
\clearpage
\setcounter{section}{8}
\refstepcounter{section}
\section*{Chapter~\thesection\quad Uniqueness up to the renormalised coupling: clause (v)\addcontentsline{toc}{section}{\protect\numberline{\thesection}Uniqueness up to the renormalised coupling: clause (v)}\label{ch:9}}
%% ---- begin inlined file: chapters/c09 ----
\begin{definition}[Radius exponents, floors on $\kappa_1$ and the fat core]
\label{9.1:not-radius-exponent-floors}
The auxiliary parameters $\mathfrak{p}$ satisfy the conditions of the frame of this chapter.
Items (i)--(iii) below give the exponent condition, the floors on $\kappa_1$ and the fat-core
condition of the frame in the form used throughout this chapter. Throughout, $d=4$, and
$T=\log g^{-2}$ is the degree variable.

(i) \emph{Exponents.} The exponent condition of the frame includes the bracket of degree
conditions in $T$ used in this chapter, with matching exponents, as follows. Ba{\l}aban's
exponents $q_0,q_1$, which are integers greater than $1$ \cite{B-CMP119}, satisfy
\begin{equation}\label{9.1:eq-radius-exponent-floors-1}
q_\flat:=\min\{q_0,q_1\}\ \ge\ p_0+2(d+1)r_{pr}+1 .
\end{equation}
At $d=4$ the right-hand side is $p_0+10r_{pr}+1$. The condition is one-sided and is met by
taking $q_0,q_1$ large. The small number $\beta_0^{pr}$ of \cite[(2.6)]{B-CMP119} is taken with
$\beta_0^{pr}\le 1/7$. The value $1/7$ is the one used as an example in \cite[\S3]{B-CMP122b}.
The condition $\beta_0^{pr}\le 1/7$ is a ceiling on the coupling, since by \cite{B-CMP119} the
number $\beta_0^{pr}$ may be
chosen arbitrarily small if $g$ is sufficiently small. It is needed because the third condition
on the coupling ceiling in the frame, at Ba{\l}aban's exponent $-\tfrac32 p_0$, has no solution
when $\beta_0^{pr}\in[\tfrac13,\tfrac12]$. The other places where $\beta_0^{pr}$ enters, namely
the exponent $7r_{pr}$, the constant $K_u$ and the floors on $\mathsf{a}$, are monotone in it.

Wherever the bracket of degree conditions in $T$ is used, it carries the prefactor exponent
$(1+d+d\beta_0^{pr})r_{pr}\le 7r_{pr}$. By \eqref{9.1:eq-radius-exponent-floors-1}, the first
two members of the augmented form of this bracket have
degree in $T$ at most
\begin{equation}\label{9.1:eq-radius-exponent-floors-2}
7r_{pr}-(q_\flat-p_0)\ \le\ -3r_{pr}-1\ <\ 0 .
\end{equation}
Hence the admissible window in $T$ is a half-line for every $L$.

Along the present construction the regime $q:=q_0=q_1$ of the correlation theorem is in force.
Moreover the constant $c_M$ introduced with the correlation theorem is reduced by the factor
$4C'_{\mathcal{J}}C_1/C_0$, where $C'_{\mathcal{J}}$ is a constant introduced with that theorem. This factor
depends on $M$ through Ba{\l}aban's condition $C_1/C_0>2B_3\,O(1)\,M$ \cite{B-CMP109}, in which
$B_3$ is a constant of \cite{B-CMP109}. The reduction is made after $M$ is chosen, as the
order of choice permits.

(ii) \emph{Floors on $\kappa_1$.} Let $\bar\kappa$ be the maximum of the $\kappa$-parameters.
Let $c_g$ be as in (iii), and let $F_1(L,\bar\kappa)$ be the fixed function of $L$ and
$\bar\kappa$, depending on no other parameter, that is introduced with the estimate of this
chapter which requires the last member of \eqref{9.1:eq-radius-exponent-floors-3}. After $L$ and
the $\kappa$-parameters, $\kappa_1$ is chosen with
\begin{equation}\label{9.1:eq-radius-exponent-floors-3}
\kappa_1\ \ge\ \max\bigl\{3,\ 2c_g(\bar\kappa+1),\ 16\bar\kappa+1,\ 4L\bar\kappa+1,
\ 8L\bar\kappa+12,\ 2F_1(L,\bar\kappa)+10\bigr\}.
\end{equation}
Put
\begin{equation*}
\varrho:=\tfrac12\bigl(e^{\kappa_1-2}-1\bigr)^{1/2},
\qquad
\kappa_1^{\natural}:=\log\frac{\varrho-1}{4}\ \ge\ \tfrac12\kappa_1-5 .
\end{equation*}

The assumption $\tfrac12(\kappa_1-1)\ge 2L\kappa$ of
\cite{B-CMP116} is a floor on $\kappa_1$. It is one of the restrictions on the constants
$M,\kappa,\kappa_1,\alpha_0,\alpha_1,\alpha_4,\alpha_6,\gamma_2,\gamma,\varepsilon_1$ of
\cite[Lemma~3]{B-CMP116}, among which $L$ does not occur. Every rate retained from
\cite[(2.37)--(2.41)]{B-CMP116} along the present construction is at most
$\tfrac12L\bar\kappa$, so this one floor covers all of them. No condition fixes
$\kappa_1$, $M$ or $A_0/A_1$ at a value.

(iii) \emph{The fat core.} The fat-core condition of the frame consists of the following floor on
$L$ and margin $\delta$. Let $c_g:=1+\sqrt d\,3^dC_d$, where $C_d$ is the dimensional constant
entering the fat-core estimate, depending on $d$ only, and let $c^\sharp:=c_g+1$,
$\beta^\star:=\max\{\beta_I,\beta_s\}=\tfrac14$ and $\kappa_R:=c^\sharp\kappa$. One application
of the construction of \cite[\S\S1--2]{B-CMP116} per step, with the fattened core, loses a
factor $c_g^2$: the fat core enters at two places. It must deliver the rate
$(1+4\beta^\star)\kappa_R$ that the subsequent estimates of this chapter spend. The conditions are
\begin{equation}\label{9.1:eq-radius-exponent-floors-a}
\begin{gathered}
L\ \ge\ \max\bigl\{7,\ 4(1+4\beta^\star)c^\sharp c_g^2\bigr\},\\
\delta:=\delta_\star:=\frac1{10}\Bigl(1-\frac{2(1+4\beta^\star)c^\sharp c_g^2}{L}\Bigr)
\in\bigl[\tfrac1{20},\tfrac1{10}\bigr),\\
\kappa_{out}(\delta'):=(1-10\delta')\,\frac{L}{2c_g^2}\,\kappa\ \ge\ (1+4\beta^\star)\kappa_R
\quad\text{for }0<\delta'\le\delta_\star,
\end{gathered}
\end{equation}
with equality in the last inequality exactly at $\delta'=\delta_\star$.

The margin $\delta$ is this single number, not an arbitrary
$\delta\in(0,\delta_\star]$; it plays the role of the parameter $\delta$ of
\cite{B-CMP116}. The reason is that the fat-core estimate uses $1/L\le\delta/8$, which
bounds $\delta$ from below. Since $c^\sharp c_g^2\ge1$, one has
\begin{equation*}
\delta_\star\le\delta_J:=\tfrac1{10}\bigl(1-\tfrac4L\bigr).
\end{equation*}
So the same $\delta$ serves the correlation theorem, which at Ba{\l}aban's cores needs only
$\delta\le\delta_J$. The floor on $L$ depends on $d$ and $\beta^\star$ only; there is no upper
bound on $L$.

A term derived from a given term in a later step is read with Ba{\l}aban's core, with one hull
$K(X)$ per term $X$ of the expansion per step. Such terms are the chain terms $\mathbb{C}^{(i)}$
and the terms $T(Y_0)$, $V(Y)$. The fat layer, the layer by which the core is thickened, is never
applied at a link between consecutive
levels in the production part of the response recursion of this chapter; otherwise, after $n$
such links, a rate $a_k$ at level $k$ would fall
to $a_k/c_g^{\,n}$.
\end{definition}

Proof outstanding.

\smallskip\noindent\textit{Route.} A proof has to establish (a) that
\eqref{9.1:eq-radius-exponent-floors-1} with $\beta_0^{pr}\le1/7$ gives the degree bound
\eqref{9.1:eq-radius-exponent-floors-2} for the first two members of the augmented bracket of
degree conditions in $T$; (b) that the floors \eqref{9.1:eq-radius-exponent-floors-3} dominate
every rate of \cite[(2.37)--(2.41), Lemma~3]{B-CMP116} used along the present construction; and
(c) that one application of \cite[\S\S1--2]{B-CMP116} with the fattened core loses at most
$c_g^2$ and delivers $\kappa_{out}(\delta_\star)=(1+4\beta^\star)\kappa_R$ under
\eqref{9.1:eq-radius-exponent-floors-a}. It must further show that all these one-sided floors
are compatible with the order in which the auxiliary parameters are chosen after $L$.

\begin{definition}[The level-wise rate, the number of stages and the strip width. Proof outstanding]
\label{9.1:not-levelwise-rate}
Write $T := \log g^{-2}$ for the degree variable. For a coupling $g$ put
\begin{equation}\label{9.1:eq-levelwise-rate-1}
\mathsf{p}(g) := \frac{\gamma_2 A_1^2\,(\log g^{-2})^{2p_0}}{64},
\qquad
\kappa_A^{pr}(g) := \frac{\mathsf{p}(g)}{N_A(g)} .
\end{equation}
The quantities $\gamma_2$, $K'$, $C_{far}$, $c_\flat$, $K_Q$, $r_A$, $K_\flat$, $\mathsf{r}$, $V^{\times}$
and the function $N_A(\cdot)$ of the coupling, which appear here and below, are fixed independently
of this definition and are not chosen here; $C_{far}$ is the constant of the outer-layer bound.
Here $\mathsf{p}(\cdot)$ is a function of the coupling and is distinct from the number of test
functions. At level $k$ the \emph{level-wise rate} is
\begin{equation}\label{9.1:eq-levelwise-rate-2}
\kappa_A := \kappa_A(k) := \min_{i \le k} \kappa_A^{pr}(g_i).
\end{equation}
It is non-increasing in $k$, it does not depend on the age of any large-field region, and it is
fixed after $\gamma_U$; it is never a number chosen before $g$.

Put $\varphi := \mathsf{p}/N_A$, regarded as a function of $T$, and let $T_*(L,\mathfrak{p},M)$
denote a threshold beyond which $\varphi$ is increasing. On the degree threshold
$T_\gamma$ the following conditions are imposed:
\begin{equation}\label{9.1:eq-levelwise-rate-3}
\begin{gathered}
\sup_{T \ge T_\gamma} K'\, T^{(1+d+d\beta_0^{pr}) r_{pr}}
\Bigl[ T^{-(q_\flat - p_1)} + C_{far}\, T^{-(q_\flat - p_0)}
+ e^{-\varphi(T-1) + c_\flat} \Bigr] \le 1 \quad\text{and}
\\
T_\gamma - 1 \ge T_* .
\end{gathered}
\end{equation}
With
\begin{equation}\label{9.1:eq-levelwise-rate-4}
\mathsf{b} := \vartheta_P + \frac{K_Q C_{far}}{r_A} + 2 e^{-(\kappa_A(k) - c_\flat)},
\qquad
\varepsilon^{\natural} := 3 K_\flat\, \mathsf{r}\, \mathsf{N}\, V^{\times} \mathsf{b},
\end{equation}
where $\mathsf{N}$ is the number of stages fixed in \eqref{9.1:eq-levelwise-rate-6} below,
the condition $\varepsilon^{\natural} \le \tfrac14$ is imposed. The list of conditions on the level
also contains the divergence
\begin{equation}\label{9.1:eq-levelwise-rate-5}
(d+1)\, r \log T - \kappa_A(k) \to -\infty .
\end{equation}
This divergence is derived, as an implication, from a further statement of this chapter; that this
condition is not a cap in the sense of Definition~\ref{2.3:def-cap} is conditional on the
hypotheses of that statement.

The \emph{number of stages} is Ba{\l}aban's integer $N$, written $\mathsf{N}$
here because $N$ is reserved for the gauge group $\mathrm{SU}(N)$; it is chosen as
\begin{equation}\label{9.1:eq-levelwise-rate-6}
\mathsf{N} := \check{R}_k .
\end{equation}
This is the choice made in \cite{B-CMP122b}, where the integer ``satisfies the conditions (i), (ii),
with $N = R_j$'', Ba{\l}aban's $R_j$ being the scaffold radius at level $j$, written
$\check{R}_j$ here. Under it the one size relation $\mathsf{N}\ln 2 < 2\check{R}_k$ imposed on
$\mathsf{N}$ holds. Ba{\l}aban's floor
\begin{equation*}
L^{-\mathsf{N}} \le (\log g_k^{-2})^{-\nu},
\quad\text{or}\quad
\mathsf{N} \ge \frac{\nu}{\log L}\,\log\log g_k^{-2},
\end{equation*}
from \cite{B-CMP122b} is also met, where $\nu$ is the exponent of that floor. The choice
\eqref{9.1:eq-levelwise-rate-6} is a relation between two scaffold integers; no cap is placed on
$\mathsf{N}$ as a number.

The \emph{strip width} is
\begin{equation}\label{9.1:eq-levelwise-rate-7}
m_\sharp := m_{hi}(k) := \Bigl\lfloor \tfrac12 \bigl(\check{R}_k - 1\bigr) \Bigr\rfloor .
\end{equation}
Together with its floor $m_{lo}(k)$, fixed with the strips, it belongs to the scaffold integers read
from $g$ by one fixed rule, as $\check{R}_k$ is.
These two integers enter
clauses (a)--(d) of the uniqueness theorem only in the way its clause (e) permits.
\end{definition}

\smallskip
\noindent\textit{Route.}
First, the supremum condition
\eqref{9.1:eq-levelwise-rate-3} must hold, the inputs being the exponent floor on $q_\flat$ of
Definition~\ref{9.1:not-radius-exponent-floors}, for the two powers of $T$, and the growth of
$\varphi$ beyond $T_*$, for the exponential term; and the bound $\varepsilon^{\natural} \le \tfrac14$
must hold at every level $k$, that is, the sum $\mathsf{b}$ must be at most
$(12 K_\flat \mathsf{r} \check{R}_k V^{\times})^{-1}$, where $\check{R}_k \ge (\log g_k^{-2})^{r_{pr}}$
grows without bound as $g_k \to 0$, so that a proof must supply bounds at level $k$ on
$\vartheta_P$, on $K_Q C_{far}/r_A$ and on $2e^{-(\kappa_A(k)-c_\flat)}$ whose sum is at most this
level-dependent quantity, the last of them requiring an upper bound on $\check{R}_k$ in terms of
$\log g_k^{-2}$, while the two inputs just named bear on neither of the first two summands.
Second, the divergence
\eqref{9.1:eq-levelwise-rate-5} must be derived from the hypotheses of the statement from which it
follows as an implication. Third, for
$\mathsf{N} = \check{R}_k$, the size relation must follow from $\ln 2 < 2$, and Ba{\l}aban's floor from
the lower bound $\check{R}_k \ge (\log g_k^{-2})^{r_{pr}}$.

\begin{definition}[The floor chain on shells and the ratio $A_0/A_1$. Proof outstanding]
\label{9.1:not-shell-floor-chain}
The following conditions on the auxiliary parameters $\mathfrak{p}$, on the coupling ceiling
$\gamma$ and on the two quantities named in item (b), and the following symbols,
are fixed. Every condition below is one-sided. Each is imposed at the place in the order of
choice of $\mathfrak{p}$ that is indicated with it.
\begin{enumerate}
\item[(a)] The floor $(F_\kappa)$ is a lower bound on the parameter $\kappa^{\flat}$,
depending on $(d,\delta_\star,c^{\sharp})$. It is imposed before $\kappa_1$ is chosen.
\item[(b)] The floor chain on shells is headed by the following floor. In it, $g$ is not the
coupling but one of the two quantities $R^{pr}M_1$ and $g(\boldsymbol{\Omega})$, the gap of the
box-set geometry $\boldsymbol{\Omega}$:
\begin{equation}\label{9.1:eq-shell-floor-chain-1}
(F_T)^{g,\varpi}:\qquad \delta_0\, g \;\ge\; 1536\,\log L .
\end{equation}
Here $\delta_0$ is the constant multiplying $g$ in the shell floors, and $R^{pr}$ is the radius
factor multiplying $M_1$ in the first of the two quantities.
This floor is required for both of these values of $g$. The chain reads
\begin{equation}\label{9.1:eq-shell-floor-chain-2}
(F_T)^{g,\varpi} \;\Longrightarrow\; (F^{\flat})^R \;\Longrightarrow\; (F_T)
\;\Longrightarrow\; (F^{\flat}) \;\Longrightarrow\; (F).
\end{equation}
The four floors $(F^{\flat})^R$, $(F_T)$, $(F^{\flat})$ and $(F)$ after the first carry the
numerical constants $512$, $256$, $32$ and $16$, respectively. Two conditions stand outside the
chain. The first is the condition in a parameter $\alpha$ and a length $RM$, carrying the
numerical constant $96$, that is taken here at $\alpha\in\{1/32,1/64\}$, with $RM$ replaced by
$g$. The second is the floor $M_1\ge 8d$.
Every shell margin has $R_1M_1$ in place of $M_1$:
\begin{equation*}
s_n := L^n\,\eta\,R_1M_1 ,
\end{equation*}
where $\eta$ is the fixed factor of the shell margins.
Here $R_1$ is a fixed power of $L$. It is
neither the radius $R_k(g) := R(x_k)$ of level $k$, nor the window radius
$R_1^w := R_0/(2\mathsf{Q}_w+2)$.
\item[(c)] The floors at the T-site are the floor $(F_\varpi)$, whose content in terms of the
parameters is stated below together with that of $(F1)'$, and
\begin{equation*}
(F1)':\qquad \varepsilon_T\le 1,\qquad
\varepsilon_T\,C_\rho(\delta_T)\le \tfrac12\,c_\eta ,
\end{equation*}
where $\varepsilon_T$, $c_\eta$, $\delta_W$ and the function $C_\rho(\cdot)$ are the constants of
the bounds at the T-site, and $\delta_T := \delta_W/4$.
The condition $(F1)'$ does not involve $C_{far}$, which is defined in (e) below. In terms of the
parameters, these floors are a floor
\begin{equation*}
\frac{\delta_W M}{4} \;\ge\; 32\kappa_1+2\log(LM)+C(d,L)
\end{equation*}
on $M$, where $C(d,L)$ is a constant depending only on $d$ and $L$, imposed after
$(L,\kappa_1)$, together with ceilings on $\alpha_0$ and $\alpha_1$,
imposed after $M$.
\item[(d)] The condition $(F_{far})^{\gamma}$ is
\begin{equation*}
M\check{R}_{k+1} \;\ge\; M_{far}(C_{far}),
\end{equation*}
where $M_{far}(C_{far})$ is a function of the constant $C_{far}$ defined in (e) below.
It is a ceiling on $\gamma$, imposed level by level.
\item[(e)] Put $\mathsf{a} := A_0/A_1$. It is required, after $M$, that
\begin{multline}\label{9.1:eq-shell-floor-chain-3}
\mathsf{a} \;\ge\; \max\Bigl\{c_{Ax}(d,L),\; c(d)(1+\beta_0)B_3L^{d+2},\;
\tfrac{d(L-1)}{2},\\
2\sqrt2\,(1+\beta_0)B_3,\; c'(d,L)M^2\Bigr\}.
\end{multline}
The constants $c_{Ax}(d,L)$, $c(d)$ and $c'(d,L)$ depend only on the arguments shown, $B_3$ is a
fixed constant, and $\beta_0$ is the small growth number of \cite{B-CMP119}.
Only the lower bound
\eqref{9.1:eq-shell-floor-chain-3} is imposed. Then
\begin{equation*}
C_{far} := C(d,L)\,A_0/A_1 ,
\end{equation*}
with a constant $C(d,L)$ depending only on $d$ and $L$, not necessarily that of item (c).
\item[(f)] The symbols of the fat core are the following. First,
$\kappa^{fat} := \kappa/c_g$. Second,
\begin{equation*}
\bar{\varepsilon}^{fat} := E_0\,\varepsilon_1\,C_1^{fat}(d,L,M_1)\,M^{\bar q}\,
e^{C_2^{fat}(d)\kappa_1}\,e^{2\kappa}.
\end{equation*}
Third, $a := 2(L+2)^4\,O(1)\,\varepsilon_2^{fat}$, subject to $a\,e^{4L\kappa^{fat}}\le 1$.
Fourth, the constant $C_3^{fat}$ is multiplied by $\alpha_6^{-1}e^{2L\kappa^{fat}}$. Fifth,
\begin{equation*}
E^{*} := O(1)\,C_3^{fat}\,\varepsilon_1\,e^{L\kappa^{fat}} \;\le\; E_0/16 .
\end{equation*}
Here $\varepsilon_1$ is the small parameter on which these conditions are ceilings, $E_0$ is the
bound that they maintain, and $\varepsilon_2^{fat}$, $\alpha_6$, $\bar q$ and $C_i^{fat}$ are
constants of the fat core of Definition~\ref{9.1:not-radius-exponent-floors}, the last depending
only on the arguments shown.
These conditions are ceilings on $\varepsilon_1$, imposed after $(L,\kappa,\kappa_1,M)$ in
the order of \cite[Lemma~3]{B-CMP116}. Sixth, the ceilings of the flat bound in its
$\varpi$-form are multiplied by a constant $c_{loc}$. In the fat core, the condition of item (b) in
$\alpha$ and $RM$ is read at $\alpha = 1/32$.
\item[(g)] Distinct symbols written alike.
\begin{itemize}
\item The symbol $c_1$ has four meanings, according to position, and none stands for
another:
\begin{itemize}
\item the polylogarithmic exponent $c_1 := p_0+q_\flat+1$, with $p_0$ the fixed exponent, as in
$\mathsf{l}_r = (\log\theta_r)^{c_1}$ and in $T^{c_1+2q_\flat+\mathfrak{a}}$;
\item the slowness constant of the random-walk bounds, as in $C_T^w := c_1L^3$ and in the
constants $C_4$, $C_{sl}^{\approx}$ and $K_o^L$ of those bounds;
\item the activity constant $c_1\in\{e^{-p_0(g_k)},\alpha^{1/3}\}$ of
\cite[(1.97)--(1.99)]{B-CMP122b}, where here $p_0(g_k)$ denotes $(\log g_k^{-2})^{p_0}$, the
degree function of \cite{B-CMP122b}, without the factor $A_0$ of the degree function
$p_0(\cdot)$ of the glossary; it is always written
\begin{equation*}
c_1^{\flat}(g_\ell) := O(1)\max\bigl\{e^{-p_0(g_\ell)},(\alpha^{\sharp}_\ell)^{1/3}\bigr\},
\end{equation*}
with $p_0(g_\ell) = (\log g_\ell^{-2})^{p_0}$ in the same sense;
\item the $c_1$ inside the constant $C_1$ of \cite{B-CMP116}, which is not the auxiliary
parameter $C_1$ of $\mathfrak{p}$.
\end{itemize}
\item The crossing length $\mathsf{d}_m := 2\check{R}_m/L-c_\times$, with $c_\times$ a fixed
constant, is the thickness of
zone $m$ at the next scale. It enters the crossing weights
$\mathsf{w}_{\ell,n} := e^{-\mathsf{d}_{\ell-1}}$. The operator size is written
$\mathfrak{d}^{op}_n := C_{cl}^T(\log x_n)^{c_2}\vartheta_T^{\,n}$, with $c_2$ a fixed exponent
of the operator-size bound, distinct from the constant $c_2$ of the glossary (with $c_5 = 2c_2$).
In the zone bounds every
$e^{-\mathsf{d}_m}$ carries the crossing length $\mathsf{d}_m$, never the operator size
$\mathfrak{d}^{op}_n$; both are distinct from the separation $\mathsf{d}$ of supports.
\item The classical rate at the T-site is $\vartheta_w := L^{-\theta_A/8}$. The inequality
$\vartheta_T\le\vartheta_2$ need not hold (it fails for $\theta_A>3/4$); the only comparison
used is $\vartheta_T<\vartheta_2^{1/4}$, through
\begin{equation*}
\vartheta_\times := \vartheta_w\vartheta_2^{-1/4} = L^{-(3-\theta_A)/32} < 1,
\qquad
e_T := \frac{3-\theta_A}{5\theta_A-3} > 1 .
\end{equation*}
In the zone bounds at the T-site, every power $\vartheta_2^{3n/4}$ is replaced by
$\vartheta_\times^{\,n}$. The exponent $-3$ of the zone bound for the operator size is replaced
by $-\min\{3,e_T\}$.
\item The counting constant $K_{\mathbf{B}}$, also written $K(d)$, of the terminal terms of the
expansion is distinct from the rate $\kappa_{\mathbf{B}} := \kappa$ at which those terms are
stored (the storage rate $\kappa_T$ evaluated at $\kappa$). For first-group terms, the
production rate of the
$\mathbf{R}$-step is $(1+\tfrac12\beta_s)\kappa = \tfrac98\kappa\ge\kappa_{\mathbf{B}}$.
\end{itemize}
\end{enumerate}
\end{definition}

Proof outstanding.

\smallskip
\noindent\emph{Route.}
A proof would have to show that the head floor \eqref{9.1:eq-shell-floor-chain-1}, taken at
$g=R^{pr}M_1$ and at $g=g(\boldsymbol{\Omega})$, yields in turn each floor of the chain
\eqref{9.1:eq-shell-floor-chain-2}, with the constants $512$, $256$, $32$ and $16$ of those
floors. It would further have to show that $(F_\varpi)$ and $(F1)'$ control the T-site with no
condition involving $C_{far}$, because, by two statements on the cubes at the T-site, no cube of
the strong class lies in the far region.
Finally, it would have to show that the floors and
ceilings of items (a)--(f) can all be met together in the order of choice stated with them.

\begin{definition}[Index exponents, rate letters and the window for $\rho$]\label{9.1:not-rho-window}
\leavevmode
\begin{itemize}
\item[(a)] \emph{Numbers fixed before $L$.} We fix $\theta_A\in\mathopen{]}3/5,1\mathclose{[}$, $\beta_I=1/5$, $\alpha=1/32$ and
\begin{equation}\label{9.1:eq-rho-window-1}
\theta_1:=\frac{5\theta_A-3}{2},\qquad \theta_2:=\frac{\theta_1}{4},\qquad \theta':=\frac{\theta_2}{2}.
\end{equation}
These are numbers, chosen before $L$. The number $\alpha$ here is an index number of this chapter
and is not one of Ba{\l}aban's radii $\alpha$.
\item[(b)] \emph{The index at which hybrid and healing terms are estimated.} We put
$\beta^{\flat}:=\beta_H-\varepsilon_H/2$, where $\beta_H$ is the index attached to the hybrid and
healing terms in this chapter and $\varepsilon_H>0$ fixes the spacing of the ladder of indices below
$\beta_H$ used in the next sentences.
The hybrid terms and the healing terms (the terms that arise when a large-field region is healed, in
the sense of the large-field vocabulary fixed with the two-point witness theorem) are estimated once,
at the index $\beta^{\flat}$, and the index $\theta$ delivered by these estimates is unchanged. The at
most eleven layer links are estimated on rungs of the ladder of indices $\beta_H-i\varepsilon_H/32$,
$i=1,2,\dots$, and the rungs used all lie above $\beta^{\flat}=\beta_H-16\varepsilon_H/32$; that is,
they have $i\le 15$.
\item[(c)] \emph{Rate letters.} We put $\vartheta_2:=L^{-\theta_2}$, the real one-move index, and
$\vartheta_C:=L^{-\theta'}$, the rate of the corollary on the chain terms; since $\theta'=\theta_2/2$,
we have $\vartheta_C=\vartheta_2^{1/2}$. The
estimates of the data away from the production block $P$ take at most eight successive halvings of the
exponent; this gives the rate $\vartheta_f:=\vartheta_2^{1/512}=\vartheta_C^{2^{-8}}$. One square root
taken in the M\"obius inversion at the terminal term replaces a rate $\hat\vartheta$ by
$\hat\vartheta^{\natural}:=\hat\vartheta^{1/2}$, and the family $\mathfrak{F}^V$ of vertex pairs costs
one joint square root, $\vartheta\mapsto\vartheta^{1/2}$. The list of losses in the exponent (the
passage from $\vartheta_2$ to $\vartheta_C$, the halvings, the M\"obius root and the joint root) is
therefore $(\tfrac12,2^{-8},\tfrac12,\tfrac12)$. Let $\mathfrak{i}$ denote the finite list of index
exponents delivered by the estimates of this chapter; its members are $\theta_2/2048$, the exponent
$\theta_2\beta_I/(1+\beta_I)$ multiplied by the roots taken with it (at the terminal stage only), and
the remaining delivered index exponents, and its least member is $\theta_2/2048$. We put
\begin{equation}\label{9.1:eq-rho-window-2}
\theta_{dl}:=\min\mathfrak{i}=\frac{\theta_2}{2048},\qquad
\vartheta_{dl}:=L^{-\theta_{dl}}=\vartheta_f^{1/4}.
\end{equation}
No further square root of $\vartheta_{dl}$ is taken: $\vartheta_{dl}$ is the delivered rate, in the
sense that the classical index rates satisfy
$\hat\vartheta^{\natural}_i\le C\vartheta_f^{i/4}=C\vartheta_{dl}^{\,i}$ with a constant $C$.
\item[(d)] \emph{The window for $\rho$.} Let $q$ denote the kernel rate of the two-run comparison of
this chapter, a number in $\mathopen{]}0,1\mathclose{[}$, distinct from Ba{\l}aban's auxiliary
parameter $q$ and from the integers $q_0,q_1$ of
Definition~\ref{9.1:not-radius-exponent-floors}. Let $\theta_\chi>0$ be the exponent of the estimate
for $\chi$, a number fixed before $L$, and
$\vartheta_\chi:=L^{-\theta_\chi}$; and let $\vartheta_\sharp>0$ be the rate letter of this chapter
whose power $L^{-\vartheta_\sharp}$ enters the window below, fixed before $\rho$. The letter $\rho$ is
any number with
\begin{equation}\label{9.1:eq-rho-window-a}
\rho\in\bigl[\max\{q^{1/2},\,L^{-\vartheta_\sharp},\,\vartheta_{dl},\,\vartheta_\chi\},\,1\bigr[,
\qquad \rho>\vartheta_{dl}.
\end{equation}
In particular $q\le\rho^2$; this is a floor on $\rho$, and it is preserved when $\rho$ is increased. The window \eqref{9.1:eq-rho-window-a} is non-empty for every $L$.
\item[(e)] \emph{What the choice of $\rho$ leaves unchanged.} The constants $c_\rho$, $C_T$ and
$T_\Pi$ that enter the two-run estimate with the power-law floor do not depend on the length $n$ of the
run; the ceiling conditions used with the two-run estimate with the power-law
floor are chosen after $\rho$; the power-law rates do not depend on $\rho$; and the choice of $\rho$
moves no rate, floor or display fixed before it.
\end{itemize}
\end{definition}

\begin{proof}
\emph{The exponent $\theta_{dl}$.} Since $\theta_A>3/5$, we have $5\theta_A-3>0$, so $\theta_1$, $\theta_2$ and $\theta'$ in \eqref{9.1:eq-rho-window-1} are positive. The product of the losses is
\begin{equation*}
\tfrac12\cdot 2^{-8}\cdot\tfrac12\cdot\tfrac12=2^{-11}=\tfrac{1}{2048},
\end{equation*}
and $\theta_2\cdot 2^{-11}=\theta_2/2048$ is the member of $\mathfrak{i}$ obtained from $\theta_2$
through the list of losses; that this member is the smallest member of $\mathfrak{i}$, and hence
equals $\theta_{dl}$ in \eqref{9.1:eq-rho-window-2}, is the statement about the delivered exponents
$\mathfrak{i}$ made in (c). The equality
$\vartheta_{dl}=\vartheta_f^{1/4}$ holds because
\begin{equation*}
\vartheta_f^{1/4}=\bigl(L^{-\theta_2/512}\bigr)^{1/4}=L^{-\theta_2/2048}=L^{-\theta_{dl}}.
\end{equation*}
For the same reason $\vartheta_f^{i/4}=\vartheta_{dl}^{\,i}$ for every $i$, which is the equality in the last sentence of (c). Since $\theta_{dl}>0$ and $L>1$, we have $0<\vartheta_{dl}<1$.

\emph{The window.} The left end point of \eqref{9.1:eq-rho-window-a} is the maximum of four numbers,
each of which is less than $1$ for every $L>1$: $q^{1/2}<1$ because the kernel rate satisfies
$0<q<1$, and $L^{-\vartheta_\sharp}$, $\vartheta_{dl}$ and $\vartheta_\chi$ are less than $1$ because
$\vartheta_\sharp$, $\theta_{dl}$ and $\theta_\chi$ are positive (for $q$, $\vartheta_\sharp$ and
$\theta_\chi$ by (d), for $\theta_{dl}$ by the first paragraph). Hence the
interval
\eqref{9.1:eq-rho-window-a} contains points strictly above $\vartheta_{dl}$, and the window is
non-empty for every $L$. From $\rho\ge q^{1/2}$ we obtain $q\le\rho^2$, and if $\rho$ is increased
this inequality persists.

\emph{Why the floor is $q\le\rho^2$.} The sums in \eqref{9.6:eq-two-run-estimate-a} and the
quantities $\hat{\mathfrak{r}}_r$ and $c_\rho$ are estimated, and the ceiling conditions used with the
two-run estimate with the power-law floor are stated, under the condition $q\le\rho^2$. In those sums
there occurs the series $\sum_{s\ge r}q^{s-r}e_s$, where $n$ is the length of the run,
$e_s=\rho^{n-s}+\Pi_n$ and $\Pi_n$ is the
power-law floor of the two-run estimate. For $q<\rho$ the geometric series give
\begin{equation*}
\sum_{s\ge r}q^{s-r}\rho^{n-s}=\rho^{n-r}\Bigl(1-\frac{q}{\rho}\Bigr)^{-1},\qquad
\sum_{s\ge r}q^{s-r}\Pi_n=\Pi_n(1-q)^{-1}\le\Pi_n\Bigl(1-\frac{q}{\rho}\Bigr)^{-1},
\end{equation*}
the last inequality because $q\le q/\rho$; hence the series is bounded by $e_r$ times
$(1-q/\rho)^{-1}$. Under
$q\le\rho^2$ we have $q/\rho\le\rho<1$, so this factor is at most $(1-\rho)^{-1}$. If instead
$q/\rho>\rho$ were allowed, the factor $(1-q/\rho)^{-1}$ would be unbounded as $\rho$ decreases to
$q$, and the three estimates indexed by $\tfrac12$, $\tfrac34$ and $\tfrac14+\tfrac14$ in the proof of
the two-run estimate with the power-law floor would not close.

\emph{Consequences.} The age rate satisfies $q_a\le q^2$, by its choice in this chapter, and hence
$q_a\le\rho^4$ by $q\le\rho^2$. Since $\rho<1$, this gives $q_a<\rho$, and together with
$\rho>\vartheta_{dl}$ from \eqref{9.1:eq-rho-window-a} we obtain $\rho>\max\{q_a,\vartheta_{dl}\}$.
Next, let $N\ge1$, $h\ge0$ and $k$ be the integers that enter the estimate of the small terms of the
comparison, which are related there by $k-1=(N-1)+h$ (in this paragraph $N$ and $k$ denote the
integers of that estimate, not the rank of the gauge group and not a level). Then $q\le\rho^2$ and
$\vartheta_{dl}<\rho$
give
\begin{align*}
(N+1)q^{N-1}\vartheta_{dl}^{\,h}
&\le(N+1)\rho^{2(N-1)+h}
=(N+1)\rho^{N-1}\rho^{k-1}
\le C'_\rho\,\rho^{k-1},\\
C'_\rho&:=\sup_{N\ge1}(N+1)\rho^{N-1},
\end{align*}
and $C'_\rho<\infty$ because $\rho<1$. This is the bound used in the estimate of the small terms of
the comparison.

\emph{The last rung.} A rung $\beta_H-i\varepsilon_H/32$ lies above
$\beta^{\flat}=\beta_H-\varepsilon_H/2$ exactly when $i\varepsilon_H/32<\varepsilon_H/2$, that is,
since $\varepsilon_H>0$, exactly when $i<16$, that is, for the integer $i$, exactly when $i\le15$:
\begin{equation*}
i\le15\ \Longrightarrow\ \frac{i\varepsilon_H}{32}\le\frac{15\varepsilon_H}{32}
<\frac{16\varepsilon_H}{32}=\frac{\varepsilon_H}{2},\qquad
i\ge16\ \Longrightarrow\ \frac{i\varepsilon_H}{32}\ge\frac{\varepsilon_H}{2}.
\end{equation*}
Hence the two descriptions of the rungs used in (b), that they lie above $\beta^{\flat}$ and that they
have $i\le15$, coincide.

\emph{Independence.} The letter $\rho$ is chosen after all the rates and floors it is compared with in
\eqref{9.1:eq-rho-window-a}. Therefore its choice moves none of them. The ceiling conditions used
with the two-run estimate with the power-law floor are imposed after $\rho$. The constants $c_\rho$,
$C_T$, $T_\Pi$ do not depend on the length $n$ of the run, by their definitions.
\end{proof}

\begin{definition}[The order in which the constants are chosen]\label{9.1:not-constant-order}
The constants of this chapter are chosen in the order \eqref{9.1:eq-constant-order-a}. Its groups
(o1)--(o9), with (o4$'$) between (o4) and (o5), are specified below.

An arrow $a\to b$ means that $b$ is chosen after $a$ and may depend on $a$. A condition listed
inside a group is a lower bound (a floor) or an upper bound (a ceiling) in the sense of
Notation~\ref{2.3:not-stages-first}, and it reads only letters chosen before the letter it bounds.
Constants named below whose definition is given later in this chapter are listed here by their
symbols only: each is defined at its first use, and the list fixes only the place at which it is
chosen.
Likewise, a condition named below by a parenthesised tag, such as $(\mathrm{F}_{cut})$,
$(\mathrm{SL})^\flat$ or $(\gamma2)$, is displayed at the place in this chapter where it is used;
the list records only the letter that the condition bounds and the group in which it is imposed.

Throughout, $d=4$, $T:=\log g^{-2}$ and $T_\gamma:=\log\gamma^{-2}$, where $\gamma$ denotes a
candidate value of the coupling ceiling $\gamma_U$ of (o8). In the members of (o8) and in the
definition of $c_a^{GR}$ in (o9), a bare $C$ denotes the constant of the bound in which the member
is used, chosen before $\gamma$.
\begin{equation}\label{9.1:eq-constant-order-a}
\begin{aligned}
&\text{(o1) numbers}\;\to\;\text{(o2) } L,\ \delta_\star(L),\ (d,L)\text{-functions},\ \gamma^{\circ}\\
&\to\;\text{(o3) } r_{pr},\,p_0,\,p_1,\,q_\flat,\,\kappa_0,\,c_1\;\to\;\text{(o4) }
\kappa^\flat,\,\kappa,\,\kappa_1,\,M_1,\,M\\
&\to\;\text{(o4$'$) }\mathsf{a},\,C_0,\,C_1,\,\mathsf{r}_1\;\to\;\text{(o5)}\;\to\;\text{(o6)}
\;\to\;\text{(o7) } r_J\\
&\to\;\text{(o8) } \gamma_U\;\to\;\text{(o9) } g\le\gamma_U .
\end{aligned}
\end{equation}
\begin{enumerate}
\item[(o1)] \emph{Numbers chosen before $L$.}
\begin{itemize}
\item The index exponent $\theta_A\in\,]3/5,1[$, also written $\theta$, and the exponents
$\theta_1,\theta_2,\theta'$ of Notation~\ref{9.1:not-rho-window}.
\item The numbers $\beta_H$, $\varepsilon_H$, $\beta^\flat:=\beta_H-\varepsilon_H/2$ and $\beta_0^H$
of Notation~\ref{9.1:not-rho-window};
the values $\beta_{pr}=\beta_I:=1/5$, $\beta_s:=1/4$, $\beta^\star:=\beta_s$; and $\alpha=1/32$.
\item The schedule number $\theta_S\in[1/2,8/9]$. One value is fixed in this window, and the
correlation theorem, the uniqueness theorem and their given inputs, labelled \textup{(K)} and
\textup{(J1)}--\textup{(J5)} where they are stated, are all read at this same value. From it one
sets
$\theta_\rho:=\tfrac14\min\{\theta_S,1-\theta_S\}$.
\item The exponent $e_T:=(3-\theta_A)/(5\theta_A-3)$.
\item The number $\beta_0^{pr}:=1/7$.
\item The numbers $C_\Sigma:=8+6\beta_I$, $K_u:=27/4$, $C_S:=4$ and $m_L:=15$.
\item The loss list $\mathfrak{i}:=(1/2,2^{-8},1/2,1/2)$ and $\theta_{dl}:=\theta_2/2048$.
\item The exponent $\theta_\chi$, defined later in this chapter.
\item The dimensional constants $c_d$, $C_d$, $c_g=c_H:=1+\sqrt{d}\,3^dC_d$,
$c^{\sharp}:=c_g+1$, $\mathsf{c}_d:=2\cdot 3^d$, $\mathsf{k}_d$,
$C_{fib}:=2e^{2\mathsf{k}_d}$, and $\beta_v:=\beta/(2\mathsf{c}_d)$ for the step's value
$\beta\in\{\beta_I,\beta_s\}$.
\item The value $\beta_\sharp:=1/20$ and the number $199^d3^{d+1}$.
\item The $d$-numbers
\[
K^{\circ}(d):=C_{fib}\sup_{y\ge1}y^de^{-(y-3)},\qquad K_{\circ}(d):=3^{d+1}e^8C_{fib}.
\]
Wherever leaves are counted through a cell, $K(d)$ is replaced by $K(d)+K_\circ+K^\circ$.
\item The scaffold parameter $k_*$, a number independent of $K$; it is the lower bound in
$n_c\ge k_*$ for the scaffold integer $n_c$.
\end{itemize}
\item[(o2)] \emph{The blocking factor $L\ge L_0$, odd.}
\begin{itemize}
\item Finitely many lower bounds on $L$ enter $L_0$, each reading (o1) only. They are:
\begin{itemize}
\item the bounds $7$ and $28$;
\item the bound $(\mathrm{F}_\star)$ $L\ge 8c^{\sharp}c_g^2$;
\item the bound $L\ge 80+4c^{\sharp}c_g^2$, which is equivalent to $1/L\le\delta_\star/8$;
\item the bound $L\ge 5$, so that four level-$m$ cubes lie in one level-$(m+1)$ cube;
\item the bound $L\ge 18(d+1)$, which re-routes walk junctions inside an island;
\item the bounds $L\ge3\ell_\circ^2$, $\ell_\circ\ge2$ and $L^{1-\theta_A}\ge2$;
\item two further lower bounds stated with their users;
\item the bound $(1+Cq\log L)^{1/2}L^{-3q}\le1$ of the correlation theorem, with $C$ and $q$ the
constants so denoted there.
\end{itemize}
\item Then $\delta:=\delta_\star(L)$.
\item Then the $(d,L)$-functions. Each of them is valid for every $M_1\ge M_1^{thr}(d,L)$. They
are:
\begin{itemize}
\item the operator-budget rate
$\vartheta_\times:=\vartheta_w\vartheta_2^{-1/4}=L^{-(3-\theta_A)/32}$;
\item $\delta_0$, $\delta_1:=\delta_0/32$, and the rate $\delta_K$ of the same kind; all three are
$(d,L)$-functions of this group, and $\delta_K$ is read in (o4$'$) and (o8);
\item $R_1$, $p_R$, $M_1^{thr}$ and Ba{\l}aban's constants $B$;
\item $c_0$, $c_2$ and $\lambda$;
\item the first ceiling $\gamma^{\circ}$ on the coupling, chosen before every supremum over the
coupling window;
\item the threshold $\gamma_0^X$ of the correlation theorem, then
$\gamma_0:=\gamma_0^X/c_\Phi(d,L)^2$, with $c_\Phi(d,L)$ a constant depending on $d$ and $L$
only, then $\gamma_0^T$;
\item $B_{\mathcal{A}}$ and $\delta_{\mathcal{A}}$;
\item the constants $B$, $\delta_W$ of the random-walk bound at the base point $\sigma\equiv0$;
\item $\Theta_1$, $\varpi_T$, $\Theta_\flat$ and $\tau_\flat(\cdot)$.
\end{itemize}
\end{itemize}
\item[(o3)] \emph{Exponents.} This group reads the number $\beta_0^{pr}=1/7$ of (o1). The order is
\[
r_{pr}\;\to\;p_0\;\to\;p_1\;\to\;q_\flat\;\to\;\kappa_0,\ c_1 ,
\]
with $p_0\ge5r$, where $r=r_{pr}$, $p_1<p_0$, and $q_0=q_1=:q_\flat$ subject to the floor
$(\mathrm{Q})'''$ of
Notation~\ref{9.1:not-radius-exponent-floors}, which contains $(\mathrm{Q})^A$.
The equality $q_0=q_1$, condition $(\mathrm{Q})'$, is used in the proof of the uniqueness theorem.
The constant $A_1>0$ is free.
\item[(o4)] \emph{Decay rates and the cube sizes.} The order is
$\kappa^c\to\kappa^\flat\to\kappa$, with $(\mathrm{F}_\kappa)$ imposed at $\kappa^\flat$
(Notation~\ref{9.1:not-shell-floor-chain}).

The floors on $\kappa$ are:
\begin{itemize}
\item $\kappa\ge10\log L$;
\item $\kappa/10\ge c_d+2\sqrt d(2+4\beta_{pr})$, read at the split spare
$\kappa/20\ge c_d+2\sqrt d(2+4\beta_{pr})$ where the values of the chain terms
$\mathbb{C}^{(i)}$ are placed;
\item $(\mathrm{F}_\kappa)$ $\kappa\ge c^{\sharp}C(d)/\delta_\star$; this floor reads
$(d,\delta_\star,c^\sharp)$ only, and it is imposed as a floor on $\kappa^\flat$, in the form
stated in Notation~\ref{9.1:not-shell-floor-chain}, before $\kappa_1$;
\item $(\mathrm{F}_\sharp)$ $\beta_\sharp\kappa\ge2(\mathsf{k}_d+1)$;
\item $\kappa\ge C(d)c^{\sharp}$, for sums at rate $a_\flat/4$;
\item $(\mathrm{F}_\sigma)^z$ $\sigma_*\ge c_d+3$, where $\sigma_*\in\{\beta\kappa,\kappa/(4c^{\sharp})\}$
is the leaf rate of the step; this is one unit more than $c_d+2$, the extra unit
being spent on the crossing weight;
\item $(\mathrm{F}_\sigma)^{\circ}$ $\sigma_*\ge\mathsf{k}_d+4$, which accommodates the margin
$\mathsf{k}_d+2$ of the relative count together with one unit for the moat crossing.
\end{itemize}
The storage rate is $\kappa_{\mathbf{B}}:=\kappa$, and it is subject to no ceiling.

After $\kappa$ come $\kappa_R:=c^{\sharp}\kappa$ and $a_\flat:=(1-\beta_{pr})\kappa/c^{\sharp}$,
then $\bar\kappa$, then $\kappa_1$ with
\[
\kappa_1\ge\max\{3,\ 2c_g(\bar\kappa+1),\ 16\bar\kappa+1,\ 4L\bar\kappa+1,\
8L\bar\kappa+12,\ 2F_1(L,\bar\kappa)+10\},
\]
then $\kappa_1^{\natural}$, then $B_W:=B(d,L)e^{16\kappa_1}$, then $x_*$.

Next comes $M_1$, after $\kappa_1$. Its floors are:
\begin{itemize}
\item the thresholds;
\item the chain headed by $(\mathrm{F}_T)^{g,\varpi}$ of
Notation~\ref{9.1:not-shell-floor-chain}, which implies $(\mathrm{F}\flat)^R$ and the further
floors of that chain, imposed on $R^{pr}M_1$ and on the gap $g(\boldsymbol{\Omega})$;
\item the floor $(\mathrm{F}_\omega)$;
\item $M_1\ge8d$.
\end{itemize}

Then $\mathsf{m}_{15}$, then $M=L^{m'}$. Writing $\tilde\delta_1$ for the constant $\delta_1$ of
\cite{B-CMP116}, the floors on $M$ are:
\begin{itemize}
\item $\tilde\delta_1M\ge4\kappa_1$;
\item $\delta_1M\ge8\log\ell_\circ$ and $e^{-\delta_1M/2}\le(2e\Upsilon)^{-1}$;
\item $\delta_1M\ge2(q_\flat-p_1+1)\log L/r_{pr}$;
\item $\delta_0M>48\kappa_1$;
\item $M\ge4LR_1M_1$, since every shell margin reads $R_1M_1$, and $M\ge8LM_1$;
\item the floors $(\mathrm{F}_{cut})$, $(\mathrm{F}_\varpi)$, $(\mathrm{F}1)'$ and
$6K_1^{TOL}O(1)e^{-\delta_0M/3}\le\mathfrak{b}^{tol}$;
\item the pair-radius condition $(\mathrm{R}\text{-}\mathrm{b})$, the floors $(\mathrm{C}2)^+$ and
$(\mathrm{F}_M)$, and the $M$-floors of two further conditions stated with their user.
\end{itemize}
$M$ is chosen after $\kappa_R$.

Then $R^{pr}$ and $R_0^{pr}M_0$, then $\gamma_2:=\mathsf{c}_2(LM)^{-4}$, then $c_\eta$.
\item[(o4$'$)] \emph{The ratio $\mathsf{a}=A_0/A_1$.} It is bounded below by a fixed multiple of
$M^2$ and by the further lower bounds of Notation~\ref{9.1:not-shell-floor-chain}. The order is
then
\[
\mathsf{a}\;\to\;C_{far}\;\to\;M_{far}:=1+\tfrac{2}{\delta_K}\log(16K_1^{TS}C_{far})\;\to\;C_0
\;\to\;C_1\ge2B_3O(1)MC_0\;\to\;\mathsf{r}_1:=C_1/C_0 .
\]
\item[(o5)] \emph{Classical constants, first group.} None of these reads $c_M$, $\varepsilon_1$ or any
later letter. They are the following.
\begin{itemize}
\item Every constant depending on $(d,L,\theta,M_1)$ only, $C_R^{off}$ among them.
\item $B_\natural(d,L,\theta,M_1,R^{sh})$, where $R^{sh}$ is the parameter of the shell sequences.
$R^{sh}$ is distinct from the radius $R:=4r_J$ of (o7) and from $\check R_k$.
\item $C_w$; the constant $B_\natural$ depends on it through the bound $2C_w$ on the local
Lipschitz constant of the Landau chart.
\item $C_s$, $c_1$, $C_T^w:=c_1L^3$, $B^{\mathsf{s}}(d,L,M_1,\kappa_1)$, $C_\flat'$ (which contains
$C_sL^2$) and $K_\Lambda'(M)$.
\item The constants
\[
C_4:=16(1+68dC_sC_T^w)\max\{C_s,B^{\mathsf{s}}\}c_1,\qquad
C_{sl}^{\approx}:=(8/\beta_I)\,c_1\Upsilon,\quad \Upsilon:=3(1+\beta_0^{pr})^3L.
\]
The factor $\max\{C_s,B^{\mathsf{s}}\}$ in $C_4$ appears because two of the cases of its estimate,
those at dilated or cut vertices, read $B^{\mathsf{s}}$ in place of $C_s$.
\item $C_{dp}$, $C_{dp}'$, $C_W$, $C_{\mathcal{J}}'$ and $C_T^{TL}$.
\item The profile constant $c_\nabla$, for which
$\ell_\nu|\nabla h|\le c_\nabla\cdot(\text{sup-profile of }h)$ on the contours used here.
\item The constant $C_4^{\sharp}:=2^9dB^{\mathsf{s}}C_sc_1L^8C_4$.
\item $\hat C_1$, the constant $C_1$ of \cite[(1.20)]{B-CMP116}.
\item After $K_\Lambda'$:
\begin{itemize}
\item $K_0$, the constant $O(1)B_3^2B_5M^6$ of \cite[(1.87)]{B-CMP122a};
\item $K^\flat:=\max\{c_QK_0,\,1.1(K_{83}+K_uK_0)\}$ and $K_\Lambda^\flat:=\max\{K_\Lambda',K^\flat\}$;
\item $K_{\mathfrak{K}}:=80d(K_0+K_\Lambda^\flat+B^{\mathsf{s}}C_sc_1K_\Lambda'+1)$;
\item $C_4^{\mathfrak{K}}:=9C_4+2B_\natural c_1(K_{\mathfrak{K}}+3)/K_\Lambda'$;
\item $K_o^{\mathfrak{K}}:=2^{20}B_\natural c_1^2C_TB^{\mathsf{s}}C_sC_1$;
\item $n_{\mathfrak{K}}:=(100\sqrt d\,M/M_2+3)^d$.
\end{itemize}
\end{itemize}

These are followed by $K_o:=2^{18}B_\natural C_{sl}^{\approx}\mathsf{r}_1$, a single value. Then
comes the $\alpha$-chain:
\begin{itemize}
\item $(\mathrm{IF}\text{-}\mathrm{c})$ $C_\sharp'(M\alpha_0+\alpha_1)\le\tfrac12$, which does not read
$\kappa_1$;
\item the primed form $(\mathrm{IF}\text{-}\mathrm{c})'$ of the preceding condition, and the
condition $(\mathrm{F}1)^\alpha$, both conditions of this chain on the radii $\alpha_0$, $\alpha_1$,
displayed where they are used;
\item $M_1\alpha_0\le a_E$.
\end{itemize}
The chain continues:
\[
E_0\;\to\;C_V:=4^{C_d}e^{\kappa_*}\;\to\;\mathsf{c}:=3C_V+5\ (\text{the class constant})
\;\to\;c_M .
\]
Here $c_M$ carries the divisors $140dC_{\mathcal{J}}'\mathsf{r}_1$ and $K_o(1+c_\nabla)/(35d)$,
each applied once and at this place, and it satisfies
\[
2^{18}B_\natural c_1C_\flat'c_M\le\theta_S\beta .
\]
Then comes $K_o^L:=2^{17}B_\natural C_T^wc_1^2B^{\mathsf{s}}/(\theta_\rho\beta)$, the constant of the
$(L)$-circle, that is, of the dilation circle $\mathsf{w}_L(p)$ of the room attached to the
condition $(L)$; this circle, which $K_o^L$ defines, is a quantity of (o9). Then
$\alpha_4$ and $\alpha_6$.

Then $\varepsilon_1$, with the ceilings
\begin{itemize}
\item $a\,e^{4L\kappa_{fat}}\le1$, where $a$ is the constant so denoted in the fat-core bound and
$\kappa_{fat}$ is the decay rate kept in that bound;
\item the ceiling carried by the constant $C_3^{fat}$, which contains
$\alpha_6^{-1}e^{2L\kappa_{fat}}$;
\item $E^*\le E_0/16$;
\item the ceilings of a further lemma of this chapter, displayed with that lemma;
\item $\varepsilon_1\le a_1R_1M_1/M$.
\end{itemize}

The group ends with
\[
\begin{aligned}
&E^*\to B_0=B_0^{\mathbf{B}}\to\varepsilon_V^*\to\bar C_{\mathsf{W}},\ \bar C_\sharp\to K_R\to
K_R^\sharp:=\max\{2K_R,K_R^o\},\\
&K_R^o:=\frac{2^{15}B_\natural C_T^wc_1^2B^{\mathsf{s}}\hat C_1}{(1-\theta_S)\beta},
\end{aligned}
\]
and then $c_V$, $c_A$, which are chosen before $\gamma$; on the run of the uniqueness theorem the
estimate of the correlation theorem at the site of the production block $(P)$ is read at the pair
radii of the condition
$(\mathrm{R}\text{-}\mathrm{b})$. The constant $K_R^o$ is used on the run of
the uniqueness theorem only; the circle $\mathsf{w}_T(k)$ which it defines is a quantity of (o9).
\item[(o6)] \emph{Classical constants, second group.} These may read $c_M$. They are chosen in the
following order.
\begin{itemize}
\item $C_i$, $C_e$, $C_P$, $C_h$, $C_h'$, $B_0^{\mathsf{s}}$, $C_{lay}$, $C_F'$,
$C^{\natural}_{LR}$, $c_b$, $y_0^{\mathsf{T}}$ (which contains $C_R^{row}$), $C^P$, $C_{IC}$,
$C_{\mathbb{C}}$, $C_\rho^{net}$, $C_\times'$, $C_\times''$, $C_{TL}$, $c_*^\sharp$, $C_\rho^\sharp$,
$C_{lay}^\times$ with its profile number $\mathsf{r}_{lay}$, $C_{TL}^\sharp$,
$C_{rim}=128C_X^+$, $K_{\mathbb{C}}$, $C_L''$ and $w^{\mathbf{B}}$.
\item Then $E_1:=4(S_\varsigma+1)E_0+4S_{\mathbf{B}}B_0$.
\item Then $E^\sharp_T$, $\mathfrak{m}_G^T$ and $C^\natural_{cl}$.
\item Then
\begin{align*}
&K_Y^{\circ}:=(K(d)+K_\circ(d)+K^\circ(d))\max K_{site},\qquad C_L''':=(1+K_Y^\circ)C_L'',\qquad
C_{cl}^T,\\
&C_T^\natural:=\frac{2E^\sharp_T(1+\bar\varepsilon^{pr}+C_F)C_{cl}^T}{E_0\mathfrak{m}_G^T},\qquad
C_T^{\mathbf{B}},\qquad E_1^z:=(5/4+4C_T^{\mathbf{B}})E_1+(5/4)B_0,\\
&C_{\mathbf{B}}^v:=2(K(d)+K_\circ+K^\circ)E_1^z/E_0 .
\end{align*}
The constant $C_{\mathbf{B}}^v$ enters
$\mathsf{c}^\star:=\mathsf{c}+C_{SM}+2+C_z+C_{\mathbf{B}}^v$.
\item The $g$-free constants of the strip: $\mathsf{K}^0$, enlarged to
$\mathsf{K}^0+K^\circ(B_0+2O(1))$ in the island format, $\mathsf{K}_{cell}^0$, $C_0^\sharp$,
$C_\sharp^0$, $A_H$ and $\alpha_*$. The bound
$\mathsf{K}_{cell}(k)\le\mathsf{K}_{cell}^0+3(N_0(g_k)+3)C_0$ is a property of the scaffold and not
a constant; only $\mathsf{K}_{cell}^0$ is chosen here. The constants $C_s$ and $C_w$ are chosen in
(o5) only.
\item Then the kernel rate $q_K$, which is distinct from the exponent $q_\flat$ of (o3).
\item Then $q_a:=e^{-c_a^{age}}$ with
\[
q_a\le\min\{q_K^2,\vartheta_\chi^2,\vartheta_\chi/2\};
\]
this is a floor on $c_a^{age}$, which is never bounded above.
\item Then $\vartheta_{dl}:=L^{-\theta_{dl}}$, and then
\[
\rho\in\bigl[\max\{q_K^{1/2},L^{-\vartheta_\sharp},\vartheta_{dl},\vartheta_\chi\},1\bigr[,
\qquad \rho>\vartheta_{dl}.
\]
\item Then $C_\rho'$, $c_\rho$, $C_T$ and $C_\sigma$.
\item Then the following constants, which read only
$(q_a,\rho,c_0,c_1,k_*,\mathfrak{a},\kappa_0,\lambda,c_2)$ and neither $C_E$ nor $C_\sigma$:
\begin{itemize}
\item $C_\Sigma^{TC}:=\sum_m(m+1)\hat{\mathsf{l}}_m^2q_a^{(m-1)_+}\rho^{-m}$;
\item $C^{(\mathfrak{s})}$, which contains $C_{\mathfrak{a}}$ (a supremum over $m$ of $\rho^{3m}$
times an $m$-dependent factor displayed where $C^{(\mathfrak{s})}$ is introduced),
$4C_\varsigma\sum_a(1+a)^3\rho^a$, and $C_1^{own}$, which contains
$q_a^{-k_*-1}$;
\item $K_\omega$.
\end{itemize}
\item Then
\[
C_\sigma^z:=(2+4C_T^{\mathbf{B}})(C_\sigma^\natural+1)+(4E_0^{-1}C_\Sigma^{TC}+1)C^{(\mathfrak{s})}.
\]
\item Then
\[
C_E^\natural:=\max\{C_V+1,\,2C_\sigma,\,2C_X^\star,\,2C_\sigma^z\},\qquad
C_E:=\tfrac87\,C_E^\natural .
\]
\item Then $C_{env}$, and then $C_{env}^z$.
\end{itemize}
The constants $C_{nc}$ and $C_{test}$ are not chosen here; they are chosen in (o9). The constant
$C_{TC}^\star$ is chosen after the two runs $\mathsf{A}$, $\mathsf{B}$.
\item[(o7)] \emph{The localisation radius.} The order is
\[
r_J\ \text{free}\;\to\;R:=4r_J\;\to\;\varrho_2,\ t_\circ,\ R_0 .
\]
\item[(o8)] \emph{The coupling ceiling $\gamma_U\le\gamma^\circ$.} Every member is either
level-wise, in the form of a supremum over $T\ge T_\gamma$, or monotone in $T$ and stated as such.
The members are the following.
\begin{enumerate}
\item[(i)] The ceilings $(\mathrm{F}_{far})^\gamma$, $(\mathrm{F}2)^<$, $(\mathrm{F}3)$ and
$(\mathrm{SL})^\flat$, and $T^{p_0-p_1}\ge4\Theta_1$.
\item[(ii)] The ceilings attached to two further displays of Ba{\l}aban's construction, and
$g(T^{\mathbb{C}})^3(LM)^4\le\varepsilon_V^*$.
\item[(iii)] $C\varrho^2e^{-\delta_KM\check R}\le1$, together with the three ceilings
$\gamma^{TS}_1$, $\gamma^{TS}_2$, $\gamma^{TS}_3$, the ceiling $(\mathrm{SL}^2)[2E_1]$ and
$\varepsilon^\natural\le\tfrac14$.
\item[(iv)] The ceilings
\begin{gather*}
4C_\Sigma c_a^{GR}\le1,\qquad C_c\varepsilon+CC_\Sigma\bar{\bar u}\le\mathfrak{r}_0,\qquad
C_h'(C^P+1)\bar{\bar u}\le\varepsilon_D,\qquad 6K_1O(\alpha^\dagger)\le\mathfrak{b}^{tol},\\
C_{lay}^{11}c\cdot\sup\mathsf{u}\le\min\{\mathfrak{r}_0,a_{R'}\},\qquad
a_{IC}\le\mathfrak{r}_e,\qquad 400LMR_*\tilde\alpha\le1 .
\end{gather*}
\item[(v)] The ceilings of two further conditions stated with their user, the ceiling on the
prefactor of the
$\kappa/10$-floor of (o4) through $\mathsf{w}_L^{-1}$, the ceilings of the one-move index, the
ceilings on $a_{rip}$ and $\mathfrak{r}_0$, and one further table ceiling stated with its user.
\item[(vi)] The ceiling on $\gamma$ under which \cite[(2.6)]{B-CMP119} holds at
$\beta_0=\beta_0^{pr}=1/7$.
\item[(vii)] The strip ceilings. First, of degree $-\infty$,
\[
(\mathrm{c}\text{-}\sharp)'\qquad \sup_{T\ge T_\gamma}C_{fib}(1+\mathsf{K}^0+\check R(g)C_0^\sharp)
e^{-\frac14\beta_\sharp\kappa(\check R(g)-4)}\le\alpha_* .
\]
Next, $\varepsilon_{br}\le1$ and
$e^{-\frac12p_0(g_k)}100^d\check R_k^{-d}\le\tfrac14\beta_v\kappa$. Next, the ceiling $(\gamma2)$,
and
\[
(\gamma3)\qquad -2(1+\beta_0^{pr})^{-1}p_0(g_k)+2+2\kappa(101\check R_k)^d\le-\tfrac32p_0(g_k).
\]
Finally, $(\gamma4)$, which concerns the factor $2^d$ at $|Z_k|$, and
\[
(N_0+3)C_0^\sharp e^{-\kappa\check R_k/40}\le1,\qquad \check R_k e^{-\kappa\check R_k}\le1 .
\]
\item[(viii)] The tolerance ceilings: $(\mathrm{c}\text{-}\mathrm{TL})^\rho$ of
\eqref{9.1:eq-tolerance-rooms-a}, and
\[
4C(C_e+C_\Sigma)\bar{\bar u}\le\mathfrak{r}^{cl},\qquad
C_W(C_e+C_\Sigma+3)\bar{\bar u}\le2,\qquad 300d\,y_0^{\mathsf{T}}\bar{\bar u}\le1 .
\]
Under the last of these, the $\bar{\bar u}$-terms of $C_\rho^{net}$ are $g$-free.
\item[(ix)] The room ceilings:
\begin{align*}
&(\mathrm{c}\text{-}\mathrm{L})\qquad 2^{15}B_\natural c_1C_4K_\Lambda'(1+\mathsf{r}_1)
e^{-\delta_1M(N_0-2)/2}\le\theta_\rho\beta,\\
&(\mathrm{c}\text{-}\mathcal{A})\qquad 2^{16}B_\natural c_1C_T^wC_4K_\Lambda'(1+\mathsf{r}_1)
4^{-(N_0-2)}\le1,\\
&(\mathrm{c}\text{-}\mathrm{TL})^\sharp\qquad 8C_{TL}^\sharp\vartheta_f^{n_c}\le1 .
\end{align*}
In $(\mathrm{c}\text{-}\mathrm{L})$ and $(\mathrm{c}\text{-}\mathcal{A})$, $N_0$ enters from below
through $L^{N_0-1}=R_{k_0+1}\ge(\log g^{-2})^{r_{pr}}$.
\item[(x)] The $\theta$-conditions of the correlation theorem, with $K_R$ replaced by $K_R^\sharp$ on
the run of the uniqueness theorem. These are $\theta_k$ and $\theta^{\mathbb{C}}_k$, level-wise,
each a supremum over $T\ge T_\gamma$. Their $g$-free member $K_R^\sharp(\Theta_\flat+1)c_A<\tfrac14$
is met by the choice of $c_A$ in (o5).
\item[(xi)] The two conditions of the $(L)$-circle: the one bounding $K_o^L\theta_j'(g)$, and
$K_o^LK_\Lambda'(1+\mathsf{r}_1)e^{-\delta_1M(N_0-2)/2}\le\theta_j'(g)$. In addition, the counting
condition on the circle inside $Z''^{(1)}_k$, at the degree in $T$ of the tower ceiling
$(\mathrm{c}\text{-}\Lambda^{re})^\sharp$ of (xiv),
and $2C_4\mathsf{s}_L\le\min\{c_4,a_{R'}\}$.
\item[(xii)] The bound
$(\alpha^\sharp)^{1/3}\le\min\{e^{-(1+\beta_s)\kappa2^d},g_k^{\kappa_0}\}$, imposed in the form
below, where $R(T)$ denotes the level radius $R_k$ of (o9) read as a function of $T=T_k$:
\begin{multline*}
(\mathrm{c}\text{-}\sharp)'''\quad \sup_{T\ge T_\gamma}
\bigl[O(1)C_{fib}(1+\mathsf{K}^0+K^\circ(B_0+2O(1))+\check R(T)C_0^\sharp)\bigr]^{1/3}\\
\times e^{-\frac1{12}\beta_\sharp\kappa(R(T)-4)}e^{\kappa_0T/2}T^{c_1+2q_\flat+\mathfrak{a}}\le1 .
\end{multline*}
\item[(xiii)] The ceiling $(\mathrm{SL}^2)[2E_1]$ read at $(2E_1^z,C_L''')$, which is called
$(\gamma\text{-}\mathbf{B}3)$. Also the ceiling $(\mathrm{SL})^\natural$ and
$A_0^2(\log X)^{-2(q_\flat-p_0)}\le1$.
\item[(xiv)] The tower ceilings:
\begin{align*}
&(\mathrm{c}\text{-}\Lambda)^\sharp\qquad 2^{20}B_\natural c_1C_T\cdot L^4C_4^\sharp K_\Lambda'
(1+\mathsf{r}_1)\ell_\circ^{-2N_0}\le1,\\
&(\mathrm{c}\text{-}\Lambda^{re})^\sharp\qquad \sup_{T\ge T_\gamma}2^{21}B_\natural c_1C_T
\Bigl[\frac{B^{\mathsf{s}}C_sc_1C_1(1+R(T)^{1/2})\,\delta'(T)}{\alpha_0(T)}
+\frac{C_{seam}\,\varepsilon(T_h)}{\alpha_0(T_h)}\Bigr]\le1,\\
&2C_4^\sharp\mathsf{s}_L\le\min\{c_4,a_{R'}\},\qquad (\mathrm{Ad})\quad
\|\mathrm{Ad}_\varpi\|\le2 ,
\end{align*}
the last being the bound $(\mathrm{Ad})$ on the adjoint action $\mathrm{Ad}_\varpi$.
\item[(xv)] The ceilings of the large-field block resummation:
\begin{itemize}
\item $(\gamma\text{-}\mathbf{B}1)$: $\mathfrak{S}(\gamma^{-2})\le1$, where $\mathfrak{S}(\cdot)$ is
the function so denoted in the large-field block resummation; this ceiling is step-uniform and of
degree $-\infty$;
\item $(\gamma\text{-}\mathbf{B}2)$:
\[
\begin{gathered}
\frac{E_1^z}{E_0}(1+\bar\varepsilon^{pr}+C_F)\bar{\mathsf{d}}(\gamma)\le\tfrac14\mathfrak{m}_G^T,\\
\bar{\mathsf{d}}(\gamma):=2C_{cl}^T(C^\natural_{cl})^{-4}\bigl(\log(\gamma^{-2}-2c_2)\bigr)^{-3c_2};
\end{gathered}
\]
\item $(\gamma\text{-}\mathbf{B}4)$:
\[
\begin{gathered}
\bar a\le\tfrac1{40},\qquad \epsilon_1\le\tfrac1{64},\qquad
(1+\epsilon_0)16\bar c_1/E_0\le\tfrac1{16},\\
S_{\mathbf{B}}\bar c_1\bigl[4E_0^{-1}E_1^zC_q+4+2(9C_\chi^\sharp/\vartheta+2)\bigr]\le B_0/4,
\end{gathered}
\]
where $\vartheta$ is the rate so denoted in the block resummation;
\item $(\gamma\text{-}\mathbf{B}5)$:
\[
\epsilon_{\mathbf{B}}K_\Lambda\le4\Lambda,\qquad
K_\Lambda:=4E_0^{-1}C_\Sigma^{TC}\cdot A_E/C_E^\natural ,
\]
where $\Lambda$ is the constant of the perturbation lemma for the system used in the proof below;
\item $(\check R_9)$ $\check R(\gamma)\ge9$.
\end{itemize}
\item[(xvi)] The ceilings under which the uniqueness theorem uses the inheritance statement. These
are:
\begin{itemize}
\item the family $(\mathrm{Q}^+)^A$ with the prefactor $K'$ of its first member replaced by $24K'$,
in sum form, with no exponent and no new floor on $q_\flat$;
\item $K_{min}\vartheta\le\mathsf{w}_0$, with $\vartheta$ the rate so denoted in the inheritance
statement, and $K_\Delta\gamma_0\le\tfrac12$;
\item $600\,\mathsf{w}_*\cdot\sup\alpha\le1$, where $\mathsf{w}_*$ is a constant of the inheritance
statement and $\alpha$ is the radius of that statement (not the number $\alpha=1/32$ of (o1));
\item the remaining ceilings of the same list, each a level-wise supremum over $T\ge T_\gamma$ or a
floor already imposed, among them the floor $C_1/C_0$ of (o4$'$).
\end{itemize}
\item[(xvii)] The ceilings of the pocket core:
\begin{align*}
&(\mathrm{c}\text{-}\mathfrak{K})\qquad 2^{20}B_\natural^2c_1^2\,\Xi_{\mathfrak{K}}\le1,\qquad
\Xi_{\mathfrak{K}}:=\Xi'+\mathfrak{l}_*+2^9L^{-R_k/2},\\
&(\mathrm{c}\text{-}\Lambda)^{\mathfrak{K}}\qquad 2^{22}B_\natural^2c_1^2C_T(1+\mathsf{r}_1)\\
&\qquad\qquad\times
\bigl[K_{\mathfrak{K}}(1+C_4)\ell_\circ^{-2N_0}+B^{\mathsf{s}}C_sC_4K_\Lambda'L^{-3(N_0-1)}\bigr]\le1 .
\end{align*}
Also $(\mathrm{c}\text{-}\mathrm{L})^{\mathfrak{K}}$ and $(\mathrm{c}\text{-}\mathcal{A})^{\mathfrak{K}}$,
which are $(\mathrm{c}\text{-}\mathrm{L})$ and $(\mathrm{c}\text{-}\mathcal{A})$ with $C_4$ replaced by
$C_4^{\mathfrak{K}}$. Finally, the counting condition
\[
1+\mathsf{w}_L(p)\ge\bigl[K_o^{\mathfrak{K}}(R_k^{1/2}+2)\delta_k'/\alpha_{0,k}\bigr]^{-1}
\quad\text{on bracket cubes meeting }\mathfrak{K}\cup Z''_k .
\]
\end{enumerate}
\item[(o9)] \emph{Quantities reading the coupling.} The order begins
\[
g\le\gamma_U\;\to\;g_k,\ T_k\;\to\;R_k,\ N_0,\ \check R_k ,
\]
and then the following quantities are chosen.
\begin{itemize}
\item $m_{hi}(k)$, $m_\sharp:=m_{hi}(k)$, and the scaffold floor $m_{lo}(k)$. The last is an integer
of the step, read after $\gamma$ and never a number chosen before it.
\item $r_A$, $\varpi_\gamma$, $\varpi_\gamma^\vee$, $N_A$, $N_Z$, the function $\mathsf{p}(g)$ and
the rates $\kappa_A^{pr}$, $\kappa_A(k)$ of Notation~\ref{9.1:not-levelwise-rate},
$\theta_k:=\theta(g_k)$ and $\theta_j'(g)$.
\item The circles $\mathsf{w}_T(k)$ and $\mathsf{w}_L(p)$. They are defined level-wise by the $(T)$-
and $(L)$-rooms through $K_R^\sharp$ and $K_o^L$ of (o5).
\item $c_a^{GR}(g):=CC_T\bar{\bar u}(g)$.
\item The alignment index $n_c^0(\gamma)\le n_c$. It is admitted only under the factor
$(q_a/\rho)^{n-\mathsf{w}}$ inside the own-source term, where the supremum over the run length is
finite, and monotonically inside the alignment bound.
\item The island sets $\mathbb{W}_m$, $\mathbb{I}_m$, $d^\circ$ and $\operatorname{fil}_m$. These
are sets of the label; no number attached to them enters any constant.
\end{itemize}
The order then ends
\[
\mathsf{w}(g),\ n_2(g),\ \ell_r,\ C_{nc}\;\to\;m_0\ge0\;\to\;C_{test}(N_w)\;\to\;n\ge n_0(N_w).
\]
\end{enumerate}
Nothing returns to $M$, $\kappa_1$, $\kappa$ or $L$: every letter reads earlier letters only.
\end{definition}

\begin{proof}
We justify, in the order of the list, the assertions made in it.

\emph{Step 1: the numbers of \textup{(o1)}.}

The window $[1/2,8/9]$ for $\theta_S$ is non-empty and lies inside $]0,1[$.

Its lower end comes from the inheritance statement. That statement runs at the spend
$\mathsf{w}_0=\min\{\tfrac12,\mathsf{c}_e\}/12$ with $\mathsf{c}_e=\theta_S$. This equals $1/24$
exactly when $\theta_S\ge\tfrac12$.

Its upper end comes from the increment lemma. Its clause at the first shell asks for
$(9/32)\theta_S\le\tfrac14$, which is $\theta_S\le8/9$. This bound comes from the one-third
reserve at depth $s=1$, through the identity $\tfrac14\beta=\tfrac12\beta\cdot2^{-1}$.

The two summation inequalities of the increment lemma, among them
$\sum_l\iota_l\le(3/16)\theta_S\beta2^{-(k-m)}$, are absolute inequalities for
$\theta_S\ge2/9$. They therefore hold on the whole window. They read nothing of $g$, and they are
not members of $\gamma_U$.

Two constants depend on $\theta_S$: the constant $K_R$ of the correlation theorem is proportional
to $(1-\theta_S)^{-1}$, and the floor on $C_0$ of the schedule lemma is proportional to
$1/((1-\theta_S)\beta)$. Since $\theta_S\le8/9<1$, both are finite on the window.

For $\theta_\rho$: since $\theta_S\ge\tfrac12$ we have $\min\{\theta_S,1-\theta_S\}=1-\theta_S$.
Moreover $1-\theta_S\in[1/9,1/2]$. Hence
\[
\theta_\rho=\tfrac14(1-\theta_S)\in[1/36,1/8].
\]

For $e_T$: since $\theta_A>3/5$, the denominator $5\theta_A-3$ is positive. Then $e_T>1$ is
equivalent to $3-\theta_A>5\theta_A-3$, that is, to $\theta_A<1$. At $\theta_A=3/4$,
\[
e_T=(9/4)/(3/4)=3.
\]

Next, $C_\Sigma=8+6/5=9.2\le10$ and $\Upsilon=3(8/7)^3L$.

The number $\beta_0^{pr}=1/7$ is admissible as a number chosen before $L$. In \cite{B-CMP119}
it is stated that ``$\beta_0>0$ can be chosen arbitrarily small, if $g$ is sufficiently small''.
This holds once $g\le\gamma_U$, because member (vi) of (o8) requires \cite[(2.6)]{B-CMP119} at
$\beta_0=1/7$. The same value is the example given in \cite{B-CMP122b}: ``for example, take
$\beta_0=1/7$''. The number $\beta_0^{pr}$ is read by the exponent of $(\mathrm{Q})^A$ in (o3), by
the floors on $\mathsf{a}$ in (o4$'$) and by $\Upsilon$ in (o5), all before (o8).

The number $k_*$ is fixed in (o1) with the scaffold's numbers. It is read by $C_{nc}$ and by
$C_1^{own}$. The number $k_*$ enters $C_1^{own}$, which enters $C^{(\mathfrak{s})}$, which in turn
enters $C_\sigma^z$ and hence $C_E^\natural$; all four are constants of (o6), chosen after $k_*$.
Hence no constant of (o6) reads a later letter through $k_*$.

\emph{Step 2: the lower bounds on $L$.}

At $d=4$ we have $\sqrt d\,3^d=2\cdot81$, and the list uses $c_g\ge1+2\cdot81$. Under this bound,
$(\mathrm{F}_\star)$ $L\ge8c^\sharp c_g^2$ implies the following floors on $L$:
\begin{itemize}
\item $3c_H$;
\item the corresponding floor of the correlation theorem;
\item $5c^{\sharp2}$, because $8c_g^2\ge5(c_g+1)$;
\item $(35/8)c_Hc^\sharp$;
\item $2K_u=13.5$;
\item $81$, which gives $K_uL^{-1/2}\le\tfrac34$.
\end{itemize}

Moreover $c^\sharp c_g^2\ge20$. Hence $4c^\sharp c_g^2\ge80$, and so
$8c^\sharp c_g^2\ge80+4c^\sharp c_g^2$. The condition $1/L\le\delta_\star/8$ is equivalent to
$L\ge80+4c^\sharp c_g^2$. Finally, $18(d+1)=90\le8c^\sharp c_g^2$.

Each of these floors reads (o1) only.

\emph{Step 3: uniformity in $M_1$.}

For Ba{\l}aban's constants, uniformity in $M_1$ is part of the published statements:
\cite[Theorem~3.1]{B-CMP99b} states: ``There exist positive constants $M_1,\delta_0,a_0,B_0$
dependent on $d$ and $L$ only
\dots\ such that for $M\ge M_1$ and for an arbitrary configuration $U$ satisfying the regularity
condition (3.35) with $M\alpha_0\le a_0$''. In that statement the letter $M$ is the parameter
written $M_1$ here.

In \cite{B-CMP102b} one reads: ``The constants $a_0,a_1,B_3$ depend on $d$ and $L$ only''.

The constant $B_\natural$ is uniform in the shell sequence. By \cite{B-CMP99b}, ``the constants in
the formulations of both theorems do not depend on the sequence $\{\Omega_j\}$''.

For the constants of (o2) introduced in this book, for instance $\Theta_1$, nothing is required:
\begin{itemize}
\item $\Theta_1$ contains the random-walk bound at the base point $\sigma\equiv0$, which reads a
partition.
\item Were one of these constants to read $M_1$, it would move, together with $\Theta_1$,
$\varpi_T$ and $\Theta_\flat$, to a place after $M_1$.
\item No floor on $\kappa$, $\kappa_1$ or $M_1$ reads them. Their readers, namely
$(\mathrm{F}_\varpi)$, $\mathfrak{b}^\flat$, $A_\flat$, $c_A$ and $\gamma$, are all later.
\end{itemize}
So the list is acyclic either way.

\emph{Step 4: $M_1$ after $\kappa_1$, and the storage rate.}

$M_1$ is placed after $\kappa_1$ because the random-walk bound contains the term
$(\delta_0-\kappa_1M_1^{-1}\cdots)|\omega|$, with $|\omega|$ the length of the walk $\omega$ in
that bound, which floors $M_1$ by $\kappa_1$. The uniformity of
Step~3 allows this placement.

For the storage rate, $\beta_s=\tfrac14$ gives
\[
(1+\tfrac12\beta_s)\kappa-\kappa_{\mathbf{B}}=\kappa/8 .
\]
Since $\beta_\sharp=1/20\le\beta_s/2$, $(\mathrm{F}_\sharp)$ gives
\[
\kappa/8\ge\beta_\sharp\kappa\ge2(\mathsf{k}_d+1)\ge\mathsf{k}_d+1 .
\]
So the requirement on
$\kappa_{\mathbf{B}}$ is met by a floor already imposed, and $\kappa_{\mathbf{B}}$ needs no ceiling.

\emph{Step 5: the floors on $M$.}

From $\tilde\delta_1M\ge4\kappa_1$ and $\kappa_1\ge4L\bar\kappa+1$ we obtain
\[
\delta_0M\ge\delta_1M\ge\kappa_1>4L\bar\kappa\ge\kappa_{out},
\]
with $\kappa_{out}$ as in \eqref{9.1:eq-radius-exponent-floors-a}. By \cite{B-CMP116}, the
walk-made pieces of $E^{(k+1)}$ have rate $\delta_0M$. By this chain they also carry the fresh
rate. The conditions $\delta_1M\ge8\log\ell_\circ$ and $e^{-\delta_1M/2}\le(2e\Upsilon)^{-1}$ of the
list are implied by the floors already imposed on $M$.

The floor $\delta_1M\ge2(q_\flat-p_1+1)\log L/r_{pr}$ is justified as follows.
\begin{itemize}
\item Since $L^{N_0-1}=R_{k_0+1}\ge(\log g^{-2})^{r_{pr}}$, the left side of the second condition of
(xi) of (o8),
$K_o^LK_\Lambda'(1+\mathsf{r}_1)e^{-\delta_1M(N_0-2)/2}\le\theta_j'(g)$, decays in $T$ at least like
$T^{-r_{pr}\delta_1M/(2\log L)}$.
\item The right side is proportional to $T^{-(q_\flat-p_1)}$, with $q_\flat$ free upward.
\item The set of admissible $T$ is therefore a half-line if and only if
$r_{pr}\delta_1M/(2\log L)\ge q_\flat-p_1$.
\item The floor gives $r_{pr}\delta_1M/(2\log L)\ge q_\flat-p_1+1$, which is more.
\end{itemize}
This floor is on $M$, chosen after $q_\flat$, and it reads $(d,L,q_\flat)$ only.

\emph{Step 6: $K_o$ and $c_M$.}

$K_o$ has the single value $2^{18}B_\natural C_{sl}^\approx\mathsf{r}_1$. It is used in the clause
\[
K_o\cdot\sup(|h|+\ell|\nabla h|)\le a^0 .
\]
The conversion of the gradient member into the sup-profile, at the constant $c_\nabla$, is carried
by the divisor $K_o(1+c_\nabla)/(35d)$ of $c_M$, and not by $K_o$. So the factor $1+c_\nabla$
enters once.

The ceilings on $\varepsilon_1$ read $E_0$, $\hat C_1$ and $\mathsf{c}$, all chosen earlier in
(o5). The bound $a\,e^{4L\kappa_{fat}}\le1$ reads the factor $\exp 5\kappa$ of \cite{B-CMP116} at the
rate kept here.

\emph{Step 7: $C_E=\tfrac87C_E^\natural$.}

The closure of the five-unknown system is carried out at $C_E^\natural$. Eliminating the
$\mathbf{B}$-part returns this system with its envelope constant at most
$C_E^\natural/(1-2\epsilon_{\mathbf{B}})$, and with $\Lambda_r$ unchanged. Under
$\epsilon_{\mathbf{B}}\le1/16$, the bound supplied with $(\gamma\text{-}\mathbf{B}4)$, this is
\[
\frac{C_E^\natural}{1-2\epsilon_{\mathbf{B}}}\le\frac{C_E^\natural}{7/8}=\tfrac87C_E^\natural=C_E .
\]

That $\Lambda_r$ is unchanged is seen as follows.
\begin{itemize}
\item The extra count of the eliminated $\mathbf{B}$-part in the $R$-entries, with $C_t$ and $A_E$
the constants of that count, is at most
\[
C_t\theta_r4E_0^{-1}e^{-p_0(g_r)}\epsilon_{\mathbf{B}}A_EC_\Sigma^{TC}\mathsf{l}_r^2e_r
\le\epsilon_{\mathbf{B}}K_\Lambda\mathfrak{r}_rC_E^\natural e_r .
\]
Here the inequality $e^{-p_0(g_r)}\mathsf{l}_r^2\le\theta_r^{-\kappa_0/2}$, proved in a corollary of
this chapter, is used, together with $q_\flat\ge p_0+1$.
\item This is compared with
$E_r^R=4C_E^\natural\hat{\mathfrak{r}}_re_r\ge64C_E^\natural\mathfrak{r}_r\Lambda_re_r$, where $e_r$
is as in \eqref{9.1:eq-rate-constants-a}.
\item So the perturbation lemma for the system applies at $\epsilon:=\epsilon_{\mathbf{B}}K_\Lambda/(64\Lambda)$.
By $(\gamma\text{-}\mathbf{B}5)$, $\epsilon\le4\Lambda/(64\Lambda)=1/16$. The lemma leaves $\Lambda$
unchanged and replaces $C_E^\natural$ by
\[
C_E^\natural/(1-2\epsilon)\le\tfrac87C_E^\natural .
\]
\end{itemize}

The quantity $\epsilon_{\mathbf{B}}$ reads $\gamma$. Hence $(\gamma\text{-}\mathbf{B}5)$ is a ceiling
on $\gamma$ of negative degree, free of the step and of the depth. The ratio $A_E/C_E^\natural$
reads $(L,\mathfrak{p},q_K,\rho,\mathsf{A},E_0,r_J)$ only.

Two forms of the same kind, with the polylogarithm written out, are
\[
E_0^{-1}2^{\kappa_0/2+c_1+2q_\flat}C_\Sigma^{TC}(A_E/C_E)\le4\Lambda(\log\gamma^{-2})^{2p_0+q_\flat+1},
\]
and
\[
4E_0^{-1}C_\Sigma^{TC}(A_E/C_E)(\log\gamma^{-2})^{c_1-2q_\flat-c_3}\le2\Lambda ,
\]
where $c_3$ is the constant entering the exponent of the polylogarithm in the second form.
They agree in kind because $c_3+2q_\flat-c_1=2p_0+q_\flat+1>0$.

Every constant of this chapter that reads $C_E$ --- among them
$c':=2c_2+\mathsf{A}C_E[(1-\rho)^{-2}+T_\Pi]$ of the uniqueness argument, with $T_\Pi$ the constant
so written there, $C_1^w$, $C_*$ and $\Pi_n$ --- is therefore valid as displayed with
$C_E=\tfrac87C_E^\natural$; no display changes form, and only the one letter $C_E^\natural$ grows by
the factor $\tfrac87$.

The order within (o6) is as follows. $C_\Sigma^{TC}$, $C^{(\mathfrak{s})}$ and $K_\omega$ read only
$(q_a,\rho,c_0,c_1,k_*,\mathfrak{a},\kappa_0,\lambda,c_2)$, so they follow $\rho$. $C_\sigma^z$ reads
them, so it follows them. $C_E^\natural$ reads $C_\sigma^z$, so it follows $C_\sigma^z$. The
constants $C_T^\natural$ and $C_L'''$ read $E_0$ and $C_L''$, so they are placed after $C_L''$, $E_1$,
$E^\sharp_T$, $\mathfrak{m}_G^T$ and $C^\natural_{cl}$.

\emph{Step 8: the members of \textup{(o8)}.}

$(\gamma3)$ at $101$ implies its form at $100$. Indeed $(100\check R_k)^d\le(101\check R_k)^d$ and
$\kappa>0$, so the left side at $100$ is at most the left side at $101$.

$(\check R_9)$ holds once $T_\gamma\ge3$: then $\check R\ge T^{r_{pr}}$ with $r_{pr}\ge2$ gives
$\check R\ge9$. The condition is monotone in $T$.

In $(\mathrm{c}\text{-}\sharp)'''$ we have $R(T)\ge T^{r_{pr}}$ with $r_{pr}\ge2$. So
$e^{-\frac1{12}\beta_\sharp\kappa R(T)}$ dominates
$e^{\kappa_0T/2}T^{c_1+2q_\flat+\mathfrak{a}}$, and the supremum has degree $-\infty$. Since
$T=\log g^{-2}$, we have $e^{-\kappa_0T/2}=g^{\kappa_0}$, so the condition delivers the factor
$g^{\kappa_0}$ in
\[
c_1^\flat(g)\le2O(1)g^{\kappa_0}(\log g^{-2})^{-c_1-2q_\flat-\mathfrak{a}} .
\]
The condition $(\mathrm{c}\text{-}\sharp)'$ likewise has degree $-\infty$.

In $(\gamma\text{-}\mathbf{B}2)$ the left member is a function of $\gamma$ alone. The index $n_c$
depends on the cutoff, by \eqref{9.1:eq-rate-constants-a}. The quantity
$\mathsf{d}_*(\gamma):=2(C_{cl}^T/C^\natural_{cl})\vartheta_2^{3n_c/4}$ is bounded by
$\bar{\mathsf{d}}(\gamma)$, using the alignment bound at $n=n_c$ and
$x_{n_c}\ge\gamma^{-2}-2c_2$. The function $\bar{\mathsf{d}}(\gamma)$ is a negative power of
$\log\gamma^{-2}$ and is monotone. With $\vartheta_\times$ in place of $\vartheta_2^{3/4}$ the same
bound holds a fortiori, since $e_T\ge1$.

In $(\mathrm{c}\text{-}\Lambda^{re})^\sharp$ the two terms have degrees $r_{pr}/2+p_1-q_\flat$ and
$p_0-q_\flat$. Both are negative under $(\mathrm{Q})'$ and $(\mathrm{Q})^A$; this is the window
``$(1/2)\nu\le p_0-p_1-1$'' of \cite{B-CMP122a}. Moreover $x_h\ge x_k$. No new floor arises.

In $(\mathrm{c}\text{-}\mathfrak{K})$ the members of $\Xi_{\mathfrak{K}}$ are:
\begin{itemize}
\item $\mathfrak{b}_*$, of degree $2r_{pr}+p_0-q_\flat<0$;
\item $\varsigma_h/\alpha_{0,h}$;
\item $(R_k^{1/2}+2)\delta_k'/\alpha_{0,k}$, of degree $r_{pr}/2+p_1-q_\flat<0$;
\item $R_k^{1/2}L^{-R_k/2}$ and $L^{-4(R_k-1)}$.
\end{itemize}
Each is a negative power of $\log g^{-2}$ or exponentially small in $R_k$. The counting condition of
(xvii) is at least $1$ under $(\mathrm{c}\text{-}\mathfrak{K})$.

The ceilings of (xvi) read neither $\check R_k$ nor a depth nor an age. In
$(\mathrm{c}\text{-}\mathrm{L})$,
$(\mathrm{c}\text{-}\mathcal{A})$, $(\mathrm{c}\text{-}\Lambda)^\sharp$ and
$(\mathrm{c}\text{-}\Lambda)^{\mathfrak{K}}$, the integer $N_0$ enters from below through
$L^{N_0-1}\ge(\log g^{-2})^{r_{pr}}$.

\emph{Step 9: nothing returns.}

The letters that might otherwise read later letters are placed as follows.
\begin{itemize}
\item The age constant has two symbols. $c_a^{age}$ is floored in (o6) and never bounded above.
$c_a^{GR}(g)$ is a quantity of (o9), bounded above by $4C_\Sigma c_a^{GR}\le1$ in (o8).
\item $\beta_0^{pr}$ and $k_*$ are numbers of (o1).
\item $\mathsf{w}_T$ and $\mathsf{w}_L$ are level-wise quantities of (o9), read only by members of
(o8) that are suprema over $T\ge T_\gamma$.
\item $C_s$ and $C_w$ are chosen once, in (o5); only $\mathsf{K}_{cell}^0$ is a constant.
\item $R^{sh}$, $R$ and $\check R_k$ are three distinct quantities.
\item $\hat C_1$ precedes $E_0$ and $K_R^o$.
\item $C_{nc}$ and $C_{test}$ are chosen in (o9), and $C_{TC}^\star$ after the two runs
$\mathsf{A}$, $\mathsf{B}$.
\item The constants of (o6) are placed as in Step~7.
\item The floors on $\kappa$ stated in (o4) read only $(d,\beta,\delta_\star,c^\sharp)$ and $L$.
\end{itemize}
With these placements, every letter of \eqref{9.1:eq-constant-order-a} reads only earlier letters.
In particular no condition after (o4) bounds $M$, $\kappa_1$, $\kappa$ or $L$, and the list contains
no cap in the sense of Definition~\ref{2.3:def-cap}.
\end{proof}

\begin{definition}[The tuned family on a fixed torus: runs, tube and standing hypotheses]
\label{9.1:def-tuned-family}
Fix the torus $\mathbb{T}_m$ and a bare point $X=g^{-2}$. Throughout this chapter we use the
following objects of the tuned-run setting:
\begin{itemize}
\item the class $\mathfrak{S}$ of tuned runs;
\item the prescribed scaffold $\Theta_X$ at the bare point $X$;
\item its ramp points $\theta_r$, the depth-$r$ inverse couplings of $\Theta_X$;
\item the step shifts $b_{(r)}$;
\item the notion of an \emph{alive} run;
\item the \emph{room} of a real centre.
\end{itemize}
These objects, and one further condition on them, are fixed in the construction of the
scaffold and are used here as fixed there.

\emph{Runs.} Let $n$ be a length and $\zeta$ a value of the running shift. The \emph{run} of
length $n$ started at $\zeta$ is the sequence $(\zeta_{(r)})_{0\le r\le n}$ defined by the
descending recursion
\begin{equation}\label{9.1:eq-tuned-family-1}
\zeta_{(n)}:=\zeta,\qquad \zeta_{(r)}:=\zeta_{(r+1)}+b_{(r)}-c_0\quad(0\le r<n).
\end{equation}
Its couplings are the numbers
\begin{equation}\label{9.1:eq-tuned-family-2}
x_{(r)}:=\theta_r-\zeta_{(r)}.
\end{equation}

\emph{Tube.} The tube of a run is given by radii $\varrho_r$, fixed in the construction of the
scaffold, that obey
\begin{equation}\label{9.1:eq-tuned-family-3}
\varrho_r\lesssim\theta_r^{1/2}(\log\theta_r)^{p_0},
\end{equation}
that is, $\varrho_r$ is at most a constant multiple of the right-hand side.

\emph{Constants of the correlation theorem.} The quantities $\mathsf{r}_k$ and $e^{J}_k$ of
the correlation theorem are re-evaluated at a lowered value of the small parameter
$\varepsilon_1$, on which they depend; the same
lowered value is used for each of the two runs compared in
Definition~\ref{9.6:def-comparison-chain}.

\emph{Standing hypotheses.} Throughout this chapter the runs of class $\mathfrak{S}$ on
$\mathbb{T}_m$, their step data and their sourced integrals satisfy the following hypotheses
(J1)--(J5). Hypotheses (J1), (J2), (J4) and (J5) are taken by statement from the correlation
theorem, read for runs of class $\mathfrak{S}$; (J3) is the
statement of the theorem on the real alive run of the prescribed scaffold $\Theta_X$. We call
these runs, with their step data and sourced integrals, the \emph{tuned family} on
$\mathbb{T}_m$.
\begin{itemize}
\item[(J1)] The step data are defined on unit balls and are holomorphic there. They do not
depend on the torus. For every argument $v$, the functionals $\beta_r[\cdot]$ of the step data
satisfy
\begin{equation*}
|\beta_r[v]|\le\mathsf{A}\,\|v\|.
\end{equation*}
The letter $\mathsf{A}$ in (J1) is a constant; it is not the run index $\mathsf{A}$ of
Definition~\ref{9.6:def-comparison-chain}. The constant $\mathsf{A}$, together with a
companion constant $C_t$ that is not used below, is a fixed constant.
\item[(J2)] The running-shift statement holds relative to the class $\mathfrak{S}$. The
numbers $a_u$ occurring in it satisfy
\begin{equation*}
a_u\lesssim\theta_u^{-3/2}(\log\theta_u)^{p_0}+\frac{\theta_u^{-(\kappa_0-4)/2}}{R},
\end{equation*}
where $R$ is the ball radius and $\kappa_0$ is the exponent entering the running-shift
statement, together with
\begin{equation*}
\sum_u a_u\le\tfrac13.
\end{equation*}
(The $a_u$ of (J2) are the numbers of the running-shift statement and are unrelated to the
$a_n$ of (J3).)
\item[(J3)] (Existence and uniqueness of the real alive run.) Let $n$ be a length and $t$ a
real number with $|t|\le r_J$, where $r_J$ is the radius of the correlation theorem. Then there
is exactly one real value $\zeta_*^n(t)$ of $\zeta$
whose run of length $n$, defined by \eqref{9.1:eq-tuned-family-1}, is alive and satisfies
$\zeta_{(0)}=t$. Along this run
\begin{equation*}
|\zeta_{(s)}|+2R\le\varrho_s
\end{equation*}
at every depth $s$. We put
\begin{equation}\label{9.1:eq-tuned-family-a}
a_n:=\theta_n-\zeta_*^n(0).
\end{equation}
\item[(J4)] Let a real centre have room $2r_J$, and let $w\in\mathcal{W}[r_J]$ be a complex
source. Then the sourced integral factors as
\begin{equation*}
Z=e^{\mathcal{E}}\tilde Z,
\end{equation*}
with exponent
\begin{equation*}
\mathcal{E}=\sum_{\mathsf{w}\le r<n}\ \sum_{X}e^{(r)}(X;w).
\end{equation*}
Here $\mathsf{w}=\mathsf{w}(g)$ is the stop depth, and $X$ runs over regions; in (J4) the
letter $X$ denotes a region, not the bare point. Each $e^{(r)}(X;\cdot)$ is analytic and
vanishes at $w=0$. It depends on $w$ only through its restriction
$w{\upharpoonright}X^{+}$ to the hull $X^{+}$ of $X$. It obeys
\begin{equation*}
|e^{(r)}(X;w)|\le C_e\,e^{-(1-\delta)\kappa\,d(X)},\qquad \kappa:=\kappa^{\flat},
\end{equation*}
where $C_e$ and $\delta$ are constants and, in (J4), the
letter $\kappa$ stands for the decay rate $\kappa^{\flat}$.
The factor $\tilde Z$ satisfies
\begin{equation*}
|\tilde Z|\le A,\qquad \tilde Z(0)\ge\frac1B,
\end{equation*}
where $A$ and $B$ are the constants of the setting of clause~(v) of
Theorem~\ref{3.1:thm-main}.
\item[(J5)] The sourced integral $Z$ and the Schwinger functions $S_{\mathsf{p}}$ are
scaffold-blind: they do not depend on the scaffold.
\end{itemize}
\end{definition}

\begin{definition}[Runs of consecutive lengths, one common source and the transplant]
\label{9.1:def-run-pair-transplant}
Let $\mathsf{A}$ be a run of length $n$ and $\mathsf{B}$ a run of length $n+1$ of the tuned family
of Definition~\ref{9.1:def-tuned-family}. At each depth $s$ the two runs use one and the same unit
lattice. Levels are indexed by the index of $\mathsf{A}$: depth $s$ is level $j=n-s$. The finest level
of $\mathsf{A}$ is $j=0$ and the finest level of $\mathsf{B}$ is $j=-1$. We write
$\mathsf{c}\in\{\mathsf{A},\mathsf{B}\}$ for the run index. For $z\in\mathbb{C}$ and $r>0$,
$D(z,r)$ is the open disc of radius $r$ about $z$ and $\bar D(z,r)$ is its closure. The number
$r_J$ is the parameter radius of Definition~\ref{9.1:def-tuned-family}.

\smallskip\noindent\emph{Lists.} A list of a run carries the structure of the lists of the
correlation theorem:
\begin{itemize}
\item it is pointed, in the sense of the correlation theorem;
\item its frozen slots are holomorphic in the complex centre $\hat{w}$;
\item its local slots are holomorphic in the source and depend on the source only through its
restriction to $Y^{+}$, in the sense of the correlation theorem;
\item its current slot is measured in the norm on current slots of lists.
\end{itemize}
Competitor lists, hybrid lists and the comparison families $N_r$ carry the same structure. For a
domain class $c$, the norm $\|\cdot\|^{c}$ of a list includes the supremum over the closed
parameter set $\bar{\mathfrak{P}}$ of the list. Here and below the italic letter $c$ and the
symbol $\flat$ denote domain classes; the italic $c$ is unrelated to the run index $\mathsf{c}$.

\smallskip\noindent\emph{Pair lists.} A \emph{pair list} is parametrised through the matching map
$\mathfrak{n}_n$. Its $\mathsf{A}$-member is taken at the parameter $v$ and its $\mathsf{B}$-member
at $\mathfrak{n}_n(v)$, and for every $v\in D(0,r_J)$ the two are matched at depth $0$. The
existence of $\mathfrak{n}_n$ follows from Rouch\'e's theorem alone. There are two kinds of
pair list.
\begin{itemize}
\item A \emph{frozen pair} has parameter set $\mathfrak{P}=D(0,r_J)$.
\item A \emph{local pair} about a real number $u$ with $|u|\le r_J/4$ is defined only for those
$n$ with $\Upsilon_n\le\tfrac14$, the restriction under which the derivative bound for the
matching map used below gives $|\mathfrak{n}_n'|\le\tfrac83$. Its radius as a list, that
is, the radius of the class $\mathcal{W}[r^{pr}_s]$ of sources on which it is defined, is
\begin{equation}\label{9.1:eq-run-pair-transplant-1}
r^{pr}_s:=\tfrac38\,r_J ,
\end{equation}
and the bound on the runs of the tuned family for sources of radius $R$ is applied with
$R=\tfrac34 r_J$ as well. For a source $W\in\mathcal{W}[r^{pr}_s]$, the $\mathsf{A}$-member is
taken at $u+W$. The source of the $\mathsf{B}$-member is
\begin{equation}\label{9.1:eq-run-pair-transplant-2}
\mathfrak{n}_n\circ(u+W)-\mathfrak{n}_n(u),
\end{equation}
which belongs to the class $\mathcal{W}^{\mathsf{B}}[r_J]$ of sources of radius $r_J$ for the
run $\mathsf{B}$, as shown below.
\end{itemize}

\smallskip\noindent\emph{One common source.} The local slots of runs, of competitors and of hybrids
are functions of one continuum source $W$, and each lattice samples $W$ at the barycentres of its
own plaquettes. For $y$ a barycentre of a plaquette of $\mathsf{A}$, let
$p^{\mathsf{A}}_y$ denote the plaquette of $\mathsf{A}$ with barycentre $y$. In formulas, for
$\mathsf{c}\in\{\mathsf{A},\mathsf{B}\}$:
\begin{equation}\label{9.1:eq-run-pair-transplant-b}
\begin{gathered}
\text{a local slot of }\mathsf{c}\text{ read on a lattice of }\mathsf{c}\text{ is a function of }
\bigl(W(x(p))\bigr)_{p},\\
p\text{ running over the plaquettes of that lattice, with one source }W
\text{ common to }\mathsf{A}\text{ and }\mathsf{B};\\
p^{\mathsf{B}}_y:=p^{\mathsf{A}}_y .
\end{gathered}
\end{equation}
No difference between a source sampled on $\mathsf{B}$ and a source sampled on $\mathsf{A}$
occurs. The sampling enters only through the weights; this is used here by statement.

\smallskip\noindent\emph{The transplant.} Let $F$ be a level-$j$ family of $\mathsf{B}$ on the
domain class $\flat$. The \emph{transplant} $\mathfrak{T}F$ is the family given by the transplant
construction for families. It has the following properties, in which $C_V$ denotes the constant
of the transplant construction for families:
\begin{equation}\label{9.1:eq-run-pair-transplant-a}
\begin{aligned}
&\text{(T1)}\quad \|\mathfrak{T}F\|^{c}\le C_V\,\|F\|^{\flat},\\
&\text{(T2)}\quad \sum_{Y\text{ from }Y'}(\mathfrak{T}F)\bigl(Y,U^{\mathsf{A}}_j(V),
J^{\mathsf{A}}_j(V)\bigr)
=F\bigl(Y',U^{\mathsf{B}}_j(V),J^{\mathsf{B}}_j(V)\bigr),\\
&\text{(T3)}\quad \beta^{\mathsf{A}}[\mathfrak{T}F]=\beta^{\mathsf{B}}[F],\\
&\text{(T4)}\quad \text{the first output of }\mathsf{B}\text{ occupies an extra slot at level }0
\text{ of }\mathsf{A}.
\end{aligned}
\end{equation}
These are to be read together with the following statements.
\begin{itemize}
\item[(T1)] $\mathfrak{T}F$ is an $\mathsf{A}$-family on the domain class $c$ and obeys the norm
bound (T1) of~\eqref{9.1:eq-run-pair-transplant-a}. It has the properties (f1) and (f2) required
of level-$j$ families in the definition of families of a run, and depends linearly on $F$. It
carries the parameters of $F$ and depends on $W$ only
through its restriction to $Y^{+}$.
\item[(T2)] Identity (T2) of~\eqref{9.1:eq-run-pair-transplant-a} holds
\begin{itemize}
\item for every real regular configuration $V$ on the whole torus of depth $s$,
\item for every torus, and
\item for every $Y'$,
\end{itemize}
the sum running over the sets $Y$ arising from $Y'$. Here $U^{\mathsf{c}}_j(V)$ and
$J^{\mathsf{c}}_j(V)$ are the level-$j$ configuration and current of the run $\mathsf{c}$
determined by $V$.
\item[(T3)] Consequently $\mathfrak{T}F$ has the background form $(f3)^{\mathrm{bg}}$ of property
(f3) in the definition of families of a run, and the coupling-increment functionals satisfy (T3)
of~\eqref{9.1:eq-run-pair-transplant-a}. The background equivalence
$\simeq^{\mathrm{bg}}$ of that definition passes through $\mathfrak{T}$, and the vacuum sums of
$\mathfrak{T}F$ and of $F$ agree.
\item[(T4)] This is the statement displayed in~\eqref{9.1:eq-run-pair-transplant-a}; it is taken
by statement from the lemma placing the first output of $\mathsf{B}$ in an extra slot at level
$0$ of $\mathsf{A}$.
\end{itemize}
\end{definition}

\begin{proof}
We verify the three claims made in the definition: the inclusion of barycentres used
in~\eqref{9.1:eq-run-pair-transplant-b}, the properties of the local pair, and the membership
of the source~\eqref{9.1:eq-run-pair-transplant-2} in $\mathcal{W}^{\mathsf{B}}[r_J]$.

\emph{Barycentres.} In physical units the finest lattice of $\mathsf{B}$, of spacing $L^{-n-1}$,
refines the finest lattice of $\mathsf{A}$, of spacing $L^{-n}$, by the odd factor
$L$. Let $a$ be the spacing of the coarser lattice. A plaquette of the coarser lattice spanned
at its corner $x$ by the directions $e_\mu,e_\nu$ has barycentre
$x+\tfrac a2(e_\mu+e_\nu)$. Consider the plaquette of the finer lattice with corner
$x+\tfrac{L-1}{2}\tfrac aL(e_\mu+e_\nu)$, which is a lattice point because $L$ is odd. Its
barycentre is
\begin{equation*}
x+\Bigl(\frac{L-1}{2}+\frac12\Bigr)\frac aL\,(e_\mu+e_\nu)=x+\frac a2\,(e_\mu+e_\nu).
\end{equation*}
Hence the plaquette barycentres of $\mathsf{A}$ are among those of $\mathsf{B}$. The
definition $p^{\mathsf{B}}_y:=p^{\mathsf{A}}_y$ in~\eqref{9.1:eq-run-pair-transplant-b}
therefore refers to a barycentre $y$ at which both lattices sample $W$.

\emph{The local pair.} Let $|u|\le r_J/4$.

First, the disc of radius $2r^{pr}_s=\tfrac34 r_J$ about $u$ lies in the parameter disc. Indeed,
if $|z-u|<\tfrac34 r_J$, then
\begin{equation*}
|z|\le|u|+|z-u|<\tfrac14 r_J+\tfrac34 r_J=r_J,
\end{equation*}
so $D(u,2r^{pr}_s)\subset D(0,r_J)$. This is why the bound on the runs of the tuned family for
sources of radius $R$ is applied with $R=\tfrac34 r_J$ as well.

Second, we locate the values of $u+W$. Let $W\in\mathcal{W}[r^{pr}_s]$. By the bound
$r_j\le 1.6\,r$ of the correlation theorem, the values of a source of class $\mathcal{W}[r]$ are
at most $1.6\,r$ in modulus. Applied with $r=r^{pr}_s$, this gives, at every point,
\begin{equation*}
|u+W|\le \tfrac14 r_J+1.6\cdot\tfrac38\,r_J=0.85\,r_J ,
\end{equation*}
so the values of $u+W$ lie in the closed disc $\bar D(0,0.85\,r_J)$.

Third, we bound the derivative of the matching map. On the closed disc $\bar D(0,0.85\,r_J)$,
the derivative bound $|\mathfrak{n}_n'|\le 1+\tfrac{20}{3}\Upsilon_n$ for the matching map on that
disc and the hypothesis $\Upsilon_n\le\tfrac14$ of the local pair give
\begin{equation*}
|\mathfrak{n}_n'|\le 1+\tfrac{20}{3}\Upsilon_n\le 1+\tfrac{5}{3}=\tfrac83 .
\end{equation*}

Fourth, we bound the Lipschitz constant of the source~\eqref{9.1:eq-run-pair-transplant-2}.
Take two points $x,y$. The segment joining $u+W(x)$ to $u+W(y)$ lies in the closed disc
$\bar D(0,0.85\,r_J)$, because that disc is convex. Integrating $\mathfrak{n}_n'$ along this
segment gives
\begin{equation*}
\bigl|\mathfrak{n}_n(u+W(x))-\mathfrak{n}_n(u+W(y))\bigr|\le\tfrac83\,|W(x)-W(y)| .
\end{equation*}
The constant $\mathfrak{n}_n(u)$ does not affect differences. Hence the Lipschitz constant
of~\eqref{9.1:eq-run-pair-transplant-2} is at most $\tfrac83$ times that of $W$. Since $W$ is
of class $\mathcal{W}[r^{pr}_s]$, the Lipschitz constant of $W$ is at most
$r^{pr}_s/r_J=\tfrac38$ times the Lipschitz bound of the class $\mathcal{W}[r_J]$. The
Lipschitz constant of~\eqref{9.1:eq-run-pair-transplant-2} is therefore at most
$\tfrac83\cdot\tfrac38=1$ times the Lipschitz bound of the class $\mathcal{W}[r_J]$.

Fifth, the supremum of this source is at most $0.875\,r_J$; this bound is used here by
statement. Together with the
Lipschitz bound just obtained, this places~\eqref{9.1:eq-run-pair-transplant-2} in
$\mathcal{W}^{\mathsf{B}}[r_J]$, which is the assertion made in the definition of the local
pair.

\emph{The transplant.} Properties (T1)--(T3) in~\eqref{9.1:eq-run-pair-transplant-a} and the
statements accompanying them are those established for the transplant construction for
families; they are used here by statement. The inputs of that construction, as listed with it,
are: its two input properties, which follow from two conditions listed there; the hypothesis
fixing the setting of this chapter; and property $(f3)^{\mathrm{bg}}$ in the definition of
families of a run. Property (T4) is taken by statement from the lemma placing the first output
of $\mathsf{B}$ in an extra slot at level $0$ of $\mathsf{A}$.
\end{proof}

\begin{definition}[Totals, components, reading units and the distance between runs]
\label{9.1:def-totals-units-distance}
Let $\mathsf{A}$ and $\mathsf{B}$ be the two runs of Definition~\ref{9.1:def-run-pair-transplant},
of lengths $n$ and $n+1$, with one common unit lattice at each depth $s$ and levels indexed by
$j=n-s$. Throughout, $\mathsf{m}\in\{\mathsf{A},\mathsf{B}\}$ denotes a run.

\emph{Representation.} For a list $G$ and a function $\varphi$ of the common unit-lattice data we
write $G\approx\varphi$ if the total of $G$ on backgrounds, taken component by component in the
sense of the statement representing totals of lists component by component, is $\varphi$;
for pointed objects this is required for each base point $y$. For frozen pair lists the relation
$G\approx\varphi$ is required up to an additive constant, on every torus. For local pair lists the
constants are included in all members, the frozen members, which are read at
$\hat{w}=W(p_y)$, included; here $p_y$ is the point attached to the base point $y$ of the pointing
of Definition~\ref{9.1:def-run-pair-transplant}, and $W(p_y)$ is the value of the source $W$ there.
With this convention the representation $\approx$ holds exactly, with zero additive constant, for
the production differences $M_r^{\bullet}$ and for transplanted
families $\mathfrak{T}F$, and the additive constant in the display of the corollary fixing the
reading units which expresses the totals in reading units is zero.

\emph{Components and reading units.} The blocks are $\bullet\in\{P,A,R\}$, and their components
are $c\in\{\mathsf{P},Q\}$ for the block $P$, $c\in\{P^{\flat},D\}$ for the block $A$, and
$c\in\{R\}$ for the block $R$. In the reading units, taken as stated in the corollary fixing the
reading units, the blocks are
\begin{equation}\label{9.1:eq-totals-units-distance-1}
\begin{gathered}
P=\bigl(E_0^{-1}\mathsf{P},\ C_t\theta_s Q\bigr),\qquad
A=\mathsf{l}_s^{-1}E_0^{-1}\bigl(P^{\flat},D\bigr),\qquad
R=C_t\theta_s\bigl(R'+R''\bigr),\\
\mathsf{l}_s:=(\log\theta_s)^{c_1},\qquad c_1:=p_0+q_\flat+1,
\end{gathered}
\end{equation}
where $\mathsf{P}$ is read at radii independent of $g$ and $A$ is pointed; $E_0$ is the constant of
the clause $\|E^{(j)}\|^{\flat}\le E_0$ of the structure statement for the lists, which the lists
of Definition~\ref{9.1:def-run-pair-transplant} carry, and $\|\cdot\|^{\flat}$ is the norm on list
members in which that clause is stated (not the norm $\|\cdot\|_{\flat}$ of sources); $C_t$ is the
constant which, together with
$\theta_s$, weights the components $Q$ and $R$ in the reading units. For a block $\bullet$,
$\|\cdot^{\bullet}\|$ denotes the maximum of the norms of its components.

\emph{The defect.} The component
\begin{equation}\label{9.1:eq-totals-units-distance-2}
D:=E(\cdot\,;W)-E_{W(p_y)}
\end{equation}
is the defect appearing in a clause of the correlation theorem: the member $E$ at the source $W$
minus the same member read at the constant source $W(p_y)$. It vanishes at constant sources, and
its level-$j$ members satisfy
\begin{equation}\label{9.1:eq-totals-units-distance-a}
\|D^{(j)}\|^{\flat}\le 2E_0 .
\end{equation}

\emph{Units.} The units are the bounds of the class $\mathfrak{K}^c$ of list bounds, as fixed by
the statement defining the class $\mathfrak{K}^c$ and its bounds, expressed in reading units:
\begin{equation}\label{9.1:eq-totals-units-distance-3}
\mathfrak{u}_s^{\mathsf{P}}:=1,\qquad
\mathfrak{u}_s^{Q}:=C_t\theta_s\mathfrak{m}_{(s)},\qquad
\mathfrak{u}_s^{P^{\flat}}:=\mathsf{l}_s^{-1},\qquad
\mathfrak{u}_s^{D}:=2\mathsf{l}_s^{-1},\qquad
\mathfrak{u}_s^{R}:=\mathfrak{r}_s .
\end{equation}
Here the per-slot $\mathsf{P}$-bound of $\mathfrak{K}^c$ divided by $E_0$ is a fixed ratio, and the
$\mathsf{P}$-unit is normalised to $1$ in reading units; $\mathfrak{m}_{(s)}$ is the depth-$s$
quantity entering the $Q$-bound of $\mathfrak{K}^c$, taken as the larger of its values for the runs
$\mathsf{A}$ and $\mathsf{B}$.

\emph{Totals and the distance between runs.} For a run $\mathsf{m}$ and a component $c$,
$\varphi_s^{\mathsf{m},c}$ is the total of the depth-$s$ component $c$ of $\mathsf{m}$ at
$U_j^{\mathsf{m}}(V)$, and $\varphi_n^{\mathsf{A},c}:=0$. We put
\begin{equation}\label{9.1:eq-totals-units-distance-4}
\mathsf{D}_s^c:=\inf\bigl\{\|G\|^c:\ G\approx\varphi_s^{\mathsf{B},c}-\varphi_s^{\mathsf{A},c}\bigr\}
\le (C_V+1)\,\mathfrak{u}_s^c .
\end{equation}
The coupling distance is
\begin{equation}\label{9.1:eq-totals-units-distance-5}
d_s:=\sup_{\bar{\mathfrak{P}}}\bigl|\zeta^{\mathsf{B}}_{(s)}-\zeta^{\mathsf{A}}_{(s)}\bigr|
\quad(s\le n),\qquad d_s:=0\quad(s>n),
\end{equation}
and its weighted tail is
\begin{equation}\label{9.1:eq-totals-units-distance-6}
d^{\sharp}_{r+1}:=\sum_{i\ge0}L^{-i/2}\bigl(d_{r+1+i}+d_{r+2+i}\bigr)\ \ge\ d_{r+1}.
\end{equation}

\emph{Rules.} Currents are never passed through the production differences $M_r^{\bullet}$; the
transplant $\mathfrak{T}$ is applied from the finer run to the coarser run only; the distances
$\mathsf{D}_s^c$ are never transplanted.
\end{definition}

\begin{proof}
\emph{The defect vanishes at constant sources.} If the source $W$ is constant, it equals $W(p_y)$
everywhere, so $E(\cdot\,;W)=E_{W(p_y)}$ and the difference
\eqref{9.1:eq-totals-units-distance-2} is zero.

\emph{The bound \eqref{9.1:eq-totals-units-distance-a}.} By the clause of the structure statement
for the lists, which the lists of Definition~\ref{9.1:def-run-pair-transplant} carry, every member
$E^{(j)}$, frozen or local, satisfies $\|E^{(j)}\|^{\flat}\le E_0$. The member
$E(\cdot\,;W)$ is local, and the member $E_{W(p_y)}$ is the frozen member read at
$\hat{w}=W(p_y)$. Applying the clause to each of them and using the triangle inequality for
$\|\cdot\|^{\flat}$ gives
\begin{equation*}
\|D^{(j)}\|^{\flat}\le\|E^{(j)}(\cdot\,;W)\|^{\flat}+\|E^{(j)}_{W(p_y)}\|^{\flat}\le E_0+E_0=2E_0 .
\end{equation*}
In the reading units of \eqref{9.1:eq-totals-units-distance-1} the component $D$ is multiplied by
$\mathsf{l}_s^{-1}E_0^{-1}$, so the bound \eqref{9.1:eq-totals-units-distance-a} becomes
$2\mathsf{l}_s^{-1}$, the value of the unit $\mathfrak{u}_s^{D}$ of
\eqref{9.1:eq-totals-units-distance-3}.

\emph{The bound in \eqref{9.1:eq-totals-units-distance-4}.} The inequality
$\mathsf{D}_s^c\le(C_V+1)\mathfrak{u}_s^c$ is taken by statement. The bounds of the class
$\mathfrak{K}^c$, expressed in reading units, are the units
\eqref{9.1:eq-totals-units-distance-3}. For $c=Q$ the unit is formed with the larger of the two
runs' values $\mathfrak{m}_{(s)}$, so that one unit serves both runs. For $s=n$ one has
$\varphi_n^{\mathsf{A},c}=0$, so that $\mathsf{D}_n^c$ is the infimum of $\|G\|^c$ over
$G\approx\varphi_n^{\mathsf{B},c}$.

\emph{The inequality in \eqref{9.1:eq-totals-units-distance-6}.} Every $d_s$ is non-negative, being
a supremum of absolute values or zero. Hence every term of the series is non-negative. Since
$d_s=0$ for $s>n$, only the finitely many terms with $r+1+i\le n$ can be non-zero, so the sum is
finite. Dropping all terms with $i\ge1$, and the summand $d_{r+2}$ of the term $i=0$, gives
\begin{equation*}
d^{\sharp}_{r+1}\ \ge\ L^{0}\bigl(d_{r+1}+d_{r+2}\bigr)\ \ge\ d_{r+1}. \qedhere
\end{equation*}
\end{proof}

\begin{definition}[Hybrid lists and the class constant]\label{9.1:def-hybrid-lists}
Let $\mathsf{A}$ and $\mathsf{B}$ be the two runs compared, $\mathsf{A}$ of length $n$ and
$\mathsf{B}$ of length $n+1$. We use the
totals, the components $c$ of the blocks $P$, $A$, $R$, the reading units $\mathfrak{u}_s^c$ and
the component norms $\|\cdot\|^c$ of Definition~\ref{9.1:def-totals-units-distance}. We write
$\mathfrak{s}^{\mathsf{A}}$ for the list of $\mathsf{A}$, $\mathfrak{e}_r^{\mathsf{A},\circ}$ for its
depth-$r$ slot of block $\circ\in\{P,A,R\}$, and $\zeta^{\mathsf{B}}_{(s)}$ for the running shift of
$\mathsf{B}$ at depth $s$. Fix a depth $r$ and a block $\bullet\in\{P,A,R\}$.

\emph{The order of the blocks.} Inside depth $r$ the production is a finite sequence of sub-steps,
namely the sub-steps of the sourced step of the correlation theorem.
Besides the slots $s>r$, each sub-step reads only those depth-$r$ blocks that are produced before it.
This defines an order $\prec$ on $\{P,A,R\}$.

For a block $\circ\in\{P,A,R\}$ we write $(\circ)$, that is $(P)$, $(A)$ or $(R)$, for the response
inequality of that block; the letter in parentheses names an inequality, the letter without
parentheses a block.
What follows does not depend on which order the production uses. The reason is that the response
inequalities $(A)$ and $(R)$ contain every same-depth entry, and an entry that is void only enlarges a
right-hand member.

\emph{Competitors.} Let $C_V$ be the constant of the uniform bound on frozen lists used in the proof
below, fixed in the order of choice before the class
constant below. A family of \emph{competitors} is a family
$(G_s^c)_{r_{\mathrm{in}}\le s\le n}$,
indexed by the components $c$ and by the depths $s$ from the lowest depth $r_{\mathrm{in}}$ of the
comparison up to $n$,
with
\begin{equation}\label{9.1:eq-hybrid-lists-1}
\|G_s^c\|^c\le (C_V+1)\,\mathfrak{u}_s^c .
\end{equation}
We write $G_s=(G_s^c)_c$ for the family over the components at depth $s$.
For a block $\circ$ we write $G_r^{\circ}$ for the part of $G_r$ belonging to $\circ$.

\emph{The hybrid list.} The \emph{hybrid list} $\mathfrak{s}^{\natural}_{r,\bullet}$ has the slots
\begin{equation}\label{9.1:eq-hybrid-lists-2}
\mathfrak{s}^{\natural}_{r,\bullet}
:=(\mathfrak{s}^{\mathsf{A}}_s+G_s)_{s>r}\;\oplus\;
\{\mathfrak{e}_r^{\mathsf{A},\circ}+G_r^{\circ}:\ \circ\prec\bullet\},
\end{equation}
and the $\zeta$-slots $\zeta^{\mathsf{B}}_{(s)}$ for $s\le n$. Here $\oplus$ denotes concatenation of
lists, and the slot $n$ consists of $G_n$ alone.

\emph{The production difference.} Let $F_r^{\mathsf{A},\bullet}$ be the depth-$r$ block $\bullet$
produced by $\mathsf{A}$ from its own list. Let $\mathsf{F}_r^{\mathsf{A},+,\bullet}$ be the
production map of $\mathsf{A}$ at depth $r$ for the block $\bullet$, extended by the extra slot of
the hybrid list. We put
\begin{equation}\label{9.1:eq-hybrid-lists-3}
M_r^{\bullet}:=\mathsf{F}_r^{\mathsf{A},+,\bullet}\bigl(\mathfrak{s}^{\natural}_{r,\bullet}\bigr)
-F_r^{\mathsf{A},\bullet}.
\end{equation}

\emph{The class constant.} Let $\mathfrak{K}^c$ be the class of list bounds on the component $c$.
The \emph{class constant} $C_{\mathfrak{K}}$ is chosen after $C_V$ and before the constants
$\varepsilon_1$, $c_\varrho$, $C_\sigma$, $C_E$ and $\gamma$, which are chosen later in the order of
choice of Definitions~\ref{2.3:not-stages-first}--\ref{2.3:not-stages-last};
it enters only through upper bounds.

The response inequalities $(P)$, $(A)$ and $(R)$ are stated for lists in $(C_V+3)\,\mathfrak{K}^c$.
They are used on the class $C_{\mathfrak{K}}\mathfrak{K}^c$:
\begin{equation}\label{9.1:eq-hybrid-lists-a}
\text{$(P)$, $(A)$, $(R)$, stated on $(C_V+3)\,\mathfrak{K}^c$, hold on
$C_{\mathfrak{K}}\,\mathfrak{K}^c$.}
\end{equation}
The hybrid list lies in this class:
$\mathfrak{s}^{\natural}_{r,\bullet}\in C_{\mathfrak{K}}\,\mathfrak{K}^c$.
The totals, the $b$-slots and the $\zeta$-slots of $\mathfrak{s}^{\natural}_{r,\bullet}$ are those of
$\mathsf{B}$.
\end{definition}

\begin{proof}
\emph{Totals and slots.} That the totals, $b$-slots and $\zeta$-slots of
$\mathfrak{s}^{\natural}_{r,\bullet}$ are those of $\mathsf{B}$ is part (a) of the statement on the
transplant $\mathfrak{T}F$ of $\mathsf{B}$-families, applied to the hybrid list
\eqref{9.1:eq-hybrid-lists-2}.

\emph{Frozen slots.} We write $E_0$ for the unit in which $\mathfrak{K}^c$ bounds a frozen entry,
$\mathsf{E}$ for the constant in which it bounds the entries $\mathsf{P}^{(\le j)}$ of the
$\mathsf{P}$-component, for each index $j$ of these entries, measured in their norms $|\cdot|_j$,
and $\mathfrak{m}$ for the quantity,
depending on the coupling, that enters these bounds beside $E_0$. The coupling is taken below the
ceiling $(C_V+1)\mathfrak{m}\le E_0$.
Each frozen slot is at most $(C_V+2)$ times the bound of $\mathfrak{K}^c$. For a
frozen entry $E$ and the entries $\mathsf{P}^{(\le j)}$, the uniform bound on frozen lists, in
its second clause, gives
\begin{equation}\label{9.1:eq-hybrid-lists-4}
\|E\|\le (C_V+2)(E_0+\mathfrak{m})\le (C_V+3)E_0,
\qquad
|\mathsf{P}^{(\le j)}|_j\le 2(C_V+3)\,\mathsf{E}.
\end{equation}

\emph{The local member.} For a component $c$ we write $G^{c}$ for the component $c$ of a
competitor, and $E^{\mathsf{A}}(\cdot;W)$ for the local entry of $\mathsf{A}$, pointed at $y$ in the
sense of Definition~\ref{9.1:def-totals-units-distance}, as a function of its argument $W$; the
subscript $W(p_y)$ on the bracket denotes evaluation of the frozen members at $\hat w=W(p_y)$.
The local member is
\begin{equation*}
E^{\mathsf{A}}(\cdot;W)+\bigl[G^{P^{\flat}}+G^{Q}\bigr]_{W(p_y)}+G^{D}.
\end{equation*}
It is at most
$E_0+(C_V+1)(3E_0+\mathfrak{m})$.

Under the ceiling $(C_V+1)\mathfrak{m}\le E_0$,
\begin{equation}\label{9.1:eq-hybrid-lists-5}
E_0+(C_V+1)(3E_0+\mathfrak{m})
\le E_0+3(C_V+1)E_0+E_0=(3C_V+5)E_0 .
\end{equation}
With \eqref{9.1:eq-hybrid-lists-4}, this places $\mathfrak{s}^{\natural}_{r,\bullet}$ in
$C_{\mathfrak{K}}\mathfrak{K}^c$.

\emph{Use of the response inequalities on $C_{\mathfrak{K}}\mathfrak{K}^c$.} Display
\eqref{9.1:eq-hybrid-lists-a} is the lemma fixing the class constant: the class constant
$C_{\mathfrak{K}}$, fixed after $C_V$ and before $\varepsilon_1$, $c_\varrho$, $C_\sigma$, $C_E$ and
$\gamma$, enters the proof of the response inequalities at finitely many places and only through
upper bounds. These places are as follows. The reading of the
response inequalities on $C_{\mathfrak{K}}\mathfrak{K}^c$ rests on item (v) of the statement
establishing the response inequalities $(P)$, $(A)$, $(R)$, whose proof is the proof of part (a)
of the transplant statement. That proof
meets the class
constant only through ceilings on $\gamma$ and on $\varepsilon_1$ and through one factor of
$C_\sigma$. All of these are chosen after $C_V$; since $C_{\mathfrak{K}}$ is chosen after $C_V$ and
before $\varepsilon_1$, $C_\sigma$ and $\gamma$, these ceilings are taken with $C_{\mathfrak{K}}$.
Frozen lists need only the
constant $C_V+3$, by \eqref{9.1:eq-hybrid-lists-4}. Hence, by the lemma fixing the class constant,
$(P)$, $(A)$ and $(R)$, stated on
$(C_V+3)\mathfrak{K}^c$, hold on $C_{\mathfrak{K}}\mathfrak{K}^c$. They therefore apply to
$\mathfrak{s}^{\natural}_{r,\bullet}$.
\end{proof}

\begin{definition}[The rates $q$, $q_a$, $\rho$ and the constants of the two-run estimate]
\label{9.1:not-rate-constants}
In the two-run estimate $n$ denotes the cutoff, the length of the shorter of the two runs
compared (the other has length $n+1$), and $r$ a depth. We use
$X=g^{-2}$ and the
inverse couplings $\theta_r=X+c_0r$ of the prescribed scaffold, the scaffold radius
$\check{R}_{(r)}$ at depth $r$, the index letters $\theta_2$, $\vartheta_{dl}$, $\vartheta_\chi$ of
Notation~\ref{9.1:not-rho-window}, and the index exponents $\hat{\alpha}$, $\beta_1$, $\beta_2$
of the same notation.

\emph{Three letters $q$.} They have different roles.
\begin{itemize}
\item[(a)] The \emph{kernel rate} is $q:=q_K$, the rate of the kernel bound in the sharpened
kernel corollary. The subscript is part of the name and does not refer to the number $K$ of
renormalisation steps, and this $q$ is not Ba{\l}aban's auxiliary parameter $q$ of the tuple
$\mathfrak{p}$. It is chosen in the interval
\begin{equation}\label{9.1:eq-rate-constants-1}
q\in\bigl]\max\{L^{-\hat{\alpha}},\,L^{-\beta_1},\,L^{-1/2}\},\,1\bigr[ .
\end{equation}
\item[(b)] The \emph{age rate} is $q_a:=e^{-c_a^{age}}$, where $c_a^{age}$ is the age
constant. It is subject to $q_a\le q^2$.
\item[(c)] The letter $q_\flat$ denotes the radius exponent. It is an exponent, not a rate.
\end{itemize}

\emph{The window letter.} We put $\vartheta_\sharp:=\min\{\beta_2-\beta_1,\tfrac12\beta_1\}$.
The letter $\vartheta_R$ denotes the rate delivered by the supply estimate; it equals
$\vartheta_{dl}=L^{-\theta_2/2048}$. The letter $\rho$, also written $\vartheta_1$, is chosen
in the window
\begin{equation*}
\rho\in\bigl[\max\{q^{1/2},\,L^{-\vartheta_\sharp},\,\vartheta_{dl},\,\vartheta_\chi\},\,1\bigr[,
\qquad \rho>\vartheta_{dl}.
\end{equation*}
This is the window \eqref{9.1:eq-rho-window-a} of Notation~\ref{9.1:not-rho-window}. The
hypotheses placed on $\rho$ are monotone in $\vartheta_1$.

\emph{The envelope functions and constants.} We put
\begin{equation*}
\mathfrak{r}_r:=C_t\,\theta_r^{1-\kappa_0/2},\qquad
\Lambda_r:=\Lambda\,(\log\theta_r)^{c_3},\quad c_3:=3p_0+2,\qquad
\ell_r:=2+\log_L\bigl(M\check{R}_{(r)}\bigr).
\end{equation*}
Together with these we put the quantities in the first three lines of the following display;
its last line is the form in which the alignment index $n_c$ may be used, explained below:
\begin{equation}\label{9.1:eq-rate-constants-a}
\begin{gathered}
e_r:=\rho^{\,n-r}+\Pi_n,\qquad \mathsf{n}_r:=r+1+(1-\rho)^{-1},\qquad
c_\rho:=\frac{5}{\rho(1-\rho)},\\
\hat{\mathfrak{r}}_r:=2\mathfrak{r}_r\Lambda_r\Bigl(\mathsf{A}\Bigl[\frac{\mathsf{n}_r}{1-\rho}
+(1-\rho)^{-2}\Bigr]+5(1-\rho)^{-1}+8\Bigr),\qquad
\Pi_n:=\frac{C_T\,\hat{\mathfrak{r}}_n}{1-\rho^{1/2}},\\
C_E:=\tfrac{8}{7}\,C_E^{\natural},\qquad
C_E^{\natural}:=\max\{C_V+1,\,2C_\sigma,\,2C_X^{\star},\,2C_\sigma^{z}\},\\
q_a^{\,n-r-n_c}\,e^{-p_0(g_{(r)})}\le C\,\mathfrak{r}_r\Lambda_r\,\rho^{\,n-r}.
\end{gathered}
\end{equation}
Since $0<\rho<1$, the numbers $1-\rho$ and $1-\rho^{1/2}$ are positive, so the quantities in the
first three lines of
\eqref{9.1:eq-rate-constants-a} are finite. Here $\mathsf{A}$ is the constant that bounds the
coupling-increment functional at depth $r$ in absolute value by $\mathsf{A}$ times the norm of
its argument;
$C_V$ is the constant in the bound $(C_V+1)\mathfrak{u}^c_s$ on competitors in
Definition~\ref{9.1:def-hybrid-lists}; $C_t$ and $\kappa_0$ are the constant and the exponent
of the weight $\mathfrak{r}_r$ of the rows of the response component $R$; $\Lambda$ is the
prefactor of $\Lambda_r$; and $C_T$ is the prefactor of the floor $\Pi_n$. The letters
$C_\sigma^{z}$ and $C_X^{\star}$ are constants of clause (b) of the consistency corollary of
the zone-resolved bound. They are fixed in Notation~\ref{9.1:not-constant-order}, before
$C_E^{\natural}$.

\emph{The constant $C_\sigma$.} Let $\varrho_k=c_\varrho L^{\vartheta_\sharp k}$ be the radius
of the dial domain at level $k$, with prefactor $c_\varrho$; this $\varrho_k$ is the dial
radius, not Ba{\l}aban's radius $\varrho$ of the tuple $\mathfrak{p}$ nor his
$\varrho_2=\varrho/2$; here the subscript is always a level. Let $k_\sharp$ be the
least $k$ with $\varrho_k\ge 2$. We put
\begin{equation*}
C_\sigma:=\max\Bigl\{\frac{7\,(E^{*}+E^{\sharp}+K_2')}{c_\varrho E_0\,\rho},\,
(C_V+1)\,\rho^{-k_\sharp}\Bigr\},
\end{equation*}
where $E_0$, $E^{*}$, $E^{\sharp}$ and $K_2'$ are constants fixed before $C_\sigma$.
The factor $7$ is $3+4$. The term $3$ comes from the frozen component $P^{\flat}$. The term
$4$ comes from clause (c) of the two-run localisation estimate, applied to the difference
$N^E(\cdot;W)-N^E_{W(p_y)}$ of two functionals of that estimate.

\emph{Convention for $C_E$.} The following proofs are carried out with $C_E^{\natural}$ in
place of $C_E$:
\begin{itemize}
\item the inductive proof of the envelope proposition;
\item the two-run localisation estimate and the sharpened step bound;
\item the proof of the uniqueness theorem.
\end{itemize}
Every occurrence of $C_E$ in these proofs, in the inductive hypothesis of the induction
proving the envelope proposition, and in the statements in which $\hat{\mathfrak{r}}_r$ occurs,
stands for $C_E^{\natural}$.

\emph{Where the alignment index may enter.} The alignment index $n_c$ depends on the cutoff;
it is of order $\log\log n$ by clause (i) of the statement on the three rates $q$, $q_a$,
$q_\flat$ and the alignment index, which is the statement from which the three letters $q$
above are taken.
It may be used only in the form of the inequality in the last line of
\eqref{9.1:eq-rate-constants-a},
where $C$ is the constant of clause (iii) of the small-field bound on which the response rows
rest; here $p_0(\cdot)$ is the degree function, $p_0(g')=A_0(\log g'^{-2})^{p_0}$, and
$g_{(r)}$ is the coupling at depth $r$. This form may appear in the rows of the component $R'$
and in the surplus of the
concluding summation. It never enters the constants $\Lambda$, $C_\sigma$, $C_E$, $c'$, nor the
constant $C_{TC}$ of the comparison theorem. It may in addition enter at two further places:
\begin{itemize}
\item[(i)] inside the own-source term $\mathfrak{s}_j$ of a zone $j$ of the zone-resolved bound
(there the run length is written $\mathfrak{n}$), which contains, for the levels $\ell$ of that
zone, the terms $c_1^{\flat}(g_\ell)\,q_a^{\,\ell-n_c}$ (this $\ell$ is an index, unrelated to
$\ell_r$);
\item[(ii)] monotonically, through the lower bound $n_c\ge n_c^0(\gamma)$, inside the
alignment condition.
\end{itemize}
It does not enter the condition of Notation~\ref{9.1:not-constant-order} whose left member is
$\bar{\mathsf{d}}(\gamma)$.

The alignment statement, which defines $n_c$, carries an exponent $c_2$ subject to
$c_2\ge 2p_0$. That exponent is written $c_2^{IS}$.

\emph{Assertions.}
\begin{itemize}
\item[(A1)] The rates satisfy
\begin{equation}\label{9.1:eq-rate-constants-2}
q_a\le q^2<q\le\rho^2<\rho .
\end{equation}
Moreover $\rho>\vartheta_R$. Hence $\rho$ lies strictly above $q$, $q_a$ and $\vartheta_R$.
\item[(A2)] $C_E^{\natural}$ is at least each of $C_V+1$, $2C_\sigma$, $2C_X^{\star}$ and
$2C_\sigma^{z}$.
\item[(A3)] Every displayed estimate in the three proofs listed under the convention for $C_E$
above, including clauses (a)--(d) of the uniqueness theorem, is the statement proved, when
read with $C_E$ as in \eqref{9.1:eq-rate-constants-a}.
\item[(A4)] At the first of the two further places, the constant $C^{(\mathfrak{s})}$ that
the terms $\mathfrak{s}_j$ contribute does not depend on the cutoff.
\item[(A5)] The inequality in the last line of \eqref{9.1:eq-rate-constants-a} holds.
\end{itemize}
\end{definition}

\begin{proof}
(A1). By \eqref{9.1:eq-rate-constants-1} we have $0<L^{-1/2}<q<1$, so $q^2<q$. With the
condition $q_a\le q^2$ this gives $q_a\le q^2<q$. The window gives $\rho\ge q^{1/2}$, hence
$\rho^2\ge q$. Since $0<\rho<1$, also $\rho^2<\rho$. This proves
\eqref{9.1:eq-rate-constants-2}. In particular $q_a<\rho$ and $q<\rho$. The window also
gives $\rho>\vartheta_{dl}$, and $\vartheta_{dl}=\vartheta_R$. The inequality $q\le\rho^2$
obtained here is the form in which the lower bound $\rho\ge q^{1/2}$ is used in the induction
proving the envelope proposition.

(A2). This holds by the definition of $C_E^{\natural}$ as a maximum.

(A3). First, carrying out the listed proofs with $C_E^{\natural}$ in place of $C_E$ is
admissible.
\begin{itemize}
\item Those proofs require their constant to be at least $\max\{C_V+1,2C_\sigma\}$.
\item The zone-resolved bound requires it to be at least $2C_X^{\star}$ and $2C_\sigma^{z}$.
\end{itemize}
Both requirements hold by (A2).

Next, clause (b) of the consistency corollary of the zone-resolved bound returns the
conclusions of those proofs with envelope constant at most
\begin{equation*}
\frac{C_E^{\natural}}{1-2\epsilon_{\mathbf{B}}}\le\tfrac{8}{7}\,C_E^{\natural}=C_E .
\end{equation*}
Here $\epsilon_{\mathbf{B}}$ is the small constant of that bound. The inequality is
equivalent to $1-2\epsilon_{\mathbf{B}}\ge 7/8$, that is, to $\epsilon_{\mathbf{B}}\le 1/16$.
It is provided by two of the conditions imposed on the zone-resolved bound. Hence each
displayed estimate listed in (A3), read with $C_E$ as in \eqref{9.1:eq-rate-constants-a}, is
the statement proved.

(A4). At the first of the two further places, the remark accompanying $\mathfrak{s}_j$ takes a
supremum over the run length $\mathfrak{n}$. That supremum is the constant $C_1^{own}$. It is
finite and depends only on
\begin{equation*}
\bigl(q_a,\,k_*,\,\mathfrak{a},\,\kappa_0,\,\lambda,\,c_2,\,\rho,\,C^{\natural}_{cl}\bigr),
\end{equation*}
the factor $(C^{\natural}_{cl})^{\mathfrak{a}/c_2}$ being carried along; here $c_2$ is not the
exponent $c_2^{IS}$ of the alignment statement. The constant $C^{(\mathfrak{s})}$ is bounded in
terms of $C_1^{own}$ and these letters only, so it does not depend on the cutoff. It enters
$C_\sigma^{z}$, and
$C_\sigma^{z}$ enters $C_E^{\natural}$; therefore the alignment index does not reach
$C_\sigma^{z}$ or $C_E^{\natural}$ through $\mathfrak{s}_j$.

At the second place, $n_c$ enters only through $n_c\ge n_c^0(\gamma)$, and the alignment
condition is monotone in $n_c$. The condition with left member $\bar{\mathsf{d}}(\gamma)$ does
not contain $n_c$.

(A5). The inequality in the last line of \eqref{9.1:eq-rate-constants-a} is supplied by clause
(iii) of the small-field bound on which the response rows rest, together with $q_a<\rho$ from
(A1).
\end{proof}

\begin{definition}[Transplanted one-lattice response bounds and the exact shift identity]
\label{9.1:def-transplant-response}
Let $\mathsf{A}$ and $\mathsf{B}$ be two runs, of lengths $n$ and $n+1$, compared in this chapter.
Let $r<n$. Let the order $\prec$ of the sub-steps inside depth $r$, the hybrid lists
$\mathfrak{s}^{\natural}_{r,\bullet}$ for $\bullet\in\{P,A,R\}$, and the competitors
$(G^{c}_s)_{r_0\le s\le n}$ be as in Definition~\ref{9.1:def-hybrid-lists}. The norms
$\|\cdot\|$ are those of Definition~\ref{9.1:def-hybrid-lists}.

We use the following notation.
\begin{itemize}
\item $G_s^{P}$, $G_s^{A}$ and $G_s^{R}$ are the entries of a competitor $G$ in the blocks $P$, $A$
and $R$ at depth $s$.
\item $G_s^{\mathsf{P}}$ and $G_s^{Q}$ are the components $\mathsf{P}$ and $Q$ of the entry $G_s^{P}$.
\item $M_r^{\bullet}$ is the production difference of block $\bullet$ at depth $r$, formed on the
hybrid list $\mathfrak{s}^{\natural}_{r,\bullet}$.
\end{itemize}

The remaining letters are the following.
\begin{itemize}
\item $q$ is the kernel rate of Notation~\ref{9.1:not-rate-constants}.
\item $g_{(s+1)}$ is the coupling of depth $s+1$.
\item $d_{s+1}$ is the coupling distance between the two runs at depth $s+1$, and $d^{\sharp}_{r+1}$
is the weighted tail of the coupling distances $(d_s)_{s\ge r+1}$.
\item $b^{\mathsf{A}}_s$ and $b^{\mathsf{B}}_s$ are the step shifts of the two runs at depth $s$.
\item $\beta^{\mathsf{A}}[\cdot]$ is the step functional of the run $\mathsf{A}$, and the constant
$\mathsf{A}$ is the one in the bound $|\beta^{\mathsf{A}}[\cdot]|\le\mathsf{A}\|\cdot\|$ for this
functional; in row (b) below, the factor $\mathsf{A}$ on the right is this constant, not the run.
\item $\alpha_r,\alpha_r^{A},\varkappa,\varkappa_A,\varkappa'_r,\varkappa''_r,\mathfrak{r}_r,\Lambda_r$
are the constants of the rows (P), (A) and (R), and $\ell_r$ is the length of the window
$r<s<r+\ell_r$ in row (R); item (a) below fixes them.
\end{itemize}

For a competitor $G$ and the coupling distances $d=(d_{s})$ we put
\begin{equation}\label{9.1:eq-transplant-response-1}
\begin{split}
B_r(G,d):={}&\sum_{r\le s<n}q^{s-r}g^2_{(s+1)}d_{s+1}
+\sum_{s\ge r}q^{s-r}\bigl(\|G_s^{P}\|+\|G_s^{A}\|\bigr)\\
&+\sum_{s>r}\|G_s^{R}\|.
\end{split}
\end{equation}

The pair $\mathsf{A}$, $\mathsf{B}$ satisfies the \emph{transplanted one-lattice response bounds and
the exact shift identity} if, for every $r<n$ and every competitor, conditions (a) and (b) below
hold. The relations referred to in them are the following.
\begin{equation}\label{9.1:eq-transplant-response-a}
\begin{aligned}
&\text{(P)}\quad \|M_r^{P}\|\le \alpha_r d^{\sharp}_{r+1}
+\varkappa\sum_{s>r}\bigl(q^{s-r}\|G_s^{P}\|+\|G_s^{R}\|\bigr)
+\varkappa''_r\sum_{s>r}q^{s-r}\|G_s^{A}\|;\\
&\text{(A)}\quad \|M_r^{A}\|\le \alpha_r^{A}d^{\sharp}_{r+1}
+\varkappa_A\Bigl[\|G_r^{P}\|
+\sum_{s>r}\Bigl(q^{s-r}\bigl(\|G_s^{P}\|+\|G_s^{A}\|\bigr)+\|G_s^{R}\|\Bigr)\Bigr]\\
&\phantom{\text{(A)}\quad \|M_r^{A}\|\le{}}+\|G_r^{R}\|;\\
&\text{(R)}\quad \|M_r^{R}\|\le \mathfrak{r}_r\Lambda_r B_r(G,d)
+\varkappa'_r\sum_{r<s<r+\ell_r}\|G_s^{R}\|;\\
&\text{(b)}\quad b^{\mathsf{B}}_s-b^{\mathsf{A}}_s
=\beta^{\mathsf{A}}[G_s^{\mathsf{P}}]+\beta^{\mathsf{A}}[G_s^{Q}],
\qquad |b^{\mathsf{B}}_s-b^{\mathsf{A}}_s|\le\mathsf{A}\|G_s^{P}\|.
\end{aligned}
\end{equation}

\begin{itemize}
\item[(a)] The production differences obey the rows (P), (A) and (R) of
\eqref{9.1:eq-transplant-response-a}, with $B_r(G,d)$ as in
\eqref{9.1:eq-transplant-response-1}. The coefficients are those supplied by the one-lattice
Lipschitz step, by its closeness condition with the constant $C_L^{\sharp}$ replaced by
$4C_L^{\sharp}$, by its statements concerning row (A), two runs and the class constant, and by its
corollary fixing the kernel rate $q$. They are taken on the set $c$ on which the competitors
$(G^{c}_s)$ and their norms are read in Definition~\ref{9.1:def-hybrid-lists}, and on the class on
the tube on which the two runs are compared. These coefficients do not depend on $K$, $n$ or $r$.
The running-shift entries of the lists are carried by the letter $\mathsf{z}^{\sharp}$, the same
as in the one-lattice Lipschitz step; on lists with local members the step at depth $r$ reads
every earlier coupling. For frozen pairs the rows (P) and (A) hold with $d_{r+1}$ in place of
$d^{\sharp}_{r+1}$; since $d_{r+1}\le d^{\sharp}_{r+1}$, the displayed rows, with the one letter
$d^{\sharp}_{r+1}$, cover these pairs as well.
\item[(b)] The identity in row (b) of \eqref{9.1:eq-transplant-response-a} holds exactly, for every
competitor $(G^{c}_s)_{r_0\le s\le n}$ of Definition~\ref{9.1:def-hybrid-lists} and at each of its
depths $s$. The bound in the same row holds as well. In row (b) the components
$G_s^{\mathsf{P}}$ and $G_s^{Q}$ of the $P$-entry stand in the identity, and the whole entry
$G_s^{P}$ stands in the bound.
\end{itemize}
\end{definition}

\begin{definition}[The two-run comparison families and their bounds]\label{9.1:def-comparison-families}
Let $C_\sigma$ and $\rho$ be constants. We use the pairs of runs $(\mathsf{A},\mathsf{B})$, of lengths $n$ and
$n+1$, the hybrid lists $\mathfrak{s}^{\natural}_{r,\bullet}$ and the norms $\|G_s^{P}\|$, $\|G_s^{A}\|$,
$\|G_s^{R}\|$ used in Definition~\ref{9.1:def-transplant-response}. The letter $\bullet$ ranges over the
production blocks $P$, $A$, $R$. The letters $R'$, $R''$ denote the two components of the block $R$ whose
norms are bounded in \eqref{9.1:eq-comparison-families-2} and \eqref{9.1:eq-comparison-families-a}.
We write
\[
e_r=\rho^{\,n-r}+\Pi_n ,
\]
where $\Pi_n$ is the power-law floor attached to the length $n$.
The letter $c$ denotes the set on which the list bounds of list-structured families are read. For a family
$F$ read component by component on $c$ we write $F^{c}$.
Let $r_\ast$ denote the matching depth of the pairs of runs compared in this chapter. We say that the
\emph{two-run comparison condition with constants $C_\sigma$, $\rho$} holds if the following is true for
every pair $(\mathsf{A},\mathsf{B})$ alive to depth $r_\ast$ and matched at depth $r_\ast$, for every $r$
with $r_\ast\le r<n$, for every block letter $\bullet$ and every hybrid list
$\mathfrak{s}^{\natural}_{r,\bullet}$, and on all domains, wrapping ones included:
there is a list-structured family $N_r^{\bullet}$ of the run $\mathsf{A}$ on the set $c$, the
\emph{comparison family at depth $r$}. Component by component on $c$ it is given by
\begin{equation}\label{9.1:eq-comparison-families-1}
N_r^{c}\approx \varphi_r^{\mathsf{B},c}
-\varphi^{c}\bigl[\mathsf{F}_r^{\mathsf{A},+,\bullet}(\mathfrak{s}^{\natural}_{r,\bullet})\bigr],
\end{equation}
where $\varphi_r^{\mathsf{B},c}$ is the family of the run $\mathsf{B}$ at depth $r$ read on $c$,
$\varphi^{c}[\cdot]$ is a production read on $c$, $\mathsf{F}_r^{\mathsf{A},+,\bullet}$ is
the production of the run $\mathsf{A}$ with the extra slot, and $\approx$ is the relation between
list-structured families of this chapter, read component by component.
Its components obey
\begin{equation}\label{9.1:eq-comparison-families-2}
\|N_r^{P}\|\le C_\sigma e_r,\qquad \|N_r^{A}\|\le C_\sigma e_r,\qquad
\|N_r^{R''}\|\le \mathfrak{r}_r\Lambda_r C_\sigma e_r ,
\end{equation}
and
\begin{equation}\label{9.1:eq-comparison-families-a}
\|N_r^{R'}\|\le \mathfrak{r}_r\Lambda_r\bigl[(C_\sigma+1)e_r+B_r(G,d)\bigr]
+\varkappa'_r\sum_{r<s<r+\ell_r}\|G_s^{R}\| .
\end{equation}
Here $\mathfrak{r}_r$, $\Lambda_r$ are the weights and $\varkappa'_r$ is the coupling constant of the
inequality $(R)$ of Definition~\ref{9.1:def-transplant-response}, $\ell_r$ is the length of the summation
window, and $B_r(G,d)$ is the quantity so denoted in the inequality \textup{(R)} of
Definition~\ref{9.1:def-transplant-response}, read at depth $r$, which
depends on the families $G=(G_s^{\bullet})_s$ and on the coupling distances $d_s$; in $B_r(G,d)$ the
letter $d$ abbreviates the tuple $(d_s)_s$.
\end{definition}

The bounds \eqref{9.1:eq-comparison-families-2} are the content of
the production theorem, its flat variant and its corollary, together with the localisation of the response.
The bound \eqref{9.1:eq-comparison-families-a} holds as an implication from the
response condition of the comparison at the inequalities carrying the total loads of the zones, through a
remark on the corresponding corollary. In that response condition the entries carrying the two-cutoff
difference of the boundary terms are
supplied zone by zone by the zone-resolved volume and filament bound, the proposition on the boundary
terms of the $\mathbf{R}$-step and the elimination of these terms.

\begin{definition}[Ceilings on the coupling for the two-run estimate]\label{9.1:def-coupling-ceilings}
Let $\rho$ be the contraction rate, and let $c_\rho$, $\alpha_r$, $\alpha_r^A$, $\mathsf{A}$, $\mathsf{n}_r$,
$\varkappa$, $\varkappa_A$, $\varkappa'_r$, $\varkappa''_r$, $\hat{\mathfrak{r}}_r$, $\Pi_n$,
$C_T$ and $\ell_r$ be the constants and depth-indexed weights of the two-run estimate, all as
fixed in Notation~\ref{9.1:not-rate-constants}. In this definition $\mathsf{A}$ denotes a constant
and not the run $\mathsf{A}$ of Definition~\ref{9.6:def-comparison-chain}. The letters
$\varkappa$, $\varkappa_A$, $\varkappa'_r$, $\varkappa''_r$ denote coupling constants of the
response bounds. They are distinct from the decay rate $\varkappa$ of the two-point witness theorem.
Here $n$ is the length of the shorter of the two runs compared, and $r$, $s$ range over depths.
The \emph{ceilings on the coupling for the two-run estimate} are the following
conditions:
\begin{equation}\label{9.1:eq-coupling-ceilings-a}
\begin{aligned}
&\text{(a)} && c_\rho\,\alpha_r\,\mathsf{A}\,\mathsf{n}_r \le \tfrac18;\\
&\text{(b)} && \varkappa \le \tfrac{1-\rho}{32};\\
&\text{(b$'$)} && 4\varkappa''_r \le \tfrac{3}{32}\,(1-\rho);\\
&\text{(c)} && c_\rho\,\alpha_r^A\,\mathsf{A}\,\mathsf{n}_r \le \tfrac14;\\
&\text{(d)} && \varkappa_A\Bigl(9+\frac{5}{1-\rho}\Bigr) \le \tfrac14;\\
&\text{(e)} && \hat{\mathfrak{r}}_r \le \tfrac1{16},\qquad
\sum_r \hat{\mathfrak{r}}_r \le 1,\qquad \Pi_n \le 1;\\
&\text{(f)} && \hat{\mathfrak{r}}_s \le C_T\,\rho^{-(n-s)/2}\,\hat{\mathfrak{r}}_n
\quad (s\le n),\qquad
\hat{\mathfrak{r}}_s \le 2\,(1+s-r)\,\hat{\mathfrak{r}}_r \quad (s\ge r);\\
&\text{(g$'$)} && 4\,\varkappa'_r\,\ell_r^{2}\,\rho^{-\ell_r} \le \tfrac14.
\end{aligned}
\end{equation}
In the sense of Definition~\ref{2.3:def-order-table}, none of the conditions
\eqref{9.1:eq-coupling-ceilings-a} reads any of the constants $C_V$, $C_\sigma$, $C_E$ of the
two-run estimate. The constant $c_\rho$ reads $(L,\mathfrak{p})$ only and carries no
polylogarithmic factor. In condition~(c), one inverse logarithm of the coupling remains available.
Indeed, with $X=g^{-2}$ and $\lambda=c_0/2$,
the factor $\alpha_r^A\,\mathsf{A}\,\mathsf{n}_r$ of the left member of~(c) satisfies
$\alpha_r^A\,\mathsf{A}\,\mathsf{n}_r\le C\,\mathsf{A}/(\lambda\log X)$
with a constant $C$ independent of $X$.
\end{definition}

\begin{lemma}[Ambient regularity and the pair law for glued minimisers]
\label{9.1:lem-glued-minimisers}
We work in the setting of the seam proposition. Its objects are:
\begin{itemize}
\item the two runs $\mathsf{A}$ and $\mathsf{B}$;
\item a vertex $\sigma'$ of the localisation;
\item the regions $W$, $\Lambda$, $Y_1$ and Ba{\l}aban's determining sets $\mathbf{B}''_k$,
$\mathbf{B}(Y_1)$;
\item the box set $\boldsymbol{\Omega}(W)$ and the region $\Omega_0(W)$;
\item Ba{\l}aban's configurations $U_{1,2}$ and $U^0_k$, which enter the response
\cite[(1.54)]{B-CMP122b};
\item the glued configuration $\tilde U_{1,2}$;
\item the minimiser $\mathfrak{U}(\sigma')=U(\mathcal{P}_W)$ of the problem $\mathcal{P}_W$;
\item the minimiser $\Xi_1(\sigma')$ on $\mathbf{B}^{\sim}(Y_1)$ with $\sigma'$-localised
layer reads;
\item the minimisers $\Xi_2$, $\Xi_3$ of the earlier members of the localisation chain.
\end{itemize}
For a pair of corresponding objects of the two runs, $\mathsf{e}$ denotes their two-run
error and $\mathsf{u}$ their single-run size. We assume the following.
\begin{itemize}
\item The hypotheses of the seam proposition. These are:
\begin{itemize}
\item the natural-interpolation statement and the current statement;
\item the mismatch lemma;
\item the localisation statements of the small-field machinery, local, on $\Lambda$ and on
the rim;
\item the axial-Landau frame statement and the case lemma for carried bounds;
\item the second clause of the tilde construction;
\item \cite[Theorem~1, (181), Proposition~9]{B-CMP102b};
\item the first $\rho$-condition at the lower bound on $L$, the second $\rho$-condition and
the zero-range condition;
\item the two floors of the order of choice \eqref{9.1:eq-constant-order-a} used by the seam
proposition, printed there as $(\mathrm{F}^{\flat})$ and $(\mathrm{F}'')$;
\item the inequalities
\begin{equation*}
M\ge 4LR_1M_1,\qquad M_1\ge 8d ,
\end{equation*}
where $R_1$ is a radius of the seam proposition (not the radius $R_1$ of the parameter
window in $\hat w$).
\end{itemize}
\item The layer lemma, read on the H\"older ladder: at H\"older exponents that decrease from
$\beta_H$ by $\varepsilon_H/32$ per link of the localisation chain, all of them at least
$\beta^\flat=\beta_H-\varepsilon_H/2>\theta$, as fixed in the order of choice
(Notation~\ref{9.1:not-constant-order}).
\item The ambient-index comparison lemma: the comparison lemma of the seam proposition under
its hypotheses, with its local-index bounds on the net mismatch $\hat{\mathfrak{d}}$ and on
the profile $\mathfrak{s}$ replaced by
\begin{align}
|\hat{\mathfrak{d}}(y)|&\le c_{\mathfrak{d}}\,\mathsf{u}\,L^{-\mathsf{y}(y)}
\vartheta^{m(y)},\label{9.1:eq-glued-minimisers-1}\\
\mathfrak{s}&\le c_S\,\mathsf{u}\,L^{-2\mathsf{y}},\label{9.1:eq-glued-minimisers-2}
\end{align}
for constants $c_{\mathfrak{d}}$, $c_S$.
For a pair problem of the seam proposition it gives the two-run error at the ambient index,
\begin{equation}\label{9.1:eq-glued-minimisers-3}
\mathsf{e}\le C_F'(c_{\mathfrak{d}}+c_S)\,\mathsf{u}\,\vartheta^m ,
\end{equation}
$C_F'$ being the constant of that lemma,
modulo one lattice gauge transformation $\Gamma$, for which, by the gauge clause of that
lemma and under that clause's hypotheses,
\begin{equation}\label{9.1:eq-glued-minimisers-6}
|\Gamma v-1|\le C_F'(c_{\mathfrak{d}}+c_S)\,\mathsf{u}\,L^{-\mathsf{y}}\vartheta^m ,
\end{equation}
where $v$ is the site frame assigned by the frame lemma of the seam proposition.
\item Clauses (a)--(c) of the regularity theorem for $U_{1,2}$ at the
scale of the label.
\item The two cut statements.
\item The ceiling on $\gamma$:
\begin{equation}\label{9.1:eq-glued-minimisers-4}
C_{\mathrm{lay}}^{11}\cdot c_b\cdot\sup\mathsf{u}\le\min\{\mathfrak{r}_0,a_{R'}\}.
\end{equation}
Here $c_b$ is the constant of the single-run bound $|b|\le c_b\,\mathsf{u}$ on the data used
in the proof below; $\mathfrak{r}_0$ and $a_{R'}$ are radii, $\mathfrak{r}_0$ that of the
bases to which the layer lemma applies and $a_{R'}$ that of the Landau chart.
\end{itemize}

Let $m$ be the level of $\mathbf{B}''_k$ on the region. Let $\mathsf{y}_W$ and $\mathsf{y}$
be the youths with respect to $\boldsymbol{\Omega}(W)\vee\mathbf{B}''_k$ and to
$\mathbf{B}(Y_1)\cup\mathbf{B}''_k$ respectively. Then at every vertex $\sigma'$ the
following hold.
\begin{itemize}
\item[(a)] $\mathfrak{U}(\sigma')$ is a diagonal entry; in particular it carries its own
bounded current. It is ambient-regular: it satisfies, at the ambient scale $\ell_m$, the field
clauses of the H\"older class $\mathsf{C}_{\beta^\flat}(\mathfrak{r}')$ of the layer lemma,
of exponent $\beta^\flat$ and radius $\mathfrak{r}'$. It has tame profile
\begin{equation*}
\mathfrak{s}_{\mathfrak{U}(\sigma')}\le\mathfrak{r}'L^{-2\mathsf{y}_W}.
\end{equation*}
The pair $(\mathfrak{U}(\sigma')^{\mathsf{B}},\mathfrak{U}(\sigma')^{\mathsf{A}})$ obeys the
pair law at ambient index
\begin{equation*}
\mathsf{e}\le C^\natural_1\,\mathsf{u}\,\vartheta^m
\end{equation*}
on $\boldsymbol{\Omega}(W)\vee\mathbf{B}''_k$ over $\Omega_0(W)$. The slot, the seam, the
cut collars, the core, the rim and the collar are included. The law holds modulo one
lattice gauge transformation $\Gamma_W$ with
\begin{equation*}
|\Gamma_W v-1|\le C^\natural_1\,\mathsf{u}\,L^{-\mathsf{y}_W}\vartheta^m ,
\end{equation*}
where $v$ is the site frame assigned by the frame lemma of the seam proposition.
\item[(b)] $\Xi_1(\sigma')$ is a diagonal entry and ambient-regular, with
$\mathfrak{s}\le\mathfrak{r}''L^{-2\mathsf{y}}$. It obeys the pair law
\begin{equation*}
\mathsf{e}\le C^\natural_{\mathrm{LR}}\,\mathsf{u}\,\vartheta^m
\end{equation*}
on $\mathbf{B}(Y_1)\cup\mathbf{B}''_k$ over all of $Y_1$. The law holds modulo a gauge
transformation $\Gamma(\sigma')$ with
\begin{equation*}
\Gamma(\sigma')v_{\mathrm{AL}}=1+O(C^\natural_{\mathrm{LR}}\mathsf{u}\vartheta^m)
\quad\text{on the slot,}
\end{equation*}
where $v_{\mathrm{AL}}$ is the axial-Landau frame on the slot.
The constants $C^\natural_1$, $C^\natural_{\mathrm{LR}}$, $C_{\mathrm{lay}}$ depend on
$(d,L,\theta,M_1,M)$ only. The quantities $\mathfrak{r}'=C_{\mathrm{lay}}\mathfrak{r}_1$ and
$\mathfrak{r}''=C_{\mathrm{lay}}\mathfrak{r}'$, where $\mathfrak{r}_1$ is the base radius
\eqref{9.1:eq-glued-minimisers-5} fixed in the proof of (a), are profiles bounded by
$C\cdot\sup\mathsf{u}$; they read $\gamma$ from above only, in the sense of
Definition~\ref{2.3:def-order-table}, so that the conditions they give on $\gamma$ are
ceilings. None of the constants depends on $k$, $K$, $h$, $N$, $N_0$, $\check R_k$ or
$|\Lambda|$, and none depends on a lower bound on $g$.
\item[(c)] Consequently the layer corollary holds in the following form.
\begin{itemize}
\item $\Xi_3$ and $\Xi_2$, at the first six members of the localisation chain, are
controlled by clauses (a) and (b) of the stage lemma, the frame lemma and the ambient-index
comparison lemma.
\item $\mathfrak{U}(\sigma')$ and $\Xi_1(\sigma')$, at the seventh to ninth members, are
controlled by the layer lemma, the frame lemma and the ambient-index comparison lemma, in
the format of a fresh pair, that is, of two minimisers of exactly paired problems compared
directly, with no common parent configuration.
\item The tenth member is controlled through the member $0'$ of the chain.
\item The seventh member uses a single-run statement only, at Ba{\l}aban's expansion centre.
No pair law is required of $U_{1,2}$ or $U^0_k$, and none is used there.
\item The eighth member is an $\ell_2$-link of the localisation chain, taken between a fresh
pair.
\item The parameter $\lambda_e$ of the layer corollary equals $\infty$.
\item The quantity $K^{\mathrm{rip}}$ of the layer corollary delivers the factor
$\vartheta^{\iota(c)}$ at every $p\in\mathfrak{F}^{V,\mathbb{R}}$, rims included.
\end{itemize}
\end{itemize}
\end{lemma}

\begin{proof}
\emph{The H\"older ladder.} The layer lemma takes a base in a H\"older class
$\mathsf{C}_\beta$ and returns an output in a smaller H\"older class. The base used in (a)
below already contains $\Xi_2$, which is an output of the layer lemma, and the base used
in (b) is the output of (a). The lemma is therefore read at a H\"older exponent that
decreases along the links of the localisation chain, by $\varepsilon_H/32$ per link. The
constant $B_0'(\varepsilon)$ of the stability statement only enlarges $C_{\mathrm{lay}}$.
The chain has at most eleven links, and $11\varepsilon_H/32<\varepsilon_H/2$. Hence every
object of the chain lies in $\mathsf{C}_{\beta^\flat}$,
$\beta^\flat=\beta_H-\varepsilon_H/2>\theta$.

\emph{Statement (a), single-run part.} The layer lemma is stated for an arbitrary
ambient-regular base $\mathfrak{U}$. It has two hypotheses:
\begin{itemize}
\item the \emph{field hypothesis}: the field clauses of $\mathsf{C}_\beta(\mathfrak{r})$ hold
at the ambient scale;
\item the \emph{current hypothesis}: $\ell^3|J_{\mathfrak{U}}|\le\mathfrak{r}$, where
$J_{\mathfrak{U}}$ is the current of the base and $\ell$ the lattice spacing at the ambient
scale.
\end{itemize}
Criticality of the base is not among them. The proof of the layer lemma treats the current
of the base as an order-zero source $\mathfrak{f}_{\mathfrak{U}}[\phi]=\langle\phi,
J_{\mathfrak{U}}\rangle$, for test fields $\phi$, of size $\mathfrak{r}$ times the third
power of the youth factor, that is, of $L$ raised to minus the youth of the set on which the
lemma is applied. We take as base
$\mathfrak{U}:=\tilde U_{1,2}$ on $Y:=\Omega_0(W)$, and we put
$\mathfrak{S}_2:=\boldsymbol{\Omega}(W)\vee\mathbf{B}''_k$.

We check both hypotheses for $\tilde U_{1,2}$, in three places. Let $\bar\Lambda^+$ be the
region of the seam proposition inside which $\tilde U_{1,2}$ is glued from the two sides. Off
$\bar\Lambda^+$, $\tilde U_{1,2}$ equals $\Xi_2(\sigma'_3,\sigma'_4)$. Indeed, by the first geometric
property of the seam proposition, on $W\setminus\Lambda$ the configuration
$\tilde U_{1,2}$ is its side-one member. There $\Xi_2$ is ambient-regular with bounded
current. This is clause (iii) of the layer lemma applied at the induction step of the layer
corollary for $\Xi_2$, at the fifth and sixth members of the chain.

Inside $\bar\Lambda^+$ and off the cut collar $\mathcal{L}$, the seam proposition defines
$\tilde U_{1,2}$ as $\Xi_2(\sigma'_3,\sigma'_4)$ on side one and as $U_2$ on side two. Here
$U_2$ is Ba{\l}aban's minimiser \cite[(9)--(10)]{B-CMP102b}, and $\Xi_2$ is the side-one
member. (The glued configuration coincides with Ba{\l}aban's $U_1$ only at the undecoupled
vertex.) There $\tilde U_{1,2}$ is ambient-regular, as off $\bar\Lambda^+$.

On $\mathcal{L}$, clause (a) of the regularity theorem for $U_{1,2}$ at the scale of the
label gives $U_{1,2}\in\mathsf{C}_\beta(a_{\mathrm{reg}})$, its
H\"older member and its own current included; $a_{\mathrm{reg}}$ is the radius of that
theorem. At the seam this is combined with the two cut statements.

Hence the field and current hypotheses hold with
\begin{equation}\label{9.1:eq-glued-minimisers-5}
\mathfrak{r}:=\max\{\mathfrak{r}'_{\Xi_2},\,B_3\varepsilon_1,\,a_{\mathrm{reg}}\}
=:\mathfrak{r}_1 ,
\end{equation}
where $\mathfrak{r}'_{\Xi_2}$ is the ambient radius of $\Xi_2$ from the fifth and sixth
members, and $B_3\varepsilon_1$ is the radius of the $B_3\varepsilon_1$-class of the seam
proposition.

\emph{Data.} At the original sites of youth zero inside $W$ we put
$b:=D^0\cdot(\mathrm{rd}\,\tilde U_{1,2})^{-1}$. In $\Lambda$ these data consist of the
shift data, restored by the argument of \cite[(1.54)]{B-CMP122b}, and the remaining data
$D^0$ of $U_{1,2}$. The single-run bound of the layer corollary controls $b$ as follows.
We have $b=Q_m(\tilde U_{1,2},\eta E)$, $Q_m(\tilde U_{1,2},\cdot)$ denoting the level-$m$
averaging map of the data at the base $\tilde U_{1,2}$, where
\begin{equation*}
e^{i\eta E}:=U^0_k\tilde U_{1,2}^{-1}
\end{equation*}
is the response of \cite[(1.54)]{B-CMP122b}, $\eta$ being its parameter and $E$ its field.
Therefore
\begin{equation*}
|b|\le 18d\ell\sup|E|\le c_b\,\mathsf{u}
\end{equation*}
by the mismatch lemma and the first row of the table of the cut construction. After the
ceiling \eqref{9.1:eq-glued-minimisers-4}, $c_b\,\mathsf{u}\le\mathfrak{r}_1$. At the rim and
collar sites we put $b:=0$. The datum there is $M^r(\tilde U_{1,2})$, the base's own averages.

Put $\Xi_b:=e^{i\eta\mathsf{X}_b}\tilde U_{1,2}$ with
$\mathsf{X}_b:=\mathfrak{X}^{\mathrm{res}}(b;\tilde U_{1,2})$ on $\mathfrak{S}_2$, where
$\mathfrak{X}^{\mathrm{res}}$ is the map of the layer lemma. Then
$U(\mathcal{P}_W)$ and $\Xi_b$ have the same data. The map $\mathfrak{X}^{\mathrm{res}}$
carries the full force $\mathfrak{f}$. Hence the source $G''_kJ_{1,2}$ of
\cite[(1.54)]{B-CMP122b} is contained in the layer lemma's source allowance
$\mathfrak{r}_1L^{-3\mathsf{y}_W}$ and is not an extra term.
By the criticality clause of the healing lemma, $\Xi_b$ is critical and lies in the space
\cite[(6)]{B-CMP102b}. Hence it lies in the minimal orbit, and
$\mathfrak{U}(\sigma')\sim\Xi_b$ by \cite[Theorem~1]{B-CMP102b}.

Clauses (i)--(iv) of the layer lemma, read on the ladder, now give the single-run
conclusions. First, $\mathfrak{U}(\sigma')$ is critical on $\mathfrak{S}_2$, hence diagonal
by clause (a) of the diagonal-entry theorem. That theorem does not require criticality in
the pair sense, but here the configuration is even critical. Second, $\mathfrak{U}(\sigma')$
is ambient-regular with $\mathfrak{r}':=C_{\mathrm{lay}}\mathfrak{r}_1$. Third,
$\mathfrak{s}_{\mathfrak{U}(\sigma')}\le\mathfrak{r}'L^{-2\mathsf{y}_W}$. This is the
profile bound \eqref{9.1:eq-glued-minimisers-2} for $\mathcal{P}_W$.

The current hypothesis of the layer lemma reads a second-order member of the base. That
member is derived in the single run, at the vertex: the current statement applied to the
ruled fifth member of the base, together with clause (a) of the regularity theorem for
$U_{1,2}$ at the scale of the label and \cite[(10)]{B-CMP102b}, gives it. A member derived in
this way is admissible in the current hypothesis under the condition on derived members of
the base that is stated with the layer lemma. Alternatively, the same conclusion holds with
the loss $\varepsilon/8$ stated with the layer lemma.

\emph{Statement (a), two-run part.} We apply the frame lemma of the seam proposition to
$\mathcal{P}_W$. The sites fall into the four classes that the seam proposition uses for its
first clause.
\begin{itemize}
\item[($\alpha$)] Original outer sites are common to both runs; $v:=1$, by the first clause
of the frame lemma.
\item[($\beta$)] At slot sites, clause (ii) of the axial-Landau frame statement gives
$v:=v_{\mathrm{AL}}$ and $\hat{\mathfrak{d}}:=\mathfrak{d}_{\mathfrak{s}}$. The
local-index form of \eqref{9.1:eq-glued-minimisers-1} holds with index $\iota=k$. Since
$k\ge m$, \eqref{9.1:eq-glued-minimisers-1} holds a fortiori.
\item[($\gamma$)] At core sites the data are common integration variables, by the seam
proposition; $v:=1$ and $\hat{\mathfrak{d}}=0$.
\item[($\delta$)] Rim and collar sites have youth $\mathsf{y}_W>0$. By the first geometric
property of the seam proposition they all lie in $W\setminus\Lambda^{\sim}$, which is
contained in side one; here $\Lambda^{\sim}$ is a set fixed in the seam proposition, distinct
from the region $\bar\Lambda^+$ above. The datum there is $M^r(\Xi_2)$.
\end{itemize}
For the class ($\delta$), the seam proposition used the three cases of the case lemma on the
bounds carried with $\Xi_2$ at local index. Now $\Xi_2^{\mathsf{A}}$ is ambient-regular, by
clause (iii) of the layer lemma at the induction step for $\Xi_2$. Moreover the pair
$(\Xi_2^{\mathsf{B}},\Xi_2^{\mathsf{A}})$ obeys the pair law at ambient index, by the
layer corollary at the fifth and sixth members. So the mismatch lemma applies directly at
$\Xi_2^{\mathsf{A}}$. Its hypothesis on the field at scale $\ell_r$ on $P(y)$ is given by
the first clause of the uniformity statement. It yields
\begin{equation*}
\begin{split}
|\hat{\mathfrak{d}}(y)|
&\le 2\ell_r\sup_{P(y)}|A_{\Xi_2}|
=2L^{-\mathsf{y}_W(y)}\,\ell_m\sup|A_{\Xi_2}|\\
&\le 2C^\natural_{\mathrm{LR},2}\,\mathsf{u}(y)\,L^{-\mathsf{y}_W(y)}\vartheta^m ,
\end{split}
\end{equation*}
where $\ell_r=L^{-\mathsf{y}_W(y)}\ell_m$ is the lattice spacing at the local index $r$ of
$y$, $P(y)$ is the block of plaquettes at $y$ read by the mismatch lemma, and $A_{\Xi_2}$ is
the field at $\Xi_2$ that the mismatch lemma reads. The last inequality is the pair law for
$(\Xi_2^{\mathsf{B}},\Xi_2^{\mathsf{A}})$ given by the layer corollary at the fifth and sixth
members, $C^\natural_{\mathrm{LR},2}$ being its constant.
Put $v:=\tilde\lambda_{\Xi_2}^{-1}$, $\tilde\lambda_{\Xi_2}$ being the gauge transformation
attached to $\Xi_2$, times the factors $\mathsf{w}_{\mathrm{site}}$ of the
dilation circles at the sites. It is single-valued up to discrepancies between these factors
of size at most $32d\ell a$, by the second clause of the frame lemma, $a$ being the constant
of that clause. No further case distinction is needed.

Thus \eqref{9.1:eq-glued-minimisers-1} holds for $\mathcal{P}_W$ with the choice
$c_{\mathfrak{d}}:=2C^\natural_{\mathrm{LR},2}+c_{\mathrm{AL}}$, where $c_{\mathrm{AL}}$ is
the constant of the slot mismatch bound of the axial-Landau frame statement. By the first
geometric property of the seam proposition, no rim point is near the slot or the core.
Hence there is no young junction with $v_{\mathrm{AL}}$. The bounds
\eqref{9.1:eq-glued-minimisers-3} and \eqref{9.1:eq-glued-minimisers-6} of the ambient-index
comparison lemma now give
\begin{equation*}
\mathsf{e}\le C_F'(c_{\mathfrak{d}}+c_S)\,\mathsf{u}\,\vartheta^m=:C^\natural_1\,\mathsf{u}\,
\vartheta^m,
\qquad
|\Gamma_Wv-1|\le C^\natural_1\,\mathsf{u}\,L^{-\mathsf{y}_W}\vartheta^m .
\end{equation*}
This proves (a).

\emph{Statement (b), single-run part.} We apply the layer lemma with base
$\mathfrak{U}:=\mathfrak{U}(\sigma')$, extended by $\tilde U_{1,2}$ off $\Omega_0(W)$. On
$\Omega_0(W)$ this base is ambient-regular with bounded current by (a). Off $\Omega_0(W)$ it
has these properties by the verification of hypotheses in the single-run part of (a). The
set is $\mathbf{B}(Y_1)\cup\mathbf{B}''_k$.

In the proof of (b) the letter $C$ denotes a layer site, not a constant. At the layer
sites $C$ we put $b:=0$. The datum there is $M^p(\mathfrak{U}(\sigma'))(C)$,
the base's own averages; here and below $p$ is a layer index, not a plaquette. By the second
geometric property of the seam proposition, all
layer sites lie outside the region $Z$ fixed in that property. At the original sites of
$Y_1^0$, $b$ is defined as in the
data paragraph of (a), site by site according to the classes ($\alpha$)--($\gamma$), and
$|b|\le c_b\,\mathsf{u}\le\mathfrak{r}'$.

Let $\Xi_b$ now be built from the base $\mathfrak{U}(\sigma')$ and these data exactly as in
(a). As in (a), $\Xi_b$ is critical and $\Xi_1(\sigma')\sim\Xi_b$. Clauses (iii) and (iv) of
the layer lemma then give that $\Xi_1(\sigma')$ is critical, hence diagonal. It is
ambient-regular with
$\mathfrak{r}'':=C_{\mathrm{lay}}\mathfrak{r}'$, and
$\mathfrak{s}\le\mathfrak{r}''L^{-2\mathsf{y}}$.
This is \eqref{9.1:eq-glued-minimisers-2} for $\Xi_1(\sigma')$.

\emph{Statement (b), two-run part.} We apply the frame lemma with the classes
($\alpha$)--($\gamma$) on $Y_1^0$. No auxiliary parent configuration occurs here, so the
error compared is that of the pair itself. We add one further class.
\begin{itemize}
\item[($\delta'$)] At layer sites the datum is $M^p(\mathfrak{U}(\sigma'))(C)$.
\end{itemize}
The layer lemma distinguishes two kinds of layer site, according to the position of $P(C)$
relative to the $(p+1)$-st layer region $\Omega^W_{p+1}$ of $W$.
\begin{itemize}
\item In the outer cases, $P(C)\not\subset\Omega^W_{p+1}$. Then
$M^p(\mathfrak{U})(C)=M^p(\Xi_2)(C)$ by the identification lemma for layer averages.
\item In the inner case, $P(C)\subset\Omega^W_{p+1}$.
\end{itemize}
Both are now direct applications of the mismatch lemma, at an ambient-regular
$\mathsf{A}$-field whose pair law at ambient index is known: for $\Xi_2$ by the fifth and
sixth members of the chain, and for $\mathfrak{U}(\sigma')$ by (a). Hence
\begin{equation*}
|\hat{\mathfrak{d}}(C)|\le 2L^{-\mathsf{y}(C)}\mathsf{e}(\cdot)
\le 2\max\{C^\natural_{\mathrm{LR},2},C^\natural_1\}\,\mathsf{u}\,L^{-\mathsf{y}}\vartheta^m .
\end{equation*}
A young junction is a layer point shared by an inner-case and an outer-case site. It is a
rim point of $W$ of youth $\mathsf{y}_W\ge\mathsf{y}-1$, so
$L^{-\mathsf{y}_W}\le L\cdot L^{-\mathsf{y}}$. At such a point the frame jump is bounded by
(a):
\begin{equation*}
|\Gamma_Wv-1|\le C^\natural_1\,\mathsf{u}\,L^{-\mathsf{y}_W}\vartheta^m
\le LC^\natural_1\,\mathsf{u}\,L^{-\mathsf{y}}\vartheta^m .
\end{equation*}
This jump is weighted, so \eqref{9.1:eq-glued-minimisers-1} holds with the choice
\begin{equation*}
c_{\mathfrak{d}}:=2\max\{C^\natural_{\mathrm{LR},2},C^\natural_1\}+LC^\natural_1
+c_{\mathrm{AL}}.
\end{equation*}
The bound \eqref{9.1:eq-glued-minimisers-3} of the ambient-index comparison lemma gives (b)
with $C^\natural_{\mathrm{LR}}:=C_F'(c_{\mathfrak{d}}+c_S)$. On the slot,
$\Gamma(\sigma')v_{\mathrm{AL}}=1+O(C^\natural_{\mathrm{LR}}\mathsf{u}\vartheta^m)$ follows
from the second clause of the frame lemma, as in the proof of the second clause of the seam
proposition.

\emph{Statement (c).} The constants compound along the localisation chain at most eleven
times, eleven being the maximal number of links of the chain, and $11\le m_L$, where $m_L$ is
the constant of that name in the setting of the layer lemma. The members
are controlled as follows.
\begin{itemize}
\item The first six members are controlled by clauses (a) and (b) of the stage lemma, the
frame lemma and the ambient-index comparison lemma, as in the proof of the layer corollary.
\item The seventh to ninth members are controlled by (a) and (b).
\item The tenth member is controlled through the member $0'$ of the chain, from the critical
parent $\Xi_1(\sigma')$, by the proposition of the seam construction on the member $0'$.
\end{itemize}
The ceiling \eqref{9.1:eq-glued-minimisers-4} keeps the layer lemma applicable at depth
eleven: $\mathsf{u}$ tends to zero with $\gamma$, while $C_{\mathrm{lay}}$ grows at most
like $L^6$.

It remains to check the two further conditions in the definition of a diagonal entry.
The entry condition at the ambient scale follows from clause (iii) of the layer lemma
and clause (a) of the diagonal-entry theorem. The level condition for the entry's set holds
because $m=\iota_N\ge\iota_\nu$, where $\iota_N$ and $\iota_\nu$ are the local indices
compared in that condition.
\end{proof}

\begin{remark}
The proof uses the listed hypotheses together with the auxiliary statements of the setting
of the seam proposition and of the layer lemma that it names; no further assumption enters.
The
first $\rho$-condition enters only where the seam proposition and the cut statements use it:
at the lower bound on $L$, and through one statement of the cut construction and one further
input of the seam proposition. We note two points about the proof.

The force of the layer lemma stays of order zero across the seam, although
$\tilde U_{1,2}$ is glued and not critical. The force
$\mathfrak{f}_{\mathfrak{U}}[\phi]=\langle\phi,J_{\mathfrak{U}}\rangle$ reads the current of
the base pointwise. On $\mathcal{L}$, clause (a) of the regularity theorem for $U_{1,2}$ at
the scale of the label bounds
$\ell_h^3|J(U_{1,2})|\le a_{\mathrm{reg}}$ in the units of the label. The first cut
statement, that $J_{1,2}$ restricted to ${\Omega''}^{\sim}_{h+1}$ is equal to $J_1$, makes
it the current of a single run there. The seam plaquette variables are exactly $1$.

Statement (b) reads no cut stratum: by the seam proposition, no cut site is read by any of
the compared objects.
\end{remark}

\begin{definition}[Hulls, filing and the strip width, with rates and floors]\label{9.1:not-hull-conventions}
Throughout this chapter $d$ denotes the dimension ($d=4$), $\kappa$, $\kappa_1$, $\delta_1$, $M$ are Ba{\l}aban's auxiliary parameters, and $\beta_{pr}=\beta_I$, $\beta_s$, $\beta^{\star}=\beta_s$ and $\delta_\star$ are the numbers fixed in Notation~\ref{9.1:not-constant-order} and Notation~\ref{9.1:not-radius-exponent-floors}; we write $\beta:=\beta_{pr}=\beta_I$. The constants, sets, rates, floors and ceilings used in this chapter are fixed in the display
\begin{equation}\label{9.1:eq-hull-conventions-a}
\begin{aligned}
&c_H=c_g:=1+\sqrt{d}\,3^dC_d,\qquad c^{\sharp}:=c_H+1,\qquad \kappa_R:=c^{\sharp}\kappa,
 \qquad a_\flat:=\frac{(1-\beta_{pr})\kappa}{c^{\sharp}},\\
&K(X):=(X^{+})_0,\qquad (\mathrm{F1})\quad Y_1^{+}(X):=K(X)\cup C(K(X)),\\
&Z^{\sharp}:=Z^{(m_\sharp)},\qquad C^{\sharp}:=C^{(m_\sharp)},\qquad
 m_\sharp:=m_{hi}(k):=\bigl\lfloor\tfrac12(\check{R}_k-1)\bigr\rfloor,\\
&A_H:=e^{\frac74 c_H(4+3^d)\kappa},\qquad \mathsf{c}_d:=2\cdot 3^d,\qquad
 \mathsf{k}_d:=1+\mathsf{c}_d\bigl(2d\log 3+(1+3^d)\log 2\bigr),\\
&C_{fib}:=2e^{2\mathsf{k}_d},\qquad \beta_v:=\frac{\beta}{2\mathsf{c}_d},\\
&\beta_\sharp:=\tfrac{1}{20},\qquad b:=(1+2\beta_s)\kappa=\tfrac32\kappa,\qquad
 a:=b+\beta_\sharp\kappa,\\
&(F_\star)\quad L\ge 8c^{\sharp}c_g^2,\qquad (F_L)\quad L\ge 3c_H,\\
&L\ge \tfrac{35}{8}c_Hc^{\sharp},\qquad \kappa\ge 10\log L,\\
&(F_\kappa)\quad \kappa\ge \frac{c^{\sharp}C(d)}{\delta_\star},\qquad
 \frac{a_\flat}{4}\ge \frac{c^{\sharp}C(d)}{\delta_\star},\qquad
 \frac{\kappa}{20}\ge c_d+2\sqrt{d}\,(2+4\beta_{pr}),\\
&(F_\sharp)\quad \beta_\sharp\kappa\ge 2(\mathsf{k}_d+1),\qquad
 \kappa_1\ge 4L\bar{\kappa}+1,\qquad \delta_1M\ge\kappa_1,\\
&(\gamma_\sharp)\quad \sup_{T\ge T_\gamma}C_{fib}\bigl(1+\mathsf{K}^0+N(g)C_0^{\sharp}\bigr)
 e^{-\frac14\beta_\sharp\kappa(\check{R}(g)-4)}\le\alpha_*,\\
&\varepsilon_{br}\le 1,\qquad e^{-\frac12p_0(g_k)}\,100^d\,\check{R}_k^{-d}\le\tfrac14\beta_v\kappa,\\
&(\gamma2)\quad \text{the ceiling on the coupling in the correlation theorem, imposed}\\
&\qquad\quad \text{at one estimate of \cite{B-CMP122b}},\\
&(\gamma3)\quad \text{the ceiling }(\gamma3)\text{ of }\eqref{9.1:eq-constant-order-a}\\
&\qquad\quad \text{at the exponent }-\tfrac32p_0(g_k),\text{ with }\beta_0^{pr}\le\tfrac17,\\
&(\gamma4)\quad 2^dc_{vol}+(2^d-1)\mathsf{K}_{cell}(k)\le C_{\log}\log g_k^{-2}.
\end{aligned}
\end{equation}
The objects in it have the following meaning.

\emph{Hull constants.} The constant $c_H$ (also written $c_g$) is the constant of the hull-length bound: for sets $Y\subset Y'\subset Y^{+}$ one has $d(Y')\le c_H(d(Y)+1)$, where, by that bound, the constant $C_d\ge0$ depends on $d$ only. The constants $c^{\sharp}$ and $\kappa_R$ are as displayed. The number $a_\flat$ is the rate at which the chain terms $\mathbb{C}^{(i)}$ filed below their top level are controlled on the route through the collar.

\emph{Hulls and filing.} For a set $X$ at level $j$, $X^{+}$ is $X$ together with one layer of cubes of Ba{\l}aban's partition $\pi_j$, and $X_0$ is the union of the cubes of side $M$ (in current lattice units) meeting $X$. Two sets \emph{meet} if they share a cube; with this convention no connector cubes occur. The hull of $X$ is $K(X):=(X^{+})_0$. The filing rule is $(\mathrm{F1})$ of \eqref{9.1:eq-hull-conventions-a}, where $C(K(X))$ is the union of the components of the large-field region $Z$ meeting $K(X)$.
The components are those meeting $K(X)$, not merely those meeting $X$: otherwise the boundary of $Y_1$ could cut a component of $Z$, and the clause of \cite[(1.61)]{B-CMP122b} requiring whole components would be lost. In \cite[(1.66)]{B-CMP122b} the set $Y_0$ is replaced by $K(X)$; all derived objects keep Ba{\l}aban's cores and his set $Y_1$. The hull parameters used in the sums of the correlation theorem over the hulls $K(X)$, a hull constant and a number of hull layers, are taken to be $c_H$ and $1$; only the filing rule $(\mathrm{F1})$ is used in this chapter, and the parameters $(c_H^3,3)$ belong to the proof of the correlation theorem.

\emph{The strip.} The large-field region at level $k$ is $Z=Z_k=\Lambda_k^{c}$. The components of $\Lambda_k^{c}$ are the classes of the equivalence relation generated by ``intersect or touch'', contact at a corner included, as in \cite{B-CMP122b}. There is one strip width, $m_\sharp=m_{hi}(k)$. For a large-field region $Z$, $Z^{(m)}$ is $Z$ together with $m$ layers of cubes of side $M$ of the current lattice, and for a component $C$ of $\Lambda_k^{c}$, $C^{(m)}$ is $C$ together with $m$ layers of cubes of side $M$ of the current lattice; $Z^{\sharp}$, $C^{\sharp}$ are as displayed. Ba{\l}aban's enlarged region $Z^{\sim}$ is read through the dictionary $\mathfrak{D}^{\sharp}$ to the strip, except at one place of \cite{B-CMP122b} and at \cite[(2.30), (2.41)]{B-CMP119}, where the tilde keeps its meaning as the collar $Z^{\boxplus}$ of thickness $M\check{R}_k$, that is, $\check{R}_k$ layers of cubes of side $M$. The strip lemma and the strip theorem of this chapter hold for every integer width $1\le m_\sharp\le m_{hi}(k)$; their smallness conclusions are asserted at $m_\sharp=m_{hi}(k)$ only.

\emph{Margin and output radius exponent.} The margin of the radius exponents is $\delta:=\delta_\star$. The output radius exponent satisfies
\begin{equation}\label{9.1:eq-hull-conventions-1}
\kappa_{out}\ge (1+4\beta^{\star})\kappa_R=2c^{\sharp}\kappa,
\end{equation}
the inequality being part of \eqref{9.1:eq-radius-exponent-floors-a}.

\emph{Combinatorial constants.} $A_H$ is the constant of the bracket estimates; it is compensated by $\mathsf{w}_L^{-1}$, and the products $A_H\mathsf{w}_L^{-1}$ are absorbed into the constant $C_\sharp$ of the strip estimates, a number not to be confused with the set $C^{\sharp}$. The constants $\mathsf{c}_d$, $\mathsf{k}_d$, $C_{fib}$ and $\beta_v$ are as displayed.

\emph{Rates.} The output rate is $b$ and the input rate is $a$; the strip rate $\beta_\sharp=\tfrac1{20}$ is the only strip rate used. In the bounds of the correlation theorem the localised terms, indexed by the sets $X$ whose hulls and filing are fixed above, enter at rate $\tfrac74\kappa$ and the chain terms at rate $\tfrac85\kappa$. The spare $\tfrac{\kappa}{10}$ of the chain terms is split: $\tfrac{\kappa}{20}$ sums their domains and $\tfrac{\kappa}{20}$ is the rate of the strip. The two halves are used once each: the strip rate $\beta_\sharp\kappa$ and the rate at which the domains are summed add up to the spare $\tfrac{\kappa}{10}$. Where the uniqueness argument's own output and input rates, $(1+2\beta_I)\kappa$ and $(1+3\beta_I)\kappa$, are used in place of $b$ and $a$, their difference is at least $\beta_\sharp\kappa$, and the strip statements apply a fortiori.

\emph{Floors.} The floors in \eqref{9.1:eq-hull-conventions-a} are one-sided. Their right-hand sides depend on $(d,\beta)$ only unless further symbols appear in them. The floor $(F_\star)$ implies $(F_L)$ and $L\ge\tfrac{35}{8}c_Hc^{\sharp}$; the two weaker forms are displayed as well. Here $C(d)$ and $c_d$ are constants depending on $d$ only: $C(d)$ enters the floor $(F_\kappa)$, and $c_d$, a letter separate from $\mathsf{c}_d$, enters the floor on $\kappa/20$; $\bar{\kappa}$ is a number fixed before $\kappa_1$ in the order of choice of Notation~\ref{9.1:not-constant-order}. $(F_\kappa)$ is imposed also with $\kappa$ replaced by the rate $a_\flat/4$, as displayed, and the floor on $\kappa/20$ is the one required by the summation of the domains of the chain terms at the rate $\kappa/20$ obtained from the split of the spare. In addition, $M$ is chosen after $\kappa_1$ and $\kappa_R$.

\emph{Ceilings.} The ceilings are members of the list defining $\gamma_U$ in \eqref{9.1:eq-constant-order-a}. They are the strip ceiling $(\gamma_\sharp)$, $\varepsilon_{br}\le1$, the bound on $e^{-\frac12p_0(g_k)}100^d\check{R}_k^{-d}$, and $(\gamma2)$--$(\gamma4)$. In $(\gamma_\sharp)$, $T=\log g^{-2}$ is the degree variable and $T_\gamma$ its threshold, so that the supremum is over the couplings $g$ with $\log g^{-2}\ge T_\gamma$; $\check{R}(g)$ is the scaffold radius at coupling $g$, $N(g)$ is the number of stages of the chain terms at coupling $g$ and $C_0^{\sharp}$ the constant of one stage. In $(\gamma4)$, $c_{vol}$ is the volume constant of the modified bound, taken at $|Z_k\cap X|$, and $C_{\log}\log g_k^{-2}\,|Z_k|$ is Ba{\l}aban's term, $C_{\log}$ being the constant written $O(1)$ in it in \cite{B-CMP122b}; the inequality is a condition on $g_k$, namely the ceiling under which the modified volume term
\begin{equation*}
\bigl(2^dc_{vol}+(2^d-1)\mathsf{K}_{cell}(k)\bigr)\,|Z_k\cap X|
\end{equation*}
stays under Ba{\l}aban's term $C_{\log}\log g_k^{-2}\,|Z_k|$ for every $X$.
\end{definition}

\begin{proof}
We verify the identities and implications stated above. By Notation~\ref{9.1:not-constant-order}, $\beta^{\star}=\beta_s=\tfrac14$ and $\beta_{pr}=\beta_I=\tfrac15$.

Since $1+4\cdot\tfrac14=2$ and $\kappa_R=c^{\sharp}\kappa$, we have $(1+4\beta^{\star})\kappa_R=2c^{\sharp}\kappa$. This is the equality in \eqref{9.1:eq-hull-conventions-1}. Likewise $1+2\cdot\tfrac14=\tfrac32$, so $b=\tfrac32\kappa$, and $a-b=\beta_\sharp\kappa$ by definition.

The split of the spare is the identity
\begin{equation*}
\tfrac85\kappa-\tfrac32\kappa=\tfrac1{10}\kappa=\tfrac1{20}\kappa+\tfrac1{20}\kappa .
\end{equation*}
For the output and input rates of the uniqueness argument,
\begin{equation*}
(1+3\beta_I)\kappa-(1+2\beta_I)\kappa=\beta_I\kappa=\tfrac15\kappa\ge\tfrac1{20}\kappa=\beta_\sharp\kappa .
\end{equation*}

By the hull-length bound the constant $C_d$ is non-negative; then $c_H\ge1$, and $c^{\sharp}=c_H+1\ge2$. Hence
\begin{equation*}
8c^{\sharp}c_g^2=8c^{\sharp}c_H^2\ge 16c_H^2\ge 3c_H ,
\end{equation*}
so $(F_\star)$ implies $(F_L)$. Furthermore
\begin{equation*}
8c^{\sharp}c_H^2=\tfrac{35}{8}c_Hc^{\sharp}\cdot\tfrac{64}{35}c_H\ge\tfrac{35}{8}c_Hc^{\sharp},
\end{equation*}
because $\tfrac{64}{35}c_H\ge1$. So $(F_\star)$ implies $L\ge\tfrac{35}{8}c_Hc^{\sharp}$ as well.
\end{proof}

\begin{proposition}[The correlation step read through leaf hulls]\label{9.1:prop-hull-step}
Consider a single run of the renormalisation step $k\to k+1$ used in the proof of the correlation
theorem, at the three sub-steps of that step that carry the sites (T), (P) and (L),
read in the following way, which we call the \emph{hull-read step}.
\begin{itemize}
\item Every leaf $F$ is read through its hull
\begin{equation}\label{9.1:eq-hull-step-1}
\tilde F(X;P):=F\bigl(X;(U_j,J_j)\bigl(X^{+},M_{\mathfrak{B}_j(X^{+})}P\bigr){\upharpoonright}X\bigr).
\end{equation}
In this display $X^{+}$ is the enlarged region of the leaf $X$ whose part $(X^{+})_0$ is the core
$K(X)$, $P$ is the chart point, $(U_j,J_j)$ are the configuration and the current at level $j$
determined from $X^{+}$ and the data $M_{\mathfrak{B}_j(X^{+})}P$, and ${\upharpoonright}X$ is
restriction to $X$.
\item The cores $K(X)$ of \eqref{9.1:eq-hull-conventions-a} are used at the three sites of the
step: the small-field site (T), the site (P) and the large-field site (L).
\item At the second sub-step of the site (L), the filing rule of
\eqref{9.1:eq-hull-conventions-a} and the assignment $Y_0\mapsto K(X)$ are used.
\item The derived objects are formed on Ba{\l}aban's cores and on his sets $Y_1$.
\item The last resummation is carried out over the strip union $Z^{\sharp}$.
\item The site (T) is run at the fat cores of hypothesis (2) below, as in the run at that site in
the proof of the uniqueness theorem.
\end{itemize}
Assume the following.
\begin{enumerate}
\item[(1)] \emph{Containment at the three sites.}
\begin{itemize}
\item At the site (P) the containment hypothesis is taken in the form of the third and fourth rows
of the second containment theorem of this chapter. The hull-read step runs its site (P) on the pair
radii $r^{\bullet}_{\square}$ (these give a smaller circle in the same domain of holomorphy, and
the piece at the small-field site attached to the box $\square$ does not depend on the circle; the
activity constants $C_2$, $\vartheta_P^{\bullet}(\square)$ and $\vartheta_P^{0}$ of the second
containment theorem obey
\begin{equation*}
2e^{C_2\kappa_1}\vartheta_P^{\bullet}(\square)\le 2e^{C_2\kappa_1}\vartheta_P^{0},
\end{equation*}
so the degree of the activity bound is unchanged). This reading is made under the floor
\begin{equation*}
\tfrac18\delta_1 M\ge(1+3\beta_I)\kappa+c_d,
\end{equation*}
which is assumed. On these radii the containment is supplied by the rate statement and by the
containment statement on the region $\mathfrak{O}$.
\item On the region $Z^{\mathrm{on}}$, over the seeds on which it is stated, the first containment
theorem of this chapter, read with the table $\mathsf{T}^{\flat}=\mathsf{T}$, is used as it stands.
\item At the sites (T) and (L), the containment statements of this chapter at those two sites are
assumed.
\end{itemize}
\item[(2)] The fat-core hypothesis, at the fat cores at which the site (T) is run.
\item[(3)] The flat first-part bound, that is, the flat form of clause (a) of the resummation
proposition in the proof of the correlation theorem. It is a consequence of the flat
zeroth-order averaging bound and of the localised response bound of this chapter.
\item[(4)] The leaf form of the tail bound.
\item[(5)] The component hypothesis, on the components fixed in the component statement of this
section, with the rows labelled (C) of the large-field bound.
\item[(6)] The flat form of \cite[(1.89)]{B-CMP122b}, in the form of the flat corollary of this
chapter.
\item[(7)] The Koteck\'y--Preiss step \cite{KP86} of the resummation, used by statement.
\item[(8)] The scaffold rule, which fixes $\pi$, $Z$, $\check{R}_k$, the cores $K_v$ and
$m_{\sharp}$, one rule for each pair of runs.
\item[(9)] The rows labelled (C) of the correlation theorem.
\end{enumerate}
Assume also the floors and coupling ceilings of \eqref{9.1:eq-hull-conventions-a} and
\eqref{9.1:eq-refiled-smallness-a}. Then the following hold.
\begin{itemize}
\item[(a)] \emph{Hull reads.}
\begin{itemize}
\item The site identities used in the correlation theorem hold with $\tilde F$ in place of
$F\circ\mathsf{R}_p$; this is algebraic exactness.
\item At every natural vertex, $\tilde F=F$.
\item The vertex-weight lemma of the correlation theorem holds as stated for the pieces of
$\tilde F$, with the entry for the weight $\mathfrak{N}(v)$ in the bracket display of that theorem
supplied by the fat-core hypothesis (2).
\end{itemize}
\item[(b)] \emph{Fresh rate for every family, uniformly in the potential.} Consider every
level-$(k+1)$ family created at the site (T). These are:
\begin{itemize}
\item the family $E$, including its $\log Z$-pieces made by random walks;
\item the family $D$;
\item the boundary families and pinned kernels of \cite[(3.26)--(3.47)]{B-CMP119};
\item the family $R''$;
\item on the flat family $\Sigma^{\flat}$ at the small-field site, its members
$\mathfrak{K}(X,\mathcal{S})$, $\mathbf{B}^{sl}$ and the terms free of the fluctuation field
$\psi$.
\end{itemize}
Each such family carries the factor
\begin{equation*}
e^{-\kappa_{out}d_{k+1}},\qquad \kappa_{out}\ge 2c^{\sharp}\kappa .
\end{equation*}
This holds for all potentials of the fat-core format in the ball of the potential family and for
all operators in the operator ball of the step. After the operation $\mathbf{R}$ these families
are stored at rate $\kappa$.
\item[(c)] \emph{Brackets at the site (L).} The boundary lemma of the correlation theorem, rows
one to three, five and the boundary row of the bracket display of that theorem, and the bracket
resummation of clause (c) of the resummation proposition keep the margins of the correlation
theorem. These margins are, with $\delta=\delta_\star$ the parameter of
Lemma~\ref{9.1:lem-bracket-margins}, fixed in Lemma~\ref{9.1:lem-fresh-rate} under the floor of
that lemma and lying in $[\tfrac1{20},\tfrac1{10})$:
\begin{itemize}
\item $(\tfrac{5}{12}-\tfrac14\delta)\kappa d_j$ and $\tfrac{17}{60}\kappa d_j$ for $j<k$;
\item $(\tfrac14-\tfrac14\delta)\kappa d_k$ for $j=k$;
\item $\tfrac12\delta\kappa d_j$ at the first sub-step of the site (L);
\item for lower chain leaves (pieces $\mathbb{C}^{(i)}$ filed below their top level), the margin
$\tfrac12 a_{\flat}d_{\ell}(X)$, which equals the third bound of the payment statement for chain
pieces.
\end{itemize}
Constants multiplied by $A_H$ are compensated by $\mathsf{w}_L^{-1}$.
\item[(d)] \emph{Values.}
\begin{itemize}
\item The strips $C^{\sharp}$ are pairwise non-touching boxes. Each contains exactly one component
$C$ of $\Lambda_k^c$, each lies inside Ba{\l}aban's collar of width $M\check{R}_k$, and
\begin{equation*}
|C^{\sharp}|\le\min\{2^d|C|,(101\check{R}_k-1)^d\},
\end{equation*}
where $|C|$ and $|C^{\sharp}|$ denote the numbers of cubes in $C$ and in $C^{\sharp}$.
\item The re-filing along $\operatorname{fil}_{\sharp}$ is linear, and
$d_{Z^{\sharp}}(\operatorname{fil}_{\sharp}Y_4)\le d_Z(Y_4)$.
\item Every exterior $Y_4$ has $d_Z(Y_4)\ge m_{\sharp}-1$.
\item Every first-part value and every subtracted copy is either grouped inside a strip
$C^{\sharp}$, where it falls in the volume class obtained from \cite[(1.69)]{B-CMP122b} by passing
to the strips and counted per cell, or obeys
\begin{equation}\label{9.1:eq-hull-step-2}
\begin{gathered}
|V^{\sharp}(Y)|\le\alpha^{\sharp}e^{-b\,d_{Z^{\sharp}}(Y)},\qquad
\mathsf{K}_{val}(k)\le\mathsf{K}^0+N(g_k)C_0^{\sharp},\\
\alpha^{\sharp}:=C_{fib}\bigl(\varepsilon_{br}+\mathsf{K}_{val}(k)\bigr)
e^{-\frac12\beta_{\sharp}\kappa(m_{\sharp}-1)}\le\alpha_*,
\end{gathered}
\end{equation}
with
\begin{equation*}
(\alpha^{\sharp})^{1/3}\le\min\bigl\{e^{-(1+\beta_s)\kappa 2^d},\,g_k^{\kappa_0}\bigr\};
\end{equation*}
here $N(g_k)$ and $C_0^{\sharp}$ are the function of the coupling and the constant of
Proposition~\ref{9.1:prop-value-cell-constants}, and $\kappa_0$ is the exponent of
Corollary~\ref{9.1:cor-refiled-smallness}.
\item $\alpha^{\sharp}\to0$ as $\gamma\to0$, uniformly in $k$. This is the smallness of $\alpha$
stated in \cite{B-CMP122b}: ``$\alpha$ \dots\ arbitrarily small for $\gamma$ small enough''.
\end{itemize}
\item[(e)] \emph{Ba{\l}aban's resummation under one dictionary.} Under the one dictionary
$\mathfrak{D}^{\sharp}$ to the strips the following hold:
\begin{itemize}
\item Ba{\l}aban's displays \cite[(1.71)--(1.100)]{B-CMP122b};
\item the remark numbered (3) accompanying the correlation theorem;
\item clauses (c)--(e) of the resummation proposition;
\item the modulus lemma of the correlation theorem.
\end{itemize}
Ba{\l}aban's displayed term $2\kappa(100R_k)^d$ of \cite{B-CMP122b}, read with
$R_k=\check{R}_k$, still absorbs the strip. For the
remaining clauses of the resummation proposition: clause (a) is hypothesis (3), clause (b) holds as
stated, and clause (f) follows from clause (b) above.
\item[(f)] \emph{Left-over chain terms; counting of units.}
\begin{itemize}
\item The chain terms $\mathbb{C}^{(i)}$ left over from earlier steps are stored at rate
$a_{\flat}$ below their top level and at rate $(1+3\beta_{pr})\kappa$ at it.
\item They are never read again at their own level.
\item The hull of a leaf of the same scale at a level $m<k$ is paid for by the rate of the
output.
\item The rate $a_{\flat}$ enters only the four bounds of the payment statement for chain pieces
and clauses (c) and (d) above.
\end{itemize}
\item[(g)] \emph{Pairs of runs.}
\begin{itemize}
\item The two runs of a pair share the scaffold: $\pi$, $Z$, $\check{R}_k$, the cores $K_v$ and
$m_{\sharp}$, by hypothesis (8), one rule for each pair.
\item Differences and affine arrays are filed alike.
\item For a pair no family arises beyond the families of the two runs, their differences and
their affine arrays, and all of these are filed by the same map $\operatorname{fil}_{\sharp}$.
\end{itemize}
\end{itemize}
Consequently the hull-read step satisfies the hull form of the relative correlation bounds. In the
proof of the uniqueness theorem, this input is thereby reduced to two things: the floors and
coupling ceilings of \eqref{9.1:eq-hull-conventions-a} and \eqref{9.1:eq-refiled-smallness-a}, and
hypotheses (1)--(9). In this proposition the correlation theorem is used as stated, read with
Ba{\l}aban's sets $Y_1$; of the present section it uses only the component statement.
\end{proposition}

\begin{proof}
Each clause is the conclusion of statements of this section. Each is applied to the hull-read step
under hypotheses (1)--(9) and under the floors and coupling ceilings of
\eqref{9.1:eq-hull-conventions-a} and \eqref{9.1:eq-refiled-smallness-a}, with the hull
parameters fixed in Notation~\ref{9.1:not-hull-conventions}.

\emph{Clause (a).} This is Lemma~\ref{9.1:lem-hull-exactness} (exactness and analyticity of hull
reads), applied to the leaves read through the hulls \eqref{9.1:eq-hull-step-1}. That lemma gives
three things: the site identities with $\tilde F$ in place of $F\circ\mathsf{R}_p$; the equality
$\tilde F=F$ at the natural vertices; and the vertex-weight lemma of the correlation theorem for the
pieces of $\tilde F$. The entry for $\mathfrak{N}(v)$ is supplied by the fat-core hypothesis (2).

\emph{Clause (b).} This is Lemma~\ref{9.1:lem-fresh-rate} (the fresh rate at the small-field site
for every potential), applied at the site (T) of the step $k\to k+1$ run at the fat cores of
hypothesis (2). It covers every family listed in clause (b), all potentials of the fat-core format
in the ball, and all operators in the operator ball. It yields the rate $\kappa_{out}\ge
2c^{\sharp}\kappa$ and storage at rate $\kappa$ after the operation $\mathbf{R}$.

\emph{Clause (c).} This is Lemma~\ref{9.1:lem-bracket-margins} (brackets at the large-field site
keep their margins), applied at the site (L) of the hull-read step. It gives the margins listed in
clause (c), including the margin $\frac12 a_{\flat}d_{\ell}(X)$ for lower chain leaves, and the
compensation of the constants multiplied by $A_H$ by $\mathsf{w}_L^{-1}$.

\emph{Clause (d).} Five statements are used.
\begin{itemize}
\item Lemma~\ref{9.1:lem-strip-geometry} (geometry of the strips) gives the geometric assertions.
These are: the strips $C^{\sharp}$ are pairwise non-touching boxes; each holds exactly one
component $C$ of $\Lambda_k^c$ and lies inside Ba{\l}aban's collar; the bound on
$|C^{\sharp}|$ holds; the inequality
$d_{Z^{\sharp}}(\operatorname{fil}_{\sharp}Y_4)\le d_Z(Y_4)$ holds; and every exterior $Y_4$ has
$d_Z(Y_4)\ge m_{\sharp}-1$.
\item Theorem~\ref{9.1:thm-strip-resummation} (the strip resummation) gives the linearity of the
re-filing. It also gives the alternative between grouping inside a strip, in the volume class per
cell, and the exponential bound with rate $b$ in $d_{Z^{\sharp}}$.
\item Lemma~\ref{9.1:lem-column-bound} (from value units to domains: a column bound off the
large-field region) supplies the constant $\mathsf{K}_{val}(k)$ for the first-part values and
subtracted copies at the site (L).
\item Proposition~\ref{9.1:prop-value-cell-constants} (bounds on the value and cell constants)
gives the bound $\mathsf{K}_{val}(k)\le\mathsf{K}^0+N(g_k)C_0^{\sharp}$ in
\eqref{9.1:eq-hull-step-2}.
\item Corollary~\ref{9.1:cor-refiled-smallness} (smallness of the re-filed values at maximal strip
width) gives the bound \eqref{9.1:eq-hull-step-2} with the stated $\alpha^{\sharp}\le\alpha_*$. It
also gives the bound on $(\alpha^{\sharp})^{1/3}$, and the convergence $\alpha^{\sharp}\to0$ as
$\gamma\to0$ uniformly in $k$.
\end{itemize}

\emph{Clause (e).} The assertions on Ba{\l}aban's displays \cite[(1.71)--(1.100)]{B-CMP122b}, on
the remark numbered (3) accompanying the correlation theorem, and on clauses (c)--(e) of the
resummation proposition are Lemma~\ref{9.1:lem-strip-combinatorics} (Ba{\l}aban's combinatorics
under the strip translation). That lemma is read with the values of clause (d) and with
Ba{\l}aban's term $2\kappa(100R_k)^d$ of \cite{B-CMP122b}, where $R_k=\check{R}_k$. The assertion
on the modulus lemma is
Corollary~\ref{9.1:cor-strip-modulus} (the modulus bound under strips). Clause (a) of the
resummation proposition is hypothesis (3) itself. Its clause (b) holds as stated. Its clause (f)
follows from clause (b) of the present proposition.

\emph{Clause (f).} The first two assertions are the statement that chain pieces are never read
again at their own level, applied to the left-over terms $\mathbb{C}^{(i)}$. The third assertion
is the payment statement for chain pieces, which pays the hull of a same-scale leaf at a level
$m<k$ by the rate of the output. The fourth assertion records where $a_{\flat}$ occurs among the
statements used here: in the four bounds of that payment statement, and in clauses (c) and (d)
above, through Lemma~\ref{9.1:lem-bracket-margins} and the statements listed for clause (d).

\emph{Clause (g).} This is Corollary~\ref{9.1:cor-run-differences} (differences of two runs on
one scaffold). It is applied to the two runs of a pair, which share the scaffold by the scaffold
rule of hypothesis (8). The corollary shows that differences and affine arrays are filed by the
same map $\operatorname{fil}_{\sharp}$ as the values of one run: the difference of the re-filed
values of the two runs is the $\operatorname{fil}_{\sharp}$-sum of the difference of their value
families, and the re-filed value of an affine array is the $\operatorname{fil}_{\sharp}$-sum of
that array. Hence every family that arises for a pair is a family of one of the two runs, a
difference or an affine array, filed by $\operatorname{fil}_{\sharp}$.

\emph{The final assertion.} By clauses (a)--(g), the hull form of the relative correlation bounds
holds for the hull-read step. In the proof of the uniqueness theorem it is therefore reduced to two
things: the floors and coupling ceilings of \eqref{9.1:eq-hull-conventions-a} and
\eqref{9.1:eq-refiled-smallness-a}, and hypotheses (1)--(9).
\end{proof}

In what follows $v=(j,S_v,f_v)$ is a leaf unit, with level $j$, set $S_v$ and function $f_v$. By the defining conditions of a leaf unit, $f_v$ is analytic on its domain $\mathfrak{D}_v$ and invariant under the complexified gauge group. We write $X$ for the domain of $S_v$ and $X^{+}$ for its hull in the sense of Notation~\ref{9.1:not-hull-conventions}; see also \eqref{9.1:eq-hull-conventions-a}.

The leaf unit $v$ is read at a site. At this site, Ba{\l}aban's response map $\mathsf{R}_p(\mathsf{t},\mathsf{s})$ depends on two kinds of interpolation parameter: the bracket parameter $\mathsf{t}$, and the localisation parameters $\mathsf{s}=(\mathsf{s}(\Delta))_{\Delta}$, with $\Delta$ ranging over the cubes to which the localisation attaches a parameter. We write $\mathbb{1}$ for the vertex at which every $\mathsf{s}(\Delta)$ equals $1$. The subscript in $\mathsf{R}_p$ is part of the name of this map; it is not the piece index $p$ used below.

The sites are of three kinds, denoted (T), (P) and (L).
\begin{itemize}
\item At (T), the T-step of \cite[\S3]{B-CMP109} re-solves its variational problem.
\item At (L), the localisation of \cite[(1.59)--(1.63)]{B-CMP122b} operates.
\item At (P), as at (T), the identities \cite[(3.7)]{B-CMP109} and \cite[(1.10), (1.23)]{B-CMP116} are used; the radius attached to (P) is the pair radius $r^{\bullet}_{\square}$ introduced below.
\end{itemize}

Let $\mathfrak{B}_j(X^{+})$ be the determining set, on $X^{+}$, of the current label, truncated at level $j$. Equivalently, writing $\boldsymbol{\Omega}(X^{+})$ for the box-set geometry of the current label on $X^{+}$, the set $\mathfrak{B}_j(X^{+})$ is $\boldsymbol{\Omega}(X^{+})$ truncated at level $j$. Let $M_{\mathfrak{B}_j(X^{+})}$ be the block average over it; this is a polynomial map. For a block datum $\mathsf{v}$ on $X^{+}$, let $(U_j,J_j)(X^{+},\mathsf{v})$ be the minimiser on $X^{+}$ with data $\mathsf{v}$, together with its current, as given by \cite[Theorem~1, Proposition~9]{B-CMP102b}.

The \emph{hull read} of $v$ is
\begin{equation}\label{9.1:eq-hull-exactness-2}
g_v(\mathsf{t},\mathsf{s}):=f_v\Big((U_j,J_j)\big(X^{+},\,
M_{\mathfrak{B}_j(X^{+})}\mathsf{R}_p(\mathsf{t},\mathsf{s})\big){\upharpoonright}X\Big).
\end{equation}

Among the statements of this book used below is a lemma that bounds the pieces $v_p$ into which a site identity splits a term; we call it the Cauchy-estimate lemma. In it, $p$ indexes the pieces, not plaquettes. The piece $p$ carries a weight $\eta_p$ and a size $m(p)$, and the estimate is taken on closed polydiscs in $(\mathsf{t},\mathsf{s})$.

The Cauchy estimate of that lemma is taken on a circle of radius $\mathsf{w}_{\mathrm{site}}$, the dilation circle at the site:
\begin{itemize}
\item $\mathsf{w}_{\mathrm{site}}=\mathsf{w}_T(k)$, fixed through the constant $K_R^{\sharp}$ of the T-site, at a (T)-site of level $k$;
\item $\mathsf{w}_{\mathrm{site}}=\mathsf{w}_L(p)$, the room of the bracket circle, fixed through the circle constant $K_o^{L}$, at (L);
\item $\mathsf{w}_{\mathrm{site}}=r^{\bullet}_{\square}$, the pair radius, at (P).
\end{itemize}
The conditions that the order of choice places on these radii are those in item~(o8) of \eqref{9.1:eq-constant-order-a}.

Finally, $\mathfrak{T}[0](X)$ denotes the target tube of slices for $X$. By its definition, membership of a datum $\mathsf{v}$ in $\mathfrak{T}[0](X)$ gives the following two properties, which are all that is used below:
\begin{itemize}
\item $\mathsf{v}$ lies within the radius $r_c$ of \cite[Proposition~9]{B-CMP102b} of a real regular datum (this radius does not depend on the coupling);
\item the slice point of $\mathsf{v}$, that is the restricted minimiser $(U_j,J_j)(X^{+},\mathsf{v}){\upharpoonright}X$ at which $f_v$ is read, lies in the subdomain $\mathfrak{D}_v[0]\subset\mathfrak{D}_v$.
\end{itemize}

\begin{lemma}[Exactness and analyticity of hull reads]\label{9.1:lem-hull-exactness}
Let $v$, the site, $\mathsf{R}_p$ and $g_v$ be as above.
\begin{enumerate}
\item[(a)] \emph{Algebraic exactness.} Consider every site identity of Ba{\l}aban's construction that is used in the proof of the correlation theorem:
\begin{itemize}
\item \cite[(3.7)]{B-CMP109}, \cite[(1.10)]{B-CMP116} and \cite[(1.23)]{B-CMP116}, at (T) and (P);
\item \cite[(1.59), (1.60)]{B-CMP122b}, and the $\mathsf{s}$-localisations of \cite[(1.61)--(1.63)]{B-CMP122b}, at (L).
\end{itemize}
Each of them is either of the form
\begin{equation}\label{9.1:eq-hull-exactness-1}
g(1,\mathbb{1})-g(0,\cdot)=\sum_{p}\prod_{\Delta}\int_0^1 d\mathsf{s}(\Delta)\,
\partial_{\mathsf{s}(\Delta)}\,\partial_{\mathsf{t}}\big|_{0}\,g,
\end{equation}
where $g$ is a function of $(\mathsf{t},\mathsf{s})$, not the coupling, and $g(0,\cdot)$ denotes the value of $g$ at $\mathsf{t}=0$, as a function of $\mathsf{s}$,
or one of the Taylor or Mayer variants of \eqref{9.1:eq-hull-exactness-1} quoted in the Cauchy-estimate lemma. It is an identity in $(\mathsf{t},\mathsf{s})$. It is valid for every function $g=G\circ\mathsf{R}_p$ that is holomorphic on the polydisc; the identity \cite[(3.4)]{B-CMP109}, which is used at (T) together with these, is a chain rule and is likewise valid for every such $g$. Hence, wherever $g_v$ is holomorphic on the polydisc, in particular under the hypothesis \eqref{9.1:eq-hull-exactness-3} of~(c), it holds with $g_v$ in place of $f_v\circ\mathsf{R}_p$.
\item[(b)] \emph{Endpoints.} Let $(\mathsf{t},\mathsf{s})$ be a vertex with the following property: for some region $W\supseteq X^{+}$, the restriction $\mathsf{R}_p(\mathsf{t},\mathsf{s}){\upharpoonright}W$ is, at real data, the minimiser of a box-set geometry $\boldsymbol{\Omega}^{\bullet}(W)$ such that $\boldsymbol{\Omega}^{\bullet}(W)(X^{+})$ equals $\boldsymbol{\Omega}(X^{+})$ truncated at $j$. At every such vertex,
\begin{equation*}
g_v(\mathsf{t},\mathsf{s})=f_v\big(\mathsf{R}_p(\mathsf{t},\mathsf{s}){\upharpoonright}X\big).
\end{equation*}
Consider the natural-vertex route, that is, the route on which the interpolation vertices are those of the natural-interpolation lemma. On it, every vertex has this property, under the three hypotheses from which the vertex lemma for hull reads is derived: the natural standing hypothesis, the junction hypothesis and the $\rho$-form of the tube condition. Consequently, under these hypotheses, the expansion with hull reads coincides, term by term, with the expansion with the direct reads $f_v\circ\mathsf{R}_p$ at the same vertices.

In Ba{\l}aban's interpolation format \cite{B-CMP109,B-CMP122b}, the two endpoints of each bracket have this property. Indeed, the configurations $U'_k$, $U^0_k$, $U_{k,1}$, $U_{k,\Lambda}$ and $U_k$ are minimisers of their determining sets by \cite[\S1, (1.51), (1.53), (1.56)]{B-CMP122b}; and at (T) the variational problem is re-solved \cite[(3.3)--(3.4)]{B-CMP109}. Therefore, by~(a), the value of each bracket is unchanged. Only the decomposition of the bracket into pieces differs from Ba{\l}aban's.
\item[(c)] \emph{Cauchy bounds.} Assume the containment hypothesis
\begin{equation}\label{9.1:eq-hull-exactness-3}
M_{\mathfrak{B}_j(X^{+})}\mathsf{R}_p(\mathsf{t},\mathsf{s})\in\mathfrak{T}[0](X)
\quad\text{on the closed polydiscs of the Cauchy-estimate lemma.}
\end{equation}
Then $g_v$ is holomorphic on the polydisc of the Cauchy-estimate lemma. Moreover $|g_v-f_v(1)|\le\mathfrak{N}(v)$, where $\mathfrak{N}(v)$ denotes the supremum of $|g_v-f_v(1)|$ over that polydisc. Finally, the conclusion of the Cauchy-estimate lemma holds without change for the pieces $v_p$ of $g_v$:
\begin{equation}\label{9.1:eq-hull-exactness-4}
|v_p|\le 2\,\mathsf{w}_{\mathrm{site}}^{-1}\,\eta_p\,e^{-(\kappa_1-1)m(p)}\,\mathfrak{N}(v),
\end{equation}
and the pieces are analytic in every parameter.
\item[(d)] \emph{Sizes.} $\mathfrak{N}(v)$ is the entry for $v$ in the size table of the correlation theorem, read under the natural standing hypothesis on the value units $\mathcal{V}$ and the second clause of the tube condition in its form for leaf units.
\end{enumerate}
\end{lemma}

\begin{proof}
(a) We treat identities of the form \eqref{9.1:eq-hull-exactness-1} first. Each is obtained by three operations:
\begin{itemize}
\item the fundamental theorem of calculus, applied in each variable $\mathsf{s}(\Delta)$ and in $\mathsf{t}$;
\item the product rule;
\item finitely many algebraic rearrangements of the resulting finite sums and products.
\end{itemize}
None of these operations uses what $g$ is. They need only $g$ to be a function $G\circ\mathsf{R}_p$ holomorphic on the polydisc. The Taylor and Mayer variants quoted in the Cauchy-estimate lemma are obtained from the same operations, and so are subject to the same remark. The identity \cite[(3.4)]{B-CMP109}, used at (T) together with these, is a chain rule, so it too is valid for every such $g$.

Unwinding \eqref{9.1:eq-hull-exactness-2}, the hull read is $g_v=G_v\circ\mathsf{R}_p$ with
\begin{equation*}
G_v(\cdot)=f_v\Big((U_j,J_j)\big(X^{+},M_{\mathfrak{B}_j(X^{+})}\,\cdot\,\big){\upharpoonright}X\Big).
\end{equation*}
Hence each identity holds with $g_v$ in place of $f_v\circ\mathsf{R}_p$, wherever $g_v$ is holomorphic on the polydisc; clause~(c) provides this holomorphy under \eqref{9.1:eq-hull-exactness-3}.

(b) Let $\hat{\mathfrak{B}}$ be a determining set, $\mathfrak{U}$ a configuration, and $f$ a holomorphic function invariant under the complexified gauge group. Write $(U,J)_{\hat{\mathfrak{B}}}(\mathsf{v})$ for the minimiser with determining set $\hat{\mathfrak{B}}$ and data $\mathsf{v}$, together with its current, so that $(U,J)_{\mathfrak{B}_j(X^{+})}(\mathsf{v})=(U_j,J_j)(X^{+},\mathsf{v})$; and write $\mathfrak{U}{\upharpoonright}\cdot$ for the restriction of $\mathfrak{U}$ to the region carrying $\hat{\mathfrak{B}}$. We use the following identity; we refer to it as the conclusion of the lemma on re-minimisation of minimisers:
\begin{equation}\label{9.1:eq-hull-exactness-5}
f\big((U,J)_{\hat{\mathfrak{B}}}(M_{\hat{\mathfrak{B}}}\mathfrak{U})\big)
=f(\mathfrak{U}{\upharpoonright}\cdot\,).
\end{equation}
We use it under the hypotheses of that lemma, in the situation in which it is applied here: $\mathfrak{U}$ is, at real data, the minimiser of a geometry whose restriction to the region carrying $\hat{\mathfrak{B}}$ is $\hat{\mathfrak{B}}$. That lemma is proved by the identity theorem on the totally real slice over the unique critical orbit of \cite{B-CMP102b}.

We apply \eqref{9.1:eq-hull-exactness-5} with
\begin{equation*}
f=f_v\circ(\,\cdot\,{\upharpoonright}X),\qquad
\hat{\mathfrak{B}}=\mathfrak{B}_j(X^{+}),\qquad
\mathfrak{U}=\mathsf{R}_p(\mathsf{t},\mathsf{s}){\upharpoonright}W .
\end{equation*}
By its definition above, $\mathfrak{B}_j(X^{+})$ is the determining set of the current label on $X^{+}$ truncated at level $j$, that is, $\boldsymbol{\Omega}(X^{+})$ truncated at $j$.
The situation described after \eqref{9.1:eq-hull-exactness-5} obtains here, for three reasons:
\begin{itemize}
\item $f$ is holomorphic and invariant under the complexified gauge group, since $f_v$ is so by the defining conditions of a leaf unit, and restriction to $X$ is a coordinate projection that commutes with the action of that group;
\item at the vertex in question, $\mathsf{R}_p(\mathsf{t},\mathsf{s}){\upharpoonright}W$ is, at real data, the minimiser of $\boldsymbol{\Omega}^{\bullet}(W)$;
\item the restriction of $\boldsymbol{\Omega}^{\bullet}(W)$ to $X^{+}$ is $\boldsymbol{\Omega}(X^{+})$ truncated at $j$, which is $\hat{\mathfrak{B}}$.
\end{itemize}
The left side of \eqref{9.1:eq-hull-exactness-5} is then $g_v(\mathsf{t},\mathsf{s})$, since $(U,J)_{\mathfrak{B}_j(X^{+})}=(U_j,J_j)(X^{+},\cdot)$, and the right side is $f_v(\mathsf{R}_p(\mathsf{t},\mathsf{s}){\upharpoonright}X)$, since the region carrying $\mathfrak{B}_j(X^{+})$ is $X^{+}\supseteq X$. This is the identity asserted in~(b).

The configurations listed in~(b), with the citations given there from \cite{B-CMP122b} and \cite{B-CMP109}, are minimisers of their determining sets at the endpoints of Ba{\l}aban's brackets, and at (T) the variational problem is re-solved. Thus the two endpoints of each bracket satisfy the hypothesis of~(b). Consequently, by the identity just proved, the endpoint values of $g_v$ and of $f_v\circ\mathsf{R}_p$ agree. By~(a), the value of each bracket is the difference of its endpoint values, for either function. The value of each bracket is therefore unchanged, and only its decomposition into pieces differs.

On the natural-vertex route, we assume the three hypotheses from which the vertex lemma for hull reads is derived: the natural standing hypothesis, the junction hypothesis and the $\rho$-form of the tube condition. Under them, the assertion for every vertex has two parts:
\begin{itemize}
\item each vertex satisfies the hypothesis of~(b): this follows from clauses (e) and (f) of the natural-interpolation lemma together with clause~(a) of the vertex lemma for hull reads;
\item no inserted or cut stratum of level at most $j$ passes through $X$: this is clause~(b) of the vertex lemma for hull reads.
\end{itemize}
Hence the identity of~(b) holds at every vertex, and the expansion with hull reads coincides term by term with the expansion with the direct reads $f_v\circ\mathsf{R}_p$ at the same vertices.

(c) By \cite[Proposition~9]{B-CMP102b}, extended to complex data by the analyticity statement for the box minimiser at complex data, the map
\begin{equation*}
\mathsf{v}\longmapsto(U_j,J_j)(X^{+},\mathsf{v})
\end{equation*}
is analytic on the ball of radius $r_c$ about each real regular datum.

Fix $(\mathsf{t},\mathsf{s})$ in the closed polydisc of the Cauchy-estimate lemma. Apply \eqref{9.1:eq-hull-exactness-3} and the first of the two properties of membership in $\mathfrak{T}[0](X)$. They show that the datum $M_{\mathfrak{B}_j(X^{+})}\mathsf{R}_p(\mathsf{t},\mathsf{s})$ lies in such a ball. The map $\mathsf{R}_p$ is holomorphic and the block average $M_{\mathfrak{B}_j(X^{+})}$ is polynomial. Hence the composite
\begin{equation*}
(\mathsf{t},\mathsf{s})\longmapsto
(U_j,J_j)\big(X^{+},M_{\mathfrak{B}_j(X^{+})}\mathsf{R}_p(\mathsf{t},\mathsf{s})\big)
\end{equation*}
is holomorphic on the polydisc. Restriction to $X$ forgets the variables outside $X$; it is a coordinate projection and preserves holomorphy, so the composite remains holomorphic after restriction.

By the second of the two properties of membership in $\mathfrak{T}[0](X)$, the slice point of this datum lies in $\mathfrak{D}_v[0]\subset\mathfrak{D}_v$, where $f_v$ is analytic; this slice point is the argument of $f_v$ in \eqref{9.1:eq-hull-exactness-2}. Therefore $g_v$, as defined in \eqref{9.1:eq-hull-exactness-2}, is holomorphic on the polydisc. The bound $|g_v-f_v(1)|\le\mathfrak{N}(v)$ holds by the definition of $\mathfrak{N}(v)$ as the supremum over the polydisc.

The proof of the Cauchy-estimate lemma is a Cauchy estimate on the circle of radius $\mathsf{w}_{\mathrm{site}}$. Nothing in it depends on the site or on the term. It therefore applies unchanged to the pieces $v_p$ into which the identities of~(a) split $g_v$, and it yields \eqref{9.1:eq-hull-exactness-4} and, as that lemma asserts, the analyticity of the pieces in every parameter.

(d) This is the content of clauses (a)--(c) of the size lemma for leaf reads. That lemma uses two inputs:
\begin{itemize}
\item the size lemmas of the correlation theorem, applied at the vertices whose box (clause~(c1) of those lemmas) lies in the regular region of the chart point intersected with the interior of a set fixed in the size lemma for leaf reads;
\item the boundary size lemma for hulls, in which the box of the size lemma of the correlation theorem is replaced by the first-order members on $X^{+}$; the boundary lemma holds under the second clause of the tube condition, in its form for leaf units.
\end{itemize}
\end{proof}

\begin{remark}[What remains to be compared after the hull reads]
By Lemma~\ref{9.1:lem-hull-exactness}, the object of the comparison of the hull-read expansion with Ba{\l}aban's is the resummation geometry with cores $K(X)$ of \eqref{9.1:eq-hull-conventions-a}. Among its constituents are the filings $Y_0\mapsto K(X)$ and $Y_1\mapsto Y_1^{+}$ and the class boundary $Z^{\sharp}$. The geometry is taken at the fixed sizes $\mathfrak{N}(v)$.
\end{remark}

\begin{lemma}[The fresh rate at the small-field site for every potential]\label{9.1:lem-fresh-rate}
Let $\mathfrak{c}$ be a value of the hull parameter; in this chapter it is read at
$\mathfrak{c}=c_H$ (Notation~\ref{9.1:not-hull-conventions}). Put
\begin{equation}\label{9.1:eq-fresh-rate-1}
\delta_\star^{(\mathfrak{c})}:=\frac{1}{10}\Bigl(1-\frac{2(1+4\beta^{\star})(\mathfrak{c}+1)c_g^2}{L}\Bigr),
\qquad
\kappa_R^{(\mathfrak{c})}:=(\mathfrak{c}+1)\kappa .
\end{equation}
At $\mathfrak{c}=c_H$ we have $\mathfrak{c}+1=c^{\sharp}$, and these are the parameter $\delta_\star$ and the
rate $\kappa_R$ of Notations~\ref{9.1:not-hull-conventions}
and~\ref{9.1:not-radius-exponent-floors}. Assume the floor
\begin{equation}\label{9.1:eq-fresh-rate-2}
(F_\star)\qquad L\ \ge\ 4(1+4\beta^{\star})(\mathfrak{c}+1)c_g^2 ,
\end{equation}
so that $\delta_\star^{(\mathfrak{c})}\in[1/20,1/10)$; at $\mathfrak{c}=c_H$ this is the floor
$(F_\star)$ of \eqref{9.1:eq-hull-conventions-a}. At the small-field site of the renormalisation step
from level $k$ to level $k+1$, run the expansion in the form of the fat-core cluster-expansion
proposition with $\delta:=\delta_\star^{(\mathfrak{c})}$. Choose $\kappa$ after $L$ and
$\delta_\star^{(\mathfrak{c})}$, subject to the floor $(F_\kappa)$ of
\eqref{9.1:eq-hull-conventions-a} and to the lower bounds on $\delta\kappa$ required in that
proposition. Choose $\kappa_1\ge 4L\kappa+1$, and choose $M$ and $\varepsilon_1$ as there. Let
$\kappa_{out}$ be the output rate of this expansion.

Then
\begin{equation}\label{9.1:eq-fresh-rate-3}
\kappa_{out}\ \ge\ (1+4\beta^{\star})\kappa_R^{(\mathfrak{c})}\ \ge\ 2(\mathfrak{c}+1)\kappa ,
\end{equation}
with equality in the first inequality at $\delta=\delta_\star^{(\mathfrak{c})}$, and every family of
terms created at this site at level $k+1$ obeys the bound assigned to it in the
correlation theorem, with $e^{-\kappa d_{k+1}(X)}$ replaced by $e^{-\kappa_{out}d_{k+1}(X)}$
for every localisation region $X$. In
detail:
\begin{itemize}
\item[(i)] the family $E^{(k+1)}$ of localised terms created at this site, in its frozen and in its
local members, obeys the fat-core form of \cite[(2.41)]{B-CMP116} at the rate $\kappa_{out}$;
\item[(ii)] the kernels of the expectation values and the boundary families $\mathfrak{B}^{(k+1)}$
of \cite[(3.26)--(3.47)]{B-CMP119} (denoted $B^{(k+1)}$ in the proof of the correlation theorem)
obey their bounds at the rate $\kappa_{out}$. They keep the
prefactors they carry there, namely the constant $B_0$ of \cite[(2.42)]{B-CMP119} and
$O(p_0^3(g_k))\,g_k$, where $p_0(\cdot)$ is the degree function;
\item[(iii)] the marked sub-sum $R''^{(k+1)}$ of the proposition on the marked sub-sum in the proof of
the correlation theorem is bounded by
$\tfrac12 g_{k+1}^{\kappa_0}e^{-\kappa_{out}d_{k+1}(X)}$, with $\kappa_0$ the exponent of the power
of $g_{k+1}$ in the bound of that proposition;
\item[(iv)] the difference $D^{(k+1)}:=E^{(k+1)}(\cdot\,;w)-E^{(k+1)}_{w(p_y)}$, that is, the local
family of (i) at the source $w$ minus the same family at the source $w(p_y)$ of the correlation
theorem, is bounded by the
factor of the Schwarz-factor lemma of the correlation theorem times $e^{-\kappa_{out}d_{k+1}(X)}$;
\item[(v)] some pieces of $E^{(k+1)}$ and $\mathfrak{B}^{(k+1)}$ are not outputs of
\cite[(2.14)--(2.41)]{B-CMP116}. These are the pieces of $\log Z^{(k)}$ of
\cite[(1.1.6)--(1.1.7)]{B-CMP116}, the second integral of \cite[(3.37)]{B-CMP119}, which carries only
the Gaussian action, and the operators of \cite[(3.37)]{B-CMP119} represented by random-walk
expansions. The first of these decay at the rate $\delta_0M$ per cube, the other two at the rate
$\kappa_1$ per cube, and all three carry $e^{-\kappa_{out}d_{k+1}(X)}$ under $\delta_1M\ge\kappa_1$;
here $\delta_0M$ is the rate that replaces $\kappa$ in the bound \cite[(1.1.18)]{B-CMP116} for the
localized pieces of $\log Z^{(k)}$, and $\delta_1$ is the constant of \cite{B-CMP116} in the
condition $\delta_1M\ge\kappa_1$ quoted in the proof;
\item[(vi)] statements (i)--(v) hold, with the same constants, for all potentials in the ball of the
cluster-expansion lemma of the correlation theorem that are of the fat-core form of Ba{\l}aban's
format of potentials, and for all operators in the operator ball of the tolerance statement
for operators;
\item[(vii)] on the flat family at the small-field site, the members of $\Sigma^{\flat}$ carry the same
rate $\kappa_{out}$.
\end{itemize}
\end{lemma}

\begin{proof}
The proofs of (i)--(vi) use neither the bound on terms created at an earlier level nor the convention
that fixes $\kappa_{out}$ for the flat family. The proof of (vii) uses that convention only as a
statement about an arbitrary rate that is extracted for fresh terms from the untilted expansion,
that is, from the expansion of the flat family with the tilt described in the proof of (vii)
switched off.

In this proof only, write $r:=(L/2c_g^2)\kappa$ for the reference rate. By the output-rate lemma in
the parameter $\delta$, the output rate of the
fat-core expansion is $(1-10\delta)r$ for every $\delta\in(0,1/10)$. At
$\delta=\delta_\star^{(\mathfrak{c})}$, definition \eqref{9.1:eq-fresh-rate-1} gives
\begin{equation}\label{9.1:eq-fresh-rate-4}
(1-10\delta_\star^{(\mathfrak{c})})\,r
=\frac{2(1+4\beta^{\star})(\mathfrak{c}+1)c_g^2}{L}\cdot\frac{L}{2c_g^2}\,\kappa
=(1+4\beta^{\star})\kappa_R^{(\mathfrak{c})} .
\end{equation}
This is the equality case of the first inequality in \eqref{9.1:eq-fresh-rate-3}.

\emph{(i).} The fat-core cluster-expansion proposition passes from its fat-core small-field hypothesis,
through its fat-core lemma, to the fat-core form of \cite[(2.41)]{B-CMP116}. The output rate there
is $(1-10\delta)r$, and at $\delta=\delta_\star^{(\mathfrak{c})}$ this rate is
$(1+4\beta^{\star})\kappa_R^{(\mathfrak{c})}$ by \eqref{9.1:eq-fresh-rate-4}. This reproduces the
proposition's output statement at $\delta=\delta_\star$.

The proof of that proposition uses $\delta$ in two ways only. First, it uses the membership
$\delta\in[1/20,1/10)$, which enters the lower bounds on $\delta\kappa$ imposed on $\kappa$, the sums of
\cite[(2.27)]{B-CMP116} taken at the rate $\delta r$, and the sum over the sets $Y_0'$ of
\cite[\S2]{B-CMP116}. Second, it uses the definition $(1-10\delta)r$ of the
output rate. It therefore applies at $\delta=\delta_\star^{(\mathfrak{c})}$ as soon as
$\delta_\star^{(\mathfrak{c})}\ge1/20$.

By \eqref{9.1:eq-fresh-rate-1}, this condition reads
\begin{equation*}
\delta_\star^{(\mathfrak{c})}\ge\frac1{20}
\iff 1-\frac{2(1+4\beta^{\star})(\mathfrak{c}+1)c_g^2}{L}\ge\frac12
\iff L\ge4(1+4\beta^{\star})(\mathfrak{c}+1)c_g^2 ,
\end{equation*}
which is $(F_\star)$ in \eqref{9.1:eq-fresh-rate-2}. The bound $\delta_\star^{(\mathfrak{c})}<1/10$
holds because the subtracted term in \eqref{9.1:eq-fresh-rate-1} is positive. This proves (i).

\emph{(ii).} At this site, the cluster-expansion lemma of the correlation theorem is the hereditary
cluster-expansion lemma, applied with its four hypotheses taken by statement. Every output of it comes
from the displays \cite[(2.14)--(2.41)]{B-CMP116} run on the activities of the input family, including
the pinned kernels of \cite[(3.47)]{B-CMP119}, the boundary kernels and the kernels of $R''$.

Three statements of \cite[\S3]{B-CMP119} are used here.
\begin{itemize}
\item It states: ``The expressions in the expectation values in (3.26), (3.28), (3.37) are expanded
and localized in the same way as the corresponding expressions in the actions''.
\item Of the pinned factor it states: ``Although different, it satisfies the same bound as the other
factors''.
\item It states that the resulting terms ``satisfy the bounds (2.42), with the constant $B_0$ replaced
by $O(p_0^3(g_k))$ \dots\ the multiplication by $g_k$ yields small bounds''.
\end{itemize}

The input family enters the expansion in two places only.
\begin{itemize}
\item[($\alpha$)] The bound on the activity of the potential $V$ is used at
\cite[(2.18)]{B-CMP116}; there the term $\tau(Y)$ of that display, indexed by the set $Y$, has its
size $|\tau(Y)|$ fixed by this activity bound. After \cite[(2.20)]{B-CMP116} no
potential occurs, and the displays \cite[(2.27)--(2.41)]{B-CMP116} see the input only through
the parameter $\varepsilon_2=2\bar\varepsilon'/\alpha_4$, formed from Ba{\l}aban's constants
$\bar\varepsilon'$ and $\alpha_4$ of \cite[\S2]{B-CMP116}, and monotonically.
\item[($\beta$)] One pinned factor occurs, and it is bounded like the others.
\end{itemize}

The decay exponent is produced in \cite[(2.27)--(2.41)]{B-CMP116}, by the following steps:
\begin{itemize}
\item the sums of \cite[(2.27)]{B-CMP116} at the rate $\delta r$;
\item the sums over pairs $(Y_0,P)$;
\item \cite[(2.36)]{B-CMP116};
\item the step at \cite[(2.32)]{B-CMP116}, which is the third of the three places singled out in the
proof of the fat-core cluster-expansion proposition;
\item \cite[(2.39)--(2.41)]{B-CMP116}.
\end{itemize}
None of these steps reads which input family is expanded. They are the three places singled out in the
proof of the fat-core cluster-expansion proposition together with its core-blind steps, and that
proposition states them for inputs of each of the three formats in which it is formulated, the third
being the fat-core form of Ba{\l}aban's format of potentials, that is, for a format and not for a
particular family. Hence each input
family carries $e^{-\kappa_{out}d_{k+1}(X)}$, with its own prefactor. This proves (ii).

\emph{(iii).} The last clause of the cluster-expansion lemma of the correlation theorem bounds the
marked sub-sum, by the cluster-expansion corollary invoked there, by
\begin{equation*}
K_{CE}\,C_A^{-1}\,\bigl(\text{activity of }V^{m}\bigr)\,e^{-\kappa d_{k+1}(X)} ,
\end{equation*}
where $K_{CE}$ is the constant of that corollary, $C_A$ the constant of the cluster-expansion lemma,
and $V^{m}$ the potential of the marked sub-sum.
In this bound $e^{-\kappa d_{k+1}(X)}$ is the output exponential of the same displays
\cite[(2.27)--(2.41)]{B-CMP116}. By the argument of (ii) it is therefore
$e^{-\kappa_{out}d_{k+1}(X)}$. The arithmetic of sizes in the proposition on the marked sub-sum does
not involve the rate, so it yields
$\tfrac12 g_{k+1}^{\kappa_0}e^{-\kappa_{out}d_{k+1}(X)}$. This proves (iii).

\emph{(iv).} $D^{(k+1)}$ is the difference of the local family of (i) at $w$ and at $w(p_y)$. By (i),
both terms are bounded by $E_0e^{-\kappa_{out}d_{k+1}(X)}$, where $E_0$ is the constant of
\cite[(1.1.18)]{B-CMP116}. The Schwarz-factor lemma of the
correlation theorem supplies a Schwarz factor on that bound. This proves (iv).

\emph{(v).} For the pieces of $\log Z^{(k)}$, \cite{B-CMP116} states: ``The expression localized in
$X$ satisfies the bound (1.1.18) with $\kappa$ replaced by $\delta_0M$, and with an absolute constant
instead of $E_0$ \dots\ we take $M$ sufficiently large, so that $\delta_0M\ge\kappa$''. The lower bound
chosen there yields only the rate $\kappa$. But these pieces of $\log Z^{(k)}$ are the parts of order
one of $E^{(k+1)}$ and $\mathfrak{B}^{(k+1)}$, with the constants $\tfrac12E_0$ and $B_0$, so they
must themselves carry the rate $\kappa_{out}$; this is secured by \eqref{9.1:eq-fresh-rate-5} below.

Of the second integral of \cite[(3.37)]{B-CMP119}, \cite[\S3]{B-CMP119} states that it contains
``only the Gaussian action''; its decay comes from $\kappa_1$. The operators in the first integral
``are represented by the generalized random walk expansions'', which produce terms satisfying
\cite[(2.42)]{B-CMP119}.

We use the chain
\begin{equation}\label{9.1:eq-fresh-rate-5}
\delta_0M\ \ge\ \delta_1M\ \ge\ \kappa_1\ >\ 4L\kappa\ \ge\ 8\kappa_{out} .
\end{equation}
The first inequality amounts to $\delta_0\ge\delta_1$, since $M>0$. The second inequality is the
condition ``for $\delta_1M\ge\kappa_1$'' of
\cite{B-CMP116}. The third
follows from $\kappa_1\ge4L\kappa+1$, which is the condition $\tfrac12(\kappa_1-1)\ge2L\kappa$ used in
\cite[\S2]{B-CMP116}. For the last, the output rate is $\kappa_{out}=(1-10\delta)(L/2c_g^2)\kappa$ with
$0<1-10\delta\le1$ and $c_g\ge1$, so $\kappa_{out}\le\tfrac12L\kappa$.

By \eqref{9.1:eq-fresh-rate-5}, both per-cube rates $\delta_0M$ and $\kappa_1$ exceed
$8\kappa_{out}$, and under this margin these pieces carry $e^{-\kappa_{out}d_{k+1}(X)}$. The condition
$\delta_1M\ge\kappa_1$ is a lower bound on $M$, chosen after $\kappa_1$ as in Ba{\l}aban's order of
choice; it is the lower bound on $M$ entered as item (o4) of this chapter's list of floors and
ceilings chosen after $L$. This proves (v).

\emph{(vi).} The proof of (i)--(iv) reads a potential only through its activity bound at
\cite[(2.18)]{B-CMP116}, and an operator only through the ball of \cite[(2.16)]{B-CMP116}. The proof of
(v) uses only the quoted statements of \cite{B-CMP116} and \cite{B-CMP119} and the chain
\eqref{9.1:eq-fresh-rate-5}, and reads no further property of the potential or of the operators.
Hence the argument holds verbatim for every potential of the fat-core
format in the ball of the cluster-expansion lemma and for every operator in the operator ball of the
tolerance statement, with the same constants. This is (vi).

\emph{(vii).} Clause (c) of the flat-family bound states
\begin{equation*}
\sum_{\mathcal{S}}e^{\kappa_A^{pr}(g_k)|\mathcal{S}|}\,
\mathfrak{N}\bigl(\mathfrak{K}(X,\mathcal{S})\bigr)\ \le\ \tfrac18B_0\,e^{-\kappa_{out}d_{k+1}(X)} ,
\end{equation*}
with $\kappa_{out}$ the rate of the untilted expansion; here $\mathfrak{N}$ is the size functional
applied to the members $\mathfrak{K}(X,\mathcal{S})$ of the flat family, and $\kappa_A^{pr}(g_k)$ is the
coupling-dependent rate weighting the size $|\mathcal{S}|$ of the filing $\mathcal{S}$. By the
convention that fixes
$\kappa_{out}$ for
the flat family, $\kappa_{out}$ may be any rate extracted for fresh terms from the untilted expansion,
for instance $(1+4\beta^{\star})\kappa_R$. In the argument establishing that convention, the displays
\cite[(2.27)--(2.30)]{B-CMP116} and \cite[(2.35)--(2.41)]{B-CMP116} are used unchanged. Moreover, by
the structural statement on these displays used with that convention, the displays
\cite[(2.27)--(2.41)]{B-CMP116} read a term only through \cite[(2.26)]{B-CMP116} and deliver the pair
consisting of the constant $E^{*}$ produced by them and the rate $\kappa_{out}$.

In this proof only, $x$ denotes the tilt parameter of the flat family (not a coupling letter of the
glossary): $\mathfrak{K}(\cdot\,;x)$ is the member at tilt $x$ and $\mathfrak{K}(\cdot\,;0)$ the
untilted member. The tilt changes the factor indexed by $P$ in the sums over pairs $(Y_0,P)$, and the
filing $\mathcal{S}$. It never changes the steps at which the
fat-core cluster-expansion proposition produces the exponent. That proposition does not mention the
flat family, and need not, because its proof concerns only those steps.

Therefore the following carry $\kappa_{out}$ on the run used for the uniqueness theorem:
\begin{itemize}
\item the members $\mathfrak{K}(X,\mathcal{S})$;
\item the differences $\mathbf{B}^{sl}=\mathfrak{K}(\cdot\,;x)-\mathfrak{K}(\cdot\,;0)$;
\item the terms that do not contain the field variable $\psi$ of the flat family's expansion, which are
reassigned by
\cite[(2.36)]{B-CMP116}, with the remnant quadratic forms coupling the variables of pairs of cubes,
left by that reassignment, carrying the factor $e^{-\delta_1M\check{R}_{k+1}}$, where
$\check{R}_{k+1}$ is Ba{\l}aban's radius at level $k+1$.
\end{itemize}
This proves (vii).

The inputs of this proof are:
\begin{itemize}
\item the fat-core cluster-expansion proposition, together with definition \eqref{9.1:eq-fresh-rate-1}
of $\delta_\star^{(\mathfrak{c})}$;
\item the cluster-expansion lemma, the proposition on the marked sub-sum and the Schwarz-factor lemma
of the correlation theorem;
\item the statements of \cite{B-CMP116} and \cite{B-CMP119} quoted above;
\item clause (c) of the flat-family bound;
\item the convention fixing $\kappa_{out}$ for the flat family, and the stated property of
\cite[(2.27)--(2.41)]{B-CMP116}.
\end{itemize}
None of these uses the present lemma.
\end{proof}

\begin{remark}
We record four consequences and one caution.
\begin{itemize}
\item[(1)] In \cite{B-CMP119} the same bound is required of all the families expanded at the site.
\item[(2)] After the $\mathbf{R}$-step, the stored norm is the norm at level $k+1$ of the correlation
theorem, at a rate $\kappa\le\kappa_{out}$. The bounds that use these terms later are monotone in the
rate.
\item[(3)] In clause (c) of the flat-family bound, ``$\kappa_{out}$ the rate of the untilted
expansion'' does not depend on the route along which the terms are reached.
\item[(4)] At the small-field site, the sums in $\delta\kappa$ are controlled because the hull of every
stored input lies inside the fat core. The inclusion $(X^+)_0\subset X_0^{\sim}$ shows that the fat
core $X_0^{\sim}$ contains the hull $K(X)=(X^+)_0$ of Notation~\ref{9.1:not-hull-conventions}, so that
one layer serves both the relative bound and the fat core at each small-field site.
The factor $c_g$ at the steps where hulls are compared is
compensated by the rescaled rate $\kappa/c_g$ and by the rescaling by $L$. At the
preliminary-integration site of the step $k\to k+1$, the stored inputs are collared and the derived
chain terms $\mathbb{C}^{(i)}$ lie on Ba{\l}aban's cores; there, in the estimate that prices the
collared route at that site, the hull of a term created at the same level costs $(a_k+1)c_H$ out of
the rate $(1+4\beta)c^{\sharp}\kappa$, with $a_k$ the level-$k$ coefficient of the hull cost and
$\beta$ the parameter of that pricing estimate; this rate is supplied by the case $\mathfrak{c}=c_H$
of the lemma. The pricing of the hull parameter used along the correlation theorem does not enter at
that site. The case $j=k$ of the
bound on terms created at an
earlier level $j$ is a corollary of the lemma.
\end{itemize}
The parameter $\delta$ does not enter lower bounds only. The estimate \cite[(2.27)]{B-CMP116} at level
$k+1$ costs the exponential of five times the rate kept there, and that rate is at least
$\kappa_{out}$; in \cite{B-CMP116} this factor is printed as $\exp 5\kappa$, at the rate $\kappa$.
Accordingly, the smallness conditions of the cluster expansion that contain $e^{-5\kappa}$ become two
conditions: the constant of the corresponding smallness condition of the fat-core cluster-expansion
proposition, multiplied by $e^{4L\kappa/c_g}$, is at most one; and a condition on
$\alpha_6^{-1}e^{2L\kappa/c_g}$, with $\alpha_6$ Ba{\l}aban's constant of \cite[\S2]{B-CMP116}. These
are upper bounds on $\varepsilon_1$, chosen after $(L,\kappa,\kappa_1,M)$, entered as item (o5) of
this chapter's list of floors and ceilings chosen after $L$.
\end{remark}

Throughout what follows, $(\mathfrak{c},\mathfrak{h}_0)$ is one of the two pairs $(c_H,1)$ and $(c_H^3,3)$, with $c_H$, $c^{\sharp}=c_H+1$ and $a_\flat$ as in Notation~\ref{9.1:not-hull-conventions}. The letters $\delta_\star$ and $\kappa_{out}$ are those of Lemma~\ref{9.1:lem-fresh-rate}, and we put $\delta:=\delta_\star$, the value of the parameter $\delta$ at which that lemma runs the step.

Fix the current level $k$ and a level $j\le k$. Let $X\in\mathbf{D}_j$ be a region entering the large-field site (L) of the sourced expansion $J^{\mathfrak{S},+}$ used in the proof of the correlation theorem. Let $X_0$ and $X^{+}$ be its filing and its collar, and let $K(X)=(X^{+})_0$ be its hull. We write $K^{fil}(X)$ for the core of $X$ at the second weight (L2) of that site. It is $K^{fil}(X)=K(X)$ when $(\mathfrak{c},\mathfrak{h}_0)=(c_H,1)$, and $K^{fil}(X)=K(X)^{+2}$ when $(\mathfrak{c},\mathfrak{h}_0)=(c_H^3,3)$. We assume throughout that these cores satisfy
\begin{equation}\label{9.1:eq-bracket-margins-a}
\begin{gathered}
X_0\subset K^{fil}(X),\qquad
d_k\bigl(K^{fil}(X)\bigr)\le\mathfrak{c}\,\bigl(d_k(X_0)+\mathfrak{h}_0\bigr),\\
d_k(X_0)\le L^{-(k-j)}\,d_j(X).
\end{gathered}
\end{equation}

For the two pairs this holds for the following reasons. First, $K(X)\subset X_0^{+}$, because a $\pi_j$-cube (a block cube of level $j$) adjacent to $X$ lies in, or is $\infty$-adjacent to, one of the $M$-cubes making up $X_0$. The hull comparison that defines $c_H$ in Notation~\ref{9.1:not-hull-conventions} states that $Y\subset Y'\subset Y^{+}$ implies $d(Y')\le c_H(d(Y)+1)$. Applied once, it gives the middle inequality for the pair $(c_H,1)$. Applied three times, it gives the bound $c_H^3d_k(X_0)+c_H^3+c_H^2+c_H$. Since $c_H\ge1$, we have $c_H^3+c_H^2+c_H\le 3c_H^3$, and this is the pair $(c_H^3,3)$. The last inequality is the comparison of lengths across levels used in the proof of the correlation theorem and in \cite{B-CMP109}.

\begin{lemma}[Brackets at the large-field site keep their margins]
\label{9.1:lem-bracket-margins}
Assume \eqref{9.1:eq-bracket-margins-a}. Assume also the following conditions on the constants:
\begin{itemize}
\item $(F_L)$: $L\ge 3\mathfrak{c}$;
\item $(F_\star)$ of \eqref{9.1:eq-hull-conventions-a};
\item $(F_\kappa)$ of \eqref{9.1:eq-hull-conventions-a}: $\kappa\ge(\mathfrak{c}+1)C(d)/\delta_\star$, where $C(d)$ is the dimensional constant of the hull summation estimate for the fat core;
\item $\kappa_1\ge 4L\kappa+1$;
\item $\kappa\ge 10\log L$.
\end{itemize}
Let the inputs at the site (L) carry the following rates:
\begin{itemize}
\item inputs at level $k$ carry the rate $\kappa_{out}\ge 2(\mathfrak{c}+1)\kappa$ of Lemma~\ref{9.1:lem-fresh-rate};
\item inputs at levels $j<k$ carry the rate $\kappa$;
\item chain terms $\mathbb{C}^{(i)}$ at their top level carry the rate $(1+3\beta_{pr})\kappa$;
\item chain pieces filed below their top level, at a level $\ell<k$, carry the rate $a_\flat$.
\end{itemize}
Then the following statements hold at the site (L) of $J^{\mathfrak{S},+}$:
\begin{itemize}
\item the boundary lemma of the proof of the correlation theorem;
\item the first, second, third and fifth bounds, and the boundary bound, of the list of bounds at that site;
\item the bracket resummation of clause (c) of the resummation proposition of the same proof.
\end{itemize}
They hold with the same margins, namely
\begin{equation}\label{9.1:eq-bracket-margins-1}
\begin{gathered}
\bigl(\tfrac{5}{12}-\tfrac14\delta\bigr)\kappa d_j,\quad
\tfrac{17}{60}\,\kappa d_j\quad(j<k),\qquad
\bigl(\tfrac14-\tfrac14\delta\bigr)\kappa d_k\quad(j=k),\\
\tfrac12\delta\kappa d_j\ \text{ under the decoupling weight at (L1)}.
\end{gathered}
\end{equation}
The constants $C(\kappa)$ and $C_\partial'$ and the bounded part of $C_\sharp$ of the proof of the correlation theorem, and the far-unit constant of the remark on far units in that proof, are each multiplied by
\begin{equation}\label{9.1:eq-bracket-margins-2}
A_H=e^{(7/4)\mathfrak{c}(4+3^d)\kappa}.
\end{equation}
This enlargement is met by the condition $\mathsf{w}_L^{-1}\le\theta'$. Along the run of the uniqueness theorem, $\mathsf{w}_L$ is the radius of the dilation circle of the bracket at the large-field site. In the corresponding entry of condition (o8) in \eqref{9.1:eq-constant-order-a}, $K_o^L$ is the circle constant of the bracket at the large-field site, $K_\Lambda'$ is the frame constant, $\mathsf{r}_1=C_1/C_0$, and $\delta_1$ and $N_0$ are the decay rate and the integer fixed in that entry; the entry reads
\begin{equation}\label{9.1:eq-bracket-margins-3}
K_o^L\theta',\qquad K_o^LK_\Lambda'(1+\mathsf{r}_1)\,e^{-\delta_1M(N_0-2)/2}\le\theta'.
\end{equation}
\end{lemma}

\begin{proof}
\emph{The weight at (L2).} The weight \cite[(1.66)]{B-CMP122b} requires of the reading unit a factor $e^{(1+3\beta_s)\kappa d_k(Y_0)}$. At the site (L) the domain $Y_0$ is replaced by $K^{fil}(X)$. The coefficient $1+3\beta_s$ is at most $7/4$, since $\beta_s\le\frac14$. Hence, by \eqref{9.1:eq-bracket-margins-a}, the required factor is at most $A_1e^{(7/4)\mathfrak{c}\kappa d_k(X_0)}$, where
\begin{equation*}
A_1:=e^{(7/4)\mathfrak{c}\mathfrak{h}_0\kappa}\le A_H,
\end{equation*}
since $\mathfrak{h}_0\le 3<4+3^d$.

Case $j<k$. By the last inequality of \eqref{9.1:eq-bracket-margins-a} and by $\mathfrak{c}/L\le\frac13$, which is $(F_L)$,
\begin{equation*}
\tfrac74\,\mathfrak{c}\kappa d_k(X_0)
\le\tfrac74\,\mathfrak{c}L^{-(k-j)}\kappa d_j(X)
\le\tfrac{7}{12}\,\kappa d_j(X).
\end{equation*}
The arithmetic of the proof of the correlation theorem at this point is $\frac74L^{-(k-j)}\le\frac{7}{12}$. It therefore holds unchanged, with the condition $L\ge3$ replaced by $L\ge3\mathfrak{c}$. The same is true of that proof's tail estimates, which are indexed by an integer $\nu\ge2$ and in which the ratio $t=\mathfrak{c}/L$ enters in place of $1/L$, with $t\le\frac13$ by $(F_L)$:
\begin{itemize}
\item the inequality $1-\frac74t-\frac16\ge\frac14$ in the case $\nu=2$ holds for $t\le\frac13$;
\item the bound $\frac{7}{36}$ in the cases $\nu\ge3$ now takes the form $\frac74\mathfrak{c}/L^2\le\frac{7}{36}$, which holds because $L^2\ge 9\mathfrak{c}^2\ge 9\mathfrak{c}$.
\end{itemize}
Since $1-\frac{7}{12}=\frac{5}{12}$, the margins \eqref{9.1:eq-bracket-margins-1} for $j<k$ are those of the correlation theorem.

Chain pieces below their top level. Let $X\in\mathbb{D}_\ell$ be a chain piece filed below its top level, at a level $\ell<k$, carrying the rate $a_\flat$. We use once the inequality $d_\ell(X)\ge L^{k-\ell}d_k(X_0)$ of \cite{B-CMP109}. Since $k-\ell\ge1$, it gives
\begin{equation*}
a_\flat d_\ell(X)\ge\tfrac12a_\flat d_\ell(X)+\tfrac12a_\flat L\,d_k(X_0).
\end{equation*}
With $a_\flat=(1-\beta_{pr})\kappa/c^{\sharp}=\frac45\kappa/c^{\sharp}$, that is $1-\beta_{pr}=\frac45$, we have
\begin{equation*}
\tfrac12a_\flat L\ge\tfrac74\,\mathfrak{c}\kappa
\iff L\ge\tfrac{35}{8}\,\mathfrak{c}c^{\sharp}.
\end{equation*}
The right-hand inequality is implied by $(F_\star)$, which gives $L\ge 8(\mathfrak{c}+1)c_H^2$. Indeed $c_H\ge2$, by the definition $c_H=1+\sqrt d\,3^dC_d$ in Notation~\ref{9.1:not-hull-conventions}, so $c_H^2\ge c_H+1=c^{\sharp}$ and $8(\mathfrak{c}+1)\ge\frac{35}{8}\mathfrak{c}$; hence
\begin{equation*}
8(\mathfrak{c}+1)c_H^2\ge\tfrac{35}{8}\,\mathfrak{c}(c_H+1).
\end{equation*}
Thus the second summand pays the factor $e^{(7/4)\mathfrak{c}\kappa d_k(X_0)}$.

The remainder $\frac12a_\flat d_\ell(X)$ replaces $\frac12\delta\kappa d_j$ in the sums over $X$ and in the anchoring of the decoupling lemma and of the bracket resummation. These steps of the proof of the correlation theorem use only some positive rate under a floor on $\kappa$. With the rate $\frac12a_\flat$, the constant $C(\kappa)$ becomes $C(a_\flat/4)$, and the floor becomes $\kappa\ge C(d)c^{\sharp}$, which has the shape of $(F_\kappa)$. This is the bound (d3) of the payment lemma for chain pieces, at general $\mathfrak{c}$.

Equivalently, $(F_\star)$ gives
\begin{equation*}
\tfrac74\,\mathfrak{c}/L<\tfrac{7}{32c_H^2}\le\tfrac{2}{5c^{\sharp}},
\end{equation*}
where the second inequality is $35(c_H+1)\le 64c_H^2$ and holds for $c_H\ge2$. Together with $d_k(X_0)\le L^{-(k-\ell)}d_\ell(X)$ this leaves the margin $\tfrac{2}{5c^{\sharp}}\kappa d_\ell(X)$.

Chain pieces at their top level carry the rate $(1+3\beta_{pr})\kappa\ge\kappa$, so the arithmetic of the case $j<k$ applies to them without change.

Case $j=k$. Here $X_0=X$. The supply is $2(\mathfrak{c}+1)\kappa d_k$, coming from $\kappa_{out}\ge2(\mathfrak{c}+1)\kappa$, and the demand is $\frac74\mathfrak{c}\kappa(d_k+\mathfrak{h}_0)$. The margin is therefore
\begin{equation*}
\bigl(2(\mathfrak{c}+1)-\tfrac74\mathfrak{c}\bigr)\kappa d_k
=\bigl(2+\tfrac14\mathfrak{c}\bigr)\kappa d_k\ge\tfrac14\kappa d_k
\ge\bigl(\tfrac14-\tfrac14\delta\bigr)\kappa d_k,
\end{equation*}
and the constant is $A_1\le A_H$.

The chain terms $\mathbb{C}^{(n)}$ of this step, at rate $(8/5)\kappa$, are the pieces at level $k$ only. They are formed with Ba{\l}aban's cores and enter through the interface of the proof of the correlation theorem itself, so the hulls $K^{fil}$ do not enter them. The pieces of this step at levels $m<k$ carry the rate $a_\flat$. For these, the form of the weight \cite[(1.66)]{B-CMP122b} is the bound (d2) of the payment lemma for chain pieces. Its condition is $(1-\frac15)^2L\ge\frac74c^{\sharp}$, and it is implied by $(F_\star)$.

\emph{The weight at (L1).} The weight $\mathsf{W}$ of the decoupling lemma reads $e^{(1-\frac12\delta)\kappa d_k(\bar S_v)}$, where the $S_v$ are the decoupled sets of that lemma and the $\bar S_v$ their closures. At the site (L) the set $\bar S_v$ is replaced by the hull $K(S_v)$.

Case $j<k$. By \eqref{9.1:eq-bracket-margins-a} with the pair $(c_H,1)$, and by $L\ge 3\mathfrak{c}\ge3c_H$,
\begin{equation*}
e^{(1-\frac12\delta)\kappa d_k(K(S_v))}\le e^{c_H\kappa}\,e^{\frac13\kappa d_j}.
\end{equation*}
This leaves at least $\frac12\delta\kappa d_j$.

Case $j=k$. Since $\mathfrak{c}\ge c_H$ and $1-\frac12\delta\le1$,
\begin{equation*}
2(\mathfrak{c}+1)\kappa d_k-\bigl(1-\tfrac12\delta\bigr)c_H\kappa(d_k+1)
\ge(c_H+2)\kappa d_k-c_H\kappa.
\end{equation*}
In both cases the constant is $e^{c_H\kappa}\le A_H$. The decay factor $\eta_p$ of the decoupling lemma, taken to the first decoupled cube, can only improve, as in the fifth step of the fat-core construction.

\emph{Counting.} Fix a hull $K$. By the hull comparison of Notation~\ref{9.1:not-hull-conventions}, the number of regions $X$ with $K^{fil}(X)=K$ is at most $2^{N(K)}$, where $N(K)$ is the number of $j$-cubes in $K$. In the hull summation estimate for the fat core, the last factor sums the at most $2^{N(K)}$ domains $Y_0$ with a given hull $K$, and this estimate holds under $(F_\kappa)$.

For the remark on far units in the proof of the correlation theorem: a far unit consists of at most $3^d$ cubes of the current level, so by \eqref{9.1:eq-bracket-margins-a} its hull satisfies $d_k\le\mathfrak{c}(3^d+\mathfrak{h}_0)$. The factor $e^{3^d\kappa}$ of that remark is therefore replaced by a factor at most $A_H$.

\emph{The term $T(Y_0)$ at (L2).} The term $T(Y_0)$ attached to the domain $Y_0$ is formed with Ba{\l}aban's core, and its exponent is that of the decoupling lemma. That exponent consists of the factor $e^{-a_sd_j(X)}$ of the region, with $a_s$ the rate of the region in the decoupling lemma, and a factor $e^{-(\kappa_1-1)}$ for each decoupled cube. The weight at (L2) asks for $(1+3\beta_s)\kappa d_k(Y_0)$, and, by the first clause of the fat-core comparison,
\begin{equation*}
d_k(Y_0)\le d_k(K(X))+\sqrt d\,m(p),
\end{equation*}
where $m(p)$ is the number of decoupled cubes of the term. The part carried by the region is paid as in the treatment of the weight at (L2) above. Each decoupled cube pays $\frac74\sqrt d\,\kappa\le\frac12(\kappa_1-1)$, which follows from $\kappa_1\ge4L\kappa+1$ because $\frac12(\kappa_1-1)\ge2L\kappa$.

No other bound in the proof of the correlation theorem reads a core. The marking in clause (iii) of the boundary lemma and the cores in its clause (i) are on the side of $X$. For far points $y$ of the large-field region $Z$ one common enlarged localisation region $Y_1^{+}=C$ is kept. This is legitimate because $X\subset\square^{\sim\nu}$, the enlarged cube containing $X$, implies
\begin{equation*}
X^{+}\subset\square^{\sim\nu+1}\subset\operatorname{int}C,
\end{equation*}
and hence $K(X)\subset C$.

\emph{The majorant of the box in the bracket resummation.} Along the run of the uniqueness theorem, the majorant of the box of the modulus lemma in the bracket resummation is the boundary form of the modulus lemma that accompanies the hull summation estimate for the fat core, applied to the first-order members on the collar $X^{+}$; these members are controlled by the transfer bound for the regions of the expansion.
\end{proof}

\begin{definition}[Cubes, tree lengths and the window of strip widths]\label{9.1:def-strip-window}
Fix a level $k$, $0\le k\le K$.

\emph{Cubes, distance, layers.} In this section a \emph{cube} is a closed cube of side $M$ of
the partition $\pi_k$ of the level-$k$ lattice into cubes of side $M$ used in
\cite{B-CMP122b}, and all lengths are measured in units of $M$. For unions $A$, $B$ of
cubes, $\operatorname{dist}(A,B)$ is their distance in the maximum norm; it is an integer,
and $\operatorname{dist}(A,B)=0$ exactly when $A$ and $B$ touch or overlap. $N_k(A)$ is the
number of cubes in $A$; it is always written with its argument and is not the $N_k$ of
Notation~\ref{0.2:not-glossary}. Unions of cubes are compared as families of cubes: ``$A$ meets $B$''
means that $A$ and $B$ have a cube in common. We put $R:=\check{R}_k$, where $\check{R}_k$ is
the integer $R_k$ of \cite{B-CMP122b} for which the large-field regions at level $k$ are
unions of cubes of side $M\check{R}_k$; throughout this section $R$ abbreviates $\check{R}_k$
and is not the ball radius of Definition~\ref{1.3:def-curvature-field}. For a union of cubes
$A$ and an integer $m\ge 1$ ($m$ here counts layers and is not the torus index), $A^{(m)}$ is
the union of the
cubes $\Delta$ with $\operatorname{dist}(\Delta,A)\le m-1$, and $A^{(0)}:=A$; thus $A^{(m)}$
is $A$ with $m$ layers of cubes added.

\emph{The layer index.} Let $Z$ be the union of components described in condition (G0) of
\eqref{9.1:eq-strip-window-a} below. For a
cube $\Delta$ put $\operatorname{lay}(\Delta):=0$ if $\Delta\subset Z$, and
$\operatorname{lay}(\Delta):=1+\operatorname{dist}(\Delta,Z)$ otherwise. Touching cubes
$\Delta$, $\Delta'$ satisfy
$\lvert\operatorname{lay}(\Delta)-\operatorname{lay}(\Delta')\rvert\le 1$. Indeed,
$\Delta=\Delta'+e$ with $e\in\{-1,0,1\}^d$; if $y\in\Delta'$ and $z\in Z$ realise
$\operatorname{dist}(\Delta',Z)$, then $y+e\in\Delta$ and
$\lvert y+e-z\rvert_\infty\le\lvert y-z\rvert_\infty+1$, so that
$\operatorname{dist}(\Delta,Z)\le\operatorname{dist}(\Delta',Z)+1$, and symmetrically.
If neither cube lies in $Z$, this is the assertion. If $\Delta\subset Z$, then
$\operatorname{dist}(\Delta',Z)=0$, so $\operatorname{lay}(\Delta')\le 1$; the case
$\Delta'\subset Z$ is symmetric.

\emph{Domains and tree lengths.} $\mathbf{D}_k$ is the class of localisation domains at
level $k$, and $d_k(\cdot)$ is the tree length, both in the sense of \cite{B-CMP122b}. For
unions of cubes $W$ and $Y$, $d_W(Y)$ is the least length $\lvert T\rvert$ of
a tree graph $T$ contained in $Y$ that intersects every cube of $Y\setminus W$, and
$d_W(Y):=0$ if $Y\subset W$. For $Y\in\mathbf{D}_k$ and $W=Z$ this is the relative tree length
$d_{k,Z}(Y)$ of \cite{B-CMP122b}.

\emph{The large-field geometry} (G0). Let $\Lambda_k$ be the small-field region at level
$k$, and let $\Lambda_k^c$ be its complement. The \emph{components} of $\Lambda_k^c$ are the
classes of its cubes under the equivalence relation generated by ``intersect, or touch''.
Touching includes contact along a face of any dimension; in particular, two cubes meeting
only at a corner lie in the same component. $Z$ is a union of components $C$ of
$\Lambda_k^c$. Each component of $\Lambda_k^c$ is a box with sides $a_1R,\dots,a_dR$, where
$a_i\in\{1,\dots,100\}$. Distinct components of $\Lambda_k^c$ do not touch, and they are at
distance at least $R$ from each other.

\emph{The window and the strips.} Let $m_\sharp\in\mathbb{N}$ lie in the window displayed
below. Put $Z^{\sharp}:=Z^{(m_\sharp)}$, and for each component $C$ of $Z$ let
$C^{\sharp}:=C^{(m_\sharp)}$ be the \emph{strip} of $C$. Put $d_\sharp:=d_{Z^\sharp}$.
The filing map $\operatorname{fil}_{\sharp}$ adds to $Y_4\in\mathbf{D}_k$ every strip that
meets $Y_4$. A domain $Y_4\in\mathbf{D}_k$ is \emph{interior} if $Y_4\subset Z^{\sharp}$,
and \emph{exterior} otherwise. In formulas,
\begin{equation}\label{9.1:eq-strip-window-a}
\begin{aligned}
&\text{(G0)}\quad Z \text{ is a union of components } C \text{ of } \Lambda_k^c,\ \text{the
classes of cubes of } \Lambda_k^c\\
&\phantom{\text{(G0)}\quad}\text{under ``intersect, or touch'', corners included};\\
&\phantom{\text{(G0)}\quad}\text{each component } C \text{ of } \Lambda_k^c \text{ is a box with
sides } a_1R,\dots,a_dR,\ a_i\in\{1,\dots,100\},\\
&\phantom{\text{(G0)}\quad}\operatorname{dist}(C,C')\ge R
\ \text{ for distinct components } C\ne C' \text{ of } \Lambda_k^c;\\
&1\le m_\sharp\le m_{hi}(k):=\Bigl\lfloor\tfrac12(R-1)\Bigr\rfloor,\qquad
Z^{\sharp}:=Z^{(m_\sharp)},\quad C^{\sharp}:=C^{(m_\sharp)},\quad d_\sharp:=d_{Z^\sharp};\\
&\operatorname{fil}_{\sharp}(Y_4):=Y_4\cup\bigcup\bigl\{C^{\sharp}:\ C^{\sharp}
\text{ meets } Y_4\bigr\};\qquad
Y_4 \text{ interior}\iff Y_4\subset Z^{\sharp};\\
&d_W(Y):=\begin{cases}
\min\bigl\{\lvert T\rvert:\ T\subset Y \text{ a tree graph meeting every cube of }
Y\setminus W\bigr\}, & Y\not\subset W,\\
0, & Y\subset W.
\end{cases}
\end{aligned}
\end{equation}

\emph{Filed domain classes.} The two classes of filed domains are
\begin{equation*}
\mathbf{D}_k^{Z}:=\bigl\{Y_4\in\mathbf{D}_k:\ Y_4\cap Z \text{ is a non-empty union of
components of } Z\bigr\},
\end{equation*}
and $\mathbf{D}_k^{\sharp}$, defined in the same way with $Z^{\sharp}$ in place of $Z$ and
the strips $C^{\sharp}$ in place of the components $C$:
\begin{equation*}
\mathbf{D}_k^{\sharp}:=\bigl\{Y_4\in\mathbf{D}_k:\ Y_4\cap Z^{\sharp} \text{ is a
non-empty union of strips } C^{\sharp}\bigr\}.
\end{equation*}

\emph{Combinatorial constants.} We put
\begin{equation*}
\mathsf{c}_d:=2\cdot 3^d,\qquad
\mathsf{k}_d:=1+\mathsf{c}_d\bigl(2d\log 3+(1+3^d)\log 2\bigr),\qquad
C_{fib}:=2e^{2\mathsf{k}_d}.
\end{equation*}

\emph{The strip width used below.} The lemma and the theorem on strips proved below in this
section hold for every $m_\sharp$ in the window \eqref{9.1:eq-strip-window-a}; in every later
use of the strips, $m_\sharp$ is taken equal to $m_{hi}(k)$.
\end{definition}

\begin{remark}[The collar of the large-field region]
The collar $Z^{\sim}$ of $Z$ in \cite{B-CMP122b} is used here as $Z^{\boxplus}$, the set $Z$
with one layer of cubes of side $MR$ added; the reason is the following. In
\cite{B-CMP122b} the operation $\sim$ is fixed by the cubes ``in terms of which the sets are
defined'', and $Z$ is a union of cubes of side $MR$; a collar by cubes of side $M$ is used at
one place, where the word ``here'' marks it as a local convention. Further, the terms of the
second group in \cite[(1.100)]{B-CMP122b} are localised in the set $\Lambda_k^{\sim-1}$ of
\cite[(2.30)]{B-CMP119}, which is obtained from $\Lambda_k$ by removing one layer of cubes
of side $MR$; hence the operation $\sim$ in the collar of \cite{B-CMP122b} adds one layer of
cubes of side $MR$. Carried out on this collar, the mechanism of \cite{B-CMP122b} gives the
lower bound $d_{k,Z}(Y)\ge R-2$ for the relative tree length of every domain $Y$ to which
that mechanism is applied in \cite{B-CMP122b}, and hence a bound, by a constant multiple of
$e^{-\beta\kappa(R-2)}$, on the factor that \cite{B-CMP122b} attaches to such a domain through
its relative tree length; here $\beta$ and $\kappa$ are the constants of \cite{B-CMP122b} in
the exponent of that factor.
A collar by cubes of side $M$ also occurs in \cite{B-CMP122b}, in the notation
$\square^{\sim n}$, which is accompanied there by the sentence ``The meaning of the symbol
$X^n$ should be obvious''. The strips $C^{\sharp}$, the filing map $\operatorname{fil}_{\sharp}$
and the classes of Definition~\ref{9.1:def-strip-window} are defined without reference to
$Z^{\boxplus}$, and the theorem on strips below does not depend on whether the collar of $Z$
is taken to be $Z^{\boxplus}$ or $Z^{(1)}$.
\end{remark}

\begin{remark}[Three test configurations]
The following three configurations concern only the bounds that \cite{B-CMP122b} states for
the terms involved: they show what these bounds imply, and no lower bound on any individual
term is asserted. Lengths are in units of
$M$. Let $C$ be a component of $Z$, let $\Delta_0\subset C$ be a cube at a face of $C$, and
let $\Delta_1$, $\Delta_2$ be the next two cubes along the outward normal to that face.

(a) \emph{The collar $Z^{(1)}$.} Here the collar of $Z$ is taken to be $Z^{(1)}$, and the set
$Y_1$ that \cite{B-CMP122b} associates with a domain is taken without a hull. Let
$X:=\Delta_0\cup\Delta_1\cup\Delta_2$, created at step $j=k$, carry a term of the family
$\mathbf{B}^{(k)}$ or of the family $E^{(k)}$ of \cite{B-CMP122b}. Then $d_k(X)=1$, and the
bound stated in \cite{B-CMP122b} for the term is $B_0e^{-(1+4\beta)\kappa}$, respectively
$E_0e^{-(1+4\beta)\kappa}$, with the constants $B_0$, $E_0$ and $\beta$ of
\cite{B-CMP122b}. The associated set is $Y_1=C\cup\Delta_1\cup\Delta_2$; it is not
contained in $Z^{(1)}$, and $d_{k,Z}(Y_1)=0$, since a single point of
$\Delta_1\cap\Delta_2$ is a tree graph of length zero in $Y_1$ meeting the cubes $\Delta_1$,
$\Delta_2$ of $Y_1\setminus Z$. Under this reading, \cite[(1.68)]{B-CMP122b} and
\cite[(1.93)]{B-CMP122b}, taken together at this term, call for a bound by
$e^{-6d(1+\beta)\kappa}$. Since the exponent $(1+4\beta)\kappa$ of the bound stated for the
term is smaller than $6d(1+\beta)\kappa$, under this reading the bound by
$e^{-6d(1+\beta)\kappa}$ is not a consequence of the stated bound; this concerns the set
$Y_1$ itself, taken without a hull, that is, before the hull convention of
\eqref{9.1:eq-hull-conventions-a} is applied.

(b) \emph{Strip width.} Let $X:=\Delta_0\cup\Delta_1$. Its hull $K(X)$ of
\eqref{9.1:eq-hull-conventions-a} contains $\Delta_2$. With the collar $Z^{(1)}$, that is,
relative to $Z^{(1)}=Z^{\sharp}$ at $m_\sharp=1$, the domain $K(X)$ is exterior; it is
interior, that is
$K(X)\subset Z^{\sharp}$, as soon as $m_\sharp\ge 2$.

(c) \emph{The collar $Z^{\boxplus}$.} Condition (G0) of \eqref{9.1:eq-strip-window-a}
admits two components $C_1$, $C_2$ of $Z$ at distance exactly $R$; the conditions that
\cite{B-CMP122b} places on large-field regions exclude only that distinct components touch.
Such a pair arises, for instance, from two boxes of cubes of
side $MR$ at distance $2a+1$ in units of $MR$, for some integer $a\ge 0$, each then enlarged
by $a$ layers of cubes of side $MR$. Since $C_1$ and
$C_2$ are unions of cubes of side $MR$ of one partition, some cube of side $MR$ touches both
of them; in the configuration just described, $C_1^{\boxplus}\cap C_2^{\boxplus}$ contains
the whole layer of cubes between them. For such a pair $C_1^{\boxplus}$ and
$C_2^{\boxplus}$ therefore have a cube of side $MR$ in common, and $Z^{\boxplus}$ has a
single component containing $C_1$ and $C_2$. A term whose domain contains $C_1\cup C_2$ and
is contained in $C_1^{\boxplus}\cup C_2^{\boxplus}$ is summed, by \cite{B-CMP122b}, into the
component $C_1^{\boxplus}\cup C_2^{\boxplus}$ of $Z^{\boxplus}$; this lies within the scope
of \cite[(1.69)]{B-CMP122b}. In \cite[(1.71)]{B-CMP122b} the correspondence
$X\leftrightarrow X^{\sim}$ pairs the components of $Z$ with those of its collar. Read with
the collar $Z^{\boxplus}$, this correspondence is not one-to-one for such a pair, by what was
just shown; for the same reason,
read with the collar $Z^{\boxplus}$, the expressions ``$V(X^{\sim})$'', ``factorize in
components of $Z$'' and ``equal to one of the components of $Z^{\sim}$'' of
\cite{B-CMP122b}, and the step of \cite{B-CMP122b} that uses the correspondence, have no
single component of the collar to refer to for each of $C_1$ and $C_2$. For a chain of
components of $Z$, each at distance exactly $R$ from the next, the argument above, applied
to each consecutive pair, shows that the collars of all members of the chain lie in one
component of $Z^{\boxplus}$; configurations containing such chains lie outside the first of
the classes of large-field configurations distinguished in the domain convention of this
chapter.
\end{remark}

\begin{lemma}[Geometry of the strips]\label{9.1:lem-strip-geometry}
Work at level $k$ with the cubes, the distance $\operatorname{dist}$, the cube count $N(\cdot)$, the
radius $R$, the hulls $A^{(m)}$, the layer function $\operatorname{lay}$ and the tree lengths $d_W$ of
Definition~\ref{9.1:def-strip-window}, and with the large-field region $Z=\Lambda_k^c$, its
components $C$, the strip width $m_\sharp$, the strips $C^{\sharp}=C^{(m_\sharp)}$, their union
$Z^{\sharp}$, Ba{\l}aban's collar $C^{\boxplus}$ of $C$, the filing map
$\operatorname{fil}_\sharp$, the classes $\mathbf{D}_k$, $\mathbf{D}_k^Z$, $\mathbf{D}_k^{\sharp}$,
the constants $\mathsf{c}_d=2\cdot 3^d$, $\mathsf{k}_d$, $C_{fib}$ and the lengths $d_Z$ and
$d_\sharp=d_{Z^\sharp}$, as fixed in \eqref{9.1:eq-hull-conventions-a} and
\eqref{9.1:eq-strip-window-a}. Assume the hypothesis on the components of $\Lambda_k^c$ displayed in
\eqref{9.1:eq-strip-window-a} and that $m_\sharp$ lies in the window of strip widths fixed there.
Then the following hold.
\begin{itemize}
\item[\textup{(S1)}] The strips are boxes. For distinct components $C\neq C'$ of $\Lambda_k^c$,
\begin{equation}\label{9.1:eq-strip-geometry-a}
\operatorname{dist}(C^{\sharp},C'^{\sharp})\ge R-2m_\sharp\ge 1,
\qquad
\operatorname{dist}(C^{\sharp},C')\ge R-m_\sharp .
\end{equation}
The strips are the components of $Z^{\sharp}$, and each contains exactly one component of
$\Lambda_k^c$. Moreover $Z^{\sharp}\setminus Z\subset\Lambda_k$, $C^{\sharp}\subset C^{\boxplus}$, and
\begin{equation*}
N(C^{\sharp})\le\min\{2^dN(C),\,(101R-1)^d\}.
\end{equation*}
\item[\textup{(S2)}] If $W\subset W'$, $Y\subset Y'$ and $Y'\setminus W'\subset Y\setminus W$, then
$d_{W'}(Y')\le d_W(Y)$.
\item[\textup{(S3)}] The map $\operatorname{fil}_\sharp$ sends $\mathbf{D}_k^Z$ into
$\mathbf{D}_k^{\sharp}$. For $Y_4\in\mathbf{D}_k^Z$ one has
$\operatorname{fil}_\sharp(Y_4)\setminus Z^{\sharp}=Y_4\setminus Z^{\sharp}$ and
$d_\sharp(\operatorname{fil}_\sharp(Y_4))\le d_Z(Y_4)$; an interior $Y_4$ lies in a single strip
$C^{\sharp}$, and $C^{\sharp}=\operatorname{fil}_\sharp(Y_4)$.
\item[\textup{(S4)}] An exterior $Y_4\in\mathbf{D}_k^Z$ has $d_Z(Y_4)\ge m_\sharp-1$.
\item[\textup{(S5)}] (a) A tree graph of length $\ell$ intersects at most $\mathsf{c}_d(\ell+1)$
cubes. (b) Let $\mathsf{I}$ be a family of pairwise non-touching unions of cubes, called
\emph{islands}, let $W$ be the union of the islands, and let $\Delta_*$ be a cube. For every integer
$n$, the number of $Y\in\mathbf{D}_k$ such that $Y\cap W$ is a union of islands,
$\Delta_*\subset Y\setminus W$ and $d_W(Y)<n$ is at most $e^{\mathsf{k}_d(n+1)}$.
(c) If $\lambda_0\ge\mathsf{k}_d+1$, the sum of $e^{-\lambda_0 d_W(Y)}$ over all
$Y\in\mathbf{D}_k$ such that $Y\cap W$ is a union of islands and $\Delta_*\subset Y\setminus W$ is at
most $C_{fib}$.
\end{itemize}
\end{lemma}

\begin{proof}
\emph{Proof of \textup{(S1)}.} We first show that for unions of cubes $A$, $B$ and integers
$m,m'\ge 0$,
\begin{equation}\label{9.1:eq-strip-geometry-1}
\operatorname{dist}(A^{(m)},B^{(m')})\ge\operatorname{dist}(A,B)-m-m'.
\end{equation}
If $m\ge1$, a cube $\Delta$ of $A^{(m)}$ satisfies $\operatorname{dist}(\Delta,A)\le m-1$ and has
$\infty$-diameter $1$, so every point of $\Delta$ lies at $\infty$-distance at most $m$ from $A$;
for $m=0$ one has $A^{(0)}=A$ and every point of $\Delta$ lies in $A$. The same holds for $B^{(m')}$
with $m'$, and the triangle inequality for the
$\infty$-norm gives \eqref{9.1:eq-strip-geometry-1}.

By the hypothesis of \eqref{9.1:eq-strip-window-a}, distinct components $C\neq C'$ of $\Lambda_k^c$
are at distance at least $R$.
Applying \eqref{9.1:eq-strip-geometry-1} with $A=C$, $B=C'$, $m=m'=m_\sharp$ gives
$\operatorname{dist}(C^{\sharp},C'^{\sharp})\ge R-2m_\sharp$, and $R-2m_\sharp\ge1$ by the window.
Applying it with $m=m_\sharp$, $m'=0$ gives $\operatorname{dist}(C^{\sharp},C')\ge R-m_\sharp$,
and $R-m_\sharp\ge R-2m_\sharp\ge1$. This is \eqref{9.1:eq-strip-geometry-a}.

Write $a_1R,\dots,a_dR$ for the side lengths of the box $C$ in units of $M$, as fixed in
\eqref{9.1:eq-strip-window-a}, with integers $1\le a_i\le 100$. The cubes at distance at most
$m_\sharp-1$ from the box $C$ are those
of the box obtained by adding $m_\sharp$ layers of cubes on each face, so $C^{\sharp}$ is the box of
sides $a_iR+2m_\sharp$. Being a box, each strip is wall-connected (any two of its cubes are joined by
a wall-chain in it, that is a sequence of cubes in which consecutive cubes share a face of dimension
$d-1$); two distinct strips are at distance at least $1$, so no cube of one touches a cube of the
other. As $Z^{\sharp}$ is the union of the strips, the strips are exactly its components. Each strip
$C^{\sharp}$ contains $C$, and shares no cube with any other component $C'$, since
$\operatorname{dist}(C^{\sharp},C')\ge1$; so it contains exactly one component of $\Lambda_k^c$. A cube
of $Z^{\sharp}$ not in $Z=\Lambda_k^c$ is a cube of $\Lambda_k$, which is the inclusion
$Z^{\sharp}\setminus Z\subset\Lambda_k$.

For the sides, $2m_\sharp\le R-1$ by the window, and $1\le a_i\le100$ give
\begin{equation}\label{9.1:eq-strip-geometry-2}
a_iR+2m_\sharp\le(a_i+1)R-1<\min\{2a_i,101\}R .
\end{equation}
Hence
\begin{equation*}
N(C^{\sharp})=\prod_{i=1}^{d}(a_iR+2m_\sharp)\le\prod_{i=1}^{d}2a_iR=2^dN(C);
\end{equation*}
and since the sides
$a_iR+2m_\sharp$ are integers smaller than $101R$, $N(C^{\sharp})\le(101R-1)^d$. Finally
$m_\sharp<R$, because $R-2m_\sharp\ge1$, and therefore $C^{\sharp}\subset C^{\boxplus}$ by the
definition of the collar in \eqref{9.1:eq-hull-conventions-a}.

\emph{Proof of \textup{(S2)}.} A tree graph enters the minimum defining $d_W(Y)$ when it lies in $Y$
and intersects every cube of $Y\setminus W$. Let $\mathcal{T}$ be such a tree of least length
$d_W(Y)$. It lies in $Y\subset Y'$ and intersects every cube of $Y\setminus W$, hence every cube of
$Y'\setminus W'\subset Y\setminus W$. So $\mathcal{T}$ enters the minimum defining $d_{W'}(Y')$, and
$d_{W'}(Y')\le d_W(Y)$.

\emph{Proof of \textup{(S3)}.} Let $Y_4\in\mathbf{D}_k^Z$; $\operatorname{fil}_\sharp(Y_4)$ is $Y_4$
together with every strip that shares a cube with $Y_4$. Each added strip is wall-connected and shares
a cube with the wall-connected $Y_4$, so the union is wall-connected and
$\operatorname{fil}_\sharp(Y_4)\in\mathbf{D}_k$. A cube of $Y_4$ lying in $Z^{\sharp}$ lies in some
strip; that strip meets $Y_4$ and is therefore added; and distinct strips share no cube by
\textup{(S1)}. Hence $\operatorname{fil}_\sharp(Y_4)\cap Z^{\sharp}$ is the union of the added strips.
It is non-empty, since $Y_4$ contains some component $C$ of $\Lambda_k^c$ and then $C^{\sharp}$ is
added. Thus $\operatorname{fil}_\sharp(Y_4)$ is a member of $\mathbf{D}_k$ meeting $Z^{\sharp}$ in a
non-empty union of strips, that is, $\operatorname{fil}_\sharp(Y_4)\in\mathbf{D}_k^{\sharp}$. The added
strips lie in $Z^{\sharp}$, so nothing off $Z^{\sharp}$ is added, and
$\operatorname{fil}_\sharp(Y_4)\setminus Z^{\sharp}=Y_4\setminus Z^{\sharp}$.

We now apply \textup{(S2)} with $W=Z$, $W'=Z^{\sharp}$, $Y=Y_4$, $Y'=\operatorname{fil}_\sharp(Y_4)$.
Here $Z\subset Z^{\sharp}$ because $C\subset C^{\sharp}$ for every component $C$;
$Y_4\subset\operatorname{fil}_\sharp(Y_4)$ by construction; and
\begin{equation*}
\operatorname{fil}_\sharp(Y_4)\setminus Z^{\sharp}=Y_4\setminus Z^{\sharp}\subset Y_4\setminus Z .
\end{equation*}
Hence $d_\sharp(\operatorname{fil}_\sharp(Y_4))=d_{Z^\sharp}(\operatorname{fil}_\sharp(Y_4))\le d_Z(Y_4)$.

Let $Y_4$ be interior, so that its cubes lie in $Z^{\sharp}$. Since distinct strips do not touch and
$Y_4$ is wall-connected, all cubes of $Y_4$ lie in one strip $C^{\sharp}$. This strip is then the only
strip meeting $Y_4$, and $\operatorname{fil}_\sharp(Y_4)=Y_4\cup C^{\sharp}=C^{\sharp}$.

\emph{Proof of \textup{(S4)}.} Let $Y_4\in\mathbf{D}_k^Z$ be exterior. It contains a cube of $Z$ (it
contains a component of $Z$) and a cube $\Delta$ not in $Z^{\sharp}$. Such a $\Delta$ is not in
$C^{(m_\sharp)}$ for any component $C$, so $\operatorname{dist}(\Delta,C)\ge m_\sharp$ for every $C$,
hence $\operatorname{dist}(\Delta,Z)\ge m_\sharp$, and $\Delta\not\subset Z$; thus
$\operatorname{lay}(\Delta)\ge m_\sharp+1$. Take a wall-chain in $Y_4$ from a cube of $Z$, of layer $0$,
to $\Delta$. Consecutive cubes of the chain touch, so their layers differ by at most $1$; therefore the
chain passes cubes $\Delta_1,\Delta_2\subset Y_4\setminus Z$ with $\operatorname{lay}(\Delta_1)=1$
and $\operatorname{lay}(\Delta_2)=m_\sharp+1$.

Let $\mathcal{T}$ be any tree graph entering the minimum defining $d_Z(Y_4)$. It intersects every cube
of $Y_4\setminus Z$, so it contains points $p_1\in\Delta_1$ and $p_2\in\Delta_2$. Since
$\operatorname{dist}(\Delta_1,Z)=0$ and $\Delta_1$ has $\infty$-diameter $1$,
$\operatorname{dist}(p_1,Z)\le1$; since $\operatorname{dist}(\Delta_2,Z)=m_\sharp$,
$\operatorname{dist}(p_2,Z)\ge m_\sharp$. The path in $\mathcal{T}$ from $p_1$ to $p_2$ therefore has
length at least $|p_1-p_2|_\infty$, and
\begin{equation*}
|p_1-p_2|_\infty\ge\operatorname{dist}(p_2,Z)-\operatorname{dist}(p_1,Z)\ge m_\sharp-1 ;
\end{equation*}
the length of $\mathcal{T}$ is at least that of the path. Taking the minimum over $\mathcal{T}$
gives $d_Z(Y_4)\ge m_\sharp-1$.

\emph{Proof of \textup{(S5)(a)}.} Let the tree graph have length $\ell$. Traversing every edge once
in each direction gives a closed walk (an Euler tour) of length $2\ell$ whose image is the whole tree.
Cut it into at most $2\ell+1$ arcs of length at most $1$. An arc has $\infty$-diameter at most $1$, so
its projection on each coordinate axis is an interval of length at most $1$, which meets at most $3$ of
the closed unit intervals of the cube grid on that axis; hence an arc meets at most $3^d$ cubes. The
tree therefore meets at most $3^d(2\ell+1)\le2\cdot3^d(\ell+1)=\mathsf{c}_d(\ell+1)$ cubes.

\emph{Proof of \textup{(S5)(b)}.} If $n\le0$ there is no such $Y$, as $d_W\ge0$. Let $n\ge1$, put
$\nu:=\mathsf{c}_d(n+1)$, and let
$Y$ be counted. Choose a tree $\mathcal{T}$ of least length $d_W(Y)<n$ entering the minimum defining
$d_W(Y)$, and let $\hat T$ be the family of cubes it intersects. By (a),
\begin{equation*}
|\hat T|\le\mathsf{c}_d(d_W(Y)+1)\le\nu .
\end{equation*}
The family $\hat T$ contains $\Delta_*$,
since $\Delta_*\subset Y\setminus W$. It is touching-connected: if $\hat T$ split into two non-empty
subfamilies no cube of one of which touches a cube of the other, their unions would be two disjoint
closed sets, each meeting $\mathcal{T}$, covering the connected set $\mathcal{T}$, which is
impossible.

A touching-connected family of $m$ cubes containing $\Delta_*$ is determined by a depth-first walk from
$\Delta_*$ on a spanning tree of its touching graph; the walk has $2(m-1)$ steps, each to one of the
fewer than $3^d$ cubes touching the current one. Such families therefore number at most
$3^{2d(m-1)}\le3^{2dm}$, and summing over $1\le m\le\nu$,
\begin{equation*}
\sum_{1\le m\le\nu}3^{2dm}\le3^{2d\nu}\sum_{j\ge0}3^{-2dj}\le2\cdot3^{2d\nu}
\end{equation*}
choices of $\hat T$. The cubes of $Y\setminus W$ form a subfamily of $\hat T$, since $\mathcal{T}$
intersects all of them and $|\hat T|\le\nu$: at most $2^\nu$ choices.

Since $Y\cap W$ is a union of islands and the islands share no cube, $Y$ is the union of
$Y\setminus W$ and the islands contained in $Y$. Each such island touches a cube of $Y\setminus W$:
a wall-chain in $Y$ from a cube of the island to $\Delta_*\not\subset W$ leaves the island, and its
first cube off the island touches a cube of the island, so it lies in no other island (islands do not
touch); being in $Y$ and in no island, it lies in $Y\setminus W$. A cube has fewer than $3^d$ touching
neighbours, and each neighbour lies in at most one island. Hence, given $Y\setminus W$, the islands
contained in $Y$ form a subfamily of fewer than $3^d\nu$ islands: at most $2^{3^d\nu}$ choices. As $Y$
is recovered from $Y\setminus W$ and these islands, the number of counted $Y$ is at most
\begin{equation*}
2\cdot3^{2d\nu}\,2^{\nu}\,2^{3^d\nu}
=\exp\bigl(\log2+\nu\,(2d\log3+(1+3^d)\log2)\bigr).
\end{equation*}
With $\nu=\mathsf{c}_d(n+1)$ and $n+1\ge1$, the exponent satisfies
\begin{align*}
\log2+\nu\bigl(2d\log3+(1+3^d)\log2\bigr)
&\le(n+1)\bigl(\log2+\mathsf{c}_d(2d\log3+(1+3^d)\log2)\bigr)\\
&\le\mathsf{k}_d(n+1),
\end{align*}
the last inequality by the choice of $\mathsf{k}_d$ in \eqref{9.1:eq-hull-conventions-a}.

\emph{Proof of \textup{(S5)(c)}.} For each $Y$ in the sum let $n\ge1$ be the integer with
$n-1\le d_W(Y)<n$; then $e^{-\lambda_0d_W(Y)}\le e^{-\lambda_0(n-1)}$, and by (b) at most
$e^{\mathsf{k}_d(n+1)}$ of the $Y$ have $d_W(Y)<n$. Using $\lambda_0-\mathsf{k}_d\ge1$,
\begin{equation*}
\sum_Y e^{-\lambda_0d_W(Y)}
\le\sum_{n\ge1}e^{\mathsf{k}_d(n+1)-\lambda_0(n-1)}
=e^{2\mathsf{k}_d}\sum_{n\ge1}e^{-(\lambda_0-\mathsf{k}_d)(n-1)}
\le e^{2\mathsf{k}_d}(1-e^{-1})^{-1},
\end{equation*}
and the right-hand side is at most $C_{fib}$ by the choice of $C_{fib}$ in
\eqref{9.1:eq-hull-conventions-a}.
\end{proof}

\begin{theorem}[The strip resummation]\label{9.1:thm-strip-resummation}
Assume \textup{(G0)} and the window of strip widths, as fixed in \eqref{9.1:eq-strip-window-a},
and the floor \textup{(F$_\sharp$)} of \eqref{9.1:eq-hull-conventions-a}, that is
$\beta_\sharp\kappa\ge 2(\mathsf{k}_d+1)$. Let $b>0$ be the output rate and $a:=b+\beta_\sharp\kappa$
the input rate fixed in \eqref{9.1:eq-hull-conventions-a}.

Let $\mathfrak{v}=\{\mathfrak{v}(Y_4)\}_{Y_4\in\mathbf{D}_k^Z}$ be any family of numbers, or of
functions, in which case $|\cdot|$ denotes the sup norm. Put
\begin{equation*}
\|\mathfrak{v}\|_a:=\sup\bigl\{e^{a\,d_Z(Y_4)}|\mathfrak{v}(Y_4)| :
Y_4\in\mathbf{D}_k^Z \text{ exterior}\bigr\}.
\end{equation*}
For $Y\in\mathbf{D}_k^{\sharp}$ put
\begin{equation*}
V^{\sharp}(Y):=\sum_{Y_4:\ \operatorname{fil}_\sharp(Y_4)=Y}\mathfrak{v}(Y_4).
\end{equation*}
Then
\begin{equation}\label{9.1:eq-strip-resummation-a}
\begin{gathered}
\text{(a)}\quad \sum_{Y_4\in\mathbf{D}_k^Z}\mathfrak{v}(Y_4)=\sum_{Y\in\mathbf{D}_k^{\sharp}}V^{\sharp}(Y),\\
\text{(b)}\quad |V^{\sharp}(Y)|\le C_{fib}\,e^{-\frac12\beta_\sharp\kappa(m_\sharp-1)}\,
\|\mathfrak{v}\|_a\,e^{-b\,d_\sharp(Y)}\qquad\text{if } Y\setminus Z^{\sharp}\neq\emptyset,\\
\text{(c)}\quad V^{\sharp}(C^{\sharp})=\sum_{Y_4\subset C^{\sharp}}\mathfrak{v}(Y_4)
\quad\text{for every strip } C^{\sharp}.
\end{gathered}
\end{equation}
More precisely:
\begin{enumerate}
\item[(a)] \emph{(identity and locality)} the identity \textup{(a)} of
\eqref{9.1:eq-strip-resummation-a} holds; the map $\mathfrak{v}\mapsto V^{\sharp}$ is linear;
$V^{\sharp}(Y)$ involves only values $\mathfrak{v}(Y_4)$ with $Y_4\subset Y$, and which of them
enter is decided by the components of $Z$ inside $Y$;
\item[(b)] for every $Y\in\mathbf{D}_k^{\sharp}$ with $Y\setminus Z^{\sharp}\neq\emptyset$, the
bound \textup{(b)} of \eqref{9.1:eq-strip-resummation-a} holds;
\item[(c)] for every strip $C^{\sharp}$, the identity \textup{(c)} of
\eqref{9.1:eq-strip-resummation-a} holds.
\end{enumerate}
Clauses \textup{(b)} and \textup{(c)} are the counterparts, for the strips $C^{\sharp}$, of
\cite[(1.68)]{B-CMP122b} and \cite[(1.69)]{B-CMP122b}.
\end{theorem}

\begin{proof}
Throughout we use the geometry of the strips, Lemma~\ref{9.1:lem-strip-geometry}, with its
items \textup{(S1)}--\textup{(S5)} as displayed in \eqref{9.1:eq-strip-geometry-a}.

\emph{(a)} By \textup{(S3)}, $\operatorname{fil}_\sharp$ is a map from $\mathbf{D}_k^Z$ into
$\mathbf{D}_k^{\sharp}$. Hence the sets
$\{Y_4\in\mathbf{D}_k^Z:\operatorname{fil}_\sharp(Y_4)=Y\}$, $Y\in\mathbf{D}_k^{\sharp}$, partition
$\mathbf{D}_k^Z$. Summing $\mathfrak{v}$ first over each such fibre and then over $Y$ gives the
identity \textup{(a)} of \eqref{9.1:eq-strip-resummation-a}. Linearity of
$\mathfrak{v}\mapsto V^{\sharp}$ holds because each $V^{\sharp}(Y)$ is a sum of values of
$\mathfrak{v}$ over an index set that does not depend on $\mathfrak{v}$.

For locality, let $Y_4$ lie in the fibre of $Y$, so $Y\supseteq Y_4$. A strip not contained in
$Y$ shares no cube with $Y$, hence none with $Y_4$. Therefore $\operatorname{fil}_\sharp(Y_4)$ is
computed from the strips inside $Y$. By \textup{(S1)}, these are the components of $Z^{\sharp}$
inside $Y$, each holding exactly one component of $\Lambda_k^c$. So
$\operatorname{fil}_\sharp(Y_4)$ is computed from the components of $Z$ inside $Y$.

\emph{(b)} Fix $Y\in\mathbf{D}_k^{\sharp}$ with $Y\setminus Z^{\sharp}\neq\emptyset$. Let
$\Delta_*$ be the first cube of $Y\setminus Z^{\sharp}$ in the lexicographic order. Let $Y_4$ be
an element of the fibre of $Y$, so that $\operatorname{fil}_\sharp(Y_4)=Y$.

By \textup{(S3)},
\begin{equation*}
Y_4\setminus Z^{\sharp}=\operatorname{fil}_\sharp(Y_4)\setminus Z^{\sharp}=Y\setminus Z^{\sharp}\ni\Delta_* .
\end{equation*}
In particular $Y_4\not\subset Z^{\sharp}$. An interior element lies in one strip by \textup{(S3)},
and the strips are components of $Z^{\sharp}$ by \textup{(S1)}. Hence $Y_4$ is not interior, so it
is exterior in the sense of \eqref{9.1:eq-strip-window-a}. Consequently
\begin{equation*}
|\mathfrak{v}(Y_4)|\le e^{-a\,d_Z(Y_4)}\|\mathfrak{v}\|_a .
\end{equation*}
Again by \textup{(S3)}, $d_Z(Y_4)\ge d_\sharp(\operatorname{fil}_\sharp Y_4)=d_\sharp(Y)$. By
\textup{(S4)}, $d_Z(Y_4)\ge m_\sharp-1$.

Split $a=b+\tfrac12\beta_\sharp\kappa+\tfrac12\beta_\sharp\kappa$. The quantities $b$ and
$\beta_\sharp\kappa$ are positive; for $\beta_\sharp\kappa$ this follows from
\textup{(F$_\sharp$)}. The two lower bounds on $d_Z(Y_4)$ therefore give
\begin{equation*}
e^{-a\,d_Z(Y_4)}\le e^{-b\,d_\sharp(Y)}\cdot e^{-\frac12\beta_\sharp\kappa(m_\sharp-1)}\cdot
e^{-\frac12\beta_\sharp\kappa\,d_Z(Y_4)} .
\end{equation*}
Summing over the fibre, we obtain
\begin{equation*}
|V^{\sharp}(Y)|\le \|\mathfrak{v}\|_a\,e^{-b\,d_\sharp(Y)}\,e^{-\frac12\beta_\sharp\kappa(m_\sharp-1)}
\sum_{Y_4:\ \operatorname{fil}_\sharp(Y_4)=Y}e^{-\frac12\beta_\sharp\kappa\,d_Z(Y_4)} .
\end{equation*}
Every element of this fibre contains the cube $\Delta_*$ and satisfies
$Y_4\setminus Z^{\sharp}=Y\setminus Z^{\sharp}$. We apply \textup{(S5)(c)} to this sum, with the
islands taken to be the components of $Z$ and with rate $\lambda_0:=\tfrac12\beta_\sharp\kappa$.
By \textup{(F$_\sharp$)}, $\lambda_0\ge\mathsf{k}_d+1$. The application bounds the last sum by
$C_{fib}$, which is the bound \textup{(b)} of \eqref{9.1:eq-strip-resummation-a}.

\emph{(c)} By definition, $V^{\sharp}(C^{\sharp})$ is the sum of $\mathfrak{v}(Y_4)$ over the
$Y_4\in\mathbf{D}_k^Z$ with $\operatorname{fil}_\sharp(Y_4)=C^{\sharp}$. By the filing statement
\textup{(S3)}, these are exactly the $Y_4\in\mathbf{D}_k^Z$ with $Y_4\subset C^{\sharp}$. This is
the identity \textup{(c)} of \eqref{9.1:eq-strip-resummation-a}.
\end{proof}

\begin{definition}[The state before the last resummation]\label{9.1:def-resummation-input}
Fix a level $k$ and the large-field region $Z$ at level $k$. We use the cubes, the families
$\mathbf{D}_k$ and $\mathbf{D}_k^Z$, the layered sets $C^{(m)}$ and the tree length $d_Z$ of
Definition~\ref{9.1:def-strip-window}, the window of strip widths
\eqref{9.1:eq-strip-window-a} with its upper end $m_{hi}(k)$, and, with the weight $a$ of the
strip resummation (Theorem~\ref{9.1:thm-strip-resummation}), the norm
\begin{equation*}
  \|\mathfrak{v}\|_a:=\sup\bigl\{e^{a\,d_Z(Y_4)}|\mathfrak{v}(Y_4)|\colon Y_4\text{ exterior}\bigr\},
\end{equation*}
where exterior is meant in the sense of Definition~\ref{9.1:def-strip-window}. As in
Definition~\ref{9.1:def-strip-window}, unions of cubes are compared as families of cubes, and
two of them meet if they share a cube.

\emph{Pre-filing.} The pre-filing is the map
\begin{equation*}
  Y_4\longmapsto Y_4\cup\bigcup\{Z'\colon Z'\text{ a connected component of }Z
  \text{ meeting }Y_4\}.
\end{equation*}
It adds nothing off $Z$ and does not raise $d_Z$. The sets $Y'_i$ of \cite{B-CMP122b} may meet
$Z$ partially; the pre-filing regroups the terms indexed by them, and the count of this
regrouping is a volume factor governed by $\kappa_1$, carried by the bounds that establish the
state defined next.

\emph{The state.} Let $\mathsf{K}_{val}(k)$ and $\mathsf{K}_{cell}(k)$ be two numbers of the
prescribed scaffold at step $k$; they depend on $k$ and are not taken to be independent of $g$.
We say that the \emph{state before the last resummation} holds at step $k$ if the following
holds. The sum of the first part of \cite[(1.49)]{B-CMP122b} with its first two terms removed,
of the terms of \cite[(1.59), (1.60)]{B-CMP122b} indexed by sets $Y_0$, and of the copies
subtracted in Ba{\l}aban's description of the first part of \cite[(1.49)]{B-CMP122b}, is, after
the pre-filing, equal to $\sum_{Y_4\in\mathbf{D}_k^Z}\mathfrak{v}(Y_4)$, where
$\mathfrak{v}=\mathfrak{v}^{br}+\mathfrak{v}^{val}$, each $\mathfrak{v}(Y_4)$ is analytic on the
complex tube over $Y_4$ jointly with the complex sources in the domain $\mathcal{W}_k(Y_4)$
attached to $Y_4$ at level $k$, and
\begin{equation}\label{9.1:eq-resummation-input-a}
  \begin{aligned}
    &\|\mathfrak{v}^{br}\|_a\le\varepsilon_{br},
      &&\text{(i)}\\
    &\|\mathfrak{v}^{val}\|_a\le\mathsf{K}_{val}(k),
      &&\text{(ii)}\\
    &\sum_{Y_4\subset C^{(m)}}|\mathfrak{v}(Y_4)|
      \le\mathsf{K}_{cell}(k)\sum_j\bigl|\Gamma^0_j\cap C^{(m)}\bigr|
      \quad\text{for every }C\text{ and every }m\le m_{hi}(k),
      &&\text{(iii)}
  \end{aligned}
\end{equation}
with $\varepsilon_{br}\to0$ as $\gamma\to0$; here $\Gamma^0_j$ is the set of cells of zone $j$.
By the definition of $\|\cdot\|_a$, the bound \textup{(ii)} says that
$|\mathfrak{v}^{val}(Y_4)|\le\mathsf{K}_{val}(k)\,e^{-a\,d_Z(Y_4)}$ for every exterior $Y_4$; in
\cite{B-CMP122b} these terms are bounded by a constant times an exponential factor, and
\textup{(ii)} records that constant as $\mathsf{K}_{val}(k)$.

\emph{The two constants.} Let $\alpha_*$ be the least of the finitely many ceilings on $\alpha$
that are independent of $g$ in \cite[(1.93)--(1.99)]{B-CMP122b}, in the form in which they are
used in this chapter; it satisfies
\begin{equation}\label{9.1:eq-resummation-input-1}
  \alpha_*^{1/3}\le\min\Bigl\{e^{-(1+\beta)\kappa 2^d},\
  \frac{\beta_v\kappa}{4C_{fib}},\ c_{KP}(d,\beta,\kappa),\ \frac{B_0}{O(1)}\Bigr\},
\end{equation}
where $c_{KP}(d,\beta,\kappa)$ and $B_0$ are the constants entering two of the ceilings in
\cite[(1.93)--(1.99)]{B-CMP122b}, the first depending only on $d$, $\beta$, $\kappa$, and $O(1)$
stands for a constant independent of $g$. We put
\begin{equation}\label{9.1:eq-resummation-input-2}
  m_{lo}(k):=1+\Bigl\lceil\frac{2}{\beta_\sharp\kappa}
  \log\frac{C_{fib}\bigl(1+\mathsf{K}_{val}(k)\bigr)}{\alpha_*}\Bigr\rceil .
\end{equation}
Like $\mathsf{K}_{val}(k)$ and $\mathsf{K}_{cell}(k)$, the integer $m_{lo}(k)$ is a number of the
prescribed scaffold at step $k$; it is fixed after $\gamma$ in the order of choice.
\end{definition}

\begin{proof}
Each set adjoined to $Y_4$ by the pre-filing is a connected component of $Z$, so every cube of
the image that is not a cube of $Y_4$ is a cube of $Z$: the pre-filing adds nothing off $Z$.
That it does not raise $d_Z$ is part of the regrouping statement for the terms indexed by the
sets $Y'_i$, and is used here by statement. The number $\alpha_*$ is the least of
finitely many numbers, so it
is well defined. Since $2^d\ge 2d$ for every integer $d\ge1$, we have
\begin{equation*}
  e^{-(1+\beta)\kappa 2^d}\le e^{-(1+\beta)\kappa\, 2d},
\end{equation*}
so the first entry of \eqref{9.1:eq-resummation-input-1} implies the same condition with
$2d$ in place of $2^d$ in the exponent.
\end{proof}

\begin{corollary}[Smallness of the re-filed values at maximal strip width]
\label{9.1:cor-refiled-smallness}
Assume the state before the last resummation, Definition~\ref{9.1:def-resummation-input},
with $\varepsilon_{br}\le 1$, and let the strip width satisfy $m_{lo}(k)\le m_\sharp\le m_{hi}(k)$,
the window of strip widths of \eqref{9.1:eq-strip-window-a}.
Let $V^{\sharp}$ be the family obtained from the input family $\mathfrak{v}$ of
\eqref{9.1:eq-resummation-input-a} by the strip resummation,
Theorem~\ref{9.1:thm-strip-resummation}. Then, for every $Y\in\mathbf{D}_k^{\sharp}$,
$V^{\sharp}(Y)$ is analytic on the tube over $Y$,
and for every such $Y$ with $Y\setminus Z^{\sharp}\neq\emptyset$ the bound in the first line of
the following display holds, with $\alpha^{\sharp}$ defined in its second line, where it is also
asserted that $\alpha^{\sharp}\le\alpha_*$; its third line is a condition on $\gamma$, used below:
\begin{equation}\label{9.1:eq-refiled-smallness-a}
\begin{aligned}
&|V^{\sharp}(Y)|\le \alpha^{\sharp}e^{-b\,d_\sharp(Y)},\\
&\alpha^{\sharp}:=C_{fib}\bigl(\varepsilon_{br}+\mathsf{K}_{val}(k)\bigr)
e^{-\frac12\beta_\sharp\kappa(m_\sharp-1)}\le\alpha_*,\\
&\sup_{g\le\gamma}C_{fib}\bigl(1+\mathsf{K}^0+N(g)C_0^{\sharp}\bigr)
e^{-\frac14\beta_\sharp\kappa(R(g)-4)}\le\alpha_*
\end{aligned}
\end{equation}
Here $\mathsf{K}^0$ and $C_0^{\sharp}$ are constants and
$N(\cdot)$ is the function of the coupling in the bound
\eqref{9.1:eq-value-cell-constants-a} on the value constant,
$\mathsf{K}_{val}(k)\le\mathsf{K}^0+N(g_k)C_0^{\sharp}$. For every strip $C^{\sharp}$
\begin{equation}\label{9.1:eq-refiled-smallness-1}
|V^{\sharp}(C^{\sharp})|\le\mathsf{K}_{cell}(k)\sum_j|\Gamma^0_j\cap C^{\sharp}|.
\end{equation}
At $m_\sharp=m_{hi}(k)$,
\begin{equation*}
\alpha^{\sharp}\le C_{fib}\bigl(1+\mathsf{K}_{val}(k)\bigr)e^{-\frac14\beta_\sharp\kappa(R_k-4)},
\end{equation*}
where $R_k=R(g_k)$ is the radius at level $k$ that fixes $m_{hi}(k)$ in
\eqref{9.1:eq-strip-window-a}, $R(\cdot)$ denotes this radius as a function of the coupling,
and $r$ is the exponent in its lower bound $R(g)\ge(\log g^{-2})^r$.
The right-hand side tends to $0$ as $\gamma\to0$, uniformly in $k$. At $m_\sharp=m_{hi}(k)$ one
has moreover, provided the exponent $r$ satisfies $r\ge2$,
$(\alpha^{\sharp})^{1/3}\le g_k^{\kappa_0}$ for $\gamma$ small, where
$\kappa_0$ is the coupling exponent fixed in the coupling-power bound on the re-filed amplitude
of this chapter.
The third line of \eqref{9.1:eq-refiled-smallness-a} is a condition on $\gamma$. The window is
non-empty at every $k$, and $\alpha^{\sharp}\le\alpha_*$ at $m_\sharp=m_{hi}(k)$, as soon as
this condition holds.
That condition is a single ceiling on $\gamma$, in the sense of
Definition~\ref{2.3:not-stages-first}, of degree $-\infty$; with $T=\log g^{-2}$ and
$T_\gamma=\log\gamma^{-2}$ its left member is the supremum over $T\ge T_\gamma$ of the same
expression. It enters the order of choice of the constants at the place listed in
\eqref{9.1:eq-constant-order-a}.
\end{corollary}

\begin{proof}
\emph{Analyticity.} By definition $V^{\sharp}(Y)$ is the sum of the $\mathfrak{v}(Y_4)$ over
the $Y_4$ with $\operatorname{fil}_\sharp(Y_4)=Y$; this is a finite sum, and by the locality
clause~(a) of Theorem~\ref{9.1:thm-strip-resummation} every $Y_4$ occurring in it satisfies
$Y_4\subset Y$. By Definition~\ref{9.1:def-resummation-input}, each $\mathfrak{v}(Y_4)$ is
analytic on the tube over $Y_4$. By the restriction property of the tubes used at every
resummation of this chapter, a function of the variables over $Y_4$ that is analytic on the
tube over $Y_4$, read as a function of the variables over $Y\supset Y_4$, is analytic on the
tube over $Y$. A finite sum of functions
analytic on the
tube over $Y$ is analytic there, so $V^{\sharp}(Y)$ is analytic on the tube over $Y$.

\emph{The bounds in the window.} Write $\mathfrak{v}=\mathfrak{v}^{br}+\mathfrak{v}^{val}$ as in
\eqref{9.1:eq-resummation-input-a}. The input bounds there give
\begin{equation*}
\|\mathfrak{v}\|_a\le\|\mathfrak{v}^{br}\|_a+\|\mathfrak{v}^{val}\|_a
\le\varepsilon_{br}+\mathsf{K}_{val}(k).
\end{equation*}
Clause~(b) of
Theorem~\ref{9.1:thm-strip-resummation} states, for $Y\setminus Z^{\sharp}\neq\emptyset$,
\begin{equation*}
|V^{\sharp}(Y)|\le C_{fib}e^{-\frac12\beta_\sharp\kappa(m_\sharp-1)}\|\mathfrak{v}\|_a
e^{-b\,d_\sharp(Y)}.
\end{equation*}
Inserting the bound on $\|\mathfrak{v}\|_a$ gives the bound in the first line of
\eqref{9.1:eq-refiled-smallness-a}, with $\alpha^{\sharp}$ as in its second line. The strip
bound \eqref{9.1:eq-refiled-smallness-1} is clause~(c) of
Theorem~\ref{9.1:thm-strip-resummation}, applied to the input family of
Definition~\ref{9.1:def-resummation-input} with its cell constant $\mathsf{K}_{cell}(k)$.

\emph{The upper end of the window.} At $m_\sharp=m_{hi}(k)=\lfloor\frac12(R_k-1)\rfloor$, see
\eqref{9.1:eq-strip-window-a}, one has
$m_\sharp-1=\lfloor\frac12(R_k-1)\rfloor-1$, and
\begin{equation*}
\lfloor\tfrac12(R_k-1)\rfloor-1\ge\tfrac12(R_k-4).
\end{equation*}
Together with $\varepsilon_{br}\le1$ this gives
\begin{equation*}
\alpha^{\sharp}\le C_{fib}\bigl(1+\mathsf{K}_{val}(k)\bigr)e^{-\frac14\beta_\sharp\kappa(R_k-4)}.
\end{equation*}
Moreover, since the running couplings satisfy $g_k\le\gamma$,
\begin{equation*}
R_k\ge(\log g_k^{-2})^r\ge(\log\gamma^{-2})^r,
\end{equation*}
so that the exponential factor in this bound is bounded uniformly in $k$:
\begin{equation*}
e^{-\frac14\beta_\sharp\kappa(R_k-4)}\le e^{\beta_\sharp\kappa}
e^{-\frac14\beta_\sharp\kappa(\log\gamma^{-2})^r}.
\end{equation*}
Bound the value constant by
$\mathsf{K}_{val}(k)\le\mathsf{K}^0+N(g_k)C_0^{\sharp}$, from
\eqref{9.1:eq-value-cell-constants-a}, and use $R_k=R(g_k)$. The right-hand side of the
last bound on $\alpha^{\sharp}$ is then at most the left member of the third line of
\eqref{9.1:eq-refiled-smallness-a}. Hence the condition in that line gives
$\alpha^{\sharp}\le\alpha_*$ at $m_\sharp=m_{hi}(k)$, for every $k$.

\emph{Finiteness and decay of the supremum.} Put $c=\frac14\beta_\sharp\kappa$ and
$T=\log g^{-2}$. The function $N(\cdot)$ grows at most like $T^r$: $N(g)\le C'T^r$ with a
constant $C'$ independent of $g$. On the other hand $R(g)\ge T^r$ gives
\begin{equation*}
e^{-c(R(g)-4)}\le e^{4c}e^{-cT^r}.
\end{equation*}
Since $T^re^{-cT^r}\to0$ as $T\to\infty$, the product
$\bigl(1+\mathsf{K}^0+N(g)C_0^{\sharp}\bigr)e^{-c(R(g)-4)}$ is bounded on $T\ge T_\gamma$, so the
supremum over $T\ge T_\gamma$ in the third line of \eqref{9.1:eq-refiled-smallness-a} is finite
and tends to $0$
as $T_\gamma\to\infty$, that is, as $\gamma\to0$. Since the bound at $m_{hi}(k)$ depends on $k$
only through $g_k\le\gamma$, the convergence to $0$ is uniform in $k$.

Finally, the bound $(\alpha^{\sharp})^{1/3}\le g_k^{\kappa_0}$ for $\gamma$ small and $r\ge2$ is
the coupling-power bound on the re-filed amplitude, applied to the bound on $\alpha^{\sharp}$ at
$m_\sharp=m_{hi}(k)$ obtained above.
\end{proof}

\begin{lemma}[From value units to domains: a column bound off the large-field region]
\label{9.1:lem-column-bound}
Let $Z$ be the large-field region and $k$ the level of
Definition~\ref{9.1:def-resummation-input}. Let $\mathcal{V}$ be the finite set of value units of the
expansion described there. These are:
\begin{itemize}
\item the constituents of the first part whose expansion is trivial (all $Y'_i=\emptyset$);
\item the subtracted copies;
\item the values of the terms $\mathbf{B}$;
\item the values of the chain terms $\mathbb{C}^{(i)}$.
\end{itemize}
Each $v\in\mathcal{V}$ has a value $\mathrm{val}_v$ and a domain $K_v\in\mathbf{D}_k$ which meets $Z$.
Here $d_k$ denotes the tree length of level-$k$ domains, and $d_Z$ the relative tree length of
\eqref{9.1:eq-strip-window-a} with $W=Z$. The number $a$ is the exponent of the weighted norms in the
bounds \eqref{9.1:eq-resummation-input-a}. Put
\begin{align*}
\|v\| &:= \sup|\mathrm{val}_v|, \qquad
Y_1(v) := K_v\cup\{\text{components of } Z \text{ meeting } K_v\},\\
\mathfrak{v}^{val}(Y_1) &:= \sum_{v:\,Y_1(v)=Y_1}\mathrm{val}_v ,
\end{align*}
and
\begin{equation}\label{9.1:eq-column-bound-1}
\mathcal{V}^{ext} := \{v\in\mathcal{V}:\ K_v\setminus Z\neq\emptyset\},\qquad
\mathsf{K}_{col}(k) := \sup_{\Delta\not\subset Z}
\sum_{v\in\mathcal{V}^{ext},\ K_v\supset\Delta} e^{a\,d_k(K_v)}\|v\| .
\end{equation}
The supremum runs over the level-$k$ cubes $\Delta$ not contained in $Z$. Then the value bound
among the bounds \eqref{9.1:eq-resummation-input-a} of Definition~\ref{9.1:def-resummation-input} holds
with $\mathsf{K}_{val}(k):=\mathsf{K}_{col}(k)$.
\end{lemma}

\begin{proof}
Fix $v\in\mathcal{V}$. The set $Y_1(v)$ arises from $K_v$ by adjoining the components of $Z$ that
meet $K_v$. A component meets $K_v$ when it shares a cube with $K_v$, so no connecting cubes are
adjoined. Hence $Y_1(v)$ belongs to $\mathbf{D}_k^Z$. Every adjoined cube lies in $Z$, so
\begin{equation}\label{9.1:eq-column-bound-2}
Y_1(v)\setminus Z = K_v\setminus Z .
\end{equation}

Now fix $Y_1$ with $Y_1\setminus Z\neq\emptyset$. Let $\Delta_*$ be the first cube of
$Y_1\setminus Z$ in a fixed enumeration of the level-$k$ cubes. Let $v$ lie in the fibre
$\{Y_1(v)=Y_1\}$.

By \eqref{9.1:eq-column-bound-2} we have $K_v\setminus Z = Y_1\setminus Z\neq\emptyset$, so
$v\in\mathcal{V}^{ext}$. Moreover
\[
\Delta_*\subset Y_1\setminus Z = K_v\setminus Z\subset K_v .
\]

Take an optimal tree of $K_v$, that is, one realising $d_k(K_v)$. It lies in $K_v\subset Y_1$ and
intersects every cube of $K_v$. Since $K_v\supseteq Y_1\setminus Z$, it intersects in particular every
cube of $Y_1\setminus Z$. By the definition of the relative tree length,
\begin{equation}\label{9.1:eq-column-bound-3}
d_Z(Y_1)\le d_k(K_v) .
\end{equation}

At every point of the domain of $\mathfrak{v}^{val}(Y_1)$ we now estimate as follows:
\begin{align*}
e^{a\,d_Z(Y_1)}\,|\mathfrak{v}^{val}(Y_1)|
&\le \sum_{v:\,Y_1(v)=Y_1} e^{a\,d_Z(Y_1)}\|v\|
\le \sum_{v:\,Y_1(v)=Y_1} e^{a\,d_k(K_v)}\|v\| \\
&\le \sum_{v\in\mathcal{V}^{ext},\ K_v\supset\Delta_*} e^{a\,d_k(K_v)}\|v\|
\le \mathsf{K}_{col}(k).
\end{align*}
The steps are justified as follows.
\begin{itemize}
\item The first inequality is the triangle inequality together with $|\mathrm{val}_v|\le\|v\|$.
\item The second follows from \eqref{9.1:eq-column-bound-3}, since the weight $e^{a\,d}$ is nondecreasing in $d$.
\item The third holds because the fibre is contained in the index set of the right-hand sum, as shown above, and all terms are nonnegative.
\item The fourth holds because $\Delta_*$, a cube of $Y_1\setminus Z$, is not contained in $Z$, so it is admissible in the supremum \eqref{9.1:eq-column-bound-1}.
\end{itemize}
Taking the supremum over the domain gives the value bound with
$\mathsf{K}_{val}(k)=\mathsf{K}_{col}(k)$. The argument involves no sum over auxiliary cubes and
requires no gain in the decay rate $a$; see the remark below.
\end{proof}

\begin{remark}[Why the column supremum is anchored off the large-field region]
The constant of Lemma~\ref{9.1:lem-column-bound} is defined level by level. The supremum defining
it runs only over cubes not contained in $Z$ and over value units whose domains reach outside $Z$.
Two features of this constant cannot be improved: it cannot be chosen uniformly in the coupling,
and the supremum cannot be extended to units with $K_v\subset Z$. Observation (a) bears on the
first; observations (b) and (c) bear on the second.

(a) The set $\mathcal{V}$ contains the values of the chain terms. \cite{B-CMP122b} states that the term
$\mathbf{B}''_k$ is the sum of the new boundary terms created by the preliminary integrations. The
value produced at one stage has norm bounded by a constant, which we denote $C_0^{\sharp}$. Summed
over $i$ stages, the bound is $iC_0^{\sharp}$, where $i$ is at most the number of preliminary
integrations.

\cite{B-CMP122a} states that the bounds from the corresponding terms cumulate and that there are at
most $N_0$ such terms, $N_0$ being an integer of Ba{\l}aban's construction there; in unweighted form
the bound is $3(N_0+3)C_0$. \cite[p.~192]{B-CMP122a} states that these terms are multiplied by $g_j$
and that their bounds involve Ba{\l}aban's constants $\delta'_j$ (the bound $4N_0^2\delta'_k$ is
displayed there). In the present bookkeeping these factors are attached to the bracketed terms in
the fluctuation field of the expansion of Definition~\ref{9.1:def-resummation-input}, and not to the
values $\mathrm{val}_v$; the values carry no factor $\mathsf{w}_L^{-1}$. Hence $\mathsf{K}_{val}$
cannot be chosen uniformly in the coupling: its supremum over small couplings $g$ is infinite.

(b) A supremum over all value units, including those with $K_v\subset Z$, would range over interior
units as well. The part of $Z$ forming its zone of level $k-s$, $s$ steps below level $k$, fills
level-$k$ cubes, each holding $(ML^s)^d$ unit cells, and each cell contributes a value of order one.
Such a constant would therefore satisfy $\mathsf{K}_{col}\gtrsim(M\check{R}_k)^d$.

(c) A chain term $\mathbb{C}^{(i)}$ placed at level $k$ decays at exactly the rate $a$ of the input
bounds \eqref{9.1:eq-resummation-input-a}. Its gain in rate is $\kappa_c=0<d$.
\end{remark}

\begin{proposition}[Bounds on the value and cell constants]\label{9.1:prop-value-cell-constants}
Let $k$ be the level at which the inputs of \eqref{9.1:eq-resummation-input-a} are read, and let
$\mathsf{K}_{val}(k)$ and $\mathsf{K}_{cell}(k)$ be the constants of its value input and of its
volume input. Let $N(g_k)$ be the bound on the number of stages of the chain of preliminary
integrations at level $k$. Let $C_0^{\sharp}$ be the bound, for one stage, on the norm
$\|\cdot\|_a$, at the rate $a$ of \eqref{9.1:eq-hull-conventions-a}, of the chain terms of that
stage. Let $N_0(g_k)$ be the stage count that enters part (a) of the payment bound for the chain
terms, and let $C_0$ be the constant of that part. Assume the floors and rates of
\eqref{9.1:eq-hull-conventions-a}, in particular the floor $L\ge 3c_H$, the hypotheses of
Lemma~\ref{9.1:lem-fresh-rate}, and the condition, of ceiling type on $\gamma$,
\begin{equation}\label{9.1:eq-value-cell-constants-1}
\bigl(N_0(g_k)+3\bigr)\,C_0^{\sharp}\,e^{-\kappa\check{R}_k/40}\le 1 .
\end{equation}
Then
\begin{equation}\label{9.1:eq-value-cell-constants-a}
\mathsf{K}_{val}(k)\le \mathsf{K}^0+N(g_k)\,C_0^{\sharp},
\qquad
\mathsf{K}_{cell}(k)\le \mathsf{K}^0_{cell}+3\bigl(N_0(g_k)+3\bigr)C_0,
\end{equation}
where the constant $\mathsf{K}^0$ does not depend on $g$.
\end{proposition}

\begin{proof}
\emph{Reduction.} Throughout, $Z$ is the large-field region at level $k$. The sets
$\mathcal{V}^{ext}\subset\mathcal{V}$, the sets $K_v$ and the norms $\|v\|$ are those of
Lemma~\ref{9.1:lem-column-bound}, and
\begin{equation*}
a=\tfrac{8}{5}\kappa-\tfrac{1}{20}\kappa=\tfrac{31}{20}\kappa .
\end{equation*}
By Lemma~\ref{9.1:lem-column-bound} the value input holds with
$\mathsf{K}_{val}(k):=\mathsf{K}_{col}(k)$, where
\begin{equation*}
\mathsf{K}_{col}(k)=\sup_{\Delta\not\subset Z}\ \sum_{v\in\mathcal{V}^{ext},\,K_v\supset\Delta}
e^{a\,d_k(K_v)}\,\|v\| .
\end{equation*}
We fix a cube $\Delta\not\subset Z$ and bound the contribution of each kind of unit to this sum. A
unit that is not in $\mathcal{V}^{ext}$ contributes nothing.

\emph{Step \textup{(i)}: the chain terms $\mathbb{C}^{(n)}$ of the present step.}
The domains $X$ of these terms are described in \cite[(1.63)]{B-CMP122a}. In the notation of that
paper, $X\subset\Omega^c_{m+1}$ for some $m\ge j+1$, $X\cap\Omega_m\neq\emptyset$,
$X\in\mathbb{D}_m$, and $X\cap(Z_{j+1}\setminus Z'_{j+1})\neq\emptyset$, where $\Omega_m$,
$Z'_{j+1}$, $\mathbb{D}_m$ and the index $j$ carry their meaning in \cite[(1.63)]{B-CMP122a},
and $Z_{j+1}=\Lambda_{j+1}^c$; in this proof the letter $m$ denotes the level of a piece, which
does not exceed $k$, and we distinguish the pieces with $m<k$ from those with $m=k$.

Consider first the pieces with $m<k$. Such a piece lies in $\Lambda_k^c$, is connected and meets
$Z$. Hence $X\subset Z$, and the unit is not in $\mathcal{V}^{ext}$.

Consider next the pieces with $m=k$. For the fixed cube $\Delta$, the sum over $X\supset\Delta$ of
$e^{a\,d_k(X)}\|\mathbb{C}^{(n)}(X)\|$ is the norm $\|\mathbb{C}^{(n)}\|_a$ of the chain term at
rate $a$. This norm is bounded by $C_0^{\sharp}$ for each stage, by the payment bound for the chain
terms, part (d2). There are at most $N(g_k)$ stages, where
\begin{equation*}
N(g_k)\le O(1)\bigl(\log g_k^{-2}\bigr)^{r}
\end{equation*}
by part (d) of the flat stage-count statement, $r$ being the fixed exponent of that statement.
The pieces with $m=k$ therefore contribute at most $N(g_k)\,C_0^{\sharp}$.

\emph{Step \textup{(ii)}: the left-over terms $\mathbb{C}^{(i)}$ of a chain stopped at level $k_1$
in $\bar c$.} Here $\bar c$ denotes the region in which the chain was stopped; it lies in a
component of $Z$, as recalled below. These terms are treated as leaves, with the hulls of
Notation~\ref{9.1:not-hull-conventions} and \eqref{9.1:eq-hull-conventions-a}.

Consider first the levels $\ell<k_1$. Every piece of the stopped chain below its stopping level
satisfies $X\subset\bar c$, by the statement that such pieces stay in $\bar c$. The region $\bar c$
lies at least ten layers of width $M\check{R}_k$ inside its component of $Z$
\cite{B-CMP122a}. Hence $K(X)\subset Z$, and the unit is not in $\mathcal{V}^{ext}$.

Consider next the level $\ell=k_1$, where the rate is $(1+3\beta_{pr})\kappa$. A unit of
$\mathcal{V}^{ext}$ of this kind joins $\bar c$ to $Z^c$, so
\begin{equation*}
d_{k_1}(X)\ge \check{R}_k\,L^{k-k_1}.
\end{equation*}
The price of passing to the hull is governed by $a\,c_H/L$. The floor $L\ge 3c_H$ of
\eqref{9.1:eq-hull-conventions-a} gives $c_H/L\le\frac13$, hence
\begin{equation*}
a\,\frac{c_H}{L}\le\frac{31}{20}\cdot\frac13\,\kappa=\frac{31}{60}\,\kappa .
\end{equation*}
After this price, more than $\kappa\,d_{k_1}(X)$ of the decay is left. Since
$d_{k_1}(X)\ge\check{R}_k$, a part $\kappa\check{R}_k/40$ of it pays for the sum of the stage
bounds of the stopped chain, which is at most $(N_0(g_k)+3)C_0^{\sharp}$: by the ceiling
\eqref{9.1:eq-value-cell-constants-1}, this sum times $e^{-\kappa\check{R}_k/40}$ is at most $1$.
With it the constant of these units is $1$. The rest of the decay sums over $X$.

\emph{Step \textup{(iii)}: stored and fresh leaves, and subtracted copies.} We use the constants
$c_H$, $c^{\sharp}=c_H+1$ and $\kappa_R=c^{\sharp}\kappa$ of
Notation~\ref{9.1:not-hull-conventions}. For sets $Y\subset Y'\subset Y^{(1)}$, where $Y^{(1)}$ is
the fattened set of Notation~\ref{9.1:not-hull-conventions} (written $Y^{+}$ there), the tree
length satisfies
\begin{equation*}
d(Y')\le c_H\bigl(d(Y)+1\bigr).
\end{equation*}
We write $\bar X$ for the smallest $\mathbf{D}_k$-domain containing $X$. For each kind of unit we
list three things: the set $K_v$, the decay it supplies, and the statement from which that decay is
taken.

\emph{Any unit, read in the core of the construction used in the proof of the correlation
theorem.} Here $K_v=\bar X$. The weight $e^{(1+3\beta_s)\kappa\,d_k(\bar X)}$ is summed,
and the factor $e^{-\frac14\delta\kappa\,d_j(X)}$ is kept, where $\delta=\delta_\star$ is the number
fixed in Lemma~\ref{9.1:lem-fresh-rate}, so the column rate is
$\kappa_c=\frac14\delta\kappa$. This is parts (i) and (ii) of the boundary-term lemma used in the
proof of the correlation theorem, rows one to five of its table.

\emph{A stored leaf created at level $j<k$.} Here $K_v$ is either the hull $(X^{+})_0$ of
\eqref{9.1:eq-hull-conventions-a} or the fat core $\bar X^{(1)}$, so that
$\bar X\subset K_v\subset\bar X^{(1)}$. By \cite[(1.31)]{B-CMP116},
\begin{equation*}
d_j(X)\ge L^{k-j}\,d_k(\bar X).
\end{equation*}
The floor $L\ge 3c_H$ gives, for the fraction $1-7c_H/(4L)$ of the rate that remains after the
hull's price,
\begin{equation*}
1-\frac{7c_H}{4L}\ge 1-\frac{7}{12}=\frac{5}{12}.
\end{equation*}
The column rate is $\kappa_c=\delta\kappa/(4c_H)$, and the bound carries the factor $A_H$. This is
parts (1) and (2) of the hull lemma used in the proof of the correlation theorem, under
$L\ge3c_H$.

\emph{A fresh leaf of level $k$.} The set $K_v$ is as for stored leaves, and the supplied decay is
\begin{equation*}
e^{-(1+4\beta^{\star})\kappa_R\,d_k(X)} .
\end{equation*}
This is part (3) of the same hull lemma. Its input is the fresh-rate lemma,
Lemma~\ref{9.1:lem-fresh-rate}, which holds for every family created at the small-field site.

\emph{The stored terms $R'$ and $\mathbf{B}'$ of the large-field steps, and left-over terms
$\mathbb{C}^{(i)}$, taken as leaves.} The set $K_v$ is as for stored leaves. The terms $R'$ and
$\mathbf{B}'$ are supplied at rate $\kappa$. A left-over term $\mathbb{C}^{(i)}$ is supplied at rate
$(1+3\beta_{pr})\kappa$ at its top level and, below it, at rate $a_{\flat}$ of
Notation~\ref{9.1:not-hull-conventions}, under the floor $\frac12 a_{\flat}L\ge\frac74 c_H\kappa$
of \eqref{9.1:eq-hull-conventions-a}. This is parts (b)--(d) of the payment bound for the
chain terms, the rate statement for pieces filed below their top level, and Step (ii) above.

\emph{Derived chain terms $\mathbb{C}^{(n)}$.} Here $K_v=\bar X$. The rates are
$a_k=(1+3\beta)\kappa$ at level $k$ and $a_m=(1-\beta)\kappa/c^{\sharp}$ at the levels
$m<k$ below the top, with the number $\beta$ of the flat derived-chain bound, under the floor
\begin{equation*}
(1-\beta)\,a_m L\ge(1+3\beta)\kappa
\end{equation*}
of \eqref{9.1:eq-hull-conventions-a}. This is the flat derived-chain bound.

\emph{Subtracted copies.} Here $K_v$ is the set $K$ of the leaf that is copied, and the supplied
decay is that leaf's. This is the statement on subtracted copies.

\emph{The bracket terms $T(Y_0)$.} These enter the bracket input of
\eqref{9.1:eq-resummation-input-a}, with
\begin{equation*}
\varepsilon_{br}=C_{\sharp}\,O(1)\,A_H\,\mathsf{w}_L^{-1},
\end{equation*}
where $C_{\sharp}$ is the constant of the weight statement used in the proof of the correlation
theorem. This is that weight statement, together with the
bracket resummation at the large-field site (Lemma~\ref{9.1:lem-bracket-margins}) and the hull
lemma.

The factor $e^{-\frac14\delta\kappa\,d_j}$ kept in the first kind of unit sums over the anchor
points $y$ of the units, by part (ii) of the boundary-term lemma and by the closing sentence of the
argument for that part. The anchor points $y$ in zones
$n<k$ lie at least $2\check{R}_k$ cubes inside $Z$ \cite{B-CMP119}. Hence a unit of
$\mathcal{V}^{ext}$ anchored at such a point has
\begin{equation*}
d_j(X)\ge\bigl(2\check{R}_k-1\bigr)L^{k-j},
\end{equation*}
and the resulting decay beats the counting factor $(ML^{k-j})^d$.

\emph{Step \textup{(iv)}: the volume input.} The bound per unit cell of zone $n$, for any cell, is
the last column of the table of boundary terms. This is part (iv) of the boundary-term lemma, rows
one to five. To it is added the unweighted sum of part (a) of the payment bound for the chain
terms; this is the origin of the summand $3(N_0(g_k)+3)C_0$ in the second inequality of
\eqref{9.1:eq-value-cell-constants-a}. The last-column bounds give $\mathsf{K}^0_{cell}$, and with
the added sum this is the second inequality of \eqref{9.1:eq-value-cell-constants-a}.

\emph{Conclusion.} By Step (i), the pieces with $m=k$ of the chain terms of the present step
contribute at most $N(g_k)C_0^{\sharp}$ to $\mathsf{K}_{col}(k)$. The pieces with $m<k$, and the
left-over pieces below their stopping level, are not in $\mathcal{V}^{ext}$. All remaining units
are summed by the decay listed in Steps (ii) and (iii), with constants that do not depend on $g$.
Their contribution is $\mathsf{K}^0$. Since $\mathsf{K}_{val}(k)=\mathsf{K}_{col}(k)$, this is the
first inequality of \eqref{9.1:eq-value-cell-constants-a}.
\end{proof}

\begin{definition}[Translation of the printed collar into strips]\label{9.1:def-collar-dictionary}
We work at level $k$ with the cubes, the cube count $N(\cdot)$, the tree lengths $d_W$ relative to a
union of cubes $W$, and the strips of Definition~\ref{9.1:def-strip-window}. Let $Z$ be a
large-field region at level $k$, and let $\{X_j\}$ and $\{Y_i\}$ be its components of first class
and of second class respectively, in the sense of \cite{B-CMP122b}. For the first-class components
$X_{j_h}$ entering $Z'$ in \cite[(1.71)--(1.72), (1.90)--(1.99)]{B-CMP122b} we put
\begin{equation*}
  Z' := \bigcup_i Y_i \,\cup\, \bigcup_h X_{j_h},
  \qquad
  N_h := N\bigl(X^{\sharp}_{j_h}\bigr).
\end{equation*}

The \emph{dictionary} $\mathfrak{D}^{\sharp}$ is the following list of substitutions, made in
\cite[(1.71)--(1.72), (1.90)--(1.99)]{B-CMP122b}.
\begin{enumerate}
  \item[(a)] Ba{\l}aban's collared regions, written $Z^{\sim}$, $X_j^{\sim}$, $Y_i^{\sim}$ in
  \cite[(1.71)--(1.72), (1.90)--(1.99)]{B-CMP122b} and denoted $Z^{\boxplus}$, $X_j^{\boxplus}$,
  $Y_i^{\boxplus}$ here, the collar of each of the regions $Z$, $X_j$, $Y_i$ being taken as in
  Definition~\ref{9.1:def-strip-window}, are replaced by the strips $Z^{\sharp}$, $X_j^{\sharp}$,
  $Y_i^{\sharp}$.
  \item[(b)] The tree lengths $d_{k,Z}$, $d_{k,Z'}$ and $d_{k,\cup Y_i}$ are replaced by $d_{\sharp}$,
  by $d_{Z'^{\sharp}}$ and by the \emph{output length}
  \begin{equation*}
    d_{out} := d_{\bigcup_i Y_i^{\sharp}} ,
  \end{equation*}
  respectively.
  \item[(c)] The volume $|X_{j_h}|$ is replaced by $|X^{\sharp}_{j_h}|$.
  \item[(d)] The coefficient $\beta/(3\cdot 2^{d})$, where $\beta$ is the constant of
  \cite{B-CMP122b} appearing in these displays, is replaced by
  \begin{equation*}
    \beta_v := \frac{\beta}{2\mathsf{c}_d},
  \end{equation*}
  where $\mathsf{c}_d = 2\cdot 3^{d}$ is the combinatorial constant of the strips.
\end{enumerate}
In \cite[(1.67)--(1.68), (1.92)]{B-CMP122b} the region $Z$ is read as its collared region
$Z^{\boxplus}$, so that at \cite[(1.92)]{B-CMP122b} the substitution (a) replaces it by the strip
$Z^{\sharp}$; this reading is used so that the surface terms carried by the collar cubes of the
second-class components $Y_i$ are among those bounded at these displays.

The dictionary $\mathfrak{D}^{\sharp}$ is not applied at one place in \cite[\S3]{B-CMP122b}, nor
in \cite[(2.30), (2.41)]{B-CMP119}; at those places the tilde continues to denote the collars
$Z^{\boxplus}$, $X_j^{\boxplus}$, $Y_i^{\boxplus}$ of (a).

The \emph{strip reading} of \cite[(1.90)--(1.100)]{B-CMP122b} is the argument of those displays,
carried out line by line for a large-field region $Z$ and its collar, with the strips of (a) in
place of the collars, the substitutions (b)--(d) made, and $\{X_j\}$, $\{Y_i\}$, $Z'$ and $N_h$ as
above; it is the collar form of the cluster-expansion step for these displays, and it is carried
out in the two lemmas that state that collar form.
\end{definition}

\begin{lemma}[Ba{\l}aban's combinatorics under the strip translation]\label{9.1:lem-strip-combinatorics}
Let $V^{\sharp}$ satisfy the conclusions of Corollary~\ref{9.1:cor-refiled-smallness} with
$b=(1+2\beta)\kappa$. Here $\beta$ is the letter of the renormalisation step at hand; in the slots
of the tuned family in which the lemma is applied it is the letter $\beta_s$. Assume condition
$(F_\sharp)$ of \eqref{9.1:eq-hull-conventions-a}, which gives
\begin{equation*}
\tfrac12\beta\kappa\ \ge\ \kappa/10\ \ge\ 4(\mathsf{k}_d+1).
\end{equation*}
Assume further:
\begin{itemize}
\item condition (G0) of \eqref{9.1:eq-strip-window-a};
\item the corollary giving the flat form of the bound \cite[(1.89)]{B-CMP122b};
\item the step of the argument of \cite[(1.90)--(1.100)]{B-CMP122b} for a general collar at
which the convergence criterion for polymer expansions of Koteck\'y and Preiss \cite{KP86} is
applied; this step is used here for the collar translated into strips by
Definition~\ref{9.1:def-collar-dictionary};
\item the first counting bound of the large-field expansion.
\end{itemize}
Then the displays \cite[(1.71)--(1.100)]{B-CMP122b} hold under the dictionary of
Definition~\ref{9.1:def-collar-dictionary}. The term $2\kappa(100R_k)^d$, in which $R_k$ is the
radius $\check{R}_k$, enters them as it is displayed in \cite[p.~389]{B-CMP122b}.
\end{lemma}

\begin{proof}
\emph{Conventions.} Throughout the proof, $\alpha^{\sharp}$ is the constant of
Corollary~\ref{9.1:cor-refiled-smallness}. With $b=(1+2\beta)\kappa$ that corollary gives
\begin{equation*}
|V^{\sharp}(Y)|\le\alpha^{\sharp}e^{-(1+2\beta)\kappa\,d_\sharp(Y)}
\qquad\text{whenever } Y\setminus Z^{\sharp}\neq\emptyset .
\end{equation*}
In the bound of \cite[(1.92)]{B-CMP122b}, read under the dictionary, this constant takes the
place of the constant written $\alpha$ there. We write $N(\cdot)$ for the number of cubes of a
union of cubes. By Definition~\ref{9.1:def-collar-dictionary}, $\beta_v=\beta/(2\mathsf{c}_d)$, and hence
\begin{equation*}
\beta_v\mathsf{c}_d=\tfrac12\beta .
\end{equation*}
We go through the places at which the collar enters \cite[(1.71)--(1.72)]{B-CMP122b} and the
argument of \cite[(1.90)--(1.100)]{B-CMP122b}, carried out for a general collar, in the order of
those displays.

\emph{The factorisation, \cite[(1.71)--(1.72)]{B-CMP122b}.} By part~(S1) of the geometry of the
strips (Lemma~\ref{9.1:lem-strip-geometry}), the strip $X^{\sharp}$ about a component $X$ of $Z$
contains exactly one component of $\Lambda_k^c$, namely $X$. Hence $V^{\sharp}(X^{\sharp})$
depends on the integration variables of $X$ only. The factorisation in components of $Z$ used
at \cite[(1.71)]{B-CMP122b} therefore holds.

By parts~(a) and~(c) of the strip resummation (Theorem~\ref{9.1:thm-strip-resummation}; see
\eqref{9.1:eq-strip-resummation-a}), the terms that are equal to one of the components are
exactly the terms of the interior class. Each of them is acted on by a single operation
$\mathbf{T}'_k(X)$.

\emph{The polymer sum, \cite[(1.90)--(1.92)]{B-CMP122b}.} Let $\mathbf{D}$ be the set of
domains $Y$ entering the term considered. Put
\begin{equation*}
X'_0:=\bigcup_{Y\in\mathbf{D}}Y\ \cup\ \bigcup_j X_j^{\sharp},
\end{equation*}
and let $X'$ be formed from it as in \cite[(1.90)]{B-CMP122b}. Here the domains $Y_i$, the
components $X_j$ and the sub-family $X_{j_h}$ are those of \cite[(1.90)--(1.92)]{B-CMP122b}. Let
$Z'=\bigcup_iY_i\cup\bigcup_hX_{j_h}$ be the set displayed in \cite[p.~388]{B-CMP122b}, and
$Z'^{\sharp}$ its translation into strips under Definition~\ref{9.1:def-collar-dictionary}.

The operation $\mathbf{T}'_k(X_j)$ acts on $e^{V^{\sharp}(Y)}-1$ if and only if $X_j\subset Y$.
Among the values $\mathfrak{v}(Y_4)$ that make up $V^{\sharp}(Y)$ there may be some with
$X_j\not\subset Y_4$; on those $\mathbf{T}'_k(X_j)$ acts trivially. By part~(a) of
Theorem~\ref{9.1:thm-strip-resummation}, $V^{\sharp}(Y)$ is decided by the components of $Z$
inside $Y$.

Consequently the bracket of \cite[(1.72)]{B-CMP122b} is a polymer sum of the kind treated in
the argument of \cite[(1.90)--(1.100)]{B-CMP122b} for a general collar. Two polymers are
compatible if and only if their parts
outside $\bigcup_iY_i^{\sharp}$ neither intersect nor touch.

Finally, for $Y\subset X'$ the set $Y\cap Z^{\sharp}$ consists of strips of $Z'$. Therefore
\begin{equation*}
d_\sharp(Y)=d_{Z'^{\sharp}}(Y).
\end{equation*}

\emph{The tree bound, \cite[(1.93)]{B-CMP122b}.} Let $N_h$ be the number of cubes of the box
$X^{\sharp}_{j_h}$, and recall $d_{out}=d_{\bigcup_iY_i^{\sharp}}$. We claim that
\begin{equation}\label{9.1:eq-strip-combinatorics-1}
\begin{aligned}
d_{out}(X')&\le\sum_{Y\in\mathbf{D}}\bigl(d_\sharp(Y)+2\sqrt d\bigr)
+\sum_h\bigl(N_h+2\sqrt d\bigr),\\
N\Bigl(X'\setminus\bigcup_iY_i^{\sharp}\Bigr)
&\le\sum_{Y\in\mathbf{D}}\mathsf{c}_d\bigl(d_\sharp(Y)+1\bigr)+\sum_hN_h .
\end{aligned}
\end{equation}

We first form a connected set out of the following pieces:
\begin{itemize}
\item for each $Y\in\mathbf{D}$, an optimal tree $\mathcal{T}_Y$ for $d_\sharp(Y)$;
\item for each $h$, a Hamiltonian path through the centres of the cubes of $X^{\sharp}_{j_h}$.
This set is a box by part~(S1) of Lemma~\ref{9.1:lem-strip-geometry}, so such a path exists,
and its length is less than $N_h$.
\end{itemize}

On the vertex set consisting of the domains $Y\in\mathbf{D}$ and the indices $h$, define a graph
as follows.
\begin{itemize}
\item $Y\sim Y'$ if a cube of $Y\setminus Z^{\sharp}$ and a cube of $Y'\setminus Z^{\sharp}$
coincide or touch.
\item $Y\sim h$ if a cube of $Y\setminus Z^{\sharp}$ touches $X^{\sharp}_{j_h}$. In particular,
$Y\sim h$ holds whenever $X^{\sharp}_{j_h}\subset Y$: the set $Y$ is connected and the strips do
not touch one another.
\item $h\sim h'$ never holds, by part~(S1) of Lemma~\ref{9.1:lem-strip-geometry}.
\end{itemize}
The set $X'$ is a single modified component in the sense of \cite[(1.90)]{B-CMP122b} if and only
if this graph is connected.

Take a spanning tree of the graph. Along each of its edges, join the two pieces through a common
point of two touching cubes; this costs a path of length at most $2\sqrt d$ inside $X'$ per edge.
A spanning tree has one edge fewer than the graph has vertices; hence the joining paths have
total length at most $2\sqrt d$ for each element of $\mathbf{D}$ plus $2\sqrt d$ for each
index $h$.

The union of the trees, the paths and the joining paths intersects every cube of
\begin{equation*}
X'\setminus\bigcup_iY_i^{\sharp}
=\bigcup_{Y\in\mathbf{D}}\bigl(Y\setminus Z^{\sharp}\bigr)\ \cup\ \bigcup_hX^{\sharp}_{j_h}.
\end{equation*}
It is therefore admissible for $d_{out}(X')$, and the first line of
\eqref{9.1:eq-strip-combinatorics-1} follows. The count of cubes in the second line is
part~(S5)(a) of Lemma~\ref{9.1:lem-strip-geometry}.

We now split $1+2\beta=(1+\beta)+\tfrac12\beta+\tfrac12\beta$ and
$\alpha^{\sharp}=(\alpha^{\sharp})^{1/3}(\alpha^{\sharp})^{2/3}$. We assert
\begin{equation}\label{9.1:eq-strip-combinatorics-a}
\begin{split}
\prod_{Y\in\mathbf{D}}\alpha^{\sharp}e^{-(1+2\beta)\kappa\,d_\sharp(Y)}
&\le e^{-(1+\beta)\kappa\,d_{out}(X')}\,
e^{-\beta_v\kappa\,N(X'\setminus\bigcup_iY_i^{\sharp})}\\
&\quad\times\prod_he^{(1+\beta+\beta_v)\kappa(N_h+2\sqrt d)}
\prod_{Y\in\mathbf{D}}(\alpha^{\sharp})^{2/3}e^{-\frac12\beta\kappa\,d_\sharp(Y)} .
\end{split}
\end{equation}

To prove \eqref{9.1:eq-strip-combinatorics-a}, insert \eqref{9.1:eq-strip-combinatorics-1} into
the first two factors on the right, and use $\beta_v\mathsf{c}_d=\frac12\beta$. The right member
then dominates the product, over $Y\in\mathbf{D}$, of the factors
\begin{equation*}
(\alpha^{\sharp})^{2/3}e^{-(1+2\beta)\kappa\,d_\sharp(Y)}
e^{-(1+\beta)\kappa 2\sqrt d-\frac12\beta\kappa},
\end{equation*}
times the product, over $h$, of the factors $e^{\beta_v\kappa 2\sqrt d}$. The left member has no
factor indexed by $h$, and each factor indexed by $h$ on the right is at least one: the balance
for each $h$ is $-\beta_v\kappa2\sqrt d\le0$. So only the factors indexed by $Y$ need comparing,
and for each $Y$ it suffices that
\begin{equation*}
\tfrac13\log(\alpha^{\sharp})^{-1}\ \ge\ (1+\beta)\kappa2\sqrt d+\tfrac12\beta\kappa .
\end{equation*}
At $d=4$ we have $(1+\beta)\kappa2d=(1+\beta)\kappa2\sqrt d+4(1+\beta)\kappa$, and
$4(1+\beta)\kappa\ge\frac12\beta\kappa$. Hence this condition is implied by the smallness
condition
\begin{equation*}
(\alpha^{\sharp})^{1/3}\le e^{-(1+\beta)\kappa 2d},
\end{equation*}
which is one of the standing smallness conditions on $\alpha^{\sharp}$ of this section.

\emph{The sum over polymers, \cite[(1.94)]{B-CMP122b}.} Every admissible $Y\subset X'$ has a cube
in $Y\setminus Z^{\sharp}$, and $Y\setminus Z^{\sharp}\subset X'\setminus Z'^{\sharp}$. Apply
part~(S5)(c) of Lemma~\ref{9.1:lem-strip-geometry} with the strips as islands, which do not
touch one another by part~(S1), and with rate $\lambda_0:=\frac12\beta\kappa$. This gives
\begin{equation*}
\sum_{Y\subset X'}(\alpha^{\sharp})^{1/3}e^{-\frac12\beta\kappa\,d_\sharp(Y)}
\ \le\ C_{fib}\,(\alpha^{\sharp})^{1/3}N\bigl(X'\setminus Z'^{\sharp}\bigr).
\end{equation*}
The constant $O(1)$ of \cite[(1.94)]{B-CMP122b} is thus $C_{fib}$.

The exponential of \cite[(1.94)]{B-CMP122b} and the last exponential of
\cite[(1.92)]{B-CMP122b} together cost a factor $e^{2C_{fib}(\alpha^{\sharp})^{1/3}N}$, with $N$
the number of cubes just obtained. Since $\bigcup_iY_i^{\sharp}\subset Z'^{\sharp}$, we have
$N(X'\setminus Z'^{\sharp})\le N(X'\setminus\bigcup_iY_i^{\sharp})$. By the smallness condition
$(\alpha^{\sharp})^{1/3}\le\beta_v\kappa/(4C_{fib})$, another of the standing smallness conditions
on $\alpha^{\sharp}$ of this section,
\begin{equation*}
2C_{fib}(\alpha^{\sharp})^{1/3}N\ \le\ \tfrac12\beta_v\kappa\,
N\Bigl(X'\setminus\bigcup_iY_i^{\sharp}\Bigr).
\end{equation*}
This is absorbed by half of the factor $e^{-\beta_v\kappa N(X'\setminus\bigcup_iY_i^{\sharp})}$
in \eqref{9.1:eq-strip-combinatorics-a}. The constant written $c_0$ in \cite[(1.94)]{B-CMP122b}
equals one exactly when $\mathbf{D}=\emptyset$, that is, exactly when $X'=X_j^{\sharp}$.

\emph{The product over boxes, \cite[(1.95)]{B-CMP122b}.} By part~(S1) of
Lemma~\ref{9.1:lem-strip-geometry}, $N_h+2\sqrt d\le(101R)^d$, where the radius $R$ of that part
is $\check{R}_k$, the radius written $R_k$ in \cite[p.~389]{B-CMP122b}. At $d=4$ we have
$\mathsf{c}_d=2\cdot3^4=162$. The letter $\beta$ satisfies $\beta\le\frac14$, a standing
condition of this section; this gives
$\beta_v\le\frac1{1296}$, and $(1.01)^4<1.041$, so
\begin{equation*}
(1+\beta+\beta_v)(1.01)^d<1.26\cdot1.041<2 .
\end{equation*}
Each factor indexed by $h$ in \eqref{9.1:eq-strip-combinatorics-a} is therefore bounded by
\begin{equation*}
e^{(1+\beta+\beta_v)\kappa(N_h+2\sqrt d)}
\le e^{(1+\beta+\beta_v)(1.01)^d\kappa(100\check{R}_k)^d}
\le e^{2\kappa(100\check{R}_k)^d}.
\end{equation*}
Consequently the inequalities of \cite[p.~389]{B-CMP122b} that lead to
\cite[(1.95)]{B-CMP122b} hold as they stand there, with the term $2\kappa(100R_k)^d$ as
displayed; the factor $(1.01)^d$ coming from $(101\check{R}_k)^d=(1.01)^d(100\check{R}_k)^d$ is
absorbed into the coefficient $2$.

The last of these inequalities has the term $-\frac32\,p_0(g_k)$ in its right member, where
$p_0(\cdot)$ is the degree function of the coupling. This agrees with the parenthetical remark
on the same page, which states that for $\beta_0=\frac17$ (here $\beta_0$ is Ba{\l}aban's small
growth number, written $\beta_0^{pr}$ in this book) the inequality means
\begin{equation*}
-\tfrac14p_0(g_k)+1+2\kappa(100R_k)^d\le0 .
\end{equation*}
In that remark the coefficient $\frac14$ of $-p_0(g_k)$ is $2(1+\beta_0)^{-1}-c$ at
$\beta_0=\frac17$, with $c$ the coefficient of $-p_0(g_k)$ in the right member of the
inequality. Since $2\cdot(\frac87)^{-1}=\frac74$, the equation $\frac74-c=\frac14$ gives
$c=\frac32$.

This is also the value the subsequent steps require. Those steps extract the factor
$e^{-p_0(g_k)}$ before the sum; it is the constant written
$c_1=\exp(-p_0(g_k))$ in \cite[(1.100)]{B-CMP122b}. At the value $q=1$ of the exponent written
$q$ in \cite[(1.95)]{B-CMP122b} this factor is available because \cite[(1.95)]{B-CMP122b}
carries $e^{-q\frac32p_0(g_k)}$, of which \cite[(1.96)]{B-CMP122b} spends
$e^{-q\frac12p_0(g_k)}$. Since $\frac32-\frac12=1$, a coefficient smaller than $\frac32$ would not
leave the factor $e^{-p_0(g_k)}$.

With the coefficient $\frac32$ the inequality can be met: its coefficient of $-p_0(g_k)$ is
$2(1+\beta_0^{pr})^{-1}-\frac32$, which is positive exactly when $\beta_0^{pr}<\frac13$, and
Ba{\l}aban's small growth number $\beta_0^{pr}$ (written $\beta_0$ in \cite[p.~389]{B-CMP122b})
satisfies $\beta_0^{pr}\le\frac17<\frac13$.

\emph{The counting of boxes, \cite[(1.96)--(1.98)]{B-CMP122b}.} The objects counted in
\cite[(1.96)]{B-CMP122b} are the boxes $X$, not their strips. The factor
$100^d(MR_k)^{-d}|\cdot|$ therefore stands as it is displayed there. Moreover, by a smallness
condition on the coupling ceiling $\gamma$ in the order of choice
\eqref{9.1:eq-constant-order-a},
\begin{equation*}
e^{-\frac12p_0(g_k)}\,100^d\check{R}_k^{-d}\le\tfrac14\beta_v\kappa .
\end{equation*}
The step at which the convergence criterion of \cite{KP86} is applied is that of the argument of
\cite[(1.90)--(1.100)]{B-CMP122b} for a general collar, used for the collar translated into
strips by Definition~\ref{9.1:def-collar-dictionary}. It applies because the activity constant
$e^{-p_0(g_k)}$ (written $c_1$ in \cite[(1.100)]{B-CMP122b}) does not exceed the constant
$c_{KP}$ of that criterion in the form used at that step.

\emph{The two groups, \cite[(1.99)--(1.100)]{B-CMP122b}.} Following \cite[(1.99)]{B-CMP122b}, the
second group consists of the $X$ disjoint from $(Z_k\cup\bigcup_iY_i)^{\boxplus}$. Each strip lies
in the corresponding collar, by part~(S1) of Lemma~\ref{9.1:lem-strip-geometry}. Hence for $X$ in
the second group $d_{out}(X)=d_k(X)$, and $X$ lies in the domain required in
\cite[(1.100)]{B-CMP122b}. With the activity constant $e^{-p_0(g_k)}$ this gives
\cite[(1.100)]{B-CMP122b}, together with the localisation of \cite[(2.30)]{B-CMP119}.

The condition defining the first group in \cite[(2.41)]{B-CMP119} is $X\cap Z_j^{\sim}\neq\emptyset$,
and \cite[p.~262]{B-CMP119} gives its alternative form. The first group is therefore the one
defined with Ba{\l}aban's collar: at this place $Z_j^{\sim}$ is the collar $Z_j^{\boxplus}$, as
Definition~\ref{9.1:def-collar-dictionary} prescribes.

The first group is assigned to $\mathbf{B}'$. Its size is measured relative to the collar in
either form of that condition. By part~(S2) of Lemma~\ref{9.1:lem-strip-geometry}, measuring
relative to the narrower strip instead of the collar only increases the tree length, and hence
strengthens the bound.

The set $Z'$ of \cite[p.~388]{B-CMP122b}, introduced above, is displayed there without a collar;
under Definition~\ref{9.1:def-collar-dictionary} the tree lengths are taken relative to its
strips $Z'^{\sharp}$.

What the argument above gives is the bound in $d_{out}=d_{\bigcup_iY_i^{\sharp}}$. Apply part~(S2)
of Lemma~\ref{9.1:lem-strip-geometry} with
\begin{equation*}
W=\bigcup_iY_i^{\sharp}\ \subset\ W'=\bigcup_iY_i^{\boxplus}
\qquad\text{and}\qquad
Y=Y'=X .
\end{equation*}
The inclusion $W\subset W'$ is part~(S1), and $X\setminus W'\subset X\setminus W$. This gives
\begin{equation*}
d_{\bigcup_iY_i^{\boxplus}}(X)\le d_{out}(X).
\end{equation*}
Hence the bound obtained implies the bound relative to the collar. That is the format of the
induction hypothesis at \cite[(2.41)--(2.42)]{B-CMP119}, read with
Definition~\ref{9.1:def-collar-dictionary}.
\end{proof}

\begin{corollary}[The modulus bound under strips]\label{9.1:cor-strip-modulus}
Let $\mathfrak{M}_k(X)$ be the majorant of the modulus lemma used in the proof of the
correlation theorem. Let $c_{\mathfrak{M}}$ be the constant of order one that multiplies
$|Z_k\cap X|$ in $\mathfrak{M}_k(X)$. Put
\begin{equation}\label{9.1:eq-strip-modulus-1}
c^{\sharp}_{\mathfrak{M}}(k):=2^d\,c_{\mathfrak{M}}+(2^d-1)\,\mathsf{K}_{cell}(k).
\end{equation}
Then that modulus lemma holds with the value $V(X^{\sim};w)$ replaced by the strip value
$V^{\sharp}(X^{\sharp};w)$, and with $c_{\mathfrak{M}}$ replaced by
$c^{\sharp}_{\mathfrak{M}}(k)$. Both constants lie under the term
$O(1)\log g_k^{-2}\,|Z_k|$ of the large-field bound at level $k$ in Ba{\l}aban's
construction, the bounds \cite[(1.80)--(1.88)]{B-CMP122b} being used as stated there:
\begin{equation}\label{9.1:eq-strip-modulus-2}
c_{\mathfrak{M}}\,|Z_k\cap X|\le O(1)\log g_k^{-2}\,|Z_k|,
\qquad
c^{\sharp}_{\mathfrak{M}}(k)\,|Z_k\cap X|\le O(1)\log g_k^{-2}\,|Z_k| .
\end{equation}
Both bounds hold by condition $(\gamma 4)$ of \eqref{9.1:eq-hull-conventions-a}, where the
quantity $N_0$ read in that condition is proportional to $\log\log g^{-2}$.
\end{corollary}

\begin{proof}
The list of factors in the modulus lemma does not involve the strips. The same holds for
the remaining items of that lemma, apart from the bound on the modulus of the exponential
of the value. These statements therefore carry over unchanged when
$V(X^{\sim};w)$ is replaced by $V^{\sharp}(X^{\sharp};w)$.

It remains to treat the bound on the modulus of the exponential of the value.
For complex $v$ we have $|e^{v}|=e^{\operatorname{Re}v}\le e^{|v|}$.
Corollary~\ref{9.1:cor-refiled-smallness} gives, on the tube of
Corollary~\ref{9.1:cor-refiled-smallness} over $X^{\sharp}$,
\begin{equation*}
|V^{\sharp}(C^{\sharp})|\le\mathsf{K}_{cell}(k)\sum_j\bigl|\Gamma^0_j\cap C^{\sharp}\bigr| .
\end{equation*}
Applied on $X^{\sharp}$, this yields, on that tube,
\begin{equation}\label{9.1:eq-strip-modulus-3}
\bigl|e^{V^{\sharp}(X^{\sharp};w)}\bigr|
\le\exp\Bigl[\mathsf{K}_{cell}(k)\sum_j\bigl|\Gamma^0_j\cap X^{\sharp}\bigr|\Bigr].
\end{equation}

By (S1) of Lemma~\ref{9.1:lem-strip-geometry}, $X^{\sharp}\setminus X\subset\Lambda_k$, and
$\Lambda_k\subset\Omega_k$ for the region $\Omega_k$ of Ba{\l}aban's construction at
level $k$, so that on $X^{\sharp}\setminus X$ the determining set of the earlier steps is the
unit lattice. Item (S1) also gives the cube count $N(C^{\sharp})\le 2^dN(C)$ for each strip
$C^{\sharp}$, $C$ denoting the component of $\Lambda_k^{c}$ it contains. Together these yield
\begin{equation}\label{9.1:eq-strip-modulus-4}
\sum_j\bigl|\Gamma^0_j\cap X^{\sharp}\bigr|
\le\sum_j\bigl|\Gamma^0_j\cap X\bigr|+(2^d-1)\,|X| .
\end{equation}
In \cite{B-CMP122a} the region $Z_k$ is identified with $Z$. Thus $|X|=|Z_k\cap X|$. Inserting
\eqref{9.1:eq-strip-modulus-4} into \eqref{9.1:eq-strip-modulus-3} gives
\begin{equation}\label{9.1:eq-strip-modulus-5}
\bigl|e^{V^{\sharp}(X^{\sharp};w)}\bigr|
\le\exp\Bigl[\mathsf{K}_{cell}(k)\sum_j\bigl|\Gamma^0_j\cap X\bigr|
+(2^d-1)\,\mathsf{K}_{cell}(k)\,|Z_k\cap X|\Bigr].
\end{equation}
The second term in the exponent of \eqref{9.1:eq-strip-modulus-5} is
$(2^d-1)\mathsf{K}_{cell}(k)\,|Z_k\cap X|$: it is the increase by
$(2^d-1)\mathsf{K}_{cell}(k)$ of the constant at $|Z_k\cap X|$ in $\mathfrak{M}_k(X)$
which enters $c^{\sharp}_{\mathfrak{M}}(k)$ in \eqref{9.1:eq-strip-modulus-1}.
Since $c_{\mathfrak{M}}\ge 0$ and $2^d\ge 1$, the constant
$c_{\mathfrak{M}}+(2^d-1)\mathsf{K}_{cell}(k)$ does not exceed $c^{\sharp}_{\mathfrak{M}}(k)$;
hence the bound on the modulus of the exponential of the value holds with
$c^{\sharp}_{\mathfrak{M}}(k)$ in place of $c_{\mathfrak{M}}$.

For \eqref{9.1:eq-strip-modulus-2}: the exponent of the large-field bound at level $k$
contains a term $O(1)\log g_k^{-2}|Z_k|$. The bounds \cite[(1.80)--(1.88)]{B-CMP122b} are
used here as stated there. They are read through the theorem
and corollary bounding the members of the flat family $\Sigma^{\flat}$. In that reading they treat this
$O(1)$ symbolically, through inequalities involving the degree function $p_0(\cdot)$ of the
coupling. Condition $(\gamma 4)$ of \eqref{9.1:eq-hull-conventions-a}, with $N_0$ proportional
to $\log\log g^{-2}$, places both $c_{\mathfrak{M}}|Z_k\cap X|$ and
$c^{\sharp}_{\mathfrak{M}}(k)|Z_k\cap X|$ under that term.
\end{proof}

\begin{remark}[The strip lemmas relative to old live islands]\label{9.1:rem-live-islands}
We work on the run of the uniqueness theorem, at the large-field operation $\mathbf{R}_k$ of
level $k$. Let $\mathbb{I}_k^{-}$ be the set of islands at level $k$, in the sense of the
definition of islands of this chapter and of the geometric lemma on islands that accompanies it.
Let $Z^{\sharp}(k)$ be
the union of the strips at level $k$, as in Lemma~\ref{9.1:lem-strip-geometry}. The union of the
\emph{old live islands} is
\begin{equation}\label{9.1:eq-live-islands-1}
  \mathbb{W}' := \bigcup\bigl\{ I\in\mathbb{I}_k^{-} : I\not\subset Z^{\sharp}(k)\bigr\}.
\end{equation}
These are the closures, enlarged by four layers, of the birth strips of second-class components
created at earlier steps that are still unhealed at level $k$.
Assume the following.
\begin{itemize}
\item[(H)] The re-reading lemma for old live islands of this chapter holds. It re-reads
the strip machinery under three substitutions, of the domains, of the tree lengths and of the
anchor conditions, together with an extension of the connectivity, namely the following.
\begin{itemize}
\item The domains are those of the class $\mathbf{D}_k^{Z,\circ}$ of that lemma, whose members
$Y_4$ have $Y_4\cap\mathbb{W}'$ a union of islands (possibly empty).
\item The tree lengths relative to $Z$ and to $Z^{\sharp}$ are replaced by
$d_{Z\cup\mathbb{W}'}$ and $d_{Z^{\sharp}\cup\mathbb{W}'}$. The output length $d_{out}$ is
replaced by
\begin{equation}\label{9.1:eq-live-islands-2}
  d_{out}^{\circ} := d_{\bigcup_\iota Y_\iota^{\sharp}\cup\mathbb{W}'}.
\end{equation}
\item Every anchor condition of the form ${\cdot}\not\subset Z$, for cubes $\Delta$ and for the
sets $K_v$ alike, is read with $Z\cup\mathbb{W}'$ in place of $Z$.
\item The connectivity off the excluded regions used in \cite[\S3]{B-CMP122b} is extended to
$\mathbb{W}'$, that is, it is taken off $\mathbb{W}'$ as well.
This is possible because no integration variable of the step lives in an island of $\mathbb{W}'$
\cite[\S1]{B-CMP122a}.
\end{itemize}
\item[(I)] The input \eqref{9.1:eq-resummation-input-a} holds with the relative norms,
as follows. The relative value bound, the value part of \eqref{9.1:eq-resummation-input-a}, is
the relative form of
Lemma~\ref{9.1:lem-column-bound}, with $\mathsf{K}^0$ replaced by
$\mathsf{K}^0+K^{\circ}(B_0+2O(1))$, where $B_0$ is the constant of that name in the re-reading
lemma for old live islands. The island-volume bound is unchanged. The relative island-boundary
bound is the output of the site at which the value units of
Lemma~\ref{9.1:lem-column-bound} are taken; that output forms part of the relative
exponentiation statement of this chapter.
\end{itemize}
Then the following statements hold verbatim under these substitutions:
\begin{itemize}
\item (S1)--(S5) of Lemma~\ref{9.1:lem-strip-geometry};
\item parts (a)--(c) of Theorem~\ref{9.1:thm-strip-resummation};
\item Corollary~\ref{9.1:cor-refiled-smallness};
\item Lemma~\ref{9.1:lem-column-bound};
\item items (u1)--(u7) of Lemma~\ref{9.1:lem-strip-combinatorics}, display
\eqref{9.1:eq-strip-combinatorics-a}, with the dictionary of
Definition~\ref{9.1:def-collar-dictionary};
\item Corollary~\ref{9.1:cor-strip-modulus}.
\end{itemize}
They hold with the same constants: $C_{fib}$, $\mathsf{k}_d$, $\alpha^{\sharp}$, $\alpha_*$,
$m_{lo}$, the condition of Corollary~\ref{9.1:cor-refiled-smallness} under which the window of
strip widths is non-empty, $\beta_v$, the
comparison of $2\kappa(100\check{R}_k)^d$ against $-\tfrac32 p_0$, and the Koteck\'y--Preiss
constant $c_{KP}$ \cite{KP86}. Here $X$ denotes a region, not the inverse coupling $g^{-2}$: for
every region
$X$ on which the boundary term $\mathbf{R}'^{(k)}(X)$ of the $\mathbf{R}$-step is defined,
\begin{equation}\label{9.1:eq-live-islands-3}
  \bigl|\mathbf{R}'^{(k)}(X)\bigr| \le O(1)\,c_1\,e^{-\frac98\kappa\, d_{out}^{\circ}(X)},
\end{equation}
with Ba{\l}aban's constant $c_1$ of \cite{B-CMP122b}. The second group of terms, those bounded in
\cite[(1.100)]{B-CMP122b}, contains no island and keeps that bound in the length $d_k$.
\end{remark}

\begin{proof}
By (H), the re-reading lemma for old live islands re-reads the proofs of
Lemma~\ref{9.1:lem-strip-geometry}, Theorem~\ref{9.1:thm-strip-resummation},
Corollary~\ref{9.1:cor-refiled-smallness}, Lemma~\ref{9.1:lem-column-bound},
Lemma~\ref{9.1:lem-strip-combinatorics} and Corollary~\ref{9.1:cor-strip-modulus} locus by locus.
It uses the substitutions listed in (H).

Only the following loci of those proofs read a tree length or an anchor condition:
\begin{itemize}
\item the chain of walls in (S4) of Lemma~\ref{9.1:lem-strip-geometry};
\item the cube $\Delta_*$ in part (b) of Theorem~\ref{9.1:thm-strip-resummation};
\item the set $Y_1$ of Lemma~\ref{9.1:lem-column-bound}, and item (ii) of the bounds
\eqref{9.1:eq-value-cell-constants-a} that supply it;
\item items (u2)--(u4) and (u7) of Lemma~\ref{9.1:lem-strip-combinatorics}, display
\eqref{9.1:eq-strip-combinatorics-a};
\item the Koteck\'y--Preiss step \cite{KP86}.
\end{itemize}
Each of these loci is re-read in that lemma with the substituted length, the substituted anchor
and the extended connectivity. The remaining steps of those proofs read neither a length nor an
anchor, so they carry over unchanged, and with them the constants listed above.

The re-read proofs take the relative input (I). Its value part is the relative value bound,
proved in that lemma with $\mathsf{K}^0$ replaced by $\mathsf{K}^0+K^{\circ}(B_0+2O(1))$;
the island-volume bound is unchanged; the relative island-boundary bound is the output of the
site at which the value units of Lemma~\ref{9.1:lem-column-bound} are taken. The conclusions of
the
re-read statements therefore hold verbatim. In the final output length $d_{out}^{\circ}$ of
\eqref{9.1:eq-live-islands-2}, they give \eqref{9.1:eq-live-islands-3}. The second group of terms
of \cite[(1.100)]{B-CMP122b} contains no island, so its bound is read in $d_k$ as before.

Finally, suppose $\mathbb{W}'=\emptyset$. Then the substitutions are identities, since
$d_{Z\cup\emptyset}=d_Z$, $d_{Z^{\sharp}\cup\emptyset}=d_{Z^{\sharp}}$,
$d_{out}^{\circ}=d_{out}$ and $Z\cup\emptyset=Z$, and the condition on $Y_4\cap\mathbb{W}'$ is
vacuous, so that $\mathbf{D}_k^{Z,\circ}=\mathbf{D}_k^{Z}$. So for labels with no live island
nothing changes.
\end{proof}

\begin{corollary}[Differences of two runs on one scaffold]\label{9.1:cor-run-differences}
Let $\mathsf{A}$ and $\mathsf{B}$ be two runs that share the scaffold at equal depth, in the sense of the
definition of a common scaffold for two compared runs: they have the same
scaffold datum $\pi$, the same large-field region $Z$, the same radius $R_k$, the same cores $K_v$ and the same
strip width $m_{\sharp}$. Let $\mathfrak{v}^{\mathsf{A}}$, $\mathfrak{v}^{\mathsf{B}}$ be their value
families on $\mathbf{D}_k^Z$, and $V^{\sharp,\mathsf{A}}$, $V^{\sharp,\mathsf{B}}$ their re-filed values in
the sense of the strip resummation (Theorem~\ref{9.1:thm-strip-resummation}). Put
$\delta\mathfrak{v}:=\mathfrak{v}^{\mathsf{B}}-\mathfrak{v}^{\mathsf{A}}$. Then, for
$Y\in\mathbf{D}_k^{\sharp}$,
\begin{equation}\label{9.1:eq-run-differences-1}
V^{\sharp,\mathsf{B}}(Y)-V^{\sharp,\mathsf{A}}(Y)
=\sum_{\operatorname{fil}_{\sharp}(Y_4)=Y}\delta\mathfrak{v}(Y_4).
\end{equation}
Let $\lambda$, $\varepsilon$ be complex numbers (in particular, $(\lambda,\varepsilon)$ may lie in the dial
domain $\mathbb{D}^{\star}$), and let $W=\{W(Y_4)\}_{Y_4\in\mathbf{D}_k^Z}$ be a mark, that is, one of the
families $W^{(z)}$ attached to a class $z$ in the first-variation statement.
For the affine array
\begin{equation}\label{9.1:eq-run-differences-2}
\mathfrak{v}^{\langle\lambda,\varepsilon\rangle}
:=(1-\lambda)\,\mathfrak{v}^{\mathsf{A}}+\lambda\,\mathfrak{v}^{\mathsf{B}}+\varepsilon W,
\end{equation}
the re-filed value $V^{\sharp\langle\lambda,\varepsilon\rangle}$ is its $\operatorname{fil}_{\sharp}$-sum:
\begin{equation}\label{9.1:eq-run-differences-3}
V^{\sharp\langle\lambda,\varepsilon\rangle}(Y)
=\sum_{\operatorname{fil}_{\sharp}(Y_4)=Y}\mathfrak{v}^{\langle\lambda,\varepsilon\rangle}(Y_4),
\qquad Y\in\mathbf{D}_k^{\sharp}.
\end{equation}
Consequently:
\begin{enumerate}
\item[(i)] parts (b) and (c) of Theorem~\ref{9.1:thm-strip-resummation} (the bounds
in~\eqref{9.1:eq-strip-resummation-a}) hold for the difference~\eqref{9.1:eq-run-differences-1} and
for the $\operatorname{fil}_{\sharp}$-sum of the mark $W$, with $\|\delta\mathfrak{v}\|_a$,
respectively $\|W\|_a$, in place of $\|\mathfrak{v}\|_a$ and with the same factor
$C_{fib}e^{-\frac12\beta_{\sharp}\kappa(m_{\sharp}-1)}$, which is at most $1$ for
$m_{\sharp}\ge m_{lo}$; in particular, for $Y\setminus Z^{\sharp}\neq\emptyset$,
\begin{equation}\label{9.1:eq-run-differences-4}
\bigl|V^{\sharp,\mathsf{B}}(Y)-V^{\sharp,\mathsf{A}}(Y)\bigr|
\le C_{fib}\,e^{-\frac12\beta_{\sharp}\kappa(m_{\sharp}-1)}\,\|\delta\mathfrak{v}\|_a\,
e^{-b\,d_{\sharp}(Y)};
\end{equation}
\item[(ii)] the balls defined by the bounds (b) and (c) of
Theorem~\ref{9.1:thm-strip-resummation} are convex balls in a weighted sup norm, so that the fourth
clause of the statement on the dial domain $\mathbb{D}^{\star}$ applies without change;
\item[(iii)] $\operatorname{fil}_{\sharp}$ changes the domain to which a term is assigned and never
its anchor: a mark anchored at a cell of the zone $n(z)$ stays anchored there; hence the bounds stated
per unit of $|\Gamma^0_{n(z)}\cap Y|$ in the second and fourth clauses of the zone-resolved bound, the
assignment of terms to zones, and the index $\vartheta^n$ of every zone $n$ are unchanged;
\item[(iv)] no constant in (i)--(iii) depends on the length of the run, on the size parameters $N$ and
$N_0$, on ages of large-field regions, or on the cutoff.
\end{enumerate}
\end{corollary}

\begin{proof}
By part (a) of Theorem~\ref{9.1:thm-strip-resummation}, the map $\mathfrak{v}\mapsto V^{\sharp}$,
$V^{\sharp}(Y)=\sum_{\operatorname{fil}_{\sharp}(Y_4)=Y}\mathfrak{v}(Y_4)$, is linear, and which
$Y_4$ contribute to $V^{\sharp}(Y)$ is decided by the components of $Z$ inside $Y$. The fibres
$\{Y_4:\operatorname{fil}_{\sharp}(Y_4)=Y\}$ are therefore determined by the scaffold alone, and by
hypothesis the scaffold of $\mathsf{A}$ is that of $\mathsf{B}$. So the two re-filed values are sums of
$\mathfrak{v}^{\mathsf{A}}$, respectively $\mathfrak{v}^{\mathsf{B}}$, over the same fibres, and
subtracting the two sums gives~\eqref{9.1:eq-run-differences-1}. In the same way, the map applied to
the array~\eqref{9.1:eq-run-differences-2} sums it over the same fibres, and by linearity
\begin{equation*}
V^{\sharp\langle\lambda,\varepsilon\rangle}
=(1-\lambda)\,V^{\sharp,\mathsf{A}}+\lambda\,V^{\sharp,\mathsf{B}}+\varepsilon\,V^{\sharp}[W],
\end{equation*}
where $V^{\sharp}[W]$ is the $\operatorname{fil}_{\sharp}$-sum of $W$; this
is~\eqref{9.1:eq-run-differences-3}.

(i) Theorem~\ref{9.1:thm-strip-resummation} applies to any family of numbers (or functions, in sup
norm) indexed by $\mathbf{D}_k^Z$. We apply it to $\delta\mathfrak{v}$ and to $W$; by the identities just
proved, the resulting $\operatorname{fil}_{\sharp}$-sums are the difference of re-filed values and
$V^{\sharp}[W]$. The factor in parts (b) and (c) involves only $C_{fib}$, $\beta_{\sharp}$, $\kappa$ and
$m_{\sharp}$, which are common to both runs, so it is the same factor. This
gives~\eqref{9.1:eq-run-differences-4} and its analogues.

(ii) Each bound in (b) and (c) requires a weighted absolute value to be at most a fixed radius for
every $Y$, that is, a weighted supremum to be at most that radius. Such a set is a ball of a sup norm
and is therefore convex. The affine combinations~\eqref{9.1:eq-run-differences-2} are summed over the
same fibres, so the fourth clause of the statement on $\mathbb{D}^{\star}$ is used with no change.

(iii) By~\eqref{9.1:eq-run-differences-1} and~\eqref{9.1:eq-run-differences-3},
$\operatorname{fil}_{\sharp}$ only decides to which $Y$ the term $\mathfrak{v}(Y_4)$ is assigned. The
term itself, and hence the cell at which a mark is anchored, is not altered. Quantities counted per
anchor cell, the zone filing and the indices $\vartheta^n$ are therefore unchanged.

(iv) The constants appearing in (i)--(iii) are $C_{fib}$, $\beta_{\sharp}$, $\kappa$ and $b$ of
Theorem~\ref{9.1:thm-strip-resummation} and the lower end $m_{lo}$ of the window of strip widths
in~\eqref{9.1:eq-strip-window-a}; no further constant is introduced, the scaffold data being common
to both runs. As fixed in Theorem~\ref{9.1:thm-strip-resummation} and
in~\eqref{9.1:eq-strip-window-a}, none of these constants depends on the length of the run, on the
size parameters $N$ and $N_0$, on ages of large-field regions or on the cutoff, which is (iv).
\end{proof}

\begin{remark}
The strip resummation does not change the family of terms being summed: by part (a) of
Theorem~\ref{9.1:thm-strip-resummation}, the last resummation regroups the terms of one exponent
under the same integrals. The right-hand side of \cite[(1.72)]{B-CMP122b}, that is $\mathbf{R}\rho_k$, is
therefore the same function for every strip width. Where no second-class component is involved, the
bracket of \cite[(1.72)]{B-CMP122b}, and hence the sum $\sum_X R'^{(k)}(X)$ of its localised terms
$R'^{(k)}(X)$, is unchanged as well. So are the
frozen whole-lattice totals, the $b_j$ and the $g_k$. What changes with the width is the localisation of
$R'$ and $\mathbf{B}'$. Near a second-class component $Y_i$, the division between
$\mathbf{T}'_k(Y_i)$ and $\mathbf{B}'$ changes too. This is a choice made in fixing the scaffold: on the
collar region of the scaffold the localisation used here is the one fixed with the scaffold.

The re-filed value $V^{\sharp}(Y)$ is a derived quantity and carries no collar. The value
$V^{\sharp}(Y_i^{\sharp})$, kept in a second-class $\mathbf{T}'$ as in \cite[(1.102)]{B-CMP122b}, lives in
$Y_i^{\sharp}\subset Y_i^{\boxplus}$, inside the ten layers that the next step adds, by the statement
fixing the layers added at each step. It is never treated as a term of the small-field expansion.

The cores $K_v$ enter only through the core-inclusion hypothesis of
Corollary~\ref{9.1:cor-refiled-smallness}. The inequality
\begin{equation*}
\Bigl[\bigl(\tfrac12 c_g+2\bigr)d_k(X)-\tfrac32 c_g-3(1+\beta)2d\Bigr]\kappa<0
\qquad\text{for } d_k(X)\le 3
\end{equation*}
(here $d=4$ is the dimension) holds and does not constrain the strip resummation. It holds because,
the constants $c_g$, $\beta$ and $\kappa$ being positive, for $d_k(X)\le 3$ the bracket is at most
\begin{equation*}
\tfrac32 c_g+6-\tfrac32 c_g-24(1+\beta)=6-24(1+\beta)<0.
\end{equation*}
The large-field regions $X$ with $d_k(X)\le 3$ have $K(X)\subset Z^{(5)}$, which is interior in the
sense of~\eqref{9.1:eq-strip-window-a} for $m_{\sharp}\ge 5$.

No statement is made in this section about the rims, relative to $Z^{\sharp}$, of the entries of the
second kind in the list of domains of this chapter. In particular, nothing is asserted about the
collars of these entries relative to the boundaries of their own regions when the strip width is of
order $R_k/2$.
\end{remark}

\begin{definition}[Cell and point types for reads at a leaf]\label{9.1:not-cell-point-types}
The small-field box sets $\boldsymbol{\Omega}(\cdot)$ of a region, with their determining sets, level
functions, distances and uniformity properties, and the two levels of a label (its level at the
reading and its level before the chain of preliminary integrations) are used as fixed, by statement,
in the notation for small-field box sets and in the notation for labels and their levels. We write
$\beta:=\beta_I=1/5$, the number fixed in
Notation~\ref{9.1:not-constant-order}. Until two runs are compared, a single run is fixed.

\emph{The leaf and its box set.} For a leaf $(F,X,j)$, that is, a term $F$ with its region $X$ and
its level $j$, we put
\begin{equation*}
\mathfrak{S}_X:=\boldsymbol{\Omega}(X^{+}),
\end{equation*}
the box set of the hull $X^{+}$ of $X$, formed over the label current at the reading and truncated at
level $j$. We write $\mathfrak{B}$ for the determining set of $\mathfrak{S}_X$,
$\lambda_{\mathfrak{S}_X}$ for its level function, $d_{\mathfrak{S}_X}$ for its distance, which
satisfies $d_{\mathfrak{S}_X}\ge d_0$, where $d_0$ denotes the distance of the notation for the
small-field box sets, and
\begin{equation*}
\ell_{\mathfrak{S}_X}:=L^{\lambda_{\mathfrak{S}_X}}\eta ,
\end{equation*}
where $\eta$ is the spacing of the finest lattice of the fixed run.

\emph{Levels.} At a point $x$, $\iota(x)$ denotes the level of the label at the reading and $n_0(x)$ its
level before the chain of preliminary integrations; the two levels coincide, $\iota(x)=n_0(x)$, except
at the points treated in the production block. The level at $x$ of $F$'s own units is denoted
$n_F(x)$. It satisfies $n_F(x)\le j\wedge n_0(x)$, with equality unless
the label has been healed at $x$ after the step at which $F$ was created. By the first uniformity
property of the small-field box sets,
\begin{equation}\label{9.1:eq-cell-point-types-1}
\lambda_{\mathfrak{S}_X}(x)=j\wedge\iota(x)\ge n_F(x),\qquad x\in X .
\end{equation}

\emph{Weighted norms and smeared profiles.} For a lattice field $A$ (the letter is generic here) and a
cell $y$ of $\mathfrak{B}$, and for a
function $f$ on $\mathfrak{B}$ (in the applications a nonnegative profile), we put
\begin{equation}\label{9.1:eq-cell-point-types-2}
|A|_y:=\ell_{\mathfrak{S}_X}(y)\sup_{\tilde{\Delta}(y)}|A|,
\qquad
f^{\approx}(x):=\sum_{y\in\mathfrak{B}}e^{-2\delta_1 d_{\mathfrak{S}_X}(x,y)}f(y),
\end{equation}
where $\tilde{\Delta}(y)$ is the enlarged cube of the cell $y$ and $\delta_1>0$ is the decay constant of
the notation for the small-field box sets.

\emph{Cell types.} For a cell $y$ of $\mathfrak{B}$ we write $\iota(y)$ for the level of the label on
$y$. A cell $y$ of $\mathfrak{B}$ of level $l$ is of type
\begin{itemize}
\item[(lab)] if $l=\iota(y)\le j$;
\item[(top)] if $l=j<\iota(y)$;
\item[(dist)] if $l<j\wedge\iota(y)$.
\end{itemize}

\emph{Point types.} A point $x$ of $X$ is of type
\begin{itemize}
\item[$(=)$] if $\iota(x)=n_F(x)$;
\item[$(\mathrm{up})$] if $\iota(x)>n_F(x)$, that is, if the label has been raised at $x$ or
$n_0(x)>n_F(x)$.
\end{itemize}
Since $\iota(x)\ge j\wedge\iota(x)\ge n_F(x)$ for every $x\in X$ by
\eqref{9.1:eq-cell-point-types-1}, the two types exhaust $X$.

\emph{Slowness} (sl). Here $\nu$ is a fixed stage of the stage tables and $\nu'$ ranges over the stages
not earlier than $\nu$; $\iota_{\nu'}$ is the label level current at stage $\nu'$,
$\bar u^{\mathrm{loc}}_{\nu'}$ are the local radii and $a^0_{\nu'}$ the shares at stage $\nu'$, and $m$
denotes a level. The slowness property, the bound on the variation of the shares and the estimate
\eqref{9.1:eq-cell-point-types-3} recorded below are the content of the slowness lemma for the share
profiles below; they are stated here for reference. For every $\nu'$, the local radii
$\bar u^{\mathrm{loc}}_{\nu'}$
and the profiles $L^{a\iota_{\nu'}}$ with $|a|\le 3$ are slow, in the sense of the slowness condition
under which the random-walk bounds are stated. This is the slowness under which the mean-value
bound of the random-walk expansion and the regularity estimate of the background-field bounds apply.
The random-walk bounds state that the condition (d1) of those bounds and this slowness condition
persist, and that they persist also for the profiles $L^{am}$.

The shares vary by at most $(4/\beta)2^{|\Delta\iota|}$, where $\Delta\iota$ is the difference between
the label levels at the cells compared. From this, from the conditions (d1) and (d2) under which the
random-walk bounds are stated, and from the floor $(\mathrm{F}^{\flat})^{g}$, one of the lower bounds
imposed after $L$ in the order in which the constants are chosen
(Notation~\ref{9.1:not-constant-order}), that lemma gives, for every profile
\begin{equation*}
\mathsf{f}:=\ell_{\mathfrak{S}_X}\,a^0_{\nu'},\qquad \nu'\ge\nu ,
\end{equation*}
the bound
\begin{equation}\label{9.1:eq-cell-point-types-3}
\mathsf{f}^{\approx}\le C_{sl}^{\approx}\,\mathsf{f},
\qquad
C_{sl}^{\approx}(d,L,M_1):=\frac{8}{\beta}\,c_1\Upsilon ,
\end{equation}
where $c_1$ is the constant of condition (d1) of the random-walk bounds, and
$\Upsilon:=3(1+\beta_0^{pr})^{3}L$, with $\beta_0^{pr}$ Ba{\l}aban's small growth number.
\end{definition}

\begin{lemma}[Slowness of the local radii across zones]\label{9.1:lem-radii-slowness}
We use the objects fixed in Notation~\ref{9.1:not-cell-point-types}:
\begin{itemize}
\item the triple $(F,X,j)$;
\item the box set $\mathfrak{S}=\boldsymbol{\Omega}(X^+)$, taken over the label current at the read and
truncated at $j$, and its determining set $\mathfrak{B}$;
\item the level function $\lambda_{\mathfrak{S}}$, the unit $\ell_{\mathfrak{S}}$ and the distance
$d_{\mathfrak{S}}\ge d$;
\item the label level $\iota$ and $\beta=\beta_I=1/5$.
\end{itemize}

A \emph{zone} of the first family is a level set of the label. A zone of the second family is, inside
one raised level $j$, one of the old layers $\Sigma_m$, $k_0<m\le j$, of \cite[(1.66)]{B-CMP122a}.

Fix a stage $\nu$. For $\nu'\ge\nu$ let $\bar u^0_{\nu'}$ be the local radii at stage $\nu'$ and
$a^0_{\nu'}$ the local radii at stage $\nu'$ weighted by their shares, so that $a^0_{\nu'}$ and
$\bar u^0_{\nu'}$ differ by the share factor. Put
\begin{equation*}
\Upsilon:=3(1+\beta_0^{pr})^3L .
\end{equation*}
For a non-negative function $\mathsf{f}$ on $\mathfrak{B}$ let
\begin{equation*}
\mathsf{f}^{\approx}(x):=\sum_{y\in\mathfrak{B}}e^{-2\delta_1 d_{\mathfrak{S}}(x,y)}\,\mathsf{f}(y)
\end{equation*}
be $\mathsf{f}$ smeared with the weight $e^{-2\delta_1 d_{\mathfrak{S}}}$.

Assume the crossing bound and the summability bound for the distance $d_{\mathfrak{S}}$ on
$\mathfrak{B}$, the shell floor $\delta_1RM_1\ge 16\log L$, and $\delta_1M\ge 4\kappa_1$. Here
$\delta_1$ is the decay rate and $c_1$ the constant of the summability bound, $R$ and $R_1$ are the
shell radii of the crossing bound and of the shells of the label respectively, and $k_0$ is the lower
index of the old layers of \cite[(1.66)]{B-CMP122a}. Then the
following hold.
\begin{itemize}
\item[(a)] For every $\nu'\ge\nu$ the profile $\mathsf{f}:=\ell_{\mathfrak{S}}a^0_{\nu'}$ satisfies,
on $\mathfrak{B}$,
\begin{equation}\label{9.1:eq-radii-slowness-1}
\mathsf{f}^{\approx}\le \frac{8}{\beta}\,c_1\Upsilon\,\mathsf{f}.
\end{equation}
\item[(b)] For every $\nu'\ge\nu$, with $\mathsf{f}$ as in (a) and the share-free profile
$\mathsf{f}_0:=\ell_{\mathfrak{S}}\bar u^0_{\nu'}$, one has, for all $x,y\in\mathfrak{B}$,
\begin{equation}\label{9.1:eq-radii-slowness-a}
\begin{gathered}
\mathsf{f}(y)\le \frac{4}{\beta}\,2\Upsilon\,e^{\delta_1 d_{\mathfrak{S}}(x,y)/4}\,\mathsf{f}(x),\\
\mathsf{f}_0(y)\le 2\Upsilon\,e^{\delta_1 d_{\mathfrak{S}}(x,y)/4}\,\mathsf{f}_0(x).
\end{gathered}
\end{equation}
Hence the quantity $\mathfrak{s}_u=2\ell_{\mathfrak{S}}\bar u^0_\nu$ satisfies
\begin{equation}\label{9.1:eq-radii-slowness-2}
\mathfrak{s}_u(y)\le 2\cdot 2\Upsilon\,e^{\delta_1 d_{\mathfrak{S}}(x,y)/4}\,
(\ell_{\mathfrak{S}}\bar u^0_\nu)(x).
\end{equation}
\item[(c)] On the level sets of the label at levels $\iota>j$ one has $\lambda_{\mathfrak{S}}=j$, and
the crossing bound does not apply to them. These zones have width at least $R_1M_1L^{\iota-j}$ in
units of level $j$. Hence crossing one of them also costs at least $RM_1$ in the distance
$d_{\mathfrak{S}}$.
\end{itemize}
\end{lemma}

\begin{proof}
For $x,y\in\mathfrak{B}$ let $s(x,y)$ be the number of zone boundaries, of either family, between $x$
and $y$. We write $s=s(x,y)$.

\emph{Step 1: ratio.} Across one zone boundary the dimensionless profile
$\ell_{\mathfrak{S}}\bar u^0_{\nu'}$ changes by at most a factor $\Upsilon$. There are two cases.
\begin{itemize}
\item Across the boundary of an old layer the change is by at most the factor $\ell_\circ^2\le L/3$.
Here
$\ell_\circ$ is Ba{\l}aban's integer $L_0$ whose powers enter \cite[(1.64), (1.66)]{B-CMP122a}; it is not
the constant $L_0(N)$ of this book. The inequality $\ell_\circ^2\le L/3$ is taken from
\cite{B-CMP122a}.
\item Across a zone boundary on cells raised since the birth of $F$ the change is by at most the
factor $L$. This
is because the reads are taken over the label current at the read (Notation~\ref{9.1:not-cell-point-types}).
\end{itemize}
In both cases two further factors enter: a numerical factor $c\le 3$ of \cite{B-CMP122a}, and the factor
$(1+\beta_0^{pr})^3$, which covers the coefficients $\alpha$ of \cite[(2.6)]{B-CMP119} and the numerical
coefficient of the bound. Since $\ell_\circ^2\le L/3\le L$, the factor across one boundary is therefore
at most
\begin{equation*}
\max\{\ell_\circ^2,L\}\cdot 3\cdot(1+\beta_0^{pr})^3\le 3(1+\beta_0^{pr})^3L=\Upsilon .
\end{equation*}

The profile $\mathsf{f}$ differs from $\mathsf{f}_0$ by the share factor. The shares vary by at most a
factor $(4/\beta)2^{|\Delta\iota|}$, where $\Delta\iota$ is the change of the label level; this is the
share bound, read on the set $I_{\nu'}$ attached to stage $\nu'$, where the label level $\iota$ takes
the value $j'$ attached to that stage. The share factor is absorbed into
$(4/\beta)2^{s}$, and multiplying
the factors over the $s$ boundaries gives
\begin{equation}\label{9.1:eq-radii-slowness-3}
\mathsf{f}(y)\le\frac{4}{\beta}(2\Upsilon)^{s}\mathsf{f}(x),\qquad
\mathsf{f}_0(y)\le\Upsilon^{s}\mathsf{f}_0(x).
\end{equation}
The second bound carries no share factor.

\emph{Step 2: distance.} The $s$ boundaries between $x$ and $y$ enclose $(s-1)_+$ full zones. We bound
the cost of crossing each.
\begin{itemize}
\item \emph{First family.} A zone of this family costs $d_{\mathfrak{S}}\ge RM_1$ to cross, by the
crossing bound. The level sets of levels $\iota>j$ are not covered by the crossing bound; for them the
same bound is part (c), whose proof at the end uses neither (a) nor (b).
\item \emph{Second family.} A zone of this family is not a shell of $\boldsymbol{\Omega}(X^+)$, so the
crossing bound says nothing about it. Its width is at least $2M\check R_mL^{m-j}$ in units of level $j$,
where $\check R_i$ are Ba{\l}aban's radii ($R_i$ in \cite{B-CMP122a}). Moreover
\begin{equation*}
L^{-(j-m)}\check R_m\ge 1\qquad\text{for }k_0<m\le j\le k .
\end{equation*}
This follows from two facts. First, the integer $N_0$ of \cite[p.~192]{B-CMP122a} is minimal with
$L^{-N_0+1}\check R_{k-N_0+1}=1$. Second, $\check R_i/\check R_{i+1}\in\{1,L\}$; in the words of
\cite{B-CMP122a}, ``in some steps the number $R_k$ decreases by the factor $L^{-1}$''. Hence the width of
such a zone is at least $2M\ge M$. Here $d$ is the distance of \cite{B-CMP96}, scaled to the unit
lattice, and $d_{\mathfrak{S}}\ge d$.
\end{itemize}
Consequently
\begin{equation}\label{9.1:eq-radii-slowness-4}
d_{\mathfrak{S}}(x,y)\ge\mathfrak{d}_{sh}\,(s-1)_+,\qquad \mathfrak{d}_{sh}:=\min\{RM_1,M\}.
\end{equation}

\emph{Step 3: sum, and proof of (a).} Two inequalities are used here and in the proof of (b). First,
$L^4\ge 2\Upsilon$. Indeed, with the value $\beta_0^{pr}=1/7$ of Ba{\l}aban's small growth number one
has $1+\beta_0^{pr}=8/7$, so $2\Upsilon=6(8/7)^3L$, and $L^4\ge 2\Upsilon$ is
the inequality $L^3\ge 6(8/7)^3$. It holds because $6(8/7)^3=3072/343<9$, while $L^3\ge 27$ for the
odd blocking factor $L>1$. Second, from the conditions on $\kappa$ and $\kappa_1$ in the order of
choice of the constants (Notation~\ref{2.3:not-stages-middle}) follow the inequalities
\begin{equation*}
\kappa_1\ge 4L\kappa,\qquad \kappa\ge 10\log L,
\end{equation*}
and with them
\begin{equation*}
4L\kappa\ge 40L\log L\ge\log(2\Upsilon),
\end{equation*}
the last inequality because $2\Upsilon\le L^4$ gives $\log(2\Upsilon)\le 4\log L\le 40L\log L$.

We now show $e^{-\delta_1\mathfrak{d}_{sh}}\le(2\Upsilon)^{-1}$.
\begin{itemize}
\item If $\mathfrak{d}_{sh}=RM_1$: the shell floor gives $e^{-\delta_1RM_1}\le L^{-16}$. Also
$L^{16}\ge 6(1+\beta_0^{pr})^3L=2\Upsilon$, since $L^{16}\ge L^4\ge 2\Upsilon$.
\item If $\mathfrak{d}_{sh}=M$: by hypothesis $\delta_1M\ge 4\kappa_1$. Further $4\kappa_1\ge 16L\kappa$,
and $16L\kappa\ge 4L\kappa\ge\log(2\Upsilon)$ by the two inequalities above.
\end{itemize}
So no new floor arises.

Next we split the weight as
\begin{equation*}
e^{-2\delta_1 d_{\mathfrak{S}}}
=e^{-\delta_1 d_{\mathfrak{S}}}\cdot e^{-\delta_1 d_{\mathfrak{S}}/2}\cdot e^{-\delta_1 d_{\mathfrak{S}}/2},
\end{equation*}
so that the weight is used only once. Each factor is treated as follows.
\begin{itemize}
\item Since $\Upsilon\ge 1$, we have $(2\Upsilon)^{s}\le 2\Upsilon(2\Upsilon)^{(s-1)_+}$.
\item By \eqref{9.1:eq-radii-slowness-4} and the bound just shown, the first factor satisfies
\begin{equation*}
e^{-\delta_1 d_{\mathfrak{S}}(x,y)}(2\Upsilon)^{(s-1)_+}
\le\bigl(e^{-\delta_1\mathfrak{d}_{sh}}\,2\Upsilon\bigr)^{(s-1)_+}\le 1 .
\end{equation*}
\item The middle factor is at most one.
\item The third factor is summed by the summability bound, which gives, for every $x\in\mathfrak{B}$,
\begin{equation*}
\sum_{y\in\mathfrak{B}}e^{-\delta_1 d_{\mathfrak{S}}(x,y)/2}\le c_1 .
\end{equation*}
\end{itemize}
Inserting the first bound of \eqref{9.1:eq-radii-slowness-3} into the definition of
$\mathsf{f}^{\approx}$ therefore gives
\begin{equation*}
\mathsf{f}^{\approx}(x)\le\frac{4}{\beta}\,2\Upsilon\,\mathsf{f}(x)
\sum_{y\in\mathfrak{B}}e^{-\delta_1 d_{\mathfrak{S}}(x,y)/2}
\le\frac{4}{\beta}\,2\Upsilon\,c_1\,\mathsf{f}(x).
\end{equation*}
This is \eqref{9.1:eq-radii-slowness-1}.

\emph{Proof of (b).} Step~1 gives \eqref{9.1:eq-radii-slowness-3} and Step~2 gives
\eqref{9.1:eq-radii-slowness-4}. We now show $e^{-\delta_1\mathfrak{d}_{sh}/4}\le(2\Upsilon)^{-1}$.
\begin{itemize}
\item If $\mathfrak{d}_{sh}=RM_1$: the shell floor gives $\delta_1RM_1/4\ge 4\log L$, and $L^4\ge
2\Upsilon$ was shown at the beginning of Step~3.
\item If $\mathfrak{d}_{sh}=M$: one has
\begin{equation*}
\delta_1M/4\ge\kappa_1\ge 4L\kappa\ge 40L\log L\ge\log(2\Upsilon).
\end{equation*}
The first inequality is the hypothesis $\delta_1M\ge 4\kappa_1$. The middle two follow from the
conditions on $\kappa$ and $\kappa_1$ of the order of choice, as recalled at the beginning of
Step~3. The last was shown at the
beginning of Step~3.
\end{itemize}
Again no new floor arises. Hence, by \eqref{9.1:eq-radii-slowness-4},
\begin{equation*}
(2\Upsilon)^{(s-1)_+}\le e^{\delta_1\mathfrak{d}_{sh}(s-1)_+/4}\le e^{\delta_1 d_{\mathfrak{S}}(x,y)/4}.
\end{equation*}
Using $\Upsilon\ge 1$ this gives
\begin{gather*}
(2\Upsilon)^{s}\le 2\Upsilon\,e^{\delta_1 d_{\mathfrak{S}}(x,y)/4},\\
\Upsilon^{s}\le\Upsilon\,(2\Upsilon)^{(s-1)_+}\le\Upsilon\,e^{\delta_1 d_{\mathfrak{S}}(x,y)/4}
\le 2\Upsilon\,e^{\delta_1 d_{\mathfrak{S}}(x,y)/4}.
\end{gather*}
Inserted into \eqref{9.1:eq-radii-slowness-3}, these bounds give \eqref{9.1:eq-radii-slowness-a}. The
second inequality of \eqref{9.1:eq-radii-slowness-a}, taken with $\nu'=\nu$ and multiplied by $2$, is
\eqref{9.1:eq-radii-slowness-2}, since $\mathfrak{s}_u=2\ell_{\mathfrak{S}}\bar u^0_\nu$.

\emph{Proof of (c).} A zone boundary of level $\iota>j$ is a shell of the current label. By the
paired-shell property of the label, its shells are exactly paired, with margins of $R_1M_1$ units of
their own level. The margin is therefore
\begin{equation*}
R_1M_1\ \text{level-}\iota\text{ units}=R_1M_1L^{\iota-j}\ \text{level-}j\text{ units}
\ge R_1M_1\ \text{level-}j\text{ units}.
\end{equation*}
On $X^+$ the distance $d_{\mathfrak{S}}$ is measured in units of level $\lambda_{\mathfrak{S}}\le j$.
Hence such a crossing costs at least $RM_1$, as used in Step~2.
\end{proof}

\begin{definition}[Stage tables and their room increments]\label{9.1:def-stage-tables}
Each of the spaces of \cite[(1.64)--(1.68)]{B-CMP122a} is given by bounds on seven
quantities. In each of these spaces the seven bounds carry one and the same coefficient,
multiplying the seven entries of \cite[(1.64)--(1.68)]{B-CMP122a}. We call
these quantities the \emph{members}. They are read on pairs
\begin{equation*}
(\mathbf{U},\mathbf{J})=(e^{i\eta A'}U,\mathbf{J}),\qquad U\ \text{real},
\end{equation*}
in the notation of \cite[(1.64)--(1.68)]{B-CMP122a}.

\emph{Free members.} These are the following five quantities:
\begin{equation*}
E_1=|\partial\mathbf{U}-1|,\quad E_3=|\mathbf{J}|,\quad E_5=|\partial U-1|,\quad
E_6=|A'|,\quad E_7=|\nabla_U A'|.
\end{equation*}

\emph{Composite members.} These are $E_2$ and $E_4$. The member $E_2$ is read on the
plaquette variables, and $E_4$ on the current, of the auxiliary configuration that
\cite{B-CMP122a} forms from $\mathbf{U}$ and writes $U_{p,X}(\,\cdot\,)$.

\emph{Tables.} A \emph{table} is a family $\mathsf{T}=(\mathsf{T}^a(x))_{a=1,\dots,7}$ of
position-dependent entries, given in absolute units, that is, as bounds on the members
themselves. For a region $X$ we let
$\mathfrak{U}[\mathsf{T}](X)$ be the set of pairs $(\mathbf{U},\mathbf{J})$ such that
$E_a<\mathsf{T}^a$ on $X$ for $a=1,\dots,7$ and such that $U$ satisfies the condition on the
real configuration over the cubes of $X$ that enters the spaces of
\cite[(1.64)--(1.68)]{B-CMP122a}.

\emph{Stage tables.} The index $\nu$ of a table, called a \emph{stage}, is the index of the
local radii $\bar u_\nu$ recalled below; it is
unrelated to the stages of the order of choice in Notation~\ref{2.3:not-stages-first}. The
table of stage $\nu$ is
\begin{equation}\label{9.1:eq-stage-tables-1}
\mathsf{T}_\nu:=\mathsf{c}_\nu\,\bar u_\nu ,
\end{equation}
where $\mathsf{c}_\nu$ is the coefficient of the space of \cite[(1.64)--(1.68)]{B-CMP122a}
used at stage $\nu$. The local radii
$\bar u^a_\nu$, the shares $a^a_\nu$ and the set $\mathsf{A}(x)$ of raising stages are those
fixed in the definition of the local radii and their shares. The weights $e_6=1$, $e_7=2$
and $e_3=3$ are attached to the members $E_6$, $E_7$ and $E_3$ respectively.

\emph{The ratio of the radii.} Let $\alpha_{0,k}$ and $\alpha_{1,k}$ be the radii at level
$k$ among Ba{\l}aban's auxiliary parameters (Notation~\ref{2.1:not-prefix}). By
\cite[(2.28)]{B-CMP119} the two radii are given by one expression in the level-$k$ coupling
that differs between them only in the constant, $C_0$ or $C_1$, and in the exponent, $q_0$ or
$q_1$. We put $\mathsf{r}_1:=C_1/C_0$. Under the condition $q_0=q_1$ on the exponents,
imposed in the order of choice for the correlation theorem, this number is also the
ratio $\alpha_{1,k}/\alpha_{0,k}$ at every level $k$.

\emph{Room increments.} The room increment of stage $\nu$ is
$\mathrm{rm}_\nu:=\mathsf{T}_\nu-\mathsf{T}_{\nu+1}$. For a stage $n_e>\nu$ we put
\begin{equation}\label{9.1:eq-stage-tables-2}
\mathsf{G}:=\mathsf{T}_\nu-\mathsf{T}_{n_e}=\sum_{\nu\le i<n_e}\mathrm{rm}_i .
\end{equation}

\emph{The summation constant.} With $\beta=1/5$ as in
Notation~\ref{9.1:not-cell-point-types}, the constant of the summation lemma used below
has the value $8+6\beta$; we put
\begin{equation}\label{9.1:eq-stage-tables-3}
C_\Sigma:=8+6\beta=9.2\le 10 .
\end{equation}
\end{definition}

\begin{proof}
Three assertions are contained in the definition.

By \cite[(2.28)]{B-CMP119}, the radius $\alpha_{0,k}$ is given by an expression in the
level-$k$ coupling carrying the constant $C_0$ as a factor and the exponent $q_0$, and the
radius $\alpha_{1,k}$ by the same expression carrying the constant $C_1$ and the exponent
$q_1$. Under the condition $q_0=q_1$ the two expressions differ only in the constant factor,
so that everything else cancels in the quotient and
$\alpha_{1,k}/\alpha_{0,k}=C_1/C_0=\mathsf{r}_1$ for every $k$.

The equality in \eqref{9.1:eq-stage-tables-2} is a telescoping sum. By the definition of the
room increments,
\begin{equation*}
\sum_{i=\nu}^{n_e-1}\mathrm{rm}_i=\sum_{i=\nu}^{n_e-1}\bigl(\mathsf{T}_i-\mathsf{T}_{i+1}\bigr)
=\mathsf{T}_\nu-\mathsf{T}_{n_e}.
\end{equation*}
Here every intermediate table $\mathsf{T}_i$ with $\nu<i<n_e$ occurs once with each sign. The
identity holds entrywise, for each $a$ and each $x$.

Finally, $\beta=1/5$ gives $6\beta=6/5$. Hence
\begin{equation*}
8+6\beta=\tfrac{46}{5}=9.2\le 10,
\end{equation*}
which is \eqref{9.1:eq-stage-tables-3}.
\end{proof}

\begin{lemma}[Each stage returns a room]\label{9.1:lem-stage-room}
Let the stage tables $\mathsf{T}_\nu=\mathsf{c}_\nu\bar u_\nu$, their room increments $\mathrm{rm}_\nu$, the
shares $a_\nu$ and, at each point $x$, the set $\mathsf{A}(x)$ of raising stages be as in
Definition~\ref{9.1:def-stage-tables}. We fix a member of the tables and drop its index from
$\mathsf{T}_\nu$, $\bar u_\nu$ and $a_\nu$. Assume the lower bounds on
the constants fixed in the
stage-summation statement. Then, in every member of the tables and at every point $x$,
\begin{equation}\label{9.1:eq-stage-room-a}
\mathrm{rm}_\nu=a_\nu\quad\text{for }\nu\notin\mathsf{A}(x),\qquad
\mathrm{rm}_\nu\ge\tfrac{1}{80}\,a_\nu\quad\text{for }\nu\in\mathsf{A}(x).
\end{equation}
\end{lemma}

\begin{proof}
Fix a member and a point $x$. Let $k$ be the current level, the upper limit of the sums below, and
let $h$ be the level such that stage $\nu$ sits at level $j:=h+\nu$ and stage $\nu+1$ at level
$j+1$, as in the stage-summation statement. Let $\beta$ be the constant appearing in the
coefficients of \cite[(1.64), (1.66)]{B-CMP122a}.

\emph{Stages off $\mathsf{A}(x)$.} For $\nu\notin\mathsf{A}(x)$ the stage-summation statement gives two
facts. First, the units are the same at stages $\nu$ and $\nu+1$. Second, the coefficient
$\mathsf{c}_\nu$ drops by exactly $\beta\,2^{-|\iota-j-1|}$ from stage $\nu$ to stage $\nu+1$, where
$\iota$ is the level index of the term of $\mathsf{c}_\nu$ singled out in the stage-summation
statement. With unchanged units
and this exact drop of the coefficient, the room increment is the share, $\mathrm{rm}_\nu=a_\nu$.
At the last stage, on the domain $\Omega_k$ of Ba{\l}aban's construction \cite{B-CMP122a}, the
contribution labelled $3$ in the stage-summation statement is lost; the coefficient then drops by
more than $\beta\,2^{-|\iota-j-1|}$, and the room increment at that stage is at least $a_\nu$.

\emph{Stages in $\mathsf{A}(x)$.} Let $\nu\in\mathsf{A}(x)$. By the
stage-summation statement, the coefficients of the tables at stages $\nu$ and $\nu+1$ are both
given by \cite[(1.64)]{B-CMP122a} at $m=j$ or by \cite[(1.66)]{B-CMP122a}, read at the levels $j$
and $j+1$ respectively; see \cite[pp.~190--191]{B-CMP122a}. At stage $\nu$ this coefficient is
\begin{equation}\label{9.1:eq-stage-room-1}
\mathsf{c}_\nu
=1-\beta\sum_{i=h+1}^{j}2^{-(j-i)}-\beta\sum_{q=j+1}^{k}2^{-(q-j)} .
\end{equation}
The two geometric sums in \eqref{9.1:eq-stage-room-1} are
\begin{equation*}
\sum_{i=h+1}^{j}2^{-(j-i)}=2-2^{-(j-h-1)},\qquad
\sum_{q=j+1}^{k}2^{-(q-j)}=1-2^{-(k-j)}.
\end{equation*}
Substituting them into \eqref{9.1:eq-stage-room-1} gives
\begin{equation}\label{9.1:eq-stage-room-2}
\mathsf{c}_\nu
=1-\beta\bigl(2-2^{-(j-h-1)}\bigr)-\beta\bigl(1-2^{-(k-j)}\bigr)\ge 1-3\beta .
\end{equation}

Replacing $j$ by $j+1$ (with $h$ and $k$ fixed) in the equality of \eqref{9.1:eq-stage-room-2}
gives $\mathsf{c}_{\nu+1}$. Subtracting the
two expressions yields
\begin{equation*}
\mathsf{c}_{\nu+1}-\mathsf{c}_\nu
=\beta\bigl(2^{-(k-j)}-2^{-(j-h)}\bigr)\le\beta\,2^{-(k-j)}\le\tfrac{\beta}{2},
\end{equation*}
where the last inequality holds because $j+1\le k$. The local radii of
Definition~\ref{9.1:def-stage-tables} satisfy
$\bar u_{\nu+1}\le\tfrac34\,\bar u_\nu$. Together with \eqref{9.1:eq-stage-room-2} these give
\begin{equation*}
\begin{aligned}
\mathsf{T}_{\nu+1}
&=\mathsf{c}_{\nu+1}\bar u_{\nu+1}
\le\tfrac34\Bigl(\mathsf{c}_\nu+\tfrac{\beta}{2}\Bigr)\bar u_\nu \\
&\le\tfrac34\Bigl(1+\frac{\beta}{2(1-3\beta)}\Bigr)\mathsf{T}_\nu
\le\tfrac{15}{16}\,\mathsf{T}_\nu .
\end{aligned}
\end{equation*}
Hence the room increment satisfies $\mathrm{rm}_\nu\ge\frac1{16}\mathsf{T}_\nu$. Using
\eqref{9.1:eq-stage-room-2} once more,
\begin{equation*}
\mathrm{rm}_\nu\ge\frac{1-3\beta}{16}\,\bar u_\nu\ge\frac{1}{40}\,\bar u_\nu .
\end{equation*}
The last inequality and the final inequality of the preceding chain are each equivalent to
$\beta\le\frac15$. Both inequalities, and the passage from the local radius $\bar u_\nu$ to the
share $a_\nu$, rest on the numerical checks \eqref{9.1:eq-room-checks-a}: by these checks the bound
$\mathrm{rm}_\nu\ge\bar u_\nu/40$ yields $\mathrm{rm}_\nu\ge a_\nu/80$, which is the second bound of
\eqref{9.1:eq-stage-room-a}.
\end{proof}

The lower bound $1-3\beta$ in \eqref{9.1:eq-stage-room-2} cannot be replaced by the smaller value
$1-4\beta$, the infimum of the table coefficients, which is attained only at stages off
$\mathsf{A}(x)$: with $1-4\beta$ in place of $1-3\beta$, the factor
$\frac34\bigl(1+\beta/(2(1-4\beta))\bigr)$ equals $\frac98>1$ at $\beta=\frac15$.

\begin{definition}[The own tube of a stored term and its dilations]\label{9.1:def-own-tube}
Let $F$ be a stored term created at step $j$, either by Ba{\l}aban's operation $\mathbf{R}$ (an
\emph{$\mathbf{R}$-born} term) or by the operation $\mathbf{T}$ (a \emph{$\mathbf{T}$-born} term).
For a point $x$ we write $n_F(x)\le j$ for the level of the layer of $F$ at the point $x$, and we
abbreviate $n_F=n_F(x)$ where the point is fixed. Tables, their members $E_a$, the weights $e_a$
and the domains $\mathfrak{U}[\mathsf{T}]$ are those of Definition~\ref{9.1:def-stage-tables}.
Below, $\alpha_{a,n}$ denotes the radius attached to the member $E_a$ at level $n$, and $\ell_n$
the scale of level $n$ entering the tables.

\emph{Own units.} The \emph{own units} of $F$ form the table $\mathsf{u}_F=(\mathsf{u}_F^a)$ with
entries
\begin{equation*}
\mathsf{u}_F^a:=\alpha_{a,n_F}\,\ell_{n_F}^{-e_a}.
\end{equation*}

\emph{Own coefficient.} For $n\le j$ put
\begin{equation*}
f_{j,n}:=1-\beta\bigl(1-2^{-(j-n)}\bigr),
\end{equation*}
where $\beta$ is the fixed small number entering the table coefficients (it is not the coupling
increment $\beta_{k+1}$).
Let $\theta_S\in\,]0,1[$ be a fixed number. The coefficient $\mathsf{c}_F^+$ is the function of the
layer $n=n_F(x)$ given as follows. If $F$ is $\mathbf{R}$-born, then $\mathsf{c}_F^+=f_{j,n}$. If $F$
is $\mathbf{T}$-born, then on the layers $n<j$ the coefficient $\mathsf{c}_F^+$ is the one supplied
by \cite[(3.43)]{B-CMP119} read layer by layer, namely
\begin{equation*}
\mathsf{c}_F^+=f_{j,n}+\theta_S\,\beta\,2^{-(j-n)},
\end{equation*}
and on its top layer $n=j$ we set $\mathsf{c}_F^+=1+\beta$.

\emph{Own tube.} The \emph{own tube} of $F$ is the domain
\begin{equation*}
\mathfrak{D}_F:=\mathfrak{U}\bigl[\mathsf{c}_F^+\mathsf{u}_F\bigr].
\end{equation*}
It is the tube on which the correlation theorem produces and bounds $F$.

\emph{Dilations.} Put
\begin{equation}\label{9.1:eq-own-tube-a}
\begin{gathered}
\varrho_F(x):=\theta_\rho\,\beta\,2^{-(j-n_F(x))},\qquad
\theta_\rho:=\tfrac14\min\{\theta_S,\,1-\theta_S\},\\
\mathfrak{D}_F[\varsigma]:=\mathfrak{U}\bigl[\bigl(\mathsf{c}_F^+-(1-\varsigma)\varrho_F\bigr)\mathsf{u}_F\bigr],
\qquad \varsigma\in[0,1].
\end{gathered}
\end{equation}
The function $\varrho_F$ is the \emph{room} of $F$, the parameter $\varsigma$ is the \emph{room
fraction}, and $\mathfrak{D}_F[\varsigma]$ is the \emph{dilation} of the own tube at room fraction
$\varsigma$. Setting $\varsigma=1$ in \eqref{9.1:eq-own-tube-a} removes the term in $\varrho_F$, so
$\mathfrak{D}_F[1]=\mathfrak{D}_F$.
\end{definition}

\begin{lemma}[The schedule of tables returns a room]\label{9.1:lem-schedule-room}
Let $F$ be a stored term of level $j$ in the sense of Definition~\ref{9.1:def-own-tube}, produced
either at an $\mathbf{R}$-step ($\mathbf{R}$-born) or at a $T$-step ($T$-born). Let
$\mathsf{c}_F^{+}$ be its own coefficient,
$\mathsf{u}_F$ its own units and $\varrho_F$ its room, as in
Definition~\ref{9.1:def-own-tube} and \eqref{9.1:eq-own-tube-a}, and let
$f_{j,n}=1-\beta(1-2^{-(j-n)})$ be as there. For a level $k$ we write $f_{k,\cdot}\mathsf{u}$
for the table whose member $a$ at a point of layer $n$ is $f_{k,n}\mathsf{u}^a_n$.

Let $\mathsf{T}$ be the table in which a production at a step later than the one that produced
$F$ places, in its containment statement, the points at which it reads $F$; this later step is
$\mathbf{R}_k$ or $T_k$ with $k\ge j$, where $k>j$ if the step is $\mathbf{R}_k$ and $F$ is
$\mathbf{R}$-born. That is:
\begin{itemize}
\item[(a)] at the step $\mathbf{R}_k$, $\mathsf{T}$ is a stage table:
  \begin{itemize}
  \item on the set $Z^{\mathrm{on}}$ on which the stage tables are in force, it is the stage
    table $\mathsf{T}_\nu$ of Definition~\ref{9.1:def-stage-tables}, with
    $\mathsf{T}_\nu\le f_{k,\cdot}\mathsf{u}$; at stage $0$ this is the table of the set
    $\tilde{U}^c_k$ with which the stages begin;
  \item on $\mathfrak{O}:=\Omega_k\cup Z^c$, where $\Omega_k$ is the small-field region of level
    $k$ of Ba{\l}aban's construction, it is the two-region table $\mathsf{T}^{\flat}_\nu$
    of the accounting statement on $\mathfrak{O}$ of this chapter,
    with the constant $c$ of Ba{\l}aban's construction on $\Omega_k$ replaced by $1$;
  \end{itemize}
\item[(b)] at the step $T_k$, $\mathsf{T}$ is obtained from the level-$k$ table
  $f_{k,\cdot}\mathsf{u}$ by using its response allotment to one half only: at layer $n$ the
  difference
  \begin{equation*}
  f_{k,n}-f_{k+1,n}=\beta2^{-(k+1-n)}
  \end{equation*}
  is split into the part $\theta_S\beta2^{-(k+1-n)}$ reserved for $T$-born terms and the
  response allotment $(1-\theta_S)\beta2^{-(k+1-n)}$, of which one half is used, so that the
  coefficient of $\mathsf{T}$ at layer $n$ is
  \begin{align*}
  &f_{k+1,n}+\theta_S\beta2^{-(k+1-n)}+\tfrac12(1-\theta_S)\beta2^{-(k+1-n)}\\
  &\qquad=f_{k,n}-\tfrac12(1-\theta_S)\beta2^{-(k+1-n)}.
  \end{align*}
\end{itemize}
Then at the $(=)$-points of layers $n<j$, at those of layer $n=j$ if $k>j$, and at those of layer
$n=j=k$ when the reading is at $T_k$,
\begin{equation}\label{9.1:eq-schedule-room-1}
\mathsf{T}\le\bigl(\mathsf{c}_F^{+}-\varrho_F\bigr)\mathsf{u}_F .
\end{equation}
\end{lemma}

\begin{proof}
Item (a) uses the two-region table on $\mathfrak{O}$. The computation
below involves only the coefficients $f_{j,n}$. It applies to the two-region table through
the accounting statement on $\mathfrak{O}$ of this chapter. At the $T$-step, the use of one
half of the response allotment in item (b) is the one provided by the accounting statement at the
$T$-step of this chapter, with its constant $K_R^{\sharp}$; no further condition on the coupling
ceiling $\gamma_U$ is needed for it.

Fix a $(=)$-point $x$ of one of the three kinds listed before
\eqref{9.1:eq-schedule-room-1}. Let $n=n_F(x)$ be its layer, and put $s:=j-n\ge0$. At
such a point the own units of Definition~\ref{9.1:def-own-tube} are
$\mathsf{u}^a_F=\alpha_{a,n}\ell_{n}^{-e_a}$, taken at the layer $n=n_F(x)$; these are the units
$\mathsf{u}^a_n$ of layer $n$ in which the table is expressed, so the table and the right side of
\eqref{9.1:eq-schedule-room-1} are both expressed
in the units of layer $n$. By Definition~\ref{9.1:def-own-tube} the room there is
$\varrho_F(x)=\theta_\rho\beta2^{-s}$, with $\theta_\rho=\tfrac14\min\{\theta_S,1-\theta_S\}$.
Hence $\theta_\rho\le\tfrac14\theta_S$, $\theta_\rho\le\tfrac14(1-\theta_S)$ and
$\theta_\rho\le\tfrac18$. It therefore suffices to show that the coefficient of $\mathsf{T}$ lies
below $\mathsf{c}_F^{+}$ by at least $\theta_S\beta2^{-s}$, or by at least $\tfrac12\beta2^{-s}$,
or by at least $\tfrac14(1-\theta_S)\beta2^{-s}$.

For $k\ge j\ge n$ the definition of $f$ gives the identity
\begin{equation}\label{9.1:eq-schedule-room-2}
f_{j,n}-f_{k,n}=\beta\bigl(2^{-(j-n)}-2^{-(k-n)}\bigr)=\beta2^{-s}\bigl(1-2^{-(k-j)}\bigr)\ge0 .
\end{equation}

The coefficient of $\mathsf{T}$ at layer $n$ is compared with $f_{k,n}$ as follows, for every
$F$. On $Z^{\mathrm{on}}$ in item (a) it is at most $f_{k,n}$ by
$\mathsf{T}_\nu\le f_{k,\cdot}\mathsf{u}$. On $\mathfrak{O}$ in item (a), the comparison of the
coefficients $f_{j,n}$, $f_{k,n}$ made below is transferred to $\mathsf{T}^{\flat}_\nu$ by
clause (c) of the accounting statement on $\mathfrak{O}$ of this chapter. In item (b) the
coefficient is at most $f_{k,n}$ because it equals
\begin{align*}
&f_{k+1,n}+\theta_S\beta2^{-(k+1-n)}+\tfrac12(1-\theta_S)\beta2^{-(k+1-n)}\\
&\qquad=f_{k,n}-\tfrac12(1-\theta_S)\beta2^{-(k+1-n)}\le f_{k,n}.
\end{align*}
We treat three cases in turn.

\emph{$T$-born terms.} On layers $n<j$, Definition~\ref{9.1:def-own-tube} gives
$\mathsf{c}_F^{+}=f_{j,n}+\theta_S\beta2^{-s}$. By \eqref{9.1:eq-schedule-room-2},
\begin{equation*}
\mathsf{c}_F^{+}-f_{k,n}=\bigl(f_{j,n}-f_{k,n}\bigr)+\theta_S\beta2^{-s}\ge\theta_S\beta2^{-s}.
\end{equation*}
On the top layer $n=j$ we have $s=0$ and $\mathsf{c}_F^{+}=1+\beta$. Since $f_{k,j}\le1$ and
$\theta_S<1$, this gives $\mathsf{c}_F^{+}-f_{k,j}\ge\beta\ge\theta_S\beta$.

\emph{$\mathbf{R}$-born terms read at a level $k>j$.} Here $\mathsf{c}_F^{+}=f_{j,n}$, and the
coefficient of $\mathsf{T}$ is at most $f_{k,n}$. By \eqref{9.1:eq-schedule-room-2} and $k-j\ge1$,
\begin{equation*}
f_{j,n}-f_{k,n}=\beta2^{-s}\bigl(1-2^{-(k-j)}\bigr)\ge\tfrac12\beta2^{-s}.
\end{equation*}
This covers layers $n<j$ and, with $s=0$, layer $n=j$.

\emph{$\mathbf{R}$-born terms read at the step $T_j$ of their own level.} Again
$\mathsf{c}_F^{+}=f_{j,n}$. By item (b), the table at $T_j$ uses one half of its
response allotment. The difference $f_{j,n}-f_{j+1,n}=\beta2^{-(s+1)}$ splits into the
$T$-born allotment $\theta_S\beta2^{-(s+1)}$ and the response allotment
$(1-\theta_S)\beta2^{-(s+1)}$. Accordingly, the coefficient of $\mathsf{T}$ at layer $n$ is
\begin{equation*}
f_{j+1,n}+\theta_S\beta2^{-(s+1)}+\tfrac12(1-\theta_S)\beta2^{-(s+1)} .
\end{equation*}
The margin is therefore
\begin{align*}
&f_{j,n}-f_{j+1,n}-\theta_S\beta2^{-(s+1)}-\tfrac12(1-\theta_S)\beta2^{-(s+1)}\\
&\qquad=\beta2^{-(s+1)}\bigl(1-\theta_S-\tfrac12(1-\theta_S)\bigr)
=\tfrac14(1-\theta_S)\beta2^{-s},
\end{align*}
which is check (k3) in \eqref{9.1:eq-room-checks-a}. With $s=0$ this is the case $n=j=k$ at
$T_k$.

In each case the margin is at least $\theta_\rho\beta2^{-s}=\varrho_F(x)$, which proves
\eqref{9.1:eq-schedule-room-1}.
\end{proof}

\begin{remark}[Readings not covered by the schedule]
Two readings lie outside Lemma~\ref{9.1:lem-schedule-room}.

The first is a fresh term ($j=k$) read on its top layer at $\mathbf{R}_k$. For this reading the
containment is supplied by the top-layer bound for fresh terms at the $\mathbf{R}$-step of their
own level. A reading of an $\mathbf{R}$-born term at the step $\mathbf{R}_j$ that produced it is
not a reading at a later step either, so the lemma does not apply to it.

The second concerns the chain terms $\mathbb{C}^{(i)}$ left over from the preliminary integrations.
They are placed in their tables by the bound for the left-over chain terms of the preliminary
integrations, not by the computation in the proof above. Their room, also written $\varrho_F$, is
not the one of Definition~\ref{9.1:def-own-tube}; it is fixed in this chapter together with the
accounting statements on $\mathfrak{O}$ and at the $T$-step.
\end{remark}

\begin{definition}[The chart about the entry's real centre and the slice map]\label{9.1:def-slice-chart}
Let the leaf $(F,X,j)$, the region $X^{+}$, the box set $\mathfrak{S}=\boldsymbol{\Omega}(X^{+})$ and the scale
$\eta$ be as in Notation~\ref{9.1:not-cell-point-types}. We write $G=\mathrm{SU}(N)$ and $G^{c}$ for its
complexification. A $G$-gauge transformation is a gauge transformation with values in $G$, and a
$G^{c}$-gauge transformation is one with values in $G^{c}$.

\emph{The real centre.} The \emph{entry} of a shape is its last diagonal object at real centre data.
Let $U_c$ denote the real centre of the entry of the shape. By the centre hypothesis of the theorem on
charts of shapes, along a chain the centre is the real minimiser of the seed and does not change.
Consequently one real field $U_c$ serves every point of the shape. Its $G$-part is one and the same at
every point, by the statement fixing the $G$-part of the centre along a chain.

\emph{The chart.} Every point $P$ of the shape has a configuration of the form
\begin{equation*}
\mathbf{U}=e^{i\eta A'}U_c
\end{equation*}
over the single real field $U_c$. This chart is the chart of the theorem on charts of shapes. The field
$A'$ carries all moves of the point, their real parts
included. By the criticality statement for real centres, the restriction $U_c{\upharpoonright}X^{+}$ is
critical for $\mathfrak{S}$.

\emph{The response and the slice map.} Let $J$ be the current map, and let $H$ be the linear operator of
the Landau-gauge response construction, with $H^{\dagger}$ as in that construction; this $H$ is not the
Hamiltonian $H$ of Notation~\ref{0.2:not-glossary}. We put $\mathsf{m}_c:=H^{\dagger}J(U_c)$. Since the base
$U_c{\upharpoonright}X^{+}$ is critical, the sourced-response statement and the tilted naturality
statement give that $\mathsf{m}_c$ is the Landau response of the box minimiser, the quantity denoted
$\mathcal{H}$ in \cite{B-CMP102b} and $\mathsf{W}(B)$ in the theorem on charts of shapes. We define
\begin{equation}\label{9.1:eq-slice-chart-a}
\begin{gathered}
B[\mathbf{U}]:=\mathbf{Q}^{\mathfrak{S}}_{U_c}(\eta A'),\qquad
\mathbf{V}(B):=e^{i\eta\mathfrak{X}(B)}U_c,\qquad
\mathfrak{X}(B):=\mathfrak{X}_{\mathfrak{S}}(B;U_c,\mathsf{m}_c),\\
S(\mathbf{U}):=\bigl(\mathbf{V}(B[\mathbf{U}]),\,J(\mathbf{V}(B[\mathbf{U}]))\bigr){\upharpoonright}X,\\
\tilde F(X;P)=F\bigl(X;S(\mathbf{U})\bigr).
\end{gathered}
\end{equation}
Here $\mathbf{Q}^{\mathfrak{S}}_{U_c}$ is the map of the chart construction on $\mathfrak{S}$ based at
$U_c$, and $\mathfrak{X}_{\mathfrak{S}}(\,\cdot\,;U_c,\mathsf{m}_c)$ is the Landau coordinate map on
$\mathfrak{S}$ at the base $U_c$ with response $\mathsf{m}_c$. The pair $S(\mathbf{U})$ is read in the
decomposition
\begin{equation*}
(U,A')=\bigl(U_c,\mathfrak{X}(B)\bigr),\qquad B=B[\mathbf{U}].
\end{equation*}
We call $S$ the \emph{slice map} and $\tilde F$ the \emph{slice} of $F$.

\emph{The meaning of the last line of \eqref{9.1:eq-slice-chart-a}.} At real points the identity
$\tilde F(X;P)=F(X;S(\mathbf{U}))$ holds modulo $G$-gauge transformations, by the statement at real
points of the theorem on charts of shapes; and the class of $J$ is invariant under $G^{c}$-gauge
transformations, by the gauge-covariance statement for the current. Off the real shell the same line is
the definition of $\tilde F$, that is, $\tilde F:=F\circ S$ there, in accordance with the statement on
slices off the real shell and the centre hypothesis of the theorem on charts of shapes. Thus $\tilde F$
is the slice of $P$ built from the chart. Besides it there is the slice of $P$ built from the data of
the point, and the identity theorem equates these two slices. It does not identify $\tilde F$ with $F$
composed with restriction to $X$.
\end{definition}

\begin{definition}[Admissible chart data and the two-cutoff seminorm]\label{9.1:def-admissible-data}
Let $\mathsf{m}\in\{\mathsf{A},\mathsf{B}\}$ be one of the two runs compared, let $F$ be a stored
term and $F^{\mathsf{m}}$ its instance in the run $\mathsf{m}$, let $X$ be a large-field region,
and let $\varsigma$ be a room fraction. We use the chart of
Definition~\ref{9.1:def-slice-chart}. Its centre $U_c$ is a real, regular configuration, critical
for the box set $\mathfrak{S}=\boldsymbol{\Omega}(X^{+})$ of the region. Its chart points are the
configurations
\[
  \mathbf{V}(B)=e^{i\eta\mathfrak{X}(B)}U_c ,
\]
and $J(\mathbf{V}(B))$ is the current of $\mathbf{V}(B)$. We write $r_c$ for the radius
of the chart and $M_{\mathfrak{B}}$ for the data map, which assigns to a configuration its data.
We use the own tube $\mathfrak{D}^{\mathsf{m}}_F[\varsigma]$ of $F$ in the
run $\mathsf{m}$, dilated by $\varsigma$, as in~\eqref{9.1:eq-own-tube-a} of
Definition~\ref{9.1:def-own-tube}.

\emph{Data.} By the statement on the descent of block data modulo block-centre gauge
transformations, data are taken modulo block-centre gauge transformations with values in the
complexification $G^{c}$ of the gauge group: two data that differ by such a gauge transformation
are identified, and a class of data is still written with the letter of any of its
representatives. By part~(d) of the theorem on the data of chart points (a theorem about the
chart of Definition~\ref{9.1:def-slice-chart}, not part of that definition), for a
complex configuration $\mathbf{U}=e^{i\eta A'}U_c$ on the chart, with chart coordinate
$B[\mathbf{U}]=\mathbf{Q}^{\mathfrak{S}}_{U_c}(\eta A')$ as in
Definition~\ref{9.1:def-slice-chart}, the data of its chart point satisfy
\[
  M_{\mathfrak{B}}\big(\mathbf{V}(B[\mathbf{U}])\big)=(\mathrm{rd}\,\mathbf{U})^{\mathsf{g}},
  \qquad \mathsf{g}=\mathsf{w}[\mathfrak{X}]\,\mathsf{w}[A']^{-1}\in G^{c},
\]
where $\mathfrak{X}=\mathfrak{X}(B[\mathbf{U}])$ is the field of
Definition~\ref{9.1:def-slice-chart}, and $\mathsf{w}[\mathfrak{X}]$, $\mathsf{w}[A']$ are the
$G^{c}$-valued block-centre gauge transformations that the chart associates with the fields
$\mathfrak{X}$ and $A'$ (the bracketed letter $\mathsf{w}[\cdot]$ is not the stop depth
$\mathsf{w}(g)$). The gauge transformation $\mathsf{g}$ is of the size of the chart radius.
Consequently the reading $\mathrm{rd}\,\mathbf{U}$ of $\mathbf{U}$ belongs to the class defined
below only modulo $G^{c}$-valued block-centre gauge transformations.

\emph{Admissible data and the chart form of a term.} We define the two objects of the following
display:
\begin{equation}\label{9.1:eq-admissible-data-a}
\begin{aligned}
  \mathfrak{T}^{\mathsf{m}}[\varsigma](X)
  &:=\Big\{M_{\mathfrak{B}}(\mathbf{V}(B)) :\ |B|\le r_c,\
     \big(\mathbf{V}(B),J(\mathbf{V}(B))\big){\upharpoonright}X
     \in\mathfrak{D}^{\mathsf{m}}_F[\varsigma]\Big\},\\
  \mathsf{S}_j^{+}[F]
  &:=\sup_X e^{\kappa_T d_j(X)}
     \sup\Big\{\big|\hat F^{\mathsf{B}}(X;\mathsf{v}\oplus c)-\hat F^{\mathsf{A}}(X;\mathsf{v})\big| :\
     \mathsf{v}\in\mathfrak{T}^{\mathsf{A}}[\tfrac78](X),\\
  &\hphantom{:=\sup_X e^{\kappa_T d_j(X)}\sup\Big\{}\
     \mathsf{v}\oplus c\in\mathfrak{T}^{\mathsf{B}}[\tfrac78](X)\Big\}.
\end{aligned}
\end{equation}
The elements of $\mathfrak{T}^{\mathsf{m}}[\varsigma](X)$ are the \emph{admissible data} of the run
$\mathsf{m}$ at room fraction $\varsigma$ on $X$; they are classes of data modulo $G^{c}$-valued
block-centre gauge transformations, in the sense above.

On these data the term is read through the chart: $\hat F^{\mathsf{m}}(X;\cdot)$ is the function
$F^{\mathsf{m}}$ there, typed on the chart as the function of the chart coordinate $B$, for $B$ as
in the first line of~\eqref{9.1:eq-admissible-data-a}:
\begin{equation}\label{9.1:eq-admissible-data-1}
  \hat F^{\mathsf{m}}(X;B)
  :=F^{\mathsf{m}}\Big(X;\big(\mathbf{V}(B),J(\mathbf{V}(B))\big){\upharpoonright}X\Big).
\end{equation}
In the second line of~\eqref{9.1:eq-admissible-data-a}, $\hat F^{\mathsf{m}}(X;\mathsf{v})$ denotes
the value~\eqref{9.1:eq-admissible-data-1} at a chart point $B$ whose data
$M_{\mathfrak{B}}(\mathbf{V}(B))$ represent $\mathsf{v}$.

\emph{The two-cutoff seminorm.} The quantity $\mathsf{S}_j^{+}[F]$ of the second line
of~\eqref{9.1:eq-admissible-data-a} is the \emph{two-cutoff seminorm} of $F$ at level $j$. In it:
\begin{itemize}
  \item the outer supremum runs over large-field regions $X$, and $d_j(X)$ is the tree length of
  $X$ at level $j$;
  \item $\kappa_T$ is the fixed rate with which the tree length $d_j(X)$ enters the weight
  $e^{\kappa_T d_j(X)}$;
  \item the run $\mathsf{A}$ has length $n$ and the run $\mathsf{B}$ length $n+1$; $c$ is the
  additional datum of the longer run $\mathsf{B}$, and $\mathsf{v}\oplus c$, the concatenation of
  $\mathsf{v}$ with $c$, is a datum of $\mathsf{B}$;
  \item the inner supremum runs over all pairs $(\mathsf{v},c)$ subject to the two stated
  membership conditions.
\end{itemize}
\end{definition}

In this subsection $F$ denotes a stored term and $\Phi$ an output of a later step. Each has its own tube $\mathfrak{D}$ with a layer-wise coefficient $\mathsf{c}^+$. The proposition below shows that one pair of dilation fractions cannot make the read of $F$ at every seed of the dilated tube of $\Phi$ lie in the dilated tube of $F$ on layers deep below the level of $F$.

In Proposition~\ref{9.1:prop-dilation-failure} and its proof we use the following conventions. For a stored term or an output $G$ with own tube $\mathfrak{D}_G=\mathfrak{U}[\mathsf{c}_G^+\mathsf{u}_G]$ (Definition~\ref{9.1:def-own-tube}) and for $\theta\in\,]0,1]$, we write $\theta\mathfrak{D}_G:=\mathfrak{U}[\theta\,\mathsf{c}_G^+\mathsf{u}_G]$. This scaling of the radius by a fraction $\theta$ is to be distinguished from the dilations $\mathfrak{D}_G[\varsigma]$ of \eqref{9.1:eq-own-tube-a}, which change the coefficient additively. We write $\theta\mathfrak{D}_G{\upharpoonright}X$ for the restriction of this tube to a region $X$. If the production of an output reads a stored term $F$ of region $X$, and $P$ is a seed of that output, then $\mathrm{rd}\,P$ denotes the configuration on $X$ at which the production evaluates $F$; we call it the read of $F$ at $P$. For the operation $\mathbf{R}_k$ of Ba{\l}aban's step at level $k$, the layers of the region it treats are called zones and are indexed by their levels $m\le k$. The letter $h<k$ denotes the level that separates the zones $m\le h$ of $\mathbf{R}_k$ from the zones $h<m\le k$. On a zone $m\le h$ the decrements of the stage tables accumulate over the levels $i$ with $h<i\le k$ (see the last part of the proof of Proposition~\ref{9.1:prop-dilation-failure}).

\begin{proposition}[Failure of uniform dilation fractions on deep layers]
\label{9.1:prop-dilation-failure}
Let $F$ be a $\mathbf{T}$-born boundary term of level $j$ with region $X$, so that
$X\cap Z_j\neq\emptyset$ \cite[(2.41)]{B-CMP119}. Assume that $X$ has a point of class $(=)$
in a layer $n\le j-3$. Let $\Phi$ be a $\mathbf{T}$-born output of level $k+1>j$ whose
production reads $F$. Then
\begin{equation}\label{9.1:eq-dilation-failure-1}
\bigl(1-\tfrac{\beta}{8}\bigr)\mathfrak{D}_\Phi{\upharpoonright}X\not\subset
\bigl(1-\tfrac{\beta}{2}\bigr)\mathfrak{D}_F .
\end{equation}
Consequently, with seed fraction $1-\beta/8$ and read fraction $1-\beta/2$, two requirements
imposed jointly on $\Phi$ and on $F$ cannot both be met. The first is that the read of $F$ at
every seed of $(1-\beta/8)\mathfrak{D}_\Phi$ lie in $(1-\beta/2)\mathfrak{D}_F$. The second,
the hypothesis made when the two runs are compared on dilated tubes that the seed fraction
exceed the read fraction, holds for these fractions; hence it is the first that fails. The same
failure occurs for the operation $\mathbf{R}_k$, in the
zones $m\le h-2$, when the term read is a $\mathbf{T}$-born term of level $k$.
\end{proposition}

\begin{proof}
Put $s:=j-n$. Then $s\ge 3$.

We first reduce the claim to the restriction map. Let $P$ be a seed of $(1-\beta/8)\mathfrak{D}_\Phi$. A seed is a complex solution. By the nesting statement for complex solutions, its restriction to $X$ is the re-solution from its own averages on $X$, holomorphically in the datum. By the statement that at a seed the read of a stored term is its restriction,
\[
\mathrm{rd}\,P=P{\upharpoonright}X .
\]
It therefore suffices to exhibit a seed $P$ of $(1-\beta/8)\mathfrak{D}_\Phi$ whose restriction $P{\upharpoonright}X$ lies outside $(1-\beta/2)\mathfrak{D}_F$: this gives \eqref{9.1:eq-dilation-failure-1} and, the read at $P$ being this restriction, the failure of the first requirement.

Next we compare coefficients on layer $n$. There all seven members of a table carry one common coefficient times the entries (Definition~\ref{9.1:def-stage-tables}). On layers of level earlier than $k+1$ the units are the layer's own \cite[(2.34)--(2.39)]{B-CMP119}. We make the comparison at the point of class $(=)$ of the hypothesis. By definition of the class $(=)$, at level $k+1$ this point still lies in layer $n$ and this layer is not raised at that point, so $\mathsf{u}^a_\Phi=\mathsf{u}^a_F=\mathsf{u}^a_n$ there. The comparison of the two tubes on layer $n$ is therefore a comparison of their coefficients.

By Definition~\ref{9.1:def-own-tube}, every coefficient $\mathsf{c}^+$ is at least $1-\beta$. The left-hand coefficient, that of $(1-\beta/8)\mathfrak{D}_\Phi$, therefore satisfies
\begin{equation}\label{9.1:eq-dilation-failure-2}
\bigl(1-\tfrac{\beta}{8}\bigr)\mathsf{c}_\Phi^+\ \ge\ \bigl(1-\tfrac{\beta}{8}\bigr)(1-\beta).
\end{equation}
Also by Definition~\ref{9.1:def-own-tube}, the $\mathbf{T}$-born term $F$ has on the layer $n<j$ the coefficient
\[
\mathsf{c}_F^+=1-\beta+(1+\theta_S)\beta 2^{-s},\qquad \theta_S\in\,]0,1[ .
\]
Since $\theta_S<1$, the right-hand coefficient satisfies
\begin{equation}\label{9.1:eq-dilation-failure-3}
\bigl(1-\tfrac{\beta}{2}\bigr)\mathsf{c}_F^+\ <\ \bigl(1-\tfrac{\beta}{2}\bigr)\bigl(1-\beta+2\beta 2^{-s}\bigr).
\end{equation}
Expanding the products shows that the inequality
\[
\bigl(1-\tfrac{\beta}{8}\bigr)(1-\beta)>\bigl(1-\tfrac{\beta}{2}\bigr)\bigl(1-\beta+2\beta2^{-s}\bigr)
\]
is equivalent to $\tfrac{3\beta}{8}(1-\beta)>2\beta\bigl(1-\tfrac{\beta}{2}\bigr)2^{-s}$. This in turn is equivalent to
\begin{equation}\label{9.1:eq-dilation-failure-4}
2^{-s}<\frac{3(1-\beta)}{16(1-\beta/2)} .
\end{equation}
For $\beta=1/5$ the right-hand side of \eqref{9.1:eq-dilation-failure-4} equals $1/6$; this is the first of the numerical checks \eqref{9.1:eq-room-checks-a}. Since $s\ge3$, we have $2^{-s}\le 1/8<1/6$. By \eqref{9.1:eq-dilation-failure-2} and \eqref{9.1:eq-dilation-failure-3}, the left-hand coefficient strictly exceeds the right-hand one on layer $n$.

We now construct the witness. It is a configuration at which the member $E_6$ reaches the left-hand radius on layer $n$, while the members $E_7$, $E_1$, $E_3$ and the healing field are smaller by the reduction factors stated below. In the chart $\mathbf{U}=e^{i\eta A'}U_c$ of Definition~\ref{9.1:def-slice-chart}, take
\[
U_c=1,\qquad A'=i\tau\chi e_\mu,\qquad \tau>0 .
\]
Here $\mu$ is a lattice direction, and $e_\mu$ is a fixed field of unit size in the direction $\mu$ with values in a one-dimensional, hence abelian, subalgebra of the Lie algebra of the gauge group. The function $\chi$ is a plateau function on layer $n$ which varies only on the scale of its width $b_\chi$, where $b_\chi\ge M\check{R}_n\ell_n$ and $\check{R}_n$ is the scaffold radius at level $n$.

On the plateau, $E_6=\tau$. The members $E_7$, $E_1$, $E_3$ and the healing field are smaller, the reduction factors being $\ell_n/b_\chi$, $\mathsf{r}_1\ell_n/b_\chi$ and $(\ell_n/b_\chi)^2$. Since $b_\chi\ge M\check{R}_n\ell_n$, these factors are at most $1/(M\check{R}_n)$, $\mathsf{r}_1/(M\check{R}_n)$ and $1/(M\check{R}_n)^2$; they tend to $0$ as $\gamma\to0$. Hence, for $\gamma$ small, the members of the healed configuration other than $E_6$ lie inside their entries of $(1-\beta/8)\mathfrak{D}_\Phi$ on $X$, and $\tau$ can be taken up to the $E_6$-radius of $(1-\beta/8)\mathfrak{D}_\Phi$ on layer $n$, namely $(1-\beta/8)\mathsf{c}_\Phi^+\mathsf{u}^6_n$, the configuration remaining a seed of $(1-\beta/8)\mathfrak{D}_\Phi$.

By the coefficient comparison, we may choose
\[
\bigl(1-\tfrac{\beta}{2}\bigr)\mathsf{c}_F^+\mathsf{u}^6_n\le\tau<\bigl(1-\tfrac{\beta}{8}\bigr)\mathsf{c}_\Phi^+\mathsf{u}^6_n .
\]
The restriction of this configuration to $X$ then has $E_6=\tau$ on layer $n$. This is not below the $E_6$-entry of $(1-\beta/2)\mathfrak{D}_F$, which gives \eqref{9.1:eq-dilation-failure-1}.

Finally, consider the zones $m\le h$ of $\mathbf{R}_k$, and let $F_k$ be a $\mathbf{T}$-born stored term of level $k$ read by $\mathbf{R}_k$. Write $\tilde u_{k,m}$ for the coefficient, on zone $m$, of the domain $\tilde U^c_k$ of the $\mathbf{R}_k$ step at level $k$. Both comparisons below are made against this coefficient. By the stage tables (Definition~\ref{9.1:def-stage-tables}), the coefficient of an entry of the table $\mathsf{T}_N$ of the last stage $N$ lies below $\tilde u_{k,m}$ by
\[
\beta\sum_{i=h+1}^{k}2^{-(i-m)}=\beta2^{-(h-m)}\bigl(1-2^{-(k-h)}\bigr)<\beta2^{-(h-m)} .
\]
The coefficient of $\mathfrak{D}_{F_k}$ exceeds $\tilde u_{k,m}$ by the amount below, the term $F_k$ being $\mathbf{T}$-born of level $k$:
\[
\theta_S\beta2^{-(k-m)}=\theta_S\beta2^{-(h-m)}2^{-(k-h)} .
\]
The total gap is therefore
\[
\beta2^{-(h-m)}\bigl(1-(1-\theta_S)2^{-(k-h)}\bigr)<\beta2^{-(h-m)} .
\]
For $h-m\ge2$ this is less than $\beta/4$, and $\beta/4<(\beta/2)(1-\beta)$ because $1-\beta>1/2$.

On the other hand, the dilation by $1-\beta/2$ lowers the coefficient of $\mathfrak{D}_{F_k}$ by $(\beta/2)\mathsf{c}_{F_k}^+\ge(\beta/2)(1-\beta)$. Hence in these zones already the undilated entry of $\mathsf{T}_N$ exceeds the coefficient of $(1-\beta/2)\mathfrak{D}_{F_k}$. The factor $1-\beta/8$ carries as well. Write $\mathsf{c}_N$ for the coefficient of the entry of $\mathsf{T}_N$ on zone $m$; by the total gap above, $\mathsf{c}_N>\mathsf{c}_{F_k}^+-\beta/4$ for $h-m\ge2$. Since $\mathsf{c}_{F_k}^+\ge1-\beta$,
\begin{align*}
\bigl(1-\tfrac{\beta}{8}\bigr)\mathsf{c}_N-\bigl(1-\tfrac{\beta}{2}\bigr)\mathsf{c}_{F_k}^+
&>\bigl(1-\tfrac{\beta}{8}\bigr)\bigl(\mathsf{c}_{F_k}^+-\tfrac{\beta}{4}\bigr)
-\bigl(1-\tfrac{\beta}{2}\bigr)\mathsf{c}_{F_k}^+\\
&=\tfrac{3\beta}{8}\,\mathsf{c}_{F_k}^+-\bigl(1-\tfrac{\beta}{8}\bigr)\tfrac{\beta}{4}
\ \ge\ \beta\bigl(\tfrac18-\tfrac{11\beta}{32}\bigr),
\end{align*}
which is positive for $\beta<4/11$, in particular for $\beta=1/5$. The coefficient comparison that gave \eqref{9.1:eq-dilation-failure-1} on layer $n$ thus holds in the zones $m\le h-2$ with $(1-\beta/8)\mathsf{c}_N$ in place of the left-hand coefficient.
\end{proof}

\begin{remark}[Additivity of the room]
\cite[p.~276]{B-CMP119} states:
\begin{quote}
The analyticity domains become smaller after each step, but the difference is very small and
exponentially decreasing in the number of steps \dots\ this difference is still much greater
than bounds on this function.
\end{quote}
In the present notation, the room between successive domains is additive. It is of order $\beta2^{-s}$ at depth $s$ by Definition~\ref{9.1:def-own-tube}, and not a fixed fraction of the radius.

More sharply, on a layer of depth $s$ both coefficients $\mathsf{c}_\Phi^+$ and $\mathsf{c}_F^+$ tend to $1-\beta$ as $s\to\infty$. Let $\theta_{\mathrm{seed}},\theta_{\mathrm{read}}\in\,]0,1[$ be fractions independent of the depth such that $\theta_{\mathrm{seed}}\mathfrak{D}_\Phi{\upharpoonright}X\subset \theta_{\mathrm{read}}\mathfrak{D}_F$ on all such layers. Then the resulting inequality of coefficients forces $\theta_{\mathrm{seed}}\le \theta_{\mathrm{read}}$. The margin needed along the comparison path of the two runs, however, requires $\theta_{\mathrm{seed}}>\theta_{\mathrm{read}}$. No uniform pair of fractions serves both purposes.

The same comparison can be made within one run. {\itshape Let $F$ be stored at level $j$, with coefficient $\mathsf{c}_F^+$ on a layer $m<j$. Let $x$ be central in an $M$-cube of $X$, in the middle of layer $m$, such that at a later step $x$ still lies in layer $m$ and this layer has not been raised. Let $\Phi$ be an output of that step with region $Y\supseteq X^+$ and coefficient $\mathsf{c}_\Phi\ge1-\beta$ there. If $(1-\beta/8)(1-\beta)>(1-\beta/2)\mathsf{c}_F^+$, then some seed of $(1-\beta/8)\mathfrak{D}_\Phi$ has its read of $F$ outside $(1-\beta/2)\mathfrak{D}_F$.}

Sketch of the argument. All seven members carry one prefactor. On layers of an earlier level the units are the layer's own \cite[(2.34)--(2.39)]{B-CMP119}, without improvement of regularity \cite{B-CMP119}. As in the proof above, $\mathrm{rd}\,P=P{\upharpoonright}X$ exactly. Scale a datum supported in the block of $x$ until the first of the seed's members to do so reaches $\theta$ times its radius in $\mathfrak{D}_\Phi$, with $\theta\uparrow1-\beta/8$. The decay bound for the response of the $\mathbf{T}$-step damps the response away from that block by a factor at least $L^8$, against ratios of units at most $L^3$. Let $E_a$ be the member that first reaches $\theta$ times its radius, and let $\mathsf{u}^a_m$ be its unit on layer $m$. At $x$ its size is
\[
\theta\,\mathsf{c}_\Phi\mathsf{u}^a_m\ \ge\ \theta(1-\beta)\mathsf{u}^a_m .
\]
For $\theta$ close to $1-\beta/8$ this exceeds $(1-\beta/2)\mathsf{c}_F^+\mathsf{u}^a_m$, by the hypothesis $(1-\beta/8)(1-\beta)>(1-\beta/2)\mathsf{c}_F^+$. So the read of $F$ at this seed, which is its restriction, lies outside $(1-\beta/2)\mathfrak{D}_F$.

For $\mathbf{R}$-born $F$ and $j-m\ge3$, Definition~\ref{9.1:def-own-tube} gives $\mathsf{c}_F^+=1-\beta(1-2^{-(j-m)})\le1-\tfrac34\beta$. Moreover,
\begin{equation}\label{9.1:eq-dilation-failure-5}
\bigl(1-\tfrac{\beta}{8}\bigr)(1-\beta)-\bigl(1-\tfrac{\beta}{2}\bigr)\bigl(1-\tfrac34\beta\bigr)
=\tfrac{\beta}{8}(1-2\beta)>0
\end{equation}
for every $\beta<1/2$. Hence the hypothesis $(1-\beta/8)(1-\beta)>(1-\beta/2)\mathsf{c}_F^+$ of the statement above holds for every such $\beta$.

Within one run, the counterpart of the proposition is that, for stored terms, the coefficient $1+\beta$ is not available uniformly. On the zones $m$ with $h<m\le k$, and for stored terms of a single scale, the random-walk bound for the stored terms continues to hold.

Suppose instead that the rooms were measured with the yardstick $\mathfrak{s}_u:=2\ell\sup\bar u$, twice the length $\ell$ of a cell times the largest local radius $\bar u$, that is, with the radius itself. Then the comparison path of the two runs, along which the chain terms are bounded by $C_{\mathbb{C}}\vartheta^{n/2}$ with a constant $C_{\mathbb{C}}$ and a rate $\vartheta$, would need
\[
C_{\mathbb{C}}\vartheta^{n/2}\le c\beta2^{-(h-n)}
\]
for a numerical constant $c$, with $n$ the layer, as in the proposition. This inequality is an upper bound on $h-n$ in terms of the constants. The comparison of the two runs in this chapter dispenses with this yardstick, not with the room.
\end{remark}

\begin{lemma}[Local Lipschitz bound for the box response]\label{9.1:lem-box-response-lipschitz}
Let $\mathfrak{S}$ be admissible, with $\ell_{\mathfrak{S}}$ as in
Notation~\ref{9.1:not-cell-point-types}. Let $U_c$ be a real configuration satisfying
\cite[(3.35)]{B-CMP99b} and critical for $\mathfrak{S}$, with $|\mathsf{m}_c|\le\mathfrak{r}_0$,
where $\mathfrak{r}_0$ is the radius admitted for the centre datum $\mathsf{m}_c$ in the
regularity properties of the chart.
Let $\mathfrak{X}(B)=\mathfrak{X}_{\mathfrak{S}}(B;U_c,\mathsf{m}_c)$ and
$\mathbf{V}(B)=e^{i\eta\mathfrak{X}(B)}U_c$ be as in Definition~\ref{9.1:def-slice-chart},
display~\eqref{9.1:eq-slice-chart-a}, and write $J(\mathbf{V})$ for the current of the
configuration $\mathbf{V}$, the left member of the Euler--Lagrange equation of the chart. Recall
that $(\cdot)^{\approx}$ is the smeared majorant of rate $2\delta_1$ fixed with the chart: for a
nonnegative function $f$ on the lattice,
\begin{equation*}
f^{\approx}(y)=\sum_{z}e^{-2\delta_1 d(y,z)}f(z),
\end{equation*}
where $d(y,z)$ is the distance entering the kernel bounds of the chart and $\delta_1>0$ is the
decay rate fixed with the chart.
Let $B^1,B^2$ be complex fields with $\sup|B^i|\le a_{R'}$ for $i=1,2$, let
$\mathsf{b}:=|B^2-B^1|$ (pointwise), and let $E:=\mathfrak{X}(B^2)-\mathfrak{X}(B^1)$. Then each of
\begin{equation*}
\ell_{\mathfrak{S}}|E|,\qquad \ell_{\mathfrak{S}}^2|\nabla E|,\qquad
\ell_{\mathfrak{S}}^3\bigl|J(\mathbf{V}(B^2))-J(\mathbf{V}(B^1))\bigr|
\end{equation*}
is bounded by $B_\natural\,\mathsf{b}^{\approx}$, where $B_\natural=B_\natural(d,L,\theta,M_1,R)$;
in this dependence the argument $d$ is the dimension, and $\theta$ and $R$ are the parameters of
the regularity properties of the chart on which its kernel bounds depend (the letter $R$ here is a
parameter of the chart, not the ball radius of Definition~\ref{1.3:def-curvature-field}).
\end{lemma}

\begin{proof}
Throughout, $\ell$ stands for $\ell_{\mathfrak{S}}$, and $C$ denotes constants depending only on
the quantities listed for $B_\natural$. If $B^1=B^2$, then $E=0$ and the two currents coincide,
so there is nothing to prove; we assume $\mathsf{b}\not\equiv0$, so that $\mathsf{b}^{\approx}>0$
at every point.

\emph{The chart as a fixed point.} We use the chart in the fixed-point form that accompanies its
regularity properties: for $i=1,2$,
\begin{equation}\label{9.1:eq-box-response-lipschitz-1}
A'_i:=\mathsf{A}_0^{-1}\mathfrak{X}(B^i)=H_0B^i-\mathcal{G}_0\,d\mathcal{N}(A'_i),
\end{equation}
where $\mathsf{A}_0$ is the coordinate map of the chart taking the chart field $A'$ to the Landau
coordinate $\mathfrak{X}$, $H_0$ and $\mathcal{G}_0$ are its two
kernels, and $d\mathcal{N}$ is a field which is analytic in $A'$ on the ball of $C^1$-fields. It
depends on $A'$ through local functions and through the kernels $H_0\mathfrak{D}_0$ and $\Theta$,
whose decay rates are at least $8\delta_1$; here $\mathfrak{D}_0$ is the operator of the chart
composed with $H_0$ in the first kernel, whose adjoint $\mathfrak{D}_0^{\dagger}$ enters below, and
$\Theta$ is a kernel of the chart, distinct from the scaffold $\Theta_X$. Put $G:=A'_2-A'_1$ and
\begin{equation*}
n(G)(y):=\sup_{\tilde\Delta(y)}\max\{\ell|G|,\ \ell^2|\nabla G|\},\qquad
f^{(8)}(y):=\sum_{z}e^{-8\delta_1 d(y,z)}f(z),
\end{equation*}
where $\tilde\Delta(y)$ is the cell neighbourhood of $y$ used in the seminorms of
Notation~\ref{9.1:not-cell-point-types}.

\emph{The nonlinear term.} By the hypothesis $\sup|B^i|\le a_{R'}$ and the regularity properties
of the chart, both $A'_1$ and $A'_2$ lie in the ball of $C^1$-fields on which
$d\mathcal{N}$ is analytic. The Cauchy estimates for $d\mathcal{N}$ on this ball, applied to the
difference of its values at $A'_2$ and $A'_1$, give, with $\varepsilon_3$ the smallness parameter
of the chart,
\begin{equation}\label{9.1:eq-box-response-lipschitz-2}
\ell^3\bigl|d\mathcal{N}(A'_2)-d\mathcal{N}(A'_1)\bigr|
\le C(\varepsilon_3+\mathfrak{r}_0)\,n(G)^{(8)}+\bigl|\mathfrak{D}_0^{\dagger}\Theta(B^2-B^1)\bigr|.
\end{equation}
The nonlocal dependence of $d\mathcal{N}$ on $A'$ passes only through $H_0\mathfrak{D}_0$ and
$\Theta$. For this reason the first term of \eqref{9.1:eq-box-response-lipschitz-2} is smeared
with the kernel $e^{-8\delta_1 d(\cdot,\cdot)}$ of rate $8\delta_1$, and not with the kernel
$e^{-2\delta_1 d(\cdot,\cdot)}$ of $\mathsf{b}^{\approx}$. The distinction is necessary. If the
contraction term were smeared at the rate $2\delta_1$ of $\mathsf{b}^{\approx}$, the argument below
would break down, because the convolution of two kernels of rate $2\delta_1$ is not bounded by a
multiple of one such kernel: the ratio grows linearly in the distance. The quotient $\mathsf{N}$
defined below would then be bounded only in terms of the diameter of the lattice. The explicit
dependence of $d\mathcal{N}$ on $B$, which is the last term of
\eqref{9.1:eq-box-response-lipschitz-2}, enters the constant $C_w$ below.

\emph{The kernel bound.} Subtracting the two equations \eqref{9.1:eq-box-response-lipschitz-1}
gives
\begin{equation*}
G=H_0(B^2-B^1)-\mathcal{G}_0\bigl(d\mathcal{N}(A'_2)-d\mathcal{N}(A'_1)\bigr).
\end{equation*}
We apply the kernel bounds for $H_0$ and $\mathcal{G}_0$ stated with the regularity
properties of the chart, with their constant $C_w$, together with \cite[(3.42)]{B-CMP99b}, and
insert \eqref{9.1:eq-box-response-lipschitz-2}. This yields, for every $x$,
\begin{equation}\label{9.1:eq-box-response-lipschitz-3}
n(G)(x)\le C_w\sum_{y}e^{-8\delta_1 d(x,y)}
\bigl[\mathsf{b}+C(\varepsilon_3+\mathfrak{r}_0)\,n(G)^{(8)}\bigr](y).
\end{equation}

\emph{Closing the bound.} Put $\mathsf{N}:=\sup_x n(G)(x)/\mathsf{b}^{\approx}(x)$. This is
finite because the lattice is finite. We use the following kernel-sum estimate for smeared
profiles: kernel sums of $\mathsf{b}^{\approx}$ return $c_1\mathsf{b}^{\approx}$, that is,
\begin{equation*}
\sum_y e^{-8\delta_1 d(x,y)}\mathsf{b}^{\approx}(y)\le c_1\,\mathsf{b}^{\approx}(x),
\end{equation*}
with a constant $c_1$ fixed below. It rests on the triangle inequality
$d(x,z)\le d(x,y)+d(y,z)$, which gives
\begin{equation*}
e^{-8\delta_1 d(x,y)}e^{-2\delta_1 d(y,z)}\le e^{-6\delta_1 d(x,y)}e^{-2\delta_1 d(x,z)};
\end{equation*}
summing over $z$ against $\mathsf{b}(z)$ and then over $y$, one obtains the estimate with any
constant at least $\sup_x\sum_y e^{-6\delta_1 d(x,y)}$, and we let $c_1$ be such a constant.
By the definition of $f^{(8)}$, this estimate reads $(\mathsf{b}^{\approx})^{(8)}\le
c_1\mathsf{b}^{\approx}$. Hence
$n(G)^{(8)}\le\mathsf{N}(\mathsf{b}^{\approx})^{(8)}\le c_1\mathsf{N}\mathsf{b}^{\approx}$. Since
$e^{-8\delta_1 d(x,y)}\le e^{-2\delta_1 d(x,y)}$, we also have
$\sum_y e^{-8\delta_1 d(x,y)}\mathsf{b}(y)\le\mathsf{b}^{\approx}(x)$. We insert these bounds into
\eqref{9.1:eq-box-response-lipschitz-3}, divide by $\mathsf{b}^{\approx}(x)$ and take the
supremum over $x$. This gives
\begin{equation*}
\mathsf{N}\le C_w+C_w c_1^2\,C(\varepsilon_3+\mathfrak{r}_0)\,\mathsf{N}.
\end{equation*}
The ceiling on $\varepsilon_3+\mathfrak{r}_0$ imposed in the regularity properties of the chart
does not depend on the coupling, and it is the one that makes
$C_wc_1^2C(\varepsilon_3+\mathfrak{r}_0)\le\frac12$.
Since $\mathsf{N}<\infty$, it follows that $\mathsf{N}\le2C_w$, that is,
\begin{equation}\label{9.1:eq-box-response-lipschitz-4}
\ell|G|,\ \ell^2|\nabla G|\le n(G)\le 2C_w\,\mathsf{b}^{\approx}.
\end{equation}

\emph{The Landau coordinate.} By \cite[(46), (73)]{B-CMP102b}, with $D_0$ the map appearing
there,
\begin{equation*}
E=G-H_0\bigl(D_0(A'_2)-D_0(A'_1)\bigr).
\end{equation*}
Through this identity, the bounds on $\ell|E|$ and $\ell^2|\nabla E|$ are read off from
\eqref{9.1:eq-box-response-lipschitz-4}.

\emph{The current.} The base $U_c$ is critical, so the parameter $\varsigma$ of the tilted form
of the natural-interpolation lemma vanishes. Indeed, item (c) of the tilted natural-interpolation
lemma gives
$\varsigma=0$ at real $B$, and $\varsigma$ is analytic in $B$ on the polydisc
$\sup|B|\le a_{R'}$; hence $\varsigma$ vanishes identically there. The Euler--Lagrange equation
therefore reads
\begin{equation*}
J(\mathbf{V}(B))=d\mathbf{Q}\bigl(\eta\mathfrak{X}(B)\bigr)^{\dagger}\mu_B,
\qquad
\mu_B(\xi)=d\Sigma(\mathfrak{X})\bigl[d\mathsf{A}_0H_0\delta_\xi\bigr],
\end{equation*}
where $\xi$ runs over the points at which the current is read, $\Sigma$ is the functional of the
chart whose differential enters the Euler--Lagrange equation, and $\delta_\xi$ is the unit field
supported at $\xi$.
The multiplier $\mu_B$ is analytic in $A'$, and its linear part is $(\Theta B)(\xi)$. Writing
$\mu_i:=\mu_{B^i}$, this gives $|\mu_2-\mu_1|\le C\mathsf{b}^{\approx}$ by
\eqref{9.1:eq-box-response-lipschitz-4}, since $\Theta$ has rate at least $8\delta_1$. Finally,
\begin{align*}
J(\mathbf{V}(B^2))-J(\mathbf{V}(B^1))
&=\bigl(d\mathbf{Q}(\eta\mathfrak{X}(B^2))-d\mathbf{Q}(\eta\mathfrak{X}(B^1))\bigr)^{\dagger}\mu_2\\
&\quad+d\mathbf{Q}(\eta\mathfrak{X}(B^1))^{\dagger}(\mu_2-\mu_1).
\end{align*}
The map $d\mathbf{Q}$ is local and analytic. Hence the first term is bounded through the bound on
$E$ just obtained, and the second through $|\mu_2-\mu_1|\le C\mathsf{b}^{\approx}$. This gives the
bound on $\ell^3|J(\mathbf{V}(B^2))-J(\mathbf{V}(B^1))|$. The constant $B_\natural$ is built
from $2C_w$ and the constants $C$ of the last two steps; the constants $C_w$ and $c_1$ come from
the regularity properties of the chart and the kernel-sum estimate and depend only on
$d,L,\theta,M_1,R$, so that $B_\natural=B_\natural(d,L,\theta,M_1,R)$.

The rate $2\delta_1$ of $\mathsf{b}^{\approx}$ equals $\delta_0/16$, the decay rate of the field
members in item (c) of the theorem on the field members of the chart, whose rate is $\delta_0$;
the argument above does not use the hypothesis from which item (c) of the theorem on the field
members of the chart is derived.
\end{proof}

\begin{lemma}[Slice displacement]\label{9.1:lem-slice-displacement}
Let $X\subseteq X^{+}\subseteq\hat X$, and let
\begin{equation*}
U_e=e^{i\eta A_e}U_c\quad\text{on }\hat X
\end{equation*}
be written in the decomposition of Definition~\ref{9.1:def-slice-chart} about the real centre $U_c$.
Let $\mathfrak{B}$ denote the box set $\mathfrak{S}=\boldsymbol{\Omega}(X^{+})$ of $X^{+}$, over which
the chart and the slice map of Definition~\ref{9.1:def-slice-chart} are defined. Assume
that $U_e$ is critical, at complex data about the real centre, for a set of which $\mathfrak{B}$ is a
refinement, in the Landau chart of that set. Examples are a seed $U(W,\cdot)$ with
$W\supseteq X^{+}$, and any object of the second kind produced by the diagonal step of the chain. Let
\begin{equation*}
\mathbf{U}=e^{i\eta\mathcal{Y}}U_e,
\end{equation*}
with $\mathcal{Y}$ complex and otherwise arbitrary. Assume
\begin{equation}\label{9.1:eq-slice-displacement-1}
\ell_{\mathfrak{S}}\,\sup\bigl(|A_e|+|\mathcal{Y}|\bigr)\le c_4,\qquad
|B[\,\cdot\,]|\le a_{R'},\qquad
C_w|A_e|_c\le 2,
\end{equation}
where $B[\,\cdot\,]$ is evaluated at $U_e$ and at $\mathbf{U}$, and the last inequality holds on every
cube $c$. Here $|\cdot|_c$ is the modulus on the cube $c$, that is, its supremum over $c$, and
$|\mathcal{Y}|_{\cdot}$ is the function $c\mapsto|\mathcal{Y}|_c$; $c_4$ is the smallness constant of
the chart of Definition~\ref{9.1:def-slice-chart};
and $C_w$ is the constant of clause~(5) of the classical entry condition
\eqref{9.1:eq-entry-conditions-a}, which bounds the remainder in the first-order expansion of the box
averages in the move used in the proof below. It is unrelated to the constant $C_W(s)$ of clause~(iv)
of Theorem~\ref{3.1:thm-main}.
Let $S$ be the slice map of \eqref{9.1:eq-slice-chart-a}. Then the following hold.
\begin{enumerate}
\item[(a)] $S(U_e)=(U_e,J(U_e)){\upharpoonright}X$. The equality holds as configurations, with the
first component carried by the decomposition $(U_c,A_e)$: its real centre is $U_c$ and its chart field
is $A_e$.
\item[(b)] $S(\mathbf{U})=\bigl(e^{i\eta(A_e+E)}U_c,\;J(U_e)+\mathsf{J}\bigr)$. On $X$ the field
$E$ and the current $\mathsf{J}$ satisfy
\begin{equation}\label{9.1:eq-slice-displacement-2}
\ell_{\mathfrak{S}}|E|,\quad \ell_{\mathfrak{S}}^{2}|\nabla E|,\quad
\ell_{\mathfrak{S}}^{3}|\mathsf{J}|\;\le\;4B_{\natural}\bigl(|\mathcal{Y}|_{\cdot}\bigr)^{\approx}.
\end{equation}
They are holomorphic in every parameter in which $(A_e,\mathcal{Y})$ are holomorphic. The member
$E_1$ of Definition~\ref{9.1:def-stage-tables} moves by at most
\begin{equation*}
2\eta^{2}|\nabla E|\,(1+C\varepsilon_3).
\end{equation*}
Here $\varepsilon_3$ is one of Ba{\l}aban's small parameters and $C$ is a constant depending only on
the parameters of the chart.
\item[(c)] The following composite members of $S(\mathbf{U})$ are its own members $E_1$ and $E_3$:
the members $E_2(p)$ and $E_4(p)$ of Definition~\ref{9.1:def-stage-tables} for $p\le j$, where $j$
is the birth label, and the pairs formed from $S(\mathbf{U})$ that enter
\cite[(iv)]{B-CMP109}.
\end{enumerate}
\end{lemma}

\begin{proof}
\emph{Part} (a). We use the nesting property of the Landau chart: a Landau gauge and a chart point
restrict to every refinement of the set for which they are defined. This property applies in both
cases of the hypothesis. For a seed over $W$, the box set of $W$ truncated to $X^{+}$ equals the box
set $\boldsymbol{\Omega}(X^{+})$, and truncation refines. For an object produced by the diagonal step
of the chain, the label of a later stage is coarser, so $\mathfrak{B}$ again refines its set.

Write $\mathbf{V}_W$ for the chart map of the seed over its box set $\boldsymbol{\Omega}(W)$, with
coordinate $\beta$, and $\mathbf{V}_{\mathfrak{S}}=\mathbf{V}$ for the map of
\eqref{9.1:eq-slice-chart-a}. The nesting property gives, for every complex $\beta$,
\begin{equation}\label{9.1:eq-slice-displacement-3}
\mathbf{V}_W(\beta){\upharpoonright}X^{+}=\mathbf{V}_{\mathfrak{S}}(B'),
\end{equation}
with $B'$ the box averages of the restricted field. Take $\beta$ to be the chart coordinate of $U_e$.
Then
\begin{equation*}
B'=\mathbf{Q}^{\mathfrak{S}}_{U_c}(\eta A_e)=B[U_e],
\end{equation*}
and \eqref{9.1:eq-slice-displacement-3} shows
that $U_e$ restricted to $X^{+}$ is the chart point $\mathbf{V}(B[U_e])$. Its chart field is
$\mathfrak{X}(B[U_e])=A_e$, about the same real centre $U_c$. By the definition of the slice map in
\eqref{9.1:eq-slice-chart-a}, this gives $S(U_e)=(U_e,J(U_e)){\upharpoonright}X$, as configurations and
in the decomposition $(U_c,A_e)$.

\emph{Part} (b). The chart field $A'$ of $\mathbf{U}$ is given by
\begin{equation*}
e^{i\eta A'}=e^{i\eta\mathcal{Y}}e^{i\eta A_e}.
\end{equation*}
We apply the first-order expansion of the box averages in the move at the complex background $U_e$.
It gives that $B[\mathbf{U}]-B[U_e]$ equals $Q(U_e,\eta\mathcal{Y})$ plus a remainder whose modulus
on each cube $c$ is at most $C_w|\mathcal{Y}|_c|A_e|_c$. Here
$Q(U_e,\cdot)$ is the linearisation of the box averages at the background $U_e$.

Next we apply the complex-background clause of the bound for these linearised averages. It gives
$|Q(U_e,\eta\mathcal{Y})(c)|\le 2|\mathcal{Y}|_c$. Combining the two estimates with the third
hypothesis in \eqref{9.1:eq-slice-displacement-1}, we get on every cube $c$
\begin{equation}\label{9.1:eq-slice-displacement-4}
\begin{aligned}
\bigl|B[\mathbf{U}]-B[U_e]\bigr|(c)
&\le \bigl(2+C_w|A_e|_c\bigr)|\mathcal{Y}|_c\\
&\le 4|\mathcal{Y}|_c .
\end{aligned}
\end{equation}
The hypothesis $C_w|A_e|_c\le 2$ enters the proof only at this point. It does not follow from the
condition $\ell_{\mathfrak{S}}\sup(|A_e|+|\mathcal{Y}|)\le c_4$, because $C_w$ is a constant of the
expansion kept symbolic. At every use of the lemma the hypothesis is met by the smallness condition
in clause~(5) of the classical entry condition \eqref{9.1:eq-entry-conditions-a},
\begin{equation*}
C_w\,(C_e+C_\Sigma+3)\,\bar{\bar u}\le 2,
\end{equation*}
where $C_e$ and $C_\Sigma$ are the
constants of that condition and $\bar{\bar u}$ is the global supremum of the local radii, since $A_e$
lies in the tube fixed there.

We now apply Lemma~\ref{9.1:lem-box-response-lipschitz} with
\begin{equation*}
B^{1}=B[U_e],\qquad B^{2}=B[\mathbf{U}].
\end{equation*}
Both satisfy $\sup|B^{i}|\le a_{R'}$ by the second hypothesis in
\eqref{9.1:eq-slice-displacement-1}. The remaining hypotheses of that lemma, namely that
$\mathfrak{S}$ is admissible, that $U_c$ is real, regular and critical for $\mathfrak{S}$, and the
bound on $\mathsf{m}_c$, are those under which the chart of Definition~\ref{9.1:def-slice-chart} is
set up. By \eqref{9.1:eq-slice-displacement-4},
\begin{equation*}
\mathsf{b}=|B^{2}-B^{1}|\le 4|\mathcal{Y}|_{\cdot}.
\end{equation*}
By part (a),
\begin{equation*}
\mathfrak{X}(B^{1})=A_e,\qquad \mathbf{V}(B^{1})=U_e{\upharpoonright}X^{+},\qquad
J(\mathbf{V}(B^{1}))=J(U_e){\upharpoonright}X^{+}.
\end{equation*}
Define
\begin{equation*}
E:=\mathfrak{X}(B[\mathbf{U}])-A_e,\qquad
\mathsf{J}:=J\bigl(\mathbf{V}(B[\mathbf{U}])\bigr)-J(U_e){\upharpoonright}X^{+}.
\end{equation*}
Then $\mathbf{V}(B[\mathbf{U}])=e^{i\eta(A_e+E)}U_c$, and the slice map gives the representation of
$S(\mathbf{U})$ stated in (b).

Lemma~\ref{9.1:lem-box-response-lipschitz} bounds $\ell_{\mathfrak{S}}|E|$,
$\ell_{\mathfrak{S}}^{2}|\nabla E|$ and $\ell_{\mathfrak{S}}^{3}|\mathsf{J}|$ by
$B_{\natural}\mathsf{b}^{\approx}$. The smearing $f\mapsto f^{\approx}$ is smearing with the positive
kernel $e^{-2\delta_1 d_{\mathfrak{S}}}$, where $\delta_1$ is the decay rate of the kernels of the
chart and $d_{\mathfrak{S}}$ is the
scaled distance attached to the box set $\mathfrak{S}$, so it is monotone and positively homogeneous.
Hence
$B_{\natural}\mathsf{b}^{\approx}\le 4B_{\natural}(|\mathcal{Y}|_{\cdot})^{\approx}$, which is
\eqref{9.1:eq-slice-displacement-2}.

\emph{Part} (c). By the construction of the slice map, $S(\mathbf{U})$ is critical for $\mathfrak{B}$.
On $X$ the label $\lambda_{\mathfrak{S}}$ of the box set $\mathfrak{S}$ equals $j\wedge\iota$, the
minimum of the birth label $j$ and the index $\iota$ of the labelling of box sets on $X$.
By the monotonicity of the
labels, the following sets refine $\mathfrak{B}$ on $X$:
\begin{itemize}
\item for $p\le j$, the box set $\mathbb{B}_p(X)$ of level $p$ on $X$ together with the sets
$\Gamma_i$, $i<p$, attached to the earlier steps $i$, that is, $\mathbb{B}_p(X)\cup\{\Gamma_i\}_{i<p}$;
here $\Gamma_i$ is unrelated to the walk operator $\Gamma$;
\item the maximal sequences of \cite{B-CMP109} on $X$.
\end{itemize}

The second half of the nesting property, applied to each of these refinements, shows the following.
The configurations from which $E_2(p)$, $E_4(p)$ and the pairs of $S(\mathbf{U})$ entering
\cite[(iv)]{B-CMP109} are
computed are $S(\mathbf{U})$ itself. So these members are its own $E_1$ and $E_3$.

It remains to compare thresholds. Since the threshold of $E_4(p)$ in the stage table of
Definition~\ref{9.1:def-stage-tables} is at least that of $E_3$, the bound on $E_3$ of
$S(\mathbf{U})$ implies the one required of $E_4(p)$.
The identification holds up to the boundary of $X$, because $X$ has its margin inside $X^{+}$.
\end{proof}

\begin{remark}
The criticality of $U_c{\upharpoonright}X^{+}$ for $\mathfrak{S}$ invoked in
Definition~\ref{9.1:def-slice-chart} is the case $A_e=0$ of part (a), obtained by the same nesting
argument.
Part (a) can also be seen as follows. The configuration $U_e$ is critical for a coarser set; the
refinement $\mathfrak{B}$ imposes more constraints, so $U_e$ is also critical for $\mathfrak{B}$, and
the uniqueness part of \cite[Theorem~1]{B-CMP102b} identifies $U_e{\upharpoonright}X^{+}$ with the
chart point $\mathbf{V}(B[U_e])$. In part (c), since labels only rise along the chain,
$\lambda_{\mathfrak{S}}$ does not exceed the label of the entry, and the refinements used there hold
at the points of class $(=)$ and at those of class $(\mathrm{up})$.
In part (b), the difference $E-\mathcal{Y}$ between the displacement of the chart field and the move is
never estimated by itself. Only $E$ and $\mathsf{J}$ enter, and their bounds
\eqref{9.1:eq-slice-displacement-2} carry the profile $(|\mathcal{Y}|_{\cdot})^{\approx}$ of the move
rather than a constant in a sup norm independent of $g$.
The current slot $\mathbf{J}$ of the chart point $P=(\mathbf{U},\mathbf{J})$ is not used here. At values
$\mathsf{s}\neq\mathbb{1}$ of the interpolation parameter, where bounds of this kind do not hold for the
operator $D^{*}D$, the covariant Laplacian of the background, no quantity of second order is required.
\end{remark}

\begin{definition}[The classical and single-run entry conditions]\label{9.1:def-entry-conditions}
We use the conventions of Notation~\ref{9.1:not-cell-point-types}, Definition~\ref{9.1:def-stage-tables}
and Definition~\ref{9.1:def-slice-chart}. Both conditions below concern one run at a time. The
\emph{entry condition} of a shape, with complex parameter polydisc $\mathfrak{P}_{\mathsf{c}}$, is
the conjunction of the classical entry condition and the single-run entry condition.

\emph{The classical entry condition.} Its radii do not depend on the coupling. It requires, on
$4\mathfrak{P}_{\mathsf{c}}$, the polydisc with the same centre as $\mathfrak{P}_{\mathsf{c}}$ and
four times its radii, the entry data included:
\begin{itemize}
\item[(1)] all compositions of moves and restrictions are defined and holomorphic, on every box;
\item[(2)] every point has the form $(e^{i\eta A'}U_c,\mathbf{J})$, with $U_c$ the real centre of
Definition~\ref{9.1:def-slice-chart}, and
\begin{equation}\label{9.1:eq-entry-conditions-1}
\ell|A'|,\quad \ell^2|\nabla A'|,\quad \ell^{2+\beta_H}[\nabla A']_{\beta_H},\quad \ell^3|\mathbf{J}|
\;\le\;\mathfrak{r}^{cl}:=\tfrac14\min\{\mathfrak{r}_0,\,a_{R'},\,c_4,\,r_c\};
\end{equation}
these are the bounds required by the domain $\mathcal{D}_{op}$ of the operator statements and by
the condition accompanying it, and by the two local inverse-function conditions of this chapter;
here $\ell$ denotes the local length $\ell_{\mathfrak{S}}$ at the point read, as in
Notation~\ref{9.1:not-cell-point-types},
$\beta_H$ is the H\"older exponent of the seminorm $[\,\cdot\,]_{\beta_H}$,
and $\mathfrak{r}_0$, $a_{R'}$, $c_4$, $r_c$ are the radii of these conditions;
\item[(3)] the move fields $h^{\mathsf{m}}_\mu$ and the currents $\mathcal{J}_\mu$ obey the size
bounds in the first line of \eqref{9.1:eq-entry-conditions-a} below, in which $a^0_\nu$ and
$a^J_\nu$ are shares at stage $\nu$ in the sense of Definition~\ref{9.1:def-stage-tables},
$\ell_\nu$ is the length scale of stage $\nu$, and $\sup_{P(x)}$ is the supremum over the set
$P(x)$ attached to the point $x$ on which the bound at $x$ reads the move fields (this $P(x)$ is
unrelated to the site labelled (P) and to the chart point $P$ below);
\item[(4)] $|A_{p_0}|\le C_e\,\bar u$ at complex entries, where $A_{p_0}$ is the field of the
entry and $C_e$ is the entry constant, the factor by which the entry field is bounded in units of
$\bar u$; the subscript $p_0$ labels the entry field and is unrelated to the exponent $p_0$ of the
degree function (Notation~\ref{0.2:not-glossary});
\item[(5)] the constants satisfy
\begin{equation}\label{9.1:eq-entry-conditions-2}
C_w\,(C_e+C_\Sigma+3)\,\bar{\bar u}\le 2 ,
\end{equation}
where $C_w$ denotes the constant of the regularity statements of this chapter,
and $C_\Sigma$ is the constant that enters beside $C_e$ in this
inequality and in the coupling-ceiling condition \eqref{9.1:eq-entry-conditions-3} below.
\end{itemize}

\emph{The single-run entry condition.} It is imposed on $\mathfrak{P}_{\mathsf{c}}$ only, and it
consists of the following two clauses.
\begin{itemize}
\item[(e)] Every entry, written in the form $(U_c,A_e)$ by its real centre and its Landau
coordinate, lies in the table of its stage or site, in the following sense. Seeds lie in their
tables by the definition of the seed classes; stage backgrounds lie in theirs by clause
(R3) of the theorem on stage backgrounds; entries at a site lie in the stage table $\mathsf{T}_N$
of Definition~\ref{9.1:def-stage-tables} (here $N$ is the stage index of that table, not the rank
of the gauge group) by clause (R4) of the same theorem and by the clause of the correlation theorem
that places its configurations in the space \cite[(1.65)]{B-CMP122b}, and this membership is
required only
on the first-class part $Z''^{(1)}_k$ of the large-field region. At the site labelled (L), one of
the sites at which entries are read, and in runs entered from (L), the
entry is the root seed, every link is an increment, and the tables are the tables
$\mathsf{T}^{\flat}$. At the site labelled (P), another such site (the label is unrelated to the
chart point $P$ below), over seeds whose entries
lie in $\mathsf{T}^{\flat}_{n_e}$, where $n_e$ is the level of the entry ($\mathsf{T}^{\flat}_{n_e}$
equals $\mathsf{T}_{n_e}$ on the part $Z^{on}$ of the large-field region, and on the region
$\Omega_k$ of Ba{\l}aban's construction at level $k$ it is
obtained from $\mathsf{T}_{n_e}$ by replacing the table coefficient $\mathsf{c}_{n_e}$ of
Definition~\ref{9.1:def-stage-tables} by $1$), reads obey condition (o1) of
\eqref{9.1:eq-constant-order-a} by the theorem that reads obey condition (o1). That reads obey (o1)
is therefore not a clause of this condition.
\item[(s)] Points lie in the stage spaces of Definition~\ref{9.1:def-stage-tables}. This is a
clause of the condition, because the theorem
that reads obey (o1) gives the free members, and the composite members of index $p\le j$ of the
slice $S(\mathbf{U})$ (here $j$ is the level of the leaf of
Notation~\ref{9.1:not-cell-point-types}, and $p$ is the index of a composite member, not a
plaquette), but gives nothing for the composite members of index $p>j$, nor for the composite
members $E_2$, $E_4$ (Definition~\ref{9.1:def-stage-tables}) of the chart point $P$ itself; these
are the members addressed by the proposition on composite members that accompanies the theorem on
stage backgrounds. The clause is used under clause (R2) of the theorem on stage backgrounds,
together with the companion corollary of that theorem and its accompanying reading. Since
the leaf's domain $\mathfrak{D}_F$ reads only the composite members of index $p\le j$, the variant
$(\mathrm{o}1)^\rho$ of condition (o1) that is read on $\mathfrak{D}_F$ is not affected.
\end{itemize}

\emph{Displayed size conditions and constants.} The size clause (3), the membership clause (e) of
the single-run entry condition, the constant $K_o$, and the condition $K_u\le L/2$ required of the
constant $K_u$ of this chapter are
\begin{equation}\label{9.1:eq-entry-conditions-a}
\begin{gathered}
K_o\sup_{P(x)}\bigl(|h^{\mathsf{m}}_\mu|+\ell_\nu|\nabla h^{\mathsf{m}}_\mu|\bigr)\le a^0_\nu(x),
\qquad \sup|\mathcal{J}_\mu|\le a^J_\nu,\\
(U_c,A_e)\ \text{lies in the table of its stage or site},\\
K_o:=2^{18}\,B_\natural\,C_{sl}^{\approx}\,\mathsf{r}_1,\qquad K_u\le L/2 .
\end{gathered}
\end{equation}

Clauses (1)--(4) of the classical entry condition hold under the coupling-ceiling condition
\begin{equation}\label{9.1:eq-entry-conditions-3}
4C\,(C_e+C_\Sigma)\,\bar{\bar u}\le\mathfrak{r}^{cl}
\end{equation}
together with the regularity statements, the local inverse-function statements, the rate
statement and the flat gradient bound of this chapter, with $c_M$ reduced as described in the
proof below; in \eqref{9.1:eq-entry-conditions-3}, $C$ denotes the constant of this condition,
which comes with the regularity and local inverse-function statements. Clause (5) is a condition
on $\bar{\bar u}$ of the same kind as \eqref{9.1:eq-entry-conditions-3}, and is required as it
stands.
\end{definition}

\begin{proof}
Clauses (1)--(4) of the classical entry condition are obtained from the coupling-ceiling condition
\eqref{9.1:eq-entry-conditions-3}, the regularity statements and the local inverse-function
statements of this chapter, together with the rate statement, which enters through the size
clause (3). The conditions \eqref{9.1:eq-entry-conditions-2} and \eqref{9.1:eq-entry-conditions-3}
are two smallness conditions on $\bar{\bar u}$; the first is clause (5) itself.

For the size clause (3) we use the bound of the rate statement for the majorant $\mathsf{H}$ of the
move fields entering that clause,
\begin{equation}\label{9.1:eq-entry-conditions-4}
\mathsf{H}\le C_\flat\,\bigl(c_M\,\mathsf{b}_j/\vartheta_P^0\bigr)\,e^{-3\delta_1 d/8},
\end{equation}
where $C_\flat$ is the constant of the rate statement, $\mathsf{b}_j$ its profile at level $j$,
$\vartheta_P^0$ its contraction rate, $\delta_1$ its decay rate, and $d$ the length in its decay
factor $e^{-3\delta_1 d/8}$.
The right-hand side of \eqref{9.1:eq-entry-conditions-4} is linear in $c_M$. A size bound that
holds with the factor $35d$ in the place of $K_o$ therefore holds with $K_o$ after one further
shrinking of $c_M$: $c_M$ is replaced by $c_M\cdot 35d/K_o$, that is, divided by $K_o/(35d)$.

It remains to treat the gradient member $\ell_\nu|\nabla h^{\mathsf{m}}_\mu|$ in the first line of
\eqref{9.1:eq-entry-conditions-a}. Let $\mathsf{q}$ denote the profile majorant of
$\ell|h^{\mathsf{m}}_\mu|$ on the contour profiles used here. Then
\begin{equation}\label{9.1:eq-entry-conditions-5}
\ell_\nu|\nabla h^{\mathsf{m}}_\mu|\le c_\nabla\,\mathsf{q},
\end{equation}
with $c_\nabla$ the constant of the flat gradient bound.
Indeed, the chart norm of the regularity statements,
\begin{equation*}
n(G)=\sup\max\{\ell|G|,\ \ell^2|\nabla G|\},
\end{equation*}
is the quantity that the regularity statements and the local Lipschitz property
$\natural\mathrm{Lip}^{\mathrm{loc}}$ of the Landau chart bound, and the bound
\eqref{9.1:eq-entry-conditions-5} is supplied by the flat gradient bound. Consequently the
two-member size clause holds wherever
\begin{equation*}
K_o\,(1+c_\nabla)\,\mathsf{q}\le a^0_\nu .
\end{equation*}
The factor $1+c_\nabla$ is absorbed where $K_o$ enters condition (o5) of the order of choice
\eqref{9.1:eq-constant-order-a}; it amounts to one more fixed shrinking of $c_M$ and introduces no
new kind of constant.
\end{proof}

\begin{theorem}[Reads land in rooms at the preliminary site]\label{9.1:thm-preliminary-reads}
Assume the following:
\begin{itemize}
\item the classical entry condition of Definition~\ref{9.1:def-entry-conditions}, at the
constant $K_o$ and with its gradient-size member;
\item clause (e) of the single-run entry condition of
Definition~\ref{9.1:def-entry-conditions};
\item slowness of the local radii across zones, Lemma~\ref{9.1:lem-radii-slowness};
\item the containment input for the stage tables on which Lemma~\ref{9.1:lem-schedule-room}
rests, at the sites of the containment of clause (v) below at which it is required;
\item the floors of the summation statement for the stage tables;
\item the floor $(F_u)$ of \eqref{9.1:eq-entry-conditions-a}.
\end{itemize}
Let $P=(\mathbf{U},\mathbf{J})$ be a point on $\mathfrak{P}_{\mathsf{c}}$ of a shape at the
preliminary site: entry at stage $n_e$, moves at stages $n_e>\nu_1>\dots>\nu_r\ge\nu$, any complex
$(b,t,\mathsf{s})$, any boxes. Let $P$ be read at stage $\nu$ by a stored term $(F,X,j)$. Write
$\mathbf{U}=e^{i\eta\mathcal{Y}}U_e$ on $X^{+}$, where $U_e$ is the last diagonal object of the
entry; links of the first kind attached after the entry are counted among the moves. Then:
\begin{enumerate}
\item[(i)] $\tilde F(X;P)=F(X;S(\mathbf{U}))$, and this is holomorphic in $\Pi$, the complex
parameters of the shape; at $r=0$ one has exactly $S(\mathbf{U})=P_0{\upharpoonright}X$, where
$P_0:=(U_e,J(U_e))$ is the configuration of the entry with its own current.
\item[(ii)] At every point of class $(=)$ --- in every zone including those of index at most $h$,
at the stages of $\mathsf{A}(x)$, and on raised histories --- in every free member
\begin{equation}\label{9.1:eq-preliminary-reads-1}
E_a(S(\mathbf{U}))\le \mathsf{T}^a_{n_e}+\tfrac18\mathsf{G}^a
=\mathsf{T}^a_\nu-\tfrac78\mathsf{G}^a ,
\end{equation}
and the point $P$ itself satisfies the same bound.
\item[(iii)] At points of class $(\mathrm{up})$, in the own units of $F$,
\begin{equation}\label{9.1:eq-preliminary-reads-2}
E_a(S(\mathbf{U}))\le\Bigl(\tfrac34+\tfrac1{200}\Bigr)\mathsf{u}_F^a .
\end{equation}
\item[(iv)] $E_5$ and the cube clause are those of the entry; the composite members are those of
part (c) of Lemma~\ref{9.1:lem-slice-displacement}.
\item[(v)] The slice point lies, in the free members and in the composite members of index
$p\le j$, in the table of the point; and $\mathrm{rd}^{\mathsf{m}}\mathbf{U}^{\mathsf{m}}\in
\mathfrak{T}^{\mathsf{m}}[0](X)$, as a class in the sense of
Definition~\ref{9.1:def-admissible-data}. This is the containment condition (o1) in the order of
constants \eqref{9.1:eq-constant-order-a}, in the form used at this site.
\end{enumerate}
The constant $K_o$ depends only on $d$, $L$, $\theta$, $M_1$, $R$, $M$ and $C_1/C_0$; nothing in
(i)--(v) depends on $g$, $k$, $K$, $N$, $N_0$, $\check{R}_k$, $h$, $r$, $n(\mathfrak{s})$, on the
ages of the large-field regions, on the depth or on the volume.
\end{theorem}

\begin{proof}
\emph{Clause (i).} The identity $\tilde F(X;P)=F(X;S(\mathbf{U}))$ and its holomorphy in $\Pi$ are
the representation \eqref{9.1:eq-slice-chart-a} of the re-read term through the slice map, in the
chart about the entry's real centre of Definition~\ref{9.1:def-slice-chart}, under the classical
entry condition. At $r=0$ there is no move, so $\mathbf{U}=U_e$, and part (a) of the slice
displacement lemma (Lemma~\ref{9.1:lem-slice-displacement}) gives
$S(U_e)=(U_e,J(U_e)){\upharpoonright}X$ as configurations, which is $P_0{\upharpoonright}X$.

For the remaining clauses we bound the displacement of each member produced by the moves. Put
\begin{equation*}
\mathsf{a}(y):=\ell_{\mathfrak{S}}(y)\sum_{\nu\le i<n_e}\ \sup_{\tilde\Delta(y)}a^0_i ,
\end{equation*}
the sum running over all stages between the read and the entry, and the supremum over the set
$\tilde\Delta(y)$ of points about $y$. The direction of the stages is that of
\cite{B-CMP122a}, where Ba{\l}aban's domain at stage $n+1$ is mapped into the one at stage $n$: a
later stage carries a smaller table, which is the ordering $n_e>\nu_1>\dots>\nu_r\ge\nu$ of the
shape. The composite move $\mathcal{Y}$
is obtained from the individual moves by the Baker--Campbell--Hausdorff formula; with the size
member of the classical entry condition (Definition~\ref{9.1:def-entry-conditions}) this gives
$|\mathcal{Y}|_y\le 2\mathsf{a}(y)/K_o$. Lemma~\ref{9.1:lem-slice-displacement} is applied, under
the classical entry condition, to $U_e$ and $\mathbf{U}=e^{i\eta\mathcal{Y}}U_e$. Combining the
bound on
$|\mathcal{Y}|_y$ with the pointwise form \eqref{9.1:eq-radii-slowness-a} of
Lemma~\ref{9.1:lem-radii-slowness}, which transports the radii from the points of
$\tilde\Delta(y)$ to $x$, and with part (b) of Lemma~\ref{9.1:lem-slice-displacement}, we obtain,
for every member $E_a$, $E_1$ included,
\begin{equation}\label{9.1:eq-preliminary-reads-a}
\ell_{\mathfrak{S}}(x)^{e_a}\cdot\bigl(\text{the move of }E_a\text{ at }x\bigr)
\le 16\,B_\natural C_{sl}^{\approx}\,\frac{\mathsf{a}(x)}{K_o},
\end{equation}
where the move of $E_a$ at $x$ is the displacement which part (b) of
Lemma~\ref{9.1:lem-slice-displacement} attaches to $E_a$, and $e_a$ is the power of the local
length carried by $E_a$.

\emph{Clause (ii).} We first show that at points of class $(=)$
\begin{equation}\label{9.1:eq-preliminary-reads-3}
\ell_\nu^{1-e_a}\sum_{\nu\le i<n_e}a^0_i\le 400\,\mathsf{r}_1\,\mathsf{G}^a ;
\end{equation}
for the member $E_1$ the left-hand side of \eqref{9.1:eq-preliminary-reads-3} is taken with the
factor $\eta^2$ that part (b) of Lemma~\ref{9.1:lem-slice-displacement} attaches to the move of
$E_1$. Let $i_1\ge\nu$ be the first stage of $\mathsf{A}(x)$ below $n_e$ met going from the read
$\nu$ towards the entry, if there is one.

For $i<i_1$ no stage of $\mathsf{A}(x)$ lies in $[\nu,i]$, so the level at $x$ is that of the
read, and
\begin{equation*}
\ell_\nu^{1-e_a}a^0_i\le\mathsf{r}_1a^a_i\le 80\,\mathsf{r}_1\,\mathrm{rm}^a_i
\end{equation*}
by Lemma~\ref{9.1:lem-stage-room} (each stage returns a room); the factor $\mathsf{r}_1=C_1/C_0$
enters for the members measured in $\alpha_0$.

The moves of the stages $i\ge i_1$ are made in coarser units than those of the read, and taken
term by term their one-step rooms in the members $E_7$ and $E_3$ are too small by the factor
$L^{(e_a-1)(\iota_i-\iota_\nu)}$, where $\iota_i$ denotes the level at which stage $i$ acts. Their
sum is nevertheless controlled: by the share-summation bound of the summation statement for the
stage tables, the sum of their shares is at
most $C_\Sigma\bar u^0_{i_1}(x)$; moreover $\ell_{i_1}(x)=\ell_\nu(x)$, and
$\bar u^a_{i_1}\le 40\,\mathrm{rm}^a_{i_1}$ by Lemma~\ref{9.1:lem-stage-room}. Hence these moves
are charged to the first raising room $\mathrm{rm}^a_{i_1}$, which is a term of $\mathsf{G}^a$:
their contribution to the left side of \eqref{9.1:eq-preliminary-reads-3} is at most
$40\,C_\Sigma\,\mathsf{r}_1\,\mathrm{rm}^a_{i_1}$; this is the check (k4) of
\eqref{9.1:eq-room-checks-a}. For the member $E_7$, a move of a coarse stage re-solved at the fine
level costs $|h^{\mathsf{m}}_\mu|/\ell_\nu$, against the room
$\mathrm{rm}^7_{i_1}\ge\alpha_1\ell_\nu^{-2}/40$ of the first raising stage.

By Lemma~\ref{9.1:lem-stage-room} the rooms $\mathrm{rm}^a_i$, $\nu\le i<n_e$, are non-negative,
and
\begin{equation*}
\mathsf{G}^a=\sum_{\nu\le i<n_e}\mathrm{rm}^a_i=\mathsf{T}^a_\nu-\mathsf{T}^a_{n_e}
\end{equation*}
is the total room between the table of the read and the table of the entry; the two
contributions charge disjoint rooms. Adding them,
\begin{equation*}
\ell_\nu^{1-e_a}\sum_{\nu\le i<n_e}a^0_i
\le \max\{80,\,40C_\Sigma\}\,\mathsf{r}_1\sum_{\nu\le i\le i_1}\mathrm{rm}^a_i
\le 368\,\mathsf{r}_1\,\mathsf{G}^a\le 400\,\mathsf{r}_1\,\mathsf{G}^a ,
\end{equation*}
since $80\vee 40C_\Sigma=368$. If no stage $i_1$ exists, only the first contribution is present and
the bound holds with $80$. This proves \eqref{9.1:eq-preliminary-reads-3}.

At a point of class $(=)$ one has $\ell_{\mathfrak{S}}=\ell_\nu$, so
$\ell_{\mathfrak{S}}^{-e_a}\mathsf{a}(x)$ is the left side of
\eqref{9.1:eq-preliminary-reads-3} read with the shares $a^0_i$ taken as their suprema over
$\tilde\Delta(x)$ (and, for $E_1$, with the factor $\eta^2$), so that, by
\eqref{9.1:eq-preliminary-reads-a} and \eqref{9.1:eq-preliminary-reads-3}, the move of $E_a$ at
$x$ satisfies
\begin{equation*}
\bigl(\text{the move of }E_a\text{ at }x\bigr)
\le \frac{6400\,B_\natural C_{sl}^{\approx}\,\mathsf{r}_1\,\mathsf{G}^a}{K_o}
<\tfrac18\,\mathsf{G}^a ,
\end{equation*}
where $16\cdot 400=6400$, and the last inequality uses the floor on $K_o$ in the numerical checks
\eqref{9.1:eq-room-checks-a}, by which $2^{18}B_\natural C_{sl}^{\approx}\mathsf{r}_1\le K_o$,
together with $6400<2^{18}/8$. The entry itself satisfies
$E_a(S(U_e))\le\mathsf{T}^a_{n_e}$ by clause (e) of the single-run entry condition. Adding the move
gives the inequality in \eqref{9.1:eq-preliminary-reads-1}; the equality there is the identity
$\mathsf{G}^a=\mathsf{T}^a_\nu-\mathsf{T}^a_{n_e}$ displayed above. For the
point $P$ itself, the members are bounded term by term by the sizes of the classical entry
condition, the gradient $\nabla h^{\mathsf{m}}_\mu$ of the move of stage $i$ being measured at the
scale $\ell_i$ of its own stage; the member $E_3$ is bounded by the current bound of the small-field
machine together with clause (e).

\emph{Clause (iii).} Since the tables are non-increasing in the stage, $\mathsf{T}_{n_e}\le\bar
u_\nu$. Next, $\ell_{n_F}^{e_a}\bar u^a_\nu\le\frac34\alpha_{a,n_F}$, where $n_F$ is the level of
the stored term $F$, $n_0\ge n_F$ is the level to which the point was raised, and $\iota\ge n_0$ is
the level of the unit being compared, as in the comparison of units across levels. If $n_0=n_F$
(the point was raised in the chain), this follows from the part of the summation statement for
the stage tables which fixes how the units and the printed coefficients change from stage to
stage, the power of $\ell_\circ$ being at most $\ell_\circ^{2(\iota-n_0)}$ as in
the proof of the comparison of units across levels. If $n_0>n_F$, put $u:=n_0-n_F\ge1$ and
$v:=\iota-n_0$; the comparison of units across levels gives
\begin{equation*}
3L^{-(\iota-n_F)}\ell_\circ^{2(\iota-n_0)}\alpha_{a,\iota}
\le K_uL^{-u}3^{-v}(u+v)^{1/2}\alpha_{a,n_F}
\le\frac{K_u}{L}\,\alpha_{a,n_F},
\end{equation*}
and by the floor $(F_u)$ of \eqref{9.1:eq-entry-conditions-a} this is at most
$\frac34\alpha_{a,n_F}$. Hence, by clause (e) of the single-run entry condition, and since in the
own units of $F$ the bound $\ell_{n_F}^{e_a}\bar u^a_\nu\le\frac34\alpha_{a,n_F}$ reads
$\bar u^a_\nu\le\frac34\mathsf{u}^a_F$,
\begin{equation*}
E_a(S(U_e))\le\mathsf{T}^a_{n_e}\le\bar u^a_\nu\le\tfrac34\,\mathsf{u}^a_F .
\end{equation*}

A smaller unit lowers a member, by the monotonicity of the table members in the unit. Hence the
move, measured in the units of level $n_F$, is at most $\ell_{n_F}/\ell_{\mathfrak{S}}$ times the
right side of \eqref{9.1:eq-preliminary-reads-a}. Using the share-summation bound of the summation
statement, and then the comparison of units across levels at $n_F\le n_0$ together with the check
(k5) of \eqref{9.1:eq-room-checks-a}, we obtain
\begin{equation*}
\frac{\ell_{n_F}}{\ell_{\mathfrak{S}}}\cdot
\frac{16B_\natural C_{sl}^{\approx}\mathsf{a}(x)}{K_o}
\le\frac{16B_\natural C_{sl}^{\approx}}{K_o}\,C_\Sigma\,\ell_{n_F}\bar u^0_\nu
\le\frac{16\,C_\Sigma K_u}{2^{18}}\,\alpha_{0,n_F}.
\end{equation*}
Since $16C_\Sigma K_u=993.6<2^{18}/200$, the move is at most $\frac1{200}$ in the own units of
$F$. Adding it to the bound $\frac34$ of the table gives \eqref{9.1:eq-preliminary-reads-2}.

\emph{Clause (iv).} The composite members of $S(\mathbf{U})$ are those described by part (c) of
Lemma~\ref{9.1:lem-slice-displacement}, applied to $U_e$ and $\mathbf{U}$ as in the derivation of
\eqref{9.1:eq-preliminary-reads-a}; $E_5$ and the cube clause are read at the entry.

\emph{Clause (v).} At points of class $(=)$, Lemma~\ref{9.1:lem-schedule-room} (the schedule of
tables returns a room), with the containment input for the stage tables where it is required,
gives $\mathsf{T}_\nu\le(\mathsf{c}_F^{+}-\varrho_F)\mathsf{u}_F$, so that
\eqref{9.1:eq-preliminary-reads-1} places the slice point in the table of the point, with the room
$\varrho_F$ of the own tube \eqref{9.1:eq-own-tube-a} to spare. At points of class
$(\mathrm{up})$,
\begin{equation*}
\tfrac34+\tfrac1{200}=0.755\le 0.775=\tfrac45-\tfrac1{40}\le\mathsf{c}_F^{+}-\varrho_F ,
\end{equation*}
the last inequality being the check (k6) of \eqref{9.1:eq-room-checks-a}; so
\eqref{9.1:eq-preliminary-reads-2} places the slice point in the table as well. In the chart of
Definition~\ref{9.1:def-slice-chart} this is the statement
$\mathrm{rd}^{\mathsf{m}}\mathbf{U}^{\mathsf{m}}\in\mathfrak{T}^{\mathsf{m}}[0](X)$ of
\eqref{9.1:eq-admissible-data-a}, as a class modulo block-centre gauges.
\end{proof}

\begin{proposition}[Reads at the small-field, large-field and box sites]\label{9.1:prop-site-reads}
Entries, increments, rooms and reads are understood as in
Theorem~\ref{9.1:thm-preliminary-reads}, and $X^{+}$ is the hull of the region read, as in
Lemma~\ref{9.1:lem-slice-displacement}. Items~(a)--(c) concern the site of the $\mathbf{T}$-step
and the link site, item~(d) the box site and item~(e) the single-scale site.
At every site below, the \emph{radius} of a read is the
Cauchy radius of the site that remains after the collar. Let $K_o$ be the circle constant of
the reads, let $k$ be a level, let $\mathsf{t}$ be the bracket parameter and $\mathsf{s}$ the
dilation, and let $B'$ be the variable of the $\mathbf{T}$-step.
\begin{enumerate}
\item[(a)] \emph{The site of the $\mathbf{T}$-step from level $k$ to level $k+1$.}
The entry is
the seed of the $\mathbf{R}$-born output $\Phi$, with its slice in the domain
$\mathfrak{D}_\Phi$. There is one increment,
\begin{equation*}
(\mathsf{t}\hat{\chi}_\square+\cdots)\,\mathbb{H}_k(\mathsf{s},B'),
\end{equation*}
where $\mathbb{H}_k(\mathsf{s},B')$ is the $\mathbf{T}$-step term of level $k$ read at the
dilation $\mathsf{s}$ and the variable $B'$, and $\hat{\chi}_\square$ the localising function of
the cube $\square$ that multiplies the bracket parameter $\mathsf{t}$.
The room at the points of class $(=)$ on the layers $n\le k$ is
\begin{equation}\label{9.1:eq-site-reads-1}
\tfrac12\,(1-\theta_S)\,\beta\,2^{-(k+1-n)}\,\mathsf{u},
\end{equation}
where $\beta$, $\theta_S$ and the unit $\mathsf{u}$ of the term read are those of
Theorem~\ref{9.1:thm-preliminary-reads},
and the room at the points of class $(\mathrm{up})$ on the small-field region $\Omega_{k+1}$ of
the step is given by scaling, as in
clause~(iii) of Theorem~\ref{9.1:thm-preliminary-reads}. The radius, that of the circle in the
bracket parameter $\mathsf{t}$ on which the read is taken, is $\mathsf{w}_T/K_o$. The
read at this site lands in its room undilated, and on $\Omega_{k+1}$ also on the
$\mathsf{t}$-circle of that radius.
\item[(b)] Assume the containment hypothesis at the $\mathbf{T}$-site: the containment of the
$\mathbf{T}$-row of the correlation theorem holds in the form ``seed plus shift does not exceed
the room''; this hypothesis is established with the continuation of these reads later in this
chapter. Then the read of~(a) lands in its room on the $\mathsf{t}$-circle on the layers $n\le k$.
\item[(c)] \emph{The link site.} The entry is the last object of the chain to whose
objects Lemma~\ref{9.1:lem-slice-displacement} applies; the increments are the link terms of
that chain and the shift in $\mathsf{t}$; the rooms are those of the chain-term row of the
correlation theorem; the radius is $\mathsf{w}_L/K_o$, that is, the radius written
$c_J\,\theta'^{-1}$, in that statement's own letters, in the proposition attached to the
correlation theorem at this site. Assume the containment hypothesis at the
link site: the containment of the chain-term row of the correlation theorem holds in the form
``entry plus increments does not exceed these rooms''; this hypothesis is established with the
continuation of these reads later in this chapter. Then the read at the link site lands in its
rooms.
\item[(d)] \emph{The box site.} For every box $Z\supseteq X^{+}$ the entry, the increments and
the room are those of the site at which the box is read, and no radius enters. The read is the
one given by the slice displacement lemma, Lemma~\ref{9.1:lem-slice-displacement}, with
$\mathfrak{S}:=\boldsymbol{\Omega}(Z)$.
\item[(e)] \emph{The single-scale site.} The entry is a configuration $\mathbf{U}$ of the
single-scale configuration class at level $j$ over the region $Y$ that is read at this
site; there are no increments; the room is the
ratio of item~(c2) of the correlation theorem; the radius is the Cauchy radius attached to that
class. The read at this site is taken up later in this chapter.
\end{enumerate}
\end{proposition}

\begin{proof}
The slice displacement lemma, Lemma~\ref{9.1:lem-slice-displacement}, does not depend on the site
at which it is applied. We apply it at the $\mathbf{T}$-site with $r=1$.

The majorant of the response is exponential in the distance
$d(\,\cdot\,,\operatorname{supp}B')$ to the support of the $\mathbf{T}$-step variable $B'$; this
is the exponential majorant bound for the response, together with \cite[(1.11)]{B-CMP116} and
item~(8) of the operator theorem at free current. Hence the majorant varies slowly.

\emph{The undilated read.} The requirement that $K_o$ times the response be at most one half of
the room allotted to it is a $\gamma$-ceiling condition, that is, a condition which holds once the
coupling lies below a ceiling fixed through $\gamma$, and its degree is $p_0-q<0$. This is the
line of item~(ii) of the size statement for the response that sets this allotment; it is the
undilated assertion of~(a).

\emph{The $\mathsf{t}$-circle at a single scale.} The number $\mathsf{w}_T$ is defined by the
room. The definition goes through the normalisation of \cite[(3.15)]{B-CMP109}, in which
Ba{\l}aban's parameter $r$ times a maximum is set equal to
$\tfrac14\alpha_2$, through item~(5) of the operator theorem at free current, and through
\cite[(1.22)]{B-CMP116}. Dividing $\mathsf{w}_T$ by $K_o$ gives the radius of~(a). At that radius
the shift together with the slice move fits into the room that the shift alone filled at
$\mathsf{w}_T$. The rows of the correlation theorem then carry the factor
$K_o\mathsf{w}_T^{-1}$, and
\begin{equation}\label{9.1:eq-site-reads-2}
K_o\,\mathsf{w}_T^{-1}\le K_o\,\theta_k ,
\end{equation}
with $\theta_k$ the inverse coupling at level $k$. Thus the rows carry a factor of at most
$K_o\theta_k$, and that the rows absorb this factor is again a $\gamma$-ceiling
condition. On the small-field region $\Omega_{k+1}$, where the room is the scaled room of
clause~(iii) of Theorem~\ref{9.1:thm-preliminary-reads}, this is the assertion of~(a) on the
$\mathsf{t}$-circle at the single scale. This proves~(a).

For~(b) and~(c) the conclusion is, in the form stated, the assumed containment: the
containment hypothesis at the $\mathbf{T}$-site reads, on the layers $n\le k$, ``seed plus shift
does not exceed the room'', and the containment hypothesis at the link site reads ``entry plus
increments does not exceed these rooms''. Part~(d) is
Lemma~\ref{9.1:lem-slice-displacement} with $\mathfrak{S}:=\boldsymbol{\Omega}(Z)$, which
applies because that lemma does not depend on the site. Part~(e) describes the read at the
single-scale site, which is taken up later in this chapter.
\end{proof}

\begin{definition}[Net mismatch of the two runs' block data in a real frame]\label{9.1:def-net-mismatch}
Let $\mathsf{A}$ and $\mathsf{B}$ be the two runs compared in this chapter, with complex configurations
$\mathbf{U}$ of run $\mathsf{A}$ and $\mathbf{U}'$ of run $\mathsf{B}$. Let $\mathsf{h}$ be the real
frame defined as follows. At the centre of each labelled cell, $\mathsf{h}$ is the real gauge frame
that the real-frame construction of this chapter fixes at that centre. At every other point,
$\mathsf{h}=1$. For a vertex pair $p\in\mathfrak{F}^V$ read by a term $F$, write $\mathsf{h}_-$ and
$\mathsf{h}_+$ for the frame factors at the two vertices of the pair, in the order in which they
enter \eqref{9.1:eq-net-mismatch-1}, and put
\begin{equation*}
\mathsf{v}:=\mathrm{rd}^{\mathsf{A}}\,\mathbf{U}.
\end{equation*}
On the paired strata, the \emph{net mismatch} $\mathfrak{d}^{net}$ at $p$ is defined by
\begin{equation}\label{9.1:eq-net-mismatch-1}
e^{i\mathfrak{d}^{net}}:=\mathsf{h}_-^{-1}\,\bigl(\mathrm{rd}^{\mathsf{B}}\,\mathbf{U}'\bigr)\,
\mathsf{h}_+\,\mathsf{v}^{-1}.
\end{equation}
The letter $\mathsf{a}$ has the meaning it has in Theorem~\ref{9.1:thm-preliminary-reads}; that it is
common to the two runs is the content of the common-parameter statement for the two runs. For the index
$j$ and the region $X$ of the two-cutoff seminorm $\mathsf{S}_j^{+}$ of
Definition~\ref{9.1:def-admissible-data}, with $n_X$ as in that definition, we put
\begin{equation}\label{9.1:eq-net-mismatch-2}
\vartheta^{\flat}:=\vartheta^{(j\wedge n_X)/2}.
\end{equation}
By the pair of bounds of the comparison statement for the two runs' readings, applied at the vertex
pair $p$, the net mismatch obeys
\begin{equation}\label{9.1:eq-net-mismatch-3}
\bigl|\mathfrak{d}^{net}\bigr|\le\bigl(|\mathfrak{d}|+2\,|\log\mathsf{w}|\bigr)\,(1+C c_4),
\end{equation}
where $\mathfrak{d}$ and $\mathsf{w}$ are the two quantities bounded by that pair of bounds (here
$\mathsf{w}$ denotes that quantity, not the stop depth $\mathsf{w}(g)$), and $C$ and $c_4$ are the
constants of that pair of bounds.
\end{definition}

\begin{proof}
Inequality \eqref{9.1:eq-net-mismatch-3} is the pair of bounds of the comparison statement for the two
runs' readings, applied at $p$ to the readings $\mathsf{v}=\mathrm{rd}^{\mathsf{A}}\,\mathbf{U}$ and
$\mathrm{rd}^{\mathsf{B}}\,\mathbf{U}'$ of \eqref{9.1:eq-net-mismatch-1}.

At the cells where the two runs are handed the same datum, $\mathfrak{d}$ and $\mathsf{w}$ are not
separately small, so the right side of \eqref{9.1:eq-net-mismatch-3} is not small there; at those
cells the two factors compensate in the product \eqref{9.1:eq-net-mismatch-1}.
\end{proof}

\begin{lemma}[The share law at common cells]\label{9.1:lem-share-law}
Assume the following:
\begin{itemize}
\item the classical entry condition \eqref{9.1:eq-entry-conditions-a} of
Definition~\ref{9.1:def-entry-conditions};
\item the per-step increment bound for the moves, and the growth bound relative to the stage
table $\mathsf{T}$;
\item the composition identity for composite moves;
\item the estimates on the averaging map at complex configurations.
\end{itemize}
Let $p=p_r\in\mathfrak{F}^V$ be a run over a seed $p_0$ of stage $n_e$. Let $c$ be a labelled cell of
level $l$. Then at every point of $\mathfrak{P}_{\mathsf{c}}$,
\begin{equation}\label{9.1:eq-share-law-1}
|\mathfrak{d}^{net}(c)|\le C_\rho^{net}\,\vartheta^{l/2}\,\mathsf{a}(c).
\end{equation}
Here $\mathfrak{d}^{net}$ is the net mismatch of Definition~\ref{9.1:def-net-mismatch}, and
$\mathsf{a}(c)$ is the share majorant of the moves at $c$. The constant $C_\rho^{net}$ is independent
of $r$, of the indices $\mathsf{N}$ and $\mathsf{N}_0$ (letters distinct from the rank $N$ of the
gauge group), of $\check{R}_k$, $k$, $g$ and of the depth.
\end{lemma}

\begin{proof}
We write $\mathcal{M}$ for one step of the block-averaging map of
Definition~\ref{1.5:def-block-average}, and $\mathcal{M}^{l}$ for its $l$-fold iterate. For a
field $A$ we write $|A|_c$ for its norm on the cell $c$, in the units of the entry condition
\eqref{9.1:eq-entry-conditions-a}. Two cases
occur, according to the level index $\iota_{n_e}(c)$ of $c$ at the entry stage $n_e$.

\emph{Case $\iota_{n_e}(c)>l$.} In this case $c$ was raised at a stage before the entry stage, so
$c$ is not a cell of the seed's set. Clause (d) of the existence statement for the moves, in the
form it takes at the corollary on vertex pairs, then applies without change and gives
\eqref{9.1:eq-share-law-1}, since the share $\mathfrak{s}_u(c)$ entering it satisfies
\begin{equation*}
\mathfrak{s}_u(c)=2\ell_\nu\,\bar u^0_\nu=\ell_\nu\,a^0_{i_1}\le\mathsf{a}(c),
\end{equation*}
where $\bar u^0_\nu$ is the share and $a^0_{i_1}$ the size bound of the entry condition
\eqref{9.1:eq-entry-conditions-a}, read at the stages $\nu$ and $i_1$ concerned, and $\ell_\nu$ is
the length unit at the stage $\nu$.
This holds because at such a cell the room is a fixed fraction of the radius, by
Lemma~\ref{9.1:lem-stage-room}.

\emph{Case $\iota_{n_e}(c)=l$.} We introduce the following notation.
\begin{itemize}
\item For a field $A$ and a configuration $U$, put
$m_c(A;U):=\mathcal{M}^{l}\bigl(e^{i\eta A}U\bigr)(c)$.
\item Let $\mathbf{U}_s$ and $\mathbf{U}'_s$ be the configurations of the runs $\mathsf{A}$ and
$\mathsf{B}$ at the member $p_s$ of the run, and put $W_s:=\mathcal{M}(\mathbf{U}'_s)$.
\item Put $\mathfrak{a}:=A_{p_0}$; in the frame fixed below, $W_0=e^{i\eta\mathfrak{a}}\mathbf{U}_0$.
\item Write $\mathbf{U}_r=e^{i\eta\mathcal{Y}}\mathbf{U}_0$, with $\mathcal{Y}$ the composite move.
\item Write $W_r=e^{i\eta Y}W_0$, with $Y:=Y[\mathcal{Y}';\mathbf{U}'_0]$.
\end{itemize}
The shells are monotone, so $c$ is a cell of the seed's own set $\mathfrak{B}(W)$, the set of cells
of the seed. Hence the two
runs are handed the same datum $\mathsf{v}_0$ at $c$, complex data included.

\emph{The frame.} By the representation theorem in the gauge $Q_j=0$, where $Q_j$ is the gauge
condition of that theorem, a gauge transform
$W^{\mathsf{u}}$ of $W$ is written $W^{\mathsf{u}}=e^{i\eta A}U$ with $Q_j(U,\eta A)=0$. By the
estimates on the averaging map,
\begin{equation*}
\mathcal{M}(W^{\mathsf{u}})=\mathsf{q}\,V\,\mathsf{q}^{-1},\qquad
\mathsf{q}=\mathsf{q}[A;U]\neq 1,
\end{equation*}
while $\mathcal{M}(W)=V$, the block field of $W$; here $\mathsf{q}$ is a group element determined
by $A$ and $U$. Applied to $W_0$, represented in the frame as $e^{i\eta\mathfrak{a}}\mathbf{U}_0$,
this gives, at block centres and modulo stabilisers, that the real frame
$\mathsf{h}$ of Definition~\ref{9.1:def-net-mismatch} is $\mathsf{q}[\mathfrak{a};\mathbf{U}_0]$,
and $|\mathsf{h}-1|\le 16d\,|\mathfrak{a}|_c$.

Two identities must therefore be read in different frames. The identity
$W_0=e^{i\eta\mathfrak{a}}\mathbf{U}_0$ holds in the frame. The identity
$m_c(0;W_0)=m_c(0;\mathbf{U}_0)$ holds outside it; read in the frame it becomes
$m_c(0;W_0)=\mathsf{h}_-\mathsf{v}_0\mathsf{h}_+^{-1}$.

Without the frame in Definition~\ref{9.1:def-net-mismatch} the seed alone would contribute
$\log[\mathsf{h}_-\mathsf{v}_0\mathsf{h}_+^{-1}\mathsf{v}_0^{-1}]$, of order
$32d\,|\mathfrak{a}|_c$ in units of the radius and independent of the moves, which only a cap
(Definition~\ref{2.3:def-cap}), bounding $\vartheta^{l}$ times $2$ to the power of the depth by a
small constant, could absorb.

\emph{The net mismatch at $c$.} Write $m_c(\mathcal{Y};\mathbf{U}_0)=e^{iy_0}\mathsf{v}_0$, and put
\begin{equation*}
\gamma_0:=\mathsf{h}_-\mathsf{v}_0\mathsf{h}_+^{-1}\mathsf{v}_0^{-1},\qquad
\varpi:=\mathsf{h}_-^{-1}\gamma_0 .
\end{equation*}
With the frame $\mathsf{h}$ in the definition, the three estimates $(\alpha)$--$(\gamma)$ below give
\begin{equation}\label{9.1:eq-share-law-2}
e^{i\mathfrak{d}^{net}(c)}
=e^{i\operatorname{Ad}_{\varpi}y_0}\,e^{-iy_0}
\bigl(1+O(\Xi)+O(|Y-\mathcal{Y}|_c)\bigr).
\end{equation}
In particular $\mathfrak{d}^{net}(c)$ vanishes at $r=0$. The factor
$e^{i\operatorname{Ad}_{\varpi}y_0}e^{-iy_0}$ differs from $1$ by at most
\begin{equation}\label{9.1:eq-share-law-3}
\bigl(|\mathsf{h}-1|+|\gamma_0-1|\bigr)\cdot 2|\mathcal{Y}|_c
\le 100d\,|\mathfrak{a}|_c\,|\mathcal{Y}|_c .
\end{equation}
This term is bilinear in $\mathfrak{a}$ and $\mathcal{Y}$, and so is measured in share units.

$(\alpha)$ By the estimates on the averaging map and Cauchy's inequality, the average
$m_c(Y;W_0)$ differs from $m_c(\mathcal{Y};W_0)$ by at most $68d\,|Y-\mathcal{Y}|_c$.

$(\beta)$ Put
\begin{equation*}
\begin{split}
\Xi(\mathcal{Y},\mathfrak{a})
:={}&\log\bigl[m_c(\mathcal{Y};e^{i\eta\mathfrak{a}}\mathbf{U}_0)\,
m_c(\mathcal{Y};\mathbf{U}_0)^{-1}\bigr]\\
&-\log\bigl[m_c(0;e^{i\eta\mathfrak{a}}\mathbf{U}_0)\,m_c(0;\mathbf{U}_0)^{-1}\bigr].
\end{split}
\end{equation*}
This function is analytic on the bidisc in the variables $(\mathcal{Y},\mathfrak{a})$. It vanishes
at $\mathcal{Y}=0$ and at $\mathfrak{a}=0$. The Schwarz lemma, applied once in each variable as in
the proof of the second-order estimate on the averaging map, gives
\begin{equation*}
|\Xi|\le C_W\,|\mathcal{Y}|_c\,|\mathfrak{a}|_c ,
\end{equation*}
with $C_W$ the constant of the estimates on the averaging map
(Definition~\ref{9.1:def-entry-conditions}).
The subtracted term
$\log[m_c(0;e^{i\eta\mathfrak{a}}\mathbf{U}_0)\,m_c(0;\mathbf{U}_0)^{-1}]$ is $\log\gamma_0$,
which the frame $\mathsf{h}$ has removed.

$(\gamma)$ The composition identity for the composite move reads
\begin{equation*}
e^{i\eta A_{p_r}}=e^{i\eta Y}\,e^{i\eta\mathfrak{a}}\,e^{-i\eta\mathcal{Y}} .
\end{equation*}
It gives, with $C_B$ the constant of the composition identity,
\begin{equation}\label{9.1:eq-share-law-4}
|Y-\mathcal{Y}|\le 2|A_{p_r}-\mathfrak{a}|
+2C_B\,\eta\,\bigl(|Y|+|\mathcal{Y}|\bigr)|\mathfrak{a}| .
\end{equation}

On all of $4\mathfrak{P}_{\mathsf{c}}$ we have
\begin{equation*}
|Y|_c+|\mathcal{Y}|_c\le 3\mathsf{a}(c).
\end{equation*}

\emph{Real parameters.} Let $\Pi$ be a real point of $4\mathfrak{P}_{\mathsf{c}}$. By the per-step
increment bound for the moves, with its constant $c_b$, and the growth bound relative to
$\mathsf{T}$, with its constant $c_a^{GR}$ and its factor $C^P$, both bounds being stated in a
unit smaller than $\ell_\nu a^0_{\nu_s}$, which accounts for the factor $\tfrac12$,
\begin{equation}\label{9.1:eq-share-law-5}
\ell_\nu|A_{p_r}-\mathfrak{a}|
\le\sum_s\ell_\nu|A_{p_s}-A_{p_{s-1}}|
\le\tfrac12 c_*\sum_s\ell_\nu\,a^0_{\nu_s}\,\vartheta^{\iota_{\nu_s}} .
\end{equation}
Here $c_*:=c_b+c_a^{GR}C^P$, $\nu_s$ is the stage of the member $p_s$, and $\iota_{\nu_s}$ is the
level index at that stage. Moreover, the following two bounds hold, where $P(c)$ is the set attached
to $c$ in the entry condition \eqref{9.1:eq-entry-conditions-a} and $y^{\mathsf{T}}_0$ is
the constant of the stage table $\mathsf{T}$:
\begin{equation*}
\iota_{\nu_s}\ge l\ \text{on}\ P(c),\qquad
|\mathfrak{a}|_c\le y^{\mathsf{T}}_0\,\bar{\bar u}\,\vartheta^{l}.
\end{equation*}
Since $\iota_{\nu_s}\ge l$ on $P(c)$ and $\vartheta<1$, and since the share majorant
$\mathsf{a}(c)$ majorises $\sum_s\ell_\nu a^0_{\nu_s}$, the bound \eqref{9.1:eq-share-law-5} gives
\begin{equation*}
2|A_{p_r}-\mathfrak{a}|_c\le c_*\,\vartheta^{l}\mathsf{a}(c).
\end{equation*}
With $\eta\le1$ and the two bounds $|\mathfrak{a}|_c\le y^{\mathsf{T}}_0\bar{\bar u}\vartheta^{l}$
and $|Y|_c+|\mathcal{Y}|_c\le3\mathsf{a}(c)$, the bound \eqref{9.1:eq-share-law-4} then gives
\begin{equation*}
|Y-\mathcal{Y}|_c\le\bigl(c_*+6C_B\,y^{\mathsf{T}}_0\bar{\bar u}\bigr)\vartheta^{l}\mathsf{a}(c).
\end{equation*}
We insert these bounds into the three contributions to
\eqref{9.1:eq-share-law-2}.
\begin{itemize}
\item The bilinear term \eqref{9.1:eq-share-law-3} is at most
\begin{equation*}
100d\cdot y^{\mathsf{T}}_0\bar{\bar u}\vartheta^{l}\cdot 3\mathsf{a}(c)
=300d\,y^{\mathsf{T}}_0\bar{\bar u}\,\vartheta^{l}\mathsf{a}(c).
\end{equation*}
\item The term $\Xi$ of $(\beta)$ is at most $3C_W\,y^{\mathsf{T}}_0\bar{\bar u}\,\vartheta^{l}\mathsf{a}(c)$.
\item The term $(\alpha)$, which is $68d$ times the last bound for $|Y-\mathcal{Y}|_c$, is at most
\begin{equation*}
68d\,(c_*+6C_B\,y^{\mathsf{T}}_0\bar{\bar u})\,\vartheta^{l}\mathsf{a}(c).
\end{equation*}
\end{itemize}
Hence
\begin{equation}\label{9.1:eq-share-law-6}
|\mathfrak{d}^{net}(c)|\le C_\rho^{net,\mathbb{R}}\,\vartheta^{l}\,\mathsf{a}(c),
\end{equation}
where
\begin{equation*}
C_\rho^{net,\mathbb{R}}:=68d\bigl(c_*+6C_B\,y^{\mathsf{T}}_0\bar{\bar u}\bigr)
+3C_W\,y^{\mathsf{T}}_0\bar{\bar u}+300d\,y^{\mathsf{T}}_0\bar{\bar u}.
\end{equation*}

\emph{A priori bound on $4\mathfrak{P}_{\mathsf{c}}$.} Apply the estimates on the averaging map in
each of the two runs, starting from the common datum $\mathsf{v}_0$. This gives, on all of
$4\mathfrak{P}_{\mathsf{c}}$,
\begin{equation}\label{9.1:eq-share-law-7}
|\mathfrak{d}^{net}(c)|\le 35d\bigl(|Y|_c+|\mathcal{Y}|_c\bigr)\le 105d\,\mathsf{a}(c).
\end{equation}
This bound is already in share units.

\emph{Complex parameters.} The map $\Pi\mapsto\mathfrak{d}^{net}(c)$ is holomorphic on
$4\mathfrak{P}_{\mathsf{c}}$ by the classical entry condition
\eqref{9.1:eq-entry-conditions-a} and by the analyticity of the averaging and large-field
operations at complex configurations. We then proceed by
the complexification corollary for the vertex-pair family: we apply the two-constant estimate, with
its constant $C_{\mathbb{C}}$, to each component of $\mathfrak{d}^{net}(c)$,
taken in a basis of the Lie algebra $\mathfrak{g}$ of $\mathrm{SU}(N)$. Its two constants are the
real bound \eqref{9.1:eq-share-law-6} and the a priori bound \eqref{9.1:eq-share-law-7}. This gives
\eqref{9.1:eq-share-law-1} on $\mathfrak{P}_{\mathsf{c}}$, with
\begin{equation*}
C_\rho^{net}:=\dim\mathfrak{g}\cdot\max\Bigl\{C_\rho^{net,\mathbb{R}},\,
\bigl(105d\,C_\rho^{net,\mathbb{R}}\bigr)^{1/2},\,40d\,C_{\mathbb{C}}\Bigr\}.
\end{equation*}
\end{proof}

\begin{remark}[Chains of links]
Lemma~\ref{9.1:lem-share-law} concerns a single run over a seed. In a chain of links there are two
kinds of link.
\begin{itemize}
\item A link of the second kind keeps its parent's data at labelled cells, by the corresponding
clause of the statement that contains the per-step increment bound for the moves.
\item A link of the first kind is a move.
\end{itemize}
Over an entry of the second type, the argument of the proof above is the one that applies. There
the links of the first kind since the root are counted among the moves, and their shares are given
by the
accumulation bound for links. This section does not carry out that argument link by link, and no
statement about chains is made here; the chain case is taken up in the section of this chapter
that contains the accumulation bound for links.
\end{remark}

\begin{lemma}[Damping beats depth]\label{9.1:lem-damping-depth}
Let $(F,X,j)$ be a leaf in the sense of Notation~\ref{9.1:not-cell-point-types}, and let
$\mathfrak{S}$, $\lambda_{\mathfrak{S}}$, $d_{\mathfrak{S}}$, $\eta$ and $n_F$ be as there. Let
$\delta_1>0$ be the decay rate in $d_{\mathfrak{S}}$ with which the profiles of this chapter are
smeared. Let $x$ be a point of class \textup{(=)} and let $y$ be a cell of type \textup{(top)}
or \textup{(dist)}, both in the sense of Notation~\ref{9.1:not-cell-point-types}, and put
$s:=j-n_F(x)$. Then
\begin{equation}\label{9.1:eq-damping-depth-1}
  e^{-\delta_1 d_{\mathfrak{S}}(x,y)/2}\le C_{dp}\,4^{-s},
\end{equation}
where the constant $C_{dp}=C_{dp}(L,\delta_1,M_1)$ depends only on $L$, $\delta_1$ and $M_1$.
\end{lemma}

\begin{proof}
By Notation~\ref{9.1:not-cell-point-types}, $n_F(x)\le j$, so $s\ge 0$. We write
$(u)_+:=\max\{u,0\}$. We treat the two types of $y$ separately.

\emph{Cells of type \textup{(top)}.} For such a cell $\lambda_{\mathfrak{S}}(y)=j$, and for
the point $x$ of class \textup{(=)} we have $\lambda_{\mathfrak{S}}(x)=n_F(x)=j-s$. The two
levels thus differ by $s$, so every path from $x$ to $y$ meets level
$j=\lambda_{\mathfrak{S}}(x)+s$. The level-crossing lower bound for the scaled distance gives
$d_{\mathfrak{S}}(x,y)\ge\mathsf{g}\,(s-1)_+$, where $\mathsf{g}>0$ is the gain in scaled
distance per level crossed in that bound. The floor condition on $\delta_1\mathsf{g}$ gives
$e^{-\delta_1\mathsf{g}/2}\le L^{-8}$. Combining the two,
\begin{equation*}
  e^{-\delta_1 d_{\mathfrak{S}}(x,y)/2}\le L^{-8(s-1)_+}.
\end{equation*}
For $s=0$ the right side is $1\le 4$. For $s\ge1$, since $L^{8}\ge 4$, it is at most
$4^{-(s-1)}=4\cdot 4^{-s}$. Hence the left side of \eqref{9.1:eq-damping-depth-1} is at most
$4\cdot4^{-s}$.

\emph{Cells of type \textup{(dist)}.} By the separation property of cells of type
\textup{(dist)},
\begin{equation*}
  |x-y|\ge 27LM_1L^{j}\eta .
\end{equation*}
Let $\gamma$ be any path from $x$ to $y$, write $d_{\mathfrak{S}}(\gamma)$ for its scaled
length, and let $j-s+t$, $t\ge0$, be the highest level it meets; in particular it reaches at
most level $j-s+t$. We bound $d_{\mathfrak{S}}(\gamma)$ from below in two ways and then pass to
the infimum over $\gamma$, which is $d_{\mathfrak{S}}(x,y)$.

First, the level-crossing lower bound for the scaled distance gives
$d_{\mathfrak{S}}(\gamma)\ge \mathsf{g}\,(t-1)_+$.

Second, the scaled length is at least the Euclidean length measured in the largest unit met
along the path \cite[(2.46)]{B-CMP96}. By Notation~\ref{9.1:not-cell-point-types} that unit is
$\ell_{\mathfrak{S}}=L^{\lambda_{\mathfrak{S}}}\eta\le L^{j-s+t}\eta$. Therefore
\begin{equation*}
  d_{\mathfrak{S}}(\gamma)\ge \frac{27LM_1L^{j}\eta}{L^{j-s+t}\eta}=27LM_1L^{s-t}.
\end{equation*}

Either $t\ge s/2$ or $t<s/2$.

If $t\ge s/2$, then $(t-1)_+\ge s/2-1$. The first bound, combined with
$e^{-\delta_1\mathsf{g}/2}\le L^{-8}$ exactly as in the case of type \textup{(top)}, gives
\begin{equation*}
  e^{-\delta_1 d_{\mathfrak{S}}(\gamma)/2}\le L^{-8(s/2-1)}=L^{-4s+8}.
\end{equation*}
If $t<s/2$, then $L^{s-t}>L^{s/2}$, and the second bound gives
\begin{equation*}
  e^{-\delta_1 d_{\mathfrak{S}}(\gamma)/2}\le \exp\bigl(-\tfrac{27}{2}\delta_1LM_1L^{s/2}\bigr).
\end{equation*}
In both cases the right side does not depend on $\gamma$, so taking the infimum over $\gamma$
yields
\begin{equation}\label{9.1:eq-damping-depth-2}
  e^{-\delta_1 d_{\mathfrak{S}}(x,y)/2}
  \le\max\Bigl\{L^{-4s+8},\ \exp\bigl(-\tfrac{27}{2}\delta_1LM_1L^{s/2}\bigr)\Bigr\}.
\end{equation}

It remains to check that $4^s$ times the right side of \eqref{9.1:eq-damping-depth-2} is
bounded in $s$. The blocking factor $L$ is an odd integer larger than one, so $L\ge3$.

For the first term,
\begin{equation*}
  4^{s}L^{-4s+8}=L^{8}\,(4/L^{4})^{s}\le L^{8}.
\end{equation*}
For the second term, put $a:=\tfrac{27}{2}\delta_1LM_1>0$. The inequality
$e^{u}\ge u^{4}/24$ for $u\ge0$ gives
\begin{equation*}
  4^{s}\exp(-aL^{s/2})\le 24\,a^{-4}\,(4/L^{2})^{s}\le 24\,a^{-4}.
\end{equation*}
Hence \eqref{9.1:eq-damping-depth-1} holds for both types of $y$ if one takes
\begin{equation*}
  C_{dp}=\max\bigl\{4,\ L^{8},\ 24\,\bigl(\tfrac{27}{2}\delta_1LM_1\bigr)^{-4}\bigr\},
\end{equation*}
which depends only on $L$, $\delta_1$ and $M_1$.
\end{proof}

\begin{corollary}[Two-cutoff tolerance in rooms]\label{9.1:cor-tolerance-rooms}
Let $\mathsf{A}$ and $\mathsf{B}$ be the two runs being compared. Assume:
\begin{itemize}
\item Theorem~\ref{9.1:thm-preliminary-reads}, in each of the two runs $\mathsf{A}$ and $\mathsf{B}$;
\item Lemma~\ref{9.1:lem-share-law};
\item the bound on the net mismatch at the cells of the entry that are not common to the two
runs, in its clause (d), and the statement at the small-field site at which that clause is read;
in the form used
below, they give at every cell $y$ of kind (top) or (dist) (the cells of the entry not common
to the two runs; the cells of kind (lab) are those at which the moves act) and of level $l$
\begin{equation}\label{9.1:eq-tolerance-rooms-1}
|\mathfrak{d}^{net}(y)|\le (1+L)\,40\,d\,C_{\mathbb{C}}\,\vartheta^{l/2}\,\mathfrak{s}_u(y),
\end{equation}
where $\mathfrak{s}_u(y)$ is the yardstick of the cell $y$ and $C_{\mathbb{C}}$ is the constant of
the chain terms in that bound;
\item the floor on the constant $\omega$ of the screening estimate used in the proof below.
\end{itemize}
Put
\begin{equation}\label{9.1:eq-tolerance-rooms-a}
\begin{aligned}
C_\times' &:= 160\,d\,(1+L)\,C\,B_\natural\,\mathsf{r}_1\,c_1\,C_{dp}^{1/2}\,\omega^{-1}C_{\mathbb{C}},\\
C_\times'' &:= 640\,d\,(1+L)\,B_\natural\,C_{\mathbb{C}}\,\omega^{-1}c_1\,\Upsilon\,\mathsf{r}_1,\\
C_{TL} &:= 2^{11}B_\natural\,C_{sl}^{\approx}\,C_\rho^{net}\,\mathsf{r}_1
+\frac{10\max\{C_\times',C_\times''\}}{\theta_\rho\beta},
\qquad 8\,C_{TL}\,\vartheta^{n_c/2}\le 1,
\end{aligned}
\end{equation}
where $C:=2\Upsilon$, with $\Upsilon=3(1+\beta_0^{pr})^3L$ as in
Lemma~\ref{9.1:lem-radii-slowness}; $\vartheta$ is the
contraction rate built on $\vartheta_2$, and $n_c$ is the classical index fixed by
$C^\natural_{cl}$ and $\vartheta_2^{1/4}$. The last inequality in \eqref{9.1:eq-tolerance-rooms-a}
is one of the coupling ceilings of the order of choice \eqref{9.1:eq-constant-order-a}, and is
assumed; that it is a ceiling on the coupling is shown at the end of the proof. Then at every
vertex pair $p\in\mathfrak{F}^V$ over a seed at the preliminary site
$(P)$, on an aligned label, with $\mathsf{v}$ the data of the seed and
$\mathsf{v}\oplus c^{\mathsf{B}}$ the data of the seed completed in run $\mathsf{B}$, the
following hold.
\begin{itemize}
\item[(a)] $\mathsf{v}\in\mathfrak{T}^{\mathsf{A}}[0]$ and
$\mathsf{v}\oplus c^{\mathsf{B}}\in\mathfrak{T}^{\mathsf{B}}[\tfrac18]$; in particular both lie in
the domain of the two-cutoff seminorm $\mathsf{S}^+$ of \eqref{9.1:eq-admissible-data-a}.
\item[(b)] For every complex $s$ with $|s|\le(8C_{TL}\vartheta^{\flat})^{-1}$, where in this
corollary $\vartheta^\flat$ stands for $\vartheta^{(j\wedge n_X)/2}$, the data
$\mathsf{v}_s:=e^{is\mathfrak{d}^{net}}(\mathsf{v}\oplus c^{\mathsf{B}})$ lie in
$\mathfrak{T}^{\mathsf{B}}[\tfrac14]$; consequently
\begin{equation}\label{9.1:eq-tolerance-rooms-2}
\bigl|\hat F^{\mathsf{B}}(\mathrm{rd}^{\mathsf{B}}\mathbf{U}')
-\hat F^{\mathsf{B}}(\mathsf{v}\oplus c^{\mathsf{B}})\bigr|
\le 16\,C_{TL}\,\vartheta^{(j\wedge n_X)/2}\,\mathfrak{n}\,e^{-\kappa_T d_j(X)} ,
\end{equation}
where $j$ is the level of the tree length $d_j(X)$, $n_X$ is a lower bound for the levels
$n_F(x)$ of the points $x$ read (see the proof), and
$\mathfrak{n}$ and the decay factor $e^{-\kappa_T d_j(X)}$, with rate $\kappa_T$ and with
$d_j(X)$ as in Notation~\ref{0.2:not-glossary}, are the factors with which the three bounds
named next are stated. Hence the large-field bound accompanying the bound on the net mismatch
named in the hypotheses, the large-field bound of the statement at the small-field site named
there, and the growth bound accompanying the definition of the classical index $n_c$ hold with
the constant $C_{LF}'C_{\mathbb{C}}$ replaced by $16C_{TL}$. The bound
\eqref{9.1:eq-tolerance-rooms-2} uses neither the constant of the $\natural$-weight bound, nor
the gauge path of the
block-centre gauge statement, nor the seminorm $|\cdot|_\natural$, nor the yardstick
$\mathfrak{s}_u$ on cells of kind (lab).
\item[(c)] Let $\Phi$ be an output. If $\Phi$ is born at a $\mathbf{T}$-step, a seed of $\Phi$ is
a datum of $\mathfrak{T}_\Phi[\tfrac78]$: its slice lies
in $\mathfrak{D}_\Phi$, the tube on which the correlation theorem runs the production of $\Phi$;
the single-run entry condition of
\eqref{9.1:eq-entry-conditions-a}, in its clause (e), at the seeds of $\Phi$ is this assertion.
For an output born at an $\mathbf{R}$-step the corresponding assertion rests on
\cite[(1.65)]{B-CMP122b} and is not part of this corollary.
Every stored term $F$ read in the production of $\Phi$ is read at points of
$\mathfrak{D}_F[0]$. The room fraction $\varsigma$ is counted, for each term, in that term's own
room, so that the fractions do not compound.
\end{itemize}
\end{corollary}

\begin{proof}
Throughout, $\mathbf{U}'$ is the point of run $\mathsf{B}$ at the vertex pair $p$ and
$S^{\mathsf{B}}$ is the slice of this point. At $s=1$ one has
$\mathsf{v}_1=\mathrm{rd}^{\mathsf{B}}\mathbf{U}'$, and at $s=0$ one has
$\mathsf{v}_0=\mathsf{v}\oplus c^{\mathsf{B}}$.

\emph{Reduction to a bound on the mismatch.} We apply Lemma~\ref{9.1:lem-box-response-lipschitz}
in run $\mathsf{B}$, the data increments being given by $B=(1/i)\log V'$ as in
\cite[Proposition~9]{B-CMP102b}, together with the corresponding item of the operator theorem at
free current. It gives that the slice of $\mathsf{v}_s$ lies within
$|1-s|\,D$ of $S^{\mathsf{B}}$, componentwise, where for each table member $a$, with $e_a$ the
exponent of $\ell_{\mathfrak{S}}$ attached to $a$,
\begin{equation}\label{9.1:eq-tolerance-rooms-3}
D^a:=2\,\ell_{\mathfrak{S}}^{-e_a}\,B_\natural\,\bigl(|\mathfrak{d}^{net}|\bigr)^{\approx},
\end{equation}
the profile $(\cdot)^{\approx}$ being smeared with $e^{-2\delta_1 d_{\mathfrak{S}}}$. We bound
$D$ separately at $(=)$-points and at $(\mathrm{up})$-points, by the kind of cell contributing.
These are the two classes of points of Theorem~\ref{9.1:thm-preliminary-reads}: at a
$(=)$-point $x$ the level $\lambda_{\mathfrak{S}}(x)$ equals $n_F(x)$, the level at which the own
units $\mathsf{u}_F$ of the stored term $F$ are taken at $x$, and at a
$(\mathrm{up})$-point the current label $\iota(x)$ exceeds $n_F(x)$.

\emph{$(=)$-points, cells of kind (lab).} Such a cell has level $l\ge j\wedge n_X$. By
Lemma~\ref{9.1:lem-share-law}, $|\mathfrak{d}^{net}(c)|\le C_\rho^{net}\vartheta^{l/2}\mathsf{a}(c)$;
by Lemma~\ref{9.1:lem-radii-slowness} and by (acc) in \eqref{9.1:eq-preliminary-reads-a},
the contribution of these cells to $D^a$ is at most
$800\,B_\natural C_{sl}^{\approx}C_\rho^{net}\mathsf{r}_1\,\vartheta^\flat\,\mathsf{G}^a$.

\emph{$(=)$-points, cells of kind (top) or (dist).} On these cells we use
\eqref{9.1:eq-tolerance-rooms-1}. By the pointwise form \eqref{9.1:eq-radii-slowness-a} of
Lemma~\ref{9.1:lem-radii-slowness}, read with the constant $C=2\Upsilon$ that enters $C_\times'$,
\begin{equation}\label{9.1:eq-tolerance-rooms-4}
\mathfrak{s}_u(y)\le 2\cdot 2\Upsilon\cdot e^{\delta_1 d_{\mathfrak{S}}(x,y)/4}\,
(\ell_{\mathfrak{S}}\bar u^0_\nu)(x).
\end{equation}
We split the smearing weight as
\begin{equation*}
e^{-2\delta_1 d_{\mathfrak{S}}}=e^{-\delta_1 d_{\mathfrak{S}}}\cdot e^{-\delta_1 d_{\mathfrak{S}}/4}
\cdot e^{-\delta_1 d_{\mathfrak{S}}/4}\cdot e^{-\delta_1 d_{\mathfrak{S}}/2},
\end{equation*}
with $d_{\mathfrak{S}}=d_{\mathfrak{S}}(x,y)$. The first factor screens: by the screening
property (d1) of the
scaled distance $d_{\mathfrak{S}}$, which is also an input of Lemma~\ref{9.1:lem-radii-slowness},
it decays in
$|\lambda_{\mathfrak{S}}(x)-l|$, where $\lambda_{\mathfrak{S}}(x)=n_F\ge n_X$, and together with
the floor on $\omega$ this gives
\begin{equation*}
e^{-\delta_1 d_{\mathfrak{S}}}\vartheta^{l/2}\le\omega^{-1}\vartheta^{(j\wedge n_X)/2}
=\omega^{-1}\vartheta^\flat .
\end{equation*}
The second factor
cancels the growth $e^{\delta_1 d_{\mathfrak{S}}/4}$ in \eqref{9.1:eq-tolerance-rooms-4}. The third
factor is the
square root of the bound of Lemma~\ref{9.1:lem-damping-depth}: since
$e^{-\delta_1 d_{\mathfrak{S}}/2}\le C_{dp}4^{-s}$ with $s=j-n_F$, one has
$e^{-\delta_1 d_{\mathfrak{S}}/4}\le C_{dp}^{1/2}2^{-s}$. The fourth factor is summed over $y$ by
the summability property (d2) of $d_{\mathfrak{S}}$.
Neither the constant $K_u$ nor a further item of the operator theorem at free current is used
here. These cells contribute at most
\begin{equation*}
C_\times'\,\vartheta^\flat\,2^{-s}\,\mathsf{u}_F^a
=\frac{C_\times'}{\theta_\rho\beta}\,\vartheta^\flat\,\varrho_F\,\mathsf{u}_F^a ,
\end{equation*}
the constant $C_\times'$ of \eqref{9.1:eq-tolerance-rooms-a} containing $2\Upsilon$ as its
constant $C$, and the equality being $\varrho_F=\theta_\rho\beta\,2^{-s}$, with $\varrho_F$ as in
\eqref{9.1:eq-own-tube-a}. Since
$800\le 2^{11}$ and $C_\times'/(\theta_\rho\beta)\le C_{TL}$ by \eqref{9.1:eq-tolerance-rooms-a},
\begin{equation}\label{9.1:eq-tolerance-rooms-5}
D\le C_{TL}\,\vartheta^\flat\,(\mathsf{G}+\varrho_F\mathsf{u}_F)\quad\text{at $(=)$-points,}
\qquad
S^{\mathsf{B}}\le(\mathsf{c}_F^{+}-\varrho_F)\,\mathsf{u}_F-\tfrac78\,\mathsf{G},
\end{equation}
the second bound by clause (ii) of Theorem~\ref{9.1:thm-preliminary-reads} and by
Lemma~\ref{9.1:lem-schedule-room}. Put $R:=(8C_{TL}\vartheta^\flat)^{-1}$, which satisfies
$R\ge1$. For
$|s|\le R$ one has $|1-s|\le 2R$, so $|1-s|\,D\le\frac14(\mathsf{G}+\varrho_F\mathsf{u}_F)$: the
slice of $\mathsf{v}_s$ lies in $\mathfrak{D}_F[\tfrac14]$. At $s=0$ the factor $|1-s|$ equals $1$
and $C_{TL}\vartheta^\flat\le\frac18$, which gives the fraction $\frac18$.

\emph{$(\mathrm{up})$-points.} Here $S^{\mathsf{B}}\le 0.755\,\mathsf{u}_F$ by clause (iii) of
Theorem~\ref{9.1:thm-preliminary-reads} and (k6) in \eqref{9.1:eq-room-checks-a}. On cells of kind
(lab) the argument of the $(=)$-case applies, with the share majorant bounded by
$\mathsf{a}\le C_\Sigma K_u\alpha_1$ in the units of $F$, $C_\Sigma$ being the summation
constant; since $2C_\Sigma K_u\le125$, the
contribution is at most $125\,B_\natural C_{sl}^{\approx}C_\rho^{net}\mathsf{r}_1\vartheta^\flat
\mathsf{u}_F$. On cells of kind (top) or (dist) we use the same split of
$e^{-2\delta_1 d_{\mathfrak{S}}}$ as at
$(=)$-points: the first factor screens, the second cancels the slowness growth, the third is
bounded by $1$, and the fourth is summed by the summability property (d2). Now the
slowness factor delivers, by the
pointwise form \eqref{9.1:eq-radii-slowness-a}, the profile at $x$ in the current units of $x$,
free of shares:
\begin{equation*}
\begin{aligned}
D&\le 2B_\natural(1+L)40dC_{\mathbb{C}}\,\omega^{-1}c_1\cdot2\cdot2\Upsilon\cdot\vartheta^\flat
\cdot\ell_{\mathfrak{S}}(x)^{-e_a}(\ell_{\mathfrak{S}}\bar u^0_\nu)(x)\\
&=\frac{C_\times''}{2\mathsf{r}_1}\,\vartheta^\flat\cdot
\ell_{\mathfrak{S}}(x)^{-e_a}(\ell_{\mathfrak{S}}\bar u^0_\nu)(x).
\end{aligned}
\end{equation*}
The remaining conversion is the unit conversion of clause (iii) of
Theorem~\ref{9.1:thm-preliminary-reads}, the same one that gave
$S^{\mathsf{B}}\le0.755\,\mathsf{u}_F$. The local radius $\bar u^0_\nu(x)$ at the current label
$\iota=\iota(x)>n_F(x)$ is of the type
$\ell_\circ^{2(\iota-n_0)}\alpha_{0,\iota}\ell_\iota^{-1}$, where $n_0\ge n_F$ is the level from
which the power of $\ell_\circ$ is counted; writing it in this form,
\begin{equation*}
\begin{aligned}
\ell_{\mathfrak{S}}(x)^{-e_a}(\ell_{\mathfrak{S}}\bar u^0_\nu)(x)
&=\Bigl(\frac{\ell_{n_F}}{\ell_{\mathfrak{S}}}\Bigr)^{e_a-1}
\frac{\ell_{n_F}}{\ell_\iota}\,\ell_{n_F}^{-e_a}\,\ell_\circ^{2(\iota-n_0)}\alpha_{0,\iota}\\
&\le \ell_{n_F}^{-e_a}\,L^{-(\iota-n_F)}\ell_\circ^{2(\iota-n_0)}\alpha_{0,\iota}\\
&\le\tfrac34\,\alpha_{0,n_F}\,\ell_{n_F}^{-e_a}\le\tfrac34\,\mathsf{u}_F^a(x).
\end{aligned}
\end{equation*}
The first inequality holds because $n_F\le\lambda_{\mathfrak{S}}\le\iota$, so that both length
ratios are at most $1$ ($e_a\ge1$). The second is proved in the two cases of clause (iii) of
Theorem~\ref{9.1:thm-preliminary-reads}. If $n_0>n_F$, put $u:=n_0-n_F\ge1$ and $v:=\iota-n_0$;
the displayed line of that conversion reads
\begin{equation*}
3L^{-(\iota-n_F)}\ell_\circ^{2(\iota-n_0)}\alpha_{a,\iota}
\le K_uL^{-u}3^{-v}(u+v)^{1/2}\alpha_{a,n_F}\le\frac{K_u}{L}\,\alpha_{a,n_F},
\end{equation*}
and the floor $(F_u)$ of \eqref{9.1:eq-entry-conditions-a}, $K_u\le L/2$, gives the factor
$\frac16\le\frac34$. If $n_0=n_F$ (the
point was raised in the chain), put $v:=\iota-n_0=\iota-n_F\ge1$; the power of $\ell_\circ$ given
by the unit rule (s2) of the summation floors, which are a hypothesis of
Lemma~\ref{9.1:lem-stage-room}, and the comparison between the radii
$\alpha_{0,\cdot}$ and $\alpha_{a,\cdot}$, give, with
$\ell_\circ^2\le L/3$,
\begin{equation*}
\Bigl(\frac{\ell_\circ^2}{L}\Bigr)^{v}\tfrac94\,v^{1/2}
\le\sup_{v\ge1}3^{-v}\tfrac94\,v^{1/2}=\tfrac34 .
\end{equation*}
The last inequality of the chain is $\alpha_{0,n_F}\le\alpha_{a,n_F}$, by the same comparison
between the radii. Therefore on cells of kind (top) or (dist)
\begin{equation*}
D\le\frac{3C_\times''}{8\mathsf{r}_1}\,\vartheta^\flat\mathsf{u}_F\le C_\times''\vartheta^\flat
\mathsf{u}_F,
\end{equation*}
since $\mathsf{r}_1\ge1$. Lemma~\ref{9.1:lem-damping-depth} and the depth
$s=j-n_F$ are not used on these cells: the third factor $e^{-\delta_1 d_{\mathfrak{S}}/4}$, which
at $(=)$-points was the square root of the bound of Lemma~\ref{9.1:lem-damping-depth}, is here
bounded by $1$. The
constant $K_u$ enters only through the floor $(F_u)$ of \eqref{9.1:eq-entry-conditions-a} in the
displayed line above, whose supremum over $(u,v)$ is attained at $(1,0)$, so that the bound is
free of the depth, by (k5) in \eqref{9.1:eq-room-checks-a}; no bound here depends on
$\iota(x)-n_F(x)$. Hence at
$(\mathrm{up})$-points
\begin{equation*}
D\le\bigl(125\,B_\natural C_{sl}^{\approx}C_\rho^{net}\mathsf{r}_1+C_\times''\bigr)\vartheta^\flat
\mathsf{u}_F\le\frac{C_{TL}}{10}\,\vartheta^\flat\mathsf{u}_F ,
\end{equation*}
because $125\le2^{11}/10$ and, as $\theta_\rho\beta\le\frac1{40}$,
$10C_\times''/(\theta_\rho\beta)\ge400\,C_\times''\ge C_\times''$; the $(=)$-case needs
$10C_\times'/(\theta_\rho\beta)$, and both are served by the maximum in
\eqref{9.1:eq-tolerance-rooms-a}. For $|s|\le R$,
\begin{equation*}
|1-s|\,D\le2R\cdot\frac{C_{TL}}{10}\,\vartheta^\flat\mathsf{u}_F=\frac{\mathsf{u}_F}{40}
\le\bigl(\mathsf{c}_F^{+}-\tfrac34\varrho_F-0.755\bigr)\mathsf{u}_F ,
\end{equation*}
the last inequality by (k7) in \eqref{9.1:eq-room-checks-a}:
$0.8-0.01875-0.755\ge\frac1{40}$.

\emph{The radius condition.} The data $\mathsf{v}_s$ stay in the ball of radius $r_c$ required by
Definition~\ref{9.1:def-admissible-data}: the yardstick $\mathfrak{s}_u$ carries the factor
$L^{l-\iota}$, and $\vartheta^{l/2}L^{l-\iota}\le\vartheta^{\iota/2}\le\vartheta^\flat$, so that
$R\sup|\mathfrak{d}^{net}|\le C'(C_{\mathbb{C}}+C_\rho^{net})\bar{\bar u}/C_{TL}$, with a constant
$C'$ distinct from the constant $C=2\Upsilon$ above. Together with the
two cases above this gives $\mathsf{v}_s\in\mathfrak{T}^{\mathsf{B}}[\tfrac14]$ for $|s|\le R$, and
$\mathsf{v}\oplus c^{\mathsf{B}}\in\mathfrak{T}^{\mathsf{B}}[\tfrac18]$.

\emph{The bound \eqref{9.1:eq-tolerance-rooms-2}.} Consider the map
$s\mapsto\hat F^{\mathsf{B}}(\mathsf{v}_s)-\hat F^{\mathsf{B}}(\mathsf{v}_0)$ on the closed disc
$|s|\le R$, on which all $\mathsf{v}_s$ lie in $\mathfrak{T}^{\mathsf{B}}[\tfrac14]$, hence in
the domain of $\mathsf{S}^+$. It is holomorphic in $s$, it vanishes at $s=0$, and its modulus is
at most twice the supremum of $|\hat F^{\mathsf{B}}|$ on $\mathfrak{T}^{\mathsf{B}}[\tfrac14]$.
The Schwarz lemma on the open disc $|s|<R$ bounds its modulus by $2|s|/R$ times this supremum,
and by continuity the same bound holds on $|s|=R$. At $s=1\le R$, where
$\mathsf{v}_1=\mathrm{rd}^{\mathsf{B}}\mathbf{U}'$, this is
\begin{equation*}
\begin{aligned}
\bigl|\hat F^{\mathsf{B}}(\mathrm{rd}^{\mathsf{B}}\mathbf{U}')
-\hat F^{\mathsf{B}}(\mathsf{v}\oplus c^{\mathsf{B}})\bigr|
&\le\frac2R\,\sup_{\mathfrak{T}^{\mathsf{B}}[1/4]}|\hat F^{\mathsf{B}}|\\
&=16\,C_{TL}\,\vartheta^\flat\,\sup_{\mathfrak{T}^{\mathsf{B}}[1/4]}|\hat F^{\mathsf{B}}| .
\end{aligned}
\end{equation*}
With the bound for this supremum that enters the two-cutoff seminorm of
\eqref{9.1:eq-admissible-data-a}, and with $\vartheta^\flat=\vartheta^{(j\wedge n_X)/2}$, this
gives \eqref{9.1:eq-tolerance-rooms-2}.

\emph{Run $\mathsf{A}$ and clause (c).} The membership $\mathsf{v}\in\mathfrak{T}^{\mathsf{A}}[0]$
is clause (v) of Theorem~\ref{9.1:thm-preliminary-reads} in run $\mathsf{A}$. In clause (c), the
assertion that the stored terms $F$ entering the production of $\Phi$ are read at points of
$\mathfrak{D}_F[0]$ is Theorem~\ref{9.1:thm-preliminary-reads} together with
Lemma~\ref{9.1:lem-schedule-room}. The corresponding assertion at the seeds of an output born at an
$\mathbf{R}$-step is not part of this corollary; it rests on \cite[(1.65)]{B-CMP122b}.

\emph{The ceiling.} The last condition in \eqref{9.1:eq-tolerance-rooms-a} is a ceiling on the
coupling: on aligned labels
$\vartheta^{n_c}\le(\log x_{n_c})^{-c_2}/C^\natural_{cl}$, where $x_{n_c}$ is the inverse squared
coupling at the level $n_c$ of the label, by clause (D) of the classical comparison of the coupling
with the contraction rates on aligned labels.
\end{proof}

\begin{remark}[No condition bounds a difference of levels]
No condition in the argument above bounds a difference of levels. Each mismatch is measured
against the room of its own origin: the moves against the stage rooms of which they are shares,
with ratio $\vartheta^{l/2}$ and independently of the depth; the cells of the entry that are not
common to the two runs against $\varrho_F$, whose factor $2^{-s}$ is paid by the factor $4^{-s}$
of Lemma~\ref{9.1:lem-damping-depth} at $(=)$-points, and against the complement of
$\frac34\mathsf{u}_F$ at $(\mathrm{up})$-points through the unit conversion of clause (iii) of
Theorem~\ref{9.1:thm-preliminary-reads}, where no depth enters. With $n$ the length of run
$\mathsf{A}$, $h$ the level from which the stage tables are counted and $j$ the level of the tree
length $d_j(X)$, the differences $h-n$, $j-n$ and $\iota-n_F$, and the index $N$ of the stage
tables, occur in no condition. The data $\mathsf{v}\oplus c^{\mathsf{B}}$
are compared inside run $\mathsf{B}$ and measured from the point of run $\mathsf{B}$ itself; no
passage from the membership of the slice of run $\mathsf{A}$ to that of run $\mathsf{B}$, which
would cost a factor $\vartheta^n2^s$, is used,
and the fraction $\tfrac18$ of clause (a) and the further fraction $\tfrac18$ allowed to the data
$\mathsf{v}_s$ in clause (b) add to the fraction $\tfrac14$ of clause (b).
\end{remark}

\begin{proposition}[Single-scale leaves: the old polydisc lies in the admissible data]
\label{9.1:prop-single-scale-leaves}
Let $X$ be the region of a single-scale leaf. This means a leaf of one of the kinds $E$,
$P^{\flat}$, $Q$, $R$, $\hat{A}$ for which $X\subset\Lambda_j$, where $\Lambda_j$ is the
complement of the large-field region $Z_j$ at level $j$, and which has no lower layer. Let
$X^{+}$ be the hull of $X$, and let $F$ be the term stored at the leaf. Then, in each of the two
runs $\mathsf{m}\in\{\mathsf{A},\mathsf{B}\}$,
\begin{equation}\label{9.1:eq-single-scale-leaves-1}
\mathbb{D}_j(X^{+};\alpha_j^s)\subset\mathfrak{T}^{\mathsf{m}}[7/8](X).
\end{equation}
Here $\mathbb{D}_j(X^{+};\alpha_j^s)$ is the localisation domain of level $j$ on $X^{+}$ at the
parameter $\alpha_j^s$, a polydisc of data; we call it the old polydisc. The set
$\mathfrak{T}^{\mathsf{m}}[7/8](X)$ is the class of admissible
chart data of Definition~\ref{9.1:def-admissible-data}, see
\eqref{9.1:eq-admissible-data-a}, at room fraction $\varsigma=7/8$.
\end{proposition}

\begin{proof}
Fix a run $\mathsf{m}$; the argument is the same for $\mathsf{m}=\mathsf{A}$ and
$\mathsf{m}=\mathsf{B}$. Let $\mathsf{v}\in\mathbb{D}_j(X^{+};\alpha_j^s)$. Of the conditions
of Definition~\ref{9.1:def-admissible-data}, see \eqref{9.1:eq-admissible-data-a}, the one at
issue here is that the slice $(\mathbf{V}(B),J(\mathbf{V}(B))){\upharpoonright}X$ attached to
$\mathsf{v}$ lie in the dilated own tube $\mathfrak{D}^{\mathsf{m}}_F[7/8]$
\eqref{9.1:eq-own-tube-a} of $F$; we verify it.

We use the leaf form of clause (b) of the two-cutoff tolerance bound, which is stated on the hull
$X^{+}$. On a single-scale leaf it coincides with the inclusion of the slices into the own tube
for the region $X$, read on $X^{+}$; since the leaf has no lower layer, this is a single-scale
statement. Both follow from the clause of the correlation theorem that places the slice of every
datum of the level-$j$ polydisc in the own tube at three quarters of its radii, together with a
further clause of the same theorem entering that estimate. These are combined with
\cite[Lemma~4]{B-CMP109}, which
controls the ratios of the radii entering at level $j$, and with the operator theorem at free
current. What the bound asserts is the following: for every
$\mathsf{v}\in\mathbb{D}_j(X^{+};\alpha_j^s)$, the slice attached to $\mathsf{v}$ lies in the
own tube of $F$ at three quarters of its radii. In the own units of $F$ these radii carry the
coefficient $\tfrac34\mathsf{c}_F^{+}$.

The radii of the dilated tube $\mathfrak{D}^{\mathsf{m}}_F[7/8]$ carry the coefficient
$\mathsf{c}_F^{+}-\varrho_F/8$, where $\varrho_F$ is the room of the own tube
\eqref{9.1:eq-own-tube-a}. We have
\begin{equation}\label{9.1:eq-single-scale-leaves-2}
\tfrac34\,\mathsf{c}_F^{+}\le\mathsf{c}_F^{+}-\tfrac18\,\varrho_F ,
\end{equation}
an inequality equivalent to $\varrho_F\le 2\,\mathsf{c}_F^{+}$, which holds since the room
$\varrho_F$ of the own tube is at most twice the coefficient $\mathsf{c}_F^{+}$.
Hence the tube at three quarters of the radii is contained in
$\mathfrak{D}^{\mathsf{m}}_F[7/8]$. The slice attached to $\mathsf{v}$ therefore lies in
$\mathfrak{D}^{\mathsf{m}}_F[7/8]$, and $\mathsf{v}\in\mathfrak{T}^{\mathsf{m}}[7/8](X)$.
This proves \eqref{9.1:eq-single-scale-leaves-1}.
\end{proof}

\begin{remark}[Statements evaluated on the old polydisc]
Consider the two-cutoff seminorm $\mathsf{S}_j^{+}[F]$ of
\eqref{9.1:eq-admissible-data-a}, which is formed over the admissible data. By
\eqref{9.1:eq-single-scale-leaves-1}, on single-scale leaves it dominates the seminorm formed
at the data of the old polydisc $\mathbb{D}_j(X^{+};\alpha_j^s)$. Consequently, every statement
that evaluates differences at data of that polydisc applies without change.
Items (a)--(c) below list the statements that evaluate differences at data of the polydisc,
and how each is covered; item (d) lists statements whose conclusion is a bound by a supremum
over seeds, which Proposition~\ref{9.1:prop-single-scale-leaves} does not cover.
\begin{itemize}
\item[(a)] The statements of this chapter evaluated at the following objects, among them the
estimate on the hull $X^{+}$ and the estimate on the region $X$, are covered: complex
configurations $\mathbf{U}\in\mathcal{U}_j^c(Y)$, the
class of complex configurations on $Y$ at level $j$ attached to a set $c$, in which the radius
enters only through the ratio $\alpha^c/\alpha^s$; the data $\beta^{\mathsf{m}}(\sigma')$ at the
localisation parameter $\sigma'$, which lie in $\mathbb{D}_j$ and are multi-affine in
$\mathsf{s}$ on $|\mathsf{s}|\le e^{\kappa_1'}$; and families on the whole lattice. For
all of these the proposition applies directly. The family $G^{+}$ at which they are evaluated
lives on one lattice and is never a seed. Ratios of the radii are therefore unrestricted, and no
ladder of levels is climbed.
\item[(b)] The leaf form of clause (b) of the two-cutoff tolerance bound places the slice of
every $\mathsf{v}\in\mathbb{D}_j$ in the own tube at three quarters of the radii. On single-scale
leaves it is the clause of the correlation theorem used in the proof above. For leaves with a
lower layer at a level $l<j$, the form of this bound carrying the weight $L^{j-l}$ is not used.
Every statement of this chapter that uses this bound for such leaves, the statement on
block-centre gauges invoked in Definition~\ref{9.1:def-admissible-data} among them, uses
instead clause (b) of Corollary~\ref{9.1:cor-tolerance-rooms}, see
\eqref{9.1:eq-tolerance-rooms-a}.
\item[(c)] Consider the statements of this chapter evaluated at comb data about the
configuration $V_\circ=1$ with the cube $\square_0$ contained in the core, among them the
bounds on comb data for single pieces and for groups of pieces. For these the
proposition applies, together with the inclusion stated next in this section. For pieces of
the second class, the trivial bound is used together with clause (b) of
Corollary~\ref{9.1:cor-tolerance-rooms}.
\item[(d)] The statements of this chapter whose conclusion is a bound by a supremum over seeds,
among them the statements of the one-lattice comparison, are not evaluated at data of the
polydisc. They are covered, for outputs $\Phi$ born at a $\mathbf{T}$-step, by clause (c) of
Corollary~\ref{9.1:cor-tolerance-rooms}: their proofs fix a seed as a parameter and run the
production of the correlation theorem on the domain $\mathfrak{D}_\Phi$ of the output $\Phi$.
Outputs born at an $\mathbf{R}$-step are treated in the section on $\mathbf{R}$-born outputs.
\end{itemize}
\end{remark}

\begin{remark}[The analyticity-radius constant is numerical]\label{9.1:rem-analyticity-constant}
The constant $C_S$ in the two-cutoff tolerance in rooms, display~\eqref{9.1:eq-tolerance-rooms-a},
is the number
\begin{equation}\label{9.1:eq-analyticity-constant-1}
C_S := 4 .
\end{equation}
Accordingly the constant $C_{rim}$ is $C_{rim} := 8C_S^2C_{X^+}$, where $C_{X^+}$ is the constant
carried by the estimates on the hull $X^+$ of the leaf. The constant of the corollary that reads the
two-cutoff seminorm
$\mathsf{S}^+$ on the hull $X^+$ is taken at $C_S^2C_{X^+}$. The constant $K_{\mathsf{l}}$, which
carries $C_S$ in the same way, is likewise taken at $C_S^2C_{X^+}$. The constant $C_S$ imposes no
condition on any parameter. It does not depend on the constants $C_e$ and $C_\Sigma$ of the
tolerance estimates in rooms, nor on $B_\natural$.
\end{remark}

\begin{proof}
The constant $C_S$ enters in two places. It appears in clause~(b) of
display~\eqref{9.1:eq-tolerance-rooms-a}, where it has the standing of a condition on the constants;
since $C_S$ is the number $4$, that clause imposes through $C_S$ no condition on any parameter. It
also measures the loss of radius in the Schwarz-lemma
estimate that compares the two runs on the datum of run $\mathsf{B}$. In that estimate the room
$2\ell\sup\bar u$, formed from the supremum of the local radii $\bar u$ with the multiplier $\ell$
fixed in that estimate, is taken as the majorant of the data, and the datum of run $\mathsf{B}$ is
majorised by $(1+2C_*)\mathfrak{s}_u$, a multiple of the data scale $\mathfrak{s}_u$ of that
estimate. This costs the factor $(1+2C_*)^2$ on $S_u^{-2}$, where $S_u$
is the radius of that estimate.

In the comparison made here, the term $P^2$ formed from the production block $P$ is read only
at the datum common to the two runs $\mathsf{A}$ and $\mathsf{B}$: by the statement that the
comparison law is void on this route, the mismatch $\mathfrak{d}$ between the data of the two runs
vanishes, $\mathfrak{d}=0$. Hence the datum of run $\mathsf{B}$ is not majorised separately, and
the factor $(1+2C_*)^2$ does not arise. We put
\begin{equation}\label{9.1:eq-analyticity-constant-2}
\psi(\sigma) := \bigl[\hat F^{\mathsf{B}}-\hat F^{\mathsf{A}}\bigr]
\bigl(X;e^{i\sigma(\beta^{\mathsf{A}}\oplus c)}\bigr),
\end{equation}
where $\hat F^{\mathsf{m}}$ is the chart function of run $\mathsf{m}$ of
Definition~\ref{9.1:def-admissible-data}. The argument is the one-parameter family of data
obtained by exponentiating the vertex pair $\beta^{\mathsf{A}}\oplus c$ formed from the data
coordinate $\beta^{\mathsf{A}}$ of run $\mathsf{A}$ and the vertex entry $c$.
We also put
\begin{equation}\label{9.1:eq-analyticity-constant-3}
S_u^\natural := \sup\Bigl\{S:\ e^{i\sigma(\beta^{\mathsf{A}}\oplus c)}
\text{ lies in the domain of } \mathsf{S}^+ \text{ for all } |\sigma|\le S\Bigr\},
\end{equation}
with $\mathsf{S}^+$ the two-cutoff seminorm of Definition~\ref{9.1:def-admissible-data}.

Two inputs bound $S_u^\natural$ from below:
\begin{itemize}
\item Proposition~\ref{9.1:prop-single-scale-leaves}: on single-scale leaves, the polydisc
$\mathbb{D}_j(X^+;\alpha_j^s)$ of radius $\alpha_j^s$ on the hull $X^+$ lies in the admissible
data $\mathfrak{T}^{\mathsf{m}}[7/8]$ for both runs.
\item The clause of the correlation theorem that supplies the constant $c_3$, applied at
$2\mathfrak{a}_n$ in each of the two runs. That the majorant $M_J$ furnished by that clause is the
same for the two runs and does not depend on the source coefficients $z$ is a separate input, the
statement that asserts this common, $z$-free form of the majorant.
\end{itemize}
Together they give
\begin{equation}\label{9.1:eq-analyticity-constant-4}
S_u^\natural \ \ge\ \frac{c_3c_s\alpha_j^\flat}{2M_J}\ =\ \frac{S_J}{4}
\qquad\text{at } c_s=\tfrac12 .
\end{equation}
Here $c_3$ is the constant and $\alpha_j^\flat$ the level-$j$ radius of that clause of the
correlation theorem. The fraction $c_s$ enters that clause; the bound is used at the value
$c_s=\tfrac12$, which is the value fixed for it in the statement that introduces $c_s$. Further,
\begin{equation*}
S_J := \frac{\alpha_{3,j}}{\sup|B^{ax}|},
\end{equation*}
where $\alpha_{3,j}$ is the analyticity radius of the analyticity lemma accompanying the correlation
theorem and $B^{ax}$ the field of that lemma; $S_J$ is this radius measured in units of
$\sup|B^{ax}|$. The quantity $\mathfrak{N}_J(u)$ is
computed through $S_J$.

By~\eqref{9.1:eq-analyticity-constant-4} and~\eqref{9.1:eq-analyticity-constant-1},
\begin{equation*}
(S_u^\natural)^{-1}\le 4S_J^{-1}=C_S\,S_J^{-1},
\qquad\text{hence}\qquad
(S_u^\natural)^{-2}\le C_S^2\,S_J^{-2}.
\end{equation*}
This is the factor $C_S^2$ carried by $C_{rim}=8C_S^2C_{X^+}$. The same factor is carried by the
constant of the corollary that reads $\mathsf{S}^+$ on the hull $X^+$, and by $K_{\mathsf{l}}$.

The case $S_u^\natural\le 1$ falls under the trivial alternative in the statement of the
correlation theorem.

The lower bound~\eqref{9.1:eq-analyticity-constant-4} contains none of $C_e$, $C_\Sigma$,
$B_\natural$. Therefore $C_S=4$ is a fixed number rather than a condition, and it reads none of these
constants.
\end{proof}

\begin{remark}[Level-free ratio of the run radii]\label{9.1:rem-radii-ratio}
In both runs $\mathsf{A}$, $\mathsf{B}$ the two radii at level $j$, indexed by $a\in\{0,1\}$ as in
\cite[(2.28)]{B-CMP119} and as used in the operator theorem at free current, are
\begin{equation}\label{9.1:eq-radii-ratio-1}
  \alpha^{\mathrm{run}}_{a,j}:=\alpha_{a,j}=g_j\,C_a\,(\log x_j)^{q},
\end{equation}
and we put $\alpha^{\flat}_j:=\min_a\alpha_{a,j}=\alpha_{0,j}$. Then
\begin{equation}\label{9.1:eq-radii-ratio-2}
  \frac{\alpha^{\mathrm{run}}_{1,j}}{\alpha^{\flat}_j}=\frac{C_1}{C_0}=\mathsf{r}_1
  \qquad\text{at every level } j .
\end{equation}
This ratio enters the constant $K_o$ of the entry conditions~\eqref{9.1:eq-entry-conditions-a} and
the constant $C_{TL}$ of Corollary~\ref{9.1:cor-tolerance-rooms}. Through these two constants it
enters ceilings on $\gamma$. Radii at two different levels $\iota>j$ occur only in the
quotient $\alpha_{a,\iota}/\alpha_{a,j}$, in which the two radii carry the same index $a$, so that
the ratio \eqref{9.1:eq-radii-ratio-2} does not enter.
The constant $C_{\mathrm{rim}}$ of this
chapter does not depend on the ratio. The equality $q_0=q_1$ of Ba{\l}aban's two exponents is used.
\end{remark}

\begin{proof}
\cite{B-CMP119} takes $q_0$, $q_1$ to be integers greater than $1$ and imposes no relation between
them. We use the radii of \cite[(2.28)]{B-CMP119} under the condition $q_0=q_1=:q$, a condition of
the correlation theorem, recorded with the radius exponents in
Notation~\ref{9.1:not-radius-exponent-floors}; under it the radii take the form
\eqref{9.1:eq-radii-ratio-1}.

We divide the radius with $a=1$ by the one with $a=0$. The factors $g_j$ cancel, and so do the
factors $(\log x_j)^q$, because the exponent is the same. What remains is $C_1/C_0=\mathsf{r}_1$.
The second equality in the definition of $\alpha^{\flat}_j$ amounts to $C_0\le C_1$, that is
$\mathsf{r}_1\ge1$, because the two radii share the factor $g_j(\log x_j)^q$; since
$\alpha^{\flat}_j=\alpha_{0,j}$, this gives \eqref{9.1:eq-radii-ratio-2}. As a function of
$\log x_j$ the ratio has degree $0$. It therefore takes the same value at every level.

The constants $K_o$ and $C_{TL}$ contain the ratio only through the number $\mathsf{r}_1$, which
does not depend on the level. In the definition of $C_{TL}$ in
Corollary~\ref{9.1:cor-tolerance-rooms}, $\mathsf{r}_1$ appears explicitly, as a factor of the
first summand, and also inside the constant $C_\times'$. The
conditions in which $K_o$ and $C_{TL}$ appear are ceilings on $\gamma$.

In the quotient $\alpha_{a,\iota}/\alpha_{a,j}$ numerator and denominator carry the same index
$a$, so the constants $C_a$ cancel in it and \eqref{9.1:eq-radii-ratio-2} does not enter there; nor
does the constant $C_{\mathrm{rim}}$ read it.

The condition $q_0=q_1$ cannot be dropped. If instead the two radii carried exponents $q_0\ne q_1$,
$\alpha_{a,j}=g_jC_a(\log x_j)^{q_a}$, the same division would give
\begin{equation}\label{9.1:eq-radii-ratio-3}
  \frac{\alpha_{1,j}}{\alpha_{0,j}}=\mathsf{r}_1\,(\log x_j)^{\,q_1-q_0},
\end{equation}
which depends on $j$. The constant $K_o$ would then depend on the level. Thus the equality
$q_0=q_1$ is used here, as it is in the single-run entry conditions of
Definition~\ref{9.1:def-entry-conditions}.
\end{proof}

\begin{remark}[Numerical checks for the room estimates]\label{9.1:rem-room-checks}
The constants of the room estimates of this section are chosen in the order of
Notation~\ref{9.1:not-constant-order}. Every condition imposed at a step of that order is
one-sided, and it involves only constants chosen at that step or earlier.

First come the numbers $\beta_I$, $\theta_S$, $\theta_\rho$, $K_u$, $C_\Sigma$, $C_S$.
Then $L$ is chosen, subject to a lower bound that reads only these numbers and is among the
hypotheses of Theorem~\ref{9.1:thm-preliminary-reads}. Then come
$\delta_1$, $\kappa$, $\kappa_1$. Then $M_1$ is chosen, subject to two lower bounds, one of
which is a hypothesis of Corollary~\ref{9.1:cor-tolerance-rooms}; after it comes $M$.
Then $\mathsf{r}_1=C_1/C_0$ is chosen, subject to the lower bound
$\mathsf{r}_1\ge 2B_3\,c'M$, where $B_3$ and $c'$ are constants that do not depend on $M$.
Then come
$B_\natural$, $C_{sl}^{\approx}$, $C_{dp}$;
then $K_o$; then $c_M$, $\mathsf{w}_T$, $\mathsf{w}_L$; and then $c_b$, $C^P$,
$y_0^{\mathsf{T}}$, $C_{\mathbb{C}}$, $C_\rho^{net}$, $C_\times'$, $C_\times''$, $C_{TL}$.

Last comes the coupling ceiling $\gamma$, subject to upper bounds only. These are the
ceilings on $\gamma$ that enter the hypotheses of Theorem~\ref{9.1:thm-preliminary-reads};
the ceiling involving $C_{TL}$ fixed in Corollary~\ref{9.1:cor-tolerance-rooms}, which is
among the conditions \eqref{9.1:eq-constant-order-a}; the condition
$\varepsilon^{\natural}\le 1/4$; and the ceilings on $\gamma$ that the rows in $\theta$ of the
conditions attached to the dilation circles $\mathsf{w}_T$, $\mathsf{w}_L$ impose through the
constants $K_R^{\sharp}$ and $K_o^{L}$. No step after the choice of $L$ imposes a condition
on $L$, and no step after the choice of $M$ imposes a condition on $M$.

With $\beta=\beta_I=1/5$ the following hold for every integer $s$. Here $s$ is not the depth
below the cutoff but an integer parameter, equal in (k1) to the gap $j-n$ between a level $j$
and a layer $n$ as in Proposition~\ref{9.1:prop-dilation-failure}; and (k3) is an identity in
$\theta_S$ and $s$:
\begin{equation}\label{9.1:eq-room-checks-a}
\begin{aligned}
&\text{(k1)}\quad (1-\beta)\tfrac{3\beta}{8}>2\beta\bigl(1-\tfrac{\beta}{2}\bigr)2^{-s}
\iff 2^{-s}<\tfrac{2.4}{14.4}=\tfrac16;\\
&\text{(k2)}\quad \Bigl(1+\tfrac{\beta/2}{1-3\beta}\Bigr)\tfrac34
=\tfrac54\cdot\tfrac34=\tfrac{15}{16},\qquad
\tfrac{2/5}{16}=\tfrac1{40}=\tfrac2{80};\\
&\text{(k3)}\quad 1-\theta_S-\tfrac12(1-\theta_S)=\tfrac12(1-\theta_S),
\ \text{and the same identity multiplied by }2^{-s-1};\\
&\text{(k4)}\quad 40C_\Sigma\le 368\le 400,\qquad 16\cdot 400=6400<2^{18}/8;\\
&\text{(k5)}\quad 16\cdot 9.2\cdot 6.75<994<2^{18}/200;\\
&\text{(k6)}\quad 0.755\le 0.775;\\
&\text{(k7)}\quad 0.8-\tfrac34\cdot 0.025-0.755=0.02625\ge\tfrac1{40},\qquad
2C_\Sigma K_u<125,\qquad 1250\le 2^{11}.
\end{aligned}
\end{equation}
\end{remark}

\begin{proof}
\emph{(k1).} Since $\beta>0$, the left inequality is equivalent to
\[
(1-\beta)\tfrac38>2\bigl(1-\tfrac\beta2\bigr)2^{-s}.
\]
Since $1-\beta/2>0$, this is in turn equivalent to
\[
2^{-s}<\frac{3(1-\beta)}{16(1-\beta/2)}.
\]
At $\beta=1/5$ we have $3(1-\beta)=2.4$ and $16(1-\beta/2)=14.4$, and $2.4/14.4=1/6$.

\emph{(k2).} At $\beta=1/5$ we have $1-3\beta=2/5$ and $\beta/2=1/10$. Hence
\[
\frac{\beta/2}{1-3\beta}=\frac14,\qquad 1+\frac14=\frac54,\qquad
\frac54\cdot\frac34=\frac{15}{16}.
\]
Moreover $(2/5)/16=2/80=1/40$.

\emph{(k3).} The identity is linear in $1-\theta_S$. Multiplying it by $2^{-s-1}$ gives
\[
(1-\theta_S)2^{-s-1}-\tfrac12(1-\theta_S)2^{-s-1}=\tfrac12(1-\theta_S)2^{-s-1}
=(1-\theta_S)2^{-s-2}.
\]

\emph{(k4).} The number $C_\Sigma$ is fixed among the numbers chosen first, before $L$, in
Notation~\ref{9.1:not-constant-order}, with $C_\Sigma\le 9.2$; that is, $40C_\Sigma\le 368$.
Further, $368\le 400$, and
\[
16\cdot 400=6400<32768=2^{18}/8.
\]

\emph{(k5).} We have $9.2\cdot 6.75=62.1$ and $16\cdot 62.1=993.6<994$. Also
$2^{18}/200=1310.72>994$.

\emph{(k6).} This is the comparison of the two decimals $0.755$ and $0.775$.

\emph{(k7).} We have $\tfrac34\cdot 0.025=0.01875$, so
\[
0.8-0.01875-0.755=0.02625\ge 0.025=\tfrac1{40}.
\]
The numbers $C_\Sigma$ and $K_u$ are fixed among the numbers chosen first, before $L$, in
Notation~\ref{9.1:not-constant-order}, with $2C_\Sigma K_u<125$.
Finally, $1250\le 2048=2^{11}$.
\end{proof}

\begin{definition}[Constants and radii across levels for site-blind estimates]
\label{9.1:not-site-blind-constants}
Throughout the site-blind estimates of this chapter the following letters have fixed meanings.
A term $F$ of the expansion is produced at a level $j$, either by Ba{\l}aban's operation
$\mathbf{T}$ (a T-step; we then call $F$ \emph{T-born}) or by Ba{\l}aban's operation
$\mathbf{R}$ (we then call $F$ \emph{$\mathbf{R}$-born}); it is read on the layers $n\le j$,
and $n_F$ denotes the layer at which $F$ is read in the estimate at hand.
\begin{enumerate}
\item[(a)] \emph{Room fraction and charts.} The letter $\varsigma$ denotes the room fraction.
The own tube of $F$ at room fraction $\varsigma$ is $\mathfrak{D}_F[\varsigma]$, and the
admissible data class on the chart is $\mathfrak{T}[\varsigma]$; the two-cutoff seminorm
$\mathsf{S}^{+}$ is taken on the tube at room fraction $7/8$. The local Lipschitz property
$\natural\mathrm{Lip}^{\mathrm{loc}}$ of the Landau chart is used with kernels of rate
$8\delta_1$, where $\delta_1$ is the decay rate of the random-walk estimates of item~(e).
The slice map $\hat F$ is taken on the chart. The net mismatch
$\mathfrak{d}^{net}$ is always taken with the real frame $\mathsf{h}$ removed.
\item[(b)] \emph{Ceiling on the radii.} The local radii $\bar u$ are taken so small that
their global supremum $\bar{\bar u}$ obeys the ceiling
\begin{equation}\label{9.1:eq-site-blind-constants-1}
C_W\,(C_e+C_\Sigma+3)\,\bar{\bar u}\le 2 ,
\end{equation}
which is an upper bound on the radii only; here $C_W$, $C_e$ and $C_\Sigma$ are the constants
entering this ceiling in the site-blind estimates (this $C_W$ is not the constant $C_W(s)$ of
the two-point witness theorem); the numerical factor $4$ in the second clause
of the small-data estimate is used under the ceiling \eqref{9.1:eq-site-blind-constants-1}.
\item[(c)] \emph{Numbers.} We put $\beta:=\beta_I=1/5$. The schedule number $\theta_S$ is the
one fixed with the room schedule of the site-blind estimates, where it is taken in the
interval $0<\theta_S<1$, and
\begin{equation}\label{9.1:eq-site-blind-constants-2}
\theta_\rho:=\tfrac14\min\{\theta_S,\,1-\theta_S\}\le\tfrac18 .
\end{equation}
For levels $n\le j$ we put
\begin{equation}\label{9.1:eq-site-blind-constants-3}
f_{j,n}:=1-\beta\bigl(1-2^{-(j-n)}\bigr).
\end{equation}
\item[(d)] \emph{Top layers.} The table coefficient of a term $F$ born at level $j$ is one
number $\mathsf{c}_F^{+}$, attached to $F$ and used on all its layers $n\le j$; it is fixed
at the top layer $n=j$ by
\begin{equation*}
\mathsf{c}_F^{+}:=1+\beta\quad\text{if $F$ is T-born},\qquad
\mathsf{c}_F^{+}:=1\quad\text{if $F$ is $\mathbf{R}$-born},
\end{equation*}
the value $1+\beta$ for T-born terms being the one of \cite[(3.43)]{B-CMP119} and the
top-layer coefficient of the room schedule. With this
choice, for every T-born $F$ born at level $j$ and on every layer $n\le j$,
\begin{equation}\label{9.1:eq-site-blind-constants-4}
\mathsf{c}_F^{+}\ge f_{j,n}+\theta_S\,\beta\,2^{-(j-n)} ,
\end{equation}
and the room of every term is
\begin{equation}\label{9.1:eq-site-blind-constants-5}
\varrho_F=\theta_\rho\,\beta\,2^{-(j-n_F)} .
\end{equation}
The letter $\iota$ denotes the level of the label at which the term is read.
\item[(e)] \emph{Distance and random-walk conditions.} The distance $d(\cdot,\cdot)$ and the
decay rate $\delta_1$ are those of the random-walk estimates; $d(\cdot,\cdot)$ is a path
metric. Conditions (d1), (d2) and $(\mathrm{T}^{w})$ below are the conditions so tagged in
the random-walk estimates. Condition (d1) is used with the gap $d_{\mathrm{gap}}$, which
satisfies the floor on the gap
\begin{equation}\label{9.1:eq-site-blind-constants-6}
e^{-\delta_1 d_{\mathrm{gap}}}\le L^{-16};
\end{equation}
condition (d2) reads, with the constant $c_1$ of those estimates,
\begin{equation}\label{9.1:eq-site-blind-constants-7}
\sum_y e^{-\delta_1 d(x,y)/2}\le c_1 \quad\text{for every } x;
\end{equation}
and condition $(\mathrm{T}^{w})$ holds with constant $C_T^{w}=c_1L^3$.
Conditions (d1) and $(\mathrm{T}^{w})$ persist on refinements, also for the power
$L^{\mathrm{am}}$ of $L$ considered in the refinement statement of the random-walk estimates,
by that statement. For a nonnegative profile $f$, the
smearings $f^{\sim}$ and $f^{\approx}$ of the site-blind estimates are formed with exponential
weights of decay rates $\delta_1$ and $2\delta_1$ respectively, and $f^{\approx}\le f^{\sim}$
pointwise.
\item[(f)] \emph{Radii across levels.} For $a\in\{0,1\}$ and levels $n'>n$,
\begin{equation}\label{9.1:eq-site-blind-constants-8}
\alpha_{0,n}\le\alpha_{1,n},\qquad
\alpha_{a,n'}\le\tfrac94\,(n'-n)^{1/2}\,\alpha_{a,n},\qquad
\alpha_{a,n}\le 2\,\alpha_{a,n'} .
\end{equation}
The radii enter the site-blind estimates only through $\alpha_{0,n}$, $\alpha_{1,n}$ and the
ratio $\mathsf{r}_1:=C_1/C_0$ of their constants. The first inequality compares the two radii
at a common level; since they have the same dependence on $g_n$ and differ only by their
constant factors $C_0$ and $C_1$, it is the inequality $\mathsf{r}_1\ge1$, taken as fixed with
the choice of these constants. The second inequality is a consequence of
\cite[(2.6)--(2.8)]{B-CMP119}; the third is a consequence of the lower bound $\frac12$ in the
reference-window condition (the reference window being fixed in
Definition~\ref{1.2:def-flow-constants}) under which the correlation theorem is applied.
\end{enumerate}
\end{definition}

\begin{proof}
The bound in \eqref{9.1:eq-site-blind-constants-2}: for every real $\theta_S$ one has
$\min\{\theta_S,1-\theta_S\}\le\frac12\bigl(\theta_S+(1-\theta_S)\bigr)=\frac12$, since the
minimum of two numbers does not exceed their mean; multiplying by $\frac14$ gives
$\theta_\rho\le\frac18$.

The bound \eqref{9.1:eq-site-blind-constants-4}: let $F$ be T-born at level $j$ and let
$n\le j$. By \eqref{9.1:eq-site-blind-constants-3},
\begin{equation*}
f_{j,n}+\theta_S\,\beta\,2^{-(j-n)}=1-\beta+(1+\theta_S)\,\beta\,2^{-(j-n)} .
\end{equation*}
Since $j-n\ge0$ we have $2^{-(j-n)}\le1$, and since $\beta>0$ and $1+\theta_S>0$,
\begin{equation*}
1-\beta+(1+\theta_S)\,\beta\,2^{-(j-n)}\le 1-\beta+(1+\theta_S)\,\beta=1+\theta_S\,\beta
\le 1+\beta ,
\end{equation*}
the last step because $\theta_S<1$. By item~(d), the coefficient $\mathsf{c}_F^{+}$ of a
T-born term is the single number $1+\beta$, used on every layer $n\le j$; hence
\eqref{9.1:eq-site-blind-constants-4} holds on every layer $n\le j$.

The comparison of the smearings: for every $t\ge0$ one has $e^{-2\delta_1 t}\le e^{-\delta_1 t}$,
since $\delta_1>0$; so the weight of decay rate $2\delta_1$ does not exceed the weight of decay
rate $\delta_1$. Each smearing is formed from the nonnegative profile $f$ with its weight
entering as a nonnegative multiplicative factor, so a smaller weight gives a smaller smearing.
Hence $f^{\approx}\le f^{\sim}$ pointwise.
\end{proof}

\begin{lemma}[The site-blind slice lemma]\label{9.1:lem-site-blind-slice}
Let $(F,X,j)$ be a leaf, that is, a term $F$ together with its region $X$ and the level index
$j$ carried by the leaf, which in general differs from the level $n_F$ of $F$ used below. Put
$\mathfrak{S}=\boldsymbol{\Omega}(X^{+})$, the truncated determining set of the leaf, and let
$\mathbf{U}=e^{i\eta\mathcal{Y}}U_e$ on $X^{+}$ be the complex configuration of the slice-map lemma:
$U_e$ is the configuration from which the move starts, $J(U_e)$ its current, $\mathcal{Y}$ the
composite move and $\eta$ its parameter. Let $\mathsf{m}\ge 0$, measured in the units attached to the
labels of the expansion, satisfy
\begin{equation}\label{9.1:eq-site-blind-slice-1}
\ell_\iota\sup_{\tilde{\Delta}(y)}|\mathcal{Y}|\le \mathsf{m}(y),\qquad
\mathsf{m}^{\approx}\le C_{\mathsf{m}}\,\mathsf{m},
\end{equation}
where $\iota$ is the level index of the site, $\ell_\iota$ the length factor attached to it,
$\tilde{\Delta}(y)$ the neighbourhood of the point $y$ over which the supremum is taken, and
$\mathsf{m}^{\approx}$ the smeared profile of $\mathsf{m}$;
and let $\mathsf{T}_e\le\mathsf{T}_r$ be tables, with $\mathsf{G}:=\mathsf{T}_r-\mathsf{T}_e$, such that
the following hold.
\begin{itemize}
\item[(E)] $E_a(U_e,J(U_e))\le\mathsf{T}_e^a$ on $X$ for the free members $E_a$ of the table; for the
member $E_5$ and for the cube clause of $U_c$ the condition is that they hold inside the own tube
$\mathfrak{D}_F$ of $F$.
\item[(A)] At points of class $(=)$, $2^7B_\natural C_{\mathsf{m}}\mathsf{m}\le\ell_\iota^{e_a}\mathsf{G}^a$,
where $e_a$ is the scaling exponent of the member $E_a$.
\item[(S)] At points of class $(=)$, $\mathsf{T}_r\le(\mathsf{c}_F^{+}-\varrho_F)\mathsf{u}_F$.
\item[(U)] At points of class $(\mathrm{up})$, $\mathsf{T}_e\le\frac34\mathsf{u}_F$ and
$2^7B_\natural C_T^{w}C_{\mathsf{m}}\mathsf{m}\le\alpha_{0,\iota}/200$.
\end{itemize}
Then $\tilde{F}(X;P)=F(X;S(\mathbf{U}))$, where $P$ is the chart point at which the re-read term
$\tilde{F}$ is evaluated; at $\mathcal{Y}=0$ the slice is exactly
$S(\mathbf{U})=(U_e,J(U_e)){\upharpoonright}X$; moreover
\begin{equation}\label{9.1:eq-site-blind-slice-2}
E_a(S(\mathbf{U}))\le
\begin{cases}
\mathsf{T}_e^a+\tfrac18\mathsf{G}^a & \text{at points of class } (=),\\[2pt]
\bigl(\tfrac34+\tfrac1{200}\bigr)\mathsf{u}_F^a & \text{at points of class } (\mathrm{up}),
\end{cases}
\end{equation}
and thus $\mathrm{rd}\,\mathbf{U}\in\mathfrak{T}[0](X)$.
\end{lemma}

\begin{proof}
We apply parts (a) and (b) of the slice-map lemma to $\mathbf{U}$. They bound the move of each
member $E_a$, that is, the change of $E_a$ in
passing from $(U_e,J(U_e))$ to $S(\mathbf{U})$, by
\begin{equation}\label{9.1:eq-site-blind-slice-3}
16B_\natural\,\ell_{\mathfrak{S}}^{-e_a}\,\bigl(|\mathcal{Y}|_{\cdot}\bigr)^{\approx},
\end{equation}
with $\ell_{\mathfrak{S}}$ the length factor attached to $\mathfrak{S}$;
for the member $E_1$ this rests on the gradient bound $2\eta^2|\nabla E|(1+C\varepsilon_3)$ on its move,
with $C$ a constant and $\varepsilon_3$ a small constant, established in the proof of the stage-table
theorem of this chapter and used here as stated there. Here
the size of the move at $y$ is
\begin{equation*}
|\mathcal{Y}|_y=\frac{\ell_{\mathfrak{S}}}{\ell_\iota}(y)\cdot\ell_\iota\sup_{\tilde{\Delta}(y)}|\mathcal{Y}|,
\end{equation*}
and the second factor is at most $\mathsf{m}(y)$ by the first inequality
in~\eqref{9.1:eq-site-blind-slice-1}. Hence
$|\mathcal{Y}|_y\le(\ell_{\mathfrak{S}}/\ell_\iota)(y)\,\mathsf{m}(y)$. The smearing
$(\cdot)^{\approx}$ has non-negative weights, so it preserves this inequality.

\emph{Points of class $(=)$.} There $\ell_{\mathfrak{S}}=\ell_\iota=\ell_{n_F}$, with $n_F$ the level
of $F$, so the ratio
$\ell_{\mathfrak{S}}/\ell_\iota$ equals $1$. By the second inequality
in~\eqref{9.1:eq-site-blind-slice-1}, the bound~\eqref{9.1:eq-site-blind-slice-3} is at most
$16B_\natural C_{\mathsf{m}}\ell_\iota^{-e_a}\mathsf{m}$. Since $16=2^7/8$, hypothesis (A) gives
\begin{equation*}
16B_\natural C_{\mathsf{m}}\ell_\iota^{-e_a}\mathsf{m}
=\tfrac18\,\ell_\iota^{-e_a}\,2^7B_\natural C_{\mathsf{m}}\mathsf{m}
\le\tfrac18\mathsf{G}^a .
\end{equation*}
Adding this to the bound $E_a(U_e,J(U_e))\le\mathsf{T}_e^a$ of hypothesis (E) gives the first line
of~\eqref{9.1:eq-site-blind-slice-2}.

\emph{Points of class $(\mathrm{up})$.} Here $(\ell_{\mathfrak{S}}/\ell_\iota)\mathsf{m}$ is a product of
powers $L^{b\lambda'}$ of $L$, in which $b$ and $\lambda'$ are the exponents carried by the length
factors, and of $\mathsf{m}$. Its smearing is
controlled by the product bound with
constant $C_T^{w}$ fixed in Notation~\ref{9.1:not-site-blind-constants}. Measured in the units of
$F$, the move at the point $x$ of class $(\mathrm{up})$ is therefore at most
\begin{equation*}
16B_\natural C_T^{w}C_{\mathsf{m}}\,\frac{\ell_{\mathfrak{S}}}{\ell_\iota}\,\mathsf{m}(x)
\le L^{-u}\,\frac{\alpha_{0,\iota}}{1600},\qquad u:=\iota-(j\wedge\iota).
\end{equation*}
Here we used $16=2^7/8$ and the second condition of hypothesis (U); the ratio
$\ell_{\mathfrak{S}}/\ell_\iota$ at the point, read in the units of $F$, supplies the factor $L^{-u}$.
By the comparison of radii
across levels in Notation~\ref{9.1:not-site-blind-constants},
\begin{equation*}
L^{-u}\alpha_{0,\iota}\le\tfrac94\,\alpha_{0,j\wedge\iota}.
\end{equation*}
Next we pass from level $j\wedge\iota$ down to level $n_F$, and put $v:=(j\wedge\iota)-n_F$. By
part~(ii) of the statement comparing units across levels, the smaller unit gives a factor
$L^{-v}$. This is set against the factor $\frac94v^{1/2}$. Since $v^{1/2}\le L^{v}$, the product of
the two factors is at most $\frac94$. The move is therefore at most
$\bigl(\tfrac94\bigr)^2\tfrac1{1600}\,\mathsf{u}_F^a$, and
\begin{equation*}
\Bigl(\frac94\Bigr)^2\frac1{1600}=\frac{81/16}{1600}=\frac{81}{25600}<\frac{128}{25600}=\frac1{200}.
\end{equation*}
Adding this to $E_a(U_e,J(U_e))\le\mathsf{T}_e^a\le\frac34\mathsf{u}_F^a$, which follows from
(E) and the first condition of (U), gives the second line of~\eqref{9.1:eq-site-blind-slice-2}.

\emph{Composite members.} These are covered by part (c) of the slice-map lemma.

\emph{Conclusion.} At points of class $(=)$, we have $\mathsf{T}_e^a+\frac18\mathsf{G}^a\le\mathsf{T}_r^a$
because $\mathsf{G}\ge0$. By (S), $\mathsf{T}_r\le(\mathsf{c}_F^{+}-\varrho_F)\mathsf{u}_F$. At points
of class $(\mathrm{up})$, the coefficient is $\frac34+\frac1{200}=0.755$, and
$0.755\le\mathsf{c}_F^{+}-\varrho_F$ by the arithmetic in the proof of the stage-table theorem of this
chapter. With~\eqref{9.1:eq-site-blind-slice-2} this gives
$\mathrm{rd}\,\mathbf{U}\in\mathfrak{T}[0](X)$.
\end{proof}

Lemma~\ref{9.1:lem-site-blind-slice} reads the move $\mathcal{Y}$ only through
$\sup_{\tilde{\Delta}(y)}|\mathcal{Y}|$. The stage-table theorem of this chapter is its special case
$\mathsf{T}_e=\mathsf{T}_{n_e}$, $\mathsf{T}_r=\mathsf{T}_\nu$ and $\mathsf{m}=2\ell\sum_i a^0_i/K_o$, where
$\mathsf{T}_{n_e}$ and $\mathsf{T}_\nu$ are the stage tables at the levels $n_e$ and $\nu$, $a^0_i$ are
the shares, and $\ell$ and $K_o$ are the length factor and the constant with which that theorem is
stated.

\begin{lemma}[Decay of a profile sourced at the top levels]\label{9.1:lem-top-source-decay}
Let $\mathcal{S}$ be a set of cells, each of label level at least $k_{\mathcal{S}}$, let
$\mathsf{m}^0$ be a constant, and call
\begin{equation}\label{9.1:eq-top-source-decay-1}
\mathsf{m}(x):=\mathsf{m}^0\,e^{-\delta_1 d(x,\mathcal{S})/4}
\end{equation}
the \emph{profile} sourced at $\mathcal{S}$, where $d(\cdot,\cdot)$ denotes the distance between
cells, with $\delta_1$, $c_1$ and the radii $\alpha_{0,m}$ as in
Notation~\ref{9.1:not-site-blind-constants}, and with $(\cdot)^{\approx}$ and $(\cdot)^{\sim}$ the
two smeared majorants of functions on cells fixed for the site-blind estimates of this chapter.
Then:
\begin{itemize}
\item[(i)] $\mathsf{m}$ is slow, that is, it satisfies the slow-variation condition imposed on
functions on cells in the site-blind estimates of this chapter, and
\begin{equation*}
\mathsf{m}^{\approx}\le\mathsf{m}^{\sim}\le c_1\,\mathsf{m};
\end{equation*}
moreover, for every function $b$ on cells,
\begin{equation}\label{9.1:eq-top-source-decay-2}
(|b|)^{\sim}\le c_1\,\Bigl(\sup|b|\Bigr)\,e^{-\delta_1 d(\cdot,\operatorname{supp} b)/4}.
\end{equation}
\item[(ii)] At every cell $x$ of label level $m\le k_{\mathcal{S}}$, with $s:=k_{\mathcal{S}}-m$,
\begin{equation}\label{9.1:eq-top-source-decay-3}
\frac{\mathsf{m}(x)}{\alpha_{0,m}}\le\frac{9}{2}\,2^{-s}\,\frac{\mathsf{m}^0}{\alpha_{0,k_{\mathcal{S}}}}.
\end{equation}
\end{itemize}
\end{lemma}

\begin{proof}
Write $(t)_+:=\max\{t,0\}$ and $1\vee t:=\max\{1,t\}$.

(i) For cells $x,y$ the triangle inequality for $d$ gives
$d(y,\mathcal{S})\ge d(x,\mathcal{S})-d(x,y)$. By \eqref{9.1:eq-top-source-decay-1},
\begin{equation*}
\mathsf{m}(y)\le\mathsf{m}(x)\,e^{\delta_1 d(x,y)/4},
\end{equation*}
and from this bound the slowness of $\mathsf{m}$ follows. Combining the factor
$e^{\delta_1 d(x,y)/4}$ of the last display with the decay weight $e^{-\delta_1 d(x,y)}$ of the
smearing gives
\begin{equation*}
e^{-\delta_1 d(x,y)}\,e^{\delta_1 d(x,y)/4}=e^{-3\delta_1 d(x,y)/4}\le e^{-\delta_1 d(x,y)/2},
\end{equation*}
and condition (d2) of Notation~\ref{9.1:not-site-blind-constants} then yields
$\mathsf{m}^{\approx}\le\mathsf{m}^{\sim}\le c_1\mathsf{m}$.
For \eqref{9.1:eq-top-source-decay-2}, consider the profile sourced at $\operatorname{supp}b$
with $\mathsf{m}^0$ replaced by $\sup|b|$. This profile dominates $|b|$ pointwise: off
$\operatorname{supp}b$ the function $|b|$ vanishes, and on $\operatorname{supp}b$ the exponential
equals one. Since the smearing $(\cdot)^{\sim}$ is monotone, $(|b|)^{\sim}$ is at most the
$(\cdot)^{\sim}$-majorant of this profile, and the bound just proved, applied to this profile,
gives \eqref{9.1:eq-top-source-decay-2}.

(ii) Let $x$ have label level $m\le k_{\mathcal{S}}$ and put $s=k_{\mathcal{S}}-m$. By condition
(d1) of Notation~\ref{9.1:not-site-blind-constants} and a lower bound on $L$ from the order of
choice of the constants,
\begin{equation*}
\mathsf{m}(x)\le\mathsf{m}^0\,L^{-4(s-1)_+}.
\end{equation*}
Combined with the comparison $(\alpha)$ of the radii $\alpha_{0,m}$ and
$\alpha_{0,k_{\mathcal{S}}}$ in Notation~\ref{9.1:not-site-blind-constants}, this shows that for
\eqref{9.1:eq-top-source-decay-3} it suffices to show
\begin{equation}\label{9.1:eq-top-source-decay-4}
2^{s}\,(1\vee s)^{1/2}\,L^{-4(s-1)_+}\le 2 \qquad (s\ge 0).
\end{equation}
We verify \eqref{9.1:eq-top-source-decay-4} case by case.
\begin{itemize}
\item For $s=0$ the left member equals $1$.
\item For $s=1$ the left member equals $2$.
\item For $s=2$ the left member equals $4\sqrt{2}\,L^{-4}$, which is less than $1$ since
$L$ is an odd integer greater than one.
\end{itemize}
For $s\ge 2$, the ratio of the left member at $s+1$ to its value at $s$ is
\begin{equation*}
2\Bigl(\frac{s+1}{s}\Bigr)^{1/2}L^{-4}\le 2\sqrt{3/2}\,L^{-4}<1,
\end{equation*}
the last inequality for the same reason.
So for $s\ge 3$ the left member is smaller than its value at $s=2$, and
\eqref{9.1:eq-top-source-decay-4} holds for all $s\ge 0$.

Part~(ii) is an estimate stated in \cite{B-CMP119} for a profile sourced at the top levels,
written here as an inequality.
\end{proof}

\begin{definition}[The classical condition in profile form]\label{9.1:def-profile-condition}
The condition defined here concerns one machine; it is a classical condition, and every radius
entering it is independent of the coupling $g$. The \emph{classical condition in profile form} is
the classical condition for one machine together with the ceiling condition, the size clause of
the former being read with the increment's own profile, namely:
\begin{itemize}
\item[(pr)] every increment $h$ of every site, with source datum $b$ and dilation
$\tau_t$ (the subscript $t$ being the parameter indexing the dilation at the site), whose
dilation parameter is $\mathsf{s}$, obeys on the contour of the site
the bound
\begin{equation}\label{9.1:eq-profile-condition-1}
\ell_\iota |h|(y) \le B^{\mathsf{s}}\, \bar\tau(y)\, (|b|)^{\sim}(y),
\end{equation}
where the dilation factor $B^{\mathsf{s}} = B^{\mathsf{s}}(d, L, M_1, \kappa_1)$ depends on
$d$, $L$, $M_1$ and $\kappa_1$ only.
\end{itemize}
Here $\ell_\iota|h|(y)$ is the size of $h$ at $y$ read in label units at label level $\iota$,
$\bar\tau(y)$ is the weight at $y$ on the contour induced by the dilation $\tau_t$, and
$(\cdot)^{\sim}$ is the majorant operation of Lemma~\ref{9.1:lem-top-source-decay}\,(i).
The constants and radii across levels are those of
Notation~\ref{9.1:not-site-blind-constants}.

Condition (pr) is a standing hypothesis. At natural vertices it is supplied by the regularity
statement together with clause (d) of the tilted natural-interpolation statement, whose
constants are free of $Z$. The bound $|\mathsf{s}| \le e^{\kappa_1}$ on the dilation parameter is
supplied at the site labelled (T) by clause (d2) of the homogeneity statement, and at the sites
labelled (P) and (L) by clauses (e) and (f) of the interpolation statement and by clause (i) of
its twisted form together with the separately stated condition (c) of the interpolation
statement.

At the site (P), when the two pair radii coincide, $r^{\times}_{\square} = r^{\bullet}_{\square}$,
condition (pr) is supplied by the following rate bound, whose left member is the size of $h_\mu$
at the box level $\lambda_Z$:
\begin{equation}\label{9.1:eq-profile-condition-2}
\ell_{\lambda_Z} |h_\mu|(y) \le \mathsf{H}^0_\nu\, e^{-3\delta_1 d(y, R_\nu)/8},
\qquad
\mathsf{H}^0_\nu := \frac{C_\flat\, c_M\, \mathsf{b}_j}{\vartheta^0_P}
\le C_\flat'\, c_M\, \alpha_{0, j_\nu},
\end{equation}
which is linear in $c_M$. In \eqref{9.1:eq-profile-condition-2} the size is read in box units,
at the box level $\lambda_Z$ (an integer level, not the constant $\lambda$ of the setting of
clause (v)), and the rate $3\delta_1/8$ refers to box units; the distance $d(y,R_\nu)$ is the same
in \eqref{9.1:eq-profile-condition-2} and in the label-unit bound
\eqref{9.1:eq-profile-condition-4} below; only the size is converted. Read in label units at
label level $\iota$, the bound becomes \eqref{9.1:eq-profile-condition-4}, which is the form
(pr) takes at the site (P); its decay rate $\delta_1/4$ is that of the decay profile
$\mathsf{h}_\nu$ attached to the index $\nu$ (an object distinct from the real frame $\mathsf{h}$
and from the normalised sourced integral $\mathsf{h}^{\mathsf{c}}$).
On the enlarged parameter polydisc $4\mathfrak{P}_{\mathsf{c}}$ the profiles $\mathsf{h}_\nu$ are
replaced by $4\mathsf{h}_\nu$, by the statement of this chapter that bounds the profiles on
$4\mathfrak{P}_{\mathsf{c}}$.
Here $h_\mu$ is the move field with index $\mu$, an increment in the sense of (pr); $\nu$ is the
index, attached to the move $h_\mu$, that labels the region $R_\nu$, the level $j_\nu$ and the
constant $\mathsf{H}^0_\nu$; $R_\nu$ is
the region from which the distance in the decay is measured (a
region, not the ball radius $R$); $j_\nu$ is the level attached to $\nu$;
$\mathsf{b}_j := C_1\delta'_j$, where $\delta'_j$ is the level-$j$ parameter of the tube-rate
statement of this chapter;
$\vartheta^0_P$ is the quantity $\vartheta$ fixed by the tube-rate statement,
namely $K_\vartheta$ times the quotient of $\delta'_j$ by the tube radius at level $j$, where
$K_\vartheta$ is the constant of that statement; $C_\flat$ is a constant independent of $g$; and
$C_\flat'$ is the constant fixed in the derivation below.
\end{definition}

\begin{proof}[Derivation of the right inequality of \eqref{9.1:eq-profile-condition-2} and of
the label-unit rate]
We take the rate bound \eqref{9.1:eq-profile-condition-2}, left inequality, in the form stated
by the rate statement at the site (P), with $\mathsf{H}^0_\nu$ defined as displayed. We show in
two steps how it yields the right inequality of \eqref{9.1:eq-profile-condition-2} and the rate
$\delta_1/4$ in label units.

\emph{The constant.} By the tube-rate statement, the denominator
$\vartheta^0_P$ is $K_\vartheta$ times the quotient of $\delta'_j$ by the tube radius at level
$j$, while the numerator contains $\mathsf{b}_j = C_1\delta'_j$. The factor $\delta'_j$ therefore
cancels, and $\mathsf{b}_j/\vartheta^0_P$ is $C_1/K_\vartheta$ times the tube radius at level $j$;
this quotient is a multiple of $\alpha_{*,j}$ by a factor independent of $g$. Multiplying by
$C_\flat c_M$ gives the bound $\mathsf{H}^0_\nu \le C_\flat' c_M \alpha_{0,j_\nu}$, where
$C_\flat'$ is chosen once so as to be at least $L^2$ times $C_\flat C_1/K_\vartheta$ times this
$g$-free factor, and so as to contain the factor $C_s L^2$ introduced below. Since $c_M$ enters
$\mathsf{H}^0_\nu$ only as a prefactor, this bound is linear in $c_M$.

\emph{The units.} The rate $3\delta_1/8$ in \eqref{9.1:eq-profile-condition-2} is the rate of the
supplying bound, with sizes and distances in box units. The size at label level $\iota$ exceeds
the size at the box level $\lambda_Z$ by at most the factor $L^{\iota-\lambda_Z}$; this factor is
the cost of reading the bound in label units. By condition (d3) of the domain convention and the
floor $(\mathrm{F}^{\flat})$, one of the floors on exponents fixed after $L$ in the order of
choice, this factor obeys
\begin{equation*}
L^{\iota-\lambda_Z} \le L^2\, e^{\delta_1 d(y, R_\nu)/16}.
\end{equation*}
Hence $\ell_\iota|h_\mu|(y)$ is bounded by $L^{\iota-\lambda_Z}$ times
$\ell_{\lambda_Z}|h_\mu|(y)$, and so by $L^{\iota-\lambda_Z}$ times the right side of the left
inequality of \eqref{9.1:eq-profile-condition-2}; using $d(y,R_\nu) \ge 0$, we obtain
\begin{align}
\ell_\iota |h_\mu|(y)
&\le L^{\iota-\lambda_Z}\,\ell_{\lambda_Z}|h_\mu|(y)
\le L^{\iota-\lambda_Z}\,\mathsf{H}^0_\nu\, e^{-3\delta_1 d(y, R_\nu)/8} \notag\\
&\le L^2\, \mathsf{H}^0_\nu\, e^{-(3/8 - 1/16)\delta_1 d(y, R_\nu)}
= L^2\, \mathsf{H}^0_\nu\, e^{-5\delta_1 d(y, R_\nu)/16} \notag\\
&\le L^2\, \mathsf{H}^0_\nu\, e^{-\delta_1 d(y, R_\nu)/4},
\label{9.1:eq-profile-condition-3}
\end{align}
because $5/16 \ge 1/4$. The rate $\delta_1/4$ is the one used by the profile $\mathsf{h}_\nu$.
The constant $C_\flat'$ contains the product $C_s L^2$, where $C_s$ is the constant that
accompanies the factor $L^2$ in this passage from box units to label units. By the choice of
$C_\flat'$ in the first step, the bound obtained there holds
with the additional factor $L^2$, that is, $L^2\,\mathsf{H}^0_\nu \le C_\flat'\,c_M\,
\alpha_{0,j_\nu}$; inserting this in \eqref{9.1:eq-profile-condition-3} gives, in label units,
\begin{equation}\label{9.1:eq-profile-condition-4}
\ell_\iota |h_\mu|(y) \le C_\flat'\, c_M\, \alpha_{0,j_\nu}\, e^{-\delta_1 d(y, R_\nu)/4},
\end{equation}
and the bound stays linear in $c_M$.
\end{proof}

\begin{lemma}[Slice displacement between non-critical points]\label{9.1:lem-noncritical-displacement}
Let $U_c$ be the centre of the chart, let $\eta$ be the chart parameter, and let $S(\mathbf{U})$
denote the slice of a configuration $\mathbf{U}$ of the chart. Let
$\mathbf{U}_1=e^{i\eta A_1}U_c$ be any point of the chart, with chart coordinate $A_1$, such that
\begin{equation*}
C_W\,|A_1|_c\le 2 ,
\end{equation*}
where $C_W$ and $|\cdot|_c$ are the constant and the seminorm through which the coordinate $A_1$
enters the estimate of the chart between two data; this $C_W$ is a constant of the chart and is
distinct from the constant $C_W(s)$ of the two-point witness theorem. Let $\mathcal{Y}$ be a
composite move and let $\mathbf{U}_2=e^{i\eta\mathcal{Y}}\mathbf{U}_1$. No criticality is assumed
of $\mathbf{U}_1$ or $\mathbf{U}_2$. Then every member of $S(\mathbf{U}_2)-S(\mathbf{U}_1)$ is
bounded by
\begin{equation}\label{9.1:eq-noncritical-displacement-1}
4\,B_{\natural}\,\bigl(|\mathcal{Y}|_{\cdot}\bigr)^{\approx},
\end{equation}
where $B_{\natural}$ is the constant of the local Lipschitz property
$\natural\mathrm{Lip}^{\mathrm{loc}}$ of the Landau chart, $|\mathcal{Y}|_{\cdot}$ is the profile
of the size of the move $\mathcal{Y}$, and $(\cdot)^{\approx}$ is the smeared majorant of a
profile.
\end{lemma}

\begin{proof}
We follow the proof of part~(b) of the slice-displacement bound at a critical point $U_e$, and
observe that no step of it uses the criticality of $U_e$. That proof bounds the members of the
difference of two slices by two inputs only, each applied to the two data $\mathbf{U}_1$ and
$\mathbf{U}_2$ of the chart. The first input is the estimate of the chart between two data,
which carries the factor $2+C_W|A_1|_c$. The second is the local Lipschitz property
$\natural\mathrm{Lip}^{\mathrm{loc}}$ of the Landau chart between two data, which carries the
constant $B_{\natural}$ and the smeared majorant $(|\mathcal{Y}|_{\cdot})^{\approx}$ of the move
$\mathcal{Y}$. Each input is a statement about an arbitrary pair of points of the chart; neither
refers to a critical point. Both therefore apply to the pair
$(\mathbf{U}_1,\mathbf{U}_2)$ of the statement. This gives the bound
$(2+C_W|A_1|_c)\,B_{\natural}(|\mathcal{Y}|_{\cdot})^{\approx}$ for every member of
$S(\mathbf{U}_2)-S(\mathbf{U}_1)$.

By the hypothesis $C_W|A_1|_c\le2$ we have
\begin{equation*}
2+C_W|A_1|_c\le 4 ;
\end{equation*}
the factor $4$ is thus the factor $2+C_W|A_1|_c$ of the chart estimate under the ceiling of the
proof of part~(b). Inserting this into the preceding bound gives
\eqref{9.1:eq-noncritical-displacement-1}.

The hypothesis holds on the complex parameter polydisc $\mathfrak{P}_{\mathsf{c}}$. There the
$A$-part of every object of the chain read on $\mathfrak{P}_{\mathsf{c}}$ is tube-small, in the
sense that
\begin{equation*}
C_W\,(C_e+C_\Sigma+3)\,\bar{\bar u}\le 2 ,
\end{equation*}
where $\bar{\bar u}$ is the global supremum of the local radii and $C_e$, $C_\Sigma$ are constants
of the chart on $\mathfrak{P}_{\mathsf{c}}$.
\end{proof}

The criticality of $U_e$ enters the slice-displacement bound through its part~(a) alone, where it
makes the hand-over at the entry exact.

\begin{definition}[Entry and room tables at the small-field site]\label{9.1:def-small-field-tables}
Fix a renormalisation step from level $k$ to level $k+1$. Let $\Phi$ be a term born in the
$\mathbf{T}$-operation of level $k+1$ and localised in the region $Y$. The \emph{seed} of
$\Phi$ is
\begin{equation*}
U_e:=U\bigl(\boldsymbol{\Omega}(Y^{+}),\mathsf{v}\bigr),
\end{equation*}
the configuration determined on the box set $\boldsymbol{\Omega}(Y^{+})$ by the data
$\mathsf{v}$. Its slice lies in $\mathfrak{D}_\Phi[7/8]$, the own tube of $\Phi$ at room
fraction $7/8$. The leaves of $\Phi$, in the sense of the hereditary description of the leaves
of $\mathbf{T}$-born terms, are triples $(F,X,j)$ of a term $F$, a region $X$ in which $F$ is
localised, and a level $j$. Every leaf has $j\le k$ and satisfies
\begin{equation*}
X^{+}\subset K(X)\subset Y,
\end{equation*}
where $K(X)=(X^{+})_0$ is the leaf hull. These inclusions are clauses (b) and (d) of the
hereditary description of the leaves of $\mathbf{T}$-born terms.

\emph{The entry table.} Let $\mathsf{u}_\Phi$ be the own units of $\Phi$, with unit
$\mathsf{u}_n$ on the layer of level $n$. Here $f_{k,n}$ and $f_{k+1,n}$ denote the
coefficients, on the layer of level $n$, of the small-field domains of levels $k$ and $k+1$;
$\beta$ (without subscript, and distinct from the coupling increment $\beta_{k+1}$) is the fixed
number entering these coefficients; and $\Omega_{k+1}$ is the small-field region of level
$k+1$. Let $\mathsf{c}_\Phi^{+}$ be the table coefficient taking the following values:
\begin{equation}\label{9.1:eq-small-field-tables-1}
\mathsf{c}_\Phi^{+}=
\begin{cases}
F_n:=f_{k+1,n}+\theta_S\,\beta\,2^{-(k+1-n)} & \text{on the layer of level } n\le k,\\[2pt]
1+\beta & \text{on } \Omega_{k+1}.
\end{cases}
\end{equation}
These are the coefficients of the post-step domain $\mathfrak{U}^{\mathrm{post}}$. The
\emph{entry table} is
\begin{equation*}
\mathsf{T}_e:=\mathsf{c}_\Phi^{+}\,\mathsf{u}_\Phi .
\end{equation*}

\emph{The room table and the room increment.} The \emph{room table} is defined layer by layer
by
\begin{equation}\label{9.1:eq-small-field-tables-2}
\mathsf{T}_r:=\Bigl(f_{k,n}-\tfrac12(1-\theta_S)\,\beta\,2^{-(k+1-n)}\Bigr)\mathsf{u}_n .
\end{equation}
The \emph{room increment} is $\mathsf{G}_T:=\mathsf{T}_r-\mathsf{T}_e$. On the layer of level
$n\le k$, subtracting $F_n\mathsf{u}_n$ from \eqref{9.1:eq-small-field-tables-2} gives
\begin{equation*}
\mathsf{G}_T=\Bigl(f_{k,n}-F_n-\tfrac12(1-\theta_S)\,\beta\,2^{-(k+1-n)}\Bigr)\mathsf{u}_n .
\end{equation*}
The allotment of the post-step domain on that layer, $f_{k,n}-F_n$, equals
$(1-\theta_S)\,\beta\,2^{-(k+1-n)}$; by the definition of $F_n$ in
\eqref{9.1:eq-small-field-tables-1} this is the property
$f_{k,n}-f_{k+1,n}=\beta\,2^{-(k+1-n)}$ of the coefficients of the post-step domain
$\mathfrak{U}^{\mathrm{post}}$, which is taken here by statement. Hence
\begin{equation}\label{9.1:eq-small-field-tables-3}
\mathsf{G}_T=\tfrac12(1-\theta_S)\,\beta\,2^{-(k+1-n)}\,\mathsf{u}_n .
\end{equation}
Thus the room increment is one half of the allotment $f_{k,n}-F_n$.

\emph{The increment.} Let $\mathfrak{g}^{\mathbb{C}}$ be the complexification of the Lie algebra
of $\mathrm{SU}(N)$. Let $\mathsf{B}$ be a $\mathfrak{g}^{\mathbb{C}}$-valued field on the free
bonds of $\Omega_{k+1}$. The datum on the rim of $\Omega_{k+1}$ is counted among these bonds,
in the convention of clause (d3) of the statement that fixes the free bonds.

Let $\{\zeta_\square\}$ be the partition of unity subordinate to the cover by cubes $\square$,
whose cubes are ordered; $\square'<\square$ denotes this order. Let $\mathsf{t}$ be the
interpolation parameter and $\mathsf{s}$ the dilation. The \emph{increment} is, in the form of
\cite[(3.4), (3.15)]{B-CMP109} and \cite[(1.10), (1.22)]{B-CMP116},
\begin{equation}\label{9.1:eq-small-field-tables-4}
\mathcal{Y}=\Bigl[\sum_{\square'<\square}\zeta_{\square'}+\mathsf{t}\,\zeta_\square\Bigr]
\mathbb{H}_k(\mathsf{s},B').
\end{equation}
Here $\mathbb{H}_k(\mathsf{s},B')$ is Ba{\l}aban's field $\mathbb{H}_k$ of those displays, taken
at dilation $\mathsf{s}$ and datum $B'$, for every datum $B'$ with
\begin{equation}\label{9.1:eq-small-field-tables-5}
|B'|\le C_1\,g_k\,|\mathsf{B}| .
\end{equation}
In this bound $C_1$ is the constant of \cite[(1.20)]{B-CMP116} and $g_k$ is the running
coupling of level $k$. This $C_1$ is not the one entering the ratio $\mathsf{r}_1=C_1/C_0$,
which is read with the constant $C_1$ of \cite[(2.28)]{B-CMP119}.
\end{definition}

\begin{lemma}[Accommodation at the small-field site]\label{9.1:lem-small-field-accommodation}
Let the step $k\to k+1$, the output $\Phi$, the seed $U_e$, the triples $(F,X,j)$ with $j\le k$,
the tables $\mathsf{T}_e$, $\mathsf{T}_r$, $\mathsf{G}_T$ and the increment
\begin{equation*}
\mathcal{Y}=\Bigl[\sum_{\square'<\square}\zeta_{\square'}+\mathsf{t}\,\zeta_{\square}\Bigr]
\mathbb{H}_k(\mathsf{s},B'),\qquad |B'|\le C_1 g_k|\mathsf{B}|,
\end{equation*}
be as in Definition~\ref{9.1:def-small-field-tables}, and let
$\mathbf{U}=e^{i\eta\mathcal{Y}}U_e$ on $X^{+}$ be as in
Lemma~\ref{9.1:lem-site-blind-slice}. Put
\begin{equation}\label{9.1:eq-small-field-accommodation-1}
K_R^{o}:=\frac{2^{15}B_{\natural}C_T^{w}c_1^{2}B^{\mathsf{s}}C_1}{(1-\theta_S)\beta}.
\end{equation}
Assume $|\mathsf{s}(\Delta)|\le e^{\kappa_1}$ and
\begin{equation}\label{9.1:eq-small-field-accommodation-2}
(1+|\mathsf{t}|)\,K_R^{o}\,g_k\sup|\mathsf{B}|\le\alpha_{0,k}.
\end{equation}
Then hypotheses \textup{(E)}, \textup{(A)}, \textup{(S)} and \textup{(U)} of
Lemma~\ref{9.1:lem-site-blind-slice} hold with $C_{\mathsf{m}}=c_1$ and
\begin{equation}\label{9.1:eq-small-field-accommodation-3}
\mathsf{m}_T:=\mathsf{m}_T^{0}\,e^{-\delta_1 d(\cdot,\Omega_{k+1})/4},\qquad
\mathsf{m}_T^{0}:=(1+|\mathsf{t}|)\,B^{\mathsf{s}}c_1C_1\,g_k\sup|\mathsf{B}|,
\end{equation}
so that Lemma~\ref{9.1:lem-site-blind-slice} gives: on every layer $n\le k$,
$E_a(S(\mathbf{U}))\le\mathsf{T}_e^{a}+\tfrac18\mathsf{G}_T^{a}$; on $\Omega_{k+1}$ and on the
layers $n>n_F$, $E_a(S(\mathbf{U}))\le(\tfrac34+\tfrac1{200})\mathsf{u}_F$.
\end{lemma}

\begin{proof}
\emph{Hypothesis \textup{(E)}.} It holds by the definition of the seed tube
$\mathfrak{T}_\Phi[7/8]$, in which the seed $U_e$ lies by
Definition~\ref{9.1:def-small-field-tables}, read with $\mathsf{T}_e=\mathsf{c}_\Phi^{+}\mathsf{u}_\Phi$.

\emph{The majorant.} Since the $\zeta_{\square'}$ belong to a partition of unity,
$0\le\sum_{\square'<\square}\zeta_{\square'}\le1$ and $0\le\zeta_\square\le1$, whence
$|\sum_{\square'<\square}\zeta_{\square'}+\mathsf{t}\zeta_\square|\le1+|\mathsf{t}|$. The size clause
(pr) of Definition~\ref{9.1:def-profile-condition}, applied to the increment
$\mathbb{H}_k(\mathsf{s},B')$ with source datum $B'$ under the hypothesis
$|\mathsf{s}(\Delta)|\le e^{\kappa_1}$, and part~(i) of
Lemma~\ref{9.1:lem-top-source-decay} applied to $b=B'$ give on the contour
\begin{equation*}
\begin{split}
\ell_\iota|\mathcal{Y}|(y)&\le(1+|\mathsf{t}|)B^{\mathsf{s}}(|B'|)^{\sim}(y)
\le(1+|\mathsf{t}|)B^{\mathsf{s}}c_1\sup|B'|\,e^{-\delta_1 d(y,\operatorname{supp}B')/4}\\
&\le(1+|\mathsf{t}|)B^{\mathsf{s}}c_1C_1g_k\sup|\mathsf{B}|\,e^{-\delta_1 d(y,\Omega_{k+1})/4}
=\mathsf{m}_T(y),
\end{split}
\end{equation*}
the third step using $|B'|\le C_1g_k|\mathsf{B}|$ and
$\operatorname{supp}B'\subset\Omega_{k+1}$. Thus $\mathsf{m}_T$ is the profile of
Lemma~\ref{9.1:lem-top-source-decay} with $\mathcal{S}=\Omega_{k+1}$ and
$\mathsf{m}^0=\mathsf{m}_T^0$, and part~(i) of that lemma gives
$\mathsf{m}_T^{\approx}\le c_1\mathsf{m}_T$, i.e.\ $C_{\mathsf{m}}=c_1$.

\emph{Hypothesis \textup{(A)}.} By \eqref{9.1:eq-small-field-accommodation-1}--\eqref{9.1:eq-small-field-accommodation-3},
\begin{equation}\label{9.1:eq-small-field-accommodation-4}
\frac{\mathsf{m}_T^0}{\alpha_{0,k}}\le\frac{B^{\mathsf{s}}c_1C_1}{K_R^o}
=\frac{(1-\theta_S)\beta}{2^{15}B_{\natural}C_T^{w}c_1}.
\end{equation}
The sources lie at label level at least $k$. At an $(=)$-point $x$ of label level $n\le k$
put $s:=k-n$; part~(ii) of Lemma~\ref{9.1:lem-top-source-decay} with $k_{\mathcal{S}}=k$,
then \eqref{9.1:eq-small-field-accommodation-4} and $2^7\cdot\tfrac92=576\le2^{13}$, give
\begin{equation*}
\begin{split}
2^7B_{\natural}c_1\mathsf{m}_T(x)
&\le2^7\cdot\tfrac92\,B_{\natural}c_1\,2^{-s}\,\frac{\mathsf{m}_T^0}{\alpha_{0,k}}\,\alpha_{0,n}
\le\frac{2^{13}}{2^{15}}\,\frac{(1-\theta_S)\beta}{C_T^{w}}\,2^{-s}\alpha_{0,n}\\
&\le\tfrac14(1-\theta_S)\beta\,2^{-s}\alpha_{0,n},
\end{split}
\end{equation*}
the last step with $C_T^{w}\ge1$. By Definition~\ref{9.1:def-small-field-tables},
$\mathsf{G}_T=\tfrac12(1-\theta_S)\beta\,2^{-(s+1)}\mathsf{u}_n=\tfrac14(1-\theta_S)\beta\,2^{-s}\mathsf{u}_n$,
and the right-hand side is the least of the members $\ell_\iota^{e_a}\mathsf{G}_T^{a}$. Since
$C_{\mathsf{m}}=c_1$, this is \textup{(A)}.

\emph{Hypothesis \textup{(S)}.} It follows from the three cases of the room schedule, the third
case being used also at $s=0$, i.e.\ at $n=j=k$, where it reads
$1-\tfrac14(1-\theta_S)\beta\le1-\theta_\rho\beta$.

\emph{Hypothesis \textup{(U)}.} At $(\mathrm{up})$-points, the bound of the entry table by the
radii and the bound of the radii by the units $\mathsf{u}_F$ give
\begin{equation*}
\mathsf{T}_e\le\tfrac65\,\alpha_{a,\iota}\ell_\iota^{-e_a}
\le\tfrac65\cdot\tfrac94\,L^{-1}\mathsf{u}_F\le\tfrac34\,\mathsf{u}_F,
\end{equation*}
the last step being $27/(10L)\le3/4$. For the size clause, the chain of \textup{(A)} with
$C_T^{w}$ kept, \eqref{9.1:eq-small-field-accommodation-4}, $(1-\theta_S)\beta\le1/5$ and
$2^{-s}\le1$ give
\begin{equation*}
2^7B_{\natural}C_T^{w}c_1\mathsf{m}_T(x)
\le\frac{\tfrac92(1-\theta_S)\beta}{256}\,2^{-s}\alpha_{0,n}
\le\frac{9/2}{1280}\,\alpha_{0,n}<\frac{\alpha_{0,n}}{200}.
\end{equation*}
Hence \textup{(U)} holds, and Lemma~\ref{9.1:lem-site-blind-slice} applies.
\end{proof}

\begin{remark}[Accommodation under homogeneity in the source]
\label{9.1:rem-homogeneous-accommodation}
In this remark $\mathsf{B}$ is the field variable at the small-field site, $\mathsf{t}$ is the
bracket parameter and $\mathsf{s}(\Delta)$ is the dilation parameter at the cell $\Delta$ of the
site in the notation of Lemma~\ref{9.1:lem-small-field-accommodation}, the slice move is the
passage from $\mathbf{U}$ to its slice $S(\mathbf{U})$ used in that lemma, and
$K_R^o$ is the constant defined in that lemma. The letter $K_R$ is the constant of the
small-field allotment $K_Rg_k\sup|\mathsf{B}|$ in the correlation theorem. We put
\begin{equation}\label{9.1:eq-homogeneous-accommodation-1}
\mathfrak{w}_k:=\|C\|\,g_k\sup|\mathsf{B}|,\qquad K_R^{\sharp}:=\max\{2K_R,\,K_R^o\},
\end{equation}
where $\|C\|$ is the norm fixed in the homogeneity clause of the correlation theorem. Assume the
homogeneity clause of the correlation theorem, in its dilation part. Then, with
Lemma~\ref{9.1:lem-small-field-accommodation} (accommodation at the small-field site), the
following hold.
\begin{itemize}
\item[(i)] Suppose that $|\mathsf{s}(\Delta)|\le e^{\kappa_1}$ and
\[
(1+|\mathsf{t}|)K_R^{\sharp}g_k\sup|\mathsf{B}|\le\alpha_{0,k}.
\]
Then the conclusions of Lemma~\ref{9.1:lem-small-field-accommodation} hold, and the inequalities
of the response estimate at the small-field site, which is stated on the polydisc of
\cite[(3.43)]{B-CMP119} only, hold there
with right members linear in $\mathfrak{w}_k$; at the dilated points with
$|\mathsf{s}(\Delta)|\le e^{\kappa_1}$ and after the slice move they hold in the same form, with
the constant $C_H$ of that response estimate replaced by $C_H e^{16\kappa_1}$. The undilated
point uses at most half of the allotment $K_R^{\sharp}g_k\sup|\mathsf{B}|$; in the statement of
this chapter whose hypothesis this is, that hypothesis is therefore met under~(i) and is not a
separate ceiling condition.
\item[(ii)] The replacement of $K_R$ by $K_R^{\sharp}$ is made only in the sourced run of the
correlation theorem on the class $\mathfrak{S}$ of tuned runs that is used for the uniqueness
theorem. It enters the definitions of the
quantities of the correlation theorem
formed from $K_R$: the dilation circle $\mathsf{w}_T$, the numbers $\theta_k$ and
$\theta^{\mathbb{C}}_k$, and the constants $c_V$ and $c_A$; these definitions are made before
$\gamma$. The statement of the correlation theorem is unchanged.
\item[(iii)] On the polydisc $|\mathsf{B}|<T^{\mathbb{C}}$, of radius $T^{\mathbb{C}}$, of the
weak slot of the correlation theorem, the condition
\begin{equation}\label{9.1:eq-homogeneous-accommodation-2}
K_R^{\sharp}(\Theta_{\flat}+1)\,c_A<\tfrac14,
\end{equation}
in which $\Theta_{\flat}$ is the constant of the correlation theorem entering it, does not
involve the coupling, and it is met by the choice of $c_A$.
\item[(iv)] In the quotient $\mathsf{w}_T/K_o$ formed in the statement of this chapter that
bounds the room at the dilation circle, the
dilation circle $\mathsf{w}_T$ is the one defined with $K_R^{\sharp}$ as in~(ii).
\end{itemize}
In the remainder of this chapter Lemma~\ref{9.1:lem-small-field-accommodation} is used in the
conjoined form of~(i), with $K_R^{\sharp}$ in place of $K_R^o$ and of $K_R$.
\end{remark}

\begin{proof}
By the hypothesis, the inequalities of the response estimate at the small-field site, on the
polydisc of \cite[(3.43)]{B-CMP119} on which it is stated, have right members linear in
$\mathfrak{w}_k$. A dilation with $|\mathsf{s}(\Delta)|\le e^{\kappa_1}$, and likewise the slice
move, multiplies the common profile of these right members by a factor that is absorbed by
replacing $C_H$ with $C_He^{16\kappa_1}$; by linearity in $\mathfrak{w}_k$ the inequalities
persist in the same form at the dilated points and after the slice move.

By \eqref{9.1:eq-homogeneous-accommodation-1} we have $K_R^{\sharp}\ge K_R^o$. Since
$g_k\sup|\mathsf{B}|\ge0$, the hypothesis of~(i) implies that of
Lemma~\ref{9.1:lem-small-field-accommodation}:
\[
(1+|\mathsf{t}|)K_R^{\sharp}g_k\sup|\mathsf{B}|\le\alpha_{0,k}
\ \Longrightarrow\
(1+|\mathsf{t}|)K_R^{o}g_k\sup|\mathsf{B}|\le\alpha_{0,k}.
\]
Together with $|\mathsf{s}(\Delta)|\le
e^{\kappa_1}$, this gives the conclusions of that lemma.

Also by \eqref{9.1:eq-homogeneous-accommodation-1}, $K_R^{\sharp}\ge2K_R$, so the allotment
$K_Rg_k\sup|\mathsf{B}|$ that the correlation theorem assigns to the undilated point satisfies
\[
K_Rg_k\sup|\mathsf{B}|\le\tfrac12K_R^{\sharp}g_k\sup|\mathsf{B}|.
\]
Thus the undilated point uses at most half of the allotment
$K_R^{\sharp}g_k\sup|\mathsf{B}|$. This is the hypothesis that the undilated point uses at most
half of the allotment, in the statement of this chapter whose hypothesis it is; it holds by
\eqref{9.1:eq-homogeneous-accommodation-1} and needs no separate ceiling. This completes~(i).

For (ii), the constant $K_R$ is fixed in the correlation theorem, and the constant $K_R^o$
depends only on $(d,L,\theta,M_1,R,\kappa_1;\theta_S)$.
Consequently $K_R^{\sharp}$ is available before the definitions of $\mathsf{w}_T$, $\theta_k$,
$\theta^{\mathbb{C}}_k$, $c_V$ and $c_A$, which precede $\gamma$. The
correlation theorem is used with its statement unchanged; the replacement is confined to the run
of the correlation theorem on $\mathfrak{S}$ that is used for the uniqueness theorem.

For (iii), the left member of \eqref{9.1:eq-homogeneous-accommodation-2} contains no coupling.
The constant $c_A$ is defined after $K_R^{\sharp}$ has been fixed, and
\eqref{9.1:eq-homogeneous-accommodation-2} is met by the choice of $c_A$, in the same way as at
the place of the correlation theorem where $c_A$ is chosen.

Clause~(iv) is the definition made in~(ii).
\end{proof}

\begin{lemma}[The untouched region and the middle shells]\label{9.1:lem-untouched-region}
Consider Ba{\l}aban's large-field operation $\mathbf{R}$ at step $k$, carried out in stages, in the
construction of \cite[(1.10)--(1.17)]{B-CMP122a}. From that construction we use:
\begin{itemize}
\item the decreasing domains $\Omega_j$, and their enlargements $\Omega^{\sim4}_j\supset\Omega_j$
\cite[\S1]{B-CMP122a};
\item the domain $\Lambda_k$ and the region $Z''_k$;
\item the integers $k_0<k$ and $N_0$, related by $k_0=k-N_0+1$;
\item the radii $\check{R}_j$, written $R_j$ in \cite{B-CMP122a}.
\end{itemize}
Let $Z$ be the union of the components processed at step $k$. Every point $x$ carries a label
assigned by that construction, in particular a level $m(x)$.
Let $d(\cdot,\cdot)$ be the distance of the random-walk estimates of this chapter (not the dimension
$d=4$), and $\mathsf{g}$ the constant of the first property of the distance $d$ (the constant that
appears in the inherited level-distance inequality). By its definition in this chapter, $d$ is a path
distance: $d(x,y)$ is the
infimum, over paths from $x$ to $y$, of the length of the path, and the length of a path is at least
the sum of the $d$-distances between the end points of any family of segments of it which pairwise
have no point in common other than a common end point.
For each stage $\nu$, let $R_\nu$ be the region of that stage and $j_\nu$ the step at which it was
created. Put
\begin{equation}\label{9.1:eq-untouched-region-1}
\begin{gathered}
\mathfrak{O}:=\Omega_k\cup Z^c,\qquad Z^{\mathrm{on}}:=Z\cap\Omega_k^c,\qquad
Z^{\mathrm{mid}}:=Z^{\mathrm{on}}\setminus Z''_k,\\
\Lambda:=\bigl(\Omega^{\sim4}_{k_0+1}\bigr)^c\cap Z,\qquad
S_j:=Z\cap(\Omega_j\setminus\Omega_{j+1})\quad(k_0<j<k).
\end{gathered}
\end{equation}
We call the sets $S_j$, $k_0<j<k$, the \emph{middle shells}. A \emph{traversal} of $S_j$ by a path
is a segment of the path that lies in $S_j$ and has one end point adjacent to $\Omega_{j+1}$ and the
other end point adjacent to $\Omega_j^c$. A set is said to be \emph{enclosed} by $S_j$ if every path
from a point of $\mathfrak{O}$ to a point of the set contains a traversal of $S_j$; and $S_j$ is said
to be at least a given number of level-$k$ units \emph{wide} if the two end points of every
traversal of $S_j$ are at distance $d$ at least that number of level-$k$ units.
Assertions (i)--(v) concern only points of $\mathfrak{O}$ and of those components of $Z$ that are
not excluded from $Z$ in the course of step $k$. Then:
\begin{enumerate}
\item[(i)] no stage and no application of $\mathbf{R}$ changes the label of a point of
$\mathfrak{O}$; the level $m(x)\le k$ of $x\in\mathfrak{O}$ is the same at every stage;
\item[(ii)] $Z''_k\subset\Lambda$, $Z^{\mathrm{on}}=Z''_k\cup Z^{\mathrm{mid}}$ and
$S_j\subset Z^{\mathrm{mid}}$ for $k_0<j<k$; at the last stage of step $k$ every point of $S_j$ has
level $k$, and each $S_j$ is at least $2M$ level-$k$ units wide;
\item[(iii)] the components of $\Lambda_k^c$ are separated by $\Lambda_k$, and
$\Lambda_k\subset\Omega_k$;
\item[(iv)] for $x\in\mathfrak{O}$ of level $m$ and $y\in R_\nu$,
\begin{equation}\label{9.1:eq-untouched-region-2}
d(x,y)\ \ge\ \mathsf{g}\bigl[(k-j_\nu-2)_+ + (k-m-1)_+\bigr];
\end{equation}
\item[(v)] for $x\in\mathfrak{O}$ of level $m$ and every set $A\subset Z$ enclosed by every middle
shell,
\begin{equation}\label{9.1:eq-untouched-region-3}
d(x,A)\ \ge\ 2M(N_0-2)+\mathsf{g}\,(k-m-1)_+ ;
\end{equation}
this applies to $A=\Lambda$, which is enclosed by every middle shell by \cite[(1.73)]{B-CMP122a},
and to the set $\bar{\Lambda}$ of the construction of the step, which, like $\Lambda$, is enclosed
by every middle shell.
\end{enumerate}
\end{lemma}

\begin{proof}
\emph{The region $\mathfrak{O}$ and the excluded components.} Assertion (i) rests on the
construction of the step in \cite[(1.10)--(1.17)]{B-CMP122a}. In that construction neither a stage
nor $\mathbf{R}$ changes the label of a point of $\Omega_k\cup Z^c$.

In the course of the step, \cite[\S1]{B-CMP122a} states:
\begin{quotation}
we exclude from $Z$ the components with the large field functions $1-\chi^{(k+1)}$
\end{quotation}
The remaining region is again denoted $Z$ there. Let $W$ be a component excluded in this way; here,
as in the definition of $\mathfrak{O}$, $Z$ is the union of all components processed at step $k$,
$W$ among them. During step $k$ such a component can be counted neither with $\mathfrak{O}$, since
the stages up to and including its exclusion have changed its level, nor with the surviving part of
$Z^{\mathrm{on}}$, since no later stage raises its level. Such a component therefore forms
a third region; it is treated separately, with its own sequence of stages, later in this chapter.
This is why (i)--(v) concern only $\mathfrak{O}$ and the surviving part of $Z$.

\emph{The decomposition of $Z^{\mathrm{on}}$.} By the construction of \cite[\S1]{B-CMP122a},
$Z''_k\subset\Lambda$. Since $k_0+1\le k$, the domains $\Omega_j$ decrease, and
$\Omega_{k_0+1}\subset\Omega^{\sim4}_{k_0+1}$ by \cite[\S1]{B-CMP122a}, we have
\begin{equation*}
\Omega_k\subset\Omega_{k_0+1}\subset\Omega^{\sim4}_{k_0+1}.
\end{equation*}
Hence $\Lambda\subset Z\cap\Omega_k^c=Z^{\mathrm{on}}$, and so $Z''_k\subset Z^{\mathrm{on}}$. The
identity $Z^{\mathrm{on}}=Z''_k\cup Z^{\mathrm{mid}}$ then follows from the definition of
$Z^{\mathrm{mid}}$ in \eqref{9.1:eq-untouched-region-1}.

Now let $k_0<j<k$. Since $j+1\le k$, we have $\Omega_k\subset\Omega_{j+1}$. Therefore
$\Omega_j\setminus\Omega_{j+1}\subset\Omega_k^c$, and so $S_j\subset Z^{\mathrm{on}}$. Moreover
$S_j\subset\Omega_j$, while, by the inclusion $\Omega_{k_0+1}\subset\Omega^{\sim4}_{k_0+1}$ used
above and since $j\ge k_0+1$,
\begin{equation*}
\Lambda\subset\Omega_{k_0+1}^c\subset\Omega_j^c .
\end{equation*}
Hence $S_j\cap\Lambda=\emptyset$, and in particular $S_j\cap Z''_k=\emptyset$. This gives
$S_j\subset Z^{\mathrm{mid}}$. By the construction of \cite[(1.10)--(1.17)]{B-CMP122a}, at the last
stage of step $k$ every point of $Z^{\mathrm{on}}$, in particular of $S_j$, carries level $k$.

\emph{Width of the shells.} By the construction of the domains in
\cite[(1.10)--(1.17)]{B-CMP122a}, the shell $S_j$ is at least $2M\check{R}_jL^{-(k-j)}$ level-$k$
units wide, and
\begin{equation}\label{9.1:eq-untouched-region-4}
2M\check{R}_jL^{-(k-j)}\ \ge\ 2M .
\end{equation}
Inequality \eqref{9.1:eq-untouched-region-4} holds by the definition of $N_0$ in
\cite[\S1]{B-CMP122a}, which reads:
\begin{quotation}
Take the smallest positive integer $N_0$ such, that $L^{-N_0+1}MR_{k-N_0+1}=M$
\end{quotation}
where $R_j$ is the radius written $\check{R}_j$ in this book.
This completes (ii). For (iii), $\Lambda_k\subset\Omega_k$ holds by the construction of
\cite[\S1]{B-CMP122a}. Distinct components of $\Lambda_k^c$ are separated by $\Lambda_k$.

\emph{Proof of \eqref{9.1:eq-untouched-region-2}.} Inequality \eqref{9.1:eq-untouched-region-2} is
the inherited level-distance inequality, stated for the distance $d$; we prove it as follows. Let
$x\in\mathfrak{O}$ have level $m$, and let $y\in R_\nu$. By the construction of the stage regions
$R_\nu$, we have $R_\nu\subset\Omega^c_{j_\nu+2}$.

Every path from $x$ to $y$ passes through a point $z$ of level $k$. Indeed, the components of
$\Lambda_k^c$ are separated by $\Lambda_k\subset\Omega_k$. Moreover,
$\Lambda=(\Omega^{\sim4}_{k_0+1})^c\cap Z$ is enclosed, in the sense fixed above, by every shell
$S_j$ \cite[(1.73)]{B-CMP122a}.

Fix one such point $z$. The first property of the distance $d$, applied twice, bounds the
$d$-distance between the end points of each of the two segments of the path:
\begin{itemize}
\item to the segment from $x$ to $z$, across the level gap from $m$ to $k$, which yields
$\mathsf{g}(k-m-1)_+$;
\item to the segment from $z$ to $y$, across the level gap from $k$ to a point of
$\Omega^c_{j_\nu+2}$, which yields $\mathsf{g}(k-j_\nu-2)_+$.
\end{itemize}
Since every path from $x$ to $y$ passes through such a point $z$, and the two segments have only $z$
in common, the path property of $d$ stated after its definition shows that $d(x,y)$ is at least
the sum of the two bounds. This gives \eqref{9.1:eq-untouched-region-2}.

\emph{Proof of \eqref{9.1:eq-untouched-region-3}.} Neither side of
\eqref{9.1:eq-untouched-region-3} depends on the stage: the level $m$ of $x$ does not, by (i), and
neither $d$ nor $A$ does. We may therefore argue at the last stage of step $k$, at which, by (ii),
every point of every middle shell has level $k$. Consider a path from $x$ to a point of $A$. By the
same separation it
passes through a point of level $k$; let $z$ be the first such point on the path. The segment from
$x$ to $z$ contributes $\mathsf{g}(k-m-1)_+$, by the first property of the distance $d$. The middle
shells $S_j$, $k_0<j<k$, enclose $A$ by hypothesis, so the path contains a traversal of each of
them,
in the sense fixed after \eqref{9.1:eq-untouched-region-1}. Every point of a traversal lies in a
middle shell and therefore has level $k$; since $z$ is the first point of level $k$ on the path,
every traversal lies in the part of the path from $z$ on, and meets the segment from $x$ to $z$ at
most in $z$. The shells are pairwise disjoint, because the sets $\Omega_j\setminus\Omega_{j+1}$ are
pairwise disjoint for the decreasing domains $\Omega_j$; hence the traversals of distinct shells
are pairwise disjoint. There are $k-k_0-1=N_0-2$ shells, since $k_0=k-N_0+1$. By (ii) and
\eqref{9.1:eq-untouched-region-4}, each shell is at least $2M$ level-$k$ units wide, so the two end
points of each traversal are at distance $d$ at least $2M$. The segment from $x$ to $z$ and the
$N_0-2$ traversals are thus segments of one path which pairwise have no point in common other than
a common end point, and by the path property of $d$ the length of the path is at least
\begin{equation*}
2M(N_0-2)+\mathsf{g}\,(k-m-1)_+ .
\end{equation*}
Taking the infimum over all paths from $x$ to points of $A$ gives
\eqref{9.1:eq-untouched-region-3}.
\end{proof}

\begin{remark}[Failure of a common entry table on the untouched region]
\label{9.1:rem-entry-failure}
Let $k$ be a step and let $\Phi$ be a term produced by Ba{\l}aban's operation $\mathbf{R}$ at
step $k$. Write $\varrho_\Phi$ for the room of $\Phi$ and $\mathfrak{T}_\Phi[7/8]$ for its
seed tube, whose members are called seeds. Let
$x\in\mathfrak{O}$ be a point of level $m$, where $\mathfrak{O}$ is the
untouched region of Lemma~\ref{9.1:lem-untouched-region}.

The remaining objects are as follows:
\begin{itemize}
\item $N\ge1$ is the stage index (not the $N$ of $\mathrm{SU}(N)$), and $h$ denotes the level
$k-N$;
\item $\mathsf{T}_N$ is the stage-$N$ table, and $\mathsf{c}_N(m)$ is its coefficient at
level $m$; the entries of a seed at $x$ are the quantities at $x$ that the stage tables bound,
and $E_6$ is one of the members of these tables;
\item $f_{k,\cdot}$ is the target profile of step $k$, and $\mathsf{u}$ is the unit in which
the table entries at the level read are measured;
\item $\bar u=\bar u(x)$ is the local radius at $x$, and $a^0_\nu$ are the shares of the
stages $\nu$;
\item $h_\mu$ are the move fields, indexed by $\mu$ (not to be confused with the level $h$),
$K_o$ is the move constant, and $\beta$, $\theta_\rho$, $\theta_S$ are the table and
schedule constants; for a term $F$ of the expansion, $\mathfrak{D}_F$ denotes its own tube, the
domain of configurations on which $F$ is read; a leaf is a term $F$ of the expansion, and the
tube of a leaf $F$ is its own tube $\mathfrak{D}_F$.
\end{itemize}
We call \emph{common-target clause} the clause of the theorem of this chapter that fixes common
entry tables at the $\mathbf{R}$-site.
\begin{itemize}
\item[(i)] One has $\mathsf{c}_N(m)<f_{k,m}-\tfrac18\varrho_\Phi$, so that there are seeds of
$\mathfrak{T}_\Phi[7/8]$ whose entries at $x$ of the second kind listed in that clause
lie outside $\mathsf{T}_N$.
\item[(ii)] Suppose that entries are placed in $f_{k,\cdot}\,\mathsf{u}$ and that the moves are
known only through the bound $K_o|h_\mu|\le a^0_\nu$. Then landing in the tube of
a leaf of level $k$ produced by the operation $\mathbf{T}$, at a level $m\le h$, requires
$K_o\ge c\,2^N$ for a constant $c>0$ independent of $N$.
\end{itemize}
The point of the remark is that neither the common-target clause, nor the condition of that
theorem which
is verified by means of the room increments $\mathrm{rm}_\nu$ of the stages, covers
$\mathfrak{O}$. The region
$\mathfrak{O}$ is therefore treated separately in the statements that follow.
\end{remark}

\begin{proof}
(i) Reading the stage table at level $m$ from \cite[(1.64)--(1.65)]{B-CMP122a}, as in the
read-off of that table in this chapter, gives
\begin{equation}\label{9.1:eq-entry-failure-1}
\mathsf{c}_N(m)=f_{k,m}-\beta\sum_{i=h+1}^{k}2^{-|m-i|}.
\end{equation}
Since $h=k-N<k$, the sum in \eqref{9.1:eq-entry-failure-1} contains the term $i=k$. That term
equals $2^{-|m-k|}=2^{-(k-m)}$, because $m\le k$ for a point of $\mathfrak{O}$ by
Lemma~\ref{9.1:lem-untouched-region}. All terms of the sum are positive, so the
sum is at least $2^{-(k-m)}$. Moreover, since the schedule constant satisfies $\theta_\rho<1$,
one has $2^{-(k-m)}>\tfrac18\theta_\rho 2^{-(k-m)}$.
Multiplying by $\beta>0$ and inserting the result into \eqref{9.1:eq-entry-failure-1}, we obtain
\begin{equation}\label{9.1:eq-entry-failure-2}
f_{k,m}-\mathsf{c}_N(m)\ \ge\ \beta\,2^{-(k-m)}\ >\ \tfrac18\,\theta_\rho\,\beta\,2^{-(k-m)}.
\end{equation}
This is the inequality of~(i), the room $\varrho_\Phi$ being read at level $m$, where the room
schedule gives $\varrho_\Phi\le\theta_\rho\beta2^{-(k-m)}$.

Each seed of $\mathfrak{T}_\Phi[7/8]$ is compared with the root of the chain at $x$ through a
chain of intermediate objects, each of which is read through its slice. By the slice bounds for
chain objects of this chapter, the slice of every chain object lies
within $\tfrac1{128}\varrho_\Phi\,\mathsf{u}$ of the root of the chain at $x$. Hence, if the
root entry of that kind at $x$ is at least $(f_{k,m}-\tfrac18\varrho_\Phi)\mathsf{u}$, the
corresponding entry stays at least $(f_{k,m}-\tfrac{17}{128}\varrho_\Phi)\mathsf{u}$ along the
whole chain. Since $\varrho_\Phi\le\theta_\rho\beta2^{-(k-m)}$ and $\theta_\rho<1$,
\eqref{9.1:eq-entry-failure-2} gives
\begin{equation*}
f_{k,m}-\mathsf{c}_N(m)\ \ge\ \beta\,2^{-(k-m)}\ >\ \tfrac{17}{128}\,\theta_\rho\,\beta\,
2^{-(k-m)}\ \ge\ \tfrac{17}{128}\,\varrho_\Phi .
\end{equation*}
Thus every seed of $\mathfrak{T}_\Phi[7/8]$ whose root entry of that kind at $x$ lies at or
above $(f_{k,m}-\tfrac18\varrho_\Phi)\mathsf{u}$ keeps that entry above
$\mathsf{c}_N(m)\mathsf{u}$, that is, outside $\mathsf{T}_N$, along the whole chain.

It remains to exhibit such a seed of $\mathfrak{T}_\Phi[7/8]$. It is obtained in the same way
as in the abelian plateau
witness, namely from an abelian plateau in the table member $E_6$.

(ii) At $x$, which lies at a level $m\le h$, the shares of the stages satisfy, by their
definition,
\begin{equation*}
\sum_\nu a^0_\nu(x)\ \ge\ \beta\,2^{-(h+1-m)}\,\bar u .
\end{equation*}
The whole margin of the leaf at level $m$ is $\theta_S\beta 2^{-(k-m)}\mathsf{u}$. When the
moves are known only through $K_o|h_\mu|\le a^0_\nu$, the bound this provides on
the composite move at $x$ is $K_o^{-1}\sum_\nu a^0_\nu(x)$, and this bound lies within the
margin only if $K_o$ is at least the ratio of $\sum_\nu a^0_\nu(x)$ to the margin. Since the
local radius $\bar u(x)$ is at least the unit $\mathsf{u}$ at the level read, the lower bound on
the shares gives for this ratio at least the quotient
of the coefficient $\beta2^{-(h+1-m)}$ of that lower bound by the coefficient
$\theta_S\beta2^{-(k-m)}$ of the margin, namely
\begin{equation}\label{9.1:eq-entry-failure-3}
\frac{\beta\,2^{-(h+1-m)}}{\theta_S\,\beta\,2^{-(k-m)}}
=\frac{2^{k-h-1}}{\theta_S}
=\frac{2^{N-1}}{\theta_S},
\end{equation}
where the last equality uses $k-h=N$. Landing in the tube of the leaf therefore requires
$K_o\ge 2^{N-1}/\theta_S$, that is, $K_o\ge c\,2^N$ with $c=1/(2\theta_S)$, a constant
independent of $N$.
\end{proof}

\begin{definition}[Two-region tables with response-sized rooms]\label{9.1:def-two-region-tables}
Work at an $\mathbf{R}$-site at step $k$, with the untouched region $\mathfrak{O}$ and the level
function $m=m(x)\le k$ of Lemma~\ref{9.1:lem-untouched-region}. Let $\mathsf{T}_\nu$,
$\nu\le N$, be the stage tables, indexed by the stage $\nu$. Here $N$ is the index of the last
stage at the site.
The following objects are those fixed together with the stage tables $\mathsf{T}_\nu$:
\begin{itemize}
\item the coefficient $f_{k,m}$, the schedule letter $\theta_S$, and the table number $\beta$
that enters the table coefficient $1-4\beta$;
\item the unit $\mathsf{u}^a_m$ of table member $a$ at level $m$, and the exponent $e_a$ of
member $a$;
\item the factor $\ell_m$;
\item the index $j_\nu$ and the set $R_\nu$ attached to stage $\nu$;
\item the constants $C'_{\flat}$, $c_M$ and $c_1$, the decay rate $\delta_1$, and the local
Lipschitz constant $B_{\natural}$.
\end{itemize}
Moreover, $\alpha_{0,j}$ is Ba{\l}aban's radius at level $j$, one of the radii among the
auxiliary parameters $\mathfrak{p}$, and $d(\cdot,R_\nu)$ is the distance to $R_\nu$. The tables,
their units and the responses defined below have one member for each index $a$. Where $a$ is
not written, as in \eqref{9.1:eq-two-region-tables-2} and below, it is suppressed, and every
identity holds member by member. The \emph{flat tables} $\mathsf{T}^{\flat}_\nu$ are defined
separately on two regions: the region $Z^{\mathrm{on}}$, formed from the processed components
$Z$ of Lemma~\ref{9.1:lem-untouched-region}, and the region $\mathfrak{O}$.
\begin{itemize}
\item[(a)] On $Z^{\mathrm{on}}$ we put $\mathsf{T}^{\flat}_\nu:=\mathsf{T}_\nu$.
\item[(b)] On $\mathfrak{O}$ we define the profile $\mathsf{h}_\nu$ and the response
$\mathsf{g}^a_\nu$ of stage $\nu$ by
\begin{equation}\label{9.1:eq-two-region-tables-1}
\mathsf{h}_\nu:=C'_{\flat}\,c_M\,\alpha_{0,j_\nu}\,e^{-\delta_1 d(\cdot,R_\nu)/4},
\qquad
\mathsf{g}^a_\nu:=2^8 B_{\natural}\,c_1\,\ell_m^{-e_a}\,\mathsf{h}_\nu .
\end{equation}
For $\nu\le N$ we put
\begin{equation}\label{9.1:eq-two-region-tables-2}
\mathsf{T}^{\flat}_\nu:=\Bigl(f_{k,m}+\tfrac12\theta_S\,\beta\,2^{-(k-m)}\Bigr)\mathsf{u}_m
+\sum_{\nu\le i<N}\mathsf{g}_i ,
\end{equation}
the sum being empty at $\nu=N$.
\end{itemize}
The coefficient $c=3$, which \cite[(1.68)]{B-CMP122a} attaches to $m=k$, does not appear
in $\mathsf{T}^{\flat}_\nu$ (see the proof below). For $\nu<N$ we write
$\mathrm{rm}^{\flat}_\nu:=\mathsf{T}^{\flat}_\nu-\mathsf{T}^{\flat}_{\nu+1}$ for the room
increment of stage $\nu$. It is formed from the flat tables in the same way as the room
increment $\mathrm{rm}_\nu$ is formed from the stage tables. For each $i<N$ the chain terms
$\mathbb{C}^{(i)}$ of the preliminary integration are considered on the domain
$\mathfrak{U}[\mathsf{T}^{\flat}_i]$ of the table $\mathsf{T}^{\flat}_i$. This domain is
contained in the flat domain $\mathfrak{D}^{\flat}_i$ on which these terms are represented; this
inclusion is proved in the lemma on the inclusion of table domains in flat domains.
\end{definition}

\begin{proof}[Proof of the room identity and of the omission of $c=3$]
\emph{Room identity.} On $\mathfrak{O}$ the first term of \eqref{9.1:eq-two-region-tables-2}
depends on $k$ and $m$ but not on $\nu$. Subtract \eqref{9.1:eq-two-region-tables-2} at
$\nu+1$ from \eqref{9.1:eq-two-region-tables-2} at $\nu$. That term cancels, and of the two
sums only the summand $i=\nu$ survives. Hence
\begin{equation}\label{9.1:eq-two-region-tables-3}
\mathrm{rm}^{\flat}_\nu
=\sum_{\nu\le i<N}\mathsf{g}_i-\sum_{\nu+1\le i<N}\mathsf{g}_i
=\mathsf{g}_\nu\qquad\text{on }\mathfrak{O}.
\end{equation}
Thus on $\mathfrak{O}$ the room of a stage is exactly its response $\mathsf{g}_\nu$. This is
the device of \cite{B-CMP119}, where the difference of two consecutive bounds is, in
Ba{\l}aban's words, ``still much greater than bounds on this function''.

\emph{Omission of $c=3$.} In \cite[(1.68)]{B-CMP122a} the coefficient is $c=1$ for $m<k$
and $c=3$ for $m=k$, and it enters at the stages $\nu<N$. The case $m=k$ is the region
$\Omega_k$, and $\Omega_k\subset\mathfrak{O}$ by the definition of $\mathfrak{O}$ in
Lemma~\ref{9.1:lem-untouched-region}. Two facts make the coefficient unnecessary there.

First, consider the table condition that the stage tables are built to satisfy. Its clause (i)
on $\Omega_k$ itself holds by \cite[(1.64)(i)]{B-CMP122a} taken at $m=j=k$, and in that bound
no coefficient $c$ occurs. At $j=k$ its clause (iii) is empty.

Second, by the proposition of this chapter that no stage before the last consumes $c=3$ on
$\Omega_k$, no bound read at a stage $\nu<N$ involves the coefficient $c=3$ on $\Omega_k$. That
proposition is needed only for the flat tables at the stages $\nu<N$: at the stage $\nu=N$ the
coefficient $c$ of \cite[(1.68)]{B-CMP122a} does not enter, so at $\nu=N$ nothing is omitted
in \eqref{9.1:eq-two-region-tables-2}.

Consequently $c=3$ may be dropped from the tables on $\mathfrak{O}$.
\end{proof}

\begin{lemma}[Accommodation on the untouched region]\label{9.1:lem-untouched-accommodation}
Let $k$ be a step at which the operation $\mathbf{R}$ acts, let $\mathfrak{O}$ be the untouched region
of Lemma~\ref{9.1:lem-untouched-region}, let $\mathsf{T}^{\flat}_\nu$ be the two-region tables of
Definition~\ref{9.1:def-two-region-tables}, and let a point of $\mathfrak{O}$ of label level $m\le k$ be
given. For each stage $i$ let $h_i$ be the move made at stage $i$, let $j_i$ be its level, and let
$\mathsf{h}_i$ be its profile majorant. Assume
\begin{equation}\label{9.1:eq-untouched-accommodation-1}
2^{18}B_\natural\, c_1\, C_\flat'\, c_M \le \theta_S\beta .
\end{equation}
Here $B_\natural$ is the constant of hypothesis (A) of Lemma~\ref{9.1:lem-site-blind-slice}, $c_1$ is
the constant of Lemma~\ref{9.1:lem-top-source-decay}(i), and $C_\flat'$, $c_M$ are the constants of the
size bound on the profile majorants $\mathsf{h}_i$ and of the floor condition on the exponents, which
together give, at a point $x$ of label level $m$,
\begin{equation*}
\mathsf{h}_i(x)\le C_\flat' c_M\,\alpha_{0,j_i}\,L^{-4(k-j_i-2)_+}\,L^{-4(k-m-1)_+}.
\end{equation*}
Then the following hold.
\begin{itemize}
\item[(a)] Let $N$ be the number of stages, so that the stages are $0\le i<N$ and $\mathsf{T}^a_N$ is
the last stage table. The room increments $\mathsf{g}^a_i$ of the tables satisfy
\begin{equation*}
\sum_{i<N}\mathsf{g}^a_i \le \tfrac{1}{16}\,\theta_S\beta\, 2^{-(k-m)}\,\mathsf{u}^a_m .
\end{equation*}
\item[(b)] In the notation of Lemma~\ref{9.1:lem-site-blind-slice} ($\ell_\iota$, $e_a$,
$\tilde\Delta(y)$, $C_{\mathsf{m}}$), let $\nu<n_e$ be stages and consider the moves $h_i$ made at the
stages in $[\nu,n_e)$.
Then the profile $\mathsf{m}=2\sum_i\mathsf{h}_i$ satisfies the conditions on $\mathsf{m}$ of
Lemma~\ref{9.1:lem-site-blind-slice} with $C_{\mathsf{m}}=c_1$, and hypothesis (A) of that lemma holds
on $\mathfrak{O}$ with $\mathsf{T}_e=\mathsf{T}^{\flat}_{n_e}$ and $\mathsf{T}_r=\mathsf{T}^{\flat}_\nu$.
\item[(c)] At the points of class $(=)$ in $\mathfrak{O}$, for every stored term $F$ read at the
operation $\mathbf{R}$ of step $k$,
\begin{equation}\label{9.1:eq-untouched-accommodation-2}
\mathsf{T}^{\flat}_0 \le (\mathsf{c}_F^{+}-\varrho_F)\,\mathsf{u}_F
-\tfrac{3}{16}\,\theta_S\beta\, 2^{-(k-m)}\,\mathsf{u}_m .
\end{equation}
\end{itemize}
None of the number $N$ of stages, the integer $N_0$, the history, the depth and the coupling $g$
enters these bounds.
\end{lemma}

\begin{proof}
Put $s:=k-m$. For each move $i$ let $j_i$ be its level and $\mathsf{h}_i$ its profile majorant, as
in the statement.

(a) By the bound recalled after \eqref{9.1:eq-untouched-accommodation-1}, written with $s=k-m$,
\begin{equation*}
\mathsf{h}_i(x)\le C_\flat' c_M\,\alpha_{0,j_i}\,L^{-4(k-j_i-2)_+}\,L^{-4(s-1)_+}.
\end{equation*}
The levels $j_i$ are distinct integers not exceeding $k$. Hence each value $t\ge1$ of
$(k-j_i-2)_+$ is taken by at most one
$i$, and these terms contribute at most $\sum_{t\ge1}L^{-4t}\le 1$; the terms with exponent zero
are those with $j_i\in\{k-2,k-1,k\}$, at most three. Therefore
$\sum_i L^{-4(k-j_i-2)_+}\le 4$. By the comparison of the radii $\alpha_{0,\cdot}$ between levels,
\begin{equation*}
\alpha_{0,j_i}\le 2\alpha_{0,k}\le \tfrac92\,(1\vee s)^{1/2}\,\alpha_{0,m}.
\end{equation*}
Combining these three bounds,
\begin{equation*}
\sum_i\mathsf{h}_i(x)\le 18\, C_\flat' c_M\,(1\vee s)^{1/2}L^{-4(s-1)_+}\,\alpha_{0,m}.
\end{equation*}
Let $q(s):=2^s(1\vee s)^{1/2}L^{-4(s-1)_+}$. Then $q(0)=1$, $q(1)=2$ and $q(2)=4\sqrt2\,L^{-4}<1$,
because $L$ is an odd integer greater than one, so that $L\ge3$ and $L^{4}\ge81>4\sqrt2$. For $s\ge2$,
\begin{equation*}
\frac{q(s+1)}{q(s)}=2\Bigl(\frac{s+1}{s}\Bigr)^{1/2}L^{-4}<1 .
\end{equation*}
Thus $q(s)\le2$ for all $s\ge0$, and
\begin{equation}\label{9.1:eq-untouched-accommodation-3}
\sum_i\mathsf{h}_i\le 36\, C_\flat' c_M\, 2^{-s}\,\alpha_{0,m}.
\end{equation}
The room increments are obtained from the profiles through the identity
$2^7B_\natural c_1\cdot 2\sum_i\mathsf{h}_i=\ell_\iota^{e_a}\sum_i\mathsf{g}^a_i$, by the definition of
the room increments $\mathsf{g}^a_i$ in terms of the profiles.
Hence, by \eqref{9.1:eq-untouched-accommodation-3} and \eqref{9.1:eq-untouched-accommodation-1},
\begin{align*}
\sum_i\mathsf{g}^a_i
&\le 2^8\cdot36\; B_\natural c_1 C_\flat' c_M\, 2^{-s}\,\ell_\iota^{-e_a}\alpha_{0,m}\\
&\le \frac{2^8\cdot 36}{2^{18}}\;\theta_S\beta\,2^{-s}\,\ell_\iota^{-e_a}\alpha_{0,m}.
\end{align*}
Since $2^8\cdot36\cdot16=147456\le 2^{18}$, the factor is at most $\tfrac1{16}$. This is (a),
since the unit of member $a$ at level $m$ is $\mathsf{u}^a_m=\ell_\iota^{-e_a}\alpha_{0,m}$ by the
definition of the units of the tables.

(b) By the rate bound on the moves, the majorant satisfies $\ell_\iota|h_i|\le\mathsf{h}_i$. The
composite move
$\mathcal{Y}$ is assembled from the moves $h_i$ by the Baker--Campbell--Hausdorff formula, which
costs at most a factor $2$. Hence
\begin{equation*}
\ell_\iota\sup_{\tilde\Delta(y)}|\mathcal{Y}|\le 2\sum_i\mathsf{h}_i(y)=\mathsf{m}(y).
\end{equation*}

For slowness, apply Lemma~\ref{9.1:lem-top-source-decay}(i) to each profile $\mathsf{h}_i$. This
gives $\mathsf{h}_i^{\approx}\le\mathsf{h}_i^{\sim}\le c_1\mathsf{h}_i$. The smearing
$(\cdot)^{\approx}$ is taken with the weight $e^{-2\delta_1 d_{\mathfrak{S}}}$, with $\delta_1$ as in
Lemma~\ref{9.1:lem-top-source-decay}. The label of a later
stage is coarser, so the distance $d(\cdot,\mathcal{S})$ to its source set is smaller, and slowness
with respect to $d(\cdot,\mathcal{S})$ passes to $d_{\mathfrak{S}}$.
Summing over $i$, which is allowed because smearing is linear, gives $\mathsf{m}^{\approx}\le c_1\mathsf{m}$.
So the requirement on $\mathsf{m}$ holds with $C_{\mathsf{m}}=c_1$.

Finally, at the points of class $(=)$ in $\mathfrak{O}$,
\begin{equation*}
2^7B_\natural C_{\mathsf{m}}\mathsf{m}=2^7B_\natural c_1\cdot 2\sum_i\mathsf{h}_i
=\ell_\iota^{e_a}\sum_i\mathsf{g}^a_i .
\end{equation*}
Here the sum runs over the stages in $[\nu,n_e)$, and the right-hand side is at most
$\ell_\iota^{e_a}\mathsf{G}^a$ with $\mathsf{G}=\mathsf{T}^{\flat}_\nu-\mathsf{T}^{\flat}_{n_e}$, the room
increments of the stages in $[\nu,n_e)$ making up the difference of the two tables. This is
hypothesis (A).

(c) Write $\mathsf{u}$ for the unit at the point read, and $f_{j,m}$ for the table coefficient
attached to step $j$ at a point of label level $m$. The unit at the point read is
$\mathsf{u}=\mathsf{u}_m$, and $\mathsf{u}_m\le\mathsf{u}_F$. By the definition of the tables
$\mathsf{T}^{\flat}_\nu$ in Definition~\ref{9.1:def-two-region-tables}, together with part (a), at the
points of class $(=)$ in $\mathfrak{O}$,
\begin{equation*}
\mathsf{T}^{\flat}_0\le\bigl(f_{k,m}+\tfrac{9}{16}\,\theta_S\beta\,2^{-s}\bigr)\mathsf{u}.
\end{equation*}
We distinguish two cases according to where $F$ was created.

First, suppose $F$ was created by the operation $\mathbf{T}$ at step $j=k$. Then, by the choice of the
room coefficient of a term created by $\mathbf{T}$ at step $k$,
\begin{equation*}
\mathsf{c}_F^{+}-\varrho_F\ge f_{k,m}+\tfrac34\,\theta_S\beta\,2^{-s}.
\end{equation*}

Second, suppose $F$ was stored at a step $j<k$. Then
\begin{equation*}
f_{j,m}-f_{k,m}-\varrho_F\ge\bigl(\tfrac12-\tfrac18\bigr)\beta\,2^{-(j-m)}\ge\tfrac34\,\beta\,2^{-s}.
\end{equation*}
The last inequality holds because $j\le k-1$ gives $2^{-(j-m)}\ge 2\cdot2^{-s}$, and
$\tfrac38\cdot2=\tfrac34$. For a term stored at a step $j<k$ the table coefficient is
$\mathsf{c}_F^{+}=f_{j,m}$, and the schedule letter satisfies $\theta_S\le1$; so in this case, too,
$\mathsf{c}_F^{+}-\varrho_F\ge f_{k,m}+\tfrac34\,\theta_S\beta\,2^{-s}$.

In either case, inserting the lower bound into the bound on $\mathsf{T}^{\flat}_0$ and using
$\tfrac34-\tfrac9{16}=\tfrac3{16}$ gives \eqref{9.1:eq-untouched-accommodation-2}.
\end{proof}

\begin{definition}[A component excluded mid-step and its table]\label{9.1:def-excluded-component}
We work at the $\mathbf{R}$-site of step $k$. Here $Z$ denotes the processed components,
$\mathfrak{O}=\Omega_k\cup Z^c$ is the untouched region of Lemma~\ref{9.1:lem-untouched-region},
$\mathsf{T}_\nu$ are the stage tables, and $\mathsf{T}^{\flat}_\nu$ are the two-region tables of
Definition~\ref{9.1:def-two-region-tables}. Further, $N$ is the number of stages of the step,
$R_i$ is the collar of stage $i$ and $\mathsf{g}_i$ its load, and $\mathsf{u}_\ell$ is the unit at
level $\ell$, as in Lemma~\ref{9.1:lem-untouched-region} and
Lemma~\ref{9.1:lem-untouched-accommodation}. Finally, $m=m(x)\le k$ is the label of
Lemma~\ref{9.1:lem-untouched-region}, $f_{k,m}$ is the table coefficient at level $k$ and label
$m$, and $f_{k,m}\mathsf{u}_m$ is the leading term of the table on $\mathfrak{O}$ of
Definition~\ref{9.1:def-two-region-tables}. The letters $\theta_S$ and $\beta=\beta_I$ are as in
Lemma~\ref{9.1:lem-untouched-accommodation}.

Let $W$ be a component that leaves $Z$ after the first $n^+$ stages of step $k$. Thus $W$ is a
new large-field region, excluded from $Z$ in the course of step $k$ as in the construction of
\cite{B-CMP122a}. The triple $(W,k,n^+)$ is the arrest of $W$, and $(\ell,\varpi)$ and $F^{(k)}$
denote its frozen structure and its frozen factor at level $k$, in the sense of the definition of
arrests given later in this chapter. By the theorem on arrested components stated there, under
the hypotheses of that theorem, from the arrest until its dissolution every space of the
construction uses on $W$ the frozen structure $(\ell,\varpi)$ and the factor $F^{(k)}$.

On $W\cap\Omega_k^c$ we put
\begin{equation}\label{9.1:eq-excluded-component-1}
\mathsf{T}^{\flat}_\nu:=\mathsf{T}_\nu\quad(\nu\le n^+),\qquad
\mathsf{T}^{\flat}_\nu:=F^{(k)}\varpi\,\mathsf{u}_\ell+\sum_{\nu\le i<N}\mathsf{g}_i\quad(\nu>n^+).
\end{equation}
Here the loads $\mathsf{g}_i$ are those of the surviving collars $R_i$ with $i>n^+$ only, read in the frozen structure $(\ell,\varpi)$ of $W$. This table has the following properties.
\begin{itemize}
\item[(i)] Condition (G) holds for the loads $\mathsf{g}_i$, $i>n^+$, through
the frozen shells of $W$.
\item[(ii)] Parts (a) and (b) of Lemma~\ref{9.1:lem-untouched-accommodation} hold on $W$ at every stage
$\nu>n^+$, with $F^{(k)}\varpi\,\mathsf{u}_\ell$ in place of $f_{k,m}\mathsf{u}_m$.
\item[(iii)] At the reading (L) of the frozen chain terms $\mathbb{C}^{(i)}(X)$ with
$X\cap W\neq\emptyset$, the root term lies under $F^{(k)}\varpi\,\mathsf{u}_\ell$.
\item[(iv)] For readings at stages $\nu\le n^+<n_e$, where $n_e$ is the end stage of the moves as
in part (b) of Lemma~\ref{9.1:lem-untouched-accommodation}, Lemma~\ref{9.1:lem-site-blind-slice} is
applied with the tables
\begin{equation}\label{9.1:eq-excluded-component-2}
\mathsf{T}_e:=\mathsf{T}^{\flat}_{n_e},\qquad
\mathsf{T}_r:=\mathsf{T}^{\flat}_{n^{+}+1}+(\mathsf{T}_\nu-\mathsf{T}_{n^+}),
\end{equation}
and hypotheses (A) and (S) of that lemma hold for these tables.
\end{itemize}
With the table \eqref{9.1:eq-excluded-component-1}, the rows of type (P) of the theorem bounding
the $\mathbf{R}$-born output $\Phi$ (see Definition~\ref{9.1:def-two-region-tables}) read
$\mathfrak{O}$ as $\mathfrak{O}$ together with $W$ at the stages $\nu>n^+$; nothing else in that
bound changes in form.
\end{definition}

\begin{proof}[Proof of properties (i)--(iv) of Definition~\ref{9.1:def-excluded-component}]
We first explain why neither of the two regions of Definition~\ref{9.1:def-two-region-tables}
covers $W$, so that a separate table is needed on $W$; the proof of (i)--(iv) follows.

Suppose first that $Z$ denotes the surviving set. Then $W\subset\mathfrak{O}$. But the stages $\le n^+$ have re-levelled $W$, and the collars $R_i$ of the stages $i\le n^+$ meet $W$. Hence condition (G) and part (a) of Lemma~\ref{9.1:lem-untouched-accommodation} fail on $W$.

Suppose next that $Z$ denotes the initial class. Then $W$ is never raised to level $k$, and the
trichotomy of points by which the tables are read fails. Moreover, for $n^+<N-1$, the root term
fills the level $F^{(k)}-\tfrac18\varrho_\Phi$, where $\varrho_\Phi$ is the room of $\Phi$, and
this level exceeds the stage coefficient $\varphi^{(N)}$. So the obstruction to the bound for the
$\mathbf{R}$-born output $\Phi$ discussed with Definition~\ref{9.1:def-two-region-tables} recurs
on $W$.

Readings on $W$ do occur. Indeed, by the definition of the chain terms, the chain term $\mathbb{C}^{(n)}(X')$ of a surviving component with $n>n^+$ reads the frozen chain terms $\mathbb{C}^{(i)}(X)$ with $X\cap W\neq\emptyset$, and so does the reading (L) of property (iii). A table on $W$ is therefore required, and \eqref{9.1:eq-excluded-component-1} is that table.

\emph{Property (i).} By the geometry of Lemma~\ref{9.1:lem-untouched-region}, for $i>n^+$ the
collar $R_i$ is contained in a surviving component. This component is separated from $W$ by
$\Lambda_k$ and by the shells of $W$. Hence condition (G) holds for these loads $\mathsf{g}_i$
through the frozen shells of $W$.

\emph{Property (ii).} At a stage $\nu>n^+$ the table \eqref{9.1:eq-excluded-component-1} on $W$ has the form of the table on $\mathfrak{O}$: a leading term plus the tail $\sum_{\nu\le i<N}\mathsf{g}_i$ of the loads. The only difference is that the leading term $f_{k,m}\mathsf{u}_m$ is replaced by $F^{(k)}\varpi\,\mathsf{u}_\ell$. The proof of parts (a) and (b) of Lemma~\ref{9.1:lem-untouched-accommodation} applies unchanged with this replacement, and gives those parts on $W$ at $\nu>n^+$.

\emph{Property (iii).} By the theorem on arrested components given later in this chapter, from
the arrest to its dissolution every space of the construction uses on $W$ the frozen structure
$(\ell,\varpi)$ and the factor $F^{(k)}$. So the frozen structure is the structure of the space on
$W$, and the root term at the reading (L) lies under $F^{(k)}\varpi\,\mathsf{u}_\ell$.

\emph{Property (iv), hypothesis (A).} Both $n^{+}+1$ and $n_e$ exceed $n^+$, so both tables are
given by the case $\nu>n^+$ of \eqref{9.1:eq-excluded-component-1}. The leading terms
$F^{(k)}\varpi\,\mathsf{u}_\ell$ therefore cancel, and the room of
Lemma~\ref{9.1:lem-site-blind-slice} is
\begin{equation*}
\mathsf{G}=\mathsf{T}_r-\mathsf{T}_e
=(\mathsf{T}_\nu-\mathsf{T}_{n^+})+\sum_{n^{+}+1\le i<n_e}\mathsf{g}_i .
\end{equation*}
The first summand is the room of the stage tables over the stages $j$ with $\nu\le j<n^+$. For the moves at these stages, hypothesis (A) is supplied by the accommodation statement for the stage tables $\mathsf{T}_\nu$, by which these tables satisfy hypothesis (A) of Lemma~\ref{9.1:lem-site-blind-slice} for the moves at their own stages. The second summand is the room over the stages $i$ with $n^+<i<n_e$. For the moves at these stages, hypothesis (A) is supplied by part (b) of Lemma~\ref{9.1:lem-untouched-accommodation}, which holds on $W$ at $\nu>n^+$ by property (ii). Adding the two gives (A).

\emph{Property (iv), hypothesis (S).} Inserting the case $\nu>n^+$ of
\eqref{9.1:eq-excluded-component-1}, at $\nu=n^{+}+1$, into \eqref{9.1:eq-excluded-component-2}
gives
\begin{equation*}
\mathsf{T}_r=\mathsf{T}_\nu+\sum_{i>n^+}\mathsf{g}_i
-\bigl(\mathsf{T}_{n^+}-F^{(k)}\varpi\,\mathsf{u}_\ell\bigr).
\end{equation*}
The bracket $\mathsf{T}_{n^+}-F^{(k)}\varpi\,\mathsf{u}_\ell$ is non-negative by the lemma comparing the stage table $\mathsf{T}_{n^+}$ at the stage of arrest with the frozen factor, which gives $\mathsf{T}_{n^+}\ge F^{(k)}\varpi\,\mathsf{u}_\ell$ on $W$. Hence
\begin{equation*}
\mathsf{T}_r\le\mathsf{T}_\nu+\sum_{i>n^+}\mathsf{g}_i
\le\Bigl(f_{k,\ell}+\tfrac{1}{16}\theta_S\beta\,2^{-(k-\ell)}\Bigr)\varpi\,\mathsf{u}_\ell .
\end{equation*}
The first inequality drops the non-negative bracket. In the second, the loads are bounded by
part (a) of Lemma~\ref{9.1:lem-untouched-accommodation}, valid on $W$ by property (ii), which
gives the term $\tfrac{1}{16}\theta_S\beta\,2^{-(k-\ell)}\varpi\,\mathsf{u}_\ell$; its exponent is
the exponent $k-m$ of that part with the label $m$ replaced by $\ell$ as in property (ii).
The arithmetic of part (c) of Lemma~\ref{9.1:lem-untouched-accommodation} therefore applies to this bound and gives (S).

It remains to treat the points raised at the stages $j$ with $\nu\le j<n^+$. There the sum
$\sum_{i>n^+}\mathsf{g}_i$, expressed in the units $\mathsf{u}_F$ of the stored term $F$ of part
(c) of Lemma~\ref{9.1:lem-untouched-accommodation} (not the frozen factor $F^{(k)}$), is at most
$\tfrac{1}{12}\cdot\tfrac23\cdot\tfrac94=\tfrac18$ of the margin
$\tfrac34\theta_S\beta\,2^{-(k-n_F)}$ of the schedule bound for the stored terms $F$, where $n_F$
is the level of the stored term $F$. This uses $\ell_\circ^2\le L/3$. So at these points the loads
are absorbed by that margin, and (S) holds there as well.
\end{proof}

\begin{definition}[The substitution chain at the large-field site and its sources]
\label{9.1:def-substitution-chain}
Throughout, $k$ is the level and we work at the large-field site.

\emph{Rows and their order.} The \emph{substitution chain} at the large-field site is the sequence
of the rows $0$, $3$, $2$, $1$, $0'$ of the row decomposition of the large-field step; a row is
one of the substitutions of which that step is composed, and each row is carried out in the three
stages (a)--(c) below. The rows are taken in this order, which is the order in which the corresponding substitutions
are made in \cite{B-CMP122b}.

\emph{Structure of a row.} Each row has the structure fixed by the stage statement of the
large-field step; in the notation of that statement it reads as follows. It starts from a
critical point $P_0$ on a set $\mathfrak{S}_0$, and a set $\mathfrak{S}_1$ refining
$\mathfrak{S}_0$ is given. The parameter $\eta$, the quantities $\mathbb{H}$, $\mathbb{H}_W$,
$\mathbb{H}_2$ in the exponents, and the regions $W$ and $Y$ are those of the stage statement.
The stages of the row are:
\begin{itemize}
\item[(a)] $P_1=e^{i\eta\mathbb{H}}P_0$;
\item[(b)] $\mathfrak{U}=e^{i\eta\mathbb{H}_W}P_0$ on $W$;
\item[(c)] $\Xi=e^{i\eta\mathbb{H}_2}P_1$, the minimiser on
$\mathfrak{S}_2=\mathbf{B}(Y)\cup\mathfrak{S}_1$ with layer data $\mathrm{rd}\,\mathfrak{U}$, the
reading of $\mathfrak{U}$; here $\mathbf{B}(Y)$ is Ba{\l}aban's determining set of the region $Y$.
\end{itemize}
The gauge frame is carried unchanged through the three stages, in accordance with the
carried-frame convention: no axial gauge condition is imposed anew.

\emph{Links.} The links of the chain are of two kinds, called the first kind and the second kind.
Stages (a) and (c) end in objects of the second kind; stage (b), and the dilated or cut vertices,
are links of the first kind.

\emph{Sources.} The sources $b_i$ of the chain are the following:
\begin{itemize}
\item the sources of \cite[(1.58)]{B-CMP122b} and \cite[(1.64)]{B-CMP122b};
\item the sources of \cite[(1.57)]{B-CMP122b}, together with the cut labelled $1$;
\item the sources of side $1$ and of the cut;
\item the sources of \cite[(1.51)]{B-CMP122b}.
\end{itemize}
All of them are carried by sites of $\bar{\Lambda}$ \cite{B-CMP122b}.

\emph{Size of the sources.} Let $\alpha_{0,k}$ and $\alpha_{1,k}$ be Ba{\l}aban's radii at level
$k$, and let $\mathsf{r}_1=C_1/C_0$ and the frame constant $K_\Lambda'(M)$, which depends on $M$,
be as in Notation~\ref{9.1:not-site-blind-constants}. We put
\begin{equation}\label{9.1:eq-substitution-chain-1}
\mathsf{s}_L:=K_\Lambda'(M)\bigl(\alpha_{0,k}+\alpha_{1,k}\bigr)
=K_\Lambda'(M)\,(1+\mathsf{r}_1)\,\alpha_{0,k},
\end{equation}
the second expression being the first rewritten with $\alpha_{1,k}=\mathsf{r}_1\alpha_{0,k}$. By
the bound on source sizes at seeds, one has at seeds, on the real supports of the sources,
\begin{equation}\label{9.1:eq-substitution-chain-2}
\sup|b_i|\le\mathsf{s}_L\qquad\text{for each of the sources } b_i \text{ listed above.}
\end{equation}

\emph{Typing of the chain.} At the large-field site the entry of the chain is the root $q_0$.
It is a seed, and its real centre $U_c$ is never changed along the chain, by the statement that
the real centre of a seed is held fixed. Every link of the chain, whether of the first kind or of
the second kind, is an increment.
\end{definition}

\begin{lemma}[Link-by-link bound over the root seed]\label{9.1:lem-link-bound}
Assume the classical condition in profile form (Definition~\ref{9.1:def-profile-condition}), the
hypothesis on the set $\bar{\Lambda}$, and the class statement at the
ambient scale. Let $q_0$ be a root seed with complex data, $\sigma'$ a vertex, and $q$ an object
of the substitution chain of Definition~\ref{9.1:def-substitution-chain}. Then there is a field
$\mathcal{Z}_q$ such that
\begin{equation}\label{9.1:eq-link-bound-1}
q=e^{i\eta\mathcal{Z}_q}q_0,\qquad
\ell_\iota|\mathcal{Z}_q|\le\mathsf{m}_L:=C_4\,\mathsf{s}_L\,e^{-\delta_1 d(\cdot,\bar{\Lambda})/4},
\end{equation}
\begin{equation}\label{9.1:eq-link-bound-2}
C_4:=16\,(1+68\,d\,C_sC_T^w)\max\{C_s,B^{\mathsf{s}}\}\,c_1 .
\end{equation}
Here $d(\cdot,\bar{\Lambda})$ is the distance to the set $\bar{\Lambda}$ attached to the root,
which is the least of the
distances occurring in the chain; as in Lemma~\ref{9.1:lem-top-source-decay}, the letter $d$
with two arguments denotes a distance, while $d$ standing alone, as in $68\,d$, is the
dimension $d=4$. The unit $\ell_\iota$ is the field unit at the level $\iota$ of the point at
which the bound is read. The constant $C_s$ is the constant
of the regularity lemma for increments, which serves the natural vertices. At the dilated or cut
link vertices of substitutions (a) and (b) the profile constant is $B^{\mathsf{s}}$, the constant
of the
size clause of Definition~\ref{9.1:def-profile-condition}; in substitution (c) the constant is
$C_s$.
Further, $c_1$
is the constant of Lemma~\ref{9.1:lem-top-source-decay}(i), and $C_T^w$ is the constant of the
transfer bound of the random-walk chain. The field $\mathcal{Z}_q$ is holomorphic in all
parameters, and $\mathcal{Z}_q=0$ at vanishing sources.
\end{lemma}

\begin{proof}
The argument has two parts. We first bound the increment contributed by one member of the chain,
and then compose the members.

\emph{Substitutions (a) and (b).} The increment of substitution (a) is
$\mathbb{H}=\mathfrak{X}_{\mathfrak{S}_1}(b;P_0,\mathsf{m})$, the increment produced by the
coordinate map $\mathfrak{X}$ on the determining set $\mathfrak{S}_1$ from the source datum $b$
at the base $P_0$, read with a profile $\mathsf{m}$ in the sense of
Lemma~\ref{9.1:lem-top-source-decay}. The increment $\mathbb{H}_W$ of substitution (b) is the
same construction on $\boldsymbol{\Omega}(W)\vee\mathfrak{S}_1$. We apply the
regularity lemma for increments with complex source datum $b$. The increment is analytic in the
base $P_0$; the regularity lemma is used together with the tilted natural-interpolation lemma. The
lemma bounds the increment, in the units of its own level, by $C_s$ times the majorant
$(|b|)^{\sim}$ of Lemma~\ref{9.1:lem-top-source-decay}(i). By that lemma this majorant satisfies
\begin{equation*}
(|b|)^{\sim}\le c_1\sup|b|\,e^{-\delta_1 d(\cdot,\operatorname{supp}b)/4}.
\end{equation*}
The distance attached to the root is the least one, so that
$d(\cdot,\operatorname{supp}b)\ge d(\cdot,\bar{\Lambda})$ for the source data $b$ of the chain.
The decay factor is therefore at most
$e^{-\delta_1 d(\cdot,\bar{\Lambda})/4}$, and at natural vertices
\begin{equation}\label{9.1:eq-link-bound-3}
\ell_\iota|\mathbb{H}|,\ \ell_\iota|\mathbb{H}_W|\ \le\ \mathsf{m}'
:=C_s\,c_1\sup|b|\,e^{-\delta_1 d(\cdot,\bar{\Lambda})/4}
\end{equation}
in the units of the increment's own level.

By the stage lemma of the random-walk chain, the increment of substitution (b) is damped at
the rim of $W$. Consequently the bound of \eqref{9.1:eq-link-bound-3} for $\mathbb{H}_W$ holds
also in the ambient units at the rim of $W$.

At dilated or cut vertices of substitutions (a) and (b) we use instead the size clause of
Definition~\ref{9.1:def-profile-condition}. At these vertices the profile constant is
$B^{\mathsf{s}}$, and the clause gives the bound of \eqref{9.1:eq-link-bound-3} with $C_s$
replaced by $B^{\mathsf{s}}$. In every
case the increment of substitutions (a) and (b) is at most
\begin{equation}\label{9.1:eq-link-bound-4}
\max\{C_s,B^{\mathsf{s}}\}\,c_1\sup|b|\,e^{-\delta_1 d(\cdot,\bar{\Lambda})/4}.
\end{equation}

\emph{Substitution (c).} By the criticality clause of the nesting lemma, applied to complex data,
the
restriction $P_1{\upharpoonright}Y$ is critical on $\mathfrak{S}_2$. Let $C$ be a cube of the
layer. The layer-source lemma bounds the layer datum $b_{\mathrm{lay}}$ of substitution (c), which
is read
from $\mathfrak{U}$. Here $p$ and $m$ denote the level of the layer cube $C$ and the level of
substitution (c), $\ell_p$ and $\ell_m$ the field units at these levels, and $P(C)$ the set of
points over
which the layer datum at $C$ is read. Together with the bounds of substitutions (a) and (b) for
$\mathbb{H}$ and $\mathbb{H}_W$, the layer-source lemma gives
\begin{equation}\label{9.1:eq-link-bound-5}
|b_{\mathrm{lay}}(C)|\le 68\,d\,\ell_p\sup_{P(C)}\bigl(|\mathbb{H}_W|+|\mathbb{H}|\bigr)
\le 136\,d\,L^{p-m}\,\mathsf{m}'(C).
\end{equation}
The factor $136=2\cdot 68$ accounts for the two increments $\mathbb{H}_W$ and $\mathbb{H}$.

Now apply the regularity lemma on $\mathfrak{S}_2$ with the datum \eqref{9.1:eq-link-bound-5},
together with the transfer bound of the random-walk chain at the value $a=1$ of its parameter.
This
gives
\begin{equation}\label{9.1:eq-link-bound-6}
\ell_m|\mathbb{H}_2|\le 136\,d\,C_sC_T^w\,\mathsf{m}'.
\end{equation}
Since $C_s\le\max\{C_s,B^{\mathsf{s}}\}$, the right side of \eqref{9.1:eq-link-bound-6} is at most
\begin{equation}\label{9.1:eq-link-bound-7}
2\cdot 68\,d\,C_sC_T^w\max\{C_s,B^{\mathsf{s}}\}\,c_1\sup|b|\,e^{-\delta_1 d(\cdot,\bar{\Lambda})/4}.
\end{equation}

\emph{Composition.} Definition~\ref{9.1:def-substitution-chain} lists the members of the chain in
the order $0,3,2,1,0'$, which is Ba{\l}aban's order of substitution. Member $0$ is the root $q_0$
itself and contributes no increment. Each of the members $3,2,1,0'$ contributes one increment, of
one of the kinds bounded above.

Increments at cut vertices are tame by the tameness clause of the statement on cut increments at
the small-field site. The raw core, in the sense of the clause on the raw core of
the same statement, lies beyond the seam fixed there, by member $1$ of the chain and that clause.

By \eqref{9.1:eq-link-bound-4} and \eqref{9.1:eq-link-bound-7}, and since
$1\le 2(1+68\,d\,C_sC_T^w)$, each of the four increments is
bounded in the units $\ell_\iota$ by
\begin{equation*}
2\,(1+68\,d\,C_sC_T^w)\max\{C_s,B^{\mathsf{s}}\}\,c_1\sup|b|\,e^{-\delta_1 d(\cdot,\bar{\Lambda})/4}.
\end{equation*}
Their sum is at most four times this common bound, that is, at most
\begin{equation*}
8\,(1+68\,d\,C_sC_T^w)\max\{C_s,B^{\mathsf{s}}\}\,c_1\sup|b|\,e^{-\delta_1 d(\cdot,\bar{\Lambda})/4}.
\end{equation*}
Combining the four exponentials into the single
exponential $e^{i\eta\mathcal{Z}_q}$ by the Baker--Campbell--Hausdorff composition costs a further
factor $2$, which gives the factor $16$ of $C_4$. For the source data of the chain one has
$\sup|b|\le\mathsf{s}_L$, where
$\mathsf{s}_L$ is the seed size of the root. Hence
$q=e^{i\eta\mathcal{Z}_q}q_0$ with
\begin{equation*}
\ell_\iota|\mathcal{Z}_q|\le C_4\,\mathsf{s}_L\,e^{-\delta_1 d(\cdot,\bar{\Lambda})/4}=\mathsf{m}_L,
\end{equation*}
with $C_4$ as in \eqref{9.1:eq-link-bound-2}. This is \eqref{9.1:eq-link-bound-1}.

\emph{Holomorphy and vanishing.} The regularity lemma is applied with complex source data and
yields increments analytic in the base, and the criticality clause of the nesting lemma is
applied with complex data; each increment is holomorphic in all parameters, and the
Baker--Campbell--Hausdorff composition of holomorphic fields is holomorphic. When the source
data $b$ of the chain vanish, each increment
vanishes by \eqref{9.1:eq-link-bound-3}, \eqref{9.1:eq-link-bound-4} and
\eqref{9.1:eq-link-bound-6}, hence $\mathcal{Z}_q=0$.

Finally, the re-minimisations
\cite[(1.61)--(1.63)]{B-CMP122b} introduce no additional error here. The error term of order
$\alpha_0$ in the re-minimisation comparison statement, where $\alpha_0$ is one of
Ba{\l}aban's radii $\alpha$ among the auxiliary parameters $\mathfrak{p}$,
concerns generic points of the tube on which that statement is made. The root is a slice point,
and every renewed minimisation on a refined determining set returns it exactly, since the
determining sets along the chain only refine (the
refinement property of the random-walk chain).
\end{proof}

\begin{definition}[Coupling ceilings and the graded-region condition]\label{9.1:def-graded-ceilings}
Let $q$ range over the chain objects of the substitution chain of
Definition~\ref{9.1:def-substitution-chain}, presented over the root seed as in
Lemma~\ref{9.1:lem-link-bound}, and let $C_4$ be the constant of that lemma. Let $N_0$ be the
integer defined by $L^{N_0-1}=\check{R}_{k_0+1}$. Here $k_0$ is the level fixed with the large-field
site of the chain, and $\check{R}_{k_0+1}$ is the scaffold radius at level $k_0+1$, written
$R_{k_0+1}$ in Ba{\l}aban's construction; it is not the ball radius $R$ of
Notation~\ref{0.2:not-glossary}. The remaining letters are the following:
\begin{itemize}
\item $B_{\natural}$ is the constant of the local Lipschitz property of the Landau chart;
\item $K_\Lambda'=K_\Lambda'(M)$ is the frame constant in terms of which the seed size
$\mathsf{s}_L$ of the layer profile of Lemma~\ref{9.1:lem-link-bound} is given;
\item $\mathsf{r}_1=C_1/C_0$;
\item $\theta_\rho$ is the parameter of the schedule of the room;
\item $c_1$ and $C_T^{w}$ are the constants entering the definition of $C_4$ in
Lemma~\ref{9.1:lem-link-bound}, and $\delta_1$ is the decay rate of the layer profile of that
lemma;
\item $M=L^{m'}$ is Ba{\l}aban's auxiliary parameter of Notation~\ref{0.2:not-glossary};
\item $\beta$ is the fixed small number of the stage tables, the one entering their coefficient
$1-4\beta$; it is not the coupling increment $\beta_{k+1}$.
\end{itemize}

\emph{Ceiling conditions.} The \emph{first} and \emph{second ceiling conditions}, together called
the \emph{coupling ceilings}, are the inequalities
\begin{align}
&2^{15}B_{\natural}\,c_1\,C_4\,K_\Lambda'\,(1+\mathsf{r}_1)\,
e^{-\delta_1 M(N_0-2)/2}\le\theta_\rho\,\beta,
\label{9.1:eq-graded-ceilings-1}\\
&2^{16}B_{\natural}\,c_1\,C_T^{w}\,C_4\,K_\Lambda'\,(1+\mathsf{r}_1)\,
4^{-(N_0-2)}\le 1 .
\label{9.1:eq-graded-ceilings-2}
\end{align}
These inequalities are conditions on $\gamma$, the ceiling of the running couplings $g_k$.
By the definition of $N_0$ and the lower bound on the scaffold radii of the prescribed scaffold,
\begin{equation}\label{9.1:eq-graded-ceilings-3}
L^{N_0-1}=\check{R}_{k_0+1}\ge(\log g^{-2})^{r},
\end{equation}
where $r$ denotes the fixed exponent of that lower bound; this $r$ is distinct from the depth
letter $r$ of this chapter.
By \eqref{9.1:eq-graded-ceilings-3}, the left members of \eqref{9.1:eq-graded-ceilings-1} and
\eqref{9.1:eq-graded-ceilings-2} are bounded by constants independent of the couplings times
negative powers of $\log g_k^{-2}$. For the second this is
the computation
\begin{equation*}
4^{-(N_0-2)}=4\,\bigl(L^{N_0-1}\bigr)^{-\log 4/\log L}\le 4\,(\log g^{-2})^{-r\log 4/\log L},
\end{equation*}
the inequality being \eqref{9.1:eq-graded-ceilings-3}; the first is treated in the same way, since
\begin{equation*}
e^{-\delta_1 M(N_0-2)/2}=e^{\delta_1 M/2}\bigl(L^{N_0-1}\bigr)^{-\delta_1 M/(2\log L)}
\le e^{\delta_1 M/2}(\log g^{-2})^{-r\delta_1 M/(2\log L)} .
\end{equation*}
Both conditions are therefore met at every level $k$ with $g_k\le\gamma$ once $\gamma$ is small
enough.

\emph{The table condition off the core.} The \emph{table condition off the core} is a bound on a
quantity of a
single construction, not on a difference between the two runs $\mathsf{A}$, $\mathsf{B}$ of
lengths $n$ and $n+1$ compared in this chapter. It concerns the regions $X$ of the
terms of the expansion that do not meet the raw core of the large-field site of the chain; here $X$
denotes such a region and not the inverse squared coupling. It is the following statement:
for every such $X$, every chain object $q$ and every table member $a$, one has on $Z''_k\cap X$,
where $Z''_k$ is the level-$k$ large-field region of Ba{\l}aban's construction
entering the present step,
\begin{equation*}
E_a\bigl(S(q)\bigr)\le\tfrac14\,\mathsf{T}^a_N .
\end{equation*}
Here $E_a$ is the $a$-th member of the stage tables of room increments, $\mathsf{T}^a_N$ is its
threshold in the table of stage $N$, the letter $N$ being here an index of the stage tables and
not the $N$ of $\mathrm{SU}(N)$, and $S(q)$ is the slice of $q$.

Two statements used elsewhere in this chapter bear on this condition. The
lemma on the growth of the stage-table thresholds states that the thresholds $\mathsf{T}^a_N$ are
at least a factor comparable to $\ell_\circ^{2(N_0-1)}$ times the corresponding level-$k$ sizes,
this factor tending to infinity as $N_0\to\infty$; here $\ell_\circ$ is the integer that
\cite{B-CMP122a} denotes $L_0$, which is not the constant $L_0=L_0(N)$ of hypothesis (H2) in
Notation~\ref{2.2:not-hypotheses}. The theorem bounding the stage-table members bounds the members
it treats by one sixteenth of their thresholds.

Regions $X$ that meet the raw core are excluded from the condition. On such a region the
object $q$ leaves the chart of the root seed, so its slice $S(q)$, which is taken in that chart,
is undefined there, and with it the left member of the condition. This definition names the table
condition off the core and does not assert that it holds. It is replaced, in the estimates of
this chapter, by the graded-region condition defined later in the chapter.
\end{definition}

\begin{lemma}[Slices of chain objects on the three regions]\label{9.1:lem-chain-slices}
Let $\Phi$ be the output of the $\mathbf{R}$-operation at step $k$, with its seeds in the seed
tube $\mathfrak{T}_\Phi[7/8]$. Let $q$ be a chain object, with root seed $q_0$ and chain field
$\mathcal{Z}_q$, so that $q=e^{i\eta\mathcal{Z}_q}q_0$ as in Lemma~\ref{9.1:lem-link-bound}.
Write $E_a$ for the members of the table and $S(\cdot)$ for the slice. Let $\mathfrak{O}$ be the
untouched region of Lemma~\ref{9.1:lem-untouched-region}, and let $\mathcal{A}$ and
$Z''^{(1)}_k$ be the two further regions of the $\mathbf{R}$-site at step $k$.
\begin{enumerate}
\item[(i)] On $\mathfrak{O}$, at every point of label level $m$,
\begin{equation}\label{9.1:eq-chain-slices-1}
E_a(S(q))\le\Bigl(f_{k,m}-\tfrac{15}{128}\varrho_\Phi\Bigr)\mathsf{u}^a_m ,
\end{equation}
where $f_{k,m}$ denotes the coefficient of the step-$k$ reading at label level $m$.
\item[(ii)] On $\mathcal{A}$,
\begin{equation}\label{9.1:eq-chain-slices-2}
E_a(q_0)\le\min\Bigl\{\tfrac{5}{12}\,\mathsf{T}^a_N,\ \tfrac16\,\mathsf{u}^a_F\Bigr\},
\qquad
E_a(S(q))\le E_a(q_0)+\tfrac{1}{960}\,\mathsf{u}^a_F .
\end{equation}
Here $\mathsf{u}_F$ are the own units of the term $F$ whose slice is read. At $(=)$-points of
a stage-$\nu$ reading, the last term is read as $\mathsf{u}_F/960\le\mathrm{rm}_{i_1}/24$,
where $\mathrm{rm}_{i_1}$ is the first raising room.
\item[(iii)] On $Z''^{(1)}_k$ the slice $S(q)$ obeys the refined slice bound on
$Z''^{(1)}_k$: at aligned points $E_a(S(q))\le\tfrac{3}{32}\mathsf{T}^a_N$, and at older
points $E_a(S(q))\le\tfrac{3}{32}\mathsf{u}^a_F$.
\end{enumerate}
\end{lemma}

\begin{proof}
Let $(F,X,j)$ be the leaf carrying the term $F$ whose slice is read, let
$\mathfrak{S}=\boldsymbol{\Omega}(X^+)$ be its box set, and let $\eta$, $\ell_\iota$ and $e_a$ be
as in Lemma~\ref{9.1:lem-site-blind-slice}.
We use three inputs throughout.
\begin{itemize}
\item Part (a) of the slice-displacement statement, applied at the root seed $q_0$. Its
hypothesis holds because the set $\mathfrak{B}$ on which part (a) is applied refines the
seed's set $\boldsymbol{\Omega}(R)$ of the root seed.
\item Part (b) of the same statement, applied with move field $\mathcal{Y}:=\mathcal{Z}_q$. This
bounds the members of $S(q)-S(q_0)$ through the smeared majorant
$(|\mathcal{Z}_q|_\cdot)^{\approx}$.
\item Part (i) of Lemma~\ref{9.1:lem-top-source-decay}, applied to the profile
$\mathsf{m}_L$ of Lemma~\ref{9.1:lem-link-bound}.
\end{itemize}

\emph{The region $\mathfrak{O}$.} On $\mathfrak{O}$ no stage changes the label
(Lemma~\ref{9.1:lem-untouched-region}); we argue at $(=)$-points.
By Lemma~\ref{9.1:lem-link-bound},
\begin{equation*}
|\mathcal{Z}_q|_y\le\frac{\ell_{\mathfrak{S}}}{\ell_\iota}(y)\,\mathsf{m}_L(y)\le\mathsf{m}_L(y)
\end{equation*}
for every point $y$. The second inequality holds because $\lambda_{\mathfrak{S}}=j\wedge\iota\le\iota$,
where $\lambda_{\mathfrak{S}}$ is the level attached to $\mathfrak{S}$, $j$ is the level of the
leaf and $\iota$ is the level index of $\ell_\iota$.
By Lemma~\ref{9.1:lem-top-source-decay}\,(i) it follows that
$(|\mathcal{Z}_q|_\cdot)^{\approx}\le c_1\mathsf{m}_L$. The displacement given by part (b) is
therefore at most $16B_\natural c_1\ell_m^{-e_a}\mathsf{m}_L(x)$. This bound does not contain
$C_T^w$, exactly as in the $(=)$-case of Lemma~\ref{9.1:lem-site-blind-slice}.

We now combine three facts: the geometry of the untouched region $\mathfrak{O}$ and of the middle
shells given in Lemma~\ref{9.1:lem-untouched-region}, the
bound of Lemma~\ref{9.1:lem-top-source-decay}\,(ii) for $\mathsf{m}_L(x)/\alpha_{0,m}$ (with $s$
as there), and the first ceiling of Definition~\ref{9.1:def-graded-ceilings},
\begin{equation*}
2^{15}B_\natural c_1C_4K_\Lambda'(1+\mathsf{r}_1)e^{-\delta_1M(N_0-2)/2}\le\theta_\rho\beta .
\end{equation*}
Lemma~\ref{9.1:lem-top-source-decay}\,(ii) contributes the factor $\tfrac92\,2^{-s}$; the
constants $B_\natural c_1$ of the displacement bound and $C_4K_\Lambda'(1+\mathsf{r}_1)$ of the
profile $\mathsf{m}_L$ are those of the first ceiling; and the geometric input on
$\mathfrak{O}$ is the one that accounts for the factor $e^{-\delta_1M(N_0-2)/2}$ of that
ceiling. Together they bound the displacement by
\begin{equation}\label{9.1:eq-chain-slices-3}
\tfrac{1}{128}\,\theta_\rho\beta\,2^{-s}\alpha_{0,m}\ell_m^{-e_a},
\end{equation}
since $16\cdot\frac92\cdot128=9216\le2^{15}$. No $C_T^w$ enters here. Adding
\eqref{9.1:eq-chain-slices-3} to the reading of $q_0$ given by part (a) yields
\eqref{9.1:eq-chain-slices-1}.

\emph{The region $\mathcal{A}$: the bound on $E_a(q_0)$.} On $Z$ the new label has level $k$,
and $\mathsf{c}_\Phi^+=1$. Hence
\begin{equation*}
E_a(q_0)\le\alpha_{a,k}\ell_k^{-e_a}.
\end{equation*}
On the strip $S_m$ of label level $m$ the stage-$N$ table is
\begin{equation*}
\mathsf{T}^a_N=\mathsf{c}\,\ell_\circ^{2(k-m)}\alpha_{a,k}\ell_k^{-e_a}.
\end{equation*}
The coefficient $\mathsf{c}$ of the stage-$N$ table satisfies $\mathsf{c}\ge1-2\beta=3/5$ by the
lower bound on that coefficient, and $\ell_\circ^2\ge4$. Since $\frac35\cdot4=\frac{12}{5}$, this
gives $E_a(q_0)\le\frac{5}{12}\mathsf{T}^a_N$. The other member of the minimum,
$E_a(q_0)\le\mathsf{u}^a_F/6$, follows from the size condition on the seeds of the tube
$\mathfrak{T}_\Phi[7/8]$ together with the normalisation of the own units $\mathsf{u}_F$ of
$F$, both used as stated.

\emph{The region $\mathcal{A}$: the displacement.} An increment of level $k$ that is re-solved
at a finer level is not small in the units of $\mathsf{T}_N$. This is the situation of the
members $E_7$ and $E_3$ in the accounting of the table members, and it is why the
second bound in \eqref{9.1:eq-chain-slices-2} is stated in the own units of $F$.

On $\mathcal{A}$, and only there, the constant $C_T^w$ enters the displacement bound; the
second ceiling of Definition~\ref{9.1:def-graded-ceilings} carries it. Measured in the units
of $F$, the displacement gains the factor
\begin{equation*}
\Bigl(\frac{\ell_{n_F}}{\ell_{\mathfrak{S}}}\Bigr)^{e_a}\frac{\ell_{\mathfrak{S}}}{\ell_k}
\le L^{-u},\qquad u:=k-n_F\ge k-m ,
\end{equation*}
where $n_F$ is the level of $F$.
We use $(9/4)u^{1/2}L^{-u}\le(9/4)4^{-u}$. Put $t:=(m-k_0-2)_+$, with $k_0$ the level of
Definition~\ref{9.1:def-graded-ceilings}; then $d(x,\bar\Lambda)\ge2Mt$, where $\bar\Lambda$
is the set from whose distance the profile $\mathsf{m}_L$ of Lemma~\ref{9.1:lem-link-bound}
decays, so the decay factor of $\mathsf{m}_L$ at $x$ is at most
$e^{-\delta_1Mt/2}$. Because $\delta_1M\ge4\kappa_1$, we have
\begin{equation*}
e^{-\delta_1Mt/2}\,4^{-(N_0-2-t)}\le4^{-(N_0-2)} .
\end{equation*}
The constants collected in this way are those of the second ceiling of
Definition~\ref{9.1:def-graded-ceilings},
\begin{equation*}
2^{16}B_\natural c_1C_T^wC_4K_\Lambda'(1+\mathsf{r}_1)4^{-(N_0-2)}\le1 .
\end{equation*}
Since $16\cdot\frac94\cdot960=34560\le2^{16}$, the displacement is at most $\mathsf{u}^a_F/960$.
Adding it to the reading of $q_0$ gives the second bound in \eqref{9.1:eq-chain-slices-2}.

\emph{Reading at $(=)$-points.} By the lower bound on the raising rooms, the first raising
room obeys $\mathrm{rm}_{i_1}\ge\bar u_{i_1}/40$; this is the input for the reading
$\mathsf{u}_F/960\le\mathrm{rm}_{i_1}/24$ stated in (ii), where $960=24\cdot40$.

\emph{The region $Z''^{(1)}_k$.} Assertion (iii) is the refined slice bound on
$Z''^{(1)}_k$, restricted to chain objects $q$; it is used here as stated.
\end{proof}

\begin{definition}[Analyticity circles at the large-field site]\label{9.1:def-large-field-circles}
At the large-field site, a bracket $p$ (in this definition $p$ ranges over brackets, not
plaquettes) of the substitution chain of
Definition~\ref{9.1:def-substitution-chain}, localised by $\zeta_\square$ to the cube
$\square$, shifts the configuration by
$\mathsf{t}\,\zeta_\square\,\mathbb{H}^{(i)}(\mathsf{s})$, with $\zeta_\square$ and
$\mathbb{H}^{(i)}(\mathsf{s})$ as in that definition. Here $\mathsf{t}$ is the bracket
parameter and $\mathsf{s}$ the dilation. Its \emph{unit majorant} on $\tilde\square$ is
\begin{equation}\label{9.1:eq-large-field-circles-1}
  \mathsf{m}_p:=B^{\mathsf{s}}c_1\sup|b_i|\,e^{-\delta_1 d(\cdot,\bar\Lambda)/4},
\end{equation}
where the $b_i$ are the sources of the chain and $d(\cdot,\bar\Lambda)$ is the distance to
$\bar\Lambda$, the small-field region of the step, a union of cells of label level at
least $k$. Hence the function $\mathsf{m}_p$ of \eqref{9.1:eq-large-field-circles-1} is a
profile in the
sense of Lemma~\ref{9.1:lem-top-source-decay}, with $\mathcal{S}=\bar\Lambda$,
$k_{\mathcal{S}}=k$ and $\mathsf{m}^0=B^{\mathsf{s}}c_1\sup|b_i|$.

The radius is fixed by the room, in the sense of \cite[(3.15)]{B-CMP109} and of the room
convention of this chapter at the large-field site. The index
$(\mathrm{L})$ marks the large-field site. The \emph{circle radius} of the bracket is
\begin{equation}\label{9.1:eq-large-field-circles-2}
  \mathsf{w}_{(\mathrm{L})}(p):=\sup\bigl\{\varpi\in\mathbb{R}:\
  2^{8}B_\natural C_T^{\mathrm{w}}c_1(1+\varpi)\,\mathsf{m}_p\le\mathsf{r}^{(\mathrm{L})}
  \text{ on }\tilde\square\cap X^{+}\bigr\}.
\end{equation}
The room $\mathsf{r}^{(\mathrm{L})}$ is prescribed on three regions of the site: the region
$\mathfrak{O}$, on which it is graded by the label level $m$ of the point; the region
$\mathfrak{A}$, on which it is fixed at level $k$; and the region $Z''^{(1)}_k$ attached to
the region $Z''_k$ at level $k$. Namely,
\begin{equation*}
  \mathsf{r}^{(\mathrm{L})}=
  \begin{cases}
    \tfrac1{16}\,\theta_\rho\beta\,2^{-(k-m)}\alpha_{0,m} & \text{on }\mathfrak{O},\\[2pt]
    \alpha_{0,k}/400 & \text{on }\mathfrak{A};
  \end{cases}
\end{equation*}
on $Z''^{(1)}_k$, at the points of the first class it is $\tfrac14\mathsf{T}^{a}_N$, one
quarter of the threshold of the stage table of Definition~\ref{9.1:def-graded-ceilings}
for the table member $a$ read at the point, and at the points of the second class it is
measured in the own units $\mathsf{u}^a_F$ of the term $F$ being bracketed.
On $Z''^{(1)}_k$ the circle of the sourced run of the correlation theorem on the side of
the uniqueness theorem is divided by the constant $K_o^{(\mathrm{L})}$ of
\eqref{9.1:eq-large-field-circles-3} below, as in the room statement at the large-field
site. We put
\begin{equation}\label{9.1:eq-large-field-circles-3}
  K_o^{(\mathrm{L})}:=\frac{2^{17}B_\natural C_T^{\mathrm{w}}c_1^{2}B^{\mathsf{s}}}{\theta_\rho\beta}.
\end{equation}
We call a constant $K$ a \emph{circle constant} of the bracket if
$1+\mathsf{w}_{(\mathrm{L})}(p)\ge[K\sup|b_i|/\alpha_{0,k}]^{-1}$. For cubes
$\square\subset\mathfrak{O}$ the
circle constant of the bracket, as it enters the sourced run, is
$K_o^{(\mathrm{L})}e^{\delta_1M(N_0-2)/2}$, that is, $K_o^{(\mathrm{L})}$ with the further
factor $e^{\delta_1M(N_0-2)/2}$; for the other cubes it is $K_o^{(\mathrm{L})}$. For every
cube $\square$ with $\tilde\square\cap X^{+}\subset\mathfrak{O}\cup\mathfrak{A}$ the
following bound holds with the constant \eqref{9.1:eq-large-field-circles-3}; it is proved
below:
\begin{equation}\label{9.1:eq-large-field-circles-4}
  1+\mathsf{w}_{(\mathrm{L})}(p)\ge\Bigl[\frac{K_o^{(\mathrm{L})}\sup|b_i|}{\alpha_{0,k}}\Bigr]^{-1}.
\end{equation}

Inside $Z''_k$ the source of a bracket is that of the last substitution of the chain, of
size $C_1\delta'_k$; the brackets of the three intermediate substitutions have
$X\subset Z^c$, that is, their localisation set $X$ lies in the complement of the
large-field region $Z$, by the localisation statement for these substitutions. The size
$C_1\delta'_k$ is the one fixed in the re-reading of the last substitution, and the
arithmetic of that re-reading gives
\begin{equation}\label{9.1:eq-large-field-circles-5}
  1+\mathsf{w}_{(\mathrm{L})}(p)\ge\Bigl[\frac{K_o^{(\mathrm{L})}\cdot\tfrac94 C_1\delta'_kN^{1/2}}
  {\alpha_{0,k}}\Bigr]^{-1}.
\end{equation}
This is a circle of the same
degree in $\log g_k^{-2}$ as the coupling condition on the re-read small-field region.

The large-field entries of the sourced run of the correlation theorem on the side of
the uniqueness theorem carry the following factors. The last substitution, where
$\sup|b_i|\le C_1\delta'_k$, carries the factor
$K_o^{(\mathrm{L})}\theta'$, with $\theta'$ one of the index numbers fixed before $L$. The
three intermediate substitutions, whose brackets have
$X\subset Z^c$, carry the factor
\begin{equation*}
  K_o^{(\mathrm{L})}K_\Lambda'(1+\mathsf{r}_1)e^{-\delta_1M(N_0-2)/2}.
\end{equation*}
The conditions that these two factors be small are coupling ceilings in the sense of
Definition~\ref{9.1:def-graded-ceilings}. In
the room statement at the large-field site, the quotient $\mathsf{w}_{(\mathrm{L})}/K_o$ is
read with
$\mathsf{w}_{(\mathrm{L})}=\mathsf{w}_{(\mathrm{L})}(p)$ of
\eqref{9.1:eq-large-field-circles-2} and
$K_o=K_o^{(\mathrm{L})}$ of \eqref{9.1:eq-large-field-circles-3}. These circles are used
only in that sourced run; the correlation theorem is used as stated.
\end{definition}

\begin{proof}
We prove \eqref{9.1:eq-large-field-circles-4} for cubes $\square$ with
$\tilde\square\cap X^{+}\subset\mathfrak{O}\cup\mathfrak{A}$.

First suppose $\sup|b_i|=0$. Then $\mathsf{m}_p=0$ and every $\varpi$ satisfies the
condition in \eqref{9.1:eq-large-field-circles-2}. Hence $\mathsf{w}_{(\mathrm{L})}(p)=+\infty$ and
\eqref{9.1:eq-large-field-circles-4} holds, its right side being read as $+\infty$.

Now let $\sup|b_i|>0$ and put
\begin{equation*}
  \varpi_0:=\frac{\alpha_{0,k}}{K_o^{(\mathrm{L})}\sup|b_i|}-1.
\end{equation*}
We show that $\varpi_0$ satisfies the condition in
\eqref{9.1:eq-large-field-circles-2} for every cube $\square$ with
$\tilde\square\cap X^{+}\subset\mathfrak{O}\cup\mathfrak{A}$. Then
$\mathsf{w}_{(\mathrm{L})}(p)\ge\varpi_0$, which is \eqref{9.1:eq-large-field-circles-4}.
Note that $1+\varpi_0>0$.

Take first a point $x\in\tilde\square\cap X^+$ in $\mathfrak{O}$, of label level
$m\le k$. The cells of $\bar\Lambda$ have label level at least $k$, so clause
(ii) of Lemma~\ref{9.1:lem-top-source-decay} applies to the function $\mathsf{m}_p$ of
\eqref{9.1:eq-large-field-circles-1}, a profile with $\mathcal{S}=\bar\Lambda$ and
$k_{\mathcal{S}}=k$. It gives
\begin{equation*}
  \mathsf{m}_p(x)\le\tfrac92\,2^{-(k-m)}\alpha_{0,m}\,
  \frac{B^{\mathsf{s}}c_1\sup|b_i|}{\alpha_{0,k}}.
\end{equation*}
Insert this bound and the definition of $\varpi_0$. Since
$2^{8}\cdot\tfrac92=9\cdot2^{7}$, the left member of the condition at $x$ satisfies
\begin{align*}
  2^{8}B_\natural C_T^{\mathrm{w}}c_1(1+\varpi_0)\mathsf{m}_p(x)
  &\le\frac{9\cdot2^{7}B_\natural C_T^{\mathrm{w}}c_1^{2}B^{\mathsf{s}}}{K_o^{(\mathrm{L})}}\,
   2^{-(k-m)}\alpha_{0,m}
   =\frac{9}{2^{10}}\,\theta_\rho\beta\,2^{-(k-m)}\alpha_{0,m}.
\end{align*}
The right side is at most $\tfrac1{16}\theta_\rho\beta2^{-(k-m)}\alpha_{0,m}$, because
$9/2^{10}\le2^{6}/2^{10}$.

Take next a point $x\in\tilde\square\cap X^{+}$ lying in $\mathfrak{A}$. Since
$e^{-\delta_1d(x,\bar\Lambda)/4}\le1$, we
have $\mathsf{m}_p(x)\le B^{\mathsf{s}}c_1\sup|b_i|$. Therefore
\begin{equation*}
  2^{8}B_\natural C_T^{\mathrm{w}}c_1(1+\varpi_0)\mathsf{m}_p(x)
  \le\frac{2^{8}\theta_\rho\beta}{2^{17}}\,\alpha_{0,k}.
\end{equation*}
It remains to see that $2^{-9}\theta_\rho\beta\le1/400$. The product $\theta_\rho\beta$ is
at most $1$ by the choice of the room fraction $\theta_\rho$ and of the exponent $\beta$ in
the order of choice of the constants; hence the right side is at most $2^{-9}\alpha_{0,k}$,
which is less than $\alpha_{0,k}/400$.

By the assumption on $\square$, every point of $\tilde\square\cap X^{+}$ lies in
$\mathfrak{O}$ or in $\mathfrak{A}$, so $\varpi_0$ satisfies the condition in
\eqref{9.1:eq-large-field-circles-2}, and \eqref{9.1:eq-large-field-circles-4} follows.
\end{proof}

\begin{definition}[Units, radii and scaffold integers in the graded region]\label{9.1:not-graded-units}
Throughout this section $k$ denotes the coarsest level in play. The following
letters are fixed in addition to those of Notation~\ref{9.1:not-site-blind-constants}.

\emph{Growth number and minimal coefficient.} We keep $\beta:=\beta_I=1/5$, as in
Notation~\ref{9.1:not-site-blind-constants}, and put
\begin{equation}\label{9.1:eq-graded-units-1}
\varphi_{\min}:=1-4\beta=\tfrac15 .
\end{equation}
Every coefficient of the exponents in part (i) of the bound of the large-field step treated in
the remark of this chapter on the table coefficients is at least $1-4\beta$. This lower bound is
the content of that remark, and it is used here as stated there.

\emph{Length units and the units of the table members.} Let $\eta>0$ be the spacing of the
finest lattice. For each level $n$ we put $\ell_n:=L^n\eta$, so that $\ell_k=L^k\eta$ is the
unit of the coarsest level. For a table member $E_a$ and a level $n$ we put
\begin{equation}\label{9.1:eq-graded-units-2}
\mathsf{u}^a_n:=\alpha_{a,n}\,\ell_n^{-e_a}.
\end{equation}
Here $\alpha_{a,n}\in\{\alpha_{0,n},\alpha_{1,n}\}$ is the radius carried by the member $a$,
and the choice is fixed with the member. The exponent $e_a\in\{1,2,3\}$ is fixed with the
member; it is given by
\begin{equation*}
e_6=1,\qquad e_1=e_5=e_7=2,\qquad e_3=3 .
\end{equation*}
The powers of $\eta$ are common to all levels. They cancel in every ratio of units of one
member at two levels, and they are suppressed.

\emph{Ba{\l}aban's integer $\ell_\circ$.} We write $\ell_\circ$ for the integer denoted $L_0$
in \cite{B-CMP122a}; it is distinct from the constant $L_0$ of hypothesis (H2). The integer
$\ell_\circ$ satisfies the one-sided conditions
\begin{equation}\label{9.1:eq-graded-units-3}
\ell_\circ^2\ge 4,\qquad L\ge 3\ell_\circ^2 .
\end{equation}
The second of them is one of the lower bounds on $L$ in the order of choice
(Lemma~\ref{2.4:lem-stage-zero-data}).

\emph{Radii across levels.} As in Notation~\ref{9.1:not-site-blind-constants}, for $n'>n$
and $i\in\{0,1\}$,
\begin{equation*}
\alpha_{0,n}\le\alpha_{1,n},\qquad
\alpha_{i,n'}\le\tfrac94\,(n'-n)^{1/2}\,\alpha_{i,n},\qquad
\alpha_{i,n}\le 2\,\alpha_{i,n'} .
\end{equation*}
The middle inequality is \cite[(2.6)--(2.8)]{B-CMP119}.

\emph{Ba{\l}aban's small parameters at level $k$.} With $x_k=g_k^{-2}$ we put
\begin{equation}\label{9.1:eq-graded-units-4}
\begin{aligned}
\varepsilon_k&:=A_0\,g_k(\log x_k)^{p_0},\\
\delta'_k&:=A_1\,g_k(\log x_k)^{p_1},\\
\alpha_{i,k}&:=C_i\,g_k(\log x_k)^{q}\qquad(i=0,1).
\end{aligned}
\end{equation}
The sources of these three quantities are as follows.
\begin{itemize}
\item $\varepsilon_k$ is the quantity of \cite[(2.4)]{B-CMP119}.
\item $\delta'_k$ is the parameter $\delta'_j=g_jA_1(\log g_j^{-2})^{p_1}$ of
\cite{B-CMP122a} at $j=k$. There $p_1<p_0$.
\item $\alpha_{i,k}$ are the radii of \cite[(2.28)]{B-CMP119}. There $q_0,q_1$ are integers
greater than $1$, and $C_0,C_1$ are sufficiently large positive numbers.
\end{itemize}
We take $q:=q_0=q_1$; it obeys the one-sided floor
\begin{equation}\label{9.1:eq-graded-units-5}
q\ge p_0+\bigl(1+d+d\beta_0^{pr}\bigr)r_{pr}+1 ,
\end{equation}
where $d=4$ is the dimension. This floor is one of the floors on the exponents fixed after $L$
in the order of choice (Notation~\ref{2.3:not-stages-middle} and~\ref{2.3:not-stages-last}); it
is used here as fixed there.
The constant $C_1$ is that of \cite[(1.20)]{B-CMP116}, and the ratio $\mathsf{r}_1=C_1/C_0$ of
the radii is that of Notation~\ref{9.1:not-site-blind-constants}. Ba{\l}aban's radii occur in
this section only through the numbers $\alpha_{i,n}$ and the ratio $\mathsf{r}_1$.

\emph{Scaffold integers.} For $x>1$ let $\check{R}(x)$ be the least power of $L$ that is at
least $(\log x)^{r_{pr}}$; this is the integer of \cite[(2.5)]{B-CMP119}. It is written with a
check to keep it apart from the letter $R$ of the glossary. For every level $j$ we put
$\check{R}_j:=\check{R}(x_j)$ with $x_j=g_j^{-2}$, and
\begin{equation*}
N:=\check{R}_k=\check{R}(x_k),\qquad h:=k-N .
\end{equation*}
From here on in this section the letter $N$ denotes the integer $\check{R}_k$, not the rank of
the gauge group. Likewise, in this section $h$ denotes the level $k-N$, not a history. The
integer $N$ is the value of $N$ in the two conditions, labelled (i) and (ii) in
\cite{B-CMP122a}, on a large-field region $Z$ under which the operation $\mathbf{R}$ is applied.
The operation $\mathbf{R}$ is applied to a region $Z$ if and only if $Z$ satisfies those two
conditions with this value of $N$ and, in addition, a nesting condition.

Let $N_0$ be the smallest positive integer such that
\begin{equation}\label{9.1:eq-graded-units-6}
L^{-N_0+1}\,M\,\check{R}_{k-N_0+1}=M ,
\end{equation}
as in \cite{B-CMP122a}. We put $k_0:=k-N_0$, so that $k-k_0=N_0$ and
$L^{N_0-1}=\check{R}_{k-N_0+1}$.

\emph{Letters fixed elsewhere in this chapter.} The following letters keep the meanings fixed
in Notation~\ref{9.1:not-site-blind-constants}:
\begin{itemize}
\item the read level $\iota$ and the level $n_F$;
\item the path metric (written $d(\cdot,\cdot)$, with two arguments, in
Lemma~\ref{9.1:lem-link-bound}) with its conditions, the constants $c_1$ and
$C_T^{w}=c_1L^3$, and the smeared profile $f^{\approx}$.
\end{itemize}
In this section the following letters have these meanings:
\begin{itemize}
\item the slice map $S(\cdot)$ and the table members $E_a$;
\item the truncated determining set $\mathfrak{S}:=\boldsymbol{\Omega}(X^{+})$ of a large-field
region $X$, with the weight $\ell_{\mathfrak{S}}(y)$ at a point $y$ that enters the weighted
size below and the quantity $\lambda_{\mathfrak{S}}$ attached to $\mathfrak{S}$;
\item the two classes of points $(=)$ and $(\mathrm{up})$, and the constant $B_\natural$ of the
local Lipschitz property of the Landau chart;
\item the weighted size $|\mathcal{Y}|_y:=\ell_{\mathfrak{S}}(y)\sup_{\tilde{\Delta}(y)}|\mathcal{Y}|$,
where $\tilde{\Delta}(y)$ is the cell attached to the point $y$.
\end{itemize}
The substitution chain is that of Definition~\ref{9.1:def-substitution-chain}. The root seed
$q_0$, the field $\mathcal{Z}_q$, the profile $\mathsf{m}_L$, the constant $C_4$ and the
dilation factor $B^{\mathsf{s}}$ are those of Lemma~\ref{9.1:lem-link-bound}, together with the
condition on the region under which that lemma is stated and the profile bound invoked there,
through which $B^{\mathsf{s}}$ enters, with
\begin{equation*}
\mathsf{s}_L:=K_\Lambda'\,(1+\mathsf{r}_1)\,\alpha_{0,k}.
\end{equation*}
\end{definition}

\begin{proof}
We prove four relations.

First, \eqref{9.1:eq-graded-units-1}: with $\beta=1/5$ we have $1-4\beta=1-4/5=1/5$.

Second, the cancellation of the powers of $\eta$. By \eqref{9.1:eq-graded-units-2} and
$\ell_n=L^n\eta$, for one member $a$ and two levels $n,n'$,
\begin{equation*}
\frac{\mathsf{u}^a_n}{\mathsf{u}^a_{n'}}
=\frac{\alpha_{a,n}\,L^{-e_an}\eta^{-e_a}}{\alpha_{a,n'}\,L^{-e_an'}\eta^{-e_a}}
=\frac{\alpha_{a,n}}{\alpha_{a,n'}}\,L^{e_a(n'-n)} .
\end{equation*}
The right-hand side does not contain $\eta$.

Third, the inequality $\ell_\circ^2\le L/3$, which appears in the hypothesis of
\cite{B-CMP122a} together with the condition $\beta\le1/4$, met by the value $\beta=1/5$ fixed
above, holds: dividing the second inequality of \eqref{9.1:eq-graded-units-3} by $3$ gives
$\ell_\circ^2\le L/3$.

Fourth, the relation between $N_0$ and the scaffold integers. Since $M>0$, dividing
\eqref{9.1:eq-graded-units-6} by $M$ gives $L^{-N_0+1}\check{R}_{k-N_0+1}=1$. This is
$L^{N_0-1}=\check{R}_{k-N_0+1}$. Finally, $k-k_0=N_0$ is the definition of $k_0$.
\end{proof}

\begin{definition}[Rims, bulk, tower and core of the graded region]\label{9.1:def-graded-geometry}
We work at an application of Ba{\l}aban's operation $\mathbf{R}$ at step $k$
(Definition~\ref{1.5:def-densities}), with $Z$ the processed components \cite[\S1]{B-CMP122a}. Let $h$
and $N_0$ be the integers of \cite[(1.10)--(1.12)]{B-CMP122a}, and put $k_0:=k-N_0$. We assume
$h+1\le k_0<k$. Here $h$ denotes this integer, not a history.

\emph{The graded sequence.} Following \cite[p.~179, (1.10)]{B-CMP122a},
\begin{equation}\label{9.1:eq-graded-geometry-1}
Z''_k=\bigl(\Omega^{\sim5}_{k_0+1}\bigr)^c\cap Z,\qquad
Z''_{k_0+1}=\bigl(\Omega^{\sim7}_{k_0+1}\bigr)^c\cap Z .
\end{equation}
Here $\Omega_{k_0+1}\subset\Omega^{\sim5}_{k_0+1}\subset\Omega^{\sim7}_{k_0+1}$ are the enlargements of
$\Omega_{k_0+1}$ in the notation of \cite[(1.10)]{B-CMP122a}.
These two sets are completed to a sequence $Z''_k,Z''_{k-1},\dots,Z''_{k_0+1}$ whose complements form an
admissible sequence of domains based on partitions into $M$-cubes of the corresponding scales: $Z''_k$ and
$Z''_{k-1}$ are separated by one layer of $M$-cubes, $Z''_{k-1}$ and $Z''_{k-2}$ by one layer of
$L^{-1}M$-cubes, and so on \cite[(1.10)--(1.12)]{B-CMP122a}. Further, for
$j=k_0,\dots,h+1$, $Z''_j$ is the intersection with $Z$ of the complement of the domain attached to level
$j$ in \cite[(1.11)]{B-CMP122a}, and $Z''_j=Z_j$ for $j=h,\dots,1$. The sequence is
decreasing, $Z''_k\supset Z''_{k-1}\supset\cdots\supset Z''_{h+1}$. The sets $\Omega''_j$ of this chapter
are $\Omega''_j=(Z''_j)^c$ for $j\ge h+1$ and $\Omega''_j=\Omega_j$ for $j\le h$,
complements being taken inside $Z$ with the convention of \cite[\S1]{B-CMP122a}. The integration domain of
\cite[(1.73)]{B-CMP122a} is
\begin{equation*}
\Lambda=\bigl(\Omega^{(4)}_{k_0+1}\bigr)^c\cap Z ,
\end{equation*}
with $\Omega^{(4)}_{k_0+1}$ as there. The set $\Lambda$ is obtained by adding one layer of $M$-cubes to
$Z''_k$ and is contained in a cube of size $100M$; all field variables on it are integrated out
\cite[(1.73)]{B-CMP122a}. Put $\Lambda_0:=\Lambda\cap\Omega''_{h+1}$. The fluctuation
field $V$ is defined on $\mathbf{B}_k\cap\Lambda\cap\Omega''_{h+1}$; on $\Lambda_0$ one sets
$V''=V'V_0$, $V_0=M_{\mathfrak{B}_0}(U_0)$, with $V_0$ defined on the bond set $\mathfrak{B}_0$, which
contains every bond meeting $\partial\Lambda$; $V$, $V'$, $U_0$, the map $M_{\mathfrak{B}_0}$ and
$\mathfrak{B}_0$ are as in \cite[(1.81)]{B-CMP122a} (the letter $M$ with the subscript $\mathfrak{B}_0$
denotes that map, not the number $M$). The characteristic function of
\cite[(1.82)]{B-CMP122a} is
\begin{equation*}
\chi'=\chi\bigl(\{|B'(b)|<\delta'_k\ \text{for}\ b\in\mathfrak{B}_0\}\bigr),
\end{equation*}
with $B'$ and $\delta'_k$ as there. The domains
\begin{equation*}
\Lambda\supset Z''_k\supset Z''_{k-1}\supset\cdots\supset Z''_{h+1}\supset(\Omega''_{h+1})^c
\end{equation*}
are rectangular parallelepipeds \cite[\S1]{B-CMP122a}, and a gauge is fixed in the differences of
consecutive members, intersected with the lattice of the corresponding scale. By the lemma of this chapter
on the nesting of these domains, the sets
\begin{equation*}
\Lambda\supset Z''_k\cap Z\supset\cdots\supset Z''_{h+1}\cap Z\supset R_{\mathrm{in}}
\end{equation*}
form one nested chain of parallelepipeds; its innermost member $R_{\mathrm{in}}$ is a parallelepiped
contained in $Z''_{h+1}\cap Z$.

\emph{The level function.} Levels of determining sets and labels are those of
\cite[(1.14)--(1.17)]{B-CMP122a}; the stages and the stage-level function are those of the stage
construction of the determining sets at step $k$. Let $N_{\mathrm{st}}$ be the number of stages of the
construction of the determining sets at step $k$. For
$x\in Z''_k$, let $m(x)$ be the level at $x$ of the determining set
$\mathbf{B}''_k=\mathbf{B}_k^{(N_{\mathrm{st}})}$, that is, the value at $x$, after the last stage, of the
stage-level function of the stage construction of the determining sets at step $k$. Here $m(\cdot)$
denotes this level function, not the torus exponent.

\emph{The pieces.} Here $\Omega''^{\sim}_{h+1}$ is the enlargement of $\Omega''_{h+1}$ by one
$\sim$-layer, of width $M\ell_h$, with $\ell_h$ the length unit at level $h$
(Notation~\ref{9.1:not-graded-units}), used in \cite[(1.52)--(1.53)]{B-CMP122b}. We call
\begin{align*}
&\mathfrak{K}:=\bigl(\Omega''^{\sim}_{h+1}\bigr)^c\cap Z\ \text{the \emph{core}},\\
&Z''_k\setminus Z''_{k_0+1}\ \text{the \emph{rims}},\qquad
Z''_{k_0+1}\setminus Z''_{k_0}\ \text{the \emph{bulk}},\\
&Z''_{k_0}\setminus\mathfrak{K}\ \text{the \emph{tower}},\qquad
Z''^{(1)}_k:=Z''_k\setminus\mathfrak{K}\ \text{the \emph{first side}}.
\end{align*}
By \cite[(1.52)--(1.53)]{B-CMP122b}, the configuration $U'_k$ there is broken into two independent parts,
supported in the domains $\Omega''^{\sim}_{h+1}$ and $(\Omega''^{\sim}_{h+1})^c$ respectively, by boundary
conditions at $\partial\Omega''^{\sim}_{h+1}$; the configuration is cut there. The core is the second of
these domains intersected with $Z$, and the first side is the part of $Z''_k$ in the first. In the
estimates of this chapter across $\partial\Omega''^{\sim}_{h+1}$, which we call the seam, the
configuration on $\mathfrak{K}$ is held fixed, and $\mathfrak{K}$ is the core beyond the seam to which
those estimates refer. Since
$\Omega''_{h+1}\subset\Omega''^{\sim}_{h+1}$, since the sequence is decreasing, and since
$\Lambda$ is obtained from $Z''_k$ by adding one layer of $M$-cubes,
\begin{equation*}
\mathfrak{K}\subset(\Omega''_{h+1})^c\cap Z=Z''_{h+1}\subset Z''_{k_0}\subset Z''_k\subset\Lambda .
\end{equation*}
Hence the rims, the bulk, the tower and the core partition $Z''_k$, which we call the
\emph{graded region}.

\emph{Properties.} For $x\in Z''_k$ and $h\le j\le k-1$,
\begin{equation}\label{9.1:eq-graded-geometry-2}
m(x)=j\ \text{ if }x\in Z''_{j+1}\setminus Z''_j,\qquad
h\le m(x)\le k-1,\qquad k-m(x)\le N_{\mathrm{st}} .
\end{equation}
On the rims $k_0<m(x)<k$; on the bulk $m(x)=k_0$; on the tower $h\le m(x)<k_0$. The level of the label
entering the step is at most $k_0$ on $Z''_k$. Moreover,
\begin{equation}\label{9.1:eq-graded-geometry-3}
Z''^{(1)}_k=Z''_k\cap\Omega''^{\sim}_{h+1},
\end{equation}
and the innermost $\sim$-layer of $Z''^{(1)}_k$ has $m=h$.
\end{definition}

\begin{proof}
\emph{The level function.} Let $x\in Z''_k$. By \cite[(1.14)--(1.17)]{B-CMP122a}, in the notation there
(in particular for the set $\Gamma_h$ and the superscript $(j)$), the part of the label in $Z''_k$ is the
union of the blocks of $\Gamma'_j=(Z''_{j+1}\setminus Z''_j)^{(j)}$, $h+1\le j\le k-1$, of level $j$, and
of the blocks of $\Gamma_h\cap Z''_{h+1}$, of level $h$; and by the same equations these blocks keep their
levels through the $N_{\mathrm{st}}$ stages. Let $h+1\le j\le k-1$. A point of $Z''_{j+1}\setminus Z''_j$ lies in a block
of $\Gamma'_j$, so $m(x)=j$ there. Let now $j=h$ and $x\in Z''_{h+1}\setminus Z''_h$. Since the sequence is
decreasing, $x$ lies in no shell $Z''_{i+1}\setminus Z''_i$ with $i\ge h+1$, hence in no block of level
$i\ge h+1$; so $x$ lies in a block of $\Gamma_h\cap Z''_{h+1}$, and $m(x)=h$. At every point of $Z''_k$ the
level is one of those carried by these blocks, so $h\le m(x)\le k-1$. Finally, $m(x)$ is the value after
the last stage of the stage-level function of the stage construction of the determining sets at step $k$,
and the statement that introduces that function gives $k-m(x)\le N_{\mathrm{st}}$. This proves
\eqref{9.1:eq-graded-geometry-2}. In particular, the blocks of $\mathbf{B}''_k$ lying in
$Z''_k\cap\Omega''^{\sim2}_{h+1}$, where $\Omega''^{\sim2}_{h+1}\supset\Omega''^{\sim}_{h+1}$ is the
enlargement in the notation of \cite[(1.10)]{B-CMP122a}, have levels between $h$ and $k-1$.

\emph{Rims and bulk.} The sequence is decreasing. Hence a point of $Z''_k\setminus Z''_{k_0+1}$ lies in
$Z''_{j+1}\setminus Z''_j$ for exactly one $j$ with $k_0+1\le j\le k-1$. By
\eqref{9.1:eq-graded-geometry-2}, $m=j$ there, so $k_0<m<k$ on the rims. The bulk is the shell with
$j=k_0$, so $m=k_0$ there by \eqref{9.1:eq-graded-geometry-2}.

\emph{Tower.} Let $x\in Z''_{k_0}\setminus\mathfrak{K}$. Then $x\in Z''_{k_0}$, and $Z''_{k_0}$ is
disjoint from every shell $Z''_{j+1}\setminus Z''_j$ with $j\ge k_0$. By the description of the blocks
above, the levels $j\ge k_0$ are carried only by blocks $\Gamma'_j$ lying in such shells. Therefore
$m(x)<k_0$, and together with $m(x)\ge h$ from \eqref{9.1:eq-graded-geometry-2} this gives
$h\le m(x)<k_0$.

\emph{The entering label.} By \eqref{9.1:eq-graded-geometry-1},
\begin{equation*}
Z''_k\subset\bigl(\Omega^{\sim5}_{k_0+1}\bigr)^c\subset\Omega^c_{k_0+1}.
\end{equation*}
So no point of $Z''_k$ lies in $\Omega_{k_0+1}$. Since the points at which the label entering the step has
level at least $k_0+1$ lie in $\Omega_{k_0+1}$ \cite[\S1]{B-CMP122a}, the level of the label entering the
step is at most $k_0$ on $Z''_k$.

\emph{The first side.} Since $Z''_k\subset Z$ by \eqref{9.1:eq-graded-geometry-1}, a point $x\in Z''_k$
lies in $\mathfrak{K}=(\Omega''^{\sim}_{h+1})^c\cap Z$ if and only if $x\notin\Omega''^{\sim}_{h+1}$.
Hence $Z''_k\setminus\mathfrak{K}=Z''_k\cap\Omega''^{\sim}_{h+1}$, which is
\eqref{9.1:eq-graded-geometry-3}. The innermost $\sim$-layer of $Z''^{(1)}_k$ is the single layer
$\Omega''^{\sim}_{h+1}\setminus\Omega''_{h+1}$ inside $Z$. It is contained in
$(\Omega''_{h+1})^c\cap Z=Z''_{h+1}$, and $Z''_{h+1}\subset Z''_k$. On $Z''_{h+1}$ no shell
$Z''_{j+1}\setminus Z''_j$ with $j\ge h+1$ is met, so no block of level $j\ge h+1$ is met. Hence $m=h$ on
this layer, by \eqref{9.1:eq-graded-geometry-2}.
\end{proof}

\begin{definition}[The last-stage table on the graded region]\label{9.1:def-last-stage-table}
We use the graded region of Definition~\ref{9.1:def-graded-geometry}: the levels $h$, $k_0$, $k$, the
nested sets $\Omega''_m$ and $Z''_m$, and the level $m(x)$ of the last-stage label at $x$. We also
use the units, radii and integers of Notation~\ref{9.1:not-graded-units}: the length unit $\ell_m$ at
level $m$ (Ba{\l}aban's $L^m\eta$), and the radii $\alpha_{0,m}$, $\alpha_{1,m}$. We write
$\ell_\circ$ for the positive integer that \cite[(1.64)--(1.65)]{B-CMP122a} denotes $L_0$; it is
unrelated to the constant $L_0$ of hypothesis (H2) in Theorem~\ref{3.1:thm-main}.

Let $N$ denote the last stage of the construction, the stage at which the step index $j$ of
\cite[(1.64)]{B-CMP122a} equals the top level $k$. The seven members of the table are the seven
entries of Ba{\l}aban's inductive description \cite[(1.64)--(1.68)]{B-CMP122a}, denoted here
$E_1,\dots,E_7$, and $\ell_m^{-e_a}$ is the power of $\ell_m$ in which the radius of $E_a$ is
measured. For
$a\in\{1,\dots,5\}$ put $\alpha_{a,m}:=\alpha_{0,m}$, and for $a\in\{6,7\}$ put
$\alpha_{a,m}:=\alpha_{1,m}$. The unit of member $a$ at level $m$ is
$\mathsf{u}^a_m:=\alpha_{a,m}\ell_m^{-e_a}$. Write $(y)_+:=\max\{0,y\}$.

The \emph{last-stage table} on $Z''_k$ is the stage table $\mathsf{T}_N$ at the last stage, read at
the points of $Z''_k$: by the definition of the stage tables, its entry $\mathsf{T}^a_N(x)$ for
member $a$ at $x$ is the coefficient of \cite[(1.64)]{B-CMP122a} at $x$ (denoted $F$ there, and
$\varphi^{(N)}_m$ below) times the local radius
of member $a$ at $x$, which is the printed entry of \cite[(1.64)--(1.68)]{B-CMP122a} with its
coefficient dropped. For each member $a$ and each $x\in Z''_k$ it is given by
\begin{equation}\label{9.1:eq-last-stage-table-1}
\begin{split}
\mathsf{T}^a_N(x)
&=\varphi^{(N)}_m\,\ell_\circ^{2(m-k_0)_+}\,\alpha_{a,m}\ell_m^{-e_a}
=\varphi^{(N)}_m\,\ell_\circ^{2(m-k_0)_+}\,\mathsf{u}^a_m,\\
m&=m(x),\qquad \varphi^{(N)}_m\ \ge\ \varphi_{\min}=\tfrac15 .
\end{split}
\end{equation}
Here $\varphi^{(N)}_m$ is the coefficient of \cite[(1.64)]{B-CMP122a} at stage $N$, recalled in
\eqref{9.1:eq-last-stage-table-3} below, and $\varphi_{\min}=1-4\beta_I$, where $\beta_I=\tfrac15$ is
the constant denoted $\beta$ in \cite[(1.64)]{B-CMP122a}, so that $\varphi_{\min}=\tfrac15$.

Every configuration of the domain $\mathfrak{U}[\mathsf{T}_N]$ attached to this table satisfies at
$x\in Z''_k$ the following
\emph{cube clause}. Let $m=m(x)$. For every cube $\square\ni x$ of the size $CM\ell_m$ with
$\square\subset\Omega''_m$ there is a gauge in which the vector potential $A$ of the configuration
satisfies
\begin{equation}\label{9.1:eq-last-stage-table-2}
\ell_m|A|,\ \ \ell_m^2|\nabla A|\ <\ \ell_\circ^{2(m-k_0)_+}\,B\,C\,M\,\alpha_{0,m}.
\end{equation}
Here $B$ and $C$ are Ba{\l}aban's constants of \cite[(1.65)]{B-CMP122a}.

The assertion attached to this definition has two parts. First, at every $x\in Z''_k$ the entries
of the stage table $\mathsf{T}_N$ are the numbers \eqref{9.1:eq-last-stage-table-1}, and the clause
\eqref{9.1:eq-last-stage-table-2} is implied by the clause that \cite[(1.64)--(1.65)]{B-CMP122a}
prescribes at stage $N$, its bound coinciding with that of \cite[(1.65)]{B-CMP122a} when
$m(x)>k_0$. Second, the
constant denoted $c$ in \cite[(1.64), (1.68)]{B-CMP122a} equals $1$ throughout $Z''_k$; the value
$c=3$ of the clause $m=k$ in \cite[(1.68)]{B-CMP122a} does not occur there.
\end{definition}

\begin{proof}
\emph{Region of the point.} At stage $N$ we have $j=k$. Let $x\in Z''_k$ and $m=m(x)$.

By Definition~\ref{9.1:def-graded-geometry}, $m(x)=j'$ on $Z''_{j'+1}\setminus Z''_{j'}$, and
$h\le m(x)\le k-1$. Hence $x\in Z''_{m+1}\setminus Z''_m$.

A point lies in $Z''_i$ exactly when it does not lie in $\Omega''_i$, by
Definition~\ref{9.1:def-graded-geometry}. Therefore
\begin{equation*}
Z''_{m+1}\setminus Z''_m=\Omega''_m\setminus\Omega''_{m+1}.
\end{equation*}
The set $Z''_{m+1}\setminus Z''_m$ is the part of the large-field region $Z$ (Ba{\l}aban's $X$ in
\cite[(1.64)]{B-CMP122a}) lying in $\Omega''_m\setminus\Omega''_{m+1}$, so
\begin{equation*}
Z''_{m+1}\setminus Z''_m=(\Omega''_m\setminus\Omega''_{m+1})\cap Z .
\end{equation*}
For $m=1,\dots,j-1$ these sets form the region of \cite[(1.64)]{B-CMP122a} on which it gives the
entries recalled below.

Since $m\le k-1=j-1$, the point $x$ lies in this region. It does not lie in Ba{\l}aban's set
$(\Omega''_j\setminus\Omega_{k_0+1})\cap Z$ of \cite[(1.64)]{B-CMP122a}, on which the entries for
$m=j$ are given, because that set is contained in $\Omega''_j=\Omega''_k$, while $x\notin\Omega''_k$.

\emph{Entries.} On this region, \cite[(1.64)]{B-CMP122a} prescribes three things. The first is the
coefficient, denoted $F$ in \cite[(1.64)]{B-CMP122a}; with its step index $j$ equal to $k$ it is
\begin{equation}\label{9.1:eq-last-stage-table-3}
\varphi^{(N)}_m:=1-\beta_I\sum_{i=h+1}^{k}2^{-|m-i|}-\beta_I\sum_{i'=m+1}^{k}2^{-(i'-m)}.
\end{equation}
The second is the weight factor, equal to
$\ell_\circ^{2\max\{0,m-k_0\}}=\ell_\circ^{2(m-k_0)_+}$ (denoted by a separate letter in
\cite[(1.64)]{B-CMP122a}). The third is the value $1$ of the constant denoted $c$ there.

By the definition of the stage tables, the entry of member $a$ at $x$ is the coefficient
$\varphi^{(N)}_m$ times the local radius of member $a$ at $x$, and the local radius is the printed
entry of \cite[(1.64)--(1.68)]{B-CMP122a} with its coefficient dropped, that is, the weight factor
times the radius of the member. By the list of the members' radii in
Notation~\ref{9.1:not-graded-units}, the radii of the seven members are
$\alpha_{0,m}$ for $E_1,\dots,E_5$ and $\alpha_{1,m}$ for $E_6$
and $E_7$. They are measured in the powers $\ell_m^{-e_a}$, that is, the radius of member $a$ is
$\alpha_{a,m}\ell_m^{-e_a}=\mathsf{u}^a_m$. Multiplying coefficient, weight factor and radius gives the
first equality in \eqref{9.1:eq-last-stage-table-1}. The second equality is the definition of
$\mathsf{u}^a_m$.

\emph{Cube clause.} \cite[(1.65)]{B-CMP122a} requires, on cubes of the size $CM\ell_m$, a gauge in
which $\ell_m|A|$ and $\ell_m^2|\nabla A|$ are bounded by $BCM\alpha_{0,m}$ times a power of
$\ell_\circ$; there $\ell_m$ is written $L^m\eta$, $\ell_\circ$ is written $L_0$, and the power of
$\ell_\circ$ there is $\ell_\circ^{2(m-k_0)}$. The weight on the right of
\eqref{9.1:eq-last-stage-table-2} is the weight factor of \cite[(1.64)]{B-CMP122a} at level $m$,
$\ell_\circ^{2\max\{0,m-k_0\}}=\ell_\circ^{2(m-k_0)_+}$. For $m>k_0$ it coincides with the power
$\ell_\circ^{2(m-k_0)}$ of \cite[(1.65)]{B-CMP122a}. For $m\le k_0$ it equals $1$; since
$\ell_\circ\ge1$, the bound
$\ell_\circ^{2(m-k_0)}BCM\alpha_{0,m}$ of \cite[(1.65)]{B-CMP122a} is then at most
$BCM\alpha_{0,m}$, so a gauge satisfying the condition of \cite[(1.65)]{B-CMP122a} satisfies
\eqref{9.1:eq-last-stage-table-2}. Reading the condition at
$m=m(x)$ for the cubes $\square\ni x$ with $\square\subset\Omega''_m$ therefore gives
\eqref{9.1:eq-last-stage-table-2}.

\emph{Lower bound on the coefficient.} The bound $\varphi^{(N)}_m\ge 1-4\beta_I$ for the coefficients
\eqref{9.1:eq-last-stage-table-3} follows from two geometric sums. For the first sum,
\begin{equation*}
\sum_{i=h+1}^{k}2^{-|m-i|}\le\sum_{i\in\mathbb{Z}}2^{-|m-i|}=1+2\sum_{i\ge1}2^{-i}=3 .
\end{equation*}
For the second sum,
\begin{equation*}
\sum_{i'=m+1}^{k}2^{-(i'-m)}<\sum_{i\ge1}2^{-i}=1 .
\end{equation*}
Hence, with $\beta_I=\tfrac15$,
\begin{equation*}
\varphi^{(N)}_m>1-4\beta_I=\varphi_{\min}=\tfrac15,
\end{equation*}
which is the last assertion of \eqref{9.1:eq-last-stage-table-1}.

In particular, for $m\le h$ every index $i$ of the first sum satisfies $i>m$. That sum is then
strictly less than $\sum_{i\ge1}2^{-i}=1$, as is the second sum for every $m$; so both sums are $<1$,
and there $\varphi^{(N)}_m>1-2\beta_I$.

\emph{The constant $c$.} The value $c=3$ of the constant $c$ in \cite[(1.68)]{B-CMP122a} belongs to
the clause $m=k$. Since $m(x)\le k-1$ for every $x\in Z''_k$, that clause does not occur on $Z''_k$,
and $c=1$ throughout.
\end{proof}

\begin{definition}[The root seed, the chain, aligned and older points]\label{9.1:def-aligned-points}
Let $\Phi$ be the output of Ba{\l}aban's operation $\mathbf{R}$ at the large-field site of
Definition~\ref{9.1:def-substitution-chain}, let $X$ be its region (in this section $X$ denotes a
region, as in \cite{B-CMP122b}), and let $\mathsf{u}_\Phi$, $\mathsf{c}_\Phi^{+}$ and
$\varrho_\Phi$ be its own units, its table coefficient and its room; the level and the own units
assigned to a region form its \emph{label}.
Let $Z$ be the large-field region of the step, and let $Z''_k$, $\Omega''_k$, $\Lambda$,
$\Lambda^{\sim}$, the sets $Y_i$ (among them $Y_3$), $Y_3^0$, $\mathbf{B}_k$, $\mathbf{B}''_k$,
$\mathbf{B}_1$, $\mathfrak{B}_0$ be the sets of \cite{B-CMP122b} at this step. Let $Z''^{(1)}_k$
denote the region inside which the sets of the chain are compared in item~\textup{(ii)}, let
$m(\cdot)$, $h$ and $k_0$ be as in Definition~\ref{9.1:def-graded-geometry}, and let $N_0$ denote
the integer, accompanying $k_0$, for which $k-k_0\ge N_0$.
\begin{enumerate}
\item[(i)] \emph{The root seed.} At the large-field site the chain is entered at the \emph{root
seed}
\begin{equation*}
q_0=e^{i\eta A_e}U_c ,
\end{equation*}
where $U_c$ is its real centre and $A_e$ its Landau coordinate; $q_0$ is a datum of the seed tube
$\mathfrak{T}_\Phi[7/8]$ of $\Phi$, and its slice on $X$ satisfies
\begin{equation}\label{9.1:eq-aligned-points-1}
S(q_0)=q_0{\upharpoonright}X\in\mathfrak{D}_\Phi[7/8]
=\mathfrak{U}\bigl[(\mathsf{c}_\Phi^{+}-\tfrac18\varrho_\Phi)\,\mathsf{u}_\Phi\bigr].
\end{equation}
\item[(ii)] \emph{The chain.} The \emph{chain objects} are the configurations
$q=e^{i\eta\mathcal{Z}_q}q_0$ of Lemma~\ref{9.1:lem-link-bound}, built by rows $0,3,2,1,0'$ of
Definition~\ref{9.1:def-substitution-chain}, which are the substitutions
\cite[(1.63), (1.58)/(1.64), (1.57), (1.54), (1.51)]{B-CMP122b} in this order, and by the layer
re-minimisations \cite[(1.61)--(1.63)]{B-CMP122b}. Inside $Z''^{(1)}_k$ row~$3$ is solved on
$\mathbf{B}_k(Y_3)=\mathbf{B}(Y_3)\cup\mathbf{B}_k$, and rows $2$, $1$, $0'$ re-solve on
$\mathbf{B}_1\supseteq\mathbf{B}''_k{\upharpoonright}Z''_k$, on $\mathbf{B}''_k$ and on
$\mathfrak{B}_0$ respectively.
\item[(iii)] \emph{Aligned and older points.} Let $F$ be the leaf of level $j$ read at the
large-field site, let $\mathfrak{S}$ be its truncated determining set, $\lambda_{\mathfrak{S}}(x)$
the level of $\mathfrak{S}$ at $x$ and $n_F(x)$ the level of the units of $F$ at $x$. On $X$,
\begin{equation}\label{9.1:eq-aligned-points-2}
\lambda_{\mathfrak{S}}(x)=\min\{j,m(x)\}.
\end{equation}
A point $x\in X$ is an \emph{aligned point} (an \textup{(a)}-point) if
$\lambda_{\mathfrak{S}}(x)=m(x)$, and an \emph{older point} (a \textup{(b)}-point) if
$\lambda_{\mathfrak{S}}(x)<m(x)$; at an older point we put $u:=k-n_F(x)$.
\end{enumerate}
With these definitions the following hold.
\begin{enumerate}
\item[(1)] The own units of $\Phi$ on $Z$ are of level $k$, and $\mathsf{c}_\Phi^{+}=1$ there.
\item[(2)] Inside $Z''^{(1)}_k$ the label read at the large-field site, the sets of rows $2$,
$1$, $0'$, and the last-stage table of Definition~\ref{9.1:def-last-stage-table} all live at the
graded level $m(x)$; only row~$3$ and the root seed are of level $k$.
\item[(3)] Every point of a leaf of level $j=k$ is aligned; these leaves are the chain terms
$\mathbb{C}^{(i)}_k$ of this step and the left-over terms, which are assigned level $j=k$ at this
step. Every point of the equal-level class $(=)$ is aligned.
\item[(4)] At an older point, $j<m(x)\le k-1$, so $F$ is an old leaf; then
$n_F(x)\le\min\{j,\lambda_{\mathrm{old}}(x)\}\le k_0$, where $\lambda_{\mathrm{old}}(x)$ is the
level carried at $x$ by the label of the old leaf $F$, and $u\ge N_0$; moreover
$n_F(x)\le\lambda_{\mathfrak{S}}(x)=j$.
\end{enumerate}
\end{definition}

\begin{proof}[Proof of the assertions in Definition~\ref{9.1:def-aligned-points}]
\emph{(1).} By \cite{B-CMP122b}, $R'^{(k)}(X)$ is an analytic function of the variables
$(\mathbf{U},\mathbf{J})$ defined on the space $U^c_k(X,\alpha_{0,k},\alpha_{1,k})$. This space is
built on the new sequence of regions produced by $\mathbf{R}$, in which $Z$ no longer appears
\cite{B-CMP122b}. Hence the label given to
$\Phi$ has level $k$ on $Z$, so that its own units there are of level $k$; and a label of level
$k$ created at this step carries table coefficient $\mathsf{c}^{+}=1$, by the definition of the
room coefficients. Thus $\mathsf{c}_\Phi^{+}=1$ on $Z$.

\emph{The seed slice.} The membership \eqref{9.1:eq-aligned-points-1} follows from the
definitions of the own tube and of the room of a term, from the statement that the slice of a
datum of the seed tube $\mathfrak{T}_\Phi[7/8]$ lies in the own tube $\mathfrak{D}_\Phi[7/8]$,
and from clause~(a) of the statement on slice domains.

\emph{The chain objects.} By Lemma~\ref{9.1:lem-link-bound} every chain object has the form
$q=e^{i\eta\mathcal{Z}_q}q_0$. All sources of the rows are supported in $\Lambda$, as stated
with their list in Definition~\ref{9.1:def-substitution-chain}.

\emph{(2).} Row $3$: by \cite[(1.63)]{B-CMP122b}, with $\mathbf{B}_k$ the set of the new
sequence \cite{B-CMP122b}, the set $\mathbf{B}_k(Y_3)=\mathbf{B}(Y_3)\cup\mathbf{B}_k$ is of level
$k$ on $\Lambda$. Its layers $\mathbf{B}(Y_3)$ lie at $\partial Y_3^0$, off the interior
$(\Lambda^{\sim})^{\circ}$: by the geometric statements of the junction one has
$Z''_k\subset\Lambda$, and $\Lambda$ lies in the interior of the enclosing set fixed there, so
that no boundary $\partial Y_i$ cuts a stratum of lower
level. Rows $2$, $1$, $0'$: they re-solve on $\mathbf{B}_1$, $\mathbf{B}''_k$, $\mathfrak{B}_0$,
graded at the levels $m(\cdot)$ of Definition~\ref{9.1:def-graded-geometry}, by
\cite[(1.52)]{B-CMP122b}. For row~$2$, in the notation of \cite[(1.57), p.~371]{B-CMP122b}, where
$U_{\mathbf{B}_1}(\cdot,\cdot,\cdot)$ is the minimiser on $\mathbf{B}_1$ with the three listed data
and $V''$, $M^k$, $M^{\cdot}$, $U_{k,\Lambda}$, $\mathbf{H}_{\mathbf{B}_1}$, $u_1$ are as defined
there, the minimiser is
\begin{align*}
U_1&=U_{\mathbf{B}_1}\Bigl(V''{\upharpoonright}_{\Lambda^c},\,
M^k(U_{k,\Lambda}){\upharpoonright}_{\Lambda\cap\Omega''_k},\,
\bigl(M^{\cdot}(U_{k,\Lambda})\bigr)^{-1}M^{\cdot}(U_{k,\Lambda}){\upharpoonright}_{Z''_k}\Bigr)\\
&=\Bigl(\exp i\eta\,\mathbf{H}_{\mathbf{B}_1}\bigl(-\tfrac1i\log
M^{\cdot}(U_{k,\Lambda}){\upharpoonright}_{Z''_k}\bigr)\,U_{k,\Lambda}\Bigr)^{u_1^{-1}} ;
\end{align*}
for row $0'$ the minimisation on $\mathfrak{B}_0$ is \cite[(1.50)/(1.51)]{B-CMP122b}, together
with \cite[(1.81)]{B-CMP122a}. The last-stage table of Definition~\ref{9.1:def-last-stage-table}
is evaluated at $m=m(x)$. Therefore inside $Z''^{(1)}_k$ the label read at the large-field site,
the sets of rows $2$, $1$, $0'$ and the last-stage table all live at the graded level $m(x)$,
while row~$3$ and, by (1), the root seed are of level $k$.

\emph{The level function.} Identity \eqref{9.1:eq-aligned-points-2} holds by the definition of the
level of a truncated determining set and the statement computing that level on $X$.

\emph{(3).} For a leaf of level $j=k$, \eqref{9.1:eq-aligned-points-2} gives
$\lambda_{\mathfrak{S}}(x)=\min\{k,m(x)\}=m(x)$, because $m(x)\le k-1$ by
Definition~\ref{9.1:def-graded-geometry}. The leaves of level $k$ are the chain terms
$\mathbb{C}^{(i)}_k$ of this step and the left-over terms, which are assigned level $j=k$ at this
step. The points of the equal-level class $(=)$ are included among the aligned points.

\emph{(4).} If $\lambda_{\mathfrak{S}}(x)<m(x)$, then $\min\{j,m(x)\}<m(x)$ by
\eqref{9.1:eq-aligned-points-2}, so $j<m(x)\le k-1$; thus $j<k$ and $F$ is an old leaf, and
$\lambda_{\mathfrak{S}}(x)=j$. For this old leaf
$n_F(x)\le\min\{j,\lambda_{\mathrm{old}}(x)\}\le k_0$, so that
\begin{equation*}
u=k-n_F(x)\ge k-k_0\ge N_0
\end{equation*}
by the choice of $N_0$ above. Finally $n_F(x)\le j=\lambda_{\mathfrak{S}}(x)$.
\end{proof}

\begin{lemma}[Graded thresholds exceed the top-level unit]\label{9.1:lem-graded-thresholds}
Put
\begin{equation}\label{9.1:eq-graded-thresholds-1}
\Gamma_\Lambda:=\frac{4}{45}\,\ell_\circ^{2N_0}.
\end{equation}
Let $x\in Z''_k$, put $m=m(x)$, and let $a$ be any member of the last-stage table of
Definition~\ref{9.1:def-last-stage-table}. Here the region $Z''_k$, the level $m(x)$, the integers
$k_0$ and $N_0$, the units $\ell_j$, the radii $\alpha_{a,j}$ with the exponents $e_a$, and the
regions $\Omega''_j$ are those of Notation~\ref{9.1:not-graded-units}, and $\ell_\circ$ is
the integer of \cite{B-CMP122a} denoted there by the letter $L_0$ (distinct from the constant
$L_0(N)$ of this book). Then the following hold.
\begin{enumerate}
\item[(i)] The table exceeds the level-$k$ unit by the factor $\Gamma_\Lambda$:
\begin{equation}\label{9.1:eq-graded-thresholds-2}
\mathsf{T}^a_N(x)\ \ge\ \Gamma_\Lambda\,L^{(e_a-1)(k-m(x))}\,\mathsf{u}^a_k(x)
\ \ge\ \Gamma_\Lambda\,\mathsf{u}^a_k(x).
\end{equation}
\item[(ii)] Let $\square\ni x$ be a cube of the size $CM\ell_m$ with $\square\subset\Omega''_m$, as in
the cube clause of Definition~\ref{9.1:def-last-stage-table}, with $A$ the field of a gauge and
$B$, $C$, $M$ as in that cube clause. The cube clause at $x$ is implied
by any gauge on a cube $\square'\supseteq\square$ of the size $CM\ell_k$ with
\begin{equation}\label{9.1:eq-graded-thresholds-3}
\ell_k|A|,\ \ell_k^2|\nabla A|\ \le\ BCM\alpha_{0,k},
\end{equation}
with the margin $\Gamma_\Lambda$: on $\square$ such a gauge satisfies
\begin{equation}\label{9.1:eq-graded-thresholds-4}
\ell_m|A|,\ \ell_m^2|\nabla A|\ \le\
\Gamma_\Lambda^{-1}\,\ell_\circ^{2(m-k_0)_+}\,BCM\alpha_{0,m}.
\end{equation}
\end{enumerate}
The number $\Gamma_\Lambda$ is free of $k$, $h$, the stage index $N$ of $\mathsf{T}^a_N$ and $m$, and
$\Gamma_\Lambda\to\infty$ as $\gamma\to0$ with $L$, $\ell_\circ$ and $r_{pr}$ fixed; indeed
\begin{equation}\label{9.1:eq-graded-thresholds-5}
\ell_\circ^{2N_0}\ \ge\ (\log x_k)^{2r_{pr}\log_L\ell_\circ}.
\end{equation}
\end{lemma}

\begin{proof}
Fix $x\in Z''_k$ and the member $a$. Then $m=m(x)\le k-1$; put $s:=k-m\ge1$. We use
$k-k_0=N_0$ (Notation~\ref{9.1:not-graded-units}), so that $(m-k_0)_+=(N_0-s)_+$, and
$s\ge N_0$ exactly when $m\le k_0$. Define
\begin{equation}\label{9.1:eq-graded-thresholds-6}
f(s):=s^{-1/2}\,L^{s}\,\ell_\circ^{2(N_0-s)_+}.
\end{equation}

\emph{The lower bound $f(s)\ge\ell_\circ^{2N_0}$ for all $s\ge1$.} Suppose first $s\ge N_0$, that
is $m\le k_0$. Then $f(s)=s^{-1/2}L^s$, and $f$ is increasing in $s$ there, because
\begin{equation*}
\frac{f(s+1)}{f(s)}=L\Bigl(\frac{s}{s+1}\Bigr)^{1/2}\ \ge\ \frac{L}{\sqrt2}\ >\ 1,
\end{equation*}
where the last inequality uses the condition $L\ge 3\ell_\circ^2$ on the blocking factor, which
gives $L\ge3$. Hence, using $L\ge3\ell_\circ^2$ and $3^{N_0}\ge N_0^{1/2}$,
\begin{equation*}
f(s)\ \ge\ f(N_0)=N_0^{-1/2}L^{N_0}\ \ge\ N_0^{-1/2}\,3^{N_0}\,\ell_\circ^{2N_0}\ \ge\ \ell_\circ^{2N_0}.
\end{equation*}
Suppose next $1\le s\le N_0-1$ (this range is empty if $N_0=1$). Then, again by
$L\ge3\ell_\circ^2$ and by $s^{-1/2}3^s\ge1$ for $s\ge1$,
\begin{equation*}
f(s)=s^{-1/2}\,\ell_\circ^{2N_0}\Bigl(\frac{L}{\ell_\circ^2}\Bigr)^{s}
\ \ge\ s^{-1/2}\,3^{s}\,\ell_\circ^{2N_0}\ \ge\ \ell_\circ^{2N_0}.
\end{equation*}
In both cases $f(s)\ge\ell_\circ^{2N_0}$.

\emph{Proof of \textup{(i)}.} By the formula for $\mathsf{T}^a_N$ in
Definition~\ref{9.1:def-last-stage-table},
\begin{equation*}
\frac{\mathsf{T}^a_N(x)}{\mathsf{u}^a_k(x)}
=\varphi^{(N)}_m\,\ell_\circ^{2(m-k_0)_+}\,\frac{\alpha_{a,m}}{\alpha_{a,k}}
\Bigl(\frac{\ell_k}{\ell_m}\Bigr)^{e_a},
\qquad \frac{\ell_k}{\ell_m}=L^{k-m}=L^{s}.
\end{equation*}
The ratio bound for the radii (Notation~\ref{9.1:not-graded-units}) gives
$\alpha_{a,m}/\alpha_{a,k}\ge\frac49(k-m)^{-1/2}=\frac49 s^{-1/2}$, and
Definition~\ref{9.1:def-last-stage-table} gives $\varphi^{(N)}_m\ge\varphi_{\min}=\frac15$.
Inserting these bounds and $(m-k_0)_+=(N_0-s)_+$, and then the lower bound on $f$,
\begin{equation*}
\frac{\mathsf{T}^a_N(x)}{\mathsf{u}^a_k(x)}
\ \ge\ \frac15\cdot\frac49\,f(s)\,L^{(e_a-1)s}
\ \ge\ \frac{4}{45}\,\ell_\circ^{2N_0}\,L^{(e_a-1)(k-m)}
=\Gamma_\Lambda\,L^{(e_a-1)(k-m)}.
\end{equation*}
This is the first inequality of \eqref{9.1:eq-graded-thresholds-2}. The second follows because
$e_a\ge1$ for every member $a$ (the exponents $e_a$ being fixed in
Notation~\ref{9.1:not-graded-units}), so that $L^{(e_a-1)(k-m)}\ge1$.

\emph{Proof of \textup{(ii)}.} Restrict the gauge from $\square'$ to $\square\subset\square'$. By
$\ell_m=L^{-(k-m)}\ell_k=L^{-s}\ell_k$ and \eqref{9.1:eq-graded-thresholds-3}, on $\square$
\begin{equation*}
\ell_m|A|=L^{-s}\,\ell_k|A|\le L^{-s}\,BCM\alpha_{0,k},
\qquad
\ell_m^2|\nabla A|=L^{-2s}\,\ell_k^2|\nabla A|\le L^{-2s}\,BCM\alpha_{0,k}.
\end{equation*}
Since $L^{-2s}\le L^{-s}$, both left members are at most $L^{-s}BCM\alpha_{0,k}$. Their ratio to
$\ell_\circ^{2(m-k_0)_+}BCM\alpha_{0,m}$ is therefore at most
\begin{equation*}
\frac{\alpha_{0,k}}{\alpha_{0,m}}\,L^{-s}\,\ell_\circ^{-2(N_0-s)_+}
\ \le\ \frac94\,s^{1/2}L^{-s}\,\ell_\circ^{-2(N_0-s)_+}
=\frac{9/4}{f(s)},
\end{equation*}
by the ratio bound for the radii with $a=0$. By the lower bound on $f$,
\begin{equation*}
\frac{9/4}{f(s)}\ \le\ \frac94\,\ell_\circ^{-2N_0}\ \le\ \frac{45}{4}\,\ell_\circ^{-2N_0}
=\Gamma_\Lambda^{-1},
\end{equation*}
which is \eqref{9.1:eq-graded-thresholds-4}.

The gauge used here is available on the cubes on which the cube clause is imposed. The cube
families of \cite[(1.65), (1.67), (1.69)]{B-CMP122a} consist of cubes of the size
$CML^m\eta$, $\eta$ the lattice spacing of \cite{B-CMP122a}, that is of the size $CM\ell_m$ in the
present units, in the regions stated there; and $C\le2$ by \cite[\S1]{B-CMP122a}. A cube of the
size $CM\ell_m$ inside $Z''_k$ lies in a cube of the size $CM\ell_k$ inside the region
$\Lambda^{\sim}$ of the $\mathbf{R}$-step, and $\Lambda^{\sim}$ is contained in the box on which
the cube clause of the output $\Phi$ of the $\mathbf{R}$-step holds; by the theorem on the output
of the $\mathbf{R}$-step, the cubes on which later stages impose their cube clauses are sub-cubes
of those cubes, and the gauge is restricted to them.

\emph{Growth of $\Gamma_\Lambda$.} Let $\check{R}_{k_0+1}$ be the scaffold radius at level
$k_0+1$ (Ba{\l}aban's $R_{k_0+1}$). The identity $L^{N_0-1}=\check{R}_{k_0+1}$ is taken from
\cite{B-CMP122a}, the bound $\check{R}_{k_0+1}\ge(\log x_{k_0+1})^{r_{pr}}$ from
\cite[(2.5)]{B-CMP119}, and the last step below compares the couplings at the levels $k_0+1$
and $k$ as in \cite[(2.6)]{B-CMP119}:
\begin{equation*}
L^{N_0-1}=\check{R}_{k_0+1}\ \ge\ (\log x_{k_0+1})^{r_{pr}}\ \ge\ (\log x_k)^{r_{pr}}.
\end{equation*}
Since Ba{\l}aban's integer $\ell_\circ$ of \cite{B-CMP122a} satisfies $\ell_\circ>1$, the exponent
$2\log_L\ell_\circ$ is positive, and raising to this power gives
\begin{equation*}
\ell_\circ^{2N_0}\ \ge\ \ell_\circ^{2(N_0-1)}
=\bigl(L^{N_0-1}\bigr)^{2\log_L\ell_\circ}
\ \ge\ (\log x_k)^{2r_{pr}\log_L\ell_\circ},
\end{equation*}
which is \eqref{9.1:eq-graded-thresholds-5}. As $\gamma\to0$ with $L$, $\ell_\circ$ and $r_{pr}$
fixed, $x_k\to\infty$, so the right member tends to infinity.

\emph{Independence.} The definition \eqref{9.1:eq-graded-thresholds-1} involves only $\ell_\circ$
and $N_0$. The factor $(k-m)^{-1/2}$ of the ratio bound for the radii is absorbed by
$L^{k-m}$ into $f(s)$, and the bound $f(s)\ge\ell_\circ^{2N_0}$ holds for every $s\ge1$; so the
infimum over $m$ is taken after this absorption, and nothing in $\Gamma_\Lambda$ depends on $k$,
$h$, $N$ or $m$.
\end{proof}

Lemma~\ref{9.1:lem-graded-thresholds} is the form, for the last-stage table of
Definition~\ref{9.1:def-last-stage-table} and the ratio bound for the radii, of the comparison
according to which every threshold imposed on the region of the later stages exceeds the
level-$k$ unit by a factor of the order of $\ell_\circ^{2(N_0-1)}$, and of the corresponding
comparison for gauges. It compares the table with the level-$k$ unit; it makes no assertion about
the rotations along the comparison chain of Definition~\ref{9.6:def-comparison-chain}.

The second half of the lemma below uses the inequality $\ell_\circ^{2N_0}\ge360$. This inequality is implied by the ceiling on $\gamma$ that the order of choice (Notation~\ref{2.3:not-stages-last}) imposes before $\gamma$ is fixed. The left member of that ceiling contains the factor $\ell_\circ^{-2N_0}$. By Lemma~\ref{9.1:lem-graded-thresholds}, $\ell_\circ^{2N_0}\ge(\log x_k)^{2r_{pr}\log_L\ell_\circ}$, so that left member is at most a constant multiple of a negative power of $\log x_k$.

\begin{lemma}[The root seed's slice in the graded region]\label{9.1:lem-root-slice}
Let $\Phi$ be an output of Ba{\l}aban's operation $\mathbf{R}$ at the present step, at level $k$. Let $q_0=e^{i\eta A_e}U_c$ be a root seed of $\Phi$ in the sense of Definition~\ref{9.1:def-aligned-points}, so that $q_0$ is a datum of the seed tube $\mathfrak{T}_\Phi[7/8]$. Let $(F,X,j)$ be a term as in Lemma~\ref{9.1:lem-site-blind-slice} with $X^{+}\cap\mathfrak{K}=\emptyset$, where $\mathfrak{K}$ is the core; in particular $X\cap\mathfrak{K}=\emptyset$. Then for every $x\in Z\cap X$ and every table member $a$,
\begin{equation}\label{9.1:eq-root-slice-1}
E_a(S(q_0))(x)<\bigl(1-\tfrac18\theta_\rho\beta\bigr)\mathsf{u}^a_k(x)\le\mathsf{u}^a_k(x).
\end{equation}
Moreover, $E_5=|\partial U_c-1|$ and the cube gauges of $U_c$ are the same at every object $q$ of the chain over $q_0$.

Assume in addition $\ell_\circ^{2N_0}\ge360$. Then the following hold.
\begin{itemize}
\item For every $x\in Z''_k\cap X$ and every $a$ we have $E_a(S(q_0))(x)\le\mathsf{T}^a_N(x)/32$.
\item If $x$ is moreover a point of class (b) of Definition~\ref{9.1:def-aligned-points}, then also $E_a(S(q_0))(x)\le\mathsf{u}^a_F(x)/32$.
\item $U_c$ satisfies the cube clause of Definition~\ref{9.1:def-last-stage-table} at every point of $Z''_k\cap X$ with margin at least $32$. That is, on each cube of that clause a gauge exists in which $\ell_m|A|$ and $\ell_m^2|\nabla A|$ are at most $32^{-1}\ell_\circ^{2(m-k_0)_+}BCM\alpha_{0,m}$.
\end{itemize}
\end{lemma}

\begin{proof}
\emph{The first display.} By part (a) of the slice decomposition used in Lemma~\ref{9.1:lem-site-blind-slice}, writing $q_0$ in the coordinates $(U_c,A_e)$ and letting $U_e$ denote the configuration of Lemma~\ref{9.1:lem-site-blind-slice} whose slice on $X$ is $q_0$, with $J(U_e)$ its current, we have
\[
S(q_0)=(U_e,J(U_e)){\upharpoonright}X=q_0{\upharpoonright}X .
\]
By Definition~\ref{9.1:def-aligned-points}, a seed of an output $\Phi$ is a datum of $\mathfrak{T}_\Phi[7/8]$ and has its slice in $\mathfrak{D}_\Phi$; thus we obtain
\[
(q_0,J(q_0)){\upharpoonright}X\in\mathfrak{D}_\Phi[7/8]
=\mathfrak{U}\bigl[(\mathsf{c}_\Phi^+-\tfrac18\varrho_\Phi)\mathsf{u}_\Phi\bigr](X).
\]
Membership in this domain bounds all seven members $E_a$ and includes the cube clause of $U_c$. This is the same reading as in hypothesis (E) of Lemma~\ref{9.1:lem-site-blind-slice}.

On $Z$ the units of $\Phi$ are those of level $k$, and $\mathsf{c}_\Phi^+=1$ there, by Definition~\ref{9.1:def-aligned-points}. The room schedule $\varrho_F=\theta_\rho\beta\,2^{-(j-n_F)}$, where $n_F$ denotes the level of the units of the term $F$, taken at $j=n_\Phi=k$, gives $\varrho_\Phi=\theta_\rho\beta$. This yields the first inequality of \eqref{9.1:eq-root-slice-1}. The second follows from $\theta_\rho\beta\ge0$.

\emph{Units of the stage-$N$ table.} Let $x\in Z''_k\cap X$. By the first inequality of Lemma~\ref{9.1:lem-graded-thresholds}, $\mathsf{T}^a_N(x)\ge\Gamma_\Lambda\mathsf{u}^a_k(x)$ with $\Gamma_\Lambda=\frac4{45}\ell_\circ^{2N_0}$. Hence
\[
\mathsf{u}^a_k\le\Gamma_\Lambda^{-1}\mathsf{T}^a_N
=\tfrac{45}{4}\ell_\circ^{-2N_0}\mathsf{T}^a_N\le\tfrac1{32}\mathsf{T}^a_N ,
\]
where the last step uses $\ell_\circ^{2N_0}\ge360=\frac{45}{4}\cdot32$. Together with \eqref{9.1:eq-root-slice-1} this gives $E_a(S(q_0))(x)\le\mathsf{T}^a_N(x)/32$.

\emph{Units of $F$ at points of class (b).} At such a point the term is old in the sense of Definition~\ref{9.1:def-aligned-points}: $n_F\le k_0$, and $u:=k-n_F\ge N_0$. The units of $F$ are those of level $n_F$. The comparison bound for the coefficients $\alpha_{a,n}$ of the units $\mathsf{u}^a_n=\alpha_{a,n}\ell_n^{-e_a}$ (Definition~\ref{9.1:def-last-stage-table}) states that $(\alpha_{a,k}/\alpha_{a,n_F})L^{-e_au}\le\frac94u^{1/2}L^{-u}$; this is the first estimate below:
\begin{align*}
\frac{\mathsf{u}^a_k}{\mathsf{u}^a_F}
&=\frac{\alpha_{a,k}}{\alpha_{a,n_F}}\,L^{-e_au}
\le\tfrac94\,u^{1/2}L^{-u}
\le\tfrac94\,N_0^{1/2}L^{-N_0}\\
&\le\tfrac94\,\ell_\circ^{-2N_0}\le\tfrac1{32}.
\end{align*}
The three later steps are justified as follows.
\begin{itemize}
\item Second step. The ratio of consecutive values of $u^{1/2}L^{-u}$ equals $((u+1)/u)^{1/2}L^{-1}$, which is at most $\sqrt2/L<1$. So the function decreases for $u\ge1$, and $u\ge N_0$.
\item Third step. Since $L\ge3\ell_\circ^2$, we have $L^{N_0}\ge(3\ell_\circ^2)^{N_0}=3^{N_0}\ell_\circ^{2N_0}$. Moreover $3^{N_0}\ge N_0^{1/2}$, so $L^{N_0}\ge N_0^{1/2}\ell_\circ^{2N_0}$.
\item Last step. The inequality $\frac94\le\frac1{32}\ell_\circ^{2N_0}$ is equivalent to $\ell_\circ^{2N_0}\ge72$, which holds since $\ell_\circ^{2N_0}\ge360$.
\end{itemize}
Combining with \eqref{9.1:eq-root-slice-1} gives $E_a(S(q_0))(x)\le\mathsf{u}^a_F(x)/32$.

\emph{$E_5$ and the cube clause.} Part (iv) of the theorem on the chain over the root seed (the chain of Definition~\ref{9.1:def-aligned-points}) states that $E_5$ and the cube clause at every chain object are those of the entry $q_0$. Part (b) of the slice decomposition used in Lemma~\ref{9.1:lem-site-blind-slice} states that for every move $\mathcal{Y}$ the slice $S(e^{i\eta\mathcal{Y}}U_e)$ has unitary part $U_c$; only the Landau coordinate and the current are shifted. Thus the unitary part of $S(q)$ is $U_c$ at every object $q$ of the chain over $q_0$. Consequently neither $E_5=|\partial U_c-1|$ nor the gauges of $U_c$ change along the chain.

The cube gauges of $U_c$ given by membership in $\mathfrak{D}_\Phi[7/8]$ are those of level $k$. On each cube of size $CM\ell_k$ of the cube clause of $\mathfrak{D}_\Phi[7/8](X)$ there is a gauge with
\[
\ell_k|A|,\ \ell_k^2|\nabla A|\le\bigl(1-\tfrac18\theta_\rho\beta\bigr)BCM\alpha_{0,k}\le BCM\alpha_{0,k}.
\]
These are gauges of the kind required in the second assertion of Lemma~\ref{9.1:lem-graded-thresholds}. That assertion transfers them to the stage-$N$ cubes of size $CM\ell_m$ at $x$ ($m=m(x)$), with margin $\Gamma_\Lambda=\frac4{45}\ell_\circ^{2N_0}$. This margin is at least $32$, since $\ell_\circ^{2N_0}\ge360$ and $\frac4{45}\cdot360=32$.
\end{proof}

\begin{lemma}[Top-level rotations are small in the graded region]\label{9.1:lem-top-rotations}
Let $X$ be a region with $X^{+}\cap\mathfrak{K}=\emptyset$, where $\mathfrak{K}$ is the core region.
Let $\mathbf{U}_1=e^{i\eta A_1}U_c$ be a point of the chart on $X^{+}$, in the notation of
Lemma~\ref{9.1:lem-noncritical-displacement}, and let $\mathbf{U}_2=e^{i\eta\mathcal{Y}}\mathbf{U}_1$.
Assume the hypotheses of Lemma~\ref{9.1:lem-noncritical-displacement},
\begin{equation*}
\ell_{\mathfrak{S}}\sup\bigl(|A_1|+|\mathcal{Y}|\bigr)\le c_4,\qquad |B[\cdot]|\le a_{R'},
\end{equation*}
where $c_4$ is the chart radius of that lemma;
and assume a level-$k$ absolute majorant: for a given number $\mathsf{m}\ge0$,
\begin{equation}\label{9.1:eq-top-rotations-1}
\sup_{\tilde{\Delta}(y)}|\mathcal{Y}|\le \frac{\mathsf{m}}{\ell_k}
\qquad\text{for all } y\in\mathfrak{B},
\end{equation}
where $\mathfrak{B}$ is the set of cells $y$ over which the smeared majorant $(\cdot)^{\approx}$
is summed and $\tilde{\Delta}(y)$ is the enlarged cell of $y$.
Then at every $x\in Z''^{(1)}_k\cap X$ the $a$-th member of $S(\mathbf{U}_2)-S(\mathbf{U}_1)$,
which we write $\delta E_a(x)$, is at most
\begin{equation}\label{9.1:eq-top-rotations-2}
16B_{\natural}c_1C_T\,\mathsf{m}\,\ell_{\mathfrak{S}}(x)^{1-e_a}\ell_k^{-1},
\end{equation}
and hence, with the tables and units of Definitions~\ref{9.1:def-last-stage-table}
and~\ref{9.1:def-aligned-points}:
\begin{itemize}
\item[(a)] at aligned points $x$, the member is at most
$180B_{\natural}c_1C_T(\mathsf{m}/\alpha_{a,k})\,\ell_\circ^{-2N_0}\,\mathsf{T}^a_N(x)$;
\item[(b)] at older points $x$, the member is at most
$36B_{\natural}c_1C_T(\mathsf{m}/\alpha_{a,k})\,\ell_\circ^{-2N_0}\,\mathsf{u}^a_F(x)$.
\end{itemize}
In particular, let $\mathsf{m}=L^4C_4^{\sharp}\mathsf{s}_L$, where $C_4^{\sharp}$ is the constant
fixed in Definition~\ref{9.1:def-aligned-points} and
$\mathsf{s}_L=K_\Lambda'(1+\mathsf{r}_1)\alpha_{0,k}$; here $\alpha_{a,k}\ge\alpha_{0,k}$.
Assume the ceiling condition on the top-level seed size in the form
\begin{equation}\label{9.1:eq-top-rotations-3}
2^{20}B_{\natural}c_1C_T\,L^4C_4^{\sharp}K_\Lambda'(1+\mathsf{r}_1)\le \ell_\circ^{2N_0}.
\end{equation}
Then the bounds in (a) and (b) are at most $180\cdot2^{-20}\mathsf{T}^a_N(x)$ and
$36\cdot2^{-20}\mathsf{u}^a_F(x)$ respectively; in particular both are at most $\tfrac{1}{2880}$ of
$\mathsf{T}^a_N(x)$, respectively of $\mathsf{u}^a_F(x)$.
\end{lemma}

\begin{proof}
\emph{The first bound.} By Lemma~\ref{9.1:lem-noncritical-displacement}, which applies under the
hypotheses assumed, the members of $S(\mathbf{U}_2)-S(\mathbf{U}_1)$ are at most
$4B_{\natural}(|\mathcal{Y}|_{\cdot})^{\approx}$ in the weights $\ell_{\mathfrak{S}}|E|$,
$\ell_{\mathfrak{S}}^2|\nabla E|$, $\ell_{\mathfrak{S}}^3|\mathsf{J}|$, where $E$ is the
field-difference coordinate and $\mathsf{J}$ the current coordinate in which that lemma weights the
members. Moreover the member $E_1$
moves by at most $2\eta^2|\nabla E|(1+C\varepsilon_3)$, with the constant $C$ and the small
parameter $\varepsilon_3$ of that lemma. From these two facts, arguing as in the
proof of Lemma~\ref{9.1:lem-chain-slices}, where the same move is bounded, the move of $E_a$ at
$x$ satisfies
\begin{equation}\label{9.1:eq-top-rotations-4}
\delta E_a(x)\le 16B_{\natural}\,\ell_{\mathfrak{S}}(x)^{-e_a}\,(|\mathcal{Y}|_{\cdot})^{\approx}(x).
\end{equation}
For composite members the same bound follows from the clause on composites of
Lemma~\ref{9.1:lem-noncritical-displacement}.

We now bound the smeared majorant. By definition $|\mathcal{Y}|_y=\ell_{\mathfrak{S}}(y)
\sup_{\tilde{\Delta}(y)}|\mathcal{Y}|$, so \eqref{9.1:eq-top-rotations-1} gives
$|\mathcal{Y}|_y\le\ell_{\mathfrak{S}}(y)\mathsf{m}/\ell_k$. Hence
\begin{align*}
(|\mathcal{Y}|_{\cdot})^{\approx}(x)
&=\sum_{y\in\mathfrak{B}}e^{-2\delta_1 d_{\mathfrak{S}}(x,y)}|\mathcal{Y}|_y
\le\frac{\mathsf{m}}{\ell_k}\sum_{y\in\mathfrak{B}}e^{-2\delta_1 d_{\mathfrak{S}}(x,y)}
\ell_{\mathfrak{S}}(y)\\
&\le c_1C_T\,\ell_{\mathfrak{S}}(x)\,\frac{\mathsf{m}}{\ell_k}.
\end{align*}
The last inequality uses four facts:
\begin{itemize}
\item the comparison of the local lengths,
$\ell_{\mathfrak{S}}(y)\le\ell_{\mathfrak{S}}(x)L^{|\lambda_{\mathfrak{S}}(x)-\lambda_{\mathfrak{S}}(y)|}$;
\item the comparison of level differences with distance,
$L^{|\lambda_{\mathfrak{S}}(x)-\lambda_{\mathfrak{S}}(y)|}\le L^2e^{\delta_1d_{\mathfrak{S}}(x,y)/8}$,
valid under the floor condition under which that comparison is stated;
\item the summation bound
$\sum_{y}e^{-\delta_1 d_{\mathfrak{S}}(x,y)/2}\le c_1$;
\item the summation constant of the random-walk estimates, $C_T=c_1L^3$.
\end{itemize}
Combined, they give
\begin{align*}
\sum_{y\in\mathfrak{B}}e^{-2\delta_1 d_{\mathfrak{S}}(x,y)}\ell_{\mathfrak{S}}(y)
&\le L^2\ell_{\mathfrak{S}}(x)\sum_{y\in\mathfrak{B}}e^{-(2-1/8)\delta_1 d_{\mathfrak{S}}(x,y)}
\le L^2\ell_{\mathfrak{S}}(x)\sum_{y\in\mathfrak{B}}e^{-\delta_1 d_{\mathfrak{S}}(x,y)/2}\\
&\le c_1L^2\,\ell_{\mathfrak{S}}(x)\le c_1C_T\,\ell_{\mathfrak{S}}(x),
\end{align*}
the last step since $C_T=c_1L^3\ge L^2$.
Inserting this into \eqref{9.1:eq-top-rotations-4} gives \eqref{9.1:eq-top-rotations-2}. This is
the bound obtained on the second of the three regions in the proof of
Lemma~\ref{9.1:lem-chain-slices},
with the level there replaced by $k$ and the majorant profile there replaced by the number
$\mathsf{m}$.

\emph{Part (a).} At an aligned point, $\ell_{\mathfrak{S}}(x)=\ell_m$ with $m=m(x)$; put $s:=k-m$.
By Definition~\ref{9.1:def-last-stage-table},
$\mathsf{T}^a_N(x)=\varphi^{(N)}_m\ell_\circ^{2(m-k_0)_+}\alpha_{a,m}\ell_m^{-e_a}$. Dividing
\eqref{9.1:eq-top-rotations-2} by this gives
\begin{equation*}
\frac{\delta E_a(x)}{\mathsf{T}^a_N(x)}
\le 16B_{\natural}c_1C_T\,\frac{\mathsf{m}}{\alpha_{a,m}}\,
\frac{\ell_m/\ell_k}{\varphi^{(N)}_m\ell_\circ^{2(m-k_0)_+}}.
\end{equation*}
We use $\varphi^{(N)}_m\ge\varphi_{\min}=1/5$ (Definition~\ref{9.1:def-last-stage-table}),
$\ell_m/\ell_k=L^{-s}$, and the comparison of the $\ell_\circ$-weights made in the proof of
Lemma~\ref{9.1:lem-graded-thresholds}, where the same quotient is estimated, which gives
$\ell_\circ^{-2(m-k_0)_+}\le\ell_\circ^{-2(N_0-s)_+}$.

If $s=0$, then $m=k$ and $\alpha_{a,m}=\alpha_{a,k}$, so the right side is at most
$16\cdot5\,B_{\natural}c_1C_T(\mathsf{m}/\alpha_{a,k})\,\ell_\circ^{-2N_0}$, which is at most the
bound in (a).

If $s\ge1$, we use in addition the growth bound for the radii $\alpha_{a,\cdot}$ across levels,
$\alpha_{a,k}\le\tfrac94 s^{1/2}\alpha_{a,m}$ over the $s$ levels from $m$ to $k$, and obtain
\begin{equation*}
\frac{\delta E_a(x)}{\mathsf{T}^a_N(x)}
\le 16\cdot5\cdot\tfrac94\,B_{\natural}c_1C_T\,\frac{\mathsf{m}}{\alpha_{a,k}}\,
s^{1/2}L^{-s}\ell_\circ^{-2(N_0-s)_+}
=\frac{180B_{\natural}c_1C_T\,\mathsf{m}/\alpha_{a,k}}{f(s)},
\end{equation*}
where $f(s):=s^{-1/2}L^{s}\ell_\circ^{2(N_0-s)_+}$. The proof of
Lemma~\ref{9.1:lem-graded-thresholds} also gives $f(s)\ge\ell_\circ^{2N_0}$. This yields (a).

\emph{Part (b).} At an older point, by Definition~\ref{9.1:def-aligned-points},
$\ell_{\mathfrak{S}}(x)=\ell_j$ with $n_F\le j<m=m(x)$, and $u:=k-n_F\ge N_0$. The unit there is
$\mathsf{u}^a_F=\alpha_{a,n_F}\ell_{n_F}^{-e_a}$. Dividing \eqref{9.1:eq-top-rotations-2} by it gives
\begin{align*}
\frac{\delta E_a(x)}{\mathsf{u}^a_F(x)}
&\le16B_{\natural}c_1C_T\,\frac{\mathsf{m}}{\alpha_{a,n_F}}
\Bigl(\frac{\ell_{n_F}}{\ell_j}\Bigr)^{e_a-1}\frac{\ell_{n_F}}{\ell_k}
\le 16B_{\natural}c_1C_T\,\frac{\mathsf{m}}{\alpha_{a,k}}\,
\frac{\alpha_{a,k}}{\alpha_{a,n_F}}\,L^{-u}\\
&\le 36B_{\natural}c_1C_T\,\frac{\mathsf{m}}{\alpha_{a,k}}\,u^{1/2}L^{-u}
\le 36B_{\natural}c_1C_T\,\frac{\mathsf{m}}{\alpha_{a,k}}\,\ell_\circ^{-2N_0}.
\end{align*}
In the first step, $(\ell_{n_F}/\ell_j)^{e_a-1}\le1$ because $n_F\le j$ and $e_a\ge1$, and
$\ell_{n_F}/\ell_k=L^{-u}$. The second step uses $\alpha_{a,k}\le\tfrac94u^{1/2}\alpha_{a,n_F}$,
the growth bound for the radii $\alpha_{a,\cdot}$ across levels used in part (a), and
$16\cdot\tfrac94=36$.

For the third step, $u\mapsto u^{1/2}L^{-u}$ is decreasing for $u\ge1$, since the ratio of
consecutive values is $((u+1)/u)^{1/2}L^{-1}\le\sqrt2/L<1$. Hence for $u\ge N_0$ it is at most
$N_0^{1/2}L^{-N_0}$. From the lower bound $L\ge3\ell_\circ^2$, one of the one-sided lower bounds on
$L$ in the order of choice, and $3^{N_0}\ge N_0^{1/2}$ we get
\begin{equation*}
L^{N_0}\ge(3\ell_\circ^2)^{N_0}\ge N_0^{1/2}\ell_\circ^{2N_0},
\end{equation*}
so $N_0^{1/2}L^{-N_0}\le\ell_\circ^{-2N_0}$. This yields (b).

\emph{The particular case.} Since $\alpha_{a,k}\ge\alpha_{0,k}$,
\begin{equation*}
\frac{\mathsf{m}}{\alpha_{a,k}}\le\frac{L^4C_4^{\sharp}K_\Lambda'(1+\mathsf{r}_1)\alpha_{0,k}}
{\alpha_{0,k}}=L^4C_4^{\sharp}K_\Lambda'(1+\mathsf{r}_1).
\end{equation*}
By \eqref{9.1:eq-top-rotations-3}, the factor
$B_{\natural}c_1C_T(\mathsf{m}/\alpha_{a,k})\ell_\circ^{-2N_0}$ in (a) and (b) is at most $2^{-20}$.
Hence the bound in (a) is at most $180\cdot2^{-20}\mathsf{T}^a_N(x)$ and the bound in (b) is at
most $36\cdot2^{-20}\mathsf{u}^a_F(x)$. Since $180\cdot2880=518400<2^{20}$ and $36<180$, both are at
most $\tfrac{1}{2880}$ times $\mathsf{T}^a_N(x)$, respectively $\mathsf{u}^a_F(x)$.

The bounds of Lemma~\ref{9.1:lem-top-rotations}, like the bound for the entry of the chain on
$Z''_k$, are the bounds of Lemma~\ref{9.1:lem-chain-slices} on the second of its three regions
moved inside $Z''_k$. The exponential decay factor available on that region is absent here.
Instead, the units
of every older term lie at level at most $k_0$, so that $k-n_F\ge N_0$, and for terms of level
$k$ every point is an aligned point.
\end{proof}

\begin{remark}[The quarter-threshold condition does not follow from the chain bounds]
\label{9.1:rem-quarter-threshold}
We use the following notation. The graded region is $Z''_k$. A point $x$ of $Z''_k$ has a graded level $m=m(x)$, and $\ell_m$ is the length scale of that level. A \emph{bulk point} is a point with $m=k_0$. The label scale of a chain object is $\ell_\iota$, and the scale of the box set $\mathfrak{S}$ is $\ell_{\mathfrak{S}}$. At a point of class~(a) one has $\ell_{\mathfrak{S}}=\ell_\iota=\ell_m$; at a point of class~(b), $\ell_{\mathfrak{S}}=\ell_j<\ell_m$. A third class of points of the same classification is denoted $(\mathrm{up})$.

For a table member $E_a$, the number $e_a$ is its scaling exponent. Its radius at level $n$ is $\alpha_{a,n}$, with $\alpha_{a,k_0}=\mathsf{r}_1\alpha_{0,k_0}$ for $a=6,7$ and $\alpha_{3,k_0}=\alpha_{0,k_0}$. The seed size is $\mathsf{s}_L=K_\Lambda'(1+\mathsf{r}_1)\alpha_{0,k}$.

The last-stage table $\mathsf{T}^a_N$ is that of Definition~\ref{9.1:def-last-stage-table}. The graded-region condition of Definition~\ref{9.1:def-graded-ceilings}, which we call the quarter-threshold condition, requires, for every leaf $X$ missing the raw core, $E_a(S(q))\le\frac14\mathsf{T}^a_N$ on $Z''_k\cap X$ for every chain object $q$.
\begin{itemize}
\item[(i)] Let $x\in Z''_k$ be a point of class~(a). Bound the move of $E_a$ at $x$ by the link-by-link bound (Lemma~\ref{9.1:lem-link-bound}) together with the slice displacement between non-critical points (Lemma~\ref{9.1:lem-noncritical-displacement}). At every bulk point, the resulting majorant exceeds $\frac14\mathsf{T}^a_N(x)$ by a factor at least $32B_\natural c_1C_4K_\Lambda'\ge 32$ for $a=6,7$, and at least $32B_\natural c_1C_4K_\Lambda'(1+\mathsf{r}_1)\ge32(1+\mathsf{r}_1)$ for $a=3$. The lower bounds $32$ and $32(1+\mathsf{r}_1)$ contain none of $L$, $M_1$, $M$, $\kappa_1$, $N_0$, $\gamma$. Consequently the quarter-threshold condition is not derivable by these majorants from the link-by-link bound, the seed-size relation $\mathsf{s}_L=K_\Lambda'(1+\mathsf{r}_1)\alpha_{0,k}$ and the slice displacement bound.
\item[(ii)] At a point of class~(b), the move given by Lemma~\ref{9.1:lem-noncritical-displacement} exceeds the scale $\ell_m^{-e_a}$ of $\mathsf{T}^a_N$ by the factor $L^{(e_a-1)(m-j)}$ when $e_a\ge2$. This holds even when the chain field is small at level $k$. There the threshold $\frac14\mathsf{T}^a_N$ is measured in the wrong units. The appropriate units are those of the term $F$ itself, as in part~(b) of Lemma~\ref{9.1:lem-top-rotations}.
\end{itemize}
\end{remark}

\begin{proof}
\emph{Part} (i). Let $q$ be a chain object, written $q=e^{i\eta\mathcal{Z}_q}q_0$ as in Lemma~\ref{9.1:lem-link-bound}. That lemma gives, in label units,
\begin{equation*}
\ell_\iota|\mathcal{Z}_q|\le\mathsf{m}_L=C_4\mathsf{s}_L e^{-\delta_1 d(\cdot,\bar\Lambda)/4}.
\end{equation*}
Here $\delta_1$ is the decay rate and $d(\cdot,\bar\Lambda)$ the distance to the region $\bar\Lambda$ of Lemma~\ref{9.1:lem-link-bound}; this distance vanishes on $Z''_k$, so there $\ell_\iota|\mathcal{Z}_q|\le C_4\mathsf{s}_L$.

Now let $x\in Z''_k$ be a point of class~(a) with graded level $m$, so that $\ell_{\mathfrak{S}}=\ell_\iota=\ell_m$. By Lemma~\ref{9.1:lem-noncritical-displacement}, read on the equal-scale line of the slice bound (the case $\ell_{\mathfrak{S}}=\ell_\iota=\ell_m$), the move of $E_a$ at $x$ is at most $16B_\natural c_1\,\ell_m^{1-e_a}\sup|\mathcal{Z}_q|$. Inserting $\sup|\mathcal{Z}_q|\le C_4\mathsf{s}_L/\ell_m$, the majorant of the move is $16B_\natural c_1C_4\mathsf{s}_L\ell_m^{-e_a}$.

Definition~\ref{9.1:def-last-stage-table} gives
\begin{equation*}
\mathsf{T}^a_N(x)=\varphi^{(N)}_m\ell_\circ^{2(m-k_0)_+}\alpha_{a,m}\ell_m^{-e_a}.
\end{equation*}
Substituting this and $\mathsf{s}_L=K_\Lambda'(1+\mathsf{r}_1)\alpha_{0,k}$, the factor $\ell_m^{-e_a}$ cancels and we obtain
\begin{equation}\label{9.1:eq-quarter-threshold-1}
\frac{16B_\natural c_1C_4\mathsf{s}_L\ell_m^{-e_a}}{\tfrac14\mathsf{T}^a_N(x)}
=\frac{64B_\natural c_1C_4K_\Lambda'(1+\mathsf{r}_1)}{\varphi^{(N)}_m\ell_\circ^{2(m-k_0)_+}}
\cdot\frac{\alpha_{0,k}}{\alpha_{a,m}}.
\end{equation}

At a bulk point we have $m=k_0$, so the weight is $\ell_\circ^{0}=1$. We use four facts:
\begin{itemize}
\item for $a=6,7$, $\alpha_{a,k_0}=\mathsf{r}_1\alpha_{0,k_0}$;
\item $\alpha_{0,k_0}\le2\alpha_{0,k}$, by the comparison of the radii at levels $k_0$ and $k$;
\item $\varphi^{(N)}_m\le1$;
\item $(1+\mathsf{r}_1)/\mathsf{r}_1\ge1$.
\end{itemize}
Together they bound the right-hand side of \eqref{9.1:eq-quarter-threshold-1} below by
\begin{equation*}
\frac{64B_\natural c_1C_4K_\Lambda'(1+\mathsf{r}_1)\alpha_{0,k}}{\mathsf{r}_1\alpha_{0,k_0}}
\ge\frac{64B_\natural c_1C_4K_\Lambda'(1+\mathsf{r}_1)}{2\mathsf{r}_1}
\ge32B_\natural c_1C_4K_\Lambda'.
\end{equation*}
Since $B_\natural,c_1,K_\Lambda'\ge1$ and $C_4\ge16$, the last member is at least $32$.

For $a=3$ the radius is $\alpha_{3,k_0}=\alpha_{0,k_0}\le2\alpha_{0,k}$. The same steps, without the factor $\mathsf{r}_1$, give the lower bound $32B_\natural c_1C_4K_\Lambda'(1+\mathsf{r}_1)$, which is at least $32(1+\mathsf{r}_1)$ since $B_\natural,c_1,K_\Lambda'\ge1$ and $C_4\ge16$.

The lower bounds $32$ and $32(1+\mathsf{r}_1)$ contain none of $L$, $M_1$, $M$, $\kappa_1$, $N_0$, $\gamma$. Hence no choice of these parameters in the order of choice makes the ratio \eqref{9.1:eq-quarter-threshold-1} at most $1$. The implication from Lemma~\ref{9.1:lem-link-bound}, the seed-size relation $\mathsf{s}_L=K_\Lambda'(1+\mathsf{r}_1)\alpha_{0,k}$ and Lemma~\ref{9.1:lem-noncritical-displacement} to the quarter-threshold condition of Definition~\ref{9.1:def-graded-ceilings} therefore fails, as an implication, at every bulk point.

\emph{Part} (ii). Let $x$ be a point of class~(b), so that $\ell_{\mathfrak{S}}=\ell_j<\ell_m$. The re-solved slice carries the gradient $\nabla A'$ of its vector potential $A'$ and a current at the scale of $\mathfrak{S}$. The move of Lemma~\ref{9.1:lem-noncritical-displacement} is therefore $\ell_j^{1-e_a}\sup|\mathcal{Z}_q|$. Since $\mathsf{T}^a_N$ is proportional to $\ell_m^{-e_a}$, this move exceeds $\mathsf{T}^a_N$ by $L^{(e_a-1)(m-j)}$ for $e_a\ge2$, even when $\mathcal{Z}_q$ is small at level $k$. In other words, a level-$k$ increment re-solved at a finer level is not small in the units of $\mathsf{T}_N$; this is the case for $E_7$ and $E_3$.

So the threshold $\frac14\mathsf{T}^a_N$ is measured in the wrong units at points of class~(b). The right units are those of the term $F$ itself, in which part~(b) of Lemma~\ref{9.1:lem-top-rotations} measures the move, namely a factor $L^{-u}$ with $u\ge N_0$. These are the units already used at the points of class $(\mathrm{up})$.
\end{proof}

\begin{lemma}[The last-stage table against a leaf's weighted unit]\label{9.1:lem-weighted-unit}
Let $F$ be a leaf of level $n_F$ (Definition~\ref{9.1:def-aligned-points}), and use the units
$\mathsf{u}^a_n=\alpha_{a,n}\ell_n^{-e_a}$ of
Notation~\ref{9.1:not-graded-units} and the last-stage table $\mathsf{T}^a_N$ of
Definition~\ref{9.1:def-last-stage-table}. Define the \emph{weighted unit} of $F$, carrying the weight of
\cite[(1.64)(i)]{B-CMP122a}, by
\begin{equation}\label{9.1:eq-weighted-unit-1}
  \bar u^a_F:=\ell_\circ^{2(n_F-k_0)_+}\,\alpha_{a,n_F}\,\ell_{n_F}^{-e_a}.
\end{equation}
Since the leaf's own unit is $\mathsf{u}^a_F=\mathsf{u}^a_{n_F}=\alpha_{a,n_F}\ell_{n_F}^{-e_a}$, we have
$\bar u^a_F=\ell_\circ^{2(n_F-k_0)_+}\mathsf{u}^a_F$; in particular $\bar u^a_F=\mathsf{u}^a_F$ whenever
$n_F\le k_0$. This is the case for old leaves, which have $n_F\le k_0$, and for the leaves of the
chain terms $\mathbb{C}^{(i)}_k$ (Definition~\ref{9.1:def-aligned-points}) of the step at level $k$
with $i\le N-N_0$: their level is $n_F=h+i$ with $h:=k-N$, so that $n_F\le k-N_0=k_0$
(Notation~\ref{9.1:not-graded-units}).
At every point $x$ of the graded region with $m:=m(x)>n_F$ (these include the
$(\mathrm{up})$-$(a)$-points of $F$, at which $m(x)>n_F$),
\begin{equation}\label{9.1:eq-weighted-unit-2}
  \mathsf{T}^a_N(x)\le \tfrac34\,\bar u^a_F .
\end{equation}
If moreover $F$ is the leaf of a chain term $\mathbb{C}^{(i)}_k$ of the step at level $k$, so that
$n_F=h+i$, put
\begin{equation*}
  \mathsf{T}^a_i(x):=\varphi^{(i)}_{n_F}\,\ell_\circ^{2(n_F-k_0)_+}\,\mathsf{u}^a_{n_F}
  =\varphi^{(i)}_{n_F}\,\bar u^a_F ,
\end{equation*}
the value at $x$ of $F$'s stage-$i$ table, that is, of the table of
Definition~\ref{9.1:def-last-stage-table} at stage $i$ read at the level $n_F$. Then
\begin{equation}\label{9.1:eq-weighted-unit-3}
  \mathsf{T}^a_N(x)\le \mathsf{T}^a_i(x).
\end{equation}
\end{lemma}

\begin{proof}
We use the following facts about the units, radii and scaffold integers of
Notation~\ref{9.1:not-graded-units} and about the stage coefficients of
Definition~\ref{9.1:def-last-stage-table}; the inequality $3\ell_\circ^2\le L$ in (iii) is the
standing condition relating the integer $\ell_\circ$ and the blocking factor $L$:
\begin{itemize}
\item[(i)] $\alpha_{a,m}/\alpha_{a,n}\le \tfrac94(m-n)^{1/2}$ whenever $n<m$;
\item[(ii)] $\ell_m=L^{m-n}\ell_n$ for $n\le m$, $e_a\ge1$ and $\ell_\circ\ge1$;
\item[(iii)] $\varphi^{(N)}_m\le 1$ and $3\ell_\circ^2\le L$;
\item[(iv)] at the step at level $k$, the levels of the graded region satisfy $m(x)\le k-1$.
\end{itemize}
The integer $r$ below is local to this proof. Fix a point $x$ with $m:=m(x)>n_F$, and
put $r:=m-n_F$, so that $r\ge1$. By
Definition~\ref{9.1:def-last-stage-table},
$\mathsf{T}^a_N(x)=\varphi^{(N)}_m\ell_\circ^{2(m-k_0)_+}\alpha_{a,m}\ell_m^{-e_a}$. Dividing by
\eqref{9.1:eq-weighted-unit-1} gives
\begin{equation}\label{9.1:eq-weighted-unit-4}
  \frac{\mathsf{T}^a_N(x)}{\bar u^a_F}
  =\varphi^{(N)}_m\,\ell_\circ^{2[(m-k_0)_+-(n_F-k_0)_+]}\,
   \frac{\alpha_{a,m}}{\alpha_{a,n_F}}\,\Bigl(\frac{\ell_m}{\ell_{n_F}}\Bigr)^{-e_a}.
\end{equation}
We bound the four factors in turn.

Since $(\,\cdot\,-k_0)_+$ is nondecreasing and $1$-Lipschitz,
\begin{equation*}
  0\le (m-k_0)_+-(n_F-k_0)_+\le m-n_F=r .
\end{equation*}
Since $\ell_\circ\ge1$ by (ii), the second factor is therefore at most $\ell_\circ^{2r}$. By (i) with
$n=n_F$, the third factor is at most $\tfrac94 r^{1/2}$. By (ii),
$\ell_m/\ell_{n_F}=L^r$ and $e_a\ge1$, so the fourth factor is $L^{-e_ar}\le L^{-r}$. By (iii), the first
factor is at most $1$. Inserting these four bounds in \eqref{9.1:eq-weighted-unit-4} and then using
$\ell_\circ^2/L\le\tfrac13$ from (iii), we obtain
\begin{equation}\label{9.1:eq-weighted-unit-5}
  \frac{\mathsf{T}^a_N(x)}{\bar u^a_F}
  \le \tfrac94\,r^{1/2}\Bigl(\frac{\ell_\circ^2}{L}\Bigr)^{r}
  \le \tfrac94\,r^{1/2}\,3^{-r}.
\end{equation}

It remains to show that $r^{1/2}3^{-r}\le\tfrac13$ for every integer $r\ge1$. At $r=1$ the two sides are
equal. For $r\ge1$ the ratio of the $(r+1)$-st term to the $r$-th term is
$\bigl((r+1)/r\bigr)^{1/2}/3$, which is at most $\sqrt2/3<1$. Hence the sequence $r^{1/2}3^{-r}$ is
decreasing, and it is bounded by its value $\tfrac13$ at $r=1$. The right side of
\eqref{9.1:eq-weighted-unit-5} is thus at most $\tfrac94\cdot\tfrac13=\tfrac34$, which proves
\eqref{9.1:eq-weighted-unit-2}.

For old leaves $\bar u^a_F=\mathsf{u}^a_F$, so \eqref{9.1:eq-weighted-unit-2} reads
$\mathsf{T}^a_N(x)\le\tfrac34\mathsf{u}^a_F$.

For \eqref{9.1:eq-weighted-unit-3}, let $F$ be the leaf of the chain term $\mathbb{C}^{(i)}_k$, so that
$n_F=h+i<m$. We retain the factor $\varphi^{(N)}_m$ in \eqref{9.1:eq-weighted-unit-4}. The remaining
three factors are bounded exactly as above, with product at most $\tfrac34$. This gives
\begin{equation*}
  \mathsf{T}^a_N(x)\le\tfrac34\,\varphi^{(N)}_m\,\bar u^a_F\le\varphi^{(N)}_m\,\bar u^a_F .
\end{equation*}
The stage coefficients form a descending sequence \cite{B-CMP122a}: each stage coefficient
$\varphi^{(\nu)}_n$ is determined by a defining quantity, and a larger defining quantity yields a
smaller coefficient. The defining quantity of $\varphi^{(N)}_m$ is
\begin{equation*}
  4-2^{-(m-h-1)}-2^{1-(k-m)},
\end{equation*}
and that of $\varphi^{(i)}_{n_F}$ is
\begin{equation*}
  3-2^{-(n_F-h-1)}-2^{-(k-n_F)} .
\end{equation*}
Their difference is
\begin{align*}
  &\bigl[4-2^{-(m-h-1)}-2^{1-(k-m)}\bigr]-\bigl[3-2^{-(n_F-h-1)}-2^{-(k-n_F)}\bigr]\\
  &\qquad=\bigl[1-2^{1-(k-m)}\bigr]+\bigl[2^{-(n_F-h-1)}-2^{-(m-h-1)}\bigr]+2^{-(k-n_F)} .
\end{align*}
The first bracket is nonnegative because $k-m\ge1$ by (iv), the second is positive because $m>n_F$,
and the last term is positive. Hence the defining quantity of $\varphi^{(N)}_m$ is at least that of
$\varphi^{(i)}_{n_F}$ for $m>n_F$, whence $\varphi^{(N)}_m\le\varphi^{(i)}_{n_F}$. Consequently
\begin{equation*}
  \mathsf{T}^a_N(x)\le\varphi^{(i)}_{n_F}\,\bar u^a_F=\mathsf{T}^a_i(x),
\end{equation*}
which is \eqref{9.1:eq-weighted-unit-3}.
\end{proof}

\begin{definition}[The chain objects read in the graded region]\label{9.1:def-read-objects}
Let $Z$ be the large-field region at level $k$ and $\mathbf{B}''_k$ the determining set of
Ba{\l}aban's construction at the large-field site \cite[\S1]{B-CMP122b}; let $Y_1$, $Y_2$ and the
set $\mathbf{B}_1(Y_2)$ be those of the chain of Definition~\ref{9.1:def-substitution-chain}.
Here $\mathbf{B}''_k(Y_1)$ and $\mathbf{B}_1(Y_2)$ are the determining sets on which the objects
listed below are defined, and the boundaries $\partial Y_1$, $\partial Y_2$ carry the layer data
of the chain. We write $\bar\Lambda^{+}$ for a fixed region of Ba{\l}aban's construction at the
large-field site at level $k$. A reading takes
place \emph{inside} if it takes place in $\bar\Lambda^{+}$, and \emph{off} $\bar\Lambda^{+}$
otherwise. We call $\bar\Lambda^{+}$ the \emph{graded region}. For a term of the expansion, $X$
denotes its localisation domain in Ba{\l}aban's sense. We say that a term of the expansion is
\emph{read at} an object if the term is evaluated at that object.

The objects of the substitution chain of Definition~\ref{9.1:def-substitution-chain} are:
\begin{enumerate}
\item[\textup{(1)}] the configuration $U_{k,\Lambda}$ of \cite[(1.56)]{B-CMP122b};
\item[\textup{(2)}] the configuration $\mathfrak{U}_4$;
\item[\textup{(3)}] the minimiser $\Xi_3$;
\item[\textup{(4)}] the configuration $U_1$;
\item[\textup{(5)}] the configuration $\mathfrak{U}_3$ on $\mathbf{B}_1(Y_2)$;
\item[\textup{(6)}] the minimiser $\Xi_2$ on $\mathbf{B}_1(Y_2)$;
\item[\textup{(7)}] the configurations $U_{1,2}$ and $U^0_k$;
\item[\textup{(8)}] the configuration $\mathfrak{U}(\sigma')$ on $\mathbf{B}''_k(Y_1)$;
\item[\textup{(9)}] the minimiser $\Xi_1(\sigma')$ on $\mathbf{B}''_k(Y_1)$, with
$\sigma'$-localised layer data;
\item[\textup{(10)}] the configuration $U''_k$ of \cite[(1.59), (1.60)]{B-CMP122b} and the
vertices of the boxes.
\end{enumerate}
We put
\begin{equation}\label{9.1:eq-read-objects-1}
\mathfrak{Q}'' := \bigl\{\, U^0_k,\ \mathfrak{U}(\sigma'),\ \Xi_1(\sigma'),\ U''_k
\text{ and the box vertices} \,\bigr\}.
\end{equation}
The clause of the statement on the chain objects, which concerns readings inside the graded
region and is formulated for every object of the chain, is used on $\mathfrak{Q}''$ only.
\end{definition}

\begin{proof}
The terms with $X\cap Z\neq\emptyset$ are read at $U''_k$ \cite[(1.59), (1.60)]{B-CMP122b},
which together with the box vertices is item (10), and at the function of
\cite[(1.61)]{B-CMP122b} on $\mathbf{B}''_k(Y_1)$. We now identify the arguments of the latter.

By the comparison statement at the large-field site, no term is read on a set cut as in
Definition~\ref{9.1:def-substitution-chain}: the objects compared there, namely
$\mathfrak{U}(\sigma')$, $\Xi_1(\sigma')$ and the third compared object $P'$ of that statement,
live on sets based on $\mathbf{B}''_k$. The classification of the readings uses in addition
clause (0)(i) of the small-field statement at the large-field site. Of these three, 
$\mathfrak{U}(\sigma')$ and
$\Xi_1(\sigma')$ are items (8) and (9), while $P'$ is not among the objects (1)--(10). The
function of \cite[(1.61)]{B-CMP122b} lives on $\mathbf{B}''_k(Y_1)$, and the chain objects
defined on $\mathbf{B}''_k(Y_1)$ are items (8) and (9); since no term is read on a cut set, no
other chain object enters there. Hence, among the objects (1)--(10), the terms with
$X\cap Z\neq\emptyset$ are read through the function of \cite[(1.61)]{B-CMP122b} at items (8)
and (9), and through \cite[(1.59), (1.60)]{B-CMP122b} at item (10). The set $\mathfrak{Q}''$
contains in addition $U^0_k$ of item (7). Thus the objects reached so far, together with
$U^0_k$, are exactly the members of $\mathfrak{Q}''$ in \eqref{9.1:eq-read-objects-1}.

We next treat the objects (1)--(6). By the same comparison statement, the moves of the chain that
produce $U_1$ and $\Xi_2$ on $\mathbf{B}_1(Y_2)$, that is, items (4) and (6), are read off
$\bar\Lambda^{+}$ only. More generally, the objects (1)--(6) fall into two classes. Either they
are read by terms with $X\subset Z^c$, as stated in \cite[\S1]{B-CMP122b}; or they enter through
layer data at the boundaries $\partial Y_i$, which, by the geometric statement on the layer data
at the large-field site, lie off $\bar\Lambda^{+}$.

The objects (1)--(7) have one further entrance, as kernel and source of the operator
$\mathbf{H}''_k$ of the large-field site; this entrance is treated under the classical member
bound and not under the clause on the chain objects.

Hence, by the three preceding paragraphs, the clause on the chain objects is used on
$\mathfrak{Q}''$ only. At the objects (4)--(6) the requirement of that clause, formulated for
every object of the chain, does not hold: it fails at the witness configuration of the
counter-example proposition for the chain objects.
\end{proof}

\begin{lemma}[Frames drop out of the block averages]\label{9.1:lem-frames-drop}
Let $q_0$ be the root seed of Definition~\ref{9.1:def-aligned-points}, $\eta$ the parameter in the
exponent of $q_0$ there, and $\mathbf{U}$ the complex configuration on which $q_0$ depends. Write
$G=\mathrm{SU}(N)$ and $G^{c}$ for its complexification. For a configuration $U$ and a site map
$v$ with values in $G^{c}$ put $U^{v}(b):=v(b_-)\,U(b)\,v(b_+)^{-1}$ for every bond $b$ with initial site
$b_-$ and final site $b_+$.

Let $\mathbb{H}_3:=\mathfrak{X}_{\mathbf{B}_k}(b_3;q_0)$ be the response over $\mathbf{B}_k$, at the
seed $q_0$, to the source $b_3$ of the member numbered~$3$ of the substitution chain of
Definition~\ref{9.1:def-substitution-chain}. Let $\mathbb{H}_Z$ be the response, formed as in
step~(c) of that definition, to the layer data of $\mathbf{B}_k(Z)$ at $\partial Z$. The
\emph{carried presentation}, in the form of \cite[(1.43)]{B-CMP122b}, is the field
$\mathcal{W}_0$ given by
\begin{equation}\label{9.1:eq-frames-drop-1}
e^{i\eta\mathcal{W}_0}:=e^{i\eta\mathbb{H}_Z}\,e^{i\eta\mathbb{H}_3},\qquad
U_0^{\mathrm{car}}(\mathbf{U}):=e^{i\eta\mathcal{W}_0}\,q_0 .
\end{equation}
Let $\mathrm{fr}$ belong to the admissible frame class $\mathfrak{F}_G:=\{\mathrm{AL},\flat\}$. Here
$\mathrm{AL}$ is the axial--Landau frame of \cite{B-CMP122a}, and $\flat$ is the flat frame. Let
$U_0^{\mathrm{fr}}(\mathbf{U})$ be the field $U_0(\mathbf{U})$ of \cite[(1.81)]{B-CMP122a}, written in
the gauge frame $\mathrm{fr}$ of the substitution chain.

Assume, with $V_0$, $\mathfrak{B}_0$, $M^{j}$ and $\pi(c)$ as below:
\begin{enumerate}
\item[(a)] any two points of the chart of complex configurations that lie over one and the same
$\mathbf{B}_k(Z)$-datum, taken modulo gauge transformations at the block centres, lie on one
critical orbit of the $G^{c}$-valued gauge transformations;
\item[(b)] for each pair of points as in~(a), the gauge transformation carrying one to the other
depends holomorphically on all parameters and takes values in $G$ at real points, by the
identity-theorem argument by which the minimiser $\Xi$ of step~(c) of
Definition~\ref{9.1:def-substitution-chain} is identified with the object $\Xi_b$ compared with it;
\item[(c)] in the flat frame, $V_0(c)=M^{\pi(c)}(U_0^{\flat})(c)$ for every bond $c$ of
$\mathfrak{B}_0$.
\end{enumerate}

Then there is a site map $\varpi$ on $Z$ with values in $G^{c}$ such that
$U_0^{\mathrm{fr}}=(U_0^{\mathrm{car}})^{\varpi}$. The map $\varpi$ is holomorphic in all
parameters and takes values in $G$ at real points. Let $V_0$ be the datum of block averages on the
bonds of $\mathfrak{B}_0$ in \cite[(1.81)]{B-CMP122a}. Here $M^{j}$ denotes Ba{\l}aban's
block-averaging map of Definition~\ref{1.5:def-block-average} to level $j$. Now let $c$ be any bond of
$\mathfrak{B}_0$: a
bond of the tower, of the rind or of the interface, or a bond crossing $\partial\Lambda$. Let
$\pi(c)$ be its level and $c_-$, $c_+$ its endpoints. Then
\begin{equation}\label{9.1:eq-frames-drop-2}
V_0(c)=\varpi(c_-)\,M^{\pi(c)}(U_0^{\mathrm{car}})(c)\,\varpi(c_+)^{-1}.
\end{equation}
\end{lemma}

\begin{proof}
Write $M_k$ for the block-averaging map of Definition~\ref{1.5:def-block-average} to level $k$, and
$M_{\mathbb{B}''}$ for the map of averages on the bond set $\mathbb{B}''$ of \cite[(1.81)]{B-CMP122a},
which assigns to each bond of $\mathbb{B}''$ the block average of the configuration on that bond.
Both $U_0^{\mathrm{fr}}(\mathbf{U})$ and $U_0^{\mathrm{car}}(\mathbf{U})$ are points of the chart of
complex configurations over
one and the same $\mathbf{B}_k(Z)$-datum, taken modulo gauge transformations at the block centres.
This datum has three parts:
\begin{itemize}
\item the block averages $M_k\mathbf{U}$ off $\Lambda$;
\item the variables $V_\Lambda(\mathbf{U})$ on $\Lambda$;
\item Ba{\l}aban's data on $\partial Z$.
\end{itemize}
By hypothesis~(a) the two fields lie on one critical orbit. Hence
$U_0^{\mathrm{fr}}=(U_0^{\mathrm{car}})^{\varpi}$ for a site map $\varpi$ on $Z$ with values
in~$G^{c}$. By hypothesis~(b), applied to this pair, $\varpi$ is holomorphic in all parameters
and $G$-valued at real points.

For the second assertion we use that the datum $V_0$ on the bonds of $\mathfrak{B}_0$ is the
collection of block averages of a single configuration. In the axial frame of
\cite[(1.81)]{B-CMP122a} this is the relation $V_0=M_{\mathbb{B}''}(U_0)$; in the two frames of
$\mathfrak{F}_G$ it reads as follows. In the axial--Landau frame,
\cite{B-CMP122a} reads $V_0$ as the block averages $M^{j}(U_0^{(\mathrm{AL})})$ at every level $j$.
In the flat frame the same reading is hypothesis~(c).
Thus, for
$\mathrm{fr}\in\mathfrak{F}_G$ and every bond $c$ of $\mathfrak{B}_0$ of level $\pi(c)$,
\begin{equation}\label{9.1:eq-frames-drop-3}
V_0(c)=M^{\pi(c)}(U_0^{\mathrm{fr}})(c)=M^{\pi(c)}\bigl((U_0^{\mathrm{car}})^{\varpi}\bigr)(c).
\end{equation}

It remains to move the gauge transformation through the block average. Write $\overline{V}$ for
the block average of a configuration $V$ and $\bar v$ for the restriction of a gauge transformation
$v$ to the sites of the block lattice. By \cite[(11)]{B-CMP98}, the block average of a
gauge-transformed configuration is the gauge transform of the block average; in this notation,
\begin{equation*}
\bigl(\overline{V^{v}}\bigr)_c=\bar v(c_-)\,\overline{V}_c\,\bar v^{-1}(c_+).
\end{equation*}
We apply this with $v=\varpi$ and the configuration $U_0^{\mathrm{car}}$, averaged to level
$\pi(c)$. The endpoints $c_-$, $c_+$ of $c$ are points of the block lattice of that level. Hence
\begin{equation*}
M^{\pi(c)}\bigl((U_0^{\mathrm{car}})^{\varpi}\bigr)(c)
=\varpi(c_-)\,M^{\pi(c)}(U_0^{\mathrm{car}})(c)\,\varpi(c_+)^{-1}.
\end{equation*}
Inserting this into \eqref{9.1:eq-frames-drop-3} gives \eqref{9.1:eq-frames-drop-2}.
\end{proof}

\begin{definition}[The gauge-rotated presentation of a chain object]
\label{9.1:def-rotated-presentation}
Let $\mathfrak{F}_G:=\{\mathrm{AL},\flat\}$ consist of the axial--Landau frame $\mathrm{AL}$ and the
flat frame $\flat$ of Lemma~\ref{9.1:lem-frames-drop}, and let $\mathrm{fr}\in\mathfrak{F}_G$. The
presentation defined below is
used in these two frames only; it is not used in the axial frame of that lemma. In the axial frame
the seam-crossing level-$h$ plaquettes carry a further source, driven by the seed, not scaled, and
of the size of the tube, for which no bound of the required rescaled form is available; this is why
the presentation is restricted to $\mathfrak{F}_G$. Let $\varpi$ be the
$G^c$-valued site map of Lemma~\ref{9.1:lem-frames-drop}, so that
\begin{equation}\label{9.1:eq-rotated-presentation-1}
U_0^{\mathrm{fr}}=(U_0^{\mathrm{car}})^{\varpi},\qquad
U_0^{\mathrm{car}}=e^{i\eta\mathcal{W}_0}q_0 .
\end{equation}
Here $\eta$, $\mathcal{W}_0$ and $q_0$ are as in Lemma~\ref{9.1:lem-frames-drop}, and, with the
gauge-transformation convention of Lemma~\ref{1.1:lem-invariance},
$U^{\varpi}(b)=\varpi(x)U(b)\varpi(y)^{-1}$ for a bond $b$ from $x$ to $y$. The identity
\eqref{9.1:eq-rotated-presentation-1} is used on the bonds of $\mathfrak{B}_0$ only.

Let $q\in\mathfrak{Q}''$ be one of the chain objects of Definition~\ref{9.1:def-read-objects}, with
the regions $\Lambda^{\sim}$, $\bar\Lambda$, $\bar\Lambda^{+}$ and the core $\mathfrak{K}$ of the
chain objects of that definition. It is
a chart point over a set $\mathfrak{S}_q\supseteq\mathbf{B}''_k{\upharpoonright}\Lambda^{\sim}$ with
data $\mathsf{D}_q$ consisting of:
\begin{itemize}
\item the configuration $V''$ off $\bar\Lambda$;
\item $e^{i\mathsf{b}}V_0$ on $\mathfrak{B}_0$, where $\mathsf{b}$ is either $0$ or the datum $g_kB$
of the vertex at which $q$ is taken;
\item $V''$ on the core $\mathfrak{K}$;
\item the layer and rim data off $\bar\Lambda^{+}$.
\end{itemize}
We write \emph{side~$1$} for the $U$-side of the junction; it contains the sites of $\mathfrak{S}_q$
off the core. In the frame $\mathrm{fr}$ we write
$q=e^{i\eta\mathcal{Y}_q}U_0^{\mathrm{fr}}$ on side~$1$, where $\mathcal{Y}_q$ is the presentation
field of $q$.
\begin{enumerate}
\item[(a)] The \emph{relative seam datum} of $q$ is the following quantity on the cut strata of the
junction:
\begin{equation}\label{9.1:eq-rotated-presentation-2}
\mathsf{b}_\Sigma:=\frac1i\log\bigl[\mathrm{rd}\,q\cdot(\mathrm{rd}\,U_0^{\mathrm{fr}})^{-1}\bigr].
\end{equation}
This is the bracket of \cite{B-CMP122a}. In the frame $\mathrm{fr}$ the factors of
$\mathrm{rd}\,q\cdot(\mathrm{rd}\,U_0^{\mathrm{fr}})^{-1}$ lie within $C(\varepsilon_h+\delta'_k)$,
where $\varepsilon_h$ and $\delta'_k$ are the small-field parameters of levels $h$ and $k$, and
$CL^{-N}\mathsf{s}_L$ of the identity; here the exponent $N$ is an integer index, not the rank of
$\mathrm{SU}(N)$. The size bounds imposed on the seam datum, namely that the seam datum is bounded
in size by $C_{\mathsf{b}}\varepsilon_h$, with $C_{\mathsf{b}}$ a constant, and by $CL^{-N}\mathsf{s}_L$,
and the bound of the same kind on the
variation of the site map $\varpi$ along the boundary $\partial\mathfrak{K}$ of the core, are bounds on
$\mathsf{b}_\Sigma$, that is,
they are taken relative to $U_0^{\mathrm{fr}}$. The junction bound in the flat frame bounds the
presented object and not the map $\varpi$.
\item[(b)] The \emph{gauge-rotated presentation} of $q$ is
\begin{equation}\label{9.1:eq-rotated-presentation-3}
q^{\sharp}:=\bigl(q{\upharpoonright}\text{side }1\bigr)^{\varpi^{-1}} .
\end{equation}
Thus $q^{\sharp}$ is a configuration on side~$1$ only; $\varpi$ is not extended to the core.
\end{enumerate}
Among the conditions defining the coupling ceiling $\gamma_U$ in the order of choice
(Notations~\ref{2.3:not-stages-first}, \ref{2.3:not-stages-middle} and~\ref{2.3:not-stages-last})
we include
\begin{equation}\label{9.1:eq-rotated-presentation-4}
2C_4^{\sharp}\mathsf{s}_L\le\min\{c_4,a_{R'}\}.
\end{equation}
Here $C_4^{\sharp}$ is the constant of the bound on the presentation of the seed, and $c_4$ and
$a_{R'}$ are radii of the chart.
The presentation has the following properties.
\begin{enumerate}
\item[(i)] On the bonds $b$ of side~$1$ at which \eqref{9.1:eq-rotated-presentation-1} is used,
\begin{equation}\label{9.1:eq-rotated-presentation-5}
q^{\sharp}(b)=e^{i\eta\operatorname{Ad}_{\varpi(x)^{-1}}\mathcal{Y}_q(b)}\,
e^{i\eta\mathcal{W}_0(b)}\,q_0(b),
\end{equation}
where $x$ is the initial site of $b$.
\item[(ii)] Let $\mathfrak{T}$ denote the tube of the chart, whose points are data taken modulo
block-centre gauge transformations; it is not the separation window
$\mathfrak{T}(\mathsf{d};\vartheta)$ of Notation~\ref{0.2:not-glossary}. On side~$1$,
$\mathrm{rd}\,q^{\sharp}=(\mathrm{rd}\,q)^{[\varpi^{-1}]}$, where
$(\mathrm{rd}\,q)^{[\varpi^{-1}]}$ denotes the reading $\mathrm{rd}\,q$ acted on by the gauge
transformation that $\varpi^{-1}$ induces on readings, a block-centre gauge transformation in the
sense of the definition of the tube $\mathfrak{T}$; and this is the
same point of the tube $\mathfrak{T}$ as the reading of $q$ on side~$1$.
\item[(iii)] $\tilde F(X;q^{\sharp})=\tilde F(X;q)$ for the regions $X$ of the graded-region condition
of Definition~\ref{9.1:def-graded-ceilings}.
\item[(iv)] The admissibility condition $\mathrm{rd}\,\mathbf{U}\in\mathfrak{T}[0](X)$ of the
definition of the target tube $\mathfrak{T}[0](X)$ may be checked at $q^{\sharp}$.
\item[(v)] The slot response entering the presentation of the seed is bounded by
\cite[(1.65)]{B-CMP116}, at a cost $(B^{\mathsf{s}})^2C_T$, where $C_T$ is a constant of the
$\mathbf{T}$-step; this cost is not included in $C_4^{\sharp}$.
The datum of the response field $\mathbb{H}_Z$ of Lemma~\ref{9.1:lem-frames-drop} at $\partial Z$ is
not a source of the response bound on the small-field region $\Lambda$. Both contributions are paid
for by the factor
$e^{-\delta_1M(N_0-2)/2}$ under the first ceiling condition of
Definition~\ref{9.1:def-graded-ceilings}.
\end{enumerate}
\end{definition}

\begin{proof}
\emph{The bonds used.} We first record the range of bonds. \cite{B-CMP122a} states that bonds
intersecting the boundary $\partial\Omega''$ of the region $\Omega''$ of the $\mathbf{R}$-step do not
belong to $\mathfrak{B}_0$. For this reason \eqref{9.1:eq-rotated-presentation-1} is used on
$\mathfrak{B}_0$ only.

\emph{Property (i).} Let $b$ be a bond of $\mathfrak{B}_0$ on side~$1$, from $x$ to $y$. On $b$ we have
$q(b)=e^{i\eta\mathcal{Y}_q(b)}U_0^{\mathrm{fr}}(b)$. By \eqref{9.1:eq-rotated-presentation-1},
\[
q(b)=e^{i\eta\mathcal{Y}_q(b)}\,\varpi(x)U_0^{\mathrm{car}}(b)\varpi(y)^{-1}.
\]
Applying the gauge transformation $\varpi^{-1}$ multiplies on the left by $\varpi(x)^{-1}$ and on the
right by $\varpi(y)$. This gives
\[
q^{\sharp}(b)=\varpi(x)^{-1}e^{i\eta\mathcal{Y}_q(b)}\varpi(x)\,U_0^{\mathrm{car}}(b).
\]
We now use the identity $g\,e^{i\eta Y}g^{-1}=e^{i\eta\operatorname{Ad}_gY}$ with $g=\varpi(x)^{-1}$,
together with $U_0^{\mathrm{car}}=e^{i\eta\mathcal{W}_0}q_0$. This yields
\eqref{9.1:eq-rotated-presentation-5}.

The field $\operatorname{Ad}_{\varpi^{-1}}\mathcal{Y}_q$ may jump across faces of level-$k$ blocks,
where $\varpi$ jumps. These jumps enter no estimate, because the size bound used for presentation
fields reads only $\sup|\mathcal{Y}_q|$.

\emph{Property (ii).} The data of $q^{\sharp}$ are the data of $\mathsf{D}_q$ on side~$1$ transformed
by $\varpi^{-1}$; no datum on the core is transformed, since $\varpi$ is not extended there. Since
the block-averaging map of Definition~\ref{1.5:def-block-average} is covariant under site-wise gauge
transformations, also under $G^c$-valued ones as in clause (iii) of the gauge-covariance lemma for
the re-read terms, the reading of
$q^{\sharp}=(q{\upharpoonright}\text{side }1)^{\varpi^{-1}}$ on side~$1$ is therefore
$(\mathrm{rd}\,q)^{[\varpi^{-1}]}$. The points of $\mathfrak{T}$ are data taken modulo block-centre gauge
transformations, by the definition of the tube $\mathfrak{T}$. Hence $\mathrm{rd}\,q^{\sharp}$ and
the reading of $q$ on side~$1$ are the same point of $\mathfrak{T}$.

\emph{Property (iii).} By the representation lemma for the re-read terms $\tilde F$,
$\tilde F(X;q)$ depends on $q$ only through the restriction of the configuration to $X^{+}$. By
clause (iii) of the gauge-covariance lemma for these terms,
it is invariant under $G^c$-valued gauge transformations. Both the representation and the
gauge-covariance statement are local on $X^{+}$.

For the regions $X$ of the graded-region condition, the localisation neighbourhoods
$\tilde\Delta(y)$ of the sites $y$ at which $\tilde F(X;\cdot)$ reads the configuration satisfy
$\tilde\Delta(y)\cap\mathfrak{K}=\emptyset$. Hence only values of $q$ on side~$1$ enter
$\tilde F(X;q)$. On side~$1$, $q^{\sharp}$ is the gauge transform of $q$ by $\varpi^{-1}$, so
$\tilde F(X;q^{\sharp})=\tilde F(X;q)$. For the same reason neither an extension of $\varpi$ to the
core nor a choice of seam point is needed.

\emph{Property (iv).} The admissibility condition $\mathrm{rd}\,\mathbf{U}\in\mathfrak{T}[0](X)$ is a
statement about the class of $\mathrm{rd}\,\mathbf{U}$ in $\mathfrak{T}$ only. By (ii), the readings
of $q$ and $q^{\sharp}$ on side~$1$ define the same class. The condition may therefore be checked at
$q^{\sharp}$.

\emph{Property (v).} The slot response is bounded by \cite[(1.65)]{B-CMP116}, with the cost
$(B^{\mathsf{s}})^2C_T$. The boundary $\partial Z$ lies at distance at least $2M(N_0-2)$ from the
region on which the response bound on $\Lambda$ is taken. Both the slot
response and the datum of $\mathbb{H}_Z$ at $\partial Z$ are therefore paid for by
$e^{-\delta_1M(N_0-2)/2}$, under the first ceiling condition of
Definition~\ref{9.1:def-graded-ceilings}.
\end{proof}

\begin{lemma}[Splitting the rotation in the graded region]\label{9.1:lem-split-rotation}
Let the chain objects be presented in a frame $\mathrm{fr}\in\mathfrak{F}_G=\{\mathrm{AL},\flat\}$, the
axial--Landau frame or the flat frame. Here $\ell_m$ is the length unit of label level $m$, so that
$\ell_m/\ell_k=L^{-(k-m)}$; $m(y)$ is the label level of a point $y$, and a bound on a field at a point of
label level $m$ is \emph{in own units} if it bounds $\ell_m$ times the field. The seam $\mathcal{L}$ separates
side~$1$ from side~$2$, $Z''^{(1)}_k$ is the part of $Z''_k$ on side~$1$, and $\tilde\Delta(y)$ is the cube
attached to $y$ at its label level. The integer $h<k$ is the level of the cut along the seam, fixed in
\cite[\S1]{B-CMP122a}, and $\Omega''^{\sim}_{h+1}$ is the region of \cite[\S1]{B-CMP122a} at level $h+1$
along whose boundary the cut strata lie. Further, $\eta$ is the lattice spacing in Ba{\l}aban's
parametrisation $e^{i\eta A}$ of bond variables near the identity, with $A$ a Lie-algebra-valued field;
$\delta_1$ is the decay rate of the kernels, and $C_s$, $c_1$ and $C_T$ are the regularity, profile-smearing
and weighted-summation constants. The numbers $\delta'_k$ and $\varepsilon_h$ are Ba{\l}aban's radii at
levels $k$ and $h$, and $C_1$ is Ba{\l}aban's constant of the glossary, one of the auxiliary parameters; it
enters through the bound $|\mathsf{b}|\le C_1\delta'_k$ on the datum $\mathsf{b}$ on $\mathfrak{B}_0$.
Assume the following.
\begin{enumerate}
\item[(a)] The classical condition in profile form, Definition~\ref{9.1:def-profile-condition}, with its
profile clause. Every increment of every site, with source datum $b$ and dilation $\tau_t$, is bounded at each
point $y$ of the contour of the site, in the units $\ell_\iota$ of that contour, by
$B^{\mathsf{s}}\bar\tau(y)\,(|b|)^{\sim}(y)$.
\item[(b)] The size bound inside $\bar\Lambda$: the level-$k$ data through which the seed enters the chain
vary over the shape by at most $\mathsf{s}_L$, and the sources of the chain inside $\bar\Lambda$ are
majorised by $\mathsf{s}_L$.
\item[(c)] The corollary placing the data of the chain objects in the admissible class at the ambient scale.
\item[(d)] The link-by-link bound, Lemma~\ref{9.1:lem-link-bound}, with its constant $C_4$, for the chain
objects of a fixed union of two families of chain objects; these are the chain objects to which it is applied
in the step on layer and rim increments below.
\item[(e)] The seam statements: the interface bound for the cut increments, with constant $C_{IC}$; the seam
proposition, on the seam $\mathcal{L}$ between the two sides; and the second cut bound at complex seeds.
\item[(f)] The ceiling condition $\|\operatorname{Ad}_{\bar\varpi}\|\le 2$, where $\bar\varpi$ is the
rotation of Definition~\ref{9.1:def-rotated-presentation}.
\item[(g)] The seam bound: on the cut strata of $\mathbf{B}_h(\Omega''^{\sim}_{h+1})$ the relative seam
datum
\begin{equation*}
\mathsf{b}_\Sigma:=\tfrac1i\log\bigl[\operatorname{rd}q\cdot(\operatorname{rd}U_0^{\mathrm{fr}})^{-1}\bigr]
\end{equation*}
of each $q\in\mathfrak{Q}''$ produces a response on side~$1$ whose size in own units is at most
$C_{seam}\,\varepsilon_h\,e^{-\delta_1 d(\cdot,\mathcal{L})/4}$, with a constant $C_{seam}$.
\end{enumerate}
For $q\in\mathfrak{Q}''$ let $q^{\sharp}$ be the presentation of
Definition~\ref{9.1:def-rotated-presentation}. At the points read here, where $\tilde\Delta(y)$ does not meet
the core $\mathfrak{K}$, it equals $(q{\upharpoonright}\text{side }1)^{\varpi^{-1}}$, where $\varpi$ is the
site map of Lemma~\ref{9.1:lem-frames-drop}. Then for every root seed $q_0$, every vertex of
$\mathfrak{P}_{\mathsf{c}}$, every $q\in\mathfrak{Q}''$ and every $y$ with
$\tilde\Delta(y)\subset Z''^{(1)}_k$,
\begin{equation}\label{9.1:eq-split-rotation-1}
q^{\sharp}=e^{i\eta\mathcal{Z}^{\sharp}}q_0,\qquad
\mathcal{Z}^{\sharp}=\mathcal{Z}^{(k)}+\mathcal{Z}^{\mathrm{re}},
\end{equation}
\begin{gather}
\sup_{\tilde\Delta(y)}|\mathcal{Z}^{(k)}|\le \frac{C_4^{\sharp}\,\mathsf{s}_L}{\ell_k},\qquad
C_4^{\sharp}:=2^{9}\,d\,B^{\mathsf{s}}C_s c_1 L^{8}C_4,\label{9.1:eq-split-rotation-2}\\
\ell_{m(y)}\sup_{\tilde\Delta(y)}|\mathcal{Z}^{\mathrm{re}}|\le \mathsf{m}^{\mathrm{re}\sharp}(y)
:=4B^{\mathsf{s}}C_s c_1 C_1\delta'_k+2C_{seam}\,\varepsilon_h\,e^{-\delta_1 d(y,\mathcal{L})/4}.
\label{9.1:eq-split-rotation-3}
\end{gather}
\end{lemma}

\begin{proof}
Every contribution to $q^{\sharp}$ is assigned to $\mathcal{Z}^{(k)}$ if it is bounded absolutely by a
constant times $\mathsf{s}_L/\ell_k$. It is assigned to $\mathcal{Z}^{\mathrm{re}}$ if it is bounded in own
units by a constant times $\delta'_k$ or times $\varepsilon_h e^{-\delta_1 d(\cdot,\mathcal{L})/4}$. The points
read satisfy $\tilde\Delta(y)\subset Z''^{(1)}_k$, and their cubes do not meet the core:
$\tilde\Delta(y)\cap\mathfrak{K}=\emptyset$. So the
rotation is needed only on side~$1$ and off $Z$. On side~$1$ the rotation $\bar\varpi$ of
Definition~\ref{9.1:def-rotated-presentation} is $\varpi$. Write $q=e^{i\eta\mathcal{Y}_q}U_0^{\mathrm{fr}}$
on side~$1$, and let $\mathcal{W}_0$ be the field of the carried presentation
$U_0^{\mathrm{car}}=e^{i\eta\mathcal{W}_0}q_0$ of Lemma~\ref{9.1:lem-frames-drop}. Then, bond by bond,
\begin{equation*}
q^{\sharp}=e^{i\eta\operatorname{Ad}_{\varpi^{-1}}\mathcal{Y}_q}\,e^{i\eta\mathcal{W}_0}\,q_0 ,
\end{equation*}
with hypothesis~(f) controlling the adjoint action. The jumps of $\varpi$ across $k$-faces cost nothing,
since only $\sup|\mathcal{Y}_q|$ is read.

\emph{The base.} Put $P:=e^{i\eta\mathbb{H}_3}q_0$, where
$\mathbb{H}_3=\mathfrak{X}_{\mathbf{B}_k}(b_3;q_0)$ is the response on $\mathbf{B}_k$ to the source $b_3$
of the third stage of the expansion. This is the first object read in
Definition~\ref{9.1:def-read-objects}. It is defined globally, is critical on $\mathbf{B}_k$, and is of level
$k$ on $\Lambda^{\sim}$. The set $\mathfrak{S}_q$ of $q$ refines $\mathbf{B}_k$ on side~$1$. By the nesting
property of the minimisers, which holds at complex data, $P$ is critical on $\mathfrak{S}_q\cap\text{side }1$
with its own readings there as data. By the definition of the carried presentation in
Lemma~\ref{9.1:lem-frames-drop}, $U_0^{\mathrm{car}}=e^{i\eta\mathbb{H}_Z}P$ on $Z$, where $\mathbb{H}_Z$ is
the response of the $\partial Z$-layer of $\mathbf{B}_k(Z)$. By Lemma~\ref{9.1:lem-frames-drop}, the rotated
data of $q$ satisfy
\begin{equation*}
\mathsf{D}_q^{[\bar\varpi^{-1}]}=\operatorname{rd}\bigl(e^{i\eta\mathbb{H}_Z}P\bigr)
\ \text{on }\mathfrak{B}_0\text{ at }\mathsf{b}=0,\qquad
\mathsf{D}_q^{[\bar\varpi^{-1}]}=\operatorname{rd}P\ \text{off }\bar\Lambda .
\end{equation*}
We prove the second identity. Off $\bar\Lambda$ one has $Q_k(\eta\mathbb{H}_3)=0$. So the data $V''$ there are
the readings of $P$ in a block-centre gauge; we denote this block-centre gauge by $\varpi_3$, and off $Z$ we put
$\bar\varpi:=\varpi_3^{-1}$. We denote by $\varpi_Z$ the gauge factor across the junction between
$\bar\varpi=\varpi$ on side~$1$ and $\bar\varpi=\varpi_3^{-1}$ off $Z$. This junction is a level-$k$ datum
increment, which enters as a further source; the size of the increment is
\begin{equation}\label{9.1:eq-split-rotation-4}
|\varpi_Z-1|\le 16\,d\,\ell_k\sup|\mathbb{H}_Z| .
\end{equation}
It is absorbed at the level-$k$ bonds crossing the junction, by the statement of the seam analysis that
absorbs datum increments at such bonds. By the bound $\ell_k|\mathbb{H}_Z|\le B^{\mathsf{s}}C_sc_1\mathsf{s}_L$
recorded in the next paragraph, it is assigned to $\mathcal{Z}^{(k)}$. Hence
$q^{\sharp}=e^{i\eta\mathbb{H}^{\sharp}}P$ on side~$1$, where $\mathbb{H}^{\sharp}$ is the response to four
kinds of sources. They all lie on side~$1$ and away from the seam $\mathcal{L}$: $\mathbb{H}_3$ and the data
increments on $\mathfrak{B}_0$, the rotated datum $\mathsf{b}$ on $\mathfrak{B}_0$, and the layer and rim
increments. The seam carries a fifth source, which is read relatively in the last step.

\emph{The third-stage response and the data increments on $\mathfrak{B}_0$.} $\mathbb{H}_3$ is a response on
a level-$k$ set. By the regularity of the response map together with hypothesis~(a), and by hypothesis~(b),
$\ell_k|\mathbb{H}_3|\le B^{\mathsf{s}}C_sc_1\mathsf{s}_L$; it is assigned to $\mathcal{Z}^{(k)}$. The
increments of the data on $\mathfrak{B}_0$,
$\operatorname{rd}(e^{i\eta\mathbb{H}_Z}P)(\operatorname{rd}P)^{-1}$, are block readings of a rotation. By the
reading lemma for block averages, at a site of level $\pi$ they are bounded by
\begin{equation*}
34\,d\,\ell_\pi\sup_\Lambda|\mathbb{H}_Z| .
\end{equation*}
The response $\mathbb{H}_Z$ obeys the bound of the same form as $\mathbb{H}_3$,
$\ell_k|\mathbb{H}_Z|\le B^{\mathsf{s}}C_sc_1\mathsf{s}_L$, by the regularity of the response map together
with hypothesis~(a). It is in fact far smaller, since $\partial Z$ lies at distance at least $2M(N_0-2)$, where
$N_0$ is an integer of Ba{\l}aban's construction of the large-field region $Z$.
These increments carry the factor $\ell_\pi$. Their response, by hypothesis~(a) and the weighted summation
bound at exponent $a=1$, with constant $C_T$, is sized by the bound for the response to these data in
Ba{\l}aban's expansion and carries the constant $(B^{\mathsf{s}})^2C_T$; this constant is not part of
$C_4^{\sharp}$. Nor is the datum of $\mathbb{H}_Z$ at $\partial Z$ among the sources bounded in
hypothesis~(b). Both are absorbed by the factor $e^{-\delta_1M(N_0-2)/2}$, under the corresponding lower bound
on $L$ in the order of choice (Notations~\ref{2.3:not-stages-first}, \ref{2.3:not-stages-middle}
and~\ref{2.3:not-stages-last}), because $\partial Z$ lies at distance at least $2M(N_0-2)$; so this response
is assigned to $\mathcal{Z}^{(k)}$.

\emph{The rotated datum on $\mathfrak{B}_0$.} The fluctuation variable satisfies $|B'|<\delta'_k$ on
$\mathfrak{B}_0$ by \cite[(1.82)]{B-CMP122a}, and the datum $\mathsf{b}$ of the vertex, the $g_kB$-datum,
obeys $|\mathsf{b}|\le C_1\delta'_k$. By hypothesis~(f),
$|\operatorname{Ad}_{\bar\varpi^{-1}}\mathsf{b}|\le 2C_1\delta'_k$ on $\mathfrak{B}_0$, the dilations being
contained in $B^{\mathsf{s}}$. By hypothesis~(a) its response is at most
$B^{\mathsf{s}}C_sc_1\cdot 2C_1\delta'_k$ in own units; it is assigned to $\mathcal{Z}^{\mathrm{re}}$.

\emph{Layer and rim increments.} At sites $y'\notin\bar\Lambda^{+}$ these are relative readings of chain
objects against the base. Let $p$ be the level of the layer cell containing $y'$. For a chain object $q'$ read
there let $\mathcal{Z}_{q'}$ be its field in Lemma~\ref{9.1:lem-link-bound}. By the reading lemma for block
averages and by hypothesis~(d), for the chain objects of the union of families named there,
\begin{equation*}
|b(y')|\le 68\,d\,\ell_{p}\sup\bigl(|\mathcal{Z}_{q'}|+|\mathcal{W}_0|\bigr)\le 136\,d\,C_4\,\mathsf{s}_L .
\end{equation*}
Let $x\in Z''_k$ have label level $m$. Every path from $x$ to $y'$ crosses the rind, which is of level $k$. By
the distance bound for such crossings, $d(x,y')\ge c_{\mathrm{cr}}M_1(k-m-1)_+$, where $c_{\mathrm{cr}}$ is
the constant of that bound. The flat exponent floor of the same order of choice gives
$e^{-\delta_1c_{\mathrm{cr}}M_1/2}\le L^{-8}$. For $m\le k$ we have $(k-m-1)_+\ge k-m-1$, hence
\begin{equation*}
e^{-\delta_1 d(x,y')/2}\le L^{-8(k-m-1)_+}\le L^{8}L^{-8(k-m)} .
\end{equation*}
Since $\ell_m/\ell_k=L^{-(k-m)}\ge L^{-8(k-m)}$ for $m\le k$, hypothesis~(a) and the decay bound for the
kernel give, for the response $\mathbb{H}$ to these increments,
\begin{equation*}
\ell_m|\mathbb{H}|(x)\le 136\,d\,B^{\mathsf{s}}C_sc_1C_4\,\mathsf{s}_L\,L^{8}\,\frac{\ell_m}{\ell_k},
\end{equation*}
that is,
\begin{equation*}
|\mathbb{H}|(x)\le \frac{136\,d\,B^{\mathsf{s}}C_sc_1C_4L^{8}\mathsf{s}_L}{\ell_k}.
\end{equation*}
This contribution is assigned to $\mathcal{Z}^{(k)}$. In particular the second and first stages of the
expansion act on $Z''_k$ only through these increments.

\emph{The seam.} By the identification of the minimiser in the random-walk expansion, $q^{\sharp}$ restricted
to side~$1$ is the minimiser on $\mathfrak{S}_q\cap\text{side }1$ joined with the cut strata of
$\mathbf{B}_h(\Omega''^{\sim}_{h+1})$. Its data are its own readings there; this is Ba{\l}aban's device in
\cite[\S1]{B-CMP122a}. The seam datum is read relatively, as the bracket $\mathsf{b}_\Sigma$ of
hypothesis~(g), which is Ba{\l}aban's bracket in \cite[\S1]{B-CMP122a}. In the frame $\mathrm{fr}$ its factors
lie within $C(\varepsilon_h+\delta'_k)$ and within $CL^{-(k-h)}\mathsf{s}_L$ of $1$, for a constant $C$; here
$k-h$ is the number of steps from the level $h$ of the seam to $k$. The three sizes entering the bracket are
read in this relative form:
\begin{itemize}
\item the core data $V_h$ on $\operatorname{supp}\chi_{h,1/2}$ are at most
$C_{IC}\varepsilon_hL^{-\mathsf{y}}$, by the interface bound of hypothesis~(e) on side~$2$ together with the
bound on the core data;
\item the readings of the base there are within $CL^{-(k-h)}\mathsf{s}_L$, by the reading lemma for block
averages;
\item the variation of the frame factor along $\partial\mathfrak{K}$ is of the same kind.
\end{itemize}
The last two are assigned to $\mathcal{Z}^{(k)}$. By hypothesis~(g) the response is at most
$C_{seam}\varepsilon_he^{-\delta_1d(\cdot,\mathcal{L})/4}$ in own units; it is assigned to
$\mathcal{Z}^{\mathrm{re}}$.

\emph{Composition.} The response map $\mathfrak{X}$ is not additive in its data. The difference between the
responses with and without one source is controlled by the local Lipschitz property
$\natural\mathrm{Lip}^{\mathrm{loc}}$ of the Landau chart. The composition of the exponentials by the
Baker--Campbell--Hausdorff formula costs a factor $2$. By \eqref{9.1:eq-split-rotation-2}, the cross terms
satisfy
\begin{equation*}
\eta\,|\mathcal{Z}^{(k)}|\,|\mathcal{Z}^{\mathrm{re}}|\le
\frac{\eta}{\ell_k}\,C_4^{\sharp}\,\mathsf{s}_L\,|\mathcal{Z}^{\mathrm{re}}| .
\end{equation*}
Under the coupling ceiling $\gamma_U$, which includes the condition
$2C_4^{\sharp}\mathsf{s}_L\le\min\{c_4,a_{R'}\}$ on the radii $c_4$ and $a_{R'}$ of the Landau chart, the
$\ell$-scaled size $C_4^{\sharp}\mathsf{s}_L$ of the first factor is
at most $1$. So the cross terms are charged to $\mathcal{Z}^{\mathrm{re}}$ and are covered by the factor $2$.
The contributions assigned to $\mathcal{Z}^{\mathrm{re}}$ give, with the factor $2$,
\begin{equation*}
2\bigl(2B^{\mathsf{s}}C_sc_1C_1\delta'_k+C_{seam}\varepsilon_he^{-\delta_1d(y,\mathcal{L})/4}\bigr)
=\mathsf{m}^{\mathrm{re}\sharp}(y),
\end{equation*}
which is \eqref{9.1:eq-split-rotation-3}. The contributions assigned to $\mathcal{Z}^{(k)}$ are the following:
$\mathbb{H}_3$, the data increments on $\mathfrak{B}_0$, the junction \eqref{9.1:eq-split-rotation-4}, the
layer and rim increments, and the seam readings of the base and of the frame. Collected with the factor $2$
under the constant $C_4^{\sharp}$, they give \eqref{9.1:eq-split-rotation-2}, and
\eqref{9.1:eq-split-rotation-1} holds.
\end{proof}

\begin{remark}
The proof uses none of the following: the pinned unit data of the cut, the second-stage datum of
\cite[(1.57)]{B-CMP122b}, the axial--Landau potential $A_0$, or any frame outside $\mathfrak{F}_G$. Consider
the plateau configuration of the proposition that constructs the abelian plateau seed. There one has
$\mathcal{W}_0=0$, and $\mathcal{Z}^{\sharp}$ is the sum of the contribution of the rotated datum on
$\mathfrak{B}_0$ and the contribution of the seam. The identity \eqref{9.1:eq-split-rotation-1} is the
counterpart, in the present notation, of the statement in the proof of \cite[(1.89)]{B-CMP122a} that the
corresponding configuration is a gauge transform of $U_0$. The seam step enters only through
hypothesis~(g), and no extension of $\bar\varpi$ to the core is used. At complex seeds
$\operatorname{Ad}_{\bar\varpi}$ takes the real datum $\mathsf{b}$ on $\mathfrak{B}_0$ into the complexified
Lie algebra; the lemma estimates sizes only. Hypothesis~(f) is one of the two forms in which the ceiling on
the coupling enters, the other being the smallness of $C_4^{\sharp}\mathsf{s}_L$ used for the cross terms: the
non-unitary part of $\bar\varpi$ is at most $K(M)\alpha_{1,k}$, with $K(M)$ a constant depending on $M$ and
$\alpha_{1,k}$ Ba{\l}aban's radius.
\end{remark}

\begin{proposition}[The graded-region slice bound]\label{9.1:prop-graded-slice}
Let $k$ be the level of the step, and let $h$, $N$, $N_0$ and the level functions $m(\cdot)$,
$n_F(\cdot)$, $\lambda_{\mathfrak{S}}(\cdot)$ be as fixed earlier in this chapter; here $N$ is the
index of the last stage, whose table is $\mathsf{T}_N$, and not the $N$ of $\mathrm{SU}(N)$.
Assume the following.
\begin{itemize}
\item[(1)] The hypotheses of Lemma~\ref{9.1:lem-split-rotation} (splitting the rotation in the graded
region); in particular its conclusion is available. Thus for every root seed $q_0$, every vertex
of $\mathfrak{P}_{\mathsf{c}}$, every $q\in\mathfrak{Q}''$ and every $y$ with
$\tilde\Delta(y)\subset Z''^{(1)}_k$, that lemma provides the object
$q^\sharp=e^{i\eta\mathcal{Z}^\sharp}q_0$ with $\mathcal{Z}^\sharp=\mathcal{Z}^{(k)}+\mathcal{Z}^{\mathrm{re}}$,
the level-$k$ bound on $\mathcal{Z}^{(k)}$ and the majorant $\mathsf{m}^{\mathrm{re}\sharp}$ of
$\mathcal{Z}^{\mathrm{re}}$.
\item[(2)] The standing hypothesis of this chapter on slices.
\item[(2$'$)] The standing hypothesis of this chapter on the units of a leaf.
\item[(3)] The two ceilings
\begin{align}
2^{20}B_\natural c_1C_T\,L^4C_4^\sharp K_\Lambda'(1+\mathsf{r}_1)\,\ell_\circ^{-2N_0}&\le 1,
\label{9.1:eq-graded-slice-1}\\
2^{21}B_\natural c_1C_T\Bigl[B^{\mathsf{s}}C_sc_1C_1(1+N^{1/2})\frac{\delta'_k}{\alpha_{0,k}}
+C_{\mathrm{seam}}\frac{\varepsilon_h}{\alpha_{0,h}}\Bigr]&\le 1,
\label{9.1:eq-graded-slice-2}
\end{align}
where $C_{\mathrm{seam}}$ is the constant with which the seam bound of this chapter holds.
These are ceilings in the sense of the order of choice (Notation~\ref{2.3:not-stages-first}),
imposed before $\gamma$ is chosen.
\end{itemize}
Let $(F,X,j)$ be a leaf. Let $\mathfrak{S}=\boldsymbol{\Omega}(X^+)$ be taken over a stage label
$\iota_\nu\le m$, which may be any such label (for $\nu=N$ this is the reading at the last stage),
and truncated at $j$, and assume $\tilde\Delta(y)\cap\mathfrak{K}=\emptyset$ for every
$y\in\mathfrak{B}$. Let $q_0$ be a root seed of the $\mathbf{R}$-born output $\Phi$, let
$q\in\mathfrak{Q}''$ at any vertex of $\mathfrak{P}_{\mathsf{c}}$, and let
$x\in Z''^{(1)}_k\cap X$.

Then $S(q^\sharp)$ is defined and holomorphic, and it equals $q_0{\upharpoonright}X$ at vanishing
sources. The quantity $E_5$ and the cube clause are those of the entry given by
Lemma~\ref{9.1:lem-root-slice}. For every free member $a$:
\begin{itemize}
\item[(a)] if $\lambda_{\mathfrak{S}}(x)=m(x)$, then
\begin{equation}\label{9.1:eq-graded-slice-3}
E_a(S(q^\sharp))(x)\le \tfrac{1}{960}\,\mathsf{T}^a_N(x);
\end{equation}
\item[(b)] if $n_F(x)\le\lambda_{\mathfrak{S}}(x)<m(x)$ (an old leaf), then
\begin{equation}\label{9.1:eq-graded-slice-4}
E_a(S(q^\sharp))(x)\le E_a(q_0)(x)+\tfrac{1}{1440}\,\mathsf{u}^a_F(x),
\end{equation}
and
\begin{equation}\label{9.1:eq-graded-slice-5}
E_a(q_0)(x)\le\min\bigl\{2^{-16}\mathsf{T}^a_N(x),\,2^{-18}\mathsf{u}^a_F(x)\bigr\}.
\end{equation}
\end{itemize}
\end{proposition}

\begin{proof}
Throughout we write $m=m(x)$ and $n_F=n_F(x)$. The proof follows the proof of
Lemma~\ref{9.1:lem-chain-slices} (slices of chain objects on the three regions). The difference is
that the displacement field is the field $\mathcal{Z}^\sharp$ of
Lemma~\ref{9.1:lem-split-rotation}, with its two parts estimated separately.

\emph{Applicability.} The slice-displacement lemma of this chapter, of which
Lemma~\ref{9.1:lem-noncritical-displacement} is the form between non-critical points, applies to
$q_0$ and $q^\sharp$. Indeed, $\ell_{\mathfrak{S}}\le\ell_m$; the fields $A_e$ and
$\mathcal{Z}^\sharp$ have size at most $2C_4^\sharp\mathsf{s}_L$ in their own units; and the
ceiling $2C_4^\sharp\mathsf{s}_L\le\min\{c_4,a_{R'}\}$, one of the ceilings imposed before $\gamma$
in the order of choice, gives
\begin{equation*}
\ell_{\mathfrak{S}}\sup\bigl(|A_e|+|\mathcal{Z}^\sharp|\bigr)
\le\ell_m\sup\bigl(|A_e|+|\mathcal{Z}^\sharp|\bigr)\le 2C_4^\sharp\mathsf{s}_L\le c_4 .
\end{equation*}
The slice-displacement lemma gives that
$S(q^\sharp)$ is defined and holomorphic and equals $q_0{\upharpoonright}X$ at vanishing sources,
exactly as in the proof of Lemma~\ref{9.1:lem-chain-slices}. By Lemma~\ref{9.1:lem-root-slice}, the
unitary part of the slice is $U_c$ at every object over $q_0$, so $E_5$ and the cube clause are those
of the entry.

\emph{Decomposition.} Apply Lemma~\ref{9.1:lem-noncritical-displacement} (slice displacement between
non-critical points) with $\mathbf{U}_1=q_0$ and $\mathcal{Y}=\mathcal{Z}^\sharp$. It bounds the
members of $S(q^\sharp)-S(q_0)$ by $4B_\natural(|\mathcal{Z}^\sharp|_\cdot)^{\approx}$; read in the
member $a$ at $x$, as in the site-blind slice lemma (Lemma~\ref{9.1:lem-site-blind-slice}), this
gives the contribution $16B_\natural\ell_{\mathfrak{S}}(x)^{-e_a}(|\mathcal{Z}^\sharp|_\cdot)^{\approx}(x)$.
The seminorm $|\cdot|_y$ and the smearing $(\cdot)^{\approx}$ are
sublinear. Hence the majorant splits into a part carried by $\mathcal{Z}^{(k)}$, whose value in the
member $a$ at $x$ we denote $\mu^{(k)}_a(x)$, and a part carried by $\mathcal{Z}^{\mathrm{re}}$,
denoted $\mu^{\mathrm{re}}_a(x)$. We obtain
\begin{equation}\label{9.1:eq-graded-slice-6}
E_a(S(q^\sharp))(x)\le E_a(q_0)(x)+\mu^{(k)}_a(x)+\mu^{\mathrm{re}}_a(x).
\end{equation}
Here the entry $E_a(q_0)(x)$ is the member $a$ at $x$ of the slice of the root seed, as in
Lemma~\ref{9.1:lem-chain-slices}.

\emph{The entry.} By Lemma~\ref{9.1:lem-root-slice}, $E_a(q_0)(x)<\mathsf{u}^a_k(x)$, and we use it
with $\ell_\circ^{2N_0}\ge 2^{20}$, which holds by the lower bound on $N_0$ in the order of choice.
By Lemma~\ref{9.1:lem-graded-thresholds} (graded thresholds
exceed the top-level unit), $\mathsf{T}^a_N(x)\ge\Gamma_\Lambda\mathsf{u}^a_k(x)$ with
$\Gamma_\Lambda=\tfrac{4}{45}\ell_\circ^{2N_0}$, that is,
$\mathsf{u}^a_k\le\tfrac{45}{4}\ell_\circ^{-2N_0}\mathsf{T}^a_N$. At the points of case (b), the
comparison of the level-$k$ unit with the leaf's unit at such points, on which the last bound of
Lemma~\ref{9.1:lem-root-slice} rests, gives
$\mathsf{u}^a_k\le\tfrac94\ell_\circ^{-2N_0}\mathsf{u}^a_F$. Since $\tfrac{45}{4}<2^4$ and
$\tfrac94<2^2$, these bounds give
\begin{equation*}
E_a(q_0)(x)\le \tfrac{45}{4}\,2^{-20}\,\mathsf{T}^a_N(x)\le 2^{-16}\,\mathsf{T}^a_N(x).
\end{equation*}
At the points of case (b) they give in addition
\begin{equation*}
E_a(q_0)(x)\le\tfrac94\,2^{-20}\,\mathsf{u}^a_F(x)\le 2^{-18}\,\mathsf{u}^a_F(x).
\end{equation*}
This is \eqref{9.1:eq-graded-slice-5}.

\emph{The level-$k$ part.} A level-$k$ absolute majorant holds on all of $\mathfrak{B}$:
\begin{equation*}
\sup_{\tilde\Delta(y)}|\mathcal{Z}^{(k)}|\le \mathsf{m}/\ell_k \quad\text{for every } y\in\mathfrak{B},
\qquad \mathsf{m}:=L^4C_4^\sharp\mathsf{s}_L .
\end{equation*}
The bound rests on three facts.
\begin{itemize}
\item Inside $Z''^{(1)}_k$, Lemma~\ref{9.1:lem-split-rotation} gives the bound
$C_4^\sharp\mathsf{s}_L/\ell_k$, which is at most $\mathsf{m}/\ell_k$ since $L^4\ge1$.
\item Off $Z''_k$, the sources of $\mathcal{Z}^{(k)}$ are of level $k$ and carry the profile of the
chain-field bound assumed in Lemma~\ref{9.1:lem-split-rotation}.
\item On a shell of label level $k'<k$, the decay factor of this profile satisfies
$e^{-\delta_1 d(y,\bar\Lambda)/4}\le L^4L^{-4(k-k')}$, where $\bar\Lambda$ is the set from which
the distance in this profile is measured.
\end{itemize}
Together they give the level-$k$ majorant above, the factor $L^4$ in $\mathsf{m}$ coming from the
third.
Lemma~\ref{9.1:lem-top-rotations} (top-level rotations are small in the graded region) applied with
this $\mathsf{m}$ bounds $\mu^{(k)}_a(x)$ in two cases:
\begin{itemize}
\item at the points of case (a), by
$180B_\natural c_1C_T(\mathsf{m}/\alpha_{a,k})\ell_\circ^{-2N_0}\mathsf{T}^a_N(x)$;
\item at the points of case (b), by
$36B_\natural c_1C_T(\mathsf{m}/\alpha_{a,k})\ell_\circ^{-2N_0}\mathsf{u}^a_F(x)$.
\end{itemize}
We have $\mathsf{s}_L=K_\Lambda'(1+\mathsf{r}_1)\alpha_{0,k}$ and $\alpha_{a,k}\ge\alpha_{0,k}$.
Hence, by the first ceiling \eqref{9.1:eq-graded-slice-1}, these bounds are at most
$180\cdot2^{-20}\mathsf{T}^a_N(x)$ and $36\cdot2^{-20}\mathsf{u}^a_F(x)$ respectively. Since
$180\cdot2880=518400<2^{20}$, both are at most $\tfrac{1}{2880}$ of $\mathsf{T}^a_N(x)$,
respectively of $\mathsf{u}^a_F(x)$.

\emph{The remainder part.} The majorant $\mathsf{m}^{\mathrm{re}\sharp}$ of
Lemma~\ref{9.1:lem-split-rotation} is slow, by part (i) of
Lemma~\ref{9.1:lem-top-source-decay} (decay of a profile sourced at the top levels). We therefore
estimate $\mu^{\mathrm{re}}_a$ by the same line as in the proof of
Lemma~\ref{9.1:lem-chain-slices}.

In case (a), $\ell_{\mathfrak{S}}=\ell_m$. Writing $\varphi=\varphi^{(N)}_m$ for the stage
coefficient, that line gives
\begin{equation*}
\frac{\mu^{\mathrm{re}}_a(x)}{\mathsf{T}^a_N(x)}
\le\frac{16B_\natural c_1C_T\,\mathsf{m}^{\mathrm{re}\sharp}}{\varphi\,\alpha_{a,m}}.
\end{equation*}
Insert the definition of $\mathsf{m}^{\mathrm{re}\sharp}$ from Lemma~\ref{9.1:lem-split-rotation},
and use $\varphi\ge\tfrac15$, $\alpha_{a,m}\ge\alpha_{0,m}$,
$\alpha_{0,k}\le\tfrac94N^{1/2}\alpha_{0,m}$ and $\alpha_{0,h}\le2\alpha_{0,m}$; the last two hold
by the comparison of the radii $\alpha_{0,\cdot}$ at different levels, the first of them together
with $k-m\le N$. This gives
\begin{equation}\label{9.1:eq-graded-slice-7}
\frac{\mu^{\mathrm{re}}_a(x)}{\mathsf{T}^a_N(x)}
\le B_\natural c_1C_T\Bigl[720B^{\mathsf{s}}C_sc_1C_1N^{1/2}\frac{\delta'_k}{\alpha_{0,k}}
+320C_{\mathrm{seam}}\frac{\varepsilon_h}{\alpha_{0,h}}\Bigr].
\end{equation}
Since $N^{1/2}\le1+N^{1/2}$ and $320\le720$, the right side of \eqref{9.1:eq-graded-slice-7} is at
most $720$ times the left member of \eqref{9.1:eq-graded-slice-2} divided by $2^{21}$, hence at most
$720\cdot2^{-21}$. Since $720\cdot2880=2073600<2^{21}$, this is at most $\tfrac1{2880}$.

In case (b), the line of the proof of Lemma~\ref{9.1:lem-chain-slices} at its points of type (b) is
used, with the length $\ell_j$ replaced by $\ell_{\lambda_{\mathfrak{S}}(x)}$. In the part of
$\mathsf{m}^{\mathrm{re}\sharp}$ carrying $\delta'_k$, the factor $\tfrac94N^{1/2}$ and the factor
coming from $\varphi$ are replaced by
\begin{equation*}
\tfrac94(k-n_F)^{1/2}L^{-(m-n_F)}\le\tfrac94\,(N^{1/2}+1),
\end{equation*}
since $(k-n_F)^{1/2}\le(k-m)^{1/2}+(m-n_F)^{1/2}\le N^{1/2}+(m-n_F)^{1/2}$, while
$L^{-(m-n_F)}\le1$ and $u^{1/2}L^{-u}\le1$ for $u=m-n_F\ge1$.
In the seam part, the factor is
\begin{equation*}
\tfrac94\bigl((h-n_F)_+\vee1\bigr)^{1/2}L^{-(m-n_F)}\le\tfrac94 .
\end{equation*}
The corresponding factors of case (a) are $5\cdot\tfrac94N^{1/2}$ and $5\cdot2$. Since
$\tfrac94(N^{1/2}+1)\le\tfrac{45}{4}(1+N^{1/2})$ and $\tfrac94\le10$, the same computation as in
case (a) gives that $\mu^{\mathrm{re}}_a(x)/\mathsf{u}^a_F(x)$ is bounded by
\begin{equation*}
B_\natural c_1C_T\Bigl[720B^{\mathsf{s}}C_sc_1C_1(1+N^{1/2})\frac{\delta'_k}{\alpha_{0,k}}
+320C_{\mathrm{seam}}\frac{\varepsilon_h}{\alpha_{0,h}}\Bigr].
\end{equation*}
By \eqref{9.1:eq-graded-slice-2}, this is at most $720\cdot2^{-21}\le\tfrac1{2880}$.

\emph{Conclusion.} In case (a), \eqref{9.1:eq-graded-slice-6} and the three bounds above give
\begin{equation*}
E_a(S(q^\sharp))(x)\le\Bigl(\tfrac{45}{4}\,2^{-20}+\tfrac{2}{2880}\Bigr)\mathsf{T}^a_N(x).
\end{equation*}
Since $\tfrac{45}{4}\cdot2880=32400<2^{20}$, the first term is below $\tfrac1{2880}$. The total is
therefore at most $\tfrac{3}{2880}\mathsf{T}^a_N(x)=\tfrac1{960}\mathsf{T}^a_N(x)$, which is
\eqref{9.1:eq-graded-slice-3}. In case (b), \eqref{9.1:eq-graded-slice-6} and the bounds on
$\mu^{(k)}_a(x)$ and $\mu^{\mathrm{re}}_a(x)$ give
\begin{equation*}
E_a(S(q^\sharp))(x)\le E_a(q_0)(x)+\tfrac{2}{2880}\mathsf{u}^a_F(x),
\end{equation*}
which is \eqref{9.1:eq-graded-slice-4}. Together with \eqref{9.1:eq-graded-slice-5}, proved above,
this completes case (b).
\end{proof}

\begin{remark}[Size of the two ceilings]
The left members of the ceilings \eqref{9.1:eq-graded-slice-1} and \eqref{9.1:eq-graded-slice-2}
are bounded by constants times negative powers of $\log x_k$. In the first, by
Lemma~\ref{9.1:lem-graded-thresholds},
$\ell_\circ^{-2N_0}\le(\log x_k)^{-2r_{pr}\log_L\ell_\circ}$. In the second, $N=\check{R}_k$, the
scaffold radius at level $k$ (Ba{\l}aban's radius $R_k$), and
$\check{R}_k\le L(\log x_k)^{r_{pr}}$ by \cite[(2.5)]{B-CMP119}; the ratios
$\delta'_k/\alpha_{0,k}$ and $\varepsilon_h/\alpha_{0,h}$ are constant multiples of
$(\log x_k)^{p_1-q}$ and $(\log x_h)^{p_0-q}$, and $x_h\ge x_k$. The exponents are then
$r_{pr}/2+p_1-q$ and $p_0-q$; here $p_0$, $p_1$ and $q$ are exponents among the auxiliary
parameters $\mathfrak{p}$ (Notation~\ref{2.1:not-prefix}), and $q$ is not the chain object $q$
above. Both exponents are negative under the one-sided floors on the exponents $p_0$, $p_1$, $q$ and
$r_{pr}$ fixed in the order of choice, which belong to the window of exponents used in
\cite{B-CMP122a}. No size determined by the seed occurs in the second ceiling.
\end{remark}

\begin{theorem}[The slice theorem at every site]\label{9.1:thm-slice-theorem}
Let $\mathfrak{O}$ and $\mathfrak{A}$ be the regions of the stage structure on whose union the field
bounds are stated, let $\mathfrak{K}$ be the raw core, let
$Z''^{(1)}_k:=Z''_k\setminus\mathfrak{K}$, where $Z''_k$ is the region of Ba{\l}aban's large-field
construction at level $k$, and let $Z^{\mathrm{on}}$ be the second region of a run at the site (P),
beside $\mathfrak{O}$ and, when present, a third region $W$; its parts in $\mathfrak{A}$ and in
$Z''^{(1)}_k$ are treated separately in the proof. Let $\Lambda$ denote the small-field region.
The statement distinguishes three sites, described by their data in the list below: the site (T),
at which the seed is moved by the dilated shift of the $\mathbf{T}$-step; the site (L), at which
the terminal read is made from the root by the links and the $\mathsf{t}$-shift; and the site (P),
at which runs of moves are read at a stage $\nu$.
Let $N$ be the
last stage of the stage structure (a stage index, not the $N$ of $\mathrm{SU}(N)$), and let $N_0$ be
the offset subtracted from $N$ in that structure. The point classes $({=})$ and $(\mathrm{up})$ are
those of the first slice theorem, the statement whose clause (e), parts (ii) and (iii) and
accounting clause are used below. The index $a$ of the table members $E_a$ and of $\mathsf{T}^a_N$ ranges over
the members of the stage table; it is unrelated to clause (a) below and to the $(a)$-points.
The $(a)$-points and $(b)$-points are the points at which clause (a),
respectively clause (b), of the sharpened $\Lambda$-estimate of hypothesis (6) below applies. Assume
the following.
\begin{enumerate}
\item[(1)] The standing conditions of this chapter in their closed form, together with their
profile condition and their slice condition.
\item[(2)] The floors of the summation bound.
\item[(3)] The unit condition.
\item[(4)] The response bound on $\Lambda$.
\item[(5)] The corollary that enters the sharpened $\Lambda$-estimate.
\item[(6)] The sharpened $\Lambda$-estimate, with its clauses (a) and (b), stated on gauge frames
$\mathrm{fr}\in\mathfrak{F}_G$ and derived from the response bound on $\Lambda$, the field bounds on
$\mathfrak{O}\cup\mathfrak{A}$, the statements at the junction surface $S_c$ (the seam) and a
further condition of that estimate, stated with it. Its seam step, and the complementary core
estimate for leaves with $X^{+}\cap\mathfrak{K}\neq\emptyset$, are separate statements; they are
not proved here and are used only where stated.
\item[(7)] The specialisation $c\mapsto 1$ of the parameter $c$; this $c$ is not the component
index of the class of list bounds.
\item[(8)] The arrest statement on the stage tables (the one carrying the table coefficient
$\varphi_{\mathrm{arrest}}$).
\item[(9)] The ceiling conditions: the ceiling $K_R^{\sharp}$; the ceiling on $c_M$ relative to
$\mathfrak{O}$; the ceiling attached to the site (L) and its circle $\mathsf{w}_L$; the ceiling on
$\mathfrak{A}$; the two sharpened ceilings on $\Lambda$; and
\begin{equation*}
2C_4^{\sharp}\mathsf{s}_L\le\min\{c_4,a_{R'}\},
\end{equation*}
where $\mathsf{s}_L$ is the seed size, the one size in which the response bound on $\Lambda$ is
written; $C_4^{\sharp}$ is a constant independent of the level $k$, taking the place of the
level-dependent constants of that letter in the standing conditions; $a_{R'}$ is a radius of the
seed's Landau and data coordinates; and $c_4$ is the constant of that name fixed with the
conditions of this chapter, so that $c_4$ and $a_{R'}$ are the two bounds below which
$2C_4^{\sharp}\mathsf{s}_L$ must lie.
\item[(10)] The dilation circles at the sites, and the circle condition in the units of $F$ at
$(b)$-points.
\end{enumerate}
Let $P=(\mathbf{U},\mathbf{J})$ be a point of the complex parameter polydisc
$\mathfrak{P}_{\mathsf{c}}$ of a shape at any site, over a seed of its output
$\Phi$ with slice in $\mathfrak{D}_\Phi[7/8]$, and read by a leaf $(F,X,j)$ of the expansion, with
term $F$ and region $X$, such that $X$ misses
the raw core $\mathfrak{K}$ and, on $Z''_k$, for every element $y$ of the family $\mathfrak{B}$ of
the truncated determining set $\boldsymbol{\Omega}(X^{+})$ the set $\tilde\Delta(y)$ attached to $y$
is disjoint from
$\mathfrak{K}$. (The leaves failing the last condition are those of the complementary core estimate
named in (6).) Then $\tilde F(X;P)=F(X;S(\mathbf{U}))$; this function is
holomorphic in $P$ on $\mathfrak{P}_{\mathsf{c}}$ and exact at vanishing increments; and
$\mathrm{rd}^{\mathsf{m}}\mathbf{U}^{\mathsf{m}}\in\mathfrak{T}^{\mathsf{m}}[0](X)$, where the
superscript $\mathsf{m}$ is the run index, $\mathrm{rd}^{\mathsf{m}}\mathbf{U}^{\mathsf{m}}$ is the
reading of the configuration of run $\mathsf{m}$, and $\mathfrak{T}^{\mathsf{m}}[0](X)$ is the
admissible data class on the chart of that run with room fraction $0$ on $X$. Precisely, with
the site's initial configuration $U_e$ and initial table $\mathsf{T}_e$, its composite move
$\mathcal{Y}$, and the table $\mathsf{T}_r$ against which points of class $({=})$ are read,
\begin{equation}\label{9.1:eq-slice-theorem-1}
E_a(S(\mathbf{U}))\le
\begin{cases}
\mathsf{T}_e+\tfrac16(\mathsf{T}_r-\mathsf{T}_e) & \text{at points of class } ({=}),\\[2pt]
\bigl(\tfrac34+\tfrac1{200}\bigr)\mathsf{u}_F & \text{at points of class } (\mathrm{up}),
\end{cases}
\end{equation}
where the sites and their data are as follows.
\begin{itemize}
\item The site (T): $U_e$ is the seed, $\mathsf{T}_e=\mathsf{c}_\Phi^{+}\mathsf{u}_\Phi$,
$\mathcal{Y}$ is the dilated shift, and at $({=})$
\begin{equation*}
\mathsf{T}_r=\bigl(f_{k,n}-\tfrac12(1-\theta_S)\beta\,2^{-(k+1-n)}\bigr)\mathsf{u},
\end{equation*}
the $f_{k,n}$ being the coefficients of the stage tables, $\beta$ being the fixed fraction that
enters the table coefficients (as in the table coefficient $\varphi_{\min}=1-4\beta$), and $n$ being
the level of the site (T) (a level, not the order of a Schwinger function).
\item The site (L), terminal read: $U_e$ is the root, the configuration in whose chart the slice
$S(\mathbf{U})$ is read; $\mathsf{T}_e$ is
$(f_{k,m}-\frac18\varrho_\Phi)\mathsf{u}$ on $\mathfrak{O}$, where $m$ is the level of the read (a
level, not the torus exponent), and is at most $\mathsf{T}_N/32$ on
$Z''^{(1)}_k$ by the bound of the stage structure that holds on leaves with
$X^{+}\cap\mathfrak{K}=\emptyset$; $\mathcal{Y}$ consists of all links and
the $\mathsf{t}$-shift. The reference table at $({=})$ is $f_{k,m}\mathsf{u}$ on $\mathfrak{O}$; on
$Z''^{(1)}_k$ it is $\frac12\mathsf{T}_N$ at $(a)$-points, the chain contributing at most
$\frac{3}{32}\mathsf{T}_N$ and the bracket at most the rest; at $(b)$-points the read is treated as
at the points of class $(\mathrm{up})$ on $\mathfrak{A}$, and the bound there is
\begin{equation*}
E_a(S(\mathbf{U}))\le\bigl(\tfrac{3}{32}+c_{\mathrm{br}}+\tfrac1{200}\bigr)\mathsf{u}_F,
\end{equation*}
where $c_{\mathrm{br}}$ is the coefficient of $\mathsf{u}_F$ contributed by the bracket.
\item A run at the site (P) over alternative (ii) of clause (e) of the first slice theorem:
$U_e$ is the root,
$\mathsf{T}_e=\mathsf{T}^{\flat}_N$, where $\mathsf{T}^{\flat}_\nu$ denotes the flat stage table at
stage $\nu$ (a table distinct from the stage table $\mathsf{T}_\nu$), $\mathcal{Y}$ consists of the
links followed by the moves at
the stages in $[\nu,N\mathclose[$, and $\mathsf{T}_r=\mathsf{T}^{\flat}_\nu$, $\nu$ being the stage
at which the run is read.
\item A run at the site (P) over alternative (i) of clause (e) of the first slice theorem:
$U_e$ is the seed,
$\mathsf{T}_e=\mathsf{T}^{\flat}_{n_e}$, the flat stage table at the seed's stage $n_e$,
$\mathcal{Y}$
consists of the moves, and $\mathsf{T}_r=\mathsf{T}^{\flat}_\nu$.
\item A box $Z\supseteq X^{+}$ for which the third box condition holds: $U_e$, $\mathsf{T}_e$,
$\mathcal{Y}$ and $\mathsf{T}_r$ are those of the underlying site.
\end{itemize}
No constant in this statement depends on $g$, $k$, $K$, $N$, $N_0$, $\check{R}_k$, on the level
$h:=k-N$, on $r:=m-n_F$, the difference between the level $m$ of the read and the level $n_F$ of
$F$, on the
ages of large-field regions, on the depth, or on the volume.
\end{theorem}

\begin{proof}
All bounds are obtained from the basic slice bound and from part (b$'$) of the slice-domination
statement, adding the majorants of the separate contributions. The items of the list above rest on
the following statements. At the site (T): the accounting statement at (T) and the schedule
comparison. At the site (L): the field bounds of the stage structure, the entry-and-link lemma and
the circle condition of hypothesis (10); part (c) of the accounting statement on $\mathfrak{O}$; the
schedule comparison; and the sharpened $\Lambda$-estimate of hypothesis (6). For the runs at (P), over
either alternative: on $\mathfrak{O}$, the accounting statement on $\mathfrak{O}$; on
$Z^{\mathrm{on}}$, the entry-and-link lemma on $\mathfrak{A}$, the sharpened $\Lambda$-estimate of
hypothesis (6) on $Z''^{(1)}_k$, and parts (ii) and (iii) of the first slice theorem.
At a box $Z\supseteq X^{+}$: the slice-domination
statement at $\mathfrak{S}:=\boldsymbol{\Omega}(Z)$. We now carry out the site-specific steps.

\emph{The site (L) on $\mathfrak{O}$.} Here the room is $\mathsf{G}=\frac18\varrho_\Phi\mathsf{u}$
and the profile of the site, denoted $\mathsf{m}$ here (a profile, unrelated to the run index
$\mathsf{m}$ of the statement), is
\begin{equation*}
\mathsf{m}=2\bigl(\mathsf{m}_L+(1+|\mathsf{t}|)\mathsf{m}_p\bigr),
\end{equation*}
where $\mathsf{m}_L$ is the layer profile, $\mathsf{t}$ is the bracket parameter, and $\mathsf{m}_p$
is the second profile, weighted by $1+|\mathsf{t}|$. The ceiling attached to the site (L) in
hypothesis (9), together with the definition of the circle
$\mathsf{w}_L$, gives condition (A) of the accounting statements. Condition (S) of the same
statements follows
from $f_{k,m}\mathsf{u}\le\mathsf{T}^{\flat}_0$ and part (c) of the accounting statement on
$\mathfrak{O}$.

\emph{Terminal reads on $\mathfrak{A}$.} A terminal read on $\mathfrak{A}$ has points of class
$(\mathrm{up})$ only, since the level $k$ of the label exceeds the level $n_F$ of $F$. The
entry-and-link lemma and part (b$'$) of the slice-domination statement give there the bound
$\bigl(\tfrac16+\tfrac1{960}+\tfrac1{200}\bigr)\mathsf{u}_F$, which lies below
$\bigl(\tfrac34+\tfrac1{200}\bigr)\mathsf{u}_F$ because $\tfrac16+\tfrac1{960}<\tfrac34$.

\emph{Terminal reads on $Z''^{(1)}_k=Z''_k\setminus\mathfrak{K}$.} We use the sharpened
$\Lambda$-estimate of hypothesis (6) and part (b$'$) of the slice-domination statement, for
leaves with $\tilde\Delta(y)\cap\mathfrak{K}=\emptyset$ for every $y\in\mathfrak{B}$; the other
leaves are those of the complementary core estimate. At $(a)$-points,
\begin{equation*}
\mathsf{T}_N\le\tfrac34\bar u_F\le(\mathsf{c}_F^{+}-\varrho_F)\bar u_F,
\end{equation*}
where the comparison is made, by the schedule comparison, with the table of $F$ in
Ba{\l}aban's construction, whose coefficient is $\mathsf{c}_F^{+}$, and where
$\bar u_F:=\ell_\circ^{2(n_F-k_0)_+}\mathsf{u}_F$, with $k_0$ the reference level of this weight;
$\bar u_F$ equals $\mathsf{u}_F$ for old leaves. At
$(b)$-points the units of $F$ are used as on $\mathfrak{A}$.

\emph{Runs on $Z^{\mathrm{on}}$.} On $\mathfrak{A}$ the entry-and-link lemma places the entry in
$\mathsf{T}_N$, and the links cost at most $\mathsf{G}/24$ at points of class $({=})$. On
$Z''^{(1)}_k$ we apply the sharpened $\Lambda$-estimate at the stage label $\iota_\nu\le m$ of the
read, the stage label $\iota_\nu(x)$ of a point $x$ being its label at stage $\nu$, and
distinguish two kinds of points of class $({=})$: a point $x$ is raised if
$\iota_\nu(x)=n_F<m(x)$, where $m(x)$ is the level of the read at $x$, so that the point was raised
at some stage $i_1\in[\nu,N\mathclose[$; it is unraised otherwise.

At an unraised point, clause (a) bounds the total by $\mathsf{T}^a_N/960\le\mathsf{T}_e$; the links
cost nothing, and the moves are covered by the accounting clause of the first slice theorem.

At a raised point, raised at the stage $i_1\in[\nu,N\mathclose[$, clause (b) gives the entry bound
$\mathsf{T}_N$ and the link bound
\begin{equation*}
\frac{\mathsf{u}_F}{1440}\le\frac{\mathrm{rm}_{i_1}}{24}\le\frac{\mathsf{G}}{24}.
\end{equation*}
Here we used the statement on the room increments, which gives
\begin{equation*}
\mathsf{G}\ge\mathrm{rm}_{i_1}\ge\frac{\bar u_{i_1}}{40}\ge\frac{\mathsf{u}_F}{40},
\end{equation*}
where $\bar u_{i_1}$ is the local unit at stage $i_1$,
and hence $\mathrm{rm}_{i_1}/24\ge\mathsf{u}_F/960\ge\mathsf{u}_F/1440$. In it,
$\bar u_{i_1}\ge\mathsf{T}_\nu$ on $Z^{\mathrm{on}}$, because the unit is unchanged until the first
raise and the coefficients in part (a) of the flat increment statement are at most $1$.

From this point on, the proof of parts (ii) and (iii) of the first slice theorem applies as
written, with $\frac18$ replaced by $\frac16=\frac18+\frac1{24}$.

\emph{Runs on $\mathfrak{O}$.} The layer profiles $\mathsf{m}_L$ of the links are covered by the
margin $\bigl(\frac12\theta_S+\frac18\theta_\rho\bigr)\beta\,2^{-s}$ of $\mathsf{T}^{\flat}_N$, where
$s$ is the difference of levels carried by this margin (not the depth below the cutoff), through
the ceiling attached to the site (L) in hypothesis (9), since
\begin{equation*}
\tfrac{9}{256}+\tfrac1{16}=\tfrac{25}{256}\le\tfrac18.
\end{equation*}
The moves are covered by the room terms $\sum_i\mathsf{g}_i$ of part (b) of the accounting statement
on $\mathfrak{O}$. The two contributions are added; part (b$'$) of the slice-domination statement
covers the non-critical intermediate objects; and part (c) of the accounting statement on
$\mathfrak{O}$ concludes.

\emph{The third region.} On the third region $W$ of a run, if there is one, the bound is read from
the table for the third region in the statement on the regions of a run.

\emph{Points of class $(\mathrm{up})$ on $\mathfrak{O}$.} These are handled by scaling, as in the
accounting statement at the site (T).
\end{proof}

\begin{definition}[Normalisation of the paired runs at their own centres]
\label{9.1:not-centre-normalisation}
In this section we fix the pair of runs compared at the last lattice, the normalisation of their
sourced integrals, the labels, slots, classes and marks of their expansions, the loads and the dial
domain, and the floors on $n$ under which all of this is used.

\medskip\noindent\emph{The pair at the last lattice.}
Let $n$ be such that $\Upsilon_n\le\frac14$, let $\hat w\in\mathcal{N}'_w$, and put
$u:=\operatorname{Re}\hat w$. We use the local pair of runs of the uniqueness argument. The run
$\mathsf{A}$ has length $n$ and real centre $\zeta_*^n+u$, where $\zeta_*^n$ is the unique real alive
run of length $n$. The run $\mathsf{B}$ has length $n+1$ and centre $\zeta_*^{n+1}+\mathfrak{n}_n(u)$,
where $\mathfrak{n}_n$ is the matching map between consecutive cutoffs. Their continuum sources are
\begin{equation}\label{9.1:eq-centre-normalisation-1}
W^{\mathsf{A}}(z):=i\operatorname{Im}\hat w+\sum_l z_lf_l,\qquad
W^{\mathsf{B}}(z):=\mathfrak{n}_n\circ\bigl(u+W^{\mathsf{A}}(z)\bigr)-\mathfrak{n}_n(u).
\end{equation}
Levels are indexed by the index of $\mathsf{A}$, $j=n-s$ at depth $s$; the finest level belonging to
$\mathsf{B}$ alone is $j=-1$. The flow stops at the depth $\mathsf{w}=\mathsf{w}(g)$ fixed by the lemma
on the stop depth. The last lattice $\mathbb{T}_w$ has $N_w=2L^{m_0+\mathsf{w}}$ sites per side; its
configuration space is $\mathcal{V}_w:=G^{\mathrm{bonds}(\mathbb{T}_w)}$, $G=\mathrm{SU}(N)$, with the
product of Haar measures; and
\begin{equation*}
k:=n-\mathsf{w}
\end{equation*}
is the stop level. In this section $k$ always denotes $n-\mathsf{w}$. The coefficients
$z=(z_l)_l$ range over the polydisc $P_{UV}\supset P_{IR}$, on which
$\|W^{\mathsf{A}}(z)\|_{\flat}<R_0<r^{pr}_s$ as in the uniqueness argument; here $r^{pr}_s$ is the
radius of the admissible class of complex sources at depth $s$. The zero residual source
is the point $\hat w=u$ real, $z=0$, of the same family: there $\operatorname{Im}\hat w=0$, so
$W^{\mathsf{A}}(0)=0$ by \eqref{9.1:eq-centre-normalisation-1}, and then
$W^{\mathsf{B}}(0)=\mathfrak{n}_n(u)-\mathfrak{n}_n(u)=0$.

\medskip\noindent\emph{Normalisation at the centre.}
For $\mathsf{c}\in\{\mathsf{A},\mathsf{B}\}$ let $c^{\mathsf{c}}$ be the centre of the run, so that
$c^{\mathsf{A}}=u$ and $c^{\mathsf{B}}=\mathfrak{n}_n(u)$, and let $v^{\mathsf{c}}:=W^{\mathsf{c}}(z)$
be its residual source. Write $Z^{\mathsf{c}}(c^{\mathsf{c}};v^{\mathsf{c}})$ for the sourced integral
of the run $\mathsf{c}$ at the stop, at centre $c^{\mathsf{c}}$ and residual source $v^{\mathsf{c}}$,
and $e^{(j),\mathsf{c}}(X;\cdot)$ for its boundary exponent on the localisation domain $X$ of level
$j$; $\mathbf{D}^{(j)}$ is the set of localisation domains of level $j$, including those that wrap
around the torus. We put
\begin{equation}\label{9.1:eq-centre-normalisation-a}
\begin{aligned}
Z^{\mathsf{c}}(c^{\mathsf{c}};v^{\mathsf{c}})
&=e^{\mathcal{E}^{\mathsf{c}}(v^{\mathsf{c}})}\,\tilde Z^{\mathsf{c}}(v^{\mathsf{c}}),\\
\mathcal{E}^{\mathsf{c}}(v)
&:=\sum_j\sum_{X\in \mathbf{D}^{(j)}}
\bigl[e^{(j),\mathsf{c}}(X;c^{\mathsf{c}}+v)-e^{(j),\mathsf{c}}(X;c^{\mathsf{c}})\bigr],\\
\tilde Z^{\mathsf{c}}(v)
&=\sum_{h\in H^{\mathsf{c}}}\tau^{\mathsf{c}}_h(v)(\mathcal{V}_w).
\end{aligned}
\end{equation}
Thus $\mathcal{E}^{\mathsf{c}}$ collects the vacuum responses of all localisation domains, subtracted
at the run's own centre $c^{\mathsf{c}}$. The expansion of $\tilde Z^{\mathsf{c}}$ over the label set
$H^{\mathsf{c}}$ (defined below) is the one given by the correlation theorem and its torus bound;
$\tau^{\mathsf{c}}_h(v)$ is the $h$-th term
$\chi\,\mathbf{T}^{w}e^{A^{w}-\mathcal{E}^{\mathsf{c}}}$ of that expansion, read as a complex Borel
measure on $\mathcal{V}_w$; here $\chi$ is the product of the characteristic functions of the label,
$\mathbf{T}^{w}$ the large-field operation at the stop and $A^{w}$ the effective action at the stop.
On second-class components the measure keeps the factor $\delta_{G_i}$ that
Ba{\l}aban's construction \cite{B-CMP122b} attaches to the $i$-th second-class component.

The norm $\|\cdot\|$ of such a measure is the total variation in the label's chart. Let
$\Psi^{\mathsf{c}}_h$ be the label's chart substitution. With $V'$ the chart variable, it is given by
\begin{equation*}
\begin{aligned}
V&=V'\,V^{\mathsf{c}}_{\Lambda_i}&&\text{on a live second-class component }\Lambda_i,\\
V''&=V'\,M_{\mathfrak{B}''}(U^{\mathsf{c}}_0)&&\text{on ejected components}.
\end{aligned}
\end{equation*}
Here $V^{\mathsf{c}}_{\Lambda_i}$ is the field that the run $\mathsf{c}$ assigns to the live
second-class component $\Lambda_i$; the ejected components are those that the operation $\mathbf{R}$
removes, and $V''$ is the configuration on such a component; and $M_{\mathfrak{B}''}(U^{\mathsf{c}}_0)$
is the field on the bond set $\mathfrak{B}''$ of the operation $\mathbf{R}$ determined by the run's
background configuration $U^{\mathsf{c}}_0$.
Thus $\Psi^{\mathsf{c}}_h$ is a triangular bondwise right translation, and, by part (a) of the lemma
establishing this substitution, it preserves the product
Haar measure. We define $\|\tau^{\mathsf{c}}_h\|$ as the total variation of the measure
$(\Psi^{\mathsf{c}}_h)^*\tau^{\mathsf{c}}_h$ on the chart space, whose image under
$\Psi^{\mathsf{c}}_h$ is $\tau^{\mathsf{c}}_h$. Equivalently, it is $\int|\cdot|\,d\mu_h$ on the space
$(\Omega_h,\mu_h)$ of the label's common inner variables. This space is common to the two runs;
accordingly, for a difference, $\|\tau^{\mathsf{A}}_h-\tau^{\mathsf{B}}_h\|$ is the total variation of
$(\Psi^{\mathsf{A}}_h)^*\tau^{\mathsf{A}}_h-(\Psi^{\mathsf{B}}_h)^*\tau^{\mathsf{B}}_h$ on it. The
bounded Borel functions $\varphi_h$
integrated against these measures are functions on the chart space. In what follows, ``on
$\mathcal{V}_w$'' and ``Borel functions on $\mathcal{V}_w$'', said of the measures $\tau^{\mathsf{c}}_h$,
are to be read in this sense.

The chart is used for the following reason. In \cite{B-CMP122b} the factor $\delta_{G_i}$ is a
function of the chart variable $V'$, while $V_{\Lambda_i}^{\mathsf{c}}$ depends on the run. Hence, on
$\mathcal{V}_w$ itself, the measures $\tau^{\mathsf{A}}_h$ and $\tau^{\mathsf{B}}_h$ of a label with a
live second-class component are carried by different graphs and are mutually singular. On
$\mathcal{V}_w$ the total variation of $\tau^{\mathsf{A}}_h-\tau^{\mathsf{B}}_h$ is therefore the sum of
the total variations of $\tau^{\mathsf{A}}_h$ and of $\tau^{\mathsf{B}}_h$, and the sum over the labels
of these total variations does not become small when the two runs are compared, however large $n$
is. In the chart, by contrast, there is a reference
measure independent of the dial parameter of the dial domain below, and the holomorphy in that
parameter is read with respect to it. Total masses do not depend on the chart: $\Psi^{\mathsf{c}}_h$
is a bijection, so the total mass of $(\Psi^{\mathsf{c}}_h)^*\tau^{\mathsf{c}}_h$ equals
$\tau^{\mathsf{c}}_h(\mathcal{V}_w)$.

By the torus bound of the correlation theorem, the label $h$ carries, besides the constant of
Ba{\l}aban's construction, the term
\begin{equation}\label{9.1:eq-centre-normalisation-2}
-\sum_j\sum_{X\in \mathbf{D}^{(j)},\,X\not\subset\Lambda_j(h)}
\bigl[e^{(j),\mathsf{c}}(X;c^{\mathsf{c}}+v)-e^{(j),\mathsf{c}}(X;c^{\mathsf{c}})\bigr].
\end{equation}
Here $\Lambda_j(h)$ is the region of $h$ at level $j$ whose complement is the large-field region of $h$
at level $j$ (see the paragraph on labels below).
These are the responses that are subtracted globally in $\mathcal{E}^{\mathsf{c}}$ but are not
produced inside the large-field regions of $h$. By the torus bound of the correlation theorem, the
one-point pieces
$\log[\mathfrak{z}(a_j(y))/\mathfrak{z}(x_j)]$, with $\mathfrak{z}$ the one-point gauge weight,
$a_j(y)$ the position-dependent inverse coupling of level $j$ at $y$
(Definition~\ref{1.5:def-domains}) and $x_j=g_j^{-2}$, belong
to $e^{(j)}$ on the cube of $y$. Consequently the $\mathfrak{z}$-pieces occur in $\mathcal{E}^{\mathsf{c}}$
for every cube. They occur again, with the opposite sign, inside $\tau^{\mathsf{c}}_h$ for the cubes
outside $\Lambda_j(h)$, through \eqref{9.1:eq-centre-normalisation-2}.

There is one normalisation per run and centre: nothing that depends on $h$ or on the source is
rescaled. The zero source is the point $v=0$ of the same objects. In particular
$\mathcal{E}^{\mathsf{c}}(0)=0$ by \eqref{9.1:eq-centre-normalisation-a}, so that
$Z^{\mathsf{c}}(c^{\mathsf{c}};0)=\tilde Z^{\mathsf{c}}(0)$.

\medskip\noindent\emph{Labels.}
$H^{\mathsf{c}}$ is the set of admissible sequences $h=\{\Omega_j,\Lambda_j,S_j\}_j$ of the run
$\mathsf{c}$ at the stop, where $\Omega_j$, $\Lambda_j$, $S_j$ are the regions of level $j$ that
describe, in Ba{\l}aban's construction, the large-field history of the effective density
(Definition~\ref{1.5:def-densities}; here $\Omega_j$ is a region, not the vacuum), together with the
classes of the events of the operation $\mathbf{R}$. On
the prescribed scaffold $\Theta_X$ at the bare point $X=g^{-2}$, with depth-$s$ inverse coupling
$\theta_s=X+c_0s$, every threshold, integer, partition and tree is read from $\theta_s$. These are
common to the two runs, by the statements on the prescribed scaffold on which the uniqueness argument
rests.

The \emph{first index} of a label is the level of its first large-field region. A label is
\emph{common} if its first index is at least $1$. It is \emph{aligned} if its first index exceeds
$n_c$ and $\varpi(h)\le\varpi_0$, where $n_c$ is the alignment cut level, subject to $n_c\ge k_*$
with $k_*$ the threshold level of the statements on the scaffold, $\varpi(h)$ is the alignment
functional of the label $h$, and $\varpi_0=\frac1{80}$. The \emph{age} at
the stop of a live component $X$ of $h$, created at level $j(X)$, is $a(X)$. Thus
\begin{equation}\label{9.1:eq-centre-normalisation-b}
H_{al}:=\{h:\ \text{first index of }h>n_c,\ \varpi(h)\le\varpi_0\},\qquad a(X):=k-j(X).
\end{equation}
The set $H_{al}$ is one set for both runs, and the pairing $\pi$ of the labels of the two runs is the
identity on $H_{al}$.

Regions only grow until $\mathbf{R}$ removes them, by item (4) of the lemma on geometric forgetting of
old regions. Hence a label of first index at most $n_c$ that is alive at the stop has a live component
of age at least $k-n_c$. The labels of $\mathsf{B}$ that are not labels of $\mathsf{A}$ have first
index $0$, so a live component created there has age $k$. No label has first index $-1$. Since
$k-n_c\le k$, the lower bound $k-n_c$ on the age covers these labels as well.

\medskip\noindent\emph{Slots, classes and marks.}
At the last lattice only the outermost stage of the assembly is present: no inner stage of the
assembly remains open, and every piece entering the assembly is terminal, a \emph{leaf}. Thus
\begin{equation*}
\tau^{\mathsf{c}}_h=\mathcal{P}_h[x^{\mathsf{c}}]
\end{equation*}
is one assembly formula $\mathcal{P}_h$, carrying no run index, applied to the array $x^{\mathsf{c}}$ of
the run's stored numbers. The entries of $x^{\mathsf{c}}$ are the uses $(j,k',\cdot)$ of stored units,
a use of a stored unit being its reading at a level $k'\in[j,k]$ with anchor in zone $j$; here zone
$j$ is the set of cells attached to level $j$, and the anchor of a use is the cell at which the unit
is filed. The stored units are:
\begin{itemize}
\item the list slots $E^{(j)}$, $P^{\flat(j)}$, $Q^{(j)}$, $R^{(j)}$, $D^{(j)}$;
\item the stored boundary terms $\mathbf{B}^{(j)}$, $\mathbf{B}'^{(j)}$;
\item the zone-resolved stored derived terms;
\item the $\zeta$-, $b$- and weight entries;
\item the pieces that do not depend on $z$ but depend on the run, namely the vacuum constants, the
$\mathfrak{z}$-pieces described after \eqref{9.1:eq-centre-normalisation-2}, and the pieces
$\mathsf{s}^0$, $D^0$ of the real object.
\end{itemize}
The \emph{class} of a use with first two entries $(j,k')$ is the pair $z:=(j,k')$, and we put
$n(z):=j$. Here and below the letter $z$ denotes a class; it is not to be confused with the tuple
$z=(z_l)_l$ of source coefficients, which from here on enters only through the residual sources
$W^{\mathsf{c}}(z)$. Each slot has exactly one class; for the zone-resolved stored derived terms this
holds because their marks are split by zone before they are filed. The \emph{mark} of the class $z$
is the part of the difference of the two arrays carried by the slots of that class:
\begin{equation}\label{9.1:eq-centre-normalisation-c}
\begin{gathered}
z:=(j,k'),\qquad n(z):=j,\\
W^{(z)}:=(x^{\mathsf{B}}-x^{\mathsf{A}})\cdot\mathbb{1}_{\{\text{slots of class }z\}},\qquad
\sum_zW^{(z)}=x^{\mathsf{B}}-x^{\mathsf{A}}.
\end{gathered}
\end{equation}
The identity on the right holds because each slot has exactly one class, so that the indicator
functions $\mathbb{1}_{\{\text{slots of class }z\}}$ sum to $1$ on every slot. Entries of both kinds,
the unknown entries and the classical entries (measured below by the loads with superscript $u$ and
with superscript $c$ respectively), lie in the one mark of their class.

\medskip\noindent\emph{Loads and the dial domain.}
Let $\rho_{\cdot}$ be the cross-scale ratio of the correlation theorem, which satisfies
$\rho_{\cdot}^2\le3$, and $\hat\alpha$ the exponent of the correlation theorem entering its weights.
For integers $\nu\ge1$ put
\begin{equation*}
\varsigma_\nu:=\bigl[\nu^3(1+\rho_{\cdot})L^{-\hat\alpha(\nu-1)}+L^{-(\nu-1)/2}\bigr]\rho_{\cdot}^2 .
\end{equation*}
Write $s(i):=n-i$ for the depth of level $i$. On the scaffold $\Theta_X$ the weights $w_{j,i}$ of
the correlation theorem satisfy $w_{j,i}\le3\theta_{s(i)}/\theta_{s(j)}$, and we put
$w^{\mathbf{B}}_{j,i}:=L^{-(j-i)/2}$. With $q_\flat$ the fixed exponent of the logarithmic factors
(one of the auxiliary exponents chosen after $L$), put
\begin{equation*}
\mathfrak{g}_j:=g_j^2(\log x_j)^{2q_\flat}\le2\theta_s^{-1}(\log\theta_s)^{2q_\flat},
\qquad s=s(j).
\end{equation*}
Let $\mathsf{z}_i:=\sup|\zeta^{\mathsf{B}}_{(i)}-\zeta^{\mathsf{A}}_{(i)}|$ and
$\mathsf{b}_i:=\sup|b^{\mathsf{B}}_{(i)}-b^{\mathsf{A}}_{(i)}|$, the suprema being taken over the
parameter set of the pair and over the torus, and put
\begin{equation*}
\mathsf{z}^\sharp_j:=\sum_{i\le j}L^{-(j-i)}\bigl(\mathsf{z}_i+\mathsf{z}_{i-1}\bigr).
\end{equation*}
For a family $F$ of stored terms of level $i$, with decay rate $\kappa_F$, put
\begin{equation*}
\mathsf{S}^+_i[F]:=\sup_Xe^{\kappa_Fd(X)}\,\sup\bigl|\hat F^{\mathsf{B}}-\hat F^{\mathsf{A}}\bigr|,
\end{equation*}
where $X$ runs over the localisation domains of level $i$, $d(X)$ is the tree length of $X$, $\hat F$
is the value of $F$ under the slice map, and the inner supremum is over the slice data of the leaf
hull $X^+$, the hull on which the leaf attached to $X$ is read; every domain other than the whole
torus is read on its own shell, domains that wrap around the torus included.
$\mathsf{S}^{+,(j)}_\ell[\cdot]$ is the zone-$j$-resolved form of this
seminorm, defined together with the zone-resolved bound on the boundary terms; the zone-resolved
boundary loads below are taken from that definition:
\begin{equation*}
\begin{aligned}
\mathsf{y}_j&:=\sum_{\ell\ge j}\Bigl[\mathsf{w}_{\ell,j}\mathsf{S}^{+,(j)}_\ell[\mathbf{B}]
+\mathsf{S}^{+,(j)}_\ell[\mathbf{B}'^{res}]\Bigr],\\
\mathsf{y}'_j&:=\sum_{\ell\le j}w^{\mathbf{B}}_{j,\ell}\,\mathsf{S}^+_\ell[\mathbf{B}'^{own}],
\end{aligned}
\end{equation*}
where $\mathsf{w}_{\ell,j}$ are the crossing weights (not the stop depth $\mathsf{w}$) and
$\mathbf{B}'^{res}$, $\mathbf{B}'^{own}$ are the resolved and the own part of the boundary terms
$\mathbf{B}'$ born in the operation $\mathbf{R}$, as defined with the zone-resolved bound on the
boundary terms. Let $\hat\vartheta_i$ be the classical index rate of level $i$, and put
$\hat\vartheta^\natural_i:=\hat\vartheta_i^{1/2}$; then
$\hat\vartheta^\natural_i\le C_\vartheta\vartheta_{dl}^i$ with
$C_\vartheta:=\sqrt2\,(16C_WC_*c_3)^{1/4}$, where $C_W$, $C_*$, $c_3$ are the constants of the bound
on the classical index rates. With $E_0$ and $B_0$ the constants of the list bounds
and of the boundary-term bounds, $\kappa_0$ the exponent of the $R$-weights, $C^\natural_{cl}$ the
classical constant, $c_2$ the exponent of the logarithmic factor in the classical term and
$\vartheta_2$ as below, the loads of zone $j$ are
\begin{equation}\label{9.1:eq-centre-normalisation-d}
\begin{aligned}
{[\![\delta]\!]}^u_j&:=\sum_{i\le j}\varsigma_{1+j-i}\mathsf{S}^+_i[E]
+\sum_{i\le j}w_{j,i}\mathsf{S}^+_i[R]+\bigl[\mathsf{y}_j+\mathsf{y}'_j\bigr]\\
&\qquad+\mathfrak{g}_j\Bigl[\mathsf{z}^\sharp_j+\sum_{i\le j}L^{-(j-i)/2}\mathsf{b}_i\Bigr],\\
{[\![\delta]\!]}^c_j&:=\sum_{i\le j}\bigl(\varsigma_{1+j-i}E_0+w_{j,i}g_i^{\kappa_0}
+w^{\mathbf{B}}_{j,i}B_0\bigr)\hat\vartheta^\natural_i\\
&\qquad+E_0\min\bigl\{1,C^\natural_{cl}(\log x_j)^{c_2}\vartheta_2^{j/4}\bigr\},\\
{[\![\delta]\!]}^{\natural+}_j&:=[\![\delta]\!]^u_j+[\![\delta]\!]^c_j .
\end{aligned}
\end{equation}
The part of $[\![\delta]\!]^u_j$ formed by its three summands other than
$\mathsf{y}_j+\mathsf{y}'_j$ is denoted $[\![\delta]\!]^{(5)}_j$. Put
\begin{equation*}
E_1:=4(S_\varsigma+1)E_0+4S_{\mathbf{B}}B_0,\qquad
E_1^z:=\Bigl(\tfrac54+4C_T^{\mathbf{B}}\Bigr)E_1+\tfrac54B_0,
\end{equation*}
with $S_\varsigma$ and $S_{\mathbf{B}}$ the constants attached to the weights $\varsigma_\nu$ and
$w^{\mathbf{B}}_{j,i}$, and $C_T^{\mathbf{B}}$ the coefficient of the zone-resolved bound on the
boundary terms. The superscript $z$ in $E_1^z$ is a fixed part of the symbol: by its definition,
$E_1^z$ does not depend on the class $z$. By the a-priori bound on the loads, stated with the
zone-resolved boundary loads and under the conditions stated there,
\begin{equation*}
[\![\delta]\!]^{\natural+}_j\le E_1^z .
\end{equation*}
In what follows the radius of the dial domain, the total dial at which the slot bounds are used and
the further constants of the slot bounds that are formed from $E_1$ are all formed with $E_1^z$, and
the slot bounds are used with the
constant $C_L''':=(1+K_Y)C_L''$, where $C_L''$ is the constant of the slot bounds and $K_Y$ the leaf
constant. The
first-variation bound of the class $z$ reads
$\mathfrak{B}_z/\Lambda_z=\mathfrak{B}_z[\![\delta]\!]^{\natural+}_{n(z)}/E_1^z$, with
$\mathfrak{B}_z$ the modulus bound and $\Lambda_z$ as below. Since $E_1^z\ge E_1$, the factor
$(E_1^z)^{-1}$ that this bound delivers is majorised by $E_1^{-1}$; the bounds that follow are stated
with the factor $E_1^{-1}$.

Rates. Let $q$ be the kernel rate and $q_a:=e^{-c_a^{age}}$ the age rate; let $\vartheta_\chi$
be the index rate of the estimate on the characteristic functions $\chi$ of the labels. Put
$\vartheta_2:=L^{-\theta_2}$ and
$\vartheta_{dl}:=L^{-\theta_2/2048}$. The rates are subject to
\begin{equation}\label{9.1:eq-centre-normalisation-f}
\begin{gathered}
q_a\le q^2,\qquad q_a\le\min\{\vartheta_\chi^2,\vartheta_\chi/2\},\\
\rho\in\bigl[\max\{q^{1/2},L^{-\vartheta_\sharp},\vartheta_{dl},\vartheta_\chi\},1\bigr),
\qquad \rho>\vartheta_{dl}.
\end{gathered}
\end{equation}
The second condition on $q_a$ is needed by the random-walk bounds for the labels; it is a lower bound
on $c_a^{age}$, imposed after $q$, the index number $\theta_\chi$ of that estimate and $L$ are
fixed; no upper bound on $c_a^{age}$ is imposed. Put
\begin{equation*}
e_s:=\rho^{n-s}+\Pi_n,\qquad \mathsf{n}_s:=s+1+(1-\rho)^{-1},\qquad
\mathsf{l}_s:=(\log\theta_s)^{c_1},
\end{equation*}
with $c_1=p_0+q_\flat+1\ge\max\{c_2,q_\flat\}$. The envelope functions $\hat{\mathfrak{r}}_s$ and
$\Pi_n$ (the power-law floor), the constants $C_T$, $C_E$, and the constant $\mathsf{A}$ of the
bound on the coupling steps $\beta_r$ (a constant, not the run $\mathsf{A}$) are those of the
uniqueness argument.

Dials. Attach to every class $z'$ a complex number $\lambda_{z'}$, its dial. The dial domain is
\begin{equation}\label{9.1:eq-centre-normalisation-3}
\mathbb{D}^{\star}:=\bigcup_z\Bigl\{(\lambda_{z'})_{z'}:\ \exists\,\lambda\in[0,1]\ \text{with}\
\lambda_{z'}=\lambda\ (z'\ne z),\ |\lambda_z-\lambda|\,[\![\delta]\!]^{\natural+}_{n(z)}\le E_1^z
\Bigr\}.
\end{equation}
At a point of $\mathbb{D}^{\star}$ in the member of index $z$, with $\varepsilon:=\lambda_z-\lambda$,
the dialled array is, by the identity in \eqref{9.1:eq-centre-normalisation-c},
\begin{equation*}
x^{\mathsf{A}}+\sum_{z'}\lambda_{z'}W^{(z')}
=x^{\mathsf{A}}+\lambda\bigl(x^{\mathsf{B}}-x^{\mathsf{A}}\bigr)+\varepsilon W^{(z)},
\end{equation*}
and
\begin{equation}\label{9.1:eq-centre-normalisation-4}
|\varepsilon|\le\Lambda_z:=E_1^z\big/[\![\delta]\!]^{\natural+}_{n(z)}\ \ge1,
\end{equation}
the lower bound $\Lambda_z\ge1$ following from the a-priori bound on the loads. The slots of class $z$
sit at total dial at most $2E_1^z$, and the slot bounds are used at that dial. The zone-resolved
boundary leaves, the left-over chain terms $\mathbb{C}^{(i)}$ and the value units $V(X^\sharp)$
attached to the strip $X^\sharp$ of a domain $X$ are the exception to the rule of one class per slot:
they are slots belonging to several classes, and they are treated in the class lemma for the
zone-resolved boundary terms. The age-weighted load sum is
\begin{equation}\label{9.1:eq-centre-normalisation-5}
\Sigma^{q_a}_k:=\sum_{j\le k}(k-j+1)\,q_a^{(k-j-1)_+}\,[\![\delta]\!]^{\natural+}_j .
\end{equation}
The activity constant on the dial is
\begin{equation}\label{9.1:eq-centre-normalisation-6}
\mathsf{c}^{\star}:=C_{\mathfrak{K}}+C_{SM}+2+C_z+C^v_{\mathbf{B}},\qquad
C^v_{\mathbf{B}}:=2\bigl(K(d)+K_\circ+K^\circ\bigr)E_1^z/E_0,
\end{equation}
where $C_{\mathfrak{K}}$ is the constant of the class of list bounds, $C_{SM}$ and $C_z$ are the
constants of the activity bound ($C_z$ is a single constant; its subscript is not the class letter
$z$), and $K(d)$, $K_\circ$, $K^\circ$ are the resummation constants; the
constant $C^v_{\mathbf{B}}$ does not depend on $n$. Let $E_+^\sharp$ be the constant of the torus
bound of the correlation theorem, and $E_+^{\sharp\sharp}$ the same constant evaluated with activity
constant $\mathsf{c}^{\star}$, with the slot bounds at total dial $2E_1^z$ and constant $C_L'''$, and
with the doubled unit. Put
\begin{equation}\label{9.1:eq-centre-normalisation-7}
A^\sharp:=e^{E_+^{\sharp\sharp}N_w^4},\qquad A:=e^{E_+^\sharp N_w^4}\le A^\sharp ;
\end{equation}
$A$ and $A^\sharp$ are numbers, distinct from the effective action $A^{w}$ and from the Wilson
action. The constant $B$ and the product $\mathsf{Q}_w:=AB$ are
those of the normalisation of the uniqueness argument.

\medskip\noindent\emph{Floors on $n$.}
The comparison theorem for the two runs is stated for $n\ge n_1$, that is, for $n$ with
$\Upsilon_l\le\frac14$ for all $l\ge n$ --- then the pair exists and $|\mathfrak{n}_n'|\le\frac32$,
by the lemma on the matching map --- and for $k=n-\mathsf{w}\ge1$. The uniqueness argument uses it
under the floor $n\ge n_0(N_w)$, where
\begin{equation}\label{9.1:eq-centre-normalisation-e}
\begin{aligned}
n_0(N_w):=\max\Bigl\{&2\mathsf{w}+2,\ n_2(g),\\
&\text{the least }n\text{ with }\Upsilon_l\le\min\bigl\{\tfrac14,\upsilon_w/(2\varrho_2)\bigr\}
\ \text{for all }l\ge n\Bigr\},
\end{aligned}
\end{equation}
and $n_2(g)$ is the least $n$ such that, for every $l\ge n$, at $r=\lfloor l/2\rfloor$,
\begin{equation*}
l/2>\max\{\ell_r,\mathsf{N}(\theta_r)\}.
\end{equation*}
Here $\ell_r$ and $\mathsf{N}(\theta_r)$ are the numbers of levels that the depth-$r$ constructions
of the uniqueness argument, namely its index sets, bands and ladders, require beneath depth $r$, the
second read at the ramp point $\theta_r$.
Thus, beneath every depth $r\le n/2$, both runs have more than $\ell_r$ and more than
$\mathsf{N}(\theta_r)$ levels, so the depth-$r$ index sets, bands and ladders of the two runs are the
same sets. These floors on $n$ depend on $g$ and $N_w$ only, and no constant used in what follows
depends on $n$, $k$, $K$ or $m_0$.
\end{definition}

\begin{definition}[Zones and the stored boundary terms of the two steps]
\label{9.1:not-zones-boundary-terms}
The following notation is used in the zone-resolved estimates for the boundary terms in the
comparison of the two runs.

\smallskip\noindent
(a) \emph{The two runs and their levels.}
The runs $\mathsf{A}$ and $\mathsf{B}$ are those of
Notation~\ref{9.1:not-centre-normalisation}. Here the length of $\mathsf{A}$ is written
$\mathfrak{n}$, and the length of $\mathsf{B}$ is $\mathfrak{n}+1$. This leaves the letter $n$
free for zones. Both runs are carried on one and the same prescribed scaffold $\Theta_X$, whose
subscript is the inverse coupling $X=g^{-2}$. In $\Theta_X$ and in $\theta_s=X+c_0s$ below, the
letter $X$ is this inverse coupling; in every other place in this notation $X$ and $X'$ denote
localisation domains. Levels $\ell$, $j$ are indexed by the index of $\mathsf{A}$, and
satisfy $\ell,j\le k'$, where $k'$ is the level at which a term is read (the \emph{reading
level}). The depth of level $\ell$ is $s(\ell):=\mathfrak{n}-\ell$. At depth $s$ we write
$\theta_s=X+c_0s$ for the inverse coupling of the scaffold. We put
\begin{equation}\label{9.1:eq-zones-boundary-terms-1}
\mathsf{l}_s:=(\log\theta_s)^{c_1},\qquad e_s:=\rho^{\mathfrak{n}-s}+\Pi_{\mathfrak{n}}.
\end{equation}
Here $\rho$ is the contraction letter of clause~(v), taken in its window, and $\Pi_{\mathfrak{n}}$
is the power-law floor.

\smallskip\noindent
(b) \emph{Regions, zones, large-field regions.}
The regions of a label form the decreasing chain of \cite{B-CMP119}:
\begin{equation}\label{9.1:eq-zones-boundary-terms-2}
\Omega_1\supset\Lambda_1\supset\Omega_2\supset\Lambda_2\supset\cdots .
\end{equation}
For $n<k'$ the \emph{zone} $n$ is $\Omega_n\setminus\Omega_{n+1}$. The top zone $k'$ is
$\Omega_{k'}$. The large-field region at level $j$ is
\begin{equation}\label{9.1:eq-zones-boundary-terms-3}
Z_j:=\Lambda_j^{c}.
\end{equation}
This is the large-field region $\Lambda^c_k$ that each term of the expansion carries in
\cite{B-CMP122a} and \cite{B-CMP119}. The components of $Z_j$ are the classes of the
equivalence relation generated by the relation ``intersect or touch'' among its cubes.

\smallskip\noindent
(c) \emph{The stored boundary terms of the $\mathbf{T}$-step.}
In \cite[(2.41)]{B-CMP119} the boundary term $B^{(\ell)}(X)$ is indexed by domains $X$ of the
class $\mathbf{D}_\ell$ of localisation domains of level $\ell$ of \cite{B-CMP119}
with $X\cap\Omega_\ell\neq\emptyset$ and $X\cap Z_\ell^{\sim}\neq\emptyset$.
Here $Z_\ell^{\sim}$ is the enlargement of $Z_\ell$ by its collar used there. The terms so indexed
include the terms $E^{(\ell)}_0(\Lambda_\ell,X,y)$ with $X\subset\Lambda_\ell$ and $y$ a point of
the rim $\Lambda_\ell\setminus\Lambda_\ell^0$, where $\Lambda^0_\ell\subset\Lambda_\ell$ is the
inner region of \cite{B-CMP119}. Such terms only touch the collar of $Z_\ell$.
In \cite{B-CMP119} the rim point is written $z$; we write $y$, the letter $z$ being reserved here
for classes (see~(e)).

In the comparison of the two runs one works with the \emph{comparison list}: the list of terms
of the two runs that are compared term by term. By the lemma that fixes the comparison list,
after each step the terms with $X\subset\Lambda_j$ are terms of the comparison list. By the
\emph{$\mathbf{T}$-site} we mean the terms produced by the $\mathbf{T}$-step with
$X\subset\Lambda_\ell$; the \emph{bound at the $\mathbf{T}$-site} is the zone-by-zone bound for
these terms, whose constants are listed in~(h). The terms $E^{(\ell)}_0(\Lambda_\ell,X,y)$ above
have $X\subset\Lambda_\ell$ and are therefore terms of the comparison list; they are assigned to
the zone of their rim point $y$, with the room weight $\varsigma$, and they are bounded by the
bound at the $\mathbf{T}$-site, whose constants do not depend on $y$.
The
\emph{stored boundary terms} $\mathbf{B}^{(\ell)}(X)$ are therefore, by definition, the remaining
terms: those with
\begin{equation}\label{9.1:eq-zones-boundary-terms-4}
X\in\mathbf{D}_\ell,\qquad X\cap\Omega_\ell\neq\emptyset,\qquad X\not\subset\Lambda_\ell .
\end{equation}
In particular, a term $E^{(\ell)}_0(\Lambda_\ell,X,y)$ with $y$ in the rim whose only
neighbouring component of $Z_\ell$ is healed at the $\mathbf{R}$-step of level $\ell$ itself has
$X\subset\Lambda_\ell$; it is a term of the comparison list, not a stored boundary term, and it
contributes nothing to the estimates for stored boundary terms.
The stored boundary terms are produced by the $\mathbf{T}$-step $\ell-1\to\ell$ in two ways:
\begin{itemize}
\item[(i)] as the terms of \cite[(3.47)]{B-CMP119} that are not contained in $\Lambda_\ell$;
\item[(ii)] as the one expectation value of that step which goes wholly to $B^{(\ell)}$. As
described in \cite{B-CMP119} (with $k+1=\ell$), the first of the two expectation values of that
step, the integral over $t$ of the expectation value of the boundary terms, contributes to
$B^{(k+1)}$ only; the second contributes to both the terms $R^{(k+1)}$ of \cite{B-CMP119} and
$B^{(k+1)}$.
Item~(ii) is the first of these.
\end{itemize}
They obey \cite[(2.42)]{B-CMP119}:
\begin{equation}\label{9.1:eq-zones-boundary-terms-5}
|\mathbf{B}^{(\ell)}(X)|\le B_0\,e^{-\kappa d_\ell(X)} .
\end{equation}
In the comparison of the two runs these terms are stored in the form in which the correlation
theorem stores them, as sums $\sum_{\mathcal{S}}\mathbf{B}^{(\ell)}(X,\mathcal{S})$ over the index
$\mathcal{S}$ of that form, with the background datum $A^{\flat}$ at which that form evaluates
them, at the fresh decay rate
$\kappa_{out}\ge 2c^{\sharp}\kappa$ supplied for the output of the $\mathbf{T}$-step in the
relative small-field format. The bound in which they are used as boundary inputs requires only a
rate $\kappa_T\ge\kappa$; we write $\kappa_{\mathbf{B}}:=\kappa$ for the rate at which they enter
below.

\smallskip\noindent
(d) \emph{The stored boundary terms of the $\mathbf{R}$-step.}
The symbol $\mathbf{B}'^{(\ell)}(X')$ denotes the first group of the terms of
\cite[(1.98)]{B-CMP122b}; we write these terms $\mathbf{R}'^{(\ell)}(X')$, and
$\mathbf{R}'^{(\ell)}(X,(\mathbf{U},\mathbf{J}))$ when their arguments $\mathbf{U}$, $\mathbf{J}$
are displayed. These are the terms whose localisation domains intersect the large-field
region that survives the step. Each term of this group satisfies
\cite[(1.99), p.~390]{B-CMP122b}:
\begin{equation}\label{9.1:eq-zones-boundary-terms-6}
|\mathbf{R}'^{(\ell)}(X,(\mathbf{U},\mathbf{J}))|
\le O(1)\,c_1\exp\Bigl(-\bigl(1+\tfrac12\beta\bigr)\kappa\,
d_{\ell,\cup_iY_i}(X)\Bigr).
\end{equation}
In \cite[(1.98)--(1.99)]{B-CMP122b} the domain of these terms is written $X'$ in their definition
and in the case distinction below, and $X$ in the bound; we keep both letters as printed there.
The quantities in this bound are as follows.
\begin{itemize}
\item $O(1)$ is the absolute constant of \cite[(1.99)]{B-CMP122b}.
\item $\beta$ is the parameter in the exponent of \cite[(1.99)]{B-CMP122b}; in the comparison of
the two runs it is taken at the step letter $\beta_s$ of the relative small-field machine, whose
value is $\beta_s=\frac14$.
\item The $Y_i$ are the regions of the $\mathbf{R}$-step entering \cite[(1.98)]{B-CMP122b}.
\item $d_{\ell,\cup_iY_i}(X)$ is the relative size of \cite[(1.67)]{B-CMP122b}: the length of a
shortest tree in $X$ that is required to meet only the cubes of $X\setminus\bigcup_iY_i$.
\item $c_1$ is Ba{\l}aban's constant of \cite[(1.99)]{B-CMP122b}:
\begin{equation}\label{9.1:eq-zones-boundary-terms-7}
c_1=
\begin{cases}
e^{-p_0(g_\ell)} & \text{if $X'$ does not intersect $\bigcup_iY_i$},\\
\alpha^{1/3} & \text{in the remaining cases,}
\end{cases}
\end{equation}
with $p_0(\cdot)$ the degree function.
\end{itemize}
This $c_1$ is Ba{\l}aban's constant and is not the exponent $c_1$ in
\eqref{9.1:eq-zones-boundary-terms-1}. In \cite[\S1]{B-CMP122b} the relative size is replaced by
$d_\ell(X)$ (Ba{\l}aban's $d_k(X)$) only for the second group of terms of
\cite[(1.98)]{B-CMP122b}; for the first group the bound \eqref{9.1:eq-zones-boundary-terms-6}
stays in the relative length.

A term $\mathbf{B}'^{(\ell)}$ is \emph{plain} if its domain meets no live island of the
large-field region. Every statement below
in which the size of a term $\mathbf{B}'^{(\ell)}$ is measured in the plain length $d_\ell$ is
asserted for plain terms only; for these $d_{\ell,\cup_iY_i}=d_\ell$, and in the comparison of
the two runs the output length $d_{out}$ of the strip format (the format in which the output of
the relative small-field step is written in that comparison) equals $d_\ell$. For the remaining
terms, those whose domain contains a live strip of a second-class component of the large-field
region (in the division of its components into a first and a second class used in the comparison
of the two runs), which is exactly the case
$c_1=\alpha^{1/3}$ of \eqref{9.1:eq-zones-boundary-terms-7}, the same displays are used in the
relative length $d_{\ell,\cup_iY_i}$ and the relative sup-norm $\mathsf{S}^{\circ}$, under the
relative expansion statement.

In the relative small-field machine, which carries the two runs in their comparison,
Ba{\l}aban's radius $\alpha$ is replaced by $\alpha^{\sharp}_\ell$, and
\eqref{9.1:eq-zones-boundary-terms-6} holds with $c_1$ given by
\eqref{9.1:eq-zones-boundary-terms-7} at this replacement. By the corollary bounding this radius
in the relative small-field format, it satisfies
\begin{equation}\label{9.1:eq-zones-boundary-terms-8}
(\alpha^{\sharp}_\ell)^{1/3}\le\min\bigl\{e^{-(1+\beta_s)\kappa 2^d},\,g_\ell^{\kappa_0}\bigr\},
\end{equation}
where $\kappa_0$ is the exponent of that corollary. We put
\begin{equation}\label{9.1:eq-zones-boundary-terms-9}
c_1^{\flat}(g_\ell):=O(1)\,\max\bigl\{e^{-p_0(g_\ell)},\,(\alpha^{\sharp}_\ell)^{1/3}\bigr\},
\end{equation}
with $O(1)$ the constant of \eqref{9.1:eq-zones-boundary-terms-6}.

\smallskip\noindent
(e) \emph{Two-cutoff letters.}
For a family $F$ of terms carried by both runs, with decay rate $\kappa_F$, the two-cutoff
seminorm $\mathsf{S}_j^{+}[F]$ is defined as follows. It is the supremum over domains $X$ of
level $j$ of $e^{\kappa_F d_j(X)}$ times the supremum, over the data $\mathsf{v}$ on the slice of
the hull $X^{+}$ of $X$ and over the local label (the label data on $X^{+}$), of
$|\hat F^{\mathsf{B}}-\hat F^{\mathsf{A}}|$. Here
$\hat F^{\mathsf{c}}$ is the slice map of run $\mathsf{c}$.

The terms of the comparison list are of three kinds, which we call units, pieces and brackets;
they are fixed with the comparison list. A unit or a piece may carry a point, as the terms
$E^{(\ell)}_0(\Lambda_\ell,X,y)$ of~(c) carry the rim point $y$.
A term read at level $k'$ carries a class $z=(n,k')$, with $n(z):=n$ its anchor zone at the
reading level. A unit or a piece that carries no point is anchored in the lowest zone it meets,
and a bracket in the lower of its zones. The two-cutoff differences are of two types: the
classical type, denoted $c$, and the other type, denoted $u$.
The loads of zone $n$ are $[\![\delta]\!]^u_n$ and $[\![\delta]\!]^c_n$, and
\begin{equation}\label{9.1:eq-zones-boundary-terms-10}
[\![\delta]\!]^{\natural+}_n:=[\![\delta]\!]^u_n+[\![\delta]\!]^c_n,\qquad
w^{\mathbf{B}}_{n,j}:=L^{-(n-j)/2}.
\end{equation}
With $S_{\mathbf{B}}$ and $S_\varsigma$ the summation constants of the weights
$w^{\mathbf{B}}$ and of the room weight $\varsigma$, and $E_0$ the constant of the bound at the
$\mathbf{T}$-site (not to be confused with the terms $E^{(\ell)}_0(\cdot)$ of~(c)), we put
\begin{equation}\label{9.1:eq-zones-boundary-terms-11}
E_1:=4(S_\varsigma+1)E_0+4S_{\mathbf{B}}B_0,\qquad
\Lambda_z:=E_1/[\![\delta]\!]^{\natural+}_{n(z)} .
\end{equation}
Here $\Lambda_z$ is the radius of class $z$ on the dial domain $\mathbb{D}^{\star}$; it is indexed
by a class and is not one of the regions of~(b). The age rate $q_a$ is the rate of geometric
forgetting of old regions.

\smallskip\noindent
(f) \emph{Aligned labels.}
A label is \emph{aligned} if its first index exceeds $n_c$. Then every zone present in it has
index at least $n_c$. The index $n_c$ is determined by the constant $C^{\natural}_{cl}$ and the
rate $\vartheta_2^{1/4}$ so that the classical index is small from $n_c$ on, that is,
$C^{\natural}_{cl}(\log x_n)^{c_2}\vartheta_2^{n/4}\le1$ for every $n\ge n_c$. Since every zone
present in an aligned label has index at least $n_c$, this holds in the form
\begin{equation}\label{9.1:eq-zones-boundary-terms-a}
C^{\natural}_{cl}\,(\log x_n)^{c_2}\,\vartheta_2^{n/4}\le1
\qquad\text{for every zone $n$ present in an aligned label.}
\end{equation}
Here $\vartheta_2:=L^{-\theta_2}$ with $\theta_2=(5\theta_A-3)/8$. The index letter $\theta_A$
ranges in the window $(3/5,1)$. The numbers $\theta_A$ and $\theta_2$ are index letters fixed
before $L$ in the order of choice; they are unrelated to the inverse couplings $\theta_s$ of~(a),
and $\theta_2$ is not the value of $\theta_s$ at $s=2$.

\smallskip\noindent
(g) \emph{The rate at the $\mathbf{T}$-site.}
The letter $\vartheta_T$ denotes the rate, per zone, of the bounds at the $\mathbf{T}$-site:
they hold with the factor $\vartheta_T^{\,n}$ in zone $n$, that is,
in zone $n$ they are at most a constant times $\vartheta_T^{\,n}$. This rate is the classical
rate $\vartheta_w$ of the statement that delivers the bounds at the $\mathbf{T}$-site, taken at
the index letter $\theta_A$. We put
\begin{equation}\label{9.1:eq-zones-boundary-terms-b}
\begin{gathered}
\vartheta_T:=\vartheta_w=L^{-\theta_A/8},\qquad
\vartheta_\times:=\vartheta_w\,\vartheta_2^{-1/4}=L^{-(3-\theta_A)/32},\\
e_T:=\frac{3-\theta_A}{5\theta_A-3}.
\end{gathered}
\end{equation}
No inequality between $\vartheta_T$ and $\vartheta_2$ is assumed. Where the zone-resolved
estimates compare the rate $\vartheta_T^{\,n}$ of the bounds at the $\mathbf{T}$-site with the
classical rate $\vartheta_2^{n/4}$, they do so through $\vartheta_\times$ and the exponent $e_T$.

\smallskip\noindent
(h) \emph{Constants of the $\mathbf{T}$-site.}
For a coupling $g'$ put $\theta(g'):=C_\theta(\log g'^{-2})^{p_0-q}$, where $C_\theta$ is the
constant fixed with the correlation theorem, and write $\bar\theta:=\theta(\gamma)$. The letter
$\theta_s$ of~(a) is the scaffold's inverse coupling and is not a value of $\theta(\cdot)$; for
the values of $\theta(\cdot)$ at the couplings $g_i$ we write $\theta(g_i)$ in full. The other
constants are:
\begin{itemize}
\item $\varkappa_i$ is the coefficient of level $i$ in the bound at the $\mathbf{T}$-site:
$\varkappa_i=K_{CE}C_L\theta(g_i)$, where $K_{CE}$ and $C_L$ are the two constant factors of this
coefficient in the one-lattice form of that bound; the coefficient is evaluated at the value
$C''_L$ of the constant $C_L$ and the value $\kappa_1^{\natural}$ of the parameter $\kappa_1$
fixed with that bound.
\item $\sigma^T_n$ is defined, for an index $n$ (a level, equivalently the zone of that index), by
\begin{equation}\label{9.1:eq-zones-boundary-terms-12}
\sigma^T_n:=C_TE_0\,(\log x_n)^{c_2}\,\vartheta_w^{\,n},
\end{equation}
where $C_T$ is the constant of the zone coefficient of the bound at the $\mathbf{T}$-site.
\item $\mathfrak{m}_G^T$ is the dial budget at the $\mathbf{T}$-site. The bounds there hold for
dial parameters whose modulus, multiplied by the per-zone size of those bounds (in zone $n$ a
constant times $\vartheta_T^{\,n}$), is at most $\mathfrak{m}_G^T$.
\item $E^{\sharp}_T$ is the constant of the boundary output of the summation lemma for these
bounds on that dial; it obeys that lemma's bound with the volume factor $\exp(O(1)\alpha_5|Z|)$,
where $\alpha_5$ is the small constant fixed with the summation lemma, replaced by
$e^{2d\alpha^{*}|Z|}$, where $\alpha^{*}$ is the constant that replaces $\alpha_5$ on the dial of
the $\mathbf{T}$-site.
\item $C_{cl}^T$ is the constant of the zone-$n$ operator bound at the $\mathbf{T}$-site, fixed
with that bound.
\end{itemize}

\smallskip\noindent
(i) \emph{Crossing length.}
A \emph{crossing} of zone $m$ is a domain of level-$(m+1)$ cubes that meets both $\Omega_m^{c}$
and $\Omega_{m+1}$. We write $\mathsf{d}_m$ for the tree length $d_{m+1}$ of a crossing
of zone $m$, that is, the length of a shortest tree joining its cubes. The constant
$c_\times=c_\times(d)$ bounds the contribution of the two end cubes of such a tree. The scaffold
radius at level $m$ is $\check R_m$, which is Ba{\l}aban's $R_m$.

\smallskip\noindent
With this notation the following hold.
\begin{itemize}
\item[(P1)] If $n<j$, zone $n$ lies in $Z_j$.
\item[(P2)] Every stored boundary term $\mathbf{B}^{(\ell)}(X)$ satisfies $X\cap Z_\ell\neq\emptyset$.
Hence some level-$\ell$ cube of $X$ lies in $Z_\ell$.
\item[(P3)] If the domain $X$ of \eqref{9.1:eq-zones-boundary-terms-6} meets no $Y_i$, then
$d_{\ell,\cup_iY_i}(X)=d_\ell(X)$. In the comparison of the two runs,
\begin{equation}\label{9.1:eq-zones-boundary-terms-13}
|\mathbf{R}'^{(\ell)}(X,(\mathbf{U},\mathbf{J}))|
\le c_1^{\flat}(g_\ell)\exp\Bigl(-\bigl(1+\tfrac12\beta\bigr)\kappa\,
d_{\ell,\cup_iY_i}(X)\Bigr).
\end{equation}
At $\beta=\beta_s$ the rate is $(1+\frac12\beta_s)\kappa=\frac98\kappa\ge\kappa_{\mathbf{B}}$.
\item[(P4)] For every $\theta_A\in(3/5,1)$:
\begin{itemize}
\item[$\cdot$] $\vartheta_\times<1$ and $e_T>1$, with $e_T=3$ at $\theta_A=\frac34$;
\item[$\cdot$] $\vartheta_\times=(\vartheta_2^{1/4})^{e_T}$;
\item[$\cdot$] $\vartheta_w\le\vartheta_2$ holds if and only if $\theta_A\le\frac34$.
\end{itemize}
\item[(P5)] Let $n$ be a zone present in an aligned label. Then, for every constant $C\ge0$,
\begin{equation}\label{9.1:eq-zones-boundary-terms-14}
\vartheta_\times^{\,n}\le\bigl(C^{\natural}_{cl}(\log x_n)^{c_2}\bigr)^{-e_T},\qquad
C(\log x_n)^{c_2}\vartheta_T^{\,n}\le\frac{C}{C^{\natural}_{cl}}\,
\vartheta_\times^{\,n},
\end{equation}
and in particular
\begin{equation}\label{9.1:eq-zones-boundary-terms-15}
\sigma^T_n\le\frac{C_TE_0}{C^{\natural}_{cl}}\,\vartheta_\times^{\,n}.
\end{equation}
If $n_0$ is the lowest zone present in an aligned label, then
\begin{equation}\label{9.1:eq-zones-boundary-terms-16}
\sum_{n\ge n_0}\vartheta_\times^{\,n}
\le(1-\vartheta_\times)^{-1}\bigl(C^{\natural}_{cl}(\log x_{n_0})^{c_2}\bigr)^{-e_T}.
\end{equation}
Since $n_0\ge n_c$, moreover
\begin{equation}\label{9.1:eq-zones-boundary-terms-17}
\sum_{n\ge n_0}\vartheta_\times^{\,n}
\le(1-\vartheta_\times)^{-1}\vartheta_\times^{\,n_c}
\le(1-\vartheta_\times)^{-1}\bigl(C^{\natural}_{cl}(\log x_{n_c})^{c_2}\bigr)^{-e_T}.
\end{equation}
These inequalities are the form in which \eqref{9.1:eq-zones-boundary-terms-a} is used in the
zone-resolved estimates, for the operator sizes (with $C=C^T_{cl}$, giving the factor
$(C^T_{cl}/C^{\natural}_{cl})\vartheta_\times^{\,n}$) and the coefficients at the
$\mathbf{T}$-site, and for their sums over zones.
\item[(P6)] For every zone $m<k'$, every crossing of zone $m$ has $d_{m+1}$-length
$\mathsf{d}_m\ge 2\check R_m/L-c_\times$.
\end{itemize}
\end{definition}

\begin{proof}
\emph{(P1).} Let $n<j$, so that $n+1\le j$. By the chain \eqref{9.1:eq-zones-boundary-terms-2},
\[
\Lambda_j\subset\Omega_j\subset\Omega_{n+1}.
\]
Zone $n$ is $\Omega_n\setminus\Omega_{n+1}$, so it is contained in $\Omega_{n+1}^c$. Hence
\[
\Omega_n\setminus\Omega_{n+1}\subset\Omega_{n+1}^c\subset\Lambda_j^c=Z_j.
\]

\emph{(P2).} By \eqref{9.1:eq-zones-boundary-terms-4} a stored boundary term has
$X\not\subset\Lambda_\ell$. So $X$ contains a point outside $\Lambda_\ell$, that is, a point of
$Z_\ell$. In Ba{\l}aban's construction \cite{B-CMP119} the domains in $\mathbf{D}_\ell$ and the
regions of a label are unions of level-$\ell$ cubes. Therefore a level-$\ell$ cube of $X$ lies in
$X\cap Z_\ell$.

\emph{(P3).} Suppose $X$ meets no $Y_i$. Then $X\setminus\bigcup_iY_i=X$. The requirement in
(d), that the tree meet only the cubes of $X\setminus\bigcup_iY_i$, then becomes the requirement
that it meet the cubes of $X$. This is the definition of $d_\ell(X)$, so
$d_{\ell,\cup_iY_i}(X)=d_\ell(X)$.

For \eqref{9.1:eq-zones-boundary-terms-13}, we replace $\alpha$ by $\alpha^{\sharp}_\ell$ in
\eqref{9.1:eq-zones-boundary-terms-7}. Then $c_1$ is one of $e^{-p_0(g_\ell)}$ and
$(\alpha^{\sharp}_\ell)^{1/3}$, hence at most their maximum. Multiplying by the constant $O(1)$ of
\eqref{9.1:eq-zones-boundary-terms-6} gives $O(1)\,c_1\le c_1^{\flat}(g_\ell)$ by
\eqref{9.1:eq-zones-boundary-terms-9}. Inserting this in \eqref{9.1:eq-zones-boundary-terms-6}
gives \eqref{9.1:eq-zones-boundary-terms-13}.

For the last assertion, since $\beta_s=\frac14$ we have
$(1+\frac12\beta_s)\kappa=(1+\frac18)\kappa=\frac98\kappa$, and
$\frac98\kappa\ge\kappa=\kappa_{\mathbf{B}}$.

\emph{(P4).} Write $\vartheta_2^{-1/4}=L^{(5\theta_A-3)/32}$. Then
\begin{align*}
\vartheta_w\,\vartheta_2^{-1/4}
&=L^{-4\theta_A/32+(5\theta_A-3)/32}\\
&=L^{(\theta_A-3)/32}=L^{-(3-\theta_A)/32}.
\end{align*}
This is the value in \eqref{9.1:eq-zones-boundary-terms-b}. Since $\theta_A<3$ and $L>1$, it is
smaller than $1$.

On the window we have $5\theta_A-3>0$. The inequality $3-\theta_A>5\theta_A-3$ is equivalent to
$6>6\theta_A$, that is, to $\theta_A<1$. Hence $e_T>1$. At $\theta_A=\frac34$,
\[
e_T=\frac{9/4}{3/4}=3 .
\]
Next, $\vartheta_2^{1/4}=L^{-(5\theta_A-3)/32}$. Therefore
\[
(\vartheta_2^{1/4})^{e_T}=L^{-(5\theta_A-3)e_T/32}=L^{-(3-\theta_A)/32}=\vartheta_\times .
\]
Finally, $\vartheta_w\le\vartheta_2$ means $L^{-\theta_A/8}\le L^{-(5\theta_A-3)/8}$. Since
$L>1$, this is equivalent to $\theta_A\ge5\theta_A-3$, that is, to $\theta_A\le\frac34$.

\emph{(P5).} Let $n$ be a zone present in an aligned label. By
\eqref{9.1:eq-zones-boundary-terms-a},
\[
0<\vartheta_2^{n/4}\le\bigl(C^{\natural}_{cl}(\log x_n)^{c_2}\bigr)^{-1}.
\]
By (P4), $\vartheta_\times^{\,n}=(\vartheta_2^{n/4})^{e_T}$ with $e_T>0$. Raising the last
inequality to the power $e_T$ gives the first inequality of
\eqref{9.1:eq-zones-boundary-terms-14}.

By \eqref{9.1:eq-zones-boundary-terms-b}, $\vartheta_T=\vartheta_\times\vartheta_2^{1/4}$, hence
$\vartheta_T^{\,n}=\vartheta_\times^{\,n}\vartheta_2^{n/4}$. Let $C\ge0$. Multiplying by
$C(\log x_n)^{c_2}$ and using \eqref{9.1:eq-zones-boundary-terms-a} again gives
\[
C(\log x_n)^{c_2}\vartheta_T^{\,n}
\le C(\log x_n)^{c_2}\vartheta_\times^{\,n}
\bigl(C^{\natural}_{cl}(\log x_n)^{c_2}\bigr)^{-1}
=\frac{C}{C^{\natural}_{cl}}\,\vartheta_\times^{\,n}.
\]
This is the second inequality of \eqref{9.1:eq-zones-boundary-terms-14}. With $C=C_TE_0$, and
with $\vartheta_w=\vartheta_T$ in \eqref{9.1:eq-zones-boundary-terms-12}, its left-hand side is
$\sigma^T_n$, which gives \eqref{9.1:eq-zones-boundary-terms-15}.

For \eqref{9.1:eq-zones-boundary-terms-16}, we have $0<\vartheta_\times<1$ by (P4). The
geometric series therefore gives
\[
\sum_{n\ge n_0}\vartheta_\times^{\,n}=(1-\vartheta_\times)^{-1}\vartheta_\times^{\,n_0}.
\]
Applying the first inequality of \eqref{9.1:eq-zones-boundary-terms-14} at the zone $n_0$ bounds
the right-hand side as claimed.

For \eqref{9.1:eq-zones-boundary-terms-17}, the zone $n_0$ is present in an aligned label, so
$n_0\ge n_c$ by~(f). Since $0<\vartheta_\times<1$, this gives
$\vartheta_\times^{\,n_0}\le\vartheta_\times^{\,n_c}$, and with the geometric series above the
first inequality of \eqref{9.1:eq-zones-boundary-terms-17} follows. For the second, the defining
property of $n_c$ in~(f), taken at $n=n_c$, gives
\[
0<\vartheta_2^{n_c/4}\le\bigl(C^{\natural}_{cl}(\log x_{n_c})^{c_2}\bigr)^{-1}.
\]
By (P4), $\vartheta_\times^{\,n_c}=(\vartheta_2^{n_c/4})^{e_T}$ with $e_T>0$; raising the last
inequality to the power $e_T$ gives
$\vartheta_\times^{\,n_c}\le\bigl(C^{\natural}_{cl}(\log x_{n_c})^{c_2}\bigr)^{-e_T}$, and
multiplying by $(1-\vartheta_\times)^{-1}>0$ gives the second inequality of
\eqref{9.1:eq-zones-boundary-terms-17}.

\emph{(P6).} By the chain \eqref{9.1:eq-zones-boundary-terms-2}, $\Omega_{m+1}\subset\Lambda_m$.
Hence zone $m$ contains $\Omega_m\setminus\Lambda_m$. A crossing of zone $m$ contains a cube
meeting $\Omega_m^c$ and a cube meeting $\Omega_{m+1}\subset\Lambda_m$. The union of a shortest
tree joining its cubes with the two end cubes of that tree is connected and joins a point of
$\Omega_m^{c}$ to a point of $\Lambda_m$; it therefore traverses the shell
$\Omega_m\setminus\Lambda_m$ from the boundary of $\Omega_m$ to the boundary of $\Lambda_m$.

By the construction of the regions in \cite{B-CMP119}, for $\Omega_m\supset\Lambda_m$ the
distance between their boundaries is at least $2M\check R_m$ in level-$m$ units. Passing to the
next step, lengths are rescaled by $L^{-1}$ \cite{B-CMP122a}. So in level-$(m+1)$ units this
distance is at least $2M\check R_m/L$. The tree length $d_{m+1}$ of \cite{B-CMP119} is measured
in cubes of side $M$ of level $m+1$, the cubes of which Ba{\l}aban's localisation domains are
unions; in this length the distance is at least $2\check R_m/L$.

Hence the $d_{m+1}$-length of the tree, increased by the contribution of its two end cubes, is at
least this distance. That contribution is at most $c_\times=c_\times(d)$. Hence
$\mathsf{d}_m\ge2\check R_m/L-c_\times$.
\end{proof}

\begin{lemma}[Zone split of the two-run differences at the fluctuation step]\label{9.1:lem-zone-split}
Let $\mathsf{A}$ and $\mathsf{B}$ be the two runs, of lengths $\mathfrak{n}$ and $\mathfrak{n}+1$, on the one
scaffold $\Theta_X$ of Notation~\ref{9.1:not-zones-boundary-terms}, with zones as fixed there. For a
set $Y$ of cubes, $j_{\min}(Y)$ denotes the lowest zone met by $Y$. Fix an aligned label (a label
in $H_{al}$) and its $\mathbf{T}$-step $i\to i+1$. For $\mathsf{m}\in\{\mathsf{A},\mathsf{B}\}$ the
objects of run $\mathsf{m}$ at the $\mathbf{T}$-site are the following.
\begin{itemize}
\item[(z1)] The fluctuation covariance $\mathsf{D}=C^{(i)}(Z_0,\sigma)^{-1}$, the walk operator
$\Gamma=\Gamma_i(Z_0,\sigma)$ and the first form $\mathsf{F}=\Gamma^{\top}\mathsf{D}^{-1}\Gamma$, as
kernels over the bond sets of the $\mathbf{T}$-site.
\item[(z2)] The potentials $\Psi$, $\Phi$, the Jacobians and the factor $\sigma(g_iB)$, which are
localised families over polymers $Y$.
\item[(z3)] The normaliser $\mathcal{L}^{(i+1)}$ together with the share of $\mathbf{B}$ in it,
localised as a sum $\sum_{X''}\mathcal{L}_R(X'')$ over domains $X''$.
\end{itemize}
For each of these objects, written $\mathsf{z}^{\mathsf{m}}$, we write
$\mathsf{z}^{\mathsf{m}}=\sum_\omega\mathsf{z}^{\mathsf{m}}_\omega$ as the sum of its pieces $\omega$
(described in (h1) below), with support $\operatorname{supp}\omega$. For a piece $\omega$ let
$n(\omega)$ be the lowest zone met by $\operatorname{supp}\omega$, let
$\delta\mathsf{z}_\omega:=\mathsf{z}^{\mathsf{B}}_\omega-\mathsf{z}^{\mathsf{A}}_\omega$, and put
\begin{equation}\label{9.1:eq-zone-split-1}
\delta\mathsf{z}^{(n)}:=\sum_{\omega:\ n(\omega)=n}\delta\mathsf{z}_\omega ,
\end{equation}
filed in the class of zone $n$ at level $i$. For a domain $X$, $\delta\mathsf{z}^{(n)}{\upharpoonright}X$
denotes the sub-sum of \eqref{9.1:eq-zone-split-1} over the pieces with
$\operatorname{supp}\omega\subset X$. Put $\vartheta_T:=\vartheta_w$ and
\begin{equation}\label{9.1:eq-zone-split-2}
\mathfrak{d}^{op}_n:=C_{cl}^T(\log x_n)^{c_2}\,\vartheta_T^{\,n},
\end{equation}
where the constant $C_{cl}^T$ collects $C_{cl}$, $C_K$, $C_Fe^{16\kappa_1}$ and the one-run
resummation constants of the three objects (z1)--(z3); in particular $C_{cl}^T\ge C_Fe^{16\kappa_1}$
and $C_{cl}^T\ge2C_{res}C_Fe^{16\kappa_1}$, with $C_{res}$ the one-run resummation constant of
(h5$'$). Let $n_c$ be the index fixed with the zone
construction; it enters below only through hypothesis (h11) and the limit in (iv).

Assume the following.
\begin{itemize}
\item[(h1)] (The walk-decoupling lemma.) Each kernel in (z1) of run $\mathsf{m}$ is the sum of its
$\mathsf{s}$-decoupled random-walk pieces $\omega$, each with connected support
$\operatorname{supp}\omega$, produced by the cluster expansion of \cite[(2.5)--(2.13)]{B-CMP116},
and the sum over pieces converges absolutely;
the objects (z2) are localised families over polymers. Both runs are taken on the one scaffold and
on one interpolation scheme, so that the index set of pieces is the same for $\mathsf{A}$ and
$\mathsf{B}$.
\item[(h2)] (The locality lemma for outputs.) An output with domain $X$ reads only the localised
terms $V(Y)$ of the effective action with $Y\subset X$.
\item[(h3)] (The filing rule for unpointed units.) An unpointed unit is filed in the lowest zone it
meets; the dial $\lambda_{j,k}$ of the dial proposition is filed by the same rule.
\item[(h4)] (The operator lemma for the two runs at the $\mathbf{T}$-site, and the
natural-interpolation lemma.) The operators of (z1) are those of the list of clause (b) of the
family lemma. For an operator $\mathbb{O}$ of that list and a bond pair $(b,b')$ with its set
$\mathsf{K}$ of cubes, which contains the fixed core $\mathsf{K}_0$ of the interpolation scheme (on
which $\mathsf{s}:=1$), and for a finite set $\mathsf{W}$ of cubes of the label's partition, put
\begin{equation}\label{9.1:eq-zone-split-3}
\mathfrak{o}^{\mathsf{m}}_{\mathsf{K}}(\mathsf{W})(b,b')
:=\sum_{\mathsf{W}'\subset\mathsf{W}}(-1)^{|\mathsf{W}\setminus\mathsf{W}'|}\,
\mathbb{O}^{\mathsf{m}}\bigl[\mathfrak{B}(\mathsf{K}\cup\mathsf{W}')\bigr](b,b'),
\end{equation}
where $\mathfrak{B}(Y)$ is the set of bonds of $Y$. Then, with $\mathsf{s}_\square$ the
interpolation parameter of the cube $\square$,
\begin{equation}\label{9.1:eq-zone-split-4}
\mathbb{O}^{\mathsf{m}}\langle\mathsf{s}\rangle(b,b')
=\sum_{\mathsf{W}}\ \prod_{\square\in\mathsf{W}\setminus\mathsf{K}_0}\mathsf{s}_\square\cdot
\mathfrak{o}^{\mathsf{m}}_{\mathsf{K}}(\mathsf{W})(b,b').
\end{equation}
\item[(h5)] (The zone-extension lemma, and the mixed-difference lemma.) Clauses (a), (c) and (d4)
of the family lemma hold for polymers $Y$ with $\vartheta_w^{k}$ replaced by
$\vartheta_w^{j_{\min}(Y)}$; clause (b) of the family lemma is not among the clauses so extended.
Clause (b) of the mixed-difference lemma holds in the mixed-size form stated at the end of the
zone-extension lemma, with normalising constant $2^{|\mathsf{W}|}C_0B_0$ and decay exponent
$\tfrac14\delta_0M$.
\item[(h5$'$)] (The majorant of the proof of clause (b) of the family lemma.) That proof attaches
to every set $\mathsf{W}$ and bond pair $(b,b')$ a length
$d^{\circ}=d^{\circ}(\mathsf{K},\mathsf{W};b,b')\ge0$, and shows, under the lower bound on $M$
assumed in the family lemma, that for every operator of the list of (h4), the two factors of the
walk operator described in the proof below included, the kernel
\begin{equation*}
(b,b')\longmapsto\sum_{\mathsf{W}}e^{\kappa_1|\mathsf{W}\setminus\mathsf{K}_0|}\,
C\,2^{|\mathsf{W}|}\,e^{-\delta|b-b'|/4}\,e^{-\frac18\delta_0Md^{\circ}}
\end{equation*}
has norm at most $C_Fe^{16\kappa_1}$ in the norm $\|\cdot\|_{\delta/4}$ of (iii), which is
submultiplicative; in this series the factor $2^{|\mathsf{W}|}$ is absorbed by
$e^{-\frac18\delta_0Md^{\circ}}$, so that no factor exponential in the number of cubes remains. The
one-run resummation of the same proof gives
$\|\mathbb{O}^{\mathsf{m}}\langle\mathsf{s}\rangle\|_{\delta/4}\le C_{res}$ on the polydisc
$|\mathsf{s}_\square|\le e^{\kappa_1}$ for these operators.
\item[(h6)] (The kernel-difference corollary.) The difference of the two runs' kernels at the
$\mathbf{T}$-site is bounded by $C_K\vartheta_K^{\,j_{\min}}$ in the norm of decay rate
$\delta_0/32$, where $\vartheta_K$ is the rate of that corollary and $j_{\min}$ is the lowest zone
met by the support of the kernel.
\item[(h7)] (Clause (c) of the family lemma, with (h5).) In zone $n$ the activities of the
differences of the objects (z2) are at most $C_{cl}\vartheta_w^{\,n}(\log x_n)^{c_2}\bar{\varepsilon}^{pr}$
in Ba{\l}aban's activity norm.
\item[(h8)] (The normaliser-localisation corollary.) The normaliser on a region $R$ localises as
$\sum_X\mathcal{L}_R(X)$, with $\mathcal{L}_R$ natural and equal to $\mathcal{L}^{nat}$ inside $R$;
its clause (e4) gives the zone-$n$ difference $\delta\hat{\mathcal{L}}_R$ the relative bound
$C_F\vartheta_w^{\,n}$, and the same for the boundary terms.
\item[(h9)] (The one-run comparison proposition at the $\mathbf{T}$-site.) Its one-run bound (a),
which is \cite[(2.16)]{B-CMP116}, places $\mathsf{D}^{\mathsf{m}}$ and $\Gamma^{\mathsf{m}}$ within
$\varepsilon_{pr}$, in the norm of the tolerance lemma (h10), of one real reference
$\mathsf{C}_{ref}^{-1}$ (and the corresponding reference for
$\Gamma$), and the constant $\alpha^{*}$ of that proposition satisfies
$2d(LM)^4\alpha^{*}\le\mathsf{b}_{cl}$.
\item[(h10)] (The tolerance lemma and its complex form.) The norm $\|\cdot\|$ of the tolerance lemma
is dominated by the norm $\|\cdot\|_{\delta/4}$ of (iii). An operator with the block and symmetry
structure of the tolerance lemma whose differences $\mathsf{E}_{\mathsf{D}}$, $\mathsf{E}_\Gamma$
from the real reference satisfy
$\|\mathsf{E}_{\mathsf{D}}\|,\|\mathsf{E}_\Gamma\|\le\varepsilon_{tol}$ lies in the ball of the
tolerance lemma; the complex form of that lemma requires one such ball along the segment considered.
\item[(h11)] (The zone construction, for the zones of Notation~\ref{9.1:not-zones-boundary-terms},
and the window on the running couplings fixed with the zones.) The zones present in an aligned
label have index $n\ge n_c$, and $\log x_n\ge1$ for every zone present; the rates satisfy
$\vartheta_w\le\vartheta_2$; and, by that window,
$\log x_{n_c}\ge\log(\gamma^{-2}-2c_2)$.
\end{itemize}
Then for every zone $n\le i$ the following hold.
\begin{itemize}
\item[(i)] $\sum_n\delta\mathsf{z}^{(n)}=\mathsf{z}^{\mathsf{B}}-\mathsf{z}^{\mathsf{A}}$ for each of
the objects (z1)--(z3).
\item[(ii)] (Locality.) A $\mathbf{T}$-site output with domain $X$ reads only pieces with
$\operatorname{supp}\omega\subset X$; hence $\delta\mathsf{z}^{(n)}{\upharpoonright}X\equiv0$ unless
$X$ meets zone $n$.
\item[(iii)] (Local index.) Write $\|\cdot\|_{\delta'}$ for the norm of a kernel over bond pairs
$(b,b')$ with weight $e^{\delta'|b-b'|}$, and $\delta$ for the decay rate of the family lemma. Then
$\|\delta\mathsf{D}^{(n)}{\upharpoonright}X\|_{\delta/4}$ and
$\|\delta\Gamma^{(n)}{\upharpoonright}X\|_{\delta/4}$ are at most $\mathfrak{d}^{op}_n$; the class-$n$
pieces of (z2) are potentials of activity at most $\mathfrak{d}^{op}_n\bar{\varepsilon}^{pr}$ in
Ba{\l}aban's activity norm; and the class-$n$ normaliser pieces obey
\begin{equation}\label{9.1:eq-zone-split-5}
|\delta\mathcal{L}^{(n)}(X'')|\le\frac{\mathfrak{d}^{op}_n}{C_{cl}^T}\,C_F\,\bar{\mathcal{L}}(X''),
\end{equation}
where $\bar{\mathcal{L}}(X'')$ is the one-run localised bound of the normaliser at $X''$.
\item[(iv)] (Budget.) For every zone $n$ present,
$\mathfrak{d}^{op}_n\le(C_{cl}^T/C^{\natural}_{cl})\,\vartheta_\times^{\,n}$, so that
\begin{equation}\label{9.1:eq-zone-split-6}
\sum_{n\ \text{present}}\mathfrak{d}^{op}_n\le(1-\vartheta_\times)^{-1}
\frac{C_{cl}^T}{C^{\natural}_{cl}}\,\vartheta_\times^{\,n_c}=:\mathfrak{d}^{op}_{*},
\end{equation}
and $\mathfrak{d}^{op}_{*}\to0$ as $\gamma\to0$ (here $n_c\asymp\log\log X$).
\end{itemize}
Finally, under the ceiling on $\gamma$
\begin{equation}\label{9.1:eq-zone-split-a}
\frac{E_1^{z}}{E_0}\bigl(1+\bar{\varepsilon}^{pr}+C_F\bigr)\,\mathfrak{d}^{op}_{*}
\le\tfrac14\,\mathfrak{m}_G^T ,
\end{equation}
and provided $\varepsilon_{pr}+\mathfrak{d}^{op}_{*}+\tfrac12\mathfrak{m}_G^T\le\varepsilon_{tol}$,
both runs' $\mathbf{T}$-site operators, every convex combination of them, and every further complex
perturbation of relative size at most $\tfrac12\mathfrak{m}_G^T$ lie in the ball of the tolerance
lemma about one real reference: the one-ball hypothesis of the complex form of the tolerance lemma
holds along the whole segment for the $\mathbf{B}$-outputs of aligned labels.
\end{lemma}

\begin{proof}
\emph{(i).} By (h1) the pieces $\omega$ form the same index set for both runs, and
$\mathsf{z}^{\mathsf{m}}=\sum_\omega\mathsf{z}^{\mathsf{m}}_\omega$ converges absolutely. Hence
$\mathsf{z}^{\mathsf{B}}-\mathsf{z}^{\mathsf{A}}=\sum_\omega\delta\mathsf{z}_\omega$, again absolutely
convergent. Each $\omega$ has exactly one lowest zone $n(\omega)$ met by its support, so regrouping
the absolutely convergent sum according to the value of $n(\omega)$ gives
\begin{equation*}
\mathsf{z}^{\mathsf{B}}-\mathsf{z}^{\mathsf{A}}
=\sum_n\ \sum_{\omega:\ n(\omega)=n}\delta\mathsf{z}_\omega=\sum_n\delta\mathsf{z}^{(n)} .
\end{equation*}

\emph{(ii).} By the locality lemma for outputs (h2), an output with domain $X$ reads only the
localised terms $V(Y)$ with $Y\subset X$. The walks and the decoupled cubes of a cluster-expansion
term lie in its domain,
since in the expansion of \cite[(2.9)--(2.13)]{B-CMP116} the set $X$ of a polymer is the hull of the
supports of its objects. Hence an output with domain $X$ reads only pieces with
$\operatorname{supp}\omega\subset X$. If such a piece has $n(\omega)=n$, then
$\operatorname{supp}\omega$ meets zone $n$ and lies in $X$, so $X$ meets zone $n$. Therefore, if $X$
does not meet zone $n$, the sub-sum defining $\delta\mathsf{z}^{(n)}{\upharpoonright}X$ is empty and
$\delta\mathsf{z}^{(n)}{\upharpoonright}X\equiv0$.

\emph{(iii), the operators (z1).} The bound is obtained from the proof of clause (b) of the family
lemma together with clause (a) in the form of the zone-extension lemma (h5), applied to the M\"obius
pieces \eqref{9.1:eq-zone-split-3}. By (h4) the operators of (z1) are those of the list of clause (b)
of the family lemma, and each such $\mathbb{O}^{\mathsf{m}}\langle\mathsf{s}\rangle$ has the
expansion \eqref{9.1:eq-zone-split-4}. The index set $\{\mathsf{W}\}$, the finite sets of cubes of
the label's partition, is common to both runs, because both runs stand on the one scaffold and the
one interpolation scheme of (h1). For a pair $(\mathsf{K},\mathsf{W})$ put
$n(\mathsf{K},\mathsf{W}):=j_{\min}(\mathsf{K}\cup\mathsf{W})$, let
$\delta\mathfrak{o}_{\mathsf{K}}(\mathsf{W}):=\mathfrak{o}^{\mathsf{B}}_{\mathsf{K}}(\mathsf{W})
-\mathfrak{o}^{\mathsf{A}}_{\mathsf{K}}(\mathsf{W})$, and define
\begin{equation}\label{9.1:eq-zone-split-7}
\delta\mathbb{O}^{(n)}\langle\mathsf{s}\rangle
:=\sum_{\mathsf{W}:\ j_{\min}(\mathsf{K}\cup\mathsf{W})=n}\ \prod_{\square\in\mathsf{W}\setminus\mathsf{K}_0}
\mathsf{s}_\square\cdot\delta\mathfrak{o}_{\mathsf{K}}(\mathsf{W}).
\end{equation}
By (i), applied to the pieces of \eqref{9.1:eq-zone-split-4}, these sum over $n$ to
$\mathbb{O}^{\mathsf{B}}\langle\mathsf{s}\rangle-\mathbb{O}^{\mathsf{A}}\langle\mathsf{s}\rangle$.

Clause (a) of the family lemma in the form of the zone-extension lemma, together with clause (b) of
the mixed-difference lemma in the mixed-size form stated at the end of the zone-extension lemma,
with normalising constant $2^{|\mathsf{W}|}C_0B_0$ and decay exponent $\tfrac14\delta_0M$ (both in
(h5)), gives, termwise, the bound of the proof of clause (b) of the family lemma with the level $k$
replaced by $j_{\min}$:
\begin{equation}\label{9.1:eq-zone-split-8}
|\delta\mathfrak{o}_{\mathsf{K}}(\mathsf{W})(b,b')|
\le C\,2^{|\mathsf{W}|}\,\vartheta_w^{\,j_{\min}(\mathsf{K}\cup\mathsf{W})}\,
e^{-\delta|b-b'|/4}\,e^{-\frac18\delta_0Md^{\circ}},
\end{equation}
where $d^{\circ}$ is the length of (h5$'$). On the polydisc $|\mathsf{s}_\square|\le e^{\kappa_1}$
one has
$\prod_{\square\in\mathsf{W}\setminus\mathsf{K}_0}|\mathsf{s}_\square|\le
e^{\kappa_1|\mathsf{W}\setminus\mathsf{K}_0|}$. Hence the sum of the bounds
\eqref{9.1:eq-zone-split-8}, multiplied by $\prod|\mathsf{s}_\square|$, over the sets $\mathsf{W}$
with $j_{\min}=n$ is at most $\vartheta_w^{\,n}$ times a sub-sum of the series of nonnegative terms
of (h5$'$), because on this sub-family the factor $\vartheta_w^{\,j_{\min}}$ equals
$\vartheta_w^{\,n}$; by (h5$'$) that series has norm at most $C_Fe^{16\kappa_1}$. The restriction to
a domain $X$ passes to a further sub-sum. Hence
\begin{equation}\label{9.1:eq-zone-split-9}
\|\delta\mathsf{D}^{(n)}{\upharpoonright}X\|_{\delta/4}\le C_Fe^{16\kappa_1}\vartheta_w^{\,n}
\qquad\text{on the polydisc},
\end{equation}
and the same bound holds for $\delta\mathbb{O}^{(n)}$ for every operator $\mathbb{O}$ to which
(h5$'$) applies. As recorded in (h5$'$), the factor $2^{|\mathsf{W}|}$ is absorbed by
$e^{-\frac18\delta_0Md^{\circ}}$ under the lower bound on $M$ assumed in the family lemma, so that no
factor exponential in the number of cubes remains. For the walk operator, which is the product
$\Gamma=\Gamma_1\Gamma_2$ of $\Gamma_1$, the
block between $Z_0$ and its complement $Z_0^c$ of the $\mathsf{s}$-interpolated operator of
Ba{\l}aban's fluctuation step that enters $\Gamma$ (interpolated as in
\cite[(2.5)--(2.13)]{B-CMP116}), and $\Gamma_2$, the $\mathsf{s}$-interpolated square root of the
covariance, we put
\begin{equation*}
\delta\Gamma^{(n)}:=\delta\Gamma_1^{(n)}\,\Gamma_2^{\mathsf{B}}+\Gamma_1^{\mathsf{A}}\,\delta\Gamma_2^{(n)},
\end{equation*}
which by (i) for $\Gamma_1$ and $\Gamma_2$ sums over $n$ to
$\Gamma_1^{\mathsf{B}}\Gamma_2^{\mathsf{B}}-\Gamma_1^{\mathsf{A}}\Gamma_2^{\mathsf{A}}
=\Gamma^{\mathsf{B}}-\Gamma^{\mathsf{A}}$. The differenced factors $\delta\Gamma_1^{(n)}$,
$\delta\Gamma_2^{(n)}$ obey the bound \eqref{9.1:eq-zone-split-9}, since (h5$'$) applies to
$\Gamma_1$ and $\Gamma_2$; the undifferenced factors $\Gamma_2^{\mathsf{B}}$, $\Gamma_1^{\mathsf{A}}$
obey the one-run bound $C_{res}$ of (h5$'$). By the submultiplicativity of the norm in (h5$'$),
with the restriction to $X$ acting on the differenced factors,
\begin{equation*}
\|\delta\Gamma^{(n)}{\upharpoonright}X\|_{\delta/4}\le2C_{res}C_Fe^{16\kappa_1}\vartheta_w^{\,n}
\qquad\text{on the polydisc}.
\end{equation*}
Since $j_{\min}(\mathsf{K}\cup\mathsf{W})$ does not depend on the
order of $(b,b')$, the pieces \eqref{9.1:eq-zone-split-7} keep the symmetry in $(b,b')$; hence the
block and symmetry clauses of the tolerance lemma hold for $\mathsf{D}^\lambda+\eta\,\delta\mathsf{D}^{(n)}$,
where $\mathsf{D}^\lambda:=\lambda\mathsf{D}^{\mathsf{B}}+(1-\lambda)\mathsf{D}^{\mathsf{A}}$,
$\lambda\in[0,1]$ (no relation to the dial $\lambda_{j,k}$ of (h3)), and $\eta\in\mathbb{C}$. Since
$C_{cl}^T\ge C_Fe^{16\kappa_1}$ and $C_{cl}^T\ge2C_{res}C_Fe^{16\kappa_1}$, since
$\vartheta_T=\vartheta_w$, and since $(\log x_n)^{c_2}\ge1$ for every zone present by (h11), the
bounds just obtained for $\delta\mathsf{D}^{(n)}$ and $\delta\Gamma^{(n)}$ give the first assertion
of (iii) with $\mathfrak{d}^{op}_n$ of \eqref{9.1:eq-zone-split-2}.

\emph{(iii), the objects (z2) and (z3).} For (z2), hypothesis (h7), which is clause (c) of the family
lemma read in zone $n$ through the zone-extension lemma (h5) and the filing rule (h3), bounds the
class-$n$ activities by $C_{cl}\vartheta_w^{\,n}(\log x_n)^{c_2}\bar{\varepsilon}^{pr}$; as
$C_{cl}^T$ collects $C_{cl}$ and $\vartheta_T=\vartheta_w$, this is
$\mathfrak{d}^{op}_n\bar{\varepsilon}^{pr}$ by \eqref{9.1:eq-zone-split-2}. For (z3), clause (e4) of
the normaliser-localisation corollary (h8) gives the class-$n$ normaliser pieces, and the same for
the boundary terms, the relative bound $C_F\vartheta_w^{\,n}$ against the one-run localised bound
$\bar{\mathcal{L}}(X'')$; since
$\mathfrak{d}^{op}_n/C_{cl}^T=(\log x_n)^{c_2}\vartheta_T^{\,n}\ge\vartheta_w^{\,n}$, because
$\vartheta_T=\vartheta_w$ and $\log x_n\ge1$ for every zone present by (h11), this gives
\eqref{9.1:eq-zone-split-5}. The constant $C_K$ of the kernel-difference corollary (h6) and the
one-run resummation constants of the three objects enter only through $C_{cl}^T$.

\emph{(iv).} By \eqref{9.1:eq-zone-split-2} and the definition of $\vartheta_\times$ in
\eqref{9.1:eq-zones-boundary-terms-b}, which gives $\vartheta_\times=\vartheta_w\vartheta_2^{-1/4}<1$,
\begin{equation*}
\vartheta_T^{\,n}(\log x_n)^{c_2}
=\Bigl[C^{\natural}_{cl}(\log x_n)^{c_2}\vartheta_2^{n/4}\Bigr]\cdot
\frac{\vartheta_T^{\,n}}{C^{\natural}_{cl}\vartheta_2^{n/4}}
\le\frac{\vartheta_\times^{\,n}}{C^{\natural}_{cl}},
\end{equation*}
because the bracket is at most $1$ by \eqref{9.1:eq-zones-boundary-terms-a} for every zone $n$
present, and $\vartheta_T^{\,n}\vartheta_2^{-n/4}=\vartheta_\times^{\,n}$. Multiplying by
$C_{cl}^T$ gives $\mathfrak{d}^{op}_n\le(C_{cl}^T/C^{\natural}_{cl})\vartheta_\times^{\,n}$. By (h11)
the zones present have $n\ge n_c$, and since $0<\vartheta_\times<1$,
\begin{equation*}
\sum_{n\ge n_c}\vartheta_\times^{\,n}\le(1-\vartheta_\times)^{-1}\vartheta_\times^{\,n_c};
\end{equation*}
this gives \eqref{9.1:eq-zone-split-6}. For the limit,
\begin{align*}
\vartheta_\times^{\,n_c}=\vartheta_w^{\,n_c}\vartheta_2^{-n_c/4}
&\le\vartheta_2^{3n_c/4}
\le\bigl(C^{\natural}_{cl}(\log x_{n_c})^{c_2}\bigr)^{-3}\\
&\le\bigl(C^{\natural}_{cl}(\log(\gamma^{-2}-2c_2))^{c_2}\bigr)^{-3}.
\end{align*}
Here the first inequality is the rate inequality $\vartheta_w\le\vartheta_2$ of (h11); the second is
\eqref{9.1:eq-zones-boundary-terms-a} at $n=n_c$, which bounds the bracket above by $1$ there, so that
$\vartheta_2^{n_c/4}\le(C^{\natural}_{cl}(\log x_{n_c})^{c_2})^{-1}$, raised to the third power; and the
third is the window $\log x_{n_c}\ge\log(\gamma^{-2}-2c_2)$ of (h11). Hence,
so that
\begin{equation}\label{9.1:eq-zone-split-10}
\mathfrak{d}^{op}_{*}\le(1-\vartheta_\times)^{-1}\frac{C_{cl}^T}{C^{\natural}_{cl}}
\bigl(C^{\natural}_{cl}(\log(\gamma^{-2}-2c_2))^{c_2}\bigr)^{-3},
\end{equation}
whose right member depends on $\gamma$ and on no zone index, and tends to $0$ as $\gamma\to0$. In
particular \eqref{9.1:eq-zone-split-a} holds whenever the same inequality holds with the right member
of \eqref{9.1:eq-zone-split-10} in place of $\mathfrak{d}^{op}_{*}$, a condition on $\gamma$ alone.
The ceiling \eqref{9.1:eq-zone-split-a} is thus a condition of the same kind as the ceiling
$\mathfrak{d}(\gamma):=(E_1/E_0)(\log\gamma^{-2})^{-c_2}\le\mathfrak{m}_G$ of the one-run comparison
proposition at the $\mathbf{T}$-site: by \eqref{9.1:eq-zone-split-10} its left member is bounded by a
function of $\gamma$ alone of negative degree in $\log\gamma^{-2}$.

\emph{The last assertion.} It is the hypothesis of the tolerance lemma (h10),
$\|\mathsf{E}_{\mathsf{D}}\|,\|\mathsf{E}_\Gamma\|\le\varepsilon_{tol}$, with
\begin{equation*}
\mathsf{E}_{\mathsf{D}}:=\bigl(\mathsf{D}^\lambda+\eta\,\delta\mathsf{D}^{(n)}\bigr)-\mathsf{C}_{ref}^{-1},
\end{equation*}
and $\mathsf{E}_\Gamma$ defined in the same way from $\Gamma^\lambda$. Write
\begin{equation*}
\mathsf{D}^\lambda-\mathsf{C}_{ref}^{-1}=(\mathsf{D}^{\mathsf{A}}-\mathsf{C}_{ref}^{-1})
+\lambda(\mathsf{D}^{\mathsf{B}}-\mathsf{D}^{\mathsf{A}}).
\end{equation*}
By the triangle inequality,
\begin{equation*}
\|\mathsf{E}_{\mathsf{D}}\|\le\|\mathsf{D}^{\mathsf{A}}-\mathsf{C}_{ref}^{-1}\|
+\lambda\sum_{n'\ \text{present}}\|\delta\mathsf{D}^{(n')}\|+\|\eta\,\delta\mathsf{D}^{(n)}\|
\le\varepsilon_{pr}+\mathfrak{d}^{op}_{*}+\tfrac12\mathfrak{m}_G^T\le\varepsilon_{tol}.
\end{equation*}
Here the sum over zones present uses (i), the terms of absent zones being empty sums; the first term
is bounded by the one-run bound (a) of (h9), which is \cite[(2.16)]{B-CMP116}; the second by (iii)
and the budget \eqref{9.1:eq-zone-split-6}, the norm of the tolerance lemma being dominated by
$\|\cdot\|_{\delta/4}$ by (h10); the third by the assumption on the relative size of the
perturbation; and the last inequality is the assumed one. The same chain bounds
$\|\mathsf{E}_\Gamma\|$. Together with $2d(LM)^4\alpha^{*}\le\mathsf{b}_{cl}$ from (h9) and the block
and symmetry clauses established in (iii), all hypotheses of the tolerance lemma hold, uniformly in
$\lambda\in[0,1]$; this is the one-ball hypothesis of its complex form along the whole segment, for
the $\mathbf{B}$-outputs of aligned labels.
\end{proof}

\begin{remark}
Beyond its hypotheses, the lemma contributes only the filing \eqref{9.1:eq-zone-split-1} and two
summations, the regrouping of (i) and the sum over zones of (iv).
\end{remark}

\begin{definition}[Zone-resolved members of the stored boundary terms]\label{9.1:def-resolved-members}
We use the two runs $\mathsf{A}$ and $\mathsf{B}$, of lengths $\mathfrak{n}$ and $\mathfrak{n}+1$, built on
the one scaffold $\Theta_X$ at the bare point $X=g^{-2}$ (in items (a)--(d) the letter $X$ otherwise
denotes a region); the levels indexed by the index of $\mathsf{A}$; the regions
$\Omega_1\supset\Lambda_1\supset\Omega_2\supset\Lambda_2\supset\cdots$ of a label; the zones; and the
large-field regions $Z_j=\Lambda_j^c$. All of these are as in
Notation~\ref{9.1:not-zones-boundary-terms}. The set of aligned labels is $H_{al}$.

We write $\mathbf{B}^{(i+1)}(X)$ for a stored boundary term produced by the fluctuation step
$i\to i+1$; we call it \emph{T-born}. We write $\mathbf{B}'^{(i)}(X')$ for a stored boundary term
produced by the step $\mathbf{R}$ at level $i$; we call it \emph{R-born}.

A \emph{class} is a pair $z=(n,k')$ consisting of a zone $n$ and a reading level $k'$. We write
$n(z)=n$ and call $n$ the zone of $z$.

A stored boundary term that is read, as an input, in a later production or at the top level is
called a \emph{read term}. For a read term $u$ and a run $\mathsf{m}\in\{\mathsf{A},\mathsf{B}\}$, we
write $\mathrm{rd}^{\mathsf{m}}P^{\mathsf{m}}$ for the datum at which run $\mathsf{m}$ reads $u$.

\begin{enumerate}
\item[(a)] \emph{Production arrays and marks.} Let the label be aligned. Let the term be a T-born
$\mathbf{B}^{(i+1)}(X)$, or, respectively, an R-born $\mathbf{B}'^{(i)}(X')$; in the latter case $X$
stands for $X'$ throughout this item.

The two runs are built on one scaffold. The productions of the two runs satisfy the closure
property stated in the first clause of the closure statement for these productions. By these two
facts the term is one formula $\mathcal{P}_X$, the same for both runs, applied to the array
$x^{\mathsf{m}}=(x^{\mathsf{m}}_u)_u$ of the values, in run $\mathsf{m}$, of the slots $u$ of its
production. The slots are of three kinds.
\begin{itemize}
\item The set $U^{dir}$ of \emph{direct uses}. It consists of the uses of the stored units $E$,
$P^{\flat}$, $Q$, $R$, $D$ of levels $\le i$ that have pieces in $X$, both kinds of piece being
counted in the mark. It also contains the $\zeta$-entries, the $b$-entries, the weight entries and
the stored constants of the third group that the load $[\![\delta]\!]^{(5)}_n$ of item~(d) carries.
All of these are the uses listed in \eqref{9.1:eq-centre-normalisation-c}.
\item The set $U^{lf}$ of \emph{read terms in $X$}. It consists of the stored derived terms read in
$X$: the terms $\mathbf{B}^{(\ell)}$ and $\mathbf{B}'^{(\ell)}$ with $\ell\le i$, and the left-over
chain terms $\mathbb{C}$ and $V(X^{\sharp})$ of the fourth group of stored quantities of the
correlation theorem.
\item The set $U^{\mathsf{z}}$ of the \emph{$z$-free pieces} inside $X$.
  \begin{itemize}
  \item At the fluctuation step these are the pieces of Lemma~\ref{9.1:lem-zone-split}.
  \item At the sub-steps of the $\mathbf{R}$-step that produce $\mathbf{B}'$ (the operation
  $\mathbf{T}'$ and the Mayer expansion among them) these are the $z$-free pieces and the chain
  pieces of those sub-steps, together with the kernel classes of the load statement for the
  productions.
  \end{itemize}
\end{itemize}
The formula reproduces the term of each run: at every common slice datum $\mathsf{v}$ and every local
label,
\begin{equation*}
\mathcal{P}_X[x^{\mathsf{A}}]=\mathbf{B}^{\mathsf{A}}(X;\mathsf{v}),\qquad
\mathcal{P}_X[x^{\mathsf{B}}]=\mathbf{B}^{\mathsf{B}}(X;\mathsf{v}),
\end{equation*}
where for seeds the stored values are those of the storage rule for seeds.

The \emph{mark} of a class $z=(n,k')$ is the array $W^{(z)}$ defined slot by slot as follows.
\begin{itemize}
\item For $u\in U^{dir}\cup U^{\mathsf{z}}$, the $u$-entry of $W^{(z)}$ is
$x^{\mathsf{B}}_u-x^{\mathsf{A}}_u$ if $u$ is of class $z$, and $0$ otherwise.
\item For $u\in U^{lf}$, the $u$-entry of $W^{(z)}$ is the sum of those among the stored class-$n$
member $G_u^{(n)}$, the reading part $G_u^{rd}$ and the own member $G_u^{own}$ of item~(c) to which
item~(c) assigns the class $z$. Each of them is assigned exactly one class.
\end{itemize}
In formula, entry by entry,
\begin{equation}\label{9.1:eq-resolved-members-1}
\bigl(W^{(z)}\bigr)_u=
\begin{cases}
x^{\mathsf{B}}_u-x^{\mathsf{A}}_u, & u\in U^{dir}\cup U^{\mathsf{z}}\text{ of class }z,\\
0, & u\in U^{dir}\cup U^{\mathsf{z}}\text{ not of class }z,\\
\text{the sum of those of }G_u^{(n)},\,G_u^{rd},\,G_u^{own}\text{ of class }z, & u\in U^{lf}.
\end{cases}
\end{equation}
The marks sum to the two-run difference:
\begin{equation}\label{9.1:eq-resolved-members-2}
\sum_z W^{(z)}=x^{\mathsf{B}}-x^{\mathsf{A}}.
\end{equation}

\item[(b)] \emph{Members.} Put $x^{\lambda}:=(1-\lambda)x^{\mathsf{A}}+\lambda x^{\mathsf{B}}$ for
$\lambda\in[0,1]$. Let $\mathcal{Z}$ denote the set of classes present in the production. For a
T-born term and $z=(n,\cdot)\in\mathcal{Z}$ we define the member
\begin{equation}\label{9.1:eq-resolved-members-3}
\delta\mathbf{B}^{(i+1),(n)}(X;\mathsf{v}):=\mathfrak{V}^{(z)}
:=\int_0^1 d\lambda\,\frac{d}{d\epsilon}\Big|_{\epsilon=0}
\mathcal{P}_X\bigl[x^{\lambda}+\epsilon W^{(z)}\bigr].
\end{equation}
This is the first-variation formula for the productions.

For an R-born term $\mathbf{B}'^{(i)}(X')$ we split $\mathcal{Z}$ as a disjoint union
$\mathcal{Z}=\mathcal{Z}_{surv}\sqcup\mathcal{Z}_{own}$.

The healed components of $Z_i=\Lambda_i^c$ are its components of the first kind in the sense of
\cite{B-CMP122b}; their set is written
$Z(i)$. A class $z=(n,\cdot)$ belongs to $\mathcal{Z}_{surv}$ if and only if two conditions hold:
\begin{itemize}
\item $n\le i-1$;
\item some slot of class $z$ of $\mathcal{P}_{X'}$ lies outside every component in $Z(i)$, outside
every live component of the second kind, and outside every live component arrested in mid-run.
\end{itemize}
The set $\mathcal{Z}_{own}$ consists of all remaining classes. It contains, in particular, the
following three kinds of class.
\begin{itemize}
\item[(b1)] Every class all of whose slots lie inside the healed components. Their slots lie
inside what the operations $\mathbf{T}'_i$ contained in $R'^{(i)}(X')$ sum and integrate over. That
is, they lie inside the sequences $\{\Omega_j\cap X_j,\,Z_j\cap X_j\}$ belonging to the components,
over which the sums of \cite[(1.71)--(1.72)]{B-CMP122b} run; the terms of $R'^{(i)}(X')$ contain at
least one large-field region connected with one of the operations $\mathbf{T}'$
\cite{B-CMP122b}. The preparatory integrations of these components, as in \cite{B-CMP122a}, are
included. Since $Z_i=\Lambda_i^c$, such a component contains its core and every zone layer around
it.
\item[(b2)] The then-top class, of zone $n=i$.
\item[(b3)] Every class whose slots lie inside a component that $\mathbf{R}_i$ opened and did not
heal. Such a component is of one of two kinds.
  \begin{itemize}
  \item A component $Y$ of the second kind, as in \cite{B-CMP122b}. Such a $Y$ keeps the factor
  $\mathbf{T}_k(Y)$ of \cite[(1.102)]{B-CMP122b}, in which $\mathbf{T}'_k(Y)$ enters as a factor;
  here $\mathbf{T}'_k(Y)$ sums over the own sequence of $Y$ of
  \cite[(1.71)]{B-CMP122b} and acts also on the last exponential in \cite[(1.72)]{B-CMP122b}.
  \item A component arrested in mid-run, as in \cite{B-CMP122a}.
  \end{itemize}
  Classes of this kind are not in $\mathcal{Z}_{surv}$, for two reasons. First, their zones $n<i$
  leave every determining set after $\mathbf{R}_i$ \cite{B-CMP122a}. Second, they cannot be dialled
  by the splitting through which the surviving classes are dialled, because
  $\mathbf{T}'_i(Y)$ stands outside $R'^{(i)}(X')$. They are carried in the label of the term,
  which consists of the history together with the outcomes of the events. The lemma on geometric
  forgetting of old regions is stated for every component, of any history, and so applies to them.
\end{itemize}
For $z=(n,\cdot)\in\mathcal{Z}_{surv}$, so that $n\le i-1$, we put
$\delta\mathbf{B}'^{(i),(n)}:=\mathfrak{V}^{(z)}$, with $\mathfrak{V}^{(z)}$ given by
\eqref{9.1:eq-resolved-members-3} with $\mathcal{P}_{X'}$ in place of $\mathcal{P}_X$.

The second index of a class $z=(n,\cdot)$ is fixed. At each step, for each read term, at most one
class of zone $n$ having a slot outside the components, that is, at most one class of zone $n$ in
$\mathcal{Z}_{surv}$, is produced, for two reasons: the classes produced at a level $k$ have reading
level $k'=k$; and in \cite[(1.71)--(1.72)]{B-CMP122b} the sum over sequences stands inside
$\mathbf{T}'$ for components of either kind. Hence $\delta\mathbf{B}'^{(i),(n)}$ is one member for
each pair (read term, $n$), and no factor $(i-n+1)$ counting classes of zone $n$ enters the
coefficients $\varkappa^{\mathbf{B}'}_i$ of the zone-resolved bound for R-born terms.

The own member is defined by
\begin{equation}\label{9.1:eq-resolved-members-4}
\delta\mathbf{B}'^{(i),own}:=\mathbf{B}'^{(i),\mathsf{B}}-\mathbf{B}'^{(i),\mathsf{A}}
-\sum_{z\in\mathcal{Z}_{surv}}\mathfrak{V}^{(z)}.
\end{equation}

\emph{Stored data.} The stored data are
$(\mathbf{B}^{\mathsf{A}},\mathbf{B}^{\mathsf{B}};(\delta\mathbf{B}^{(n)})_n)$, together with
$\delta\mathbf{B}'^{own}$ for R-born terms. They are functions of the common datum; the parameters
$\lambda$ and $\epsilon$ are not stored.

\emph{Norms.} Put $\kappa_{\mathbf{B}}:=\kappa$. This is the rate at which stored boundary terms are
consumed as read terms: the input bound for read terms in the productions holds at a rate
$\kappa_T\ge\kappa$ and is used here at $\kappa$. The rate $\kappa_{\mathbf{B}}$ lies below the
production rates $\kappa_{out}$ and $(1+\tfrac12\beta_s)\kappa$. We set
\begin{equation}\label{9.1:eq-resolved-members-5}
\mathsf{S}^{+,(n)}_{\ell}[\mathbf{B}]:=\sup_X e^{\kappa_{\mathbf{B}}d_{\ell}(X)}
\sup_{\mathsf{v},\,\text{label}}\bigl|\delta\mathbf{B}^{(\ell),(n)}(X;\mathsf{v})\bigr|.
\end{equation}
The norms $\mathsf{S}^{+,(n)}_{\ell}[\mathbf{B}'^{res}]$ and
$\mathsf{S}^{+}_{\ell}[\mathbf{B}'^{own}]$ are defined in the same way, with
$\delta\mathbf{B}'^{(\ell),(n)}$, respectively $\delta\mathbf{B}'^{(\ell),own}$, in place of
$\delta\mathbf{B}^{(\ell),(n)}$, for those $\mathbf{B}'$ whose assigned domain is plain.

For those $\mathbf{B}'$ whose assigned domain is island-blind, the norms are the
relative sup-norms
\begin{equation*}
\mathsf{S}^{\circ,(n)}_{\ell}:=\sup_Y e^{\kappa_{\mathbf{B}}d^{\circ}_{\ell}(Y)}\sup|\cdot|,
\end{equation*}
where $d^{\circ}_{\ell}$ is the length off the live islands that the assigned domain contains; this
is Ba{\l}aban's relative length $d_{k,\cup Y_i}$. In the statements of this chapter that follow,
$\mathsf{S}$ denotes the supremum over both kinds of norm, the island-blind part entering the loads
as numbers.

\item[(c)] \emph{Classes of the marks of a read term at a reading level $k'\ge\ell$.} Let
$u=\mathbf{B}^{(\ell)}(X_u)$, or $u=\mathbf{B}'^{(\ell)}(X_u)$, be a read term, read in a
production or at the top level at the reading level $k'$. Its mark is
\begin{equation*}
\sum_n G_u^{(n)}+G_u^{rd}\quad\Bigl[\,+\,G_u^{own}\text{ for R-born }u\Bigr],
\end{equation*}
with the following three parts.
\begin{itemize}
\item The stored class-$n$ member, of class $(n,k')$:
\begin{equation*}
G_u^{(n)}:=\delta\mathbf{B}^{(\ell),(n)}(X_u;\mathrm{rd}^{\mathsf{A}}P^{\mathsf{A}}),
\end{equation*}
and $G_u^{(n)}:=\delta\mathbf{B}'^{(\ell),(n)}(X_u;\mathrm{rd}^{\mathsf{A}}P^{\mathsf{A}})$ when $u$
is R-born.
\item The reading part, which is classical and of class $(n_{X_u},k')$, where $n_{X_u}$ is the
lowest zone met by $X_u$:
\begin{equation*}
G_u^{rd}:=\mathbf{B}^{(\ell),\mathsf{B}}(X_u;\mathrm{rd}^{\mathsf{B}}P^{\mathsf{B}})
-\mathbf{B}^{(\ell),\mathsf{B}}(X_u;\mathrm{rd}^{\mathsf{A}}P^{\mathsf{A}}).
\end{equation*}
Its affine part in $g^2$ lies in the class $(n_{X_u},k')$ and is accounted for by the last entry
of the classical load $[\![\delta]\!]^c_n$ of item~(d) and by the entries
$w^{\mathbf{B}}B_0\hat{\vartheta}^{\natural}$ of that load.
\item For R-born $u$, the own member, of class $(i_u\vee\ell,k')$, where $i_u$ is the zone of the
block of the use:
\begin{equation*}
G_u^{own}:=\delta\mathbf{B}'^{(\ell),own}(X_u;\mathrm{rd}^{\mathsf{A}}P^{\mathsf{A}}).
\end{equation*}
\end{itemize}

\item[(d)] \emph{Loads.} The five-unknown load of zone $n$ is the load $[\![\delta]\!]^u_n$ of the
two-run differences with its sum over boundary terms removed. It consists of the
entries of the five unknowns:
\begin{equation*}
\begin{split}
[\![\delta]\!]^{(5)}_n:=&\sum_{j\le n}\varsigma_{1+n-j}\,\mathsf{S}^+_j[E]
+\sum_{j\le n}w_{n,j}\,\mathsf{S}^+_j[R]\\
&+\mathfrak{g}_n\Bigl[\mathsf{z}^{\sharp}_n+\sum_{j\le n}L^{-(n-j)/2}\,\mathsf{b}_j\Bigr].
\end{split}
\end{equation*}
Here $w_{n,j}$ are the weights with which the entries of the stored unit $R$ enter this load.
The crossing weights are
\begin{equation}\label{9.1:eq-resolved-members-6}
\mathsf{w}_{\ell,n}:=1\quad(\ell\le n+1),\qquad
\mathsf{w}_{\ell,n}:=e^{-\mathsf{d}_{\ell-1}}\quad(\ell\ge n+2),
\end{equation}
with $\mathsf{d}_m$ the crossing length of zone $m$ of
Notation~\ref{9.1:not-zones-boundary-terms}. With them we put
\begin{equation*}
\begin{gathered}
\mathsf{y}_n^{(\le k')}:=\sum_{n\le\ell\le k'}\Bigl[\mathsf{w}_{\ell,n}
\,\mathsf{S}^{+,(n)}_{\ell}[\mathbf{B}]+\mathsf{S}^{+,(n)}_{\ell}[\mathbf{B}'^{res}]\Bigr],\\
\mathsf{y}'_n:=\sum_{\ell\le n}w^{\mathbf{B}}_{n,\ell}\,\mathsf{S}^+_{\ell}[\mathbf{B}'^{own}].
\end{gathered}
\end{equation*}

The classical load $[\![\delta]\!]^c_n$ is the load of the classical two-run differences, without
change. Its entries for the read boundary terms are kept, and so is its last entry, with its
constant $E_0$,
\begin{equation*}
E_0\min\bigl\{1,\;C^{\natural}_{cl}(\log x_n)^{c_2}\vartheta_2^{n/4}\bigr\}.
\end{equation*}
The stores for the left-over chain terms $\mathbb{C}$ and for $V(X^{\sharp})$ are unchanged.

The total loads are
\begin{equation*}
\mathfrak{a}_n:=[\![\delta]\!]^{(5)}_n+\mathsf{y}'_n+[\![\delta]\!]^c_n,\qquad
[\![\delta]\!]^{\natural+,(k')}_n:=\mathfrak{a}_n+\mathsf{y}_n^{(\le k')},
\end{equation*}
and $[\![\delta]\!]^{\natural+}_n:=\sup_{k'}[\![\delta]\!]^{\natural+,(k')}_n$.

The dial domain $\mathbb{D}^{\star}$ and the class radii
\begin{equation*}
\Lambda_z:=E_1^z/[\![\delta]\!]^{\natural+}_{n(z)}
\end{equation*}
are formed with the constants
\begin{equation*}
E_1^z:=\Bigl(\tfrac54+4C_T^{\mathbf{B}}\Bigr)E_1+\tfrac54B_0,\qquad
C_L''':=(1+K_Y)\,C_L''.
\end{equation*}
Here $E_1$ and $B_0$ are the constants of the lemma that fixes $E_1$ for the dial construction. The
$\gamma$-ceiling condition under which $\mathbb{D}^{\star}$ is defined, stated there with the
constants $2E_1$ and $C_L''$, is required with the substitution
\begin{equation}\label{9.1:eq-resolved-members-b}
2E_1\ \longmapsto\ 2E_1^z,\qquad C_L''\ \longmapsto\ C_L'''.
\end{equation}
The $\gamma$-ceiling condition read with \eqref{9.1:eq-resolved-members-b} consists
of $\gamma$-ceilings of the same form as those of the condition stated with $2E_1$ and $C_L''$.

\item[(e)] \emph{The weight sum.} In this item $X=g^{-2}$ denotes the inverse reference coupling,
not a region, and $\lambda$ denotes the slope of the reference window estimate
$x_i-x_j\ge\lambda(j-i)-2c_2$, not the interpolation parameter of item~(b). For $X-2c_2>1$ put
\begin{equation}\label{9.1:eq-resolved-members-7}
\mathfrak{S}(X):=e^{c_{\times}}\sum_{t\ge0}
\exp\Bigl(-\frac{2}{L}\bigl(\log(X-2c_2+\lambda t)\bigr)^{r_{pr}}\Bigr),
\end{equation}
where $c_{\times}$ is the crossing-length constant. We require
\begin{equation}\label{9.1:eq-resolved-members-a}
\mathfrak{S}(\gamma^{-2})\le1 .
\end{equation}
Then
\begin{equation}\label{9.1:eq-resolved-members-8}
\mathfrak{S}:=\sup_n\sum_{\ell\ge n+2}\mathsf{w}_{\ell,n}\le\mathfrak{S}(X);
\end{equation}
$\mathfrak{S}(X)$ is finite for $r_{pr}>1$, is decreasing in $X=g^{-2}$ and tends to $0$ as
$X\to\infty$; the bound \eqref{9.1:eq-resolved-members-8} is uniform in the length $\mathfrak{n}$ of
the run, in the zone $n$ and in the number of steps; and under \eqref{9.1:eq-resolved-members-a},
for $g\le\gamma$,
\begin{equation*}
\sum_{\ell\ge n}\mathsf{w}_{\ell,n}\le3 .
\end{equation*}
\end{enumerate}
\end{definition}

\begin{proof}[Proof of the assertions contained in items (a), (b) and (e)]
\emph{The identity \eqref{9.1:eq-resolved-members-2}.} We compare the two sides of
\eqref{9.1:eq-resolved-members-2} slot by slot.

Let first $u\in U^{dir}\cup U^{\mathsf{z}}$. Each direct slot carries exactly one class, by the
assignment of classes to the inputs of the productions. For the $z$-free slots the same holds
because of the zone split of Lemma~\ref{9.1:lem-zone-split}(i). By
\eqref{9.1:eq-resolved-members-1} the $u$-entry of $W^{(z)}$ is
therefore $x^{\mathsf{B}}_u-x^{\mathsf{A}}_u$ for the one class $z$ of $u$, and $0$ for every other
class. Summing over $z$ gives $x^{\mathsf{B}}_u-x^{\mathsf{A}}_u$.

Let now $u\in U^{lf}$ be a read term of level $\ell\le i$. Item~(c) assigns to each of
$G_u^{(n)}$, $G_u^{rd}$ and $G_u^{own}$ exactly one class. Hence the sum over $z$ of the $u$-entries
of $W^{(z)}$ is
\begin{equation*}
\sum_n G_u^{(n)}+G_u^{rd}\ \bigl[+G_u^{own}\bigr].
\end{equation*}
The value of the slot in run $\mathsf{m}$ is
$x^{\mathsf{m}}_u=\mathbf{B}^{(\ell),\mathsf{m}}(X_u;\mathrm{rd}^{\mathsf{m}}P^{\mathsf{m}})$. By the
definition of $G_u^{rd}$ we can therefore write
\begin{equation*}
\begin{split}
x^{\mathsf{B}}_u-x^{\mathsf{A}}_u
&=\bigl[\mathbf{B}^{(\ell),\mathsf{B}}(X_u;\mathrm{rd}^{\mathsf{B}}P^{\mathsf{B}})
-\mathbf{B}^{(\ell),\mathsf{B}}(X_u;\mathrm{rd}^{\mathsf{A}}P^{\mathsf{A}})\bigr]\\
&\quad+\bigl[\mathbf{B}^{(\ell),\mathsf{B}}-\mathbf{B}^{(\ell),\mathsf{A}}\bigr]
(X_u;\mathrm{rd}^{\mathsf{A}}P^{\mathsf{A}})\\
&=G_u^{rd}+\bigl[\mathbf{B}^{(\ell),\mathsf{B}}-\mathbf{B}^{(\ell),\mathsf{A}}\bigr]
(X_u;\mathrm{rd}^{\mathsf{A}}P^{\mathsf{A}}).
\end{split}
\end{equation*}
For R-born $u$ the same holds with $\mathbf{B}'$ in place of $\mathbf{B}$. The remaining equality,
\begin{equation*}
\sum_n G_u^{(n)}\ \bigl[+G_u^{own}\bigr]
=\bigl[\mathbf{B}^{(\ell),\mathsf{B}}-\mathbf{B}^{(\ell),\mathsf{A}}\bigr]
(X_u;\mathrm{rd}^{\mathsf{A}}P^{\mathsf{A}}),
\end{equation*}
is the summation identity for the marks of read terms stated below in this chapter, applied at the
level $\ell$ of $u$. Since $\ell\le i$, the identity \eqref{9.1:eq-resolved-members-2} for a term
produced at level $i+1$ uses that summation identity only at levels $\le i$. The argument is
therefore an induction in the level.

Combining the two cases, every entry of $\sum_zW^{(z)}$ equals the corresponding entry of
$x^{\mathsf{B}}-x^{\mathsf{A}}$, which is \eqref{9.1:eq-resolved-members-2}.

\emph{The decomposition of an R-born term.} Rearranging \eqref{9.1:eq-resolved-members-4} gives
\begin{equation*}
\mathbf{B}'^{(i),\mathsf{B}}-\mathbf{B}'^{(i),\mathsf{A}}
=\sum_{z\in\mathcal{Z}_{surv}}\mathfrak{V}^{(z)}+\delta\mathbf{B}'^{(i),own}.
\end{equation*}
This holds identically, whatever the classes in $\mathcal{Z}_{own}$ contribute. Since each
$z\in\mathcal{Z}_{surv}$ is the unique class of its zone $n$ in $\mathcal{Z}_{surv}$, the sum over
$\mathcal{Z}_{surv}$ can be written as the sum of $\delta\mathbf{B}'^{(i),(n)}$ over those
$n\le i-1$ whose class lies in $\mathcal{Z}_{surv}$.

\emph{The weight sum.} Fix $n$ and a level $\ell\ge n+2$, and let $t\ge0$ be the depth of the level
$\ell-1$. Two lower bounds enter.
\begin{itemize}
\item By the reference window estimate $x_i-x_j\ge\lambda(j-i)-2c_2$ used with the correlation
theorem, the inverse coupling at the level $\ell-1$ satisfies
\begin{equation*}
x_{\ell-1}\ge X+\lambda t-2c_2 .
\end{equation*}
\item By the definition of the scaffold radius $R(x)$ used with the correlation theorem,
\begin{equation*}
\check{R}_{\ell-1}\ge(\log x_{\ell-1})^{r_{pr}} .
\end{equation*}
\end{itemize}
Inserting these two lower bounds into $\mathsf{w}_{\ell,n}=e^{-\mathsf{d}_{\ell-1}}$ gives
\begin{equation*}
\mathsf{w}_{\ell,n}\le e^{c_{\times}}
\exp\Bigl(-\frac{2}{L}\bigl(\log(X-2c_2+\lambda t)\bigr)^{r_{pr}}\Bigr).
\end{equation*}
For $X-2c_2>1$ the right-hand side is a decreasing function of $X+\lambda t$, because
$y\mapsto(\log y)^{r_{pr}}$ is increasing for $y>1$. Distinct levels $\ell$ have distinct depths
$t$, and all terms are non-negative. Summing over $\ell\ge n+2$ and taking the supremum over $n$
therefore gives
\begin{equation*}
\sup_n\sum_{\ell\ge n+2}\mathsf{w}_{\ell,n}\le\mathfrak{S}(X),
\end{equation*}
which is \eqref{9.1:eq-resolved-members-8}.

We now establish four properties of $\mathfrak{S}(X)$.

\emph{Finiteness.} Let $r_{pr}>1$; this holds because $r_{pr}\ge2$ by the condition fixing $r_{pr}$
in the order of choice. Then, for all sufficiently large $t$,
\begin{equation*}
\frac{2}{L}\bigl(\log(X-2c_2+\lambda t)\bigr)^{r_{pr}}
\ge\frac{2}{L}\bigl(\log(\lambda t)\bigr)^{r_{pr}}\ge2\log t ,
\end{equation*}
since $(\log(\lambda t))^{r_{pr}}/\log t\to\infty$ as $t\to\infty$. So the terms of
\eqref{9.1:eq-resolved-members-7} are eventually bounded by $t^{-2}$, and $\mathfrak{S}(X)$ is
finite.

\emph{Monotonicity.} Each term of \eqref{9.1:eq-resolved-members-7} is decreasing in $X$, so
$\mathfrak{S}(X)$ is decreasing in $X=g^{-2}$.

\emph{Limit.} Fix $X_0$ with $X_0-2c_2>1$. For $X\ge X_0$ the terms are dominated by those of
$\mathfrak{S}(X_0)$, and each term tends to $0$ as $X\to\infty$. By dominated convergence,
$\mathfrak{S}(X)\to0$ as $X\to\infty$. In particular \eqref{9.1:eq-resolved-members-a} is a
ceiling on $\gamma$ depending only on $L$, $\lambda$, $c_2$, $c_{\times}$ and $r_{pr}$, which is met
by every sufficiently small $\gamma$.

\emph{Uniformity.} The right-hand side of \eqref{9.1:eq-resolved-members-8} depends only on $X$,
$L$, $\lambda$, $c_2$, $c_{\times}$ and $r_{pr}$. The bound is therefore uniform in the length
$\mathfrak{n}$ of the run, in the zone $n$ and in the number of steps.

\emph{Conclusion.} Let $g\le\gamma$, so that $X=g^{-2}\ge\gamma^{-2}$. Since $\mathfrak{S}(X)$ is
decreasing, \eqref{9.1:eq-resolved-members-8} and \eqref{9.1:eq-resolved-members-a} give
$\mathfrak{S}\le\mathfrak{S}(\gamma^{-2})\le1$. With $\mathsf{w}_{n,n}=\mathsf{w}_{n+1,n}=1$ from
\eqref{9.1:eq-resolved-members-6} we obtain
\begin{equation*}
\sum_{\ell\ge n}\mathsf{w}_{\ell,n}=\mathsf{w}_{n,n}+\mathsf{w}_{n+1,n}
+\sum_{\ell\ge n+2}\mathsf{w}_{\ell,n}\le1+1+1=3 .
\end{equation*}
\end{proof}

\begin{lemma}[Per-cell domination of the resolved boundary members]
\label{9.1:lem-per-cell-domination}
Let $k'$ be a reading level and consider an aligned label, with its regions
$\Omega_1\supset\Lambda_1\supset\Omega_2\supset\Lambda_2\supset\cdots$, its zones
(zone $n$ is $\Omega_n\setminus\Omega_{n+1}$ for $n<k'$, and the top zone is $\Omega_{k'}$) and the
large-field regions $Z_j=\Lambda_j^c$, as in Notation~\ref{9.1:not-zones-boundary-terms}. Let
$(n,k')$ be a class, and let $\Gamma^0_n$ be the set of cells of zone $n$, that is, of the level-$n$
cubes of zone $n$.

For $\ell\ge n$ and $X\in\mathbf{D}_\ell$ let $\delta\mathbf{B}^{(\ell),(n)}(X)$ and
$\delta\mathbf{B}'^{(\ell),(n)}(X)$ be the class-$(n,k')$ members of the stored boundary terms of
level $\ell$ produced by a $\mathbf{T}$-step and by an $\mathbf{R}$-operation respectively, as
resolved in Definition~\ref{9.1:def-resolved-members}. Let $d_\ell(X)$ be the tree length of $X$,
$\kappa_{\mathbf{B}}$ the decay rate of the stored boundary terms, and let $\mathsf{w}_{\ell,n}$ and
$\mathsf{S}^{+,(n)}_\ell[\cdot]$ be the crossing weights and the weighted sup-norms of the class-$n$
differences between the two runs $\mathsf{A}$ and $\mathsf{B}$ at level $\ell$ that accompany the
zone-resolved bound \eqref{9.1:eq-resolved-members-a}; $\mathbf{B}'^{\mathrm{res}}$
denotes the resolved stored terms produced by an $\mathbf{R}$-operation, whose class-$n$ members are
the $\delta\mathbf{B}'^{(\ell),(n)}(X)$. Put
\begin{equation}\label{9.1:eq-per-cell-domination-1}
\mathsf{x}^{(n)}_\ell:=
\begin{cases}
\mathsf{w}_{\ell,n}\,\mathsf{S}^{+,(n)}_\ell[\mathbf{B}]
+\mathsf{S}^{+,(n)}_\ell[\mathbf{B}'^{\mathrm{res}}], & \ell\ge n,\\
0, & \ell<n.
\end{cases}
\end{equation}
The boundary load of class $(n,k')$ is
$\mathsf{y}_n^{(\le k')}=\sum_{\ell\ge n}\mathsf{x}^{(n)}_\ell\in[0,+\infty]$. Let $K(d)$ be
the resummation constant of the resummation lemma of the correlation theorem, and $c_d$ the constant
of that lemma for which, for every level-$\ell$ cube $\hat c$, the sum of $e^{-(c_d+1)d_\ell(X)}$ over
the $X\in\mathbf{D}_\ell$ containing $\hat c$ is at most $K(d)$. Assume that the scaffold radii
$\check{R}_j$ of consecutive levels satisfy $\check{R}_\ell\ge\check{R}_{\ell-1}/L$ for every $\ell$.
Then the following hold.
\begin{itemize}
\item[(a)] \emph{Values.} For every cell $c\in\Gamma^0_n$,
\begin{equation}\label{9.1:eq-per-cell-domination-2}
\sum_{\ell\ge n}\ \sum_{\substack{X\in\mathbf{D}_\ell\\ X\supset c}}
e^{(\kappa_{\mathbf{B}}-c_d-2)d_\ell(X)}
\Bigl(\bigl|\delta\mathbf{B}^{(\ell),(n)}(X)\bigr|
+\bigl|\delta\mathbf{B}'^{(\ell),(n)}(X)\bigr|\Bigr)
\le \bigl(K(d)+K_{\circ}(d)+K^{\circ}(d)\bigr)\,\mathsf{y}_n^{(\le k')} .
\end{equation}
Every non-zero member has such a cell $c\subset X$, with one exception. Let $X\subset\Omega_\ell$
carry a term produced by a $\mathbf{T}$-step, and consider its member of class $n=\ell-1$. At
production this was the top member, the zone of class $n$ then being $\Omega_{\ell-1}$. At later
reading levels $X$ contains no cell of $\Gamma^0_{\ell-1}=\Omega_{\ell-1}\setminus\Omega_\ell$.
Such a member is anchored at a level-$\ell$ cube of $X\cap Z_\ell$, which is a cell of
$\Gamma^0_\ell$. Consequently, for every region $Y$ over whose cells the activity sums of (b) are
taken, the values of the members are charged on the dial domain $\mathbb{D}^{\star}$ at most
\begin{equation*}
|\varepsilon|\bigl(K(d)+K_{\circ}(d)+K^{\circ}(d)\bigr)\,\mathsf{y}_n^{(\le k')}
\end{equation*}
per cell of $(\Gamma^0_n\cup\Gamma^0_{n+1})\cap Y$,
where $\varepsilon$ is the dial parameter. Since the per-block activity bound of hypothesis~(i)
of~(b) charges per unit of $\sum_j|\Gamma^0_j\cap Y|$, the phrase ``per cell of
$\Gamma^0_n\cap Y$'' is read as ``per cell of $Y$ met by $X$''.
\item[(b)] \emph{Activities.} Consider the per-cell and per-block activity sums for the flat marks of
stored boundary terms of sizes $\mathsf{x}_\ell e^{-\kappa_{\mathbf{B}}d_\ell(X)}$ that are bounded by
\begin{equation}\label{9.1:eq-per-cell-domination-3}
K_{site}\Bigl[\sum_{\ell\le i}w^{\mathbf{B}}_{i,\ell}\,\mathsf{x}_\ell
+\sum_{\ell>i}\varpi_{i,\ell}\,\mathsf{x}_\ell\Bigr],
\end{equation}
where $i$ is the zone of the block and $\varpi_{i,\ell}\le1$ is the crossing factor that the bound
in question attaches to level $\ell$. Such bounds are supplied by the following statements, which are
assumed as hypotheses:
\begin{itemize}
\item[(i)] the summation of boundary terms at finer zones, together with the resummation lemma of
the correlation theorem, by which the sum over $X$ of the marks is the norm;
\item[(ii)] the activity bound for flat marks on aligned labels;
\item[(iii)] the rate statement for these activity sums, with rate exponent $\sigma_*$ of the sums,
$\sigma_*\in\{\beta\kappa,\kappa/(4c^{\sharp})\}$, and floor $\sigma_*-1\ge c_d+1$;
\item[(iv)] part~(b) of the bound for the weights $w^{\mathbf{B}}_{i,\ell}$.
\end{itemize}
Assume moreover the floor
\begin{equation}\label{9.1:eq-per-cell-domination-a}
\sigma_*\ge c_d+3
\end{equation}
on $\kappa$. Then every such bound holds for the class-$(n,k')$ members, with $\mathsf{x}_\ell$
replaced by $\mathsf{x}^{(n)}_\ell$ of \eqref{9.1:eq-per-cell-domination-1} and the sums run at the
rate $\kappa_{\mathbf{B}}-1$. Each such sum is therefore at most $K_{site}\,\mathsf{y}_n^{(\le k')}$.
We put
\begin{equation*}
K_Y:=\bigl(K(d)+K_{\circ}(d)+K^{\circ}(d)\bigr)\cdot\max K_{site},
\end{equation*}
the maximum being taken over the constants $K_{site}$ of the
bounds supplied by hypotheses (i)--(iv).
\item[(c)] \emph{Reading parts.} Consider the reading parts $G_u^{rd}$ of the stored derived terms of
coarser levels $\ell'$, $n<\ell'\le k'$. Definition~\ref{9.1:def-resolved-members} assigns them to
class $(n,k')$, and the zone-by-zone two-run bound \eqref{9.1:eq-fluctuation-boundary-bound-a}
counts them among the entries of $[\![\delta]\!]^c$. These parts, each summed over $\ell'$ with the
crossing weight $\mathsf{w}_{\ell',n}$ of the term of level $\ell'$ that carries it (its support
being that term's $X$), satisfy, for every cell $c\in\Gamma^0_n$,
\begin{equation}\label{9.1:eq-per-cell-domination-4}
\begin{split}
\sum_{n<\ell'\le k'}\ \sum_{\substack{X\in\mathbf{D}_{\ell'}\\ X\supset c}}
e^{(\kappa_{\mathbf{B}}-c_d-2)d_{\ell'}(X)}\,\bigl|G_u^{rd}(X)\bigr|
&\le 3\bigl(K(d)+K_{\circ}(d)+K^{\circ}(d)\bigr)\,C B_0\,\vartheta^{n/4}\\
&\le \frac{3\bigl(K(d)+K_{\circ}(d)+K^{\circ}(d)\bigr)\,C B_0}{E_0\,C^{\natural}_{cl}}\,
[\![\delta]\!]^c_n ,
\end{split}
\end{equation}
where $C$, $B_0$, $E_0$, $C^{\natural}_{cl}$ and the rate $\vartheta$ are the constants of the
zone-resolved bound \eqref{9.1:eq-resolved-members-a}.
For the stored terms produced by an $\mathbf{R}$-operation, the corresponding sum is controlled
through $\sum_{\ell'}c_1^{\flat}(g_{\ell'})$. No age $k'-n$ enters.
\end{itemize}
Here $K_{\circ}(d)$ and $K^{\circ}(d)$ are the resummation constants of the island-blind format,
whose terms are charged to the cells of their own island $I$.
\end{lemma}

\begin{proof}
\emph{Support.} Let $\delta\mathbf{B}^{(\ell),(n)}(X)$ or $\delta\mathbf{B}'^{(\ell),(n)}(X)$ be a
non-zero member. The zone-by-zone
two-run bound for boundary terms of the fluctuation step,
\eqref{9.1:eq-fluctuation-boundary-bound-a}, in its support part and applied inductively in the
level, gives $X\cap(\text{zone }n)\ne\emptyset$, with zone $n$ as it stood at the production of the
term. At reading level $k'$ this yields a cell of $\Gamma^0_n$ contained in $X$, except in the case
treated under (a).

\emph{Crossing.} Let $X\in\mathbf{D}_\ell$ carry a term produced by a $\mathbf{T}$-step. By
\cite[(2.41)]{B-CMP119}, $X\cap\Omega_\ell\ne\emptyset$. Suppose $n\le\ell-2$. The set $X$, being a
member of $\mathbf{D}_\ell$, is a connected union of cubes, and by \cite[(2.42)]{B-CMP119} its
$d_\ell(X)$ is a tree length.
Hence $X$ contains a path from a point of zone $n$ to a point of $\Omega_\ell$.

Since $n+1\le\ell-1$ and the regions decrease, zone $n$ is contained in
$\Omega_{n+1}^c\subset\Omega_{\ell-1}^c$.
On the other side, $\Omega_\ell\subset\Lambda_{\ell-1}\subset\Omega_{\ell-1}$. Consider the part of
the path between its last entrance into $\Omega_{\ell-1}$ and its arrival in $\Lambda_{\ell-1}$. This
part crosses $\Omega_{\ell-1}\setminus\Lambda_{\ell-1}$, which is contained in zone $\ell-1$ because
$\Omega_\ell\subset\Lambda_{\ell-1}$. By the thickness property of the zones, which is available for
the zones determined by a renormalization transformation, zone $\ell-1$ has
thickness at least $2M\check{R}_{\ell-1}$ at scale $\ell-1$. Therefore
\begin{equation}\label{9.1:eq-per-cell-domination-5}
d_\ell(X)\ge\mathsf{d}_{\ell-1},
\end{equation}
where $\mathsf{d}_{\ell-1}$ is the crossing length of zone $\ell-1$ of
Notation~\ref{9.1:not-zones-boundary-terms}.

A second derivation of \eqref{9.1:eq-per-cell-domination-5} does not use the thickness property; it
is needed for the following reason. The domains made by an $\mathbf{R}$-operation are excluded from
that property: Ba{\l}aban's paper \cite{B-CMP119} speaks of ``The domains $\Omega_j$, $\Lambda_j$,
which are determined by the $j$-th renormalization transformation, but not by the R-operation'', and
it also states ``we admit the
possibility that $\Lambda_j=\Omega_j$ for some indices $j$. Such a situation may arise as a result
of an R-operation''.

Accordingly, suppose zone $\ell-1$ is a zone of the history of the label, that is, a zone whose
domains $\Omega_{\ell-1}$, $\Lambda_{\ell-1}$ are made by an $\mathbf{R}$-operation rather than by
the $(\ell-1)$-th renormalization transformation, lying inside a component of
the second class or an arrested component. Then $\Lambda_{\ell-1}=\Omega_{\ell-1}$ is possible, and
the set $\Omega_{\ell-1}\setminus\Lambda_{\ell-1}$ used above may be empty. In that case the crossing
is taken from the construction
of the $(\ell-1)$-th renormalization transformation itself. Every class $n\le\ell-2$, whether a zone
or a zone of the history, lies in $Z_{\ell-1}=\Lambda_{\ell-1}^c$. Let $Z'_{\ell-1}$ be the union of
the cubes of size $LM\check{R}_\ell$ meeting $Z_{\ell-1}$. Then $\Omega_\ell$ is cut out of the
complement of $Z'^{\sim4}_{\ell-1}$, which is $Z'_{\ell-1}$ enlarged by four layers of such cubes.
The construction \cite[(3.5)]{B-CMP119} adds one, one and two further layers, eight layers of
cubes of size $LM\check{R}_\ell$ in all. These are made by the $\mathbf{T}$-step at every step.

Hence a term produced by a $\mathbf{T}$-step on $X\in\mathbf{D}_\ell$ with $X\cap\Omega_\ell\ne\emptyset$
and a class-$n$ member satisfies
\begin{equation*}
d_\ell(X)\ \ge\ 5\check{R}_\ell-c_\times\ \ge\ 5\check{R}_{\ell-1}/L-c_\times\ \ge\ \mathsf{d}_{\ell-1},
\end{equation*}
where $c_\times$ is a fixed correction constant of this count of layers.
The middle inequality is the hypothesis $\check{R}_\ell\ge\check{R}_{\ell-1}/L$ of the lemma on
consecutive scaffold radii. The crossing weights and their sum
$\mathfrak{S}$ are unchanged, and \eqref{9.1:eq-resolved-members-a} holds,
whichever of the two derivations of \eqref{9.1:eq-per-cell-domination-5} is used. Starting
the path instead at a point of $X\cap\Lambda_\ell^c$ gives the same conclusion, with the crossing
length $\mathsf{d}_{\ell-1}$ taken as the smaller of the two thicknesses of zone $\ell-1$ obtained in
the two derivations above.

\emph{The crossing weight.} Let $X$ carry a non-zero class-$n$ member of a term produced by a
$\mathbf{T}$-step. We split off one unit of the rate,
$e^{-\kappa_{\mathbf{B}}d_\ell(X)}=e^{-d_\ell(X)}e^{-(\kappa_{\mathbf{B}}-1)d_\ell(X)}$.
\begin{itemize}
\item For $\ell\le n+1$ we have $\mathsf{w}_{\ell,n}=1$, and $e^{-d_\ell(X)}\le1$ because
$d_\ell(X)\ge0$.
\item For $\ell\ge n+2$, \eqref{9.1:eq-per-cell-domination-5} gives
$e^{-d_\ell(X)}\le e^{-\mathsf{d}_{\ell-1}}$, and $e^{-\mathsf{d}_{\ell-1}}\le\mathsf{w}_{\ell,n}$ by
the definition of the crossing weights that accompany the zone-resolved bound
\eqref{9.1:eq-resolved-members-a}.
\end{itemize}
In both cases
\begin{equation}\label{9.1:eq-per-cell-domination-6}
e^{-\kappa_{\mathbf{B}}d_\ell(X)}\le\mathsf{w}_{\ell,n}\,e^{-(\kappa_{\mathbf{B}}-1)d_\ell(X)} .
\end{equation}

\emph{Proof of (a).} Let $\ell\ge n$ and $c\subset X$. Since $c$ is a level-$n$ cube and $\ell\ge n$,
$X$ contains the unique level-$\ell$ cube $\hat c\supseteq c$. The sum over $X\supset c$ is therefore
at most the sum over $X\supset\hat c$. The resummation constant of the resummation lemma of the
correlation theorem, which rests on the tree count \cite[(2.29)]{B-CMP116}, gives
\begin{equation}\label{9.1:eq-per-cell-domination-7}
\sum_{\substack{X\in\mathbf{D}_\ell\\ X\supset\hat c}}e^{-(c_d+1)d_\ell(X)}\le K(d).
\end{equation}

The class-$n$ members of level $\ell$ on $X$ are bounded by their weighted sup-norms at the rate
$\kappa_{\mathbf{B}}$:
\begin{equation*}
\bigl|\delta\mathbf{B}^{(\ell),(n)}(X)\bigr|
\le\mathsf{S}^{+,(n)}_\ell[\mathbf{B}]\,e^{-\kappa_{\mathbf{B}}d_\ell(X)},
\qquad
\bigl|\delta\mathbf{B}'^{(\ell),(n)}(X)\bigr|
\le\mathsf{S}^{+,(n)}_\ell[\mathbf{B}'^{\mathrm{res}}]\,e^{-\kappa_{\mathbf{B}}d_\ell(X)} .\end{equation*}
For the first member we apply \eqref{9.1:eq-per-cell-domination-6}. For the second we use
$e^{-(c_d+2)d_\ell(X)}\le e^{-(c_d+1)d_\ell(X)}$. Together these give
\begin{equation*}
e^{(\kappa_{\mathbf{B}}-c_d-2)d_\ell(X)}
\Bigl(\bigl|\delta\mathbf{B}^{(\ell),(n)}(X)\bigr|+\bigl|\delta\mathbf{B}'^{(\ell),(n)}(X)\bigr|\Bigr)
\le\mathsf{x}^{(n)}_\ell\,e^{-(c_d+1)d_\ell(X)} .
\end{equation*}
Summing over $X\supset\hat c$ and using \eqref{9.1:eq-per-cell-domination-7}, the $\ell$-th summand
of \eqref{9.1:eq-per-cell-domination-2} is at most
\begin{equation*}
K(d)\bigl[\mathsf{w}_{\ell,n}\mathsf{S}^{+,(n)}_\ell[\mathbf{B}]
+\mathsf{S}^{+,(n)}_\ell[\mathbf{B}'^{\mathrm{res}}]\bigr]=K(d)\,\mathsf{x}^{(n)}_\ell .
\end{equation*}
Summing over $\ell\ge n$ gives \eqref{9.1:eq-per-cell-domination-2} with the constant $K(d)$ for the
terms outside the island-blind format; the terms of that format are treated below.

By the support step, every non-zero member has a cell of zone $n$ in $X$, with the exception named
in (a). For that exception, let $X\subset\Omega_\ell$ carry a term produced by a $\mathbf{T}$-step
with member of class $n=\ell-1$. By the construction of $\mathbf{B}^{(\ell)}$,
$X\not\subset\Lambda_\ell$ and $X\cap Z_\ell$ contains a level-$\ell$ cube; the same anchoring is
used in \cite{B-CMP119}, which applies ``this inequality with $|\Gamma_n\setminus\Lambda_k|$ for
$j=n$'' (the letters $\Gamma_n$, $\Lambda_k$ and $j$ being Ba{\l}aban's there). Moreover
$X\cap Z_\ell\subset\Omega_\ell\setminus\Lambda_\ell\subset\Omega_\ell\setminus\Omega_{\ell+1}$, so
this level-$\ell$ cube of $X\cap Z_\ell$ is a cell of $\Gamma^0_\ell=\Gamma^0_{n+1}$. The
estimate of the $\ell$-th summand just given applies with $\hat c$ this anchor cube.

Each non-zero member is therefore charged to a cell of $(\Gamma^0_n\cup\Gamma^0_{n+1})$ contained in
$X$, at most $\bigl(K(d)+K_{\circ}(d)+K^{\circ}(d)\bigr)\mathsf{y}_n^{(\le k')}$ per cell, the terms
of the island-blind format being included as shown below. On the dial domain the values enter with
the factor $\varepsilon$, which gives the charge
$|\varepsilon|\bigl(K(d)+K_{\circ}(d)+K^{\circ}(d)\bigr)\mathsf{y}_n^{(\le k')}$ per cell of
$(\Gamma^0_n\cup\Gamma^0_{n+1})\cap Y$. Since the per-block activity bound of hypothesis~(i) of~(b)
charges per unit of $\sum_j|\Gamma^0_j\cap Y|$, the cells of $Y$ met by $X$ carry this charge.

\emph{Proof of (c).} The reading parts $G_u^{rd}$ of the stored derived terms of levels
$\ell'\in\,]n,k']$ have the supports of their terms. The crossing argument reads only the support
$X$, so it gives each of these parts the crossing weight $\mathsf{w}_{\ell',n}$ of the term of level
$\ell'$ that carries it. Summed over $\ell'>n$, they obey the chain
\eqref{9.1:eq-per-cell-domination-4} by \eqref{9.1:eq-resolved-members-a}. For the stored terms
produced by an $\mathbf{R}$-operation the sum is controlled through
$\sum_{\ell'}c_1^{\flat}(g_{\ell'})$. Neither bound involves the age $k'-n$.

\emph{Island-blind terms.} Each term of the island-blind format is charged to the cells of its own
island $I$. The number $N_\ell(I)$ of level-$\ell$ cubes of $I$ satisfies
$N_\ell(I)\le\sum_j|\Gamma^0_j\cap I|$ by part~(a) of the island-blind counterpart of the present
lemma. The resummation constants $K_{\circ}(d)$ and $K^{\circ}(d)$ of that format then add to $K(d)$
in every bound above, which gives the constant $K(d)+K_{\circ}(d)+K^{\circ}(d)$ of (a), of (c) and
of $K_Y$.

\emph{Proof of (b).} The sums of hypotheses (i)--(iv) read of a stored term only three things: its
support, its rate and its size constant. The class-$n$ member has the same support as the term. By
the rate corollaries for the cluster expansion and for the first variation, it has the rate
$\kappa_{\mathbf{B}}$ of its production. One unit of that rate is converted into
$\mathsf{w}_{\ell,n}$ by \eqref{9.1:eq-per-cell-domination-6}. The member is thus a flat mark of size
$\mathsf{x}^{(n)}_\ell e^{-(\kappa_{\mathbf{B}}-1)d_\ell(X)}$, and the sums run at the rate
$\kappa_{\mathbf{B}}-1$.

With the sums run at the rate $\kappa_{\mathbf{B}}-1$, the floor $\sigma_*-1\ge c_d+1$ of
hypothesis~(iii) is replaced by $\sigma_*\ge c_d+3$, which is the floor
\eqref{9.1:eq-per-cell-domination-a} assumed in~(b).
Because $\sigma_*\in\{\beta\kappa,\kappa/(4c^{\sharp})\}$ by hypothesis~(iii),
\eqref{9.1:eq-per-cell-domination-a} is a lower bound on $\kappa$ depending on $d$ and $\beta$ only.

With $\mathsf{x}^{(n)}$ in place of $\mathsf{x}$, the right member \eqref{9.1:eq-per-cell-domination-3}
is bounded using $w^{\mathbf{B}}_{i,\ell}\le1$ and $\varpi_{i,\ell}\le1$:
\begin{equation*}
K_{site}\Bigl[\sum_{\ell\le i}w^{\mathbf{B}}_{i,\ell}\,\mathsf{x}^{(n)}_\ell
+\sum_{\ell>i}\varpi_{i,\ell}\,\mathsf{x}^{(n)}_\ell\Bigr]
\le K_{site}\sum_{\ell\ge n}\mathsf{x}^{(n)}_\ell=K_{site}\,\mathsf{y}_n^{(\le k')} .
\end{equation*}
Here $\mathsf{x}^{(n)}_\ell=0$ for $\ell<n$.

The quantity $\mathsf{y}_n^{(\le k')}$ may be $+\infty$. The lemma is an inequality in $[0,+\infty]$,
and the finiteness of $\mathsf{y}_n^{(\le k')}$ is a separate matter, proved where the boundary loads
are shown to be finite.
\end{proof}

\begin{lemma}[Interpolation with resolved boundary terms and the age of live components]
\label{9.1:lem-resolved-interpolation}
Let $\mathsf{A}$ and $\mathsf{B}$ be the two runs compared in this chapter, and let $\lambda$ and
$\lambda_z$, with $z=(n,k')$ running over the classes, be the parameters of the dial domain
$\mathbb{D}^{\star}$. We write $\varepsilon=\lambda_z-\lambda$ and $Z_j=\Lambda_j^c$. Let $X$ be the
domain of the output term, $\mathbf{B}^{(\ell)}(X)$ or $\mathbf{B}'^{(\ell)}(X)$ as in
Definition~\ref{9.1:def-resolved-members}, and let $X^{\sim}$ be its enlargement.

Assume the following.
\begin{itemize}
\item[(A)] The hypotheses of the dial interpolation proposition are in force, so that its
conclusions hold on $\mathbb{D}^{\star}$. Among these hypotheses are the natural-zone interpolation
proposition, from which the dial interpolation proposition is derived, and the bounds on those items
of the assembly that read only sizes and holomorphy of derived objects, in the form in which the
dial interpolation proposition uses them.
\item[(B)] The activity condition of the dial domain holds at total size $2E_1^z$ with constant
$C_L'''$. In what follows this is called \emph{the activity condition}.
\item[(C)] The lemma on geometric forgetting of old regions is available; it is stated for every
component of any history, and it applies to a component whose age satisfies its age hypothesis.
\item[(D)] The zone split of the two-run differences at the fluctuation step
(Lemma~\ref{9.1:lem-zone-split}) holds, together with the display
\eqref{9.1:eq-zone-split-a}.
\end{itemize}

Then the following hold at the $\mathbf{T}$-site producing the stored terms $\mathbf{B}$, at the two
further sites of the assembly at which the dial interpolation proposition is applied, and at the top
storey of the assembly. In each case the arrays are those of parts (a) and (c) of
Definition~\ref{9.1:def-resolved-members}.
\begin{itemize}
\item[(i)] The ends are exact.
\item[(ii)] Let $E_0$ be the constant in the lower bound on the classical loads of aligned zones
used in the proof below, and let
\begin{equation}\label{9.1:eq-resolved-interpolation-1}
C_{\mathbf{B}}^v:=\frac{2\,(K(d)+K_{\circ}+K^{\circ})\,E_1^z}{E_0}.
\end{equation}
On $\mathbb{D}^{\star}$ every derived object is holomorphic in $\varepsilon$ and continuous in
$\lambda$. It obeys the bounds of the sharp interpolation proposition with the class constant
$\mathsf{c}$ replaced by
\begin{equation*}
\mathsf{c}+C_{SM}+2+C_z+C_{\mathbf{B}}^v ,
\end{equation*}
where $C_{SM}$ and $C_z$ are the constants of the dial interpolation proposition. The conclusion of
geometric forgetting of old regions holds there for the derived objects.
\item[(iii)] Let $z=(n,k')$. For a component $C$ write $j(C)$ for the step at which $C$ was created
(its first large-field index), so that its age at level $k'$ is $\operatorname{age}(C)=k'-j(C)$.
The parameter $\lambda_z$ occurs, that is, its coefficient is not identically zero, only in labels
of the following kind. Their history has zone $n$, at some level $k''\le k'$, inside a component
$C$ of $\Lambda_{k'}^c$ that is live at $k'$ and meets $X^{\sim}$; that is, $C$ is of one of the
three kinds:
\begin{itemize}
\item[--] $C$ is untouched;
\item[--] $C$ is second-class: for the region $Y$ of the operation $\mathbf{T}_i(Y)$ one has
$Y\subset Z_i^{\mathrm{new}}=Z_i\cup\bigcup_\iota Y_\iota$, with the factor
$\chi^c_k(Y^{\sim-6})$ of \cite[(1.102)]{B-CMP122b}, and the sequence $\omega_Y$ (the summation
index of $\mathbf{T}_i(Y)$) is kept in the label;
\item[--] $C$ is arrested.
\end{itemize}
Every such component satisfies
\begin{equation}\label{9.1:eq-resolved-interpolation-a}
j(C)\le n+1,\qquad \operatorname{age}(C)=k'-j(C)\ge k'-n-1 .
\end{equation}
Suppose a component $C$ containing the large-field component in which the slot sits is healed at
some $m\le k'$. Then the stored term is a term of the first part of the $\mathbf{R}$-operation and
has been absorbed into $R'^{(m)}$: in \cite[\S1]{B-CMP122b} the terms whose localisation domains
intersect the region $Z$ are absorbed. Hence the second inequality of
\eqref{9.1:eq-resolved-interpolation-a} holds at every reading level.

The age hypothesis of geometric forgetting of old regions is therefore satisfied, and that lemma
supplies, for every class $z$, the factor
\begin{equation}\label{9.1:eq-resolved-interpolation-2}
(k'-n+1)\,q_a^{(k'-n-1)_+},
\end{equation}
in the sequence form of the age bound for zone classes that accompanies the dial interpolation
proposition.
For slots of stored terms in the island-blind format the same holds with the relative lengths and
norms of that format. There the relative-format analogue of the present lemma takes its place,
since the relative sup-ball is convex.
\end{itemize}
\end{lemma}

\begin{proof}
\emph{Proof of} (i). The marks of the classes sum to the difference of the two arrays,
\begin{equation*}
\sum_z W^{(z)}=x^{\mathsf{B}}-x^{\mathsf{A}} .
\end{equation*}
The derived objects of both runs are produced by one assembly formula applied to the array of slot
values, on the common variables. This is item (c1) of the dial interpolation proposition, available
by hypothesis~(A). Hence the ends are exact.

\emph{Proof of} (ii). We follow the proof of the dial interpolation proposition item by
item. Only the slots listed below behave differently.

Slots of direct uses and slots carrying the entries $\mathsf{z}_i$ have one class each. Their total
dial is at most $2E_1^z$.

A slot $u$ of a resolved stored term $\mathbf{B}$ carries the value
\begin{equation*}
(1-\lambda)\,x_u^{\mathsf{A}}+\lambda\,x_u^{\mathsf{B}}+\varepsilon\,G_u^{(n)} ,
\end{equation*}
and it has one class. In each run the stored term obeys \cite[(2.42)]{B-CMP122b} if it is
$\mathbf{T}$-born and \cite[(1.99)]{B-CMP122b} if it is $\mathbf{R}$-born. The sup-ball in which
the correlation theorem reads stored derived terms is convex. Hence the convex combination
$(1-\lambda)x_u^{\mathsf{A}}+\lambda x_u^{\mathsf{B}}$ lies in that ball.

Consider the increment $\varepsilon G_u^{(n)}$. Part (b) of Lemma~\ref{9.1:lem-per-cell-domination}
gives the first inequality below, and $\mathsf{y}_n\le[\![\delta]\!]^{\natural+}_n$ gives the second:
\begin{equation*}
\text{activity per cell}\ \le\ |\varepsilon|\,K_Y\,\mathsf{y}_n
\ \le\ |\varepsilon|\,K_Y\,[\![\delta]\!]^{\natural+}_n\ \le\ K_Y\,E_1^z .
\end{equation*}
This lies inside the activity condition of hypothesis~(B), at the constant $C_L'''$.

Part (a) of Lemma~\ref{9.1:lem-per-cell-domination} gives, per cell of $\Gamma^0_n\cap Y$, with $Y$
the region of that lemma,
\begin{equation*}
\text{value per cell}\ \le\ |\varepsilon|\,K(d)\,\mathsf{y}_n\ \le\ K(d)\,E_1^z .
\end{equation*}
This lies inside the per-cell budget of the class constant $\mathsf{c}$, which item (iv) of the
natural-format bound for stored derived terms charges per unit of $|\Gamma^0_{n(z)}\cap Y|$.

At the $\mathbf{T}$-site, operator slots again carry a convex combination of the two runs' values.
It lies in the tolerance ball of the tolerance lemma of the $\mathbf{T}$-site. The class-$n$ pieces,
at $|\varepsilon|\le\Lambda_z$, move this combination by at most $|\varepsilon|\,\mathfrak{d}^{op}_n$.
By part (iii) of Lemma~\ref{9.1:lem-zone-split},
\begin{equation*}
\mathfrak{d}^{op}_n=C_{cl}^T(\log x_n)^{c_2}\vartheta_T^n .
\end{equation*}
On aligned zones we have
\begin{equation*}
[\![\delta]\!]^c_n\ \ge\ E_0\,C^{\natural}_{cl}\,(\log x_n)^{c_2}\,\vartheta_2^{n/4} .
\end{equation*}
By part (iv) of Lemma~\ref{9.1:lem-zone-split}, this lower bound and
\eqref{9.1:eq-zone-split-a},
\begin{align*}
|\varepsilon|\,\mathfrak{d}^{op}_n
&\ \le\ \frac{E_1^z}{[\![\delta]\!]^c_n}\,\mathfrak{d}^{op}_n
\ \le\ \frac{E_1^z}{E_0}\,\frac{C_{cl}^T}{C^{\natural}_{cl}}\,\vartheta_T^{\,n}\,\vartheta_2^{-n/4}\\
&\ \le\ \frac{E_1^z}{E_0}\,\frac{C_{cl}^T}{C^{\natural}_{cl}}\,\vartheta_\times^{\,n}
\ \le\ \tfrac14\,\mathfrak{m}_G^T .
\end{align*}
Here $\vartheta_\times$ is the classical index rate that dominates the quotient
$\vartheta_T\vartheta_2^{-1/4}$ of the two rates displayed above.
The tolerance lemma and the holomorphy proposition of the $\mathbf{T}$-site therefore give
holomorphy at these points, and the bounds of the summation lemma of the correlation theorem hold
there.

Among the class pieces at the $\mathbf{T}$-site, those that enter as potentials have activity
\begin{equation*}
|\varepsilon|\,\mathfrak{d}^{op}_n\,\bar{\varepsilon}^{pr}
\ \le\ \frac{E_1^z}{E_0\,C^{\natural}_{cl}}\,\bar{\varepsilon}^{pr} .
\end{equation*}
This is the factor admitted by the activity condition. The remaining class pieces enter additively
and need no such factor.

The remaining items of the proof of the dial interpolation proposition, and geometric forgetting
of old regions, read only sizes and holomorphy of derived objects. They therefore hold as in that
proof.

\emph{Proof of} (iii). A slot of class $(n,k')$ is of one of three kinds.
\begin{itemize}
\item[--] A direct use anchored in zone $n$. This is the definition of the classes of direct uses
in the sharp interpolation proposition.
\item[--] A class-$n$ piece free of $\mathsf{z}$, with support in $X$ meeting zone $n$. This is part
(ii) of Lemma~\ref{9.1:lem-zone-split}.
\item[--] A resolved member $G_u^{(n)}\neq0$ with $X_u\subset X$ meeting zone $n$. This is the
support property of the resolved members stated in
Lemma~\ref{9.1:lem-per-cell-domination}: every non-zero member has a cell of $\Gamma^0_n$ inside
its domain.
\end{itemize}
The members $G_u^{rd}$ are assigned to a zone met by $X_u$. The members $G_u^{own}$ are assigned to
the zone of the block.

Now
\begin{equation*}
\text{zone }n\ \subset\ \Omega_{n+1}^c\ \subset\ \Lambda_{n+1}^c=Z_{n+1}\ \subset\ Z_{k'} .
\end{equation*}
Let $C$ be the component of $Z_{k'}=\Lambda_{k'}^c$ containing the piece of zone $n$ that $X$
meets. Then $C$ contains points outside $\Omega_{n+1}$. So $C$ meets $Z_{n+1}$, and its first
large-field index satisfies $j(C)\le n+1$, as in item (d) of the lemma on geometric forgetting of
old regions.

Next, $C$ is live at $k'$. Suppose $\mathbf{R}_m$ had healed $C$ at some $m<k'$. Then every term of
the action whose domain meets $\tilde C\supseteq C$ would have been absorbed into $R'^{(m)}$: the
activities $F(X')$ of \cite[(1.90)--(1.91)]{B-CMP122b} collect every term linked to the domains of
the $\mathbf{T}'$-operations. But the slot is anchored in zone $n\subset C$, or
supported on a set meeting it, or is a member whose domain meets it. Such a slot would then not
exist at $k'$, a contradiction. Hence
\begin{equation*}
\operatorname{age}(C)=k'-j(C)\ge k'-n-1 ,
\end{equation*}
which is \eqref{9.1:eq-resolved-interpolation-a}. This is the hypothesis of geometric forgetting
of old regions.

That lemma is applicable here. It is stated for every component of any history; the step of its
proof that requires $Z_j\cap C\neq\emptyset$ for $j(C)\le j\le k'$ is satisfied throughout the
event considered; and \cite[(1.79)]{B-CMP122b} is stated for the $\mathbf{T}$-operation connected
with an arbitrary large-field region. It therefore supplies the factor
\eqref{9.1:eq-resolved-interpolation-2}.

Two kinds of class require a separate word.

\emph{The top member of a $\mathbf{T}$-born term.} Let $n=\ell-1$ be the top member of a
$\mathbf{T}$-born $\mathbf{B}^{(\ell)}$, whose slots lie in the then top region $\Omega_{\ell-1}$.
Here $X_u\cap Z_\ell\neq\emptyset$ by \cite[(2.41)]{B-CMP122b}. This gives directly a component $C$
with $j(C)\le\ell=n+1$.

\emph{The member $G_u^{own}$ of an $\mathbf{R}$-born term.} Let $G_u^{own}$ belong to a
$\mathbf{B}'^{(\ell)}$-term and be assigned to the class $(i_u\vee\ell,k')$.
\begin{itemize}
\item[--] If $i_u\ge\ell$, the block's own zone $i_u$ meets $X_u$, and the argument above applies.
\item[--] If $i_u<\ell$, then $X_u$ meets the old zone $i_u$. After $\mathbf{R}_\ell$ this zone
exists only inside a surviving component, because healed components have left the label. So a
component $C$ with $j(C)\le i_u+1\le\ell+1$ is at hand.
\end{itemize}
\end{proof}

\begin{theorem}[Zone-by-zone two-run bound for boundary terms of the fluctuation step]
\label{9.1:thm-fluctuation-boundary-bound}
Let $\mathsf{A}$ and $\mathsf{B}$ be the two runs being compared, fix an aligned label, and assume the following.
\begin{itemize}
\item[(h1)] The hypotheses of part~(a) of the two-run comparison theorem at the site of the fluctuation step. These are the one-formula property of the stored terms, the class hypothesis of that theorem, the coupling condition $g\le\gamma_R$, where $\gamma_R$ is the bound on the coupling required by that theorem, the interpolation construction together with its weight condition at its grade, and the tolerance condition for the operators at the site of the fluctuation step, as supplied by the tolerance lemma.
\item[(h2)] Lemma~\ref{9.1:lem-zone-split} (zone split of the two-run differences at the fluctuation step).
\item[(h3)] The zone-resolved members of Definition~\ref{9.1:def-resolved-members}, formed for the stores of levels $\le i$.
\item[(h4)] The smallness condition for the activities at the site of the fluctuation step, read at the constants $(2E_1^z,C_L''')$.
\item[(h5)] The display \eqref{9.1:eq-zone-split-a} of Lemma~\ref{9.1:lem-zone-split}.
\item[(h6)] The marked-potential corollary of the cluster expansion, and the first-variation formula, in the forms recalled in the proof below.
\end{itemize}
Consider the fluctuation step $\mathbf{T}$ from level $i$ to level $i+1$ of the aligned label, an output
$\mathbf{B}^{(i+1)}(X)$ as stored by this step, and a zone $n\le i$. Write
$\delta\mathbf{B}^{(i+1),(n)}(X;\mathsf{v})$ for the member of zone $n$ of the two-run difference of this output, in the
sense of Definition~\ref{9.1:def-resolved-members}. Then
\begin{equation}\label{9.1:eq-fluctuation-boundary-bound-a}
\begin{aligned}
&\text{(a)}\quad \sum_{n\le i}\delta\mathbf{B}^{(i+1),(n)}(X;\mathsf{v})
=\mathbf{B}^{(i+1),\mathsf{B}}(X;\mathsf{v})-\mathbf{B}^{(i+1),\mathsf{A}}(X;\mathsf{v}),\\
&\text{(b)}\quad \delta\mathbf{B}^{(i+1),(n)}(X)\equiv0\quad\text{unless $X$ meets zone $n$},\\
&\text{(c)}\quad \mathsf{S}^{+,(n)}_{i+1}[\mathbf{B}]\le
\varkappa^{\mathbf{B}}_i\,[\![\delta]\!]^{\natural+,(i)}_n+\sigma^{T,(n)} .
\end{aligned}
\end{equation}
Identity (a) holds at every datum $\mathsf{v}$ common to the two runs and at every local label. Each member is holomorphic in $\mathsf{v}$ on the tube of the term.

In (b), zone $n$ is understood at the production of the term. For the member $n=i$, the top zone at that production, it is the region $\Omega_i$. At a later reading level, the anchor of the member is a cube of $X\cap Z_{i+1}$, by clause (b) of Lemma~\ref{9.1:lem-per-cell-domination}. For history zones inside second-class or arrested components, (b) is read in the sequence form of \eqref{9.1:eq-resolved-interpolation-a}.

In (c), $\mathsf{S}^{+,(n)}_{i+1}[\mathbf{B}]$ is the weighted supremum norm of the two-run differences of zone $n$,
namely the supremum of
$e^{\kappa_{\mathbf{B}}d_{i+1}(X)}|\delta\mathbf{B}^{(i+1),(n)}(X;\mathsf{v})|$ over data $\mathsf{v}$, local labels and domains $X$. The coefficients are
\begin{equation}\label{9.1:eq-fluctuation-boundary-bound-1}\varkappa^{\mathbf{B}}_i:=(1+K_Y^{\circ})\varkappa_i=K_{CE}C_L'''\theta_i
\le\bar{\varkappa}:=(1+K_Y^{\circ})\sup_i\varkappa_i .
\end{equation}
Here $\varkappa_i$ is the coefficient of part~(a) of the two-run comparison theorem, proportional to $\theta_i$, and $K_Y^{\circ}$ is the constant of Lemma~\ref{9.1:lem-per-cell-domination} in the island-blind format. The constant $K_{CE}$ is that of the marked-potential corollary of the cluster expansion recalled in the proof. The constants $E_1^z$ and $C_L'''$ are those with which the smallness condition (h4) is applied: the activity bound is taken at $2E_1^z$, and $C_L'''$ is the weight constant of the per-cell budget, which is also the constant of the per-cell bound of the marked activity in the proof. The second coefficient is
\begin{equation}\label{9.1:eq-fluctuation-boundary-bound-2}
\sigma^{T,(n)}:=C_T^{\natural}E_0(\log x_n)^{c_2}\vartheta_T^n,
\qquad
C_T^{\natural}:=\frac{2E^{\sharp}_T(1+\bar{\varepsilon}^{pr}+C_F)C_{cl}^T}{E_0\,\mathfrak{m}_G^T}.
\end{equation}
This is the quantity $\sigma^T_k$ of part~(a) of the two-run comparison theorem, with the level $k$ of the output replaced by the zone $n$ of the class and nothing else changed. Here $E_0$ is the constant of part~(a) of the two-run comparison theorem that enters $\sigma^T_k$; it is also the constant of the lower bound \eqref{9.1:eq-zones-boundary-terms-a} on the classical load of aligned zones. The constant $C_F$ is the constant by which the normaliser moves in clause (iii) of Lemma~\ref{9.1:lem-zone-split}. On aligned zones,
\begin{equation}\label{9.1:eq-fluctuation-boundary-bound-3}
\sigma^{T,(n)}\le C_T^{\mathbf{B}}\vartheta_\times^n[\![\delta]\!]^c_n\le C_T^{\mathbf{B}}[\![\delta]\!]^c_n,
\qquad C_T^{\mathbf{B}}:=C_T^{\natural}/C^{\natural}_{cl}.
\end{equation}
The members carry the rate $\kappa_{out}\ge\kappa_{\mathbf{B}}$ of the production. For outputs that read a unit of the island-blind format, the decay factor is $e^{-\kappa_{out}d^{\circ}}$, by the relative exponential bound of the fluctuation step, as in the island-blind version of the present theorem.

The load $[\![\delta]\!]^{\natural+,(i)}_n$ is the total load of zone $n$ of Definition~\ref{9.1:def-resolved-members}(d). It is formed with the stores of levels $\le i$ for the terms $\mathbf{B}$ and of levels $<i$ for the terms $\mathbf{B}'$, since the production reads only stores produced before it.
\end{theorem}

\begin{proof}
\emph{Proof of} (a). By Definition~\ref{9.1:def-resolved-members}(a), the term $\mathbf{B}^{(i+1)}(X)$ is one formula $\mathcal{P}_X$ applied to the array $x^{\mathsf{c}}$ of slot values of its production, $\mathsf{c}\in\{\mathsf{A},\mathsf{B}\}$. By the same clause, the marks $W^{(z)}$ of the classes $z$ add up to the difference of the two arrays:
\begin{equation*}
\sum_z W^{(z)}=x^{\mathsf{B}}-x^{\mathsf{A}} .
\end{equation*}
Let $x^{\lambda}:=x^{\mathsf{A}}+\lambda(x^{\mathsf{B}}-x^{\mathsf{A}})$, $\lambda\in[0,1]$, be the convex combination of the two arrays. We apply the first-variation formula with $F:=\mathcal{P}_X$, $x^0:=x^{\mathsf{A}}$ and $x^1:=x^{\mathsf{B}}$. It states that
\begin{equation}\label{9.1:eq-fluctuation-boundary-bound-4}
F(x^1)-F(x^0)=\sum_z\mathfrak{V}^{(z)},\qquad
\mathfrak{V}^{(z)}:=\int_0^1\frac{d}{d\varepsilon}\Big|_{\varepsilon=0}
F\big(x^{\lambda}+\varepsilon W^{(z)}\big)\,d\lambda ,
\end{equation}
provided that two conditions hold.
\begin{itemize}
\item[(i)] At each $\lambda\in[0,1]$ there is joint holomorphy near $0$ in the finitely many interpolation variables $\varepsilon_z$ ($z$ ranging over the classes of the zones $n\le i$) and in the variable $\tau$.
\item[(ii)] There is holomorphy in $\varepsilon_z$, together with a bound, on $|\varepsilon_z|\le\Lambda_z$.
\end{itemize}
Both conditions follow from clause (c2)$^{\star}$ of Lemma~\ref{9.1:lem-resolved-interpolation}, applied at the site of the fluctuation step. That clause states that on $\mathbb{D}^{\star}$ every derived object is holomorphic in $\varepsilon=\lambda_z-\lambda$ and continuous in $\lambda$, with bounds. Since there are finitely many classes, the same majorant serves a neighbourhood of $0$ in all the $\varepsilon_z$ at once.

By clause (c1) of Lemma~\ref{9.1:lem-resolved-interpolation} the ends are exact, that is,
$\mathcal{P}_X[x^{\mathsf{c}}]=\mathbf{B}^{(i+1),\mathsf{c}}(X;\mathsf{v})$ for $\mathsf{c}\in\{\mathsf{A},\mathsf{B}\}$. The left side of \eqref{9.1:eq-fluctuation-boundary-bound-4} is therefore
$\mathbf{B}^{(i+1),\mathsf{B}}(X;\mathsf{v})-\mathbf{B}^{(i+1),\mathsf{A}}(X;\mathsf{v})$. Grouping the terms $\mathfrak{V}^{(z)}$ by the zone $n$ of $z$ gives the members $\delta\mathbf{B}^{(i+1),(n)}(X;\mathsf{v})$ and the identity (a).

Holomorphy in $\mathsf{v}$ follows from three facts.
\begin{itemize}
\item The entries $x^{\mathsf{c}}(\mathsf{v})$ and $W^{(z)}(\mathsf{v})$ are holomorphic on the data tube. For the stored units and their marks, this is the data-holomorphy statement that accompanies the two-run comparison theorem. For the members of stored derived terms, it is identity (a) at the levels $\le i$, used inductively. For the operators, holomorphy in the parameters is part of the tolerance lemma.
\item The function $F$ is holomorphic.
\item The operation $\int_0^1 d\lambda\,(d/d\varepsilon)_{\varepsilon=0}$ preserves holomorphy, because the majorant of clause (c2)$^{\star}$ of Lemma~\ref{9.1:lem-resolved-interpolation} is uniform in $\lambda$.
\end{itemize}

\emph{Proof of} (b). Suppose that $X$ does not meet zone $n$. We show that $W^{(z)}=0$ for every class $z=(n,\cdot)$ of zone $n$.
\begin{itemize}
\item The direct uses of class $(n,\cdot)$ are anchored in zone $n$, and only pieces inside $X$ are read. This is the localisation lemma of the correlation expansion: an output with domain $X$ reads only $V(Y)$ with $Y\subset X$. It is also the locality clause of the weighted-size bound of the interpolation. Hence no direct use of class $(n,\cdot)$ enters.
\item The class-$n$ pieces $\delta\mathsf{z}^{(n)}{\upharpoonright}X$ of the zone split of Lemma~\ref{9.1:lem-zone-split} vanish, by its clause (ii).
\item Let $X_u\subset X$ be the domain of a stored derived term read in $X$ whose member $G_u^{(n)}$ of zone $n$ is non-zero. Then $X_u$ meets zone $n$. This is statement (b) at the level $\ell\le i$ of $X_u$ for terms born at fluctuation steps. For terms born at $\mathbf{R}$-steps, it is the corresponding vanishing statement for those terms. The argument is an induction on the level. Since $X_u\subset X$ and $X$ does not meet zone $n$, no such member occurs.
\end{itemize}
Hence $W^{(z)}=0$ for $z=(n,\cdot)$, and by \eqref{9.1:eq-fluctuation-boundary-bound-4} $\mathfrak{V}^{(z)}=0$. Therefore $\delta\mathbf{B}^{(i+1),(n)}(X)\equiv0$.

\emph{Proof of} (c). Fix the class $z=(n,\cdot)$ and split its direction as
$W^{(z)}=W^{(z),mk}+W^{(z),op}$. The operator direction $W^{(z),op}$ consists of the class-$n$ pieces of Lemma~\ref{9.1:lem-zone-split}: the operator pieces, the potential pieces and the normaliser pieces. The mark direction $W^{(z),mk}$ is the rest of class $z$. It contains three kinds of entry:
\begin{itemize}
\item the direct marks, with both of their parts;
\item the members $G^{(n)}$ of the stored derived terms;
\item the reading parts $G^{rd}$ and own members $G^{own}$ that Definition~\ref{9.1:def-resolved-members}(c) assigns to this class.
\end{itemize}
By clause (c2)$^{\star}$ of Lemma~\ref{9.1:lem-resolved-interpolation}, $F$ is holomorphic near $x^{\lambda}$, so its derivative is linear in the direction:
\begin{equation*}
\begin{aligned}
\frac{d}{d\varepsilon}\Big|_{0}F\big(x^{\lambda}+\varepsilon W^{(z)}\big)
&=\frac{d}{d\varepsilon}\Big|_{0}F\big(x^{\lambda}+\varepsilon W^{(z),mk}\big)
+\frac{d}{d\varepsilon}\Big|_{0}F\big(x^{\lambda}+\varepsilon W^{(z),op}\big)\\
&=:\mathfrak{d}^{mk}(\lambda)+\mathfrak{d}^{op}(\lambda).
\end{aligned}
\end{equation*}
By \eqref{9.1:eq-fluctuation-boundary-bound-4}, therefore,
$\mathfrak{V}^{(z)}=\int_0^1(\mathfrak{d}^{mk}(\lambda)+\mathfrak{d}^{op}(\lambda))\,d\lambda$. We bound the two derivatives separately.

\smallskip
\emph{Step $(\alpha)$: the mark direction.} At the base point $x^{\lambda}$, the data of the fluctuation step are convex combinations of the two runs' unit values, values of stored derived terms, operators, potentials and normaliser. These data are admissible for the localisation lemma of the correlation expansion, with activity below its threshold. This is clause (c2)$^{\star}$ of Lemma~\ref{9.1:lem-resolved-interpolation} at $\varepsilon=0$, which rests on convexity and on the smallness condition (h4) read at $2E_1^z$.

The perturbation $\varepsilon W^{(z),mk}$ enters the piece bound of the correlation expansion linearly in the unit values and in the values of the stored derived terms. It therefore enters as a marked potential
$V^{m,z}:=V[W^{(z),mk}]$.

We now apply the marked-potential corollary of the cluster expansion, in the normalisation of the budget of the interpolation, in which the activity of a marked potential $V^m$ is denoted $\mathsf{W}^m$. It states: if the potential is $V=V^u+\varepsilon V^m$, with $V^u$ the unmarked part, then on the disc in $\varepsilon$ where the total activity stays below the activity threshold $\nu_0$ the marked sub-sum $t_1-t_0$ is at most $K_{CE}\mathsf{W}^me^{-\kappa d_{i+1}(X)}$, where $d_{i+1}(X)$ is the tree length at level $i+1$ of the domain $X$ of the output. Its proof applies Schwarz's lemma on that disc. Hence the derivative at $\varepsilon=0$ obeys
\begin{equation*}
\Big|\frac{d}{d\varepsilon}\Big|_{\varepsilon=0}(t_1-t_0)\Big|\le K_{CE}\mathsf{W}^me^{-\kappa d_{i+1}(X)} .
\end{equation*}
The corollary applies to the output that produces the boundary term, because the localisation lemma lists among its outputs the pinned kernels of \cite[(3.47)]{B-CMP119} and the kernels of $\mathbf{B}$ and of $R''$. These outputs carry the site's own rate, the multiple $\mathsf{r}(\delta)\kappa\ge2\kappa$ of the rate $\kappa$ of the expansion, which is the rate $\kappa_{out}$ of the run by the rate lemma for the site of the fluctuation step.

The activity of $V^{m,z}$ is measured per cell of $\Gamma^0_n$. This is the budget of the weighted-size bound of the interpolation and of the smallness condition (h4): for every region $Y$, the budget is per unit of $\sum_j|\Gamma^0_j\cap Y|$, as in clause (a) of Lemma~\ref{9.1:lem-per-cell-domination} and in clause (c2)$^{\star}$ of Lemma~\ref{9.1:lem-resolved-interpolation}. The cells of the site that carry activity lie off every live island, by clause (d) of the island-relative geometry of the fluctuation step. There the island-blind version of clause (b) of Lemma~\ref{9.1:lem-per-cell-domination} bounds the island-blind part with the constant $K^{\circ}$. Per such cell the activity $\mathsf{W}^m$ of $V^{m,z}$ satisfies
\begin{equation}\label{9.1:eq-fluctuation-boundary-bound-5}
\begin{aligned}
\mathsf{W}^m
&\le\theta_iC_L''\cdot[\text{weighted sizes of the direct class-$n$ marks}]\\
&\qquad+\theta_iC_L''\cdot[\text{members of stored derived terms}]\\
&\le\theta_iC_L''\big([\![\delta]\!]^{(5)}_n+[\![\delta]\!]^c_n+\mathsf{y}'_n\big)
+\theta_iC_L''K_Y^{\circ}\mathsf{y}_n^{(\le i)}\\
&\le\theta_iC_L'''[\![\delta]\!]^{\natural+,(i)}_n .
\end{aligned}
\end{equation}
The individual terms are justified as follows.
\begin{itemize}
\item The direct marks are bounded by the weighted-size bound of the interpolation: their interpolated parts have, per cell of zone $n$, weighted size at most $C_L^{\sharp}|\varepsilon|[\![\delta]\!]$, where $C_L^{\sharp}$ is the constant of that bound and its interpolation parameter plays the role of $\varepsilon$ here; in this and the following quoted bounds, $[\![\delta]\!]$, $[\![\delta]\!]_n$ and $[\![\delta]\!]^T_k$ denote the load of the quoted statement, in its notation. This bound carries the constant $C_X$ of the two-run comparison theorem and the constant $C_L''$ of the interpolation construction. The classical parts of the direct marks are the entries of $[\![\delta]\!]^c_n$, since by Definition~\ref{9.1:def-resolved-members}(a) the marks carry both parts inside the mark. The own members are the entries of $\mathsf{y}'_n$, by the assignment of own members with the crossing weights $w^{\mathbf{B}}$.
\item The members $G^{(n)}$ are bounded by clause (b) of Lemma~\ref{9.1:lem-per-cell-domination}.
\item The factor $\theta_i$ majorises the prefactor $\mathsf{w}_{site}^{-1}\eta_p$ of the piece bound of the correlation expansion,
\begin{equation*}
|v_p|\le2\mathsf{w}_{site}^{-1}\eta_pe^{-(\kappa_1-1)m(p)}\mathfrak{N}(v),
\end{equation*}
where $v_p$ is the contribution of the piece $p$, $\eta_p$ its weight, $\mathfrak{N}(v)$ the norm of the input $v$, and $m(p)$ the length of $p$ governing its decay, all in the notation of that bound.
This prefactor is common to every piece at the site, the stored derived terms read there included. It is at most $\theta_i$ by the activity bound of the interpolation, which gives
\begin{equation*}
\mathsf{W}\le\mathsf{W}_J+|\varepsilon|\mathsf{w}^{-1}C_L[\![\delta]\!]_n,\qquad \mathsf{w}^{-1}\le\theta,
\end{equation*}
where $\mathsf{W}$ is the activity at the site, $\mathsf{W}_J$ the activity bound of the correlation theorem, $C_L$ the constant of the activity bound of the interpolation, $\varepsilon$ the interpolation parameter as above, and $\theta=\theta_i$ at the site; and, for the marked part,
$\mathsf{W}^m\le\theta_kC_L[\![\delta]\!]^T_k$ (here $k=i$).
\item For a member of a stored derived term, $\mathfrak{N}(v)$ is twice its stored size, by the entry of the correlation theorem's expansion that lists the stored derived terms.
\item The weighted-size bound of the interpolation also contains a maximum over adjacent zones. That maximum concerns the sum over all classes near one cell, whereas $V^{m,z}$ holds one class only.
\end{itemize}
Inserting the bound \eqref{9.1:eq-fluctuation-boundary-bound-5} on $\mathsf{W}^m$ into the bound $K_{CE}\mathsf{W}^me^{-\kappa d_{i+1}(X)}$ of the corollary, at the rate $\kappa_{out}$ of the site, gives, uniformly in $\lambda$,
\begin{equation*}
|\mathfrak{d}^{mk}(\lambda)|\le K_{CE}C_L'''\theta_i[\![\delta]\!]^{\natural+,(i)}_n
e^{-\kappa_{out}d_{i+1}(X)}
=\varkappa^{\mathbf{B}}_i[\![\delta]\!]^{\natural+,(i)}_ne^{-\kappa_{out}d_{i+1}(X)} .
\end{equation*}

\smallskip
\emph{Step $(\beta)$: the operator direction.} Put $G(\mu):=\mathcal{P}_X[x^{\lambda}+\mu W^{(z),op}]$. By clause (iii) of Lemma~\ref{9.1:lem-zone-split}, the perturbation $\mu W^{(z),op}$ does three things:
\begin{itemize}
\item it moves the operators at the site by at most $|\mu|C_{cl}^T(\log x_n)^{c_2}\vartheta_T^n$ in $\|\cdot\|_\delta$;
\item it adds potentials of activity at most $|\mu|C_{cl}^T(\log x_n)^{c_2}\vartheta_T^n\bar{\varepsilon}^{pr}$;
\item it moves the normaliser inside its additive class by at most $|\mu|C_{cl}^T(\log x_n)^{c_2}\vartheta_T^nC_F$ times the one-run bound.
\end{itemize}
Define the radius $R_n$ (used in this proof only) by
\begin{equation*}
R_n:=\frac{\tfrac12\mathfrak{m}_G^T}{(1+\bar{\varepsilon}^{pr}+C_F)\,C_{cl}^T(\log x_n)^{c_2}\vartheta_T^n},
\end{equation*}
and let $|\mu|\le R_n$. Then three things hold.
\begin{itemize}
\item The operators stay in the ball of the tolerance lemma. The base operators lie within $\varepsilon_{pr}+\mathsf{d}_*$ of the centre of that ball, in the notation and by the statement of clause (iv) of Lemma~\ref{9.1:lem-zone-split}; the displacement by $\mu W^{(z),op}$ adds at most $\frac12\mathfrak{m}_G^T$.
\item The added activity is at most $\frac12\mathfrak{m}_G^T\bar{\varepsilon}^{pr}$. This lies inside the factor of the smallness condition (h4) and inside the class of potential pieces of the two-run comparison theorem.
\item The normaliser stays within the class of normaliser pieces of that theorem.
\end{itemize}
The tolerance lemma states that the output at the site is holomorphic in the interpolation and decoupling parameters of the expansion at the site and in the remaining parameters, and obeys the tolerance bounds, so that the output-producing steps of the correlation expansion and the localisation lemma hold unchanged under the interpolation. Hence $G$ is holomorphic on $|\mu|\le R_n$, with $|G|\le E^{\sharp}_Te^{-\kappa_{out}d_{i+1}(X)}$ there. Cauchy's estimate at $\mu=0$ gives
\begin{equation*}
\begin{aligned}
|G'(0)|&\le E^{\sharp}_Te^{-\kappa_{out}d_{i+1}(X)}/R_n\\
&=\frac{2E^{\sharp}_T(1+\bar{\varepsilon}^{pr}+C_F)C_{cl}^T(\log x_n)^{c_2}\vartheta_T^n}
{\mathfrak{m}_G^T}\,e^{-\kappa_{out}d_{i+1}(X)}\\
&=\sigma^{T,(n)}e^{-\kappa_{out}d_{i+1}(X)} .
\end{aligned}
\end{equation*}
The last equality is \eqref{9.1:eq-fluctuation-boundary-bound-2}. No load and no radius $\Lambda_z$ enters Step $(\beta)$: the radius $R_n$ is the tolerance lemma's own.

The operator difference of an old zone, of size $C_{cl}^T(\log x_n)^{c_2}\vartheta_T^n$, is assigned to the class of zone $n$. Its size is then the classical load of that class, and it is controlled by the age factor of \eqref{9.1:eq-resolved-interpolation-a} at every reading level. Since $\mathfrak{d}^{op}(\lambda)=G'(0)$,
\begin{equation*}
|\mathfrak{d}^{op}(\lambda)|\le\sigma^{T,(n)}e^{-\kappa_{out}d_{i+1}(X)} .
\end{equation*}

\smallskip
\emph{Conclusion of} (c). Add the bounds of Steps $(\alpha)$ and $(\beta)$ and integrate in $\lambda$ over $[0,1]$. This gives
\begin{equation*}
|\mathfrak{V}^{(z)}|\le\big(\varkappa^{\mathbf{B}}_i[\![\delta]\!]^{\natural+,(i)}_n+\sigma^{T,(n)}\big)e^{-\kappa_{out}d_{i+1}(X)}.
\end{equation*}
Multiply by the weight $e^{\kappa_{\mathbf{B}}d_{i+1}(X)}$, which is at most $e^{\kappa_{out}d_{i+1}(X)}$ because $\kappa_{\mathbf{B}}\le\kappa_{out}$. Taking the supremum over $\mathsf{v}$, local labels and $X$ gives (c).

\emph{The comparison} \eqref{9.1:eq-fluctuation-boundary-bound-3}. By \eqref{9.1:eq-zones-boundary-terms-b}, $\vartheta_T=\vartheta_w$ and $\vartheta_w^n=\vartheta_2^{n/4}\vartheta_\times^n$. Hence
\begin{equation*}
\begin{aligned}
\sigma^{T,(n)}&=C_T^{\natural}E_0(\log x_n)^{c_2}\vartheta_w^n
=C_T^{\natural}E_0(\log x_n)^{c_2}\vartheta_2^{n/4}\cdot\vartheta_\times^n\\
&\le(C_T^{\natural}/C^{\natural}_{cl})\,[\![\delta]\!]^c_n\,\vartheta_\times^n .
\end{aligned}
\end{equation*}
The inequality holds because, on aligned zones, \eqref{9.1:eq-zones-boundary-terms-a} gives
\begin{equation*}
[\![\delta]\!]^c_n\ge E_0\min\{1,C^{\natural}_{cl}(\log x_n)^{c_2}\vartheta_2^{n/4}\}
=E_0C^{\natural}_{cl}(\log x_n)^{c_2}\vartheta_2^{n/4}.
\end{equation*}
This is the first inequality of \eqref{9.1:eq-fluctuation-boundary-bound-3}. The second follows from it because the rate letter $\vartheta_\times$ is a negative power of $L$, hence at most $1$.
\end{proof}

\begin{remark}
Five remarks on Theorem~\ref{9.1:thm-fluctuation-boundary-bound}.
\begin{itemize}
\item[(r1)] At $n=i$, the top zone, the bound (c) of \eqref{9.1:eq-fluctuation-boundary-bound-a} is the bound $\varkappa_i[\![\delta]\!]+\sigma^T_i$ that part~(a) of the two-run comparison theorem gives for the kernels of $\mathbf{B}$, in the notation of that theorem. Theorem~\ref{9.1:thm-fluctuation-boundary-bound} extends that bound, class by class, to the remaining terms of \cite[(3.47)]{B-CMP119}. The one expectation value that passes wholly into $\mathbf{B}^{(i+1)}$ in \cite{B-CMP119} is an output of the localisation lemma of the same kind.
\item[(r2)] The marked-potential corollary of the cluster expansion bounds by the activity of the marked potential. That activity is a supremum over the cells of $X$, which include cells of adjacent and older zones and earlier stored terms $\mathbf{B}$. In the proof, $V^{m,z}$ holds one class only. The per-cell size of that class is its load wherever the cell lies (Lemma~\ref{9.1:lem-per-cell-domination}, and the weighted-size bound of the interpolation).
\item[(r3)] The top members, the reading parts and the own members of the stored derived terms inside $X$ belong to their own classes, and are differentiated when those classes are interpolated. Nothing is counted twice, nothing is lost, and the sum over $z$ returns the full difference, by (a) of \eqref{9.1:eq-fluctuation-boundary-bound-a}.
\item[(r4)] The members carry the rate $\kappa_{out}$ of the production; the accounting of decay rates elsewhere in the chapter is unaffected, and the only rate proper to the present bounds is that of the per-cell domination \eqref{9.1:eq-per-cell-domination-a}.
\item[(r5)] Nothing in the bound (c) depends on the run's length $\mathfrak{n}$, the reading level $k'$, $K$, $\check{R}$, any age, $m_0$, the torus, or the two integers written $N$ and $N_0$ in the setting of the two-run comparison theorem. The quantities $\theta_i$ and $x_n$ depend on $g$ only.
\end{itemize}
\end{remark}

\begin{proposition}[Two-run bound for boundary terms of the large-field step]
\label{9.1:prop-large-field-boundary}
Let $\mathsf{A}$ and $\mathsf{B}$ be the two runs compared, let $i$ be a level, and assume the following.
\begin{itemize}
\item[(1)] The hypotheses of the dial theorem at the site of the operation $\mathbf{R}$, parts (A)--(D), namely: the resolution statement for the production arrays; the load bound for the production block $P$; the fourth response estimate in its natural form; the interpolation statement on the dial domain $\mathbb{D}^{\star}$; geometric forgetting of old regions (listed separately as (2)); and the three conditions of the storage rule that govern the $\mathbf{R}$-site. The conclusions (A)--(D) of that theorem are assumed as well, and the clause of its part (D)(e) stating that its bound holds also for the history part of $\mathbf{B}'$ is used as stated.
\item[(2)] Clauses (d) and (e) of the age-dial lemma for the operations $\mathbf{T}'$, and geometric forgetting of old regions.
\item[(3)] Lemma~\ref{9.1:lem-resolved-interpolation}.
\item[(4)] Definition~\ref{9.1:def-resolved-members} for the stored terms of all levels $\le i$.
\item[(5)] The second-order load condition, read at the pair $(2E_1^z,C_L''')$ formed from the load constant $E_1^z$ and the constant $C_L'''$ of that condition.
\item[(6)] The strip bound: $(\alpha^{\sharp})^{1/3}\le g_k^{\kappa_0}$ for $\gamma$ small.
\end{itemize}
Let $\mathbf{B}'^{(i)}(X')$ be a boundary term of an aligned label produced by the operation $\mathbf{R}_i$ after the fluctuation step at level $i$, with top zone $i$. Assume that its stored domain is of the plain class. For each class $z$ write $n(z)$ for the zone of $z$. Write $\mathcal{Z}_{surv}$ and $\mathcal{Z}_{own}$ for the surviving and own class sets of Definition~\ref{9.1:def-resolved-members}. Write $\delta\mathbf{B}'^{(i),(n(z))}$ for the resolved member of the class $z\in\mathcal{Z}_{surv}$, and $\delta\mathbf{B}'^{(i),\mathrm{own}}$ for the own member. Write $\delta\mathbf{B}'^{(i),(n)}$ for the resolved member of a surviving class $z$ with $n(z)=n$, put $\mathbf{B}'^{\mathrm{res}}:=\sum_{z\in\mathcal{Z}_{surv}}\delta\mathbf{B}'^{(i),(n(z))}$, and let $\mathbf{B}'^{\mathrm{own}}$ denote the own part. Finally, with the load constant $E_1^z$ of (5), put
\begin{equation*}
\varkappa^{\mathbf{B}'}_i:=\frac{2c_1^{\flat}(g_i)}{E_1^z}.
\end{equation*}
Then the following hold.
\begin{equation}\label{9.1:eq-large-field-boundary-a}
\begin{aligned}
&\sum_{z\in\mathcal{Z}_{surv}}\delta\mathbf{B}'^{(i),(n(z))}+\delta\mathbf{B}'^{(i),\mathrm{own}}
=\mathbf{B}'^{(i),\mathsf{B}}-\mathbf{B}'^{(i),\mathsf{A}},\\
&\mathsf{S}^{+,(n)}_i[\mathbf{B}'^{\mathrm{res}}]\le\varkappa^{\mathbf{B}'}_i\,
[\![\delta]\!]^{\natural+,(i)}_n,\\
&\bar{\varkappa}':=\sup_n\sum_{i\ge n}\varkappa^{\mathbf{B}'}_i
\le\frac{4O(1)}{E_1^z}\sum_{t\ge0}(X-2c_2+\lambda t)^{-\kappa_0/2},\\
&\mathsf{S}_i^{+}[\mathbf{B}'^{\mathrm{own}}]\le c_1^{\flat}(g_i)\Bigl[4E_0^{-1}\sum_{n\le i}(i-n+1)\,
q_a^{(i-n-1)_+}[\![\delta]\!]^{\natural+,(i)}_n\\
&\hphantom{\mathsf{S}_i^{+}[\mathbf{B}'^{\mathrm{own}}]\le c_1^{\flat}(g_i)\Bigl[}
+4q_a^{\,i-n_c}+2\bigl(9C_\chi^{\sharp}/\vartheta+2\bigr)\vartheta^{i}\Bigr].
\end{aligned}
\end{equation}
\begin{itemize}
\item[(a)] The first line of \eqref{9.1:eq-large-field-boundary-a} holds at every datum, and all its members are holomorphic in the data $\mathsf{v}$.
\item[(b)] Let $n\le i-1$ and suppose $\delta\mathbf{B}'^{(i),(n)}\not\equiv0$. Read zone $n$ of the label in the sequence form used in \eqref{9.1:eq-resolved-interpolation-a}. Then $X'$ meets zone $n$ of the label after $\mathbf{R}_i$ at a point lying
 \begin{itemize}
 \item outside the healed components of $\Lambda_i^c$, collars included,
 \item outside every live second-class or arrested component,
 \item inside a surviving untouched component.
 \end{itemize}
 Hence, at every later level $k'$ at which the term is read and as long as it is stored, some live component $C$ of $Z_{k'}$ with $j(C)\le n+1$ meets $X'^{\sim}$ (the enlargement of $X'$ used in \eqref{9.1:eq-resolved-interpolation-a}), as in \eqref{9.1:eq-resolved-interpolation-a}. The class of the top zone $n=i$ belongs to the own member. It is stored by the storage rule at the zone of the block, whose age is the age of the block (Lemma~\ref{9.1:lem-resolved-interpolation}).
\item[(c)] For every zone $n\le i-1$ the second line of \eqref{9.1:eq-large-field-boundary-a} holds. Its decay rate in the output length is the rate of \cite[(1.99)]{B-CMP122b}, namely $(1+\tfrac12\beta_s)\kappa=\tfrac98\kappa$, and $\tfrac98\kappa\ge\kappa_{\mathbf{B}}$. The third line holds, where $X=g^{-2}$ is the inverse reference coupling (not a region) and $\lambda=c_0/2$. Since $\kappa_0\ge7$, where $\kappa_0$ is the exponent of the coupling power in the bounds of (e), its right side is a convergent series, and it tends to $0$ as $\gamma\to0$.
\item[(d)] The fourth line of \eqref{9.1:eq-large-field-boundary-a} holds. It is the display of part (D)(e) of the dial theorem with $e^{-p_0(g_i)}$ replaced by $c_1^{\flat}(g_i)$; $E_0$, $n_c$ and $C_\chi^{\sharp}$ are the load constant, the index and the constant that enter that display.
\item[(e)] Put $T:=\log g^{-2}$ and let $T_\gamma$ be its value at the coupling ceiling $\gamma$. Let $R(T)$ and $N(T)$ be the scaffold radius $\check{R}_k$ and the count $N(g_k)$ of the strip bound, written as functions of $T$, and let $B_0$ and $C_0^{\sharp}$ be the constants of the value family under the strip. In \eqref{9.1:eq-large-field-boundary-b}, in the bound below and in the absorption inequality, the exponent $c_1$ is the polylogarithm exponent of $\mathsf{l}_s=(\log\theta_s)^{c_1}$; it is not Ba{\l}aban's constant $c_1$ of \cite[(1.97)]{B-CMP122b}, which the strip reading turns into $c_1^{\flat}$ and which is the meaning of $c_1$ wherever it multiplies an activity bound. Suppose that
\begin{equation}\label{9.1:eq-large-field-boundary-b}
\begin{split}
\sup_{T\ge T_\gamma}\Bigl[O(1)\,C_{fib}\bigl(1+\mathsf{K}^0+K^{\circ}(B_0+2O(1))+N(T)C_0^{\sharp}\bigr)\Bigr]^{1/3}
&\\
\times e^{-\frac{1}{12}\beta_\sharp\kappa(R(T)-4)}\,e^{\kappa_0T/2}\,
T^{c_1+2q_\flat+\mathfrak{a}}&\le1 .
\end{split}
\end{equation}
Then for every level $\ell$
\begin{equation*}
c_1^{\flat}(g_\ell)\le2O(1)\,g_\ell^{\kappa_0}(\log x_\ell)^{-c_1-2q_\flat-\mathfrak{a}}.
\end{equation*}
Consequently, every use of the display of part (D)(e) of the dial theorem that absorbs polylogarithmic factors through the inequality
\begin{equation*}
e^{-p_0(g_k)}\le g_k^{\kappa_0}(\log x_k)^{-c_1-2q_\flat}
\end{equation*}
(the polylogarithm-absorption inequality of the age-dial lemma) does so with $c_1^{\flat}$ in place of $e^{-p_0}$.
\end{itemize}
\end{proposition}

\begin{proof}
Throughout, $\mathcal{P}_{X'}$ is the formula by which $\mathbf{B}'^{(i)}(X')$ is produced from the array of its stored slot values (Definition~\ref{9.1:def-resolved-members}). The array $x^{\lambda}$ is the convex combination of the arrays of the two runs, and $W^{(z)}$ is the mark of the class $z$.

\emph{Proof of (a).} The argument is the one for the first display of Theorem~\ref{9.1:thm-fluctuation-boundary-bound}, see \eqref{9.1:eq-fluctuation-boundary-bound-a}. The production of the operation $\mathbf{R}$ takes the place of the production of the fluctuation step.

By parts (A)--(C) of the dial theorem at the $\mathbf{R}$-site, the production of $R'$ and of $\mathbf{B}'$ is one formula applied to the array of stored numbers at a physical base. This formula is holomorphic on $\mathbb{D}^{\star}$. For the numbers produced at the step, this is the interpolation statement on $\mathbb{D}^{\star}$ in hypothesis (1). For the stored boundary terms of earlier steps that the production reads, it is Lemma~\ref{9.1:lem-resolved-interpolation}.

Part (c) of the first-variation lemma, applied with the class set $\mathcal{Z}=\mathcal{Z}_{surv}\sqcup\mathcal{Z}_{own}$, writes the difference $\mathbf{B}'^{(i),\mathsf{B}}-\mathbf{B}'^{(i),\mathsf{A}}$ at every datum as the sum over $z\in\mathcal{Z}$ of the first-variation members of the classes. Each of these members is holomorphic in $\mathsf{v}$. The members of the classes in $\mathcal{Z}_{own}$ are summed, by definition, into the single member $\delta\mathbf{B}'^{(i),\mathrm{own}}$. This is the first line of \eqref{9.1:eq-large-field-boundary-a}.

\emph{Proof of (b).} Let $z\in\mathcal{Z}_{surv}$. By Definition~\ref{9.1:def-resolved-members}, this means that $n=n(z)\le i-1$ and that some slot of class $z$ lies outside the healed components. Such a slot is a use at step $i$ of one of three kinds:
\begin{itemize}
\item a direct use anchored in zone $n$, by the anchoring rule for direct uses;
\item a use supported on a piece meeting zone $n$, by the zone split of pieces together with that anchoring rule;
\item a stored boundary term meeting zone $n$, by the localisation statement among the displays \eqref{9.1:eq-fluctuation-boundary-bound-a}.
\end{itemize}
In each case the use lies inside $X'$. Indeed, in \cite{B-CMP122b} the term $R'(X)$ depends on the background field $U_k$ restricted to $X$, and the support statement of the dial theorem places both kinds of stored units inside the mark. The use also lies in the part of the label that $\mathbf{R}_i$ leaves standing.

Every zone $n\le i-1$ lies in $Z_i=\Lambda_i^c$, hence in some component of it. Since the slot lies outside the healed components, its piece of zone $n$ lies in a surviving component. The label keeps the surviving components after $\mathbf{R}_i$. In \cite{B-CMP119} the large-field regions are changed after the $\mathbf{R}$-operation: the healed components, and everything that $\mathbf{T}'_i$ integrated in them, leave the label, as stated in \cite{B-CMP122a} and in the removal clause of the storage rule. Moreover, zone $n=\Omega_n\setminus\Omega_{n+1}$ near a surviving component is a fixed set from then on.

Therefore the argument of \eqref{9.1:eq-resolved-interpolation-a} in Lemma~\ref{9.1:lem-resolved-interpolation} applies at every later level $k'$. The surviving component is alive as long as the term is stored, and its creation step satisfies $j\le n+1$. This proves (b).

The class of the top zone is not resolved, for the following reason. Its slots lie in $\Omega_i$, which later splits into zone $i$ and $\Omega_{i+1}$. Whether $X'$ still meets zone $i$ at a later level is not decided at production. Storage at the zone of the block, by the storage rule, always carries the age of the block. The cost of this choice is that no age factor is available at production. That cost is nil, because the weight of the term $n=i$ in the fourth line of \eqref{9.1:eq-large-field-boundary-a} is $(i-i+1)q_a^{0}=1$.

\emph{Proof of (c).} Apply part (a) of the first-variation lemma with class radius
\begin{equation*}
\Lambda_z:=E_1^z\big/[\![\delta]\!]^{\natural+,(i)}_n ,
\end{equation*}
which satisfies $\Lambda_z\ge1$ by the a priori load bound, and with
\begin{equation}\label{9.1:eq-large-field-boundary-1}
\mathfrak{B}_z:=\sup_{|\varepsilon|\le\Lambda_z}\bigl|\mathcal{P}_{X'}[x^{\lambda}+\varepsilon W^{(z)}]\bigr|
\le2O(1)\,c_1^{\flat}(g_i)\,e^{-\frac98\kappa d_{out}(X')} .
\end{equation}

The bound \eqref{9.1:eq-large-field-boundary-1} is obtained as follows. On $\mathbb{D}^{\star}$ the activities \cite[(1.91)]{B-CMP122b} obey \cite[(1.97)]{B-CMP122b}, hence the bound \cite[(1.98)--(1.99)]{B-CMP122b}. This is how part (D)(e) of the dial theorem is proved. There, by the holomorphy statement of that theorem, the activities \cite[(1.91)]{B-CMP122b} obey \cite[(1.97)]{B-CMP122b} with $2c_1$ on the domain $\mathbb{D}^{\natural+}$ of that proof; here it is used in its form on $\mathbb{D}^{\star}$. Two replacements are made:
\begin{itemize}
\item Ba{\l}aban's $c_1$ is read as $c_1^{\flat}$ under the strip, the strip reading replacing $\alpha$ by $\alpha^{\sharp}$;
\item the factor $2$ is the factor of the bound $2\mathfrak{B}_X$ in clause (e) of the age-dial lemma.
\end{itemize}
The dialled activities obey \cite[(1.97)]{B-CMP122b} with $2c_1$ in the relative length. The cluster-expansion step then delivers the rate $(1+\tfrac12\beta_s)\kappa=\tfrac98\kappa$ of \cite[(1.99)]{B-CMP122b} in the output length $d_{out}=d_{\cup_\iota Y_\iota^{\sharp}}$. For the plain stored domains considered here, $d_{out}(X')=d_i(X')$, and \eqref{9.1:eq-large-field-boundary-1}, with the input as read above, holds with $\tfrac98\kappa\ge\kappa_{\mathbf{B}}$.

Part (a) of the first-variation lemma therefore gives
\begin{equation*}
\bigl|\mathfrak{V}^{(z)}\bigr|\le\frac{\mathfrak{B}_z}{\Lambda_z}
=\frac{2O(1)\,c_1^{\flat}(g_i)}{E_1^z}\,[\![\delta]\!]^{\natural+,(i)}_n\,
e^{-\frac98\kappa d_{out}(X')} .
\end{equation*}
With $\varkappa^{\mathbf{B}'}_i=2c_1^{\flat}(g_i)/E_1^z$, this is the second line of \eqref{9.1:eq-large-field-boundary-a} at the stated rate.

No age factor exists for a surviving zone at production, and none is claimed. The factor in the age of the surviving component is supplied at the later levels by \eqref{9.1:eq-resolved-interpolation-a}.

For the third line, we first bound $c_1^{\flat}$. Here $c_1^{\flat}$ is Ba{\l}aban's $c_1$ under the strip reading, which replaces $\alpha$ by $\alpha^{\sharp}$; through this reading, (e) or the strip bound $(\alpha^{\sharp})^{1/3}\le g_k^{\kappa_0}$ of hypothesis (6) gives
\begin{equation*}
c_1^{\flat}(g_i)\le2O(1)\,g_i^{\kappa_0}=2O(1)\,x_i^{-\kappa_0/2}.
\end{equation*}
By the reference lower bound on the running inverse couplings, $x_i\ge X-2c_2+\lambda t$ at depth $t$, where $X=g^{-2}$ is the inverse reference coupling (not a region) and $\lambda=c_0/2$ is the slope of that lower bound (not the interpolation variable of the dial). Hence, at depth $t$,
\begin{equation*}
\varkappa^{\mathbf{B}'}_i\le\frac{4O(1)}{E_1^z}(X-2c_2+\lambda t)^{-\kappa_0/2},
\end{equation*}
and the sum over $i\ge n$ is bounded by the series over $t\ge0$. The series converges for $\kappa_0>2$, in particular for $\kappa_0\ge7$. It decreases in $X$ and tends to $0$ as $X\to\infty$. Since $X=g^{-2}$ with $g$ below the ceiling $\gamma$, this gives $\bar{\varkappa}'\to0$ as $\gamma\to0$.

\emph{Proof of (d).} We treat first the class of the top zone, then the inner classes, then the remaining steps.

The class of the top zone $n=i$ has $a:=(i-i-1)_+=0$. For this class there is no $\mu$-dial and $\mathfrak{B}_z=2\mathfrak{B}_{X'}$, where $\mathfrak{B}_{X'}$ is the modulus bound of the dial theorem, which bounds $|\mathbf{B}'(X';\lambda,\mu)|$ on the dial. Part (a) of the first-variation lemma is applied at $\Lambda_z$. The weight of this class in the display is $(i-i+1)q_a^{0}=1$. This is the term $n=k$ of the display of part (D)(e) of the dial theorem, taken without change.

Now let $z$ belong to the inner class set $\mathcal{Z}_{inner}\subset\mathcal{Z}_{own}$ of Definition~\ref{9.1:def-resolved-members}, the classes all of whose slots lie inside healed components of $\Lambda_i^c$. Every slot of class $z$ lies inside a healed component $X_j$ of $\Lambda_i^c$, namely in $\mathbf{T}'_i(X_j)$. A slot of class $(n,\cdot)$ there forces zone $n$ to be present inside $X_j^{\sim}$. Hence $j(X_j)\le n+1$, and the age of $X_j$ at its healing is at least $i-n-1$. This is clause (d) of the age-dial lemma, whose region is the component operated on by $\mathbf{T}'$.

Put $a:=(i-n-1)_+$ and dial the part $\mathbf{T}'_{i,\ge a}$ as in clause (e) of the age-dial lemma, that is, use
\begin{equation*}
\mathbf{T}'_{i,<a}+\mu\,\mathbf{T}'_{i,\ge a},\qquad|\mu|\le q^{-a}.
\end{equation*}
Geometric forgetting of old regions keeps \cite[(1.92)]{B-CMP122b}. The sum in the definition \cite[(1.71)]{B-CMP122b} acts also on the last exponential in \cite[(1.72)]{B-CMP122b}, so Ba{\l}aban's exponent $V(Y)$ of \cite[(1.72)]{B-CMP122b} is split together with it. Then $\mathbf{B}'(\cdot\,;\mu=0)$ does not contain $\lambda_z$. This is the statement, in the proof of clause (e) of the age-dial lemma, that $R'(\cdot,0)$ does not contain the dial variable $\lambda_{n,k'}$ of the class $(n,k')$, and it is what the dial theorem means by the history part.

The rest of the proof of part (D)(e) of the dial theorem is used without change, with $c_1^{\flat}$ in place of $e^{-p_0}$. It consists of:
\begin{itemize}
\item the pointwise bound $|\mathbf{B}'(X';\lambda,\mu)|\le\mathfrak{B}_{X'}$ on the dial;
\item the first-derivative estimate in $\mu$;
\item the first-variation lemma over the inner classes, each with class radius $\Lambda_z$;
\item the supremum over seeds;
\item the non-aligned inner sequences and the remainder $\Delta^{\chi}$, which give the last two terms of the display.
\end{itemize}
The replacement of $e^{-p_0}$ by $c_1^{\flat}$ does not affect the age factors. The surplus $q_a^{a+1}$ supplied by geometric forgetting of old regions sits in the first exponential of \cite[(1.79)]{B-CMP122b}, as in the proof of the age-dial lemma. It is independent of which constant, Ba{\l}aban's $c_1$ or $c_1^{\flat}$, appears in \cite[(1.97)]{B-CMP122b}. With the clause of part (D)(e) on the history part of $\mathbf{B}'$, used as stated in hypothesis (1), this gives the fourth line of \eqref{9.1:eq-large-field-boundary-a}.

\emph{Proof of (e).} The bound on $c_1^{\flat}$ is obtained from \eqref{9.1:eq-large-field-boundary-b} by the mechanism of the strip bound $(\alpha^{\sharp})^{1/3}\le g_k^{\kappa_0}$ for $\gamma$ small, hypothesis (6). In the derivation of that bound:
\begin{itemize}
\item $\alpha^{\sharp}\le C_{fib}(1+\mathsf{K}_{val})\,e^{-\frac14\beta_\sharp\kappa(\check{R}_k-4)}$;
\item $\check{R}_k\ge(\log g_k^{-2})^{r_{pr}}$ with $r_{pr}\ge2$;
\item $\mathsf{K}_{val}(k)\le\mathsf{K}^0+N(g_k)C_0^{\sharp}$ (the supremum form of the value bound under the strip).
\end{itemize}
The left side of \eqref{9.1:eq-large-field-boundary-b} carries, besides these members, the member $K^{\circ}(B_0+2O(1))$. Its supremum runs over all $T\ge T_\gamma$, that is, over all couplings below the ceiling. At $T=\log g_\ell^{-2}$ one has $e^{-\kappa_0T/2}=g_\ell^{\kappa_0}$ and $T=\log x_\ell$. In these variables \eqref{9.1:eq-large-field-boundary-b} is read as the stated bound on $c_1^{\flat}(g_\ell)$.

The second assertion of (e) rests on the form of the first. The majorant $2O(1)g_\ell^{\kappa_0}(\log x_\ell)^{-c_1-2q_\flat-\mathfrak{a}}$ is the right side of the absorption inequality $g_\ell^{\kappa_0}(\log x_\ell)^{-c_1-2q_\flat}$ multiplied by $2O(1)(\log x_\ell)^{-\mathfrak{a}}$.
\end{proof}

\begin{remark}[Classes of the summed sequences]
Every term $R'^{(k)}(X)$, those of the first group included, contains at least one large-field region connected with one of the operations $\mathbf{T}'$. Its first-variation classes therefore include zones of the summed sequences, and these zones are absent from every later label. If such classes were stored as resolved members, the localisation \eqref{9.1:eq-resolved-interpolation-a} would fail for them.

These are the classes of $\mathcal{Z}_{inner}\subset\mathcal{Z}_{own}$. They are kept out of the resolved members and bounded as a whole, together with the class of the top zone, by the fourth line of \eqref{9.1:eq-large-field-boundary-a}; their age factors are used at production. The resolved members are those of the old surviving classes $n\le i-1$, for which \eqref{9.1:eq-resolved-interpolation-a} holds at every later level.

A zone met both inside a healed history and outside it belongs to $\mathcal{Z}_{surv}$. The whole class is then resolved, with its age factor supplied at the later levels and none at production. By the first line of \eqref{9.1:eq-large-field-boundary-a}, no slot is counted twice and none is lost.

The own member is one function of the data per stored term. It is stored by the storage rule at the zone of the block and enters the loads as the boundary entry $\mathsf{y}'$ with the crossing weight $w^{\mathbf{B}}$. Its display is the display of part (D)(e) of the dial theorem with $e^{-p_0}$ replaced by $c_1^{\flat}$. Under \eqref{9.1:eq-large-field-boundary-b}, $c_1^{\flat}$ obeys the absorption inequality $g^{\kappa_0}(\log x)^{-c_1-2q_\flat}$ in the same way as $e^{-p_0}$.
\end{remark}

In what follows $n$, $n'$ and $\ell$ denote zones of the aligned label under consideration. The aligned zones form the finite set $\{n_c,\dots,k\}$, where $n_c$ is the lowest aligned zone and $k$, the stop level, is the top zone. We use the following objects.
\begin{itemize}
\item The zone-resolved members of the stored boundary terms of Definition~\ref{9.1:def-resolved-members}.
\item The crossing weights $\mathsf{w}_{\ell,n}$. By \eqref{9.1:eq-resolved-members-a} their sums satisfy
\begin{equation*}
\sum_{\ell}\mathsf{w}_{\ell,n}\le 3 .
\end{equation*}
\item The weights $w^{\mathbf{B}}_{n,\ell}=L^{-(n-\ell)/2}$, $\ell\le n$, of the crossing-weight statement for the
stored boundary terms.
\item The norm $\mathsf{S}^{+,(n)}_\ell[\,\cdot\,]$, the weighted supremum norm of the difference between the
two runs $\mathsf{A}$ and $\mathsf{B}$ of the member of zone $n$ of a store produced at level $\ell$.
\item The norm $\mathsf{S}^{+}_\ell[\,\cdot\,]$, the same norm of a whole member.
\end{itemize}

The stores of the fluctuation step and of the large-field step are written $\mathbf{B}^{(\ell)}$ and $\mathbf{B}'^{(\ell)}$. The residual members and the own member of $\mathbf{B}'^{(\ell)}$ are written $\mathbf{B}'^{(\ell),\mathrm{res}}$ and $\mathbf{B}'^{(\ell),\mathrm{own}}$. The half-steps are written $\mathbf{T}_\ell$ and $\mathbf{R}_\ell$.

We use the loads of zone $n$ attached to the zone-resolved members of
Definition~\ref{9.1:def-resolved-members}, which are bounded a priori in
\eqref{9.1:eq-apriori-loads-a}. They are the following.
\begin{itemize}
\item $[\![\delta]\!]^{(5)}_n$ is the part carried by the entries of the five unknowns.
\item $[\![\delta]\!]^{c}_n$ is the part carried by the classical entries.
\item $[\![\delta]\!]^{\natural+}_n$ is the total load.
\item $[\![\delta]\!]^{\natural+,(\ell)}_n$ is the total load as read at level $\ell$. It satisfies
$[\![\delta]\!]^{\natural+,(\ell)}_n\le[\![\delta]\!]^{\natural+}_n$.
\end{itemize}

The boundary loads are the following two.
\begin{itemize}
\item $\mathsf{y}_n$ collects the members of zone $n$ of the stores produced after level $n$. A store
$\mathbf{B}^{(\ell)}$, $\ell\ge n+1$, contributes $\mathsf{w}_{\ell,n}\mathsf{S}^{+,(n)}_\ell[\mathbf{B}^{(\ell)}]$.
A store $\mathbf{B}'^{(\ell)}$, $\ell\ge n+1$, contributes
$\mathsf{S}^{+,(n)}_\ell[\mathbf{B}'^{(\ell),\mathrm{res}}]$.
\item $\mathsf{y}'_n$ is given by
\begin{equation*}
\mathsf{y}'_n=\sum_{\ell\le n}w^{\mathbf{B}}_{n,\ell}\,\mathsf{S}^+_\ell\bigl[\mathbf{B}'^{(\ell),\mathrm{own}}\bigr].
\end{equation*}
\end{itemize}
For a level $k'$ we write $\mathsf{y}_n^{(\le k')}$ for the part of $\mathsf{y}_n$ carried by the stores
produced up to level $k'$.

We abbreviate by $\mathfrak{a}_n$ the sum in the first formula below; the second formula is the
decomposition of the total load of Definition~\ref{9.1:def-resolved-members} into its five-unknown,
classical and boundary parts:
\begin{equation}\label{9.1:eq-load-elimination-1}
\mathfrak{a}_n:=[\![\delta]\!]^{(5)}_n+[\![\delta]\!]^{c}_n+\mathsf{y}'_n,
\qquad
[\![\delta]\!]^{\natural+}_n=\mathfrak{a}_n+\mathsf{y}_n .
\end{equation}

\emph{Order of production.} In the renormalization transformation of \cite{B-CMP119}, each step is the composition of the operations $\mathbf{T}$ and $\mathbf{R}$, applied in this order. Hence, for a fixed zone $n$, the stores carrying a member resolved to zone $n$ appear in the half-step order
\begin{equation*}
\mathbf{T}_{n+1},\ \mathbf{R}_{n+1},\ \mathbf{T}_{n+2},\ \mathbf{R}_{n+2},\ \dots
\end{equation*}
Here $\mathbf{T}_{n+1}$ produces $\mathbf{B}^{(n+1)}$, whose top member at that moment is the member of zone $n$. The half-step $\mathbf{R}_{n+1}$ produces $\mathbf{B}'^{(n+1)}$, with members of zones $\le n$. Each production reads only stores produced before it.

For fixed $n$, the quantity $\mathfrak{a}_n$ reads levels $\le n$ only. Indeed, the youngest member of $\mathsf{y}'_n$, namely $\mathbf{B}'^{(n),\mathrm{own}}$, is produced at $\mathbf{R}_n$, before $\mathbf{T}_{n+1}$. Consequently, every production reads a load of zone $n$ that is at most
\begin{equation}\label{9.1:eq-load-elimination-2}
\mathfrak{a}_n+u ,
\end{equation}
where $u$ is the partial sum of $\mathsf{y}_n$ over the stores produced before it.

\begin{proposition}[Elimination of the boundary-term loads]\label{9.1:prop-load-elimination}
Assume the conclusions of Theorem~\ref{9.1:thm-fluctuation-boundary-bound} and of
Proposition~\ref{9.1:prop-large-field-boundary}, applied inductively along the half-steps. Wherever clause~(a) of the first-variation bound is used, the a priori bound on the loads and the unit dial radius, \eqref{9.1:eq-apriori-loads-a}, supplies $\Lambda_z\ge 1$.

The constants are as follows.
\begin{itemize}
\item $\bar{\varkappa}$ bounds the coefficients $\varkappa^{\mathbf{B}}_i$ of the two-run bound of
\eqref{9.1:eq-fluctuation-boundary-bound-a} for the members of the fluctuation-step stores.
\item $\bar{\varkappa}'$ bounds the sum $\sum_\ell\varkappa^{\mathbf{B}'}_\ell$ of the coefficients of the
bound of \eqref{9.1:eq-large-field-boundary-a} for the residual members
$\mathbf{B}'^{(\ell),\mathrm{res}}$.
\item $C_T^{\mathbf{B}}$, the source term $\sigma^{T,(n)}$ of the two-run bound for the member of zone $n$,
and the rate $\vartheta_\times$ are those of
\eqref{9.1:eq-fluctuation-boundary-bound-a} and \eqref{9.1:eq-zones-boundary-terms-b}.
\item $E_0$, $c_1^{\flat}(g_\ell)$, $q_a$, $C_\chi^{\sharp}$ and $\vartheta$ are the constants of the bound
of \eqref{9.1:eq-large-field-boundary-a} for the own member $\mathbf{B}'^{(\ell),\mathrm{own}}$.
\end{itemize}

Put
\begin{equation*}
\bar{a}:=3\bar{\varkappa}+\bar{\varkappa}',\qquad \epsilon_0:=e^{\bar{a}}\bar{a},\qquad
\bar{c}_1:=\sup_\ell c_1^{\flat}(g_\ell).
\end{equation*}
For $n'\le n$ put
\begin{equation*}
\mathfrak{k}(n,n'):=4E_0^{-1}\sum_{\ell:\ n'\le\ell\le n}w^{\mathbf{B}}_{n,\ell}\,c_1^{\flat}(g_\ell)\,
(\ell-n'+1)\,q_a^{(\ell-n'-1)_+}\ \ge 0 ,
\end{equation*}
and put $\mathfrak{k}(n,n'):=0$ for $n'>n$. Further put
\begin{align*}
\mathfrak{s}_n&:=\sum_{\ell\le n}w^{\mathbf{B}}_{n,\ell}\,c_1^{\flat}(g_\ell)
\Bigl[4q_a^{\ell-n_c}+2\bigl(9C_\chi^{\sharp}/\vartheta+2\bigr)\vartheta^{\ell}\Bigr],\\
\mathfrak{K}&:=(1+\epsilon_0)\,\mathfrak{k},\\
\mathfrak{a}^0_n&:=(1+\epsilon_0)\bigl([\![\delta]\!]^{(5)}_n+[\![\delta]\!]^{c}_n+\mathfrak{s}_n\bigr)
+3e^{\bar{a}}C_T^{\mathbf{B}}[\![\delta]\!]^{c}_n .
\end{align*}

For the envelope, let $\mathsf{l}_s$, $e_s$, $\hat{\mathsf{l}}_m$, $\rho$, the kernel rate $q$, the load weights $\varsigma_\nu$ and $C_\varsigma$ be those of \eqref{9.1:eq-load-envelope-a}. Let $s(n)$ be the depth of zone $n$, so that $s(n')=s(n)+n-n'$. Put
\begin{equation*}
\omega_n:=\mathsf{l}_{s(n)}^2\,e_{s(n)},\qquad
\|\psi\|_\omega:=\sup_n |\psi_n|/\omega_n .
\end{equation*}
Write $\|\mathfrak{M}\|_\omega$ for the norm of a matrix $\mathfrak{M}$ on the aligned zones as an operator for $\|\cdot\|_\omega$. Put
\begin{align*}
K_\omega&:=4C_\varsigma\sum_{\xi,\eta\ge0}(1+\xi)^3(\eta+1)\,\hat{\mathsf{l}}_{\xi+\eta}^{2}\,
\rho^{\xi}\rho_\eta ,\\
\rho_0&:=1,\qquad \rho_\eta:=\rho^{3\eta-4}\ \ (\eta\ge1),\\
\epsilon_1&:=(1+\epsilon_0)\,4E_0^{-1}\bar{c}_1K_\omega .
\end{align*}

The matrix $\mathfrak{K}$ on the finitely many aligned zones is non-negative and lower-triangular in age:
its entry $(n,n')$ vanishes unless $n'\le n$. Its diagonal entries are at most
$(1+\epsilon_0)\,16\bar{c}_1/E_0$, and this number is less than $1$ for $\gamma$ sufficiently small. Hence
\begin{equation*}
\mathfrak{G}:=\sum_{\iota\ge0}\mathfrak{K}^\iota=(I-\mathfrak{K})^{-1}
\end{equation*}
exists and is non-negative and lower-triangular.

Then the following hold for every aligned zone $n$:
\begin{equation}\label{9.1:eq-load-elimination-a}
\begin{aligned}
&\text{(i)}\quad \mathsf{y}_n\le e^{\bar{a}}\bigl(\bar{a}\,\mathfrak{a}_n
+3C_T^{\mathbf{B}}\vartheta_\times^{n}[\![\delta]\!]^{c}_n\bigr)\\
&\qquad\qquad\le\epsilon_0\mathfrak{a}_n+3e^{\bar{a}}C_T^{\mathbf{B}}[\![\delta]\!]^{c}_n ,\\
&\text{(ii)}\quad \mathsf{y}_n+\mathsf{y}'_n\le\epsilon_0\bigl([\![\delta]\!]^{(5)}_n+[\![\delta]\!]^{c}_n\bigr)
+3e^{\bar{a}}C_T^{\mathbf{B}}[\![\delta]\!]^{c}_n\\
&\qquad\qquad\qquad +(1+\epsilon_0)\sum_{n'\le n}\mathfrak{k}(n,n')[\![\delta]\!]^{\natural+}_{n'}
+(1+\epsilon_0)\mathfrak{s}_n ,\\
&\text{(iii)}\quad [\![\delta]\!]^{\natural+}_n\le\sum_{n'\le n}\mathfrak{G}(n,n')\,\mathfrak{a}^0_{n'} ,\\
&\text{(iv)}\quad \|\mathfrak{K}\|_\omega\le\epsilon_1,\qquad
\|\mathfrak{G}\|_\omega\le(1-\epsilon_1)^{-1},\\
&\qquad\qquad [\![\delta]\!]^{\natural+}_n\le(1-\epsilon_1)^{-1}\|\mathfrak{a}^0\|_\omega\,\omega_n .
\end{aligned}
\end{equation}

The following qualifications belong to the statement.
\begin{itemize}
\item In (i), the bound holds uniformly in the number of later steps. Stores in the island-blind format add
terms of the same form, with the same coefficients.
\item The source $\mathfrak{a}^0$ in (iii) contains only the entries of the five unknowns and the classical
entries of the loads.
\item In (iv), $\epsilon_1<1$ for $\gamma$ sufficiently small.
\item $\epsilon_0$, $\epsilon_1$ and $\bar{c}_1$ tend to $0$ as $\gamma\to0$.
\item None of them depends on $\mathfrak{n}$ (the length of the run), $k'$, $K$, $N$, $N_0$,
$\check{R}$ (the scaffold radius), an age, or $m_0$.
\end{itemize}
\end{proposition}

\begin{proof}
\emph{Clause (i).} Fix $n$. Let $u_m$, $m=0,1,\dots$, be the successive partial sums of
$\mathsf{y}_n^{(\le k')}$. The sums are taken along the half-steps after which a store with a member of
zone $n$ exists, and $u_0=0$. By \eqref{9.1:eq-load-elimination-2}, the load of zone $n$ read at the
production following the $m$-th of these half-steps is at most $\mathfrak{a}_n+u_m$.

Consider first a half-step $\mathbf{T}$ producing $\mathbf{B}^{(\ell)}$ with $\ell\ge n+1$. It adds
$\mathsf{w}_{\ell,n}\mathsf{S}^{+,(n)}_\ell[\mathbf{B}^{(\ell)}]$ to the partial sum. By the two-run bound
of \eqref{9.1:eq-fluctuation-boundary-bound-a} for the member of zone $n$ of a fluctuation-step store,
applied with the load it read, this increment satisfies
\begin{equation*}
\mathsf{w}_{\ell,n}\,\mathsf{S}^{+,(n)}_\ell[\mathbf{B}^{(\ell)}]
\le\mathsf{w}_{\ell,n}\bigl[\varkappa^{\mathbf{B}}_{\ell-1}(\mathfrak{a}_n+u_m)+\sigma^{T,(n)}\bigr].
\end{equation*}

Consider next a half-step $\mathbf{R}$ producing $\mathbf{B}'^{(\ell)}$ with $\ell\ge n+1$. It adds
$\mathsf{S}^{+,(n)}_\ell[\mathbf{B}'^{(\ell),\mathrm{res}}]$. By the bound of
\eqref{9.1:eq-large-field-boundary-a} for the residual members this is at most
$\varkappa^{\mathbf{B}'}_\ell(\mathfrak{a}_n+u_m)$.

Stores in the island-blind format contribute terms of the same form with the same coefficients, by the
island-blind reading of the stores. Hence in all cases
\begin{equation*}
u_{m+1}\le(1+\mu_m)\,u_m+\mu_m\,\mathfrak{a}_n+\mu'_m ,
\end{equation*}
where the coefficients are
\begin{equation*}
\mu_m\in\{\mathsf{w}_{\ell,n}\varkappa^{\mathbf{B}}_{\ell-1},\ \varkappa^{\mathbf{B}'}_\ell\},\qquad
\mu'_m\in\{\mathsf{w}_{\ell,n}\sigma^{T,(n)},\ 0\}.
\end{equation*}

We bound the sums of the coefficients. Since $\varkappa^{\mathbf{B}}_{\ell-1}\le\bar{\varkappa}$, the bound $\sum_\ell\mathsf{w}_{\ell,n}\le3$ of \eqref{9.1:eq-resolved-members-a} and the bound $\sum_{\ell\ge n}\varkappa^{\mathbf{B}'}_\ell\le\bar{\varkappa}'$ of Proposition~\ref{9.1:prop-large-field-boundary} give
\begin{equation*}
\sum_m \mu_m\le\bar{\varkappa}\sum_\ell\mathsf{w}_{\ell,n}+\sum_{\ell\ge n}\varkappa^{\mathbf{B}'}_\ell
\le3\bar{\varkappa}+\bar{\varkappa}'=\bar{a}.
\end{equation*}
Again by $\sum_\ell\mathsf{w}_{\ell,n}\le3$, and by the comparison of $\sigma^{T,(n)}$ with the classical
load contained in the two-run bound of \eqref{9.1:eq-fluctuation-boundary-bound-a}, with the rate
$\vartheta_\times$ of \eqref{9.1:eq-zones-boundary-terms-b},
\begin{equation*}
\sum_m \mu'_m\le3\,\sigma^{T,(n)}\le3\,C_T^{\mathbf{B}}\vartheta_\times^{n}[\![\delta]\!]^{c}_n .
\end{equation*}

We show by induction on $\bar m$ that
\begin{equation}\label{9.1:eq-load-elimination-3}
u_{\bar m}\le\sum_{m<\bar m}(\mu_m\mathfrak{a}_n+\mu'_m)\prod_{m<m'<\bar m}(1+\mu_{m'}) .
\end{equation}
For $\bar m=0$ both sides vanish. If the bound holds for $\bar m$, the recursion gives
\begin{align*}
u_{\bar m+1}&\le(1+\mu_{\bar m})\sum_{m<\bar m}(\mu_m\mathfrak{a}_n+\mu'_m)\prod_{m<m'<\bar m}(1+\mu_{m'})
+\mu_{\bar m}\mathfrak{a}_n+\mu'_{\bar m}\\
&=\sum_{m<\bar m+1}(\mu_m\mathfrak{a}_n+\mu'_m)\prod_{m<m'<\bar m+1}(1+\mu_{m'}),
\end{align*}
since the empty product for $m=\bar m$ equals one.

Each product in \eqref{9.1:eq-load-elimination-3} satisfies
\begin{equation*}
\prod(1+\mu_{m'})\le e^{\sum_m \mu_m}\le e^{\bar{a}} .
\end{equation*}
Therefore, for every $\bar m$,
\begin{equation*}
u_{\bar m}\le e^{\sum_m \mu_m}\Bigl(\sum_m \mu_m\,\mathfrak{a}_n+\sum_m \mu'_m\Bigr)
\le e^{\bar{a}}\bigl(\bar{a}\,\mathfrak{a}_n+3C_T^{\mathbf{B}}\vartheta_\times^{n}[\![\delta]\!]^{c}_n\bigr).
\end{equation*}
For the last half-step up to level $k'$ one has $u_{\bar m}=\mathsf{y}_n^{(\le k')}$, and the bound is
independent of $\bar m$ and of $k'$. This gives the first inequality of (i), uniformly in the number of later
steps. The second inequality of (i) uses $e^{\bar{a}}\bar{a}=\epsilon_0$ and $\vartheta_\times^{n}\le1$.

\emph{Clause (ii).} The load $\mathsf{y}_n$ is bounded by (i). For $\mathsf{y}'_n$ we use the bound of
\eqref{9.1:eq-large-field-boundary-a} for the own member $\mathbf{B}'^{(\ell),\mathrm{own}}$. Multiplied by
$w^{\mathbf{B}}_{n,\ell}$ and summed over $\ell\le n$, the part of its right side carrying the loads
contributes
\begin{equation*}
\begin{split}
&\sum_{\ell\le n}w^{\mathbf{B}}_{n,\ell}\,c_1^{\flat}(g_\ell)\,4E_0^{-1}\\
&\qquad\times\sum_{n'\le\ell}(\ell-n'+1)\,q_a^{(\ell-n'-1)_+}[\![\delta]\!]^{\natural+,(\ell)}_{n'}\\
&\qquad\le\sum_{n'\le n}\mathfrak{k}(n,n')\,[\![\delta]\!]^{\natural+}_{n'} .
\end{split}
\end{equation*}
Here the finite sums over $\ell\le n$ and $n'\le\ell$ have been exchanged into sums over $n'\le n$ and $n'\le\ell\le n$. We also used $[\![\delta]\!]^{\natural+,(\ell)}_{n'}\le[\![\delta]\!]^{\natural+}_{n'}$. The remaining parts of its right side contribute exactly $\mathfrak{s}_n$. Thus
\begin{equation}\label{9.1:eq-load-elimination-4}
\mathsf{y}'_n\le\sum_{n'\le n}\mathfrak{k}(n,n')\,[\![\delta]\!]^{\natural+}_{n'}+\mathfrak{s}_n .
\end{equation}

Now insert $\mathfrak{a}_n=[\![\delta]\!]^{(5)}_n+[\![\delta]\!]^{c}_n+\mathsf{y}'_n$ from \eqref{9.1:eq-load-elimination-1} into the second inequality of (i). This gives
\begin{equation*}
\mathsf{y}_n\le\epsilon_0\bigl([\![\delta]\!]^{(5)}_n+[\![\delta]\!]^{c}_n\bigr)+\epsilon_0\,\mathsf{y}'_n
+3e^{\bar{a}}C_T^{\mathbf{B}}[\![\delta]\!]^{c}_n .
\end{equation*}
Adding $\mathsf{y}'_n$ gives the term $(1+\epsilon_0)\mathsf{y}'_n$. Bounding it by $(1+\epsilon_0)$ times the right side of \eqref{9.1:eq-load-elimination-4} yields (ii).

\emph{Clause (iii).} By \eqref{9.1:eq-load-elimination-1} and then (ii),
\begin{align*}
[\![\delta]\!]^{\natural+}_n&=\mathfrak{a}_n+\mathsf{y}_n
\le[\![\delta]\!]^{(5)}_n+[\![\delta]\!]^{c}_n+(\mathsf{y}_n+\mathsf{y}'_n)\\
&\le\mathfrak{a}^0_n+\bigl(\mathfrak{K}\,[\![\delta]\!]^{\natural+}\bigr)_n .
\end{align*}
Indeed, $[\![\delta]\!]^{(5)}_n+[\![\delta]\!]^{c}_n$ plus the first two terms and the last term of (ii) add up to $\mathfrak{a}^0_n$. The remaining term of (ii) is $\sum_{n'\le n}\mathfrak{K}(n,n')[\![\delta]\!]^{\natural+}_{n'}$.

We turn to the diagonal of $\mathfrak{k}$. For $n'=n$, the only summand in $\mathfrak{k}(n,n)$ is $\ell=n$. There the factor $(\ell-n'+1)q_a^{(\ell-n'-1)_+}$ equals $(0+1)q_a^0=1$, and $w^{\mathbf{B}}_{n,n}=L^0=1$. Hence
\begin{equation*}
\mathfrak{k}(n,n)=4E_0^{-1}c_1^{\flat}(g_n)\le4\bar{c}_1/E_0 ,
\end{equation*}
and therefore
\begin{equation*}
\mathfrak{K}(n,n)\le(1+\epsilon_0)\,4\bar{c}_1/E_0\le(1+\epsilon_0)\,16\bar{c}_1/E_0<1 .
\end{equation*}
Off the diagonal, $\mathfrak{K}(n,n')$ is non-zero only for $n'<n$. So $\mathfrak{K}$ is non-negative and lower-triangular on the finite set of aligned zones. Its eigenvalues are its diagonal entries, so its spectral radius is its largest diagonal entry and is less than one.

Write $\mathfrak{K}=\mathfrak{D}+\mathfrak{N}$, with $\mathfrak{D}$ diagonal and $\mathfrak{N}$ strictly lower-triangular. The matrix $I-\mathfrak{D}$ is diagonal with entries $1-\mathfrak{K}(n,n)>0$. So $(I-\mathfrak{D})^{-1}$ exists and is non-negative. Since
\begin{equation*}
(I-\mathfrak{D})\bigl(I-(I-\mathfrak{D})^{-1}\mathfrak{N}\bigr)=I-\mathfrak{K},
\end{equation*}
and $(I-\mathfrak{D})^{-1}\mathfrak{N}$ is non-negative and strictly lower-triangular on the $k-n_c+1$
aligned zones, hence nilpotent (its $(k-n_c+1)$-th power vanishes), we obtain
\begin{equation*}
(I-\mathfrak{K})^{-1}=\bigl(I-(I-\mathfrak{D})^{-1}\mathfrak{N}\bigr)^{-1}(I-\mathfrak{D})^{-1}
=\sum_{\iota=0}^{k-n_c}\bigl((I-\mathfrak{D})^{-1}\mathfrak{N}\bigr)^\iota(I-\mathfrak{D})^{-1}\ge0 .
\end{equation*}
This inverse is lower-triangular.

It equals $\sum_\iota\mathfrak{K}^\iota$. Write $\mathfrak{G}_{\le\tau}:=\sum_{\iota\le\tau}\mathfrak{K}^\iota$ for $\tau\ge0$.
Then $(I-\mathfrak{K})\mathfrak{G}_{\le\tau}=I-\mathfrak{K}^{\tau+1}$, so
\begin{equation*}
\mathfrak{G}_{\le\tau}=(I-\mathfrak{K})^{-1}-(I-\mathfrak{K})^{-1}\mathfrak{K}^{\tau+1}
\le(I-\mathfrak{K})^{-1} ,
\end{equation*}
because $(I-\mathfrak{K})^{-1}\ge0$ and $\mathfrak{K}^{\tau+1}\ge0$. The partial sums
$\mathfrak{G}_{\le\tau}$ are entrywise non-decreasing in $\tau$ and bounded by $(I-\mathfrak{K})^{-1}$, so
they converge to a limit $\mathfrak{G}$. The terms $\mathfrak{K}^{\tau+1}$ of the convergent non-negative
series then tend to zero, and $(I-\mathfrak{K})\mathfrak{G}=I$. By uniqueness of the inverse,
$\mathfrak{G}=(I-\mathfrak{K})^{-1}$. Thus $\mathfrak{G}$ exists, equals $(I-\mathfrak{K})^{-1}$, and is
non-negative and lower-triangular.

The load vector $[\![\delta]\!]^{\natural+}=([\![\delta]\!]^{\natural+}_n)_n$ is non-negative, and it is finite by the a priori bound on the loads and the unit dial radius, \eqref{9.1:eq-apriori-loads-a}. The inequality above reads $[\![\delta]\!]^{\natural+}\le\mathfrak{a}^0+\mathfrak{K}[\![\delta]\!]^{\natural+}$, that is, $(I-\mathfrak{K})[\![\delta]\!]^{\natural+}\le\mathfrak{a}^0$. Applying the non-negative matrix $(I-\mathfrak{K})^{-1}$ gives $[\![\delta]\!]^{\natural+}\le\mathfrak{G}\,\mathfrak{a}^0$, which is (iii).

\emph{Clause (iv).} Let $\psi$ be any vector on the aligned zones. Since $\mathfrak{K}\ge0$ and $|\psi_{n'}|\le\|\psi\|_\omega\,\omega_{n'}$, exchanging the finite sums back as in the proof of (ii) gives
\begin{equation*}
\begin{split}
|(\mathfrak{K}\psi)_n|&\le\|\psi\|_\omega\,(1+\epsilon_0)\,4E_0^{-1}\bar{c}_1\\
&\qquad\times\sum_{\ell\le n}\sum_{n'\le\ell}w^{\mathbf{B}}_{n,\ell}\,(\ell-n'+1)\,
q_a^{(\ell-n'-1)_+}\,\omega_{n'} ,
\end{split}
\end{equation*}
where also $c_1^{\flat}(g_\ell)\le\bar{c}_1$ was used. Put $\xi:=n-\ell$ and $\eta:=\ell-n'$. Then
$s(n')=s(n)+\xi+\eta$.

By the monotonicity of $e_s$ and of $\mathsf{l}_s$ in the depth,
\eqref{9.1:eq-load-envelope-a},
\begin{equation*}
\omega_{n'}\le\hat{\mathsf{l}}_{\xi+\eta}^{2}\,\rho^{-(\xi+\eta)}\,\omega_n .
\end{equation*}
By the comparison of the weights $w^{\mathbf{B}}_{n,\ell}$ with the load weights $\varsigma$ in the
crossing-weight statement for the stored boundary terms, and by the bound on the load weights
$\varsigma_\nu$ in \eqref{9.1:eq-load-envelope-a},
\begin{equation*}
w^{\mathbf{B}}_{n,\ell}=L^{-\xi/2}\le4\varsigma_{1+\xi}\le4C_\varsigma(1+\xi)^3q^\xi .
\end{equation*}
Hence the $(\xi,\eta)$-summand of the double sum is at most
\begin{equation*}
4C_\varsigma(1+\xi)^3(\eta+1)\,\hat{\mathsf{l}}_{\xi+\eta}^{2}\,(q/\rho)^\xi\,q_a^{(\eta-1)_+}
\rho^{-\eta}\,\omega_n .
\end{equation*}

We now use the two inequalities between the kernel rate $q$, the age rate $q_a$ and $\rho$, namely
$q/\rho\le\rho$ and $q_a/\rho\le\rho^3$. They give
$(q/\rho)^\xi\le\rho^\xi$. For $\eta\ge1$ they give
\begin{equation*}
q_a^{\eta-1}\rho^{-\eta}=(q_a/\rho)^{\eta-1}\rho^{-1}\le\rho^{3(\eta-1)-1}=\rho^{3\eta-4}=\rho_\eta .
\end{equation*}
For $\eta=0$ the factor $q_a^{0}\rho^{0}$ equals $1=\rho_0$. Since $\rho<1$, the factor
$\rho^\xi\rho_\eta$ is $\rho^{-4}\rho^{\xi+3\eta}$ for $\eta\ge1$, and the polynomial factors and
$\hat{\mathsf{l}}_{\xi+\eta}^2$ are summable against $\rho^{\xi+3\eta}$. Hence $K_\omega<\infty$.
Extending the sum to all $\xi,\eta\ge0$ bounds the double sum by $K_\omega\,\omega_n$.

Therefore $\|\mathfrak{K}\psi\|_\omega\le\epsilon_1\|\psi\|_\omega$, that is, $\|\mathfrak{K}\|_\omega\le\epsilon_1$. It follows that $\|\mathfrak{K}^\iota\|_\omega\le\epsilon_1^\iota$, and, since $\epsilon_1<1$,
\begin{equation*}
\|\mathfrak{G}\|_\omega\le\sum_{\iota\ge0}\epsilon_1^\iota=(1-\epsilon_1)^{-1} .
\end{equation*}
Finally, $[\![\delta]\!]^{\natural+}\le\mathfrak{G}\,\mathfrak{a}^0$ from (iii) gives
\begin{equation*}
[\![\delta]\!]^{\natural+}_n\le(\mathfrak{G}\,\mathfrak{a}^0)_n\le\|\mathfrak{G}\,\mathfrak{a}^0\|_\omega\,\omega_n
\le(1-\epsilon_1)^{-1}\|\mathfrak{a}^0\|_\omega\,\omega_n .
\end{equation*}

\emph{Smallness.} The constant $\bar{\varkappa}$ is bounded by a constant multiple of the small parameter
$\theta(\gamma)$ of the correlation theorem. The constant $\bar{\varkappa}'$ is bounded by a constant
multiple of $\sum_\ell c_1^{\flat}(g_\ell)$, by Proposition~\ref{9.1:prop-large-field-boundary}. By
\eqref{9.1:eq-large-field-boundary-b}, with $\kappa_0$ the exponent of $\gamma$ in that bound,
\begin{equation*}
\bar{c}_1\le2\,O(1)\,\gamma^{\kappa_0} .
\end{equation*}
Hence $\epsilon_0=e^{\bar{a}}\bar{a}$, $\epsilon_1$ and $\bar{c}_1$ tend to $0$ as $\gamma\to0$. In
particular the inequalities $(1+\epsilon_0)16\bar{c}_1/E_0<1$ and $\epsilon_1<1$ used in (iii) and (iv)
hold for $\gamma$ sufficiently small. None of the bounds used above involves $\mathfrak{n}$, $k'$, $K$, $N$, $N_0$,
$\check{R}$, an age, or $m_0$.
\end{proof}

\begin{remark}
No fixed point is solved for in the proof of
Proposition~\ref{9.1:prop-load-elimination}. Clause (i) is an induction on the step, and clause (iii) is
the inversion of an explicit triangular matrix. The loads were finite beforehand, by
\eqref{9.1:eq-apriori-loads-a}.

The uniformity of (i) in the number of steps comes from two summability facts. The first is the bound
$\sum_\ell\mathsf{w}_{\ell,n}\le3$ of \eqref{9.1:eq-resolved-members-a}. The second is the bound
$\sum_\ell\varkappa^{\mathbf{B}'}_\ell\le\bar{\varkappa}'$ of
Proposition~\ref{9.1:prop-large-field-boundary}. The crossing of zone $\ell-1$ costs a factor
$e^{-\mathsf{d}_{\ell-1}}$, and, by the lower bound on the crossing lengths of the zones, the crossing
length satisfies, with a constant $c_\times$ and with $x_{\ell-1}=g_{\ell-1}^{-2}$ the running inverse
coupling of Definition~\ref{1.2:def-couplings},
\begin{equation*}
\mathsf{d}_{\ell-1}\ge2(\log x_{\ell-1})^{r_{pr}}/L-c_\times ,
\end{equation*}
which grows towards the ultraviolet. No bound on ages or on the number of steps enters.
\end{remark}

\begin{lemma}[A priori bound on the loads and the unit dial radius]\label{9.1:lem-apriori-loads}
Let $\bar a$ and $\epsilon_0=e^{\bar a}\bar a$ be the constants of the elimination of the
boundary-term loads, Proposition~\ref{9.1:prop-load-elimination}. Let $\epsilon_1$,
the $\mathbf{R}$-born coefficient $\bar c_1$, the load-weight constant $S_\varsigma$, the
crossing-weight constant $S_{\mathbf{B}}$, and the constants $E_0$, $E_1$, $B_0$,
$C_\chi^{\sharp}$, $C_T^{\mathbf{B}}$ be those of the zone-resolved members
\eqref{9.1:eq-resolved-members-a}--\eqref{9.1:eq-resolved-members-b} of this section, and let
$\vartheta$ be the rate letter of Notation~\ref{0.2:not-glossary}. Put
\begin{equation*}
C_q:=\sum_{m\ge 0}(m+1)\,q_a^{(m-1)_+},\qquad (m-1)_+:=\max\{m-1,0\}.
\end{equation*}
Define $E_1^z$ by the first line of \eqref{9.1:eq-apriori-loads-a} and assume the remaining four
conditions, which are ceilings on the coupling $g$:
\begin{equation}\label{9.1:eq-apriori-loads-a}
\begin{gathered}
E_1^z:=\Bigl(\tfrac54+4C_T^{\mathbf{B}}\Bigr)E_1+\tfrac54 B_0,\\
\bar a\le\tfrac1{40},\qquad \epsilon_1\le\tfrac1{64},\qquad
(1+\epsilon_0)\,\frac{16\,\bar c_1}{E_0}\le\frac1{16},\\
S_{\mathbf{B}}\,\bar c_1\Bigl[4E_0^{-1}E_1^zC_q+4+2\Bigl(\frac{9C_\chi^{\sharp}}{\vartheta}+2\Bigr)\Bigr]
\le\frac{B_0}{4}.
\end{gathered}
\end{equation}
Then $\epsilon_0\le 1/38$ and $\epsilon_{\mathbf{B}}:=\epsilon_0+2\epsilon_1\le 1/16$. Moreover,
every load read by a production step, or by a step that reads the stored terms, obeys
\begin{equation}\label{9.1:eq-apriori-loads-1}
[\![\delta]\!]^{\natural+,(k')}_n\le E_1^z
\end{equation}
for every zone $n$ and every level $k'$ at which it is read. Consequently
\begin{equation*}
\Lambda_z=\frac{E_1^z}{[\![\delta]\!]^{\natural+}_{n(z)}}\ge 1\qquad\text{on }\mathbb{D}^{\star}.
\end{equation*}
\end{lemma}

\begin{proof}
\emph{The two small constants.} For $0\le x<1$ one has $e^{x}\le 1/(1-x)$; with
$x=1/40$ this gives $e^{1/40}\le 40/39\le 1.03$. Hence, by $\bar a\le 1/40$,
\begin{equation*}
\epsilon_0=e^{\bar a}\bar a\le\frac{1.03}{40}\le\frac1{38},\qquad
\epsilon_{\mathbf{B}}=\epsilon_0+2\epsilon_1\le\frac1{38}+\frac1{32}=\frac{35}{608}
\le\frac{38}{608}=\frac1{16},
\end{equation*}
where $\epsilon_1\le 1/64$ was used. The same estimate gives $3e^{\bar a}\le 3.09\le 4$, and
$e^{\bar a}\bar a\le 1/38\le 1/4$.

\emph{Induction.} We argue by induction along the half-steps, that is, along the
fluctuation step and the large-field step of each renormalisation step, in their order. The
first stored term of an aligned label reads no member of a stored boundary term, so the
induction hypothesis is vacuous at the start. Suppose that every load read so far obeys
\eqref{9.1:eq-apriori-loads-1}. Lemma~\ref{9.1:lem-resolved-interpolation},
Theorem~\ref{9.1:thm-fluctuation-boundary-bound} and
Proposition~\ref{9.1:prop-large-field-boundary} read the radius of the dial domain
$\mathbb{D}^{\star}$ at $E_1^z$ and their load hypothesis at $2E_1^z$ (see the remark following
this proof); by the induction hypothesis these hypotheses hold, so the three statements apply to
the current production, the first-variation bound being used at $\Lambda_z\ge 1$ where it enters.

\emph{The loads other than the boundary loads.} The five-unknown part and the classical part of
the load of zone $n$ satisfy
\begin{equation}\label{9.1:eq-apriori-loads-2}
[\![\delta]\!]^{(5)}_n+[\![\delta]\!]^c_n\le E_1 .
\end{equation}
This is the count of part (a) of the lemma bounding the crossing weights
$w^{\mathbf{B}}_{n,j}$, carried out with $\mathsf{S}^+[\mathbf{B}]\le 2B_0$,
$\hat{\vartheta}^{\natural}\le 2$ and the normalisation \eqref{9.1:eq-centre-normalisation-d}:
the two remaining families of five-unknown entries of the load of zone $n$ (those of the table
members and of the response row), the entries coming from the running shift $\zeta$ and from
the step shifts $\mathsf{b}$,
together with the classical entries other than those attached to the leaves of the stored
boundary terms, stay within $4(S_\varsigma+1)E_0$, as in the corresponding entry bounds of the
one-lattice count and of the normalisation at the centres; and the classical entries attached to
those leaves obey
\begin{equation*}
\sum_j w^{\mathbf{B}}_{n,j}\,B_0\,\hat{\vartheta}^{\natural}_j\le 2S_{\mathbf{B}}B_0
\le 4S_{\mathbf{B}}B_0 .
\end{equation*}

\emph{The own load of the large-field step.} By the bound on the own member in
\eqref{9.1:eq-large-field-boundary-a}, applied with the loads it reads bounded by $E_1^z$
(induction hypothesis), and by the last condition in \eqref{9.1:eq-apriori-loads-a},
\begin{equation*}
\begin{split}
\mathsf{y}'_n&\le S_{\mathbf{B}}\sup_{\ell\le n}\mathsf{S}^+_\ell[\mathbf{B}'^{\,own}]\\
&\le S_{\mathbf{B}}\,\bar c_1\Bigl[4E_0^{-1}E_1^zC_q+4
+2\Bigl(\frac{9C_\chi^{\sharp}}{\vartheta}+2\Bigr)\Bigr]\le\frac{B_0}{4}.
\end{split}
\end{equation*}
Together with \eqref{9.1:eq-apriori-loads-2} this gives, for the source term of the elimination,
\begin{equation}\label{9.1:eq-apriori-loads-3}
\mathfrak{a}_n\le E_1+\frac{B_0}{4}.
\end{equation}

\emph{The boundary load.} By item (i) of Proposition~\ref{9.1:prop-load-elimination},
\eqref{9.1:eq-load-elimination-a}, then by
$e^{\bar a}\bar a\le 1/4$ and $3e^{\bar a}\le 4$, by \eqref{9.1:eq-apriori-loads-3}, and by
$[\![\delta]\!]^c_n\le E_1$ from \eqref{9.1:eq-apriori-loads-2},
\begin{equation*}
\mathsf{y}_n\le e^{\bar a}\bar a\,\mathfrak{a}_n+3e^{\bar a}C_T^{\mathbf{B}}[\![\delta]\!]^c_n
\le\frac14\Bigl(E_1+\frac{B_0}{4}\Bigr)+4C_T^{\mathbf{B}}E_1 .
\end{equation*}

\emph{The total load.} Adding \eqref{9.1:eq-apriori-loads-3} and using $B_0\ge 0$,
\begin{equation*}
\begin{split}
[\![\delta]\!]^{\natural+}_n=\mathfrak{a}_n+\mathsf{y}_n
&\le\frac54\Bigl(E_1+\frac{B_0}{4}\Bigr)+4C_T^{\mathbf{B}}E_1\\
&=\Bigl(\frac54+4C_T^{\mathbf{B}}\Bigr)E_1+\frac5{16}B_0\le E_1^z .
\end{split}
\end{equation*}
This is \eqref{9.1:eq-apriori-loads-1} for the current production, which completes the
induction. The bound $\Lambda_z\ge 1$ on $\mathbb{D}^{\star}$ follows by dividing
\eqref{9.1:eq-apriori-loads-1}, at the zone $n(z)$, by $[\![\delta]\!]^{\natural+}_{n(z)}$.
\end{proof}

\begin{remark}
The constant $E_1^z$ replaces $E_1$ wherever the radius of $\mathbb{D}^{\star}$ is read and
wherever the load condition at twice the load bound is read, namely in item (d) of the
definition of the zone-resolved stores and in \eqref{9.1:eq-resolved-members-b}. It is fixed at
the same stage of the order of choice as the constants of the correlation theorem, of the
tolerance lemma at the fluctuation step and of the normalisation
\eqref{9.1:eq-centre-normalisation-d}, before the coupling ceiling $\gamma$; its constituent
$C_T^{\mathbf{B}}$ is
\begin{equation*}
C_T^{\mathbf{B}}=\frac{2\,E^{\sharp}_T\,C_{cl}^T\,(1+\bar{\varepsilon}^{pr}+C_F)}
{E_0\,\mathfrak{m}_G^T\,C^{\natural}_{cl}},
\end{equation*}
where $E^{\sharp}_T$, $C^T_{cl}$, $\bar{\varepsilon}^{pr}$, $C_F$ and $\mathfrak{m}^T_G$ are
constants of the fluctuation step that enter the two-run bound of
Theorem~\ref{9.1:thm-fluctuation-boundary-bound}, and $C^{\natural}_{cl}$ is the classical constant
of the interpolation lemma, Lemma~\ref{9.1:lem-resolved-interpolation}, and of the statements it
takes as inputs.
\end{remark}

\begin{corollary}[The envelope bound without a boundary-term hypothesis]
\label{9.1:cor-envelope-corollary}
Assume the hypotheses of Proposition~\ref{9.1:prop-load-elimination} and of
Lemma~\ref{9.1:lem-apriori-loads}, so that their conclusions hold.
Assume also the following.
\begin{itemize}
\item[(h0)] The conditions on the resolved boundary members and on the loads: the zone-resolved
bound \eqref{9.1:eq-resolved-members-a}, the two further conditions on the resolved boundary
members stated with it, and the ceilings on $g$ of \eqref{9.1:eq-apriori-loads-a}. The two-run bound
\eqref{9.1:eq-large-field-boundary-b} for the boundary terms of the
large-field step. The per-cell domination \eqref{9.1:eq-per-cell-domination-a}.
\end{itemize}
Assume moreover:
\begin{itemize}
\item[(h1)] the five conclusions of the five-unknown response proposition, at the constants of
part~(b) below;
\item[(h2)] part~(a) of the conversion corollary, which puts the response inequalities in which a
load enters into kernel form;
\item[(h3)] the crossing-weight lemma for the weights $w^{\mathbf{B}}_{n,\ell}$;
\item[(h4)] for some $\Lambda>0$ with $\Lambda\le\Lambda_r$ for every depth $r$, the ceiling
\begin{equation}\label{9.1:eq-envelope-corollary-a}
\epsilon_{\mathbf{B}}K_\Lambda\le 4\Lambda,
\qquad
K_\Lambda:=\frac{4E_0^{-1}C_\Sigma^{TC}A_E}{C_E^{\natural}},
\qquad
A_E:=C^{(E)}+C^{(R)}+C^{(\zeta)},
\end{equation}
holds, where $\epsilon_{\mathbf{B}}:=\epsilon_0+2\epsilon_1$, as in
Lemma~\ref{9.1:lem-apriori-loads}; $E_0$ is the constant appearing in the ceilings of that lemma;
$C_E^{\natural}$ is the constant of part~(b)(iii) below; and $\Lambda_r$ is the weight of the
$R$-inequality at depth $r$.
\end{itemize}
Assume finally that the loads are those of their defining display in this chapter (the loads
$[\![\delta]\!]^u_j$, $[\![\delta]\!]^c_j$, $[\![\delta]\!]^{\natural+}_j$ of zone $j$), with $E_1$
replaced by the radius $E_1^z$ of \eqref{9.1:eq-apriori-loads-a}, and that $g\le\gamma_U$. Then
the following hold.

\smallskip
\noindent(a) For every aligned zone $j$, with $s:=\mathfrak{n}-j$,
\begin{equation}\label{9.1:eq-envelope-corollary-1}
[\![\delta]\!]^{\natural+}_j\le C_{env}^z\,\mathsf{l}_s^2\,e_s,
\end{equation}
where
\begin{equation}\label{9.1:eq-envelope-corollary-2}
C_{env}^z:=(1-\epsilon_1)^{-1}\Bigl[(1+\epsilon_0)\bigl(C^{(E)}+C^{(R)}+C^{(\zeta)}+C^{(c)}
+C^{(\mathfrak{s})}\bigr)+3e^{\bar a}C_T^{\mathbf{B}}C^{(c)}\Bigr].
\end{equation}
Here $C^{(E)},C^{(R)},C^{(\zeta)}$ are the constants of the bound on the five-unknown part of the
load in \eqref{9.1:eq-load-envelope-b}. The constant $C^{(c)}$ is that of the bound on the classical
load there, and $C^{(\mathfrak{s})}$ is the own-source constant of \eqref{9.1:eq-own-source-a}.
Consequently, at the stop level $k$,
\begin{equation}\label{9.1:eq-envelope-corollary-3}
\Sigma^{q_a}_k\le C_\Sigma^{TC}\,C_{env}^z\,\mathsf{l}_{\mathsf{w}}^2\,e_{\mathsf{w}} .
\end{equation}
The rate bound $\varepsilon_n\le C_{TC}\,e_{\mathsf{w}}$ of the two-cutoff comparison theorem holds
with $C_{env}^z$ in place of its envelope constant $C_{env}$. In particular this rate bound uses no
envelope hypothesis on the stored boundary terms.

\smallskip
\noindent(b) Consider the kernel-form system of response inequalities on the scaffold together with the
response inequalities of its $R'$-part.
\begin{itemize}
\item[(i)] No stored boundary term, whether produced by a $\mathbf{T}$-step or by an
$\mathbf{R}$-step, enters the $T$-inequalities, that is, the $P$- and $A$-inequalities of the
kernel-form system. At every level $k'$ these bound the outputs of the $T$-site on sets
$X\subset\Lambda_{k'+1}$, where $\Lambda_{k'+1}$ is the small-field region (not to be confused with
the weights $\Lambda_r$), namely $E^{(k'+1)}$, $P^{\flat}$, $D$, $Q$ and the
whole-lattice part of $R''$. Hence the load $[\![\delta]\!]^{\natural T}_{k'}$ entering them is the
boundary-free load of the kernel-form system, and their coupling constants
$\varkappa,\varkappa_A,\varkappa''$ are the same as in that system.
\item[(ii)] The loads enter the $R$-inequality and the $R'$-part only through the age-weighted sum
$\Sigma^{q_a}_{k'}\bigl[[\![\delta]\!]^{\natural+}\bigr]$ at the level $k'$ of the step, by
part~(e) of the conversion corollary and part~(D)(e) of the closing corollary of the response
system. The conversion corollary puts these inequalities into kernel form. In doing so, the constant
$C_\sigma^{\natural}$ of the index and classical sources is replaced by
\begin{equation}\label{9.1:eq-envelope-corollary-4}
C_\sigma^z:=(2+4C_T^{\mathbf{B}})(C_\sigma^{\natural}+1)
+\bigl(4E_0^{-1}C_\Sigma^{TC}+1\bigr)C^{(\mathfrak{s})} .
\end{equation}
With $(M_{ij})$ the kernel and $E=(E_j)$ the envelope of the system, as in
Lemma~\ref{9.6:lem-maximum-principle-slack}, the
evaluation $\sum_jM_{ij}E_j$ at the index $i$ of the $R$-inequality at depth $r$ is enlarged by at
most $\epsilon_{\mathbf{B}}K_\Lambda\,\mathfrak{r}_r\,C_E^{\natural}\,e_r$. The quantity
$\Lambda_r$ is not changed.
\item[(iii)] Take as the constant of the five-unknown response proposition
\begin{equation*}
C_E^{\natural}:=\max\{C_V+1,\,2C_X^{\star},\,2C_\sigma^z\},
\end{equation*}
with $C_V$ and $C_X^{\star}$ the constants of that proposition. Its side conditions are ceilings
on $g$ or properties of $\hat{\mathfrak{r}}$; they hold at $C_\sigma^z$ and $C_E^{\natural}$ for
$\gamma$ small, since none of their forms changes. Then
Lemma~\ref{9.6:lem-maximum-principle-slack} applies with
$\epsilon:=\epsilon_{\mathbf{B}}K_\Lambda/(64\Lambda)\le
1/16$, and it returns the five conclusions of that proposition at the constant
\begin{equation*}
\frac{C_E^{\natural}}{1-2\epsilon}\le\frac{8}{7}\,C_E^{\natural}=:C_E ,
\end{equation*}
namely the bounds on the three response quantities $g^P,g^A,g^R$ of that proposition,
\begin{equation*}
g^P\le C_E\,e_r,\qquad g^A\le 4C_E\,e_r,\qquad g^R\le 4C_E\,\hat{\mathfrak{r}}_r\,e_r,
\end{equation*}
and the $d$-inequality with the constant $C_E$.
\item[(iv)] The whole-lattice part of $R''$ is covered by~(i). Under
\eqref{9.1:eq-envelope-corollary-a}, by the same estimate as in~(ii), the boundary part of the
$R'$-part is at most $\epsilon_{\mathbf{B}}C_E^{\natural}\hat{\mathfrak{r}}_re_r$.
\end{itemize}
In particular the system of the five-unknown response proposition has no sixth unknown: the stored
boundary terms never enter the index set $I$ of Lemma~\ref{9.6:lem-maximum-principle-slack}, and
from its substitution step on, the proof of that proposition uses only the three response
inequalities, the $d$-inequality, the boundary entries and the side conditions.
\end{corollary}

\begin{proof}[Proof of part (a)]
The bounds $(E)$, $(R)$, $(\zeta)$ of \eqref{9.1:eq-load-envelope-b} and the bound on the classical
load there read
\begin{equation*}
[\![\delta]\!]^{(5)}_j\le\bigl(C^{(E)}+C^{(R)}+C^{(\zeta)}\bigr)\mathsf{l}_s^2e_s,
\qquad
[\![\delta]\!]^{c}_j\le C^{(c)}\mathsf{l}_se_s .
\end{equation*}
They are derived from hypothesis~(h1) and use no bound on the stored boundary terms; no bound on
the boundary part of the load enters the present argument. By
Remark~\ref{9.1:rem-own-source}, the own sources satisfy
$\mathfrak{s}_j\le C^{(\mathfrak{s})}e_s$, which is \eqref{9.1:eq-own-source-a}. Since
$\mathsf{l}_s\ge1$, we have $e_s\le\mathsf{l}_se_s\le\mathsf{l}_s^2e_s$, so each of the three
bounds just listed is at most its constant times $\mathsf{l}_s^2e_s$.

We now apply part~(iv) of Proposition~\ref{9.1:prop-load-elimination}. By the three bounds and
the inequality $e_s\le\mathsf{l}_s^2e_s$, the source term $\mathfrak{a}^0$ of that proposition
satisfies
\begin{equation*}
\|\mathfrak{a}^0\|_\omega\le(1+\epsilon_0)\bigl(C^{(E)}+C^{(R)}+C^{(\zeta)}+C^{(c)}
+C^{(\mathfrak{s})}\bigr)+3e^{\bar a}C_T^{\mathbf{B}}C^{(c)} ,
\end{equation*}
and part~(iv) then gives \eqref{9.1:eq-envelope-corollary-1} with the constant
\eqref{9.1:eq-envelope-corollary-2}.

The argument giving the bound on the age-weighted sum in \eqref{9.1:eq-load-envelope-b} applies
verbatim with \eqref{9.1:eq-envelope-corollary-1} as its input and gives
\eqref{9.1:eq-envelope-corollary-3}. The rate bound $\varepsilon_n\le C_{TC}e_{\mathsf{w}}$ of the
two-cutoff comparison theorem, in which the envelope constant $C_{env}$ enters, then holds with
$C_{env}^z$ in its place; no envelope hypothesis on the stored boundary terms has been used.
\end{proof}

\begin{proof}[Proof of part (b)]
The passage to kernel form in the conversion corollary is linear and monotone in the entries of the
loads, since the proof of the kernel-form corollary of the response system, on which it rests, uses
only the right-hand members of the inequalities. Upper bounds for the loads may therefore be
inserted entry by entry. By part~(iii) of Proposition~\ref{9.1:prop-load-elimination}, every load is
bounded by the resolvent $\mathfrak{G}$ applied to its five-unknown and classical entries. In
checking the hypotheses of Lemma~\ref{9.6:lem-maximum-principle-slack}, the inequalities are
evaluated on the candidate envelope $E$, that is, one bounds the sums $\sum_jM_{ij}E_j$.

\emph{Proof of (i).} Fix a level $k'$ and a set $X\subset\Lambda_{k'+1}$. For the outputs
$E^{(k'+1)}$, $P^{\flat}$, $D$, \cite{B-CMP119} states that the integral over the interpolation
parameter $t$ of the boundary terms contributes to $B^{(k'+1)}$ only, and that $E^{(k'+1)}$ comes
from the regular terms of the effective action only. For an arbitrary output on $X$ we argue by
locality, by the localisation lemma of the correlation theorem. The action on which the
$\mathbf{T}$-step at level $k'+1$ acts satisfies the inductive hypothesis of \cite{B-CMP122b} in
the current labelling, and the stored terms linked to healed components have been absorbed, by
\eqref{9.1:eq-resolved-interpolation-a}. Hence a stored boundary term could enter the output on $X$
only if the part $X_u$ of $X$ carrying it met $(Z_\ell^{cur})^{\sim}$ for some $\ell\le k'$, where
$Z_\ell^{cur}$ is the current large-field region created at step $\ell$. But, with
$\Omega_{\ell+1}$ the region of \cite[(2.1), p.~254]{B-CMP119},
$Z_\ell^{\sim}\cap\Omega_{\ell+1}=\emptyset$, since by \cite{B-CMP119} at least six layers of
$LM\check{R}_{\ell+1}$-cubes separate them, and $\Lambda_{k'+1}\subset\Omega_{\ell+1}$ by
\cite[(2.1), p.~254]{B-CMP119}. The argument uses $\Omega_{\ell+1}$ and not the sets $\Lambda_\ell$,
since healing can enlarge the sets $\Lambda_j$ \cite{B-CMP119}. Therefore
$[\![\delta]\!]^{\natural T}_{k'}$ is the boundary-free load of the kernel-form system, and the
coupling constants $\varkappa,\varkappa_A,\varkappa''$ of the $P$- and $A$-inequalities are
unchanged.

\emph{Proof of (ii).} Evaluated on $E$, the five-unknown part of the loads is the explicit function
$\Phi(E)$ of $E$ given by the bounds $(E)$, $(R)$, $(\zeta)$ of \eqref{9.1:eq-load-envelope-b}, and
$\Phi(E)\le A_E\,\omega$ with $A_E$ as in \eqref{9.1:eq-envelope-corollary-a}. Let $c$ denote the
classical load and $\mathrm{id}$ the identity. By part~(iv) of
Proposition~\ref{9.1:prop-load-elimination}, the eliminated boundary part evaluated on $E$ is at most
\begin{equation*}
(\mathfrak{G}-\mathrm{id})\bigl(\Phi(E)+c\bigr)
+\mathfrak{G}\bigl[\epsilon_0\bigl(\Phi(E)+c\bigr)+(1+\epsilon_0)\mathfrak{s}
+3e^{\bar a}C_T^{\mathbf{B}}c\bigr],
\end{equation*}
whose $\omega$-norm is at most
\begin{equation*}
(1-\epsilon_1)^{-1}\Bigl[(\epsilon_0+\epsilon_1+\epsilon_0\epsilon_1)A_E
+\eta_c\,C^{(c)}
+(1+\epsilon_0)C^{(\mathfrak{s})}\Bigr],
\end{equation*}
where $\eta_c$ is the coefficient of the classical load that results from the bounds on
$\mathfrak{G}$ and $\mathfrak{G}-\mathrm{id}$ in part~(iv) of
Proposition~\ref{9.1:prop-load-elimination}; it is the sum of
$\epsilon_0+\epsilon_1+3e^{\bar a}C_T^{\mathbf{B}}$ and of terms of higher order in
$\epsilon_0,\epsilon_1$. Accordingly the eliminated part is bounded at zone $s$ by
\begin{equation*}
\epsilon_{\mathbf{B}}A_E\,\omega+\bigl(1+4C_T^{\mathbf{B}}\bigr)C^{(c)}\mathsf{l}_se_s
+2C^{(\mathfrak{s})}e_s .
\end{equation*}
For the first term we used
$(\epsilon_0+\epsilon_1+\epsilon_0\epsilon_1)/(1-\epsilon_1)\le\epsilon_0+2\epsilon_1
=\epsilon_{\mathbf{B}}$: multiplying out, this inequality becomes
$2\epsilon_1(\epsilon_0+\epsilon_1)\le\epsilon_1$, which holds because
$2(\epsilon_0+\epsilon_1)\le1$, and the latter holds since $\epsilon_0\le1/38$ and
$\epsilon_1\le1/64$ by Lemma~\ref{9.1:lem-apriori-loads}. For the middle term,
$(1-\epsilon_1)^{-1}\eta_c\le1+4C_T^{\mathbf{B}}$: the terms
$\epsilon_0+\epsilon_1+3e^{\bar a}C_T^{\mathbf{B}}$ of $\eta_c$ contribute at most
\begin{equation*}
\frac{64}{63}\Bigl(\frac{1}{38}+\frac{1}{64}\Bigr)+\frac{64}{63}\,3e^{1/40}C_T^{\mathbf{B}}
\le\frac{1}{16}+\frac{13}{4}\,C_T^{\mathbf{B}},
\end{equation*}
by $\bar a\le1/40$ and the same two bounds of Lemma~\ref{9.1:lem-apriori-loads}, and the terms of
higher order in $\epsilon_0,\epsilon_1$ are absorbed in the remaining margin up to
$1+4C_T^{\mathbf{B}}$. For the last term, $(1+\epsilon_0)/(1-\epsilon_1)\le2$
by the same two bounds.

The load enters the $R$-inequality at depth $r$ through the prefactor
$C_t\theta_r4E_0^{-1}e^{-p_0(g_r)}$ multiplying the age-weighted sum, and the age-weighted sum of
$\omega$ is at most $C_\Sigma^{TC}\mathsf{l}_r^2e_r$ by the bound on age-weighted sums in
\eqref{9.1:eq-load-envelope-b}. Hence the extra contribution of the eliminated boundary part to the
kernel evaluation of the $R$-inequality satisfies
\begin{equation*}
C_t\theta_r4E_0^{-1}e^{-p_0(g_r)}\cdot\epsilon_{\mathbf{B}}A_EC_\Sigma^{TC}\mathsf{l}_r^2e_r
\le\epsilon_{\mathbf{B}}K_\Lambda\,\mathfrak{r}_rC_E^{\natural}e_r .
\end{equation*}
Here the definition of $K_\Lambda$ in \eqref{9.1:eq-envelope-corollary-a} has been inserted, and
$C_t\theta_re^{-p_0(g_r)}\mathsf{l}_r^2\le\mathfrak{r}_r$ has been used, with $C_t$ the constant of
the $R'$-prefactor; the latter rests on
$e^{-p_0(g_r)}\mathsf{l}_r^2\le\theta_r^{-\kappa_0/2}$, which the kernel-form corollary of the
response system provides, with its exponent $\kappa_0$, since $q_\flat\ge p_0+1$. The envelope
entry of $g^R$ at depth $r$ is
$E^R_r=4C_E^{\natural}\hat{\mathfrak{r}}_re_r\ge64C_E^{\natural}\mathfrak{r}_r\Lambda_re_r$.
Since $\Lambda\le\Lambda_r$ by (h4),
\begin{equation*}
\epsilon_{\mathbf{B}}K_\Lambda\,\mathfrak{r}_rC_E^{\natural}e_r
\le\frac{\epsilon_{\mathbf{B}}K_\Lambda}{64\Lambda_r}E^R_r
\le\frac{\epsilon_{\mathbf{B}}K_\Lambda}{64\Lambda}E^R_r=\epsilon E^R_r ,
\end{equation*}
and $\epsilon\le4\Lambda/(64\Lambda)=1/16$ by \eqref{9.1:eq-envelope-corollary-a}. The unperturbed
kernel evaluation of the $R$-inequality is
$2\mathfrak{r}_r\Lambda_rB_r(E)\le C_E^{\natural}\hat{\mathfrak{r}}_re_r=\tfrac14E^R_r$, with
$B_r(E)$ the kernel functional of the $R$-inequality evaluated on $E$, by the
definition of $\hat{\mathfrak{r}}_r$, with $\Lambda_r$ unchanged; so the perturbed evaluation is at
most $(\tfrac14+\epsilon)E^R_r\le(\tfrac12+\epsilon)E^R_r$.

The remaining two terms of the eliminated part join the sources. The classical terms join the index
sources. The own sources $\mathfrak{s}$ reach the $R$-inequality only through the prefactor
$C_t\theta_r4E_0^{-1}e^{-p_0(g_r)}\le4E_0^{-1}\mathfrak{r}_r$, and could reach the
$T$-inequalities only with a factor $\varkappa\le1$; by~(i) they do not reach them. The constant of
the index and classical sources therefore becomes $C_\sigma^z$ of
\eqref{9.1:eq-envelope-corollary-4}. The source term of the $R$-inequality at depth $r$ then
satisfies
\begin{equation*}
\sigma=\mathfrak{r}_r\Lambda_r(2C_\sigma^z+1)e_r
\le\frac{\hat{\mathfrak{r}}_r}{16}\,2C_E^{\natural}e_r\le\tfrac12E^R_r ,
\end{equation*}
which needs only $C_E^{\natural}\ge2C_\sigma^z$, one of the members of the maximum defining
$C_E^{\natural}$ in~(iii).

\emph{Proof of (iii).} The side conditions of the five-unknown response proposition are ceilings on
$g$ and properties of $\hat{\mathfrak{r}}$; their forms do not change, so they hold at the new
constants for $\gamma$ small. By~(i), the $P$- and $A$-inequalities keep their kernel evaluations, so
the hypotheses of Lemma~\ref{9.6:lem-maximum-principle-slack} hold there with $\epsilon=0$; by~(ii),
they hold for the $R$-inequality with $\epsilon=\epsilon_{\mathbf{B}}K_\Lambda/(64\Lambda)$, the
sources being at most half the envelope; the boundary entries are those of the five-unknown response
proposition. Lemma~\ref{9.6:lem-maximum-principle-slack}, applied with the envelope at the constant
$C_E^{\natural}$, gives the unknowns at most $E/(1-2\epsilon)$, and $1/(1-2\epsilon)\le8/7$ since
$\epsilon\le1/16$. This is the conclusion of~(iii) with $C_E=\tfrac87C_E^{\natural}$.

\emph{Proof of (iv).} The whole-lattice part of $R''$ is an output of the $T$-site on a set
$X\subset\Lambda_{k'+1}$, so~(i) applies to it. For the $R'$-part, the estimate of~(ii) gives, under
\eqref{9.1:eq-envelope-corollary-a}, the bound
$\epsilon_{\mathbf{B}}C_E^{\natural}\hat{\mathfrak{r}}_re_r$ on its boundary part.

Finally, the stored boundary terms have been eliminated through parts~(iii) and~(iv) of
Proposition~\ref{9.1:prop-load-elimination}, so they are not among the unknowns: the index set $I$
of Lemma~\ref{9.6:lem-maximum-principle-slack} is that of the five-unknown response proposition.
\end{proof}

\smallskip
\noindent The following consequences concern the proof of the two-cutoff comparison theorem and the
statements on which it depends; we continue the lettering of
Corollary~\ref{9.1:cor-envelope-corollary}.

\smallskip
\noindent(c) In the proof of the two-cutoff comparison theorem, the following assertion holds with
the resolved stored terms: the inner history of a large-field region removed by an
$\mathbf{R}$-step is summed inside the stored terms $R'$ and $\mathbf{B}'$ of that step, whose
difference between the two runs is the difference of a stored term assigned to another class;
hence the label $h$ has a live component $X$ at the stop depth with $j(X)\le j+1$, where $j(X)$ is
the step at which $X$ was created and $j$ is the zone fixed at that place of the proof. Indeed, the
difference of the inner history between the two runs is the own part $\mathbf{B}'^{own}$ of the
stored term of that $\mathbf{R}$-step, assigned to the larger of the level of that term and the
block zone, where it is present. The surviving classes consist of the resolved members of
\eqref{9.1:eq-resolved-members-a}, and for these \eqref{9.1:eq-resolved-interpolation-a} gives the
live component directly. The other steps of that proof, and its first display, hold as stated there.

\smallskip
\noindent(d) In the closing lemma of the two-cutoff comparison, the interpolation statement is used
in its resolved form \eqref{9.1:eq-resolved-interpolation-a}, and the summation of boundary terms
at finer zones in its resolved form. This adds one further kind of term occupying slots of several
classes, the resolved boundary members, beside the left-over chain terms $\mathbb{C}$ and the terms
$V(X^{\sharp})$; the resolved boundary members are admissible on the dial domain
$\mathbb{D}^{\star}$ exactly as those terms are, one complex class at a time.

\smallskip
\noindent(e) The load-propagation statement, the zone-splitting estimate of the response comparison
and part~(b) of the crossing-weight lemma hold for the loads of
Corollary~\ref{9.1:cor-envelope-corollary}, with the constant $K_{\mathsf{l}}'$ of the
load-propagation statement replaced by $K_{\mathsf{l}}'(1+K_Y)$ and the constant $K_{\mathbb{C}}$
attached to the chain terms $\mathbb{C}$ changed accordingly. These constants enter only the
ceiling constant $C_L''$, which becomes $C_L'''$, that is, only ceilings on $\gamma$; they do not
enter the coefficient $4E_0^{-1}e^{-p_0}$.

\smallskip
\noindent(f) In place of part~(a) of the crossing-weight lemma, Lemma~\ref{9.1:lem-apriori-loads} is
used, with $E_1^z$ in place of $E_1$ in the definition of
$\mathbb{D}^{\star}$ and in its smallness condition. Part~(c) of the crossing-weight lemma, the
bound $w^{\mathbf{B}}\le 4\varsigma$, is used in this section only for the boundary loads
$\mathsf{y}'$.

\begin{remark}[The own-source bound absorbs the alignment logarithm]\label{9.1:rem-own-source}
Let $\mathfrak{n}$ be the length of the run, let $j$ be a zone and put $s:=\mathfrak{n}-j$; assume
$X\ge e$.
Let $\mathfrak{s}_j$ be the own-source quantity of zone $j$ in the two-run bound for boundary
terms of the large-field step, display~\eqref{9.1:eq-large-field-boundary-a}. Then
\begin{equation}\label{9.1:eq-own-source-a}
\mathfrak{s}_j\le C^{(\mathfrak{s})}e_s,\qquad
C^{(\mathfrak{s})}:=\bigl(C_1^{own}+C_2^{own}\bigr)\,4C_\varsigma\sum_{a\ge 0}(1+a)^3\rho^{a},
\end{equation}
where $C_\varsigma$ is the constant of the bound on the crossing weights $w^{\mathbf{B}}_{j,\ell}$
used in the proof, and $C_1^{own}$ and $C_2^{own}$ are the constants defined in
\eqref{9.1:eq-own-source-2} and \eqref{9.1:eq-own-source-3} below. The constant $C_1^{own}$
depends only on $(q_a,k_*,\mathfrak{a},\kappa_0,\lambda,c_2,\rho,C^{\natural}_{cl})$ and on the
constant of the large-field smallness bound, and not on $\mathfrak{n}$; here $k_*$ and
$\mathfrak{a}$ are the offset and the exponent of the bound on the alignment level,
$C^{\natural}_{cl}$ is the constant of that bound, and $\kappa_0$ is the exponent of the
large-field smallness.
The constant $C^{\natural}_{cl}$ is fixed before $C^{(\mathfrak{s})}$, and none of
$C_1^{own}$, $C_2^{own}$, $C_\varsigma$, $\rho$ depends on $\mathfrak{n}$ or on $X$; hence
$C^{(\mathfrak{s})}$ depends on no cutoff. Thus the polylogarithmic factor produced by the
alignment level is absorbed by the large-field smallness $c_1^{\flat}(g_\ell)$ of each level
$\ell\le j$ that contributes to $\mathfrak{s}_j$, and it need not be placed in the constant
$C_{nc}$ fixed after $g$.
\end{remark}

\begin{proof}
Fix a level $\ell\le j$ and put
\begin{equation*}
s':=\mathfrak{n}-\ell\ \ge\ s=\mathfrak{n}-j,\qquad a:=s'-s=j-\ell\ \ge 0 .
\end{equation*}
At the level $\ell$, the quantity $\mathfrak{s}_j$ collects two terms with the crossing weight
$w^{\mathbf{B}}_{j,\ell}$. The first is $4c_1^{\flat}(g_\ell)\,q_a^{\ell-n_c}$, where $n_c$ is
the alignment level. The second is
\begin{equation*}
2\bigl(9C^{\sharp}_\chi/\vartheta+2\bigr)\,c_1^{\flat}(g_\ell)\,\vartheta^{\ell}.
\end{equation*}
This term comes from the two-run bound \eqref{9.1:eq-large-field-boundary-a}, where it is the
last term, $C^{\sharp}_\chi$ being its constant, and it carries the
classical (test) rate $\vartheta$ at the level $\ell$. We bound each term by a multiple of
$e_{s'}$ and then sum over $\ell$. Throughout we use $\rho<1$.

\emph{First term.} The bound on the alignment level states
\begin{equation}\label{9.1:eq-own-source-1}
n_c\le k_*+1+\frac{\log\bigl(C^{\natural}_{cl}\,(\log(x_0+2c_2))^{c_2}\bigr)}{\log\vartheta^{-1}},
\qquad x_0\le X+3\lambda\mathfrak{n}+2c_2 ,
\end{equation}
where $k_*$ is the fixed offset in that bound and $x_0$ is the inverse coupling at the level $0$
of the run. We use it with the exponent $\mathfrak{a}$ of the alignment-level bound, for which
\begin{equation*}
q_a^{-\log y/\log\vartheta^{-1}}\le y^{\mathfrak{a}/c_2}\qquad (y\ge1).
\end{equation*}
Since $q_a<1$ (shown below), $q_a^{-n_c}$ is increasing in $n_c$. Replacing $C^{\natural}_{cl}$
by $\max\{1,C^{\natural}_{cl}\}$ in \eqref{9.1:eq-own-source-1} only enlarges its right-hand
side, because $\log\vartheta^{-1}>0$; we use the bound on $n_c$ in this weaker form. We apply the
inequality above with $y=\max\{1,C^{\natural}_{cl}\}(\log(x_0+2c_2))^{c_2}$, which is at least
$1$: the lower bound on the running inverse couplings used below, applied with $\ell$ replaced
by $0$, that is at depth $\mathfrak{n}$, gives $x_0\ge X+\lambda\mathfrak{n}-2c_2$, so
$x_0+2c_2\ge X+\lambda\mathfrak{n}\ge e$ and $\log(x_0+2c_2)\ge1$.
We also use $x_0+2c_2\le X+3\lambda\mathfrak{n}+4c_2$,
and note that $\ell=\mathfrak{n}-s'$. Together these give
\begin{equation*}
\begin{split}
q_a^{\ell-n_c}=q_a^{\mathfrak{n}-s'-n_c}
&\le \bigl(\max\{1,C^{\natural}_{cl}\}\bigr)^{\mathfrak{a}/c_2}\,q_a^{-k_*-1}\\
&\quad\times\bigl(\log(X+3\lambda\mathfrak{n}+4c_2)\bigr)^{\mathfrak{a}}\,q_a^{\mathfrak{n}-s'} .
\end{split}
\end{equation*}
We have $q_a\le q^2\le\rho^4$, so $q_a<1$. Hence
\begin{equation*}
q_a^{\mathfrak{n}-s'}\le\rho^{4(\mathfrak{n}-s')}=\rho^{\mathfrak{n}-s'}\,\rho^{3(\mathfrak{n}-s')}.
\end{equation*}

Next we treat the logarithm. Expanding the product gives
\begin{align*}
(X+3\lambda s'+4c_2)\bigl(1+3\lambda(\mathfrak{n}-s')\bigr)
&=X+3\lambda\mathfrak{n}+4c_2\\
&\quad+3\lambda(\mathfrak{n}-s')\,(X+3\lambda s'+4c_2-1).
\end{align*}
The last summand is nonnegative, so
\begin{equation*}
X+3\lambda\mathfrak{n}+4c_2\le (X+3\lambda s'+4c_2)\bigl(1+3\lambda(\mathfrak{n}-s')\bigr).
\end{equation*}
Since $X\ge e$, we have $\log(X+3\lambda s'+4c_2)\ge 1$. Therefore
\begin{align*}
\log(X+3\lambda\mathfrak{n}+4c_2)
&\le \log(X+3\lambda s'+4c_2)+\log\bigl(1+3\lambda(\mathfrak{n}-s')\bigr)\\
&\le \log(X+3\lambda s'+4c_2)\,\bigl(1+\log(1+3\lambda(\mathfrak{n}-s'))\bigr).
\end{align*}
Raising to the power $\mathfrak{a}\ge0$ and inserting the factor $2^{\mathfrak{a}}\ge1$, we obtain
\begin{equation*}
\bigl(\log(X+3\lambda\mathfrak{n}+4c_2)\bigr)^{\mathfrak{a}}
\le 2^{\mathfrak{a}}\bigl(\log(X+3\lambda s'+4c_2)\bigr)^{\mathfrak{a}}
\bigl(1+\log(1+3\lambda(\mathfrak{n}-s'))\bigr)^{\mathfrak{a}} .
\end{equation*}
Let $C_{\mathfrak{a}}$ be the supremum over integers $b\ge0$ of
$\rho^{3b}(1+\log(1+3\lambda b))^{\mathfrak{a}}$. It is finite because $\rho<1$, and
\begin{equation*}
\rho^{3(\mathfrak{n}-s')}\bigl(1+\log(1+3\lambda(\mathfrak{n}-s'))\bigr)^{\mathfrak{a}}\le C_{\mathfrak{a}}.
\end{equation*}

It remains to compare $X+3\lambda s'+4c_2$ with the coupling of the level $\ell$. The
reference-window lower bound on the running inverse couplings in the correlation theorem gives
\begin{equation*}
x_\ell\ \ge\ x_{\mathfrak{n}}+\lambda s'-2c_2\ \ge\ X+\lambda s'-2c_2 .
\end{equation*}
Hence $3x_\ell+10c_2\ge 3X+3\lambda s'+4c_2$, and so
\begin{equation*}
X+3\lambda s'+4c_2\le 3x_\ell+10c_2 .
\end{equation*}
Now use the large-field smallness of the level $\ell$. Let $A_\flat$ be the constant of the
large-field smallness bound, which states $c_1^{\flat}(g_\ell)\le 2A_\flat\,x_\ell^{-\kappa_0/2}$
at every level with $x_\ell\ge1$, $\kappa_0>0$ being its exponent; $A_\flat$ depends neither on
the level nor on the run. This gives
\begin{equation*}
c_1^{\flat}(g_\ell)\bigl(\log(X+3\lambda s'+4c_2)\bigr)^{\mathfrak{a}}
\le 2A_\flat\,x_\ell^{-\kappa_0/2}\bigl(\log(3x_\ell+10c_2)\bigr)^{\mathfrak{a}}
\le 2A_\flat\,C'_{own}.
\end{equation*}
Here $C'_{own}$ is the supremum of $x^{-\kappa_0/2}(\log(3x+10c_2))^{\mathfrak{a}}$ over
$x\ge1$, applied at $x=x_\ell$; and $x_\ell\ge1$ throughout the reference window.
It is finite because the power $x^{-\kappa_0/2}$ dominates the logarithm as
$x\to\infty$. Define
\begin{equation}\label{9.1:eq-own-source-2}
C_1^{own}:=8A_\flat\,2^{\mathfrak{a}}\,\bigl(\max\{1,C^{\natural}_{cl}\}\bigr)^{\mathfrak{a}/c_2}\,
q_a^{-k_*-1}\,C_{\mathfrak{a}}\,C'_{own}.
\end{equation}
Collecting the displays of this step, we get
\begin{equation*}
4c_1^{\flat}(g_\ell)\,q_a^{\ell-n_c}\ \le\ C_1^{own}\rho^{\mathfrak{n}-s'}\ \le\ C_1^{own}e_{s'},
\end{equation*}
where the last inequality holds by the definition of the envelope $e_{s'}$. By
\eqref{9.1:eq-own-source-2}, $C_1^{own}$ depends on
$(q_a,k_*,\mathfrak{a},\kappa_0,\lambda,c_2,\rho,C^{\natural}_{cl})$ and on $A_\flat$ only,
and not on $\mathfrak{n}$.

\emph{Second term.} The classical and test rates are powers of $L$, each at most
$\vartheta_{dl}=L^{-\theta_2/2048}$. The window
for $\rho$ requires $\rho\ge\max\{\vartheta_{dl},\vartheta_\chi\}$. Hence
$\vartheta\le\vartheta_{dl}\le\rho$. Since $\ell=\mathfrak{n}-s'$,
\begin{equation*}
\vartheta^{\ell}\le\rho^{\mathfrak{n}-s'}\le e_{s'} .
\end{equation*}
Define
\begin{equation}\label{9.1:eq-own-source-3}
C_2^{own}:=4\bigl(9C^{\sharp}_\chi/\vartheta+2\bigr)A_\flat .
\end{equation}
Since $x_\ell\ge1$, the bound of the first step gives
$c_1^{\flat}(g_\ell)\le 2A_\flat x_\ell^{-\kappa_0/2}\le 2A_\flat$. Then
\begin{equation*}
2\bigl(9C^{\sharp}_\chi/\vartheta+2\bigr)\,c_1^{\flat}(g_\ell)\,\vartheta^{\ell}\le C_2^{own}e_{s'}.
\end{equation*}

\emph{Sum.} By the monotonicity of the envelope in \eqref{9.1:eq-load-envelope-a},
$e_{s'}\le\rho^{-a}e_s$. The crossing weights satisfy
$w^{\mathbf{B}}_{j,\ell}\le 4C_\varsigma(1+a)^3q^{a}$, and $q/\rho\le\rho$. So the contribution
of the level $\ell$ to $\mathfrak{s}_j$ is at most
\begin{equation*}
4C_\varsigma(1+a)^3q^{a}\,\bigl(C_1^{own}+C_2^{own}\bigr)\rho^{-a}e_s
\le\bigl(C_1^{own}+C_2^{own}\bigr)\,4C_\varsigma(1+a)^3\rho^{a}e_s .
\end{equation*}
Distinct levels $\ell\le j$ give distinct values $a=j-\ell\ge0$. Summing over $\ell\le j$
therefore gives
\begin{equation*}
\mathfrak{s}_j\le\bigl(C_1^{own}+C_2^{own}\bigr)\,4C_\varsigma\sum_{a\ge0}(1+a)^3\rho^{a}\,e_s
=C^{(\mathfrak{s})}e_s,
\end{equation*}
and the series converges because $\rho<1$. This is \eqref{9.1:eq-own-source-a}.
\end{proof}

\begin{definition}[Cubes, relative tree length and strips at the large-field step]
\label{9.1:not-cubes-strips}
\emph{Levels.} The letters $i$, $\ell$, $m$, $k'$ denote levels. The $m$-th renormalisation step
consists of the operation $\mathbf{T}_m$ followed by the operation $\mathbf{R}_m$, in this order,
as in \cite{B-CMP119}.

\emph{Cubes at level $m$.} Let $\pi_m$ be the partition of the level-$m$ lattice into cubes of
side $M$ used by Ba{\l}aban's construction at level $m$ \cite{B-CMP119}. At level $m$ a
\emph{cube} is a closed $M$-cube of $\pi_m$, and lengths are measured in units of $M$. For unions
of cubes $A$, $B$ we write $\operatorname{dist}_m(A,B)$ for their $\infty$-distance (maximum over
coordinates) measured in cubes; it is an integer, and $\operatorname{dist}_m(A,B)=0$ when $A$ and
$B$ touch or overlap. For a set $A$, $\mathsf{cl}_m(A)$ is the union of the level-$m$ cubes meeting
$A$, and $N_m(A)$ is the number of level-$m$ cubes in $\mathsf{cl}_m(A)$. The partitions are nested:
the partition at each level is chosen, in the words of \cite{B-CMP119}, ``compatible with all the
previous partitions''.
Consequently
\begin{equation*}
\mathsf{cl}_{m'}\circ\mathsf{cl}_m=\mathsf{cl}_{m'}\qquad\text{for } m'\ge m .
\end{equation*}
For a set $X$ we put $\bar X^{(m)}:=\mathsf{cl}_m(X)$; it is the smallest domain of the class
$\mathbf{D}_m$ containing $X$. Here $\mathbf{D}_m$ is the class of localisation domains at level
$m$ and $d_m$ the tree length on it, both as fixed in the lemma on the localisation
domains at level $m$ and their tree length.

\emph{Relative tree length.} Let $W$ be a union of level-$m$ cubes and $Y$ a domain at level $m$.
For a tree graph $\tau$ write $|\tau|$ for its length (in units of $M$). We put
\begin{equation}\label{9.1:eq-cubes-strips-1}
d_{m,W}(Y):=
\begin{cases}
\min_\tau|\tau| , & Y\not\subset W,\\[2pt]
0, & Y\subset W,
\end{cases}
\end{equation}
the minimum being taken over tree graphs $\tau\subset Y$ that meet every level-$m$ cube
$\Delta\subset Y$ with $\Delta\not\subset W$. This is the length of \cite[(1.67)]{B-CMP122b} with the
set $Z$ there replaced by $W$; in \cite[(1.67)]{B-CMP122b} it is described as ``the length of a
shortest tree graph contained in $Y$ and intersecting all $M$-cubes of the components of
$Y\setminus Z$''. The length of the tree is counted wherever the tree runs, inside $W$ included;
only the requirement that the tree visit the cubes of $W$ is removed.

\emph{Regions at the operation $\mathbf{R}_m$.} With $\Lambda_m$ as in
Notation~\ref{9.1:not-zones-boundary-terms}, $Z_m=\Lambda_m^c$ is the large-field region at
level $m$, as there and as in \cite{B-CMP122a} and \cite{B-CMP119}. We write $Z(m)\subset Z_m$ for
the union of the components of $Z_m$ of the first two classes of the classification at the operation
$\mathbf{R}_m$ in \cite{B-CMP122a}, and $\{X_j\}$, $\{Y_\iota\}$ for its components of the first and
of the second class respectively, as in \cite{B-CMP122b}.

\emph{Geometry of the large-field region.} Throughout, the large-field region is assumed to have
the following box geometry, a standing convention of the strip construction: each component of
$\Lambda_m^c$ is a box of sides $a_\nu\check{R}_m$, $\nu=1,\dots,d$, with
$a_\nu\in\{1,\dots,100\}$, and distinct components are at distance at least $\check{R}_m$ from each
other. Here, for every level $k$, $\check{R}_k$ is Ba{\l}aban's radius $R_k$ of \cite{B-CMP122a} at
level $k$, evaluated at the real running coupling $g_k$.

\emph{Strips.} Put
\begin{equation*}
m_\sharp:=m_{hi}(m)=\bigl\lfloor \tfrac12(\check{R}_m-1)\bigr\rfloor ,
\end{equation*}
where $m_{hi}(m)$ is the upper strip index of the dictionary to the strip. For a component $C$ of
the large-field region, its strip is $C^{\sharp}:=C^{(m_\sharp)}$, the set attached to $C$ at
index $m_\sharp$ by the dictionary to the strip; there also the union $Z^{\sharp}$ of the strips,
the filing map $\operatorname{fil}_\sharp$ to strip domains and the tree length $d_\sharp$
relative to the strips are defined. The output length is the relative tree length
\eqref{9.1:eq-cubes-strips-1}, at the level at which it is read, with $W$ the union of the strips of
the second-class components:
\begin{equation*}
d_{out}:=d_{\bigcup_\iota Y_\iota^{\sharp}} .
\end{equation*}
The combinatorial constants are
\begin{equation}\label{9.1:eq-cubes-strips-2}
C_{fib}:=2e^{2\mathsf{k}_d},\qquad
\mathsf{k}_d:=1+\mathsf{c}_d\bigl(2d\log 3+(1+3^d)\log 2\bigr),\qquad
\mathsf{c}_d:=2\cdot 3^d .
\end{equation}

\emph{Rates.} The decay rate of the stored boundary terms is $\kappa_{\mathbf{B}}:=\kappa$: by the
bound for the terms stored by the operation $\mathbf{T}$, their decay rate is at least $\kappa$,
and it is used at the value $\kappa$. With the step letter $\beta_s=\frac14$, the combination step
under the strip delivers the rate $(1+2\beta_s)\kappa$; its factor is split as
\begin{equation*}
1+2\beta_s=(1+\beta_s)+\tfrac12\beta_s+\tfrac12\beta_s ,
\end{equation*}
and the step of \cite[(1.98)]{B-CMP122b} then gives \cite[(1.99)]{B-CMP122b} at the rate
$(1+\frac12\beta_s)\kappa$; this is the rate written $(1+\frac12\beta)\kappa$ in
\cite[(1.99)]{B-CMP122b}. Hence the operation $\mathbf{R}$ produces its terms of the first group,
recalled after \eqref{9.1:eq-cubes-strips-6}, at the rate
\begin{equation*}
\bigl(1+\tfrac12\beta_s\bigr)\kappa=\tfrac98\,\kappa .
\end{equation*}
By the lemma on the output rate at the $\mathbf{T}$-site,
the output rate there satisfies $\kappa_{out}\ge 2(\mathfrak{c}+1)\kappa$, where $\mathfrak{c}$ is
the constant of that lemma.

\emph{The coefficient of the $\mathbf{R}$-step.} With $p_0(\cdot)$ the degree function and
$O(1)$ the numerical constant of \cite[(1.99)]{B-CMP122b}, we put
\begin{equation*}
c_1^{\flat}(g_\ell):=O(1)\max\bigl\{e^{-p_0(g_\ell)},\,(\alpha^{\sharp}_\ell)^{1/3}\bigr\},
\end{equation*}
where $\alpha^{\sharp}_\ell$ is the constant $\alpha^{\sharp}$ of the value family under the strip,
taken at level $\ell$.

\emph{The radius from step to step.} The radius satisfies
\begin{equation}\label{9.1:eq-cubes-strips-3}
\check{R}_{m+1}\ge \check{R}_m/L .
\end{equation}
Indeed, in Ba{\l}aban's construction the radius is kept from one step to the next except that,
in the words of \cite{B-CMP122a}, ``in some steps the number $R_k$ decreases by the factor
$L^{-1}$''.

\emph{The bounds of the operation $\mathbf{R}$ in \cite{B-CMP122b}.} We recall the displays of
\cite[pp.~389--391]{B-CMP122b} that the above notation serves; in them Ba{\l}aban's level $k$ plays
the role of $m$, and $\beta$, $\kappa$, $\alpha$, $c_0$, $c_1$ are the parameters and constants of
\cite{B-CMP122b}; in particular $c_0$, $c_1$ there are not the constants $c_0$, $c_1$ of
Notation~\ref{9.1:not-zones-boundary-terms}. In these displays $X$, $X'$, $Y_i$ are localisation
domains, and $i$ indexes the second-class components, as in \cite{B-CMP122b}. For the quantities
$F(X')$ of \cite[(1.95)]{B-CMP122b}, that display reads
\begin{equation}\label{9.1:eq-cubes-strips-4}
\begin{split}
|F(X')|\le{}&\sum_q\ \sideset{}{'}\sum_{\{X_{j_1},\dots,X_{j_q}\}}
\exp\bigl(-q\tfrac32\,p_0(g_k)\bigr)\,c_0\,
\exp\Bigl(-(1+\beta)\kappa\, d_{k,\cup_i Y_i}(X')\Bigr)\\
&\qquad\times\exp\Bigl(-\tfrac{1}{6\cdot 2^d}\,\beta\kappa M^{-d}\,
\bigl|X'\setminus\textstyle\bigcup_i Y_i\bigr|\Bigr),
\end{split}
\end{equation}
and \cite[(1.97)]{B-CMP122b} states
\begin{equation}\label{9.1:eq-cubes-strips-5}
|F(X')|\le c_1\exp\Bigl(-(1+\beta)\kappa\, d_{k,\cup_i Y_i}(X')\Bigr),
\end{equation}
``where $c_1=\exp(-p_0(g_k))$ if $X'$ does not intersect $\bigcup Y_i$, and $c_1=\alpha^{1/3}$ in the
remaining cases''. There, each domain $X'$ either contains one of the $Y_i$ or is disjoint from
$\bigcup_i Y_i$ \cite[pp.~389--391]{B-CMP122b}. After the step of \cite[(1.98)]{B-CMP122b} the
terms $\mathbf{R}'^{(k)}(X,(\mathbf{U},\mathbf{J}))$ are bounded in \cite[(1.99)]{B-CMP122b}:
\begin{equation}\label{9.1:eq-cubes-strips-6}
\bigl|\mathbf{R}'^{(k)}\bigl(X,(\mathbf{U},\mathbf{J})\bigr)\bigr|
\le O(1)\,c_1\exp\Bigl(-\bigl(1+\tfrac12\beta\bigr)\kappa\, d_{k,\cup_iY_i}(X)\Bigr),
\end{equation}
where $d_{k,\cup_iY_i}$ is the relative tree length \eqref{9.1:eq-cubes-strips-1} at level $k$ with
$W=\bigcup_iY_i$. The terms are then divided into two groups \cite[(1.99)--(1.100)]{B-CMP122b}.
The first group consists of ``all the terms with the localization domains $X$ intersecting the large
field region $Z_k^{\sim}\cup\bigcup Y_i^{\sim}$''; these ``are new boundary terms $\dots$
$\mathbf{B}'^{(k)}(X)$''. For the second group only, \cite{B-CMP122b} states: ``By their definition
the linear size of the domain $X$ in (1.99) can be replaced by $d_k(X)$, and $c_1$ is equal to
$\exp(-p_0(g_k))$'', which gives \cite[(1.100)]{B-CMP122b}.

The operation $\mathbf{T}$ at level $k$ is given in \cite[(1.102)]{B-CMP122b} by
\begin{equation*}
\mathbf{T}_k(Y)=\chi_k^c\bigl(Y^{\sim-6}\bigr)\,\chi_{k,\Lambda}\,
\delta_G\bigl(V_kV_\Lambda^{-1}\bigr)\,\mathbf{T}'_k(Y),
\end{equation*}
in the notation of \cite{B-CMP122b}, and there ``The new large field region is equal to
$Z_k\cup\bigcup Y_i$''.

In the alternative representation of \cite[(1.103)--(1.104), p.~391]{B-CMP122b}, ``the first
product over $h$ is replaced by the product of $e^3$ over indices $i$ such that $Y_i\subset X'$, the
number $\beta$ is replaced by $(1/4)\beta$, and $\alpha$ by $O(1)c_1$ $\dots$ We have the inequality
(1.93), but with the linear size $d_k(X'\setminus\cup Y_i)$ $\dots$ This requires a comment, because
in this bound there are no tree decay exponential factors for the domains $Y_i$. In fact $\dots$
such factors are not needed, because there are no summations over these domains, they are fixed.''

Thus in both representations of \cite{B-CMP122b} a first-group term whose domain contains a
second-class component $Y_i$ is bounded with tree decay off $Y_i$ only, and the use of such terms
at later steps rests on the fact that no summation runs over the domains $Y_i$, which are fixed.
This is the alternative inductive form described in \cite{B-CMP119}: ``we may resum all the terms
with localization domains intersecting a given component of the large field region $\dots$ we
obtain the bound (2.42) with $d_j(X)$ replaced by $d_j(X\setminus Z_j)$. All constructions and proofs
of the procedure can be carried on with these inductive assumptions.''
\end{definition}

We use the level-$m$ cubes, the cube distance $\operatorname{dist}_m$ (an integer, equal to $0$ when
two sets touch or overlap), the closure $\mathsf{cl}_m$ in level-$m$ cubes and the strips of
Notation~\ref{9.1:not-cubes-strips}. By $\operatorname{dist}$ without subscript we mean the
$\infty$-distance in level-$m$ lattice units. In these units a level-$m$ cube has side $M$ and a
level-$(m+1)$ cube has side $LM$. For a union $A$ of level-$m$ cubes and an integer $n\ge0$ we write
$A^{(n)}$ for the union of the level-$m$ cubes at cube distance at most $n$ from $A$, that is, $A$
with $n$ layers of level-$m$ cubes added.

\emph{Levels.} Level $m$ refers to the geometry of the $m$-th effective density. Step $m$ consists
of the renormalisation transformation $T_m$ of \cite{B-CMP119} followed by Ba{\l}aban's large-field
operation $\mathbf{R}_m$. The transformation $T_m$ produces domains $\Omega_m\supset\Lambda_m$
\cite{B-CMP119}, and we put $Z_m:=\Lambda_m^c$. We write $Z(m)\subset Z_m$ for the healable class at
$\mathbf{R}_m$ and $Z^{\sharp}(m)$ for the union of its strips. By \cite{B-CMP122b}, each component
of $Z(m)$ is of first class (it is healed by $\mathbf{R}_m$) or of second class. After $\mathbf{R}_m$
the large-field region is
\begin{equation*}
Z_m^{+}:=\bigl(\text{the unhealed part of }Z_m\bigr)\cup\bigcup_\iota Y_\iota ,
\end{equation*}
where the $Y_\iota$ are the new regions. \cite[p.~390]{B-CMP122b} states that the new large-field
region equals $Z_k\cup\bigcup_i Y_i$, $k$ being the level there; healed components are not part of
it.

The transformation $T_{m+1}$ is applied to a term whose large-field region, $Z_k=\Lambda_k^c$ in the
notation of \cite[\S3]{B-CMP119}, is $Z_m^{+}$. The construction proceeds as follows:
\begin{itemize}
\item[(1)] it forms the union $Z'$ of the $LM\check R_{m+1}$-cubes of the level-$m$ lattice that
meet $Z_m^{+}$ (Ba{\l}aban's $R_{m+1}$ is $\check R_{m+1}$);
\item[(2)] it surrounds $Z'$ by four layers of such cubes, which gives the domain $Z'^{\sim4}$;
\item[(3)] it introduces a decomposition of unity for $(Z'^{\sim4})^c$ and sums over domains
$P_{m+1}\subset(Z'^{\sim4})^c$ \cite[\S3]{B-CMP119}.
\end{itemize}
By \cite[(3.5)]{B-CMP119},
\begin{equation*}
\Omega_{m+1}\subset(Z'^{\sim5})^c .
\end{equation*}
Since $Z'^{\sim4}\subset Z'^{\sim5}$, in particular $\Omega_{m+1}\subset(Z'^{\sim4})^c$; only these
four layers are used below. Moreover $\Lambda_{m+1}\subset\Omega_{m+1}$, as at every level.

In level-$m$ units one layer of $LM\check R_{m+1}$-cubes has thickness $LM\check R_{m+1}$. This is
at least $M\check R_m$, by the one-step comparison of scaffold radii
$\check R_m\le L\check R_{m+1}$. Since $Z_m^{+}\subset Z'$, we obtain
\begin{equation}\label{9.1:eq-island-geometry-a}
\operatorname{dist}\bigl(\Omega_{m+1},Z_m^{+}\bigr)\ \ge\ 4LM\check R_{m+1}\ \ge\ 4M\check R_m
\quad\text{(level-$m$ units)}.
\end{equation}
Thus $\Omega_{m+1}$ is at least $4\check R_m$ level-$m$ cubes away from $Z_m^{+}$.

\emph{Births, strips, islands.} A \emph{birth} is a pair $(C,i)$ in which $C$ is a second-class
component of $Z(i)$ at $\mathbf{R}_i$ \cite{B-CMP122b}. Its \emph{birth strip} is
\begin{equation*}
s(C,i):=C^{\sharp}=C^{(m_\sharp(i))},\qquad
m_\sharp(i)=\bigl\lfloor\tfrac12(\check R_i-1)\bigr\rfloor .
\end{equation*}
By the strip geometry of the large-field step, $s(C,i)$ is a box of level-$i$ cubes with
$N_i(s(C,i))\le(101\check R_i-1)^d$, and it lies inside the collar $C^{\boxplus}$.

The birth is \emph{live at level $m\ge i$} if at no $\mathbf{R}_{m'}$ with $i<m'\le m$ the component
of $Z_{m'}$ containing $C$ was of first class (that is, healed; it then belongs to $Z(m')$). We put
\begin{equation*}
\mathfrak{S}_m:=\bigcup\{s(C,i):\ i\le m,\ (C,i)\ \text{live at }m\},\qquad
\mathbb{W}_m:=\bigl[\mathsf{cl}_m(\mathfrak{S}_m)\bigr]^{(4)}.
\end{equation*}
So $\mathbb{W}_m$ is the closure of $\mathfrak{S}_m$ in level-$m$ cubes with four layers of level-$m$
cubes added. The \emph{islands at level $m$}, forming the family $\mathbb{I}_m$, are the
touch-components of $\mathbb{W}_m$. They are pairwise non-touching by definition. An island may be a
cluster of the closures of several strips; no bound on its size or on the number of strips in it is
claimed or used.

We also put
\begin{equation*}
\mathbb{W}_m^{-}:=\bigl[\mathsf{cl}_m\bigl(\textstyle\bigcup\{s(C,i):\ i\le m-1,\ (C,i)
\ \text{live at }m\}\bigr)\bigr]^{(4)}\subset\mathbb{W}_m .
\end{equation*}
These are the islands seen by the outputs of $T_m$, before the births of $\mathbf{R}_m$. We write
$\mathbb{I}_m^{-}$ for the touch-components of $\mathbb{W}_m^{-}$.

The four spare layers of $\mathbb{W}_m$ absorb the one-layer fattenings of the domain at which a term
is evaluated, performed at the three places of part (d) below: the fattened core at the
renormalisation transformation, the hull at the resummation, and the one-layer shell at the
preparatory integrations. Throughout we assume
\begin{equation}\label{9.1:eq-island-geometry-b}
L\ge5,\qquad\qquad \check R_m\ge9\quad\text{for every level }m .
\end{equation}
The first condition ensures that four consecutive layers of level-$m$ cubes have total width
$4/L<1$ level-$(m+1)$ cube; it is a lower bound on $L$ weaker than those already imposed in the
order of choice. The second is a
lower bound on the scaffold radius. Since $\check R_m\ge T_\gamma^{\,r_{pr}}$ by the lower bound on
the scaffold radius in the resummation proposition of the large-field step, with $T_\gamma$ the
degree variable at the coupling ceiling and $r_{pr}$ the exponent in that lower bound, it holds as
soon as $T_\gamma\ge3$. Both conditions are one-sided: the first is a lower bound on $L$, the second a
lower bound on $T_\gamma$, that is, an upper bound on the coupling ceiling; neither is a cap in the
sense of Definition~\ref{2.3:def-cap}.

\emph{Stored terms.} Stored boundary terms of level $\ell$ are measured in the island-blind format
$\mathbf{D}_\ell^{\mathbb{I}}$ of this chapter, by a relative tree length
$d^{\circ}_\ell$. This length does not
charge a union of islands of $\mathbb{I}_\ell$, the term's \emph{blind set}; we say the term is
\emph{blind over} these islands. The localisation domain of a stored term is the domain assigned to
it by the filing map $\operatorname{fil}_\ell$ of that format.

\begin{lemma}[Moat and position of the islands around live strips]\label{9.1:lem-island-geometry}
Assume \eqref{9.1:eq-island-geometry-b}. Let $(C,i)$ be a birth and $s:=s(C,i)$, live at the levels
named below. Let $I$ denote an island of $\mathbb{I}_m$.
\begin{itemize}
\item[(a)] \emph{Moat.} For every $m\ge i$ at which $(C,i)$ is live,
\begin{equation*}
\operatorname{dist}\bigl(\mathsf{cl}_m(s),\Omega_{m+1}\bigr)\ge 3LM\check R_{m+1}
\quad\text{(level-$m$ units)}.
\end{equation*}
That is, $\mathsf{cl}_m(s)$ is at least $3L\check R_{m+1}\ge3\check R_m$ level-$m$ cubes from
$\Omega_{m+1}$. Hence, for every $I\in\mathbb{I}_m$,
\begin{equation*}
\operatorname{dist}_m(I,\Omega_{m+1})\ \ge\ 3\check R_m-4\ \ge\ 23 .
\end{equation*}
In level-$(m+1)$ cubes,
\begin{equation*}
\operatorname{dist}_{m+1}\bigl(\mathsf{cl}_{m+1}(I),\Omega_{m+1}\bigr)\ \ge\ 3\check R_{m+1}-2 .
\end{equation*}
In particular no island meets any of the following three sets: $\Lambda_{m+1}\subset\Omega_{m+1}$;
the region $Y_0\subset\Lambda_{m+1}$ of the decoupling lemma at $T_{m+1}$; the support of the
fluctuation field of step $m+1$ \cite{B-CMP119}.
\item[(b)] \emph{Inside.} For $m\ge i+1$,
\begin{equation*}
\mathsf{cl}_m(s)^{(4)}\subset\Omega_m^c\subset Z_m .
\end{equation*}
This set lies inside the component $\mathcal{C}_m(C)$ of $Z_m$ containing $C$, at distance at least
$3\check R_m-5$ cubes from $\Lambda_m$. Consequently:
\begin{itemize}
\item[(b1)] an island all of whose births are at levels $\le m-1$ lies inside one live component of
$Z_m$, at depth at least $3\check R_m-5$;
\item[(b2)] an island containing strips born at $m$ is the union of the sets $s'^{(4)}$, $s'$
running over those strips, and of the
older islands inside their components. There are several such strips if the strips of adjacent
second-class components lie within $8$ cubes of one another. The set $s(C',m)^{(4)}$ contributed by
each new strip lies in the collar $C'^{\boxplus}$ of its component $C'$, because
$m_\sharp+4\le\check R_m$ by \eqref{9.1:eq-island-geometry-b};
\item[(b3)] from level $m+1$ on, all components met by one island are one component of $Z_{m+1}$:
their $4LM\check R_{m+1}$-fattenings overlap, by \eqref{9.1:eq-island-geometry-a}. So the births of
an island are healed together or not at all.
\end{itemize}
\item[(c)] \emph{The partition at $\mathbf{R}_m$.} Let $C'$ be a component of $Z(m)$.
\begin{itemize}
\item[(c1)] Let $I\in\mathbb{I}_m^{-}$ be an island all of whose births have their components
contained in a component $C''$ of $Z_m$ that is not in $Z(m)$. Then
$I\subset C''$, by (b), and $\operatorname{dist}_m(I,C')\ge\check R_m$. Every cube of $I$ has layer
index $\operatorname{lay}\ge\check R_m+1>m_\sharp+1$ with respect to $Z(m)$, in the sense of
Notation~\ref{9.1:not-cubes-strips}. Moreover
\begin{equation*}
\operatorname{dist}_m(I,C'^{\sharp})\ \ge\ \check R_m-m_\sharp\ \ge\ \tfrac12(\check R_m+1)\ \ge\ 5 .
\end{equation*}
\item[(c2)] Let $I\in\mathbb{I}_m^{-}$ lie inside $C'$, and let $C'$ be of first class. Then
$I\subset C'\subset C'^{\sharp}=X^{\sharp}_j$, the healed box. Every stored term whose localisation
domain contains $I$ meets $C'$, and is therefore one of the terms that $\mathbf{R}_m$ absorbs into the
activities of \cite[(1.91)]{B-CMP122b}. After $\mathbf{R}_m$ the births inside $C'$ are not live,
and no stored term is blind over them.
\item[(c3)] Let $I\in\mathbb{I}_m^{-}$ lie inside $C'$, and let $C'$ be of second class. Then
$I\subset C'\subset s(C',m)$, and $I$ is swallowed by the island of $\mathbb{I}_m$ that contains
$s(C',m)^{(4)}$.
\end{itemize}
\item[(d)] \emph{Cells of the three places never lie in an island.}
\begin{itemize}
\item[(d1)] At the renormalisation transformation $T_{m+1}$, each of the following lies in
$\Omega_{m+1}$: every decoupling cube $\Delta\subset Y_0$ of the decoupling lemma, every cube
carrying fluctuation field, and every origin of a
random-walk piece. Hence each of them is at least $3\check R_m-4$ cubes from $\mathbb{W}_m$, by (a).
\item[(d2)] At the resummation of $\mathbf{R}_m$, the exterior accounting of the resummation
proposition of the large-field step is anchored at cubes off $Z(m)$. This accounting consists of the
proposition's two exterior input bounds, taken in the norm $\|\cdot\|_a$, with the suprema of the
filing taken over the anchor cubes $\Delta\not\subset Z(m)$. The relative resummation lemma moves
these anchors off $Z(m)\cup\mathbb{W}'_m$, where
\begin{equation*}
\mathbb{W}'_m:=\bigcup\{I\in\mathbb{I}_m^{-}:\ I\not\subset Z^{\sharp}(m)\}.
\end{equation*}
Such an anchor exists on every exterior domain, by that lemma. The interior accounting of the
proposition is a value budget per unit cell: for every unit cell $C_{cell}$,
\begin{equation*}
\sum_{Y_4\subset C_{cell}}|\mathfrak{v}(Y_4)|\ \le\ \mathsf{K}_{cell}(k)\sum_j
|\Gamma^0_j\cap C_{cell}| ,
\end{equation*}
where $k=m$ is the level of the step and $Y_4$ runs over the domains of the value family
$\mathfrak{v}$ of the proposition. This budget is served by the island-blind value lemma.
\item[(d3)] At the preparatory integrations of $\mathbf{R}_m$ (with $k=m$), the integrations run over
the variables localised in $Z_{j+1}\setminus Z''_{j+1}$, for $j=k-N_{\mathrm{pr}},\dots,k-1$
\cite{B-CMP122a}. Here $N_{\mathrm{pr}}$ (the integer $N$ of \cite{B-CMP122a}) is the number of
preceding renormalisation steps in \cite[condition~(ii)]{B-CMP122a}. The shell
$Z_{j+1}\setminus Z''_{j+1}$ is one (scaled) layer of
$M$-cubes at $\partial Z_{j+1}$ inside the components $C'$ of $Z(m)$ \cite{B-CMP122a}. An island
lying in another component of $Z_{j+1}$ is off the shell: by (c1), at least $\tfrac12(\check R_m+1)$
cubes away, when $j+1=m$ and that component is not in $Z(m)$, and by (b) and the separation of
components otherwise. Consider an island inside
$C'$ that belongs to $\mathbb{I}_{j+1}$ if $j+1<m$ and to $\mathbb{I}_m^{-}$ if $j+1=m$; every
birth of it is at a level $i'\le m-N_{\mathrm{pr}}-1$. Therefore, at level $j+1\ge i'+1$, for each
of its strips $s'$ the fattened closure $\mathsf{cl}_{j+1}(s')^{(4)}$ lies at least
$3\check R_{j+1}-5\ge2$ level-$(j+1)$ cubes inside
$Z_{j+1}$, off the shell, and so the whole island lies off the shell. Hence a cell of the preparatory
integrations at scale $j+1$ is off every
island of $\mathbb{I}_{j+1}$ if $j+1<m$, and off every island of $\mathbb{I}_m^{-}$ if $j+1=m$, the
islands of $\mathbb{I}_m^{-}$ being those seen by the outputs of $T_m$, before the births of
$\mathbf{R}_m$. The treatment by the preparatory integrations of a stored term younger
than the shell at which it is evaluated (of level $\ell>j+1$) is not part of this statement; it is
governed by the class assignment of stored terms in the relative expansion at the preparatory
integrations.
\end{itemize}
\item[(e)] \emph{Absorbed or alive.} A stored term blind over an island exists at level $m$ only if
all births of that island are live at $m$. From the level after the youngest birth on, these births
share one component, by (b). At the first $\mathbf{R}_{m'}$ that heals this component, the stored
term is absorbed, by (c2): its localisation domain contains the island and hence meets the
component. While the term exists, $\mathsf{cl}_m$ of its blind set lies in $\mathbb{W}_m$. Indeed,
let $\ell$ be its level, let $m>\ell$, let $\mathcal{B}$ be its blind set, and let
$\mathfrak{S}'_\ell$ be the union of the strips of $\mathfrak{S}_\ell$ whose births are live at $m$;
then
\begin{equation*}
\mathsf{cl}_m(\mathcal{B})\ \subset\ \mathsf{cl}_m\bigl([\mathsf{cl}_\ell(\mathfrak{S}'_\ell)]^{(4)}\bigr)
\ \subset\ [\mathsf{cl}_m(\mathfrak{S}_m)]^{(1)}\ \subset\ \mathbb{W}_m ,
\end{equation*}
since four layers of level-$\ell$ cubes have total width less than one level-$m$ cube for $m>\ell$
and $L\ge5$.
\item[(f)] \emph{Healing and re-birth dominate.} Consider the $\mathbf{R}_{m'}$ that heals (first
class) or re-strips (second class) the component $\mathcal{C}$ of $C$. Every stored term blind over
an island $I\subset\mathcal{C}$ and meeting $Z(m')^{\boxplus}$ is resummed by the resummation
proposition of the large-field step, with input norm relative to $Z(m')\supseteq\mathcal{C}\supseteq I$.
This is the length $d_{k,Z}$ in $\|\cdot\|_a$, with $k=m'$ and $Z=Z(m')$. The term's own relative
decay dominates this norm, by the strip geometry of the large-field step: for the domain $Y_1$ of the
term and the hull $K_v$ of its value unit, in
the sense of the resummation proposition, the tree length of $Y_1$ relative to
$Z(m')\cup\mathbb{W}'_{m'}$ is at most the island-blind length $d^{\circ}_\ell(K_v)$, $\ell$ being
the level of the term:
\begin{equation*}
d_{m',Z(m')\cup\mathbb{W}'_{m'}}(Y_1)\ \le\ d^{\circ}_\ell(K_v).
\end{equation*}
Inside the combinatorial lemma of the resummation, a healed box
$X^{\sharp}_j\supseteq\mathcal{C}^{\sharp}\supseteq I$ is crossed by a Hamiltonian path that is paid
for by \cite[(1.89)]{B-CMP122b}. Hence nothing blind survives a healing. At a re-birth the old island
lies inside the new strip, by (c3). Blindness over an island therefore costs nothing at the step
where the island's component is operated on.
\end{itemize}
\end{lemma}

\begin{proof}
\emph{(a)} We argue by induction on $m\ge i$, over the levels at which $(C,i)$ is live.

\emph{Base case $m=i$.} The strip $s$ is a union of level-$i$ cubes, so $\mathsf{cl}_i(s)=s$, and
$s$ lies within $m_\sharp=m_\sharp(i)$ level-$i$ cubes of $C$. Since $C$ is of second class, it is not
healed, and $C\subset Z_i^{+}$. By \eqref{9.1:eq-island-geometry-a} at level $i$,
\begin{align*}
\operatorname{dist}(\Omega_{i+1},s)
&\ge 4LM\check R_{i+1}-m_\sharp M
\ \ge\ 4LM\check R_{i+1}-\tfrac12\check R_i M\\
&\ge \bigl(4L-\tfrac12L\bigr)M\check R_{i+1}
\ \ge\ 3LM\check R_{i+1}.
\end{align*}
Here we used $m_\sharp\le\frac12(\check R_i-1)\le\frac12\check R_i$ and, in the third inequality,
$\check R_i\le L\check R_{i+1}$.

\emph{Inductive step from $m$ to $m+1$.} Let $(C,i)$ be live at $m+1$, and assume the bound at $m$.
The bound at $m$ gives $\mathsf{cl}_m(s)\subset\Omega_{m+1}^c$. Since $\Lambda_{m+1}\subset
\Omega_{m+1}$, this set lies in $\Lambda_{m+1}^c=Z_{m+1}$.

The set $\mathsf{cl}_m(s)$ is connected, because $s$ is a box, and it contains $C$. So it lies in the
component of $Z_{m+1}$ containing $C$. By liveness at $m+1$, this component is not of first class at
$\mathbf{R}_{m+1}$, so $\mathbf{R}_{m+1}$ does not remove it, and $\mathsf{cl}_m(s)\subset
Z_{m+1}^{+}$.

The partitions are nested, so $\mathsf{cl}_{m+1}(s)=\mathsf{cl}_{m+1}(\mathsf{cl}_m(s))$. This set
lies within one level-$(m+1)$ cube of $\mathsf{cl}_m(s)$, that is, within $M$ level-$(m+1)$ units.
Hence \eqref{9.1:eq-island-geometry-a} at level $m+1$ gives
\begin{equation*}
\operatorname{dist}\bigl(\Omega_{m+2},\mathsf{cl}_{m+1}(s)\bigr)\ \ge\ 4LM\check R_{m+2}-M
\ \ge\ 3LM\check R_{m+2},
\end{equation*}
because $L\check R_{m+2}\ge1$. This completes the induction. Dividing by the side $M$ of a level-$m$
cube, and using $L\check R_{m+1}\ge\check R_m$, gives the statement in level-$m$ cubes.

\emph{The island bound.} Every $I\in\mathbb{I}_m$ lies in $\mathbb{W}_m=[\mathsf{cl}_m
(\mathfrak{S}_m)]^{(4)}$, and $\mathfrak{S}_m$ is the union of the strips of births $(C,i)$ with
$i\le m$ that are live at $m$. To each such strip the bound just proved applies. Therefore
\begin{equation*}
\operatorname{dist}_m(I,\Omega_{m+1})\ \ge\ \min_s\operatorname{dist}_m\bigl(\mathsf{cl}_m(s),
\Omega_{m+1}\bigr)-4\ \ge\ 3\check R_m-4\ \ge\ 23,
\end{equation*}
where the last inequality is $\check R_m\ge9$ from \eqref{9.1:eq-island-geometry-b}.

\emph{The bound in level-$(m+1)$ cubes.} The distance $3LM\check R_{m+1}$ level-$m$ units equals
$3\check R_{m+1}$ level-$(m+1)$ cubes. From this we subtract one cube for the closure
$\mathsf{cl}_{m+1}$. We subtract one more cube for the four level-$m$ layers, whose rescaled width
$4/L$ is less than one level-$(m+1)$ cube. This gives $3\check R_{m+1}-2$.

\emph{The final sentence of (a).} It follows because $\Lambda_{m+1}\subset\Omega_{m+1}$. The region
$Y_0$ of the decoupling lemma lies in $\Lambda_{m+1}$, since at the place where the decoupling lemma
is used the expansion lemma is run on the integral over $\Lambda_{m+1}$. The fluctuation field is
localised in $\Lambda_{m+1}$, hence in $\Omega_{m+1}$ \cite{B-CMP119}.

\emph{(b)} Let $m\ge i+1$. Statement (a) at level $m-1$ gives
\begin{equation*}
\operatorname{dist}\bigl(\mathsf{cl}_{m-1}(s),\Omega_m\bigr)\ \ge\ 3LM\check R_m
\quad\text{(level-$(m-1)$ units)}.
\end{equation*}
This is $3\check R_m$ level-$m$ cubes. The set $\mathsf{cl}_m(s)^{(4)}$ lies within five level-$m$
cubes of $\mathsf{cl}_{m-1}(s)$: one cube for the closure and four for the layers. Hence it lies in
$\Omega_m^c$, with $3\check R_m-5\ge22$ cubes to spare, and it lies in $\Lambda_m^c=Z_m$. Since
$\Lambda_m\subset\Omega_m$, its distance from $\Lambda_m$ is at least $3\check R_m-5$ cubes. It is a
connected union of level-$m$ cubes containing $C$, so it lies in the component $\mathcal{C}_m(C)$ of
$Z_m$ containing $C$.

\emph{Consequence (b1).} Let all births of an island be at levels $\le m-1$. By the first part of (b),
the set $\mathsf{cl}_m(s')^{(4)}$ of each of their strips $s'$ lies inside one component of $Z_m$.
These components are
not of first class, by liveness. Distinct components of $Z_m$ are at least $\check R_m$ cubes apart,
by the separation of components at the large-field step. Two such closures can therefore touch only
if they lie in the same component. A touch-component is then contained in one component, at depth at
least $3\check R_m-5$.

\emph{Consequence (b2).} Consider a cluster containing new strips. An old closure in $\mathbb{W}_m$
belongs to a birth that is live at $m$. So its component is not of first class at $\mathbf{R}_m$:
either it is not in $Z(m)$, or it is a second-class component of $Z(m)$.
\begin{itemize}
\item Old closures inside other live components, those not in $Z(m)$, are at least
$3\check R_m-5$ deep in them. Distinct components are at least $\check R_m$ apart, by the separation
of components. A new strip's four layers reach at most $m_\sharp+4\le\frac12(\check R_m-1)+4$ cubes
from its component $C'$, and this is less than $\check R_m$ because $\check R_m\ge9$; for
$\check R_m=9$ one has $m_\sharp=4$ and
\begin{equation*}
m_\sharp+4=8<9 .
\end{equation*}
So such old closures cannot touch a new strip's four layers.
\item The remaining old islands lie inside second-class components $C'$ of $Z(m)$, and
$C'\subset s(C',m)$.
\end{itemize}
Hence a cluster containing new strips contains, besides them, only islands inside their own
components $C'\subset s(C',m)$. The strips of two second-class components can share a cluster only
if they are within $8$ cubes of each other; this is the case in which their four-layer fattenings
touch.

\emph{Consequence (b3).} In the situation just described, the two components lie in $Z_m^{+}$. Each
strip lies within $m_\sharp$ cubes of its component, so the two components are at most
$2m_\sharp+8\le\check R_m+7<8\check R_m$ cubes apart. By \eqref{9.1:eq-island-geometry-a} at level
$m$, the complement $\Omega_{m+1}^c$ contains both components fattened by
$4LM\check R_{m+1}\ge4M\check R_m$. These fattenings overlap, so their union is a connected subset of
$\Omega_{m+1}^c\subset Z_{m+1}$. It merges the two components into one component of $Z_{m+1}$.
If an island meets more than two components, we apply this to each pair of new strips whose
four-layer fattenings touch; since the island is a touch-component, these pairs link all its strips,
and all the components it meets lie in one component of $Z_{m+1}$. At each later level $m'\ge m+1$ at
which the births are live, this component is not healed by $\mathbf{R}_{m'}$, so it lies in
$Z_{m'}^{+}$; by \eqref{9.1:eq-island-geometry-a} at level $m'$ its $4LM\check R_{m'+1}$-fattening,
which is connected, lies in $\Omega_{m'+1}^c\subset Z_{m'+1}$, so the components remain in one
component of $Z_{m'+1}$.

\emph{(c1)} The island $I$ consists of closures of births at levels $\le m-1$ that are live at $m$,
and the components of these births lie in $C''$.
By (b), $I\subset C''$. By the separation of components, $\operatorname{dist}_m(C'',C')\ge\check R_m$,
hence $\operatorname{dist}_m(I,C')\ge\check R_m$, and this holds for every component $C'$ of $Z(m)$.
So every cube of $I$ has layer index at least $\check R_m+1$ with respect to $Z(m)$. This exceeds
$m_\sharp+1$, since $m_\sharp\le\frac12(\check R_m-1)<\check R_m$. The strip $C'^{\sharp}$ lies within
$m_\sharp$ cubes of $C'$, so
\begin{equation*}
\operatorname{dist}_m(I,C'^{\sharp})\ \ge\ \check R_m-m_\sharp\ \ge\ \check R_m-\tfrac12(\check R_m-1)
\ =\ \tfrac12(\check R_m+1)\ \ge\ 5,
\end{equation*}
where the last inequality uses $\check R_m\ge9$.

\emph{(c2)} \cite{B-CMP122b} states that each term determines the localisation domain
$X'_0=\bigcup_{Y\in D}Y\cup\bigcup_jX_j$, with $D$ and the sets $X_j$ as there. Every term of the
action whose localisation domain is
linked to the $\mathbf{T}'$-domain of a first-class component is a factor of some activity $F(X')$ of
\cite[(1.91)]{B-CMP122b}, and it is absorbed into the output of $\mathbf{R}_m$. A localisation domain
containing $I\subset C'$ meets $C'$, and $C'\subset C'^{\boxplus}$. Hence every stored term whose
localisation domain contains $I$ is such a term. After $\mathbf{R}_m$ the births inside $C'$ are not
live, by the definition of liveness, and no stored term is blind over them.

\emph{(c3)} By hypothesis $I\subset C'$. Since $C'\subset C'^{(m_\sharp(m))}=s(C',m)$, the island $I$
lies in $s(C',m)^{(4)}$. Since $(C',m)$ is a birth live at $m$, we have
$s(C',m)\subset\mathfrak{S}_m$, and $s(C',m)$ is a union of level-$m$ cubes, so
\begin{equation*}
s(C',m)^{(4)}=\bigl[\mathsf{cl}_m(s(C',m))\bigr]^{(4)}\subset\mathbb{W}_m .
\end{equation*}
This set is connected, so it lies in one island of $\mathbb{I}_m$, which therefore contains $I$, in
accordance with (b).

\emph{(d)} We treat the three places in turn.

\emph{The renormalisation transformation.} The decoupling cubes of the decoupling lemma are cubes
$\Delta\subset Y_0$, and at the place where that lemma is used the expansion lemma is run on the
integral over $\Lambda_{m+1}$; so these
cubes lie in $\Lambda_{m+1}\subset\Omega_{m+1}$. In \cite{B-CMP119} the fluctuation field is localised
in $\Lambda_{k+1}$, hence in $\Omega_{k+1}$, with $k=m$. The generalised random walk expansions start
at the cubes of the fluctuation field \cite{B-CMP119}. All three kinds of cells therefore lie in
$\Omega_{m+1}$, and the distance bound follows from the island bound in (a).

\emph{The resummation.} The anchoring and the value budget are those of the resummation proposition
of the large-field step and the relative resummation lemma, as stated in (d2). Every island in
$\mathbb{W}'_m$ lies, by (b), inside a component of $Z_m$. That component is not in $Z(m)$: an island
of $\mathbb{I}_m^{-}$ inside a component $C'$ of $Z(m)$ lies in $C'^{\sharp}$ or in $s(C',m)$, hence
in $Z^{\sharp}(m)$, by (c2) and (c3), and so does not belong to $\mathbb{W}'_m$. Hence, by (c1),
every island in $\mathbb{W}'_m$ is at least $\frac12(\check R_m+1)\ge5$ cubes from every strip of
$Z(m)$.

\emph{The preparatory integrations.} The integration variables and the shell are those of
\cite{B-CMP122a}, as stated in (d3). Let the shell lie in the component $C'$ of $Z(m)$. An island
lying in another component of $Z_m$ is off the shell. If that component is not in $Z(m)$, the island
is at least $\frac12(\check R_m+1)$ cubes away, by (c1). If that component is in $Z(m)$, the island
lies inside it by (b), and distinct components are at least $\check R_m$ cubes apart by the
separation of components. At a scale $j+1<m$, let an island of $\mathbb{I}_{j+1}$ lie in a component
of $Z_{j+1}$ other than the one containing the part of the shell under consideration. Its old closures
lie inside their component by (b), and the four layers of a strip born at $j+1$ reach at most
$m_\sharp(j+1)+4<\check R_{j+1}$ cubes from its component, as in the proof of (b2). Since distinct
components of $Z_{j+1}$ are at least $\check R_{j+1}$ cubes apart, the island keeps a positive cube
distance from the shell.

Now fix $j$ with $m-N_{\mathrm{pr}}\le j\le m-1$, and let an island lie inside a component $C'$ of
$Z(m)$, belonging to $\mathbb{I}_{j+1}$ if $j+1<m$ and to $\mathbb{I}_m^{-}$ if $j+1=m$. In either
case its births are at levels $\le m-1$. \cite[condition~(ii)]{B-CMP122a} requires
that in the preceding $N_{\mathrm{pr}}$ renormalisation steps no new large-field regions were created
inside this component. A second-class event at a level $i'$ creates such a region: by
\cite[p.~390]{B-CMP122b} the new large-field region is $Z_k\cup\bigcup_iY_i$ (compare the factor
$\chi^c_k(Y^{\sim-6})$ of \cite[(1.102)]{B-CMP122b}). Condition~(ii) therefore excludes births
inside $C'$ at the levels $m-N_{\mathrm{pr}},\dots,m-1$. Since the births of the island are at levels
$\le m-1$, every birth $(C_1,i')$ of the island satisfies $i'\le m-N_{\mathrm{pr}}-1$. For each such
birth let $s':=s(C_1,i')$.

Since $j\ge m-N_{\mathrm{pr}}$, we have $j+1\ge i'+1$. So (b) at level $j+1$ gives, for each of
these strips $s'$,
\begin{equation*}
\mathsf{cl}_{j+1}(s')^{(4)}\subset Z_{j+1},\qquad
\operatorname{dist}_{j+1}\bigl(\mathsf{cl}_{j+1}(s')^{(4)},\Lambda_{j+1}\bigr)\ \ge\
3\check R_{j+1}-5\ \ge\ 2 .
\end{equation*}
The shell $Z_{j+1}\setminus Z''_{j+1}$ is one layer of cubes at $\partial Z_{j+1}$. Each of these
closures therefore lies off the shell, hence so does the whole island, and no cell of the preparatory
integrations at scale $j+1$ lies in such an island.

\emph{(e)} Let a stored term have level $\ell$. Its blind set is a union of islands of
$\mathbb{I}_\ell$, so all births of each of these islands are live at $\ell$. Suppose the term exists
at a level $m\ge\ell$ and some birth of one of these islands is not live at $m$; then $m>\ell$. By
(b), all births of the island lie in one component from the level after the youngest birth on. By
the definition of liveness, that component was healed at some $\mathbf{R}_{m'}$ with
$\ell<m'\le m$. At the first such $\mathbf{R}_{m'}$, the localisation domain of the
stored term contains the island, so it meets the component, and the term is absorbed, by (c2). So
the term does not exist at level $m$. This proves the first sentence.

For the second sentence, let the stored term have level $\ell\le m$ and exist at $m$. If $m=\ell$,
its blind set lies in $[\mathsf{cl}_\ell(\mathfrak{S}_\ell)]^{(4)}=\mathbb{W}_\ell$ by definition. Let
$m>\ell$. The blind set $\mathcal{B}$ is a union of islands of $\mathbb{I}_\ell$. Each of them is a
touch-component of $\mathbb{W}_\ell$, hence a union of sets $[\mathsf{cl}_\ell(s'')]^{(4)}$, $s''$
running over its own strips. By the first sentence, the births of these strips are live at $m$, so
the strips lie in $\mathfrak{S}'_\ell$, and $\mathfrak{S}'_\ell\subset\mathfrak{S}_m$. This gives the
first inclusion of the display. The partitions are nested
(Notation~\ref{9.1:not-cubes-strips}), so $\mathsf{cl}_m\circ\mathsf{cl}_\ell=\mathsf{cl}_m$. Four
layers of level-$\ell$ cubes have total width $4L^{-(m-\ell)}$ in units of level-$m$ cubes, and this
is less than one because $L\ge5$ and $m>\ell$. Hence the closure in level-$m$ cubes of
$[\mathsf{cl}_\ell(\mathfrak{S}'_\ell)]^{(4)}$ lies within one level-$m$ layer of
$\mathsf{cl}_m(\mathfrak{S}'_\ell)\subset\mathsf{cl}_m(\mathfrak{S}_m)$. This gives the second
inclusion, and $[\mathsf{cl}_m(\mathfrak{S}_m)]^{(1)}\subset\mathbb{W}_m$ by definition.

\emph{(f)} The norm of the resummation proposition of the large-field step is
\begin{equation*}
\|\mathfrak{v}\|_a:=\sup_{Y_4\setminus Z\neq\emptyset}e^{a\,d_{k,Z}(Y_4)}|\mathfrak{v}(Y_4)| .
\end{equation*}
Its input length is therefore the tree length $d_{k,Z}$ relative to $Z=Z(m')$, with $k=m'$, and
$Z(m')\supseteq\mathcal{C}\supseteq I$. The domination of this length by the term's own relative
decay is the comparison of relative tree lengths supplied by the strip geometry of the large-field
step:
\begin{equation*}
d_{m',Z(m')\cup\mathbb{W}'_{m'}}(Y_1)\ \le\ d^{\circ}_\ell(K_v),
\end{equation*}
with $Y_1$, $K_v$ and $\ell$ as in (f).

The crossing of a healed box $X^{\sharp}_j\supseteq\mathcal{C}^{\sharp}\supseteq I$ by a Hamiltonian
path is paid for by \cite[(1.89)]{B-CMP122b}, through the Hamiltonian-path bounds of the large-field
step. These set the factors
$\exp\bigl(2\kappa(100R_k)^d\bigr)$ of the product taken in \cite{B-CMP122b} against the inequality
\begin{equation*}
-2(1+\beta_0)^{-1}p_0(g_k)+1+2\kappa(100R_k)^d\ \le\ -\tfrac32\,p_0(g_k)
\end{equation*}
stated in \cite[\S1]{B-CMP122b}, with $\beta_0$ Ba{\l}aban's constant of that place (not a coupling
increment), $R_k=\check R_k$ and
$p_0(\cdot)$ the degree function. Hence no
blind stored term survives a healing. At a re-birth the old island lies inside the new strip, by
(c3). These steps are carried out in detail in the relative resummation lemma.
\end{proof}

We note three features of Lemma~\ref{9.1:lem-island-geometry}.

None of the statements (a)--(f) depends on the number of strips in an island, on the size of an
island, on $N_{\mathrm{pr}}$, on an age, or on the number of steps. The scaffold radius enters only
through distances of the form $3\check R-O(1)$, which are used as ``$\ge1$'' or ``$\ge23$'', and
through \eqref{9.1:eq-island-geometry-b}.

The four layers of Ba{\l}aban's construction are what (a) uses. A live strip is separated from every
later fluctuation region by a moat several strip-widths wide, since
$m_\sharp\le\frac12\check R_m<3L\check R_{m+1}$. In \cite{B-CMP122b} the same property is expressed
by the statement that the variables inside a second-class region $Y_i$ do not participate in the
operations connected with the next renormalisation transformations.

The rule by which the inner zones of second-class components, and of components arrested in the
sense of the arrest theorem, leave every determining set of Ba{\l}aban's construction concerns the
class to which a stored term is assigned, not the length in which it is measured. It is
compatible with the statements above and does not affect them.

We fix a level $\ell$. We use the level-$m$ cubes and the following objects of Notation~\ref{9.1:not-cubes-strips}:
\begin{itemize}
\item the distances $\operatorname{dist}_m$ counted in cubes;
\item the cube counts $N_m(\cdot)$;
\item the closures $\mathsf{cl}_m$;
\item the tree lengths $d_m(\cdot)$;
\item the relative tree lengths $d_{m,W}(\cdot)$;
\item the scaffold radii $\check{R}_m$;
\item the crossing lengths $\mathsf{d}_m$.
\end{itemize}

We also use the union $\mathbb{W}_\ell$ of the live strips at level $\ell$ and the set $\mathbb{I}_\ell$ of its islands, as in Lemma~\ref{9.1:lem-island-geometry}, and the class $\mathbf{D}_\ell$ of localisation domains of level $\ell$.

\emph{The island-blind format.} Put
\begin{equation*}
\mathbf{D}_\ell^{\mathbb{I}}:=\{Y\in\mathbf{D}_\ell:\ Y\cap\mathbb{W}_\ell
\text{ is a union of islands of }\mathbb{I}_\ell\},
\end{equation*}
where the union may be empty.

A \emph{store of level $\ell$ in format} is a family $F=\{F(Y)\}_{Y\in\mathbf{D}_\ell^{\mathbb{I}}}$. It has one function for each domain and for each kind of stored term, in the same way as the stored boundary terms $\mathbf{B}^{(\ell)}$, $\mathbf{B}'^{(\ell)}$ of the $\mathbf{T}$- and $\mathbf{R}$-steps. Its support satisfies $\operatorname{supp}F(Y)\subset Y$, and inside the islands it may be smaller than $Y$. We speak of the store $Y$ when we mean the term stored on the domain $Y$.

A domain $Y$, and the store on it, is \emph{blind} if $Y\cap\mathbb{W}_\ell\neq\emptyset$, and \emph{plain} otherwise. The \emph{excused set} of $Y$ is $Y\cap\mathbb{W}_\ell$. The \emph{relative length} of $Y$ is
\begin{equation*}
d^{\circ}_\ell(Y):=d_{\ell,\mathbb{W}_\ell}(Y);
\end{equation*}
for plain $Y$ it equals the tree length $d_\ell(Y)$.

\emph{Filing.} For a connected union $X$ of level-$\ell$ cubes put
\begin{equation*}
\operatorname{fil}_\ell(X):=X\cup\bigcup\{I\in\mathbb{I}_\ell:\ I\cap X\neq\emptyset\}
\in\mathbf{D}_\ell^{\mathbb{I}} .
\end{equation*}
Then $\operatorname{fil}_\ell(X)\setminus\mathbb{W}_\ell=X\setminus\mathbb{W}_\ell$. Filing is an indexing of the stored function: the function is indexed by the domain $\operatorname{fil}_\ell(X)$, and its support stays what it was, exactly as for the filing map $\operatorname{fil}_{\sharp}$ of the strip format.

A store $Y$ of level $\ell$ is read at a level $m\ge\ell$ through
\begin{equation*}
\bar Y^{(m)}:=\mathsf{cl}_m(Y),\qquad
d^{\circ}_m(Y):=d_{m,\mathbb{W}_m}\bigl(\operatorname{fil}_m(\bar Y^{(m)})\bigr).
\end{equation*}

\emph{Norms.} The one-machine size of a store $F$ in format is
\begin{equation*}
\mathsf{N}^{\circ}_\ell[F]:=\sup_{Y}e^{r_F d^{\circ}_\ell(Y)}\sup|F(Y)| .
\end{equation*}
Here $r_F$ is the decay rate assigned to the kind of term stored:
\begin{itemize}
\item $r_F=\tfrac98\kappa$ for the $\mathbf{R}$-born terms of the first group, by the relative bound for $\mathbf{R}$-born terms;
\item $r_F=\kappa_{out}$ at the site of the $\mathbf{T}$-step;
\item $r_F=a_k$ or $a_m$ at the collar site, where $a_k$ and $a_m$ are the decay rates, indexed by the levels $k$ and $m$, of the bound on the terms stored at the collar site.
\end{itemize}

For the two-machine members we put
\begin{equation*}
\mathsf{S}^{\circ,(n)}_\ell[F]:=\sup_{Y}e^{\kappa_{\mathbf{B}}d^{\circ}_\ell(Y)}
\sup_{\mathsf{v},\,\mathrm{label}}\bigl|\delta F^{(\ell),(n)}(Y;\mathsf{v})\bigr|,
\qquad \kappa_{\mathbf{B}}=\kappa .
\end{equation*}
Here $\delta F^{(\ell),(n)}(Y;\mathsf{v})$ is the two-machine difference of the term stored on $Y$, its member of index $n$, at data $\mathsf{v}$, and the supremum runs over the data and the labels. For the sub-family $\mathbf{B}'^{own}$ of the stored boundary terms $\mathbf{B}'^{(\ell)}$ of the $\mathbf{R}$-step, the seminorm $\mathsf{S}^{\circ}_\ell[\mathbf{B}'^{own}]$ is defined by the same formula.

For plain $Y$ these are the norms already used for plain domains. The stores and statements for plain domains are unchanged, and $\mathsf{S}$ below denotes the supremum over plain and blind domains together.

\begin{lemma}[Domain counts in the island-blind length]\label{9.1:lem-blind-counts}
Let $\ell$ be a level. The following statements hold.
\begin{enumerate}
\item[(i)] \textup{(Count at level $\ell$.)}
Let $\Delta_\ell$ be a level-$\ell$ cube with $\Delta_\ell\not\subset\mathbb{W}_\ell$, and let $\lambda_0\ge\mathsf{k}_d+1$. Then
\begin{equation}\label{9.1:eq-blind-counts-a}
\sum_{Y\in\mathbf{D}_\ell^{\mathbb{I}},\ Y\supset\Delta_\ell}
e^{-\lambda_0 d^{\circ}_\ell(Y)}\le C_{fib}.
\end{equation}
\item[(ii)] \textup{(Site cells.)}
Let $m\ge\ell$, and let $\Delta$ be a level-$m$ cube with $\operatorname{dist}_m(\Delta,\mathbb{W}_m)\ge1$. By Lemma~\ref{9.1:lem-island-geometry}\textup{(d)} this hypothesis is met by every cell $\Delta$ of the site of the $\mathbf{T}$-step, which has $\operatorname{dist}_m(\Delta,\mathbb{W}_m)\ge 3\check{R}_m-4$, and by every exterior anchor cell and every collar cell $\Delta$, which have $\operatorname{dist}_m(\Delta,\mathbb{W}_m)\ge1$. Let $\lambda\ge\mathsf{k}_d+2$. Summing over the blind stores $Y\in\mathbf{D}_\ell^{\mathbb{I}}$ of level $\ell$ still present at $m$,
\begin{equation*}
\sum_{\substack{Y\in\mathbf{D}_\ell^{\mathbb{I}}\ \mathrm{blind},\\ \bar Y^{(m)}\supset\Delta}}
e^{-\lambda d^{\circ}_\ell(Y)}\le K^{\circ}(d)
:=C_{fib}\sup_{y\ge1}y^d e^{-(y-3)} .
\end{equation*}
\item[(iii)] \textup{(An island's own cells.)}
Let $I\in\mathbb{I}_\ell$ and $\lambda_0\ge\mathsf{k}_d+1$. Then
\begin{equation*}
\sum_{Y\in\mathbf{D}_\ell^{\mathbb{I}},\ Y\supseteq I}e^{-\lambda_0 d^{\circ}_\ell(Y)}
\le 1+3^d C_{fib}N_\ell(I),
\qquad
N_\ell(I)\le\sum_j\bigl|\Gamma^0_j\cap I\bigr| .
\end{equation*}
Here $\Gamma^0_j$ is the set of cells of zone $j$ of the current determining set, and $|\Gamma^0_j\cap I|$ is the number of those cells lying in $I$.
\item[(iv)] \textup{(Crossing.)}
Let $Y$ be a blind store of level $\ell$ born at the $\mathbf{T}$-step, that is, with $Y\cap\Omega_\ell\neq\emptyset$ (see \cite[(2.41)]{B-CMP119}). Then
\begin{equation*}
d^{\circ}_\ell(Y)\ge 3\check{R}_\ell-8\ge\mathsf{d}_{\ell-1}-8,
\end{equation*}
where $\mathsf{d}_{\ell-1}=2\check{R}_{\ell-1}/L-c_\times$ is the crossing length of Notation~\ref{9.1:not-cubes-strips}.
\item[(v)] \textup{(Levels.)}
Let $Y$ be a store of level $\ell$ present at a level $m\ge\ell$. Then $\operatorname{fil}_m(\bar Y^{(m)})\in\mathbf{D}_m^{\mathbb{I}}$ and
\begin{equation*}
d^{\circ}_m(Y)\le L^{-(m-\ell)}d^{\circ}_\ell(Y).
\end{equation*}
More generally, let $W'\supset Y\cap\mathbb{W}_\ell$ be a set, treated as excused, with $\mathsf{cl}_m(W')\subset\mathbb{W}_m$, and let $T$ be any tree in $Y$ meeting every level-$\ell$ cube of $Y\setminus W'$. Then $d^{\circ}_m(Y)\le L^{-(m-\ell)}|T|$, where $|T|$ is the length of $T$ in level-$\ell$ units.
\end{enumerate}
\end{lemma}

\begin{proof}
\emph{(i)} This is the relative domain count of the strip format, in its parts (b) and (c), applied with excused set $W:=\mathbb{W}_\ell$ and islands $\mathbb{I}_\ell$. The cube $\Delta_*$ of that count is taken to be $\Delta_\ell$; it lies in $Y\setminus\mathbb{W}_\ell$ for every $Y$ in the sum.

The count applies in this setting because its proof uses only one property of the islands: they are pairwise non-touching unions of cubes. Precisely, the fewer than $3^d$ neighbours of a cube lie in fewer than $3^d$ islands. No bound on the size or on the make-up of an island enters that proof. The islands of $\mathbb{I}_\ell$ are pairwise non-touching unions of cubes, so the count gives \eqref{9.1:eq-blind-counts-a}.

\emph{(ii)} Let $Y$ be a blind store of level $\ell$ present at $m$. By Lemma~\ref{9.1:lem-island-geometry}\textup{(e)}, its excused set satisfies $\mathsf{cl}_m(Y\cap\mathbb{W}_\ell)\subset\mathbb{W}_m$.

Since $\operatorname{dist}_m(\Delta,\mathbb{W}_m)\ge1$, the cube $\Delta$ is not contained in $\mathbb{W}_m$. Hence $\Delta$ is not one of the cubes of $\mathsf{cl}_m(Y\cap\mathbb{W}_\ell)$, and so $\Delta$ does not meet $Y\cap\mathbb{W}_\ell$. Consequently none of the $L^{d(m-\ell)}$ level-$\ell$ sub-cubes $\Delta_\ell\subset\Delta$ lies in $Y\cap\mathbb{W}_\ell$.

Such a $\Delta_\ell$ also lies in none of the other islands of $\mathbb{W}_\ell$. A sub-cube inside another island $I'\in\mathbb{I}_\ell$ that is live at $m$ would make $\Delta$ meet $I'$; this would put $\Delta$ inside $\mathsf{cl}_m(I')\subset\mathbb{W}_m$, contrary to $\operatorname{dist}_m(\Delta,\mathbb{W}_m)\ge1$. Islands whose birth is no longer live carry no stores at $m$, and they may be deleted from $\mathbb{W}_\ell$ for the purpose of this count.

Thus every $\Delta_\ell\subset\Delta$ satisfies $\Delta_\ell\not\subset\mathbb{W}_\ell$. Moreover, $\bar Y^{(m)}\supset\Delta$ if and only if $Y$ contains some level-$\ell$ sub-cube $\Delta_\ell\subset\Delta$.

Fix such a $Y$ and such a $\Delta_\ell\subset Y$. Since $Y$ is blind, $Y\cap\mathbb{W}_\ell$ is a non-empty union of islands, so $Y\supseteq I$ for some $I\in\mathbb{I}_\ell$. This island satisfies
\begin{equation*}
\mathsf{cl}_m(I)\subset\mathsf{cl}_m(Y\cap\mathbb{W}_\ell)\subset\mathbb{W}_m .
\end{equation*}
Therefore
\begin{equation*}
\operatorname{dist}_\ell(\Delta_\ell,I)\ge L^{m-\ell}\operatorname{dist}_m(\Delta,\mathbb{W}_m)
\ge L^{m-\ell}=:y .
\end{equation*}

The relative tree of $Y$ meets every level-$\ell$ cube of $Y\setminus\mathbb{W}_\ell$. So it contains a point of $\Delta_\ell$, and also a point of a cube of $Y\setminus\mathbb{W}_\ell$ that touches $I$. Such a cube exists for the following reason. First, $Y\supsetneq I$, because $\Delta_\ell\subset Y$ and $\Delta_\ell\not\subset\mathbb{W}_\ell$. Second, $Y$ is connected, so there is a chain of cubes inside $Y$ from $I$ to $\Delta_\ell$. Third, islands do not touch, so the first cube on such a chain that is not in $I$ is off $\mathbb{W}_\ell$ and touches $I$.

Since that cube touches $I$ and $\operatorname{dist}_\ell(\Delta_\ell,I)\ge y$, the length of the relative tree is at least $y-3$. That is,
\begin{equation*}
d^{\circ}_\ell(Y)\ge y-3 .
\end{equation*}

Put $\lambda_0:=\lambda-1$; then $\lambda_0\ge\mathsf{k}_d+1$. For each $Y$ in the sum, write $e^{-\lambda d^{\circ}_\ell(Y)}=e^{-(\lambda-\lambda_0)d^{\circ}_\ell(Y)}e^{-\lambda_0 d^{\circ}_\ell(Y)}$ and use $d^{\circ}_\ell(Y)\ge y-3$. Count each $Y$ once for each sub-cube $\Delta_\ell\subset\Delta$ that it contains; this includes every $Y$ of the sum at least once. For each $\Delta_\ell$, apply \eqref{9.1:eq-blind-counts-a} of part (i), which is legitimate since $\Delta_\ell\not\subset\mathbb{W}_\ell$. This gives
\begin{align*}
\sum_{\substack{Y\ \mathrm{blind},\\ \bar Y^{(m)}\supset\Delta}}e^{-\lambda d^{\circ}_\ell(Y)}
&\le\sum_{\Delta_\ell\subset\Delta}\ \sum_{Y\in\mathbf{D}_\ell^{\mathbb{I}},\ Y\supset\Delta_\ell}
e^{-(\lambda-\lambda_0)(y-3)}e^{-\lambda_0 d^{\circ}_\ell(Y)}\\
&\le L^{d(m-\ell)}e^{-(\lambda-\lambda_0)(y-3)}C_{fib}
= y^d e^{-(y-3)}C_{fib}\le K^{\circ}(d).
\end{align*}
Here we used $\lambda-\lambda_0=1$ and $L^{d(m-\ell)}=y^d$, and the last inequality holds because $y\ge1$. At $m=\ell$ we have $y=1$, and the bound reads $e^2C_{fib}\le K^{\circ}(d)$.

\emph{(iii)} The domain $Y=I$ contributes a single term. Its excused set is all of $I$, so $d^{\circ}_\ell(I)=0$ and the term equals $1$.

Now let $Y\supsetneq I$. As in part (ii), $Y$ is connected and islands do not touch. Hence $Y$ contains a cube $\Delta_*$ with $\Delta_*\not\subset\mathbb{W}_\ell$ that touches $I$. Each such cube touches one of the $N_\ell(I)$ cubes of $I$, and a cube touches fewer than $3^d$ other cubes. So there are at most $3^dN_\ell(I)$ such cubes. Assign each $Y\supsetneq I$ to one such cube $\Delta_*\subset Y$ and apply part (i) at each $\Delta_*$. This gives
\begin{equation*}
\sum_{Y\in\mathbf{D}_\ell^{\mathbb{I}},\ Y\supseteq I}e^{-\lambda_0 d^{\circ}_\ell(Y)}
\le 1+\sum_{\Delta_*}\ \sum_{Y\in\mathbf{D}_\ell^{\mathbb{I}},\ Y\supset\Delta_*}
e^{-\lambda_0 d^{\circ}_\ell(Y)}
\le 1+3^dC_{fib}N_\ell(I).
\end{equation*}

For the cells, take a level-$\ell$ cube of $I$. It lies in some zone $j\le\ell$ or in $\Omega_\ell$. In either case the cells of the current determining set there have side at most $1$ in level-$\ell$ lattice units, the side in zone $j$ being
\begin{equation*}
L^{-(\ell-j)}\le 1\le M,
\end{equation*}
by the scales of the determining sets in \cite{B-CMP119}; and a level-$\ell$ cube has side $M\ge1$ in these units. So the cube contains at least one cell. The cubes of $I$ are distinct, so distinct cubes of $I$ contain distinct cells, and therefore $N_\ell(I)\le\sum_j|\Gamma^0_j\cap I|$.

\emph{(iv)} The store $Y$ contains an island $I\in\mathbb{I}_\ell^-$, one of the old live islands, and $Y$ meets $\Omega_\ell$. Let $I_0$ be the un-layered closure of $I$. Lemma~\ref{9.1:lem-island-geometry}\textup{(a)} at $m=\ell-1$ gives, in level-$\ell$ cubes,
\begin{equation*}
\operatorname{dist}_\ell\bigl(\mathsf{cl}_\ell(I_0),\Omega_\ell\bigr)\ge 3\check{R}_\ell-2,
\qquad\text{and hence}\qquad
\operatorname{dist}_\ell(I,\Omega_\ell)\ge 3\check{R}_\ell-6 .
\end{equation*}

The relative tree of $Y$ joins a cube of $Y\setminus\mathbb{W}_\ell$ touching $I$, found as in part (ii), to a cube of $Y$ meeting $\Omega_\ell$. Its length is therefore at least $3\check{R}_\ell-8$, which is the first inequality.

For the second inequality, the step relation of the scaffold radii gives
\begin{equation*}
3\check{R}_\ell\ge 3\check{R}_{\ell-1}/L\ge\mathsf{d}_{\ell-1},
\end{equation*}
where $\mathsf{d}_{\ell-1}=2\check{R}_{\ell-1}/L-c_\times$. Subtracting $8$ gives $3\check{R}_\ell-8\ge\mathsf{d}_{\ell-1}-8$.

\emph{(v)} Write $\bar Y=\bar Y^{(m)}$. By Lemma~\ref{9.1:lem-island-geometry}\textup{(e)}, $\mathsf{cl}_m(Y\cap\mathbb{W}_\ell)\subset\mathbb{W}_m$.

By the definition of filing, $\operatorname{fil}_m(\bar Y)\setminus\mathbb{W}_m=\bar Y\setminus\mathbb{W}_m$. Also, by construction, $\operatorname{fil}_m(\bar Y)\cap\mathbb{W}_m$ is a union of islands of $\mathbb{I}_m$, so $\operatorname{fil}_m(\bar Y)\in\mathbf{D}_m^{\mathbb{I}}$.

Let $Q$ be a level-$m$ cube of $\operatorname{fil}_m(\bar Y)\setminus\mathbb{W}_m=\bar Y\setminus\mathbb{W}_m$. Then $Q$ meets $Y$ in at least one level-$\ell$ cube not contained in $Y\cap\mathbb{W}_\ell$. Otherwise $Q\subset\mathsf{cl}_m(Y\cap\mathbb{W}_\ell)\subset\mathbb{W}_m$, contradicting $Q\subset\bar Y\setminus\mathbb{W}_m$.

Let $T$ be an optimal relative tree of $Y$: a tree in $Y$ meeting all level-$\ell$ cubes of $Y\setminus\mathbb{W}_\ell$, whose length in level-$\ell$ units is $d^{\circ}_\ell(Y)$. Then $T$ meets the level-$\ell$ cube just found, and hence meets $Q$. Thus $T$ meets every level-$m$ cube of $\operatorname{fil}_m(\bar Y)\setminus\mathbb{W}_m$, and $T\subset Y\subset\operatorname{fil}_m(\bar Y)$. So $T$ bounds $d_{m,\mathbb{W}_m}(\operatorname{fil}_m(\bar Y))=d^{\circ}_m(Y)$ by its length in level-$m$ units, which is $L^{-(m-\ell)}$ times its length in level-$\ell$ units. Hence $d^{\circ}_m(Y)\le L^{-(m-\ell)}d^{\circ}_\ell(Y)$.

The general form is proved by the same argument with $W'$ in place of $Y\cap\mathbb{W}_\ell$ and the given tree $T$ in place of the optimal one. A level-$m$ cube $Q$ of $\bar Y\setminus\mathbb{W}_m$ meets $Y$ in a level-$\ell$ cube not contained in $W'$, since otherwise $Q\subset\mathsf{cl}_m(W')\subset\mathbb{W}_m$. The tree $T$ meets that cube, hence $Q$, and $T\subset Y\subset\operatorname{fil}_m(\bar Y)$. Therefore $d^{\circ}_m(Y)\le L^{-(m-\ell)}|T|$.
\end{proof}

\begin{lemma}[The large-field resummation with island-blind units]\label{9.1:lem-blind-resummation}
Let $m$ be a step (in this lemma, as in Lemma~\ref{9.1:lem-island-geometry}, the letter $m$ indexes
a level of the renormalisation group, not the torus exponent), and consider the large-field operation
$\mathbf{R}_m$ of that step with its large-field region $Z_m$, its healable class $Z(m)$ with
first-class components $X_j$ and second-class components $Y_\iota$, the strips $X_j^{\sharp}$,
$Y_\iota^{\sharp}$ and the set $Z^{\sharp}(m)$ of the strip lemma listed in \textup{(i)} below, and the
island unions $\mathbb{W}_m^{-}\subset\mathbb{W}_m$ and island families $\mathbb{I}_m^{-}$,
$\mathbb{I}_m$ of Lemma~\ref{9.1:lem-island-geometry}. Put
\begin{equation*}
\mathbb{W}' := \bigcup\bigl\{ I\in\mathbb{I}_m^{-} : I\not\subset Z^{\sharp}(m)\bigr\}.
\end{equation*}
By Lemma~\ref{9.1:lem-island-geometry}(b),(c) these are exactly the islands lying inside live
components that do not belong to $Z(m)$; every other island of $\mathbb{I}_m^{-}$ lies inside a
strip of $Z^{\sharp}(m)$, namely inside a healed box or inside a new second-class strip. For a set
$W$ of level-$k$ cubes, $d_W$ denotes the level-$k$ tree length relative to $W$, in which a tree has
to meet only the cubes of a domain that do not lie in $W$; $d_{k,Z}$ is this length for $W=Z$, and
$\mathbf{D}_k$ is the class of localisation domains at level $k$. The statements listed in
\textup{(i)} below are written at a generic level $k$ with large-field region $Z$; they are applied
here at $k=m$. Assume the following.
\begin{itemize}
\item[(i)] Clauses (c)--(e) of the resummation proposition for the large-field operation, in the
form of its constituent statements: the strip lemma with its clauses (S1)--(S5), the
strip-resummation theorem with its clauses (a)--(c), its corollary on value units, the filing lemma
for value units, the dictionary $\mathfrak{D}^{\sharp}$ to the strip forms, the combination lemma with
its clauses (u1)--(u7), and the corollary on the modified bound.
\item[(ii)] Lemma~\ref{9.1:lem-island-geometry} and Lemma~\ref{9.1:lem-blind-counts}.
\item[(iii)] The island-blind input hypothesis, consisting of the following parts. Its volume part is
the volume part of the input hypothesis of the resummation proposition, unchanged. Its bracket part
is that the relative norm of the brackets is at most $\varepsilon_{br}$; this is the output of the
site (L) of Proposition~\ref{9.1:prop-blind-site-expansions}, whose proof is outstanding. Its value
inputs are the following two. First, the chain units $\mathbb{C}^{(i)}$ of the step obey the bound of
the chain-unit corollary; those that read an island-blind store are the units produced at the site
(P) of the same Proposition~\ref{9.1:prop-blind-site-expansions}, whose proof is outstanding. Here an
\emph{island-blind store} is a stored value of $\mathbf{B}$ or $\mathbf{B}'$ kept in the island-blind
format: a filed domain $Y\in\mathbf{D}_\ell^{\mathbb{I}}$, an excused set, and the relative length
$d^{\circ}_\ell$. Second, each stored value of $\mathbf{B}$ or $\mathbf{B}'$ entering the filing lemma
is either plain, bounded as in the filing lemma with the constant $B_0$ of that lemma, or an
island-blind store; an island-blind $\mathbf{B}'$ of level $j<k$ carries the rate $\tfrac98\kappa$ in
$d^{\circ}_j$, and an island-blind store born at a $\mathbf{T}$-step carries the rate $\kappa_{out}$,
with $\kappa_{out}\ge2(\mathfrak{c}+1)\kappa$.
\item[(iv)] For the dialled form only: clause (D)(e) of the dial lemma on the dial domain
$\mathbb{D}^{\star}$, and the linearity corollary for the filing maps.
\end{itemize}
The value part of the island-blind input hypothesis, the bound of the value family in the norm
$\|\cdot\|_a^{\circ}$ defined below, is not assumed: it is derived in the proof from the value inputs
of \textup{(iii)}.
In the statements listed in \textup{(i)}, read through the dictionary $\mathfrak{D}^{\sharp}$ to the
strips, make the substitutions
\begin{equation}\label{9.1:eq-blind-resummation-1}
\begin{aligned}
&\mathbf{D}_k^{Z}\mapsto\mathbf{D}_k^{Z,\circ},\qquad
\mathbf{D}_k^{\sharp}\mapsto\mathbf{D}_k^{\sharp,\circ};\\
&d_{k,Z}\mapsto d_{Z\cup\mathbb{W}'},\qquad d_\sharp\mapsto d_{Z^{\sharp}\cup\mathbb{W}'},\qquad
d_{out}\mapsto d_{out}^{\circ};\\
&\|\cdot\|_a\mapsto\|\cdot\|_a^{\circ},\qquad
\mathsf{K}_{val}(k)\mapsto\mathsf{K}_{val}^{\circ}(k);\qquad
\textstyle\bigcup_\iota Y_\iota\mapsto\bigcup_\iota Y_\iota^{\sharp}\cup\mathbb{W}',
\end{aligned}
\end{equation}
which are read as follows.
\begin{itemize}
\item[] \emph{Domains.} With $Y_4$ denoting a localisation domain, as in the resummation proposition,
put
\begin{equation*}
\begin{aligned}
\mathbf{D}_k^{Z,\circ} := \bigl\{ Y_4\in\mathbf{D}_k :\ & Y_4\cap Z \text{ is a non-empty union
of components of } Z,\\
& Y_4\cap\mathbb{W}' \text{ is a union of islands}\bigr\}.
\end{aligned}
\end{equation*}
The pre-filing of the input hypothesis, which replaces $Y_4$ by $Y_4$ together with the components
of $Z$ meeting $Y_4$, is replaced by the same operation followed by the adjunction of the islands of
$\mathbb{W}'$ meeting $Y_4$; this adds nothing outside $Z\cup\mathbb{W}'$ and does not raise
$d_{Z\cup\mathbb{W}'}$, for the same reason as in the original pre-filing. The class
$\mathbf{D}_k^{\sharp}$ is replaced by $\mathbf{D}_k^{\sharp,\circ}$ in the same way. The filing
map $\operatorname{fil}_\sharp$ is unchanged (it adds strips). The islands in (S5) are the strips of
$Z^{\sharp}$ together with the islands of $\mathbb{W}'$; they are pairwise non-touching by (S1) and
Lemma~\ref{9.1:lem-island-geometry}(c1).
\item[] \emph{Lengths and norms.} The lengths are replaced as in \eqref{9.1:eq-blind-resummation-1},
with
\begin{equation*}
d_{out}^{\circ} := d_{\bigcup_\iota Y_\iota^{\sharp}\cup\mathbb{W}'},\qquad
\|\mathfrak{v}\|_a^{\circ} := \sup_{Y_4\not\subset Z^{\sharp}}
e^{a\, d_{Z\cup\mathbb{W}'}(Y_4)}\,|\mathfrak{v}(Y_4)|.
\end{equation*}
A domain $Y_4$ is \emph{exterior} if $Y_4\not\subset Z^{\sharp}$ (unchanged). Every anchor
condition $\Delta\not\subset Z$ becomes $\Delta\not\subset Z\cup\mathbb{W}'$, and every anchor
``the first cube of $Y\setminus Z^{\sharp}$'', ``a cube in $Y\setminus Z^{\sharp}$'' becomes the
first cube, respectively a cube, of $Y\setminus(Z^{\sharp}\cup\mathbb{W}')$. Finally
\begin{equation*}
\mathsf{K}_{val}^{\circ}(k) := \sup_{\Delta\not\subset Z\cup\mathbb{W}'}
\sum_{v\in\mathcal{V}^{ext},\,K_v\ni\Delta} e^{a\, d^{\circ}_k(K_v)}\,\|v\|,
\end{equation*}
where $d^{\circ}_k$ is the relative length of the input store (Lemma~\ref{9.1:lem-blind-counts}(v))
for island-blind value units, and $d_k$ for plain ones.
\item[] \emph{Connectivity.} In \cite{B-CMP122b} the connectivity of the resummation reads:
\begin{quote}
At first we take components of $X'_0\setminus\bigcup Y_i$ and we consider two such components as
connected, if they are contained in one localization domain $Y$ of a term in the product. This
means that each component has a correct tree graph decay factor\,[\dots].
\end{quote}
Here $\bigcup_\iota Y_\iota$ is replaced by $\bigcup_\iota Y_\iota^{\sharp}\cup\mathbb{W}'$. That is:
in (u2), two polymers are compatible if and only if their parts outside
$\bigcup_\iota Y_\iota^{\sharp}\cup\mathbb{W}'$ neither meet nor touch; in (u3), the links ``cubes of
$Y\setminus Z^{\sharp}$ and $Y'\setminus Z^{\sharp}$ coincide or touch'' and ``a cube of
$Y\setminus Z^{\sharp}$ touches $X^{\sharp}_{j_h}$'' are read with $Z^{\sharp}\cup\mathbb{W}'$ in
place of $Z^{\sharp}$.
\end{itemize}
Then:
\begin{itemize}
\item[(a)] The strip lemma (S1)--(S5), the strip-resummation theorem (a)--(c), its corollary on
value units, the filing lemma for value units, the combination lemma (u1)--(u7) and the corollary on
the modified bound hold verbatim after the substitutions \eqref{9.1:eq-blind-resummation-1}, with the
same constants: $C_{fib}$, $\mathsf{k}_d$, $\alpha^{\sharp}$, $\alpha_*$, $m_{lo}$, the smallness
condition on the constants of the resummation proposition, $\beta_v$, the comparison of
$2\kappa(100\check{R}_m)^d$ against $-\tfrac32 p_0(g_m)$, and $c_{KP}$.
\item[(b)] For every $X$, the localised term $\mathbf{R}'^{(m)}(X,(\mathbf{U},\mathbf{J}))$ produced
by $\mathbf{R}_m$, in the notation of \cite[(1.99)]{B-CMP122b}, satisfies
\begin{equation}\label{9.1:eq-blind-resummation-2}
\bigl|\mathbf{R}'^{(m)}(X,(\mathbf{U},\mathbf{J}))\bigr|
\le O(1)\, c_1\, e^{-(1+\frac12\beta_s)\kappa\, d_{out}^{\circ}(X)},
\end{equation}
with $c_1=e^{-p_0(g_m)}$ if $X\cap\bigcup_\iota Y_\iota^{\sharp}=\emptyset$ (whether or not $X$
contains islands of $\mathbb{W}'$), and $c_1=(\alpha^{\sharp}_m)^{1/3}$ otherwise. This bound has the
form of \cite[(1.99)]{B-CMP122b}, the relative length of that display being taken relative also to
the old live islands inside $X$.
\item[(c)] The second group, formed by the $X$ with
$X\cap(Z_m\cup\bigcup_\iota Y_\iota)^{\boxplus}=\emptyset$, contains no island, because by
Lemma~\ref{9.1:lem-island-geometry}(b) the islands lie inside live components, which are contained
in $Z_m$; for it the bound \cite[(1.100)]{B-CMP122b} in $d_m(X)$ holds with $c_1=e^{-p_0(g_m)}$.
\item[(d)] Every $X$ of the first group (the remaining ones) is filed as
$\operatorname{fil}_m(X)\in\mathbf{D}_m^{\mathbb{I}}$ with respect to the island family
$\mathbb{I}_m$ of Lemma~\ref{9.1:lem-island-geometry} present after $\mathbf{R}_m$, obtained by
clustering the strips born at $\mathbf{R}_m$ and the islands of $\mathbb{W}'$; its excused set is
$\operatorname{fil}_m(X)\cap\mathbb{W}_m$, and
$d^{\circ}_m(\operatorname{fil}_m X)\le d_{out}^{\circ}(X)$.
\item[(e)] On the dial domain $\mathbb{D}^{\star}$, the statements (a)--(d) hold with $2c_1$ in place
of $c_1$ for the dialled activities of the dial lemma in hypothesis \textup{(iv)}; in particular
\begin{equation}\label{9.1:eq-blind-resummation-a}
\bigl|\mathbf{R}'^{(m)}(X,(\mathbf{U},\mathbf{J}))\bigr|
\le 2\,O(1)\,c_1\,e^{-(1+\frac12\beta_s)\kappa\, d_{out}^{\circ}(X)}
\quad\text{on } \mathbb{D}^{\star}.
\end{equation}
\end{itemize}
\end{lemma}

\begin{proof}
The substitutions \eqref{9.1:eq-blind-resummation-1} act only at the places where the proofs of the
statements in hypothesis (i) read a length or an anchor. We go through these places in order; at
every other place those proofs apply as they stand. Throughout the proof $k=m$ is the step under
consideration; we write $k$ where the statements in hypothesis (i) do, and $m$ where the objects of
Lemma~\ref{9.1:lem-island-geometry} do.

\emph{The strip lemma, clauses (S1)--(S3).} These clauses concern the strips only, and the
substitutions do not touch them; they hold unchanged.

\emph{Clause (S4), the escape bound by at least $m_\sharp-1$.} Its proof uses a chain of cubes across
the wall of a strip, with a cube $\Delta_1$ in layer $1$ and a cube $\Delta_2$ in layer $m_\sharp+1$,
in the layer numbering of the strip lemma.
By Lemma~\ref{9.1:lem-island-geometry}(c1) the cubes of $\mathbb{W}'$ lie in layers at least
$\check{R}_m+1$, so $\Delta_1$ and $\Delta_2$ are not in $\mathbb{W}'$. Hence the admissible tree
still contains points of $\Delta_1$ and $\Delta_2$, and these cubes lie in
$Y_4\setminus(Z\cup\mathbb{W}')$; the escape bound follows as before.

\emph{Clause (S5).} By the domain substitution the islands entering (S5) are the strips of
$Z^{\sharp}$ and the islands of $\mathbb{W}'$, and they are pairwise non-touching, which is the
property of the island family that the proof of (S5) invokes; (S5) therefore holds verbatim.

\emph{The strip-resummation theorem, clause (b).} Let $Y$ be exterior. We claim that the first cube
$\Delta_*$ of $Y\setminus(Z^{\sharp}\cup\mathbb{W}')$ exists. Indeed $Y\not\subset Z^{\sharp}$,
$Y$ contains a strip $C^{\sharp}$, and $Y$ is connected. If $Y\setminus Z^{\sharp}$ were contained in
$\mathbb{W}'$, then, leaving $C^{\sharp}$ along $Y$, one would enter a cube that touches a strip of
$Z^{\sharp}$ and lies in $\mathbb{W}'$; this is excluded, since
$\operatorname{dist}_m(\mathbb{W}',C'^{\sharp})\ge 5$ for every strip $C'^{\sharp}$ of $Z^{\sharp}$ by
Lemma~\ref{9.1:lem-island-geometry}(c1). Now let $Y_4$ belong to the fibre of
$\operatorname{fil}_\sharp$ over $Y$. Since $\operatorname{fil}_\sharp$ adds strips only,
\begin{equation*}
Y_4\setminus(Z^{\sharp}\cup\mathbb{W}') = Y\setminus(Z^{\sharp}\cup\mathbb{W}')\ni\Delta_*.
\end{equation*}
By (S2), $Y_4\subset Y$ and $Y\setminus(Z^{\sharp}\cup\mathbb{W}')\subset
Y_4\setminus(Z\cup\mathbb{W}')$; therefore
\begin{equation*}
d_{Z\cup\mathbb{W}'}(Y_4)\ge d_{Z^{\sharp}\cup\mathbb{W}'}(Y),
\end{equation*}
and by (S4) also $d_{Z\cup\mathbb{W}'}(Y_4)\ge m_\sharp-1$. Clause (S5)(c) then sums the fibre,
anchored at $\Delta_*$, as in the original proof.

\emph{Clauses (a), (c) of the strip-resummation theorem, its corollary on value units, and the
corollary on the modified bound.} Their proofs read no length; they hold unchanged.

\emph{The filing lemma for value units.} For a value unit $v$ with domain $K_v$ put
\begin{equation*}
\begin{aligned}
Y_1(v) := K_v &\cup\{\text{components of } Z \text{ meeting } K_v\}\\
&\cup\{\text{islands of } \mathbb{W}' \text{ meeting } \bar{K}_v\},
\end{aligned}
\end{equation*}
with $\bar{K}_v$ the set associated with $K_v$ in Lemma~\ref{9.1:lem-blind-counts}. Then
$Y_1(v)\in\mathbf{D}_k^{Z,\circ}$, and
$Y_1\setminus(Z\cup\mathbb{W}')=K_v\setminus(Z\cup\mathbb{W}')$, which contains the cube $\Delta_*$
for exterior $Y_1$ by the argument given for clause (b). Fix an exterior
$Y_1\in\mathbf{D}_k^{Z,\circ}$, and let $v$ range over the value units with $Y_1(v)=Y_1$. If $v$ is
plain, an optimal tree of $K_v$ meets every cube of $K_v$, and
$K_v\supseteq Y_1\setminus(Z\cup\mathbb{W}')$; hence, by the same argument as in the filing lemma,
\begin{equation*}
d_{Z\cup\mathbb{W}'}(Y_1)\le d_k(K_v).
\end{equation*}
If $v$ is an island-blind store of level $\ell\le k$, apply Lemma~\ref{9.1:lem-blind-counts}(v) with
$W'$ the union of the excused set of $v$ and of $K_v\cap Z$: its relative tree lies in
$K_v\subset Y_1$ and meets every cube of $\bar{K}_v\setminus\mathbb{W}_k$, a set containing
$Y_1\setminus(Z\cup\mathbb{W}')$; hence
\begin{equation*}
d_{Z\cup\mathbb{W}'}(Y_1)\le d^{\circ}_k(K_v).
\end{equation*}
In both cases $e^{a\,d_{Z\cup\mathbb{W}'}(Y_1)}\le e^{a\,d^{\circ}_k(K_v)}$, with $d^{\circ}_k$ read as
$d_k$ for plain units. Since every such $K_v$ contains the anchor
$\Delta_*\not\subset Z\cup\mathbb{W}'$, the value family $\mathfrak{v}^{val}$ obeys
\begin{align*}
e^{a\,d_{Z\cup\mathbb{W}'}(Y_1)}\,|\mathfrak{v}^{val}(Y_1)|
&\le\sum_{v:\,Y_1(v)=Y_1}e^{a\,d^{\circ}_k(K_v)}\,\|v\|\\
&\le\sum_{v\in\mathcal{V}^{ext},\,K_v\ni\Delta_*}e^{a\,d^{\circ}_k(K_v)}\,\|v\|
\le\mathsf{K}_{val}^{\circ}(k),
\end{align*}
and taking the supremum over exterior $Y_1$ gives
\begin{equation*}
\|\mathfrak{v}^{val}\|_a^{\circ}\le\mathsf{K}_{val}^{\circ}(k).
\end{equation*}
This is the value part of the island-blind input hypothesis.

It remains to bound $\mathsf{K}_{val}^{\circ}(k)$ independently of $g$; the value units entering it
are of two kinds, given by the value inputs of hypothesis (iii).

First, the chain units $\mathbb{C}^{(i)}$ of this step: the pieces with $i<k$ lie inside $Z$; those
with $i=k$ are plain outputs of the chain-unit corollary, unless they read an island-blind store, in
which case they are the units of the site (P) of
Proposition~\ref{9.1:prop-blind-site-expansions}, whose proof is outstanding; in either case they are
bounded by hypothesis (iii).

Second, the stored values of the boundary terms $\mathbf{B}$ and $\mathbf{B}'$. For island-blind ones
the anchored sum at a cube $\Delta\not\subset Z\cup\mathbb{W}'$ is bounded by
Lemma~\ref{9.1:lem-blind-counts}(ii) at the anchor, applied with the island union present when the
filing lemma is applied; this union is $\mathbb{W}_k^{-}$, because the strips born at $\mathbf{R}_k$
do not yet exist at that point. The hypothesis of that clause,
$\operatorname{dist}_k(\Delta,\mathbb{W}_k^{-})\ge 1$, holds: $\Delta\notin\mathbb{W}'$, and every other
island of $\mathbb{I}_k^{-}$ lies at least $3\check{R}_k-5$ cubes inside a component
$C'\subset Z(k)$ by Lemma~\ref{9.1:lem-island-geometry}(b), while $\Delta\not\subset Z$; even a cube
$\Delta$ in the collar $Z^{\sharp}\setminus Z$ is at least $3\check{R}_k-5\ge 22$ cubes from such an
island. In this
sum the crossing estimate of the filing lemma, in which the factor $e^{-\kappa\check{R}_kL^{k-j}}$
compensates the count $L^{d(k-j)}$, is replaced by the factor $y^de^{-(y-3)}$ of
Lemma~\ref{9.1:lem-blind-counts}(ii). As for the rates: by hypothesis (iii) an island-blind
$\mathbf{B}'$ of level $j<k$ carries the rate $\tfrac98\kappa$ in $d^{\circ}_j$, and
$d^{\circ}_j\ge L\,d^{\circ}_k$, so in level-$k$ units its rate is
\begin{equation*}
\tfrac98\kappa L\ge\tfrac{27}{8}\kappa\ge a+\tfrac{\kappa}{20},
\end{equation*}
using $L\ge3$ and $a=\tfrac{31}{20}\kappa$; by the same hypothesis an island-blind store born at a
$\mathbf{T}$-step carries the rate $\kappa_{out}$, and
\begin{equation*}
\kappa_{out}\ge 2(\mathfrak{c}+1)\kappa\ge a+\tfrac{\kappa}{20}.
\end{equation*}
Hence the island-blind contributions join the constant $\mathsf{K}^0$,
with $K^{\circ}(d)$ of Lemma~\ref{9.1:lem-blind-counts}(ii) in place of the counts $e^{c_d}$ of the
filing lemma:
\begin{equation*}
\mathsf{K}^0\mapsto\mathsf{K}^0+K^{\circ}(d)\bigl(B_0+2\,O(1)\bigr),
\end{equation*}
where $B_0$ is the constant of the filing lemma bounding the stored values of $\mathbf{B}$ and
$\mathbf{B}'$. The new constant is independent of $g$, and the smallness condition of the resummation
proposition is unchanged in form.

\emph{The combination lemma, clause (u1).} Unchanged.

\emph{Clause (u2).} The weight of a family under $\prod_j\mathbf{T}'_m(X_j)$ factorises over the
components of the connectivity of \eqref{9.1:eq-blind-resummation-1}: two domains overlapping only
inside $\bigcup_\iota Y_\iota^{\sharp}\cup\mathbb{W}'$ share no integration variable of this step.
Indeed, $\mathbf{T}'_m(X_j)$ acts in $X_j^{\sharp}\subset Z^{\sharp}$, which is disjoint from
$\mathbb{W}'$ by Lemma~\ref{9.1:lem-island-geometry}(c1) and from the second-class strips. Inside a
second-class strip nothing is integrated at step $m$, since the product in
\cite[(1.72)]{B-CMP122b} runs over the first class only (see also \cite[(1.102)]{B-CMP122b}). Inside a
component not in $Z(m)$ nothing is operated, in accordance with the sentence of \cite{B-CMP122a}:
\begin{quote}
[\dots]\,we do not make any changes in the operation $\mathbf{T}_k(Z_k\cap Z^c)$\,[\dots].
\end{quote}
This is Ba{\l}aban's own reason for the connectivity quoted above, applied to one more kind of fixed
region.

\emph{Clause (u3).} We claim
\begin{equation*}
d_{out}^{\circ}(X')\le\sum_{Y}\bigl(d_{Z^{\sharp}\cup\mathbb{W}'}(Y)+2\sqrt d\,\bigr)
+\sum_h\bigl(N_h+2\sqrt d\,\bigr),
\end{equation*}
the first sum running over the domains $Y$ of the family and the second over the healed boxes
$X^{\sharp}_{j_h}$ crossed; here, as in clause (u3), $X'$ is the polymer, $\mathcal{T}_Y$ the
relative tree of $Y$ and $N_h$ the length of the Hamiltonian path through the healed box
$X^{\sharp}_{j_h}$ in clause (u3).
The links are at cubes outside $Z^{\sharp}\cup\mathbb{W}'$, and these
cubes are reached by the relative trees $\mathcal{T}_Y$, which meet every cube of
$Y\setminus(Z^{\sharp}\cup\mathbb{W}')$. The Hamiltonian paths through the healed boxes
$X^{\sharp}_{j_h}$ are unchanged: a healed box containing an old island is crossed whole, since
nothing is excused inside a box. The union of these trees and paths is connected and meets every cube
of $X'\setminus(\bigcup_\iota Y_\iota^{\sharp}\cup\mathbb{W}')$, which proves the claim.

\emph{Clause (u4).} The splitting $1+2\beta=(1+\beta)+\tfrac12\beta+\tfrac12\beta$, with $\beta$ as
in clause (u4), is kept. In the
count, the statement that every admissible $Y\subset X'$ has a cube in $Y\setminus Z^{\sharp}$ is
replaced by the statement that it has a cube in $Y\setminus(Z^{\sharp}\cup\mathbb{W}')$, shown in the
step for clause (b) of the strip-resummation theorem. Clause (S5)(c) with the enlarged island family
then bounds the count by
\begin{equation*}
C_{fib}\,\alpha^{1/3}N_m\bigl(X'\setminus(Z'^{\sharp}\cup\mathbb{W}')\bigr)
\le C_{fib}\,\alpha^{1/3}N_m\bigl(X'\setminus Z'^{\sharp}\bigr),
\end{equation*}
where $Z'$ denotes the set so named in \cite[(1.93)]{B-CMP122b} and $Z'^{\sharp}$ the corresponding
set in the strip form of clause (u4).
Ba{\l}aban's volume factor counts cubes outside $\bigcup Y_i$, the present one cubes outside
$\bigcup_\iota Y_\iota^{\sharp}$. The cubes of $\mathbb{W}'\cap X'$ are not counted in favour of
$N_m(X'\setminus\bigcup_\iota Y_\iota^{\sharp})$ in either case: they lie inside $Z_m$, where the last
exponential of \cite[(1.95)]{B-CMP122b} is not produced. That factor comes from the factor
\begin{equation*}
\exp\bigl(O(1)\alpha^{1/3}M^{-d}|X'\setminus Z'|\bigr)
\end{equation*}
of \cite[(1.93)]{B-CMP122b} and from the volume factors of the components \cite{B-CMP122b}, and an
island of $\mathbb{W}'$ is in neither: no $\mathbf{T}'$ acts on it, and no value unit
$v\in\mathcal{V}$ has a domain $K_v$ meeting it except island-blind stores, whose own bound is in the
relative length.

\emph{Clauses (u5), (u6).} They concern healed boxes only and are unchanged, with the comparison
against $-\tfrac32 p_0(g_m)$ as in \cite{B-CMP122b}.

\emph{Clause (u7).} The second group is disjoint from $(Z_m\cup\bigcup_\iota Y_\iota)^{\boxplus}$,
which contains every live component and hence $\mathbb{W}'$; therefore $d_{out}^{\circ}=d_m$ there,
which gives conclusion (c). For the first group, the term $\mathbf{B}'$ is filed by
$\operatorname{fil}_m$, and Lemma~\ref{9.1:lem-blind-counts}(v) is applied with
\begin{equation*}
W' := \bigcup_\iota Y_\iota^{\sharp}\cup(\mathbb{W}'\cap X')\subset\mathbb{W}_m;
\end{equation*}
the inclusion holds because the strips born at $\mathbf{R}_m$ contain the $Y_\iota^{\sharp}$, and
$\mathbb{W}'\subset\mathbb{W}_m^{-}\subset\mathbb{W}_m$. This gives conclusion (d).

\emph{The Koteck\'y--Preiss step} \cite{KP86} of the resummation proposition, with $c_1\le c_{KP}$:
the sums are anchored at cubes of a polymer outside $\bigcup_\iota Y_\iota^{\sharp}\cup\mathbb{W}'$.
Every polymer has such a cube, since it contains an exterior $Y$, which has one by the step for
clause (b). Clause (S5)(c) then applies, and the step is unchanged in form.

With these places accounted for, the statements in hypothesis (i) hold after the substitutions
\eqref{9.1:eq-blind-resummation-1} with the same constants, which is conclusion (a), and the bound of
the resummation proposition becomes \eqref{9.1:eq-blind-resummation-2}, which is conclusion (b).

\emph{The dialled form (e).} By clause (D)(e) of the dial lemma, the activities of
\cite[(1.91)]{B-CMP122b} obey \cite[(1.97)]{B-CMP122b} with $2c_1$ on $\mathbb{D}^{\star}$. By the
linearity corollary for the filing maps, $\operatorname{fil}_\sharp$ and $\operatorname{fil}_m$ are
linear, they move domains and never anchors, and the sup-balls of the strip forms of
\cite[(1.68), (1.69)]{B-CMP122b} are convex. Hence the argument above applies on the dial with $2c_1$
in place of $c_1$, which gives \eqref{9.1:eq-blind-resummation-a}.
\end{proof}

\begin{proposition}[Large-field boundary terms with island-blind domains]
\label{9.1:prop-blind-boundary-production}
Assume the following.
\begin{itemize}
\item[(i)] The hypotheses of Proposition~\ref{9.1:prop-large-field-boundary} (the two-run
bound for boundary terms of the large-field step) hold. These are the hypotheses of parts
(A)--(D) of the theorem on the dial domain $\mathbb{D}^{\star}$, taken at the
$\mathbf{R}$-site, together with the further hypotheses listed in that proposition, one of
which is imposed for levels at most $i$.
\item[(ii)] The conclusions of Lemma~\ref{9.1:lem-blind-resummation} (the large-field
resummation with island-blind units) hold, including the relative bound among the displays
\eqref{9.1:eq-blind-resummation-a}. These conclusions are used in place of clauses (c)--(e)
of the resummation chain for the ordinary localisation domains.
\end{itemize}
Let $\mathbf{B}'^{(i)}(X')$ be a boundary term of the first group of an aligned label. Assume
that its domain $\operatorname{fil}_i(X')$ under the filing map is island-blind. This is the
case when $X'$ contains a second-class strip $s(Y_\iota,i)$, or $X'$ contains an island of
$\mathbb{W}'_i$, or $X'$ merely meets $\mathbb{W}_i$. In the last case $\operatorname{fil}_i$
excuses the islands that $X'$ meets, and the resulting bound is weaker and remains true. Let
$n\le i-1$ be a zone. Then the following hold.
\begin{itemize}
\item[(a)] The statements (zs1$'$) and (zs2$'$) of the display
\eqref{9.1:eq-large-field-boundary-a} hold as in
Proposition~\ref{9.1:prop-large-field-boundary}. Here the classes whose slots lie inside
second-class or arrested components are counted in $\mathcal{Z}_{own}$.
\item[(b)] In the relative weighted sup-norm, whose weight is
$e^{\kappa_{\mathbf{B}}d^{\circ}_i(\operatorname{fil}_iX')}$,
\begin{equation}\label{9.1:eq-blind-boundary-production-1}
\mathsf{S}^{\circ,(n)}_i\bigl[\mathbf{B}'^{res}\bigr]
\le \varkappa^{\mathbf{B}'}_i\,[\![\delta]\!]^{\natural+,(i)}_n ,
\qquad
\varkappa^{\mathbf{B}'}_i:=\frac{2\,c_1^{\flat}(g_i)}{E_1^{z}} .
\end{equation}
This is the same coefficient $\varkappa^{\mathbf{B}'}_i$ as in
Proposition~\ref{9.1:prop-large-field-boundary}. The rate is that of
Lemma~\ref{9.1:lem-blind-resummation}, namely
\begin{equation*}
\tfrac98\,\kappa\ \ge\ \kappa_{\mathbf{B}} ,
\end{equation*}
measured in the island-blind length $d^{\circ}_i$. Consequently the constant
$\bar{\varkappa}'$ of Proposition~\ref{9.1:prop-large-field-boundary} is unchanged.
\item[(c)] The quantity $\mathsf{S}^{\circ}_i\bigl[\mathbf{B}'^{own}\bigr]$ is bounded, in the
relative sup-norm, by the display of clause (D)(e) of the theorem on the dial domain taken at
$c_1^{\flat}$. This is the bound (own) of \eqref{9.1:eq-large-field-boundary-a}, verbatim.
\end{itemize}
\end{proposition}

\begin{proof}
We take the three conclusions in turn.

\emph{Conclusion \textup{(a)}.} The proofs of (zs1$'$) and (zs2$'$) given for
Proposition~\ref{9.1:prop-large-field-boundary} use no length of a domain. They therefore
apply without change when $\operatorname{fil}_i(X')$ is island-blind. The classes whose slots
lie inside second-class or arrested components are placed in $\mathcal{Z}_{own}$ by the
statement that fixes the class sets $\mathcal{Z}_{surv}$ and $\mathcal{Z}_{own}$ for these
boundary terms.

\emph{Conclusion \textup{(b)}.} We follow the proof of the bound (res) of
Proposition~\ref{9.1:prop-large-field-boundary}, with one change. On the dial domain
$\mathbb{D}^{\star}$ we apply part (a) of the first-variation bound with class radius
\begin{equation*}
\Lambda_z=\frac{E_1^{z}}{[\![\delta]\!]^{\natural+,(i)}_n}
\end{equation*}
and modulus bound
\begin{equation*}
\mathfrak{B}_z:=\sup_{|\varepsilon|\le\Lambda_z}
\bigl|\mathcal{P}_{X'}\bigl[x^{\lambda}+\varepsilon W^{(z)}\bigr]\bigr| .
\end{equation*}
By the relative bound among the displays \eqref{9.1:eq-blind-resummation-a} of
Lemma~\ref{9.1:lem-blind-resummation},
\begin{equation}\label{9.1:eq-blind-boundary-production-2}
\mathfrak{B}_z\le 2\,O(1)\,c_1^{\flat}(g_i)\,e^{-\frac98\kappa\, d^{\circ}_{out}(X')} .
\end{equation}
Two readings enter this bound. First, Ba{\l}aban's constant $c_1$ of
\cite[(1.99)]{B-CMP122b} is read as $c_1^{\flat}$ under the strip, the radius $\alpha$ being
replaced by $\alpha^{\sharp}$.
Second, the factor $2$ is the one supplied by clause (e) of the lemma from which $c_1^{\flat}$
is taken. Part (a) of the first-variation bound then gives, for the first-variation member of
class $z$,
\begin{equation}\label{9.1:eq-blind-boundary-production-3}
\bigl|\mathfrak{V}^{(z)}(X')\bigr|
\le \frac{2\,O(1)\,c_1^{\flat}(g_i)}{E_1^{z}}\,[\![\delta]\!]^{\natural+,(i)}_n\,
e^{-\frac98\kappa\, d^{\circ}_{out}(X')} .
\end{equation}
The right-hand side of \eqref{9.1:eq-blind-boundary-production-3} equals the right-hand side
of \eqref{9.1:eq-blind-boundary-production-2} divided by $\Lambda_z$.

Next we compare weights. By the last clause of Lemma~\ref{9.1:lem-blind-resummation}, the
island-blind length of the assigned domain is at most the island-blind output length of
$X'$. Since moreover $\kappa_{\mathbf{B}}\le\frac98\kappa$, we obtain
\begin{equation}\label{9.1:eq-blind-boundary-production-4}
e^{\kappa_{\mathbf{B}}d^{\circ}_i(\operatorname{fil}_iX')}
\le e^{\kappa_{\mathbf{B}}d^{\circ}_{out}(X')}
\le e^{\frac98\kappa\, d^{\circ}_{out}(X')} .
\end{equation}
Multiplying \eqref{9.1:eq-blind-boundary-production-3} by the weight and using
\eqref{9.1:eq-blind-boundary-production-4}, we find
\begin{equation*}
e^{\kappa_{\mathbf{B}}d^{\circ}_i(\operatorname{fil}_iX')}\bigl|\mathfrak{V}^{(z)}(X')\bigr|
\le \frac{2\,O(1)\,c_1^{\flat}(g_i)}{E_1^{z}}\,[\![\delta]\!]^{\natural+,(i)}_n .
\end{equation*}
Taking the supremum that defines $\mathsf{S}^{\circ,(n)}_i$ gives
\eqref{9.1:eq-blind-boundary-production-1}.

The convergence of $\sum_t c_1^{\flat}(g_t)$ used for
Proposition~\ref{9.1:prop-large-field-boundary} is unchanged. Hence $\bar{\varkappa}'$ is
unchanged.

\emph{Conclusion \textup{(c)}.} The proof of clause (D)(e) of the theorem on the dial domain
bounds the term pointwise on the dial:
\begin{equation*}
|\mathbf{B}'(X';\lambda,\mu)|\le\mathfrak{B}_{X'} .
\end{equation*}
Here $\mathfrak{B}_{X'}$ is the relative bound among the displays
\eqref{9.1:eq-blind-resummation-a} of Lemma~\ref{9.1:lem-blind-resummation}. The remaining
ingredients of that proof use no length of a domain. These ingredients are the age factors,
geometric forgetting of old regions, the $\mu$-dial and the remainder term $\Delta^{rest}$ of
that proof. The
display of clause (D)(e), taken at $c_1^{\flat}$, therefore bounds
$\mathsf{S}^{\circ}_i\bigl[\mathbf{B}'^{own}\bigr]$ in the relative sup-norm.
\end{proof}

\begin{remark}[The island-blind weight of boundary terms produced by the large-field operation]
Let a boundary term produced by $\mathbf{R}_i$ have an island-blind domain. It is stored with
the weight $e^{\kappa_{\mathbf{B}}d^{\circ}}$, where $d^{\circ}$ is the length off the live
islands that the domain contains; this is Ba{\l}aban's relative length $d_{i,Z}$ of
Notation~\ref{0.2:not-glossary}, evaluated at the domain, with $Z$ the union of the live
islands contained in the domain. Its constant is that of
Proposition~\ref{9.1:prop-large-field-boundary}. So are its summability over levels, its
splitting into the parts $\mathbf{B}'^{res}$ and $\mathbf{B}'^{own}$, its filing, and its
place, as numbers, in the loads $\mathsf{y}_n$ and $\mathsf{y}'_n$.

No count is taken at production, because $\mathsf{S}^{\circ}$ is a supremum. In particular
no power of $\check{R}$ multiplies $c_1^{\flat}$. The summability condition on the constants
$c_1^{\flat}(g_t)$ used in the proof of Proposition~\ref{9.1:prop-large-field-boundary} is
not changed. Anchoring a sum at a cube of some level $\ell$ inside a strip would cost a factor
\begin{equation*}
2d\,(101\,\check{R}_\ell+1)^{d-1} ;
\end{equation*}
the sums of the island-blind format are never so anchored.
\end{remark}

\begin{proposition}[Site expansions with island-blind inputs. Proof outstanding]
\label{9.1:prop-blind-site-expansions}
Consider any renormalisation step and any one of the three sites \textup{(T)}, \textup{(L)},
\textup{(P)} at which the correlation theorem performs its expansions: the site \textup{(T)} of the
expansion at a $\mathbf{T}$-step; the site \textup{(P)} of the preliminary integrations, which
produce the chain terms $\mathbb{C}^{(i)}$; and the site \textup{(L)}, whose expansion produces the
input bounds on the bracket terms and on the volume terms. By the expansion statement of a site we
mean the following.
\begin{itemize}
\item[\textup{(T)}] The cluster expansion lemma of the site \textup{(T)}, together with its
corollary on the marked sub-sum and the tolerance lemma, and clauses (i)--(vii) of the
$\mathbf{T}$-step expansion lemma.
\item[\textup{(L)}] The expansion lemma of the site \textup{(L)}, together with the summation of
boundary terms at finer zones, the boundary lemma and the bracket lemma. These together produce the
input bound on the bracket terms $\mathfrak{v}^{br}$ and the input bound on the volume terms.
\item[\textup{(P)}] Part (a) of the expansion proposition of the site \textup{(P)} and its
corollary, which concern the chain terms $\mathbb{C}^{(i)}$, with the norm $\|\cdot\|^{A}$ of that
corollary taken at $a_k$ and at $a_m$; and the remaining chain terms $\mathbb{C}^{(i)}$ of arrested
chains.
\end{itemize}
Then the expansion statement of the site holds also when some of its input terms are island-blind
stored terms, defined next; its outputs are then as in (a), (b) below, and (c) lists the terms that
are never island-blind.

An island-blind stored term of level $\ell$ is a stored term in the island-blind format, given by
the following data:
\begin{itemize}
\item its filed domain $Y$ lies in $\mathbf{D}_\ell^{\mathbb{I}}$, and its excused set is
$Y\cap\mathbb{W}_\ell$;
\item its size in one run is measured in the norm $\mathsf{N}^{\circ}_\ell$;
\item its differences between the two runs $\mathsf{A}$, $\mathsf{B}$ are measured by
$\mathsf{S}^{\circ}$.
\end{itemize}
Such a term is read at level $m$ through the filing map $\operatorname{fil}_m$ and through the
one-layer hull, the fat core and the collar of the site. These lie inside the four spare layers of
$\mathbb{W}_m$. Consequently the hull inequality, for the one-layer enlargement $Y^{(1)}$ of $Y$,
\begin{equation*}
Y\subset Y'\subset Y^{(1)}\ \Longrightarrow\ d(Y')\le c_g\bigl(d(Y)+1\bigr)
\end{equation*}
holds, by its own proof carried out off $\mathbb{W}_m$, in the form
\begin{equation*}
d_{m,W}(Y')\le c_g\bigl(d^{\circ}(Y)+1\bigr),
\end{equation*}
where $d_{m,W}$ denotes the tree length relative to $\mathbb{W}_m$. We write $\mathbb{W}$, without
index, for the set $\mathbb{W}_m$ at the level $m$ read.

The outputs are as follows.
\begin{itemize}
\item[(a)] An output whose term reads no island-blind input term is unchanged.
\item[(b)] Let an output's term read an island-blind input term. Such an output, with domain $X$,
is filed under $\operatorname{fil}(X)$, and $\operatorname{fil}(X)$ lies in $\mathbf{D}^{\mathbb{I}}$
of the output's level. The output obeys the bound of the site's expansion statement with every tree
length $d_j(\cdot)$ replaced by the island-blind length $d^{\circ}_j(\cdot)$, as follows.
\begin{itemize}
\item At \textup{(T)}, the factor $e^{-\kappa_{out}d_{k+1}(X)}$ becomes
$e^{-\kappa_{out}d^{\circ}_{k+1}(X)}$. This applies in the outputs of the cluster expansion lemma,
in the marked sub-sum of its corollary, and in clauses (ii) and (vi) of the $\mathbf{T}$-step
expansion lemma.
\item At \textup{(L)}, the bound reads $\|\mathfrak{v}^{br}\|_a^{\circ}\le\varepsilon_{br}$.
\item At \textup{(P)}, the bound reads $\|\mathbb{C}^{(n)}\|^{A,\circ}\le C_0^{\sharp}$, with $d_m$
replaced by $d^{\circ}_m$ in the weight of the norm.
\end{itemize}
The constants and the rates are those of the site's statement; the floors are those of the site's
statement with the per-cell domination \eqref{9.1:eq-blind-domination-a} added.
\item[(c)] The following are never island-blind: the $E$-, $D$- and $R''$-families, the Gaussian
part of the second expectation, and the $\mathbf{R}'$-terms of the second group. The domain of each
of them lies in $\Lambda_{k+1}$ or off $(Z\cup\bigcup_\iota Y_\iota)^{\boxplus}$, a set containing
$\mathbb{W}$.
\end{itemize}
Finally, for an island-blind potential the activity hypothesis of the cluster expansion lemma of
the site \textup{(T)} is its anchored relative activity at the cells of the site. By
Lemma~\ref{9.1:lem-blind-counts}\,(ii) (domain counts in the island-blind length), this activity
is at most $K^{\circ}(d)$, $d=4$, times its sup-norm.
\end{proposition}

Proof outstanding.

\emph{Route.} We first explain why the outputs must themselves be island-blind. The argument
is written for a $\mathbf{T}$-step.

Let $s=s(C,i)$ be the birth strip of a birth $(C,i)$, and suppose the expansion at a
$\mathbf{T}$-step reads an island-blind stored term whose domain $X'$ satisfies $X'\supseteq s$.
The expansion then produces outputs with
$X\supseteq X'\supseteq s$. A bound for these outputs in the full length would need the tree
through the strip. At the next level that tree has length at least
\begin{equation*}
\frac{N(s)}{\mathsf{c}_d}\gtrsim\frac{(\check{R}_i/L)^d}{\mathsf{c}_d},
\end{equation*}
where $N(s)$ is the number of cubes of $s$ at the next level.
Against this stand two factors. The first is the small factor
$(\alpha^{\sharp})^{1/3}\asymp e^{-\beta_\sharp\kappa\check{R}_i/12}$ carried by the stored term.
The second is the
factor $e^{-O(\kappa_1\check{R})}$ of the moat. Both are only linear in $\check{R}$ in the
exponent, so the length cannot be paid for young strips. For ages
$\ge\log_L(101\check{R}_i)$ the deficit is $O(1)$. A remedy by age, however, would require the
bounds to read $\check{R}$, and the format does not provide this.

Hence the outputs must themselves be island-blind. This is the alternative inductive assumption of
\cite{B-CMP119}, which states: ``All constructions and proofs of
the procedure can be carried on with these inductive assumptions.'' Clause (ii) of the
$\mathbf{T}$-step expansion lemma states that the estimates of \cite[(2.27)--(2.41)]{B-CMP116}
produce the output exponential for inputs of the format of (1.42), (1.43) and (1.36) of the same
paper, the last read on the fat core of the site; this is a format in the full length.

A proof of Proposition~\ref{9.1:prop-blind-site-expansions} would repeat the expansion of
\cite[\S2]{B-CMP116}, and its analogues at \textup{(P)} and \textup{(L)}, for a single run, that is
without comparing two runs, with island-blind inputs, and would establish the following five
points.
The geometry these points use is proved in
Lemma~\ref{9.1:lem-island-geometry} (moat and position of the islands around live strips) and
Lemma~\ref{9.1:lem-blind-counts}.

(i) \emph{Volume factors.} Every decoupling cube, every fluctuation bond and every origin of a
random walk lies in the region $\Omega_{m+1}$ of Lemma~\ref{9.1:lem-island-geometry}. Each lies at
distance at least $3\check{R}_m-4$ cubes from
every island, by Lemma~\ref{9.1:lem-island-geometry}\,(a),(d). The volume factors of
\cite{B-CMP116} count fluctuation volume only; in them, in the notation of \cite[\S2]{B-CMP116},
$|Z|$ is the fluctuation volume. These factors are:
\begin{itemize}
\item the factor $\exp O(1)\alpha_5|Z|$ over fluctuation bonds in \cite[(2.25)--(2.26)]{B-CMP116},
with the constant $\alpha_5$ of that paper;
\item the sums over the walk cubes $Z_0\setminus Y_0$ in \cite[(2.31)--(2.34)]{B-CMP116}, the sets
of cubes $Z_0$ and $Y_0$ being those so denoted there;
\item the last exponential in \cite[(2.37)]{B-CMP116}.
\end{itemize}
The fluctuation volume lies off $\mathbb{W}$ and is bounded by the relative tree length. Indeed,
\cite[(2.30)]{B-CMP116} in the form
\begin{equation*}
M^{-d}|Z\setminus\mathbb{W}|\le\mathsf{c}_d\bigl(d^{\circ}(Z)+1\bigr)
\end{equation*}
is the volume bound of the island-blind format. Moreover, the inequality
$2d_k\ge Ld_{k+1}$ of \cite[(2.36)]{B-CMP116} holds for the relative lengths, with
$\mathbb{W}_{k+1}\supseteq\mathsf{cl}_{k+1}(\mathbb{W}_k)$; this is
Lemma~\ref{9.1:lem-blind-counts}\,(v).

(ii) \emph{Domain sums.} Three kinds of domain sums occur:
\begin{itemize}
\item the sums of \cite[(2.29)]{B-CMP116}, of which \cite{B-CMP116} says ``Inequalities of this
type were proved many times'';
\item the sums over $n$ components after \cite[(2.36)]{B-CMP116};
\item the sums carrying the weight $e^{(1-\frac12\delta)\kappa d(\bar S_v)}$ of the expansion lemma
of the site \textup{(L)}, with the parameter $\delta$ and the sets $\bar S_v$ of that lemma.
\end{itemize}
All of them are anchored at polymer cubes. When anchored off $\mathbb{W}$, they are bounded by
Lemma~\ref{9.1:lem-blind-counts}\,(i),(ii) at $\lambda_0\ge\mathsf{k}_d+1$. This replaces the
factor $e^{c_d(n+1)}$ of \cite[(1.26)]{B-CMP116}, with the constant $c_d$ of that paper; the margin
is contained in the condition $\delta\kappa\ge2(\mathsf{k}_d+2)$.

A polymer cube inside an island occurs only as a cube of a random walk that crossed the moat. The
terms attached there are enumerated by their true supports containing that cube. Note that the
domain a stored term receives from the filing map may exceed its support inside islands.
Attachment is by support, while counting is by the filing domain. The count is
Lemma~\ref{9.1:lem-blind-counts}\,(ii), taken from the last off-island cube of the walk. To it is
added the walk's own enumeration, $3^d$ per step, which is paid by $e^{-(\kappa_1-1)}$.

(iii) \emph{Tree composition.} The composition inequality of \cite[(2.27)]{B-CMP116},
\begin{equation*}
\sum_D\bigl(d(Y)+5\bigr)\ge d(Y_0)+5,
\end{equation*}
holds in relative form when the Mayer links and the decoupling links sit at cubes off
$\mathbb{W}$, since each relative tree reaches such cubes. Links through the fluctuation field sit
off $\mathbb{W}$ by (i).

A junction at an island cube $c$, reached by a walk from $\Omega_{m+1}$, is re-routed inside the
island to the nearest off-island cube of the support of the term, at most $\mathfrak{j}_m$ cubes
away, where
\begin{equation*}
\mathfrak{j}_m:=d\cdot(101\check{R}_m+8)+1 .
\end{equation*}
Here $d=4$ is the dimension. This bound holds because the strips of an island have closures of
side at most
\begin{equation*}
101\check{R}_iL^{-(m-i)}+10\le101\check{R}_m+10,
\end{equation*}
since $L\check{R}_{j+1}\ge\check{R}_j$ for every $j\ge i$ by
Lemma~\ref{9.1:lem-island-geometry}\,(a), whence $\check{R}_iL^{-(m-i)}\le\check{R}_m$. A clustered
island is re-routed strip by strip along the walk, which visits those strips.

The re-routing is paid at the storage rate $\kappa_{\mathbf{B}}+\sigma_*$, not at
$\kappa_{out}$, by the moat factor $e^{-(\kappa_1-1)(3\check{R}_m-4)}$ of that walk. The required
inequality, and the conditions from which it follows, are
\begin{align*}
(\kappa_1-1)(3\check{R}_m-4)\ge(\kappa_{\mathbf{B}}+\sigma_*)\,\mathfrak{j}_m
&\ \Longleftarrow\ 4L\kappa\,(3\check{R}_m-4)\ge2\kappa\bigl(d\cdot(101\check{R}_m+8)+1\bigr)\\
&\ \Longleftarrow\ L\ge18(d+1).
\end{align*}
These implications use $\kappa_1\ge4L\kappa+1$ (from the bracket lemma), $\sigma_*\le\kappa$, and
$\check{R}_m\ge9$, which is \eqref{9.1:eq-island-geometry-b}. The floor $L\ge18(d+1)$ depends on $d$
only. It is implied by the floor $L\ge4(1+4\beta^{\star})(c_g+1)c_g^2$ with $c_g>31$.

Finally, a Mayer link between two terms that overlap only inside an island is declared a non-link,
since the two terms share no integration variable there. This is the device of
\cite{B-CMP122b}, also used in the relative localisation statement. With this convention, no
junction without a walk occurs inside an island.

(iv) \emph{The sites \textup{(P)} and \textup{(L)}.}
\begin{itemize}
\item At \textup{(P)}, the points above hold in the same way with the preparatory covariances of
\cite[\S1]{B-CMP122a}. The collar cells lie off the islands of their own scale, by
Lemma~\ref{9.1:lem-island-geometry}\,(d).
\item At \textup{(L)}, the pieces of the expansion lemma of the site \textup{(L)} attach at
decoupling cubes in the region of the site \textup{(L)}: inside the collars of $Z(m)^{\sharp}$ and
at near points. These cubes lie
off $\mathbb{W}'$ by Lemma~\ref{9.1:lem-island-geometry}\,(c1). Inside a first-class component
nothing is excused, by Lemma~\ref{9.1:lem-island-geometry}\,(c2).
\end{itemize}

(v) \emph{Hull, fat core and collar.} Each is one layer, and all lie inside the four spare layers
of $\mathbb{W}_m$. For each of them the hull inequality holds in the relative form displayed in the
statement, by the tour-and-neighbours argument of its proof carried out off $\mathbb{W}_m$.

\emph{The alternative representation.} The representation
\cite[(1.103)--(1.104)]{B-CMP122b} integrates all the second-class histories in one step, with the
factors $\mathbf{T}'_k(Y_i)1$ brought in front, and re-exponentiates the terms of the first group.
Reading the activities of \cite[p.~391]{B-CMP122b}, one finds in them a factor $e^3$ for each
$Y_i\subset X'$, because the normalised operations $(\mathbf{T}'1)^{-1}\mathbf{T}'$ are of order
one, not small. One also finds the length $d_k(X'\setminus\bigcup_iY_i)$.

Hence the new terms $\mathbf{B}'$ are still relative. Moreover, the small factors
$\mathbf{T}'_k(Y_i)1$ are density factors shared by all stored terms over $Y_i$, and are not
available term by term. How the density is shared between the factors $\mathbf{T}'_k(Y_i)$ and the
terms $\mathbf{B}'$ is part of the choice of scaffold, so the scaffold could be chosen according to
\cite[(1.104)]{B-CMP122b}. That choice would not remove the need for
Proposition~\ref{9.1:prop-blind-site-expansions}, and it is not adopted in this book.

Throughout the lemma below and its proof we work in the setting of the per-cell domination of the resolved boundary members (Lemma~\ref{9.1:lem-per-cell-domination}), under its hypotheses.

Stores of both kinds are admitted:
\begin{itemize}
\item those born at the $\mathbf{T}'$-step, whose class-$n$ members of level $\ell$ and support $X$ are the $\delta\mathbf{B}^{(\ell),(n)}(X)$;
\item those born at the $\mathbf{R}$-step, whose class-$n$ members are the $\delta\mathbf{B}'^{(\ell),(n)}(X)$.
\end{itemize}
The loads $\mathsf{y}_n$, $\mathsf{y}'_n$ and $[\![\delta]\!]^{\natural+}_n$ of zone $n$ are those of Lemma~\ref{9.1:lem-per-cell-domination}. An island-blind store (blind, for short) of level $\ell$ has a domain $X\in\mathbf{D}_\ell^{\mathbb{I}}$ and relative length $d^{\circ}_\ell(X)$. Its size enters only through the weighted sup-norms $\mathsf{S}^{\circ,(n)}_\ell[\mathbf{B}]$ and $\mathsf{S}^{\circ,(n)}_\ell[\mathbf{B}'^{res}]$, which are numbers.

Among the estimates on which Lemma~\ref{9.1:lem-per-cell-domination} and its consequences rest, only the following read the length of a support, and they are therefore the only ones that change when island-blind stores are admitted:
\begin{itemize}
\item clauses (a) and (b) of that lemma, together with its crossing weight;
\item the ball in clause (c2)$^{\star}$ of the zone-resolved cluster lemma;
\item the factor $e^{-\kappa_{out}d}$ in clauses $(\alpha)$ and $(\beta)$ of the zone-resolved volume and filament bound.
\end{itemize}
The elimination of the $\mathbf{B}$-terms, and the statements derived from it, read only sizes; like every other estimate that reads sizes only, they are unchanged. The lemma below is the counterpart of clauses (a) and (b) of Lemma~\ref{9.1:lem-per-cell-domination} for island-blind stores.

\begin{lemma}[Per-cell domination for island-blind stores]\label{9.1:lem-blind-domination}
Assume the following two statements.
\begin{itemize}
\item The site expansion at the site (T) holds when island-blind stores are among its input units. This is the conclusion of Proposition~\ref{9.1:prop-blind-site-expansions} at (T); the proof of that proposition is outstanding.
\item The per-cell and per-block activity sums for the marks of the stored $\mathbf{B}$-terms hold at the sites (T), (L) and (P) of Proposition~\ref{9.1:prop-blind-site-expansions}. These are the sums on which clause (b) of Lemma~\ref{9.1:lem-per-cell-domination} rests, namely: the bound of the correlation theorem for the large-field boundary terms, with the statement that the sum over the supports $X$ is the norm itself; the per-cell activity bound for those marks; the per-block fibre count at the site (P), with its margin $\sigma_*$; and the bound on the sums $w^{\mathbf{B}}_{n,\ell}$ of the crossing weights. Each is a sum, over the stored terms whose support contains a given cell $\Delta$, of $e^{-\kappa_{\mathbf{B}}d}$ times the size of the term, and holds by the geometry of the supports alone.
\end{itemize}
Then, at reading level $k'$, on an aligned label and for the class $(n,k')$, the following hold.

\smallskip\noindent
(a) \emph{Values.} Let $Y$ be a resummation or output domain: a strip $C^{\sharp}$, a domain of the form $V(Y')$ as in \cite[(1.69)]{B-CMP122b}, or a component of the region of the operation $\mathbf{T}'$. Put $K_{\circ}(d):=3^{d+1}e^{8}C_{fib}$, and let $\mathsf{y}_n^{\circ}$ be the blind part of $\mathsf{y}_n$. Then the island-blind class-$n$ members supported in $Y$ satisfy
\begin{equation}\label{9.1:eq-blind-domination-1}
\sum_{\ell\ge n}\ \sum_{\substack{X\in\mathbf{D}_\ell^{\mathbb{I}}\ \text{blind}\\ X\subset Y}}
\Bigl(|\delta\mathbf{B}^{(\ell),(n)}(X)|+|\delta\mathbf{B}'^{(\ell),(n)}(X)|\Bigr)
\le K_{\circ}(d)\,\mathsf{y}_n^{\circ,(\le k')}\sum_j|\Gamma^0_j\cap Y| .
\end{equation}
Consider next the members anchored at a cell $c\in\Gamma^0_n$ with $c\not\subset\mathbb{W}_{k'}$. They obey the per-cell inequality of Lemma~\ref{9.1:lem-per-cell-domination}(a) with $K(d)$ replaced by $K^{\circ}(d)$, the constant of Lemma~\ref{9.1:lem-blind-counts}(ii). Consequently the value budget of the large-field resummation, which charges per unit of $\sum_j|\Gamma^0_j\cap Y|$, is met with $K(d)$ replaced by $K(d)+K_{\circ}(d)+K^{\circ}(d)$.

\smallskip\noindent
(b) \emph{Activities.} Let $\Delta$ be a site cell of level $k'$ of one of the following kinds:
\begin{itemize}
\item at (T), a cell $\Delta\subset Y_0$, where $Y_0\subset\Lambda_{k'+1}$ is the region of the site expansion at (T);
\item at (L), an exterior anchor, that is, a cell with $\Delta\not\subset Z\cup\mathbb{W}'$;
\item at (P), a collar cell of scale at least the level of the store.
\end{itemize}
The activity sums of the second hypothesis hold for the island-blind class-$(n,k')$ members, with the following changes:
\begin{itemize}
\item the support length $d(X)$ is replaced by the relative length $d^{\circ}_\ell(X)$;
\item the sizes are
\begin{equation*}
\mathsf{x}^{\circ,(n)}_\ell:=\mathsf{w}_{\ell,n}\,\mathsf{S}^{\circ,(n)}_\ell[\mathbf{B}]
+\mathsf{S}^{\circ,(n)}_\ell[\mathbf{B}'^{res}] ;
\end{equation*}
\item the rate is $\kappa_{\mathbf{B}}-1$.
\end{itemize}
These sums are bounded by $K^{\circ}(d)\,K_{site}\,\mathsf{y}_n^{\circ,(\le k')}$. Clause (b) is asserted under the floor
\begin{equation}\label{9.1:eq-blind-domination-a}
\sigma_*\ \ge\ \mathsf{k}_d+4 ,
\end{equation}
where $\sigma_*$, the margin of those activity sums, is one of $\beta\kappa$ and $\kappa/(4c^{\sharp})$, $\beta$ being the fixed fraction of $\kappa$ that the per-block fibre count at the site (P) retains as its margin. This floor is a condition on $\kappa$ depending on the dimension $d$ only; it is imposed beside the floor \eqref{9.1:eq-per-cell-domination-a} of Lemma~\ref{9.1:lem-per-cell-domination}. Finally, the constant $K_Y$ of Lemma~\ref{9.1:lem-per-cell-domination} is replaced by
\begin{equation*}
K_Y^{\circ}:=\bigl(K(d)+K_{\circ}(d)+K^{\circ}(d)\bigr)\max K_{site} ,
\end{equation*}
the maximum being taken over the sites (T), (L) and (P).
\end{lemma}

\begin{proof}
\emph{(a), grouping by islands.} Let $X\in\mathbf{D}_\ell^{\mathbb{I}}$ be a blind domain of level $\ell$ with $X\subset Y$. Then $X$ contains an island $I\in\mathbb{I}_\ell$ with $I\subset Y$. Fix an order on $\mathbb{I}_\ell$ and assign to $X$ the first island $I\subset X$. This groups the domains by the pair $(\ell,I)$, each $X$ lying in exactly one group.

Write $\delta^{(n)}(X)$ for $\delta\mathbf{B}^{(\ell),(n)}(X)$ or $\delta\mathbf{B}'^{(\ell),(n)}(X)$. Write $\mathsf{S}^{\circ,(n)}_\ell$ for $\mathsf{S}^{\circ,(n)}_\ell[\mathbf{B}]$ or $\mathsf{S}^{\circ,(n)}_\ell[\mathbf{B}'^{res}]$ accordingly. By the definition of the weighted sup-norm, for fixed $(\ell,I)$,
\begin{equation*}
\sum_{X\supseteq I}|\delta^{(n)}(X)|\le\mathsf{S}^{\circ,(n)}_\ell
\sum_{X\supseteq I}e^{-\kappa_{\mathbf{B}}d^{\circ}_\ell(X)} .
\end{equation*}

\emph{(a), stores born at the $\mathbf{R}$-step.} For the $\mathbf{B}'^{res}$-stores we apply Lemma~\ref{9.1:lem-blind-counts}(iii) with $\lambda_0:=\kappa_{\mathbf{B}}=\kappa$. We need $\kappa\ge\mathsf{k}_d+2$. This holds because the lower bound on $\kappa$ under which the correlation theorem holds contains $\kappa/20\ge2(\mathsf{k}_d+1)$, whence $\kappa\ge40(\mathsf{k}_d+1)\ge\mathsf{k}_d+2$. Let $N_\ell(I)$ be the number of level-$\ell$ cubes of $I$. Lemma~\ref{9.1:lem-blind-counts}(iii) gives
\begin{equation*}
\sum_{X\supseteq I}|\delta^{(n)}(X)|
\le\bigl(1+3^dC_{fib}N_\ell(I)\bigr)\mathsf{S}^{\circ,(n)}_\ell[\mathbf{B}'^{res}]
\le3^{d+1}C_{fib}\sum_j|\Gamma^0_j\cap I|\,\mathsf{S}^{\circ,(n)}_\ell[\mathbf{B}'^{res}].
\end{equation*}

\emph{(a), stores born at the $\mathbf{T}'$-step.} For these we first take out the crossing. Let $n\le\ell-2$. A non-zero class-$n$ member sits on a $\mathbf{T}'$-born $X$ with $d^{\circ}_\ell(X)\ge\mathsf{d}_{\ell-1}-8$, by Lemma~\ref{9.1:lem-blind-counts}(iv). Hence
\begin{equation*}
e^{-\kappa_{\mathbf{B}}d^{\circ}_\ell(X)}
\le e^{8}\,\mathsf{w}_{\ell,n}\,e^{-(\kappa_{\mathbf{B}}-1)d^{\circ}_\ell(X)} ,
\end{equation*}
with the crossing weight $\mathsf{w}_{\ell,n}$ of Lemma~\ref{9.1:lem-per-cell-domination}. For $\ell\le n+1$ we have $\mathsf{w}_{\ell,n}=1$, and the same inequality holds trivially, since $e^{-\kappa_{\mathbf{B}}d^{\circ}_\ell(X)}\le e^{-(\kappa_{\mathbf{B}}-1)d^{\circ}_\ell(X)}$ and $e^{8}\ge1$; so it holds for all $\ell\ge n$. Then we apply Lemma~\ref{9.1:lem-blind-counts}(iii) with $\lambda_0:=\kappa_{\mathbf{B}}-1$, which satisfies $\lambda_0\ge\mathsf{k}_d+1$ by $\kappa\ge\mathsf{k}_d+2$. The group of $(\ell,I)$ therefore contributes at most
\begin{equation*}
e^{8}\,\mathsf{w}_{\ell,n}\,3^{d+1}C_{fib}\sum_j|\Gamma^0_j\cap I|\,\mathsf{S}^{\circ,(n)}_\ell[\mathbf{B}] .
\end{equation*}

\emph{(a), summation.} The islands are disjoint, so
\begin{equation*}
\sum_{I\subset Y}\sum_j|\Gamma^0_j\cap I|\le\sum_j|\Gamma^0_j\cap Y| .
\end{equation*}
Summing the two bounds over $I\subset Y$ and over $\ell\ge n$ gives \eqref{9.1:eq-blind-domination-1} with
\begin{equation*}
\mathsf{y}_n^{\circ}=\sum_\ell\Bigl[\mathsf{w}_{\ell,n}\mathsf{S}^{\circ,(n)}_\ell[\mathbf{B}]
+\mathsf{S}^{\circ,(n)}_\ell[\mathbf{B}'^{res}]\Bigr] .
\end{equation*}

\emph{(a), the per-cell inequality.} Let $c\in\Gamma^0_n$ with $c\not\subset\mathbb{W}_{k'}$. For each $\ell\ge n$, $c$ lies in a unique level-$\ell$ cube $\hat c$, and $\hat c\not\subset\mathbb{W}_\ell$. This follows from Lemma~\ref{9.1:lem-island-geometry}(e): the level-$k'$ closure of the live part of $\mathbb{W}_\ell$ is contained in $\mathbb{W}_{k'}$. We apply Lemma~\ref{9.1:lem-blind-counts}(i) at $\hat c$, taking out the crossing as above. This yields the per-cell inequality of Lemma~\ref{9.1:lem-per-cell-domination}(a) with $K(d)$ replaced by $e^{8}C_{fib}\le K^{\circ}(d)$.

\emph{(b).} The activity sums of the second hypothesis are sums over the supports $X\ni\Delta$ of $e^{-\kappa_{\mathbf{B}}d(X)}$ times the size of the term on $X$. Splitting this exponential factor, we write them as
\begin{equation*}
\sum_{X\ni\Delta}e^{-(\kappa_{\mathbf{B}}-\sigma_*)d(X)}
\Bigl(e^{-\sigma_*d(X)}\cdot\bigl(\text{size of the term on }X\bigr)\Bigr).
\end{equation*}
As in Lemma~\ref{9.1:lem-per-cell-domination}(b), they are bounded through the count of supports through $\Delta$ at margin $\sigma_*-1$, with one unit kept for the crossing; the split-off factor $e^{-\sigma_*d(X)}$ supplies this margin and this unit.

For island-blind stores of level $\ell\le k'$, the count through the level-$k'$ cell $\Delta$ is Lemma~\ref{9.1:lem-blind-counts}(ii). Its hypothesis $\operatorname{dist}_{k'}(\Delta,\mathbb{W}_{k'})\ge1$ holds at every listed cell by Lemma~\ref{9.1:lem-island-geometry}(d), as follows.
\begin{itemize}
\item At (T) the cell $\Delta$ lies in the region $Y_0$ of the site expansion of the first hypothesis, and the distance is at least $3\check{R}_{k'}-4$.
\item At (L), for exterior anchors, it follows from the second relation of \eqref{9.1:eq-blind-resummation-a}.
\item At (P), let $j+1$ be the scale of the collar, so that $\ell\le j+1$ for the stores considered. The collar does not meet the islands of $\mathbb{I}_{j+1}$, and the level-$\ell$ sub-cubes of $\Delta$ do not meet $\mathbb{W}_\ell$.
\end{itemize}
The count is applied with $\lambda=\sigma_*-2$, and $\lambda\ge\mathsf{k}_d+2$ by \eqref{9.1:eq-blind-domination-a}. The remaining unit gives the crossing weight through Lemma~\ref{9.1:lem-blind-counts}(iv), exactly as in (a).
\end{proof}

\begin{remark}
Clause (a)'s value budget is the one place where a sum over island-blind stores is anchored inside an island. There the island pays with its own cells, since $N_\ell(I)\le\sum_j|\Gamma^0_j\cap I|$. This bound has the form of $|V(Y)|\le O(1)\sum_j|\Gamma^0_j\cap Y|$ in \cite[(1.69)]{B-CMP122b}, and of the charge per unit of $\sum_j|\Gamma^0_j\cap Y|$ in the value budget of the large-field resummation.

By Lemma~\ref{9.1:lem-island-geometry}(d), no activity sum is anchored inside an island. Hence neither the boundary size $|\partial I|$ of an island nor its cube count $N_\ell(I)$ is ever paid by a stored term in an activity sum. This matters because neither of these two numbers is bounded by any function of the scaffold radii $\check{R}_m$: an island may cluster many strips.
\end{remark}

\begin{lemma}[Interpolation with island-blind boundary terms]\label{9.1:lem-blind-interpolation}
Assume the following.
\begin{itemize}
\item[(i)] The hypotheses of Lemma~\ref{9.1:lem-resolved-interpolation} hold.
\item[(ii)] The expansion lemma at the $\mathbf{T}$-site holds with blind stores among its input
units, in the form stated in Proposition~\ref{9.1:prop-blind-site-expansions}, whose proof is
outstanding.
\end{itemize}
Then the conclusions of Lemma~\ref{9.1:lem-resolved-interpolation} hold when some leaf slots (slots
carrying stored boundary terms) are blind, that is, carry stores in the island-blind format, with filed
domains in the class $\mathbf{D}_\ell^{\mathbb{I}}$, that is used as input in
Proposition~\ref{9.1:prop-blind-site-expansions}, with the following readings.
\begin{itemize}
\item[(a)] In clause (c2)$^\star$ of Lemma~\ref{9.1:lem-resolved-interpolation}, the requirement that
the leaf of each of the two machines obey the leaf bound stated for that machine in clause
(c2)$^\star$ reads, for a blind leaf, as membership in the relative sup-ball
\begin{equation}\label{9.1:eq-blind-interpolation-1}
\Bigl\{\, x \;:\; \sup_X e^{\kappa_{\mathbf{B}} d^{\circ}(X)}\,|x(X)|
\le \text{the size of the leaf in one machine} \,\Bigr\}.
\end{equation}
\item[(b)] The activity per cell is at most $|\varepsilon|\,K_Y^{\circ}\,\mathsf{y}_n$, where
$\varepsilon=\lambda_z-\lambda$ is the variable of the domain $\mathbb{D}^{\star}$ of
Lemma~\ref{9.1:lem-resolved-interpolation}.
\item[(c)] For every resummation or output domain $Y$ as in
Lemma~\ref{9.1:lem-blind-domination}\textup{(a)}, the values of the members supported in $Y$ are at
most
\begin{equation*}
|\varepsilon|\,\bigl(K(d)+K_{\circ}+K^{\circ}\bigr)\,\mathsf{y}_n\sum_j|\Gamma^0_j\cap Y|.
\end{equation*}
\item[(d)] Clause (c3)$^z$ of Lemma~\ref{9.1:lem-resolved-interpolation} holds unchanged, read for
second-class and arrested inner zones in the sequence form of the zone bound for such zones.
\end{itemize}
The constant $C_{\mathbf{B}}^v$ of clause (c2)$^\star$ of
Lemma~\ref{9.1:lem-resolved-interpolation} takes the value
\begin{equation}\label{9.1:eq-blind-interpolation-2}
C_{\mathbf{B}}^v = 2\,\bigl(K(d)+K_{\circ}+K^{\circ}\bigr)\,\frac{E_1^z}{E_0},
\end{equation}
with $E_1^z$ and $E_0$ as in Lemma~\ref{9.1:lem-resolved-interpolation}.
\end{lemma}

\begin{proof}
By hypothesis (ii), the expansion at the $\mathbf{T}$-site is available with blind stores among its
input units; by hypothesis (i), the argument of Lemma~\ref{9.1:lem-resolved-interpolation} applies, and
we indicate the places where a blind leaf slot enters it.

(a) For a blind leaf, the leaf bound is the relative sup-ball
\eqref{9.1:eq-blind-interpolation-1}. This set is convex: if $x_u^{\mathsf{A}}$ and $x_u^{\mathsf{B}}$
both lie in it and $\lambda\in[0,1]$, then for every $X$
\begin{equation*}
\begin{split}
&e^{\kappa_{\mathbf{B}} d^{\circ}(X)}\,
\bigl|(1-\lambda)x_u^{\mathsf{A}}(X)+\lambda x_u^{\mathsf{B}}(X)\bigr|\\
&\qquad\le (1-\lambda)\,e^{\kappa_{\mathbf{B}} d^{\circ}(X)}|x_u^{\mathsf{A}}(X)|
+\lambda\, e^{\kappa_{\mathbf{B}} d^{\circ}(X)}|x_u^{\mathsf{B}}(X)|,
\end{split}
\end{equation*}
and the right-hand side is at most the size of the leaf in one machine. Hence the convex combination
$(1-\lambda)x_u^{\mathsf{A}}+\lambda x_u^{\mathsf{B}}$ of the entries of the slot-value arrays of the
two runs stays in the ball, which is what clause (c2)$^\star$ requires of the leaf along the
interpolation.

(b) The bound $|\varepsilon|K_Y^{\circ}\mathsf{y}_n$ on the activity per cell follows from part~(b) of
the per-cell domination for island-blind stores, Lemma~\ref{9.1:lem-blind-domination}.

(c) The bound of item (c) on the values of the members supported in $Y$ follows from part~(a) of
Lemma~\ref{9.1:lem-blind-domination}, applied to that domain $Y$.

(d) Clause (c3)$^z$ reads the place where zone $n$ is met, not a length, so it is unchanged; for
second-class and arrested inner zones it is read with the sequence form of the zone bound for such
zones, as stated in item (d).

Finally, with the value constant $K(d)+K_{\circ}+K^{\circ}$ of item (c), the constant
$C_{\mathbf{B}}^v$ of clause (c2)$^\star$ is \eqref{9.1:eq-blind-interpolation-2}.
\end{proof}

\begin{theorem}[Zone-by-zone bound for fluctuation outputs reading island-blind units]
\label{9.1:thm-blind-fluctuation-bound}
Work in the setting of Theorem~\ref{9.1:thm-fluctuation-boundary-bound}, the zone-by-zone
two-run bound for boundary terms of the fluctuation step. As hypothesis, assume the part of
Proposition~\ref{9.1:prop-blind-site-expansions}, site expansions with island-blind inputs,
whose proof is outstanding, that concerns the site of the fluctuation step $\mathbf{T}$
(below, the T-site).

Consider the $\mathbf{T}$-step (fluctuation step) $i\to i+1$ of an aligned label and an output
$\mathbf{B}^{(i+1)}(X)$ whose expansion reads an island-blind input unit. Its domain is filed,
in the format of that proposition, as
$X:=\operatorname{fil}_{i+1}(\bar{X})\in\mathbf{D}_{i+1}^{\mathbb{I}}$, where $\bar{X}$ is the
domain of the output before filing and $\operatorname{fil}_{i+1}$ is the filing map of level
$i+1$ of that proposition.
Then the following hold.
\begin{enumerate}
\item[(a)] The first two assertions of
Theorem~\ref{9.1:thm-fluctuation-boundary-bound} hold for this output as stated there.
\item[(b)] For each zone $n$ as in Theorem~\ref{9.1:thm-fluctuation-boundary-bound},
\begin{equation}\label{9.1:eq-blind-fluctuation-bound-1}
\mathsf{S}^{\circ,(n)}_{i+1}[\mathbf{B}]
\le \varkappa^{\mathbf{B}}_i\,[\![\delta]\!]^{\natural+,(i)}_n+\sigma^{T,(n)} .
\end{equation}
Here $\mathsf{S}^{\circ,(n)}_{i+1}[\mathbf{B}]$ is the zone-$n$ seminorm
$\mathsf{S}^{+,(n)}_{i+1}[\mathbf{B}]$ of Theorem~\ref{9.1:thm-fluctuation-boundary-bound} of
the two-run differences of the outputs of level $i+1$, in which the supremum over domains is
taken with the weight $e^{\kappa_{\mathbf{B}}d^{\circ}_{i+1}(X)}$ of the island-blind length
$d^{\circ}_{i+1}$ in place of the weight of the plain length.
The numbers $\varkappa^{\mathbf{B}}_i$ and $\sigma^{T,(n)}$ are the coefficients of
Theorem~\ref{9.1:thm-fluctuation-boundary-bound}; $\sigma^{T,(n)}$ is unchanged, and
$\varkappa^{\mathbf{B}}_i$ is the coefficient $\varkappa^{\mathbf{B}}_i$ of that theorem,
proportional to $C_L'''\,\theta_i$, in which the constant $C_L'''$ of that theorem is evaluated
with $K_Y^{\circ}$ in place of $K_Y$.
\item[(c)] If, for some zone $n\le i-1$, the zone-$n$ member
$\delta\mathbf{B}^{(i+1),(n)}(X)$ of the decomposition in the first assertion of
Theorem~\ref{9.1:thm-fluctuation-boundary-bound} for this output is not identically zero,
then
\begin{equation*}
d^{\circ}_{i+1}(X)\ge\mathsf{d}_i-8 .
\end{equation*}
\end{enumerate}
\end{theorem}

\begin{proof}
Assertion (a) is taken over verbatim from Theorem~\ref{9.1:thm-fluctuation-boundary-bound};
nothing further is to be proved for it.

For (b) we go through the two parts of the proof of
Theorem~\ref{9.1:thm-fluctuation-boundary-bound}: the part that bounds the activity of the
marked potential, and the part that bounds the remaining term. In each part the length of the
output domain is replaced by the island-blind length $d^{\circ}_{i+1}(X)$.

\emph{The marked potential.} The activity per cell of the marked potential $V^{m,z}$ is read at
the cells of the T-site. By item~(d) of the relative geometry of sites and live islands,
these cells lie off $\mathbb{W}_i$. Hence part (b) of Lemma~\ref{9.1:lem-blind-domination},
the per-cell domination for island-blind stores, applies at each of them. It bounds the
activity per cell by
\begin{equation*}
\theta_i\,C_L'''\,[\![\delta]\!]^{\natural+,(i)}_n ,
\end{equation*}
with the constant $K_Y^{\circ}$ in place of $K_Y$.

The corollary on the interpolated expansion at the T-site carries an output factor
$e^{-\kappa\,d_{i+1}(X)}$, where $d_{i+1}(X)$ is the plain length at level $i+1$ of the output
domain $X$ used there.
This factor is the
output exponential of the expansion lemma at the T-site. For outputs that read
an island-blind unit, the T-site part of Proposition~\ref{9.1:prop-blind-site-expansions},
whose proof is outstanding, supplies this exponential in the form
\begin{equation*}
e^{-\kappa_{out}\,d^{\circ}_{i+1}(X)} .
\end{equation*}

\emph{The remaining term.} The radius fixed at the T-site does not depend on any length of the
output domain, and neither does the crossing length $\mathsf{d}_n$. The term written $G$ in
this part of the proof of Theorem~\ref{9.1:thm-fluctuation-boundary-bound} is bounded by the
T-site part of Proposition~\ref{9.1:prop-blind-site-expansions}, whose proof is outstanding.
That part asserts, for outputs reading an island-blind unit, item~(vi) of the bound lemma for
the expansion at the T-site, in its clause for island-blind inputs. This clause
gives
\begin{equation*}
|G|\le E^{\sharp}_T\,e^{-\kappa_{out}\,d^{\circ}_{i+1}(X)} .
\end{equation*}

\emph{Conclusion of} (b). The quantity $\mathsf{S}^{\circ,(n)}_{i+1}[\mathbf{B}]$ is a
supremum over domains with the weight $e^{\kappa_{\mathbf{B}}d^{\circ}_{i+1}(X)}$. This weight
satisfies
\begin{equation*}
e^{\kappa_{\mathbf{B}}d^{\circ}_{i+1}(X)}\le e^{\kappa_{out}\,d^{\circ}_{i+1}(X)} ;
\end{equation*}
since $d^{\circ}_{i+1}(X)\ge0$, this is the inequality $\kappa_{\mathbf{B}}\le\kappa_{out}$
between the two rates. So after multiplication by the weight, the bounds of the two parts hold
with the same coefficients as in the proof of
Theorem~\ref{9.1:thm-fluctuation-boundary-bound}. This gives
\eqref{9.1:eq-blind-fluctuation-bound-1}, with $\varkappa^{\mathbf{B}}_i$ the coefficient of
Theorem~\ref{9.1:thm-fluctuation-boundary-bound} and $C_L'''$ evaluated at $K_Y^{\circ}$.

\emph{Assertion} (c). This is item~(iv) of Lemma~\ref{9.1:lem-blind-counts}, the domain
counts in the island-blind length, applied at level $i+1$.
\end{proof}

\begin{remark}
The proposition that eliminates the stored boundary terms $\mathbf{B}$, and the two lemmas
and the corollary of this section that are stated together with it and read the coefficients
listed below, remain valid without change when island-blind outputs are present.

These statements use the following quantities only as numbers:
$\varkappa^{\mathbf{B}}$, $\varkappa^{\mathbf{B}'}$, $\sigma^{T}$, $c_1^{\flat}$,
the crossing weights $\mathsf{w}_{\ell,n}$, $w^{\mathbf{B}}$, $q_a$, the constants $E_0$,
$E_1$, $B_0$ of those statements, and the loads. The island-blind pieces add terms of the same
form to $\mathsf{y}_n$ and $\mathsf{y}'_n$, with the same coefficients.

The following are used exactly as before: the bound
$\sum_{\ell}\mathsf{w}_{\ell,n}\le3$; the numbers $\bar{\varkappa}'$, $\bar{c}_1$,
$\mathfrak{k}$, $\mathfrak{s}$; the four smallness conditions on the coupling imposed for the
boundary terms; the condition on the strip constants; and the constant $E_1^z$ of
Theorem~\ref{9.1:thm-fluctuation-boundary-bound}. The one change is that
$K_Y$ is replaced by $K_Y^{\circ}$ inside $C_L'''$ and inside the third of those four
smallness conditions.

No scaffold radius $\check{R}$, neither of the two further size parameters of the geometric
preliminaries of this chapter, no age, no number of steps and no torus
size enters these statements. In the geometric preliminaries of this chapter, $\check{R}$
appears only inside distances used in the form $\ge1$ or $\ge23$, and in one further
condition of those preliminaries involving $\check{R}$.
\end{remark}

We use four elementary facts. They hold for $\mathsf{w}\le s\le s'\le n$, $m\ge 0$ and $\nu\ge 1$. In the last line the indices are the level $i=n-s'$ and the zone $j=n-s$, and the run lies on the scaffold $\Theta_X$.
\begin{equation}\label{9.1:eq-load-envelope-a}
\begin{gathered}
e_{s'}\le\rho^{-(s'-s)}e_s,\qquad e_s\le e_{s'},\qquad \mathsf{n}_s\le\mathsf{n}_{s'};\\
\mathsf{l}_{s+m}\le\hat{\mathsf{l}}_m\,\mathsf{l}_s,\qquad
\hat{\mathsf{l}}_m:=\bigl(1+\log(1+c_0m)\bigr)^{c_1};\\
\varsigma_\nu\le C_\varsigma\,\nu^3q^{\nu-1},\qquad C_\varsigma:=18;\\
g_{(s)}^2\le 2\theta_s^{-1},\qquad \log x_{(s)}\le 2\log\theta_s,\qquad
w_{j,i}\le C^\theta_w\,\frac{\theta_{s'}}{\theta_s},\qquad C^\theta_w:=8.
\end{gathered}
\end{equation}

We justify them line by line.

For the first line, recall $e_r=\rho^{n-r}+\Pi_n$ with $0<\rho<1$. Then $\rho^{n-s'}=\rho^{-(s'-s)}\rho^{n-s}$, and $\Pi_n\le\rho^{-(s'-s)}\Pi_n$; adding the two gives $e_{s'}\le\rho^{-(s'-s)}e_s$. Since $n-s'\le n-s$, we have $\rho^{n-s}\le\rho^{n-s'}$, hence $e_s\le e_{s'}$. The inequality $\mathsf{n}_s\le\mathsf{n}_{s'}$ holds because, by its definition, $\mathsf{n}_r$ is non-decreasing in the depth $r$.

For the second line, recall $\mathsf{l}_s=(\log\theta_s)^{c_1}$ with $\theta_s=X+c_0s$ and $\log\theta_s\ge 1$. Since $\theta_s\ge 1$,
\[
\theta_{s+m}=\theta_s+c_0m\le\theta_s(1+c_0m).
\]
Hence, using $\log\theta_s\ge1$,
\[
\log\theta_{s+m}\le\log\theta_s+\log(1+c_0m)\le\log\theta_s\,\bigl(1+\log(1+c_0m)\bigr).
\]
Raising this to the power $c_1$ gives the second line.

The third line follows from the definition of the load weights $\varsigma_\nu$: the exponent $\hat\alpha$ in that definition satisfies $L^{-\hat\alpha}<q$, also $L^{-1/2}<q$, and the squares of the remaining constants of that definition are at most $3$.

For the fourth line, on $\Theta_X$ we have $x_{(s)}=\theta_s-\zeta_{(s)}$, and the tube condition gives $|\zeta_{(s)}|\le\varrho_s\le\theta_s/2$. Therefore $\theta_s/2\le|x_{(s)}|\le 3\theta_s/2$. The lower bound gives $|g_{(s)}^2|=|x_{(s)}|^{-1}\le 2\theta_s^{-1}$. Since $\log\theta_s\ge1$ forces $\theta_s\ge 3/2$, the upper bound gives $|x_{(s)}|\le\theta_s^2$, that is, $\log|x_{(s)}|\le2\log\theta_s$. For real $x_{(s)}$ these are the first two inequalities of the fourth line; in general those inequalities are read with moduli. The weights $w_{j,i}$, with which the entries of level $i$ enter the loads of zone $j$, obey $w_{j,i}\le C^\theta_w\theta_{s'}/\theta_s$ with $C^\theta_w=8$ by the weight bound stated with the correlation theorem.

\begin{lemma}[The envelope of the loads]\label{9.1:lem-load-envelope}
Let $g\le\gamma_U$, and assume the following.
\begin{itemize}
\item[(a)] The five conclusions of the response proposition for the paired runs hold at the constants of the envelope bound without a boundary-term hypothesis, together with the remark on the constituent-seeing components of that response system read at a constant $C_X^\star\ge1$. That is, the response system is run at a constant $C_E^\natural$ and returns its bounds at a constant at most $C_E^\natural/(1-2\epsilon)$, where $0\le\epsilon\le\frac1{16}$; this is at most
\[
C_E:=\tfrac87\,C_E^\natural .
\]
The weights $\Lambda_r$ are unchanged under the ceiling \eqref{9.1:eq-envelope-corollary-a}, and the coefficients $\varkappa$, $\varkappa_A$, $\varkappa''$ of the fluctuation-step rows are unchanged.
\item[(b)] Part (a) of the two-run single-step bound holds, and the bound on the crossing weights $w^{\mathbf{B}}$ holds.
\item[(c)] The normalisation \eqref{9.1:eq-centre-normalisation-d} holds, read at the radius $E_1^z$ of \eqref{9.1:eq-apriori-loads-a}.
\item[(d)] The following hold under the four coupling ceilings on $g$ that they carry, the last of which is \eqref{9.1:eq-apriori-loads-a}, and under the two further conditions on the constants under which they are stated:
  \begin{itemize}
  \item the zone-by-zone two-run bound for boundary terms of the fluctuation step (Theorem~\ref{9.1:thm-fluctuation-boundary-bound});
  \item the two-run bound for boundary terms of the large-field step (Proposition~\ref{9.1:prop-large-field-boundary});
  \item the elimination of the boundary-term loads (Proposition~\ref{9.1:prop-load-elimination});
  \item the a priori bound on the loads (Lemma~\ref{9.1:lem-apriori-loads}).
  \end{itemize}
\end{itemize}
Consider the local pair of Notation~\ref{9.1:not-centre-normalisation} at any $z\in P_{UV}$, or the frozen pair at any $\hat{w}$ in the disc $D(0,r_J)$ of radius $r_J$ about the origin. Let $j\le k$ be a zone present in an aligned label (an aligned zone in the sense of Proposition~\ref{9.1:prop-load-elimination}), and put $s:=n-j$. Then
\begin{equation}\label{9.1:eq-load-envelope-1}
[\![\delta]\!]^{\natural+}_j\le C^z_{env}\,\mathsf{l}_s^2\,e_s ,
\end{equation}
with $C^z_{env}$ as in \eqref{9.1:eq-load-envelope-b}, and hence
\begin{equation}\label{9.1:eq-load-envelope-2}
\Sigma^{q_a}_k\le C^{TC}_\Sigma\,C^z_{env}\,\mathsf{l}_{\mathsf{w}}^2\,e_{\mathsf{w}},\qquad
C^{TC}_\Sigma:=\sum_{m\ge0}(m+1)\,\hat{\mathsf{l}}_m^2\,q_a^{(m-1)_+}\rho^{-m}<\infty .
\end{equation}
The constants $C^{(E)}$, $C^{(R)}$, $C^{(\zeta)}$, $C^{(c)}$ of \eqref{9.1:eq-load-envelope-b} and $C^{TC}_\Sigma$ depend only on
\[
(L,\mathfrak{p},q,q_a,\rho;\,C_E,\mathsf{A},E_0,B_0,C_\zeta,C_\vartheta,C_t,C_T,
C^\natural_{cl},c_0).
\]
The constant $C^z_{env}$ reads in addition the constants $\bar a$, $\epsilon_0$, $C^{\mathbf{B}}_T$ of Proposition~\ref{9.1:prop-load-elimination}, the constant $\epsilon_1$ of Lemma~\ref{9.1:lem-apriori-loads} and the constant $C^{(\mathfrak{s})}$ of Remark~\ref{9.1:rem-own-source}. None of these constants depends on $n$, $k$, $j$, $s$, $N_w$, $m_0$, the torus or the position.
\end{lemma}

\begin{proof}
The proof has four steps.

\emph{Step (1): the unknowns at depth $s'$.} Fix a level $i=n-s'\le j$, so that $s'\ge s$.

The response proposition for the paired runs, assumed in hypothesis (a), applies with starting depth zero. Indeed, the local pair is matched at depth zero for every frozen value, by the pinning statement for the runs on $\Theta_X$, and the frozen pair likewise. Apply it, together with the remark on its constituent-seeing components read at the constant $C_X^\star$. For the components $\mathsf{K}^P_{s'}$, $\mathsf{K}^A_{s'}$, $\mathsf{K}^R_{s'}$ of the rows $P$, $A$, $R$ this gives
\[
C_X^\star\,\mathsf{K}^P_{s'}\le E^P_{s'},\quad C_X^\star\,\mathsf{K}^A_{s'}\le E^A_{s'},\quad
C_X^\star\,\mathsf{K}^R_{s'}\le E^R_{s'},\qquad C_X^\star\ge1,
\]
where
\[
E^P=C_Ee_{s'},\qquad E^A=4C_Ee_{s'},\qquad E^R=4C_E\hat{\mathfrak{r}}_{s'}e_{s'} .
\]
Since $C_X^\star\ge1$, each component is at most the corresponding right-hand side.

In the reading units of the response rows, the $A$-component is $\mathsf{l}^{-1}(P^\flat;D)/E_0$, the $P$-component contains $C_t\theta_{s'}Q$, and the $R$-component is $C_t\theta_{s'}(R'+R'')$. The three bounds therefore read
\[
\mathsf{S}_i^+[P^\flat],\ \mathsf{S}_i^+[D]\le 4C_EE_0\mathsf{l}_{s'}e_{s'},\qquad
\mathsf{S}_i^+[Q]\le\frac{C_Ee_{s'}}{C_t\theta_{s'}},\qquad
\mathsf{S}_i^+[R]\le\frac{4C_E\hat{\mathfrak{r}}_{s'}e_{s'}}{C_t\theta_{s'}} .
\]
Now $E=P^\flat+Q+D$, and $D=0$ for the frozen pair. By the triangle inequality for $\mathsf{S}_i^+$,
\begin{equation}\label{9.1:eq-load-envelope-3}
\mathsf{S}_i^+[E]\le C_E\bigl(8E_0\mathsf{l}_{s'}+(C_t\theta_{s'})^{-1}\bigr)e_{s'} .
\end{equation}

For the couplings, the coupling bound in part (b) of the transplanted coupling hypothesis of that response proposition, combined with (a), gives
\[
\mathsf{b}_i\le\mathsf{A}\,g^P_{s'}\le\mathsf{A}C_Ee_{s'} .
\]
Here $g^P_{s'}$ is the coupling entry of the $P$-row at depth $s'$. The same bound also gives
\[
d_{s'}\le\mathsf{A}C_E\mathsf{n}_{s'-1}e_{s'-1}\le\mathsf{A}C_E\mathsf{n}_{s'}e_{s'},
\]
where the second inequality uses the first line of \eqref{9.1:eq-load-envelope-a}.

For the $\zeta$-fields of the local pair, the running-shift comparison lemma for the local pair gives
\[
\sup_x\bigl|\zeta^{\mathsf{B}}_{(s')}(x)-\zeta^{\mathsf{A}}_{(s')}(x)\bigr|
\le d_{s'}+\Bigl[\tfrac12\mathsf{A}C_E(1-\rho)^{-1}+r_J\Bigr]e_{s'} .
\]
This holds where the profile $\varphi$ equals $1$, which is the case down to depth $s'\ge\mathsf{w}$. In a transition layer the field is the convex combination of its values at two consecutive depths. With the first line of \eqref{9.1:eq-load-envelope-a} this gives, for every $s'\in[\mathsf{w},n]$,
\begin{equation}\label{9.1:eq-load-envelope-4}
\mathsf{z}_{s'}\le C_\zeta(1+s')e_{s'},
\end{equation}
where
\[
C_\zeta:=\rho^{-1}\Bigl[\mathsf{A}C_E\bigl(2+(1-\rho)^{-1}\bigr)+\tfrac12\mathsf{A}C_E(1-\rho)^{-1}+r_J\Bigr].
\]

The stored boundary terms need no bound in this step: their two-cutoff difference enters the loads only through the entries $\mathsf{y}_j+\mathsf{y}'_j$ of \eqref{9.1:eq-centre-normalisation-d}, which are eliminated after Step (3).

\emph{Step (2): the $u$-load of zone $j$.} Throughout, the level $i\le j$ corresponds to the depth $s'=s+m$, with $m=j-i\ge0$.

Line (E). By the third line of \eqref{9.1:eq-load-envelope-a} with $\nu=1+m$, $\varsigma_{1+m}\le C_\varsigma(m+1)^3q^m$. In \eqref{9.1:eq-load-envelope-3} we use three bounds:
\begin{itemize}
\item $\mathsf{l}_{s+m}\le\hat{\mathsf{l}}_m\mathsf{l}_s$, from the second line of \eqref{9.1:eq-load-envelope-a};
\item $(C_t\theta_{s'})^{-1}\le1$;
\item $e_{s+m}\le\rho^{-m}e_s$, from the first line of \eqref{9.1:eq-load-envelope-a}.
\end{itemize}
This gives
\[
\sum_{i\le j}\varsigma_{1+j-i}\mathsf{S}_i^+[E]
\le\sum_{m\ge0}C_\varsigma(m+1)^3q^m\,C_E\bigl(8E_0\hat{\mathsf{l}}_m\mathsf{l}_s+1\bigr)\rho^{-m}e_s .
\]
Using $\hat{\mathsf{l}}_m\ge1$ and $\mathsf{l}_s\ge1$, this is at most $C^{(E)}\mathsf{l}_se_s$, with $C^{(E)}$ as in \eqref{9.1:eq-load-envelope-b}. The constant is finite: $q/\rho\le\rho<1$, and $\hat{\mathsf{l}}_m$ grows only like a power of $\log m$.

Line (R). By the fourth line of \eqref{9.1:eq-load-envelope-a} and the bound on $\mathsf{S}_i^+[R]$ from Step (1),
\begin{align*}
\sum_{i\le j}w_{j,i}\mathsf{S}_i^+[R]
&\le\sum_{s'\ge s}C^\theta_w\frac{\theta_{s'}}{\theta_s}\cdot
\frac{4C_E\hat{\mathfrak{r}}_{s'}e_{s'}}{C_t\theta_{s'}}
=\frac{4C^\theta_wC_E}{C_t\theta_s}\sum_{s'\ge s}\hat{\mathfrak{r}}_{s'}e_{s'}\\
&\le\frac{4C^\theta_wC_E}{C_t\theta_s}\bigl(\hat{\mathfrak{r}}_se_s+2\Pi_n\bigr)
\le\frac{4C^\theta_wC_E}{C_t\theta_s}\Bigl(\frac1{16}+2\Bigr)e_s
\le C^{(R)}e_s .
\end{align*}
The second line uses four facts:
\begin{itemize}
\item the summation bound on the $R$-row weights, $\sum_{s'>s}\hat{\mathfrak{r}}_{s'}e_{s'}\le2\Pi_n$;
\item $\Pi_n\le e_s$;
\item the smallness condition $4\hat{\mathfrak{r}}_s\le\frac14$ on the $R$-row weights;
\item $C_t\theta_s\ge1$.
\end{itemize}
The flat remnant sum thus costs no factor beyond a constant.

The boundary-term loads $\mathsf{y}_j+\mathsf{y}'_j$ are not bounded in this step; they are eliminated by Proposition~\ref{9.1:prop-load-elimination} after Step (3).

Line ($\zeta$). For the $\zeta$-load, the weight $\mathfrak{g}_j$ is bounded through the fourth line of \eqref{9.1:eq-load-envelope-a}, and the entries through Step (1): $\mathsf{z}^\sharp_j$ through \eqref{9.1:eq-load-envelope-4}, and $\mathsf{b}_i\le\mathsf{A}C_Ee_{s+m}$, together with the first line of \eqref{9.1:eq-load-envelope-a}. This gives
\begin{align*}
&\mathfrak{g}_j\Bigl[\mathsf{z}^\sharp_j+\sum_{i\le j}L^{-(j-i)/2}\mathsf{b}_i\Bigr]\\
&\quad\le\frac{2(\log\theta_s)^{2q_\flat}}{\theta_s}\Bigl[\sum_{m\ge0}L^{-m}C_\zeta(3+2s+2m)\rho^{-m-1}\\
&\qquad\qquad+\sum_{m\ge0}L^{-m/2}\mathsf{A}C_E\rho^{-m}\Bigr]e_s .
\end{align*}

We now simplify the right-hand side using four facts:
\begin{itemize}
\item $3+2s+2m\le(2s+3)(1+m)$ and $2s+3\ge1$;
\item $(\log\theta_s)^{2q_\flat}\le\mathsf{l}_s^2$, because $c_1\ge q_\flat$ and $\log\theta_s\ge1$;
\item $L^{-m}\rho^{-m-1}=\rho^{-1}(L^{-1}/\rho)^m$ and $L^{-m/2}\rho^{-m}=(L^{-1/2}/\rho)^m$;
\item $(2s+3)/\theta_s\le3+2/c_0$, because $\theta_s=X+c_0s\ge1+c_0s$ and $2s+3\le(3+2/c_0)(1+c_0s)$.
\end{itemize}
With these, the right-hand side is at most
\begin{align*}
&2\mathsf{l}_s^2\,\frac{2s+3}{\theta_s}\Bigl[C_\zeta\rho^{-1}\sum_{m\ge0}(1+m)\Bigl(\frac{L^{-1}}{\rho}\Bigr)^m\\
&\qquad\qquad+\mathsf{A}C_E\sum_{m\ge0}\Bigl(\frac{L^{-1/2}}{\rho}\Bigr)^m\Bigr]e_s
\le C^{(\zeta)}\mathsf{l}_s^2e_s .
\end{align*}
Both ratios $L^{-1}/\rho\le L^{-1/2}/\rho$ are at most $\rho<1$, since $L^{-1/2}<q\le\rho^2$. Hence $C^{(\zeta)}<\infty$.

\emph{Step (3): the $c$-load.} The classical index rates obey
\[
\hat{\vartheta}^\natural_i\le C_\vartheta\vartheta_{dl}^{\,i}=C_\vartheta\vartheta_{dl}^{\,n-s'}\le C_\vartheta\rho^{n-s'},
\]
since $\vartheta_{dl}\le\rho$. We bound the four entries in turn.

Line ($E_0$). By the third line of \eqref{9.1:eq-load-envelope-a},
\[
E_0\sum_{i\le j}\varsigma_{1+j-i}\hat{\vartheta}^\natural_i
\le E_0C_\vartheta C_\varsigma\rho^{n-s}\sum_{m\ge0}(m+1)^3\Bigl(\frac q\rho\Bigr)^m
=C^aE_0\rho^{n-s}\le C^aE_0e_s ,
\]
with $C^a$ as in \eqref{9.1:eq-load-envelope-b}.

Line ($B_0$). The same computation with $4\varsigma B_0$ in place of $\varsigma E_0$ gives the bound $4C^aB_0e_s$.

Line ($g^{\kappa_0}$). Here $g_i$ is the coupling $g_{(s')}$ at depth $s'$, and $\kappa_0$ is the exponent with which it enters this entry. By the fourth line of \eqref{9.1:eq-load-envelope-a}, $w_{j,i}\le C^\theta_w\theta_{s'}/\theta_s$ and $g_i^{\kappa_0}\le 2^{\kappa_0/2}\theta_{s'}^{-\kappa_0/2}$. Therefore
\[
\sum_{i\le j}w_{j,i}g_i^{\kappa_0}\hat{\vartheta}^\natural_i
\le\frac{C^\theta_wC_\vartheta2^{\kappa_0/2}}{\theta_s}\sum_{s'\ge s}\theta_{s'}^{1-\kappa_0/2}\rho^{n-s'}
=\frac{C^\theta_wC_\vartheta2^{\kappa_0/2}}{C_t\theta_s}\sum_{s'\ge s}\mathfrak{r}_{s'}\rho^{n-s'} ,
\]
where the equality expresses the sum through the $R$-row weights $\mathfrak{r}_{s'}$. By the condition on the $R$-row weights,
\[
\mathfrak{r}_{s'}\le\frac{\hat{\mathfrak{r}}_{s'}}{16}\le\frac{C_T}{16}\rho^{-(n-s')/2}\hat{\mathfrak{r}}_n .
\]
Hence the last sum is at most
\[
\frac{C_T\hat{\mathfrak{r}}_n}{16}\sum_{s'}\rho^{(n-s')/2}\le\frac{\Pi_n}{16}\le\frac{e_s}{16},
\]
and the line is at most $C^b(C_t\theta_s)^{-1}e_s$. This is the one entry that carries no index smallness at old levels. It is paid for by the decay of $\mathfrak{r}$ itself, through the floor $\Pi_n$, by the same mechanism as the summation bound used in line (R).

Line (last). By the fourth line of \eqref{9.1:eq-load-envelope-a}, $\log x_j\le2\log\theta_s$. Also $(\log\theta_s)^{c_2}\le\mathsf{l}_s$, because $c_2\le c_1$ and $\log\theta_s\ge1$. Therefore
\[
E_0\min\bigl\{1,C^\natural_{cl}(\log x_j)^{c_2}\vartheta_2^{j/4}\bigr\}
\le E_0C^\natural_{cl}2^{c_2}\mathsf{l}_s\bigl(\vartheta_2^{1/4}\bigr)^{n-s}
\le2^{c_2}C^\natural_{cl}E_0\mathsf{l}_se_s ,
\]
using $\vartheta_2^{1/4}=L^{-\theta_2/4}\le\vartheta_{dl}\le\rho$ and $\rho^{n-s}\le e_s$.

\emph{Assembly of Steps (2) and (3).} Add the lines (E), (R), ($\zeta$) of Step (2) and the four lines of Step (3), using $\mathsf{l}_s\ge1$ and $(C_t\theta_s)^{-1}\le1$. This gives the bound on the five-unknown and classical loads:
\begin{equation}\label{9.1:eq-load-envelope-5}
[\![\delta]\!]^{(5)}_j+[\![\delta]\!]^c_j\le\bigl(C^{(E)}+C^{(R)}+C^{(\zeta)}\bigr)\mathsf{l}_s^2e_s
+C^{(c)}\mathsf{l}_se_s .
\end{equation}
The constants of Steps (2), (3) and (4), and the constant of the lemma, are
\begin{equation}\label{9.1:eq-load-envelope-b}
\begin{aligned}
C^{(E)}&:=C_EC_\varsigma(8E_0+1)\sum_{m\ge0}(m+1)^3\hat{\mathsf{l}}_m\Bigl(\frac q\rho\Bigr)^m,\qquad
C^{(R)}:=4C^\theta_wC_E\Bigl(\frac1{16}+2\Bigr),\\
C^{(\zeta)}&:=2\Bigl(3+\frac2{c_0}\Bigr)\Bigl[C_\zeta\rho^{-1}\sum_{m\ge0}(1+m)\Bigl(\frac{L^{-1}}{\rho}\Bigr)^m
+\mathsf{A}C_E\sum_{m\ge0}\Bigl(\frac{L^{-1/2}}{\rho}\Bigr)^m\Bigr],\\
C^a&:=C_\vartheta C_\varsigma\sum_{m\ge0}(m+1)^3\Bigl(\frac q\rho\Bigr)^m,\qquad
C^b:=\frac{C^\theta_wC_\vartheta2^{\kappa_0/2}}{16},\\
C^{(c)}&:=(5C^a+1)E_0+4C^aB_0+C^b+2^{c_2}C^\natural_{cl}E_0,\\
C^{TC}_\Sigma&=\sum_{m\ge0}(m+1)\,\hat{\mathsf{l}}_m^2\,q_a^{(m-1)_+}\rho^{-m},\\
C^z_{env}&:=(1-\epsilon_1)^{-1}\Bigl[(1+\epsilon_0)\bigl(C^{(E)}+C^{(R)}+C^{(\zeta)}+C^{(c)}
+C^{(\mathfrak{s})}\bigr)\\
&\qquad\qquad\qquad+3e^{\bar a}C^{\mathbf{B}}_TC^{(c)}\Bigr].
\end{aligned}
\end{equation}
Here $\bar a$, $\epsilon_0=e^{\bar a}\bar a$ and $C^{\mathbf{B}}_T$ are the constants of Proposition~\ref{9.1:prop-load-elimination}, and $\epsilon_1$ is that of Lemma~\ref{9.1:lem-apriori-loads}.

The two-cutoff difference of the stored boundary terms enters the loads only through the entries $\mathsf{y}_j+\mathsf{y}'_j$ of \eqref{9.1:eq-centre-normalisation-d}. Proposition~\ref{9.1:prop-load-elimination}, parts (ii)--(iv) of \eqref{9.1:eq-load-elimination-a}, bounds these entries by the five-unknown and classical entries of the same and older zones; no sixth unknown occurs. In the envelope norm $\omega_j:=\mathsf{l}_s^2e_s$, they are at most $\epsilon_{\mathbf{B}}$ times the five-unknown and classical loads, plus the own sources $\mathfrak{s}_j$, uniformly in the number of steps.

Now apply Proposition~\ref{9.1:prop-load-elimination}, part (iv) of \eqref{9.1:eq-load-elimination-a}, in the envelope norm $\omega_j=\mathsf{l}_s^2e_s$. Its inputs are:
\begin{itemize}
\item the five-unknown and classical bound \eqref{9.1:eq-load-envelope-5};
\item the own-source bound $\mathfrak{s}_j\le C^{(\mathfrak{s})}e_s$ of Remark~\ref{9.1:rem-own-source}, with $C^{(\mathfrak{s})}$ as in \eqref{9.1:eq-own-source-a};
\item $\mathsf{l}_s\ge1$ and $e_s\le\mathsf{l}_s^2e_s$, which express all entries in that norm.
\end{itemize}
This proves \eqref{9.1:eq-load-envelope-1} with the constant $C^z_{env}$ of \eqref{9.1:eq-load-envelope-b}.

\emph{Step (4): the weighted sum.} Put $m:=k-j=s-\mathsf{w}\in[0,k]$, so that the zone $k-m$ has depth $\mathsf{w}+m$. By \eqref{9.1:eq-load-envelope-1},
\begin{align*}
\Sigma^{q_a}_k&=\sum_{m}(m+1)q_a^{(m-1)_+}[\![\delta]\!]^{\natural+}_{k-m}
\le C^z_{env}\sum_m(m+1)q_a^{(m-1)_+}\mathsf{l}_{\mathsf{w}+m}^2e_{\mathsf{w}+m}\\
&\le C^z_{env}\,\mathsf{l}_{\mathsf{w}}^2e_{\mathsf{w}}\sum_{m\ge0}(m+1)\hat{\mathsf{l}}_m^2q_a^{(m-1)_+}\rho^{-m}
=C^{TC}_\Sigma C^z_{env}\,\mathsf{l}_{\mathsf{w}}^2e_{\mathsf{w}} .
\end{align*}
The second inequality uses $\mathsf{l}_{\mathsf{w}+m}\le\hat{\mathsf{l}}_m\mathsf{l}_{\mathsf{w}}$ and $e_{\mathsf{w}+m}\le\rho^{-m}e_{\mathsf{w}}$, from \eqref{9.1:eq-load-envelope-a}. This is \eqref{9.1:eq-load-envelope-2}.

It remains to show $C^{TC}_\Sigma<\infty$. For $m\ge2$ the summand equals
\[
(m+1)\hat{\mathsf{l}}_m^2\rho^{-1}\Bigl(\frac{q_a}{\rho}\Bigr)^{m-1}\le(m+1)\hat{\mathsf{l}}_m^2\rho^{3m-4},
\]
because $q_a\le q^2\le\rho^4$. Since $\rho<1$ and $\hat{\mathsf{l}}_m$ grows only like a power of $\log m$, the series converges.

All constants in \eqref{9.1:eq-load-envelope-b} and $C^{TC}_\Sigma$ are built from the listed parameters and from constants that do not involve $n$, $k$, $j$, $s$, $N_w$, $m_0$, the torus or the position.
\end{proof}

We note two points about the lemma.

First, depths $s>n/2$ are covered. The response proposition for the paired runs, assumed in hypothesis (a), holds along the whole chain $0\le s\le n$ with depth-free constants, with its boundary row at $s=n$. The output of $\mathsf{B}$ at its private finest level has no partner in $\mathsf{A}$ and enters as the extra slot $G_n$. No comparison of vacua at depth $r>n/2$ is used: the comparison is between terms, whose content in old zones is suppressed by their age.

Second, $\mathsf{l}_{\mathsf{w}}=(\log\theta_{\mathsf{w}})^{c_1}$ depends on $g$ through $\mathsf{w}=\mathsf{w}(g)$ and $X=g^{-2}$. It does not depend on $n$ or on $N_w$.

\begin{lemma}[Comparison of aligned terms at common ramps]\label{9.1:lem-aligned-terms}
Work in the setting of Notation~\ref{9.1:not-centre-normalisation}. There $\mathsf{A}$ and
$\mathsf{B}$ are the two runs, $\mathcal{V}_w$ is the configuration space of the last lattice
$\mathbb{T}_w$ with its Haar measure, and $k$ is the stop level. Assume the following.
\begin{itemize}
\item The interpolation statement on the dial domain $\mathbb{D}^{\star}$, with its clauses (c1),
(c2)$^{\star}$ and (c3), in the form with resolved boundary terms given by
Lemma~\ref{9.1:lem-resolved-interpolation}. This form includes the identification of live
components \eqref{9.1:eq-resolved-interpolation-a}.
\item The first-variation statement on $\mathbb{D}^{\star}$, with its clauses (a) and (c).
\item Geometric forgetting of old regions, together with its corollary on old live components.
\item The slot hypothesis of Lemma~\ref{9.1:lem-resolved-interpolation}, read at total dial
$2E_1$.
\item $g\le\gamma_R$, where $\gamma_R>0$ is a coupling ceiling; the lemma uses $\gamma_R$ only
through this inequality.
\end{itemize}
For an aligned label $h\in H_{al}$, let $\tau_h^{\mathsf{B}}[t^{\mathsf{A}}]$ denote the term of
$\mathsf{B}$ formed with the objects of $\mathsf{B}$ and with the ramps $t^{\mathsf{A}}$ of
$\mathsf{A}$ in the common variables. This is the term at the end $(\lambda,\varepsilon)=(1,0)$ of
the interpolation of clause (c1).

Let $\varphi=(\varphi_h)_{h\in H_{al}}$ be a family of Borel functions on $\mathcal{V}_w$ with
$|\varphi_h|\le 1$, each $\varphi_h$ being read in the chart of the label $h$. Then for every
$z\in P_{UV}$, and also at the zero source,
\begin{equation}\label{9.1:eq-aligned-terms-1}
\Bigl|\sum_{h\in H_{al}}\Bigl[\int\varphi_h\,d\tau_h^{\mathsf{B}}[t^{\mathsf{A}}]
-\int\varphi_h\,d\tau_h^{\mathsf{A}}\Bigr]\Bigr|
\le E_1^{-1}\,\Sigma^{q_a}_k\cdot A^{\sharp}.
\end{equation}
Consequently
\begin{equation}\label{9.1:eq-aligned-terms-2}
\sum_{h\in H_{al}}\bigl\|\tau_h^{\mathsf{B}}[t^{\mathsf{A}}]-\tau_h^{\mathsf{A}}\bigr\|
\le E_1^{-1}\,\Sigma^{q_a}_k\,A^{\sharp}.
\end{equation}
\end{lemma}

\begin{proof}
Throughout the proof the point of $P_{UV}$ in the statement is fixed. From now on the
letter $z$ denotes a class $z=(j,k')$, as in the marks $W^{(z)}$. We write $\|\cdot\|$ for the
total variation of a complex measure on $\mathcal{V}_w$.

\emph{The function.} For an array $x$ of slot values let $\tau_h[x;t^{\mathsf{A}}]:=\mathcal{P}_h[x]$
be the assembly formula of the label $h$ applied to $x$ at the ramps $t^{\mathsf{A}}$, as in
\eqref{9.1:eq-centre-normalisation-c}. Put
\[
F(x):=\sum_{h\in H_{al}}\int\varphi_h\,d\tau_h[x;t^{\mathsf{A}}].
\]
On the $z$-fibre of $\mathbb{D}^{\star}$ (here and below $z=(j,k')$ is a class) define
\begin{equation}\label{9.1:eq-aligned-terms-4}
x^{\lambda,\varepsilon}:=x^{\mathsf{A}}+\lambda\bigl(x^{\mathsf{B}}-x^{\mathsf{A}}\bigr)
+\varepsilon W^{(z)},\qquad
\mathsf{G}_z(\lambda,\varepsilon):=F\bigl(x^{\lambda,\varepsilon}\bigr).
\end{equation}
We first establish the regularity of each term. By clause (c2)$^{\star}$, for every Borel set
$E\subset\mathcal{V}_w$ the number $\tau_h[x^{\lambda,\varepsilon};t^{\mathsf{A}}](E)$ is
holomorphic in $\varepsilon$ and continuous in $\lambda$. To see this, note three facts.
\begin{itemize}
\item The measure $\tau_h[x^{\lambda,\varepsilon};t^{\mathsf{A}}]$ is the $\mathcal{V}_w$-marginal
of an integrand that $\mathcal{P}_h$ assembles from the array; its value on $E$ is the number in
question.
\item That integrand is entire in the entries of the array.
\item The integrand is dominated, uniformly in $(\lambda,\varepsilon)$, by the majorant of the
moduli argument in the proof of the correlation theorem, taken at the activity constant
$\mathsf{c}^{\star}$.
\end{itemize}
Holomorphy of the marginal then follows from Morera's theorem under this majorant, as in that
moduli argument.

Next we pass this regularity to $\mathsf{G}_z$. Since $|\varphi_h|\le1$, dominated convergence
carries holomorphy in $\varepsilon$ and continuity in $\lambda$ over to
$\int\varphi_h\,d\tau_h[x^{\lambda,\varepsilon};t^{\mathsf{A}}]$. Using the summability
\eqref{9.1:eq-aligned-terms-5} below, dominated convergence also carries it over to the sum over
labels, and hence to $\mathsf{G}_z$.

The hypothesis of the first-variation statement on the classes is also met. At fixed $n$ the
classes $z=(j,k')$ with $j\ge0$, $k'\ge j$ and $k'\le k$ are finitely many. Moreover, joint
holomorphy near the origin in the class parameters $(\varepsilon_z)_z$ together with the
interpolation parameter of that statement holds wherever the total dial of every slot stays at
most $2E_1$. This uses the same majorant.

\emph{The modulus bound, in moduli form.} For an integer $a\ge0$ and $h\in H_{al}$, let $m_a(h)$
be the number of live components of $h$ whose age is at least $a$. We claim the following. For
$(\lambda,\varepsilon)\in\mathbb{D}^{\star}$, every integer $a\ge0$ and every $\mu\in\mathbb{C}$
with $|\mu|\le q_a^{-a}$,
\begin{equation}\label{9.1:eq-aligned-terms-3}
\sum_{h\in H_{al}}|\mu|^{m_a(h)}\,
\bigl\|\tau_h[x^{\lambda,\varepsilon};t^{\mathsf{A}}]\bigr\|\le A^{\sharp}.
\end{equation}

We prove the claim in three steps.

First, the proof of the upper member of the torus bound of the correlation theorem bounds
$\sum_h\|\tau_h\|$ by the moduli of the factors. Here $\sum_h\|\tau_h\|$ is the sum of the total
variations, integrated over $\mathcal{V}_w$. Of the objects that enter the terms, that proof uses
only the following:
\begin{itemize}
\item the bounds \cite[(2.43)--(2.49)]{B-CMP119} and \cite[(1.68)--(1.69)]{B-CMP122b}, in their
form with the complex source;
\item the right member of \cite[(1.79)]{B-CMP122b} as printed, in the form in which the moduli
argument uses it;
\item the raw Wilson strength;
\item \cite[(1.97)]{B-CMP122b}.
\end{itemize}

Second, clause (c2)$^{\star}$ supplies all of these for the array $x^{\lambda,\varepsilon}$ on
$\mathbb{D}^{\star}$ at the activity constant $\mathsf{c}^{\star}$. In doing so:
\begin{itemize}
\item the ramps are those of $\mathsf{A}$, with the majorants of the random-walk bound on the ramps;
\item the raw Wilson strength and the gains are those of $\mathsf{A}$;
\item the data $\mathsf{s}^0$ and $D^0$ are real;
\item the reading unit is the doubled one.
\end{itemize}

Third, geometric forgetting of old regions holds on $\mathbb{D}^{\star}$ by clause (c2)$^{\star}$.
It leaves on each live component of age $a'$ an unused factor $q_a^{a'+1}$. For $a'\ge a$ this
factor pays for $|\mu|$:
\[
|\mu|\,q_a^{a'+1}\le q_a\le 1.
\]
This is the statement of the corollary on old live components. It is applied under the sum over
labels, with the weights $|\varphi_h|\le1$ inside the moduli. This proves
\eqref{9.1:eq-aligned-terms-3}.

Two consequences of \eqref{9.1:eq-aligned-terms-3} will be used.
\begin{itemize}
\item[(i)] Take $a=0$ and $\mu=1$. Since $|\varphi_h|\le1$ gives
$|\int\varphi_h\,d\tau_h|\le\|\tau_h\|$, we obtain on $\mathbb{D}^{\star}$
\begin{equation}\label{9.1:eq-aligned-terms-5}
|\mathsf{G}_z(\lambda,\varepsilon)|
\le\sum_{h\in H_{al}}\bigl\|\tau_h[x^{\lambda,\varepsilon};t^{\mathsf{A}}]\bigr\|
\le A^{\sharp}.
\end{equation}
\item[(ii)] Take $|\mu|=q_a^{-a}$. Each label with $m_a(h)\ge1$ carries a factor
$|\mu|^{m_a(h)}\ge q_a^{-a}$. Hence
\begin{equation}\label{9.1:eq-aligned-terms-a}
\sum_{h\in H_{al}:\ m_a(h)\ge1}
\bigl\|\tau_h[x^{\lambda,\varepsilon};t^{\mathsf{A}}]\bigr\|\le q_a^{a}A^{\sharp}.
\end{equation}
\end{itemize}
Neither a Schwarz-type step nor any cancellation between different labels is used.

\emph{The $\varepsilon$-sensitive part.} Fix a class $z=(j,k')$. By clause (c3), the variable
$\varepsilon$ enters $\tau_h[x^{\lambda,\varepsilon};t^{\mathsf{A}}]$ only if the zone $j$ is
present in $h$ at level $k'$.

Suppose $j\le k-2$. Then this presence forces a large-field region $Z_{j+1}$ of $h$ that is alive
at level $k'$. The label $h$ is a label of the final density, so its large-field regions only grow
until the operation $\mathbf{R}$ removes them. Two cases arise.
\begin{itemize}
\item If the region has been removed by an $\mathbf{R}$-step, its inner history is summed inside
the stored boundary terms $R'$ and $\mathbf{B}'$ of that $\mathbf{R}$-step. This holds because the
sum in the definition \cite[(1.71)]{B-CMP122b} acts also on the last exponential in
\cite[(1.72)]{B-CMP122b}. The pair-difference of these stored terms is that of the own part
$\mathbf{B}'^{\,own}$. That part is attributed to the larger of the level of the stored boundary
term and the block zone, and this zone is present.
\item If the region survives, the classes carrying it are resolved members, for which
\eqref{9.1:eq-resolved-interpolation-a} gives the live component directly; see clause (c) of
Corollary~\ref{9.1:cor-envelope-corollary}.
\end{itemize}
Hence $h$ has, at the stop, a live component $C$ created at a step $j(C)\le j+1$. At the stop
level $k$ the age of $C$ is $k-j(C)\ge k-j-1$, so
\[
m_{(k-j-1)_+}(h)\ge1.
\]

Let $c_{\lambda,z}$ be the sum of $\int\varphi_h\,d\tau_h[x^{\lambda,\varepsilon};t^{\mathsf{A}}]$
over those aligned labels $h$ in which the class $z$ does not occur. By clause (c3) it does not
depend on $\varepsilon$.

On the $z$-fibre of $\mathbb{D}^{\star}$, that is for $|\varepsilon|\le\Lambda_z$ with
$\Lambda_z=E_1/[\![\delta]\!]^{\natural+}_j$ and $\lambda$ in its range,
\eqref{9.1:eq-aligned-terms-a} with $a=(k-j-1)_+$ gives
\begin{equation}\label{9.1:eq-aligned-terms-6}
\bigl|\mathsf{G}_z(\lambda,\varepsilon)-c_{\lambda,z}\bigr|
\le\sum_{h\in H_{al}:\ z\text{ occurs in }h}
\bigl\|\tau_h[x^{\lambda,\varepsilon};t^{\mathsf{A}}]\bigr\|
\le q_a^{(k-j-1)_+}A^{\sharp}=:\mathfrak{B}_z.
\end{equation}
For $j\ge k-1$ the exponent is $0$, and \eqref{9.1:eq-aligned-terms-5} gives the same bound.

\emph{The first variation.} We apply the first-variation statement with
\[
x^0:=x^{\mathsf{A}},\qquad x^1:=x^{\mathsf{B}},\qquad W:=\sum_z W^{(z)}.
\]
By \eqref{9.1:eq-centre-normalisation-c}, $W=x^1-x^0$. Clause (c) of that statement, and then its
clause (a), give
\[
F(x^{\mathsf{B}})-F(x^{\mathsf{A}})=\sum_z\mathfrak{V}^{(z)},\qquad
\mathfrak{V}^{(z)}=\int_0^1 d\lambda\,
\frac{\partial}{\partial\varepsilon}\mathsf{G}_z(\lambda,\varepsilon)\Big|_{\varepsilon=0},
\]
with
\[
|\mathfrak{V}^{(z)}|\le\frac{\mathfrak{B}_z}{\Lambda_z}
=q_a^{(k-j-1)_+}A^{\sharp}\,\frac{[\![\delta]\!]^{\natural+}_j}{E_1}.
\]
The equality is the definition of $\mathfrak{B}_z$ in \eqref{9.1:eq-aligned-terms-6} together with
that of $\Lambda_z$.

Now we sum over the classes. For fixed $j$ the bound does not depend on $k'$, and $k'$ ranges over
the $k-j+1$ integers from $j$ to $k$. Summing first over $k'$ and then over $j$ gives
\[
\bigl|F(x^{\mathsf{B}})-F(x^{\mathsf{A}})\bigr|
\le E_1^{-1}\sum_{0\le j\le k}(k-j+1)\,q_a^{(k-j-1)_+}\,[\![\delta]\!]^{\natural+}_j\cdot A^{\sharp}
=E_1^{-1}\,\Sigma^{q_a}_k\,A^{\sharp}.
\]
The last equality is the definition of the age-weighted load sum,
\[
\Sigma^{q_a}_k=\sum_{0\le j\le k}(k-j+1)\,q_a^{(k-j-1)_+}\,[\![\delta]\!]^{\natural+}_j .
\]

It remains to identify the two values of $F$. By clause (c1),
\[
F(x^{\mathsf{A}})=\sum_{h\in H_{al}}\int\varphi_h\,d\tau_h^{\mathsf{A}},
\qquad
F(x^{\mathsf{B}})=\sum_{h\in H_{al}}\int\varphi_h\,d\tau_h^{\mathsf{B}}[t^{\mathsf{A}}].
\]
This proves \eqref{9.1:eq-aligned-terms-1}.

\emph{Total variation.} The bound \eqref{9.1:eq-aligned-terms-1} is uniform in the family
$\varphi$. Put
\[
\nu_h:=\tau_h^{\mathsf{B}}[t^{\mathsf{A}}]-\tau_h^{\mathsf{A}}.
\]
By \eqref{9.1:eq-aligned-terms-5}, applied at the two ends of the interpolation, each $\nu_h$ is a
complex measure of finite total variation. Let $d\nu_h=e^{i\psi_h}\,d|\nu_h|$ be its polar
decomposition, with $\psi_h$ Borel. Take $\varphi_h:=e^{-i\psi_h}$, which is Borel with
$|\varphi_h|=1$. Then $\int\varphi_h\,d\nu_h=\|\nu_h\|$, so the left member of
\eqref{9.1:eq-aligned-terms-1} becomes $\sum_{h\in H_{al}}\|\nu_h\|$. This is
\eqref{9.1:eq-aligned-terms-2}.

\emph{The zero source.} The zero source is the point $\hat w=u$ with vanishing source coefficients
of the same family. The argument above applies there without change.
\end{proof}

\begin{remark}[The constants in Lemma~\ref{9.1:lem-aligned-terms}]
\leavevmode
\begin{itemize}
\item[(a)] The coefficient of $E_1^{-1}$ in \eqref{9.1:eq-aligned-terms-1} is $1$. This is because
clause (a) of the first-variation statement is applied directly, with the modulus bound
$\mathfrak{B}_z$ taken from \eqref{9.1:eq-aligned-terms-a}. Thus no Schwarz-type step enters,
which would cost a factor $2$, and no factor $2$ from the corollary on old live components enters
either. The weaker bound with $4E_1^{-1}$ in place of $E_1^{-1}$ follows a fortiori; later
estimates of this chapter use \eqref{9.1:eq-aligned-terms-1} in that weaker form.
\item[(b)] The classical part of each mark $W^{(z)}$ is paid for by the classical load
$[\![\delta]\!]^c_j$. This load is contained in the total load $[\![\delta]\!]^{\natural+}_j$ that
defines the class radius $\Lambda_z$.
\end{itemize}
\end{remark}

\begin{lemma}[The ramp comparison and unaligned common labels]\label{9.1:lem-ramp-comparison}
Let the two runs $\mathsf{A}$ and $\mathsf{B}$, the last lattice $\mathbb{T}_w$ of side $N_w$ with
configuration space $\mathcal{V}_w$, and the stop level $k$ be as in
Notation~\ref{9.1:not-centre-normalisation}. For $\mathsf{m}\in\{\mathsf{A},\mathsf{B}\}$ let
$H^{\mathsf{m}}$ be the label set of run $\mathsf{m}$, let $\tau_h^{\mathsf{m}}$, $h\in H^{\mathsf{m}}$,
be the $h$-th term of its normalised sourced integral $\tilde{Z}^{\mathsf{m}}$, a complex measure on
$\mathcal{V}_w$ with total variation norm $\|\tau_h^{\mathsf{m}}\|$, let $t^{\mathsf{m}}$ be the ramps
of run $\mathsf{m}$, let $H_{al}$ be the set of aligned labels, and let
$\tau_h^{\mathsf{B}}[t^{\mathsf{A}}]$ be the term of $\mathsf{B}$ with the objects of $\mathsf{B}$ and
the ramps of $\mathsf{A}$ in the common variables, as in Lemma~\ref{9.1:lem-aligned-terms}. Write
$H^{\mathsf{m}}_{nc}$ for the set of labels $h\in H^{\mathsf{m}}$ that are common to the two runs,
are not aligned, and have first index larger than the cut index $n_c$ of this chapter. Assume:
\begin{itemize}
\item[(a)] the ramp-weight statements of this chapter, including the common-slot statement (the slots
are common to the two runs on common labels), the label-by-label modulus bound, the doubling persistence
bound, the weight statement, the proposition of the ramp-weight statements from which the move corollary
is drawn, and the move corollary with its clauses (C1) and (C2), these two clauses in the form derived
from the smallness hypothesis of this chapter;
\item[(b)] $g\le\gamma_{\mathrm{ramp}}$, where $\gamma_{\mathrm{ramp}}$ is the coupling ceiling of the
statements in (a);
\item[(c)] the standing hypothesis of this chapter in the form in which the statements in (a) are
applied here.
\end{itemize}
Then for every $z\in P_{UV}$ and at the zero source
\begin{equation}\label{9.1:eq-ramp-comparison-1}
\sum_{h\in H_{al}}\bigl\|\tau_h^{\mathsf{B}}-\tau_h^{\mathsf{B}}[t^{\mathsf{A}}]\bigr\|
+\sum_{\mathsf{m}\in\{\mathsf{A},\mathsf{B}\}}\ \sum_{h\in H^{\mathsf{m}}_{nc}}
\bigl\|\tau_h^{\mathsf{m}}\bigr\|
\le C_{test}\,\vartheta_\chi^{k}\,A^{\sharp},
\end{equation}
where
\begin{equation}\label{9.1:eq-ramp-comparison-2}
C_{test}:=2e\bigl(80\,C_\chi N_w^4+9\,C_\chi^{\sharp}\vartheta_\chi^{-1}+2\bigr)
\end{equation}
is the constant of the move corollary in the form of its clause (C1); here $C_\chi$ is the constant of
the stop-level slot bound and $C_\chi^{\sharp}$ that of the admitted multiplier range in the statements
in (a). The constant $C_{test}$ is a number depending only on $d$,
$\mathfrak{p}$, $A_0/A_1$, $C_\chi$ and $N_w$. Only the form of the bound
\eqref{9.1:eq-ramp-comparison-1}, a number times $\vartheta_\chi^{k}A^{\sharp}$, is used later in this
chapter.
\end{lemma}

\begin{proof}
The moves of the move corollary are made inside one run at a time, and on all live components of that
run at once. Three of them are used here: the move (m3) in $\mathsf{B}$, the move (m2) in $\mathsf{B}$
and the move (m4) in $\mathsf{A}$. The move (m1) is not used in this lemma; it enters the estimate of
the remaining terms.

\emph{The move \textup{(m3)}.} It acts in $\mathsf{B}$ on the aligned labels. For a complex parameter
$\mathsf{s}$ it replaces the ramps by $t^{\mathsf{A}}+\mathsf{s}(t^{\mathsf{B}}-t^{\mathsf{A}})$, and at
the same time
it interpolates the small slots of the production block $P$ at the stop level. Write $i(\iota)$ for the
level read by a slot $\iota$, and $\eta^{\mathsf{A}}_\iota$, $\eta^{\mathsf{B}}_\iota$ for the values of
the slot $\iota$ in the two runs. The slots concerned read $V_k$ only, that is $i(\iota)=k$, and there are
at most $N_w^4$ of them. Each interpolated slot obeys
\begin{equation}\label{9.1:eq-ramp-comparison-3}
\bigl|\eta^{\mathsf{A}}_\iota+\mathsf{s}(\eta^{\mathsf{B}}_\iota-\eta^{\mathsf{A}}_\iota)\bigr|
\le\bigl(1+20\,C_\chi\vartheta_\chi^{k}|\mathsf{s}|\bigr)\,
\mathbf{1}\bigl\{f^{\mathsf{B}}<(1+\varpi_0)c\bigr\},
\end{equation}
where the indicator restricts to the configurations on which the quantity $f^{\mathsf{B}}$ attached to
the slot in $\mathsf{B}$ lies below $(1+\varpi_0)$ times the threshold $c$.
Hence the product of the moduli of these at most $N_w^4$ slots is bounded by the product of the
indicators of \eqref{9.1:eq-ramp-comparison-3} times
\begin{equation*}
\bigl(1+20\,C_\chi\vartheta_\chi^{k}|\mathsf{s}|\bigr)^{N_w^4}
\le\exp\bigl(20\,C_\chi N_w^4\vartheta_\chi^{k}|\mathsf{s}|\bigr)\le e
\qquad\text{for }|\mathsf{s}|\le\bigl(20\,C_\chi N_w^4\vartheta_\chi^{k}\bigr)^{-1}.
\end{equation*}
Write $\tau_h[\mathsf{s}]$ for the term of $\mathsf{B}$ with label $h\in H_{al}$ after this move. At
$\mathsf{s}=0$ the
ramps and the interpolated stop-level slot values are those of $\mathsf{A}$ and the objects those of
$\mathsf{B}$; this is the term $\tau_h[0]=\tau_h^{\mathsf{B}}[t^{\mathsf{A}}]$ of
Lemma~\ref{9.1:lem-aligned-terms}. At $\mathsf{s}=1$ they are
those of $\mathsf{B}$, so $\tau_h[1]=\tau_h^{\mathsf{B}}$.

\emph{The moves \textup{(m2)} and \textup{(m4)}.} The move (m2) acts in $\mathsf{B}$, the move (m4) in
$\mathsf{A}$. Each multiplies the terms with common non-aligned labels, in particular those with labels
in $H^{\mathsf{m}}_{nc}$, by a complex number $\mu$,
\begin{equation*}
|\mu|\le\bigl(4\,C_\chi^{\sharp}\vartheta_\chi^{k}\bigr)^{-1},
\end{equation*}
which is the range admitted by the clause of the doubling persistence bound on multiplication of the
common non-aligned labels.

\emph{Persistence of the upper bound.} Each of the three moves is followed by the increment estimate on
the disc of its parameter that accompanies the move corollary. That estimate uses the following fact:
the upper bound of the correlation theorem persists when the majorant of each component is doubled,
since $2\le e^{\mathsf{V}}$ for the exponent $\mathsf{V}$ attached to each component, by the lower bound
on these exponents stated in this chapter. We write $\mathcal{M}^{\mathsf{m}}$ for this persisting
upper bound, that is, the upper bound of the correlation theorem for run $\mathsf{m}$ with the majorant
of each component doubled. This persistence is the
content of the doubling persistence bound together with the label-by-label modulus bound, and it is a
bound label by label on moduli. Indeed, write the integrand of a label as a product $\prod_\iota v_\iota$
over its components, with factors $v^{\mathsf{A}}_\iota$, $v^{\mathsf{B}}_\iota$ in the two runs, the
components being ordered. The telescoping identity
\begin{equation*}
\prod_\iota v^{\mathsf{B}}_\iota-\prod_\iota v^{\mathsf{A}}_\iota
=\sum_\iota\bigl(v^{\mathsf{B}}_\iota-v^{\mathsf{A}}_\iota\bigr)
\prod_{\sigma<\iota}v^{\mathsf{B}}_\sigma\prod_{\sigma>\iota}v^{\mathsf{A}}_\sigma
\end{equation*}
and the triangle inequality give
\begin{equation}\label{9.1:eq-ramp-comparison-4}
\Bigl|\prod_\iota v^{\mathsf{B}}_\iota-\prod_\iota v^{\mathsf{A}}_\iota\Bigr|
\le\sum_\iota\bigl|v^{\mathsf{B}}_\iota-v^{\mathsf{A}}_\iota\bigr|
\prod_{\sigma<\iota}\bigl|v^{\mathsf{B}}_\sigma\bigr|\prod_{\sigma>\iota}\bigl|v^{\mathsf{A}}_\sigma\bigr|,
\end{equation}
and each summand on the right of \eqref{9.1:eq-ramp-comparison-4} is an admissible filling in the class
$\mathcal{U}[\varpi_0]$ of admissible fillings with parameter $\varpi_0$. For such a filling $\omega$ let
$\mathsf{V}(\omega)$ be its exponent, and let $\mathsf{e}$ be the coefficient fixed with the weight
statement. The sum over the fillings is controlled by $\sum_\omega e^{-\mathsf{e}\mathsf{V}(\omega)}\le1$,
by clause (ii) of the weight statement and by geometric forgetting of old regions. Since the bound is a
bound on moduli taken label by label,
it continues to hold when a Borel weight $\varphi_h$ with $|\varphi_h|\le1$ is inserted in the
integrand of each label $h$: the bounded function $F$ admitted in the doubling persistence bound
contains $\varphi_h$ together with the last exponential of \cite[(1.72)]{B-CMP122b}, the slots being
functions of the integration variables only. The resulting bound is uniform in the family
$\varphi=(\varphi_h)_h$.

\emph{Total variation from the increments.} Consequently the increments given by the disc estimate
for the move (m3) bound
\begin{equation*}
\Bigl|\sum_{h\in H_{al}}\int\varphi_h\,d\bigl(\tau_h[1]-\tau_h[0]\bigr)\Bigr|
\end{equation*}
by the quantity that this estimate assigns to the move (m3),
uniformly in all families of Borel functions $\varphi_h$ on $\mathcal{V}_w$ with $|\varphi_h|\le1$.
For a finite family of complex measures $(\nu_h)_h$ on $\mathcal{V}_w$ one has
\begin{equation*}
\sum_h\|\nu_h\|=\sup_{\varphi}\Bigl|\sum_h\int\varphi_h\,d\nu_h\Bigr|,
\end{equation*}
the supremum over all such families, as in the passage from the first to the second bound of
Lemma~\ref{9.1:lem-aligned-terms}: the inequality $\le$ is attained by taking for $\varphi_h$ the
complex conjugate of the density of $\nu_h$ with respect to $|\nu_h|$, and $\ge$ follows from
$|\int\varphi_h\,d\nu_h|\le\|\nu_h\|$. Applied to
$\nu_h=\tau_h[1]-\tau_h[0]=\tau_h^{\mathsf{B}}-\tau_h^{\mathsf{B}}[t^{\mathsf{A}}]$, the supremum over
$\varphi$ gives the first sum in \eqref{9.1:eq-ramp-comparison-1}. For the moves (m2) and (m4), the
terms of the labels in $H^{\mathsf{m}}_{nc}$ are homogeneous in the multiplier $\mu$, and the disc
estimate on $|\mu|\le(4C_\chi^{\sharp}\vartheta_\chi^{k})^{-1}$ gives, by the same supremum over
$\varphi$,
\begin{equation*}
\sum_{h\in H^{\mathsf{m}}_{nc}}\bigl\|\tau_h^{\mathsf{m}}\bigr\|
\le4\,C_\chi^{\sharp}\vartheta_\chi^{k}\,\mathcal{M}^{\mathsf{m}},
\qquad\mathsf{m}\in\{\mathsf{A},\mathsf{B}\}.
\end{equation*}

\emph{Objects and slots.} In each move the objects are those of the run in which the move is made, at
its own real centre; their modulus $A$ satisfies $A\le A^{\sharp}$ by the bound on the object modulus at
the real centre. The
slots are functions of the integration variables only, and on common labels they are common to the two
runs by the common-slot statement; therefore each move is an operation inside a single run, as required
by the move corollary. Adding the bound for (m3) and the bounds for (m2) and (m4), read with these
objects and $A\le A^{\sharp}$, with the constant of the move corollary in the form of its clause (C1),
which is $C_{test}$ of \eqref{9.1:eq-ramp-comparison-2}, gives \eqref{9.1:eq-ramp-comparison-1}.
\end{proof}

\begin{lemma}[The old unpaired labels]\label{9.1:lem-unpaired-labels}
Let the runs $\mathsf{A}$ and $\mathsf{B}$, the stop depth $\mathsf{w}=\mathsf{w}(g)$, the stop level
$k:=n-\mathsf{w}$, the last lattice $\mathbb{T}_w$ with configuration space $\mathcal{V}_w$, the label
sets $H^{\mathsf{A}}$, $H^{\mathsf{B}}$ and the terms $\tau_h^{\mathsf{c}}$, $\mathsf{c}\in\{\mathsf{A},\mathsf{B}\}$,
be as in Notation~\ref{9.1:not-centre-normalisation}. Each $\tau_h^{\mathsf{c}}$ is a complex measure
on $\mathcal{V}_w$, and $\|\tau_h^{\mathsf{c}}\|$ is its total variation. For a label $h$ write $j(h)$
for its first index. Let $n_c$ be an integer, a level counted by the index of $\mathsf{A}$ and
depending on the cutoff, for which hypothesis (c) below holds; of $n_c$ only (c) and (d) are used.
Assume $g\le\gamma_R$, where $\gamma_R$ is a threshold on the reference coupling $g$, and:
\begin{itemize}
\item[(a)] in each of the runs $\mathsf{A}$ and $\mathsf{B}$: geometric forgetting of old regions,
together with its corollary bounding the labels that carry a live component of prescribed age;
\item[(b)] in each of the runs $\mathsf{A}$ and $\mathsf{B}$: the torus bound of the correlation theorem,
in its modulus form;
\item[(c)] the centre-depth bound for the level $n_c$, read with the constants $C^{\natural}_{cl}$ and
$\vartheta_2^{1/4}$:
\begin{equation}\label{9.1:eq-unpaired-labels-3}
n_c\le k_*+1+\frac{\log\bigl(C^{\natural}_{cl}\,(\log(x_0+2\bar c_2))^{\bar c_2}\bigr)}
{\log\vartheta_2^{-1/4}} ,
\end{equation}
where $x_0$ is the bare inverse coupling of the run $\mathsf{A}$, $k_*$ is the integer entering this
bound, which does not depend on $n$, and $\bar c_2$ is the constant entering this bound, distinct from
the constant $c_2$ with $c_5=2c_2$ of the glossary;
\item[(d)] the rule fixing the level $n_c$: $n_c$ depends on the cutoff and enters the comparison of
the two runs only through the bound of the present lemma.
\end{itemize}
Then for every $z\in P_{UV}$ and at the zero source
\begin{equation}\label{9.1:eq-unpaired-labels-1}
\sum_{\substack{h\in H^{\mathsf{A}}\\ j(h)\le n_c}}\|\tau_h^{\mathsf{A}}\|
+\sum_{\substack{h'\in H^{\mathsf{B}}\\ j(h')\le n_c}}\|\tau_{h'}^{\mathsf{B}}\|
\le 4q_a^{\,k-n_c}A\le 4q_a^{\,k-n_c}A^{\sharp},
\end{equation}
and
\begin{equation}\label{9.1:eq-unpaired-labels-2}
4q_a^{\,k-n_c}\le C_{nc}\,\rho^{\,n-\mathsf{w}}\le C_{nc}\,e_{\mathsf{w}},
\end{equation}
where the constant $C_{nc}$, the explicit constant \eqref{9.1:eq-unpaired-labels-4} of the proof,
depends on $(L,\mathfrak{p},q_a,\rho,k_*,\mathsf{w}(g),X)$ only, and not on $n$ and not on $N_w$.
\end{lemma}

\begin{proof}
\emph{Age.} Let $\mathsf{c}\in\{\mathsf{A},\mathsf{B}\}$. By the normalisation of the two runs at their
own centres, \eqref{9.1:eq-centre-normalisation-b}, levels are counted by $\mathsf{A}$'s index and the
stop is at level $k$. Hence a label of first index at most $n_c$ that is alive at the stop has a live
component of age at least $k-n_c$. In $\mathsf{B}$ this includes the labels born at $\mathsf{B}$'s private
finest level $-1$; these have age $k+1\ge k-n_c$. Put
\begin{equation*}
a:=k-n_c ,
\end{equation*}
and note that, in the notation $m_a(h)$ of the corollary in hypothesis (a), $m_a(h)\ge1$ holds exactly
when the label $h$ has a live component of age at least $a$. The labels of $H^{\mathsf{c}}$ with
$j(h)\le n_c$ then
satisfy $m_a(h)\ge1$, so their sum in
\eqref{9.1:eq-unpaired-labels-1} is at most the sum of $\|\tau_h^{\mathsf{c}}\|$ over the labels with
$m_a(h)\ge1$.

In each run separately, under hypotheses (a) and (b) for that run, we apply the corollary of
geometric forgetting of old regions at $a=k-n_c$, in the form in which
\eqref{9.1:eq-aligned-terms-a} bounds the total variations of
the individual terms; no interpolated array of slot values enters, the interpolation parameter taking
only its two end values $0$ and $1$. This gives
\begin{equation*}
\sum_{\substack{h\in H^{\mathsf{c}}\\ m_a(h)\ge1}}\|\tau_h^{\mathsf{c}}\|
\le q_a^{\,a}A\le 2q_a^{\,a}A ;
\end{equation*}
we keep the factor $2$, which is the form in which the corollary is stated.
Adding the two runs yields the bound $4q_a^{\,k-n_c}A$, which is the first inequality in
\eqref{9.1:eq-unpaired-labels-1}; the second holds because $A^{\sharp}$ majorises $A$.

\emph{Arithmetic.} The level $n_c$ depends on the cutoff. By hypothesis (d) it enters only through the
bound just proved, and we now replace $q_a^{\,k-n_c}$ by a quantity whose constant reads neither $n$
nor $N_w$. Since $k=n-\mathsf{w}$,
\begin{equation*}
k-n_c=n-\mathsf{w}-n_c .
\end{equation*}
Hypothesis (c) bounds $n_c$ by \eqref{9.1:eq-unpaired-labels-3}. Put
\begin{equation*}
\mathfrak{a}:=\frac{\bar c_2\log q_a^{-1}}{\log\vartheta_2^{-1/4}}
=\frac{4\bar c_2\log q_a^{-1}}{\log\vartheta_2^{-1}} .
\end{equation*}
Since $q_a<1$, the map $t\mapsto q_a^{-t}$ is increasing. For every $Y>0$,
\begin{equation*}
q_a^{-\log Y/\log\vartheta_2^{-1/4}}
=\exp\Bigl(\frac{\log q_a^{-1}\,\log Y}{\log\vartheta_2^{-1/4}}\Bigr)=Y^{\mathfrak{a}/\bar c_2}.
\end{equation*}
Applying this with $Y=C^{\natural}_{cl}(\log(x_0+2\bar c_2))^{\bar c_2}$, the bound
\eqref{9.1:eq-unpaired-labels-3} gives
\begin{equation*}
q_a^{-n_c}\le q_a^{-k_*-1}\,(C^{\natural}_{cl})^{\mathfrak{a}/\bar c_2}\,
(\log(x_0+2\bar c_2))^{\mathfrak{a}} .
\end{equation*}
The factor $(C^{\natural}_{cl})^{\mathfrak{a}/\bar c_2}$ is carried into $C_{nc}$ below.

The bound on the bare inverse coupling of the run $\mathsf{A}$ along the reference flow gives
\begin{equation*}
x_0\le X+3\lambda n+2\bar c_2 ,\qquad X=g^{-2},\quad \lambda=c_0/2 .
\end{equation*}
Hence $\log(x_0+2\bar c_2)\le\log(X+3\lambda n+4\bar c_2)$, and the power $\mathfrak{a}>0$ preserves
this inequality. Now $q_a^{\,k}=q_a^{\,n-\mathsf{w}}=\rho^{\,n-\mathsf{w}}(q_a/\rho)^{n-\mathsf{w}}$, and
since $n\ge\mathsf{w}$ the term at $n'=n$ is at most the supremum over $n'\ge\mathsf{w}$ below, so
\begin{align*}
q_a^{\,k-n_c}
&=\rho^{\,n-\mathsf{w}}\,(q_a/\rho)^{n-\mathsf{w}}\,q_a^{-n_c}\\
&\le\rho^{\,n-\mathsf{w}}\,(C^{\natural}_{cl})^{\mathfrak{a}/\bar c_2}\,q_a^{-k_*-1}
\sup_{n'\ge\mathsf{w}}\,(q_a/\rho)^{n'-\mathsf{w}}\bigl(\log(X+3\lambda n'+4\bar c_2)\bigr)^{\mathfrak{a}}\\
&=\tfrac14\,C_{nc}\,\rho^{\,n-\mathsf{w}},
\end{align*}
where
\begin{equation}\label{9.1:eq-unpaired-labels-4}
C_{nc}:=4\,(C^{\natural}_{cl})^{\mathfrak{a}/\bar c_2}\,q_a^{-k_*-1}
\sup_{n'\ge\mathsf{w}}\,(q_a/\rho)^{n'-\mathsf{w}}
\bigl(\log(X+3\lambda n'+4\bar c_2)\bigr)^{\mathfrak{a}} .
\end{equation}
The supremum is finite. Indeed, by the choice of the contraction rate $\rho$ relative to the age rate
$q_a$, one has $q_a/\rho\le\rho^{3}<1$, so the factor $(q_a/\rho)^{n'-\mathsf{w}}$ decays
geometrically in $n'$. It therefore beats the power $\mathfrak{a}$ of $\log(X+3\lambda n'+4\bar c_2)$,
which grows like a power of $\log n'$, and the product tends to $0$ as $n'\to\infty$. The supremum reads
only $q_a$, $\rho$, $\mathsf{w}(g)$, $X=g^{-2}$, $\lambda$, $\bar c_2$ and $\mathfrak{a}$; the remaining
factors of $C_{nc}$ read $q_a$, $k_*$, $C^{\natural}_{cl}$ and $\mathfrak{a}$. Therefore $C_{nc}$ reads
neither $n$ nor $N_w$. This proves the first inequality in \eqref{9.1:eq-unpaired-labels-2}.

Finally, $e_{\mathsf{w}}=\rho^{\,n-\mathsf{w}}+\Pi_n$ with the non-negative power-law floor $\Pi_n$, so
$\rho^{\,n-\mathsf{w}}\le e_{\mathsf{w}}$. This is the second inequality in
\eqref{9.1:eq-unpaired-labels-2}.
\end{proof}

\begin{theorem}[The two-cutoff comparison of terms at the last lattice]
\label{9.1:thm-last-lattice-comparison}
Let the run $\mathsf{A}$ of length $n$ and the run $\mathsf{B}$ of length $n+1$, their residual
sources $W^{\mathsf{A}}(z)$ and $W^{\mathsf{B}}(z)$, the stop depth $\mathsf{w}=\mathsf{w}(g)$, the
stop level $k:=n-\mathsf{w}$, the last lattice $\mathbb{T}_w$ with $N_w$ sites per side and its
configuration space $\mathcal{V}_w$ be as in Notation~\ref{9.1:not-centre-normalisation}. For
$\mathsf{m}\in\{\mathsf{A},\mathsf{B}\}$, $H^{\mathsf{m}}$ is the label set of the run $\mathsf{m}$,
and $\tau^{\mathsf{m}}_h$, $h\in H^{\mathsf{m}}$, is the $h$-th term of its normalised sourced
integral $\tilde{Z}^{\mathsf{m}}$, a complex measure on $\mathcal{V}_w$. These terms are those of
the expansions on the dial domain $\mathbb{D}^{\star}$, at the total loads
$[\![\delta]\!]^{\natural+}$, the activity constant $\mathsf{c}^{\star}$ and the value
$\kappa_1^{\natural}$ of $\kappa_1$, with the objects of Lemma~\ref{9.1:lem-aligned-terms}.
$H_{al}$ is the set of
aligned labels, and $\|\cdot\|$ is the total-variation norm of complex measures on $\mathcal{V}_w$,
as in Lemmas~\ref{9.1:lem-aligned-terms}, \ref{9.1:lem-ramp-comparison}
and~\ref{9.1:lem-unpaired-labels}. Assume:
\begin{itemize}
\item[(a)] the standing hypotheses of the two expansions, as fixed in this chapter (the standing
statements on which the two runs rest and the supplementary conditions attached to them), insofar
as they are not among the hypotheses named in (b)--(e); among them are the correlation theorem,
the age lemma, the
statement that supplies the exponent $\theta_\chi$ of hypothesis~(d), and the second index
condition $(\iota')$ of \eqref{9.1:eq-centre-normalisation-f};
\item[(b)] the hypotheses of the comparison of aligned terms at common ramps
(Lemma~\ref{9.1:lem-aligned-terms}), of the ramp comparison
(Lemma~\ref{9.1:lem-ramp-comparison}), of the bound for the old unpaired labels
(Lemma~\ref{9.1:lem-unpaired-labels}) and of the envelope of the loads
(Lemma~\ref{9.1:lem-load-envelope});
\item[(c)] $g\le\gamma_U$, and $n\ge n_1$ with $k=n-\mathsf{w}\ge1$, where $n_1$ is the lower
threshold on the length $n$ of the run $\mathsf{A}$;
\item[(d)] the index condition $(\iota)$ of \eqref{9.1:eq-centre-normalisation-f}, namely
$\vartheta_\chi\le\rho$. This is the fourth member of the window display
\eqref{9.1:eq-centre-normalisation-b}, with the exponent $\theta_\chi$ attached to the rate
$\vartheta_\chi$ among the index exponents $\mathfrak{i}$ of that display, that is,
$\theta_\chi\ge\theta_{dl}$;
\item[(e)] for labels with a live second-class strip, the conclusion of
Proposition~\ref{9.1:prop-blind-site-expansions}, whose proof is outstanding.
\end{itemize}
With $\theta_\chi$ among the index exponents $\mathfrak{i}$, the window
$[\max\{q^{1/2},L^{-\vartheta_\sharp},\vartheta_{dl}\},1)$ for $\rho$ is non-empty for every $L$.
Then the following holds for every $\hat{w}\in\mathcal{N}'_w$ and every $z\in P_{UV}$ (in particular
every $z\in P_{IR}\subset P_{UV}$), and also at the zero residual source. The normalisation is
\eqref{9.1:eq-centre-normalisation-a}: each run is taken at its own real centre, with its own
vacuum responses subtracted at that centre. The paired label set is $H_p:=H_{al}$, and the pairing
map $\pi$ is the identity. Writing $\tau^{\mathsf{A}}_h$ for $\tau^{\mathsf{A}}_h(W^{\mathsf{A}}(z))$
and $\tau^{\mathsf{B}}_{h}$ for $\tau^{\mathsf{B}}_{h}(W^{\mathsf{B}}(z))$ wherever the argument
$W^{\mathsf{m}}(z)$ is suppressed, we have
\begin{equation}\label{9.1:eq-last-lattice-comparison-1}
\begin{split}
&\sum_{h\in H_{al}}\bigl\|\tau^{\mathsf{B}}_h(W^{\mathsf{B}}(z))
-\tau^{\mathsf{A}}_h(W^{\mathsf{A}}(z))\bigr\|
+\sum_{h\in H^{\mathsf{A}}\setminus H_{al}}\|\tau^{\mathsf{A}}_h\|\\
&\qquad+\sum_{h'\in H^{\mathsf{B}}\setminus H_{al}}\|\tau^{\mathsf{B}}_{h'}\|
\;\le\;\varepsilon_n A^{\sharp},
\end{split}
\end{equation}
where
\begin{equation}\label{9.1:eq-last-lattice-comparison-a}
\begin{gathered}
\varepsilon_n:=4E_1^{-1}\Sigma^{q_a}_k+4q_a^{\,k-n_c}+C_{test}\vartheta_\chi^{\,k}
\;\le\;C_{TC}(g,N_w)\,e_{\mathsf{w}},\\
C_{TC}:=\frac{4C_\Sigma^{TC}C_{env}^z\,\mathsf{l}_{\mathsf{w}}^2}{E_1}+C_{nc}+C_{test}.
\end{gathered}
\end{equation}
Here $\Sigma^{q_a}_k$ is the age-weighted load sum and $E_1$ the constant of
Lemma~\ref{9.1:lem-aligned-terms}; $C_{test}$ is the constant of Lemma~\ref{9.1:lem-ramp-comparison};
$n_c$ is the cut index of Lemma~\ref{9.1:lem-unpaired-labels} (the labels bounded there are those
whose first index is at most $n_c$) and $C_{nc}$ is the constant of that lemma; $C_\Sigma^{TC}$,
$C_{env}^z$, $E_1^z$, the envelope function $\mathsf{l}_{\mathsf{w}}$ and
$e_{\mathsf{w}}=\rho^{\,n-\mathsf{w}}+\Pi_n$ are those of Lemma~\ref{9.1:lem-load-envelope}, with
$\rho$ the window letter of \eqref{9.1:eq-centre-normalisation-b} and $\Pi_n$ the power-law floor.
The rate bound in \eqref{9.1:eq-last-lattice-comparison-a} also holds with $E_1^z$ in place of
$E_1$ in both places:
\begin{equation*}
4(E_1^z)^{-1}\Sigma^{q_a}_k+4q_a^{\,k-n_c}+C_{test}\vartheta_\chi^{\,k}
\le\Bigl(\frac{4C_\Sigma^{TC}C_{env}^z\,\mathsf{l}_{\mathsf{w}}^2}{E_1^z}+C_{nc}+C_{test}\Bigr)
e_{\mathsf{w}}.
\end{equation*}
The constant $C_{TC}=C_{TC}(g,N_w)$ does not depend on $n$, $k$, $K$, $m_0$ or the position.
\end{theorem}

\begin{proof}
Fix $\hat{w}\in\mathcal{N}'_w$ and $z\in P_{UV}$, or take the zero residual source. Every
statement applied below is applied at the same $\hat{w}$ and $z$ (respectively at the zero source),
in the normalisation \eqref{9.1:eq-centre-normalisation-a}. Hypothesis~(b) supplies the hypotheses
of each lemma used.

For $\mathsf{m}\in\{\mathsf{A},\mathsf{B}\}$, let $H^{\mathsf{m}}_{\mathrm{old}}$ be the set of labels
$h\in H^{\mathsf{m}}$ whose first index, in the level indexing of
Notation~\ref{9.1:not-centre-normalisation}, is at most $n_c$. Let $H^{\mathsf{m}}_{\mathrm{cn}}$ be
the set of labels $h\in H^{\mathsf{m}}$ that are common to both runs, have first index $>n_c$, and
are not aligned. Put
\begin{equation*}
S_{\mathrm{cn}}:=\sum_{\mathsf{m}\in\{\mathsf{A},\mathsf{B}\}}\;\sum_{h\in H^{\mathsf{m}}_{\mathrm{cn}}}
\|\tau^{\mathsf{m}}_h\|,
\end{equation*}
the second sum on the left-hand side of the bound of Lemma~\ref{9.1:lem-ramp-comparison}.

\emph{Aligned labels.} For $h\in H_{al}$, let $\tau^{\mathsf{B}}_h[t^{\mathsf{A}}]$ be the hybrid
term of Lemma~\ref{9.1:lem-aligned-terms}. It is the term of $\mathsf{B}$ with the objects of
$\mathsf{B}$ and the ramps of $\mathsf{A}$ in the common variables. By the triangle inequality for
the total-variation norm,
\begin{equation*}
\|\tau^{\mathsf{B}}_h-\tau^{\mathsf{A}}_h\|
\le\|\tau^{\mathsf{B}}_h-\tau^{\mathsf{B}}_h[t^{\mathsf{A}}]\|
+\|\tau^{\mathsf{B}}_h[t^{\mathsf{A}}]-\tau^{\mathsf{A}}_h\|
\qquad(h\in H_{al}).
\end{equation*}
We sum over $h\in H_{al}$. The ramp comparison (Lemma~\ref{9.1:lem-ramp-comparison}) states
\begin{equation*}
\sum_{h\in H_{al}}\|\tau^{\mathsf{B}}_h-\tau^{\mathsf{B}}_h[t^{\mathsf{A}}]\|+S_{\mathrm{cn}}
\le C_{test}\vartheta_\chi^{\,k}A^{\sharp},
\end{equation*}
so the first sum is at most $C_{test}\vartheta_\chi^{\,k}A^{\sharp}-S_{\mathrm{cn}}$. The comparison
of aligned
terms at common ramps (Lemma~\ref{9.1:lem-aligned-terms}) states
\begin{equation*}
\sum_{h\in H_{al}}\|\tau^{\mathsf{B}}_h[t^{\mathsf{A}}]-\tau^{\mathsf{A}}_h\|
\le E_1^{-1}\Sigma^{q_a}_kA^{\sharp}.
\end{equation*}
The right-hand side bounds a sum of norms, so it is non-negative. It is therefore at most
$4E_1^{-1}\Sigma^{q_a}_kA^{\sharp}$. Together,
\begin{equation}\label{9.1:eq-last-lattice-comparison-2}
\sum_{h\in H_{al}}\|\tau^{\mathsf{B}}_h-\tau^{\mathsf{A}}_h\|
\le\bigl(C_{test}\vartheta_\chi^{\,k}+4E_1^{-1}\Sigma^{q_a}_k\bigr)A^{\sharp}-S_{\mathrm{cn}} .
\end{equation}

\emph{Unaligned labels.} For $\mathsf{m}\in\{\mathsf{A},\mathsf{B}\}$ the complement of the aligned
labels splits as the disjoint union
\begin{equation*}
H^{\mathsf{m}}\setminus H_{al}=H^{\mathsf{m}}_{\mathrm{old}}\sqcup H^{\mathsf{m}}_{\mathrm{cn}}.
\end{equation*}
A label of $\mathsf{B}$ whose first index is $-1$, the finest level private to $\mathsf{B}$, lies in
$H^{\mathsf{B}}_{\mathrm{old}}$. The bound for the old unpaired labels
(Lemma~\ref{9.1:lem-unpaired-labels}) states
\begin{equation*}
\sum_{h\in H^{\mathsf{A}}_{\mathrm{old}}}\|\tau^{\mathsf{A}}_h\|
+\sum_{h'\in H^{\mathsf{B}}_{\mathrm{old}}}\|\tau^{\mathsf{B}}_{h'}\|
\le4q_a^{\,k-n_c}A\le4q_a^{\,k-n_c}A^{\sharp}.
\end{equation*}
The sums over $H^{\mathsf{A}}_{\mathrm{cn}}$ and $H^{\mathsf{B}}_{\mathrm{cn}}$ together make up
$S_{\mathrm{cn}}$. Hence
\begin{equation}\label{9.1:eq-last-lattice-comparison-3}
\sum_{h\in H^{\mathsf{A}}\setminus H_{al}}\|\tau^{\mathsf{A}}_h\|
+\sum_{h'\in H^{\mathsf{B}}\setminus H_{al}}\|\tau^{\mathsf{B}}_{h'}\|
\le4q_a^{\,k-n_c}A^{\sharp}+S_{\mathrm{cn}} .
\end{equation}

\emph{Total.} Add \eqref{9.1:eq-last-lattice-comparison-2} and
\eqref{9.1:eq-last-lattice-comparison-3}. The term $S_{\mathrm{cn}}$ cancels, and the left-hand side
of
\eqref{9.1:eq-last-lattice-comparison-1} is at most
\begin{equation*}
\bigl(4E_1^{-1}\Sigma^{q_a}_k+4q_a^{\,k-n_c}+C_{test}\vartheta_\chi^{\,k}\bigr)A^{\sharp}
=\varepsilon_nA^{\sharp}.
\end{equation*}
This proves \eqref{9.1:eq-last-lattice-comparison-1} with $\varepsilon_n$ as defined in
\eqref{9.1:eq-last-lattice-comparison-a}.

\emph{The rate.} We bound the three summands of $\varepsilon_n$ separately.

First, the envelope of the loads (Lemma~\ref{9.1:lem-load-envelope}) gives the last of the bounds
\eqref{9.1:eq-load-envelope-b}, in its form with the constant $C_{env}^z$. It yields
\begin{equation*}
4E_1^{-1}\Sigma^{q_a}_k
\le\frac{4C_\Sigma^{TC}C_{env}^z\,\mathsf{l}_{\mathsf{w}}^2}{E_1}\,e_{\mathsf{w}},
\end{equation*}
and the same bound holds with $E_1^z$ in place of $E_1$ in both places. For labels with a live
second-class strip this bound rests on hypothesis~(e), that is, on
Proposition~\ref{9.1:prop-blind-site-expansions}, whose proof is outstanding.

Second, Lemma~\ref{9.1:lem-unpaired-labels} also states
\begin{equation*}
4q_a^{\,k-n_c}\le C_{nc}\rho^{\,n-\mathsf{w}}\le C_{nc}e_{\mathsf{w}}.
\end{equation*}

Third, $k=n-\mathsf{w}$, and $\vartheta_\chi\le\rho$ by hypothesis~(d), while $\vartheta_\chi>0$,
being a power of $L$. Moreover
$\rho^{\,n-\mathsf{w}}\le e_{\mathsf{w}}$, where $e_{\mathsf{w}}=\rho^{\,n-\mathsf{w}}+\Pi_n$.
Therefore
\begin{equation*}
C_{test}\vartheta_\chi^{\,k}=C_{test}\vartheta_\chi^{\,n-\mathsf{w}}
\le C_{test}\rho^{\,n-\mathsf{w}}\le C_{test}e_{\mathsf{w}}.
\end{equation*}

Adding the three bounds gives
\begin{equation*}
\varepsilon_n\le\Bigl(\frac{4C_\Sigma^{TC}C_{env}^z\,\mathsf{l}_{\mathsf{w}}^2}{E_1}
+C_{nc}+C_{test}\Bigr)e_{\mathsf{w}}=C_{TC}\,e_{\mathsf{w}}.
\end{equation*}
With the envelope bound in its form carrying $E_1^z$ in both places, the same addition, with the
second and third bounds unchanged, gives the bound with $E_1^z$ in place of $E_1$ in both places
stated after \eqref{9.1:eq-last-lattice-comparison-a}.

The constants $C_\Sigma^{TC}$, $C_{env}^z$, $E_1$, $E_1^z$, $C_{nc}$ and $C_{test}$, together with
$\mathsf{l}_{\mathsf{w}}$, are those of Lemmas~\ref{9.1:lem-aligned-terms},
\ref{9.1:lem-ramp-comparison}, \ref{9.1:lem-unpaired-labels} and~\ref{9.1:lem-load-envelope},
attributed after \eqref{9.1:eq-last-lattice-comparison-a}. None of them depends on $n$, $k$, $K$,
$m_0$ or the position, and $N_w$ enters only through $C_{test}$; this is the uniformity asserted
for $C_{TC}(g,N_w)$.
\end{proof}

\begin{remark}[Use of the bound; volume dependence and summability]
\emph{Use of the bound.} The bound \eqref{9.1:eq-last-lattice-comparison-1} is the comparison
hypothesis, read through the chart substitutions $\Psi^{\mathsf{m}}_h$ of the labels, at the last
lattice in the proof of the uniqueness theorem, with paired
label set $H_p=H_{al}$, pairing map the identity, and the majorant $A^{\sharp}$, which dominates the
constant $A$ of Lemma~\ref{9.1:lem-unpaired-labels}, in place of $A$. It is used there only through
the total masses of the complex measures on its left-hand side.

\emph{Volume dependence and summability.} The side $N_w$ of the last lattice enters $C_{TC}$ only
through the term of $C_{test}$ that carries $N_w^4$. The majorant $A^{\sharp}$ carries the factor
$e^{E_+^{\sharp\sharp}N_w^4}$. Thus the rate $e_{\mathsf{w}}$ in
\eqref{9.1:eq-last-lattice-comparison-a} is free of the volume, while the constant is not; this is
part~(d) of the corollary on the envelope functions. Part~(c) of the same corollary
gives
\begin{equation*}
\sum_n\bigl(\rho^{\,n-\mathsf{w}}+\Pi_n\bigr)<\infty\qquad\text{for }\kappa_0>6 ,
\end{equation*}
for the exponent $\kappa_0$ of that corollary.
Since $C_{TC}$ does not depend on $n$, \eqref{9.1:eq-last-lattice-comparison-a} then gives
$\sum_n\varepsilon_n<\infty$.
\end{remark}

\begin{corollary}[The normalised sourced integrals of the two runs]\label{9.1:cor-sourced-integrals}
Assume the hypotheses of Theorem~\ref{9.1:thm-last-lattice-comparison}. Let $\varepsilon_n$ and
$C_{TC}$ be as in \eqref{9.1:eq-last-lattice-comparison-a}. Let $A$, $A^{\sharp}$, $E_+^{\sharp}$,
$E_+^{\sharp\sharp}$, $N_w$ and $e_{\mathsf{w}}$ be the quantities of
Theorem~\ref{9.1:thm-last-lattice-comparison}. For every $\hat{w}\in\mathcal{N}'_w$ and every
$z\in P_{IR}$, with each run normalised at its own centre as in
\eqref{9.1:eq-centre-normalisation-a}, one has
\begin{equation}\label{9.1:eq-sourced-integrals-1}
\bigl|\tilde{Z}^{\mathsf{B}}(W^{\mathsf{B}}(z))-\tilde{Z}^{\mathsf{A}}(W^{\mathsf{A}}(z))\bigr|
\le \varepsilon_n^{\star}A ,
\end{equation}
and the same bound holds at the zero residual source. Here
$\varepsilon_n^{\star}:=\varepsilon_nA^{\sharp}/A$; it satisfies
\begin{equation}\label{9.1:eq-sourced-integrals-2}
\varepsilon_n^{\star}\le C_{TC}^{\star}(N_w;g)\,e_{\mathsf{w}},
\qquad
C_{TC}^{\star}:=e^{(E_+^{\sharp\sharp}-E_+^{\sharp})N_w^4}\,C_{TC}.
\end{equation}
The constant $C_{TC}^{\star}$ does not depend on $n$.
\end{corollary}

\begin{proof}
Fix $\hat{w}\in\mathcal{N}'_w$ and $z\in P_{IR}$. Since $P_{IR}\subset P_{UV}$,
Theorem~\ref{9.1:thm-last-lattice-comparison} applies at $z$, and it also applies at the zero
residual source.

For each run $\mathsf{c}\in\{\mathsf{A},\mathsf{B}\}$ we use the following objects, as in
Theorem~\ref{9.1:thm-last-lattice-comparison}:
\begin{itemize}
\item $\tilde{Z}^{\mathsf{c}}$ is the normalised sourced integral of the run, evaluated at the
residual source $W^{\mathsf{c}}(z)$;
\item for each label $h$ of the label set $H^{\mathsf{c}}$, the term $\tau_h^{\mathsf{c}}$ is a
complex measure on the configuration space $\mathcal{V}_w$ of the last lattice
$\mathbb{T}_w$;
\item $H_{al}$ is the set of aligned labels, and the pairing of labels is the identity.
\end{itemize}
The representation of the normalised sourced integrals at the last lattice as total masses of their
terms, split along the aligned labels, gives
\begin{equation}\label{9.1:eq-sourced-integrals-3}
\begin{split}
\tilde{Z}^{\mathsf{B}}-\tilde{Z}^{\mathsf{A}}
={}&\sum_{h\in H_{al}}\bigl[\tau_h^{\mathsf{B}}-\tau_h^{\mathsf{A}}\bigr](\mathcal{V}_w)
-\sum_{h\in H^{\mathsf{A}}\setminus H_{al}}\tau_h^{\mathsf{A}}(\mathcal{V}_w)\\
&+\sum_{h'\in H^{\mathsf{B}}\setminus H_{al}}\tau_{h'}^{\mathsf{B}}(\mathcal{V}_w).
\end{split}
\end{equation}

A complex measure $\mu$ on $\mathcal{V}_w$ satisfies $|\mu(\mathcal{V}_w)|\le\|\mu\|$, where
$\|\mu\|$ is its total variation. We apply this to each summand of
\eqref{9.1:eq-sourced-integrals-3} and use the triangle inequality. The result bounds
$|\tilde{Z}^{\mathsf{B}}-\tilde{Z}^{\mathsf{A}}|$ by the sum of the total variations on the
left-hand side of \eqref{9.1:eq-last-lattice-comparison-a}. By
Theorem~\ref{9.1:thm-last-lattice-comparison}, that sum is at most $\varepsilon_nA^{\sharp}$. With
$\varepsilon_n^{\star}=\varepsilon_nA^{\sharp}/A$ as in the statement, we obtain
\begin{equation*}
\bigl|\tilde{Z}^{\mathsf{B}}(W^{\mathsf{B}}(z))-\tilde{Z}^{\mathsf{A}}(W^{\mathsf{A}}(z))\bigr|
\le\varepsilon_nA^{\sharp}=\varepsilon_n^{\star}A .
\end{equation*}
This is \eqref{9.1:eq-sourced-integrals-1} at $z$.

The same argument applies at the zero residual source, for two reasons. First, the normalisation
\eqref{9.1:eq-centre-normalisation-a} uses the same objects at $z$ and at the zero source: each run
is taken at its own real centre, with its own vacuum responses subtracted there. Second, it rescales
nothing that depends on the label $h$ or on the source. Hence the decomposition
\eqref{9.1:eq-sourced-integrals-3}, the bound by total variations and the bound of
Theorem~\ref{9.1:thm-last-lattice-comparison} hold there with the same constants.

It remains to prove \eqref{9.1:eq-sourced-integrals-2}. Theorem~\ref{9.1:thm-last-lattice-comparison}
bounds the sum of total variations by $\varepsilon_nA^{\sharp}$, with the majorant $A^{\sharp}$,
which is independent of $n$ and satisfies $A^{\sharp}\ge A$, in place of $A$; the ratio
$A^{\sharp}/A$ is absorbed into the constant, which becomes $C_{TC}^{\star}$. By the definitions
$A=e^{E_+^{\sharp}N_w^4}$ and $A^{\sharp}=e^{E_+^{\sharp\sharp}N_w^4}$ of
Theorem~\ref{9.1:thm-last-lattice-comparison},
\begin{equation*}
\frac{A^{\sharp}}{A}=e^{(E_+^{\sharp\sharp}-E_+^{\sharp})N_w^4}.
\end{equation*}
Together with $\varepsilon_n\le C_{TC}e_{\mathsf{w}}$ from
\eqref{9.1:eq-last-lattice-comparison-a}, this gives
\begin{equation*}
\varepsilon_n^{\star}=\frac{A^{\sharp}}{A}\,\varepsilon_n
\le e^{(E_+^{\sharp\sharp}-E_+^{\sharp})N_w^4}\,C_{TC}\,e_{\mathsf{w}}
=C_{TC}^{\star}(N_w;g)\,e_{\mathsf{w}} .
\end{equation*}
The factor $A^{\sharp}/A$ does not depend on $n$, and neither does $C_{TC}$. Hence $C_{TC}^{\star}$
does not depend on $n$.
\end{proof}

In the statements of this chapter that apply Corollary~\ref{9.1:cor-sourced-integrals} to compare
the increments of the two runs, the constant written $C_{TC}(N_w;g)$ is $C_{TC}^{\star}$ of
\eqref{9.1:eq-sourced-integrals-2} (the undecorated $C_{TC}$ of this corollary is the constant of
\eqref{9.1:eq-last-lattice-comparison-a}; the two differ by the factor $A^{\sharp}/A$). In
particular, the constant of those statements is
\begin{equation*}
C_*(N_w):=2C_{TC}^{\star}\,\mathsf{Q}_w(1+\mathsf{Q}_w),
\end{equation*}
where $\mathsf{Q}_w=AB$ is the product of the volume constants $A$ and $B$ of the setting of
clause~(v).

\begin{remark}[Centre normalisation, length floor and gauge-weight pieces]
\label{9.1:rem-centre-normalisation-b}
Let the runs $\mathsf{A}$, $\mathsf{B}$, the stop depth $\mathsf{w}=\mathsf{w}(g)$ and the stop level
$k=n-\mathsf{w}$ be as in Notation~\ref{9.1:not-centre-normalisation}. For
$\mathsf{c}\in\{\mathsf{A},\mathsf{B}\}$ write $c^{\mathsf{c}}$ for the real centre of the run
$\mathsf{c}$ fixed there. Three features of the normalisation
\eqref{9.1:eq-centre-normalisation-a} and of the floor \eqref{9.1:eq-centre-normalisation-e}
enter the two-cutoff comparison of terms at the last lattice
(Theorem~\ref{9.1:thm-last-lattice-comparison}).
\begin{enumerate}
\item[(1)] \emph{Centre normalisation.} In \eqref{9.1:eq-centre-normalisation-a} each
$\tilde{Z}^{\mathsf{c}}$ is the representation given by the correlation theorem at the run's own
real centre $c^{\mathsf{c}}$, with $\mathcal{E}^{\mathsf{c}}$ subtracted at $c^{\mathsf{c}}$.
Inside $\tau_h^{\mathsf{c}}$ the compensating factors and the vacuum constants that do not depend
on $z$ lie in the dial domain $\mathbb{D}^{\star}$ and are counted in the total load
$[\![\delta]\!]^{\natural+}$. The normalisation admits exactly one free positive constant for each
run and each centre.
\item[(2)] \emph{Length floor.} The floor $n\ge n_0(N_w)$ in
\eqref{9.1:eq-centre-normalisation-e}, of which $n\ge n_2(g)$ is a part, is a hypothesis of the
statement in which Theorem~\ref{9.1:thm-last-lattice-comparison} is applied, not of that theorem.
Theorem~\ref{9.1:thm-last-lattice-comparison} requires $n\ge n_1$, $k\ge 1$ and the common
scaffold $\Theta_X$, and its constant $C_{TC}$ does not depend on $n_2$.
\item[(3)] \emph{Gauge-weight pieces.} The terms $\log[\mathfrak{z}(a_j(y))/\mathfrak{z}(x_j)]$
occur in $\mathcal{E}$ on every cube $y$. With the opposite sign they occur inside
$\tau_h^{\mathsf{c}}$ on the cubes outside $\Lambda_j(h)$. There, moving along the dial changes
their exponent on each cube by at most $CE_1$, which is one unit of the activity constant
$\mathsf{c}^{\star}$.
\end{enumerate}
No hypothesis beyond those of Theorem~\ref{9.1:thm-last-lattice-comparison} enters.
\end{remark}

\begin{proof}
Here $\mathfrak{z}$ is the one-point gauge weight, $x_j$ the inverse coupling at level $j$,
$\zeta_j$ the running shift at level $j$, $\varepsilon_j$ the small-field threshold at level $j$,
and $\Lambda_j(h)$ the small-field region of the label
$h$ at level $j$. The letter $\lambda\in[0,1]$ denotes the interpolation parameter of the dial,
that is, of the convex combination $x^{\lambda}$ of slot arrays; it is not the constant $c_0/2$ of
the ladder. The symbol $\delta\zeta_j$ denotes the difference of $\zeta_j$ between
the two runs.

\emph{Item \textup{(1)}.} By \eqref{9.1:eq-centre-normalisation-a} and the normalisation
statement for the sourced integrals of the two runs, each $\tilde{Z}^{\mathsf{c}}$
is the correlation theorem's representation at $c^{\mathsf{c}}$, and $\mathcal{E}^{\mathsf{c}}$ is
subtracted at the same point $c^{\mathsf{c}}$.

A subtraction made once for each length, rather than once for each run and centre, would not do.
It would put $\tilde{Z}^{\mathsf{c}}(0)$ outside the interval $[1/B,A]$, and the two-cutoff
comparison between the two lengths would then fail.

Inside $\tau_h^{\mathsf{c}}$ the subtraction is compensated by the factors
\begin{equation*}
-\sum_{Y\not\subset\Lambda_j(h)}
\bigl[e^{(j),\mathsf{c}}(Y;c^{\mathsf{c}}+v)-e^{(j),\mathsf{c}}(Y;c^{\mathsf{c}})\bigr],
\end{equation*}
where $e^{(j),\mathsf{c}}(Y;\cdot)$ is the level-$j$ boundary exponent of the run $\mathsf{c}$ on
the region $Y$, and $v$ is the displacement of its source argument from the real centre
$c^{\mathsf{c}}$.
\begin{itemize}
\item These factors are the values of the slice at the data value $\mathsf{v}=1$ of stored local
families of class $(j,k)$, namely the families of boundary exponents in the table of stored local
families.
\item The vacuum constants that do not depend on $z$ occupy slots of the entries $\mathsf{z}$. This
is the first clause of the slot assignment for vacuum constants.
\end{itemize}
Hence both kinds of term lie inside the dial domain $\mathbb{D}^{\star}$ and are counted in
$[\![\delta]\!]^{\natural+}$.

The three normalisation conditions on the representation leave exactly one free positive constant
for each run and each centre. This proves item~(1).

\emph{Item \textup{(2)}.} The floor $n\ge n_0(N_w)$ of \eqref{9.1:eq-centre-normalisation-e}
contains $n\ge n_2(g)$ among its parts. It is a hypothesis of the statement that consumes
Theorem~\ref{9.1:thm-last-lattice-comparison} at the last lattice, where it serves the second
clause of that statement's transfer bound.

Of the length, Theorem~\ref{9.1:thm-last-lattice-comparison} asks only $n\ge n_1$ and
$k=n-\mathsf{w}\ge1$, together with the reading of both runs on the common scaffold $\Theta_X$; its
remaining hypotheses do not involve $n_2$.
On the common scaffold $\Theta_X$ the labels whose first index exceeds the common-scaffold level
$n_c$ are the same sets in both runs. Here $n_c\ge k_*$, where $k_*$ is the level parameter of the
class $\mathfrak{S}$ of tuned runs to which both runs belong, and above the level $k_*$ the label
sets of two runs of that class coincide.

The number $n_2$ enters none of these requirements, and the constant $C_{TC}$ does not depend on
$n_2$. This proves item~(2).

\emph{Item \textup{(3)}.} The term $\log[\mathfrak{z}(a_j(y))/\mathfrak{z}(x_j)]$ lies in
$\mathcal{E}$ on every cube $y$:
\begin{itemize}
\item at depths $r\le n/2$ by part (b) of the localisation corollary;
\item at larger depths by the ultraviolet bound at the last lattice.
\end{itemize}
Neither of these is part of the proof of Theorem~\ref{9.1:thm-last-lattice-comparison}.

With the opposite sign the same term lies inside $\tau_h^{\mathsf{c}}$ on the cubes outside
$\Lambda_j(h)$. This follows from the torus bound of the correlation theorem, in its form at the
last lattice, together with the clause of that bound which places the gauge-weight factor, with the
opposite sign, on the cubes outside $\Lambda_j(h)$.

On those cubes the argument $a_j(y)=x_j-\zeta_j(p_y)$, where $p_y$ is the plaquette of the cube $y$
at which the running shift is read, is a weight slot of class $(j,k)$. It is
the gauge-fixing entry of the table of stored local families. Its difference between the two runs
is the part of the load entry $\mathsf{z}_j$ governed by the matching estimate.

The dial therefore moves the exponent on each cube by at most
\begin{equation}\label{9.1:eq-centre-normalisation-b-1}
\begin{aligned}
|\lambda\,\delta\zeta_j|\cdot\sup|(\log\mathfrak{z})'|
&\le \frac{E_1}{\mathfrak{g}_j}\cdot C\varepsilon_j^2\\
&\le CE_1A_0^2(\log x_j)^{2p_0-2q_\flat}\\
&\le CE_1 .
\end{aligned}
\end{equation}
The three inequalities are justified as follows.
\begin{itemize}
\item \emph{First inequality.} The factor $|\lambda\,\delta\zeta_j|$ is at most
$E_1/\mathfrak{g}_j$. The factor $\sup|(\log\mathfrak{z})'|$ is at most $C\varepsilon_j^2$, by the
gauge-fixing clause of the correlation theorem's lemma on the small-field factor. That clause gives
$1-\operatorname{Re}\operatorname{tr}U(\partial p)\le C\varepsilon_j^2$ for the plaquettes $p$ of
the cube, on the support of $\chi_j$.
\item \emph{Second inequality.} This is the last clause of the hypothesis of
Theorem~\ref{9.1:thm-last-lattice-comparison} that is stated with the constant $2E_1$.
\item \emph{Third inequality.} Since $q_\flat\ge p_0$, the exponent $2p_0-2q_\flat$ is
non-positive, and $x_j\ge X$, where $X=g^{-2}$; hence
\begin{equation*}
A_0^2(\log x_j)^{2p_0-2q_\flat}\le A_0^2(\log X)^{-2(q_\flat-p_0)}\le 1
\end{equation*}
by the ceiling $A_0^2(\log X)^{-2(q_\flat-p_0)}\le 1$. This ceiling is a condition on $\gamma_U$
alone once $A_0$ is fixed, and $A_0$ is chosen before $\gamma_U$; it therefore holds under the
hypothesis $g\le\gamma_U$ of Theorem~\ref{9.1:thm-last-lattice-comparison}.
\end{itemize}

The first inequality needs the bound $|(\log\mathfrak{z})'|\le C\varepsilon_j^2$ of the
gauge-fixing clause. The bound $|(\log\mathfrak{z})'|\le2$ of the localisation corollary would
not suffice, since the factor $\varepsilon_j^2$ is what the second inequality consumes.

Thus the dial changes the exponent of each gauge-weight piece by at most $CE_1$ per cube, which is
one unit of $\mathsf{c}^{\star}$. This proves item~(3).
\end{proof}

\begin{definition}[Units, radii and degrees on the core]\label{9.1:not-core-units}
In this definition the subscripts $n$, $n'$, $k$ are levels. The following letters are
used with the meanings stated.

\emph{Coefficients.} We put $\beta=1/5$ and
\begin{equation*}
\varphi_{\min}:=1-4\beta=\tfrac15 .
\end{equation*}
Every coefficient of \cite[(1.64)(i)]{B-CMP122a} is at least $1-4\beta$; this is the content
of a remark of this book on those coefficients. On a live
piece on which the arrest
theorem has acted (an \emph{arrested} piece), by the arrest estimate for such pieces,
$\varphi_{\min}$ is replaced by $\varphi_{\mathrm{arrest}}:=1-4.51\beta$.

\emph{Units.} Put $\ell_n:=L^n\eta$, with $\eta$ the factor common to all levels;
$\ell_k$ is the unit of level $k$, and we set $\ell_k=1$ where convenient.
For a table member $a$ (one of the members $E_a$ of the stage tables recalled in the final
paragraph below), with radius $\alpha_{a,n}$ at level $n$, the unit of $a$ at level $n$ is
\begin{equation*}
\mathsf{u}^a_n:=\alpha_{a,n}\,\ell_n^{-e_a},\qquad e_6=1,\quad e_1=e_5=e_7=2,\quad e_3=3,
\end{equation*}
where the powers of $\eta$ common to all levels are suppressed.

\emph{Radii and degrees.} With $x_n=g_n^{-2}$, the constants $A_0,A_1,C_0,C_1$ and the
exponents $p_0,p_1$ among Ba{\l}aban's auxiliary parameters $\mathfrak{p}$, we put
\begin{equation*}
\varepsilon_n=A_0\,g_n(\log x_n)^{p_0},\qquad
\delta'_n=A_1\,g_n(\log x_n)^{p_1},\qquad p_1<p_0,
\end{equation*}
the last inequality as in \cite{B-CMP122a}, and, as in \cite[(2.28)]{B-CMP119},
\begin{equation*}
\alpha_{i,n}=C_i\,g_n(\log x_n)^{q},\qquad i=0,1.
\end{equation*}
The radii and the common exponent $q$ satisfy
\begin{equation}\label{9.1:eq-core-units-a}
\begin{gathered}
\alpha_{0,n}\le\alpha_{1,n},\qquad
\alpha_{a,n'}\le\tfrac94\,(n'-n)^{1/2}\,\alpha_{a,n}\ \ (n'>n),\qquad
\alpha_{a,n}\le2\,\alpha_{a,n'},\\
q:=q_0=q_1\ \ge\ p_0+\bigl(1+d+d\beta_0^{pr}\bigr)r_{pr}+1 ,
\end{gathered}
\end{equation}
where $q_0,q_1$ are the exponents of $\alpha_{0,n},\alpha_{1,n}$, $r_{pr}$ is the exponent
of the radius $R(x)$ below, and $\beta_0^{pr}$ is Ba{\l}aban's small growth number.

\emph{The radius $R(x)$ and depths.} For $x\ge1$ let $R(x)$ be the least power $L^j$, $j$ a
non-negative integer, with $L^j\ge(\log x)^{r_{pr}}$, and put $R_n:=R(x_n)$; the sequence
$(R_n)$ is non-increasing in $n$ on the levels $n$ with $x_n\ge1$.
We take $N=R_k$, the value prescribed for $N$ by the construction of the core in this
chapter, and
put $h:=k-N$ and $k_0:=k-N_0$, with $N_0$ the integer so denoted in
\cite{B-CMP122a}. Wherever this definition is used, $N$ denotes the integer $R_k$, $h$ the
level $k-N$ and $R$ the function just defined; neither the integer $N$ of the gauge group
$\mathrm{SU}(N)$, nor a history $h$, nor the radius $R$ of the balls supporting the observables
is meant by these letters there.

\emph{Seed size and one constant of Ba{\l}aban.} With $\mathsf{r}_1:=C_1/C_0$ we put
\begin{equation*}
\tilde\alpha:=\alpha_{0,k}+\alpha_{1,k}=(1+\mathsf{r}_1)\,\alpha_{0,k},\qquad
\mathsf{s}_L:=K_\Lambda'(M)\,\tilde\alpha ,
\end{equation*}
where $K_\Lambda'(M)$ is the frame constant of the statement of this chapter that constructs
the seed and bounds its size by $\mathsf{s}_L$. Finally, $K_0$ denotes the
constant $O(1)B_3^2B_5M^6$ of \cite[(1.87)]{B-CMP122a}, with $B_3,B_5$ the constants so named
there.

\emph{Further letters.} The following letters are not defined in this definition; each keeps
the meaning fixed at its introduction elsewhere in this chapter: the local Lipschitz property
$\natural\mathrm{Lip}^{\mathrm{loc}}$, the slice $S(\cdot)$, the tubes $\mathfrak{T}[\varsigma]$,
the domains $\mathfrak{D}_F[\varsigma]$, the table members $E_a$ with their point classes $(=)$
and $(\mathrm{up})$, the stage tables $\mathsf{T}_\nu$ and their room increments
$\mathrm{rm}_\nu$; the layer profile $\mathsf{m}_L$, the constant $C_4$ and the table of the
theorem in which these two are fixed; the constant $c_1$ bounding the sum of the decay weights,
the constant
$C_T=c_1L^3$, and the lower bound on distances used with that sum; and the remaining letters
of the layer construction. Among these are the constant $B_\natural$ of the local Lipschitz
property, the dilation factor $B^{\mathsf{s}}$, and a further constant $C_s$. Each
of $B_\natural$, $c_1$, $C_s$, $B^{\mathsf{s}}$, $C_T$, $C_4$, $K_\Lambda'$ depends on
$d$, $L$, $M_1$, $M$, $\kappa_1$ and the number $\theta$ fixed with these letters only.
\end{definition}

\begin{proof}
The value of $\varphi_{\min}$ is arithmetic: $1-4\cdot\frac15=\frac15$.

The three inequalities in the first line of \eqref{9.1:eq-core-units-a} are
\cite[(2.6)--(2.8)]{B-CMP119}, read for the radii $\alpha_{a,n}$. In the second line the two
exponents of $\alpha_{0,n}$ and $\alpha_{1,n}$ are one and the same number $q$; the lower bound
on $q$ is not imposed in this definition: it is the condition on $q$ under which two
statements of this book, on which the estimates of this chapter depend, are stated, and it is
recorded here from them.

For the identity defining $\tilde\alpha$: since $\alpha_{0,k}$ and $\alpha_{1,k}$ carry the same
factor $g_k(\log x_k)^{q}$, the definition of $\alpha_{i,k}$ gives
\begin{equation*}
\alpha_{1,k}=C_1\,g_k(\log x_k)^q=\frac{C_1}{C_0}\,C_0\,g_k(\log x_k)^q
=\mathsf{r}_1\,\alpha_{0,k},
\end{equation*}
and therefore $\alpha_{0,k}+\alpha_{1,k}=(1+\mathsf{r}_1)\alpha_{0,k}$.

For the monotonicity of $(R_n)$: by \cite[(2.6)]{B-CMP119} the couplings $g_n$ increase with
$n$, so $x_n=g_n^{-2}$ and $\log x_n$ do not increase with $n$. Consider the levels $n$ with
$x_n\ge1$, on which $R_n$ is defined; the exponent $r_{pr}$ is non-negative (it is fixed in the
order of choice of the constants, Notations~\ref{2.3:not-stages-first},
\ref{2.3:not-stages-middle} and~\ref{2.3:not-stages-last}), so the quantity
$(\log x_n)^{r_{pr}}$ does not increase with
$n$ on these levels. The least power of
$L$ that is at least a given number is a non-decreasing function of that number, so
$R_n=R(x_n)$ does not increase with $n$ on these levels.
\end{proof}

\begin{definition}[The core region, its collar and its deep part]\label{9.1:def-core-collar}
\emph{Setting.} Consider a site of Ba{\l}aban's operation $\mathbf{R}$ at step $k$, on a connected
component $Z$ of the large-field region that satisfies the two conditions of \cite[\S1]{B-CMP122a}:
(i) $Z$ is contained in a cube of size $100M\check{R}_k$; (ii) in the preceding $\bar N$
renormalization steps no new large-field regions were created inside $Z$, and the previous regions
contained in $Z$ satisfy condition (i) on the corresponding scales. Here $\check{R}_k$ is
Ba{\l}aban's radius at level $k$ (written $R_k$ in \cite{B-CMP122a}), and $\bar N$ is the number of
steps written $N$ in \cite{B-CMP122a}. The large-field regions inside $Z$ form the chain
\begin{equation*}
Z=Z_k\supset Z_{k-1}\supset\cdots\supset Z_{h+1}\supset Z_h\supset\cdots,
\end{equation*}
the last dots standing for the older regions, if any. In this definition the letters $h$, $j$ and $m$
index levels of Ba{\l}aban's construction; $h$ is the level of \cite[(1.10)--(1.12)]{B-CMP122a} at
and below which the regions $Z_j$ and the domains $\Omega_j$ are left unmodified, as recalled below.
With $\Omega_m$ the small-field domains of
\cite{B-CMP122a}, the zone $(\Omega_m\setminus\Omega_{m+1})\cap Z$ carries label level $m$. In the
notation of \cite[(1.10)--(1.12)]{B-CMP122a}, with $N_0$ as in \cite[(1.10)]{B-CMP122a},
$Z''_j=(\Omega^{\sim5}_j)^c\cap Z$ for
$j=k-N_0,\dots,h+1$, $Z''_j=Z_j$ and $\Omega''_j=\Omega_j$ for $j=h,\dots,1$, and all components of
the domains $Z_j$ are rectangular parallelepipeds for $j=h,\dots,k$. For $j=k-N_0,\dots,h+1$ the
domains $\Omega''_j$ are the modified small-field domains of \cite[(1.10)--(1.12)]{B-CMP122a}. For
a small-field domain, the superscript $\sim$ denotes the domain changed by one boundary layer in the
sense of \cite{B-CMP122a}, and $\sim2$, $\sim4$, $\sim5$ denote the domain changed by two, four and
five such layers. The label
level $m(x)$ of a
point $x$, in the convention of this chapter, is $m(x)=j$ on $Z''_{j+1}\setminus Z''_j$, and
$h\le m(x)\le k-1$ on $Z''_k$.

The set $\Lambda=(\Omega^{\sim4}_{k_0+1})^c\cap Z$, with $k_0$ the level fixing the domain
$\Omega_{k_0+1}$ from which $\Lambda$ is cut in \cite[\S1]{B-CMP122a}, is contained in a cube of
size $100M$ \cite[\S1]{B-CMP122a}, and
\begin{equation}\label{9.1:eq-core-collar-1}
\Lambda_0:=\Lambda\cap\Omega''^{\sim2}_{h+1}
\end{equation}
carries the variables of the step $\mathbf{R}$: by \cite[(1.81)]{B-CMP122a}, $V''=V'V_0$ on
$\Lambda_0$. Here $U_0$ are the bond variables of the step $\mathbf{R}$ and $V_0$ their average,
both as in \cite[(1.81)]{B-CMP122a}, the average being taken on the bond set
$\mathfrak{B}_0=\mathfrak{B}''_k\cap\Lambda_0$, with $\mathfrak{B}''_k$ the bond set of that
display; bonds intersecting $\partial\Lambda$ belong to $\mathfrak{B}_0$, and bonds
intersecting the boundary of the modified small-field domain named in \cite[(1.81)]{B-CMP122a}
do not. By the lemma of this chapter on the prepared data of the step $\mathbf{R}$ (its item on
the bond set), the boundary meant in
\cite[(1.81)]{B-CMP122a} is the surface
\begin{equation}\label{9.1:eq-core-collar-2}
S_c:=\partial\,\Omega''^{\sim2}_{h+1}.
\end{equation}
After the $\mathbf{T}$-step, $V'=e^{ig_kB}$, with $B$ the variable of the $\mathbf{T}$-step, in the
form stated in \cite[(1.50)]{B-CMP122b}, where the configuration $U''_k$ is $U'_k$ evaluated at
$V''$, at $\exp(ig_kB)$ times the average of $U_0(\mathbf{U})$, with $\mathbf{U}$ the configuration
of that display, and at $V''$, on the three sets of \cite[(1.50)]{B-CMP122b}, the second being a
part of $\Lambda\cap\Omega''$;
and on the support of the characteristic function $\chi'$ of \cite[(1.82)]{B-CMP122a}, at real $B$,
one has $|g_kB|\le C_1\delta'_k$, with $\delta'_k$ the constant of that display. The cut of
\cite[(1.52)--(1.53)]{B-CMP122b} breaks the configuration $U''_k$ into two independent parts by
boundary conditions at the boundary $\partial\Omega''^{\sim}_{h+1}$, the \emph{seam}. By the lemma
of this chapter on the cut at the seam, the seam lies one $\sim$-layer, of width $M\ell_h$, outside
$S_c$, with $\ell_h$ the length scale at level $h$, and the set
$(\Omega''^{\sim}_{h+1})^c\cap\Omega''^{\sim2}_{h+1}$ has width $M\ell_h$ and is of level $h$ for
the class $\mathbf{B}_2$ of determining sets of that lemma.

\emph{Definition.} The \emph{core} $\mathfrak{K}$, its \emph{collar} $\mathfrak{K}_1$ and its
\emph{deep part} $\mathfrak{K}_2$ are
\begin{equation}\label{9.1:eq-core-collar-a}
\begin{gathered}
\mathfrak{K}:=(\Omega''^{\sim}_{h+1})^c\cap Z=\mathfrak{K}_1\cup\mathfrak{K}_2,\\
\mathfrak{K}_1:=\mathfrak{K}\cap\Omega''^{\sim2}_{h+1},\qquad
\mathfrak{K}_2:=(\Omega''^{\sim2}_{h+1})^c\cap Z.
\end{gathered}
\end{equation}
The collar $\mathfrak{K}_1$ has label level $h$; its cells are $\mathfrak{B}_0$-cells, that is, cells
whose bonds belong to $\mathfrak{B}_0$, and its data
are $e^{ig_kB}V_0$. On the deep part $\mathfrak{K}_2$ the data are the real configuration $V_h$,
common to the two runs $\mathsf{A}$ and $\mathsf{B}$ compared in this chapter, on its zone-$h$ part
$\mathfrak{K}_2\cap\Omega_h$,
and the older common configurations $V_m$, $m<h$, on the zones of $Z_h$ itself.

\emph{The guard.} By \cite[(1.88)]{B-CMP122a}, in the form used in the lemma of this chapter on the
prepared data of the step $\mathbf{R}$ (its item on the guard function), on $Z$ after $\mathbf{R}$
one has the function
$\chi_{h,1/2}\bigl((\Omega''^{\sim}_{h+1})^c\cap\Omega_h\bigr)$. It is a product over the
$LM_2\check{R}_h$-cubes $\square$ for the $L^{-h}$-scale that intersect the domain
$(\Omega''^{\sim}_{h+1})^c$ and are contained in $\Omega_h$. Each factor is the characteristic
function of a condition of the type
\begin{equation}\label{9.1:eq-core-collar-3}
\bigl|U_{h,\square}\bigl((1,V_h),\partial p\bigr)-1\bigr|
<\tfrac12(\cdots)\,\varepsilon_h\,(L^{k-h}\eta)^2,
\end{equation}
the omitted factor being as in \cite[(1.88)]{B-CMP122a}, with $\varepsilon_h$ and $\eta$ the
small-field parameters of that display, for the local function $U_{h,\square}$ whose
data are $1$ on the cells on the side of $\Omega''^{\sim2}_{h+1}$ and $V_h$ inside. The symbol
$(1,V_h)$ means the configuration
$(1{\upharpoonright}_{\Omega''^{\sim2}_{h+1}},\,V_h{\upharpoonright}_{(\Omega''^{\sim2}_{h+1})^c})$.
By the flat scaffold comparison of this chapter, the group-valued factor $U_{h,\square}(1,V_h)$ does
not depend on the seed variable and has curvature less than
$\tfrac12\varepsilon_h(L^{k-h}\eta)^2$.

\emph{The frame hypothesis on the core.} The uniqueness theorem of this chapter is stated under the
following hypothesis, referred to as the frame hypothesis on the core: the restriction
$V''{\upharpoonright}_{\mathfrak{K}}$, which consists of $V_h$ and the older $V_j$ on the deep part,
and
the variable $B_0$ before passage to the axial-Landau frame are common integration parameters of
the two runs $\mathsf{A}$ and $\mathsf{B}$, with the $\varepsilon$-regularity condition of the core
holding at their sites.

\emph{Sizes.} Since $\mathfrak{K}\subset Z_{h+1}$, and since by conditions (i) and (ii) $Z_{h+1}$ is
contained in a cube of $100M\check{R}_{h+1}$ level-$(h+1)$ units, in the units of
Notation~\ref{9.1:not-core-units} $\mathfrak{K}$ is at most
\begin{equation}\label{9.1:eq-core-collar-4}
K_h:=100\sqrt{d}\,ML\check{R}_{h+1}\le 100\sqrt{d}\,ML\check{R}_h
\end{equation}
level-$h$ units across. A deeper zone $m<h$ of $\mathfrak{K}$ lies in $Z_{m+1}$, which by the second
part of condition (ii) (the previous regions satisfy condition (i) on the corresponding scales) is
contained in a cube of
\begin{equation}\label{9.1:eq-core-collar-5}
K_m:=100\sqrt{d}\,ML\check{R}_{m+1}
\end{equation}
level-$m$ units. In level-$k$ units the core $\mathfrak{K}$ has diameter at most $K_hL^{-\bar N}$,
and $K_hL^{-\bar N}<1$, indeed much smaller than $1$, by the diameter bound for the component $Z$.
Finally, every path from a cell of
zone $m'$ to a point of zone $m$ crosses the $|m-m'|$ zone boundaries lying between them.
\end{definition}

\begin{definition}[The objects read on the core and their data]\label{9.1:def-core-data}
Fix a step $k$, an $\mathbf{R}$-site of step $k$ on a component $Z$, and the level $h<k$,
the domain $\Lambda$, the bond set $\mathfrak{B}_0$ and the core
$\mathfrak{K}=\mathfrak{K}_1\cup\mathfrak{K}_2$ of
Definition~\ref{9.1:def-core-collar}, see \eqref{9.1:eq-core-collar-a}.
Let $\zeta$ denote the seed variable and $B$ the variable of the $\mathbf{T}$-step. Let
$U''_k(\zeta,B)$ be the configuration of \cite[(1.50)]{B-CMP122b}, let
$U^0_k(\zeta)$ be the same configuration at $B=0$, and put
\begin{equation*}
\tilde{\alpha}:=\alpha_{0,k}+\alpha_{1,k}=(1+\mathsf{r}_1)\,\alpha_{0,k},
\qquad \mathsf{r}_1=C_1/C_0 .
\end{equation*}
\begin{itemize}
\item[(a)] \emph{The reading set.} A term $(F,X,j)$ of the expansion, $X$ its localisation
region, with $X\cap Z\neq\emptyset$ is a \emph{first-part term}. The objects at which
first-part terms are read form the set
\begin{equation*}
\mathfrak{Q}'':=\bigl\{U^0_k(\zeta),\ U''_k(\zeta,B),\ \Xi_1(\sigma')\bigr\}
\cup\{\text{box and bracket vertices over these}\},
\end{equation*}
where $\Xi_1(\sigma')$ is the minimiser on
$\mathbf{B}^{\sim}(Y_1):=\mathbf{B}(Y_1)\cup\mathbf{B}''_k$ described in item (c), with
$Y_1$ the region of \cite[(1.59)--(1.61)]{B-CMP122b}.
\item[(b)] \emph{The data of the minimisers.} Each
$P\in\{U^0_k(\zeta),U''_k(\zeta,B)\}$ is the minimiser $U_{\mathbf{B}''_k}$ with the
following data. Off $\Lambda$ the data are $\mathsf{v}$, the seed's own data. On
$\mathfrak{B}_0$ (the rind, whose bonds are of level $k$; the towers, whose bonds $c$ have a
level $\pi=\pi(c)$; and the collar $\mathfrak{K}_1$) the data are
$e^{ig_kB}V_0^{\mathrm{fr}}(\zeta)$, with $B=0$ for $U^0_k$, where
\begin{equation*}
V_0(\zeta)=M_{\mathbf{B}''_k}\bigl(U_0(\zeta)\bigr),\qquad
U_0(\zeta)=U_{k,Z}\bigl(\mathsf{v}{\upharpoonright}(Z\setminus\Lambda)\,\big|\,V_\Lambda\bigr),
\end{equation*}
with $V_\Lambda$ the data on $\Lambda$ of \cite[(1.79), (1.81)]{B-CMP122a}.
For a frame $\mathrm{fr}$ we write $U_0^{\mathrm{fr}}(\zeta)$ for $U_0(\zeta)$ in the gauge
$\mathrm{fr}$, and $V_0^{\mathrm{fr}}(\zeta):=M_{\mathbf{B}''_k}\bigl(U_0^{\mathrm{fr}}(\zeta)\bigr)$.
On the deep core $\mathfrak{K}_2$ the data are the common real variables $V_h$ and $V_m$,
$m<h$.
\item[(c)] \emph{The minimiser $\Xi_1(\sigma')$.} It is the minimiser on
$\mathbf{B}^{\sim}(Y_1)$ with the same interior data as in (b). On the layer cells at
$\partial Y_1^0$ of \cite[(1.59)--(1.61)]{B-CMP122b}, its data are the averages localised
with the localisation parameter $\sigma'$, read on those cells.
\item[(d)] \emph{The frame.} Throughout we take $\mathrm{fr}:=\mathrm{AL}$, the
axial--Landau frame. The class $\mathfrak{F}_G$ consists of the gauge frames whose gauge
transformation $\gamma'$ equals $\bar u_0$ on $\partial^{+}$; here $\bar u_0$ is the
boundary gauge factor, and $\partial^{+}$ is the boundary set to which the compensating
gauge transformations recalled in the proof below are restricted.
\end{itemize}
With these conventions the data have the following raw sizes, for every $\zeta$ in the
complex tube of seeds:
\begin{equation}\label{9.1:eq-core-data-a}
\begin{aligned}
&\text{(i)}\quad \bigl|V_0^{\mathrm{AL}}(\zeta)(c)-1\bigr|\le 34\,d\,K_0\,\tilde{\alpha}\,L^{\pi-k},\\
&\text{(ii)}\quad \bigl|\log V_0^{\mathrm{AL}}(\zeta)(c)\bigr|\le 2K_\Lambda^{\flat}\,\tilde{\alpha}.
\end{aligned}
\end{equation}
Bound (i) holds for every tower bond $c$ of level $\pi=\pi(c)$ (both of whose level-$\pi$
blocks lie in the domain $Z''_k$ of \cite[(1.10)--(1.12)]{B-CMP122a}), the bonds of the
collar, for which $\pi=h$, included. Bound (ii)
holds for the level-$k$ bonds of the rind, for those of the level-$k$ bond set
$\mathfrak{i}$ of the flat-frame data bound (on these bonds that bound's datum is
$V_f\bar u_0^{-1}$; see the proof), and for those crossing $\partial\Lambda$. The constants
in (ii) are
\begin{equation*}
K^{\flat}:=\max\{c_QK_0,\ 1.1\,(K_f+K_uK_0)\},\qquad
K_\Lambda^{\flat}:=\max\{K'_\Lambda,\ K^{\flat}\}.
\end{equation*}
Here $K_0$ is the constant with which \cite[(1.87)]{B-CMP122a}, with $\tilde{\alpha}$ in
place of $\varepsilon_k$, reads $\ell_k|A_0(\zeta)|\le K_0\tilde{\alpha}$ on $Z''_k$ (see
\eqref{9.1:eq-core-data-1} below), with $A_0$ the vector potential of
\cite[(1.86)--(1.87)]{B-CMP122a} (not the auxiliary parameter $A_0$ of the glossary);
$K'_\Lambda=K'_\Lambda(M)$ is the constant of the seed size $K'_\Lambda(M)\tilde{\alpha}$;
$K_f$ is the constant of the bound $|\log V_f|\le K_f\tilde{\alpha}$ and $K_u$ that of
the bound $|\log\bar u_0|\le K_uK_0\tilde{\alpha}$, both recalled in the proof, $V_f$ being
the field entering the flat-frame datum; and $c_Q$ is the constant of the flat-frame data
bound with which its member $c_QK_0$ of $K^{\flat}$ is written; it enters here only through
$K^{\flat}$.
\end{definition}

\begin{proof}
\emph{The reading set.} By the reading statement for first-part terms, whose form in
Ba{\l}aban's construction is given in \cite[(1.59)--(1.61)]{B-CMP122b}, a first-part term
is read at the objects of $\mathfrak{Q}''$. This is the first part in the sense of
\cite[(1.59)--(1.61)]{B-CMP122b}. The second part of \cite[\S1]{B-CMP122b} consists of the
terms with $X\subset Z^c$. By the statement that no term is read on a cut set, a first-part
term is never read at the configurations $U_{k,\Lambda}$, $U_1$, $U_{1,2}$ of
\cite[\S1]{B-CMP122b}, or at the interpolating
configurations of \cite[(1.54)]{B-CMP122b}. These configurations serve the second-part
terms. They enter first-part terms only through averages read at the far boundary layer
$Y_1\setminus Y_1^0$ \cite[(1.59)--(1.61)]{B-CMP122b}.

\emph{The data off $\Lambda$.} By \cite[(1.50)]{B-CMP122b}, each
$P\in\{U^0_k,U''_k\}$ is the minimiser $U_{\mathbf{B}''_k}$ of its data.
\cite[\S1]{B-CMP122b} states two facts. First, the determining set $\mathbf{B}''_k$
restricted to $\Lambda^c$, and the field $V''{\upharpoonright}\Lambda^c$, coincide with the
corresponding determining set and field for the new configuration $U_k$. Second, the
averages $M_{\cdot}(U_k)$, with $M_{\cdot}$ the block-averaging operation and its
determining set suppressed from the notation, and the field $V''$ are replaced by the
corresponding averages $M_{\cdot}(\mathbf{U})$. By the identification of the seed,
$\mathbf{U}$ is the seed $\zeta$ on
the seed's region $Y_{\mathrm{out}}^{+}$. Hence the data off $\Lambda$ are $\mathsf{v}$,
the seed's own data.

\emph{The data on $\mathfrak{B}_0$ and on $\mathfrak{K}_2$.} On $\mathfrak{B}_0$ the data of
$U''_k$ are $e^{ig_kB}V_0^{\mathrm{fr}}(\zeta)$, with
$V_0^{\mathrm{fr}}(\zeta)=M_{\mathbf{B}''_k}(U_0^{\mathrm{fr}}(\zeta))$. Here
$U_0(\zeta)$ is the configuration $U_{k,Z}$ with data
$\mathsf{v}{\upharpoonright}(Z\setminus\Lambda)$ and $V_\Lambda$; see
\cite[(1.79), (1.81)]{B-CMP122a} and \cite[(1.43)]{B-CMP122b}. Setting $B=0$ gives the data
of $U^0_k$. On the deep core $\mathfrak{K}_2$ the data are the common real variables
$V_h$, and $V_m$ with $m<h$, as in Definition~\ref{9.1:def-core-collar}.

\emph{The minimiser $\Xi_1(\sigma')$.} This is the minimiser on $\mathbf{B}^{\sim}(Y_1)$
described in \cite[(1.59)--(1.61)]{B-CMP122b}. It has the interior data just described and,
on the layer cells at $\partial Y_1^0$, the averages localised with the parameter
$\sigma'$. By the distance lemma for the layer cells, these cells lie at distance at least
$M-3M_1$ from $\Lambda$.

\emph{The frame.} \cite[\S1]{B-CMP122a} states that the configuration $U_0$ can be chosen
in the axial--Landau gauge from the beginning. By \cite[(1.86)--(1.87)]{B-CMP122a},
$U_0^{u_0}=U_0^{(\mathrm{AL})}=\exp(i\eta A_0)$ inside $\Lambda$, with $u_0$ the gauge
transformation taking $U_0$ to the axial--Landau gauge and $\eta$ the spacing of the lattice
on which \cite[(1.86)--(1.87)]{B-CMP122a} are written. On $Z''_k$, \cite[(1.87)]{B-CMP122a}
gives
\begin{equation*}
|A_0|,\ |\nabla^{\eta}A_0|,\ |\partial^{\eta*}\partial^{\eta}A_0|
<O(1)\,B_3^2B_5M^6\varepsilon_k,
\end{equation*}
with $B_3$, $B_5$ constants of \cite{B-CMP122a} and $\nabla^{\eta}$, $\partial^{\eta}$,
$\partial^{\eta*}$ its lattice derivatives. \cite[\S1]{B-CMP122b} states that the analytic
extension $A_0=H_{k,\Lambda}\bigl(\tfrac1i\log M_{\cdot}(U_0(\mathbf{U}))\bigr)$, with
$H_{k,\Lambda}$ the operator of \cite{B-CMP122b}, satisfies
\cite[(1.87)]{B-CMP122a} with $\alpha_{0,k}+\alpha_{1,k}=\tilde{\alpha}$ in place of
$\varepsilon_k$. Therefore
\begin{equation}\label{9.1:eq-core-data-1}
\begin{gathered}
U_0^{\mathrm{AL}}(\zeta)=e^{i\eta A_0(\zeta)}\ \text{on all of }
\Lambda\supset Z''_k\supset\mathfrak{K},\\
\ell_k\,|A_0(\zeta)|\le K_0\,\tilde{\alpha}\ \text{on }Z''_k,
\end{gathered}
\end{equation}
for every $\zeta$ in the tube, with $\ell_j$ the lattice spacing at level $j$ and $K_0$ the
constant with which the bound of \cite[(1.87)]{B-CMP122a}, with $\tilde{\alpha}$ in place
of $\varepsilon_k$, is written in this normalisation. Off
$\Lambda$, $U_0^{\mathrm{AL}}(\zeta)$ is the block-axial configuration $U_0(\zeta)$ of
\cite[(1.81)]{B-CMP122a}, and we put $u_0:=1$ there.

Both the axial--Landau frame and the flat frame $\flat$ lie in $\mathfrak{F}_G$. By the
frame lemma, part (b), they give the same operation $\mathbf{T}'$. \cite[\S1]{B-CMP122a}
states that changing the gauge of $U_0$ in the integral \cite[(1.76)]{B-CMP122a} requires
the corresponding compensating gauge transformations of the variables $V_h$, restricted to
$\partial^{+}$. Thus $V_h$ is co-rotated on $\partial^{+}$ when the frame changes, and this
is the defining condition of $\mathfrak{F}_G$. The block-axial frame of
\cite[(1.81)]{B-CMP122a} belongs to $\mathfrak{F}_G$ if and only if $\bar u_0$ is constant
on $\partial^{+}$.

\emph{Bound \textup{(i)}.} Let $c$ be a tower bond of level $\pi$; both of its level-$\pi$
blocks lie in $Z''_k$. At $c$, the average $M_{\mathbf{B}''_k}$ is the level-$\pi$ block
average $M^{\pi}$ over these blocks. By \eqref{9.1:eq-core-data-1}, $U_0^{\mathrm{AL}}(\zeta)$
equals $e^{i\eta A_0(\zeta)}$ on these blocks, so
\begin{equation*}
V_0^{\mathrm{AL}}(\zeta)(c)=M^{\pi}\bigl(e^{i\eta A_0(\zeta)}\bigr)(c).
\end{equation*}
We apply the lemma on block averages of configurations close to the identity, with base
configuration $1$, to the configuration $e^{i\eta A_0(\zeta)}$. Its size is
controlled by the second relation in \eqref{9.1:eq-core-data-1}. The hypothesis of that
lemma, $\ell_\pi K_0\tilde{\alpha}/\ell_k\le c_4$ with $c_4$ the constant of the lemma, is
satisfied in the present setting, and the lemma gives
\begin{equation*}
|V_0^{\mathrm{AL}}(\zeta)(c)-1|\le 34\,d\,K_0\,\tilde{\alpha}\,L^{\pi-k}
\end{equation*}
for every $\zeta$ in the tube, which is (i).

On tower bonds $V_0^{\flat}=V_0^{\mathrm{AL}}$: by the frame lemma, part (a), both blocks
of a tower bond lie in $Z''_k$, where the gauged configuration equals $e^{i\eta A_0}$.
Hence bound (i) and the tower-bond case of the flat-frame data bound, stated there with
right-hand side $K^{\flat}\ell_{\pi(c)}\tilde{\alpha}$, estimate the same quantity.

\emph{Bound \textup{(ii)}.} This bound is taken by statement from the flat-frame data
bound, together with the bound on the boundary gauge factor. We recall the content used,
as it appears in the proof of the flat-frame data bound.
\begin{itemize}
\item On $\mathfrak{i}$ the flat datum is $V_f(\zeta)\bar u_0(\zeta)^{-1}$; on the rind it
is $V_f(\zeta)$.
\item $|\log V_f|\le K_f\tilde{\alpha}$.
\item $|\log\bar u_0|\le K_uK_0\tilde{\alpha}$, by the bound on the boundary gauge factor.
\item By the frame lemma, parts (a) and (b), $V_0^{\mathrm{AL}}$ differs from
$V_0^{\flat}$ on these bonds by factors $\bar u_0$ only.
\end{itemize}
Bound (ii),
$|\log V_0^{\mathrm{AL}}(\zeta)(c)|\le 2K_\Lambda^{\flat}\tilde{\alpha}$ on the level-$k$
bonds of the rind, of $\mathfrak{i}$, and on those crossing $\partial\Lambda$, is the
flat-frame data bound as stated, with $K^{\flat}$ and $K_\Lambda^{\flat}$ as in the
statement, transferred to the axial--Landau frame by the frame lemma.
\end{proof}

\begin{definition}[The seed, the chart centre and criticality]\label{9.1:def-seed-criticality}
Let $\mathfrak{K}=\mathfrak{K}_1\cup\mathfrak{K}_2$ be the core, its collar and its deep part, and let $h$ be
the label level, all as in Definition~\ref{9.1:def-core-collar}. Let the units, radii and degrees on the
core be those of Notation~\ref{9.1:not-core-units}, and let the objects read on the core and their data be
those of Definition~\ref{9.1:def-core-data}.
\begin{itemize}
\item[(a)] \emph{The output $\Phi$ and its label.} $\Phi$ denotes the output of Ba{\l}aban's operation
$\mathbf{R}$ at level $k$. Its label is the new determining set $\mathbf{B}_k$, which is of level $k$ on
$Z$, by the labelling rule of the stage tables and by the second clause of the statement on the labels
of the healed kernels. The terms of $\Phi$, written $R'^{(k)}(X)$ in \cite{B-CMP122b}, are, in the form
stated there, analytic functions of the variables
$(\mathbf{U},\mathbf{J})$ defined on the space of complex configurations over $X$ of radii
$\alpha_{0,k}$, $\alpha_{1,k}$, written $U^c_k(X,\alpha_{0,k},\alpha_{1,k})$ there.
A \emph{seed} is a configuration $\zeta$ on the region $R$ (the region $Y_{out}^{+}$ of the
construction that produces the seed's set $\boldsymbol{\Omega}(R)$), the variable of the objects read on
the core in Definition~\ref{9.1:def-core-data}; $\boldsymbol{\Omega}(R)$ denotes the seed's set, the box
set produced by that construction, and $\mathfrak{T}_\Phi[7/8]$ the seed tube of $\Phi$.
\item[(b)] \emph{The root and the chart centre.} Let $U_c$ be the real, regular configuration supplied by
the initial stage of the construction of the seed's set, critical, in the sense of the
variational problem of \cite[Theorem~1]{B-CMP102b} with the sub-averages over $\boldsymbol{\Omega}(R)$ as
data, for the seed's set $\boldsymbol{\Omega}(R)$ over the label of $\Phi$; at that stage the level
function $\lambda_R$ of this set takes one and the same level, the one fixed by that construction, at
every point of $Y_4$, and hence of $Y_1$ and of $X$, where $Y_4\supseteq Y_1\supseteq X$ are regions of
that construction. We call $U_c$
the \emph{chart centre}. Let $\beta_e$ be a complex data coordinate with $|\beta_e|\le r_c$, and let
$A_e:=\mathfrak{X}_{\boldsymbol{\Omega}(R)}(\beta_e;U_c)$ be its Landau coordinate. With $\eta$ the
spacing of the lattice on which the configurations live, the \emph{root} is
\begin{equation*}
q_0:=e^{i\eta A_e}\,U_c ;
\end{equation*}
it is a datum of the seed tube $\mathfrak{T}_\Phi[7/8]$, for every seed.
\item[(c)] \emph{Sub-averages.} With $M_{\mathbf{B}''_k}(\cdot)$ Ba{\l}aban's average over the determining
set $\mathbf{B}''_k$ of the minimisers $U^0_k$, $U''_k$ of Definition~\ref{9.1:def-core-data} (the
letter $M$ in this average denotes the averaging operation, not the number $M=L^{m'}$ of the cube
clause in item (C3) of the lemma that follows), we put
\begin{equation*}
\mathsf{v}'':=M_{\mathbf{B}''_k}(q_0),\qquad \mathsf{v}''_c:=M_{\mathbf{B}''_k}(U_c).
\end{equation*}
For a bond $c$ we put $|A_e|_c:=\ell_c\sup_{P(c)}|A_e|$, where $\ell_c$ and $P(c)$ are the length unit
and the region attached to the bond $c$ in the averaging lemma.
\item[(d)] \emph{Tables on the core.} For $x\in\mathfrak{K}$, $m=m(x)$ is the zone index of $x$; the
stage-$N$ table entries are $\mathsf{T}^a_N(x)$, the units $\mathsf{u}^a_m(x)$, and the stage
coefficients $\varphi^{(N)}_m$. Here $k_0$ is the integer so denoted in \cite{B-CMP122a}, and the label
level satisfies $h\le k_0$; in the form
stated in \cite{B-CMP122a}, the conditions of \cite[(1.64)(i)]{B-CMP122a} for $m\le k_0$ carry no weight
factor built from the integer $\ell_\circ$ of that paper, and this is the role of $k_0$ in what follows.
We write $B$, $C$ for the constants so written in the cube
clause of the stage-table identity; in item (C3) of the lemma that follows these letters denote only
those constants.
\end{itemize}
\end{definition}

\begin{lemma}[Criticality of the chart centre and the root, and the tables on the core]
Let the objects be as in Definition~\ref{9.1:def-seed-criticality}. The following hold.
\begin{itemize}
\item[(C1)] Each of $\mathbf{B}''_k{\upharpoonright}R$, $\mathbf{B}^{\sim}(Y_1){\upharpoonright}R$, and the
set $\mathfrak{S}=\boldsymbol{\Omega}(X^+)$ of a first-part term $(F,X,j)$ (taken over
$\mathbf{B}''_k$ and truncated at $j$) refines $\boldsymbol{\Omega}(R)$. Hence $U_c$ and $q_0$ are
critical for all three sets, with their own sub-averages as data; in particular the first clause of the
entry lemma applies at the entry point $q_0$ on $\mathfrak{S}$.
\item[(C2)] Modulo block-centre gauges,
\begin{equation}\label{9.1:eq-seed-criticality-1}
\mathsf{v}''=e^{i\beta''_e}\,\mathsf{v}''_c,\qquad |\beta''_e(c)|\le 4\,|A_e|_c ,
\end{equation}
and for the raw representative, with trivial block-centre gauge transformations,
\begin{equation}\label{9.1:eq-seed-criticality-2}
|\beta''_e(c)|\le 34\,d\,|A_e|_c .
\end{equation}
\item[(C3)] Every $x\in\mathfrak{K}$ has $m(x)\le h\le k_0$, and
\begin{equation}\label{9.1:eq-seed-criticality-3}
\mathsf{T}^a_N(x)=\varphi^{(N)}_m\,\mathsf{u}^a_m(x)\ \ge\ \tfrac15\,\mathsf{u}^a_m(x),\qquad m=m(x),
\end{equation}
together with the cube clause of the stage-table identity, on cubes of size $CM\ell_m$ with bound
$BCM\alpha_{0,m}$. On zone $0$ (inside $Z_1$, if present) no table clause is imposed, and a first-part
term carries no dependence subject to analyticity: it does not depend on the fields $(\mathbf{U},\mathbf{J})$
and the background potential restricted to any component of $X\cap\Omega_0$ that does not connect to
$\Omega_j$ \cite{B-CMP119}.
\item[(C4)] On a live arrested piece $W_{\mathfrak{a}}\subset\mathfrak{K}$ (by the statement on live
arrested pieces of the healing construction, its arrest level $k_{\mathfrak{a}}$ satisfies
$k_{\mathfrak{a}}\le h$, and $W_{\mathfrak{a}}\subset Z_h$) the
structure $(m,1,\varphi^{(N)}_m)$ is replaced by the frozen triple
$(\ell(x),\omega_{\mathrm{arrest}}(x),F^{(k)}(x))$ of the frozen-structure clauses of the arrest theorem,
where $\ell(x)$ denotes the localisation length at $x$, $\omega_{\mathrm{arrest}}(x)$ the frozen weight
and $F^{(k)}(x)$ the frozen coefficient, and $F^{(k)}\ge\varphi_{\mathrm{arrest}}$, the table coefficient
of the arrest theorem. The
domain $\mathfrak{D}_F$ of the frozen terms is the one fixed in the domain assignment for frozen terms
that accompanies the seed-tube proposition. Every
ratio of table entries used in what follows is unchanged, with $1/5$ replaced by
$\varphi_{\mathrm{arrest}}$.
\end{itemize}
\end{lemma}

\begin{proof}
\emph{The root is a datum of the seed tube.} By the definition of the seed tube
$\mathfrak{T}_\Phi[7/8]$ and by the first clause of the seed-tube proposition, which applies to any
seed, the configuration $q_0=e^{i\eta A_e}U_c$ with $|\beta_e|\le r_c$ is a datum of
$\mathfrak{T}_\Phi[7/8]$. The criticality of $U_c$ for $\boldsymbol{\Omega}(R)$ over the label of $\Phi$
is part of the initial stage of the construction of the seed's set, which also fixes the values of
$\lambda_R$ on $Y_4\supseteq Y_1\supseteq X$.

\emph{Proof of \textup{(C1)}.} On $Z$ the three sets are, against level $k$, respectively graded, graded
with far layers, and graded and truncated. Off $Z$ all three equal $\mathbf{B}_k$ up to layers and
truncation: by \cite{B-CMP122b}, the determining set $\mathbf{B}''_k$ and its field, written $V''$ in
\cite{B-CMP122b}, restricted to $\Lambda^c$, coincide with the corresponding determining set
and field for the new configuration $U_k$. The construction of the seed's set has the
property that the sets it produces only refine. Hence each of the three sets refines
$\boldsymbol{\Omega}(R)$. By the second clause of the nesting lemma for critical configurations, applied
to the refinements just listed, $U_c$ is critical for the three sets, with its own sub-averages as data.
The root $q_0$ is obtained from $U_c$ through the complex data coordinate $\beta_e$; since that clause
holds for every complex data coordinate, it applies to $q_0$ as well and
gives criticality of $q_0$ for the three sets, with its own sub-averages as data. In particular, since
$q_0$ is critical for $\mathfrak{S}$, the first clause of the entry lemma applies at the entry point
$q_0$ on $\mathfrak{S}$.

\emph{Proof of \textup{(C2)}.} The averaging lemma, applied at base $U_c$ to $q_0=e^{i\eta A_e}U_c$,
together with its second clause, gives $\mathsf{v}''=e^{i\beta''_e}\mathsf{v}''_c$ modulo block-centre
gauges, with the bound \eqref{9.1:eq-seed-criticality-1}. For the raw representative, with trivial
block-centre gauge transformations, the same two clauses give the bound
\eqref{9.1:eq-seed-criticality-2}; this is the representative used in the lemma bounding the raw data
of the seed.

\emph{Proof of \textup{(C3)}.} By Definition~\ref{9.1:def-core-collar}, every point of $\mathfrak{K}$ lies
on a zone of index at most $h$: the collar $\mathfrak{K}_1$ is at label level $h$, and the deep part
$\mathfrak{K}_2$ consists of its zone-$h$ part and of zones of index $m<h$. So a point
$x\in\mathfrak{K}$ has $m(x)\le h$, and $h\le k_0$ by item (d). We apply the
stage-table identity. This identity rests on \cite[(1.64)(i)]{B-CMP122a}, which is imposed
on $(\Omega''_m\setminus\Omega''_{m+1})\cap X$ for $m=1,\dots,j-1$; in the form stated in
\cite{B-CMP122a}, it carries no $\ell_\circ$-weight for $m\le k_0$. It gives
\eqref{9.1:eq-seed-criticality-3}. The cube clause of
the same identity gives the bound $BCM\alpha_{0,m}$ on cubes of size $CM\ell_m$.

For zone $0$: in \cite{B-CMP119} the corresponding conditions are imposed for $n=1,\dots,j-1$ only, so
zone $0$ carries no clause. Moreover \cite{B-CMP119} states that if a component of $X\cap\Omega_0$ does
not connect to $\Omega_j$, the function does not depend on the fields $(\mathbf{U},\mathbf{J})$ and the
background potential, written $A$ in \cite{B-CMP119},
restricted to this component. Hence zone $0$ carries no dependence subject to analyticity. Under the
standing hypothesis on the core, the sites of zone $0$ satisfy the regularity condition with parameter
$\varepsilon$ of the definition of regular configurations in any case.

\emph{Proof of \textup{(C4)}.} By the statement on live arrested pieces of the healing construction, a
live arrested piece $W_{\mathfrak{a}}\subset\mathfrak{K}$ has arrest level $k_{\mathfrak{a}}\le h$ and
$W_{\mathfrak{a}}\subset Z_h$.
On it, by the frozen-structure clauses of the arrest theorem, the table structure
$(m,1,\varphi^{(N)}_m)$ is replaced by the frozen $(\ell(x),\omega_{\mathrm{arrest}}(x),F^{(k)}(x))$, with
$F^{(k)}\ge\varphi_{\mathrm{arrest}}$. The weight
$\omega_{\mathrm{arrest}}(x)$ multiplies both sides of
every ratio of table entries, so every such ratio is unchanged. The lower bound
$F^{(k)}\ge\varphi_{\mathrm{arrest}}$ then takes the place of the lower bound $1/5$ in
\eqref{9.1:eq-seed-criticality-3}.
\end{proof}

\begin{definition}[Regularity of the common core variables]\label{9.1:def-core-regularity}
Let $\mathfrak{K}=\mathfrak{K}_1\cup\mathfrak{K}_2$ be the core region, with its collar and its deep
part, of Definition~\ref{9.1:def-core-collar}, let $Z$ be the large-field region containing it and
$h$ its label level.

The \emph{common real core variables} are the following variables, which enter the two compared runs
$\mathsf{A}$ and $\mathsf{B}$ as common integration parameters:
\begin{itemize}
\item the real variable $V_h$ on the zone-$h$ part $\mathfrak{K}_2\cap\Omega_h$ of the deep part;
\item the older real variables $V_m$, $m<h$, on the zones of $Z_h$.
\end{itemize}

They are called \emph{regular} if they lie in the supports of the effective density's own local
small-field functions, in the following sense.

Fix a zone $m\ge 1$ of $\mathfrak{K}$, or the frozen zone of a live piece. Consider every cube
$\square$ of the density's partition on the scale of that zone which meets $\mathfrak{K}$. These are
the cubes of side $LM_2R_m\ell_m$, where $R_m$ is the number $R_j$ determining the cubes of
\cite[(1.3)]{B-CMP122a}, taken at $j=m$, and $\ell_m$ is the scale factor of zone $m$, so that
$\eta^2\ell_m^{-2}$ is the quantity $(L^{k-j}\eta)^2$ of \cite[(1.3)]{B-CMP122a} at $j=m$,
$k$ being the current level. For
$m=h$ on
$\mathfrak{K}\cap\Omega_h$ they are the $LM_2R_h$-cubes of the characteristic function
$\chi_{h,1/2}$ of \cite[(1.88)]{B-CMP122a}; there the local function is evaluated on the
configuration $(1,V_h)$ of \cite[(1.88)]{B-CMP122a}, equal to $1$ on $\Omega''^{\sim 2}_{h+1}$ and to
$V_h$ on its complement.

For each such cube $\square$ the local function $U_{m,\square}(V_m)$ of \cite[(1.3)]{B-CMP122a}
(with $j=m$) is required to satisfy
\begin{equation}\label{9.1:eq-core-regularity-a}
\bigl|U_{m,\square}(V_m,\partial p)-1\bigr|\;\le\;c_\chi\,\varpi\,\varepsilon_m\,\eta^2\,\ell_m^{-2}
\qquad\text{for every plaquette } p\subset\square .
\end{equation}
Here $\varepsilon_m$ and $\eta$ are the quantities denoted by these letters in the condition of
\cite[(1.3)]{B-CMP122a}, taken at zone $m$, and $c_\chi\le 3$. The factor $\varpi$ is the weight
carried by a live piece, with $\varpi=1$ off live pieces.
\end{definition}

\begin{remark}
Condition \eqref{9.1:eq-core-regularity-a} is the quantitative form of the small-field condition
which the hypothesis on common core variables requires at the sites of these variables. Its form is
read off at four places: three in Ba{\l}aban's construction and one in the proof of the lemma
representing the live arrested pieces.

\emph{Zones of the most recent steps.} For these zones, \cite[(1.3)]{B-CMP122a} contains the
characteristic functions
\begin{equation*}
\chi\bigl(|U_{j,\square}(V_j,\partial p)-1|<\varepsilon_j(L^{k-j}\eta)^2\bigr).
\end{equation*}
Here $j$ is the zone index, written $m$ above, $k$ is the current level, and
$(L^{k-j}\eta)^2=\eta^2\ell_j^{-2}$.
Ba{\l}aban specifies their cubes as follows \cite[\S1]{B-CMP122a}: the cubes $\square$ in
\cite[(1.3)]{B-CMP122a} are the $LM_2R_j$-cubes, and the cubes $\square'$ in
\cite[(1.4), (1.7)--(1.9)]{B-CMP122a} are the $LM_2R_{j+1}$-cubes.

\emph{The zone-$h$ part after $\mathbf{R}$.} After the operation $\mathbf{R}$, of the characteristic
functions of these recent zones only that of zone $h$ survives, and only on the part of zone $h$
lying in $\mathfrak{K}$. It does so through $\chi_{h,1/2}$, that is, with the bound at half strength
\cite[(1.88)]{B-CMP122a}.

\emph{Inside $Z_h$.} Here the same functions occur one induction level down: the previous regions
satisfy the inductive condition on their own scales \cite[\S1, condition~(ii)]{B-CMP122a}. The
operation $\mathbf{T}_h(Z_h)$ is the inductive operation, with its own partition of unity
$\{\zeta_\square\}$ and its own characteristic functions $\chi$ \cite[(2.19)--(2.21)]{B-CMP119}.
Moreover, \cite[\S1]{B-CMP122a} states,
as part of the inductive assumption, that all the characteristic functions introduced there depend on
the field variables localised in the corresponding components.

\emph{Live arrested pieces.} On a live arrested piece $W$ with frozen zone $m$, the proof of the
lemma representing the live arrested pieces carries the bound that the plaquette values of the local
function of the piece at the level of its arrest, evaluated on the restriction $V{\upharpoonright}_W$
of the variables to $W$, differ from $1$ at every point of the piece by less than
\begin{equation*}
c\,\varpi\,\varepsilon_m L^{-2m},
\end{equation*}
with a constant $c$; this is the form in which \eqref{9.1:eq-core-regularity-a} enters on the frozen
zone of a live piece.
\end{remark}

\begin{remark}[What the reading of a term consumes: cellwise data increments]
\label{9.1:rem-cellwise-increments}
Let $(F,X,j)$ be a term of the expansion on the core $\mathfrak{K}$ of
Definition~\ref{9.1:def-core-data}. Let it be read at an object $q\in\mathfrak{Q}''$ over a seed
$q_0$ in the chart of the shape, and at a vertex of the complex parameter polydisc
$\mathfrak{P}_{\mathsf{c}}$ of the shape.
\begin{enumerate}
\item[(a)] \emph{What is needed.} The bound for terms read on the chart --- in its plain form,
its accumulated form and its form along the run --- and the dissolution bound of the low-order
scheme all need, for such a term, the single input
\begin{equation}\label{9.1:eq-cellwise-increments-1}
\mathrm{rd}\,q\in\mathfrak{T}[0](X).
\end{equation}
By the definition of the target tube $\mathfrak{T}[0](X)$ this means the following. Let the slice
$S(q)$ be the $\mathfrak{S}$-critical configuration that has the $\mathfrak{S}$-data of $q$,
written in the Landau chart about the real centre of the shape. Then $S(q)$ lies in
$\mathfrak{D}_F[0]$. By the membership criterion for $\mathfrak{D}_F[0]$, this amounts to two
conditions on the entry table plus the moves. At the $(=)$-points this sum is at most a fraction
of the gap at that point. At the $(\mathrm{up})$-points it is at most
$(\tfrac34+\tfrac1{200})\mathsf{u}_F$.
\item[(b)] \emph{What is read.} Let $q=e^{i\eta\mathcal{Y}}q_0$ be any presentation of $q$ over
$q_0$, with presentation field $\mathcal{Y}$ and presentation parameter $\eta$. For every member
and every point $x$, part (b$'$) of the move lemma gives: the move of the member at $x$ is at most
\begin{equation}\label{9.1:eq-cellwise-increments-2}
16\,B_\natural\,\ell_{\mathfrak{S}}(x)^{-e_a}\,\bigl(|\mathcal{Y}|_{\cdot}\bigr)^{\approx}(x).
\end{equation}
Here $\ell_{\mathfrak{S}}(x)^{-e_a}$ is the position-dependent weight of the move lemma, with
$\ell_{\mathfrak{S}}$ and the exponent $e_a$ fixed in its statement; $|\mathcal{Y}|_c$ is the
supremum of $|\mathcal{Y}|$ over the cell $c$, $|\mathcal{Y}|_{\cdot}$ is the function
$c\mapsto|\mathcal{Y}|_c$, and $(\cdot)^{\approx}$ is the smeared majorant.
Its proof uses the presentation only through the cellwise data increment
\begin{equation}\label{9.1:eq-cellwise-increments-3}
|B[q]-B[q_0]|(c)\le 4\,|\mathcal{Y}|_c .
\end{equation}
The local Lipschitz property $\natural\mathrm{Lip}^{\mathrm{loc}}$ of the Landau chart converts
such data increments into the Landau field.
\item[(c)] \emph{Representatives may be chosen freely.} Consider the hypothesis of
$\natural\mathrm{Lip}^{\mathrm{loc}}$: the data representatives $B^1$ and $B^2$ are complex with
$\sup|B^i|\le a_{R'}$. In it, each $B^i$ may be any representative of its data orbit in that
ball. The two representatives may be chosen independently of each other, and complex block-centre
gauges with values in the complexified gauge group $G^{c}$ are allowed.
\item[(d)] \emph{Consequences on the core.} Each minimiser-type object $P\in\mathfrak{Q}''$ is
handled directly through its data; $q$ is never presented as a rotation of another configuration.
Write $P^{\mathrm{re}}:=P(\zeta=U_c,\ B=0)$ for the real object attached to $P=P(\zeta,B)$,
where $U_c$ is the real centre; here $\zeta$ is the seed variable and $B$ the variable of the
$\mathbf{T}$-step, while, as in item~(b), $B[\cdot]$ denotes the block data of a configuration.
This has three consequences.
\begin{enumerate}
\item[(i)] No change of frame is needed at complex seeds.
\item[(ii)] On deep-core cells no bracket term across the seam is needed.
\item[(iii)] The only new estimate is real and common. It is the bound on
$B[P^{\mathrm{re}}]-B[q_0]$ on the core.
\end{enumerate}
\end{enumerate}
\end{remark}

\begin{proof}
\emph{Item}~(a). This is the definition of the target tube $\mathfrak{T}[0](X)$, combined with
the membership criterion for $\mathfrak{D}_F[0]$.

\emph{Item}~(b). The bound \eqref{9.1:eq-cellwise-increments-2} is part (b$'$) of the move lemma,
applied to the given presentation $q=e^{i\eta\mathcal{Y}}q_0$. Its proof estimates the data of $q$
against those of $q_0$ cell by cell. That cellwise estimate is \eqref{9.1:eq-cellwise-increments-3},
which is the data-increment bound of the tower estimate. The conversion into the Landau field is
the statement of $\natural\mathrm{Lip}^{\mathrm{loc}}$.

\emph{Item}~(c). By the definition of the target tube, the chart point is determined by the data
$M_{\mathfrak{B}}(\mathbf{V}(B))$, the image of $\mathbf{V}(B)$ under the block-data map
$M_{\mathfrak{B}}$, modulo block-centre gauges, and so is its Landau coordinate.
This is the form in which the lemma on the dependence of the data on block-centre gauges applies.
Take two representatives of one data orbit. Their critical configurations are gauge-equivalent,
because \cite[Theorem~1]{B-CMP102b} states that the critical orbit with given data is unique.
Hence, by the chart proposition on critical orbits, they have the same Landau representative
about the centre. It follows that each $B^i$ in the hypothesis of $\natural\mathrm{Lip}^{\mathrm{loc}}$
may be replaced by any representative of its data orbit in the ball, independently for $i=1,2$.
Block-centre gauges with values in $G^{c}$ are admitted by the third item of the same lemma.

\emph{Item}~(d). Every minimiser-type $P\in\mathfrak{Q}''$ has prescribed data
on the cells of $\mathbf{B}''_k$ (respectively of $\mathbf{B}^{\sim}(Y_1)$), by
Definition~\ref{9.1:def-core-data}. By Definition~\ref{9.1:def-seed-criticality}, the seed $q_0$
and the centre are critical for the same set. The input \eqref{9.1:eq-cellwise-increments-1} for
such a $P$ is then obtained in four steps.
\begin{itemize}
\item The increment $|B[P]-B[q_0]|$ is bounded cell by cell from the data, by the cellwise data
lemma.
\item $\natural\mathrm{Lip}^{\mathrm{loc}}$ is applied about the centre on that same set. By
item~(c) the representatives may be chosen freely. This gives a presentation field
$\mathcal{Y}_P$ with cellwise sup bounds.
\item The link steps of the chain are composed by the composition rule for presentations, in
the same way as in the composition that yields the chain field $\mathcal{Z}^{4F}$.
\item The bound \eqref{9.1:eq-cellwise-increments-2} of the move lemma is applied to the term
$(F,X,j)$.
\end{itemize}
The entry table is the root entry table of the move lemma. By part (a) of the move lemma and by
the proposition on the root coefficient, its coefficient is exactly $1$ at every point. Hence
nothing changes at the points of the two remaining point classes --- the class carried by the tube
and the class carried by the field --- except for one further source that is exponentially small.

\emph{Item}~(d)(i). Let $P^{\mathrm{re}}=P(\zeta=U_c,\ B=0)$ be the real object of item~(d).
On a cell $c$ of $\mathfrak{B}_0$, the data of $P(\zeta,B)$ and the data of $P^{\mathrm{re}}$
are both data in the axial-Landau frame $\mathrm{AL}$. Their ratio is
\begin{equation}\label{9.1:eq-cellwise-increments-4}
e^{ig_kB}\,V_0^{\mathrm{AL}}(\zeta)(c)\,V_0^{\mathrm{AL}}(U_c)(c)^{-1}.
\end{equation}
Here $V_0^{\mathrm{AL}}(\zeta)(c)$ denotes the field $V_0$ that the $\mathbf{R}$-step carries on
the bonds of $\mathfrak{B}_0$, taken in the cell $c$, written in the frame $\mathrm{AL}$ and
evaluated at the seed $\zeta$; $V_0^{\mathrm{AL}}(U_c)(c)$ is the same object at $U_c$.
This ratio is estimated raw and bond by bond, at every point of the tube, by the raw sizes
\eqref{9.1:eq-core-data-a}. On tower bonds it lies within
\begin{equation}\label{9.1:eq-cellwise-increments-5}
B^{\mathsf{s}}C_1\delta'_k+68\,d\,K_0\,\tilde{\alpha}\,L^{\pi-k}
\end{equation}
of $1$; this is the tower estimate at base $\mathbb{1}$, applied, with $\tilde{\alpha}$, to the
bound with constant $K_0$ from which the raw sizes \eqref{9.1:eq-core-data-a} on tower bonds are
derived. Here $B^{\mathsf{s}}$ is the dilation factor, $C_1$ is one of
Ba{\l}aban's auxiliary parameters $\mathfrak{p}$, $\delta'_k$ is the level-$k$ constant fixed
with the tube, $K_0$ is the constant of the bound underlying \eqref{9.1:eq-core-data-a},
$\tilde{\alpha}=\alpha_{0,k}+\alpha_{1,k}$, and $\pi$ is
the tower level entering the bound. On the level-$k$ bonds it lies within
\begin{equation}\label{9.1:eq-cellwise-increments-6}
B^{\mathsf{s}}C_1\delta'_k+4K_\Lambda^{\flat}\tilde{\alpha}
\end{equation}
of $1$, where $K_\Lambda^{\flat}$ is a frame constant. Thus no change of frame is applied at
complex seeds: neither the frame map $\varpi$, nor the adjoint action $\operatorname{Ad}$, nor
the three steps of the frame-map construction is applied there.

\emph{Item}~(d)(ii). On deep-core cells, the core-frame hypothesis says that the data of $P$ and
the data of $P^{\mathrm{re}}$ are the same common real variables. Hence, exactly,
\begin{equation}\label{9.1:eq-cellwise-increments-7}
B[P]-B[P^{\mathrm{re}}]=0\quad\text{on every deep-core cell.}
\end{equation}
The cut strata of the core construction are auxiliary cells of the sets attached to the two
configurations $U_1$ and $U_2$ of that construction. By the convention that fixes which cells a
term reads, they are read by no term. They are
also cells of none of $\mathbf{B}''_k$, $\mathbf{B}^{\sim}(Y_1)$ and $\mathfrak{S}$.

\emph{Item}~(d)(iii). The only estimate that is new is that of
$B[P^{\mathrm{re}}]-B[q_0]$ on the core. It measures how far the block data of the old field are
from those of the smooth centre. It is real and common. It is bounded by the lemma that bounds
the block data of the old field against those of the centre, using the cubes of the density
itself. The frame enters at one place only:
through the raw size of the actual data on the collar, measured against the reference value $1$
of the guard factor.
\end{proof}

Let $U_c$ be the real centre of the shape (Definition~\ref{9.1:def-seed-criticality}). It is one real field, common to both constructions compared on the core. This is the content of the definition of the seed, of the identification of the real centre in the tube, and of the identification of its healing field.

The chart, the tables and the term $F$ are covariant under $G$-valued gauge transformations, by the gauge-covariance statement for the chart and the tables. In \cite{B-CMP119} the expansions are stated to be ``invariant with respect to $G$-valued gauge transformations''. The block averages are covariant by \cite[(11)]{B-CMP98}.

We may therefore fix the representative of $U_c$ once. Every object below is then taken in the correspondingly regauged representative. This includes Ba{\l}aban's data $M_{\cdot}(\mathbf{U})$, which are averages of the seed and hence covariant.

At the real seed two inputs combine: the representation, at the real seed, of the configuration
$U_0$ of \cite[(1.79)]{B-CMP122a} as a perturbation of the seed by a healing field, and the
description of the layer along $\partial Z$ at the real seed.
Together they give $U_0(U_c)=e^{i\eta\mathcal{W}_0}U_c$, up to gauge on $Z$, with
\begin{equation}\label{9.1:eq-centre-representative-1}
\ell_k\sup|\mathcal{W}_0|\le 2B^{\mathsf{s}}C_{sc_1}\mathsf{s}_L .
\end{equation}
Indeed, the source of the healing field is bounded by $\mathsf{s}_L$ by the source bound on $\Lambda$,
and $\partial Z$ lies at distance at least $2M(N_0-2)$, $N_0$ being the integer fixed with the radii
on the core in Notation~\ref{9.1:not-core-units}. Here $B^{\mathsf{s}}$ is the dilation factor and
$\mathsf{s}_L$ the seed size (Notation~\ref{9.1:not-core-units}); $C_{sc_1}$ is the constant of the
description of the layer along $\partial Z$ that bounds the healing field by its source; $\eta$ is
the factor in the exponential parametrisation $\exp i\eta A_0$ of \cite[(1.86)]{B-CMP122a}; and
$\ell_k$ is the unit at level $k$ (Notation~\ref{9.1:not-core-units}).

Let $\mathcal{W}_0'$ be $\mathcal{W}_0$ transported into the axial--Landau representative. Then $|\mathcal{W}_0'|=|\mathcal{W}_0|$, because the map relating the two representatives is $G$-valued at the real seed, by the $G$-valuedness of the frame change at real configurations. Hence the gauge copy
\begin{equation*}
\mathring{U}_c:=e^{-i\eta\mathcal{W}_0'}\,U_0^{\mathrm{AL}}(U_c)
\end{equation*}
represents $U_c$. Here $U_0^{\mathrm{AL}}(U_c)$ is the configuration $U_0$ of \cite[(1.79)]{B-CMP122a}, at the real seed, in the axial--Landau frame of Definition~\ref{9.1:def-core-data}. We take $\mathring{U}_c$ as the centre of the chart and from now on write $U_c$ for it.

On each set used below we put $\mathsf{v}_c:=M_{\cdot}(U_c)$. On a cell $c$ of level $\lambda$,
$\mathsf{v}_c(c)=M^{\lambda}(U_c)(c)$ is the $\lambda$-level block average of
Definition~\ref{1.5:def-block-average}.

Let $D^0$ be the data of the real object of Definition~\ref{9.1:def-core-data}:
\begin{itemize}
\item[] on $\mathfrak{B}_0$, $D^0=V_0^{\mathrm{AL}}(U_c)$;
\item[] off $\Lambda$, $D^0=\mathsf{v}_c$, since at $\zeta=U_c$ Ba{\l}aban's exterior data are
$M_{\cdot}(U_c)$ \cite[\S1]{B-CMP122b};
\item[] on $\mathfrak{K}_2$, $D^0$ is the common real data $V_h,V_m$.
\end{itemize}
Finally put
\begin{equation*}
a_\Lambda:=K_0\tilde{\alpha}+2B^{\mathsf{s}}C_{sc_1}\mathsf{s}_L .
\end{equation*}

\begin{lemma}[Distance of the centre's representative from the real data]
\label{9.1:lem-centre-representative}
With $U_c$, $\mathsf{v}_c$, $D^0$ and $a_\Lambda$ as above, put
\begin{equation*}
\mathfrak{d}_\lambda:=68\,d\,B^{\mathsf{s}}C_{sc_1}\mathsf{s}_L\,L^{\lambda-k},
\qquad
\mathfrak{d}'_\lambda:=34\,d\,a_\Lambda\,L^{\lambda-k}.
\end{equation*}
In $\mathfrak{d}_\lambda$ the exponent $\lambda-k$ is read as $0$ off $Z$. Then the following hold:
\begin{equation}\label{9.1:eq-centre-representative-a}
\begin{gathered}
\bigl|\log[D^0(c)\,\mathsf{v}_c(c)^{-1}]\bigr|\le\mathfrak{d}_\lambda,\\
|\mathsf{v}_c(c)-1|\le\mathfrak{d}'_\lambda ,
\end{gathered}
\end{equation}
for the following ranges of cells.
\begin{itemize}
\item[(a)] The first line of \eqref{9.1:eq-centre-representative-a} holds for every cell $c$ of $\mathbf{B}''_k$ of level $\lambda$.
\item[(b)] The second line of \eqref{9.1:eq-centre-representative-a} holds for every cell $c$ of level $\lambda$ both of whose blocks lie inside $\Lambda$; these include, in particular, the tower cells, the cells of the collar, and every cell of $\mathfrak{K}$.
\end{itemize}
\end{lemma}

\begin{proof}
\emph{Part (a) off $\Lambda$.} There both $D^0(c)$ and $\mathsf{v}_c(c)$ are data of type $\mathsf{v}_c$ of the same point of the gauge orbit. In fact they are equal, since $D^0=\mathsf{v}_c$ off $\Lambda$ by definition. Hence the left side of the first bound is $0$, and the bound holds, so that off $\Lambda$ one could take $\mathfrak{d}_\lambda=0$.

\emph{Part (a) on $\mathfrak{B}_0$.} On this set $D^0(c)=V_0^{\mathrm{AL}}(U_c)(c)$, which is the
average of $U_0^{\mathrm{AL}}(U_c)$ at $c$ (Definition~\ref{9.1:def-core-data}); in
$V_0^{\mathrm{AL}}(U_c)$ and $U_0^{\mathrm{AL}}(U_c)$ the argument $U_c$ is the real centre, at
which these objects are evaluated. We apply the perturbation lemma for block averages with base
$U_0^{\mathrm{AL}}(U_c)$ and field $-\mathcal{W}_0'$.

The base is regular. It is a gauge copy of the configuration $U_0$ of \cite[(1.79)]{B-CMP122a}, whose plaquettes obey \cite[(1.80)]{B-CMP122a} with $\tilde{\alpha}$ in place of the radius there \cite{B-CMP122b}.

The size of the field satisfies
\begin{equation*}
\ell_\lambda\sup|\mathcal{W}_0'|\le 2B^{\mathsf{s}}C_{sc_1}\mathsf{s}_L\,L^{\lambda-k}\le c_4 .
\end{equation*}
The first inequality follows from \eqref{9.1:eq-centre-representative-1}, from $|\mathcal{W}_0'|=|\mathcal{W}_0|$, and from the relation between the units at levels $\lambda$ and $k$ (Notation~\ref{9.1:not-core-units}). Here $c_4$ is the smallness threshold of the perturbation lemma.

Since the centre of the chart is $\mathring{U}_c=e^{-i\eta\mathcal{W}_0'}U_0^{\mathrm{AL}}(U_c)$, with
$U_c$ the real centre in the argument of $U_0^{\mathrm{AL}}$, the perturbation lemma then gives
\begin{align*}
\mathsf{v}_c(c)=M^{\lambda}(\mathring{U}_c)(c)
&=M^{\lambda}\bigl(e^{-i\eta\mathcal{W}_0'}U_0^{\mathrm{AL}}(U_c)\bigr)(c)\\
&=\mathsf{q}_-\,e^{i\mathcal{Q}_\lambda}\,D^0(c)\,\mathsf{q}_+^{-1},
\end{align*}
with factors $\mathsf{q}_\pm$ and an exponent $\mathcal{Q}_\lambda$ furnished by the perturbation lemma, which satisfy
\begin{equation*}
|\mathcal{Q}_\lambda|\le 2\ell_\lambda\sup|\mathcal{W}_0'|,\qquad
|\mathsf{q}_\pm-1|\le 16\,d\,\ell_\lambda\sup|\mathcal{W}_0'|.
\end{equation*}
It follows that
\begin{equation*}
D^0(c)\,\mathsf{v}_c(c)^{-1}
=\bigl(D^0(c)\,\mathsf{q}_+\,D^0(c)^{-1}\bigr)\,e^{-i\mathcal{Q}_\lambda}\,\mathsf{q}_-^{-1}.
\end{equation*}
Since $D^0(c)$ is $G$-valued, $\operatorname{Ad}_{D^0(c)}$ is an isometry. Hence the first factor lies at the same distance from $1$ as $\mathsf{q}_+$. Therefore
\begin{equation*}
\bigl|\log[D^0(c)\,\mathsf{v}_c(c)^{-1}]\bigr|
\le(2+32d)\cdot 1.01\cdot\ell_\lambda\sup|\mathcal{W}_0'|
\le(4+64d)\cdot 1.01\cdot B^{\mathsf{s}}C_{sc_1}\mathsf{s}_L\,L^{\lambda-k};
\end{equation*}
here the factor $1.01$ accounts for the passage from the sizes of the three factors $D^0(c)\mathsf{q}_+D^0(c)^{-1}$, $e^{-i\mathcal{Q}_\lambda}$ and $\mathsf{q}_-^{-1}$ to the logarithm of their product, all three being close to the identity because $\ell_\lambda\sup|\mathcal{W}_0'|\le c_4$.
For $d=4$ we have $(4+64d)\cdot1.01=262.6\le 272=68d$. So the right side is at most $\mathfrak{d}_\lambda$, which proves (a).

\emph{Part (b).} Let $c$ have both blocks inside $\Lambda$. By Definition~\ref{9.1:def-core-data}, the axial--Landau configuration is $U_0^{\mathrm{AL}}(U_c)=e^{i\eta A_0(U_c)}$ on all of $\Lambda\supset Z''_k\supset\mathfrak{K}$, with
\begin{equation*}
\ell_k|A_0(U_c)|\le K_0\tilde{\alpha}.
\end{equation*}
Hence there
\begin{equation*}
\mathring{U}_c=e^{-i\eta\mathcal{W}_0'}\,e^{i\eta A_0(U_c)}
=e^{i\eta\,\mathrm{BCH}(-\mathcal{W}_0',A_0)},
\end{equation*}
where $\mathrm{BCH}(\mathcal{X}_1,\mathcal{X}_2)$ denotes the Baker--Campbell--Hausdorff field, defined by
$e^{i\eta\mathcal{X}_1}e^{i\eta\mathcal{X}_2}=e^{i\eta\,\mathrm{BCH}(\mathcal{X}_1,\mathcal{X}_2)}$.

The two fields have $\ell_k$-sizes at most $2B^{\mathsf{s}}C_{sc_1}\mathsf{s}_L$ (by \eqref{9.1:eq-centre-representative-1}) and $K_0\tilde{\alpha}$. Their sum is $a_\Lambda$, and
\begin{equation*}
\ell_k\sup|\mathrm{BCH}(-\mathcal{W}_0',A_0)|\le 1.01\,a_\Lambda .
\end{equation*}

We apply the perturbation lemma for block averages at the base $\mathbb{1}$, with the field
$\mathrm{BCH}(-\mathcal{W}_0',A_0)$. The averages of the base $\mathbb{1}$ are $1$, so in the
representation used in part (a) the datum $D^0(c)$ is replaced by $1$:
\begin{equation*}
\mathsf{v}_c(c)=\mathsf{q}_-\,e^{i\mathcal{Q}_\lambda}\,\mathsf{q}_+^{-1},
\end{equation*}
where now $|\mathcal{Q}_\lambda|$ and $|\mathsf{q}_\pm-1|$ obey the bounds of part (a) with
$\ell_\lambda\sup|\mathcal{W}_0'|$ replaced by $\ell_\lambda\sup|\mathrm{BCH}(-\mathcal{W}_0',A_0)|$.
By the algebraic identity
\begin{equation*}
\mathsf{q}_-e^{i\mathcal{Q}_\lambda}\mathsf{q}_+^{-1}-1
=(\mathsf{q}_--1)\,e^{i\mathcal{Q}_\lambda}\mathsf{q}_+^{-1}
+(e^{i\mathcal{Q}_\lambda}-1)\,\mathsf{q}_+^{-1}+(\mathsf{q}_+^{-1}-1),
\end{equation*}
$|\mathsf{v}_c(c)-1|$ is bounded by the sum $16d+2+16d=2+32d$ of the three coefficients times
$\ell_\lambda\sup|\mathrm{BCH}(-\mathcal{W}_0',A_0)|$, up to a factor $1.01$ that absorbs the
factors $|e^{i\mathcal{Q}_\lambda}\mathsf{q}_+^{-1}|$ and $|\mathsf{q}_+^{-1}|$ and the passage from
$|\mathsf{q}_+-1|$ to $|\mathsf{q}_+^{-1}-1|$, all within one per cent of their values at the
identity under the smallness condition $c_4$ of the perturbation lemma. Together with the relation
between the units at levels $\lambda$ and $k$ (Notation~\ref{9.1:not-core-units}) this gives
\begin{align*}
|\mathsf{v}_c(c)-1|
&\le(2+32d)\cdot1.01\cdot\ell_\lambda\sup|\mathrm{BCH}(-\mathcal{W}_0',A_0)|\\
&\le(2+32d)\cdot1.01^2\,a_\Lambda\,L^{\lambda-k}
\le 34\,d\,a_\Lambda\,L^{\lambda-k}=\mathfrak{d}'_\lambda ,
\end{align*}
since for $d=4$ we have $(2+32d)\cdot1.01^2=132.613\le136=34d$.

This holds on $\mathfrak{K}_2$ as well. There $A_0$ is defined and bounded on all of $Z''_k\supset\mathfrak{K}$ \cite[(1.86)--(1.87)]{B-CMP122a}, and $\mathring{U}_c$ is smooth. Thus the centre's representative has no jumps across faces of $k$-blocks on the core. This is unlike the configuration $U_0$ in the axial gauge of \cite[(1.81)]{B-CMP122a}, whose jumps across faces of $k$-blocks are described by the layer profile $\mathsf{m}_L$. This proves (b).
\end{proof}

We use the following objects:

\begin{itemize}
\item the core region $\mathfrak{K}=\mathfrak{K}_1\cup\mathfrak{K}_2$ of Definition~\ref{9.1:def-core-collar} (see \eqref{9.1:eq-core-collar-a}), with its collar $\mathfrak{K}_1$, its deep core $\mathfrak{K}_2$ and its zones;
\item the objects read on the core and their raw size bounds \eqref{9.1:eq-core-data-a} (Definition~\ref{9.1:def-core-data});
\item the regularity \eqref{9.1:eq-core-regularity-a} of the common core variables (Definition~\ref{9.1:def-core-regularity}), with its constants $c_\chi\le 3$, $\varepsilon_m$, $\eta$, $\ell_m$ and its weight $\varpi$ ($\varpi=1$ off live pieces);
\item the centre's representative $U_c$ and its block averages $\mathsf{v}_c$ from Lemma~\ref{9.1:lem-centre-representative};
\item the two bounds \eqref{9.1:eq-centre-representative-a} of that lemma. The first bounds $|\log[D^0(c)\mathsf{v}_c(c)^{-1}]|$ by $\mathfrak{d}_\lambda$ on every cell $c$ of $\mathbf{B}''_k$ of level $\lambda$. The second bounds $|\mathsf{v}_c(c)-1|$ by $\mathfrak{d}'_\lambda$ on every cell of level $\lambda$ with both blocks inside $\Lambda$; this includes all of $\mathfrak{K}\subset\Lambda$.
\end{itemize}

We write $G=\mathrm{SU}(N)$ and $G^{\mathbb{C}}$ for its complexification. For a cell-bond $c$ we write $c_-$, $c_+$ for the centres of its initial and final blocks. For a gauge transformation $v$ on sites, $U^v$ denotes the transformed configuration (Lemma~\ref{1.1:lem-invariance}); for data $V$ on cell-bonds and a map $v$ on block centres we put $V^v(c)=v(c_-)V(c)v(c_+)^{-1}$. We write $\mathcal{M}^m$ for the block-averaging map of Definition~\ref{1.5:def-block-average} to level $m$.

\emph{The real object.} Let $V$ denote the common real core variables: $V_h$ on $\mathfrak{K}_2\cap\Omega_h$, and $V_m$ on zone $m$ of $Z_h$. Let
\begin{equation*}
U^{0,\mathrm{re}}:=U^0_k(U_c;B=0)
\end{equation*}
be the configuration determined, in the frame $\mathrm{AL}$, by the following data on $\mathbf{B}''_k$:
\begin{itemize}
\item $\mathsf{v}_c$ off $\Lambda$;
\item $V_0^{\mathrm{AL}}(U_c)$ on $\mathfrak{B}_0$;
\item $V$ on $\mathfrak{K}_2$.
\end{itemize}

Let $\mathbf{B}''_k{\upharpoonright}R$ denote the restriction of Ba{\l}aban's determining set $\mathbf{B}''_k$ to the region $R$ (here $R$ is a region of the lattice, not the ball radius of the glossary). The datum of $U^{0,\mathrm{re}}$ on a cell-bond $c$ of $\mathbf{B}''_k{\upharpoonright}R$ is $D^0(c)$; these are the data $D^0$ of Lemma~\ref{9.1:lem-centre-representative}.

A \emph{block-centre gauge} is a map $g$ from cell centres to $G^{\mathbb{C}}$. For such $g$ we put
\begin{equation}\label{9.1:eq-core-gauge-1}
B^0_g(c):=\frac{1}{i}\log\bigl[g(c_-)\,D^0(c)\,g(c_+)^{-1}\,\mathsf{v}_c(c)^{-1}\bigr].
\end{equation}
This is a representative of the data coordinate of $U^{0,\mathrm{re}}$ in the chart centred at $U_c$, the centre fixed in Lemma~\ref{9.1:lem-centre-representative}.

\emph{Constants.} Let $K_m$ denote the bound on the linear extent of zone $m$ of $\mathfrak{K}$, in units of $\ell_m$, fixed together with the core region; it satisfies $K_m\le 100\sqrt{d}\,MLR_m$. Let $K_0$ and $\tilde\alpha=\alpha_{0,k}+\alpha_{1,k}$ be the constants of the raw size bounds \eqref{9.1:eq-core-data-a}; the constant $K_0$ is unrelated to the zone constants $K_m$. We put
\begin{equation}\label{9.1:eq-core-gauge-2}
\begin{aligned}
a_m&:=2dLM_2R_m\,c_\chi\varpi\,\varepsilon_m\qquad(\text{zone } m\le h \text{ of }\mathfrak{K}),\\
n_{\mathfrak{K}}&:=\bigl(100\sqrt{d}\,M/M_2+3\bigr)^d,\qquad
\mathfrak{q}_\lambda:=34dK_0\tilde\alpha L^{\lambda-k},\\
\varsigma_h&:=2dK_h(2a_h+2\mathfrak{q}_h)+4d(2n_{\mathfrak{K}}+1)LM_2R_ha_h .
\end{aligned}
\end{equation}
Since the density's scale-$m$ cubes have side $LM_2R_m$, the bound on $K_m$ shows that $n_{\mathfrak{K}}$ is at least the number of these cubes meeting zone $m$ of $\mathfrak{K}$.

The number $\varsigma_h$ is the gauge drift on the collar. Its first term comes from paths inside the collar and its second from the matching of cubes; both are explained in the proof of clause (c) below.

\begin{lemma}[A real common block gauge on the core]\label{9.1:lem-core-gauge}
Make the following three assumptions:
\begin{itemize}
\item the common core variables satisfy the regularity \eqref{9.1:eq-core-regularity-a};
\item the gauge frame for $U_0$ is the axial--Landau frame, $\mathrm{fr}=\mathrm{AL}$;
\item the ceiling $a_m\le c_4$ holds for $m\le h$, as part of the coupling ceilings \eqref{9.1:eq-core-ceilings-a}; here $c_4$ is the smallness constant of the averaging lemma for small perturbations of a regular configuration recalled in the proof.
\end{itemize}

Then there is a real, that is $G$-valued, block-centre gauge $g$ on the centres of $\mathbf{B}''_k{\upharpoonright}R$, with $g=1$ at every centre outside $\mathfrak{K}\cup\mathfrak{K}_1$, such that
\begin{equation}\label{9.1:eq-core-gauge-a}
\begin{aligned}
&\text{(a)}\quad B^0_g(c)=0, &&\qquad \text{(b)}\quad |B^0_g(c)|\le\mathfrak{d}_\lambda,\\
&\text{(c)}\quad |B^0_g(c)|\le\mathfrak{d}_h+2\varsigma_h, &&\qquad \text{(d)}\quad |B^0_g(c)|\le\mathfrak{b}_m,
\end{aligned}
\end{equation}
where the bounds hold on the following cells.
\begin{itemize}
\item Clause (a) holds on every cell off $\Lambda$.
\item Clause (b) holds on every cell of $\mathfrak{B}_0$ of level $\lambda$ outside the collar. These are the rind, the set $\mathfrak{i}$, the bonds crossing $\partial\Lambda$ and the towers; $g=1$ at their centres, and clause (b) is the first bound of \eqref{9.1:eq-centre-representative-a}.
\item Clause (c) holds on the collar cells and on the junction cells of the collar with the towers. In addition, $|g-1|\le\varsigma_h$ at every collar centre.
\item Clause (d) holds on three kinds of cells: the cells of $\mathfrak{K}_2$ of zone $m$ whose two end centres are zone-$m$ centres; the cells crossing $S_c$; and the junction cells between zones $m$ and $m+1$.
\end{itemize}

Here
\begin{equation}\label{9.1:eq-core-gauge-3}
\mathfrak{b}_m:=\mathfrak{d}'_m+2(a_m+a_{m+1})
+4d(2n_{\mathfrak{K}}+1)LM_2\bigl(R_ma_m+R_{m+1}a_{m+1}\bigr).
\end{equation}
For $m=h$ one reads $a_{h+1}:=0$ and adds $\varsigma_h$. In particular
\begin{equation}\label{9.1:eq-core-gauge-4}
\mathfrak{b}_m\le 50d^2n_{\mathfrak{K}}(LM_2)^2R_m^2\,c_\chi\varpi\,(\varepsilon_m\vee\varepsilon_{m+1})
+\mathfrak{d}'_m+\mathbb{1}_{m=h}\,\varsigma_h .
\end{equation}
The gauge $g$ is built from the common real data and $U_c$ only. It is the same for both runs $\mathsf{A}$ and $\mathsf{B}$, and it does not depend on $\zeta$, $B$, $\mathsf{t}$, $\mathsf{s}$, $\sigma'$.
\end{lemma}

\begin{proof}
\textit{Clause (a).} The data of $U^{0,\mathrm{re}}$ off $\Lambda$ are Ba{\l}aban's exterior data of the seed. At $\zeta=U_c$ these are the block averages of $U_c$, that is, $\mathsf{v}_c$. Hence $D^0(c)=\mathsf{v}_c(c)$ on every cell $c$ off $\Lambda$. Moreover $g=1$ at the centres of these cells. By \eqref{9.1:eq-core-gauge-1}, $B^0_g(c)=\frac{1}{i}\log 1=0$.

\textit{Clause (b).} On these cells $g=1$ at both end centres, so $B^0_g(c)=\frac{1}{i}\log[D^0(c)\mathsf{v}_c(c)^{-1}]$. The first bound of \eqref{9.1:eq-centre-representative-a} bounds this by $\mathfrak{d}_\lambda$.

Moreover, on tower bonds the first raw size bound of \eqref{9.1:eq-core-data-a}, taken at $\zeta=U_c$, gives
\begin{equation*}
|D^0(c)-1|\le 34dK_0\tilde\alpha L^{\lambda-k}=\mathfrak{q}_\lambda .
\end{equation*}
This bound is used below only on the collar, where $\lambda=h$.

\textit{Clause (d): one cube.} Let $\square$ be a cube on scale $m$, of the kind described in the regularity \eqref{9.1:eq-core-regularity-a}, that meets $\mathfrak{K}$. Let $U_\square:=U_{m,\square}(V)$ be its local function on $\square^{\sim}$.

For $m=h$ on $\mathfrak{K}\cap\Omega_h$, $U_\square$ is the guard's local function $U_{h,\square}((1,V_h))$ \cite{B-CMP122a}. Its data are $1$ on the cells of $\square^{\sim}$ on the ${\Omega''}^{\sim2}_{h+1}$ side and $V_h$ on the others, bonds crossing $S_c$ included.

By \eqref{9.1:eq-core-regularity-a}, the plaquette variables of spacing $\eta$ of $U_\square$ are within $c_\chi\varpi\varepsilon_m\eta^2\ell_m^{-2}$ of $1$ on $\square^{\sim}$.

Fix the axial gauge $u_\square$ of $U_\square$ on $\square^{\sim}$ from a corner. It is taken along the comb tree, that is, the coordinate-ordered tree of \cite[(1.15)]{B-CMP99a}. In this gauge $U_\square^{u_\square}=1$ on the tree bonds. The variable of a non-tree bond is the ordered product of the plaquette variables tiling its comb strip, and there are at most $d(LM_2R_m+2)\ell_m/\eta$ of them. Hence $U_\square^{u_\square}=e^{i\eta a_\square}$ with
\begin{equation}\label{9.1:eq-core-gauge-5}
\ell_m\sup|a_\square|\le d(LM_2R_m+2)\,c_\chi\varpi\varepsilon_m\le a_m .
\end{equation}

The data of $U_\square$ on the cells of $\square^{\sim}$ are $V$. Indeed, a local function returns its data as averages; this is so for the functions $U_{j,\square}(V_j)$ of \cite[(1.3)]{B-CMP122a} by \cite[Theorem~1]{B-CMP102b}.

We now use two facts.

First, \cite[(11)]{B-CMP98} states the covariance of block averages:
\begin{equation*}
\mathcal{M}^m(U^v)(c)=\bar v(c_-)\,\mathcal{M}^m(U)(c)\,\bar v(c_+)^{-1},
\end{equation*}
where $\bar v$ denotes the values of $v$ at the block centres.

Second, we use the averaging lemma for small perturbations of a regular configuration, at the base configuration identically equal to $1$. In that setting it reads as follows. For a real field $a$ with $\ell_m\sup|a|\le c_4$, one has
\begin{equation*}
\mathcal{M}^m(e^{i\eta a})(c)=\xi(c_-)e^{iQ(c)}\xi(c_+)^{-1},\qquad
|Q|\le 2\ell_m\sup|a|,\qquad |\xi-1|\le 16d\,\ell_m\sup|a| .
\end{equation*}
It applies here because $\ell_m\sup|a_\square|\le a_m\le c_4$, by \eqref{9.1:eq-core-gauge-5} and the assumed ceiling.

Together these give, on every cell-bond $c$ of $\square^{\sim}$,
\begin{equation*}
V(c)=\mathcal{M}^m(U_\square)(c)
=\bar u_\square(c_-)^{-1}\,\xi(c_-)\,e^{iQ_\square(c)}\,\xi(c_+)^{-1}\,\bar u_\square(c_+),
\end{equation*}
with $|Q_\square|\le 2a_m$ and $|\xi-1|\le 16da_m$. Here $\bar u_\square$ and $\xi$ are centre values, and they are real because $V$ and $U_\square$ are real.

Put $g_\square:=\xi^{-1}\bar u_\square$ on the centres of the cells of $\square^{\sim}$. Substituting the last display gives
\begin{equation}\label{9.1:eq-core-gauge-6}
g_\square(c_-)\,V(c)\,g_\square(c_+)^{-1}=e^{iQ_\square(c)},\qquad |Q_\square(c)|\le 2a_m,
\end{equation}
for every cell-bond $c$ of $\square^{\sim}$. This covers interior cells, cells crossing $S_c$ and collar-side cells of a guard cube alike.

On a collar-side cell of a guard cube, the guard's datum is $1$. There the actual datum $D^0(c)$ is within $\mathfrak{q}_h$ of $1$, by the first raw size bound of \eqref{9.1:eq-core-data-a} at $\zeta=U_c$ in the frame $\mathrm{AL}$. Writing
\begin{equation*}
g_\square(c_-)D^0(c)g_\square(c_+)^{-1}
=e^{iQ_\square(c)}\,\operatorname{Ad}_{g_\square(c_+)}\bigl(D^0(c)\bigr),
\end{equation*}
with $\operatorname{Ad}$ an isometry since $g_\square$ is $G$-valued, we obtain the same display as \eqref{9.1:eq-core-gauge-6} with $|Q|\le 2a_m+\mathfrak{q}_h(1+2a_m)$.

Finally, let $S\in G$ be constant. Replacing $g_\square$ by $S g_\square$ keeps \eqref{9.1:eq-core-gauge-6}, with $Q\mapsto\operatorname{Ad}_SQ$ and the norm unchanged.

\textit{Clause (d): matching the cubes.} For each zone $m$ let $\mathfrak{Q}_m$ be the set of scale-$m$ cubes meeting zone $m$ of $\mathfrak{K}$; then $|\mathfrak{Q}_m|\le n_{\mathfrak{K}}$.

Call two cubes of $\mathfrak{Q}_m\cup\mathfrak{Q}_{m+1}$ \emph{adjacent} if their enlargements share a layer of zone-$m$ cells. The scale-$m$ cubes tile zone $m$. The scale-$(m+1)$ cubes lie in $(\Omega^{\sim}_{m+1}\setminus\Omega_{m+2})\cap Z$ \cite[\S1]{B-CMP122a}, so they overlap zone $m$ in one layer.

On an overlap of adjacent cubes $\square\in\mathfrak{Q}_m$ and $\square'\in\mathfrak{Q}_{m'}$, $m'\in\{m,m+1\}$, both local functions carry the same data $V$. Solving \eqref{9.1:eq-core-gauge-6} for $V(c)$ in both cubes and equating gives
\begin{equation*}
g_\square(c_-)^{-1}e^{iQ_\square(c)}g_\square(c_+)=g_{\square'}(c_-)^{-1}e^{iQ_{\square'}(c)}g_{\square'}(c_+).
\end{equation*}
So $r:=g_\square g_{\square'}^{-1}$ satisfies $r(c_+)=e^{-iQ_\square(c)}\,r(c_-)\,e^{iQ_{\square'}(c)}$ along every cell-bond $c$ of the overlap.

A path $y'\to y$ inside one overlap, between centres $y$ and $y'$ of that overlap, has at most $d(LM_2R_m+2)$ overlap bonds. Hence
\begin{equation*}
|\log[r(y)r(y')^{-1}]|\le(2a_m+2a_{m'})\cdot d(LM_2R_m+2)=2d(LM_2R_m+2)(a_m+a_{m'}).
\end{equation*}

In each zone $m$ choose a spanning tree $T_m$ of the adjacency graph of $\mathfrak{Q}_m$. Choose also one adjacency between the root of $T_m$ and a cube of $\mathfrak{Q}_{m+1}$. For $m=h$ this is an adjacency between the root $\square_0\in\mathfrak{Q}_h$, a guard cube, and the collar (see below).

Going down from zone $h$, replace each cube's $g_\square$ by $S_\square g_\square$. Here $S_\square\in G$ is the constant that makes $r=1$ at one centre of the overlap of the cube with its parent.

Now take two cubes adjacent in the geometry, in zones $m$ and $m'\in\{m,m+1\}$. The tree path between them has at most $2n_{\mathfrak{K}}+1$ steps. Each step contributes a drift of at most $2d(LM_2R_{m''}+2)\cdot 2a_{m''}$ with $m''\in\{m,m+1\}$. So at every centre of their overlap
\begin{equation}\label{9.1:eq-core-gauge-7}
|\log r|\le 4d(2n_{\mathfrak{K}}+1)LM_2\bigl(R_ma_m+R_{m+1}a_{m+1}\bigr).
\end{equation}
Here $R_{m+1}\le R_m$ is used, and the summand $2$ in $LM_2R_{m''}+2$ is absorbed in the factor $4d$.

Fix a rule that assigns to every centre lying in some $\square^{\sim}$, $\square\in\bigcup_m\mathfrak{Q}_m$, one such cube, its \emph{assigned cube}. Define $g:=S_\square g_\square$ on the centres assigned to $\square$, and $g:=1$ elsewhere.

Let $c$ be a cell-bond whose two ends are assigned to one cube. Since $D^0=V$ there, \eqref{9.1:eq-core-gauge-1} and \eqref{9.1:eq-core-gauge-6} give $B^0_g(c)=\frac{1}{i}\log[e^{iQ(c)}\mathsf{v}_c(c)^{-1}]$. The second bound of \eqref{9.1:eq-centre-representative-a}, $|\mathsf{v}_c-1|\le\mathfrak{d}'_m$ on $\Lambda\supset\mathfrak{K}$, then yields
\begin{equation*}
|B^0_g(c)|\le 2a_m+\mathfrak{d}'_m(1+2a_m).
\end{equation*}

Now let $c$ have ends assigned to adjacent cubes, $\square\ni c_-$ and $\square'\ni c_+$. This covers the same zone, a zone junction, or a bond crossing $S_c$ whose ends are assigned to two guard cubes. With $g_\square$ and $g_{\square'}$ already multiplied by their constants, $g_\square$ extends to $c_+\in\square^{\sim}$, and
\begin{equation*}
B^0_g(c)=\frac{1}{i}\log\bigl[g_\square(c_-)V(c)g_\square(c_+)^{-1}\cdot r(c_+)\cdot\mathsf{v}_c(c)^{-1}\bigr],
\end{equation*}
with $r=g_\square g_{\square'}^{-1}$ at $c_+$. Hence $|B^0_g(c)|\le 2a_m+|\log r|+\mathfrak{d}'_m$. By \eqref{9.1:eq-core-gauge-7} this is the bound \eqref{9.1:eq-core-gauge-3}.

\textit{Clause (c) and the root.} The domain $\square_0^{\sim}$ of the root guard cube contains two kinds of cells:
\begin{itemize}
\item collar cells, whose datum is $1$ for the guard and $D^0$ for $U^{0,\mathrm{re}}$;
\item bonds crossing $S_c$, whose datum is $V_h$ for both.
\end{itemize}

Normalise $S_{\square_0}$ so that $g(y_\times)=1$ at one collar centre $y_\times\in\square_0^{\sim}$ adjacent to the towers.

Join a collar centre $y$ to $y_\times$ by a path of at most $2dK_h$ collar cell-bonds. Along a collar bond $c'$, with the guard's datum $1$, the display \eqref{9.1:eq-core-gauge-6} reads $g(c'_-)\cdot 1\cdot g(c'_+)^{-1}=e^{iQ(c')}$. Hence
\begin{equation*}
|\log[g(c'_+)g(c'_-)^{-1}]|\le|Q(c')|\le 2a_h+\mathfrak{q}_h(1+2a_h).
\end{equation*}
Summing along the path, $|g(y)-1|\le 2dK_h(2a_h+2\mathfrak{q}_h)$ for collar centres assigned to $\square_0$.

A collar centre assigned to another guard cube adds that cube's matched drift $|\log r|\le 4d(2n_{\mathfrak{K}}+1)LM_2R_ha_h$, by \eqref{9.1:eq-core-gauge-7} with $a_{h+1}=0$. Altogether $|g(y)-1|\le\varsigma_h$ at every collar centre; the two terms of $\varsigma_h$ are exactly these two contributions.

On a collar cell $c$ we write
\begin{equation*}
B^0_g(c)=\frac{1}{i}\log\Bigl[g(c_-)\cdot\bigl(D^0\mathsf{v}_c^{-1}\bigr)(c)\cdot
\operatorname{Ad}_{\mathsf{v}_c(c)}\bigl(g(c_+)^{-1}\bigr)\Bigr].
\end{equation*}
Since $\mathsf{v}_c(c)\in G$, $\operatorname{Ad}_{\mathsf{v}_c(c)}$ is an isometry. The first bound of \eqref{9.1:eq-centre-representative-a} then gives $|B^0_g(c)|\le\mathfrak{d}_h+2\varsigma_h$, the products of the small quantities being absorbed in the rounding.

On a junction cell between the collar and the towers, $g=1$ at the tower end, and the same bound follows. On a bond crossing $S_c$ and on the first core cells, the one-cube or two-cube bound above applies, which is at most $\mathfrak{b}_h$.

\textit{Reality and commonness.} Every ingredient of $g$ is real and common to both runs: $V$, $U_\square$, $u_\square$, the conjugators $\xi$ at real fields, the constants $S_\square$, and $U_c$. The variables $\zeta$, $B$, $\mathsf{t}$, $\mathsf{s}$, $\sigma'$ do not occur.

\textit{The bound \eqref{9.1:eq-core-gauge-4}.} Put $\bar\varepsilon:=\varepsilon_m\vee\varepsilon_{m+1}$. By \eqref{9.1:eq-core-gauge-2} and $R_{m+1}\le R_m$,
\begin{equation*}
2(a_m+a_{m+1})\le 8dLM_2R_m\,c_\chi\varpi\bar\varepsilon .
\end{equation*}
Using also $2n_{\mathfrak{K}}+1\le 3n_{\mathfrak{K}}$,
\begin{equation*}
4d(2n_{\mathfrak{K}}+1)LM_2(R_ma_m+R_{m+1}a_{m+1})
\le 48d^2n_{\mathfrak{K}}(LM_2)^2R_m^2\,c_\chi\varpi\bar\varepsilon .
\end{equation*}
Moreover $8dLM_2R_m\le 2d^2n_{\mathfrak{K}}(LM_2)^2R_m^2$, since $n_{\mathfrak{K}}\ge 3^d$ and the side $LM_2R_m$ of the cubes is at least $1$. Adding these to $\mathfrak{d}'_m$ and, for $m=h$, to $\varsigma_h$ gives \eqref{9.1:eq-core-gauge-4}.
\end{proof}

\begin{remark}[Role of the gauge frame and size of the core increments]\label{9.1:rem-frame-size}
We use the objects and constants of Notation~\ref{9.1:not-core-units} and of
Lemma~\ref{9.1:lem-core-gauge}: the datum $D^0(c)$ on a cell-bond $c$, the data coordinate
$B^0_g(c)$, the zone field sizes $a_m$ (defined through the constant $c_\chi$ and the
small-field sizes $\varepsilon_m$), the cube count $n_{\mathfrak{K}}$, the sizes
$\mathfrak{q}_\lambda$, the collar gauge drift $\varsigma_h$ (in which the constant $K_h$
enters), and the core increments $\mathfrak{b}_m$, $m\le h$.
\begin{itemize}
\item[(a)] \emph{Where the frame enters.} The gauge frame $\mathrm{fr}$ enters only in the
statement that the actual collar datum $D^0$ is within $\mathfrak{q}_h$ of $1$; this is the
bound of Lemma~\ref{9.1:lem-core-gauge} on the cells of level $\lambda$ outside the collar, taken
at $\lambda=h$. In the axial--Landau frame $\mathrm{AL}$ the datum $D^0(c)$ is the value at $c$
of the $h$-fold block average of Definition~\ref{1.5:def-block-average}, applied to the
exponentiated field of that frame. In the axial frame this
statement fails at the junction surface $S_c$, and the condition on the core stated below in
this section (which contains the ceiling $a_m\le c_4$, $m\le h$, assumed in
Lemma~\ref{9.1:lem-core-gauge}) cannot be met. Every frame $\mathrm{fr}$ of the admissible class
$\mathfrak{F}_G$ gives the same datum $D^0$ on the collar cells adjacent to $S_c$ and the same
operation $\mathbf{T}'$; everything is therefore stated for the frame $\mathrm{AL}$, as in
\cite{B-CMP122a}.
\item[(b)] \emph{Size.} For every zone $m\le h$ of the core,
\begin{equation}\label{9.1:eq-frame-size-1}
\frac{\mathfrak{b}_m}{\alpha_{0,m}}
\le \frac{50d^2 n_{\mathfrak{K}}(LM_2)^2R_m^2c_\chi(\varepsilon_m\vee\varepsilon_{m+1})}{\alpha_{0,m}}
+\frac{\mathfrak{d}'_m}{\alpha_{0,m}}+\mathbf{1}_{\{m=h\}}\,\frac{\varsigma_h}{\alpha_{0,h}} .
\end{equation}
The first term on the right is at most a constant multiple of
$(\log x_m)^{2r_{pr}+p_0-q}$, and the condition on the exponent $q$ in
\eqref{9.1:eq-core-units-a} gives
\begin{equation}\label{9.1:eq-frame-size-2}
2r_{pr}+p_0-q\le -1-3r_{pr}<0 .
\end{equation}
With $k-m\ge R_k$, the second term satisfies
\begin{equation}\label{9.1:eq-frame-size-3}
\begin{split}
\frac{\mathfrak{d}'_m}{\alpha_{0,m}}
&=34d\bigl(K_0+2B^{\mathsf{s}}C_{sc_1}K_\Lambda'\bigr)(1+\mathsf{r}_1)
\frac{\alpha_{0,k}}{\alpha_{0,m}}L^{-(k-m)}\\
&\le 77d\bigl(K_0+2B^{\mathsf{s}}C_{sc_1}K_\Lambda'\bigr)(1+\mathsf{r}_1)(k-m)^{1/2}L^{-(k-m)},
\end{split}
\end{equation}
and the third satisfies
\begin{equation}\label{9.1:eq-frame-size-4}
\frac{\varsigma_h}{\alpha_{0,h}}\le C(\log x_h)^{2r_{pr}+p_0-q}
+CK_0(1+\mathsf{r}_1)R_k^{1/2}R_hL^{-R_k},
\end{equation}
where $C$ is a constant; the two terms are, respectively, of negative degree in $\log x_h$ and
exponentially small in $R_k$. No constant in
\eqref{9.1:eq-frame-size-1}--\eqref{9.1:eq-frame-size-4}, $C$ included, depends on $R_k$, $h$,
$k$ or on ages;
the radii $R_m$, $R_h$, $R_k$ enter only through the displayed quantities, which stand in the
left-hand sides of the conditions on the core in which they are used.
\end{itemize}
\end{remark}

\begin{proof}
(a) In the axial frame the datum $D^0(c)$ is the value at $c$ of the same averaging operation
applied to the configuration $U_0$ written in that frame. The holonomy lower bound, applied on
its plateau, states two things: the logarithm of this datum is the axial holonomy, and the
modulus of that logarithm is at least $\alpha_{0,k}$ times a fixed positive factor times the
excess of the plateau parameter, which is at least $M$, over a fixed constant.
Then $\mathfrak{b}_h$ is at least one half of that factor times $M\alpha_{0,k}$,
less a small term, so the condition on the
core cannot be met: the statement that the collar datum is within $\mathfrak{q}_h$ of $1$ fails
at the junction surface $S_c$ in that frame. For the frame class $\mathfrak{F}_G$, the statement
on admissible frames fixes the boundary datum of the frame on the outer boundary that it
specifies, independently of the choice of frame in $\mathfrak{F}_G$; hence every
$\mathrm{fr}\in\mathfrak{F}_G$ yields the same datum $D^0$ on the collar cells adjacent to $S_c$,
and the same operation $\mathbf{T}'$. This is why the frame $\mathrm{AL}$ may be used
throughout, as in \cite{B-CMP122a}.

(b) Lemma~\ref{9.1:lem-core-gauge} bounds $\mathfrak{b}_m$ by
$\mathfrak{d}'_m+\mathbf{1}_{\{m=h\}}\varsigma_h$ plus the product of the live-piece weight of
that lemma, which equals $1$ off live pieces, with
\begin{equation*}
50d^2n_{\mathfrak{K}}(LM_2)^2R_m^2c_\chi(\varepsilon_m\vee\varepsilon_{m+1}) .
\end{equation*}
Dividing by $\alpha_{0,m}$ (with
$\alpha_{0,m}=\alpha_{0,h}$ when $m=h$), the live-piece weight cancels on live pieces against the
same weight in the live-piece unit of Notation~\ref{9.1:not-core-units}; this gives
\eqref{9.1:eq-frame-size-1}.

For the first term we use the comparison of $\varepsilon_{m+1}$ with $\varepsilon_m$ in
\cite[(2.8)]{B-CMP119}, which contributes the factor $(1+\beta_0^{119})^2$, with
$\beta_0^{119}$ the constant of \cite[(2.8)]{B-CMP119}. By the radii and units of
Notation~\ref{9.1:not-core-units}, the first term is at most
\begin{equation*}
C(\log x_m)^{2r_{pr}}\cdot\frac{A_0}{C_0}(\log x_m)^{p_0-q}\cdot(1+\beta_0^{119})^2 ,
\end{equation*}
with $C$ a constant. Its degree in $\log x_m$ is $2r_{pr}+p_0-q$. Under the condition on the
exponent $q$ in \eqref{9.1:eq-core-units-a} this degree obeys \eqref{9.1:eq-frame-size-2}.

For the second term, the bound $\mathfrak{d}'_m$ gives the equality in
\eqref{9.1:eq-frame-size-3}, in which $\tilde\alpha=(1+\mathsf{r}_1)\alpha_{0,k}$ has been used.
The inequality in \eqref{9.1:eq-frame-size-3} is the bound
$\alpha_{0,k}/\alpha_{0,m}\le\frac{77}{34}(k-m)^{1/2}$ on the units $\alpha_{0,\cdot}$ of
Notation~\ref{9.1:not-core-units}, for $k-m\ge R_k$.

For the third term, insert the definitions of $a_h$ and $\mathfrak{q}_h$ into that of
$\varsigma_h$ (all three in Lemma~\ref{9.1:lem-core-gauge}), and use the bound
$K_h\le 100\sqrt{d}\,MLR_h$ that enters the cube count $n_{\mathfrak{K}}$. Then
$\varsigma_h/\alpha_{0,h}$ is bounded by a sum, with constant coefficients, of terms
$R_h^2(\varepsilon_h/\alpha_{0,h})$ and $K_hK_0(\tilde\alpha/\alpha_{0,h})L^{-(k-h)}$. The first
is bounded as the first term above. For the second,
$\tilde\alpha/\alpha_{0,h}=(1+\mathsf{r}_1)\alpha_{0,k}/\alpha_{0,h}$; with $k-h\ge R_k$, the
bound on $\alpha_{0,k}/\alpha_{0,m}$ above taken at $m=h$, the bound on $K_h$, and the decrease
of $t\mapsto t^{1/2}L^{-t}$ for $t\ge 1$ bound it by
$CK_0(1+\mathsf{r}_1)R_k^{1/2}R_hL^{-R_k}$. This is \eqref{9.1:eq-frame-size-4}.

The constants $C$, $d$, $n_{\mathfrak{K}}$, $L$, $M_2$, $c_\chi$, $K_0$, $B^{\mathsf{s}}$,
$C_{sc_1}$, $K_\Lambda'$, $\mathsf{r}_1$, $A_0$, $C_0$, $\beta_0^{119}$ appearing here depend on
none of $R_k$, $h$, $k$ or the ages. The radii $R_m$, $R_h$, $R_k$ enter the bounds only where
displayed, never through the constants.
\end{proof}

In what follows we fix the objects on which the increments are read.

Let $\zeta=e^{i\eta A_e}U_c$ be a seed in the chart of the shape, with
\begin{equation*}
|A_e|\le\tilde{\alpha}/\ell_k\quad\text{on } Z .
\end{equation*}
This is one of the bounds that define the data of the seed tube $\mathfrak{T}_\Phi[7/8]$, the units of the $\mathbf{R}$-born output $\Phi$ being those of level $k$ on $Z$. Off $Z$ the units are those of the label, and only cells in $Z$ carry non-zero increments below.

Let $\mathsf{s}$ be a vertex of the decoupling vertex set $\mathfrak{P}_{\mathsf{c}}$, so that $|\mathsf{s}(\Delta)|\le e^{\kappa_1}$ for every cube $\Delta$. The decoupling bound for the prepared data carries the factor $B^{\mathsf{s}}$ on the sources.

Let $B$ be a variable of the $\mathbf{T}$-step in the support of $\chi'$, where, by \cite[(1.82)]{B-CMP122a}, $\chi'$ is the characteristic function of the set
\begin{equation*}
\{|B'(b)|<\delta'_k \text{ for } b\in\mathfrak{B}_0\}.
\end{equation*}
Here $B'$ is the variable of \cite[(1.82)]{B-CMP122a}, which is the present $B$.
Let $B^{(\mathsf{s})}$ be its dilation. By the bound on the dilated $\mathbf{T}$-step variable in the preparation of the data, together with \cite[(1.82)]{B-CMP122a},
\begin{equation}\label{9.1:eq-data-increments-1}
|g_kB^{(\mathsf{s})}(c)|\le B^{\mathsf{s}}C_1\delta'_k\quad\text{on }\mathfrak{B}_0 .
\end{equation}
Let $\sigma'\in[0,1]^{\mathrm{cubes}}$ (one localisation parameter per cube). Let $\mathsf{t}$ lie on the bracket circle of \eqref{9.1:eq-core-ceilings-a}.

The objects $P$ read, and the sets $\mathfrak{S}_P$ on which their data are read, are the following.
\begin{itemize}
\item For $P\in\{U^0_k(\zeta),\,U''_k(\zeta,B)\}$ put $\mathfrak{S}_P:=\mathbf{B}''_k{\upharpoonright}R$.
\item For $P=\Xi_1(\sigma')$, the minimiser on $\mathbf{B}^{\sim}(Y_1)$, put $\mathfrak{S}_P:=\mathbf{B}^{\sim}(Y_1){\upharpoonright}R$.
\end{itemize}
We write $P^{\mathrm{re}}:=U^{0,\mathrm{re}}:=U^0_k(U_c;B=0)$ for the real object of Lemma~\ref{9.1:lem-core-gauge}. In the notation of that lemma, $D^0(c)$ is its datum and $B^0_g(c)$ is its data coordinate in a block-centre gauge $g$.

For each of the objects $P$, $P^{\mathrm{re}}$, $q_0$ and a cell $c$ of $\mathfrak{S}_P$ with ends $c_-$, $c_+$, the data coordinate about $U_c$ is
\begin{equation*}
B[\cdot](c)=\tfrac1i\log\bigl[g(c_-)\,(\text{datum on }c)\,g(c_+)^{-1}\,\mathsf{v}_c(c)^{-1}\bigr],
\end{equation*}
where $g$ is a block-centre gauge (the representative). The representative is free, and we fix it as follows.
\begin{itemize}
\item For $P$ and for $P^{\mathrm{re}}$ we use the real gauge $g$ of Lemma~\ref{9.1:lem-core-gauge}. It equals $1$ at every centre outside $\mathfrak{K}\cup\mathfrak{K}_1$, and by construction $B[P^{\mathrm{re}}]=B^0_g$.
\item For the seed object $q_0$ (Definition~\ref{9.1:def-seed-criticality}) we use the raw representative $g\equiv1$.
\end{itemize}
Writing $M^{\lambda(c)}(\cdot)(c)$ for the block average on a cell $c$ of level $\lambda(c)$ (Ba{\l}aban's averaging operator, not the constant $M$), this gives
\begin{equation}\label{9.1:eq-data-increments-2}
B[q_0](c)=\tfrac1i\log\bigl[M^{\lambda(c)}(q_0)(c)\,\mathsf{v}_c(c)^{-1}\bigr]=:\beta''_e(c).
\end{equation}
On cells of level $\lambda(c)$ inside $Z$ it satisfies
\begin{equation}\label{9.1:eq-data-increments-3}
|\beta''_e(c)|\le34d\,|A_e|_c\le34d\,\tilde{\alpha}L^{\lambda(c)-k}.
\end{equation}
This follows from the averaging lemma about a regular base, applied at the base $U_c$ with the field $A_e$, where $|A_e|\le\tilde{\alpha}/\ell_k$ on $Z$. The element $\mathsf{v}_c(c)\in G$ acts isometrically under $\mathrm{Ad}$. Finally we put
\begin{equation*}
\mathsf{b}_P(c):=|B[P](c)-B[q_0](c)|.
\end{equation*}

\begin{lemma}[Data increments of the read objects against the seed]\label{9.1:lem-data-increments}
Assume the hypotheses of Lemma~\ref{9.1:lem-core-gauge}, the decoupling bound and the chain-displacement statement. Put
\begin{equation*}
K_{\mathfrak{K}}:=80d\bigl(K_0+K^{\flat}_\Lambda+B^{\mathsf{s}}C_{sc_1}K'_\Lambda+1\bigr).
\end{equation*}
Here $C_{sc_1}$ is the constant in the bound $\mathfrak{d}_\lambda$ of \eqref{9.1:eq-centre-representative-a}. Then the following hold.
\begin{equation}\label{9.1:eq-data-increments-a}
\begin{aligned}
\text{(a)}\ \ &\mathsf{b}_P(c)=0 \qquad\text{on $\mathbf{B}''_k$-cells $c$ off $\Lambda$;}\\
\text{(b)}\ \ &\mathsf{b}_P(c)\le\mathbb{1}_{P\neq U^0_k}B^{\mathsf{s}}C_1\delta'_k
 +2\max\bigl\{34dK_0\tilde{\alpha}L^{\lambda-k},\,2K^{\flat}_\Lambda\tilde{\alpha}\mathbb{1}_{\lambda=k}\bigr\}\\
&\qquad\quad+\mathfrak{d}_\lambda+34d\tilde{\alpha}L^{\lambda-k}+\mathbb{1}_{\lambda=h}2\varsigma_h\\
&\qquad\le B^{\mathsf{s}}C_1\delta'_k+K_{\mathfrak{K}}\tilde{\alpha}L^{\lambda-k}
 +\mathbb{1}_{\lambda=h}2\varsigma_h
 \qquad\text{on $\mathfrak{B}_0$-cells $c$ of level $\lambda$;}\\
\text{(c)}\ \ &\mathsf{b}_P(c)\le\mathfrak{b}_m+34d\tilde{\alpha}L^{m-k}\le2\mathfrak{b}_m
 \qquad\text{on the cells listed in item (c) below (zone $m$);}\\
\text{(d)}\ \ &\mathsf{b}_P(c)\le4|\mathcal{Z}^{\mathrm{ch}}_{U^0_k}|_c
 +\bigl(\text{the $\sigma'$-decoupling increment at $c$}\bigr)\\
&\qquad\le8\ell_c\,\mathsf{m}_L(c)/\ell_{\iota(c)}\qquad\text{on the layer cells of (d) below.}
\end{aligned}
\end{equation}
The clauses are made precise as follows.
\begin{itemize}
\item[(a)] On these cells $P$ and $q_0$ have the same datum and the same representative.
\item[(b)] On the rind, on the set $\mathfrak{i}$ and on the bonds crossing $\partial\Lambda$ the level is $\lambda=k$; here $\mathfrak{i}$ is the family of $\mathfrak{B}_0$-cells of level $k$ listed, with the rind and the $\partial\Lambda$-crossing bonds, among the $\mathfrak{B}_0$-cells outside the collar in Lemma~\ref{9.1:lem-core-gauge}. On the towers the level is the tower level; on the collar $\lambda=h$. The exponent is $0$ at $\lambda=k$, and products of small terms are absorbed in the rounding.
\item[(c)] The cells are the cells of zone $m$ of $\mathfrak{K}_2$, the cells crossing the junction surface $S_c$, and the zone junctions. The bound is independent of $P$, $\zeta$, $B$, $\mathsf{s}$. The second inequality holds because $34d\tilde{\alpha}L^{m-k}\le\mathfrak{d}'_m\le\mathfrak{b}_m$.
\item[(d)] Here $P=\Xi_1(\sigma')$. The layer cells of $\mathbf{B}(Y_1)$ lie at distance at least $M-3M_1$ from $\Lambda$ by the layer geometry. The chain field $\mathcal{Z}^{\mathrm{ch}}_{U^0_k}$ is the displacement of the chain object from the seed at the layer, as in the chain-displacement statement; the $\sigma'$-decoupling increment is that of the decoupling bound; the constants $C_4$, $\delta_1$, the set $\bar{\Lambda}$ and the scales $\ell_c$, $\ell_{\iota(c)}$ are those of the chain-displacement statement; $\operatorname{dist}(\cdot,\bar{\Lambda})$ denotes the distance to $\bar{\Lambda}$; and
\begin{equation*}
\mathsf{m}_L=C_4\mathsf{s}_L\,e^{-\delta_1\operatorname{dist}(\cdot,\bar{\Lambda})/4}.
\end{equation*}
\item[(e)] Under the coupling ceilings \eqref{9.1:eq-core-ceilings-a}, all representatives lie in the $a_{R'}$-ball: $|B[q_0]|\le r_c$ and $\mathsf{b}_P\le10^{-1}a_{R'}$.
\end{itemize}
\end{lemma}

\begin{proof}
\emph{Clause (a).} Off $\Lambda$ the exterior data of $P$ are the block averages of $\zeta$, and these are those of $q_0$. This is the exterior part of \eqref{9.1:eq-core-data-a}, which is Ba{\l}aban's identification of the data outside $\Lambda$ \cite{B-CMP122b}, together with the exterior-data identity of the averaging lemma. So $P$ and $q_0$ have the same datum on every $\mathbf{B}''_k$-cell off $\Lambda$. They also have the same representative there: the representative is $g=1$ for both on these cells. Hence $B[P](c)=B[q_0](c)$ and the difference is $0$.

\emph{Clause (b).} Since $B[P^{\mathrm{re}}]=B^0_g$, we have
\begin{equation}\label{9.1:eq-data-increments-4}
B[P]-B[q_0]=\bigl(B[P]-B[P^{\mathrm{re}}]\bigr)+B^0_g-B[q_0],
\end{equation}
and we bound the three terms separately.

\emph{First term.} The datum of $P$ on $c$ is $e^{ig_kB^{(\mathsf{s})}(c)}V_0^{\mathrm{AL}}(\zeta)(c)$. For $P=U^0_k(\zeta)$, which is taken at $B=0$, the factor $e^{ig_kB^{(\mathsf{s})}(c)}$ equals $1$; this is the indicator $\mathbb{1}_{P\neq U^0_k}$. The datum of $P^{\mathrm{re}}$ is $V_0^{\mathrm{AL}}(U_c)(c)$. Both are taken in the same frame $\mathrm{AL}$ with the same representative $g$, so the ratio is
\begin{equation*}
e^{ig_kB^{(\mathsf{s})}(c)}\cdot V_0^{\mathrm{AL}}(\zeta)(c)\cdot V_0^{\mathrm{AL}}(U_c)(c)^{-1}.
\end{equation*}
By parts (i) and (ii) of \eqref{9.1:eq-core-data-a}, at every point of the seed tube each factor $V_0^{\mathrm{AL}}(\cdot)(c)$ is close to $1$:
\begin{itemize}
\item within $34dK_0\tilde{\alpha}L^{\lambda-k}$ on tower and collar bonds;
\item within $2K^{\flat}_\Lambda\tilde{\alpha}$ on bonds of level $k$.
\end{itemize}
Together with \eqref{9.1:eq-data-increments-1}, this gives
\begin{equation*}
|\log(\text{ratio})|\le\mathbb{1}_{P\neq U^0_k}B^{\mathsf{s}}C_1\delta'_k
+2\max\bigl\{34dK_0\tilde{\alpha}L^{\lambda-k},\,2K^{\flat}_\Lambda\tilde{\alpha}\mathbb{1}_{\lambda=k}\bigr\}
\end{equation*}
up to products of these small terms. The phrase ``same representative'' implies a conjugation by the real element $g(c_-)$. Off the collar $g=1$, so the conjugation is an isometry. On the collar $|g-1|\le\varsigma_h$ by the collar estimate in \eqref{9.1:eq-core-gauge-a}, and the conjugation costs at most $2\varsigma_h$.

\emph{Second term.} We use the estimates of Lemma~\ref{9.1:lem-core-gauge} collected in \eqref{9.1:eq-core-gauge-a}.
\begin{itemize}
\item On $\mathfrak{B}_0$-cells of level $\lambda$ outside the collar (the rind, $\mathfrak{i}$, the bonds crossing $\partial\Lambda$, the towers), $g=1$ and $|B^0_g(c)|\le\mathfrak{d}_\lambda$. This is \eqref{9.1:eq-centre-representative-a}.
\item On collar cells and on the junction cells of the collar with the towers, $|B^0_g(c)|\le\mathfrak{d}_h+2\varsigma_h$. In the sum the term $2\varsigma_h$ is counted once.
\end{itemize}

\emph{Third term.} We have $|B[q_0](c)|=|\beta''_e(c)|\le34d\tilde{\alpha}L^{\lambda-k}$ by \eqref{9.1:eq-data-increments-3}.

Adding the three bounds gives the first inequality of (b). For the second inequality:
\begin{itemize}
\item $\mathbb{1}_{P\neq U^0_k}\le1$.
\item $2\max\{x,y\}\le2x+2y$. Moreover $\mathbb{1}_{\lambda=k}\neq0$ only where $L^{\lambda-k}=1$.
\item By \eqref{9.1:eq-centre-representative-a}, $\mathfrak{d}_\lambda=68dB^{\mathsf{s}}C_{sc_1}\mathsf{s}_LL^{\lambda-k}$, where the seed size is $\mathsf{s}_L=K'_\Lambda\tilde{\alpha}$.
\end{itemize}
Hence everything except $\mathbb{1}_{P\neq U^0_k}B^{\mathsf{s}}C_1\delta'_k$ and $\mathbb{1}_{\lambda=h}2\varsigma_h$ is at most
\begin{equation*}
\bigl(68dK_0+4K^{\flat}_\Lambda+68dB^{\mathsf{s}}C_{sc_1}K'_\Lambda+34d\bigr)\tilde{\alpha}L^{\lambda-k}
\le K_{\mathfrak{K}}\tilde{\alpha}L^{\lambda-k}.
\end{equation*}
The last inequality holds because $d\ge1$ and all constants are non-negative. Products of small terms are absorbed in the rounding.

\emph{Clause (c).} On these cells the datum of $P$ is the common real core variable $V$, for every $P$, $\zeta$, $B$, $\mathsf{s}$. This follows from \eqref{9.1:eq-core-data-a} and the common-core statement. The dilation $\mathsf{s}$ and the variable $B$ enter only through the variables on $\mathfrak{B}_0$ \cite{B-CMP122a}, so they leave the core variables unchanged. So the datum of $P$ equals that of $P^{\mathrm{re}}$, with the same representative $g$.

In \eqref{9.1:eq-data-increments-4} the first term is therefore exactly $0$. The second is at most $\mathfrak{b}_m$ by the core estimate of Lemma~\ref{9.1:lem-core-gauge} in \eqref{9.1:eq-core-gauge-a}. The third is at most $34d\tilde{\alpha}L^{m-k}$ by \eqref{9.1:eq-data-increments-3} with $\lambda=m$. None of these bounds depends on $P$, $\zeta$, $B$, $\mathsf{s}$, since $g$ is built from the common real data and $U_c$ only.

For the second form, $34d\tilde{\alpha}L^{m-k}\le\mathfrak{d}'_m$ with $\mathfrak{d}'_m$ the bound of \eqref{9.1:eq-centre-representative-a}. Also $\mathfrak{d}'_m\le\mathfrak{b}_m$, because $\mathfrak{b}_m$ is $\mathfrak{d}'_m$ plus non-negative terms by \eqref{9.1:eq-core-gauge-a}.

\emph{Clause (d).} The layer data of $\Xi_1(\sigma')$ are reads of the chain object at the layer, localised by $\sigma'$. The localisation is through the localisation kernel of \cite{B-CMP122b}, whose set $\sigma_0$ is disjoint from $\Lambda$. By the chain-displacement statement, which holds at complex seeds, the chain object at the layer is $e^{i\eta\mathcal{Z}^{\mathrm{ch}}_{U^0_k}}q_0$. By the averaging lemma's bound for a configuration shifted by a field, its reads therefore differ from those of $q_0$ by at most $4|\mathcal{Z}^{\mathrm{ch}}_{U^0_k}|_c$.

The $\sigma'$-interpolation moves the reads by at most as much again, by the decoupling bound for kernels decoupled in $\mathsf{s}$. So the $\sigma'$-decoupling increment at $c$ is at most $4|\mathcal{Z}^{\mathrm{ch}}_{U^0_k}|_c$. By the chain-displacement statement (its bounds for the layer reads of the chain), the displacement of the chain from the seed at the layer obeys $|\mathcal{Z}^{\mathrm{ch}}_{U^0_k}|_c\le\ell_c\mathsf{m}_L(c)/\ell_{\iota(c)}$. Together these give the bound $8\ell_c\mathsf{m}_L(c)/\ell_{\iota(c)}$.

\emph{Clause (e).} By the definition of the seed tube, $|B[q_0]|\le r_c$. By (a)--(d), under the ceilings \eqref{9.1:eq-core-ceilings-a}, every other representative differs from $B[q_0]$ by $\mathsf{b}_P\le10^{-1}a_{R'}$.

\emph{What does not enter.} The bounds \cite[(1.54)]{B-CMP122b} and \cite[(1.57)]{B-CMP122b} of Ba{\l}aban's construction (the first of them carries the localisation kernel used in clause (d)), and the configurations $U_1$, $U_2$ and the cut of Ba{\l}aban's construction, do not enter the proof above. One further bound of Ba{\l}aban's localisation in \cite{B-CMP122b} enters, but only at the real seed, inside the size constant $a_\Lambda$ (the choice of the centre's representative $\mathring{U}_c$). The complex seed enters the proof above only through Ba{\l}aban's bound $|A_0(\zeta)|\le K_0\tilde{\alpha}/\ell_k$ \cite{B-CMP122a}, where $A_0(\zeta)$ is the field of his construction attached to $\zeta$ (a field, not the auxiliary constant $A_0$), and through $|A_e|\le\tilde{\alpha}/\ell_k$. The complex variable $B$ enters only through $\chi'$ and the factor $B^{\mathsf{s}}$.
\end{proof}

\begin{corollary}[Every read object as a small rotation of the seed]\label{9.1:cor-seed-rotation}
Under the hypotheses and in the setting of Lemma~\ref{9.1:lem-data-increments}, let
$q\in\mathfrak{Q}''$ be an object at which a first-part term is read, written as
$q=e^{i\eta\mathcal{Y}_{lk}}P$. Here $P$ is one of the configurations whose data increments
$\mathsf{b}_P$ against the seed $q_0$ are bounded in Lemma~\ref{9.1:lem-data-increments}, and
$\mathfrak{S}_P$ is the set over which $P$ is a chart point. The field $\mathcal{Y}_{lk}$ is the
composition of the link moves, namely
\begin{itemize}
\item the $\sigma'$-decoupling differences on $\Lambda$;
\item the $\mathsf{t}$-bracket: $\mathsf{t}$ times the weight of the cube $\square$ in
Ba{\l}aban's partition of unity, applied to the operator $\mathbb{H}^{(0')}(\mathsf{s})$ of the
leaf-reading theorem at dilation $\mathsf{s}$;
\item the dilation corrections;
\item the box re-presentations,
\end{itemize}
in the form fixed by the three link displays of the leaf-reading theorem and by item~(b) of the
chain-field bounds.
Then the following hold.
\begin{enumerate}
\item[(a)] On the seed's region $R$ (Definition~\ref{9.1:def-seed-criticality}) we have
$q=e^{i\eta\mathcal{Y}_q}q_0$. The field $\mathcal{Y}_q$ is holomorphic in
$(\zeta,B,\mathsf{t},\mathsf{s},\sigma')$. Up to the layer reads bounded in the layer item of
Lemma~\ref{9.1:lem-data-increments}, $\mathcal{Y}_q=0$ at
$(\zeta,B,\mathsf{t},\mathsf{s},\sigma')=(U_c,0,0,\mathbb{1},1)$.
\item[(b)] For every cell $y$ of $R$,
\begin{equation}\label{9.1:eq-seed-rotation-1}
\ell_\iota(y)\sup_{\widetilde{\Delta}(y)}|\mathcal{Y}_q|\le\mathsf{y}_q(y).
\end{equation}
\item[(c)] On cells $c\subset\bar\Lambda^{+}$,
\begin{equation}\label{9.1:eq-seed-rotation-2}
\ell_\iota\sup_{\widetilde{\Delta}(c)}|\mathcal{Y}_{lk}|\le\mathsf{m}_{br}(c)+\mathsf{m}_{lk}(c).
\end{equation}
\end{enumerate}
In (b) and (c) the majorants are
\begin{equation}\label{9.1:eq-seed-rotation-a}
\begin{split}
\mathsf{y}_q(y)&:=1.1\,B_\natural(\mathsf{b}_P)^{\approx}(y)
+\ell_\iota(y)\sup_{\widetilde{\Delta}(y)}|\mathcal{Y}_{lk}|,\\
\mathsf{m}_{br}(c)&:=(1+|\mathsf{t}|)B^{\mathsf{s}}C_{sc_1}C_1\delta'_k\,
e^{-\delta_1 d(c,\mathfrak{B}_0\cap\widetilde{\square})/4},\\
\mathsf{m}_{lk}(c)&:=8B^{\mathsf{s}}C_{sc_1}C_4\mathsf{s}_L\,e^{-\delta_1 d(c,(\bar\Lambda)^c)/4}.
\end{split}
\end{equation}
Here $(\mathsf{b}_P)^{\approx}$ is the smeared majorant of $\mathsf{b}_P$, taken over the cells
of $\mathfrak{S}_P$; $\ell_\iota(y)$ and $\widetilde{\Delta}(y)$ are the length scale and the
reading neighbourhood attached to the cell $y$ in the leaf-reading theorem;
$\widetilde{\square}$ is the bracket cube of the first-part term read; and $C_{sc_1}$ is the
response constant of the decay of responses.

The profile $\mathsf{m}_{br}$ comes from $g_kB$ on $\mathfrak{B}_0\cap\widetilde{\square}$, which
drives the bracket and the dilations. The profile $\mathsf{m}_{lk}$ comes from the
$\sigma_0$-cubes disjoint with $\Lambda$ and from the far layers, which drive the
$\sigma'$-links. The second profile decays into $\Lambda$ from its complement. Consequently at
$\mathfrak{K}$ it carries the factor $L^{-4(\nu_{\mathfrak{K}}-1)}$, where $\nu_{\mathfrak{K}}$
is the integer attached to the core $\mathfrak{K}$ by its zone structure, and at a tower
cell of level $\pi$ (an integer level) it carries the factor $L^{-4(k-\pi-1)_+}$.

Box vertices are not links. A box object on $Z\supseteq X^{+}$ is the slice on
$\boldsymbol{\Omega}(Z)$ of its parent, the slice being taken at
$\mathfrak{S}:=\boldsymbol{\Omega}(Z)$ as in the last row of the table of the leaf-reading
theorem. It is therefore read through the same $\mathcal{Y}_q$.
\end{corollary}

\begin{proof}
\emph{Chart points.} Write $P=\mathbf{V}_{\mathfrak{S}_P}(B[P])$ and
$q_0=\mathbf{V}_{\mathfrak{S}_P}(B[q_0])$. Both are chart points over $U_c$ on $\mathfrak{S}_P$,
for the following reasons.
\begin{itemize}
\item By Definition~\ref{9.1:def-seed-criticality}, both $P$ and $q_0$ are critical for
$\mathfrak{S}_P$, and $\mathfrak{S}_P$ is a refinement of $\boldsymbol{\Omega}(R)$.
\item The critical orbit with given data is unique \cite[Theorem~1]{B-CMP102b}.
\item By the identity theorem for holomorphic functions, the same holds at complex data.
\end{itemize}
Hence, with $\mathfrak{X}$ the Landau coordinate map about $U_c$,
\begin{equation*}
P=e^{i\eta\mathfrak{X}(B[P])}U_c,\qquad
q_0=e^{i\eta\mathfrak{X}(B[q_0])}U_c=e^{i\eta A_e}U_c .
\end{equation*}
Put $E_P:=\mathfrak{X}(B[P])-\mathfrak{X}(B[q_0])$. Then $P=e^{i\eta\mathcal{Y}_P}q_0$, where
$\mathcal{Y}_P$ is given by the Baker--Campbell--Hausdorff series:
\begin{equation*}
\mathcal{Y}_P:=\operatorname{BCH}\bigl(\mathfrak{X}(B[P]),-\mathfrak{X}(B[q_0])\bigr)
=E_P+\tfrac{\eta}{2}[E_P,A_e]+\cdots .
\end{equation*}

\emph{The Landau field of $P$ against the seed.} We apply the local Lipschitz property
of the Landau chart about $U_c$ on $\mathfrak{S}_P$. Its
hypotheses hold, for three reasons.
\begin{itemize}
\item $U_c$ is a real regular configuration, critical for $\mathfrak{S}_P$.
\item $|\mathsf{m}_c|\le\mathfrak{r}_0$, with $\mathfrak{r}_0$ the bound of the leaf-reading
theorem.
\item The two data $B^1=B[P]$ and $B^2=B[q_0]$ satisfy $\sup|B^i|\le a_{R'}$. This holds because,
by the last assertion of Lemma~\ref{9.1:lem-data-increments}, all representatives used lie in
the $a_{R'}$-ball.
\end{itemize}
The property gives
\begin{equation*}
\ell_\iota|E_P|\le B_\natural(\mathsf{b}_P)^{\approx},\qquad
\ell_\iota^2|\nabla E_P|\le B_\natural(\mathsf{b}_P)^{\approx}.
\end{equation*}
From the series for $\mathcal{Y}_P$ we obtain, with a constant $C$,
\begin{equation*}
|\mathcal{Y}_P|\le|E_P|\bigl(1+C\ell_\iota|A_e|\bigr)\le1.1\,|E_P| .
\end{equation*}

\emph{The bracket and the dilations.} We bound the link moves by applying the response estimate
that is among the hypotheses of Lemma~\ref{9.1:lem-data-increments}
term by term, each term at its own sources.

For the bracket and the dilations, the source is $g_kB$ restricted to
$\mathfrak{B}_0\cap\widetilde{\square}$. To see this, consider the bracket display of the
leaf-reading theorem. The brackets of its lower rows have $X\subset Z^c$. The bracket cube
$\widetilde{\square}$ of a first-part term, however, may meet $\Lambda$; hence the set
$\mathfrak{B}_0$ enters.

At the source, the size in its own units is at most $(1+|\mathsf{t}|)B^{\mathsf{s}}C_1\delta'_k$.
The response to it is at most $B^{\mathsf{s}}C_{sc_1}$ times this, with the profile given by
the decay of responses, part~(i). This is $\mathsf{m}_{br}$.

\emph{The $\sigma'$-links.} These have their sources at two kinds of place.
\begin{itemize}
\item At the $\sigma_0$-cubes. These are disjoint with $\Lambda$ in each of the four
localisations of \cite{B-CMP122b}, two of which are on \cite[p.~374]{B-CMP122b}.
\item At the layers of $\partial Y_i^0$. These lie at distance at least $M-3M_1$ from $\Lambda$,
by the geometric separation lemma for the layers of $\partial Y_i^0$.
\end{itemize}
Their arguments are at most $C_4\mathsf{s}_L$, by items~(b) and~(c) of the chain-field bounds.
The decay of responses, part~(i), then gives the stated profile $\mathsf{m}_{lk}$. This proves
\eqref{9.1:eq-seed-rotation-2}.

\emph{Composition.} Since $q=e^{i\eta\mathcal{Y}_{lk}}e^{i\eta\mathcal{Y}_P}q_0$, the field
$\mathcal{Y}_q$ is the Baker--Campbell--Hausdorff composite of the field of $P$ and the fields
of the links. All $\ell$-scaled sizes are at most $10^{-1}$, so the composition contributes a
factor at most $1.1$. Together with the two preceding paragraphs this gives
\eqref{9.1:eq-seed-rotation-1}.

\emph{Exactness at the real point.} At $(\zeta,B,\mathsf{t},\mathsf{s},\sigma')=
(U_c,0,0,\mathbb{1},1)$ all increments vanish, with one exception: the reads on the layer cells
bounded in Lemma~\ref{9.1:lem-data-increments}. By the nesting of the $\sigma'$-localised reads,
these vanish at $\sigma'=1$.
\end{proof}

\begin{lemma}[Pointwise moves of the table members from cellwise increments]
\label{9.1:lem-pointwise-moves}
Let $(F,X,j)$ be a term $F$ of the resummation together with its region $X$ and its
truncation level $j$, and let $X^{+}$ be the hull of $X$. Let
$\mathfrak{S}=\boldsymbol{\Omega}(X^{+})$ be the box set of $X^{+}$ over
Ba{\l}aban's determining set $\mathbf{B}''_k$, truncated at the level $j$. Let $q$ be one
of the objects of $\mathfrak{Q}''$ at which the first-part terms are read, written as
$q=e^{i\eta\mathcal{Y}_q}q_0$, in the presentation of
Corollary~\ref{9.1:cor-seed-rotation}, with $P$, $\mathfrak{S}_P$, $\mathsf{y}_q$,
$\mathsf{m}_{br}$ and $\mathsf{m}_{lk}$ as there; $\tilde{\Delta}(y)$ and
$\bar{\Lambda}^{+}$ denote the enlarged cell at $y$ and the region of that corollary.
For a table member $E_a$ with exponent
$e_a$, write $\operatorname{move}_a(x)$ for the move of $E_a$ at $x\in X$. Put
\begin{equation}\label{9.1:eq-pointwise-moves-1}
\sigma(x):=\min\Bigl\{1,\;L\,\frac{\ell_{\mathfrak{S}}}{\ell_\iota}(x)\Bigr\},
\qquad
\mathsf{b}^{\sim}(x):=\sum_{c}e^{-\delta_1 d(x,c)}\,\mathsf{b}_P(c),
\end{equation}
where the sum runs over the cells $c$ of $\mathfrak{S}_P$ and $\mathsf{b}_P(c)$ are the
data increments of Lemma~\ref{9.1:lem-data-increments}. Then, for every $x\in X$,
\begin{equation}\label{9.1:eq-pointwise-moves-a}
\begin{split}
\operatorname{move}_a(x)
&\le 18B_{\natural}^{2}c_1\,\sigma(x)\,\ell_{\mathfrak{S}}(x)^{-e_a}\,
\mathsf{b}^{\sim}(x)\\
&\quad+16B_{\natural}c_1C_T\,\sigma(x)\,\ell_{\mathfrak{S}}(x)^{-e_a}\,
(\mathsf{m}_{br}+\mathsf{m}_{lk})(x).
\end{split}
\end{equation}
Here $\ell_{\mathfrak{S}}$ and $\ell_\iota$ are the length functions formed from the level
functions $\lambda_{\mathfrak{S}}$ and $\iota$, $\delta_1$ is the decay rate of the cell
weights, $c_1$ is the constant of the summation estimate and $C_T$ that of the transfer
estimate for slowly decaying profiles, both among the random-walk estimates.
\end{lemma}

\begin{proof}
\emph{The move bound at the seed.} The entry $q_0$ of the presentation of $q$ is critical
for $\mathfrak{S}$ in the sense of Definition~\ref{9.1:def-seed-criticality}, by the
criticality statement for the seed and its box set. Hence
part~(a) of the move bound for the table members applies at $q_0$, and its part~(b$'$)
gives, for $x\in X$,
\begin{equation}\label{9.1:eq-pointwise-moves-2}
\operatorname{move}_a(x)\le 16B_{\natural}\,\ell_{\mathfrak{S}}(x)^{-e_a}\,
\bigl(|\mathcal{Y}_q|_{\cdot}\bigr)^{\approx}(x),
\qquad
|\mathcal{Y}_q|_y:=\ell_{\mathfrak{S}}(y)\sup_{\tilde{\Delta}(y)}|\mathcal{Y}_q|,
\end{equation}
where $(\cdot)^{\approx}$ is the smearing with the weight $e^{-2\delta_1 d}$ over the
cells of $\mathfrak{S}$.

\emph{Passage to the majorant.} Corollary~\ref{9.1:cor-seed-rotation} gives
$\ell_\iota(y)\sup_{\tilde{\Delta}(y)}|\mathcal{Y}_q|\le\mathsf{y}_q(y)$;
multiplying by $\ell_{\mathfrak{S}}(y)/\ell_\iota(y)$ gives
\begin{equation}\label{9.1:eq-pointwise-moves-3}
|\mathcal{Y}_q|_y\le\frac{\ell_{\mathfrak{S}}}{\ell_\iota}(y)\,\mathsf{y}_q(y),
\qquad
\mathsf{y}_q(y)=1.1B_{\natural}(\mathsf{b}_P)^{\approx}(y)
+\ell_\iota(y)\sup_{\tilde{\Delta}(y)}|\mathcal{Y}_{lk}|,
\end{equation}
the second identity being the definition of $\mathsf{y}_q$ in
\eqref{9.1:eq-seed-rotation-a}, with $(\mathsf{b}_P)^{\approx}$ taken over the cells of
$\mathfrak{S}_P$.

\emph{Transport of the ratio.} Since $\lambda_{\mathfrak{S}}$ and $\iota$ are level
functions, by the floor condition on the level functions, the comparison of level
functions at two points among the random-walk
estimates gives
\begin{equation}\label{9.1:eq-pointwise-moves-4}
\frac{\ell_{\mathfrak{S}}}{\ell_\iota}(y)
\le L\,e^{\delta_1 d(x,y)/16}\,\frac{\ell_{\mathfrak{S}}}{\ell_\iota}(x).
\end{equation}

\emph{The increment term.} Inserting the first summand of $\mathsf{y}_q$ from
\eqref{9.1:eq-pointwise-moves-3} into \eqref{9.1:eq-pointwise-moves-2}, the double
smearing carries the weight $e^{-2\delta_1 d(x,y)-2\delta_1 d(y,c)}$, with $y$ a cell of
$\mathfrak{S}$ and $c$ a cell of $\mathfrak{S}_P$. By the triangle inequality
$d(x,c)\le d(x,y)+d(y,c)$,
\begin{equation*}
d(x,c)+\tfrac12 d(x,y)\le\tfrac32 d(x,y)+d(y,c)\le 2d(x,y)+2d(y,c),
\end{equation*}
so this weight is at most $e^{-\delta_1 d(x,c)}e^{-\delta_1 d(x,y)/2}$. Together with
the factor $e^{\delta_1 d(x,y)/16}$ of \eqref{9.1:eq-pointwise-moves-4}, the weight left
in $y$ is $e^{-7\delta_1 d(x,y)/16}$, and the sum over $y$ is controlled by the
summation estimate among the random-walk estimates, which supplies the constant $c_1$.
The remaining sum over $c$ is
$\mathsf{b}^{\sim}(x)$ of \eqref{9.1:eq-pointwise-moves-1}. The numerical factor is
$16\cdot 1.1\,B_{\natural}\cdot B_{\natural}\le 18B_{\natural}^{2}$. Together with the
factor $L\,(\ell_{\mathfrak{S}}/\ell_\iota)(x)$ of \eqref{9.1:eq-pointwise-moves-4},
this gives the first line of \eqref{9.1:eq-pointwise-moves-5} below.

\emph{The link term.} On the cells $c\subset\bar{\Lambda}^{+}$,
Corollary~\ref{9.1:cor-seed-rotation} bounds the second summand of $\mathsf{y}_q$:
\begin{equation*}
\ell_\iota\sup_{\tilde{\Delta}(c)}|\mathcal{Y}_{lk}|
\le\mathsf{m}_{br}(c)+\mathsf{m}_{lk}(c).
\end{equation*}
The profiles $\mathsf{m}_{br}$ and $\mathsf{m}_{lk}$ of \eqref{9.1:eq-seed-rotation-a}
decay only at the rate $\delta_1/4$, slower than the weight $e^{-2\delta_1 d}$ of the
smearing. Their smearing against the weight
$e^{-2\delta_1 d(x,y)}$, multiplied by the factor $e^{\delta_1 d(x,y)/16}$ of
\eqref{9.1:eq-pointwise-moves-4}, is estimated by the transfer estimate for slowly
decaying profiles among the random-walk estimates, together with the summation estimate;
this gives the factor $c_1C_T$ in front of
$(\mathsf{m}_{br}+\mathsf{m}_{lk})(x)$. With the constant $16B_{\natural}$ of
\eqref{9.1:eq-pointwise-moves-2} and the factor $L\,(\ell_{\mathfrak{S}}/\ell_\iota)(x)$
of \eqref{9.1:eq-pointwise-moves-4}, this is the second line of
\eqref{9.1:eq-pointwise-moves-5} below.

\emph{Conclusion.} In both terms, the ratio at $y$ from
\eqref{9.1:eq-pointwise-moves-3} has been transported to $x$ by
\eqref{9.1:eq-pointwise-moves-4}, and it contributes the factor
$L\,(\ell_{\mathfrak{S}}/\ell_\iota)(x)$. Adding the two terms gives, for $x\in X$,
\begin{equation}\label{9.1:eq-pointwise-moves-5}
\begin{split}
\operatorname{move}_a(x)
&\le 18B_{\natural}^{2}c_1\,L\,\frac{\ell_{\mathfrak{S}}}{\ell_\iota}(x)\,
\ell_{\mathfrak{S}}(x)^{-e_a}\,\mathsf{b}^{\sim}(x)\\
&\quad+16B_{\natural}c_1C_T\,L\,\frac{\ell_{\mathfrak{S}}}{\ell_\iota}(x)\,
\ell_{\mathfrak{S}}(x)^{-e_a}\,(\mathsf{m}_{br}+\mathsf{m}_{lk})(x).
\end{split}
\end{equation}
On the cells $y$ of $\mathfrak{S}$ the two length functions satisfy
$\ell_{\mathfrak{S}}\le\ell_\iota$, by the comparison of the length functions of
$\lambda_{\mathfrak{S}}$ and $\iota$ on the box set; hence
\eqref{9.1:eq-pointwise-moves-3} gives $|\mathcal{Y}_q|_y\le\mathsf{y}_q(y)$ without
the transport \eqref{9.1:eq-pointwise-moves-4}. We run the two summations above with
this bound. In the increment term the weight left in $y$ is now
$e^{-\delta_1 d(x,y)/2}$, and since
$e^{-\delta_1 d(x,y)/2}\le e^{-7\delta_1 d(x,y)/16}$ the summation estimate applies with
the same constant $c_1$. In the link term the factor $e^{\delta_1 d(x,y)/16}\ge1$ is now
absent, so the transfer estimate and the summation estimate apply with the same constant
$c_1C_T$. This gives \eqref{9.1:eq-pointwise-moves-5} with the factor
$L\,(\ell_{\mathfrak{S}}/\ell_\iota)(x)$ replaced by $1$ in both lines. Both bounds hold
at every $x\in X$, and in each of them the two lines carry the same factor; the smaller
of the two right-hand sides is therefore the right-hand side with the factor
$\sigma(x)$ of \eqref{9.1:eq-pointwise-moves-1}, which is
\eqref{9.1:eq-pointwise-moves-a}.
\end{proof}

\begin{lemma}[The smeared increment on the core, the tower cells and far cells]
\label{9.1:lem-smeared-increment}
Assume the hypotheses of Lemma~\ref{9.1:lem-data-increments}, and let $\mathsf{b}^{\sim}$ be the
decay-smeared increment of Lemma~\ref{9.1:lem-pointwise-moves},
$\mathsf{b}^{\sim}(x)=\sum_c e^{-\delta_1 d(x,c)}\mathsf{b}_P(c)$, the sum running over the cells
of $\mathfrak{S}_P$. Let $h$ be the level of the collar $\mathfrak{K}_1$, and let $R_k$ denote the
depth of the core, the number of levels separating the collar from level $k$, so that $k-h=R_k$ and
$R_k$ is a positive integer. Assume that this depth satisfies $\tfrac94(1+2L)\le3L^{R_k/2}$. Let
$\Lambda$ be
the region of Lemma~\ref{9.1:lem-core-gauge}, let $\bar{\Lambda}$ and $\bar{\Lambda}^{+}$ be the
enlarged regions attached to it, and let $Z''^{(1)}_k$ be the tower region, the region of the tower
cells. Let $x$ be a point whose label has level $m$ (written $\pi$ in (ii)). Then:
\begin{equation}\label{9.1:eq-smeared-increment-a}
\begin{aligned}
&\text{(i) if } x\in\mathfrak{K}, \text{ with } m\le h: &
&\mathsf{b}^{\sim}(x)\le c_1\alpha_{0,m}\,\Xi';\\
&\text{(ii) if } x\in Z''^{(1)}_k, \text{ with } \pi\ge h: &
&\mathsf{b}^{\sim}(x)\le c_1\bigl[4B^{\mathsf{s}}C_1\delta'_k
+3K_{\mathfrak{K}}\tilde{\alpha}L^{\min\{\pi+1,k\}-k}\bigr]\\
& & &\phantom{\mathsf{b}^{\sim}(x)\le}
+c_1\alpha_{0,\pi}\,\Xi';\\
&\text{(iii) if } x\notin\bar{\Lambda}^{+}: &
&\mathsf{b}^{\sim}(x)\le c_1e^{-\delta_1 d(x,\bar{\Lambda})}\sup_{\bar{\Lambda}^{+}}\mathsf{b}_P
+8c_1\mathsf{m}_L(x)\\
& & &\phantom{\mathsf{b}^{\sim}(x)}
\le c_1(K_{\mathfrak{K}}+3)\tilde{\alpha}\,e^{-\delta_1 d(x,\bar{\Lambda})}+8c_1\mathsf{m}_L(x);\\
&\text{where} & &\mathfrak{b}_*:=\sup_{1\le m'\le h}\frac{\mathfrak{b}_{m'}}{\alpha_{0,m'}},\\
& & &\Xi':=14\,\mathfrak{b}_*+10\,\frac{\varsigma_h}{\alpha_{0,h}}
+7B^{\mathsf{s}}C_1\bigl(R_k^{1/2}+2\bigr)\frac{\delta'_k}{\alpha_{0,k}}\\
& & &\phantom{\Xi':=}+3K_{\mathfrak{K}}(1+\mathsf{r}_1)R_k^{1/2}L^{-R_k/2},\\
& & &C_4^{\mathfrak{K}}:=9C_4+\frac{2B_{\natural}c_1(K_{\mathfrak{K}}+3)}{K_\Lambda'}.
\end{aligned}
\end{equation}
If, moreover, the coupling ceilings \eqref{9.1:eq-core-ceilings-a} hold, then the second inequality
in (iii) holds and, at the far cells
$x\notin\bar{\Lambda}^{+}$, the majorant of the presentation at $x$ satisfies the far-field bound
assumed in Lemma~\ref{9.1:lem-data-increments}, with the constant $C_4$ multiplying
$\mathsf{m}_L$ there replaced by $C_4^{\mathfrak{K}}$.
\end{lemma}

\begin{proof}
\emph{Preliminaries.} By the decay estimate for cells of different levels, a cell $c$ of level $m'$
satisfies $e^{-\delta_1 d(x,c)/2}\le L^{-8(|m-m'|-1)_+}$, and $\sum_c e^{-\delta_1 d(x,c)/2}\le c_1$.
Write $e^{-\delta_1 d}=e^{-\delta_1 d/2}\,e^{-\delta_1 d/2}$. If $\mathcal{C}$ is a class of cells of
$\mathfrak{S}_P$ and $\mathsf{b}_P(c)\le\mu(m')$ for every $c\in\mathcal{C}$ of level $m'$, then,
bounding the first factor by the decay estimate and summing the second over the cells of each level,
\begin{equation}\label{9.1:eq-smeared-increment-4}
\sum_{c\in\mathcal{C}}e^{-\delta_1 d(x,c)}\mathsf{b}_P(c)
\le\sum_{m'}L^{-8(|m-m'|-1)_+}\mu(m')\sum_{\substack{c\in\mathcal{C}\\ \text{level } m'}}
e^{-\delta_1 d(x,c)/2}
\le c_1\sum_{m'}L^{-8(|m-m'|-1)_+}\mu(m').
\end{equation}
Bounds obtained in this way for disjoint classes add. By \eqref{9.1:eq-core-units-a},
$\alpha_{0,m'}\le\frac94(1\vee|m-m'|)^{1/2}\alpha_{0,m}$ for either order of $m,m'$. For every
$L\ge2$,
\begin{equation}\label{9.1:eq-smeared-increment-5}
\sum_{\Delta\in\mathbb{Z}}\tfrac94(1\vee|\Delta|)^{1/2}L^{-8(|\Delta|-1)_+}\le7,
\qquad
\sum_{\Delta\ge0}L^{-8(\Delta-1)_+}\le3.
\end{equation}
In the first sum the terms with $|\Delta|\le1$ give $\frac{27}{4}$ and the others less than
$\frac14$. In the second the terms $\Delta=0,1$ give $2$ and the others less than $1$. Below we also
use that $L$, being odd and larger than one, is at least $3$.

\emph{(i).} Let $x\in\mathfrak{K}$, of level $m\le h$. We treat the classes of cells in turn.

Cells of $\mathfrak{K}$. By \eqref{9.1:eq-data-increments-a}, on a cell of the core of level $m'$
one has $\mathsf{b}_P(c)\le2\mathfrak{b}_{m'}$. By the definition of $\mathfrak{b}_*$ in
\eqref{9.1:eq-smeared-increment-a} and the comparison of units,
\begin{equation*}
2\mathfrak{b}_{m'}\le2\mathfrak{b}_*\alpha_{0,m'}
\le2\mathfrak{b}_*\tfrac94(1\vee|m-m'|)^{1/2}\alpha_{0,m}.
\end{equation*}
By \eqref{9.1:eq-smeared-increment-4} and the first sum in \eqref{9.1:eq-smeared-increment-5}, these
cells contribute at most $14c_1\mathfrak{b}_*\alpha_{0,m}$.

Cells of $\mathfrak{B}_0$. These have levels $\lambda$ with $h\le\lambda\le k$ (here $\lambda$ is a
level index), the collar being at level $h$. By \eqref{9.1:eq-data-increments-a},
\begin{equation*}
\mathsf{b}_P(c)\le B^{\mathsf{s}}C_1\delta'_k+K_{\mathfrak{K}}\tilde{\alpha}L^{\lambda-k}
+\mathbb{1}_{\lambda=h}2\varsigma_h .
\end{equation*}
Since $h=k-R_k$, for every level $m\le h$ one has $k-m=R_k+(h-m)\ge R_k$ and, the square root being
subadditive, $(k-m)^{1/2}\le R_k^{1/2}+(h-m)^{1/2}$. In addition, for every integer $j\ge0$,
$j^{1/2}L^{-8(j-1)_+}\le2$.

For the last term, only $\lambda=h$ occurs. Write $2\varsigma_h=2\alpha_{0,h}\,\varsigma_h/\alpha_{0,h}$
and use $\alpha_{0,h}\le\frac94(1\vee(h-m))^{1/2}\alpha_{0,m}$. This gives the contribution
\begin{equation*}
c_1L^{-8(h-m-1)_+}2\varsigma_h
\le\tfrac92\,(1\vee(h-m))^{1/2}L^{-8(h-m-1)_+}\,c_1\alpha_{0,m}\frac{\varsigma_h}{\alpha_{0,h}}
\le10\,c_1\alpha_{0,m}\frac{\varsigma_h}{\alpha_{0,h}} .
\end{equation*}

For the first term we sum over the levels $\lambda\ge h$. If $h=m$, the second sum in
\eqref{9.1:eq-smeared-increment-5} gives $\sum_{\lambda\ge h}L^{-8(\lambda-m-1)_+}\le3$. If $h>m$,
then $(\lambda-m-1)_+=(h-m-1)+(\lambda-h)$, and the same sum gives
$\sum_{\lambda\ge h}L^{-8(\lambda-m-1)_+}\le3L^{-8(h-m-1)}$. Hence in both cases this sum is at most
$3L^{-8(h-m-1)_+}$. Write $\delta'_k=\alpha_{0,k}\,\delta'_k/\alpha_{0,k}$ and use
$\alpha_{0,k}\le\frac94(1\vee(k-m))^{1/2}\alpha_{0,m}$ together with
\begin{equation*}
L^{-8(h-m-1)_+}(1\vee(k-m))^{1/2}
\le\max\bigl\{1,\;R_k^{1/2}+(h-m)^{1/2}L^{-8(h-m-1)_+}\bigr\}
\le R_k^{1/2}+2 .
\end{equation*}
The first term therefore contributes at most
\begin{equation*}
\tfrac{27}{4}\bigl(R_k^{1/2}+2\bigr)\,c_1\alpha_{0,m}B^{\mathsf{s}}C_1\frac{\delta'_k}{\alpha_{0,k}},
\end{equation*}
which is at most the third term of $c_1\alpha_{0,m}\Xi'$.

For the second term write $\tilde{\alpha}=(1+\mathsf{r}_1)\alpha_{0,k}$. Since $k-m\ge R_k$ and
$R_k$ is a positive integer, the comparison of units gives
$\alpha_{0,k}\le\frac94(k-m)^{1/2}\alpha_{0,m}$. Since $m\le h\le\lambda$, the terms of the sum over
$\lambda$ are $L^{m-k}$ for $\lambda=m$ (which occurs only if $h=m$) and $L^{m+1-k}L^{-7j}$ for
$\lambda=m+1+j$, $j\ge0$; as $\sum_{j\ge0}L^{-7j}\le2$,
\begin{equation*}
\sum_{h\le\lambda\le k}L^{\lambda-k}L^{-8(\lambda-m-1)_+}\le(1+2L)L^{-(k-m)}.
\end{equation*}
The second term therefore contributes at most
\begin{equation*}
c_1K_{\mathfrak{K}}(1+\mathsf{r}_1)\alpha_{0,m}
\cdot\tfrac94(1+2L)L^{-(k-m)/2}\cdot(k-m)^{1/2}L^{-(k-m)/2}.
\end{equation*}
The last factor is bounded through
\begin{equation*}
(k-m)^{1/2}L^{-(k-m)/2}\le R_k^{1/2}L^{-R_k/2}.
\end{equation*}
This holds because $k-m\ge R_k\ge1$ and $s\mapsto s^{1/2}L^{-s/2}$ is decreasing on $[1,\infty)$: its
logarithmic derivative $\frac1{2s}-\frac{\log L}{2}$ is negative there, since $\log3>1$. The middle
factor is at most $\frac94(1+2L)L^{-R_k/2}$, because $k-m\ge R_k$, and by the assumption
$\tfrac94(1+2L)\le3L^{R_k/2}$ on the depth of the core this is at most $3$. The result is the
fourth term
$3c_1K_{\mathfrak{K}}(1+\mathsf{r}_1)R_k^{1/2}L^{-R_k/2}\alpha_{0,m}$ of
$c_1\alpha_{0,m}\Xi'$.

Cells off $\Lambda$. By \eqref{9.1:eq-data-increments-a}, $\mathsf{b}_P(c)=0$ on these cells, so they
contribute nothing.

Cells of the layers. Their contribution is contained in the rounding up of the numerical
coefficients of $\Xi'$, and gives no separate term.

Adding the four non-zero contributions gives $\mathsf{b}^{\sim}(x)\le c_1\alpha_{0,m}\Xi'$.

\emph{(ii).} Let $x\in Z''^{(1)}_k$, its label being of level $\pi\ge h$; thus $h\le\pi\le k$. On the
cells of $\mathfrak{B}_0$ the terms $B^{\mathsf{s}}C_1\delta'_k$ and
$K_{\mathfrak{K}}\tilde{\alpha}L^{\lambda-k}$ of \eqref{9.1:eq-data-increments-a}, summed over the
levels $h\le\lambda\le k$ with the decay factors $L^{-8(|\pi-\lambda|-1)_+}$ of
\eqref{9.1:eq-smeared-increment-4}, contribute at most
\begin{equation*}
c_1\Bigl[B^{\mathsf{s}}C_1\delta'_k\sum_{h\le\lambda\le k}L^{-8(|\pi-\lambda|-1)_+}
+K_{\mathfrak{K}}\tilde{\alpha}\sum_{h\le\lambda\le k}L^{\lambda-k}L^{-8(|\pi-\lambda|-1)_+}\Bigr].
\end{equation*}
The first sum is two-sided in $\lambda-\pi$ and is at most
\begin{equation*}
\sum_{\Delta\in\mathbb{Z}}L^{-8(|\Delta|-1)_+}=3+2\sum_{j\ge1}L^{-8j}\le4 .
\end{equation*}
In the second sum, the levels $\lambda=\pi-i$ with $i\ge0$ give at most
\begin{equation*}
L^{\pi-k}\Bigl(1+\sum_{i\ge1}L^{-i}L^{-8(i-1)}\Bigr)\le L^{\pi-k}\bigl(1+2L^{-1}\bigr),
\end{equation*}
and the levels $\lambda=\pi+1+j$ with $j\ge0$, present only if $\pi<k$, give
\begin{equation*}
\sum_{j\ge0}L^{\pi+1+j-k}L^{-8j}=L\,L^{\pi-k}\sum_{j\ge0}L^{-7j}\le L^{\pi-k}\bigl(L+2L^{-6}\bigr).
\end{equation*}
If $\pi=k$, only the first group occurs and, since $L\ge3$, the second sum is at most
$1+2L^{-1}\le3=3L^{\min\{\pi+1,k\}-k}$. If $\pi<k$, the two groups together give at most
\begin{equation*}
L^{\pi+1-k}\bigl(L^{-1}+2L^{-2}+1+2L^{-7}\bigr)\le3L^{\pi+1-k},
\end{equation*}
again since $L\ge3$. Hence in both cases the second sum is at most $3L^{\min\{\pi+1,k\}-k}$. These
cells therefore contribute at most
\begin{equation*}
c_1\bigl[4B^{\mathsf{s}}C_1\delta'_k+3K_{\mathfrak{K}}\tilde{\alpha}L^{\min\{\pi+1,k\}-k}\bigr].
\end{equation*}

The core and collar cells have levels at most $h$. They contribute terms carrying decay factors at
most $L^{-8(\pi-h-1)_+}$: the collar term $2\varsigma_h$ with exactly this factor, and the core
increments $2\mathfrak{b}_{m'}$ of level $m'$ with the factor $L^{-8(\pi-m'-1)_+}$. We compare them
with $\alpha_{0,\pi}$
as follows. First, by the comparison of units read from level $\pi$,
\begin{equation*}
2\mathfrak{b}_{m'}\le2\mathfrak{b}_*\tfrac94(1\vee|\pi-m'|)^{1/2}\alpha_{0,\pi},
\end{equation*}
and the first sum in \eqref{9.1:eq-smeared-increment-5} gives at most $14c_1\mathfrak{b}_*\alpha_{0,\pi}$.
Second, summing over the collar cells,
\begin{equation*}
c_1L^{-8(\pi-h-1)_+}2\varsigma_h
\le\tfrac92\,(1\vee(\pi-h))^{1/2}L^{-8(\pi-h-1)_+}c_1\alpha_{0,\pi}\frac{\varsigma_h}{\alpha_{0,h}}
\le10\,c_1\alpha_{0,\pi}\frac{\varsigma_h}{\alpha_{0,h}} .
\end{equation*}
All terms of $\Xi'$ are nonnegative, so these contributions are at most $c_1\alpha_{0,\pi}\Xi'$.
Cells off $\Lambda$ contribute nothing and the layer cells are treated as in (i). This
proves (ii).

\emph{(iii).} Let $x\notin\bar{\Lambda}^{+}$. By \eqref{9.1:eq-data-increments-a},
$\mathsf{b}_P(c)=0$ on the cells off $\Lambda$, so these contribute nothing. The remaining cells of
$\bar{\Lambda}^{+}$ contribute at most
$c_1e^{-\delta_1 d(x,\bar{\Lambda})}\sup_{\bar{\Lambda}^{+}}\mathsf{b}_P$. The layer cells contribute,
by the layer bound of \eqref{9.1:eq-data-increments-a}, at most $8c_1\mathsf{m}_L(x)$, their increments
being carried by the layer profile $\mathsf{m}_L$. This is the first inequality.

For the second, the ceilings \eqref{9.1:eq-core-ceilings-a} give $\sup_m\mathfrak{b}_m\le\tilde{\alpha}$
and $2\varsigma_h\le\tilde{\alpha}$. Inserted in the bounds of \eqref{9.1:eq-data-increments-a}, they
give $\sup_{\bar{\Lambda}^{+}}\mathsf{b}_P\le(K_{\mathfrak{K}}+3)\tilde{\alpha}$.

Finally, under the ceilings \eqref{9.1:eq-core-ceilings-a}, the bound (iii) on $\mathsf{b}^{\sim}(x)$
is inserted into the bound of Lemma~\ref{9.1:lem-pointwise-moves} for the moves of the table members
at $x$, which is linear in $\mathsf{b}^{\sim}(x)$ and in the profiles $\mathsf{m}_{br}$,
$\mathsf{m}_{lk}$. This gives, at the far cells, the far-field bound assumed in
Lemma~\ref{9.1:lem-data-increments} with $C_4$ replaced by $C_4^{\mathfrak{K}}$ of
\eqref{9.1:eq-smeared-increment-a}, which is the last assertion of the lemma.
\end{proof}

\begin{definition}[Coupling ceilings and the bracket circle near the core]\label{9.1:def-core-ceilings}
We use the units, radii and degrees on the core $\mathfrak{K}$ fixed in Notation~\ref{9.1:not-core-units},
the sizes of the core increments of Remark~\ref{9.1:rem-frame-size}, and the core majorant $\Xi'$,
the supremum $\mathfrak{b}_*$ and the far-cell constant $C_4^{\mathfrak{K}}$ of the smeared
increment on the core, the tower cells and far cells (Lemma~\ref{9.1:lem-smeared-increment},
display~\eqref{9.1:eq-smeared-increment-a}). In particular, the levels $h$ and $k$, the integers
$R_h$, $R_k$ and $N_0$, the quantities $\delta'_k$ and $\alpha_{0,k}$, the set $Z''_k$ and the
constants $B_\natural$, $B^{\mathsf{s}}$, $C_s$, $C_T$, $C_1$, $C_4$, $c_1$, $K_\Lambda'$,
$K_{\mathfrak{K}}$ and $\mathsf{r}_1$ are those of Notation~\ref{9.1:not-core-units}. Put
\begin{equation}\label{9.1:eq-core-ceilings-1}
\begin{aligned}
\Xi_{\mathfrak{K}}&:=\Xi'+\mathfrak{l}_*+2^{9}L^{-R_k/2},\\
\mathfrak{l}_*&:=8B^{\mathsf{s}}C_sC_4C_TK_\Lambda'(1+\mathsf{r}_1)\,(R_k^{1/2}+2)\,L^{-4(R_k-1)},\\
K_o^{\mathfrak{K}}&:=2^{20}B_\natural c_1^2C_TB^{\mathsf{s}}C_sC_1 .
\end{aligned}
\end{equation}
The \emph{coupling ceilings near the core} are the conditions (a)--(c) below; they are ceilings on
$\gamma$ in the sense of Notation~\ref{2.3:not-stages-first}. The \emph{bracket circle near the core}
is, at every plaquette $p$ of a bracket cube $\tilde{\square}$ meeting $\mathfrak{K}\cup Z''_k$, the
circle of the room analysis of the seed and of the remark on the bracket circle in the analysis away
from the core, with room $\mathsf{w}_L(p)$; it is required to satisfy the lower bound (d):
\begin{equation}\label{9.1:eq-core-ceilings-a}
\begin{aligned}
&\text{(a)}\quad 2^{20}B_\natural^2c_1^2\,\Xi_{\mathfrak{K}}\le 1;\\
&\text{(b)}\quad 2^{22}B_\natural^2c_1^2C_T(1+\mathsf{r}_1)
\Bigl[K_{\mathfrak{K}}(1+C_4)\,\ell_\circ^{-2N_0}
+B^{\mathsf{s}}C_sC_4K_\Lambda'L^{-3(N_0-1)}\Bigr]\le 1;\\
&\text{(c)}\quad \text{two ceilings of the room analysis of the seed, each with $C_4$ replaced}\\
&\phantom{\text{(c)}}\quad \text{by } C_4^{\mathfrak{K}};\\
&\text{(d)}\quad 1+\mathsf{w}_L(p)\ge
\Bigl[K_o^{\mathfrak{K}}\,(R_k^{1/2}+2)\,\delta'_k/\alpha_{0,k}\Bigr]^{-1}.
\end{aligned}
\end{equation}
Here $\ell_\circ$ is the integer written $L_0$ in \cite{B-CMP122a}, renamed $\ell_\circ$ here to avoid
confusion with $L_0=L_0(N)$. The bracket cubes $\tilde{\square}$ are those of the room analysis of the
seed.
Condition (b) is of the same kind as the
corresponding ceiling of the analysis away from the core, its first term arising from that ceiling
under the substitution of $4B_\natural c_1K_{\mathfrak{K}}$ for the product $C_4^{(k)}K_\Lambda'$
occurring there, in the form with $K_{\mathfrak{K}}$ in place of $K_{\mathfrak{K}}(1+C_4)$;
the second term in square brackets is the tower term carried by the links of the
localisation parameter $\sigma'$. Since $C_4>0$ and $K_{\mathfrak{K}}>0$
(Notation~\ref{9.1:not-core-units}), (b) implies the
same inequality with $K_{\mathfrak{K}}$ in place of $K_{\mathfrak{K}}(1+C_4)$. Each of (a)--(c) is
imposed uniformly in the levels: its left member is replaced by its supremum over the levels at which it
is read, as in the level-wise supremum convention of the room analysis of the seed.

In the one-sided order of choice of the present construction, which inserts its constants into the order
of Notations~\ref{2.3:not-stages-first}, \ref{2.3:not-stages-middle} and~\ref{2.3:not-stages-last} after
$\mathsf{r}_1$ and before $\gamma$, conditions (a)--(c) and the lower bound (d) are imposed after
$K_{\mathfrak{K}}$, $C_4^{\mathfrak{K}}$, $K_o^{\mathfrak{K}}$ and the cube count $n_{\mathfrak{K}}$ are
fixed, beside the
other ceilings of the room analysis and of the analysis away from the core, and before $\gamma$; no
condition among them returns to $M$ or to $L$.
\end{definition}

\begin{lemma}[The conditions near the core are ceilings]
Conditions (a)--(c) of Definition~\ref{9.1:def-core-ceilings} are ceilings in the sense of
Notation~\ref{2.3:not-stages-first}: every member of their left sides is a negative power of
$\log x_j$ at a level $j$ involved, or is exponentially small in $R_k$ or in $N_0$, with coefficients
depending on the dimension $d=4$ and on $L,M,M_1,M_2,\kappa_1$ only; none of them is a floor, and none
is a cap in the sense of
Definition~\ref{2.3:def-cap}. Under (a), the right member of (d) is at least
$7B_\natural/(C_TC_s)$.
\end{lemma}

\begin{proof}
\emph{The form of the left members.} By Remark~\ref{9.1:rem-frame-size}, for each level $m'$ with
$1\le m'\le h$ the quotient $\mathfrak{b}_{m'}/\alpha_{0,m'}$ is bounded by a power of $\log x_{m'}$ of
degree $2r_{pr}+p_0-q<0$, plus terms carrying the factor $L^{-(k-m')/2}$; hence so is
$\mathfrak{b}_*$, the supremum of these quotients over $1\le m'\le h$. By the size bound for the collar
gauge drift $\varsigma_h$, the
quotient $\varsigma_h/\alpha_{0,h}$ is bounded by the sum of a power of $\log x_h$ of the same degree
$2r_{pr}+p_0-q$ and a term $R_hR_k^{1/2}L^{-R_k}$. The term of $\Xi'$ containing $\delta'_k$, and the
quantity in square brackets in (d), carry the quotient $\delta'_k/\alpha_{0,k}$, a power of $\log x_k$
of degree
$r_{pr}/2+p_1-q<0$, the degree of the $\delta'$-member of the strengthened re-localisation ceiling of
the localisation analysis. The remaining members are
exponentially small: in $\Xi'$ the term $3K_{\mathfrak{K}}(1+\mathsf{r}_1)R_k^{1/2}L^{-R_k/2}$, in
$\mathfrak{l}_*$ the factor $L^{-4(R_k-1)}$, the term $2^9L^{-R_k/2}$ of $\Xi_{\mathfrak{K}}$, and in (b)
the factors $\ell_\circ^{-2N_0}$ and $L^{-3(N_0-1)}$. The conditions (c) are of the same kind as the
ceilings they modify, only the constant $C_4^{\mathfrak{K}}$ being new. The terms of $\mathfrak{b}_*$
carrying the factor $L^{-(k-m')/2}$ are exponentially small in $R_k$ as well, since $k-m'\ge R_k$ for
$m'\le h$, as used in the proof of Lemma~\ref{9.1:lem-smeared-increment}; the term
$R_hR_k^{1/2}L^{-R_k}$ carries the factor $L^{-R_k}$. The coefficients of all these members are
products of numerical factors and of the constants listed at the beginning of the definition; these
constants are fixed in the order of choice before $\gamma$, and they depend on
the dimension $d=4$ and on $L,M,M_1,M_2,\kappa_1$ only. The numbers $R_k$, $N_0$, $h$ and
$k$ occur in left members only; as prefactors they occur only through $R_k^{1/2}$ and
$R_hR_k^{1/2}$, each multiplied either by a factor exponentially small in $R_k$ or by a negative power
of $\log x_k$. This is within the range allowed in \cite{B-CMP122a}, where the integer playing the role
of $R_k$ is allowed to be a positive power of $\log g_k^{-2}$. Hence (a)--(c) bound no letter of the
order of choice from below and are not caps in the sense of Definition~\ref{2.3:def-cap}: they are
ceilings, neither floors nor caps.

\emph{The bracket circle.} All terms of $\Xi'$ are non-negative, the sizes $\mathfrak{b}_{m'}$,
$\varsigma_h$, $\delta'_k$ and the constants of Notation~\ref{9.1:not-core-units} being non-negative and
$\alpha_{0,m'}$, $\alpha_{0,h}$, $\alpha_{0,k}$ positive; so are $\mathfrak{l}_*$ and
$2^9L^{-R_k/2}$; therefore
\begin{equation*}
\Xi_{\mathfrak{K}}\ \ge\ \Xi'\ \ge\ 7B^{\mathsf{s}}C_1(R_k^{1/2}+2)\,\delta'_k/\alpha_{0,k},
\end{equation*}
the last member being the $\delta'$-member of $\Xi'$ in display~\eqref{9.1:eq-smeared-increment-a}. By (a),
\begin{equation*}
2^{20}\cdot 7\,B_\natural^2c_1^2B^{\mathsf{s}}C_1(R_k^{1/2}+2)\,\delta'_k/\alpha_{0,k}\le 1 .
\end{equation*}
By the definition of $K_o^{\mathfrak{K}}$ in~\eqref{9.1:eq-core-ceilings-1}, the quantity in square
brackets in (d) equals
$C_TC_s/(7B_\natural)$ times the left member of the last inequality; the constants $B_\natural$, $C_T$,
$C_s$ and the quantity $\delta'_k$ being positive (Notation~\ref{9.1:not-core-units}), both are
positive. Hence
\begin{equation*}
K_o^{\mathfrak{K}}(R_k^{1/2}+2)\,\delta'_k/\alpha_{0,k}\le \frac{C_TC_s}{7B_\natural},
\end{equation*}
and the right member of (d) is at least $7B_\natural/(C_TC_s)$.
\end{proof}

\begin{proposition}[Slices of the read objects on leaves meeting the core]
\label{9.1:prop-core-leaf-slices}
Assume the following.
\begin{itemize}
\item The closed form of the matching hypothesis of the reading theorem for large-field terms,
together with its composition rule for links, its slice hypothesis, its unit hypothesis and the
floors of its summation estimate.
\item The radius condition on the large-field region $\Lambda$, the kernel corollary, and the
far-cell presentation of the read objects.
\item The hypothesis on the core $\mathfrak{K}$, together with the regularity of the common core
variables, Definition~\ref{9.1:def-core-regularity}, display~\eqref{9.1:eq-core-regularity-a}.
\item $\mathrm{fr}=\mathrm{AL}$, the axial--Landau frame (any frame $\mathrm{fr}\in\mathfrak{F}_G$
may be used instead).
\item The ceilings and circles of Definition~\ref{9.1:def-core-ceilings},
display~\eqref{9.1:eq-core-ceilings-a}: the \emph{core ceiling}
\begin{equation*}
2^{20}B_\natural^2c_1^2\,\Xi_{\mathfrak{K}}\le 1,\qquad
\Xi_{\mathfrak{K}}=\Xi'+\mathfrak{l}_*+2^{9}L^{-R_k/2},
\end{equation*}
with $\Xi'$ the core majorant of Lemma~\ref{9.1:lem-smeared-increment} and $\mathfrak{l}_*$ the
link term of Definition~\ref{9.1:def-core-ceilings}; the \emph{large-field ceiling near the core},
in the form
\begin{equation*}
\begin{split}
&2^{22}B_\natural^2c_1^2C_T(1+\mathsf{r}_1)\\
&\qquad\cdot\bigl[K_{\mathfrak{K}}(1+C_4)\,\ell_\circ^{-2N_0}
+B^{\mathsf{s}}C_sC_4K'_\Lambda L^{-3(N_0-1)}\bigr]\le1;
\end{split}
\end{equation*}
the two far-point ceilings of Definition~\ref{9.1:def-core-ceilings},
display~\eqref{9.1:eq-core-ceilings-a}, which are those of the reading theorem with $C_4$
replaced by $C_4^{\mathfrak{K}}$; and
the \emph{bracket circles} on bracket cubes meeting $\mathfrak{K}\cup Z''_k$, with constant
$K_o^{\mathfrak{K}}$.
\end{itemize}
Let $(F,X,j)$ be a leaf with $X^{+}\cap\mathfrak{K}\neq\emptyset$ (for what follows,
equivalently: some site $y$ of the leaf's site set $\mathfrak{B}$ has
$\tilde{\Delta}(y)\cap\mathfrak{K}\neq\emptyset$; these
are exactly the leaves outside the scope of the slice statement for leaves meeting the large-field
region away from the core). Frozen left-overs on live pieces $W_{\mathfrak{a}}\subset\mathfrak{K}$
are included, with the letters $(\ell,\mathfrak{w},F)$ of the arrest construction, $\mathfrak{w}$
being its live-piece weight. Let
$q_0=e^{i\eta A_e}U_c$ be any seed of $\mathfrak{T}_\Phi[7/8]$ of the term $\Phi$ born at the
operation $\mathbf{R}$ (complex, in the chart of the shape, $U_c$ the centre of the chart in its
representative $\mathring{U}_c$); fix a vertex of $\mathfrak{P}_{\mathsf{c}}$; and let
$q\in\mathfrak{Q}''$, with its parameters $\zeta$, $B$ (on the dilated support of the cut-off
function $\chi'$ that carries $B$), $\sigma'$, and $\mathsf{t}$ (in the bracket circle).

Then $\tilde F(X;q)=F(X;S(q))$, where $S$ is holomorphic in all parameters and
$S(q)=q_0{\upharpoonright}X$ exactly at vanishing increments; the fifth member $E_5$ and the cube
clause of $S(q)$ are those of the entry, that is, of $U_c$; and in every free member $a$
\begin{equation}\label{9.1:eq-core-leaf-slices-a}
\begin{aligned}
\text{(a)}\quad & E_a(S(q))(x)\le \tfrac{1}{16}\,\mathsf{T}^a_N(x),\\
\text{(b)}\quad & E_a(S(q))(x)\le E_a(q_0)(x)+\tfrac{1}{1440}\,\mathsf{u}^a_F(x)
\le \tfrac{1}{16}\,\mathsf{u}^a_F(x),
\end{aligned}
\end{equation}
where (a) holds at the $(=)$-points of $X\cap\mathfrak{K}$ and at the points of the first class of
$X\cap Z''^{(1)}_k$, and (b) holds at the $(\mathrm{up})$-points of $X\cap\mathfrak{K}$ (old leaf,
$j<m(x)$) and at the points of the second class of $X\cap Z''^{(1)}_k$. More precisely:
\begin{itemize}
\item[(a)] The entry contributes at most $\mathsf{T}_N/32$, and the chain, the $\sigma'$-links and
the bracket together at most $\mathsf{T}_N/960$. Hence $\mathrm{rd}\,q$ meets the terminal row of
the reading theorem at the points of its first class, $\mathsf{T}_r=\tfrac12\mathsf{T}_N$, and the
requirement of its run row that the links be at most $\mathsf{G}/24$ (in the run-row form of the
link condition); and the schedule inequality gives
\begin{equation*}
\tfrac12\mathsf{T}_N\le\mathsf{T}_N\le(\mathsf{c}_F^{+}-\varrho_F)\,\mathsf{u}_F,
\end{equation*}
since, with $f_{k,m}$ the schedule coefficient of zone $m$ at level $k$,
$\mathsf{T}_N(x)=\varphi^{(N)}_m\mathsf{u}_m\le f_{k,m}\mathsf{u}_m$ on $\mathfrak{K}$ in
every zone, while on live pieces its live-piece form applies.
\item[(b)] There $E_a(q_0)\le\min\{\mathsf{T}^a_N,\mathsf{u}^a_F\}/32$, so that the bound (b),
which is at most $\tfrac1{16}\mathsf{u}^a_F\le(\tfrac34+\tfrac1{200})\mathsf{u}_F$, is of the raised
form of the reading theorem; at raised points of a run the links cost at most
\begin{equation*}
\tfrac{1}{1440}\,\mathsf{u}_F\le\tfrac{1}{24}\,\mathrm{rm}_{i_1}
\end{equation*}
by the room bound, as the run row of the reading theorem requires.
\item[(c)] At $x\in X\setminus\bar{\Lambda}^{+}$ (the points of the class $\mathcal{A}$, the points
of the tube class and the points of $Z^{c}$) the bounds of the reading theorem for its terminal row
hold verbatim with $C_4$ replaced by $C_4^{\mathfrak{K}}$ inside the slice lemma on the large-field
domain and inside the two far-point ceilings, which thereby become those of
Definition~\ref{9.1:def-core-ceilings}, display~\eqref{9.1:eq-core-ceilings-a}: the seed is handed
over with coefficient $1$, and the core and the collar add one exponentially decayed source.
\end{itemize}
Consequently $\mathrm{rd}\,q\in\mathfrak{T}[0](X)$ for every such leaf, so the slice condition of
the terminal row holds on leaves meeting $\mathfrak{K}$. The points of side one of leaves meeting
$\mathfrak{K}$ form the case $x\in Z''^{(1)}_k$ of (a)--(b). The dissolution case
$\Lambda\supset W_{\mathfrak{a}}$ of the schedule inequality on live pieces is covered by (a)--(b)
with the letters of the arrest construction, since $W_{\mathfrak{a}}\subset Z_h\subset\mathfrak{K}$.
No constant in these statements depends on $g$, $k$, $h$, $N$, $N_0$, the radii $R_\cdot$, the
ages, the depth or the volume.
\end{proposition}

\begin{proof}
\emph{Representation, holomorphy, exactness.} The representation
$\tilde F(X;q)=F(X;S(q))$ and the holomorphy of $S$ in all parameters follow from the
representation formula for slices and the move lemma for table members. They also use the
holomorphic dependence of the tube data on the parameters, and the covariance of the critical orbits
under complex $G^{c}$-valued block-centre gauges. Exactness, $S(q)=q_0{\upharpoonright}X$ at
vanishing increments, is Corollary~\ref{9.1:cor-seed-rotation}: there every read object is
presented as $q=e^{i\eta\mathcal{Y}_q}q_0$, with the majorant
\eqref{9.1:eq-seed-rotation-a}, and $\mathcal{Y}_q$ vanishes at vanishing increments.

\emph{The fifth member and the cube clause.} By the second part of the move lemma for table
members, the $G$-part of every slice is the centre $U_c$. At a point $x\in\mathfrak{K}$ of zone
$m$ the cube clause asks for cubes of size $CM\ell_m$ on which
\begin{equation*}
\ell_m|A|,\ \ell_m^2|\nabla A|<\varphi^{(N)}_m\,BCM\,\alpha_{0,m}.
\end{equation*}
The level-$k$ cube gauge of $U_c$ is a datum of $\mathfrak{T}_\Phi$: it is the cube clause at
level $k$ on $Z$, with bound $(1-\tfrac18\theta_\rho\beta)BCM\alpha_{0,k}$. Restricted to a
sub-cube of size $CM\ell_m$ it gives the ratio
\begin{equation}\label{9.1:eq-core-leaf-slices-1}
\frac{\alpha_{0,k}}{\alpha_{0,m}}\,L^{-(k-m)}
\le\tfrac94\,(k-m)^{1/2}L^{-(k-m)}
\le 3L^{-R_k/2}
\le\tfrac15 .
\end{equation}
The first inequality is \eqref{9.1:eq-core-units-a}. For the second, every zone $m\le h$ of
$\mathfrak{K}$ has $k-m\ge R_k$, as in Lemma~\ref{9.1:lem-smeared-increment}. Moreover
$s^{1/2}L^{-s/2}\le1$ for $s\ge1$, because $s\le L^{s}$. Hence, with $s=k-m$,
\begin{equation*}
\tfrac94\,s^{1/2}L^{-s}\le\tfrac94\,L^{-s/2}\le\tfrac94\,L^{-R_k/2}.
\end{equation*}
The third inequality follows from the last member of $\Xi_{\mathfrak{K}}$ in the core ceiling,
which gives $2^{20}\cdot2^{9}L^{-R_k/2}\le1$. This is the corresponding line of the threshold
arithmetic of the slice lemma on the large-field domain, continued to the zones $m\le h$. Likewise,
\begin{equation*}
E_5=|\partial U_c-1|<\alpha_{0,k}\ell_k^{-2}\le 3L^{-R_k/2}\,\mathsf{T}^5_N .
\end{equation*}
By the second part of the move lemma, the first member $E_1$ moves by at most
$2\eta^2|\nabla E|(1+C\varepsilon_3)$, with $E$ and $C\varepsilon_3$ as in that lemma, and this is
contained in the move bound below.

\emph{The entry.} By the definition of the tube and the first display of the entry bound of the
slice lemma on the large-field domain, $E_a(q_0)(x)<\mathsf{u}^a_k(x)$ on $Z$. At
$x\in\mathfrak{K}$ of zone $m$ we use $\mathsf{T}^a_N\ge\tfrac15\alpha_{a,m}\ell_m^{-e_a}$ and
\eqref{9.1:eq-core-units-a}, and obtain
\begin{equation*}
\frac{\mathsf{u}^a_k}{\mathsf{T}^a_N}
\le 5\,\frac{\alpha_{a,k}}{\alpha_{a,m}}\,L^{-e_a(k-m)}
\le 12\,(k-m)^{1/2}L^{-(k-m)}
\le 12\,L^{-R_k/2}
\le 2^{-5}.
\end{equation*}
The third step is the argument of \eqref{9.1:eq-core-leaf-slices-1}. The last step uses
$2^{20}\cdot2^{9}L^{-R_k/2}\le1$ from the last member of the core ceiling, since
$12\cdot2^{-29}\le2^{-5}$. At $(\mathrm{up})$-points the same argument gives
\begin{equation*}
\frac{\mathsf{u}^a_k}{\mathsf{u}^a_F}
\le\tfrac94\,(k-j)^{1/2}L^{-(k-j)}
\le 2^{-5}.
\end{equation*}
On $Z''^{(1)}_k$ the entry bound of the slice lemma on the large-field domain applies unchanged.
Its ceiling is implied by the large-field ceiling near the core. Thus the entry is at most
$\mathsf{T}_N/32$ at the $(=)$-points of $\mathfrak{K}$, and at most
$\min\{\mathsf{T}^a_N,\mathsf{u}^a_F\}/32$ at the $(\mathrm{up})$-points.

\emph{Moves at an $(=)$-point $x\in\mathfrak{K}$.} Here $\ell_{\mathfrak{S}}=\ell_m$ and
$\sigma(x)=1$. Moreover $\mathsf{T}^a_N\ge\tfrac15\alpha_{a,m}\ell_m^{-e_a}$ and
$\alpha_{a,m}\ge\alpha_{0,m}$. On a live piece the factor $\tfrac15$ is replaced by
$\varphi_{\mathrm{arrest}}\mathfrak{w}$, and $\mathfrak{b}$ and $\mathsf{u}$ carry the same weight
$\mathfrak{w}$. We treat the chain, the $\sigma'$-links and the bracket in turn.

The chain part comes from the first term of \eqref{9.1:eq-pointwise-moves-a}
(Lemma~\ref{9.1:lem-pointwise-moves}) and from part~(i) of \eqref{9.1:eq-smeared-increment-a}
(Lemma~\ref{9.1:lem-smeared-increment}), which gives
$\mathsf{b}^{\sim}(x)\le c_1\alpha_{0,m}\Xi'$. Together they bound it by
\begin{equation*}
18B_\natural^2c_1\,\ell_m^{-e_a}\,c_1\alpha_{0,m}\Xi'
\le 18B_\natural^2c_1\cdot5\cdot c_1\Xi'\,\mathsf{T}^a_N
=90B_\natural^2c_1^2\,\Xi'\,\mathsf{T}^a_N .
\end{equation*}

For the $\sigma'$-links, a path from a point of zone $m$ of $\mathfrak{K}$ to $(\bar{\Lambda})^{c}$
crosses the strata $m+1,\dots,k-1$: the zones $m+1,\dots,h$ of $\mathfrak{K}$, and then the strata
$h+1,\dots,k-1$. By the distance property of the level functions, therefore,
\begin{equation*}
\mathsf{m}_{lk}(x)\le 8B^{\mathsf{s}}C_sc_1C_4\,\mathsf{s}_L\,L^{-4(k-m-1)},
\end{equation*}
with $\mathsf{s}_L=K'_\Lambda\tilde{\alpha}$ and $\tilde{\alpha}=(1+\mathsf{r}_1)\alpha_{0,k}$.
The second term of \eqref{9.1:eq-pointwise-moves-a} then bounds their move by
\begin{equation*}
\begin{split}
&80B_\natural c_1C_T\cdot 8B^{\mathsf{s}}C_sc_1C_4K'_\Lambda(1+\mathsf{r}_1)
\frac{\alpha_{0,k}}{\alpha_{0,m}}L^{-4(k-m-1)}\,\mathsf{T}^a_N\\
&\qquad\le 180B_\natural^2c_1^2\,\mathfrak{l}_*\,\mathsf{T}^a_N ,
\end{split}
\end{equation*}
using $\alpha_{0,k}/\alpha_{0,m}\le\tfrac94(k-m)^{1/2}$, the definition of $\mathfrak{l}_*$, and
\begin{equation*}
\begin{split}
(k-m)^{1/2}L^{-4(k-m-1)}
&\le\bigl(R_k^{1/2}+(h-m)^{1/2}\bigr)L^{-4(R_k-1)}L^{-4(h-m)}\\
&\le(R_k^{1/2}+2)\,L^{-4(R_k-1)}.
\end{split}
\end{equation*}
Here $(k-m)^{1/2}\le R_k^{1/2}+(h-m)^{1/2}$, as in Lemma~\ref{9.1:lem-smeared-increment};
$k-m-1\ge(R_k-1)+(h-m)$, because $k-h\ge R_k$ (the case $m=h$ of $k-m\ge R_k$); and
$(h-m)^{1/2}L^{-4(h-m)}\le1$ for every integer $h-m\ge0$.
Since $90\le180$ and $\Xi'+\mathfrak{l}_*\le\Xi_{\mathfrak{K}}$, the core ceiling bounds chain
and $\sigma'$-links together by
\begin{equation*}
180B_\natural^2c_1^2(\Xi'+\mathfrak{l}_*)\,\mathsf{T}^a_N
\le180\cdot2^{-20}\,\mathsf{T}^a_N<\tfrac{1}{5800}\,\mathsf{T}^a_N ,
\end{equation*}
because $180\cdot5800=1044000<2^{20}=1048576$.

For the bracket,
\begin{equation*}
\mathsf{m}_{br}(x)\le(1+|\mathsf{t}|)\,B^{\mathsf{s}}C_sc_1C_1\,\delta'_k\,L^{-4(h-m-1)_+}.
\end{equation*}
We use $\alpha_{0,k}/\alpha_{0,m}\le\tfrac94(k-m)^{1/2}$ together with
$(k-m)^{1/2}\le R_k^{1/2}+(h-m)^{1/2}$, as in Lemma~\ref{9.1:lem-smeared-increment}, and
$(h-m)^{1/2}L^{-4(h-m-1)_+}\le2$ for $L\ge2$ (the maximum over $h-m$ is attained at $h-m=1$);
and we bound $1+|\mathsf{t}|$ by $1+\mathsf{w}_L$ for $\mathsf{t}$ in the bracket circle of
Definition~\ref{9.1:def-core-ceilings}.
Its move is then at most
\begin{equation*}
\begin{split}
&80B_\natural c_1C_T\,B^{\mathsf{s}}C_sc_1C_1(1+|\mathsf{t}|)\tfrac94(k-m)^{1/2}
L^{-4(h-m-1)_+}\frac{\delta'_k}{\alpha_{0,k}}\,\mathsf{T}_N\\
&\qquad\le180B_\natural c_1^2C_TB^{\mathsf{s}}C_sC_1(1+\mathsf{w}_L)(R_k^{1/2}+2)
\frac{\delta'_k}{\alpha_{0,k}}\,\mathsf{T}_N\\
&\qquad\le180\cdot2^{-20}\,\mathsf{T}_N<\tfrac{1}{5800}\,\mathsf{T}_N .
\end{split}
\end{equation*}
The last line uses the bracket circle of Definition~\ref{9.1:def-core-ceilings}, whose constant is
$K_o^{\mathfrak{K}}=2^{20}B_\natural c_1^2C_TB^{\mathsf{s}}C_sC_1$.

Adding up, the links are less than $\mathsf{T}_N/2900<\mathsf{T}_N/960$. With the entry, the total
is less than $\mathsf{T}_N/32+\mathsf{T}_N/2900<\mathsf{T}_N/16$. This proves (a) on
$\mathfrak{K}$.

The $(=)$-points are the points with $j\ge m(x)$, since there the label level
$\lambda_{\mathfrak{S}}$ and the unit level $\iota$ satisfy
$\lambda_{\mathfrak{S}}=j\wedge m=m=\iota$; this is the identification of the label level with the
unit level, in the point-class convention of the slice lemma on the large-field domain. At these
points the table is $\mathsf{T}_N$, by the identification of the table in that lemma. That
identification rests on \cite[(1.64)(i)]{B-CMP122a}, stated there for $m=1,\dots,j-1$ and
carrying no $\ell_\circ$-weight for $m\le k_0$, where $k_0$ is the level above which the table
carries the weight $\ell_\circ^{2(\pi-k_0)}$ at a level $\pi>k_0$.

\emph{Moves at an $(\mathrm{up})$-point $x\in\mathfrak{K}$.} Here the leaf is old,
$j<m(x)\le h$. We have $n_F(x)=j$, where $n_F(x)$ is the level of $x$ for the term $F$: the point
$x\in\Omega_m\subset\Omega_j$, with $\Omega_m$ the domain of level $m$ in the decreasing sequence
of domains, was a level-$j$ point at
the birth of $F$ and has been coarsened since, as in the definition of raised points. Hence
$\ell_{\mathfrak{S}}(x)=\ell_j$ and $\sigma(x)\le L\cdot L^{-(m-j)}$. The first term of
\eqref{9.1:eq-pointwise-moves-a} with part~(i) of \eqref{9.1:eq-smeared-increment-a}, and
\eqref{9.1:eq-core-units-a}, give
\begin{equation*}
\begin{split}
\frac{\operatorname{move}_a(x)}{\mathsf{u}^a_F}
&\le18B_\natural^2c_1\,L\cdot L^{-(m-j)}\,\frac{c_1\alpha_{0,m}\Xi'}{\alpha_{a,j}}
\le18\cdot\tfrac94\,B_\natural^2c_1^2\,L(m-j)^{1/2}L^{-(m-j)}\,\Xi'\\
&\le41B_\natural^2c_1^2\,\Xi'\le41\cdot2^{-20}.
\end{split}
\end{equation*}
Here $Ls^{1/2}L^{-s}\le1$ for $s\ge1$ and $L\ge2$, and the last step is the core ceiling. The
$\sigma'$-links and the bracket carry the same factors $\sigma$ and $\tfrac94(m-j)^{1/2}$ and the
same change of unit, the $\sigma'$-link profile carrying the factor $L^{-4(k-m-1)}$ obtained above.
As on the $(=)$-points, they are therefore bounded together by
$2\cdot180\cdot2^{-20}\,\mathsf{u}^a_F$. Hence the total of the links is at most
\begin{equation*}
401\cdot2^{-20}\,\mathsf{u}^a_F<\tfrac{1}{2600}\,\mathsf{u}^a_F<\tfrac{1}{1440}\,\mathsf{u}^a_F ,
\end{equation*}
since $401\cdot2600=1042600<2^{20}$. With the entry, which is at most $\mathsf{u}_F/32$, the total
is less than $\mathsf{u}_F/16$. This proves (b) on $\mathfrak{K}$.

\emph{Moves at $x\in Z''^{(1)}_k$ of level $\pi$.} At points of the first class
$\ell_{\mathfrak{S}}=\ell_\pi$. At points of the second class $\ell_{\mathfrak{S}}=\ell_j<\ell_\pi$
and $u=k-n_F\ge N_0$; this is the point-class convention of the slice lemma on the large-field
domain. We split the bound of part~(ii) of \eqref{9.1:eq-smeared-increment-a},
\begin{equation*}
\mathsf{b}^{\sim}(x)\le c_1\bigl[3B^{\mathsf{s}}C_1\delta'_k+3K_{\mathfrak{K}}\tilde{\alpha}
L^{\pi-k}\bigr]+c_1\alpha_{0,\pi}\,\Xi',
\end{equation*}
and treat its terms in turn.

The term $3K_{\mathfrak{K}}\tilde{\alpha}L^{\pi-k}$ is absolute at level $k$ when read as a field
majorant. Indeed Corollary~\ref{9.1:cor-seed-rotation} gives
$\sup|\mathcal{Y}|\le1.1B_\natural(\mathsf{b})^{\approx}/\ell_\iota$, so
\begin{equation*}
\sup|\mathcal{Y}|
\le1.1B_\natural c_1\cdot3K_{\mathfrak{K}}\tilde{\alpha}L^{\pi-k}/\ell_\pi
\le4B_\natural c_1K_{\mathfrak{K}}\tilde{\alpha}/\ell_k
=:\mathsf{m}^{\mathfrak{K}}/\ell_k .
\end{equation*}
The link lemma of the slice lemma on the large-field domain, applied with
$\mathsf{m}=\mathsf{m}^{\mathfrak{K}}$, bounds the resulting move by
$180B_\natural c_1C_T(\mathsf{m}^{\mathfrak{K}}/\alpha_{a,k})\ell_\circ^{-2N_0}\mathsf{T}^a_N$ at
points of the first class. At points of the second class the bound is the same quantity with $36$
in place of $180$ and $\mathsf{u}^a_F$ in place of $\mathsf{T}^a_N$. Since
$\tilde{\alpha}/\alpha_{a,k}\le1+\mathsf{r}_1$, both are at most
\begin{equation*}
720B_\natural^2c_1^2C_TK_{\mathfrak{K}}(1+\mathsf{r}_1)\ell_\circ^{-2N_0}
\cdot(\mathsf{T}^a_N\ \text{resp.}\ \mathsf{u}^a_F)
<\tfrac{1}{5800}\,(\mathsf{T}^a_N\ \text{resp.}\ \mathsf{u}^a_F)
\end{equation*}
under the large-field ceiling near the core, since $720\cdot5800=4176000<2^{22}$.

The term $3B^{\mathsf{s}}C_1\delta'_k$ and the core and collar term $c_1\alpha_{0,\pi}\Xi'$ are
measured in the own unit at level $\pi$. Since $\pi\ge h$ we have $(k-\pi)^{1/2}\le R_k^{1/2}$,
and $3\cdot\tfrac94<7$, so the first member of the bracket below is at most $\Xi'$. Their ratio to
$\mathsf{T}_N$ is therefore at most
\begin{equation*}
90B_\natural^2c_1^2\Bigl[3B^{\mathsf{s}}C_1\tfrac94(k-\pi)^{1/2}\frac{\delta'_k}{\alpha_{0,k}}
+\Xi'\Bigr]
\le90B_\natural^2c_1^2\cdot2\,\Xi_{\mathfrak{K}}<\tfrac{1}{5800},
\end{equation*}
by the core ceiling. At points of the second class, the factor $\sigma\le L\cdot L^{-(\pi-j)}$ and
the change of unit give the second-class arithmetic of the slice lemma on the large-field domain,
with the same outcome in the units of $\mathsf{u}_F$.

For the $\sigma'$-links, $\mathsf{m}_{lk}$ at level $\pi$ carries the factor
$L^{-4(k-\pi-1)_+}$. For $\pi\le k_0$ this factor is at most $L^{-4(N_0-1)}$, that is,
$k-\pi\ge N_0$. For $k_0<\pi<k$ the
weight $\ell_\circ^{2(\pi-k_0)}$ of the table multiplies the threshold, and the threshold function
$f$ of the threshold arithmetic of the slice lemma on the large-field domain satisfies
$f(s)\ge\ell_\circ^{2N_0}$. For $\pi\le k_0$ the ratio is therefore at most
\begin{equation*}
80B_\natural c_1C_T\cdot8B^{\mathsf{s}}C_sc_1C_4K'_\Lambda(1+\mathsf{r}_1)\tfrac94(k-\pi)^{1/2}
L^{-4(k-\pi-1)},
\end{equation*}
and for $k_0<\pi<k$ it is at most the same expression multiplied by $\ell_\circ^{-2N_0}$. Now
$s^{1/2}L^{-(s-1)}\le1$ for every integer $s\ge1$ and $L\ge2$, because
$L^{s-1}\ge2^{s-1}\ge s\ge s^{1/2}$. With $s=k-\pi$ this gives
\begin{equation*}
(k-\pi)^{1/2}L^{-4(k-\pi-1)}\le L^{-3(k-\pi-1)}\le L^{-3(N_0-1)}\quad(\pi\le k_0),
\end{equation*}
and $(k-\pi)^{1/2}L^{-4(k-\pi-1)}\le1$ for $k_0<\pi<k$. In both ranges, therefore, the ratio is at
most
\begin{equation*}
1440B_\natural c_1^2C_TB^{\mathsf{s}}C_sC_4K'_\Lambda(1+\mathsf{r}_1)
\max\{L^{-3(N_0-1)},\ell_\circ^{-2N_0}\}<\tfrac{1}{2900}.
\end{equation*}
In the $L$-branch this is the second member of the large-field ceiling near the core
($1440\cdot2900<2^{22}$). In the $\ell_\circ$-branch we use
$B^{\mathsf{s}}C_sC_4K'_\Lambda\le K_{\mathfrak{K}}C_4$, and the branch is controlled by the first
member of that ceiling, $K_{\mathfrak{K}}(1+C_4)\ell_\circ^{-2N_0}$.

The bracket is treated as on $\mathfrak{K}$, with $L^{-4(h-m-1)_+}$ replaced by $1$. By the bracket
circle it contributes less than $\mathsf{T}_N/5800$.

Totals: at points of the first class the links are less than
\begin{equation*}
\Bigl(\tfrac{3}{5800}+\tfrac{1}{2900}\Bigr)\mathsf{T}_N=\tfrac{5}{5800}\,\mathsf{T}_N
=\tfrac{1}{1160}\,\mathsf{T}_N<\tfrac{1}{960}\,\mathsf{T}_N ,
\end{equation*}
and at points of the second class, with the constants of the second-class arithmetic, less than
$\mathsf{u}_F/1450\le\mathsf{u}_F/1440$. Adding the entry, which is at most $\mathsf{T}_N/32$,
respectively $E_a(q_0)$, gives (a) and (b) on $Z''^{(1)}_k$.

\emph{Points $x\notin\bar{\Lambda}^{+}$.} By part~(iii) of \eqref{9.1:eq-smeared-increment-a}, the
presentation $q=e^{i\eta\mathcal{Y}_q}q_0$ obeys the display of the far-cell presentation with
$C_4^{\mathfrak{K}}$ in place of $C_4$, and the links are those of that presentation. The proof of
the slice lemma on the large-field domain at the points of the tube class and of the class
$\mathcal{A}$, and the proof of the terminal row of the reading theorem, read nothing else of the
chain. They use the slice bound with $\mathsf{m}=2(\mathsf{m}_L+(1+|\mathsf{t}|)\mathsf{m}_p)$,
where $\mathsf{m}_L$ is the layer profile and $\mathsf{m}_p$ the profile multiplying
$1+|\mathsf{t}|$ in the far-cell presentation, the auxiliary gauge and
dilation estimates of those proofs, and the two far-point ceilings, which with
$C_4^{\mathfrak{K}}$ in place of $C_4$ are those of Definition~\ref{9.1:def-core-ceilings},
display~\eqref{9.1:eq-core-ceilings-a}. These hold verbatim with
$C_4^{\mathfrak{K}}$, which proves (c).

\emph{Schedule inequality and rooms.} At the $(=)$-points of $\mathfrak{K}$,
$\mathsf{T}_N=\varphi^{(N)}_m\mathsf{u}_m$ with
\begin{equation*}
\varphi^{(N)}_m=f_{k,m}-\beta\sum_{i=h+1}^{k}2^{-(i-m)}<f_{k,m}.
\end{equation*}
This is the formula for the table coefficients, with \cite[pp.~190--191]{B-CMP122a} and the
corresponding display in the proof of the reading theorem. The schedule inequality applies at
$(=)$-points of layers $n<j$, and at $n=j$ when $k>j$. It therefore gives
\begin{equation*}
\tfrac12\mathsf{T}_N\le(\mathsf{c}_F^{+}-\varrho_F)\,\mathsf{u}_F
\end{equation*}
for every stored $F$ read at $\mathbf{R}_k$ with $k>j$. For $j=k$ (the chain terms
$\mathbb{C}^{(i)}_k$ of the present step, and the left-overs whose level $k_a$ equals $k$) the
points of $\mathfrak{K}$ have $n=m<k=j$, which is the first case of the schedule inequality. On
live pieces its live-piece form applies, with its room coefficient $\varrho_F$. At
$(\mathrm{up})$-points and at points of the second class, the room inequality at raised points
gives
\begin{equation*}
\tfrac18\,\mathsf{u}_F\le0.755\,\mathsf{u}_F\le(\mathsf{c}_F^{+}-\varrho_F)\,\mathsf{u}_F .
\end{equation*}
Together with (a)--(c), these are the conditions of the slice bound at the points of $X$, so
$\mathrm{rd}\,q\in\mathfrak{T}[0](X)$.
\end{proof}

\begin{remark}[The same bounds on leaves missing the core]\label{9.1:rem-leaves-missing-core}
Assume the hypotheses of Proposition~\ref{9.1:prop-core-leaf-slices} except the condition
$X^{+}\cap\mathfrak{K}\neq\emptyset$, and let $(F,X,j)$ be a leaf all of whose cells miss the
core $\mathfrak{K}$. Such a leaf lies outside the scope of
Proposition~\ref{9.1:prop-core-leaf-slices}. For it the inequality of the slice statement for
leaves not meeting the core holds in the form in which that statement gives it, namely as
follows. The point classes $(=)$, $(\mathrm{up})$, $(\mathrm{a})$, $(\mathrm{b})$ are those of
Proposition~\ref{9.1:prop-core-leaf-slices}, taken here on the cells of the leaf. At the points
of the first class (the $(=)$- and $(\mathrm{a})$-points) there is a link budget of the same
form as in item~(a) of Proposition~\ref{9.1:prop-core-leaf-slices}: the chain, link and bracket
contributions together are at most $\mathsf{T}_N/960$. At the points of the second class (the
$(\mathrm{up})$- and $(\mathrm{b})$-points),
\begin{equation}\label{9.1:eq-leaves-missing-core-1}
E_a(S(q))(x)\le E_a(q_0)(x)+\mathsf{u}^a_F(x)/1440 .
\end{equation}
Here $q_0$, $q$ and the slice map $S$ are as in Proposition~\ref{9.1:prop-core-leaf-slices}.
In the derivation of both bounds, the ceiling on the core and the ceiling on $\Lambda$ near the
core, the first two conditions of \eqref{9.1:eq-core-ceilings-a}, take the place of the two
ceilings on $\Lambda$, the additional hypothesis and the seam constant under which the slice
statement for leaves not meeting the core is stated; the additional hypothesis and the seam
constant do not occur here, and the other hypotheses are those assumed above. The two separate
estimates of that statement, made under its two ceilings on $\Lambda$, reappear here as the two
terms $K_{\mathfrak{K}}\tilde{\alpha}L^{\lambda-k}$ and $B^{\mathsf{s}}C_1\delta'_k$ of the bound
on the cells of $\mathfrak{B}_0$ in \eqref{9.1:eq-data-increments-a}, with $\lambda$ the level
of the cell, as there.
\end{remark}

\begin{proof}
The proof of the data-increment bounds of Lemma~\ref{9.1:lem-data-increments}, displayed in
\eqref{9.1:eq-data-increments-a}, nowhere uses the condition
$X^{+}\cap\mathfrak{K}\neq\emptyset$. Neither do the pointwise moves
\eqref{9.1:eq-pointwise-moves-a} nor the bounds on the smeared increment of
Lemma~\ref{9.1:lem-smeared-increment}, displayed in \eqref{9.1:eq-smeared-increment-a}. All
three therefore hold for the leaf $(F,X,j)$.

Since every cell of the leaf misses $\mathfrak{K}$, the bound of
Lemma~\ref{9.1:lem-smeared-increment} for points of the core, its item~(i), does not enter. Only
its items~(ii) and~(iii) enter: the bounds on the tower cells and on the far cells. These two
items, under the first two conditions of \eqref{9.1:eq-core-ceilings-a}, give the inequality of
the slice statement for leaves not meeting the core in the form stated above: at the points of
the first class, the bound $\mathsf{T}_N/960$ on the sum of the chain, link and bracket
contributions and, at the points of the second class, \eqref{9.1:eq-leaves-missing-core-1}.
\end{proof}

\begin{remark}[The junction surface, the cut surface and their cost]\label{9.1:rem-junction-cut}
Two surfaces are attached to the core region
$\mathfrak{K}=\mathfrak{K}_1\cup\mathfrak{K}_2$ of \eqref{9.1:eq-core-collar-a}
(Definition~\ref{9.1:def-core-collar}). The \emph{junction surface}
$S_c:=\partial\Omega''^{\sim 2}_{h+1}$ is the surface where the variables on $\mathfrak{B}_0$
stop and the variables $V_h$ begin; it separates the collar $\mathfrak{K}_1$ from the deep part
$\mathfrak{K}_2$. The \emph{cut surface} $\partial\Omega''^{\sim}_{h+1}$ is the
surface at which the configuration is cut in \cite[(1.52)--(1.53)]{B-CMP122b}; it is the outer
boundary of $\mathfrak{K}$ and lies one $\sim$-layer outside $S_c$. We call these two surfaces
the \emph{seams}. For these two surfaces the following hold.
Here, as in the proof of Lemma~\ref{9.1:lem-data-increments}, $B[P](c)$ (not to be confused with
the $\mathbf{T}$-step variable $B$) denotes the data
coordinate of a read object $P$ at a cell $c$, taken in the gauge
$g$ of Lemma~\ref{9.1:lem-core-gauge}; in particular $B[P^{\mathrm{re}}]=B^0_g$.
\begin{itemize}
\item[(a)] At the real point, in the axial--Landau frame $\mathrm{AL}$, the junction surface costs
at most $\mathfrak{q}_h+2a_h$ per bond, up to products of these two small quantities, and no term
of another kind enters. In the axial frame of \cite[(1.81)]{B-CMP122a}, on the real plateau
configuration of the axial plateau proposition, the logarithm of the collar datum is the axial
holonomy $\mathsf{m}_{\mathrm{ax}}$ of that proposition, which satisfies
\begin{equation*}
|\log\mathsf{m}_{\mathrm{ax}}|\ge(t_{\mathrm{ax}}-\mathfrak{c}_{\mathrm{ax}})\,
f_{\mathrm{ax}}\,\alpha_{0,k},
\end{equation*}
where $t_{\mathrm{ax}}\ge M$, $f_{\mathrm{ax}}$ and $\mathfrak{c}_{\mathrm{ax}}$ are the constants
of that proposition; there the junction cost is at least
$\frac12 f_{\mathrm{ax}}M\alpha_{0,k}$. Consequently $\mathfrak{b}_h$ is at least this large,
and the smallness condition on the core increments $\mathfrak{b}_m$ stated in
Remark~\ref{9.1:rem-frame-size}, which contains the ceiling hypothesis of
Lemma~\ref{9.1:lem-core-gauge}, cannot be met: the set on which it holds is empty. Among
the estimates of this section, the gauge frame enters only through the collar datum $D^0$ of the
real object, in its comparison with the datum $1$ of the guard function quoted in
\eqref{9.1:eq-core-regularity-a}.
\item[(b)] At complex seeds $\zeta$ and complex $B,\mathsf{t},\mathsf{s},\sigma'$, every read
object $P$ of Definition~\ref{9.1:def-core-data} has, on the cells of $\mathfrak{K}_2$ and on
the bonds crossing $S_c$, the same data as $P^{\mathrm{re}}$, namely the common real $V$. Hence
$B[P]-B[P^{\mathrm{re}}]$ vanishes identically on the cells of $\mathfrak{K}_2$ and on the bonds
crossing $S_c$, and it is explicit on the collar $\mathfrak{K}_1$.
The cut strata of \cite[(1.52)]{B-CMP122b} are auxiliary cells of the determining sets of the two
configurations into which the cut separates the configuration, and cells of none of
$\mathbf{B}''_k$, $\mathbf{B}^{\sim}(Y_1)$, $\boldsymbol{\Omega}(X^{+})$.
No bracket is taken on the cut strata, no constant attached to the seams enters, and no
conjugation correction is needed.
\item[(c)] Let $m\le h$. At a point of zone $m$ of $\mathfrak{K}$ the two seams produce four
increments, namely the four terms of
\begin{equation}\label{9.1:eq-junction-cut-1}
\begin{split}
\mathfrak{b}_m
&+B^{\mathsf{s}}C_1\delta'_k\,L^{-8(h-m-1)_+}\\
&+K_{\mathfrak{K}}\tilde{\alpha}\,L^{h-k}\,L^{-8(h-m-1)_+}
+\varsigma_h ,
\end{split}
\end{equation}
where $(x)_+:=\max\{x,0\}$: the old field against the smooth centre, the $\mathbf{T}$-step
variable and the seed, each reaching zone $m$ from the collar inward with the factor
$L^{-8(h-m-1)_+}$, and the gauge drift of the collar.
\end{itemize}
\end{remark}

\begin{proof}
\emph{(a).} We use Lemma~\ref{9.1:lem-core-gauge}, namely its bounds on the core and on the
collar in \eqref{9.1:eq-core-gauge-a}, applied on the root cube. Let $\square_0$ be the root guard
cube (a cube of the guard function quoted in \eqref{9.1:eq-core-regularity-a})
of the matching in the proof of Lemma~\ref{9.1:lem-core-gauge}, the
cube normalised at a
collar centre adjacent to the towers. Its local function $U_{h,\square_0}((1,V_h))$ has data $1$
on the cells on the collar side. On the cells inside $S_c$, and on the bonds crossing $S_c$, its
data are $V_h$ \cite[text of (1.88)]{B-CMP122a}.

The axial gauge of this local function presents all of these data as $e^{iQ}$ with
$|Q|\le 2a_h$; this is the display for one cube in the proof of
Lemma~\ref{9.1:lem-core-gauge}. The actual collar datum of the real object $U^{0,\mathrm{re}}$
(Definition~\ref{9.1:def-core-data}) is $D^0$. In the
frame $\mathrm{AL}$ it differs from the guard's datum $1$ by at most
\begin{equation*}
\mathfrak{q}_h=34dK_0\tilde{\alpha}L^{h-k}.
\end{equation*}
This is the raw size bound (i) of \eqref{9.1:eq-core-data-a}, taken at $\zeta=U_c$. Therefore, on
the collar-side cells of $\square_0$ the actual datum is presented as $e^{iQ}$ with
$|Q|\le 2a_h+\mathfrak{q}_h$, up to products of small terms, while on the bonds crossing $S_c$ the
guard's data and the actual data agree. Thus the junction costs at most
$\mathfrak{q}_h+2a_h$ per bond, and no further term enters.

In the axial frame of \cite[(1.81)]{B-CMP122a}, the collar datum is the block average at level
$h$ (Definition~\ref{1.5:def-block-average}) of $U_0$ taken in that axial gauge, written
$U_0^{\mathrm{ax}}$. On
the real plateau configuration of the axial plateau proposition, the logarithm of this datum is
the axial holonomy $\mathsf{m}_{\mathrm{ax}}$ of that proposition, and by that proposition
\begin{equation*}
|\log\mathsf{m}_{\mathrm{ax}}|\ge(t_{\mathrm{ax}}-\mathfrak{c}_{\mathrm{ax}})\,
f_{\mathrm{ax}}\,\alpha_{0,k},
\end{equation*}
with $t_{\mathrm{ax}}\ge M$.
The junction cost is then at least $\frac12 f_{\mathrm{ax}}M\alpha_{0,k}$, so that
$\mathfrak{b}_h\ge\frac12 f_{\mathrm{ax}}M\alpha_{0,k}$ up to small terms. Hence the smallness
condition on the core increments of Remark~\ref{9.1:rem-frame-size} cannot be satisfied at $S_c$
in that frame.

In Ba{\l}aban's construction, a change of the gauge of $U_0$ rotates $V_h$ with it on the
boundary layer of the cells of $V_h$ at $S_c$ \cite[\S1]{B-CMP122a}.
For this reason, in the frames $\mathrm{AL}$ and $\flat$ the guard's data and the
collar data fit together.

\emph{(b).} We use Lemma~\ref{9.1:lem-data-increments}, in its bounds on the deep core and on
the cells of $\mathfrak{B}_0$ in \eqref{9.1:eq-data-increments-a}. For every read object $P$, the
data on the cells of $\mathfrak{K}_2$ and on the bonds crossing $S_c$ are the common real $V$, and
they are identical for $P(\zeta,B,\dots)$ and for $P^{\mathrm{re}}$. The reason is that the
transformations of the integrand act on the bonds of $\mathfrak{B}_0$ only
\cite[\S1]{B-CMP122a}.

The parameters therefore enter only through the data $e^{ig_kB}V_0^{\mathrm{AL}}(\zeta)$ on
$\mathfrak{B}_0$. Since the data on the cells of $\mathfrak{K}_2$ and on the bonds crossing $S_c$
are common, $B[P]-B[P^{\mathrm{re}}]=0$ there identically. On the collar $\mathfrak{K}_1$ the
parameters enter through $e^{ig_kB}V_0^{\mathrm{AL}}(\zeta)$, which by (i) and (ii) of
\eqref{9.1:eq-core-data-a} is raw-close to $V_0^{\mathrm{AL}}(U_c)$, and the difference is
explicit: it is the logarithm of the ratio of $e^{ig_kB}V_0^{\mathrm{AL}}(\zeta)$ to
$V_0^{\mathrm{AL}}(U_c)$, conjugated by the real gauge value $g(c_-)$ of
Lemma~\ref{9.1:lem-core-gauge}, and this conjugation costs at most $2\varsigma_h$, as in the proof
of Lemma~\ref{9.1:lem-data-increments}. No bracket on the cut strata, no constant attached to the
seams and no conjugation correction at the seams is needed.

The cut strata of \cite[(1.52)]{B-CMP122b} are auxiliary cells of the determining sets
$\mathbf{B}_1$, $\mathbf{B}_2$. These are the sets of the two configurations $U_1$, $U_2$ into
which the cut at $\partial\Omega''^{\sim}_{h+1}$ separates the configuration. By the reading rule
for cut sets, no term is read on a cut set. The cut strata are cells of none of
$\mathbf{B}''_k$, $\mathbf{B}^{\sim}(Y_1)$, $\boldsymbol{\Omega}(X^{+})$, by the filing rule for
strata, under which auxiliary strata are not entered into these sets while label strata are.
Moreover, by the local-level rule for strata, strata carrying data $1$ do not lower the local level.

Near the seams, the read object $U^0_k$ of Definition~\ref{9.1:def-core-data} depends on a
complex seed through the response of the global minimiser to
the data increments on the collar and on the towers; no bracket at the seams enters this response.
The response is written as a field by the local Lipschitz property
$\natural\mathrm{Lip}^{\mathrm{loc}}$ of the Landau chart about $U_c$, in
Corollary~\ref{9.1:cor-seed-rotation}.

\emph{(c).} The four terms of \eqref{9.1:eq-junction-cut-1} arise as follows.
\begin{itemize}
\item $\mathfrak{b}_m$ is the increment of the old field against the smooth centre. It is real and
common (Lemma~\ref{9.1:lem-core-gauge}), and its bound is of degree $2r_{pr}+p_0-q$ in the degree
variable $T$.
\item $B^{\mathsf{s}}C_1\delta'_k\,L^{-8(h-m-1)_+}$ is the contribution of the variable of the
$\mathbf{T}$-step, from the collar inward.
\item $K_{\mathfrak{K}}\tilde{\alpha}L^{h-k}L^{-8(h-m-1)_+}$ is the contribution of the seed,
which enters through the tower data of the collar.
\item $\varsigma_h$ is the gauge drift of the collar, bounded in
Lemma~\ref{9.1:lem-core-gauge}. Its two terms are of the types $\check{R}_h^2\varepsilon_h$ and
$\check{R}_hK_0\tilde{\alpha}L^{h-k}$.
\end{itemize}

The seed's contribution follows the route of the estimate of the seed's readings at level $h$ in
the core region: the seed reaches
$\mathfrak{K}$ only through the readings
at level $h$, which that estimate bounds by a constant multiple of $L^{h-k}\mathsf{s}_L$,
$\mathsf{s}_L$ the seed size. By the same estimate these readings
are damped layer by layer from the collar inward, and this damping gives the factor
$L^{-8(h-m-1)_+}$. By that estimate, too, the term $g_kB$ enters with $\delta'_k$ from the collar
inward and is damped in the same way.

The margin at the real point comes from the regularity \eqref{9.1:eq-core-regularity-a} of the
common core variables together with Lemma~\ref{9.1:lem-core-gauge}. It is not taken from
\cite[(1.96)]{B-CMP122a}. That display bounds, on the whole domain
$(\Omega''^{\sim}_{h+1})^c\cap\Omega_h$, the deviation from $1$ of the plaquette variables of the
local functions by $\frac12\varepsilon_h(L^{k-h}\eta)^2$ times the factor printed there; this
bound is the guard itself, which \eqref{9.1:eq-core-regularity-a} quotes.
\end{proof}

\begin{definition}[The bare datum on the torus and the stop depth]\label{9.1:not-stop-depth}
Throughout this chapter $g>0$ is the reference coupling, $c_0$ and $c_2$ are the constants of the
coupling flow listed in the glossary (Notation~\ref{0.2:not-glossary}), and $\mathfrak{p}$ is the
tuple of Ba{\l}aban's auxiliary parameters.

\emph{Letters.} We put
\begin{equation}\label{9.1:eq-stop-depth-1}
X:=g^{-2},\qquad \theta_r:=X+c_0r\quad(r\ge0),\qquad \lambda:=\frac{c_0}{2}.
\end{equation}
The leading coupling increment $\beta^0_l$ at level $l$ is used in the form
\begin{equation}\label{9.1:eq-stop-depth-2}
\beta^0_l=c_0+s_{l-1}-s_l,\qquad |s_l|\le c_2 ,
\end{equation}
where $(s_l)$ is a sequence of numbers bounded in modulus by $c_2$; this is the form given by the
statement on the leading coupling increments, which this chapter uses as stated.

\emph{Depth and length.} The \emph{depth} $r\ge0$ of a lattice is the number of levels by which it
is finer than the lattice of spacing $\ell_g$, the reference spacing attached to the reference
coupling $g$. A \emph{run} of length $n$ is the sequence of
renormalisation steps started from the bare lattice of spacing $L^{-n}\ell_g$; its \emph{length}
$n$ is its cutoff, and the bare lattice lies at depth $n$.

\emph{The bare datum.} For integers $m_0\ge0$ and $n$, and a bare inverse coupling $x$, the
\emph{bare datum} $(m_0+n,x)$ is the weight
\begin{equation*}
\exp\bigl[-A(x-w,U)\bigr]
\end{equation*}
on the periodic lattice with $2L^{m_0+n}$ sites per side, where $A(x-w,U)$ is the weighted Wilson
action of Definition~\ref{1.1:def-wilson-action} and $w$ is the complex source. Its
spacing is $L^{-n}\ell_g$, so its side is $2L^{m_0}\ell_g$. The lattice at depth $r$ of this torus
is one and the same lattice for all lengths $n\ge r$.

\emph{The stop depth and the constants of the last lattice.} We put
$\hat B:=\max\{B^{\sharp},3\lambda\}$, a number distinct from the constant $\bar B=C_\beta+1$ of the
glossary, and, for test functions $f,f_1,\dots,f_{\mathsf{p}}$
($\mathsf{p}\ge2$), with all lengths measured in units of $\ell_w$, and with $\ell_w$ and $N_w$ as
defined in \eqref{9.1:eq-stop-depth-3} below,
\begin{equation}\label{9.1:eq-stop-depth-a}
\begin{gathered}
\mathsf{w}(g):=\min\Bigl\{m\ge0:\ L^{t}\ge L^{12}M\bigl(\log(X+\hat B(t+1))\bigr)^{r_{pr}}
\ \text{for all integers } t\ge m\Bigr\},\\
\|f\|_{\sharp}:=\max\{\|f\|_\infty,\,40M\operatorname{Lip}f\},\qquad
\mathsf{d}:=\min_{i\ne j}\operatorname{dist}(\operatorname{supp}f_i,\operatorname{supp}f_j),\\
A:=e^{E_+^{\sharp}N_w^4},\qquad B:=e^{C^w_{low}N_w^4},\qquad \mathsf{Q}_w:=AB .
\end{gathered}
\end{equation}
The numbers $A$ and $B$ in \eqref{9.1:eq-stop-depth-a} are constants of the last lattice; in
particular $A$ is not the action of the bare datum.
We also write $D:=\operatorname{diam}\operatorname{supp}f_1$. The run stops at depth
$\mathsf{w}:=\mathsf{w}(g)$. In subscripts the stop depth $\mathsf{w}$ is abbreviated to $w$; this
subscript is unrelated to the complex source $w$. The \emph{last lattice} $\mathbb{T}_w$ is the
lattice at depth $\mathsf{w}$, and we put
\begin{equation}\label{9.1:eq-stop-depth-3}
\ell_w:=L^{-\mathsf{w}}\ell_g,\qquad N_w:=2L^{m_0+\mathsf{w}},\qquad
\mathcal{V}_w:=\mathrm{SU}(N)^{\mathrm{bonds}},
\end{equation}
so that $\mathbb{T}_w$ has $N_w$ sites per side, and $\mathcal{V}_w$ is its configuration space, the
product of copies of $\mathrm{SU}(N)$ over the bonds of $\mathbb{T}_w$. Every length considered in
this chapter satisfies
\begin{equation}\label{9.1:eq-stop-depth-4}
n>\mathsf{w},\qquad n\ge2 .
\end{equation}
The $n-\mathsf{w}\ge1$ renormalisation steps from depth $n$ to depth $\mathsf{w}$ are run on the
torus. At the depths $r<\mathsf{w}$ only numbers attached to the whole lattice, frozen at the stop,
are carried, in the form in which the correlation theorem states them. In
\eqref{9.1:eq-stop-depth-a}, $E_+^{\sharp}$ and $C^w_{low}$ are the exponents, per site of
$\mathbb{T}_w$ (which has $N_w^4$ sites), of the constants $A$ and $B$; both are evaluated at the
inverse coupling $X_w^{+}$ of the last lattice. The number $\mathsf{Q}_w$ depends on
$(N_w,g,L,\mathfrak{p})$ only. The lemma in which $X_w^{+}$ is fixed and the theorem on the last
lattice in which the constants $A$ and $B$ are fixed are taken in the form in which they are stated.

\emph{The torus exponent.} No upper bound is placed on $m_0$. The integer $m_0$ enters
$\mathsf{Q}_w$, the radius $\upsilon_w$ of the parameter window, the constant $C_*(N_w)$ and the
threshold length $n_0$ only through $N_w$, and no upper bound is imposed on any of these four;
this is the content of the lemma on the dependence of the last-lattice constants on $m_0$, which we
use as stated.

\emph{The clock.} For a large-field region $Z$ created at step $j$ and read at level $l$, the
number $K_j(Z)$ attached to $Z$, called its \emph{clock}, is read truncated at the age $l-j$, that
is, as $\tilde K_j:=\min\{K_j(Z),l-j\}$. The truncated clock enters in the form given by the lemma
on the truncated clock.

\emph{Conditions on the auxiliary parameters.} Ba{\l}aban's free integers and exponents, entries of
$\mathfrak{p}$, are chosen after $L$, subject to the one-sided floors
\begin{equation}\label{9.1:eq-stop-depth-5}
\begin{gathered}
\delta_0RM_1\ge16\log L,\qquad 4\delta_0R_0M_0\ge(6+\theta)\log L,\qquad
\delta_{\flat}RM_1\ge(3+\theta)\log L,\quad \delta_{\flat}:=\delta_0/16,\\
\kappa_0\ge10,\qquad \delta_1M\ge\kappa_1^{+}:=L^d\kappa_1^{\sharp},\\
q_\flat:=\min\{q_0,q_1\}\ge\max\{p_0+3r_{pr}+2,\ p_1+8r_{pr}+1\}.
\end{gathered}
\end{equation}
Here $R$, $R_0$, $R_1$, $M_0$, $M_1$, $\delta_0$, $\delta_1$, $\kappa_0$, $q_0$, $q_1$ and $p_1$ are
such entries of $\mathfrak{p}$, while $\theta$ and $\kappa_1^{\sharp}$ are fixed numbers entering the
right sides; in particular $R$, $R_0$, $R_1$ are
neither the ball radius $R$ of the glossary nor radii of the parameter window. The letter $\theta$
without subscript denotes the fixed number just named, not one of the numbers $\theta_r$
of \eqref{9.1:eq-stop-depth-1}.
Further we require $r_{pr}\ge2$ (the letter $\nu:=r_{pr}$ is also used for it), and $p_1$ is taken
in the window
\begin{equation}\label{9.1:eq-stop-depth-6}
\tfrac12\bigl(p_0+(d+5)r_{pr}\bigr)<p_1\le p_0-1-\tfrac12 r_{pr} .
\end{equation}
This condition on $p_1$ is two-sided. It is a window in which $p_1$ is chosen once $p_0$ and
$r_{pr}$ are fixed; it bounds no other quantity. It is made non-empty by the floor on the degree
exponent
\begin{equation}\label{9.1:eq-stop-depth-7}
p_0\ge(d+6)r_{pr}+4 .
\end{equation}
Finally we impose the floors, the lower bound $L\ge4c_g$ on $L$, and the relation $M=R_1M_1$
(an identity, not a floor) displayed here:
\begin{equation}\label{9.1:eq-stop-depth-8}
\begin{gathered}
\kappa_1\ge\max\{3,\,2c_g(\kappa+1)\},\qquad L\ge4c_g,\qquad \delta_0M\ge8c_d(\kappa_1+1),\\
R\ge5LR_1,\qquad M=R_1M_1 .
\end{gathered}
\end{equation}
In \eqref{9.1:eq-stop-depth-8}, $c_d$ is the numerical constant, indexed by the dimension $d$, that
sets the floor on $\delta_0M$ in proportion to $\kappa_1+1$; it is fixed with the statements from
which this floor is taken, and the floor is one-sided in $\delta_0M$.

\emph{The source coefficients and the correlation theorem.} The radius of the correlation theorem
is split into two conditions: a parameter radius $r_s$, and the matching condition
\begin{equation*}
\begin{gathered}
|c(p)|\,\mathfrak{g}_m\le\varpi\quad\text{on zone } m,\\
\mathfrak{g}_j:=\frac{(\log\mathsf{x}_j)^{2q_\flat}}{\mathsf{x}_j},\qquad
\mathsf{x}_j:=\theta_{n-j},
\end{gathered}
\end{equation*}
where $c(p)$ is the source coefficient at the plaquette $p$, $\varpi$ is the matching radius, and
the zones, indexed by the level $m$, are the level-indexed regions of the splitting of the radius,
fixed where that splitting is stated; on
the class $\mathfrak{S}$ of tuned runs the matching radius is $\varpi_{\mathfrak{S}}:=\varpi_\natural$,
the number fixed by the conditions on the matching radius that accompany this splitting of the
radius.
For a healed label the coupling profile $\varphi_j$ is reset on $Z\cup\mathsf{R}$, where $Z$ is the
large-field region of the label and $\mathsf{R}$ is the ring about $Z$ that is specified together
with the healed label; without this reset
a ring about the region would keep an old coupling at all later steps, which would amount to an
upper bound on $n$. In the correlation theorem the factor $\sigma_0$ is replaced by a factor
$\sigma_0^{\sharp}$ that does not depend on $g$ and is fixed once in the correlation theorem, in
both of the settings it compares; and the
condition on the couplings in the setting is read as $x_j\ge X$, never with the bare inverse
coupling $x_0$.
\end{definition}

\begin{proof}[Proof that $\mathsf{w}(g)$ is well defined and that the window for $p_1$ is non-empty]
Let $\mathcal{M}_g$ be the set of integers $m\ge0$ such that the inequality in the first line of
\eqref{9.1:eq-stop-depth-a} holds for every integer $t\ge m$.

The set $\mathcal{M}_g$ is closed upward. Indeed, if $m\in\mathcal{M}_g$ and $m'\ge m$, then every
integer $t\ge m'$ also satisfies $t\ge m$, so the inequality holds at $t$.

The set $\mathcal{M}_g$ is non-empty. Since $\hat B\ge3\lambda=\tfrac32c_0>0$, the number
$X+\hat B(t+1)$ tends to infinity linearly in $t$. Hence, for all large $t$, the number
$\log(X+\hat B(t+1))$ is positive and of order $\log t$, and the right side
$L^{12}M(\log(X+\hat B(t+1)))^{r_{pr}}$ grows like a fixed power of $\log t$. Since $L>1$, the left
side $L^t$ grows geometrically in $t$. The inequality therefore holds for all $t$ beyond some
integer $t_0$, and this $t_0$ belongs to $\mathcal{M}_g$.

A non-empty set of non-negative integers has a least element, so $\mathsf{w}(g)=\min\mathcal{M}_g$
exists.

For the window \eqref{9.1:eq-stop-depth-6}, assume \eqref{9.1:eq-stop-depth-7} and put
\begin{equation*}
p_1:=p_0-1-\lceil r_{pr}/2\rceil .
\end{equation*}
The right-hand inequality of \eqref{9.1:eq-stop-depth-6} holds because
$\lceil r_{pr}/2\rceil\ge r_{pr}/2$. For the left-hand inequality, note that, since
$\lceil x\rceil<x+1$ for every real $x$, we have $\lceil r_{pr}/2\rceil<r_{pr}/2+1$. Hence
\begin{equation*}
p_1>p_0-2-\tfrac12r_{pr}.
\end{equation*}
It therefore suffices that
\begin{equation*}
p_0-2-\tfrac12r_{pr}\ge\tfrac12p_0+\tfrac12(d+5)r_{pr},
\end{equation*}
that is,
\begin{equation*}
\tfrac12p_0\ge2+\tfrac12(d+6)r_{pr},
\end{equation*}
and this is \eqref{9.1:eq-stop-depth-7} multiplied by $\tfrac12$.
So this $p_1$ lies in the window, and the window is non-empty.
\end{proof}

\begin{definition}[Integer floors after $L$ and the two lengths unbounded in the cutoff]
\label{9.1:not-integer-floors}
We work in the setting of Notation~\ref{9.1:not-stop-depth}, with $d=4$. The words \emph{floor}
and \emph{ceiling} are used as in Notation~\ref{2.3:not-stages-first}, and \emph{cap} as in
Definition~\ref{2.3:def-cap}.

\emph{Integers and radii.} Two free integers of Ba{\l}aban's construction are chosen after $L$
and are written with a superscript: $R^{pr}$ and $R_0^{pr}$. They enter the floors below together
with the integers $M_0$ and $M_1$. The
letter $M$ denotes the side of the cubes of the partitions used for the runs. The quotient
\begin{equation*}
R_1^{116}:=M/M_1
\end{equation*}
is a derived quantity and not a free integer. In \cite{B-CMP116} the integer here written
$R_1^{116}$ is a power of $L$; the least admissible value of that integer in \cite{B-CMP102b} is
another number, not used here.

These integers are distinct from the \emph{complex radii}
\begin{equation*}
R:=4r_J,\qquad R_0:=\tfrac{3}{64}\,r_J,\qquad R_1:=\frac{R_0}{2\mathsf{Q}_w+2}.
\end{equation*}
Here $r_J$ is the radius of the complex sources in the correlation theorem, and $\mathsf{Q}_w$ is
the constant of \eqref{9.1:eq-stop-depth-a}; the radius $R$ is not the ball radius $R$ of
Notation~\ref{0.2:not-glossary}. In the order of choice, $M_1$, $M$, $R^{pr}$ and
$R_0^{pr}M_0$ are among the integers, and $r_J$, $R$, $R_0$, $R_1$ are among the radii.

\emph{Floors on integers and exponents after $L$.} With Ba{\l}aban's constants $\delta_0$ and
$\kappa_0$, a fixed exponent $\theta$, the constants $\delta_1$ and $\kappa_1^{\sharp}$, fixed in
another statement of this chapter of which the third line below is a hypothesis, and
$\delta_\flat:=\delta_0/16$, we impose the following conditions.
\begin{equation}\label{9.1:eq-integer-floors-a}
\begin{gathered}
\delta_0R^{pr}M_1\ge 16\log L,\qquad 4\delta_0R_0^{pr}M_0\ge(6+\theta)\log L,\\
\delta_\flat R^{pr}M_1\ge(3+\theta)\log L,\qquad \kappa_0\ge 10,\\
\delta_1M\ge\kappa_1^{+}:=L^{d}\kappa_1^{\sharp},\\
q_\flat:=\min\{q_0,q_1\}\ge\max\{p_0+3r_{pr}+2,\;p_1+8r_{pr}+1\},\\
r_{pr}\ge 2,\qquad \nu:=r_{pr},\qquad
\tfrac12\bigl(p_0+(d+5)r_{pr}\bigr)<p_1\le p_0-1-\tfrac12 r_{pr}.
\end{gathered}
\end{equation}
The exponents $q_0$ and $q_1$ are integers greater than $1$, as in \cite[(2.28)]{B-CMP119}.

The chain of inequalities in the last line of \eqref{9.1:eq-integer-floors-a} is Ba{\l}aban's
two-sided window for $p_1$; the same line fixes $r_{pr}\ge2$ and $\nu:=r_{pr}$. The window is
non-empty whenever
\begin{equation*}
p_0\ge(d+6)r_{pr}+4,
\end{equation*}
with the choice $p_1:=p_0-1-\lceil r_{pr}/2\rceil$, and in this sense it is a floor on $p_0$. The
third line is the floor on $M$ required by the statement in which $\delta_1$ and
$\kappa_1^{\sharp}$ are fixed. Further lower
bounds on $q_\flat$ are added later, by the bracket estimates only; they accumulate with the one
above.

\emph{Fat core.} We impose
\begin{equation*}
\kappa_1\ge\max\{3,\,2c_g(\kappa+1)\},\qquad L\ge 4c_g,\qquad
\delta_0M\ge 8c_d(\kappa_1+1),
\end{equation*}
where $c_d$ is a constant depending on the dimension $d$ only.

\emph{Thickness.} The condition $R^{pr}M_1\ge 5LM$ is imposed on the run sequences only. The
condition $2M\check{R}_j\ge 5LM$ is a ceiling on $g$. The shells
$\boldsymbol{\Omega}(Z)$, of width $(2L-3)M_1$, do not involve the condition $R^{pr}M_1\ge 5LM$.
Ba{\l}aban's own instance of this condition is \cite[\S F]{B-CMP102b}, where the cube side is
$M'R_1M_1$, with $M'$ an integer of that section and $R_1$ there Ba{\l}aban's integer and not the
radius $R_1$ above. The partition of unity used for the runs requires only that the region
$\Lambda_{i+1}$ have width at least $1.5\,\ell_i$, with $\ell_i$ the length scale of level $i$ of
that partition.

\emph{Further floors read by the chain estimates.} We further impose
\begin{equation}\label{9.1:eq-integer-floors-1}
\begin{gathered}
\beta\le\tfrac14,\qquad \ell_\circ^2\le L/3,\qquad L\ge 3,\qquad
\delta M/M_1\ge 2+\theta'\log L,\qquad \kappa\ge10\log L,\\
L>2(1+4\beta)(c^{\sharp})^2,\quad c^{\sharp}:=c_g+1,\qquad
\kappa\ge c^{\sharp}C(d)/\delta,\qquad L\ge 7,\\
\delta_0R^{pr}M_1\ge 256\log L,\qquad \delta_0R^{pr}M_1\ge 32\log L,\qquad L\ge5,\qquad
M_1\ge 8d.
\end{gathered}
\end{equation}
Here $\beta$ is the number $\beta$ of \cite{B-CMP122a}, not the coupling increment $\beta_{k+1}$;
$\delta$ is a further constant of Ba{\l}aban's construction beside $\delta_0$, $\theta'$ is one of
the index numbers of this chapter fixed before $L$, $C(d)$
is a constant depending on $d$ only, and $\ell_\circ$ is the integer written $L_0$ in
\cite{B-CMP122a}, renamed here to keep it apart from the constant $L_0=L_0(N)$ of (H2). The first
three conditions are conditions of \cite{B-CMP122a}, read through the
admissibility of the prepared large-field data, printed without proof in
\cite[pp.~178--179]{B-CMP122a}; cited as printed. The floor
$\delta M/M_1\ge 2+\theta'\log L$ is implied by the floor
\begin{equation*}
(M/M_1)\min\{\delta_0,\delta\}\ge 8\kappa\sqrt d
\end{equation*}
of the statement of this chapter that fixes the rate used below as $\vartheta_P$, together with
the floor $\kappa\ge 10\log L$, which is imposed here. The remaining floors are hypotheses of
statements of this chapter: the two floors involving $c^{\sharp}$ are hypotheses of the lemma in
which the constant $c^{\sharp}$ is introduced; $L\ge7$ is a hypothesis of a further lemma; the
floor $\delta_0R^{pr}M_1\ge 256\log L$ is a hypothesis of the statement that also imposes the
third ceiling of \eqref{9.1:eq-integer-floors-3} below; and the floors
$\delta_0R^{pr}M_1\ge 32\log L$, $L\ge5$ and $M_1\ge8d$ are hypotheses of the statement in which
they are read together with
\cite[(2.59)]{B-CMP96}. The floors with $16\log L$ and $32\log L$ are
implied by the one with $256\log L$ and are listed separately because different statements read
them. The conditions of \eqref{9.1:eq-integer-floors-1} are conditions on $L$, on integers chosen
after $L$ and on the constants $\beta$, $\kappa$ and $\ell_\circ$; none of them involves $g$, $n$
or the chain lengths.

\emph{The two lengths unbounded in the cutoff.} Let $\mathsf{x}_k=\theta_{n-k}$ be the inverse
coupling at level $k$ of the run, where $\theta_r=X+c_0r$ is the ladder of
Notation~\ref{9.1:not-stop-depth}; neither the exponent $\theta$ of
\eqref{9.1:eq-integer-floors-a} nor the function $\theta(g)$ below is this ladder. The chain of
preliminary integrations has two lengths, $N_0$ and $\mathsf{N}:=\check{R}_k$, and by the lemma
giving the growth of these lengths they satisfy
\begin{equation}\label{9.1:eq-integer-floors-2}
N_0\asymp\log_L\log\mathsf{x}_k,\qquad \mathsf{N}=\check{R}_k\asymp(\log\mathsf{x}_k)^{r_{pr}}.
\end{equation}
Both are unbounded in the cutoff. In this chapter the quantities $N^{pr}$, $N_0$ and $\check{R}_k$
occur only in the following two ways. (Here $N^{pr}$ denotes Ba{\l}aban's integer $N$ of
\cite{B-CMP122a}, written with a superscript to keep it apart from the $N$ of $\mathrm{SU}(N)$; it
satisfies \eqref{9.1:eq-integer-floors-4}.)
\begin{itemize}
\item[(i)] They occur in ceilings on the coupling of the kind used in Ba{\l}aban's papers. These
ceilings are monotone in $g$ and imposed after every constant:
\begin{equation}\label{9.1:eq-integer-floors-3}
\begin{gathered}
\varepsilon_P^N:=3\cdot 2^{2d}K_P\,N^{pr}\,V_P\,\vartheta_P\le\tfrac12,\qquad
C_\Sigma c_a^{GR}(\gamma)\le\tfrac12,\\
N^{pr}\alpha(g)\le a_T,\qquad (N_0+C)\,c_1\theta(g)\le\log 2.
\end{gathered}
\end{equation}
The first ceiling is a hypothesis of the proposition of this chapter on the terms of the chain of
preliminary integrations; the second is a hypothesis of a further statement accompanying that
proposition; the third is imposed in the statement that also
imposes the floor $\delta_0R^{pr}M_1\ge256\log L$ of \eqref{9.1:eq-integer-floors-1}; the fourth
is a hypothesis of the statement of this chapter that is used, at some places, in place of the
statement imposing the second ceiling, and there it replaces the second ceiling. Each of these
statements fixes the letters of its ceiling. The first carries a constant $K_P$, a factor $V_P$
and a rate $\vartheta_P$, bounded in
the proof below; the superscript $N$ of $\varepsilon_P^N$ records its dependence on $N^{pr}$ and
is not the $N$ of $\mathrm{SU}(N)$. The second carries constants $C_\Sigma$ and
$c_a^{GR}(\gamma)$, the latter
depending on $\gamma$. The third carries a function $\alpha(g)$ of the coupling and its threshold
$a_T$. The fourth carries a fixed constant $C$ of that statement, not one of the generic
constants $C(\cdot)$ of this definition, a constant $c_1$ and a function $\theta(g)$ of the
coupling. The length $N_0$ also occurs in Ba{\l}aban's own bound
\cite[p.~192]{B-CMP122a} comparing
$N_0\delta'_j$ with $4N_0\delta'_k$, where $N_0=O(\log g_k^{-2})$.
\item[(ii)] They occur as polylogarithmic factors $(\log\mathsf{x}_k)^{c_P}$ on the difference
constants $C_W$ and $C_*$ of the response bounds and the bracket estimates, and on the constant
$C_{TC}$ of the comparison theorem, and nowhere else. They never occur on the value constant
$C_0^{\sharp}$, which is independent of $g$ and bounded by $C(L,\mathfrak{p})$. They never occur
on $C_\sigma$ or $C_E$, whose bounds leave no room for a power of $\log\mathsf{x}_k$, by the lemma
of this book bounding these two constants. On the
constant $c_3$ of $\Lambda$ they occur only through the replacement $c_3\mapsto c_3+c_h$ at the
level at which $c_3$ is read.
\end{itemize}
An argument that needs a polylogarithmic factor on a value constant requires, in addition, the
conditions
\begin{equation*}
c_P+(d+1)r_{pr}<p_0,\qquad c_P<q_\flat-p_1-8r_{pr}-1,
\end{equation*}
imposed before $p_0$. In this chapter no such factor is used.

The first ceiling of \eqref{9.1:eq-integer-floors-3} is a ceiling on $g$, imposed after $L$, $M$,
$\kappa_1$ and $K_P$; it introduces neither a new floor nor a cap.
\end{definition}

\begin{proof}
\emph{The window for $p_1$.} Put $p_1:=p_0-1-\lceil r_{pr}/2\rceil$. Since
$\lceil r_{pr}/2\rceil\ge r_{pr}/2$, we have $p_1\le p_0-1-\tfrac12r_{pr}$. Since
$2\lceil r_{pr}/2\rceil\le r_{pr}+1$, we have
\begin{equation*}
2p_1-p_0-(d+5)r_{pr}\ge p_0-3-(d+6)r_{pr}.
\end{equation*}
The right-hand side is at least $1$ when $p_0\ge(d+6)r_{pr}+4$. Hence
$\tfrac12(p_0+(d+5)r_{pr})<p_1$, and the window is non-empty.

\emph{The floors of \eqref{9.1:eq-integer-floors-1}.} Since $L\ge3$ we have $\log L>0$, so
$\delta_0R^{pr}M_1\ge 256\log L$ implies $\delta_0R^{pr}M_1\ge 16\log L$; the latter is the first
floor of \eqref{9.1:eq-integer-floors-a}. Since $\min\{\delta_0,\delta\}\le\delta$ and
$\sqrt d=2$, the floor $(M/M_1)\min\{\delta_0,\delta\}\ge 8\kappa\sqrt d$ together with
$\kappa\ge10\log L$ gives
\begin{equation*}
\delta M/M_1\ge(M/M_1)\min\{\delta_0,\delta\}\ge 16\kappa\ge 160\log L,
\end{equation*}
and from this lower bound the floor $\delta M/M_1\ge 2+\theta'\log L$ is obtained. Each condition
in \eqref{9.1:eq-integer-floors-1}
involves only $L$, $\beta$, $\ell_\circ$, $\delta$, $\delta_0$, $\theta'$, $\kappa$, $c_g$, $d$,
$M$, $M_1$ and $R^{pr}$. None of these is $g$, $n$, $N^{pr}$, $N_0$ or $\check{R}_k$.

\emph{The ceiling $\varepsilon_P^N\le\frac12$.} We bound each factor of $\varepsilon_P^N$ in turn.
By \cite{B-CMP122a},
\begin{equation}\label{9.1:eq-integer-floors-4}
N^{pr}\le O(1)\,(\log g_k^{-2})^{\nu},
\end{equation}
and $\nu=r_{pr}$ by the last line of \eqref{9.1:eq-integer-floors-a}. By the lemma giving
\eqref{9.1:eq-integer-floors-2}, the length $\mathsf{N}=\check{R}_k$ obeys a bound of the same
form, a constant times $(\log g_k^{-2})^{r_{pr}}$; hence the estimate below also holds with
$N^{pr}$ replaced by $\mathsf{N}=\check{R}_k$. Next,
$V_P\le C(100L\check{R}_k)^d$; this contributes a factor $L^{O(d)}$ times
$(\log g_k^{-2})^{d\,r_{pr}}$. Finally, $\vartheta_P$ is the rate fixed in the
statement imposing the floor $(M/M_1)\min\{\delta_0,\delta\}\ge 8\kappa\sqrt d$, a constant
multiple of
$(\log g_k^{-2})^{p_1-q_*}$, where $q_*:=\min\{q_0,q_1\}$; by the definition of $q_\flat$ in
\eqref{9.1:eq-integer-floors-a}, $q_*=q_\flat$. Multiplying these bounds with
the constant $3\cdot2^{2d}K_P$, we obtain
\begin{equation*}
\varepsilon_P^N\le C'L^{O(d)}(\log g_k^{-2})^{e},\qquad e:=\nu+d\,r_{pr}-(q_\flat-p_1),
\end{equation*}
where the constant $C'$ is fixed once $M$, $\kappa_1$, $K_P$ and the parameters chosen before them
are fixed.

By the second member of the maximum in the fourth line of \eqref{9.1:eq-integer-floors-a}, we have
$q_\flat-p_1\ge 8r_{pr}+1$. With $\nu=r_{pr}$ and $d=4$ this gives
\begin{equation}\label{9.1:eq-integer-floors-5}
e\le r_{pr}+4r_{pr}-(8r_{pr}+1)=-3r_{pr}-1<0.
\end{equation}
Hence $\varepsilon_P^N\le\frac12$ holds as soon as $g$ is below a threshold depending on $L$, $M$,
$\kappa_1$ and $K_P$. This is a ceiling on $g$ imposed after these constants. It asks nothing new
of $L$ or of the integers, so it is not a new floor, and it bounds nothing from above except the
coupling, so it is not a cap.

The weaker condition $q_\flat\ge p_1+1$ of the statement fixing $\vartheta_P$ would not suffice
when the length is $\mathsf{N}=\check{R}_k$, which grows like $(\log g_k^{-2})^{r_{pr}}$: the same
count then gives only $e\le 5r_{pr}-1$, which is positive. The
floor on $q_\flat$ in \eqref{9.1:eq-integer-floors-a} gives \eqref{9.1:eq-integer-floors-5}.
\end{proof}

\begin{definition}[An intermediate exponent floor and the thickness of run sequences]
\label{9.1:not-run-thickness}
All conditions of Notation~\ref{9.1:not-integer-floors} are kept, and are used in the form
displayed in \eqref{9.1:eq-integer-floors-a}. This includes the lower bounds on Ba{\l}aban's
integers $M_1$, $M$, $R^{pr}$, $R^{pr}_0M_0$ and the bound $\kappa_0\ge 10$; the derived integer
$R_1^{116}:=M/M_1$ (the $R_1$ of \cite{B-CMP116}); the width condition
$\operatorname{width}(\Lambda_{n+1})\ge 1.5\,\ell_n$; and $\delta_0M\ge 8c_d(\kappa_1+1)$. As
there, the complex radii $r_J,R,R_0,R_1$ are letters distinct from these integers. There are
exactly two exceptions, (a) and (b) below. The words \emph{floor}, \emph{ceiling} and
\emph{cap} have the meaning of Notation~\ref{2.3:not-stages-first} and
Definition~\ref{2.3:def-cap}. Throughout, $T$ denotes the degree variable
$\log x_k=\log g_k^{-2}$ at the level $k$ under discussion.

\smallskip\noindent\textup{(a)} \emph{The intermediate exponent floor.} The lower bound on
$q_\flat$ in \eqref{9.1:eq-integer-floors-a}, namely
$q_\flat\ge\max\{p_0+3r_{pr}+2,\;p_1+2d\,r_{pr}+1\}$, is replaced by
\begin{equation}\label{9.1:eq-run-thickness-a}
q_\flat:=\min\{q_0,q_1\}\ \ge\ \max\bigl\{p_0+3r_{pr}+2,\ p_1+2(d+1)r_{pr}+1\bigr\}.
\end{equation}
At $d=4$ the second member equals $p_1+10r_{pr}+1$. In \cite[(2.28)]{B-CMP119} the exponents
$q_0,q_1$ are required to be integers greater than $1$, and
\eqref{9.1:eq-run-thickness-a} is a further lower bound on them. It is a floor. It is chosen
after $r_{pr}$, $p_0$, $p_1$ and before $\gamma$. No condition links $q_\flat$ with
$\kappa_1$, $M_1$, $M$, $R^{pr}$ or with the two further parameters $\kappa^{c}$,
$\kappa^{\flat}$.
Further lower bounds on $q_\flat$, wherever they occur, are added to
\eqref{9.1:eq-run-thickness-a}. The bound \eqref{9.1:eq-run-thickness-a} replaces the previous
one, in the sum form of the multiplicities described in the proof, in each of the following:
the weight estimate, in which the weight $\vartheta_P$, of order $T^{-(q_\flat-p_1)}$, is fixed,
and clause (c) of the theorem resting on it; the correlation theorem, and one further statement
used alongside it; the closure of a further theorem; and the two statements that carry the half
power $\vartheta_P^{\natural}$ of item (e). The lower bound
\begin{equation*}
p_1+(1+d+d\beta_0^{pr})r_{pr}+1
\end{equation*}
that the correlation theorem imposes on its own exponent is contained in the previous bound on
$q_\flat$, and hence in
\eqref{9.1:eq-run-thickness-a}.

\smallskip\noindent\textup{(b)} \emph{Thickness of run sequences.} Let
$\mathsf{m}_{15}:=R_1M_1$ be the product of the two integers $R_1$ and $M_1$ of
\cite{B-CMP102b}; this $R_1$ is distinct both from $R_1^{116}$ and from the complex radius $R_1$.
Only of the domains $\boldsymbol{\Omega}$ of run sequences do we require, for the thickness
functional of the box-set geometry of $\boldsymbol{\Omega}$,
\begin{equation}\label{9.1:eq-run-thickness-1}
\text{thickness of }\boldsymbol{\Omega}\ \ge\ 5L\,\mathsf{m}_{15},\qquad L\ge 5 .
\end{equation}
This replaces the thickness condition $R^{pr}M_1\ge 5LM$ of
Notation~\ref{9.1:not-integer-floors}. A condition with $5LM$ in
place of $5L\mathsf{m}_{15}$ does not hold on domains produced by the operation $\mathbf{R}$.
The shells $\boldsymbol{\Omega}(Z)$, of width $(2L-3)M_1$, never read either condition. For the
fat core we require
\begin{equation}\label{9.1:eq-run-thickness-2}
L\ \ge\ \max\{7,\ 4(1+4\beta_I)c_g^2\},\qquad
\kappa_1\ \ge\ \max\{3,\ 2c_g(\kappa+1)\},
\end{equation}
where $c_g$ depends on $d$ only and $\beta_I$ is the number of item (d). The first bound
accounts for two successive fattenings,
whose effects compound. To these we add the further floors of the fat-core lemma, with
$\kappa_1$ chosen after $(L,\kappa)$ and $M$ after $\kappa_1$.

\smallskip\noindent\textup{(c)} \emph{The two lengths of the chain.} The two lengths
$N_0\asymp\log_L\log x_k$ and $\mathsf{N}=\check{R}_k\asymp(\log x_k)^{r_{pr}}$ of the chain are
unbounded in the cutoff. Besides these we use Ba{\l}aban's number $N$ of stages, a letter
distinct from the rank $N$ of $\mathrm{SU}(N)$; its growth with the cutoff is governed by the
bound recalled in (c-i). The lengths $N$, $N_0$ and $\check{R}_k$ occur in exactly three ways.

\noindent\textup{(c-i)} They occur in ceilings on $\gamma$ of the kind used by Ba{\l}aban.
These ceilings are monotone in $g$ and are imposed after every constant; they all lie beneath
one further ceiling on $\gamma$, in which the rate $\rho$ enters, and which is stated in supply
form and is not the coupling ceiling $\gamma_U$. Concretely they are
\begin{equation}\label{9.1:eq-run-thickness-3}
\varepsilon_P^N:=3\cdot 2^{2d}K_P\,N\,V_P\,\vartheta_P\le\tfrac12,\qquad
C_\Sigma\,c_a^{GR}(\gamma)\le\tfrac12,\qquad N\,\alpha(g)\le a_T,
\end{equation}
together with the family of conditions, within the scope of item (a), in which $\vartheta_P$
multiplies multiplicities of the chain, at full and at half power, and the window condition at
half power, whose multiplicity has degree $d\,r_{pr}+0^+$ in $T$, $0^+$ denoting an
arbitrarily small positive excess. In
\eqref{9.1:eq-run-thickness-3}, the first inequality is the hypothesis of the proposition whose
smallness parameter is $\varepsilon_P^N$, in which $K_P$ is a constant and $V_P$ a volume
factor of degree $d\,r_{pr}$ in $T$; the second bounds the coefficient $c_a^{GR}(\gamma)$, in
which the Cauchy quotient $\theta'_j(g)$ of item (e) enters, against the constant
$C_\Sigma=8+6\beta_I$ of item (d), and is the hypothesis of the estimate to which that
coefficient belongs; the third is the hypothesis of one further lemma; there $\alpha(g)$ is a
function of the coupling and $a_T$ a constant. They also include the
estimate of \cite[p.~192]{B-CMP122a}, in which $N_0^2$ makes only a logarithmic
contribution, and the bound $N\le O(1)(\log g_k^{-2})^{\nu}$, with $\nu=r_{pr}$ by
\eqref{9.1:eq-integer-floors-a}, stated in \cite{B-CMP122a}, in which $N$
enters through a positive power of $\log g_k^{-2}$.

\noindent\textup{(c-ii)} They occur as polylogarithms $(\log x_j)^{c_P}$, and only on the
difference constants $C_W$ (not the constant $C_W(s)$ of the glossary,
Notation~\ref{0.2:not-glossary}) and
$C_*^{\mathsf{m}}$ in the component $R'$ of the production block $R$ of the response estimates
and in the estimates of the brackets (the terms carrying
the bracket parameter $\mathsf{t}$), and on the constant $C_{TC}$ of the comparison theorem. Each
such factor at level $j$ is
compensated by a factor at the same level $j$. For a term of the first class this factor is
its share of $e^{-p_0(g_j)}$, with $p_0(\cdot)$ the degree function. For an ejected or
second-class term it is the factor of the large-field event. No such factor ever falls on
the value constant $C_0^{\sharp}$ of \cite[(1.70)]{B-CMP122a} or on the constants $C_\sigma$,
$C_E$, whose estimates leave no room for a power of $\log x_j$. A polylogarithm on a value
constant would require
\begin{equation*}
c_P+(d+1)r_{pr}<p_0\quad\text{and}\quad
c_P<\tfrac12(q_\flat-p_1)-(d+1)r_{pr}-\tfrac12 ,
\end{equation*}
and none is used.

\noindent\textup{(c-iii)} They occur in quantities none of which bounds $N_0$, $\mathsf{N}$,
$n$, $L$, $m_0$ or an age from above. These are:
\begin{itemize}
\item the numbers $\mathsf{w}(g)$ of \eqref{9.1:eq-stop-depth-a} and $n_2(g)$, which depend on
$g$ only and are floors on $n$;
\item the range $\ell_r$, a band at the term's own level, compensated by $\varkappa'_r$;
\item the zones $n\le k-\mathsf{N}$, and the ages in the exponent of the rate. These are
controlled by geometric forgetting of old regions together with
$\rho>\max\{q_a,\vartheta^{1/2}\}$, which give, for the age sum $S_k$ of the age-sum lemma,
\begin{equation*}
S_k\le q_a^{-1}(1-q_a/\rho)^{-2}\rho^k ,
\end{equation*}
with a constant free of $k$, $\mathsf{N}$ and the ages;
\item the budget of that forgetting, of degree $(d+2)r_{pr}<p_0$ in $T$, as in
\eqref{9.1:eq-integer-floors-a};
\item the clock statement;
\item the component $R''$ of the production block $R$ of the response estimates;
\item the radius-split lemma of the tree bound;
\item clause (e) of the estimate whose floor on $L$ is $L>2(1+4\beta_I)c^{\sharp2}$;
\item the thickness condition \eqref{9.1:eq-run-thickness-1};
\item the line of complex radii;
\item the confined number $n_c$;
\item the window in $N$ of \cite{B-CMP122a}.
\end{itemize}
In the estimates written separately for each class mark $z$, the number $N$ of stages enters
with the factor $N^{-1}$ at each level separately.

\smallskip\noindent\textup{(d)} \emph{The number $\beta$, the weights, and the floors after
$L$.} We put
\begin{equation*}
\beta:=\beta_I=\tfrac15
\end{equation*}
wherever the number $\beta$ of \cite{B-CMP122a} is read. These places are the rooms of
\cite[(1.64)--(1.70)]{B-CMP122a} and \cite[(1.24), (1.89), (1.96)]{B-CMP122a}. The value
$\beta=\frac14$ given as an example in \cite{B-CMP122a} is the endpoint at which the
coefficient
\begin{equation*}
1-4\beta+\beta\bigl(2^{-(m-h-1)}+2^{-(j-m)}+2^{-(k-m)}\bigr)
\end{equation*}
of \cite[(1.64)]{B-CMP122a} is at most $2^{-\lfloor N/2\rfloor}$, where $h$, $m$, $j$, $k$ are
the step indices and $N$ is the number so denoted in \cite[(1.64)]{B-CMP122a}; this is the
content of the lemma on the endpoint
$\beta=\frac14$. With $\beta_I$ the number
$(1-4\beta_I)^{-1}=5$ enters as a constant. Accordingly the following quantities are read at
$\beta_I$: the bound $L>2(1+4\beta_I)c^{\sharp 2}$; the number $2(1+4\beta_I)c_g^2/(1-\beta_I)$;
$(1+3\beta_I)\kappa$; $(1-\beta_I)\kappa/c^{\sharp}$; $C_\Sigma=8+6\beta_I$; and the quantities
$m_H$ and $A_H$ of the estimate whose floor on $L$ is $L>2(1+4\beta_I)c^{\sharp 2}$.
The coefficient $\frac34$ of the summation estimate in which $\beta_0^{119}$ enters is the value
of $\frac13(1+\beta_0^{119})^2$ at
$\beta_0^{119}=\frac12$, and $\frac13(1+\beta_0^{119})^2\le\frac34$ for
$0\le\beta_0^{119}\le\frac12$; here $\beta_0^{119}$ is the small
number of \cite[(2.6)]{B-CMP119}, a different letter. We write $\ell_\circ$ for the integer
denoted $L_0$ in \cite{B-CMP122a}. It satisfies $\ell_\circ^2\le L/3$ and is distinct from the
constant $L_0$. The weights are
\begin{equation*}
\omega_W:=e^{-\delta_0R^{pr}M_1/32},\qquad \omega_T:=\vartheta .
\end{equation*}
Here $\omega_T$ is the weight of the norm $|\cdot|_\omega$, and the associated constant is
$C_2=4C_Dc\,\omega_W^{-4}\max\{(\vartheta L)^{\pm2}\}$, in which $C_D$ and $c$ are the two
constants of the estimate for the norm $|\cdot|_\omega$ from which $C_2$ is formed. The
conditions $\omega_W\le\vartheta/2$ and $\omega_WL^3\le\vartheta/2$, used with the weight
$\omega_W$, follow from the floor $\delta_0R^{pr}M_1\ge256\log L$ below. The floors on $L$ and
on integers
chosen after $L$ are:
\begin{equation}\label{9.1:eq-run-thickness-5}
\begin{gathered}
\delta M/M_1\ge 2+\theta'\log L,\qquad L>2(1+4\beta_I)c^{\sharp2},\qquad
\kappa\ge c^{\sharp}C(d)/\delta,\qquad L\ge 7,\\
\delta_0R^{pr}M_1\ge 256\log L,\qquad \delta_0R^{pr}M_1\ge 32\log L,\qquad M_1\ge 8d,\\
\delta_0R^{pr}M_1\ge 512\log L .
\end{gathered}
\end{equation}
Here $\theta'$ is the index letter of item (e). The first of these is implied by the hypothesis
$(M/M_1)\min\{\delta_0,\delta\}\ge 8\kappa\sqrt d$ of the weight estimate of item (a),
together with $\kappa\ge10\log L$. The last is the same as
$\delta_1R^{pr}M_1\ge 16\log L$ with
$\delta_1:=\delta_0/32$. It is used with \cite[(2.59)]{B-CMP96}, with the parameter $\alpha$
there equal to $\frac1{64}$, and
$L\ge5$. The floors also include the lower bound on $R^{pr}M_1$ that provides one shell of
width $R^{pr}M_1$ from the surface $S_b$, and the bound $\beta_0^{H}>\theta_A$, with $\theta_A$
of item (e), on the free
H\"older exponent denoted $\beta_0$ in \cite[(1.36)]{B-CMP99a} and \cite[(3.43)]{B-CMP99b},
written $\beta_0^{H}$ here to keep it apart from $\beta_0^{119}$ and $\beta_0^{pr}$. This floor
is compatible with a ceiling on $\beta_0^{H}$ of the form $\beta_0^{H}<1$, since
$]\theta_A,1[\ne\emptyset$. None of these floors reads $g$, $n$ or $N$.

\smallskip\noindent\textup{(e)} \emph{Index letters.} Let $\theta_A\in\,]\frac35,1[$ be the
exponent of the index theorem on which $\theta_1$ and $\theta_R$ below are built, and put
$\vartheta_A:=L^{-\theta_A}$. The remaining index letters are
\begin{equation}\label{9.1:eq-run-thickness-6}
\theta_1:=\tfrac12(5\theta_A-3),\qquad \theta_R:=\tfrac14\theta_1,\qquad
\vartheta_R:=L^{-\theta_R},\qquad \theta':=\tfrac12\theta_R .
\end{equation}
Here $\theta_R$ is the one-move index. The letter $\theta'\in\,]0,\frac18[$ is the classical
index on polydiscs, the index entering the rate $\vartheta_C$; it arises from one square root,
and it is the $\theta'$ of
\eqref{9.1:eq-run-thickness-5}. Any fixed $\theta'>0$ serves, since $\rho$ is chosen strictly
above $\vartheta_R$, the kernel rate $q$ and $q_a$ afterwards. The Cauchy quotient
$\theta'_j(g)$ in
$c_a^{GR}$ reads $g$; it is a ceiling letter and never an index. Neither $\vartheta_P$, the
weight of the weight estimate of item (a), of order
$(\log x_k)^{-(q_\flat-p_1)}=T^{-(q_\flat-p_1)}$, nor
its half power $\vartheta_P^{\natural}$ is ever an index. Four
letters $C_*$ occur: $C_*^{GR}:=2(c_0+C_\Sigma c_b)$, with $c_b$ a constant of the estimate to
which $c_a^{GR}$ belongs; $C_*^{\kappa}$, the corresponding constant of the $\kappa$-dependent
form of that estimate; $C_*^{\mathsf{m}}$ of
(c-ii); and $C_*(N_w)$. No inequality links them.
\end{definition}

\begin{proof}
\emph{Item (a).} At $d=4$ we have $2(d+1)r_{pr}+1=10r_{pr}+1$. Since $2(d+1)r_{pr}+1$ exceeds
$2d\,r_{pr}+1$, the bound \eqref{9.1:eq-run-thickness-a} implies the previous bound, whose
second member is $p_1+2d\,r_{pr}+1=p_1+8r_{pr}+1$. The lower bound that the correlation theorem
imposes on its exponent, displayed in item (a), is contained in the previous bound as soon as
$(1+d+d\beta_0^{pr})r_{pr}+1\le 2d\,r_{pr}+1$, for which $\beta_0^{pr}\le(d-1)/d=\frac34$ at
$d=4$ suffices; it is then contained in \eqref{9.1:eq-run-thickness-a} as well.

The previous member $2d\,r_{pr}+1$ suffices against one multiplicity
$\mathsf{l}^2\asymp\check{R}^d$, of degree $d\,r_{pr}$ in $T$, at the half power
$\vartheta_P^{\natural}\asymp\vartheta_P^{1/2}$. Such a term has degree at most
$d\,r_{pr}-\frac12(2d\,r_{pr}+1)=-\frac12$. The product $N\cdot V_P$, however, has degree
$(d+1)r_{pr}$. Under the previous bound its degree at half power is
\begin{equation*}
(d+1)r_{pr}-\tfrac12(2d\,r_{pr}+1)=r_{pr}-\tfrac12>0 .
\end{equation*}
The corresponding ceiling on $g$ would therefore fail as $g\to0$, so it would be a cap on $n$.
The multiplicities add and do not multiply: the total is
\begin{equation*}
3\cdot4^d\bigl[(N_0+3)\mathsf{l}^2+N\cdot V_P+o(1)\bigr]C_0 ,
\end{equation*}
the contributions entering through the links between consecutive levels being assigned to the
level at which they are read. Hence the degree of the total in $T$ is the
largest degree of its summands, and it suffices that each summand have negative degree. The
computation below shows this under \eqref{9.1:eq-run-thickness-a}.

\emph{Item (c).} Put $\mu_\flat:=q_\flat-p_1$. Then $\vartheta_P\asymp T^{-\mu_\flat}$, and
\eqref{9.1:eq-run-thickness-a} gives $\mu_\flat\ge 2(d+1)r_{pr}+1=10r_{pr}+1$. The quantity
$\varepsilon_P^N$ in \eqref{9.1:eq-run-thickness-3} carries $N\cdot V_P$, of degree
$(d+1)r_{pr}$, at full power. The half-power members of the family of item (a) carry
$N\cdot V_P$ at half power. The
window condition of (c-i) carries a multiplicity of degree $d\,r_{pr}+0^+$ at half power.
Their degrees in $T$ are therefore
\begin{equation}\label{9.1:eq-run-thickness-4}
\begin{aligned}
\varepsilon_P^N:&\quad (d+1)r_{pr}-\mu_\flat\le -(d+1)r_{pr}-1=-5r_{pr}-1,\\
\text{half power:}&\quad (d+1)r_{pr}-\tfrac12\mu_\flat\le -\tfrac12,\\
\text{window:}&\quad d\,r_{pr}+0^+-\tfrac12\mu_\flat\le -r_{pr}-\tfrac12+0^+ .
\end{aligned}
\end{equation}
All three are negative. The first member of \eqref{9.1:eq-run-thickness-3}, the half-power
members of the family of item (a) and the window condition are thus bounds of the form
$C\,T^{-e}\le\mathrm{const}$ with $e>0$, where $C$ depends
on $(L,M,\kappa_1,K_P)$. Such a bound holds for $g$ small enough. It is therefore a ceiling on
$g$, imposed after $(L,M,\kappa_1,K_P)$, and not a cap. The second and third members of
\eqref{9.1:eq-run-thickness-3} are ceilings on $g$ by their statements in (c-i), being monotone
in $g$. The only new floor is
\eqref{9.1:eq-run-thickness-a}.

The quantities of (c-iii) are floors on $n$, bands at a single level, or bounds whose
constants are free of $k$, $\mathsf{N}$ and the ages, as stated there. Hence none of them
bounds $N_0$, $\mathsf{N}$, $n$, $L$, $m_0$ or an age from above.

\emph{Item (d).} From $\beta_I=\frac15$ we get $1-4\beta_I=\frac15$, so
$(1-4\beta_I)^{-1}=5$. With $\beta_0^{119}=\frac12$,
\begin{equation*}
\tfrac13\bigl(1+\beta_0^{119}\bigr)^2=\tfrac13\cdot\tfrac94=\tfrac34 ,
\end{equation*}
and since $\frac13(1+\beta_0^{119})^2$ increases with $\beta_0^{119}\ge0$, it is at most
$\frac34$ for $0\le\beta_0^{119}\le\frac12$.
The hypothesis $(M/M_1)\min\{\delta_0,\delta\}\ge 8\kappa\sqrt d$ gives
$\delta M/M_1\ge 8\kappa\sqrt d=16\kappa$ at $d=4$, and with $\kappa\ge10\log L$ this is at
least $160\log L$. Since $\theta'<\frac18$ by item (e) and $L\ge7$, so that $\log L>1$,
\begin{equation*}
160\log L\ \ge\ 2+\log L\ >\ 2+\theta'\log L ,
\end{equation*}
which is the first member of \eqref{9.1:eq-run-thickness-5}.
The bound $\delta_0R^{pr}M_1\ge 512\log L$ implies the bounds with $256\log L$ and
$32\log L$. It also implies $\delta_0R^{pr}M_1\ge16\log L$ of
\eqref{9.1:eq-integer-floors-a}. Since $\omega_W=e^{-\delta_0R^{pr}M_1/32}$, the bound with
$256\log L$ gives $\omega_W\le e^{-8\log L}=L^{-8}$, hence $\omega_W\le\omega_WL^3\le L^{-5}$.
The conditions $\omega_W\le\vartheta/2$ and
$\omega_WL^3\le\vartheta/2$ therefore follow from the
bound with $256\log L$ whenever $\vartheta\ge 2L^{-5}$. Every member of
\eqref{9.1:eq-run-thickness-5} involves only $L$,
integers chosen after $L$, $d$, and constants fixed before $L$. None involves $g$, $n$ or $N$.

\emph{Item (e).} From $\theta_A\in\,]\frac35,1[$ we get $5\theta_A-3\in\,]0,2[$. Hence
$\theta_1\in\,]0,1[$, $\theta_R\in\,]0,\frac14[$ and $\theta'\in\,]0,\frac18[$.
\end{proof}

\begin{definition}[The adopted exponent floor and one-sided bounds on $\kappa_1$]
\label{9.1:not-exponent-floor}
Throughout, $d=4$. All letters, floors, ceilings and index letters of
Notation~\ref{9.1:not-run-thickness}, of the displays \eqref{9.1:eq-run-thickness-a} and
\eqref{9.1:eq-integer-floors-a}, and of the setting \eqref{9.1:eq-stop-depth-a} are kept,
with the changes and additions listed in (a)--(f) below. The conditions that these items
introduce are collected, in letters introduced in (a)--(f) below, in
\begin{equation}\label{9.1:eq-exponent-floor-a}
\begin{gathered}
(\mathrm{Q})''':\ q_\flat\ge p_0+2(d+1)r_{pr}+1,\qquad
(\mathrm{Q})^{A}:\ q_*\ge p_0+(1+d+d\beta_0^{pr})r_{pr}+1,\\
\kappa_1\ge\max\{3,\ 2c_g(\kappa+1),\ 16\kappa+1,\ 4L\kappa+1\},\\
(\mathrm{F}^{\flat})^{R}\Longrightarrow(\mathrm{F}_T)\Longrightarrow(\mathrm{F}^{\flat})
\Longrightarrow(\mathrm{F}),\\
\rho:=\vartheta_1:=\max\{q^{1/2},\ L^{-\vartheta_\sharp},\ \vartheta_{dl}^{1/2}\}.
\end{gathered}
\end{equation}

\emph{(a) The exponent floor.} Let $r_{pr}$, $p_0$, $p_1$ be the exponents of Ba{\l}aban's
large-field construction, with $p_0>p_1$ \cite{B-CMP122a}, and let $q_0$, $q_1$ be the two
exponents of Ba{\l}aban's convergent renormalization expansions \cite{B-CMP119}; $q_0$ and
$q_1$ are integers greater than $1$, and the construction leaves them free upward. Put
$q_\flat:=\min\{q_0,q_1\}$. In place of the floor
$(\mathrm{Q})''$ of \eqref{9.1:eq-run-thickness-a} we impose the floor $(\mathrm{Q})'''$ of
\eqref{9.1:eq-exponent-floor-a}, which at $d=4$ reads
\begin{equation}\label{9.1:eq-exponent-floor-1}
q_\flat\ge p_0+10\,r_{pr}+1 .
\end{equation}
Each stage of the construction integrates two variables; we write $A_j$ for the second, and
$r_A$ for the radius at which the bound on $A_j$ is proved. We write $q_*$ for the exponent
for which $r_A$ is proportional to $(\log g^{-2})^{q_*-p_0}$.
We impose the floor $(\mathrm{Q})^{A}$ of \eqref{9.1:eq-exponent-floor-a}; it is contained
in $(\mathrm{Q})'''$, in the sense that its right side is at most the right side of
$(\mathrm{Q})'''$.
Both are floors on exponents: they are imposed after
$(r_{pr},p_0,p_1)$ and before $\gamma$. Neither of them reads any of the parameters
$\kappa^{c}$, $\kappa^{\flat}$, $\kappa_1$, $M_1$, $M$, $R^{pr}$, and none of these
parameters is constrained by them. The floor $(\mathrm{Q})'''$ implies $(\mathrm{Q})''$. The
floor $(\mathrm{Q})'''$ replaces $(\mathrm{Q})''$ also where the correlation theorem reads it,
and $(\mathrm{Q})^{A}$ is added to the exponent floors read there. On run sequences the volume
factor $V$ has degree $d\,r_{pr}$ in $T:=\log g^{-2}$, because the scaffold radius at level
$k$ satisfies $\check{R}_k\le 2^{r_{pr}}L\,R_k$, where $R_k$ is the radius at level $k$ of
Ba{\l}aban's construction.

The floor $(\mathrm{Q})'''$ is a sufficient condition for what follows, not a necessary one.
The floor $(\mathrm{Q})''$ would suffice in either of two other cases: if
$\Delta:=p_0-p_1\le 5r_{pr}$, or if the radius $r_A$ carried a power of $g^{-1}$. Because
$r_A$ is proportional to $(\log g^{-2})^{q_*-p_0}$, the floor
$(\mathrm{Q})'''$ is the one adopted.

\emph{(b) One-sided bounds on $\kappa_1$.} Every condition on $\kappa_1$ is a lower bound;
$\kappa_1$ is never fixed by an equality. The parameter $\kappa_1$ is chosen after $(L,\kappa)$
subject to the second line of \eqref{9.1:eq-exponent-floor-a}. The member $4L\kappa+1$ there
is one of the restrictions of \cite[Lemma~3]{B-CMP116}. The condition $\kappa_1\ge16\kappa+1$
enters only as a member of this maximum. No upper
bound on $L$ is imposed in terms of $\kappa_1$, or in any other way.

\emph{(c) Thickness and the fat core.} For a box set $\boldsymbol{\Omega}$ we write
$\operatorname{gap}(\boldsymbol{\Omega})$ for its gap, the quantity bounded below in the
thickness condition of \eqref{9.1:eq-run-thickness-a}. The thickness condition
$(\mathrm{THK})^{N}$ of \eqref{9.1:eq-run-thickness-a} is kept,
\begin{equation}\label{9.1:eq-exponent-floor-2}
\operatorname{gap}(\boldsymbol{\Omega})\ge 5L\,\mathsf{m}_{R},
\end{equation}
and it is required of run sequences only. Here $\mathsf{m}_{R}$ is the product $R_1M_1$ of
the integers $R_1$ and $M_1$ of \cite{B-CMP116} fixed with
\eqref{9.1:eq-run-thickness-a}; this $R_1$ is not the radius $R_1$ of the parameter window in
$\hat{w}$, and $\mathsf{m}_{R}$ is a letter distinct from
$\mathsf{m}:=q_\flat-p_1$ and from $C_*^{\mathsf{m}}$. The condition is stated with
$\mathsf{m}_{R}$ rather than with $M$ because the bound with $5LM$ in place of
$5L\mathsf{m}_{R}$ fails on domains produced by the operation $\mathbf{R}$.
The fat-core floor on $L$ is
\begin{equation}\label{9.1:eq-exponent-floor-3}
L\ge\max\{7,\ 4(1+4\beta^{\star})c_g^2\},\qquad \beta^{\star}:=\max\{\beta_I,\beta_s\},
\end{equation}
where $\beta_s$ is one quarter of the margin exponent of the correlation theorem.

\emph{(d) Ceilings on $g$.} The list of ceilings on $g$ of \eqref{9.1:eq-run-thickness-a}
is changed in three respects.
\begin{itemize}
\item[(d1)] The first addition is the ceiling $(\mathrm{Q}^{+})^{A}$, of the form
\begin{equation}\label{9.1:eq-exponent-floor-4}
K_\flat\,\mathsf{r}\,N\,V\,\bigl[\vartheta+K_Q/r_A+2e^{-(\kappa_A-c_\flat)}\bigr]\le\tfrac14 .
\end{equation}
Here $N$ is the length of Ba{\l}aban's construction that is bounded by a positive power of
$\log g_k^{-2}$; it is unbounded in the cutoff, it has degree $r_{pr}$ in $T$ in these
ceilings, and it is distinct from the length $N_0$ of \eqref{9.1:eq-run-thickness-a}, from
the length $\mathsf{N}$ introduced below and from
the rank $N$ of $\mathrm{SU}(N)$. Further, $V$ is the volume factor of item (a);
$K_\flat$, $K_Q$, $c_\flat$ are constants and $\mathsf{r}$ is a factor of the bound on $A_j$
at radius $r_A$, and $\kappa_A$ is the exponent of its exponential term.
Its member $N\mathsf{r}V/r_A$ has degree $(d+1)r_{pr}-(q_\flat-p_0)$ in $T$, and this degree
is at most $-5r_{pr}-1$ under $(\mathrm{Q})'''$. Since the skin bound gives
$\kappa_A^{pr}\asymp p_0(g)^2$, where $p_0(\cdot)$ is the degree function, the quantity
\[
e^{-\kappa_A^{pr}(g)}(\log g^{-2})^{(1+d+d\beta_0^{pr})r_{pr}}
\]
tends to $0$ as $g\to0$; the condition on $\kappa_A$ that this quantity be small is therefore
a ceiling on $g$.
\item[(d2)] The second addition consists of the ceilings
$C_*^{GR}\bar{\bar u}\le a_{R'}$ and $C_\Sigma\bar{\bar u}\le a_{rip}$, where
$\bar{\bar u}$ is the global supremum of the local radii and $a_{R'}$, $a_{rip}$ are the two
radii against which it is compared.
\item[(d3)] The ceilings on $g$ are subsumed under the contraction ceiling associated with
$\rho$ in \eqref{9.1:eq-run-thickness-a}, and not under the coupling ceiling $\gamma_U$. The
exceptions are the family $(\mathrm{Q}^{+})$ and two ceilings of Ba{\l}aban: the bound of
\cite[p.~192]{B-CMP122a} in which $N_0^2$ makes only a logarithmic contribution, and the
bound \cite[(1.94)]{B-CMP122a} of $N$ by a positive power of $\log g_k^{-2}$. These are among
the ceilings contained in the coupling ceiling $\gamma_U$.
\end{itemize}

The value-polylogarithm condition of \eqref{9.1:eq-run-thickness-a} is not imposed. It would
be read only if the variant whose constant is $C_*^{\kappa}$ were used; that variant is not
used here. No condition imposed here bounds the lengths $N$, $N_0$ and $\mathsf{N}$ from
above. We put $\mathsf{N}:=\check{R}_k$, the scaffold radius at the level $k$
read. It is never the radius at the step $j$ that created a large-field region; that choice
would force an upper bound on the age $k-j$ of the region. Only the age $k-j$ of a
large-field region created at step $j$ and read at level $k$ is used, and it is unbounded.
The age enters through three counts of elapsed steps, and none of them is bounded from above.

\emph{(e) Floors on $L$ and on integers chosen after $L$.} The choices and floors on $L$ and on
the integers chosen after $L$ that accompany \eqref{9.1:eq-run-thickness-a} and
\eqref{9.1:eq-integer-floors-a} are kept: the choices $\beta:=\beta_I=1/5$, the weights
$\omega_W$ and $\omega_T$ and the integer $\ell_\circ$, and the floors $(\mathrm{F}_M)$,
$(\mathrm{F}_H)$, $(\mathrm{F}_\kappa)$, $L\ge7$, $(\mathrm{F}_T)$, $(\mathrm{F}^{\flat})$,
the floor on $RM_1$, and the floor $\beta_{\mathrm{h}}>\theta_A$ on a free H\"older exponent
$\beta_{\mathrm{h}}$ chosen later. They are kept with the following
changes. The four floors of the third line of \eqref{9.1:eq-exponent-floor-a} read
\begin{equation}\label{9.1:eq-exponent-floor-5}
\begin{aligned}
(\mathrm{F}^{\flat})^{R}&:\ \delta_0R^{pr}M_1\ge512\log L, &\qquad
(\mathrm{F}_T)&:\ \delta_0R^{pr}M_1\ge256\log L,\\
(\mathrm{F}^{\flat})&:\ \delta_0R^{pr}M_1\ge32\log L, &
(\mathrm{F})&:\ \delta_0R^{pr}M_1\ge16\log L.
\end{aligned}
\end{equation}
The floor $M_1\ge 8d$ is a separate condition and is not part of this chain. The floor
$(\mathrm{F}^{\flat})^{R}$ is imposed twice: once on $R^{pr}M_1$, for run sequences, and once
on the gap, for shells and for strata produced by the operation $\mathbf{R}$. The second form
is
\begin{equation}\label{9.1:eq-exponent-floor-6}
\delta_0\operatorname{gap}(\boldsymbol{\Omega})\ge 512\log L,
\end{equation}
a floor on $M_1$ imposed after $L$. The floor $(\mathrm{F}_\kappa)$ is imposed after
$\kappa_1$. To these floors we add the following.
\begin{itemize}
\item[(e1)] The floors $\delta M/M_1\ge3$, $\delta_1M\ge4\kappa_1$ and $L\ge8$, with $M$
chosen after $\kappa_1$.
\item[(e2)] The floor $(\mathrm{C2})^{+}$, namely $\delta_0M/16\ge 6\log L+c$, with an
additive constant $c$.
\item[(e3)] The ceiling $400LMR_*\tilde{\alpha}\le1$ on $g$.
\item[(e4)] The floor $S_u\ge C_S$ on the parameter $S_u$, where
$C_S:=1+40d(C_e+C_\Sigma)$, and the constant $C_{rim}:=8C_S^2C_X^{+}$; in the norm $\|F\|$,
in $\mathsf{S}^{+}$ and in the large-field bound, $\kappa$ is replaced by $\kappa_T$.
\end{itemize}

\emph{(f) Index letters.} The letter $\theta_A$ denotes a number in ${]}3/5,1{[}$, fixed before
$L$. In particular the floor $L^{1-\theta_A}\ge2$, and every other floor on $L$, reads only
numbers fixed before $L$. The condition $\beta_{\mathrm{h}}>\theta_A$ remains a floor on an
exponent chosen later. We put
\begin{equation}\label{9.1:eq-exponent-floor-7}
\begin{gathered}
\theta_1:=\tfrac12(5\theta_A-3),\qquad \theta_2:=\theta_1/4,\qquad
\vartheta_2:=L^{-\theta_2},\\
\theta':=\theta_2/2,\qquad \vartheta_C:=L^{-\theta'},\qquad
\vartheta_{dl}:=L^{-\theta'/2}.
\end{gathered}
\end{equation}
These letters have the following roles.
\begin{itemize}
\item $\vartheta_2$ is the one-move index at real parameters.
\item $\vartheta_C$ is the index on complex polydiscs, obtained with one square root.
\item $\vartheta_{dl}$ is the rate delivered by the index term of the response rows, one
square root of $\vartheta_C$.
\end{itemize}
A further halving, that is, $\vartheta_{dl}$ replaced by $L^{-\theta'/4}$, is made exactly
when the dilation parameter $\mathsf{s}$ is paid for through a total-variation estimate;
unless that alternative estimate is used, $\vartheta_{dl}$ has the value $L^{-\theta'/2}$ of
\eqref{9.1:eq-exponent-floor-7}.
The contraction letter $\rho$, also written $\vartheta_1$, is defined in the last line of
\eqref{9.1:eq-exponent-floor-a}. This is the choice of $\rho$ made in the setting of
clause~(v), with the rate that enters that choice read as $\vartheta_{dl}$. That choice also
requires $\rho>q_a$, and then $\rho>\max\{q,q_a,\vartheta_{dl}\}$. The letters
$\theta'_j(g)$, $\vartheta_P$,
$\vartheta_P^{\natural}$ and the four constants $C_*$ are as in
\eqref{9.1:eq-run-thickness-a}. The letters $\theta_A,\theta_1,\theta_2,\theta'$ are numbers
fixed before $L$. The letters $\vartheta_2,\vartheta_C,\vartheta_{dl},\rho$ are fixed once
$L$ and $q$ are.
\end{definition}

\begin{proof}
\emph{Step 1: $(\mathrm{Q})'''$ implies $(\mathrm{Q})''$.} The floor $(\mathrm{Q})''$ asks
that $q_\flat$ be at least each of $p_1+2(d+1)r_{pr}+1$ and $p_0+3r_{pr}+2$. For the first
member we use $p_0>p_1$ \cite{B-CMP122a}, the inequality recorded with (w6) in
\eqref{9.1:eq-integer-floors-a}:
\[
q_\flat\ge p_0+2(d+1)r_{pr}+1>p_1+2(d+1)r_{pr}+1 .
\]
For the second member, the difference of the two right sides is
\[
\bigl(p_0+2(d+1)r_{pr}+1\bigr)-\bigl(p_0+3r_{pr}+2\bigr)=(2d-1)r_{pr}-1,
\]
which is $7r_{pr}-1$ at $d=4$; it is nonnegative because the exponent $r_{pr}$ of
Ba{\l}aban's construction \cite{B-CMP122a} satisfies $(2d-1)r_{pr}\ge1$, that is,
$r_{pr}\ge1/7$ at $d=4$.
Hence
\[
q_\flat\ge p_0+2(d+1)r_{pr}+1\ge p_0+3r_{pr}+2 .
\]
At $d=4$
the right side of $(\mathrm{Q})'''$ is $p_0+10r_{pr}+1$, which is
\eqref{9.1:eq-exponent-floor-1}.

\emph{Step 2: $(\mathrm{Q})^{A}$ and $(\mathrm{Q})'''$.} The two floors have the same
constant term $p_0+1$. The coefficient of $r_{pr}$ is $1+d+d\beta_0^{pr}$ in
$(\mathrm{Q})^{A}$ and $2(d+1)$ in $(\mathrm{Q})'''$. The difference of the two coefficients
is
\[
2(d+1)-(1+d+d\beta_0^{pr})=d+1-d\beta_0^{pr}.
\]
This difference is positive because the small growth number $\beta_0^{pr}$ of
Ba{\l}aban's construction satisfies $0<\beta_0^{pr}<1$ \cite{B-CMP122a}, so that
$d+1-d\beta_0^{pr}>1$. Then the
right side of $(\mathrm{Q})^{A}$ is at most that of $(\mathrm{Q})'''$, which is the
containment stated in item (a). Both floors
are imposed in \eqref{9.1:eq-exponent-floor-a}.

\emph{Step 3: why the floor is placed on $p_0$.} Each stage integrates two variables. The
second, $A_j$, is cut by a condition in Ba{\l}aban's large-field step whose parameter
$\delta$ carries the degree function $p_0(\cdot)$, and not $p_1$. Suppose the bound on $A_j$
is proved at a radius $r_A$ drawn from the spaces of \cite[(2.28)]{B-CMP119}, which allow
\[
r_A\le (C/A_1)\,T^{q_\flat-p_0},
\]
and that $r_A$ is chosen of this maximal order.
Then the ceiling $N\mathsf{r}V/r_A$ has degree
\begin{equation}\label{9.1:eq-exponent-floor-8}
(d+1)r_{pr}-(q_\flat-p_0)
\end{equation}
in $T$. Here $N$ and $V$ together contribute $(d+1)r_{pr}$, and on run sequences $V$ has
degree $d\,r_{pr}$ by the bound on $\check{R}$ in item (a).

Under $(\mathrm{Q})''$ we only know
$q_\flat-p_0\ge 2(d+1)r_{pr}+1-\Delta$. The degree \eqref{9.1:eq-exponent-floor-8} is
therefore at most $\Delta-(d+1)r_{pr}-1=\Delta-5r_{pr}-1$, and this bound is negative exactly
when $\Delta<5r_{pr}+1$, in particular when $\Delta\le5r_{pr}$.

Under $(\mathrm{Q})'''$ we have $q_\flat-p_0\ge 2(d+1)r_{pr}+1$, so
\eqref{9.1:eq-exponent-floor-8} is at most $-(d+1)r_{pr}-1=-5r_{pr}-1$. The same inequality
covers the half-power row:
\[
(d+1)r_{pr}-\tfrac12(q_\flat-p_0)\le(d+1)r_{pr}-(d+1)r_{pr}-\tfrac12=-\tfrac12 .
\]
This also gives the degree bound stated for the member $N\mathsf{r}V/r_A$ of
\eqref{9.1:eq-exponent-floor-4}.

\emph{Step 4: the degrees of the ceilings.} Put $\mathsf{m}_0:=q_\flat-p_0$. By
$(\mathrm{Q})'''$, $\mathsf{m}_0\ge10r_{pr}+1$. Put $\mathsf{m}:=q_\flat-p_1=\mathsf{m}_0+\Delta$,
and note $\mathsf{m}>\mathsf{m}_0$ because $\Delta=p_0-p_1>0$. At $d=4$ the four ceilings
have the following degrees in $T$.
\begin{align*}
&\text{the ceiling }\varepsilon_P^N:&
5r_{pr}-\mathsf{m}&<5r_{pr}-\mathsf{m}_0\le-5r_{pr}-1,\\
&\text{the family }(\mathrm{Q}^{+})^{\natural,N}:&
5r_{pr}-\tfrac12\mathsf{m}&<5r_{pr}-\tfrac12\mathsf{m}_0\le-\tfrac12,\\
&\text{the slot of }A_j:&
5r_{pr}-\mathsf{m}_0&\le-5r_{pr}-1,\\
&\text{the slot of }A_j\text{ at half power}:&
5r_{pr}-\tfrac12\mathsf{m}_0&\le-\tfrac12 .
\end{align*}
The half-power slot of $A_j$ enters only if the bound on $A_j$, read with the half-power
weight $\vartheta_P^{\natural}$, which is of the order of $\vartheta_P^{1/2}$, involves the
factor $r_A^{-1}$. All four degrees
are negative. Each of these quantities therefore
tends to $0$ as $g\to0$, so each condition is a ceiling on $g$. It is imposed after
$(L,M,\kappa_1,K_P,\kappa_A)$ and bounds no parameter from above. The only new floors are
$(\mathrm{Q})'''$ and $(\mathrm{Q})^{A}$, and both are imposed in
\eqref{9.1:eq-exponent-floor-a}.

\emph{Step 5: the floor chain.} All four floors in \eqref{9.1:eq-exponent-floor-5} have the
same left side $\delta_0R^{pr}M_1$. Since $L\ge8$ we have $\log L>0$, and hence
\[
512\log L\ge256\log L\ge32\log L\ge16\log L .
\]
Consequently $(\mathrm{F}^{\flat})^{R}$ implies $(\mathrm{F}_T)$, $(\mathrm{F}_T)$ implies
$(\mathrm{F}^{\flat})$, and $(\mathrm{F}^{\flat})$ implies $(\mathrm{F})$; this is the third
line of \eqref{9.1:eq-exponent-floor-a}. None of these floors involves $M_1\ge 8d$, which is
therefore imposed separately. The second form \eqref{9.1:eq-exponent-floor-6} of
$(\mathrm{F}^{\flat})^{R}$ is a floor on $M_1$ given $L$. Both forms, and the further floors
listed in item (e), read only $L$, $\kappa_1$ and integers chosen after $L$; none of them
reads $g$, the cutoff, or $N$.

\emph{Step 6: the bounds on $\kappa_1$ are one-sided, and nothing bounds $L$.}
Ba{\l}aban's conditions constrain $\kappa_1$ only from below: \cite{B-CMP116} requires
$\kappa_1$ sufficiently large, and \cite[Lemma~3]{B-CMP116} imposes the restriction
$\kappa_1\ge4L\kappa+1$. Its other occurrences, such as $\delta_1M\ge\kappa_1$, are conditions
on later parameters, monotone in $\kappa_1$, and are met by choosing $M$ after $\kappa_1$.
Every quantity used later reads $\kappa_1$ monotonically:
\begin{itemize}
\item the constants of the form $e^{C_2\kappa_1}$ are fixed before the ceilings on $g$;
\item $M$ is chosen after $\kappa_1$, subject to $\delta_1M\ge4\kappa_1$ and
$\delta_0M\ge8c_d(\kappa_1+1)$.
\end{itemize}
Now fix $L$ and $\kappa$ satisfying their floors. Choose $\kappa_1$ at least the maximum in
\eqref{9.1:eq-exponent-floor-a}, and then choose $M$ and every later parameter. Each
condition is met, because raising $\kappa_1$ raises the lower bounds imposed on the later
parameters and lowers none of the admissible ranges. Hence no condition produces an upper
bound on $L$, and $\kappa_1$ is fixed by inequalities only.

The fat-core floor \eqref{9.1:eq-exponent-floor-3} is a floor on $L$ whichever of $\beta_I$
and $\beta_s$ is the larger. Indeed $\beta^{\star}$ and $c_g$ are numbers fixed before $L$:
$c_g$ depends on $d$ only.

\emph{Step 7: the index letters.} Since $3/5<\theta_A<1$ we have $0<5\theta_A-3<2$, hence
$0<\theta_1<1$. Consequently
\[
0<\theta'=\theta_2/2=\theta_1/8<\tfrac18 .
\]
All four numbers $\theta_A,\theta_1,\theta_2,\theta'$ are determined by $\theta_A$, which is
fixed before $L$. Since $L>1$, each of $\vartheta_2,\vartheta_C,\vartheta_{dl}$ in
\eqref{9.1:eq-exponent-floor-7} lies in ${]}0,1{[}$; each is determined by $L$ and $\theta_A$.
Moreover
\[
\vartheta_{dl}=L^{-\theta'/2}=\vartheta_C^{1/2}.
\]
By the definition of $\rho$, we have $\rho\ge\vartheta_{dl}^{1/2}>\vartheta_{dl}$ because
$0<\vartheta_{dl}<1$, and $\rho\ge q^{1/2}>q$ for the rate $q\in{]}0,1{[}$. The inequality
$\rho>q_a$ is a requirement of the choice of $\rho$ in the setting of clause~(v), and is
carried from there. The three inequalities give $\rho>\max\{q,q_a,\vartheta_{dl}\}$, as
required by the lemma bounding the age-weighted sum $S_k$ by
$q_a^{-1}(1-q_a/\rho)^{-2}\rho^k$, with a prefactor independent of $k$, of $\mathsf{N}$ and of
the ages. In
particular the case $\rho<\vartheta_{dl}$ does not occur. The remaining conditions of the
choice of $\rho$ in the setting of clause~(v) are monotone in $\rho$, and the conditions read
after $\rho$ are ceilings on $g$. Hence the constants $c_\rho$, $C_T$, $T_\Pi$ of the choice
of $\rho$ do not depend
on the cutoff $n$, and no upper bound on any parameter and no condition reading the cutoff
arises. Finally, $\rho$ is a function of $q$, $L$, $\theta_A$ and $\vartheta_\sharp$:
$\vartheta_{dl}$ is a function of $L$ and of $\theta_A$, which is fixed before $L$, and
$\vartheta_\sharp$ is a rate fixed with $L$. So $\rho$ is fixed once $L$ and $q$ are.
\end{proof}

\begin{definition}[The natural-vertex hypothesis]\label{9.1:def-natural-vertex}
In this definition $\mathbf{U}$ is the complex configuration and $\mathbf{J}=J(\mathbf{U})$ is the
complex current attached to it. The letters $\Delta$, $\Delta_a$, $\Delta^{(2)}$, $H$ and $Q$
denote operators of \cite{B-CMP102b}, with $Q^*$ the adjoint of $Q$; in this definition $H$ is this
operator of \cite{B-CMP102b}, not the Hamiltonian of Notation~\ref{0.2:not-glossary}. The operators
$G_1$, $H_1$ and $(QG_1Q^*)^{-1}$ are those of \cite[(180)]{B-CMP102b}, where
\begin{equation*}
G_1=(\Delta_a-\Delta^{(2)})^{-1}.
\end{equation*}
The letter $\mathbf{G}_1$ denotes the corresponding operator of the natural-interpolation lemma, and
$T^{\mathrm{nat}}=T^{\mathrm{nat}}(Y_0)$ is the natural-interpolation term of that lemma, with
$\delta T^{\mathrm{nat}}$ its variation. The parameter that appears in the exponents of the kernel
bounds below is written $N_{\mathrm{nat}}$, to distinguish it from the $N$ of $\mathrm{SU}(N)$.

The \emph{natural-vertex hypothesis} is the conjunction of the following four conditions.
\begin{enumerate}
\item Clauses (a), (b) and (d) of the base statement on natural vertices hold.

\item The flat statements hold, and so does the accompanying corollary. The flat statements are the
flat lemma, the flat proposition, the flat localisation statement and the flat form of clause (c) of
the base statement. Clause (a) of the flat lemma is the identity
\begin{equation*}
\mathbf{G}_1\Delta H=H-H_1+\mathbf{G}_1\Delta^{(2)}H .
\end{equation*}
Clauses (a), (d) and (v) of the flat lemma hold at $\mathbf{J}=J(\mathbf{U})$. Off $\mathbf{J}$ the
flat system is, by definition, the extension from $\mathbf{J}$.

\item Either the kernel bounds described below hold, or the by-pass statement holds. The by-pass
statement is the estimate \cite[(3.9)]{B-CMP109} taken at complex values of the parameter
$\mathsf{t}$ of the natural-interpolation lemma.

The kernel bounds concern the first case of \cite[(3.5)]{B-CMP109}, together with
\cite[(3.21)]{B-CMP109}, and only these. In \cite[p.~271]{B-CMP109} the condition
$\operatorname{dist}\ge M$ is measured on the length scale $\eta$ of \cite{B-CMP109};
\cite{B-CMP119} measures it on the same scale. In what follows $\kappa$ and $\kappa_1$ are
Ba{\l}aban's auxiliary parameters, $X$ is a large-field region (not the inverse coupling $g^{-2}$)
and $d_j(X)$ is as in Notation~\ref{0.2:not-glossary}; the constants $E_0$, $\alpha_2$, $B_3$,
$B_0$ are constants of Ba{\l}aban's construction in \cite{B-CMP109}; the number $\sigma_j$ carries
the same index $j$ as $d_j(X)$; the numbers $\sigma_j$, $\theta$ and the quantity $B'$ enter as
parameters of the bounds below; and $\delta$ is the number by which the floor below lowers the rate
$2\kappa$ to $2(1-\delta)\kappa$. Put
\begin{equation*}
\mathsf{N}_X=E_0\,e^{-\kappa d_j(X)}\,6\alpha_2^{-1}B_3\cdot 8B_0^{\mathsf{t}}\,|B'| .
\end{equation*}
The kernel bounds are
\begin{equation}\label{9.1:eq-natural-vertex-a}
\begin{aligned}
|T^{\mathrm{nat}}|&\le\bigl(2e^{-(\kappa_1-1)}\bigr)^{N_{\mathrm{nat}}/2}\,\sigma_j^{1/2}\,
\mathsf{N}_X,\\
|\delta T^{\mathrm{nat}}|&\le 2\,\theta^{1-\beta'-\gamma'}\,\sigma_j^{\beta'}\,
e^{-[\gamma'(\kappa_1-1)-(1-\gamma')\log 2]\,N_{\mathrm{nat}}}\,\mathsf{N}_X .
\end{aligned}
\end{equation}
The first bound in \eqref{9.1:eq-natural-vertex-a} rests on \cite[(190)]{B-CMP102b}, taken at
vertices, and on the flat form of clause (c) of the base statement. The second bound holds for
every $\beta'>0$ and every $\gamma'>0$. Its right-hand side is twice the weighted geometric mean,
with exponents $1-\beta'-\gamma'$, $\beta'$ and $\gamma'$ (which sum to one), of
$\theta\,2^{N_{\mathrm{nat}}}\mathsf{N}_X$, $\sigma_j\,2^{N_{\mathrm{nat}}}\mathsf{N}_X$ and
$e^{-(\kappa_1-1)N_{\mathrm{nat}}}\mathsf{N}_X$, namely
\begin{equation*}
\begin{split}
&2\,\bigl(\theta\,2^{N_{\mathrm{nat}}}\mathsf{N}_X\bigr)^{1-\beta'-\gamma'}\,
\bigl(\sigma_j\,2^{N_{\mathrm{nat}}}\mathsf{N}_X\bigr)^{\beta'}\\
&\qquad\times\bigl(e^{-(\kappa_1-1)N_{\mathrm{nat}}}\mathsf{N}_X\bigr)^{\gamma'} .
\end{split}
\end{equation*}
In the natural-vertex hypothesis the parameters satisfy, moreover, the floor
\begin{equation*}
\gamma'(\kappa_1-1)-(1-\gamma')\log 2\ \ge\ 2(1-\delta)\kappa .
\end{equation*}
In the second case of \cite[(3.5)]{B-CMP109} the kernel bounds are replaced by
\cite[(1.18)]{B-CMP109} alone.

\item The zone statement holds at the zone's own parameter $\hat{\varepsilon}_n$, in the form given
by the zone proposition, which derives it from the conjunction of the following:
\begin{itemize}
\item the halving of \cite[\S F]{B-CMP102b}, applied locally on both kinds of two-level cube;
\item the thickness floor, one of the floors on exponents fixed after $L$;
\item the reference statement;
\item the condition $a_5$;
\item the H\"older bounds of \cite[(3.133)]{B-CMP99b}, used as stated there.
\end{itemize}
Taken in this form, at $\hat{\varepsilon}_n$, the zone statement lets the small-field statement
keep $c_2^{\mathrm{IS}}\ge 2p_0$, where $c_2^{\mathrm{IS}}$ is the constant of the small-field
statement and $p_0$ is the exponent of the degree function $p_0(\cdot)$.
\end{enumerate}

The statements named in these four conditions are stated in terms of the following objects and
bounds, among other inputs.
\begin{itemize}
\item Bounds on the operators $G_1$, $H_1$ and $(QG_1Q^*)^{-1}$ of \cite[(180)]{B-CMP102b}, outside
the set $\mathfrak{O}$. They consist of three bounds, the third in a strengthened form.
\item The bounds \cite[(189)--(190)]{B-CMP102b}, taken at vertices.
\item The vertex statement, the response statement and its corollary.
\end{itemize}
\end{definition}

\begin{lemma}[Non-linear natural interpolation under the natural-vertex hypothesis]
\label{9.1:lem-nonlinear-interpolation}
We use the letters of the setting of the natural-interpolation lemma for the decoupled
operators, whose items are lettered (a)--(g). Fix one of the non-linear cases of that lemma,
namely the case of the first kind or one of the seven cases of the second kind. In the linear
cases the operator interpolation is kept, and those cases are not treated here.

The case comes with its data $(\boldsymbol{\Omega},U_0,\mathsf{B},S)$ and the following objects.
\begin{itemize}
\item A core $K$ that contains Ba{\l}aban's core and is fattened. In the cases of the second
kind the fattening is by the union $\bigcup\Lambda^{\sim}$, by the fattening lemma.
\item The set $\sigma_0=\pi\setminus\mathsf{K}$ of the cubes of the collection $\pi$ of the
setting that do not lie in the cube set $\mathsf{K}$ of the setting. In the cases of the second
kind it consists of unit $M$-cubes, which is a choice made here. In the case of the first kind
it consists of cubes at their proper scales, by the zone-scale statement.
\item The boxes $Z=K\cup Y'$.
\end{itemize}
On each $\boldsymbol{\Omega}(Z)$ the operators are taken at $(U_0{\upharpoonright}Z,\,J(U_0))$.
They are decoupled in the parameters $\mathsf{t}$ by Ba{\l}aban's random walks, with
$|\mathsf{t}(\Delta)|\le e^{\kappa_1+1}$.

Put $A':=H_1(\mathsf{t})\mathsf{B}+\mathcal{A}$, where the field $\mathcal{A}$ is subject to
\begin{equation}\label{9.1:eq-nonlinear-interpolation-1}
\mathcal{A}=-\mathfrak{j}_Z(\mathsf{t})-\mathbf{G}_1(\mathsf{t})V'_{\flat}(\mathsf{t})(A')
+(H-H_1)(\mathsf{t})D(\mathsf{t})(A'),
\end{equation}
and put
\begin{equation}\label{9.1:eq-nonlinear-interpolation-2}
\mathbb{H}^{\flat}_Z(\mathsf{t}):=H_1(\mathsf{t})\bigl(\mathsf{B}-D(\mathsf{t})(A')\bigr)
-\mathbf{G}_1(\mathsf{t})V'_{\flat}(\mathsf{t})(A')-\mathfrak{j}_Z(\mathsf{t}).
\end{equation}
Here $\mathfrak{j}:=0$, except in the second and the fifth of the seven cases of the second
kind. In those two cases
\begin{equation}\label{9.1:eq-nonlinear-interpolation-3}
\mathfrak{j}_Z(\mathsf{t}):=\mathbf{G}_1(\mathsf{t})(1-\theta)Q^{*}\theta_1
H_{\boldsymbol{\Omega}_1(Z)}(\mathsf{t})^{T}J_1+\mathbf{G}_1(\mathsf{t})\theta J_{1,2},
\end{equation}
as in the decoupled-operator construction, with each factor decoupled.
Here $H$, $H_1$, $\mathbf{G}_1$ are the operators and $D(\mathsf{t})$ is the map of the
decoupled-operator construction, $Q^{*}$ is the adjoint of the linearised block-averaging map,
$\theta$ and $\theta_1$ are the localising functions of that construction, and $J_1$, $J_{1,2}$
are its currents.

Assume:
\begin{itemize}
\item[(i)] the natural-vertex hypothesis, Definition~\ref{9.1:def-natural-vertex};
\item[(ii)] the strengthened vertex statement;
\item[(iii)] the zone-scale statement.
\end{itemize}
Then items (a)--(g) of the natural-interpolation lemma hold, with the following two
substitutions.
\begin{itemize}
\item The natural interpolation $\langle\mathsf{s}\rangle$ on $(\mathfrak{X},\natural)$ is
replaced by two objects. The first consists of the families $f_Z(\mathsf{t})$ of analytic
functionals of
$(\mathbb{H}^{\flat}_Z,\mathfrak{L}^{\flat}_Z){\upharpoonright}K^{\circ}$ that are bounded by
$\mathsf{N}$ on their range. The second is
\begin{equation}\label{9.1:eq-nonlinear-interpolation-4}
T^{nat}(Y_0):=\sum_{\mathsf{Y}'\subset\mathsf{Y}_0}(-1)^{|\mathsf{Y}_0\setminus\mathsf{Y}'|}
f_{K\cup Y'}(\mathbb{1}).
\end{equation}
\item The constant $B^{\natural}$ of the natural-interpolation lemma is replaced by
$(B_0^{\mathsf{t}})^3$, where $B_0^{\mathsf{t}}$ is the constant of the bounds on the decoupled
operators, taken at $\kappa_1+1$.
\end{itemize}
\end{lemma}

\begin{proof}
We take the items in order.

\emph{Item (a).} Clause (0) of the flat multiscale bounds is part of the natural-vertex
hypothesis (Definition~\ref{9.1:def-natural-vertex}). It contains the following count, for
each random walk $\omega$ of the decoupling, of length $\ell_\omega$, with $S(\omega)$ the set
of cubes attached to $\omega$ and $c_d$, $C_d$ the constants of those bounds:
\begin{equation*}
|S(\omega)|\le c_d+C_d\,\ell_\omega/M .
\end{equation*}
This count holds for unit cubes in any zone, because the scaled length is bounded below by the
Euclidean length \cite[(2.46)]{B-CMP96}. For cubes at their proper scales it holds by the
zone-scale statement, hypothesis (iii).

Clauses (0)--(ii) of the flat multiscale bounds see neither the stratum of $\mathsf{B}$ nor
$\mathfrak{j}$. Indeed, the norm entering them is the supremum over all pairs $(j,c)$ of a
level $j$ and a component $c$.
Moreover there is one geometry, and the bounds carry no loss near its boundary, as in the
decoupled-operator construction.

Write $\Phi^{\mathsf{t}}_{\flat}$ for the map
\begin{equation*}
\mathcal{A}\longmapsto-\mathbf{G}_1(\mathsf{t})V'_{\flat}(\mathsf{t})(A')
+(H-H_1)(\mathsf{t})D(\mathsf{t})(A'),
\qquad A'=H_1(\mathsf{t})\mathsf{B}+\mathcal{A}.
\end{equation*}
Then \eqref{9.1:eq-nonlinear-interpolation-1} reads
$\mathcal{A}=-\mathfrak{j}_Z(\mathsf{t})+\Phi^{\mathsf{t}}_{\flat}(\mathcal{A})$. Under these
clauses, in the norm $|\cdot|_1$ of \cite[(3.31)]{B-CMP109} and with the radius
$\varepsilon_2$ of the flat multiscale bounds, $\Phi^{\mathsf{t}}_{\flat}$ maps
$\{|\mathcal{A}|_1\le\varepsilon_2\}$ into
$\{|\mathcal{A}|_1\le\tfrac14\varepsilon_2\}$.

Let $\mathcal{A}_0$ be a point with $|\mathcal{A}_0|_1\le\tfrac12\varepsilon_2$. The ball of
radius $\tfrac12\varepsilon_2$ about $\mathcal{A}_0$ lies in the $\varepsilon_2$-ball, and on it
$|\Phi^{\mathsf{t}}_{\flat}|_1\le\tfrac14\varepsilon_2$. The Cauchy estimate therefore gives
\begin{equation*}
\|D\Phi^{\mathsf{t}}_{\flat}(\mathcal{A}_0)\|\le
\frac{\tfrac14\varepsilon_2}{\tfrac12\varepsilon_2}=\frac12 .
\end{equation*}
So $\|D\Phi^{\mathsf{t}}_{\flat}\|\le\tfrac12$ on the $\tfrac12\varepsilon_2$-ball.

Now let $\mathfrak{j}$ be analytic with $|\mathfrak{j}|_1\le\tfrac14\varepsilon_2$. On the
$\tfrac12\varepsilon_2$-ball,
\begin{equation*}
|-\mathfrak{j}+\Phi^{\mathsf{t}}_{\flat}(\mathcal{A})|_1
\le\tfrac14\varepsilon_2+\tfrac14\varepsilon_2=\tfrac12\varepsilon_2 .
\end{equation*}
Hence $-\mathfrak{j}+\Phi^{\mathsf{t}}_{\flat}$ maps that ball into itself. The ball is convex and
the derivative is bounded by $\tfrac12$ there, so the map is a contraction with constant
$\tfrac12$. It therefore has a unique fixed point $\mathcal{A}$, and this fixed point is analytic
in the variables in which $\mathfrak{j}$ and $\Phi^{\mathsf{t}}_{\flat}$ are analytic, as the
uniform limit of the analytic iterates. From
$\mathcal{A}=-\mathfrak{j}+\Phi^{\mathsf{t}}_{\flat}(\mathcal{A})$ and the contraction bound,
\begin{equation*}
|\mathcal{A}|_1\le|\mathfrak{j}|_1+|\Phi^{\mathsf{t}}_{\flat}(0)|_1+\tfrac12|\mathcal{A}|_1 .
\end{equation*}
Rearranging gives
\begin{equation}\label{9.1:eq-nonlinear-interpolation-5}
|\mathcal{A}|_1\le2\bigl(|\mathfrak{j}|_1+|\Phi^{\mathsf{t}}_{\flat}(0)|_1\bigr).
\end{equation}
Under the smallness ceilings on the radii $\varepsilon_2$ and $\varepsilon_3$ of the flat
multiscale bounds, \eqref{9.1:eq-nonlinear-interpolation-2} then gives
\begin{equation}\label{9.1:eq-nonlinear-interpolation-6}
|\mathbb{H}^{\flat}_Z|_1\le4B^{\mathsf{t}}_0|\mathsf{B}|+2|\mathfrak{j}|_1 .
\end{equation}
These bounds hold in the norms of \cite[(3.31)]{B-CMP109}. The same holds for the component
$\mathfrak{L}^{\flat}$ in the slot of the current $J'$ on $K^{\circ}$, by clause (iv) of the flat
multiscale bounds. When $\mathfrak{j}=\mathbf{G}_1\varphi_{\mathfrak{j}}$, this component acquires
the additional term $-\mathfrak{L}_{\mathbf{G}_1}(\mathsf{t})\varphi_{\mathfrak{j}}$. All
constants are independent of $Z$ and $\boldsymbol{\Omega}$.

For the current $\mathfrak{j}$ of \eqref{9.1:eq-nonlinear-interpolation-3} we have
\begin{equation*}
|\mathfrak{j}|_1\le(B^{\mathsf{t}}_0)^2\bigl[C B^{\mathsf{t}}_0|J_1|_{(-3)}+|J_{1,2}|_{(-3)}\bigr]
\le\tfrac14\varepsilon_2 .
\end{equation*}
The first inequality uses the bounds of clause (0) of the flat multiscale bounds on the
transposed operators, together with the bound of $Q^{*}$ by a constant, used as stated. The
second inequality holds by the choice of the exponent $\gamma$. So the hypothesis
$|\mathfrak{j}|_1\le\tfrac14\varepsilon_2$ used above is met.

\emph{Item (b).} Set $\mathsf{t}=\mathbb{1}$. On the configuration's own currents, the identity
(d) of the flat system holds; there $\mathcal{A}=-\mathbf{G}_1\mathbf{J}-\mathbf{G}_1V'$, as in the
decoupled-operator construction. Combining this with the definition of $\mathfrak{j}$ and clause
(c) of that construction, equation \eqref{9.1:eq-nonlinear-interpolation-1} is the box equation
\cite[(111)]{B-CMP102b}. By the uniqueness statements \cite[Propositions~5 and~6]{B-CMP102b},
$e^{i\eta\mathbb{H}^{\flat}_Z(\mathbb{1})}U_0{\upharpoonright}Z$ is the natural vertex. At $Z=T$,
where $T$ is the largest box of the setting, which contains every box $Z$, it is the function
of Ba{\l}aban's construction.

\emph{Item (c).} The four steps of the flat localisation statement remain valid with
$\mathfrak{j}$ present, for three reasons.
\begin{itemize}
\item The kernels involved are decoupled kernels and their transposes, and $\theta$,
$\theta_1$, $Q^{*}$ are local.
\item No interpolation parameter is switched on $\mathcal{L}$, nor on
$\operatorname{supp}\theta_1\subset Z''_k\subset K^{\circ}$.
\item $\boldsymbol{\Omega}_1(Z)$ and $\boldsymbol{\Omega}''(Z)$ are contained in boxes
determined by the same $Z$. This follows from the second and the strengthened third of the
properties of the operators listed among the inputs of the natural-vertex hypothesis
(Definition~\ref{9.1:def-natural-vertex}).
\end{itemize}

Hence the decomposition clauses (a)--(d) of Definition~\ref{9.1:def-natural-vertex} apply to
every $f_Z$ of the statement. This covers the terms, the readings
$\operatorname{rd}_{S_v}$ on sets $S_v\subset K^{\circ}$, and the block-local functionals
$u$. They give
\begin{equation}\label{9.1:eq-nonlinear-interpolation-7}
|T^{nat}(Y_0)|\le e^{-(\kappa_1-1)N}\mathsf{N},\qquad
\sum_{Y_0}T^{nat}(Y_0)=f_T(\mathbb{1}),
\end{equation}
and $T^{nat}(Y_0)=0$ unless $K\cup Y_0$ is wall-connected. Here $N=N(Y_0)$ is the size of
$Y_0$ that enters the decay factors of the bounds \eqref{9.1:eq-natural-vertex-a}.

Suppose now that $S$ lies off the component of $x$. Then $\mathsf{B}=0$ and $\mathfrak{j}=0$
there, so the fixed point of item (a) is $0$. This is the vanishing clause of Ba{\l}aban's
construction.

\emph{Item (d).} This is the Schwarz-lemma step of the natural-interpolation lemma. It passes
from the radius $e^{\kappa_1+1}$ in the parameters $\mathsf{t}$ to the radius $e^{\kappa_1}$ and
gives the decay $e^{-n(x,S)}$. The readings at levels $j<k$ are the raw readings of
\cite[(1.60)]{B-CMP122b}, and they give the bounds \eqref{9.1:eq-natural-vertex-a}.

\emph{Items (e) and (f).} These items concern only the vertices. Their proofs in the
natural-interpolation lemma apply without change, with the following inputs:
\begin{itemize}
\item item (b) above, in place of the vertex-structure statement used there;
\item clause (c) of the decoupled-operator construction;
\item items (a)--(d) of the vertex statement.
\end{itemize}

\emph{Item (g).} Item (g) follows from item (a). Readings directly up to the boundary
$\partial X$ are not supplied by this lemma; they are supplied by the statement on reading
units of this chapter.
\end{proof}

\begin{theorem}[The tube theorem at every site]\label{9.1:thm-tube-theorem}
Assume the following:
\begin{itemize}
\item the closed form of the hypothesis on the two machines, together with its preparation
clause, its slice clause, the floors of the summation lemma, the unit clause and the frame
relation;
\item the kernel corollary;
\item the frame estimate off the healer's core $\mathfrak{K}$, on the family of frames
$\mathfrak{F}_G$, on the region
$Z^{\prime\prime(1)}_{k}:=Z''_k\setminus\mathfrak{K}$;
\item the arrest statement;
\item the conditions on the constants: the bound on $K^{\sharp}_R$, the condition on the regular
region with the constant $c_M$, the condition on the link constants, the condition on the
aged region $\mathcal{A}$, the sharpened condition on the constants of the frame estimate and
the sharpened condition on the constants of its restricted form, and the inequality
\begin{equation*}
2C_4^{\sharp}\mathsf{s}_L\le\min\{c_4,a_R'\};
\end{equation*}
\item the circles fixed with the sites;
\item the stage-room inequality.
\end{itemize}
Let $P=(\mathbf{U},\mathbf{J})$ be a point on the pair tube $\mathfrak{P}_{\mathsf{c}}$ of a shape
at any site, lying over a seed of the output $\Phi$ of that shape, the slice of the seed lying in
$\mathfrak{D}_{\Phi}[7/8]$. Let $P$ be read by a constituent term $(F,X,j)$ such that $X$ does not
meet the raw core and such that, on $Z''_k$,
\begin{equation*}
\tilde{\Delta}(y)\cap\mathfrak{K}=\emptyset\qquad\text{for every box }y\in\mathfrak{B}
\text{ of the box tower }\boldsymbol{\Omega}(X^{+}).
\end{equation*}
(The terms for which the second of these two conditions fails are those treated by the estimate
on the healer's core.) Then
$\tilde{F}(X;P)=F(X;S(\mathbf{U}))$, where $\tilde{F}(X;P)$ is the value of $F$ read at the pair
point $P$ and $S(\mathbf{U})$ is the configuration determined by the pair variable $\mathbf{U}$,
both as in the definition of the pair tube; this function is holomorphic and exact at vanishing
increments; and
\begin{equation*}
\mathrm{rd}^{\mathsf{m}}\mathbf{U}^{\mathsf{m}}\in\mathfrak{T}^{\mathsf{m}}[0](X).
\end{equation*}
Precisely, at the (up)-points one has
$E_a(S(\mathbf{U}))\le(\tfrac34+\tfrac1{200})\mathsf{u}_F$; at the (=)-points,
\begin{equation*}
E_a(S(\mathbf{U}))\le\mathsf{T}_e+\tfrac16\bigl(\mathsf{T}_r-\mathsf{T}_e\bigr).
\end{equation*}
Here the (up)-points and the (=)-points are the two classes of points of the definition of the
tube $\mathfrak{T}^{\mathsf{m}}[\varsigma](X)$, and $E_a$ is the quantity which the tables of that
definition bound. The starting configuration $U_e$ of the increments, the entry table
$\mathsf{T}_e$, the
increment generators $\mathcal{Y}$ and the reference table $\mathsf{T}_r$ at (=)-points are, site
by site, as follows. Below, $\mathsf{N}$ is the number of stages; $n_F$ is the level of the term
$F$; $f_{k,n}$ and $f_{k,m}$ are the layer coefficients of the layers read, indexed by their
levels $n$ and $m$; $\nu$ is the stage at which a run is read and $n_e$ its entry stage;
(a)-points and (b)-points are
the points of $Z^{\prime\prime(1)}_{k}$ covered by clause~(a), respectively clause~(b), of the
frame estimate off the healer's core (in a run, the unraised, respectively the raised,
(=)-points); and the bracket term of that estimate, which together with the chain term makes
up the reference value at (a)-points, is written $\mathsf{b}\,\mathsf{u}_F$ at (b)-points.
\begin{equation}\label{9.1:eq-tube-theorem-a}
\begin{array}{|p{0.17\textwidth}|p{0.63\textwidth}|}
\hline
site & data \\
\hline
(T) &
$U_e$: the seed; $\mathsf{T}_e=\mathsf{c}_{\Phi}^{+}\mathsf{u}_{\Phi}$.\newline
$\mathcal{Y}$: the dilated shift.\newline
$\mathsf{T}_r$: $\bigl(f_{k,n}-\tfrac12(1-\theta_S)\beta\,2^{-(k+1-n)}\bigr)\mathsf{u}$.\newline
Inputs: the accounting lemma at the T-step; the schedule lemma. \\
\hline
(L), terminal &
$U_e$: the root; $\mathsf{T}_e$: on $\mathfrak{O}$, $(f_{k,m}-\tfrac18\varrho_{\Phi})\mathsf{u}$;
on $Z^{\prime\prime(1)}_{k}$, at most $\mathsf{T}_{\mathsf{N}}/32$.\newline
$\mathcal{Y}$: all links, and the $\mathsf{t}$-shift.\newline
$\mathsf{T}_r$: on $\mathfrak{O}$, $f_{k,m}\mathsf{u}$; on $Z^{\prime\prime(1)}_{k}$ at (a)-points,
$\tfrac12\mathsf{T}_{\mathsf{N}}$, the chain taking at most $\tfrac{3}{32}\mathsf{T}_{\mathsf{N}}$
and the bracket at most the rest; at (b)-points, as at the (up)-points of $\mathcal{A}$: at most
$(\tfrac{3}{32}+\mathsf{b}+\tfrac1{200})\mathsf{u}_F$.\newline
Inputs: the fourfold tube lemma, the entry lemma, the stage-room inequality; clause (c) of the
accounting lemma on $\mathfrak{O}$; the schedule lemma; the entry estimate on
$Z^{\prime\prime(1)}_{k}$ for the bound $\mathsf{T}_e\le\mathsf{T}_{\mathsf{N}}/32$; the frame
estimate off the healer's core. \\
\hline
(P)-run over case (e-ii) &
$U_e$: the root; $\mathsf{T}_e=\mathsf{T}^{\flat}_{\mathsf{N}}$.\newline
$\mathcal{Y}$: links, then moves in $[\nu,\mathsf{N}[$.\newline
$\mathsf{T}_r=\mathsf{T}^{\flat}_{\nu}$.\newline
Inputs: on $\mathfrak{O}$, the accounting lemma on $\mathfrak{O}$; on $Z^{\mathrm{on}}$, the entry
lemma and clauses (ii), (iii) of the first tube theorem. \\
\hline
(P)-run over case (e-i) &
$U_e$: the seed; $\mathsf{T}_e=\mathsf{T}^{\flat}_{n_e}$.\newline
$\mathcal{Y}$: moves.\newline
$\mathsf{T}_r=\mathsf{T}^{\flat}_{\nu}$.\newline
Inputs: as for case (e-ii). \\
\hline
box $Z\supseteq X^{+}$, case (S3) &
$U_e$, $\mathsf{T}_e$, $\mathcal{Y}$, $\mathsf{T}_r$: those of the site.\newline
Inputs: the site-descent lemma at the box tower $\mathfrak{S}:=\boldsymbol{\Omega}(Z)$. \\
\hline
\end{array}
\end{equation}
None of the constants of this theorem depends on $g$, $k$, $K$, $\mathsf{N}$, $N_0$,
$\check{R}_k$, the level $h$ of the healer's core, $r$, on the ages of large-field regions, on
the depth or on the volume.
\end{theorem}

\begin{proof}
At every site we apply the tube-membership criterion together with clause (b$'$) of the
site-descent lemma, the majorants of the increments of the several kinds being added. It
remains to verify the hypotheses of the criterion region by region.

\emph{The site} (L) \emph{on} $\mathfrak{O}$. We take the gap $\mathsf{G}=\tfrac18\varrho_{\Phi}
\mathsf{u}$ and the majorant $\mathsf{m}=2\bigl(\mathsf{m}_L+(1+|\mathsf{t}|)\mathsf{m}_p\bigr)$.
The two hypotheses of the tube-membership criterion then read
\begin{equation}\label{9.1:eq-tube-theorem-b}
\begin{gathered}
\text{(A)}\quad \text{clause (A) of the tube-membership criterion holds}\\
\text{with the majorant } \mathsf{m} \text{ and the gap } \mathsf{G};\\
\text{(S)}\quad f_{k,m}\mathsf{u}\le\mathsf{T}^{\flat}_{0}.
\end{gathered}
\end{equation}
Clause (A) follows from the condition on the link constants and the definition of the link
weight $\mathsf{w}_L$. Clause (S) is clause (c) of the accounting lemma on $\mathfrak{O}$.

\emph{The region} $\mathcal{A}$. A terminal read on $\mathcal{A}$ has (up)-points only, since
there the label level $k$ exceeds $n_F$. By the entry lemma and clause (b$'$) of the
site-descent lemma the read is bounded there by $(\tfrac16+\tfrac1{960}+\tfrac1{200})
\mathsf{u}_F$. This lies below the bound $(\tfrac34+\tfrac1{200})\mathsf{u}_F$ of the statement,
since
\begin{equation*}
\tfrac16+\tfrac1{960}=\tfrac{161}{960}<\tfrac{720}{960}=\tfrac34.
\end{equation*}

\emph{The region} $Z^{\prime\prime(1)}_{k}=Z''_k\setminus\mathfrak{K}$. Here we use the frame
estimate off the healer's core together with clause (b$'$) of the site-descent lemma, on the
terms for which $\tilde{\Delta}(y)\cap\mathfrak{K}=\emptyset$ for every $y\in\mathfrak{B}$. This
is the hypothesis of the theorem; the other terms are those of the estimate on the healer's
core. At (a)-points,
\begin{equation*}
\mathsf{T}_{\mathsf{N}}\le\tfrac34\bar{\mathsf{u}}_F\le(\mathsf{c}_F^{+}-\varrho_F)
\bar{\mathsf{u}}_F,
\end{equation*}
the comparison being with the table of thresholds carried by $F$, by the schedule lemma. Here
$\bar{\mathsf{u}}_F:=\ell_{\circ}^{2(n_F-k_0)_{+}}\mathsf{u}_F$, which equals $\mathsf{u}_F$ for
old terms. At (b)-points the units of $F$ are used exactly as on $\mathcal{A}$.

\emph{Runs on} $Z^{\mathrm{on}}$. On $\mathcal{A}$, the entry lemma places the entry in
$\mathsf{T}_{\mathsf{N}}$, and at (=)-points the links cost at most $\mathsf{G}/24$.

On $Z^{\prime\prime(1)}_{k}$ we apply the frame estimate off the healer's core at the stage label
$\iota_{\nu}\le m$ of the read. At unraised (=)-points, its clause (a) bounds the total by
$\mathsf{T}^{a}_{\mathsf{N}}/960\le\mathsf{T}_e$; the links cost nothing there, and the moves are
paid by the accounting clause of the first tube theorem.

At raised (=)-points, that is, where $\iota_{\nu}(x)=n_F<m(x)$, $m(x)$ being the layer level at
$x$, the point was raised at some
stage $i_1\in[\nu,\mathsf{N}[$. Clause (b) of the frame estimate off the healer's core gives an
entry at most $\mathsf{T}_{\mathsf{N}}$ and a cost of the links bounded by
\begin{equation}\label{9.1:eq-tube-theorem-1}
\frac{\mathsf{u}_F}{1440}\le\frac{\mathrm{rm}_{i_1}}{24}\le\frac{\mathsf{G}}{24}.
\end{equation}
Both inequalities in \eqref{9.1:eq-tube-theorem-1} rest on the lower bound for rooms, which gives
\begin{equation}\label{9.1:eq-tube-theorem-2}
\mathsf{G}\ge\mathrm{rm}_{i_1}\ge\frac{\bar{\mathsf{u}}_{i_1}}{40}\ge\frac{\mathsf{u}_F}{40}.
\end{equation}
Here $\bar{\mathsf{u}}_{i_1}$ is the unit at the stage $i_1$, and
$\bar{\mathsf{u}}_{i_1}\ge\mathsf{T}_{\nu}$ on $Z^{\mathrm{on}}$, because the unit is
unchanged until the first raise and the coefficients in clause (a) of the flat increment lemma
are at most $1$. By \eqref{9.1:eq-tube-theorem-2}, $\mathrm{rm}_{i_1}/24\ge\mathsf{u}_F/960\ge
\mathsf{u}_F/1440$, which is the first inequality in \eqref{9.1:eq-tube-theorem-1}. The second
follows from $\mathsf{G}\ge\mathrm{rm}_{i_1}$.

From this point on, the proof of clauses (ii) and (iii) of the first tube theorem applies, with
the fraction $\tfrac18$ replaced by $\tfrac16$; this yields the bound $\mathsf{T}_e+\tfrac16
(\mathsf{T}_r-\mathsf{T}_e)$ at (=)-points.

On $\mathfrak{O}$, the link majorants $\mathsf{m}_L$ are paid from the margin
$(\tfrac12\theta_S+\tfrac18\theta_{\rho})\beta\,2^{-s}$ of $\mathsf{T}^{\flat}_{\mathsf{N}}$, where
$s$ is here the number of levels between the layer read and the current level,
through the condition on the link constants. This uses
\begin{equation*}
\tfrac{9}{256}+\tfrac1{16}=\tfrac{25}{256}\le\tfrac{32}{256}=\tfrac18.
\end{equation*}
The moves are paid by the stage room increments $\sum_i\mathsf{g}_i$ (clause (b) of the
accounting lemma on $\mathfrak{O}$). The two payments are added, and clause (b$'$) of the
site-descent lemma covers the non-critical intermediate objects. Clause (c) of the accounting
lemma on $\mathfrak{O}$ then concludes.

On the third region $W$, if present, the bound is read off the table of thresholds on the third
region. The (up)-points of $\mathfrak{O}$ are treated by scaling, as in the proof of the
accounting lemma at the T-step.
\end{proof}

\begin{remark}[Entry and increments within the rooms at link sites]\label{9.1:rem-link-accounting}
We use the regions in which a term $(F,X,j)$ reads a point $P=(\mathbf{U},\mathbf{J})$ of the pair
tube $\mathfrak{P}_{\mathsf{c}}$, in the setting of Theorem~\ref{9.1:thm-tube-theorem}: the regular region
$\mathfrak{O}$, the aged region $\mathcal{A}$, the large-field shells $Z''_k$, the healer's core
$\mathfrak{K}$, the top layer $\Omega_k$, and
\begin{equation*}
Z''^{(1)}_k:=Z''_k\setminus\mathfrak{K}.
\end{equation*}
Here $\mathbf{U}$ and $\mathbf{J}$ are the pair variables of $P$ as in Theorem~\ref{9.1:thm-tube-theorem},
and $S=S(\mathbf{U})$ is the configuration at which $F$ is
evaluated in that theorem.
\begin{enumerate}
\item[(i)] The accounting at the link site, that is, the assertion that the entry value plus the increments is
bounded by the rooms, is the conjunction of three statements: the fourfold-tube statement, the
entry-and-link statement, and a further condition listed among the hypotheses of
Theorem~\ref{9.1:thm-tube-theorem}. This conjunction is proved on
$\mathfrak{O}\cup\mathcal{A}$. On $Z''^{(1)}_k$ it is supplemented by the entry statement on the
large-field shells.
\item[(ii)] At link sites clause \textup{(c-i)} of the general site conditions bears only on a component of
$P$ that the estimate used at sites never reads, and clause \textup{(c-ii)} is void.
\item[(iii)] The accounting required on the top layer $\Omega_k$ is
the instance $s=0$ of clause~(c) of the regular-region accounting, $s$ being the index in which that
clause is stated. It is not among the separately named hypotheses of the one-site tube statement used in
the proof of Theorem~\ref{9.1:thm-tube-theorem}.
\end{enumerate}
\end{remark}

\begin{proof}
(i) The accounting at the link site is defined as the conjunction of the three statements listed in
item~(i). On $\mathfrak{O}\cup\mathcal{A}$ the fourfold-tube statement and the entry-and-link statement
hold, and the further condition is a hypothesis
of Theorem~\ref{9.1:thm-tube-theorem}. All three conjuncts are therefore available on
$\mathfrak{O}\cup\mathcal{A}$, and their conjunction, which is the bound of the entry value plus the
increments by the rooms, holds there. On $Z''^{(1)}_k$ the supplementary conjunct is the entry statement
on the large-field shells, used here by statement; it is taken as a consequence of the
fourfold-tube statement in its large-field form and three further statements on large-field regions.

(ii) Clause \textup{(c-i)} is a condition on the component $\mathbf{J}$ of $P$ over the whole tube. The
estimate by which the proof of Theorem~\ref{9.1:thm-tube-theorem} bounds the read at a site never reads
the component of $P$ on which clause \textup{(c-i)} bears, namely $\mathbf{J}$.
Hence clause \textup{(c-i)} constrains nothing that is used.

Clause \textup{(c-ii)} concerns the collar of width $4M_1\xi_j$ at $\partial X\cap\partial Y_1$, where
$\xi_j$ is the length attached to level $j$ in the general site conditions, and $Y_1$, the outermost
region of the chain displayed below, is that of the general site conditions. Along the
chain of enlargements used at link sites, with $K(X)$ the cube hull of $X$ and $X^{+}$ the enlargement
of $X$, one has
\begin{equation*}
Y_1\supseteq K(X)\supseteq X^{+},
\end{equation*}
by clause~(b) of the statement that provides this chain of enlargements. On $X$ the current entering
the read is that of $S$ itself. No stratum finer than $j$
passes through $X$. Both facts are the content of the geometric statement accompanying the general site
conditions together with clause~(b) of the companion statement on the boundary of the enlarged region.
Consequently the situation to which clause
\textup{(c-ii)} refers does not occur at a link site, and the clause is void.

(iii) We have $\Omega_k\subset\mathfrak{O}$. The numerical input is
\begin{equation}\label{9.1:eq-link-accounting-1}
1+\tfrac{9}{16}\,\theta_S\beta\le 1+\beta-\theta_{\rho}\beta .
\end{equation}
Since $\beta>0$, \eqref{9.1:eq-link-accounting-1} is equivalent to
\begin{equation*}
\tfrac{9}{16}\,\theta_S+\theta_{\rho}\le 1 ,
\end{equation*}
a condition on the schedule fractions $\theta_S$ and $\theta_{\rho}$ alone, which holds by the choice of
these fractions.
Only one of the two machines enters the read on $\Omega_k$, by the lemma asserting this for the top
layer. Therefore the accounting required on $\Omega_k$ is clause~(c) of the regular-region accounting
read at $s=0$, and so it is not a separate named hypothesis of the one-site tube statement.
\end{proof}

\begin{definition}[Arrests, the frozen structure and the diminished layer factor]
\label{9.1:def-arrest-factor}
Throughout this definition $\beta$ denotes the layer constant, and for it we put
\begin{equation}\label{9.1:eq-arrest-factor-1}
\varphi_{\mathrm{ARR}}:=1-4.51\,\beta .
\end{equation}

\emph{Arrests.} An \emph{arrest} is a triple $\mathfrak{a}=(W_{\mathfrak{a}},k_a,n^{+}_{a})$ in which
$W_{\mathfrak{a}}$ is a component that leaves the large-field region $Z$ after $n^{+}_{a}$ letters
of its history, $k_a$ is the step of the arrest (the frozen terms below are stored at its end),
and $n^{+}_{a}$ is that number of letters. Here $h_a$ denotes the level from which the
$n^{+}_{a}$ letters of the arrest are counted. To an arrest we attach
\begin{equation}\label{9.1:eq-arrest-factor-a}
j^{+}_{a}:=h_a+n^{+}_{a},\qquad
I(\mathfrak{a}):=\{h_a+1,\,h_a+2,\,\dots,\,j^{+}_{a}+1\}.
\end{equation}

\emph{Frozen terms.} The stage-$i$ chain terms $\mathbb{C}^{(i)}_{k_a}(X)$ with $i\le n^{+}_{a}$ are
\emph{frozen}: they are stored at the end of step $k_a$. Inside step $k_a$ they are not stored
but derived, in the unfolded form given by the unfolding identity for the chain terms.

\emph{Validity of the frozen structure.} Assume the hypotheses of the arrest theorem, namely
\begin{itemize}
\item $\beta=\beta_I=\tfrac15$;
\item $\mathsf{N}=R_k$, where $\mathsf{N}$ is the stage count;
\item the ratio condition $\delta_1M/M_1\ge 2$ on the auxiliary parameter $\delta_1$ of
Ba{\l}aban's construction;
\item $L\ge 8$;
\item the equality $q_1=q_0$ of the two auxiliary parameters $q_1$ and $q_0$ of
Ba{\l}aban's construction;
\item $C_0\ge \tfrac52\,A_0/\varphi_{\mathrm{ARR}}$;
\item the two further conditions of the hypothesis list of the arrest theorem, as stated with
that theorem.
\end{itemize}
Under these hypotheses the arrest theorem states the following, which is used here as stated
there: from the arrest $\mathfrak{a}$ to its dissolution, in the sense of the last clause of the
arrest theorem, every attempt space of the construction uses on $W_{\mathfrak{a}}$ the structure
$(\ell,w)$ frozen at the arrest, namely the layer label $\ell(x)$ of a point $x\in W_{\mathfrak{a}}$
and the weight $w$ by which the units are multiplied on $W_{\mathfrak{a}}$ (this $w$ is not the
complex source of the glossary), with the collar parameter $c\equiv 1$, and the diminished layer
factor
\begin{equation}\label{9.1:eq-arrest-factor-b}
F^{(k')}(x):=f_{k',\ell(x)}-\beta\sum_{\mathfrak{a}'\ \text{live},\ x\in W_{\mathfrak{a}'}}
\mathfrak{h}_{\mathfrak{a}'}(\ell(x)),
\qquad
\mathfrak{h}_{\mathfrak{a}}(\ell):=\sum_{i\in I(\mathfrak{a})}2^{-|\ell-i|},
\end{equation}
where the sum runs over the arrests $\mathfrak{a}'$ with $x\in W_{\mathfrak{a}'}$ that are
\emph{live} at level $k'$, that is, for which $k'$ lies between the arrest and its dissolution,
and $f_{k',\ell}$ are the layer coefficients. The factor satisfies
$F^{(k')}(x)\ge\varphi_{\mathrm{ARR}}$, and the frozen structure and the factor do not depend on
the complex source variables $z_i$. Moreover $\varphi_{\mathrm{ARR}}>0$ at $\beta=\beta_I$, and
the one size relation among the hypotheses of the arrest theorem, $\mathsf{N}\ln 2<2R_k$,
holds at $\mathsf{N}=R_k$.

\emph{Substitution on live pieces.} On live pieces (a live piece is a part of $W_{\mathfrak{a}}$,
for an arrest $\mathfrak{a}$ that is live at the level considered), in the generic table
formulas for terms born in an $\mathbf{R}$-step and in the constructions of the tables that use
them, one substitutes
\begin{equation}\label{9.1:eq-arrest-factor-2}
f_{j,n}\;\longmapsto\;F^{(j)}(x),\qquad \mathsf{u}_n\;\longmapsto\;w\,\mathsf{u}_{\ell}.
\end{equation}
On live pieces, at fixed structure, every difference $F^{(k')}(x)-F^{(k'')}(x)$ equals
$f_{k',\ell(x)}-f_{k'',\ell(x)}$, and likewise every difference of two coefficients at the same
point belonging to two stages is unchanged by the subtraction of the sums of the
$\mathfrak{h}_{\mathfrak{a}}$; neither of the two estimates that use the coefficient itself rather
than a difference of coefficients, recalled in the proof below, is applied to a coefficient
diminished by some $\mathfrak{h}_{\mathfrak{a}}$.
\end{definition}

\begin{proof}
The conclusion of the paragraph on the validity of the frozen structure --- the use of the frozen
structure and of the factor \eqref{9.1:eq-arrest-factor-b} from the arrest to its dissolution,
the lower bound $F^{(k')}(x)\ge\varphi_{\mathrm{ARR}}$ and the independence from the complex
source variables --- is the statement of the arrest theorem, used here as stated there. We check
the remaining assertions in the order in which they are stated.

\emph{Finiteness.} The set $I(\mathfrak{a})$ in \eqref{9.1:eq-arrest-factor-a} consists of
the $n^{+}_{a}+1$ integers from $h_a+1$ to $h_a+n^{+}_{a}+1$, so
$\mathfrak{h}_{\mathfrak{a}}(\ell)$ is a finite sum of positive terms. The sum over live
arrests in \eqref{9.1:eq-arrest-factor-b} is finite because only finitely many arrests occur
on the finite lattice in the $K$ steps. Hence $F^{(k')}(x)$ is well defined.

\emph{The numerical facts.} At $\beta=\tfrac15$,
\begin{equation*}
\varphi_{\mathrm{ARR}}=1-\frac{4.51}{5}=1-0.902=0.098>0 .
\end{equation*}
More generally $\varphi_{\mathrm{ARR}}>0$ if and only if $\beta<100/451$. Since $4\cdot100<451$
we have $100/451<\tfrac14$. Consequently $\varphi_{\mathrm{ARR}}>0$ fails at $\beta=\tfrac14$,
and the lower bound $F^{(k')}\ge\varphi_{\mathrm{ARR}}$ of the arrest theorem is used only at
$\beta=\beta_I$, in particular never at the value $\beta_s=\tfrac14$ of the layer constant used
in other parts of the construction. The size relation at $\mathsf{N}=R_k$ reads
$R_k\ln 2<2R_k$, which holds since $R_k>0$ and $\ln 2<2$.

\emph{Differences of coefficients.} By \eqref{9.1:eq-arrest-factor-b},
$\mathfrak{h}_{\mathfrak{a}}(\ell)$ depends only on $\ell$ and on the set
$I(\mathfrak{a})$, which is fixed at the arrest. It depends on neither the level $k'$ nor the
stage. Let two levels $k'$ and $k''$ be considered at the same point $x$ of a live piece, at
fixed structure, so that $\ell(x)$ and the set of live arrests containing $x$ are the same for
both levels. Then the two subtracted sums in \eqref{9.1:eq-arrest-factor-b} coincide, and
\begin{equation*}
F^{(k')}(x)-F^{(k'')}(x)=f_{k',\ell(x)}-f_{k'',\ell(x)} .
\end{equation*}
For the same reason, if two coefficients attached to two stages at the same point $x$ of a live
piece are both diminished by the same sum
$\beta\sum_{\mathfrak{a}'}\mathfrak{h}_{\mathfrak{a}'}(\ell(x))$, that sum does not depend on
the stage and cancels in their difference.
Thus every difference of coefficients used in the schedule lemma, in the accumulation lemma
for the thresholds $\mathsf{T}$, in the third clause of the accumulation lemma for the regular
region $\mathfrak{O}$, and in the closing estimate of the construction of the tables, is the
same after the substitution \eqref{9.1:eq-arrest-factor-2} as before it.

\emph{The estimates that use the coefficient itself.} The points at which a constituent term is
estimated again at a later level are divided into $(=)$-points and $(\mathrm{up})$-points, in
the sense fixed by the theorem that places each such later estimate in a threshold table, every
point being of exactly one of the two kinds. For a constituent term $F$ (here $F$ denotes the
term, not the factor of \eqref{9.1:eq-arrest-factor-b}) we write $\mathsf{c}_F^{+}$ for its top
coefficient. Only two estimates use the coefficient itself rather than a difference of
coefficients:
\begin{itemize}
\item the lower bound $1-3\beta$ on the layer coefficient on the set $\mathsf{A}(x)$ attached to
$x$ in the construction of the rooms $\mathrm{rm}_{\nu}$, used in that construction;
\item the lower bound $\mathsf{c}_F^{+}\ge\tfrac45$ for the top coefficient at
$(\mathrm{up})$-points, stated with the construction of the tables.
\end{itemize}
Neither is used on a live piece, for the following reasons.

For the first: $\mathsf{A}(x)\cap W_{\mathfrak{a}}=\emptyset$. Indeed, the last clause of the
arrest theorem, together with the lemma that places $W_{\mathfrak{a}}$ in the healer's core
$\mathfrak{K}$, keeps $\mathsf{A}(x)$ off $W_{\mathfrak{a}}$.

For the second, consider a point of a live piece. It is an $(=)$-point, not an
$(\mathrm{up})$-point, for every constituent term of type $\mathbb{B}$ born after the arrest.
A constituent term of type $\mathbb{B}$ born before the arrest carries no summand
$\mathfrak{h}_{\mathfrak{a}}$. So the lower bound
$\mathsf{c}_F^{+}\ge\tfrac45$ at $(\mathrm{up})$-points is never applied to a coefficient
diminished by $\mathfrak{h}_{\mathfrak{a}}$.
\end{proof}

\begin{definition}[Domain and margin function of a left-over term]\label{9.1:def-leftover-margin}
Let $\mathfrak{a}=(W_{\mathfrak{a}},k_a,n^{+}_a)$ be an arrest, with $j^{+}_a$, the frozen
structure $(\ell,w)$ on $W_{\mathfrak{a}}$ and the layer constant $\beta$ as in
Definition~\ref{9.1:def-arrest-factor}. Let $F=\mathbb{C}^{(i)}_{k_a}(X)$ be one of the frozen
stage-$i$ chain terms of $\mathfrak{a}$, $i\le n^{+}_a$, stored at the end of step $k_a$. We call
such an $F$ a \emph{left-over term} of $\mathfrak{a}$ and attach to it the following objects.
\begin{enumerate}
\item[(a)] The \emph{domain} of $F$ is
\begin{equation*}
\mathfrak{D}_F:=\mathfrak{U}[\mathsf{T}^{\flat}_i],
\end{equation*}
the set of pairs below the table $\mathsf{T}^{\flat}_i$, in the sense of the definition of the
tables of thresholds by stage and of the sets of pairs $\mathfrak{U}[\mathsf{T}]$ below a table,
both taken at level $k=k_a$. The level
of $F$ is $j:=k_a$.
\item[(b)] The \emph{margin function} $\varrho_F$ of $F$ is defined by
\begin{equation}\label{9.1:eq-leftover-margin-a}
\varrho_F(x):=
\begin{cases}
\theta_{\rho}\beta\,\max\bigl\{2^{-|\ell(x)-j^{+}_a-1|},\,2^{-(k_a-\ell(x))}\bigr\},
& x\in W_{\mathfrak{a}}\cap Z^{\mathrm{on}},\\[4pt]
\theta_{\rho}\beta\,2^{-(k_a-m(x))},
& x\in \mathfrak{O}_a\cup\bigl(Z^{\mathrm{on}}_{k_a}\setminus W_{\mathfrak{a}}\bigr).
\end{cases}
\end{equation}
Here $\theta_{\rho}$ is the schedule fraction of the margin, $\ell(x)$ is the level of the frozen
structure at $x$, and $m(x)$ is the layer level of the point $x$ off the frozen piece, that is, the
index $m$ of the standard layer coefficient $f_{k_a,m}$ that is read at $x$ there.
\end{enumerate}
If $W_{\mathfrak{a}}$ contains a second arrested component, then in the first line of
\eqref{9.1:eq-leftover-margin-a} the number $j^{+}_a$ is, on each live piece, the one of the
arrest to which that live piece belongs; every live piece carries its own value of $j^{+}$.

On $\mathfrak{O}_a$ and on $Z^{\mathrm{on}}_{k_a}\setminus W_{\mathfrak{a}}$ the second line of
\eqref{9.1:eq-leftover-margin-a} is the generic margin function of a term of level $k_a$ born in an
$\mathbf{R}$-step, in the sense of the definition of the generic margin functions, and $F$ is
treated there as any such term. On $W_{\mathfrak{a}}\cap Z^{\mathrm{on}}$ the maximum is taken so
that $\varrho_F(x)\ge\theta_{\rho}\beta\,2^{-(k_a-\ell(x))}$, the standard depth factor that the
later depth bounds convert. The first member alone would not have this property: with $h_a$ as in
Definition~\ref{9.1:def-arrest-factor}, so that $j^{+}_a=h_a+n^{+}_a$, at $\ell(x)=k_a-1$ and
$n^{+}_a=1$ the first member equals
\begin{equation*}
\theta_{\rho}\beta\,2^{-|k_a-h_a-3|},
\end{equation*}
which, when $k_a-h_a=\mathsf{N}$ with $\mathsf{N}$ the stage count, is
$\theta_{\rho}\beta\,2^{-(\mathsf{N}-3)}$, whereas the depth bounds produce the factor $4^{-s}$ at
depth $s=k_a-\ell(x)=1$. With the maximum the constant $C_{\mathrm{TL}}^{\sharp}$ of the bound in
which $\varrho_F$ is converted does not depend on $\mathsf{N}$.
\end{definition}

\begin{lemma}[Schedule bound for later reads of a left-over term]\label{9.1:lem-leftover-schedule}
Let $\mathfrak{a}=(W_{\mathfrak{a}},k_a,n^{+}_a)$ be an arrest, with $j^{+}_a$ and the diminished
layer factor $F^{(k')}$ of \eqref{9.1:eq-arrest-factor-b} as in
Definition~\ref{9.1:def-arrest-factor}. Let $F=\mathbb{C}^{(i)}_{k_a}(X)$, $i\le n^{+}_a$, be one
of its frozen terms, with the domain $\mathfrak{D}_F$ and the margin function $\varrho_F$ of
\eqref{9.1:eq-leftover-margin-a} (Definition~\ref{9.1:def-leftover-margin}). We write
$\mathsf{c}_F^{+}$ for the top coefficient of the table attached to $F$ and $\mathsf{u}_F$ for
its unit. Assume the following two hypotheses.
\begin{itemize}
\item[(a)] The coefficient comparison for arrests at fixed structure: for every $k'\ge k_a$,
with the frozen structure $(\ell,w)$ held fixed on $W_{\mathfrak{a}}$, at every point $x$ of a
live piece of $\mathfrak{a}$, with $\ell=\ell(x)$,
\begin{equation}\label{9.1:eq-leftover-schedule-1}
\varphi^{(n^{+}_a),(k_a)}_{\ell}-F^{(k')}(x)
\ge\beta\,2^{-|\ell-j^{+}_a-1|}+\beta\bigl(2^{-(k_a-\ell)}-2^{-(k'-\ell)}\bigr),
\end{equation}
and the same bound holds for the enclosing attempt spaces; here
$\varphi^{(n^{+}_a),(k_a)}_{\ell}=\varphi^{(n^{+}_a)}_{\ell}$ is the layer coefficient at
level $\ell$, the superscript $(k_a)$ recording only the step at which it is read.
\item[(b)] The core statement: for every constituent term whose enlarged domain $X^{+}$ meets
the healer's core $\mathfrak{K}$, at the real centre of the core the slice lies in
$\mathfrak{D}_F[0]$, the read is placed in the table $\tfrac12\mathsf{T}_N$, and it is compared
with the printed table of $F$, whose unit is $\bar{\mathsf{u}}_F$, through
$\tfrac12\mathsf{T}_N\le\tfrac12\mathsf{T}_i$; and the inequality
$\theta_{\rho}\beta\le\tfrac12\varphi^{(i)}$ holds.
\end{itemize}
Let $\mathsf{T}$ be the table in which the tube theorem
(Theorem~\ref{9.1:thm-tube-theorem}) places a later read of $F$: a read at a $T$-step $T_{k'}$
with $k'\ge k_a$, or at an $\mathbf{R}$-step $\mathbf{R}_{k'}$ with $k'>k_a$, at any stage,
including the site \textup{(L)}. Then at the points which the tube theorem labels $(=)$,
\begin{equation}\label{9.1:eq-leftover-schedule-a}
\mathsf{T}\le\bigl(\mathsf{c}_F^{+}-\varrho_F\bigr)\mathsf{u}_F ,
\end{equation}
and at all other points the read is placed as at the points which the tube theorem labels
$(\mathrm{up})$.
\end{lemma}

\begin{proof}
\emph{The tables.} The tube theorem (Theorem~\ref{9.1:thm-tube-theorem}) places the reads in
the following tables. At a step
$k'>k_a$, by part~(a) of the accumulation bound on the regular region $\mathfrak{O}$ and by the
table of reads on the third region, every table satisfies, at a point $x$ of level $\ell$,
\begin{equation}\label{9.1:eq-leftover-schedule-2}
\mathsf{T}\le\Bigl(F^{(k')}(x)+\tfrac{9}{16}\,\theta_S\beta\,2^{-(k'-\ell)}\Bigr)\mathsf{u}_F .
\end{equation}
At the $T$-step $T_{k_a}$ the reads are placed, by the accumulation bound at the $T$-step, in its
reference table
\begin{equation}\label{9.1:eq-leftover-schedule-3}
\mathsf{T}_r=\Bigl(F^{(k_a)}(x)-\tfrac14(1-\theta_S)\,\beta\,2^{-(k_a-\ell)}\Bigr)\mathsf{u}_F .
\end{equation}
Inside step $k_a$ the frozen terms are derived and unfolded, as in
Definition~\ref{9.1:def-arrest-factor}. The diminished coefficient $F^{(k_a)}$ entering
\eqref{9.1:eq-leftover-schedule-3} is therefore the one of \eqref{9.1:eq-arrest-factor-b},
computed within the single run of step $k_a$. On $W_{\mathfrak{a}}$ every space from the arrest
to its dissolution uses the frozen structure $(\ell,w)$ of
Definition~\ref{9.1:def-arrest-factor}, so the tables \eqref{9.1:eq-leftover-schedule-2} and
\eqref{9.1:eq-leftover-schedule-3} are read there in the unit $\mathsf{u}_F$ of $F$. In both
cases the bound
\eqref{9.1:eq-leftover-schedule-a} says that the difference between $\mathsf{c}_F^{+}$ and the
coefficient of the table, which we call the room above the table, is at least $\varrho_F$. We
use throughout the following inequalities between the schedule fractions $\theta_S$ and
$\theta_{\rho}$, which are conditions on the choice of these fractions:
\begin{equation*}
\theta_S\le1,\qquad \theta_{\rho}\le\tfrac18\ \ (\text{so }\theta_{\rho}\le1),\qquad
\theta_{\rho}\le\tfrac14(1-\theta_S),\qquad \tfrac12\theta_S\ge2\theta_{\rho}.
\end{equation*}

\emph{On the live piece} $W_{\mathfrak{a}}\cap Z^{\mathrm{on}}$. Here
\eqref{9.1:eq-leftover-margin-a} gives
\begin{equation*}
\varrho_F(x)=\theta_{\rho}\beta\max\bigl\{2^{-|\ell-j^{+}_a-1|},\,2^{-(k_a-\ell)}\bigr\},
\qquad \ell=\ell(x).
\end{equation*}
By the descent property of the top coefficients, $\mathsf{c}_F^{+}\ge\varphi^{(n^{+}_a)}_{\ell}$.
Combined with hypothesis~(a), that is with \eqref{9.1:eq-leftover-schedule-1}, this gives
\begin{equation}\label{9.1:eq-leftover-schedule-4}
\mathsf{c}_F^{+}-F^{(k')}(x)
\ge\beta\,2^{-|\ell-j^{+}_a-1|}+\beta\bigl(2^{-(k_a-\ell)}-2^{-(k'-\ell)}\bigr),
\qquad k'\ge k_a .
\end{equation}

\emph{Case $k'>k_a$.} By \eqref{9.1:eq-leftover-schedule-2} and
\eqref{9.1:eq-leftover-schedule-4}, the room above the table is at least
\begin{equation*}
\beta\,2^{-|\ell-j^{+}_a-1|}+\beta\,2^{-(k_a-\ell)}
-\beta\,2^{-(k'-\ell)}-\tfrac{9}{16}\,\theta_S\beta\,2^{-(k'-\ell)} .
\end{equation*}
Since $2^{-(k'-\ell)}=2^{-(k_a-\ell)}\,2^{-(k'-k_a)}$, this expression equals
\begin{equation*}
\beta\,2^{-|\ell-j^{+}_a-1|}
+\beta\,2^{-(k_a-\ell)}\Bigl[1-\bigl(1+\tfrac{9}{16}\theta_S\bigr)2^{-(k'-k_a)}\Bigr].
\end{equation*}
As $k'-k_a\ge1$ and $\theta_S\le1$, the bracket satisfies
\begin{equation*}
1-\bigl(1+\tfrac{9}{16}\theta_S\bigr)2^{-(k'-k_a)}
\ge 1-\tfrac12\cdot\tfrac{25}{16}=1-\tfrac{25}{32}=\tfrac{7}{32}.
\end{equation*}
Hence the room is at least
\begin{equation}\label{9.1:eq-leftover-schedule-5}
\beta\,2^{-|\ell-j^{+}_a-1|}+\tfrac{7}{32}\,\beta\,2^{-(k_a-\ell)} .
\end{equation}
Each of the two summands of \eqref{9.1:eq-leftover-schedule-5} dominates the corresponding
member of the maximum defining $\varrho_F$. The first does so because $\theta_{\rho}\le1$. The
second does so because $\theta_{\rho}\le\tfrac18<\tfrac{7}{32}$. Both summands are nonnegative,
so their sum is at least
\begin{equation*}
\theta_{\rho}\beta\max\bigl\{2^{-|\ell-j^{+}_a-1|},\,2^{-(k_a-\ell)}\bigr\}=\varrho_F(x).
\end{equation*}
This is \eqref{9.1:eq-leftover-schedule-a}.

\emph{Case $k'=k_a$, the step $T_{k_a}$.} By \eqref{9.1:eq-leftover-schedule-3} and by
\eqref{9.1:eq-leftover-schedule-4} at $k'=k_a$, where the second summand vanishes, the room above
the table is at least
\begin{equation*}
\beta\,2^{-|\ell-j^{+}_a-1|}+\tfrac14(1-\theta_S)\,\beta\,2^{-(k_a-\ell)} .
\end{equation*}
Since $\theta_{\rho}\le1$ and $\theta_{\rho}\le\tfrac14(1-\theta_S)$, each summand again
dominates its member of the maximum, and the sum is at least $\varrho_F(x)$.

\emph{Off the live piece}, that is on $\mathfrak{O}_a$ and on
$Z^{\mathrm{on}}_{k_a}\setminus W_{\mathfrak{a}}$. Let $m=m(x)$ be the level at $x$. There the
tables \eqref{9.1:eq-leftover-schedule-2} and \eqref{9.1:eq-leftover-schedule-3} are likewise
read in the unit $\mathsf{u}_F$ of $F$, and the comparison is between coefficients. There
\eqref{9.1:eq-leftover-margin-a} gives $\varrho_F(x)=\theta_{\rho}\beta\,2^{-(k_a-m)}$, and the top
coefficient satisfies
\begin{equation}\label{9.1:eq-leftover-schedule-6}
\mathsf{c}_F^{+}\ge f_{k_a,m}+\tfrac12\,\theta_S\beta\,2^{-(k_a-m)} .
\end{equation}
Off the live piece no replacement of the coefficients is made (the replacement
$f_{j,n}\mapsto F^{(j)}(x)$ of Definition~\ref{9.1:def-arrest-factor} is made on live pieces
only), so the coefficient read in \eqref{9.1:eq-leftover-schedule-2} and
\eqref{9.1:eq-leftover-schedule-3} in place of $F^{(k')}(x)$ is $f_{k',m}$.

At $T_{k_a}$, since $\theta_S\le1$, the coefficient of the table \eqref{9.1:eq-leftover-schedule-3}
is at most $f_{k_a,m}$. By \eqref{9.1:eq-leftover-schedule-6} the room is therefore at least the
margin $\tfrac12\theta_S\beta\,2^{-(k_a-m)}$. Because $\tfrac12\theta_S\ge2\theta_{\rho}$, this
margin is at least $2\theta_{\rho}\beta\,2^{-(k_a-m)}\ge\varrho_F(x)$.

For $k'>k_a$, by the definition of the layer coefficients $f_{k,m}$,
\begin{equation*}
f_{k_a,m}-f_{k',m}\ge\tfrac12\,\beta\,2^{-(k_a-m)}\ge\beta\,2^{-(k'-m)} .
\end{equation*}
The second inequality holds because $2^{k'-k_a}\ge2$. By \eqref{9.1:eq-leftover-schedule-2} and
\eqref{9.1:eq-leftover-schedule-6}, the room above the table is at least
\begin{equation*}
\bigl(f_{k_a,m}-f_{k',m}\bigr)-\tfrac{9}{16}\,\theta_S\beta\,2^{-(k'-m)}
+\tfrac12\,\theta_S\beta\,2^{-(k_a-m)} .
\end{equation*}
By the preceding display the first term is at least $\beta\,2^{-(k'-m)}$, which, since
$\theta_S\le1$, is at least $\tfrac{9}{16}\theta_S\beta\,2^{-(k'-m)}$. The room is therefore at
least $\tfrac12\theta_S\beta\,2^{-(k_a-m)}$. As at $T_{k_a}$, this is at least
$2\theta_{\rho}\beta\,2^{-(k_a-m)}\ge\varrho_F(x)$.

\emph{Dissolution.} Let $k$ be the step at which the dissolution is read. At the dissolution
the region $\Lambda$ of the arrest property of Definition~\ref{9.1:def-arrest-factor} contains
$W_{\mathfrak{a}}$.
Moreover $W_{\mathfrak{a}}$ lies in the healer's core $\mathfrak{K}$, for the following reason.
The arrest creates a new large-field region inside $W_{\mathfrak{a}}$ at step $k_a$.
\cite[(ii)]{B-CMP122a} states that ``in the preceding $N$ renormalization
steps no new large field regions were created inside this component''; here $N$ is the stage
count $\mathsf{N}=R_k$ of Definition~\ref{9.1:def-arrest-factor}. Put $h:=k-\mathsf{N}$, a level
(in this paragraph $h$ does not denote a history). Hence $k_a\le h$, and
\begin{equation*}
W_{\mathfrak{a}}\subset Z_{k_a}\subset Z_h\subset\mathfrak{K}.
\end{equation*}
This is the clause of the arrest property which, through the lemma on the location of arrested
components, puts $W_{\mathfrak{a}}$ in the healer's core, read at levels $k\ge k_a+\mathsf{N}$.

Every frozen term $\mathbb{C}^{(i)}_{k_a}(X)$ has $X\cap W_{\mathfrak{a}}\neq\emptyset$
\cite[(1.63)]{B-CMP122a}. Since $X\subset X^{+}$, the enlarged domain $X^{+}$ therefore meets
$\mathfrak{K}$. The reads at the dissolution are thus governed by the core statement of
hypothesis~(b). They are not governed by the estimate on the frames of $\mathfrak{F}_G$, which
concerns the region $Z''_k\setminus\mathfrak{K}$.

By hypothesis~(b), at the real centre of the core the slice lies in
$\mathfrak{D}_F[0]$, the table of the read is $\mathsf{T}=\tfrac12\mathsf{T}_N$, and its
comparison with the printed table of $F$ is
\begin{equation*}
\tfrac12\,\mathsf{T}_N\le\tfrac12\,\mathsf{T}_i
\le\bigl(\mathsf{c}_F^{+}-\varrho_F\bigr)\bar{\mathsf{u}}_F .
\end{equation*}
The last inequality is implied by $\theta_{\rho}\beta\le\tfrac12\varphi^{(i)}$. Here
$\bar{\mathsf{u}}_F$ is the unit of the printed table of $F$.

The second-class terms $V(X^{\sharp})$ carry the standard letters, by the remark on
second-class terms of the arrest construction: they are born in an $\mathbf{R}$-step.
\end{proof}

\begin{definition}[The inherited domain and the stage-increment sums]\label{9.1:not-inherited-domain}
Fix a level $k$, a stage $i\le\mathsf{N}$ and a domain $X$ of a constituent term. The objects below
are evaluated at
points $y$ of the shell $m=m(y)$ in the units $r_a(y)=\mathsf{u}^a_m(y)$ attached to the thresholds
indexed by $a$. We put $s:=k-m$.

\emph{Ambient convention.} A pair $(\mathbf{U},\mathbf{J})$ belonging to a set of pairs attached to $X$
is given on the enlarged domain
\begin{equation}\label{9.1:eq-inherited-domain-1}
X^{+}:=X\cup\bigcup_{x\in X}B(x),
\end{equation}
where $B(x)$ is a ball about $x$ whose radius is less than the side of one $M$-cube of the local
scale. The pair obeys the conditions imposed on it at every point of $X^{+}$. Every set of pairs below
is understood with this convention, which we call the ambient convention.

\emph{The three factors.} The conditions below bear on the part of $X$ in the region
$Z^{\mathrm{on}}$ of processed onions and on the part of $X$ in the tower region $\mathfrak{T}$; the
top layer $\Omega_k$ is the shell $m=k$. The shells $m<k$ of the tower region $\mathfrak{T}$ are
shells of untouched onions. For each stage $l$ let $\iota_l(y)$ and $\iota^1_l(y)$
denote the stage-$l$ increments of the thresholds at $y$, and let $\mathbf{r}^T_a$ denote the radius
of a constituent term $T$ at the index $a$. We write $E_a$ for the entry of a pair that is compared
with the thresholds indexed by $a$.
\begin{itemize}
\item[(1)] $\mathfrak{L}_i(X\cap Z^{\mathrm{on}})$ is the set of pairs that satisfy, on
$X\cap Z^{\mathrm{on}}$, the lettered conditions of Ba{\l}aban's construction
\cite[(1.64)--(1.69)]{B-CMP122a}, taken with the stage-$i$ coefficients.
\item[(2)] $\mathfrak{I}_i(X\cap\mathfrak{T})$ is the set of pairs that satisfy, on
$X\cap\mathfrak{T}$, the clause proper to each constituent term $T$ of the stage-$i$ chain term
$\mathbb{C}^{(i)}_k(X)$, $T$ being a term of class $\mathbb{B}$, $\mathbb{E}$ or
$\mathbb{R}$. The threshold $\theta_{T,a}$ of
that clause is diminished to $\theta^{(i)}_{T,a}$ as follows. For $T$ of class $\mathbb{B}$ it is
diminished by $\sum_{l\le i}\iota_l\,\mathbf{r}^T_a$. For $T$ of class $\mathbb{E}$ or $\mathbb{R}$
it is diminished by
\begin{equation}\label{9.1:eq-inherited-domain-2}
\tfrac13\bigl(\theta_{T,a}-\theta_{\mathrm{src},a}\bigr)+\sum_{l\le i}\iota^1_l\,\mathbf{r}^T_a .
\end{equation}
Here $\theta_{\mathrm{src},a}$ is the source threshold at the index $a$, the threshold over which the
room of $T$ in (c) below is measured. For a term of class $\mathbb{E}$ or $\mathbb{R}$ and
the composite entries $a\in\{2,4\}$, the following is required at the nested triples, with
subdomain $X'\subset X$, over which the chain sums are taken; here $p$ denotes the level at which
the composite entry is formed and $j$ the level of the term (for a left-over term $j=k_a$, as in
Definition~\ref{9.1:def-leftover-margin}). If $p\le j$, the
composite entries of a pair are taken to be its own entries of the types $E_1$ and $E_3$, and these
are required to lie below the same diminished thresholds. If $p>j$, the condition imposed is that the
points lie in the stage spaces.
\item[(3)] $\mathfrak{N}^{(i)}(X)$ is the set of pairs that satisfy Ba{\l}aban's direct conditions,
understood with the ambient convention. On $\mathfrak{T}$ the conditions are taken with the factor
$1+2\beta-\sum_{l\le i}\iota^1_l$, and on $Z^{\mathrm{on}}$ with the factor $5w_{*}$, where $w_*$ is
a constant with $w_*\ge1$ that is not smaller than the coefficients of the lettered conditions.
\end{itemize}
The \emph{inherited domain} at stage $i$ is
\begin{equation}\label{9.1:eq-inherited-domain-b}
\mathfrak{D}^{\flat}_i(X):=\mathfrak{N}^{(i)}(X)\cap\mathfrak{L}_i\bigl(X\cap Z^{\mathrm{on}}\bigr)
\cap\mathfrak{I}_i\bigl(X\cap\mathfrak{T}\bigr).
\end{equation}

\emph{Constants.} Let $\beta$ be the layer constant, and put $\mathsf{w}_0:=1/24$. The schedule fraction
$\theta_S$ is an auxiliary parameter chosen once, in the window $\tfrac12\le\theta_S\le\tfrac89$, at
the place in the order of choice of the constants where the schedule fractions are fixed. This
window is a condition on the choice of $\theta_S$; it
adds no member to $\gamma_U$. The entry coefficient is $\mathsf{c}_e:=\theta_S$.

\emph{Properties used.} We use the following statements of this book, recorded here in the present
notation: the theorem on the domain and bound of the two-slot chain terms at stage $i$, whose bound
gives (a) below and whose estimate of the increment sums gives (b) below; and the corollary of that
theorem together with the accompanying statement on the room over the source at real slots, which
give (c) below.
\begin{itemize}
\item[(a)] Consider the two-slot stage-$i$ chain term $\mathbb{C}^{(i)}_k(X,\mathcal{S})$. It is a
stored term whose datum in the first slot lies in $\mathfrak{D}^{\flat}_i(X)$; its datum in the
second slot lies
in $\mathfrak{D}_{\mathcal{S}}(X)$, the domain attached to the second slot $\mathcal{S}$. On
$\mathfrak{D}^{\flat}_i(X)\times\mathfrak{D}_{\mathcal{S}}(X)$ it is holomorphic and jointly
$\mathcal{G}$-invariant. It satisfies
\begin{equation}\label{9.1:eq-inherited-domain-3}
\bigl\|\mathbb{C}^{(i)}_k\bigr\|^{A,\flat}_{(i)}\le C_0^{\sharp},
\end{equation}
where $C_0^{\sharp}$ is a constant of that theorem and $\|\cdot\|^{A,\flat}_{(i)}$ is the weighted
supremum norm of the two-slot chain terms at stage $i$, whose weights are the decay rates of those
terms. In particular
$\mathbb{C}^{(i)}_k$ is bounded on $\mathfrak{D}^{\flat}_i(X)\times\mathfrak{D}_{\mathcal{S}}(X)$ by
$C_0^{\sharp}e^{-a_md_m(X)}$, where $a_md_m(X)$ is the decay exponent attached to $X$ in that
theorem.
\item[(b)] Let $n_W(y)$ be the decay exponent that the bound on the increment sums carries at $y$,
indexed by the shells of untouched onions; on such a shell $m=k-s<k$ it equals $s-1$. At every point
$y$ of such a shell, and at
every point $y\in\Omega_k$, the stage-increment sums obey
\begin{equation}\label{9.1:eq-inherited-domain-a}
\begin{aligned}
\sum_{l\le i}\iota_l(y)&\le\tfrac12\mathsf{w}_0\beta e^{-n_W(y)}=\tfrac1{48}\beta e^{-(s-1)}
&&\text{(shell $m=k-s<k$, untouched)},\\
\sum_{l\le i}\iota_l(y)&\le\tfrac13\mathsf{w}_0\beta=\tfrac{\beta}{72}
&&\text{(on $\Omega_k$)},\\
\sum_{l\le i}\iota^1_l(y)&\le2\mathsf{w}_0\beta=\tfrac{\beta}{12}. &&
\end{aligned}
\end{equation}
These inequalities hold for every admissible choice of the constants; they are not conditions of the
order of choice and involve no member of $\gamma_U$.
\item[(c)] Take the slots real and $\mathsf{c}_e=\theta_S$. The room $r_T$ of a constituent term $T$
over its source, measured in the units $r_a(y)$, is at least $\tfrac12\beta\theta_S2^{-(k-m)}$ at
points $y$ of $\mathfrak{T}$ in the shell $m$, and at least $\tfrac12\beta$ on $\Omega_k$.
\end{itemize}
The domain of a left-over term in Definition~\ref{9.1:def-leftover-margin} is the set
$\mathfrak{U}[\mathsf{T}^{\flat}_i](X)$ of pairs below the table $\mathsf{T}^{\flat}_i$. Its inclusion,
understood with the ambient convention, in $\mathfrak{D}^{\flat}_i(X)$ is the subject of the
lemma asserting $\mathfrak{U}[\mathsf{T}^{\flat}_i](X)\subset\mathfrak{D}^{\flat}_i(X)$, under the
hypotheses stated there; on $\mathfrak{D}^{\flat}_i(X)$ the bound is the
one given by (a).
\end{definition}

\begin{lemma}[Inclusion of the table domain in the inherited domain]\label{9.1:lem-domain-inclusion}
Let $\theta_S\in[\tfrac12,\tfrac89]$, $\mathsf{w}_0=\tfrac1{24}$ and $\mathsf{c}_e=\theta_S$. Consider the
site of the operation $\mathbf{R}_k$ of step $k$, where $k=k_a$ when the term read is a left-over
term of an arrest $\mathfrak{a}=(W_{\mathfrak{a}},k_a,n^{+}_a)$ in the sense of
Definition~\ref{9.1:def-arrest-factor}. Let $i\le\mathsf{N}$ be a stage and let $X$ be a leaf
domain that does not meet the raw core; when $k=k_a$, assume moreover that
$X\cap W_{\mathfrak{a}}=\emptyset$, where $W_{\mathfrak{a}}$ is the component excluded in mid-run at
that step. For $y\in X\cap\mathfrak{T}$ on the shell $m=m(y)$ put $s:=k-m$, and write every threshold
at $y$ as its coefficient in the unit $r_a(y)=\mathsf{u}^a_m(y)$. Assume the following.
\begin{enumerate}
\item[(i)] The stage-increment sums of Notation~\ref{9.1:not-inherited-domain},
\eqref{9.1:eq-inherited-domain-a}: on an untouched shell $m=k-s<k$, that is, a shell of an onion
lying outside the processed onions $Z^{\mathrm{on}}$,
\begin{equation*}
\sum_{l\le i}\iota_l(y)\le\tfrac12\mathsf{w}_0\beta e^{-(s-1)}=\tfrac1{48}\beta e^{-(s-1)};
\end{equation*}
on $\Omega_k$ this sum is at most $\tfrac13\mathsf{w}_0\beta=\tfrac{\beta}{72}$; everywhere
$\sum_{l\le i}\iota^1_l(y)\le2\mathsf{w}_0\beta=\tfrac{\beta}{12}$; and on $\Omega_k$ the
$\iota^1$-increments subtracted from the clause of a newborn $\mathbb{E}^{(k)}$- or
$\mathbb{R}^{(k)}$-term are at most $\sum_l\tfrac23\mathsf{w}_0\rho^{(l)}_k$, where the
$\rho^{(l)}_k$, one for each stage $l$, are the top-layer quantities of
Notation~\ref{9.1:not-inherited-domain} entering these increments, with $\sum_l\rho^{(l)}_k<2\beta$.
\item[(ii)] The bound for scale drops of Notation~\ref{9.1:not-inherited-domain}: a constituent $T$
whose own top shell $\ell_T$ lies below $m$, $\ell_T<m$ (a scale drop), has
$\theta^{(i)}_{T,a}\ge\tfrac32\,r_a$.
\item[(iii)] The room table of the constituents on the untouched shells and on the top layer
$\Omega_k$: for
$\ell_T\ge m$ one has $\mathbf{r}^T_a=r_a$; an older leaf $\mathbf{B}^{(j')}$ or $\mathbf{B}'^{(j')}$,
$m\le j'\le k-1$, has room
\begin{equation*}
r_T=f_{j',m}-f_{k,m}=\beta2^{-s}(2^{k-j'}-1)
\end{equation*}
over the level-$k$ coefficient; an $\mathbb{E}^{(m)}$- or $\mathbb{R}^{(m)}$-term with $\ell_T=m<k$ has
$\theta_{T,a}=r_a$ (coefficient $1$ on its own top shell), $\theta_{\mathrm{src},a}=f_{k,m}r_a$ and
$r_T=\beta(1-2^{-s})$; and an $\mathbb{E}^{(k)}$- or
$\mathbb{R}^{(k)}$-term at $m=k$ has $r_T=\beta$.
\item[(iv)] The layer-wise schedule with $\mathsf{c}_e=\theta_S$: the newborn leaf $\mathbf{B}^{(k)}$
born at the step $T_k$ has, on a shell $m<k$, the threshold
$\theta_{T,a}=(f_{k,m}+\theta_S\beta2^{-s})r_a$, and on the top layer $\Omega_k$ the factor $1+\beta$,
$\theta_{T,a}=(1+\beta)r_a$. A newborn $\mathbb{E}^{(k)}$- or $\mathbb{R}^{(k)}$-term carries the same
top-layer factor: on $\Omega_k$ it has $\theta_{T,a}=(1+\beta)r_a$ and
$\theta_{\mathrm{src},a}=f_{k,k}r_a=r_a$, so that $\theta_{T,a}-\theta_{\mathrm{src},a}=\beta r_a$, the
room $r_T=\beta$ of (iii). The top-layer factor
is that of \cite[(3.16)]{B-CMP109}; in
\cite{B-CMP119} the new terms of the type $\mathbb{E}$ are analytic at $(1+\beta)\alpha$.
\item[(v)] The diminished layer factor on live pieces: for a live piece of a left-over term of an
arrest $\mathfrak{a}'=(W_{\mathfrak{a}'},k_{a'},n^{+}_{a'})$, read at step $k'>k_{a'}$, at a point of
shell $\ell$, in the frozen structure of $\mathfrak{a}'$,
\begin{equation*}
\varphi^{(n^{+}_{a'})}_{\ell}-F^{(k')}\ge\beta2^{-(k_{a'}-\ell)}-\beta2^{-(k'-\ell)},
\end{equation*}
with $F^{(k')}$ as in \eqref{9.1:eq-arrest-factor-b}.
\end{enumerate}
The set $\mathfrak{U}[\mathsf{T}^{\flat}_i](X)$ is read with the same conventions as
$\mathfrak{D}^{\flat}_i(X)$ in Notation~\ref{9.1:not-inherited-domain}: a pair belongs to it if it is
given on $X^{+}=X\cup\bigcup_{x\in X}B(x)$, where $B(x)$ has radius less than one $M$-cube of the local
scale, and obeys the table there; the entry $E_5$ and the cube clause are those of the source term, as
in the definition of $\mathfrak{U}[\mathsf{T}]$, and on $\mathfrak{T}$ they do not depend on the stage;
the composite entries formed at summation triples of level $p\le j$, $j$ as in
Definition~\ref{9.1:def-leftover-margin}, are the
pair's own entries of the types $E_1$ and $E_3$, and those formed at level $p>j$ are read under the
rider that the points lie in the stage spaces. Then
\begin{equation}\label{9.1:eq-domain-inclusion-1}
\mathfrak{U}[\mathsf{T}^{\flat}_i](X)\subset\mathfrak{D}^{\flat}_i(X),
\end{equation}
and every slice that the tube theorem (Theorem~\ref{9.1:thm-tube-theorem}), the room-accumulation
statement on the regular region, the entry statement, the reading of the third region, and the schedule
bound for later reads of a left-over term (Lemma~\ref{9.1:lem-leftover-schedule}) place in the tube
$\mathfrak{T}_F[\varsigma](X)$ of $F=\mathbb{C}^{(i)}_{k_a}(X)$ is such a pair. Precisely, at every
$y\in X\cap\mathfrak{T}$ and for every index $a\notin\{5,8\}$,
\begin{equation}\label{9.1:eq-domain-inclusion-2}
\mathsf{T}^{\flat}_i(y)\le\mathsf{T}^{\flat}_0(y)\le f_{k,m}+\tfrac{9}{16}\theta_S\beta2^{-s}
<\theta^{(i)}_{T,a}(y)
\end{equation}
for every constituent $T$ of $\mathfrak{I}_i$ with $(y,a)\in\operatorname{Cl}(T)$. The last inequality
may hold with equality, and then only for $\theta_S=\tfrac89$, $s=1$ and $T$ an $\mathbb{E}$- or
$\mathbb{R}$-term; this suffices, because $\mathfrak{U}[\mathsf{T}]$ and $\mathfrak{I}_i$ are both
defined by strict inequalities $E_a<\text{threshold}$, so that $\mathsf{T}\le\theta$ gives
$\{E_a<\mathsf{T}\}\subset\{E_a<\theta\}$. The margins in the last inequality are at least
$\tfrac14\theta_S\beta2^{-s}$ for the newborn leaf $\mathbf{B}^{(k)}$ on a shell $m<k$;
$\tfrac{11}{24}\beta2^{-s}$ for older leaves and for live left-overs of arrests $\mathfrak{a}'$ with
$k_{a'}<k$;
\begin{equation*}
\Bigl[\tfrac23(1-2^{-s})-\tfrac1{12}-\tfrac{9}{16}\theta_S2^{-s}\Bigr]\beta\ge0
\end{equation*}
for $\mathbb{E}^{(m)}$- and $\mathbb{R}^{(m)}$-terms with $m<k$; and $\tfrac1{18}\beta$ on
$\Omega_k\cap\mathfrak{T}$, that is, for the newborn leaf $\mathbf{B}^{(k)}$ and the newborn
$\mathbb{E}^{(k)}$- and $\mathbb{R}^{(k)}$-terms there. The direct clauses of
$\mathfrak{N}^{(i)}$ exceed the table by at least $\tfrac5{12}\beta$. No condition of the family
$\gamma_U$ is used: each of these inequalities is absolute, valid for every
$\theta_S\in[\tfrac29,\tfrac89]$.
\end{lemma}

\begin{proof}
\emph{Scope.} The case $k=k_a$ with $X\cap W_{\mathfrak{a}}\neq\emptyset$ is not covered by the
present lemma: the bound for live left-overs used below applies only at
$k>k_a$; at $k=k_a$ and stages $i>n^{+}_a$ on $W_{\mathfrak{a}}$ itself, the diminished layer factor of
Definition~\ref{9.1:def-arrest-factor} leaves against the sum of the stage increments only the notch
$\beta2^{-|\ell-j^{+}_a-1|}$, which the inequality \eqref{9.1:eq-domain-inclusion-a} below does not
control when the notch is deep. The inclusion on $W_{\mathfrak{a}}$ at its own arrest step is the
content of the increment statement.

Throughout, $y\in X\cap\mathfrak{T}$ lies on the shell $m=m(y)$, $s=k-m$, and thresholds at $y$ are
coefficients in the unit $r_a(y)$. By \eqref{9.1:eq-inherited-domain-b}, $\mathfrak{D}^{\flat}_i(X)$ is
the intersection of $\mathfrak{N}^{(i)}(X)$, the set defined by the direct conditions at the factor
$1+2\beta-\sum_l\iota^1_l$ on $\mathfrak{T}$ and at the factor $5w_{*}$ on $Z^{\mathrm{on}}$, where
$w_{*}\ge1$ is the constant of the direct conditions on the processed onions in
Notation~\ref{9.1:not-inherited-domain}; of
$\mathfrak{L}_i(X\cap Z^{\mathrm{on}})$, the set defined by the clauses
\cite[(1.64)--(1.69)]{B-CMP122a} on the
processed onions; and of $\mathfrak{I}_i(X\cap\mathfrak{T})$, the set defined by the clause of each
constituent $T$ diminished in the stage. For a leaf the diminished threshold is
\begin{equation*}
\theta^{(i)}_{T,a}=\theta_{T,a}-\sum_{l\le i}\iota_l\,\mathbf{r}^T_a,
\end{equation*}
and for an $\mathbb{E}$- or $\mathbb{R}$-term it is
\begin{equation}\label{9.1:eq-domain-inclusion-3}
\theta^{(i)}_{T,a}=\theta_{T,a}-\tfrac13(\theta_{T,a}-\theta_{\mathrm{src},a})
-\sum_{l\le i}\iota^1_l\,\mathbf{r}^T_a,
\end{equation}
the composite entries $a\in\{2,4\}$ being subject to the composite clause. We treat the three sets in
turn, and then the extension to $X^{+}$ and the composite entries.

\emph{(a) The processed onions $Z^{\mathrm{on}}$.} There the table is, by the definition of the flat
table on the processed onions, $\mathsf{T}^{\flat}_\nu:=\mathsf{T}_\nu$, the table of
\cite[(1.64)--(1.69)]{B-CMP122a} at
stage $\nu$ multiplied by the unit $\bar{\mathsf{u}}$, and it is decreasing in $\nu$
\cite{B-CMP122a}. Hence the restriction of $\mathfrak{U}[\mathsf{T}^{\flat}_i](X)$ to
$X\cap Z^{\mathrm{on}}$ is exactly the set defined by the clauses of $\mathfrak{L}_i$. The direct
entries of such a pair are below $5w_{*}$ times the unit, because $w_{*}\ge1$ and the coefficients of
those clauses are at most $1$. So the clauses of $\mathfrak{N}^{(i)}$ on $Z^{\mathrm{on}}$ hold.

\emph{(b) The region $\mathfrak{T}$: the table.} By the definition of the flat table on
$\mathfrak{T}$, the stage-$i$ table lies below the stage-$0$ table,
and
\begin{equation*}
\mathsf{T}^{\flat}_0=\bigl(f_{k,m}+\tfrac12\theta_S\beta2^{-s}\bigr)\mathsf{u}_m
+\sum_{0\le l<\mathsf{N}}\mathsf{g}_l .
\end{equation*}
By part (a) of the room-accumulation statement on the regular region, the room increments sum to at
most $\tfrac1{16}\theta_S\beta2^{-s}\mathsf{u}_m$, so that
\begin{equation}\label{9.1:eq-domain-inclusion-4}
\mathsf{T}^{\flat}_i\le\mathsf{T}^{\flat}_0\le\bigl(f_{k,m}+\tfrac{9}{16}\theta_S\beta2^{-s}\bigr)
\mathsf{u}_m .
\end{equation}
This holds also at $s=0$, where $f_{k,k}=1$. On the third region $W$ the same bound holds, at a point
of shell $\ell$ of $W$, with the coefficient $F^{(k)}w\,\mathsf{u}_\ell$ in place of
$f_{k,m}\mathsf{u}_m$, where $F^{(k)}$ is as in \eqref{9.1:eq-arrest-factor-b} and $w$ is the factor
with which the reading of the third region states its table there. We note for later use that
$f_{k,m}\le1$: for $s=0$ this is $f_{k,k}=1$, and for $s\ge1$ the older-leaf entry of
hypothesis (iii) at $j'=m$, together with $f_{m,m}=1$, gives
\begin{equation*}
1-f_{k,m}=f_{m,m}-f_{k,m}=\beta(1-2^{-s})\ge0 .
\end{equation*}
These are the first two inequalities of \eqref{9.1:eq-domain-inclusion-2}. It remains to compare the
right side of \eqref{9.1:eq-domain-inclusion-4} with the diminished thresholds of the constituents.

\emph{(c) The region $\mathfrak{T}$: the constituents.} We use hypothesis (iii) (in particular
$\mathbf{r}^T_a=r_a$ for $\ell_T\ge m$) and the two sums of hypothesis (i).

\emph{(c1) The newborn leaf $\mathbf{B}^{(k)}$ born at the step $T_k$, on a shell $m<k$.} By
hypothesis (iv), $\theta_{T,a}=(f_{k,m}+\theta_S\beta2^{-s})r_a$. Since $\mathbf{r}^T_a=r_a$ and the
$\iota$-sum is bounded by hypothesis (i),
\begin{equation*}
\theta^{(i)}_{T,a}\ge\theta_{T,a}-\sum_{l\le i}\iota_lr_a
\ge\bigl(f_{k,m}+\theta_S\beta2^{-s}-\tfrac1{48}\beta e^{-(s-1)}\bigr)r_a .
\end{equation*}
Now
\begin{equation}\label{9.1:eq-domain-inclusion-a}
\tfrac1{48}e^{-(s-1)}\le\tfrac{3}{16}\theta_S2^{-s}
\iff\theta_S\ge\tfrac29\bigl(\tfrac2e\bigr)^{s-1},
\end{equation}
as one sees by multiplying both sides by $\tfrac{16}{3}2^{s}$; the right-hand condition follows from
$\theta_S\ge\tfrac29$, since $s\ge1$ and $2/e<1$. Hence
$\theta^{(i)}_{T,a}\ge(f_{k,m}+\tfrac{13}{16}\theta_S\beta2^{-s})r_a$, which exceeds the right side of
\eqref{9.1:eq-domain-inclusion-4} by the margin $\tfrac14\theta_S\beta2^{-s}$.

\emph{(c2) Older leaves $\mathbf{B}^{(j')}$, $\mathbf{B}'^{(j')}$ with $m\le j'\le k-1$, and live
left-overs of arrests $\mathfrak{a}'$ with $k_{a'}<k$.} For the older leaves hypothesis (iii) gives
\begin{equation*}
r_T=f_{j',m}-f_{k,m}=\beta2^{-s}(2^{k-j'}-1)\ge\beta2^{-s},
\end{equation*}
because $k-j'\ge1$. For a live piece of a left-over term of such an arrest $\mathfrak{a}'$, hypothesis
(v) at $k'=k>k_{a'}$ gives, at a point of shell $\ell=m$, in the frozen structure of $\mathfrak{a}'$,
\begin{equation*}
\varphi^{(n^{+}_{a'})}_{\ell}-F^{(k)}\ge\beta2^{-(k_{a'}-\ell)}-\beta2^{-(k-\ell)}
\ge\beta2^{-(k-\ell)},
\end{equation*}
the second inequality being $2^{-(k_{a'}-\ell)}\ge2^{-(k-1-\ell)}$, that is $k_{a'}\le k-1$; this is
the same lower bound $r_T\ge\beta2^{-s}$. In both cases, by hypothesis (i),
$\theta^{(i)}_{T,a}\ge(f_{k,m}+\beta2^{-s}-\tfrac1{48}\beta e^{-(s-1)})r_a$. Further
\begin{equation*}
\tfrac{9}{16}\theta_S2^{-s}+\tfrac1{48}e^{-(s-1)}\le\bigl(\tfrac12+\tfrac1{24}\bigr)2^{-s},
\end{equation*}
because $\tfrac{9}{16}\theta_S\le\tfrac{9}{16}\cdot\tfrac89=\tfrac12$ and
$\tfrac1{48}e^{-(s-1)}2^{s}=\tfrac1{24}(2/e)^{s-1}\le\tfrac1{24}$. Hence $\theta^{(i)}_{T,a}$ exceeds the
right side of \eqref{9.1:eq-domain-inclusion-4} by at least
$(1-\tfrac12-\tfrac1{24})\beta2^{-s}=\tfrac{11}{24}\beta2^{-s}$.

\emph{(c3) $\mathbb{E}^{(m)}$- and $\mathbb{R}^{(m)}$-terms with $\ell_T=m<k$.} By hypothesis (iii),
$\theta_{T,a}=r_a$, $\theta_{\mathrm{src},a}=f_{k,m}r_a$ and $r_T=\beta(1-2^{-s})$; by (b),
$\theta_{T,a}-\theta_{\mathrm{src},a}=(1-f_{k,m})r_a=\beta(1-2^{-s})r_a$. By
\eqref{9.1:eq-domain-inclusion-3}, with $\mathbf{r}^T_a=r_a$ and the $\iota^1$-sum of hypothesis (i),
\begin{equation*}
\theta^{(i)}_{T,a}=\theta_{T,a}-\tfrac13(\theta_{T,a}-\theta_{\mathrm{src},a})
-\sum_{l\le i}\iota^1_l\mathbf{r}^T_a
\ge\bigl(f_{k,m}+\tfrac23\beta(1-2^{-s})-\tfrac{\beta}{12}\bigr)r_a .
\end{equation*}
Comparison with \eqref{9.1:eq-domain-inclusion-4} requires
\begin{equation}\label{9.1:eq-domain-inclusion-5}
\tfrac{9}{16}\theta_S2^{-s}\le\tfrac23(1-2^{-s})-\tfrac1{12}.
\end{equation}
At $s=1$ this reads $\tfrac{9}{32}\theta_S\le\tfrac13-\tfrac1{12}=\tfrac14$, that is
$\theta_S\le\tfrac89$; this is the origin of the upper end of the interval for $\theta_S$, and at
$\theta_S=\tfrac89$, $s=1$ equality holds. For $s\ge2$ the right side of
\eqref{9.1:eq-domain-inclusion-5} is at least its value $\tfrac5{12}$ at $s=2$ and the left side is at
most its value $\tfrac{9}{64}\theta_S$ at $s=2$, and $\tfrac5{12}\ge\tfrac{9}{64}\theta_S$. The margin is
the difference of the two sides of \eqref{9.1:eq-domain-inclusion-5} times $\beta$, which is the
expression displayed in the statement.

\emph{(c4) Scale drops, $\ell_T<m$.} By hypothesis (ii),
$\theta^{(i)}_{T,a}\ge\tfrac32r_a>(1+\beta)r_a$, and $(1+\beta)r_a$ is at least the right side of
\eqref{9.1:eq-domain-inclusion-4}, because $f_{k,m}\le1$ by (b) and $\tfrac{9}{16}\theta_S\le\tfrac12$.

\emph{(c5) $\Omega_k\cap\mathfrak{T}$, where $m=k$ and $s=0$.} By (b), at $s=0$ the room increments sum
to at most $\tfrac1{16}\theta_S\beta$ and the table is at most $1+\tfrac{9}{16}\theta_S\beta$. For the
newborn leaf $\mathbf{B}^{(k)}$ born at $T_k$, on its top layer, hypothesis (iv) gives
$\theta_{T,a}=(1+\beta)r_a$; subtracting the $\iota$-sum, at most $\tfrac{\beta}{72}$ on $\Omega_k$ by
hypothesis (i), gives $\theta^{(i)}_{T,a}\ge(1+\tfrac{71}{72}\beta)r_a$, which exceeds
$1+\tfrac12\beta\ge1+\tfrac{9}{16}\theta_S\beta$ by at least
$(\tfrac{71}{72}-\tfrac12)\beta=\tfrac{35}{72}\beta\ge\tfrac1{18}\beta$. For the newborn
$\mathbb{E}^{(k)}$- and
$\mathbb{R}^{(k)}$-terms, hypothesis (iv) gives $\theta_{T,a}=(1+\beta)r_a$ and
$\theta_{\mathrm{src},a}=f_{k,k}r_a=r_a$, so that $\theta_{T,a}-\theta_{\mathrm{src},a}=\beta r_a$; this is
the room $r_T=\beta$ of hypothesis (iii), and it is the top-layer factor $1+\beta$
of the layer-wise schedule. The subtracted stage increments are, by hypothesis (i), at
most
\begin{equation*}
\sum_l\tfrac23\mathsf{w}_0\rho^{(l)}_k\le\tfrac43\mathsf{w}_0\beta=\tfrac{\beta}{18},
\end{equation*}
because $\sum_l\rho^{(l)}_k<2\beta$. By \eqref{9.1:eq-domain-inclusion-3}, with these values of
$\theta_{T,a}$ and $\theta_{\mathrm{src},a}$,
\begin{equation*}
\theta^{(i)}_{T,a}\ge\bigl(1+\beta-\tfrac13\beta-\tfrac{\beta}{18}\bigr)r_a
=\bigl(1+\tfrac{11}{18}\beta\bigr)r_a\ge\bigl(1+\tfrac12\beta\bigr)r_a
\ge\bigl(1+\tfrac{9}{16}\theta_S\beta\bigr)r_a,
\end{equation*}
the last step by $\theta_S\le\tfrac89$. The margin is $(\tfrac{11}{18}-\tfrac{9}{16}\theta_S)\beta$,
which is at least $(\tfrac{11}{18}-\tfrac12)\beta=\tfrac19\beta$, attained at $\theta_S=\tfrac89$ and
larger below; in particular it is at least $\tfrac1{18}\beta$. No other constituents occur on $\Omega_k$
outside (c4): an older $\mathbb{E}^{(j')}$- or $\mathbb{R}^{(j')}$-term, $j'<k$, on $\Omega_k$ has
$\ell_T=j'<m=k$ and is a scale drop, treated in (c4) with $\theta^{(i)}_{T,a}\ge\tfrac32r_a$; and older
leaves with $\ell_T\ge k$ do not exist. So on $\Omega_k$ the constituents are exactly the newborn leaf
$\mathbf{B}^{(k)}$, the newborn $\mathbb{E}^{(k)}$- and $\mathbb{R}^{(k)}$-terms, and the scale drops.

\emph{(c6) The indices $a\in\{5,8\}$.} For these indices the clause of the table is that of the source
term, as in the definition of $\mathfrak{U}[\mathsf{T}]$. The constituents' thresholds for these indices
are never incremented in the stage, $\theta^{(i)}_{T,a}=\theta_{T,a}\ge\theta_{\mathrm{src},a}$, and the
direct clause of $\mathfrak{N}^{(i)}$ for them is at $(1+2\beta)r_a$; both are met by a pair obeying the
source term's clause.

Parts (c1)--(c6) give the last inequality of \eqref{9.1:eq-domain-inclusion-2} with the margins of
the statement, hence the clauses of $\mathfrak{I}_i$ on $X\cap\mathfrak{T}$, by the strict-inequality
remark in the statement.

\emph{(d) The set $\mathfrak{N}^{(i)}$ and the extension to $X^{+}$.} For the direct entries on
$X\cap\mathfrak{T}$, the factor of $\mathfrak{N}^{(i)}$ exceeds the table coefficient of
\eqref{9.1:eq-domain-inclusion-4}, which is at most $1+\beta+\tfrac{9}{16}\theta_S\beta$ by $f_{k,m}\le1$,
by
\begin{equation*}
\Bigl(1+2\beta-\sum_l\iota^1_l\Bigr)-\bigl(1+\beta+\tfrac{9}{16}\theta_S\beta\bigr)
\ge2\beta-\tfrac{\beta}{12}-\beta-\tfrac{9}{16}\theta_S\beta\ge\tfrac5{12}\beta>0,
\end{equation*}
using hypothesis (i) and $\tfrac{9}{16}\theta_S\le\tfrac12$.

On $X^{+}\setminus X$ we have $X^{+}\setminus X\subset K(X)\subset Y$, where $Y$ denotes the region on
which the entry bound of the tube $\mathfrak{T}_{\Phi}$ holds by its definition in
Theorem~\ref{9.1:thm-tube-theorem}, and where $K(X)$ is the union of the $M$-cubes meeting the cube
enlargement of $X$ about which that theorem forms its box tower; $K(X)$ contains
$X\cup\bigcup_{x\in X}B(x)=X^{+}$ by two properties of the cube hull $K(X)$ stated with its
definition, since each $B(x)$ has radius less than one $M$-cube of the local
scale. The pairs of $\mathfrak{U}[\mathsf{T}^{\flat}_i](X)$ are, by hypothesis, given on $X^{+}$ and
obey the table there. The slices of the tube theorem are such pairs. Indeed, the slice $S(q)$ of
Theorem~\ref{9.1:thm-tube-theorem} is by definition the restriction of the chart configuration
$(\mathbf{V}(B[q]),\mathbf{J}(\mathbf{V}(B[q])))$ of that theorem, which is defined on the whole box
tower, with boxes $\mathfrak{B}$, about the cube enlargement of $X$, and this tower contains
$K(X)\supseteq X^{+}$ because of its margins of size $R_1M_1$. Moreover the bounds of
Theorem~\ref{9.1:thm-tube-theorem} are
pointwise: the entry bound $E_a\le\mathsf{T}_e$ holds on the region $Y\supset X^{+}$, by the definition
of the tube $\mathfrak{T}_{\Phi}$ (for the root of the tube, on the box tower about the cube
enlargement of $Y$); the
profile bounds of the moves (the local Lipschitz bounds, the difference bounds, the gap bound, the link
and cube-size conditions, and the bound through the constant $K^{\sharp}_R$ of that theorem) hold at
every point of the chart's
domain; and the unit of $F$ is continued to $y\in X^{+}\setminus X$ by
$n_F(y):=\min\{j,m(y)\}$. Therefore the table inequality \eqref{9.1:eq-domain-inclusion-4}, and with
it the comparisons (b)--(c), hold on $X^{+}$ with the same margins.

\emph{(e) The composite entries.} Let $T$ be an $\mathbb{E}$- or $\mathbb{R}$-term and $a\in\{2,4\}$,
at a summation triple $(X',p)$ with $X'\subset X$ and of level $p$, at which the composite clause is
imposed. For $p\le j$ the
composite entries of a pair in the chart are its own entries of the types $E_1$ and $E_3$, because the
boxes on which the composite entries are formed refine the boxes $\mathfrak{B}$ of the tower; hence they
lie below the same thresholds by (c3). For $p>j$ the composite clause is not established here: what is
used is only the rider that the points of the chart lie in the stage spaces, supplied by the regularity
theorem for the chart configurations and by its corollary on the composite entries, used here as
stated there; the composite clause at $p>j$ itself is the content of the composite-box statement and is
not proved here.

\emph{(f) Chain-term constituents.} The constituents $\mathbb{C}^{(i')}(S,\mathcal{S}')$ with $i'<i$
and $S\subset X$ are not subject to a clause of the table: they are treated within the theorem cited
in Notation~\ref{9.1:not-inherited-domain}, whose part (a) bounds the chain terms on
$\mathfrak{D}^{\flat}_i(X)$; nothing is asked of the table for
them beyond (a)--(e).

By (a), (c) and (d), a pair of $\mathfrak{U}[\mathsf{T}^{\flat}_i](X)$ satisfies the clauses of
$\mathfrak{L}_i$, $\mathfrak{I}_i$ and $\mathfrak{N}^{(i)}$, with (e) for the composite entries and (f)
for the chain terms; this is \eqref{9.1:eq-domain-inclusion-1}. By (d) the slices placed in
$\mathfrak{T}_F[\varsigma](X)$ are such pairs. No condition of the family $\gamma_U$ is used; the
inequalities in (c1) need $\theta_S\ge\tfrac29$ and those in (c2)--(c5) and (d) need
$\theta_S\le\tfrac89$.
\end{proof}

\begin{remark}[Analyticity and exponential bound of left-over terms]\label{9.1:rem-leftover-bound}
Work at the $\mathbf{R}$-site of step $k=k_a$, under the hypotheses of
Lemma~\ref{9.1:lem-domain-inclusion}: $\theta_S\in[\tfrac12,\tfrac89]$, $\mathsf{w}_0=1/24$,
$\mathsf{c}_e=\theta_S$; a stage $i\le\mathsf{N}$; and a domain $X$ missing the raw core, with
$X\cap W_{\mathfrak{a}}=\emptyset$ for the component $W_{\mathfrak{a}}$ of the arrest $\mathfrak{a}$ at
step $k_a$. Let $F=\mathbb{C}^{(i)}_{k_a}(X)$ be the left-over
term with domain $\mathfrak{D}_F=\mathfrak{U}[\mathsf{T}^{\flat}_i](X)$ of
Definition~\ref{9.1:def-leftover-margin}. This domain is read with the same conventions as the inherited
domain:
\begin{itemize}
\item the pair is given on the enlargement $X^{+}$ and satisfies there the strict inequalities that
define $\mathfrak{U}[\mathsf{T}^{\flat}_i]$;
\item the entries $E_5$ and the cube clause of $\mathfrak{U}[\mathsf{T}]$ are those of its definition,
with the thresholds of the source term;
\item the composite entries are those specified in Lemma~\ref{9.1:lem-domain-inclusion}.
\end{itemize}
Assume the following two statements.
\begin{itemize}
\item[(1)] The bound stated in Definition~\ref{9.1:not-inherited-domain}: $\mathbb{C}^{(i)}_{k_a}(X)$ is
holomorphic on the inherited domain $\mathfrak{D}^{\flat}_i(X)$ of \eqref{9.1:eq-inherited-domain-b}, and it
is bounded there by $C_0^{\sharp}e^{-a_m d_m(X)}$, with the constant $C_0^{\sharp}$, the rate $a_m$ and
the distance $d_m(X)$ of Definition~\ref{9.1:not-inherited-domain}.
\item[(2)] The inclusion \eqref{9.1:eq-domain-inclusion-a} of Lemma~\ref{9.1:lem-domain-inclusion}.
\end{itemize}
Then $F$ is holomorphic on $\mathfrak{D}_F$, and it is bounded there by $C_0^{\sharp}e^{-a_m d_m(X)}$.
When $X$ meets $W_{\mathfrak{a}}$, the inclusion (2) is not provided by
Lemma~\ref{9.1:lem-domain-inclusion}; the comparison on $W_{\mathfrak{a}}$ at its own arrest step is the
subject of the increment statement, and that case is not covered by this remark.

The window $\theta_S\in[\tfrac12,\tfrac89]$, imposed on the schedule fraction $\theta_S$ among this
chapter's conditions on the schedule fractions and assumed in
Lemma~\ref{9.1:lem-domain-inclusion}, has the following origin.
\begin{itemize}
\item The lower endpoint $\tfrac12$ is fixed together with $\mathsf{w}_0=1/24$ in the hypotheses under
which the bound (1) is stated (Definition~\ref{9.1:not-inherited-domain}).
\item The upper endpoint $\tfrac89$ comes from the case of the $\mathbb{E}$- and $\mathbb{R}$-terms with
top shell of level $m<k$ in the proof of Lemma~\ref{9.1:lem-domain-inclusion}; here and below, as in that
lemma, $m$ denotes the level of a shell of $X\cap\mathfrak{T}$. On the first shell below level $k$,
$m=k-1$ (depth $s:=k-m=1$), that case requires $\tfrac{9}{32}\theta_S\le\tfrac14$, that is,
$\theta_S\le\tfrac89$.
\end{itemize}

No condition is placed in $\gamma_U$. At a point $y$ of $X\cap\mathfrak{T}$ lying on an untouched shell
$m<k$, with $s=k-m$, the bound on the stage-increment sums stated in
Definition~\ref{9.1:not-inherited-domain} and taken as a hypothesis in
Lemma~\ref{9.1:lem-domain-inclusion} gives
\begin{equation*}
\sum_{l\le i}\iota_l(y)\le\tfrac{1}{48}\,\beta\,e^{-(s-1)}\le\tfrac{3}{16}\,\theta_S\,\beta\,2^{-s};
\end{equation*}
the second inequality is equivalent to $\theta_S\ge\tfrac29(2/e)^{s-1}$, and so holds as soon as
$\theta_S\ge\tfrac29$, because $(2/e)^{s-1}\le1$ for $s\ge1$.
This is the inequality used in the case of the newborn terms $\mathbf{B}^{(k)}$ on a shell $m<k$ in the
proof of Lemma~\ref{9.1:lem-domain-inclusion}. It holds identically for every $\theta_S\ge\tfrac29$, hence
throughout the window, and is therefore not one of the conditions collected in $\gamma_U$.

The bound (1) is used as stated in Definition~\ref{9.1:not-inherited-domain}, under the hypotheses
stated there. For the composite entries of level $p$ above $j=k_a$ (in the notation of
Definition~\ref{9.1:def-leftover-margin} and Lemma~\ref{9.1:lem-domain-inclusion}), the inclusion (2) is
used only in the form given by Lemma~\ref{9.1:lem-domain-inclusion}, namely under the condition that the
points lie in the stage spaces.
\end{remark}

\begin{proof}
By (2), read with the conventions listed in the statement, we have
$\mathfrak{D}_F=\mathfrak{U}[\mathsf{T}^{\flat}_i](X)\subset\mathfrak{D}^{\flat}_i(X)$. By (1),
$F=\mathbb{C}^{(i)}_{k_a}(X)$ is holomorphic on $\mathfrak{D}^{\flat}_i(X)$. Its restriction to the subset
$\mathfrak{D}_F$ is therefore holomorphic. The bound $C_0^{\sharp}e^{-a_m d_m(X)}$ of (1) holds on all of
$\mathfrak{D}^{\flat}_i(X)$, so it holds on $\mathfrak{D}_F$.
\end{proof}

\begin{lemma}[The share law independent of the site]\label{9.1:lem-share-law-bis}
Let $p_0$ be a seed pair in the first frame, and put $\mathfrak{a}:=A_{p_0}$, where $A_p$ denotes
the quantity attached to a vertex pair $p$ through whose deviation $A_{p_r}-A_{p_0}$ the identity
relating the increment generators of the two machines bounds the difference of these generators
(see \eqref{9.1:eq-share-law-3-bis} in the proof below); its changes along a presentation are
controlled in \eqref{9.1:eq-share-law-a} below. Let $p_r$ be a vertex pair presented over $p_0$ by
increments of arbitrary sites, and let $p_0,p_1,\dots,p_r$ be the natural vertices of this
presentation, $p_s$ being reached from $p_{s-1}$ by the $s$-th increment. The increments have
increment parameter $\eta$ and increment generators $\mathcal{Y}$ and $Y$ of the two machines.
Write $\mathbf{U}$ and $\mathbf{U}'$ for the pair variables read by the first and by the second
machine, and $M(\cdot)$ for the map through which the pair variable $\mathbf{U}'$ of the second
machine is acted on by its generator $Y$. The presentation then reads
\begin{equation}\label{9.1:eq-share-law-1-bis}
\mathbf{U}_r=e^{i\eta\mathcal{Y}}\mathbf{U}_0,\qquad
M(\mathbf{U}'_r)=e^{i\eta Y}M(\mathbf{U}'_0).
\end{equation}
Let $\mathsf{m}_s$ be majorants of the increments common to the two machines, in the sense of
the first majorant condition, with
$\ell_{\iota}|h^{\mathsf{m}}_s|\le\mathsf{m}_s$ on $4\mathfrak{P}_{\mathsf{c}}$, put
\begin{equation*}
\mathsf{a}^{\sharp}:=2\sum_s\sup_{\tilde{\Delta}(\cdot)}\mathsf{m}_s ,
\end{equation*}
and assume that, for source majorants $\mathsf{d}_s$ of the increments (not to be confused with
the separation of supports, which is also denoted $\mathsf{d}$ elsewhere in this book), at real
points $\Pi$ of the fourfold pair tube $4\mathfrak{P}_{\mathsf{c}}$ and with $p_0,\dots,p_r$ the
natural vertices of the presentation, one has, with a constant $c_{*}^{\sharp}$,
\begin{equation}\label{9.1:eq-share-law-a}
\ell_{\iota}\bigl|A_{p_s}-A_{p_{s-1}}\bigr|
\le\tfrac12\,c_{*}^{\sharp}\,\vartheta^{\iota}\bigl[\mathsf{m}_s+\mathsf{d}_s\bigr].
\end{equation}
Put $\mathsf{d}:=\sum_s\mathsf{d}_s$. Then at every cell $c$, of level $l$, of the seed's own set
of cells, on $\mathfrak{P}_{\mathsf{c}}$,
\begin{equation}\label{9.1:eq-share-law-b}
\bigl|\mathfrak{d}^{\mathrm{net}}(c)\bigr|
\le C_{\rho}^{\sharp}\,\vartheta^{l/2}\bigl[\mathsf{a}^{\sharp}+\mathsf{d}\bigr](c),
\end{equation}
where the constant $C_{\rho}^{\sharp}$ does not depend on $r$, $\mathsf{N}$, $N_0$,
$\check{R}_k$, $k$, $g$ or the depth.
\end{lemma}

\begin{proof}
The proof of the share law at the original site uses the run of increments only through
four facts. We establish each of them under the present hypotheses. In what follows $|\cdot|_c$
is the norm at the cell $c$ used in that share law, and $P(c)$ is the set attached to $c$ in
that share law, on which the stage labels $\iota$ of the increments are at least the level $l$
of $c$.

\emph{First fact.} At the cell $c$ both machines are handed one and the same datum. Hence,
once the common part $\mathsf{h}$ of that datum is taken out, $\mathfrak{d}^{\mathrm{net}}(c)$
vanishes at $\mathcal{Y}=Y=0$. The two bounds used in the first step of the proof of the share
law at the original site, together with the one-datum bound (the bound for the net share at a
cell at which both machines are handed one datum), then give
\begin{equation}\label{9.1:eq-share-law-2-bis}
\bigl|\mathfrak{d}^{\mathrm{net}}(c)\bigr|
\le 68d\,|Y-\mathcal{Y}|_c
+\bigl(C_W+300d\bigr)|\mathfrak{a}|_c\bigl(|Y|_c+|\mathcal{Y}|_c\bigr).
\end{equation}
Here $C_W$ is the constant of the two-machine comparison lemma, which is applied together with
two uses of the Schwarz lemma.

\emph{Second fact.} The identity relating the two increment generators to the pair
variables gives
\begin{equation}\label{9.1:eq-share-law-3-bis}
|Y-\mathcal{Y}|
\le 2\bigl|A_{p_r}-\mathfrak{a}\bigr|+2C_B\,\eta\bigl(|Y|+|\mathcal{Y}|\bigr)|\mathfrak{a}|,
\end{equation}
where $C_B$ is the constant of that identity.

\emph{Third fact: the bound at real arguments.} Let $\Pi$ be real. Since
$\mathfrak{a}=A_{p_0}$, the difference telescopes:
\begin{equation*}
A_{p_r}-\mathfrak{a}=\sum_{s=1}^{r}\bigl(A_{p_s}-A_{p_{s-1}}\bigr).
\end{equation*}
Each summand is bounded by \eqref{9.1:eq-share-law-a}. On $P(c)$ one has $\iota\ge l$, so
that the factor $\vartheta^{\iota}$ in \eqref{9.1:eq-share-law-a} is bounded by
$\vartheta^{l}$. For the seed,
\begin{equation*}
|\mathfrak{a}|_c\le y_0^{\mathsf{T}}\,\bar{\bar{\mathsf{u}}}\,\vartheta^{l},
\end{equation*}
by the bound of the seed pair by $y_0^{\mathsf{T}}\bar{\bar{\mathsf{u}}}$. For the
generators,
\begin{equation*}
|Y|_c+|\mathcal{Y}|_c\le3\,\mathsf{a}^{\sharp}(c).
\end{equation*}
We insert the telescoped form of \eqref{9.1:eq-share-law-a} and these two bounds into
\eqref{9.1:eq-share-law-3-bis}, and the result into \eqref{9.1:eq-share-law-2-bis}. This gives, at
real $\Pi$,
\begin{equation}\label{9.1:eq-share-law-4-bis}
\bigl|\mathfrak{d}^{\mathrm{net}}(c)\bigr|
\le C^{\mathbb{R}}\,\vartheta^{l}\bigl[\mathsf{a}^{\sharp}+\mathsf{d}\bigr](c),
\end{equation}
with a constant $C^{\mathbb{R}}$ formed from $d$, $c_{*}^{\sharp}$, $C_W$, $C_B$, $\eta$ and
$y_0^{\mathsf{T}}\bar{\bar{\mathsf{u}}}$. The sum over $s=1,\dots,r$ of the right-hand sides of
\eqref{9.1:eq-share-law-a} is bounded, by the definitions of $\mathsf{a}^{\sharp}$ and
$\mathsf{d}$, by $\tfrac12c_{*}^{\sharp}\vartheta^{l}[\mathsf{a}^{\sharp}+\mathsf{d}](c)$, so that
$C^{\mathbb{R}}$ does not depend on $r$.

\emph{Fourth fact: the a priori bound and holomorphy.} On the whole of
$4\mathfrak{P}_{\mathsf{c}}$, and because both machines start from the common datum,
\begin{equation}\label{9.1:eq-share-law-5-bis}
\bigl|\mathfrak{d}^{\mathrm{net}}(c)\bigr|\le105d\,\mathsf{a}^{\sharp}(c).
\end{equation}
Moreover, $\mathfrak{d}^{\mathrm{net}}(c)$ is holomorphic in $\Pi$ on
$4\mathfrak{P}_{\mathsf{c}}$. This holds under the holomorphy hypothesis of the tube theorem
at every site (Theorem~\ref{9.1:thm-tube-theorem}). The passage from $4\mathfrak{P}_{\mathsf{c}}$ to
$\mathfrak{P}_{\mathsf{c}}$ uses a single Cauchy estimate, at the cost of one factor $2C$, where
$C$ is the constant of the Cauchy bound that accompanies the identity of the second fact. This
estimate comes in addition
to the Cauchy estimate already spent in the corollary on runs of the share law at the original
site.

None of the four facts refers to the site at which an increment of the presentation
\eqref{9.1:eq-share-law-1-bis} is made. With \eqref{9.1:eq-share-law-2-bis},
\eqref{9.1:eq-share-law-3-bis}, \eqref{9.1:eq-share-law-4-bis} and \eqref{9.1:eq-share-law-5-bis} in
hand, the argument of the share law at the original site applies without change. It yields
\eqref{9.1:eq-share-law-b} on $\mathfrak{P}_{\mathsf{c}}$, with a constant $C_{\rho}^{\sharp}$
built from the constants of \eqref{9.1:eq-share-law-2-bis}--\eqref{9.1:eq-share-law-5-bis} and the
one Cauchy factor, none of which involves $r$, $\mathsf{N}$, $N_0$, $\check{R}_k$, $k$, $g$ or
the depth.
\end{proof}

\begin{definition}[Link-by-link majorants and the layer source]\label{9.1:not-link-majorants}
The share law independent of the site, Lemma~\ref{9.1:lem-share-law-bis}, takes as its hypothesis
the inequality \eqref{9.1:eq-share-law-a} for each increment $p_{s-1}\to p_s$ of a presentation.
In that inequality a majorant $\mathsf{m}_s$ and a source term $\mathsf{d}_s$ are attached to the
$s$-th increment. We fix them according to the kind of the increment:
\begin{itemize}
\item[(P)] for a move at stage $\nu$, $\mathsf{m}_s=\min\{\ell a^{0}_{\nu}/K_{o},\,\mathsf{h}_{\nu}\}$
and $\mathsf{d}_s=0$, where $\ell$ (a letter distinct from the $\ell(n,R,g)$ of the glossary),
$a^{0}_{\nu}$ and $K_{o}$ are the stage-$\nu$ quantities whose quotient $\ell a^{0}_{\nu}/K_{o}$
majorises the move in the move bound, and $\mathsf{h}_{\nu}$ is its move majorant;
\item[(T)] for a shift, $\mathsf{m}_s=\mathsf{m}_T$ and $\mathsf{d}_s=0$;
\item[(a,b)] for the dilated and the cut links of kind $(\ell1)$ of stages (a) and (b) of the
staged construction of the links, $\mathsf{m}_s\in\{\mathsf{m}_L,(1+|\mathsf{t}|)\mathsf{m}_p\}$,
and $\mathsf{d}_s=0$;
\item[(c)] for the links of kind $(\ell2)$ of stage (c), carried by its three pieces
$\Xi_3,\Xi_2,\Xi_1$, $\mathsf{m}_s=\mathsf{m}_L$ and $\mathsf{d}_s=\mathsf{d}_{\mathrm{lay}}$, where
\begin{equation}\label{9.1:eq-link-majorants-a}
\mathsf{d}_{\mathrm{lay}}:=C\,\mathsf{u}\,e^{-\delta_1 d(\cdot,\mathrm{lay})/4}.
\end{equation}
Here $\mathsf{u}$ is the slow local radius of the random-walk construction,
$d(\cdot,\mathrm{lay})$ is the distance to the layer strata of stage (c), $\delta_1$ is the decay
rate of the factor $e^{-\delta_1 d}$ carried by the source terms of the relative bounds (B)
and (C), and $C$ is the constant produced by the case of stage (c) in the proof below.
\end{itemize}
We call $\mathsf{d}_{\mathrm{lay}}$ the \emph{layer source}.

Assume the following statements:
\begin{itemize}
\item the relative bounds (B) and (C) for the increment generator;
\item the move bound, and the bound with constant $C^{\natural}$ of the same statement;
\item the rate bound;
\item clause E(b) of the entry bounds, which allows the relative bounds to be applied over the
entries;
\item clauses (a)--(c) of the staged construction of the links;
\item the corollary on links and clause (b) of the initial-condition statement;
\item clause (e) of the naturality of the tilt;
\item the level--distance condition (d1), which, for a source term of level $q$ and index
$\vartheta^{q}$ seen at a point $x$ at distance $d(x)$ from it, bounds
$\vartheta^{q}e^{-\delta_1 d(x)/4}$ by $L^{4}\vartheta^{\lambda(x)}$, where $\lambda(x)$ is the
level attached to $x$;
\item properties (f3) and (N), by which the background increment of a shift is bounded, and
property $(\Lambda)$, which with (f3) bounds the background increment of a link;
\item the identity comparing the generators with the pair variables.
\end{itemize}
Then, with these choices of $\mathsf{m}_s$ and $\mathsf{d}_s$, the inequality
\eqref{9.1:eq-share-law-a} holds for every increment, at real points $\Pi$ of the fourfold pair
tube $4\mathfrak{P}_{\mathsf{c}}$ and at natural vertices, as required in
Lemma~\ref{9.1:lem-share-law-bis}.
\end{definition}

\begin{proof}
The bounds (B) and (C) are relative bounds: in their notation,
\begin{equation*}
\bigl|Y[\tau\mathbb{H}';\mathbf{U}']-\tau\mathbb{H}\bigr|_x
\le C\bar{\tau}\bigl\{\mathsf{P}_b\Psi_p+\mathsf{P}_{\delta b}\bigr\}
+CL^{-\lambda}\mathsf{H}\bigl(M^{-1}+\mathsf{H}\bigr),
\qquad \mathsf{H}=\bar{\tau}\mathsf{P}_b ,
\end{equation*}
and the factor $\Psi_p$ satisfies
\begin{equation*}
\Psi_p=\mathsf{T}^{\sim}_p+\mathfrak{r}_0L^{-\theta\lambda}
\le C\bigl(C^{P}\bar{\bar{\mathsf{u}}}+\mathfrak{r}_0\bigr)\vartheta^{\lambda}.
\end{equation*}
Here and below $\lambda(x)$ and $m(y)$ are levels attached to the points $x$ and $y$, and the
exponents $p$ and $m$ in $L^{p-m}$ are levels (the letter $m$ does not denote the torus exponent
there). In the two displays above, $\tau$ is the factor multiplying $\mathbb{H}'$ and
$\mathbb{H}$; $\bar{\tau}$ is the factor which, multiplied by $\mathsf{P}_b$, gives $\mathsf{H}$;
$\lambda$ is the level exponent carried by the decay factor $\vartheta^{\lambda}$ and by
$L^{-\lambda}$; $\theta$ is the fraction of that exponent which appears in the remainder
$\mathfrak{r}_0L^{-\theta\lambda}$ of $\Psi_p$; $\mathsf{T}^{\sim}_p$ is the entry at $p$ of the
threshold table, the other summand of $\Psi_p$; and $C^{P}$ is the coefficient of the unit
$\bar{\bar{\mathsf{u}}}$ in the bound on $\Psi_p$. The letter $C$ denotes a constant, not
necessarily the same at each occurrence; in \eqref{9.1:eq-link-majorants-a} it is the constant
produced by the case of stage (c) below.
In each of the four cases below the bound on the generator is followed by the identity
comparing the generators with the pair variables. This passes from the generator to the
difference $A_{p_s}-A_{p_{s-1}}$ that appears in \eqref{9.1:eq-share-law-a}.

The source terms are converted as follows. Consider a source term of level $q$ carrying the index
$\vartheta^{q}$, seen at a point $x$ at distance $d(x)$ from it. By condition (d1) it satisfies
\begin{equation}\label{9.1:eq-link-majorants-1}
\vartheta^{q}e^{-\delta_1 d(x)/4}\le L^{4}\vartheta^{\lambda(x)} .
\end{equation}
Accordingly we split the decay factor as
\begin{equation*}
e^{-\delta_1 d(x)}=e^{-\delta_1 d(x)/4}\,e^{-\delta_1 d(x)/2}\,e^{-\delta_1 d(x)/4}.
\end{equation*}
The first quarter converts the index by \eqref{9.1:eq-link-majorants-1}. The half is kept for the
sum over the source terms. The last quarter remains in the profile, and it is the exponent in
\eqref{9.1:eq-link-majorants-a}.

\emph{Moves.} The move bound is the relative bound (C) combined with the rate bound. It gives
\eqref{9.1:eq-share-law-a} with $\mathsf{m}_s=\min\{\ell a^{0}_{\nu}/K_{o},\mathsf{h}_{\nu}\}$ and
no source term.

\emph{Shifts.} We apply the relative bound (C) on the pair of box towers
$(\boldsymbol{\Omega}'(Z),\boldsymbol{\Omega}(Z))$. Clause (e) of the naturality of the tilt
allows this application. The background increment is $\delta b=(\mathrm{Ad}_g-1)b$, where $b$ is
the background and $\mathrm{Ad}_g$ is the adjoint action of a group element $g$ (here $g$ is not
the reference coupling), and it is
bounded by properties (f3) and (N). This gives $\mathsf{m}_s=\mathsf{m}_T$ and no source term.

\emph{Dilated and cut $(\ell1)$-links, stages (a) and (b).} We apply the relative bounds (B) and
(C) over the entries, as clause E(b) of the entry bounds allows. The background increment at a
link point $y$ satisfies
\begin{equation}\label{9.1:eq-link-majorants-2}
|\delta b(y)|\le c_{\delta}\,\mathsf{s}_L\,\vartheta^{m(y)} ,
\end{equation}
where $c_{\delta}$ is the constant of this bound.
This follows from properties (f3) and $(\Lambda)$, the corollary on links, and clause (b) of the
initial-condition statement. The result is, for each such link, one of the majorants
$\mathsf{m}_L$, $(1+|\mathsf{t}|)\mathsf{m}_p$, with no source term.

\emph{The $(\ell2)$-links of stage (c), pieces $\Xi_3,\Xi_2,\Xi_1$.} We apply the relative bound
(B) on $\mathfrak{S}_2$. The layer background and its increment obey
\begin{equation*}
|b_{\mathrm{lay}}|\le 136\,d\,L^{p-m}\mathsf{m}'
\quad\text{and}\quad
|\delta b_{\mathrm{lay}}(\cdot)|\le C L^{p-m}\mathsf{u}\,\vartheta^{\lambda_0},
\end{equation*}
where the factor $136\,d$ is the coefficient of the bound on the layer background in clause (c)
of the staged construction (the letter $d$ there is not the distance $d(x)$ used above),
$\lambda_0$ is the level exponent in the bound on the increment, and the increment
$\delta b_{\mathrm{lay}}(\cdot)$ is evaluated at the arguments of clause (c).
These two bounds come from clause (c) of the staged construction and the bound with constant
$C^{\natural}$ of the statement containing the move bound. With these two bounds the relative
bound (B) on
$\mathfrak{S}_2$ yields \eqref{9.1:eq-share-law-a} with $\mathsf{m}_s=\mathsf{m}_L$ and with a
source term proportional to the slow local radius $\mathsf{u}$. After the conversion by
\eqref{9.1:eq-link-majorants-1} and the splitting of $e^{-\delta_1 d(x)}$, this source term
becomes the layer source \eqref{9.1:eq-link-majorants-a}.
\end{proof}

The layer source $\mathsf{d}_{\mathrm{lay}}$ of \eqref{9.1:eq-link-majorants-a}, entering the link-by-link majorants of
Notation~\ref{9.1:not-link-majorants}, is measured in radius units. It is a source of distance type: it decays
with the distance $d(\cdot,\mathrm{lay})$ to the layer strata. The following lemma converts that decay into
decay in the depth.

\begin{lemma}[Decay in depth of the distance to the layer strata]\label{9.1:lem-layer-depth-decay}
Let $X$ be a domain whose cube hull is cut at level $j$, and let $Y_i$ be the enlargement of that hull
relative to which the layer strata are placed. Let
$x\in X$, let $F$ be the term read at $x$, of level $n_F(x)$, and suppose that $x$ is an $(=)$-point of the
share law in which $\mathsf{d}_{\mathrm{lay}}$ enters as a source. Let
\begin{equation*}
s=j-n_F(x)
\end{equation*}
be the depth of $x$.
Let $d(x,\mathrm{lay})$ be the distance from $x$ to the layer strata, measured in units of level $n_F(x)$,
and let $\delta_1$ be the constant in the exponent of the layer source \eqref{9.1:eq-link-majorants-a}. Then
\begin{equation}\label{9.1:eq-layer-depth-decay-a}
e^{-\delta_1 d(x,\mathrm{lay})/8}\le C_{\mathrm{dp}}'\,4^{-s}.
\end{equation}
\end{lemma}

\begin{proof}
We first place the layer strata relative to the hull.

\emph{Strata of level above $j$.} A hull cut at level $j$ does not see layer strata of level $p>j$; this is
the rim statement for hulls. Only strata of level $p\le j$ therefore enter $d(x,\mathrm{lay})$.

\emph{Strata of level at most $j$.} By the geometric statement on layer strata, a layer stratum of level
$p\le j$ lies within distance
\begin{equation}\label{9.1:eq-layer-depth-decay-1}
4L^{p}M_1\eta\le 4L^{j}M_1\eta
\end{equation}
of $Y_i^{c}$, where $\eta$ denotes the spacing of the lattice on which the hull and the strata are measured.

\emph{The hull and $Y_i^c$.} By two clauses of the hull-distance statement,
\begin{equation}\label{9.1:eq-layer-depth-decay-2}
\operatorname{dist}(X,Y_i^{c})\ge ML^{j}\eta .
\end{equation}

\emph{Case $p<j$.} Let $\Sigma$ be a layer stratum of level $p<j$, let $d(x,\Sigma)$ denote the distance from
$x$ to $\Sigma$ in units of level $n_F(x)$, and let $y\in\Sigma$. By
\eqref{9.1:eq-layer-depth-decay-1}, there is a point $z\in Y_i^{c}$ with
\begin{equation*}
|y-z|\le 4L^{j}M_1\eta .
\end{equation*}
Since $x\in X$, \eqref{9.1:eq-layer-depth-decay-2} gives $|x-z|\ge ML^{j}\eta$. The triangle inequality then
yields
\begin{equation*}
|x-y|\ge (M-4M_1)L^{j}\eta .
\end{equation*}
Since $y\in\Sigma$ was arbitrary, expressing this in units of level $n_F(x)$ and using $j-n_F(x)=s$ gives
\begin{equation*}
d(x,\Sigma)\ge (M-4M_1)L^{s}.
\end{equation*}
The second case of the depth-decay bound applies to a distance bounded below in this way. It gives
\begin{equation*}
e^{-\delta_1 d(x,\Sigma)/8}\le C_{\mathrm{dp}}'\,4^{-s}.
\end{equation*}

\emph{Case $p=j$.} Let $\Sigma$ be a layer stratum of level $j$, the stratum at the hull's own level, and let
$d(x,\Sigma)$ be as above. The lower bound on the box-tower distance bounds $d_{\mathfrak{S}}$ from below by
its coefficient times $(s-1)_{+}$, exactly as in the first case of the depth-decay bound. As in that case,
this yields
\begin{equation*}
e^{-\delta_1 d(x,\Sigma)/8}\le L^{-2(s-1)_{+}}.
\end{equation*}
It remains to compare $L^{-2(s-1)_{+}}$ with $4\cdot 4^{-s}$. Since the blocking factor $L$ of
Definition~\ref{1.1:def-lattice} is an odd integer greater than one, $L^{2}\ge 4$.
\begin{itemize}
\item If $s\ge 1$, then
\begin{equation*}
L^{-2(s-1)}\le 4^{-(s-1)}=4\cdot 4^{-s}.
\end{equation*}
\item If $s\le 0$, then $(s-1)_{+}=0$ and
\begin{equation*}
L^{-2(s-1)_{+}}=L^{0}=1\le 4^{-s}\le 4\cdot 4^{-s}.
\end{equation*}
\end{itemize}
In both sub-cases
\begin{equation*}
e^{-\delta_1 d(x,\Sigma)/8}\le L^{-2(s-1)_{+}}\le 4\cdot 4^{-s}.
\end{equation*}

\emph{Conclusion.} The distance $d(x,\mathrm{lay})$ is the minimum of $d(x,\Sigma)$ over the layer strata
$\Sigma$ of level $p\le j$, strata of level $p>j$ not entering. Hence $e^{-\delta_1 d(x,\mathrm{lay})/8}$ is
the maximum of $e^{-\delta_1 d(x,\Sigma)/8}$ over these strata, and each of these is bounded by the bound of
its case. This covers every level $p\le j$. This completes the proof.
\end{proof}

\begin{corollary}[Tube landing and the share-phase bound at every site]\label{9.1:cor-tube-landing}
Assume the following.
\begin{enumerate}
\item[(1)] The tube theorem at every site, Theorem~\ref{9.1:thm-tube-theorem}, holds in both
machines $\mathsf{A}$ and $\mathsf{B}$.
\item[(2)] The share law independent of the site, Lemma~\ref{9.1:lem-share-law-bis}, holds. Its
hypothesis \eqref{9.1:eq-share-law-a} is satisfied by every increment, with the majorants
$\mathsf{m}_s$ and the sources $\mathsf{d}_s$ of the table of link-by-link majorants in
Notation~\ref{9.1:not-link-majorants}.
\item[(3)] The exterior estimate for cells of type \textup{(top)} and \textup{(dist)} holds; its
clause~(d) is the one used below.
\item[(4)] The bound on the rider quantity $\mathsf{p}$ on which that exterior estimate is read
holds. In runs it is supplied by the corollary of the share law for runs. Off the sites of type
\textup{(P)} it is the inequality $\|\mathsf{p}\|\le x_{*}$ at $\vartheta_f$; this inequality
holds whenever the two rider conditions, the share-law rider and the square-root rider on
clause~(a) of the single-site tube landing, hold at the yardstick $\mathfrak{s}_u$ of the exterior
estimate.
\item[(5)] The screening estimate, by which the sources read through the exterior estimate are
screened, holds.
\end{enumerate}
For a vertex pair read by a term $F$ on $X$ at level $j$, put
$\vartheta^{\sharp}:=\vartheta^{(j\wedge n_X)/2}$ in runs and
$\vartheta^{\sharp}:=\vartheta_f^{\,j\wedge n_X}$ off the sites of type \textup{(P)}. Put
\begin{equation}\label{9.1:eq-tube-landing-a}
C^{\sharp}_{\mathrm{TL}}:=2^4C^{\sharp}_{\rho}
+\frac{10\,C'_{\times}+C^{\times}_{\mathrm{lay}}}{\theta_{\rho}\beta},
\qquad
8\,C^{\sharp}_{\mathrm{TL}}\,\vartheta_f^{\,n_c}\le 1 .
\end{equation}
Here $C^{\times}_{\mathrm{lay}}$ is the constant \eqref{9.1:eq-tube-landing-3} fixed in the proof
below. Assume the inequality in \eqref{9.1:eq-tube-landing-a}; it is a ceiling on $\gamma$, in the
sense of the word fixed in Notation~\ref{2.3:not-stages-first}.

Then the following holds at every vertex pair of every site on an aligned label, read by a term
$F$ on $X$ at level $j$, where we put $\mathsf{v}:=\mathrm{rd}^{\mathsf{A}}\mathbf{U}$; the letters
of the clauses are those of the single-site tube landing, whose clause~(c) is not restated here.
\begin{enumerate}
\item[(a)] $\mathsf{v}\in\mathfrak{T}^{\mathsf{A}}[0]$ and
$\mathsf{v}\oplus c^{\mathsf{B}}\in\mathfrak{T}^{\mathsf{B}}[\tfrac18]$.
\item[(b)] Let $|s|\le(8C^{\sharp}_{\mathrm{TL}}\vartheta^{\sharp})^{-1}$. Then
$e^{is\mathfrak{d}^{\mathrm{net}}}(\mathsf{v}\oplus c^{\mathsf{B}})\in\mathfrak{T}^{\mathsf{B}}[\tfrac14]$,
and consequently
\begin{equation}\label{9.1:eq-tube-landing-1}
\bigl|\hat{F}^{\mathsf{B}}(\mathrm{rd}^{\mathsf{B}}\mathbf{U}')
-\hat{F}^{\mathsf{B}}(\mathsf{v}\oplus c^{\mathsf{B}})\bigr|
\le 16\,C^{\sharp}_{\mathrm{TL}}\,\vartheta^{\sharp}\,\mathfrak{n}\,e^{-\kappa_T d_j(X)} .
\end{equation}
\item[(d)] The seeds of an output $\Phi$ fill $\mathfrak{T}_{\Phi}[\tfrac78]$. All reads of its
production land in $\mathfrak{T}_F[\tfrac14]$. This includes the outputs $\Phi$ born at an
$\mathbf{R}$-step and the left-over terms.
\end{enumerate}
\end{corollary}

\begin{proof}
We follow the argument of the single-site tube-landing statement, with $\mathsf{G}$ now the gap of
the site at hand. Let $F$ be the term read at the vertex pair, on $X$ at level $j$. Put
\begin{equation}\label{9.1:eq-tube-landing-2}
D^a:=2\,\ell_{\mathfrak{S}}^{-e_a}\,B_{\natural}\,\bigl(|\mathfrak{d}^{\mathrm{net}}|\bigr)^{\approx}.
\end{equation}
Here $(\cdot)^{\approx}$ is the operation denoted in this way in the single-site tube-landing
statement.
We bound the contributions to $D^a$ class by class, over the cells of the seed's own set.

\emph{Cells of type \textup{(lab)}.} At a site of type \textup{(T)} these are all cells of type
\textup{(lab)} of the seed's set. At a site of type \textup{(L)} they are the cells of type
\textup{(lab)} lying on $\mathfrak{O}$, since, by the first line of the table of local radii of the
random-walk construction, the quantity $\lambda_R$ of that construction equals $m$ on its region
$Y_4$. On all these cells, hypothesis~(2) lets us
apply Lemma~\ref{9.1:lem-share-law-bis}, which gives
\begin{equation*}
|\mathfrak{d}^{\mathrm{net}}(c)|\le C^{\sharp}_{\rho}\,\vartheta^{l/2}\,[\mathsf{a}^{\sharp}+\mathsf{d}](c).
\end{equation*}
Clause (A) of the tube theorem, displayed in \eqref{9.1:eq-tube-theorem-b}, reads
\begin{equation*}
2^7B_{\natural}\,c_1\,\mathsf{a}^{\sharp}\le 2\,\ell^{e_a}\,\mathsf{G}.
\end{equation*}
Hence the part of $D^a$ coming from $\mathsf{a}^{\sharp}$, the share part, is at most
$C^{\sharp}_{\rho}\vartheta^{\sharp}\mathsf{G}/16$.

The table of Notation~\ref{9.1:not-link-majorants} gives $\mathsf{d}_s=0$ for every increment
except the third stage of the link construction, whose source is $\mathsf{d}_{\mathrm{lay}}$. We
show that the part of $D^a$ coming from $\mathsf{d}_{\mathrm{lay}}$ is at most
$(C^{\times}_{\mathrm{lay}}/\theta_{\rho}\beta)\,\vartheta^{\sharp}\varrho_F\mathsf{u}_F$.

By \eqref{9.1:eq-link-majorants-a},
$\mathsf{d}_{\mathrm{lay}}=C_{\mathsf{d}}\,\mathsf{u}\,e^{-\delta_1 d(\cdot,\mathrm{lay})/4}$.
Here $\mathsf{u}$ is the slowly varying local radius of the random-walk construction, and
$C_{\mathsf{d}}$ is the constant of the third stage of the link construction in the table of
Notation~\ref{9.1:not-link-majorants}, which is the same constant as in the bound on the layer
increments quoted in that table. We split the factor
$e^{-\delta_1 d}$ into an index part, a sum part and a
profile part, with exponents $\frac14$, $\frac12$ and $\frac14$, as in
Notation~\ref{9.1:not-link-majorants}. The index part is converted by the inequality
$\vartheta^q e^{-\delta_1 d/4}\le L^4\vartheta^{\lambda(x)}$ stated there. We then apply the
smoothing bound with constant $c_1$. This gives
\begin{equation*}
\begin{split}
2\,\ell^{-e_a}B_{\natural}\,\bigl(C^{\sharp}_{\rho}\vartheta^{\lambda}\mathsf{d}_{\mathrm{lay}}\bigr)^{\approx}(x)
&\le 2B_{\natural}C^{\sharp}_{\rho}c_1C_{\mathsf{d}}L^4\cdot\vartheta^{\sharp}\\
&\qquad\cdot e^{-\delta_1 d(x,\mathrm{lay})/8}\cdot\ell_{\mathfrak{S}}^{-e_a}\mathsf{u}(x).
\end{split}
\end{equation*}
Let $x$ be a $(=)$-point of depth $s=j-n_F(x)$. The decay in depth of the distance to the layer
strata, Lemma~\ref{9.1:lem-layer-depth-decay} in the form \eqref{9.1:eq-layer-depth-decay-a}, turns
the exponential factor into
\begin{equation*}
e^{-\delta_1 d(x,\mathrm{lay})/8}\le C'_{\mathrm{dp}}4^{-s}=C'_{\mathrm{dp}}2^{-s}\cdot 2^{-s}.
\end{equation*}
Further, $\ell_{\mathfrak{S}}^{-e_a}\mathsf{u}(x)\le\mathsf{r}_{\mathrm{lay}}\mathsf{u}_F^a(x)$.
Here $\mathsf{r}_{\mathrm{lay}}$ is the ratio of the local radius of the random-walk construction to
the unit of $F$ at the same point. It is a quotient of Ba{\l}aban's radii $\alpha$ and powers of
$L_0$ at a single level. Under the standing hypothesis on units of the single-site analysis it does
not depend on $g$. It also does not contain the factor $\mathsf{r}_1\ell_{\circ}^{2N_0}$, because at
$(=)$-points both quantities are read at the level $n_F$ of $x$ itself. We put
\begin{equation}\label{9.1:eq-tube-landing-3}
C^{\times}_{\mathrm{lay}}:=2B_{\natural}C^{\sharp}_{\rho}c_1C_{\mathsf{d}}L^4C'_{\mathrm{dp}}\,
\mathsf{r}_{\mathrm{lay}};
\end{equation}
this constant does not depend on $N$. At the $(=)$-point $x$ the level $\ell(x)$ entering the
margin function is $n_F(x)$, so one of the two factors $2^{-s}$ equals $2^{-(j-\ell(x))}$, and
$\theta_{\rho}\beta\,2^{-(j-\ell(x))}\le\varrho_F(x)$: for a left-over term
$F=\mathbb{C}^{(i)}_{k_a}(X)$ on $W_{\mathfrak{a}}\cap Z^{\mathrm{on}}$ one has $j=k_a$, and this
is the second member of the maximum in \eqref{9.1:eq-leftover-margin-a}; in every other case
$\varrho_F$ is the margin function of the single-site analysis. The other factor $2^{-s}$ is
bounded by $1$. The three displays then give the bound
$(C^{\times}_{\mathrm{lay}}/\theta_{\rho}\beta)\,\vartheta^{\sharp}\varrho_F\mathsf{u}_F$ for the
part coming from $\mathsf{d}_{\mathrm{lay}}$.

\emph{Cells of type \textup{(top)} and \textup{(dist)}; cells of $Z^{\mathrm{on}}$.} On these cells we
use clause~(d) of the exterior estimate of hypothesis~(3). It is read with the bound of
hypothesis~(4). It is screened by the index inequality of Notation~\ref{9.1:not-link-majorants}
quoted above and by the screening estimate of hypothesis~(5). We combine it with the single-site
depth bound. A cell of $Z^{\mathrm{on}}$ reaches a $(=)$-point of $\mathfrak{O}$ only through a
factor $e^{-2\delta_1 d}$. The geometric bound and the second clause of the double-decay bound then
bound these contributions by $(C'_{\times}/\theta_{\rho}\beta)\,\vartheta^{\sharp}\varrho_F\mathsf{u}_F$.
At $(=)$-points of $Z^{\mathrm{on}}$ we use the single-site share law, in its two cases.

\emph{Comparison with the placement in machine $\mathsf{B}$.} Adding the contributions, with
$C^{\sharp}_{\mathrm{TL}}$ as in \eqref{9.1:eq-tube-landing-a}, gives
\begin{equation*}
D^a\le C^{\sharp}_{\mathrm{TL}}\,\vartheta^{\sharp}\,(\mathsf{G}+\varrho_F\mathsf{u}_F).
\end{equation*}
This is to be compared with the placement $S^{\mathsf{B}}$ of the read in machine $\mathsf{B}$,
which satisfies
\begin{equation*}
S^{\mathsf{B}}\le(\mathsf{c}_F^{+}-\varrho_F)\,\mathsf{u}_F-\tfrac56\,\mathsf{G}.
\end{equation*}
For $|s|\le(8C^{\sharp}_{\mathrm{TL}}\vartheta^{\sharp})^{-1}$ the first display gives
$|s|\,D^a\le\frac18(\mathsf{G}+\varrho_F\mathsf{u}_F)$. The $(\mathrm{up})$-points and the region
$Z''_k$ have rooms that are fixed fractions, and there the radius yardstick serves, as in the
single-site argument.

The first inclusion of (a) is the conclusion of the tube theorem in machine $\mathsf{A}$, since
$\mathsf{v}=\mathrm{rd}^{\mathsf{A}}\mathbf{U}$. The second inclusion of (a) and the inclusion of (b)
follow from the comparison just made. The bound \eqref{9.1:eq-tube-landing-1} follows from the
inclusion of (b) by the Schwarz lemma in the variable $s$, on the disc
$|s|\le(8C^{\sharp}_{\mathrm{TL}}\vartheta^{\sharp})^{-1}$.

\emph{Clause (d).} The seeds and the reads of every output, including the outputs born at an
$\mathbf{R}$-step, are placed by the tube theorem, Theorem~\ref{9.1:thm-tube-theorem}, together
with the schedule bound and its flat form. The later reads of a left-over term are placed by the
schedule bound for later reads of a left-over term, Lemma~\ref{9.1:lem-leftover-schedule}, in the
form \eqref{9.1:eq-leftover-schedule-a}. Together these give (d).
\end{proof}

\begin{remark}[Yardsticks at cells of type \textup{(lab)}]
At cells of type \textup{(lab)} with additive room the yardstick must be a share, and there the
argument above does not use the two rider conditions. In particular, the yardstick $\mathfrak{s}_u$ is
not replaced at such cells by the share $\mathsf{a}(c)$ of the single-site statement. The rider
conditions enter only through the bound off the sites of type \textup{(P)} in hypothesis~(4), at the
yardstick $\mathfrak{s}_u$ of the exterior estimate.
\end{remark}

\begin{proposition}[The delivered tube contains the loaded tube]\label{9.1:prop-delivered-tube}
Assume the hypotheses of the tube theorem at every site
(Theorem~\ref{9.1:thm-tube-theorem}), of the schedule bound for later reads of a left-over term
(Lemma~\ref{9.1:lem-leftover-schedule}) and of tube landing at every site
(Corollary~\ref{9.1:cor-tube-landing}). Let $\Phi$ be an output and let $(F,X,j)$ be a leaf, in
the sense of Theorem~\ref{9.1:thm-tube-theorem}, that reads $\Phi$ at the same scale; we call
this a \emph{hand-over} $\Phi\to F$.
\begin{enumerate}
\item[(i)] In Theorem~\ref{9.1:thm-tube-theorem} an entry is any seed of
$\mathfrak{T}_{\Phi}[7/8]$. The only size conditions on an entry are those of the definition of
the tube: its slice lies in the tube, and the size clause $|B|\le r_c$ of that definition holds,
where the constant $r_c$ of that definition does not depend on $g$. No
condition on the radius of an entry occurs. In particular, for an $\mathbf{R}$-born output
$\Phi$, the seeds of $\mathfrak{T}_{\Phi}[7/8]$ on $\mathfrak{O}$ that lie beyond the table
$\mathsf{T}_N$ are entries.
\item[(ii)] The seed is handed to every read exactly, with coefficient $1$. The constants
$B_{\natural}$, $C_T^{w}$, $c_1$, $B^{\mathsf{s}}$, $\mathsf{r}_1$ and $C_4$ multiply increments
only. They are paid either by $c_M$, $K_R^{\sharp}$ and $\mathsf{w}_L(p)$, through dilations, or
by the two ceilings on $\gamma$ among the hypotheses of
Theorem~\ref{9.1:thm-tube-theorem} that govern the link increments and the increments on
$\mathcal{A}$. None of them restricts a domain of delivery. Here $c_M$ and $K_R^{\sharp}$ are
constants occurring among the hypotheses of Theorem~\ref{9.1:thm-tube-theorem}, and
$\mathsf{w}_L(p)$ is the link weight.
\item[(iii)] At every hand-over $\Phi\to F$, with the fraction $\varsigma$ of a tube counted in
the room of $F$ itself, whose coefficient at $(=)$-points (Theorem~\ref{9.1:thm-tube-theorem}) is
$\mathsf{c}_F^{+}-\varrho_F$, the delivered tube contains the loaded tube:
\begin{equation*}
\mathfrak{T}_F[7/8]\ \supseteq\ \mathfrak{T}_F[1/4].
\end{equation*}
The coefficient of the load at $(=)$-points and the room coefficient are as follows; in them $s$
denotes the number of levels by which the point read lies below the level of $F$.
\begin{enumerate}
\item[(a)] For a $T$-born leaf $F$ of level $k$, read at $\mathbf{R}_k$ on $\mathfrak{O}$, the
coefficient of the load at $(=)$-points is at most
\begin{equation*}
f_{k,m}+\tfrac{9}{16}\,\theta_S\beta\,2^{-s},
\end{equation*}
and the room coefficient satisfies
\begin{equation*}
\mathsf{c}_F^{+}-\varrho_F\ \ge\ f_{k,m}+\tfrac34\,\theta_S\beta\,2^{-s}.
\end{equation*}
\item[(b)] For the same hand-over on $Z^{\mathrm{on}}$, with $\mathsf{G}$ the gap of the site,
the room coefficient is as in (a), and the load at $(=)$-points is bounded by
\begin{equation*}
\mathsf{T}_{\nu}-\tfrac56\,\mathsf{G}\le f_{k,m}.
\end{equation*}
\item[(c)] For stage $n+1$ reading $\mathbb{C}^{(i)}$ with $i\le n$, the terms are read unfolded,
with $\mathsf{T}^{\flat}_n\le\mathsf{T}^{\flat}_i$, rooms $\mathrm{rm}_{\nu}$ and per-stage room
increments $\mathsf{g}_{\nu}$; no separate room coefficient enters.
\item[(d)] For an $\mathbf{R}$-born leaf $F$ of level $k$, read at $T_k$, the coefficient of the
load at $(=)$-points is at most
\begin{equation*}
f_{k,n}-\tfrac14\,(1-\theta_S)\beta\,2^{-s},
\end{equation*}
and the room coefficient is
\begin{equation*}
\mathsf{c}_F^{+}-\varrho_F=f_{k,n}-\theta_{\rho}\beta\,2^{-s}.
\end{equation*}
\item[(e)] For a left-over term $F$ read at $T_{k_a}$, the coefficient of the load at
$(=)$-points is at most the diminished layer factor $F^{(k_a)}$
of~\eqref{9.1:eq-arrest-factor-b}, and at most $f_{k_a,m}$. The room is given by the domain and
margin function of Definition~\ref{9.1:def-leftover-margin}
(see~\eqref{9.1:eq-leftover-margin-a}).
\end{enumerate}
\item[(iv)] The pair of fractions $(1/4,7/8)$ depends on no generation: no count of steps enters it.
\end{enumerate}
\end{proposition}

\begin{proof}
Theorem~\ref{9.1:thm-tube-theorem}, Lemma~\ref{9.1:lem-leftover-schedule} and
Corollary~\ref{9.1:cor-tube-landing} are proved as implications from their named hypotheses,
which are assumed here. Every appeal to them below is made under these hypotheses.

\emph{Parts (i) and (ii).} Both restate Theorem~\ref{9.1:thm-tube-theorem} and the seed bound of
the tube construction.

Theorem~\ref{9.1:thm-tube-theorem} takes as its point $P$ any point on the pair tube
$\mathfrak{P}_{\mathsf{c}}$ of a shape at any site that lies over a seed of the output $\Phi$
with slice at fraction $7/8$. The hypotheses of that theorem impose on this seed no condition
other than those of the definition of the tube, namely that the slice lies in the tube and that
the size clause $|B|\le r_c$ of that definition holds, its constant $r_c$ being independent of
$g$. This gives (i). The last sentence of (i) is the special case of an $\mathbf{R}$-born output:
the hypotheses of Theorem~\ref{9.1:thm-tube-theorem} do not restrict entries on $\mathfrak{O}$
to points of $\mathsf{T}_N$.

For (ii), the seed-delivery statement, part (a), hands the seed to the read with coefficient~$1$.
The proof of Theorem~\ref{9.1:thm-tube-theorem} proceeds additively in the majorants of the
increments. In it the constants $B_{\natural}$, $C_T^{w}$, $c_1$, $B^{\mathsf{s}}$,
$\mathsf{r}_1$ and $C_4$ occur only as factors of increments. These products are absorbed either
by $c_M$, $K_R^{\sharp}$ and the link weight $\mathsf{w}_L(p)$, through dilations, or by the two
ceilings on $\gamma$ named in (ii), all of which are hypotheses of
Theorem~\ref{9.1:thm-tube-theorem}. Since none of these constants multiplies the seed, none of
them restricts the set of points at which the seed is delivered.

\emph{Part (iii).} Each load bound and each room coefficient is supplied by the statement named
for its case.
\begin{itemize}
\item[(a)] For a $T$-born leaf of level $k$ read at $\mathbf{R}_k$ on $\mathfrak{O}$, both bounds
are part (c) of the accounting of thresholds on the regular region $\mathfrak{O}$.
\item[(b)] On $Z^{\mathrm{on}}$, the load bound is clause (ii) of the tube statement on runs,
together with the entry-and-link bound and the schedule bound against the printed table of a
leaf. The room coefficient is the same as in (a).
\item[(c)] For stage $n+1$ on $\mathbb{C}^{(i)}$ with $i\le n$, the reading of the unfolded terms
against $\mathsf{T}^{\flat}_n\le\mathsf{T}^{\flat}_i$, with rooms $\mathrm{rm}_{\nu}$ and
increments $\mathsf{g}_{\nu}$, is the tube statement on runs together with part (b) of the
accounting of thresholds on $\mathfrak{O}$.
\item[(d)] For an $\mathbf{R}$-born leaf of level $k$ read at $T_k$, the load bound and the value
of the room coefficient are the accounting of thresholds at the $T$-site together with the
schedule bound against the printed table of a leaf.
\item[(e)] For a left-over term read at $T_{k_a}$, the load is bounded by the diminished layer
factor of Definition~\ref{9.1:def-arrest-factor}. By
Lemma~\ref{9.1:lem-leftover-schedule}, in the form~\eqref{9.1:eq-leftover-schedule-a}, a
later read satisfies $\mathsf{T}\le(\mathsf{c}_F^{+}-\varrho_F)\mathsf{u}_F$ at $(=)$-points,
with $\varrho_F$ as in Definition~\ref{9.1:def-leftover-margin}.
\end{itemize}
Along the pair path from the seed to the read, parts (a) and (b) of
Corollary~\ref{9.1:cor-tube-landing} apply, with the constant and ceiling
of~\eqref{9.1:eq-tube-landing-a}. By (a), the read value
$\mathsf{v}=\mathrm{rd}^{\mathsf{A}}\mathbf{U}$ lies in $\mathfrak{T}^{\mathsf{A}}[0]$ and
$\mathsf{v}\oplus c^{\mathsf{B}}$ lies in $\mathfrak{T}^{\mathsf{B}}[1/8]$. By (b), for every real
phase $\eta$ with $|\eta|\le(8C_{\mathrm{TL}}^{\sharp}\vartheta^{\sharp})^{-1}$, where
$\vartheta^{\sharp}$ is as in Corollary~\ref{9.1:cor-tube-landing}, the rotated values
$e^{i\eta\mathfrak{d}^{\mathrm{net}}}(\mathsf{v}\oplus c^{\mathsf{B}})$ lie in
$\mathfrak{T}^{\mathsf{B}}[1/4]$.

Combined with the case list, this shows that the loads at every hand-over land in
$\mathfrak{T}_F[1/4]$, counted in the room of $F$, while by (i) the delivered seeds are all of
$\mathfrak{T}_{\Phi}[7/8]$. This is the asserted containment.

\emph{Part (iv).} None of the constants and none of the ceilings on $\gamma$ on which
Theorem~\ref{9.1:thm-tube-theorem} and the statements named in (iii) depend reads a count of
steps. The fractions $1/4$ and $7/8$ are therefore the same at every generation.
\end{proof}

\begin{definition}[Lists, loads, uses and the dialled production]\label{9.2:not-lists-loads-uses}
Throughout this section we work on one lattice with one scaffold, in the setting of the one-lattice
Lipschitz step.

\emph{Lists and loads.} Let $\mathfrak{s}$ and $\tilde{\mathfrak{s}}$ be two lists in the admissible
class $\mathfrak{K}_k$ at level $k$, and put $\delta := \tilde{\mathfrak{s}} - \mathfrak{s}$. For each
zone $n \le k$ (zones in the sense of the setting of the one-lattice Lipschitz step),
$[\![\delta]\!]_n$ denotes the load of $\delta$ at zone $n$. The load is taken in the form fixed in
the extension proposition; the features of that form which go beyond the basic definition of the load
play no role in what follows. With
the load constants $E_0$ and $S_\varsigma$ we put
\begin{equation*}
E_1 := 2(S_\varsigma + 1)E_0 .
\end{equation*}
Since the load constants satisfy $S_\varsigma \ge 0$ and $E_0 \ge 0$, we have $E_1 \ge 2E_0$; and the
load bound of the setting of the one-lattice Lipschitz step gives $[\![\delta]\!]_n \le E_1$ for
differences of lists in $\mathfrak{K}_k$.

\emph{The age ratio and the weighted load sum.} Let $c_a \ge \log 2$ be the constant fixed after $L$
and before $\gamma$ in the order of choice, and put $q := e^{-c_a}$. Since
$c_a \ge \log 2$, we have $q \le e^{-\log 2} = \tfrac12$. The weighted load sum of $\delta$ at level $k$ is
\begin{equation}\label{9.2:eq-lists-loads-uses-1}
\Sigma^q_k[\delta] := \sum_{n \le k} (k-n+1)\, q^{(k-n-1)_+}\, [\![\delta]\!]_n ,
\end{equation}
where $(k-n-1)_+ = \max\{k-n-1, 0\}$.

\emph{Derived objects.} The derived objects are the following:
\begin{itemize}
\item the boundary terms $\mathbf{B}^{(j)}$ produced by the operation $\mathbf{T}$, the frozen terms
$\mathbf{C}$ of arrested components, the terms $V(Y)$, one for each region $Y$, produced by the
production formulas, and the operations $\mathbf{T}_{k'}$ and $\mathbf{T}'_{k'}(X)$;
\item the activities of \cite[(1.91)]{B-CMP122b};
\item the rest terms $R'$ and $R''$ of the second group and the boundary terms $\mathbf{B}'$ of the
first group of \cite[(1.98)]{B-CMP122b}, and the outputs of the operation $\mathbf{T}$.
\end{itemize}
They are given by the production formulas of the one-lattice Lipschitz step.
In these formulas, after each step $j$, the terms with $X \subset \Lambda_j$, where $\Lambda_j$ is the
small-field region after step $j$ in Ba{\l}aban's construction, are entries of the list.
Every derived object is a composition of holomorphic maps of the values of the \emph{units} of the
production, namely the entries of the first, second, third and fifth rows and of the row marked
$\partial$ of the table of units of the one-lattice Lipschitz step.

\emph{Uses.} A \emph{use} is a triple $u = (k', n, v)$ consisting of a level $k'$, a zone $n$ and a
value index $v$, so that $f_v$ is a value of a unit. The zone of a value is defined as follows: if
$\mathfrak{b}$ is the argument at which the value $f_v(\mathfrak{b})$ is taken, the zone of
$f_v(\mathfrak{b})$ is the zone of the anchor in the determining set of $\mathfrak{b}$, anchors and
determining sets being those of the setting of the one-lattice Lipschitz step.

\emph{The dialled production.} Let $\lambda = (\lambda_{n,k'})$ be complex dial variables, one for each
zone $n$ and level $k'$. We write $f_v[\mathfrak{s}]$ for the value $f_v$ computed from the
list $\mathfrak{s}$, and $f_v[\delta]$ for the corresponding value attached to the difference $\delta$,
in the notation of the one-lattice Lipschitz step. The \emph{dialled production} is the production in
which, in every use $(k', n, v)$, the value $f_v[\mathfrak{s}]$ is replaced by
\begin{equation}\label{9.2:eq-lists-loads-uses-2}
f_v[\mathfrak{s}] + \lambda_{n,k'}\, f_v[\delta] .
\end{equation}
The dials range over the polydisc
\begin{equation}\label{9.2:eq-lists-loads-uses-3}
\mathbb{D}_\delta := \bigl\{ \lambda \,:\, |\lambda_{n,k'}|\, [\![\delta]\!]_n \le E_1
\ \text{for all } n, k' \bigr\} .
\end{equation}

\emph{Stored derived objects.} The fourth row of the table of units, omitted from the list of units
above, consists of the sup-norms of derived objects of lower level or lower stage; this input is
supplied by the induction hypothesis of the one-lattice Lipschitz step on the level and the stage. The
derived objects of this kind that are stored from earlier steps are the following.
\begin{itemize}
\item The boundary terms $\mathbf{B}^{(m)}(X_F)$ produced by the operation $\mathbf{T}$, in the
notation of \cite[(2.40)--(2.42)]{B-CMP119}.
\item The boundary terms $\mathbf{B}'^{(m)}(X')$ of the first group of \cite[(1.98)]{B-CMP122b},
produced by the operation $\mathbf{R}$.
\item The frozen terms $\mathbf{C}$ of arrested components in \cite{B-CMP122a}.
\end{itemize}
These stored derived objects are recomputed, at every later level and for every label (the index
carried by a term of Ba{\l}aban's expansion), from the dialled values
\eqref{9.2:eq-lists-loads-uses-2} of the units, and they enter the later formulas through their
values only.
\end{definition}

\begin{definition}[The per-use polydisc and switching off a set of uses]
\label{9.2:def-per-use-polydisc}
We use the lists $\mathfrak{s},\tilde{\mathfrak{s}}\in\mathfrak{K}_k$, their difference
$\delta=\tilde{\mathfrak{s}}-\mathfrak{s}$, the loads $[\![\delta]\!]_n$, the constant $E_1$,
the uses $(k',n,v)$, the dialled production and the polydisc $\mathbb{D}_\delta$ of
Definition~\ref{9.2:not-lists-loads-uses}.

\emph{Use-copies.} The dialled uses are indexed individually. A \emph{use-copy} $\upsilon$ is one
occurrence of a use $(k'(\upsilon),n(\upsilon),v)$ inside the fully expanded production formula
of a given object. Two occurrences of the same unit value in two distinct terms of the formula
(for instance in two stored derived objects of lower level) are two distinct use-copies.
Each use-copy $\upsilon$ carries its own dial variable $\lambda_\upsilon\in\mathbb{C}$.
In the use $\upsilon$ this variable takes the place of the variable $\lambda_{n(\upsilon),k'(\upsilon)}$
of the dialled production.

\emph{The per-use polydisc.} We put
\begin{equation}\label{9.2:eq-per-use-polydisc-1}
\mathbb{D}^{\mathrm{use}}_\delta
:=\prod_{\upsilon}\bigl\{\lambda_\upsilon\in\mathbb{C}:
|\lambda_\upsilon|\,[\![\delta]\!]_{n(\upsilon)}\le E_1\bigr\}.
\end{equation}
The polydisc $\mathbb{D}_\delta$ is identified with the diagonal of
$\mathbb{D}^{\mathrm{use}}_\delta$, that is, with the image of $\mathbb{D}_\delta$ under the map
$\lambda\mapsto(\lambda_{n(\upsilon),k'(\upsilon)})_\upsilon$; it consists of the points of
$\mathbb{D}^{\mathrm{use}}_\delta$ whose coordinate $\lambda_\upsilon$ depends on $\upsilon$ only
through the pair $(n(\upsilon),k'(\upsilon))$.

\emph{Validity on the per-use polydisc.} The estimates of the proposition bounding the dialled
production on $\mathbb{D}_\delta$ hold verbatim on $\mathbb{D}^{\mathrm{use}}_\delta$; this is the
extension proposition (stated; proof incomplete at the step marked $(\ast)$).

The mechanism of that proposition is the following: the proof of these estimates charges each
use separately. The units entering the production are affine in the list. The dialled part of
each such unit in a use-copy $\upsilon$ has, per cell of the zone $n=n(\upsilon)$, weighted size
at most $C_L^{\sharp}\,|\lambda_\upsilon|\,[\![\delta]\!]_{n(\upsilon)}\le C_L'$. The remaining
entries of the table of estimates in that proof read only the sizes of their inputs. Sup-norms
of derived objects of lower level or stage are bounded by the induction hypothesis. Thus giving
two use-copies of the same pair $(n,k')$ different values inside their common radius changes no
entry of that table of estimates.

\emph{Switching off a set of uses.} Let $S$ be a set of use-copies and let
$c\in\mathbb{D}^{\mathrm{use}}_\delta$. We define $c\lceil_S$ to be the point with coordinates
\begin{equation}\label{9.2:eq-per-use-polydisc-2}
(c\lceil_S)_\upsilon:=
\begin{cases}
0, & \upsilon\in S,\\
c_\upsilon, & \upsilon\notin S.
\end{cases}
\end{equation}
This point again lies in $\mathbb{D}^{\mathrm{use}}_\delta$. Indeed, for $\upsilon\in S$ the
defining condition in \eqref{9.2:eq-per-use-polydisc-1} reads
$0\cdot[\![\delta]\!]_{n(\upsilon)}=0\le E_1$. For $\upsilon\notin S$ the coordinate is that
of $c$. For an object $\psi$ we put
\begin{equation*}
\psi\lceil_S:=\psi(c\lceil_S).
\end{equation*}

\emph{Diagonal use of the disc estimate.} The disc estimate of the one-lattice Lipschitz step is
applied, as for $\mathbb{D}_\delta$, to the diagonal variables $\lambda_z$ with $z=(n,k')$. For
each zone $n\le k$ there are $k-n+1$ values of $k'$. In the remainder of this section every
majorant and every structural statement is asserted, and holds, at every point of
$\mathbb{D}^{\mathrm{use}}_\delta$, not only on the diagonal $\mathbb{D}_\delta$.
\end{definition}

\begin{definition}[The age factor and the leaf majorants, taken by statement]
\label{9.2:not-age-factor-majorants}
Throughout, $\mathrm{e}$ denotes Euler's number, and $q=\mathrm{e}^{-c_a}\le\frac12$, the lists $\mathfrak{s},
\tilde{\mathfrak{s}}\in\mathfrak{K}_k$, their difference $\delta$, the loads $[\![\delta]\!]_n$, the
constants $E_0\le\frac12 E_1$ and the smallness allowance on the dialled production are those of
Notation~\ref{9.2:not-lists-loads-uses}. The diagonal polydisc $\mathbb{D}_\delta$ and the per-use
polydisc $\mathbb{D}^{\mathrm{use}}_\delta$ are those of Definition~\ref{9.2:def-per-use-polydisc}.
The following statements are used in the form stated. The statements derived from them here (the
tail-summed bound in (b), the value of $c_1$ for plain terms in (d), the value of
$\mathsf{c}'$ in (e), the consequence drawn in (f), the inclusion in (i) and the lower bound in (k))
are proved below; the two facts on which part (g2) of (g) rests are recorded there as well.

\begin{enumerate}
\item[(a)] \emph{Schwarz-type estimates.} For $\sigma>0$ let $D(0,\sigma)\subset\mathbb{C}$ be the
disc of radius $\sigma$ about $0$, and let $\mathfrak{B}>0$. If $f$ is holomorphic on
$D(0,\sigma)$ with $\sigma\ge1$ and $|f|\le\mathfrak{B}$ there, then $|f(1)|\le\mathfrak{B}/\sigma$
when $f(0)=0$, and $|f(1)-f(0)|\le2\mathfrak{B}/\sigma$ in general. If $G$ is holomorphic on
$D(0,\sigma_1)\times D(0,\sigma_2)$ with $\sigma_1,\sigma_2\ge2$ and $|G|\le\mathfrak{B}$ there, then
\begin{equation*}
|G(1,1)-G(1,0)-G(0,1)+G(0,0)|\le\frac{4\mathfrak{B}}{\sigma_1\sigma_2}.
\end{equation*}
These are the Schwarz-type lemmas stated and proved with the one-lattice Lipschitz step, used here
in the form stated.

\item[(b)] \emph{Age.} Let $X$ be a component of the large-field region read at level $k$, for any
history. For an admissible sequence $\omega$ in the sum \cite[(1.71)]{B-CMP122b} put
\begin{equation*}
j(X)(\omega):=\min\{i:\ Z_i\cap X\ne\emptyset\ \text{in}\ \omega\},
\end{equation*}
the index of a first large-field region contained in $X$ in the sense of \cite{B-CMP122b}, and call
$a:=k-j(X)$ the \emph{age} of $\omega$ in $X$. Let $\mathbf{T}'_{k,a}(X)$ be the part of the sum
over sequences in \cite[(1.71)]{B-CMP122b} with age $a$ (for an arbitrary component the sequence
belongs to the label, the datum attached to the component that carries its history, and $a$ is the
label's age), and put
$\mathbf{T}'_{k,\ge a}(X):=\sum_{a'\ge a}\mathbf{T}'_{k,a'}(X)$.
Geometric forgetting of old regions states: under the smallness allowance, on the product of
$\mathbb{D}_\delta$ with the tube of configurations and the auxiliary domain $\mathcal{W}_k$ of the
one-lattice Lipschitz step, for every component $X$ (any history), every $a\ge0$ and every
bounded $F$,
\begin{equation*}
\begin{split}
|\mathbf{T}'_{k,a}(X)F|
&\le\mathrm{e}\,\exp\bigl(-2p_0(g_{k-a})\bigr)\,q^{a+1}\sup|F|\\
&\le\mathrm{e}\,\exp\bigl(-2(1+\beta_0)^{-1}p_0(g_k)\bigr)\,q^{a+1}\sup|F|,
\end{split}
\end{equation*}
with $p_0(\cdot)$ the degree function. Its proof is carried out with the entropy constant
$C_{\mathfrak{d}}$ of the theorem it rests on equal to $\mathsf{c}$, under the condition
$p_0\ge6r+1$ on the exponent of the degree function, and with the two-class statement as a
hypothesis, which enters through the events of the
second class. The
induction in that proof holds at any value of the entropy constant, its degree inequalities being
followed by a ceiling on $\gamma$ chosen after that constant; in particular it holds with
$\mathsf{c}$ replaced by the constant $\mathsf{c}'$ of (e), under a ceiling on $\gamma$ chosen
after $\mathsf{c}'$. The bound is obtained by majorising the summands and summing, the summations
over the admissible sequences being replaced by the factors $\exp\bigl(O(1)(MR_j)^{-d}|Z_j|\bigr)$
of \cite{B-CMP122b}, with $R_j$ the radius of that paper; hence the part of the sum over any
sub-family of the sequences of age $a$ obeys the same majorant. We put
\begin{equation}\label{9.2:eq-age-factor-majorants-a}
\mathfrak{E}_k:=\exp\bigl(-2(1+\beta_0)^{-1}p_0(g_k)\bigr),\qquad
|\mathbf{T}'_{k,\ge a}(X)F|\le\mathrm{e}\,\mathfrak{E}_k\,q^{a}\sup|F|,
\end{equation}
the inequality holding for the part of $\mathbf{T}'_{k,\ge a}(X)$ over every sub-family of
sequences of age at least $a$.

\item[(c)] \emph{Holomorphy of the dialled production.} The holomorphy proposition for the dialled
production states that every derived object of the dialled production of
Notation~\ref{9.2:not-lists-loads-uses} is holomorphic on $\mathbb{D}_\delta$ and, by
Definition~\ref{9.2:def-per-use-polydisc}, on $\mathbb{D}^{\mathrm{use}}_\delta$, and satisfies
there, under the smallness allowance, the bounds of the correlation theorem:
\cite[(2.42)]{B-CMP119} for $\mathbf{B}^{(j)}$; the weighted form of
\cite[(1.69)]{B-CMP122b} with the term $(\mathsf{c}-c_a)\sum_j|\Gamma^0_j\cap Y|$, $Y$ being the
localisation domain of \cite[(1.69)]{B-CMP122b};
\begin{equation*}
|\mathbf{T}'_k(X)F|\le\mathrm{e}\,\sup_\omega\Pi^{[\mathsf{c}-c_a]}_\omega\,\sup|F|,
\end{equation*}
where $\Pi^{[\mathsf{c}-c_a]}_\omega$ is the product of factors which the correlation theorem's
bound on $\mathbf{T}'_k(X)$ assigns to the admissible sequence $\omega$, taken at the constant
$\mathsf{c}-c_a$;
\cite[(1.97)]{B-CMP122b} with its coefficient $c_1$ replaced by $2c_1$; and
$|R'^{(k)}(X)|\le\mathrm{e}^{-p_0(g_k)}\mathrm{e}^{-\kappa d_k(X)}$. In the estimates
\cite[(1.92)--(1.100)]{B-CMP122b} the dialled objects enter multiplicatively: the terms $V$
through their sup-bounds and $\mathbf{T}'$ through \cite[(1.89)]{B-CMP122b} times $\mathrm{e}$.
The replacement of the summand $+2$ by $+3$ in the corresponding estimate of \cite{B-CMP122b},
which is already part of the smallness allowance, covers the factor $2\mathrm{e}$ of one splitting
in the auxiliary dial $\mu$, one splitting at a time; and the factor $\mathrm{e}^{-p_0/2}$ left
over at \cite[(1.96)]{B-CMP122b}, $p_0$ denoting there the value of the degree function
$p_0(\cdot)$ in that estimate, absorbs the factor $2\,O(1)$.

\item[(d)] \emph{Leaves and their majorants.} By the same argument, for a term of the first group
of \cite[(1.98)]{B-CMP122b} the bound \cite[(1.99)]{B-CMP122b},
\begin{equation*}
|R'^{(k)}(X,(U,J))|\le O(1)\,c_1\exp\Bigl(-\bigl(1+\tfrac12\beta\bigr)\kappa\,
d_{k,\bigcup_iY_i}(X)\Bigr),
\end{equation*}
holds on the dialled polydisc with $c_1$ replaced by $2c_1$, $(U,J)$ being the arguments of
$R'^{(k)}$ in \cite[(1.99)]{B-CMP122b}; here $c_0$, $c_1$ and $\beta$ are the
coefficients and the exponent of \cite[(1.99)]{B-CMP122b}, used in that sense in (c), (d), (j) and
the proof below only. Consider such a first-group
term, with domain $X'$, produced at level $j$ (the step called $k$ in \cite[(1.99)]{B-CMP122b}); we
say it is \emph{plain} if $X'\cap\bigcup_hY_h=\emptyset$, and a \emph{strip} term if $X'$ contains
one of the domains $Y_h$. Write $d^{\mathrm{rel}}_j(X')$ for the relative size $d_{k,\bigcup_iY_i}(X')$ of
\cite[(1.99)]{B-CMP122b}, read at the level $j$ of production, and put
\begin{equation}\label{9.2:eq-age-factor-majorants-b}
\mathfrak{B}'_{X'}:=C_{99}\,c_1(X')\,
\mathrm{e}^{-(1+\frac12\beta)\kappa\, d^{\mathrm{rel}}_j(X')},\qquad
\mathfrak{B}_X:=\mathrm{e}^{-p_0(g_k)}\,\mathrm{e}^{-\kappa d_k(X)},
\end{equation}
where $C_{99}$ is twice the absolute constant written $O(1)$ in \cite[(1.99)]{B-CMP122b}.
Here $\mathfrak{B}'_{X'}$ is the bound which the holomorphy proposition delivers for the term of
\cite[(1.98)]{B-CMP122b} with domain $X'$ on the polydisc enlarged by one disc in the auxiliary dial
$\mu$ (covered by the summand $+3$ of (c)) and one disc in the auxiliary dial $\nu$ (covered by
the enlarged allowance of (e)); $\mathfrak{B}_X$ is the bound for the second group,
\cite[(1.100)]{B-CMP122b}. For a plain term, $c_1(X')=\mathrm{e}^{-p_0(g_j)}$ and
$d^{\mathrm{rel}}_j(X')=d_j(X')$, the latter identification being taken with the majorant as
stated.

\item[(e)] \emph{The enlarged allowance.} The enlarged allowance is the smallness allowance with
$C_L'$ replaced by $2C_L'$ at both of its occurrences, in
$\mathsf{c}=C_{\sharp}O(1)+1+C_L'+c_a$ and in the factor $(C_{\sharp}+C_L'/E_0)$ of the activity
constants; that is, one more unit $C_L'$ of the per-cell allowance and of the bracket budget:
\begin{equation}\label{9.2:eq-age-factor-majorants-c}
\begin{gathered}
C_{\sharp}O(1)+1+2C_L'+c_a=\mathsf{c}',\qquad \mathsf{c}':=\mathsf{c}+C_L',\\
(C_{\sharp}+C_L'/E_0)\ \text{replaced by}\ (C_{\sharp}+2C_L'/E_0).
\end{gathered}
\end{equation}
The replacement of $+2$ by $+3$ of (c) is already part of the smallness allowance and covers one
splitting in $\mu$ at a time. Each clause of the enlarged allowance is a ceiling on $\gamma$ chosen
after $L$, $M$, $\kappa$, the constant $\kappa_0$ of the smallness allowance of
Notation~\ref{9.2:not-lists-loads-uses}, the exponents, $r$ and $c_a$; the resulting ceiling on
$\gamma$ is written $\gamma_L'$, and it satisfies $\gamma_L'\le\gamma_L$, where $\gamma_L$ is the
ceiling on $\gamma$ of the smallness allowance of Notation~\ref{9.2:not-lists-loads-uses}.

\item[(f)] \emph{Budget of stored families.} The contribution of a family of stored boundary terms
and frozen terms (the $\mathbf{B}$- and $\mathbf{C}$-terms, whose piles are counted at each term's
own level) to the budget of a site is linear in the family and is measured by its pile per native
cell, in the attribution of the pile lemma; the conversion from cubes to cells is the volume part of
the volume and filament bound. Precisely, if, for some number $\varpi\ge0$, the values of the
family have pile per native cell at most $\varpi C_L'$ on the reading domain, then the family adds
at most the first quantity in \eqref{9.2:eq-age-factor-majorants-d} per unit of
$\sum_l|\Gamma^{\natural}_l\cap Y|$ to the constant of the weighted form of
\cite[(1.69)]{B-CMP122b}, and at most the second to the site's activity:
\begin{equation}\label{9.2:eq-age-factor-majorants-d}
\varpi C_L',\qquad C_A\mathsf{w}_{\mathrm{site}}^{-1}\varpi C_L'.
\end{equation}
Hence a family with $\varpi\le1$ is covered by the enlarged allowance. This is the way in which
the dialled values of Notation~\ref{9.2:not-lists-loads-uses} enter: they add at most $C_L'$ per
unit of $\sum_j|\Gamma^0_j\cap Y|$, without a site-weight factor $\mathsf{w}_{\mathrm{site}}^{-1}$,
and the brackets of the
activity gain $C_A\mathsf{w}_{\mathrm{site}}^{-1}C_L'$.

\item[(g)] \emph{Anchoring of stored terms.} Let $Z_j:=\Lambda_j^c$, $Z_j^{\boxplus}$ and $\Omega_j$
denote the unprimed sets of the current label. (g1) For a component $C$ of the current $Z_j$,
$j<k$, one has $(C^{\boxplus})^{(1)}\subset Z_{j+1}$ (current), $(\cdot)^{(1)}$ denoting the
enlargement operation on regions of the notation of the correlation theorem. (g2) Every stored
boundary-type term
$B$ of level $j\le k$ still stored at the current level has a cube of its localisation domain $A_B$
in $(Z_j^{\boxplus})^{(1)}$, in the current label. Between the label at the production of a term and
the current label the inclusion $(Z_j^{\boxplus})^{(1)}\subset Z_{j+1}$ fails across a component
healed in between.

\item[(h)] \emph{Localisation.} An output with domain $X$ reads only terms $V(Y)$ with
$Y\subset X$: the domain of an object contains the domains of the objects and units it reads.

\item[(i)] \emph{Filament inclusion.} Let $X$ be a component of $Z_{k'}=\Lambda_{k'}^c$, components
being the classes of $M\check{R}_{k'}$-cubes of $\Lambda_{k'}^c$ under the equivalence relation
generated by adjacency of cubes. Then $X^{\boxplus}\setminus X\subset
\Lambda_{k'}$; hence every point of $X^{\boxplus}$ lying in $\Lambda_i^c$ for some $i\le k'$ lies in
$X$. Moreover $X^{\sim}\subset X^{\boxplus}$, the enlargement $X^{\sim}$ being by one
$M\check{R}$-layer \cite{B-CMP119}.

\item[(j)] \emph{Format of the densities, by statement.}
\begin{itemize}
\item[(j1)] The density after $k$ steps is a sum over admissible sequences (histories) of terms: in
\cite[(2.18)--(2.22)]{B-CMP119} the operation is
\begin{equation*}
\mathbf{T}_k(X)=\prod_j\mathbf{T}^{(j)}(Z_{j+1}\cap X)\,\mathbf{T}^{(\bar m)}(Z_{\bar m}\cap X),
\end{equation*}
$\bar m$ the last index of the sequence, the operation of level $j$ acting on $Z_{j+1}\cap X$; when
$\mathbf{T}^{(j)}$ is changed
by an $\mathbf{R}$-operation the general form \cite[(2.21)]{B-CMP119} is preserved and the
characteristic functions are changed. Each term carries its own stored objects,
$\mathbf{B}_k=\sum_j\mathbf{B}^{(j)}(U_k,A,\{S_i\})$ \cite[(2.40)--(2.41)]{B-CMP119}, $A$ and
$\{S_i\}$ being the arguments of \cite[(2.40)]{B-CMP119}, the constant of the bound there
depending also on $\{\Omega_j\}$, $\{\Lambda_j\}$. Inside $\mathbf{T}'_k(X)$ of
$\mathbf{R}_k$, \cite[(1.71)]{B-CMP122b}, the sum over admissible sequences
$\{\Omega_j\cap X,Z_j\cap X\}$ acts on everything built along them: in \cite{B-CMP122b} the terms
$V(Y,U_k)$ depend on the sequence restricted to the components of $Z$ contained in $Y$, the sum in
\cite[(1.71)]{B-CMP122b} acts also on the last exponential in \cite[(1.72)]{B-CMP122b}, and
$V(X^{\sim})$ is a sum of many terms, mainly boundary terms from all the renormalisation steps,
depending on the configuration \cite[(1.61)--(1.63)]{B-CMP122b} localised in $X^{\sim}$. The
reading of this format used in this chapter, and nothing finer, is the following. Write
\begin{equation}\label{9.2:eq-age-factor-majorants-e}
\mathbf{T}'_k(X)=\sum_\omega\mathbf{T}'_{k,\omega}(X).
\end{equation}
The $\omega$-summand contains exactly the stored objects produced along the history
$\omega{\upharpoonright}$, the label's history completed inside $X$ by $\omega$; a stored object
whose production required a component $C$ of $Z_e$ (respectively a point $y$ in zone $n$ at level
$k'$) is absent from the $\omega$-summands in which $C\notin Z_e^{\omega}$ (respectively
$y\notin\Gamma_n^{\omega}$ at level $k'$).
\item[(j2)] Anchoring: the statement (g).
\item[(j3)] Absorption: at a completed $\mathbf{R}_s$ with operated class $Z^{(s)}$, the first part
consists of the terms of the rest of the effective action whose localisation domains intersect the
region $Z$, the second part of the remaining terms, whose localisation domains are disjoint from $Z$
\cite{B-CMP122b}; the first part is expanded and resummed into the $V(Y)$ and thence into the
activities \cite[(1.90)--(1.91)]{B-CMP122b} and into $R'^{(s)}$, $\mathbf{B}'^{(s)}$, which contain
at least one large-field region connected with one of the $\mathbf{T}'$-operations.
\item[(j4)] Deletion: by \cite[(1.33)]{B-CMP122b}, $\Gamma_j=\Gamma''_j\cap Z^c$ for $j<k$ and
$\Gamma_k=\Gamma''_k\cup Z^{(k)}$; whole components of $\Lambda_i^c$ and of
$\Omega^{\natural c}_i$ with $i<k''$ lying inside $Z$ are deleted, $k''$ denoting the level at
which the deletion is performed; arrests of components raise indices inside the arrested
component and leave $\Lambda$ unchanged \cite{B-CMP122a}; and the new large-field region is
$Z_k\cup\bigcup_iY_i$ \cite{B-CMP122b}.
\item[(j5)] Growth and separation: each renormalisation step adds at least ten layers of
$MR_k$-cubes \cite[(i),(ii)]{B-CMP122a}; the domains $\Omega_j$, $\Lambda_j$ are unions of
$MR_j$-cubes and
the distance between their boundaries is at least $2MR_j$ \cite{B-CMP119};
by the separation statement for the native strata of the scaffold,
$\operatorname{dist}_\infty(\Omega^{\natural}_{i+1},\Lambda_i^c)\ge8M\check{R}_{i+1}\ell_{i+1}$;
and $\check{R}_{j+1}\le\check{R}_j\le L\check{R}_{j+1}$, $\log x_j\le\check{R}_j$, under the
one-sided hypothesis on the scaffold radii, closed under passage to prefixes, which holds on the
scaffold used here.
\item[(j6)] First group and classes, as in \cite{B-CMP122b}: regions of the first class $X_i$ are
healed, regions of the second class $Y_i$ are kept with the complementary characteristic function
written $\chi^c$ in \cite{B-CMP122b};
$X'_0$ is the union of the regions $X_j$ of the first class with a family of domains $Y$ specified
in \cite{B-CMP122b}, and the clusters are the components of
$X'_0\setminus\bigcup_iY_i$; the coefficients $c_0$, $c_1$ are those recalled in (d); a term of
\cite[(1.98)]{B-CMP122b} is of the first group if and only if its domain $X'$ meets
$Z^{\sim}\cup\bigcup_iY_i^{\sim}$ of the new label, and these are the new boundary terms
$\mathbf{B}'^{(k)}(X)$.
\end{itemize}

\item[(k)] \emph{Lower flow bound.} The lower member of the reference window
(Definition~\ref{1.2:def-flow-constants}) is used in the following form, taken by statement, in
which $\lambda$ is the rate and $b_+$ the first-passage overshoot of the flow: for $0\le i<j\le K$,
\begin{equation*}
\lambda(j-i)-2c_2\le x_i-x_j,\qquad x_K\in\,]g^{-2},g^{-2}+b_+].
\end{equation*}
Hence, for every $e\le K$,
\begin{equation*}
x_e\ge g^{-2}-2c_2+\lambda(K-e),
\end{equation*}
and the degree function $p_0(g')=A_0(\log x')^{p_0}$, $x'=g'^{-2}$, is increasing in $x'$ on
$x'\ge1$.
\end{enumerate}
\end{definition}

\begin{proof}
\emph{The tail-summed bound in (b).} Let $a\ge0$ and fix a sub-family of sequences of age at least
$a$. Its part of $\mathbf{T}'_{k,\ge a}(X)F$ is the sum over $a'\ge a$ of the parts over its
sequences of age $a'$. Each of these is the part of $\mathbf{T}'_{k,a'}(X)F$ over a sub-family of
the sequences of age $a'$, so by the last sentence of (b) before \eqref{9.2:eq-age-factor-majorants-a}
it obeys the second majorant of geometric forgetting of old regions with $a'$ in place of $a$, that
is, it is at most $\mathrm{e}\,\mathfrak{E}_k\,q^{a'+1}\sup|F|$; the factor $\mathfrak{E}_k$ does
not depend on $a'$. Since $0<q\le\frac12<1$, the series converges and
\begin{equation*}
\sum_{a'\ge a}q^{a'+1}=\frac{q^{a+1}}{1-q}\le q^{a},
\end{equation*}
the last inequality being equivalent to $q\le1-q$, that is, to $q\le\frac12$. Summing the bounds
gives \eqref{9.2:eq-age-factor-majorants-a}.

\emph{Plain terms in (d).} The coefficient $c_1$ of \cite[(1.99)]{B-CMP122b}
equals $\exp(-p_0(g_k))$ if $X'$ does not intersect $\bigcup_hY_h$, and $\alpha^{1/3}$ in the
remaining cases, with $\alpha$ the parameter so denoted in \cite[(1.99)]{B-CMP122b}; the case with
the coefficient $c_0=\alpha^{1/3}$ arises only when the domain $X'$
contains one of the domains $Y_h$, which is the case of strip terms. A plain term has
$X'\cap\bigcup_hY_h=\emptyset$, so the first alternative applies; reading the step of
\cite[(1.99)]{B-CMP122b} as the level $j$ of production gives $c_1(X')=\mathrm{e}^{-p_0(g_j)}$.

\emph{The value of $\mathsf{c}'$ in (e).} Replacing $C_L'$ by $2C_L'$ in
$\mathsf{c}=C_{\sharp}O(1)+1+C_L'+c_a$ gives
$C_{\sharp}O(1)+1+2C_L'+c_a=\mathsf{c}+C_L'=\mathsf{c}'$, which is the first line of
\eqref{9.2:eq-age-factor-majorants-c}; the second line is the same replacement in the factor of the
activity constants.

\emph{The consequence in (f).} If $\varpi\le1$, then by \eqref{9.2:eq-age-factor-majorants-d} the
family adds at most $C_L'$ per unit of $\sum_l|\Gamma^{\natural}_l\cap Y|$ to the constant of the
weighted form of \cite[(1.69)]{B-CMP122b} and at most $C_A\mathsf{w}_{\mathrm{site}}^{-1}C_L'$ to the
site's activity. By \eqref{9.2:eq-age-factor-majorants-c} the enlarged allowance provides exactly
one further unit $C_L'$ in $\mathsf{c}$ and one further unit $C_L'$ in the factor of the activity
constants, entering in the same way as the dialled values recalled at the end of (f). Hence the
family is covered by the enlarged allowance.

\emph{The facts underlying part (g2) of (g).} Part (g2) is taken by statement; it rests on the
following two facts. At its production, a stored term $B=\mathbf{B}'^{(j)}$ born of an
$\mathbf{R}$-operation meets the large-field region of the label created by that operation, the set
written $Z_k^{\sim}\cup\bigcup_iY_i^{\sim}$ in \cite{B-CMP122b}, $k$ there being the step of that
operation, the enlargement $\sim$ there being $\boxplus$. At each later completed
$\mathbf{R}$-operation with operated region $Z'$, if $B$ remains stored it belongs to the second
part, so that $A_B\cap Z'=\emptyset$ by (j3); a term meeting $Z'$ belongs to the first part and is
no longer stored after that operation.

\emph{The inclusion (i).} Let $Q\subset X^{\boxplus}\setminus X$ be an $M\check{R}_{k'}$-cube. Since
$X^{\boxplus}$ is the enlargement of $X$ by one layer, $Q$ is adjacent to $X$. If $Q\subset
\Lambda_{k'}^c$, then $Q$ would lie in the class of $X$, that is, in $X$, which is not the
case; hence $Q\subset\Lambda_{k'}$, and $X^{\boxplus}\setminus X\subset\Lambda_{k'}$. Next, for
$i\le k'$ one has \cite{B-CMP119}
\begin{equation*}
\Lambda_i\supset\Omega_{i+1}\supset\Lambda_{i+1}\supset\dots\supset\Lambda_{k'},
\end{equation*}
hence $\Lambda_i^c\subset\Lambda_{k'}^c$. A point of $X^{\boxplus}$
lying in $\Lambda_i^c$ with $i\le k'$ therefore lies in $\Lambda_{k'}^c$, so it is not in
$X^{\boxplus}\setminus X\subset\Lambda_{k'}$, and it lies in $X$.

\emph{The lower bound (k).} Let $e<K$. The first inequality of (k) with $i=e$, $j=K$ and the
inclusion $x_K\in\,]g^{-2},g^{-2}+b_+]$ give
\begin{equation*}
x_e\ge x_K+\lambda(K-e)-2c_2>g^{-2}-2c_2+\lambda(K-e).
\end{equation*}
For $e=K$ the claim reads $x_K\ge g^{-2}-2c_2$, which follows from $x_K>g^{-2}$ and $c_2\ge0$.
Finally
$x'\mapsto\log x'$ is increasing and non-negative on $x'\ge1$, and $t\mapsto A_0t^{p_0}$ is
increasing on $t\ge0$ since $A_0>0$ and $p_0>0$; hence $p_0(g')=A_0(\log x')^{p_0}$ is increasing
in $x'$ on $x'\ge1$.
\end{proof}

\begin{definition}[Object occurrences, witnessing components and absorption at a healing]
\label{9.2:def-witnessing-components}
Fix a history and a time $\mathsf{t}$, that is, a level, read between two half-steps. For a level
$j$ write $Z_j=\Lambda_j^{c}$ for the large-field region of level $j$ of the history, the
complement of the domain $\Lambda_j$ in Ba{\l}aban's sequence of domains
$\Omega_j\supset\Lambda_j\supset\Omega_{j+1}$ \cite{B-CMP119}. For a region $X$, $X^{\boxplus}$
and $X^{\sim}$ are Ba{\l}aban's enlargements of $X$, and $Y^{(1)}$ is the enlargement of a set
$Y$ used in the anchoring statement for stored boundary terms. The phrases \emph{alive at
$\mathsf{t}$}, \emph{healed by $\mathbf{R}_s$} (said of a component of the first class of
$\mathbf{R}_s$), \emph{second-class} and \emph{arrested} are used in the sense of the absorption
and deletion rules of the fixed history. For the internal sum over admissible sequences $\omega$
of the operated component in $\mathbf{T}'_k(X)$, they are read in the history
$\omega{\upharpoonright}$, which completes the fixed history inside $X$ by $\omega$.

\begin{enumerate}
\item[(a)] \emph{Object occurrences.} An \emph{object occurrence} $\psi$ is a stored derived
object of some level $e(\psi)$, with localisation domain $X_\psi$. It is of one of three kinds:
\begin{itemize}
\item a $\mathbf{T}$-born boundary term $\mathbf{B}^{(e)}(X_\psi)$;
\item an $\mathbf{R}$-born boundary term of the first group $\mathbf{B}'^{(e)}(X_\psi)$;
\item a frozen term $\mathbf{C}$.
\end{itemize}
The occurrence is either the stored term itself, or the copy of its formula inside a later
object that read its value. The domain of an object contains the domains of the objects it
reads (localisation of reading). Hence nested occurrences have nested domains.

\item[(b)] \emph{Witnessing components.} The witnessing components of $\psi$, with $e=e(\psi)$,
are defined according to its kind.
\begin{itemize}
\item For an $\mathbf{R}$-born boundary term of level $e$: the components $C$ of $Z_e$, in the
history at the birth of $\psi$ (after $\mathbf{R}_e$), with
$(C^{\boxplus})^{(1)}\cap X_\psi\neq\emptyset$. By part~(ii) of the anchoring statement, in the
case of a term at its birth, there is at least one.
\item For a $\mathbf{T}$-born $\mathbf{B}^{(e)}(X_\psi)$: the components $C$ of $Z_e$ with
$(C^{\boxplus})^{(1)}\cap X_\psi\neq\emptyset$. There is at least one: the domains of the
boundary terms $\mathbf{B}^{(j)}$ satisfy $X\cap Z_j^{\sim}\neq\emptyset$ by
\cite[(2.41)]{B-CMP119}, so $X_\psi\cap Z_e^{\sim}\neq\emptyset$; since
$X^{\sim}\subset X^{\boxplus}$, the domain $X_\psi$ meets $Z_e^{\boxplus}$, and hence meets
$C^{\boxplus}\subset(C^{\boxplus})^{(1)}$ for some component $C$ of $Z_e$.
\item For a frozen term $\mathbf{C}$: its arrested component.
\end{itemize}
We call $\psi$ \emph{live-witnessed at $\mathsf{t}$} if some witnessing component of $\psi$ is
alive at $\mathsf{t}$, that is, has not been deleted by a healing at any time
$\le\mathsf{t}$.

\item[(c)] \emph{Geometric witnesses.} Let $\upsilon$ be a use-copy (one use of a stored value
by a later object, carrying its own dial variable) of class $z=(n,k')$, the class being indexed
by a zone $n$ and a level $k'$, with $n<k'$, and let $y_\upsilon$ be its anchor point, the point
at which the use is anchored. Since $\Lambda_{n+1}\subset\Omega_{n+1}$, at level $k'$ in its
history
\begin{equation}\label{9.2:eq-witnessing-components-1}
y_\upsilon\in\Omega_n\setminus\Omega_{n+1}\subset\Lambda_{n+1}^{c}=Z_{n+1}.
\end{equation}
Let $C_\upsilon$ be the component of $\Lambda_{n+1}^{c}$ containing $y_\upsilon$. If $C_\upsilon$
is alive at $\mathsf{t}$, we call $\upsilon$ \emph{geometrically witnessed} at $\mathsf{t}$, with
witness level $n+1$.

Components of $\Lambda^{c}$ are deleted only by healings, and an arrest does not change the sets
$\Lambda$. Consequently, for a use-copy anchored inside a component of $\Lambda_{k'}^{c}$ alive
at $k'$ (untouched, second-class or arrested), the geometric witness is the component
$C_\upsilon$ of $Z_{n+1}=\Lambda_{n+1}^{c}$, read in the sets $\Lambda$, which arrests leave
unchanged. Use-copies of the top zone ($n=k'$) have no geometric witness.
\end{enumerate}

With these definitions the following three facts hold.

\emph{Fact~A.} If $\psi$ is a stored term itself, alive at $\mathsf{t}$, then $\psi$ is
live-witnessed at $\mathsf{t}$.

\emph{Fact~B.} Let $C$ be a witnessing component of $\psi$ healed by $\mathbf{R}_s$ with
$s>e(\psi)$, and let $\Psi$ be any stored object containing an occurrence of $\psi$, with
localisation domain $X_\Psi$. Then, $Z^{(s)}$ denoting the operated class of $\mathbf{R}_s$,
\begin{equation}\label{9.2:eq-witnessing-components-a}
X_\Psi\cap Z^{(s)}\neq\emptyset;
\end{equation}
consequently $\Psi$ belongs to the first part at $\mathbf{R}_s$ and is absorbed into the terms
$R'^{(s)}$, $\mathbf{B}'^{(s)}$ produced by $\mathbf{R}_s$.

\emph{Fact~C.} An $\mathbf{R}$-born boundary term alive at $\mathsf{t}$ has copied, at its
birth, the values of only such objects as are still live-witnessed at $\mathsf{t}$.
\end{definition}

\begin{proof}
\emph{Fact~A.} Let $\psi$ be a stored term of level $e$, alive at $\mathsf{t}$.

By part~(ii) of the anchoring statement, $X_\psi$ has a cube (its anchor cube) lying in
$(Z_e^{\boxplus})^{(1)}$. Here $Z_e$ is the region of the history current at $\mathsf{t}$, not
of the history at the birth of $\psi$. The anchor cube therefore lies in $(C^{\boxplus})^{(1)}$
for a component $C$ of the current $Z_e$, and so $(C^{\boxplus})^{(1)}\cap X_\psi\neq\emptyset$.

The component $C$ belongs to the history current at $\mathsf{t}$. It has therefore not been
deleted by a healing at any time $\le\mathsf{t}$, so it is alive at $\mathsf{t}$. Components of
$\Lambda^{c}$ are deleted only by healings, and arrests leave the sets $\Lambda$ unchanged, as
recalled in part~(c). Hence $C$ is a component of $Z_e$ at the birth of $\psi$, and it is a
witnessing component of $\psi$. Thus $\psi$ is live-witnessed at $\mathsf{t}$.

\emph{Fact~B.} Let $\Psi$ be a stored object containing an occurrence of $\psi$. By the
localisation of reading in part~(a), $X_\Psi\supset X_\psi$. By part~(b), $X_\psi$ contains a
cube of $(C^{\boxplus})^{(1)}$.

Let $X^{(s)}$ be the descendant of $C$ in the operated class $Z^{(s)}$ of $\mathbf{R}_s$, that
is, the region healed by $\mathbf{R}_s$ that contains $C$. Then
\begin{equation}\label{9.2:eq-witnessing-components-2}
(C^{\boxplus})^{(1)}\subset X^{(s)}\in Z^{(s)}.
\end{equation}
This inclusion follows from two facts. First, part~(i) of the anchoring statement gives
$(D^{\boxplus})^{(1)}\subset Z_{j+1}$ for a component $D$ of $Z_j$ with $j$ below the current
level. Second, each renormalization step adds at least ten layers of cubes to the large-field
regions \cite{B-CMP122a}. By the first fact, applied with $D=C$ and $j=e$,
$(C^{\boxplus})^{(1)}$ is contained in the large-field region of the next level; by the second,
the large-field regions containing $C$ at the successive levels up to $s$ only grow, so that
$(C^{\boxplus})^{(1)}$ stays inside the region containing $C$ at each of these levels, and in
particular inside the descendant $X^{(s)}$ of $C$ healed by $\mathbf{R}_s$.

Consequently $X_\Psi$ contains a cube of $X^{(s)}$ and meets $Z^{(s)}$; this is
\eqref{9.2:eq-witnessing-components-a}. At a completed operation $\mathbf{R}_s$ with operated
class $Z^{(s)}$, the terms of the rest of the effective action are split as follows
\cite{B-CMP122b}:
\begin{itemize}
\item the first part consists of the terms whose localisation domains intersect $Z^{(s)}$;
\item the second part consists of the terms whose localisation domains are disjoint from
$Z^{(s)}$.
\end{itemize}
The first part is expanded and resummed into the activities and into the terms $R'^{(s)}$,
$\mathbf{B}'^{(s)}$ produced by $\mathbf{R}_s$. Hence $\Psi$ belongs to the first part at
$\mathbf{R}_s$ and is absorbed into the terms $R'^{(s)}$, $\mathbf{B}'^{(s)}$ produced by
$\mathbf{R}_s$.

\emph{Fact~C.} Let $\varphi$ be an $\mathbf{R}$-born boundary term alive at $\mathsf{t}$, and let
$\psi$ be an object whose value $\varphi$ copied at its birth. Then $\varphi$ contains an
occurrence of $\psi$.

Suppose $\psi$ is not live-witnessed at $\mathsf{t}$. Then every witnessing component of $\psi$
has been deleted by a healing at a time $\le\mathsf{t}$. Let the last of them, $C$, be healed
by $\mathbf{R}_s$, so that $s\le\mathsf{t}$. Its deletion is a healing after the birth of $\psi$,
which took place after $\mathbf{R}_{e(\psi)}$, so $s>e(\psi)$.

Fact~B applied to $C$ and to the stored object $\varphi$ shows that $\varphi$ is absorbed into
the terms $R'^{(s)}$, $\mathbf{B}'^{(s)}$ produced by $\mathbf{R}_s$. Then $\varphi$ is not
alive at $\mathsf{t}\ge s$, a contradiction. Hence $\psi$ is live-witnessed at $\mathsf{t}$.
\end{proof}

\begin{definition}[The own-witness sets of a stored object]\label{9.2:def-own-witness-sets}
We use the object-occurrences and their witnessing components of
Definition~\ref{9.2:def-witnessing-components}. Let $\upsilon$ be a use-copy of class $z=(n,k')$,
$n$ the zone and $k'$ the level of the use. If $n<k'$, the anchor point of $\upsilon$ lies in
$\Lambda_{n+1}^c=Z_{n+1}$, the large-field region of level $n+1$, at level $k'$ of its history.
We write
$C_\upsilon$ for the component of $\Lambda_{n+1}^c$ that contains this anchor point, and we call
$\upsilon$ \emph{geometrically witnessed at a time $t$} when $C_\upsilon$ is alive at $t$; its
witness level is then $n+1$. A use-copy with $n=k'$ lies in the top zone at its level and has no
geometric witness. A use-copy inside a formula is \emph{direct} if it does not lie inside a nested
leaf-formula, a nested $\mathbf{T}$-born object or a nested frozen term. For each stored object,
that is, each object-occurrence of Definition~\ref{9.2:def-witnessing-components} and each
second-group term $R'^{(s)}(X)$ as in~(c) below, and for each class $z=(n,k')$, we define a set
of use-copies $\mathcal{W}(\cdot,z)$, the \emph{own-witness set of class $z$}. The set is defined
once, at the birth of the object, and does not change afterwards. There are three cases.

\begin{enumerate}
\item[(a)] \emph{$\mathbf{T}$-born and frozen objects.} Let $\psi$ be of level $e$, either a
$\mathbf{T}$-born object $\mathbf{B}^{(e)}(X_\psi)$, born at the $\mathbf{T}$-operation of step
$e$, or a frozen term $\mathbf{C}$, born at the arrest of its component. Then
$\mathcal{W}(\psi,z)$ is the set of direct use-copies of class $z$ of $\psi$ with $n=k'$, that is,
those in the top zone at the level of the use; for these $k'\in\{e-1,e\}$. Every other use-copy
read by $\psi$ is either geometrically witnessed (the case $n<k'$), or lies inside an older
occurrence whose value $\psi$ copies, and such an occurrence keeps its own set.

\item[(b)] \emph{$\mathbf{R}$-born leaves.} Let $\varphi=\mathbf{B}'^{(s)}(X')$ be born at
$\mathbf{R}_s$. Let $X_1,\dots,X_p$ be the healed (first-class) components of the operated
family and $Y_1,\dots,Y_m$ the second-class components with which $\varphi$ is born at
$\mathbf{R}_s$. Put
\begin{equation}\label{9.2:eq-own-witness-sets-a}
\mathcal{W}(\varphi,z):=(\mathrm{A1})\cup(\mathrm{A2})\cup(\mathrm{B})\cup(\mathrm{G}),
\end{equation}
where $(\mathrm{A1})$, $(\mathrm{A2})$, $(\mathrm{B})$, $(\mathrm{G})$ are the following sets of
use-copies of class $z$.
\begin{itemize}
\item[(A1)] The use-copies of class $z$ made in the histories of the $X_i$. They are of two
kinds.
\begin{itemize}
\item[$\cdot$] Those made by the inner operations $\mathbf{T}^{(k')}$, $k'<s$, of
$\mathbf{T}'_s(X_i)$ (see \cite[(1.71)]{B-CMP122b} and \cite[(2.22)]{B-CMP119}).
\item[$\cdot$] Those made by the terms of $V(X_i^{\sim})$, and of $V(Y)$ for regions $Y$
containing a component of $X_i^{\sim}$, that belong to the old zones of $X_i$; here
$V(X_i^{\sim})$ and $V(Y)$ denote the terms attached to the regions $X_i^{\sim}$ (the
enlargement of $X_i$) and $Y$ in \cite{B-CMP122b}, distinct from the block field $V_k$. The old
zones are read off from the old determining set $\{\Gamma^0_j\}$ of \cite{B-CMP122b}.
\end{itemize}
Equivalently, $(\mathrm{A1})$ is the set of direct use-copies of class $z$ of the formula of
$\varphi$ whose anchor lies in $\bigcup_i X_i$ and whose level satisfies $k'<s$.

\item[(A2)] The direct use-copies of class $z$ of the formula of $\varphi$ with $k'=n=s$. These
are the uses in the top zone at production, and they may lie anywhere in $X'$. By the definition
of $(\mathrm{A1})$, every direct use of level $s$ inside $\bigcup_i X_i$ is of this kind, since
there the new determining set is $\{Z^{(s)}\}\subset\Gamma_s$ \cite[(1.33)]{B-CMP122b},
$\Gamma_s$ being the determining set of level $s$.

\item[(B)] The union
\begin{equation*}
\bigcup_{\psi}\mathcal{W}(\psi,z),
\end{equation*}
over all object-occurrences $\psi$ nested in the formula of $\varphi$ all of whose witnessing
components are among the $X_i$ or descend to them. These are the occurrences that lose
their last witness at $\mathbf{R}_s$.

\item[(G)] The geometrically witnessed use-copies of class $z$ inside the formula of $\varphi$,
at any depth of nesting, whose component $C_\upsilon$ is one of the $X_i$ or descends to one of
them. These use-copies lose their geometric witness at $\mathbf{R}_s$.
\end{itemize}
The following do not enter $\mathcal{W}(\varphi,z)$: the nested occurrences that keep a live
witness, namely a surviving component (an unoperated component, one of the $Y_h$, or an arrested
component); and the geometrically witnessed use-copies whose component $C_\upsilon$ survives.

\item[(c)] \emph{Second-group terms.} Let $R'^{(s)}(X)$ be a second-group term of
\cite[(1.98)]{B-CMP122b}, with $X\subset\Lambda_s^{\mathrm{new}}$, the part of $\Lambda_s$ over
which the second-group terms $R'^{(s)}(X)$ are indexed; here $\Lambda_s$ is the complement of the
large-field region $Z_s$ of level $s$, so that $\Lambda_s^{c}=Z_s$.
We put
\begin{equation*}
\mathcal{W}\bigl(R'^{(s)}(X),z\bigr):=(\mathrm{A1})\cup(\mathrm{A2})\cup(\mathrm{B})\cup(\mathrm{G}),
\end{equation*}
where the four sets are formed as in~(b) for the cluster of regions and components from which
$R'^{(s)}(X)$ arises, with $R'^{(s)}(X)$ in place of $\varphi$.
\end{enumerate}
\end{definition}

\begin{lemma}[Partition of uses by witness, and the age forced at a reader]
\label{9.2:lem-witness-partition}
We use the object-occurrences, their levels $e(\psi)$ and domains $X_\psi$, the witnessing
components, the geometric witness $C_\upsilon$ of a use-copy $\upsilon$, and the words
\emph{alive}, \emph{live-witnessed} and \emph{geo-witnessed}, as fixed in
Definition~\ref{9.2:def-witnessing-components}. We use the own-witness sets
$\mathcal{W}(\psi,z)$ of Definition~\ref{9.2:def-own-witness-sets}, and the ages of
admissible sequences as fixed in Notation~\ref{9.2:not-age-factor-majorants}. Let
$z=(n,k')$ be a class.
\begin{itemize}
\item[(W1)] At every time $t$ of the history, let $\upsilon$ be a use-copy of class $z$ inside
the formula of an object alive at $t$. Then exactly one of the following holds:
\begin{itemize}
\item $\upsilon$ is geo-witnessed at $t$;
\item $\upsilon\in\mathcal{W}(\psi,z)$ for exactly one object-occurrence $\psi$, and that
$\psi$ is live-witnessed at $t$.
\end{itemize}
The \emph{witness level} $w(\upsilon)$ is given by the first line of
\eqref{9.2:eq-witness-partition-a}.
\item[(W2)] Let $k$ be at least the level of $t$, and let $X$ be a component of $Z_k$ (the
\emph{reader}: the stored objects enter $\mathbf{T}'_k(X)$ and are read there at level $k$). Let
$\omega$ be one of the admissible sequences of $X$. Let $\upsilon$ be a use-copy of class $z$
present in the $\omega$-summand of
\begin{equation*}
\mathsf{G}=\mathbf{T}'^{\,\lambda}_k(X)F,
\end{equation*}
that is, of $\mathbf{T}'_k(X)F$ as a function of the dial variables. The use-copy may be
present in one of two ways:
\begin{itemize}
\item directly, through the inner operations of $X$ itself, through the direct terms of
the sum $V(X^{\sim})$ of boundary terms of \cite[\S1]{B-CMP122b}, or through the constants of
the summand, each of which is attached to an index $j$ with $Z_j\cap X\neq\emptyset$ in
$\omega$;
\item inside an object-occurrence lying in $X^{\sim}$.
\end{itemize}
Then $\operatorname{age}(\omega)=k-j(X)(\omega)$ obeys the second line of
\eqref{9.2:eq-witness-partition-a}.
\end{itemize}
Here $j(X)(\omega)$ is the least index $i$ with $Z_i\cap X\neq\emptyset$ in the history
completed inside $X$ by $\omega$.
\begin{equation}\label{9.2:eq-witness-partition-a}
\begin{gathered}
w(\upsilon):=
\begin{cases}
n+1 & \text{if $\upsilon$ is geo-witnessed at $t$},\\
e(\psi) & \text{if $\upsilon\in\mathcal{W}(\psi,z)$},
\end{cases}
\\[4pt]
\operatorname{age}(\omega)\ge
\begin{cases}
k-w(\upsilon) & \text{if $\upsilon\in\mathcal{W}(\psi,z)$},\\
(k-n-1)_+ & \text{if $\upsilon$ is geo-witnessed or a direct use of $X$.}
\end{cases}
\end{gathered}
\end{equation}
\end{lemma}

\begin{proof}
Throughout, $p$ denotes the anchor of $\upsilon$. Recall that $Z_l=\Lambda_l^c$ for every
level $l$.

Three properties of Ba{\l}aban's construction are used below.
\begin{itemize}
\item \emph{Filling property of the enlargements.} Let $Y$ be a component of $Z_l$. Then
every point of $Y^{\boxplus}$ that lies in $\Lambda_i^c$ for some $i\le l$ lies in $Y$.
Moreover $Y^{\sim}\subset Y^{\boxplus}$ (see \cite[\S2]{B-CMP119}).
\item \emph{First part of the anchor property.} Let $\psi$ be an object-occurrence of level
$e$, and let $C$ be a witnessing component of $\psi$ that is alive. If $\psi$ is a leaf with
anchor cube $Q$, then $Q\subset X_\psi\cap(C^{\boxplus})^{(1)}$ and $Q\subset Z_{e+1}$. If
$\psi$ is a carrier, then $X_\psi\cap(C^{\boxplus})^{(1)}\neq\emptyset$ and this set lies in
$Z_{e+1}$ (Definition~\ref{9.2:def-witnessing-components}).
\item \emph{Absorption at a completed $\mathbf{R}_s$.} Let $Z^{(s)}$ be the operated class
of $\mathbf{R}_s$. Every term of the effective action whose localisation domain meets
$Z^{(s)}$ is expanded and resummed into the first-group terms
$\mathbf{B}'^{(s)}(X')$ and the second-group terms $R'^{(s)}(\cdot)$ of $\mathbf{R}_s$
\cite[\S1]{B-CMP122b}.
\end{itemize}
We also use the reading \eqref{9.2:eq-age-factor-majorants-e} of the sum over admissible
sequences. By this reading, a stored object whose production required a component $C$ of
$Z_e$ is absent from every $\omega$-summand in which $C\notin Z_e^{\omega}$. Likewise, an
object whose production required a point $p$ in zone $n$ at level $k'$ is absent from every
$\omega$-summand in which $p\notin\Gamma_n^{\omega}$ at $k'$. Here $Z_e^{\omega}$ is the
level-$e$ large-field region, and $\Gamma_n^{\omega}$ the zone-$n$ determining set, of the
history completed inside $X$ by $\omega$.

\emph{Proof of (W1).} We argue by induction along the history.

\emph{Creation of $\upsilon$.} The use-copy $\upsilon$ is created at some level $k'$. This
happens either in the production of some object $\psi_0$, or directly by an inner operation.
\begin{itemize}
\item If $n<k'$, the anchor $p$ lies in $\Lambda_{n+1}^c$. So the component $C_\upsilon$ of
$\Lambda_{n+1}^c$ containing $p$ exists and is alive, and $\upsilon$ is geo-witnessed.
\item If $n=k'$, then $\upsilon\in\mathcal{W}(\psi_0,z)$. For an $\mathbf{R}$-born object
this is the part (A2) of \eqref{9.2:eq-own-witness-sets-a}. For a $\mathbf{T}$-born or
frozen object (a carrier, in the terminology of
Definition~\ref{9.2:def-witnessing-components}), this is the clause of
Definition~\ref{9.2:def-own-witness-sets} for carriers: their own-witness set of class $z$
consists of their direct use-copies of class $z$ with $n=k'$. Moreover, by
Definition~\ref{9.2:def-witnessing-components}, $\psi_0$ is live-witnessed at its birth.
\end{itemize}

\emph{A $\mathbf{T}$-step or an arrest.} Such a step deletes no component of any
$\Lambda^c_i$. The objects it creates have own-witness sets that consist only of their own
top-zone uses. Hence no use-copy already present changes its witness, and the alternative of
(W1) that held before the step holds after it.

\emph{A completed $\mathbf{R}_s$.} Let $Z^{(s)}$ be the operated class of $\mathbf{R}_s$.
Three cases arise.
\begin{itemize}
\item[(a)] Let $\upsilon$ be geo-witnessed, with $C_\upsilon$ contained in (descending to) one
of the components of $Z^{(s)}$ healed at this step. Then
\begin{equation*}
p\in C_\upsilon\subset Z^{(s)}.
\end{equation*}
So every object containing $\upsilon$ has localisation domain meeting $Z^{(s)}$. By the
absorption property, each such
object is absorbed, occurrence by occurrence, into some first- or second-group term
$\varphi$ of $\mathbf{R}_s$. Then $\upsilon$ is geo-witnessed inside the formula of
$\varphi$, and its component $C_\upsilon$ descends to one of the healed components. So
$\upsilon\in\mathcal{W}(\varphi,z)$ by the part (G) of
\eqref{9.2:eq-own-witness-sets-a}.
\item[(b)] Let $\upsilon\in\mathcal{W}(\psi,z)$, where the last live witnessing component of
$\psi$ is healed at this step. By \eqref{9.2:eq-witnessing-components-a}, every object
containing that occurrence of $\psi$ is absorbed into some $\varphi$. Since the component
healed now was the last live witnessing component of $\psi$, the occurrence $\psi$ nested in
the formula of $\varphi$ loses its last witness at $\mathbf{R}_s$. So
$\upsilon\in\mathcal{W}(\varphi,z)$ by the part (B) of
\eqref{9.2:eq-own-witness-sets-a}.
\item[(c)] In every other case, $\upsilon$ keeps its witness. This witness is either
$C_\upsilon$ or a live witnessing component of $\psi$, and it stays alive.
\end{itemize}

\emph{Uniqueness.} Each of the cases above transfers membership in an own-witness set, and
none duplicates it. The reason is that the four parts (A1), (A2), (B), (G) of
\eqref{9.2:eq-own-witness-sets-a} are pairwise disjoint by construction:
\begin{itemize}
\item the parts (A1) and (A2) consist of direct use-copies of the formula of $\varphi$;
\item the part (B) consists of use-copies lying in the own-witness sets of nested
occurrences;
\item the part (G) consists of geo-witnessed copies, which belonged to no own-witness set
before the step.
\end{itemize}
So the occurrence $\psi$ of the second alternative in (W1) is unique at every time. This
completes the induction.

\emph{Proof of (W2).} First note that $X$ is a component of $Z_k$. Hence
$Z_k\cap X=X\neq\emptyset$, so $j(X)(\omega)\le k$ and $\operatorname{age}(\omega)\ge 0$. It
follows that the bound $(k-n-1)_+$ needs proof only when $n+1\le k$. Likewise the bound
$k-w(\upsilon)$ needs proof only when $w(\upsilon)<k$. We treat three cases.

\emph{Direct uses of $X$.} Let $\upsilon$ be a direct use of $X$ of a non-top zone, $n<k'$, at
its level $k'\le k$. In the history completed inside $X$ by $\omega$, zone $n$ lies in
\begin{equation*}
\Omega^c_{n+1}\subset Z_{n+1}.
\end{equation*}
The anchor $p$ lies in zone $n$, hence in $Z_{n+1}=\Lambda^c_{n+1}$, and in
$X^{\sim}\subset X^{\boxplus}$. Since $n+1\le k$, the filling property, applied to the
component $X$ of $Z_k$ with $i=n+1$, gives $p\in X$.
Hence $Z_{n+1}\cap X\neq\emptyset$ in $\omega$. So $j(X)(\omega)\le n+1$, and
\begin{equation*}
\operatorname{age}(\omega)\ge k-n-1.
\end{equation*}

For the constants of the summand the reasoning is similar. A constant attached to an index
$j$ with $Z_j\cap X\neq\emptyset$ in $\omega$ is assigned the zone $n:=j$. Then
$j(X)(\omega)\le j$, and
\begin{equation*}
\operatorname{age}(\omega)\ge k-j=k-n\ge k-n-1.
\end{equation*}

Consider finally a top-zone use ($n=k'$) made by the inner operation $\mathbf{T}^{(k')}$. By
\cite[(2.22)]{B-CMP119}, this operation lives on $Z_{k'+1}\cap X$. In the $\omega$-summand it is
therefore present only if $Z^{\omega}_{k'+1}\cap X\neq\emptyset$. Then
$j(X)(\omega)\le k'+1=n+1$, and again $\operatorname{age}(\omega)\ge k-n-1$.

\emph{Geo-witnessed $\upsilon$ inside an occurrence lying in $X^{\sim}$.} Here $p\in
C_\upsilon$, with $C_\upsilon$ alive. Since $n+1\le k$,
\begin{equation*}
C_\upsilon\subset\Lambda^c_{n+1}\subset\Lambda^c_k.
\end{equation*}
Also $p\in X^{\sim}\subset X^{\boxplus}$. Apply the filling property to the component $X$
of $Z_k$, with $i=n+1\le k$. It gives $p\in X$. Now $C_\upsilon$ is connected, is contained
in $\Lambda^c_k=Z_k$, and meets the component $X$ of $Z_k$ at $p$. Therefore
$C_\upsilon\subset X$.

By the reading \eqref{9.2:eq-age-factor-majorants-e}, the occurrence is present in the
$\omega$-summand only if $p\in\Gamma_n^{\omega}$ at $k'$. This means
\begin{equation*}
C_\upsilon\subset Z^{\omega}_{n+1}\cap X.
\end{equation*}
Hence $j(X)(\omega)\le n+1$, and $\operatorname{age}(\omega)\ge k-n-1$.

\emph{$\upsilon\in\mathcal{W}(\psi,z)$.} Here $\psi$ is live-witnessed by a component $C$ of
$Z_e$ that is alive, where $e:=e(\psi)\le k$ and $w(\upsilon)=e$. If $e=k$ there is nothing to
prove, so let $e<k$. Since the occurrence lies in $X^{\sim}$, we have $X_\psi\subset
X^{\sim}$.

Suppose first that $\psi$ is a leaf. Let $Q$ be its anchor cube, so that
\begin{equation*}
Q\subset X_\psi\cap(C^{\boxplus})^{(1)}.
\end{equation*}
By the first part of the anchor property, which applies because $C$ is alive,
$Q\subset Z_{e+1}$. Also $Q\subset X_\psi\subset X^{\sim}\subset X^{\boxplus}$. Since
$e+1\le k$, the filling property gives $Q\subset X$.

The sets $C$ and $Q$ lie in one component of $Z_{e+1}$. Because $e+1\le k$,
\begin{equation*}
Z_{e+1}\subset Z_k.
\end{equation*}
So this component is a connected subset of $Z_k$ that meets $X$ (in $Q$). It is therefore
contained in $X$, and in particular $C\subset X$.

Suppose next that $\psi$ is a carrier. By the first part of the anchor property, the set
$X_\psi\cap(C^{\boxplus})^{(1)}$ is non-empty and lies in $Z_{e+1}$. The argument of the two
preceding paragraphs applies to this set in place of $Q$. Again $C\subset X$.

By the reading \eqref{9.2:eq-age-factor-majorants-e}, the occurrence is present in the
$\omega$-summand only if $C\subset Z_e^{\omega}$. Since $C\subset X$, this gives
\begin{equation*}
j(X)(\omega)\le e=w(\upsilon),
\end{equation*}
that is, $\operatorname{age}(\omega)\ge k-w(\upsilon)$. This completes the proof of (W2).
\end{proof}

\begin{lemma}[Presence of own-witness uses inside the healed components]\label{9.2:lem-presence-healed}
Consider the step $\mathbf{R}_s$, and let $X_1,\dots,X_p$ be the healed components operated by $\mathbf{T}'_s$ at that step. Let $\varphi$ be a term of the first or of the second group of \cite[(1.98)]{B-CMP122b}. For $i=1,\dots,p$ let $\omega_i$ range over the admissible sequences of $X_i$ in the sum \cite[(1.71)]{B-CMP122b}. Write $Z^{\omega_i}_j$ for the large-field region of index $j$ of the history completed inside $X_i$ by $\omega_i$, and put
\begin{equation*}
\operatorname{age}_s(\omega_i):=s-j(X_i)(\omega_i),
\end{equation*}
where $j(X_i)(\omega_i)$ is the first index of $X_i$ along $\omega_i$ (Notation~\ref{9.2:not-age-factor-majorants}). The own-witness sets are those of Definition~\ref{9.2:def-own-witness-sets}, with the parts (A1), (B), (G) of \eqref{9.2:eq-own-witness-sets-a}.
\begin{itemize}
\item[(P1)] Let $\upsilon\in\mathrm{(A1)}$ be a use-copy of class $(n,k')$, and let $X_i$ be the component containing its anchor. Then $\upsilon$ is present in the $(\omega_1,\dots,\omega_p)$-summand of the cluster of $\varphi$ only if the condition in line (P1) of \eqref{9.2:eq-presence-healed-a} holds; consequently $\operatorname{age}_s(\omega_i)\ge s-n-1$.
\item[(P2)] Let $\upsilon\in\mathrm{(G)}$ be a use-copy of class $(n,k')$ with $C_\upsilon\subset X_i$. Then $\upsilon$ is present in that summand only if $C_\upsilon\subset Z^{\omega_i}_{n+1}$, the condition in line (P2); this forces $\operatorname{age}_s(\omega_i)\ge s-n-1$.
\item[(P3)] Let $\psi\in\mathrm{(B)}$ be an occurrence of level $e$ all of whose witnessing components lie in $\bigcup_i X_i$, and let $C\subset X_{i_0}$ be one of them. Then $\psi$ is present in that summand only if the condition in line (P3) of the display holds, and then $\operatorname{age}_s(\omega_{i_0})\ge s-e$.
\end{itemize}
The three conditions and age bounds are
\begin{equation}\label{9.2:eq-presence-healed-a}
\begin{aligned}
&\text{(P1)}\quad
\begin{cases}
Z^{\omega_i}_{n+1}\cap X_i\neq\emptyset, & n<k',\\
Z^{\omega_i}_{k'+1}\cap X_i\neq\emptyset, & n=k'<s,
\end{cases}
\qquad \operatorname{age}_s(\omega_i)\ge s-n-1;\\
&\text{(P2)}\quad C_\upsilon\subset Z^{\omega_i}_{n+1},
\qquad \operatorname{age}_s(\omega_i)\ge s-n-1;\\
&\text{(P3)}\quad C\subset Z^{\omega_{i_0}}_{e}\cap X_{i_0},
\qquad \operatorname{age}_s(\omega_{i_0})\ge s-e.
\end{aligned}
\end{equation}
\end{lemma}

\begin{proof}
\emph{Proof of (P1).} We use the localisation of uses that is established in the proof of the theorem of this chapter named the one-lattice Lipschitz step. We apply it at level $s$ to the component $X_i$ that contains the anchor of $\upsilon$. This $X_i$ is the component operated by $\mathbf{T}'_s$, and we take it with its own history.

That localisation states the following. A use inside $X_i$ in zone $n$, a subset of $\Lambda_{n+1}^c$, at a level $k'<s$ exists only in those sequences $\omega_i$ whose zone $n$ at level $k'$ is present inside (that is, meets) the one-layer enlargement $X_i^{\boxplus}$ of $X_i$.

By the filament part of the volume and filament bound, applied at level $s$, the part of zone $n$ met inside $X_i^{\boxplus}$ lies inside $X_i$. Hence the $\omega_i$-summand carries $\upsilon$ only if zone $n$ of the history $\omega_i$ meets $X_i$. For $n<k'$ this is the condition $Z^{\omega_i}_{n+1}\cap X_i\neq\emptyset$ of \eqref{9.2:eq-presence-healed-a}.

For $n=k'<s$ the use is made by the inner operation $\mathbf{T}^{(k')}$ of $X_i$. By \cite[(2.22)]{B-CMP119} the inner operation of level $k'$ lives on $Z_{k'+1}\cap X_i$. So the $\omega_i$-summand carries $\upsilon$ only if $Z^{\omega_i}_{k'+1}\cap X_i\neq\emptyset$.

\emph{Proof of (P2) and (P3).} We use the representation \eqref{9.2:eq-age-factor-majorants-e} of the density as a sum over admissible histories. Under it, the $\omega$-summand of $\mathbf{T}'_k(X)$ contains exactly the stored objects produced along the label's history completed inside $X$ by $\omega$. A stored object whose production required a component $C$ of $Z_e$ is absent from the $\omega$-summands in which $C$ is not a component of $Z_e^{\omega}$.

In the case (P2) we apply this reading to the object containing the copy $\upsilon$. That object was produced along a history in which $C_\upsilon$ is a component of $Z_{n+1}$.

In the case (P3) we apply it to the occurrence $\psi$ itself or, if $\psi$ is present only through a copy carried by a term of the first part of the rest of the effective action, to that term. Here the first part, at a completed $\mathbf{R}$-step, consists of the terms of the rest of the effective action whose localisation domains intersect the operated region \cite[\S1]{B-CMP122b}; it is distinct from the first group of \cite[(1.98)]{B-CMP122b}. The production history of that term contains that of $\psi$. In either case the object was produced along a history in which $C$ is a component of $Z_e$.

Inside the sum \cite[(1.71)]{B-CMP122b}, such a requirement is a condition on the sequence $\omega_{i_0}$ alone (in the case (P2), $i_0=i$). Indeed, \cite[\S1]{B-CMP122b} states that the terms localised in a region $Y$ depend on the sequence $\{\Omega_j,Z_j\}$ restricted to the components of $Z$ contained in $Y$, and that the sum in \cite[(1.71)]{B-CMP122b} acts also on the last exponential in \cite[(1.72)]{B-CMP122b}. Inside $X_{i_0}$ this restricted sequence is $\{\Omega_j\cap X_{i_0},\,Z_j\cap X_{i_0}\}$, that is, $\omega_{i_0}$. Hence the object is present in the $(\omega_1,\dots,\omega_p)$-summand only if $C_\upsilon\subset Z^{\omega_i}_{n+1}$, respectively only if $C\subset Z^{\omega_{i_0}}_e$.

Finally, in (P2), $C_\upsilon\subset X_i$ by hypothesis. In (P3), a witnessing component that lies in $\bigcup_i X_i$ is contained in a single $X_{i_0}$, because large-field regions grow into their descendants, by the growth property of the large-field regions in Notation~\ref{9.2:not-age-factor-majorants}: the witnessing component $C$ descends to a component contained in $X_{i_0}$. This justifies the choice $C\subset X_{i_0}$ and gives the intersection with $X_{i_0}$ in line (P3). This gives the conditions of (P2) and (P3) in \eqref{9.2:eq-presence-healed-a}.

\emph{The age bounds.} In each case the age bound restates the condition through the first index $j(X_i)(\omega_i)$ of Notation~\ref{9.2:not-age-factor-majorants}. In (P1) and (P2) the history $\omega_i$ has a large-field region of index $n+1$ in $X_i$, so $j(X_i)(\omega_i)\le n+1$. Therefore $\operatorname{age}_s(\omega_i)\ge s-n-1$. In (P3) it has one of index $e$ in $X_{i_0}$, so $j(X_{i_0})(\omega_{i_0})\le e$. Therefore $\operatorname{age}_s(\omega_{i_0})\ge s-e$.
\end{proof}

\begin{remark}[Why uses are grouped by witness level]\label{9.2:rem-witness-level-grouping}
The auxiliary dials $\nu_e$ are not introduced one per level of nesting of occurrences. They are introduced one per witness level. The content of a stored object is grouped at its witness level, the level of the innermost enclosing occurrence that witnesses it. This grouping coincides with the static membership fixed by the own-witness sets of Definition~\ref{9.2:def-own-witness-sets}. With it, both the step at an operation $\mathbf{R}_s$ that produces a new term and the transport of that term to the level at which it is read are exact. Grouping instead by the level of nesting, that is, at the level of the enclosing occurrence itself, would require of that occurrence a decay that it does not have.
\end{remark}

\begin{proof}
Let $\psi'$ be an occurrence of level $e'$ nested in an occurrence $\psi$ of level $e>e'$. Suppose that the anchor of $\psi'$ is live at the birth of $\psi$ and is healed later, at the operation $\mathbf{R}_s$.

While the anchor of $\psi'$ was live, $\psi$ copied the content of $\psi'$ by value. Hence the top-zone and healed content of $\psi'$ carries no decay at level $e$.

At $\mathbf{R}_s$, the content of $\psi'$ passes to the term $\varphi_s$ that absorbs the healed region. We write $X^{(s)}$ for this region, with healed components $X_1,\dots,X_p$. The content passes through the part of the own-witness sets of $\varphi_s$ (Definition~\ref{9.2:def-own-witness-sets}) that collects nested occurrences. It does not pass to $\psi$.

Since the anchor of $\psi'$ is healed at $\mathbf{R}_s$, the witnessing components of $\psi'$ lie in $X_1\cup\dots\cup X_p$; let $C$ be one of them and $X_{i_0}$ the healed component containing it. Apply part (P3) of Lemma~\ref{9.2:lem-presence-healed} to the occurrence $\psi'$, of level $e'$. In the notation of that lemma, with $(\omega_1,\dots,\omega_p)$ ranging over the admissible sequences of the healed components, it shows that $\psi'$ is present in the $(\omega_1,\dots,\omega_p)$-summand of the expansion of $\varphi_s$ only if
\begin{equation*}
\operatorname{age}_s(\omega_{i_0})\ge s-e'.
\end{equation*}
Therefore the split of the age dial $\mu$ of $\varphi_s$ at the threshold $s-e'$ pays, for the content of $\psi'$, the age of $X^{(s)}$.

The occurrence $\psi$ copies by value, so it never acquires an age factor for the content of $\psi'$. Consequently that content must be grouped at $e'$, its witness level before this healing. This is the level of the innermost live-witnessed occurrence enclosing it, that is, the static membership of Definition~\ref{9.2:def-own-witness-sets}. With this grouping the step at $\mathbf{R}_s$ producing $\varphi_s$ and the transport of $\varphi_s$ to the level at which it is read are both exact.

Grouping the content at $e$ would supply too small a factor. It would ask $\psi$ for a decay at level $e$ which, by the value-copying described above, $\psi$ does not have.
\end{proof}

\begin{lemma}[The per-cell pile of first-group leaves over all levels]\label{9.2:lem-leaf-pile}
We use the following notation. For a level $l$, $\ell_l=L^l\eta$ is the spacing of the level-$l$ lattice, where $\eta=L^{-K}$ is the bare spacing of Definition~\ref{1.1:def-lattice}. For a set $S$ of the torus, $|S|_l$ is the number of level-$l$ cells contained in $S$. The partition $\pi_l$ of the torus consists of the cubes of side $M\ell_l$; it is chosen compatible with all the previous partitions, as in \cite{B-CMP119}. The native strata $\Gamma^{\natural}_l$ partition the torus into unions of level-$l$ cells, and a point has \emph{native index} $l$ if it lies in $\Gamma^{\natural}_l$. We put
\begin{equation}\label{9.2:eq-leaf-pile-1}
K_0:=\sup_{j,\ \square\in\pi_j}\ \sum_{D\in\mathbf{D}_j,\ D\supset\square}e^{-(\kappa-1)d_j(D)} ,
\end{equation}
where $\mathbf{D}_j$ is the class of localisation domains of level $j$. By the splitting of the sum given in \cite{B-CMP119}, $K_0$ is a finite number depending only on the dimension $d=4$ and on $\kappa$.

Fix a history and a time, let $k\le K$ be the current level at that time, and let $Z_e$ be the large-field region of level $e$ of the current label. Put $P_e:=(Z_e^{\boxplus})^{(1)}$, the set obtained by the operation $(\cdot)^{(1)}$ from the one-layer enlargement $Z_e^{\boxplus}$ of $Z_e$, formed in this current label; it is the region of the current label in which the stored boundary-type terms of level $e$ have their anchor cubes (first item below).

A \emph{first-group leaf} of level $e$ is an $\mathbf{R}$-born first-group boundary term $\mathbf{B}'^{(e)}(X_\psi)$ of \cite[(1.98)]{B-CMP122b}; it is \emph{plain} if its domain $X_\psi$ does not meet the union $\bigcup_hY_h$ of the regions $Y_h$ of \cite[(1.99)]{B-CMP122b}. Let $\mathcal{L}_e$ be the set of plain first-group leaves of level $e$ that are alive (stored and not yet ceased) at the fixed time. They are counted as stored family members, and not as copies of their formula inside later objects, in the sense of Definition~\ref{9.2:def-witnessing-components}. Each $\psi\in\mathcal{L}_e$ has two attributes:
\begin{itemize}
\item an anchor cube $\square_\psi\in\pi_e$ with $\square_\psi\subset X_\psi\cap P_e$, supplied by the anchor property of stored boundary-type terms in the current label;
\item a dial majorant, which is the leaf majorant \eqref{9.2:eq-age-factor-majorants-b} of Notation~\ref{9.2:not-age-factor-majorants} for a plain leaf:
\begin{equation}\label{9.2:eq-leaf-pile-2}
\mathfrak{B}_\psi=C_{99}\,e^{-p_0(g_e)}\,e^{-(1+\frac12\beta)\kappa d_e(X_\psi)},
\end{equation}
with the exponent $\beta$ of \cite[(1.99)]{B-CMP122b}.
\end{itemize}

The constant $C_L$ is a budget constant of the allowance, chosen at least as large as the constant that converts sup-norm bounds on stored terms into per-cell value contributions; since this condition on $C_L$ is a lower bound only, we may and do take $C_L\ge1$. The constant $C_L'$ is the budget constant in whose units per-native-cell piles are measured in \eqref{9.2:eq-age-factor-majorants-d}. Assume the following.
\begin{enumerate}
\item[(1)] (From the volume and filament bound.) Every point of $P_e$ has native index at most $e$.
\item[(2)] (From the volume and filament bound.) For every $\square\in\pi_e$ all of whose points have native index at most $e$,
\begin{equation}\label{9.2:eq-leaf-pile-3}
\sum_{l\le e}L^{-4(e-l)}\,|\square\cap\Gamma^{\natural}_l|_l=M^4 .
\end{equation}
\item[(3)] (The lower member of the reference window, Definition~\ref{1.2:def-flow-constants}.) For $0\le i<j\le K$ one has $\lambda(j-i)-2c_2\le x_i-x_j$, and $x_K>X$, where $X=g^{-2}$ and $\lambda>0$ is the level-wise rate of Notation~\ref{9.1:not-levelwise-rate}.
\item[(4)] (Relative count for strip leaves.) For the first-group leaves that are not plain (strip leaves, whose domain meets $\bigcup_hY_h$), the relative count of first-group leaves holds; that is, the bound of the form \eqref{9.2:eq-leaf-pile-a} holds for the strip leaves of level $e$, in which $d_e$ is replaced by the relative length of \cite[(1.99)]{B-CMP122b} and the factor $e^{-p_0(g_e)}$ by the value $\alpha^{1/3}$ that the constant $c_1$ of \cite[(1.99)]{B-CMP122b} takes for such leaves, and the corresponding bound summed over levels holds as well. We use it as stated.
\end{enumerate}
Then for every level $e$ and every set $S$ of the torus,
\begin{equation}\label{9.2:eq-leaf-pile-a}
\begin{split}
\sum_{\psi\in\mathcal{L}_e,\ \square_\psi\subset S}\mathfrak{B}_\psi
&\le K_0M^{-4}\,C_{99}\,e^{-p_0(g_e)}\sum_{l\le e}L^{-4(e-l)}\,
|S\cap\Gamma^{\natural}_l|_l\\
&\le \varepsilon_e\,C_L'\sum_{l}|S\cap\Gamma^{\natural}_l|_l ,
\qquad
\varepsilon_e:=\frac{C_LK_0M^{-4}C_{99}}{C_L'}\,e^{-p_0(g_e)} .
\end{split}
\end{equation}
Moreover, uniformly in the history, the time, $k$ and $K$,
\begin{equation}\label{9.2:eq-leaf-pile-b}
\begin{gathered}
\varepsilon_{\mathrm{tot}}:=\sup\sum_{e}\varepsilon_e
\le\frac{C_LK_0M^{-4}C_{99}}{C_L'}\,\mathfrak{c}(X),\\
\mathfrak{c}(X):=\sum_{t\ge0}\exp\bigl(-A_0[\log(X-2c_2+\lambda t)]^{p_0}\bigr).
\end{gathered}
\end{equation}
Here the supremum is over histories, times, $k$ and $K$. The series $\mathfrak{c}(X)$ converges and decreases to $0$ as $X\uparrow\infty$ ($\gamma\downarrow0$).

In the accounting of \eqref{9.2:eq-age-factor-majorants-d}, the level-$e$ family of plain leaves therefore has per-native-cell pile at most $\varepsilon_e$. For strip leaves we write $\varepsilon^{\mathrm{strip}}_e$ and $\varepsilon^{\mathrm{strip}}_{\mathrm{tot}}$ for the constants of the relative count in hypothesis~(4), which is used as stated.
\end{lemma}

\begin{proof}
\emph{Grouping by anchor cubes.} The leaves of $\mathcal{L}_e$ are members of the family $(\mathbf{B}'^{(e)}(D))_{D\in\mathbf{D}_e}$ indexed by domains, $\psi\leftrightarrow\mathbf{B}'^{(e)}(X_\psi)$ with $X_\psi\in\mathbf{D}_e$; thus at most one leaf of $\mathcal{L}_e$ has a given domain. Fix $\square\in\pi_e$. Every $\psi$ with $\square_\psi=\square$ satisfies $X_\psi\supset\square$. We have $(1+\frac12\beta)\kappa\ge\kappa-1$ and $d_e(D)\ge0$, so each exponential factor satisfies $e^{-(1+\frac12\beta)\kappa d_e(D)}\le e^{-(\kappa-1)d_e(D)}$. Together with \eqref{9.2:eq-leaf-pile-2} and the definition \eqref{9.2:eq-leaf-pile-1} of $K_0$ this gives
\begin{equation}\label{9.2:eq-leaf-pile-4}
\sum_{\psi:\ \square_\psi=\square}\mathfrak{B}_\psi
\le C_{99}\,e^{-p_0(g_e)}\sum_{D\in\mathbf{D}_e,\ D\supset\square}
e^{-(1+\frac12\beta)\kappa d_e(D)}
\le C_{99}\,e^{-p_0(g_e)}\,K_0 .
\end{equation}

\emph{Counting the cubes.} If $\psi\in\mathcal{L}_e$ and $\square_\psi\subset S$, then $\square_\psi\subset S\cap P_e$, because $\square_\psi\subset P_e$. It therefore suffices to count the cubes $\square\in\pi_e$ with $\square\subset S\cap P_e$. We do this as in the count of the volume and filament bound. Let $\square$ be such a cube. By hypothesis~(1) every point of $\square\subset P_e$ has native index at most $e$, so hypothesis~(2) applies to $\square$. Dividing \eqref{9.2:eq-leaf-pile-3} by $M^4$ expresses the number $1$ as a weighted cell count of $\square$. The cubes of $\pi_e$ are pairwise disjoint, and those counted lie in $S$. Hence, for each $l$, the counts $|\square\cap\Gamma^{\natural}_l|_l$ of the distinct cubes add up to at most $|S\cap\Gamma^{\natural}_l|_l$. Summing over the cubes therefore gives
\begin{equation}\label{9.2:eq-leaf-pile-5}
\begin{split}
\#\{\square\in\pi_e:\ \square\subset S\cap P_e\}
&=\sum_{\square}M^{-4}\sum_{l\le e}L^{-4(e-l)}\,|\square\cap\Gamma^{\natural}_l|_l\\
&\le M^{-4}\sum_{l\le e}L^{-4(e-l)}\,|S\cap\Gamma^{\natural}_l|_l .
\end{split}
\end{equation}
Summing \eqref{9.2:eq-leaf-pile-4} over these cubes and using \eqref{9.2:eq-leaf-pile-5} gives the first inequality of \eqref{9.2:eq-leaf-pile-a}.

\emph{The second inequality.} By the definition of $\varepsilon_e$,
\begin{equation*}
\varepsilon_eC_L'=C_LK_0M^{-4}C_{99}\,e^{-p_0(g_e)}.
\end{equation*}
For $l\le e$ we have $L^{-4(e-l)}\le1$, and the terms with $l>e$ added on the right are nonnegative. Hence the middle member of \eqref{9.2:eq-leaf-pile-a} is at most
\begin{equation*}
K_0M^{-4}C_{99}\,e^{-p_0(g_e)}\sum_{l}|S\cap\Gamma^{\natural}_l|_l
=C_L^{-1}\,\varepsilon_eC_L'\sum_{l}|S\cap\Gamma^{\natural}_l|_l ,
\end{equation*}
and since $C_L\ge1$ this is at most the right member. The factor $C_L$ converts sup-norm bounds into per-cell value contributions. In the accounting of \eqref{9.2:eq-age-factor-majorants-d} this is the statement that the level-$e$ family of plain leaves has per-native-cell pile at most $\varepsilon_e$.

\emph{The sum over levels.} Leaves alive at the current time have level $e\le k\le K$. By the definition of $\varepsilon_e$, $\sum_e\varepsilon_e$ is $C_LK_0M^{-4}C_{99}/C_L'$ times $\sum_{e\le k}e^{-p_0(g_e)}$, where $p_0(g_e)=A_0(\log x_e)^{p_0}$.

We first bound $x_e$ from below. For $e<K$, hypothesis~(3) with $i=e$ and $j=K$ gives
\begin{equation*}
x_e\ge x_K+\lambda(K-e)-2c_2>X-2c_2+\lambda(K-e).
\end{equation*}
For $e=K$ the same lower bound is $x_K>X$ together with $c_2\ge0$. Hence
\begin{equation}\label{9.2:eq-leaf-pile-6}
x_e\ge X-2c_2+\lambda(K-e)\qquad\text{for every }e\le K.
\end{equation}
The degree function $p_0(\cdot)$ is increasing in $x=g^{-2}$. So \eqref{9.2:eq-leaf-pile-6} gives
\begin{equation*}
e^{-p_0(g_e)}\le\exp\bigl(-A_0[\log(X-2c_2+\lambda(K-e))]^{p_0}\bigr).
\end{equation*}
We re-index by $t:=K-e\ge0$. Distinct levels give distinct $t$, so $\sum_{e\le k}e^{-p_0(g_e)}\le\mathfrak{c}(X)$. This bound does not depend on the history, the time, $k$ or $K$, which proves \eqref{9.2:eq-leaf-pile-b}.

\emph{Convergence and monotonicity of the series.} The $t$-th term of $\mathfrak{c}(X)$ is positive and nonincreasing in $X$, since $A_0>0$ and the degree function is increasing. Since the fixed exponent of the degree function $p_0(\cdot)$ of Definition~\ref{1.5:def-domains} satisfies $p_0>1$, and $\lambda>0$, we have $A_0(\log y)^{p_0}\ge2\log y$ for all large $y$. Hence the $t$-th term is at most $(X-2c_2+\lambda t)^{-2}$ for all large $t$, and the series converges. Each term tends to $0$ as $X\uparrow\infty$. For $X$ above a fixed value $X_1$, each term is dominated by the corresponding term at $X_1$. By dominated convergence, $\mathfrak{c}(X)$ decreases to $0$.

\emph{Strip leaves.} For strip leaves, hypothesis~(4) supplies the constants $\varepsilon^{\mathrm{strip}}_e$ and $\varepsilon^{\mathrm{strip}}_{\mathrm{tot}}$, and nothing further is proved here. For plain leaves, by contrast, the sum of $e^{-p_0(g_e)}$ over all levels converges by itself, and no comparison of these factors with a power of the coupling is needed.
\end{proof}

\begin{theorem}[The carry lemma. Stated; proof incomplete at the step marked $(\ast)$]
\label{9.2:thm-carry}
Let $g\le\gamma_L'$, the ceiling on the coupling under which the allowance
\eqref{9.2:eq-age-factor-majorants-c} holds. Let
$\mathfrak{s},\tilde{\mathfrak{s}}\in\mathfrak{K}_k$ be two lists on one lattice and one
scaffold, with difference $\delta$, and let the production formulas be dialled on the per-use
polydisc $\mathbb{D}^{\mathrm{use}}_\delta$ of Definition~\ref{9.2:def-per-use-polydisc}.
Assume the following statements:
\begin{itemize}
\item the holomorphy and majorant statement on $\mathbb{D}^{\mathrm{use}}_\delta$ of
Definition~\ref{9.2:def-per-use-polydisc}, read at the allowance
\eqref{9.2:eq-age-factor-majorants-c}, together with the extension proposition and the budget
bound it contains;
\item the lemma on geometric forgetting of old regions, in the form stated in
Notation~\ref{9.2:not-age-factor-majorants}, at the enlarged entropy constant $\mathsf{c}'$ of
\eqref{9.2:eq-age-factor-majorants-c}, together with the standing conditions listed with this
lemma in that notation, including the condition named there explicitly as an input of its proof;
\item the lemmas of the correlation theorem, taken by statement, including both clauses of its
anchoring statement and the locality of its summation lemma;
\item the volume and filament bound;
\item the format \cite[(1.92)--(1.99)]{B-CMP122b} of the effective densities, taken by
statement;
\item for strip leaves, the relative-length bound for first-group leaves, with its pile
$\varepsilon^{\mathrm{strip}}$.
\end{itemize}
Let $\varepsilon_e$, $\varepsilon_{\mathrm{tot}}$ be the piles of
\eqref{9.2:eq-leaf-pile-a}--\eqref{9.2:eq-leaf-pile-b}, and let
$\varepsilon^{\mathrm{strip}}_e$, $\varepsilon^{\mathrm{strip}}_{\mathrm{tot}}$ be the
corresponding piles of strip leaves. Assume the one-sided condition
\begin{equation}\label{9.2:eq-carry-a}
\begin{gathered}
16\,\varepsilon^{\mathrm{all}}_{\mathrm{tot}}\le 1,\qquad
\varepsilon^{\mathrm{all}}_{\mathrm{tot}}:=\varepsilon_{\mathrm{tot}}
+\varepsilon^{\mathrm{strip}}_{\mathrm{tot}},\qquad
C_*:=\frac{4}{1-4\varepsilon^{\mathrm{all}}_{\mathrm{tot}}}\ \Bigl(\le\frac{16}{3}\Bigr),\\
\varepsilon^{(\mathrm{all})}_e:=\varepsilon_e+\varepsilon^{\mathrm{strip}}_e
\qquad\Bigl(\text{so that }\sum_e\varepsilon^{(\mathrm{all})}_e
\le\varepsilon^{\mathrm{all}}_{\mathrm{tot}}\Bigr).
\end{gathered}
\end{equation}
For a leaf-formula occurrence $\varphi$ of level $j$ and $z=(n,k')$ with $k'\le j$, and for a
level $e$, put
\begin{equation}\label{9.2:eq-carry-b}
\mathfrak{d}(\varphi;z):=\varphi-\varphi\lceil_{\mathcal{W}(\varphi,z)},\qquad
R_e(z):=\bigl[C_*\,\varepsilon^{(\mathrm{all})}_e\,q^{(e-n-1)_+}\bigr]^{-1}\ (\ge 3),
\end{equation}
with $\mathcal{W}(\varphi,z)$ as in Definition~\ref{9.2:def-own-witness-sets}; here
$\mathfrak{d}(\varphi;z)$ is a holomorphic function on $\mathbb{D}^{\mathrm{use}}_\delta$. Then the
following hold.
\begin{enumerate}
\item[(i)] (Occurrence classes.) For every first-group leaf $\mathbf{B}'^{(j)}(X')$ of
\cite[(1.98)]{B-CMP122b} and every $z=(n,k')$ with $k'\le j$, each use-copy of class $z$ inside
its formula is exactly one of the following.
\begin{enumerate}
\item[($\alpha$)] It is geo-witnessed ($n<k'$, its anchor in a live component of
$\Lambda_{n+1}^c$, that is, in the zone of a surviving component; this includes the zones of
live second-class and of arrested components, read through $\Lambda$); or it is a top-zone copy
($n=k'$) owned by a live-witnessed carrier (a $\mathbf{T}$-born boundary term, a frozen term
$\mathbf{C}$ or an outer leaf) of level $n$ or $n+1$, or owned by a live-witnessed
nested leaf-occurrence of level $e$ with $e\le n+1$.
\item[($\beta$)] It is a member of $\mathcal{W}(\mathbf{B}'^{(j)}(X'),z)$, see
\eqref{9.2:eq-own-witness-sets-a}: a use-copy of the own histories of the healed components, of
the top class, or one of the use-copies that lost their last witness at $\mathbf{R}_j$.
\item[($\gamma$)] It is a member of $\mathcal{W}(\psi,z)$ for a nested leaf-occurrence $\psi$
of level $e$, $n+2\le e<j$, live-witnessed at the present time (an older leaf anchored at a
survivor and copied by value).
\end{enumerate}
\item[(ii)] (Inner decay.) For every plain first-group leaf-formula occurrence $\varphi$ of level
$j$ and every $z=(n,k')$ with $k'\le j$, at every point of $\mathbb{D}^{\mathrm{use}}_\delta$,
\begin{equation}\label{9.2:eq-carry-c}
\sup\bigl|\mathfrak{d}(\varphi;z)\bigr|
=\sup\bigl|\varphi-\varphi\lceil_{\mathcal{W}(\varphi,z)}\bigr|
\le C_*\,\mathfrak{B}'_{X'}\,q^{(j-n-1)_+},
\end{equation}
with $\mathfrak{B}'_{X'}$ as in \eqref{9.2:eq-age-factor-majorants-b}. The same bound holds for
every second-group term $R'^{(j)}(X)$ of \cite[(1.98)]{B-CMP122b}, with $\mathfrak{B}_X$ in
place of $\mathfrak{B}'_{X'}$; for a strip leaf the same display holds with
$\mathfrak{B}'_{X'}$ formed with a relative length $d_{j,Z}(X')$, as in the relative-length bound
for first-group leaves, in place of $d_j(X')$, as an implication from the pile
$\varepsilon^{\mathrm{strip}}$ of that bound. The supremum is over
$\mathbb{D}^{\mathrm{use}}_\delta$, all other use-copies taking arbitrary admissible values; the
constant $C_*$ depends on no level, age, depth, $k$, $K$, torus or position.
\item[(iii)] (Transport to the reading level.) For every component $X$ of $Z_k$ (of any
history), every bounded $F$ and every $z=(n,k')$, with
$\mathsf{G}:=\mathbf{T}'^{\lambda}_k(X)F$,
\begin{equation}\label{9.2:eq-carry-d}
\mathfrak{B}_z:=\sup_{\mathbb{D}_\delta}\bigl|\mathsf{G}-\mathsf{G}\lceil_{\lambda_z=0}\bigr|
\le 2e\,(1+C_*\varepsilon^{\mathrm{all}}_{\mathrm{tot}})\,\mathfrak{E}_k\,
q^{(k-n-1)_+}\sup|F|,
\end{equation}
where the prefactor $e$ is Euler's number. Hence
\begin{equation}\label{9.2:eq-carry-1}
\begin{split}
\bigl|(\mathbf{T}'^{\tilde{\mathfrak{s}}}_k(X)-\mathbf{T}'^{\mathfrak{s}}_k(X))F\bigr|
&\le e\,(1+C_*\varepsilon^{\mathrm{all}}_{\mathrm{tot}})\,E_0^{-1}\,\mathfrak{E}_k
\sup|F|\cdot\Sigma^q_k[\delta]\\
&\le 2e\,E_0^{-1}\,e^{-2(1+\beta_0)^{-1}p_0(g_k)}\sup|F|\cdot\Sigma^q_k[\delta],
\end{split}
\end{equation}
which is the bound of clause (d) of the one-lattice Lipschitz step as stated there. The sum over
leaves enters through one further unit of the allowance; this is the allowance in the form
\eqref{9.2:eq-age-factor-majorants-c}.
\item[(iv)] The same holds inside clause (e) of the one-lattice Lipschitz step (its second-group
member as stated there, its first-group member resolved by zones), inside the proposition on the
tail above a zone (with its bound unchanged and its clause replaced by the condition that no
use-copy of alternative ($\alpha$) of clause (i) has its zone below $n$), inside items (3)--(4) of
the alignment proposition (for the aligned part at every depth), and for the first index
$h(\omega)$ of clause (f) of the one-lattice Lipschitz step, over the full history. These four
members are stated in the four statements that follow the proof of clauses (ii) and (iii).
\end{enumerate}
\end{theorem}

Clauses (i)--(iii) are proved as an implication from the statements listed in the hypotheses.

\begin{proof}
\emph{The constants.} By \eqref{9.2:eq-carry-a}, $4\varepsilon^{\mathrm{all}}_{\mathrm{tot}}\le
\frac14$, hence
\begin{equation*}
C_*\le\frac{4}{3/4}=\frac{16}{3},\qquad
C_*\varepsilon^{\mathrm{all}}_{\mathrm{tot}}
=\frac{4\varepsilon^{\mathrm{all}}_{\mathrm{tot}}}{1-4\varepsilon^{\mathrm{all}}_{\mathrm{tot}}}
\le\frac{1/4}{3/4}=\frac13 .
\end{equation*}
Thus $1+C_*\varepsilon^{\mathrm{all}}_{\mathrm{tot}}\le\frac43$. Since the piles are
non-negative, $\varepsilon^{(\mathrm{all})}_e\le\varepsilon^{\mathrm{all}}_{\mathrm{tot}}$, and
since $q\le\frac12$, we get
$C_*\varepsilon^{(\mathrm{all})}_e q^{(e-n-1)_+}\le\frac13$, that is $R_e(z)\ge3$, as stated
in \eqref{9.2:eq-carry-b}.

\emph{Clause (i).} This is the first part of Lemma~\ref{9.2:lem-witness-partition}, read inside
the formula of $\mathbf{B}'^{(j)}(X')$. Every class-$z$ use-copy there is either geo-witnessed or
lies in $\mathcal{W}(\psi,z)$ for exactly one live-witnessed object-occurrence $\psi$. The
geo-witnessed copies, together with those whose $\psi$ is a carrier (of level $n$ or $n+1$, by
Definition~\ref{9.2:def-own-witness-sets}) or a nested leaf of level at most $n+1$, form
($\alpha$); those with $\psi=\mathbf{B}'^{(j)}(X')$ form ($\beta$); the remaining $\psi$ are
nested leaves of level $e$ with $n+2\le e<j$, and the use-copies they own form ($\gamma$).

\emph{Clauses (ii) and (iii): the device.} Clause (ii) and the bound \eqref{9.2:eq-carry-d} of
clause (iii) are proved in the lemmas that follow this theorem; \eqref{9.2:eq-carry-1} is derived
from \eqref{9.2:eq-carry-d} below. The device is the following. At production, the sequence sums
of the healed components are split by an auxiliary dial $\mu$ at the age threshold fixed for this
split by the lemma on the split threshold of the healed components' sequence sums.
The nested occurrences that lose their witness are frozen by one auxiliary dial $\nu_e$ for each
witness level $e$, and are unfrozen by the second Schwarz-type disc lemma.

At the reading level $k$, the sequence sum of $X$ is split at the threshold $k-e$ for each
witness level $e$. The live-witnessed occurrences of that level are dialled by $\nu_e$ and
undialled by the first Schwarz-type disc lemma.

In both steps, holomorphy and the majorant on the auxiliary discs come from the holomorphy and
majorant statement on $\mathbb{D}^{\mathrm{use}}_\delta$ of
Definition~\ref{9.2:def-per-use-polydisc}, read at the allowance
\eqref{9.2:eq-age-factor-majorants-c}. It is applied through three inputs: the budget bound of
the extension proposition; Lemma~\ref{9.2:lem-leaf-pile}; and, by induction on the level, the
inner decay \eqref{9.2:eq-carry-c} at the levels $e<j$ of the nested leaves.

The carriers are the $\mathbf{T}$-born boundary terms $\mathbf{B}^{(\cdot)}$, the frozen terms
$\mathbf{C}$ and the outer leaves. They are recomputed from the dialled values, and therefore stay
inside their own majorants. Each of them is bounded on the whole dial, so the
Schwarz-type bounds apply through them. This suffices because the own-witness variation
$\mathfrak{d}(\varphi;z)$ of each leaf is dialled where the leaf is produced.

\emph{From \eqref{9.2:eq-carry-d} to \eqref{9.2:eq-carry-1}.} The third Schwarz-type disc lemma
is applied in the variables $\lambda_z$ on discs of radius $E_1/[\![\delta]\!]_n$, and for each
$n$ with the $k-n+1$ values of $k'$. It turns \eqref{9.2:eq-carry-d} into the first inequality of
\eqref{9.2:eq-carry-1}. There the factor $2$ of \eqref{9.2:eq-carry-d} is absorbed by
$E_1\ge2E_0$, and
\begin{equation*}
\sum_{n\le k}(k-n+1)\,q^{(k-n-1)_+}[\![\delta]\!]_n=\Sigma^q_k[\delta]
\end{equation*}
by definition. The second inequality of \eqref{9.2:eq-carry-1} follows from
$1+C_*\varepsilon^{\mathrm{all}}_{\mathrm{tot}}\le\frac43\le2$ and
$\mathfrak{E}_k=e^{-2(1+\beta_0)^{-1}p_0(g_k)}$.

$(\ast)$ \emph{Clause (iv).} Its four members are the analogues of the transport of clause (iii)
in the following places:
\begin{itemize}
\item clause (e) of the one-lattice Lipschitz step, with the second-group member as stated there
and the first-group member resolved by zones;
\item the proposition on the tail above a zone, with its bound unchanged and its clause replaced
by the condition that no use-copy of alternative ($\alpha$) of clause (i) has its zone below $n$;
\item items (3)--(4) of the alignment proposition, for the aligned part at every depth;
\item the first index $h(\omega)$ of clause (f) of the one-lattice Lipschitz step, over
the full history.
\end{itemize}
At this step the proof is incomplete: the four members are stated, in the form in which they are
used, in the four statements that follow the proof of clauses (ii) and (iii); they are not proved
there.
\end{proof}

Throughout what follows, $\mathbb{D}^{\mathrm{use}}_\delta$ is the per-use polydisc of Definition~\ref{9.2:def-per-use-polydisc}, and $c\lceil_S$ and $\psi\lceil_S$ are as there. The constant $C_*$ and the condition on the piles are those of \eqref{9.2:eq-carry-a}. The own-witness variation $\mathfrak{d}(\varphi;z)$ and the radius $R_e(z)$ are those of \eqref{9.2:eq-carry-b}. The inductive bound is \eqref{9.2:eq-carry-c}. All three displays belong to the carry lemma (Theorem~\ref{9.2:thm-carry}), whose proof is incomplete at one step. For an object-occurrence $\psi$ and a class $z=(n,k')$, $\mathcal{W}(\psi,z)$ is the own-witness set of class $z$ of $\psi$ used in \eqref{9.2:eq-carry-b}.

\emph{The interpolated value.} Let $\psi$ be a leaf-formula occurrence of level $e$, let $z=(n,k')$ with $k'\le e$, let $c\in\mathbb{D}^{\mathrm{use}}_\delta$ and let $\nu\in\mathbb{C}$. We put
\begin{equation}\label{9.2:eq-interpolation-budget-1}
\psi^{[\nu]}(c):=\psi\bigl(c\lceil_{\mathcal{W}(\psi,z)}\bigr)+\nu\,\mathfrak{d}(\psi;z)(c)
=(1-\nu)\,\psi\lceil_{\mathcal{W}(\psi,z)}(c)+\nu\,\psi(c).
\end{equation}
The second equality is the definition $\mathfrak{d}(\psi;z)=\psi-\psi\lceil_{\mathcal{W}(\psi,z)}$ of \eqref{9.2:eq-carry-b}, inserted and regrouped. At $\nu=1$ the interpolated value is the true value $\psi(c)$. At $\nu=0$ it is the value with the own-witness content of class $z$ switched off.

\emph{Dialling.} Let $\mathcal{S}$ be a set of leaf-formula occurrences of level $e$ and let $\nu_e\in\mathbb{C}$. To \emph{dial $\mathcal{S}$ by $\nu_e$} means the following. In every production formula of the dialled expansion that reads the value of some $\psi\in\mathcal{S}$ as a stored term, that value is replaced by $\psi^{[\nu_e]}$. Every object produced later is recomputed from the replaced values by its own production formula, so that for every label the objects are again those the production formulas give.

By \emph{the holomorphy-and-majorant proposition} we mean the item of Notation~\ref{9.2:not-age-factor-majorants} giving \eqref{9.2:eq-age-factor-majorants-b}: every derived object of the dialled production is holomorphic on $\mathbb{D}^{\mathrm{use}}_\delta$ and obeys there its majorant, $\mathfrak{B}'_{X'}$ for a first-group term and $\mathfrak{B}_X$ for a second-group term (both of \eqref{9.2:eq-age-factor-majorants-b}). It holds under the enlarged allowance \eqref{9.2:eq-age-factor-majorants-c}. For sub-sums of $\mathbf{T}'_k(X)F$ of age at least $a$ we use in addition the tail-summed age bound \eqref{9.2:eq-age-factor-majorants-a}, the consequence of geometric forgetting of old regions stated in the same Notation.

\emph{The age split.} Let $X$ be a component, $k$ a level, $a\ge0$ and $F$ a bounded function on which $\mathbf{T}'_k(X)$ acts, as in \eqref{9.2:eq-age-factor-majorants-a}. Write $\mathbf{T}'_{k,<a}(X)$ for the part of the sum over admissible sequences in \cite[(1.71)]{B-CMP122b} with age below $a$, and $\mathbf{T}'_{k,\ge a}(X)$ for the part with age at least $a$. The age split of $X$ replaces the full sum $\mathbf{T}'_k(X)=\mathbf{T}'_{k,<a}(X)+\mathbf{T}'_{k,\ge a}(X)$ by
\begin{equation}\label{9.2:eq-interpolation-budget-2}
\mathbf{T}'_{k,<a}(X)+\mu\,\mathbf{T}'_{k,\ge a}(X),\qquad |\mu|\le q^{-a},
\end{equation}
where $\mu$ is an auxiliary complex variable.

\begin{lemma}[Inner decay: the auxiliary interpolation and its budget]
\label{9.2:prf-interpolation-budget}
Let $e$ be a level and $z=(n,k')$ with $k'\le e$. Assume the pile condition \eqref{9.2:eq-carry-a}. Let $\mathcal{S}$ be a set of leaf-formula occurrences of level $e$, each of them plain or strip in the sense of Notation~\ref{9.2:not-age-factor-majorants}. Assume that every member of $\mathcal{S}$ satisfies the inductive bound \eqref{9.2:eq-carry-c} at level $e$ at every point of $\mathbb{D}^{\mathrm{use}}_\delta$. Dial $\mathcal{S}$ by $\nu_e$, $|\nu_e|\le R_e(z)$. Then the following hold.
\begin{enumerate}
\item[(i)] \emph{Budget.} At every site that reads members of $\mathcal{S}$ directly, the replaced family $\psi\mapsto\nu_e\,\mathfrak{d}(\psi;z)$, $\psi\in\mathcal{S}$, obeys the first line of \eqref{9.2:eq-interpolation-budget-a}. Its per-native-cell pile is at most $C_L'$, i.e.\ $\pi\le1$ in the budget rule \eqref{9.2:eq-age-factor-majorants-d}, hence inside the allowance \eqref{9.2:eq-age-factor-majorants-c}. The other summand $\psi\lceil_{\mathcal{W}(\psi,z)}$ of \eqref{9.2:eq-interpolation-budget-1} is a point value of $\psi$ on $\mathbb{D}^{\mathrm{use}}_\delta$. It therefore lies in the ball in which the holomorphy-and-majorant proposition admits values of stored terms.
\item[(ii)] \emph{Propagation at no cost.} Let $\xi$ be an object produced later that reads members of $\mathcal{S}$, directly or through intermediate objects. After recomputation, $\xi$ is a holomorphic function of $\nu_e$ on $|\nu_e|\le R_e(z)$. There it obeys the majorant $\mathfrak{B}_\xi$ that the holomorphy-and-majorant proposition assigns to it (second line of \eqref{9.2:eq-interpolation-budget-a}). No amplification factor accrues along chains of carriers, in the sense of clause~(i) of the carry lemma (Theorem~\ref{9.2:thm-carry}, whose proof is incomplete at one step), of any length.
\item[(iii)] \emph{Conclusion.} Assume that at any moment only one auxiliary variable (a dial $\nu_e$ as above, or the variable $\mu$ of an age split \eqref{9.2:eq-interpolation-budget-2}) ranges over a disc, and that all other interpolations are at $\nu\in\{0,1\}$. Let $\mathfrak{f}$ be any function to which the Schwarz-type disc estimates are applied in the inner-decay argument. Then $\mathfrak{f}$ is holomorphic on the disc of the active variable ($|\nu_e|\le R_e(z)$, respectively $|\mu|\le q^{-a}$) and is bounded there by the majorant $\mathfrak{B}_{\mathfrak{f}}$ of the object it interpolates (third and fourth lines of \eqref{9.2:eq-interpolation-budget-a}). That majorant is $\mathfrak{B}'_{X'}$ for a first-group term, $\mathfrak{B}_X$ for a second-group term, and $\mathrm{e}\,\mathfrak{E}_kq^a\sup|F|$ for an age-$\ge a$ sub-sum of $\mathbf{T}'_k(X)F$, $F$ bounded, where $\mathrm{e}=\exp(1)$.
\end{enumerate}
\begin{equation}\label{9.2:eq-interpolation-budget-a}
\begin{gathered}
|\nu_e\,\mathfrak{d}(\psi;z)|\le R_e(z)\,C_*\,q^{(e-n-1)_+}\,\mathfrak{B}_\psi
=\mathfrak{B}_\psi/\varepsilon^{(\mathrm{all})}_e\qquad(\psi\in\mathcal{S}),\\
\sup_{|\nu_e|\le R_e(z)}|\xi|\le\mathfrak{B}_\xi,\\
\sup_{|\nu_e|\le R_e(z)}|\mathfrak{f}|\le\mathfrak{B}_{\mathfrak{f}}\quad\text{respectively}\quad
\sup_{|\mu|\le q^{-a}}|\mathfrak{f}|\le\mathfrak{B}_{\mathfrak{f}},\\
\mathfrak{B}_{\mathfrak{f}}\in\bigl\{\mathfrak{B}'_{X'},\ \mathfrak{B}_X,\ \mathrm{e}\,\mathfrak{E}_k\,q^a\sup|F|\bigr\}.
\end{gathered}
\end{equation}
\end{lemma}

\begin{proof}
\emph{Part (i).} Let $\psi\in\mathcal{S}$.
\begin{itemize}
\item By the inductive bound \eqref{9.2:eq-carry-c} for $\psi$ at level $e$, and since $|\nu_e|\le R_e(z)$, we get the inequality in the first line of \eqref{9.2:eq-interpolation-budget-a}.
\item By the definition of $R_e(z)$ in \eqref{9.2:eq-carry-b}, namely $R_e(z)=[C_*\varepsilon^{(\mathrm{all})}_eq^{(e-n-1)_+}]^{-1}$, the product $R_e(z)C_*q^{(e-n-1)_+}$ equals $1/\varepsilon^{(\mathrm{all})}_e$. This gives the equality in the first line of \eqref{9.2:eq-interpolation-budget-a}.
\item For the piles we use Lemma~\ref{9.2:lem-leaf-pile} applied to the sub-family of plain members of $\mathcal{S}$, which is contained in $\mathcal{L}_e$. For the strip members we use the relative bound for strip leaves. The majorants $\mathfrak{B}_\psi$ have per-native-cell pile at most $\varepsilon_e$ in units of $C_L'$. Since the values are at most $\mathfrak{B}_\psi/\varepsilon^{(\mathrm{all})}_e$ by the first line of \eqref{9.2:eq-interpolation-budget-a}, their pile is at most $(1/\varepsilon_e)\cdot\varepsilon_e=1$ in units of $C_L'$.
\item By the budget rule \eqref{9.2:eq-age-factor-majorants-d}, for a region $Y$, a family of stored terms with per-native-cell pile at most $\pi C_L'$ on the reading domain adds at most $\pi C_L'$ per unit of $\sum_l|\Gamma^{\natural}_l\cap Y|$ to the constant of the weighted inductive bound entering the budget rule \eqref{9.2:eq-age-factor-majorants-d}, and at most $C_A\mathsf{w}_{\mathrm{site}}^{-1}\pi C_L'$ to the activity of the site.
\item With $\pi\le1$ this is the one further unit of the per-cell allowance and of the bracket budget that the enlarged allowance \eqref{9.2:eq-age-factor-majorants-c} provides over the unenlarged one, through the replacement of $C_L'$ by $2C_L'$.
\item The other summand is $\psi\lceil_{\mathcal{W}(\psi,z)}(c)=\psi(c\lceil_{\mathcal{W}(\psi,z)})$. By Definition~\ref{9.2:def-per-use-polydisc}, the point $c\lceil_{\mathcal{W}(\psi,z)}$ lies again in $\mathbb{D}^{\mathrm{use}}_\delta$. So this summand is a point value of $\psi$ on $\mathbb{D}^{\mathrm{use}}_\delta$, and it lies in the ball in which the holomorphy-and-majorant proposition admits stored values.
\end{itemize}

\emph{Part (ii).} We follow the induction along the production steps by which the holomorphy-and-majorant proposition is proved, now with $\nu_e$ as an additional variable on the disc $|\nu_e|\le R_e(z)$.
\begin{itemize}
\item At the members of $\mathcal{S}$ themselves, each replaced value $\psi^{[\nu_e]}$ is affine in $\nu_e$ by \eqref{9.2:eq-interpolation-budget-1}. Its coefficients are holomorphic on $\mathbb{D}^{\mathrm{use}}_\delta$.
\item At the production of $\xi$, two kinds of input enter. The members of $\mathcal{S}$ read directly cost the one unit of part~(i). The intermediate objects enter, as stored terms, only through their sup-norms.
\item By the induction, each intermediate object obeys its own majorant of the holomorphy-and-majorant proposition on the disc. Its sup-norm on the disc is therefore bounded by the same majorant that enters the induction without the dial, so these stored terms are charged exactly as before.
\item Hence $\xi$ is holomorphic in $\nu_e$ on $|\nu_e|\le R_e(z)$ and obeys $\mathfrak{B}_\xi$ there. This is the second line of \eqref{9.2:eq-interpolation-budget-a}.
\item Every link of a chain of carriers is such an intermediate object and enters the next link through its sup-norm. So no factor is picked up from one link to the next, whatever the length of the chain.
\end{itemize}

\emph{Part (iii).} First consider the variables that do not range over a disc. Every interpolation other than the single active one is at $\nu\in\{0,1\}$. By \eqref{9.2:eq-interpolation-budget-1}, $\psi^{[1]}=\psi(c)$ and $\psi^{[0]}=\psi(c\lceil_{\mathcal{W}(\psi,z)})$. Both are values at genuine points of $\mathbb{D}^{\mathrm{use}}_\delta$ and cost nothing.

If the active variable is a dial $\nu_e$, parts~(i) and~(ii) apply. Under \eqref{9.2:eq-carry-a} the disc $|\nu_e|\le R_e(z)$ contains $0$ and $1$, since $R_e(z)\ge 3$.

If the active variable is the variable $\mu$ of an age split, the sum over admissible sequences of one component $X$ is replaced by \eqref{9.2:eq-interpolation-budget-2}. We use this split as follows.
\begin{itemize}
\item Geometric forgetting of old regions, in the form stated in Notation~\ref{9.2:not-age-factor-majorants}, keeps the bound \cite[(1.92)]{B-CMP122b} for \eqref{9.2:eq-interpolation-budget-2}. Moreover, since the sum over sequences in the definition \cite[(1.71)]{B-CMP122b} acts also on the last exponential in \cite[(1.72)]{B-CMP122b}, that last exponential factor of \cite[(1.72)]{B-CMP122b} is split together with the sequence sum.
\item The sequences of age below $a$ form a sub-family of age $\ge0$, and the sequences of age at least $a$ form a sub-family of age $\ge a$. The tail-summed age bound \eqref{9.2:eq-age-factor-majorants-a}, which is the consequence of geometric forgetting of old regions stated there, applies to each sub-family. It gives
\begin{equation}\label{9.2:eq-interpolation-budget-3}
\begin{aligned}
|\mathbf{T}'_{k,<a}(X)F|+q^{-a}|\mathbf{T}'_{k,\ge a}(X)F|
&\le \mathrm{e}\,\mathfrak{E}_k\sup|F|+q^{-a}\,\mathrm{e}\,\mathfrak{E}_kq^a\sup|F|\\
&=2\,\mathrm{e}\,\mathfrak{E}_k\sup|F|.
\end{aligned}
\end{equation}
\item So the split doubles the input $\mathrm{e}\,\mathfrak{E}_k$ of that component to \cite[(1.92)]{B-CMP122b}.
\item This factor $2$ is covered by the replacement of the summand $2$ by $3$ in the constant of \cite[\S1]{B-CMP122b} that carries this summand. That replacement is already contained in the allowance \eqref{9.2:eq-age-factor-majorants-c} and covers one split at a time. Accordingly only one $\mu$ is active at a time.
\end{itemize}

In every case, therefore, the function $\mathfrak{f}$ is produced on a polydisc consisting of $\mathbb{D}^{\mathrm{use}}_\delta$ enlarged by at most one $\mu$-disc, covered by the summand $3$, and at most one $\nu$-disc, covered by the allowance \eqref{9.2:eq-age-factor-majorants-c}. This is the polydisc on which the majorants stated in Notation~\ref{9.2:not-age-factor-majorants} hold:
\begin{itemize}
\item $\mathfrak{B}'_{X'}$ of \eqref{9.2:eq-age-factor-majorants-b} for the first-group term with domain $X'$;
\item $\mathfrak{B}_X$ of \eqref{9.2:eq-age-factor-majorants-b} for the second-group term with domain $X$;
\item $\mathrm{e}\,\mathfrak{E}_kq^a\sup|F|$ of \eqref{9.2:eq-age-factor-majorants-a} for an age-$\ge a$ sub-sum of $\mathbf{T}'_k(X)F$.
\end{itemize}
Holomorphy on these discs follows from part~(ii) for $\nu_e$. For $\mu$ it follows in the same way as in part~(ii): the split \eqref{9.2:eq-interpolation-budget-2} is affine in $\mu$ at the component $X$, and every object produced later from it is holomorphic on $|\mu|\le q^{-a}$ by the induction of the holomorphy-and-majorant proposition, the doubled input being covered by the summand $3$. This is the third and fourth lines of \eqref{9.2:eq-interpolation-budget-a}.
\end{proof}

\emph{The induction hypothesis.} Lemma~\ref{9.2:prf-interpolation-budget} is used in the proof of clause~(ii) of the carry lemma (Theorem~\ref{9.2:thm-carry}, whose proof is incomplete at one step). That proof proceeds by induction on the level $j$. Its induction hypothesis at level $j$ is the bound \eqref{9.2:eq-carry-c} for every leaf-formula occurrence of every level $e<j$, at every point of $\mathbb{D}^{\mathrm{use}}_\delta$. For plain occurrences it is what the induction has established at the levels $e<j$; for strip occurrences it is an input of the carry lemma taken by statement.

For the base of the induction, take the lowest level carrying a leaf. There the part labelled (B) of the own-witness set $\mathcal{W}(\psi,z)$ is empty; the label refers to the decomposition of $\mathcal{W}(\psi,z)$ into parts given in the definition of the own-witness sets, on which the definition of $\mathfrak{d}(\psi;z)$ in \eqref{9.2:eq-carry-b} rests. Together with the estimate of the own-witness variation $\mathfrak{d}(\psi;z)$ carried out in the proof of clause~(ii) of the carry lemma (Theorem~\ref{9.2:thm-carry}, whose proof is incomplete at one step), this gives \eqref{9.2:eq-carry-c} with the constant $4$. This constant is at most $C_*$, because $C_*=4/(1-4\varepsilon^{\mathrm{all}}_{\mathrm{tot}})\ge4$ under \eqref{9.2:eq-carry-a}.

There is no circularity. The leaves of one level $j$ are produced simultaneously by the one operation $\mathbf{R}_j$ of step $j$, and none of them reads another.

\begin{lemma}[Inner decay: the lump term by one age split]\label{9.2:prf-lump-term}
Let $\varphi=\mathbf{B}'^{(j)}(X')$ be a first-group term of \cite[(1.98)]{B-CMP122b}, born at
$\mathbf{R}_j$ with healed components $X_1,\dots,X_p$. The same letters are used for a
second-group term $R'^{(j)}(X)$. Let $z=(n,k')$ with $k'\le j$, and put
\begin{equation*}
a:=(j-n-1)_+ .
\end{equation*}
For each level $e$ let $\mathcal{S}_e$ be the set of nested leaf-occurrences of level $e$ in
the formula of $\varphi$ all of whose witnessing components lie in $\bigcup_iX_i$. Put
$\mathcal{W}^{(e)}:=\bigcup\{\mathcal{W}(\psi,z):\psi\in\mathcal{S}_e\}$. Let
$\mathcal{W}^{\mathrm{lump}}$ be the union of three kinds of set:
\begin{itemize}
\item the parts \textup{(A1)}, \textup{(A2)} and \textup{(G)} of $\mathcal{W}(\varphi,z)$ in
\eqref{9.2:eq-own-witness-sets-a};
\item the sets $\mathcal{W}(\psi,z)$ of the nested carriers $\psi$, of any level, that lose
their last witness at $\mathbf{R}_j$;
\item the sets $\mathcal{W}(\psi,z)$ of the nested leaf-occurrences $\psi$ of level
$e\le n+1$ that lose their last witness at $\mathbf{R}_j$.
\end{itemize}
The own-witness set of $\varphi$ is decomposed as
\begin{equation}\label{9.2:eq-lump-term-1}
\mathcal{W}(\varphi,z)=\mathcal{W}^{\mathrm{lump}}\sqcup\bigsqcup_{e=n+2}^{j-1}\mathcal{W}^{(e)} .
\end{equation}
For an admissible sequence $\omega_i$ of $X_i$ write
$\operatorname{age}_j(\omega_i):=j-j(X_i)(\omega_i)$, as in Lemma~\ref{9.2:lem-presence-healed}.
\begin{enumerate}
\item[\textup{(i)}] Every member of $\mathcal{W}^{\mathrm{lump}}$ other than the members of
\textup{(A2)} is absent
from each summand of the cluster of $\varphi$ in which $\operatorname{age}_j(\omega_i)<a$ for all
$i$.
\item[\textup{(ii)}] Let $c$ be a point of $\mathbb{D}^{\mathrm{use}}_\delta$ \textup{(}in the
application, the ambient point of the enclosing proof\textup{)}. Put
\begin{equation*}
c':=c\lceil_{\bigsqcup_{e=n+2}^{j-1}\mathcal{W}^{(e)}} ,
\end{equation*}
the point $c$ with the dials of the groups $\mathcal{W}^{(e)}$, $n+2\le e\le j-1$, of
\eqref{9.2:eq-lump-term-1} set to $0$.
In the cluster of $\varphi$ replace every $\mathbf{T}'_j(X_i)$ by
\begin{equation}\label{9.2:eq-lump-term-2}
\mathbf{T}'_{j,<a}(X_i)+\mu\,\mathbf{T}'_{j,\ge a}(X_i),\qquad |\mu|\le q^{-a},
\end{equation}
and write $\varphi(\,\cdot\,;\mu)$ for the resulting function. Put
\begin{equation}\label{9.2:eq-lump-term-3}
\Phi(\mu):=\varphi(c';\mu)-\varphi\bigl(c'\lceil_{\mathcal{W}^{\mathrm{lump}}};\mu\bigr).
\end{equation}
Then, with $\mathfrak{B}'_{X'}$ the first-group majorant of
\eqref{9.2:eq-age-factor-majorants-b} \textup{(}for a second-group term $R'^{(j)}(X)$ read
$\mathfrak{B}_X$ in place of $\mathfrak{B}'_{X'}$ throughout\textup{)},
\begin{equation}\label{9.2:eq-lump-term-a}
|\Phi(1)|\le 2\,\mathfrak{B}'_{X'}\,q^{a}.
\end{equation}
\end{enumerate}
\end{lemma}

\begin{proof}
\emph{Proof of (i).} For $a=0$ there is nothing to show, since ages are non-negative; assume
$a\ge1$, so that $a=j-n-1$. We treat the members of $\mathcal{W}^{\mathrm{lump}}$ kind by kind. In
each case we show the following: if the member is present in the summand indexed by
$(\omega_1,\dots,\omega_p)$, then some $X_i$ has $\operatorname{age}_j(\omega_i)\ge a$. Such a
summand is therefore not one in which all ages are below $a$.

\emph{The copies in (A1) and in (G).} Let $\upsilon$ be a copy in (A1) of class $(n,k')$.
Lemma~\ref{9.2:lem-presence-healed}\,(P1), applied at $\mathbf{R}_j$, says that $\upsilon$ is present
only if $\operatorname{age}_j(\omega_i)\ge j-n-1$. Here $X_i$ is the component that contains the
anchor of $\upsilon$. Now let $\upsilon$ be a copy in (G), with $C_\upsilon\subset X_i$.
Lemma~\ref{9.2:lem-presence-healed}\,(P2) gives the same lower bound
$\operatorname{age}_j(\omega_i)\ge j-n-1$ for that $X_i$. Since $a=j-n-1$, both bounds read
$\operatorname{age}_j(\omega_i)\ge a$.

\emph{The carriers.} Let $\psi$ be a nested carrier of level $e$ that loses its last witness at
$\mathbf{R}_j$. Let that witness be a component $C\subset X_{i_0}$. By
Definition~\ref{9.2:def-own-witness-sets}, the own-witness set $\mathcal{W}(\psi,z)$ of a carrier
consists of top-zone use-copies. So it is non-empty only for classes $(n,k')$ with $k'=n$ and
$n\in\{e-1,e\}$. In particular $e\le n+1$. A member of $\mathcal{W}(\psi,z)$ is present only
where $\psi$ is present. Lemma~\ref{9.2:lem-presence-healed}\,(P3) gives
\begin{equation*}
\operatorname{age}_j(\omega_{i_0})\ge j-e\ge j-n-1 ,
\end{equation*}
and hence $\operatorname{age}_j(\omega_{i_0})\ge a$.

\emph{The leaves of level at most $n+1$.} Let $\psi$ be a nested leaf-occurrence of level
$e\le n+1$ that loses its last witness at $\mathbf{R}_j$, in a component contained in
$X_{i_0}$. Lemma~\ref{9.2:lem-presence-healed}\,(P3) again gives
\begin{equation*}
\operatorname{age}_j(\omega_{i_0})\ge j-e\ge j-n-1 ,
\end{equation*}
which is $a$.

\emph{The part (A2).} By \eqref{9.2:eq-own-witness-sets-a}, the part (A2) consists of the direct
use-copies of class $z$ with $k'=n=j$. It is therefore non-empty only if $n=k'=j$. In that
case $a=(j-j-1)_+=0$. This proves (i).

\emph{Proof of (ii).} We first check that the two points in \eqref{9.2:eq-lump-term-3} are
admissible. Setting the dials of a set of use-copies to $0$ leaves a point of
$\mathbb{D}^{\mathrm{use}}_\delta$. So $c'$ and $c'\lceil_{\mathcal{W}^{\mathrm{lump}}}$ are
points of $\mathbb{D}^{\mathrm{use}}_\delta$. Each of the two terms in \eqref{9.2:eq-lump-term-3}
is therefore the term of \cite[(1.98)]{B-CMP122b} belonging to the cluster of $\varphi$,
evaluated at a point of $\mathbb{D}^{\mathrm{use}}_\delta$, with a single $\mu$-split
\eqref{9.2:eq-lump-term-2}.

We next bound $\Phi$ on its disc. By the auxiliary interpolation and its budget, established
earlier in the inner-decay argument, each of the two
terms of \eqref{9.2:eq-lump-term-3} is holomorphic in $\mu$ on $|\mu|\le q^{-a}$ and bounded
there by the majorant $\mathfrak{B}'_{X'}$ of the term it interpolates, the doubling of each
$X_i$'s input under \eqref{9.2:eq-lump-term-2} being absorbed in that majorant there.
Since $\Phi$ is the difference of these two terms, $\Phi$ is holomorphic on $|\mu|\le q^{-a}$ and
\begin{equation}\label{9.2:eq-lump-term-4}
|\Phi(\mu)|\le 2\,\mathfrak{B}'_{X'}\qquad (|\mu|\le q^{-a}).
\end{equation}

At $\mu=1$ the replacement \eqref{9.2:eq-lump-term-2} gives back
$\mathbf{T}'_{j,<a}(X_i)+\mathbf{T}'_{j,\ge a}(X_i)=\mathbf{T}'_j(X_i)$. Hence
\begin{equation*}
\Phi(1)=\varphi(c')-\varphi\bigl(c'\lceil_{\mathcal{W}^{\mathrm{lump}}}\bigr).
\end{equation*}

\emph{The case $a\ge1$.} At $\mu=0$ each $\mathbf{T}'_j(X_i)$ is replaced by
$\mathbf{T}'_{j,<a}(X_i)$. So only the summands with $\operatorname{age}_j(\omega_i)<a$ for all $i$
remain. Since $a\ge1$ forces $n\le j-2$, the part (A2) is empty. By (i), no member of
$\mathcal{W}^{\mathrm{lump}}$ is present in the remaining summands.

By the definition of $\lceil$, the points $c'$ and $c'\lceil_{\mathcal{W}^{\mathrm{lump}}}$ differ
only in the dials indexed by $\mathcal{W}^{\mathrm{lump}}$. Since the remaining summands contain no
use-copy of $\mathcal{W}^{\mathrm{lump}}$, they do not read those dials. Hence the two terms of
\eqref{9.2:eq-lump-term-3} coincide at $\mu=0$, and
\begin{equation*}
\Phi(0)=0 .
\end{equation*}
The passage from $c$ to $c'$ matters here: with the groups $\mathcal{W}^{(e)}$,
$n+2\le e\le j-1$, switched off, the only own-witness dials of $\varphi$ still live at $c'$ are,
by \eqref{9.2:eq-lump-term-1}, those of $\mathcal{W}^{\mathrm{lump}}$, and by (i) these occur only
in summands in which $\operatorname{age}_j(\omega_i)\ge a$ for some $i$.

Since $q\le\frac12$ and $a\ge1$, the radius of the disc satisfies $q^{-a}\ge2\ge1$; in particular
the point $\mu=1$ lies in the disc. We apply the Schwarz-type disc lemma for a function
vanishing at the centre to $\Phi$ on the disc
$|\mu|\le q^{-a}$: $\Phi$ is holomorphic there, bounded by $2\mathfrak{B}'_{X'}$ by
\eqref{9.2:eq-lump-term-4}, and $\Phi(0)=0$. It gives
\begin{equation*}
|\Phi(\mu)|\le 2\,\mathfrak{B}'_{X'}\,\frac{|\mu|}{q^{-a}}=2\,\mathfrak{B}'_{X'}\,|\mu|\,q^{a}
\qquad (|\mu|\le q^{-a}),
\end{equation*}
and at $\mu=1$ this is
\begin{equation*}
|\Phi(1)|\le 2\,\mathfrak{B}'_{X'}\,q^{a}.
\end{equation*}

\emph{The case $a=0$.} The disc is then $|\mu|\le q^0=1$ and contains $\mu=1$. Hence
\eqref{9.2:eq-lump-term-4} gives
\begin{equation*}
|\Phi(1)|\le 2\,\mathfrak{B}'_{X'}=2\,\mathfrak{B}'_{X'}\,q^{a}.
\end{equation*}

In both cases \eqref{9.2:eq-lump-term-a} holds.
\end{proof}

\begin{proof}[Proof of clause (ii) of the carry lemma, continued: inner decay, one pair of dials per witness level]
We keep the notation of the preceding parts of this proof of the carry lemma
(Theorem~\ref{9.2:thm-carry}). Thus
$\varphi=\mathbf{B}'^{(j)}(X')$ is a first-group leaf formed at $\mathbf{R}_j$ with healed components
$X_1,\dots,X_p$ (for a second-group term $R'^{(j)}(X)$ the same letters are used), $z=(n,k')$ with
$k'\le j$, and $a:=(j-n-1)_+$. The own-witness set $\mathcal{W}(\varphi,z)$ is the disjoint union of the
lump $\mathcal{W}^{\mathrm{lump}}$ and the high groups $\mathcal{W}^{(e)}$, $n+2\le e\le j-1$, where
$\mathcal{W}^{(e)}$ is the union of the sets $\mathcal{W}(\psi,z)$ over the set $\mathcal{S}_e$ of nested
leaf-occurrences $\psi$ of level $e$ all of whose witnessing components lie in $\bigcup_{i'}X_{i'}$.
The point $c$ is the ambient point of $\mathbb{D}^{\mathrm{use}}_\delta$, and
$c':=c\lceil_{\bigsqcup_{e=n+2}^{j-1}\mathcal{W}^{(e)}}$ is the point with all high groups switched
off, as in the lump estimate \eqref{9.2:eq-lump-term-a}.

\emph{The high groups.} Let $n+2\le e_1<e_2<\dots<e_r\le j-1$ be the levels $e$ with
$\mathcal{S}_e\neq\emptyset$; then $r\le j-n-2$, but $r$ enters no constant below. For $1\le i\le r$ let
\begin{equation*}
c_i:=c\lceil_{\mathcal{W}^{(e_{i+1})}\sqcup\dots\sqcup\mathcal{W}^{(e_r)}},
\end{equation*}
the point at which the groups above $e_i$ are switched off and the groups of level at most $e_i$ are
left on. Then $c_r=c$; $c_1\lceil_{\mathcal{W}^{(e_1)}}=c'$, all high groups being off; and, unwinding
the definition, $c_i\lceil_{\mathcal{W}^{(e_i)}}=c_{i-1}$ for $2\le i\le r$. Hence the difference
telescopes:
\begin{equation}\label{9.2:eq-witness-level-dials-1}
\varphi(c)-\varphi(c')=\sum_{i=1}^{r}\bigl[\varphi(c_i)-\varphi(c_i\lceil_{\mathcal{W}^{(e_i)}})\bigr],
\end{equation}
the sum being empty, and $c=c'$, when $r=0$.

Fix $i$ and write $e:=e_i$. Let $G_e(\mu,\nu)$ be the term $\varphi$ of its cluster, that is, the
formula of $\varphi$ itself, evaluated at the ambient point $c_i$, with two modifications:
\begin{itemize}
\item[(the $\mu$-split)] every $\mathbf{T}'_j(X_{i'})$ in the cluster is split as
$\mathbf{T}'_{j,<j-e}(X_{i'})+\mu\,\mathbf{T}'_{j,\ge j-e}(X_{i'})$, with $|\mu|\le q^{-(j-e)}$;
since $e\le j-1$ and $q\le\frac12$, this radius is at least $2$;
\item[(the $\nu$-interpolation)] the set $\mathcal{S}_e$ is dialled by $\nu$ through the auxiliary
interpolation $\psi\mapsto\psi^{[\nu]}$, whose budget is \eqref{9.2:eq-interpolation-budget-a},
with $|\nu|\le R_e(z)$, the radius of \eqref{9.2:eq-carry-b}.
\end{itemize}
We establish four properties of $G_e$.

First, $G_e(1,1)=\varphi(c_i)$: at $\mu=1$ the $\mu$-split restores $\mathbf{T}'_j(X_{i'})$, and
at $\nu=1$ the $\nu$-interpolation restores every $\psi\in\mathcal{S}_e$ at $c_i$.

Second, $G_e(1,0)=\varphi(c_i\lceil_{\mathcal{W}^{(e)}})$. At $\nu=0$ each $\psi\in\mathcal{S}_e$ takes
the value $\psi\lceil_{\mathcal{W}(\psi,z)}(c_i)$. Since $\mathcal{W}(\psi,z)$ consists of uses inside
$\psi$, switching it off affects $\varphi$ only through the value of $\psi$; this includes the copies of
$\psi$ inside carriers nested in the cluster, which are recomputed from that value, as arranged in the
construction of the auxiliary interpolation whose budget is \eqref{9.2:eq-interpolation-budget-a}.

Third, $G_e$ is holomorphic on the bidisc $\{|\mu|\le q^{-(j-e)}\}\times\{|\nu|\le R_e(z)\}$ and
satisfies $|G_e|\le\mathfrak{B}'_{X'}$ there. This is the budget
\eqref{9.2:eq-interpolation-budget-a}: its first bound is applied together with the inductive
hypothesis \eqref{9.2:eq-carry-c} at level $e$ for the occurrences $\psi\in\mathcal{S}_e$, at the
ambient point $c_i$, which is a genuine point of $\mathbb{D}^{\mathrm{use}}_\delta$; its second bound is
applied as it stands; and the allowance $+3$ for one $\mu$-split in
\eqref{9.2:eq-interpolation-budget-a} covers the split of the $\mathbf{T}'_j(X_{i'})$.

Fourth, $G_e(0,\nu)$ does not depend on $\nu$. At $\mu=0$ only the summands remain in which every
healed component satisfies $\operatorname{age}_j(\omega_{i'})<j-e$. Let $\psi\in\mathcal{S}_e$, and let
$C\subset X_{i_0}$ be its last witnessing component. By clause (P3) of the lemma on presence of
own-witness uses inside the healed components (Lemma~\ref{9.2:lem-presence-healed}; see
\eqref{9.2:eq-presence-healed-a}), applied at $\mathbf{R}_j$, $\psi$ is present in the
$(\omega_1,\dots,\omega_p)$-summand only if $\operatorname{age}_j(\omega_{i_0})\ge j-e$. Hence $\psi$ is
absent from every summand that remains at $\mu=0$, and this holds in each of the ways in which $\psi$
could enter. It is absent as a stored member read in the cluster. It is absent as a copy inside a
nested carrier or inside an outer occurrence, because the production histories of these contain that
of $\psi$, by \eqref{9.2:eq-age-factor-majorants-e}. Finally, a nested occurrence that contains a copy
of $\psi$ while keeping a live witness of its own does not exist: its domain would contain $X_\psi$,
which contains a cube of $(C^{\boxplus})^{(1)}\subset X_{i_0}$, so by the absorption property
\eqref{9.2:eq-witnessing-components-a} of Definition~\ref{9.2:def-witnessing-components} it is
absorbed at $\mathbf{R}_j$; it is therefore nested, not live. So at $\mu=0$ the dial $\nu$ acts on
nothing, and $G_e(0,1)=G_e(0,0)$.

By the first two properties, the $i$-th summand of \eqref{9.2:eq-witness-level-dials-1} equals
$G_e(1,1)-G_e(1,0)$; by the fourth, this equals the mixed second difference
$G_e(1,1)-G_e(1,0)-G_e(0,1)+G_e(0,0)$. We apply to it the Schwarz-type bound for holomorphic functions
of two variables on a bidisc, which bounds this mixed second difference by four times the bound
$\mathfrak{B}'_{X'}$ of the third property divided by the product of the radii $q^{-(j-e)}$ and
$R_e(z)$. Inserting the definition \eqref{9.2:eq-carry-b} of $R_e(z)$ gives
\begin{equation}\label{9.2:eq-witness-level-dials-a}
\begin{split}
\bigl|\varphi(c_i)-\varphi(c_i\lceil_{\mathcal{W}^{(e)}})\bigr|
&=\bigl|G_e(1,1)-G_e(1,0)-G_e(0,1)+G_e(0,0)\bigr|
\le\frac{4\,\mathfrak{B}'_{X'}}{q^{-(j-e)}\,R_e(z)}\\
&=4\,\mathfrak{B}'_{X'}\,q^{j-e}\,C_*\,\varepsilon^{(\mathrm{all})}_e\,q^{e-n-1}
=4\,C_*\,\varepsilon^{(\mathrm{all})}_e\,\mathfrak{B}'_{X'}\,q^{j-n-1}.
\end{split}
\end{equation}
Here $(e-n-1)_+=e-n-1$ because $e\ge n+2$, and $j-n-1=a$ because $j>e\ge n+2$.

\emph{Closure.} Since $\mathcal{W}(\varphi,z)$ is the disjoint union of $\mathcal{W}^{\mathrm{lump}}$ and
the high groups, $c'\lceil_{\mathcal{W}^{\mathrm{lump}}}=c\lceil_{\mathcal{W}(\varphi,z)}$. Moreover the
lump function $\Phi$ of \eqref{9.2:eq-lump-term-a}, taken at $\mu=1$, where the split restores
$\mathbf{T}'_j(X_{i'})$, equals $\varphi(c')-\varphi(c'\lceil_{\mathcal{W}^{\mathrm{lump}}})$. Hence
\begin{equation*}
\mathfrak{d}(\varphi;z)(c)=\varphi(c)-\varphi(c\lceil_{\mathcal{W}(\varphi,z)})
=\bigl[\varphi(c)-\varphi(c')\bigr]+\Phi(1).
\end{equation*}
We bound the first bracket by \eqref{9.2:eq-witness-level-dials-1} and
\eqref{9.2:eq-witness-level-dials-a}, summed over $i$, and the second term by the lump bound
$|\Phi(1)|\le2\,\mathfrak{B}'_{X'}q^a$ of \eqref{9.2:eq-lump-term-a}. We then bound the partial sum of
the piles by $\varepsilon^{\mathrm{all}}_{\mathrm{tot}}$, and write $q^a=q^{(j-n-1)_+}$:
\begin{equation}\label{9.2:eq-witness-level-dials-2}
\begin{split}
|\mathfrak{d}(\varphi;z)(c)|
&\le|\varphi(c)-\varphi(c')|+|\Phi(1)|
\le\mathfrak{B}'_{X'}\,q^{a}\Bigl[2+4C_*\sum_{e=n+2}^{j-1}\varepsilon^{(\mathrm{all})}_e\Bigr]\\
&\le\mathfrak{B}'_{X'}\,q^{(j-n-1)_+}\bigl[2+4C_*\,\varepsilon^{\mathrm{all}}_{\mathrm{tot}}\bigr]
\le C_*\,\mathfrak{B}'_{X'}\,q^{(j-n-1)_+}.
\end{split}
\end{equation}
The last inequality is $2+4C_*\varepsilon^{\mathrm{all}}_{\mathrm{tot}}\le C_*$. This is equivalent to
$C_*(1-4\varepsilon^{\mathrm{all}}_{\mathrm{tot}})\ge2$, which holds for
$C_*=4/(1-4\varepsilon^{\mathrm{all}}_{\mathrm{tot}})$ of \eqref{9.2:eq-carry-a}: indeed the left-hand
side then equals $4$. Thus \eqref{9.2:eq-witness-level-dials-2} is \eqref{9.2:eq-carry-c} at level $j$
at the arbitrary ambient point $c$.

The constant $C_*$ is the same at every level. The number $r$ of groups, the nesting depth of the
occurrences in $\mathcal{S}_e$, the ages $j-e$ and the ages of the healed components enter no constant.
They enter only exponents of $q$, and these are produced by the lemma on geometric forgetting of old
regions, which holds for every component $X$, whatever its history, and every $a\ge0$; it is applied
here to the healed components' own sequence sums inside the first group. No age is capped. This
completes the proof of clause (ii).
\end{proof}

\begin{proof}[Proof of the inner decay for second-group terms and strip clusters]
\label{9.2:prf-second-group-terms}
Let $j$ be a level and let $R'^{(j)}(X)$ be a term of the second group of the expansion
\cite[(1.98)]{B-CMP122b}, that is, a term whose domain satisfies $X\subset\Lambda_j^{\mathrm{new}}$.
Here $\Lambda_j^{\mathrm{new}}$ is the small-field region of step $j$, and its complement
$\Lambda_j^{c,\mathrm{new}}=Z_j\cup\bigcup_h Y_h$ is the large-field region of the step; in
\cite[\S3]{B-CMP122b} these are the terms whose domains are disjoint from that region.
Let $z=(n,k')$ be a class, and let $\mathcal{W}(R'^{(j)}(X),z)$ be the own-witness
set of Definition~\ref{9.2:def-own-witness-sets}, namely the union
\eqref{9.2:eq-own-witness-sets-a} formed for the cluster of the term. We show the following.
\begin{itemize}
\item[(a)] $R'^{(j)}(X)$ has no class-$z$ content outside $\mathcal{W}(R'^{(j)}(X),z)$. Its
own-witness variation is therefore its full $\lambda_z$-variation, and
\begin{equation}\label{9.2:eq-second-group-terms-a}
\sup\bigl|R'^{(j)}(X)-R'^{(j)}(X)\lceil_{\lambda_z=0}\bigr|
\le C_*\,\mathfrak{B}_X\,q^{(j-n-1)_+}.
\end{equation}
Here $R'^{(j)}(X)\lceil_{\lambda_z=0}$ is the term with every dial of class $z$ set to $0$. The
supremum is over $\mathbb{D}^{\mathrm{use}}_\delta$, with all other use-copies at arbitrary
admissible values. The factor $\mathfrak{B}_X$ is the second-group majorant of
\eqref{9.2:eq-age-factor-majorants-b}, and $C_*$ is the fixed-point constant of the inner-decay
bound for first-group leaves; it satisfies $C_*\le 16/3$ and depends on no level, age, depth,
$k$, $K$, torus or position.
\item[(b)] Let $X'$ be a domain containing one of the domains $Y_h$, and consider a cluster with
domain $X'$ that is counted in the term of index zero of the sum that Ba{\l}aban forms for such a
domain in \cite[\S3]{B-CMP122b}; that is, none of the components $X_i$ of the region $X'_0$
formed there is among the constituents of the cluster. For such a cluster the own-witness set
$\mathcal{W}$ consists of (A2) alone, the invariance \eqref{9.2:eq-carry-c} of the carry lemma
(Theorem~\ref{9.2:thm-carry}, whose proof is incomplete at one step) holds trivially, and the
cluster is a strip leaf, that is, a leaf whose constant in that sum is $\alpha^{1/3}$.
\end{itemize}

\emph{Part (a): the bound on the own-witness variation.} The argument that proves the inner decay
for a first-group leaf $\mathbf{B}'^{(j)}(X')$, with one pair of dials per witness level, applies
without change to $R'^{(j)}(X)$, with the majorant $\mathfrak{B}_X$ in place of
$\mathfrak{B}'_{X'}$. It gives
\begin{equation*}
\sup\bigl|R'^{(j)}(X)-R'^{(j)}(X)\lceil_{\mathcal{W}(R'^{(j)}(X),z)}\bigr|
\le C_*\,\mathfrak{B}_X\,q^{(j-n-1)_+}.
\end{equation*}

\emph{Part (a): no class-$z$ content outside the own-witness set.} We exclude in turn the three
kinds of class-$z$ content that could lie outside $\mathcal{W}(R'^{(j)}(X),z)$.

First, suppose there were a geo-witnessed copy, that is, one with $n<k'$. Its anchor $p$ would lie
in a live component of $\Lambda_{n+1}^c$. Such a component is contained in
$Z_j\cup\bigcup_h Y_h=\Lambda_j^{c,\mathrm{new}}$. On the other hand, the anchor lies in
$X\subset\Lambda_j^{\mathrm{new}}$, which is disjoint from $\Lambda_j^{c,\mathrm{new}}$. This is
impossible.

Second, suppose there were a nested occurrence $\psi$ with a surviving witness $C$. Then
$X\supseteq X_\psi$, where $X_\psi$ is the domain of $\psi$, and $X_\psi$ contains a cube meeting
$C^{\boxplus}$. Moreover $C^{\boxplus}$ is contained in the region obtained by applying the
enlargement $(\,\cdot\,)^{\sim}$ to the descendant of $C$. This
would make $X$ a first-group domain, which it is not.

Third, suppose there were a carrier $\psi$ with a surviving witnessing component $C$. Then $C$
would meet the domain $X_\psi\subset X$ of $\psi$, so that $X$ would meet the surviving component
$C$; as in the second case, this would make $X$ a first-group domain, which it is not.

Hence every dial of class $z$ on which $R'^{(j)}(X)$ depends belongs to
$\mathcal{W}(R'^{(j)}(X),z)$. Setting the dials of that set to zero is therefore the same as
setting every $\lambda_z$ to zero:
\begin{equation*}
R'^{(j)}(X)\lceil_{\mathcal{W}(R'^{(j)}(X),z)}=R'^{(j)}(X)\lceil_{\lambda_z=0}.
\end{equation*}
Inserting this into the display of the first part of (a) gives
\eqref{9.2:eq-second-group-terms-a}. This display has the form of the bound of the one-lattice
Lipschitz step, with the constant $4$ there replaced by $C_*$, and $C_*\le 16/3$.

\emph{Part (b).} Consider a cluster as in (b). Its domain $X'$ contains one of the domains $Y_h$.
By \cite[\S3]{B-CMP122b}, the sum that Ba{\l}aban forms over its summation index for such a domain
then begins with the index zero, and the constant that Ba{\l}aban denotes $c_0$ in that sum takes
the value $\alpha^{1/3}$; this constant is not the constant $c_0$ of
$\underline{\lambda}=c_0/4$. The cluster considered is counted in the term of index zero.

By hypothesis, none of the components $X_i$ of $X'_0$ is among the constituents of the cluster.
By the fact stated in \cite[\S3]{B-CMP122b} that every component $X_i$ of $X'_0$ that meets the
cluster belongs to it, the cluster therefore meets no $X_i$; in particular it contains no healed
component $X_i$.

Therefore no healed component's dial enters the own-witness set, and (A1) and (G) are empty. For
every witness level $e$ the set $\mathcal{S}_e$ is empty as well: an occurrence all of whose
witnesses die has a cube in some $X_i$; that $X_i$ would then meet the cluster and hence belong
to it, contrary to the hypothesis of (b).

Hence $\mathcal{W}$ consists of (A2) alone. For the same reason the invariance
\eqref{9.2:eq-carry-c} of the carry lemma (Theorem~\ref{9.2:thm-carry}, whose proof is incomplete
at one step) holds trivially for such a cluster. Since the constant $c_0$ of that sum equals
$\alpha^{1/3}$, such clusters are strip leaves. The inner-decay bound for strip leaves, with the
majorant $\mathfrak{B}'$ taken in the relative length and obtained as an implication from the
relative leaf bound, therefore applies to them.

The live-witnessed nested content such a cluster copies belongs to one of the occurrence classes
$(\alpha)$, $(\gamma)$ of part (i) of the carry lemma (Theorem~\ref{9.2:thm-carry}, whose proof
is incomplete at one step). That content is dialled at its own witness level, not at the level
of the cluster.

\emph{Scope of the argument.} The inductive argument establishes no decay for geo-witnessed
content, nor for the live-witnessed nested content of a first-group leaf at the level where the
leaf is produced; in neither case does such decay exist. No age factor exists for a surviving
zone at the level of production, and none is asserted. What such content costs at the level of
production is paid instead, at the steps that read it, by the class-$z$ clause of the estimate
made at those steps.

The statement obtained by applying the lemma on geometric forgetting of old regions to a component
whose history passed through a second-class event is a separate statement, the second-class
forgetting statement; it enters unchanged as a named input of another statement of this book, and
plays no part in the present induction.
\end{proof}

\begin{proof}[Proof of clause \textup{(iii)} of the carry lemma (Theorem~\ref{9.2:thm-carry}):
transport at the reading step $\mathbf{T}'_k(X)$ and the Lipschitz clause \textup{(d)}]
\label{9.2:prf-reader-transport}
Throughout this proof $\mathrm{e}$ denotes Euler's number, while $e$ denotes a level. We prove
the following. For every component $X$ of $Z_k$ (for any history), every bounded $F$ and every
$z=(n,k')$, with $\mathsf{G}:=\mathbf{T}^{\prime\lambda}_k(X)F$, the quantity
$\mathfrak{B}_z:=\sup_{\mathbb{D}_\delta}|\mathsf{G}-\mathsf{G}\lceil_{\lambda_z=0}|$ obeys the bound
\eqref{9.2:eq-carry-d}, which is clause (iii) of the carry lemma (Theorem~\ref{9.2:thm-carry},
stated; proof incomplete). The proof of clause (iii) given here is complete; the appeals it makes
to other clauses of that theorem are marked where they occur. From the bound we derive clause (d)
of the one-lattice Lipschitz step with its constant unchanged.

\emph{Set-up.} Write
\begin{equation*}
\mathsf{G}(\lambda):=\mathbf{T}^{\prime\lambda}_k(X)F=\sum_{\omega}\mathsf{G}_\omega(\lambda),
\end{equation*}
the sum running over the admissible sequences $\omega$ of $X$. This function is holomorphic on
$\mathbb{D}_\delta$ by the proposition of this chapter that establishes holomorphy in the dials
together with the majorants of all objects of the expansion. Let $z=(n,k')$ and put
$b:=(k-n-1)_+$. Fix $\lambda\in\mathbb{D}_\delta$, and let $c$ be the corresponding diagonal point
of $\mathbb{D}^{\mathrm{use}}_\delta$. For each $\omega$ let $U_\omega$ be the set of class-$z$
use-copies of the $\omega$-summand. Then
\begin{equation*}
\mathsf{G}_\omega(\lambda)=\mathsf{G}_\omega(c),\qquad
\mathsf{G}_\omega\lceil_{\lambda_z=0}(\lambda)=\mathsf{G}_\omega(c\lceil_{U_\omega}).
\end{equation*}

We partition $U_\omega$ by Lemma~\ref{9.2:lem-witness-partition}, clause (W1) of
\eqref{9.2:eq-witness-partition-a}, read in the history $\omega{\upharpoonright}$. For a level $e$
let $\mathcal{L}_e(\omega)$ be the set of leaf-occurrences of level $e$ in the $\omega$-summand that
are live-witnessed in $\omega{\upharpoonright}$. These are the occurrences read directly by the term
$V(X^{\sim})$ in the last exponential of the definition of $\mathbf{T}'_k(X)$ in \cite{B-CMP122b},
and those inside carriers or outer leaves lying in $X^{\sim}$. Put
\begin{equation*}
W^{(e)}_\omega:=\bigcup\bigl\{\mathcal{W}(\psi,z):\psi\in\mathcal{L}_e(\omega)\bigr\}.
\end{equation*}
Let $T_\omega$ be the rest of $U_\omega$. It consists of $X$'s own direct uses, the geometrically
witnessed copies, the own sets of live-witnessed carriers (of levels $n$ and $n+1$), and the own sets
of live-witnessed leaf-occurrences of level $e\le n+1$. By (W1), every $\upsilon\in U_\omega$ is
either geometrically witnessed, or lies in $\mathcal{W}(\psi,z)$ for exactly one live-witnessed
occurrence $\psi$. Hence the sets $W^{(e)}_\omega$, $n+2\le e\le k-1$, and $T_\omega$ are pairwise
disjoint; we put $W_\omega:=\bigsqcup_{e=n+2}^{k-1}W^{(e)}_\omega$, and
\begin{equation}\label{9.2:eq-reader-transport-1}
U_\omega=T_\omega\sqcup\bigsqcup_{e=n+2}^{k-1}W^{(e)}_\omega .
\end{equation}
Levels $e\ge k$ do not occur. The leaves of level $k$ are outputs of $\mathbf{R}_k$ itself, which
are produced after $\mathbf{T}'_k$, and the boundary terms $\mathbf{B}^{(k)}$ produced by the
operation $\mathbf{T}$ are carriers.

By clause (W2) of \eqref{9.2:eq-witness-partition-a}, a member of $W^{(e)}_\omega$ has witness level
$w(\upsilon)=e$, so it forces $\operatorname{age}(\omega)\ge k-e$. Each member of $T_\omega$ forces
$\operatorname{age}(\omega)\ge b$. For $X$'s own direct uses and for geometrically witnessed copies
this is the second alternative of (W2). For a member of the own set of a live-witnessed carrier or
leaf of level $e\le n+1$, (W2) gives $\operatorname{age}(\omega)\ge k-e\ge k-n-1$. Since ages are
non-negative, this implies $\operatorname{age}(\omega)\ge b$.

Inserting the intermediate point $c\lceil_{W_\omega}$ in each summand, we have
\begin{equation}\label{9.2:eq-reader-transport-2}
\begin{split}
\mathsf{G}(\lambda)-\mathsf{G}\lceil_{\lambda_z=0}(\lambda)
&=\sum_\omega\bigl[\mathsf{G}_\omega(c)-\mathsf{G}_\omega(c\lceil_{W_\omega})\bigr]\\
&\quad+\sum_\omega\bigl[\mathsf{G}_\omega(c\lceil_{W_\omega})
-\mathsf{G}_\omega(c\lceil_{U_\omega})\bigr].
\end{split}
\end{equation}
We call the second sum the $T$-part and the first sum the $W$-part.

\emph{(a) The $T$-part.} By \eqref{9.2:eq-reader-transport-1}, the points $c\lceil_{W_\omega}$ and
$c\lceil_{U_\omega}$ differ only on $T_\omega$. By the preceding paragraph, $T_\omega$ is empty
unless $\operatorname{age}(\omega)\ge b$. So the summands of the $T$-part with
$\operatorname{age}(\omega)<b$ vanish, and the sum runs over the sub-family of sequences of age at
least $b$. Each of the two resulting sub-sums is a sub-sum of $\mathbf{T}'_{k,\ge b}(X)F$ whose
summands are evaluated at points of $\mathbb{D}^{\mathrm{use}}_\delta$ depending on $\omega$. We
apply the tail-summed form \eqref{9.2:eq-age-factor-majorants-a} of geometric forgetting of old
regions, as stated in Notation~\ref{9.2:not-age-factor-majorants}, whose majorant holds for every
sub-family of sequences of age at least $b$, at the enlarged entropy constant $\mathsf{c}'$ of
\eqref{9.2:eq-age-factor-majorants-c}. It bounds each sub-sum by
$\mathrm{e}\,\mathfrak{E}_kq^b\sup|F|$. Hence
\begin{equation}\label{9.2:eq-reader-transport-3}
\Bigl|\sum_\omega\bigl[\mathsf{G}_\omega(c\lceil_{W_\omega})
-\mathsf{G}_\omega(c\lceil_{U_\omega})\bigr]\Bigr|
\le 2\,\mathrm{e}\,\mathfrak{E}_k\,q^{b}\sup|F| .
\end{equation}
In the proof of clause (d) of the one-lattice Lipschitz step the analogous supremum is bounded by
exactly this tail sum:
\begin{align*}
&2\sum_{a\ge(k-n-1)_+}\mathrm{e}\,\exp\bigl(-2(1+\beta_0)^{-1}p_0(g_k)\bigr)\,q^{a+1}\sup|F|\\
&\qquad\le 2\,\mathrm{e}\,\exp\bigl(-2(1+\beta_0)^{-1}p_0(g_k)\bigr)\,q^{(k-n-1)_+}\sup|F| .
\end{align*}
Here it bounds the $T$-part only, and the $W$-part is added in (b) and (c). That step argues that
the dial $\lambda_{n,k'}$ enters only sequences in which zone $n$ is present inside $X^{\sim}$, the
other sequences cancelling. For the members of $T_\omega$ this cancellation is the content of
clause (W2); the members of $W_\omega$ are transported in (b).

\emph{(b) The $W$-part: transport.} Enumerate the levels $e$ of $[n+2,k-1]$ as
$e_1<e_2<\dots<e_r$. As in the telescope leading to \eqref{9.2:eq-witness-level-dials-a}, we define
ambient points $\omega$-wise by
\begin{equation*}
c_i:=c\lceil_{W^{(e_{i+1})}_\omega\sqcup\dots\sqcup W^{(e_r)}_\omega},\qquad 1\le i\le r .
\end{equation*}
Then $c_r=c$, we have $c_i\lceil_{W^{(e_i)}_\omega}=c_{i-1}$ for $i\ge2$, and
$c_1\lceil_{W^{(e_1)}_\omega}=c\lceil_{W_\omega}$. Therefore
\begin{equation*}
\begin{split}
&\sum_\omega\bigl[\mathsf{G}_\omega(c)-\mathsf{G}_\omega(c\lceil_{W_\omega})\bigr]\\
&\qquad=\sum_{i=1}^r\sum_\omega\bigl[\mathsf{G}_\omega(c_i)
-\mathsf{G}_\omega(c_i\lceil_{W^{(e_i)}_\omega})\bigr].
\end{split}
\end{equation*}

Fix $i$ and put $e:=e_i$. For $\nu\in\mathbb{C}$ define
\begin{equation*}
\Phi_e(\nu):=\sum_\omega\mathsf{G}_\omega\bigl(c_i;\ \mathcal{L}_e(\omega)\text{ dialled by }\nu\bigr),
\end{equation*}
where dialling is meant in the sense of the auxiliary interpolation whose budget and propagation are
\eqref{9.2:eq-interpolation-budget-a}. Each $\psi\in\mathcal{L}_e(\omega)$ is read through its
interpolated value $\psi^{[\nu]}$, and later objects are recomputed. At $\nu=1$ every $\psi$ takes its
true value. At $\nu=0$ every $\psi$ takes the value $\psi\lceil_{\mathcal{W}(\psi,z)}(c_i)$, and the
union of these own sets over $\mathcal{L}_e(\omega)$ is $W^{(e)}_\omega$. So $\Phi_e(1)$ and
$\Phi_e(0)$ are the two consecutive telescope points $\sum_\omega\mathsf{G}_\omega(c_i)$ and
$\sum_\omega\mathsf{G}_\omega(c_i\lceil_{W^{(e)}_\omega})$.

If $\mathcal{W}(\psi,z)=\emptyset$, the interpolated value of $\psi$ is its true value. Hence the
$\omega$-summand depends on $\nu$ only if $\mathcal{L}_e(\omega)$ contributes a use-copy, that is, only
if $W^{(e)}_\omega\ne\emptyset$. By (W2) this forces $\operatorname{age}(\omega)\ge k-e$. Therefore
\begin{equation*}
\Phi_e(\nu)-\Phi_e(0)=\Phi^{\ge}_e(\nu)-\Phi^{\ge}_e(0),
\end{equation*}
where $\Phi^{\ge}_e$ is the restriction of the sum to the sub-family of sequences of age at least
$k-e$.

The function $\Phi^{\ge}_e$ is holomorphic on $|\nu|\le R_e(z)$, with $R_e(z)$ as in
\eqref{9.2:eq-carry-b}, and there
\begin{equation*}
|\Phi^{\ge}_e(\nu)|\le\mathrm{e}\,\mathfrak{E}_k\,q^{k-e}\sup|F| .
\end{equation*}
This holds for the following reasons.
\begin{itemize}
\item By the budget clause of \eqref{9.2:eq-interpolation-budget-a}, the dialled stored leaves of
level $e$ read by the term $V(X^{\sim})$ and by the inner data cost one unit of the enlarged
allowance \eqref{9.2:eq-age-factor-majorants-c}.
\item By the propagation clause of \eqref{9.2:eq-interpolation-budget-a}, the carriers and outer
leaves containing copies are recomputed inside their own majorants.
\item Geometric forgetting of old regions in the tail-summed form
\eqref{9.2:eq-age-factor-majorants-a}, taken together with the bound on modified terms that enters
the proof of the one-lattice Lipschitz step, bounds the age-$\ge(k-e)$ sub-sum at the constant
$\mathsf{c}'$.
\item The invariance clause \eqref{9.2:eq-carry-c} of the carry lemma
(Theorem~\ref{9.2:thm-carry}, stated; proof incomplete) is used at level $e$ at the
ambient point $c_i$, which is a genuine point of $\mathbb{D}^{\mathrm{use}}_\delta$.
\end{itemize}

We apply the one-variable Schwarz bound to $\Phi^{\ge}_e$ on the disc of radius $R_e(z)$. By
\eqref{9.2:eq-carry-b} together with Lemma~\ref{9.2:lem-leaf-pile},
$R_e(z)^{-1}=C_*\varepsilon^{(\mathrm{all})}_e\,q^{(e-n-1)_+}$, and $(e-n-1)_+=e-n-1$ because
$e\ge n+2$. This gives
\begin{equation}\label{9.2:eq-reader-transport-a}
\begin{split}
|\Phi_e(1)-\Phi_e(0)|
&\le\frac{2\,\mathrm{e}\,\mathfrak{E}_kq^{k-e}\sup|F|}{R_e(z)}
=2\,\mathrm{e}\,\mathfrak{E}_k\sup|F|\cdot C_*\varepsilon^{(\mathrm{all})}_e\,q^{(k-e)+(e-n-1)}\\
&=2\,\mathrm{e}\,\mathfrak{E}_k\sup|F|\cdot C_*\varepsilon^{(\mathrm{all})}_e\,q^{k-n-1}.
\end{split}
\end{equation}
Here $k-n-1=b$. Indeed, the range $[n+2,k-1]$ is non-empty only for $n\le k-3$, and then
$k-n-1\ge2$.

\emph{(c) Summation and the Lipschitz clause (d).} We insert \eqref{9.2:eq-reader-transport-3} and
\eqref{9.2:eq-reader-transport-a}, summed over $i$, into \eqref{9.2:eq-reader-transport-2}. The sum
of the $\varepsilon^{(\mathrm{all})}_e$ over the levels $e$ is bounded by
$\varepsilon^{\mathrm{all}}_{\mathrm{tot}}$ of Lemma~\ref{9.2:lem-leaf-pile}. The condition
\eqref{9.2:eq-carry-a} of the carry lemma (Theorem~\ref{9.2:thm-carry}, stated; proof
incomplete) gives $1+C_*\varepsilon^{\mathrm{all}}_{\mathrm{tot}}\le 4/3$. Taking the supremum over
$\lambda\in\mathbb{D}_\delta$ we obtain
\begin{equation}\label{9.2:eq-reader-transport-4}
\mathfrak{B}_z\le 2\,\mathrm{e}\,\mathfrak{E}_k\,q^{b}\sup|F|\,
\bigl(1+C_*\varepsilon^{\mathrm{all}}_{\mathrm{tot}}\bigr)
\le\frac{8\mathrm{e}}{3}\,\mathfrak{E}_k\,q^{b}\sup|F| .
\end{equation}
This is \eqref{9.2:eq-carry-d}.

At the point of $\mathbb{D}_\delta$ with all dials equal to $1$, resp.\ $0$, the function
$\mathsf{G}$ takes the value $\mathbf{T}^{\prime\tilde{\mathfrak{s}}}_k(X)F$, resp.\
$\mathbf{T}^{\prime\mathfrak{s}}_k(X)F$; we write $\mathsf{G}(1)$, $\mathsf{G}(0)$ for these values,
so that $\mathsf{G}(1)-\mathsf{G}(0)$ is the difference of clause (d). We bound it by the
several-variable Schwarz bound over the diagonal variables $\lambda_z$, whose radii in
$\mathbb{D}_\delta$ are $\varrho_z=E_1/[\![\delta]\!]_n\ge1$; it gives
$|\mathsf{G}(1)-\mathsf{G}(0)|\le\sum_z\mathfrak{B}_z/\varrho_z$. The bound
\eqref{9.2:eq-reader-transport-4} depends on $z=(n,k')$ only through $n$, and for each $n$ there are
$k-n+1$ values of $k'$. We then use $E_1=2(S_\varsigma+1)E_0\ge 2E_0$ and the definition of the
weighted load sum $\Sigma^q_k[\delta]$. This yields
\begin{align*}
|\mathsf{G}(1)-\mathsf{G}(0)|
&\le\sum_z\mathfrak{B}_z\,[\![\delta]\!]_n/E_1
\le\frac{8\mathrm{e}}{3}E_1^{-1}\mathfrak{E}_k\sup|F|
\sum_n(k-n+1)\,q^{(k-n-1)_+}[\![\delta]\!]_n\\
&\le\frac{4\mathrm{e}}{3}E_0^{-1}\mathfrak{E}_k\sup|F|\,\Sigma^q_k[\delta]\\
&\le 2\,\mathrm{e}\,E_0^{-1}\exp\bigl(-2(1+\beta_0)^{-1}p_0(g_k)\bigr)\sup|F|\cdot\Sigma^q_k[\delta].
\end{align*}
This is clause (d) of the one-lattice Lipschitz step with its constant unchanged. It holds under the
enlarged allowance \eqref{9.2:eq-age-factor-majorants-c} and the condition \eqref{9.2:eq-carry-a}.

\emph{On the arithmetic of the age factors.} Let $j=j(X)(\omega)$ be the index of the first
large-field region contained in $X$ along $\omega$, so that $\operatorname{age}(\omega)=a:=k-j$.
The inequality
\begin{equation*}
q^{k-j+1}q^{(j-n-1)_+}\ge q\cdot q^{(k-n-1)_+}
\end{equation*}
belongs to the per-age form of geometric forgetting of old regions, in which the age-$a$ part carries
the factor $q^{a+1}$. In the tail-summed form \eqref{9.2:eq-age-factor-majorants-a} used in (a) and
(b), the additional factor $q$ is spent on $1/(1-q)\le2$, since
$\sum_{a'\ge a}q^{a'+1}=q^{a+1}/(1-q)\le q^{a}$ for $q\le\tfrac12$. The bounds above are therefore
unaffected.
\end{proof}

\begin{proposition}[Zone-resolved variation of the first-group boundary terms. Stated; proof incomplete at the step marked $(\ast)$]
\label{9.2:prop-boundary-term-variation}
Let $\mathfrak{s}$ and $\tilde{\mathfrak{s}}$ be two lists in the admissible class $\mathfrak{K}_k$ and
$\delta$ their difference. Let $[\![\delta]\!]_n$ be the load of $\delta$ at the zone $n$, and let
\begin{equation*}
\Sigma^q_k[\delta]:=\sum_{n\le k}(k-n+1)\,q^{(k-n-1)_+}\,[\![\delta]\!]_n
\end{equation*}
be the weighted load sum, with $q$ the age ratio and $E_0$, $E_1=2(S_\varsigma+1)E_0$ the load constants.
For a pair $z=(n,k')$ of a zone $n$ and a level $k'$ with $n\le k'\le k$ (a \emph{class}), write
$\lambda_z$ for its dial variable and $\psi\lceil_{\lambda_z=0}$ for an object $\psi$ with this dial
set to $0$. Second-group terms are the rest terms $R'^{(k)}(X)$ of Ba{\l}aban's expansion
\cite[(1.98)]{B-CMP122b}, and first-group boundary terms are the terms
$\mathbf{B}'^{(k)}(X')$, also called first-group leaves, plain or strip as in
Lemma~\ref{9.2:lem-leaf-pile}; a superscript $\mathfrak{s}$ or $\tilde{\mathfrak{s}}$ records the list at
which a term is evaluated. We put
\begin{equation*}
\|\delta R'^{(k)}\|^{\flat}:=\sup_X e^{\kappa d_k(X)}\,
\bigl|R'^{(k),\tilde{\mathfrak{s}}}(X)-R'^{(k),\mathfrak{s}}(X)\bigr|,
\end{equation*}
and
\begin{gather*}
\delta R'^{(k)}(X):=R'^{(k),\tilde{\mathfrak{s}}}(X)-R'^{(k),\mathfrak{s}}(X),\\
\delta\mathbf{B}'^{(k)}(X'):=\mathbf{B}'^{(k),\tilde{\mathfrak{s}}}(X')-\mathbf{B}'^{(k),\mathfrak{s}}(X');
\end{gather*}
$\|\delta\mathbf{B}'^{(k)}\|^{\flat}$ is defined as $\|\delta R'^{(k)}\|^{\flat}$, with $X'$ in place of
$X$. The region $\Lambda_{n+1}$ is the small-field region of level $n+1$ of Ba{\l}aban's construction,
and its complement $\Lambda_{n+1}^c$ is the large-field region of that level; a component of
$\Lambda_{n+1}^c$ is \emph{live} if it still carries its large-field factor at the level read.

Assume the carry lemma, Theorem~\ref{9.2:thm-carry} (stated; proof incomplete), with its displays
\eqref{9.2:eq-carry-a}, \eqref{9.2:eq-carry-b}, \eqref{9.2:eq-carry-c}; the partition of uses by
witness, Lemma~\ref{9.2:lem-witness-partition}, in the form \eqref{9.2:eq-witness-partition-a}; the
per-cell pile of first-group leaves, Lemma~\ref{9.2:lem-leaf-pile}, with the pile
$\varepsilon^{(\mathrm{all})}_e$ of first-group leaves of level $e$; the bound
\eqref{9.2:eq-second-group-terms-a} for second-group terms; the interpolation budget
\eqref{9.2:eq-interpolation-budget-a}; the transport estimate \eqref{9.2:eq-reader-transport-a}; the
two-list estimate over the diagonal variables; and the Schwarz-type estimate on a disc. Let
$\mathfrak{B}_X$ and $\mathfrak{B}'_{X'}$ be the majorants of \eqref{9.2:eq-age-factor-majorants-b}.
Then the following hold.
\begin{itemize}
\item[(a)] For the second-group terms,
\begin{equation}\label{9.2:eq-boundary-term-variation-1}
\|\delta R'^{(k)}\|^{\flat}\le 4E_0^{-1}e^{-p_0(g_k)}\,\Sigma^q_k[\delta].
\end{equation}
\item[(b)] For every plain first-group leaf $\mathbf{B}'^{(k)}(X')$ and every class $z=(n,k')$,
\begin{equation}\label{9.2:eq-boundary-term-variation-a}
\begin{split}
\sup\bigl|\mathbf{B}'^{(k)}(X')&-\mathbf{B}'^{(k)}(X')\lceil_{\lambda_z=0}\bigr|
\le\mathfrak{B}'_{X'}\Bigl[C_*q^{(k-n-1)_+}
+2C_*\!\!\sum_{e}\varepsilon^{(\mathrm{all})}_e q^{e-n-1}\\
&+2\cdot\mathbf{1}\{X'\text{ meets a live component }C\text{ of }\Lambda_{n+1}^c
\text{ with }n<k',\text{ or }n=k'\\
&\qquad\text{and }X'\text{ holds a live-witnessed occurrence of level}\le n+1\}\Bigr],
\end{split}
\end{equation}
where the sum runs over the levels $e$ with $n+2\le e<k$ such that $X'$ copies a live leaf of level
$e$, and the supremum is over the per-use polydisc $\mathbb{D}^{\mathrm{use}}_\delta$ of the carry
lemma, all other use-copies being held at arbitrary admissible values. Every term beyond the first is
witnessed: at any later term of level $k''$ whose formula contains it (a \emph{reader}), it occurs
only in
sequences of age at least $k''-e$, respectively $(k''-n-1)_+$, and its decay in the age is obtained
there through \eqref{9.2:eq-reader-transport-a}.
\item[(c)] As a two-list bound,
\begin{equation}\label{9.2:eq-boundary-term-variation-2}
\begin{split}
|\delta\mathbf{B}'^{(k)}(X')|\le\mathfrak{B}'_{X'}E_1^{-1}\Bigl[&C_*\Sigma^q_k[\delta]
+2C_*\sum_{e}\varepsilon^{(\mathrm{all})}_e\,\Sigma^q_e[\delta]{\upharpoonright}\\
&+2\sum_{n\in\mathcal{N}(X')}(k-n+1)[\![\delta]\!]_n\Bigr].
\end{split}
\end{equation}
Here the middle sum runs over the levels $e$ of the live leaves that $X'$ copies, and
$\mathcal{N}(X')$ is the set of zones of the live components met by $X'$. The sum
\begin{equation*}
\Sigma^q_e[\delta]{\upharpoonright}:=\sum_{n\le e}(k-n+1)\,q^{(e-n-1)_+}\,[\![\delta]\!]_n
\end{equation*}
carries the age weights of the production level $e$ of the copied leaf, not those of $k$. The
constant $c_1$ used below is the small factor in Ba{\l}aban's bound on the boundary terms
\cite{B-CMP122b}, equal to $e^{-p_0(g_k)}$ for plain leaves, and equal to $\alpha^{1/3}$ for strip
leaves under the sharpened condition on strip leaves, in which case $\mathfrak{B}'_{X'}$ is taken in
the relative length of the strip leaf. The last term carries no factor $q$.
\item[(d)] $(\ast)$ $\|\delta\mathbf{B}'^{(k)}\|^{\flat}\le 4E_0^{-1}c_1\,\Sigma^q_k[\delta]$, with
$c_1$ as in (c): $c_1=e^{-p_0(g_k)}$ for plain leaves, and $c_1=\alpha^{1/3}$ for strip leaves under
the sharpened condition on strip leaves, the length $d_k(X')$ in the flat norm being replaced by the
relative length of the strip leaf.
\end{itemize}
\end{proposition}

\begin{proof}
\emph{Clause (a).} Let $X$ be the domain of a second-group term and $z=(n,k')$ a class. The inner
decay for second-group terms, \eqref{9.2:eq-second-group-terms-a}, states that for a second-group
term the own-witness variation is the full $\lambda_z$-variation, since such a term has no class-$z$
content outside its own witness set $\mathcal{W}(R'^{(k)}(X),z)$, and that
\begin{equation*}
\sup\bigl|R'^{(k)}(X)-R'^{(k)}(X)\lceil_{\lambda_z=0}\bigr|\le C_*\,\mathfrak{B}_X\,q^{(k-n-1)_+},
\end{equation*}
the supremum being over the per-use polydisc $\mathbb{D}^{\mathrm{use}}_\delta$ of the carry lemma,
all other use-copies held at arbitrary admissible values.
We apply the two-list estimate over the diagonal variables to these class-wise variations. It gives
\begin{equation*}
|\delta R'^{(k)}(X)|\le\sum_{z}C_*\,\mathfrak{B}_X\,q^{(k-n-1)_+}\,[\![\delta]\!]_n/E_1 .
\end{equation*}
For a fixed $n$ there are $k-n+1$ classes $z=(n,k')$, one for each $k'$ with $n\le k'\le k$. The sum
over $z$ is therefore
\begin{equation*}
\sum_{n\le k}(k-n+1)\,q^{(k-n-1)_+}\,[\![\delta]\!]_n=\Sigma^q_k[\delta].
\end{equation*}
By \eqref{9.2:eq-age-factor-majorants-b} and $E_1=2(S_\varsigma+1)E_0$, the factor in front of the
load sum satisfies
\begin{equation*}
C_*\,\mathfrak{B}_X/E_1\le\tfrac{C_*}{2}\,E_0^{-1}e^{-p_0(g_k)}e^{-\kappa d_k(X)},
\end{equation*}
and the bound becomes
\begin{equation*}
|\delta R'^{(k)}(X)|\le\tfrac{C_*}{2}\,E_0^{-1}e^{-p_0(g_k)}e^{-\kappa d_k(X)}\,\Sigma^q_k[\delta].
\end{equation*}
The constant $C_*$ of \eqref{9.2:eq-carry-a} (Theorem~\ref{9.2:thm-carry}; stated, proof incomplete)
satisfies $C_*\le 16/3$, so $C_*/2\le 8/3\le 4$.
Multiplying by $e^{\kappa d_k(X)}$ and taking the supremum over $X$ gives
\eqref{9.2:eq-boundary-term-variation-1}.

\emph{Clause (b).} Let $\varphi=\mathbf{B}'^{(k)}(X')$ be a plain first-group leaf and fix
$z=(n,k')$. Consider any class-$z$ use-copy in the formula of $\varphi$. By part (W1) of
Lemma~\ref{9.2:lem-witness-partition}, in the form \eqref{9.2:eq-witness-partition-a}, it is either
geo-witnessed, or it belongs to
$\mathcal{W}(\psi,z)$ for exactly one object-occurrence $\psi$, and that $\psi$ is live-witnessed.
Here $\psi$ is either $\varphi$ itself or a nested occurrence of a level $e<k$. Setting $\lambda_z=0$
on these three groups of use-copies in turn splits the full $\lambda_z$-variation of $\varphi$ into
three parts, which we call the own part, the flat part and the nested part. The
own part is the own-witness variation $\mathfrak{d}(\varphi;z)$ of \eqref{9.2:eq-carry-b}
(Theorem~\ref{9.2:thm-carry}; stated, proof incomplete). The flat part is
the variation of the geo-witnessed content together with the live-witnessed nested content of levels
$e\le n+1$. The nested part is the variation of the live-witnessed nested content of levels $e$ with
$n+2\le e<k$.

The own part obeys \eqref{9.2:eq-carry-c} at level $k$, by the carry lemma,
Theorem~\ref{9.2:thm-carry}, whose proof is incomplete at one step. Thus
\begin{equation*}
\sup|\mathfrak{d}(\varphi;z)|\le C_*\,\mathfrak{B}'_{X'}\,q^{(k-n-1)_+}.
\end{equation*}

For the nested part, fix a level $e$ with $n+2\le e<k$ of a live leaf copied into $\varphi$. We
interpolate
once in the auxiliary dial $\nu_e$ of \eqref{9.2:eq-interpolation-budget-a}. On the disc of radius
$R_e=R_e(z)$ of \eqref{9.2:eq-carry-b} (Theorem~\ref{9.2:thm-carry}; stated, proof incomplete),
$\varphi$ stays bounded by $\mathfrak{B}'_{X'}$. The
Schwarz-type estimate on a disc of radius $R_e$ bounds the variation by
\begin{equation*}
2\,\mathfrak{B}'_{X'}/R_e=2C_*\,\varepsilon^{(\mathrm{all})}_e\,\mathfrak{B}'_{X'}\,q^{(e-n-1)_+},
\end{equation*}
where $\varepsilon^{(\mathrm{all})}_e$ is the per-cell pile of Lemma~\ref{9.2:lem-leaf-pile}. These
bounds are summed over the levels $e$ with $n+2\le e<k$ of the live leaves copied into $\varphi$; for
these levels $(e-n-1)_+=e-n-1$.

For the flat part, both $\varphi$ and $\varphi\lceil_{\lambda_z=0}$ are bounded by
$\mathfrak{B}'_{X'}$, so this part is at most $2\mathfrak{B}'_{X'}$, whatever the levels of the
occurrences whose content it contains; this is why the nested content of levels $e\le n+1$ is counted
in the flat part and not in the nested part. It carries no age factor. It vanishes unless this group
of use-copies
is non-empty. By the classification of occurrences in clause (i) of the carry lemma,
Theorem~\ref{9.2:thm-carry} (stated; proof incomplete), a geo-witnessed class-$z$ copy occurs only for
$n<k'$ and has its anchor in a live component $C$ of $\Lambda_{n+1}^c$, which then meets $X'$; for
$n=k'$ the copies of this group are top-zone copies owned by live-witnessed occurrences of level $n$
or $n+1$. Hence the flat part is present only when the indicator in
\eqref{9.2:eq-boundary-term-variation-a} equals one.

Adding the three bounds gives \eqref{9.2:eq-boundary-term-variation-a}. Now let $k''$ be the level of
a later reader. The nested content of level $e$ has witness level $e$, and the geo-witnessed content
has witness level $n+1$. By part (W2) of Lemma~\ref{9.2:lem-witness-partition}, in the form
\eqref{9.2:eq-witness-partition-a}, such a use-copy occurs at
that reader only in sequences of age at least $k''-e$, respectively $(k''-n-1)_+$. The decay in that
age is obtained at the reader by \eqref{9.2:eq-reader-transport-a}.

\emph{Clause (c).} We apply the two-list estimate over the diagonal variables to
\eqref{9.2:eq-boundary-term-variation-a}, summing over the $k-n+1$ classes $z=(n,k')$ for each $n$,
and divide by $E_1$. The terms are treated one by one.
\begin{itemize}
\item The first term gives $C_*\Sigma^q_k[\delta]$, exactly as in clause (a). With
$c_1=e^{-p_0(g_k)}$ in $\mathfrak{B}'_{X'}$, this is the bound of (a) for the own part.
\item In the second term, fix a level $e$. The zones contributing to it have $n\le e-2$, and for
them $e-n-1=(e-n-1)_+$. Their contribution is
\begin{equation*}
\sum_{n\le e-2}(k-n+1)\,q^{e-n-1}[\![\delta]\!]_n\le\Sigma^q_e[\delta]{\upharpoonright},
\end{equation*}
since the remaining summands of $\Sigma^q_e[\delta]{\upharpoonright}$ are non-negative.
\item In the third term, each zone $n\in\mathcal{N}(X')$ contributes $(k-n+1)[\![\delta]\!]_n$, with
no factor $q$.
\end{itemize}
This gives \eqref{9.2:eq-boundary-term-variation-2}.

\emph{Clause (d).} $(\ast)$ The argument above gives the form of (d) for the own part of
$\delta\mathbf{B}'^{(k)}$: this is the first term of \eqref{9.2:eq-boundary-term-variation-2}. For
the nested part it gives only the production-level weights $\Sigma^q_e[\delta]{\upharpoonright}$.
For the flat part it yields only the flat term
\begin{equation*}
2\sum_{n\in\mathcal{N}(X')}(k-n+1)\,[\![\delta]\!]_n,
\end{equation*}
without the factor $q^{(k-n-1)_+}$; a bound of this term by a constant times $\Sigma^q_k[\delta]$,
whose zone-$n$ summand carries that factor, is not established here. The proof of (d) stops at this
step.
\end{proof}

\begin{remark}[The step marked $(\ast)$ and the use of the first-group boundary terms]
The proof of clause (d) is incomplete at the step marked $(\ast)$; the obstruction comes from the
following leaves. Consider a first-group leaf one of whose
polymer factors contains a unit term lying deep inside a live component of $\Lambda^c_{n+1}$. For such
a leaf the geo-witnessed content at the zones of that component has, at the level $k$ of production,
no bound carrying a factor in the age: the argument above bounds it only through the flat term, and
tail terms of this kind, with domain $Y$, occur with the weight $e^{-\kappa d_k(Y)}$ in the length
$d_k(Y)$ of their domain only \cite[(1.68)]{B-CMP122b}. For a
difference $\delta$ loaded on one such zone $n$ with $k-n$ large, the variation of such a leaf is
therefore controlled only by $2E_1^{-1}(k-n+1)[\![\delta]\!]_n\,\mathfrak{B}'_{X'}$, while clause (d)
requires the bracket of \eqref{9.2:eq-boundary-term-variation-2} to be bounded by a constant times
$\Sigma^q_k[\delta]$, whose zone-$n$ summand carries the factor $q^{(k-n-1)_+}$; the argument above
does not yield the bound of clause (d) for these leaves. Clause (d) is not used in what follows: the
first-group boundary terms are
derived quantities and enter only through the two bounds from which they are derived, never through the
number $\|\delta\mathbf{B}'^{(k)}\|^{\flat}$. After $\mathfrak{B}'_{X'}$ is inserted as in the proof
of clause (a), the first term of \eqref{9.2:eq-boundary-term-variation-2} gives a bound of the form
\eqref{9.2:eq-boundary-term-variation-1} with $e^{-p_0(g_k)}$ replaced by $c_1$. This first term is
the contribution of the own part, and the remaining two, from the nested and the flat part, form the
residual part of this bound. Finally, the re-filing
of the expansion leaves
$\mathbf{B}^{(k)}$ and $\mathbf{B}'^{(k)}$ unchanged. Hence, writing $\hat{\mathbf{B}}'^{(k)}$ for the
first-group boundary terms of the re-filed expansion, the copy of
clause (b) for $\hat{\mathbf{B}}'^{(k)}$ holds with \eqref{9.2:eq-boundary-term-variation-a} word for
word.
\end{remark}

\begin{proposition}[The tail of the rest terms above a threshold zone. Stated; proof incomplete at the step marked $(\ast)$]
\label{9.2:prop-rest-term-tail}
Let $k$ be a level, let $n$ be a zone, and let $R'^{(k)}$ be the rest terms of the second group at level $k$. Write $R'^{(k)}_{>n}$ for the part of $R'^{(k)}$ that the age split on the first index $j(X)(\omega)$ of the operated components assigns to the summands in which every operated component's sequence $\omega$ has first index greater than $n$. Then the following hold.
\begin{enumerate}
\item[\textup{(a)}] In the flat norm,
\begin{equation}\label{9.2:eq-rest-term-tail-1}
\bigl\|R'^{(k)}-R'^{(k)}_{>n}\bigr\|^{\flat}
\le 2\,e^{-(1-\beta_0)\,p_0^{\sharp}(n)}\,q^{\,k-n}.
\end{equation}
Here $p_0^{\sharp}(n)$ is the degree at zone $n$ in the age split of the one-lattice Lipschitz step; $\beta_0$ is the exponent of the age lemma comparing the degree function at an earlier level, $p_0(g_{k-a})$ for an age $a\ge1$, with $p_0(g_k)$; and $q=e^{-c_a}\le\tfrac12$, with $c_a$ the age constant of the one-lattice Lipschitz step, is its age ratio.
\item[\textup{(b)}] The tail $R'^{(k)}_{>n}$ contains no direct use-copy, of type \textup{(A1)} or \textup{(G)} in the classification \eqref{9.2:eq-own-witness-sets-a}, of a zone $n'<n$, and no live-witnessed content of witness level at most $n$. The use-copies of zones $n'<n$ that $R'^{(k)}_{>n}$ does carry lie in own-witness sets of occurrences of witness level $e>n$. Each such own-witness set obeys the invariant \eqref{9.2:eq-carry-c} of the carry lemma (Theorem~\ref{9.2:thm-carry}, whose proof is incomplete at one step) at level $e$, with inner factor
\begin{equation}\label{9.2:eq-rest-term-tail-a}
C_*\,q^{(e-n'-1)_+}\le C_*\,q^{\,e-n'-1}\le C_*\,q^{\,n-n'} .
\end{equation}
These copies appear only in summands in which the sequence $\omega$ of the operated component concerned has first index at most $e$.
\end{enumerate}
\end{proposition}

Stated; the proof below is incomplete at the step marked $(\ast)$.

\begin{proof}
Throughout, Lemma~\ref{9.2:lem-presence-healed} is applied at the level $s=k$ of $\mathbf{R}_k$, with $X_1,\dots,X_p$ the healed components and $\omega_i$ ranging over the admissible sequences of $X_i$; accordingly we write $k$ for $s$.

\emph{The bound \eqref{9.2:eq-rest-term-tail-1}.} Split $R'^{(k)}$ according to the first index $j(X)(\omega)$ of the operated components. The summands in which every operated component's sequence has first index greater than $n$ make up $R'^{(k)}_{>n}$. The summands in which at least one operated component's sequence has first index at most $n$ make up $R'^{(k)}-R'^{(k)}_{>n}$. This is a pure age split on that first index. $(\ast)$ The bound \eqref{9.2:eq-rest-term-tail-1} for the second part is used here as stated in the one-lattice Lipschitz step, in the age split of its rest terms. Its derivation is not carried out in this proposition.

\emph{Absence of direct copies below $n$.} Let $\upsilon$ be a use-copy of type (A1) or (G) of a zone $n'<n$. Apply (P1), respectively (P2), of Lemma~\ref{9.2:lem-presence-healed} with $n'$ in place of $n$ (see \eqref{9.2:eq-presence-healed-a}). The copy $\upsilon$ is present in the summand of an admissible sequence $\omega_i$ only if
\[
\operatorname{age}_k(\omega_i)=k-j(X_i)(\omega_i)\ge k-n'-1,
\]
that is, only if $j(X_i)(\omega_i)\le n'+1\le n$. Hence $\upsilon$ lives only in summands in which at least one operated component's sequence has first index at most $n$. These belong to $R'^{(k)}-R'^{(k)}_{>n}$, not to $R'^{(k)}_{>n}$.

\emph{Absence of live-witnessed content of level at most $n$.} Let $\psi$ be an occurrence of the class (B) of \eqref{9.2:eq-own-witness-sets-a} of witness level $e\le n$. We use, without deriving it here, that every occurrence of class (B) met in $R'^{(k)}$ has all its witnessing components among $X_1,\dots,X_p$; this is the hypothesis of (P3) of Lemma~\ref{9.2:lem-presence-healed}. Let $C\subset X_{i_0}$ be one of these components. By (P3) of Lemma~\ref{9.2:lem-presence-healed}, $\psi$ is present only if $\operatorname{age}_k(\omega_{i_0})\ge k-e$, that is, only if $j(X_{i_0})(\omega_{i_0})\le e\le n$. As in the preceding paragraph, $\psi$ does not occur in $R'^{(k)}_{>n}$.

\emph{The remaining copies of zones below $n$.} By the two preceding paragraphs, a use-copy of a zone $n'<n$ carried by $R'^{(k)}_{>n}$ is neither of type (A1) or (G) nor part of live-witnessed content of witness level at most $n$. Such a copy lies in an own-witness set $\mathcal{W}(\psi,z)$ of an occurrence $\psi$ whose witness level $e$ exceeds $n$. By (P3) of Lemma~\ref{9.2:lem-presence-healed}, applied to $\psi$ under the assumption on its witnessing components used in the preceding paragraph, such a copy appears only in summands in which the sequence $\omega_i$ of the operated component concerned has first index at most $e$.

Each such own-witness set obeys \eqref{9.2:eq-carry-c} at level $e$, by the carry lemma (Theorem~\ref{9.2:thm-carry}), whose proof is incomplete at one step. With $e$ in place of $j$ and $n'$ in place of $n$, the inner factor is $C_*q^{(e-n'-1)_+}$. Since $e>n>n'$, we have $e-n'-1\ge n-n'\ge1$. The positive part is therefore the number itself, which gives the first inequality in \eqref{9.2:eq-rest-term-tail-a}. The second inequality follows from $e-n'-1\ge n-n'$ and $q\le1$.
\end{proof}

Where $R'^{(k)}_{>n}$ enters a later estimate only through a majorant, the fact that part (b) permits use-copies of zones below $n$ in $R'^{(k)}_{>n}$ changes nothing there.

\begin{lemma}[Aligned parts of nested leaves and the unmatched term. Stated; proof incomplete
at the step marked $(\ast)$]\label{9.2:lem-aligned-parts}
Let $k_*$ be the alignment level of the cross-cutoff proposition. An admissible sequence
$\omega$ of a component $X$ is called
\emph{aligned} if its first index satisfies $j(X)(\omega)\ge k_*$. Let $\varphi$ be a
leaf-formula occurrence of level $j$ of a first-group leaf $\mathbf{B}'^{(j)}(X')$, and let
$X_1,X_2,\dots$ be the healed components of its cluster. The \emph{aligned part}
$\varphi^{\mathrm{AL}}$ of $\varphi$ is defined recursively in the level: it is the cluster of
$\varphi$ in which
\begin{itemize}
\item[(a)] for every $X_i$, the sum over the admissible sequences of $X_i$ is restricted to
the aligned sequences, and
\item[(b)] every nested occurrence $\psi$ that thereby loses its witness is replaced by its
aligned part $\psi^{\mathrm{AL}}$.
\end{itemize}
Put
\begin{equation}\label{9.2:eq-aligned-parts-1}
C^{\mathrm{AL}}:=\frac{1}{1-4\varepsilon^{\mathrm{all}}_{\mathrm{tot}}}\le\frac43 .
\end{equation}
Then the following hold.
\begin{itemize}
\item[(i)] At every point of $\mathbb{D}^{\mathrm{use}}_\delta$,
\begin{equation}\label{9.2:eq-aligned-parts-2}
|\varphi-\varphi^{\mathrm{AL}}|\le C^{\mathrm{AL}}\,\mathfrak{B}'_{X'}\,q^{(j-k_*)_+}.
\end{equation}
\item[(ii)] Under the hypotheses of the cross-cutoff proposition, its display holds with its
last member $4\exp(-p_0(g_k))\,q^{k-k_*}$ replaced by
\begin{equation}\label{9.2:eq-aligned-parts-3}
\bigl(1+C^{\mathrm{AL}}\varepsilon^{\mathrm{all}}_{\mathrm{tot}}\bigr)\,4\exp(-p_0(g_k))\,q^{k-k_*},
\qquad
1+C^{\mathrm{AL}}\varepsilon^{\mathrm{all}}_{\mathrm{tot}}\le 1+\tfrac{1}{12}<\tfrac43 ,
\end{equation}
and the factor $1+C^{\mathrm{AL}}\varepsilon^{\mathrm{all}}_{\mathrm{tot}}$ is absorbed where that
proposition absorbs its constants, in its bound $\le\mathfrak{r}_r\Lambda C_\sigma\vartheta_1^k$,
where $\Lambda$ and $C_\sigma$ denote constants of that proposition that may be enlarged (the
letter $\Lambda$ is not a region here).
\end{itemize}
\end{lemma}

Stated; the proof below is incomplete at the step marked $(\ast)$.

\begin{proof}
Here $\varepsilon^{(\mathrm{all})}_e$ is the per-cell pile of leaves of level $e$ denoted so
in the notation of \eqref{9.2:eq-leaf-pile-a}, $\varepsilon^{\mathrm{all}}_{\mathrm{tot}}$ is its sum
over the levels $e$, and $\mathfrak{B}'_{X'}$, $\mathfrak{B}_X$ are the majorants of
\eqref{9.2:eq-age-factor-majorants-b}. The need for the aligned part is the following: an
aligned sequence of $X$ may carry occurrences whose own-witness content predates $k_*$, so that
restricting the sequences of $X$ alone to first index $\ge k_*$ does not remove every
contribution from before $k_*$. The aligned part removes them recursively.

\emph{Proof of \textup{(i)}.} We argue by induction on the level $j$: assume
\eqref{9.2:eq-aligned-parts-2} for every leaf-formula occurrence of level $e<j$. The bound is
obtained by the same two stages as the own-witness bound in the proof of the carry lemma,
Theorem~\ref{9.2:thm-carry} (whose proof is incomplete at one step), with ``switching off the
own-witness set'' replaced by ``restricting the sequence sum of every nested healed component to
first index $\ge k_*$''. Accordingly $\varphi-\varphi^{\mathrm{AL}}$ is written as the
difference produced by the restriction (a) of the own sequence sums of the $X_i$, plus,
telescoped over the levels $e<j$ of the nested occurrences that lose their witness, the
differences produced by the replacement (b) of the occurrences of level $e$ by their aligned
parts.

First stage. An admissible sequence of $X_i$ that is not aligned has first index $<k_*$, hence
age $>j-k_*$ at level $j$. By the lemma on geometric forgetting of old regions, applied to the
own sequence sums of the $X_i$, the unaligned sub-sum of the $X_i$ is at most
$\mathfrak{B}'_{X'}q^{(j-k_*)_+}$.

Second stage. Fix a level $e<j$ of nested occurrences $\psi$ that lose their witness. We use the
two dials of the proof step on inner decay (one pair of dials per
witness level): the dial $\mu$ splits every large-field operation of the cluster at the age
threshold $j-e$, with $|\mu|\le q^{-(j-e)}$; the dial $\nu_e$ acts on the differences
$\psi-\psi^{\mathrm{AL}}$ of the occurrences of level $e$. By the induction hypothesis
\eqref{9.2:eq-aligned-parts-2} at level $e$ and the per-cell pile
\eqref{9.2:eq-leaf-pile-a}, these differences have per-cell size at most
$C^{\mathrm{AL}}\varepsilon^{(\mathrm{all})}_e q^{(e-k_*)_+}$, so that $\nu_e$ runs over the disc of
radius $\bigl(C^{\mathrm{AL}}\varepsilon^{(\mathrm{all})}_e q^{(e-k_*)_+}\bigr)^{-1}$. With these
radii the mixed-difference estimate \eqref{9.2:eq-witness-level-dials-a} bounds the contribution of
the level $e$ by
\begin{equation*}
4\,\mathfrak{B}'_{X'}\,q^{j-e}\,C^{\mathrm{AL}}\varepsilon^{(\mathrm{all})}_e\,q^{(e-k_*)_+}
\le 4\,C^{\mathrm{AL}}\varepsilon^{(\mathrm{all})}_e\,\mathfrak{B}'_{X'}\,q^{(j-k_*)_+},
\end{equation*}
where we used $q\le\frac12$ and $(j-e)+(e-k_*)_+\ge(j-k_*)_+$; the latter holds because for
$e\ge k_*$ the left side equals $j-k_*$ and is $\ge0$, while for $e<k_*$ it equals
$j-e>j-k_*$ and is $\ge0$.

Closure. Summing the first stage and the second stage over $e$, and using
$\sum_e\varepsilon^{(\mathrm{all})}_e\le\varepsilon^{\mathrm{all}}_{\mathrm{tot}}$,
\begin{equation*}
|\varphi-\varphi^{\mathrm{AL}}|
\le\mathfrak{B}'_{X'}\,q^{(j-k_*)_+}\bigl(1+4C^{\mathrm{AL}}\varepsilon^{\mathrm{all}}_{\mathrm{tot}}\bigr)
\le C^{\mathrm{AL}}\,\mathfrak{B}'_{X'}\,q^{(j-k_*)_+},
\end{equation*}
since the definition \eqref{9.2:eq-aligned-parts-1} gives
$C^{\mathrm{AL}}(1-4\varepsilon^{\mathrm{all}}_{\mathrm{tot}})=1$, that is, the fixed-point relation
$1+4C^{\mathrm{AL}}\varepsilon^{\mathrm{all}}_{\mathrm{tot}}=C^{\mathrm{AL}}$. This is
\eqref{9.2:eq-aligned-parts-2} at level $j$ and closes the induction. The constant
$C^{\mathrm{AL}}$ does not depend on $j$. The bound $C^{\mathrm{AL}}\le\frac43$ in
\eqref{9.2:eq-aligned-parts-1} is equivalent to
$\varepsilon^{\mathrm{all}}_{\mathrm{tot}}\le\frac1{16}$.

\emph{Proof of \textup{(ii)}.} $(\ast)$ This part rests on the proof of the cross-cutoff
proposition, used as it stands, and on two inputs of that proof which are not proved here and
are used unchanged: the zone-wise input of its affine interpolation step, which is stated on the
zones $\ge k_*-1$, and a second input of that proof. Two steps of the proof of the cross-cutoff
proposition are modified.

In the step that runs the affine interpolation dial, the dial is run on the fully aligned parts,
that is, on the terms in which the own sequence sum of $X$ is restricted to the aligned
sequences and the nested occurrences are treated as in (a) and (b) of the definition of the
aligned part; at $\lambda\equiv0$ the dial gives the fully aligned part of the expansion of the
first run $\mathsf{A}$ of Definition~\ref{9.6:def-comparison-chain}. All use-copies of the fully
aligned parts lie in zones $\ge k_*-1$, where the zone-wise input of the interpolation step is
stated.

In the step that collects the unmatched term, this term now collects, in each of the two
expansions compared, the own unaligned sequences of $X$, bounded in that proof by
$2\cdot2\mathfrak{B}_Xq^{k-k_*}$, and the nested remainders $\varphi-\varphi^{\mathrm{AL}}$. The
nested remainders are transported to the level $k$ by the transport estimate
\eqref{9.2:eq-reader-transport-a}, with the per-cell size of the dialled differences now given by
part~(i). In the notation of that display, with $n$ the zone of the class considered there and
$C_*$ as in \eqref{9.2:eq-carry-a}, the contribution of the level $e$ is therefore that of
\eqref{9.2:eq-reader-transport-a} with the factor $C_*\varepsilon^{(\mathrm{all})}_e q^{e-n-1}$
replaced by $C^{\mathrm{AL}}\varepsilon^{(\mathrm{all})}_e q^{(e-k_*)_+}$; that is, it is
\begin{equation*}
2\,\exp(1)\,\mathfrak{E}_k\,q^{k-e}\,
C^{\mathrm{AL}}\varepsilon^{(\mathrm{all})}_e\,q^{(e-k_*)_+}
\end{equation*}
times the remaining factors of \eqref{9.2:eq-reader-transport-a}, where
$\exp(1)\,\mathfrak{E}_k$ is the majorant of that display built on the age factor
\eqref{9.2:eq-age-factor-majorants-a}. Since $(k-e)+(e-k_*)_+\ge k-k_*$ (for $e\ge k_*$ both
sides are equal, and for $e<k_*$ the left side is $k-e>k-k_*$) and $q\le\frac12$, this
contribution carries the power $q^{k-k_*}$ of the own term of $X$; with its remaining factors
compared with those of that own term, which is part of the step marked $(\ast)$, it is at most
$C^{\mathrm{AL}}\varepsilon^{(\mathrm{all})}_e$ times that own term
$2\cdot2\mathfrak{B}_Xq^{k-k_*}$. Summing over $e$ and using
$\sum_e\varepsilon^{(\mathrm{all})}_e\le\varepsilon^{\mathrm{all}}_{\mathrm{tot}}$, the nested
remainders contribute at most $C^{\mathrm{AL}}\varepsilon^{\mathrm{all}}_{\mathrm{tot}}$ times the
own term of $X$. Hence the last member $4\exp(-p_0(g_k))\,q^{k-k_*}$ of the display of the
cross-cutoff proposition is multiplied by $1+C^{\mathrm{AL}}\varepsilon^{\mathrm{all}}_{\mathrm{tot}}$,
which is \eqref{9.2:eq-aligned-parts-3}; with $C^{\mathrm{AL}}\le\frac43$ and
$\varepsilon^{\mathrm{all}}_{\mathrm{tot}}\le\frac1{16}$ one has
$C^{\mathrm{AL}}\varepsilon^{\mathrm{all}}_{\mathrm{tot}}\le\frac1{12}$.

The factor is absorbed where the cross-cutoff proposition absorbs its constants, in the
bound $\le\mathfrak{r}_r\Lambda C_\sigma\vartheta_1^k$, since the constants $\Lambda$ and
$C_\sigma$ of that proposition are free to be enlarged; the remaining steps of the proof of that
proposition are not modified here. The variant of hypothesis (H3) of
Definition~\ref{2.2:not-hypotheses} in which the cross-cutoff proposition uses it keeps the same
form.
\end{proof}

\begin{remark}[The first index over the full history. Stated; proof incomplete at the step marked $(\ast)$]
\label{9.2:rem-full-history-index}
Clause (f) of the one-lattice Lipschitz step is a pointwise two-sided ratio bound at real data. Fix a
label and an admissible sequence $\omega$ completing its history. Let $h(\omega)$ be the first index
entering clause (f). Let $h^{\mathrm{full}}(\omega)$ be the least first index over the full history of
the label, the healed components whose first-group boundary terms $\mathbf{B}'^{(e)}$ are carried by
$\omega$ included.

Clause (f) is used with its first index replaced as follows.
\begin{itemize}
\item[(1)] The index $h(\omega)$ is replaced by $h^{\mathrm{full}}(\omega)$. Equivalently, the load
$[\![\delta]\!]_\omega$ of the list difference $\delta$ read by the label is replaced by
\begin{equation}\label{9.2:eq-full-history-index-1}
\sum_{n\ge h^{\mathrm{full}}(\omega)-1}(k-n+1)\,[\![\delta]\!]_n .
\end{equation}
\item[(2)] In the volume member, the count of $\mathbf{B}$-terms is the sum of two parts.
\begin{itemize}
\item[(a)] The first part is the count on the native strata $\Gamma^{\natural}_l$ and on the flattened
region $\Lambda_k^{\flat}$, with constant
\begin{equation}\label{9.2:eq-full-history-index-2}
C_P := C_L+2K_0c_LM^{-4}B_0/E_0 ,
\end{equation}
where $B_0$ is the size constant of the boundary terms and $c_L$ is the constant of the volume bound
on the flattened region.
\item[(b)] The second part is the contribution $2\varepsilon^{\mathrm{all}}_{\mathrm{tot}}C_L'/E_0$ of
the first-group boundary terms.
\end{itemize}
\end{itemize}
\end{remark}

\begin{proof}
\emph{The derived-term member.} The derived-term member of clause (f) uses the Schwarz-type bound on a
disc with the trivial majorant $\mathfrak{B}_u=2\mathfrak{B}_B$ in each dial variable $u$,
$\mathfrak{B}_B$ being the majorant on a disc of the boundary term. This gives
\begin{equation}\label{9.2:eq-full-history-index-3}
|\delta B^{(j)}(X)|\le 2B_0E_0^{-1}e^{-\kappa d_j(X)}\,[\![\delta]\!]_\omega .
\end{equation}
The bound \eqref{9.2:eq-full-history-index-3} is flat: its only weight is $e^{-\kappa d_j(X)}$, and it
carries no factor depending on the age.

\emph{Zones touched by the stored objects.} Let $\tau_\omega$ be the term of the label. An object stored
in $\tau_\omega$ touches the zones $n\ge h(\omega)-1$ of the live history. It also touches the zones of
the healed histories nested in it.

This follows from Lemma~\ref{9.2:lem-witness-partition}, with both parts of
\eqref{9.2:eq-witness-partition-a}: the own-witness sets $\mathcal{W}(\psi,z)$ there occur at every
witness level. Since a use-copy in the formula of an object stored in $\tau_\omega$ may belong to the
own-witness set of an object-occurrence of any witness level $e\le k$, these use-copies reach the zones
of every healed history nested in that object, not only the zones $n\ge h(\omega)-1$ of the live
history.

The load read by the label must therefore reach the first index of every history nested in the stored
objects, healed ones included. This is the replacement of $h(\omega)$ by $h^{\mathrm{full}}(\omega)$,
that is, of $[\![\delta]\!]_\omega$ by the sum \eqref{9.2:eq-full-history-index-1}.

\emph{The volume member.} The first part of the count of $\mathbf{B}$-terms is the one given by the
volume bound on the native strata and on the flattened region $\Lambda_k^{\flat}$, whose constant is
$C_P$ of \eqref{9.2:eq-full-history-index-2}. The second part comes from the first-group boundary terms
carried by $\omega$, the healed components included; their sum over all levels is bounded by
Lemma~\ref{9.2:lem-leaf-pile}, and this gives the added contribution
$2\varepsilon^{\mathrm{all}}_{\mathrm{tot}}C_L'/E_0$.

\emph{The form of the bound.} After this replacement, every member of clause (f) still has flat weights
and is a product of loads and volumes. The form of each member is therefore the same as before. The
statements that use clause (f) read $\mathsf{H}_\omega$ as an exponent, in the normalised form in which
the first index enters; the first index there is now $h^{\mathrm{full}}(\omega)$.

$(\ast)$ It remains to show that every other estimate in the proof of clause (f) holds unchanged with
$h^{\mathrm{full}}(\omega)$ in place of $h(\omega)$; this is not done here.
\end{proof}

\begin{remark}[Example: an old healed region beside a young live one]\label{9.2:rem-old-beside-young}
We test the carry lemma (Theorem~\ref{9.2:thm-carry}), whose proof is incomplete at one step, on
the following configuration. Let $h_1,n,j_D,j,k$ be levels with
\begin{equation*}
h_1<n+1<j_D\le j,\qquad j+2\le k ,
\end{equation*}
and no bound on their differences. The data are the following.
\begin{itemize}
\item[(a)] $X_1$ is a large-field region created at step $h_1$. Zone $n$ is a zone of its
history, and the operation $\mathbf{R}_j$ heals it with first-class outcome; at level $j$ it
satisfies conditions (i) and (ii) of \cite{B-CMP122a}.
\item[(b)] $D$ is a young large-field component created at step $j_D$. It is live, lies beside
$X_1$, its level-$j$ component is distinct from that of $X_1$, and it does not belong to the
class operated by $\mathbf{R}_j$.
\item[(c)] $\varphi:=\mathbf{B}'^{(j)}(X')$ is a first-group boundary term of
\cite[(1.98)]{B-CMP122b}. Its cluster $X'$ contains $\mathbf{T}'_j(X_1)$ and meets $D^{\sim}$.
After $\mathbf{R}_j$ it is stored, with witness the live component $D$.
\item[(d)] $z=(n,k')$ with $k'\le j$.
\item[(e)] $k\ge j+2$ is a level at which the component $X$ of $Z_k$ containing $D$ also
contains $X'$. Under one of the lower bounds on $L$ of the order of choice
(Lemma~\ref{2.4:lem-stage-zero-data}) the component of $D$ engulfs $X'$ within two steps, so
that $k=j+2$ may be taken; alternatively, without appeal to that bound, one takes for $k$ the
first level at which $X'$ is engulfed. In the current label $Z_i\cap X=\emptyset$ for $i<j_D$,
so that the first index of $X$ is $j(X)=j_D>n+1$, and every sequence $\omega$ of $X$ has
\begin{equation*}
\operatorname{age}(\omega)\le k-j_D<k-n-1 .
\end{equation*}
We put
$\mathsf{G}:=\mathbf{T}'^{\lambda}_k(X)F$ with $F$ bounded, and
\begin{equation*}
\mathfrak{B}_{n,k'}:=\sup_{\mathbb{D}_\delta}\bigl|\mathsf{G}-\mathsf{G}\lceil_{\lambda_z=0}\bigr|,
\end{equation*}
as in \eqref{9.2:eq-carry-d}.
\end{itemize}
We assume that $\varphi$ is plain, that is, $X'\cap\bigcup_h Y_h=\emptyset$, where the $Y_h$
are the second-class components of the class operated by $\mathbf{R}_j$. In this configuration
the operated class near $D$ has $X_1$ of the first class. If a second-class $Y_h$ were met, the
argument below applies with $\mathfrak{B}'_{X'}$ read in the relative length and with
$\varepsilon^{\mathrm{strip}}_e$ in place of $\varepsilon_e$; nothing else changes.

Under the smallness condition \eqref{9.2:eq-carry-a}, and with $\mathfrak{E}_k$ as in
\eqref{9.2:eq-age-factor-majorants-a}, the following hold.
\begin{itemize}
\item[(i)] Every use-copy of class $z$ in $\varphi$ is own-witness content of witness level
at least $n+2$. Moreover, for every sequence $\omega$ of $X$, the part $T_\omega$ of the
partition used in the proof of the transport estimate \eqref{9.2:eq-reader-transport-a} is
empty; this is the precise content of the statement that zone $n$ is present inside $X^{\sim}$
in no sequence of $X$.
\item[(ii)] We have
\begin{equation}\label{9.2:eq-old-beside-young-1}
\sup_{\mathbb{D}^{\mathrm{use}}_\delta}\bigl|\mathfrak{d}(\varphi;z)\bigr|
\le C_*\,\mathfrak{B}'_{X'}\,q^{\,j-n-1},
\end{equation}
the supremum being that of \eqref{9.2:eq-carry-c}, uniformly in the age $j-h_1$ of $X_1$.
\item[(iii)] We have
\begin{equation}\label{9.2:eq-old-beside-young-2}
\mathfrak{B}_{n,k'}\le \tfrac{2e}{3}\,\mathfrak{E}_k\,q^{(k-n-1)_+}\sup|F| .
\end{equation}
The exponent is assembled as $k-n-1=(j-n-1)+(k-j)$.
\item[(iv)] Suppose $\varphi$ is read instead through the exponentiation of $\mathbf{R}_k$.
Then either $D$ is healed at level $k$, and $\varphi$ falls under the proof of
\eqref{9.2:eq-witness-level-dials-a} at level $k$ in sequences of age at least $k-j$; or
$\varphi$ keeps its witness and is among the witnessed terms of
\eqref{9.2:eq-boundary-term-variation-a} (Proposition~\ref{9.2:prop-boundary-term-variation},
whose proof is incomplete at one step). For every threshold zone $n'>n$ the tail
$R'^{(k)}_{>n'}$ carries the zone-$n$ content of $\varphi$ only in sub-sequences of first index
at most $j$: not at all if $n'\ge j$, and with inner factor $C_*q^{j-n-1}\le C_*q^{n'-n}$ if
$n'<j$, in the form required by \eqref{9.2:eq-rest-term-tail-a}
(Proposition~\ref{9.2:prop-rest-term-tail}, whose proof is incomplete at one step).
\end{itemize}
\end{remark}

\begin{proof}
\emph{(i) Classes.} Every zone of $D$ is at least $j_D-1$. Since $j_D\ge n+2$, this exceeds
$n$. Hence no use-copy of class $z$ in $\varphi$ is geometrically witnessed inside $D$. The
use-copies of class $z$ in $\varphi$ are therefore of the following kinds.

The first kind, of type (A1) in \eqref{9.2:eq-own-witness-sets-a}, comes from the history of
$X_1$. It consists of the inner operations $\mathbf{T}^{(k')}$ of $\mathbf{T}'_j(X_1)$ and of
the terms $V(X_1^{\sim})$ and $V(Y)$, $Y$ a region of the neighbourhood of $X_1$, in the old
zone $n$ of $X_1$; here and below $V(\cdot)$ denotes the interaction terms of
\cite{B-CMP122b}, which enter the expansion exponentiated, as $e^{V(X^{\sim})}$ under the
operation $\mathbf{T}'_k(X)$. The inner operations
$\mathbf{T}^{(k')}$ of $\mathbf{T}'_j(X_1)$ and the terms $V(Y)$ of the neighbourhood of $X_1$
read units in zone $n$ inside $X_1$ at levels $k'\in[n,j]$.

There may also be copies of types (B) and (G). These are older stored boundary terms and
carriers (the $\mathbf{T}$-born boundary terms $\mathbf{B}^{(\cdot)}$) of the neighbourhood of
$X_1$. Their witnesses
were ancestors of $X_1$ or earlier-healed neighbours in the lineage of $X_1$, nested to any
depth. In particular they include the older terms $\mathbf{B}'^{(j')}(\cdot)$, $j'<j$, of
components healed earlier beside the lineage of $X_1$.

The last kind consists of copies, read by the terms $V(Y)$ of $X'$, of live-witnessed older
stored boundary terms anchored at $D$. These have levels $e\in[j_D,j)$, so their witness level
is $e\ge j_D\ge n+2$. This is case $(\gamma)$ of part (i) of the carry lemma
(Theorem~\ref{9.2:thm-carry}, whose proof is incomplete at one step). That part rests on
Lemma~\ref{9.2:lem-witness-partition}, clause (W1) of \eqref{9.2:eq-witness-partition-a}.

The copies of the first two kinds lie in $\mathcal{W}(\varphi,z)$ and have witness level $j$.
Those of the last kind lie in $\mathcal{W}(\psi,z)$ for occurrences $\psi$ of level $e$. Thus
all content of class $z$ in $\varphi$ is own-witness content of witness level at least $n+2$.
None of it is of the kinds collected in $T_\omega$ in the proof of the transport estimate
\eqref{9.2:eq-reader-transport-a}, namely direct uses of $X$, geometrically witnessed copies,
and the own-witness sets of live-witnessed carriers of levels $n,n+1$ and of live-witnessed
occurrences of level at most $n+1$. Thus $\varphi$
contributes nothing to $T_\omega$. The set $T_\omega$ collects the class-$z$ use-copies of these
kinds in the $\omega$-summand of $\mathsf{G}$. By Lemma~\ref{9.2:lem-witness-partition}, clause
(W2), every member of $T_\omega$ forces $\operatorname{age}(\omega)\ge(k-n-1)_+=k-n-1$. By
datum (e), every sequence $\omega$ of $X$ satisfies
\begin{equation*}
\operatorname{age}(\omega)\le k-j_D<k-n-1 .
\end{equation*}
Hence $T_\omega=\emptyset$ for every sequence $\omega$ of $X$.

\emph{(ii) Inner decay.} We apply the argument that proves \eqref{9.2:eq-witness-level-dials-a}
to $\varphi$. Its conclusion is \eqref{9.2:eq-carry-c} with $(j-n-1)_+=j-n-1$, and
\begin{equation*}
j-n-1\ge j_D-n-1\ge 1 .
\end{equation*}

The first ingredient is the $\mu$-split of $\mathbf{T}'_j(X_1)$ at age $a=j-n-1$, which catches
the copies of type (A1). Consider a use in zone $n$ inside $X_1$ at level $k'$. By
Lemma~\ref{9.2:lem-presence-healed}, clause (P1) of \eqref{9.2:eq-presence-healed-a}, applied
at $\mathbf{R}_j$ to the component $X_1$ operated by $\mathbf{T}'_j$, such a use is present only
in sequences $\omega_1$ of $X_1$ that satisfy
\begin{equation*}
Z^{\omega_1}_{n+1}\cap X_1\neq\emptyset,\qquad j(X_1)(\omega_1)\le n+1,\qquad
\operatorname{age}_j(\omega_1)\ge j-n-1 .
\end{equation*}
Here $\operatorname{age}_j(\omega_1):=j-j(X_1)(\omega_1)$. This is the statement that the dial
$\lambda_{n,k'}$ occurs only in sequences of $X_1$ in which zone $n$ is present inside $X_1$,
which clause (P1) proves for the operated component.

The second ingredient concerns the nested content of type (B) in the neighbourhood of $X_1$ with
witness level $e\in[n+2,j-1]$. This content is frozen by the auxiliary dial $\nu_e$. By
Lemma~\ref{9.2:lem-presence-healed}, clause (P3), it lives in sequences of $X_1$ of age at least
$j-e$. The estimate \eqref{9.2:eq-witness-level-dials-a} then returns
\begin{equation*}
4C_*\,\varepsilon_e\,q^{(j-e)+(e-n-1)}=4C_*\,\varepsilon_e\,q^{\,j-n-1}.
\end{equation*}
The sum over $e$ closes at the constant $C_*$ of \eqref{9.2:eq-carry-a}, which gives
\eqref{9.2:eq-old-beside-young-1}.

The age $j-h_1$ of $X_1$ is arbitrary. It is not bounded, not compared with anything and not
spent. What is spent is the factor $q^{\operatorname{age}_j(\omega_1)}$ supplied by geometric
forgetting of old regions. It is spent on the sub-family of sequences of $X_1$ that reach zone
$n$; the resulting bound is uniform in the age $j-h_1$ of $X_1$. The gap factor
$e^{-(1+\frac12\beta)\kappa d_j(X')}$ of
$\mathfrak{B}'_{X'}$, where $\beta$ is the exponent in \cite[(1.99)]{B-CMP122b}, is a prefactor
throughout and pays for nothing.

\emph{(iii) Transport.} By the construction of the terms $V(\cdot)$ in \cite{B-CMP122b},
$\varphi$ is a term of $V(X^{\sim})$, and it is
live-witnessed by $D$. We claim that in every sequence $\omega$ of $X$ carrying $\varphi$,
\begin{equation}\label{9.2:eq-old-beside-young-3}
D\subset Z^{\omega}_j\cap X,\qquad j(X)(\omega)\le j,\qquad \operatorname{age}(\omega)\ge k-j .
\end{equation}
This follows from Lemma~\ref{9.2:lem-witness-partition}, clause (W2), once its hypotheses are
checked. First, $D$ is a current component of $Z_j$. Next, by clause (ii) of the anchoring
statement for stored boundary terms in the correlation theorem, the anchor cube
$\square_\varphi$ of $\varphi$ satisfies
$\square_\varphi\subset X'\cap(D^{\boxplus})^{(1)}$, where $(D^{\boxplus})^{(1)}$ is the region
attached in that statement to the enlargement $D^{\boxplus}$ of $D$. By clause (i) of that
anchoring statement, $(D^{\boxplus})^{(1)}\subset Z_{j+1}$. Finally, $\square_\varphi\subset X^{\boxplus}$,
so the filling property of the enlargement (an anchor cube of a stored boundary term lying in
$X^{\boxplus}$ lies in $X$) gives $\square_\varphi\subset X$.

We now apply the transport estimate \eqref{9.2:eq-reader-transport-a} to the group $e=j$. The
dial $\nu_j$ dials $\varphi$, together with the other level-$j$ stored boundary terms anchored
at live components inside $X$. The function $\Phi^{\ge}_j$ is the sub-sum over sequences of age
at least $k-j$. By \eqref{9.2:eq-old-beside-young-3} and geometric forgetting of old regions it
is holomorphic on the disc $|\nu|\le R_j(z)$ and satisfies there
\begin{equation*}
|\Phi^{\ge}_j(\nu)|\le e\,\mathfrak{E}_k\,q^{k-j}\sup|F| .
\end{equation*}
The Schwarz-type estimate on the disc of radius $R_j(z)$ then gives
\begin{equation*}
|\Phi_j(1)-\Phi_j(0)|\le \frac{2e\,\mathfrak{E}_k\,q^{k-j}\sup|F|}{R_j(z)} ,
\end{equation*}
and with $1/R_j(z)=C_*\varepsilon_j\,q^{\,j-n-1}$ this contribution is
\begin{equation*}
2e\,\mathfrak{E}_k\sup|F|\cdot C_*\varepsilon_j\,q^{(k-j)+(j-n-1)}
=2e\,\mathfrak{E}_k\sup|F|\cdot C_*\varepsilon_j\,q^{\,k-n-1}.
\end{equation*}
The older stored boundary terms anchored at $D$, of levels $e\in[j_D,j)$, give in the same way
$2e\,\mathfrak{E}_k\sup|F|\,C_*\varepsilon_e\,q^{k-n-1}$.

By part (i), $T_\omega=\emptyset$ for every $\omega$, so the contribution of the sets $T_\omega$
to the transport argument vanishes. Summing over $e$, and using
$\sum_e\varepsilon_e\le\varepsilon_{\mathrm{tot}}\le\varepsilon^{\mathrm{all}}_{\mathrm{tot}}$
with $\varepsilon_{\mathrm{tot}}$ as in \eqref{9.2:eq-leaf-pile-b}, gives
\begin{align*}
\mathfrak{B}_{n,k'}
&\le 2e\,\mathfrak{E}_k\,q^{k-n-1}\sup|F|\cdot C_*\sum_e\varepsilon_e
\le 2e\,\mathfrak{E}_k\,q^{k-n-1}\sup|F|\cdot C_*\varepsilon^{\mathrm{all}}_{\mathrm{tot}}\\
&\le \tfrac{2e}{3}\,\mathfrak{E}_k\,q^{(k-n-1)_+}\sup|F| .
\end{align*}
In the last step we used $C_*\varepsilon^{\mathrm{all}}_{\mathrm{tot}}\le\frac13$, which is the
condition $1+C_*\varepsilon^{\mathrm{all}}_{\mathrm{tot}}\le\frac43$ under
\eqref{9.2:eq-carry-a}, and $(k-n-1)_+=k-n-1$, since $k\ge j+2>n+1$. This proves
\eqref{9.2:eq-old-beside-young-2}. The bound lies inside the bound
$2e\,\mathfrak{E}_k\,q^{(k-n-1)_+}\sup|F|$ for $\mathfrak{B}_{n,k'}$ used in clause (d) of the
one-lattice Lipschitz step, and inside \eqref{9.2:eq-carry-d}.

Geometric forgetting of old regions applied to $X$ alone, whose first index is $j_D$, supplies
only $q^{k-j_D+1}$. Since $j_D\ge n+2$,
\begin{equation*}
q^{k-j_D+1}\ge q^{(k-n-1)_+},
\end{equation*}
and this is strictly weaker as soon as $j_D>n+2$.
The exponent $k-n-1$ is assembled instead as the sum of two ages. The first is $j-n-1$, the age
of the sequences of $X_1$ carrying zone $n$, paid when $\varphi$ is produced. The second is
$k-j$, the age of the sequences of $X$ since $D\in Z_j$, paid at level $k$. The youth of $D$
($j_D>n+1$) and the antiquity of $X_1$ ($h_1\le n$, possibly much smaller) play no role.

\emph{(iv) Reading through the exponentiation of $\mathbf{R}_k$.} Suppose the term $\varphi$ is
read by the exponentiation of $\mathbf{R}_k$, which links $\mathbf{T}'_k(X)$ to a term $V(Y)$
containing $\varphi$. There are two cases.

Suppose first that $D$, as an ancestor of $X$, is healed at level $k$. Then $\varphi$ loses its
witness and belongs to the family $\mathcal{S}_j$ of the proof of
\eqref{9.2:eq-witness-level-dials-a}, which applies at level $k$. By
Lemma~\ref{9.2:lem-presence-healed}, clause (P3), $\varphi$ lives in sequences with
$\operatorname{age}_k\ge k-j$, and the estimate \eqref{9.2:eq-witness-level-dials-a} applies.

Suppose instead that $D$ survives $\mathbf{R}_k$ as a second-class component, or that the
cluster is anchored elsewhere. Then $\varphi$ keeps its witness. The bound
\eqref{9.2:eq-boundary-term-variation-a} of Proposition~\ref{9.2:prop-boundary-term-variation},
whose proof is incomplete at one step, counts it among its witnessed terms.

Consider finally the tail $R'^{(k)}_{>n'}$ of Proposition~\ref{9.2:prop-rest-term-tail}, whose
proof is incomplete at one step, at level $k$, for a threshold zone $n'>n$. It carries the
zone-$n$ content of $\varphi$ only in sub-sequences in which $D\in Z_j$, that is, of first index
at most $j$.
If $n'\ge j$, this content is therefore absent from $R'^{(k)}_{>n'}$. If $n'<j$, it is present
with inner factor $C_*q^{j-n-1}$. Since $n'<j$ gives $j-n-1\ge n'-n$, and $q<1$,
\begin{equation*}
C_*\,q^{\,j-n-1}\le C_*\,q^{\,n'-n}.
\end{equation*}
This is the form required by \eqref{9.2:eq-rest-term-tail-a}.

At no point is it used, for $\varphi$, that the dial $\lambda_{n,k'}$ occurs in the sequences
of $X$ only where zone $n$ is present inside $X^{\sim}$. That property is used only where it is
proved: for the kinds collected in $T_\omega$, by Lemma~\ref{9.2:lem-witness-partition}, clause
(W2), in part (i), and, in the analogous form for the component $X_1$ operated by
$\mathbf{T}'_j$ (zone $n$ present inside $X_1$), by Lemma~\ref{9.2:lem-presence-healed}, clause
(P1), in part (ii).
\end{proof}

\begin{definition}[Cube partitions, regions, components and fattening]\label{9.2:not-cubes-regions}
Throughout this chapter one cutoff and one scaffold (the radii $\check{R}_k$) are fixed, together
with the two runs compared, whose large-field labels are common to the two lists of labels and are
never interpolated between them; $d=4$. Let $\eta>0$
be the spacing of level~$0$. For $i\ge 0$, $\ell_i:=L^i\eta$ is the spacing of level $i$; larger $i$ is
coarser.

\emph{Lengths.} Lengths are measured in the $\ell^\infty$ metric of the torus, written $\mathbb{T}$
(the torus of Definition~\ref{1.1:def-lattice}; its index is fixed throughout this chapter and
suppressed, so that $m$ denotes a level); $|x-y|_\infty$ is the $\ell^\infty$-distance of two points
$x,y$ of $\mathbb{T}$. For $x\in\mathbb{T}$ and $\emptyset\ne B\subset\mathbb{T}$ we put
$\operatorname{dist}_\infty(x,B):=\inf_{y\in B}|x-y|_\infty$, and
$\operatorname{dist}_\infty(\cdot,\emptyset):=+\infty$; we write $(x)_+:=\max\{x,0\}$. For a
continuous path $\gamma\colon[\alpha,\beta]\to\mathbb{T}$,
\begin{equation*}
\operatorname{len}(\gamma):=\sup\Bigl\{\sum_{l=1}^{n}|\gamma(t_l)-\gamma(t_{l-1})|_\infty :
\ n\ge 1,\ \alpha=t_0<t_1<\dots<t_n=\beta\Bigr\}.
\end{equation*}
The $\ell^\infty$ metric is a \emph{length metric}: two points are joined by a path of
$\ell^\infty$-length equal to their distance. Moreover, for $\alpha\le\alpha'<\beta'\le\beta$,
\begin{equation*}
\operatorname{len}\bigl(\gamma|_{[\alpha',\beta']}\bigr)\ \ge\ |\gamma(\beta')-\gamma(\alpha')|_\infty ,
\end{equation*}
and the lengths of consecutive pieces of $\gamma$ add to at most $\operatorname{len}(\gamma)$.

\emph{Partitions.} A partition of the torus into closed cubes is a family of closed cubes covering the
torus with pairwise disjoint interiors. For each level $m$, $\mathcal{Q}_m$ is the partition of the
torus into closed cubes of side $M\check{R}_m\ell_m$, and $\pi_m$ is the partition into closed
\emph{$M$-cubes}, of side $M\ell_m$. All these partitions and the finer ones are
\emph{compatible}: for $m\le m'$, of two of these families of levels $m$ and $m'$, every cube of the
coarser is a union of cubes of the finer. Partitions of this kind,
each compatible with all the previous ones, are used in \cite{B-CMP119}.

\emph{Regions and components.} A \emph{region} is a finite union of closed cubes of one of these
families; we write $\mathcal{S}$ for the set of its cubes and $|\mathcal{S}|$ for their union, the
\emph{carrier}, which is closed. A region is \emph{connected} if its carrier is connected. The
\emph{components} of $\mathcal{S}$ are the classes of the equivalence relation on $\mathcal{S}$
generated by ``two cubes of $\mathcal{S}$ meet (share a point)''; each component is itself a
region. There are finitely many components; their carriers are compact, pairwise disjoint and at
positive $\ell^\infty$-distance from one another; in particular a region is connected if and only
if it has exactly one component. The following consequences are used.
\begin{itemize}
\item[(c1)] If $\mathcal{A}$ is a region, $\mathcal{D}$ a component of $\mathcal{A}$, and
$C$ a connected set, then
\begin{equation}\label{9.2:eq-cubes-regions-a}
C\subset|\mathcal{A}|,\quad C\cap|\mathcal{D}|\ne\emptyset\ \Longrightarrow\ C\subset|\mathcal{D}|.
\end{equation}
\item[(c2)] (Interior rule.) If $Q$ is a cube of $\mathcal{Q}_m$ (or of $\pi_m$) and
$\mathcal{T}$ is a region of a compatible family coarser than or equal to that of $Q$, then
\begin{equation}\label{9.2:eq-cubes-regions-b}
\operatorname{int}Q\cap|\mathcal{T}|\ne\emptyset\ \Longrightarrow\ Q\subset|\mathcal{T}|.
\end{equation}
\item[(c3)] (Granularity.) If $\mathcal{U}$ is a union of whole components of a region
$\mathcal{T}$ and $\mathcal{S}$ is a region with $|\mathcal{S}|\subset|\mathcal{T}|$, then for every
component $\mathcal{C}$ of $\mathcal{S}$,
\begin{equation}\label{9.2:eq-cubes-regions-c}
|\mathcal{C}|\cap|\mathcal{U}|\ne\emptyset\ \Longrightarrow\ |\mathcal{C}|\subset|\mathcal{U}|,
\end{equation}
and the other components of $\mathcal{S}$ have carriers disjoint from $|\mathcal{U}|$.
\end{itemize}

\emph{Fattening.} For a region $\mathcal{S}$ of the family $\mathcal{Q}_m$, let
$\mathcal{S}^{\sim 0}:=\mathcal{S}$, let $\mathcal{S}^{\sim 1}$ be the set of cubes of
$\mathcal{Q}_m$ meeting $|\mathcal{S}|$, and
$\mathcal{S}^{\sim n}:=(\mathcal{S}^{\sim(n-1)})^{\sim 1}$ for $n\ge 2$. Fattening is monotone:
$\mathcal{S}^{\sim(n-1)}\subset\mathcal{S}^{\sim n}$ for $n\ge1$, and $\mathcal{S}\subset\mathcal{T}$
implies $\mathcal{S}^{\sim n}\subset\mathcal{T}^{\sim n}$; moreover
$(\mathcal{S}\cup\mathcal{T})^{\sim n}=\mathcal{S}^{\sim n}\cup\mathcal{T}^{\sim n}$ for regions
$\mathcal{S},\mathcal{T}$ of $\mathcal{Q}_m$, and every $Q\in\mathcal{S}^{\sim n}$ heads a chain
$Q=Q_n,Q_{n-1},\dots,Q_0\in\mathcal{S}$ of successively meeting cubes with
$Q_t\in\mathcal{S}^{\sim t}$. For a closed set $A$, let $Z'_A$ be the set of cubes of
$\mathcal{Q}_m$ meeting $A$. (Transfer.) If moreover $Q$ meets $A$, then
\begin{equation}\label{9.2:eq-cubes-regions-d}
Q_{n-t}\in(Z'_A)^{\sim t}\quad(0\le t\le n),\qquad\text{in particular}\quad
Q_0\in\mathcal{S}\cap(Z'_A)^{\sim n}.
\end{equation}
If $A$ is the carrier of a connected region of a compatible finer family, then $Z'_A$ and every
$(Z'_A)^{\sim n}$ are connected.
\end{definition}

\begin{proof}
\emph{Lengths.} Given $x,y$, choose lifts $\tilde x,\tilde y$ in $\mathbb{R}^4$ with
$|\tilde y-\tilde x|_\infty=|x-y|_\infty$, and let $\gamma(t)$ be the image of
$\tilde x+t(\tilde y-\tilde x)$, $t\in[0,1]$. For any partition,
$|\gamma(t_l)-\gamma(t_{l-1})|_\infty\le(t_l-t_{l-1})|\tilde y-\tilde x|_\infty$, so every sum is at
most $|x-y|_\infty$; by the triangle inequality every sum is at least $|x-y|_\infty$. Hence
$\operatorname{len}(\gamma)=|x-y|_\infty$. The lower bound on
$\operatorname{len}(\gamma|_{[\alpha',\beta']})$ is the sum over the partition
$\alpha'=t_0<t_1=\beta'$. If $\alpha\le\tau\le\beta$, a partition of $[\alpha,\tau]$ followed by one of
$[\tau,\beta]$ is a partition of $[\alpha,\beta]$, so the two sums add to at most
$\operatorname{len}(\gamma)$; taking suprema gives
\begin{equation*}
\operatorname{len}\bigl(\gamma|_{[\alpha,\tau]}\bigr)+\operatorname{len}\bigl(\gamma|_{[\tau,\beta]}\bigr)
\le\operatorname{len}(\gamma),
\end{equation*}
and induction on the number of pieces gives the general case.

\emph{Components.} They are finitely many because $\mathcal{S}$ is finite. Carriers are finite
unions of closed cubes of the compact torus, hence compact. If $x\in|\mathcal{C}|\cap|\mathcal{C}'|$,
then $x\in Q\cap Q'$ with $Q\in\mathcal{C}$, $Q'\in\mathcal{C}'$, so $Q,Q'$ meet and
$\mathcal{C}=\mathcal{C}'$; thus distinct carriers are disjoint, and, being compact, at positive
distance. Each carrier is connected: two cubes $Q,Q'$ of a component $\mathcal{C}$ are joined by a
chain $Q=P_0,\dots,P_r=Q'$ of cubes of $\mathcal{S}$ with $P_{i-1}\cap P_i\ne\emptyset$, all in
$\mathcal{C}$; by induction on $i$, $P_0\cup\dots\cup P_i$ is connected (a union of two connected
sets that meet), so any two points of $|\mathcal{C}|$ lie in a connected subset of $|\mathcal{C}|$.
Hence a region with exactly one component is connected; and a region with two or more components
is not, since its carrier is then the union of the carrier of one component and the union of the
carriers of the others, two disjoint non-empty closed sets.

(c1). Let $F$ be the union of the carriers of the components of $\mathcal{A}$ other than
$\mathcal{D}$. Then $|\mathcal{A}|=|\mathcal{D}|\cup F$ with both sets closed and disjoint, so
$C\cap|\mathcal{D}|$ and $C\cap F$ are disjoint, relatively closed in $C$ and cover $C$. The first
is non-empty and $C$ is connected, so $C\cap F=\emptyset$ and $C\subset|\mathcal{D}|$.

(c2). If two distinct cubes $Q,Q'$ of one partition satisfied $\operatorname{int}Q\cap Q'\ne\emptyset$,
then, since $Q'$ is the closure of its interior and $\operatorname{int}Q$ is open,
$\operatorname{int}Q\cap\operatorname{int}Q'\ne\emptyset$, contrary to the definition of a
partition. Now $\operatorname{int}Q$ meets some $T\in\mathcal{T}$; by compatibility $T$ is a union of
cubes of the family of $Q$, so $\operatorname{int}Q$ meets one of them, $Q'\subset T$. Hence $Q'=Q$ and
$Q\subset T\subset|\mathcal{T}|$.

(c3). $|\mathcal{C}|$ is connected and contained in $|\mathcal{S}|\subset|\mathcal{T}|$. If it meets
$|\mathcal{U}|$, it meets $|\mathcal{D}|$ for a component $\mathcal{D}$ of $\mathcal{T}$ contained
in $\mathcal{U}$, and \eqref{9.2:eq-cubes-regions-a} with $\mathcal{A}=\mathcal{T}$ and
$C=|\mathcal{C}|$ gives $|\mathcal{C}|\subset|\mathcal{D}|\subset|\mathcal{U}|$. Otherwise
$|\mathcal{C}|\cap|\mathcal{U}|=\emptyset$.

\emph{Fattening.} $\mathcal{Q}_m$ is finite, so each $\mathcal{S}^{\sim n}$ is a region of
$\mathcal{Q}_m$. Every cube meets itself, so $\mathcal{S}^{\sim(n-1)}\subset\mathcal{S}^{\sim n}$;
and $\mathcal{S}\subset\mathcal{T}$ gives $|\mathcal{S}|\subset|\mathcal{T}|$, hence
$\mathcal{S}^{\sim1}\subset\mathcal{T}^{\sim1}$ and, by induction,
$\mathcal{S}^{\sim n}\subset\mathcal{T}^{\sim n}$. Since
$|\mathcal{S}\cup\mathcal{T}|=|\mathcal{S}|\cup|\mathcal{T}|$, a cube meets it exactly when it meets
$|\mathcal{S}|$ or $|\mathcal{T}|$, so
$(\mathcal{S}\cup\mathcal{T})^{\sim1}=\mathcal{S}^{\sim1}\cup\mathcal{T}^{\sim1}$; applying this to
$\mathcal{S}^{\sim(n-1)}$ and $\mathcal{T}^{\sim(n-1)}$ gives the identity for every $n$ by
induction. Chains: for $n=0$ take $Q_0=Q$. If $n\ge1$ and $Q\in\mathcal{S}^{\sim n}$, then $Q$ meets
$|\mathcal{S}^{\sim(n-1)}|$, hence meets some $Q_{n-1}\in\mathcal{S}^{\sim(n-1)}$, and the chain
for $Q_{n-1}$ given by induction extends to one for $Q$.

(Transfer.) For $t=0$, $Q_n=Q$ meets $A$, so $Q_n\in Z'_A=(Z'_A)^{\sim0}$. If
$Q_{n-t}\in(Z'_A)^{\sim t}$ with $t<n$, then $Q_{n-t-1}$ meets $Q_{n-t}$, hence meets
$|(Z'_A)^{\sim t}|$, so $Q_{n-t-1}\in(Z'_A)^{\sim(t+1)}$. At $t=n$ this gives
$Q_0\in(Z'_A)^{\sim n}$, and $Q_0\in\mathcal{S}$ by construction.

Connectedness. Each point of $A$ lies in a cube of $\mathcal{Q}_m$, which then meets $A$; so
$A\subset|Z'_A|$ and $|Z'_A|=A\cup\bigcup_{Q\in Z'_A}Q$ is a union of the connected set $A$ with
connected cubes each meeting $A$, hence connected. For $n\ge1$,
$(Z'_A)^{\sim n}=Z'_{B}$ with $B=|(Z'_A)^{\sim(n-1)}|$, by the definition of fattening; $B$ is
connected by induction, and the same argument with $B$ in place of $A$ shows that
$(Z'_A)^{\sim n}$ is connected.
\end{proof}

\begin{definition}[The frozen step and its stored boundary terms]\label{9.2:not-frozen-step}
Throughout this section we work with one renormalisation step $\mathbf{T}$ from level $k$ to
level $k+1$. This is the frozen step, taken on the large-field labels common to the two runs
of Definition~\ref{9.2:not-cubes-regions}. We put
\begin{equation*}
  \Omega := \Omega_{k+1},
\end{equation*}
the fluctuation domain of this step in the sense of \cite{B-CMP119}. It is a region in the
sense of Definition~\ref{9.2:not-cubes-regions}.

If $\Omega=\emptyset$, every quantity introduced below in this section is trivial; in
particular the per-unit response factor $\mathsf{r}_{\varpi}$, introduced later in this
section, vanishes identically. We therefore assume $\Omega\neq\emptyset$ from now on.

The estimates of this section read the stored boundary terms $\mathbf{B}^{(j)}$ of the
levels $j\le k$. Ba{\l}aban's treatment of these terms is described in \cite{B-CMP119} as
follows:
\begin{quotation}
The expression determined by the boundary terms $B_k$ is much simpler to deal with. To its
terms we apply the formulas (3.6), (3.7) [I], and next we localize the obtained expressions
\end{quotation}
Here the reference [I] of \cite{B-CMP119} is \cite{B-CMP109}. The paper \cite{B-CMP119}
also states that
\begin{quotation}
the integral over $t$ of the expectation value of the boundary terms \ldots{} contributes to
$B^{(k+1)}$ only
\end{quotation}
By \cite{B-CMP119}, the stored levels are the integers $j$ with $1\le j\le k$. For such a
level $j$ we write
\begin{equation*}
  s := k-j \ \ge 0
\end{equation*}
for the number of steps that separate the stored level $j$ from the current level $k$. In
this section the letter $s$ carries this meaning only.
\end{definition}

\begin{definition}[Step-made and current large-field sets; local scale]\label{9.2:def-current-sets}
We work with the frozen step $k\to k+1$ on the common large-field label, as in
Notation~\ref{9.2:not-frozen-step}, and with the cube partitions $\mathcal{Q}_m$, regions and
closures of Notation~\ref{9.2:not-cubes-regions}.

\emph{Time.} The operations that make the label are linearly ordered; the points of this order are
called \emph{stages}, or times.

The $\mathbf{T}$-step $m-1\to m$ occupies the stage $t_{\mathbf{T}}(m)$. It makes two
sets:
\begin{itemize}
\item the action domain $\Lambda_m^{\mathbf{T}}$, which is a region of $\mathcal{Q}_m$;
\item the determining domain $\Omega_m^{\mathbf{T}}$ of that step.
\end{itemize}
In \cite{B-CMP119} it is stated that the domains $\Omega_j,\Lambda_j$, which are determined by the
$j$-th renormalization transformation and not by the $\mathbf{R}$-operation, are unions of
$MR_j$-cubes (in the notation of \cite{B-CMP119}); with the scaffold's radii these are the cubes of
$\mathcal{Q}_j$ (cubes of side $M\check{R}_j\ell_j$; Notation~\ref{9.2:not-cubes-regions}). The
native determining sequence $\{\Omega_n^{\natural}\}$ of the label is, by the statement on the
native determining sequence, the sequence made by the $\mathbf{T}$-steps with whole healed
components deleted; these are the reset classes $E$ of \eqref{9.2:eq-attempt-windows-a}.

An $\mathbf{R}$-attempt (completed or not; see \eqref{9.2:eq-attempt-windows-a}) at step $k_a$ acts
on $\rho_{k_a}$ after the time $t_{\mathbf{T}}(k_a)$ and
before the $\mathbf{T}$-step $k_a\to k_a+1$ reads the label. This is the order in which the step of
Notation~\ref{9.2:not-frozen-step} reads the large-field operation at step $k+1$: the outputs
$R'^{(k+1)}$, $\mathbf{B}'^{(k+1)}$ of $\mathbf{R}$ at step $k+1$, on the components to which it
applies at that step, are read near $\partial\Lambda_{k+1}\cup\mathbf{B}'$, and carry the prefactor
$e^{-p_0(g_k)}$, with $p_0(\cdot)$ the degree function. The stage \emph{now} is the reading of
the label by the $\mathbf{T}$-step $k\to k+1$.

\emph{Step-made sets.} Put
\begin{equation}\label{9.2:eq-current-sets-1}
Z_m^{\mathbf{T}}:=\operatorname{cl}\bigl((\Lambda_m^{\mathbf{T}})^{c}\bigr),
\end{equation}
a region of $\mathcal{Q}_m$, namely the cubes of $\mathcal{Q}_m$ not in $\Lambda_m^{\mathbf{T}}$.

The action domains $\Lambda_i$ of the label and its native determining sequence
$\{\Omega_n^{\natural}\}$, with whole healed components deleted, satisfy, by the statement on the
native determining sequence, the nesting
\begin{equation}\label{9.2:eq-current-sets-a}
\Omega^{\natural}_{i+1}\subset\Lambda_i\subset\Omega^{\natural}_i,
\qquad \Lambda_i\ \text{a union of cubes of}\ \mathcal{Q}_i .
\end{equation}
In this definition \eqref{9.2:eq-current-sets-a} is read at the making time $t_{\mathbf{T}}(m)$
only: with $i=m$ it gives $\Lambda_m^{\mathbf{T}}\subset\Omega_m^{\mathbf{T}}$, cube-wise in
$\mathcal{Q}_m$, in the sense of \eqref{9.2:eq-cubes-regions-b}; with $i=m-1$ it gives
$\Omega_m^{\mathbf{T}}\subset\Lambda_{m-1}$, where $\Lambda_{m-1}$ is the action domain of level
$m-1$ current at $t_{\mathbf{T}}(m)$.

\emph{Current sets.} Let $E$ denote the reset class of a completed $\mathbf{R}$-operation (see
\eqref{9.2:eq-attempt-windows-a}) and $k_E$ its step. The \emph{current set} of level $m$ is
\begin{equation}\label{9.2:eq-current-sets-2}
Z_m:=Z_m^{\mathrm{cur}}:=\operatorname{cl}\Bigl(|Z_m^{\mathbf{T}}|\setminus\bigcup_{E}|E|\Bigr),
\end{equation}
where the union is over the reset classes $E$ of completed $\mathbf{R}$-operations with
$k_E\ge m$ processed so far. Each $|E|$ is closed, being the carrier of a union of whole
components (see \eqref{9.2:eq-attempt-windows-a}), and so is the finite union $\bigcup_E|E|$.
Each class $E$ with $k_E\ge m$ is a union of whole components of a region of $\mathcal{Q}_{k_E}$,
so $|E|$ is a union of cubes of $\mathcal{Q}_m$ by \eqref{9.2:eq-cubes-regions-b}. A cube of
$\mathcal{Q}_m$ that is not among the cubes making up $\bigcup_E|E|$ meets that set only on its
boundary, since the cubes of $\mathcal{Q}_m$ have disjoint interiors; hence it is not contained in
$\bigcup_E|E|$, and it lies in the closure of its part outside that set. Therefore
$Z_m^{\mathrm{cur}}$ is the region of $\mathcal{Q}_m$ consisting of the cubes of
$Z_m^{\mathbf{T}}$ not contained in $\bigcup_E|E|$, and the right-hand side of
\eqref{9.2:eq-current-sets-2} is its carrier.
Here and below $Z_m$ and $Z_m^{(t)}$ are read as subsets of the torus, a region being identified
with its carrier $|\cdot|$ wherever it is compared with sets of other levels.
At a stage $t\ge t_{\mathbf{T}}(m)$ of the label's history, the current set $Z_m^{(t)}$ is defined
by \eqref{9.2:eq-current-sets-2} with the union taken over the reset classes $E$ with $k_E\ge m$
processed before $t$. Thus $Z_m=Z_m^{(t)}$ when $t$ is the stage now.

The letters $Z_j=\operatorname{cl}(\Lambda_j^{c})$, $Z_j^{\boxplus}$ (the weak-anchor set of
level $j$), $\Lambda_j$ and $\Omega_j$
without further decoration, with $\Lambda_j$ the current action domain, always denote the sets of
the current label
(the re-made sets of \eqref{9.2:eq-attempt-windows-a} carry primes). Arrested or excluded attempts
(see \eqref{9.2:eq-attempt-windows-a}) leave all $\Lambda_j$ unchanged.

The healed classes may be accounted for in two ways: either the current action domain of level $m$
is taken to be
\begin{equation*}
\Lambda_m:=\Lambda_m^{\mathbf{T}}\cup\bigcup_{E}|E|,
\end{equation*}
the union being the one in \eqref{9.2:eq-current-sets-2}, or $\Lambda_m^{\mathbf{T}}$ is kept and
the classes are deleted from $Z$. The set $Z_m^{\mathrm{cur}}$ is the same either way:
\begin{equation*}
\operatorname{cl}(\Lambda_m^{c})=Z_m^{\mathrm{cur}},
\end{equation*}
by the closure identity proved below. In what follows the undecorated $\Lambda_m$ denotes the first
of the two, so that $Z_m=\operatorname{cl}(\Lambda_m^{c})$. At the making time
$t_{\mathbf{T}}(m)$ no class with $k_E\ge m$ has been processed, since a completed operation at
step $k_E$ acts after $t_{\mathbf{T}}(k_E)\ge t_{\mathbf{T}}(m)$; there
$\Lambda_m=\Lambda_m^{\mathbf{T}}$, and this is the reading of \eqref{9.2:eq-current-sets-a} used
below.
The set $Z_m^{\mathrm{cur}}$ is the only large-field set of level $m$ used below.

\emph{Nesting.} For $i<k$ and every stage $t\ge t_{\mathbf{T}}(i+1)$,
\begin{equation}\label{9.2:eq-current-sets-b}
Z_i^{(t)}\subset Z_{i+1}^{(t)};\qquad\text{in particular}\quad Z_i\subset Z_{i+1}\ \text{now.}
\end{equation}

\emph{Current cubes and local scale.} For $n\le k$, the \emph{current stratum}
$\Omega_n\setminus\Omega_{n+1}$ is a union of closed cubes of compatible families. These are called
the \emph{current cubes of scale $n$}. The cubes of $\Omega=\Omega_{k+1}$ are the current cubes of
scale $k+1$. Write $\Omega_n^{\mathrm{cur}}:=\Omega_n$ for the current determining domain of level
$n$; these sets are closed. The \emph{local scale} of a point $x$ is
\begin{equation}\label{9.2:eq-current-sets-3}
n(x):=\max\{\,n\le k+1:\ x\in\Omega_n^{\mathrm{cur}}\,\}.
\end{equation}
The \emph{native scale} $n^{\natural}(x)$ is that of the statement on the native determining
sequence; one has $n^{\natural}(x)\le n(x)$.
\end{definition}

\begin{proof}
We prove four statements in turn: a closure identity, the inclusion at the making time, the nesting
\eqref{9.2:eq-current-sets-b}, and the comparison of the two scales.

\emph{A closure identity.} Let $G$ be a subset of the torus and $\mathcal{F}$ a closed subset.
Every point of $\operatorname{cl}(G)\setminus\mathcal{F}$ has a neighbourhood disjoint from
$\mathcal{F}$ and is therefore a limit of points of $G\setminus\mathcal{F}$; hence
$\operatorname{cl}(G)\setminus\mathcal{F}\subset\operatorname{cl}(G\setminus\mathcal{F})$, while
$G\setminus\mathcal{F}\subset\operatorname{cl}(G)\setminus\mathcal{F}$. Taking closures,
\begin{equation}\label{9.2:eq-current-sets-5}
\operatorname{cl}(G\setminus\mathcal{F})
=\operatorname{cl}\bigl(\operatorname{cl}(G)\setminus\mathcal{F}\bigr).
\end{equation}
With $G:=(\Lambda_m^{\mathbf{T}})^{c}$ and $\mathcal{F}:=\bigcup_E|E|$, the union of
\eqref{9.2:eq-current-sets-2}, which is closed, one has $\Lambda_m^{c}=G\setminus\mathcal{F}$ and
$|Z_m^{\mathbf{T}}|=\operatorname{cl}(G)$, so \eqref{9.2:eq-current-sets-5} gives
$\operatorname{cl}(\Lambda_m^{c})=Z_m^{\mathrm{cur}}$.

\emph{Inclusion at the making time.} Read \eqref{9.2:eq-current-sets-a} at the making time
$t_{\mathbf{T}}(m)$ with $i=m$. At that time two identifications hold:
\begin{itemize}
\item the action domain of level $m$ is $\Lambda_m^{\mathbf{T}}$;
\item the native determining domain of level $m$ is $\Omega_m^{\mathbf{T}}$: the native
sequence is the one made by the $\mathbf{T}$-steps with healed classes deleted, and at the time
$t_{\mathbf{T}}(m)$ the $\mathbf{T}$-step $m-1\to m$ has just made $\Omega_m^{\mathbf{T}}$, while
no operation of the label that comes after it in the linear order has yet acted, so that no healed
class has been deleted from it; thus $\Omega_m^{\natural}=\Omega_m^{\mathbf{T}}$ at that time.
\end{itemize}
The right-hand inclusion of \eqref{9.2:eq-current-sets-a} therefore gives
\begin{equation}\label{9.2:eq-current-sets-4}
\Lambda_m^{\mathbf{T}}\subset\Omega_m^{\mathbf{T}} .
\end{equation}
Both sides are unions of cubes of $\mathcal{Q}_m$, by \eqref{9.2:eq-current-sets-a} and the
statement of \cite{B-CMP119} recalled above. Hence \eqref{9.2:eq-current-sets-4} holds cube-wise in
$\mathcal{Q}_m$, in the sense of \eqref{9.2:eq-cubes-regions-b}.

\emph{The nesting \eqref{9.2:eq-current-sets-b} at its birth.} Fix $i<k$.

By \eqref{9.2:eq-attempt-windows-a}, the $\mathbf{T}$-step $i\to i+1$ makes
$\Omega_{i+1}^{\mathbf{T}}$, $\Lambda_{i+1}^{\mathbf{T}}$, $Z_{i+1}^{\mathbf{T}}$ and nothing at
lower levels. Write $t_{\mathbf{T}}(i+1)^{-}$ for the stage just before $t_{\mathbf{T}}(i+1)$. Then
the set of level~$i$ at $t_{\mathbf{T}}(i+1)$ is the then-current set
$Z_i^{(t_{\mathbf{T}}(i+1)^{-})}$, and the set of level $i+1$ at that stage is
$Z_{i+1}^{\mathbf{T}}$, since a completed operation with $k_E\ge i+1$ acts after
$t_{\mathbf{T}}(k_E)\ge t_{\mathbf{T}}(i+1)$, so that none has yet been processed.

The inclusion $Z_i^{(t_{\mathbf{T}}(i+1)^{-})}\subset Z_{i+1}^{\mathbf{T}}$ is proved in the
statement on the ten-layer collars of the current sets from the ten-layer collar
\eqref{9.2:eq-ten-layer-collars-c}, read on $Z_i^{(t_{\mathbf{T}}(i+1)^{-})}$, and
\eqref{9.2:eq-current-sets-a} read at the making time $t_{\mathbf{T}}(i+1)$, namely
\eqref{9.2:eq-current-sets-4} with $m=i+1$ and the left-hand inclusion of
\eqref{9.2:eq-current-sets-a} at index $i$; we use it from there. Thus
\eqref{9.2:eq-current-sets-b} holds at $t=t_{\mathbf{T}}(i+1)$.

\emph{Persistence under deletions.} By \eqref{9.2:eq-attempt-windows-a}, for
$t\ge t_{\mathbf{T}}(i+1)$ nothing changes $Z_i$ or $Z_{i+1}$ except the deletion of reset classes
$E$ of completed $\mathbf{R}$-operations. Each deletion removes $E$ from $Z_m^{\mathrm{cur}}$ for
every $m\le k_E$ at which $E$ was large-field. We proceed by induction over the completed
operations processed after $t_{\mathbf{T}}(i+1)$, in their linear order. Suppose
$Z_i^{(t)}\subset Z_{i+1}^{(t)}$, and let the next such operation have class $E$.
Since it is processed after $t_{\mathbf{T}}(i+1)$, and an operation at step $k_E$ is processed
before the $\mathbf{T}$-step $k_E\to k_E+1$ reads the label, we have $k_E\ge i+1>i$. By the
description of $Z_m^{(t)}$ through \eqref{9.2:eq-current-sets-2}, $E$ is therefore deleted from
$Z_{i+1}$ and from $Z_i$ as well. For $m\in\{i,i+1\}$ write
$Z_m^{(t)}=\operatorname{cl}(G_m)$, with $G_m$ the set $|Z_m^{\mathbf{T}}|$ less the classes
processed before $t$; after the deletion the set of level $m$ is
$\operatorname{cl}(G_m\setminus|E|)$. Since $|E|$ is closed, the closure identity
\eqref{9.2:eq-current-sets-5}, applied with $G=G_m$ and $\mathcal{F}=|E|$, gives
\begin{equation*}
\operatorname{cl}(G_m\setminus|E|)=\operatorname{cl}\bigl(Z_m^{(t)}\setminus|E|\bigr).
\end{equation*}
Set difference with $|E|$ and closure both preserve inclusions, so
\begin{equation*}
\operatorname{cl}\bigl(Z_i^{(t)}\setminus|E|\bigr)
\subset\operatorname{cl}\bigl(Z_{i+1}^{(t)}\setminus|E|\bigr) .
\end{equation*}
Hence \eqref{9.2:eq-current-sets-b} holds at every stage $t\ge t_{\mathbf{T}}(i+1)$.

Finally, the stage now --- the reading of the label by the $\mathbf{T}$-step $k\to k+1$ --- comes
after $t_{\mathbf{T}}(i+1)$, because $i+1\le k$; hence $Z_i\subset Z_{i+1}$ now.

\emph{Native and local scale.} The native determining sequence is the one made by the
$\mathbf{T}$-steps, with whole healed components deleted. The inequality
$n^{\natural}(x)\le n(x)$ is part of the statement on the native determining sequence, and is used
here as stated there.
\end{proof}

\begin{lemma}[The ten-layer collars of the current sets]\label{9.2:lem-ten-layer-collars}
We use the cube partitions $\mathcal{Q}_i$ into cubes of side $M\check{R}_i\ell_i$, the regions, their
carriers $|\cdot|$, their components and the fattening $\mathcal{S}^{\sim n}$ in $\mathcal{Q}_i$ of
Notation~\ref{9.2:not-cubes-regions}, the $\ell^{\infty}$-distance $\operatorname{dist}_{\infty}$ of the
torus, and the step-made sets $Z_m^{\mathbf{T}}$, $\Lambda_m^{\mathbf{T}}$, $\Omega_m^{\mathbf{T}}$, the stages
$t$, $t_{\mathbf{T}}(m)$, $t^{-}$, the sets $Z_i^{(t)}$ and the current sets $Z_i=Z_i^{\mathrm{cur}}$ of
Definition~\ref{9.2:def-current-sets}. Let $k$ be the level of the term; the current set
$Z_k=Z_k^{\mathrm{cur}}$ is the set $Z_k^{(t_{\mathbf{T}}(k+1)^{-})}$ read by the step $k\to k+1$, and
$\Omega:=\Omega_{k+1}^{\mathbf{T}}$ is the determining domain of that step, assumed non-empty. Assume the
following.
\begin{itemize}
\item[(a)] \emph{Ten layers per step.} For every $i\in[1,k]$ and every component $D$ of the set
$Z_{i-1}^{(t_{\mathbf{T}}(i)^{-})}$ read by the step $i-1\to i$, that step produces a set $Z_i^{\mathbf{T}}$
containing $D$ together with at least ten layers of cubes of $\mathcal{Q}_i$ around $D$. This is the
sentence ``Each renormalization step adds at least ten layers of $MR_k$-cubes'' of \cite{B-CMP122a} (in
the notation of \cite{B-CMP122a}, whose $MR_k$-cubes are the cubes of $\mathcal{Q}_k$ here), in the
componentwise form in which it is used in the proof of the ancestry clause of the correlation theorem.
\item[(b)] \emph{The eight-layer step collar.} For every $m\le k+1$, let $Z'_{m-1}$ be the set of cubes of
$\mathcal{Q}_m$ meeting $Z_{m-1}^{(t_{\mathbf{T}}(m)^{-})}$, the large-field region which the term carries
into the step $m-1\to m$. Then, with $\sim$ the fattening in $\mathcal{Q}_m$,
\begin{equation}\label{9.2:eq-ten-layer-collars-b}
\begin{gathered}
\text{no cube of }\Omega_m^{\mathbf{T}}\text{ lies in }(Z'_{m-1})^{\sim 8},\\
\operatorname{dist}_{\infty}\bigl(|\Omega_m^{\mathbf{T}}|,\,Z_{m-1}^{(t_{\mathbf{T}}(m)^{-})}\bigr)
\ge 8M\check{R}_m\ell_m .
\end{gathered}
\end{equation}
This is the collar of \cite[(3.2)--(3.5)]{B-CMP119}, stated there for the native large-field region:
``The term has a large field region $Z_k=\Lambda_k^c$, and we denote by $Z_k'$ the union of all
$LMR_{k+1}$-cubes intersecting the region $Z_k$''; ``We surround $Z_k'$ by four layers of the
$LMR_{k+1}$-cubes''; then ``we surround the union $P''_{k+1}\cup {Z_k'}^{\sim 4}$ by a layer of
$LMR_{k+1}$-cubes'', followed by a further layer and by two layers of such cubes, $4+1+1+2=8$ layers in
all (the symbols $Z_k$, $Z_k'$, $P''_{k+1}$ are those of \cite{B-CMP119}, whose $LMR_{k+1}$-cubes are the
cubes of $\mathcal{Q}_{k+1}$ here). The form \eqref{9.2:eq-ten-layer-collars-b}, for the large-field region
carried by the term, in which whole healed components are deleted, is taken by statement from the
correlation theorem.
\item[(c)] \emph{Domains of the step}, \eqref{9.2:eq-current-sets-a}: for every $m$,
$\Lambda_m^{\mathbf{T}}\subset\Omega_m^{\mathbf{T}}$, and both are unions of cubes of $\mathcal{Q}_m$.
\item[(d)] \emph{The deletion rule} of \eqref{9.2:eq-attempt-windows-a}. After $t_{\mathbf{T}}(i)$ the set
$Z_i^{(t)}$ changes only at the stages $t_E$ of completed large-field operations $\mathbf{R}$. At such a
stage the reset class $E$, of step $k_E$, is deleted from $Z_i^{(t_E^{-})}$ for every $i\le k_E$, and $E$ is
a union of whole components of $Z_{k_E}^{(t_E^{-})}$. A completed operation resets its whole component to
one level: \cite[(1.33)]{B-CMP122b} reads ``We define $\mathbf{B}_k$ as equal to $\mathbf{B}''_k$ on $Z^c$,
and to $\{Z^{(k)}\}$ on $Z$'' (the symbols $\mathbf{B}_k$, $\mathbf{B}''_k$, $Z$, $Z^{(k)}$ being those of
\cite{B-CMP122b}). A preparation of an operation re-makes the primed sets $\Omega''_l(a)$, $Z''_l(a)$
inside its own component and leaves the unprimed $\Lambda_j$ unchanged \cite{B-CMP122a}. These last two
sentences are used in the form in which they appear in the proof of the ancestry clause of the
correlation theorem.
\end{itemize}
Then the following hold.
\begin{itemize}
\item[(i)] For every $i\in[1,k]$ and every component $D$ of the current set $Z_{i-1}$,
\begin{equation}\label{9.2:eq-ten-layer-collars-a}
\operatorname{dist}_{\infty}(D,\,Z_i^{c})\ge 10M\check{R}_i\ell_i ,
\end{equation}
and hence $\operatorname{dist}_{\infty}(Z_{i-1},Z_i^{c})\ge 10M\check{R}_i\ell_i$ for every $i\in[1,k]$, in
particular for every $i\in[j+1,k]$ whenever $0\le j<k$.
\item[(ii)] For every $m\le k+1$, no cube of $\Omega_m^{\mathbf{T}}$ meets
$Z_{m-1}^{(t_{\mathbf{T}}(m)^{-})}$, and
\begin{equation}\label{9.2:eq-ten-layer-collars-c}
\begin{gathered}
Z_{m-1}^{(t_{\mathbf{T}}(m)^{-})}\subset\operatorname{cl}\bigl((\Lambda_m^{\mathbf{T}})^{c}\bigr)
=Z_m^{\mathbf{T}},\qquad
Z_{m-1}^{(t)}\subset Z_m^{(t)}\quad(t\ge t_{\mathbf{T}}(m)),\\
|\Omega_m^{\mathbf{T}}|\cap Z_i^{(t)}=\emptyset\quad(i\le m-1,\ t\ge t_{\mathbf{T}}(m)),
\qquad\text{in particular}\quad |\Omega_m^{\mathbf{T}}|\cap Z_i^{\mathrm{cur}}=\emptyset .
\end{gathered}
\end{equation}
\item[(iii)] For every $m\le k+1$,
\begin{equation}\label{9.2:eq-ten-layer-collars-d}
\text{every cube of }(Z'_{m-1})^{\sim 8}\text{ is a cube of }Z_m^{\mathbf{T}} .
\end{equation}
\item[(iv)] At $m=k+1$,
\begin{equation}\label{9.2:eq-ten-layer-collars-e}
\begin{gathered}
\operatorname{dist}_{\infty}(Z_k,|\Omega|)\ge 8M\check{R}_{k+1}\ell_{k+1}=8LM\check{R}_{k+1}\ell_k,
\qquad |\Omega|\cap Z_k=\emptyset,\\
\emptyset\neq|\Omega|\subset Z_k^{c}\subset Z_i^{c}\quad\text{for every }i\le k .
\end{gathered}
\end{equation}
\end{itemize}
\end{lemma}

\begin{proof}
\emph{Proof of \eqref{9.2:eq-ten-layer-collars-a}.} Fix $i\in[1,k]$ and a component $D$ of the current set
$Z_{i-1}$.

First, $D$ is a component of $Z_{i-1}^{(t)}$ for every stage $t\ge t_{\mathbf{T}}(i-1)$. Indeed, by (d) a
preparation does not change the unprimed sets, and the only changes of $Z_{i-1}^{(t)}$ after its birth are
deletions of classes $E$ with $k_E\ge i-1$. Such a class is a union of whole components of
$Z_{k_E}^{(t_E^{-})}$, and $Z_{k_E}^{(t_E^{-})}\supseteq Z_{i-1}^{(t_E^{-})}$ by the nesting
\eqref{9.2:eq-current-sets-b}. So, by \eqref{9.2:eq-cubes-regions-c}, each deletion removes whole components
of the then $Z_{i-1}$ and leaves the others untouched. A component of the set which remains after all
deletions, in particular $D$, was therefore a component of $Z_{i-1}^{(t)}$ at every earlier stage
$t\ge t_{\mathbf{T}}(i-1)$. In particular, $D$ is a component of the set $Z_{i-1}^{(t_{\mathbf{T}}(i)^{-})}$
read by the step $i-1\to i$.

Hence, by (a), that step made $Z_i^{\mathbf{T}}\supseteq D$ together with at least ten layers of cubes of
side $M\check{R}_i\ell_i$ around $D$. By Notation~\ref{9.2:not-cubes-regions}, $D$ is the carrier of a
connected region of the finer family $\mathcal{Q}_{i-1}$, which is compatible with $\mathcal{Q}_i$. So $D$
together with these layers forms a connected region of $\mathcal{Q}_i$ contained in $Z_i^{\mathbf{T}}$. By
\eqref{9.2:eq-cubes-regions-a}, this region lies in one component $C_i$ of $Z_i^{\mathbf{T}}$.

We show by induction along the stages $t\ge t_{\mathbf{T}}(i)$ that $C_i\subset Z_i^{(t)}$. At
$t=t_{\mathbf{T}}(i)$ this holds because $Z_i^{(t)}=Z_i^{\mathbf{T}}$. Suppose it holds just before a stage
$t_E$ at which, by (d), a class $E$ with $k_E\ge i$ is deleted from $Z_i$. Then $E$ is a union of whole
components of $Z_{k_E}^{(t_E^{-})}$, and
\begin{equation*}
C_i\subset Z_i^{(t_E^{-})}\subset Z_{k_E}^{(t_E^{-})}
\end{equation*}
by the nesting \eqref{9.2:eq-current-sets-b}. Since $C_i$ is connected, it lies in one component of
$Z_{k_E}^{(t_E^{-})}$. By \eqref{9.2:eq-cubes-regions-c}, there are two cases.
\begin{itemize}
\item[--] $E$ contains that component, hence $C_i$, and with it $D$. Then, since $i-1\le k_E$, $D$ is
deleted from $Z_{i-1}$ as well, and $D$ is not a component of the current $Z_{i-1}$. This contradicts the
choice of $D$, so this case does not occur.
\item[--] $E$ is a union of other whole components. Then $E$ is disjoint from $C_i$, and $C_i$ remains in
$Z_i^{(t_E)}$.
\end{itemize}
This completes the induction. So $C_i$, and with it the ten layers of cubes around $D$, lies in the current
set $Z_i$. Every point of $Z_i^{c}$ therefore lies outside these ten layers of cubes of side
$M\check{R}_i\ell_i$ around $D$, and has $\ell^{\infty}$-distance at least $10M\check{R}_i\ell_i$ from $D$.
This is \eqref{9.2:eq-ten-layer-collars-a} with the current $Z_i$.

Since $Z_{i-1}$ is the union of its components, the distance of $Z_{i-1}$ to $Z_i^{c}$ is the minimum over
these components. So $\operatorname{dist}_{\infty}(Z_{i-1},Z_i^{c})\ge 10M\check{R}_i\ell_i$ for every
$i\in[1,k]$, in particular for every $i\in[j+1,k]$ whenever $0\le j<k$.

\emph{Proof of \eqref{9.2:eq-ten-layer-collars-c}.} Let $m\le k+1$. If a cube of $\Omega_m^{\mathbf{T}}$ met
$Z_{m-1}^{(t_{\mathbf{T}}(m)^{-})}$, it would be one of the cubes forming $Z'_{m-1}$. It would therefore
lie in $(Z'_{m-1})^{\sim 8}\supseteq Z'_{m-1}$, contrary to the first line of
\eqref{9.2:eq-ten-layer-collars-b}. So no cube of $\Omega_m^{\mathbf{T}}$ meets
$Z_{m-1}^{(t_{\mathbf{T}}(m)^{-})}$.

By (c), $\Lambda_m^{\mathbf{T}}\subset\Omega_m^{\mathbf{T}}$ is a union of cubes of $\mathcal{Q}_m$. So
$|\Lambda_m^{\mathbf{T}}|\cap Z_{m-1}^{(t_{\mathbf{T}}(m)^{-})}=\emptyset$, that is,
\begin{equation*}
Z_{m-1}^{(t_{\mathbf{T}}(m)^{-})}\subset(\Lambda_m^{\mathbf{T}})^{c}
\subset\operatorname{cl}\bigl((\Lambda_m^{\mathbf{T}})^{c}\bigr)=Z_m^{\mathbf{T}},
\end{equation*}
the last equality being the definition of $Z_m^{\mathbf{T}}$ in
Definition~\ref{9.2:def-current-sets}. This is the nesting of Definition~\ref{9.2:def-current-sets} at the
birth of $Z_m$, since $Z_m^{(t_{\mathbf{T}}(m))}=Z_m^{\mathbf{T}}$.

It persists for $t\ge t_{\mathbf{T}}(m)$ by the deletion rule (d). A class deleted from $Z_m$ has
$k_E\ge m>m-1$, so it is deleted from $Z_{m-1}$ at the same stage. A class deleted from $Z_{m-1}$ alone has
$k_E=m-1$ and removes points from the smaller set only. Each deletion therefore preserves the inclusion,
and $Z_{m-1}^{(t)}\subset Z_m^{(t)}$ for all $t\ge t_{\mathbf{T}}(m)$.

Now let $i\le m-1$ and $t\ge t_{\mathbf{T}}(m)$. The nesting at the levels below $m$ gives
$Z_i^{(t)}\subset Z_{m-1}^{(t)}$. By (d) the sets only shrink after their birth, so
$Z_{m-1}^{(t)}\subset Z_{m-1}^{(t_{\mathbf{T}}(m)^{-})}$. Since no cube of $\Omega_m^{\mathbf{T}}$ meets
the last set, $|\Omega_m^{\mathbf{T}}|\cap Z_i^{(t)}=\emptyset$. Taking for $t$ the current stage gives
$|\Omega_m^{\mathbf{T}}|\cap Z_i^{\mathrm{cur}}=\emptyset$.

\emph{Proof of \eqref{9.2:eq-ten-layer-collars-d}.} Let $Q$ be a cube of $(Z'_{m-1})^{\sim 8}$. By the first
line of \eqref{9.2:eq-ten-layer-collars-b}, $Q$ is not a cube of $\Omega_m^{\mathbf{T}}$, hence by (c) not a
cube of $\Lambda_m^{\mathbf{T}}\subset\Omega_m^{\mathbf{T}}$. By \eqref{9.2:eq-cubes-regions-b}, $Q$ is then a
cube of $\operatorname{cl}((\Lambda_m^{\mathbf{T}})^{c})=Z_m^{\mathbf{T}}$.

\emph{Proof of \eqref{9.2:eq-ten-layer-collars-e}.} Take $m=k+1$, so that $\Omega_m^{\mathbf{T}}=\Omega$.
The set $Z_{m-1}^{(t_{\mathbf{T}}(m)^{-})}=Z_k^{(t_{\mathbf{T}}(k+1)^{-})}$ is the current $Z_k$. The second
line of \eqref{9.2:eq-ten-layer-collars-b} gives
$\operatorname{dist}_{\infty}(Z_k,|\Omega|)\ge 8M\check{R}_{k+1}\ell_{k+1}$. Since $\ell_{k+1}=L^{k+1}\eta$
and $\ell_k=L^{k}\eta$, we have $\ell_{k+1}=L\ell_k$, and the bound equals $8LM\check{R}_{k+1}\ell_k$.

By \eqref{9.2:eq-ten-layer-collars-c} at $m=k+1$, no cube of $\Omega$ meets $Z_k$, so
$|\Omega|\cap Z_k=\emptyset$, that is, $|\Omega|\subset Z_k^{c}$. For $i\le k$ the nesting of the current
sets, \eqref{9.2:eq-current-sets-b}, gives $Z_i\subset Z_k$, hence $Z_k^{c}\subset Z_i^{c}$. Together with
$\Omega\neq\emptyset$, hence $|\Omega|\neq\emptyset$, this is the second line of
\eqref{9.2:eq-ten-layer-collars-e}.

Of the eight layers of the cube form, the first line of \eqref{9.2:eq-ten-layer-collars-b}, the applications
of this lemma later in this chapter use at most seven. The bound \eqref{9.2:eq-ten-layer-collars-e} uses the
distance form, the second line of \eqref{9.2:eq-ten-layer-collars-b}, with its factor $8$.
\end{proof}

\begin{definition}[Monotonicity and floor of the radii]\label{9.2:not-radii-floor}
Let $(\check{R}_i)_i$ be the scaffold's radii, $\check{R}_i$ being the radius at level $i$. For a
level $k$ of the scaffold we put $\check{R}^{\flat}:=\check{R}_k$. Assume the following.
\begin{enumerate}
\item[(i)] The radii are non-increasing in the level and decrease by at most the factor $L$ per
level: for every $i$,
\begin{equation}\label{9.2:eq-radii-floor-a}
\text{(a)}\quad \check{R}_{i+1}\le\check{R}_i,\qquad\qquad
\text{(b)}\quad \check{R}_i\le L\,\check{R}_{i+1}.
\end{equation}
\item[(ii)] The radius floor of the correlation theorem holds: for every level $k$ of the
scaffold, with the threshold $\check{\mathsf{T}}_k$ of the scaffold, the degree exponent
$r_{\mathrm{pr}}$ and $c_5=2c_2$,
\begin{equation}\label{9.2:eq-radii-floor-b}
\bigl(\log(x_k-c_5)\bigr)^{r_{\mathrm{pr}}}\le\check{\mathsf{T}}_k^{\,r_{\mathrm{pr}}}\le\check{R}_k .
\end{equation}
\item[(iii)] The coupling ceiling on the scaffold holds: $g_k\le\gamma$ for every level $k$ of the
scaffold.
\item[(iv)] For every level $k$ of the scaffold,
\begin{equation}\label{9.2:eq-radii-floor-1}
\check{R}_k\ge2 .
\end{equation}
\end{enumerate}
Then, for every level $k$, $\check{R}_i\ge\check{R}_k$ for every $i\le k$, and
$L\check{R}_{k+1}\ge\check{R}_k$. Each
collar between levels $j$ and $k$ is read at its own radius $\check{R}_i$, $j\le i\le k$, and
$\check{R}_i\ge\check{R}^{\flat}$. Moreover, for every $k$,
\begin{equation*}
\check{R}_k\ge\bigl(\log(\gamma^{-2}-c_5)\bigr)^{r_{\mathrm{pr}}}.
\end{equation*}
In what follows the floor \eqref{9.2:eq-radii-floor-1} is used; the floor $\check{R}_{k+1}\ge2$
enters only in the optional sharpening of the weak radius bound.
\end{definition}

\begin{proof}[Proof of the consequences]
For $i\le k$, applying \eqref{9.2:eq-radii-floor-a}(a) at the levels $i,i+1,\dots,k-1$ gives the
chain
\begin{equation*}
\check{R}_k\le\check{R}_{k-1}\le\dots\le\check{R}_i .
\end{equation*}
Taking $i=k$ in \eqref{9.2:eq-radii-floor-a}(b) gives $\check{R}_k\le L\check{R}_{k+1}$. For a
collar between $j$ and $k$, its level $i$ satisfies $i\le k$, so $\check{R}_i\ge\check{R}_k
=\check{R}^{\flat}$ by the chain above.

By (iii), $x_k=g_k^{-2}\ge\gamma^{-2}$, hence $x_k-c_5\ge\gamma^{-2}-c_5$. Since $\log$ is
increasing and $t\mapsto t^{r_{\mathrm{pr}}}$ is non-decreasing
on $[0,\infty)$, we obtain
\begin{equation*}
\bigl(\log(\gamma^{-2}-c_5)\bigr)^{r_{\mathrm{pr}}}\le\bigl(\log(x_k-c_5)\bigr)^{r_{\mathrm{pr}}}
\le\check{R}_k ,
\end{equation*}
the last inequality being \eqref{9.2:eq-radii-floor-b}. This holds for every $k$.
\end{proof}

\begin{definition}[Level-$j$ boundary families and their size]\label{9.2:def-boundary-families}
We use the cube partitions and regions of Notation~\ref{9.2:not-cubes-regions} and the step
$k\to k+1$ with its stored boundary terms of Notation~\ref{9.2:not-frozen-step}. Let $j$ be the
level of a stored boundary term, and let $\pi_j$ be the partition into cubes of side $M\ell_j$.

\emph{Domains.} $\mathbf{D}_j$ is Ba{\l}aban's family of localisation domains of level $j$
\cite[(2.41)]{B-CMP119}. The stored $\mathbf{T}$-born boundary terms of level $j$ present at the
step $k\to k+1$ are indexed by those $X\in\mathbf{D}_j$ with $X\cap\Omega_j\neq\emptyset$ and
$X\cap Z_j^{\sim}\neq\emptyset$; these are the domains over which the sum of
\cite[(2.41)]{B-CMP119} runs. Here $\Omega_j$ is the fluctuation domain of the step to level $j$
(for the step $k\to k+1$ this is $\Omega_{k+1}$ of Notation~\ref{9.2:not-frozen-step}), and
$Z_j^{\sim}$ is the fattening of the level-$j$ large-field region in the sense of
\cite[(2.41)]{B-CMP119}. Below, $X\in\mathbf{D}_j$ always ranges over these domains. Each such
$X$ is a connected union of closed cubes of $\pi_j$.

\emph{Tree length.} For $X\in\mathbf{D}_j$, $d_j(X)\ge 0$ is the tree length of $X$ of
\cite{B-CMP119}, measured in units of $M\ell_j$. By the volume and filament bound, it satisfies
\begin{equation}\label{9.2:eq-boundary-families-a}
  d_j(X)\;\ge\;\frac{\operatorname{diam}_\infty(X)}{M\ell_j}-2 ,
\end{equation}
where $\operatorname{diam}_\infty$ is the diameter in the $\ell^\infty$ metric of the torus.

\emph{Families and their size.} Let $f_j\ge 0$. A \emph{level-$j$ $\mathbf{B}$-family} of
\emph{size} $f_j$ is an assignment $X\mapsto F^{(j)}(X)$ of a value to each $X\in\mathbf{D}_j$
such that
\begin{equation}\label{9.2:eq-boundary-families-1}
  \bigl|F^{(j)}(X)\bigr|\;\le\; f_j\, e^{-\kappa d_j(X)}\qquad (X\in\mathbf{D}_j).
\end{equation}
Two families enter below: $F=\delta\mathbf{B}$, with
$f_j=\delta^{\mathbf{B}}_j:=\|\delta\mathbf{B}^{(j)}\|$, the norm, in the sense of the one-lattice
Lipschitz step, of the difference of the stored level-$j$ boundary terms along the segment
$[\mathfrak{h},\tilde{\mathfrak{h}}]\subset\mathfrak{K}_k^{+}$ of that step; and $F=\mathbf{B}$,
with $f_j=\tfrac{9}{8}B_0$, where $B_0$ is the constant in Ba{\l}aban's exponential bound on the
stored boundary terms $\mathbf{B}^{(j)}(X)$.

\emph{Overlap constant.} We put
\begin{equation}\label{9.2:eq-boundary-families-b}
  K_0^{(2)}\;:=\;\sup_{\square\in\pi_j}\;\sum_{\substack{X\in\mathbf{D}_j\\ X\supset\square}}
  e^{-\kappa d_j(X)/2}.
\end{equation}
By the volume and filament bound, $K_0^{(2)}$ is a $(d,\kappa)$-number, that is, a number
depending only on $d$ and $\kappa$. These two facts, the inequality
\eqref{9.2:eq-boundary-families-a} and the $(d,\kappa)$-dependence of the constant
\eqref{9.2:eq-boundary-families-b}, are used below under those labels.
\end{definition}

\begin{lemma}[Anchors of surviving stored boundary terms]\label{9.2:lem-anchors}
Let $j\le k$. Write $Z_j$ for the current large-field set of level $j$, and $\Lambda_j$
for the action domain made by the $\mathbf{T}$-step of
level $j$ (Definition~\ref{9.2:def-current-sets}). Let $X\in\mathbf{D}_j$ be a unit of a
level-$j$ boundary family (Definition~\ref{9.2:def-boundary-families}), that is, the
localisation domain of a stored $\mathbf{T}$-born boundary term $\mathbf{B}^{(j)}(X)$ of level
$j$ alive at the current time. Lengths are $\ell^\infty$-lengths on the torus, and the
$\pi_j$-cubes have side $M\ell_j$ (Notation~\ref{9.2:not-cubes-regions}).
\begin{enumerate}
\item[(a)] \emph{Strong anchor.} Assume the following two statements.
\begin{itemize}
\item[(i)] Every $\mathbf{T}$-born boundary term $\mathbf{B}^{(j)}(X)$ has
$X\not\subset\Lambda_j$, that is, $X$ shares a closed $\pi_j$-cube with the large-field set
of level $j$ itself, as it stood when the term was born, and not merely with a fattening
of it. This is the birth statement for stored boundary
terms used in the proof of the correlation theorem.
\item[(ii)] Let a later operation $\mathbf{R}$ reset a class $E$ of the large-field set. A
stored term whose localisation domain meets $E$ lies in the first part of that operation
$\mathbf{R}$. According to \cite{B-CMP122b}, the first part includes every term whose
localisation domain intersects the region $Z$ of the operation. A term of the first part
undergoes the operations \cite[(1.59), (1.60)]{B-CMP122b} and remains among the terms of
\cite[(1.70), (1.71)]{B-CMP122b}; in particular, after that operation it is no longer carried
as a stored $\mathbf{T}$-born boundary term.
\end{itemize}
Then there are $y_0\in X$ and $z_0\in Z_j$ (the current set) with
\begin{equation}\label{9.2:eq-anchors-a}
|y_0-z_0|_\infty\le M\ell_j .
\end{equation}
\item[(b)] \emph{Weak anchor.} Assume instead only the following statement.
\begin{itemize}
\item[(iii)] Every stored boundary-type term $F$ of level $j\le k$ alive at the current time
has, in the current label, a cube of $A_F$ in $(Z_j^{\boxplus})^{(1)}$, where the set $A_F$
attached to $F$ satisfies $A_F\subset X$, $X$ being the localisation domain of $F$; the set
$Z_j^{\boxplus}$ is $Z_j$ fattened by $\check{R}_j$ layers of $\pi_j$-cubes; and the
superscript $(1)$ adds one further layer of $\pi_j$-cubes. These readings of $A_F$,
${}^{\boxplus}$ and $(1)$ are taken as part of this hypothesis. This is the displayed anchoring
clause for stored boundary-type terms in the correlation theorem.
\end{itemize}
Then there are $y_0\in X$, $z_0\in Z_j$ and a number $a\ge0$ with
\begin{equation}\label{9.2:eq-anchors-b}
|y_0-z_0|_\infty\le(1+a)M\ell_j,\qquad a\le\check{R}_j\le L\check{R}_{j+1}.
\end{equation}
If moreover $j=k$, then $a\le\check{R}_k$.
\end{enumerate}
Under the hypotheses of (a), the bound \eqref{9.2:eq-anchors-b} holds with $a=0$.
\end{lemma}

\begin{proof}
\emph{Part (a).} Let $\mathbf{B}^{(j)}(X)$ be the stored $\mathbf{T}$-born term with domain $X$,
and consider the step producing level $j$, at which it is born. By hypothesis (i),
$X\not\subset\Lambda_j$. This is the classification of \cite[(3.26), (3.28)]{B-CMP119}, read
for that step: the small-field terms $\mathbf{E}$ and the terms denoted $R''$ in
\cite[(3.26), (3.28)]{B-CMP119} are those with $X\subset\Lambda_j$, and, as stated in
\cite{B-CMP119}, the remaining terms are included into the boundary term. So $X$ shares a
closed $\pi_j$-cube $Q$ with the large-field set of level $j$ as it stood when the term was
born, and $Q$ lies in a component of that set.

We show that $Q$ still lies in the current set $Z_j$. Suppose that a
later operation $\mathbf{R}$ had reset a class $E$ containing the cube $Q$. Then $X$ meets $E$,
since $Q$ lies in both $X$ and $E$. By statement (ii), the term $\mathbf{B}^{(j)}(X)$
would then lie in the first part of that $\mathbf{R}$, would undergo the operations
\cite[(1.59), (1.60)]{B-CMP122b}, and would remain among the terms of
\cite[(1.70), (1.71)]{B-CMP122b}. By the last sentence of statement (ii), it would
therefore not survive as a stored $\mathbf{T}$-born term. This contradicts the assumption that
$X$ is a unit, that is, the domain of such a term alive at the current time.

Put differently, a surviving term whose shared cube lay in a component of $Z_j$ healed since
its birth would have met the healed region $Z$ and would have ceased to be stored. Hence $Q$
lies in a still-current component of $Z_j$.

Now choose $y_0\in X\cap Q$ and $z_0\in Z_j\cap Q$. Both points lie in the closed $\pi_j$-cube
$Q$, whose side is $M\ell_j$, so $|y_0-z_0|_\infty\le M\ell_j$. This is
\eqref{9.2:eq-anchors-a}, and it is \eqref{9.2:eq-anchors-b} with $a=0$.

\emph{Part (b).} Let $F$ be the stored term whose localisation domain is $X$. By statement
(iii), the set $A_F\subset X$ contains a $\pi_j$-cube $Q'$ of
$(Z_j^{\boxplus})^{(1)}$. The set $Z_j^{\boxplus}$ adds $\check{R}_j$ layers of $\pi_j$-cubes
to $Z_j$, and the superscript $(1)$ adds one more layer. Therefore $Q'$ lies within
$\check{R}_j+1$ layers of $\pi_j$-cubes around $Z_j$.

Put $Z_j^{\sim 0}:=Z_j$, and call the $n$-th layer, $n\ge1$, the set of $\pi_j$-cubes of
$Z_j^{\sim n}$ not in $Z_j^{\sim(n-1)}$ (Notation~\ref{9.2:not-cubes-regions}); each such cube
meets a cube of $Z_j^{\sim(n-1)}$. We show by induction on $n$ that every point of a cube of
$Z_j^{\sim n}$ is at $\ell^\infty$-distance at most $nM\ell_j$ from $Z_j$. For $n=0$ the cube
is a cube of $Z_j$ and the distance is $0$. For $n\ge1$, a cube of the $n$-th layer is a
closed $\pi_j$-cube meeting a cube of $Z_j^{\sim(n-1)}$. Each of its points is therefore within
the side $M\ell_j$ of a point of that cube, which by induction lies within $(n-1)M\ell_j$ of
$Z_j$.

Take any $y_0\in Q'\subset X$. Since $Z_j$ is a finite union of closed cubes of the torus,
there is a $z_0\in Z_j$ at which the distance from $y_0$ to $Z_j$ is attained. By the
induction,
\begin{equation*}
|y_0-z_0|_\infty\le(\check{R}_j+1)M\ell_j ,
\end{equation*}
which is \eqref{9.2:eq-anchors-b} with $a:=\check{R}_j$.

The bound $\check{R}_j\le L\check{R}_{j+1}$ is part (b) of \eqref{9.2:eq-radii-floor-a}
(Notation~\ref{9.2:not-radii-floor}).

If $j=k$, then $\check{R}_j=\check{R}_k$, and so $a\le\check{R}_k$.
\end{proof}

\begin{definition}[Scaled distance to the fluctuation domain]\label{9.2:def-scaled-distance}
Lengths are measured in the $\ell^\infty$ metric $|\cdot|_\infty$ of the torus, as in
Notation~\ref{9.2:not-cubes-regions}. Let $j\le k$ be the levels fixed in
Definition~\ref{9.2:def-current-sets}, and for $j\le i\le k$ let $Z_i:=Z_i^{\mathrm{cur}}$
be the current large-field set of scale $i$ of that definition.

\emph{Depth.} The \emph{depth} of a point $y$ is
\begin{equation*}
\iota(y):=
\begin{cases}
\min\{\, i\in[j,k] : y\in Z_i \,\}, & y\in Z_k,\\
k+1, & y\notin Z_k .
\end{cases}
\end{equation*}
For $\iota(y)\le i\le k$ we write $A_i(y)$ for the connected component of $Z_i$ that contains
$y$. The depth of a non-empty set $X$ is $\iota(X):=\max_{y\in X}\iota(y)$. Since $\iota$ takes
values in the finite set $\{j,\dots,k+1\}$, $\iota(X)$ is finite and the maximum is attained.

\emph{Scaled distance.} Let $\Omega$, the fluctuation domain, be the set of scale $k+1$ of the
current determining
sequence of Definition~\ref{9.2:def-current-sets} at level $k$. Ba{\l}aban's construction in
\cite{B-CMP122a} measures distances to the determining domain by a scaled distance defined in
terms of $M_1$-cubes on the corresponding scales. In this book we do not use that definition
directly. Instead we use the following reading of it, which counts $M$-cubes of the local scale,
and which is the definition of $d_{\mathrm{sc}}$ throughout what follows.

A point $x$ is said to \emph{lie only in current cubes of scale $\le m$} if every cube of the
current sets of Definition~\ref{9.2:def-current-sets} that contains $x$ has scale at most $m$.

A \emph{path} is a continuous map $\Gamma$ from an interval $[\alpha_0,\beta_0]$,
$\alpha_0\le\beta_0$, into the torus; for $\alpha_0=\beta_0$ it is a single point, and the path
is called \emph{degenerate}.
A \emph{scaled length} is a function $\ell_{\mathrm{sc}}$ that assigns to every path
$\Gamma$ a number $\ell_{\mathrm{sc}}(\Gamma)\ge 0$ and satisfies
the conditions (D1) and (D2) below; (D3) is not a condition on $\ell_{\mathrm{sc}}$ but records
the standing property of the current cubes at level $k$ that is used together with them.
Conditions (D1) and (D2) are required for every path $\Gamma$ on $[\alpha_0,\beta_0]$, every
partition $\alpha_0=t_0<t_1<\dots<t_n=\beta_0$ (when $\alpha_0<\beta_0$), every sub-path
$\Gamma|_{[\alpha,\beta]}$ with $\alpha_0\le\alpha<\beta\le\beta_0$, every $S\ge 0$, and every
scale $m$:
\begin{equation}\label{9.2:eq-scaled-distance-a}
\begin{aligned}
&\text{(D1)}\quad \ell_{\mathrm{sc}}(\Gamma)\ \ge\ \sum_{l=1}^{n}
  \ell_{\mathrm{sc}}\bigl(\Gamma|_{[t_{l-1},t_l]}\bigr);\\
&\text{(D2)}\quad \ell_{\mathrm{sc}}\bigl(\Gamma|_{[\alpha,\beta]}\bigr)\ \ge\ \frac{S}{M\ell_m}
  \quad\text{if } |\Gamma(\beta)-\Gamma(\alpha)|_\infty\ge S \text{ and each } \Gamma(t),\
  \alpha\le t<\beta,\\
&\qquad\quad \text{lies only in current cubes of scale}\le m;\\
&\text{(D3)}\quad \text{every point } x\notin\Omega \text{ lies only in current cubes of
  scale}\le k .
\end{aligned}
\end{equation}
Condition (D3) expresses that $\Omega$ is the only set of scale $k+1$ at level $k$.

A path $\Gamma$ on $[\alpha_0,\beta_0]$ is called \emph{admissible from $y$} if three things
hold: $\Gamma(\alpha_0)=y$, $\Gamma(\beta_0)\in\Omega$, and $\Gamma(t)\notin\Omega$ for
$\alpha_0\le t<\beta_0$. In other words,
it is a path from $y$ to $\Omega$ stopped at its first point of $\Omega$. If $y\in\Omega$, the
paths admissible from $y$ are exactly the degenerate paths at $y$, since a non-degenerate one
would violate the last condition at $t=\alpha_0$; if $y\notin\Omega$, every path admissible from
$y$ is non-degenerate. The \emph{scaled
distance} of $y$ to $\Omega$ is
\begin{equation}\label{9.2:eq-scaled-distance-1}
d_{\mathrm{sc}}(y,\Omega):=\inf\{\,\ell_{\mathrm{sc}}(\Gamma) : \Gamma \text{ admissible from } y\,\},
\qquad \inf\emptyset:=+\infty .
\end{equation}

Suppose the scaled length is counted in $M_1$-cubes, that is, (D2) holds with $M_1$ in place of
$M$. Since $M_1\le M$, we have $S/(M_1\ell_m)\ge S/(M\ell_m)$, so (D2) as displayed also holds.
Thus (D2) with $M$ is the weaker of the two requirements.

\emph{Upper-bound clause.} In addition to (D1)--(D3), the reading includes the following
one-sided upper bound. If $x$ and $x'$ lie in one local $M$-cube, that is, $x'$ lies in the cube
of the partition $\pi_n$ into $M\ell_n$-cubes that contains $x$, where $n=n(x)$ is the local
scale of $x$ of Definition~\ref{9.2:def-current-sets}, then
\begin{equation}\label{9.2:eq-scaled-distance-3}
d_{\mathrm{sc}}(x,\Omega)\ \le\ d_{\mathrm{sc}}(x',\Omega)+1 .
\end{equation}
In other words, passing between points of one local $M$-cube costs at most one unit of
$d_{\mathrm{sc}}$. This clause is not derived from (D1)--(D3); it is part of the reading, taken
by statement. The bound \eqref{9.2:eq-scaled-distance-3} is kept separate from
\eqref{9.2:eq-scaled-distance-a}. It is used only in bounds on constants, never in a lower bound
on $d_{\mathrm{sc}}$.

Write $\operatorname{dist}_\infty(y,\Omega):=\inf_{z\in\Omega}|y-z|_\infty$. Then for every
point $y$,
\begin{equation}\label{9.2:eq-scaled-distance-2}
d_{\mathrm{sc}}(y,\Omega)\ \ge\ \frac{\operatorname{dist}_\infty(y,\Omega)}{M\ell_k},
\end{equation}
and this is (D2) at $m=k$.
\end{definition}

\begin{proof}[Proof of \eqref{9.2:eq-scaled-distance-2}]
Suppose first that no path is admissible from $y$. Then the left side of
\eqref{9.2:eq-scaled-distance-2} is $+\infty$ by the convention $\inf\emptyset=+\infty$ in
\eqref{9.2:eq-scaled-distance-1}, and there is nothing to prove.

Suppose next that $y\in\Omega$. Then $\operatorname{dist}_\infty(y,\Omega)=0$. The degenerate
path at $y$ is admissible from $y$, so the set in \eqref{9.2:eq-scaled-distance-1} is not empty,
and the left side is an infimum of numbers $\ell_{\mathrm{sc}}(\Gamma)\ge 0$. So it is
non-negative, and the inequality holds.

Now let $y\notin\Omega$, and let $\Gamma$ on $[\alpha_0,\beta_0]$ be admissible from $y$; as
noted above, $\alpha_0<\beta_0$. Put $S:=\operatorname{dist}_\infty(y,\Omega)$.

Since $\Gamma(\beta_0)\in\Omega$, the definition of $\operatorname{dist}_\infty$ gives
\begin{equation*}
|\Gamma(\beta_0)-\Gamma(\alpha_0)|_\infty=|\Gamma(\beta_0)-y|_\infty\ \ge\ S .
\end{equation*}

By admissibility, $\Gamma(t)\notin\Omega$ for $\alpha_0\le t<\beta_0$. By (D3) in
\eqref{9.2:eq-scaled-distance-a}, each such point $\Gamma(t)$ therefore lies only in current
cubes of scale $\le k$.

Now apply (D2) to the sub-path $\Gamma|_{[\alpha_0,\beta_0]}=\Gamma$, with $\alpha=\alpha_0$,
$\beta=\beta_0$ and $m=k$.
Its two hypotheses were just verified, so
\begin{equation*}
\ell_{\mathrm{sc}}(\Gamma)\ \ge\ \frac{S}{M\ell_k}.
\end{equation*}
This bound holds for every path admissible from $y$. Taking the infimum over all such $\Gamma$
in \eqref{9.2:eq-scaled-distance-1} gives \eqref{9.2:eq-scaled-distance-2}.
\end{proof}

We work at the frozen step $k\to k+1$ of Notation~\ref{9.2:not-frozen-step}. Its fluctuation
domain is $\Omega=\Omega_{k+1}\neq\emptyset$. We take a level-$j$ boundary family $F$, with values
$F^{(j)}(X)$ on the localisation domains $X\in\mathbf{D}_j$, its constant $f_j$ and the tree
length $d_j(X)$, as in Definition~\ref{9.2:def-boundary-families}. The scaled distance
$d_{\mathrm{sc}}(\cdot,\Omega)$ is that of Definition~\ref{9.2:def-scaled-distance}.

\emph{The response profile.} Let $v$ be a unit of the expansion of the correlation theorem, and let
$S_v$ be its read set. Let $p=(\square,Y_0)$ be a response piece of $v$, in the sense of the
Cauchy-estimate lemma of the correlation theorem. Here $\square$ is the block of the piece, and
$Y_0$ is the set that labels it together with $\square$ in that theorem. We measure the piece in the
unit $\mathsf{U}:=B_HK_Rg_k\sup|\mathsf{B}|$. This unit is fixed in the correlation theorem and is
used here as fixed there. The profile of the piece is
\begin{equation*}
e_p(y):=\sup_{\mathfrak{M}}\Theta_0^{\mathrm{loc}}(\delta H_p)(y)/\mathsf{U}.
\end{equation*}
In this formula $\delta H_p$ is the piece of the unit shift belonging to $p$. The functional
$\Theta_0^{\mathrm{loc}}$ and the set $\mathfrak{M}$ over which the supremum runs are those of the
correlation theorem. By the response statement at the step, $e_p$ obeys three bounds:
\begin{itemize}
\item[(a)] $e_p\le\eta_p$ on the read set $S_v$. Here $\eta_p$ is the correlation theorem's decay
factor between $\square$ and $S_v$. It is a function of the pair $(\square,S_v)$ with
$0<\eta_p\le1$, and it equals $1$ at every touching block.
\item[(b)] For a constant $C_b\ge1$ and a rate $\delta_1^{\mathrm{sc}}>0$,
\begin{equation*}
e_p(y)\le C_b\,e^{-\delta_1^{\mathrm{sc}}(d_{\mathrm{sc}}(y,\Omega)-1)}.
\end{equation*}
\item[(c)] A third bound, on the logarithm of the law factor on the Cauchy disc. It is not used
here and is not reproduced.
\end{itemize}
Only (b) enters below.

\emph{The Cauchy radius.} The units of the summation of boundary terms at finer zones have as
localisation domains the domains $X\in\mathbf{D}_j$ of the stored level-$j$ boundary terms of
Definition~\ref{9.2:def-boundary-families}. For every such unit the correlation theorem takes the
Cauchy radius $\rho_A:=\mathsf{w}/\eta_p^{1/2}$ in the rescaled amplitude $\mathsf{t}$. At the
frozen step $\mathsf{w}=\mathsf{w}_T$, and the Cauchy-estimate lemma gives
$\mathsf{w}_T^{-1}\le\theta_k$, where $\theta_k$ is the activity scale of the step. This is the
same inequality that enters the activity bound of the one-lattice Lipschitz step. There an activity
is $C_U\theta_k$ times the weight of its kind of term times its size, $C_U$ being the constant of
that step.

\emph{The response letters.} Fix $\varpi\in\,]0,1]$ and put
\begin{equation}\label{9.2:eq-response-factor-a}
\begin{aligned}
\rho(y)&:=\min\bigl\{1,\;C_b\,e^{-\delta_1^{\mathrm{sc}}(d_{\mathrm{sc}}(y,\Omega)-1)}\bigr\},\\
\mathsf{r}_{\varpi}(X)&:=\sup_{y\in X}\rho(y)^{\varpi}=\rho(y^{\circ})^{\varpi},\\
\delta_T&:=\min\{\varpi\delta_1^{\mathrm{sc}},\,\kappa/2\},
\qquad
C_w:=C_b^{\varpi}e^{\varpi\delta_1^{\mathrm{sc}}},\\
w(X)&:=e^{-(\kappa/2)d_j(X)}\,\mathsf{r}_{\varpi}(X),
\qquad
E'(X):=d_j(X)+d_{\mathrm{sc}}(y^{\circ},\Omega),\\
\Pi_j(\Delta)&:=\sum_{X\in\mathbf{D}_j:\,X\cap\Delta\neq\emptyset}
|F^{(j)}(X)|\,\mathsf{r}_{\varpi}(X),
\qquad
\Lambda_L:=\log(2L^4).
\end{aligned}
\end{equation}

The conventions behind these letters are as follows.
\begin{itemize}
\item $\Delta$ is a box.
\item $y^{\circ}\in X$ is a point at which $d_{\mathrm{sc}}(\cdot,\Omega)$ takes its least value
on $X$. It is fixed once for each unit. If the infimum is not attained, $y^{\circ}$ is chosen as in
the remark on non-attained minimising points.
\item The function $t\mapsto\min\{1,C_be^{-\delta_1^{\mathrm{sc}}(t-1)}\}$ is continuous and
non-increasing. Hence $\rho(y)$ is non-increasing in $d_{\mathrm{sc}}(y,\Omega)$. Since
$u\mapsto u^{\varpi}$ is increasing, this gives the equality in the definition of
$\mathsf{r}_{\varpi}$.
\item The same definition gives, for every set $S$,
\begin{equation*}
\mathsf{r}_{\varpi}(S)=\min\bigl\{1,C_be^{-\delta_1^{\mathrm{sc}}
(\inf_Sd_{\mathrm{sc}}(\cdot,\Omega)-1)}\bigr\}^{\varpi}.
\end{equation*}
In particular $\mathsf{r}_{\varpi}(S)\le\mathsf{r}_{\varpi}(X)$ when $S\subset X$.
\item Since $\varpi,\delta_1^{\mathrm{sc}},\kappa>0$, we have $\delta_T>0$.
\end{itemize}

Finally, $C_A^{(1/2)}$ denotes the one-unit piece-sum constant of the correlation theorem at Cauchy
radius $\mathsf{w}/\eta_p^{1/2}$. The rate $\kappa''>0$ is the decay rate of the receding class.

\begin{proposition}[The per-unit response factor at the step. Stated; proof incomplete at the step
marked $(\ast)$]\label{9.2:prop-response-factor}
Let $v$ be a marked unit of the summation of boundary terms at finer zones at the frozen step,
\emph{marked} in the sense of the correlation theorem's expansion. Assume that its localisation
domain is $X_v=X\in\mathbf{D}_j$ and that its size satisfies
$\mathfrak{N}(v)\le f_je^{-\kappa d_j(X)}$. Let $\hat{\Delta}=\square\in\pi_k$ be a block of the
block-sum lemma, and let $\tilde{\square}\supset\square$ be its fattened block.

After the Cauchy estimate at $\rho_A=\mathsf{w}_T/\eta_p^{1/2}$ and the allocation of the decay (b)
of the datum by the complex-type allocation lemma of the correlation theorem, the contribution of
$v$ to the activity that the block-sum lemma reads at $\hat{\Delta}$ is at most
\begin{equation}\label{9.2:eq-response-factor-b}
\begin{cases}
C_A^{(1/2)}\,\theta_k\,f_je^{-\kappa d_j(X)}\,\mathsf{r}_{\varpi}(X)
& \text{if } X\cap\tilde{\square}\neq\emptyset,\\[2pt]
C_A^{(1/2)}\,\theta_k\,f_je^{-\kappa d_j(X)}\,\mathsf{r}_{\varpi}(X)\,
e^{-\kappa''\operatorname{dist}_\infty(X,\square)/(M\ell_k)}
& \text{if } X\cap\tilde{\square}=\emptyset,
\end{cases}
\end{equation}
where $\operatorname{dist}_\infty$ is the distance in the maximum norm.
\begin{itemize}
\item[(i)] The first line is the case of the \emph{touching class}.
\item[(ii)] The second line is the case of the \emph{receding class}.
\item[(iii)] If the read set $S_v$ is not $X$ itself, both bounds hold with $S_v$ in place of $X$.
\end{itemize}
Consequently, the marked level-$j$ $\mathbf{B}$-activity at $\hat{\Delta}$ coming from the
touching class is at most
\begin{equation}\label{9.2:eq-response-factor-1}
C_U\,\theta_k\,C_A^{(1/2)}\,\Pi_j(\tilde{\square}).
\end{equation}
\end{proposition}

The proof of the correlation theorem displays the allocation of the decay (b) to the units of the
summation of boundary terms at finer zones only once for each level $j$, as an aggregate summed over
all these units. The factor $\mathsf{r}_{\varpi}(X)$ per unit is read from that aggregate; it is not
part of the statement of the correlation theorem.

\begin{proof}
Let $v$, $X$, $\square$, $\tilde{\square}$ be as in the statement. Let $p$ be a piece of $v$ read
at $\square$.

\emph{The Cauchy estimate.} We apply the Cauchy-estimate lemma of the correlation theorem to the
amplitude $\mathsf{t}$ on the disc of radius $\rho_A=\mathsf{w}_T/\eta_p^{1/2}$. The reciprocal of
this radius is $\eta_p^{1/2}\mathsf{w}_T^{-1}$. By $\mathsf{w}_T^{-1}\le\theta_k$ it is at most
$\theta_k\eta_p^{1/2}$. The estimate therefore charges the piece at most
\begin{equation*}
\theta_k\mathfrak{N}(v)\eta_p^{1/2}\le\theta_k\mathfrak{N}(v),
\end{equation*}
using $\eta_p\le1$. We call $\theta_k\mathfrak{N}(v)\eta_p^{1/2}$ the \emph{positional charge} of
the piece. It depends on $(\square,S_v)$ only through $\eta_p$.

$(\ast)$ \emph{The allocation.} The complex-type allocation lemma of the correlation theorem assigns
to the pieces of the units of the summation of boundary terms at finer zones a part of the decay (b)
of their datum. We take this allocation per unit, with the fraction $\varpi$ of the decay rate. That
is, the piece $p$ is taken to carry, at a read point $y\in S_v$, the factor
\begin{equation*}
\min\bigl\{1,C_be^{-\delta_1^{\mathrm{sc}}(\ell_{\mathrm{sc}}(y)-1)}\bigr\}^{\varpi},
\end{equation*}
where $\ell_{\mathrm{sc}}(y)$ is the scaled length over which the datum of $p$ decays at rate
$\delta_1^{\mathrm{sc}}$ when read at $y$. The minimum with $1$ expresses that no allocation can
exceed the positional charge. As explained after the statement, the proof of the correlation
theorem displays this allocation per level, not per unit.

\emph{The exponent.} The datum of $p$ is localised at $\square$ and built from the $\mathbf{R}$-born
boundary terms $\mathbf{B}'$, which live on $\Omega$. By the input used in the proof of (b), it
carries the factor $e^{-\delta_1^{\mathrm{sc}}d_{\mathrm{sc}}(\square,\Omega)}$. Hence the length
$\ell_{\mathrm{sc}}(y)$ at a read point $y\in S_v$ is the scaled length of the chain
$\Omega\to\square\to y$, and by the triangle inequality for $d_{\mathrm{sc}}$ it satisfies
\begin{equation*}
\ell_{\mathrm{sc}}(y)\;\ge\;d_{\mathrm{sc}}(y,\Omega)\;\ge\;\inf_{S_v}d_{\mathrm{sc}}(\cdot,\Omega).
\end{equation*}

\emph{The per-unit factor.} Since
$t\mapsto\min\{1,C_be^{-\delta_1^{\mathrm{sc}}(t-1)}\}^{\varpi}$ is non-increasing, the last two
steps bound the factor of $v$ by
\begin{equation*}
\min\bigl\{1,\,C_b\,e^{-\delta_1^{\mathrm{sc}}(\inf_{S_v}d_{\mathrm{sc}}(\cdot,\Omega)-1)}\bigr\}^{\varpi}
=\mathsf{r}_{\varpi}(S_v).
\end{equation*}
The equality is the identity recorded after the definition of the response letters. When
$S_v\subset X$, the same identity gives $\mathsf{r}_{\varpi}(S_v)\le\mathsf{r}_{\varpi}(X)$. Thus
$\mathsf{r}_{\varpi}(X)$ is a conservative charge, and any finer allocation only lowers the factor.

\emph{The three cases.} We sum over the pieces of $v$ at Cauchy radius $\mathsf{w}/\eta_p^{1/2}$
with the one-unit piece-sum constant $C_A^{(1/2)}$, and use $\mathfrak{N}(v)\le f_je^{-\kappa
d_j(X)}$.
\begin{itemize}
\item For the touching class this gives case (i).
\item For the receding class, $X\cap\tilde{\square}=\emptyset$, and the pieces read at $\square$
carry in addition the distance factor
$e^{-\kappa''\operatorname{dist}_\infty(X,\square)/(M\ell_k)}$ of that class. This gives case (ii).
\item If the read set $S_v$ is not $X$, the same argument with $S_v$ in place of $X$ gives case
(iii).
\end{itemize}

\emph{The touching class in sum.} Let $v'$ range over the marked units of level $j$ whose domains
$X_{v'}$ meet $\tilde{\square}$. The size of such a unit is
$\mathfrak{N}(v')=|F^{(j)}(X_{v'})|$, which is at most $f_je^{-\kappa d_j(X_{v'})}$ by
Definition~\ref{9.2:def-boundary-families}. For each such unit, the activity bound of the
one-lattice Lipschitz step contributes the factor $C_U\theta_k$, the weight
$C_A^{(1/2)}\mathsf{r}_{\varpi}(X_{v'})$ and the size $\mathfrak{N}(v')=|F^{(j)}(X_{v'})|$.
Summing over the domains $X=X_{v'}\in\mathbf{D}_j$ with $X\cap\tilde{\square}\neq\emptyset$, the
definition of $\Pi_j$ in \eqref{9.2:eq-response-factor-a} gives \eqref{9.2:eq-response-factor-1}.
\end{proof}

\begin{remark}[On the fraction $\varpi$]
The proposition is stated for every $\varpi\in\,]0,1]$. The proof of the correlation theorem
displays two values:
\begin{itemize}
\item $\varpi=1/2$, the half rate that accompanies the Cauchy radius $\mathsf{w}/\eta_p^{1/2}$;
\item $\varpi=1/4$, the quarter of the decay in the aggregate allocation over a level.
\end{itemize}
The fraction $\varpi$ enters $\delta_T$, $C_w$ and $\mathsf{r}_{\varpi}$ as displayed in
\eqref{9.2:eq-response-factor-a}.
\end{remark}

\begin{definition}[Attempt windows; thin and thick levels]\label{9.2:def-attempt-windows}
We use the following from earlier in this section:
\begin{itemize}
\item[--] the cube partitions $\mathcal{Q}_m$, regions, carriers $|\mathcal{S}|$, components and fattenings
$\mathcal{S}^{\sim n}$ of Notation~\ref{9.2:not-cubes-regions};
\item[--] the step-made sets, the current sets, the linear order of the operations that make the large-field
label, the stage $t_{\mathbf{T}}(m)$ of the renormalisation step $m-1\to m$ and the local scale $n(x)$ of
Definition~\ref{9.2:def-current-sets};
\item[--] the scaled distance $d_{\mathrm{sc}}$ of Definition~\ref{9.2:def-scaled-distance}.
\end{itemize}
As in Definition~\ref{9.2:def-current-sets}, a position in the linear order of operations is a \emph{stage};
for a stage $t$, $t^{-}$ is the stage immediately preceding it, and $Z_m^{(t)}$, $\Omega_m^{(t)}$ are the
large-field region and the determining domain at level $m$ as they stand at $t$. A cube of $\mathcal{Q}_m$
has \emph{scale} $m$.

The integer $k$ is fixed, and the label is considered up to the stage at which the renormalisation step
$k\to k+1$ takes it as input; we call this the \emph{present stage}. An \emph{attempt} is an application of
Ba{\l}aban's large-field operation $\mathbf{R}$ \cite{B-CMP122a,B-CMP122b}. An attempt is either completed
or non-completed; a non-completed attempt is one that was arrested or excluded, these being two of the
outcomes of $\mathbf{R}$ in \cite{B-CMP122a,B-CMP122b}. A further outcome, a second-class outcome, comes
with a region $W$, its \emph{frame}, and a step $k_W$; in \textup{(S4a)} it is treated like a non-completed
attempt with class $W$ and step $k_W$.

\emph{Standing hypotheses.} The history of the label is assumed to have the structure (S0)--(S6) below.
Where a hypothesis paraphrases an operation defined in Ba{\l}aban's papers, the place is given.

\smallskip
\noindent\textup{(S0)} \emph{Operations.} The current label arises from the $\mathbf{T}$-made sets by the
following three kinds of operations, and by no others.
\begin{itemize}
\item[(a)] \emph{$\mathbf{T}$-steps.} The step $m-1\to m$ makes $\Omega_m^{\mathbf{T}}$,
$\Lambda_m^{\mathbf{T}}$ and $Z_m^{\mathbf{T}}$, and nothing at lower levels \cite{B-CMP119}.
\item[(b)] \emph{Completed $\mathbf{R}$-operations.} A completed $\mathbf{R}$-operation at step $k_E$
resets a class $E$.
\begin{itemize}
\item[--] $E$ is a union of whole components of $Z_{k_E}$ as it stands when the operation is processed.
For each application of $\mathbf{R}$, the set denoted $Z$ in \cite{B-CMP122a}, defined there as the union
of a class of components of the large-field region, is the class on which that application acts: it is $E$
for a completed operation and $W_a$ for a non-completed attempt $a$.
\item[--] The whole of each such component is reset to one level \cite[(1.33)]{B-CMP122b}.
\item[--] $E$ is deleted from $Z_m^{\mathrm{cur}}$ for every $m\le k_E$ at which it was large-field. Hence a
component of a lower-level large-field region that lay in the part so taken away is no longer current.
\end{itemize}
For the frame $W$ of a second-class outcome at step $k_W$ one has instead $W\subset Z_{k_W}$: the frame is
not deleted from the large-field region. This case enters through \textup{(S4a)}.
\item[(c)] \emph{Preparations of non-completed attempts.} The preparation of a non-completed attempt $a$
(its class $W_a$, its step $k_a$ and its levels $h_a<j_a$ are the data fixed in \textup{(S1)})
re-makes the sets $\Omega''_l(a)$ and $Z''_l(a)$, for $h_a<l\le j_a$, inside its class $W_a$ only
\cite{B-CMP122a}. These sets are frozen thereafter: of the components on which the operation is not
completed, \cite{B-CMP122a} says ``we exclude them from the region $Z$. They are not changed by the
subsequent operations'', and this applies to $W_a$ and to the sets re-made inside it. On $W_a$,
\begin{equation*}
\Omega_l\cap W_a=W_a\setminus Z''_l(a)\qquad(h_a<l\le j_a).
\end{equation*}
The preparation leaves unchanged all domains $\Lambda_j$, the functions denoted $\varphi_j$ in
\cite{B-CMP122a}, the small-field terms $\mathbf{E}^{(j)}$, the large-field terms $\mathbf{R}^{(j)}$ and the
previously stored boundary terms $\mathbf{B}^{(j)}$ \cite{B-CMP122a}. Above $j_a$ nothing is re-made inside
$W_a$: the sets of the levels $l>j_a$ there are those current before $a$ is processed
\cite[(1.15)--(1.16)]{B-CMP122a}.
\end{itemize}

\emph{Attempts in force.} Let $a$ be a non-completed attempt, processed at a stage $t_a$, and let $t>t_a$.
We say that $a$ is \emph{in force at $t$} if $|W_a|$ is not contained in the union of the classes $E$ reset
by completed $\mathbf{R}$-operations processed after $t_a$ and before $t$. We say that $a$ is \emph{in force
now} if it is in force at the present stage. The content of this definition is as follows.
\begin{itemize}
\item[--] On the part of $|W_a|$ outside those classes, the re-made sets $\Omega''_l(a)$ and $Z''_l(a)$ are
the label's at $t$.
\item[--] On the part of $|W_a|$ inside those classes, they are deleted together with $E$.
\item[--] An attempt not in force at $t$ has no re-made set left in the label at $t$.
\end{itemize}
This is part of the definition, not a further condition.

\emph{Exhaustiveness.} Let $l\le k$ and let $t$ be any stage. For $t<t_{\mathbf{T}}(l)$ the determining
domain at level $l$ is not yet made (it is first made at $t_{\mathbf{T}}(l)$,
Definition~\ref{9.2:def-current-sets}), so the region $\Omega_l^{(t)}$ is empty and there is nothing to
assert; let therefore $t\ge t_{\mathbf{T}}(l)$. By (a), no
$\mathbf{T}$-step after the step $l-1\to l$ makes anything at level $l$; by (a)--(c), the determining domain
at level $l$ at stage $t$ is therefore $\Omega_l^{\mathbf{T}}$ modified only by operations processed before
$t$ of the two remaining kinds. A completed $\mathbf{R}$-operation acts at level $l$ only if $l\le k_E$, and
only inside its class $E$, by (b). A preparation of $a$ acts at level $l$ only if $h_a<l\le j_a$, and only
inside $W_a$, by (c); the region it leaves at level $l$ is $W_a\setminus Z''_l(a)$, whose carrier is
$U_l(a)$ below, and by the definition of an attempt in force these re-made sets are present in the label at
$t$ only if $a$ is in force at $t$, and then only outside the classes reset after $t_a$ and before $t$.
\textup{(S0)} asserts in addition that every cube $Q$ of the region $\Omega_l^{(t)}$ satisfies at least one
of the following; only this exhaustiveness is used.
\begin{itemize}
\item[(n)] $Q\subset|\Omega_l^{\mathbf{T}}|$.
\item[(r)] $Q\subset|E|$ for a reset class $E$ processed before $t$ with $k_E\ge l$.
\item[(p)] $Q\subset U_l(a):=|W_a\setminus Z''_l(a)|$ for a non-completed attempt $a$ processed before
$t$, in force at $t$, with $h_a<l\le j_a$.
\end{itemize}

\emph{The large-field regions.} Nothing else changes $Z_m$: for $t\ge t_{\mathbf{T}}(m)$,
\begin{equation}\label{9.2:eq-attempt-windows-a}
|Z_m^{(t)}|=|Z_m^{\mathbf{T}}|\setminus\bigcup\bigl\{\,|E|\;:\;E\text{ a reset class processed before }t,
\ k_E\ge m\,\bigr\}.
\end{equation}
The sets $W_a$ and $E$ are unions of whole components of $Z_{k_a}$ and of $Z_{k_E}$, respectively, as these
regions stand when $a$, respectively the operation resetting $E$, is processed.

\smallskip
\noindent\textup{(S1)} \emph{Data of a non-completed attempt.} A non-completed attempt $a$ has the
following data.
\begin{itemize}
\item[--] Its step $k_a$, at which it acts on the effective density $\rho_{k_a}$ at coupling $x_{k_a}$.
\item[--] Its clock $N_a$, and $h_a:=k_a-N_a$.
\item[--] Ba{\l}aban's integer $N_0(a)$ \cite{B-CMP122a}: ``Take the smallest positive integer $N_0$ such,
that $L^{-N_0+1}MR_{k-N_0+1}=M$ \dots\ We assume that $N>N_0$'', taken at $k=k_a$, $N=N_a$. Here $R_j$ are
the radii of \cite{B-CMP122a}, and a cube of side $L^{-(k-j)}MR_j$ there is a cube of $\mathcal{Q}_j$ here
(see \textup{(S2)}).
\item[--] $k_0(a):=k_a-N_0(a)$.
\item[--] The last re-made level $j_a\le k_a-1$, so that $n_a:=j_a-h_a<N_a$ integrations have been
performed.
\item[--] The \emph{window} $I_a:=[h_a,k_a-1]$, which consists of $N_a$ levels.
\item[--] The \emph{thin part} $\Theta_a:=[k_0(a)+1,k_a-1]$.
\end{itemize}
The levels of $\Theta_a$, the top $N_0(a)-1$ levels of the window, are the thin strata of $a$. These strata
are separated by one layer of $M$-cubes, one layer of $L^{-1}M$-cubes,
and so on \cite{B-CMP122a}, and each has $d_{\mathrm{sc}}$-thickness at least $1$. The levels $\le k_0(a)$
keep their own scales. Inside $W_a$ the collars (Lemma~\ref{9.2:lem-ten-layer-collars}) of the earlier
domains $\Lambda_i$, $k_0(a)<i<j_a$, lie
outside the union of all the re-made strata of $a$, in the part of $W_a$ where the current index is $j_a$.

\smallskip
\noindent\textup{(S2)} \emph{The re-made strata.} Inside $W_a$,
\begin{equation*}
Z''_l(a)=W_a\setminus\bigl(\Omega_l^{(a)}\bigr)^{\sim5}\qquad\bigl(h_a<l\le k_0(a)\bigr).
\end{equation*}
This is \cite[(1.11)]{B-CMP122a}, ``$Z''_j=(\Omega_j^{\sim5})^c\cap Z$ for $j=k-N_0,\dots,h+1$'', written
in the notation of \cite{B-CMP122a}, where $k$ is the step of the attempt, and taken at $k=k_a$, $h=h_a$ and
$Z=W_a$. In the same notation \cite[(1.10)]{B-CMP122a} states
\begin{equation*}
Z''_{k-N_0+1}=(Z'_{k-N_0})^{\sim}=(\Omega^{\sim7}_{k-N_0+1})^c\cap Z .
\end{equation*}
At the level $k_0(a)+1$ the hypothesis is taken in either of the following two readings of that display:
\begin{equation*}
Z''_{k_0+1}(a)=\bigl(Z'_{k_0}\bigr)^{\sim1}\quad\text{(reading A)},\qquad
Z''_{k_0+1}(a)=W_a\setminus\bigl(\Omega^{(a)}_{k_0+1}\bigr)^{\sim7}\quad\text{(reading B)}.
\end{equation*}
Here $k_0=k_0(a)$, and $Z'_{k_0}$ is the set so denoted in \cite[(1.10)]{B-CMP122a}; it is not the set
$Z'_A$ of Notation~\ref{9.2:not-cubes-regions}. Reading A is the
middle member of that display, with the unindexed fattening $(\,\cdot\,)^{\sim}$ taken as one layer of
cubes; reading B is its right member, with $Z=W_a$. The display asserts that the two members agree; both
readings are carried in what follows.

The sets $Z''_l(a)$ increase with $l$ in the thin range $l\ge k_0(a)+1$. Moreover the preparation changes
nothing at the levels $j\le h_a$: there $\Omega''_j(a)=\Omega^{(a)}_j$. There is one cube family: the sets
$\Omega_j^{\sim n}$ are unions of $L^{-(k-j)}MR_j$-cubes,
where $k$ is the step of the attempt and lengths are measured in units of the spacing $\ell_k$ of that step;
that is, they are unions of the cubes of side $M\check{R}_j\ell_j$ of $\mathcal{Q}_j$, the family of
Notation~\ref{9.2:not-cubes-regions} at scale $j$. The current sets entering these formulae are
$\Omega^{(a)}:=\Omega^{(t_a^{-})}$, those current just before $a$ is processed. The fattening radii $7$ in
\cite[(1.10)]{B-CMP122a} and $5$ in \cite[(1.11)]{B-CMP122a} are used exactly as printed.

\smallskip
\noindent\textup{(S3)} \emph{Attempts in force at a point.} For a point $y$, $\mathfrak{A}(y)$ is the set of
non-completed attempts $a$ that are in force now and satisfy $y\in W_a$.

\smallskip
\noindent\textup{(S4)} \emph{Disjoint windows.} \textup{(S4a)} Let $a\ne a'$ be two non-completed attempts
whose classes have meeting carriers, $|W_a|\cap|W_{a'}|\ne\emptyset$; arrested, excluded and second-class
attempts are treated alike. Then
\begin{equation*}
I_a\cap I_{a'}=\emptyset,\qquad\text{and}\qquad k_{a'}<k_a\ \Longrightarrow\ k_{a'}\le h_a .
\end{equation*}
This rests on the statement that an arrest creates a new large-field region, which gives $h_b\ge k_a$ for a
later attempt $b$
whose class has carrier meeting $|W_a|$. A second-class outcome is likewise a creation \cite{B-CMP122b}, so
that $k_W\le h_a$ in that case too. \textup{(S4a)} is a statement about the data of the two attempts, taken
at the stage at which the step of the later attempt takes the label as input. It therefore remains valid
after later resets. \textup{(S4b)} Two distinct non-completed attempts of one step have classes with
disjoint carriers, since both
windows contain that step less one; this is derived from \textup{(S4a)} below, as consequence \textup{(ii)}.

\smallskip
\noindent\textup{(S5)} \emph{The last window.} $k_a\le k$ for every attempt in force at the present stage,
and the bound is sharp: the operation that acts on $\rho_k$ before the step $k\to k+1$ takes the label as
input ($\mathbf{R}_k$ in the indexing of \cite{B-CMP122a,B-CMP122b}) has step $k$, so that if this
application is not completed it is an attempt in force with $k_a=k$. Its remainders are new boundary terms
near $\partial\Lambda_{k+1}$ and among the terms $\mathbf{B}'$, and in the comparison of the two runs
(Definition~\ref{9.6:def-comparison-chain}) they are
read by the term collecting the boundary terms near $\partial\Lambda_{k+1}\cup\mathbf{B}'$.
In the words of \cite{B-CMP122a}, these are ``some new boundary terms only, i.e., some new terms in
$\mathbf{B}_k$''. In that indexing the arrested pieces and frames entering at step $k$ are those of steps
$\le k-1$, one index down, which agrees with the present indexing.

\smallskip
\noindent\textup{(S6)} \emph{The clock.} $N_a>N_0(a)$ \cite{B-CMP122a}, and
\begin{equation*}
2N_0(a)\le N_a,\qquad N_0(a)\le N_0^{\sharp}(x_{k_a}),\qquad\text{where}\quad
N_0^{\sharp}(x):=2+r_{\mathrm{pr}}\log_L\log\bigl((1+B^{\sharp})x\bigr).
\end{equation*}
The condition $2N_0^{\sharp}(x)\le(\log(x-c_5))^{r_{\mathrm{pr}}}$, that is
\begin{equation*}
4+2r_{\mathrm{pr}}\log_L\log\bigl((1+B^{\sharp})x\bigr)\le\bigl(\log(x-c_5)\bigr)^{r_{\mathrm{pr}}},
\end{equation*}
is the form of the condition $2N_0\le N$ on live large-field attempts in which only the coupling enters.

\smallskip
\noindent\emph{Thin and thick levels.} Let $j$ be a level with $j+1\le k$. For $y\in Z_i$ with
$j+1\le i\le k$, let $A_i(y)$ be the component of $Z_i$ containing $y$. Level $i$ is \emph{thin} for $y$ if
some $x\in A_i(y)$ has $n(x)>i$, and \emph{thick} otherwise. The \emph{thick indicator} is
\begin{equation}\label{9.2:eq-attempt-windows-b}
c_i(y):=\begin{cases}0,&\text{if }n(x)>i\text{ for some }x\in A_i(y)\quad(\text{level }i\text{ thin for }y),\\
1,&\text{if }n(x)\le i\text{ for all }x\in A_i(y)\quad(\text{level }i\text{ thick for }y).\end{cases}
\end{equation}
If level $i$ is thick for $y$, then every current cube meeting $A_i(y)$ has scale $\le i$.

\smallskip
\noindent\emph{The two-sided comparison of couplings.} For the running inverse couplings $x_i=g_i^{-2}$ of
Definition~\ref{1.2:def-couplings} we use the two-sided comparison of the running couplings, with constants
$\lambda>0$ and $c_2\ge0$, namely
\begin{equation}\label{9.2:eq-attempt-windows-c}
\lambda(j'-i')-2c_2\;\le\;x_{i'}-x_{j'}\;\le\;3\lambda(j'-i')+2c_2\qquad\text{for }0\le i'<j'\le K.
\end{equation}
The constants $c_5$ and $B^{\sharp}$ enter here as $c_5=2c_2\ge0$ and $B^{\sharp}=3\lambda+2c_2>0$. The
degree exponent satisfies $r_{\mathrm{pr}}\ge2$.

\smallskip
\noindent\emph{Consequences.} The following are proved below.
\begin{itemize}
\item[(i)] Every window satisfies $\max I_a\le k-1$; in particular level $k$ lies in no window.
\item[(ii)] \textup{(S4b)}: two distinct non-completed attempts of one step have classes with disjoint
carriers.
\item[(iii)] If level $i$ is thick for $y$, then every current cube meeting $A_i(y)$ has scale $\le i$.
\end{itemize}
\end{definition}

\begin{proof}[Proof of the consequences]
\emph{Consequence \textup{(i)}.} Let $a$ be an attempt in force at the present stage. By \textup{(S5)},
$k_a\le k$, hence $k_a-1\le k-1$. By \textup{(S1)}, $I_a=[h_a,k_a-1]$, so $\max I_a=k_a-1\le k-1$, and
$k\notin I_a$.

\emph{Consequence \textup{(ii)}.} Let $a\ne a'$ be non-completed attempts with $k_a=k_{a'}$. By
\textup{(S1)}, $N_0(a)$ is a positive integer. By \textup{(S6)}, $N_a>N_0(a)\ge1$, so $N_a\ge2$ and
$h_a=k_a-N_a\le k_a-1$. Hence $k_a-1\in[h_a,k_a-1]=I_a$, and in the same way $k_{a'}-1=k_a-1\in I_{a'}$.
Thus $I_a\cap I_{a'}\ne\emptyset$. If $|W_a|$ and $|W_{a'}|$ met, \textup{(S4a)} would give
$I_a\cap I_{a'}=\emptyset$, a contradiction. So $|W_a|\cap|W_{a'}|=\emptyset$.

\emph{Consequence \textup{(iii)}.} Let level $i$ be thick for $y$. By \eqref{9.2:eq-attempt-windows-b},
$n(x)\le i$
for every $x\in A_i(y)$. Let $Q$ be a current cube of scale $m$ meeting $A_i(y)$. Cubes are closed, so the
meeting is a point $x\in Q\cap A_i(y)$. A point of a current cube of scale $m$ has local scale $n(x)\ge m$
(Definition~\ref{9.2:def-current-sets}). Hence
\begin{equation*}
m\le n(x)\le i .
\end{equation*}
\end{proof}

\begin{theorem}[The pile bound for stored boundary terms]\label{9.2:thm-pile-bound}
Work on one lattice and one scaffold, at the frozen step $k\to k+1$ of
Notation~\ref{9.2:not-frozen-step}, on any common large-field label in the sense of
Notation~\ref{9.2:not-cubes-regions}. Untouched labels, labels reset elsewhere, labels with
second-class frames and labels carrying live arrested or excluded attempts are all allowed.
Let $j\le k$, put $s:=k-j$, and let $F$ be a level-$j$ boundary family in the sense of
Definition~\ref{9.2:def-boundary-families}, with units $X\in\mathbf{D}_j$ and size $f_j$.
Let $\varpi\in\,]0,1]$; the values $\varpi=1/2$ and $\varpi=1/4$ are the ones displayed below.
Take the profile $\rho$, the per-unit response factor $\mathsf{r}_{\varpi}$, the joint exponent
$E'$, the pile per box $\Pi_j$ and the letters $\delta_T$, $\Lambda_L$, $C_w$, $w(X)$ as fixed in
\eqref{9.2:eq-response-factor-a}. Only the definitions collected in that display are used here;
Proposition~\ref{9.2:prop-response-factor}, stated with proof incomplete at one step, is not
invoked. We recall
\begin{gather*}
\delta_T=\min\{\varpi\delta_1^{\mathrm{sc}},\kappa/2\},\qquad
\Lambda_L=\log(2L^4),\\
C_w=C_b^{\varpi}e^{\varpi\delta_1^{\mathrm{sc}}},\qquad
w(X)=e^{-(\kappa/2)d_j(X)}\,\mathsf{r}_{\varpi}(X),
\end{gather*}
and put $\check{R}^{\flat}:=\check{R}_k$.
The following floors are assumed throughout: $\check{R}_k\ge2$, $L\ge3$ and
$\kappa\ge10\log L$. On labels with non-completed attempts in force, condition (S6) of
Definition~\ref{9.2:def-attempt-windows} is assumed as well.

Assume the following.
\begin{itemize}
\item The ten-layer collars \eqref{9.2:eq-ten-layer-collars-a} of the current sets and the
eight-layer step collar \eqref{9.2:eq-ten-layer-collars-e} hold
(Lemma~\ref{9.2:lem-ten-layer-collars}).
\item Every unit of $F$ carries the anchor \eqref{9.2:eq-anchors-a} of
Lemma~\ref{9.2:lem-anchors}. A unit that carries only the weak anchor
\eqref{9.2:eq-anchors-b} is treated in the final paragraph below, which replaces (i) for it.
\item The nesting $\Omega^{\natural}_{i+1}\subset\Lambda_i\subset\Omega^{\natural}_i$ holds,
with $\Lambda_i$ the label's action domains (Definition~\ref{9.2:def-current-sets}), each a
union of cubes of $\mathcal{Q}_i$.
\item The following hold: the monotonicity of the radii and the radius floor of
Notation~\ref{9.2:not-radii-floor}; the bound $g_k\le\gamma$ on the scaffold; the constant
$K_0^{(2)}$ of \eqref{9.2:eq-boundary-families-b}, together with the counting bound for
units on which it rests.
\item The scaled distance $d_{\mathrm{sc}}$ is that of
Definition~\ref{9.2:def-scaled-distance}.
\item On labels with live arrested or excluded attempts, the attempt windows have the
structure (S0)--(S6) of Definition~\ref{9.2:def-attempt-windows}. The count bound for live
attempts, with its function $N_0^{\sharp}$, holds, and so does the reference bound on the
running couplings with constants $\lambda$, $c_2$ and $c_5=2c_2$.
\end{itemize}
The conditions used below are the following, all one-sided:
\begin{equation}\label{9.2:eq-pile-bound-a}
\begin{aligned}
&(\mathrm{a}_n)\colon\ n\,\delta_T\check{R}^{\flat}\ge\Lambda_L,\qquad n\in\{10,5,4\};\\
&(\mathrm{b})\colon\ 5\,\delta_T\check{R}_k\ge N_0^{\sharp}(2x_k)\,\Lambda_L;\\
&(\mathrm{c})\colon\ 8B^{\sharp}N_0^{\sharp}(y)\le y\quad\text{for every }
y\ge\gamma^{-2}-c_5;\\
&(\mathrm{d})\colon\ (1+B^{\sharp})\,y\ge\max\{e,2^{r_{\mathrm{pr}}}\}\quad\text{for every }
y\ge\gamma^{-2}-c_5.
\end{aligned}
\end{equation}
Condition $(\mathrm{a}_4)$ implies $(\mathrm{a}_5)$, and $(\mathrm{a}_5)$ implies
$(\mathrm{a}_{10})$. At $\varpi=1/2$, condition $(\mathrm{a}_4)$ is the pair of floors
\begin{equation*}
e^{2\delta_1^{\mathrm{sc}}\check{R}_k}\ge2L^4,\qquad e^{2\kappa\check{R}_k}\ge2L^4 .
\end{equation*}
At both $\varpi=1/2$ and $\varpi=1/4$ this pair implies $(\mathrm{a}_{10})$.

Let $\sigma\in\{1/M,1,\hat{\sigma}\}$, let $\Delta$ be a box of side $\sigma M\ell_k$ anywhere
on the torus, let $c$ and $C_{\mathrm{pile}}(\sigma)$ be as in (ii) below, and let $U(X)$ be the
set of thin levels of (i) below. Then, for every unit $X$ of $F$,
\begin{equation}\label{9.2:eq-pile-bound-b}
\begin{aligned}
w(X)&\le C_w\,e^{-(13/2)\delta_T\check{R}_k}\,
e^{-10\delta_T\check{R}^{\flat}(s-\#U(X))},\\
\Pi_j(\Delta)&\le C_{\mathrm{pile}}(\sigma)\,e^{-c\,\delta_T\check{R}_k}\,2^{-s}f_j,\\
\sum_{j\le k}\Pi_j(\Delta)&\le 2\,C_{\mathrm{pile}}(\sigma)\,e^{-c\,\delta_T\check{R}_k}
\max_{j\le k}f_j,\\
w(X)&\le C_w\,e^{-(11/2)\delta_T\check{R}_k}\,(2L^4)^{-s}
\quad\text{(weak anchor, grades ($\alpha$), ($\beta$))}.
\end{aligned}
\end{equation}
The first line holds without any of the conditions \eqref{9.2:eq-pile-bound-a}. The second and
third lines hold under those conditions as specified in (ii). The fourth line holds for a
weakly anchored unit of grade ($\alpha$) or ($\beta$), under the hypotheses and conditions
stated in the final paragraph. More precisely:
\begin{enumerate}
\item[(i)] \emph{Per unit.} Let $y^{\circ}$ be the point attached to $X$ in
\eqref{9.2:eq-response-factor-a}, and let $\iota(y^{\circ})$ be its depth. Let
$U(X)\subset[j+1,k]$ be the set of levels thin for $y^{\circ}$ (in the sense of
Definition~\ref{9.2:def-attempt-windows}) at or above $\max\{\iota(y^{\circ}),j+1\}$. Put
$P(X):=[j+1,k]\setminus U(X)$. Then $w(X)\le C_w e^{-\delta_T E'(X)}$, and
\begin{equation}\label{9.2:eq-pile-bound-1}
E'(X)\ \ge\ 10\sum_{i\in P(X)}\check{R}_i+8L\check{R}_{k+1}-3
\ \ge\ 10\check{R}^{\flat}\bigl(s-\#U(X)\bigr)+\tfrac{13}{2}\check{R}_k .
\end{equation}
Hence the first line of \eqref{9.2:eq-pile-bound-b} holds. Moreover:
\begin{enumerate}
\item[($\alpha$)] Suppose $U(X)=\emptyset$. This holds in particular whenever no non-completed
attempt in force acts on a class containing a point of $X$. Then, under $(\mathrm{a}_{10})$,
\begin{equation}\label{9.2:eq-pile-bound-2}
w(X)\le C_w\,e^{-(13/2)\delta_T\check{R}_k}\,(2L^4)^{-s}.
\end{equation}
\item[($\beta$)] Suppose $y^{\circ}$ has no exceptional attempt window (in the sense fixed in
the proof of (i) below). A sufficient condition
is that no attempt $a$ in force on a class containing a point of $X$ has both $h_a\le j$ and
$j\le k_a-2$. Then, under $(\mathrm{a}_5)$, the bound \eqref{9.2:eq-pile-bound-2} holds.
\item[($\gamma$)] For every unit on every label, under $(\mathrm{a}_4)$, (b), (c) and (d),
\begin{equation*}
w(X)\le C_w\,e^{-(3/2)\delta_T\check{R}_k}\,(2L^4)^{-s}.
\end{equation*}
\end{enumerate}
At $s=0$ no condition of \eqref{9.2:eq-pile-bound-a} is used, and
$w(X)\le C_w e^{-(13/2)\delta_T\check{R}_k}$. No upper bound on the age $u:=k-k_{a^{*}}$ of an
exceptional window $a^{*}$ is used: this age is split, never bounded.
\item[(ii)] \emph{Pile per box.} The three values of $\sigma$ correspond to a level-$k$ cell
($\sigma=1/M$), a $\pi_k$-cube ($\sigma=1$) and the fattened block of the block-sum lemma
($\sigma=\hat{\sigma}$). Put
\begin{equation*}
C_{\mathrm{pile}}(\sigma):=(\sigma+2)^4K_0^{(2)}C_w .
\end{equation*}
For $\sigma\le1$ this gives $C_{\mathrm{pile}}(\sigma)\le81K_0^{(2)}C_w$. The exponent is
$c:=13/2$ if every unit meeting $\Delta$ is of grade ($\alpha$) or ($\beta$), and $c:=3/2$ in
general. For weakly anchored units the two values are $11/2$ and $3/2$ respectively. With this
choice, the second line of \eqref{9.2:eq-pile-bound-b} holds under the corresponding conditions
of \eqref{9.2:eq-pile-bound-a}.
\item[(iii)] \emph{Level sum.} Given, for every $j\le k$, a level-$j$ boundary family with size
$f_j$ as above, under the same conditions the third line of
\eqref{9.2:eq-pile-bound-b} holds, and its right-hand side tends to zero as
$\gamma\downarrow0$. This is the form of decay, per level $j$ and summed over the shell index
$s=k-j$, that appears in the correlation theorem.
\end{enumerate}

\emph{Weakly anchored units.} Suppose a unit carries the weak anchor \eqref{9.2:eq-anchors-b}.
Suppose further that the anchor excess $a$ of \eqref{9.2:eq-anchors-b}
satisfies $a\le L\check{R}_{j+1}$ when $s\ge1$, and $a\le\check{R}_k$ when $s=0$. Then the
grades ($\alpha$), ($\beta$), ($\gamma$) hold under the same conditions, with the exponent
$13/2,\ 13/2,\ 3/2$ of $\delta_T\check{R}_k$ replaced by $11/2,\ 11/2,\ 3/2$ respectively; for
grades ($\alpha$) and ($\beta$) this is the fourth line of \eqref{9.2:eq-pile-bound-b}.
More precisely, the exponents are $13/2,\ 8,\ 3$ when $s\ge1$, and $11/2$ when $s=0$. They are
$13/2,\ 13/2,\ 3/2$ throughout if the floor $\check{R}_{k+1}\ge2$ is also assumed.

Every constant above depends only on $C_b$, $\delta_1^{\mathrm{sc}}$, $\varpi$, $\sigma$, and on
$\kappa$ through $K_0^{(2)}$. No constant depends on $K$, $k$, $j$, the ages, the attempt data
$N_a$ and $N_0(a)$, the radii $\check{R}_i$, the torus or the position. No upper bound on $L$,
$M$, $K$, $k$, $j$, the ages, $N_a$, $N_0(a)$ or a nesting depth is used or implied.
\end{theorem}

The proof follows.

\begin{lemma}[Additivity of sup-norm distances across nested sets]\label{9.2:lem-nested-distances}
Distances are taken in the $\ell^\infty$ metric of the torus of
Notation~\ref{9.2:not-cubes-regions}. For a point $y$ and sets $A,C$ we write
$\operatorname{dist}_\infty(y,A)$ for the distance from $y$ to $A$, and
$\operatorname{dist}_\infty(A,C)$ for the distance between $A$ and $C$.
\begin{itemize}
\item[(a)] Let $A_0\subset A_1\subset\cdots\subset A_m$ be closed sets of the torus, let
$z\in A_0$, and let $y\notin A_m$.
\item[(b)] Let $A\subset B$ be sets of the torus and let $y$ be a point.
\end{itemize}
Then, in case (a) and in case (b) respectively,
\begin{equation}\label{9.2:eq-nested-distances-a}
\begin{aligned}
&\text{(a)}\quad |z-y|_\infty \ge \sum_{i=1}^{m}\operatorname{dist}_\infty(A_{i-1},A_i^c)
+\operatorname{dist}_\infty(y,A_m),\\
&\text{(b)}\quad \operatorname{dist}_\infty(y,B^c)\ge \operatorname{dist}_\infty(A,B^c)
-\operatorname{dist}_\infty(y,A).
\end{aligned}
\end{equation}
\end{lemma}

\begin{proof}
(a) By Notation~\ref{9.2:not-cubes-regions} the $\ell^\infty$ metric of the torus is a length
metric. Thus there is a path $\omega$ in the torus, parametrised by $[0,1]$, with $\omega(0)=z$,
$\omega(1)=y$ and $\operatorname{len}(\omega)=|z-y|_\infty$. Here $\operatorname{len}$ is the
length defined there, which satisfies
$\operatorname{len}(\omega|_{[\alpha',\beta']})\ge|\omega(\beta')-\omega(\alpha')|_\infty$, and
the lengths of consecutive pieces of a path add.

For $i=0,\dots,m$ put $S_i:=\omega^{-1}(A_i)$. This set is closed, since $A_i$ is closed and
$\omega$ is continuous. It contains $0$, since $z\in A_0\subset A_i$. Being a non-empty closed
subset of $[0,1]$, it has a largest element; put
\begin{equation*}
\tau_i:=\max S_i,\qquad u_i:=\omega(\tau_i)\in A_i .
\end{equation*}
Since $y\notin A_m$ and $A_m\supseteq A_i$, we have $y\notin A_i$, hence $1\notin S_i$ and
$\tau_i<1$. By maximality, $\omega(t)\notin A_i$ for every $t\in(\tau_i,1]$. Letting
$t\downarrow\tau_i$ and using continuity of $\omega$ gives
$u_i\in\operatorname{cl}(A_i^c)$. The nesting $A_{i-1}\subset A_i$ gives
$S_{i-1}\subset S_i$, hence
\begin{equation*}
\tau_0\le\tau_1\le\cdots\le\tau_m<1 .
\end{equation*}

Fix $i\in\{1,\dots,m\}$. Since $u_{i-1}\in A_{i-1}$ and $u_i\in\operatorname{cl}(A_i^c)$,
\begin{equation*}
|u_{i-1}-u_i|_\infty\ge\operatorname{dist}_\infty\bigl(A_{i-1},\operatorname{cl}(A_i^c)\bigr)
=\operatorname{dist}_\infty(A_{i-1},A_i^c).
\end{equation*}
The equality holds because the distance from a point to a set equals the distance from that
point to the closure of the set. Since $u_m\in A_m$,
\begin{equation*}
|u_m-y|_\infty\ge\operatorname{dist}_\infty(y,A_m).
\end{equation*}
The intervals $[\tau_{i-1},\tau_i]$, $i=1,\dots,m$, and $[\tau_m,1]$ are consecutive pieces of
$[\tau_0,1]\subset[0,1]$. Lengths add over consecutive pieces, and the remaining piece
$[0,\tau_0]$ has non-negative length. Applying the lower bound of a piece's length by the distance
of its endpoints to each piece, we obtain
\begin{align*}
|z-y|_\infty=\operatorname{len}(\omega)
&\ge\sum_{i=1}^{m}\operatorname{len}\bigl(\omega|_{[\tau_{i-1},\tau_i]}\bigr)
+\operatorname{len}\bigl(\omega|_{[\tau_m,1]}\bigr)\\
&\ge\sum_{i=1}^{m}|u_{i-1}-u_i|_\infty+|u_m-y|_\infty .
\end{align*}
Inserting the two lower bounds of the previous paragraph gives~(a)
of~\eqref{9.2:eq-nested-distances-a}.

(b) Let $a'\in A$ and $b\in B^c$. By the triangle inequality,
\begin{equation*}
\operatorname{dist}_\infty(A,B^c)\le|a'-b|_\infty\le|a'-y|_\infty+|y-b|_\infty .
\end{equation*}
Taking the infimum over $b\in B^c$, and then the infimum over $a'\in A$, gives
\begin{align*}
\operatorname{dist}_\infty(A,B^c)
&\le|a'-y|_\infty+\operatorname{dist}_\infty(y,B^c),\\
\operatorname{dist}_\infty(A,B^c)
&\le\operatorname{dist}_\infty(y,A)+\operatorname{dist}_\infty(y,B^c).
\end{align*}
This is~(b) of~\eqref{9.2:eq-nested-distances-a}.
\end{proof}

\begin{corollary}[Distances across the nested current sets]
Let $Z_i$ be the current large-field sets of Definition~\ref{9.2:def-current-sets}. Let
$j\le b\le k$, $z\in Z_j$ and $y\notin Z_b$. Then
\begin{equation}\label{9.2:eq-nested-distances-b}
|z-y|_\infty\ge\sum_{i=j+1}^{b}10M\check{R}_i\ell_i+\operatorname{dist}_\infty(y,Z_b).
\end{equation}
\end{corollary}

\begin{proof}
Apply part~(a) of Lemma~\ref{9.2:lem-nested-distances} with $m:=b-j$ and
$A_{i'}:=Z_{j+i'}$ for $i'=0,\dots,m$. These sets are closed and nested by
Definition~\ref{9.2:def-current-sets} and~\eqref{9.2:eq-current-sets-b}. Moreover
$z\in A_0=Z_j$ and $y\notin A_m=Z_b$. This gives
\begin{equation*}
|z-y|_\infty\ge\sum_{i=j+1}^{b}\operatorname{dist}_\infty(Z_{i-1},Z_i^c)
+\operatorname{dist}_\infty(y,Z_b).
\end{equation*}
When $b=j$ the sum is empty and the claim is already proved.

Now let $j+1\le i\le b\le k$. Every point of $Z_{i-1}$ lies in some current component $D$ of
$Z_{i-1}$. Hence $\operatorname{dist}_\infty(Z_{i-1},Z_i^c)$ is the infimum, over the current
components $D$ of $Z_{i-1}$, of $\operatorname{dist}_\infty(D,Z_i^c)$. By the collar bound
\eqref{9.2:eq-ten-layer-collars-a} of Lemma~\ref{9.2:lem-ten-layer-collars},
each of these distances is at least
$10M\check{R}_i\ell_i$. Hence
\begin{equation*}
\operatorname{dist}_\infty(Z_{i-1},Z_i^c)\ge10M\check{R}_i\ell_i .
\end{equation*}
Inserting this for each $i$ into the display above gives~\eqref{9.2:eq-nested-distances-b}.
\end{proof}

\begin{lemma}[The size side: tree length across collars]\label{9.2:lem-size-side}
Assume the following statements: the distance bound \eqref{9.2:eq-nested-distances-b};
the size bound \eqref{9.2:eq-boundary-families-a} for level-$j$
boundary families of Definition~\ref{9.2:def-boundary-families}; the anchor statement,
Lemma~\ref{9.2:lem-anchors}; and the collar bound \eqref{9.2:eq-ten-layer-collars-e} of
Lemma~\ref{9.2:lem-ten-layer-collars}.

Let $X$ be a unit in the sense of Definition~\ref{9.2:def-boundary-families}, with an anchor
$(y_0,z_0)$ as in Lemma~\ref{9.2:lem-anchors} ($y_0\in X$, $z_0\in Z_j$) satisfying
$|y_0-z_0|_\infty\le(1+a)M\ell_j$ for some real number $a$, where $a=0$ under
\eqref{9.2:eq-anchors-a}. Let $y\in X$ be a
point whose depth $\iota:=\iota(y)$ lies in $[j+1,k+1]$, and put
$t:=\operatorname{dist}_\infty(y,Z_{\iota-1})$. Then
\begin{equation}\label{9.2:eq-size-side-a}
d_j(X)\ \ge\ \sum_{i=j+1}^{\iota-1}10\check{R}_iL^{i-j}+\frac{t}{M\ell_j}-3-a .
\end{equation}
If moreover $\iota=k+1$ and $y\in\Omega$, the fluctuation domain of
Definition~\ref{9.2:def-scaled-distance}, then $t\ge 8LM\check{R}_{k+1}\ell_k$, and therefore
\begin{equation}\label{9.2:eq-size-side-1}
d_j(X)\ \ge\ \sum_{i=j+1}^{k}10\check{R}_iL^{i-j}+8L\check{R}_{k+1}L^{k-j}-3-a .
\end{equation}
\end{lemma}

\begin{proof}
By the definition of the depth $\iota(y)$ of a point, $y\notin Z_{\iota-1}$; and
$\iota-1\in[j,k]$ because $\iota\in[j+1,k+1]$. Since $z_0\in Z_j$ and $y\notin Z_{\iota-1}$,
the bound \eqref{9.2:eq-nested-distances-b} applies with $z=z_0$ and $b=\iota-1$, and gives
\begin{equation}\label{9.2:eq-size-side-2}
|z_0-y|_\infty\ \ge\ \sum_{i=j+1}^{\iota-1}10M\check{R}_i\ell_i+t .
\end{equation}
Both $y_0$ and $y$ lie in $X$, so the triangle inequality and the anchor bound give
\begin{equation*}
\operatorname{diam}_\infty X\ \ge\ |y_0-y|_\infty\ \ge\ |z_0-y|_\infty-|y_0-z_0|_\infty
\ \ge\ |z_0-y|_\infty-(1+a)M\ell_j .
\end{equation*}
We divide by $M\ell_j$. Since the spacings of levels $i$ and $j$ satisfy $\ell_i/\ell_j=L^{i-j}$,
\eqref{9.2:eq-size-side-2} yields
\begin{equation*}
\frac{\operatorname{diam}_\infty X}{M\ell_j}\ \ge\
\sum_{i=j+1}^{\iota-1}10\check{R}_iL^{i-j}+\frac{t}{M\ell_j}-1-a .
\end{equation*}
The size bound \eqref{9.2:eq-boundary-families-a} of Definition~\ref{9.2:def-boundary-families},
applied to the unit $X$, states
\begin{equation*}
d_j(X)\ \ge\ \frac{\operatorname{diam}_\infty X}{M\ell_j}-2 .
\end{equation*}
Combining the last two inequalities gives \eqref{9.2:eq-size-side-a}.

For the second assertion let $\iota=k+1$ and $y\in\Omega$. Then $\iota-1=k$, so
$t=\operatorname{dist}_\infty(y,Z_k)$. Because $y\in\Omega$,
\begin{equation*}
t=\operatorname{dist}_\infty(y,Z_k)\ \ge\ \operatorname{dist}_\infty(\Omega,Z_k)
\ \ge\ 8LM\check{R}_{k+1}\ell_k ,
\end{equation*}
the last inequality being the collar bound \eqref{9.2:eq-ten-layer-collars-e} of
Lemma~\ref{9.2:lem-ten-layer-collars}. Since
\begin{equation*}
\frac{8LM\check{R}_{k+1}\ell_k}{M\ell_j}=8L\check{R}_{k+1}L^{k-j},
\end{equation*}
inserting this lower bound for $t$ into \eqref{9.2:eq-size-side-a}, whose sum now runs up to
$\iota-1=k$, gives \eqref{9.2:eq-size-side-1}.

Neither the hypotheses nor the argument above use the local scale $n(x)$, a determining set,
or a window $I_a$ of a non-completed attempt; the bounds \eqref{9.2:eq-size-side-a} and
\eqref{9.2:eq-size-side-1} therefore hold in the same form whichever window is considered.
\end{proof}

\begin{lemma}[Frozen inclusions of re-made strata]\label{9.2:lem-frozen-inclusions}
Assume the following statements:
\begin{itemize}
\item[(i)] the structure of attempt windows, clauses (S0)--(S6) of \eqref{9.2:eq-attempt-windows-a}
in Definition~\ref{9.2:def-attempt-windows};
\item[(ii)] the collar statement of Lemma~\ref{9.2:lem-ten-layer-collars} in its cube form,
\eqref{9.2:eq-ten-layer-collars-b}, \eqref{9.2:eq-ten-layer-collars-c} and
\eqref{9.2:eq-ten-layer-collars-d}, each read with the large-field sets current at the stage of
the $\mathbf{T}$-step it concerns;
\item[(iii)] the statement \eqref{9.2:eq-current-sets-a} and the nesting
\eqref{9.2:eq-current-sets-b} of the current sets of Definition~\ref{9.2:def-current-sets};
\item[(iv)] the compatibility of the cube partitions,
\eqref{9.2:eq-cubes-regions-a}--\eqref{9.2:eq-cubes-regions-d} of Notation~\ref{9.2:not-cubes-regions};
\item[(v)] the layer count of the collar in \textup{(ii)}, as stated in
Lemma~\ref{9.2:lem-ten-layer-collars}: the collar made by a $\mathbf{T}$-step consists of eight layers
of cubes of the partition of the level that step makes, the same cube family in which the
fattenings below are counted.
\end{itemize}
Write $\mathrm{now}$ for the present stage of the history of the label, the stage at which the
sets $Z_i^{\mathrm{cur}}$ are read. Let $b$ be a non-completed attempt in force at the present
stage, with class $W_b$, step $k_b$ and levels $h_b=k_b-N_b$, $k_0(b)=k_b-N_0(b)$ and $j_b$ as in
Definition~\ref{9.2:def-attempt-windows}. For $h_b<m\le j_b$ put
\begin{equation*}
U_m(b):=|W_b\setminus Z''_m(b)|,
\end{equation*}
the re-made stratum of level $m$ of $b$ inside $W_b$ that is in force; by
\eqref{9.2:eq-cubes-regions-b} it equals $\operatorname{cl}(|W_b|\setminus|Z''_m(b)|)$. Then
\begin{equation}\label{9.2:eq-frozen-inclusions-b}
\begin{aligned}
&\text{(a)}\quad U_m(b)\cap Z_{m-1}^{\mathrm{cur}}=\emptyset
\qquad\text{for } h_b<m\le\min\{k_0(b),j_b\};\\
&\text{(b)}\quad U_{k_0(b)+1}(b)\cap Z_{k_0(b)}^{\mathrm{cur}}=\emptyset
\qquad\text{if } k_0(b)+1\le j_b;\\
&\text{(c)}\quad U_m(b)\subset U_{k_0(b)+1}(b),\ \text{hence}\
U_m(b)\cap Z_{k_0(b)}^{\mathrm{cur}}=\emptyset
\qquad\text{for } k_0(b)+1\le m\le j_b,
\end{aligned}
\end{equation}
and \textup{(b)} holds under either of the two readings, A and B, of $Z''_{k_0(b)+1}(b)$ in clause
(S2) of \eqref{9.2:eq-attempt-windows-a}. Here $N_b>N_0(b)$ by clause (S6), so
$k_0(b)=k_b-N_0(b)$ exceeds $k_b-N_b=h_b$.
\end{lemma}

\begin{proof}
\emph{Preliminary facts.} We first establish four facts, for stages $t\le t'$, levels $i$, $m\le m'$
and $l>i$, for resets $E$ (classes reset by completed $\mathbf{R}$-operations, at step $k_E$ and
stage $t_E$) performed so far, and for attempts $a$, completed or not:
\begin{equation}\label{9.2:eq-frozen-inclusions-a}
\begin{aligned}
&\text{(F1)}\quad Z_m^{(t')}\subset Z_m^{(t)},\qquad
Z_m^{(t)}\subset Z_{m'}^{(t)}\ \ (t\ge t_{\mathbf{T}}(m')),\qquad
Z_i^{\mathrm{cur}}\subset Z_i^{\mathbf{T}};\\
&\text{(F2)}\quad \text{if } k_E\ge m:\ E\cap Z_m^{(t_E^-)}\ \text{is a union of whole components of }
Z_m^{(t_E^-)};\\
&\text{(F3)}\quad \text{if } i\le k_a,\ A \text{ a component of } Z_i^{\mathrm{cur}},\
|A|\cap|W_a|\neq\emptyset:\ |A|\subset|W_a|;\\
&\text{(F4)}\quad |\Omega_l^{\mathbf{T}}|\cap Z_i^{\mathrm{cur}}=\emptyset,\qquad
|E|\cap Z_i^{\mathrm{cur}}=\emptyset\ \ (k_E\ge l).
\end{aligned}
\end{equation}

(F1) The three inclusions are, in this order, the last clause of (S0) in
\eqref{9.2:eq-attempt-windows-a}, the nesting \eqref{9.2:eq-current-sets-b} of the current sets,
and \eqref{9.2:eq-current-sets-a} together with \eqref{9.2:eq-ten-layer-collars-c}.

(F2) Let $E$ be a reset at stage $t_E$ with $k_E\ge m$. We apply \eqref{9.2:eq-cubes-regions-c}
with $\mathcal{S}=Z_m^{(t_E^-)}$, $\mathcal{T}=Z_{k_E}^{(t_E^-)}$ and $\mathcal{U}=E$; the inclusion
$\mathcal{T}\supseteq\mathcal{S}$ holds by the second inclusion of (F1), and $E$ is a union of whole
components of $\mathcal{T}$ by (S0). It gives: $E\cap Z_m^{(t_E^-)}$ is a union of whole components of
$Z_m^{(t_E^-)}$, and every component of $Z_m^{(t_E^-)}$ not contained in $|E|$ is disjoint from
$|E|$. Consequently $Z_m^{(t)}$ is the union of a subfamily, decreasing in $t$, of the components of
$Z_m^{\mathbf{T}}$; and a component $A$ of $Z_m^{\mathrm{cur}}$ is a component of $Z_m^{(t)}$ for
every $t\in[t_{\mathbf{T}}(m),\mathrm{now}]$ and is disjoint from every reset class $E$ with
$k_E\ge m$ performed so far.

(F3) Let $a$ be an attempt, $i\le k_a$, and $A$ a component of $Z_i^{\mathrm{cur}}$ whose carrier
meets $|W_a|$. By (F2), $A$ is a component of $Z_i^{(t_a^-)}$. We apply
\eqref{9.2:eq-cubes-regions-c} with $\mathcal{S}=Z_i^{(t_a^-)}$, $\mathcal{T}=Z_{k_a}^{(t_a^-)}$ and
$\mathcal{U}=W_a$, the class $W_a$ playing the role of the union of whole components of
$\mathcal{T}$; since $A$ is a component of $\mathcal{S}$ whose carrier meets $|W_a|$, it follows that
$|A|\subset|W_a|$.

(F4) Let $l>i$. The first statement follows from \eqref{9.2:eq-ten-layer-collars-c} and (F1). For a
reset $E$ performed so far with $k_E\ge l$ we have $k_E>i$, so (F2) with $m=i$ shows that every
component of $Z_i^{\mathrm{cur}}$ is disjoint from $|E|$; hence $|E|\cap Z_i^{\mathrm{cur}}=\emptyset$.

\emph{Proof of \textup{(c)}.} Let $k_0(b)+1\le m\le j_b$. In this thin range
$Z''_m(b)\supseteq Z''_{k_0(b)+1}(b)$ as point sets: by clause (S2) of
\eqref{9.2:eq-attempt-windows-a} the sets $Z''_{k_0(b)+2}(b),\dots$ are built from
$Z''_{k_0(b)+1}(b)$ by adding layers, so the sets $Z''$ increase with the level. Hence
\begin{equation*}
|W_b|\setminus|Z''_m(b)|\subset|W_b|\setminus|Z''_{k_0(b)+1}(b)|,
\end{equation*}
and, taking closures and using the description of $U_m(b)$ in the statement,
$U_m(b)\subset U_{k_0(b)+1}(b)$. The second assertion of (c) follows from this inclusion and (b).

\emph{Proof of \textup{(a)} and \textup{(b)}.} Fix $m\in(h_b,k_0(b)+1]$, and put $n:=5$ if
$m\le k_0(b)$ and $n:=7$ if $m=k_0(b)+1$. Let $\Omega_m^{(b)}:=\Omega_m^{(t_b^-)}$ be the determining
domain of level $m$ at the stage $t_b^-$; it is a region of $\mathcal{Q}_m$, the one cube family of
clause (S2). By (S2), for $m\le k_0(b)$ we have
$U_m(b)=|W_b\cap(\Omega_m^{(b)})^{\sim5}|$; for $m=k_0(b)+1$, under reading B,
$U_{k_0(b)+1}(b)=|W_b\cap(\Omega_{k_0(b)+1}^{(b)})^{\sim7}|$; and under reading A,
$Z''_{k_0(b)+1}(b)=(Z'_{k_0(b)})^{\sim1}$, where $Z'_{k_0(b)}$ is the set of cubes of
$\mathcal{Q}_{k_0(b)+1}$ meeting $Z_{k_0(b)}^{(t_b^-)}\cap|W_b|$.

Under reading A, (b) is proved as follows. By (F1),
$Z_{k_0(b)}^{\mathrm{cur}}\subset Z_{k_0(b)}^{(t_b^-)}$. Hence every cube of $W_b$ meeting
$Z_{k_0(b)}^{\mathrm{cur}}$ meets $Z_{k_0(b)}^{(t_b^-)}\cap|W_b|$, so it lies in $Z'_{k_0(b)}$,
which is contained in $Z''_{k_0(b)+1}(b)$. Therefore no cube of $W_b\setminus Z''_{k_0(b)+1}(b)$
meets $Z_{k_0(b)}^{\mathrm{cur}}$, and (b) holds; nothing of $\Omega^{(b)}$ is used.

For reading B and for (a) it suffices to prove
\begin{equation}\label{9.2:eq-frozen-inclusions-1}
\text{no cube of } (\Omega_m^{(b)})^{\sim n},\ n\le 7,\ \text{meets } Z_{m-1}^{\mathrm{cur}}\cap|W_b|.
\end{equation}
Indeed, a cube of $W_b\cap(\Omega_m^{(b)})^{\sim n}$ lies in $|W_b|$, so by
\eqref{9.2:eq-frozen-inclusions-1} it does not meet $Z_{m-1}^{\mathrm{cur}}$; with the formulas for
$U_m(b)$ just recalled this is (a) for $m\le k_0(b)$ and (b) for $m=k_0(b)+1$.

To prove \eqref{9.2:eq-frozen-inclusions-1}, let $A'$ be a component of $Z_{m-1}^{\mathrm{cur}}$
whose carrier meets $|W_b|$; these components cover $Z_{m-1}^{\mathrm{cur}}\cap|W_b|$.

\emph{The component lies in the class.} Since $m-1\le k_b$, (F3) gives $|A'|\subset|W_b|$.

\emph{The component of $Z_m^{\mathbf{T}}$ containing it.} By (F1) and
\eqref{9.2:eq-ten-layer-collars-c},
\begin{equation*}
|A'|\subset Z_{m-1}^{\mathrm{cur}}\subset Z_{m-1}^{(t_{\mathbf{T}}(m)^-)}\subset Z_m^{\mathbf{T}},
\end{equation*}
and $|A'|$ is connected; by \eqref{9.2:eq-cubes-regions-a} there is a component $C_m$ of
$Z_m^{\mathbf{T}}$ with $|A'|\subset|C_m|$. It has three properties.

(i) By (F2), $A'$ is a component of $Z_{m-1}^{(t)}$ for every
$t\in[t_{\mathbf{T}}(m-1),\mathrm{now}]$; in particular of $Z_{m-1}^{(t_{\mathbf{T}}(m)^-)}$, the set
read by the $\mathbf{T}$-step from level $m-1$ to level $m$.

(ii) No reset $E$ with $k_E\ge m$ performed so far contains a cube of $C_m$, and $C_m$ is a component
of $Z_m^{(t)}$ for every $t\in[t_{\mathbf{T}}(m),\mathrm{now}]$. We argue by induction along the
stages. Let $E$ be such a reset and suppose every earlier such reset misses $C_m$. Then, by (F2),
$C_m$ is still a component of $Z_m^{(t_E^-)}$. Suppose $E$ held a cube $Q$ of $C_m$. By (F1),
$Z_m^{(t_E^-)}\subset Z_{k_E}^{(t_E^-)}$, and by \eqref{9.2:eq-cubes-regions-a} there is a component
$\hat{C}$ of $Z_{k_E}^{(t_E^-)}$ containing $|C_m|$. Now $E$ is a union of whole components of
$Z_{k_E}^{(t_E^-)}$ and meets $|\hat{C}|$ in $Q$, so $\hat{C}\subset E$ by
\eqref{9.2:eq-cubes-regions-c}, whence $|A'|\subset|C_m|\subset|E|$. Since $m-1\le k_E$, the reset
of $E$ deletes $A'$ from $Z_{m-1}$ (clause (S0), and the statement, used in the proof of
Lemma~\ref{9.2:lem-ten-layer-collars}, that a component so deleted is not current), so that $A'$ is
not current; this contradicts the choice of $A'$ as a component of $Z_{m-1}^{\mathrm{cur}}$. Hence $E$
misses $C_m$, and the induction proceeds.

(iii) If $c$ is a non-completed attempt with $k_c\ge m$ whose class contains a cube of $C_m$, then
$|C_m|\subset|W_c|$: this is the argument of (ii) with $W_c$ in place of $E$ ($C_m$ is a component
of $Z_m^{(t_c^-)}$ by (ii), lies by (F1) and \eqref{9.2:eq-cubes-regions-a} in a component of
$Z_{k_c}^{(t_c^-)}$, which $W_c$ meets and therefore contains by \eqref{9.2:eq-cubes-regions-c}).
In particular $|C_m|\subset|W_b|$: a cube $q$ of $A'$ has
$\operatorname{int}q\subset|W_b|\cap|C_m|$; by \eqref{9.2:eq-cubes-regions-b} the cube $Q$ of
$\mathcal{Q}_m$ containing $q$ is a cube of $C_m$ inside $|W_b|$; and $k_b\ge m$.

\emph{A fattened neighbourhood of the component.} Let $Z'_{A'}$ be the set of cubes of
$\mathcal{Q}_m$ meeting $|A'|$, and put $\mathcal{V}_{A'}:=(Z'_{A'})^{\sim7}$. Since $|A'|$ is
contained in $Z_{m-1}^{(t_{\mathbf{T}}(m)^-)}$, the set read by the $\mathbf{T}$-step from level
$m-1$ to level $m$ (property (i)), we have $Z'_{A'}\subset Z'_{m-1}$, where $Z'_{m-1}$ is the set of
cubes of $\mathcal{Q}_m$ meeting that set; hence $\mathcal{V}_{A'}\subset(Z'_{m-1})^{\sim8}$. So
every cube of $\mathcal{V}_{A'}$ is a cube of $Z_m^{\mathbf{T}}$, by \eqref{9.2:eq-ten-layer-collars-d}.
The region $\mathcal{V}_{A'}$ is connected and contains the cubes of $C_m$ holding the cubes of $A'$;
hence $\mathcal{V}_{A'}\subset C_m$ by \eqref{9.2:eq-cubes-regions-a}.

\emph{No cube of $\Omega_m^{(b)}$ lies in $\mathcal{V}_{A'}$.} Let $Q_0$ be a cube of
$\Omega_m^{(b)}$ and suppose $Q_0\in\mathcal{V}_{A'}$. By the exhaustiveness in (S0) at the stage
$t_b^-$, one of the following three alternatives holds.

First, $Q_0\subset|\Omega_m^{\mathbf{T}}|$. Then $Q_0$ is a cube of $\Omega_m^{\mathbf{T}}$ by
\eqref{9.2:eq-cubes-regions-b}, lying inside $(Z'_{m-1})^{\sim8}$; this is excluded by
\eqref{9.2:eq-ten-layer-collars-b}. Of the eight layers of the collar only the inequality $7\le8$,
respectively $5\le8$, is used, and the fattening and the collar are counted in the same cube
family $\mathcal{Q}_m$.

Second, $Q_0\subset|E|$ for a reset $E$ performed before $t_b$ with $k_E\ge m$ (then $k_E\le k_b$; a
class reset earlier in the same step $k_b$ is included). Then $E$ contains the cube $Q_0$ of $C_m$,
which is excluded by property (ii).

Third, $Q_0\subset U_m(c')\subset|W_{c'}|$ for a non-completed attempt $c'$ performed before $t_b$,
so $c'\neq b$, with $h_{c'}<m\le j_{c'}$, so that $k_{c'}\ge j_{c'}+1>m$. Then $W_{c'}$ contains the
cube $Q_0$ of $C_m$, and property (iii) gives $|C_m|\subset|W_{c'}|$. Together with
$|C_m|\subset|W_b|$ this shows that the carriers of $b$ and $c'$ meet, and clause (S4a) gives
$I_b\cap I_{c'}=\emptyset$. But
\begin{equation*}
\begin{gathered}
h_b<m\le k_0(b)+1\le k_b,\qquad h_{c'}<m\le j_{c'}\le k_{c'}-1,\\
\text{so}\quad m-1\in[h_b,k_b-1]\cap[h_{c'},k_{c'}-1],
\end{gathered}
\end{equation*}
a contradiction. Equivalently: if $k_{c'}<k_b$, then
\begin{equation*}
k_{c'}\le h_b<m\le j_{c'}<k_{c'},
\end{equation*}
which is absurd; if $k_{c'}=k_b$, then $b$ and $c'$ are two classes of one step with meeting
carriers, which is absurd by clause (S4b); and an attempt with $k_{c'}>k_b$ cannot have been
performed before $t_b$.

\emph{Conclusion of \eqref{9.2:eq-frozen-inclusions-1}.} Suppose a cube $Q$ of
$(\Omega_m^{(b)})^{\sim n}$, $n\le7$, met $|A'|$. The transfer property
\eqref{9.2:eq-cubes-regions-d} would give a cube
\begin{equation*}
Q_0\in\Omega_m^{(b)}\cap(Z'_{A'})^{\sim n}\subset\mathcal{V}_{A'},
\end{equation*}
which has just been excluded. Since the components $A'$ cover $Z_{m-1}^{\mathrm{cur}}\cap|W_b|$,
this proves \eqref{9.2:eq-frozen-inclusions-1}, and with it (a) and (b).

\emph{Persistence.} \cite{B-CMP122a} states that the sets $Z''_m(b)$ are not changed by the
subsequent operations. By (F1) the current sets $Z^{\mathrm{cur}}$ only shrink. By clause (S5) the
$\mathbf{T}$-steps performed later make sets only at levels greater than the present level $k$, and
$k\ge k_b$. A later attempt $c$ in force whose carrier meets $|W_b|$ has, by clause (S4a),
$h_c\ge k_b>m$, and re-makes nothing at level $m$. Hence the inclusions
\eqref{9.2:eq-frozen-inclusions-b} are not altered by the operations that follow.
\end{proof}

\begin{lemma}[Which levels can be thin]\label{9.2:lem-thin-levels}
Assume the following:
\begin{itemize}
\item[(1)] the reading \textup{(S0)}--\textup{(S6)} of the current label by attempt windows,
Definition~\ref{9.2:def-attempt-windows}, displayed in \eqref{9.2:eq-attempt-windows-a};
\item[(2)] the collar inclusions \eqref{9.2:eq-ten-layer-collars-b} and
\eqref{9.2:eq-ten-layer-collars-c} of Lemma~\ref{9.2:lem-ten-layer-collars}, in cube form and
with the sets current at the stage to which they are applied;
\item[(3)] the cube-wise description \eqref{9.2:eq-current-sets-a} of the step-made sets,
Definition~\ref{9.2:def-current-sets};
\item[(4)] the disjointness \eqref{9.2:eq-ten-layer-collars-e} of $\Omega$ from $Z_k$.
\end{itemize}
Let $y\in Z_i$ with $j+1\le i\le k$, and suppose that some $x\in A_i(y)$ has local scale
$n(x)=:l>i$. Then there is exactly one $b\in\mathfrak{A}(y)$ with $i\in\Theta_b$, and for it
\begin{equation}\label{9.2:eq-thin-levels-a}
\begin{gathered}
i\in[k_0(b)+1,\,j_b-1]\subset\Theta_b,\\
n(x')\le j_b\le k_b-1\le k-1\qquad\text{for every }x'\in A_i(y).
\end{gathered}
\end{equation}
Consequently, with thin and thick levels as in \eqref{9.2:eq-attempt-windows-b}:
\begin{itemize}
\item[(i)] a level $i$ that is thin for $y$ lies in $\Theta_b$ for a unique $b\in\mathfrak{A}(y)$;
\item[(ii)] the level $k$ is never thin;
\item[(iii)] the $\mathbf{T}$-band $\Lambda_k\setminus\Omega$ has scale $k$;
\item[(iv)] if the level $i$ is thick for $y$, every current cube meeting $A_i(y)$ has scale at
most $i$.
\end{itemize}
\end{lemma}

\begin{proof}
Throughout we use Lemma~\ref{9.2:lem-frozen-inclusions}: its inclusions \textup{(F1)}--\textup{(F4)} of
\eqref{9.2:eq-frozen-inclusions-a} and its claims \textup{(a)}, \textup{(b)}, \textup{(c)} of
\eqref{9.2:eq-frozen-inclusions-b} on the re-made strata $U_m(b)$ of an attempt $b$ in force.
Here and below $Z_m$ denotes the current set $Z_m^{\mathrm{cur}}$; since $A_i(y)$ is a current
component of $Z_i$, we have $x\in|A_i(y)|\subset Z_i^{\mathrm{cur}}$.

\emph{Step 0: $i<l\le k$.} The only set of scale $k+1$ is $\Omega$. If $l=k+1$, then
$x\in\Omega$; but
\begin{equation*}
x\in|A_i(y)|\subset Z_i\subset Z_k,
\end{equation*}
and $\Omega\cap Z_k=\emptyset$ by \eqref{9.2:eq-ten-layer-collars-e}. Hence $l\ne k+1$, so
$i<l\le k$, and, $l$ being the local scale of $x$, the point $x$ lies in a cube $Q$ of
$\Omega_l^{\mathrm{cur}}$.

We now apply \textup{(S0)} at the present stage of the label. By \textup{(S0)} the current
label arises from the $\mathbf{T}$-made sets by three kinds of operations only: $\mathbf{T}$-steps,
completed $\mathbf{R}$-operations (resets), and the preparations of non-completed attempts.
Accordingly, the cube $Q$ of $\Omega_l^{\mathrm{cur}}$ falls under one of the following three
cases.

\emph{Case \textup{(s1)}: $Q\subset|\Omega_l^{\mathbf{T}}|$.} Then
$x\in|\Omega_l^{\mathbf{T}}|\cap Z_i^{\mathrm{cur}}$. Since $i\le l-1$, this intersection is
empty by \textup{(F4)}. This case is impossible.

\emph{Case \textup{(s2)}: $Q\subset|E|$, where $E$ is the class of a reset with step
$k_E\ge l>i$.} Then $x\in|E|\cap Z_i^{\mathrm{cur}}$, and this intersection is empty by
\textup{(F4)}. This case is impossible.

\emph{Case \textup{(s3)}.} By the two preceding cases, $Q\subset U_l(b)$ for an attempt $b$ in
force with $h_b<l\le j_b$, and $x\in|W_b|$. By \textup{(S1)},
\begin{equation*}
i\le l-1\le j_b-1\le k_b-2 .
\end{equation*}
Hence $A_i(y)=A_i(x)$ is a current component of $Z_i$ with $i\le k_b$, and it meets $|W_b|$
at $x$. By \textup{(F3)}, $|A_i(y)|\subset|W_b|$. Thus $y\in W_b$, that is,
$b\in\mathfrak{A}(y)$.

Suppose first that $l\le k_0(b)$. Then $h_b<l\le\min\{k_0(b),j_b\}$, and claim \textup{(a)} of
\eqref{9.2:eq-frozen-inclusions-b}, applied with $m=l$, gives
$U_l(b)\cap Z_{l-1}^{\mathrm{cur}}=\emptyset$. Since $x\in Q\subset U_l(b)$, we get
$x\notin Z_{l-1}^{\mathrm{cur}}$. By \textup{(F1)} and $i\le l-1$ we have
$Z_{l-1}^{\mathrm{cur}}\supseteq Z_i^{\mathrm{cur}}$, so $x\notin Z_i^{\mathrm{cur}}$. This
contradicts $x\in Z_i^{\mathrm{cur}}$.

Therefore $k_0(b)+1\le l\le j_b$. Claims \textup{(b)} and \textup{(c)} of
\eqref{9.2:eq-frozen-inclusions-b} give
\begin{equation*}
x\in U_l(b)\subset U_{k_0(b)+1}(b),\qquad
U_{k_0(b)+1}(b)\cap Z_{k_0(b)}^{\mathrm{cur}}=\emptyset .
\end{equation*}
Hence $x\notin Z_{k_0(b)}^{\mathrm{cur}}$. By \textup{(F1)},
$Z_m^{\mathrm{cur}}\subset Z_{k_0(b)}^{\mathrm{cur}}$ for $m\le k_0(b)$, so $x\notin
Z_m^{\mathrm{cur}}$ for every $m\le k_0(b)$. For the depth of $x$ this means
$\iota(x)\ge k_0(b)+1$. On the other hand, $x\in Z_i$ gives $\iota(x)\le i$. Therefore
\begin{equation*}
k_0(b)+1\le i\le l-1\le j_b-1\le k_b-2 .
\end{equation*}
In particular $i\in[k_0(b)+1,\,j_b-1]$. By Definition~\ref{9.2:def-attempt-windows}, this
interval is contained in $\Theta_b$. This is the first line of \eqref{9.2:eq-thin-levels-a}.

\emph{Uniqueness.} Let $b'\in\mathfrak{A}(y)$ with $b'\ne b$, and suppose that
$i\in\Theta_{b'}$. Since $\Theta_{b'}\subset I_{b'}$ and $\Theta_b\subset I_b$, we get
$i\in I_b\cap I_{b'}$. Moreover $y\in W_b\cap W_{b'}$. This contradicts \textup{(S4a)}. Hence
$b$ is the only element of $\mathfrak{A}(y)$ whose thin part contains $i$.

\emph{Scale bound on $A_i(y)$.} Let $x'\in A_i(y)$. If $n(x')\le i$, then
$n(x')\le i\le j_b-1$. Suppose instead that $n(x')=:l'>i$. Repeating the argument from Step~0
through the first line of \eqref{9.2:eq-thin-levels-a} with $x'$ in place of $x$ gives an attempt
$b'\in\mathfrak{A}(y)$ in force with $i\in\Theta_{b'}$ and $l'\le j_{b'}$. By uniqueness,
$b'=b$, so $l'\le j_b$. In both cases $n(x')\le j_b$. Moreover $j_b\le k_b-1$ by \textup{(S1)}
and $k_b\le k$ by \textup{(S5)}, so
\begin{equation*}
n(x')\le j_b\le k_b-1\le k-1 .
\end{equation*}
This is the second line of \eqref{9.2:eq-thin-levels-a}.

\emph{Consequence \textup{(i)}.} A level $i$ that is thin for $y$ is one at which some point of
$A_i(y)$ has local scale greater than $i$. The first part of the lemma then applies and gives
the unique $b\in\mathfrak{A}(y)$ with $i\in\Theta_b$.

\emph{Consequence \textup{(ii)}.} Take $i=k$. A point $x\in A_k(y)$ with $n(x)>k$ would lie in
$\Omega$, by Step~0. But $x\in Z_k$ and $\Omega\cap Z_k=\emptyset$ by
\eqref{9.2:eq-ten-layer-collars-e}. Hence no such point exists, and the level $k$ is never thin.

\emph{Consequence \textup{(iii)}.} By \eqref{9.2:eq-current-sets-a},
$\Lambda_k^{\mathbf{T}}\subset\Omega_k^{\mathbf{T}}$ cube-wise. By \textup{(S0)}, the operations
performed after the stage $t_{\mathbf{T}}(k)$ and before the present stage are of two kinds.

The first kind are resets with $k_E\le k$. If $k_E=k$, the reset keeps $|E|\subset|\Omega_k|$.
If $k_E<k$, then by \textup{(S0)} and \textup{(F1)}
\begin{equation*}
|E|\subset Z^{(t_E^{-})}_{k_E}\subset Z^{(t_{\mathbf{T}}(k)^{-})}_{k-1},
\end{equation*}
and by \eqref{9.2:eq-ten-layer-collars-c} this set shares no cube with
$\Omega_k^{\mathbf{T}}$.

The second kind are preparations of attempts $a$, which re-make levels at most
$j_a\le k-1$.

None of these operations removes a cube of $\Omega_k^{\mathbf{T}}$. Hence every cube of
$\Lambda_k^{\mathbf{T}}$ is still a current cube of $\Omega_k$, so each of its points has local
scale at least $k$. Outside $\Omega$ the local scale is at most $k$, since $\Omega$ is the only
set of scale $k+1$. Therefore every point of $\Lambda_k\setminus\Omega$ has scale $k$.

\emph{Consequence \textup{(iv)}.} Let the level $i$ be thick for $y$, and let a current cube meet
$A_i(y)$ at a point $x$. The local scale $n(x)$ is at least the scale of the cube, since $x$
lies in it. It is at most $i$, since the level is thick. Hence the scale of the cube is at
most $i$.
\end{proof}

\begin{remark}
The proof of consequence \textup{(iii)} of Lemma~\ref{9.2:lem-thin-levels} rests on two inputs.
The first is the reading \textup{(S0)} of Definition~\ref{9.2:def-attempt-windows}, that is,
the list of the operations performed between $t_{\mathbf{T}}(k)$ and the present stage and
their effect on the cubes of $\Omega_k^{\mathbf{T}}$. The second consists of
\eqref{9.2:eq-current-sets-a}, \eqref{9.2:eq-ten-layer-collars-c} and \textup{(F1)} of
\eqref{9.2:eq-frozen-inclusions-a}.

Consequence \textup{(iii)} is not used in what follows. Two places might appear to need it: the
portion of a path crossing the $\mathbf{T}$-band, and the non-terminal points outside $\Omega$
of a path in the case $\iota=k+1$. In both, clause \textup{(D3)} of
Definition~\ref{9.2:def-scaled-distance}, displayed in \eqref{9.2:eq-scaled-distance-a}, already
gives what is needed. It supplies that every non-terminal point of the path lies only in
current cubes of scale at most $k$, and this is all that clause \textup{(D2)} there requires
with $m=k$. No bound and no floor on the auxiliary parameters $\mathfrak{p}$ depends on the
consequence.
\end{remark}

\begin{remark}
Two further situations do not affect the conclusions of Lemma~\ref{9.2:lem-thin-levels}.

\emph{Outer attempts.} Let $c''$ be an attempt in force with $k_{c''}>k_b$ whose carrier meets
$|W_b|$. By \textup{(S4a)}, $h_{c''}\ge k_b$. Lemma~\ref{9.2:lem-frozen-inclusions} then shows
that the re-made strata $U_m(c'')$, for $m\ge h_{c''}+1\ge k_b+1$, miss $Z_{k_b}^{\mathrm{cur}}$.
This set satisfies
\begin{equation*}
Z_{k_b}^{\mathrm{cur}}\supseteq|W_b|\cap Z_{k_b}^{\mathrm{cur}}\supseteq A_i(y)
\end{equation*}
for every $A_i(y)$ occurring in the lemma. Hence $c''$ raises no local scale there.

\emph{Second-class outcomes.} A second-class outcome may be treated either as a completed reset
or as the in-force re-made strata of an attempt. In the first case Case \textup{(s2)} of the
proof of the lemma applies without change, in the second Case \textup{(s3)}; so the lemma holds
in both. In both cases the condition \textup{(S4a)} for such outcomes is supplied by the
classification of second-class outcomes in \cite{B-CMP122b}.
\end{remark}

\begin{remark}
The exclusion of the case $l\le k_0(b)$ in the proof of Lemma~\ref{9.2:lem-thin-levels}, through
claim \textup{(a)} of \eqref{9.2:eq-frozen-inclusions-b}, depends on one requirement of
\textup{(S2)} in Definition~\ref{9.2:def-attempt-windows}: the five-fold fattening
$\mathcal{S}^{\sim 5}$ entering \textup{(S2)} and the collars of
\eqref{9.2:eq-ten-layer-collars-b} are taken in one and the same cube family. If the five-fold
fattening were taken in a family coarser than that of \eqref{9.2:eq-ten-layer-collars-b}, five
coarse layers could exceed eight fine ones.
\end{remark}

\begin{lemma}[The response side: scaled distance across collars]\label{9.2:lem-response-side}
Assume the following statements:
\begin{itemize}
\item[(a)] properties (D1)--(D3) of the scaled distance, \eqref{9.2:eq-scaled-distance-a}, in Definition~\ref{9.2:def-scaled-distance};
\item[(b)] the collar bounds \eqref{9.2:eq-ten-layer-collars-a} and \eqref{9.2:eq-ten-layer-collars-e} of Lemma~\ref{9.2:lem-ten-layer-collars};
\item[(c)] part~(b) of Lemma~\ref{9.2:lem-nested-distances};
\item[(d)] Lemma~\ref{9.2:lem-thin-levels}.
\end{itemize}
Let $y\notin\Omega$, and let $\Gamma\colon[0,T]\to\mathbb{T}_m$ be a continuous path with
$\Gamma(0)=y$, $\Gamma(T)\in\Omega$ and $\Gamma(\tau)\notin\Omega$ for $\tau<T$. Put $\iota:=\iota(y)$.

\emph{Case $\iota\le k$.} For $i\in[\iota,k]$ let
\[
e_i:=\inf\{\tau\in[0,T]:\Gamma(\tau)\notin Z_i\}.
\]
Then $0\le e_\iota\le\dots\le e_k<T$. Consider the consecutive portions
\[
P_\iota:=\Gamma|_{[0,e_\iota]},\qquad
P_i:=\Gamma|_{[e_{i-1},e_i]}\ \ (\iota<i\le k),\qquad
P_T:=\Gamma|_{[e_k,T]}.
\]
They satisfy:
\begin{itemize}
\item[(1)] $\ell_{\mathrm{sc}}(P_i)\ge 10\check{R}_i\,c_i(y)$ for $\iota<i\le k$;
\item[(2)] if $\iota\ge j+1$, then $\ell_{\mathrm{sc}}(P_\iota)\ge c_\iota(y)\,(10\check{R}_\iota-v)_+$, where
$v:=t/(M\ell_\iota)$ and $t:=\operatorname{dist}_\infty(y,Z_{\iota-1})$;
\item[(3)] $\ell_{\mathrm{sc}}(P_T)\ge 8L\check{R}_{k+1}$.
\end{itemize}

\emph{Case $\iota=k+1$.} The path satisfies
\[
\ell_{\mathrm{sc}}(\Gamma)\ge(8L\check{R}_{k+1}-u)_+,\qquad
u:=\operatorname{dist}_\infty(y,Z_k)/(M\ell_k).
\]

Consequently
\begin{equation}\label{9.2:eq-response-side-a}
d_{\mathrm{sc}}(y,\Omega)\ \ge\
\begin{cases}
c_\iota(y)(10\check{R}_\iota-v)_+\,\mathbf{1}[\iota\ge j+1]
+\sum_{i=\iota+1}^{k}10\check{R}_i\,c_i(y)+8L\check{R}_{k+1}, & \iota\le k,\\[2pt]
(8L\check{R}_{k+1}-u)_+, & \iota=k+1.
\end{cases}
\end{equation}
Here $\mathbf{1}[\iota\ge j+1]$ is $1$ if $\iota\ge j+1$ and $0$ otherwise.
\end{lemma}

\begin{proof}
Throughout, the sets $Z_i$ are the current large-field sets of \eqref{9.2:eq-current-sets-b}. They are nested, $Z_{i-1}\subset Z_i$. The number $c_i(y)\in\{0,1\}$ is the thick indicator of \eqref{9.2:eq-attempt-windows-b}, and $A_i(y)$ is the component of $Z_i$ that contains $y$, in the sense of \eqref{9.2:eq-cubes-regions-a}.

We first record that $\Omega$ lies in the complement of every $Z_i$ with $i\le k$. By \eqref{9.2:eq-ten-layer-collars-e},
\[
\operatorname{dist}_\infty(Z_k,\Omega)\ge 8LM\check{R}_{k+1}\ell_k>0 .
\]
Hence $\Omega\cap Z_k=\emptyset$. Since $Z_i\subset Z_k$ for $i\le k$, this gives $\Omega\subset Z_i^c$.

\emph{Exit times.} Let $\iota\le k$ and fix $i\in[\iota,k]$. Put
\[
E_i:=\{\tau\in[0,T]:\Gamma(\tau)\notin Z_i\}.
\]
Since $\Gamma(T)\in\Omega\subset Z_i^c$, we have $T\in E_i$. So $E_i$ is non-empty and $e_i=\inf E_i\in[0,T]$.

By the definition of the depth $\iota(y)$, $\Gamma(0)=y\in Z_\iota$. By nesting, $Z_\iota\subset Z_i$, so $0\notin E_i$.

For $\tau<e_i$ we have $\tau\notin E_i$, that is, $\Gamma(\tau)\in Z_i$. The point $\Gamma(e_i)$ also lies in $Z_i$:
\begin{itemize}
\item if $e_i>0$, it is the limit from the left of the points $\Gamma(\tau)$, $\tau<e_i$, by continuity of $\Gamma$, and these points lie in the closed set $Z_i$;
\item if $e_i=0$, it equals $y\in Z_i$.
\end{itemize}
Thus $e_i\notin E_i$. As $e_i$ is the infimum of $E_i$ but not a member of it, there are $\tau_n\in E_i$ with $\tau_n\downarrow e_i$. By continuity, $\Gamma(e_i)=\lim\Gamma(\tau_n)\in\operatorname{cl}(Z_i^c)$. Moreover $e_i<T$, since $T\in E_i$ and $e_i\notin E_i$.

Finally, for $\iota<i\le k$ the inclusion $Z_{i-1}\subset Z_i$ gives $E_{i-1}\supset E_i$, hence $e_{i-1}\le e_i$. This proves $0\le e_\iota\le\dots\le e_k<T$.

\emph{The portions lie in the components of $y$.} The set $\Gamma([0,e_i])$ is connected, as the continuous image of an interval. It is contained in $Z_i$ by the previous step, and it contains $y$. Hence it is contained in the component $A_i(y)$ of $Z_i$ containing $y$, by \eqref{9.2:eq-cubes-regions-a}. Consequently:
\begin{itemize}
\item every point of $P_i$, $\iota<i\le k$, lies in $A_i(y)$;
\item every point of $P_\iota$ lies in $A_\iota(y)$.
\end{itemize}
If $\Gamma$ re-enters $Z_i$ after the time $e_i$, those times belong to later portions and are never credited to level $i$.

The portions $P_\iota,P_{\iota+1},\dots,P_k,P_T$ are consecutive and together make up $\Gamma$. Therefore (D1) in \eqref{9.2:eq-scaled-distance-a} gives
\begin{equation}\label{9.2:eq-response-side-1}
\ell_{\mathrm{sc}}(\Gamma)\ge\ell_{\mathrm{sc}}(P_\iota)+\sum_{i=\iota+1}^{k}\ell_{\mathrm{sc}}(P_i)
+\ell_{\mathrm{sc}}(P_T),
\end{equation}
and all terms on the right are non-negative.

\emph{The portions $P_i$, $\iota<i\le k$.} The portion $P_i$ starts at $\Gamma(e_{i-1})\in Z_{i-1}$ and ends at $\Gamma(e_i)\in\operatorname{cl}(Z_i^c)$. The distance from a point to a set equals its distance to the closure of that set. So the sup-norm distance between the endpoints of $P_i$ is at least $\operatorname{dist}_\infty(Z_{i-1},Z_i^c)$.

The collar bound \eqref{9.2:eq-ten-layer-collars-a} at level $i$ holds for every current component of $Z_{i-1}$. Taking the infimum over these components gives
\[
\operatorname{dist}_\infty(Z_{i-1},Z_i^c)\ge 10M\check{R}_i\ell_i .
\]

Suppose $c_i(y)=1$, that is, level $i$ is thick for $y$. Every current cube containing a point of $P_i$ meets $A_i(y)$, by the previous step. Since $y\in Z_i$, the last assertion of Lemma~\ref{9.2:lem-thin-levels} shows that every such cube has scale at most $i$. We apply (D2) in \eqref{9.2:eq-scaled-distance-a} with $S=10M\check{R}_i\ell_i$ and $m=i$. This gives
\[
\ell_{\mathrm{sc}}(P_i)\ge 10\check{R}_i .
\]
If $c_i(y)=0$, then $\ell_{\mathrm{sc}}(P_i)\ge0=10\check{R}_i\,c_i(y)$. This proves~(1).

\emph{The portion $P_\iota$, $\iota\ge j+1$.} The portion $P_\iota$ starts at $y$ and ends at $\Gamma(e_\iota)\in\operatorname{cl}(Z_\iota^c)$. Its endpoint distance is therefore at least $\operatorname{dist}_\infty(y,Z_\iota^c)$.

We apply part~(b) of Lemma~\ref{9.2:lem-nested-distances} with $A=Z_{\iota-1}\subset B=Z_\iota$, and then the collar bound \eqref{9.2:eq-ten-layer-collars-a} at level $\iota\in[j+1,k]$:
\[
\operatorname{dist}_\infty(y,Z_\iota^c)
\ge\operatorname{dist}_\infty(Z_{\iota-1},Z_\iota^c)-\operatorname{dist}_\infty(y,Z_{\iota-1})
\ge 10M\check{R}_\iota\ell_\iota-t .
\]

Suppose $c_\iota(y)=1$ and $10M\check{R}_\iota\ell_\iota-t>0$. Every current cube containing a point of $P_\iota$ meets $A_\iota(y)$. By the last assertion of Lemma~\ref{9.2:lem-thin-levels}, every such cube has scale at most $\iota$. Then (D2) in \eqref{9.2:eq-scaled-distance-a}, with $m=\iota$, gives
\[
\ell_{\mathrm{sc}}(P_\iota)\ge\frac{10M\check{R}_\iota\ell_\iota-t}{M\ell_\iota}=10\check{R}_\iota-v .
\]
Otherwise, either $c_\iota(y)=0$ or $(10\check{R}_\iota-v)_+=0$. In both cases $c_\iota(y)(10\check{R}_\iota-v)_+=0$, and the bound $\ell_{\mathrm{sc}}(P_\iota)\ge0$ is the assertion. This proves~(2).

For $\iota=j$ nothing is credited to $P_\iota$; we only use $\ell_{\mathrm{sc}}(P_\iota)\ge0$.

\emph{The terminal portion $P_T$.} The portion $P_T$ starts at $\Gamma(e_k)\in Z_k$ and ends at $\Gamma(T)\in\Omega$. By \eqref{9.2:eq-ten-layer-collars-e}, its endpoint distance is at least
\[
\operatorname{dist}_\infty(Z_k,\Omega)\ge 8LM\check{R}_{k+1}\ell_k .
\]
Every non-terminal point of $P_T$ is of the form $\Gamma(\tau)$ with $\tau<T$, so it lies outside $\Omega$. By (D3) in \eqref{9.2:eq-scaled-distance-a}, it therefore lies in cubes of scale at most $k$; for the band $\Lambda_k\setminus\Omega$, Lemma~\ref{9.2:lem-thin-levels} gives scale $k$. We apply (D2) in \eqref{9.2:eq-scaled-distance-a} with $m=k$; the terminal point is exempt there. This gives
\[
\ell_{\mathrm{sc}}(P_T)\ge\frac{8LM\check{R}_{k+1}\ell_k}{M\ell_k}=8L\check{R}_{k+1},
\]
which is~(3).

\emph{The infimum over paths, case $\iota\le k$.} We insert (1)--(3) into \eqref{9.2:eq-response-side-1}, and use $\ell_{\mathrm{sc}}(P_\iota)\ge0$ when $\iota=j$. This gives
\[
\ell_{\mathrm{sc}}(\Gamma)\ge c_\iota(y)(10\check{R}_\iota-v)_+\,\mathbf{1}[\iota\ge j+1]
+\sum_{i=\iota+1}^{k}10\check{R}_i\,c_i(y)+8L\check{R}_{k+1}.
\]
The right member does not depend on $\Gamma$. Taking the infimum over $\Gamma$ in the definition of $d_{\mathrm{sc}}(y,\Omega)$ gives the first line of \eqref{9.2:eq-response-side-a}. If no such $\Gamma$ exists, then $d_{\mathrm{sc}}(y,\Omega)=+\infty$ and there is nothing to prove.

\emph{The case $\iota=k+1$.} By the definition of the depth and the hypothesis, $y\notin Z_k\cup\Omega$. Since $\Omega\cap Z_k=\emptyset$, we have $Z_k\subset\Omega^c$. The endpoint distance of $\Gamma$ is at least $\operatorname{dist}_\infty(y,\Omega)$. We apply part~(b) of Lemma~\ref{9.2:lem-nested-distances} with $A=Z_k\subset B=\Omega^c$, for which $B^c=\Omega$, and then \eqref{9.2:eq-ten-layer-collars-e}:
\[
\operatorname{dist}_\infty(y,\Omega)
\ge\operatorname{dist}_\infty(Z_k,\Omega)-\operatorname{dist}_\infty(y,Z_k)
\ge 8LM\check{R}_{k+1}\ell_k-uM\ell_k .
\]
The non-terminal points of $\Gamma$ lie outside $\Omega$. By (D3) in \eqref{9.2:eq-scaled-distance-a}, they lie in cubes of scale at most $k$.

If $8L\check{R}_{k+1}-u>0$, then (D2) in \eqref{9.2:eq-scaled-distance-a}, with $m=k$, gives
\[
\ell_{\mathrm{sc}}(\Gamma)\ge 8L\check{R}_{k+1}-u .
\]
Otherwise the bound $\ell_{\mathrm{sc}}(\Gamma)\ge0$ is the assertion. Hence $\ell_{\mathrm{sc}}(\Gamma)\ge(8L\check{R}_{k+1}-u)_+$ in all cases. The right member does not depend on $\Gamma$. Taking the infimum over $\Gamma$, with $d_{\mathrm{sc}}(y,\Omega)=+\infty$ if no such $\Gamma$ exists, gives the second line of \eqref{9.2:eq-response-side-a}.
\end{proof}

\begin{lemma}[Combined size and response exponent]\label{9.2:lem-combined-exponent}
Assume the following statements:
\begin{itemize}
\item[(i)] the size side, Lemma~\ref{9.2:lem-size-side}, with $a=0$, which is the case of anchors
as in \eqref{9.2:eq-anchors-a};
\item[(ii)] the response side, Lemma~\ref{9.2:lem-response-side};
\item[(iii)] the form of the response profile, in terms of the scaled distance
$d_{\mathrm{sc}}(y,\Omega)$ to the domain $\Omega$ of Definition~\ref{9.2:def-scaled-distance},
\begin{equation*}
\rho(y)=\min\bigl\{1,\;C_b\,e^{-\delta_1^{\mathrm{sc}}(d_{\mathrm{sc}}(y,\Omega)-1)}\bigr\},
\end{equation*}
as stated in Proposition~\ref{9.2:prop-response-factor}, whose proof is incomplete at one step;
\item[(iv)] the monotonicity of the radii, parts (a) and (b) of \eqref{9.2:eq-radii-floor-a}
(Notation~\ref{9.2:not-radii-floor});
\item[(v)] the floor $\check{R}_k\ge 2$.
\end{itemize}
Let $X$ be a unit, with its level $j$, the current level $k$, $s:=k-j$, the point $y^{\circ}\in X$,
the weight $w(X)$, the exponent $E'(X)$, the constant $C_w$, the fraction $\varpi$ and the common rate
$\delta_T=\min\{\varpi\delta_1^{\mathrm{sc}},\kappa/2\}$ as in \eqref{9.2:eq-response-factor-a}, and
put $\check{R}^{\flat}:=\check{R}_k$. Put $\iota^{\circ}:=\iota(y^{\circ})$ and $\iota:=\iota(X)$;
for $i\in[\max\{\iota^{\circ},j+1\},k]$ put $c_i^{\circ}:=c_i(y^{\circ})$, the thick indicator of
\eqref{9.2:eq-attempt-windows-b} at $y^{\circ}$; let
\begin{equation*}
U(X):=\bigl\{i\in[\max\{\iota^{\circ},j+1\},k]:\ c_i^{\circ}=0\bigr\}\subset[j+1,k],
\qquad P:=[j+1,k]\setminus U(X),
\end{equation*}
so that $\#P=s-\#U(X)$. Then
\begin{equation}\label{9.2:eq-combined-exponent-a}
w(X)\le C_w\,e^{-\delta_T E'(X)},
\end{equation}
and
\begin{equation}\label{9.2:eq-combined-exponent-b}
\begin{split}
E'(X)&\ge 10\sum_{i\in P}\check{R}_i+8L\check{R}_{k+1}-3\\
&\ge 10\check{R}^{\flat}\bigl(s-\#U(X)\bigr)+\tfrac{13}{2}\check{R}_k .
\end{split}
\end{equation}
Consequently
\begin{equation}\label{9.2:eq-combined-exponent-1}
w(X)\le C_w\,e^{-\frac{13}{2}\delta_T\check{R}_k}\,
e^{-10\delta_T\check{R}^{\flat}(s-\#U(X))},
\end{equation}
and if $s=0$, then $P=U(X)=\emptyset$ and $E'(X)\ge 8L\check{R}_{k+1}-3\ge\tfrac{13}{2}\check{R}_k$.
\end{lemma}

\begin{proof}
\emph{The bound \eqref{9.2:eq-combined-exponent-a}.} Let $y$ be any point. By hypothesis (iii),
since $\min\{1,b\}\le b$ for $b>0$ and $\varpi>0$,
\begin{equation*}
\begin{split}
\rho(y)^{\varpi}
&=\min\bigl\{1,\;C_b\,e^{-\delta_1^{\mathrm{sc}}(d_{\mathrm{sc}}(y,\Omega)-1)}\bigr\}^{\varpi}
\le\bigl(C_b\,e^{\delta_1^{\mathrm{sc}}}\bigr)^{\varpi}
e^{-\varpi\delta_1^{\mathrm{sc}}d_{\mathrm{sc}}(y,\Omega)}\\
&\le C_w\,e^{-\delta_T d_{\mathrm{sc}}(y,\Omega)} .
\end{split}
\end{equation*}
The last inequality uses $d_{\mathrm{sc}}(y,\Omega)\ge0$ together with
$\varpi\delta_1^{\mathrm{sc}}\ge\delta_T$, which holds because $\delta_T$ is the minimum of
$\varpi\delta_1^{\mathrm{sc}}$ and $\kappa/2$, and it uses that the constant $C_w$ of
\eqref{9.2:eq-response-factor-a} is at least $(C_b e^{\delta_1^{\mathrm{sc}}})^{\varpi}$. By
\eqref{9.2:eq-response-factor-a}, $w(X)\le e^{-(\kappa/2)d_j(X)}\,\mathsf{r}_{\varpi}(X)$, where the
response factor $\mathsf{r}_{\varpi}(X)=\rho(y^{\circ})^{\varpi}$ is read at the point $y^{\circ}$.
We take $y=y^{\circ}$ in the
bound just proved and use $(\kappa/2)\,d_j(X)\ge\delta_T\,d_j(X)$, valid because $d_j(X)\ge0$ and
$\kappa/2\ge\delta_T$. This gives
\begin{equation*}
w(X)\le C_w\,e^{-\delta_T(d_j(X)+d_{\mathrm{sc}}(y^{\circ},\Omega))},
\end{equation*}
and the exponent $d_j(X)+d_{\mathrm{sc}}(y^{\circ},\Omega)$ is $E'(X)$. This proves
\eqref{9.2:eq-combined-exponent-a}.

\emph{Position of $y^{\circ}$ and the depths.} The depths satisfy
$j\le\iota^{\circ}\le\iota\le k+1$. In Cases A, B and E below,
$y^{\circ}\in Z_{\iota^{\circ}}\subset Z_k$, so $y^{\circ}\notin\Omega$, and
Lemma~\ref{9.2:lem-response-side} applies at
$y^{\circ}$. In Case D, $y^{\circ}\notin\Omega$ holds by the hypothesis of that case. The radii are
non-negative, so every term $10\check{R}_ic_i^{\circ}$ and every term $(\,\cdot\,)_+$ below is
non-negative and may be dropped from a lower bound. We prove the first inequality of
\eqref{9.2:eq-combined-exponent-b} in five cases.

\emph{Case A: $\iota^{\circ}=\iota\le k$.} Suppose first that $\iota=j$. Then $X\subset Z_j$, and we
use only $d_j(X)\ge0$. Here $\max\{\iota^{\circ},j+1\}=j+1$. Apply
\eqref{9.2:eq-response-side-a} at $y^{\circ}$ and drop its term at level $\iota=j$, which is
non-negative. This gives
\begin{equation*}
E'(X)\ge\sum_{i=j+1}^{k}10\check{R}_ic_i^{\circ}+8L\check{R}_{k+1}
=10\sum_{i\in P}\check{R}_i+8L\check{R}_{k+1},
\end{equation*}
since $P=\{i\in[j+1,k]:c_i^{\circ}=1\}$ in this sub-case. A fortiori the first inequality of
\eqref{9.2:eq-combined-exponent-b} holds.

Suppose next that $\iota\ge j+1$. Put $t:=\operatorname{dist}_\infty(y^{\circ},Z_{\iota-1})$ and
$v:=t/(M\ell_\iota)$. We apply \eqref{9.2:eq-size-side-a} (Lemma~\ref{9.2:lem-size-side} with
$a=0$) at the point
$y^{\circ}\in X$, whose depth $\iota$ lies in $[j+1,k]$. We apply \eqref{9.2:eq-response-side-a}
(Lemma~\ref{9.2:lem-response-side}) at $y^{\circ}$ with this $v$. Adding the two bounds, and using
$t/(M\ell_j)=v\,\ell_\iota/\ell_j=vL^{\iota-j}$, we get
\begin{equation*}
\begin{split}
E'(X)\ge{}&\Bigl[\sum_{i=j+1}^{\iota-1}10\check{R}_iL^{i-j}+vL^{\iota-j}-3\Bigr]\\
&+\Bigl[c_\iota^{\circ}(10\check{R}_\iota-v)_+
+\sum_{i=\iota+1}^{k}10\check{R}_ic_i^{\circ}+8L\check{R}_{k+1}\Bigr].
\end{split}
\end{equation*}
Since $L\ge1$ and $v\ge0$, we have $vL^{\iota-j}\ge v$. Next,
$v+c_\iota^{\circ}(10\check{R}_\iota-v)_+\ge10\check{R}_\iota c_\iota^{\circ}$. If
$c_\iota^{\circ}=0$ this is
$v\ge0$. If $c_\iota^{\circ}=1$ it is $v+(10\check{R}_\iota-v)_+\ge10\check{R}_\iota$, which holds
because $b+(c-b)_+\ge c$ for all real $b$ and $c$. Thus the collar of level $\iota$ that $y^{\circ}$
straddles is accounted for jointly: its inner part by the size side, its outer part by the response
side. Using also $L^{i-j}\ge1$, we obtain
\begin{equation*}
\begin{split}
E'(X)&\ge\sum_{i=j+1}^{\iota-1}10\check{R}_i+\sum_{i=\iota}^{k}10\check{R}_ic_i^{\circ}
+8L\check{R}_{k+1}-3\\
&=10\sum_{i\in P}\check{R}_i+8L\check{R}_{k+1}-3 .
\end{split}
\end{equation*}
The equality holds because here $\max\{\iota^{\circ},j+1\}=\iota$, so that
\begin{equation*}
P=[j+1,\iota-1]\sqcup\{i\in[\iota,k]:c_i^{\circ}=1\}.
\end{equation*}

\emph{Case B: $\iota^{\circ}<\iota\le k$.} Pick $y^{*}\in X$ with $\iota(y^{*})=\iota$; then
$\iota\ge j+1$ because $\iota>\iota^{\circ}\ge j$. We apply \eqref{9.2:eq-size-side-a} at $y^{*}$ and
drop the term in $t$, which is non-negative. With $L^{i-j}\ge1$ this gives
\begin{equation*}
d_j(X)\ge\sum_{i=j+1}^{\iota-1}10\check{R}_i-3 .
\end{equation*}
We apply \eqref{9.2:eq-response-side-a} at $y^{\circ}$ and drop its non-negative terms of levels
$i<\iota$. This gives
\begin{equation*}
d_{\mathrm{sc}}(y^{\circ},\Omega)\ge\sum_{i=\iota}^{k}10\check{R}_ic_i^{\circ}+8L\check{R}_{k+1}.
\end{equation*}
Adding the two bounds,
\begin{equation*}
\begin{split}
E'(X)&\ge\sum_{i=j+1}^{\iota-1}10\check{R}_i+\sum_{i=\iota}^{k}10\check{R}_ic_i^{\circ}
+8L\check{R}_{k+1}-3\\
&\ge10\sum_{i\in P}\check{R}_i+8L\check{R}_{k+1}-3 .
\end{split}
\end{equation*}
The last step holds because
$P\subset[j+1,\iota-1]\cup\{i\in[\iota,k]:c_i^{\circ}=1\}$. Indeed, a level $i\in P$ with
$i\ge\iota$
lies in $[\max\{\iota^{\circ},j+1\},k]$ and not in $U(X)$, so $c_i^{\circ}=1$. The levels of
$[\max\{\iota^{\circ},j+1\},\iota-1]$ with $c_i^{\circ}=0$ are accounted for by the size side but are
counted in $U(X)$. The estimate is therefore not sharp at these levels, and it remains valid.

\emph{Case C: $\iota^{\circ}=\iota=k+1$ and $y^{\circ}\in\Omega$.} The response factor, which is at
most $1$, is not used: we use only $E'(X)\ge d_j(X)$, valid because $d_{\mathrm{sc}}\ge0$. We apply
\eqref{9.2:eq-size-side-a} at $y^{\circ}$ in its clause for $\iota=k+1$ and $y\in\Omega$. With
$L^{i-j}\ge1$ and $L^{s}\ge1$ this gives
\begin{equation*}
E'(X)\ge d_j(X)\ge\sum_{i=j+1}^{k}10\check{R}_i+8L\check{R}_{k+1}-3 .
\end{equation*}
Here $[\max\{\iota^{\circ},j+1\},k]=[k+1,k]$ is empty. Hence $U(X)=\emptyset$ and $P=[j+1,k]$, and
the right member is $10\sum_{i\in P}\check{R}_i+8L\check{R}_{k+1}-3$.

\emph{Case D: $\iota^{\circ}=\iota=k+1$ and $y^{\circ}\notin\Omega$.} Put
$t:=\operatorname{dist}_\infty(y^{\circ},Z_k)$ and write $t=uM\ell_k$, so that
$t/(M\ell_j)=uL^{s}$. Applying \eqref{9.2:eq-size-side-a} at $y^{\circ}$ gives
\begin{equation*}
d_j(X)\ge\sum_{i=j+1}^{k}10\check{R}_iL^{i-j}+uL^{s}-3 .
\end{equation*}
Applying \eqref{9.2:eq-response-side-a} at $y^{\circ}$ gives
$d_{\mathrm{sc}}(y^{\circ},\Omega)\ge(8L\check{R}_{k+1}-u)_+$. Since $L^{s}\ge1$ and $u\ge0$,
\begin{equation*}
uL^{s}+(8L\check{R}_{k+1}-u)_+\ge u+(8L\check{R}_{k+1}-u)_+\ge8L\check{R}_{k+1}.
\end{equation*}
With $L^{i-j}\ge1$ we obtain
\begin{equation*}
E'(X)\ge\sum_{i=j+1}^{k}10\check{R}_i+8L\check{R}_{k+1}-3 .
\end{equation*}
As in Case C, $U(X)=\emptyset$ and $P=[j+1,k]$.

\emph{Case E: $\iota=k+1>\iota^{\circ}$.} Pick $y^{*}\in X\setminus Z_k$ and apply
\eqref{9.2:eq-size-side-a} at $y^{*}$ with the term in $t$ dropped. With $L^{i-j}\ge1$ this gives
\begin{equation*}
d_j(X)\ge\sum_{i=j+1}^{k}10\check{R}_i-3 .
\end{equation*}
Apply \eqref{9.2:eq-response-side-a} at $y^{\circ}$. Keep only the contribution of the final portion
of the path, the one ending in $\Omega$ (the portion written $P_T$ in
Lemma~\ref{9.2:lem-response-side}), and drop the non-negative contributions of the other portions.
This gives
$d_{\mathrm{sc}}(y^{\circ},\Omega)\ge8L\check{R}_{k+1}$. Adding the two bounds,
\begin{equation*}
E'(X)\ge\sum_{i=j+1}^{k}10\check{R}_i+8L\check{R}_{k+1}-3
\ge10\sum_{i\in P}\check{R}_i+8L\check{R}_{k+1}-3,
\end{equation*}
since $P\subset[j+1,k]$ and the radii are non-negative.

\emph{The cases are exhaustive.} Since $j\le\iota^{\circ}\le\iota\le k+1$, there are three
possibilities. If $\iota\le k$, we are in Case A ($\iota^{\circ}=\iota$) or Case B
($\iota^{\circ}<\iota$). If $\iota=k+1$ and $\iota^{\circ}\le k$, we are in Case E. If
$\iota^{\circ}=\iota=k+1$, we are in Case C or Case D, according as $y^{\circ}\in\Omega$ or not. If
$s=0$, that is $k=j$, Case B cannot occur, and in Case A one has $\iota=j$. This proves the first
inequality of \eqref{9.2:eq-combined-exponent-b} for every unit $X$.

\emph{The second inequality of \eqref{9.2:eq-combined-exponent-b}.} By part (a) of
\eqref{9.2:eq-radii-floor-a}, applied repeatedly, $\check{R}_i\ge\check{R}_k=\check{R}^{\flat}$ for
every $i\le k$, in particular for every $i\in P\subset[j+1,k]$. Hence
\begin{equation*}
10\sum_{i\in P}\check{R}_i\ge10\check{R}^{\flat}\,\#P=10\check{R}^{\flat}\bigl(s-\#U(X)\bigr).
\end{equation*}
By part (b) of \eqref{9.2:eq-radii-floor-a} with $i=k$, $\check{R}_k\le L\check{R}_{k+1}$, so
$8L\check{R}_{k+1}\ge8\check{R}_k$. Finally,
$8\check{R}_k-3\ge\tfrac{13}{2}\check{R}_k$ is equivalent to $\tfrac{3}{2}\check{R}_k\ge3$, that is to
$\check{R}_k\ge2$, which is hypothesis (v). Combining the three facts,
\begin{equation*}
10\sum_{i\in P}\check{R}_i+8L\check{R}_{k+1}-3
\ge10\check{R}^{\flat}\bigl(s-\#U(X)\bigr)+\tfrac{13}{2}\check{R}_k .
\end{equation*}

\emph{Consequences.} Insert \eqref{9.2:eq-combined-exponent-b} into
\eqref{9.2:eq-combined-exponent-a}; since $\delta_T>0$, this gives \eqref{9.2:eq-combined-exponent-1}.
No condition on the radii is used beyond $\check{R}_k\ge2$. If $s=0$, the interval $[j+1,k]$ is
empty, so $P=U(X)=\emptyset$. The first inequality of \eqref{9.2:eq-combined-exponent-b} then reads
$E'(X)\ge8L\check{R}_{k+1}-3$, and $8L\check{R}_{k+1}-3\ge8\check{R}_k-3\ge\tfrac{13}{2}\check{R}_k$
by the two facts just used.
\end{proof}

\begin{remark}[Three observations on the combined exponent]
These concern Lemma~\ref{9.2:lem-combined-exponent}.
\begin{enumerate}
\item[(1)] In all five cases the levels below $\iota(X)$ are accounted for by the size side.
Therefore \eqref{9.2:eq-combined-exponent-b} also holds with $U(X)$ replaced by
$U(X)\cap[\iota(X),k]$. We keep $U(X)$ as defined above.
\item[(2)] Suppose the least value of $d_{\mathrm{sc}}(\cdot,\Omega)$ on $X$ is not attained. Then
run the five cases at an arbitrary point $y\in X$ in place of $y^{\circ}$; this produces a set
$U(y)\subset[j+1,k]$. As $y$ ranges over $X$, the right member of
\eqref{9.2:eq-combined-exponent-b} takes only finitely many values, because $U(y)\subset[j+1,k]$.
Choose $y^{\circ}$ to minimise it. Then
\begin{equation*}
\mathsf{r}_{\varpi}(X)\le C_w\,e^{-\delta_T\inf_{y\in X}d_{\mathrm{sc}}(y,\Omega)},
\end{equation*}
and the per-unit form holds with $U(X):=U(y^{\circ})$: the bound
\eqref{9.2:eq-combined-exponent-a} holds, $E'(X)$ is bounded below by the last member of
\eqref{9.2:eq-combined-exponent-b}, and hence \eqref{9.2:eq-combined-exponent-1} holds.
\item[(3)] The coefficient $10$ is the largest per-level coefficient that this argument admits. The
hypotheses are consistent with a configuration with the following properties: the distances satisfy
\begin{equation*}
\operatorname{dist}(Z_{i-1},Z_i^{c})=10M\check{R}_i\ell_i,\qquad
\operatorname{dist}(Z_k,\Omega)=8LM\check{R}_{k+1}\ell_k
\end{equation*}
exactly; all levels are thick; the local-scale condition of
Definition~\ref{9.2:def-scaled-distance} holds with
equality along an $\ell^\infty$-geodesic; and $X$ is a single cube of $\pi_j$ at the boundary of
$Z_j$. In that configuration
\begin{equation*}
E'(X)=\sum_i10\check{R}_i+8L\check{R}_{k+1}+O(1),
\end{equation*}
which \eqref{9.2:eq-combined-exponent-b} reproduces.
\end{enumerate}
\end{remark}

Throughout, $X$ is a unit in the sense of Lemma~\ref{9.2:lem-combined-exponent} (combined size and response exponent), and the attempts, their windows and thin parts are those of Definition~\ref{9.2:def-attempt-windows} (attempt windows; thin and thick levels).

\begin{lemma}[The per-window balance]\label{9.2:lem-window-balance}
Assume the following three statements:
\begin{itemize}
\item the description of the current label by attempt windows, \eqref{9.2:eq-attempt-windows-a}
in Definition~\ref{9.2:def-attempt-windows}, including its inequalities $2N_0(b)\le N_b$ and
$N_b>N_0(b)$ for every attempt $b$ (referred to below as the window model);
\item Lemma~\ref{9.2:lem-thin-levels} (which levels can be thin);
\item the bound \eqref{9.2:eq-combined-exponent-b} of Lemma~\ref{9.2:lem-combined-exponent}.
\end{itemize}
Let $X$ have levels $j\le k$, the number $s=k-j$ of levels of $[j+1,k]$ and the exponent
$E'=E'(X)$ of Lemma~\ref{9.2:lem-combined-exponent}. Let $y^{\circ}$ be its minimising point, of
depth $\iota^{\circ}$. Let $U=U(X)$ and $P=P(X)$ be its sets of thin and of thick levels; each
level of $[j+1,k]$ is either thin or thick for $y^{\circ}$, so that $[j+1,k]$ is the disjoint
union of $U$ and $P$. Let
$\mathfrak{A}(y^{\circ})$ be the set of attempts in force at $y^{\circ}$. Each
$b\in\mathfrak{A}(y^{\circ})$ carries its window $I_b$, its thin part $\Theta_b$, its clock $h_b$,
and the integers $k_b$, $j_b$, $N_b$, $N_0(b)$, $k_0(b)$ of
Definition~\ref{9.2:def-attempt-windows}. Put $m^{\circ}:=\max\{\iota^{\circ},j+1\}$.

An attempt $b\in\mathfrak{A}(y^{\circ})$ whose window meets $[j+1,k]$ is called \emph{contained}
if $h_b\ge j+1$, and \emph{straddling} if $h_b\le j$. A level of $[j+1,k]$ is \emph{free} if it
lies in no window $I_b$, $b\in\mathfrak{A}(y^{\circ})$. An attempt $b\in\mathfrak{A}(y^{\circ})$ is
\emph{exceptional} if $h_b\le j$ and $\Theta_b\cap[m^{\circ},k]\neq\emptyset$.

\smallskip\noindent\textup{(a)} The following hold; in (i)--(iii), (v), (vi) $b$ ranges over
$\mathfrak{A}(y^{\circ})$, and in (v) $a^{*}$ denotes the exceptional attempt when there is one.
\begin{equation}\label{9.2:eq-window-balance-b}
\begin{aligned}
&\text{(i)}\ \ \Theta_b\subset I_b\ \text{ and }\ \max\Theta_b\le k_b-1;\\
&\text{(ii)}\ \ I_b\cap U\subset\Theta_b;\\
&\phantom{\text{(ii)}}\ \ \text{every free level is not thin for }
  y^{\circ}\text{ and lies in }P;\\
&\phantom{\text{(ii)}}\ \ \#(U\cap I_b)\le\#\Theta_b\le N_0(b)-1\le \tfrac12 N_b-1;\\
&\text{(iii)}\ \ k\notin I_b;\ \text{if } b \text{ is contained, } I_b\subset[j+1,k-1];\\
&\phantom{\text{(iii)}}\ \ \text{if } b\text{ is straddling, } I_b\supseteq\{j,j+1\};\\
&\phantom{\text{(iii)}}\ \ \text{at most one attempt in }\mathfrak{A}(y^{\circ})
  \text{ is straddling};\\
&\text{(iv)}\ \ \text{at most one attempt in }\mathfrak{A}(y^{\circ})
  \text{ is exceptional,}\\
&\phantom{\text{(iv)}}\ \ \text{and an exceptional attempt is straddling};\\
&\text{(v)}\ \ \text{if } a^{*}\text{ exists, put } k_2:=k_{a^{*}};\
  \text{then } k_2\le k;\\
&\phantom{\text{(v)}}\ \ \text{if moreover } b\neq a^{*}\text{ and } I_b\cap[j+1,k]\neq\emptyset,\\
&\phantom{\text{(v)}}\ \ \text{then } h_b\ge k_2 \text{ and } I_b\subset[k_2,k-1];\\
&\text{(vi)}\ \ \text{if no attempt is exceptional and } b \text{ is straddling,}\\
&\phantom{\text{(vi)}}\ \ \text{then } I_b\cap[j+1,k]\subset P.
\end{aligned}
\end{equation}

\smallskip\noindent\textup{(b)} Fix $a\in\{10,5,4\}$. For $i\in P$ put
$\sigma_a(i):=10/a-1$, and for $i\in U$ put $\sigma_a(i):=-1$. For $S\subset P\cup U$ put
$\sigma_a(S):=\sum_{i\in S}\sigma_a(i)$. Assume the condition at parameter $a$ of
\eqref{9.2:eq-pile-bound-a}, which gives $10\,\delta_T\check{R}^{\flat}\ge(10/a)\Lambda_L$.
Then for every real number $c$
\begin{equation}\label{9.2:eq-window-balance-a}
\delta_T E'-s\Lambda_L-c\,\delta_T\check{R}_k\ \ge\ \sigma_a([j+1,k])\,\Lambda_L
+\Bigl(\frac{13}{2}-c\Bigr)\delta_T\check{R}_k .
\end{equation}

\smallskip\noindent\textup{(c)} With $\sigma_a$ as in \textup{(b)}, every
$b\in\mathfrak{A}(y^{\circ})$ with $I_b\subset[j+1,k]$ satisfies
\begin{equation}\label{9.2:eq-window-balance-c}
\sigma_a(I_b)\ \ge\ \Bigl(\frac{5}{a}-1\Bigr)N_b+\frac{10}{a}.
\end{equation}
At $a=10$, $5$ and $4$ this lower bound equals $1-N_b/2$, $2$ and $N_b/4+5/2$ respectively.
\end{lemma}

\begin{proof}
\emph{Part (a), item (i).} The lower end of the window $I_b$ is compared with the lower end of the
thin part $\Theta_b$ through the inequality $k_0(b)+1\ge h_b+1$. This inequality holds because
$N_0(b)\le N_b$, which in turn follows from $N_b>N_0(b)$ in the window model
\eqref{9.2:eq-attempt-windows-a}. With it, the window model gives $\Theta_b\subset I_b$, and no
element of $\Theta_b$ exceeds $k_b-1$.

\emph{Item (ii).} Let $i\in I_b$ be a level of $U$, that is, a level of $[j+1,k]$ that is thin for
$y^{\circ}$. By Lemma~\ref{9.2:lem-thin-levels}, $i\in\Theta_c$ for some
$c\in\mathfrak{A}(y^{\circ})$, and $\Theta_c\subset I_c$ by item (i). Thus $i\in I_b\cap I_c$, and
the window model \eqref{9.2:eq-attempt-windows-a} gives $c=b$. Hence $i\in\Theta_b$.

The same argument applies to any level of $[j+1,k]$ that is thin for $y^{\circ}$: such a level
lies in some $\Theta_c\subset I_c$, so it is not free. Consequently a free level is not thin, and
so it lies in $P$.

Since $I_b\cap U\subset\Theta_b$, we get $\#(U\cap I_b)\le\#\Theta_b$. The window model
\eqref{9.2:eq-attempt-windows-a} gives $\#\Theta_b\le N_0(b)-1$ and $2N_0(b)\le N_b$. These
together give the displayed chain.

\emph{Item (iii).} By the window model \eqref{9.2:eq-attempt-windows-a}, level $k$ lies in no
window $I_b$, $b\in\mathfrak{A}(y^{\circ})$. The window $I_b$ is an interval of levels whose lowest
level is $h_b$ (Definition~\ref{9.2:def-attempt-windows}).

Suppose $b$ is contained. Then $I_b$ begins at $h_b\ge j+1$, meets $[j+1,k]$ and does not contain
$k$. Hence $I_b\subset[j+1,k-1]$.

Suppose $b$ is straddling. Then $I_b$ begins at $h_b\le j$ and contains a level $\ge j+1$. Hence
$I_b\supseteq\{j,j+1\}$.

Two straddling attempts would have windows sharing the level $j$. The windows of distinct
attempts in $\mathfrak{A}(y^{\circ})$ are disjoint by the window model
\eqref{9.2:eq-attempt-windows-a}, so at most one attempt is straddling.

\emph{Item (iv).} Let $b$ be exceptional, and pick a level of $\Theta_b$ in $[m^{\circ},k]$. By
item (i) this level lies in $I_b$ and is at most $k_b-1$. Therefore
\begin{equation*}
k_b-1\ \ge\ m^{\circ}\ \ge\ j+1 .
\end{equation*}
In particular $I_b$ meets $[j+1,k]$, and $h_b\le j$. So $b$ is straddling, and
$I_b\supseteq\{j,j+1\}$ by item (iii). By item (iii) at most one attempt is straddling, hence at
most one is exceptional. This argument uses no information on the lengths of the windows.

\emph{Item (v).} By item (iv), $a^{*}$ is straddling, so by item (iii) its window $I_{a^{*}}$ is an
interval that contains $j+1\le k$ and does not contain $k$; hence all its levels are below $k$.
Its largest level is $k_2-1$, so $k_2-1\le k-1$, that is, $k_2\le k$.

The windows $I_b$ and $I_{a^{*}}$ are disjoint intervals by the window model
\eqref{9.2:eq-attempt-windows-a}, so one lies entirely
below the other. Suppose $I_b$ lay below $I_{a^{*}}$. Then $k_b-1<h_{a^{*}}\le j$, and $I_b$ could
not meet $[j+1,k]$, contrary to hypothesis. Therefore $I_b$ lies above $I_{a^{*}}$, whose largest
level is $k_2-1$. So $h_b>k_2-1$, that is, $h_b\ge k_2$.

By the proof of item (iv), $k_2-1\ge j+1$, so $h_b\ge k_2>j+1$ and $b$ is contained. Item (iii)
then gives $I_b\subset[j+1,k-1]$. Combined with $h_b\ge k_2$, this yields
$I_b\subset[k_2,k-1]$.

\emph{Item (vi).} By Definition~\ref{9.2:def-attempt-windows}, thinness of a level $i$ for
$y^{\circ}$ is read through the component $A_i(y^{\circ})$ of $Z_i$ containing $y^{\circ}$, as in
the hypothesis $y\in Z_i$ of Lemma~\ref{9.2:lem-thin-levels}; so a level that is thin for
$y^{\circ}$ is at least the depth $\iota^{\circ}$ of $y^{\circ}$. Hence $U\subset[m^{\circ},k]$, and
a level of $I_b\cap[j+1,k]$ that lies in $U$ lies in $I_b\cap[m^{\circ},k]$.

Assume that no attempt is exceptional, and let $b$ be straddling. Suppose a
level $i\in I_b\cap[m^{\circ},k]$ were in $U$. By item (ii), $i$ would lie in $\Theta_b$. Then
$\Theta_b\cap[m^{\circ},k]\neq\emptyset$, and since $h_b\le j$, the attempt $b$ would be
exceptional, which is excluded. Hence no level of $I_b\cap[j+1,k]$ lies in $U$, and since
$[j+1,k]$ is the disjoint union of $P$ and $U$, we get $I_b\cap[j+1,k]\subset P$.

\emph{Part (b).} Since $[j+1,k]$ is the disjoint union of $P$ and $U$, we have
$s=\#P+\#U$. By \eqref{9.2:eq-combined-exponent-b},
\begin{equation*}
E'\ \ge\ 10\,\check{R}^{\flat}\,(s-\#U)+\tfrac{13}{2}\check{R}_k
\ =\ 10\,\check{R}^{\flat}\,\#P+\tfrac{13}{2}\check{R}_k .
\end{equation*}
We multiply by $\delta_T>0$ and use the inequality $10\,\delta_T\check{R}^{\flat}\ge(10/a)\Lambda_L$
given by the assumed condition of \eqref{9.2:eq-pile-bound-a}. This gives
\begin{equation*}
\delta_T E'\ \ge\ 10\,\delta_T\check{R}^{\flat}\,\#P+\tfrac{13}{2}\delta_T\check{R}_k
\ \ge\ \frac{10}{a}\Lambda_L\,\#P+\tfrac{13}{2}\delta_T\check{R}_k .
\end{equation*}
Next we subtract $s\Lambda_L=(\#P+\#U)\Lambda_L$ and $c\,\delta_T\check{R}_k$. The left side of
\eqref{9.2:eq-window-balance-a} is therefore at least
\begin{equation*}
\Bigl(\bigl(\tfrac{10}{a}-1\bigr)\#P-\#U\Bigr)\Lambda_L+\Bigl(\tfrac{13}{2}-c\Bigr)\delta_T\check{R}_k
=\sigma_a([j+1,k])\,\Lambda_L+\Bigl(\tfrac{13}{2}-c\Bigr)\delta_T\check{R}_k ,
\end{equation*}
by the definition of $\sigma_a$.

\emph{Part (c).} Let $I_b\subset[j+1,k]$, and put $t:=\#(U\cap I_b)$. By item (ii) of
\eqref{9.2:eq-window-balance-b}, $t\le N_b/2-1$. The window $I_b$ consists of $N_b$ levels, each in
$P$ or in $U$, so exactly $N_b-t$ of them lie in $P$. By the definition of $\sigma_a$,
\begin{equation*}
\sigma_a(I_b)=\Bigl(\frac{10}{a}-1\Bigr)(N_b-t)-t=\Bigl(\frac{10}{a}-1\Bigr)N_b-\frac{10}{a}\,t .
\end{equation*}
Since $10/a>0$, we may insert the upper bound on $t$:
\begin{equation*}
\sigma_a(I_b)\ \ge\ \Bigl(\frac{10}{a}-1\Bigr)N_b-\frac{10}{a}\Bigl(\frac{N_b}{2}-1\Bigr)
=\Bigl(\frac{5}{a}-1\Bigr)N_b+\frac{10}{a} .
\end{equation*}
Substituting $a=10$, $5$, $4$ gives $1-N_b/2$, $2$ and $N_b/4+5/2$.

We record how \textup{(c)} changes under a weaker hypothesis.
Suppose the inequality $2N_0(b)\le N_b$
in \eqref{9.2:eq-attempt-windows-a} is replaced by the weaker
$2N_0(b)\le N_b+\varepsilon$. Then the chain in item (ii) becomes
\begin{equation*}
t\ \le\ N_0(b)-1\ \le\ \tfrac12(N_b+\varepsilon)-1,
\end{equation*}
and the computation above gives
\begin{equation*}
\sigma_a(I_b)\ \ge\ \Bigl(\frac{5}{a}-1\Bigr)N_b+\frac{10-5\varepsilon}{a}.
\end{equation*}
At $a=5$ this lower bound is $2-\varepsilon$, which is $\ge0$ for $\varepsilon\le2$. So each window
with $I_b\subset[j+1,k]$ contributes a non-negative amount, and so does any sum of such
contributions. At $a=4$ the lower bound is at least $N_b/4$ for $\varepsilon\le2$. Consequently
the two counting estimates of this section that use \textup{(c)} at $a=5$ and at $a=4$ keep their
conclusions for $\varepsilon\le2$.
\end{proof}

\begin{proof}[Proof of the per-unit bound in three grades]
\label{9.2:prf-per-unit-grades}
We use the notation of Theorem~\ref{9.2:thm-pile-bound}:
\begin{equation*}
s=k-j,\qquad \Lambda_L=\log(2L^4),\qquad
\delta_T=\min\{\varpi\delta_1^{\mathrm{sc}},\kappa/2\},\qquad
\check{R}^{\flat}=\check{R}_k .
\end{equation*}
The sets $U=U(X)$ and $P=P(X)$ are the thin and the thick levels of the unit $X$. Every level of
$[j+1,k]$ lies in exactly one of them. For $a\in\{10,5,4\}$ and $I\subset[j+1,k]$ we write
$\sigma_a(I)=\sum_{i\in I}\sigma_a(i)$, where $\sigma_a(i)=10/a-1$ for $i\in P$ and $\sigma_a(i)=-1$
for $i\in U$.

The grades are the following three cases:
\begin{itemize}
\item[(\(\alpha\))] $U=\emptyset$, under the floor at index $10$ of \eqref{9.2:eq-pile-bound-a};
\item[(\(\beta\))] there is no exceptional window, under the floor at index $5$;
\item[(\(\gamma\))] the general case, under the floor at index $4$ together with the floor on
$N_0^{\sharp}$ and the two conditions on the ray $y\ge\gamma^{-2}-c_5$, all collected in
\eqref{9.2:eq-pile-bound-a}.
\end{itemize}
We prove that in each of these cases
\begin{equation}\label{9.2:eq-per-unit-grades-1}
\delta_T E'(X)\ \ge\ s\Lambda_L+c\,\delta_T\check{R}_k,
\qquad c=\tfrac{13}{2},\ \tfrac{13}{2},\ \tfrac{3}{2}
\text{ in grades }(\alpha),\ (\beta),\ (\gamma).
\end{equation}
By the combined size and response exponent (Lemma~\ref{9.2:lem-combined-exponent}, display
\eqref{9.2:eq-combined-exponent-a}),
$w(X)\le C_we^{-\delta_TE'(X)}$. Since $e^{-s\Lambda_L}=(2L^4)^{-s}$, the inequality
\eqref{9.2:eq-per-unit-grades-1} is the bound
\begin{equation}\label{9.2:eq-per-unit-grades-2}
w(X)\ \le\ C_w\,e^{-c\,\delta_T\check{R}_k}\,(2L^4)^{-s}.
\end{equation}

Grades $(\alpha)$ and $(\beta)$, and the first branch of grade $(\gamma)$ below, use only the
following inputs: Lemma~\ref{9.2:lem-combined-exponent}, the per-window balance
(Lemma~\ref{9.2:lem-window-balance}), the statements \eqref{9.2:eq-window-balance-a} and
\eqref{9.2:eq-window-balance-b}, and the floor at the relevant index. The second branch of grade
$(\gamma)$ uses, in addition, the following statements of this section:
\begin{itemize}
\item the upper member of the reference bound \eqref{9.2:eq-attempt-windows-c} of
Definition~\ref{9.2:def-attempt-windows};
\item the bound $N_0(a)\le N_0^{\sharp}(x_{k_a})$ on the thin count of an attempt, which is part of
\eqref{9.2:eq-attempt-windows-a};
\item the floor bounding $N_0^{\sharp}(2x_k)\Lambda_L$ by $5\delta_T\check{R}_k$, the inequality
$4B^{\sharp}N_0^{\sharp}(y)\le\tfrac12 y$, and the condition giving $\log((1+B^{\sharp})y)\ge1$,
the last two for $y\ge\gamma^{-2}-c_5$, all in \eqref{9.2:eq-pile-bound-a}.
\end{itemize}

By the per-window balance \eqref{9.2:eq-window-balance-a} under the floor at index $a$ of
\eqref{9.2:eq-pile-bound-a}, the inequality \eqref{9.2:eq-per-unit-grades-1} at $c=13/2$ follows
once $\sigma_a([j+1,k])\ge0$. Hence, at $c=13/2$, it suffices to show $\sigma_a([j+1,k])\ge0$.

\emph{Grade $(\alpha)$.} Here $U=\emptyset$, so $P=[j+1,k]$ and $\sigma_{10}(i)=10/10-1=0$ for
every level $i$. Thus $\sigma_{10}\equiv0$ on $P=[j+1,k]$, and \eqref{9.2:eq-per-unit-grades-1}
holds with $c=13/2$. Directly from Lemma~\ref{9.2:lem-combined-exponent} with $\#U(X)=0$, and
since the floor at index $10$ of \eqref{9.2:eq-pile-bound-a} bounds
$e^{-10\delta_T\check{R}^{\flat}s}$ by $(2L^4)^{-s}$, we obtain
\begin{equation*}
w(X)\ \le\ C_we^{-\frac{13}{2}\delta_T\check{R}_k}\,e^{-10\delta_T\check{R}^{\flat}s}
\ \le\ C_we^{-\frac{13}{2}\delta_T\check{R}_k}\,(2L^4)^{-s}.
\end{equation*}

\emph{Grade $(\beta)$.} Assume there is no exceptional window. Then $[j+1,k]$ is the disjoint
union of three sets:
\begin{itemize}
\item the set $F$ of free levels;
\item the set $S$ of levels, inside $[j+1,k]$, of the straddling window, of which there is at most
one;
\item the contained windows $I_b$.
\end{itemize}
By \eqref{9.2:eq-window-balance-b}, $\sigma_5(F)=\#F\ge0$ and $\sigma_5(S)=\#S\ge0$. By
Lemma~\ref{9.2:lem-window-balance} at $a=5$, $\sigma_5(I_b)\ge2$ for every contained window.
Summing, $\sigma_5([j+1,k])\ge0$, which is \eqref{9.2:eq-per-unit-grades-1} with $c=13/2$.

\emph{Grade $(\gamma)$, no exceptional window.} Since $3/2\ge1$ on $P$ and both weights equal
$-1$ on $U$, we have $\sigma_4\ge\sigma_5$ level by level. The decomposition of grade $(\beta)$
therefore gives $\sigma_4([j+1,k])\ge\sigma_5([j+1,k])\ge0$. So
\eqref{9.2:eq-per-unit-grades-1} holds even with $c=13/2$.

\emph{Grade $(\gamma)$, the exceptional window.} Otherwise let $a^{*}$ be the exceptional window of
\eqref{9.2:eq-window-balance-b}, and put
\begin{equation*}
k_2:=k_{a^{*}}\le k,\qquad u:=k-k_2\ge0,\qquad
d^{*}:=\#(U\cap I_{a^{*}})=\#(U\cap\Theta_{a^{*}})\le N_0(a^{*})-1 .
\end{equation*}
Then $[j+1,k]=[j+1,k_2-1]\sqcup[k_2,k]$, with $[j+1,k_2-1]\subset I_{a^{*}}$ and
$\#[k_2,k]=u+1$.

On $[j+1,k_2-1]$ the $P$-levels have weight $3/2\ge0$. The $U$-levels have weight $-1$, and there
are $d^{*}$ of them. Hence $\sigma_4([j+1,k_2-1])\ge-d^{*}$.

On $[k_2,k]$ each level is of one of two kinds:
\begin{itemize}
\item it is free, with weight $3/2\ge\tfrac14$; or
\item it lies in a contained window $I_b\subset[k_2,k-1]$, by \eqref{9.2:eq-window-balance-b}.
For such a window $\sigma_4(I_b)\ge N_b/4$ by Lemma~\ref{9.2:lem-window-balance} at $a=4$.
\end{itemize}
Thus $\sigma_4([k_2,k])\ge\tfrac14(u+1)$. Inserting both estimates into the per-window balance
\eqref{9.2:eq-window-balance-a} under the floor at index $4$ of \eqref{9.2:eq-pile-bound-a}, and
moving the term $c\,\delta_T\check{R}_k$, we get for every $c$
\begin{equation}\label{9.2:eq-per-unit-grades-a}
\delta_TE'-s\Lambda_L-c\,\delta_T\check{R}_k\ \ge\
\bigl[\tfrac14(u+1)-d^{*}\bigr]\Lambda_L+\bigl(\tfrac{13}{2}-c\bigr)\delta_T\check{R}_k .
\end{equation}

\emph{Branch 1: $\tfrac14(u+1)\ge d^{*}$.} Both terms on the right of
\eqref{9.2:eq-per-unit-grades-a} are non-negative for every $c\le13/2$, whatever the value of $u$.

\emph{Branch 2: $u+1<4d^{*}$.} Since $u+1\ge1$, this forces $d^{*}\ge1$. Since $u$ and $d^{*}$ are
integers, it also gives $u\le4d^{*}-2$. We proceed in six steps.

Step 1. We show $x_{k_2}-x_k\le B^{\sharp}u$.
\begin{itemize}
\item If $u=0$, both sides vanish.
\item If $u\ge1$, the upper member of \eqref{9.2:eq-attempt-windows-c} gives
$x_{k_2}-x_k\le3\lambda u+2c_2$. With $B^{\sharp}=3\lambda+2c_2$ we have
$B^{\sharp}u-(3\lambda u+2c_2)=2c_2(u-1)\ge0$, because $c_2\ge0$.
\end{itemize}

Step 2. From $u\le4d^{*}$, from $d^{*}\le N_0(a^{*})-1$, and from the bound
$N_0(a^{*})\le N_0^{\sharp}(x_{k_{a^{*}}})$ of \eqref{9.2:eq-attempt-windows-a}, we get
\begin{equation*}
B^{\sharp}u\ \le\ 4B^{\sharp}d^{*}\ \le\ 4B^{\sharp}\bigl(N_0(a^{*})-1\bigr)
\ \le\ 4B^{\sharp}N_0^{\sharp}(x_{k_2}).
\end{equation*}
Here we used that the coupling of the attempt $a^{*}$ is $x_{k_{a^{*}}}=x_{k_2}$.

Step 3. We have $x_{k_2}=g_{k_2}^{-2}\ge\gamma^{-2}$, hence $x_{k_2}\ge\gamma^{-2}-c_5$. So the
inequality
$4B^{\sharp}N_0^{\sharp}(y)\le\tfrac12y$ of \eqref{9.2:eq-pile-bound-a} applies at $y=x_{k_2}$ and
gives $4B^{\sharp}N_0^{\sharp}(x_{k_2})\le\tfrac12x_{k_2}$. With Steps 1 and 2, this yields
$x_{k_2}-x_k\le\tfrac12x_{k_2}$, that is, $x_{k_2}\le2x_k$.

Step 4. By the condition in \eqref{9.2:eq-pile-bound-a}, $\log((1+B^{\sharp})y)\ge1>0$ for
$y\ge\gamma^{-2}-c_5$. Hence
$N_0^{\sharp}(y)=2+r_{\mathrm{pr}}\log_L\log((1+B^{\sharp})y)$ is defined on that ray. It is
non-decreasing there, being a composition of non-decreasing maps. Both $x_{k_2}$ and $2x_k$ lie on
the ray: the first by Step 3, the second because $2x_k\ge x_{k_2}$. Therefore
$N_0^{\sharp}(x_{k_2})\le N_0^{\sharp}(2x_k)$.

Step 5. Combining $d^{*}\le N_0(a^{*})-1\le N_0^{\sharp}(x_{k_2})-1$ with Step 4 gives
$d^{*}\le N_0^{\sharp}(2x_k)-1$. The floor on $N_0^{\sharp}$ in \eqref{9.2:eq-pile-bound-a} then gives
\begin{equation*}
(d^{*}+1)\Lambda_L\ \le\ N_0^{\sharp}(2x_k)\Lambda_L\ \le\ 5\delta_T\check{R}_k .
\end{equation*}

Step 6. At $c=3/2$, by Step 5,
\begin{equation*}
\bigl(\tfrac{13}{2}-c\bigr)\delta_T\check{R}_k-d^{*}\Lambda_L
=5\delta_T\check{R}_k-d^{*}\Lambda_L\ \ge\ \Lambda_L\ \ge\ 0 .
\end{equation*}
Since $\tfrac14(u+1)\Lambda_L\ge0$, the right-hand side of \eqref{9.2:eq-per-unit-grades-a} is
non-negative. This is \eqref{9.2:eq-per-unit-grades-1} with $c=3/2$.

When $u=0$ we have $k_2=k$, so $x_{k_2}=x_k$. The upper member of
\eqref{9.2:eq-attempt-windows-c} and the inequality $4B^{\sharp}N_0^{\sharp}(y)\le\tfrac12y$ are then
not used.

The age $u$ of the exceptional window is split between the two branches and is never bounded. When
$u$ is large, it is paid for by $\tfrac14\Lambda_L$ per level above $I_{a^{*}}$. When $u$ is small,
it is converted into $x_{k_2}\le2x_k$. No property of $N_b$ beyond $2N_0\le N_b$ enters, and no
bound on the number of windows or on a nesting depth is used.

In each grade we have proved $\delta_TE'\ge s\Lambda_L+c\,\delta_T\check{R}_k$, that is,
\eqref{9.2:eq-per-unit-grades-2}, with $c=13/2$, $13/2$, $3/2$ in grades $(\alpha)$, $(\beta)$,
$(\gamma)$. This is the per-unit bound (i) of \eqref{9.2:eq-pile-bound-b}.
\end{proof}

\begin{lemma}[Robustness of the counting to window placement]\label{9.2:lem-window-robustness}
The counting in the proof~\ref{9.2:prf-per-unit-grades} of the per-unit bound uses a window $I_b$
only through the following properties.
\begin{itemize}
\item[(a)] $I_b$ is an integer interval of $N_b\ge 2$ consecutive levels.
\item[(b)] $\max I_b\le k-1$.
\item[(c)] The levels thin for $y^{\circ}$ that are attributable to $b$ lie in a set $\Theta_b\subset I_b$
with
\begin{equation*}
\#\Theta_b\le N_0(b)-1\le \tfrac12 N_b-1 .
\end{equation*}
These levels lie among the top $N_0(b)$ levels of $I_b$; their position there is not used.
\item[(d)] The windows are pairwise disjoint.
\item[(e)] Every level thin for $y^{\circ}$ lies in some $\Theta_b$, and level $k$ is never thin
(the statement on which levels can be thin, \eqref{9.2:eq-thin-levels-a}).
\item[(f)] For the exceptional window $a^{*}$ only, and in branch~2 of grade $(\gamma)$ only: the
coupling index is $k_{a^{*}}=\max I_{a^{*}}+1$, and $N_0(a^{*})\le N_0^{\sharp}(x_{k_{a^{*}}})$.
\end{itemize}
Let the windows $I_b=[h_b,k_b-1]$ be placed one index up,
$I'_b:=[h_b+1,k_b]$. Then $k_b\notin\Theta_b$, and level $k$, if it lies in $I'_b$, is thick,
hence in $P$. The displays \eqref{9.2:eq-window-balance-b} hold verbatim. The per-window balance of
grade $(\beta)$ satisfies
\begin{equation}\label{9.2:eq-window-robustness-1}
N_b-2\,\#(U\cap I'_b)\;\ge\; N_b-2N_0(b)+2\;\ge\;2 .
\end{equation}
In grade $(\gamma)$ the levels above $I'_{a^{*}}$ number $u':=k-k_{a^{*}}$ and their surplus is at
least $\tfrac14 u'\Lambda_L$; branch~1 reads $\tfrac14 u'\ge d^{*}$, and branch~2 reads
$u'\le 4d^{*}-1$ and
\begin{equation}\label{9.2:eq-window-robustness-2}
x_{k_{a^{*}}}-x_k\;\le\; B^{\sharp}u'\;\le\; 4B^{\sharp}d^{*},
\end{equation}
with $x_{k_{a^{*}}}=x_k$ when $u'=0$. The only term lost in passing from $I_b$ to $I'_b$ is the
summand $1$ in the factor $\tfrac14(u+1)$ of grade $(\gamma)$, which that grade does not use.
\end{lemma}

\begin{proof}
Throughout, $U$ and $P$ are the thin and thick level sets $U(X)$ and $P(X)$, and
$\Lambda_L=\log(2L^4)$.

\emph{Properties (a)--(f) for $I'_b$.} We have $I'_b=I_b+1$. Hence $I'_b$ is again an integer
interval of $N_b$ consecutive levels, which is (a); and since every window is translated by the same
integer, pairwise disjointness is preserved, which is (d). By (b), $k_b-1=\max I_b\le k-1$, so
$\max I'_b=k_b\le k$; the only level of $I'_b$ that (b) would exclude is $k$, and it lies in $I'_b$
exactly when $k_b=k$. By (e), level $k$ is never thin; so if it lies in $I'_b$ it is thick, that is,
it belongs to $P$. This replaces (b). Since $\Theta_b\subset I_b$ and $\max I_b=k_b-1$, we have
$k_b\notin\Theta_b$. By (c) the thin levels attributable to $b$ lie among the top $N_0(b)$ levels
of $I_b$, that is, in $[k_b-N_0(b),\,k_b-1]$. Since $N_0(b)\le \tfrac12 N_b<N_b=k_b-h_b$, we have
$k_b-N_0(b)\ge h_b+1$, so these levels lie in $I'_b$, and (c) holds for $I'_b$ with the same set
$\Theta_b$ and the same bound on $\#\Theta_b$. Property (e) does not involve the placement of the
windows and is unchanged. In (f) the coupling index $k_{a^{*}}$ is now $\max I'_{a^{*}}$; the
inequality $N_0(a^{*})\le N_0^{\sharp}(x_{k_{a^{*}}})$ reads the same index $k_{a^{*}}$ and is unchanged.

\emph{The displays \eqref{9.2:eq-window-balance-b}.} Each of them reads a window only through its
position relative to the two levels $\{j,j+1\}$; they therefore hold verbatim for the windows $I'_b$.

\emph{Grade $(\beta)$.} The per-window balance of grade $(\beta)$ on $I'_b$ is at least
$N_b-2\,\#(U\cap I'_b)$. A thin level in $I'_b$ lies, by (e), in some $\Theta_c$, hence by the
preceding paragraph in $I'_c$; by (d), $c=b$. Thus $U\cap I'_b\subset\Theta_b$, and by (c),
$\#(U\cap I'_b)\le N_0(b)-1$. This gives the first inequality of \eqref{9.2:eq-window-robustness-1};
the second is $N_0(b)\le \tfrac12 N_b$ from (c). The resulting lower bound $2$ is the value that the
per-window balance (Lemma~\ref{9.2:lem-window-balance}) gives at $a=5$.

\emph{Grade $(\gamma)$.} The levels above $I'_{a^{*}}$ are $k_{a^{*}}+1,\dots,k$; they number
$u'=k-k_{a^{*}}$, and the count of grade $(\gamma)$ gives them a surplus of at least
$\tfrac14 u'\Lambda_L$. Branch~1 is then the case $\tfrac14 u'\ge d^{*}$. In branch~2 this fails;
since $u'$ and $d^{*}$ are integers, this means $u'\le 4d^{*}-1$. The bound
$x_{k_{a^{*}}}-x_k\le B^{\sharp}u'$ is obtained verbatim as in branch~2 of grade $(\gamma)$, now with
the coupling index $k_{a^{*}}=\max I'_{a^{*}}$ of (f). The second inequality of
\eqref{9.2:eq-window-robustness-2} follows from $u'\le 4d^{*}-1$. If $u'=0$, then $k_{a^{*}}=k$ and
$x_{k_{a^{*}}}=x_k$.

\emph{The lost term.} For $I_b=[h_b,k_b-1]$ the levels above $I_{a^{*}}$ are $k_{a^{*}},\dots,k$.
Their surplus enters grade $(\gamma)$ as $\tfrac14(u+1)\Lambda_L$, and for $I'_{a^{*}}$ it enters as
$\tfrac14 u'\Lambda_L$. Every other step above holds verbatim. So the summand $1$ in $\tfrac14(u+1)$ is
the only term lost. Grade $(\gamma)$ does not use it, as the two branches above show.
\end{proof}

\begin{proof}[Proof of clause~(ii) of Theorem~\ref{9.2:thm-pile-bound}: the pile per box]
\label{9.2:prf-pile-per-box}
Clause~(ii) is proved as an implication from three inputs, which we use and nothing else.
\begin{itemize}
\item[(a)] Clause~(i) of Theorem~\ref{9.2:thm-pile-bound}, in the following form: for every domain
$X\in\mathbf{D}_j$ with $X\cap\Delta\neq\emptyset$ there is a grade $c$ such that
\begin{equation*}
w(X)\le C_w\,e^{-c\delta_T\check{R}_k}\,(2L^4)^{-s}.
\end{equation*}
Here $w(X)$, $C_w$ and $\delta_T$ are the quantities of \eqref{9.2:eq-response-factor-a}, and $s=k-j$.
\item[(b)] The size bound \eqref{9.2:eq-boundary-families-b} for level-$j$ boundary families. For every
$\square\in\pi_j$ it bounds $\sum_{X\ni\square}e^{-\kappa d_j(X)/2}$ by $K_0^{(2)}$.
\item[(c)] The block structure of the compatible partitions, as in \cite{B-CMP119}. The partition $\pi_j$
into closed $M\ell_j$-cubes is a product partition: its cubes are the products of four closed arcs
$[n\ell,(n+1)\ell]$, one in each coordinate, where $\ell:=M\ell_j$.
\end{itemize}

Let $\Delta$ be a closed axis-parallel box of side $\sigma M\ell_k$, with $\sigma>0$. Since
$\ell_k=L^{s}\ell_j$, this side equals $\sigma L^sM\ell_j$. Let $F$ be a level-$j$ boundary family
(Definition~\ref{9.2:def-boundary-families}). The pile $\Pi_j(\Delta)$ of \eqref{9.2:eq-response-factor-a}
is the sum of the nonnegative terms $|F^{(j)}(X)|\,\mathsf{r}_{\varpi}(X)$ over the domains
$X\in\mathbf{D}_j$ that meet $\Delta$.

\emph{Step 1.} Every $X$ meeting $\Delta$ contains a closed $\pi_j$-cube meeting $\Delta$. Indeed, by
Definition~\ref{9.2:def-boundary-families}, $X$ is a connected union of closed $\pi_j$-cubes. A point of
$X\cap\Delta$ therefore lies in one of these cubes, and that cube meets $\Delta$.

\emph{Step 2.} We show that
\begin{equation*}
\#\{\square\in\pi_j:\ \square\cap\Delta\neq\emptyset\}\le(\sigma L^s+2)^4\le\bigl((\sigma+2)L^s\bigr)^4.
\end{equation*}
First fix one coordinate and put $a:=\sigma L^s$. The projection of $\Delta$ onto that coordinate is a
closed arc of length $a\ell$; lift it to $[t,t+a\ell]$. The arc $[n\ell,(n+1)\ell]$ meets it if and only if
$n\ell\le t+a\ell$ and $(n+1)\ell\ge t$. Equivalently, $n$ lies in the real interval $[t/\ell-1,\ t/\ell+a]$,
whose length is $a+1$. Such an interval holds at most $a+2$ integers. If $a\ell$ exceeds the circumference,
every arc is met; the number of arcs is then at most $a$, because the circumference is smaller than $a\ell$.
In either case at most $a+2$ of the arcs meet the projection.

Now use (c). A product of closed arcs meets the product $\Delta$ if and only if, in each coordinate, its arc
meets the projection of $\Delta$. Hence at most $(\sigma L^s+2)^4$ cubes of $\pi_j$ meet $\Delta$. Finally
$\sigma L^s+2\le(\sigma+2)L^s$, because $2\le 2L^s$.

Without using the product structure, the same count follows from volumes. The cubes of $\pi_j$ meeting
$\Delta$ lie in a box of side $\sigma M\ell_k+2M\ell_j$, and their interiors are disjoint. Each of them has
volume $(M\ell_j)^4$. Their number is therefore at most
\begin{equation*}
\Bigl(\frac{\sigma M\ell_k+2M\ell_j}{M\ell_j}\Bigr)^4=(\sigma L^s+2)^4 .
\end{equation*}

\emph{Step 3.} We split $e^{-\kappa d_j}=e^{-\kappa d_j/2}\cdot e^{-\kappa d_j/2}$ once, and the second
factor is the one belonging to $w(X)$. Since $|F^{(j)}(X)|\le f_je^{-\kappa d_j(X)}$ and
$w(X)=e^{-(\kappa/2)d_j(X)}\mathsf{r}_{\varpi}(X)$, this gives
\begin{equation*}
|F^{(j)}(X)|\,\mathsf{r}_{\varpi}(X)\le f_j\,e^{-(\kappa/2)d_j(X)}\,w(X).
\end{equation*}
By Step~1, each $X$ meeting $\Delta$ occurs at least once in the double sum over pairs $(\square,X)$ with
$\square\in\pi_j$, $\square\cap\Delta\neq\emptyset$ and $X\ni\square$. All terms are nonnegative, so
\begin{align}
\Pi_j(\Delta)
&\le f_j\Bigl[\sup_{X\cap\Delta\neq\emptyset}w(X)\Bigr]
\sum_{\square\cap\Delta\neq\emptyset}\ \sum_{X\ni\square}e^{-\kappa d_j(X)/2}\notag\\
&\le f_j\,C_we^{-c\delta_T\check{R}_k}(2L^4)^{-s}\cdot\bigl((\sigma+2)L^s\bigr)^4\cdot K_0^{(2)}\notag\\
&=(\sigma+2)^4K_0^{(2)}C_w\cdot e^{-c\delta_T\check{R}_k}\cdot 2^{-s}\cdot f_j .
\label{9.2:eq-pile-per-box-a}
\end{align}
The reasons for the two inequalities are as follows.
\begin{itemize}
\item The first inequality bounds each $w(X)$ by the supremum over the domains meeting $\Delta$.
\item The second inequality uses three bounds. The supremum is bounded by input~(a). Here $c$ is the worst
grade among the domains meeting $\Delta$, that is, the smallest; since $\delta_T\check{R}_k\ge0$, the bound
of~(a) holds for each of these domains with this $c$. The inner sum is bounded by $K_0^{(2)}$, by
input~(b). The number of cubes in the outer sum is bounded by Step~2.
\end{itemize}
The final equality holds because $L^{4s}(2L^4)^{-s}=2^{-s}$.

We put
\begin{equation*}
C_{\mathrm{pile}}(\sigma):=(\sigma+2)^4K_0^{(2)}C_w .
\end{equation*}
For $\sigma\le1$ this gives $C_{\mathrm{pile}}(\sigma)\le3^4K_0^{(2)}C_w=81\,K_0^{(2)}C_w$.

The count does not depend on where $\Delta$ lies: $\Delta$ may sit deep inside $Z_j$ or inside
$\Omega_j$. Each domain's weight $w(X)$ is taken once, at its own point $y^{\circ}$.
\end{proof}

\begin{proof}[Proof of the pile bound for stored boundary terms, continued: the sum over levels and its smallness]\label{9.2:prf-level-sum}
We prove clause~(iii) of \eqref{9.2:eq-pile-bound-b}. Fix a box $\Delta$. For every level
$j\le k$ the bound \eqref{9.2:eq-pile-per-box-a} on the pile per box reads, with the constant
$c>0$ and the level factors $f_j$ of that display,
\begin{equation}\label{9.2:eq-level-sum-1}
  \Pi_j(\Delta)\le C_{\mathrm{pile}}(\sigma)\,e^{-c\delta_T\check{R}_k}\,2^{-s}\,f_j,
  \qquad s=k-j .
\end{equation}
Here $s=k-j$ is the depth of the level $j$ below the level $k$, not the depth below the cutoff.
The constants $C_{\mathrm{pile}}(\sigma)$, $c$ and $\delta_T$ and the radius $\check{R}_k$ do not depend
on $j$. We bound $f_j$ by $\max_j f_j$ in each term of the sum over $j\le k$. The map
$j\mapsto s=k-j$ sends distinct levels $j\le k$ to distinct integers $s\ge0$, and every term
of the geometric series is non-negative. Hence \eqref{9.2:eq-level-sum-1} gives
\begin{align*}
  \sum_{j\le k}\Pi_j(\Delta)
  &\le C_{\mathrm{pile}}(\sigma)\,e^{-c\delta_T\check{R}_k}\,\max_j f_j\,\sum_{s\ge0}2^{-s}\\
  &= 2\,C_{\mathrm{pile}}(\sigma)\,e^{-c\delta_T\check{R}_k}\,\max_j f_j .
\end{align*}
This is the level sum of clause~(iii).

For its smallness we use the floor of the radii \eqref{9.2:eq-radii-floor-b}, which is part of
Notation~\ref{9.2:not-radii-floor}:
\[
  \check{R}_k\ge\bigl(\log(\gamma^{-2}-c_5)\bigr)^{r_{\mathrm{pr}}} .
\]
Since $c\delta_T>0$, the quantity $e^{-c\delta_T\check{R}_k}$ is decreasing in $\check{R}_k$. Therefore
\[
  e^{-c\delta_T\check{R}_k}\le
  \exp\Bigl(-c\delta_T\bigl(\log(\gamma^{-2}-c_5)\bigr)^{r_{\mathrm{pr}}}\Bigr),
\]
and the right-hand side tends to $0$ as $\gamma\downarrow0$.
\end{proof}

\begin{remark}[Collars with fewer layers]
The argument of the pile bound for stored boundary terms changes only in its numerical
coefficients if the collars are taken with fewer layers.

Suppose first that the collar \eqref{9.2:eq-ten-layer-collars-a} is taken with eight layers,
as \eqref{9.2:eq-ten-layer-collars-b} alone provides. Then the number $10$ is replaced by $8$
in the following places: in \eqref{9.2:eq-nested-distances-b}, in \eqref{9.2:eq-size-side-a},
in \eqref{9.2:eq-response-side-a}, in the joint inequalities and in \eqref{9.2:eq-combined-exponent-b}.
In the proof of \eqref{9.2:eq-pile-bound-b}, the floor condition on the auxiliary parameters
$\mathfrak{p}$ that is read there with ten layers is then needed with $8$ in place of $10$.

Suppose next that the step collar \eqref{9.2:eq-ten-layer-collars-e} is taken with five layers,
in the form stated in \cite{B-CMP119}. Then $8L\check{R}_{k+1}$ is replaced by
$5L\check{R}_{k+1}$, and $8\check{R}_k-3$ is replaced by $5\check{R}_k-3$. The latter satisfies
\[
  5\check{R}_k-3\ge\tfrac72\,\check{R}_k \qquad\text{whenever } \check{R}_k\ge2 ,
\]
because this inequality is equivalent to $\tfrac32\check{R}_k\ge3$. In the same way, for
$\check{R}_k\ge2$ one has $8\check{R}_k-3\ge\tfrac{13}{2}\check{R}_k$, since this is again equivalent
to $\tfrac32\check{R}_k\ge3$; accordingly the exponents $13/2$ become $7/2$.

In both cases the rate $2^{-s}$ in \eqref{9.2:eq-level-sum-1} is unchanged, so the level sum
and its smallness as $\gamma\downarrow0$ hold with the moved coefficients.
\end{remark}

\begin{remark}[Half-rate displacement of the presented argument]\label{9.2:rem-half-rate-displacement}
Let $p$ be a response piece of a unit $v$ at the step, with read set $S_v$, profile $e_p$,
decay factor $\eta_p$ satisfying $0<\eta_p\le 1$, width $\mathsf{w}>0$ and law factor
$\mathfrak{u}_p$, as in Proposition~\ref{9.2:prop-response-factor}, whose proof is incomplete
at one step. Let $\Omega$ be the fluctuation domain of the step, $d_{\mathrm{sc}}$ the scaled
distance of Definition~\ref{9.2:def-scaled-distance}, and $\delta_1^{\mathrm{sc}}$, $C_b$ the
rate and the constant of the response bounds. In this remark only we put
\begin{equation*}
R:=\frac{\mathsf{w}}{\eta_p},\qquad R':=\frac{\mathsf{w}}{\eta_p^{1/2}} .
\end{equation*}
Since $0<\eta_p\le1$ we have $\eta_p\le\eta_p^{1/2}$, hence $R'\le R$. For a law site $s$ of $v$
and the second variable $\sigma$ of the law factor, ranging over $|\sigma|\le e^{\kappa_1}$ as in
Proposition~\ref{9.2:prop-response-factor}, put
\begin{equation*}
g_s(\mathsf{t}):=\log\bigl[\mathfrak{u}_p(0,\sigma)^{-1}\mathfrak{u}_p(\mathsf{t},\sigma)\bigr](s),
\end{equation*}
where $\mathfrak{u}_p(\mathsf{t},\sigma)(s)$ takes its values in the normed algebra in which the
law factor of Proposition~\ref{9.2:prop-response-factor} takes its values; we use that this
algebra is complex, complete and unital and that its norm, written $|\cdot|$, is
submultiplicative, products and inverses being taken at the site $s$.
Here $\log a$, for an element $a$ of this algebra, denotes a logarithm of $a$, that is an element
whose exponential is $a$; hypothesis~(c) below includes that the bracket possesses the logarithm
$g_s(\mathsf{t})$ to which the bound there applies.
Thus $g_s(\mathsf{t})$ measures, in logarithmic size, the displacement of the argument
$\mathfrak{u}_p(\mathsf{t},\sigma)(s)$ presented at the law site $s$ as the amplitude moves from
$0$ to $\mathsf{t}$. We use the following hypotheses.
\begin{itemize}
\item[(a)] $e_p\le\eta_p$ on $S_v$.
\item[(b)] For every site $y$,
$e_p(y)\le C_b\,e^{-\delta_1^{\mathrm{sc}}(d_{\mathrm{sc}}(y,\Omega)-1)}$, with a constant $C_b\ge1$.
\item[(c)] On the closed disc $|\mathsf{t}|\le R$, for $|\sigma|\le e^{\kappa_1}$ and at every law
site $s$ of $v$, $|g_s(\mathsf{t})|\le\vartheta^{T}e_p(s)/\eta_p$.
\item[(d)] Holomorphy of the law factor in the amplitude, in the form in which it enters the
statement of the correlation theorem: for every law site $s$ and every $\sigma$ with
$|\sigma|\le e^{\kappa_1}$, the map $\mathsf{t}\mapsto\mathfrak{u}_p(\mathsf{t},\sigma)(s)$ is
holomorphic on the open disc $|\mathsf{t}|<R$.
\item[(e)] $0\le\vartheta^{T}<\log 2$.
\end{itemize}
Hypotheses (a)--(c) are the bounds stated in Proposition~\ref{9.2:prop-response-factor}, whose
proof is incomplete at one step; (d)
and (e) are taken by statement from the correlation theorem, (e) being the condition on the
displacement bound $\vartheta^{T}$ appearing in (c).

\smallskip
\noindent(1) Under (a) and (b), for every law site $s\in S_v$,
\begin{equation}\label{9.2:eq-half-rate-displacement-1}
e_p(s)\le\min\bigl\{\eta_p,\;C_b\,e^{-\delta_1^{\mathrm{sc}}(d_{\mathrm{sc}}(s,\Omega)-1)}\bigr\}
\le\eta_p^{1/2}\,C_b^{1/2}\,e^{-\frac12\delta_1^{\mathrm{sc}}(d_{\mathrm{sc}}(s,\Omega)-1)} .
\end{equation}

\noindent(2) Under (a)--(e), for every law site $s\in S_v$, every $\sigma$ with
$|\sigma|\le e^{\kappa_1}$ and every $\mathsf{t}$ with $|\mathsf{t}|\le R'$,
\begin{equation}\label{9.2:eq-half-rate-displacement-2}
\begin{aligned}
|g_s(\mathsf{t})|
&\le\vartheta^{T}\,\frac{e_p(s)\,|\mathsf{t}|}{\mathsf{w}}
\le\vartheta^{T}\,\frac{e_p(s)}{\eta_p^{1/2}}\\
&\le\vartheta^{T}\min\bigl\{\eta_p^{1/2},\;
C_b^{1/2}\,e^{-\frac12\delta_1^{\mathrm{sc}}(d_{\mathrm{sc}}(s,\Omega)-1)}\bigr\}.
\end{aligned}
\end{equation}

Consequently, on the disc $|\mathsf{t}|\le\rho_A=\mathsf{w}/\eta_p^{1/2}$ of
Proposition~\ref{9.2:prop-response-factor}, whose proof is incomplete at one step, on which the
Cauchy estimate in the amplitude is taken for the response pieces carrying $\eta_p^{1/2}$, the
presented argument at every law site $s\in S_v$ moves, in the size of its logarithm, by at most
\begin{equation*}
\vartheta^{T}\min\bigl\{\eta_p^{1/2},\;
C_b^{1/2}\,e^{-\frac12\delta_1^{\mathrm{sc}}(d_{\mathrm{sc}}(s,\Omega)-1)}\bigr\}.
\end{equation*}
This is the precise sense in which the displacement is paid by the distance to $\Omega$ at half
the rate $\delta_1^{\mathrm{sc}}$; equivalently, at depth $d_{\mathrm{sc}}$ the presented rotation
retains the factor $e^{-\frac12\delta_1^{\mathrm{sc}}d_{\mathrm{sc}}}$. It matches the mechanism
which \cite{B-CMP119} states in the
following words:
\begin{quote}
on the domain $\Omega^c_{k+1}$ there is no regularity improvement of the configuration
$U_{k+1}$ on the other hand we have to use a part of the analyticity domains of terms in $B_k$
for the function $H_k$. This is the reason for putting the powers of $1/2$ in the conditions
(2.36)--(2.39). The analyticity domains become smaller after each step, but the difference is
very small and exponentially decreasing in the number of steps
\end{quote}
with reference to \cite[(2.36)--(2.39)]{B-CMP119}.

The bound~\eqref{9.2:eq-half-rate-displacement-2} controls displacement, not the activity marked
by the $\Omega$-distance which the block-sum lemma uses. That Cauchy estimate, in the rescaled
amplitude $\mathsf{t}$ at the positional radius $\rho_A=\mathsf{w}/\eta_p^{1/2}$, alone gives, per
piece, the factor $\theta_k\,\mathfrak{N}(v)\,\eta_p^{1/2}$, which contains no $\Omega$-distance:
$\eta_p=1$ at touching blocks whatever the value of $d_{\mathrm{sc}}$. To turn the unused room
in~\eqref{9.2:eq-half-rate-displacement-2} into an activity factor
$e^{-\varpi\delta_1^{\mathrm{sc}}d_{\mathrm{sc}}}$ one has to spend the Lipschitz modulus, in its
configuration argument, of a stored boundary term $\mathbf{B}^{(j)}(X)$, of some level $j$ and
large-field region $X$, on the room
which \cite[(2.41)(ii)]{B-CMP119} provides for it, namely
\begin{quote}
(ii) it has an extension to an analytic function on the space $U_j(X,\alpha_0,\alpha_1)$
\end{quote}
and to decide how much of the resulting smallness is assigned to the activity and how much to
the length carried by the output. These pieces carry the square root of the decay
$e^{-\delta_1^{\mathrm{sc}}d_{\mathrm{sc}}}$ in the $\Omega$-distance, that is half of it in the
exponent; in the allocation made inside the proof of the complex-type lemma of the correlation
theorem one half of this half, a quarter of the whole, is assigned to the activity and the other
half to the length carried by the output. This is why the per-unit response
factor $\mathsf{r}_{\varpi}$ of Proposition~\ref{9.2:prop-response-factor} is carried with
$\varpi$ as a parameter and the proof of that proposition is incomplete at the step where this
allocation is needed for each unit separately; and why $\varpi=1/2$ is the natural per-unit value
on the disc of radius $\rho_A$, while $\varpi=1/4$ is the more conservative value used in the
proof of the complex-type lemma of the correlation theorem.
\end{remark}

\begin{proof}
(1) Let $s\in S_v$ be a law site. By (a), $e_p(s)\le\eta_p$; by (b) at $y=s$,
$e_p(s)\le C_b\,e^{-\delta_1^{\mathrm{sc}}(d_{\mathrm{sc}}(s,\Omega)-1)}$. Hence $e_p(s)$ is at
most the minimum of the two right-hand sides, which is the first inequality
of~\eqref{9.2:eq-half-rate-displacement-1}. For real numbers $a,b\ge0$ one has
$\min\{a,b\}^2\le ab$: if $a\le b$, then $\min\{a,b\}^2=a^2\le ab$, and symmetrically if $b\le a$.
Taking square roots, $\min\{a,b\}\le a^{1/2}b^{1/2}$. With $a=\eta_p\ge0$ and
$b=C_b\,e^{-\delta_1^{\mathrm{sc}}(d_{\mathrm{sc}}(s,\Omega)-1)}\ge0$ (recall $C_b\ge1$) this is the
second inequality
of~\eqref{9.2:eq-half-rate-displacement-1}. Thus the right-hand side
of~\eqref{9.2:eq-half-rate-displacement-1} is the geometric mean of the bounds (a) and (b); it
is the half-rate bound on $e_p$ used in the treatment of the complex-type terms of the
correlation theorem.

(2) Fix a law site $s\in S_v$ and $\sigma$ with $|\sigma|\le e^{\kappa_1}$. We first show that
$g_s$ is
holomorphic on $|\mathsf{t}|<R$. By (d), $\mathsf{t}\mapsto\mathfrak{u}_p(\mathsf{t},\sigma)(s)$ is
holomorphic there; the bracket in the definition of $g_s$,
\begin{equation*}
\Phi_s(\mathsf{t}):=\bigl[\mathfrak{u}_p(0,\sigma)^{-1}\mathfrak{u}_p(\mathsf{t},\sigma)\bigr](s),
\end{equation*}
is obtained from it by left multiplication with the element $\mathfrak{u}_p(0,\sigma)(s)^{-1}$,
which does not depend on $\mathsf{t}$, so $\Phi_s$ is holomorphic on $|\mathsf{t}|<R$. By (c), by
(a) at the site $s\in S_v$ and by (e),
\begin{equation*}
|g_s(\mathsf{t})|\le\vartheta^{T}\,\frac{e_p(s)}{\eta_p}\le\vartheta^{T}<\log 2
\qquad(|\mathsf{t}|<R).
\end{equation*}
For an element $\xi$ of the algebra with $|\xi|<\log2$, submultiplicativity gives
\begin{equation*}
|e^{\xi}-1|\le\sum_{m\ge1}\frac{|\xi|^m}{m!}=e^{|\xi|}-1<1 ,
\end{equation*}
and the double series obtained by inserting
$e^{\xi}-1=\sum_{m\ge1}\xi^m/m!$ into $\sum_{n\ge1}(-1)^{n+1}(e^{\xi}-1)^n/n$ is dominated by
$\sum_{n\ge1}(e^{|\xi|}-1)^n/n<\infty$; it may therefore be rearranged, and the identity of formal
power series $\log(e^{z})=z$ shows that the logarithmic series at $e^{\xi}$ returns $\xi$. Applied
to $\xi=g_s(\mathsf{t})$, whose exponential is $\Phi_s(\mathsf{t})$, this gives
\begin{equation*}
g_s(\mathsf{t})=\sum_{n\ge1}\frac{(-1)^{n+1}}{n}\bigl(\Phi_s(\mathsf{t})-1\bigr)^n,
\qquad |\Phi_s(\mathsf{t})-1|\le e^{\vartheta^{T}}-1<1 ,
\end{equation*}
for $|\mathsf{t}|<R$. The series therefore converges uniformly on $|\mathsf{t}|<R$; its terms are
holomorphic, being polynomials in the holomorphic function $\Phi_s$, and so $g_s$ is holomorphic
on $|\mathsf{t}|<R$. It vanishes at $\mathsf{t}=0$, since there the bracket is
$\mathfrak{u}_p(0,\sigma)^{-1}\mathfrak{u}_p(0,\sigma)$, the identity, whose logarithm is $0$. By
(c), $|g_s|\le B:=\vartheta^{T}e_p(s)/\eta_p$ on $|\mathsf{t}|<R$.

Let $\ell$ be a continuous linear functional of norm at most $1$ on the value space. Then
$h:=\ell\circ g_s$ is a scalar holomorphic function on $|\mathsf{t}|<R$ with $h(0)=0$ and
$|h|\le B$. If $B=0$ then $h\equiv0$. If $B>0$, the function $z\mapsto h(Rz)/B$ is holomorphic on
the unit disc, vanishes at $0$ and is bounded by $1$ in modulus, so by the Schwarz lemma
$|h(Rz)|\le B|z|$ for $|z|<1$, that is $|h(\mathsf{t})|\le B|\mathsf{t}|/R$ for
$|\mathsf{t}|<R$. Taking the supremum over all such $\ell$, which by the Hahn--Banach theorem
equals the norm,
\begin{equation*}
|g_s(\mathsf{t})|\le\vartheta^{T}\,\frac{e_p(s)}{\eta_p}\,\frac{|\mathsf{t}|}{R}
=\vartheta^{T}\,\frac{e_p(s)\,|\mathsf{t}|}{\mathsf{w}}\qquad(|\mathsf{t}|<R),
\end{equation*}
where the equality is $R=\mathsf{w}/\eta_p$. On the circle $|\mathsf{t}|=R$ the same inequality
reads $|g_s(\mathsf{t})|\le\vartheta^{T}e_p(s)/\eta_p$, which is (c) itself. Hence the first
inequality of~\eqref{9.2:eq-half-rate-displacement-2} holds on $|\mathsf{t}|\le R$, in
particular on $|\mathsf{t}|\le R'$, as $R'\le R$. For $|\mathsf{t}|\le R'=\mathsf{w}/\eta_p^{1/2}$
we have $|\mathsf{t}|/\mathsf{w}\le\eta_p^{-1/2}$, which gives the second inequality.

For the third, $e_p(s)/\eta_p^{1/2}\le\eta_p/\eta_p^{1/2}=\eta_p^{1/2}$ by (a), and, dividing the
outer inequality of~\eqref{9.2:eq-half-rate-displacement-1} by $\eta_p^{1/2}>0$,
\begin{equation*}
\frac{e_p(s)}{\eta_p^{1/2}}\le C_b^{1/2}\,e^{-\frac12\delta_1^{\mathrm{sc}}(d_{\mathrm{sc}}(s,\Omega)-1)}
\end{equation*}
by part (1). Thus $e_p(s)/\eta_p^{1/2}$ is at most the minimum of the two, and multiplying by
$\vartheta^{T}\ge0$ gives the last inequality of~\eqref{9.2:eq-half-rate-displacement-2}.
\end{proof}

\begin{lemma}[The weakly anchored case]\label{9.2:lem-weak-anchor}
The setting is that of the pile bound for stored boundary terms,
Theorem~\ref{9.2:thm-pile-bound}. We are given the frozen step $k\to k+1$ on a common label, a level
$j\le k$ with $s:=k-j$, and a level-$j$ family $F$ of stored boundary terms. The floors of that
theorem are in force, among them $\check{R}_k\ge 2$ and $L\ge 3$. Assume that the units $X$ of $F$
carry the weak anchor \eqref{9.2:eq-anchors-b}, and that each unit $X$ accordingly has points
$y_0\in X$ and $z_0\in Z_j$ with $|y_0-z_0|_\infty\le(1+a)M\ell_j$ for a number $a\ge 0$, so that
the size side, Lemma~\ref{9.2:lem-size-side}, applies to $X$ with this value of $a$. Assume further
that
\begin{equation}\label{9.2:eq-weak-anchor-1}
a\le L\check{R}_{j+1}\quad\text{if } s\ge 1,\qquad\qquad a\le \check{R}_k\quad\text{if } s=0 .
\end{equation}
Then the three grades $(\alpha)$, $(\beta)$, $(\gamma)$ of the per-unit bound of
Theorem~\ref{9.2:thm-pile-bound} hold under the same floors, with the constants $c$ given below: for
every unit $X$ of $F$ the joint exponent $E'(X)$ of Lemma~\ref{9.2:lem-combined-exponent} satisfies
\begin{equation}\label{9.2:eq-weak-anchor-2}
\delta_T E'(X)\ \ge\ s\Lambda_L+c\,\delta_T\check{R}_k,
\qquad\text{so that}\qquad
w(X)\ \le\ C_w\,e^{-c\,\delta_T\check{R}_k}\,(2L^4)^{-s},
\end{equation}
with $c=\tfrac{11}{2},\ \tfrac{11}{2},\ \tfrac32$ in the grades $(\alpha)$, $(\beta)$, $(\gamma)$
respectively. More precisely:
\begin{itemize}
\item if $s\ge1$, one may take $c=\tfrac{13}{2},\ 8,\ 3$;
\item if $s=0$, one may take $c=\tfrac{11}{2}$ in each grade;
\item if in addition the floor $\check{R}_{k+1}\ge 2$ holds, one may take
$c=\tfrac{13}{2},\ \tfrac{13}{2},\ \tfrac32$ for every $s$.
\end{itemize}
\end{lemma}

\begin{proof}
Inequality \eqref{9.2:eq-weak-anchor-1} is a hypothesis; that it follows from the data that define
the weak anchor \eqref{9.2:eq-anchors-b} is not proved, and the last paragraph below explains why a
bound of this form is needed.

\emph{Notation and inputs.} We write $\check{R}^{\flat}=\check{R}_k$. The following objects are
those of Lemmas~\ref{9.2:lem-size-side}, \ref{9.2:lem-response-side} and
\ref{9.2:lem-combined-exponent}:
\begin{itemize}
\item the minimising point $y^{\circ}$, with depth $\iota^{\circ}:=\iota(y^{\circ})$;
\item the depth $\iota:=\iota(X)$ of the unit;
\item the thick indicators $\theta_i\in\{0,1\}$ at $y^{\circ}$;
\item the sets $Z_i$ and $\Omega$.
\end{itemize}
We write $y^{*}$ for a point of $X$ of depth $\iota(y^{*})=\iota$. At the levels paid by the size
side we put $\hat\theta_i:=1$ for $i<\iota^{\circ}$, and $\hat\theta_i:=\theta_i$ for
$i\ge\iota^{\circ}$.

The anchor enters only the size side \eqref{9.2:eq-size-side-a}, through the term $-3-a$. The
response side \eqref{9.2:eq-response-side-a} does not read the anchor, and neither does the level
counting in the proof of the per-unit bound of Theorem~\ref{9.2:thm-pile-bound}. That counting uses
Lemma~\ref{9.2:lem-thin-levels}, on which levels can be thin, the per-window balance,
Lemma~\ref{9.2:lem-window-balance} with \eqref{9.2:eq-window-balance-b} and
\eqref{9.2:eq-window-balance-c}, and the count in three grades \eqref{9.2:eq-per-unit-grades-a}.
Besides these statements we use the
monotonicity and floor of the radii, \eqref{9.2:eq-radii-floor-a}, recorded in
Notation~\ref{9.2:not-radii-floor}: by its part (a), $\check{R}_i\ge\check{R}_k$ for $i\le k$; by its
part (b), $L\check{R}_{k+1}\ge\check{R}_k$. We also use the floors $\check{R}_k\ge2$ and $L\ge3$, and
the hypothesis \eqref{9.2:eq-weak-anchor-1}.

\emph{Elementary inequalities.} Under \eqref{9.2:eq-weak-anchor-1} we have
\begin{equation}\label{9.2:eq-weak-anchor-a}
\begin{aligned}
&10L\check{R}_{j+1}-3-a\ \ge\ 9L\check{R}_{j+1}-3\ \ge\ 10\check{R}_{j+1}
&&\text{if } s\ge1,\\
&8L\check{R}_{k+1}-3-a\ \ge\ 7\check{R}_k-3\ \ge\ \tfrac{11}{2}\check{R}_k
&&\text{if } s=0 .
\end{aligned}
\end{equation}
For the first line of \eqref{9.2:eq-weak-anchor-a}, the first inequality is $a\le L\check{R}_{j+1}$.
The second is equivalent to
$(9L-10)\check{R}_{j+1}\ge3$. This holds because $9L-10\ge17$ (as $L\ge3$) and
$\check{R}_{j+1}\ge\check{R}_k\ge2$ (part (a), as $j+1\le k$), so that
$(9L-10)\check{R}_{j+1}\ge 17\cdot2\ge3$.

For the second line, part (b) gives $8L\check{R}_{k+1}\ge 8\check{R}_k$, and
$a\le\check{R}_k$; together they give the first inequality. The second is
$\tfrac32\check{R}_k\ge3$, which is $\check{R}_k\ge2$.

If the floor $\check{R}_{k+1}\ge2$ holds, the second line improves. Since
$a\le\check{R}_k\le L\check{R}_{k+1}$, and $L\check{R}_{k+1}\ge6$ because $L\ge3$ and
$\check{R}_{k+1}\ge2$,
\begin{equation}\label{9.2:eq-weak-anchor-3}
8L\check{R}_{k+1}-3-a\ \ge\ 8L\check{R}_{k+1}-L\check{R}_{k+1}-3
=7L\check{R}_{k+1}-3\ \ge\ \tfrac{13}{2}L\check{R}_{k+1}\ \ge\ \tfrac{13}{2}\check{R}_k .
\end{equation}
Here $7L\check{R}_{k+1}-3\ge\tfrac{13}{2}L\check{R}_{k+1}$ is $L\check{R}_{k+1}\ge6$, and the last
step is part (b).

Three further facts on real numbers are needed. For $u\ge0$, $C\ge0$ and $\lambda'\ge1$,
\begin{equation}\label{9.2:eq-weak-anchor-4}
\lambda' u+(C-u)_+\ \ge\ C .
\end{equation}
Indeed, if $u\le C$ the left side is $C+(\lambda'-1)u\ge C$, and if $u>C$ it is at least
$\lambda'u\ge u>C$.

For $L>1$ and $A,B\ge0$ with $A/L\le B$,
\begin{equation}\label{9.2:eq-weak-anchor-5}
\min_{v\ge0}\bigl[\max\{0,vL-A\}+(B-v)_+\bigr]=B-\frac{A}{L},
\end{equation}
attained at $v=A/L$. To see this, note how the bracket behaves on three intervals:
\begin{itemize}
\item on $[0,A/L]$ it equals $B-v$, since $v\le A/L\le B$, and so has slope $-1$;
\item on $[A/L,B]$ it equals $vL-A+B-v$, with slope $L-1>0$;
\item for $v\ge B$ it equals $vL-A$, with slope $L$.
\end{itemize}
The bracket is thus decreasing on $[0,A/L]$ and increasing on $[A/L,\infty)$, so its minimum over
$v\ge0$ is its value $B-A/L$ at $v=A/L$.

For $0\le A\le C$,
\begin{equation}\label{9.2:eq-weak-anchor-6}
\min_{u\ge0}\bigl[\max\{0,u-A\}+(C-u)_+\bigr]=C-A .
\end{equation}
Indeed, the bracket equals $C-u\ge C-A$ on $[0,A]$, equals $C-A$ on $[A,C]$, and equals
$u-A\ge C-A$ for $u\ge C$; its minimum over $u\ge0$ is therefore $C-A$.

\emph{The cases.} We call $8L\check{R}_{k+1}$ the collar term. In each case below, $E'(X)$ is
bounded below by combining a lower bound for $d_j(X)$ from the size side with the response side at
$y^{\circ}$; the cases are distinguished by $s$, $\iota^{\circ}$ and $\iota$.
\begin{enumerate}
\item[Case 1.] $\iota=j$. The anchor is not read in this case, and
\begin{equation*}
E'(X)\ \ge\ \sum_{i=j+1}^{k}10\check{R}_i\theta_i+8L\check{R}_{k+1}.
\end{equation*}
\item[Case 2.] $s\ge1$ and $\iota^{\circ}=\iota=j+1$. Put
$v:=\operatorname{dist}_\infty(y^{\circ},Z_j)/(M\ell_{j+1})$. The size side at $y=y^{\circ}$, whose
depth is $j+1$, has an empty sum and $t=\operatorname{dist}_\infty(y^{\circ},Z_j)$; since
$\ell_{j+1}=L\ell_j$, its distance term is $t/(M\ell_j)=vL$. Together with $d_j(X)\ge0$ it gives
$d_j(X)\ge\max\{0,vL-3-a\}$, and the response side gives the remaining terms:
\begin{equation*}
E'(X)\ \ge\ \max\{0,vL-3-a\}+\theta_{j+1}\bigl(10\check{R}_{j+1}-v\bigr)_+
+\sum_{i=j+2}^{k}10\check{R}_i\theta_i+8L\check{R}_{k+1}.
\end{equation*}
If $\theta_{j+1}=0$ we use only $d_j(X)\ge0$, which does not involve $a$, so nothing is lost
relative to the case $a=0$. If
$\theta_{j+1}=1$, we apply \eqref{9.2:eq-weak-anchor-5} with $A=3+a$ and $B=10\check{R}_{j+1}$;
its hypothesis $A/L\le B$ holds, since $(3+a)/L\le 1+\check{R}_{j+1}\le10\check{R}_{j+1}$ by
$L\ge3$, $a\le L\check{R}_{j+1}$ and $\check{R}_{j+1}\ge2$. Minimising over $v\ge0$, level $j+1$ is
credited
\begin{equation*}
10\check{R}_{j+1}-\frac{3+a}{L}\ \ge\ 9\check{R}_{j+1}-\frac{3}{L}\ \ge\ \frac{17}{2}\check{R}^{\flat},
\end{equation*}
where the first inequality is $a\le L\check{R}_{j+1}$, and the second uses
$\check{R}_{j+1}\ge\check{R}^{\flat}$ and $3/L\le1\le\check{R}^{\flat}/2$. This is the only case in
which a level is credited less than $10\check{R}_i\hat\theta_i$; the collar term is not reduced.
\item[Case 3.] $s\ge2$ and $\iota^{\circ}=\iota\in[j+2,k]$. The sum of the size side at $y^{\circ}$
contains the term $10L\check{R}_{j+1}$ (index $i=j+1$), which absorbs $-3-a$ by the first line of
\eqref{9.2:eq-weak-anchor-a}; the terms with $i\ge j+2$ are at least $10\check{R}_i$, as
$L^{i-j}\ge1$. The distance term of the size side and the response term at level $\iota$ are
combined by \eqref{9.2:eq-weak-anchor-4}. This gives
\begin{equation*}
E'(X)\ \ge\ \sum_{i=j+1}^{\iota-1}10\check{R}_i+\sum_{i=\iota}^{k}10\check{R}_i\theta_i
+8L\check{R}_{k+1}.
\end{equation*}
\item[Case 4.] $s\ge1$ and $\iota^{\circ}=j<\iota=j+1$. As in Case 1: $d_j(X)\ge0$ and the response
side at $y^{\circ}\in Z_j$ give the bound of Case 1.
\item[Case 5.] $s\ge2$ and $\iota^{\circ}<\iota\in[j+2,k]$. The size side at $y^{*}$, with its
distance term $t\ge0$ dropped, the first line of \eqref{9.2:eq-weak-anchor-a} for the term $i=j+1$
and $L^{i-j}\ge1$ for the others, give
$d_j(X)\ge\sum_{i=j+1}^{\iota-1}10\check{R}_i$. The response side at $y^{\circ}$ then gives
\begin{equation*}
E'(X)\ \ge\ \sum_{i=j+1}^{\iota-1}10\check{R}_i+\sum_{i=\iota}^{k}10\check{R}_i\theta_i
+8L\check{R}_{k+1}.
\end{equation*}
\item[Case 6.] $s\ge1$ and $y^{\circ}\in\Omega$. The clause of the size side for points of
$\Omega$, with $L^{s}\ge1$, $L^{i-j}\ge1$ and the first line of \eqref{9.2:eq-weak-anchor-a},
gives
\begin{equation*}
E'(X)\ \ge\ \sum_{i=j+1}^{k}10\check{R}_i+8L\check{R}_{k+1}.
\end{equation*}
\item[Case 7.] $s=0$ and $y^{\circ}\in\Omega$. The same clause, whose sum is now empty and in
which $L^{s}=1$, and the second line of \eqref{9.2:eq-weak-anchor-a} give
\begin{equation*}
E'(X)\ \ge\ 8L\check{R}_{k+1}-3-a\ \ge\ \tfrac{11}{2}\check{R}_k ;
\end{equation*}
under the floor $\check{R}_{k+1}\ge2$, \eqref{9.2:eq-weak-anchor-3} replaces $\tfrac{11}{2}$ by
$\tfrac{13}{2}$.
\item[Case 8.] $s\ge1$, $\iota^{\circ}=k+1$ and $y^{\circ}\notin\Omega$. The size side at
$y^{\circ}$, the first line of \eqref{9.2:eq-weak-anchor-a} for its term $i=j+1$, and
\eqref{9.2:eq-weak-anchor-4} applied to its distance term together with the response side give
\begin{equation*}
E'(X)\ \ge\ \sum_{i=j+1}^{k}10\check{R}_i+8L\check{R}_{k+1}.
\end{equation*}
\item[Case 9.] $s=0$ and $y^{\circ}\notin\Omega\cup Z_k$. Here $j=k$ and $\iota(y^{\circ})=k+1$;
let $u:=\operatorname{dist}_\infty(y^{\circ},Z_k)/(M\ell_k)$ be the distance term of the size side
at $y^{\circ}$, whose sum is empty. The size side, together with $d_j(X)\ge0$, and the response side
give
\begin{equation*}
E'(X)\ \ge\ \max\{0,u-3-a\}+\bigl(8L\check{R}_{k+1}-u\bigr)_+ .
\end{equation*}
We apply \eqref{9.2:eq-weak-anchor-6} with $A=3+a$ and $C=8L\check{R}_{k+1}$; its hypothesis
$A\le C$ holds since $A\le\check{R}_k+3$ and $8L\check{R}_{k+1}\ge8\check{R}_k\ge\check{R}_k+3$ by
part (b) and $\check{R}_k\ge2$. With the second line of \eqref{9.2:eq-weak-anchor-a} this gives
\begin{equation*}
E'(X)\ \ge\ 8L\check{R}_{k+1}-3-a\ \ge\ \tfrac{11}{2}\check{R}_k ,
\end{equation*}
and $\tfrac{13}{2}\check{R}_k$ under the floor $\check{R}_{k+1}\ge2$, by
\eqref{9.2:eq-weak-anchor-3}.
\item[Case 10.] $s\ge1$ and $\iota^{\circ}\le k<\iota$. The size side at $y^{*}$, with its distance
term dropped, and the first line of \eqref{9.2:eq-weak-anchor-a}, together with the response side
at $y^{\circ}$, give
\begin{equation*}
E'(X)\ \ge\ \sum_{i=j+1}^{k}10\check{R}_i+8L\check{R}_{k+1}.
\end{equation*}
\item[Case 11.] $s=0$ and $\iota^{\circ}=k<\iota$. In this case $E'(X)\ge8L\check{R}_{k+1}$.
\end{enumerate}
These cases exhaust all possible values of $(\iota^{\circ},\iota,s)$. There are two shortfalls
relative to the case $a=0$: level $j+1$ in Case 2, at $s\ge1$, and the collar term in Cases 7
and 9, at $s=0$. They never occur together, since the first requires $s\ge1$ and the second $s=0$.

\emph{The grades for $s\ge1$.} By the cases, every level $i\in[j+1,k]$ is credited at least
$10\check{R}_i\hat\theta_i$, as in the count for $a=0$ in the proof of the per-unit bound
\eqref{9.2:eq-per-unit-grades-a}, except level $j+1$ in Case 2, which is credited at least
$9\check{R}_{j+1}-3/L\ge9\check{R}^{\flat}-1$ (as $3/L\le1$) and at least
$\tfrac{17}{2}\check{R}^{\flat}$. The collar term is at least $8\check{R}_k$ by part (b).

Grade $(\alpha)$. Under the floor of grade $(\alpha)$ in \eqref{9.2:eq-pile-bound-a},
$10\delta_T\check{R}^{\flat}\ge\Lambda_L$, the levels $i\ge j+2$ pay at least $\Lambda_L$ each, as for
$a=0$. Level $j+1$ pays at least
\begin{equation*}
\delta_T\bigl(9\check{R}^{\flat}-1\bigr)=10\delta_T\check{R}^{\flat}-\delta_T\bigl(\check{R}^{\flat}+1\bigr)
\ \ge\ \Lambda_L-\tfrac32\delta_T\check{R}_k ,
\end{equation*}
since $1\le\check{R}_k/2$ and $\check{R}^{\flat}=\check{R}_k$. The collar pays $8\delta_T\check{R}_k$.
Hence
\begin{equation*}
\delta_T E'(X)\ \ge\ (s-1)\Lambda_L+\Lambda_L-\tfrac32\delta_T\check{R}_k+8\delta_T\check{R}_k
= s\Lambda_L+\tfrac{13}{2}\delta_T\check{R}_k .
\end{equation*}

Grade $(\beta)$. Under the floor of grade $(\beta)$ in \eqref{9.2:eq-pile-bound-a},
$10\delta_T\check{R}^{\flat}\ge2\Lambda_L$. A thick level $j+1$ that is free in the count of
\eqref{9.2:eq-per-unit-grades-a} pays at least
\begin{equation*}
\tfrac{17}{2}\delta_T\check{R}^{\flat}\ \ge\ \tfrac{17}{10}\Lambda_L\ \ge\ \Lambda_L .
\end{equation*}
Consider a window $I_b\subset[j+1,k]$ of the per-window balance whose bottom is $h_b=j+1$. By
\eqref{9.2:eq-window-balance-b}, $k_0(b)-h_b=N_b-N_0(b)\ge N_0(b)\ge1$, so $j+1\le k_0(b)$; since the
thin levels of the window lie in $[k_0(b)+1,\,j_b-1]$ (Lemma~\ref{9.2:lem-thin-levels}), level
$j+1$ is thick there. In the count for $a=0$ the window receives $10\delta_T\check{R}^{\flat}(N_b/2+1)$
from its thick levels; at level $j+1$ at most $\tfrac32\delta_T\check{R}^{\flat}$ of this is lost, so
the window balances:
\begin{equation*}
10\delta_T\check{R}^{\flat}\bigl(\tfrac{N_b}{2}+1\bigr)-\tfrac32\delta_T\check{R}^{\flat}
=\bigl(5N_b+\tfrac{17}{2}\bigr)\delta_T\check{R}^{\flat}\ \ge\ 5N_b\delta_T\check{R}^{\flat}
\ \ge\ N_b\Lambda_L .
\end{equation*}
Every other window satisfies
\begin{equation*}
10\delta_T\check{R}^{\flat}\bigl(\tfrac{N_b}{2}+1\bigr)\ \ge\ (N_b+2)\Lambda_L .
\end{equation*}
A window with bottom at most $j$ meeting $[j+1,k]$ would be exceptional, and such windows are
excluded in grade $(\beta)$. The collar is not borrowed from and pays $8\delta_T\check{R}_k$ in full.
Hence $\delta_T E'(X)\ge s\Lambda_L+8\delta_T\check{R}_k$.

Grade $(\gamma)$. Under the conditions of grade $(\gamma)$ in \eqref{9.2:eq-pile-bound-a}, among them
$4\delta_T\check{R}^{\flat}\ge\Lambda_L$: if level $j+1$ is thin, it is one of the $d^{*}$ thin levels
counted in the proof of \eqref{9.2:eq-per-unit-grades-a}, and nothing differs from the case $a=0$.
If it is thick, it pays at least
\begin{equation*}
\tfrac{17}{2}\delta_T\check{R}^{\flat}\ \ge\ \tfrac{17}{8}\Lambda_L\ \ge\ \Lambda_L ,
\end{equation*}
which is its own charge. The rest of the count is that of \eqref{9.2:eq-per-unit-grades-a}, its
second branch borrowing $5\delta_T\check{R}_k-\Lambda_L$ from the collar. Hence
$\delta_T E'(X)\ge s\Lambda_L+3\delta_T\check{R}_k$.

\emph{The case $s=0$.} At $s=0$ only Cases 1, 7, 9 and 11 can occur. In Cases 1 and 11 the sums are
empty and $E'(X)\ge8L\check{R}_{k+1}\ge8\check{R}_k$ by part (b); in Cases 7 and 9,
$E'(X)\ge\tfrac{11}{2}\check{R}_k$. Hence $\delta_T E'(X)\ge\tfrac{11}{2}\delta_T\check{R}_k$ in each
grade, and $\delta_T E'(X)\ge\tfrac{13}{2}\delta_T\check{R}_k$ under the floor $\check{R}_{k+1}\ge2$.

\emph{Conclusion.} Taking in each grade the smaller of the values at $s\ge1$ and at $s=0$ gives
$\tfrac{11}{2}$, $\tfrac{11}{2}$ and $3$; the displayed $\tfrac32$ in grade $(\gamma)$ is smaller than
the value $3$ proved. Under the floor $\check{R}_{k+1}\ge2$ the minima are $\tfrac{13}{2}$,
$\tfrac{13}{2}$ and $3\ge\tfrac32$. Since $w(X)\le C_w e^{-\delta_T E'(X)}$
(Lemma~\ref{9.2:lem-combined-exponent}, as used in \eqref{9.2:eq-per-unit-grades-a}) and
$e^{-s\Lambda_L}=(2L^4)^{-s}$ by $\Lambda_L=\log(2L^4)$, each bound
$\delta_TE'(X)\ge s\Lambda_L+c\,\delta_T\check{R}_k$ gives the bound on $w(X)$ in
\eqref{9.2:eq-weak-anchor-2}. Parts (ii) and (iii) of the pile bound \eqref{9.2:eq-pile-bound-b}
read part (i) only through the factor $(2L^4)^{-s}$ and a prefactor independent of $k$ and $j$;
their proofs apply without change with the constants above.

\emph{Sharpness relative to the inputs.} If the size side and the response side hold with equality,
the credits computed above are attained. At level $j+1$, take $a=L\check{R}_{j+1}$ and
$v=\check{R}_{j+1}+3/L$: then $vL-3-a=0$, so $d_j(X)=0$ is allowed, and level $j+1$ is credited
$10\check{R}_{j+1}-v=9\check{R}_{j+1}-3/L$. At the collar, take $L\check{R}_{k+1}=\check{R}_k=2$ and
$a=2$: then $8L\check{R}_{k+1}-3-a=11=\tfrac{11}{2}\check{R}_k$.

\emph{Why the form of \eqref{9.2:eq-weak-anchor-1} matters.} If $a$ were not bounded by a fixed
multiple of $L\check{R}_{j+1}$, then the minimum over $v\ge0$ of the bracket in
\eqref{9.2:eq-weak-anchor-5} with $A=3+a$ and $B=10\check{R}_{j+1}$, which is
$\max\{0,\,10\check{R}_{j+1}-(3+a)/L\}$ (for $A/L>B$ the bracket vanishes at $v=B$), would vanish as
soon as $a\ge10L\check{R}_{j+1}-3$, and a thick level $j+1$ would have no source for its factor
$(2L^4)^{-1}$. On the other hand, any bound of $a$ by a fixed multiple of $L\check{R}_{j+1}$ with
factor less than $9$, together with a bound of $a$ at $s=0$ by a fixed multiple of $\check{R}_k$
with factor less than $5$, still gives bounds of the form \eqref{9.2:eq-weak-anchor-2}, with smaller
constants $c$.
\end{proof}

\begin{remark}[Cost of thin levels under a normalisation floor]\label{9.2:rem-thin-level-cost}
Let $y$ be a point, let $a\in\mathfrak{A}(y)$ be an attempt in force at $y$, with window data
$k_a$, $j_a$, $k_0(a)$, $N_0(a)$ as in Definition~\ref{9.2:def-attempt-windows}, and let $i\in\Theta_a$
be a thin level for $y$. Let $P_i$ be the portion at level $i$ of a path $\Gamma$ from $y$ to $\Omega$,
as in the response side lemma (Lemma~\ref{9.2:lem-response-side}). Put
\begin{equation}\label{9.2:eq-thin-level-cost-1}
\mathfrak{r}_a:=\check{R}_{k_0(a)+1}\,L^{-(N_0(a)-1)} .
\end{equation}
Then
\begin{equation}\label{9.2:eq-thin-level-cost-2}
\ell_{\mathrm{sc}}(P_i)\ \ge\ 10\,L\,\mathfrak{r}_a .
\end{equation}
In \cite{B-CMP122a} Ba{\l}aban fixes an integer $N_0$ by $L^{-N_0+1}MR_{k-N_0+1}=M$, where $R_k$
denotes the radii of his construction, distinct from the ball radius $R$ of the glossary
(Notation~\ref{0.2:not-glossary}). If $N_0(a)$ is the analogous integer for the radii
$\check{R}$, that is $L^{-N_0(a)+1}M\check{R}_{k_0(a)+1}=M$, then $\mathfrak{r}_a=1$.
The remark uses this normalisation of $N_0(a)$ only as a hypothesis and does not decide whether
it holds.
Put $\mathfrak{r}:=\inf_a\mathfrak{r}_a$, the infimum over all non-completed attempts $a$ of the
label (Definition~\ref{9.2:def-attempt-windows}). Under the floor
\begin{equation}\label{9.2:eq-thin-level-cost-3}
10\,\mathfrak{r}\,\varpi\,\delta_1^{\mathrm{sc}}\ \ge\ \Lambda_L ,
\end{equation}
which bounds from below the normalisation of $\delta_1^{\mathrm{sc}}$ relative to the decay rate of
\cite{B-CMP122a} (there $M$ is chosen large enough that Ba{\l}aban's decay constant $\delta$
satisfies $\delta M/M_1\ge 2$), every thin level would contribute at least $\Lambda_L$ to the decay
exponent along $\Gamma$, in the sense that
$\varpi\,\delta_1^{\mathrm{sc}}\,\ell_{\mathrm{sc}}(P_i)\ge\Lambda_L$, and the grades $(\beta)$ and
$(\gamma)$ into which thin levels are sorted in the count of thin levels would contain no level.
The floor \eqref{9.2:eq-thin-level-cost-3} is
not assumed anywhere in this book, and no statement of the book depends on it; whether it can be
met within the order of choice of the constants (Notation~\ref{2.3:not-stages-first},
Notation~\ref{2.3:not-stages-middle} and Notation~\ref{2.3:not-stages-last}) is not examined here.
\end{remark}

\begin{proof}
By the response side lemma (Lemma~\ref{9.2:lem-response-side},
\eqref{9.2:eq-response-side-a}), the portion $P_i$ runs inside $A_i(y)$ and its endpoints are at
distance at least $10M\check{R}_i\ell_i$; this holds on thin levels as well. Since $i\in\Theta_a$,
the lemma on which levels can be thin (Lemma~\ref{9.2:lem-thin-levels},
\eqref{9.2:eq-thin-levels-a}) applies with the unique attempt $a$: it gives
$i\in[k_0(a)+1,\,j_a-1]$, and every $x'\in A_i(y)$ has $n(x')\le j_a\le k_a-1$. So all scales met by
$P_i$ are at most $k_a-1$, and the definition of the scaled distance,
\eqref{9.2:eq-scaled-distance-a} in its second part, applied with $m=k_a-1$ converts the endpoint
distance into scaled length; as $\ell_i/\ell_{k_a-1}=L^{-(k_a-1-i)}$,
\begin{equation}\label{9.2:eq-thin-level-cost-4}
\ell_{\mathrm{sc}}(P_i)\ \ge\ 10\,\check{R}_i\,L^{-(k_a-1-i)} .
\end{equation}
By part (b) of the monotonicity and floor of the radii (Notation~\ref{9.2:not-radii-floor},
\eqref{9.2:eq-radii-floor-a}), $\check{R}_l\le L\check{R}_{l+1}$ for every $l$; iterating from
$l=k_0(a)+1$ to $l=i-1$, which is possible since $i\ge k_0(a)+1$, gives
\[
\check{R}_i\ \ge\ \check{R}_{k_0(a)+1}\,L^{-(i-k_0(a)-1)} .
\]
Inserting this in \eqref{9.2:eq-thin-level-cost-4} and adding the exponents,
$(i-k_0(a)-1)+(k_a-1-i)=k_a-k_0(a)-2$, which equals $N_0(a)-2$ since $k_a-k_0(a)=N_0(a)$ for the
window of $a$ (Definition~\ref{9.2:def-attempt-windows}). Hence
\[
\ell_{\mathrm{sc}}(P_i)\ \ge\ 10\,\check{R}_{k_0(a)+1}\,L^{-(N_0(a)-2)}\ =\ 10\,L\,\mathfrak{r}_a ,
\]
which is \eqref{9.2:eq-thin-level-cost-2}.

If $N_0(a)$ is the analogous integer for the radii $\check{R}$, that is
$L^{-N_0(a)+1}M\check{R}_{k_0(a)+1}=M$, equivalently $L^{N_0(a)-1}=\check{R}_{k_0(a)+1}$, then
\eqref{9.2:eq-thin-level-cost-1} gives
\[
\mathfrak{r}_a\ =\ \check{R}_{k_0(a)+1}\,L^{-(N_0(a)-1)}\ =\ 1 .
\]

Finally, if \eqref{9.2:eq-thin-level-cost-3} held, then by \eqref{9.2:eq-thin-level-cost-2} and
$\mathfrak{r}_a\ge\mathfrak{r}$,
\[
\varpi\,\delta_1^{\mathrm{sc}}\,\ell_{\mathrm{sc}}(P_i)\ \ge\ 10\,L\,\mathfrak{r}\,\varpi\,
\delta_1^{\mathrm{sc}}\ \ge\ L\,\Lambda_L\ \ge\ \Lambda_L ,
\]
so each thin level would contribute at least $\Lambda_L$ in the sense stated, and the grades
$(\beta)$ and $(\gamma)$ would contain no level.
\end{proof}

\begin{lemma}[Read sets of cells meeting the domain]\label{9.2:lem-read-set}
Let $\Omega$ be the fluctuation domain and $d_{\mathrm{sc}}$ the scaled distance to $\Omega$ of
Definition~\ref{9.2:def-scaled-distance}. For a set $S$ write $d_{\mathrm{sc}}(S,\Omega)$ for the
infimum of $d_{\mathrm{sc}}(y,\Omega)$ over $y\in S$. Let $X$ be a domain, and call \emph{cells}
the blocks of side $\ell_k$ of the level at which the domain $X$ is read. Assume the following two
properties of the scaled distance, which are used here by statement.
\begin{itemize}
\item[(a)] For every point $s$, every cell containing $s$ lies in the local $M$-cube of $s$.
\item[(b)] For every point $s$ and every point $\tilde{s}$ of a cell containing $s$ (by (a), such a
point lies in the local $M$-cube of $s$), passing from $s$ to $\tilde{s}$ costs at most one unit of
$d_{\mathrm{sc}}$, that is, $d_{\mathrm{sc}}(\tilde{s},\Omega)\ge d_{\mathrm{sc}}(s,\Omega)-1$.
\end{itemize}
Let $v$ be a unit of the summation of boundary terms at finer zones in the expansion of the
correlation theorem, in the case of Proposition~\ref{9.2:prop-response-factor} in which the width
is $\mathsf{w}_T$, and let $S_v$ be its read set in the sense of the correlation theorem's
localisation statement; suppose that $S_v$
is the union of the cells meeting $X$, rather than $X$ itself. Then
$d_{\mathrm{sc}}(S_v,\Omega)\ge d_{\mathrm{sc}}(X,\Omega)-1$, and the per-unit response factor
$\mathsf{r}_{\varpi}$ of \eqref{9.2:eq-response-factor-a} (a display of
Proposition~\ref{9.2:prop-response-factor}, stated; proof incomplete at the step marked $(\ast)$)
satisfies
$\mathsf{r}_{\varpi}(S_v)\le e^{\varpi\delta_1^{\mathrm{sc}}}\mathsf{r}_{\varpi}(X)$. Let $\Delta$ be
a box of side $\sigma M\ell_k$ that is a union of cells, and let $\Delta^{+}$ be $\Delta$ enlarged
by one layer of cells, of
side $(\sigma+2/M)M\ell_k$. If $S_v$ meets $\Delta$, then $X$ meets $\Delta^{+}$. Consequently,
the constants of \eqref{9.2:eq-response-factor-a} and of \eqref{9.2:eq-pile-per-box-a} are replaced
as follows:
\begin{equation}\label{9.2:eq-read-set-a}
\begin{gathered}
C_w\ \longmapsto\ C_w\,e^{\varpi\delta_1^{\mathrm{sc}}},\\
C_{\mathrm{pile}}(\sigma)\ \longmapsto\ C_{\mathrm{pile}}(\sigma+2/M)
\le (1+2/M)^4\,C_{\mathrm{pile}}(\sigma)\le 2\,C_{\mathrm{pile}}(\sigma)
\quad\text{for } M\ge 12 .
\end{gathered}
\end{equation}
Only $C_w$ and $C_{\mathrm{pile}}$ change, each by the bounded factor displayed; no floor and no
exponent changes.
\end{lemma}

\begin{proof}
By \cite[(2.41)(i)]{B-CMP119} the localised term ``depends on $U_k$, $A$ restricted to $X$''; the
lemma treats the case in which the read set $S_v$ is taken to be the union of the cells meeting
$X$ rather than $X$.

We first bound the scaled distance. Let $y\in S_v$. Then $y$ lies in a cell that meets $X$ at
some point $x$; this cell contains $x$. By (a), this cell lies in the local $M$-cube of $x$. By
(b), applied with $s=x$ and $\tilde{s}=y$, passing from $x$ to
$y$ inside that cube costs at most one unit of $d_{\mathrm{sc}}$. Hence
\begin{equation*}
d_{\mathrm{sc}}(y,\Omega)\ \ge\ d_{\mathrm{sc}}(x,\Omega)-1\ \ge\ d_{\mathrm{sc}}(X,\Omega)-1 .
\end{equation*}
Taking the infimum over $y\in S_v$ gives $d_{\mathrm{sc}}(S_v,\Omega)\ge d_{\mathrm{sc}}(X,\Omega)-1$.

Next we compare the response factors. In \eqref{9.2:eq-response-factor-a}, a display of
Proposition~\ref{9.2:prop-response-factor}, whose proof is incomplete at one step, the factor
$\mathsf{r}_{\varpi}$ of a set $S$ has the form $\min\{1,t(S)\}$, where
$t(S)=C\,e^{-\varpi\delta_1^{\mathrm{sc}}d_{\mathrm{sc}}(S,\Omega)}$ with a constant $C\ge0$ that
does not depend on $S$. The bound on the scaled distance therefore gives
\begin{equation*}
t(S_v)\le e^{\varpi\delta_1^{\mathrm{sc}}}\,t(X) .
\end{equation*}
For $a\ge1$ and $t\ge0$ one has $\min\{1,at\}\le a\min\{1,t\}$: if $t\le1$ the left side is at
most $at$, and if $t>1$ it is $1\le a$. Since $\min\{1,\cdot\}$ is non-decreasing, applying this
with $a=e^{\varpi\delta_1^{\mathrm{sc}}}$ and $t=t(X)$ gives
$\mathsf{r}_{\varpi}(S_v)\le e^{\varpi\delta_1^{\mathrm{sc}}}\mathsf{r}_{\varpi}(X)$. Wherever
$C_w$ bounds the read-set weight through $\mathsf{r}_{\varpi}(X)$, it is therefore replaced by
$C_we^{\varpi\delta_1^{\mathrm{sc}}}$. This is the first line of \eqref{9.2:eq-read-set-a}.

We turn to the box condition. Suppose $S_v$ meets $\Delta$ at a point $y$. A cell that is part of
$S_v$ and contains $y$ meets both $\Delta$ and $X$. A cell meeting $\Delta$ lies in $\Delta^{+}$,
since $\Delta^{+}$ is $\Delta$ with one layer of cells added. Hence $X$ meets $\Delta^{+}$.

The box $\Delta^{+}$ has side $(\sigma+2/M)M\ell_k$, so the pile per box of
\eqref{9.2:eq-pile-per-box-a} is now read with $\sigma+2/M$ in place of $\sigma$. In that display
$C_{\mathrm{pile}}(\sigma)$ depends on $\sigma$ only through the factor $(\sigma+2)^4$. Since
$\sigma\ge0$,
\begin{equation*}
\frac{\sigma+2+2/M}{\sigma+2}\;=\;1+\frac{2}{M(\sigma+2)}\;\le\;1+\frac2M,
\end{equation*}
and therefore $C_{\mathrm{pile}}(\sigma+2/M)\le(1+2/M)^4C_{\mathrm{pile}}(\sigma)$. For $M\ge12$,
\begin{equation*}
(1+2/M)^4\ \le\ (7/6)^4\ =\ \frac{2401}{1296}\ <\ 2 .
\end{equation*}
This is the second line of \eqref{9.2:eq-read-set-a}.

The condition $M\ge12$ is implied by the floor $M\ge L^4$ of the order of choice of the
constants (Definitions~\ref{2.3:not-stages-first}, \ref{2.3:not-stages-middle}
and~\ref{2.3:not-stages-last}), because $L$ is an odd integer greater than one, so that
$L^4\ge81$; that condition is a floor, not a cap in the sense of Definition~\ref{2.3:def-cap}.
\end{proof}

\begin{lemma}[Block reach and the shell lemma]\label{9.2:lem-shell-lemma}
Let $0\le j\le k$, let $\square=\hat{\Delta}\in\pi_k$ be the block of the block-sum lemma, and put
$s_0:=M\ell_k$, the side of the blocks of $\pi_k$.
Write $\operatorname{dist}_\infty$ for the distance on $\mathbb{T}_m$ induced by the maximum norm.
Assume the following.
\begin{itemize}
\item[(a)] \emph{Block dichotomy.} Take the block $\square$ and the fattened block
$\tilde{\square}\supset\square$ of \cite[p.~271]{B-CMP109}. Take also the case distinction
made there according to whether a level-$j$ domain $X$ satisfies $X\cap\tilde{\square}=\emptyset$ or
$X\cap\tilde{\square}\neq\emptyset$. With these, the units read at $\square$ without positional decay are
exactly those whose domain meets $\tilde{\square}$. These units form the \emph{touching class} at
$\square$, and the remaining units form the \emph{receding class}. The fattened block $\tilde{\square}$ is
a box of side $\hat{\sigma}M\ell_k$ containing $\square$, where $\hat{\sigma}$ is the fattened-block side
of the block-sum lemma. No upper bound on $\hat{\sigma}$ is assumed; conclusion (ii) records the bound
$\hat{\sigma}\le3$ in the case where $\tilde{\square}$ is $\square$ with its neighbouring blocks.
\item[(b)] \emph{Receding factor.} Let $v$ be a unit of the correlation theorem belonging to the receding
class at $\square$, let $p$ be a response piece of $v$ at $\square$, and let $S_v$ be the read set of
$v$. Then the piece carries the factor
\begin{equation*}
\eta_p=e^{-\delta_1M\operatorname{dist}(\square,S_v)} ,
\end{equation*}
and this factor is read only at the two sites at which the correlation theorem states it. Here
$\delta_1>0$ is the decay rate per lattice unit of level $k$ (distinct from the response decay rate
$\delta_1^{\mathrm{sc}}$). The quantity
$\operatorname{dist}(\square,S_v)$ is the distance of the correlation theorem's decay factor, counted in
blocks of side $s_0$, so that $\delta_1M$ is the rate per block. The number $M$ satisfies
$\delta_1M\ge\kappa_1$, which is the floor on $M$ under which the correlation theorem is stated. The
charge of the response piece $p$ at $\square$ is $\theta_k\mathfrak{N}(v)\eta_p^{1/2}$.
\end{itemize}
Then the following hold.
\begin{itemize}
\item[(i)] \emph{Shell lemma.} Let $\Delta\in\pi_k$, let $a\ge 0$ be any function of the level-$j$
domains, and put $\Pi(\Delta'):=\sum_{X\cap\Delta'\neq\emptyset}a(X)$ for $\Delta'\in\pi_k$. Then for
every $\xi>0$, with $\mu:=e^{-\xi}$,
\begin{equation}\label{9.2:eq-shell-lemma-a}
\begin{gathered}
\sum_X e^{-\xi\operatorname{dist}_\infty(X,\Delta)/s_0}\,a(X)
\le C_{\mathrm{shell}}(\xi)\,\sup_{\Delta'\in\pi_k}\Pi(\Delta'),\\
C_{\mathrm{shell}}(\xi):=\sum_{n\ge0}(2n+3)^4e^{-\xi n}\\
\le 81\sum_{n\ge0}(n+1)^4\mu^n=\frac{81\,(1+11\mu+11\mu^2+\mu^3)}{(1-\mu)^{5}} .
\end{gathered}
\end{equation}
\item[(ii)] \emph{Touching class.} The touching class at $\square$ is controlled by the bound
\eqref{9.2:eq-pile-per-box-a} on the pile per box, applied to the box $\tilde{\square}$, that is, at the
factor $(\hat{\sigma}+2)^4$. If $\tilde{\square}$ is $\square$ together with its neighbouring blocks of
$\pi_k$, then $\hat{\sigma}\le3$ and the factor is at most $5^4$. For any fixed side
$\hat{\sigma}M\ell_k$ of $\tilde{\square}$ the same bound holds with the factor $(\hat{\sigma}+2)^4$.
Moreover, $\hat{\sigma}$ enters the delivery constant $C_T$ only through this factor.
\item[(iii)] \emph{Receding class.} For a response piece at a block separated from the read set by $n$
blocks of $\pi_k$, the charge carries the factor $e^{-\frac12\delta_1Mn}$. The number
\begin{equation*}
\kappa'':=\frac14\,\delta_1M
\end{equation*}
satisfies $\kappa''\ge\frac14\kappa_1>0$. It is the rate of the positional weight
$e^{-\kappa''\operatorname{dist}_\infty(X,\square)/(M\ell_k)}$ with which the receding class enters the
pile.
\item[(iv)] \emph{Total at a block.} Let $a\ge0$ be a function of the level-$j$ domains, $a(X)$ being the
contribution of $X$ to the pile per box, such that $\sum_{X\cap B\neq\emptyset}a(X)\le\Pi_j(B)$ for every
box $B$ to which \eqref{9.2:eq-pile-per-box-a} is applied, in particular for $B=\tilde{\square}$ and for
every $B=\Delta'\in\pi_k$. Let $s$, $c$, $\delta_T$,
$C_w$ and $f_j$ be as in \eqref{9.2:eq-pile-per-box-a}, and let $K_0^{(2)}$ be the constant of
\eqref{9.2:eq-boundary-families-b}. Then
\begin{equation}\label{9.2:eq-shell-lemma-1}
\begin{aligned}
&\sum_{X\cap\tilde{\square}\neq\emptyset}a(X)
+\sum_{X\cap\tilde{\square}=\emptyset}e^{-\kappa''\operatorname{dist}_\infty(X,\square)/s_0}\,a(X)\\
&\qquad\le\bigl[(\hat{\sigma}+2)^4+81\,C_{\mathrm{shell}}(\kappa'')\bigr]
K_0^{(2)}C_w\,e^{-c\delta_T\check{R}_k}\,2^{-s}f_j .
\end{aligned}
\end{equation}
When the read sets are the sets of cells meeting the domains, the right-hand side is further multiplied
by the read-set factor $2e^{\varpi\delta_1^{\mathrm{sc}}}$ of \eqref{9.2:eq-read-set-a}.
\end{itemize}
\end{lemma}

\begin{proof}
\emph{Proof of (i).} Fix $\Delta\in\pi_k$. Every $\Delta'\in\pi_k$ can be written as
$\Delta'=\Delta+s_0\nu$ with $\nu\in\mathbb{Z}^4$. For each $\Delta'$ we fix one such $\nu$ that realises the
distance on the torus. The blocks are cubes of side $s_0$. In each coordinate direction $i$ the two
projections are therefore intervals of length $s_0$ at offset $s_0\nu_i$, and their distance is
$s_0(|\nu_i|-1)_+$. Taking the maximum over $i$ gives
\begin{equation*}
\operatorname{dist}_\infty(\Delta',\Delta)=s_0\,(|\nu|_\infty-1)_+ .
\end{equation*}
Hence $n(\Delta'):=\operatorname{dist}_\infty(\Delta',\Delta)/s_0$ is a non-negative integer.

If $n(\Delta')=n\ge1$, then $|\nu|_\infty=n+1$. If $n(\Delta')=0$, then $|\nu|_\infty\le1$. In both cases
$|\nu|_\infty\le n+1$. Since $\Delta'\mapsto \nu$ is injective, we obtain
\begin{equation*}
\#\{\Delta'\in\pi_k:n(\Delta')=n\}\le\#\{\nu\in\mathbb{Z}^4:|\nu|_\infty\le n+1\}=(2n+3)^4 .
\end{equation*}

Now let $X$ be a level-$j$ domain. It is a finite set, so some $x\in X$ realises
$\operatorname{dist}_\infty(X,\Delta)$. Let $\Delta'\in\pi_k$ be the block containing $x$. Then
\begin{equation*}
\operatorname{dist}_\infty(\Delta',\Delta)\le\operatorname{dist}_\infty(x,\Delta)
=\operatorname{dist}_\infty(X,\Delta).
\end{equation*}
Put $n_X:=\min\{n(\Delta'):\Delta'\cap X\neq\emptyset\}$. The block containing $x$ meets $X$, so
$n_X\le\operatorname{dist}_\infty(X,\Delta)/s_0$. Choose one block $\Delta'_X$ that meets $X$ and has
$n(\Delta'_X)=n_X$.

We now justify each step of the chain below.
\begin{itemize}
\item The first inequality holds because the exponential is decreasing, $a\ge0$ and
$n_X\le\operatorname{dist}_\infty(X,\Delta)/s_0$.
\item The equality comes from grouping the domains $X$ according to $\Delta'_X$.
\item The second inequality holds because every $X$ with $\Delta'_X=\Delta'$ meets $\Delta'$ and
$a\ge0$, so the inner sum is at most $\Pi(\Delta')$.
\item The last inequality is the count of blocks with $n(\Delta')=n$.
\end{itemize}
The chain reads
\begin{align*}
\sum_Xe^{-\xi\operatorname{dist}_\infty(X,\Delta)/s_0}a(X)
&\le\sum_Xe^{-\xi n_X}a(X)
=\sum_{\Delta'\in\pi_k}e^{-\xi n(\Delta')}\sum_{X:\,\Delta'_X=\Delta'}a(X)\\
&\le\sum_{\Delta'\in\pi_k}e^{-\xi n(\Delta')}\,\Pi(\Delta')
\le\sum_{n\ge0}e^{-\xi n}(2n+3)^4\sup_{\Delta'\in\pi_k}\Pi(\Delta').
\end{align*}
This is the first line of \eqref{9.2:eq-shell-lemma-a}.

We now derive the closed form. For $n\ge0$ we have $2n+3\le3n+3$, hence $(2n+3)^4\le81(n+1)^4$. This
gives the inequality $C_{\mathrm{shell}}(\xi)\le81\sum_{n\ge0}(n+1)^4\mu^n$ in
\eqref{9.2:eq-shell-lemma-a}.

Since $\xi>0$, we have $0<\mu<1$, and the series $\sum_{n\ge0}(n+1)^4\mu^n$ converges absolutely.
Multiplying it by the polynomial $(1-\mu)^5=\sum_{i=0}^5(-1)^i\binom5i\mu^i$ gives the power series
$\sum_{n\ge0}c_n\mu^n$ with
\begin{equation*}
c_n=\sum_{i=0}^{\min\{n,5\}}(-1)^i\binom5i(n+1-i)^4 .
\end{equation*}
The first coefficients are
\begin{equation*}
c_0=1,\quad c_1=16-5=11,\quad c_2=81-80+10=11,\quad c_3=256-405+160-10=1 .
\end{equation*}
At $n=4$ the term with $i=5$, namely $-\binom55(n-4)^4$, vanishes, so $c_4$ is already the full
alternating sum. Computing directly,
\begin{equation*}
c_4=625-1280+810-160+5=0 .
\end{equation*}
For $n\ge5$, the coefficient $c_n=\sum_{i=0}^5(-1)^i\binom5i(n+1-i)^4$ is the fifth backward finite
difference, at $n$, of the quartic $n\mapsto(n+1)^4$. That polynomial has degree four, so $c_n=0$.
Therefore $(1-\mu)^5\sum_{n\ge0}(n+1)^4\mu^n=1+11\mu+11\mu^2+\mu^3$, which is the equality in
\eqref{9.2:eq-shell-lemma-a}.

\emph{Proof of (ii).} By hypothesis (a), the units read at $\square$ without positional decay are exactly
those whose domain meets $\tilde{\square}$. Their contribution is therefore the pile on the box
$\tilde{\square}$, which has side $\hat{\sigma}M\ell_k$. The bound \eqref{9.2:eq-pile-per-box-a} can be
applied to a box of side $\sigma M\ell_k$. It bounds the pile on such a box by
$C_{\mathrm{pile}}(\sigma)e^{-c\delta_T\check{R}_k}2^{-s}f_j$, with
$C_{\mathrm{pile}}(\sigma)=(\sigma+2)^4K_0^{(2)}C_w$. At $\sigma=\hat{\sigma}$ this is the factor
$(\hat{\sigma}+2)^4$.

Suppose that the fattened block $\tilde{\square}$ of \cite[p.~271]{B-CMP109} is the block $\square$
together with its neighbouring blocks of $\pi_k$. Then it is a box of side at most $3M\ell_k$. Hence
$\hat{\sigma}\le3$ and $(\hat{\sigma}+2)^4\le5^4$. For any other fixed side $\hat{\sigma}M\ell_k$ the same
application holds, and only the factor $(\hat{\sigma}+2)^4$ changes.

\emph{Proof of (iii).} Let the block be separated from the read set $S_v$ by $n$ blocks of $\pi_k$. Then
$\operatorname{dist}(\square,S_v)\ge n$ in units of $s_0$. By hypothesis (b),
\begin{equation*}
\eta_p^{1/2}=e^{-\frac12\delta_1M\operatorname{dist}(\square,S_v)}\le e^{-\frac12\delta_1Mn}.
\end{equation*}
So the charge $\theta_k\mathfrak{N}(v)\eta_p^{1/2}$ carries the factor $e^{-\frac12\delta_1Mn}$.

The block sum over pieces whose constant is $C_A^{(1/2)}$ may already use this decay. To cover that case,
we split the exponent $\frac12\delta_1Mn$ in halves once. One half, $\kappa''n$ with
$\kappa'':=\frac14\delta_1M$, is kept for the positional weight of the pile. The other half is left to
that block sum. From $\delta_1M\ge\kappa_1$ we get $\kappa''\ge\frac14\kappa_1$, which is positive because
$\kappa_1>0$. This uses only $\delta_1M\ge\kappa_1$, the floor on $M$ of the correlation theorem, and no
new condition is imposed. The number $\kappa''$ is the rate of the positional weight with which the
receding class enters the pile in Proposition~\ref{9.2:prop-response-factor}, whose proof is incomplete
at one step.

\emph{Proof of (iv).} The first sum in \eqref{9.2:eq-shell-lemma-1} is at most the pile
$\Pi_j(\tilde{\square})$ on the box $\tilde{\square}$, by the hypothesis on $a$ with $B=\tilde{\square}$.
By (ii) it is at most
\begin{equation*}
(\hat{\sigma}+2)^4K_0^{(2)}C_w\,e^{-c\delta_T\check{R}_k}2^{-s}f_j .
\end{equation*}

For the second sum, enlarge it to the sum over all $X$; this is allowed because $a\ge0$. Then apply (i)
with $\Delta=\square$, with $\xi=\kappa''$ as in (iii), and with
$\Pi(\Delta'):=\sum_{X\cap\Delta'\neq\emptyset}a(X)$, which satisfies $\Pi(\Delta')\le\Pi_j(\Delta')$ for
every $\Delta'\in\pi_k$ by the hypothesis on $a$. Every block $\Delta'\in\pi_k$
is a box of side $M\ell_k$. Hence \eqref{9.2:eq-pile-per-box-a} at $\sigma=1$ gives
\begin{equation*}
\sup_{\Delta'\in\pi_k}\Pi(\Delta')\le\sup_{\Delta'\in\pi_k}\Pi_j(\Delta')
\le C_{\mathrm{pile}}(1)\,e^{-c\delta_T\check{R}_k}2^{-s}f_j ,
\qquad C_{\mathrm{pile}}(1)=81K_0^{(2)}C_w .
\end{equation*}
So the second sum is at most
\begin{equation*}
81\,C_{\mathrm{shell}}(\kappa'')K_0^{(2)}C_w\,e^{-c\delta_T\check{R}_k}2^{-s}f_j .
\end{equation*}
Adding the two bounds gives \eqref{9.2:eq-shell-lemma-1}.

Suppose now that the read sets are the sets of cells meeting the domains. Then
Lemma~\ref{9.2:lem-read-set} multiplies each contribution by at most the factor
$2e^{\varpi\delta_1^{\mathrm{sc}}}$ of \eqref{9.2:eq-read-set-a}. The total is therefore multiplied by
the same factor.

In both classes the counting is carried out in blocks of side $M\ell_k$. The factor $2^{-s}f_j$ is the same
as in \eqref{9.2:eq-pile-per-box-a}, and only the constant depends on $\hat{\sigma}$ and $\kappa''$.
\end{proof}

\begin{proposition}[Geometric weight of old boundary terms]\label{9.2:prop-old-term-weight}
Consider the frozen step $\mathbf{T}$ from level $k$ to level $k+1$ on one lattice and one scaffold, for
any label admitted in Theorem~\ref{9.2:thm-pile-bound}. Let $\mathsf{b}=(\mathbf{B}^{(j)})_{j\le k}$ be the
boundary-family slot of the one-lattice Lipschitz step, and let $[\mathfrak{h},\tilde{\mathfrak{h}}]\subset
\mathfrak{K}_k^{+}$ be its segment. Put
\begin{equation*}
F:=\bigl(\delta\mathbf{B}^{(j)}(X)\bigr)_{j\le k,\ X\in\mathbf{D}_j},
\end{equation*}
the family of the differences along $[\mathfrak{h},\tilde{\mathfrak{h}}]$ of the $\mathbf{T}$-born members
$\mathbf{B}^{(j)}$, $j\le k$, of $\mathsf{b}$. Here $X$ ranges over the localisation domains $\mathbf{D}_j$ of the
stored $\mathbf{T}$-born boundary terms of level $j$ alive when read at the step, as in
Definition~\ref{9.2:def-boundary-families}. Moreover
\begin{equation*}
|\delta\mathbf{B}^{(j)}(X,\cdot)|\le\delta^{\mathbf{B}}_j e^{-\kappa d_j(X)},\qquad
f_j:=\delta^{\mathbf{B}}_j:=\|\delta\mathbf{B}^{(j)}\|.
\end{equation*}
The differences are read as marked units of the summation of boundary terms at finer zones of the
correlation theorem, by the lemma on marked activities of the one-lattice Lipschitz step. Assume:
\begin{itemize}
\item clauses (ii) and (iii) of the pile bound for stored boundary terms,
Theorem~\ref{9.2:thm-pile-bound}, displayed in \eqref{9.2:eq-pile-bound-b};
\item the bound of the correlation theorem for the units of the summation of boundary terms at finer
zones at the step $\mathbf{T}$, with its piece-sum constant $C_A^{(1/2)}$;
\item the lemma on marked activities of the one-lattice Lipschitz step, which is linear in the unit and
bounds the marked activity with the prefactor $C_U\theta_k$;
\item clauses (R4-i)--(R4-iii) of Proposition~\ref{9.2:prop-response-factor}, displayed in
\eqref{9.2:eq-response-factor-b}, whose proof is incomplete at one step, together with the profile bounds
(a)--(c) for response pieces stated there;
\item Lemma~\ref{9.2:lem-shell-lemma}, with its receding rate $\kappa''>0$ and its fattened block side
$\hat{\sigma}$;
\item Lemma~\ref{9.2:lem-read-set}.
\end{itemize}
Then, at every block $\hat{\Delta}$ of the block-sum lemma of the one-lattice Lipschitz step, at the site of
the step $\mathbf{T}$, the marked activity of these members satisfies
\begin{equation}\label{9.2:eq-old-term-weight-1}
\text{marked activity at }\hat{\Delta}\ \le\ C_U\theta_k\sum_{j\le k}\mathfrak{a}^{T}_{k,j}\,
\delta^{\mathbf{B}}_j ,
\end{equation}
where
\begin{equation}\label{9.2:eq-old-term-weight-a}
\begin{gathered}
\mathfrak{a}^{T}_{k,j}:=C_T\,e^{-c\delta_T\check{R}_k}\,2^{-(k-j)}\le C_T\,2^{-(k-j)},\\
C_T:=n_{\mathrm{sp}}\,C_A^{(1/2)}\,\bigl[(\hat{\sigma}+2)^4+81\,C_{\mathrm{shell}}(\kappa'')\bigr]\,
2e^{\varpi\delta_1^{\mathrm{sc}}}\,K_0^{(2)}C_w .
\end{gathered}
\end{equation}
The symbols in these displays are the following.
\begin{itemize}
\item $c$ is the exponent of Theorem~\ref{9.2:thm-pile-bound} in the grade in which that theorem is
applied. It equals $13/2$, $13/2$, $3/2$ in its first, second and third grade. For weakly anchored
members it equals $11/2$, $11/2$, $3/2$.
\item $n_{\mathrm{sp}}$ is the number of kinds of boundary term stored at one level under the one norm
$\delta^{\mathbf{B}}_j$, with $n_{\mathrm{sp}}=1$ if the slot is one family.
\item $C_U$ and $\theta_k$ are the constant and the activity scale of the one-lattice Lipschitz step.
\item $\delta_T$ and $C_w$ are as in \eqref{9.2:eq-response-factor-a}, in which the fraction $\varpi$ of the
decay and the response decay rate $\delta_1^{\mathrm{sc}}$ enter; $\check{R}_k$ is the radius of the scaffold
at level $k$.
\item $C_{\mathrm{shell}}(\kappa'')$ and $\hat{\sigma}$ are as in \eqref{9.2:eq-shell-lemma-a}.
\item $K_0^{(2)}$ is as in \eqref{9.2:eq-boundary-families-b}.
\end{itemize}
The constant $C_T$ depends on $d,\kappa,\delta_1^{\mathrm{sc}},C_b,\varpi,\hat{\sigma},\kappa'',n_{\mathrm{sp}}$
and $C_A^{(1/2)}$ only. It does not depend on $K$, $k$, $j$, the age, $N$, $N_0$, the torus, or the position.

The proposition concerns the $\mathbf{T}$-born members $\mathbf{B}^{(j)}$, $j\le k$, of $\mathsf{b}$ only. The
$\mathbf{R}$-born members $\mathbf{B}'^{(j)}$, $j\le k$, are not covered by it. These members enter the slot
through the boundary term $\mathbf{B}^{(k+1)}$ near $\partial\Lambda_{k+1}\cup\mathbf{B}'$, together with
their descendants.
\end{proposition}

\begin{proof}
\emph{Reduction to one kind at one level.} By hypothesis the marked activity is linear in the unit.
The marked activity at $\hat{\Delta}$ of the family $F$ is therefore the sum, over the levels $j\le k$ and
over the $n_{\mathrm{sp}}$ kinds of boundary term stored at level $j$, of the marked activities of the
corresponding parts of $F$. Each kind at level $j$ is a level-$j$ family in the sense of
Definition~\ref{9.2:def-boundary-families}, with size $f_j=\delta^{\mathbf{B}}_j$ and decay
$e^{-\kappa d_j(X)}$. Clauses (ii) and (iii) of Theorem~\ref{9.2:thm-pile-bound} are assumed for this
family with the anchor required there, or with its weak form under the rider (R-weak) of
\eqref{9.2:eq-pile-bound-b}; in the second case $c$ takes the second list of values stated above. Fix such a
kind and a level $j$, and put $s:=k-j$. Relative to the block
$\hat{\Delta}$, divide its units into the touching class and the receding class of
Lemma~\ref{9.2:lem-shell-lemma}.

\emph{Touching class.} The touching class is estimated by (R4-i) of \eqref{9.2:eq-response-factor-b},
which is part of Proposition~\ref{9.2:prop-response-factor}, whose proof is incomplete at one step. This
reduces it to the pile $\Pi_j$ of Theorem~\ref{9.2:thm-pile-bound}, taken at side $\sigma=\hat{\sigma}$.
Clause (ii) of that theorem, in the form \eqref{9.2:eq-pile-per-box-a}, bounds this pile at
$\sigma=\hat{\sigma}$ by
\begin{equation}\label{9.2:eq-old-term-weight-2}
C_{\mathrm{pile}}(\hat{\sigma})\,e^{-c\delta_T\check{R}_k}\,2^{-s}f_j,\qquad
C_{\mathrm{pile}}(\hat{\sigma})=(\hat{\sigma}+2)^4K_0^{(2)}C_w .
\end{equation}

\emph{Receding class.} The receding class is estimated by (R4-ii) of \eqref{9.2:eq-response-factor-b},
again part of Proposition~\ref{9.2:prop-response-factor}, whose proof is incomplete at one step. It is then
estimated by Lemma~\ref{9.2:lem-shell-lemma}, applied with $\Pi=\Pi_j$ taken on $\pi_k$-cubes, which
carries the factor $C_{\mathrm{shell}}(\kappa'')$; here $\kappa''>0$ as stated in that lemma. The pile on
one $\pi_k$-cube is bounded by clause (ii) of Theorem~\ref{9.2:thm-pile-bound} at $\sigma=1$. By
\eqref{9.2:eq-pile-per-box-a}, and since $(1+2)^4=81$, this gives
\begin{equation}\label{9.2:eq-old-term-weight-3}
C_{\mathrm{shell}}(\kappa'')\,C_{\mathrm{pile}}(1)\,e^{-c\delta_T\check{R}_k}\,2^{-s}f_j,\qquad
C_{\mathrm{pile}}(1)=81\,K_0^{(2)}C_w .
\end{equation}
The sum of the two class bounds \eqref{9.2:eq-old-term-weight-2} and \eqref{9.2:eq-old-term-weight-3} is
\begin{equation*}
\bigl[(\hat{\sigma}+2)^4+81\,C_{\mathrm{shell}}(\kappa'')\bigr]K_0^{(2)}C_w\,
e^{-c\delta_T\check{R}_k}\,2^{-s}f_j .
\end{equation*}

\emph{Read set.} By (R4-iii) of \eqref{9.2:eq-response-factor-b}, part of
Proposition~\ref{9.2:prop-response-factor}, whose proof is incomplete at one step, the unit is read on its
read set rather than on its domain $X$. By Lemma~\ref{9.2:lem-read-set}, passing from $X$ to the read set
costs the factor $e^{\varpi\delta_1^{\mathrm{sc}}}$ in the per-unit response factor
$\mathsf{r}_{\varpi}$.

\emph{Sum over kinds and levels.} The lemma on marked activities of the one-lattice Lipschitz step, linear
in the unit, contributes the prefactor $C_U\theta_k$. The bound of the correlation theorem for the units of
the summation of boundary terms at finer zones at the step $\mathbf{T}$ contributes the piece-sum constant
$C_A^{(1/2)}$. The factors for one kind of boundary term at level $j$ are therefore the following:
\begin{itemize}
\item $C_U\theta_k$, from the lemma on marked activities, and $C_A^{(1/2)}$, from the correlation theorem;
\item the read-set factor $e^{\varpi\delta_1^{\mathrm{sc}}}$;
\item the bracket of the two classes;
\item $K_0^{(2)}C_w$;
\item $e^{-c\delta_T\check{R}_k}\,2^{-(k-j)}\,\delta^{\mathbf{B}}_j$.
\end{itemize}
Summing over the $n_{\mathrm{sp}}$ kinds of boundary term at each level contributes the factor
$n_{\mathrm{sp}}$, and summing over the levels $j\le k$ gives the sum in \eqref{9.2:eq-old-term-weight-1}.
The product of the factors listed above and of $n_{\mathrm{sp}}$ is
\begin{equation*}
C_U\theta_k\cdot\tfrac12\,C_T\,e^{-c\delta_T\check{R}_k}\,2^{-(k-j)}\,\delta^{\mathbf{B}}_j ,
\end{equation*}
with $C_T$ as in \eqref{9.2:eq-old-term-weight-a}. All factors of $C_T$ are non-negative, so
$\tfrac12C_T\le C_T$, and \eqref{9.2:eq-old-term-weight-1} holds with $\mathfrak{a}^{T}_{k,j}$ and $C_T$ as in
\eqref{9.2:eq-old-term-weight-a}.

\emph{The inequality for $\mathfrak{a}^{T}_{k,j}$ and the dependence of $C_T$.} We have $c>0$ and
$\delta_T>0$, and $\check{R}_k\ge2$ among the floors of Theorem~\ref{9.2:thm-pile-bound}. Hence
$e^{-c\delta_T\check{R}_k}\le1$, which gives $\mathfrak{a}^{T}_{k,j}\le C_T2^{-(k-j)}$. By its definition
\eqref{9.2:eq-old-term-weight-a}, $C_T$ is built from the following quantities only:
\begin{itemize}
\item $n_{\mathrm{sp}}$ and $C_A^{(1/2)}$;
\item $\hat{\sigma}$ and $C_{\mathrm{shell}}(\kappa'')$;
\item $\varpi$ and $\delta_1^{\mathrm{sc}}$;
\item $K_0^{(2)}$;
\item $C_w=C_b^{\varpi}e^{\varpi\delta_1^{\mathrm{sc}}}$.
\end{itemize}
This gives the stated list of dependences. The indices $k$ and $j$ enter $\mathfrak{a}^{T}_{k,j}$ only
through $e^{-c\delta_T\check{R}_k}$ and $2^{-(k-j)}$; they do not enter $C_T$.
\end{proof}

\begin{remark}[The boundary-family entry of the Lipschitz step]\label{9.2:rem-boundary-entry}
Fix one lattice and one scaffold, and consider the frozen renormalisation step $k\to k+1$ within the
one-lattice Lipschitz step on the segment $[\mathfrak{h},\tilde{\mathfrak{h}}]\subset\mathfrak{K}_k^{+}$.
Assume the following statements.
\begin{itemize}
\item[(H-a)] Proposition~\ref{9.2:prop-old-term-weight} (geometric weight of old boundary terms),
whose weight is displayed in \eqref{9.2:eq-old-term-weight-a}.
\item[(H-b)] The following statements of the one-lattice Lipschitz step.
\begin{itemize}
\item[(1)] Its boundary-family slot is
$\mathsf{b}=(\mathbf{B}^{(j)})_{j\le k}$. The norm of the slot is the weighted sup-norm
\begin{equation*}
\|\mathbf{B}^{(j)}\|:=\sup_{X\in\mathbf{D}_j}e^{\kappa d_j(X)}|\mathbf{B}^{(j)}(X,\cdot)|,
\end{equation*}
where $|\mathbf{B}^{(j)}(X,\cdot)|$ is the supremum over the analyticity domain
$U_j(X,\alpha_0,\alpha_1)$ of \cite[(2.41)(ii)]{B-CMP119}. In this norm the class of each member,
the bound of \cite[(2.40)--(2.42)]{B-CMP119}, reads $\|\mathbf{B}^{(j)}\|\le B_0$, that is,
$|\mathbf{B}^{(j)}(X,\cdot)|\le B_0e^{-\kappa d_j(X)}$ for every $X\in\mathbf{D}_j$. On
$\mathfrak{K}_k^{+}$ the constant $B_0$ of this class is replaced by $\tfrac98 B_0$.
\item[(2)] The terms $\mathbf{B}$ and $\mathbb{C}^{(i)}$ are read by the sup-norm in the fourth class
of units of the correlation theorem, the class formed in its list of units by the stored boundary
terms $\mathbf{B}$ and the complex-type terms $\mathbb{C}^{(i)}$.
\item[(3)] Marks are set at the norm of a difference: the marking lemma of the one-lattice
Lipschitz step is applied to each family with its size parameter equal to the norm of the
difference of the stored terms along the segment, as it is applied to the polymer terms with the
norm $\delta^{P}_j$ of the difference of the level-$j$ polymer activities.
\item[(4)] On the scaffold $g_k\le\gamma$ for every $k$, and the activity scale is
$\theta_k=\theta(g_k)$. Every entry of the one-lattice Lipschitz step that depends on the
coupling tends to $0$ as $\gamma\downarrow0$ at fixed $(L,M,\mathfrak{p},r)$; for the activity
scale, $\sup_{g\le\gamma}\theta(g)\to0$.
\item[(5)] Constants growing with $L$, such as $(6L)^4$ and $C_U$, sit on the small side of a
ceiling on $g$.
\item[(6)] Its lemmas that resum marked members require the activity at a block of each marked
member to lie inside the activity room $\varrho_{\Sigma}$ of the block-sum lemma; this is imposed
through the cutoff radius $R_c$ of that step, defined there as the minimum of two quantities, the
second of which is $\varrho_{\Sigma}$ divided by twice a level-$k$ quantity of that step, and it
enters as a ceiling on $g$.
\item[(7)] Its bound contains the boundary-family entry $\sum_{j\le k}2^{-(k-j)}\delta^{\mathbf{B}}_j$.
\end{itemize}
\item[(H-c)] Ba{\l}aban's description of the stored boundary terms.
\cite[(2.41)(ii), (iv), (2.42)]{B-CMP119} states that $\mathbf{B}^{(j)}(X,\cdot)$ has an extension
to an analytic function on the space $U_j(X,\alpha_0,\alpha_1)$, and that it satisfies
\begin{equation*}
|\mathbf{B}^{(j)}(X,\cdot)|\le B_0\exp(-\kappa d_j(X)).
\end{equation*}
\end{itemize}
Then the following hold.

\smallskip\noindent
\emph{(a) The entry.} The old boundary terms $\mathbf{B}^{(j)}$, namely the $\mathbf{T}$-born
members of the slot $\mathsf{b}$, enter the bound of the one-lattice Lipschitz
step as units of the fourth class of the correlation theorem. They are read at the
$\mathbf{T}$-site, that is, at the site of the renormalisation step $\mathbf{T}$ (as opposed to the
large-field operation $\mathbf{R}$), at the Cauchy radius
\begin{equation*}
\rho_A=\mathsf{w}_T/\eta_p^{1/2},
\end{equation*}
with the per-unit response factor $\mathsf{r}_{\varpi}$ of \eqref{9.2:eq-response-factor-b}, whose
proof is incomplete at one step. They
are piled per block by the pile bound (Theorem~\ref{9.2:thm-pile-bound}) together with the shell
lemma (Lemma~\ref{9.2:lem-shell-lemma}). Under the one-sided floors of
Theorem~\ref{9.2:thm-pile-bound} and the room condition \eqref{9.2:eq-boundary-entry-a}, this
supplies the weight per level of age
\begin{equation*}
\mathfrak{a}^{T}_{k,j}=C_Te^{-c\delta_T\check{R}_k}2^{-(k-j)}\le C_T2^{-(k-j)},
\end{equation*}
with the constants $c$ and $C_T$ of \eqref{9.2:eq-old-term-weight-a}.
The geometric bound \cite[(2.47)]{B-CMP119} is not used. The $\mathbf{R}$-born terms
$\mathbf{B}'^{(j)}$, $j\le k$, are folded into the slot by the definition of $\mathbf{B}^{(k+1)}$
near $\partial\Lambda^{\mathbf{T}}_{k+1}\cup\mathbf{B}'$ ($\Lambda^{\mathbf{T}}_{k+1}$ the action
domain at level $k+1$); for them this weight is not supplied here. The marking
of $\delta\mathbf{B}^{(j)}$, its groups, its units and its vertex are as for the polymer terms, by
the lemma on types and the marking device of the one-lattice Lipschitz step, and
the terms $\mathbf{B}'$ are treated with the $\mathbf{R}$-terms, as in that step.

\smallskip\noindent
\emph{(b) The entry of the increment bound.} The entry
\begin{equation}\label{9.2:eq-boundary-entry-1}
\sum_{j\le k}2^{-(k-j)}\delta^{\mathbf{B}}_j
\end{equation}
of the one-lattice Lipschitz step stands for the $\mathbf{T}$-born members of $\mathsf{b}$. For the
folded $\mathbf{R}$-born members $\mathbf{B}'^{(j)}$, $j\le k$, neither the pile bound nor this
remark supplies the marked-activity bound. With this reading the bound of the one-lattice
Lipschitz step is to be read with the following two modifications.
\begin{itemize}
\item[(1)] The constant $K_{\mathrm{CE}}C_U$ of the one-lattice Lipschitz step in front of
\eqref{9.2:eq-boundary-entry-1} becomes $K_{\mathrm{CE}}C_UC_T$
(equivalently, $C_T$ is displayed beside the sum). The factor $e^{-c\delta_T\check{R}_k}\le1$ is
dropped in \eqref{9.2:eq-boundary-entry-1} and kept in \eqref{9.2:eq-pile-bound-b}.
\item[(2)] The following conditions are added to those defining $\gamma_{2R}$:
  \begin{itemize}
  \item[$\cdot$] the grade condition of \eqref{9.2:eq-pile-bound-a} of
  Theorem~\ref{9.2:thm-pile-bound} at the grade of the family, together, at the third grade, with
  the three further conditions displayed with it;
  \item[$\cdot$] that theorem's condition on labels with non-completed attempts in force;
  \item[$\cdot$] the room condition \eqref{9.2:eq-boundary-entry-a}.
  \end{itemize}
\end{itemize}

\smallskip\noindent
\emph{(c) Size.} For every $X\in\mathbf{D}_j$,
\begin{equation}\label{9.2:eq-boundary-entry-2}
|\delta\mathbf{B}^{(j)}(X,\cdot)|\le\delta^{\mathbf{B}}_je^{-\kappa d_j(X)},
\qquad \delta^{\mathbf{B}}_j\le\tfrac94B_0\ \text{on }\mathfrak{K}_k^{+}.
\end{equation}
Thus $\delta\mathbf{B}$ is a level-$j$ boundary family in the sense of
Definition~\ref{9.2:def-boundary-families}, with $f_j=\delta^{\mathbf{B}}_j$.

\smallskip\noindent
\emph{(d) Room.} Let $\mathfrak{b}_4\in\,]0,1[$ be the share of the activity room
$\varrho_{\Sigma}$ of the block-sum lemma allotted to the marked $\mathbf{T}$-born member of the
fourth class, and impose
\begin{equation}\label{9.2:eq-boundary-entry-a}
\sup_k\;C_U\theta_k\cdot2C_T\cdot\tfrac94B_0\cdot e^{-c\delta_T\check{R}_k}
\le\mathfrak{b}_4\varrho_{\Sigma}.
\end{equation}
Then the marked member lies inside $\varrho_{\Sigma}$. Moreover, \eqref{9.2:eq-boundary-entry-a}
is a one-sided ceiling on $\gamma$, of the same kind as the ceiling $\gamma_{2R}$. It can be met for
every $L$, and none of $K$, $k$, $j$, the age, nor the attempt data $N_a$, $N_0(a)$ of any label
enters it.

\smallskip\noindent
\emph{(e) Block reach.} The block reach used in (a) and (d) is that of
Lemma~\ref{9.2:lem-shell-lemma}.

\smallskip\noindent
Beyond (c): the correlation theorem files the fourth class in its own norm. There a unit is its
orbit datum, the datum on which the correlation theorem files it, and the size of a unit is twice
the supremum $\mathfrak{N}^{\mathbb{C}}(\cdot)$ of that datum. This quantity dominates the
real-slice size $\mathfrak{N}$ with the same constants. It carries the strong weights
\begin{equation*}
e^{\kappa_A^{\mathrm{pr}}(g_\cdot)\|\mathcal{S}_\cdot\|}\ \text{on the newest slot},\qquad
e^{\kappa_A(k)\|\mathcal{S}_{j'}\|}\ \text{on the older slots},
\end{equation*}
where $\mathcal{S}$ denotes the strong set of a slot, and, by the amplitude bound of the
correlation theorem for its stored terms, it is bounded by
$B_0e^{-\kappa_{\mathrm{out}}d_j(X)}$, with the rate $\kappa_{\mathrm{out}}$ of the correlation
theorem. Writing the norm
of $\mathsf{b}$ in the one-lattice Lipschitz step as that filing norm changes none of the uses made
of $\mathsf{b}$, for three reasons:
\begin{itemize}
\item[$\cdot$] the strong sets belong to the geometry of the label and are common to the two norms;
\item[$\cdot$] every use of $\delta^{\mathbf{B}}_j$ reads it as a number;
\item[$\cdot$] $\kappa$ is replaced by $\kappa_{\mathrm{out}}$ in the size entry of $\delta_T$ and
in the $\kappa$-members of the floors, and these are conditions of the same kind.
\end{itemize}
The identification of the correlation theorem's $\mathfrak{N}^{\mathbb{C}}$-class with the class of
\cite[(2.41)(ii)--(2.42)]{B-CMP119}, and with the room in powers of $\tfrac12$ that
\cite{B-CMP119} provides for it, is not
proved. Neither the pile bound nor \eqref{9.2:eq-boundary-entry-1} uses it.
\end{remark}

\begin{proof}
\emph{Part (a).} The $\mathbf{T}$-born members of $\mathsf{b}$ are stored level-$j$ boundary terms
whose domains lie in $\mathbf{D}_j$. At the $\mathbf{T}$-site they are read as units of the fourth
class at the Cauchy radius $\rho_A=\mathsf{w}_T/\eta_p^{1/2}$; this is used as stated in the
correlation theorem for that class. Each unit carries the response factor
$\mathsf{r}_{\varpi}$ of \eqref{9.2:eq-response-factor-b}, whose bound is stated; proof incomplete
(it is a reading of the proof of the lemma $V^{\mathbb{C}}$ of the correlation theorem).

By Proposition~\ref{9.2:prop-old-term-weight} (hypothesis (H-a)), the pile of such units per block
is bounded by Theorem~\ref{9.2:thm-pile-bound}, in the form \eqref{9.2:eq-pile-bound-b}, together with
Lemma~\ref{9.2:lem-shell-lemma}. They are applied under:
\begin{itemize}
\item[$\cdot$] that the members form a level-$j$ boundary family, which is part (c) proved below;
\item[$\cdot$] the one-sided floors of Theorem~\ref{9.2:thm-pile-bound};
\item[$\cdot$] the room condition \eqref{9.2:eq-boundary-entry-a} imposed in part (d).
\end{itemize}
The conclusion is the weight \eqref{9.2:eq-old-term-weight-a}, that is,
$\mathfrak{a}^{T}_{k,j}=C_Te^{-c\delta_T\check{R}_k}2^{-(k-j)}$. Since $c\delta_T\check{R}_k\ge0$,
we have $e^{-c\delta_T\check{R}_k}\le1$, and so $\mathfrak{a}^{T}_{k,j}\le C_T2^{-(k-j)}$.

Nowhere in this chain is \cite[(2.47)]{B-CMP119} invoked. The chain applies only to the
$\mathbf{T}$-born members, because Proposition~\ref{9.2:prop-old-term-weight} is stated for those
members only. This is why the weight is not supplied for the folded $\mathbf{B}'^{(j)}$.

\emph{Part (b).} By (a), the $\mathbf{T}$-born part of the boundary-family entry is bounded by
$\sum_{j\le k}\mathfrak{a}^{T}_{k,j}\,\delta^{\mathbf{B}}_j$. Using $e^{-c\delta_T\check{R}_k}\le1$,
this is at most
\begin{equation*}
C_T\sum_{j\le k}2^{-(k-j)}\delta^{\mathbf{B}}_j,
\end{equation*}
which is \eqref{9.2:eq-boundary-entry-1} with the constant multiplied by $C_T$. The factor
$e^{-c\delta_T\check{R}_k}$ so dropped remains available in \eqref{9.2:eq-pile-bound-b}.

The three conditions listed in (b)(2) are exactly the hypotheses under which
Theorem~\ref{9.2:thm-pile-bound} and the room condition of part (d) are used in (a).
Each is a condition of the same kind as those defining $\gamma_{2R}$, so each joins $\gamma_{2R}$.

\emph{Part (c).} By (H-b)(1), the norm of $\mathsf{b}$ is the weighted sup-norm displayed there, in
which each member satisfies $\|\mathbf{B}^{(j)}\|\le B_0$. This norm is taken on the analyticity
domain $U_j(X,\alpha_0,\alpha_1)$ of \cite[(2.41)(ii)]{B-CMP119}, on which, by (H-c),
$\mathbf{B}^{(j)}(X,\cdot)$ is analytic and bounded by $B_0\exp(-\kappa d_j(X))$. On
$\mathfrak{K}_k^{+}$, hence
at every point of the segment $[\mathfrak{h},\tilde{\mathfrak{h}}]$, the constant is $\tfrac98B_0$.

By (H-b)(2), the fourth class of units is read by this sup-norm. By (H-b)(3), marks are set at the
norm of the difference, that is, $\delta^{\mathbf{B}}_j=\|\delta\mathbf{B}^{(j)}\|$ along
$[\mathfrak{h},\tilde{\mathfrak{h}}]$. Since the norm is a weighted sup-norm with weight
$e^{\kappa d_j(X)}$, it follows that
\begin{equation*}
|\delta\mathbf{B}^{(j)}(X,\cdot)|\le\delta^{\mathbf{B}}_je^{-\kappa d_j(X)}
\end{equation*}
for every $X\in\mathbf{D}_j$. By the triangle inequality applied to the difference of the values at
two points of the segment, each bounded by $\tfrac98B_0$ in this norm,
\begin{equation*}
\delta^{\mathbf{B}}_j\le\tfrac98B_0+\tfrac98B_0=\tfrac94B_0.
\end{equation*}
This is \eqref{9.2:eq-boundary-entry-2}. It is the size bound of
Definition~\ref{9.2:def-boundary-families} with $F=\delta\mathbf{B}$ and $f_j=\delta^{\mathbf{B}}_j$,
which is the hypothesis on the family in Theorem~\ref{9.2:thm-pile-bound}.

\emph{Part (d).} By (H-b)(6), the lemmas of the one-lattice Lipschitz step that resum marked
members need the marked member to lie inside the activity room $\varrho_{\Sigma}$ of the block-sum
lemma; that requirement is imposed there through a ceiling on $g$. We bound the activity at a block of the
marked $\mathbf{T}$-born member of the fourth class in three steps.
\begin{itemize}
\item[(1)] By Proposition~\ref{9.2:prop-old-term-weight}, the pile at a block of the level-$j$
units contributes at most
\begin{equation*}
\mathfrak{a}^{T}_{k,j}\delta^{\mathbf{B}}_j=C_Te^{-c\delta_T\check{R}_k}2^{-(k-j)}\delta^{\mathbf{B}}_j.
\end{equation*}
\item[(2)] By Lemma~\ref{9.2:lem-shell-lemma}, the units reaching a block are those counted in the
pile; the activity at a block of a marked member is its piled size times the activity scale
$C_U\theta_k$ of the one-lattice Lipschitz step, $\theta_k=\theta(g_k)$.
\item[(3)] By part (c), each unit has size at most $\tfrac94B_0$.
\end{itemize}
Summing over $j\le k$ and using
\begin{equation*}
\sum_{j\le k}2^{-(k-j)}\le\sum_{s\ge0}2^{-s}=2,
\end{equation*}
the activity is at most
\begin{equation*}
C_U\theta_k\cdot2C_Te^{-c\delta_T\check{R}_k}\cdot\tfrac94B_0.
\end{equation*}
Under \eqref{9.2:eq-boundary-entry-a} this is at most $\mathfrak{b}_4\varrho_{\Sigma}$, a share
$\mathfrak{b}_4<1$ of $\varrho_{\Sigma}$. The marked member therefore lies inside the room.

It remains to see that \eqref{9.2:eq-boundary-entry-a} is a ceiling on $\gamma$ of the stated kind.
We treat the two sides in turn.

\emph{The left side tends to $0$.}
\begin{itemize}
\item[$\cdot$] Since $\theta_k=\theta(g_k)$ with $g_k\le\gamma$, hypothesis (H-b)(4) gives
\begin{equation*}
\theta_k\le\sup_{g\le\gamma}\theta(g)\to0
\end{equation*}
as $\gamma\downarrow0$ at fixed $(L,M,\mathfrak{p},r)$, uniformly in $k$.
\item[$\cdot$] By \eqref{9.2:eq-radii-floor-b} of Definition~\ref{9.2:not-radii-floor},
$\check{R}_k\ge(\log(\gamma^{-2}-c_5))^{r_{\mathrm{pr}}}$. Since $c\delta_T>0$, this gives
\begin{equation*}
e^{-c\delta_T\check{R}_k}\le e^{-c\delta_T(\log(\gamma^{-2}-c_5))^{r_{\mathrm{pr}}}},
\end{equation*}
and the right-hand side is free of $k$ and decreases to $0$ as $\gamma\downarrow0$.
\end{itemize}

\emph{The right side is fixed.} The quantity $\mathfrak{b}_4\varrho_{\Sigma}$ does not depend on
$\gamma$.

\emph{The remaining factors.} The factors growing with $L$, $M$ or $\kappa$ are $C_U$ and the
factors $(\hat{\sigma}+2)^4$, $C_{\mathrm{shell}}$ and $(1+2/M)^4\le2$ (which holds for $M\ge11$)
carried by the constants of the pile bound and the shell lemma. By
(H-b)(5) they stand on the small
side of a ceiling on $g$. They are fixed before $\gamma$ and only lower the admissible $\gamma$.

Hence \eqref{9.2:eq-boundary-entry-a} is one-sided, and it holds for all $\gamma$ below a ceiling
depending on $(L,M,\mathfrak{p},r)$, so it can be met for every $L$. Since the supremum is over $k$
and the bound on its left side is uniform in $k$, none of $K$, $k$, $j$, the age, nor the attempt
data $N_a$, $N_0(a)$ of any label enters it.

\emph{Part (e).} This is Lemma~\ref{9.2:lem-shell-lemma}, applied as in (a) and (d).

\emph{The final paragraph of the statement.} The filing norm of the correlation theorem differs
from the norm of (H-b)(1) in three respects:
\begin{itemize}
\item[$\cdot$] the strong weights;
\item[$\cdot$] the factor two in front of $\mathfrak{N}^{\mathbb{C}}$;
\item[$\cdot$] the rate $\kappa_{\mathrm{out}}$ in place of $\kappa$.
\end{itemize}
The strong sets are determined by the label and are the same in both norms. Every step above uses
$\delta^{\mathbf{B}}_j$ only through the number bounding it. Replacing $\kappa$ by
$\kappa_{\mathrm{out}}$ alters only the size entry of $\delta_T$ and the $\kappa$-members of the
floors, which remain one-sided floors of the same kind. Parts (a)--(e) therefore hold in the
filing norm with $\kappa$ replaced by $\kappa_{\mathrm{out}}$ at those two places. They use only
that norm, and not its identification with the class of
\cite[(2.41)(ii)--(2.42)]{B-CMP119}.
\end{proof}

\begin{lemma}[Corrected age weight against the geometric weight]\label{9.2:lem-age-weight-comparison}
Let $0\le j\le k$, put $s:=k-j$ (in this lemma only), and let $\vartheta_{*}$ be an age rate with
$0<\vartheta_{*}<1$.
Let
\begin{equation}\label{9.2:eq-age-weight-comparison-1}
\mathfrak{w}^{\mathbf{B}}_{k,j}:=\sum_{i=j}^{k}2^{-(i-j)}\,\vartheta_{*}^{\,k-i}.
\end{equation}
This is the weight that the inner-read bound of the one-lattice Lipschitz step assigns to a stored
level-$j$ boundary term $\mathbf{B}^{(j)}(X)$ read by an inner operation of $\mathbf{R}_{k+1}$ at a
level $i\ge j$. Such a read carries the level-$i$ unit weight $2^{-(i-j)}$ and the age factor
$\vartheta_{*}^{\,k-i}$ from level $i$ to level $k$. Let
$C_{\mathbf{B}}^{v}=K_0M^{-4}c_L\ge0$ be the per-cell constant of the volume and filament bound, with
$K_0$ and $c_L$ constants fixed in that bound, and put
\begin{equation}\label{9.2:eq-age-weight-comparison-2}
\mathfrak{w}^{\mathbf{B},\mathrm{corr}}_{k,j}:=C_{\mathbf{B}}^{v}\,(k-j+2)\,\vartheta_{*}^{\,k-j}.
\end{equation}
\begin{enumerate}
\item[(i)] If $\vartheta_{*}\le\frac14$, then
$\mathfrak{w}^{\mathbf{B},\mathrm{corr}}_{k,j}\le 2C_{\mathbf{B}}^{v}\,\mathfrak{w}^{\mathbf{B}}_{k,j}$.
\item[(ii)] For every $\vartheta_{*}\in\,]0,1[$,
\begin{equation}\label{9.2:eq-age-weight-comparison-3}
C_{\mathbf{B}}^{v}\,(s+2)\Bigl(\frac{\vartheta_{*}}{2}\Bigr)^{s}
\le 2C_{\mathbf{B}}^{v}\,\mathfrak{w}^{\mathbf{B}}_{k,j}.
\end{equation}
If moreover $\vartheta_{*}\le\frac23$, then
\begin{equation*}
C_{\mathbf{B}}^{v}(s+2)\vartheta_{*}^{\,2s}\le 2C_{\mathbf{B}}^{v}\,\mathfrak{w}^{\mathbf{B}}_{k,j}.
\end{equation*}
\end{enumerate}
Assume in addition the following two hypotheses. The native scale, the zones, the charge zones and
the set $X^{\boxplus}$ attached to a region $X$, which occur in them and below, are those of the
volume and filament bound.
\begin{enumerate}
\item[(h1)] The volume part of the volume and filament bound: inside
$\mathbf{R}$ and $\mathbf{T}'$ a stored $\mathbf{B}^{(j)}$ is read per cell of its native scale, at
weight $C_{\mathbf{B}}^{v}$, in its own zone.
\item[(h2)] The filament part of the volume and filament bound together with
geometric forgetting of old regions: the stored $\mathbf{B}^{(j)}$ is read at most once
per level $i\in[j,k+1]$, and each read carries the factor $\vartheta_{*}^{\,k-j}$.
\end{enumerate}
Then:
\begin{enumerate}
\item[(iii)] The total weight, per native cell of its zone, of the inner reads of $\mathbf{B}^{(j)}$
by $\mathbf{R}_{k+1}$ is at most
$\mathfrak{w}^{\mathbf{B},\mathrm{corr}}_{k,j}$. If $\vartheta_{*}\le\frac14$, the inner-read bound of the
one-lattice Lipschitz step holds as written, with the weight
\eqref{9.2:eq-age-weight-comparison-1}, once the entry of its constant $K_R$ for boundary terms is
multiplied by $2C_{\mathbf{B}}^{v}$. The same holds for $\vartheta_{*}>\frac14$ if hypothesis (h2)
holds with the rate $\vartheta_{*}/2$ in place of $\vartheta_{*}$. For $\vartheta_{*}\le\frac23$ it also
holds if hypothesis (h2) holds with the rate $\vartheta_{*}^{2}$ in place of $\vartheta_{*}$.
\end{enumerate}
Assume finally:
\begin{enumerate}
\item[(h3)] The filament part of the volume and filament bound, as it bears on stored differences:
a stored difference enters at weight one per native cell, with the constant $C_{\mathbf{B}}^{v}$ and
the zone-age factor $\vartheta_{*}^{\,k-j}$, and it is absent unless a charge zone lies in
$X^{\boxplus}$, that is, it carries the indicator of that event; the statements that read this
constituent of the quantity $\Xi_X$ attached to the region $X$ in the bound of the one-lattice
Lipschitz step for the contribution of the renormalisation step $\mathbf{T}$, namely the inner-read
bound of (iii), the auxiliary slot of a further corollary, the proposition on old regions and the
corollary that chooses its rate subject to $\vartheta_{*}\le\min\{q,\vartheta\}$, read it in that
form and are unaffected by the replacement.
\end{enumerate}
Then:
\begin{enumerate}
\item[(iv)] In the bound of the one-lattice Lipschitz step for the contribution of the
renormalisation step $\mathbf{T}$, the term $2^{-(i-j)}\delta^{\mathbf{B}}_j$ inside $\Xi_X$ is
replaced by $C_{\mathbf{B}}^{v}\delta^{\mathbf{B}}_j$ at weight one per native cell; the
constituent of $\Xi_X$ so obtained is
\begin{equation}\label{9.2:eq-age-weight-comparison-4}
C_{\mathbf{B}}^{v}\cdot\delta^{\mathbf{B}}_j\cdot
\mathbf{1}\bigl[\text{a charge zone lies in }X^{\boxplus}\bigr],
\end{equation}
where $\mathbf{1}[\cdot]$ equals $1$ if the bracketed condition holds and $0$ otherwise.
The zone-age factor $\vartheta_{*}^{\,k-j}$ is supplied by the filament part of the volume and
filament bound and is not displayed in \eqref{9.2:eq-age-weight-comparison-4}. Neither the pile per
box $\Pi_j(\Delta)$ nor the pile constant $C_{\mathrm{pile}}(\sigma)$ enters the boundary-term entry
of the inner-read bound treated in (iii) or the constituent \eqref{9.2:eq-age-weight-comparison-4}
of $\Xi_X$ in (iv).
\end{enumerate}
\end{lemma}

\begin{proof}
\emph{Arithmetic.} Put $a_s:=(s+2)2^{-s}$ for $s\ge0$. Then $a_0=2$ and
\begin{equation*}
\frac{a_{s+1}}{a_s}=\frac{s+3}{2(s+2)}\le\frac34,
\end{equation*}
because $4(s+3)\le 6(s+2)$ is equivalent to $0\le 2s$. By induction on $s$,
$a_{s+1}\le\frac34a_s\le a_s$, and therefore $a_s\le a_0=2$ for all $s$. The first values are
$2,\ \frac32,\ 1,\ \frac58$.

All summands in \eqref{9.2:eq-age-weight-comparison-1} are non-negative. Hence
$\mathfrak{w}^{\mathbf{B}}_{k,j}$ is at least its summand $i=k$, which is $2^{-s}$. It is also at least
its summand $i=j$, which is $\vartheta_{*}^{\,s}$.

(i) Let $\vartheta_{*}\le\frac14$. Then $\vartheta_{*}^{\,s}\le 4^{-s}$, and
\begin{equation*}
(s+2)\vartheta_{*}^{\,s}\le (s+2)4^{-s}=a_s2^{-s}\le 2\cdot2^{-s}
\le 2\mathfrak{w}^{\mathbf{B}}_{k,j}.
\end{equation*}
Multiplying by $C_{\mathbf{B}}^{v}$ and using \eqref{9.2:eq-age-weight-comparison-2} gives
$\mathfrak{w}^{\mathbf{B},\mathrm{corr}}_{k,j}\le 2C_{\mathbf{B}}^{v}\mathfrak{w}^{\mathbf{B}}_{k,j}$.

(ii) For every $\vartheta_{*}\in\,]0,1[$,
\begin{equation*}
(s+2)\Bigl(\frac{\vartheta_{*}}{2}\Bigr)^{s}=a_s\vartheta_{*}^{\,s}\le 2\vartheta_{*}^{\,s}
\le 2\mathfrak{w}^{\mathbf{B}}_{k,j}.
\end{equation*}
Multiplying by $C_{\mathbf{B}}^{v}$ gives \eqref{9.2:eq-age-weight-comparison-3}.

For the second assertion we first show the following: $(s+2)\vartheta_{*}^{\,s}\le2$ for all $s\ge0$
if and only if $\vartheta_{*}\le\frac23$. Necessity follows from the case $s=1$, where the inequality
reads $3\vartheta_{*}\le2$. For sufficiency, put $a'_s:=(s+2)(\frac23)^s$. Then $a'_0=2$ and
\begin{equation*}
\frac{a'_{s+1}}{a'_s}=\frac23\cdot\frac{s+3}{s+2}\le\frac23\cdot\frac32=1.
\end{equation*}
So $a'_s\le2$ for all $s$, and $(s+2)\vartheta_{*}^{\,s}\le a'_s\le2$ when
$\vartheta_{*}\le\frac23$. For such $\vartheta_{*}$,
\begin{equation*}
(s+2)\vartheta_{*}^{\,2s}=\bigl((s+2)\vartheta_{*}^{\,s}\bigr)\vartheta_{*}^{\,s}
\le 2\vartheta_{*}^{\,s}\le 2\mathfrak{w}^{\mathbf{B}}_{k,j}.
\end{equation*}
Multiplying by $C_{\mathbf{B}}^{v}$ gives the claim.

(iii) By hypothesis (h1), each read of the stored $\mathbf{B}^{(j)}$ inside $\mathbf{R}$ or
$\mathbf{T}'$ costs $C_{\mathbf{B}}^{v}$ per native cell in its own zone. By hypothesis (h2), there is
at most one read per level $i\in[j,k+1]$, and each read carries $\vartheta_{*}^{\,k-j}$. There are
$k-j+2$ such levels. Summing the cost over them gives at most
$C_{\mathbf{B}}^{v}(k-j+2)\vartheta_{*}^{\,k-j}=\mathfrak{w}^{\mathbf{B},\mathrm{corr}}_{k,j}$.

If $\vartheta_{*}\le\frac14$, part (i) bounds this by
$2C_{\mathbf{B}}^{v}\mathfrak{w}^{\mathbf{B}}_{k,j}$. Since in the inner-read bound the weight
\eqref{9.2:eq-age-weight-comparison-1} enters through its product with the boundary-term entry of
$K_R$, and that bound is monotone in this product, the inner-read bound with the weight
\eqref{9.2:eq-age-weight-comparison-1} holds with the boundary-term entry of $K_R$ multiplied by
$2C_{\mathbf{B}}^{v}$.

If $\vartheta_{*}>\frac14$ and hypothesis (h2) holds with the rate $\vartheta_{*}/2$, the same count
gives the total $C_{\mathbf{B}}^{v}(s+2)(\vartheta_{*}/2)^{s}$. By
\eqref{9.2:eq-age-weight-comparison-3} this is at most
$2C_{\mathbf{B}}^{v}\mathfrak{w}^{\mathbf{B}}_{k,j}$, so the same constant suffices.

If $\vartheta_{*}\le\frac23$ and hypothesis (h2) holds with the rate $\vartheta_{*}^{2}$, the count
gives $C_{\mathbf{B}}^{v}(s+2)\vartheta_{*}^{\,2s}$. The second assertion of (ii) bounds this by
$2C_{\mathbf{B}}^{v}\mathfrak{w}^{\mathbf{B}}_{k,j}$.

(iv) Hypothesis (h3) is applied to the term $2^{-(i-j)}\delta^{\mathbf{B}}_j$ inside $\Xi_X$. It
replaces this term by $C_{\mathbf{B}}^{v}\delta^{\mathbf{B}}_j$ at weight one per native cell, with
the zone-age factor $\vartheta_{*}^{\,k-j}$ supplied by the filament part and not displayed in
\eqref{9.2:eq-age-weight-comparison-4}, and by the same hypothesis the replaced term is absent unless
a charge zone lies in $X^{\boxplus}$, that is, it carries the indicator of that event. This is the
form \eqref{9.2:eq-age-weight-comparison-4}.

The boundary-term entry of the inner-read bound of (iii) is given by
\eqref{9.2:eq-age-weight-comparison-1}, \eqref{9.2:eq-age-weight-comparison-2} and the constant
$C_{\mathbf{B}}^{v}$, and the boundary-term constituent of $\Xi_X$ in (iv) by
\eqref{9.2:eq-age-weight-comparison-4}. None of these contains $\Pi_j(\Delta)$ or
$C_{\mathrm{pile}}(\sigma)$, and hypotheses (h1)--(h3) introduce neither.
\end{proof}

The condition $\vartheta_{*}\le\frac14$ in part (i) of Lemma~\ref{9.2:lem-age-weight-comparison} is
one-sided. It holds for the exemplary value of
the rate used in the one-lattice Lipschitz step,
$\vartheta_{*}=\min\{L^{-\hat{\alpha}},L^{-1/2}\}^{2}$, where $\hat{\alpha}$ is an exponent fixed in
that step whose value is immaterial here. Whatever $\hat{\alpha}$ is, the minimum is at most
$L^{-1/2}$, so that
\begin{equation*}
\vartheta_{*}\le\bigl(L^{-1/2}\bigr)^{2}=L^{-1}.
\end{equation*}
The condition therefore holds whenever $L\ge4$, and this is implied by the lower bounds on $L$ of
Lemma~\ref{2.4:lem-stage-zero-data}, so that no new condition on $L$ arises.

Statements that choose their own rate, namely the corollary that chooses it subject to
$\vartheta_{*}\le\min\{q,\vartheta\}$ and the proposition on old regions,
are covered by the third sentence of (iii), that is, by hypothesis (h2) run at the rate
$\vartheta_{*}/2$, whenever their rate $\vartheta_{*}$ exceeds $\frac14$. The ceilings on the coupling
in geometric forgetting of old regions then read $\log\vartheta_{*}^{-1}+\log2$ in place of
$\log\vartheta_{*}^{-1}$. Since $\log\vartheta_{*}^{-1}$ is proportional to $\log L$ and stands on the
small side of a ceiling on $g$, the added $\log2$ is of the same kind. It introduces no cap in the
sense of Definition~\ref{2.3:def-cap}.

\begin{remark}[Scope: which boundary-family members are covered]\label{9.2:rem-covered-members}
The one-lattice Lipschitz step carries a slot $\mathsf{b}=(\mathbf{B}^{(j')})_{j'\le k}$ for the boundary families on the common labels; fix $j\le k$. The members of this slot, and the neighbouring terms listed below, stand as follows with respect to the pile bound for stored boundary terms (Theorem~\ref{9.2:thm-pile-bound}), the weakly anchored case (Lemma~\ref{9.2:lem-weak-anchor}), the geometric weight of old boundary terms (Proposition~\ref{9.2:prop-old-term-weight}) and the boundary-family entry of the Lipschitz step (Remark~\ref{9.2:rem-boundary-entry}).
\begin{enumerate}
\item[(a)] The $\mathbf{T}$-born members $\mathbf{B}^{(j)}$, $j\le k$, of the slot $\mathsf{b}$ are covered. They form the level-$j$ family $F$ of Definition~\ref{9.2:def-boundary-families}, indexed by the localisation domains of the stored $\mathbf{T}$-born boundary terms alive at the read. Under the strong anchor \eqref{9.2:eq-anchors-a} they are covered by Theorem~\ref{9.2:thm-pile-bound}. Under the weak anchor \eqref{9.2:eq-anchors-b} they are covered by Lemma~\ref{9.2:lem-weak-anchor}, subject to its hypothesis.
\item[(b)] The plain $\mathbf{R}$-born terms $\mathbf{B}'^{(j)}$, $j\le k$, of earlier steps are not covered. Theorem~\ref{9.2:thm-pile-bound} does not cover them, nor do Proposition~\ref{9.2:prop-old-term-weight} and Remark~\ref{9.2:rem-boundary-entry}. A size bound for them at their true localisation domain, of the form required by the hypothesis of Theorem~\ref{9.2:thm-pile-bound}, is not proved.
\item[(c)] The $\mathbf{T}$-born descendants which the correlation theorem carries wrapped in its relative terms, and the terms it attaches to its strips, are not covered.
\item[(d)] The frozen terms $\mathbb{C}^{(i)}_{k_1}$, $k_1<k$, of the arrested class $\mathcal{C}_{\mathrm{arr}}$ and the change-of-variables terms are not members of the slot $\mathsf{b}$ and are not in $F$.
\item[(e)] Reads inside $\mathbf{R}$ of the stored boundary terms entering the summation of boundary terms at finer zones are made per native cell of the component and are governed by the volume and filament bound; no pile enters. The inner reads of the one-lattice Lipschitz step weighted by its age weights are not $\mathbf{T}$-site reads (Lemma~\ref{9.2:lem-age-weight-comparison}).
\end{enumerate}
\end{remark}

\begin{proof}
\emph{(a).} Definition~\ref{9.2:def-boundary-families} takes as its units exactly the localisation domains $X\in\mathbf{D}_j$ of the stored $\mathbf{T}$-born boundary terms of level $j$ alive at the read. On these units, $F^{(j)}(X)=\mathbf{B}^{(j)}(X)$ (or its difference $\delta\mathbf{B}^{(j)}(X)$ along the segment $[\mathfrak{h},\tilde{\mathfrak{h}}]\subset\mathfrak{K}_k^{+}$ of the one-lattice Lipschitz step) satisfies the size bound \eqref{9.2:eq-boundary-families-a}.

Suppose the anchor \eqref{9.2:eq-anchors-a} of Lemma~\ref{9.2:lem-anchors} holds. Then the hypotheses of Theorem~\ref{9.2:thm-pile-bound} are met for this family, and its conclusions apply.

Suppose instead only \eqref{9.2:eq-anchors-b} holds. Then Lemma~\ref{9.2:lem-weak-anchor} gives the same conclusions, under its hypothesis on the anchor distance.

For these members, Proposition~\ref{9.2:prop-old-term-weight} and Remark~\ref{9.2:rem-boundary-entry} supply the weight $\mathfrak{a}^{T}_{k,j}$ with constant $C_T$ of \eqref{9.2:eq-old-term-weight-a}.

\emph{(b).} The one-lattice Lipschitz step folds the plain $\mathbf{R}$-born $\mathbf{B}'^{(j)}$ into its slot $\mathsf{b}$. Its boundary term $\mathbf{B}^{(k+1)}$ collects the terms of one of that step's classes localised near $\partial\Lambda_{k+1}$, together with $\mathbf{B}'$. So $\mathsf{b}$ lists both the $\mathbf{B}^{(j)}$ and the folded $\mathbf{B}'^{(j)}$.

The correlation theorem itself reads these terms at the $\mathbf{T}$-site. In that theorem's bound \eqref{9.2:eq-covered-members-1} below for these terms, let $\lambda_\theta$ be the exponent of the prefactor, $c_1^{\flat}(g_j)$ the coefficient, evaluated at the coupling $g_j$, and $\beta_s$, with $s=k-j$ the age as in Theorem~\ref{9.2:thm-pile-bound}, the coefficient of the additional decay; and let $\mathfrak{N}^{A,\mathbb{C}}$ be its size functional on a label $A$. At the $\mathbf{T}$-site the plain $\mathbf{B}'^{(j)}$ are read with the bound $e^{\lambda_\theta}c_1^{\flat}(g_j)$. These are the first-group terms off every strip, and every such unit is anchored in $Z_k\cup S$, with $S$ the set so denoted in that theorem.

They are nevertheless not covered by Theorem~\ref{9.2:thm-pile-bound}: the family $F$ of Definition~\ref{9.2:def-boundary-families} is indexed by the $\mathbf{T}$-born terms only, and, as stated with Theorem~\ref{9.2:thm-pile-bound}, the $\mathbf{R}$-born $\mathbf{B}'$, the frozen terms of $\mathcal{C}_{\mathrm{arr}}$ and the change-of-variables terms are not in $F$. Nor are they covered by Proposition~\ref{9.2:prop-old-term-weight} and Remark~\ref{9.2:rem-boundary-entry}, which supply the geometric weight \eqref{9.2:eq-old-term-weight-a} and the boundary-family entry of the Lipschitz step for the $\mathbf{T}$-born members only.

What the correlation theorem provides for these members is a bound on the label. Let $A$ be the label on which such a member is stored, as distinct from its localisation domain. The bound reads
\begin{equation}\label{9.2:eq-covered-members-1}
\mathfrak{N}^{A,\mathbb{C}}\bigl(\mathbf{B}'^{(j)}(A,\cdot)\bigr)
\le e^{\lambda_\theta}\,c_1^{\flat}(g_j)\,
e^{-(1+\frac14\beta_s)\kappa d_j(A)} .
\end{equation}
This bound is on the label $A$. It is not a bound in $d_j(X_F)$ at the true domain $X_F$, by which we mean the localisation domain of the member.

The hypothesis of Theorem~\ref{9.2:thm-pile-bound} would require instead the bound $|F(X)|\le f_je^{-\kappa d_j(X)}$ of \eqref{9.2:eq-boundary-families-a} at $X=X_F$. In that theorem the per-unit response factor $\mathsf{r}_{\varpi}(X)$ is a supremum over the same domain $X$. For these members such a bound at the true domain is not proved.

The entry of the one-lattice Lipschitz step that places $\mathbf{B}'$ in its response inequality concerns only the output $\mathbf{B}'^{(k+1)}$, whose difference norm $\|\delta\mathbf{B}'^{(k+1)}\|$ is bounded there. It is not a ground of this exclusion.

There is a scope difference with the slot of the one-lattice Lipschitz step, which we state without resolving it:
\begin{itemize}
\item the folded earlier $\mathbf{B}'^{(j)}$, $j\le k$, are members of the slot $\mathsf{b}$, and the differences $\delta\mathbf{B}^{(j)}$ of that slot's members appear in the boundary-family entry of the one-lattice Lipschitz step;
\item their weight is supplied neither by Theorem~\ref{9.2:thm-pile-bound} nor by Proposition~\ref{9.2:prop-old-term-weight} and Remark~\ref{9.2:rem-boundary-entry}.
\end{itemize}
This difference has no bearing on the boundary-family entry of Remark~\ref{9.2:rem-boundary-entry} for the $\mathbf{T}$-born members of (a).

\emph{(c).} In the correlation theorem the wrapped descendants are carried among the relative terms, and the members attached to its strips go with them; neither is covered here, and both are bounded by the carry lemma's statement on the relative boundary terms.

\emph{(d).} The frozen $\mathbb{C}^{(i)}_{k_1}$, $k_1<k$, which make up the arrested class $\mathcal{C}_{\mathrm{arr}}$ of the correlation theorem, and the change-of-variables terms, are neither members of the slot $\mathsf{b}$ nor in $F$. They therefore lie outside Theorem~\ref{9.2:thm-pile-bound}, as they lie outside the volume and filament bound.

The correlation theorem carries its own piles for them. For the $\mathbb{C}$'s the depth is at most $6(\mathsf{N}_0(\mathsf{T})+3)$, where $\mathsf{N}_0(\cdot)$ is the stage-count function of that theorem and $\mathsf{T}$ the degree threshold at which that theorem evaluates it. For the change-of-variables terms the pile is at most one per block.

\emph{(e).} The correlation theorem reads the stored boundary terms entering the summation of boundary terms at finer zones at two sites: at the $\mathbf{T}$-site, and inside $\mathbf{R}$. At the read inside $\mathbf{R}$, a stored term is read per native cell of its component. That read is governed by the volume and filament bound, and no pile enters.

The reads weighted by the age weights $\mathfrak{w}^{\mathbf{B}}_{k,j}$ of the one-lattice Lipschitz step are inner reads of $\mathbf{R}_{k+1}$ at levels $i\ge j$. They are not $\mathbf{T}$-site reads, as Lemma~\ref{9.2:lem-age-weight-comparison} states.
\end{proof}

\begin{lemma}[One-sidedness of the conditions on the coupling]\label{9.2:lem-one-sided-conditions}
We use the notation of the pile bound for stored boundary terms
(Theorem~\ref{9.2:thm-pile-bound}): $\varpi\in\,]0,1]$ and
\begin{equation*}
\delta_T:=\min\{\varpi\delta_1^{\mathrm{sc}},\,\kappa/2\},\qquad
\Lambda_L:=\log(2L^4),\qquad \check{R}^{\flat}:=\check{R}_k,\qquad x_k=g_k^{-2},
\end{equation*}
and we write
\begin{equation}\label{9.2:eq-one-sided-conditions-1}
N_0^{\sharp}(x):=2+r_{\mathrm{pr}}\log_L\log\bigl((1+B^{\sharp})x\bigr).
\end{equation}
The conditions on the coupling that enter that theorem, restated from \eqref{9.2:eq-pile-bound-a}, are
\begin{align}
a\,\delta_T\,\check{R}^{\flat}&\ge\Lambda_L\quad\text{for every }k,\qquad a\in\{10,5,4\},
\label{9.2:eq-one-sided-conditions-2}\\
5\,\delta_T\,\check{R}_k&\ge N_0^{\sharp}(2x_k)\,\Lambda_L\quad\text{for every }k,
\label{9.2:eq-one-sided-conditions-3}\\
y&\ge 8B^{\sharp}N_0^{\sharp}(y)\quad\text{for all }y\ge\gamma^{-2}-c_5,
\label{9.2:eq-one-sided-conditions-4}\\
\gamma^{-2}&\ge c_5+\frac{\max\{e,\,2^{r_{\mathrm{pr}}}\}}{1+B^{\sharp}}.
\label{9.2:eq-one-sided-conditions-5}
\end{align}
Assume:
\begin{itemize}
\item[(a)] $\kappa\ge 10\log L$, $L\ge3$, and $\check{R}_k\ge2$ at every level $k$;
\item[(b)] the floor of the radii, $(\log(x_k-c_5))^{r_{\mathrm{pr}}}\le\check{R}_k$ at every level $k$,
of Notation~\ref{9.2:not-radii-floor}, display~\eqref{9.2:eq-radii-floor-b};
\item[(c)] $g_k\le\gamma$ at every level $k$ (the scaffold of the one-lattice Lipschitz step, as
recorded in Notation~\ref{9.2:not-radii-floor});
\item[(d)] the condition of the correlation theorem for live regions, in the form
\begin{equation}\label{9.2:eq-one-sided-conditions-6}
4+2r_{\mathrm{pr}}\log_L\log\bigl((1+B^{\sharp})x\bigr)\le(\log(x-c_5))^{r_{\mathrm{pr}}}
\quad\text{for all }x\ge\gamma^{-2};
\end{equation}
\item[(e)] condition \eqref{9.2:eq-one-sided-conditions-5}.
\end{itemize}
Call the \emph{$\delta_1$-member} of \eqref{9.2:eq-one-sided-conditions-2} or
\eqref{9.2:eq-one-sided-conditions-3} the inequality obtained by replacing $\delta_T$ with
$\varpi\delta_1^{\mathrm{sc}}$, and its \emph{$\kappa$-member} the inequality obtained by replacing
$\delta_T$ with $\kappa/2$. Then:
\begin{enumerate}
\item Each of \eqref{9.2:eq-one-sided-conditions-2}, \eqref{9.2:eq-one-sided-conditions-3} holds
if and only if both its $\delta_1$-member and its $\kappa$-member hold.
\item Under (a), the $\kappa$-members of \eqref{9.2:eq-one-sided-conditions-2} hold for every $k$.
Under (a)--(e), the $\kappa$-member of \eqref{9.2:eq-one-sided-conditions-3} holds for every $k$.
\item Under (b) and (c), the $\delta_1$-member of \eqref{9.2:eq-one-sided-conditions-2} holds for
every $k$ provided
\begin{equation}\label{9.2:eq-one-sided-conditions-7}
\gamma^{-2}\ge c_5+\exp(T_a),\qquad
T_a:=\Bigl(\frac{\Lambda_L}{a\varpi\delta_1^{\mathrm{sc}}}\Bigr)^{1/r_{\mathrm{pr}}},
\end{equation}
and the $\delta_1$-member of \eqref{9.2:eq-one-sided-conditions-3} holds for every $k$ provided
\begin{equation}\label{9.2:eq-one-sided-conditions-8}
5\varpi\delta_1^{\mathrm{sc}}(\log(x-c_5))^{r_{\mathrm{pr}}}
\ge\Bigl[2+\frac{r_{\mathrm{pr}}}{\log L}\log\log\bigl(2(1+B^{\sharp})x\bigr)\Bigr]\Lambda_L
\quad\text{for all }x\ge\gamma^{-2}.
\end{equation}
Condition \eqref{9.2:eq-one-sided-conditions-7} is a lower bound on $\gamma^{-2}$, finite for every
$L$, fixed once $(L,a,\varpi,\delta_1^{\mathrm{sc}},r_{\mathrm{pr}},c_5)$ are fixed; condition
\eqref{9.2:eq-one-sided-conditions-8} holds whenever $\gamma^{-2}\ge x_{b'}$, for a number
$x_{b'}=x_{b'}(L,\varpi,\delta_1^{\mathrm{sc}},r_{\mathrm{pr}},B^{\sharp},c_5)$ finite for every $L$.
\item Condition \eqref{9.2:eq-one-sided-conditions-4} is a condition on $\gamma$ alone once
$(B^{\sharp},r_{\mathrm{pr}},c_5)$ are fixed, satisfied by all sufficiently small $\gamma$,
uniformly in $L\ge3$; condition \eqref{9.2:eq-one-sided-conditions-5} is a lower bound on
$\gamma^{-2}$ not involving $L$.
\item Each condition on the coupling above is one-sided and is stated in the form ``for all
$x\ge\gamma^{-2}$'' (respectively ``for all $y\ge\gamma^{-2}-c_5$''), or as a single lower bound on
$\gamma^{-2}$; hence if it holds for $\gamma$, it holds for every $\gamma'\in\,]0,\gamma]$. None of the
conditions used in this section bounds $L$, $M$, $K$, $k$, $j$, an age, the numbers $N_b$, $N_0(b)$
of a non-completed attempt $b$, or the nesting depth from above, nor the running inverse coupling
$x_k$ from above; the only upper bounds are on the reference coupling $\gamma$, each in the form of a
lower bound on $\gamma^{-2}$; no value of $L_0$ enters.
\end{enumerate}
\end{lemma}

\begin{proof}
\emph{Splitting, item (1).} Let $\alpha>0$, $A,B\ge0$, $\check{R}\ge0$ and $\Lambda\in\mathbb{R}$.
If $\alpha\min\{A,B\}\check{R}\ge\Lambda$, then $\alpha A\check{R}\ge\alpha\min\{A,B\}\check{R}\ge\Lambda$
and likewise $\alpha B\check{R}\ge\Lambda$. Conversely, $\min\{A,B\}$ equals $A$ or $B$, so if both
$\alpha A\check{R}\ge\Lambda$ and $\alpha B\check{R}\ge\Lambda$ hold, then
$\alpha\min\{A,B\}\check{R}\ge\Lambda$. We apply this with $A=\varpi\delta_1^{\mathrm{sc}}$,
$B=\kappa/2$: for \eqref{9.2:eq-one-sided-conditions-2} with $\alpha=a$,
$\check{R}=\check{R}^{\flat}$, $\Lambda=\Lambda_L$, and for \eqref{9.2:eq-one-sided-conditions-3}
with $\alpha=5$, $\check{R}=\check{R}_k$, $\Lambda=N_0^{\sharp}(2x_k)\Lambda_L$. This is item (1).

\emph{The $\kappa$-members of \eqref{9.2:eq-one-sided-conditions-2}.} By (a), $\kappa\ge10\log L>0$,
$\check{R}^{\flat}=\check{R}_k\ge2$ and $a\ge4$, so
\begin{equation*}
\tfrac{a}{2}\,\kappa\,\check{R}^{\flat}\ge 2\cdot10\log L\cdot2=40\log L
\ge\log2+4\log L=\Lambda_L,
\end{equation*}
the last inequality because $36\log L\ge\log 2$ for $L\ge2$.

\emph{The $\delta_1$-members of \eqref{9.2:eq-one-sided-conditions-2}.} Since $\Lambda_L>0$, $T_a$ is
well defined and $T_a\ge0$. Assume \eqref{9.2:eq-one-sided-conditions-7}. For $x\ge\gamma^{-2}$ we
have $x-c_5\ge\exp(T_a)$, hence $\log(x-c_5)\ge T_a\ge0$, and since $t\mapsto t^{r_{\mathrm{pr}}}$ is
increasing on $[0,\infty)$,
\begin{equation*}
(\log(x-c_5))^{r_{\mathrm{pr}}}\ge T_a^{r_{\mathrm{pr}}}=\frac{\Lambda_L}{a\varpi\delta_1^{\mathrm{sc}}}.
\end{equation*}
By (c), $x_k=g_k^{-2}\ge\gamma^{-2}$, so this applies at $x=x_k$; by (b),
$\check{R}^{\flat}=\check{R}_k\ge(\log(x_k-c_5))^{r_{\mathrm{pr}}}$. Multiplying by
$a\varpi\delta_1^{\mathrm{sc}}>0$ gives $a\varpi\delta_1^{\mathrm{sc}}\check{R}^{\flat}\ge\Lambda_L$
for every $k$. Condition \eqref{9.2:eq-one-sided-conditions-7} involves $\gamma$ and the fixed
quantities $L,a,\varpi,\delta_1^{\mathrm{sc}},r_{\mathrm{pr}},c_5$ only, and $c_5+\exp(T_a)$ is finite for
every $L$, so it is met by all sufficiently small $\gamma$.

At $\varpi=\tfrac12$ the $\delta_1$-members for $a=10,5,4$ read
$5\,\delta_1^{\mathrm{sc}}\check{R}^{\flat}\ge\Lambda_L$,
$\tfrac52\,\delta_1^{\mathrm{sc}}\check{R}^{\flat}\ge\Lambda_L$ and
$2\,\delta_1^{\mathrm{sc}}\check{R}^{\flat}\ge\Lambda_L$. Since $e^{\Lambda_L}=2L^4$, the last is
the displayed floor $e^{2\delta_1^{\mathrm{sc}}\check{R}_k}\ge2L^4$, and since
$\delta_1^{\mathrm{sc}}\check{R}^{\flat}\ge0$ it implies the other two. At $\varpi=\tfrac14$ they read
$\tfrac52$, $\tfrac54$, $1$ times $\delta_1^{\mathrm{sc}}\check{R}^{\flat}\ge\Lambda_L$; the case $a=10$
is still implied by the floor $e^{2\delta_1^{\mathrm{sc}}\check{R}_k}\ge2L^4$ because $\tfrac52\ge2$,
and the cases $a=5,4$ are the same floor with a lowered coefficient, each guaranteed by
\eqref{9.2:eq-one-sided-conditions-7}.

\emph{The $\delta_1$-member of \eqref{9.2:eq-one-sided-conditions-3}.} By
\eqref{9.2:eq-one-sided-conditions-1},
\begin{equation*}
N_0^{\sharp}(2x)=2+\frac{r_{\mathrm{pr}}}{\log L}\log\log\bigl(2(1+B^{\sharp})x\bigr),
\end{equation*}
which is the bracket in \eqref{9.2:eq-one-sided-conditions-8}. Assume
\eqref{9.2:eq-one-sided-conditions-8}. By (c), $x_k\ge\gamma^{-2}$, and by (b),
\begin{equation*}
5\varpi\delta_1^{\mathrm{sc}}\check{R}_k\ge5\varpi\delta_1^{\mathrm{sc}}(\log(x_k-c_5))^{r_{\mathrm{pr}}}
\ge N_0^{\sharp}(2x_k)\Lambda_L
\end{equation*}
for every $k$. In \eqref{9.2:eq-one-sided-conditions-8} the left side is a positive multiple of
$(\log(x-c_5))^{r_{\mathrm{pr}}}$ with $r_{\mathrm{pr}}\ge2$, while the right side grows like
$(\log L)^{-1}\log\log x$ times $\Lambda_L$; the ratio of the left side to the right side therefore
tends to infinity as $x\to\infty$, so the inequality holds for all $x$ in a ray $[x_{b'},\infty)$ with
$x_{b'}$ depending on $(L,\varpi,\delta_1^{\mathrm{sc}},r_{\mathrm{pr}},B^{\sharp},c_5)$ and finite for
every $L$. Hence \eqref{9.2:eq-one-sided-conditions-8} holds whenever $\gamma^{-2}\ge x_{b'}$; it is a
condition of the same form as \eqref{9.2:eq-one-sided-conditions-6}. The third item is proved.

\emph{The $\kappa$-member of \eqref{9.2:eq-one-sided-conditions-3}.} Assume (a)--(e). For
$y\ge\gamma^{-2}-c_5$, \eqref{9.2:eq-one-sided-conditions-5} gives
$(1+B^{\sharp})y\ge\max\{e,2^{r_{\mathrm{pr}}}\}$, that is
\begin{equation}\label{9.2:eq-one-sided-conditions-9}
\log\bigl((1+B^{\sharp})y\bigr)\ge\max\{1,\;r_{\mathrm{pr}}\log2\}.
\end{equation}
In particular $\log\log((1+B^{\sharp})y)\ge0$, so $N_0^{\sharp}(y)\ge2$, and $N_0^{\sharp}$ is
non-decreasing there, being a composition of increasing functions. We use these facts at $y=x_k$ and
$y=2x_k$, which satisfy $y\ge\gamma^{-2}-c_5$ because $x_k\ge\gamma^{-2}$ by (c) and $c_5=2c_2\ge0$.
For such $x$,
\begin{equation*}
N_0^{\sharp}(2x)-N_0^{\sharp}(x)
=r_{\mathrm{pr}}\log_L\frac{\log(2(1+B^{\sharp})x)}{\log((1+B^{\sharp})x)}
=r_{\mathrm{pr}}\log_L\Bigl(1+\frac{\log2}{\log((1+B^{\sharp})x)}\Bigr),
\end{equation*}
and by \eqref{9.2:eq-one-sided-conditions-9} the fraction is at most $1/r_{\mathrm{pr}}$. Since
$\log(1+v)\le v$ for $v\ge0$, we have $r\log(1+1/r)\le1$ for $r>0$; hence
\begin{equation*}
N_0^{\sharp}(2x)-N_0^{\sharp}(x)\le r_{\mathrm{pr}}\log_L\Bigl(1+\frac1{r_{\mathrm{pr}}}\Bigr)
=\frac{r_{\mathrm{pr}}\log(1+1/r_{\mathrm{pr}})}{\log L}\le\frac1{\log L}=\log_Le\le1,
\end{equation*}
the last step because $L\ge3\ge e$. Next, $2N_0^{\sharp}(x)$ is the left side of
\eqref{9.2:eq-one-sided-conditions-6}; applying \eqref{9.2:eq-one-sided-conditions-6} at
$x=x_k\ge\gamma^{-2}$ and then (b) gives
\begin{equation*}
2N_0^{\sharp}(x_k)\le(\log(x_k-c_5))^{r_{\mathrm{pr}}}\le\check{R}_k.
\end{equation*}
With $\kappa\ge10\log L$, the left side of the $\kappa$-member satisfies
\begin{equation*}
\tfrac{5\kappa}{2}\,\check{R}_k\ge25\log L\cdot\check{R}_k\ge50\log L\cdot N_0^{\sharp}(x_k).
\end{equation*}
For the right side, $\Lambda_L=\log2+4\log L\le5\log L$ for $L\ge2$, and $n+1\le\tfrac32n$ for
$n\ge2$, so, using $N_0^{\sharp}(x_k)\ge2$,
\begin{equation*}
N_0^{\sharp}(2x_k)\Lambda_L\le\bigl(N_0^{\sharp}(x_k)+1\bigr)\Lambda_L
\le\bigl(N_0^{\sharp}(x_k)+1\bigr)\,5\log L\le\tfrac{15}{2}\log L\cdot N_0^{\sharp}(x_k).
\end{equation*}
Since $50\ge\tfrac{15}2$ and $\log L\cdot N_0^{\sharp}(x_k)\ge0$, we obtain
$\tfrac{5\kappa}2\check{R}_k\ge N_0^{\sharp}(2x_k)\Lambda_L$ for every $k$. This completes the second
item. Without assumption (d), the $\kappa$-member of \eqref{9.2:eq-one-sided-conditions-3} is itself a
condition on $\gamma$ of the same form as \eqref{9.2:eq-one-sided-conditions-6}, uniform in $L\ge3$.

\emph{Item (4).} For $L\ge3$ and $t\ge1$, $\log_Lt=\log t/\log L\le\log t/\log3$, so the right side of
\eqref{9.2:eq-one-sided-conditions-4} is at most
$8B^{\sharp}\bigl(2+r_{\mathrm{pr}}\log\log((1+B^{\sharp})y)/\log3\bigr)$ wherever
$\log((1+B^{\sharp})y)\ge1$, a bound free of $L$. Since $\log\log$ grows more slowly than any positive
power, this bound is smaller than $y$ for all sufficiently large $y$; taking the threshold also at
least $e/(1+B^{\sharp})$, it depends on $(B^{\sharp},r_{\mathrm{pr}})$ only, and
\eqref{9.2:eq-one-sided-conditions-4} holds for all
$y\ge\gamma^{-2}-c_5$ as soon as $\gamma^{-2}$ exceeds $c_5$ plus that threshold, uniformly in
$L\ge3$. Condition \eqref{9.2:eq-one-sided-conditions-5} is a lower bound on $\gamma^{-2}$ by a number
depending on $(c_5,r_{\mathrm{pr}},B^{\sharp})$ only.

\emph{Item (5): one-sidedness.} Each of \eqref{9.2:eq-one-sided-conditions-4},
\eqref{9.2:eq-one-sided-conditions-6} and \eqref{9.2:eq-one-sided-conditions-8} asserts an
inequality for all $x\ge\gamma^{-2}$, respectively all $y\ge\gamma^{-2}-c_5$, and each of
\eqref{9.2:eq-one-sided-conditions-5}, \eqref{9.2:eq-one-sided-conditions-7} is a single lower bound
on $\gamma^{-2}$. If $0<\gamma'\le\gamma$, then $\gamma'^{-2}\ge\gamma^{-2}$, so the range of $x$
(respectively $y$) for $\gamma'$ is contained in that for $\gamma$, and a lower bound met by
$\gamma^{-2}$ is met by $\gamma'^{-2}$; hence each condition that holds for $\gamma$ holds for
$\gamma'$. By items (3) and (4) each such condition is non-empty for every $L$.

The standing floors used in this section are the following, all lower bounds. The floor
$\check{R}_k\ge2$ is used in the inequality $8\check{R}_k-3\ge\tfrac{13}2\check{R}_k$ of the combined
size and response exponent (Lemma~\ref{9.2:lem-combined-exponent}), in the weakly anchored case
(Lemma~\ref{9.2:lem-weak-anchor}, display~\eqref{9.2:eq-weak-anchor-a}), and in the $\kappa$-members
above. The floor $L\ge3$ is used above; elsewhere in this section only $L\ge1$ or $L\ge2$ is used,
except for $L\ge4$ in the bound $\vartheta_{*}\le L^{-1}\le\tfrac14$ on the age rate $\vartheta_{*}$ of
the weights of stored boundary terms; both are implied by the lower bound on $L$ in force in this
chapter. The floor $\kappa\ge10\log L$ is used in the $\kappa$-members. The floor $M\ge12$ is used only
in Lemma~\ref{9.2:lem-read-set} (read sets of cells meeting the domain), and follows from the floor
$M\ge L^4$ since $L\ge3$. Of the conditions \eqref{9.2:eq-attempt-windows-a} on attempt windows
(Definition~\ref{9.2:def-attempt-windows}), only the last is used, and only on labels at which
non-completed attempts are in force. The following conditions appear in this section but are
not used in the arguments above: a further clause of the condition of the correlation theorem for
live regions; the thin-level condition on the radius ratio $\mathfrak{r}_a$; the condition
$\delta_1M\ge\kappa_1$ of the correlation theorem, which enters Lemma~\ref{9.2:lem-shell-lemma}
(block reach and the shell lemma) as a condition of that theorem; and the displacement bound
$\vartheta^{T}<\log2$. The room condition \eqref{9.2:eq-boundary-entry-a} of
Remark~\ref{9.2:rem-boundary-entry} is of the same kind as a coupling condition of the one-lattice
Lipschitz step.

No condition of this section bounds the running inverse coupling $x_k$ from above; the only upper
bounds are on the reference coupling $\gamma$, each in the form of one of the lower bounds on
$\gamma^{-2}$ just listed. The only quantity of the type of an age that occurs, the age $u$
of the exceptional window (together with its thin count $u'$), is split, never bounded. The depth
$s=k-j$, the number of windows, the numbers $N_b$, $N_0(b)$ of a non-completed attempt $b$, the
nesting depth and the anchor excess of the weakly anchored case enter no upper bound on a quantity of
the run; the hypothesis of the weakly anchored case (Lemma~\ref{9.2:lem-weak-anchor}) bounds that
excess by a geometric quantity given together with the family. The quantities that depend on
$L$ --- $\Lambda_L$, $1/\log L$, $(\hat{\sigma}+2)^4$ and $C_U$ --- occur only on the side of a
ceiling on $\gamma$ that is dominated once $\gamma$ is small. No value of $L_0$ enters any of these
conditions. This proves item (5).

\emph{Alternative readings.} The same conclusions hold, with changed coefficients and no change in the
form of any condition, under the following alternative readings. If the collar estimate used in the
pile bound is available at eight layers only, the coefficients $a=10,5,4$ in
\eqref{9.2:eq-one-sided-conditions-2} become $8,4,\tfrac{16}5$, and
\eqref{9.2:eq-one-sided-conditions-4} is adjusted accordingly. If the collar of
the step is taken at five layers, as in \cite{B-CMP119}, the exponent $\tfrac{13}2$ in
Lemma~\ref{9.2:lem-combined-exponent} becomes $\tfrac72$, and \eqref{9.2:eq-one-sided-conditions-3}
becomes $(\tfrac72-c)\,\delta_T\check{R}_k\ge N_0^{\sharp}(2x_k)\Lambda_L$, where $c$ denotes the
constant by which that five-layer collar lowers the coefficient. If the read set is the set
of cells meeting the domain, only constants change (Lemma~\ref{9.2:lem-read-set}). If the windows are
placed one index higher, nothing changes.
\end{proof}

\begin{definition}[Parameters, running couplings and geometric constants]
\label{9.2:not-parameters-couplings}
\emph{Parameters.} The dimension is $d=4$. The blocking factor $L$ is an odd integer with
$L\ge 5$. The block sizes $M_1<M_2=L^2M_1<M$ are powers of $L$ \cite{B-CMP119}, and all
partitions of the lattice introduced below are compatible with one another \cite{B-CMP119}.
Since $M$ and $M_2$ are powers of $L$ with $M>M_2$, the quotient $M/M_2$ is a power of $L$
larger than $1$, hence at least $L$; as $L\ge 5$,
\[
M\ge LM_2\ge 4M_2 .
\]
\emph{Exponents:} $p_0\ge 1$ is the exponent of the degree function $p_0(\cdot)$, and
$r\ge 1$ is the exponent of the lengths $R_n$ introduced below.

\emph{Constants:} $A_0\ge 1$ is a sufficiently large constant, as in \cite{B-CMP119}.
$A_1>0$ and $\gamma_0>0$ are the constants of Ba{\l}aban's large-field factors
\cite[p.~381]{B-CMP122b}; as stated there, $\gamma_0$ is an absolute positive constant, and
in the large-field factor one may take $\gamma_0=1/2$. $\beta_0\in\,]0,1[$ is the flow
constant of
\cite[(2.6)--(2.7)]{B-CMP119}. $B_3$ is the constant of \cite[Theorem~1]{B-CMP102b}.
$\Theta_1=\Theta_1(d,L)\ge 1$ is the collar constant, which depends on $d$ and $L$ only.
$M_0$ is the constant of \cite[(1.101)]{B-CMP122a}. $C_\Lambda$ is the constant of the
regular extension. $C_\vartheta$ is the constant of the error term of the second-class
income bound. The decay constant $\kappa$ and the constant $C_\sharp$ of the correlation
theorem, and the dial constant $C_{L'}=C_LE_1$, enter only through the charge constant $C_0$
of the majorant, which is fixed later in this chapter.

\emph{Running couplings.} Let $k\ge 1$ be the current level; levels are integers $n$ with
$1\le n\le k$. The numbers $g_1,\dots,g_k\in\,]0,1[$ are the running couplings of the flow
of clause~(0) of Theorem~0 (Definition~\ref{1.2:def-couplings};
\cite[(2.6)--(2.7)]{B-CMP119}). The letter $\gamma$ denotes the coupling parameter with
respect to which every running coupling obeys the bound \eqref{9.2:eq-parameters-couplings-1}
below (a lower bound on $g_n^{-2}$), and on which the upper bounds on $\gamma$ required later
are placed. All numbers
defined below are formed from $g_1,\dots,g_k$ at their real values; they do not depend on
the complex coupling parameter of the correlation theorem. Put
\begin{equation}\label{9.2:eq-parameters-couplings-a}
\begin{gathered}
\Xi_n:=\log g_n^{-2},\qquad \mathsf{P}_n:=p_0(g_n)=A_0\,\Xi_n^{p_0},\qquad
\varepsilon_n:=g_n\,p_0(g_n),\qquad \mathfrak{p}_n:=\varepsilon_n/g_n,\\
\mathfrak{b}_n^2:=A_1^2\,\Xi_n^{2p_0},\\
\mathsf{Q}_n:=\frac{\gamma_0A_1^2\,\Xi_n^{2p_0}}{256\,\Theta_1^2(1+\beta_0)^2},\qquad
\mathsf{D}_n:=\frac{\mathsf{Q}_n}{2\sqrt d},\qquad
\upsilon_n:=e^{-2\mathsf{Q}_n}.
\end{gathered}
\end{equation}
Here $\varepsilon_n$ is the quantity of \cite[(2.4)]{B-CMP119}, and $\mathfrak{b}_n^2$ is the
quantity written $A_1^2p_0^2(g_n)$ in \cite[p.~381]{B-CMP122b}.

The following properties of the running couplings are taken by statement from the flow
facts of the flow of clause~(0) of Theorem~0. First, with $c_2$ the flow constant (not to be
confused with the charge numbers $c_n$ defined below), every running coupling satisfies
\begin{equation}\label{9.2:eq-parameters-couplings-1}
g_n^{-2}\ge\gamma^{-2}-2c_2 .
\end{equation}
Consequently, if $\gamma^{-2}>2c_2+1$, then $\gamma^{-2}-2c_2>1$, and since the logarithm
is increasing, \eqref{9.2:eq-parameters-couplings-1} gives
$\Xi_n\ge\log(\gamma^{-2}-2c_2)>0$ for every $n$. Moreover, in either of the two forms in
which \cite[(2.4)]{B-CMP119} writes $\varepsilon_n$ one has $\mathfrak{p}_n\ge\mathsf{P}_n$;
with the form adopted in \eqref{9.2:eq-parameters-couplings-a},
$\mathfrak{p}_n=p_0(g_n)=\mathsf{P}_n$.

\emph{Lengths.} For each level $n$, $R_n$ is the length of level $n$ fixed in the correlation
theorem; it is a power of $L$ with
\begin{equation*}
1\le R_n<L\,\Xi_n^{r},\qquad R_{n+1}\le R_n ,
\end{equation*}
the second inequality being one of the flow facts. Ba{\l}aban's own length
\cite[(2.5)]{B-CMP119}, the least power of $L$ not smaller than $\Xi_n^{r}$, is covered by
the same notation: everything below applies to it under the same letter $R_n$. Put
$\ell_n:=LR_{n+1}/R_n$ for $n<k$, and $\ell_n:=L$ for $n\ge k$ (the
constant continuation past the last step). The flow facts state, under the definition of the
lengths $R_n$ and the upper bound on $\gamma$ that they require, that $\ell_n\in\{1,L\}$
for $n<k$, and that $\ell_n=1$ implies that ($n=1$ or $\ell_{n-1}=L$) and $\ell_{n+1}=L$;
together with the definition for $n\ge k$, $\ell_n\in\{1,L\}$ for every $n$. Of $R_n$ only
the two displayed facts and this property of the ratios $\ell_n$ are used. Put
$\beta_n:=d\,(LMR_n)^d$, the number of free bonds per $MR_n$-cell available to the
refusing-bond labels.

\emph{Geometric constants.} The numbers
\begin{equation}\label{9.2:eq-parameters-couplings-b}
\begin{gathered}
\Gamma:=\sqrt d\cdot 2^{d+1}\cdot 219^{d},\qquad \sigma:=220\sqrt d,\qquad c_0:=2d+24,\\
\Gamma'_n:=\Gamma\ \text{ if }\ell_n=L,\qquad \Gamma'_n:=3\Gamma\ \text{ if }\ell_n=1,
\end{gathered}
\end{equation}
depend on $d$ only (and $\Gamma'_n$ on $\ell_n$). At $d=4$, $\sqrt d=2$ and $2^{d+1}=32$, so
$\Gamma=64\cdot 219^4$, $\sigma=440$ and $c_0=32$. Since $\ell_n\in\{1,L\}$,
\eqref{9.2:eq-parameters-couplings-b} defines $\Gamma'_n$ for every $n$.

\emph{Charge numbers and the forgetting ratio.} For a number $C>0$ (a charge constant) put
$c_n:=C M^dR_n^{d+1}$, the charge number of level $n$ (not to be confused with the flow
constant $c_2$ or with $c_0$ above).
Since $R_{n+1}\le R_n$ and $d+1>0$, one has $c_{n+1}\le c_n$.
Finally, $c_a\ge\log 2$ is a free number, and $q:=e^{-c_a}$; since $c_a\ge\log 2$,
$q\le e^{-\log 2}=\tfrac12$.
\end{definition}

\begin{definition}[Cells, coarsening, box steps and linear size]\label{9.2:def-cells-coarsening}
The lengths $R_n$, the coarsening ratios $\ell_n$, the number $M$ and the dimension $d=4$ are
those of Notation~\ref{9.2:not-parameters-couplings}.

\emph{Cells.} After the $n$-th renormalisation step the configurations live on the unit lattice
of the $L^{-n}$-scale, which is a torus. We write $\mathbb{L}_n$ for the set of the
$MR_n$-cubes of Ba{\l}aban's partition of this torus \cite{B-CMP122b}. Its elements are the
\emph{cells of level $n$}; they are indexed by $\mathbb{Z}^d$ modulo the torus of level $n$.
For $x,y\in\mathbb{L}_n$, $|x-y|_\infty$ is the sup-distance of their indices on the torus.
Two cells $x,y$ are \emph{adjacent} if $|x-y|_\infty\le 1$. The words ``connected'' and
``component'' refer to this adjacency: a component of a set $A\subset\mathbb{L}_n$ is a maximal
non-empty connected subset of $A$. For $A\subset\mathbb{L}_n$ we put
$\operatorname{dist}_\infty(y,A):=\min_{x\in A}|x-y|_\infty$,
\begin{equation*}
A^{\sim s}:=\{\,y\in\mathbb{L}_n:\ \operatorname{dist}_\infty(y,A)\le s\,\},
\end{equation*}
and we write $|A|$ for the number of cells of $A$. A set $A$ \emph{fits in $D$ cells per axis}
if $A\subset v+\{0,\dots,D-1\}^d$ for some $v$. A \emph{$100$-cube} is a set of the form
$v+\{0,\dots,99\}^d$. The \emph{realisation} of the cell $x$ is the closed cube
$x+[0,1]^d\subset\mathbb{R}^d$. Lengths are measured in level-$n$ cell units, one unit being
$MR_n$ lattice spacings of the level-$n$ lattice.

\emph{Coarsening.} The partitions of consecutive levels are nested. Each $MR_{n+1}$-cube of the
level-$(n+1)$ lattice is an $LMR_{n+1}$-cube of the level-$n$ lattice \cite{B-CMP119}. Since
$LMR_{n+1}/(MR_n)=\ell_n$, the coarsening ratio $\ell_n=LR_{n+1}/R_n$, such a cube is a block
of $\ell_n^d$
cells of $\mathbb{L}_n$. The \emph{coarsening map}
$\pi_n:\mathbb{L}_n\to\mathbb{L}_{n+1}$, given by
$x\mapsto\lfloor x/\ell_n\rfloor$ coordinatewise, assigns to $x$ the cell of level $n+1$
that contains it. For $A\subset\mathbb{L}_n$ we put $\pi_nA:=\{\pi_n(x):x\in A\}$.
Ba{\l}aban's operation that covers a region by a smallest union of $LMR_{n+1}$-cubes and then
adds ten layers of such cubes \cite{B-CMP122b} is, in this notation,
\begin{equation}\label{9.2:eq-cells-coarsening-1}
S_n(A):=(\pi_nA)^{\sim 10}\subset\mathbb{L}_{n+1}.
\end{equation}

\emph{Box steps.} A \emph{box step} on a set $Y\subset\mathbb{L}_n$ replaces one component $K$
of $Y$ that fits in a $100$-cube by the smallest rectangular parallelepiped of cells that
contains $K$ \cite[p.~269]{B-CMP119}.

\emph{Linear size.} For a finite non-empty $A\subset\mathbb{L}_n$,
\begin{equation}\label{9.2:eq-cells-coarsening-a}
\begin{split}
d'(A):=\inf\bigl\{\operatorname{length}(T):\ & T\subset\mathbb{R}^d \text{ a finite connected
union of straight segments}\\
& \text{meeting the closed cell of every } x\in A\bigr\},
\end{split}
\end{equation}
with lengths in level-$n$ cell units. We call $d'$ the \emph{unconstrained linear size}. The
\emph{constrained linear size} $d'_P(A)$ of Ba{\l}aban's construction is the same infimum with
$T$ required in addition to lie in the union of the closed cells of $A$; the infimum of the
empty set is $+\infty$, a convention that covers the sets $A$ that are not connected, for which
no such $T$ exists. For Ba{\l}aban's definition of the linear size in terms of $MR_n$-cubes
see \cite{B-CMP122b}. These two quantities satisfy
\begin{equation}\label{9.2:eq-cells-coarsening-2}
0\le d'(A)\le d'_P(A).
\end{equation}
When $A\subset\mathbb{L}_n$ and $\pi_nA\subset\mathbb{L}_{n+1}$ are compared, $d'(\pi_nA)$ is
measured in level-$(n+1)$ units. The functional $d'$ enters what follows only through the
deposit bounds, the bounds on cores, and one recursion of this chapter. No lower bound for $d'$
is used anywhere.
\end{definition}

\begin{proof}[Proof of \eqref{9.2:eq-cells-coarsening-2}]
Every set $T$ admitted in \eqref{9.2:eq-cells-coarsening-a} has length at least $0$, so
$d'(A)\ge 0$.

A set $T$ admitted in the definition of $d'_P(A)$ satisfies every requirement of
\eqref{9.2:eq-cells-coarsening-a}, and in addition it lies in the union of the closed cells of
$A$. Hence the class over which $d'_P(A)$ is the infimum is contained in the class over which
$d'(A)$ is the infimum. The infimum over a subclass is at least the infimum over the whole
class, which gives $d'(A)\le d'_P(A)$. If the subclass is empty, then $d'_P(A)=+\infty$ and the
inequality holds trivially.
\end{proof}

\begin{lemma}[Elementary properties of the linear size]\label{9.2:lem-linear-size-properties}
Let $A,B\subset\mathbb{L}_n$ be finite and non-empty, and let $d'$ be the linear size
\eqref{9.2:eq-cells-coarsening-a} of Definition~\ref{9.2:def-cells-coarsening}.
\begin{enumerate}
\item[(a)] If $A$ is connected, then $d'(A)\le\sqrt{d}\,(|A|-1)$.
\item[(b)] $d'(\pi_nA)\le d'(A)/\ell_n$, where the left member is measured in level-$(n+1)$
cell units.
\item[(c)] If $A\subset B$, then $d'(A)\le d'(B)$.
\end{enumerate}
\end{lemma}

\begin{proof}
Call a set $T\subset\mathbb{R}^d$ \emph{admissible for} a finite non-empty set of cells if it is a
finite connected union of straight segments meeting the closed cell of every element of that set.
By \eqref{9.2:eq-cells-coarsening-a}, $d'$ of the set is the infimum of $\operatorname{length}(T)$
over the sets $T$ admissible for it.

(a) The \emph{adjacency graph} of $A$ has vertex set $A$, two distinct cells being joined by an
edge when they are adjacent in the sense of Definition~\ref{9.2:def-cells-coarsening}. Since $A$
is connected, this graph is connected and therefore has a spanning tree, which has $|A|-1$ edges.
For each edge $\{x,y\}$ of the tree, join the centres $x+(\tfrac12,\dots,\tfrac12)$ and
$y+(\tfrac12,\dots,\tfrac12)$ of the closed cells of $x$ and $y$ by a straight segment. Adjacent
cells satisfy $|x-y|_\infty\le1$, so their centres differ by at most $1$ in each coordinate, and
each segment has length at most $\sqrt{d}$. Let $T$ be the union of these segments (when
$|A|=1$, the centre of the single cell, a segment of length zero). Since the tree is connected,
$T$ is connected; it meets the closed cell of every $x\in A$ at its centre; and its length is at
most $\sqrt{d}\,(|A|-1)$. Thus $T$ is admissible for $A$, and
$d'(A)\le\sqrt{d}\,(|A|-1)$.

(b) Let $T$ be admissible for $A$. For every $x\in A$, the closed block realising the cell
$\pi_n(x)$ of level $n+1$ is the union of the closed cells of the $\ell_n^d$ cells of
$\mathbb{L}_n$ that it contains, among them the closed cell $x+[0,1]^d$. Since $T$ meets the
closed cell of $x$, it meets the closed block realising $\pi_n(x)$. Hence $T$ meets the closed
cell of every element of $\pi_nA$, which is finite and non-empty, and is admissible for $\pi_nA$.
One level-$(n+1)$ cell unit is $\ell_n$ level-$n$ cell units, so the length of $T$ measured in
level-$(n+1)$ units is $\operatorname{length}(T)/\ell_n$. Therefore
$d'(\pi_nA)\le\operatorname{length}(T)/\ell_n$, and taking the infimum over all $T$ admissible
for $A$ gives $d'(\pi_nA)\le d'(A)/\ell_n$.

(c) Let $A\subset B$. A set meeting the closed cell of every element of $B$ meets the closed cell
of every element of $A$, so every set admissible for $B$ is admissible for $A$. The infimum
defining $d'(A)$ is taken over a family containing the one defining $d'(B)$, and hence
$d'(A)\le d'(B)$.
\end{proof}

\begin{definition}[The construction family and the outcomes of the operation]
\label{9.2:def-construction-outcomes}
Cells, the set $\mathbb{L}_n$ of cells of level $n$, the coarsening map $\pi_n$, Ba{\l}aban's
ten-layer operation $S_n$, box steps and the neighbourhoods $A^{\sim s}$, $Y^{\sim -s}$ of a cell
set are as in Definition~\ref{9.2:def-cells-coarsening}.

\emph{Large-field regions.} For each level $n\ge 1$, $Z_n\subset\mathbb{L}_n$ denotes the
large-field region produced by the step $n-1\to n$, read before the operation $\mathbf{R}$ of
step $n$. It is a union of cells \cite{B-CMP122b}. We write $\operatorname{comp}(Z_n)$ for the
set of its components.

\emph{The family.} The constructions considered form a family containing Ba{\l}aban's
construction (\cite[\S3]{B-CMP119}, \cite{B-CMP122a}, \cite[\S1]{B-CMP122b}). A member of the
family enters what follows only through three things: the properties (G1)--(G4) below, the
outcomes (s0)--(s2) below, and the hypotheses of the lemmas of this chapter that use it.
Ba{\l}aban's construction is the member of this family in which the set $I'$ below is empty.

\emph{(G4) The operation $\mathbf{R}$ at step $n$.} This operation is begun on those
components of $Z_n$ that are of class (i) or (ii) in the following sense of maturity
(\cite{B-CMP122a}, \cite{B-CMP122b}):
\begin{itemize}
\item[(i)] the component is contained in a cube of side $100$ cells;
\item[(ii)] during the preceding $\mathsf{N}$ steps, where the length $\mathsf{N}$ of this window
satisfies $\mathsf{N}\le R_n$ (the assumption of \cite{B-CMP122b}, ``We have assumed here that
$N\le R_k$''), its lineage (the components of earlier levels from which it descends, each a
constituent, in the sense defined below, of the next) had no creation, join or re-creation
(these events are defined below), subject to the further nesting clause of the maturity rule.
\end{itemize}
For such a component $Z_0$ there are three outcomes.
\begin{itemize}
\item[(s0)] \emph{Healed (first class).} All preparatory characteristic functions are
small-field. The integral \cite[(1.71)]{B-CMP122b} is performed, and the last decomposition of
unity of \cite[p.~378]{B-CMP122b} gives the small-field function. Then $Z_0$ leaves the
large-field region, its operation becomes a polymer activity of \cite[(1.72)]{B-CMP122b}, and
$Z_0$ belongs to no later history.
\item[(s1)] \emph{Survival after a preparatory failure.} A preparatory large-field
characteristic function of the step is present in $Z_0$ (\cite[pp.~381--383]{B-CMP122b};
in \cite{B-CMP122a} such components ``do not satisfy (ii)''). New large fields are introduced
inside $Z_0$, the integral \cite[(1.71)]{B-CMP122b} is not performed, and $Z_0$ stays.
\item[(s2)] \emph{Survival as second class.} All preparatory functions are small-field and the
integral \cite[(1.71)]{B-CMP122b} is performed. The last decomposition
\cite[p.~378]{B-CMP122b} gives the second-class function $\chi^c_n(Y^{\sim -6})$ of the new field
$V_n$ restricted to $Y$, with $Y=Z_0$ (\cite[(1.72)]{B-CMP122b}). In the words of
\cite{B-CMP122b}: ``The new large field region is equal to $Z_k\cup\bigcup Y_i$, and the
operation $\mathbf{T}_k$ for a component $Y$ of the set $\bigcup Y_i$ has the form (1.102)''.
Then $Z_0$ stays.
\end{itemize}
The phrase \emph{the integral of $\mathbf{R}$ is performed on $Z_0$ at step $n$} means that the
outcome is (s0) or (s2). The phrase \emph{$Z_0$ is re-created at step $n$} means that
$\mathbf{R}$ was begun on $Z_0$ at step $n$ and $Z_0$ survived, that is, the outcome is (s1) or
(s2). The \emph{living components of level $n$} are the components of $Z_n$ that are not healed
at $n$; we write $Z_n^{\mathrm{liv}}$ for their union.

\emph{The step $n\to n+1$.} This step receives $Z_n^{\mathrm{liv}}$ and produces $Z_{n+1}$ as
follows.
\begin{itemize}
\item[--] \emph{Creation set.} Put
\begin{equation}\label{9.2:eq-construction-outcomes-1}
\mathcal{W}_{n+1}:=P'\cup Q'\cup R'\cup I'\subset\mathbb{L}_{n+1}.
\end{equation}
This is the set of cells meeting the large-field cubes created by the tests of the step; its
cells are the \emph{creation cells} of level $n+1$. Here
$P'$ (resp.\ $Q'$, $R'$) is the set of cells of $\mathbb{L}_{n+1}$ meeting the large-field cubes
of Ba{\l}aban's type $P$ (resp.\ $Q$, $R$) of \cite[\S3]{B-CMP119} (see also the cover in
\cite{B-CMP122b}), and $I'$ is the set of cells of $\mathbb{L}_{n+1}$ meeting the cubes of the
refusing-bond type $I$, a fourth type of large-field cube created by the tests of the step,
admitted by the family and absent from Ba{\l}aban's construction. The set $I'$ is empty on
Ba{\l}aban's member of the family. The
components of $\mathcal{W}_{n+1}$ are the \emph{creation components of level $n+1$}.
\item[--] \emph{(G1) Seeds.} The seed set is
\begin{equation}\label{9.2:eq-construction-outcomes-a}
\mathfrak{s}:=(\pi_n Z_n^{\mathrm{liv}})^{\sim 10}\cup P'^{\sim 6}\cup Q'^{\sim 4}
\cup R'^{\sim 1}\cup I'^{\sim 1}\subset Z_{n+1}.
\end{equation}
The layer counts are those of \cite[\S3]{B-CMP119}, in particular \cite[(3.20)]{B-CMP119}. For
the first term, \cite{B-CMP122b} writes ``we add ten layers''.
\item[--] \emph{(G2)} The set $Z_{n+1}$ is obtained from the seed set
\eqref{9.2:eq-construction-outcomes-a} and, on members of the family whose tests within one step
are carried out in several successive phases, from the analogous seed sets of the later phases,
by finitely many unions and box steps. On Ba{\l}aban's member of the family there is a single
phase.
\item[--] \emph{(G3)} Suppose a component $Z'$ of $Z_{n+1}$ contains exactly one set
$\pi_nZ_0$ with $Z_0$ a living component of level $n$, and contains no creation cell. Then
$Z'=S_n(Z_0)$ or $Z'$ is the box of $S_n(Z_0)$. In the words of \cite{B-CMP122b}: ``Such a
domain arises as a new large field region in our procedure, if no large fields are created in
a neighborhood of $Z$, more precisely in $S(Z)\setminus Z$''.
\item[--] \emph{Constituents.} Let $Z'$ be a component of $Z_{n+1}$. Its \emph{old
constituents} are the living components $Z_0$ of level $n$ with $\pi_nZ_0\subset Z'$; their
number is $m$. Its \emph{new constituents} are the creation components $W\subset Z'$ of level
$n+1$; their number is $q'$. Here and in the lemmas of this chapter the letters $m$, $q'$ and
$N$ count constituents; they do not denote the torus exponent $m$ of $\mathbb{T}_m$ or the $N$
of $\mathrm{SU}(N)$.

Each $\pi_nZ_0$ with $Z_0$ living is connected and lies in $Z_{n+1}$, by
\eqref{9.2:eq-construction-outcomes-a}. The same holds for each creation component. Hence each
such set lies in exactly one component of $Z_{n+1}$.

Every component has at least one constituent, that is, $N:=m+q'\ge 1$. Indeed, the core
$\mathcal{C}_{n+1}(Z')$ of $Z'$ is the union of the coarsened cores of the old constituents and
of the new constituents. It is non-empty by part (i) of the lemma on cores and connectors later
in this chapter.
\item[--] \emph{Events.} Let $Z'\in\operatorname{comp}(Z_{n+1})$. The step $n\to n+1$ is a
\emph{creation-or-join} for $Z'$ if $(m,q')\neq(1,0)$.

Otherwise $Z'$ has a single old constituent $Z_0$ and no creation cell. Indeed, a creation cell
in $Z'$ lies in a creation component, and that component lies in exactly one component of
$Z_{n+1}$, namely $Z'$; it would then be a new constituent, contrary to $q'=0$. By (G3) it
follows that $Z'\in\{S_nZ_0,\ \text{the box of }S_nZ_0\}$. In this case the step is a
\emph{re-creation step} for $Z'$ if $Z_0$ was re-created at $n$, and a \emph{quiet step}
otherwise.

A component re-created at $n$ and joined at $n+1$ falls under creation-or-join.
\end{itemize}
\end{definition}

\begin{definition}[Histories, ages and cores of a large-field component]\label{9.2:def-histories-ages}
Throughout, constituents (old and new), living components, quiet steps, re-creation, the
outcomes \textup{(s0)}--\textup{(s2)} of the operation, Ba{\l}aban's ten-layer operation $S_n$,
the box step and $100$-cubes are those of the construction family
of Definition~\ref{9.2:def-construction-outcomes}, whose outcomes are collected in
\eqref{9.2:eq-construction-outcomes-a}. Let $k$ be a level, let $Z_k$ be the large-field region
at level $k$, and let $X\in\operatorname{comp}(Z_k)$ be one of its components.

\emph{The lineage of $X$.} Put $\operatorname{comp}(\mathsf{Z}_k):=\{X\}$. For $n<k$, descending
in $n$, let $\operatorname{comp}(\mathsf{Z}_n)$ be the set of old constituents of the members of
$\operatorname{comp}(\mathsf{Z}_{n+1})$; these are living components of level $n$. Let
$\mathsf{Z}_n$ be their union, and put
\begin{equation*}
j:=\min\{n\ge 1:\ \mathsf{Z}_n\neq\emptyset\}.
\end{equation*}
A member of $\operatorname{comp}(\mathsf{Z}_n)$ is, by definition, an old constituent of a member
of $\operatorname{comp}(\mathsf{Z}_{n+1})$; hence $\operatorname{comp}(\mathsf{Z}_n)\neq\emptyset$
forces $\operatorname{comp}(\mathsf{Z}_{n+1})\neq\emptyset$. The set of levels $n\le k$ with
$\mathsf{Z}_n\neq\emptyset$ is therefore closed upwards, contains $k$, and has least element $j$:
\begin{equation}\label{9.2:eq-histories-ages-a}
\mathsf{Z}_n\neq\emptyset\quad\text{exactly for } j\le n\le k.
\end{equation}

\emph{Histories.} A \emph{history} $h$ of $X$ consists of the nested sequence
$(\mathsf{Z}_n)_{j\le n\le k}$ together with the internal data of the construction inside $X$ at
each level: Ba{\l}aban's small-field domains $\Omega_n\cap X$, the types and locations of the
large-field, preparatory and second-class functions that occur, and the labels
$\Lambda_n\cap X$ and $I_n(h)$, the latter a subset of the free bonds in $\mathsf{Z}_n$. We
write $\operatorname{Hist}(X)$ for the set of admissible histories; these are the admissible
sequences of \cite[(1.71)]{B-CMP122b}. For $h\in\operatorname{Hist}(X)$ let $j(h):=j$, the index
of the first large-field region contained in $X$ in the sense of \cite{B-CMP122b}; let
$a(h):=k-j(h)$ be the \emph{age} of $h$; and call the levels $j(h)\le n\le k$ the \emph{live
levels} of $h$. Put
\begin{align*}
\operatorname{births}(h)&:=\{(n,W):\ W \text{ a new constituent of a member of }
\operatorname{comp}(\mathsf{Z}_n)\},\\
\operatorname{recr}(h)&:=\{(n,Z):\ n<k,\ Z\in\operatorname{comp}(\mathsf{Z}_n),\
Z \text{ re-created at } n\},\\
\operatorname{anc}(h)&:=\{(n,Z):\ j(h)\le n\le k,\ Z\in\operatorname{comp}(\mathsf{Z}_n)\}.
\end{align*}
For $(n,W)\in\operatorname{births}(h)$ the set $W$ is a creation component of level $n$, created
by the step from level $n-1$ to level $n$. The set $\operatorname{recr}(h)$ is the union
$\operatorname{recr}_1(h)\cup\operatorname{recr}_2(h)$, where $\operatorname{recr}_1(h)$ and
$\operatorname{recr}_2(h)$ collect the re-creations with outcome \textup{(s1)} and with outcome
\textup{(s2)} respectively.

By \eqref{9.2:eq-histories-ages-a}, $\mathsf{Z}_n\neq\emptyset$, hence
$\operatorname{comp}(\mathsf{Z}_n)\neq\emptyset$, at each live level of $h$; each of the
$k-j(h)+1$ summands below is therefore at least $1$, and
\begin{equation}\label{9.2:eq-histories-ages-b}
\#\operatorname{anc}(h)=\sum_{n=j(h)}^{k}\#\operatorname{comp}(\mathsf{Z}_n)\ \ge\ k-j(h)+1
=a(h)+1.
\end{equation}

\emph{Sub-lineages and first indices.} For $(n,Z)\in\operatorname{anc}(h)$ the sub-lineage of $Z$
is defined downwards from $Z$ in the same way: its set of components at level $n$ is $\{Z\}$,
and below level $n$ it consists of the old constituents, level by level. We write $h|Z$ for the
restriction of $h$ to the sub-lineage of $Z$, and $j(Z)$ for its first index, defined recursively
by
\begin{equation*}
j(Z):=\min\bigl(\{j(Z_0):\ Z_0 \text{ an old constituent of } Z\}\cup
\{n:\ Z \text{ has a new constituent}\}\bigr).
\end{equation*}
The set on the right is non-empty because $Z$ has a constituent
(Definition~\ref{9.2:def-construction-outcomes}). The members of
$\operatorname{comp}(\mathsf{Z}_{j(h)})$ have no old constituent, since otherwise
$\mathsf{Z}_{j(h)-1}\neq\emptyset$, contrary to \eqref{9.2:eq-histories-ages-a}; hence each of
them has at least one new constituent. Moreover $j(X)=j(h)$. Indeed, every number entering the
recursion for $j(X)$ is a level of the lineage of $X$, so $j(X)\ge j(h)$; conversely, a member
$Y$ of $\operatorname{comp}(\mathsf{Z}_{j(h)})$ is joined to $X$ by a chain
$Y=Y_{j(h)},Y_{j(h)+1},\dots,Y_k=X$ in which each $Y_n$ is an old constituent of $Y_{n+1}$, so
that $j(Y_{n+1})\le j(Y_n)$ for each $n$, while $j(Y)\le j(h)$ because $Y$ has a new
constituent; hence $j(X)\le j(h)$.

\emph{Cores.} For $(n,Z)\in\operatorname{anc}(h)$ with old constituents $Z_0^{(v)}$ and new
constituents $W_i$, the \emph{core} of $Z$ is
\begin{equation*}
\mathcal{C}_n(Z):=\bigcup_v \pi_{n-1}\,\mathcal{C}_{n-1}\bigl(Z_0^{(v)}\bigr)\ \cup\
\bigcup_i W_i ,
\end{equation*}
where $\pi_{n-1}$ is the coarsening map; at a level where $Z$ has no old constituent,
$\mathcal{C}_n(Z)=\bigcup_i W_i$. Through a quiet step or a re-creation step, from a member $Z_0$
of level $n$ to its descendant $Z'$ of level $n+1$,
$\mathcal{C}_{n+1}(Z')=\pi_n\,\mathcal{C}_n(Z_0)$.

\emph{The clock.} In \cite{B-CMP122b} the corresponding number is introduced as follows:
\begin{quote}
the number $K$ is the smallest positive integer having the property that the domain $S^K(Z)$,
considered as a domain in the lattice of the scale $L^{-(j+K)}$, satisfies the conditions (i),
(ii), with $N = R_j$
\end{quote}
In this quotation the letters $K$, $N$, $j$ and $R_j$ are Ba{\l}aban's; in the definition below
the roles of $j$ and $R_j$ are taken by the level $n$ of $Z$ and the length $R_n$, and $S$ is
Ba{\l}aban's ten-layer operation, written $S_{n+t}$ at level $n+t$. For
$(n,Z)\in\operatorname{anc}(h)$ we define the \emph{fictitious quiet evolution} of $Z$ by
\begin{equation*}
Z^{[0]}:=Z,\qquad Z^{[t+1]}:=\text{the box of } S_{n+t}\bigl(Z^{[t]}\bigr)\quad(t\ge 0),
\end{equation*}
that is, $S_{n+t}$ followed by the box step of the construction where that step applies, the
coarsening ratios $\ell_m$ being continued by $L$ for $m$ beyond the last step of the
construction. The
\emph{clock} $K_n(Z)$ is the least integer $u\ge 0$ with the following property: in the window
of the $R_n$ steps ending at level $n+u$, formed by the actual lineage of $Z$ up to level $n$ and
by the fictitious evolution after level $n$, every step is quiet, and at each of the $R_n+1$
levels of the window the lineage's region is one component fitting in a $100$-cube. This is
the pair of conditions (i), (ii) of \cite{B-CMP122b} quoted above, with Ba{\l}aban's parameter
$N$ of those conditions taken equal to $R_n$, together with
the requirement that the region be a single component at every level of the window. That such an
integer $u$ exists, so that $K_n(Z)$ is finite, is part~(ii) of the lemma on the behaviour of the
clock (decrease of the clock along quiet steps, bounds on the clock in terms of the size of the
region and of its core, and vanishing of the clock only at a re-creation): the fictitious
evolution contracts regions.
\end{definition}

\begin{definition}[The large-field integral operation split by age]\label{9.2:def-integral-by-age}
Let $k$ be a level. Let $X\in\operatorname{comp}(Z_k)$ be a connected component of the large-field
region $Z_k$ on which the integral of Ba{\l}aban's renormalization operation $\mathbf{R}$ is
performed at step $k$.

\emph{The operation.} We write $\mathbf{T}'_k(X)$ for the integral operation displayed in
\cite[(1.71)]{B-CMP122b}. It contains the following pieces, all in the notation of
\cite{B-CMP122b}:
\begin{itemize}
\item the group integrations localised in $X$;
\item the Gaussian fluctuation integrals;
\item the Wilson term
\begin{equation*}
-A\bigl((g''_k(\cdot))^{-2}\zeta_1\Gamma_X,\;U''_{k+2}\bigr),
\end{equation*}
that is, minus the weighted Wilson action of Definition~\ref{1.1:def-wilson-action} at the
configuration $U''_{k+2}$, with the weight $(g''_k(\cdot))^{-2}\zeta_1\Gamma_X$ written in that
display ($g''_k(\cdot)$ a position-dependent coupling as in Definition~\ref{1.5:def-domains});
\item the terms $V(X^{\sim})$ and $E_k(X)$ of that display;
\item the normalisation constants.
\end{itemize}
Here $\zeta_1$ is the letter so written in that display and not the complex source shift
$\zeta$ of this book, and $E_k(X)$ is the term so written there, not one of the letters
$E_+,E_-$ of the glossary.
In \cite{B-CMP122b} this operation is described as a composition of two integrations. The first
is the integration in \cite[(1.70)]{B-CMP122b}, localised in the component $X$. The second is the
integration in \cite[(1.100)]{B-CMP122a}, taken together with the $\mathbf{T}''_k$-operation
there.

\emph{Decomposition over histories.} Let $\operatorname{Hist}(X)$ be the set of admissible
histories of $X$, as in Definition~\ref{9.2:def-histories-ages}. The sum over admissible
sequences in the display of \cite[(1.71)]{B-CMP122b} writes $\mathbf{T}'_k(X)$ as a sum of
operations $\tau_h$,
one for each $h\in\operatorname{Hist}(X)$. Here $\tau_h$ is the part of the operation that belongs
to the history $h$:
\begin{equation}\label{9.2:eq-integral-by-age-1}
\mathbf{T}'_k(X)=\sum_{h\in\operatorname{Hist}(X)}\tau_h .
\end{equation}
For a family $\mathcal{A}\subset\operatorname{Hist}(X)$ of histories we put
\begin{equation}\label{9.2:eq-integral-by-age-2}
\mathbf{T}'_{k,\mathcal{A}}(X):=\sum_{h\in\mathcal{A}}\tau_h .
\end{equation}
For $h\in\operatorname{Hist}(X)$ let $a(h)=k-j(h)$ be its age
(Definition~\ref{9.2:def-histories-ages}). For an integer $a\ge 0$ we define the part of
$\mathbf{T}'_k(X)$ of age $a$ and the part of age at least $a$ by
\begin{align}
\mathbf{T}'_{k,a}(X)&:=\mathbf{T}'_{k,\{h:\,a(h)=a\}}(X),
\label{9.2:eq-integral-by-age-3}\\
\mathbf{T}'_{k,\ge a}(X)&:=\sum_{a'\ge a}\mathbf{T}'_{k,a'}(X).
\label{9.2:eq-integral-by-age-4}
\end{align}

\emph{Domain of action.} The operators \eqref{9.2:eq-integral-by-age-1}--\eqref{9.2:eq-integral-by-age-4}
act on bounded functions $F$ of the variables not integrated in \cite[(1.71)]{B-CMP122b}.
For such an $F$, $\sup|F|$ denotes the supremum of $|F|$ over its arguments.

\emph{The parameter domain $\mathfrak{P}$.} The operation is considered in two settings.
\begin{itemize}
\item[(i)] The operation is that of a single run of the correlation theorem. Then
$\mathfrak{P}$ is the product of the complex tube $\mathcal{T}$ of external fields $(U,J)$ and
the weight domain $\mathcal{W}_k$, both as fixed in that theorem.
\item[(ii)] The operation is that of setting (i) in the form in which it depends, in addition, on
dial variables, and the fixed list of conditions on the parameters under which that form is
defined is assumed. Then
\begin{equation*}
\mathfrak{P}=\mathbb{D}_\delta\times\mathcal{T}\times\mathcal{W}_k ,
\end{equation*}
where $\mathcal{T}$ is the tube of setting (i) and $\mathbb{D}_\delta$ is the polydisc in the
dial variables on which this form of the operation is defined.
\end{itemize}
Every bound stated in what follows for these operations holds pointwise in $\mathfrak{P}$. Its
constants do not depend on the point of $\mathfrak{P}$.
\end{definition}

\begin{definition}[The pre-payment numbers at a given charge constant]\label{9.2:def-prepayment-numbers}
We use the following notation.
\begin{itemize}
\item the dimension $d$, the number $M$ and the lengths $R_n$ are those of
Notation~\ref{9.2:not-parameters-couplings};
\item $\mathsf{P}_n=p_0(g_n)$ is the degree function $p_0(\cdot)$ evaluated at the running coupling
of level $n$;
\item $\mathsf{D}_n$ is the deposit attached to one large-field cube at level $n$, in its linear-size
form;
\item $d'(\cdot)$ is the unconstrained linear size;
\item histories $h$, their sets of ancestors $\operatorname{anc}(h)$, the events of a history and the
first index $j(\cdot)$ of a member are as in Definition~\ref{9.2:def-histories-ages}.
\end{itemize}

Fix a constant $C>0$ and put, for every level $n$,
\begin{equation*}
c_n:=C\,M^dR_n^{d+1}.
\end{equation*}
For a history $h$ and $(n,Z)\in\operatorname{anc}(h)$, the \emph{charge of $Z$ at level $n$} is
\begin{equation*}
\mathsf{c}_n(Z):=c_n\sqrt{d}\,|Z|,
\end{equation*}
where $|Z|$ denotes the number of cells of level $n$ of which $Z$ consists.

The \emph{pre-payment numbers} $\kappa_n(Z)=\kappa^{[C]}_n(Z;h)$, $(n,Z)\in\operatorname{anc}(h)$, are
defined recursively over the events of the history $h$. When $C$ and $h$ are fixed we write
$\kappa_n(Z)$ for $\kappa^{[C]}_n(Z;h)$; these numbers carry the argument $Z$ and are not the
auxiliary parameters $\kappa$, $\kappa_1$ of $\mathfrak{p}$. The recursion is taken according to the
event at the step $n-1\to n$ at which the member $Z$ of level $n$ arises. Three kinds of event are
distinguished.
\begin{itemize}
\item[$(\kappa1)$] Creation or join at the step $n-1\to n$. Here $Z$ has old constituents
$Z_0^{(v)}$, $1\le v\le m$, and new constituents $W_i$, $1\le i\le q'$, with $(m,q')\neq(1,0)$.
In this clause $m$ and $q'$ denote the numbers of old and new constituents, not the torus index
$m$. Put $j(v):=j(Z_0^{(v)})$.
\item[$(\kappa2)$] Quiet step. Here $Z_0$ is the member of level $n-1$ from which $Z$ arises.
\item[$(\kappa3)$] Re-creation step. Here $Z$ arises from the member $Z_0$ of level $n-1$, and
$Z_0$ is re-created at $n-1$.
\end{itemize}
In these three cases respectively we set
\begin{equation}\label{9.2:eq-prepayment-numbers-a}
\begin{aligned}
&(\kappa1)\quad \kappa_n(Z):=\sum_{v=1}^{m}\bigl[\kappa_{n-1}(Z_0^{(v)})+2\mathsf{P}_{j(v)}\bigr]
+\sum_{i=1}^{q'}\bigl[\mathsf{D}_n\bigl(d'(W_i)+1\bigr)+2\mathsf{P}_n\bigr]\\
&\qquad\qquad\qquad -\mathsf{c}_n(Z)-2\mathsf{P}_{j(Z)},\\
&(\kappa2)\quad \kappa_n(Z):=\kappa_{n-1}(Z_0)-\mathsf{c}_n(Z),\\
&(\kappa3)\quad \kappa_n(Z):=\kappa_{n-1}(Z_0)+\mathsf{P}_{n-1}-\mathsf{c}_n(Z).
\end{aligned}
\end{equation}
The recursion starts at the members with no old constituent. For these, clause $(\kappa1)$ of
\eqref{9.2:eq-prepayment-numbers-a} applies with $m=0$.

In clause $(\kappa1)$ the subtracted term $2\mathsf{P}_{j(Z)}$ cancels exactly one of the $m+q'$
linear terms, namely the terms $2\mathsf{P}_{j(v)}$, $1\le v\le m$, and the $q'$ copies of
$2\mathsf{P}_n$.

The numbers are normalised as follows. In \cite{B-CMP122b} the factor connected with a region $Z$
is written in the form
\begin{equation*}
\exp\bigl(-\kappa_j(Z)-2p_0(g_{j(Z)})\bigr).
\end{equation*}
The definition \eqref{9.2:eq-prepayment-numbers-a} follows this form. Its three clauses correspond
respectively to \cite[(1.85)]{B-CMP122b}, to \cite[(1.83)]{B-CMP122b} and to
\cite[p.~386]{B-CMP122b}.

At a re-creation step, \cite[p.~386]{B-CMP122b} records the new factor $\exp(-p_0(g_j))$ and
defines
\begin{equation*}
\kappa_{j+1}(Z)=p_0(g_j)-O(1)M^dR^{d+1}_{j+1}\,d'_{j+1}(Z).
\end{equation*}
The right-hand side of this definition consists of the term $p_0(g_j)$ and a charge term. Clause
$(\kappa3)$ of \eqref{9.2:eq-prepayment-numbers-a} contains the corresponding terms
$\mathsf{P}_{n-1}$ and $-\mathsf{c}_n(Z)$ and, in addition, the summand $\kappa_{n-1}(Z_0)$ of the
previous level.
\end{definition}

\begin{proposition}[Summation over admissible histories at a cost per cell]
\label{9.2:prop-history-summation}
Let $k$ be a level and let $X\subset\mathbb{L}_k$ be connected, where $\mathbb{L}_n$ denotes the set
of cells of level $n$ and $\pi_n$ the coarsening map from level $n$ to level $n+1$, as in
Definition~\ref{9.2:def-histories-ages}. A \emph{labelled nested sequence ending at $X$} consists
of a level $j(\eta)\le k$ and a sequence $\eta=(\mathsf{Z}_n,\sigma_n)_{j(\eta)\le n\le k}$ with
\begin{equation*}
\emptyset\neq\mathsf{Z}_n\subset\mathbb{L}_n,\qquad \mathsf{Z}_k=X,\qquad
\pi_n\mathsf{Z}_n\subset\mathsf{Z}_{n+1},\qquad \sigma_n=(\sigma_n^{\mathrm{pr}},I_n),
\end{equation*}
where $\sigma_n^{\mathrm{pr}}$ is the label datum of Ba{\l}aban's large-field step in
$\mathsf{Z}_n$ and $I_n$ is a set of
free bonds in $\mathsf{Z}_n$; there are at most $\beta_n|\mathsf{Z}_n|$ free bonds in
$\mathsf{Z}_n$. Here $\beta_n$ and $\upsilon_n$ are the numbers fixed in
Notation~\ref{9.2:not-parameters-couplings}, $s(n)$ is the number of label types per test cube at
level $n$, and we put
\begin{equation*}
c_A(n):=\bigl(L^{2d}+s(n)(LM/M_2)^d+10\bigr)\log 2+1,\qquad c_A^{\flat}(n):=c_A(n)+1 .
\end{equation*}
Assume the following two statements.
\begin{itemize}
\item[(a)] The count of label data: $\sigma_n^{\mathrm{pr}}$ takes at most
$2^{(s(n)(LM/M_2)^d+9)|\mathsf{Z}_n|}$ values.
\item[(b)] The lemma on the sum over labelled nested sequences: if $\beta_n\upsilon_n\le 1$ for
all $n\le k$, then for every function $\Psi\ge 0$ of labelled nested sequences ending at $X$
\begin{equation}\label{9.2:eq-history-summation-a}
\sum_{\eta}\Psi(\eta)\le\sup_{\eta}\Bigl[\Psi(\eta)\prod_{n=j(\eta)}^{k}
\upsilon_n^{-|I_n|}\,e^{c_A^{\flat}(n)|\mathsf{Z}_n|}\Bigr].
\end{equation}
\end{itemize}
Then every admissible history $h\in\operatorname{Hist}(X)$ is a labelled nested sequence ending
at $X$, with sets $\mathsf{Z}_n(h)$ and labels $\sigma_n(h)=(\sigma_n^{\mathrm{pr}}(h),I_n(h))$, and,
if $\beta_n\upsilon_n\le1$ for all $n\le k$, then for every family
$\mathcal{A}\subset\operatorname{Hist}(X)$ and all numbers $m(h)\ge0$
\begin{equation}\label{9.2:eq-history-summation-1}
\sum_{h\in\mathcal{A}}m(h)\le\sup_{h\in\mathcal{A}}\Bigl[m(h)\prod_{n=j(h)}^{k}
\upsilon_n^{-|I_n(h)|}\,e^{c_A^{\flat}(n)|\mathsf{Z}_n(h)|}\Bigr].
\end{equation}
\end{proposition}

It is stated in \cite{B-CMP122b} that the summations over the admissible sequences can be replaced
by the factors $\exp O(1)(MR_j)^{-d}|Z_j|$, in the notation there; the bound
\eqref{9.2:eq-history-summation-1} is the form of that summation used in this chapter, with the
label $I_n$ of free bonds carried along. The hypothesis $\beta_n\upsilon_n\le1$ for all $n\le k$
is the free-bond condition among the standing conditions on the parameters under which the
estimates of this chapter are made.

\begin{proof}
Let $h\in\operatorname{Hist}(X)$. We check the four conditions defining a labelled nested
sequence ending at $X$.

First, by \eqref{9.2:eq-histories-ages-a} the set $\mathsf{Z}_n(h)$ is non-empty exactly for $n$
in the interval $[j(h),k]$; so $h$ gives a sequence indexed by $j(h)\le n\le k$ with
$\emptyset\ne\mathsf{Z}_n(h)\subset\mathbb{L}_n$, and we set $j(\eta):=j(h)$.

Second, $\mathsf{Z}_k(h)=X$, since $h$ is a history of the component $X$.

Third, let $j(h)\le n<k$ and let $Z_0$ be a member of
$\operatorname{comp}(\mathsf{Z}_n(h))$. By Definition~\ref{9.2:def-histories-ages}, $Z_0$ is an
old constituent of a member of $\operatorname{comp}(\mathsf{Z}_{n+1}(h))$, and $\pi_nZ_0$ lies in
that member. Taking the union over all members $Z_0$ gives
$\pi_n\mathsf{Z}_n(h)\subset\mathsf{Z}_{n+1}(h)$.

Fourth, the inner sum of the integral operation $\mathbf{T}'_k(X)$ at level $n$ runs over the
traces on $X$ of the two large-field label sets $\Omega_n$, $\Lambda_n$ of Ba{\l}aban's step and
of the step's set of free bonds at level $n$. The first two traces constitute the label datum
$\sigma_n^{\mathrm{pr}}(h)$ in $\mathsf{Z}_n(h)$; we write $I_n(h)$ for the third, which is a set
of free bonds in $\mathsf{Z}_n(h)$. Hence the data of $h$ at level $n$ are of the
form $\sigma_n(h)=(\sigma_n^{\mathrm{pr}}(h),I_n(h))$.

Thus $h$ is a labelled nested sequence ending at $X$; hypothesis~(a) is the count of label data
under which the lemma in hypothesis~(b) is stated, so (b) applies to non-negative functions of
such sequences.

Now let $\mathcal{A}\subset\operatorname{Hist}(X)$ and $m(h)\ge0$ be given, and assume
$\beta_n\upsilon_n\le1$ for all $n\le k$. Define $\Psi:=m\cdot\mathbf{1}_{\mathcal{A}}$ on
admissible histories, that is $\Psi(h)=m(h)$ for $h\in\mathcal{A}$ and $\Psi(h)=0$ for
$h\in\operatorname{Hist}(X)\setminus\mathcal{A}$, and extend $\Psi$ by $0$ to the labelled nested
sequences ending at $X$ that are not admissible histories. Then $\Psi\ge0$, and by
hypothesis~(b) the bound \eqref{9.2:eq-history-summation-a} holds for $\Psi$. Its left-hand side
equals $\sum_{h\in\mathcal{A}}m(h)$, because $\Psi$ vanishes off $\mathcal{A}$. In its right-hand
side, every $\eta$ outside $\mathcal{A}$ contributes the value $0$ to the supremum, while every
$h\in\mathcal{A}$ contributes
\begin{equation*}
m(h)\prod_{n=j(h)}^{k}\upsilon_n^{-|I_n(h)|}\,e^{c_A^{\flat}(n)|\mathsf{Z}_n(h)|}\ \ge 0 .
\end{equation*}
Hence the supremum over all $\eta$ equals the supremum over $h\in\mathcal{A}$ of these numbers
(both sides being $0$ if $\mathcal{A}$ is empty), and \eqref{9.2:eq-history-summation-a} becomes
\eqref{9.2:eq-history-summation-1}.
\end{proof}

\begin{proposition}[The history-wise modulus majorant. Proof outstanding]
\label{9.2:prop-modulus-majorant}
We use the construction family and the outcomes
\eqref{9.2:eq-construction-outcomes-a}, and the notation of
Definitions~\ref{9.2:def-construction-outcomes}, \ref{9.2:def-histories-ages}
and~\ref{9.2:def-integral-by-age}. Here $\operatorname{Hist}(X)$ is the set of admissible
histories of a component $X$, and $\tau_h$ is the part of the integral operation belonging to
the history $h$. Setting~(i) is the setting of the correlation theorem. Setting~(ii) is the dial
polydisc $\mathbb{D}_\delta$ under the standing conditions of the dial.

There is a number $C_0\ge 1$, the \emph{charge constant of the majorant}, with the following
properties:
\begin{itemize}
\item in setting~(i) it depends only on $(d,L,M,M_2,\kappa,C_\sharp)$, and in setting~(ii) it
depends in addition on $C_{L'}$;
\item in the order of choice it is fixed before $c_a$ and before $\gamma$.
\end{itemize}
Further, let $k$ be a level, let $X\in\operatorname{comp}(Z_k)$ be a component on which the
integral of $\mathbf{R}$ is performed at step $k$, and let $h\in\operatorname{Hist}(X)$. Then
there is a number $F^{\natural}(h)\ge 0$, independent of $F$ and of the point of
$\mathfrak{P}$, such that clauses (a)--(c) below hold.
\begin{itemize}
\item[(a)] \emph{Majorant.} For every bounded $F$ and every point of $\mathfrak{P}$,
\begin{equation}\label{9.2:eq-modulus-majorant-a}
\lvert \tau_h F\rvert \le e\cdot \sup\lvert F\rvert\cdot F^{\natural}(h).
\end{equation}
\item[(b)] \emph{Factors.} $F^{\natural}(h)$ is a product taken over two indices: the live
levels $n$ of $h$, and the members $Z\in\operatorname{comp}(\mathsf{Z}_n)$. The factors are the
groups $(\alpha)$--$(\delta)$ of non-negative factors listed below. Each of them is a number free
of the complex coupling parameter, and each is read at the real couplings $g_n$. The real part
of every raw Wilson coefficient is at least $(1-\theta_W)$ times its real value, where
\begin{equation*}
1-\theta_W\ge \tfrac34 .
\end{equation*}
In setting~(i) this holds by clause~(c) of the Wilson-coefficient lemma of the correlation
theorem and by step~(3) of the proof of the modulus lemma of that theorem. In setting~(ii) it
holds by the
corresponding clause of the standing conditions of the dial. The groups are these.
\begin{itemize}
\item[$(\alpha)$] Let $W$ be a new constituent of $Z$, that is, $(n,W)\in\operatorname{births}(h)$.
For each such $W$, the group contains the majorants of the moduli of those large-field factors
of the step $n-1\to n$ that are located in $W$. These are, first, the factors of the
characteristic functions that create the $P$-, $Q$- and $R$-cubes of $W$
\cite{B-CMP122b}. Second, for every refusing bond of the label $I_n(h)$ located in $W$, they
include the full factor of that bond, which is at most $\upsilon_n^2$.
\item[$(\beta)$] The positive exponentials attached to the level $n$ and the member $Z$. The
first of them are Ba{\l}aban's normalization constants and counterterms, ``to which every
large field domain $Z_j$ contributes the constant $O(1)\log g_j^{-2}\lvert Z_j\rvert$''
\cite{B-CMP122b}. Next come the factors ``$\exp O(1)\lvert Z_j\cap \Omega_j\rvert$''
\cite{B-CMP122b}, and the bound \cite[(1.69)]{B-CMP122b} of the terms $V$. The group further
contains the shell and boundary terms per cell, and the enlargements in the list of standing
conditions of the correlation theorem. In setting~(ii) it also contains the per-cell factors
that the standing conditions of the dial add to the positive exponentials on
$\mathbb{D}_\delta$.
\item[$(\gamma)$] This group is present if $n<k$ and $Z$ was re-created at $n$. For kind
\textup{(s1)} of \eqref{9.2:eq-construction-outcomes-a}, it is the majorant of the modulus of
Ba{\l}aban's preparatory factor present in $Z$ at step $n$. For kind \textup{(s2)}, it is a
product of two numbers. One is the number that replaces the second-class function
$\chi^c_n(Z^{\sim-6})$ in the majorant (clause~(c3)). The other is the prefactor $e$ of the
inner performed operation $\mathbf{T}'_n(Z)$ (clause~(c2)).
\item[$(\delta)$] All remaining factors are bounded by $1$ and are omitted. These are the
characteristic functions, the $\delta$-functions, and the group-valued integrations: ``All the
integrals with respect to the group valued variables in the definition (1.71) are estimated by
1'' \cite[p.~383]{B-CMP122b}; here (1.71) is \cite[(1.71)]{B-CMP122b}. The Gaussian integrals
are also bounded by $1$, by positivity
\cite{B-CMP122b}.
\end{itemize}
The product structure follows the renormalization steps. In the words of \cite{B-CMP122b},
``(1.79) is a product of factors coming from the successive renormalization steps. The product
of factors coming from the first $j$ steps is connected with the region $Z_j$, and can be
factorized in components of $Z_j$''. Here (1.79) refers to \cite[(1.79)]{B-CMP122b}.
\item[(c)] \emph{Nesting: performed inner operations.} Let
$(n,Y)\in\operatorname{recr}_2(h)$, that is, a survival as second class at a level $n<k$. Then
\textup{(c1)}--\textup{(c3)} hold.
\begin{itemize}
\item[(c1)] The operation $\mathbf{T}'_n(Y)$ performed at step $n$ is a compositional factor
of $\tau_h$. The reason is as follows. After step $n$, the operation attached to $Y$ is
\begin{equation*}
\mathbf{T}_n(Y)=\chi^c_n(Y^{\sim-6})\,\chi_{n,\Lambda}\,\delta_G(V'_n)\,\mathbf{T}'_n(Y)
\end{equation*}
by \cite[(1.102)]{B-CMP122b}. The operations $\mathbf{T}^{(m)}$ of the later steps (in the
letters of \cite[(2.20)--(2.21)]{B-CMP119}, where $m$ is a running index of that paper) compose
with it, by \cite[(2.20)--(2.21)]{B-CMP119} and \cite[(1.72)]{B-CMP122b}; in the latter the
composed operations appear as the factor
\begin{equation*}
\Bigl(\sum_{\{\Omega^c_j,\,Z_j\}}\mathbf{T}''_k(Z_k)\Bigr),
\end{equation*}
summed over the admissible sequences; here $j$ and $k$ are running indices of
\cite[(1.72)]{B-CMP122b}, not the levels of this proposition. Now consider any later level
$k'\le k$ at which the
integral of $\mathbf{R}$ is performed on a member $X'$ of the lineage such that $X'$ contains
the descendant of $Y$. At such a level, $\mathbf{T}'_{k'}(X')$ is ``a composition of the
integration in (1.70) localized in this component, and the integration together with the
$\mathbf{T}''$-operation in (1.100) [IV]'' \cite{B-CMP122b}. Here (1.70) is
\cite[(1.70)]{B-CMP122b}, [IV] denotes
\cite{B-CMP122a}, and (1.100) is \cite[(1.100)]{B-CMP122a}. This composition contains
$\mathbf{T}_n(Y)$. The admissible sequences of $Y$ are the restrictions $h|Y$ of the admissible
histories $h$ of $X$ to the sub-lineage of $Y$. Moreover, $\tau_h$ contains
$\tau^{(Y)}_{h|Y}$, the corresponding part of $\mathbf{T}'_n(Y)$.
\item[(c2)] $F^{\natural}(h)$ contains the factor-by-factor majorant of $\tau^{(Y)}_{h|Y}$,
jointly indexed by $h|Y$. That is, the groups $(\alpha)$, $(\beta)$, $(\gamma)$ indexed by the
levels $\le n$ and the members of the sub-lineage of $Y$ are at the same time the factors of the
inner operation. In addition, the group $(\gamma)$ of $(n,Y)$ contains the prefactor of the
inner operation itself: the constant $e$ of \textup{(a)}, applied to $\mathbf{T}'_n(Y)$. In
Ba{\l}aban's notation this is the prefactor of \cite[(1.75)]{B-CMP122b}. Thus $F^{\natural}(h)$
contains one factor $e$ for each performed inner operation of the lineage, that is, one for
each element of $\operatorname{recr}_2(h)$, and no other cost.
\item[(c3)] Let $(m,Y_0)\in\operatorname{recr}_2(h)$, and let $\chi^c_m(Y_0^{\sim-6})$ be its
second-class function. The number that replaces this function is delivered by the second-class
bound of this chapter, at its clause $(\alpha)$. That clause is applied with $k'$ equal to the
first level $k'\in(m,k]$ at which the integral of $\mathbf{R}$ is performed on a member $X'$ of
the lineage such that $X'$ contains the descendant of $Y_0$. Such a level $k'$ exists, since the
integral is
performed on $X$ at level $k$. By \textup{(c2)}, that inner majorant is part of
$F^{\natural}(h)$. For the last element of $\operatorname{recr}_2(h)$, which has no later
performed inner operation, $k'=k$, and the clause is applied at the outer operation itself.
\end{itemize}
\end{itemize}
\end{proposition}

\noindent\emph{Proof outstanding.}

\noindent\emph{Route.} In setting~(i), clauses (a), (b) and (c) are to be read off
steps (1)--(5) of the proof of
the modulus lemma of the correlation theorem. That proof bounds the modulus of
$\mathbf{T}'_k(X)F$ factor by factor, under the standing conditions of the correlation theorem.
In setting~(ii), the clauses are to be read off the proof of the dial-domain modulus
proposition, which is the modulus bound for $\mathbf{T}'$ on $\mathbb{D}_\delta$. Its proof
inspects the factors of the operation one by one.

The numerical bounds for the groups $(\alpha)$, $(\beta)$, $(\gamma)$ of \textup{(b)} are the
inputs assembled in the assembly lemma of this chapter. The compositional facts in
\textup{(c1)} are those of \cite[(1.102), (1.72)]{B-CMP122b}, \cite[(2.20)--(2.21)]{B-CMP119}
and \cite[(1.100)]{B-CMP122a}, used as cited there.

Beyond these facts, a proof has to establish the following statement about the majorant built
in the proof of the modulus lemma. Consider a lineage containing operations
$\mathbf{T}'_n(Y)$ performed at levels $n<k$ with outcome \textup{(s2)}. Then the
factor-by-factor majorant of $\mathbf{T}'_k(X)F$ contains the factor-by-factor majorants of
those operations, with three properties:
\begin{itemize}
\item they are jointly indexed by the inner admissible sequences $h|Y$, or taken as blocks, in
the sense of the remark on the blockwise reading that follows;
\item exactly one prefactor $e$ is present per performed operation;
\item no other cost is incurred.
\end{itemize}
Should the proof of the modulus lemma yield an absolute prefactor $e^{\mathfrak{e}_0}$, with
$\mathfrak{e}_0>1$, per performed operation, then the condition on $\mathsf{P}_n$ in the table of
conditions of this chapter becomes $\mathsf{P}_n\ge \mathfrak{e}_0$, and nothing else changes.
Should it yield no inner prefactor, that condition is idle.

\begin{remark}[Blockwise reading of the nesting clause]
Suppose the proof of the modulus lemma bounds an inner performed operation
$\mathbf{T}'_n(Y)$ as a block, instead of jointly indexed by $h|Y$. The bound is then $e$ times
the operation's own bound, which has the form of the supremum over the inner admissible
sequences times the sequence-sum factor, with $e$ the prefactor of \cite[(1.75)]{B-CMP122b}. In
that case, nothing that is drawn from
Proposition~\ref{9.2:prop-modulus-majorant} in this chapter changes, for three reasons.
\begin{itemize}
\item The inner supremum of the inner comparison product is bounded by the same product
$D(h|Y)\,I(h|Y)\,E^{[C]}(h|Y)$ of deposit, income and charge factors as in the assembly lemma
of this chapter, at the charge constant $C$ of that lemma, that is, by that product restricted
to the sub-lineage of $Y$.
\item The inner sequence-sum factor $\exp O(1)(MR_j)^{-d}\lvert Z_j\rvert$ is carried, with its
constant made explicit, as $c_A(n)$ per cell. It is therefore one of the per-cell charges of the
counting bound of this chapter.
\item Applying the outer bound over the larger index set, that is, the decomposition of the
integral operation into the parts $\tau_h$ of Definition~\ref{9.2:def-integral-by-age} together
with \eqref{9.2:eq-modulus-majorant-a}, only weakens the bound.
\end{itemize}
\end{remark}

\begin{lemma}[The creation deposit in linear-size form]\label{9.2:lem-creation-deposit}
Let $h$ be a history in the sense of Definition~\ref{9.2:def-histories-ages}, and let the quantities
$\mathsf{P}_n$, $\mathsf{Q}_n$, $\mathsf{D}_n$ and $\upsilon_n$ be those of
Notation~\ref{9.2:not-parameters-couplings}. The hypotheses are clause (E2) of
Lemma~\ref{9.2:lem-linear-size-properties} and the creation-deposit bound, in the following form.
Assume that for every $(n,W)\in\operatorname{births}(h)$, after one factor $\upsilon_n$ per refusing bond
of $I_n(h)$ located in $W$ has been set aside, the remaining factors of the group $(\alpha)$ of the
history-wise modulus majorant $F^{\natural}(h)$ (Proposition~\ref{9.2:prop-modulus-majorant}, clause
(b), whose proof is outstanding) located in $W$ multiply to at most $e^{-\mathsf{Q}_n|W|}$.

Then, for every $(n,W)\in\operatorname{births}(h)$ such that
\begin{equation}\label{9.2:eq-creation-deposit-a}
\mathsf{Q}_n\ \ge\ 4\,\mathsf{P}_n ,
\end{equation}
one has
\begin{equation}\label{9.2:eq-creation-deposit-1}
e^{-\mathsf{Q}_n|W|}\ \le\ \exp\bigl(-\mathsf{D}_n\,(d'(W)+1)-2\,\mathsf{P}_n\bigr).
\end{equation}
Consequently, under \eqref{9.2:eq-creation-deposit-a} and after the set-aside described above, the
remaining factors of the group $(\alpha)$ located in $W$ multiply to at most the right-hand side of
\eqref{9.2:eq-creation-deposit-1}.
\end{lemma}

\begin{proof}
Let $(n,W)\in\operatorname{births}(h)$ and assume \eqref{9.2:eq-creation-deposit-a}. By
Definition~\ref{9.2:def-histories-ages}, $W$ is a finite non-empty connected subset of $\mathbb{L}_n$,
so clause (E2) of
Lemma~\ref{9.2:lem-linear-size-properties} therefore applies to $A=W$ and gives
$d'(W)\le\sqrt d\,(|W|-1)$. Rearranging, and using $\sqrt d\ge1$, we obtain
\begin{equation*}
|W|\ \ge\ 1+\frac{d'(W)}{\sqrt d}\ \ge\ \frac{d'(W)+1}{\sqrt d}.
\end{equation*}
Multiplying by $\tfrac12\mathsf{Q}_n$, which is non-negative, and using
$\mathsf{D}_n=\mathsf{Q}_n/(2\sqrt d)$ (Notation~\ref{9.2:not-parameters-couplings}), we get
\begin{equation*}
\tfrac12\,\mathsf{Q}_n|W|\ \ge\ \frac{\mathsf{Q}_n\,(d'(W)+1)}{2\sqrt d}\ =\ \mathsf{D}_n\,(d'(W)+1).
\end{equation*}
Since $|W|\ge1$, hypothesis \eqref{9.2:eq-creation-deposit-a} yields
\begin{equation*}
\tfrac12\,\mathsf{Q}_n|W|\ \ge\ \tfrac12\,\mathsf{Q}_n\ \ge\ 2\,\mathsf{P}_n.
\end{equation*}
We split $\mathsf{Q}_n|W|$ into two halves and apply the two lower bounds, one to each half:
\begin{equation*}
e^{-\mathsf{Q}_n|W|}
= e^{-\frac12\mathsf{Q}_n|W|}\,e^{-\frac12\mathsf{Q}_n|W|}
\le e^{-\mathsf{D}_n(d'(W)+1)}\,e^{-2\mathsf{P}_n}.
\end{equation*}
This is \eqref{9.2:eq-creation-deposit-1}. The final assertion of the lemma follows by combining
\eqref{9.2:eq-creation-deposit-1} with the creation-deposit bound assumed above.
\end{proof}

\begin{remark}
The creation-deposit bound assumed in Lemma~\ref{9.2:lem-creation-deposit} rests on the following
facts. The test cubes are nested in the cells, and each cell of the cover meets a bad cube. Each of the
three per-cube factors of \cite[p.~381]{B-CMP122b} is at least $\mathsf{Q}_n$ per bad cube. A refusing
bond lies in at
most two test cubes, and these two are adjacent; hence it lies in a single creation component. After the
set-aside, such a bond keeps a factor of at most
\begin{equation*}
\upsilon_n=\exp\bigl(-\gamma_0\,\mathfrak{b}_n^2/(128\,\Theta_1^2(1+\beta_0)^2)\bigr).
\end{equation*}
The refusing bonds and the bonds of the $R_n$-cubes are disjoint coordinates of one fluctuation
quadratic form,
and this form is bounded below by $\gamma_0$. Finally, the passage from the coupling $g_{n-1}$ of the
creating step to $g_n$ costs the factor $(1+\beta_0)^{-2}$ of \cite[(2.7)]{B-CMP119}, which is included
in $\mathsf{Q}_n$.
\end{remark}

\begin{remark}[The second-class income within the majorant]\label{9.2:rem-second-class-income}
Let $X\in\operatorname{comp}(Z_k)$ be a component on which the integral of $\mathbf{R}$ is performed at
step $k$, and let $h\in\operatorname{Hist}(X)$, as in the history-wise modulus majorant
(Proposition~\ref{9.2:prop-modulus-majorant}, whose proof is outstanding). Let
$(n,Y)\in\operatorname{recr}_2(h)$, so that $n\ge1$ and $Y$ is a component on which $\mathbf{R}$ was
begun at step $n$ with a second-class outcome in the sense of
Definition~\ref{9.2:def-construction-outcomes}.

The factor of the group $(\gamma)$ of $F^{\natural}(h)$ attached to the pair $(n,Y)$ is the income
$e^{-E_{\mathrm{2cl}}(n)}$ of the statement that a second-class outcome pays a tag (in dimension $d=4$,
with $M\ge 4M_2$), whose price is the raw Wilson term restricted to one set $\Delta'\subset Y^{\sim-4}$
of the level-$n$ stratum of $Y$. That statement is applied in the majorant of the operation under which
the level-$n$ variables of $Y$ are integrated, which has two forms: the integral
$\mathbf{T}'_{k'}(X')1$ of \cite[(1.71)]{B-CMP122b} at the first step $k'>n$ at which the integral of
$\mathbf{R}$ is performed on a component $X'$ containing the descendant of $Y$, and the final integral
if there is no such step. Its hypotheses include a lower bound $c_W\ge\tfrac34$ on the real part of the
Wilson coefficient over the level-$n$ stratum of $Y$ in units of $g_n^{-2}$ and a ceiling
$g_n\le g_\bullet$ on the coupling; the exponent $E_{\mathrm{2cl}}(n)$ is bounded below in terms of
$c_W$ and of a counting number $N_Y$, and under the ceiling $E_{\mathrm{2cl}}(n)\ge 2\mathsf{P}_n$.

In the use made of it in $F^{\natural}(h)$, the hypotheses of that statement hold, the income
$e^{-E_{\mathrm{2cl}}(n)}$ enters $F^{\natural}(h)$ once, and no other factor of $F^{\natural}(h)$
draws on that price.
\end{remark}

\begin{proof}
The argument has six points.

(i) \emph{The place of use.} The second-class pairs at which the statement is used inside
$F^{\natural}(h)$ are the elements $(n,Y)$ of $\operatorname{recr}_2(h)$; for each of them the use is
at the first of its two forms, the integral $\mathbf{T}'_{k'}(X')1$, with $k'\le k$, as recorded in
clause (c3) of Proposition~\ref{9.2:prop-modulus-majorant}, whose proof is outstanding. By
clause (c2) of the same proposition, the majorant of this integral, and with it the income
$e^{-E_{\mathrm{2cl}}(n)}$, is a factor of $F^{\natural}(h)$.

The statement is also used in a corollary of the age lemma; there the use is at the second form, the
final integral.

The word ``performed'' in the definition of $k'$ matters. Suppose that at some step $\mathbf{R}$ is
begun on a member and a preparatory failure occurs. Then no integral \cite[(1.71)]{B-CMP122b} is
formed at that step, and no variable of $Y$ is integrated. Ba{\l}aban's construction simply continues
with the large-field region obtained, by the ten-layer operation, from the one present before the
step, together with its whole history \cite{B-CMP122b}. Thus $k'$ is the first step at which the
integral of $\mathbf{R}$ is actually performed.

(ii) \emph{Independence of the charge constant.} Neither the exponent $E_{\mathrm{2cl}}(n)$ nor the
number $N_Y$ involves the charge constant $C_0$ of the majorant. Enlarging $C_0$ to the enlarged
charge constant $C$ therefore leaves both the lower bound on $E_{\mathrm{2cl}}(n)$ and its consequence
$E_{\mathrm{2cl}}(n)\ge 2\mathsf{P}_n$ unchanged.

(iii) \emph{The Wilson coefficient on the dial.} The statement requires $c_W\ge\tfrac34$ on the dial
polydisc $\mathbb{D}_\delta$. This holds because the smallness condition imposed on the dial keeps the
total relative shift of every raw Wilson coefficient at most one quarter. That total shift is
$\theta_W$ plus a non-negative remainder contribution.

In the complex coordinates of the correlation theorem, a lemma on the Wilson coefficient proved with
that theorem gives that the real part of the raw Wilson coefficient over the level-$n$ stratum of $Y$
is at least $(1-\theta_W)\,g_n^{-2}$.
Since $\theta_W$ is at most the total relative shift, hence at most $\tfrac14$, we may take
$c_W\ge1-\theta_W\ge\tfrac34$. This is the hypothesis on the Wilson coefficient required by the
statement that a second-class outcome pays a tag, which holds on $\mathbb{D}_\delta$.

(iv) \emph{The ceiling on the coupling.} The bound $g_j\le g_\bullet$ for all $j\le k$, in particular
for $j=n$, follows from a ceiling on $\gamma$, by the comparison of the running couplings with the
reference coupling. It is the
line labelled $(g_\bullet)$ in the list of ceilings on $\gamma$ whose minimum is $\gamma^{\sharp}$.

(v) \emph{Counting once.} A second-class survivor enters the next step as an old constituent of a
member, never as a creation component. The counting rule of the statement that a second-class outcome
pays a tag is therefore respected, and in $F^{\natural}(h)$ the income appears exactly once, in the
group $(\gamma)$ of factors attached to the pair $(n,Y)$.

(vi) \emph{No other factor of $F^{\natural}(h)$ draws on the price.} We treat three kinds of factor.

First, the gains listed in \cite[pp.~380--383]{B-CMP122b} (among them the factors of the cubes
$P_j,Q_j,R_j$ and of the preparatory functions) do not draw on the price. Nor do the other
applications of the statement that a second-class outcome pays a tag. Both facts are a clause of that
statement and are established in its proof. In particular, second-class survivals at other levels of
the same lineage draw on distinct strata.

Second, the only factors of $F^{\natural}(h)$ that lie outside Ba{\l}aban's list are the refusing-bond
factors of the tightened family (Definition~\ref{9.2:def-construction-outcomes}). Of each such factor,
one half, $\upsilon_j=e^{-2\mathsf{Q}_j}$ for a bond of a region born at level $j$, is spent in the
creation deposit of Lemma~\ref{9.2:lem-creation-deposit}, which is proved as an implication from the
statements named there. The other half is spent in the counting weight over histories that accompanies
Proposition~\ref{9.2:prop-modulus-majorant}, whose proof is outstanding.

These factors are values of a fluctuation quadratic form. As in Lemma~\ref{9.2:lem-creation-deposit},
the refusing bonds and the bonds of the $R$-cubes are disjoint coordinates of one quadratic form,
bounded below by $\gamma_0$. The form is read on the outermost layer of the seed of the region. This is
the same source as Ba{\l}aban's factor for the $R$-cubes, which \cite[pp.~380--383]{B-CMP122b} obtains
by the same positivity bound.

Third, that the refusing-bond factors do not draw on the price is a matter of where they live. The
argument is the one used, in the proof of the statement that a second-class outcome pays a tag, for
Ba{\l}aban's $Q$- and $R$-cube factors. Those factors are gains from fluctuation quadratic forms:
\cite[pp.~380--383]{B-CMP122b} obtains them ``by expanding the Wilson action locally in the approximate
fluctuation field, and using the positivity bound for the resulting quadratic form''. Such a form lives
on a domain that, at its creation, is contained in $(Z^{\sim5})^c$ \cite{B-CMP119}. If the form was
created before step $n$, the step at which $\mathbf{R}$ was begun on $Y$, and lies inside $Y$, it
contributes no factor here, its variables having been integrated inside $\mathbf{T}'_n(Y)$.

Likewise, the refusing-bond factors are placed by the same alternative. Such a factor is read on the
outermost layer, the collar, of the seed of the region at whose creation it arises. If that region is
created at step $n$ or later, then $Y\subset Z$ at its creation, and the collar lies in
$(Z^{\sim5})^c$, hence outside $Y\supset Y^{\sim-4}$. If it is created before step $n$, the region
either lies inside $Y$, and then the factor contributes nothing here, its variables having been
integrated inside $\mathbf{T}'_n(Y)$, or lies in a component of $Z_n$ other than $Y$, and is then
disjoint from $Y$. The price, by contrast, is the raw Wilson term restricted to the
set $\Delta'\subset Y^{\sim-4}$ of the level-$n$ stratum of $Y$. In every case the refusing-bond
factor and the price are disjoint.

Finally, the prefactor $e$ in clause (c2) of Proposition~\ref{9.2:prop-modulus-majorant}, whose proof
is outstanding, is a constant, not an energy. It has no price, and it is paid by income, in the lemma
of this chapter in which the prefactors of clause (c2) are absorbed.
\end{proof}

\begin{lemma}[The comparison product of a history]\label{9.2:lem-comparison-product}
Let $k$ be given, let $X\in\operatorname{comp}(Z_k)$ be a component on which the integral of
$\mathbf{R}$ is performed at step $k$, and let $\operatorname{Hist}(X)$, $j(h)$, $\mathsf{Z}_n$,
$\operatorname{anc}(h)$, $\operatorname{births}(h)$ and $\operatorname{recr}(h)$ be as in
Definition~\ref{9.2:def-histories-ages}. For $h\in\operatorname{Hist}(X)$ and $C>0$ put
\begin{gather*}
D(h):=\prod_{(n,W)\in\operatorname{births}(h)}
\exp\bigl(-\mathsf{D}_n\,(d'(W)+1)-2\mathsf{P}_n\bigr),\qquad
I(h):=\prod_{(n,Z)\in\operatorname{recr}(h)} e^{-\mathsf{P}_n},\\
E^{[C]}(h):=\exp\Bigl(\sum_{(n,Z)\in\operatorname{anc}(h)} C\sqrt{d}\,M^dR_n^{d+1}|Z|\Bigr).
\end{gather*}
Neither $D(h)$ nor $I(h)$ contains $C$.

Let $d=4$ and $M\ge 4M_2$, and assume that for every $n\le k$:
\begin{itemize}
\item condition $(S)_n$ of \eqref{9.2:eq-creation-deposit-a} holds;
\item condition $(P)_n$ holds, that is, $\mathsf{P}_n\ge 1$, the line (P) of
\eqref{9.2:eq-forgetting-conditions-a};
\item $g_n\le g_\bullet$, where $g_\bullet$ is the ceiling on the coupling in the second-class
income bound listed below;
\item $c_W\ge \tfrac34$ on $\mathfrak{P}$, where $c_W$ is a lower bound, in units of $g_n^{-2}$,
of the real part of the Wilson coefficient on the level-$n$ stratum.
\end{itemize}
Assume further the following statements as hypotheses, for every $h\in\operatorname{Hist}(X)$:
\begin{itemize}
\item clauses (b) and (c) of the history-wise modulus majorant,
Proposition~\ref{9.2:prop-modulus-majorant}, whose proof is outstanding, with its charge constant
$C_0$ and its majorant $F^{\natural}(h)$; clause (b) sorts the factors of $F^{\natural}(h)$ into
the groups $(\alpha)$ of the pairs $(n,W)\in\operatorname{births}(h)$, $(\beta)$ of the pairs
$(n,Z)\in\operatorname{anc}(h)$, $(\gamma)$ of the pairs $(n,Z)\in\operatorname{recr}(h)$, and
$(\delta)$;
\item the creation deposit in linear-size form, Lemma~\ref{9.2:lem-creation-deposit};
\item the charge bound at the constant $C_0$: for every $(n,Z)\in\operatorname{anc}(h)$, the
group $(\beta)$ of $(n,Z)$ times the counting factor $e^{c_A^{\flat}(n)|Z|}$ is at most
$\exp(C_0\sqrt{d}\,M^dR_n^{d+1}|Z|)$;
\item the income of a survival after a preparatory failure: for every
$(n,Z)\in\operatorname{recr}_1(h)$, a re-creation of the first kind of
Definition~\ref{9.2:def-histories-ages}, the group $(\gamma)$ of $(n,Z)$ is at most
$e^{-\mathsf{P}_n}$;
\item the second-class income bound in dimension four: under $d=4$, $M\ge 4M_2$,
$g_n\le g_\bullet$ and $c_W\ge\tfrac34$, the income exponent $E_{\mathrm{2cl}}(n)$ of a
second-class outcome at level $n$ satisfies
\begin{equation*}
E_{\mathrm{2cl}}(n)\ge 2\mathfrak{p}_n\ge 2\mathsf{P}_n .
\end{equation*}
\end{itemize}
Here $I_n(h)$ is the refusing-bond label of $h$ at level $n$, $\upsilon_n=e^{-2\mathsf{Q}_n}$, and
$c_A^{\flat}(n)$ is the counting constant of \eqref{9.2:eq-history-summation-a}.
Then for every $h\in\operatorname{Hist}(X)$,
\begin{equation}\label{9.2:eq-comparison-product-a}
F^{\natural}(h)\cdot\prod_{n=j(h)}^{k}\upsilon_n^{-|I_n(h)|}\,e^{c_A^{\flat}(n)|\mathsf{Z}_n|}
\;\le\; D(h)\cdot I(h)\cdot E^{[C_0]}(h).
\end{equation}
\end{lemma}

\begin{proof}
Fix $h\in\operatorname{Hist}(X)$.

\emph{Distribution over members.} We distribute the left member of
\eqref{9.2:eq-comparison-product-a} over the live levels $n$, $j(h)\le n\le k$, and over the members
$Z\in\operatorname{comp}(\mathsf{Z}_n)$. The members of one level are disjoint components of
$\mathsf{Z}_n$. Hence $|\mathsf{Z}_n|=\sum_{Z\in\operatorname{comp}(\mathsf{Z}_n)}|Z|$, and
therefore
\begin{equation*}
e^{c_A^{\flat}(n)|\mathsf{Z}_n|}=\prod_{Z\in\operatorname{comp}(\mathsf{Z}_n)}e^{c_A^{\flat}(n)|Z|}.
\end{equation*}
Every refusing bond of $I_n(h)$ lies in $\mathsf{Z}_n$. As in
Lemma~\ref{9.2:lem-creation-deposit}, it is located in exactly one creation component $W$ of level
$n$, with $(n,W)\in\operatorname{births}(h)$. This $W$ is a new constituent of a member of
$\operatorname{comp}(\mathsf{Z}_n)$, because the cells of refusing cubes are creation cells, that
is, $I'\subset\mathcal{W}_n$. For $(n,W)\in\operatorname{births}(h)$ write $I_n(h;W)$ for the set
of refusing bonds of $I_n(h)$ located in $W$. The sets $I_n(h;W)$ partition $I_n(h)$, and so
\begin{equation}\label{9.2:eq-comparison-product-2}
\prod_{n=j(h)}^{k}\upsilon_n^{-|I_n(h)|}
=\prod_{(n,W)\in\operatorname{births}(h)}\upsilon_n^{-|I_n(h;W)|}.
\end{equation}

We now sort the groups into which clause (b) of Proposition~\ref{9.2:prop-modulus-majorant},
whose proof is outstanding, divides $F^{\natural}(h)$. To each group we attach the factors of the
left member of \eqref{9.2:eq-comparison-product-a} that belong to the same pair.

\emph{Group $(\alpha)$ of $(n,W)\in\operatorname{births}(h)$, times
$\upsilon_n^{-|I_n(h;W)|}$.} For each refusing bond of $I_n(h;W)$ the group contains a full factor
$\upsilon_n^2$. One factor $\upsilon_n$ of it cancels against the corresponding factor
$\upsilon_n^{-1}$ of \eqref{9.2:eq-comparison-product-2}. This is the factor that
Lemma~\ref{9.2:lem-creation-deposit} sets aside. The other factor $\upsilon_n$ is kept among the
remaining factors of the group. Lemma~\ref{9.2:lem-creation-deposit}, applied under condition
$(S)_n$ of \eqref{9.2:eq-creation-deposit-a}, bounds the product of these remaining factors by
\begin{equation*}
\exp\bigl(-\mathsf{D}_n\,(d'(W)+1)-2\mathsf{P}_n\bigr).
\end{equation*}
The product of these bounds over $\operatorname{births}(h)$ is $D(h)$.

\emph{Group $(\beta)$ of $(n,Z)\in\operatorname{anc}(h)$, times $e^{c_A^{\flat}(n)|Z|}$.} By the
charge bound at the constant $C_0$, this product is at most
$\exp(C_0\sqrt{d}\,M^dR_n^{d+1}|Z|)$. The product of these bounds over $\operatorname{anc}(h)$ is
$E^{[C_0]}(h)$.

\emph{Group $(\gamma)$ of $(n,Z)\in\operatorname{recr}(h)$.} There are two kinds of re-creation,
fixed in Definition~\ref{9.2:def-histories-ages}.

If $(n,Z)\in\operatorname{recr}_1(h)$, a survival after a preparatory failure, then the group is at
most $e^{-\mathsf{P}_n}$ by the income of a survival after a preparatory failure.

If $(n,Z)\in\operatorname{recr}_2(h)$, a second-class outcome, then by the second and third items
of clause (c) of Proposition~\ref{9.2:prop-modulus-majorant}, whose proof is outstanding, the group
equals $e\cdot e^{-E_{\mathrm{2cl}}(n)}$. The hypotheses $d=4$, $M\ge 4M_2$, $g_n\le g_\bullet$ and
$c_W\ge\tfrac34$ of the second-class income bound hold by assumption. Hence
the second-class income bound gives $E_{\mathrm{2cl}}(n)\ge 2\mathsf{P}_n$. Since
$\mathsf{P}_n\ge 1$ by $(P)_n$, the group is at most
\begin{equation*}
e\cdot e^{-E_{\mathrm{2cl}}(n)}\le e^{1-2\mathsf{P}_n}\le e^{-\mathsf{P}_n}.
\end{equation*}
In both cases the bound is $e^{-\mathsf{P}_n}$. Its product over $\operatorname{recr}(h)$ is $I(h)$.

\emph{Group $(\delta)$.} It is at most $1$.

\emph{Conclusion.} By \eqref{9.2:eq-comparison-product-2} and the distribution over members, every
factor of the left member of \eqref{9.2:eq-comparison-product-a} has been attached to exactly one
of the groups $(\alpha)$--$(\delta)$ of $F^{\natural}(h)$. All factors involved are non-negative
numbers, by clause (b) of Proposition~\ref{9.2:prop-modulus-majorant}, whose proof is outstanding.
Multiplying the four bounds obtained above therefore gives
\eqref{9.2:eq-comparison-product-a}.
\end{proof}

\begin{remark}[The prefactor of the inner operation of a second-class survival]
A second-class survival $(n,Y)$ brings two things into $F^{\natural}(h)$. The first is the income
$e^{-E_{\mathrm{2cl}}(n)}\le e^{-2\mathsf{P}_n}$. The second is the prefactor $e$ of the inner
performed operation $\mathbf{T}'_n(Y)$. There is one such prefactor per survival, and their
product over a lineage with many such survivals is not bounded in the age.

In the proof of Lemma~\ref{9.2:lem-comparison-product} the income is split into two halves. One
half, $e^{-\mathsf{P}_n}$, is kept in the factor $I(h)$ of
\eqref{9.2:eq-comparison-product-a}; it is the income used later in the induction of this chapter.
The other half pays the prefactor, since $e\cdot e^{-\mathsf{P}_n}\le 1$ under $(P)_n$.

Suppose the inner prefactor supplied by clause (c) of
Proposition~\ref{9.2:prop-modulus-majorant}, whose proof is outstanding, were an absolute factor
$e^{\mathfrak{e}_0}$ with $\mathfrak{e}_0>1$ in place of $e$. Then only the line (P) of
\eqref{9.2:eq-forgetting-conditions-a} would change, into $\mathsf{P}_n\ge\mathfrak{e}_0$. No
other line of \eqref{9.2:eq-forgetting-conditions-a} changes, and no constant of the standing list
of constants of this chapter changes.
\end{remark}

\begin{lemma}[Live levels and the per-level charge floor]\label{9.2:lem-live-levels}
Let $k$ be a level, let $X\in\operatorname{comp}(Z_k)$ be a component of the large-field region
$Z_k$, and let $h\in\operatorname{Hist}(X)$ be an admissible history of $X$, in the sense of
Definition~\ref{9.2:def-histories-ages}. For a set $Z$ of cells of $\mathbb{L}_n$ write $|Z|$ for
the number of cells it contains, and write $\#$ for the cardinality of a finite set. Then the
following hold.
\begin{itemize}
\item[(a)] The live levels of $h$ are exactly the integers $n$ with $j(h)\le n\le k$; their number
is $a(h)+1$.
\item[(b)] For each live level $n$ of $h$, the members of $\operatorname{comp}(\mathsf{Z}_n)$ are
pairwise disjoint non-empty sets of cells,
\begin{equation}\label{9.2:eq-live-levels-1}
|\mathsf{Z}_n| \;=\; \sum_{Z\in\operatorname{comp}(\mathsf{Z}_n)} |Z|
\;\ge\; \#\operatorname{comp}(\mathsf{Z}_n) \;\ge\; 1,
\end{equation}
and $\#\operatorname{anc}(h)\ge a(h)+1$.
\item[(c)] For every $C>0$ and every $(n,Z)\in\operatorname{anc}(h)$ one has
$C\sqrt{d}\,M^dR_n^{d+1}|Z|\ge C$; hence
\begin{equation}\label{9.2:eq-live-levels-2}
\sum_{(n,Z)\in\operatorname{anc}(h)} \sqrt{d}\,M^dR_n^{d+1}|Z| \;\ge\; a(h)+1 .
\end{equation}
\end{itemize}
\end{lemma}

\begin{proof}
(a) The live levels of $h$ are described by \eqref{9.2:eq-histories-ages-a}: they are the integers
$n$ with $j(h)\le n\le k$. Their number is $k-j(h)+1$, and since the age of $h$ at level $k$ is
$a(h)=k-j(h)$ (Definition~\ref{9.2:def-histories-ages}), this number equals $a(h)+1$.

(b) Let $n$ be a live level of $h$. A member of $\operatorname{comp}(\mathsf{Z}_n)$ is a component
of $Z_n$, that is, a maximal non-empty connected set of cells of $\mathbb{L}_n$
(Definition~\ref{9.2:def-cells-coarsening}); this is the cell structure of the large-field regions
of \cite{B-CMP122b}, where the domain $Z$ at level $j$ is a union of $MR_j$-cubes. Each member is
therefore a non-empty set of cells. Two distinct members are distinct components of $Z_n$, and
distinct maximal connected sets are disjoint; hence the members are pairwise disjoint. Since
$\mathsf{Z}_n$ is the union of the members of $\operatorname{comp}(\mathsf{Z}_n)$ and this union is
disjoint, the number of cells of $\mathsf{Z}_n$ is the sum of the numbers of cells of the members,
which is the equality in \eqref{9.2:eq-live-levels-1}. Each summand satisfies $|Z|\ge 1$ because
each member is non-empty, so the sum is at least the number of summands,
$\#\operatorname{comp}(\mathsf{Z}_n)$. At a live level $n$ the set $\mathsf{Z}_n$ is non-empty, by
the definition of live levels in Definition~\ref{9.2:def-histories-ages}, so
$\#\operatorname{comp}(\mathsf{Z}_n)\ge 1$. Finally, $\#\operatorname{anc}(h)\ge a(h)+1$ is
\eqref{9.2:eq-histories-ages-b}.

(c) Let $(n,Z)\in\operatorname{anc}(h)$. The set $Z$ is non-empty, so $|Z|\ge1$; moreover $M\ge1$
and $R_n\ge1$ (Notation~\ref{9.2:not-parameters-couplings}), and $\sqrt{d}\ge1$. The product of
these four factors gives
\begin{equation*}
\sqrt{d}\,M^dR_n^{d+1}|Z|\;\ge\;1 ,
\end{equation*}
and multiplying by $C>0$ gives $C\sqrt{d}\,M^dR_n^{d+1}|Z|\ge C$. Summing the displayed inequality
over $(n,Z)\in\operatorname{anc}(h)$ shows that the left side of \eqref{9.2:eq-live-levels-2} is at
least $\#\operatorname{anc}(h)$, which is at least $a(h)+1$ by (b). This proves
\eqref{9.2:eq-live-levels-2}.
\end{proof}

The lower bound \eqref{9.2:eq-live-levels-2} is the per-level floor on the charges of the ancestors
of a history; its proof uses no property of the shape of the regions and no property of the linear
size $d'$.

\begin{remark}[Ba{\l}aban's linear size in block-cube units]\label{9.2:rem-linear-size-block-cubes}
Let $Z$ be a connected union of $MR_n$-cubes of the level-$n$ lattice, that is, of cells of
$\mathbb{L}_n$ in the sense of Definition~\ref{9.2:def-cells-coarsening}. The linear size
$d'_n(Z)$ that enters the bounds \cite[(1.79)--(1.88)]{B-CMP122b} is
\begin{equation}\label{9.2:eq-linear-size-block-cubes-a}
\begin{split}
d'_n(Z) = d'_P(Z) := {}&\text{the least length, in units of } MR_n, \text{ of a tree graph}\\
&\text{contained in } Z \text{ and meeting every closed } MR_n\text{-cube of } Z .
\end{split}
\end{equation}
Thus $d'_n$ is the constrained variant of the linear size $d'$ of
\eqref{9.2:eq-cells-coarsening-a}: the infimum defining $d'$ is taken with the connecting set
required, in addition, to lie in the union of the closed cells of $Z$.

The relative size $d^Z_k(Y)$ of \cite[(1.67)]{B-CMP122b} is a different functional: it is
the length of a shortest tree graph in $Y$ meeting the $M$-cubes of the components of
$Y\setminus Z$. It is not the functional $d'_n$ of \cite[(1.79)--(1.88)]{B-CMP122b}.
Likewise, \cite[(6.31)]{B-CMP109}, which \cite{B-CMP122b} invokes for its coarsening
inequality, is a scaling property and not a definition; it is not used in what follows.
\end{remark}

\begin{proof}
In \cite{B-CMP109} a localization domain $X$ at level $j$ is a connected union of cubes of the
partition into $M$-cubes. For such an $X$, \cite{B-CMP109} considers the class of tree graphs
contained in $X$ and intersecting all the cubes in $X$, and defines the linear size $d_j(X)$ of
$X$ as the length of a shortest graph in this class, divided by $M$. In other words, $d_j(X)$
is the least length of a tree graph that lies in $X$ and meets every $M$-cube of $X$, measured
in units of $M$.

In \cite{B-CMP122b}, the quantity $d'_j$ is stated to be the linear size defined in terms of
$MR_j$-cubes instead of $M$-cubes. Carrying out this replacement in the definition of the
preceding paragraph, with the index $j$ written as $n$, yields the following. The domain is a
connected union $Z$ of $MR_n$-cubes of the level-$n$ lattice. The admissible graphs are the
tree graphs contained in $Z$ that intersect every $MR_n$-cube of $Z$. The minimal length is
divided by $MR_n$, that is, it is measured in units of $MR_n$. This is exactly the right-hand
side of \eqref{9.2:eq-linear-size-block-cubes-a}, which proves the identity $d'_n(Z)=d'_P(Z)$.

It remains to compare with \eqref{9.2:eq-cells-coarsening-a}. There, $d'(A)$ for a finite
non-empty $A\subset\mathbb{L}_n$ is the infimum of the lengths, in level-$n$ cell units, of
finite connected unions of straight segments meeting the closed cell of every $x\in A$. The
level-$n$ cells are the $MR_n$-cubes, so the level-$n$ cell unit is $MR_n$. The definition of
$d'_P(Z)$ imposes the same meeting condition and the same unit, and adds one requirement: the
connecting set must lie in $Z$, the union of the closed cells of $Z$. Hence $d'_P$ is the
constrained variant of $d'$, as asserted.

The statements about $d^Z_k(Y)$ and \cite[(6.31)]{B-CMP109} in the remark are readings of
the two displays cited there. The first follows from the definition of $d^Z_k(Y)$ in
\cite[(1.67)]{B-CMP122b}, whose trees lie in $Y$ and need only meet the $M$-cubes of the
components of $Y\setminus Z$. The second follows from the use of \cite[(6.31)]{B-CMP109} in
\cite{B-CMP122b} as a scaling property supporting the coarsening inequality there.
\end{proof}

\begin{lemma}[Where the linear size vanishes]\label{9.2:lem-linear-size-vanishing}
Let $A\subset\mathbb{Z}^d$ be finite and non-empty, let $d'(A)$ be its linear size
\eqref{9.2:eq-cells-coarsening-a}, and let $d'_P(A)$ be the constrained variant of
Definition~\ref{9.2:def-cells-coarsening}.
\begin{itemize}
\item[(a)] If $|x-y|_\infty\le 1$ for all $x,y\in A$, then the closed cells of $A$ have a
common point, and $d'(A)=d'_P(A)=0$.
\item[(b)] Otherwise $d'_P(A)\ge d'(A)\ge 1$.
\end{itemize}
Consequently $d'$ and $d'_P$ take no value in $]0,1[$. They are $\ge 1$ as soon as $A$
contains two cells at sup-distance $\ge 2$, in particular as soon as $A\supseteq c^{\sim 1}$
for some cell $c$.
\end{lemma}

\begin{proof}
(a) Fix a coordinate $i\in\{1,\dots,d\}$. By hypothesis the integers $x_i$, $x\in A$,
differ pairwise by at most $1$. Put $v_i:=\max_{x\in A}x_i$. Then every $x_i$ lies in
$\{v_i-1,v_i\}$. Put $v:=(v_1,\dots,v_d)$. For every $x\in A$ and every $i$ we then have
$v_i\in[x_i,x_i+1]$. Since the closed cell of $x$ is the set of points whose $i$-th
coordinate lies in $[x_i,x_i+1]$ for every $i$, the point $v$ lies in every closed cell
of $A$, and hence in their union. The degenerate tree $\{v\}$ is therefore admissible in
both infima defining $d'(A)$ and $d'_P(A)$, and its length is $0$. Since both infima are
non-negative, both vanish.

(b) Suppose the hypothesis of (a) fails. Then there are $x,y\in A$ and an index $i$ with
$|y_i-x_i|\ge 2$; after exchanging $x$ and $y$ if necessary, $y_i\ge x_i+2$.

Let $T$ be a connected finite union of segments that meets the closed cell of $x$ and the
closed cell of $y$. A point $p$ of $T$ in the closed cell of $x$ satisfies
$p_i\le x_i+1$. A point $p'$ of $T$ in the closed cell of $y$ satisfies
$p'_i\ge y_i\ge x_i+2$. Since $T$ is connected, it contains a polygonal path from $p$ to
$p'$ that runs along each segment of $T$ at most once. The $i$-th coordinate is
$1$-Lipschitz along this path, so the length of the path is at least
\begin{equation*}
|p'_i-p_i|\ \ge\ (x_i+2)-(x_i+1)\ =\ 1 .
\end{equation*}
The length of $T$ is at least the length of this path, and hence at least $1$.

Every set admissible in the infimum defining $d'(A)$ meets all closed cells of $A$, in
particular those of $x$ and $y$. Hence $d'(A)\ge 1$. The inequality $d'_P(A)\ge d'(A)$ holds
because the constrained infimum is taken over a subclass of the sets admitted in the
unconstrained one.

For the consequences, note that exactly one of the hypotheses of (a) and (b) holds. In the
first case both quantities equal $0$; in the second both are $\ge 1$. Hence neither takes a
value in $]0,1[$.

If $A$ contains two cells at sup-distance $\ge 2$, the hypothesis of (a) fails, and (b)
gives the bound $\ge 1$.

Finally, let $c$ be the cell with lower corner $x$, and let $e_1$ be the first unit
vector. Then $c^{\sim 1}$ contains the cells with lower corners $x-e_1$ and $x+e_1$, which
are at sup-distance $2$. So the preceding case applies whenever $A\supseteq c^{\sim 1}$.
\end{proof}

Thus non-emptiness of $A$ alone gives no positive lower bound for $d'$. A single cube, or
the $2^d$ cubes at a vertex, have $d'=0$. This is why lower bounds on the linear size are
taken in cells.

\begin{lemma}[Linear size at least one for constructed regions; part~(2): stated; proof incomplete
at the step marked $(\ast)$]\label{9.2:lem-linear-size-lower}
Let $d'$ be the linear size of \eqref{9.2:eq-cells-coarsening-a} and $d'_P$ its constrained variant,
and let the construction family of Definition~\ref{9.2:def-construction-outcomes} be given.
\begin{enumerate}
\item[(1)] Assume conditions \textup{(G1)} and \textup{(G2)} of \eqref{9.2:eq-construction-outcomes-a}.
Then every component of $Z_{n+1}$ produced by a step $n\to n+1$ of the construction contains
$c^{\sim 1}$ for some cell $c\in\mathbb{L}_{n+1}$, and for every non-empty $Z_0\subset\mathbb{L}_n$
one has $S_n(Z_0)\supseteq c^{\sim 10}$ for some cell $c$. Hence every living component of every level
$n\ge 2$, every survivor and every member of every lineage at such a level has $d'\ge 1$ and
$d'_P\ge 1$.
\item[(2)] The vanishing of the linear size at the place of \cite{B-CMP122b} reading
``or $d'_n(S^{n-j}(Z))\le (64)^d$ if it is equal to $0$'', and in the sentence ``Let $n_0$ be the last
index $n$ such that $d'_n(Z^{(n-j)})>0$'' of the same paper, concerns the covers $Z^{(n-j)}$ of a fixed
region by the larger cubes of later levels (``the cover of $Z$ by $MR_n$-cubes in the corresponding
scale'', \cite{B-CMP122b}), which shrink into a block of $2^d$ cells, inside the estimate
\cite[(1.81)]{B-CMP122b} of the clock. The regions $S^{n-j}(Z)$ themselves contain ten layers and
satisfy part~\textup{(1)}.
\end{enumerate}
\end{lemma}

Part~(1) is proved below as an implication from conditions \textup{(G1)} and \textup{(G2)} of
\eqref{9.2:eq-construction-outcomes-a}. Part~(2) is stated; proof incomplete at the step marked
$(\ast)$.

For a cell $c$ we write $c\pm e_1$ for the two cells obtained from $c$ by a translation by one cell in
the first coordinate direction, in either sense.

\begin{proof}[Proof of \textup{(1)}]
Let $\mathfrak{s}$ be a seed set of the step, and let $y\in\mathfrak{s}$ be a cell. By \textup{(G1)},
\begin{equation*}
\mathfrak{s}=(\pi_nZ_n^{\mathrm{liv}})^{\sim 10}\cup P'^{\sim 6}\cup Q'^{\sim 4}\cup R'^{\sim 1}
\cup I'^{\sim 1}.
\end{equation*}
The $s$-neighbourhood of a set of cells is the union of the $s$-neighbourhoods of its cells. Hence if
$y\in(\pi_nZ_n^{\mathrm{liv}})^{\sim 10}$, there is a cell $c\in\pi_nZ_n^{\mathrm{liv}}$ with
\begin{equation*}
y\in c^{\sim 10}\subset(\pi_nZ_n^{\mathrm{liv}})^{\sim 10}\subset\mathfrak{s};
\end{equation*}
in the same way, if $y$ lies in $P'^{\sim 6}$, $Q'^{\sim 4}$, $R'^{\sim 1}$ or $I'^{\sim 1}$, there is a
cell $c$ of $P'$, $Q'$, $R'$ or $I'$ respectively with $y\in c^{\sim s}\subset\mathfrak{s}$, where
$s=6,4,1,1$. Thus $y\in c^{\sim s}\subset\mathfrak{s}$ for some cell $c$ and some
$s\in\{10,6,4,1\}$. The set $c^{\sim s}$ is connected and contains $y$, so the component of
$\mathfrak{s}$ containing $y$ contains $c^{\sim s}\supseteq c^{\sim 1}$. Every component of a seed set
therefore contains a set $c^{\sim 1}$.

Say that a set of cells has property $(\mathrm{P})$ if each of its components contains $c^{\sim 1}$ for
some cell $c$. A union $Y=\bigcup_i Y_i$ of sets with property $(\mathrm{P})$ has property
$(\mathrm{P})$: a component $W$ of $Y$ is non-empty, so it meets some $Y_i$ and hence some component
$W_i$ of $Y_i$; since $W_i$ is connected, contained in $Y$ and meets $W$, we have $W_i\subset W$, and
$W_i$ contains some $c^{\sim 1}$. Let $Y$ have property $(\mathrm{P})$. A box step replaces a non-empty
set $W\subset Y$ by its box $\hat W\supseteq W$, so that the set after the step is
$(Y\setminus W)\cup\hat W=Y\cup\hat W$. The box $\hat W$ is connected and meets $Y$, since it contains
$W$; hence a component of $Y\cup\hat W$ that meets $\hat W$ contains $\hat W$ and so meets $Y$. Every
component of $Y\cup\hat W$ therefore meets $Y$, and, as in the union case, contains a component of $Y$
and so some $c^{\sim 1}$; the result has property $(\mathrm{P})$. By \textup{(G2)}, $Z_{n+1}$ is
obtained from the seed sets of the step's stages by finitely many unions and box steps; by induction
over these operations,
every component of $Z_{n+1}$ contains $c^{\sim 1}$ for some cell $c\in\mathbb{L}_{n+1}$.

If $Z_0\subset\mathbb{L}_n$ is non-empty, then $\pi_nZ_0$ is non-empty, and for any
$c\in\pi_nZ_0$ one has $S_n(Z_0)=(\pi_nZ_0)^{\sim 10}\supseteq c^{\sim 10}$.

The operation $\mathbf{R}$ removes whole components or leaves them as they are. Consequently, by
Definition~\ref{9.2:def-construction-outcomes}, each of the sets named in the last sentence of
\textup{(1)}, at a level $n\ge 2$, is a component of a set produced by a step $n-1\to n$, or a set of
the form $S_{n-1}(Z_0)$ with $Z_0$ non-empty, and by the two assertions just proved it contains
$c^{\sim 1}$ for some cell $c$. The set $c^{\sim 1}$ contains
the two cells $c-e_1$ and $c+e_1$, which are at sup-distance $2$. By
Lemma~\ref{9.2:lem-linear-size-vanishing}\textup{(b)}, any set containing two cells at sup-distance at
least $2$ satisfies $d'_P\ge d'\ge 1$. This proves \textup{(1)}.
\end{proof}

\begin{proof}[Proof of \textup{(2)}]
The regions $S^{n-j}(Z)$ are formed by adding ten layers at each application of $S$
\cite{B-CMP122b}; by part~\textup{(1)} each of them contains $c^{\sim 10}\supseteq c^{\sim 1}$ for some
cell $c$, and so has $d'\ge 1$ and $d'_P\ge 1$.

$(\ast)$ That the linear size whose vanishing is considered at the two places of \cite{B-CMP122b}
quoted in \textup{(2)} is that of the covers $Z^{(n-j)}$ of the fixed region $Z$ by $MR_n$-cubes of
the later level $n$, which shrink into a block of $2^d$ cells inside the estimate
\cite[(1.81)]{B-CMP122b} of the clock, and not that of the regions $S^{n-j}(Z)$, is a reading of that
passage and is not proved here.
\end{proof}

\begin{remark}[Use of the lower bounds on the linear size]
In the bound that charges the living large-field components level by level, the count is taken in
numbers of cells: the quantity charged per level $n$ and living component $Z$ is
$\sqrt{d}\,M^dR_n^{d+1}|Z|$, and $|Z|\ge 1$ for $Z\neq\emptyset$ requires no geometry. That bound
therefore depends neither on Lemma~\ref{9.2:lem-linear-size-vanishing} nor on
Lemma~\ref{9.2:lem-linear-size-lower}, and part~\textup{(2)} is used nowhere in this book. The two
lemmas show the following. Under Ba{\l}aban's definition \eqref{9.2:eq-linear-size-block-cubes-a},
every region containing two cubes at sup-distance at least $2$ has $d'_j\ge 1$; in particular a region
of extent larger than $20MR_j$, the extent of a region being its diameter in the sup-norm, has
$d'_j\ge 1$ whenever it contains two such cubes. At a creation level the relevant fact is
part~\textup{(1)}, and not the statement of \cite{B-CMP122a} on the extent gained by one step.
Part~\textup{(2)}, whose proof is incomplete at the step marked $(\ast)$, reads the vanishing linear
size of \cite{B-CMP122b} discussed there as belonging to the estimate of the clock and not to a charged
region. The lemma makes no assertion at level $1$; the numbers of layers added in the first step are
those of \cite[\S1, (1.10), (1.20)]{B-CMP119} and do not enter it.

In \cite{B-CMP122b} the chain
\begin{equation*}
\begin{aligned}
O(1)\log g_j^{-2}\,|Z_j|
&\le O(1)\log g_j^{-2}\,(MR_j)^d\,d'_j(Z_j)\\
&\le O(1)\,M^dR_j^{d+1}\,d'_j(Z_j)
\end{aligned}
\end{equation*}
appears as the contribution of every large-field domain $Z_j$. It is used here in neither direction.
For a non-empty region inside a block of $2^d$ cells one has $d'=0$ by
Lemma~\ref{9.2:lem-linear-size-vanishing}\textup{(a)}, so a bound of the volume by the linear size
carries an additive term; the general form of such a bound is the volume--linear-size inequality
$|Z|\le 2^{d+1}(d'(Z)+1)$.
\end{remark}

\begin{proposition}[The pre-payment induction at a given charge constant]
\label{9.2:prop-prepayment-induction}
Let $C>0$ and put
\begin{equation*}
c_n:=C\,M^dR_n^{d+1},\qquad \mathsf{c}_n(Z):=c_n\sqrt{d}\,|Z| ,
\end{equation*}
so that $\mathsf{c}_n(Z)$ is the charge of $Z$ at level $n$. Let $\kappa^{[C]}$ be the pre-payment
numbers of Definition~\ref{9.2:def-prepayment-numbers}, defined by the rules $(\kappa1)$--$(\kappa3)$
of \eqref{9.2:eq-prepayment-numbers-a}. Let $\Gamma$, $\Gamma'_n$, $\sigma$, $c_0$, $\beta_0$,
$\ell_n$, $R_n$, $\mathsf{D}_n$ and $\mathsf{P}_n$ be as in
Notation~\ref{9.2:not-parameters-couplings}, and let the histories, their components, cores and
clocks be those of Definitions~\ref{9.2:def-construction-outcomes}
and~\ref{9.2:def-histories-ages}. Assume the following.
\begin{enumerate}
\item[(a)] $L\ge 5$, and the flow facts \textup{(C0a)}--\textup{(C0c)} of
Notation~\ref{9.2:not-parameters-couplings} hold.
\item[(b)] The conclusions (i)--(iii) of the lemma on cores and connectors of large-field
components hold, namely: for every history of the construction family of
Definition~\ref{9.2:def-construction-outcomes} (whose conditions \textup{(G1)}--\textup{(G4)} are
those of \eqref{9.2:eq-construction-outcomes-a}), every level $n$ and every component $Z$ of
$Z_n$, with $\mathcal{C}:=\mathcal{C}_n(Z)$,
\begin{enumerate}
\item[(b-i)] $\emptyset\neq\mathcal{C}\subset Z\subset\mathcal{C}^{\sim 109}$;
\item[(b-ii)] $\sqrt{d}\,|Z|\le\Gamma\,(d'(\mathcal{C})+1)$;
\item[(b-iii)] if $Z$ has $N$ constituents with cores $\mathcal{C}_1,\dots,\mathcal{C}_N$, where
$\mathcal{C}_a=\pi_{n-1}\mathcal{C}_{n-1}(Z_0)$ for an old constituent $Z_0$ and
$\mathcal{C}_a=W$ for a new constituent with creation cell $W$, then
\begin{equation*}
d'(\mathcal{C})+1\le\sum_{a=1}^{N}\bigl(d'(\mathcal{C}_a)+1\bigr)+(N-1)\sigma .
\end{equation*}
\end{enumerate}
\item[(c)] The conclusions (i)--(iii) of the clock lemma for large-field components hold, namely:
for every history of the same family, every level $n$ and every
living component $Z$ of level $n$,
\begin{enumerate}
\item[(c-i)] if the step $n\to n+1$ is quiet for the descendant $Z'$ of $Z$ and $K_n(Z)\ge1$,
then $K_{n+1}(Z')\le K_n(Z)-1$;
\item[(c-ii)] if $Z$ fits in $D$ cells per axis, then $K_n(Z)\le 2\log_L(D+1)+4+R_n$; and
\begin{equation*}
K_n(Z)\le\tfrac14\bigl(d'(\mathcal{C}_n(Z))+1\bigr)+c_0+R_n ;
\end{equation*}
\item[(c-iii)] if $Z$ is living after step $n$ and $K_n(Z)=0$, then $Z$ was re-created at $n$.
\end{enumerate}
\item[(d)] For all levels $n\le k$,
\begin{equation}\label{9.2:eq-prepayment-induction-a}
\begin{gathered}
\mathsf{D}_n\ge\Gamma c_n\,(9+4c_0+4R_n),\\
(1+\beta_0)^{-1}\mathsf{P}_n\ge\Gamma c_n\,
\bigl(5\sigma+4\sqrt{d}\cdot101^d+4c_0+4R_n+48\bigr).
\end{gathered}
\end{equation}
\end{enumerate}
Then for every level $k$, every $X\in\operatorname{comp}(Z_k)$, every
$h\in\operatorname{Hist}(X)$, every live level $n$ and every
$Z\in\operatorname{comp}(\mathsf{Z}_n)$,
\begin{equation}\label{9.2:eq-prepayment-induction-b}
\kappa^{[C]}_n(Z)\ge\Gamma'_n\,c_n\bigl(d'(\mathcal{C}_n(Z))+1\bigr)+4\Gamma c_n K_n(Z);
\end{equation}
in particular $\kappa^{[C]}_n(Z)\ge\Gamma c_n>0$.

The inequalities \eqref{9.2:eq-prepayment-induction-a} are used only at levels $n\le k$: at a
step $n\to n+1$ with $n+1\le k$ they are used at level $n+1$ or at level $n$, as the proof shows.
\end{proposition}

\begin{proof}
We write $\kappa_n$ for $\kappa^{[C]}_n$. In this proof the letters $N$, $m$, $q'$, $K$, $D$ and
$Q$ denote the local counts and objects introduced where they first occur, and not the constants
of the glossary. The proof is by induction on the level $n$, from
$j(h)$ upwards; at each level the members of $\operatorname{comp}(\mathsf{Z}_n)$ are treated one
by one according to the event that produces them in the construction family of
Definition~\ref{9.2:def-construction-outcomes}.

\emph{Notation for a step.} Consider a step $n\to n+1$ with $n+1\le k$ and a member
$Z'\in\operatorname{comp}(\mathsf{Z}_{n+1})$. We write $c:=c_n$ and $c':=c_{n+1}$; by \textup{(C0b)},
$c'\le c$. Since $c_n=CM^dR_n^{d+1}$ with $C>0$, the inequality $c'\le c$ is equivalent to
$R_{n+1}\le R_n$, and $c'>0$. For an old constituent $Z^{(v)}\in\operatorname{comp}(\mathsf{Z}_n)$
of $Z'$ we put $x_v:=d'(\mathcal{C}_n(Z^{(v)}))$. We put $\mathcal{C}':=\mathcal{C}_{n+1}(Z')$
and $y:=d'(\mathcal{C}')$. By (b-ii), applied at level $n+1$ to $Z'$ and multiplied by
$c'>0$,
\begin{equation}\label{9.2:eq-prepayment-induction-1}
\mathsf{c}_{n+1}(Z')=c'\sqrt{d}\,|Z'|\le\Gamma c'(y+1).
\end{equation}

\emph{The base.} Members with no old constituent --- all members of
$\operatorname{comp}(\mathsf{Z}_{j(h)})$, and any member born fresh at a later level --- fall under
the creation-or-join case below with $m=0$, and that case, when $m=0$, uses no induction
hypothesis. This is the base of the induction; for a member of
$\operatorname{comp}(\mathsf{Z}_{j(h)})$ the case below is read with $n+1$ replaced by $j(h)$.
There remain three kinds of step, which we treat in turn.

\emph{Quiet step.} Suppose $Z'$ has the single old constituent $Z_0$, no creation cell, and that
$Z_0$ was not re-created at $n$. Then $Z_0$ is living after step $n$, since it has the descendant
$Z'$. If $K_n(Z_0)$ were $0$, (c-iii) would show that $Z_0$ was re-created at $n$, which is
excluded; hence $K:=K_n(Z_0)\ge1$. By rule $(\kappa2)$ of \eqref{9.2:eq-prepayment-numbers-a},
\begin{equation*}
\kappa_{n+1}(Z')=\kappa_n(Z_0)-\mathsf{c}_{n+1}(Z').
\end{equation*}
By the induction hypothesis \eqref{9.2:eq-prepayment-induction-b} at level $n$, with
$x:=d'(\mathcal{C}_n(Z_0))$,
\begin{equation*}
\kappa_n(Z_0)\ge\Gamma'_n c\,(x+1)+4\Gamma cK .
\end{equation*}
By \eqref{9.2:eq-prepayment-induction-1}, $\mathsf{c}_{n+1}(Z')\le\Gamma c'(y+1)$. Since
$\mathcal{C}'=\pi_n\mathcal{C}_n(Z_0)$, property (E3) of
Lemma~\ref{9.2:lem-linear-size-properties} gives $y\le x/\ell_n$. By (c-i), applied to $Z_0$
(the step is quiet for its descendant $Z'$ and $K\ge1$), $K_{n+1}(Z')\le K-1$.

The right member of \eqref{9.2:eq-prepayment-induction-b} at level $n+1$ for $Z'$, namely
$\Gamma'_{n+1}c'(y+1)+4\Gamma c'K_{n+1}(Z')$, and the bound $\Gamma c'(y+1)$ for
$\mathsf{c}_{n+1}(Z')$ are increasing in $y$, and the former is increasing in $K_{n+1}(Z')$; so we
may replace $y$ by $x/\ell_n$ and $K_{n+1}(Z')$ by $K-1$. Hence
\eqref{9.2:eq-prepayment-induction-b} at level $n+1$ for $Z'$ follows from
\begin{equation*}
\Gamma'_n c\,(x+1)+4\Gamma cK-\Gamma c'\Bigl(\frac{x}{\ell_n}+1\Bigr)
\ge\Gamma'_{n+1}c'\Bigl(\frac{x}{\ell_n}+1\Bigr)+4\Gamma c'(K-1).
\end{equation*}
Since $c\ge c'$ and $K\ge0$,
\begin{equation*}
4\Gamma cK\ge4\Gamma c'K=4\Gamma c'(K-1)+4\Gamma c',
\end{equation*}
and $\Gamma'_n c(x+1)\ge\Gamma'_n c'(x+1)$. Subtracting $4\Gamma c'(K-1)$ from both sides and
dividing by $c'>0$, we see that it suffices that
\begin{equation}\label{9.2:eq-prepayment-induction-2}
\Gamma'_n(x+1)+4\Gamma\ge(\Gamma+\Gamma'_{n+1})\Bigl(\frac{x}{\ell_n}+1\Bigr).
\end{equation}
If $\ell_n=L\ge5$, then $\Gamma'_n=\Gamma$ and $\Gamma'_{n+1}\le3\Gamma$, so the right member of
\eqref{9.2:eq-prepayment-induction-2} is at most $4\Gamma(x/L+1)$; and
\begin{equation*}
\Gamma(x+1)+4\Gamma=\Gamma(x+5)\ge4\Gamma\Bigl(\frac{x}{5}+1\Bigr)
\ge4\Gamma\Bigl(\frac{x}{L}+1\Bigr),
\end{equation*}
where the first inequality holds because $x\ge\frac45x$ (as $x\ge0$) and $5\ge4$, and the second
because $L\ge5$. If $\ell_n=1$, then by \textup{(C0c)} we have $\Gamma'_n=3\Gamma$ and
$\ell_{n+1}=L$, so $\Gamma'_{n+1}=\Gamma$; the right member of
\eqref{9.2:eq-prepayment-induction-2} is $2\Gamma(x+1)$, and
$3\Gamma(x+1)+4\Gamma\ge2\Gamma(x+1)$. In both cases \eqref{9.2:eq-prepayment-induction-2} holds.
No condition on $C$ is used in this step.

\emph{Creation or join.} Suppose $Z'$ has $m$ old constituents $Z^{(v)}$ and $q'$ new
constituents with creation cells $W_1,\dots,W_{q'}$, with $(m,q')\neq(1,0)$, and put $N:=m+q'$. In
rule $(\kappa1)$ of \eqref{9.2:eq-prepayment-numbers-a} the term $-2\mathsf{P}_{j(Z')}$ cancels
one of the $N$ linear terms, and each of the remaining $N-1$ linear terms is at least
$2(1+\beta_0)^{-1}\mathsf{P}_{n+1}$: for a term $2\mathsf{P}_{j(v)}$ of an old constituent this is
\textup{(C0a)}, since every $j(v)\le n\le n+1$; for the term $2\mathsf{P}_{n+1}$ of a new
constituent it is trivial, since $\mathsf{P}_{n+1}>0$ and $\beta_0\ge0$. Thus
\begin{equation}\label{9.2:eq-prepayment-induction-3}
\begin{split}
\kappa_{n+1}(Z')\ge{}&\sum_v\kappa_n(Z^{(v)})+\sum_i\mathsf{D}_{n+1}\bigl(d'(W_i)+1\bigr)\\
&+(N-1)\cdot2(1+\beta_0)^{-1}\mathsf{P}_{n+1}-\mathsf{c}_{n+1}(Z').
\end{split}
\end{equation}
Put
\begin{equation*}
\Phi:=\Gamma'_{n+1}c'(y+1)+4\Gamma c'K_{n+1}(Z')+\mathsf{c}_{n+1}(Z').
\end{equation*}
Then \eqref{9.2:eq-prepayment-induction-b} at level $n+1$ for $Z'$ follows if the sum of the
first three terms of the right member of \eqref{9.2:eq-prepayment-induction-3} is at least
$\Phi$. By \eqref{9.2:eq-prepayment-induction-1}, $\mathsf{c}_{n+1}(Z')\le\Gamma c'(y+1)$. By the
second inequality of (c-ii) at level $n+1$, multiplied by $4\Gamma c'$,
\begin{equation*}
4\Gamma c'K_{n+1}(Z')\le\Gamma c'(y+1)+4\Gamma c'(c_0+R_{n+1}).
\end{equation*}
Hence
\begin{equation*}
\Phi\le\Gamma''c'(y+1)+4\Gamma c'(c_0+R_{n+1}),\qquad \Gamma'':=\Gamma'_{n+1}+2\Gamma ,
\end{equation*}
and $\Gamma''\le5\Gamma$ because $\Gamma'_{n+1}\le3\Gamma$, while $\Gamma''=3\Gamma$ if
$\ell_{n+1}=L$, because then $\Gamma'_{n+1}=\Gamma$. The cores of the old constituents, read at
level $n+1$, are $\pi_n\mathcal{C}_n(Z^{(v)})$, and by (E3) of
Lemma~\ref{9.2:lem-linear-size-properties} their linear sizes are at most $x_v/\ell_n$. So (b-iii)
at level $n+1$ gives
\begin{equation*}
y+1\le\sum_v\Bigl(\frac{x_v}{\ell_n}+1\Bigr)+\sum_i\bigl(d'(W_i)+1\bigr)+(N-1)\sigma .
\end{equation*}
Multiplying by $\Gamma''c'$ and using $\Gamma''\le5\Gamma$ in the last term, we obtain
\begin{equation*}
\begin{split}
\Phi\le{}&\sum_v\Gamma''c'\Bigl(\frac{x_v}{\ell_n}+1\Bigr)+\sum_i\Gamma''c'\bigl(d'(W_i)+1\bigr)\\
&+5(N-1)\sigma\Gamma c'+4\Gamma c'(c_0+R_{n+1}).
\end{split}
\end{equation*}
We bound the terms of the first two sums.

Old $v$ with $\ell_n=L$: since $\Gamma''\le5\Gamma$ and $L\ge5$,
\begin{equation*}
\begin{split}
\Gamma''c'\Bigl(\frac{x_v}{L}+1\Bigr)&\le5\Gamma c'\Bigl(\frac{x_v}{5}+1\Bigr)
=\Gamma c'(x_v+1)+4\Gamma c'\\
&\le\Gamma c(x_v+1)+4\Gamma c'
\le\kappa_n(Z^{(v)})+4\Gamma c',
\end{split}
\end{equation*}
the last inequality by the induction hypothesis \eqref{9.2:eq-prepayment-induction-b} at level
$n$ for $Z^{(v)}$, together with $\Gamma'_n\ge\Gamma$ and $K_n(Z^{(v)})\ge0$.

Old $v$ with $\ell_n=1$: then $\ell_{n+1}=L$ by \textup{(C0c)}, so $\Gamma''=3\Gamma=\Gamma'_n$, and
\begin{equation*}
\Gamma''c'(x_v+1)\le3\Gamma c(x_v+1)\le\kappa_n(Z^{(v)}),
\end{equation*}
again by \eqref{9.2:eq-prepayment-induction-b} at level $n$ and $K_n(Z^{(v)})\ge0$.

New $i$: the first inequality of \eqref{9.2:eq-prepayment-induction-a} at level $n+1$ gives
$\mathsf{D}_{n+1}\ge9\Gamma c'\ge5\Gamma c'$, because $c_0\ge0$ and $R_{n+1}\ge0$. Hence
$\mathsf{D}_{n+1}-5\Gamma c'\ge0$, and since $d'(W_i)+1\ge1$,
\begin{equation*}
\begin{split}
\Gamma''c'\bigl(d'(W_i)+1\bigr)&\le5\Gamma c'\bigl(d'(W_i)+1\bigr)\\
&=\mathsf{D}_{n+1}\bigl(d'(W_i)+1\bigr)-(\mathsf{D}_{n+1}-5\Gamma c')\bigl(d'(W_i)+1\bigr)\\
&\le\mathsf{D}_{n+1}\bigl(d'(W_i)+1\bigr)-(\mathsf{D}_{n+1}-5\Gamma c').
\end{split}
\end{equation*}

Collecting these bounds (each of the $m$ old constituents contributes at most $4\Gamma c'$ beyond
$\kappa_n(Z^{(v)})$, and each of the $q'$ new constituents contributes
$-(\mathsf{D}_{n+1}-5\Gamma c')$ beyond $\mathsf{D}_{n+1}(d'(W_i)+1)$), and comparing with the first
three terms of the right member of \eqref{9.2:eq-prepayment-induction-3}, we see that
\eqref{9.2:eq-prepayment-induction-b} at level $n+1$ for $Z'$ follows from
\begin{equation}\label{9.2:eq-prepayment-induction-c}
\begin{split}
&4m\Gamma c'+5(N-1)\sigma\Gamma c'+4\Gamma c'(c_0+R_{n+1})\\
&\qquad\le(N-1)\cdot2(1+\beta_0)^{-1}\mathsf{P}_{n+1}+q'(\mathsf{D}_{n+1}-5\Gamma c').
\end{split}
\end{equation}

If $N\ge2$: then $m\le N\le2(N-1)$ and $N-1\ge1$, so the left member of
\eqref{9.2:eq-prepayment-induction-c} is at most
\begin{equation*}
(N-1)\Gamma c'\bigl(8+5\sigma+4c_0+4R_{n+1}\bigr)\le(N-1)(1+\beta_0)^{-1}\mathsf{P}_{n+1},
\end{equation*}
the inequality by the second inequality of \eqref{9.2:eq-prepayment-induction-a} at level
$n+1$, since $c'=c_{n+1}$ and $48\ge8$ (and $4\sqrt{d}\cdot101^d\ge0$). As
$\mathsf{P}_{n+1}\ge0$, the right member of this display is at most
$(N-1)\cdot2(1+\beta_0)^{-1}\mathsf{P}_{n+1}$; and
$q'(\mathsf{D}_{n+1}-5\Gamma c')\ge0$ by the first inequality of
\eqref{9.2:eq-prepayment-induction-a} at level $n+1$, as shown above. So
\eqref{9.2:eq-prepayment-induction-c} holds.

If $N=1$: since $(m,q')\neq(1,0)$, we have $(m,q')=(0,1)$, a birth standing alone. Then
\eqref{9.2:eq-prepayment-induction-c} reads
$4\Gamma c'(c_0+R_{n+1})\le\mathsf{D}_{n+1}-5\Gamma c'$, that is,
\begin{equation*}
\mathsf{D}_{n+1}\ge\Gamma c'\bigl(5+4c_0+4R_{n+1}\bigr),
\end{equation*}
which is implied by the first inequality of \eqref{9.2:eq-prepayment-induction-a} at level $n+1$,
since $9\ge5$.

\emph{Re-creation step.} Suppose $Z'$ has the single old constituent $Z_0$, no creation cell, and
that $Z_0$ was re-created at $n$. Then the operation $\mathbf{R}$ was begun on $Z_0$ at $n$, so
$Z_0$ is of class (i) in the construction family of Definition~\ref{9.2:def-construction-outcomes}
and fits in a cube of $100$ cells per axis; and, by \textup{(G3)} in
\eqref{9.2:eq-construction-outcomes-a}, $Z'$ is either $S_nZ_0$ or the box of $S_nZ_0$. The
induction hypothesis \eqref{9.2:eq-prepayment-induction-b} at level $n$ gives
$\kappa_n(Z_0)\ge0$, and together with rule $(\kappa3)$ of \eqref{9.2:eq-prepayment-numbers-a}
this yields
\begin{equation*}
\kappa_{n+1}(Z')\ge\mathsf{P}_n-\mathsf{c}_{n+1}(Z').
\end{equation*}
Now $\mathcal{C}'=\pi_n\mathcal{C}_n(Z_0)\subset\pi_nZ_0$, because
$\mathcal{C}_n(Z_0)\subset Z_0$ by (b-i). The image under $\pi_n$ of a set fitting in $100$ cells
per axis fits in $\lceil100/\ell_n\rceil+1\le101$ cells per axis, that is, it lies in a cube $Q$
of $101$ cells per axis. The cube $Q$ is connected and has $101^d$ cells, so (E4) and (E2) of
Lemma~\ref{9.2:lem-linear-size-properties} give
\begin{equation*}
y\le d'(Q)\le\sqrt{d}\,(101^d-1)<\sqrt{d}\cdot101^d .
\end{equation*}
Further, $Z'$ fits in $101+20=121$ cells per axis: $S_n$ adds ten cells on each side, and a box
has the extents of its set. So the first inequality of (c-ii) at level $n+1$, with $D=121$,
gives
\begin{equation*}
K_{n+1}(Z')\le2\log_L122+4+R_{n+1}\le10+R_{n+1},
\end{equation*}
since $L^3\ge125>122$ gives $\log_L122<3$. By \eqref{9.2:eq-prepayment-induction-1},
$\mathsf{c}_{n+1}(Z')\le\Gamma c'(y+1)$. Therefore, with $\Phi$ defined for $Z'$ as in the
preceding case, using $\Gamma'_{n+1}\le3\Gamma$ and
$y+1\le\sqrt{d}\,(101^d-1)+1\le\sqrt{d}\cdot101^d$ (as $\sqrt{d}\ge1$),
\begin{equation*}
\begin{split}
\Phi&=\Gamma'_{n+1}c'(y+1)+4\Gamma c'K_{n+1}(Z')+\mathsf{c}_{n+1}(Z')\\
&\le4\Gamma c'(y+1)+4\Gamma c'(R_{n+1}+10)\\
&\le\Gamma c'\bigl(4\sqrt{d}\cdot101^d+4R_{n+1}+44\bigr).
\end{split}
\end{equation*}
Since $\kappa_{n+1}(Z')\ge\mathsf{P}_n-\mathsf{c}_{n+1}(Z')$,
\eqref{9.2:eq-prepayment-induction-b} at level $n+1$ for $Z'$ follows from $\mathsf{P}_n\ge\Phi$,
hence from
\begin{equation*}
\mathsf{P}_n\ge\Gamma c'\bigl(4\sqrt{d}\cdot101^d+4R_{n+1}+44\bigr).
\end{equation*}
This is implied by the second inequality of \eqref{9.2:eq-prepayment-induction-a} at level $n$,
because $c'\le c_n$, $R_{n+1}\le R_n$, $5\sigma+4c_0+48\ge44$, and
$\mathsf{P}_n\ge(1+\beta_0)^{-1}\mathsf{P}_n$ (as $\beta_0\ge0$ and $\mathsf{P}_n\ge0$).

The term $\mathsf{P}_n$ of rule $(\kappa3)$ is, for survivors of kind (s1) in
\eqref{9.2:eq-construction-outcomes-a}, the exponent of the factor $\exp(-p_0(g_n))$ of
Ba{\l}aban's induction for the numbers $\kappa$ \cite[(1.80)--(1.88), pp.~384--387]{B-CMP122b},
and, for survivors of kind (s2), the half-income that the second-class income statement provides
for such survivors; the inequality used above is the same in both cases.

The argument thus has five places: first, the lone birth, case $N=1$ of the creation-or-join
step; second, the quiet step; third, the re-creation step for survivors of kind (s1); fourth,
the join, case $N\ge2$ of the creation-or-join step; fifth, the re-creation step for survivors of
kind (s2).

\emph{Conclusion.} By induction, \eqref{9.2:eq-prepayment-induction-b} holds at every live level
$n$ for every $Z\in\operatorname{comp}(\mathsf{Z}_n)$. Finally, since $d'(\mathcal{C}_n(Z))\ge0$,
$K_n(Z)\ge0$ and $\Gamma'_n\ge\Gamma$, \eqref{9.2:eq-prepayment-induction-b} gives
\begin{equation*}
\kappa_n(Z)\ge\Gamma'_nc_n\cdot1+0\ge\Gamma c_n>0 .
\end{equation*}
\end{proof}

\begin{remark}[The five comparisons of the charge constant]\label{9.2:rem-five-comparisons}
Let $C>0$ and $c_n=CM^dR_n^{d+1}$ be as in Proposition~\ref{9.2:prop-prepayment-induction}. The
charge constant $C$ enters the proof of Proposition~\ref{9.2:prop-prepayment-induction} only
through the two conditions $(B)_C$ and $(M)_C$ of \eqref{9.2:eq-prepayment-induction-a}. These
conditions are invoked at the following five places of the case list of Ba{\l}aban's induction for
the numbers $\kappa$ \cite[pp.~384--387]{B-CMP122b}, and nowhere else. The outcomes of the
operation named in the list are those of \eqref{9.2:eq-construction-outcomes-a}.
\begin{enumerate}
\item[(a)] At a birth standing alone, that is, a creation with a single constituent, $(B)_C$ is
invoked. It compares the quadratic deposit $\mathsf{D}_n$, of degree $2p_0$ in $\Xi_n$, with
$CM^dR_n^{d+2}$, of degree $(d+2)r$, where $R_n$ is of degree $r$ in $\Xi_n$
(Notation~\ref{9.1:not-radius-exponent-floors}). This is the base case of
\cite[(1.81)--(1.82)]{B-CMP122b}, whose condition is stated there in the words ``This condition is
satisfied for $p_0$ large, and $g_1$ sufficiently small''.
\item[(b)] At a quiet step no condition on $C$ is invoked. This is the identity
\cite[(1.83)]{B-CMP122b}, which carries the same constant on both sides.
\item[(c)] At a survival after a preparatory failure, $(M)_C$ is invoked at the level $n$ of the
survival. It compares the linear income $\mathsf{P}_n$, of degree $p_0$ in $\Xi_n$, with
$CM^dR_n^{d+2}$. This is the survivor rule of \cite[p.~386]{B-CMP122b}. The income used is the
income bound for a survival after a preparatory failure.
\item[(d)] At a creation or a join with at least two constituents, $(M)_C$ is invoked at level
$n+1$ for the old constituents and for the connectors. For the new constituents $(B)_C$ is invoked
at level $n+1$; their deposits satisfy $\mathsf{D}_{n+1}\ge 5\Gamma c_{n+1}$. This is
\cite[(1.85)--(1.88)]{B-CMP122b}, with the condition stated there as ``for $p_0$ large and $\gamma$
small enough''.
\item[(e)] At a survival as second class, the case list of \cite[pp.~384--387]{B-CMP122b} contains
no case for it. The components of \cite[p.~378]{B-CMP122b} re-enter the large-field region, and the
product in \cite[(1.79)]{B-CMP122b} assigns no income to them. The income is supplied by the
statement that a second-class outcome pays a tag, which delivers a factor
\begin{equation*}
e^{-E_{\mathrm{2cl}}(n)}\le e^{-2\mathsf{P}_n}
\end{equation*}
inside the history-wise modulus majorant $F^{\natural}$ of the first integral performed later (the
nesting clause of Proposition~\ref{9.2:prop-modulus-majorant}, whose proof is outstanding). Half of
it, $e^{-\mathsf{P}_n}$, is the income of the third clause of
Definition~\ref{9.2:def-prepayment-numbers}. The other half pays the
prefactor of the inner operation (Lemma~\ref{9.2:lem-comparison-product}). The inequality of~(c)
then applies verbatim. The factor supplied by that statement is drawn on by no other factor of the
majorant.
\end{enumerate}
Both $(B)_C$ and $(M)_C$ consist of a degree inequality on the exponents, namely $2p_0>(d+2)r$ and
$p_0>(d+2)r$ respectively, followed by a ceiling on $\gamma$ that depends on $C$, hence on $c_a$;
these are the lines so named in the one-sided list of conditions on the parameters.
The incomes $\mathsf{D}_n$, $\mathsf{P}_n$, $E_{\mathrm{2cl}}(n)$ do not contain $C$. Nor do the
geometric lemmas invoked in the proof of Proposition~\ref{9.2:prop-prepayment-induction}.
Without the statement that a second-class outcome pays a tag, geometric forgetting of old regions
would not be proved for histories containing a survival as second class, and it would be proved as
an implication from the remaining inputs for all other histories.
\end{remark}

\begin{proof}
Ba{\l}aban's induction for the numbers $\kappa$ \cite[pp.~384--387]{B-CMP122b} distinguishes the
following cases.
\begin{itemize}
\item The base case of a fresh region \cite[(1.81)--(1.82)]{B-CMP122b}. Its condition is
``satisfied for $p_0$ large, and $g_1$ sufficiently small''.
\item The first case: no large field is created in the last step, $Z=S(Z_0)$ with $S$ the ten-layer
operation, and the identity \cite[(1.83)]{B-CMP122b} holds.
\item Within the first case, the survivor after a preparatory failure, treated in
\cite[p.~386]{B-CMP122b} as follows:
\begin{quotation}
Then $\kappa_j(Z_0) \ge 0$, and we do not have a better bound for it, but we have the new factor
$\exp(-p_0(g_j))$. We define $\kappa_{j+1}(Z) = p_0(g_j) - O(1)M^dR^{d+1}_{j+1}d'_{j+1}(Z)$. It
satisfies (1.80), because $Z$ is a small domain, it is contained in a cube of the size
$100MR_{j+1}$, hence $K = R_{j+1}$ for $Z$
\end{quotation}
\item The second case, the join \cite[(1.84)--(1.88)]{B-CMP122b}, whose condition is ``for $p_0$
large and $\gamma$ small enough''.
\end{itemize}

\emph{Places (a)--(d).} In the proof of Proposition~\ref{9.2:prop-prepayment-induction} these cases
are treated as follows.
\begin{itemize}
\item The base case is place~(a). There the charge $c_n$ is compared with a deposit through
$(B)_C$.
\item The survivor rule is place~(c). There $c_n$ is compared with the income $\mathsf{P}_n$
through $(M)_C$.
\item The join is place~(d). There $c_{n+1}$ is compared with incomes through $(M)_C$ for the old
constituents and the connectors, and with deposits through $(B)_C$ for the new constituents. For
the new constituents, $(B)_C$ at level $n+1$ gives
\begin{equation*}
\mathsf{D}_{n+1}\ge\Gamma c_{n+1}(9+4c_0+4R_{n+1})\ge 5\Gamma c_{n+1}.
\end{equation*}
\item The identity \cite[(1.83)]{B-CMP122b} is place~(b). It carries the same constant on both
sides, so no comparison of $C$ with a deposit or an income is made there.
\end{itemize}
These three comparisons and the identity are the first four places of the list.

\emph{Place (e).} The survivor rule has a second occupant for which \cite{B-CMP122b} writes no
factor. A component $Y$ re-included as second class \cite[p.~378]{B-CMP122b} is continued by
$Z=S(Y)$, in the same way as a preparatory survivor. However, the product in
\cite[(1.79)]{B-CMP122b} of factors $\exp(-p_0(g_j))$ is taken ``over components \dots{} for which
some large fields are created during the preparatory steps''. The case list
\cite[pp.~384--387]{B-CMP122b} has no case for such a component. A number $\kappa_{j+1}(S(Y))$
built from $\kappa_j(Y)\ge0$ alone does not satisfy \cite[(1.80)]{B-CMP122b}, whose right member is
positive, at any value of the charge constant; an income from outside that case list is therefore
required here. This is the fifth place.

Its income is that of the statement that a second-class outcome pays a tag, at $d=4$, counted once.
That statement delivers the factor $e^{-E_{\mathrm{2cl}}(n)}\le e^{-2\mathsf{P}_n}$ inside the
majorant $F^{\natural}$ of the first integral performed later. The delivery uses the nesting clause
of the history-wise modulus majorant (Proposition~\ref{9.2:prop-modulus-majorant}, whose proof is
outstanding). We split
\begin{equation*}
e^{-2\mathsf{P}_n}=e^{-\mathsf{P}_n}\cdot e^{-\mathsf{P}_n}.
\end{equation*}
The first half is the income entering the third clause of
Definition~\ref{9.2:def-prepayment-numbers}. The second half pays the prefactor $e$ of the inner
performed operation in the bound $|\tau_hF|\le e\cdot\sup|F|\cdot F^{\natural}(h)$ of
Proposition~\ref{9.2:prop-modulus-majorant}, whose proof is outstanding. That bound is applied in
the comparison product of a history, Lemma~\ref{9.2:lem-comparison-product}. Under the hypotheses
of that lemma, $\mathsf{P}_n\ge1$, hence $e\cdot e^{-\mathsf{P}_n}\le1$.

With the income $e^{-\mathsf{P}_n}$ in hand, the inequality to be verified at the survival level is
the same as that of the preparatory survivor in the re-creation step of the proof of
Proposition~\ref{9.2:prop-prepayment-induction}. That inequality is $(M)_C$, as at place~(c). The
factor supplied by the statement that a second-class outcome pays a tag is drawn on by no other
factor of the majorant. Therefore the two halves are not used elsewhere.

\emph{Where $C$ enters.} At no other step of the proof of
Proposition~\ref{9.2:prop-prepayment-induction} is $(B)_C$ or $(M)_C$ invoked.
None of the quantities $\mathsf{D}_n$, $\mathsf{P}_n$, $E_{\mathrm{2cl}}(n)$ contains $C$, and the
same holds for the geometric lemmas assumed there. Hence $C$ enters only through $(B)_C$ and
$(M)_C$ at the five places above.

\emph{Form of the conditions.} Each of $(B)_C$ and $(M)_C$ has its left member of degree $2p_0$,
respectively $p_0$, in $\Xi_n$, and a right member built from $CM^dR_n^{d+2}$, of degree $(d+2)r$.
Each is therefore a degree inequality on the exponents, $2p_0>(d+2)r$ respectively $p_0>(d+2)r$,
followed by a ceiling on $\gamma$ that depends on $C$, hence on $c_a$.
\end{proof}

\begin{lemma}[Unrolling the pre-payment numbers along a lineage]\label{9.2:lem-lineage-unrolling}
Let $X$ be a large-field component, let $h\in\operatorname{Hist}(X)$ be a history in the sense of
Definition~\ref{9.2:def-histories-ages}, with first index $j(h)$, and let $k$ be the level of its
top member $X$. Let $D(h)$, $I(h)$ and $E^{[C]}(h)$ be the deposit, income and charge products of
$h$ that appear in Lemma~\ref{9.2:lem-comparison-product}: $D(h)$ has one factor
$\exp(-\mathsf{D}_n(d'(W)+1)-2\mathsf{P}_n)$ for each birth $(n,W)\in\operatorname{births}(h)$,
$I(h)$ has one factor $e^{-\mathsf{P}_n}$ for each re-creation $(n,Z)\in\operatorname{recr}(h)$, and
$E^{[C]}(h)$ has one factor $e^{\mathsf{c}_n(Z)}$, the charge factor, for each member
$(n,Z)\in\operatorname{anc}(h)$. For $(n,Z)\in\operatorname{anc}(h)$ let $D_n(Z)$, $I_n(Z)$ and
$E^{[C]}_n(Z)$ denote these three products restricted, respectively, to the births of the members
of the sub-lineage of $Z$ (levels $\le n$); to the survivals $(m,Z_1)\in\operatorname{recr}(h)$ with
$m<n$ in the sub-lineage of $Z$; and to the members $(m,Z_1)$, $m\le n$, of the sub-lineage of $Z$.
Then for every $C>0$, every history $h$ and every $(n,Z)\in\operatorname{anc}(h)$,
\begin{equation}\label{9.2:eq-lineage-unrolling-1}
  D_n(Z)\, I_n(Z)\, E^{[C]}_n(Z)\le \exp\bigl(-\kappa^{[C]}_n(Z)-2\mathsf{P}_{j(Z)}\bigr),
\end{equation}
where $\kappa^{[C]}_n(Z)$ are the pre-payment numbers of
Definition~\ref{9.2:def-prepayment-numbers}. In particular
\begin{equation}\label{9.2:eq-lineage-unrolling-2}
  D(h)\, I(h)\, E^{[C]}(h)\le \exp\bigl(-\kappa^{[C]}_k(X)-2\mathsf{P}_{j(h)}\bigr).
\end{equation}
\end{lemma}

\begin{proof}
Fix $C>0$ and $h$. We prove \eqref{9.2:eq-lineage-unrolling-1} by induction along the lineage tree
of $Z$: a member is treated after all members of its sub-lineage at lower levels. We write
\emph{left member} for the left-hand side of \eqref{9.2:eq-lineage-unrolling-1}.

\emph{Members with no old constituent.} Let $(n,Z)\in\operatorname{anc}(h)$ have no old constituent,
so that $Z$ is made of new constituents $W_i$ born at level $n$. Its sub-lineage consists of the
member $(n,Z)$ alone, so $I_n(Z)$ is the empty product, equal to $1$, $D_n(Z)$ is the product of the
deposits of the $W_i$, and $E^{[C]}_n(Z)$ is the single charge factor of $(n,Z)$. Hence the left
member is
\begin{equation*}
  \prod_i \exp\bigl(-\mathsf{D}_n(d'(W_i)+1)-2\mathsf{P}_n\bigr)\cdot e^{\mathsf{c}_n(Z)}
  =\exp\bigl(-\kappa^{[C]}_n(Z)-2\mathsf{P}_{j(Z)}\bigr),
\end{equation*}
the equality being case~$(\kappa1)$ of \eqref{9.2:eq-prepayment-numbers-a} with $m=0$ old
constituents, together with $j(Z)=n$. Here \eqref{9.2:eq-lineage-unrolling-1} holds with equality.

\emph{From level $n$ to level $n+1$.} Let $(n+1,Z')\in\operatorname{anc}(h)$ have constituents, and
assume \eqref{9.2:eq-lineage-unrolling-1} for every old constituent of $Z'$ (these are members at
level $n$ of its sub-lineage). The sub-lineage of $Z'$ at levels $\le n+1$ consists of the
sub-lineages of its old constituents $Z^{(v)}$ together with the member $(n+1,Z')$. Unwinding the
three restricted products accordingly, the left member for $Z'$ is the product of:
the left members for the old constituents $Z^{(v)}$; a factor $e^{-\mathsf{P}_n}$ for each old
constituent re-created at level $n$ (these are the survivals $(m,Z_1)$ with $m=n$, the only ones in
$I_{n+1}(Z')$ not already in some $I_n(Z^{(v)})$); the deposits
$\exp(-\mathsf{D}_{n+1}(d'(W_i)+1)-2\mathsf{P}_{n+1})$ of the new constituents $W_i$ born at level
$n+1$; and the charge factor $e^{\mathsf{c}_{n+1}(Z')}$ of the member $(n+1,Z')$. We treat the three
cases of \eqref{9.2:eq-prepayment-numbers-a} in turn.

In case $(\kappa1)$ we bound each income factor $e^{-\mathsf{P}_n}$ by $1$, the number
$\mathsf{P}_n=A_0\Xi_n^{p_0}$ being non-negative, and apply the induction hypothesis to each left
member for $Z^{(v)}$. Writing $j(v):=j(Z^{(v)})$, the product is at most
\begin{align*}
  &\exp\Bigl(-\sum_v\bigl[\kappa^{[C]}_n(Z^{(v)})+2\mathsf{P}_{j(v)}\bigr]
   -\sum_i\bigl[\mathsf{D}_{n+1}(d'(W_i)+1)+2\mathsf{P}_{n+1}\bigr]
   +\mathsf{c}_{n+1}(Z')\Bigr)\\
  &\qquad=\exp\bigl(-\kappa^{[C]}_{n+1}(Z')-2\mathsf{P}_{j(Z')}\bigr),
\end{align*}
where the equality is the definition of $\kappa^{[C]}_{n+1}(Z')$ in case $(\kappa1)$.

In case $(\kappa2)$ the product consists of the left member for the single old constituent $Z_0$
and the charge factor $e^{\mathsf{c}_{n+1}(Z')}$. By the induction hypothesis for $Z_0$,
\begin{equation*}
  (\text{left member for }Z_0)\cdot e^{\mathsf{c}_{n+1}(Z')}
  \le\exp\bigl(-\kappa^{[C]}_n(Z_0)-2\mathsf{P}_{j(Z_0)}+\mathsf{c}_{n+1}(Z')\bigr)
  =\exp\bigl(-\kappa^{[C]}_{n+1}(Z')-2\mathsf{P}_{j(Z')}\bigr),
\end{equation*}
the equality being case~$(\kappa2)$ of \eqref{9.2:eq-prepayment-numbers-a} together with
$j(Z')=j(Z_0)$.

In case $(\kappa3)$ the product consists of the left member for the single old constituent $Z_0$,
the income factor $e^{-\mathsf{P}_n}$ of its re-creation at level $n$, and the charge factor
$e^{\mathsf{c}_{n+1}(Z')}$. By the induction hypothesis for $Z_0$ it is at most
\begin{equation*}
  \exp\bigl(-\kappa^{[C]}_n(Z_0)-2\mathsf{P}_{j(Z_0)}-\mathsf{P}_n+\mathsf{c}_{n+1}(Z')\bigr)
  =\exp\bigl(-\kappa^{[C]}_{n+1}(Z')-2\mathsf{P}_{j(Z')}\bigr),
\end{equation*}
the equality being case~$(\kappa3)$ of \eqref{9.2:eq-prepayment-numbers-a}.

This completes the induction and proves \eqref{9.2:eq-lineage-unrolling-1}. At the top of the
lineage tree, $(n,Z)=(k,X)$ and the sub-lineage of $X$ is all of $h$, so that $D_k(X)=D(h)$,
$I_k(X)=I(h)$, $E^{[C]}_k(X)=E^{[C]}(h)$ and $j(X)=j(h)$; then \eqref{9.2:eq-lineage-unrolling-1}
is \eqref{9.2:eq-lineage-unrolling-2}.
\end{proof}

\begin{corollary}[The fundamental inequality in cell units]\label{9.2:cor-fundamental-inequality}
Assume the hypotheses of Proposition~\ref{9.2:prop-prepayment-induction} (the pre-payment
induction at a given charge constant) with charge constant $C>0$. Then for every
level $k$, every $X\in\operatorname{comp}(Z_k)$ and every history $h\in\operatorname{Hist}(X)$,
\begin{equation}\label{9.2:eq-fundamental-inequality-1}
D(h)\,I(h)\,E^{[C]}(h)\;\le\;\exp\bigl(-\Gamma c_k-2\mathsf{P}_{j(h)}\bigr)
\;<\;\exp\bigl(-2p_0(g_{j(h)})\bigr).
\end{equation}
\end{corollary}

Inequality \eqref{9.2:eq-fundamental-inequality-1} is the form of \cite[(1.89)]{B-CMP122b},
with the additional term that accompanies it in \cite{B-CMP122b}, taken at the charge
constant $C$ and measured in cells.

\begin{proof}
Fix $k$, $X\in\operatorname{comp}(Z_k)$ and $h\in\operatorname{Hist}(X)$.

Lemma~\ref{9.2:lem-lineage-unrolling}, unrolling the pre-payment numbers along a lineage, holds
for every $C>0$. In its particular form at the last member $(k,X)$ of the lineage of $h$ it gives
\begin{equation}\label{9.2:eq-fundamental-inequality-2}
D(h)\,I(h)\,E^{[C]}(h)\;\le\;\exp\bigl(-\kappa^{[C]}_k(X)-2\mathsf{P}_{j(h)}\bigr).
\end{equation}

Next we apply Proposition~\ref{9.2:prop-prepayment-induction}, whose hypotheses are assumed, at
level $n=k$ with $Z=X$. This gives \eqref{9.2:eq-prepayment-induction-b} at $n=k$:
\begin{equation*}
\kappa^{[C]}_k(X)\;\ge\;\Gamma'_k c_k\bigl(d'(\mathcal{C}_k(X))+1\bigr)+4\Gamma c_k K_k(X),
\end{equation*}
together with the particular consequence that
Proposition~\ref{9.2:prop-prepayment-induction} records after it, which at $n=k$ and $Z=X$
reads $\kappa^{[C]}_k(X)\ge\Gamma c_k$. Since the
exponential is increasing, inserting this lower bound into \eqref{9.2:eq-fundamental-inequality-2}
gives the first inequality of \eqref{9.2:eq-fundamental-inequality-1}.

For the second inequality, recall that $\mathsf{P}_{j(h)}=p_0(g_{j(h)})$ by definition, where
$p_0(\cdot)$ is the degree function. Moreover $\Gamma>0$, and $c_k=CM^dR_k^{d+1}>0$ since $C$,
$M$ and $R_k$ are positive. Hence $-\Gamma c_k-2\mathsf{P}_{j(h)}<-2p_0(g_{j(h)})$. Since the
exponential is strictly increasing, the second inequality follows.
\end{proof}

\begin{definition}[The one-sided conditions of geometric forgetting]\label{9.2:def-forgetting-conditions}
Let $d=4$ and let $k\ge 1$ be a level. The \emph{one-sided conditions of geometric forgetting}
at level $k$ are the conditions listed in \eqref{9.2:eq-forgetting-conditions-a} below. They
are the hypotheses of the statement on geometric forgetting of old regions. At $d=4$ the
following values hold for $\sqrt{d}$, for the $d$-dependent constants of
Notation~\ref{9.2:not-parameters-couplings}, and for the number $\beta_n$ of free bonds per cell:
\[
\sqrt{d}=2,\qquad \Gamma=64\cdot 219^4,\qquad \sigma=440,\qquad c_0=32,
\]
\[
4\sqrt{d}\cdot 101^d=8\cdot 101^4,\qquad \beta_n=4(LMR_n)^4 .
\]
The conditions are:
\begin{equation}\label{9.2:eq-forgetting-conditions-a}
\begin{aligned}
&(\mathrm{L}) && L\ \text{odd},\ L\ge 5;\quad M_1<M_2=L^2M_1<M\ \text{powers of}\ L,
  \ \text{so}\ M\ge 4M_2;\\
&(\mathrm{E}) && r\ge 1,\quad p_0\ge (d+2)r+1=6r+1,\quad A_0\ge 1;\\
&(c_a) && c_a\ge \log 2,\quad q:=e^{-c_a},\quad C:=C_0+c_a;\\
& && \text{for every running coupling } g_n,\ n\le k,\ \text{with } \Xi_n=\log g_n^{-2}:\\
&(\mathrm{P}) && \Xi_n\ge 1;\\
&(\mathrm{S}) && \Xi_n^{p_0}\ge \frac{1024(1+\beta_0)^2\Theta_1^2A_0}{\gamma_0A_1^2};\\
&(\mathrm{B})_C && \Xi_n^{2p_0-6r}\ge
  \frac{1024(1+\beta_0)^2\Theta_1^2\,\Gamma(13+4c_0)\,C\,M^4L^6}{\gamma_0A_1^2};\\
&(\mathrm{M})_C && \Xi_n^{p_0-6r}\ge
  \frac{(1+\beta_0)\,\Gamma\bigl(5\sigma+8\cdot 101^4+4c_0+52\bigr)\,C\,M^4L^6}{A_0};\\
&(\mathrm{F}_I) && \log 4+4\log(LM)+4\log L+4r\log\Xi_n\le 2\mathsf{Q}_n;\\
&(\gamma) && \text{the ceilings on}\ \gamma\ \text{stated after this display};\\
&(g_\bullet) && g_n\le g_\bullet(L,M,M_2,B_3,A_0,r,p_0,M_0,C_\Lambda,4),
  \ \text{that is}\ \Xi_n\ge \log g_\bullet^{-2};\\
&(\mathrm{St}) && \text{the standing conditions of the setting, stated after this display.}
\end{aligned}
\end{equation}
The lines are to be read as follows. Lines $(\mathrm{L})$, $(\mathrm{E})$ and $(c_a)$ include the
further requirements stated in the next two paragraphs; these are part of the conditions.

In $(\mathrm{L})$ the floors on $L$ imposed by the correlation theorem remain in force, and all
the cell partitions involved are compatible; the scales $M_1<M_2<M$ are those of
\cite{B-CMP119}. Since $M$ and $M_2$ are powers of $L$ with $M>M_2$, we have $M\ge LM_2\ge 4M_2$,
and this inequality is also the one used in the statement that a second-class outcome pays a tag.
In $(\mathrm{E})$ the floors on the exponents imposed by the correlation theorem remain in force.
The floor on $p_0$ gives $p_0-6r\ge 1$, and hence $2p_0-6r\ge p_0+1\ge 1$.

In $(c_a)$ the exponent $c_a$ is fixed after $(L,M,C_0)$ and before $\gamma$. Here $C_0$ is the
charge constant of the history-wise modulus majorant,
Proposition~\ref{9.2:prop-modulus-majorant}, whose proof is outstanding, and it is the charge
constant at which the charge bound is applied.

Line $(\mathrm{P})$ together with $A_0\ge 1$ from $(\mathrm{E})$ gives
$\mathsf{P}_n=A_0\Xi_n^{p_0}\ge 1$. This pays the one prefactor $e$ incurred for each survival as
second class in Lemma~\ref{9.2:lem-comparison-product}.

Line $(\mathrm{S})$ is the inequality $\mathsf{Q}_n\ge 4\mathsf{P}_n$, that is, the condition
\eqref{9.2:eq-creation-deposit-a} of Lemma~\ref{9.2:lem-creation-deposit}.

Line $(\mathrm{B})_C$ implies $\mathsf{D}_n\ge \Gamma c_n(9+4c_0+4R_n)$, which is hypothesis
$(\mathrm{B})_C$ of Proposition~\ref{9.2:prop-prepayment-induction},
\eqref{9.2:eq-prepayment-induction-a}. Line $(\mathrm{M})_C$ implies
\[
(1+\beta_0)^{-1}\mathsf{P}_n\ge \Gamma c_n\bigl(5\sigma+8\cdot 101^4+4c_0+4R_n+48\bigr),
\]
which is hypothesis $(\mathrm{M})_C$ of Proposition~\ref{9.2:prop-prepayment-induction},
\eqref{9.2:eq-prepayment-induction-a}, at $d=4$. Both are hypotheses of
Proposition~\ref{9.2:prop-prepayment-induction}, which is applied here at $C=C_0+c_a$.

Line $(\mathrm{F}_I)$ implies $\beta_n\upsilon_n\le 1$, which is used in
Proposition~\ref{9.2:prop-history-summation}.

Line $(\gamma)$ consists of the ceilings on $\gamma$ under which the first and the third of the
flow conditions assumed in Proposition~\ref{9.2:prop-prepayment-induction} hold. In $(g_\bullet)$
the function $g_\bullet$ is the coupling threshold of the statement that a second-class outcome
pays a tag, taken at $d=4$.

The two settings are those in which the parameter domain $\mathfrak{P}$ is taken: in setting (i),
one run of the correlation theorem, $\mathfrak{P}$ is the analyticity tube of the external fields
times the weight domain $\mathcal{W}_k$; in setting (ii), the construction on the dial polydisc
$\mathbb{D}_\delta$, $\mathfrak{P}$ is $\mathbb{D}_\delta$ times that tube times $\mathcal{W}_k$.
In $(\mathrm{St})$, setting (i) requires the standing conditions of the correlation theorem.
Setting (ii) requires the standing conditions of the construction on the dial polydisc
$\mathbb{D}_\delta$, in particular their clause, on the raw Wilson coefficients, which bounds by
$\tfrac14$ the sum of the relative shift $\theta_W$ and the correction term of those conditions.
This clause gives $c_W\ge \tfrac34$ on $\mathbb{D}_\delta$.

In setting (ii) the standing conditions of the construction on the dial polydisc
$\mathbb{D}_\delta$ contain a clause that uses, in place of the unspecified constant $O(1)$ of
\cite{B-CMP122b}, the number
\[
\mathsf{c}:=C_\sharp O(1)+1+C_{L'}+c_a
\]
(a constant of those standing conditions, unrelated to the charges $\mathsf{c}_n(Z)$).
That clause is followed by the conditions ``$p_0$ large, and $g_1$ sufficiently small'' and
``for $p_0$ large and $\gamma$ small enough'' of \cite{B-CMP122b}. For the statement on geometric
forgetting of old regions, this clause and these conditions are replaced by two inputs: the charge
bound at $C_0$, and Proposition~\ref{9.2:prop-prepayment-induction} at $C=C_0+c_a$. In other
words, they are replaced by the lines $(\mathrm{P})$, $(\mathrm{S})$, $(\mathrm{B})_C$,
$(\mathrm{M})_C$, $(\mathrm{F}_I)$, $(\gamma)$ and $(g_\bullet)$ above, with $M\ge 4M_2$. The
floor $p_0>(d+2)r_{\mathrm{pr}}$ on the exponent used with those standing conditions, where
$r_{\mathrm{pr}}$ is the exponent in the definition of the length $R$ in \cite{B-CMP122b}, is of
the same kind as $(\mathrm{E})$; the conditions used here require $p_0\ge 6r+1$.
\end{definition}

\begin{lemma}[The conditions reduce to one coupling ceiling]\label{9.2:lem-coupling-ceiling}
Let $d=4$, let the lines $(\mathrm{L})$, $(\mathrm{E})$, $(c_a)$ of \eqref{9.2:eq-forgetting-conditions-a}
hold, and let $c_2$ be the constant of Definition~\ref{1.2:def-flow-constants} entering the lower
bound $g_n^{-2}\ge\gamma^{-2}-2c_2$ on the running inverse couplings.
\begin{itemize}
\item[(a)] Each of the lines $(\mathrm{P})$, $(\mathrm{S})$, $(\mathrm{B})_C$, $(\mathrm{M})_C$,
$(\mathrm{F}_I)$, $(g_\bullet)$ of \eqref{9.2:eq-forgetting-conditions-a}, read as a condition on a
real variable $\Xi\ge 1$ in place of $\Xi_n$, holds for all sufficiently large $\Xi$. Let $\Xi_\ast$
be the supremum of the set of $\Xi\ge1$ at which at least one of these lines fails, with
$\Xi_\ast:=1$ if this set is empty. Then $\Xi_\ast<\infty$; $\Xi_\ast$ depends only on
$L$, $M$, $M_2$, $A_0$, $A_1$, $\gamma_0$, $\beta_0$, $\Theta_1$, $B_3$, $M_0$, $C_\Lambda$, $r$, $p_0$
and $C$; and every one of these lines
holds at every $\Xi>\Xi_\ast$.
\item[(b)] There is $\gamma^\sharp>0$, depending only on the parameters listed in \textup{(a)} and on
$c_2$, such that $\gamma\le\gamma^\sharp$ implies $\log(\gamma^{-2}-2c_2)>\Xi_\ast$, and then, by the
lower bound $g_n^{-2}\ge\gamma^{-2}-2c_2$ on the running couplings in the correlation theorem, all the
lines of \textup{(a)} hold for every running coupling $g_n$, $n\le k$.
\item[(c)] The lines $(\mathrm{C0})$ and $(\mathrm{SL}/\mathrm{SL}^2)$ are ceilings on $\gamma$. Thus
every condition in \eqref{9.2:eq-forgetting-conditions-a} other than $(\mathrm{L})$, $(\mathrm{E})$,
$(c_a)$ is a ceiling on $\gamma$, chosen last, after $d$, $L$, $M$, $M_2$, $\kappa$, $\kappa_0$, the
exponents, $r$, $C_0$, $c_a$; each is monotone in the sense that a larger $C$ (a larger $c_a$) or a
larger $L$, $M$ asks only for a smaller $\gamma$. There is no upper bound on $L$, $M$, $k$, the
number of levels, ages, lifetimes, cascade lengths, numbers of constituents or the volume, and no
lower bound on $\gamma$; the clock $K_n(Z)$ and the integer index of the large-field construction
are local integers and occur in no constant.
\end{itemize}
\end{lemma}

\begin{proof}
\emph{Step 1: what the lines secure.} We first derive, at $d=4$, the inequalities for which the lines
of \eqref{9.2:eq-forgetting-conditions-a} are imposed. Recall from
\eqref{9.2:eq-parameters-couplings-a} that $c_n=CM^dR_n^{d+1}$, and that the
length satisfies $R_n\ge1$ and $R_n<L\Xi_n^r$. Since $R_n\ge1$ we have $9+4c_0+4R_n\le(13+4c_0)R_n$,
where $c_0$ is the $d$-only geometric constant of the core sandwich, hence
\begin{align*}
\Gamma c_n(9+4c_0+4R_n)&\le\Gamma(13+4c_0)CM^dR_n^{d+2}\\
&<\Gamma(13+4c_0)CM^dL^{d+2}\Xi_n^{(d+2)r}.
\end{align*}
By \eqref{9.2:eq-creation-deposit-a} the linear-size deposit is
\begin{equation*}
\mathsf{D}_n=\frac{\gamma_0A_1^2\Xi_n^{2p_0}}{512\sqrt{d}\,\Theta_1^2(1+\beta_0)^2}
=\frac{\gamma_0A_1^2\Xi_n^{2p_0}}{1024\,\Theta_1^2(1+\beta_0)^2}\quad\text{at } d=4,
\end{equation*}
so multiplying $(\mathrm{B})_C$ by $\gamma_0A_1^2\Xi_n^{6r}/(1024(1+\beta_0)^2\Theta_1^2)$ gives
\begin{equation*}
\mathsf{D}_n\ge\Gamma(13+4c_0)CM^4L^6\Xi_n^{6r}>\Gamma c_n(9+4c_0+4R_n).
\end{equation*}
For $(\mathrm{M})_C$ the argument is the same: $R_n\ge1$ gives
\begin{equation*}
5\sigma+8\cdot101^4+4c_0+4R_n+48\le(5\sigma+8\cdot101^4+4c_0+52)R_n,
\end{equation*}
so
\begin{align*}
\Gamma c_n(5\sigma+8\cdot101^4+4c_0+4R_n+48)
&<\Gamma(5\sigma+8\cdot101^4+4c_0+52)CM^4L^6\Xi_n^{6r}\\
&\le(1+\beta_0)^{-1}\mathsf{P}_n,
\end{align*}
the last inequality by multiplying $(\mathrm{M})_C$ by $A_0\Xi_n^{6r}/(1+\beta_0)$ and using
$\mathsf{P}_n=A_0\Xi_n^{p_0}$. The implication
$(\mathrm{S})\Rightarrow\mathsf{Q}_n\ge4\mathsf{P}_n$ follows from \eqref{9.2:eq-parameters-couplings-a}.
For $(\mathrm{F}_I)$: the number $\beta_n$ of free bonds per cell is
$\beta_n=4(LMR_n)^4<4(LM)^4L^4\Xi_n^{4r}$, so
\begin{equation*}
\log\beta_n<\log4+4\log(LM)+4\log L+4r\log\Xi_n\le2\mathsf{Q}_n=\log\upsilon_n^{-1},
\end{equation*}
that is, $\beta_n\upsilon_n\le1$. For $(\mathrm{P})$: $\mathsf{P}_n=A_0\Xi_n^{p_0}\ge A_0$ by
$\Xi_n\ge1$, and $A_0\ge1$ by $(\mathrm{E})$.

\emph{Step 2: each line holds for large $\Xi$.} Line $(\mathrm{P})$ reads $\Xi\ge1$ and holds on the
whole range. In $(\mathrm{S})$, $(\mathrm{B})_C$, $(\mathrm{M})_C$ the left-hand sides are the powers
$\Xi^{p_0}$, $\Xi^{2p_0-6r}$, $\Xi^{p_0-6r}$, whose exponents are $\ge1$ by $(\mathrm{E})$ (which gives
$p_0\ge6r+1$, hence $2p_0-6r\ge1$ and $p_0-6r\ge1$), and the right-hand sides are constants; a
positive power of $\Xi$ exceeds a fixed constant for all large $\Xi$. In $(\mathrm{F}_I)$ the
polynomial $2\mathsf{Q}(\Xi)$, that is $2\mathsf{Q}_n$ written as a function of $\Xi=\Xi_n$, of degree
$2p_0\ge2$ and with positive leading coefficient (by the definition of the per-bad-cube deposit
$\mathsf{Q}_n$, see \eqref{9.2:eq-parameters-couplings-a}), is compared with the constant
$\log4+4\log(LM)+4\log L$ plus $4r\log\Xi$; it exceeds this for all large $\Xi$. Line $(g_\bullet)$
reads $\Xi\ge\log g_\bullet^{-2}$, with $g_\bullet$ the function of line $(g_\bullet)$, whose arguments
are $L$, $M$, $M_2$, $B_3$, $A_0$, $r$, $p_0$, $M_0$, $C_\Lambda$ and the number $4$,
and holds for all large $\Xi$.

\emph{Step 3: the threshold $\Xi_\ast$.} By Step~2 the set of $\Xi\ge1$ at which a given line fails
is bounded above; the set at which at least one line fails is the union of six such sets, hence
bounded above, and $\Xi_\ast<\infty$. By the definition of the supremum no line fails at any
$\Xi>\Xi_\ast$. The constants entering the six lines are the numbers $\Gamma$, $\sigma$, $c_0$, which
depend on $d$ only, and the parameters listed in (a); hence so does $\Xi_\ast$. This proves (a).

\emph{Step 4: the ceiling $\gamma^\sharp$.} Put
\begin{equation}\label{9.2:eq-coupling-ceiling-1}
\gamma^\sharp:=\bigl(e^{\Xi_\ast+1}+2|c_2|\bigr)^{-1/2}>0 ,
\end{equation}
which depends only on $\Xi_\ast$ and $c_2$; the absolute value is used only to avoid a discussion of
the sign of $c_2$. If $\gamma\le\gamma^\sharp$ then
$\gamma^{-2}-2c_2\ge e^{\Xi_\ast+1}>0$, so $\log(\gamma^{-2}-2c_2)\ge\Xi_\ast+1$. By the lower
bound on the running couplings in the correlation theorem, every running coupling obeys
$g_n^{-2}\ge\gamma^{-2}-2c_2$; since the logarithm is increasing and $\Xi_\ast\ge1$,
\begin{equation*}
\Xi_n=\log g_n^{-2}\ge\log(\gamma^{-2}-2c_2)\ge\Xi_\ast+1>\Xi_\ast\ge1,\qquad n\le k.
\end{equation*}
By Step~3 every line of (a) then holds at $\Xi=\Xi_n$, for all $n\le k$. This
proves (b).

\emph{Step 5: the remaining lines and the order of choice.} The line $(\mathrm{C0})$ consists, as it
stands in \eqref{9.2:eq-forgetting-conditions-a}, of the ceilings on $\gamma$ under which the flow facts
\begin{equation*}
\mathsf{P}_m\ge(1+\beta_0)^{-1}\mathsf{P}_n\quad(m\le n\le k),\qquad \ell_n\in\{1,L\}
\end{equation*}
hold, the second together with the rule that $\ell_n=1$ forces ($n=1$ or $\ell_{n-1}=L$) and
$\ell_{n+1}=L$; here $\mathsf{P}_n$ is the value of the degree function at $g_n$ and $\ell_n$ the
coarsening ratio. The line $(\mathrm{SL}/\mathrm{SL}^2)$
consists of the standing conditions of the two settings, which are ceilings on $\gamma$. Together with
(b), every condition of \eqref{9.2:eq-forgetting-conditions-a} other than $(\mathrm{L})$, $(\mathrm{E})$,
$(c_a)$ is a ceiling on $\gamma$. The number $\gamma^\sharp$ depends only on $\Xi_\ast$ and $c_2$,
neither of which depends on $\gamma$; hence $\gamma$ is chosen last, after $d$, $L$, $M$, $M_2$,
$\kappa$, $\kappa_0$, the exponents, $r$, $C_0$ and $c_a$ (recall $C=C_0+c_a$). The constant sides of
$(\mathrm{B})_C$ and $(\mathrm{M})_C$ are nondecreasing in $C$, $L$, $M$; the constant side
$\log4+4\log(LM)+4\log L$ of $(\mathrm{F}_I)$ is nondecreasing in $L$, $M$ and does not involve $C$;
and the constant side of $(\mathrm{S})$ involves none of them. The $\Xi$-dependent sides
$\Xi^{p_0}$, $\Xi^{2p_0-6r}$, $\Xi^{p_0-6r}$ and $2\mathsf{Q}(\Xi)$ involve none of $C$, $L$, $M$ (for
$\mathsf{Q}$ by \eqref{9.2:eq-parameters-couplings-a}). For these four lines a larger $C$ (a
larger $c_a$) or a larger $L$, $M$ therefore enlarges the set of $\Xi$ at which the line fails. For
every line, such a change alters only the value of the ceiling on $\gamma$ (for the lines of (a),
through $\Xi_\ast$ and \eqref{9.2:eq-coupling-ceiling-1}); since no line bounds $C$, $L$ or $M$ from
above, the change asks only for a smaller $\gamma$. The
lines $(\mathrm{L})$, $(\mathrm{E})$, $(c_a)$ are lower bounds, and all other lines are upper bounds on
$\gamma$; none of the lines involves an upper bound on $L$, $M$, $k$, the number of levels, ages,
lifetimes, cascade lengths, numbers of constituents or the volume, or a lower bound on $\gamma$.
Finally, neither the clock $K_n(Z)$ nor the local integer index of the large-field construction
(which is not the number $K$ of renormalisation steps)
appears in any line of \eqref{9.2:eq-forgetting-conditions-a}, so neither occurs in a constant. This
proves (c).
\end{proof}

Throughout the following theorem the dimension is $d=4$ and the gauge group is $\mathrm{SU}(2)$. For $\mathrm{SU}(N)$ only the constant $\bar\gamma_0=1/N$ changes.

We fix the following data.
\begin{itemize}
\item An odd integer $L$, block sizes $M_1<M_2<M$, exponents $p_0$ and $r$, and constants $A_0,A_1,\gamma_0,\beta_0,B_3,\Theta_1$.
\item A free number $c_a\ge\log 2$; we put $q:=e^{-c_a}$, so that $q\in(0,\tfrac12]$.
\end{itemize}

For a level $k\ge1$, let $X$ be a component of the large-field region $Z_k$ of level $k$ on which the integral \cite[(1.71)]{B-CMP122b} of Ba{\l}aban's renormalization operation $\mathbf{R}$ is performed at step $k$; its past is arbitrary. We write that operation as the sum over the admissible histories $h\in\operatorname{Hist}(X)$ of $X$:
\[
\mathbf{T}'_k(X)=\sum_{h\in\operatorname{Hist}(X)}\tau_h .
\]
For an integer $a\ge0$ its part of age $a$ is
\[
\mathbf{T}'_{k,a}(X):=\sum_{h:\,a(h)=a}\tau_h .
\]
Here $a(h):=k-j(h)$, and $j(h)$ is the index of the first large-field region contained in $X$ along $h$ (see \cite{B-CMP122b} and Definition~\ref{9.2:def-integral-by-age}).

For a coupling $g$ we use the degree function
\[
p_0(g):=A_0\bigl(\log g^{-2}\bigr)^{p_0}
\]
of \cite{B-CMP119} (the degree function of the glossary, Notation~\ref{0.2:not-glossary}). We write $\Xi_n:=\log g_n^{-2}$ and $\mathsf{P}_n:=p_0(g_n)=A_0\Xi_n^{p_0}$. Finally $e:=\exp(1)$.

\begin{theorem}[Geometric forgetting of old regions]\label{9.2:thm-geometric-forgetting}
Let $k\ge1$ be a level with running couplings $g_1,\dots,g_k$. Assume the following.
\begin{enumerate}
\item[(A)] The history-wise modulus majorant, Proposition~\ref{9.2:prop-modulus-majorant}, whose proof is outstanding, with its clauses (a), (b) and (c) and its charge constant $C_0$.
\item[(B)] The following statements on admissible histories:
\begin{itemize}
\item the creation-deposit bound, that is, the bound on the deposit made when a large-field region is created;
\item the per-cell charge bound, that is, the bound on the charges per $MR_n$-cell at the charge constant $C_0$;
\item the income bound for a survival after a preparatory failure;
\item Lemma~\ref{9.2:lem-live-levels}, with its clauses (a)--(c);
\item Lemma~\ref{9.2:lem-comparison-product}.
\end{itemize}
\item[(C)] Summation over admissible histories at a cost per cell, Proposition~\ref{9.2:prop-history-summation}.
\item[(D)] The second-class income statement in dimension $d=4$: a processed component that survives as second class supplies an income.
\item[(E)] The three flow facts on the running couplings $g_1,\dots,g_k$ that are hypotheses of Proposition~\ref{9.2:prop-prepayment-induction}, and the further flow fact on the running couplings relative to the flow constants of Definition~\ref{1.2:def-flow-constants}. The first of the three yields the comparison
\[
\mathsf{P}_n\ge(1+\beta_0)^{-1}\mathsf{P}_k\qquad\text{for every level }n\le k.
\]
\item[(F)] The one-sided conditions \eqref{9.2:eq-forgetting-conditions-a} of Definition~\ref{9.2:def-forgetting-conditions}. These are:
\begin{itemize}
\item floors on $L$, on $M$, on the exponents and on $c_a$, fixed first;
\item ceilings on the coupling $\gamma$, chosen last, after $d$, $L$, $M$, $M_2$, the constants of the correlation theorem and of the construction with the complex dial variable, the exponents, and $c_a$.
\end{itemize}
\end{enumerate}
Let $X$ be any component of $Z_k$ on which the integral \cite[(1.71)]{B-CMP122b} is performed at step $k$, with any past. Let $a\ge0$ be any integer and $F$ any bounded function of the variables that \cite[(1.71)]{B-CMP122b} does not integrate. Fix any point of the parameter domain $\mathfrak{P}$, which is one of the following:
\begin{itemize}
\item in setting (i), a single run of the correlation theorem, $\mathfrak{P}$ is the analyticity tube of the external fields $U$ and $J$ of that theorem times the weight domain $\mathcal{W}_k$;
\item in setting (ii), the construction with the complex dial variable, $\mathfrak{P}$ is the dial polydisc $\mathbb{D}_\delta$ of that construction times that tube times $\mathcal{W}_k$.
\end{itemize}
Then
\begin{equation}\label{9.2:eq-geometric-forgetting-a}
\begin{split}
\bigl|\mathbf{T}'_{k,a}(X)F\bigr|
&\le e\cdot\exp\bigl(-2p_0(g_{k-a})\bigr)\cdot q^{a+1}\cdot\sup|F|\\
&\le e\cdot\exp\bigl(-2(1+\beta_0)^{-1}p_0(g_k)\bigr)\cdot q^{a+1}\cdot\sup|F|.
\end{split}
\end{equation}
Here $\sup|F|$ is the supremum of $|F|$ over its arguments. No constant in \eqref{9.2:eq-geometric-forgetting-a} depends on $k$, on the number of levels of the whole construction, on $a$, on the volume or on the position of $X$.
\end{theorem}

\begin{remark}[Scope of the bound]
In Theorem~\ref{9.2:thm-geometric-forgetting} the component $X$ ranges over the components on which \cite[(1.71)]{B-CMP122b} is performed at step $k$, with arbitrary admissible past. Two further cases are handled as follows.
\begin{itemize}
\item For the components alive at the last level, the same bound holds inside the final integral. This is the corollary of Theorem~\ref{9.2:thm-geometric-forgetting} for these components, in whose proof the second-class income statement is used inside the final integral.
\item For any other living component, the majorant of its accumulated operation is formed only inside a later operation of one of these two kinds, and the bound holds there.
\end{itemize}
This is the sense of the remarks of \cite{B-CMP122b} that \cite[(1.79)]{B-CMP122b} ``holds for the T-operation connected with an arbitrary large field region'' and that \cite[(1.80)]{B-CMP122b} ``holds quite generally for such regions''. Every application of Theorem~\ref{9.2:thm-geometric-forgetting} in this book takes place inside an operation of one of these two kinds.
\end{remark}

\begin{proof}
Assume hypotheses (A)--(F). Fix $k$, a component $X\in\operatorname{comp}(Z_k)$ on which the integral of $\mathbf{R}$ is performed at step $k$, an integer $a\ge0$, a bounded $F$, and a point of $\mathfrak{P}$. Let
\[
\mathcal{A}:=\{h\in\operatorname{Hist}(X):a(h)=a\}.
\]
For $h\in\mathcal{A}$ we have $j(h)=k-a(h)=k-a$. If $\mathcal{A}$ is empty, then $\mathbf{T}'_{k,a}(X)=0$ and there is nothing to prove. We therefore assume that $\mathcal{A}$ is non-empty.

\begin{subequations}\label{9.2:eq-geometric-forgetting-b}
\emph{Step 1: the sum over histories is replaced by a supremum.}
By definition
\[
\mathbf{T}'_{k,a}(X)F=\sum_{h\in\mathcal{A}}\tau_hF .
\]
We apply the triangle inequality. To each $h\in\mathcal{A}$, at the fixed $F$ and the fixed point of $\mathfrak{P}$, we apply clause (a) of Proposition~\ref{9.2:prop-modulus-majorant}, whose proof is outstanding; this is the bound \eqref{9.2:eq-modulus-majorant-a}. This gives
\begin{equation}\label{9.2:eq-geometric-forgetting-1}
\bigl|\mathbf{T}'_{k,a}(X)F\bigr|\le\sum_{h\in\mathcal{A}}|\tau_hF|
\le e\cdot\sup|F|\cdot\sum_{h\in\mathcal{A}}F^{\natural}(h).
\end{equation}

Next we apply the estimate \eqref{9.2:eq-history-summation-a} of Proposition~\ref{9.2:prop-history-summation}. Its hypotheses hold:
\begin{itemize}
\item $X$ is connected, being a component;
\item the line $(F_I)$ of \eqref{9.2:eq-forgetting-conditions-a} holds for all $n\le k$, by hypothesis (F);
\item the admissible histories are sequences of the form required there.
\end{itemize}
We apply it to the function $\Psi:=F^{\natural}\cdot\mathbf{1}_{\mathcal{A}}$, extended by $0$ to the sequences that are not admissible histories. The function $\Psi$ equals $F^{\natural}$ on $\mathcal{A}$ and vanishes elsewhere, and $F^{\natural}\ge0$ by Proposition~\ref{9.2:prop-modulus-majorant}, whose proof is outstanding. Hence the sum of $\Psi$ is $\sum_{h\in\mathcal{A}}F^{\natural}(h)$. Moreover, since all weighted values are non-negative and $\mathcal{A}$ is non-empty, the supremum of $\Psi$ times the weights of \eqref{9.2:eq-history-summation-a} is its supremum over $\mathcal{A}$. Therefore
\begin{equation}\label{9.2:eq-geometric-forgetting-2}
\begin{split}
\sum_{h\in\mathcal{A}}F^{\natural}(h)
&\le\sup_{h\in\mathcal{A}}\Bigl[F^{\natural}(h)\prod_{n=j(h)}^{k}
\upsilon_n^{-|I_n(h)|}e^{c_A^{\flat}(n)|\mathsf{Z}_n|}\Bigr]\\
&\le\sup_{h\in\mathcal{A}}D(h)\cdot I(h)\cdot E^{[C_0]}(h).
\end{split}
\end{equation}
The last step is the inequality \eqref{9.2:eq-comparison-product-a} of Lemma~\ref{9.2:lem-comparison-product}, applied to each $h\in\mathcal{A}$. The hypotheses of that lemma are as follows.
\begin{itemize}
\item $d=4$, $M\ge4M_2$, the lines (S) and (P), and $g_n\le g_\bullet$ for all $n\le k$: all are contained in \eqref{9.2:eq-forgetting-conditions-a}.
\item $c_W\ge\tfrac34$ on $\mathfrak{P}$. In setting (i) this holds by clause (c) of the lemma on the Wilson coefficient in the proof of the correlation theorem, together with the third step of the proof of the modulus bound of that theorem. In setting (ii) it holds by the clause of \eqref{9.2:eq-forgetting-conditions-a} that provides it.
\end{itemize}
The second-class incomes entering Lemma~\ref{9.2:lem-comparison-product} are those supplied by the second-class income statement, hypothesis (D).

\emph{Step 2: the surplus is split off before the induction.}
Put $C:=C_0+c_a$. For every $h$ we have the identity
\begin{equation}\label{9.2:eq-geometric-forgetting-3}
\begin{split}
E^{[C_0]}(h)
&=E^{[C]}(h)\cdot\exp\Bigl(-(C-C_0)\sum_{(n,Z)\in\operatorname{anc}(h)}\sqrt d\,M^dR_n^{d+1}|Z|\Bigr)\\
&=E^{[C]}(h)\cdot\exp\Bigl(-c_a\sum_{(n,Z)\in\operatorname{anc}(h)}\sqrt d\,M^dR_n^{d+1}|Z|\Bigr),
\end{split}
\end{equation}
whose second line uses $C-C_0=c_a$.

\emph{Step 3: the induction at the constant $C$.}
The hypotheses of Proposition~\ref{9.2:prop-prepayment-induction} hold at $C=C_0+c_a$:
\begin{itemize}
\item the inequalities \eqref{9.2:eq-prepayment-induction-a} hold for all $n\le k$, by \eqref{9.2:eq-forgetting-conditions-a} together with the implications between its lines recorded with Definition~\ref{9.2:def-forgetting-conditions};
\item $L\ge5$ holds by \eqref{9.2:eq-forgetting-conditions-a};
\item the flow facts of hypothesis (E) hold;
\item the hypotheses of the two lemmas on which the proof of that proposition rests hold; these are four structural conditions on the family of regions $(Z_n)$, and they are met because the regions $Z_n$ are those of Ba{\l}aban's construction.
\end{itemize}
We now apply Corollary~\ref{9.2:cor-fundamental-inequality} to each $h\in\mathcal{A}$. Since $j(h)=k-a$, it gives
\begin{equation}\label{9.2:eq-geometric-forgetting-4}
D(h)\cdot I(h)\cdot E^{[C]}(h)\le\exp\bigl(-2p_0(g_{j(h)})\bigr)=\exp\bigl(-2p_0(g_{k-a})\bigr).
\end{equation}

In the proof of Proposition~\ref{9.2:prop-prepayment-induction} the constant $C$ enters only at five places. The incomes used there are:
\begin{itemize}
\item for a component surviving after a preparatory failure, the income of hypothesis (B) for such a survival;
\item for a component surviving as second class, half of the income supplied by the second-class income statement of hypothesis (D).
\end{itemize}
Each income is counted once, in the third clause of the definition of the pre-payment numbers $\kappa^{[C]}_n(Z)$, at the step following the survival. The second-class income exponent $E_{\mathrm{2cl}}(n)$ does not depend on $C$.

\emph{Step 4: the floor per live level.}
Let $h\in\mathcal{A}$. The sum in \eqref{9.2:eq-geometric-forgetting-3} is over $\operatorname{anc}(h)$.
\begin{itemize}
\item By Lemma~\ref{9.2:lem-live-levels}(b) it has $\#\operatorname{anc}(h)\ge a(h)+1=a+1$ terms.
\item By Lemma~\ref{9.2:lem-live-levels}(c) each term satisfies $\sqrt d\,M^dR_n^{d+1}|Z|\ge1$.
\end{itemize}
Since $c_a>0$, it follows that
\begin{equation}\label{9.2:eq-geometric-forgetting-5}
\exp\Bigl(-c_a\sum_{(n,Z)\in\operatorname{anc}(h)}\sqrt d\,M^dR_n^{d+1}|Z|\Bigr)
\le e^{-c_a(a+1)}=q^{a+1}.
\end{equation}
\end{subequations}

\emph{Combination.}
For $h\in\mathcal{A}$ put
\[
x_h:=D(h)I(h)E^{[C]}(h),\qquad
y_h:=\exp\Bigl(-c_a\sum_{(n,Z)\in\operatorname{anc}(h)}\sqrt d\,M^dR_n^{d+1}|Z|\Bigr).
\]
By \eqref{9.2:eq-geometric-forgetting-3} we have $D(h)I(h)E^{[C_0]}(h)=x_hy_h$. By \eqref{9.2:eq-geometric-forgetting-5} we have $0<y_h\le q^{a+1}$.

The inequality \eqref{9.2:eq-comparison-product-a} bounds $x_hy_h$ from below by $F^{\natural}(h)$ times a product of the positive numbers $\upsilon_n^{-|I_n(h)|}$ and $e^{c_A^{\flat}(n)|\mathsf{Z}_n|}$. Since $F^{\natural}(h)\ge0$ by Proposition~\ref{9.2:prop-modulus-majorant}, whose proof is outstanding, this lower bound is non-negative. Hence $x_hy_h\ge0$, and since $y_h>0$, also $x_h\ge0$. Consequently $x_hy_h\le x_hq^{a+1}$.

Combining \eqref{9.2:eq-geometric-forgetting-1}, \eqref{9.2:eq-geometric-forgetting-2}, \eqref{9.2:eq-geometric-forgetting-3}, \eqref{9.2:eq-geometric-forgetting-5} and then \eqref{9.2:eq-geometric-forgetting-4}, we obtain
\begin{align*}
\bigl|\mathbf{T}'_{k,a}(X)F\bigr|
&\le e\cdot\sup|F|\cdot\sup_{h\in\mathcal{A}}\bigl[D(h)I(h)E^{[C]}(h)\cdot q^{a+1}\bigr]\\
&\le e\cdot\sup|F|\cdot\exp\bigl(-2p_0(g_{k-a})\bigr)\cdot q^{a+1}.
\end{align*}
This is the first inequality of \eqref{9.2:eq-geometric-forgetting-a}.

\emph{Step 5: the index of the coupling.}
Since $k-a\le k$, the flow comparison displayed in hypothesis (E) gives
\[
p_0(g_{k-a})=\mathsf{P}_{k-a}\ge(1+\beta_0)^{-1}\mathsf{P}_k=(1+\beta_0)^{-1}p_0(g_k).
\]
Since the exponential is increasing,
\[
\exp\bigl(-2p_0(g_{k-a})\bigr)\le\exp\bigl(-2(1+\beta_0)^{-1}p_0(g_k)\bigr).
\]
This gives the second inequality of \eqref{9.2:eq-geometric-forgetting-a}. It is the same comparison along the flow by which \cite{B-CMP122b} reaches \cite[(1.89)]{B-CMP122b}.

None of the constants used in Steps 1--5 depends on $k$, on the number of levels, on $a$, on the volume or on the position of $X$.
\end{proof}

\begin{corollary}[Geometric forgetting summed over ages]\label{9.2:cor-forgetting-summed}
Assume the hypotheses of Theorem~\ref{9.2:thm-geometric-forgetting}, and let $q=e^{-c_a}$ be the
ratio appearing there. Let $k\ge 1$, let $X$ be a component of $Z_k$ on which the integral
operation $\mathbf{T}'_k(X)$ is performed at step $k$, let $a\ge 0$ be an integer, let $F$ be a
bounded function of the variables that $\mathbf{T}'_k(X)$ does not integrate, and fix a point of
the parameter domain $\mathfrak{P}$. Then
\begin{equation}\label{9.2:eq-forgetting-summed-1}
\bigl|\mathbf{T}'_{k,\ge a}(X)F\bigr|
\le e\cdot e^{-2(1+\beta_0)^{-1}p_0(g_k)}\,\sup|F|\,\frac{q^{a+1}}{1-q}.
\end{equation}
If moreover $q\le \tfrac12$, then
\begin{equation}\label{9.2:eq-forgetting-summed-2}
\bigl|\mathbf{T}'_{k,\ge a}(X)F\bigr|
\le e\cdot e^{-2(1+\beta_0)^{-1}p_0(g_k)}\,q^{a}\,\sup|F|.
\end{equation}
\end{corollary}

\begin{proof}
The operation $\mathbf{T}'_{k,\ge a}(X)$ is the sum of the parts $\mathbf{T}'_{k,a'}(X)$ of
$\mathbf{T}'_k(X)$ of age $a'\ge a$. For each such $a'$, the data $k$, $X$, $a'$, $F$ and the
chosen point of $\mathfrak{P}$ satisfy the hypotheses of
Theorem~\ref{9.2:thm-geometric-forgetting}. The second inequality of
\eqref{9.2:eq-geometric-forgetting-a}, applied with $a'$ in place of $a$, therefore gives
\begin{equation*}
\bigl|\mathbf{T}'_{k,a'}(X)F\bigr|
\le e\cdot e^{-2(1+\beta_0)^{-1}p_0(g_k)}\,q^{a'+1}\,\sup|F|.
\end{equation*}
The prefactor on the right does not depend on $a'$. Since $0<q<1$, summing over $a'\ge a$ and
using the triangle inequality yields
\begin{align*}
\bigl|\mathbf{T}'_{k,\ge a}(X)F\bigr|
&\le \sum_{a'\ge a}\bigl|\mathbf{T}'_{k,a'}(X)F\bigr|\\
&\le e\cdot e^{-2(1+\beta_0)^{-1}p_0(g_k)}\,\sup|F|\sum_{a'\ge a}q^{a'+1}\\
&= e\cdot e^{-2(1+\beta_0)^{-1}p_0(g_k)}\,\sup|F|\,\frac{q^{a+1}}{1-q}.
\end{align*}
This is \eqref{9.2:eq-forgetting-summed-1}.

For \eqref{9.2:eq-forgetting-summed-2}, write
\begin{equation*}
\frac{q^{a+1}}{1-q}=q^{a}\cdot\frac{q}{1-q}.
\end{equation*}
If $q\le\tfrac12$, then $q\le 1-q$, so $q/(1-q)\le 1$. Inserting this into
\eqref{9.2:eq-forgetting-summed-1} gives \eqref{9.2:eq-forgetting-summed-2}.
\end{proof}

\begin{remark}[Remarks on the proof of geometric forgetting]\label{9.2:rem-forgetting-remarks}
The following observations concern the proof of Theorem~\ref{9.2:thm-geometric-forgetting}
(geometric forgetting of old regions) under the one-sided conditions
\eqref{9.2:eq-forgetting-conditions-a}.
\begin{enumerate}
\item[(i)] No gap between creations is counted, and no window of levels appears in the argument.
\item[(ii)] The surplus that produces the geometric factor is, for a history $h$,
\begin{equation*}
\exp\Bigl(-c_a\sum_{(n,Z)\in\operatorname{anc}(h)}\sqrt{d}\,M^dR_n^{d+1}|Z|\Bigr),
\end{equation*}
and it is at most the corresponding expression with $\sqrt{d}\,|Z|$ replaced by $d'(Z)+1$.
\item[(iii)] The positions of the internal histories cost nothing.
\item[(iv)] For $a=0$ the bound \eqref{9.2:eq-geometric-forgetting-a} holds with the factor $q^1$;
in general the exponent of $q$ is $a+1$ and not $a$.
\item[(v)] No income is counted twice.
\item[(vi)] Quiet live levels of a second-class survivor cause no loss.
\item[(vii)] The surplus survives the mergers and the preparatory case of Ba{\l}aban's construction.
\item[(viii)] Enlarging the charge constant from $C_0$ to $C$ affects only the ceiling on $\gamma$.
\item[(ix)] The order of choices is one-sided: no later choice returns to $L$, $M$ or $c_a$.
\item[(x)] The dial and complex couplings do not involve $c_a$.
\item[(xi)] No constant of \eqref{9.2:eq-geometric-forgetting-a} depends on the number of stages,
of constituents, on the clock, on $k$, on $a$ or on the volume.
\item[(xii)] The number of levels enters no constant of \eqref{9.2:eq-geometric-forgetting-a}.
\item[(xiii)] Of the length $R_n$ only the bounds $1\le R_n<L\,\Xi_n^{r}$, the conditions (C0b),
(C0c) of \eqref{9.2:eq-forgetting-conditions-a} and, by statement, the charge bound of the majorant
are used.
\end{enumerate}
\end{remark}

\begin{proof}
(i) The splitting identity in \eqref{9.2:eq-geometric-forgetting-b}, which writes
$E^{[C_0]}(h)$ as $E^{[C]}(h)$ times the surplus of (ii), is an identity on each single history;
it is applied before the induction of Proposition~\ref{9.2:prop-prepayment-induction} is invoked.
Joins and survivals change the recursion \eqref{9.2:eq-prepayment-numbers-a} of the pre-payment
numbers $\kappa^{[C]}$; that recursion runs at the constant $C$ whatever the history, and these
events never act on the factor split off. The step of Ba{\l}aban's argument in which the
pre-payment number is discarded at a survival \cite{B-CMP122b} acts inside the bound on
$D(h)\,I(h)\,E^{[C]}(h)$, not on the surplus. Hence no gap between creations needs to be counted.

(ii) The surplus is the exponential factor in the splitting identity of
\eqref{9.2:eq-geometric-forgetting-b}, since $C-C_0=c_a$. By the lower bound for the number $|Z|$
of cells of a component $Z$ in terms of its linear size $d'(Z)$, it satisfies
\begin{equation*}
\exp\Bigl(-c_a\!\!\sum_{(n,Z)\in\operatorname{anc}(h)}\!\!\sqrt{d}\,M^dR_n^{d+1}|Z|\Bigr)
\le\exp\Bigl(-c_a\!\!\sum_{(n,Z)\in\operatorname{anc}(h)}\!\!M^dR_n^{d+1}\bigl(d'(Z)+1\bigr)\Bigr),
\end{equation*}
and the right-hand side is smaller than the surplus expression used in the remark on the surplus
of the localised large-field lemma. In Ba{\l}aban's linear size, part (1) of
Lemma~\ref{9.2:lem-linear-size-lower}, a lemma whose proof is incomplete at one step, gives
$d'\ge1$ at every level $n\ge2$; this is not used in the proof of
Theorem~\ref{9.2:thm-geometric-forgetting}.

(iii) The sum over the internal histories is taken inside the supremum bound
of \eqref{9.2:eq-geometric-forgetting-b} on $\sum_{h}F^{\natural}(h)$; the only summation left to a
statement that uses Theorem~\ref{9.2:thm-geometric-forgetting} is over the places of $X$ on the
lattice of level $k$.

(iv) For $a=0$ the history has the single level $k$. The first line of
\eqref{9.2:eq-prepayment-numbers-a}, with its index $m$ equal to $0$ at level $k$ ($m$ the index of
that display, not the torus exponent), subtracts $\mathsf{c}_k(X)$, and
\eqref{9.2:eq-prepayment-induction-b} at $n=k$ gives $\kappa^{[C]}_k(X)\ge0$. The factor split off is
\begin{equation*}
\exp\bigl(-c_a\sqrt{d}\,M^dR_k^{d+1}|X|\bigr)\le e^{-c_a}=q,
\end{equation*}
because $|X|\ge1$, $R_k\ge1$ and $\sqrt{d}\,M^d\ge1$; thus \eqref{9.2:eq-geometric-forgetting-a}
holds with $q^1$. In general the charge of level $k$ belongs to $\operatorname{anc}(h)$, so the
number of terms of the surplus is at least $a+1$, whence the exponent $a+1$ and not $a$.

(v) The income bound is the source of the income in the third line of
\eqref{9.2:eq-prepayment-numbers-a} for survivals of the first kind, and the statement that a
second-class outcome pays a tag is the source of that income for survivals of the second kind.
The income of a survival appears once
in $F^{\natural}(h)$, in the group of its factors that carries the incomes of survivals, and is
used once, in the third line of \eqref{9.2:eq-prepayment-numbers-a}. A survivor joined at the next
step has its income bounded by $1$ in Lemma~\ref{9.2:lem-lineage-unrolling}, and that income is
never used. Repeated survivals of one lineage are distinct events at distinct levels, with incomes
consumed at distinct levels of the recursion \eqref{9.2:eq-prepayment-numbers-a}. A survivor is
never also counted as a creation. The surplus is split off from the charges $E$, not from the
incomes. The income statement with the per-cell charge constant $C_{\mathfrak{d}}^{\flat}$ is not
used in the proof, so no income is imported through it.

(vi) At the quiet live levels of a second-class survivor the linear size $d'$ of
\cite{B-CMP122b} may vanish. The floor used in the proof is $|Z|\ge1$, which holds for every
non-empty component whatever its shape. Independently, the vanishing at those places concerns the
covers inside the clock estimate, the subject of part (2) of Lemma~\ref{9.2:lem-linear-size-lower},
whose proof is incomplete at one step; and by part (1) of the same lemma a survivor and its
descendants contain $c^{\sim1}$ for some cell $c$ and have $d'\ge1$.

(vii) By (i): the first and third lines of \eqref{9.2:eq-prepayment-numbers-a} subtract the full
charge at the constant $C$, and Proposition~\ref{9.2:prop-prepayment-induction} absorbs it at every
event, including the mergers \cite[(1.84)--(1.88)]{B-CMP122b} and the preparatory case
\cite{B-CMP122b}.

(viii) The statement that a second-class outcome pays a tag (with $E_{\mathrm{2cl}}$ and $N_Y$),
the deposits $\mathsf{Q}_n$, $\mathsf{D}_n$ of the deposit bound, the income bound, the counting
constant $c_A^{\flat}(n)$, the core and clock lemmas used in
Proposition~\ref{9.2:prop-prepayment-induction}, and the lines (S), (P), $(F_I)$ and
(C0a)--(C0c) of \eqref{9.2:eq-forgetting-conditions-a} are all free of $C$. Only the lines
$(B)_C$ and $(M)_C$ of \eqref{9.2:eq-forgetting-conditions-a} (see
\eqref{9.2:eq-prepayment-induction-a}) depend on $C$, and monotonically: their right-hand sides
increase with $C$, so a larger $C$ asks for a smaller $\gamma$ and for nothing else, and the
condition stays one-sided.

(ix) The exponents ($p_0$, $r$ and the exponents of the correlation theorem) are chosen first;
then $L$; then $M_1$, $M_2$, $M$; then the constants chosen after $(L,M)$ ($\kappa$,
$\Theta_1$, $C_\sharp$, $C_{L'}$, and the others) and the charge constant $C_0$; then $c_a$, free
subject to $c_a\ge\log2$ (the statements that use the localised large-field lemma fix it after
$L$, for instance by $q\le\vartheta^2$); then
$C=C_0+c_a$; then $\gamma$ by \eqref{9.2:eq-forgetting-conditions-a}. No later choice returns to
$L$, $M$ or $c_a$.

(x) All factors of $F^{\natural}$ are majorants of moduli: \eqref{9.2:eq-modulus-majorant-a} comes
from the proof of the modulus lemma and from the complex-coupling proposition. The quantities
$\mathsf{D}_n$, $\mathsf{P}_n$, $E_{\mathrm{2cl}}(n)$, $N_Y$ are numbers at the real couplings. By
the line (SL/SL$^2$) of \eqref{9.2:eq-forgetting-conditions-a}, the raw Wilson coefficient keeps,
on $\mathbb{D}_\delta$, a real part at least $\tfrac34$ of its value; this is what the deposit bound
and the Wilson clause of the statement that a second-class outcome pays a tag use. Positivity of
the raw Wilson term for group-valued bond variables is used as stated in \cite{B-CMP122b}. The
per-cell loads of the dial are contained in $C_0$. None of these involves $c_a$.

(xi) The core and clock lemmas are uniform in the number of stages and of constituents. The bound
\eqref{9.2:eq-prepayment-induction-c} is linear in $N-1$, with $N$ the integer of that display (not
the $N$ of $\mathrm{SU}(N)$), and is paid by the $N-1$ surviving linear deposits. The clock
$K_n(Z)$ is a local integer bounded by part (ii)
of the clock lemma and occurs in no constant of \eqref{9.2:eq-geometric-forgetting-a}; neither do
$k$, $a$ or the volume.

(xii) In the statement that a second-class outcome pays a tag, the number of levels could enter
only through the coupling at which the Wilson energy is converted into income; its Wilson clause
takes that coupling
frozen at the processing level, so it does not enter there. In the inner prefactors, one per
second-class survival, it could enter, but these are paid survival by survival by the line (P) of
\eqref{9.2:eq-forgetting-conditions-a} and Lemma~\ref{9.2:lem-comparison-product}. It enters
nowhere else.

(xiii) The letter $R_n$ denotes the length $\check R_n$ throughout, not the ball radius $R$ of the
observables. Of it only $1\le R_n<L\,\Xi_n^{r}$ and the conditions (C0b), (C0c) of
\eqref{9.2:eq-forgetting-conditions-a} are used, together with the charge bound of the majorant,
by statement.
\end{proof}

\begin{lemma}[A difference bound for bounded holomorphic functions on a disc]
\label{9.2:lem-holomorphic-disc-bound}
Let $\varrho>1$, write $D(0,\varrho)=\{w\in\mathbb{C}:|w|<\varrho\}$ for the open disc of radius
$\varrho$ centred at $0$, and let $f$ be holomorphic on $D(0,\varrho)$ with $|f|\le\mathfrak{B}$
there, where $\mathfrak{B}\ge 0$. Then
\begin{equation}\label{9.2:eq-holomorphic-disc-bound-1}
|f(1)-f(0)|\le\frac{2\mathfrak{B}}{\varrho},
\end{equation}
and if $f(0)=0$, then
\begin{equation}\label{9.2:eq-holomorphic-disc-bound-2}
|f(1)|\le\frac{\mathfrak{B}}{\varrho}.
\end{equation}
\end{lemma}

\begin{proof}
Suppose first that $f(0)=0$. If $\mathfrak{B}=0$, then $f$ vanishes identically on $D(0,\varrho)$ and
\eqref{9.2:eq-holomorphic-disc-bound-2} holds. If $\mathfrak{B}>0$, consider the map
$z\mapsto f(\varrho z)/\mathfrak{B}$ on the open unit disc. Since $z\mapsto\varrho z$ maps the open
unit disc into $D(0,\varrho)$, this map is holomorphic on the open unit disc; its modulus is at most
$1$ there, because $|f|\le\mathfrak{B}$ on $D(0,\varrho)$, and its value at $0$ is
$f(0)/\mathfrak{B}=0$. Schwarz's lemma therefore gives $|f(\varrho z)|/\mathfrak{B}\le|z|$ for
$|z|<1$. For $|w|<\varrho$ we take $z=w/\varrho$, which lies in the open unit disc, and obtain
\begin{equation*}
|f(w)|=\mathfrak{B}\,\frac{|f(\varrho z)|}{\mathfrak{B}}\le\mathfrak{B}\,|z|
=\frac{\mathfrak{B}\,|w|}{\varrho}.
\end{equation*}
Since $\varrho>1$, the point $w=1$ belongs to $D(0,\varrho)$, and this inequality at $w=1$ is
\eqref{9.2:eq-holomorphic-disc-bound-2}.

In general, the function $f-f(0)$ is holomorphic on $D(0,\varrho)$ and vanishes at $0$. For
$|w|<\varrho$ we have
\begin{equation*}
|f(w)-f(0)|\le|f(w)|+|f(0)|\le 2\mathfrak{B},
\end{equation*}
since $0$ and $w$ both lie in $D(0,\varrho)$. The case already treated, applied to $f-f(0)$ with the
bound $2\mathfrak{B}$ in place of $\mathfrak{B}$, gives $|f(1)-f(0)|\le 2\mathfrak{B}/\varrho$; this is
\eqref{9.2:eq-holomorphic-disc-bound-1}.
\end{proof}

In the applications of Lemma~\ref{9.2:lem-holomorphic-disc-bound} in this chapter the radius is
$\varrho=q^{-a}$, where $q=e^{-c_a}$ is the geometric forgetting ratio and $a$ is an age; it
satisfies $\varrho\ge 2$ when $a\ge 1$, while for $a=0$ the bound needed is trivial, so the lemma in
the form stated, for the open disc, suffices.

\begin{corollary}[Geometric smallness of the old-age part of the final integral]
\label{9.2:cor-old-age-smallness}
We work at the last level $K'$ of the correlation theorem; the lattice of that level has $N^4$
sites. In this corollary and its proof $N$ denotes the side length of that lattice, not the rank
of the gauge group. Let $\tilde Z$ be the final quantity of the correlation theorem, regarded as a
function on the dial domain $\mathbb{D}$, and let $E_+^{\sharp}$ be the constant in the exponent
of the correlation theorem's bound of its final quantity, \eqref{9.2:eq-old-age-smallness-1}
below.
For a live component $X$ of level $K'$ let $\mathbf{T}_{K'}(X)$ be the operation of $X$ in
$\tilde Z$. Histories $h\in\operatorname{Hist}(X)$ and their ages $a(h)$ are taken with $k=K'$.
We write $\mathbf{T}_{K',a}(X)$ for the part of $\mathbf{T}_{K'}(X)$ belonging to histories of
age $a$, and $\mathbf{T}_{K',\ge a}(X)$ for the sum of these parts over all ages at least $a$.
Put $q=e^{-c_a}$. The antecedents are the following.
\begin{itemize}
\item[(A)] Geometric forgetting of old regions, Theorem~\ref{9.2:thm-geometric-forgetting},
holds at the last level. Precisely, the history-wise majorant of the operations
$\mathbf{T}_{K'}(X)$ (in the sense of Proposition~\ref{9.2:prop-modulus-majorant}, whose proof
is outstanding) is formed inside the bound \eqref{9.2:eq-old-age-smallness-1} of
antecedent~(B) below. The deposits, charges
and incomes are the same objects as in Theorem~\ref{9.2:thm-geometric-forgetting}, and the
second-class incomes are paid at the second of the two loci of the statement that a
second-class outcome pays a tag, namely at the last level, from the Wilson term of the final
action.
Consequently \eqref{9.2:eq-geometric-forgetting-a} holds for the parts $\mathbf{T}_{K',a}(X)$.
\item[(B)] The correlation theorem's bound of its final quantity on the dial domain,
\begin{equation}\label{9.2:eq-old-age-smallness-1}
|\tilde Z|\le e^{E_+^{\sharp}N^4}\qquad\text{on }\mathbb{D},
\end{equation}
holds together with its structure. In that bound each live component's operation enters
through the supremum over its histories of
\begin{equation}\label{9.2:eq-old-age-smallness-2}
e\cdot e^{-2(1+\beta_0)^{-1}p_0(g_{K'})}
\end{equation}
times the history's comparison factors (the factors of $F^{\natural}(h)$ in
Proposition~\ref{9.2:prop-modulus-majorant}(b), whose proof is outstanding), where $p_0(\cdot)$
is the degree function. Part of this antecedent, together with
\eqref{9.2:eq-old-age-smallness-1}, is the persistence under the dial of the bounds stated in
\cite[(1.92)--(1.100)]{B-CMP122b}.
\item[(C)] The difference bound for bounded holomorphic functions on a disc,
Lemma~\ref{9.2:lem-holomorphic-disc-bound}.
\end{itemize}
Let $a\ge 0$ be an integer. Write $\tilde Z=\tilde Z_{<a}+\tilde Z_{\ge a}$, where
$\tilde Z_{\ge a}$ collects the terms whose label contains a live component with a history of
age at least $a$. Then, on $\mathbb{D}$,
\begin{equation}\label{9.2:eq-old-age-smallness-3}
|\tilde Z_{\ge a}|\le 2q^{a}\,e^{E_+^{\sharp}N^4}.
\end{equation}
The factor $e^{E_+^{\sharp}N^4}$ is the bound \eqref{9.2:eq-old-age-smallness-1} of $|\tilde Z|$;
it depends on the volume and is kept explicit in \eqref{9.2:eq-old-age-smallness-3}.
\end{corollary}

\begin{proof}
Fix a point of $\mathbb{D}$. All bounds below hold at that point with constants independent of
it.

\emph{The case $a=0$.} Here $q^a=1$, and \eqref{9.2:eq-old-age-smallness-3} reads
$|\tilde Z_{\ge 0}|\le 2e^{E_+^{\sharp}N^4}$. This is the bound
\eqref{9.2:eq-old-age-smallness-1} of antecedent~(B), multiplied by two.

\emph{The case $a\ge 1$: the dial.} In the operation of every live component $X$ of level $K'$,
replace the part $\mathbf{T}_{K',\ge a}(X)$ by $\mu\,\mathbf{T}_{K',\ge a}(X)$. Here $\mu$ is a
complex number with $|\mu|<q^{-a}$. Leave the parts $\mathbf{T}_{K',a'}(X)$ with $a'<a$
unchanged. Call the resulting quantity $\tilde Z(\mu)$.

The lattice of level $K'$ is finite, so only finitely many live components occur. Hence
$\tilde Z(\mu)$ is a polynomial in $\mu$.

For $\mu=1$ nothing is changed, so $\tilde Z(1)=\tilde Z$. For $\mu=0$ every term whose label
contains a live component with a history of age at least $a$ vanishes. The remaining terms are
exactly those of $\tilde Z_{<a}$, so $\tilde Z(0)=\tilde Z_{<a}$.

\emph{The case $a\ge 1$: the bound persists under the dial.} By antecedent~(B), in the majorant
behind \eqref{9.2:eq-old-age-smallness-1} each live component's operation enters through the
supremum over its histories of the quantity \eqref{9.2:eq-old-age-smallness-2} times the
history's comparison factors.

After the replacement, consider a history $h$ of age $a(h)=a'\ge a$. By antecedent~(A), the
history-wise majorant is formed inside this bound and \eqref{9.2:eq-geometric-forgetting-a}
holds for the part of age $a'$. So $h$ carries in that supremum the quantity
$|\mu|\,e\,e^{-2p_0(g_{K'-a'})}q^{a'+1}$. By \eqref{9.2:eq-geometric-forgetting-a} and
$|\mu|<q^{-a}$,
\begin{align*}
|\mu|\cdot e\cdot e^{-2p_0(g_{K'-a'})}\,q^{a'+1}
&\le e\cdot e^{-2(1+\beta_0)^{-1}p_0(g_{K'})}\cdot q^{-a}q^{a'+1}\\
&\le e\cdot e^{-2(1+\beta_0)^{-1}p_0(g_{K'})}\cdot q .
\end{align*}
The second inequality holds because $q^{-a}q^{a'+1}=q^{a'-a}\,q$, with $a'-a\ge 0$ and $q\le 1$.

The right-hand side lies within what the bound \eqref{9.2:eq-old-age-smallness-1} uses for a
history in that supremum. Histories of age $a'<a$ are untouched by the replacement. Therefore
the majorant behind \eqref{9.2:eq-old-age-smallness-1} still applies, and
\begin{equation}\label{9.2:eq-old-age-smallness-4}
|\tilde Z(\mu)|\le e^{E_+^{\sharp}N^4}\qquad\text{for all }\mu\in\mathbb{C},\ |\mu|<q^{-a}.
\end{equation}

\emph{The case $a\ge 1$: the disc.} Define $f(\mu):=\tilde Z(\mu)-\tilde Z(0)$. As a polynomial,
$f$ is holomorphic on the open disc of radius $\varrho:=q^{-a}$ about the origin, and $f(0)=0$.

Under the conditions of Theorem~\ref{9.2:thm-geometric-forgetting}, which include a floor on
$c_a$, one has $q\le\tfrac12$. Since moreover $a\ge1$, we have $\varrho=q^{-a}\ge 2>1$.
By \eqref{9.2:eq-old-age-smallness-4}, applied
at $\mu$ and at $0$, we have $|f|\le\mathfrak{B}:=2e^{E_+^{\sharp}N^4}$ on that disc.

Lemma~\ref{9.2:lem-holomorphic-disc-bound}, in its case $f(0)=0$, gives
$|f(1)|\le\mathfrak{B}/\varrho$. Since $f(1)=\tilde Z(1)-\tilde Z(0)=\tilde Z-\tilde Z_{<a}$, we get
\begin{equation*}
|\tilde Z_{\ge a}|=|f(1)|\le\frac{2e^{E_+^{\sharp}N^4}}{q^{-a}}=2q^{a}\,e^{E_+^{\sharp}N^4},
\end{equation*}
which is \eqref{9.2:eq-old-age-smallness-3}.
\end{proof}

\begin{proposition}[Truncation of the remainder to histories of recent first index]
\label{9.2:prop-remainder-truncation}
Work in the setting of the geometric forgetting theorem (Theorem~\ref{9.2:thm-geometric-forgetting})
and of the persistence of Ba{\l}aban's bounds \cite[(1.92)--(1.100)]{B-CMP122b} under the dial
(multiplication of the terms of a family of histories by a complex parameter $\mu$; see
hypothesis (b) below), for every choice of the data admitted there. Let $n<k$ be levels, and put
\begin{equation*}
x_m:=g_m^{-2},\qquad x^{\sharp}_n:=x_n-2c_2,\qquad
p_0^{\sharp}(n):=A_0\bigl(\log x^{\sharp}_n\bigr)^{p_0}.
\end{equation*}
Let $q=e^{-c_a}$ be the ratio of Theorem~\ref{9.2:thm-geometric-forgetting}, and let $a(h)$ and $j(h)$
denote the age and the first index of an admissible history $h$. Let $R'^{(k)}$ be the term of
\cite[(1.98)]{B-CMP122b}, and let $R'^{(k)}_{>n}$ be the part of $R'^{(k)}$ built from the histories of
first index $j(h)>n$. Assume:
\begin{itemize}
\item[(a)] the first inequality of \eqref{9.2:eq-geometric-forgetting-a}, the conclusion of
Theorem~\ref{9.2:thm-geometric-forgetting};
\item[(b)] the persistence of the bounds \cite[(1.92)--(1.100)]{B-CMP122b} under the dial: if, for some
radius $\bar\mu>1$, the terms of the histories of a family are multiplied by a complex variable $\mu$
with $|\mu|<\bar\mu$, and if for $|\mu|<\bar\mu$ the modified majorants still satisfy the requirement
of \cite[(1.92)]{B-CMP122b}, then \cite[(1.92)--(1.100)]{B-CMP122b} hold for the modified quantities on
$|\mu|<\bar\mu$, and these quantities depend holomorphically on $\mu$ there;
\item[(c)] Lemma~\ref{9.2:lem-holomorphic-disc-bound};
\item[(d)] the flow fact that every running coupling satisfies $x_i\ge x_n-2c_2=x^{\sharp}_n$ for
$i\le n$;
\item[(e)] the degree comparison between levels: for $k>n$ one has $x_k\le x_n+2c_2$, and then
\begin{equation*}
p_0(g_k)\le(1+\beta_0)\,p_0^{\sharp}(n).
\end{equation*}
\end{itemize}
Then $R'^{(k)}_{>n}$ contains no use of a zone of level $<n$, and
\begin{equation}\label{9.2:eq-remainder-truncation-1}
\bigl\|R'^{(k)}-R'^{(k)}_{>n}\bigr\|^{\flat}\le 2\,e^{-(1-\beta_0)p_0^{\sharp}(n)}\,q^{k-n}.
\end{equation}
\end{proposition}

\begin{proof}
\emph{No zone below $n$.} A use of a zone of level $<n$ occurs only in histories whose first index
is $\le n$. The part $R'^{(k)}_{>n}$ is built from histories of first index $>n$ only, and therefore
contains no such use.

\emph{The majorant of the discarded histories.} The term $R'^{(k)}$ is built from the integral
operations $\mathbf{T}'_k(X)$ on the components processed at step $k$, so the first inequality of
\eqref{9.2:eq-geometric-forgetting-a} (hypothesis (a)) applies to each history entering it. Let $h$ be
a history of first index $j(h)\le n$. Its age at level $k$ satisfies $a(h)\ge k-n$. By hypothesis (d)
with $i=j(h)\le n$ we have $x_{j(h)}\ge x_n-2c_2=x^{\sharp}_n$. Since the degree function
$p_0(g')=A_0(\log g'^{-2})^{p_0}$ is increasing in $g'^{-2}$, this gives
$p_0(g_{j(h)})\ge p_0^{\sharp}(n)$. Hence the first inequality of \eqref{9.2:eq-geometric-forgetting-a},
combined with this lower bound, shows that the majorant of $h$ carries the factor
\begin{equation}\label{9.2:eq-remainder-truncation-2}
e\cdot e^{-2p_0^{\sharp}(n)}\,q^{a(h)+1}.
\end{equation}
For comparison, \cite[(1.92)]{B-CMP122b} requires only the factor
$e\cdot e^{-2(1+\beta_0)^{-1}p_0(g_k)}$.

\emph{The dial.} Multiply the terms of all histories of first index $\le n$ by a complex number $\mu$,
leaving the other terms unchanged, and write $R'^{(k)}(X;\mu)$ for the resulting quantity on a
component $X$. By construction
\begin{equation*}
R'^{(k)}(X;1)=R'^{(k)}(X),\qquad R'^{(k)}(X;0)=R'^{(k)}_{>n}(X).
\end{equation*}
Put
\begin{equation*}
\bar\mu:=q^{-(k-n)}\exp\Bigl(2p_0^{\sharp}(n)-2(1+\beta_0)^{-1}p_0(g_k)\Bigr).
\end{equation*}
By hypothesis (e), $p_0^{\sharp}(n)\ge(1+\beta_0)^{-1}p_0(g_k)$, so the exponential factor is at least
$1$. Since $k-n\ge1$, we have $q^{-(k-n)}\ge q^{-1}$, and $q^{-1}=e^{c_a}\ge2$ under the floor on $c_a$
among the conditions of Theorem~\ref{9.2:thm-geometric-forgetting}. Hence $\bar\mu\ge q^{-1}\ge 2$.

Let $|\mu|<\bar\mu$ and let $h$ be a history of first index $\le n$. Using
\eqref{9.2:eq-remainder-truncation-2}, the definition of $\bar\mu$, the bound $a(h)\ge k-n$ and
$q\le 1$, we obtain
\begin{align*}
|\mu|\cdot e\cdot e^{-2p_0^{\sharp}(n)}\,q^{a(h)+1}
&\le e\cdot e^{-2(1+\beta_0)^{-1}p_0(g_k)}\,q^{a(h)+1-(k-n)}\\
&\le e\cdot e^{-2(1+\beta_0)^{-1}p_0(g_k)}\cdot q .
\end{align*}
Since $q\le1$, the right-hand side is at most the factor required by \cite[(1.92)]{B-CMP122b}. The
histories of first index $>n$ are unchanged and meet that requirement as before. Hypothesis (b) thus
applies: the bounds \cite[(1.92)--(1.100)]{B-CMP122b} persist under the dial on $|\mu|<\bar\mu$, and
they give
\begin{equation}\label{9.2:eq-remainder-truncation-3}
\bigl|R'^{(k)}(X;\mu)\bigr|\le e^{-p_0(g_k)}\,e^{-\kappa d_k(X)}\qquad(|\mu|<\bar\mu).
\end{equation}

\emph{Application of the disc bound.} The dial variable enters $R'^{(k)}(X;\mu)$ holomorphically on the
open disc of radius $\bar\mu$, by the holomorphy clause of hypothesis (b). Apply
Lemma~\ref{9.2:lem-holomorphic-disc-bound} (hypothesis (c)) to $\mu\mapsto R'^{(k)}(X;\mu)$ with
$\varrho=\bar\mu$, which is admissible since $\bar\mu\ge2>1$. Take as bound $\mathfrak{B}$ the right-hand
side of \eqref{9.2:eq-remainder-truncation-3}. By the two identities displayed above, the value at
$\mu=1$ minus the value at $\mu=0$ is $R'^{(k)}(X)-R'^{(k)}_{>n}(X)$. The lemma gives
\begin{equation*}
\bigl|R'^{(k)}(X)-R'^{(k)}_{>n}(X)\bigr|\le \frac{2\,e^{-p_0(g_k)}}{\bar\mu}\,e^{-\kappa d_k(X)} .
\end{equation*}
This holds for every component $X$, so, by the definition of the flat norm $\|\cdot\|^{\flat}$,
\begin{equation}\label{9.2:eq-remainder-truncation-4}
\begin{split}
\bigl\|R'^{(k)}-R'^{(k)}_{>n}\bigr\|^{\flat}
&\le\frac{2\,e^{-p_0(g_k)}}{\bar\mu}\\
&=2\,q^{k-n}\exp\Bigl(-p_0(g_k)+2(1+\beta_0)^{-1}p_0(g_k)-2p_0^{\sharp}(n)\Bigr).
\end{split}
\end{equation}

\emph{The exponent.} Since $0<(1+\beta_0)^{-1}\le1$ and $p_0(g_k)\ge0$, we have
$2(1+\beta_0)^{-1}p_0(g_k)\le 2p_0(g_k)$. By hypothesis (e), which for $k>n$ rests on
$x_k\le x_n+2c_2$, we have $p_0(g_k)\le(1+\beta_0)p_0^{\sharp}(n)$. Together these give
\begin{align*}
-p_0(g_k)+2(1+\beta_0)^{-1}p_0(g_k)-2p_0^{\sharp}(n)
&\le p_0(g_k)-2p_0^{\sharp}(n)\\
&\le (1+\beta_0)\,p_0^{\sharp}(n)-2p_0^{\sharp}(n)
=-(1-\beta_0)\,p_0^{\sharp}(n).
\end{align*}
Inserting this into \eqref{9.2:eq-remainder-truncation-4} yields \eqref{9.2:eq-remainder-truncation-1}.
\end{proof}

\begin{definition}[Paired variational problems at consecutive cutoffs]\label{9.3:not-paired-variational}
We use Ba{\l}aban's variational problem in the form of \cite{B-CMP102b}, with $d=4$, and pair two
instances of it at the lattice spacings $\eta$ and $\eta/L$.

\emph{The variational problem at spacing $\eta$.} Let $\eta=L^{-k}$ and let $T_{\eta}$ be the torus
viewed as a lattice of spacing $\eta$. A decreasing sequence of domains
$\boldsymbol{\Omega}=(\Omega_0\supset\Omega_1\supset\dots\supset\Omega_k)$ in $T_{\eta}$ is
\emph{admissible} if each $\Omega_j$ is a union of blocks of side $M_1L^j\eta$ and
\begin{equation*}
(L^j\eta)^{-1}\operatorname{dist}\bigl(\Omega_j^{c},\Omega_{j+1}\bigr)>RM_1,
\end{equation*}
as in \cite[(1)]{B-CMP102b}; domains equal to $T_{\eta}$ are admitted, as in \cite{B-CMP99a}.
Here $R$ and $M_1$ are the integers of \cite{B-CMP102b}; this $R$ is the integer of
\cite[(1)]{B-CMP102b}, not the ball radius $R$ of the observables. We use the convention
$\Omega_{k+1}:=\emptyset$. Put
\begin{equation*}
\Lambda_j:=(\Omega_j\setminus\Omega_{j+1})^{(j)},\qquad
\mathfrak{B}:=\bigcup_{j=0}^{k}\Lambda_j,
\end{equation*}
the layers $\Lambda_j$ and the determining set $\mathfrak{B}$. The superscript $(j)$, as in
\cite{B-CMP102b}, indicates that $\Lambda_j$ is taken on the lattice of spacing $L^j\eta$, on which
the $j$-step averages $M^{j}(U)$ below are defined; we call $\Lambda_j$ the layer of level $j$.
The data are configurations $V$ on $\mathfrak{B}$ subject to the condition of
\cite[(7)]{B-CMP102b}, written $(7)_{\varepsilon}$:
\begin{equation}\label{9.3:eq-paired-variational-a}
|(\partial V)(p')-1|<\varepsilon,\qquad p'\in\mathfrak{B}.
\end{equation}
With $M^{j}$ Ba{\l}aban's $j$-step averaging map of \cite{B-CMP98}, $U_{\mathfrak{B}}(V)$ denotes a
minimal configuration of
\begin{equation*}
A^{\eta}(U)=\sum_{p\subset\Omega_0}\eta^{d-4}\bigl[1-\operatorname{Re}\operatorname{tr}U(\partial p)\bigr]
\end{equation*}
on the set $\{U:\ M^{j}(U)=V\text{ on }\Lambda_j,\ 0\le j\le k\}$.

\emph{Levels and norms.} The \emph{level} of a point $x\in\Omega_0$ is
$j(x):=\max\{j:\ x\in\Omega_j\}$, and $\ell:=L^{j(x)}\eta$ is the local scale at $x$. For a
configuration $U$ we put, as in \cite[(27)--(28)]{B-CMP102b},
$\mathfrak{F}_U(p):=\eta^{-2}\operatorname{Im}U(\partial p)$ and
$J_U:=D_U^{*}\mathfrak{F}_U$, where $D_U$ is the covariant derivative taking functions $A$ on bonds
to functions on plaquettes and $D_U^{*}$ is its adjoint, so that
\begin{equation*}
\langle A,J_U\rangle=\langle D_UA,\mathfrak{F}_U\rangle .
\end{equation*}
For a real $\gamma$, the level-weighted norm of \cite[(3.41)]{B-CMP99b} is
\begin{equation*}
|A|_{(\gamma)}:=\sup_{j}\ \sup_{\Omega_j\setminus\Omega_{j+1}}(L^j\eta)^{-\gamma}|A|,
\end{equation*}
and for a real parameter $\theta$, with $\nabla_UA$ the covariant lattice derivatives of $A$ (an
operator distinct from $D_U$, whose weighted size enters the following norm),
\begin{equation}\label{9.3:eq-paired-variational-b}
\|A\|_{\theta}:=\max\bigl\{|A|_{(-1-\theta)},\ |\nabla_UA|_{(-2-\theta)}\bigr\}.
\end{equation}
Thus, for $\tau>0$, $\|A\|_{\theta}\le\tau\eta^{\theta}$ holds if and only if at every level $j$
\begin{equation}\label{9.3:eq-paired-variational-1}
\ell|A|\le\tau L^{-\theta j}\quad\text{and}\quad \ell^{2}|\nabla_UA|\le\tau L^{-\theta j}.
\end{equation}
The distance $d(y,y')$ between lattice points is that of \cite[(2.1)--(2.4)]{B-CMP96}.

\emph{A lattice beneath.} Put $\eta':=\eta/L$. Let
$\boldsymbol{\Omega}'=(\Omega'_0,\dots,\Omega'_{k+1})$ be an admissible sequence in $T_{\eta'}$, with
layers $\Lambda'_j$, determining set $\mathfrak{B}'$ and data $V'$ on $\mathfrak{B}'$, such that
\begin{equation}\label{9.3:eq-paired-variational-c}
\begin{gathered}
\Omega'_{j+1}=\Omega_j,\qquad V'=V\ \text{on}\ \Lambda'_{j+1}=\Lambda_j\qquad(0\le j\le k),\\
U':=U_{\mathfrak{B}'}(V'),\qquad W:=M(U')\ \text{on}\ \Omega_0 ;
\end{gathered}
\end{equation}
the domain $\Omega'_0$ and the data $V'$ on $\Lambda'_0$ are arbitrary, subject to admissibility of
$\boldsymbol{\Omega}'$ and to \eqref{9.3:eq-paired-variational-a} for $V'$ on $\mathfrak{B}'$, and
$\Lambda'_0$ is empty when $\Omega_0$ is the whole torus. Then
\begin{equation}\label{9.3:eq-paired-variational-2}
M^{j}(W)=V\quad\text{on }\Lambda_j,\qquad 0\le j\le k.
\end{equation}

\emph{The floor.} On the integers $R$ and $M_1$, which are chosen after $L$, we impose, with
$\delta_0$ the exponential decay rate of the estimates of \cite{B-CMP102b}, the one-sided
condition, called the floor $(\mathrm{F})$,
\begin{equation}\label{9.3:eq-paired-variational-d}
\delta_0RM_1\ge 16\log L .
\end{equation}
\end{definition}

\begin{proof}
The identity $\langle A,J_U\rangle=\langle D_UA,\mathfrak{F}_U\rangle$ is the definition of the
adjoint $D_U^{*}$ applied to $J_U=D_U^{*}\mathfrak{F}_U$.

For \eqref{9.3:eq-paired-variational-1}: since the domains decrease and $\Omega_{k+1}=\emptyset$,
the sets $\Omega_j\setminus\Omega_{j+1}$, $0\le j\le k$, partition $\Omega_0$, and a point $x$ lies
in $\Omega_j\setminus\Omega_{j+1}$ exactly when $j(x)=j$, in which case $L^j\eta=\ell$. Hence
$|A|_{(-1-\theta)}\le\tau\eta^{\theta}$ means $\ell^{1+\theta}|A|\le\tau\eta^{\theta}$ at every
point of every level $j$. Since $\ell^{\theta}=L^{\theta j}\eta^{\theta}$ there, this reads
$\ell|A|\cdot L^{\theta j}\eta^{\theta}\le\tau\eta^{\theta}$, that is,
$\ell|A|\le\tau L^{-\theta j}$. In the same way $|\nabla_UA|_{(-2-\theta)}\le\tau\eta^{\theta}$
means $\ell^{2+\theta}|\nabla_UA|\le\tau\eta^{\theta}$, that is,
$\ell^{2}|\nabla_UA|\le\tau L^{-\theta j}$. A maximum of two quantities is at most
$\tau\eta^{\theta}$ if and only if each of them is; this gives the equivalence.

The equality $\Lambda'_{j+1}=\Lambda_j$ in \eqref{9.3:eq-paired-variational-c} holds because
$\Omega'_{j+1}\setminus\Omega'_{j+2}=\Omega_j\setminus\Omega_{j+1}$ for $0\le j\le k$ (with
$\Omega'_{k+2}=\emptyset=\Omega_{k+1}$), and the layer $\Lambda'_{j+1}$ is taken at the spacing
$L^{j+1}\eta'=L^{j}\eta$ at which $\Lambda_j$ is taken.

For \eqref{9.3:eq-paired-variational-2}: by \cite[(43)]{B-CMP98}, $M^{j+1}=M^{j}\circ M$. On
$\Lambda_j=\Lambda'_{j+1}$ we therefore have
\begin{equation*}
M^{j}(W)=M^{j}\bigl(M(U')\bigr)=M^{j+1}(U')=V'=V,
\end{equation*}
where the third equality is the constraint $M^{j+1}(U')=V'$ on $\Lambda'_{j+1}$ satisfied by
$U'=U_{\mathfrak{B}'}(V')$, and the last is the pairing $V'=V$ on $\Lambda'_{j+1}$ of
\eqref{9.3:eq-paired-variational-c}.
\end{proof}

\begin{definition}[Bounds on the minimisers at both depths, {\cite[Theorem~1, (9), (10)]{B-CMP102b}}]
\label{9.3:not-minimiser-bounds}
We work with the paired variational problems of Notation~\ref{9.3:not-paired-variational}.
On the coarser side, $\eta = L^{-k}$ and $\boldsymbol{\Omega}=(\Omega_0\supset\dots\supset\Omega_k)$ is admissible
in $T_{\eta}$, with determining set $\mathfrak{B}$ and data $V$ subject to
\eqref{9.3:eq-paired-variational-a}.
On the finer side, $\eta'=\eta/L$ and $\boldsymbol{\Omega}'=(\Omega'_0,\dots,\Omega'_{k+1})$ is admissible in
$T_{\eta'}$, with determining set $\mathfrak{B}'$ and data $V'$ as in
\eqref{9.3:eq-paired-variational-c}.
We write $U_{\mathfrak{B}}=U_{\mathfrak{B}}(V)$ and $U'=U_{\mathfrak{B}'}(V')$ for the two minimisers.

\emph{Constants and exponents.} Throughout this section, $C$ and $C_L$ denote constants depending only
on $d$, $L$ and $\theta$. Here $\theta$ is the exponent of the norm $\|\cdot\|_{\theta}$ of
Notation~\ref{9.3:not-paired-variational}. We put
\begin{equation}\label{9.3:eq-minimiser-bounds-1}
\beta := \frac{1+\theta}{2}, \qquad \epsilon := \frac{1-\theta}{2}, \qquad\text{so that}\qquad
\theta = \beta-\epsilon .
\end{equation}
The exponents $\beta$ and $\epsilon$ are local to this section: $\beta$ is neither the coupling
increment $\beta_{k+1}$ nor a one-loop coefficient $\beta^{0}_{l}$, and the H\"older exponent
$\epsilon$ is distinct from the parameter $\varepsilon$ of \eqref{9.3:eq-paired-variational-a}.
In the bounds below, $R_1M_1$ is the product that multiplies $\varepsilon$ in \cite[(9)]{B-CMP102b},
$R_1$ being the integer so denoted there.

\emph{Correspondence of levels.} Since $\Omega'_{j+1}=\Omega_j$, level $j$ of $\boldsymbol{\Omega}$ and level $j+1$
of $\boldsymbol{\Omega}'$ carry the same local scale
\begin{equation*}
\ell = L^{j}\eta = L^{j+1}\eta' .
\end{equation*}

\emph{The bounds.} By \cite[Theorem~1]{B-CMP102b}, applied on $T_{\eta}$ and on $T_{\eta'}$, the two minimisers
lie in Ba{\l}aban's classes of regular configurations:
\begin{equation*}
U_{\mathfrak{B}}\in\mathfrak{U}_k(\boldsymbol{\Omega},B_3\varepsilon), \qquad
U'\in\mathfrak{U}_{k+1}(\boldsymbol{\Omega}',B_3\varepsilon),
\end{equation*}
where $B_3$ is the constant of \cite[(9)]{B-CMP102b}.

For $U'$ this means the following. At every level of $\boldsymbol{\Omega}'$, $\ell$ being the local
scale of that level,
\begin{equation}\label{9.3:eq-minimiser-bounds-2}
|\mathfrak{F}_{U'}| \le B_3\varepsilon\,\ell^{-2}, \qquad
|J_{U'}| \le 2B_3\varepsilon\,\ell^{-3} .
\end{equation}
At level $j+1$ of $\boldsymbol{\Omega}'$, that is, at level $j$ of $\boldsymbol{\Omega}$, every plaquette
$p'$ of $T_{\eta'}$ at that level satisfies, with $\alpha' := |U'(\partial p')-1|$,
\begin{equation*}
\alpha' < B_3\varepsilon L^{-2j-2} .
\end{equation*}
Moreover, on each cube entering the definition of the class $\mathfrak{U}_{k+1}(\boldsymbol{\Omega}',B_3\varepsilon)$
there is a local gauge in which $U' = e^{i\eta' A'}$. In that gauge,
\begin{equation}\label{9.3:eq-minimiser-bounds-3}
\begin{gathered}
|A'| < B_3R_1M_1\varepsilon\,\ell^{-1}, \qquad |\nabla A'| < B_3R_1M_1\varepsilon\,\ell^{-2}, \\
\|\nabla A'\|_{\beta} < B_4(\beta)R_1M_1\varepsilon\,\ell^{-2-\beta} .
\end{gathered}
\end{equation}
Here $B_4(\beta)$ is the constant of \cite[(9)]{B-CMP102b}, depending on $\beta$, and $\|\nabla A'\|_{\beta}$
is the H\"older norm of exponent $\beta$ in the notation of \cite[(9)]{B-CMP102b} (not the
weighted sup norm $\|\cdot\|_{\theta}$ of Notation~\ref{9.3:not-paired-variational}).

The membership $U_{\mathfrak{B}}\in\mathfrak{U}_k(\boldsymbol{\Omega},B_3\varepsilon)$ is the same set of
conditions read on $T_{\eta}$. At level $j$ of $\boldsymbol{\Omega}$, with the same local scale
$\ell = L^{j}\eta$, every plaquette $p$ of $T_{\eta}$ at that level satisfies
\begin{equation*}
|U_{\mathfrak{B}}(\partial p)-1| < B_3\varepsilon L^{-2j} ;
\end{equation*}
the bounds \eqref{9.3:eq-minimiser-bounds-2} hold at every level of $\boldsymbol{\Omega}$ with
$U_{\mathfrak{B}}$ in place of $U'$; and on each cube entering the definition of
$\mathfrak{U}_k(\boldsymbol{\Omega},B_3\varepsilon)$ there is a local gauge in which
$U_{\mathfrak{B}} = e^{i\eta A}$, the potential $A$ satisfying the bounds
\eqref{9.3:eq-minimiser-bounds-3} in place of $A'$.

These statements are \cite[Theorem~1, (9), (10)]{B-CMP102b}, read at the two depths. The proof is
that of the cited paper and is not reproduced.
\end{definition}

\begin{lemma}[The chain lift through one averaging step]\label{9.3:lem-chain-lift}
Let $\eta=L^{-k}$ and $\eta'=\eta/L$. Let $U'$ be the minimal configuration on $T_{\eta'}$ of the paired
variational problems of Definition~\ref{9.3:not-paired-variational}, let $M$ be the one-step averaging map
of \cite[(0.11)--(0.12)]{B-CMP109}, and put $W:=M(U')$; the results of \cite{B-CMP98} on the averaging
\cite[(42)]{B-CMP98} that are used below hold for this mean, as stated in \cite{B-CMP109}. Let $\alpha'$
be the fine plaquette size of Definition~\ref{9.3:not-minimiser-bounds}. We write $C_L$ for constants
depending only on $d$, $L$ and the exponent $\theta$ of Definition~\ref{9.3:not-minimiser-bounds}, whose
value may change from line to line. For a function
$\delta A'$ on the bonds of $T_{\eta'}$ with values in $\mathfrak{g}$ we define a function on the bonds
of $T_\eta$ by
\begin{equation*}
Q_1\delta A' := (i\eta)^{-1}\frac{d}{ds}\Big[M\big(e^{is\eta'\delta A'}U'\big)W^{-1}\Big]_{s=0}.
\end{equation*}
For $y\in T_\eta$ let $B(y)$ be its block in $T_{\eta'}$, and for $x\in B(y)$ let $\Gamma_{y,x}$ be
the axial contour in $B(y)$ from $y$ to $x$. Fix the block-axial gauge $U'(\Gamma_{y,x})=1$ for
$x\in B(y)$, $y\in T_\eta$. For a bond $c$ of $T_\eta$ with initial point $c_-$ and final point $c_+$,
the \emph{corridor} of $c$ is
\begin{equation*}
B(c):=\{b' : b'_-\in B(c_-),\ b'_+\in B(c_+)\}.
\end{equation*}
In this gauge, by \cite{B-CMP98} and \cite[Lemma~1]{B-CMP99a},
\begin{equation}\label{9.3:eq-chain-lift-1}
\begin{aligned}
|U'(b')-1| &\le dL\alpha' &&\text{if both endpoints of $b'$ lie in one block},\\
|U'(b')-W(c)| &\le C_L\alpha' &&\text{if } b'\in B(c).
\end{aligned}
\end{equation}
For a function $a$ on the bonds of the domain $\Omega_1$ of the sequence $\boldsymbol{\Omega}$ of
Definition~\ref{9.3:not-paired-variational}, extended by zero to the other bonds of $T_\eta$, put
\begin{equation*}
(Ea)(b') :=
\begin{cases} L\,a(c) & \text{if } b'\in B(c),\\ 0 & \text{otherwise};\end{cases}
\end{equation*}
this is the linearisation of \cite[(1.3)]{B-CMP119}. For a plaquette $p'$ of $T_{\eta'}$ and a
plaquette $p$ of $T_\eta$ write $p'\sim p$ if the four bonds of $\partial p'$ lie in the corridors of
the four bonds of $\partial p$; for each $p$ there are $L^{d-2}$ such $p'$. Then:
\begin{itemize}
\item[(a)] $Q_1Ea=(1+\mathsf{k})a$, where $\mathsf{k}$ is bond-diagonal and
$|\mathsf{k}(c)|\le C_L\alpha'$ for every bond $c$;
\item[(b)] for every plaquette $p'$ of $T_{\eta'}$,
\begin{equation}\label{9.3:eq-chain-lift-2}
(D'_{U'}Ea)(p') = L^2(D_Wa)(p)\,\mathbb{1}_{p'\sim p} + (\mathsf{r}a)(p'),
\end{equation}
where $p$ is the plaquette with $p'\sim p$ if there is one (it is unique, since a fine bond lies in at
most one corridor) and the first term is absent otherwise; here $\mathsf{r}a=0$ unless $\partial p'$
meets a corridor, and
\begin{equation*}
|(\mathsf{r}a)(p')|\le C_L\,\eta'^{-1}\alpha'\sup|a|.
\end{equation*}
\end{itemize}
Here $D_W$ is the covariant derivative of Definition~\ref{9.3:not-paired-variational} on $T_\eta$ in the
background $W$, and $D'_{U'}$ the same operator on $T_{\eta'}$ in the background $U'$.
\end{lemma}

\begin{proof}
(a) Write $U'_s:=e^{is\eta' Ea}U'$. Since $\eta'L=\eta$, we have $\eta'(Ea)(b')=\eta\,a(c)$ for
$b'\in B(c)$, and $Ea$ vanishes on every bond whose endpoints lie in one block, because such a bond
belongs to no corridor. Every contour $\Gamma$ entering $M(U')(c)$, as described in
\cite[(42)]{B-CMP98} for the mean \cite[(0.11)--(0.12)]{B-CMP109} (see the statement above), lies in
$B(c_-)\cup B(c_+)$, crosses the
common face of the two blocks exactly once, through a bond $b'_\Gamma\in B(c)$, and meets no other
corridor. Let $\Gamma_<\subset B(c_-)$ be the part of $\Gamma$ before the crossing and $\Gamma_>$ the
part after it. Then only the variable on $b'_\Gamma$ is changed, and moving the factor
$e^{is\eta a(c)}$ past $U'(\Gamma_<)$ with the adjoint rotation $R$ gives
\begin{equation*}
\begin{aligned}
U'_s(\Gamma) &= U'(\Gamma_<)\,e^{is\eta a(c)}\,U'(b'_\Gamma)\,U'(\Gamma_>)\\
&= \exp\big(is\eta\,R(U'(\Gamma_<))a(c)\big)\,U'(\Gamma).
\end{aligned}
\end{equation*}
The contour $\Gamma_<$ consists of at most $C_L$ bonds of the block $B(c_-)$, on each of which
$|U'(b')-1|\le dL\alpha'$ by \eqref{9.3:eq-chain-lift-1}. Hence
$|R(U'(\Gamma_<))-1|\le C_L\alpha'$.

The value $M(\cdot)(c)$ is an analytic function of the contour variables $U'(\Gamma)$. At coinciding
arguments its logarithmic derivative is the arithmetic mean over the contours. At the configuration
$U'$ all loops formed from these contours are within $C_L\alpha'$ of $1$. Therefore
\begin{equation*}
(i\eta)^{-1}\frac{d}{ds}\Big[M(U'_s)(c)\,W(c)^{-1}\Big]_{s=0}
\end{equation*}
equals the arithmetic mean of the elements
$R(U'(\Gamma_<))a(c)$, up to an error at most $C_L\alpha'|a(c)|$. Each of these elements differs from
$a(c)$ by at most $C_L\alpha'|a(c)|$. This gives $(Q_1Ea)(c)=a(c)+\mathsf{k}(c)a(c)$, where
$\mathsf{k}(c)$ is a linear map on $\mathfrak{g}$ with $|\mathsf{k}(c)|\le C_L\alpha'$.

Since the contours entering $M(\cdot)(c)$ meet no corridor other than $B(c)$, only $a(c)$ enters
$(Q_1Ea)(c)$. Hence $\mathsf{k}$ is bond-diagonal.

(b) If $\partial p'$ meets no corridor, then $Ea$ vanishes on all four bonds of $\partial p'$. Then
$(D'_{U'}Ea)(p')=0$, no $p$ satisfies $p'\sim p$, and $(\mathsf{r}a)(p')=0$.

Let $\partial p'$ meet a corridor. Then one of two cases holds.

In the first case, $\partial p'$ has two opposite bonds in one corridor $B(c)$, and its other two bonds
lie inside blocks. In this case no $p$ satisfies $p'\sim p$, and
\begin{equation*}
(D'_{U'}Ea)(p') = \eta'^{-1}L\big[a(c)-R(U'(b'_0))a(c)\big],
\end{equation*}
where $b'_0$ is one of the two bonds of $\partial p'$ lying inside a block. By
\eqref{9.3:eq-chain-lift-1},
$|R(U'(b'_0))-1|\le C_L\alpha'$. So $(\mathsf{r}a)(p'):=(D'_{U'}Ea)(p')$ satisfies
$|(\mathsf{r}a)(p')|\le C_L\eta'^{-1}\alpha'\sup|a|$.

In the second case $p'\sim p$. Let $b'_i$ be the bonds of $\partial p'$ and $c_i$ the bonds of
$\partial p$, with $b'_i\in B(c_i)$. The two sides of \eqref{9.3:eq-chain-lift-2} to be compared are
\begin{equation*}
\eta'^{-1}L\sum_i \pm R(U'(b'_i))a(c_i)
\qquad\text{and}\qquad
L^2\eta^{-1}\sum_i \pm R(W(c_i))a(c_i),
\end{equation*}
with the same signs. The rotations are those of the transports entering the definitions of
$D'_{U'}$ and $D_W$. Because $\eta'=\eta/L$, the prefactors agree: $\eta'^{-1}L=L^2\eta^{-1}$. By the
corridor bound in \eqref{9.3:eq-chain-lift-1}, $|U'(b'_i)-W(c_i)|\le C_L\alpha'$ for each $i$, so
corresponding rotations differ by at most $C_L\alpha'$. Hence the difference
$(\mathsf{r}a)(p'):=(D'_{U'}Ea)(p')-L^2(D_Wa)(p)$ satisfies
$|(\mathsf{r}a)(p')|\le C_L\eta'^{-1}\alpha'\sup|a|$.

The two cases exhaust the plaquettes meeting a corridor, which proves (b).
\end{proof}

\begin{lemma}[The face-mean multiplier and the current defect]\label{9.3:lem-current-defect}
Let $U'$ be the minimiser at depth $k+1$ and $W := M(U')$, with the bounds of
Notation~\ref{9.3:not-minimiser-bounds}, and let $E$, $\mathsf{k}$, $\mathsf{r}$ and the corridors
$B(c)$ be those of the chain lift through one averaging step (Lemma~\ref{9.3:lem-chain-lift}). Write
$D' := D_{U'}$, and let $\langle\cdot,\cdot\rangle_{\eta}$, $\langle\cdot,\cdot\rangle_{\eta'}$ be the
pairings on $T_{\eta}$, $T_{\eta'}$ (weights $\eta^{d}$, $\eta'^{d}$). Put
$\widehat{E} := E(1+\mathsf{k})^{-1}$ and define the bond function $\mu$ on $T_{\eta}$ by
\begin{equation}\label{9.3:eq-current-defect-1}
\langle a,\mu\rangle_{\eta} := \langle \widehat{E}a, J_{U'}\rangle_{\eta'},
\end{equation}
where $a$ ranges over the $\mathfrak{g}$-valued bond functions on $\Omega_1$.
At a bond or plaquette of level $j$ we write $\ell = L^{j}\eta$.
We assume $\varepsilon$ small enough, depending on $(d,L,\theta)$ only, that
$C_LB_3\varepsilon \le 1/3$, where $C_L$ is the constant in the bound
$|\mathsf{k}(c)| \le C_L\alpha'$ of Lemma~\ref{9.3:lem-chain-lift}(a).
\begin{enumerate}
\item[(a)] $\langle \delta A'', J_{U'}\rangle_{\eta'} = \langle Q_1\delta A'', \mu\rangle_{\eta}$ for all
$\mathfrak{g}$-valued bond functions $\delta A''$ on $T_{\eta'}$; $|\mu| \le 3B_3\varepsilon\ell^{-3}$;
$\langle D_W\lambda,\mu\rangle_{\eta} = 0$ for all $\mathfrak{g}$-valued site functions $\lambda$ on
$T_{\eta}$; and $\langle a,\mu\rangle_{\eta} = 0$ whenever $Q_j(W)a = 0$ on $\Lambda_j$ for all $j$.
\item[(b)] On $\Omega_1$ there are a plaquette function $f$ and a bond function $f_0$ with
\begin{equation}\label{9.3:eq-current-defect-a}
J_W = \mu + D_W^{*}f + f_0,\qquad
|f| \le C_f\varepsilon\ell^{-2}L^{-\beta j},\qquad
|f_0| \le C_f\varepsilon^{2}\ell^{-3}L^{-j},
\end{equation}
for a constant $C_f$ depending on $(d,L,\theta)$ and $R_1M_1$ only.
\end{enumerate}
\end{lemma}

The decomposition \eqref{9.3:eq-current-defect-a} is asserted for the bonds of $\Omega_1$; its value
form across the collar pulled back from the coarser problem is not asserted.

\begin{proof}
Throughout, $C$ and $C_L$ denote constants depending on $(d,L,\theta)$ only, not necessarily the same
at each occurrence.

\emph{(a).} Let $M_{\mathfrak{B}'}$ denote the constraint map of the finer problem, whose level set
$\{M_{\mathfrak{B}'} = V'\}$ is the constraint set of the finer problem, on which $U'$ is the minimiser.
By \cite[(178)]{B-CMP102b}, applied at depth $k+1$, the current of the minimiser is orthogonal to
every variation $\delta A'$ in the kernel of the linearised constraint at $U'$, and this relation
involves no gauge-fixing restriction on $\delta A'$. Equivalently, $U' \in
\mathfrak{U}_{k+1}(\boldsymbol{\Omega}',B_3\varepsilon)$ is an interior minimum of the action on
$\{M_{\mathfrak{B}'} = V'\}$, so that $J_{U'} \perp \ker DM_{\mathfrak{B}'}(U')$. The constraints of
the finer problem are expressed through the maps $M^{j+1} = M^{j}\circ M$; by the chain rule,
\begin{equation*}
Q_{j+1}(U') = Q_j(W)\,Q_1 .
\end{equation*}
Hence every $\delta A''$ with $Q_1\delta A'' = 0$ lies in the kernel of the linearised constraints at
$U'$, and $J_{U'} \perp \ker Q_1$. By Lemma~\ref{9.3:lem-chain-lift}(a), $Q_1E = 1+\mathsf{k}$, so
$Q_1\widehat{E} = 1$ and $Q_1(\delta A'' - \widehat{E}Q_1\delta A'') = 0$, i.e.\
$\delta A'' - \widehat{E}Q_1\delta A'' \in \ker Q_1$. Therefore
\begin{equation*}
\langle \delta A'', J_{U'}\rangle_{\eta'}
= \langle \widehat{E}Q_1\delta A'', J_{U'}\rangle_{\eta'}
= \langle Q_1\delta A'', \mu\rangle_{\eta},
\end{equation*}
the last equality being \eqref{9.3:eq-current-defect-1} with $a = Q_1\delta A''$.

To make $\mu$ explicit, insert $(Ea)(b') = L\,a(c)$ for $b' \in B(c)$ into
\eqref{9.3:eq-current-defect-1} and use $\eta'^{d} = L^{-d}\eta^{d}$. Writing
$(1+\mathsf{k})^{-\dagger}$ for the adjoint of the bond-diagonal operator $(1+\mathsf{k})^{-1}$, we get
\begin{equation}\label{9.3:eq-current-defect-2}
\mu(c) = (1+\mathsf{k})^{-\dagger}\, L^{1-d}\sum_{b'\in B(c)} J_{U'}(b'),
\qquad \#B(c) = L^{d-1}.
\end{equation}
Thus $\mu(c)$ is the mean of the fine current over the $L^{d-1}$ fine bonds crossing the common face of
the two blocks joined by $c$, corrected by $(1+\mathsf{k})^{-\dagger}$. By
Notation~\ref{9.3:not-minimiser-bounds}, $|J_{U'}| \le 2B_3\varepsilon\ell^{-3}$, and $\ell$ is the same
for $c$ and for the fine bonds of $B(c)$, since level $j$ of the coarse problem corresponds to level
$j+1$ of the fine one and $L^{j}\eta = L^{j+1}\eta'$. The mean in \eqref{9.3:eq-current-defect-2} is
therefore at most $2B_3\varepsilon\ell^{-3}$. Since $|\mathsf{k}(c)| \le C_L\alpha' < C_LB_3\varepsilon$,
the factor $(1+\mathsf{k})^{-\dagger}$ has norm at most $(1-C_L\alpha')^{-1}$. This is at most $3/2$,
since $C_LB_3\varepsilon \le 1/3$ by the smallness of $\varepsilon$ assumed in the statement, which
gives $|\mu| \le 3B_3\varepsilon\ell^{-3}$.

For every gauge transformation $u$ of $T_{\eta'}$ the averaging map satisfies
$M(U'^{u}) = M(U')^{u|_{T_{\eta}}}$. Differentiating at $s=0$ along the one-parameter family
$u = e^{is\eta'\lambda'}$, $\lambda'$ a $\mathfrak{g}$-valued site function on $T_{\eta'}$, gives
\begin{equation*}
Q_1 D'\lambda' = D_W\bigl(\lambda'|_{T_{\eta}}\bigr).
\end{equation*}
Moreover $D'^{*}J_{U'} = 0$. Given $\lambda$ on $T_{\eta}$, choose $\lambda'$ on $T_{\eta'}$ with
$\lambda'|_{T_{\eta}} = \lambda$. By the first identity of (a),
\begin{equation*}
\langle D_W\lambda,\mu\rangle_{\eta}
= \langle Q_1D'\lambda', \mu\rangle_{\eta}
= \langle D'\lambda', J_{U'}\rangle_{\eta'}
= \langle \lambda', D'^{*}J_{U'}\rangle_{\eta'} = 0 .
\end{equation*}
Finally, let $Q_j(W)a = 0$ on $\Lambda_j$ for all $j$. Then $Q_1\widehat{E}a = a$, so
$Q_{j+1}(U')\widehat{E}a = Q_j(W)a = 0$ on $\Lambda_j$ for all $j$. Hence $\widehat{E}a$ lies in the
kernel of the linearised constraints at $U'$, and \cite[(178)]{B-CMP102b} gives
$\langle a,\mu\rangle_{\eta} = \langle \widehat{E}a, J_{U'}\rangle_{\eta'} = 0$.

\emph{(b).} Let $\delta A$ be a bond function on $\Omega_1$ and put $a := (1+\mathsf{k})^{-1}\delta A$,
so that $\widehat{E}\delta A = Ea$. Since $J_{U'} = D'^{*}\mathfrak{F}_{U'}$, we have
$\langle \delta A,\mu\rangle_{\eta} = \langle D'Ea, \mathfrak{F}_{U'}\rangle_{\eta'}$. We insert
Lemma~\ref{9.3:lem-chain-lift}(b),
\begin{equation*}
(D'Ea)(p') = L^{2}(D_Wa)(p)\mathbb{1}_{p'\sim p} + (\mathsf{r}a)(p'),
\end{equation*}
and use $\eta'^{d}L^{2} = \eta^{d}L^{2-d}$. This gives
\begin{equation}\label{9.3:eq-current-defect-3}
\langle \delta A,\mu\rangle_{\eta}
= \sum_{p}\eta^{d}\operatorname{tr}\bigl((D_Wa)(p)\,g(p)\bigr)
+ \langle \mathsf{r}a, \mathfrak{F}_{U'}\rangle_{\eta'},
\qquad
g(p) := L^{2-d}\sum_{p'\sim p}\mathfrak{F}_{U'}(p').
\end{equation}
Since exactly $L^{d-2}$ plaquettes $p'$ satisfy $p' \sim p$ (Lemma~\ref{9.3:lem-chain-lift}), $g(p)$
is a normalised mean of $\mathfrak{F}_{U'}$. On the other hand, $J_W = D_W^{*}\mathfrak{F}_W$ gives
$\langle \delta A, J_W\rangle_{\eta} = \langle D_W\delta A, \mathfrak{F}_W\rangle_{\eta}$. Put
$f := \mathfrak{F}_W - g$. Since $\delta A - a = \mathsf{k}(1+\mathsf{k})^{-1}\delta A$,
\eqref{9.3:eq-current-defect-3} yields
\begin{equation*}
\langle \delta A, J_W - \mu - D_W^{*}f\rangle_{\eta}
= \bigl\langle D_W\bigl(\mathsf{k}(1+\mathsf{k})^{-1}\delta A\bigr), g\bigr\rangle_{\eta}
- \langle \mathsf{r}a, \mathfrak{F}_{U'}\rangle_{\eta'} .
\end{equation*}
We let $f_0$ be the bond function on $\Omega_1$ representing the right side; this is
\eqref{9.3:eq-current-defect-a}. It collects two terms.

The term $\langle \mathsf{r}a, \mathfrak{F}_{U'}\rangle_{\eta'}$ contributes
$\mathsf{r}^{\dagger}\mathfrak{F}_{U'}$, composed with $(1+\mathsf{k})^{-\dagger}$.
By Lemma~\ref{9.3:lem-chain-lift}(b), $\mathsf{r}a(p')$ vanishes unless $\partial p'$ meets a corridor,
and $|(\mathsf{r}a)(p')| \le C_L\eta'^{-1}\alpha'\sup|a|$. At most $C_LL^{d-1}$ fine plaquettes
contribute to a given $c$, and the weight ratio is $\eta'^{d}/\eta^{d} = L^{-d}$. Since
$\eta'^{-1} = L\eta^{-1}$, this term is at most $C_L\eta^{-1}\alpha'|\mathfrak{F}_{U'}|$.

The term containing $\mathsf{k}(1+\mathsf{k})^{-1}$ is estimated with the crude bound
$|D_W(\mathsf{k}(1+\mathsf{k})^{-1}\delta A)| \le 4\eta^{-1}|\mathsf{k}||\delta A|$, paired against $g$.
Here $|g| \le \sup|\mathfrak{F}_{U'}|$, because $g$ is a normalised mean, and each bond lies in
$2(d-1)$ plaquettes.

We now use $|\mathsf{k}| \le C_L\alpha'$, $\alpha' < B_3\varepsilon L^{-2j-2}$,
$|\mathfrak{F}_{U'}| \le B_3\varepsilon\ell^{-2}$ and $\eta^{-1} = L^{j}\ell^{-1}$. Since
\begin{equation*}
\eta^{-1}\alpha'|\mathfrak{F}_{U'}|
\le L^{j}\ell^{-1}\cdot B_3\varepsilon L^{-2j-2}\cdot B_3\varepsilon\ell^{-2}
= B_3^{2}L^{-2}\varepsilon^{2}\ell^{-3}L^{-j},
\end{equation*}
both terms give $|f_0| \le C\varepsilon^{2}\ell^{-3}L^{-j}$.

For $f$ we use \cite[(47)--(49)]{B-CMP98}. Let $y_0$ be a corner of $p$; for $x \in B(y_0)$ let
$(p)_x$ be the family of $L^{2}$ fine plaquettes parallel to $p$ that \cite[(47)]{B-CMP98} attaches
to the point $x$. That reference gives
\begin{equation*}
W(\partial p) - 1 = \sum_{x\in B(y_0)} L^{-d}\sum_{p'\subset (p)_x}\bigl(U'(\partial p')-1\bigr)
+ O(L^{4}\alpha'^{2}).
\end{equation*}
We multiply the imaginary part by $\eta^{-2} = L^{-2}\eta'^{-2}$ and use
$\alpha'^{2} < B_3^{2}\varepsilon^{2}L^{-4j-4}$ and $\eta^{-2} = L^{2j}\ell^{-2}$. This gives
\begin{equation}\label{9.3:eq-current-defect-4}
\mathfrak{F}_W(p) = L^{-d-2}\sum_{x\in B(y_0)}\sum_{p'\subset(p)_x}\mathfrak{F}_{U'}(p')
+ O\bigl(\varepsilon^{2}\ell^{-2}L^{-2j}\bigr),
\end{equation}
again a normalised mean, over $L^{d}\cdot L^{2}$ plaquettes.

Both means, in \eqref{9.3:eq-current-defect-3} and \eqref{9.3:eq-current-defect-4}, run over plaquettes
parallel to $p$ within distance $2d\eta$ of it. Hence $|f(p)|$ is bounded by the oscillation of
$\mathfrak{F}_{U'}$ over such plaquettes, plus the cost of the parallel transports needed to compare
its values, plus the error in \eqref{9.3:eq-current-defect-4}.

If $L^{j} < 2d$ the bound on $f$ is trivial: then $L^{-\beta j} > (2d)^{-\beta}$, while
$|g| \le B_3\varepsilon\ell^{-2}$, and by \eqref{9.3:eq-current-defect-4}
$|\mathfrak{F}_W| \le B_3\varepsilon\ell^{-2} + O(\varepsilon^{2}\ell^{-2}L^{-2j})$.

Let $L^{j} \ge 2d$; then every distance $r \le 2d\eta$ between the plaquettes entering the two means
satisfies $r \le \ell$. In the local gauge of Notation~\ref{9.3:not-minimiser-bounds},
\begin{equation*}
\mathfrak{F}_{U'} = \partial A' + i[A',A'] + O\bigl(\eta'|A'||\nabla A'|\bigr).
\end{equation*}
Its oscillation over distances $r \le 2d\eta$ is therefore at most
\begin{equation*}
\begin{split}
2\|\nabla A'\|_{\beta}\,r^{\beta} + C|A'||\nabla A'|\,r
&\le 2B_4(\beta)R_1M_1\varepsilon\ell^{-2-\beta}(2d)^{\beta}\ell^{\beta}L^{-\beta j}\\
&\quad + C(R_1M_1)^{2}\varepsilon^{2}\ell^{-3}\cdot 2d\,\ell L^{-j},
\end{split}
\end{equation*}
where we used $|A'| < B_3R_1M_1\varepsilon\ell^{-1}$, $|\nabla A'| < B_3R_1M_1\varepsilon\ell^{-2}$,
$\|\nabla A'\|_{\beta} < B_4(\beta)R_1M_1\varepsilon\ell^{-2-\beta}$, with $R_1M_1$ the constant of
Notation~\ref{9.3:not-minimiser-bounds}, and $\eta = \ell L^{-j}$. The
transports cost at most
\begin{equation*}
C\eta|A'||\mathfrak{F}_{U'}| \le CR_1M_1\varepsilon^{2}\ell^{-2}L^{-j}.
\end{equation*}
The term $O(\eta'|A'||\nabla A'|)$ and the error in \eqref{9.3:eq-current-defect-4} are of the same or
smaller order. All these contributions are at most $C(R_1M_1)^{2}\varepsilon\ell^{-2}L^{-\beta j}$,
since $R_1M_1 \ge 1$, the terms of order $\varepsilon^{2}$ being so for $\varepsilon \le 1$, since
$\beta \le 1$. This gives $|f| \le C(R_1M_1)^{2}\varepsilon\ell^{-2}L^{-\beta j}$. Both bounds of (b)
hold with $C_f$ the larger of the two constants obtained for $f_0$ and for $f$.
\end{proof}

\begin{remark}
In the flat abelian case, where the covariant derivatives are the lattice exterior derivatives
$\partial'$, $\partial$ and $\mathsf{k} = 0$, one has $\mu = \partial^{*}g$ exactly. In the sum
$\sum_{b'\in B(c)}(\partial'^{*}\mathfrak{F}_{U'})(b')$, the plaquettes lying between two bonds of one
corridor cancel in pairs, and only the junction plaquettes remain.
\end{remark}

\begin{lemma}[Gauge fixing under a divergence-form current bound]\label{9.3:lem-gauge-fixing}
Let $\eta=L^{-k}$, the admissible sequence of domains $\boldsymbol{\Omega}$, the determining set
$\mathfrak{B}$, the level $j(x)$ of a point $x$ and the local scale $\ell=L^{j(x)}\eta$ be as in
Notation~\ref{9.3:not-paired-variational}, and let $\epsilon$ be the H\"older exponent of
Notation~\ref{9.3:not-minimiser-bounds}. Write $|\cdot|_{(\gamma)}$ for the level-weighted sup
norm of \cite[(3.41)]{B-CMP99b}. Let $\alpha_0,\alpha_1>0$. Let
$U_0\in\mathfrak{U}_k(\boldsymbol{\Omega},\alpha_0)$, where $\mathfrak{U}_k$ is the class of regular
configurations of \cite[Theorem~1]{B-CMP102b}, and assume that $U_0$ satisfies
\cite[(3.35)]{B-CMP99b}. Let $U$ satisfy condition \cite[(1.7)]{B-CMP99a} with $\alpha_0$, let
$UU_0^{-1}\in\mathrm{Ax}_k(\mathfrak{B},U_0)$, the axial class of \cite{B-CMP99a}, let $U$ satisfy
condition \cite[(1.66)]{B-CMP99a} with $\alpha_1$, and assume
\begin{equation}\label{9.3:eq-gauge-fixing-a}
\begin{gathered}
(i\eta^2)^{-1}D_U^{*}\partial U=j_0+D_U^{*}f,\qquad |j_0|_{(-3)}\le\alpha_0,\\
\mathsf{N}_{\epsilon}(f):=\sup_{j}\ \sup_{x:\,j(x)=j}L^{j\epsilon}\ell^{2}|f(x)|\le\alpha_0 .
\end{gathered}
\end{equation}
Then there are a constant $c_1^{\sharp}=c_1^{\sharp}(d,L,\epsilon)>0$ and a constant
$B_1^{\sharp}$ such that, if $\alpha_0+\alpha_1\le c_1^{\sharp}$, there is exactly one gauge
transformation $u$ satisfying \cite[(1.29)]{B-CMP99a} such that
$U_1:=(UU_0^{-1})^{u^{-1}}=e^{i\eta A}$ obeys \cite[(1.37)]{B-CMP99a},
\cite[(1.38)]{B-CMP99a} and
\begin{equation}\label{9.3:eq-gauge-fixing-1}
\ell|A|\le B_1^{\sharp}(\alpha_0+\alpha_1),\qquad
\ell^{2}|\nabla_{U_0}A|\le B_1^{\sharp}(\alpha_0+\alpha_1).
\end{equation}
\end{lemma}

The lemma is \cite[Theorem~4]{B-CMP99a} with the hypothesis \cite[(1.9)]{B-CMP99a} on $U$
replaced by the current condition \eqref{9.3:eq-gauge-fixing-a}.

\begin{proof}
We follow the proof of \cite[Theorem~4]{B-CMP99a} and show at each stage that the hypothesis
\cite[(1.9)]{B-CMP99a} on $U$ can be replaced by \eqref{9.3:eq-gauge-fixing-a}.

\emph{Step 1. The induction.} \cite[Theorem~4]{B-CMP99a} is proved by induction on $k$, and the
hypotheses of the theorem at $k$ pass to the corresponding hypotheses at $k-1$
\cite{B-CMP99a}. The condition \eqref{9.3:eq-gauge-fixing-a} passes to $k-1$ as well. Indeed, the
bound $\alpha_0\ell^{-3}$ on $|j_0|$ imposed by $|j_0|_{(-3)}\le\alpha_0$ decreases in $j$, and the
weight $L^{j\epsilon}\ell^{2}$ in the definition of $\mathsf{N}_{\epsilon}$ increases in $j$.

\emph{Step 2. The inductive output is used only to first order.} In the inductive step of
\cite{B-CMP99a} only the first-order part of the output of the induction hypothesis enters. The
proof there states: ``The other conditions in (1.37), (1.38), and (1.62) hold also, but we will
not need them'', and further: ``We will use only the properties (1.69), (1.73), (1.74)''
\cite{B-CMP99a}.

\emph{Step 3. The first-order bound on $A$.} \cite[Proposition~5]{B-CMP99a} assumes exactly the
properties \cite[(1.69), (1.73), (1.74)]{B-CMP99a} listed in Step 2. Its conclusions
\cite[(1.110)--(1.111)]{B-CMP99a} give, with the constants $C_1'$ and $B_1$ of \cite{B-CMP99a},
\begin{equation}\label{9.3:eq-gauge-fixing-2}
\ell|A|\le C_1'B_1(\alpha_0+\alpha_1)=:\alpha_2 .
\end{equation}

\emph{Step 4. Restoring the constants.} In \cite{B-CMP99a} the constants of the induction are
restored by \cite[Proposition~3]{B-CMP99a}, and this is the only place in the proof where the
hypothesis \cite[(1.9)]{B-CMP99a} on $U=U_1U_0$ enters. There, \cite[(1.54)]{B-CMP99a} is an
identity with remainders of type $O_1$, which the proof states is ``derived under the assumption
(0.1) for $U_0$ and (1.41) for $U_1$'' \cite{B-CMP99a}; condition \cite[(0.1)]{B-CMP99a} is a
condition on the plaquette variables only. The proof then continues: ``From (1.40) the term on the
left-hand side and the first term on the right-hand side can be estimated by
$\alpha_0(L^j\eta)^{-3}\eta^2$'' \cite{B-CMP99a}.

Under \eqref{9.3:eq-gauge-fixing-a} in place of \cite[(1.9)]{B-CMP99a}, the estimate drawn from
\cite[(1.40)]{B-CMP99a} in the sentence just quoted is replaced by the decomposition
\begin{equation}\label{9.3:eq-gauge-fixing-3}
J:=D_{U_0}^{*}D_{U_0}A=j_1+D_{U_0}^{*}f_1,\qquad f_1:=R(u)f,
\end{equation}
where $R(u)$ is the adjoint rotation by $u$, and where the current $j_1$ satisfies
\begin{equation}\label{9.3:eq-gauge-fixing-4}
|j_1|_{(-3)}\le 2\alpha_0+2d\alpha_2\alpha_0+|r|_{(-3)},
\end{equation}
with $r$ the sum of the remainder terms of type $O_1$ of \cite[(1.54)]{B-CMP99a}, which are
estimated there. The bound
\eqref{9.3:eq-gauge-fixing-4} uses that the operator $D_U^{*}-D_{U_0}^{*}$ is of order zero,
together with \eqref{9.3:eq-gauge-fixing-2}.

By \cite[(1.58)]{B-CMP99a},
\begin{equation}\label{9.3:eq-gauge-fixing-5}
A=G(U_0)J+G(U_0)\sum_{j}Q_j^{*}(L^j\eta)^{-3}B_1 ,
\end{equation}
with $G(U_0)$, $Q_j$ and $B_1$ as in \cite[(1.58)]{B-CMP99a}. Inserting
\eqref{9.3:eq-gauge-fixing-3} into \eqref{9.3:eq-gauge-fixing-5}, the only term not present in
\cite{B-CMP99a} is $G D^{*}f_1$, where we write $G=G(U_0)$ and $D^{*}=D_{U_0}^{*}$. We estimate it
and its covariant derivative.

For the term itself, \cite[(3.47)]{B-CMP99b} with $\gamma=-2$ gives
\begin{equation}\label{9.3:eq-gauge-fixing-6}
|GD^{*}f_1|_{(-1)}\le B_0|f_1|_{(-2)} ,
\end{equation}
with the constant $B_0$ of \cite[(3.47)]{B-CMP99b}. Since $f_1$ is obtained from $f$ by an adjoint
rotation, $|f_1|=|f|$ pointwise; since $L^{j\epsilon}\ge1$, the definition of
$\mathsf{N}_{\epsilon}$ in \eqref{9.3:eq-gauge-fixing-a} gives $|f_1|_{(-2)}\le
\mathsf{N}_{\epsilon}(f)\le\alpha_0$.

For the derivative, take a partition of unity $\{\zeta_{y'}\}$ with
$\zeta_{y'}\in C_0^{\infty}(\widetilde{\Delta}(y'))$ on the enlarged cubes
$\widetilde{\Delta}(y')$ and $\|\zeta_{y'}\|^{\xi'}_{1}\le4$, where $y'$ is a cube of level $j'$,
$\xi'=L^{-j'}$, and $\|\cdot\|^{\xi'}_{\epsilon}$ is the scaled H\"older norm of
\cite[(3.44)]{B-CMP99b}. Then \cite[(3.44)]{B-CMP99b} gives, for $x$ in the cube $y$ of the cover
underlying \cite[(3.44)]{B-CMP99b} that contains $x$,
\begin{equation}\label{9.3:eq-gauge-fixing-7}
|\nabla GD^{*}f_1|(x)\le\sum_{y'}B_0'(\epsilon)\,e^{-\delta_0 d(y,y')}
\bigl(\|\zeta_{y'}f_1\|^{\xi'}_{\epsilon}+|\zeta_{y'}f_1|\bigr),
\end{equation}
with the constant $B_0'(\epsilon)$ and the decay rate $\delta_0$ of \cite[(3.44)]{B-CMP99b}, and the
multi-scale distance $d(\cdot,\cdot)$ of \cite{B-CMP96}.

The H\"older norm on the right of \eqref{9.3:eq-gauge-fixing-7} is estimated by the supremum alone.
On a lattice of spacing $\eta$, distinct points are at distance at least $\eta$. Hence, for every
function $h$ on the lattice,
\begin{equation}\label{9.3:eq-gauge-fixing-b}
\|h\|^{\xi'}_{\epsilon}\le 2L^{j'\epsilon}\sup|h| ,
\end{equation}
whatever the parallel transports used to compare values of $h$ at different points. In particular
no regularity of the gauge transformation $u$ enters. Since $0\le\zeta_{y'}\le1$, the product
bound
\begin{equation*}
\|\zeta_{y'}f_1\|^{\xi'}_{\epsilon}\le\|\zeta_{y'}\|^{\xi'}_{1}\sup|f_1|+\|f_1\|^{\xi'}_{\epsilon}
\end{equation*}
holds. With $\|\zeta_{y'}\|^{\xi'}_{1}\le4$, with \eqref{9.3:eq-gauge-fixing-b} for $h=f_1$, and
with $|\zeta_{y'}f_1|\le|f_1|=|f|$ and $L^{j'\epsilon}\ge1$, the bracket in
\eqref{9.3:eq-gauge-fixing-7} is at most $(4+2+1)L^{j'\epsilon}\sup|f|$, the norms and suprema
being taken over the support of $\zeta_{y'}$. On the support of
$\zeta_{y'}$, the definition of $\mathsf{N}_{\epsilon}$ gives
\begin{equation*}
L^{j'\epsilon}|f|\le\mathsf{N}_{\epsilon}(f)(L^{j'}\eta)^{-2}.
\end{equation*}
The right side of \eqref{9.3:eq-gauge-fixing-7} is therefore bounded as follows:
\begin{equation}\label{9.3:eq-gauge-fixing-8}
\begin{split}
|\nabla GD^{*}f_1|(x)
&\le 7B_0'(\epsilon)\,\mathsf{N}_{\epsilon}(f)\sum_{y'}e^{-\delta_0 d(y,y')}(L^{j'}\eta)^{-2}\\
&\le C\ell^{-2}\alpha_0 .
\end{split}
\end{equation}
The second inequality, in which $C$ is a constant depending on $d$, $L$ and $\epsilon$, uses the
summation bounds of \cite[Lemma~2.1, (2.60)--(2.61)]{B-CMP96},
which apply under the floor \eqref{9.3:eq-paired-variational-d}, together with
$\mathsf{N}_{\epsilon}(f)\le\alpha_0$.

By \eqref{9.3:eq-gauge-fixing-6} and \eqref{9.3:eq-gauge-fixing-8}, the new term $GD^{*}f_1$
satisfies bounds of the same form as those obtained in \cite{B-CMP99a} from
\cite[(1.40)]{B-CMP99a}. By \eqref{9.3:eq-gauge-fixing-4}, the current $j_1$ enters as the
corresponding current does there. Hence the first two bounds of
\cite[(1.59)--(1.62)]{B-CMP99a} hold with $\alpha_0$ replaced by $(1+C)\alpha_0$. With these bounds
in place of the ones derived from \cite[(1.9)]{B-CMP99a}, the restoration of the constants in
\cite[Proposition~3]{B-CMP99a} proceeds as in \cite{B-CMP99a}. This gives the bounds
\eqref{9.3:eq-gauge-fixing-1} for a constant $B_1^{\sharp}$ and a threshold
$c_1^{\sharp}(d,L,\epsilon)$, and the conditions \cite[(1.37)]{B-CMP99a} and
\cite[(1.38)]{B-CMP99a}.

\emph{Step 5. Uniqueness.} The uniqueness part of the proof of \cite[Theorem~4]{B-CMP99a} uses only
three inputs: the first-order bounds, \cite[Proposition~8]{B-CMP98} and
\cite[Proposition~5]{B-CMP99a}. The first-order bounds are available by Step 3, and the hypotheses
of \cite[Proposition~5]{B-CMP99a} hold by Steps 2 and 3. Hence that argument applies without
change, and $u$ is unique.
\end{proof}

The lemma is distinct from the generalisation considered after \cite[(1.147)]{B-CMP99a}. There, in
place of \cite[(1.9)]{B-CMP99a}, the weaker H\"older regularity condition \cite[(1.147)]{B-CMP99a}
is assumed, and in the paragraph following it an analogue of \cite[Theorem~2]{B-CMP99a}, without
the conclusion \cite[(1.39)]{B-CMP99a}, is stated without proof; Lemma~\ref{9.3:lem-gauge-fixing}
is not that statement and
does not use it. The hypotheses of the two statements are not comparable. The one-step average
$W:=M(U')$ of the minimiser $U'$ at the finer depth of Notation~\ref{9.3:not-minimiser-bounds}
satisfies both. In the lemma the function $f$ is small enough that the crude lattice bound
\eqref{9.3:eq-gauge-fixing-b} on the H\"older norm suffices.

\begin{lemma}[The averaged minimiser in the gauge-fixed slice]\label{9.3:prf-averaged-gauge}
Let $\boldsymbol{\Omega}$, $\mathfrak{B}=\bigcup_j\Lambda_j$, the data $V$ satisfying the condition
$(7)_{\varepsilon}$ of Notation~\ref{9.3:not-paired-variational},
the minimiser $U_{\mathfrak{B}}(V)$, the paired data $\boldsymbol{\Omega}'$, $V'$, the finer minimiser
$U'=U_{\mathfrak{B}'}(V')$ and its one-step average $W := M(U')$ on $\Omega_0$ be as in
Notation~\ref{9.3:not-paired-variational}, so that $M^{j}(W)=V$ on $\Lambda_j$ for every $j$. Assume
that $\varepsilon$ is sufficiently small, depending on $(d,L,\theta)$ only. Then there are a gauge
transformation $\mathsf{u}$ and a Lie-algebra valued field
$A$ on the bonds of $\Omega_0$ such that
\begin{equation}\label{9.3:eq-averaged-gauge-a}
\begin{gathered}
\widehat{W} := W^{\mathsf{u}} = e^{i\eta A}\,U_{\mathfrak{B}}, \qquad
R(U_{\mathfrak{B}})D^{*}A = 0, \qquad Q_j(U_{\mathfrak{B}},\eta A) = 0 \ \text{ for all } j,\\
\|A\|_{0} \le \varepsilon_2 := C\varepsilon ,
\end{gathered}
\end{equation}
where $C$ depends on $(d,L,\theta)$ only and
$\|A\|_{0}=\max\{|A|_{(-1)},\,|\nabla_{U_{\mathfrak{B}}}A|_{(-2)}\}$.
\end{lemma}

\begin{proof}
We apply Lemma~\ref{9.3:lem-gauge-fixing} with $U_0 := U_{\mathfrak{B}}(V)$. By the
bounds of Notation~\ref{9.3:not-minimiser-bounds},
$U_0\in\mathfrak{U}_k(\boldsymbol{\Omega},B_3\varepsilon)$. Let $\lambda$ be a gauge transformation,
equal to the identity on $\mathfrak{B}$, with $W^{\lambda}U_0^{-1}\in\mathrm{Ax}_k(\mathfrak{B},U_0)$
\cite{B-CMP99a}, and take $U := W^{\lambda}$. The plaquette variables and the current of
$W^{\lambda}$ are those of $W$ conjugated pointwise by $\lambda$, and the bounds below are bounds
in norm, so we state them for $W$.

\emph{Condition \cite[(1.7)]{B-CMP99a}.} By \cite[Proposition~1]{B-CMP98} and the bound
$\alpha' < B_3\varepsilon L^{-2j-2}$ of Notation~\ref{9.3:not-minimiser-bounds}, for a plaquette $p$
of level $j$,
\begin{equation}\label{9.3:eq-averaged-gauge-1}
|W(\partial p)-1| \le L^{2}\alpha' + C_0\,(L^{2}\alpha')^{2} \le 2B_3\varepsilon L^{-2j},
\end{equation}
where $C_0$ is the constant of \cite[Proposition~1]{B-CMP98}.

\emph{Equal averages.} Since $M^{j}(W)=V=M^{j}(U_{\mathfrak{B}})$ on $\Lambda_j$ for every $j$, the
quantity $B$ of \cite[(1.31)]{B-CMP99a} vanishes, and condition \cite[(1.66)]{B-CMP99a} holds with
$\alpha_1 := 11d^{2}\alpha_0$ by \cite[(1.65)]{B-CMP99a}, whose derivation uses \cite[(1.7)]{B-CMP99a}
only.

\emph{Condition \eqref{9.3:eq-gauge-fixing-a}.} Here $\operatorname{Re}$ of a matrix is its
Hermitian part. On $\Omega_1$ the decomposition
\eqref{9.3:eq-current-defect-a} of Lemma~\ref{9.3:lem-current-defect}(b),
$J_W = \mu + D_W^{*}f + f_0$, gives \eqref{9.3:eq-gauge-fixing-a} for $U=W$ with that $f$ and
\begin{equation*}
j_0 := \mu + f_0 + (i\eta^{2})^{-1}D_W^{*}\bigl(\operatorname{Re}\partial W - 1\bigr).
\end{equation*}
Since $W(\partial p)$ is unitary,
\begin{equation*}
\operatorname{Re}W(\partial p) - 1 = -\tfrac12\,\bigl(W(\partial p)-1\bigr)^{*}\bigl(W(\partial p)-1\bigr),
\end{equation*}
so by
\eqref{9.3:eq-averaged-gauge-1}, and since $D_W^{*}$ contributes a factor $\eta^{-1}$, the last
term is bounded at a point of level $j$ by
\begin{equation*}
C\eta^{-3}\bigl(\varepsilon L^{-2j}\bigr)^{2} = C\varepsilon^{2}\ell^{-3}L^{-j},
\qquad \ell = L^{j}\eta .
\end{equation*}
Together with $|\mu|\le 3B_3\varepsilon\ell^{-3}$ and $|f_0|\le C_f\varepsilon^{2}\ell^{-3}L^{-j}$
from Lemma~\ref{9.3:lem-current-defect}, this bounds $j_0$ at these points. At points of
level $0$ and on $\Lambda_0$ we put $f:=0$ and the whole of
$(i\eta^{2})^{-1}D_W^{*}\partial W$ into $j_0$; there the current bound \cite[(1.9)]{B-CMP99a}
follows from \cite[(1.7)]{B-CMP99a} up to a factor $8d$. Altogether $|j_0|_{(-3)}\le C\varepsilon$.
Finally,
$|f|\le C_f\varepsilon\ell^{-2}L^{-\beta j}$ and $\beta-\epsilon=\theta$ give, at level $j$,
\begin{equation*}
L^{j\epsilon}\ell^{2}|f| \le C_f\varepsilon L^{-\theta j} \le C_f\varepsilon,
\qquad\text{so}\qquad \mathsf{N}_{\epsilon}(f)\le C_f\varepsilon .
\end{equation*}
Thus the hypotheses of Lemma~\ref{9.3:lem-gauge-fixing} on the field $U=W^{\lambda}$ hold with
$\alpha_0 := C\varepsilon$ and $\alpha_1 = 11d^{2}\alpha_0$, where $C$ depends on $(d,L,\theta)$ only;
taking $C\ge B_3$, also
$U_0\in\mathfrak{U}_k(\boldsymbol{\Omega},B_3\varepsilon)\subset\mathfrak{U}_k(\boldsymbol{\Omega},\alpha_0)$.
Since $\epsilon=(1-\theta)/2$, the smallness of $\varepsilon$ gives
$\alpha_0+\alpha_1\le c_1^{\sharp}(d,L,\epsilon)$.

\emph{Output.} Lemma~\ref{9.3:lem-gauge-fixing} furnishes a gauge transformation $u$; the gauge
transformation $\mathsf{u}$ of the statement is $\lambda$ followed by $u^{-1}$, and
$\widehat{W} := W^{\mathsf{u}} = e^{i\eta A}U_{\mathfrak{B}}$, with the conditions
\cite[(1.37), (1.38)]{B-CMP99a}, in the form $R(U_{\mathfrak{B}})D^{*}A=0$ and
$Q_j(U_{\mathfrak{B}},\eta A)=0$, and the bounds
$\ell|A|,\ \ell^{2}|\nabla_{U_{\mathfrak{B}}}A| \le B_1^{\sharp}(1+11d^{2})\alpha_0$. By the
definition of $\|A\|_0$ these bounds give
$\|A\|_0\le B_1^{\sharp}(1+11d^{2})\alpha_0 =: \varepsilon_2$, which is of the form $C\varepsilon$
with a constant $C$ depending on $(d,L,\theta)$ only; this is
\eqref{9.3:eq-averaged-gauge-a}.
\end{proof}

\begin{lemma}[The stability equation for the averaged minimiser]\label{9.3:lem-stability-equation}
Let $U_{\mathfrak{B}}=U_{\mathfrak{B}}(V)$, the gauge transformation $\mathsf{u}$, the configuration
$\widehat{W}=W^{\mathsf{u}}=e^{i\eta A}U_{\mathfrak{B}}$ and the potential $A$, with $RD^{*}A=0$,
$Q_j(U_{\mathfrak{B}},\eta A)=0$ and $\|A\|_0\le\varepsilon_2=C\varepsilon$, be as in
\eqref{9.3:eq-averaged-gauge-a}. Here $\|\cdot\|_0$ is the norm \eqref{9.3:eq-paired-variational-b}
at $\theta=0$. Let $f$ and $f_0$ be the defect and the remainder of \eqref{9.3:eq-current-defect-a}.

Let $\mathsf{A}(A'):=A'-HD(A')$ be the parametrisation of the slice in \cite[(47)]{B-CMP102b}, where
$H$ is the operator and $D(\cdot)$ the map by which that display parametrises the slice; the map
$D(\cdot)$ is not the covariant derivative $D_U$.

By \cite[Proposition~3]{B-CMP102b} there is a potential $A'_W$ with $A=\mathsf{A}(A'_W)$ and
$\|A'_W\|_0\le 2\varepsilon_2$. It satisfies $QA'_W=0$ and $RD^{*}A'_W=0$ by
\cite[(48), (76)]{B-CMP102b}, where $Q$ stands for the linearised constraint maps $Q_j$ of
\cite[(44)]{B-CMP102b}.

For $A'$ in the Landau slice $\{QA'=0,\ RD^{*}A'=0\}$, the functional $V(A')$ is defined by the expansion
\begin{equation}\label{9.3:eq-stability-equation-1}
A^{\eta}\bigl(e^{i\eta\mathsf{A}(A')}U_{\mathfrak{B}}\bigr)
=A^{\eta}(U_{\mathfrak{B}})+\langle A',J_{U_{\mathfrak{B}}}\rangle
+\tfrac12\langle A',\Delta_1A'\rangle+V(A')
\end{equation}
of \cite[(74), (81)]{B-CMP102b}. In this expansion $\Delta_1$ is Ba{\l}aban's fluctuation operator,
and the functional $V(\cdot)$ is distinct from the data $V$. We put $V':=\delta V/\delta A'$.

Then there are a plaquette field $f_2$ and a bond field $f_3$ such that, with
$\mathfrak{S}:=D_{U_{\mathfrak{B}}}^{*}f_2+f_3$,
\begin{equation}\label{9.3:eq-stability-equation-a}
A'_W=-\mathfrak{G}(U_{\mathfrak{B}})\bigl[V'(A'_W)-\mathfrak{S}\bigr],
\end{equation}
\begin{equation}\label{9.3:eq-stability-equation-2}
|f_2|\le 2|f|,\qquad
|f_3|_{(-3-\theta)}\le C\bigl(\varepsilon_2|f|_{(-2-\theta)}+|f_0|_{(-3-\theta)}\bigr).
\end{equation}
Here $\mathfrak{G}(U_{\mathfrak{B}})=G_1\mathfrak{P}^{*}$ is Ba{\l}aban's propagator on the slice,
formed from the operators $G_1$ and $\mathfrak{P}$ of \cite[(105)--(111)]{B-CMP102b}. Here and
below, $C$ denotes a constant, not necessarily the same at each occurrence.
\end{lemma}

\begin{proof}
Throughout, a hat denotes transport by the gauge transformation $\mathsf{u}$. Since
$\widehat{W}=W^{\mathsf{u}}$, on $\Omega_1$ the decomposition \eqref{9.3:eq-current-defect-a} of
Lemma~\ref{9.3:lem-current-defect} transports to
\begin{equation*}
J_{\widehat{W}}=\hat{\mu}+D_{\widehat{W}}^{*}\hat{f}+\hat{f}_0 .
\end{equation*}

\emph{First computation of a derivative.} Let $\delta A'$ satisfy $Q\delta A'=0$ and $RD^{*}\delta A'=0$.
Both conditions are linear, so the line $s\mapsto A'_W+s\delta A'$ stays in the Landau slice. On this
slice the expansion \eqref{9.3:eq-stability-equation-1} holds. Put
\begin{equation*}
\widehat{W}_s:=e^{i\eta\mathsf{A}(A'_W+s\delta A')}U_{\mathfrak{B}} .
\end{equation*}
Since $\mathsf{A}(A'_W)=A$, we have $\widehat{W}_0=\widehat{W}$. By \cite[(170)]{B-CMP102b},
$\langle\delta A',J_{U_{\mathfrak{B}}}\rangle=0$. Hence the linear term of
\eqref{9.3:eq-stability-equation-1} contributes nothing, and the derivative of
$s\mapsto A^{\eta}(\widehat{W}_s)$ at $s=0$ is
\begin{equation}\label{9.3:eq-stability-equation-3}
\frac{d}{ds}\Big|_{s=0}A^{\eta}(\widehat{W}_s)
=\bigl\langle\delta A',\,\Delta_1A'_W+V'(A'_W)\bigr\rangle .
\end{equation}

\emph{Second computation of the same derivative.} Let $\mathfrak{D}(A'_W)$ be the derivative of the map
$D(\cdot)$ at $A'_W$. Define the operator $\mathsf{g}$, acting bond by bond, by
\begin{equation*}
\mathsf{g}(b):=\int_0^1 e^{it\eta\operatorname{ad}A(b)}\,dt .
\end{equation*}
With $\tilde{\delta}:=\mathsf{g}\bigl(\delta A'-H\mathfrak{D}(A'_W)\delta A'\bigr)$ we have
\begin{equation*}
\widehat{W}_s=e^{is\eta\tilde{\delta}+O(s^2)}\,\widehat{W}.
\end{equation*}
The same derivative is therefore the first variation of the action at $\widehat{W}$ in the direction
$\tilde{\delta}$. By the transported decomposition above, it reads
\begin{equation}\label{9.3:eq-stability-equation-4}
\frac{d}{ds}\Big|_{s=0}A^{\eta}(\widehat{W}_s)
=\langle\tilde{\delta},J_{\widehat{W}}\rangle
=\langle\tilde{\delta},\hat{\mu}\rangle+\langle D_{\widehat{W}}\tilde{\delta},\hat{f}\rangle
+\langle\tilde{\delta},\hat{f}_0\rangle .
\end{equation}

\emph{The multiplier term vanishes.} By \cite[(48)]{B-CMP102b}, for every $s$ and every $j$,
\begin{equation*}
Q_j\bigl(U_{\mathfrak{B}},\eta\mathsf{A}(A'_W+s\delta A')\bigr)
=L^{j}\eta\,Q_j(A'_W+s\delta A')=0 .
\end{equation*}
Hence, by \cite[(1.30)--(1.31)]{B-CMP99a}, each $\widehat{W}_s$ is gauge equivalent to a configuration
whose averages are $V$, that is, to a configuration $U$ with $M^{j}(U)=V$ on $\Lambda_j$ for all $j$.
Consequently $\tilde{\delta}$ is the sum of two vectors:
\begin{itemize}
\item a vector tangent at $\widehat{W}$ to the set of configurations with these averages;
\item a gauge direction.
\end{itemize}
By Lemma~\ref{9.3:lem-current-defect}(a), transported by $\mathsf{u}$, the multiplier $\hat{\mu}$
annihilates both. It annihilates the tangent vector because that vector satisfies the linearised
constraints $Q_j(\widehat{W})a=0$ on $\Lambda_j$ for all $j$. It annihilates the gauge direction because
$\langle D_{\widehat{W}}\lambda,\hat{\mu}\rangle=0$ for all $\lambda$. Therefore
$\langle\tilde{\delta},\hat{\mu}\rangle=0$.

\emph{The equation.} Since $\tilde{\delta}$ depends linearly on $\delta A'$, we may define the bond
field $\mathfrak{S}$ by
\begin{equation}\label{9.3:eq-stability-equation-5}
\langle\delta A',\mathfrak{S}\rangle
:=\langle D_{\widehat{W}}\tilde{\delta},\hat{f}\rangle+\langle\tilde{\delta},\hat{f}_0\rangle .
\end{equation}
Equating \eqref{9.3:eq-stability-equation-3} with \eqref{9.3:eq-stability-equation-4}, and using
$\langle\tilde{\delta},\hat{\mu}\rangle=0$, we obtain
\begin{equation*}
\bigl\langle\delta A',\,\Delta_1A'_W+V'(A'_W)-\mathfrak{S}\bigr\rangle=0
\end{equation*}
for every $\delta A'$ in the tangent space $\{Q\delta A'=0,\ RD^{*}\delta A'=0\}$ of the slice. The
argument of \cite[(105)--(111)]{B-CMP102b} applies to this identity without change. In that argument
the tangent vectors are written $\delta A'=\mathfrak{P}\delta A$, and $G_1\mathfrak{P}^{*}=\mathfrak{G}$.
The argument yields \eqref{9.3:eq-stability-equation-a}.

\emph{Form of $\mathfrak{S}$.} We read off the structure of \eqref{9.3:eq-stability-equation-5} from the
following four facts.

First, the operator $D_{\widehat{W}}-D_{U_{\mathfrak{B}}}$ is of order zero and of size $O(|A|)$.

Second, at a bond $b$ of level $j$, where $\ell=L^{j}\eta$ is the local scale at $b$ and hence
$\eta\ell^{-1}=L^{-j}$, the bound $\|A\|_0\le\varepsilon_2$ gives
\begin{equation*}
|\mathsf{g}-1|\le\eta|A|\le\varepsilon_2L^{-j}.
\end{equation*}

Third, crudely,
\begin{equation*}
\bigl|D_{\widehat{W}}\bigl((\mathsf{g}-1)\varphi\bigr)\bigr|
\le 4\eta^{-1}|\mathsf{g}-1|\,|\varphi|
\le 4\varepsilon_2\ell^{-1}|\varphi| .
\end{equation*}

Fourth, the non-local piece $\langle D_{\widehat{W}}\mathsf{g}H\mathfrak{D}(A'_W)\delta A',\hat{f}\rangle$
is bounded by \cite[(46), (73)]{B-CMP102b} and \cite[(3.133)]{B-CMP99b}:
\begin{equation*}
\bigl|\langle D_{\widehat{W}}\mathsf{g}H\mathfrak{D}(A'_W)\delta A',\hat{f}\rangle\bigr|
\le\sum_b\eta^{d}|\delta A'(b)|\cdot C\varepsilon_2\ell^{-1}
\sum_{y'}e^{-\frac12\delta_0 d(b,y')}\sup_{\widetilde{\Delta}(y')}|f| .
\end{equation*}
Here $y'$ runs over the points indexing the enlarged cubes $\widetilde{\Delta}(y')$ of
\cite[(3.44), (3.47)]{B-CMP99b}, and $d(\cdot,\cdot)$ is the multi-scale distance of \cite{B-CMP96}.

These are the bounds from which $\mathfrak{S}$ takes the form $D_{U_{\mathfrak{B}}}^{*}f_2+f_3$ with
$f_2$ and $f_3$ satisfying \eqref{9.3:eq-stability-equation-2}. Namely, replacing in
$\langle D_{\widehat{W}}\tilde{\delta},\hat{f}\rangle$ the operator $D_{\widehat{W}}$ by
$D_{U_{\mathfrak{B}}}$ (first fact), the operator $\mathsf{g}$ by $1$ (second fact), and omitting the
non-local piece, leaves
$\langle D_{U_{\mathfrak{B}}}\delta A',\hat{f}\rangle=\langle\delta A',D_{U_{\mathfrak{B}}}^{*}\hat{f}\rangle$;
this is the term from which $f_2$ is read off. The field $f_3$ collects the three corrections bounded
in the first, third and fourth facts, together with the term $\langle\tilde{\delta},\hat{f}_0\rangle$.
\end{proof}

\begin{proposition}[Ba{\l}aban's second-variation and data-response bounds; printed without proof
in {\cite[(189)--(190), p.~308]{B-CMP102b}}]\label{9.3:prop-second-variation}
Let $\eta=L^{-k}$, let $\boldsymbol{\Omega}=(\Omega_0\supset\dots\supset\Omega_k)$ be an admissible
sequence of domains in $T_{\eta}$, with layers $\Lambda_j$, as in
Notation~\ref{9.3:not-paired-variational}. The bounds below are stated in
\cite[(189)--(190), p.~308]{B-CMP102b} without proof; cited as printed.

\emph{The second variation.} Write $V(A')$ for the higher-order functional of the fine potential
$A'$ in the slice parametrisation $A'-HD(A')$ of \cite[(47)]{B-CMP102b}, in which $H$ is an
operator and $D(A')$ a functional of $A'$ fixed there; it satisfies $V'(0)=0$. The letter $V$ in
\cite[(189)]{B-CMP102b} denotes this functional, not the data of
Notation~\ref{9.3:not-paired-variational}. Write
$\delta^2/\delta A'^2$ for its second variation. For $y\in\Lambda_j$, let
$\Delta(y)$ be the cube of \cite{B-CMP102b} of side $L^j\eta$ indexed by $y$, the same symbol
standing for restriction to it, and $\widetilde{\Delta}(y')$ the enlarged
cube indexed by $y'$. Let $|\cdot|_{(\gamma)}$ be the level-weighted sup norm of
\cite[(3.41)]{B-CMP99b}, $d(y,y')$ the multi-scale distance of \cite{B-CMP96} (the exponent $d$
in the second display below is the dimension, $d=4$), $\delta_0$
Ba{\l}aban's decay rate, and $\varepsilon_3$ the small parameter of \cite{B-CMP102b} that enters
the bound quoted below for potentials $A'$ satisfying \cite[(77)]{B-CMP102b}.
The bound \cite[(189)]{B-CMP102b} reads:
\begin{quote}
\begin{multline*}
\Bigl|\Delta(y)\frac{\delta^2}{\delta A'^2}V(A')\mathfrak{A}\Bigr|
\le O(1)\varepsilon_3(L^j\eta)^{-3}\exp\bigl(-\tfrac14\delta_0 d(y,y')\bigr)\\
\max\bigl\{|\mathfrak{A}|_{(-1)},|\nabla\mathfrak{A}|_{(-2)}\bigr\}
\end{multline*}
for $\operatorname{supp}\mathfrak{A}\subset\widetilde{\Delta}(y')$, $A'$ satisfying (77),
$y\in\Lambda_j$
\end{quote}
The equation number inside the quotation is that of \cite{B-CMP102b}: the condition on $A'$ is
\cite[(77)]{B-CMP102b}. Here $\mathfrak{A}$ is a test form on the fine lattice.

\emph{The response to the data.} Let $\tilde{\delta}$ denote the variation of the minimal
configuration $U_{\mathfrak{B}}(V_0)$ induced by a variation of the data $V_0$ on $\mathfrak{B}$
at a bond $c$ of $\mathfrak{B}$ at the point $y'$, of level $j'$. The first entry of
\cite[(190)]{B-CMP102b}, for $y\in\Lambda_j$, bounds the size of
$\tilde{\delta}$ at $y$; the printed right-hand side reads
\begin{quote}
\begin{equation*}
\le O(1)(L^j\eta)^{-1}(L^{j'}\eta)^{-d}\exp\bigl(-\tfrac18\delta_0 d(y,y')\bigr)
\end{equation*}
\end{quote}
Of these bounds \cite{B-CMP102b} writes:
\begin{quote}
We do not perform these calculations here
\end{quote}
\end{proposition}

\begin{theorem}[Index stationarity of the minimal configurations]\label{9.3:thm-stationarity-minimisers}
Let $\theta\in(0,1)$. There are $a_A>0$ and $C_A<\infty$, depending on $(d,L,\theta)$ only, such
that the following holds for all $k\ge 1$, all pairs of problems at consecutive cutoffs as in
\eqref{9.3:eq-paired-variational-c}, and all data $V$ satisfying
\eqref{9.3:eq-paired-variational-a} with $\varepsilon\le a_A$. Let $W:=M(U')$ be the one-step
average of the finer minimal configuration $U'$ of the pair. Then there is a gauge
transformation $\mathsf{u}$ on $\Omega_0$ such that
\begin{equation}\label{9.3:eq-stationarity-minimisers-1}
\begin{gathered}
W^{\mathsf{u}}=\exp(i\eta A)\,U_{\mathfrak{B}}(V),\qquad R(U_{\mathfrak{B}})D^{*}A=0,\qquad
Q_j(U_{\mathfrak{B}},\eta A)=0\ \text{on }\Lambda_j,\ 0\le j\le k,\\
A=0\ \text{on }\Lambda_0,\qquad \|A\|_{\theta}\le C_A\,\varepsilon\,\eta^{\theta},
\end{gathered}
\end{equation}
where $\|\cdot\|_{\theta}$ is the norm \eqref{9.3:eq-paired-variational-b}.
\end{theorem}

\begin{proof}
Throughout, $C$, $C_1$, $C_2$, $C_3$, $C_4$ denote constants depending on $(d,L,\theta)$ only,
and $\beta=(1+\theta)/2$, $\epsilon=(1-\theta)/2$ are the exponents of
Notation~\ref{9.3:not-minimiser-bounds}, so that $\theta=\beta-\epsilon$.

\emph{The gauge transformation and the stability equation.} For $\varepsilon$ small, depending
on $(d,L,\theta)$, the averaged minimiser in the gauge-fixed slice,
\eqref{9.3:eq-averaged-gauge-a}, obtained from Lemma~\ref{9.3:lem-gauge-fixing}, supplies a
gauge transformation $\mathsf{u}$ on $\Omega_0$ with
$\widehat{W}:=W^{\mathsf{u}}=e^{i\eta A}U_{\mathfrak{B}}$, $R(U_{\mathfrak{B}})D^{*}A=0$,
$Q_j(U_{\mathfrak{B}},\eta A)=0$ on $\Lambda_j$ for every $j$, and $\|A\|_0\le\varepsilon_2$,
where $\|\cdot\|_0$ denotes the norm \eqref{9.3:eq-paired-variational-b} with $\theta=0$ and
$\varepsilon_2:=C_3\varepsilon$, with $C_3=C_3(d,L,\theta)$ the constant of
\eqref{9.3:eq-averaged-gauge-a}. Write $A=\mathsf{A}(A'_W)=A'_W-HD(A'_W)$ in the slice
parametrisation \cite[(47)]{B-CMP102b}, with $\|A'_W\|_0\le2\varepsilon_2$, as in the setting of
Lemma~\ref{9.3:lem-stability-equation}; here $V(A')$ is the higher-order part of the action in
that parametrisation and $V'=\delta V/\delta A'$ (not to be confused with the data $V$). By
Lemma~\ref{9.3:lem-stability-equation}, $A'_W$ satisfies the stability equation
\eqref{9.3:eq-stability-equation-a},
\begin{equation*}
A'_W=-\mathfrak{G}\bigl[V'(A'_W)-\mathfrak{S}\bigr],\qquad
\mathfrak{G}:=\mathfrak{G}(U_{\mathfrak{B}}),\qquad
\mathfrak{S}=D_{U_{\mathfrak{B}}}^{*}f_2+f_3,
\end{equation*}
with $|f_2|\le 2|f|$ and
$|f_3|_{(-3-\theta)}\le C(\varepsilon_2|f|_{(-2-\theta)}+|f_0|_{(-3-\theta)})$; here $f$ and
$f_0$ are the defect and the remainder of \eqref{9.3:eq-current-defect-a}, which obey
$|f|\le C_f\varepsilon\ell^{-2}L^{-\beta j}$ and
$|f_0|\le C_f\varepsilon^{2}\ell^{-3}L^{-j}$ at every level $j$. It remains to prove that
$A=0$ on $\Lambda_0$ and the bound on $\|A\|_{\theta}$. This is done in four steps.

\emph{Step (1): the source term.} By \cite[Theorem~3.13]{B-CMP99b}, $\mathfrak{G}$ obeys the
bounds \cite[(3.42)--(3.47)]{B-CMP99b}, with the exception of the Laplace inequality. We use
\cite[(3.47)]{B-CMP99b} at $\gamma=-3-\theta$ and at $\gamma=-2-\theta$, both of which lie in
$[-4,4]$; with $B_0$ the constant there, this gives
\begin{equation}\label{9.3:eq-stationarity-minimisers-2}
|\mathfrak{G}f_3|_{(-1-\theta)},\ |\nabla\mathfrak{G}f_3|_{(-2-\theta)}
\le B_0|f_3|_{(-3-\theta)},\qquad
|\mathfrak{G}D^{*}f_2|_{(-1-\theta)}\le B_0|f_2|_{(-2-\theta)} .
\end{equation}
The right-hand sides are bounded by unwinding the definition of $|\cdot|_{(\gamma)}$. At level
$j$ one has $\ell=L^{j}\eta$ and $L^{-\theta j}=\eta^{\theta}(L^{j}\eta)^{-\theta}$, so
\begin{equation*}
(L^{j}\eta)^{2+\theta}|f|\le C_f\varepsilon(L^{j}\eta)^{\theta}L^{-\beta j}
=C_f\varepsilon\eta^{\theta}L^{(\theta-\beta)j}\le C_f\varepsilon\eta^{\theta},
\end{equation*}
that is, $|f|_{(-2-\theta)}\le C_f\varepsilon\eta^{\theta}$, and in the same way
\begin{equation*}
(L^{j}\eta)^{3+\theta}|f_0|\le C_f\varepsilon^{2}(L^{j}\eta)^{\theta}L^{-j}
=C_f\varepsilon^{2}\eta^{\theta}L^{(\theta-1)j}\le C_f\varepsilon^{2}\eta^{\theta}
\end{equation*}
at every level $j$, that is, $|f_0|_{(-3-\theta)}\le C_f\varepsilon^2\eta^{\theta}$.
Hence $|f_2|_{(-2-\theta)}\le 2C_f\varepsilon\eta^{\theta}$ and, since
$\varepsilon\le a_A\le1$ (we may and do take $a_A\le1$), $|f_3|_{(-3-\theta)}\le C\varepsilon\eta^{\theta}$.

For the gradient of $\mathfrak{G}D^{*}f_2$ we use \cite[(3.44)]{B-CMP99b} exactly as in the
estimate \eqref{9.3:eq-gauge-fixing-b} in the proof of Lemma~\ref{9.3:lem-gauge-fixing}: with
the partition of unity $\{\zeta_{y'}\}$, $\zeta_{y'}$ supported in the enlarged cube
$\widetilde{\Delta}(y')$ of a point $y'$ of level $j'$ and $\xi'=L^{-j'}$, the lattice H\"older
norm $\|\cdot\|^{\xi'}_{\epsilon}$ is bounded by the supremum as there, and the bounds
$|f_2|\le2|f|$ and $|f|\le C_f\varepsilon\ell^{-2}L^{-\beta j}$ give
\begin{equation*}
\begin{split}
\|\zeta_{y'}f_2\|^{\xi'}_{\epsilon}+|\zeta_{y'}f_2|
&\le 14L^{j'\epsilon}C_f\varepsilon(L^{j'}\eta)^{-2}L^{-\beta j'}\\
&=14C_f\varepsilon\eta^{\theta}(L^{j'}\eta)^{-2-\theta},
\end{split}
\end{equation*}
the equality because $L^{j'(\epsilon-\beta)}=L^{-\theta j'}=\eta^{\theta}(L^{j'}\eta)^{-\theta}$.
By \cite[(3.44)]{B-CMP99b}, at a point $x$ of level $j$ lying in the cube of $y$,
\begin{equation*}
|\nabla\mathfrak{G}D^{*}f_2|(x)\le\sum_{y'}B_0'(\epsilon)e^{-\delta_0d(y,y')}
\bigl(\|\zeta_{y'}f_2\|^{\xi'}_{\epsilon}+|\zeta_{y'}f_2|\bigr).
\end{equation*}
The ratio of the weights satisfies
$(L^{j'}\eta)^{-2-\theta}(L^{j}\eta)^{2+\theta}\le L^{3|j-j'|}$. For $y$ of level $j$ and $y'$
of level $j'$, admissibility of $\boldsymbol{\Omega}$
(Notation~\ref{9.3:not-paired-variational}) gives $d(y,y')\ge RM_1(|j-j'|-1)$, so half of the
factor $e^{-\delta_0d(y,y')}$ absorbs this ratio: by the floor
\eqref{9.3:eq-paired-variational-d}, $\delta_0RM_1\ge16\log L$, one has, for $|j-j'|\ge1$,
\begin{equation*}
e^{-\frac12\delta_0RM_1(|j-j'|-1)}L^{3|j-j'|}\le L^{-8(|j-j'|-1)}L^{3|j-j'|}
=L^{8-5|j-j'|}\le L^{3}.
\end{equation*}
The other half of the factor is summed over $y'$ by \cite[Lemma~2.1]{B-CMP96}.
Therefore $|\nabla\mathfrak{G}D^{*}f_2|_{(-2-\theta)}\le C\varepsilon\eta^{\theta}$. Collecting
these bounds with \eqref{9.3:eq-stationarity-minimisers-2} and the definition
\eqref{9.3:eq-paired-variational-b} of $\|\cdot\|_{\theta}$,
\begin{equation}\label{9.3:eq-stationarity-minimisers-3}
\|\mathfrak{G}\mathfrak{S}\|_{\theta}\le C_1\varepsilon\eta^{\theta}.
\end{equation}

\emph{Step (2): the quadratic term.} One has $V'(0)=0$ and
\begin{equation*}
V'(A')=\int_0^1\frac{\delta^2V}{\delta A'^2}(tA')\,A'\,dt,\qquad
A'=\sum_{y'}\zeta_{y'}A'.
\end{equation*}
We apply Proposition~\ref{9.3:prop-second-variation}, printed without proof in
\cite[(189)]{B-CMP102b} and cited as printed: for $y\in\Lambda_j$, a test form $\mathfrak{A}$ with
$\operatorname{supp}\mathfrak{A}\subset\widetilde{\Delta}(y')$, and $A'$ satisfying
\cite[(77)]{B-CMP102b}, the second derivative of $V$ restricted to the cube $\Delta(y)$ obeys
\begin{equation*}
\Bigl|\Delta(y)\frac{\delta^2}{\delta A'^2}V(A')\,\mathfrak{A}\Bigr|
\le O(1)\,\varepsilon_3(L^{j}\eta)^{-3}e^{-\frac14\delta_0d(y,y')}
\max\bigl\{|\mathfrak{A}|_{(-1)},|\nabla\mathfrak{A}|_{(-2)}\bigr\},
\end{equation*}
$\varepsilon_3$ being the smallness parameter of that bound. For
$\mathfrak{A}=\zeta_{y'}A'$ one has
\begin{equation*}
\max\bigl\{|\mathfrak{A}|_{(-1)},|\nabla\mathfrak{A}|_{(-2)}\bigr\}
\le C(L^{j'}\eta)^{-\theta}\|A'\|_{\theta}.
\end{equation*}
Here $A'=tA'_W$, $0\le t\le1$, satisfies $\|tA'_W\|_0\le2\varepsilon_2$, and the bound is applied
with $\varepsilon_3$ replaced by a constant multiple of $\varepsilon_2$; the ratio
$(L^{j'}\eta)^{-\theta}(L^{j}\eta)^{\theta}\le L^{\theta|j-j'|}$ of the weights is absorbed, and
the sum over $y'$ controlled, by \cite[Lemma~2.1]{B-CMP96} and the floor
\eqref{9.3:eq-paired-variational-d} as in Step~(1).
Inserting this with $A'=tA'_W$, $0\le t\le1$, into the integral representation of $V'$ and
summing over $y'$ gives
\begin{equation}\label{9.3:eq-stationarity-minimisers-4}
|V'(A'_W)|_{(-3-\theta)}\le C_2\varepsilon_2\|A'_W\|_{\theta}.
\end{equation}

\emph{Step (3): absorption.} The lattice $T_\eta$ is finite, so $\|A'_W\|_{\theta}<\infty$. By
the stability equation, the bounds \eqref{9.3:eq-stationarity-minimisers-2} applied to
$V'(A'_W)$ in place of $f_3$, \eqref{9.3:eq-stationarity-minimisers-4} and
\eqref{9.3:eq-stationarity-minimisers-3},
\begin{equation*}
\|A'_W\|_{\theta}\le B_0|V'(A'_W)|_{(-3-\theta)}+\|\mathfrak{G}\mathfrak{S}\|_{\theta}
\le B_0C_2\varepsilon_2\|A'_W\|_{\theta}+C_1\varepsilon\eta^{\theta}
\le 2C_1\varepsilon\eta^{\theta},
\end{equation*}
the last inequality once $B_0C_2\varepsilon_2\le\frac12$. Since $\varepsilon_2=C_3\varepsilon$,
this is a condition $\varepsilon\le c$ with $c=c(d,L,\theta)>0$; this condition, together with
the smallness thresholds of the averaged-gauge step \eqref{9.3:eq-averaged-gauge-a}, of
\cite[Theorem~3.13]{B-CMP99b} and of the results of \cite{B-CMP102b} used above, and with
$a_A\le1$, defines $a_A$.

\emph{Step (4): the correction $HD$.} One has
$D(A')=\int_0^1\mathfrak{D}(tA')A'\,dt$, with $\mathfrak{D}$ the derivative of $D$ in the
notation of \cite{B-CMP102b}. By \cite[(73)]{B-CMP102b} and \cite[(46)]{B-CMP102b},
$\|HD(A'_W)\|_{\theta}\le C_4\varepsilon_2\|A'_W\|_{\theta}$. Hence
\begin{equation*}
\|A\|_{\theta}\le\|A'_W\|_{\theta}+\|HD(A'_W)\|_{\theta}
\le(1+C_4\varepsilon_2)\,2C_1\varepsilon\eta^{\theta}\le C_A\varepsilon\eta^{\theta},
\end{equation*}
with $C_A:=2C_1(1+C_4C_3a_A)$, since $\varepsilon_2=C_3\varepsilon\le C_3a_A$. Finally, the
condition $A=0$ on $\Lambda_0$ is the constraint $Q_0A=0$, that is, the case $j=0$ of
$Q_j(U_{\mathfrak{B}},\eta A)=0$ on $\Lambda_j$.

When the domains $\Omega_j$ are all equal to $T_\eta$, the weights in the norms
$|\cdot|_{(\gamma)}$ are constants, and \cite[(98)]{B-CMP102b} takes the place of
\cite[(189)]{B-CMP102b} in Step~(2); the rest of the argument is unchanged.
\end{proof}

\begin{corollary}[Stationarity of level data, curvature and action]
\label{9.3:cor-stationarity-level-data}
Let $\theta\in(0,1)$, $k\ge1$, the pair, the data $V$ satisfying the condition $(7)_{\varepsilon}$
with $\varepsilon\le a_A$, the gauge transformation $\mathsf{u}$ on $\Omega_0$ and the
potential $A$ be as in Theorem~\ref{9.3:thm-stationarity-minimisers}. Let $U'$, the minimal
configuration at the finer depth on $T_{\eta'}$, and $W:=M(U')$ be as in that theorem, and let
$U_{\mathfrak{B}}=U_{\mathfrak{B}}(V)$ be the minimiser on $T_{\eta}$.
At a point $x$ write $j=j(x)$ and $\ell=L^{j}\eta$. In (ii)--(iv), $C<\infty$ denotes a constant
depending on $(d,L,\theta)$ only.
\begin{enumerate}
\item[(i)] \emph{Level data.} For $i\le j(x)$ there is a potential $A^{(i)}$ with
\begin{equation*}
\bigl(M^{i+1}(U')\bigr)^{\mathsf{u}\restriction T_{L^{i}\eta}}
=\exp\bigl(\mathrm{i}L^{i}\eta A^{(i)}\bigr)\,M^{i}(U_{\mathfrak{B}}),
\qquad
\ell\,|A^{(i)}|\le 2C_A\varepsilon L^{-\theta j},
\end{equation*}
where $\mathsf{u}\restriction T_{L^{i}\eta}$ is the restriction of $\mathsf{u}$ to the sites of the
lattice $T_{L^{i}\eta}$ of spacing $L^{i}\eta$.
\item[(ii)] \emph{Curvature.} For every plaquette $p$ of $T_\eta$
and every plaquette $p'$ of $T_{\eta'}$ parallel to $p$ and within distance $\eta$ of it,
\begin{equation*}
\bigl|R(\cdot)\mathfrak{F}_{U'}(p')-\mathfrak{F}_{U_{\mathfrak{B}}}(p)\bigr|
\le C\varepsilon\ell^{-2}L^{-\theta j},
\end{equation*}
where $R(\cdot)$ is the adjoint rotation by the value of $\mathsf{u}$ at the site at which the loop
$\partial p$ begins, so that $\mathfrak{F}_{W^{\mathsf{u}}}(p)=R(\cdot)\mathfrak{F}_{W}(p)$.
\item[(iii)] \emph{Action, relative form.} For $y\in\Lambda_j$ let $\Delta(y)$ be the cube of side
$\ell$ at $y$, so that $|\Delta(y)|=\ell^{d}=\ell^{4}$, and for a configuration $U$ on a lattice of
spacing $a$ let $A^{a}(U)|_{\Delta(y)}$ be the Wilson action of
Definition~\ref{1.1:def-wilson-action} on that lattice, each plaquette weighted by $a^{d-4}$ and
restricted to $\Delta(y)$, that is, the sum of
$a^{d-4}\bigl(1-\operatorname{Re}\operatorname{tr}U(\partial p)\bigr)$
over the plaquettes $p$ of that lattice contained in $\Delta(y)$. Then
\begin{equation*}
\bigl|A^{\eta'}(U')|_{\Delta(y)}-A^{\eta}(U_{\mathfrak{B}})|_{\Delta(y)}\bigr|
\le C\varepsilon^{2}L^{-\theta j}.
\end{equation*}
\item[(iv)] \emph{Local backgrounds.} The conclusion of
Theorem~\ref{9.3:thm-stationarity-minimisers} holds at levels $j\ge2$ for the pair
$U_{k+1,\Omega}(V)$, $U_{k,\Omega}(V)$ of local backgrounds of \cite[(2.13)--(2.16)]{B-CMP119}. In
particular the test functionals
\begin{equation*}
t_k(V):=\sup_{p}\eta^{-2}\bigl|U_{k,\square}(V,\partial p)-1\bigr|,
\end{equation*}
the supremum running over the cubes $\square$ and the plaquettes $p$ on which the local
backgrounds $U_{k,\square}(V)$ of \cite[(2.13)--(2.17)]{B-CMP119} are defined,
satisfy $|t_{k+1}(V)-t_k(V)|\le C\varepsilon L^{-\theta k}$.
\end{enumerate}
\end{corollary}

We prove parts (i)--(iii) here. The proof of part (iv) uses the refinement lemma and the
decay lemmas of this chapter and is given after them.

\begin{proof}[Proof of parts \textup{(i)}--\textup{(iii)}]
\emph{Bound on $A$.} By the definition of the level-weighted norms
$|\cdot|_{(-1-\theta)}$ and $|\cdot|_{(-2-\theta)}$ \cite[(3.41)]{B-CMP99b}, the bound
$\|A\|_{\theta}\le C_A\varepsilon\eta^{\theta}$ of
Theorem~\ref{9.3:thm-stationarity-minimisers} gives, with $\ell$ and $j$ read at the point
considered,
\begin{equation}\label{9.3:eq-stationarity-level-data-1}
\ell\,|A|\le C_A\varepsilon(\eta/\ell)^{\theta}=C_A\varepsilon L^{-\theta j},
\qquad
\ell^{2}\,|\nabla_{U_{\mathfrak{B}}}A|\le C_A\varepsilon L^{-\theta j}.
\end{equation}

\emph{Part (i).} Since $M^{i+1}=M^{i}\circ M$, we have $M^{i+1}(U')=M^{i}(W)$. Ba{\l}aban's
averaging maps are gauge covariant \cite{B-CMP98}, so
$M^{i}(W)^{\mathsf{u}\restriction T_{L^{i}\eta}}=M^{i}(W^{\mathsf{u}})$, and
$W^{\mathsf{u}}=\exp(\mathrm{i}\eta A)U_{\mathfrak{B}}$ by
Theorem~\ref{9.3:thm-stationarity-minimisers}. The constraint map $Q_i(U_{\mathfrak{B}},\cdot)$
expresses $M^{i}(\exp(\mathrm{i}\eta A)U_{\mathfrak{B}})$ as
$\exp\bigl(\mathrm{i}Q_i(U_{\mathfrak{B}},\eta A)\bigr)M^{i}(U_{\mathfrak{B}})$, and by
\cite[Proposition~4]{B-CMP98}, which is \cite[(44)]{B-CMP102b}, $Q_i(U_{\mathfrak{B}},\eta A)$ is
the sum of its linear part $L^{i}\eta\,Q_iA$ and a remainder, with
\begin{equation*}
|Q_iA|\le\sup|A|,
\qquad
\bigl|Q_i(U_{\mathfrak{B}},\eta A)-L^{i}\eta\,Q_iA\bigr|\le C'\,(L^{i}\eta|A|)^{2},
\end{equation*}
where $C'$ is the constant of \cite[Proposition~4]{B-CMP98} and the supremum runs over the block
on which the average at $x$ depends. We put
$A^{(i)}:=(L^{i}\eta)^{-1}Q_i(U_{\mathfrak{B}},\eta A)$; this is the identity of
part (i). By \eqref{9.3:eq-stationarity-level-data-1} and $L^{i}\eta\le\ell$ (as $i\le j$),
\begin{equation*}
\ell\,|A^{(i)}|\le \ell\sup|A|+C'\frac{L^{i}\eta}{\ell}\,(\ell\sup|A|)^{2}
\le C_A\varepsilon L^{-\theta j}\bigl(1+C'C_A\varepsilon\bigr),
\end{equation*}
which is at most $2C_A\varepsilon L^{-\theta j}$ when $C'C_A\varepsilon\le1$. After decreasing
$a_A$ if necessary, which leaves Theorem~\ref{9.3:thm-stationarity-minimisers} valid, we may
assume $C'C_Aa_A\le1$.

\emph{Part (ii).} Put $\widehat{W}:=W^{\mathsf{u}}=\exp(\mathrm{i}\eta A)U_{\mathfrak{B}}$. By gauge
covariance of the curvature, $\mathfrak{F}_{\widehat{W}}(p)=R(\cdot)\mathfrak{F}_W(p)$. Let $g(p)$ be
the normalised mean of the fine curvature introduced in the proof of
Lemma~\ref{9.3:lem-current-defect}, a mean of $\mathfrak{F}_{U'}$ over plaquettes parallel to $p$
within $2d\eta$ of it. Since $R(\cdot)$ is an isometry, we split
\begin{align*}
\bigl|R(\cdot)\mathfrak{F}_{U'}(p')-\mathfrak{F}_{U_{\mathfrak{B}}}(p)\bigr|
&\le\bigl|\mathfrak{F}_{U'}(p')-g(p)\bigr|+\bigl|g(p)-\mathfrak{F}_{W}(p)\bigr|\\
&\quad+\bigl|\mathfrak{F}_{\widehat{W}}(p)-\mathfrak{F}_{U_{\mathfrak{B}}}(p)\bigr|.
\end{align*}
The first term is bounded by the oscillation of $\mathfrak{F}_{U'}$ over distances of order $\eta$,
together with the cost of the parallel transports, as established in the proof of
Lemma~\ref{9.3:lem-current-defect}: at most $C\varepsilon\ell^{-2}L^{-\beta j}$ with
$\beta=(1+\theta)/2$. The second term is $|f(p)|$, where $f:=\mathfrak{F}_W-g$ is the defect of
Lemma~\ref{9.3:lem-current-defect}; by its bound in \eqref{9.3:eq-current-defect-a} it is at most
$C_f\varepsilon\ell^{-2}L^{-\beta j}$. Since $\beta\ge\theta$, both are at most
$C\varepsilon\ell^{-2}L^{-\theta j}$. For the third term, \cite[(1.47)]{B-CMP99a} gives
\begin{equation*}
\mathfrak{F}_{\widehat{W}}-\mathfrak{F}_{U_{\mathfrak{B}}}
=D_{U_{\mathfrak{B}}}A+O\bigl(|A|^{2}+|A|\,\eta\,|\mathfrak{F}_{U_{\mathfrak{B}}}|\bigr).
\end{equation*}
Here $|D_{U_{\mathfrak{B}}}A|\le C|\nabla_{U_{\mathfrak{B}}}A|\le CC_A\varepsilon\ell^{-2}L^{-\theta j}$
and $|A|^{2}\le(C_A\varepsilon)^{2}\ell^{-2}L^{-2\theta j}$ by
\eqref{9.3:eq-stationarity-level-data-1}. By \cite[Theorem~1]{B-CMP102b}, $U_{\mathfrak{B}}$ lies
in the class $\mathfrak{U}_k(\boldsymbol{\Omega},B_3\varepsilon)$ of regular configurations, so
$|\mathfrak{F}_{U_{\mathfrak{B}}}|\le C\varepsilon\ell^{-2}$, and
$|A|\eta|\mathfrak{F}_{U_{\mathfrak{B}}}|\le CC_A\varepsilon^{2}\ell^{-2}L^{-\theta j}(\eta/\ell)$.
Summing the three terms gives (ii).

\emph{Part (iii).} For a configuration $U$ on a lattice of spacing $a$ and a plaquette $p$,
\begin{equation*}
a^{d-4}\bigl(1-\operatorname{Re}\operatorname{tr}U(\partial p)\bigr)
=\tfrac12 a^{d}|\mathfrak{F}_U(p)|^{2}\bigl(1+O(a^{4}|\mathfrak{F}_U(p)|^{2})\bigr).
\end{equation*}
The plaquettes of $T_{\eta'}$ in $\Delta(y)$ are distributed among the plaquettes $p$ of $T_\eta$ in
$\Delta(y)$, $L^{d}$ parallel ones to each $p$, each within $\eta$ of $p$; and
$\eta'^{d}L^{d}=\eta^{d}$. Hence the leading terms of the difference in (iii) equal
\begin{equation*}
\sum_{p\subset\Delta(y)}\tfrac12\eta'^{d}\sum_{p'}
\bigl(|\mathfrak{F}_{U'}(p')|^{2}-|\mathfrak{F}_{U_{\mathfrak{B}}}(p)|^{2}\bigr),
\end{equation*}
where the inner sum runs over the $L^{d}$ plaquettes $p'$ assigned to $p$. As
$|\mathfrak{F}_{U'}(p')|=|R(\cdot)\mathfrak{F}_{U'}(p')|$, part (ii) and
$|\mathfrak{F}_{U_{\mathfrak{B}}}|\le C\varepsilon\ell^{-2}$ give
$\bigl||\mathfrak{F}_{U'}(p')|^{2}-|\mathfrak{F}_{U_{\mathfrak{B}}}(p)|^{2}\bigr|\le
C\varepsilon^{2}\ell^{-4}L^{-\theta j}$. There are at most $\binom{d}{2}|\Delta(y)|\eta^{-d}$
plaquettes $p$ in $\Delta(y)$, so the leading terms contribute at most
$C\varepsilon^{2}L^{-\theta j}|\Delta(y)|\ell^{-4}=C\varepsilon^{2}L^{-\theta j}$, using
$d=4$, $|\Delta(y)|\ell^{-4}=1$. By the same count for the plaquettes of $T_{\eta'}$, and since
$|\mathfrak{F}_{U'}|,|\mathfrak{F}_{U_{\mathfrak{B}}}|\le C\varepsilon\ell^{-2}$ (the first by
part (ii)), the correction terms contribute at most
$C\varepsilon^{4}(\eta/\ell)^{4}|\Delta(y)|\ell^{-4}$, which is at most
$C\varepsilon^{2}L^{-\theta j}$ since $\varepsilon^{4}\le a_A^{2}\varepsilon^{2}$ and
$(\eta/\ell)^{4}=L^{-4j}\le L^{-\theta j}$. This proves (iii).
\end{proof}

\begin{definition}[The perfect action and its coupling-free tube]\label{9.3:not-perfect-action}
We use the data of the paired variational problems at consecutive cutoffs
(Notation~\ref{9.3:not-paired-variational}): the spacing $\eta=L^{-k}$, the layers $\Lambda_j$,
the determining set $\mathfrak{B}=\bigcup_j\Lambda_j$, and, for data $V$ on $\mathfrak{B}$
satisfying \eqref{9.3:eq-paired-variational-a} for some $\varepsilon$, a minimal configuration
$U_{\mathfrak{B}}(V)$ of the action $A^{\eta}$ under the constraints $M^{j}(U)=V$ on $\Lambda_j$.
Let $a_A$ be the threshold of the index stationarity of the minimal configurations
(Theorem~\ref{9.3:thm-stationarity-minimisers}). Write $\mathfrak{g}$ for the Lie algebra of
$\mathrm{SU}(N)$ and $\mathfrak{g}^{c}$ for its complexification.

The \emph{perfect action} is $\mathcal{A}_{\mathfrak{B}}(V):=A^{\eta}(U_{\mathfrak{B}}(V))$.
For a bond $c\subset\Lambda_j$ and $\mathsf{X}\in\mathfrak{g}$, let $\mathsf{X}\delta_c$ be the
function on $\mathfrak{B}$ equal to $\mathsf{X}$ at $c$ and to $0$ elsewhere; the \emph{first
variation} of the perfect action is
\begin{equation*}
\omega_{\mathfrak{B}}(V)(c)[\mathsf{X}]
:=\frac{d}{ds}\,\mathcal{A}_{\mathfrak{B}}\bigl(e^{is\mathsf{X}\delta_c}V\bigr)\Big|_{s=0}.
\end{equation*}
We fix the radius $r_c$ in the display below, with $c_{\gamma}$ the constant fixed in the
calculus of exponentially decaying kernels, \eqref{9.3:eq-kernel-calculus-a}. For real data
$V_0$ on $\mathfrak{B}$ satisfying \eqref{9.3:eq-paired-variational-a} with $\varepsilon=a_A/2$
and for $\varrho>0$, the tube $\mathbb{D}(\varrho)$ about $V_0$ is defined in the same display:
\begin{equation}\label{9.3:eq-perfect-action-a}
\begin{gathered}
r_c:=\min\{a_A/10,\;c_{\gamma}\},\\
\mathbb{D}(\varrho):=\bigl\{V=e^{iB}V_0:\ B \text{ a } \mathfrak{g}^{c}\text{-valued function on }
\mathfrak{B},\ \sup|B|<\varrho\bigr\}.
\end{gathered}
\end{equation}
The radius $r_c$ does not depend on the coupling.
For $V=e^{iB}V_0\in\mathbb{D}(r_c)$ and bonds
$c_1,\dots,c_n$ of $\mathfrak{B}$, the \emph{$n$-th jet} of the perfect action is the iterated
derivative with respect to the values $B(c_1),\dots,B(c_n)$ of the tube coordinate $B$ at the
bonds $c_1,\dots,c_n$:
\begin{equation*}
a^{(n)}_{\mathfrak{B}}(c_1,\dots,c_n;V):=\partial_{B(c_1)}\cdots\partial_{B(c_n)}\,
\mathcal{A}_{\mathfrak{B}}(V).
\end{equation*}
For any quantity built on the determining set $\mathfrak{B}$, the prefix $\delta$ denotes its
value built on the determining set $\mathfrak{B}'$ of the paired problem at the other depth minus
its value built on $\mathfrak{B}$; thus
$\delta\mathcal{A}=\mathcal{A}_{\mathfrak{B}'}-\mathcal{A}_{\mathfrak{B}}$. We write $j_{*}$ for
the minimum of the levels
$j(\cdot)$ of the bonds $c_1,\dots,c_n$, and $\operatorname{diam}$ for the diameter of
$\{c_1,\dots,c_n\}$ measured in the multi-scale distance $d(y,y')$ of \cite{B-CMP96}.
\end{definition}

\begin{lemma}[Analyticity of the perfect action on the tube]
Let $V_0$, $r_c$ and $\mathbb{D}(r_c)$ be as in
Definition~\ref{9.3:not-perfect-action}. The minimal configuration
$V\mapsto U_{\mathfrak{B}}(V)$, the
difference $\mathcal{A}_{\mathfrak{B}}-\mathcal{A}_{\mathfrak{B}}(V_0)$ and the first variation
$\omega_{\mathfrak{B}}$ are analytic on $\mathbb{D}(r_c)$.
\end{lemma}

\begin{proof}
By \cite[Proposition~9]{B-CMP102b}, applied at the real data $V_0$, which satisfy
\eqref{9.3:eq-paired-variational-a} with $\varepsilon=a_A/2$, under the hypothesis
\cite[(172)]{B-CMP102b} of that proposition with the constants written $C_1$, $\varepsilon_1$
there (notation of that paper) satisfying $C_1\varepsilon_1\ge r_c$, the map
$V\mapsto U_{\mathfrak{B}}(V)$ is analytic on $\mathbb{D}(r_c)$.

By definition, $\mathcal{A}_{\mathfrak{B}}(V)$ is the action $A^{\eta}$ evaluated at
$U_{\mathfrak{B}}(V)$. Hence the analyticity of $U_{\mathfrak{B}}$ on $\mathbb{D}(r_c)$ gives the
analyticity of $\mathcal{A}_{\mathfrak{B}}-\mathcal{A}_{\mathfrak{B}}(V_0)$ there. The first
variation $\omega_{\mathfrak{B}}(V)(c)[\mathsf{X}]$ is the derivative of
$\mathcal{A}_{\mathfrak{B}}$ along the one-parameter family $s\mapsto e^{is\mathsf{X}\delta_c}V$
at $s=0$. It is therefore analytic on $\mathbb{D}(r_c)$ as well.
\end{proof}

\begin{lemma}[A two-constants lemma on the disc]\label{9.3:lem-two-constants}
Write $D(0,3)=\{z\in\mathbb{C}:|z|<3\}$. Let $\varphi$ be holomorphic on $D(0,3)$, and let
$0\le\mathsf{m}\le\mathsf{M}$ be constants such that $|\varphi|\le\mathsf{M}$ on $D(0,3)$ and
$|\varphi|\le\mathsf{m}$ on the diameter $\mathopen{]}-3,3\mathclose{[}$. Then
\begin{equation*}
|\varphi(i)|\le(\mathsf{m}\mathsf{M})^{1/2}.
\end{equation*}
\end{lemma}

\begin{proof}
If $\mathsf{m}=0$, then $\varphi$ vanishes on $\mathopen{]}-3,3\mathclose{[}$, so $\varphi\equiv0$ by
the identity theorem and there is nothing to prove. Assume therefore $0<\mathsf{m}\le\mathsf{M}$.

Let $D_+=\{z\in D(0,3):\operatorname{Im}z>0\}$ be the upper half-disc. The function
$\log|\varphi|$ is subharmonic on $D_+$ and bounded above by $\log\mathsf{M}$.

Put $\Phi(z)=\bigl((3+z)/(3-z)\bigr)^2$. The M\"obius map $z\mapsto(3+z)/(3-z)$ sends
$\mathopen{]}-3,3\mathclose{[}$ onto $\mathopen{]}0,\infty\mathclose{[}$. It sends the point
$3e^{it}$, $0<t<\pi$, to $i\cot(t/2)$, a point of the positive imaginary axis, and it sends the
interior point $i$ to $(3+i)/(3-i)=(4+3i)/5$, which lies in the open first quadrant. Hence it maps
$D_+$ onto the open first quadrant, and squaring maps that quadrant onto the upper half-plane.
Thus $\Phi$ maps $D_+$ onto the upper half-plane, the diameter onto
$\mathopen{]}0,\infty\mathclose{[}$, and the open semicircle onto the negative real axis.

Let $\arg$ take values in $[0,\pi]$ on the closed upper half-plane minus $0$, and put
\begin{equation}\label{9.3:eq-two-constants-1}
\varpi(z)=1-\frac{1}{\pi}\arg\Phi(z),\qquad
v(z)=\varpi(z)\log\mathsf{m}+\bigl(1-\varpi(z)\bigr)\log\mathsf{M}.
\end{equation}
The function $\varpi$ is the harmonic measure of the diameter in $D_+$. It is harmonic on $D_+$
and continuous on $\overline{D_+}\setminus\{\pm3\}$, where it equals $1$ on the diameter and $0$
on the semicircle. Consequently $v$ is harmonic on $D_+$ and bounded.

Consider the function
\begin{equation*}
u=\log|\varphi|-v .
\end{equation*}
At every boundary point of $D_+$ other than $\pm3$ we have $\limsup u\le0$: on the diameter
because $\varphi$ is continuous there with $|\varphi|\le\mathsf{m}$, and on the semicircle because
$|\varphi|\le\mathsf{M}$. For $\varepsilon>0$ set
\begin{equation*}
u_\varepsilon(z)=u(z)+\varepsilon\log\frac{|z-3|\,|z+3|}{36}.
\end{equation*}
The added term is harmonic on $D_+$ and is $\le0$ there, since $|z\mp3|\le6$; it tends to
$-\infty$ at $\pm3$. Since $u$ is bounded above, $u_\varepsilon$ is subharmonic with
$\limsup u_\varepsilon\le0$ at every boundary point, $\pm3$ included. The maximum principle
therefore gives $u_\varepsilon\le0$ on $D_+$, and letting $\varepsilon\to0$ gives
$\log|\varphi|\le v$ on $D_+$.

It remains to evaluate at $z=i$. Since $(3+i)/(3-i)=(3+i)^2/10$ has modulus $1$ and argument
$2\arctan(1/3)$, we get $\Phi(i)=e^{4i\arctan(1/3)}$, with
$4\arctan(1/3)\in\mathopen{]}0,\pi\mathclose{[}$.
Hence
\begin{equation*}
\varpi(i)=1-\frac{4}{\pi}\arctan\frac13\approx0.59 .
\end{equation*}
Moreover $\varpi(i)>\frac12$, because $\frac13<\sqrt2-1=\tan(\pi/8)$.

Since $\log\mathsf{m}\le\log\mathsf{M}$ and $\varpi(i)>\frac12$,
\begin{align*}
\log|\varphi(i)|\le v(i)
&=\log\mathsf{M}-\varpi(i)\bigl(\log\mathsf{M}-\log\mathsf{m}\bigr)\\
&\le\log\mathsf{M}-\tfrac12\bigl(\log\mathsf{M}-\log\mathsf{m}\bigr)
=\tfrac12\log(\mathsf{m}\mathsf{M}),
\end{align*}
which is the assertion.
\end{proof}

\begin{theorem}[Stationarity of the perfect action's jets]\label{9.3:thm-stationarity-jets}
Let $\theta\in(0,1)$ and let $a_A$ be the threshold of
Theorem~\ref{9.3:thm-stationarity-minimisers}. Consider paired variational problems at
consecutive cutoffs to which Theorem~\ref{9.3:thm-stationarity-minimisers} applies, with
determining sets $\mathfrak{B}$ and $\mathfrak{B}'$. We use the following objects of
Notation~\ref{9.3:not-perfect-action}:
\begin{itemize}
\item the perfect action $\mathcal{A}_{\mathfrak{B}}$ and its first variation
$\omega_{\mathfrak{B}}$;
\item the tube $\mathbb{D}(\varrho)$ around a real datum $V_0$ that satisfies the regularity
condition \eqref{9.3:eq-paired-variational-a}, $|(\partial V)(p')-1|<\varepsilon$ on
$\mathfrak{B}$, with $\varepsilon=a_A/2$;
\item the radius $r_c$ and the jets $a^{(n)}_{\mathfrak{B}}$;
\item the differences $\delta X:=X_{\mathfrak{B}'}-X_{\mathfrak{B}}$;
\item the minimal level $j_{*}$ of the arguments and the diameter $\operatorname{diam}$ in
the distance $d(\cdot,\cdot)$.
\end{itemize}
For a bond $c\subset\Lambda_j$, $|\delta\omega(V)(c)|$ denotes the norm of the linear form
$\mathsf{X}\mapsto\delta\omega(V)(c)[\mathsf{X}]$ on $\mathfrak{g}$, and
$|a^{(n)}_{\mathfrak{B}}|$, $|\delta a^{(n)}|$ denote the norms of the corresponding
multilinear forms. There are constants $C_B$ and $C_n$ ($n\ge2$), depending on
$(d,L,\theta)$ and $n$ only, such that for every $k\ge1$ and every such pair the following hold.
\begin{itemize}
\item[(a)] Let $V$ be real and satisfy \eqref{9.3:eq-paired-variational-a} with
$\varepsilon\le a_A$. Then for every bond $c\subset\Lambda_j$
\begin{equation*}
|\delta\omega(V)(c)|\le C_B\,\varepsilon\,L^{-\theta j}.
\end{equation*}
\item[(b)] Let $V\in\mathbb{D}(r_c/4)$. Then for every bond $c\subset\Lambda_j$
\begin{equation*}
|\delta\omega(V)(c)|\le C_B\,L^{-\theta j/2}.
\end{equation*}
\item[(c)] Let $V\in\mathbb{D}(r_c/8)$ and $n\ge2$. Then for all bonds
$c_1,\dots,c_n$ of $\mathfrak{B}$, with $\operatorname{diam}$ the diameter of
$\{c_1,\dots,c_n\}$,
\begin{align*}
|a^{(n)}_{\mathfrak{B}}(c_1,\dots,c_n;V)| &\le C_n\,e^{-\delta_1\operatorname{diam}},\\
|\delta a^{(n)}(c_1,\dots,c_n;V)| &\le C_n\,L^{-\theta j_{*}/4}\,
e^{-\delta_1\operatorname{diam}/2}.
\end{align*}
\end{itemize}
\end{theorem}

\begin{proof}
Throughout, $C$ denotes a finite constant depending on $(d,L,\theta)$ only, whose value may
change from one occurrence to the next. The dimension is $d=4$. For a point $x$ we write
$\ell=L^{j(x)}\eta$.

\emph{Proof of (a): the setting.} Fix a real $V$ satisfying
\eqref{9.3:eq-paired-variational-a} with $\varepsilon\le a_A$, a bond $c\subset\Lambda_j$ and
$\mathsf{X}\in\mathfrak{g}$. Let $U'$ be the minimal configuration of the finer problem and let
$W:=M(U')$ be its one-step average.

Theorem~\ref{9.3:thm-stationarity-minimisers} provides a gauge transformation $\mathsf{u}$ and
a potential $A$ with $W^{\mathsf{u}}=e^{i\eta A}U_{\mathfrak{B}}(V)$. We put
$\widehat{W}:=W^{\mathsf{u}}$ and write $U_{\mathfrak{B}}$ for $U_{\mathfrak{B}}(V)$.

As in Lemma~\ref{9.3:lem-stability-equation}, let $A'_W$ be the potential with
$A=\mathsf{A}(A'_W)$, $QA'_W=0$ and $RD^{*}A'_W=0$. Let
$\mathfrak{S}=D^{*}_{U_{\mathfrak{B}}}f_2+f_3$ be the source term of the stability equation
\eqref{9.3:eq-stability-equation-a}.

Let $B_{\mathsf{X}}:=\mathsf{X}\delta_c$ be the $\mathfrak{g}$-valued function on
$\mathfrak{B}$ that equals $\mathsf{X}$ at $c$ and vanishes elsewhere. By the definition of
$\omega_{\mathfrak{B}}$ in Notation~\ref{9.3:not-perfect-action}, $\omega_{\mathfrak{B}}(V)(c)[\mathsf{X}]$
is the derivative at $s=0$ of $\mathcal{A}_{\mathfrak{B}}(e^{isB_{\mathsf{X}}}V)$.

\emph{Proof of (a): the identity.} We show
\begin{equation}\label{9.3:eq-stationarity-jets-1}
\delta\omega(V)(c)[\mathsf{X}]=\big\langle H_1B_{\mathsf{X}},\,V'(A'_W)-\mathfrak{S}\big\rangle .
\end{equation}

Put $\delta A'_{*}:=H_1B_{\mathsf{X}}$, where $H_1=G_1Q^{*}(QG_1Q^{*})^{-1}$ is Ba{\l}aban's
minimiser map \cite[(3.129)]{B-CMP99b}. By the properties of $H_1$ stated in
\cite[(45)]{B-CMP102b}, we have $L^j\eta\,Q_j\delta A'_{*}=B_{\mathsf{X}}$ and
$RD^{*}\delta A'_{*}=0$.

Hence the line $s\mapsto A'_W+s\delta A'_{*}$ stays in the gauge-fixing slice of
Lemma~\ref{9.3:lem-stability-equation}. Recall from the proof of that lemma the identity
$Q_j(U_{\mathfrak{B}},\eta\mathsf{A}(A'))=L^j\eta\,Q_j(A')$. By this identity, the
configurations
\begin{equation*}
\widehat{W}_s:=e^{i\eta\mathsf{A}(A'_W+s\delta A'_{*})}U_{\mathfrak{B}}
\end{equation*}
have, up to a gauge transformation, the averages $e^{isB_{\mathsf{X}}}V$ on $\mathfrak{B}$.

We compute the derivative of $A^{\eta}(\widehat{W}_s)$ at $s=0$ in two ways, as in the proof of
Lemma~\ref{9.3:lem-stability-equation}.

First, on the slice the action has the expansion
\begin{equation*}
A^{\eta}\big(e^{i\eta\mathsf{A}(A')}U_{\mathfrak{B}}\big)
=A^{\eta}(U_{\mathfrak{B}})+\langle A',J_{U_{\mathfrak{B}}}\rangle
+\tfrac12\langle A',\Delta_1A'\rangle+V(A')
\end{equation*}
used there. Differentiating along the line gives
\begin{equation}\label{9.3:eq-stationarity-jets-2}
\frac{d}{ds}\Big|_{s=0}A^{\eta}(\widehat{W}_s)
=\big\langle\delta A'_{*},\,J_{U_{\mathfrak{B}}}+\Delta_1A'_W+V'(A'_W)\big\rangle .
\end{equation}
Here, unlike in Lemma~\ref{9.3:lem-stability-equation}, the term
$\langle\delta A'_{*},J_{U_{\mathfrak{B}}}\rangle$ need not vanish, since
$Q\delta A'_{*}\ne0$.

Second, we have $\widehat{W}_s=e^{is\eta\tilde{\delta}_{*}+O(s^2)}\widehat{W}$, with
\begin{equation*}
\tilde{\delta}_{*}:=\mathsf{g}\big(\delta A'_{*}-H\mathfrak{D}(A'_W)\delta A'_{*}\big),
\end{equation*}
where $\mathsf{g}$ and $\mathfrak{D}$ are as in that proof. So the same derivative equals
$\langle\tilde{\delta}_{*},J_{\widehat{W}}\rangle$.

Lemma~\ref{9.3:lem-current-defect}(b), transported by $\mathsf{u}$ (indicated by hats), gives
$J_{\widehat{W}}=\widehat{\mu}+D^{*}_{\widehat{W}}\widehat{f}+\widehat{f}_0$. Therefore the
derivative equals
\begin{equation*}
\langle\tilde{\delta}_{*},\widehat{\mu}\rangle
+\langle D_{\widehat{W}}\tilde{\delta}_{*},\widehat{f}\rangle
+\langle\tilde{\delta}_{*},\widehat{f}_0\rangle
=\langle\tilde{\delta}_{*},\widehat{\mu}\rangle+\langle\delta A'_{*},\mathfrak{S}\rangle .
\end{equation*}
The last equality is the definition of $\mathfrak{S}$ in the proof of
Lemma~\ref{9.3:lem-stability-equation}: the expression defining it is linear in the
variation, here evaluated at $\delta A'_{*}$.

We now identify the terms of \eqref{9.3:eq-stationarity-jets-2} and of the second expression.
\begin{itemize}
\item The term $\langle H_1B_{\mathsf{X}},J_{U_{\mathfrak{B}}}\rangle$ is the first-order term
of the expansion \cite[(2.8)]{B-CMP109}. Hence it equals
$\omega_{\mathfrak{B}}(V)(c)[\mathsf{X}]$.
\item The term $\langle H_1B_{\mathsf{X}},\Delta_1A'_W\rangle$ vanishes by the defining identity
of $H_1$ in \cite{B-CMP102b}, because $QA'_W=0=RD^{*}A'_W$.
\item Transporting the definition of $\mu$ in Lemma~\ref{9.3:lem-current-defect} by
$\mathsf{u}$ gives
$\langle\tilde{\delta}_{*},\widehat{\mu}\rangle
=\langle\widehat{E}\tilde{\delta}_{*},J_{U'}\rangle$. This equals
$\omega_{\mathfrak{B}'}(V)(c)[\mathsf{X}]$. Indeed, two fine variations that induce the same
variation of the averages on $\mathfrak{B}'$ differ by an element of $\ker Q(U')$ plus a gauge
direction, and $J_{U'}$ vanishes on both.
\end{itemize}

Equating the two expressions for the derivative gives
\begin{equation*}
\omega_{\mathfrak{B}}(V)(c)[\mathsf{X}]+\langle H_1B_{\mathsf{X}},V'(A'_W)\rangle
=\omega_{\mathfrak{B}'}(V)(c)[\mathsf{X}]+\langle H_1B_{\mathsf{X}},\mathfrak{S}\rangle .
\end{equation*}
This is \eqref{9.3:eq-stationarity-jets-1}.

\emph{Proof of (a): the bound.} By \eqref{9.3:eq-stationarity-jets-1} and the definition of the
adjoint,
\begin{equation*}
|\delta\omega(V)(c)[\mathsf{X}]|\le\mathrm{(I)}+\mathrm{(II)}+\mathrm{(III)},
\end{equation*}
where
\begin{equation*}
\mathrm{(I)}:=|\langle H_1B_{\mathsf{X}},V'(A'_W)\rangle|,\quad
\mathrm{(II)}:=|\langle D_{U_{\mathfrak{B}}}H_1B_{\mathsf{X}},f_2\rangle|,\quad
\mathrm{(III)}:=|\langle H_1B_{\mathsf{X}},f_3\rangle| .
\end{equation*}

We collect the pointwise bounds needed.
\begin{itemize}
\item The bound \cite[(3.133)]{B-CMP99b}, applied to $H_1$, gives
\begin{align*}
|H_1B_{\mathsf{X}}|(x)&\le C\ell^{-1}e^{-\frac12\delta_1 d(y,c)}|\mathsf{X}|,\\
|\nabla H_1B_{\mathsf{X}}|(x)&\le C\ell^{-2}e^{-\frac12\delta_1 d(y,c)}|\mathsf{X}|
\end{align*}
for $x$ in the block $\Delta(y)$. The second bound also controls
$D_{U_{\mathfrak{B}}}H_1B_{\mathsf{X}}$.
\item By Proposition~\ref{9.3:prop-second-variation} (printed without proof in
\cite[p.~308]{B-CMP102b}; cited as printed),
$|V'(A'_W)|_{(-3-\theta)}\le C\varepsilon_2\|A'_W\|_{\theta}$, where $\varepsilon_2$ is the
fixed multiple of $\varepsilon$ of Lemma~\ref{9.3:lem-stability-equation}. By the proof of
Theorem~\ref{9.3:thm-stationarity-minimisers}, $\|A'_W\|_{\theta}\le C\varepsilon\eta^{\theta}$.
Reading the level-weighted norm at a point, we get
\begin{equation*}
|V'(A'_W)|(x)\le C\varepsilon_2\varepsilon\,\eta^{\theta}\ell^{-3-\theta}.
\end{equation*}
\item Lemma~\ref{9.3:lem-stability-equation} gives $|f_2|\le2|f|$ and
\begin{equation*}
|f_3|_{(-3-\theta)}\le C\big(\varepsilon_2|f|_{(-2-\theta)}+|f_0|_{(-3-\theta)}\big).
\end{equation*}
\item Lemma~\ref{9.3:lem-current-defect}(b) gives
$|f|(x)\le C_f\varepsilon\ell^{-2}L^{-\beta j(x)}$ and
$|f_0|(x)\le C_f\varepsilon^2\ell^{-3}L^{-j(x)}$.
\end{itemize}

Since $\beta=(1+\theta)/2\ge\theta$, $\theta<1$ and $\eta^{\theta}\ell^{-\theta}=L^{-\theta j(x)}$,
we have
\begin{equation*}
L^{-\beta j(x)}\le\eta^{\theta}\ell^{-\theta},\qquad L^{-j(x)}\le\eta^{\theta}\ell^{-\theta}.
\end{equation*}
Hence
\begin{equation*}
|f_2|(x)\le2C_f\varepsilon\,\eta^{\theta}\ell^{-2-\theta},\qquad
|f_3|(x)\le C(\varepsilon_2+\varepsilon)\varepsilon\,\eta^{\theta}\ell^{-3-\theta}.
\end{equation*}

The pairings are sums over points with weight $\eta^{d}$. We group the points into the blocks
$\Delta(y)$, on each of which $\ell$ is constant and $\sum_{x\in\Delta(y)}\eta^{d}=\ell^{d}$.
With $d=4$ the powers of $\ell$ combine in each term as follows:
\begin{itemize}
\item in $\mathrm{(I)}$ and $\mathrm{(III)}$: $\ell^{4}\cdot\ell^{-1}\cdot\ell^{-3-\theta}=\ell^{-\theta}$;
\item in $\mathrm{(II)}$: $\ell^{4}\cdot\ell^{-2}\cdot\ell^{-2-\theta}=\ell^{-\theta}$.
\end{itemize}
Since $\eta^{\theta}\ell^{-\theta}=L^{-\theta j(y)}$, it follows that
\begin{equation*}
\mathrm{(I)}+\mathrm{(II)}+\mathrm{(III)}\le C\varepsilon\,(1+\varepsilon_2+\varepsilon)\,|\mathsf{X}|
\sum_{y}e^{-\frac12\delta_1 d(y,c)}L^{-\theta j(y)} .
\end{equation*}
By the summation lemma \cite[Lemma~2.1]{B-CMP96}, the sum over $y$ is bounded by
$CL^{-\theta j}$, where $j$ is the level of $c$. Since $\varepsilon\le a_A$ and
$\varepsilon_2$ is a fixed multiple of $\varepsilon$, the prefactor $1+\varepsilon_2+\varepsilon$
is bounded by a constant. This gives (a).

\emph{Proof of (b).} Let $V=e^{iB}V_0\in\mathbb{D}(r_c)$. Write
$U=e^{i\eta\mathcal{H}(B)}U_{\mathfrak{B}}(V_0)$ with the response map $\mathcal{H}$ of
\cite[Proposition~9]{B-CMP102b}. Then
\begin{equation*}
\partial\mathcal{A}_{\mathfrak{B}}/\partial B(c)=\langle\tilde{\delta},J_U\rangle,
\end{equation*}
where $\tilde{\delta}$ is the variation of $U$ induced by $\partial_{B(c)}$. By the same
proposition, the response map $\mathcal{H}$ obeys \cite[(19)--(21)]{B-CMP102b}, with smallness
parameter $B_5\varepsilon_1$ given by the constants $B_5$, $\varepsilon_1$ of that proposition.
Consequently $|J_U|\le C\ell^{-3}$.

As in the statement, $j$ is the level of $c$; for a block $y$ put $\ell(y):=L^{j(y)}\eta$. By
the first entry of \cite[(190)]{B-CMP102b}, in the form included in
Proposition~\ref{9.3:prop-second-variation} (printed without proof in
\cite[p.~308]{B-CMP102b}; cited as printed), the induced variation is bounded at a point of a
block $y$ by
\begin{equation*}
O(1)(L^{j(y)}\eta)^{-1}(L^{j}\eta)^{-d}\exp\!\big(-\tfrac18\delta_0 d(y,c)\big).
\end{equation*}
Inserting these bounds into the representation of the derivative gives, for $d=4$,
\begin{equation*}
|\omega_{\mathfrak{B}}|\le(L^{j}\eta)^{d}\sum_{y}\ell(y)^{d}\,\ell(y)^{-1}(L^{j}\eta)^{-d}
e^{-\frac18\delta_0 d(y,c)}\,C\ell(y)^{-3}\le C_{\omega}
\end{equation*}
on $\mathbb{D}(r_c)$. The same bound holds at the depth $\mathfrak{B}'$.

Now let $V=e^{i(B_1+iB_2)}V_0\in\mathbb{D}(r_c/4)$, with $B_1,B_2$ $\mathfrak{g}$-valued. Fix
$\mathsf{X}$ with $|\mathsf{X}|=1$ and put
\begin{equation*}
\varphi(z):=\delta\omega\big(e^{i(B_1+zB_2)}V_0\big)(c)[\mathsf{X}].
\end{equation*}
For $|z|<3$ we have $\sup|B_1+zB_2|<r_c/4+3r_c/4=r_c$. Hence $\varphi$ is holomorphic on
$D(0,3)$, by the analyticity on $\mathbb{D}(r_c)$ recorded in
Notation~\ref{9.3:not-perfect-action}, and $|\varphi|\le2C_{\omega}$ there.

For real $z\in\,]-3,3[$ the datum $e^{i(B_1+zB_2)}V_0$ is real and satisfies
\eqref{9.3:eq-paired-variational-a} with $\varepsilon=a_A/2+5r_c$. Since $r_c\le a_A/10$, we have
$a_A/2+5r_c\le a_A$. So (a) gives $|\varphi(z)|\le C_Ba_AL^{-\theta j}$ for real $z$.

We apply the two-constants lemma (Lemma~\ref{9.3:lem-two-constants}) with $\mathsf{M}:=2C_{\omega}$
and $\mathsf{m}:=\min\{C_Ba_AL^{-\theta j},2C_{\omega}\}$. It yields
\begin{equation*}
|\delta\omega(V)(c)[\mathsf{X}]|=|\varphi(i)|
\le(2C_{\omega}C_Ba_A)^{1/2}L^{-\theta j/2}.
\end{equation*}
After enlarging $C_B$ if necessary, which keeps (a) valid, this is (b).

\emph{Proof of (c): the second jet.} By \cite[(2.8)]{B-CMP109} and
\cite[(174)--(177)]{B-CMP102b}, the second-order coefficient of
$B\mapsto\mathcal{A}_{\mathfrak{B}}(e^{iB}V)$ is $\frac12\langle H_1B,\Delta_1H_1B\rangle$.

Let $a$ be the constant in the averaging term $Q^{*}aQ$ of
$G_1^{-1}=\Delta_1+DRD^{*}+Q^{*}aQ$. By the formula $H_1=G_1Q^{*}(QG_1Q^{*})^{-1}$, $QH_1$ is
the identity and $G_1^{-1}H_1=Q^{*}(QG_1Q^{*})^{-1}$. Moreover $RD^{*}H_1=0$. Therefore
\begin{align*}
\langle H_1B,\Delta_1H_1B\rangle
&=\langle H_1B,G_1^{-1}H_1B\rangle-\langle H_1B,DRD^{*}H_1B\rangle
-\langle QH_1B,aQH_1B\rangle\\
&=\big\langle B,\big[(QG_1Q^{*})^{-1}-a\big]B\big\rangle .
\end{align*}

The bound \cite[(3.132)]{B-CMP99b} is stated there also for the operator with $G_1$ in place of
$G$. It is valid on the complex class \cite[(3.37)--(3.38)]{B-CMP99b}, by
\cite[Theorems~3.4 and~3.12]{B-CMP99b}. It gives
\begin{equation*}
|a^{(2)}_{\mathfrak{B}}(c,c';V)|\le C_2\,e^{-\delta_1 d(c,c')}\qquad\text{on }\mathbb{D}(r_c).
\end{equation*}
The scale factors are handled as follows.
\begin{itemize}
\item For $d=4$ the scale factors cancel when $c$ and $c'$ are of equal level.
\item Otherwise they are absorbed by \cite[Lemma~2.1]{B-CMP96}.
\end{itemize}
The exponential coordinates $B$ differ from the left-invariant
coordinates in which the coefficient above is computed by bond-diagonal terms from the
derivative of the exponential map, multiplied by $\omega_{\mathfrak{B}}$. These terms do not
affect the decay.

\emph{Proof of (c): higher jets.} Let $n>2$. The form $a^{(n)}_{\mathfrak{B}}$ is symmetric in
its arguments, since partial derivatives commute. Choose two of the bonds $c_1,\dots,c_n$,
$c'$ and $c''$ say, with $d(c',c'')=\operatorname{diam}$. Then
$a^{(n)}_{\mathfrak{B}}(c_1,\dots,c_n;V)$ is the derivative of
$a^{(2)}_{\mathfrak{B}}(c',c'';\cdot)$ in the remaining $n-2$ variables, evaluated at $V$.

For $V\in\mathbb{D}(r_c/8)$, a shift of each of these variables by at most $r_c/2$ stays in
$\mathbb{D}(r_c)$, since $r_c/8+r_c/2<r_c$. The Cauchy inequalities on polydiscs of radius
$r_c/2$ therefore give
\begin{equation*}
|a^{(n)}_{\mathfrak{B}}(c_1,\dots,c_n;V)|\le C_n e^{-\delta_1\operatorname{diam}}.
\end{equation*}
The same argument applies at the depth $\mathfrak{B}'$.

\emph{Proof of (c): the differences.} The difference $\delta a^{(n)}$ is an $(n-1)$-st
derivative of $\delta\omega$. For $V\in\mathbb{D}(r_c/8)$, polydiscs of radius $r_c/8$ in the
$n-1$ variables stay in $\mathbb{D}(r_c/4)$, where (b) holds. The Cauchy inequalities and (b)
give
\begin{equation*}
|\delta a^{(n)}|\le C_nL^{-\theta j_{*}/2},
\end{equation*}
since every argument has level at least $j_{*}$. By the bound just proved, applied at both
depths, we also have $|\delta a^{(n)}|\le2C_ne^{-\delta_1\operatorname{diam}}$.

Using $\min\{x,y\}\le(xy)^{1/2}$ and enlarging $C_n$, we obtain
\begin{equation*}
|\delta a^{(n)}|\le C_nL^{-\theta j_{*}/4}e^{-\delta_1\operatorname{diam}/2}.
\end{equation*}
\end{proof}

\begin{remark}[Flat jets on the whole lattice]\label{9.3:rem-flat-jets}
Take the trivial data $V_0=1$ and the admissible sequence of domains with
$\Omega_j=T_{\eta}$ for $0\le j\le k$. For $n\ge2$ write $T_n^{(k)}$ for the $n$-th jet at $V_0=1$
(the \emph{flat jet}) of the perfect action $\mathcal{A}_{\mathfrak{B}}$ of
Theorem~\ref{9.3:thm-stationarity-jets}, where at depth $k$ the determining set $\mathfrak{B}$ is
the whole level-$k$ lattice. For a kernel
$\mathsf{K}$ of $n$ bond arguments let $\operatorname{tree}$ be the length, in the distance
$d(\cdot,\cdot)$, of a shortest tree joining the arguments, and put
$\|\mathsf{K}\|_{\delta}=\sup e^{\delta\operatorname{tree}}|\mathsf{K}|$, the supremum over all
$n$-tuples of bond arguments. Then, for every $\delta>0$ with $(n-1)\delta\le\delta_1/2$, part (c)
of Theorem~\ref{9.3:thm-stationarity-jets} yields $\sup_k\|T_n^{(k)}\|_{\delta}\le C_n$ and
$\|T_n^{(k+1)}-T_n^{(k)}\|_{\delta}\le C_nL^{-\theta k/4}$, so that all flat jets $T_n^{(k)}$
converge geometrically in $k$ in this norm, for every averaging mean
admissible in the sense of \cite{B-CMP109}.
\end{remark}

\begin{proof}
Since $V_0=1$ is $e^{iB}V_0$ with $B=0$, it lies in $\mathbb{D}(r_c/8)$, so part (c) of
Theorem~\ref{9.3:thm-stationarity-jets} applies. With $\Omega_j=T_{\eta}$ for $0\le j\le k$, every
difference $\Omega_j\setminus\Omega_{j+1}$ with $j<k$ is empty; every bond therefore has level $k$,
$\mathfrak{B}=\Lambda_k$ is the whole level-$k$ lattice, which is the lattice of unit spacing and the
same at every depth, so that $T_n^{(k)}$ and $T_n^{(k+1)}$ are kernels on the same bonds; moreover
$a^{(n)}_{\mathfrak{B}}=T_n^{(k)}$, while $j_{*}=k$. The
first bound of part (c), whose constants $C_n$ do not depend on $k$, gives
$|T_n^{(k)}|\le C_ne^{-\delta_1\operatorname{diam}}$, and the second gives
\begin{equation}\label{9.3:eq-flat-jets-1}
|T_n^{(k+1)}-T_n^{(k)}|\le C_nL^{-\theta k/4}e^{-\delta_1\operatorname{diam}/2}.
\end{equation}
For $n$ points, $\operatorname{diam}\le\operatorname{tree}\le(n-1)\operatorname{diam}$. For the
first inequality, the tree contains a path between two points at distance
$\operatorname{diam}$. For the second, the star joining one point to the other $n-1$ has $n-1$
edges, each of length at most $\operatorname{diam}$. Hence, for $(n-1)\delta\le\delta_1/2$,
\begin{equation*}
\sup_k\|T_n^{(k)}\|_{\delta}\le C_n
\qquad\text{and}\qquad
\|T_n^{(k+1)}-T_n^{(k)}\|_{\delta}\le C_nL^{-\theta k/4}.
\end{equation*}
Since $L^{-\theta/4}<1$, the second bound is summable in $k$. Thus $(T_n^{(k)})_k$ is Cauchy and
converges geometrically. Since the argument uses no property of the mean beyond the hypotheses of
Theorem~\ref{9.3:thm-stationarity-jets}, it holds for every admissible mean.
\end{proof}

\begin{lemma}[One-step data as functionals of the jets]\label{9.3:lem-one-step-data}
Consider the step $k\to k+1$ of Ba{\l}aban's small-field expansion in \cite{B-CMP109}, with
data $V^{(k)}$ and the operator $\Delta^{(k)}$ of the quadratic form of the step as in
\cite[(2.3)]{B-CMP109}, and let
$\mathcal{A}_{\mathfrak{B}}$, $\omega_{\mathfrak{B}}$ and $a^{(n)}_{\mathfrak{B}}$ be the perfect
action, its first variation and its $n$-th jets on the determining set $\mathfrak{B}$ of depth $k$
(Notation~\ref{9.3:not-perfect-action}). Let $h$ and $C^{(2)}(V;\cdot)$ be the one-step objects of
the change of variables of the step in \cite[\S2]{B-CMP109}, $G^{(2)}(V;\cdot)$ the one-step
quadratic term of the same expansion, $C(V)$ the
eliminating operator of \cite[\S2]{B-CMP109}, $\mathfrak{j}(V)$ the Jacobian attached to the
constraint $\delta(QB')$, and $\mathbb{H}_k(B')$ the non-linear response of \cite[(2.6)]{B-CMP109}.
For $V$ in the tube $\mathbb{D}(r_c)$ of Notation~\ref{9.3:not-perfect-action} and
$\mathfrak{g}^{c}$-valued bond functions $B$ on $\mathfrak{B}$, define the quadratic form
$\Delta^{[\mathfrak{B}]}(V)$ by
\begin{equation}\label{9.3:eq-one-step-data-a}
\begin{split}
\langle B,\Delta^{[\mathfrak{B}]}(V)B\rangle
:={}& a^{(2)}_{\mathfrak{B}}(V)[B,B]
-2\,\omega_{\mathfrak{B}}(V)\bigl[hC^{(2)}(V;B)\bigr]\\
&+G^{(2)}(V;B).
\end{split}
\end{equation}
Here $a^{(2)}_{\mathfrak{B}}(V)[B,B]$ is the sum over bonds $c_1,c_2$ of
$a^{(2)}_{\mathfrak{B}}(c_1,c_2;V)$ applied to $B(c_1),B(c_2)$, and, for a $\mathfrak{g}^{c}$-valued
bond function $Y$ on $\mathfrak{B}$, $\omega_{\mathfrak{B}}(V)[Y]$ is the sum over bonds $c$ of
$\omega_{\mathfrak{B}}(V)(c)[Y(c)]$, extended complex-linearly in $Y(c)$.
Then:
\begin{itemize}
\item[(a)] $\Delta^{(k)}=\Delta^{[\mathfrak{B}]}(V^{(k)})$.
\item[(b)] $h$, $C^{(2)}$, $G^{(2)}$, $C(V)$ and $\mathfrak{j}(V)$ are one-step objects: each has
range at most two $L$-blocks, is analytic in $V$, and does not depend on the index $k$. The
non-linear response $\mathbb{H}_k(B')$ is $U_k$ regarded as a function of the data $V$.
\item[(c)] Consequently every coupling-free datum of the step $k\to k+1$ is a functional of
$(\omega_{\mathfrak{B}},a^{(n)}_{\mathfrak{B}})$ given by one-step objects, the same for every $k$.
\end{itemize}
\end{lemma}

\begin{proof}
(a) We begin with \cite[(2.8)]{B-CMP109}. It expands the action of the minimal configuration with
data $V'V^{(k)}$, where $V'$ is the fluctuation datum parametrised by $B'$, as
\begin{equation*}
\begin{split}
A\bigl(U_k(V'V^{(k)})\bigr)={}&A(U_{k+1})+\langle H_1B',J\rangle\\
&+\tfrac12\langle H_1B',\Delta_1H_1B'\rangle+V(H_1B').
\end{split}
\end{equation*}
Here $H_1=G_1Q^{*}(QG_1Q^{*})^{-1}$ is the minimiser map, $J$ the current, $\Delta_1$ the
fluctuation operator appearing in \cite[(2.8)]{B-CMP109}, and $V(\cdot)$ the higher-order
functional of \cite{B-CMP109}; the last
symbol is not the datum $V$. In \cite{B-CMP109}, $U_k(V)$ is the minimal configuration with
prescribed averages $V$ on the determining set of depth $k$, that is, $U_{\mathfrak{B}}(V)$ of
Notation~\ref{9.3:not-perfect-action}. The left side is therefore the perfect action
$\mathcal{A}_{\mathfrak{B}}$ at the datum $V'V^{(k)}$, and its expansion around $V^{(k)}$ is read
through the jets $\omega_{\mathfrak{B}}(V^{(k)})$ and $a^{(n)}_{\mathfrak{B}}(V^{(k)})$.

Next, \cite[\S2]{B-CMP109} makes the change of variables
\begin{equation*}
B'=B-hD(B),
\end{equation*}
where $D$ is the non-linear map of that change and $D^{(2)}$, its part of second order in $B$,
equals $C^{(2)}$. Further, \cite[(2.11)]{B-CMP109} collects the terms of order zero in $g_k$, and
\cite{B-CMP109} states in the same place that the quadratic form in $B'$ among these terms equals
$-\frac12\langle B',\Delta^{(k)}B'\rangle$.

The right side of~\eqref{9.3:eq-one-step-data-a} at $V=V^{(k)}$ is built from the same
ingredients; its three terms are:
\begin{itemize}
\item the second jet $a^{(2)}_{\mathfrak{B}}(V^{(k)})$ of the perfect action;
\item the first variation $\omega_{\mathfrak{B}}(V^{(k)})$ composed with the second-order part
$hC^{(2)}(V^{(k)};B)$ of the change of variables, which enters the right side
of~\eqref{9.3:eq-one-step-data-a} with the factor $-2$;
\item the remaining one-step quadratic term $G^{(2)}(V^{(k)};B)$.
\end{itemize}
From the expansion \cite[(2.8)]{B-CMP109} of the perfect action at $V'V^{(k)}$, the change of
variables $B'=B-hD(B)$ with $D^{(2)}=C^{(2)}$, the identification of the quadratic form of order
zero in $g_k$ in \cite[(2.11)]{B-CMP109}, and \cite[(1.5)]{B-CMP109}, we obtain
$\Delta^{(k)}=\Delta^{[\mathfrak{B}]}(V^{(k)})$.

(b) The operators $h$, $C^{(2)}$ and $G^{(2)}$ are those of \cite[\S2]{B-CMP109}, and $C(V)$ is
the eliminating operator of \cite[\S2]{B-CMP109}, which is almost local in the sense of
\cite{B-CMP99b}. These operators and the Jacobian $\mathfrak{j}(V)$ of $\delta(QB')$ are the
one-step objects of the step: each has range at most two $L$-blocks, is analytic in $V$, and
carries no index $k$. The non-linear response $\mathbb{H}_k(B')$ of
\cite[(2.6)]{B-CMP109} is $U_k$ written as a function of $V$.

(c) The coupling-free data of the step $k\to k+1$ are built, in \cite[\S2]{B-CMP109}, from the
objects listed in (b) and from the non-linear response $\mathbb{H}_k(B')$. By (a), the quadratic
form of the step is the functional~\eqref{9.3:eq-one-step-data-a} of
$\omega_{\mathfrak{B}}$ and $a^{(2)}_{\mathfrak{B}}$, evaluated at $V^{(k)}$, with one-step
coefficients. By (b), the objects $h$, $C^{(2)}$, $G^{(2)}$, $C(V)$, $\mathfrak{j}(V)$ are
one-step objects that do not depend on $k$, and the response $\mathbb{H}_k(B')$, being $U_k$ as a
function of $V$, enters only through the perfect action $\mathcal{A}_{\mathfrak{B}}$ and its jets
$a^{(n)}_{\mathfrak{B}}$. Hence
every coupling-free datum of the step $k\to k+1$ is a functional of
$(\omega_{\mathfrak{B}},a^{(n)}_{\mathfrak{B}})$ given by one-step objects that do not depend
on $k$.
\end{proof}

\begin{lemma}[Calculus of exponentially decaying kernels]\label{9.3:lem-kernel-calculus}
Let $\mathcal{B}$ be a finite set of bonds of the unit lattice: all bonds of the torus, or the
bonds of a subset of it. Write $|b-b'|$ for the distance of the bonds $b,b'$, that is, the
Euclidean distance of their midpoints, taken periodic on the torus. Fix a finite-dimensional
complex inner product space. A \emph{kernel} is a family
$\mathsf{K}=(\mathsf{K}(b,b'))_{b,b'\in\mathcal{B}}$ of linear maps of that space. It acts on
the functions on $\mathcal{B}$ with values in that space, which carry the inner product
$\langle u_1,u_2\rangle:=\sum_b\langle u_1(b),u_2(b)\rangle$, and $\mathsf{K}_1\mathsf{K}_2$ has
kernel $\sum_{b''}\mathsf{K}_1(b,b'')\mathsf{K}_2(b'',b')$. For $\delta>0$ put
\begin{equation*}
\|\mathsf{K}\|_{\delta}:=\sup_{b,b'}|\mathsf{K}(b,b')|\,e^{\delta|b-b'|},\qquad
c_{\delta}:=\sup_{b}\sum_{b'}e^{-\delta|b-b'|/2},
\end{equation*}
with $|\cdot|$ the operator norm on that space. In $c_{\delta}$ the supremum and the sum run over
the bonds of $\mathbb{Z}^4$, and $|b-b'|$ is the Euclidean distance of the midpoints.
We write $\operatorname{Re}T\ge\gamma_0$ if
$\operatorname{Re}\langle u,Tu\rangle\ge\gamma_0\|u\|^2$ for all $u$. Let
$\mathsf{K}_1,\mathsf{K}_2$ be kernels, let $T$ be a kernel with $\operatorname{Re}T\ge\gamma_0>0$
and $\|T\|_{\delta}\le\Gamma$, let $\lambda\ge0$, and let $T'$ be a kernel with $T'+\lambda$
invertible; write $\delta T:=T'-T$ and
$\delta(T+\lambda)^{-1}:=(T'+\lambda)^{-1}-(T+\lambda)^{-1}$. Then
\begin{equation}\label{9.3:eq-kernel-calculus-a}
\begin{gathered}
(\mathrm{K}1)\qquad \|\mathsf{K}_1\mathsf{K}_2\|_{\delta/2}\le c_{\delta}\|\mathsf{K}_1\|_{\delta}
\|\mathsf{K}_2\|_{\delta},\\
(\mathrm{K}2)\qquad \|(T+\lambda)^{-1}\|_{\delta'}\le C'(\gamma_0+\lambda)^{-1},\\
T^{-1/2}=\pi^{-1}\int_0^{\infty}(T+\lambda)^{-1}\lambda^{-1/2}\,d\lambda ,\\
\delta(T+\lambda)^{-1}=-(T'+\lambda)^{-1}\,\delta T\,(T+\lambda)^{-1},
\end{gathered}
\end{equation}
where $\delta'=\delta'(\gamma_0,\Gamma,\delta)>0$ and $C'$ is an absolute constant (the proof
gives $C'=2$).
In the application, $T_k=C^{*}\Delta^{(k)}C$ is the operator of
\cite[(3.187)]{B-CMP99b}. For $\widetilde{B}$ with $e^{i\widetilde{B}}V_0$ in the tube
$\mathbb{D}(r_c)$ of \eqref{9.3:eq-perfect-action-a}, write $T(\widetilde{B})$ for $T_k$ at the
data $e^{i\widetilde{B}}V_0$. Here $\delta>0$ is a fixed decay rate at which the kernel of
$T(\widetilde{B})$ is bounded in the norm $\|\cdot\|_{\delta}$ on the tube, and $c_{\delta}$ is
evaluated at this $\delta$. On this tube
$\|T(\widetilde{B})-T(\operatorname{Re}\widetilde{B})\|_{\delta}\le C_{\mathbb{D}}r_c$, with
$C_{\mathbb{D}}$ the constant of the Cauchy estimate for the analytic map
$\widetilde{B}\mapsto T(\widetilde{B})$ in the norm $\|\cdot\|_{\delta}$. The constant
$c_{\gamma}$ of Notation~\ref{9.3:not-perfect-action} is
$c_{\gamma}:=\gamma_0/(2C_{\mathbb{D}}c_{\delta})$, so that
$C_{\mathbb{D}}c_{\delta}r_c\le\gamma_0/2$.
\end{lemma}

\begin{proof}
On the torus each bond has lifts to $\mathbb{Z}^4$, and its periodic distance to $b$ is the least
distance of a lift. Hence the sums below are bounded by the sums over $\mathbb{Z}^4$.

\emph{(K1).} By the triangle inequality,
$|b-b''|+|b''-b'|\ge|b-b'|$. Hence
\begin{align*}
|(\mathsf{K}_1\mathsf{K}_2)(b,b')|
&\le\|\mathsf{K}_1\|_{\delta}\|\mathsf{K}_2\|_{\delta}
\sum_{b''}e^{-\delta|b-b''|-\delta|b''-b'|}\\
&\le\|\mathsf{K}_1\|_{\delta}\|\mathsf{K}_2\|_{\delta}\,e^{-\delta|b-b'|/2}
\sum_{b''}e^{-\delta|b-b''|/2}.
\end{align*}
The last sum is at most $c_{\delta}$.

\emph{A norm bound.} Let $\mathsf{K}$ be a kernel. By Cauchy--Schwarz,
\begin{equation*}
|\langle u_1,\mathsf{K}u_2\rangle|\le\sum_{b,b'}|u_1(b)||\mathsf{K}(b,b')||u_2(b')|
\le\Bigl(\sup_b\sum_{b'}|\mathsf{K}(b,b')|\Bigr)^{1/2}
\Bigl(\sup_{b'}\sum_{b}|\mathsf{K}(b,b')|\Bigr)^{1/2}\|u_1\|\|u_2\|.
\end{equation*}
Thus the operator norm of $\mathsf{K}$ is at most the larger of its row and column sums.

\emph{(K2).} We use the argument of Combes and Thomas \cite{CT73}. Fix $b'$. For $a>0$ let
$T_a$ be $T$ conjugated by multiplication with $e^{a|\cdot-b'|}$, that is, the kernel with
entries $e^{a|b-b'|}T(b,b'')e^{-a|b''-b'|}$. Its entries differ from those of $T$ by
\begin{equation*}
T(b,b'')\bigl(e^{a(|b-b'|-|b''-b'|)}-1\bigr).
\end{equation*}
We have $\bigl||b-b'|-|b''-b'|\bigr|\le|b-b''|$, and $|e^{y}-1|\le|y|e^{|y|}$ for real $y$.
For $a\le\delta/2$, the entry therefore has norm at most
\begin{equation*}
\Gamma a|b-b''|e^{-\delta|b-b''|/2}\le\Gamma a\,\tfrac{4}{e\delta}\,e^{-\delta|b-b''|/4}.
\end{equation*}
This uses $xe^{-\delta x/4}\le 4/(e\delta)$. The row and column sums are then at most
$\Gamma a\,4c_{\delta/2}/(e\delta)$.

Take
\begin{equation*}
\delta':=\min\{\delta/2,\;e\delta\gamma_0/(8\Gamma c_{\delta/2})\}.
\end{equation*}
With $a=\delta'$, the norm bound gives $\|T_a-T\|\le\gamma_0/2$. Hence
\begin{equation*}
\operatorname{Re}(T_a+\lambda)\ge\gamma_0/2+\lambda\ge(\gamma_0+\lambda)/2.
\end{equation*}
It follows that $T_a+\lambda$ is invertible with
$\|(T_a+\lambda)^{-1}\|\le2(\gamma_0+\lambda)^{-1}$. Now $(T_a+\lambda)^{-1}$ is
$(T+\lambda)^{-1}$ conjugated by multiplication with $e^{a|\cdot-b'|}$, and its $(b,b')$ entry is
$e^{\delta'|b-b'|}(T+\lambda)^{-1}(b,b')$, since $|b'-b'|=0$. This entry is bounded by the
operator norm. As $\delta'$ does not depend on $b'$, (K2) holds with $C'=2$.

\emph{The integral formula in (K2).} If $Tu=zu$ with $\|u\|=1$, then
$\operatorname{Re}z=\operatorname{Re}\langle u,Tu\rangle\ge\gamma_0$. So the spectrum of $T$
lies in $\{\operatorname{Re}z\ge\gamma_0\}$, and $T^{-1/2}$ is given by the holomorphic
functional calculus with the principal branch.

For $t>0$, substitute $\lambda=tu^2$:
\begin{equation*}
\int_0^\infty\frac{\lambda^{-1/2}\,d\lambda}{t+\lambda}
=2t^{-1/2}\int_0^\infty\frac{du}{1+u^2}=\pi t^{-1/2}.
\end{equation*}
We have $|t+\lambda|\ge\operatorname{Re}t+\lambda$. Hence the integral converges locally
uniformly on $\{\operatorname{Re}t>0\}$ and is holomorphic there, so the identity extends to that
half-plane.

Represent $(T+\lambda)^{-1}$ and $T^{-1/2}$ as contour integrals of
$(z-T)^{-1}$ over a contour in $\{\operatorname{Re}z>0\}$ enclosing the spectrum. Interchanging
the $\lambda$- and $z$-integrals by Fubini gives the integral formula for $T^{-1/2}$ in
\eqref{9.3:eq-kernel-calculus-a}.

\emph{The resolvent identity in (K2).} We have
\begin{equation*}
(T'+\lambda)^{-1}-(T+\lambda)^{-1}
=(T'+\lambda)^{-1}\bigl[(T+\lambda)-(T'+\lambda)\bigr](T+\lambda)^{-1},
\end{equation*}
and the bracket equals $-\delta T$.

\emph{Application.} \cite[(3.187)]{B-CMP99b} states that $C^{*}\Delta^{(k)}C$ has ``a lower
bound $\gamma_0>0$ independent of $k$ and $U$''. The bound on the tube comes from the Cauchy
estimate for the analytic dependence of the kernel of $T$ on the tube variable.
Notation~\ref{9.3:not-perfect-action} gives $r_c\le c_{\gamma}$. With the definition of
$c_{\gamma}$ this yields $C_{\mathbb{D}}c_{\delta}r_c\le\gamma_0/2$.
\end{proof}

\begin{theorem}[Stationarity of the fluctuation covariance]\label{9.3:thm-stationarity-covariance}
Let $\theta\in(0,1)$ be as in Theorem~\ref{9.3:thm-stationarity-minimisers}. Work on the unit lattice of the step from level $k$ to level $k+1$. Either all the domains $\Omega_j$ of the admissible sequence $\boldsymbol{\Omega}$ equal the whole torus of this lattice, or only bonds of a set $\Lambda\subset\Lambda_k$ with $\operatorname{dist}(\Lambda,\Lambda_k^{c})>RM$ are considered. Here $R$ is the free integer of Ba{\l}aban's construction that enters, with his decay rate $\delta_0$, the floor $\delta_0RM_1\ge16\log L$, and $M=L^{m'}$ is the integer of the auxiliary parameters $\mathfrak{p}$. For a kernel $\mathsf{K}$ on bonds put
\begin{equation*}
\|\mathsf{K}\|_{\delta}:=\sup_{b,b'}|\mathsf{K}(b,b')|\,e^{\delta|b-b'|},
\qquad
\vartheta_C:=L^{-\theta/8}.
\end{equation*}
Consider the data of the step computed with $k$ lattices beneath the unit lattice (determining set $\mathfrak{B}$) and with $k+1$ lattices beneath it (determining set $\mathfrak{B}'$). These data are written in the axial gauge of \cite[(2.2)--(2.3)]{B-CMP109} and regarded as analytic functions of the next field $\widetilde{W}$ on its complex tube. For each such object $X$ write $\delta X:=X_{\mathfrak{B}'}-X_{\mathfrak{B}}$. The objects are:
\begin{itemize}
\item the configuration $V^{(k)}$ of \cite[(2.3)]{B-CMP109};
\item the operators $(QG_1Q^{*})^{-1}$, $\Delta^{(k)}$ and $C^{(k)}$ of \cite[(2.3)]{B-CMP109};
\item the operator $\Gamma_k$ of \cite[(2.5)--(2.6)]{B-CMP116}.
\end{itemize}
There are $\delta>0$ and $C_C<\infty$, depending on $(d,L,\theta)$ only, and in particular not on $k$, on $\boldsymbol{\Omega}$ or on the point $\widetilde{W}$ of the tube, such that for every $k$ and every $\boldsymbol{\Omega}$ as above
\begin{equation}\label{9.3:eq-stationarity-covariance-1}
\begin{gathered}
|\delta V^{(k)}|,\quad \|\delta(QG_1Q^{*})^{-1}\|_{\delta},\quad \|\delta\Delta^{(k)}\|_{\delta},\\
\|\delta C^{(k)}\|_{\delta},\quad \|\delta(C^{(k)})^{1/2}\|_{\delta},\quad \|\delta\Gamma_k\|_{\delta}
\;\le\; C_C\,\vartheta_C^{k}.
\end{gathered}
\end{equation}
Moreover, the first and second derivatives of each of these six differences with respect to $\widetilde{W}$ are bounded by $C_C\vartheta_C^{k}e^{-\delta\operatorname{diam}}$. Here $\operatorname{diam}$ is the diameter, in the distance $d(\cdot,\cdot)$, of the set formed by the arguments of the quantity and the bonds at which the derivatives are taken.
\end{theorem}

\begin{proof}
\emph{The configuration $V^{(k)}$.} By \cite[(2.3)]{B-CMP109},
\begin{equation}\label{9.3:eq-stationarity-covariance-a}
V^{(k)}(\widetilde{W})=M^{k}\bigl(U_{k+1}(\widetilde{W})\bigr).
\end{equation}
Here $M^{k}$ is the $k$-step averaging map, not a power of the integer $M$.

\emph{Real data.} Let $\widetilde{W}$ be real. We apply Theorem~\ref{9.3:thm-stationarity-minimisers} to the pair of depths $k+1$, $k+2$. We apply Corollary~\ref{9.3:cor-stationarity-level-data}\,(i) with the level taken one below the top. That clause relates $M^{i+1}(U')$, after the gauge transformation $\mathsf{u}$, to $M^{i}(U_{\mathfrak{B}})$ through the exponential of the imaginary unit times $L^{i}\eta A^{(i)}$, and it bounds $\ell|A^{(i)}|$ by $2C_A\varepsilon L^{-\theta j}$. Throughout this proof, the level letters $j$ and $j_{*}$ of the statements cited take the value $k$, the level of the unit lattice. Here $\varepsilon$ is the parameter of the condition $(7)_{\varepsilon}$ of Theorem~\ref{9.3:thm-stationarity-minimisers}. By \eqref{9.3:eq-stationarity-covariance-a}, the two versions of $V^{(k)}$ therefore agree up to a gauge transformation within $C_L\varepsilon L^{-\theta k}$, with a constant $C_L<\infty$ depending on $(d,L,\theta)$ only.

\emph{Passage to the axial gauges.} It remains to pass from $\mathsf{u}$ to the canonical axial gauges of \cite[(2.2)--(2.3)]{B-CMP109}. Let $u$ denote the gauge transformation relating the two axial representatives. It is determined by the configuration $e^{i\eta A}$, where $A$ is the potential of Theorem~\ref{9.3:thm-stationarity-minimisers}, the factor relating $W^{\mathsf{u}}$ to $U_{\mathfrak{B}}(V)$; by \cite[(1.70)--(1.72)]{B-CMP99a}
\begin{equation*}
|u-1|\le CL\sum_{n}(L^{n}\eta)\,\sup|A|.
\end{equation*}
This passage therefore costs a factor $C_L$. In the axial gauges, $|\delta V^{(k)}|\le C_L\varepsilon L^{-\theta k}$ for real $\widetilde{W}$.

\emph{Complex data.} The components of $\delta V^{(k)}$ are analytic and bounded on the $\widetilde{W}$-tube; let $\mathsf{M}$ be their bound there. Put $\mathsf{m}=C_L\varepsilon L^{-\theta k}$, the bound on real data. The two-constants lemma on the disc (Lemma~\ref{9.3:lem-two-constants}) gives $|\delta V^{(k)}|\le(\mathsf{m}\mathsf{M})^{1/2}$. Hence
\begin{equation}\label{9.3:eq-stationarity-covariance-2}
|\delta V^{(k)}|\le CL^{-\theta k/2}.
\end{equation}

\emph{Decay rates.} In the rest of this proof decay rates are written generically: the letter $\delta$ in an exponent or as the subscript of a kernel norm denotes a positive rate depending on $(d,L,\theta)$ only, which decreases at each of the finitely many steps below; the $\delta$ of the statement is the last of them. Placed in front of an object, $\delta$ denotes, as in the statement, the difference between the two depths.

\emph{Locality of $V^{(k)}$.} Applied to $V^{(k)}$, $\partial_c$ denotes differentiation in the component of $\widetilde{W}$ at the bond $c$. By Ba{\l}aban's second-variation bound (Proposition~\ref{9.3:prop-second-variation}, printed without proof in \cite[(190)]{B-CMP102b}; cited as printed) and by \cite[Proposition~5]{B-CMP98}, the derivatives of $V^{(k)}$ with respect to $\widetilde{W}$ decay:
\begin{equation*}
|\partial_cV^{(k)}(c_1)|\le Ce^{-\delta|c-c_1|}.
\end{equation*}
Let $r$ be the radius of the tube. The Cauchy estimate, applied to this bound in the variable at $c'$, gives the bound $(C/r)e^{-\delta|c-c_1|}$ for $|\partial_c\partial_{c'}V^{(k)}(c_1)|$; exchanging the roles of $c$ and $c'$ gives the bound $(C/r)e^{-\delta|c'-c_1|}$. Hence
\begin{equation*}
|\partial_c\partial_{c'}V^{(k)}(c_1)|
\le\frac{C}{r}\,e^{-\delta\max\{|c-c_1|,\,|c'-c_1|\}}
\le\frac{C}{r}\,e^{-\delta\operatorname{diam}/2}.
\end{equation*}
The last inequality holds because the diameter of $\{c,c',c_1\}$ is at most $|c-c_1|+|c'-c_1|\le2\max\{|c-c_1|,|c'-c_1|\}$. These bounds hold at each depth, hence for both terms of $\delta V^{(k)}$.

\emph{The operator $(QG_1Q^{*})^{-1}$.} Write $(QG_1Q^{*})^{-1}=a^{(2)}_{\mathfrak{B}}+a$, where $a$ is a one-step object of the step in the sense of Lemma~\ref{9.3:lem-one-step-data}: analytic in the data, of range at most two $L$-blocks, and the same functional of the data at both depths. Write $V^{(k)\prime}$ for $V^{(k)}$ at depth $\mathfrak{B}'$. Then
\begin{equation*}
\begin{split}
\delta(QG_1Q^{*})^{-1}={}&\delta a^{(2)}\bigl(V^{(k)\prime}\bigr)\\
&+\bigl\{(a^{(2)}_{\mathfrak{B}}+a)\bigl(V^{(k)\prime}\bigr)-(a^{(2)}_{\mathfrak{B}}+a)\bigl(V^{(k)}\bigr)\bigr\}.
\end{split}
\end{equation*}
The first term is bounded by the bound on $\delta a^{(n)}$ at $n=2$ in Theorem~\ref{9.3:thm-stationarity-jets}\,(c). The brace is estimated along the segment joining $V^{(k)}$ to $V^{(k)\prime}$, with \eqref{9.3:eq-stationarity-covariance-2}, exactly as the brace in the decomposition of $\delta\Delta^{(k)}$ below.

\emph{The operator $\Delta^{(k)}$: decomposition.} By the lemma on one-step data as functionals of the jets (Lemma~\ref{9.3:lem-one-step-data}), $\Delta^{(k)}=\Delta^{[\mathfrak{B}]}(V^{(k)})$, with $\Delta^{[\mathfrak{B}]}$ given by \eqref{9.3:eq-one-step-data-a}. In this formula $h$, $C^{(2)}$ and $G^{(2)}$ are one-step objects, the same at both depths. Then
\begin{equation*}
\begin{split}
\delta\Delta^{(k)}={}&\bigl[\delta a^{(2)}-2\,\delta\omega\circ hC^{(2)}\bigr]\bigl(V^{(k)\prime}\bigr)\\
&+\bigl\{\Delta^{[\mathfrak{B}]}\bigl(V^{(k)\prime}\bigr)-\Delta^{[\mathfrak{B}]}\bigl(V^{(k)}\bigr)\bigr\}.
\end{split}
\end{equation*}

\emph{The bracket.} The bracket is bounded by Theorem~\ref{9.3:thm-stationarity-jets}. Part (b) gives $|\delta\omega(c)|\le C_BL^{-\theta j/2}$. Part (c) gives
\begin{equation*}
|\delta a^{(n)}|\le C_nL^{-\theta j_{*}/4}e^{-\delta_1\operatorname{diam}/2},
\end{equation*}
applied here at the level $k$ of the unit lattice.

\emph{The brace.} We estimate the brace along the segment joining $V^{(k)}$ to $V^{(k)\prime}$. Let $\delta B(c)$ be the difference of the coordinates of the two data at the bond $c$; it is at most a constant times $\sup|\delta V^{(k)}|$. Applied to $\Delta^{[\mathfrak{B}]}$, here and below, $\partial_c$ denotes differentiation in the coordinate $B(c)$ of the data. By Theorem~\ref{9.3:thm-stationarity-jets}\,(c) for the jet terms, and by the analyticity in $V$ and the range of at most two $L$-blocks of the one-step objects $h$, $C^{(2)}$, $G^{(2)}$ (Lemma~\ref{9.3:lem-one-step-data}) for the others, $|\partial_c\Delta^{[\mathfrak{B}]}(b,b')|\le C_3e^{-\delta_1\operatorname{diam}(b,b',c)}$. Hence
\begin{equation*}
\begin{split}
\sum_c|\partial_c\Delta^{[\mathfrak{B}]}(b,b')|\,|\delta B(c)|
&\le\sum_cC_3e^{-\delta_1\operatorname{diam}(b,b',c)}\sup|\delta V^{(k)}|\\
&\le Ce^{-\delta_1|b-b'|/2}L^{-\theta k/2}.
\end{split}
\end{equation*}
In the last step we used three facts. First, \eqref{9.3:eq-stationarity-covariance-2}. Second, $\operatorname{diam}(b,b',c)\ge\frac12(|b-b'|+|b-c|)$. Third, $\sum_ce^{-\delta_1|b-c|/2}$ is bounded independently of $b$. Together these give the bound on $\|\delta\Delta^{(k)}\|_{\delta}$.

\emph{The operators $C^{(k)}$, $(C^{(k)})^{1/2}$ and $\Gamma_k$.} Let $C(V)$ be the eliminating operator of the one-step data and put $T_k:=C(V^{(k)})^{*}\Delta^{(k)}C(V^{(k)})$. Then $C^{(k)}=T_k^{-1}$, and $\Gamma_k:=(T_k)_{IE}(C^{(k)})^{1/2}$, with $(T_k)_{IE}$ the block of $T_k$ between the two bond sets $I$ and $E$ of \cite[(2.5)--(2.6)]{B-CMP116}; its kernel norms, and those of its difference between the two depths, are therefore at most the corresponding norms of $T_k$ and of $\delta T_k$ at the same rate. These are controlled by the calculus of exponentially decaying kernels (Lemma~\ref{9.3:lem-kernel-calculus}), namely (K1) and (K2) of \eqref{9.3:eq-kernel-calculus-a}. By Lemma~\ref{9.3:lem-kernel-calculus}, $T_k$ satisfies the hypotheses of (K2) on the tube, with the constant $\gamma_0$ given there.
\begin{itemize}
\item Applied to the products defining $T_k$, (K1) bounds $\delta T_k$ in kernel norm at a smaller rate. Here the variation of $C(V^{(k)})$ between the two depths is bounded by a constant times $e^{-\delta_1|b-b'|/2}L^{-\theta k/2}$ exactly as the brace, $C(V)$ being a one-step object analytic in $V$ (Lemma~\ref{9.3:lem-one-step-data}).
\item The resolvent identity of (K2) at $\lambda=0$ gives $\delta C^{(k)}=-C^{(k)}_{\mathfrak{B}'}\,\delta T_k\,C^{(k)}_{\mathfrak{B}}$. Both inverses are controlled by (K2), and the product by (K1).
\item The integral formula of (K2) gives
\begin{equation*}
\delta(C^{(k)})^{1/2}
=-\frac1\pi\int_0^\infty(T_k'+\lambda)^{-1}\,\delta T_k\,(T_k+\lambda)^{-1}\lambda^{-1/2}\,d\lambda,
\end{equation*}
where $T_k'$ is $T_k$ at depth $\mathfrak{B}'$. The integrand is bounded by $C(\gamma_0+\lambda)^{-2}\lambda^{-1/2}$ times the kernel norm of $\delta T_k$ at the rate of the preceding item, which is integrable.
\item Finally,
\begin{equation*}
\delta\Gamma_k=\delta(T_k)_{IE}\,(C^{(k)}_{\mathfrak{B}'})^{1/2}+(T_k)_{IE}\,\delta(C^{(k)})^{1/2},
\end{equation*}
and (K1) bounds both terms.
\end{itemize}
Each application of (K1) or (K2) replaces the decay rate by a smaller positive rate. Finitely many applications occur, and $\delta$ denotes the final rate.

\emph{Derivatives of each object.} The $\widetilde{W}$-derivatives are obtained by the chain rule. For instance,
\begin{equation*}
\partial_c\Delta^{(k)}(b,b')
=\sum_{c_1}\partial_{c_1}\Delta^{[\mathfrak{B}]}(V^{(k)})(b,b')\,\partial_cV^{(k)}(c_1).
\end{equation*}
This is combined with the locality of $V^{(k)}$ established after \eqref{9.3:eq-stationarity-covariance-a}, with the decay of the jets in Theorem~\ref{9.3:thm-stationarity-jets}\,(c), and, for the inverses and the square root, with (K1) and (K2). This yields bounds $Ce^{-\delta\operatorname{diam}}$ for the first two derivatives of each object at each depth.

\emph{Derivatives of the differences.} Each derivative of a difference $\delta X$ obeys two bounds:
\begin{itemize}
\item a small bound, given by the Cauchy estimate applied to the bound on $\delta X$ on the tube;
\item a decaying bound $Ce^{-\delta\operatorname{diam}}$, since both terms of $\delta X$ obey it.
\end{itemize}
The minimum of the two is at most their geometric mean.

\emph{Exponent count.} The exponent $\theta k$ of Theorem~\ref{9.3:thm-stationarity-minimisers} and Corollary~\ref{9.3:cor-stationarity-level-data}\,(i) is halved three times:
\begin{itemize}
\item by Lemma~\ref{9.3:lem-two-constants}, in \eqref{9.3:eq-stationarity-covariance-2};
\item in the bound on $\delta a^{(n)}$ of Theorem~\ref{9.3:thm-stationarity-jets}\,(c);
\item by the geometric mean.
\end{itemize}
This gives the factor $L^{-\theta k/8}=\vartheta_C^{k}$, both in \eqref{9.3:eq-stationarity-covariance-1} and for the derivatives. The decay rate is halved as well, and is renamed $\delta$.

\emph{The case $\Lambda\subset\Lambda_k$.} On bonds of $\Lambda\subset\Lambda_k$ with $\operatorname{dist}(\Lambda,\Lambda_k^{c})>RM$ the same argument applies verbatim. The bounds of Theorem~\ref{9.3:thm-stationarity-minimisers} and Theorem~\ref{9.3:thm-stationarity-jets} are uniform in $\boldsymbol{\Omega}$.
\end{proof}

\begin{theorem}[Convergence of the one-loop coefficients and offsets]\label{9.3:thm-one-loop-convergence}
Work in the setting and under the conditions of Theorem~\ref{9.3:thm-stationarity-covariance}, with
$\vartheta_C=L^{-\theta/8}$. For $l\ge1$ let $\beta^{0}_{l}$ be the one-loop coefficient at level $l$,
that is, the value $\beta[\widehat{L}^{(l)}]$ of the one-loop functional $\beta[\cdot]$ on the localised
one-loop family $\widehat{L}^{(l)}$ at level $l$.
There is a constant $C_D$, depending only on $d$, $L$ and $\theta$, such that
\begin{equation}\label{9.3:eq-one-loop-convergence-1}
\bigl|\beta^{0}_{l+1}-\beta^{0}_{l}\bigr|\le C_D\,\vartheta_C^{\,l}\qquad(l\ge1).
\end{equation}
Assume in addition that the offsets $(s_l)_{l\ge1}$ of the correlation theorem form a real sequence with
\begin{enumerate}
\item[(a)] $s_{l-1}-s_l=\beta^{0}_{l}-c_0$ for every $l\ge2$;
\item[(b)] the differences $s_i-s_j$, $i,j\ge1$, are bounded, as the statement on the offsets of the
correlation theorem asserts.
\end{enumerate}
Then $\beta^{0}_{l}\to c_0$, the limit $s_\infty:=\lim_{l\to\infty}s_l$ exists, and for every $l\ge1$
\begin{equation}\label{9.3:eq-one-loop-convergence-2}
\bigl|\beta^{0}_{l}-c_0\bigr|\le\frac{C_D\,\vartheta_C^{\,l}}{1-\vartheta_C},\qquad
\bigl|s_l-s_\infty\bigr|\le\frac{C_D\,\vartheta_C^{\,l}}{(1-\vartheta_C)^{2}}.
\end{equation}
\end{theorem}

\begin{proof}
\emph{The second derivative of the step's logarithm.} For the next field $\widetilde{W}$ on its tube put
$\varphi_l(\widetilde{W}):=\log Z^{(l)}(U_{l+1}(\widetilde{W}))$, where $Z^{(l)}$ is the fluctuation
integral of the step from level $l$. By \cite[(1.4)]{B-CMP109},
\begin{equation}\label{9.3:eq-one-loop-convergence-3}
\varphi_l(\widetilde{W})=-\tfrac12\log\det T_l+\mathfrak{j}(V^{(l)}),
\end{equation}
where $T_l=C^{*}\Delta^{(l)}C$ is the sandwiched fluctuation operator, $C=C(V^{(l)})$ the eliminating
operator, $\mathfrak{j}$ the Jacobian, and $V^{(l)}=V^{(l)}(\widetilde{W})$. For bonds $c,c'$ write
$\partial_c$ for the derivative with respect to the variable of $\widetilde{W}$ at the bond $c$ (in the
coordinates $\widetilde{W}=e^{iB}$ of the tube), and put $\Pi_l(c,c'):=\partial_c\partial_{c'}\varphi_l(1)$.
Since $\partial_c\log\det T=\operatorname{Tr}[T^{-1}\partial_cT]$ and
$\partial_{c'}T^{-1}=-T^{-1}\partial_{c'}T\,T^{-1}$, differentiating
\eqref{9.3:eq-one-loop-convergence-3} twice gives, with $T=T_l$ evaluated at $\widetilde{W}=1$,
\begin{equation}\label{9.3:eq-one-loop-convergence-4}
\Pi_l(c,c')=-\tfrac12\operatorname{Tr}\bigl[T^{-1}\partial_c\partial_{c'}T\bigr]
+\tfrac12\operatorname{Tr}\bigl[T^{-1}\partial_cT\,T^{-1}\partial_{c'}T\bigr]
+\partial_c\partial_{c'}\mathfrak{j}.
\end{equation}

\emph{Decay.} The kernel of $T^{-1}=C^{(l)}$ decays exponentially by the resolvent bound in
\eqref{9.3:eq-kernel-calculus-a} of Lemma~\ref{9.3:lem-kernel-calculus}; the kernels of $\partial_cT$,
$\partial_c\partial_{c'}T$ and of the derivatives of $\mathfrak{j}$ decay exponentially away from $c$ and
from $c'$. Each trace in \eqref{9.3:eq-one-loop-convergence-4} is thus a sum
over bonds of products of kernels decaying from $c$ and from $c'$. By the product bound in
\eqref{9.3:eq-kernel-calculus-a} of Lemma~\ref{9.3:lem-kernel-calculus} this sum converges absolutely,
and there are constants $C'$ and $\delta_3>0$, independent of $l$ and of the torus, with
\begin{equation}\label{9.3:eq-one-loop-convergence-7}
|\Pi_l(c,c')|\le C'\,e^{-\delta_3|c-c'|}\qquad(l\ge0).
\end{equation}

\emph{The difference between consecutive levels.} Let $\delta$ denote the difference of an object between
the two depths, as in Theorem~\ref{9.3:thm-stationarity-covariance}. Apply $\delta$ to each term of
\eqref{9.3:eq-one-loop-convergence-4} and telescope the products of kernels by
$\delta(\mathsf{K}_1\mathsf{K}_2)=\delta\mathsf{K}_1\,\mathsf{K}_2'+\mathsf{K}_1\,\delta\mathsf{K}_2$,
with $\mathsf{K}_2'$ the kernel at the second depth. Each resulting term contains exactly one
$\delta$-factor. That
factor is the difference $\delta$ of $C^{(l)}$, of $\Delta^{(l)}$ or of one of its first two
$\widetilde{W}$-derivatives, or of a one-step functional of $V^{(l)}$ that does not depend on the level
(the eliminating operator $C(V^{(l)})$ or $\mathfrak{j}$) or of one of its first two
$\widetilde{W}$-derivatives; a factor of the last kind is controlled through $\delta V^{(l)}$ and its
first two $\widetilde{W}$-derivatives. By
Theorem~\ref{9.3:thm-stationarity-covariance}, each such factor is at most
$C_C\vartheta_C^{\,l}e^{-\delta_2\operatorname{diam}}$, where $\delta_2>0$ is the decay rate written
$\delta$ in the bounds of that theorem and $\operatorname{diam}$ is the diameter of the set of bond
arguments. All remaining factors are bounded as in the proof of \eqref{9.3:eq-one-loop-convergence-7},
and the sums over bonds are controlled by the product
bound in \eqref{9.3:eq-kernel-calculus-a}. After enlarging $C'$ and decreasing $\delta_3$ if necessary,
we obtain
\begin{equation}\label{9.3:eq-one-loop-convergence-5}
|\delta\Pi_l(c,c')|\le C'\,\vartheta_C^{\,l}\,e^{-\delta_3|c-c'|},
\end{equation}
uniformly in the torus, for every $l\ge1$. The limit of the tori to $\mathbb{Z}^d$ in
\cite[(1.21)--(1.22)]{B-CMP109}
preserves \eqref{9.3:eq-one-loop-convergence-5} and \eqref{9.3:eq-one-loop-convergence-7}, both being
uniform in the torus. Write $\Pi_{l,\mu\nu}(x)$ for the resulting infinite-volume kernel in the notation of
the defining identity recalled below (a bond of direction $\mu$ at the origin, a bond of direction $\nu$
at $x$).

\emph{The coefficients.} By the defining identity of the one-loop coefficient in the correlation
theorem, the functional $\beta[\cdot]$ is applied to the total second derivative $\Pi_l$,
Jacobian term included, and for $l\ge0$
\begin{equation*}
\beta^{0}_{l+1}=\beta[\widehat{L}^{(l+1)}]=\sum_{x}\Pi_{l,\mu\nu}(x)\,x_\mu x_\nu .
\end{equation*}
The bound below is uniform in the directions $\mu,\nu$; a summation over $\mu,\nu$ costs a factor
$d^{2}$, absorbed in $C'$. Applied at the second depth, the same identity
expresses $\beta^{0}_{l+2}$ as the second moment of $\Pi_l$ at the second depth, that is, of the total
second derivative of the step with one more lattice beneath. Hence
$\beta^{0}_{l+2}-\beta^{0}_{l+1}=\sum_x\delta\Pi_{l,\mu\nu}(x)x_\mu x_\nu$. By
\eqref{9.3:eq-one-loop-convergence-5} and $|x_\mu x_\nu|\le|x|^2$,
\begin{equation*}
\bigl|\beta^{0}_{l+2}-\beta^{0}_{l+1}\bigr|\le C'\,\vartheta_C^{\,l}\,S,\qquad
S:=\sum_x|x|^{2}e^{-\delta_3|x|},
\end{equation*}
for every $l\ge1$, and $S$ is finite. Put $C_D:=2C'\,S\,\vartheta_C^{-1}$. Replacing $l$ by $l-1$
gives \eqref{9.3:eq-one-loop-convergence-1} for $l\ge2$. For $l=1$,
\eqref{9.3:eq-one-loop-convergence-7} gives $|\beta^{0}_{l+1}|\le C'\,S$ for every $l\ge0$, hence
$|\beta^{0}_{2}-\beta^{0}_{1}|\le 2C'\,S=C_D\vartheta_C$, which is \eqref{9.3:eq-one-loop-convergence-1}
for $l=1$.

\emph{Arithmetic.} Since $L>1$ and $\theta>0$, we have $0<\vartheta_C<1$, so by
\eqref{9.3:eq-one-loop-convergence-1} the series $\sum_l(\beta^{0}_{l+1}-\beta^{0}_{l})$ converges
absolutely. Thus $\beta^{0}_{l}\to\beta^{0}_{\infty}$ for some $\beta^{0}_{\infty}$, and
\begin{equation}\label{9.3:eq-one-loop-convergence-6}
\bigl|\beta^{0}_{l}-\beta^{0}_{\infty}\bigr|\le\sum_{l'\ge l}C_D\vartheta_C^{\,l'}
=\frac{C_D\vartheta_C^{\,l}}{1-\vartheta_C}.
\end{equation}
For $1\le i<j$, summing (a) over $l=i+1,\dots,j$ gives
\begin{equation*}
s_i-s_j=\sum_{l=i+1}^{j}(\beta^{0}_{l}-c_0)
=(j-i)(\beta^{0}_{\infty}-c_0)+\sum_{l=i+1}^{j}(\beta^{0}_{l}-\beta^{0}_{\infty}).
\end{equation*}
By \eqref{9.3:eq-one-loop-convergence-6}, the last sum is at most $C_D/(1-\vartheta_C)^2$ in absolute
value. By (b), $s_i-s_j$ is bounded, so $(j-i)(\beta^{0}_{\infty}-c_0)$ is bounded in $j$. Letting
$j\to\infty$ forces $\beta^{0}_{\infty}=c_0$. This gives $\beta^{0}_{l}\to c_0$ and, by
\eqref{9.3:eq-one-loop-convergence-6}, the first bound of \eqref{9.3:eq-one-loop-convergence-2}.

By (a) and this bound, the increments $s_{l'-1}-s_{l'}=\beta^{0}_{l'}-c_0$ are summable, so $(s_l)$ is
Cauchy and $s_\infty$ exists. Moreover
\begin{align*}
|s_l-s_\infty|&\le\sum_{l'\ge l+1}|\beta^{0}_{l'}-c_0|
\le\sum_{l'\ge l+1}\frac{C_D\vartheta_C^{\,l'}}{1-\vartheta_C}\\
&=\frac{C_D\vartheta_C^{\,l+1}}{(1-\vartheta_C)^2}\le\frac{C_D\vartheta_C^{\,l}}{(1-\vartheta_C)^2},
\end{align*}
which is the second bound of \eqref{9.3:eq-one-loop-convergence-2}.
\end{proof}

\begin{lemma}[Refinement of the determining set]\label{9.3:lem-refinement}
Let $\boldsymbol{\Omega}=(\Omega_0\supset\dots\supset\Omega_k)$ be admissible on $T_{\eta}$, and let
$V$ be data on the determining set $\mathfrak{B}=\bigcup_j\Lambda_j$ satisfying
\eqref{9.3:eq-paired-variational-a} with parameter $\varepsilon$, as in
Definition~\ref{9.3:not-paired-variational}.
Let $\widetilde{\boldsymbol{\Omega}}=(\widetilde{\Omega}_0\supset\dots\supset\widetilde{\Omega}_k)$ be
admissible on $T_{\eta}$ with $\widetilde{\Omega}_j\subseteq\Omega_j$ for every $j$. Put
$\widetilde{\Lambda}_j:=(\widetilde{\Omega}_j\setminus\widetilde{\Omega}_{j+1})^{(j)}$ and
$\widetilde{\mathfrak{B}}:=\bigcup_j\widetilde{\Lambda}_j$. Write $M_{\mathfrak{B}}$ for the map
that assigns to a configuration the averages $M^{j}(\cdot)$ on $\Lambda_j$ for every $j$, and
$M_{\widetilde{\mathfrak{B}}}$ for the map that assigns the averages $M^{j}(\cdot)$ on
$\widetilde{\Lambda}_j$ for every $j$. Let
\begin{equation*}
U:=U_{\mathfrak{B}}(V),\qquad \widetilde{V}:=M_{\widetilde{\mathfrak{B}}}(U).
\end{equation*}
Then $U$ is minimal for $(\widetilde{\mathfrak{B}},\widetilde{V})$.
\end{lemma}

\begin{proof}
Because $\widetilde{\Omega}_j\subseteq\Omega_j$ for every $j$, the averages that define
$M_{\mathfrak{B}}$ are compositions of averaging maps applied to the averages that define
$M_{\widetilde{\mathfrak{B}}}$. Hence every configuration whose averages on
$\widetilde{\mathfrak{B}}$ equal $\widetilde{V}=M_{\widetilde{\mathfrak{B}}}(U)$ has averages
$M_{\mathfrak{B}}(U)=V$ on $\mathfrak{B}$. This gives
\begin{equation}\label{9.3:eq-refinement-1}
U\in\{M_{\widetilde{\mathfrak{B}}}=\widetilde{V}\}\subset\{M_{\mathfrak{B}}=V\},
\end{equation}
where $U$ belongs to the smaller set by the definition of $\widetilde{V}$. In the words of
\cite[(150)]{B-CMP102b}, the smaller set ``is defined by more restrictive functional conditions''.

The configuration $U=U_{\mathfrak{B}}(V)$ minimises $A^{\eta}$ on $\{M_{\mathfrak{B}}=V\}$ and is
therefore a critical point of $A^{\eta}$ restricted to that set. By \eqref{9.3:eq-refinement-1},
it is then also a critical point of $A^{\eta}$ restricted to the smaller set
$\{M_{\widetilde{\mathfrak{B}}}=\widetilde{V}\}$.

By \cite[Theorem~1]{B-CMP102b}, recorded in Definition~\ref{9.3:not-minimiser-bounds},
$U\in\mathfrak{U}_k(\boldsymbol{\Omega},B_3\varepsilon)$, with the constant $B_3$ of that
theorem. This membership carries over to
$U\in\mathfrak{U}_k(\widetilde{\boldsymbol{\Omega}},B_3\varepsilon)$ because
$\widetilde{\Omega}_j\subseteq\Omega_j$ for every $j$. Indeed, since
$\widetilde{\Omega}_j\subseteq\Omega_j$, the level of every point with respect to
$\widetilde{\boldsymbol{\Omega}}$ is at most its level with respect to $\boldsymbol{\Omega}$, so the
local scale $\ell$ can only decrease and the bounds defining the class only weaken.

The data $\widetilde{V}$ satisfy condition \eqref{9.3:eq-paired-variational-a} on
$\widetilde{\mathfrak{B}}$ with $\varepsilon$ replaced by $2B_3\varepsilon$, by
\cite[Proposition~2]{B-CMP98}.

We apply \cite[Theorem~1]{B-CMP102b} to the admissible sequence $\widetilde{\boldsymbol{\Omega}}$
and the data $\widetilde{V}$, which satisfy the data condition with parameter $2B_3\varepsilon$.
It gives a minimal configuration $U_{\widetilde{\mathfrak{B}}}(\widetilde{V})$ of $A^{\eta}$ on
$\{M_{\widetilde{\mathfrak{B}}}=\widetilde{V}\}$ and the uniqueness of the critical gauge orbit of
the problem $(\widetilde{\mathfrak{B}},\widetilde{V})$ within the class of regular configurations
in which that theorem asserts uniqueness for data at parameter $2B_3\varepsilon$. This class is
defined by upper bounds that increase with the parameter, and it contains
$\mathfrak{U}_k(\widetilde{\boldsymbol{\Omega}},B_3\varepsilon)$, to which $U$ belongs by the
preceding paragraphs. Since $U$ is a critical point of this problem, it lies on the gauge orbit of
$U_{\widetilde{\mathfrak{B}}}(\widetilde{V})$. Since $A^{\eta}$ is a sum over plaquettes of
functions of $\operatorname{Re}\operatorname{tr}U(\partial p)$, it is gauge invariant, so $U$
has the same action as $U_{\widetilde{\mathfrak{B}}}(\widetilde{V})$. Hence $U$ is minimal for
$(\widetilde{\mathfrak{B}},\widetilde{V})$.
\end{proof}

\begin{lemma}[Decay of the response to local data]\label{9.3:lem-response-decay}
Let $\widetilde{\mathfrak{B}}$ be a refined determining set as in Lemma~\ref{9.3:lem-refinement}, and
write $U(V) := U_{\widetilde{\mathfrak{B}}}(V)$ for the minimal configuration with prescribed averages
$V$ on $\widetilde{\mathfrak{B}}$. Let $V^1$ be real data on $\widetilde{\mathfrak{B}}$ satisfying
the regularity condition of \cite[(7)]{B-CMP102b} with parameter $a_A/2$, so that the tube
$\mathbb{D}(r_c)$ of \eqref{9.3:eq-perfect-action-a} about $V^1$ is
defined. Let $S\subset\widetilde{\mathfrak{B}}$, and let
$V^2 = e^{iB}V^1$, where $B$ is $\mathfrak{g}^{c}$-valued, $\sup|B| < r_c$ and
$\operatorname{supp} B \subset S$; for a point $y'$ we write $B(y')$ for the value of $B$ at $y'$.
Then, up to a gauge transformation,
\begin{equation*}
U(V^2) = e^{i\eta\mathcal{H}(B)}\,U(V^1),
\end{equation*}
where $\mathcal{H}$ is the response map of the minimiser of \cite[Proposition~9]{B-CMP102b}, and
there is a constant $C$, independent of $B$, $S$ and $y$, such that at every
point $y$, with local scale $\ell = L^{j(y)}\eta$,
\begin{equation}\label{9.3:eq-response-decay-1}
\ell\,\bigl|\mathcal{H}(B)(y)\bigr|,\quad \ell^{2}\,\bigl|(\nabla\mathcal{H}(B))(y)\bigr|
\;\le\; C\sum_{y'\in S} e^{-\frac18\delta_0 d(y,y')}\,|B(y')| .
\end{equation}
\end{lemma}

\begin{proof}
For $t\in[0,1]$ we have $\sup|tB| \le \sup|B| < r_c$, so the segment $t\mapsto e^{itB}V^1$ lies in
the tube $\mathbb{D}(r_c)$, on which, by \cite[Proposition~9]{B-CMP102b} applied as in
Notation~\ref{9.3:not-perfect-action}, the minimiser $U$ is analytic and
has the representation $U(e^{itB}V^1) = e^{i\eta\mathcal{H}(tB)}U(V^1)$ up to a gauge
transformation, with $\mathcal{H}$ analytic in its argument. Taking $t=1$ gives the first
assertion. The response map of \cite[Proposition~9]{B-CMP102b} vanishes at the zero perturbation,
$\mathcal{H}(0)=0$.

Hence, at each point $y$, by the fundamental theorem of calculus along the segment and the chain
rule,
\begin{equation*}
\mathcal{H}(B)(y) = \int_0^1 \frac{d}{dt}\,\mathcal{H}(tB)(y)\,dt
= \int_0^1 \sum_{y'} \partial_{B(y')}\mathcal{H}(tB)(y)\bigl[B(y')\bigr]\,dt ,
\end{equation*}
where $\partial_{B(y')}$ denotes the derivative in the value of $B$ at $y'$, a linear map on
$\mathfrak{g}^{c}$; the same identity holds with $\nabla\mathcal{H}$ in place of $\mathcal{H}$.
Because $\operatorname{supp}B\subset S$, only points $y'\in S$ contribute to the sums. By
\cite[(190)]{B-CMP102b}, uniformly for $tB$ in the tube,
\begin{equation*}
\bigl|\partial_{B(y')}\mathcal{H}(tB)(y)\bigr| \le C\ell^{-1}e^{-\frac18\delta_0 d(y,y')},
\qquad
\bigl|\partial_{B(y')}(\nabla\mathcal{H}(tB))(y)\bigr| \le C\ell^{-2}e^{-\frac18\delta_0 d(y,y')} ;
\end{equation*}
here the first entry of \cite[(190)]{B-CMP102b} bounds the kernel of the derivative of
$\mathcal{H}$ by
\begin{equation*}
O(1)\,(L^{j(y)}\eta)^{-1}(L^{j'}\eta)^{-d}\exp\bigl(-\tfrac18\delta_0 d(y,y')\bigr),
\end{equation*}
with $j'$ the level of $y'$, and the volume factor $(L^{j'}\eta)^{-d}$ of the kernel is absorbed by
differentiating in the point value $B(y')$. Inserting these bounds into the two integrals, the
integrand is bounded uniformly in $t\in[0,1]$, and the integral over $t$ gives
\eqref{9.3:eq-response-decay-1}.
\end{proof}

\begin{proof}[Proof of Corollary~\ref{9.3:cor-stationarity-level-data}, clause~(iv)]
\label{9.3:prf-local-backgrounds}
Throughout, $\eta=L^{-k}$ and $\eta'=\eta/L$. For a point $y$ of $\Omega$ we write
$\operatorname{dist}(y,\Omega^{c})$ for its distance to $\Omega^{c}$. We write
$B_k(\Omega)\subset T_{\eta}$ and $B_{k+1}(\Omega)\subset T_{\eta'}$ for the maximal
determining sets in $\Omega$ of \cite[\S2]{B-CMP119} at the two depths. Their domain sequences are written
$\Omega_1,\Omega_2,\dots$ and $\Omega'_1,\Omega'_2,\dots$ respectively.

\emph{Geometry of the two determining sets.} By maximality of $B_k(\Omega)$, the domain $\Omega_1$ is
$\Omega$ with one layer of cubes of side $L\eta M_1$ removed. For $n\ge2$, the domain $\Omega_n$ is $\Omega$
with two layers of cubes of side $L^{n}\eta M_1$ removed. To see the second statement, recall that
admissibility requires every cube of side $L^{n}\eta M_1$ in $\Omega_n$ to lie at distance at least
$L^{n}\eta M_1$ from $\Omega_{n-1}^{c}$, and that maximality takes $\Omega_n$ as large as this allows.
Count the layers of
cubes of side $L^{n}\eta M_1$ from the boundary of $\Omega$. A cube of the $m$-th layer lies at distance
\begin{equation*}
(m-1)L^{n}\eta M_1-\bigl(\text{width of }\Omega\setminus\Omega_{n-1}\bigr)
\end{equation*}
from $\Omega_{n-1}^{c}$. The width of $\Omega\setminus\Omega_{n-1}$ is positive and smaller than
$L^{n}\eta M_1$: by the description at the previous index it is $L\eta M_1$ for $n=2$ and
$2L^{n-1}\eta M_1$ for $n\ge3$. Hence the least $m$ for which this distance is at least $L^{n}\eta M_1$
is $m=3$. By maximality, exactly the first two layers are removed.

In $B_{k+1}(\Omega)\subset T_{\eta'}$ the cubes of level $n$ have side
$L^{n}\eta'M_1=L^{n-1}\eta M_1$. The same description therefore gives the following:
$\Omega'_1$ is $\Omega$ with one layer of cubes of side $\eta M_1$ removed; $\Omega'_2$ is $\Omega$ with
two layers of cubes of side $L\eta M_1$ removed; and $\Omega'_{n+1}=\Omega_n$ for $n\ge2$. So the two
determining sets differ only in their two finest layers.

We also use the identity
\begin{equation}\label{9.3:eq-local-backgrounds-1}
M(Q_{k+1}V)=Q_kV .
\end{equation}
Here $Q_kV$ denotes the configuration on $T_{\eta}$ attached to the datum $V$ in \cite[\S2]{B-CMP119},
and $Q_{k+1}V$ the corresponding configuration on $T_{\eta'}$; this is not the constraint map
$Q_j(U,\cdot)$ of Theorem~\ref{9.3:thm-stationarity-minimisers}.
The identity holds because, in the one-step average, contours inside a block see $1$ and contours
crossing between blocks see $V(c)$.

Put
\begin{equation*}
\widetilde{\boldsymbol{\Omega}}:=(\Omega'_1,\Omega'_2,\Omega_2,\dots,\Omega_k),
\end{equation*}
and let $\widetilde{\mathfrak{B}}$ be its determining set. The sequence
$\widetilde{\boldsymbol{\Omega}}$ is admissible. It is paired with $B_{k+1}(\Omega)$ as in
\eqref{9.3:eq-paired-variational-c}, with the integer $M_1$ of \cite{B-CMP119} in the role of the
product $RM_1$ of \cite{B-CMP102b}.

\emph{Step 1.} We apply Theorem~\ref{9.3:thm-stationarity-minimisers} (stationarity of the minimal
configurations) to the pair $(\widetilde{\boldsymbol{\Omega}},B_{k+1}(\Omega))$, for data satisfying
\eqref{9.3:eq-paired-variational-a} with $\varepsilon\le a_A$. Write $V^{1}$ for the datum on
$\widetilde{\mathfrak{B}}$ that this pairing attaches to the datum of $B_{k+1}(\Omega)$; it is the
one-step average $M(Q_{k+1}V)$, so by \eqref{9.3:eq-local-backgrounds-1} its entries at level $0$ are
those of $Q_kV$. The theorem yields a
gauge transformation $\mathsf{u}$ and a potential $A$ with $\|A\|_{\theta}\le C_A\varepsilon\eta^{\theta}$.
These relate the one-step average of $U_{k+1,\Omega}(V)$ to $U_{\widetilde{\mathfrak{B}}}(V^{1})$ in the
form stated there.

\emph{Step 2.} The sequence $\widetilde{\boldsymbol{\Omega}}$ refines $B_k(\Omega)$ in the sense of
Lemma~\ref{9.3:lem-refinement} (refinement of the determining set). Indeed, indexwise,
$\widetilde{\Omega}_0=\Omega'_1\subseteq\Omega_0=\Omega$, the difference being the outer strip of width
$\eta M_1$, which consists of level-$0$ bonds of $B_k(\Omega)$, on which the configuration is
prescribed by the datum itself; $\widetilde{\Omega}_1=\Omega'_2\subseteq\Omega_1$ (two layers of
cubes of side $L\eta M_1$ removed against one); and $\widetilde{\Omega}_j=\Omega_j$ for $j\ge2$.
Lemma~\ref{9.3:lem-refinement} therefore gives
\begin{equation*}
U_{k,\Omega}(V)=U_{\widetilde{\mathfrak{B}}}(\widetilde{V}),\qquad
\widetilde{V}:=M_{\widetilde{\mathfrak{B}}}\bigl(U_{k,\Omega}(V)\bigr).
\end{equation*}
The datum $\widetilde{V}$ coincides with the datum of Step~1 except on the strip
\begin{equation*}
S:=\{y:\ L\eta M_1\le\operatorname{dist}(y,\Omega^{c})<2L\eta M_1\}.
\end{equation*}
The strip $S$ lies in the level-$0$ layer of $\widetilde{\mathfrak{B}}$, and there $\widetilde{V}$ is
$U_{k,\Omega}(V)$ itself. On $S$ the configurations $\widetilde{V}$ and $Q_kV$ have the same $L$-averages,
and both are in the block-axial gauge. By \cite[Lemma~1]{B-CMP99a},
\begin{equation}\label{9.3:eq-local-backgrounds-2}
\bigl|\widetilde{V}(Q_kV)^{-1}-1\bigr|\le 4d^{2}\cdot 2B_3L^{2}\varepsilon<r_c .
\end{equation}
The second inequality in \eqref{9.3:eq-local-backgrounds-2} is a further smallness condition on
$\varepsilon$, imposed here in addition to $\varepsilon\le a_A$. Here $r_c$ is the coupling-free tube
radius of \eqref{9.3:eq-perfect-action-a}. Since $V^{1}=Q_kV$ on $S$ by Step~1, we define $B$ on $S$
as the logarithm $-i\log\bigl(\widetilde{V}(V^{1})^{-1}\bigr)$ and put $B=0$ off $S$; for
$\varepsilon$ below this ceiling, $\sup|B|$ is bounded by a constant times the left side of
\eqref{9.3:eq-local-backgrounds-2}, and the ceiling is taken so that this bound is less than $r_c$.
Thus $\widetilde{V}=e^{iB}V^{1}$ with
$\operatorname{supp}B\subset S$ and $\sup|B|<r_c$.

\emph{Step 3.} We apply Lemma~\ref{9.3:lem-response-decay} (decay of the response to local data) on
$\widetilde{\mathfrak{B}}$, with $V^{1}$ and $V^{2}:=\widetilde{V}=e^{iB}V^{1}$. Modulo gauge,
\begin{equation*}
U_{\widetilde{\mathfrak{B}}}(\widetilde{V})=e^{i\eta\mathcal{H}(B)}\,U_{\widetilde{\mathfrak{B}}}(V^{1}).
\end{equation*}
Moreover, both $\ell|\mathcal{H}|$ and $\ell^{2}|\nabla\mathcal{H}|$ at a point $y$ are bounded by
\begin{equation*}
C\sum_{y'\in S}e^{-\frac18\delta_0 d(y,y')}|B(y')| .
\end{equation*}
For $y$ of level $j\ge2$ we have $d(y,y')\ge d(y,S)\ge RM_1(j-2)$ for every $y'\in S$, by
\cite[(2.60)]{B-CMP96}. The floor \eqref{9.3:eq-paired-variational-d} states
$\delta_0RM_1\ge16\log L$, so
\begin{equation*}
e^{-\frac18\delta_0RM_1(j-2)}\le e^{-2(j-2)\log L}=L^{-2(j-2)} .
\end{equation*}
Hence at levels $j\ge2$ every term of the response to the change of datum on $S$ carries the factor
$L^{-2(j-2)}$.

We now combine the three steps. By Step~2, $U_{k,\Omega}(V)=U_{\widetilde{\mathfrak{B}}}(\widetilde{V})$.
By Step~3, at levels $j\ge2$ this configuration differs from $U_{\widetilde{\mathfrak{B}}}(V^{1})$, modulo
gauge, by the factor $e^{i\eta\mathcal{H}(B)}$. Since $\sup_S|B|\le C\varepsilon$ by
\eqref{9.3:eq-local-backgrounds-2}, each term of the sum in Step~3 is at most
\begin{equation*}
C\varepsilon L^{-2(j-2)}\le C\varepsilon L^{4}L^{-\theta j},
\end{equation*}
the inequality because $\theta<1$. By Step~1,
$U_{\widetilde{\mathfrak{B}}}(V^{1})$ is related to $U_{k+1,\Omega}(V)$ as in
Theorem~\ref{9.3:thm-stationarity-minimisers}. Composing the two relations gives, at levels $j\ge2$,
the bound of Theorem~\ref{9.3:thm-stationarity-minimisers} for the pair
$U_{k+1,\Omega}(V)$, $U_{k,\Omega}(V)$, with a constant $C$ in place of $C_A$. These are the local
backgrounds of \cite[(2.13)--(2.16)]{B-CMP119}.

\emph{The test functionals.} For the test functionals $t_k$ of \cite[(2.17)]{B-CMP119}, we apply
clause~(ii) (curvature) of
Corollary~\ref{9.3:cor-stationarity-level-data} on cubes $\square\subset\Omega_k$. These lie at level
$j=k$, where $\ell=L^{k}\eta=1$ and $L^{-\theta j}=L^{-\theta k}$. This gives
$|t_{k+1}-t_k|\le C\varepsilon L^{-\theta k}$.
\end{proof}

\begin{remark}[Stationarity beside large-field regions]\label{9.3:rem-large-field-regions}
Fix $\theta\in(0,1)$. Let $a_A$, $C_A$ be the constants of
Theorem~\ref{9.3:thm-stationarity-minimisers}, and $C_B$ that of
Theorem~\ref{9.3:thm-stationarity-jets}. Let $h$ be a large-field history whose domains form
an admissible decreasing sequence $\boldsymbol{\Omega}=(\Omega_j)_{0\le j\le k}$ of domains
in $T_{\eta}$. Write $\varepsilon_j$ for the small-field parameter of level $j$ in \cite{B-CMP119}.
\begin{itemize}
\item[(a)] Every background of a run whose large-field history is $h$, in the sense of
\cite[(2.10)--(2.12)]{B-CMP119}, is a minimal configuration $U_{\mathfrak{B}}(V)$. Its data
satisfy condition \eqref{9.3:eq-paired-variational-a} with parameter
\begin{equation}\label{9.3:eq-large-field-regions-1}
\varepsilon=O(L^2)\max_j\varepsilon_j .
\end{equation}
\item[(b)] Write $\Omega^{h}_j$ and $\Omega^{h'}_j$ for the domains of two histories $h$ and $h'$,
so that $\Omega_j=\Omega^{h}_j$. Call $h$ and $h'$ \emph{paired} if three conditions hold:
$\Omega^{h'}_{j+1}=\Omega^{h}_j$ for each level $j$ of $h$; $\Omega^{h'}_1$ contains the support of
the data of $h$; and the data of $h$ and $h'$ coincide at equal depth. For paired histories,
Theorems~\ref{9.3:thm-stationarity-minimisers} and~\ref{9.3:thm-stationarity-jets} hold as
stated, under their hypotheses. The bound they give at a point carries the factor $L^{-\theta j}$,
where $j$ is the level of that point itself.
\item[(c)] Consider a component of a large-field region of $h$ at level $j_0$. At every level
$j\ge j_0+1$, its imprint, that is, its contribution to the bounds of (b) at points of level $j$,
is at most $e^{-\delta_0RM_1(j-j_0-1)}L^{-\theta j_0}$, and
\begin{equation}\label{9.3:eq-large-field-regions-2}
e^{-\delta_0RM_1(j-j_0-1)}L^{-\theta j_0}\le L^{\theta}L^{-\theta j}.
\end{equation}
Here $\delta_0$ is the decay rate appearing in Lemma~\ref{9.3:lem-response-decay}, and $R$,
$M_1$ are the integers of the floor condition $(\mathrm{F})$, $\delta_0RM_1\ge16\log L$ (this $R$ is
Ba{\l}aban's integer, not the ball radius $R$ of the observables); the gaps
between the consecutive domains $\Omega_{j_0+1}\supset\dots\supset\Omega_k$, each measured in
lattice units of its own level, exceed $RM_1$.
\item[(d)] At the youngest levels, $0$ and $1$, there is no smallness, and none can be obtained.
At these levels the estimate of the large-field terms pays with the factor $e^{-p_0}$; this is the
factor that the hypothesis on the large-field terms at these levels prescribes.
\item[(e)] A component of $h'$ that has no partner in $h$ is a change of determining set and of
data. Let $\widetilde{\boldsymbol{\Omega}}=(\widetilde{\Omega}_j)$ be the sequence of domains
that describes the change of set; assume that it is admissible in $T_{\eta}$ and that
$\widetilde{\Omega}_j\subseteq\Omega_j$ for every $j$, and let $\widetilde{\mathfrak{B}}$ be its
determining set. The change of set leaves $U_{\mathfrak{B}}(V)$ minimal for
$\widetilde{\mathfrak{B}}$, with induced data
$\widetilde{V}=M_{\widetilde{\mathfrak{B}}}(U_{\mathfrak{B}}(V))$. Write the change of data as
$\widetilde{V}\mapsto e^{iB}\widetilde{V}$, with $B$ supported on a set $S$ of bonds and
$\sup|B|<r_c$, and write
\begin{align*}
U_{\widetilde{\mathfrak{B}}}(e^{iB}\widetilde{V})
&=e^{i\eta\mathcal{H}(B)}U_{\widetilde{\mathfrak{B}}}(\widetilde{V})\\
&=e^{i\eta\mathcal{H}(B)}U_{\mathfrak{B}}(V)\quad\text{modulo gauge},
\end{align*}
the first equality by Lemma~\ref{9.3:lem-response-decay}, the second by
Lemma~\ref{9.3:lem-refinement}. Let $\ell=L^{j(y)}\eta$ be the local scale at a point
$y$. Then
\begin{equation*}
\ell\,|\mathcal{H}(B)(y)|,\ \ell^{2}|\nabla\mathcal{H}(B)(y)|
\le C\,e^{-\frac18\delta_0 d(y,S)}\sum_{y'\in S}|B(y')| ,
\end{equation*}
where $d(y,S)=\min_{y'\in S}d(y,y')$ is the distance from $y$ to $S$ in the multi-scale
distance $d(\cdot,\cdot)$ of \cite{B-CMP96}, and $C$ is the constant of
Lemma~\ref{9.3:lem-response-decay}.
\item[(f)] Consider the chain of minimisers of the large-field operation $\mathbf{R}$ in
\cite[(1.50)--(1.58)]{B-CMP122b}. Theorems~\ref{9.3:thm-stationarity-minimisers}
and~\ref{9.3:thm-stationarity-jets} and Lemmas~\ref{9.3:lem-refinement}
and~\ref{9.3:lem-response-decay} apply to each configuration of this chain as an object in its
own right. The configuration $U_{1,2}$ of this chain is piecewise such a minimiser, and the four
statements apply to each piece.
The expansion \cite[(1.54)--(1.55)]{B-CMP122b} taken term by term is not covered.
\end{itemize}
\end{remark}

\begin{proof}
\emph{(a).} In the notation of \cite{B-CMP119}, the gauge field variables $V_{j-1}$ are defined
on the bonds of the set $\Omega_j^{(j-1)c}$; this set, the domain of definition of the variables
$V_{j-1}$ in that paper, contains the set
$\Gamma_{j-1}$ of that paper on which $V_{j-1}$ is assumed regular:
$\Gamma_{j-1}\subset\Omega_j^{(j-1)c}$. The regularity assumed on $\Gamma_{j-1}$ is
$|\partial V_{j-1}-1|<O(L^2)\varepsilon_{j-1}$. The data are supported on a small neighbourhood of
$\Omega_1$ that includes a layer of $M_1$-cubes \cite{B-CMP119}. Outside this neighbourhood no
condition is imposed \cite{B-CMP99b}.

Each background of \cite[(2.10)--(2.12)]{B-CMP119} is built from these data with the averages
prescribed on these bonds. It is therefore a minimal configuration $U_{\mathfrak{B}}(V)$, where
$\mathfrak{B}$ is the determining set carrying the prescribed averages. At level $j-1$ the
regularity assumption just stated is condition \eqref{9.3:eq-paired-variational-a} with parameter
$O(L^2)\varepsilon_{j-1}$. Since no condition is asked outside the support, taking the maximum
over the levels gives \eqref{9.3:eq-large-field-regions-1}.

\emph{(b).} Let $h$ and $h'$ be paired. The three pairing conditions are: the domains of $h'$ are
those of $h$ shifted by one level; $\Omega^{h'}_1$ contains the support of the data of $h$; and
the data coincide at equal depth. Together they say that the determining sets and data of $h$ at
spacing $\eta$ and of $h'$ at spacing $\eta/L$ form a pair of paired variational problems at
consecutive cutoffs, in the sense of Theorem~\ref{9.3:thm-stationarity-minimisers}.

By (a), the data satisfy \eqref{9.3:eq-paired-variational-a} with the parameter
\eqref{9.3:eq-large-field-regions-1}. Whenever $\varepsilon\le a_A$, the hypotheses of
Theorem~\ref{9.3:thm-stationarity-minimisers} hold, so its conclusion holds without change. The
same applies to Theorem~\ref{9.3:thm-stationarity-jets} under its hypotheses.

In Theorem~\ref{9.3:thm-stationarity-jets}(a), the bound at a bond $c$ of level $j$ is
$C_B\varepsilon L^{-\theta j}$. In Theorem~\ref{9.3:thm-stationarity-minimisers}, the norm
$\|A\|_{\theta}=\max\{|A|_{(-1-\theta)},|\nabla_UA|_{(-2-\theta)}\}$ is built from the
level-weighted norms $|\cdot|_{(\gamma)}$ of \cite[(3.41)]{B-CMP99b}. In both cases the
weight at a point is given by that point's own level, whatever the levels of the large-field
regions of $h$.

\emph{(c).} Let the component be at level $j_0$. The domains
$\Omega_{j_0+1}\supset\dots\supset\Omega_k$ are separated by gaps larger than $RM_1$, each gap
measured in the units of its own level. This geometry screens the component. By (b), its bound at
its own level is $L^{-\theta j_0}$. Passing to level $j$ crosses $j-j_0-1$ gaps, each of width
more than $RM_1$ against decay at rate $\delta_0$. Hence its imprint at level $j$ is at most the
left side of \eqref{9.3:eq-large-field-regions-2}.

For the inequality in \eqref{9.3:eq-large-field-regions-2}, write the right side as
$L^{-\theta(j-1)}$. Dividing both sides by $L^{-\theta j_0}$, the inequality becomes
\begin{equation*}
e^{-\delta_0RM_1(j-j_0-1)}\le L^{-\theta(j-j_0-1)} .
\end{equation*}
When $j=j_0+1$ both sides equal $1$ and there is nothing to prove. When $j-j_0-1\ge1$, this is
the inequality $e^{-\delta_0RM_1}\le L^{-\theta}$ raised to the power $j-j_0-1$, so it follows from
the condition $\delta_0RM_1\ge\theta\log L$. This condition
holds: the floor condition $(\mathrm{F})$ on these parameters gives $\delta_0RM_1\ge16\log L$, and
$16\log L>\theta\log L$ because $\log L>0$ and $\theta<1$.

\emph{(d).} Write $K_0$ and $K_1$ for the kernels of the quadratic fluctuation form of the
renormalisation step at levels $0$ and $1$. We use, as an input of this proof, the statement that
these two kernels are distinct, $K_1\neq K_0$.
This is a change of the quadratic form of relative size $O(1)$ at
levels $0$ and $1$. A bound carrying a small factor at these levels therefore cannot hold.
The estimate of the large-field terms at levels $0$ and $1$ pays there with the factor
$e^{-p_0}$ instead. That factor is the form in which the hypothesis on the large-field terms at
levels $0$ and $1$ is stated.

\emph{(e).} A component of $h'$ with no partner in $h$ breaks the pairing of (b) in two ways: it
changes the determining set, and it changes the data. The change of set is handled by the
refinement lemma (Lemma~\ref{9.3:lem-refinement}). Its hypotheses hold by the assumption made in
(e): $\widetilde{\boldsymbol{\Omega}}$ is admissible in $T_{\eta}$ and
$\widetilde{\Omega}_j\subseteq\Omega_j$ for every $j$. By that lemma, the minimal configuration
$U=U_{\mathfrak{B}}(V)$ is again minimal for the refined determining set
$\widetilde{\mathfrak{B}}$, with induced data $\widetilde{V}=M_{\widetilde{\mathfrak{B}}}(U)$;
that is, $U_{\widetilde{\mathfrak{B}}}(\widetilde{V})=U_{\mathfrak{B}}(V)$ modulo gauge.
The change of data is handled by the decay lemma (Lemma~\ref{9.3:lem-response-decay}), stated
for the same refined determining set $\widetilde{\mathfrak{B}}$. For data
$e^{iB}\widetilde{V}$ with $\sup|B|<r_c$ and $\operatorname{supp}B\subset S$, that lemma gives
\begin{equation*}
U_{\widetilde{\mathfrak{B}}}(e^{iB}\widetilde{V})
=e^{i\eta\mathcal{H}(B)}U_{\widetilde{\mathfrak{B}}}(\widetilde{V})
\quad\text{modulo gauge},
\end{equation*}
and, combined with the identity just obtained from the refinement lemma, this is the display
of (e). The decay lemma bounds both
$\ell|\mathcal{H}(B)(y)|$ and $\ell^{2}|\nabla\mathcal{H}(B)(y)|$ at $y$ by
\begin{equation*}
C\sum_{y'\in S}e^{-\frac18\delta_0d(y,y')}|B(y')| .
\end{equation*}
Since $d(y,y')\ge d(y,S)$ for every $y'\in S$, this sum is at most
$C\,e^{-\frac18\delta_0 d(y,S)}\sum_{y'\in S}|B(y')|$. Together the two lemmas therefore give
the bounds of (e).

\emph{(f).} The configurations in question are $U_k''$, $U_k^{0}$, $U_1$, $U_2$ of
\cite[(1.52)--(1.53)]{B-CMP122b} and $U_{k,\Lambda}$ of \cite[(1.56)]{B-CMP122b}. Each is a
minimiser on a determining set. As stated in \cite[(1.50)--(1.58)]{B-CMP122b}, these determining
sets are joined by an equation of the earlier paper of the series cited there.

Theorems~\ref{9.3:thm-stationarity-minimisers} and~\ref{9.3:thm-stationarity-jets} and
Lemmas~\ref{9.3:lem-refinement} and~\ref{9.3:lem-response-decay} concern minimisers on
determining sets. They therefore apply to each of these configurations taken as an object.
The configuration $U_{1,2}$ of \cite[(1.50)--(1.58)]{B-CMP122b} is piecewise a minimiser of
this kind, and the statements apply to each piece.

The expansion \cite[(1.54)--(1.55)]{B-CMP122b} is of a different kind. Taken term by term, it
expands, in the notation of \cite[(1.54)--(1.55)]{B-CMP122b}, the factor
\begin{equation*}
\exp\bigl(i\eta\,\mathbb{H}(\mathfrak{G}J_{1,2},H_1Q(\eta A_0))\bigr).
\end{equation*}
The symbols $J_{1,2}$ and $A_0$ in this display are those of that paper; in particular $A_0$
there denotes the object of that paper that appears in the argument of $Q$ in this display, not
the auxiliary parameter $A_0$ of $\mathfrak{p}$.
This expansion, taken term by term, is not treated here; the four statements above are applied
only to the minimisers of the chain.
\end{proof}

\begin{proposition}[Weak stationarity of the fine-lattice propagators. Proof outstanding]
\label{9.3:prop-propagator-stationarity}
Fix $\theta\in(0,1)$ and let $a_A$ be the threshold of
Theorem~\ref{9.3:thm-stationarity-minimisers}. There are constants $C<\infty$, $\theta'>0$ and
$\delta>0$ such that the following holds for
every $k\ge1$, every pair $(\mathfrak{B},\mathfrak{B}')$ of determining sets as in
\eqref{9.3:eq-paired-variational-c}, and every $V$ satisfying the condition $(7)_{\varepsilon}$
of that theorem on $\mathfrak{B}$, $|(\partial V)(p')-1|<\varepsilon$, with $\varepsilon\le a_A$.
Let $U_{\mathfrak{B}}=U_{\mathfrak{B}}(V)$ and the finer
minimal configuration $U'$ on $T_{\eta'}$ be as in Theorem~\ref{9.3:thm-stationarity-minimisers}.
Let $\mathsf{G}$ be one of the fine-lattice propagators $G'$, $G$, $G_1$ of \cite{B-CMP99b} or
the slice propagator $\mathfrak{G}$ of Lemma~\ref{9.3:lem-stability-equation}. Let
$Q_1^{\dagger}$ be the adjoint, for the pairings $\langle\cdot,\cdot\rangle_{\eta}$ and
$\langle\cdot,\cdot\rangle_{\eta'}$, of the linearised averaging $Q_1$ of
Lemma~\ref{9.3:lem-chain-lift}, so that $Q_1^{\dagger}$ carries functions on $T_{\eta}$ to
functions on $T_{\eta'}$. Let $y_1,y_2$ be centres of enlarged cubes $\widetilde{\Delta}(y_1)$,
$\widetilde{\Delta}(y_2)$ in $T_{\eta}$, let $f_1$, $f_2$ be functions on $T_{\eta}$ supported in
them, and let their scaled sup norms and
scaled Lipschitz norms be at most one. Then
\begin{equation}\label{9.3:eq-propagator-stationarity-1}
\begin{split}
&\bigl|\langle Q_1^{\dagger}f_1,\mathsf{G}(U')\,Q_1^{\dagger}f_2\rangle_{\eta'}
-\langle f_1,\mathsf{G}(U_{\mathfrak{B}})\,f_2\rangle_{\eta}\bigr|\\
&\qquad\le C\,L^{-\theta'\min j}\,\mathsf{S}_{\mathsf{G}}(y_1,y_2)\,e^{-\delta\, d(y_1,y_2)}.
\end{split}
\end{equation}
Here $\min j$ is the least level $j(x)$ of
the points of the supports, $d(\cdot,\cdot)$ is the multi-scale distance, and
$\mathsf{S}_{\mathsf{G}}(y_1,y_2)$ denotes the scale factors that
stand in Ba{\l}aban's bounds on $G'$, $G$, $G_1$ in \cite{B-CMP99b}, and the corresponding
factors for $\mathfrak{G}$.
\end{proposition}

Proof outstanding.

\medskip
\noindent\emph{Route.}
The bound \eqref{9.3:eq-propagator-stationarity-1} enters only the term-by-term comparison of
random walk expansions whose junctions $h_{\square}$, $K(h_{\square})$ lie on the fine lattice
\cite[(3.90), (3.107)]{B-CMP99b}. The first such expansion is the localisation of
$\log Z^{(k)}$ that \cite{B-CMP116} constructs from the generalised random walk expansion for
$Z^{(k)}(U_{k+1})$. The comparison of those expansions is what bounds the norm
$\|\widehat{L}^{(j+1)}-\widehat{L}^{(j)}\|$ of the difference of the localised one-loop families.
The bound also enters
the decoupled forms of $C^{(k)}$ and of $\Gamma_k$ introduced in \cite[(2.8)]{B-CMP116}
and the expansion of \cite[(1.54)--(1.55)]{B-CMP122b}.

It suffices to prove \eqref{9.3:eq-propagator-stationarity-1} with the last factor
$e^{-\delta d(y_1,y_2)}$ replaced by $1$. Ba{\l}aban's bounds on $\mathsf{G}$ decay in
$d(y_1,y_2)$ at the rate $\delta_0 M$, with $M=L^{m'}$ the auxiliary parameter of
$\mathfrak{p}$, and
this rate satisfies $\delta_0 M\ge2\kappa$; interpolating between the bound without decay and
this decaying bound gives \eqref{9.3:eq-propagator-stationarity-1} with a smaller exponent
$\theta'$ and a rate $\delta$ smaller than $\delta_0 M$.

The stationarity theorems proved above (Theorems~\ref{9.3:thm-stationarity-minimisers},
\ref{9.3:thm-stationarity-jets} and \ref{9.3:thm-stationarity-covariance}) do not yield
\eqref{9.3:eq-propagator-stationarity-1}, for two reasons. First, the gauge-fixing form
$\|R(U_{\mathfrak{B}})D^{*}A\|^2$ is not nested under the one-step averaging
map $M$ of Ba{\l}aban's construction, an object distinct from the parameter $M$ in the rate
$\delta_0M$. Second, a sandwich of a propagator by $Q$ on $\mathfrak{B}$ is a $V$-derivative of
the minimiser and is therefore covered by item~(c) of Theorem~\ref{9.3:thm-stationarity-jets},
whereas the fine-lattice junctions are not of that kind.

Two routes are available. The first proves \eqref{9.3:eq-propagator-stationarity-1}; the second
does not prove it, but obtains directly the localisation for which it is needed, so that
\eqref{9.3:eq-propagator-stationarity-1} is bypassed.

\emph{First route.} Rerun the chain lift through one averaging step
(Lemma~\ref{9.3:lem-chain-lift}), the face-mean multiplier and the current defect
(Lemma~\ref{9.3:lem-current-defect}), and the stability equation for the averaged minimiser
(Lemma~\ref{9.3:lem-stability-equation}) for the scalar form
\begin{equation}\label{9.3:eq-propagator-stationarity-2}
\langle\psi,(\Delta_U+Q'^{*}aQ')\psi\rangle .
\end{equation}
Here $\psi$ is a scalar function, $\Delta_U$ is the covariant lattice Laplacian in the background
$U$, $Q'$ is the scalar averaging map and $a$ is the weight of the averaging term (the letter $a$
has this meaning only here). For this form the following hold.
\begin{itemize}
\item $Q'$ nests exactly under one averaging step.
\item The lift is the block-constant extension.
\item The defect is the oscillation of $\nabla\psi$.
\item The bounds \cite[(3.42)--(3.45)]{B-CMP99b} serve for $G'$.
\end{itemize}
Then pass to the operator
\begin{equation}\label{9.3:eq-propagator-stationarity-3}
I-G'Q'^{*}\bigl(Q'G'^{2}Q'^{*}\bigr)^{-1}Q'G'
\end{equation}
of \cite[(3.25)]{B-CMP99b} and to the Neumann series \cite[(3.130), (3.138)]{B-CMP99b}.

\emph{Second route.} This route avoids \eqref{9.3:eq-propagator-stationarity-1}. By the statement
that only equivalence classes of step data, for the equivalence $\simeq$, are fed
forward to the next renormalisation step, it is enough to localise $\delta\varphi_l$ on the
unit lattice. To do so, expand the resolvent $(T_l+\lambda)^{-1}$, $\lambda\ge0$, where $T_l$ is
the sandwiched fluctuation operator of Theorem~\ref{9.3:thm-stationarity-covariance} at level $l$,
in a walk expansion built from local blocks $T_{\square}$. Here $T_{\square}$ is the sandwiched
fluctuation operator $T$ of Theorem~\ref{9.3:thm-stationarity-covariance}, built on the
determining set
$B_{l+1}(\square^{\sim})$. These blocks are exactly local in $\widetilde{W}$, and they are
stationary by the following three statements:
\begin{itemize}
\item item (iv) of the stationarity of level data, curvature and action
(Corollary~\ref{9.3:cor-stationarity-level-data});
\item the stationarity of the perfect action's jets (Theorem~\ref{9.3:thm-stationarity-jets});
\item the stationarity of the fluctuation covariance
(Theorem~\ref{9.3:thm-stationarity-covariance}).
\end{itemize}
Combine this expansion with
\begin{equation}\label{9.3:eq-propagator-stationarity-4}
\log\det T=\int_0^{\infty}d\lambda\,
\operatorname{Tr}\bigl[(1+\lambda)^{-1}-(T+\lambda)^{-1}\bigr].
\end{equation}
Then introduce one holomorphic parameter that multiplies all the differences $\delta T_{\square}$,
and apply the Schwarz lemma.

A telescoping over fattened cubes $\square^{\sim n}$ does not serve in place of either route.
The reason is that $d_k(\square^{\sim n})$ grows like $n^4$, while the available decay is only
$e^{-\delta_0 M n}$.

\begin{remark}[Remarks on analytic functionals and shell width]\label{9.3:rem-shell-width}
We use the notation of Theorem~\ref{9.3:thm-stationarity-minimisers} (index stationarity of the
minimal configurations) and of Corollary~\ref{9.3:cor-stationarity-level-data} (stationarity of level
data, curvature and action). For a potential $A$ we write
$\|A\|_{0}:=\max\{|A|_{(-1)},|\nabla_U A|_{(-2)}\}$, the norm $\|A\|_{\theta}$ at $\theta=0$.
\begin{enumerate}
\item[(a)] Theorem~\ref{9.3:thm-stationarity-minimisers} concerns data $V$ with values in
$\mathrm{SU}(N)$ only, not data in the complex tube $\mathbb{D}(\varrho)$. Let $V$
satisfy its hypotheses, put $U:=U_{\mathfrak{B}}(V)$, and let $W$, $\mathsf{u}$, $A$ be as there,
so that $W^{\mathsf{u}}=e^{i\eta A}U$. Let $\alpha_1>0$, and assume $\|A\|_{0}<\alpha_1$. Let $F$ be
a gauge-invariant function of configurations, analytic and of modulus at most $1$ on the
neighbourhood of $U$ of radius $\alpha_1$ in the sense of \cite[(3.37)]{B-CMP99b}. Then
\begin{equation}\label{9.3:eq-shell-width-1}
|F(W)-F(U)|=|F(e^{i\eta A}U)-F(U)|\le \frac{2\|A\|_{0}}{\alpha_1}.
\end{equation}
A functional needed at complex data is treated by applying Lemma~\ref{9.3:lem-two-constants}
(a two-constants lemma on the disc) to the scalar holomorphic function it defines.
\item[(b)] Let $t_k(V)$, $t_{k+1}(V)$ be the test functionals of
Corollary~\ref{9.3:cor-stationarity-level-data}\,(iv), and let $\varepsilon_k>0$ be the threshold
in the small-field characteristic functions $\chi_k$. For data $V$ satisfying the hypotheses of
that clause with $\varepsilon=\varepsilon_k$, the symmetric difference of
$\{t_k<\varepsilon_k\}$ and $\{t_{k+1}<\varepsilon_k\}$ is contained in the shell
$|t_k-\varepsilon_k|\le C\varepsilon_kL^{-\theta k}$.
\end{enumerate}
\end{remark}

\begin{proof}
(a) Since $F$ is gauge invariant, $F(W)=F(W^{\mathsf{u}})=F(e^{i\eta A}U)$, which is the equality
in~\eqref{9.3:eq-shell-width-1}. If $\|A\|_{0}=0$ then $A=0$, so $e^{i\eta A}U=U$ and the inequality
holds. Otherwise put $R:=\alpha_1/\|A\|_{0}>1$. For complex $t$ with $|t|\,\|A\|_{0}<\alpha_1$,
that is $|t|<R$, the configuration $e^{it\eta A}U$ lies in the neighbourhood of $U$ of radius
$\alpha_1$; hence $\varphi(t):=F(e^{it\eta A}U)$ is holomorphic on $|t|<R$ with $|\varphi|\le1$
there. The function
$\psi(t):=(\varphi(t)-\varphi(0))/2$ is holomorphic on $|t|<R$, satisfies $|\psi|\le1$ there and
$\psi(0)=0$. The Schwarz lemma applied to $s\mapsto\psi(Rs)$ on the unit disc gives
$|\psi(t)|\le|t|/R$ for $|t|<R$. At $t=1$, admissible because $R>1$,
\begin{equation*}
|\varphi(1)-\varphi(0)|\le \frac{2}{R}=\frac{2\|A\|_{0}}{\alpha_1},
\end{equation*}
and $\varphi(0)=F(U)$, $\varphi(1)=F(e^{i\eta A}U)$.

(b) Let $V$ lie in the symmetric difference. Either $t_k<\varepsilon_k\le t_{k+1}$, and then
$0<\varepsilon_k-t_k\le t_{k+1}-t_k$; or $t_{k+1}<\varepsilon_k\le t_k$, and then
$0\le t_k-\varepsilon_k<t_k-t_{k+1}$. In both cases
$|t_k-\varepsilon_k|\le|t_{k+1}-t_k|$. By Corollary~\ref{9.3:cor-stationarity-level-data}\,(iv)
with $\varepsilon=\varepsilon_k$, $|t_{k+1}-t_k|\le C\varepsilon_k L^{-\theta k}$, which gives the
shell.
\end{proof}

\label{9.4:par-how-clause-v-consumes-the-two-cutoff-res}
Clause~(v) of Theorem~\ref{3.1:thm-main} uses the following statements of Chapter~5, exactly
as they are stated there and not proved again: the two-cutoff response statement at zeroth order
and the two routes to the two-cutoff response. Beside the statement at zeroth order it reads the
flat energy-difference bound of Chapters~5 and~6, which is used by reference only. The statement
pairing the two runs is used by its statement alone, as one of the named hypotheses from which the
two-cutoff response statement at zeroth order is proved as an implication.

Clause~(v) reads these inputs in three steps. First, the correlation theorem is applied to the
second of the two runs of Definition~\ref{9.6:def-comparison-chain}. It supplies the quantities of
that run to which the statement at zeroth order is applied. Next, these quantities are calibrated
between the two runs. Finally, the exact-pin statement of the uniqueness theorem turns the
calibrated comparison into the uniqueness asserted in clause~(v).

That the statement at zeroth order can be read on the exactly paired shells on which clause~(v)
compares the two runs, with no estimate beyond those it contains, is not proved. It is listed in
the index of Chapter~\ref{ch:13}. This reading requires at least three things:
\begin{itemize}
\item the response functionals it uses transform covariantly under the change of base
configuration;
\item their domains of analyticity still contain the re-based arguments after the radii are shrunk
by a fixed factor, a shrink incurred once per functional and not once per renormalization step;
\item the resulting constant is absorbed into the constants of that statement without a further
smallness condition.
\end{itemize}
It also requires identifications of the points at which clause~(v) reads that statement, which
are likewise not proved. If the three requirements fail, an estimate of the response in supremum
form is needed, and that estimate is not proved either.

The clauses of Theorem~\ref{3.1:thm-main} are proved in the following order.
\begin{description}
\item[(0)] Clause~(0) is proved first.
\item[(i)] Clause~(i) follows clause~(0).
\item[(iii)] Clause~(iii) is proved before clause~(ii), because clause~(ii) uses it.
\item[(ii)] Clause~(ii) is proved next. Its torus sub-clauses (ii-$\mathbb{T}$-a) and
(ii-$\mathbb{T}$-b) use the identification asserted in clause~(v), by its statement.
\item[(iv)] The proof of clause~(iv) uses the two-point witness theorem and the volume transfer of
Chapter~8. It does not use an estimate, uniform in the volume, of the never-healed old-age tail;
such an estimate is not proved in this book. This tail is the contribution of
large-field regions that were created more than the prescribed number of steps before the witness
depth (the depth below the cutoff fixed in Chapter~8 for the two-point witness theorem) and have
not been re-absorbed into the regular expansion by then.
\item[(v)] Clause~(v) is proved as an implication, in the three steps above, from the following
statements:
\begin{itemize}
\item the correlation theorem;
\item the two-cutoff response statement at zeroth order, with the two routes to the two-cutoff
response;
\item the flat energy-difference bound;
\item the reading on exactly paired shells described above, which is not proved;
\item the uniqueness theorem with its named hypotheses, among them that the shells of the two runs
are exactly paired.
\end{itemize}
\item[(vi)] Clause~(vi) is conditional on the far-sphere flux hypothesis
$(\mathrm{IR}\text{-}\mathrm{fl})_{\omega_0}$ of Theorem~\ref{3.1:thm-main}: the flux of the
lattice rotation current through far spheres vanishes along one sequence of radii. That theorem
displays beside it a sufficient condition in the form of a stress-tensor decay bound. Under this
bound, one stress-tensor insertion at distance $|x|$ from the supports of the observables
decorrelates from them faster than $|x|^{-4}$, uniformly in the cutoff.
\end{description}

Acyclicity concerns both the constants and the clauses.

Every condition on the constants displayed in the statements of Chapter~5 is a condition of the
order of choice. That order has no cycle and imposes only one-sided conditions
(Lemma~\ref{2.4:lem-no-cycle}), and its conditions can be met simultaneously
(Proposition~\ref{2.4:prop-order-consistent}). Two conditions met in the reading of clause~(v) are
not shown here to be of this kind. The first is the shrink of the radii once per functional, one
of the three requirements stated above. The second is a further condition, that no constant
compounds with the number of repetitions of the step of the two-run comparison that carries a
complex dilation parameter. If either condition failed, it would bound the volume. Neither is
proved.

Among the clauses, the one forward appeal is that of the torus sub-clauses of clause~(ii) to the
identification asserted in clause~(v). With that appeal the order above is the order of proof.
Clause~(v) is proved from the statements named in item~(v) and not from clause~(ii), so no cycle
arises.

\begin{remark}\label{9.4:rem-the-zeroth-order-response-read-on-the-sh}
The uniqueness theorem compares two runs of
Ba{\l}aban's renormalization group whose numbers of renormalization steps differ by one (the
runs of Definition~\ref{9.6:def-comparison-chain}). It
compares them on shells paired exactly between the two runs. One statement about these shells is
set down here because the deduction of clause (v) uses it. Its content is as follows.
On these shells, the two-cutoff response statement at zeroth order holds for every object stored
in the comparison, after that object is re-expressed on the common base configuration through the
gauge covariance of its arguments. This re-expression requires no estimate beyond the
implicit-function lemma that already pays every change of base on the route of the uniqueness
theorem. Its only price is a shrink of the radii of analyticity by the radius factor
$(1-C_g M\alpha_1)$ of the covariance identity, under that identity's standing condition
$C_g M\alpha_1<1$. The shrink is incurred once per object and not once per
renormalization step. The constant of the implicit-function lemma that pays for the
re-expression does not depend on $\kappa_1$. It is absorbed by the smallness condition already
among the hypotheses of the zeroth-order statement, with no further condition.
This reading is not proved. Proof outstanding.
Route: run the zeroth-order argument of the two-cutoff response statement on those shells, with
the radii shrunk by the factor $(1-C_g M\alpha_1)$ once per object. The reading is used in the
deduction of clause (v), and through the uniqueness theorem in the identification of part
(ii-$\mathbb{T}$-a) of clause (ii). Without it, that use would require instead a further estimate
in supremum form among the inputs of the uniqueness theorem, and that estimate is not proved
either. Clause by clause, in the order of Theorem~\ref{3.1:thm-main}, the position is as
follows.

\begin{description}
\item[(0)] Clause (0) is deduced without the reading.
\item[(i)] Clause (i) is deduced without the reading.
\item[(iii)] Clause (iii) is deduced without the reading.
\item[(ii)] Clause (ii) uses clause (iii). Its identification in part (ii-$\mathbb{T}$-a) uses the
uniqueness theorem, and through that theorem it uses the reading. The use of clause (iii) does
not involve the reading.
\item[(iv)] Clause (iv) is deduced without the reading.
\item[(v)] Clause (v) is deduced from the uniqueness theorem. That theorem uses, on the exactly
paired shells, the two-cutoff response statement at zeroth order together with the flat
energy-difference bound. For the objects stored in the comparison, it uses the former through the
reading, or else through the further estimate in supremum form described above. Clause (v)
therefore rests on the reading.
\item[(vi)] Clause (vi) is deduced under the far-sphere flux hypothesis
$(\mathrm{IR}\text{-}\mathrm{fl})_{\omega_0}$ of Theorem~\ref{3.1:thm-main}: the flux of the
lattice rotation current through far spheres vanishes along one sequence of radii. A sufficient
condition for it is that one stress-tensor insertion placed far away decorrelates from the
observables faster than the inverse fourth power of its distance, uniformly in the cutoff. The
reading is not used directly in the deduction of clause (vi).
\end{description}

Adding the reading creates no cycle among the deductions of this chapter. The reading is used
only in the deduction of clause (v) and, through the uniqueness theorem, in that of part
(ii-$\mathbb{T}$-a). The route
indicated for the reading rests on the zeroth-order argument of the two-cutoff response
statement, which depends on neither of these clauses. Clause (ii) uses clause (iii), and clause (iii)
is deduced without clause (ii); this is why (iii) is listed before (ii). The order in which the
constants of Theorem~\ref{3.1:thm-main} are chosen has no cycle either
(Lemma~\ref{2.4:lem-no-cycle}).
\end{remark}

\begin{proposition}[The two-cutoff smallness, at the current depth, of the two factors kept
un-integrated inside a thrown-back region, on the enlarged class, under one ceiling. Proof
outstanding]
\label{9.5:prop-the-two-cutoff-smallness-of-the-kept-fac}
Consider two runs of the renormalization group, at two values of the cutoff with their bare
couplings, carried on one common sequence of large-field regions and localization domains.
Call a component of the large-field region a thrown-back region at step $j$ if the large-field
operation acts on it at step $j$. Consider such a region that still carries its large-field
factor at a later step $k > j$. Two factors are kept un-integrated over the region. The first is
the quadratic form kept by the large-field operation, built on the region's own background and
localized to the region. The second is the sum of the
effective-action terms of the earlier steps together with the current step's localized remainders
of the action expansion, localized in the region enlarged by a collar of prescribed width around
it.

For each of these two factors, form the difference between the two runs. Measure it by its
oscillation over the enlarged class of pairs of configurations on which the uniqueness theorem
compares the two runs, one configuration for each run.
The oscillation is taken with the real field on the region itself held fixed at one value common
to both runs and lying in the supports of both runs' densities, and the supremum is then taken
over such values. Divide the sum of the two oscillations by the
fixed power of the radius parameter of the region's own step $j$ that normalizes the bound on the
kept factors.
Here the radius parameter of step $j$ is a prescribed function of the running coupling at that
step, not exceeding the corresponding radius of Ba{\l}aban's construction; it is also taken as
the number of preparatory steps before the large-field operation at that step. Then take the
supremum over all regions still carrying their large-field factor at step $k$.

The assertion is that, at the current step $k$, this supremum does not exceed the threshold on
which the two-cutoff comparison in the proof of the uniqueness theorem relies, that is, the bound
which that comparison requires of the two kept factors. More exactly, for
every pair of cutoffs and every level $k$ of the common sequence of regions,
it does not exceed that threshold under a single one-sided ceiling on the
coupling. This ceiling is placed after the constant of the bound on the kept factors in the order
of choice.
\end{proposition}

\noindent
\emph{Route.} Prove the smallness at the current depth layer by layer: split the regions
according to the step at which they were created; bound each layer by the defect of the
two-cutoff comparison at that layer, that is, by the quantity with which the induction in the
proof of the uniqueness theorem measures the difference of the two runs at that step, and which
vanishes when the two cutoffs coincide; and pay the number of steps since
creation once by the region's large-field factor, taken without normalization. The current
step's remainder terms of the action expansion enter through one further input, not proved: the
bound, by the same defect, of the difference between the two runs of the remainder terms of the
basic step of the large-field operation of which \cite{B-CMP122a} writes ``we omit details'',
localized in the enlarged region and measured and normalized as above.

\begin{proposition}[The mixed-size counting estimate. Proof outstanding]\label{9.5:prop-the-mixed-first-generation-bound}
The statement of this proposition is not written out here.
\end{proposition}

\noindent\emph{Route.} Cover the configurations of blocks of mixed sizes that are to be counted by balls of their
local scales. Arrange the balls along a depth-first circuit of the tree formed by their overlaps,
and bound the count step by step along that circuit. The argument involves no field variable and
uses two one-sided lower bounds on $M$.

Proposition~\ref{9.5:prop-the-mixed-first-generation-bound} is an input of the proof of
clause~(v). Besides this proposition and the uniqueness theorem for $N\ge 3$, three statements of
clause~(v)'s own making have outstanding proofs.

The first bounds the two-run difference, along the comparison chain of
Definition~\ref{9.6:def-comparison-chain}, of the field-free remainder terms of the basic step of
$\mathbf{R}$ for which \cite{B-CMP122a} writes ``we omit details''. These terms are taken near a
large-field region that the step leaves un-integrated, and they are split by the lowest level
that they meet.

The second bounds the two-run difference of the Gaussian factors, quadratic in the real variables
common to the two runs, that are kept inside such a region by large-field steps interrupted
there, one nested in another.

The third bounds the two-run difference of one part of each boundary term produced at a step of
$\mathbf{R}$. This part is carried by the small-field region of that step, or by large-field
regions whose large-field factor is absorbed back into the small-field format at that step.

For each clause of Theorem~\ref{3.1:thm-main}, the item below says whether its proof depends on
these four statements, namely the proposition and the three statements just described. Where that
dependence runs through other statements of this chapter, it is given in the list at the end of
this chapter.

\begin{description}
\item[(0)] The proof of clause~(0) uses none of the four statements.
\item[(i)] The proof of clause~(i) uses none of the four statements.
\item[(iii)] The inputs of the proof of clause~(iii) include the correlation theorem. That proof
uses none of the four statements.
\item[(ii)] The infinite-volume part of clause~(ii) is deduced from clause~(iii) and uses none of
the four statements. The torus sub-clause (ii-$\mathbb{T}$-a) uses statements introduced for
the proof of clause~(v); each appears in the list at the end of this chapter.
\item[(iv)] Among the clauses, the proof of clause~(iv) uses only clauses~(0) and~(iii) and the
lower bound of clause~(i). Its other inputs are the correlation theorem and the two-point
witness theorem, together with the statements on which that theorem depends. It uses neither
clause~(v) nor the four statements.
\item[(v)] All four statements are inputs of the proof of clause~(v). Among the clauses, some
statements used in that proof use clauses~(iii) and~(iv); none uses clause~(vi).
\item[(vi)] Clause~(vi) is a complete deduction from the far-sphere flux hypothesis
$(\mathrm{IR\text{-}fl})_{\omega_0}$ of that clause and from the following: clauses~(0), (i)
and~(iii), the infinite-volume part of clause~(ii), and the one-lattice Lipschitz step, on which
clause~(v) depends as well. The list at the end of this chapter states whether the proof of the
one-lattice Lipschitz step uses any of the four statements.
\end{description}

The dependence among the clauses has no cycle, for four reasons:
\begin{itemize}
\item the proof of clause~(ii) uses clause~(iii), which is why (iii) is stated before (ii);
\item the torus sub-clause of~(ii) uses intermediate statements introduced for clause~(v), not the
conclusion of clause~(v);
\item the proofs of clauses~(iv) and~(v) use only clauses stated before them;
\item the deduction of clause~(vi) does not use the conclusion of clause~(v).
\end{itemize}

All constants on which these proofs depend are chosen in one order. The conditions of the order
of choice of Definitions~\ref{2.3:not-stages-first}, \ref{2.3:not-stages-middle}
and~\ref{2.3:not-stages-last} have no cycle and are one-sided, by Lemma~\ref{2.4:lem-no-cycle}.
The proofs of clauses~(iv) and~(v) add further conditions, among them the two lower bounds on $M$
used in the route to Proposition~\ref{9.5:prop-the-mixed-first-generation-bound}. Each added
condition is of one of two kinds:
\begin{itemize}
\item a lower bound on an auxiliary constant, placed after the constants it reads;
\item an upper bound on the coupling threshold $\gamma$, chosen last.
\end{itemize}
None of them is a cap in the sense of Definition~\ref{2.3:def-cap}.

Five statements of clause~(v)'s own making have an outstanding proof, and they are the only such
statements:
\begin{itemize}
\item the uniqueness theorem for $N\ge 3$;
\item Proposition~\ref{9.5:prop-the-mixed-first-generation-bound};
\item the three statements described above.
\end{itemize}
Every other such statement is proved, proved as an implication from named statements, or stated
with proof incomplete at a marked step, as listed at the end of this chapter.

\begin{definition}[Parameters of the two-lattice comparison]\label{9.5:not-comparison-parameters}
Throughout the two-lattice comparison the dimension is $d=4$. The following parameters are
kept symbolic: none of them is given a numerical value.
\begin{itemize}
\item the blocking factor $L$, an odd integer with $L\ge 5$;
\item the auxiliary parameters $M$, $M_1$, $\kappa$, $\kappa_1$ and the parameter $R_1$;
\item the integers $N$, $N_0$, the number of steps $K$ and the level $k$;
\item $\check{R}=\check{R}_k$;
\item the running couplings $g_k$ and the parameter $\gamma$;
\item an exponent $\theta$ with $3/5<\theta<1$.
\end{itemize}
From $\theta$ we form the exponents
\begin{equation}\label{9.5:eq-comparison-parameters-1}
\theta_2:=\frac{5\theta-3}{8},\qquad
\vartheta:=L^{-\theta_2},\qquad
\beta:=\frac{1+\theta}{2}.
\end{equation}
Since $3/5<\theta<1$, we have $0<5\theta-3<2$, hence $0<\theta_2<1/4$. Moreover $1/4<3/5<\theta$,
and therefore
\begin{equation}\label{9.5:eq-comparison-parameters-2}
0<\theta_2<\theta .
\end{equation}
Finally, $\delta_0$ denotes the decay constant of the propagators in Ba{\l}aban's construction,
and we put
\begin{equation}\label{9.5:eq-comparison-parameters-3}
\delta_1:=\frac{\delta_0}{32}.
\end{equation}
\end{definition}

\begin{definition}[The two lattices, their pairs and the mismatch field]\label{9.5:def-mismatch-field}
Let $G$ be a compact Lie group with Lie algebra $\mathfrak{g}$, and let $G^{\mathbb{C}}$ and
$\mathfrak{g}^{\mathbb{C}}$ be their complexifications. Let $|\cdot|$ be an
$\operatorname{Ad}(G)$-invariant norm on $\mathfrak{g}^{\mathbb{C}}$, extended to linear maps of
$\mathfrak{g}^{\mathbb{C}}$ as the associated operator norm.

\emph{The two machines.} Two \emph{machines} work on one scaffold, that is, with one and the same
cover and one and the same set of rules for the scaffold, in the sense fixed with the parameters
of the two-lattice comparison (Notation~\ref{9.5:not-comparison-parameters}). In
Definition~\ref{9.6:def-comparison-chain} the same two machines are called the two runs.
\begin{itemize}
\item The machine $\mathsf{A}$ works on the lattice of spacing $\eta=L^{-k}$ on the torus
$\mathbb{T}_m$, written $T_\eta$ (the letter $T$ with a spacing subscript denotes a lattice, not
the degree variable of Notation~\ref{9.1:not-levelwise-rate}).
\item The machine $\mathsf{B}$ works on the lattice $T_{\eta'}$ of spacing $\eta'=\eta/L$ on the
same torus.
\end{itemize}
A \emph{configuration} on a region is a map from the bonds of that region to $G^{\mathbb{C}}$.

\emph{The averaging map.} We write $M$ for Ba{\l}aban's one-step block-averaging map from
configurations on $T_{\eta'}$ to configurations on $T_\eta$
(Definition~\ref{1.5:def-block-average}). Throughout the two-lattice comparison of this chapter
the letter $M$ denotes this map and not the cube size $L^{m'}$. The map $M$ has three properties.
\begin{itemize}
\item It is local: for a bond $x$ of $T_\eta$, $M(\mathbf{U}')(x)$ depends on $\mathbf{U}'$ only
through its values on the fine bonds of a fixed finite neighbourhood of $x$. The set of these fine
bonds is called the \emph{stencil} of $x$ and is written $\operatorname{st}(x)$.
\item It is covariant.
\item It is analytic.
\end{itemize}

\emph{Points and pairs.} A \emph{point} is a couple $P=(\mathbf{U},\mathbf{J})$ of a
configuration $\mathbf{U}$ and a \emph{slot content} $\mathbf{J}$. The slot content is a
$\mathfrak{g}^{\mathbb{C}}$-valued bond field, namely the variable $\mathbf{J}$ of
\cite{B-CMP109}. A \emph{pair} is a couple $p=(P',P)$ with the following two members:
\begin{itemize}
\item a point $P$ on a region of $T_\eta$;
\item a point $P'=(\mathbf{U}',\mathbf{J}')$ on the same region of $T_{\eta'}$.
\end{itemize}

\emph{The mismatch.} The \emph{mismatch} of the pair $p$ is the $\mathfrak{g}^{\mathbb{C}}$-valued
bond field
\begin{equation}\label{9.5:eq-mismatch-field-a}
A_p(x):=(i\eta)^{-1}\log\bigl[M(\mathbf{U}')(x)\cdot\mathbf{U}(x)^{-1}\bigr],
\end{equation}
where $\log$ is the \emph{principal logarithm}, fixed as follows. Let $\mathcal{B}$ be the union
of all open balls in $\mathfrak{g}^{\mathbb{C}}$ centred at $0$ (for the norm $|\cdot|$) which
$\exp\colon\mathfrak{g}^{\mathbb{C}}\to G^{\mathbb{C}}$ maps diffeomorphically onto open subsets
of $G^{\mathbb{C}}$; there is at least one such ball, since the differential of $\exp$ at $0$ is
the identity. These balls are nested, so $\mathcal{B}$ is an open ball centred at $0$ or all of
$\mathfrak{g}^{\mathbb{C}}$; any two of its points lie in one of the balls, so $\exp$ is
injective on $\mathcal{B}$, and it is a local diffeomorphism there, hence a diffeomorphism of
$\mathcal{B}$ onto the open neighbourhood $\exp(\mathcal{B})$ of the identity. The principal
logarithm is the inverse of this diffeomorphism. The field $A_p$ is defined at every bond $x$ of
the region at which $M(\mathbf{U}')(x)\,\mathbf{U}(x)^{-1}$ lies in $\exp(\mathcal{B})$.

\emph{The variation of the averaging map under a fine increment.} Let $h$, the \emph{increment},
be a $\mathfrak{g}^{\mathbb{C}}$-valued bond field on the fine
lattice and let $V$ be a configuration on the fine lattice, such that $V$ and the fine
configuration $e^{i\eta' h}V$ defined below both lie in the domain on which $M$ is analytic, and
such that the principal logarithm below is defined at every coarse bond $x$, that is,
$M\bigl(e^{i\eta' h}V\bigr)(x)\,M(V)(x)^{-1}$ lies in $\exp(\mathcal{B})$. We put
\begin{equation}\label{9.5:eq-mismatch-field-b}
Y[h;V](x):=(i\eta)^{-1}\log\bigl[M\bigl(e^{i\eta' h}V\bigr)(x)\cdot M(V)(x)^{-1}\bigr],
\end{equation}
where $e^{i\eta' h}V$ is the fine configuration defined by
$(e^{i\eta' h}V)(y):=e^{i\eta' h(y)}\,V(y)$ for every fine bond $y$. For such $h$ and $V$ the
field $Y[h;V]$ is local, since by the locality of $M$ its value at $x$ depends on $h$ and $V$
only through their values on $\operatorname{st}(x)$; it is analytic, since $M$ is analytic and
the exponential, products, inverses and the principal logarithm are analytic where they are
defined; and $Y[0;V]=0$, since $M(V)(x)\,M(V)(x)^{-1}$ is the identity, whose principal logarithm
is $0$. By the variation bound for the one-step averaging map, for every coarse bond $x$ it
satisfies
\begin{equation}\label{9.5:eq-mismatch-field-1}
\bigl|Y[h;V](x)\bigr|\le 34\,d\,\sup_{y\in\operatorname{st}(x)}|h(y)| .
\end{equation}
The slot contents $\mathbf{J}$ and $\mathbf{J}'$ never enter the mismatch
\eqref{9.5:eq-mismatch-field-a}.

\emph{The frame rule.} The representative of the machine $\mathsf{B}$ is rotated once, at the
root or at the seed of a presentation, by a fixed site field $g$ (in the two-lattice comparison
the letter $g$ denotes this field and not the
reference coupling). A source datum $b$ common to
both machines is read as the couple $(\operatorname{Ad}_g b,\,b)$, the first member by
$\mathsf{B}$, whose representative is rotated, and the second by $\mathsf{A}$, in the order of
the members of a pair. We write
\begin{equation}\label{9.5:eq-mismatch-field-2}
\delta b:=(\operatorname{Ad}_g-1)\,b .
\end{equation}
\end{definition}

\begin{definition}[Admissible sets, boxes, shells and fat cores]\label{9.5:def-admissible-sets}
We use the two machines $\mathsf{A}$ and $\mathsf{B}$, the scaffold on which they work, the spacing
$\eta=L^{-k}$ of the lattice of $\mathsf{A}$ and the spacing $\eta'=\eta/L$ of the lattice of
$\mathsf{B}$, all as in Definition~\ref{9.5:def-mismatch-field}.

\emph{Admissible sets.} An \emph{admissible set} $\mathfrak{S}$ is a decreasing sequence of regions
$(\Omega_j)_{j\ge0}$ with the block and spacing structure of Ba{\l}aban's sequences of regions.
An admissible set determines the following objects.
\begin{itemize}
\item The \emph{level function} $\lambda_{\mathfrak{S}}$: at a bond, $\lambda_{\mathfrak{S}}$ is the
largest $j$ such that the bond lies in $\Omega_j$.
\item The \emph{local scale} $\ell_{\mathfrak{S}}(x):=L^{\lambda_{\mathfrak{S}}(x)}\eta$.
\item The \emph{local cubes} $\tilde{\Delta}_{\mathfrak{S}}(y)$, whose side is a fixed multiple of
$\ell_{\mathfrak{S}}(y)$.
\item The \emph{scaled path metric} $d_{\mathfrak{S}}$, in which the length of a path is measured in
local units, with Ba{\l}aban's weights.
\end{itemize}
A pair $(\mathfrak{S}',\mathfrak{S})$ consists of an admissible set
$\mathfrak{S}'=(\Omega'_j)_{j\ge0}$ for $\mathsf{B}$ and an admissible set
$\mathfrak{S}=(\Omega_j)_{j\ge0}$ for $\mathsf{A}$. Such a pair is \emph{exactly paired} if
\begin{equation}\label{9.5:eq-admissible-sets-1}
\Omega'_{j+1}=\Omega_j\qquad\text{for all }j\ge0 .
\end{equation}
The outermost stratum of $\mathsf{B}$, namely $\Omega'_0\setminus\Omega'_1$, is then an extra
stratum that has no counterpart in $\mathfrak{S}$, and it lies outside $\Omega_0$.
Of these objects only four statements are used in this chapter; they are
stated later in the chapter.

\emph{Labels.} The \emph{label} at a site is the current admissible set of the scaffold at that
site. Its level function is written $\iota$ and its local scale $\ell_\iota$.

\emph{Boxes and shells.} A \emph{box} is a union $Z$ of cubes of the cover of
Definition~\ref{9.5:def-mismatch-field} at the proper scale. For
$s\ge0$ let $Z^{-[s]}$ be the set of points of $Z$ whose distance from $Z^{c}$ is at least $s$. The
\emph{shells of $Z$ over} $\mathfrak{S}$ form the admissible set
\begin{equation}\label{9.5:eq-admissible-sets-2}
\boldsymbol{\Omega}_{\mathfrak{S}}(Z):=\bigl(\Omega_n\cap Z^{-[s_n]}\bigr)_{n\ge0},
\end{equation}
where $(s_n)_{n\ge0}$ is the sequence of shell widths of the block geometry. We put
\begin{equation*}
\lambda_Z:=\lambda_{\boldsymbol{\Omega}_{\mathfrak{S}}(Z)},\qquad
\ell_Z:=L^{\lambda_Z}\eta,\qquad
d_Z:=d_{\boldsymbol{\Omega}_{\mathfrak{S}}(Z)} .
\end{equation*}
Since $\Omega_n\cap Z^{-[s_n]}\subset\Omega_n$ for every $n$, a bond of $Z$ that lies in the $n$-th
shell lies in $\Omega_n$. Hence the largest index of a shell containing it does not exceed the
largest index of a region $\Omega_n$ containing it, that is, $\lambda_Z\le\lambda_{\mathfrak{S}}$
pointwise on $Z$. Because $L>1$, it follows that $\ell_Z\le\ell_{\mathfrak{S}}$ on $Z$. Moreover,
by the block geometry, $d_Z\ge d_{\mathfrak{S}}$ on $Z$.

The \emph{full-level region} of $Z$ and the \emph{deficit} of $Z$ are
\begin{equation}\label{9.5:eq-admissible-sets-3}
Z^{=}:=\{x\in Z:\ \lambda_Z(x)=\lambda_{\mathfrak{S}}(x)\},\qquad
\operatorname{def}_Z:=\lambda_{\mathfrak{S}}-\lambda_Z .
\end{equation}
By the preceding paragraph, $\operatorname{def}_Z\ge0$ on $Z$. For a pair $(\mathfrak{S}',\mathfrak{S})$ we
write $(\boldsymbol{\Omega}'(Z),\boldsymbol{\Omega}(Z))$ for the pair of shells of $Z$ over
$\mathfrak{S}'$ and over $\mathfrak{S}$, respectively.

\emph{Fat cores.} A \emph{fat core} is a union $K$ of proper-scale cubes that is designated, for a
given unit, never to be decoupled. In this section the letter $K$ denotes a fat core and not the
number of renormalisation steps. We write $\sigma_0(K)$ for the set of
proper-scale cubes of the site that are not among the cubes making up $K$. The \emph{boxes of the
unit} are the sets
\begin{equation}\label{9.5:eq-admissible-sets-4}
Z=K\cup Y',\qquad Y'\subset\sigma_0(K)\ \text{finite},
\end{equation}
where $Y'$ is identified with the union of its cubes. We write $K^{\circ}$ for the \emph{interior}
of $K$ in the sense of the block geometry, a subset of $K$ determined by it. Of $K^{\circ}$, only
three statements of the block geometry are used in this chapter; they are
stated later in the chapter. They assert that certain sets lie in $K^{\circ}$, and that
$K^{\circ}\subset Z^{=}$ for every box $Z\supseteq K$.
\end{definition}

\begin{definition}[Local norms, exponential smearing and slow functions]\label{9.5:def-slow-functions}
Fix an admissible set $\mathfrak{S}$ in the sense of Definition~\ref{9.5:def-admissible-sets}, with
its level function $\lambda=\lambda_{\mathfrak{S}}$, its local scale $\ell=\ell_{\mathfrak{S}}$, its
local cubes $\tilde{\Delta}(y)=\tilde{\Delta}_{\mathfrak{S}}(y)$ and its scaled path metric
$d=d_{\mathfrak{S}}$. Let $\delta_1>0$ be the fixed decay constant.

\emph{Local norms.} For a bond field $f$ and a point $y$ put
\begin{equation}\label{9.5:eq-slow-functions-1}
  |f|^{\circ}_y:=\sup_{\tilde{\Delta}(y)}|f|,\qquad |f|_y:=\ell(y)\,|f|^{\circ}_y .
\end{equation}

\emph{Exponential smearing.} For a non-negative function $f$ on bonds or on sites and for $a>0$
put
\begin{equation}\label{9.5:eq-slow-functions-2}
  f^{\sim,a}(x):=\sum_{y}e^{-a\delta_1 d(x,y)}\,f(y),\qquad f^{\sim}:=f^{\sim,1}.
\end{equation}
For a datum $b$ on sites, with pointwise modulus $|b|$, we write $\mathsf{P}_b:=|b|^{\sim}$.

\emph{Slow functions.} A positive function $\Psi$ is \emph{slow of type} $(C,a)$ if
\begin{equation}\label{9.5:eq-slow-functions-3}
  \Psi(y)\le C\,e^{a\delta_1 d(x,y)}\,\Psi(x)\qquad\text{for all }x,y .
\end{equation}
A function slow of type $(L,1)$ is called \emph{slow}.

Two elementary facts are used throughout:
\begin{enumerate}
\item[(i)] for every $f\ge0$ and every $a>0$, the function $f^{\sim,a}$ satisfies
\eqref{9.5:eq-slow-functions-3} with $(C,a)$ replaced by $(1,a)$, that is, it is slow of
type $(1,a)$;
\item[(ii)] if $\Psi_1,\dots,\Psi_n$ are slow of types $(C_1,a_1),\dots,(C_n,a_n)$, then
$\prod_i\Psi_i$ is slow of type $(\prod_iC_i,\sum_ia_i)$, and $\max_i\Psi_i$ is slow of type
$(\max_iC_i,\max_ia_i)$.
\end{enumerate}
\end{definition}

\begin{proof}
(i) Let $x,y$ be points and $z$ a summation point. The triangle inequality for $d$ gives
$d(x,z)\le d(x,y)+d(y,z)$, that is, $-d(y,z)\le d(x,y)-d(x,z)$. Since $a\delta_1>0$ and
$f(z)\ge0$, multiplying by $a\delta_1$, exponentiating and summing over $z$ yields
\begin{equation*}
  f^{\sim,a}(y)=\sum_{z}e^{-a\delta_1 d(y,z)}f(z)
  \le e^{a\delta_1 d(x,y)}\sum_{z}e^{-a\delta_1 d(x,z)}f(z)
  =e^{a\delta_1 d(x,y)}f^{\sim,a}(x),
\end{equation*}
which is \eqref{9.5:eq-slow-functions-3} with $C=1$. Every weight $e^{-a\delta_1 d(x,z)}$ is
positive, so $f^{\sim,a}$ is positive as soon as $f$ is not identically zero.

(ii) Fix $x,y$. By \eqref{9.5:eq-slow-functions-3} for each $\Psi_i$, and since all $\Psi_i$
are positive, multiplying the $n$ inequalities gives
\begin{equation*}
  \prod_i\Psi_i(y)\le\Bigl(\prod_iC_i\Bigr)e^{(\sum_ia_i)\delta_1 d(x,y)}\prod_i\Psi_i(x).
\end{equation*}
For the maximum, put $C:=\max_iC_i$ and $a:=\max_ia_i$. Since $\delta_1 d(x,y)\ge0$, for
each $i$
\begin{equation*}
  \Psi_i(y)\le C_ie^{a_i\delta_1 d(x,y)}\Psi_i(x)\le C\,e^{a\delta_1 d(x,y)}\max_j\Psi_j(x),
\end{equation*}
and taking the maximum over $i$ on the left gives
$\max_i\Psi_i(y)\le C\,e^{a\delta_1 d(x,y)}\max_i\Psi_i(x)$. Both functions are positive, so
they are slow of the stated types.
\end{proof}

\begin{definition}[Regularity classes, the slot rule and the box response]\label{9.5:def-box-response}
Let $\mathfrak{S}$ be an admissible set in the sense of
Definition~\ref{9.5:def-admissible-sets}, with level function $\lambda_{\mathfrak{S}}$, local scale
$\ell=\ell_{\mathfrak{S}}$ and local cubes $\tilde{\Delta}(y)=\tilde{\Delta}_{\mathfrak{S}}(y)$. Let
$\eta$ be the spacing of the finest lattice, so that $\ell(y)=L^{\lambda(y)}\eta$. Fix
$0<\beta<1$. For a field $f$ on $\tilde{\Delta}(y)$ we write $[f]_\beta$ for its H\"older
seminorm of exponent $\beta$ over $\tilde{\Delta}(y)$.

\emph{Regularity class of points.} For $\mathfrak{r}>0$, the class
$\mathsf{C}_\beta(\mathfrak{r})$ consists of the real points $P=(\mathbf{U},\mathbf{J})$, made of a
configuration $\mathbf{U}$ and a slot content $\mathbf{J}$, with the following property: for every
$y$ there is a gauge on $\tilde{\Delta}(y)$ in which $\mathbf{U}=e^{i\eta A}$ with
\begin{equation}\label{9.5:eq-box-response-1}
\ell(y)\,|A|,\quad \ell(y)^2\,|\nabla A|,\quad \ell(y)^{2+\beta}\,[\nabla A]_\beta,
\quad \ell(y)^3\,|\mathbf{J}|\ \le\ \mathfrak{r}
\end{equation}
on $\tilde{\Delta}(y)$.

\emph{Class of increments.} For a non-negative function $\mathsf{r}$ on sites (a
\emph{profile}), the class $\operatorname{Inc}_\beta(\mathsf{r})$ consists of the real fields
$\mathfrak{h}$ such that, for every $y$,
\begin{equation}\label{9.5:eq-box-response-2}
\ell(y)\,|\mathfrak{h}|,\quad \ell(y)^2\,|\nabla\mathfrak{h}|,\quad
\ell(y)^{2+\beta}\,[\nabla\mathfrak{h}]_\beta,\quad \ell(y)^3\,|\Delta_{\mathbf{U}}\mathfrak{h}|
\ \le\ \mathsf{r}(y)
\end{equation}
on $\tilde{\Delta}(y)$, where $\Delta_{\mathbf{U}}$ is the covariant lattice Laplacian at
$\mathbf{U}$.

\emph{The slot rule.} With $D$ the covariant lattice derivative at $\mathbf{U}$, $D^*$ its adjoint,
$R_{\mathfrak{S}}$ Ba{\l}aban's projection for the set $\mathfrak{S}$ at $\mathbf{U}$, and
$F(\mathbf{U},\mathfrak{h})$ the remainder term of \cite[(3.13)]{B-CMP109}, we put
\begin{equation}\label{9.5:eq-box-response-3}
\mathcal{J}_{\mathfrak{S}}(\mathfrak{h};\mathbf{U})
:=\bigl(D^*D+D\,R_{\mathfrak{S}}\,D^*\bigr)\mathfrak{h}+F(\mathbf{U},\mathfrak{h}),
\end{equation}
which is the rule of \cite[(3.13)]{B-CMP109}.

\emph{The response.} Let $P=(\mathbf{U},\mathbf{J})$ be a point and $b$ a datum on sites. Put
$\mathsf{m}_P:=H_{\mathbf{U}}^{\dagger}\mathbf{J}$, where $H_{\mathbf{U}}^{\dagger}$ is the adjoint
operator of the external response lemma. Under the convention fixed with that lemma, the
\emph{response} $\mathbb{H}_{\mathfrak{S}}(b;P)$ is the small critical point
\begin{equation}\label{9.5:eq-box-response-4}
\mathbb{H}_{\mathfrak{S}}(b;P):=\mathfrak{X}(b;\mathbf{U},\mathsf{m}_P)
\end{equation}
of the functional $\Sigma_{\mathbf{U},\mathsf{m}_P}$ of the external response lemma on the constraint
set $\{X:\ \mathbf{Q}_{\mathbf{U}}(\eta X)=b\}$, with $\mathbf{Q}_{\mathbf{U}}$ the averaging
constraint at $\mathbf{U}$ and all operators taken for the admissible set $\mathfrak{S}$.

\emph{The vertex of a box.} For a box $Z$ in the sense of
Definition~\ref{9.5:def-admissible-sets}, with shells forming the admissible set
$\boldsymbol{\Omega}(Z)$, the \emph{vertex} is
\begin{equation}\label{9.5:eq-box-response-5}
\mathbb{H}_Z(b;P):=\mathbb{H}_{\boldsymbol{\Omega}(Z)}\bigl(b;P\restriction Z\bigr)
\quad\text{on } Z,\qquad \mathbb{H}_Z(b;P):=0\quad\text{off } Z,
\end{equation}
where the operators $D$, $R_{\boldsymbol{\Omega}(Z)}$, $\mathbf{Q}_{\mathbf{U}}$ and
$H_{\mathbf{U}}^{\dagger}$ entering \eqref{9.5:eq-box-response-4} are those of the admissible set
$\boldsymbol{\Omega}(Z)$, evaluated at the restricted point $P\restriction Z$.
The vertex $\mathbb{H}_Z(b;P)$ vanishes on the outermost stratum of $\boldsymbol{\Omega}(Z)$ and at
$b=0$; these are the two properties of the vertex stated with the external response lemma.
\end{definition}

\begin{definition}[Presented pairs and their total variation]\label{9.5:def-total-variation}
We use the classes $\operatorname{Inc}_\beta(\mathsf{r})$, the slot rule $\mathcal{J}_{\mathfrak{S}}$ and the
constant $\eta$ of Definition~\ref{9.5:def-box-response}, and the mismatch field $A_p$ of a pair of
points over exactly paired sets, \eqref{9.5:eq-mismatch-field-a}--\eqref{9.5:eq-mismatch-field-b}.

\emph{Presented points.} Let $P_0=(\mathbf{U}_0,\mathbf{J}_0)$ be a point and let $r\ge 0$ be an
integer. A \emph{presented point over} $P_0$ is a point $P=(\mathbf{U}_r,\mathbf{J}_r)$ obtained by
$r$ links,
\begin{equation}\label{9.5:eq-total-variation-1}
\mathbf{U}_s := e^{i\eta\mathfrak{h}_s}\,\mathbf{U}_{s-1},\qquad
\mathbf{J}_s := \mathbf{J}_{s-1}+\mathcal{J}_{\mathfrak{S}_s}(\mathfrak{h}_s;\mathbf{U}_{s-1}),
\qquad 1\le s\le r,
\end{equation}
where $\mathfrak{S}_s$ are the admissible sets of the $s$-th link and
$\mathfrak{h}_s\in\operatorname{Inc}_\beta(\mathsf{r}_s)$ for a non-negative profile $\mathsf{r}_s$,
the class being taken in the units of $\mathfrak{S}_s$, that is, with the local scale
$\ell_{\mathfrak{S}_s}$ and the cubes $\tilde{\Delta}_{\mathfrak{S}_s}(y)$. No further relation
among the $\mathfrak{h}_s$ is required; in particular they may be dilated, cut or glued. The
data $(\mathfrak{S}_s,\mathfrak{h}_s)_{1\le s\le r}$ form the \emph{presentation}, $r$ is its
\emph{length}, and we put
\[
\mathfrak{r}_P := \mathfrak{r}_e + C\sum_{1\le s\le r}\sup\mathsf{r}_s,
\]
with $\mathfrak{r}_e$ the class radius of the seed or entry over which the presentation is built and
$C$ the constant that accompanies the box response of Definition~\ref{9.5:def-box-response}.
A restriction $P\restriction Z$ shrinks the region on which $P$ is read and leaves
the presentation unchanged. Writing $\lambda$ for the level function of the current sets, the
\emph{level condition} on a presentation is
\begin{equation}\label{9.5:eq-total-variation-a}
\lambda_{\mathfrak{S}_s}\ \ge\ \lambda\quad\text{on the region read, for every } s\le r .
\end{equation}

\emph{Presented pairs.} A \emph{presented pair} $p$ consists of two presentations of equal length
$r$, one in each of the machines $\mathsf{A}$ and $\mathsf{B}$, over exactly paired sets
$\mathfrak{S}_s$, $\mathfrak{S}'_s$, with equal cut-off parameters in both machines, the seed or
the entry of the pair being written in the first frame; the objects of the second presentation
are written with primes. For $0\le s\le r$ the \emph{truncation} $p_s$ is the pair formed by the
first $s$ links of both presentations; thus $p_0$ is the seed or entry and $p_r=p$. The
\emph{link errors} are
\[
e_s := Y[\mathfrak{h}'_s;\mathbf{U}'_{s-1}] - \mathfrak{h}_s,\qquad 1\le s\le r ,
\]
where $Y[\mathfrak{h}'_s;\mathbf{U}'_{s-1}]$ is the field, in the units of the first presentation,
that the composition identity associates with the link $(\mathfrak{h}'_s,\mathbf{U}'_{s-1})$ of the
second presentation; the link errors enter only through the bound stated at the end of this
definition.

\emph{Total variation.} The \emph{total variation} of $p$, measured in the local norms
$|\cdot|_x$ of the current sets, is
\begin{equation}\label{9.5:eq-total-variation-b}
\mathsf{T}_p(x) := \mathsf{E}(x) + \sum_{1\le t\le r}\bigl|A_{p_t}-A_{p_{t-1}}\bigr|_x
\ \ge\ |A_{p_s}|_x\qquad(0\le s\le r),
\end{equation}
where $\mathsf{E}(x):=|A_{p_0}|_x$ when $p_0$ is a seed, and, when $p_0$ is an entry, $\mathsf{E}$
is the modulus of the entry, defined below together with the entries of pairs; seeds and entries,
the two forms that the initial pair $p_0$ may take, are introduced there. Each increment term lives
on the region of its link: the sum at $x$ runs over the links whose region contains $x$.

\emph{Derived majorants.} We put $\mathsf{S}_p := \mathsf{T}_p/\ell$, with $\ell$ the local scale of
the current sets, that is, the same sum with the
norms $|\cdot|^{\circ}$ in place of $|\cdot|$, and
\begin{equation}\label{9.5:eq-total-variation-2}
\Psi_p := \mathsf{T}_p^{\sim} + \mathfrak{r}_0\,L^{-\theta\lambda},\qquad
\Psi^B_p := \mathsf{T}_p^{\sim} + L^{-\theta\lambda},
\end{equation}
where $\mathsf{T}_p^{\sim}$ is the exponential smearing of $\mathsf{T}_p$ and $\theta$ is a fixed
exponent.

By the composition identity, the link errors of a presented pair satisfy
\[
\sum_{s\le r}|e_s|_x\ \le\ 3\,\mathsf{T}_p(x).
\]
\end{definition}

\begin{definition}[Entries and their modulus]\label{9.5:def-entries}
Let $\mathsf{E}$ be a non-negative function on sites and $C_e$ a constant. An \emph{entry of modulus}
$(\mathsf{E},C_e)$ is a real pair $p_0=(P_0,P_0')$ on exactly paired sets, with
$P_0=(\mathbf{U}_0,\mathbf{J}_0)$ and $P_0'=(\mathbf{U}_0',\mathbf{J}_0')$ (here $p_0$ denotes a pair,
not the degree function $p_0(\cdot)$ of Notation~\ref{0.2:not-glossary}), such that the following
three conditions hold.
\begin{enumerate}
\item[(e1)] $P_0,P_0'\in\mathsf{C}_\beta(\mathfrak{r}_e)$.
\item[(e2)] The mismatch field $A_{p_0}$ of the pair satisfies $|A_{p_0}|_x\le\mathsf{E}(x)$ at every
site $x$, and $\sup\mathsf{E}\le\varepsilon_D$.
\item[(e3)] On every regular fine test field $w$ (a test field, not a complex source), the two slots
$\mathbf{J}_0$ and $\mathbf{J}_0'$ pair alike up to an error of absolute value at most
\begin{equation*}
C_e\sum_y w^{\sim}(y)\,\bigl[\mathsf{E}^{\sim}+\mathfrak{r}_0L^{-\theta\lambda}\bigr](y),
\end{equation*}
where $w^{\sim}$ and $\mathsf{E}^{\sim}$ are the exponential smearings of $w$ and $\mathsf{E}$.
\end{enumerate}
The entry is \emph{diagonal} if $\mathbf{J}_0=J(\mathbf{U}_0)$ and $\mathbf{J}_0'=J(\mathbf{U}_0')$.

Let a presented pair over the entry $p_0$ be given, with links of sets $\mathfrak{S}_s$,
$1\le s\le r$ (Definition~\ref{9.5:def-total-variation}). It satisfies the level condition for
entries if, on the region read, the sets of its links and the set $\mathfrak{S}$ on which (e1) holds
all have level at least $\lambda$; that is, on the region read,
\begin{equation}\label{9.5:eq-entries-1}
\lambda_{\mathfrak{S}_s}\ \ge\ \lambda\quad(1\le s\le r)\qquad\text{and}\qquad
\lambda_{\mathfrak{S}}\ \ge\ \lambda,
\end{equation}
the first half of which is the condition \eqref{9.5:eq-total-variation-a}.
\end{definition}

\begin{definition}[Blocks, letter assignments and the letter profile]\label{9.5:def-letter-profile}
Fix a site and the admissible set $\mathfrak{S}$ that it carries, with its layers, its local scale
$\ell=\ell_{\mathfrak{S}}$, its local cubes $\tilde{\Delta}(y)=\tilde{\Delta}_{\mathfrak{S}}(y)$ and its
scaled path metric $d_{\mathfrak{S}}$, in the sense of Definitions~\ref{9.5:def-admissible-sets}
and~\ref{9.5:def-slow-functions}. We keep the subscript on $d_{\mathfrak{S}}$ to distinguish
this metric from the dimension $d$.

\emph{Blocks and the partition of unity.} The site carries Ba{\l}aban's cover $\{\square\}$ of
its layers by \emph{blocks} \cite{B-CMP109}. A block is a union of $2^d$ neighbouring $M$-cubes
of the scale of its layer. For a block $\square$ we write $\ell_\square$ for the scale of its
layer and $\tilde{\square}$ for its enlargement. The cover comes with a partition of unity
$\{\zeta_\square\}$ of class $C^2$ \cite{B-CMP109}:
\begin{equation}\label{9.5:eq-letter-profile-1}
  \operatorname{supp}\zeta_\square\subset\tilde{\square},\qquad
  \sum_{\square}\zeta_\square=1,\qquad
  \zeta_\square\ge 0 .
\end{equation}
At each point at most $2^d$ of the functions $\zeta_\square$ are non-zero. Moreover
\begin{equation}\label{9.5:eq-letter-profile-2}
  \ell_\square\,|\nabla\zeta_\square|\le\frac{C_\zeta}{M},\qquad
  \ell_\square^{2}\,|\nabla^{2}\zeta_\square|\le\frac{C_\zeta}{M^{2}} ,
\end{equation}
where $C_\zeta$ is a number. These two bounds express that $\zeta_\square$ varies at scale $M$;
they are used here as stated in \cite{B-CMP109}.

\emph{Letter assignments and cut-offs.} A \emph{letter assignment} $c$ assigns a real or complex
number $c_\square$ to each block $\square$. Its \emph{cut-off} is the function
\begin{equation*}
  \chi_c:=\sum_{\square}c_\square\,\zeta_\square .
\end{equation*}
For a point $y$ of the site we put
\begin{equation*}
  \bar{c}(y):=\max\bigl\{|c_\square| :
  \tilde{\Delta}(y)\cap\operatorname{supp}\zeta_\square\neq\emptyset\bigr\}.
\end{equation*}

\emph{The source region.} Let $\Omega^{\mathsf{s}}$ denote the site's source region. It is the
member $\Omega_{i+1}$ of the admissible set $\mathfrak{S}$ at a site of type (T), the region
$\bar{\Lambda}$ at a site of type (L), and the region $R_\nu$ at a site of type (P).

\emph{The letter profile.} Let $c_0^\chi$ and $\mathfrak{r}_0^\chi$ be numbers. Consider a letter
assignment whose letters on the strip blocks (the strip letters) obey
\begin{equation}\label{9.5:eq-letter-profile-3}
  |c_\square|\le\mathfrak{r}_0^\chi\,
  e^{\delta_1 d_{\mathfrak{S}}(\tilde{\square},\,\Omega^{\mathsf{s}})/8},
\end{equation}
and whose letters on all other blocks obey $|c_\square|\le c_0^\chi$. For such a letter
assignment the \emph{letter profile} $\bar{\tau}_\chi$ and the number $C_\Delta$ entering it are
\begin{equation}\label{9.5:eq-letter-profile-a}
  \begin{aligned}
  \bar{\tau}_\chi(y)&:=\max\Bigl\{c_0^\chi,\;
  \mathfrak{r}_0^\chi\,C_\Delta\,
  e^{\delta_1 d_{\mathfrak{S}}(y,\,\Omega^{\mathsf{s}})/8}\Bigr\},\\
  C_\Delta&:=\sup_{y}\exp\Bigl(\frac{\delta_1}{8}\,
  \operatorname{diam}_{d_{\mathfrak{S}}}\Bigl(\tilde{\Delta}(y)\cup
  \bigcup_{\tilde{\square}\ni y}\tilde{\square}\Bigr)\Bigr).
  \end{aligned}
\end{equation}
The number $C_\Delta$ depends only on $d$, $L$, $M_1$ and $M$. Finally we put
\begin{equation}\label{9.5:eq-letter-profile-5}
  C_\chi:=2^{d}\,C_\zeta\,L .
\end{equation}
\end{definition}

\begin{definition}[Terms over strips and lettered cut-offs]\label{9.5:def-lettered-cutoffs}
This definition uses the following objects of Definition~\ref{9.5:def-letter-profile}: the blocks
$\square$ of the cover of a site, their enlargements $\tilde{\square}$ and $\tilde{\square}^{4}$, the
partition of unity $\{\zeta_\square\}$, and the constant $C_\chi$. The sites and their types T, L
and P, the stages, entry chains and moves of a site of type P, the bracket and the seed of a site of
type T, the root $q_0$ and the rows of a site of type L, and the live large-field regions $Y$ with
their strips $\mathbb{I}_i(Y)$ are taken as they are defined with the sites of this chapter; $N$ is
the number of stages of Definition~\ref{9.1:not-levelwise-rate}.

\emph{Terms over strips.} A \emph{term over strips} read at a site $\mathsf{s}$ of current level $i$ is a
triple $c=(F;X,X^\theta)$ with first component $F$, taken as fixed with the sites of this chapter,
true domain $X$ and tree domain $X^\theta$. Its tree core is
$K^\theta:=K(X^\theta)$, the fat core of $X^\theta$ in the sense of
Definition~\ref{9.1:not-radius-exponent-floors}. The term carries its live large-field regions $Y$ at
level $i$, each with its strip $\mathbb{I}_i(Y)$.

\emph{Strip blocks and plain blocks.} Let $X^{+}$ be the set associated with the true domain $X$ in the
data of the site $\mathsf{s}$, taken as it is fixed with the sites of this chapter. The
\emph{strip blocks} of $c$ are the
blocks in
\begin{equation*}
  \mathcal{D}_c:=\bigl\{\square:\ \tilde{\square}\cap(X^{+}\setminus K^\theta)\neq\emptyset\bigr\}.
\end{equation*}
Every other block meeting $X^{+}$ is called \emph{plain}. For a plain block we have
$\tilde{\square}\cap X^{+}\subset K^\theta$.

\emph{Rows of a site.} $\mathsf{Rows}(\mathsf{s})$ denotes the finite ordered family of increments
of the site $\mathsf{s}$. It is given as follows.
\begin{itemize}
\item At a site of type T it consists of the single increment $\mathcal{Y}$.
\item At a site of type L it consists of the rows $3,2,1,0'$ over the root $q_0$, in this order.
\item At a site of type P at stage $\nu$ it consists of the links of the entry chain, followed by the
  moves of the stages in $[\nu,N[$.
\end{itemize}
We put $\mathcal{D}^{+}_c:=\mathcal{D}_c\times\mathsf{Rows}(\mathsf{s})$.

\emph{The core of a set of indexed strip blocks.} For $D\subset\mathcal{D}^{+}_c$ put
\begin{equation}\label{9.5:eq-lettered-cutoffs-1}
  \bar{S}_D:=[K^\theta]_i\ \cup\bigcup_{(\square,r)\in D}\bigl(\tilde{\square}^{4}\cup
  \operatorname{Cor}(\square)\bigr),
  \qquad K_D:=(\bar{S}_D)^{\sim}.
\end{equation}
Here $\operatorname{Cor}(\square)$ is the corridor of the block $\square$; $[K^\theta]_i$ is the set
determined by the tree core $K^\theta$ at the current level $i$; and $(\,\cdot\,)^{\sim}$ is the
operation on sets that turns the core seed set $\bar{S}_D$ into $K_D$. Both operations are taken as
they are fixed with the sites of this chapter. The set $K_D$ is the never-decoupled core of the unit
$v_D$ indexed by $D$. The boxes of $v_D$
are the sets $Z=K_D\cup Y'$ with $Y'\subset\sigma_0(K_D)$ finite.

\emph{Lettered cut-offs.} Let $r\in\mathsf{Rows}(\mathsf{s})$, and let
$t=(t_{\square,r})_{(\square,r)\in\mathcal{D}^{+}_c}$ be a family of \emph{strip letters}, letters in
the sense of Definition~\ref{9.5:def-letter-profile}. The
\emph{lettered cut-off} of the row $r$ is
\begin{equation}\label{9.5:eq-lettered-cutoffs-a}
  \chi^{(r)}_t:=\sum_{\square\ \mathrm{plain}}\tau^{(r)}_\square\,\zeta_\square
  +\sum_{\square\in\mathcal{D}_c}t_{\square,r}\,\zeta_\square .
\end{equation}
Here $\tau^{(r)}=(\tau^{(r)}_\square)$ are the letters that the increment $r$ carries on plain blocks.
They are given as follows.
\begin{itemize}
\item At a site of type T, $\tau^{(r)}_\square$ equals $1$ on the plain blocks before the block of the
  bracket, $\mathsf{t}$ at that block, and $0$ on the plain blocks after it.
\item At a site of type L, they are the letters fixed for the row $r$ on plain blocks in the
  definition of the rows of a site of type L.
\item At a site of type P, they are the block letters $\mathbf{1}_{A'}$ of the move $r$, that is,
  the indicator of the set $A'$ attached to that move.
\end{itemize}

\emph{The lettered configuration.} The \emph{lettered configuration}, at which the term is read, is
\begin{equation}\label{9.5:eq-lettered-cutoffs-2}
  \mathcal{C}_t:=\prod^{\mathrm{ord}}_{r\in\mathsf{Rows}(\mathsf{s})}
  e^{\,\mathrm{i}\eta\,\chi^{(r)}_t\mathfrak{h}^{(r)}}\cdot U^0 .
\end{equation}
The product is ordered along $\mathsf{Rows}(\mathsf{s})$, $\mathrm{i}=\sqrt{-1}$,
$\mathfrak{h}^{(r)}$ is the increment
field of the row $r$, and $\eta$ is taken as it is defined with the sites of this chapter. The
configuration $U^0$ depends on the type of the site: it is the seed at a site
of type T, the root $q_0$ at a site of type L, and the entry at a site of type P.

\emph{Radii.} The two-argument $d(\cdot,\cdot)$ below is the distance between sets in the metric fixed
with the sites of this chapter, and not the dimension $d=4$; $\Omega_{i+1}$ and $\bar{\Lambda}$ are
taken as they are defined with the sites. At a site of type T the letter radii are
\begin{equation*}
  r_\square:=r^0_T\,e^{\delta_1 d(\tilde{\square},\Omega_{i+1})/8}.
\end{equation*}
At a site of type L they are
\begin{equation*}
  r^L_\square:=r^0_L\,e^{\delta_1 d(\tilde{\square},\bar{\Lambda})/8}.
\end{equation*}
At a site of type P they are the pair radii $r^{\bullet}_\square$ fixed with the sites of type P. We
write $r'_\square:=r_\square/C_\chi$.
\end{definition}

\begin{remark}[The slot of a lettered configuration]
We adopt the following convention for the point read at a lettered configuration. The point read is
$(\mathcal{C}_t,\mathbf{J}_t)$. Here the slot $\mathbf{J}$ is moved at each lettered increment by the
slot rule $\mathcal{J}$ of this chapter, with the same cut field $\chi^{(r)}_t\mathfrak{h}^{(r)}$ as
the field component:
\begin{equation*}
  \mathbf{J}\longmapsto\mathbf{J}+\mathcal{J}\bigl(\chi^{(r)}_t\mathfrak{h}^{(r)};\,\cdot\,\bigr),
\end{equation*}
that is, by the rule of \cite[(3.13)]{B-CMP109}. At a site of type P this slot move is the move
$\Phi_\mu$, taken as defined with the sites of
type P. Display
\eqref{9.5:eq-lettered-cutoffs-2} concerns the field component only. The slot component at which the
term is read at $\mathcal{C}_t$ is not identified here with the slot of this convention; that
identification concerns one of the machines $\mathsf{A}$, $\mathsf{B}$ only.
\end{remark}

\begin{definition}[The three families of lettered links]\label{9.5:def-lettered-links}
Throughout, a pair is written $p=(P',P)$, the primed point first, and for increment fields
$\mathfrak{h}',\mathfrak{h}$ we write ``$p$ extended by $(\mathfrak{h}',\mathfrak{h})$'' for the
presented pair obtained from $p$ by one further link with these increment fields, in the sense of
Definition~\ref{9.5:def-total-variation}. The responses $\mathbb{H}_Z$, $\mathbb{H}'_Z$ are those
of Definition~\ref{9.5:def-box-response}. The strip blocks $\mathcal{D}_c$, the indexed strip
blocks $\mathcal{D}^{+}_c$ and the lettered cut-offs $\chi^{(r)}_t$ are those of
Definition~\ref{9.5:def-lettered-cutoffs}. Entries and diagonal entries are those of
Definition~\ref{9.5:def-entries}. For a fat core $K$ and a region $Y'$ we write $Z=K\cup Y'$ for
the corresponding box.

\emph{The family \textup{(P)}.} A \emph{move of stage $\nu$} is a quadruple $\mu=(c,A,K,Z)$
consisting of an input $c$ of stage $\nu+1$, a finite set $A$ of blocks, a fat core
\begin{equation*}
K\supseteq\bigcup_{\square\in A}\tilde{\square}^{4}
\end{equation*}
which contains the core letter $\mathsf{d}^\theta(e)$ ($K$ is not necessarily $K_c=X^{+}$), and a
box $Z=K\cup Y'$. Let $p=(P',P)$ be a parent on a region
$Y\supseteq Z$. The real parameters of the move are $\pi=(b,t)$, where $b$ is a real source datum
in the range of source data admitted in Definition~\ref{9.5:def-box-response} and
$t=(t_\square)_\square$ satisfies
\begin{equation*}
|t_\square|\le 4r^{\bullet}_\square\quad\text{for }\square\in A,\qquad
t_\square=0\quad\text{for }\square\notin A .
\end{equation*}
We put
\begin{equation}\label{9.5:eq-lettered-links-1}
\tau_t:=\sum_{\square\in A}t_\square\,\zeta_\square ,
\end{equation}
and we write $p\restriction\mathrm{strip}$ for the parent restricted to the part, lying off $K$,
of the hull of the term at which the produced pair is read; that term reads the parent on that
part as a spectator. The move $\mu$ applied to $p$ at $\pi$ produces the pair
\begin{equation}\label{9.5:eq-lettered-links-2}
\begin{aligned}
\Phi^{\times}_\mu(\pi;p):={}&\Bigl(p\text{ extended by }
\bigl(\tau_t\,\mathbb{H}'_Z\bigl((\operatorname{Ad}_g b)\oplus 0;P'\bigr),\
\tau_t\,\mathbb{H}_Z(b;P)\bigr)\Bigr)\restriction K^{\circ}\\
&\cup\bigl(p\restriction\mathrm{strip}\bigr).
\end{aligned}
\end{equation}
The system $\mathfrak{F}^{\mathbb{R},\mathbb{Y}}$ is the system of sets of real pairs generated by
the rules \textup{(F0)} (entries), \textup{(F1)} (restriction) and \textup{(F2)} (moves) of the
definition of the family $\mathfrak{F}$, in which a move of stage $\nu$ is the quadruple $\mu$
described above. Every point of $\mathfrak{F}^{\mathbb{R},\mathbb{Y}}$ has, by these rules, a
\emph{run presentation}
\begin{equation*}
p_0\to p_1\to\dots\to p_r,\qquad
p_s=\Phi^{\times}_{\mu_s}(\pi_s;p_{s-1})\restriction S_s\quad(1\le s\le r),
\end{equation*}
in which $p_0$ is an entry, $\mu_s$ is a move of stage $\nu_s$, $\pi_s$ are admissible real
parameters for $\mu_s$, $S_s$ is a region on which the restriction \textup{(F1)} is taken, $n$ is
the number of stages of Notation~\ref{9.1:not-levelwise-rate}, and the stages satisfy
$n>\nu_1>\dots>\nu_r$.

\emph{The family \textup{(T)}.} The parent is the seed pair of the step on the site,
\begin{equation*}
p_0=\bigl((U_0',J(U_0')),\,(U_0,J(U_0))\bigr),
\end{equation*}
formed by the critical minimisers on the hull at equal values of the common variables, written in
the frame fixed for seed pairs among the entries of the family $\mathfrak{F}$. The family carries
one lettered link. Its data are:
\begin{itemize}
\item a set $D\subset\mathcal{D}^{+}_c=\mathcal{D}_c\times\{\mathcal{Y}\}$;
\item a box $Z=K_D\cup Y'$;
\item a real bracket letter $\mathsf{t}$ within the range fixed for it;
\item real strip letters $t_\square$ with $|t_\square|\le r'_\square$.
\end{itemize}
With $\chi_t:=\chi^{(\mathcal{Y})}_t$ and $p_0=(P_0',P_0)$, and with $B'$ the source datum of the
step, we let $p_1$ be the pair which on $Z$ is
\begin{equation}\label{9.5:eq-lettered-links-3}
p_0\restriction Z\ \text{extended by}\
\bigl(\chi_t\,\mathbb{H}'_Z\bigl((\operatorname{Ad}_g B')\oplus 0;P_0'\bigr),\
\chi_t\,\mathbb{H}_Z(B';P_0)\bigr),
\end{equation}
and which coincides with $p_0$ off $Z$. The gluing is seamless, because the increment vanishes on
the outermost stratum of $\boldsymbol{\Omega}(Z)$.

\emph{The family \textup{(L)}.} The root $q_0$ is either a seed pair or a diagonal entry on the
region $Y_{\mathrm{out}}^{+}$, carrying its own admissible set $\mathfrak{S}_e$. The family
carries rows $r=1,\dots,m$. Here, in this definition, $m$ denotes the number of links of type
\textup{($\ell$1)} that follow
the last diagonal link of the chain, the types being those of the corollary that enumerates the
links of the chain, and $m\le m_L$. All rows are taken at one
box $Z=K_D\cup Y'$ of the unit $v_D$; each unit carries one family of decoupling parameters
$\mathsf{s}$. The partial products start from
\begin{equation*}
p^{(0)}:=q_0\restriction Z ,
\end{equation*}
and we write $p^{(r)}=(P^{(r)\prime},P^{(r)})$. For each $r=1,\dots,m$, exactly one of the
following two cases holds.
\begin{itemize}
\item[(a)] \emph{Response type.} The $r$-th row is the link
\begin{equation}\label{9.5:eq-lettered-links-4}
\begin{aligned}
p^{(r)}:=p^{(r-1)}\ \text{extended by}\ \Bigl(
&\chi^{(r)}_t\,\mathbb{H}^{(r)\prime}_Z\bigl((\operatorname{Ad}_g b_r)\oplus 0;P^{(r-1)\prime}\bigr),\\
&\chi^{(r)}_t\,\mathbb{H}^{(r)}_Z\bigl(b_r;P^{(r-1)}\bigr)\Bigr).
\end{aligned}
\end{equation}
Here $\mathbb{H}^{(r)}_Z$ is the response of the current point on the shells
$\boldsymbol{\Omega}_{\mathfrak{S}_r}(Z)$ of $Z$ over the $r$-th base set $\mathfrak{S}_r$ of the
chain; by that corollary, base sets only refine along the
chain. The source datum $b_r$ is a
real source datum in the range of source data admitted in Definition~\ref{9.5:def-box-response}.
\item[(b)] \emph{Object type.} The $r$-th row is the link
\begin{equation}\label{9.5:eq-lettered-links-5}
p^{(r)}:=p^{(r-1)}\ \text{extended by}\
\bigl(\chi^{(r)}_t\,\mathcal{Z}^{(r)\prime},\ \chi^{(r)}_t\,\mathcal{Z}^{(r)}\bigr).
\end{equation}
Here $\mathcal{Z}^{(r)}$ is the displacement to a re-solved configuration of
\cite[(1.61)--(1.63)]{B-CMP122b}, that is, a link of type \textup{($\ell$2)} of that enumeration.
It is not of response form.
\end{itemize}
The enumeration of the links of the chain in that corollary, read against the displays
\cite[(1.51), (1.54), (1.57), (1.58)]{B-CMP122b},
determines which of those displays give links of response type and which give links of object
type. The present definition does not make that assignment and does not state it.

\emph{Vertices.} In each of the three families, consider a vertex, that is, the value of the
decoupling parameters with $\mathsf{s}=\mathbf{1}$ on $Y'$ and $\mathsf{s}=0$ on
$\sigma_0(K)\setminus Y'$. At a vertex, the increment produced by the family of decoupling
parameters of the unit is $\chi\,\mathbb{H}_Z$, with $\chi$ the cut-off of the link in question
($\tau_t$, $\chi_t$ or $\chi^{(r)}_t$); this follows from the vertex identity. This
identification is used on all of $Z$.
\end{definition}

\begin{definition}[Standing hypotheses on comparison and boxes]\label{9.5:not-hypotheses-comparison}
The four statements below are used in the rest of this chapter, each only in the form restated
here. Item~(d) is the external vertex identity; it is taken as stated and is not proved again.
\begin{itemize}
\item[(a)] \emph{Response comparison on paired admissible sets.} Let two exactly paired admissible
sets be given, the pairing including their labels, their boxes $\boldsymbol{\Omega}(Z)$ and their
joins. Let $p$ be a real presented pair on their region, in the sense of
Definition~\ref{9.5:def-total-variation}, satisfying the level condition of that definition.
Assume that its two presented points $P'$ and $P$ satisfy
\begin{equation*}
\mathfrak{r}_{P'}\le\mathfrak{r}_0,\qquad \mathfrak{r}_P\le\mathfrak{r}_0,
\end{equation*}
and that $\sup\mathsf{T}_p\le\varepsilon_D$.
Let $b$ and $\delta b$ be real, with $\sup\bigl(|b|+|\delta b|\bigr)\le a_{R'}$. For
conclusion~(C) below assume moreover the following. The function
$\tau$ is real and the same in both machines. It satisfies
\begin{gather*}
|\tau|\le\bar{\tau},\qquad \ell\,|\nabla\tau|\le\bar{\tau}/M,\qquad
\ell^{2}\,|\nabla^{2}\tau|\le\bar{\tau}/M^{2},
\end{gather*}
where $\bar{\tau}$ is slow. Finally,
\begin{equation*}
\mathsf{H}:=\bar{\tau}\,\mathsf{P}_b\le a_{R'} .
\end{equation*}
Then the conclusions (A)--(D) of the lettered response theorem, stated there with the letter
profile $\bar{\tau}_\chi$ and the constant $C_\chi$, hold with $C_\chi=1$ and with $\bar{\tau}$ in
place of $C_\chi\bar{\tau}_\chi$. The constants in these conclusions do not depend on the length
$r$ of the presentation. This comparison is used below through its form over entries, item~(b).
It is proved as an implication, under a convention fixed with the lettered response theorem,
from six inputs listed with that theorem. Among them are the level condition of
Definition~\ref{9.5:def-total-variation} and a second input stated together with it; we call the
latter the companion input of the level condition.

\item[(b)] \emph{Entries.} Throughout this item the level condition for entries of
Definition~\ref{9.5:def-entries} is in force on the region read.
\begin{itemize}
\item[--] A diagonal pair $p_0$ whose points satisfy $P_0,P_0'\in\mathsf{C}_\beta(\mathfrak{r}_e)$
is an entry of modulus $(|A_{p_0}|_{\cdot},C_i)$, that is, with $|A_{p_0}|_{\cdot}$ in the role of
$\mathsf{E}$ and a constant $C_i$ in the role of the constant $C_e$ of condition (e3) of
Definition~\ref{9.5:def-entries}. No criticality is required of $p_0$.
\item[--] Over an entry, the following statements hold verbatim: the first two clauses of the
statement on presented pairs over a diagonal base point, and its third clause in its form for
entries; the statement on pairing; the response comparison of item~(a) with its conclusions
(A)--(D); the corollary stated alongside these; and clause (e) of the box lemma quoted in
item~(c). In each of
them $\mathsf{T}_p$ is read as $\mathsf{E}$ plus the total variation of
Definition~\ref{9.5:def-total-variation}, and the constants do not depend on $r$.
\item[--] A presented pair $p$ over an entry, with $\mathsf{T}_p$ so read, is then itself an entry
of modulus $(\mathsf{T}_p,C_P)$, that is, with $\mathsf{T}_p$ in the role of $\mathsf{E}$ and a
constant $C_P$ in the role of $C_e$.
\end{itemize}
Moreover, consider real dilations and natural vertices. A chain end of the second kind of the
comparison chain of Definition~\ref{9.6:def-comparison-chain} is a
presentation of at most $m_L$ links over its last diagonal link, on sets that only refine.
Suppose every object of the chain satisfies the chain law with constant $C_{LR'}$. Then the
yardstick of the chain end in $\mathsf{T}$ is at most
\begin{equation*}
y_0^{\mathsf{T}}:=(2m_L+1)\max\{c_0,\,C_{LR'}\}.
\end{equation*}
The statements of this item are proved as an implication from a representation statement named
with them, from the structure of the comparison chain, Lemma~\ref{9.6:lem-chain-structure}, taken
as it is stated there, and from the chain law for the objects of the chain.

\item[(c)] \emph{Boxes} (clauses (d) and (e) of the box lemma). Let $Z$ be a union of cubes
common to the two exactly paired admissible sets. Then $\boldsymbol{\Omega}'(Z)$ is
exactly paired with $\boldsymbol{\Omega}(Z)$. The four structural conditions required of
admissible sets hold on $\boldsymbol{\Omega}(Z)$ with constants independent of $Z$. The response
comparison of item~(a) applies to the pair $(\boldsymbol{\Omega}'(Z),\boldsymbol{\Omega}(Z))$.
These statements are proved as an implication from two inputs: the companion input of the level
condition (item~(a)), and one further input listed with the box lemma, used as stated there.

\item[(d)] \emph{Vertex identity.} Consider a unit with never-decoupled core $K$ and with
decoupling parameters $\mathsf{s}$ on $\sigma_0(K)$. For a subset $Y'\subset\sigma_0(K)$, set
$\mathsf{s}=\mathbf{1}$ on $Y'$ and $\mathsf{s}=0$ on $\sigma_0(K)\setminus Y'$. Then the
$\mathsf{s}$-decoupled operators of the unit are those of the geometry
$\boldsymbol{\Omega}(K\cup Y')$. In the operator notation,
\begin{equation*}
\mathbb{O}\langle\mathbf{1}_{Y'}\rangle=\mathbb{O}_{K\cup Y'},
\end{equation*}
where $\mathbb{O}_{K\cup Y'}$ denotes the operators of the geometry
$\boldsymbol{\Omega}(K\cup Y')$. Consequently the $\mathsf{s}$-decoupled response equals
$\mathbb{H}_{K\cup Y'}$ on $K\cup Y'$ and vanishes off $K\cup Y'$.
\end{itemize}
\end{definition}

\begin{definition}[Standing hypotheses on metric and block geometry]\label{9.5:not-hypotheses-geometry}
In the remainder of this chapter the following hypotheses are in force. They concern an admissible set
$\mathfrak{S}$, with level function $\lambda_{\mathfrak{S}}$, local scale $\ell_{\mathfrak{S}}$,
local cubes $\tilde{\Delta}_{\mathfrak{S}}(y)$ and scaled path metric $d_{\mathfrak{S}}$ as in
Definition~\ref{9.5:def-admissible-sets}, and every refinement of $\mathfrak{S}$. In particular they
apply to the label, whose level function is $\iota$ and whose scale is $\ell_\iota$, to the shell set
$\boldsymbol{\Omega}(Z)$ of a box $Z$, and to an entry's own set $\mathfrak{S}_e$. When
$\mathfrak{S}$ is fixed we write $\lambda$ for $\lambda_{\mathfrak{S}}$ and $d$ for
$d_{\mathfrak{S}}$; $(t)_+:=\max\{t,0\}$; the smearing $f^{\sim,a}$ and the profile
$\mathsf{P}_b=|b|^{\sim}$ of a datum $b$ are those of Definition~\ref{9.5:def-slow-functions}.

\emph{Metric geometry.} With $\delta_1$ the decay constant entering the smearing of
Definition~\ref{9.5:def-slow-functions}, there are constants $c_1$, $C_T$ and, for $0<a\le 1$,
$C_w(a)$, such that the following hold.
\begin{itemize}
\item[(i)] For all $x,y$ the level comparison (i) of \eqref{9.5:eq-hypotheses-geometry-a} holds.
It holds in the same form for $\boldsymbol{\Omega}(Z)$ and for $\mathfrak{S}_e$, with their own
levels and metrics in place of $\lambda_{\mathfrak{S}}$ and $d_{\mathfrak{S}}$; more generally it
persists on refinements.
\item[(ii)] If $\mathfrak{S}^{\sharp}$ refines $\mathfrak{S}$ (so that its local units are finer),
then the metric comparison (ii) holds pointwise; in particular
$d_{\boldsymbol{\Omega}(Z)}\ge d_{\mathfrak{S}}$, and we write $d_Z:=d_{\boldsymbol{\Omega}(Z)}$.
\item[(iii)] For every datum $b$ on sites the profile bound (iii) holds at every $x$.
\item[(iv)] For the local unit function $\bar{u}^{\mathrm{loc}}$ of the site at which an estimate is
read, the summation bound (iv) holds at every $x$.
\item[(v)] The kernels of Ba{\l}aban's construction entering the response, among them the external
kernel $\mathcal{G}_0$, decay at rate at least $4\delta_1$ in the metric $d$. Writing $\mathsf{K}$
for the action of any one of them on non-negative functions, the smearing bound (v) holds for
$0<a\le1$; the same bound holds with $f^{\sim,a}$ replaced on both sides by
$L^{c\lambda}f^{\sim,a}$, $|c|\le 3$, with a constant again denoted $C_w(a)$.
\end{itemize}
\begin{equation}\label{9.5:eq-hypotheses-geometry-a}
\begin{aligned}
&\mathrm{(i)}\quad
L^{2(|\lambda_{\mathfrak{S}}(x)-\lambda_{\mathfrak{S}}(y)|-1)_+}
\le e^{\delta_1 d_{\mathfrak{S}}(x,y)/8},\\
&\mathrm{(ii)}\quad
d_{\mathfrak{S}}(x,y)\le d_{\mathfrak{S}^{\sharp}}(x,y),\\
&\mathrm{(iii)}\quad
\mathsf{P}_b(x)\le c_1\,\sup|b|\;e^{-\delta_1 d(x,\operatorname{supp} b)/2},\\
&\mathrm{(iv)}\quad
\sum_{y} e^{-\delta_1 d(x,y)/2}\,\bar{u}^{\mathrm{loc}}(y)\le C_T\,\bar{u}^{\mathrm{loc}}(x),\\
&\mathrm{(v)}\quad
\mathsf{K}\bigl[f^{\sim,a}\bigr]\le C_w(a)\,f^{\sim,a}.
\end{aligned}
\end{equation}
Items (i), (ii) and (iv), including the persistence of (i) on refinements, are taken by statement
from the metric estimates for admissible sets that accompany the random-walk expansion; items (iii)
and (v) are taken by statement as well.

\emph{The cover.} The cover $\{\square\}$ of a site by blocks and the partition of unity
$\{\zeta_\square\}$ have the properties displayed in Definition~\ref{9.5:def-letter-profile}: the
$\zeta_\square$ are of class $C^2$ and satisfy the derivative bounds at scale $M$ stated there, and
the cover has the multiplicity stated there. These are properties of Ba{\l}aban's construction, used
here as stated in \cite{B-CMP109}.

\emph{Block geometry.} For a set $S$ let $S^{\mathrm{fat}}$ denote the union of $S$, its
$\tilde{\Delta}$-fattening and its stencil-fattening. For a term over strips and a set $D$ of blocks
indexing the unit $v_D$, with core $K_D$, as in Definition~\ref{9.5:def-lettered-cutoffs}, let
$K^{\theta}_{\mathrm{pl}}$ be the union of the cells of the tree core $K^{\theta}$ that belong to
plain blocks. The following hold.
\begin{itemize}
\item[(vi)] For every fat core $K$ and every set $A$ of blocks with
$K\supseteq\bigcup_{\square\in A}\tilde{\square}^{4}$, the inclusion (vi) of
\eqref{9.5:eq-hypotheses-geometry-b} holds.
\item[(vii)] For every fat core $K$ and every box $Z\supseteq K$, the inclusion (vii) holds, where
$Z^{=}$ is the full-level region of $Z$.
\item[(viii)] The cubes at which the unit $v_D$ reads an increment are the cells of
$K^{\theta}_{\mathrm{pl}}$ and the enlargements $\tilde{\square}$, $\square\in D$. Together with
their $\tilde{\Delta}$- and stencil-fattenings they lie in the interior $K_D^{\circ}$; this is the
inclusion (viii).
\end{itemize}
\begin{equation}\label{9.5:eq-hypotheses-geometry-b}
\begin{aligned}
&\mathrm{(vi)}\quad
\Bigl(\bigcup_{\square\in A}\tilde{\square}\Bigr)^{\mathrm{fat}}\subset K^{\circ},\\
&\mathrm{(vii)}\quad
K^{\circ}\subset Z^{=},\\
&\mathrm{(viii)}\quad
\Bigl(K^{\theta}_{\mathrm{pl}}\cup\bigcup_{\square\in D}\tilde{\square}\Bigr)^{\mathrm{fat}}
\subset K_D^{\circ}.
\end{aligned}
\end{equation}
The inclusions (vi)--(viii) are taken by statement from the geometry of fat cores, of the boxes
containing them and of the units $v_D$.
\end{definition}

\begin{definition}[Standing hypotheses on summation, frames and entries]\label{9.5:not-hypotheses-summation}
The statements of this section are made under the standing hypotheses (a)--(g) below. Pairs, their
mismatch fields $A_p$ and their regularity classes are those of
Definitions~\ref{9.5:def-total-variation} and~\ref{9.5:def-entries}. Sites of types (T), (L), (P),
letters, cut-offs and clause $(\mathbb{Y})^{+}$ are those of
Definition~\ref{9.5:def-lettered-cutoffs}. Moves and the real lettered family
$\mathfrak{F}^{\mathbb{R},\mathbb{Y}}$ are those of Definition~\ref{9.5:def-lettered-links}. For a stage
$\nu$ of a run we write:
\begin{itemize}
\item $I_\nu$ for the set of lattice points on which the link of stage $\nu$ acts, the stages and their
links being those of Definition~\ref{9.5:def-lettered-links};
\item $\iota_\nu$ for the level function of the label at stage $\nu$, and $\iota^e_\nu$ for its
box-capped index;
\item $\bar{u}_\nu$, $u_\nu$ and $a_\nu$ for the ambient unit function, the majorant function and the
stage majorant at stage $\nu$.
\end{itemize}
We write $\mathcal{S}$ for a finite set of stages and $\min\mathcal{S}$ for its smallest element.

\begin{enumerate}
\item[(a)] \emph{Summation.} For every point $x$,
\begin{equation}\label{9.5:eq-hypotheses-summation-a}
\begin{gathered}
\text{(s1)}\quad \mathsf{A}(x):=\{\nu : x\in I_\nu\}\ \text{is an interval of stages};\\
\text{(s2)}\quad \iota_{\nu+1}(x)=\iota_\nu(x)+1,\quad
\bar{u}_{\nu+1}(x)\le \tfrac34\,\bar{u}_\nu(x)\qquad (\nu\in\mathsf{A}(x)),\\
\phantom{\text{(s2)}}\quad \iota_{\nu+1}(x)=\iota_\nu(x),\quad
\bar{u}_{\nu+1}(x)=\bar{u}_\nu(x)\qquad (\nu\notin\mathsf{A}(x));\\
\text{(s3)}\quad \sum_{\nu'\in\mathcal{S}} a_{\nu'}(x)\,\vartheta^{\iota^e_{\nu'}(x)}
\le C_\Sigma\, u_{\min\mathcal{S}}(x)\quad\text{for every finite set $\mathcal{S}$ of stages}.
\end{gathered}
\end{equation}
Here $C_\Sigma:=8+6\beta$, and $a_\nu(x)=2\bar{u}_\nu(x)$ for $\nu\in\mathsf{A}(x)$. Moreover, with $C$
a numerical constant and $\alpha_1(\gamma)$ a number depending only on $\gamma$,
\begin{equation*}
\bar{\bar{u}}:=\sup \bar{u}^{\mathrm{loc}}\le C\,\alpha_1(\gamma)\,L(L+1)\,\check{R}.
\end{equation*}

\item[(b)] \emph{Composition identity.} Let $p$ be any parent pair, and let $p^{+}$ be the pair built from
it with the increment field $\mathfrak{h}$ and a second field $Y$; in this item the letter $Y$ denotes
this field and not a region. Let $C_B$ be the constant fixed with the composition identity. Then
\begin{equation}\label{9.5:eq-hypotheses-summation-1}
\begin{gathered}
e^{i\eta A_{p^{+}}}=e^{i\eta Y}\,e^{i\eta A_p}\,e^{-i\eta\mathfrak{h}},\qquad
|A_{p^{+}}|\le k_B\bigl(|A_p|+|Y-\mathfrak{h}|\bigr),\\
k_B:=1+C_B\,\eta\,\bigl(|Y|+|\mathfrak{h}|\bigr).
\end{gathered}
\end{equation}

\item[(c)] \emph{Frame bound.} Let a lettered link be given over a parent of its family. With $g$ the
frame rotation of the link (not the coupling) and $b$ its source datum, put
$\delta b:=(\operatorname{Ad}_g-1)b$, and let $c_b^{\mathrm{fr}}$ denote the frame constant. Then
for every $x$
\begin{equation}\label{9.5:eq-hypotheses-summation-2}
2C_{R'}\,C_\chi\,\bar{\tau}_\chi\,\mathsf{P}_{\delta b}(x)
\le \tfrac12\,\ell_\iota(x)\,a(x)\,\vartheta^{\iota^e_Z(x)}\,c_b^{\mathrm{fr}}.
\end{equation}

\item[(d)] \emph{Entry laws.}
\begin{enumerate}
\item[(d1)] A seed pair $p_0$ obeys the following bound at every point $x$ of its region, where
$c_0:=C_A C_Z$, with $C_A$ and $C_Z$ the constants fixed with the seed pairs:
\begin{equation}\label{9.5:eq-hypotheses-summation-3}
|A_{p_0}(x)|\le c_0\,\bar{u}(x)\,\vartheta^{\iota(x)}.
\end{equation}
\item[(d2)] Every chain object $q$ of type (e-ii) of the chain construction, at real parameters, obeys
$|A_q|\le C_{LR'}\,u$ on its own set, with $A_q$ read in the frame of the root.
\item[(d3)] The seed points satisfy $P_0,P_0'\in\mathsf{C}_\beta(C_c\varepsilon)$, where $\varepsilon$
is the smallness parameter of the seeds and $C_c$ a constant, both fixed with the seed class.
\end{enumerate}

\item[(e)] \emph{Format of the real lettered family.} Consider the generation rules (F0)--(F2) of
$\mathfrak{F}^{\mathbb{R},\mathbb{Y}}$ in Definition~\ref{9.5:def-lettered-links}. Rule (F2) applies a
stage-$\nu$ move $\mu=(c,A,K,Z)$ only to parents on a region containing the box $Z$, and it restricts
the output to $K^{\circ}$ together with the strip. At sites of type (P), clause $(\mathbb{Y})^{+}$ of
Definition~\ref{9.5:def-lettered-cutoffs} is, word for word, the clause for moves over the block set
$A$ of Definition~\ref{9.5:def-lettered-links}, read with the core letter $\mathsf{d}(e)$
replaced by $\mathsf{d}^\theta(e)$, where
\begin{equation*}
\mathsf{d}^\theta(e)\supseteq\bigcup_{\square\in A}\bigl(\tilde{\square}^{4}\cup\operatorname{Cor}(\square)\bigr).
\end{equation*}

\item[(f)] \emph{Holomorphy.}
\begin{enumerate}
\item[(f1)] On the complex polydiscs $\mathfrak{P}^{\mathbb{Y}}$, every lettered configuration is
defined.
\item[(f2)] On $\mathfrak{P}^{\mathbb{Y}}$, every lettered configuration is holomorphic in $b$, the
plain letters, the move parameters $t=(t_\square)$ and the seed parameters.
\item[(f3)] For every lettered link, with $\chi$ its cut-off, $\mathfrak{h}$ its increment field, $a$
the majorant of (c), and $d=4$ the dimension,
\begin{equation*}
35\,d\cdot\sup|\chi\,\mathfrak{h}|\le a .
\end{equation*}
\item[(f4)] Let $F$ be a $\mathfrak{g}^{\mathbb{C}}$-valued function,
holomorphic on $\{|z|<3\}$, with $|F|\le\mathsf{M}$ there and $|F|\le m$ on the real segment of that
disc. Then, with an absolute constant $c$,
\begin{equation*}
|F(z)|\le c\,\dim\mathfrak{g}\,(2m\mathsf{M})^{1/2}\qquad(|z|\le1).
\end{equation*}
\end{enumerate}

\item[(g)] \emph{Letter radii.} The two conditions on the parameters that give the letter radii the
lower bound $r_\square\ge2$ are imposed with $8C_\chi$ in place of $2$.
\end{enumerate}
Items (a), (b), (c), (d1) and (f4) are, respectively, the summation lemma for stage majorants, the
composition identity for pairs, the frame bound for lettered links, the mismatch bound for seed pairs
and the square-root lemma, and they are used here as stated in those lemmas. The seed class (d3) is
that of the seed construction, used as stated there. Item (e) restates part of
Definitions~\ref{9.5:def-lettered-links} and~\ref{9.5:def-lettered-cutoffs}, and (g) consists of
conditions on the auxiliary parameters.
Item (f2) is proved as an implication, from the hypotheses listed in it, in the holomorphy lemma for
lettered configurations.
The mismatch law (d2) and items (f1) and (f3) are not proved in this book.
Every statement of this section that uses them carries them as hypotheses.
\end{definition}

\begin{lemma}[Single-lattice containment of the lettered majorant. Proof outstanding]
\label{9.5:lem-majorant-containment}
Fix one of the two machines, the coarse-lattice machine $\mathsf{A}$ or the fine-lattice machine
$\mathsf{B}$, so that everything below is read on a single lattice. Let $\mathsf{s}$ be a site and $Z$
the box read at $\mathsf{s}$; at sites of type (P) this is the box $Z=K\cup Y'$ of the move of
Definition~\ref{9.5:def-lettered-links}. Let $\bar{\tau}_\chi$ be the
letter profile and $C_\chi$ the constant of Definition~\ref{9.5:def-letter-profile}, let
$\mathsf{P}_b=|b|^{\sim}$ be the smeared profile of the source datum $b$ of
Definition~\ref{9.5:def-slow-functions}, let $\iota$ be the level function of the label read at
$\mathsf{s}$ and $\ell_\iota$ its scale, let $a$ be the majorant function of the conversion and
$a_{R'}$ the ceiling constant of the theorem. Put
$\mathsf{H}_\chi:=C_\chi\,\bar{\tau}_\chi\,\mathsf{P}_b$, the \emph{lettered majorant}. Assume that
at sites of type (P) the cut-off is $\tau=\tau_t$ of
Definition~\ref{9.5:def-lettered-links} with real letters $|t_\square|\le 4r^{\bullet}_\square$, and
that at sites of type (T) and (L) the letters satisfy $|t_\square|\le r'_\square$, with
\begin{equation*}
r'_\square:=r_\square/C_\chi .
\end{equation*}
Then at every bond $x$ of $Z$
\begin{equation}\label{9.5:eq-majorant-containment-a}
\mathsf{H}_\chi(x)=C_\chi\,\bar{\tau}_\chi(x)\,\mathsf{P}_b(x)
\le\min\Bigl\{a_{R'},\ \tfrac14\,\ell_\iota(x)\,a(x)\Bigr\}.
\end{equation}
\end{lemma}

Proof outstanding.

\medskip
\noindent\emph{Route.} At sites of type (P), with $\tau=\tau_t$ and real letters
$|t_\square|\le 4r^{\bullet}_\square$, the display \eqref{9.5:eq-majorant-containment-a} is to be
obtained from the rate estimate for the majorant of the response, taken together with its accompanying
condition on the data and the sharpened increment bound of the one-lattice response theorem: these
give the majorant bounded by $\tfrac14 a_\nu^{\mathrm{loc}}$, one quarter of the local stage majorant,
for real letters $|t_\square|\le 4r^{\bullet}_\square$, once a constant of the rate
estimate has been shrunk by a factor depending only on the dimension and on $(L,M_1,M)$; the
constant $C_\chi$ of Definition~\ref{9.5:def-letter-profile} is one admissible such shrinking factor.
At sites of type (T) and (L), the display is to be obtained from
the three majorant inequalities of the lemma on strip-letter majorants: the first majorant is at most
one half of the majorant that lemma attaches to the rim of the strip (its rim majorant); the second
majorant is a prefactor times $e^{-\delta_1 d/8}$, $d$ the scaled path metric of
Definition~\ref{9.5:def-slow-functions}; and that prefactor is at most one quarter of the rim constant
of that lemma. These inequalities are to be read with the
profile $\mathsf{P}_b$ of Definition~\ref{9.5:def-slow-functions} in place of the profile used in that
lemma, and with letters up to $r'_\square=r_\square/C_\chi$. Since the majorant there is linear in the
letter radius, dividing the radius by $C_\chi$ compensates the factor $C_\chi$ in $\mathsf{H}_\chi$, and
the three inequalities then yield \eqref{9.5:eq-majorant-containment-a} at every bond of $Z$. What a
proof has to supply at (T) and (L) is this re-reading carried out in full: the three inequalities with
$\mathsf{P}_b$ in place of the lemma's own profile, together with the allotment of the constants among
them, on letters $|t_\square|\le r'_\square$.

\begin{remark}[What the comparison proof reads]\label{9.5:rem-comparison-inputs}
Consider the proof of the comparison statement whose hypotheses and conclusions (A)--(D) are
collected in Notation~\ref{9.5:not-hypotheses-comparison}. We sort what each of its steps uses
into five kinds:
\begin{itemize}
\item[($\alpha$)] the current admissible sets;
\item[($\beta$)] the presentation of the parent point, in the sense of
Definition~\ref{9.5:def-total-variation};
\item[($\gamma$)] the cut-off;
\item[($\delta$)] the box $Z$ as against its fat core $K$;
\item[($\varepsilon$)] the region over which the exponential smearing $f^{\sim}$ is summed.
\end{itemize}
For a presented pair $p$ we use the level condition of
Definition~\ref{9.5:def-total-variation}: on the region read,
$\lambda_{\mathfrak{S}_s}\ge\lambda$ for every $s\le r$. Then the following hold.
\begin{itemize}
\item The presentation, and with it the nesting of the sets $\mathfrak{S}_s$ and the level
condition, is read only in parts (a) and (b) of the class lemma and in part (ii) of the pair
estimate.
\item The cut-off is read only through a slow $C^2$ majorant.
\item The box is read only as a box, except in the conversion carried out in the rate step and
the concluding estimate of that proof, where the index at points off $K^{\circ}$ is the one
of the box.
\item The smearing $f^{\sim}$ is always summed over the whole region carrying $\mathsf{T}_p$.
\item The inclusions $K\supseteq K_A$ and $K\supseteq X^{+}$, of the fat core $K$ over the core
$K_A$ of the machine $\mathsf{A}$ and over the region $X^{+}$, are read in
none of the steps of that proof preceding the rate step.
\end{itemize}
\end{remark}

\begin{proof}
We go through the steps of the proof of the comparison statement in their order and record,
for each, what it reads of ($\alpha$)--($\varepsilon$).

\begin{enumerate}
\item[(1)] The one-machine lemmas: the bounds on the sources, on the response field
$\mathbb{H}_{\mathfrak{S}}$, stability, regularity, the auxiliary one-machine estimate, the
background-field bounds, and the Lipschitz bound for Ba{\l}aban's projection $R_{\mathfrak{S}}$.
Each is an estimate in one machine, on the current sets only. It reads the hypothesis
\cite[(3.35)]{B-CMP99b} (printed without proof in \cite[(3.35)]{B-CMP99b}; cited as printed)
and the regularity class at the current scale. On the shells
$\boldsymbol{\Omega}(Z)$ it holds with constants that do not depend on $Z$, by the standing
hypotheses on boxes of Notation~\ref{9.5:not-hypotheses-comparison}. It reads nothing of
($\beta$)--($\delta$).

\item[(2)] The nest estimates and the bounds proved from them. These are: the opening nest
estimate; the four nest inequalities for the slot content $\mathbf{J}$; part (a) of the nest
bound that follows these inequalities; the estimate that follows that bound; the divergence
bound; the comparison bound; the
nest bound for Ba{\l}aban's projection; the inversion bound; the nest $N\mathcal{W}$; and the
penalty bound. They read three things: the current exactly
paired sets; the moduli of the fields at the scale of the sets in the statement; and the
pointwise index $L^{-\beta\lambda}$ or $L^{-\theta\lambda}$ at the point of evaluation. They read
nothing of ($\gamma$) or ($\delta$).

\item[(3)] Parts (a) and (b) of the class lemma place the class of the parent point in the
current units, and here the level condition is read. Put
$\mathcal{J}_s:=\mathcal{J}_{\mathfrak{S}_s}(\mathfrak{h}_s;\mathbf{U}_{s-1})$. An increment
$\mathfrak{h}_s\in\operatorname{Inc}_\beta(\mathsf{r}_s)$, taken in the units of its link's set
$\mathfrak{S}_s$, lies in $\operatorname{Inc}_\beta(\mathsf{r}_s)$ in the current units at the
points where $\ell\le\ell_s$. These are the points where $\lambda\le\lambda_s$, with
$\lambda_s:=\lambda_{\mathfrak{S}_s}$ and $\ell_s$ the local scale of $\mathfrak{S}_s$. The level
condition gives exactly $\lambda\le\lambda_s$ on the region read. The same holds for
$\ell^3|\mathcal{J}_s|$. Part (c) of the class lemma is an identity valid for every regular
$\mathbf{U}$, and it reads no presentation.

\item[(4)] The proof of (B) has three steps.
\begin{itemize}
\item The first step rests on two inputs: the comparison bound of step (2) and a separate
proposition of the same proof.
\item The second step consists of six bracketed estimates. They use the background-field bounds,
the nest $N\mathcal{W}$, the Lipschitz bound for $R_{\mathfrak{S}}$, the auxiliary bound of the
second step, the auxiliary one-machine estimate of step (1) in the machine $\mathsf{B}$, the
divergence bound, the nest bound for Ba{\l}aban's projection, and the comparison bound
differentiated.
\item The third step is the stability estimate.
\end{itemize}
All three steps are carried out on the current sets. Of the parent point they read three things:
$\sup|\mathfrak{a}|$, with $|\mathfrak{a}|_y\le\mathsf{T}_p(y)$ on the current region; the class;
and $|\mathsf{m}'-\mathsf{m}|$, the difference of the data majorants of the two points, through
(A). By the splitting of conclusion (B) established for the comparison statement, conclusion (B)
splits as $(\mathrm{B})=(\mathrm{B})^{\mathrm{free}}\wedge(\mathrm{A})$, where
$(\mathrm{B})^{\mathrm{free}}$ denotes the conjunct of (B) other than (A); and
$(\mathrm{B})^{\mathrm{free}}$ does not read the presentation.

The third step uses that $\Psi_p$ is slow. Two facts give this. First, $\mathsf{T}_p^{\sim}$ is
slow of type $(1,1)$ for every $\mathsf{T}_p$, by part (i) of the statement that defines
slowness and its type. Second, $L^{-\theta\lambda}$ is slow of type $(L^{\theta},\theta/16)$, by
the first
inequality of \eqref{9.5:eq-hypotheses-geometry-a} on the current sets. As for
($\varepsilon$): $f^{\sim}$ is summed over the whole region carrying $\mathsf{T}_p$.

\item[(5)] Part (i) of the pair estimate reads the class and the bound
$|A_{p_s}|_{\cdot}\le\mathsf{T}_p$, where $A_{p_s}$ is the mismatch field of the pair $p_s$
formed at the $s$-th link of the presentation. It reads no earlier sets.

Part (ii) of the pair estimate reads several things: the penalty bound, the nest bound for
Ba{\l}aban's projection and the divergence bound at the sets
$(\mathfrak{S}'_s,\mathfrak{S}_s)$ of the earlier link; the Lipschitz bound for $R_{\mathfrak{S}}$;
and, a second time, the level condition. The defects of stage $s$ carry $L^{-\beta\lambda_s}$
and are counted as $L^{-\beta\lambda}$. This is admissible because $\lambda\le\lambda_s$ gives
$L^{-\beta\lambda_s}\le L^{-\beta\lambda}$. Moreover $\sum_s|e_s|\le 3\,\mathsf{T}_p$ by the
composition identity, and $\sum_s\mathsf{r}_s\le\mathfrak{r}_0$. Hence no constant depends on
the number $r$ of links.

\item[(6)] The proof of (A) combines the pair estimate with a duality argument on the current
sets. It reads part (b) of the class lemma, that is, the boundedness of $\mathbf{J}$ on the
current region.

\item[(7)] The proofs of (C) and (D) read (B), the regularity estimate, the composition identity,
and the slot rule with the same cut field. Of $\tau$ they read only $\bar{\tau}$, $\ell|\nabla\tau|$
and $\ell^2|\nabla^2\tau|$, that is, a slow $C^2$ majorant.

\item[(8)] The naturality step holds for any box $Z$, with the constants independent of $Z$
noted in step (1), by the standing hypotheses on boxes of
Notation~\ref{9.5:not-hypotheses-comparison}.

\item[(9)] The rate step is an estimate in one machine. Here, and only here in the proof, the
deficit of the box enters. Write $R_\nu\subset K^{\circ}$ for the source regions of this step,
$\nu$ indexing the stages; these are not Ba{\l}aban's projections $R_{\mathfrak{S}}$. At points
off $K^{\circ}$, the conversion of $\mathsf{H}$ to the units of the label is paid for by the
decay of $\mathsf{H}$ away from the sources $R_\nu$. This decay is furnished by the decay
hypothesis (d3) of Notation~\ref{9.5:not-hypotheses-geometry} and by the increment bound.

\item[(10)] The concluding estimate reads three things: the yardstick
$\sup\mathsf{T}_p/(\ell u)$; the weight hypothesis in
\eqref{9.5:eq-hypotheses-geometry-a}; and the decay hypothesis (d3) of step (9). The local growth
bound is applied to the total variation $\mathsf{T}_p$ without modification. The passage to
complex values costs one factor $2C$.
\end{enumerate}

Reading the list off gives each item of the remark.
\begin{itemize}
\item The presentation ($\beta$) is read only through parts (a) and (b) of the class lemma and
through part (ii) of the pair estimate: parts (a) and (b) are established in step (3), the class
of the parent point that they place in the current units is read again in step (4), and part (b)
is invoked again in step (6); part (ii) of the pair estimate enters in step (5). No other step
reads the presentation.
\item The cut-off appears in step (7) only through $\bar{\tau}$, $\ell|\nabla\tau|$,
$\ell^2|\nabla^2\tau|$, and in step (9) only through the majorant $\mathsf{H}$ of
Notation~\ref{9.5:not-hypotheses-comparison}, which carries it through $\bar{\tau}$; in every
case it is read only through a slow $C^2$ majorant.
\item The box appears in steps (1) and (8) only as a box, with constants free of $Z$, and as more
than a box only in steps (9) and (10).
\item The summation of $f^{\sim}$ enters the list at step (4), and there it is over the whole
region carrying $\mathsf{T}_p$.
\item No step among (1)--(8) reads either inclusion $K\supseteq K_A$ or $K\supseteq X^{+}$. \qedhere
\end{itemize}
\end{proof}

\begin{lemma}[Slow twice-differentiable majorant of the lettered cut-off]\label{9.5:lem-cutoff-majorant}
Let $\chi=\sum_{\square}c_\square\zeta_\square$ be the cut-off of a letter assignment
(Definitions~\ref{9.5:def-letter-profile} and~\ref{9.5:def-lettered-cutoffs}) whose letters satisfy
$|c_\square|\le c_0^\chi$ off the strip blocks and
\begin{equation*}
|c_\square|\le \mathfrak{r}_0^\chi\, e^{\delta_1 d(\tilde{\square},\Omega^{\mathsf{s}})/8}
\quad\text{on the strip blocks,}
\end{equation*}
where $d$ is the metric of the label, and let $\bar{\tau}_\chi$ be the letter profile of
\eqref{9.5:eq-letter-profile-a}. Then, in the units of any admissible set $\mathfrak{S}$ refining the
label (in particular $\boldsymbol{\Omega}(Z)$), with $\ell=\ell_{\mathfrak{S}}$ its local scale:
\begin{itemize}
\item[(a)] $|\chi(y)|\le \bar{c}(y)\le \bar{\tau}_\chi(y)$;
\item[(b)] $\ell(y)\,|\nabla\chi(y)|\le C_\chi\bar{\tau}_\chi(y)/M$ and
$\ell(y)^2\,|\nabla^2\chi(y)|\le C_\chi\bar{\tau}_\chi(y)/M^2$;
\item[(c)] $\bar{\tau}_\chi$ is slow of type $(1,1/8)$, hence slow of type $(L,1)$, for every one of these
sets, in the sense of Definition~\ref{9.5:def-slow-functions};
\item[(d)] $\bar{\tau}_\chi$ and $C_\chi$ do not depend on $|D|$, on the number of blocks with non-zero
letters, or on the number of dilated letters.
\end{itemize}
\end{lemma}

\begin{proof}
Recall from Definition~\ref{9.5:def-letter-profile} the properties of the cover and of its partition of
unity that are used: $\operatorname{supp}\zeta_\square\subset\tilde{\square}$, $\zeta_\square\ge 0$,
$\sum_\square\zeta_\square=1$, at most $2^d$ of the $\zeta_\square$ are non-zero at a point, and
\begin{equation}\label{9.5:eq-cutoff-majorant-1}
\ell_\square\,|\nabla\zeta_\square|\le C_\zeta/M,\qquad
\ell_\square^2\,|\nabla^2\zeta_\square|\le C_\zeta/M^2,
\end{equation}
with $\ell_\square$ the scale of the layer of $\square$ and $C_\zeta$ a number
(\cite{B-CMP109}). The quantity $\bar{c}(y)$ bounds $|c_\square|$ for every block $\square$ with
$\tilde{\square}$ meeting $\tilde{\Delta}(y)$.

(a) Since $\zeta_\square(y)\neq0$ forces $y\in\tilde{\square}$, every block contributing at $y$ has
$\tilde{\square}$ meeting $\tilde{\Delta}(y)$, and therefore
\begin{equation*}
|\chi(y)|\le \sum_\square |c_\square|\,\zeta_\square(y)
\le \bar{c}(y)\sum_\square\zeta_\square(y)=\bar{c}(y).
\end{equation*}
For the second inequality let $\square$ be a strip block with $\tilde{\square}$ meeting
$\tilde{\Delta}(y)$, and pick $z\in\tilde{\square}\cap\tilde{\Delta}(y)$. By the triangle inequality for $d$,
\begin{equation*}
d(\tilde{\square},\Omega^{\mathsf{s}})\le d(z,\Omega^{\mathsf{s}})
\le d(y,\Omega^{\mathsf{s}})+d(y,z)
\le d(y,\Omega^{\mathsf{s}})+\operatorname{diam}_d\bigl(\tilde{\Delta}(y)\cup\tilde{\square}\bigr),
\end{equation*}
so that, with the constant $C_\Delta$ of \eqref{9.5:eq-letter-profile-a} bounding
$e^{\delta_1\operatorname{diam}_d(\tilde{\Delta}(y)\cup\tilde{\square})/8}$,
\begin{equation}\label{9.5:eq-cutoff-majorant-2}
|c_\square|\le \mathfrak{r}_0^\chi\, C_\Delta\, e^{\delta_1 d(y,\Omega^{\mathsf{s}})/8}.
\end{equation}
For the other blocks with $\tilde{\square}$ meeting $\tilde{\Delta}(y)$ one has $|c_\square|\le c_0^\chi$.
The profile $\bar{\tau}_\chi(y)$ of \eqref{9.5:eq-letter-profile-a} dominates both $c_0^\chi$ and the
right-hand side of \eqref{9.5:eq-cutoff-majorant-2}; hence $\bar{c}(y)\le\bar{\tau}_\chi(y)$.

(b) We have $\nabla\chi=\sum_\square c_\square\nabla\zeta_\square$, and by the multiplicity of the cover
at most $2^d$ terms are non-zero near $y$; each of them belongs to a block with $y\in\tilde{\square}$, so its
letter is bounded by $\bar{c}(y)$, and by \eqref{9.5:eq-cutoff-majorant-1} each term is at most
$\bar{c}(y)C_\zeta/(M\ell_\square)$. The enlargement $\tilde{\square}$ meets at most the two layers adjacent
to the layer of $\square$, so the scale of the block's layer and the label's scale $\ell_\iota$ near $y$
differ by at most a factor $L$; and the scale of a set refining the label is at most the label's. Thus
$\ell(y)\le\ell_\iota(y)\le L\ell_\square$ for each contributing block, and
\begin{equation*}
\ell(y)\,|\nabla\chi(y)|
\le 2^d\,\bar{c}(y)\,\frac{C_\zeta}{M}\,\frac{\ell(y)}{\ell_\square}
\le \frac{2^d C_\zeta L\,\bar{c}(y)}{M}
\le \frac{C_\chi\bar{\tau}_\chi(y)}{M},
\end{equation*}
the last step by (a). The second derivative is treated in the same way:
$\nabla^2\chi=\sum_\square c_\square\nabla^2\zeta_\square$ has at most $2^d$ non-zero terms near $y$, each
at most $\bar{c}(y)C_\zeta/(M^2\ell_\square^2)$ by the second bound of \eqref{9.5:eq-cutoff-majorant-1},
and $\ell(y)^2\le L^2\ell_\square^2$, so that
\begin{equation*}
\ell(y)^2\,|\nabla^2\chi(y)|
\le \frac{2^d C_\zeta L^2\,\bar{c}(y)}{M^2}
\le \frac{C_\chi\bar{\tau}_\chi(y)}{M^2}.
\end{equation*}

(c) For all $x,y$ the triangle inequality for the path metric $d$ gives
$d(y,\Omega^{\mathsf{s}})\le d(x,y)+d(x,\Omega^{\mathsf{s}})$, hence
\begin{equation*}
e^{\delta_1 d(y,\Omega^{\mathsf{s}})/8}
\le e^{\delta_1 d(x,y)/8}\,e^{\delta_1 d(x,\Omega^{\mathsf{s}})/8}.
\end{equation*}
The maximum of a constant $c$ and a function $f$ with $f(y)\le e^{\delta_1d(x,y)/8}f(x)$ satisfies
\begin{equation*}
\max\{c,f(y)\}\le\max\bigl\{c,\,e^{\delta_1 d(x,y)/8}f(x)\bigr\}
\le e^{\delta_1 d(x,y)/8}\max\{c,f(x)\},
\end{equation*}
because $e^{\delta_1 d(x,y)/8}\ge1$. Therefore
$\bar{\tau}_\chi(y)\le e^{\delta_1 d(x,y)/8}\,\bar{\tau}_\chi(x)$, that is, $\bar{\tau}_\chi$ is slow of type
$(1,1/8)$ with respect to $d$. For a set $\mathfrak{S}$ refining the label,
$d\le d_{\mathfrak{S}}$ pointwise by the monotonicity of the metric under refinement
(Notation~\ref{9.5:not-hypotheses-geometry}, \eqref{9.5:eq-hypotheses-geometry-a}), so
$e^{\delta_1 d(x,y)/8}\le e^{\delta_1 d_{\mathfrak{S}}(x,y)/8}$ and the type only improves. Since $1\le L$
and $e^{\delta_1 d_{\mathfrak{S}}(x,y)/8}\le e^{\delta_1 d_{\mathfrak{S}}(x,y)}$, type $(1,1/8)$ implies type
$(L,1)$.

(d) This holds by construction: $\bar{\tau}_\chi$ and $C_\chi$ are fixed in
\eqref{9.5:eq-letter-profile-a} from the data listed there, and the proofs of (a)--(c) used, besides
these, only the multiplicity $2^d$ of the cover, the constant $C_\zeta$ and the factor $L$ between adjacent
layer scales; none of these involves $|D|$, the number of blocks with non-zero letters, or the number of
dilated letters.
\end{proof}

\begin{remark}
In what follows the cut-off $\chi$ enters the response bounds through the profile
$\bar{\tau}:=C_\chi\bar{\tau}_\chi$. With this choice Lemma~\ref{9.5:lem-cutoff-majorant} gives
\begin{equation*}
|\chi|\le\bar{\tau},\qquad \ell\,|\nabla\chi|\le\bar{\tau}/M,\qquad
\ell^2\,|\nabla^2\chi|\le\bar{\tau}/M^2,
\end{equation*}
and the lettered majorant is $\mathsf{H}_\chi=C_\chi\bar{\tau}_\chi\mathsf{P}_b$. The bound into which the
profile enters is linear in it; complex letters are nevertheless treated separately, by the proposition on
complex letters.
\end{remark}

\begin{lemma}[Points of a run are presented pairs over entries]\label{9.5:lem-run-points}
Let $p_0\to p_1\to\dots\to p_r$ be a run presentation in the family
$\mathfrak{F}^{\mathbb{R},\mathbb{Y}}$ of Definition~\ref{9.5:def-lettered-links}, with links
$\mu_s=(c_s,A_s,K_s,Z_s)$ of stages $\nu_1>\nu_2>\dots>\nu_r$, with the letter vectors
$t_s=(t_\square)_\square$ of the real parameters of the moves
(Definition~\ref{9.5:def-lettered-links}), and with cut-offs
$\tau_s:=\tau_{t_s}$, $1\le s\le r$. Write $K_0^{\circ}$ for the region of the entry $p_0$. For
each $s$ let $S_s^{+}$ denote the fattening of $\bigcup_{\square\in A_s}\tilde{\square}$ by the
local cubes $\tilde{\Delta}(y)$ and by the stencils of the bonds they contain. Then for every $s$:
\begin{enumerate}
\item[(i)] $Z_s\subset K_{s-1}^{\circ}\subset Z_{s-1}$;
\item[(ii)] the increment $\tau_s\mathbb{H}_{Z_s}$ and its $\mathsf{T}$-increment
$|A_{p_s}-A_{p_{s-1}}|_{\cdot}$ vanish outside $S_s^{+}$, and
$S_s^{+}\subset K_s^{\circ}\subset Z_s^{=}$;
\item[(iii)] the truncation $p_{s-1}$, read on $Z_s$, satisfies the level condition
\eqref{9.5:eq-total-variation-a} and the entry level condition of
Definition~\ref{9.5:def-entries}, as written, with respect to the current sets
$(\boldsymbol{\Omega}'(Z_s),\boldsymbol{\Omega}(Z_s))$;
\item[(iv)] $\tau_s$ has the slow twice-differentiable majorant
$C_\chi\bar{\tau}_{\tau_s}$ of Lemma~\ref{9.5:lem-cutoff-majorant}, with $c_0^\chi=0$ and with
$\mathfrak{r}_0^\chi$ the constant of the exponential bound on the letter radii
$r^{\bullet}_\square$ that enter the move rule of Definition~\ref{9.5:def-lettered-links}, whose
rate $\tfrac12\delta$ satisfies $\tfrac12\delta\le\delta_1/8$ in the metric $d$, in either of
the two cases of that bound.
\end{enumerate}
Consequently, if $p_0$ is an entry and the ceilings
$\mathfrak{r}_P,\mathfrak{r}_{P'}\le\mathfrak{r}_0$ on the class radii of $p_{s-1}$ and
$\sup\mathsf{T}_{p_{s-1}}\le\varepsilon_D$ hold, then $p_{s-1}$ on $Z_s$ satisfies every
hypothesis that the comparison statement forming the second item of
Notation~\ref{9.5:not-hypotheses-comparison} imposes on its parent at the link $\mu_s$.
\end{lemma}

\begin{proof}
(i) By the move rule of Definition~\ref{9.5:def-lettered-links}, the stage-$\nu_s$ move $\mu_s$ is
applied to a parent on a region $\hat{X}\supseteq Z_s$. This parent is $p_{s-1}$, which was
produced at stage $\nu_{s-1}>\nu_s$; by the restriction built into the move rule, together with
the restriction rule of the same definition, it lives on $K_{s-1}^{\circ}$ (and on the strip).
Hence $Z_s\subset\hat{X}\subset K_{s-1}^{\circ}$. Finally
$K_{s-1}^{\circ}\subset K_{s-1}\subset Z_{s-1}$,
since by the definition of the box in Definition~\ref{9.5:def-lettered-links} the box
$Z_{s-1}$ is the union of the fat core $K_{s-1}$ with a further set.

(ii) The letters $t_\square$ of the cut-off $\tau_s$ vanish for $\square\notin A_s$, so
\begin{equation*}
\operatorname{supp}\tau_s\subset\bigcup_{\square\in A_s}\tilde{\square}\subset S_s^{+},
\end{equation*}
and therefore $\tau_s\mathbb{H}_{Z_s}$ vanishes outside $S_s^{+}$. The fat core satisfies
$K_s\supseteq\bigcup_{\square\in A_s}\tilde{\square}^{4}$ by
Definition~\ref{9.5:def-lettered-links}; the first of the block-geometry hypotheses
\eqref{9.5:eq-hypotheses-geometry-b}, applied to $K_s$ and the blocks of $A_s$, gives
$S_s^{+}\subset K_s^{\circ}$, and the second gives $K_s^{\circ}\subset Z_s^{=}$. It remains to
treat the $\mathsf{T}$-increment. At a point $y$ it equals
\begin{equation*}
\ell(y)\sup_{\tilde{\Delta}(y)}|A_{p_s}-A_{p_{s-1}}|.
\end{equation*}
By the composition identity of Notation~\ref{9.5:not-hypotheses-summation}, which expresses
$e^{i\eta A_{p_s}}$ through $e^{i\eta A_{p_{s-1}}}$, the field denoted $Y$ there, and the
increment $\mathfrak{h}_s$ of the link, and by the locality of the block-averaging map $M$ and
of that field $Y$, the difference $A_{p_s}-A_{p_{s-1}}$ vanishes at every bond whose stencil
misses $\operatorname{supp}\tau_s$. Consequently the supremum over $\tilde{\Delta}(y)$ vanishes
unless $\tilde{\Delta}(y)$ contains a bond whose stencil meets
$\operatorname{supp}\tau_s\subset\bigcup_{\square\in A_s}\tilde{\square}$, that is, by the
definition of $S_s^{+}$, unless $y\in S_s^{+}$.

(iii) Fix $q<s$. On the region read we have $Z_s\subset K_q^{\circ}$, by (i) applied
successively to the links $q+1,\dots,s$. By the second block-geometry hypothesis
\eqref{9.5:eq-hypotheses-geometry-b}, $K_q^{\circ}\subset Z_q^{=}$, so on $Z_s$ the level of
$\boldsymbol{\Omega}(Z_q)$ is the label's level at stage $\nu_q$, namely $\iota_{\nu_q}$. The
level of the current sets satisfies, pointwise on $Z_s$,
\begin{equation*}
\lambda_{Z_s}\le\iota_{\nu_s}\le\iota_{\nu_q},
\end{equation*}
the last inequality by the second of the summation hypotheses
\eqref{9.5:eq-hypotheses-summation-a}: the label's level at a point does not decrease as the
stage index increases, and $\nu_q>\nu_s$. Thus every link set of $p_{s-1}$ has level at least
$\lambda_{Z_s}$ on $Z_s$, which is the level condition \eqref{9.5:eq-total-variation-a}. For the
entry level condition of Definition~\ref{9.5:def-entries}, the set on which the regularity
condition of that definition holds is the entry's own set, for each of the two kinds of entry
admitted by the entries rule of Definition~\ref{9.5:def-lettered-links}: the label, for a seed,
and the object's own set $\mathfrak{S}_e$, over which the later boxes are built, for an entry
carried by an object. In either case this
set has level at least $\lambda_{Z_s}$ on $Z_s$, and together with the level condition just
proved this is the entry level condition.

(iv) The profile $\bar{\tau}_{\tau_s}$ of the cut-off $\tau_s$ is the letter profile
$\bar{\tau}_\chi$ of Lemma~\ref{9.5:lem-cutoff-majorant} for the letter assignment $t_s$; at a
point $y$ it is of the exponential type
$\mathfrak{r}_0^\chi C_\Delta e^{\delta_1 d(y,\Omega^{\mathsf{s}})/8}$, with $\Omega^{\mathsf{s}}$
the source region of that lemma. This is the situation of
Lemma~\ref{9.5:lem-cutoff-majorant} with
$c_0^\chi=0$ and with $\mathfrak{r}_0^\chi$ the constant of the bound on the letter radii
$r^{\bullet}_\square$ named in the statement, and that lemma gives the slow
twice-differentiable majorant $C_\chi\bar{\tau}_{\tau_s}$.

For the final assertion, the hypotheses that the comparison statement forming the second item of
Notation~\ref{9.5:not-hypotheses-comparison} imposes on its parent are: the class bound on the
radii, $\mathfrak{r}_P,\mathfrak{r}_{P'}\le\mathfrak{r}_0$; the bound
$\sup\mathsf{T}_{p_{s-1}}\le\varepsilon_D$; and
presentation over an entry with the level condition \eqref{9.5:eq-total-variation-a} and the
entry level condition of Definition~\ref{9.5:def-entries}. The first two are assumed for
$p_{s-1}$; $p_0$ is an entry by assumption; and the two level conditions on $Z_s$ are (iii).
\end{proof}

\begin{remark}[Non-nested boxes]
Suppose a run presentation outside the format of Definition~\ref{9.5:def-lettered-links} had a
link $q$ whose box $Z_q$ does not contain the current box $Z_s$. Then on
$Z_s\setminus Z_q$ the set $\boldsymbol{\Omega}(Z_q)$ has no level, and the level condition
\eqref{9.5:eq-total-variation-a} fails as written. This does not occur in
$\mathfrak{F}^{\mathbb{R},\mathbb{Y}}$: the move rule of Definition~\ref{9.5:def-lettered-links}
forces part~(i) of Lemma~\ref{9.5:lem-run-points}. What would be needed in that case is the
separate argument for run presentations with non-nested boxes.
\end{remark}

\begin{lemma}[The single lettered step meets the comparison hypotheses]\label{9.5:lem-single-step}
Consider the $(\mathrm{T})$-family of lettered links of Definition~\ref{9.5:def-lettered-links},
with seed pair $p_0$, whose point has configurations $\mathbf{U}_0,\mathbf{U}_0'$. Let $D$ be the
index of any unit $v_D$, with core $K_D$, and let $Z=K_D\cup Y'$. Let $\mathsf{t}$ be a real
bracket letter in the range
prescribed for it in the $(\mathrm{T})$-family, and let the strip letters be real with
$|t_\square|\le r'_\square$. Let $\chi_t$ be the corresponding lettered cut-off and $B'$ the
source datum of the link. Then the following hold.
\begin{enumerate}
\item[(i)] $p_0\restriction Z$ is a diagonal pair satisfying condition (e1) of
Definition~\ref{9.5:def-entries}. Hence, by part (a) of the comparison theorem over entries,
it is an entry of modulus $(|A_{p_0}|_{\cdot},C_i)$, where $C_i$ is the constant of part (a)
of that theorem.
\item[(ii)] The condition \eqref{9.5:eq-total-variation-a} holds. So does the level condition of
Definition~\ref{9.5:def-entries}: on the region read, the sets of the links and the set on
which (e1) holds have level at least $\lambda_Z$.
\item[(iii)] $\chi_t$ has the slow twice-differentiable majorant $C_\chi\bar{\tau}_{\chi_t}$ in
the sense of Lemma~\ref{9.5:lem-cutoff-majorant}, with
\begin{equation*}
c_0^\chi:=1+|\mathsf{t}|,\qquad \mathfrak{r}_0^\chi:=r^0_T/C_\chi .
\end{equation*}
\item[(iv)] The increment $\chi_t\,\mathbb{H}_Z(B';P_0)$ is supported in
\begin{equation*}
Z\cap\bigcup\{\tilde{\square}:\ c_\square\neq0\},
\end{equation*}
where $c_\square$ is the letter of $\square$.
This set in general meets $Z\setminus Z^{=}$, since the cubes of $Y'$ and the rim of $K_D$
carry plain blocks with letter $1$. The cells at which $v_D$ reads an increment lie in
$K_D^{\circ}$, and $K_D^{\circ}\subset Z^{=}$.
\end{enumerate}
Suppose moreover that $\sup_Z|A_{p_0}|_{\cdot}\le\varepsilon_D$, and that the containment
inequality \eqref{9.5:eq-majorant-containment-a} holds at every bond of $Z$. That inequality is
the content of Lemma~\ref{9.5:lem-majorant-containment}, whose proof is outstanding. Then every
hypothesis of part (b) of the comparison theorem over entries holds at the link
$(Z,\chi_t,B')$ over $p_0\restriction Z$.
\end{lemma}

\begin{proof}
\emph{(i).} The slot contents of the seed pair are
\begin{equation*}
\mathbf{J}_0=J(\mathbf{U}_0),\qquad \mathbf{J}_0'=J(\mathbf{U}_0').
\end{equation*}
These identities hold by the construction of the seed pair, in which each slot carries the current
of its own configuration; they are exactly the defining condition of a diagonal pair in
Definition~\ref{9.5:def-entries}, so $p_0\restriction Z$ is diagonal.

Condition (e1) requires $P_0,P_0'\in\mathsf{C}_\beta(\mathfrak{r}_e)$. In the units of the label
this is the standing hypothesis on the seed class. It also holds in the units of
$\boldsymbol{\Omega}(Z)$, for the following reason. The set $\boldsymbol{\Omega}(Z)$ refines the
label (Notation~\ref{9.5:not-hypotheses-geometry}). Its local scale at
each point is therefore at most the label's local scale $\ell_\iota$. The conditions
\cite[(3.35)]{B-CMP99b} and the class $\mathsf{C}_\beta$ weaken as the scale decreases, since each
quantity in their definition carries a positive power of $\ell$
and so decreases with $\ell$. Hence (e1) holds a fortiori in the units of
$\boldsymbol{\Omega}(Z)$. Part (a) of the
comparison theorem over entries then makes $p_0\restriction Z$ an entry of modulus
$(|A_{p_0}|_{\cdot},C_i)$.

\emph{(ii).} The pair $p_0\restriction Z$ is a restriction of the seed pair, so its presentation
contains no earlier link. Condition \eqref{9.5:eq-total-variation-a} therefore involves no set of
an earlier link, and it is satisfied by the seed presentation alone.
For the same reason, the level condition of Definition~\ref{9.5:def-entries}
reduces to the set on which (e1) holds, namely the label. The label's level function satisfies
$\iota\ge\lambda_Z$ on $Z$. Both conditions therefore hold.

\emph{(iii).} We apply Lemma~\ref{9.5:lem-cutoff-majorant} to $\chi_t$, with the ranges of the
letters in the $(\mathrm{T})$-family. These ranges give the hypotheses of that lemma with
$c_0^\chi:=1+|\mathsf{t}|$ off the strip blocks and $\mathfrak{r}_0^\chi:=r^0_T/C_\chi$ on the
strip blocks. Its conclusions hold in the units of $\boldsymbol{\Omega}(Z)$. Part (a) of that
lemma gives $|\chi_t|\le\bar{c}\le\bar{\tau}_{\chi_t}$, part (b) bounds
$\ell|\nabla\chi_t|$ and $\ell^2|\nabla^2\chi_t|$ by
$C_\chi\bar{\tau}_{\chi_t}/M$ and $C_\chi\bar{\tau}_{\chi_t}/M^2$, and part (c) gives the
slowness of $\bar{\tau}_{\chi_t}$. Together these say that $C_\chi\bar{\tau}_{\chi_t}$ is a slow
twice-differentiable majorant of $\chi_t$.

\emph{(iv).} The cut-off is $\chi_t=\sum_\square c_\square\zeta_\square$, and $\zeta_\square$
vanishes off $\tilde{\square}$. Hence
\begin{equation*}
\operatorname{supp}\chi_t\subset\bigcup\{\tilde{\square}:\ c_\square\neq0\}.
\end{equation*}
The response field $\mathbb{H}_Z$ is, by its definition, a field on $Z$ that vanishes on the
outermost stratum. The product $\chi_t\,\mathbb{H}_Z(B';P_0)$ is therefore supported in the
intersection displayed in (iv). The blocks of $Y'$ and of the rim of $K_D$ are plain blocks with
letter $1$, so this set in general meets $Z\setminus Z^{=}$. The last statement of (iv) follows
from the second and third conditions of the block geometry \eqref{9.5:eq-hypotheses-geometry-b}
of Notation~\ref{9.5:not-hypotheses-geometry}. Together they place the cells read by $v_D$ in
$K_D^{\circ}$, and give $K_D^{\circ}\subset Z^{=}$.

\emph{The hypotheses of part (b) of the comparison theorem over entries.} We check them in turn.
\begin{itemize}
\item The sets $(\boldsymbol{\Omega}'(Z),\boldsymbol{\Omega}(Z))$ are exactly paired admissible,
by the standing hypothesis on boxes of Notation~\ref{9.5:not-hypotheses-comparison}.
\item The parent $p_0\restriction Z$ is an entry satisfying (e1), by (i). Its modulus is
$\mathsf{E}=|A_{p_0}|_{\cdot}$, so the assumption $\sup_Z|A_{p_0}|_{\cdot}\le\varepsilon_D$ gives
$\sup\mathsf{E}\le\varepsilon_D$. The level condition of Definition~\ref{9.5:def-entries} holds
by (ii).
\item The source datum is $b=B'$, real and in its prescribed range, with
$\delta b=(\operatorname{Ad}_g-1)B'$.
\item The cut-off is $\tau=\chi_t$. It is real because the letters are real, and by (iii) it has
the slow twice-differentiable majorant $C_\chi\bar{\tau}_{\chi_t}$.
\item It remains to bound $\mathsf{H}$. With $\bar{\tau}=C_\chi\bar{\tau}_{\chi_t}$, the quantity
$\mathsf{H}=\bar{\tau}\,\mathsf{P}_b$ is the lettered majorant $\mathsf{H}_\chi$. The assumed
inequality \eqref{9.5:eq-majorant-containment-a} gives, at every bond $x$ of $Z$,
\begin{equation*}
\mathsf{H}_\chi(x)\le\min\{a_{R'},\tfrac14\ell_\iota(x)a(x)\}\le a_{R'},
\end{equation*}
that is, $\mathsf{H}\le a_{R'}$.
\end{itemize}
All hypotheses of part (b) of the comparison theorem over entries therefore hold at the link
$(Z,\chi_t,B')$ over $p_0\restriction Z$.
\end{proof}

\begin{lemma}[Response-type rows meet the comparison hypotheses]\label{9.5:lem-response-rows}
Consider the $(L)$-family of lettered links (Definition~\ref{9.5:def-lettered-links}) at one box
$Z=K_D\cup Y'$ of the unit $v_D$. Take the parameters real and the strip letters real, with
\begin{equation*}
|t_{\square,r}|\le r^{L}_{\square}/C_\chi .
\end{equation*}
Then the following hold.
\begin{itemize}
\item[(i)] The restriction $p^{(0)}=q_0\restriction Z$ is an entry in the sense of
Definition~\ref{9.5:def-entries}.
\item[(ii)] For every row $r$, let $p^{(r-1)}$ be the truncation of the presented pair after the
rows preceding $r$. Read on $Z$, it satisfies the level condition
\eqref{9.5:eq-total-variation-a}. The sets of
the earlier response rows are the shells of the same box $Z$ over base sets
$\mathfrak{S}_{r'}$, $r'<r$, and the chain only refines these base sets. Consequently,
pointwise on $Z$,
\begin{equation*}
\lambda_{\boldsymbol{\Omega}_{\mathfrak{S}_{r'}}(Z)}\ge\lambda_{\boldsymbol{\Omega}_{\mathfrak{S}_{r}}(Z)} .
\end{equation*}
The level condition for entries in Definition~\ref{9.5:def-entries} holds in the same way,
since the base set $\mathfrak{S}_e$ of the entry is the coarsest.
\item[(iii)] The cut-off $\chi^{(r)}_t$ of row $r$ has the slow twice-differentiable majorant
$C_\chi\bar{\tau}_{\chi^{(r)}}$ of Lemma~\ref{9.5:lem-cutoff-majorant}, formed with $c_0^\chi$
equal to the bound on the plain letters of row $r$ fixed in the data of the $(L)$-family
(Definition~\ref{9.5:def-lettered-cutoffs}) and with $\mathfrak{r}_0^\chi:=r^0_L/C_\chi$.
\item[(iv)] The increment of each response row is supported in
$Z\cap\bigcup\{\tilde{\square}:\ \text{letter}\neq0\}$. In general this set meets
$Z\setminus Z^{=}$. The cells at which $v_D$ reads increments lie in $K_D^{\circ}\subset Z^{=}$.
\end{itemize}
Now let $r$ be a row of response type. Suppose that each of its predecessors is either of
response type or of object type with letters in $\{0,1\}$. Assume the ceilings for
$p^{(r-1)}$ required in Notation~\ref{9.5:not-hypotheses-comparison}. Assume also the
single-lattice containment \eqref{9.5:eq-majorant-containment-a}. Then every hypothesis
listed in Notation~\ref{9.5:not-hypotheses-comparison} holds at the link of row $r$ over
$p^{(r-1)}$, in the form in which part~(b) of the lettered comparison theorem uses it.
\end{lemma}

\begin{proof}
\emph{Clause (ii).} For an admissible set $\mathfrak{S}$, the $n$-th set of the shell
$\boldsymbol{\Omega}_{\mathfrak{S}}(Z)$ is $\Omega_n^{\mathfrak{S}}\cap Z^{-[s_n]}$, where
$\Omega_n^{\mathfrak{S}}$ is the $n$-th set of $\mathfrak{S}$ and $Z^{-[s_n]}$ is the box with
which it is intersected in the construction of the shells.

Let $r'<r$. The chain passes from $\mathfrak{S}_{r'}$ to $\mathfrak{S}_r$ by refinement. Hence
$\Omega_n^{\mathfrak{S}_{r'}}\supseteq\Omega_n^{\mathfrak{S}_r}$ for every $n$. Intersecting
both sides with $Z^{-[s_n]}$ gives, for every $n$,
\begin{equation*}
\Omega_n^{\mathfrak{S}_{r'}}\cap Z^{-[s_n]}\supseteq\Omega_n^{\mathfrak{S}_r}\cap Z^{-[s_n]} .
\end{equation*}
Thus every set of the shell of $Z$ over $\mathfrak{S}_{r'}$ contains the set of the same index
of the shell of $Z$ over $\mathfrak{S}_r$. The level function of a shell is defined through
these sets. Therefore the levels satisfy the inequality displayed in~(ii) at every point
of $Z$.

The sets of the earlier response rows are exactly these shells. So condition
\eqref{9.5:eq-total-variation-a} for $p^{(r-1)}$ on $Z$ follows from this inequality.

The entry's base set $\mathfrak{S}_e$ precedes all $\mathfrak{S}_{r'}$ in the chain. The same
argument with $\mathfrak{S}_e$ in place of $\mathfrak{S}_{r'}$ therefore gives the level
condition for entries of Definition~\ref{9.5:def-entries}.

\emph{Clauses (i), (iii), (iv).} These are proved as the corresponding clauses of
Lemma~\ref{9.5:lem-single-step}. In place of the range of the bracket letter and of the strip
letters of the $(T)$-family, we use the ranges of the parameters and of the strip letters
$t_{\square,r}$ of the $(L)$-family.

For~(i): a seed restricted to $Z$ is a diagonal pair satisfying condition (e1) of
Definition~\ref{9.5:def-entries}, and a diagonal pair satisfying (e1) is an entry. If $q_0$ is
instead a root of the second alternative for the initial pair of the $(L)$-family, it is the
diagonal link $q_{i_*}$ supplied by the corollary on entries along the chain, and the same
reasoning applies to it.

For~(iii): we apply Lemma~\ref{9.5:lem-cutoff-majorant} to $\chi^{(r)}_t$, with $c_0^\chi$ and
$\mathfrak{r}_0^\chi=r^0_L/C_\chi$ as stated. The plain letters of row $r$ obey the bound
$c_0^\chi$. That the strip letters, which obey $|t_{\square,r}|\le r^{L}_{\square}/C_\chi$,
satisfy the strip hypothesis of that lemma with $\mathfrak{r}_0^\chi=r^0_L/C_\chi$ is seen
exactly as in the proof of Lemma~\ref{9.5:lem-single-step}\,(iii), with the letter radii
$r^{L}_\square$ and $r^0_L$ of the $(L)$-family in place of $r'_\square$ and $r^0_T$.

For~(iv): as for the increment $\chi_t\mathbb{H}_Z$ in Lemma~\ref{9.5:lem-single-step}\,(iv),
the increment of a response row is the cut-off $\chi^{(r)}_t$ times the response field of the
box $Z$, which is supported in $Z$; and the cut-off
$\chi^{(r)}_t=\sum_\square c_\square\zeta_\square$ of Lemma~\ref{9.5:lem-cutoff-majorant}
vanishes off $\bigcup\{\tilde{\square}:\ c_\square\neq0\}$. The inclusion of the reading cells
of $v_D$ in $K_D^{\circ}\subset Z^{=}$ is the block geometry \eqref{9.5:eq-hypotheses-geometry-b}.

\emph{Predecessors of object type with letters in $\{0,1\}$.} Let a predecessor of row $r$ be
of object type, with letters in $\{0,1\}$. Its increment has the form $\chi\mathcal{Z}$, where
$\mathcal{Z}$ is the displacement of that row and $\chi$ is its cut-off. We check three things.

First, $\chi\mathcal{Z}\in\operatorname{Inc}_\beta$. The cut-off $\chi$ has the majorant
$C_\chi\max\{c_0^\chi,1\}$. This is Lemma~\ref{9.5:lem-cutoff-majorant} with the strip bound
$\mathfrak{r}_0^\chi e^{\delta_1 d(\tilde{\square},\Omega^{\mathsf{s}})/8}$ replaced by $1$,
which is legitimate because every letter lies in $\{0,1\}$. The membership of $\mathcal{Z}$ in
$\operatorname{Inc}_\beta$ is the typing of the links of the chain given in the corollary on
entries along the chain.

Second, the slot of this link is assigned by the slot rule of the lettered families. The link
is therefore a link of a presented pair in the sense of the definition of presented pairs and
their total variation. That
definition imposes no restriction of shape on the increment fields $\mathfrak{h}_s$ of the
links: they may be dilated, cut off and glued. Hence inserting this link leaves the
presentation structure of $p^{(r-1)}$ intact.

Third, the contribution of this link to the total variation $\mathsf{T}_{p^{(r-1)}}$ is
bounded in clause~(iii) of the bound for object-type rows.

\emph{The final assertion.} Clauses (i) and (ii), together with the preceding paragraph, show
that $p^{(r-1)}$ read on $Z$ is a presented pair starting at an entry and satisfying the level
conditions. Clause (iii) gives the slow twice-differentiable majorant of the cut-off of row
$r$. Clause (iv) gives the support of its increment and the position of the cells read by
$v_D$. The ceilings for $p^{(r-1)}$ and the containment \eqref{9.5:eq-majorant-containment-a}
are assumed. The latter is the content of Lemma~\ref{9.5:lem-majorant-containment}, whose
proof is outstanding. Together these are the hypotheses of
Notation~\ref{9.5:not-hypotheses-comparison} at the link of row $r$ over $p^{(r-1)}$.
\end{proof}

\begin{theorem}[Comparison bounds for lettered links. Stated; proof incomplete at the step marked $(\ast)$]\label{9.5:thm-lettered-comparison}
Assume the following statements:
\begin{itemize}
\item[\textup{(i)}] the one-link comparison theorem, clauses (A)--(D), under the convention stated
with it among the standing hypotheses in Notation~\ref{9.5:not-hypotheses-comparison}, together with
its form for a majorant profile $\bar{\tau}$ slow of type $(C_\chi,1/8)$, with the same
constants except that its constant $C$ (of clauses (C) and (D)) is multiplied by $C_\chi$;
\item[\textup{(ii)}] parts (d) and (e) of the tilted naturality lemma, as stated in
Notation~\ref{9.5:not-hypotheses-comparison};
\item[\textup{(iii)}] parts (a)--(c) of the entry theorem, its corollary on entries, and the level
condition at entries, as stated in Notation~\ref{9.5:not-hypotheses-comparison};
\item[\textup{(iv)}] the vertex identity, as stated in Notation~\ref{9.5:not-hypotheses-comparison};
\item[\textup{(v)}] the inequalities on levels and scaled path metrics of an admissible set and of its
refinements (in particular of a box set $\boldsymbol{\Omega}(Z)$ and of an entry's own set), the
monotonicity $d_{\mathfrak{S}}\le d_{\mathfrak{S}^{\sharp}}$ of the metric under passage to a
refinement $\mathfrak{S}^{\sharp}$ of an admissible set $\mathfrak{S}$, the weight inequality and the
kernel-summation bound, as stated in Notation~\ref{9.5:not-hypotheses-geometry};
\item[\textup{(vi)}] the second-order smoothness bounds and the scale-$M$ bounds for the partition of
unity $\{\zeta_\square\}$ of \cite{B-CMP109}, and the block geometry of the fat core $K$, its interior
$K^{\circ}$ and the cubes $\sigma_0(K)$, as stated in Notation~\ref{9.5:not-hypotheses-geometry};
\item[\textup{(vii)}] the move format of the \textup{(P)}-family in
Definition~\ref{9.5:def-lettered-links}, with its core letter $\mathsf{d}^\theta(e)$.
\end{itemize}
Then there are constants $C_{\mathsf{m}}$, $C_{R'}$, $C$ and ceilings $\mathfrak{r}_0$,
$\varepsilon_D$, $a_{R'}$, depending on $(d,L,\theta,M_1)$ only (those of the one-link comparison
theorem in \textup{(i)}, with $C$ enlarged by the factor $C_\chi^2$), and free of the number of
links, of $k$, $K$, $N$, $N_0$, $\check{R}$, of the ages of large-field regions, of the volume and
of $g$, of the box $Z$, of $|D|$, of $|\mathcal{D}^{+}_c|$, of the number of dilated blocks and of
the lengths of the corridors, such that the following holds at each site type.

Consider a term over strips $c$, a set $D\subset\mathcal{D}^{+}_c$, a box $Z$ of the families of
Definition~\ref{9.5:def-lettered-links}, and a lettered link $(Z,\chi,b)$ of that definition with
real letters in the ranges displayed there and real datum $b$ with
$\sup(|b|+|\delta b|)\le a_{R'}$. Let $p$ be a real parent of the family on a region
$\hat{X}\supseteq Z$ with $\mathfrak{r}_{P'},\mathfrak{r}_P\le\mathfrak{r}_0$ and
$\sup\mathsf{T}_p\le\varepsilon_D$, the supremum taken over the region carrying $\mathsf{T}_p$.
Assume that the lettered majorant $\mathsf{H}_\chi:=C_\chi\bar{\tau}_\chi\mathsf{P}_b$, with
$C_\chi$ the constant of Lemma~\ref{9.5:lem-cutoff-majorant} and $\bar{\tau}_\chi$ as in
\eqref{9.5:eq-letter-profile-a}, satisfies the single-lattice containment
\eqref{9.5:eq-majorant-containment-a} of Lemma~\ref{9.5:lem-majorant-containment}, whose proof is
outstanding. Read all quantities in the units of the current sets
$(\boldsymbol{\Omega}'(Z),\boldsymbol{\Omega}(Z))$, that is, with the levels $\lambda_Z$, the scales
$\ell_Z$, and the cubes and metric of $\boldsymbol{\Omega}(Z)$, and put
\begin{equation*}
\mathbb{H}:=\mathbb{H}_Z(b;P),\qquad
\mathbb{H}':=\mathbb{H}'_Z\bigl((b+\delta b)\oplus 0;P'\bigr).
\end{equation*}
Then:
\begin{itemize}
\item[\textup{(A)}] at every paired data site $y$ of $Z$,
\begin{equation*}
|\mathsf{m}_{P'}(y)-\mathsf{m}_P(y)|\le C_{\mathsf{m}}\Psi_p(y);
\end{equation*}
\item[\textup{(B)}] at every $x\in Z$, with $Y[\,\cdot\,;\mathbf{U}']$ as in clause (B) of the
one-link comparison theorem of Notation~\ref{9.5:not-hypotheses-comparison},
\begin{equation*}
\bigl|Y[\mathbb{H}';\mathbf{U}']-\mathbb{H}\bigr|_x
\le C_{R'}\bigl\{\mathsf{P}_b(x)\Psi^B_p(x)+\mathsf{P}_{\delta b}(x)\bigr\};
\end{equation*}
\item[\textup{(C)}] the extended pair $p^\chi$, that is, $p$ extended by $(\chi\mathbb{H}',\chi\mathbb{H})$,
obeys at every $x\in Z$, with $\mathsf{T}$ and $A$ as in \eqref{9.5:eq-total-variation-b} and
\eqref{9.5:eq-mismatch-field-b},
\begin{equation}\label{9.5:eq-lettered-comparison-a}
\begin{split}
\mathsf{T}_{p^\chi}(x)-\mathsf{T}_p(x)
&=|A_{p^\chi}-A_p|_x\\
&\le 2C_{R'}\bigl\{\mathsf{H}_\chi\Psi^B_p+C_\chi\bar{\tau}_\chi\mathsf{P}_{\delta b}\bigr\}(x)
+C\,L^{-\lambda_Z(x)}\mathsf{H}_\chi(x)\bigl[M^{-1}+\mathsf{H}_\chi+\mathsf{T}_p\bigr](x);
\end{split}
\end{equation}
\item[\textup{(D)}] the pair $p^\chi$ is presented: in the units of $\boldsymbol{\Omega}(Z)$, with
$C_s$ the constant of clause (D) of the one-link comparison theorem of
Notation~\ref{9.5:not-hypotheses-comparison},
\begin{equation}\label{9.5:eq-lettered-comparison-b}
\chi\mathbb{H}\in\operatorname{Inc}_\beta\bigl(4C_sC_\chi\mathsf{H}_\chi\bigr),\qquad
\chi\mathbb{H}'\in\operatorname{Inc}_\beta\bigl(4C_sC_\chi\mathsf{H}_\chi\bigr),
\end{equation}
and the class radii grow by $4CC_sC_\chi\sup\mathsf{H}_\chi$.
\end{itemize}
These assertions are made in the following cases. Clause \textup{(P)}: for the links of
$\mathfrak{F}^{\mathbb{R},\mathbb{Y}}$ over its points (the cut-off being $\chi=\tau_t$). Clause
\textup{(T)}: for the one link over $p_0$. Clause \textup{(L)}: for rows of response type over the
partial products $p^{(r-1)}$ whose earlier rows are of response type, or of object type with strip
letters in $\{0,1\}$. For rows of object type at dilated strip letters nothing is claimed.
\end{theorem}

The step marked $(\ast)$ concerns clause \textup{(P)} for runs entered at an \textup{(L)}-chain
containing a row of object type at dilated strip letters; for clauses \textup{(T)} and \textup{(L)},
and for clause \textup{(P)} at runs entered at a seed or at an \textup{(L)}-chain whose rows of
object type carry strip letters in $\{0,1\}$ only, the proof below derives the conclusions from the
hypotheses listed in the theorem.

\begin{proof}
Fix a site type, and fix a link of Definition~\ref{9.5:def-lettered-links} over a parent $p$ of the
family satisfying the ceilings $\mathfrak{r}_{P'},\mathfrak{r}_P\le\mathfrak{r}_0$ and
$\sup\mathsf{T}_p\le\varepsilon_D$. We verify, one by one, the hypotheses of the one-link comparison
theorem of Notation~\ref{9.5:not-hypotheses-comparison} for this link, read on $Z$.

\emph{The parent is presented over an entry.} At \textup{(P)} this is the statement of
Lemma~\ref{9.5:lem-run-points}, at \textup{(T)} that of Lemma~\ref{9.5:lem-single-step}, and at
\textup{(L)} that of Lemma~\ref{9.5:lem-response-rows}: in each case the parent read on $Z$ is a
presented pair over an entry, and the level condition required of the parent in
Notation~\ref{9.5:not-hypotheses-comparison} and the level condition at entries of hypothesis
\textup{(iii)} hold in
exactly the form required by the one-link comparison theorem. Where a run of the \textup{(P)}-family
is entered at an \textup{(L)}-chain containing a row of object type at dilated strip letters, the
presentation of its entry is not supplied by the statements listed in the theorem: it rests on the
relative bound \eqref{9.5:eq-relative-object-bound-a} for such rows, whose proof is outstanding; this
is the step $(\ast)$ at which the proof is incomplete.

\emph{The current sets.} By the vertex identity, together with the standing hypotheses on boxes of
Notation~\ref{9.5:not-hypotheses-comparison}, the increment of the unit at the vertex is
$\chi\mathbb{H}_Z$ on all of $Z$; hence the current sets are
$(\boldsymbol{\Omega}'(Z),\boldsymbol{\Omega}(Z))$. By the same standing hypotheses these two sets
are exactly paired admissible sets, with constants free of $Z$.

\emph{The cut-off.} The cut-off $\chi$ is real, since the letters and the parameters are real, and it
is the same function in both machines, since the parameters are equal in both. Put
$\bar{\tau}:=C_\chi\bar{\tau}_\chi$. Parts (a) and (b) of
Lemma~\ref{9.5:lem-cutoff-majorant}, applied in the units of $\boldsymbol{\Omega}(Z)$ (an admissible
set refining the label), give
\begin{equation*}
|\chi|\le\bar{\tau},\qquad
\ell_Z|\nabla\chi|\le\bar{\tau}/M,\qquad
\ell_Z^2|\nabla^2\chi|\le\bar{\tau}/M^2;
\end{equation*}
the hypotheses of Lemma~\ref{9.5:lem-cutoff-majorant} on the letters hold because the letters of
the link lie in the ranges displayed in
Definition~\ref{9.5:def-lettered-links} (at \textup{(T)} this is part (iii) of
Lemma~\ref{9.5:lem-single-step}).
By part (c) of Lemma~\ref{9.5:lem-cutoff-majorant}, $\bar{\tau}$ is slow of type
$(C_\chi\cdot 1,1/8)$. We therefore use the form of the one-link comparison theorem for profiles
slow of type $(C_\chi,1/8)$ listed in hypothesis \textup{(i)}, whose constant $C$ carries the factor
$C_\chi$; this is accounted for by the enlargement of $C$ to $CC_\chi^2$ in the statement.

\emph{The majorant.} The one-link comparison theorem's majorant is $\mathsf{H}:=\bar{\tau}\mathsf{P}_b$,
which by the choice of $\bar{\tau}$ equals $\mathsf{H}_\chi=C_\chi\bar{\tau}_\chi\mathsf{P}_b$. By the
containment \eqref{9.5:eq-majorant-containment-a}, which is a hypothesis of the theorem (it is the
content of Lemma~\ref{9.5:lem-majorant-containment}, whose proof is outstanding),
$\mathsf{H}=\mathsf{H}_\chi\le a_{R'}$.

\emph{Conclusion.} By part (b) of the entry theorem, the presented parent over its entry, with the
data just verified, meets the hypotheses of the one-link comparison theorem, and that theorem, under
the convention stated with it in Notation~\ref{9.5:not-hypotheses-comparison}, gives its
clauses (A)--(D) with its own constants. These constants are free of the number $r$ of links and of
the box $Z$. The letters of the link enter the one-link comparison theorem only through $\bar{\tau}$;
hence, by part (d) of Lemma~\ref{9.5:lem-cutoff-majorant}, the constants are also free of $|D|$, of
the number of non-zero letters and of the lengths of the corridors. Written out with
$\bar{\tau}=C_\chi\bar{\tau}_\chi$, so that $\mathsf{H}=\mathsf{H}_\chi$, the clauses (A)--(D) of
the one-link comparison theorem, in the form of hypothesis \textup{(i)}, yield the clauses (A) and
(B), the bound \eqref{9.5:eq-lettered-comparison-a} and the presentation
\eqref{9.5:eq-lettered-comparison-b} of the theorem, after enlarging the constants by factors
$C_\chi\ge 1$, the class $\operatorname{Inc}_\beta(\mathsf{r})$ being increasing in the profile
$\mathsf{r}$.
\end{proof}

\begin{remark}[Hypotheses met verbatim and hypotheses never read]
\label{9.5:rem-verbatim-hypotheses}
Six properties of the boxes, of the fat core $K$ and of the cut-offs could enter as hypotheses
of the comparison bounds for lettered links (Theorem~\ref{9.5:thm-lettered-comparison}). Their
standing is as follows at the three site types of that theorem: the points of a run of lettered
links (Lemma~\ref{9.5:lem-run-points}), the single lettered step
(Lemma~\ref{9.5:lem-single-step}) and the response-type rows
(Lemma~\ref{9.5:lem-response-rows}).
\begin{itemize}
\item[(a)] Nesting of boxes is used only through the level condition
\eqref{9.5:eq-total-variation-a}; at the points of a run it holds in the form
$Z_s\subset K_{s-1}^{\circ}$, and at the other two site types there is nothing to nest.
\item[(b)] The containments $K\supseteq K_A$ and $K\supseteq X^{+}$, with $K_A$ and $X^{+}$
as fixed in the setting of Theorem~\ref{9.5:thm-lettered-comparison}, are used by no step of
the proof of Theorem~\ref{9.5:thm-lettered-comparison}.
\item[(c)] At the points of a run the cut-off $\tau_s$ is supported in $K_s^{\circ}$; at the
single lettered step and at the response-type rows the increment $\chi\mathbb{H}_Z$ is in
general not supported in $K_D^{\circ}$; the conclusions
of Theorem~\ref{9.5:thm-lettered-comparison} hold at every $x\in Z$ nonetheless.
\item[(d)] The level condition \eqref{9.5:eq-total-variation-a} holds verbatim at all three
site types.
\item[(e)] The majorants $\bar{\tau}_\chi$ and $\Psi_p$ are slow, and the slowness of
$\bar{\tau}_\chi$ imposes no condition beyond those already present.
\item[(f)] The exponential smearing $f^{\sim}$ is taken over the whole region; no restriction
to $X^{+}$ is made.
\end{itemize}
\end{remark}

\begin{proof}
We treat the six properties in turn.

(a) The nesting of boxes enters the proof only through the level condition
\eqref{9.5:eq-total-variation-a}, which is read only at clauses (a) and (b) of the comparison
lemma and at clause (ii) of the pair lemma listed among the inputs in
Remark~\ref{9.5:rem-comparison-inputs}. At the points of a run it holds in the form
\begin{equation*}
Z_s\subset K_{s-1}^{\circ}\qquad\text{for every } s,
\end{equation*}
which is Lemma~\ref{9.5:lem-run-points}(i) and follows from the conditions that define a run
presentation in the real lettered family $\mathfrak{F}^{\mathbb{R},\mathbb{Y}}$. At the single
lettered step there is no earlier link, so no pair of boxes has to be nested. At the
response-type rows a single box $Z$ occurs.

(b) No step of the proof of Theorem~\ref{9.5:thm-lettered-comparison}, as listed in
Remark~\ref{9.5:rem-comparison-inputs}, uses $K\supseteq K_A$ or $K\supseteq X^{+}$. Hence
the facts that the core $K_D$ does not contain the union of the enlarged blocks
$\tilde{\square}^{4}$ over the blocks carrying plain letters, and that $K_D$ does not contain
$X^{+}$, do not affect that theorem. They affect only two things: the index produced by the
conversion of the comparison bounds into the target law, and the reading of the strip as a
spectator in the entry law.

(c) At the points of a run, Lemma~\ref{9.5:lem-run-points}(ii) gives
$\operatorname{supp}\tau_s\subset K_s^{\circ}$. At the single lettered step and at the
response-type rows the increment $\chi\mathbb{H}_Z$ is in general not supported in
$K_D^{\circ}$: the support lies in $Z$ and in general meets $Z\setminus Z^{=}$, by
Lemma~\ref{9.5:lem-single-step}(iv) and Lemma~\ref{9.5:lem-response-rows}(iv). The proof of
Theorem~\ref{9.5:thm-lettered-comparison} does not use this support property, so its
conclusions hold at every $x\in Z$ in all three cases. The price of the missing support
property appears in the pointwise form of the third conclusion of that theorem at points of
$Z\setminus Z^{=}$, which is read in the conversion of the comparison bounds into the target
law.

(d) The level condition \eqref{9.5:eq-total-variation-a} holds verbatim at all three site
types. At the points of a run this is Lemma~\ref{9.5:lem-run-points}(iii), at the single
lettered step Lemma~\ref{9.5:lem-single-step}(ii), and at the response-type rows
Lemma~\ref{9.5:lem-response-rows}(ii).

(e) The letter profile $\bar{\tau}_\chi$ is slow by Lemma~\ref{9.5:lem-cutoff-majorant}(c).
On the strip blocks the dilated strip letters are allowed to grow like
\begin{equation*}
e^{\delta_1 d(\cdot,\,\Omega^{\mathsf{s}})/8}.
\end{equation*}
This is the rate $\delta_1/8$ in $d(\cdot,\Omega^{\mathsf{s}})$ at which the pair radii grow
in the hypotheses of the one-lattice response bounds underlying
Theorem~\ref{9.5:thm-lettered-comparison}. The slowness of $\bar{\tau}_\chi$ therefore adds no
condition. The majorant $\Psi_p$ is slow by clause (i) of the definition of the slow majorants
$\Psi_p$ together with condition (d1) of that setting.

(f) The first three conclusions of Theorem~\ref{9.5:thm-lettered-comparison} are used in what
follows at every bond of $Z$. The laws derived from them are stated at every bond of the
regions that carry the total variation $\mathsf{T}_p$. Neither the proof nor these
applications restrict the smearing $f^{\sim}$ to $X^{+}$.
\end{proof}

\begin{lemma}[The boundary price off the full-level region]\label{9.5:lem-boundary-price}
Let $(Z,\chi,b)$ be a lettered link of the family of site type $(\mathrm{T})$ or of the family
of site type $(\mathrm{L})$, and let $x$ be a bond of $Z$. Write $\mathcal{B}(x)$ for the right
member of the comparison bound \eqref{9.5:eq-lettered-comparison-a} at the bond $x$, and
$p$, $p^{\chi}$ for the two points compared there. Then
\begin{equation}\label{9.5:eq-boundary-price-a}
\begin{aligned}
\sup_{\tilde{\Delta}_Z(x)} \bigl|A_{p^{\chi}}-A_p\bigr|
&\le \ell_Z(x)^{-1}\,\mathcal{B}(x)\\
&\le L^{\operatorname{def}_Z(x)}\,\ell_\iota(x)^{-1}\,\mathcal{B}(x).
\end{aligned}
\end{equation}
On the full-level region $Z^{=}$, where $\operatorname{def}_Z=0$, this is the bound in the
label's units $\ell_\iota(x)$ and with the label's index. Off $Z^{=}$ the pointwise bound
\eqref{9.5:eq-boundary-price-a} is weaker than the bound in the label's units by the factor
$L^{\operatorname{def}_Z(x)}$, and
the index in $\mathcal{B}(x)$ is
\begin{equation}\label{9.5:eq-boundary-price-1}
L^{-\theta\lambda_Z(x)}=L^{-\theta\iota(x)}\,L^{\theta\operatorname{def}_Z(x)};
\end{equation}
altogether, against the form with the label's index $L^{-\theta\iota(x)}$ and the label's
units $\ell_\iota(x)$, the bound \eqref{9.5:eq-boundary-price-a} carries the factor
$L^{(1+\theta)\operatorname{def}_Z(x)}$.
\end{lemma}

\begin{proof}
The first inequality of \eqref{9.5:eq-boundary-price-a} is the comparison bound
\eqref{9.5:eq-lettered-comparison-a} of Theorem~\ref{9.5:thm-lettered-comparison}, read at the
bond $x$ of $Z$.

For the second inequality we use the relation between the two local scales at $x$. The local
scale over the box is $\ell_Z(x)=L^{\lambda_Z(x)}\eta$ (Definition~\ref{9.5:def-admissible-sets}),
the label's scale is $\ell_\iota(x)=L^{\iota(x)}\eta$, and the deficit of the box is the
difference of the two levels, so that
\begin{equation}\label{9.5:eq-boundary-price-2}
\ell_Z(x)=\ell_\iota(x)\,L^{-\operatorname{def}_Z(x)},
\qquad
\lambda_Z(x)=\iota(x)-\operatorname{def}_Z(x).
\end{equation}
Hence $\ell_Z(x)^{-1}=L^{\operatorname{def}_Z(x)}\ell_\iota(x)^{-1}$, and multiplying this
identity by $\mathcal{B}(x)$ gives the second inequality of \eqref{9.5:eq-boundary-price-a}.

The second relation in \eqref{9.5:eq-boundary-price-2}, multiplied by $-\theta$ and
exponentiated in base $L$, is \eqref{9.5:eq-boundary-price-1}. On $Z^{=}$ one has
$\operatorname{def}_Z(x)=0$, so $\ell_Z(x)=\ell_\iota(x)$ and $\lambda_Z(x)=\iota(x)$, and
\eqref{9.5:eq-boundary-price-a} is the bound in the label's units with the label's index.
Off $Z^{=}$ the factor of units is $L^{\operatorname{def}_Z(x)}$ by the second inequality of
\eqref{9.5:eq-boundary-price-a}, the factor in the index is $L^{\theta\operatorname{def}_Z(x)}$ by
\eqref{9.5:eq-boundary-price-1}, and their product is
$L^{(1+\theta)\operatorname{def}_Z(x)}$.
\end{proof}

The factor $L^{(1+\theta)\operatorname{def}_Z(x)}$ of Lemma~\ref{9.5:lem-boundary-price} is
compensated by none of the inputs of Theorem~\ref{9.5:thm-lettered-comparison} at bonds of
$Z\setminus Z^{=}$ near the sources of the
site. Indeed, Lemma~\ref{9.5:lem-majorant-containment}, whose proof is outstanding, bounds the
lettered majorant $\mathsf{H}_\chi$ by $\tfrac14\ell_\iota a$, that is, in the label's units; and
the decay of the smeared profile $\mathsf{P}_b$ comes from the site's source region
$\Omega^{\mathsf{s}}$, which at site types $(\mathrm{T})$ and $(\mathrm{L})$ the box meets through
its toggled cubes, which we denote by $Y'$. At these site types, by the vanishing clause of the
statement on the pieces of the $(\mathrm{T})$- and $(\mathrm{L})$-families, a piece vanishes
unless every wall-component of the set $\bigcup_D\tilde{\square}^{4}\cup\mathsf{y}$ formed in
that statement holds a source cell; so
$\operatorname{def}_Z$ may be large where $\mathsf{P}_b$ is not small. At site type
$(\mathrm{P})$ the question does not arise: by Lemma~\ref{9.5:lem-run-points}(ii) the
$\mathsf{T}$-increment vanishes off $K_s^{\circ}\subset Z_s^{=}$, and the conversion of rates at
points off $K^{\circ}$ that enters Lemma~\ref{9.5:lem-majorant-containment} (proof outstanding)
at site type $(\mathrm{P})$ concerns the majorant only.

For increments supported off $K^{\circ}$, the form of the comparison with the label's index and
the label's units at a bond $x\notin Z^{=}$ would require
\begin{equation}\label{9.5:eq-boundary-price-b}
\mathsf{H}_\chi(x)\,L^{(1+\theta)\operatorname{def}_Z(x)}
\le \tfrac14\,\ell_\iota(x)\,a(x)\cdot\mathrm{const},
\end{equation}
where $\mathrm{const}$ denotes a number; that is, a decay of the lettered majorant that beats
the deficit of the box at $x$. Inequality \eqref{9.5:eq-boundary-price-b} is not proved here
and is not used below. Where $\mathsf{P}_b$ decays from sources inside $K^{\circ}$, it would
follow from the conversion of rates that enters Lemma~\ref{9.5:lem-majorant-containment}
(proof outstanding) at site type $(\mathrm{P})$; near sources lying in $Y'$ it fails in form.
It is not needed in what follows: every use of the pointwise bound
\eqref{9.5:eq-boundary-price-a} is at bonds of $Z^{=}$, since the cells at which $v_D$ reads
increments lie in $Z^{=}$ by \eqref{9.5:eq-hypotheses-geometry-b}; and the role of the
increments supported up to
the boundary of $Z$ as increments of the parent is carried in terms of the total variation
$\mathsf{T}$ by the lemma converting the pointwise bound into a bound on $\mathsf{T}$ with the
box-capped index $\iota^e_Z$, in which that index absorbs the deficit.

\begin{proposition}[Comparison bounds for weakly presented pairs. Stated; proof incomplete at the step marked $(\ast)$]\label{9.5:prop-weak-presentation}
Let $Z$ be a box, and let $\lambda$ be as in the level condition \eqref{9.5:eq-total-variation-a}. Let $p_0\to\dots\to p_r$ be a presentation, not necessarily in $\mathfrak{F}^{\mathbb{R},\mathbb{Y}}$, with links $t=1,\dots,r$, the link $t$ having admissible set $\mathfrak{S}_t$, box $Z_t$, increment field $\mathfrak{h}_t$ with profile $\mathsf{r}_t$, and the stencil-fattened support $S_t^{+}$ of $\mathfrak{h}_t$ (compare Lemma~\ref{9.5:lem-run-points}\textup{(ii)}). Write $\mathbf{U}_{t-1}$ for the configuration of the point $p_{t-1}$.

A pair with presentation $p_0\to\dots\to p_r$ over the box $Z$ is called \emph{weakly presented} if it satisfies every requirement of a presented pair in the sense of the definition containing \eqref{9.5:eq-total-variation-a}, except the level condition \eqref{9.5:eq-total-variation-a} itself. That condition is replaced by the following two conditions. First, for every $t$,
\begin{equation}\label{9.5:eq-weak-presentation-1}
\lambda_{\mathfrak{S}_t}\ge\lambda\qquad\text{on } S_t^{+}\cap Z .
\end{equation}
Second, for every $t$, with a constant $C_J$, the tail bound
\begin{equation}\label{9.5:eq-weak-presentation-2}
\ell(y)^3\,\bigl|\mathcal{J}_{\mathfrak{S}_t}(\mathfrak{h}_t;\mathbf{U}_{t-1})(y)\bigr|
\le C_J\,\bigl(\mathsf{r}_t\mathbf{1}_{S_t^{+}}\bigr)^{\sim,5/8}(y),
\qquad y\in Z ,
\end{equation}
holds. Then clauses \textup{(A)}--\textup{(D)} of the comparison bounds for presented pairs hold for weakly presented pairs, with the same constants up to absolute factors.
\end{proposition}

Stated; the proof below is incomplete at the step marked $(\ast)$.

\begin{proof}
We first say why the weak form is needed. Suppose the link $t$ of a presentation has a box $Z_t$ with $Z\not\subset Z_t$. Then the level condition \eqref{9.5:eq-total-variation-a} fails on $Z\setminus Z_t$, because there the shells of $Z_t$ carry no level. Consequently the comparison bounds for presented pairs, as stated, do not apply to such a presentation.

We now list the places in the proof of clauses \textup{(A)}--\textup{(D)} where the level function $\lambda_{\mathfrak{S}_t}$ of link $t$ is actually read (compare Remark~\ref{9.5:rem-comparison-inputs}). There are three of them.
\begin{enumerate}
\item The step that bounds the members on the cubes $\tilde{\Delta}(y)$ reads $\lambda_{\mathfrak{S}_t}$ only on those cubes $\tilde{\Delta}(y)$ that meet $\operatorname{supp}\mathfrak{h}_t$. Elsewhere the members bounded in that step vanish.
\item The step that bounds the slot term reads $\lambda_{\mathfrak{S}_t}$ through
\[
\ell^3\,\bigl|\mathcal{J}_{\mathfrak{S}_t}(\mathfrak{h}_t)\bigr| .
\]
Its part $D R_{\mathfrak{S}_t} D^{*}\mathfrak{h}_t$ has exponentially decaying tails over $Z_t$. The tail bound \eqref{9.5:eq-weak-presentation-2} is the hypothesis on $Z$ that this step is to use in the weak form.
\item The pairing step at stage $t$ reads $\lambda_{\mathfrak{S}_t}$ in two ways. First, it reads it through three defect bounds at stage $t$, which we call the first, the second and the third defect bound, in the order in which they are listed in the proof of the comparison bounds for presented pairs. Each defect is paired against $D^{*}\mathfrak{h}_t$, against $e_t$ or against $D^{*}e_t$. Each of these three fields is supported on the stencil-fattened support $S_t^{+}$ of $\mathfrak{h}_t$. Together with the first step, this reading is what suggests condition \eqref{9.5:eq-weak-presentation-1}, which asks for the level comparison on $S_t^{+}\cap Z$ only. Second, it reads it through the Lipschitz bound for Ba{\l}aban's projection $R_{\mathfrak{S}_t}$, which reads $\sup\lvert A_{p_{t-1}}\rvert$ on $Z_t$.
\end{enumerate}
The weak form affects only the pairing step and the two bounding steps above, the one on the cubes $\tilde{\Delta}(y)$ and the one on the slot term.

$(\ast)$ The pairing step and the first defect bound entering it are stated with smeared majorants $\mathsf{r}_t^{\sim}$ and with the index $L^{-\beta\lambda}$. The bookkeeping of scales (which field is measured in which length unit) is not displayed at points where $\lambda_{\mathfrak{S}_t}\ne\lambda$. At such points the test field is regular at the current scale, while the increment is regular at the scale of its link $t$. Passing to weakly presented pairs requires two further steps:
\begin{itemize}
\item the bookkeeping of scales at these points;
\item the conversion of the factor $L^{-\beta\lambda_{\mathfrak{S}_t}(y')}$ at a point $y'\in S_t^{+}$ into the factor $L^{-\beta\lambda_Z(c)}$ at a data site $c$, where $\lambda_Z$ is the level function of the box $Z$.
\end{itemize}
The conversion requires a comparison of the scaled path metrics $d_{Z_t}$ and $d_Z$. These two metrics are not comparable in general, and the proof stops at this step.

The proposition is not needed for the comparison bounds for lettered links (Theorem~\ref{9.5:thm-lettered-comparison}). For the presentations of the real lettered family $\mathfrak{F}^{\mathbb{R},\mathbb{Y}}$ that enter that theorem, the level condition \eqref{9.5:eq-total-variation-a} holds verbatim, site type by site type:
\begin{itemize}
\item at sites of type \textup{(P)}, the defining conditions of the family force the nesting of Lemma~\ref{9.5:lem-run-points}\textup{(i)};
\item at sites of type \textup{(L)}, there is a single box;
\item at sites of type \textup{(T)}, there is no earlier link.
\end{itemize}
\end{proof}

\begin{lemma}[Slowness of the box-capped index]\label{9.5:lem-capped-index}
Let $Z$ be a box, in the sense of Definition~\ref{9.5:def-admissible-sets}, over an admissible
set $\mathfrak{S}$ refining the label and, at a site of type \textup{(L)}, the entry's set
$\mathfrak{S}_e$. We write $\iota$ and $d$ for the level function and the scaled path metric of
the label, $\lambda_e$ and $d_{\mathfrak{S}_e}$ for those of $\mathfrak{S}_e$, and $\lambda_Z$
and $d_Z$ for those of the admissible set $\boldsymbol{\Omega}(Z)$ of
Definition~\ref{9.5:def-admissible-sets}; $Z^{=}$ and $\operatorname{def}_Z$ denote the full-level
region and the deficit function of $Z$ defined there. Let
$v=\bar{u}^{\mathrm{loc}}$ be the site's local unit, and let
\begin{equation*}
  \iota^e_Z=\min\{\iota,\lambda_Z,\lambda_e\}
\end{equation*}
be the box-capped index on $Z$ of Proposition~\ref{9.5:prop-reading-decay}. With
$\vartheta=L^{-\theta_2}$ put $w:=v\,\vartheta^{\iota^e_Z}$ on $Z$. Then the following hold.
\begin{enumerate}
\item[(a)] For $x,y\in Z$,
\begin{equation}\label{9.5:eq-capped-index-1}
  |\iota^e_Z(x)-\iota^e_Z(y)|\le 1+\frac{\delta_1 d_Z(x,y)}{16\log L}.
\end{equation}
\item[(b)] For $x\in Z$,
\begin{equation}\label{9.5:eq-capped-index-2}
  \sum_{y\in Z}e^{-\delta_1 d_Z(x,y)}\,w(y)\le L^{\theta_2}C_T\,w(x),
\end{equation}
where $C_T$ is the constant of the weighted summation bound among the standing hypotheses
\eqref{9.5:eq-hypotheses-geometry-a}.
\item[(c)] $L^{-\theta\lambda_Z}\le\vartheta^{\iota^e_Z}$ on $Z$, where $\theta$ is the exponent,
not smaller than $\theta_2$, among the parameters of the two-lattice comparison
(Notation~\ref{9.5:not-comparison-parameters}).
\item[(d)] For $y\in Z$ and $x\in Z^{=}$,
\begin{equation}\label{9.5:eq-capped-index-3}
  L^{\operatorname{def}_Z(y)}\le L^{2}e^{\delta_1 d_Z(x,y)/8};
\end{equation}
that is, the deficit at $y$ is paid for by the box distance from $y$ to the full-level region.
\end{enumerate}
\end{lemma}

\begin{proof}
We use the standing hypotheses on metric and block geometry of
Notation~\ref{9.5:not-hypotheses-geometry}, displayed in \eqref{9.5:eq-hypotheses-geometry-a}.

(a) For an admissible set and for each of its refinements, in particular for the label, for the
admissible set $\boldsymbol{\Omega}(Z)$ of a box and for an entry's own set $\mathfrak{S}_e$,
the level-difference bound in \eqref{9.5:eq-hypotheses-geometry-a} states, for such a set
$\mathfrak{S}'$ with level function $\lambda_{\mathfrak{S}'}$ and scaled path metric
$d_{\mathfrak{S}'}$,
\begin{equation*}
  L^{2(|\lambda_{\mathfrak{S}'}(x)-\lambda_{\mathfrak{S}'}(y)|-1)_+}
  \le e^{\delta_1 d_{\mathfrak{S}'}(x,y)/8}.
\end{equation*}
Taking logarithms,
$2(|\lambda_{\mathfrak{S}'}(x)-\lambda_{\mathfrak{S}'}(y)|-1)_+\log L\le\delta_1
d_{\mathfrak{S}'}(x,y)/8$, hence
\begin{equation}\label{9.5:eq-capped-index-4}
  |\lambda_{\mathfrak{S}'}(x)-\lambda_{\mathfrak{S}'}(y)|
  \le 1+\frac{\delta_1 d_{\mathfrak{S}'}(x,y)}{16\log L}.
\end{equation}
Thus $\iota$ satisfies \eqref{9.5:eq-capped-index-4} with the metric $d$, $\lambda_Z$ with the
metric $d_Z$, and $\lambda_e$ with the metric $d_{\mathfrak{S}_e}$. Since
$\boldsymbol{\Omega}(Z)$ refines both the label and $\mathfrak{S}_e$, the monotonicity of the
scaled path metric under refinement in \eqref{9.5:eq-hypotheses-geometry-a} gives
$d\le d_Z$ and $d_{\mathfrak{S}_e}\le d_Z$ pointwise. Hence each of the three functions
$f\in\{\iota,\lambda_Z,\lambda_e\}$ satisfies
\begin{equation*}
  |f(x)-f(y)|\le c(x,y):=1+\frac{\delta_1 d_Z(x,y)}{16\log L},\qquad x,y\in Z.
\end{equation*}
A minimum of functions each satisfying this bound satisfies it as well: if $f_j$ attains the
minimum at $y$, then
\begin{equation*}
  \min_i f_i(x)\le f_j(x)\le f_j(y)+c(x,y)=\min_i f_i(y)+c(x,y),
\end{equation*}
and the same argument with
$x$ and $y$ exchanged gives the reverse inequality, the function $c$ being symmetric. Applied to
$\iota^e_Z=\min\{\iota,\lambda_Z,\lambda_e\}$ this is \eqref{9.5:eq-capped-index-1}.

(b) By (a), $\iota^e_Z(y)\ge\iota^e_Z(x)-1-\delta_1 d_Z(x,y)/(16\log L)$, and therefore, since
$L^{\theta_2\delta_1 d_Z(x,y)/(16\log L)}=e^{\theta_2\delta_1 d_Z(x,y)/16}$,
\begin{equation*}
  \vartheta^{\iota^e_Z(y)}=L^{-\theta_2\iota^e_Z(y)}
  \le L^{\theta_2}e^{\theta_2\delta_1 d_Z(x,y)/16}\,\vartheta^{\iota^e_Z(x)}.
\end{equation*}
Multiplying by $e^{-\delta_1 d_Z(x,y)}v(y)$ and summing over $y\in Z$ gives
\begin{align*}
  \sum_{y\in Z}e^{-\delta_1 d_Z(x,y)}w(y)
  &\le L^{\theta_2}\vartheta^{\iota^e_Z(x)}
     \sum_{y\in Z}e^{-(1-\theta_2/16)\delta_1 d_Z(x,y)}v(y)\\
  &\le L^{\theta_2}\vartheta^{\iota^e_Z(x)}\sum_{y\in Z}e^{-\delta_1 d(x,y)/2}v(y)
   \le L^{\theta_2}C_T\,w(x).
\end{align*}
In the second inequality we used $1-\theta_2/16\ge 1/2$, that is $\theta_2\le 8$, and
$d\le d_Z$, the latter by the
monotonicity of the metric under refinement as in (a). The third inequality is the weighted
summation bound of \eqref{9.5:eq-hypotheses-geometry-a}, applied to the unit function $v$,
which bounds the last sum by $C_T\,v(x)$, followed by
$v(x)\vartheta^{\iota^e_Z(x)}=w(x)$. This is \eqref{9.5:eq-capped-index-2}.

(c) On $Z$ we write
\begin{equation*}
  L^{-\theta\lambda_Z}=\vartheta^{\lambda_Z}L^{-(\theta-\theta_2)\lambda_Z}
  \le\vartheta^{\lambda_Z}\le\vartheta^{\iota^e_Z}.
\end{equation*}
The identity is $\vartheta=L^{-\theta_2}$. The first inequality is
$L^{-(\theta-\theta_2)\lambda_Z}\le 1$, that is $(\theta-\theta_2)\lambda_Z\ge 0$, which holds
since $\theta\ge\theta_2$ and $\lambda_Z\ge 0$. The second
holds because $\iota^e_Z\le\lambda_Z$ by the definition of $\iota^e_Z$ as a minimum and
$\vartheta\le 1$.

(d) Let $y\in Z$ and $x\in Z^{=}$. At the point $x$ of the full-level region
$\lambda_Z(x)=\iota(x)$, so the deficit at $y$ can be written
\begin{equation*}
  \operatorname{def}_Z(y)=[\iota(y)-\iota(x)]+[\lambda_Z(x)-\lambda_Z(y)].
\end{equation*}
By \eqref{9.5:eq-capped-index-4} for $\iota$ with the metric $d$ and for $\lambda_Z$ with the
metric $d_Z$, and then by $d\le d_Z$,
\begin{equation*}
  \operatorname{def}_Z(y)\le 2+\frac{\delta_1[d(x,y)+d_Z(x,y)]}{16\log L}
  \le 2+\frac{\delta_1 d_Z(x,y)}{8\log L}.
\end{equation*}
Raising $L$ to this power and using $L^{\delta_1 d_Z(x,y)/(8\log L)}=e^{\delta_1 d_Z(x,y)/8}$
gives \eqref{9.5:eq-capped-index-3}.
\end{proof}

\begin{lemma}[Substitution into ambient units on the whole box]\label{9.5:lem-ambient-substitution}
Let the comparison bound \eqref{9.5:eq-lettered-comparison-a} of Theorem~\ref{9.5:thm-lettered-comparison}
hold at the lettered link $(Z,\chi,b)$ over $p$, with $p^{\chi}$, $\mathsf{T}_{p^{\chi}}$ and $\mathsf{T}_p$ as in
that theorem. Assume the containment \eqref{9.5:eq-majorant-containment-a} of the lettered majorant
(Lemma~\ref{9.5:lem-majorant-containment}, whose proof is outstanding), and the frame hypothesis among
the standing hypotheses of Notation~\ref{9.5:not-hypotheses-summation}, which supplies the constant
$c_b^{\mathrm{fr}}$ below. Let $\mathsf{y}(p)$ be the
yardstick of Proposition~\ref{9.5:prop-reading-decay}. Then
\begin{equation}\label{9.5:eq-ambient-substitution-a}
\mathsf{T}_{p^{\chi}}(x)-\mathsf{T}_p(x)
\le \tfrac12\,\ell_\iota(x)\,a(x)\,\vartheta^{\iota^e_Z(x)}\bigl[c_b+c_a\,\mathsf{y}(p)\bigr]
\qquad (x\in Z),
\end{equation}
where
\begin{equation*}
c_b:=C_{R'}+c_b^{\mathrm{fr}}+\tfrac14\,C\,C_\chi^2\bigl(M^{-1}+a_{R'}\bigr),
\qquad
c_a:=\bigl[C_{R'}L^{\theta_2}C_T+\tfrac12\,C\,C_\chi^2\bigr]\,\bar{\bar{u}},
\end{equation*}
and $c_a\to0$ as $\gamma\to0$.
\end{lemma}

\begin{proof}
Fix $x\in Z$. We write the right member of \eqref{9.5:eq-lettered-comparison-a} at $x$ as the sum of
five terms, all evaluated at $x$:
\begin{equation}\label{9.5:eq-ambient-substitution-1}
\begin{split}
&\text{(i)}\ 2C_{R'}\,\mathsf{H}_\chi\,\mathsf{T}_p^{\sim}
\;+\;\text{(ii)}\ 2C_{R'}\,\mathsf{H}_\chi\,L^{-\theta\lambda_Z}
\;+\;\text{(iii)}\ 2C_{R'}\,C_\chi\,\bar{\tau}_\chi\,\mathsf{P}_{\delta b}\\
&\quad+\;\text{(iv)}\ C\,C_\chi^2\,L^{-\lambda_Z}\,\mathsf{H}_\chi\bigl(M^{-1}+\mathsf{H}_\chi\bigr)
\;+\;\text{(v)}\ C\,C_\chi^2\,L^{-\lambda_Z}\,\mathsf{H}_\chi\,\mathsf{T}_p .
\end{split}
\end{equation}
The factor $C_\chi^2$ is the enlargement of the constant $C$ stated in
Theorem~\ref{9.5:thm-lettered-comparison}.

Put $w:=v\,\vartheta^{\iota^e_Z}$. By the definition of $\mathsf{y}(p)$ in
Proposition~\ref{9.5:prop-reading-decay} as the supremum of $\mathsf{T}_p(y)/(v(y)\vartheta^{\iota^e_Z(y)})$
over the region carrying $\mathsf{T}_p$, we have $\mathsf{T}_p(y)\le\mathsf{y}(p)\,w(y)$ at every $y$
carrying $\mathsf{T}_p$. The smear $\mathsf{T}_p^{\sim}$ at the current link is taken, by
Definition~\ref{9.5:def-total-variation}, in the metric of the shells $\boldsymbol{\Omega}(Z)$ for the
terms living on $Z$, and in the metrics of the earlier links for the entry term. Writing $d_\sharp$ for
the metric in which a given term is smeared, each of these metrics dominates the box metric $d_Z$ in
which part (b) of Lemma~\ref{9.5:lem-capped-index} is stated: $d_\sharp\ge d_Z$ in all these cases. The
proof of part (b) of Lemma~\ref{9.5:lem-capped-index} uses only the inequality $d_\sharp\ge d_Z$ and
part (a) of that lemma; hence part (b) applies in either case and gives
\begin{equation*}
\mathsf{T}_p^{\sim}(x)\le L^{\theta_2}C_T\,\mathsf{y}(p)\,w(x).
\end{equation*}
From the containment \eqref{9.5:eq-majorant-containment-a} (Lemma~\ref{9.5:lem-majorant-containment},
whose proof is outstanding) we use, at $x$, the two bounds $\mathsf{H}_\chi\le\frac14\ell_\iota a$ and
$\mathsf{H}_\chi\le a_{R'}$. We now estimate the five terms of \eqref{9.5:eq-ambient-substitution-1}.

(i) By $\mathsf{H}_\chi\le\frac14\ell_\iota a$ and the bound on $\mathsf{T}_p^{\sim}(x)$, and then by
$v\le\bar{\bar{u}}$,
\begin{equation*}
\text{(i)}
\le \tfrac12 C_{R'}\,\ell_\iota a\,L^{\theta_2}C_T\,\mathsf{y}(p)\,v(x)\,\vartheta^{\iota^e_Z(x)}
\le \tfrac12\,\ell_\iota a\,\vartheta^{\iota^e_Z}\bigl[C_{R'}L^{\theta_2}C_T\,\bar{\bar{u}}\bigr]\,
\mathsf{y}(p).
\end{equation*}

(ii) By $\mathsf{H}_\chi\le\frac14\ell_\iota a$ and by part (c) of Lemma~\ref{9.5:lem-capped-index},
which states $L^{-\theta\lambda_Z}\le\vartheta^{\iota^e_Z}$ on $Z$,
\begin{equation*}
\text{(ii)}\le \tfrac12 C_{R'}\,\ell_\iota a\,L^{-\theta\lambda_Z}
\le \tfrac12\,\ell_\iota a\,\vartheta^{\iota^e_Z}\,C_{R'} .
\end{equation*}

(iii) The frame hypothesis of Notation~\ref{9.5:not-hypotheses-summation} bounds the term (iii) at $x$:
$\text{(iii)}\le\frac12\,\ell_\iota a\,\vartheta^{\iota^e_Z}\,c_b^{\mathrm{fr}}$.

(iv) By $\mathsf{H}_\chi\le\frac14\ell_\iota a$ for the first factor $\mathsf{H}_\chi$ and
$\mathsf{H}_\chi\le a_{R'}$ inside the bracket,
\begin{equation*}
\text{(iv)}\le \tfrac14\,C\,C_\chi^2\,\ell_\iota a\,L^{-\lambda_Z}\bigl(M^{-1}+a_{R'}\bigr),
\end{equation*}
and $L^{-\lambda_Z}\le L^{-\theta\lambda_Z}\le\vartheta^{\iota^e_Z}$, the second inequality being again
part (c) of Lemma~\ref{9.5:lem-capped-index}.

(v) By $\mathsf{H}_\chi\le\frac14\ell_\iota a$, by $\mathsf{T}_p(x)\le\mathsf{y}(p)\,w(x)$, by
$L^{-\lambda_Z(x)}\le1$ and by $v\le\bar{\bar{u}}$,
\begin{equation*}
\text{(v)}\le \tfrac14\,C\,C_\chi^2\,\ell_\iota a\,L^{-\lambda_Z}\,\mathsf{y}(p)\,w(x)
\le \tfrac12\,\ell_\iota a\,\vartheta^{\iota^e_Z}\bigl[\tfrac12\,C\,C_\chi^2\,\bar{\bar{u}}\bigr]\,
\mathsf{y}(p).
\end{equation*}

Adding the bounds on (i)--(v) and using \eqref{9.5:eq-lettered-comparison-a} yields
\eqref{9.5:eq-ambient-substitution-a} at $x$.
\end{proof}

\begin{remark}[The constants and the pointwise reading of the substitution bound]
\leavevmode
\begin{enumerate}
\item The constants $c_b$ and $c_a$ of Lemma~\ref{9.5:lem-ambient-substitution} are free of $g$, $N$,
$N_0$, $k$, $K$, of ages and volume, of $r$, of the box $Z$, of $D$ and of the letters; $c_b$ is also
free of $\check{R}$, while $c_a$ depends on $\check{R}$ only through the factor $\bar{\bar{u}}$.
\item The pointwise reading of \eqref{9.5:eq-ambient-substitution-a} is the one of
Proposition~\ref{9.5:prop-reading-decay}: on $Z^{=}$ it is the reading-decay bound
\eqref{9.5:eq-reading-decay-a} in label units, with $\mathsf{y}(p)$ as yardstick; off $Z^{=}$ it carries
the factor $L^{\operatorname{def}_Z}$, as in Lemma~\ref{9.5:lem-boundary-price} and its display
\eqref{9.5:eq-boundary-price-a}.
\item For the $s$-th link $p_{s-1}\to p_s$, of stage $\nu_s$, of a run presentation $p_0\to\dots\to p_r$
of the $(\mathrm{P})$-family as in Lemma~\ref{9.5:lem-run-points}, by part (ii) of that lemma both
members of \eqref{9.5:eq-ambient-substitution-a} vanish off $S_s^{+}\subset Z_s^{=}$. Hence
\eqref{9.5:eq-ambient-substitution-a} may be divided by $\ell_\iota=\ell_Z$ wherever its members do not
vanish, and gives
\begin{equation}\label{9.5:eq-ambient-substitution-b}
\mathsf{S}_{p_s}(x)-\mathsf{S}_{p_{s-1}}(x)
\le \tfrac12\,a_{\nu_s}(x)\,\vartheta^{\iota^e(x)}
\bigl[c_b+c_a\,\mathsf{y}(p_{s-1})\bigr]\,\mathbf{1}_{S_s^{+}}(x).
\end{equation}
\end{enumerate}
\end{remark}

\begin{proposition}[Reading decay in ambient units]\label{9.5:prop-reading-decay}
Assume the following.
\begin{enumerate}
\item[(a)] The named inputs from which Theorem~\ref{9.5:thm-lettered-comparison} (comparison
bounds for lettered links) is proved as an implication. Throughout, only those site types and
clauses of that theorem are considered at which it is established from these inputs.
\item[(b)] The three summation statements over stages, together with the ambient unit functions
and majorant functions that they fix.
\item[(c)] The composition identity.
\item[(d)] The frame conditions, in the form in which they enter Lemma~\ref{9.5:lem-ambient-substitution}
(substitution into ambient units on the whole box).
\item[(e)] The entry law for moves and the entry law at seams, both in the form of the entry
corollary.
\end{enumerate}
The following objects are fixed at a given site.
\begin{itemize}
\item $u$ is the ambient unit function of the site. At a site of type (P) it is given stage by
stage by $u_\nu:=\bar{u}_\nu\vartheta^{\iota^e_\nu}$, where $\vartheta=L^{-\theta_2}$; at sites of
types (T) and (L) it is the site's unit.
\item $a$ is the one-machine majorant function of the site's increments.
\item $\ell_\iota$ is the label scale.
\item $v:=\ell_\iota\bar{u}=:\bar{u}^{\mathrm{loc}}$ is the local unit, and
$\bar{\bar{u}}:=\sup\bar{u}^{\mathrm{loc}}$.
\end{itemize}
For a box $Z$ put
\begin{equation}\label{9.5:eq-reading-decay-1}
\iota^e_Z:=
\begin{cases}
\min\{\iota,\lambda_Z,\lambda_e\} & \text{on } Z,\\
\iota^e:=\min\{\iota,\lambda_e\} & \text{off } Z,
\end{cases}
\qquad
\mathsf{y}(p):=\sup_y\frac{\mathsf{T}_p(y)}{v(y)\,\vartheta^{\iota^e_Z(y)}}.
\end{equation}
The supremum is taken over the region carrying $\mathsf{T}_p$, and each term is read in the units
of its own link.

Then there are constants $c_b\ge 0$ and $c_a=c_a(\gamma)$, with $c_a(\gamma)\to 0$ as
$\gamma\to 0$, with the following property. Both constants are independent of $g$, $N$, $N_0$,
$\check{R}$, $k$, $K$, the ages, the volume, the number of links, the box, $D$ and the letters.
For every lettered link $(Z,\chi,b)$ of the families on which
Theorem~\ref{9.5:thm-lettered-comparison} is stated, over a real parent $p$ of the family with
$\mathsf{y}(p)\le\varpi_0$, and for every $x\in Z$,
\begin{equation}\label{9.5:eq-reading-decay-a}
\mathsf{T}_{p^\chi}(x)-\mathsf{T}_p(x)
\le\tfrac12\,\ell_\iota(x)\,a(x)\,\vartheta^{\iota^e_Z(x)}\bigl[c_b+c_a\,\mathsf{y}(p)\bigr].
\end{equation}
Hence, pointwise,
\begin{equation}\label{9.5:eq-reading-decay-2}
\sup_{\tilde{\Delta}_Z(x)}\bigl|A_{p^\chi}-A_p\bigr|
\le\tfrac12\,a(x)\,L^{\operatorname{def}_Z(x)}\,\vartheta^{\iota^e_Z(x)}
\bigl[c_b+c_a\,\mathsf{y}(p)\bigr].
\end{equation}
Consider the full-level region $Z^{=}$ of the box. It contains $K_D^{\circ}$, and therefore every
bond at which the unit $v_D$ reads an increment. On $Z^{=}$ the bound \eqref{9.5:eq-reading-decay-2}
is the increment bound in ambient units, read with the yardstick $\mathsf{y}(p)$.
\end{proposition}

\begin{proof}
Fix a site type and a clause of Theorem~\ref{9.5:thm-lettered-comparison} at which that theorem is
established from the inputs assumed in (a). Fix a lettered link $(Z,\chi,b)$ of its families over
a real parent $p$ with $\mathsf{y}(p)\le\varpi_0$, and a point $x\in Z$.

First, Theorem~\ref{9.5:thm-lettered-comparison} applies at this clause, because (a) supplies its
inputs. It gives its comparison bound, clause (C), at the link $(Z,\chi,b)$ over $p$.

Second, we check the hypotheses of Lemma~\ref{9.5:lem-ambient-substitution}. That lemma requires
three things:
\begin{itemize}
\item the comparison bound at the link, obtained in the first step;
\item the one-machine containment, which is one of the inputs of
Theorem~\ref{9.5:thm-lettered-comparison} and is therefore assumed in (a);
\item the frame conditions in the form in which the lemma reads them, which are assumed in (d).
\end{itemize}
The quantity $\mathsf{y}(p)$ in the lemma is the one defined in \eqref{9.5:eq-reading-decay-1}.
Lemma~\ref{9.5:lem-ambient-substitution} therefore yields \eqref{9.5:eq-ambient-substitution-a}
at the point $x\in Z$. This is exactly \eqref{9.5:eq-reading-decay-a}. The constants $c_b$ and
$c_a$ of the proposition are those of Lemma~\ref{9.5:lem-ambient-substitution}, with the
properties stated there.

Third, we record where the hypothesis $\mathsf{y}(p)\le\varpi_0$ is used. By the definition
\eqref{9.5:eq-reading-decay-1}, every $y$ in the region carrying $\mathsf{T}_p$ satisfies
\begin{equation*}
\mathsf{T}_p(y)\le\mathsf{y}(p)\,v(y)\,\vartheta^{\iota^e_Z(y)}.
\end{equation*}
Moreover $v\le\bar{\bar{u}}$, by the definition of $\bar{\bar{u}}$. The hypothesis
$\mathsf{y}(p)\le\varpi_0$ is used only to keep $\sup\mathsf{T}_p\le\varepsilon_D$, where
$\varepsilon_D$ is the ceiling of Theorem~\ref{9.5:thm-lettered-comparison}. This is done through
the condition
\begin{equation*}
(y^{\mathsf{s}}+1)\,\bar{\bar{u}}\le\varepsilon_D
\end{equation*}
among the conditions \eqref{9.5:eq-parameter-order-a}, imposed at each site type. Here
$y^{\mathsf{s}}$ is the site type's yardstick constant.

Since $x\in Z$ was arbitrary, \eqref{9.5:eq-reading-decay-a} holds at every point of $Z$, for
every parent with $\mathsf{y}(p)\le\varpi_0$, on each of the clauses considered.
\end{proof}

\begin{lemma}[No accumulation of total variation along a run]\label{9.5:lem-no-accumulation}
Assume clause (P) of the comparison bounds for lettered links
(Theorem~\ref{9.5:thm-lettered-comparison}); the single-lattice containment of the lettered
majorant \eqref{9.5:eq-majorant-containment-a} at site type (P)
(Lemma~\ref{9.5:lem-majorant-containment}, whose proof is outstanding); the standing hypotheses
on summation, frames and entries of Notation~\ref{9.5:not-hypotheses-summation}; and the smallness
ceilings among the conditions \eqref{9.5:eq-parameter-order-a} of
Remark~\ref{9.5:rem-parameter-order}.

Let $p_0\to p_1\to\cdots\to p_r$ be a run of the real lettered family
$\mathfrak{F}^{\mathbb{R},\mathbb{Y}}$ whose links have stages $n>\nu_1>\cdots>\nu_r$, and put
$\nu_0:=n$. For a stage $\nu$, a point $p_s$ of the run and a bond $y$ put
\begin{equation}\label{9.5:eq-no-accumulation-1}
\mathsf{S}^{(\nu)}_{p_s}(y):=\mathsf{E}^{\circ}(y)
+\sum_{t=1}^{s}\ \sup_{\tilde{\Delta}_\nu(y)}\bigl|A_{p_t}-A_{p_{t-1}}\bigr|,
\qquad
\mathsf{y}_\nu(p_s):=\sup_y\frac{\mathsf{S}^{(\nu)}_{p_s}(y)}{u_\nu(y)},
\end{equation}
where $\mathsf{E}^{\circ}$ is the entry modulus, $A_p$ is the mismatch field
\eqref{9.5:eq-mismatch-field-a}, $\tilde{\Delta}_\nu(y)$ are the label's cubes at stage $\nu$,
$u_\nu$ is the stage majorant of Notation~\ref{9.5:not-hypotheses-summation}, and the supremum over
$y$ runs over the union of the regions of the links of the run, each term of the sum being taken
where it is defined. Assume $\mathsf{y}_n(p_0)\le y_0$. Then for every $s\in\{0,\dots,r\}$
\begin{equation}\label{9.5:eq-no-accumulation-a}
\mathsf{y}_{\nu_s}(p_s)\le C^{P,\mathbb{Y}}:=4\bigl(y_0+C_\Sigma c_b\bigr),
\end{equation}
where $C_\Sigma$ is the summation constant of Notation~\ref{9.5:not-hypotheses-summation} and
$c_b$ is the constant of Lemma~\ref{9.5:lem-ambient-substitution}; the constant
$C^{P,\mathbb{Y}}$ is free of $r$, $N$, $N_0$, $\check{R}$, $k$ and $g$. Moreover, along the run
the ceilings of Theorem~\ref{9.5:thm-lettered-comparison} are reproduced: every point of the run
has class radius at most $\mathfrak{r}_0$, and its total variation
\eqref{9.5:eq-total-variation-b} satisfies $\sup\mathsf{T}\le\varepsilon_D$.
\end{lemma}

\begin{proof}
We use three monotonicity facts.

(a) For every point $p$ and every bond $y$, $\mathsf{S}^{(\nu)}_p(y)$ is non-decreasing in $\nu$.
Indeed the entry term does not depend on $\nu$, and the cubes $\tilde{\Delta}_\nu(y)$ grow with
$\nu$, because the level of the label at $y$ does, by $(\mathrm{s}2)$ in
\eqref{9.5:eq-hypotheses-summation-a}; so each supremum in \eqref{9.5:eq-no-accumulation-1} is
taken over a set that grows with $\nu$.

(b) The stage majorant $u_\nu$ is non-increasing in $\nu$; this follows from $(\mathrm{s}2)$ in
\eqref{9.5:eq-hypotheses-summation-a}.

(c) Hence $\mathsf{y}_\nu(p)$ is non-decreasing in $\nu$: in \eqref{9.5:eq-no-accumulation-1}
the numerator is non-decreasing by (a) and the positive denominator is non-increasing by (b). In
particular, for $1\le t\le s$ we have $\nu_t+1\le\nu_{t-1}$, and therefore
\begin{equation}\label{9.5:eq-no-accumulation-2}
\mathsf{y}_{\nu_t+1}(p_{t-1})\le\mathsf{y}_{\nu_{t-1}}(p_{t-1}).
\end{equation}

We prove \eqref{9.5:eq-no-accumulation-a} by induction on $s$, together with the two ceilings at
$p_s$, namely $\sup\mathsf{T}_{p_s}\le\varepsilon_D$ and $\mathfrak{r}_{P_s}\le\mathfrak{r}_0$
for the class radius $\mathfrak{r}_{P_s}$ of $p_s$. For $s=0$, \eqref{9.5:eq-no-accumulation-a}
is the hypothesis $\mathsf{y}_n(p_0)\le y_0$, since $y_0\le4(y_0+C_\Sigma c_b)$; the two ceilings
at $p_0$ are the chains \eqref{9.5:eq-no-accumulation-3} and \eqref{9.5:eq-no-accumulation-4}
below with $t=1$, in which the sum in \eqref{9.5:eq-no-accumulation-4} is empty. Let
$1\le s\le r$ and suppose that \eqref{9.5:eq-no-accumulation-a} and the two ceilings hold at every
index below $s$, that is, at $p_0,\dots,p_{s-1}$.

\emph{The ceilings at the parents.} Fix $t$ with $1\le t\le s$. Since $t-1<s$, the induction
hypothesis gives $\mathsf{y}_{\nu_{t-1}}(p_{t-1})\le C^{P,\mathbb{Y}}$ and the two ceilings at
$p_{t-1}$; they arise from \eqref{9.5:eq-no-accumulation-a} at $t-1$ as follows. With the
smallness condition $C^{P,\mathbb{Y}}\bar{\bar{u}}\le\varepsilon_D$ of
\eqref{9.5:eq-parameter-order-a},
\begin{equation}\label{9.5:eq-no-accumulation-3}
\sup\mathsf{T}_{p_{t-1}}
\le\mathsf{y}_{\nu_{t-1}}(p_{t-1})\cdot\sup(\ell_\iota u)
\le C^{P,\mathbb{Y}}\,\bar{\bar{u}}
\le\varepsilon_D ,
\end{equation}
where $\ell_\iota$ is the label's scale, and the middle inequality bounds $\sup(\ell_\iota u)$ by
$\bar{\bar{u}}$ of Notation~\ref{9.5:not-hypotheses-summation}.
For the class radius $\mathfrak{r}_{P_{t-1}}$ of $p_{t-1}$ we have, writing $\mathsf{H}_{t'}$ for
the lettered majorant $\mathsf{H}_\chi$ of the link $t'$, $\mathfrak{r}_e$ for the class radius
of the entry, $C_c\varepsilon$ for the bound on $\mathfrak{r}_e$, and $a^{\mathrm{loc}}_\nu$ for
the local majorant at stage $\nu$ in terms of which the containment
\eqref{9.5:eq-majorant-containment-a} is read at site type (P) in
Lemma~\ref{9.5:lem-majorant-containment},
\begin{equation}\label{9.5:eq-no-accumulation-4}
\begin{split}
\mathfrak{r}_{P_{t-1}}
&\le\mathfrak{r}_e+4CC_sC_\chi\sum_{t'<t}\sup\mathsf{H}_{t'}
\le C_c\varepsilon+CC_\chi\sum_{t'<t}\tfrac14\sup a^{\mathrm{loc}}_{\nu_{t'}}\\
&\le C_c\varepsilon+C_1'\,\bar{\bar{u}}
\le\mathfrak{r}_0 .
\end{split}
\end{equation}
The first inequality is the comparison bound \eqref{9.5:eq-lettered-comparison-b} applied at the
links $t'<t$, whose parents $p_{t'-1}$ satisfy the two ceilings by the induction hypothesis; the
second uses $\mathfrak{r}_e\le C_c\varepsilon$ and the containment
\eqref{9.5:eq-majorant-containment-a} (by
Lemma~\ref{9.5:lem-majorant-containment}, whose proof is outstanding); the third is the summation
bound $(\mathrm{s}3)^{0}$, the form of $(\mathrm{s}3)'$ in
\eqref{9.5:eq-hypotheses-summation-a} without the factors $\vartheta^{\iota^e}$, which is part of
the standing hypotheses of Notation~\ref{9.5:not-hypotheses-summation}, and $C_1'$ denotes the
constant it produces together with the factor $CC_\chi$; the last is the smallness condition
$C_c\varepsilon+C_1'\bar{\bar{u}}\le\mathfrak{r}_0$ of \eqref{9.5:eq-parameter-order-a}.

\emph{The substitution bound at link $t$.} By Lemma~\ref{9.5:lem-run-points}, the parent
$p_{t-1}$ on the box $Z_t$ satisfies the hypotheses of clause (P) of
Theorem~\ref{9.5:thm-lettered-comparison}, the ceilings required there being
\eqref{9.5:eq-no-accumulation-3} and \eqref{9.5:eq-no-accumulation-4}. Together with the assumed
containment \eqref{9.5:eq-majorant-containment-a}, the substitution bound
\eqref{9.5:eq-ambient-substitution-b} of Lemma~\ref{9.5:lem-ambient-substitution} therefore holds
at link $t$, with the yardstick $\mathsf{y}(p_{t-1})$ formed at that link. This yardstick is
dominated by $\mathsf{y}_{\nu_t+1}(p_{t-1})$: its box-capped index $\iota^e_{Z_t}$ does not exceed
$\iota^e$, and its cubes are those of the shells $\boldsymbol{\Omega}(Z_t)$, which are finer than
the label's cubes at stage $\nu_t+1$. With \eqref{9.5:eq-no-accumulation-2} and the induction
hypothesis at $t-1$ we obtain
\begin{equation}\label{9.5:eq-no-accumulation-5}
\mathsf{y}(p_{t-1})\le\mathsf{y}_{\nu_t+1}(p_{t-1})\le\mathsf{y}_{\nu_{t-1}}(p_{t-1})
\le C^{P,\mathbb{Y}} .
\end{equation}

\emph{Summation along the run.} Fix a bond $x$. We measure every term of
$\mathsf{S}^{(\nu_s)}_{p_s}(x)$ with the cubes $\tilde{\Delta}_{\nu_s}(x)$, which are contained in
$\tilde{\Delta}_{\nu_t}(x)$ for every $t\le s$ by (a), since $\nu_s\le\nu_t$. The entry term
$\mathsf{E}^{\circ}(x)$ is $\mathsf{S}^{(n)}_{p_0}(x)$. The $t$-th increment is bounded by
\eqref{9.5:eq-ambient-substitution-b} at link $t$, with the yardstick bounded by
\eqref{9.5:eq-no-accumulation-5}, and it vanishes outside the fattened support $S_t^{+}$ by
Lemma~\ref{9.5:lem-run-points}\,(ii). Hence
\begin{align}
\mathsf{S}^{(\nu_s)}_{p_s}(x)
&\le\mathsf{S}^{(n)}_{p_0}(x)+\sum_{t=1}^{s}\tfrac12\,a_{\nu_t}(x)\,
\vartheta^{\iota^e_{\nu_t}(x)}\bigl[c_b+c_aC^{P,\mathbb{Y}}\bigr]\mathbf{1}_{S_t^{+}}(x)
\notag\\
&\le y_0\,u_n(x)+\tfrac12\bigl[c_b+c_aC^{P,\mathbb{Y}}\bigr]
\sum_{t=1}^{s}a_{\nu_t}(x)\,\vartheta^{\iota^e_{\nu_t}(x)}
\notag\\
&\le y_0\,u_{\nu_s}(x)+\tfrac12\,C_\Sigma\bigl[c_b+c_aC^{P,\mathbb{Y}}\bigr]u_{\nu_s}(x).
\label{9.5:eq-no-accumulation-6}
\end{align}
The second inequality uses the hypothesis $\mathsf{y}_n(p_0)\le y_0$ and
$\mathbf{1}_{S_t^{+}}\le1$. The third uses (b), which gives $u_n\le u_{\nu_s}$ because $n>\nu_s$,
and the summation bound $(\mathrm{s}3)'$ of \eqref{9.5:eq-hypotheses-summation-a} applied to the
finite set of stages $\mathcal{S}=\{\nu_t:1\le t\le s\}$, whose least element is $\nu_s$.

By the ceiling $4C_\Sigma c_a\le1$, one of the smallness conditions in
\eqref{9.5:eq-parameter-order-a}, we have
$\tfrac12C_\Sigma c_aC^{P,\mathbb{Y}}\le\tfrac18C^{P,\mathbb{Y}}$, so that
the chain \eqref{9.5:eq-no-accumulation-6} gives
\begin{equation*}
\mathsf{S}^{(\nu_s)}_{p_s}(x)\le u_{\nu_s}(x)\bigl[y_0+\tfrac12C_\Sigma c_b
+\tfrac18C^{P,\mathbb{Y}}\bigr].
\end{equation*}
Since $\tfrac18C^{P,\mathbb{Y}}=\tfrac12(y_0+C_\Sigma c_b)$,
\begin{align*}
y_0+\tfrac12C_\Sigma c_b+\tfrac18C^{P,\mathbb{Y}}
&=y_0+\tfrac12C_\Sigma c_b+\tfrac12\bigl(y_0+C_\Sigma c_b\bigr)
=\tfrac32\,y_0+C_\Sigma c_b\\
&\le\tfrac32\bigl(y_0+C_\Sigma c_b\bigr)=\tfrac38\,C^{P,\mathbb{Y}}
\le\tfrac34\,C^{P,\mathbb{Y}} .
\end{align*}
Thus $\mathsf{S}^{(\nu_s)}_{p_s}(x)\le\tfrac34C^{P,\mathbb{Y}}u_{\nu_s}(x)$ for every bond $x$ of
the union of the regions of the links; dividing by $u_{\nu_s}(x)$ and taking the supremum gives
$\mathsf{y}_{\nu_s}(p_s)\le\tfrac34C^{P,\mathbb{Y}}\le C^{P,\mathbb{Y}}$, which is
\eqref{9.5:eq-no-accumulation-a} at $s$.

The chains \eqref{9.5:eq-no-accumulation-3} and \eqref{9.5:eq-no-accumulation-4}, written with $s$
in place of $t-1$, give the two ceilings at $p_s$: the first uses
\eqref{9.5:eq-no-accumulation-a} at $s$, and the second uses
\eqref{9.5:eq-lettered-comparison-b} at the links $t'\le s$, whose parents $p_{t'-1}$ satisfy the
two ceilings by the induction hypothesis. This completes the induction; it proves
\eqref{9.5:eq-no-accumulation-a} for every $s\in\{0,\dots,r\}$ and shows that the ceilings are
reproduced along the whole run.
\end{proof}

We record four features of the lemma and its proof.
\begin{enumerate}
\item The recursion in the first line of the chain \eqref{9.5:eq-no-accumulation-6} is additive,
with coefficient one in front of
the earlier variation: the multiplicative factor $k_B$ of the composition identity of
Notation~\ref{9.5:not-hypotheses-summation} is contained in the last term of item (C) of
Theorem~\ref{9.5:thm-lettered-comparison}, and is converted in part (v) of
Lemma~\ref{9.5:lem-ambient-substitution}.
\item The supremum defining $\mathsf{y}_\nu$ runs over every bond of every region carrying a term
of $\mathsf{S}^{(\nu)}$, in particular over bonds off the region $X^{+}$ and off the region of the
current link. At such bonds the later increments vanish, by Lemma~\ref{9.5:lem-run-points}\,(ii),
and the bound persists by (b).
\item None of the constants depends on $r$, on the number $|\mathsf{A}(x)|$ of stages through a
bond $x$, on $N$ or on $N_0$; the parameter $\check{R}$ enters only through $\bar{\bar{u}}$, and
only inside the ceilings.
\item For the entries, $y_0=C_T'c_0$ for seeds, by the entry laws, using that $u$ varies slowly
inside a cube; and $y_0=y_0^{\mathsf{T}}$ for entries of type (e-ii) of chains without dilated rows
of object type, by the corollary on entries, given the weight law for the chain's objects. For
lettered chains the entry's yardstick is the constant $y^L$ of the lemma on lettered chains that
follows.
\end{enumerate}

\begin{lemma}[The yardstick after one step and along a chain]\label{9.5:lem-chain-yardstick}
Let $u$ be the ambient unit function, $a$ the majorant function, $\ell_\iota$ the label scale,
$v=\ell_\iota\bar{u}=\bar{u}^{\mathrm{loc}}$ the local unit, $\bar{\bar{u}}=\sup\bar{u}^{\mathrm{loc}}$, and let
$\iota^e$, $\iota^e_Z$ and the yardstick $\mathsf{y}(p)$ be as in
Proposition~\ref{9.5:prop-reading-decay}. Let $c_b$, $c_a$ be the constants of
Lemma~\ref{9.5:lem-ambient-substitution}, and put
\begin{equation}\label{9.5:eq-chain-yardstick-1}
C_a:=\sup\frac{a\,\ell_\iota}{2v}=\sup\frac{a}{2\bar{u}} .
\end{equation}
Fix a term over strips and a box $Z$ of the families of
Theorem~\ref{9.5:thm-lettered-comparison}. At the site type $(T)$ let $p_0$ be the seed and $p_1$
its image under the lettered link of clause $(T)$ of Theorem~\ref{9.5:thm-lettered-comparison}.
At the site type $(L)$ fix a chain of truncations $p^{(0)},p^{(1)},\dots,p^{(m)}$ read on $Z$,
row $r$ being the link $p^{(r-1)}\mapsto p^{(r)}$, with $m\le m_L$, where $m_L$ is a number free
of $g$. For a row $r$ of object type write $\mathcal{Z}^{(r)}$ for its displacement,
$q_{\mathrm{prev}}$ for the chain's unlettered object before the row and
$q=e^{i\eta\mathcal{Z}^{(r)}}q_{\mathrm{prev}}$ for the unlettered object after it. Assume:
\begin{enumerate}
\item[(a)] at every lettered link of the seed step at $(T)$ and of every row of response type at
$(L)$, the comparison bound \eqref{9.5:eq-lettered-comparison-a} of clause $(T)$, respectively
clause $(L)$, of Theorem~\ref{9.5:thm-lettered-comparison} holds, together with the remaining
hypotheses of Lemma~\ref{9.5:lem-ambient-substitution};
\item[(b)] the entry laws, with $c_0$ the constant of these laws: at $(T)$ the seed $p_0$
satisfies, on its region,
$\mathsf{T}_{p_0}(y)=|A_{p_0}|_y\le c_0'\,v(y)\,\vartheta^{\iota(y)}$ with $c_0':=C_T'c_0$; at
$(L)$ the root $p^{(0)}$ satisfies $\mathsf{y}(p^{(0)})\le y_0^{\mathsf{T}}$;
\item[(c)] the law of the chain objects (the response bound together with the seam bound for the
chain's unlettered objects), namely: at every row $r$ of object type at $(L)$ the mismatch fields
of the unlettered objects satisfy $|A_q|\le C_{LR'}u$ and $|A_{q_{\mathrm{prev}}}|\le C_{LR'}u$;
\item[(d)] the parameters satisfy the conditions of Remark~\ref{9.5:rem-parameter-order}.
\end{enumerate}
Then the following hold.
\begin{enumerate}
\item[(i)] At $(T)$: $\mathsf{y}(p_0)\le c_0'$, $\sup\mathsf{T}_{p_0}\le c_0'\bar{\bar{u}}\le\varepsilon_D$,
and for every $x\in Z$
\begin{equation}\label{9.5:eq-chain-yardstick-2}
\mathsf{T}_{p_1}(x)\le v(x)\,\vartheta^{\iota^e_Z(x)}\,y^T,\qquad
y^T:=c_0'+C_a\bigl(c_b+c_ac_0'\bigr).
\end{equation}
\item[(ii)] At $(L)$, suppose that every row of object type carries strip letters in $\{0,1\}$. Put
$y^{(0)}:=y_0^{\mathsf{T}}$ and, for $1\le r\le m_L$,
\begin{equation}\label{9.5:eq-chain-yardstick-3}
y^{(r)}:=
\begin{cases}
y^{(r-1)}+C_a\bigl(c_b+c_ay^{(r-1)}\bigr) & \text{if $r>m$ or row $r$ is of response type,}\\
y^{(r-1)}+C_{\mathrm{obj}} & \text{if $r\le m$ and row $r$ is of object type,}
\end{cases}
\end{equation}
where
\begin{equation*}
C_{\mathrm{obj}}:=4\bigl(c_0^\chi\vee1\bigr)C_T'C_{LR'}\bigl(1+C_\chi/M+\cdots\bigr),
\end{equation*}
the dots denoting further contributions free of $g$; thus $C_{\mathrm{obj}}$ is free of $g$.
Then for $0\le r\le m$
and every $x\in Z$, $\mathsf{T}_{p^{(r)}}(x)\le v(x)\,\vartheta^{\iota^e_Z(x)}\,y^{(r)}$; and
$y^L:=y^{(m_L)}$ is a number free of $g$.
\end{enumerate}
\end{lemma}

\begin{proof}
\emph{(i) The seed step at $(T)$.} By the entry law (b),
$\mathsf{T}_{p_0}(y)=|A_{p_0}|_y\le c_0'\,v(y)\,\vartheta^{\iota(y)}$ on the seed's region. Since
$\iota^e_Z=\min\{\iota,\lambda_Z,\lambda_e\}\le\iota$ and $0<\vartheta<1$, we have
$\vartheta^{\iota(y)}\le\vartheta^{\iota^e_Z(y)}$, so the definition of $\mathsf{y}(p_0)$ as the
supremum of $\mathsf{T}_{p_0}(y)/(v(y)\vartheta^{\iota^e_Z(y)})$ gives $\mathsf{y}(p_0)\le c_0'$.
Since $\vartheta^{\iota}\le1$ and $v=\bar{u}^{\mathrm{loc}}\le\bar{\bar{u}}$, also
$\sup\mathsf{T}_{p_0}\le c_0'\bar{\bar{u}}$, and $c_0'\bar{\bar{u}}\le\varepsilon_D$ is one of the
ceiling conditions of Remark~\ref{9.5:rem-parameter-order} (see
\eqref{9.5:eq-parameter-order-a}). Thus the seed lies below the ceilings of
Theorem~\ref{9.5:thm-lettered-comparison}, and by hypothesis (a) clause $(T)$ of that theorem
supplies the comparison bound \eqref{9.5:eq-lettered-comparison-a} at the link $p_0\mapsto p_1$.
Lemma~\ref{9.5:lem-ambient-substitution} therefore applies, and \eqref{9.5:eq-ambient-substitution-a}
gives, for $x\in Z$,
\begin{equation*}
\mathsf{T}_{p_1}(x)-\mathsf{T}_{p_0}(x)
\le\tfrac12\ell_\iota(x)a(x)\vartheta^{\iota^e_Z(x)}\bigl[c_b+c_a\mathsf{y}(p_0)\bigr].
\end{equation*}
By \eqref{9.5:eq-chain-yardstick-1},
$\tfrac12\ell_\iota a\vartheta^{\iota^e_Z}\le C_a v\vartheta^{\iota^e_Z}$; here $v=\ell_\iota\bar{u}$
gives the second expression for $C_a$, and $C_a=1$ where $a=2\bar{u}$, which is the case on the
interval of stages through $x$ by Notation~\ref{9.5:not-hypotheses-summation}. With
$\mathsf{y}(p_0)\le c_0'$ the right side is at most $C_a(c_b+c_ac_0')\,v(x)\vartheta^{\iota^e_Z(x)}$,
and by the definition of $\mathsf{y}(p_0)$,
$\mathsf{T}_{p_0}(x)\le c_0'\,v(x)\vartheta^{\iota^e_Z(x)}$. Adding the two bounds gives
\eqref{9.5:eq-chain-yardstick-2}.

\emph{(ii) Rows of response type at $(L)$.} We argue by induction on $r\le m\le m_L$. For $r=0$,
the root $p^{(0)}$ is an entry (Lemma~\ref{9.5:lem-response-rows}, part (i)), and the entry law
(b) gives $\mathsf{y}(p^{(0)})\le y_0^{\mathsf{T}}=y^{(0)}$, which by the definition of the
yardstick is the assertion at $r=0$. Let $r\ge1$ and assume
$\mathsf{T}_{p^{(r-1)}}(x)\le v(x)\vartheta^{\iota^e_Z(x)}y^{(r-1)}$ for $x\in Z$, that is,
$\mathsf{y}(p^{(r-1)})\le y^{(r-1)}$. If row $r$ is of response type, the ceilings of
Theorem~\ref{9.5:thm-lettered-comparison} are reproduced at this row by the ceiling conditions
of Remark~\ref{9.5:rem-parameter-order}, so clause $(L)$ of that theorem supplies, by hypothesis
(a), the comparison bound at the link $p^{(r-1)}\mapsto p^{(r)}$, and
Lemma~\ref{9.5:lem-ambient-substitution} together with the inequality
$\tfrac12\ell_\iota a\vartheta^{\iota^e_Z}\le C_av\vartheta^{\iota^e_Z}$ of part (i) gives, for
$x\in Z$,
\begin{equation*}
\mathsf{T}_{p^{(r)}}(x)
\le v(x)\vartheta^{\iota^e_Z(x)}\Bigl[y^{(r-1)}+C_a\bigl(c_b+c_ay^{(r-1)}\bigr)\Bigr]
=v(x)\vartheta^{\iota^e_Z(x)}\,y^{(r)} .
\end{equation*}

\emph{(iii) A row of object type at $(L)$ with strip letters in $\{0,1\}$.} At such a row the
comparison bound \eqref{9.5:eq-lettered-comparison-a} is not available, and we bound the link
directly. Write $\mathcal{Z}:=\mathcal{Z}^{(r)}$, so that $q=e^{i\eta\mathcal{Z}}q_{\mathrm{prev}}$
with $q_{\mathrm{prev}}$, $q$ the unlettered objects before and after the row; write
$\mathcal{Z}'$ for the
displacement read in the previous configuration $\mathbf{U}'_{\mathrm{prev}}$, and
$Y[\chi\mathcal{Z}';\mathbf{U}'_{\mathrm{prev}}]$, $Y[\mathcal{Z}']$ for the field $Y$ of the
composition identity of Notation~\ref{9.5:not-hypotheses-summation} at the lettered and at the
unlettered object link; $Q_1$ is the linear operator through which the displacement enters this
field. Applying the composition identity in both directions, with $k_B'$ the multiplicative
factor, of the form $1+C_B\eta(|Y|+|\mathfrak{h}|)$, that the identity carries when applied in the
reverse direction,
\begin{align*}
\bigl|Y[\chi\mathcal{Z}';\mathbf{U}'_{\mathrm{prev}}]-\chi\mathcal{Z}\bigr|
&\le|\chi|\cdot\bigl|Y[\mathcal{Z}']-\mathcal{Z}\bigr|+\bigl|[Q_1,\chi]\mathcal{Z}'\bigr|
+\text{quadratic terms},\\
\bigl|Y[\mathcal{Z}']-\mathcal{Z}\bigr|
&\le k_B'\bigl(|A_q|+|A_{q_{\mathrm{prev}}}|\bigr).
\end{align*}
By hypothesis (c), $|A_q|$ and $|A_{q_{\mathrm{prev}}}|$ are at most $C_{LR'}u$. The letters being
in $\{0,1\}$ on the strip blocks and those of the assignment on the plain blocks, the cut-off
satisfies, in the units of Lemma~\ref{9.5:lem-cutoff-majorant}, part (b),
\begin{equation*}
|\chi(y)|\le c_0^\chi\vee1,\qquad \ell(y)\,|\nabla\chi(y)|\le C_\chi\bigl(c_0^\chi\vee1\bigr)/M .
\end{equation*}
Inserting these bounds, the $\mathsf{T}$-increment of the lettered object link is at most
$v\,\vartheta^{\iota^e}\,C_{\mathrm{obj}}$, with $C_{\mathrm{obj}}$ as in part (ii) of the
statement.
Since $\iota^e=\min\{\iota,\lambda_e\}\ge\iota^e_Z$ and $\vartheta<1$, we have
$\vartheta^{\iota^e}\le\vartheta^{\iota^e_Z}$ on $Z$, so the induction hypothesis gives, for
$x\in Z$,
\begin{equation*}
\mathsf{T}_{p^{(r)}}(x)\le v(x)\vartheta^{\iota^e_Z(x)}\bigl(y^{(r-1)}+C_{\mathrm{obj}}\bigr)
=v(x)\vartheta^{\iota^e_Z(x)}y^{(r)},
\end{equation*}
in accordance with
\eqref{9.5:eq-chain-yardstick-3}, and the induction of part (ii) continues.

This completes the induction. Since $m_L$ is a number and each step of
\eqref{9.5:eq-chain-yardstick-3} adds a quantity determined by $C_a$, $c_b$, $c_a$ and
$C_{\mathrm{obj}}$, all free of $g$, the number $y^L=y^{(m_L)}$ is free of $g$.

The computation of part (iii) does not extend to rows of object type carrying dilated strip
letters: there the same computation yields, in place of $c_0^\chi\vee1$, the factor
$\bar{\tau}_\chi$ of Lemma~\ref{9.5:lem-cutoff-majorant}, which is comparable to $r^L_\square$
and is not bounded above by a number free of $g$, since the radius $r^0_L$ is defined by the room
of the second-derivative room lemma and grows with the distance from the strip to $\bar{\Lambda}$.
Such rows are accordingly excluded from part (ii) of the statement; they are taken up separately,
in the statement on rows of object type carrying dilated strip letters.
\end{proof}

\begin{proposition}[Complex strip letters at the cost of a square root]\label{9.5:prop-complex-letters}
Let $\mathsf{s}$ be one of the site types (T), (L) and (P). Write $y^{\mathsf{s}}$ for the
real-parameter constant of that site type: $y^{T}$ at (T), $y^{L}$ at (L) and
$C^{P,\mathbb{Y}}$ at (P). Let $D$ be the set of indexed strip blocks $(\square,r)$ of the
unit, and $t_{\square,r}$ the strip letters. Let $\mathfrak{P}^{\mathbb{Y}}$ be the polydisc
that is the product of four factors. The first is the range of the complex source datum $b$.
The second is the product of the discs of the complex plain letters. The third is the factor
\begin{equation*}
\prod_{(\square,r)\in D}\bigl\{\,|t_{\square,r}|\le r'_\square/4\,\bigr\}
\end{equation*}
for the strip letters. The fourth is the product of the discs of the seed's complex parameters.
Each coordinate ranges over a closed disc about a real centre; the last factor is present when
the seed carries complex parameters. Write $4\mathfrak{P}^{\mathbb{Y}}$ for the polydisc with
the same centres and four times the radii.

In addition to the standing hypotheses of Notation~\ref{9.5:not-hypotheses-summation}, assume
the following.
\begin{enumerate}
\item[(a)] The one-machine complex extension hypothesis holds, with a priori constant
$C_{\mathrm{ap}}$; its proof is outstanding. It includes the analyticity of the one-step
block-averaging map $M$. Moreover, the response field of the strip-letter response lemma is
holomorphic in the strip letters at radii $r'_\square$.
\item[(b)] The square-root lemma holds.
\item[(c)] The real-parameter bounds of Lemma~\ref{9.5:lem-no-accumulation} and
Lemma~\ref{9.5:lem-chain-yardstick} hold.
\item[(d)] The lower bound on the letter radii among the conditions
\eqref{9.5:eq-parameter-order-a} of Remark~\ref{9.5:rem-parameter-order} holds.
\end{enumerate}
Let $\Pi$ be a point of $\mathfrak{P}^{\mathbb{Y}}$, and write $p=p(\Pi)$ for the pair it
determines. Let $x\in Z^{=}$; at (P), let $x$ lie in the interior $K^{\circ}_{s}$ of the fat
core. Then
\begin{equation}\label{9.5:eq-complex-letters-1}
|A_p(x)|\le C^{\mathsf{s}}_{\mathbb{C}}\,\bar{u}(x)\,\vartheta^{\iota^e(x)/2},
\qquad
C^{\mathsf{s}}_{\mathbb{C}}:=\dim\mathfrak{g}\cdot\bigl(2C_{\mathrm{ap}}\,y^{\mathsf{s}}\bigr)^{1/2}.
\end{equation}
\end{proposition}

\begin{proof}
Fix the site type $\mathsf{s}$ and a point $x$ as in the statement. Write $\Pi$ for a point of
the parameter polydisc and $p(\Pi)$ for the pair it determines.

\emph{Holomorphy on the fourfold polydisc.} The map $\Pi\mapsto A_{p(\Pi)}(x)$ is holomorphic
on $4\mathfrak{P}^{\mathbb{Y}}$. This follows from hypothesis~(a). The response field of the
strip-letter response lemma is holomorphic in the strip letters at radii $r'_\square$. This
radius is four times the radius $r'_\square/4$ that appears in $\mathfrak{P}^{\mathbb{Y}}$. The
block-averaging map $M$ is analytic.

\emph{A priori bound.} On $4\mathfrak{P}^{\mathbb{Y}}$ the same map is bounded a priori by
\begin{equation*}
|A_{p(\Pi)}(x)|\le C_{\mathrm{ap}}\,\bar{u}(x).
\end{equation*}
Three inputs give this bound. The first is the composition identity of
Notation~\ref{9.5:not-hypotheses-summation}. It is used with the one-machine sizes
$35d\,\sup|\chi\mathfrak{h}|\le a$, the standing size condition of the one-machine setting. Here
$\chi$ is the lettered cut-off and $\mathfrak{h}$ the increment field of the link. The second
input is hypothesis~(a). The third is the summation bound of
Notation~\ref{9.5:not-hypotheses-summation}:
\begin{equation*}
\sum_{\nu'\in\mathcal{S}}a_{\nu'}\vartheta^{\iota^e_{\nu'}}\le C_\Sigma\,u_{\min\mathcal{S}}
\qquad\text{for every finite set } \mathcal{S} \text{ of stages.}
\end{equation*}

\emph{The complex line through a point.} Let $\Pi\in\mathfrak{P}^{\mathbb{Y}}$. Write
$\Pi=\Pi_1+i\Pi_2$ with $\Pi_1$ and $\Pi_2$ real, taken relative to the real centres. For
$|z|<3$ put
\begin{equation*}
\Pi(z):=\Pi_1+z\Pi_2 .
\end{equation*}
Take a coordinate with real centre $c$ and radius $\varsigma$. Write the coordinate of $\Pi$ in
it as $c+w$ with $|w|\le\varsigma$. Then the coordinate of $\Pi_1$ is $c+\operatorname{Re}w$,
and that of $\Pi_2$ is $\operatorname{Im}w$. Hence
\begin{equation*}
|\operatorname{Re}w+z\operatorname{Im}w|\le\varsigma+3\varsigma=4\varsigma .
\end{equation*}
So $z\mapsto\Pi(z)$, $|z|<3$, takes values in $4\mathfrak{P}^{\mathbb{Y}}$. Moreover
$\Pi(i)=\Pi$.

\emph{Real points of the line.} For real $z$, $\Pi(z)$ is a real parameter. Its strip letters
are real and satisfy $|t_{\square,r}|\le r'_\square$. This is the estimate just made with
$\varsigma=r'_\square/4$. Lemma~\ref{9.5:lem-cutoff-majorant} applies to every real letter
assignment in these ranges. Hence $\Pi(z)$ is a point of the real lettered family
$\mathfrak{F}^{\mathbb{R},\mathbb{Y}}$. Hypothesis~(c) then gives
\begin{equation}\label{9.5:eq-complex-letters-2}
|A_{p(\Pi(z))}(x)|\le y^{\mathsf{s}}\,\bar{u}(x)\,\vartheta^{\iota^e(x)},
\qquad z\in\mathbb{R},\ |z|<3,
\end{equation}
on $Z^{=}$, and at (P) on $K^{\circ}_{s}$.

\emph{The square root.} Expand $A_{p(\Pi(z))}(x)$ in an orthonormal basis of $\mathfrak{g}$.
Then each component is bounded in modulus by the norm, and the norm is bounded by the sum of the
moduli of the components. Each of the $\dim\mathfrak{g}$ components is a holomorphic function of
$z$ on $|z|<3$. By the a priori bound, it is bounded on that disc by
$\mathsf{M}:=C_{\mathrm{ap}}\bar{u}(x)$. By \eqref{9.5:eq-complex-letters-2}, it is bounded at
the real points by $m:=y^{\mathsf{s}}\bar{u}(x)\vartheta^{\iota^e(x)}$.

We use the square-root lemma, hypothesis~(b), in the following form. A function holomorphic on
the disc $|z|<3$ is bounded in modulus by $(2\mathsf{M}m)^{1/2}$ at $z=i$, provided it is bounded
in modulus by $\mathsf{M}$ on the disc and by $m$ on its real segment. We apply this to each
component at $z=i$. Each component of $A_{p(\Pi)}(x)$ is then bounded by
\begin{equation*}
(2\mathsf{M}m)^{1/2}=\bigl(2C_{\mathrm{ap}}y^{\mathsf{s}}\bigr)^{1/2}\,
\bar{u}(x)\,\vartheta^{\iota^e(x)/2}.
\end{equation*}
Summing over the $\dim\mathfrak{g}$ components gives \eqref{9.5:eq-complex-letters-1}.

Hypothesis~(d) plays no role in the proof of the bound \eqref{9.5:eq-complex-letters-1}. It makes
$r'_\square/4\ge 2$. The circles of the lettered merging lemma require this on the polydisc
$\mathfrak{P}^{\mathbb{Y}}$.
\end{proof}

\begin{definition}[The strip-lettered pointwise mismatch law]
Let $\mathsf{A}$ and $\mathsf{B}$ be the coarse-lattice and the fine-lattice machines. For a pair $p$ of
their configurations let $A_p$ be the mismatch field of \eqref{9.5:eq-mismatch-field-a}. Let
$P(\beta)$, $\beta\in\mathfrak{B}$, be the cells at which the pair of configurations is read. The
\emph{strip-lettered pointwise mismatch law} consists of the following two conditions, required for
every $D$ and every vertex point.
\begin{itemize}
\item[(i)] \emph{Regularity clause.} The pair formed by the lettered configurations of the two
machines satisfies the regularity clause of the external response lemma on the region $X^{+}$ of
that lemma.
\item[(ii)] \emph{Field clause.} There are a number $y\ge 1$, independent of $g$, and a fixed
power $\rho$ of $\vartheta$ such that
\begin{equation}\label{9.5:eq-mismatch-law-1}
|A_p(x)|\le y\,\bar{u}(x)\,\rho^{\iota^e(x)}
\qquad\text{for every bond $x$ of every cell $P(\beta)$, $\beta\in\mathfrak{B}$.}
\end{equation}
\end{itemize}
The \emph{mismatch law at letters in $\{0,1\}$} is the restriction of this law to real parameters
and to strip letters in $\{0,1\}$.
\end{definition}

Theorem~\ref{9.5:thm-mismatch-law} below proves the field clause (ii). At real parameters it
holds with $(y,\rho)=(y^{\mathsf{s}},\vartheta)$. On the complex polydiscs of
Proposition~\ref{9.5:prop-complex-letters} it holds with
$(y,\rho)=(C_{\mathbb{C}}^{\mathsf{s}},\vartheta^{1/2})$. The sites covered are as follows: sites
of type (P), for seed-entered runs and for (L)-entered runs under the restriction of
Lemma~\ref{9.5:lem-chain-yardstick}(iii); sites of type (T); and sites of type (L), for rows of
response type at all letters and for all rows at letters in $\{0,1\}$. At sites of type (L) the
rows of object type at dilated letters are not reached, and at sites of type (P) neither are the
(L)-entered runs without that restriction; these are the cases that require the relative bound
\eqref{9.5:eq-relative-object-bound-a} for dilated object rows, which is not proved. The
regularity clause (i) is not addressed.

\begin{theorem}[The pointwise mismatch law on read cells]\label{9.5:thm-mismatch-law}
Assume the hypotheses of Lemma~\ref{9.5:lem-no-accumulation} and of
Lemma~\ref{9.5:lem-chain-yardstick}. For the assertions on the polydisc $\mathfrak{P}^{\mathbb{Y}}$
of Proposition~\ref{9.5:prop-complex-letters}, assume in addition the hypotheses of that
proposition. For clause (b) at roots of the second entry type, assume the root entry bound
\begin{equation}\label{9.5:eq-mismatch-law-2}
|A_{p_0}(x)|\le C_{LR'}\,u(x)
\end{equation}
at the bonds $x$ described in (b); this bound is supplied by the two-machine root law together
with the seam condition. Then, at real parameters [resp.\ on $\mathfrak{P}^{\mathbb{Y}}$], for
every $D$ and every vertex box $Z$ the following hold.
\begin{itemize}
\item[(a)] Let $x$ be a bond of a cell at which $v_D$ reads an increment; such cells lie in
$K_D^{\circ}\subset Z^{=}$. Then
\begin{equation}\label{9.5:eq-mismatch-law-3}
|A_p(x)|\le y^{\mathsf{s}}\,\bar{u}(x)\,\vartheta^{\iota^e(x)}
\qquad\Bigl[\text{resp. }
|A_p(x)|\le C_{\mathbb{C}}^{\mathsf{s}}\,\bar{u}(x)\,\vartheta^{\iota^e(x)/2}\Bigr].
\end{equation}
Here the constant $y^{\mathsf{s}}$ is chosen as follows.
\begin{itemize}
\item At sites of type (P), $y^{\mathsf{s}}=C^{P,\mathbb{Y}}$. This covers seed-entered runs, and
(L)-entered runs under the restriction of Lemma~\ref{9.5:lem-chain-yardstick}(iii), where the
initial constant is $y_0:=y^L C_T'$.
\item At sites of type (T), $y^{\mathsf{s}}=y^T$.
\item At sites of type (L), $y^{\mathsf{s}}=y^L$. This covers rows of response type at all letters,
and all rows at letters in $\{0,1\}$.
\end{itemize}
\item[(b)] At every other bond $x$ of a cell $P(\beta)$, $\beta\in\mathfrak{B}$, the law is the
entry law. These are the bonds of the strip off $\bigcup_D\tilde{\square}$, where the configuration
is that of $U^0$. The entry law reads, for a seed and for a root of the second entry type
respectively,
\begin{equation}\label{9.5:eq-mismatch-law-4}
\begin{aligned}
|A_p(x)|=|A_{p_0}(x)|&\le c_0\,\bar{u}(x)\,\vartheta^{\iota(x)},\\
|A_p(x)|=|A_{p_0}(x)|&\le C_{LR'}\,u(x).
\end{aligned}
\end{equation}
\item[(c)] At every bond $x$ of $Z\setminus Z^{=}$, the law of (a) holds with $\iota^e$ replaced by
$\iota^e_Z$ and with the extra factor $L^{\operatorname{def}_Z(x)}$:
\begin{equation}\label{9.5:eq-mismatch-law-5}
\begin{gathered}
|A_p(x)|\le L^{\operatorname{def}_Z(x)}\,y^{\mathsf{s}}\,\bar{u}(x)\,\vartheta^{\iota^e_Z(x)}\\
\Bigl[\text{resp. }
|A_p(x)|\le L^{\operatorname{def}_Z(x)}\,C_{\mathbb{C}}^{\mathsf{s}}\,\bar{u}(x)\,
\vartheta^{\iota^e_Z(x)/2}\Bigr].
\end{gathered}
\end{equation}
\end{itemize}
\end{theorem}

\begin{proof}
\emph{Clause (a).} Let $x$ be a bond of a cell at which $v_D$ reads an increment. By the standing
hypotheses on block geometry, \eqref{9.5:eq-hypotheses-geometry-b} (see
Notation~\ref{9.5:not-hypotheses-geometry}), such a cell lies in $K_D^{\circ}$, and
$K_D^{\circ}\subset Z^{=}$. On $Z^{=}$ the deficit $\operatorname{def}_Z$ vanishes
(Lemma~\ref{9.5:lem-boundary-price}), so $\ell_Z=\ell_\iota$ and
$\iota^e_Z=\iota^e$ hold on $Z^{=}\supseteq K_D^{\circ}$. The mismatch field at $x$ is bounded by
the unweighted variation $\mathsf{S}_p$, and this in turn by the total variation $\mathsf{T}_p$ at
the local scale $\ell_Z$ of the box:
\begin{equation}\label{9.5:eq-mismatch-law-6}
|A_p(x)|\le \mathsf{S}_p(x)\le \frac{\mathsf{T}_p(x)}{\ell_Z(x)}
=\frac{\mathsf{T}_p(x)}{\ell_\iota(x)} .
\end{equation}
At real parameters we insert into \eqref{9.5:eq-mismatch-law-6} the bounds on the total variation
at bonds of $Z^{=}$ from two results. The first is Lemma~\ref{9.5:lem-no-accumulation}, which
applies at sites of type (P) along seed-entered runs. The second is
Lemma~\ref{9.5:lem-chain-yardstick}, which applies at sites of type (T) and (L) for the rows and
letters listed in (a), and at (P) for (L)-entered runs under its restriction (iii) with
$y_0=y^L C_T'$. Since $\iota^e_Z=\iota^e$ at $x$, and since the majorant of
Lemma~\ref{9.5:lem-chain-yardstick} is $v=\ell_\iota\bar{u}$, this gives the first inequality of
\eqref{9.5:eq-mismatch-law-3} with $y^{\mathsf{s}}=C^{P,\mathbb{Y}}$, $y^T$, $y^L$ respectively. On
$\mathfrak{P}^{\mathbb{Y}}$ the bound is that of Proposition~\ref{9.5:prop-complex-letters}, read
at the bond $x\in K_D^{\circ}\subset Z^{=}$. It is the bracketed inequality of
\eqref{9.5:eq-mismatch-law-3}.

\emph{Clause (b).} By the construction of the lettered family $\mathfrak{F}^{\mathbb{R},\mathbb{Y}}$,
at the bonds of the strip off $\bigcup_D\tilde{\square}$ both machines carry the seed configuration
$U^0$; so the configurations of $p$ and of its entry $p_0$ agree there, and $A_p=A_{p_0}$. Hence
$|A_p(x)|$ is bounded by the entry law of the entry $p_0$ of the run. For a seed, this is the entry
law for seeds, the first inequality of \eqref{9.5:eq-mismatch-law-4}. For a root of the second
entry type, it is the hypothesis \eqref{9.5:eq-mismatch-law-2}, which gives the second inequality of
\eqref{9.5:eq-mismatch-law-4}.

\emph{Clause (c).} At the bonds of $Z\setminus Z^{=}$ the argument of (a) is run at the local scale
$\ell_Z$ of the box and with the box-capped index $\iota^e_Z$. The price of this is displayed in
Lemma~\ref{9.5:lem-boundary-price}: by \eqref{9.5:eq-boundary-price-a} the passage from
$\ell_Z(x)^{-1}$ to $\ell_\iota(x)^{-1}$ costs at most the factor $L^{\operatorname{def}_Z(x)}$
(see also \eqref{9.5:eq-boundary-price-b}). This gives \eqref{9.5:eq-mismatch-law-5}.
\end{proof}

\begin{remark}[The conditions on the parameters and their order]\label{9.5:rem-parameter-order}
The parameters of the two-lattice comparison (Notation~\ref{9.5:not-comparison-parameters})
are subject to the floors and ceilings listed below and, in the order of choice given after
them, to one further condition on the radii $\alpha$ and $\varepsilon_1$. No other floor is
imposed, and each condition is imposed after $L$ has been fixed.

\emph{Floors.} First, $L\ge 5$ and $\theta>3/5$, together with the further floor that
accompanies them in Notation~\ref{9.5:not-comparison-parameters}. Second, the
floor on $\kappa_1$:
\begin{equation}\label{9.5:eq-parameter-order-1}
  \tfrac12(\kappa_1-1)\ \ge\ 2L\kappa .
\end{equation}
Third, one of two alternatives. Either the floor on $M_1$
\begin{equation}\label{9.5:eq-parameter-order-2}
  M_1\ \ge\ 128\,\delta/\delta_0 ,
\end{equation}
or the two floors on $M$
\begin{equation}\label{9.5:eq-parameter-order-3}
  \tfrac18\,\delta_1 M\ \ge\ (1+3\beta_I)\kappa+c_d,
  \qquad
  \delta_1 M\ \ge\ \kappa_1 .
\end{equation}
Fourth, the two conditions of the standing hypotheses on summation
(Notation~\ref{9.5:not-hypotheses-summation}): the floor on $L$ of those hypotheses, which
bounds $L$ from below by three times the square of the fixed base constant of
Ba{\l}aban's construction, and the bound $\beta\le \tfrac14$ on the exponent $\beta$.
Fifth, the radius condition: the letter radius $r_\square$ of
Definition~\ref{9.5:def-letter-profile} satisfies the condition $r_\square\ge 8C_\chi$
displayed in the first line of \eqref{9.5:eq-parameter-order-a} below, where
$C_\chi=2^dC_\zeta L$, so that $8C_\chi=2^{d+3}C_\zeta L$.
In clause~(T) of the comparison theorem with strip letters this is a lower bound on
$\check{R}$, of the same kind as the lower bound on $\check{R}$ among the conditions of the
strip-letter clause. In clause~(L) of that theorem it reads
\begin{equation}\label{9.5:eq-parameter-order-4}
  \varrho_L\, e^{\delta_1 M (N_0-2)/4}\ \ge\ 8C_\chi ,
\end{equation}
where $\varrho_L$ is the positive radius constant, not depending on $N_0$, of the
corresponding condition of the strip-letter clause at (L). This is
a lower bound on $N_0$ chosen after $M$ and $L$, of the same kind as that condition.

\emph{Ceilings.} These are smallness conditions, placed last and monotone in $\gamma$. The
numbers $\mathfrak{r}_0$, $\varepsilon_D$, $a_{R'}$, $\varepsilon_3$ are absolute in
$(d,L,\theta,M_1)$. With
\begin{equation}\label{9.5:eq-parameter-order-5}
  c_a\ =\ \Bigl[C_{R'}L^{\theta_2}C_T+\tfrac12\,C\,C_\chi^2\Bigr]\bar{\bar{u}},
\end{equation}
the radius condition and the ceilings (i)--(iii) read:
\begin{equation}\label{9.5:eq-parameter-order-a}
  \begin{gathered}
    r_\square\ \ge\ 8C_\chi ,\\
    \text{(i)}\ \ C_c\varepsilon+C_1'\bar{\bar{u}}\ \le\ \mathfrak{r}_0,\qquad
    \text{(ii)}\ \ (y^{\mathsf{s}}+1)\bar{\bar{u}}\ \le\ \varepsilon_D,\qquad
    \text{(iii)}\ \ 4C_\Sigma c_a\ \le\ 1 .
  \end{gathered}
\end{equation}

\emph{Order of choice.} The parameters are chosen in the following order:
\begin{enumerate}
  \item $L$;
  \item $\delta(L)$ and $\theta$;
  \item $\kappa$;
  \item $\kappa_1$, subject to the floor \eqref{9.5:eq-parameter-order-1};
  \item $M_1$, subject to \eqref{9.5:eq-parameter-order-2}, or $M$, subject to
        \eqref{9.5:eq-parameter-order-3};
  \item the radii $\alpha$ and $\varepsilon_1$, subject to the condition on $\alpha$ and
        $\varepsilon_1$ of the one-lattice response bounds;
  \item $\check{R}$ and $N_0$, from below, subject to the radius condition in the first line
        of \eqref{9.5:eq-parameter-order-a};
  \item $\gamma$, last, subject to the ceilings (i)--(iii) of
        \eqref{9.5:eq-parameter-order-a}.
\end{enumerate}
The ceiling $C_B(4+6\beta)\bar{\bar{u}}\le\log 2$ of the local growth lemma is not used here,
because the recursion for the total variation $\mathsf{T}_p$ is additive; the ceilings
(i)--(iii) are of the same kind.

The statements proved under these conditions have the following form. There is $L_0$, with
the floors $L\ge5$, $\theta>3/5$ and the accompanying floor of
Notation~\ref{9.5:not-comparison-parameters} built into its choice, such that, for every
$L\ge L_0$, every $\theta>3/5$, and every choice of the later parameters obeying the displayed
floors and
ceilings in the order just given, the conclusions hold with constants depending on
$(d,L,\theta,M_1,M)$ only. No condition bounds any of $L$, $M$, $M_1$, $\kappa$, $\kappa_1$,
$K$, $N$, $N_0$, $\check{R}$, $|\mathcal{D}^{+}_c|$, the number of links, the number of letters
or the length of a corridor from above, and no condition bounds a coupling from below.

Each condition is one-sided in the parameter chosen at its step of the order of choice, given
the parameters chosen before it.
\end{remark}

\begin{proof}
We take the stages in turn. The floor \eqref{9.5:eq-parameter-order-1} is equivalent to
$\kappa_1\ge 4L\kappa+1$, a lower bound on $\kappa_1$ in terms of $L$ and $\kappa$, which are
chosen earlier. The inequality \eqref{9.5:eq-parameter-order-2} bounds $M_1$ from below by a
quantity determined by $\delta$ and $\delta_0$, both fixed before $M_1$. Since $\delta_1>0$,
the two inequalities \eqref{9.5:eq-parameter-order-3} are equivalent to
\begin{equation*}
  M\ \ge\ \max\Bigl\{\,8\bigl((1+3\beta_I)\kappa+c_d\bigr)/\delta_1,\ \kappa_1/\delta_1\Bigr\},
\end{equation*}
a lower bound on $M$ in terms of parameters chosen before $M$. The conditions $L\ge5$,
the floor on $L$ of the summation hypotheses and $\theta>3/5$ are lower bounds. In
clause~(T) the radius condition is, by
its form, a lower bound on $\check{R}$. In clause~(L), since $\delta_1M>0$ and $\varrho_L>0$
does not depend on $N_0$, the
left-hand side of \eqref{9.5:eq-parameter-order-4} is increasing in $N_0$, while $8C_\chi$
depends on $d$, $C_\zeta$ and $L$ only. Hence \eqref{9.5:eq-parameter-order-4} holds for all
$N_0$ above a value determined by $M$, $L$, $\delta_1$, $\varrho_L$, $d$ and $C_\zeta$, which is
a lower bound on $N_0$
chosen after $M$ and $L$. Finally, the ceilings (i)--(iii) of
\eqref{9.5:eq-parameter-order-a} are imposed on $\gamma$ at the last step of the order of
choice. Their right-hand
sides $\mathfrak{r}_0$, $\varepsilon_D$ and $1$ are fixed before $\gamma$, and their left-hand
sides are monotone in $\gamma$. Each of them is therefore a one-sided condition on $\gamma$.
\end{proof}

\begin{remark}[Relative bound for dilated object rows; proof outstanding]
\label{9.5:rem-relative-object-bound}
On the lettered families, the proof of the comparison bounds for lettered links,
Theorem~\ref{9.5:thm-lettered-comparison}, goes through at every step at the site types
(P) and (T), and at (L) for rows of response type (Lemma~\ref{9.5:lem-response-rows}).
The first place at which it does not is not a step of that proof but the absence of a
comparison bound: at (L), for a row of the (L)-chain of object type that carries dilated
strip letters, no bound of that theorem covers the increment of the link.

Let $r$ be such a row. The link $p^{(r-1)}\mapsto p^{(r)}$ is a presented link: its
increment $\chi^{(r)}_t\mathcal{Z}^{(r)}$ lies in $\operatorname{Inc}_\beta$, and its
slot is given by the lettered slot rule. Theorem~\ref{9.5:thm-lettered-comparison}
therefore continues to apply to the later rows of response type, with $\Psi_p$
containing the $\mathsf{T}$-increment of this link. That increment, however, is not
shown to be small.

The comparison bound \eqref{9.5:eq-lettered-comparison-a} does not apply, because
$\mathcal{Z}^{(r)}$ is not a response of the current sets. The absolute route of
Lemma~\ref{9.5:lem-chain-yardstick}, part~(iii), uses two inputs: the composition
identity, and the absolute mismatch bound $|A_q|,|A_{q_{\mathrm{prev}}}|\le C_{LR'}u$
for the unlettered objects $q$, $q_{\mathrm{prev}}$ of the chain, taken here as a
hypothesis. With these it gives only
\begin{equation}\label{9.5:eq-relative-object-bound-1}
\mathsf{T}_{p^{(r)}}(x)-\mathsf{T}_{p^{(r-1)}}(x)
\le C\,\bar{\tau}_{\chi^{(r)}}(x)\, C_{LR'}\, v(x)\,\vartheta^{\iota^e(x)} .
\end{equation}
On the dilated strip blocks the profile satisfies
\begin{equation*}
\bar{\tau}_{\chi^{(r)}}(x)\asymp r^0_L\, e^{\delta_1 \operatorname{dist}(x,\bar{\Lambda})/8},
\end{equation*}
where $\operatorname{dist}(x,\bar{\Lambda})$ is the distance from $x$ to the source region
$\bar{\Lambda}$ of the (L)-family, measured as in the definitions of $r^0_L$ and
$\mathsf{m}_L$.
By the definition of $r^0_L$, this profile is bounded below by a $g$-free number times
$e^{\delta_1 M(N_0-2)/4}$. It is not bounded above by a $g$-free number.

The inequality that is missing is the following relative bound. Let a row of the
(L)-chain be of object type: its increment is a displacement $\mathcal{Z}_q$ to a
configuration re-solved as in \cite[(1.61)--(1.63)]{B-CMP122b}, entered as in clause
$(\ell 2)$ of the entry corollary. Primes mark the objects of the other run,
$p_{\mathrm{prev}}$ is the preceding link of the chain, and $Y[\,\cdot\,;\,\cdot\,]$ is
the map of the composition identity. The bound is required at real parameters, for every
box $Z$ of a unit $v_D$ that reads the row, and for every $x\in Z$:
\begin{equation}\label{9.5:eq-relative-object-bound-a}
\bigl|Y[\mathcal{Z}'_q;\mathbf{U}'_{\mathrm{prev}}]-\mathcal{Z}_q\bigr|_x
\le C_{\mathrm{rel}}\,\ell(x)\,\mathsf{m}_q(x)\,\vartheta^{\iota^e_Z(x)}
\bigl[c_b'+c_a'\,\mathsf{y}(p_{\mathrm{prev}})\bigr],
\end{equation}
where $\ell(x)$ is the local scale at $x$.
Here $\mathsf{m}_q$ is the majorant of the row's object, so that $\ell\,\mathsf{m}_q$
majorises $\ell\,|\mathcal{Z}_q|$; for a row of the (L)-chain it is the object majorant
\begin{equation*}
\mathsf{m}_L=C_4\,\mathsf{s}_L\,e^{-\delta_1 \operatorname{dist}(\cdot,\bar{\Lambda})/4}
\end{equation*}
of the (L)-family, with the constant $C_4$ and the decoupling parameters $\mathsf{s}_L$
of the statement that fixes that majorant.
The constants $C_{\mathrm{rel}}$ and $c_b'$ are free of $g$, and $c_a'\to 0$.

Proof outstanding.

Route. A proof would have to establish \eqref{9.5:eq-relative-object-bound-a} in the
relative form displayed. The inequality belongs among the bounds of the two-lattice
response theorem for Ba{\l}aban's objects of object type, together with the seam
statement. It must be proved there in relative form. The
absolute mismatch bound $|A_q|\le C_{LR'}u$ is the absolute form of
\eqref{9.5:eq-relative-object-bound-a}, and it does not imply it.

Suppose \eqref{9.5:eq-relative-object-bound-a} holds. Carry out the computation of
Lemma~\ref{9.5:lem-chain-yardstick}, part~(iii), with
\eqref{9.5:eq-relative-object-bound-a} in place of the absolute mismatch bound, and with
the inequality
\begin{equation*}
r^L_\square\,\ell\,\mathsf{m}_q\le \tfrac14\,\mathsf{r}^L ,
\end{equation*}
which holds by the definition of $r^0_L$. This computation bounds the
$\mathsf{T}$-increment of the lettered object link by a $g$-free multiple of
$v\,\vartheta^{\iota^e_Z}$. The clauses of Theorem~\ref{9.5:thm-mismatch-law} at the
site types (L) and (P) then extend to all rows and to all entries. This consequence is
stated here only under the hypothesis \eqref{9.5:eq-relative-object-bound-a}.
\end{remark}

\begin{definition}[The two factors kept inside a thrown-back region]\label{9.5:def-kept-factors}
Let $m$ be a step of the renormalisation group, that is, a level in the sense of the glossary
(written $k$ in \cite{B-CMP122b}; in this definition $m$ is not the torus exponent), and let $Z$ be
the large-field region on whose
components Ba{\l}aban's large-field operation $\mathbf{R}_m$ acts. By \cite[p.~378]{B-CMP122b},
these components are split by a decomposition of unity into two classes.
\begin{itemize}
\item[(a)] The first class $\{X_i\}$ consists of the components whose new field is small; they
are healed.
\item[(b)] The second class $\{Y_i\}$ consists of the components carrying the second, large-field
function. In the words of \cite[p.~378]{B-CMP122b}, they ``are included again into the large field
domain of the new action''. We call such a component $Y$ a \emph{thrown-back region} created at
step~$m$.
\end{itemize}
For a component $X$, \cite[(1.71)]{B-CMP122b} defines $\mathbf{T}'_m(X)$ as a sum over the
component's own sequences $\{\Omega_j\cap X,\ Z_j\cap X\}$, built from the small-field regions
$\Omega_j$ and the large-field regions $Z_j$ of \cite{B-CMP122b}. Each summand is an integral, over
the fluctuation field $B$ and the older integration variables, of an exponential whose exponent
contains a quadratic form in $B$ and the term $V(X^{\sim},U_m)$. By the sentence following
\cite[(1.71)]{B-CMP122b}, these two are the only ingredients of that definition that depend on the
background configuration $(U,J)$. By \cite[(1.102)]{B-CMP122b}, for a thrown-back region $Y$ the
operation $\mathbf{T}_m(Y)$ keeps $\mathbf{T}'_m(Y)$ in the density.

Consequently, as long as $Y$ is live, the density carries, under an integral sign, two functions of
the background on a neighbourhood of $Y$. These are the \emph{kept factors} of $Y$:
\begin{equation}\label{9.5:eq-kept-factors-a}
\mathfrak{K}(Y):=\bigl\{\,V(Y^{\sim}),\ \mathsf{Q}_Y\,\bigr\},
\qquad
\mathsf{Q}_Y:=-\tfrac{1}{2}\,\bigl\langle DH''_{1,m,Y}B,\ \zeta\, DH''_{1,m,Y}B\bigr\rangle .
\end{equation}
Here $V(Y^{\sim})$ is the term $V(X^{\sim},U_m)$ of \cite[(1.71)]{B-CMP122b} taken for $X=Y$,
and $\mathsf{Q}_Y$ is the quadratic form in $B$ of the same formula, taken for $X=Y$. The
operator $DH''_{1,m,Y}$ and the function $\zeta$ are those of \cite[(1.71)]{B-CMP122b}; this
$\zeta$ is Ba{\l}aban's letter and is distinct from the complex source $\zeta$ of the glossary. We
call $V(Y^{\sim})$ the \emph{kept potential} and $\mathsf{Q}_Y$ the \emph{kept quadratic form} of
$Y$. In the notation of the statement of this chapter that fixes $V^{\sharp}$, the kept potential
is written $V^{\sharp}(Y^{\sharp})$, with $Y^{\sharp}$ the region attached to $Y$ there; in the
notation of the bound of the quadratic form, the kept quadratic form is written $\mathfrak{q}_Y$.
\end{definition}

Both members of $\mathfrak{K}(Y)$ depend on the background, which every operation after step $m$
changes; so every later step at which $Y$ is still live must control them.

\cite[(1.104)]{B-CMP122b} gives an alternative representation of these terms. Of that
representation, \cite[(1.104)]{B-CMP122b} states: ``These functions depend only on the new field
variables $V_k$ restricted to $Y_i$, therefore they do not participate in operations connected with
next renormalization transformations''. The representation \cite[(1.104)]{B-CMP122b} and this
sentence are not used here; in the representation by \cite[(1.71), (1.102)]{B-CMP122b}, which is
the one used throughout, both members of $\mathfrak{K}(Y)$ are functions of the background.

The size of the kept potential is controlled by \cite[(1.69)]{B-CMP122b}: for $Y$ a component of
$Z^{\sim}$,
\begin{equation}\label{9.5:eq-kept-factors-1}
|V(Y)|\le O(1)\sum_j\bigl|\Gamma^{0}_j\cap Y\bigr| ,
\end{equation}
with the regions $\Gamma^{0}_j$ of \cite[(1.69)]{B-CMP122b}. The right-hand side of
\eqref{9.5:eq-kept-factors-1} is a volume counted in old cells, and it therefore grows with the
nesting depth. What a later step needs, however, is not the size of a kept factor but the
difference of two of its values taken at one and the same value of the inner integration variables.

\begin{definition}[Two runs on one scaffold and the two units of distance]\label{9.5:not-distance-units}
\emph{The two runs.} We compare two runs $\mathsf{A}$ and $\mathsf{B}$ of Ba{\l}aban's
renormalisation steps, of lengths $\mathfrak{n}$ and $\mathfrak{n}+1$ respectively, on one
scaffold $\Theta_X$, the system of regions, strips, cells and covers, listed below, that the
two runs share. Levels $\ell,j,m,k'$ are numbered by the index of the run $\mathsf{A}$,
and the depth of the level $\ell$ is $s(\ell):=\mathfrak{n}-\ell$. The zones, the sets
$Z_j:=\Lambda_j^c$, where $\Lambda_j$ is the small-field region of level $j$, and their
components are those of the scaffold. The two runs carry one label $\mathcal{L}$, and
$\mathcal{L}$ is \emph{aligned}: every zone present in it has an index $n\ge n_c$, and at the
index $n$ of every zone present the alignment condition
\begin{equation}\label{9.5:eq-distance-units-1}
C^{\natural}_{cl}\,(\log x_n)^{c_2}\,\vartheta_2^{\,n/4}\le 1
\end{equation}
holds; here $n_c$ is the threshold index of aligned labels and $C^{\natural}_{cl}$ is the
constant of the alignment condition.
A \emph{thrown-back region} $(Y,m)$ is a component $Y$ of the second class of the
large-field operation $\mathbf{R}$ at step $m$, that is, a component placed in the second of
the classes into which the operation $\mathbf{R}$ at step $m$ sorts the components on which it
acts. Its strip $Y^{\sharp}$ consists of
$m_{\sharp}:=\lfloor\frac12(\check{R}_m-1)\rfloor$ layers, where $\check{R}_m$ is the width
of a region created at step $m$. The region is live until the operation $\mathbf{R}$ is
completed on the component that contains it. All objects of the scaffold are the same for
the two runs: the thrown-back regions, the strips, the cells, the covers, the families of
cells, the index sets of the lists of units and of the pieces, the radii, $\check{R}$, the
two integers $\mathcal{N}$ and $\mathcal{N}_0$ attached to the scaffold, the distance
$d^{\wedge}$ and the far factors $\vartheta(\square)$ introduced below, and the constants
$K_A$ of the bound at the pair site.

\emph{Hypotheses of the comparison.} Both runs follow the scheme of the run used in the
uniqueness theorem, in which every kept term is read through its hull, by the hull-averaging
read map $\mathrm{rd}^{\mathsf{m}}$, and in which the pair site is taken at the pair radii
$r^{\bullet}_\square$ defined below. We
assume
\begin{equation}\label{9.5:eq-distance-units-c}
\begin{gathered}
\text{the run hypotheses and the one-family dictionary stated below}\\
\text{hold in the run }\mathsf{m},\quad \mathsf{m}\in\{\mathsf{A},\mathsf{B}\}.
\end{gathered}
\end{equation}
The \emph{run hypotheses} are these: in the run $\mathsf{m}$, at every step, the hypotheses
of the one-run bound on the kept potential $V^{\sharp}(Y^{\sharp})$ hold; so do the
hypotheses of the bounds on the kept terms at the later sites, the hypotheses under which the
kept quadratic form $\mathsf{Q}_Y$ of Definition~\ref{9.5:def-kept-factors} is defined and
bounded, and the six conditions of the statement that bounds the kept terms of the run; each
of them is taken with the decay rates of the run $\mathsf{m}$ substituted for the rates in its
statement. Consequently, by the one-run bound on the kept potential, its conclusion holds in
each run $\mathsf{m}$. The \emph{one-family
dictionary} consists of the following three sentences.
\begin{enumerate}
\item[(a)] The configurations presented to a unit at the \textup{(L)}-site, namely the
restriction $\mathrm{chain}(m)\upharpoonright X$ to the region $X$ of the chain
$\mathrm{chain}(m)$ of configurations at step $m$, and the objects $\mathfrak{Q}''$ of the
interrupted run, are the same objects at natural vertices and at real data, their holomorphic
continuations in the seed and in the decoupling parameters $\vec{\mathsf{s}}$ agree, and the
core $\mathsf{C}$ lies in the union of the region $Z''_m$ of step $m$ of the run of the
uniqueness theorem and the region carrying the kept factors; the seed and the natural
vertices are those of that run.
\item[(b)] The vertices $\vec{\mathsf{s}}\in\{0,1\}$ of the families at the
\textup{(L)}-, \textup{(P)}- and \textup{(T)}-sites are natural vertices.
\item[(c)] For a term over strips, the tree domain $X^{\theta}$ of a term of the first kind,
with its distances $d^{\theta}_j(X)$, and the label $A_F$ of a term of the second kind, with
its distances $d_j(A_F)$, denote the same set and the same distances; these are level-$j$
distances of the kinds $d^{\theta}(X)$ and $d_j(X)$ of Notation~\ref{0.2:not-glossary}.
\end{enumerate}

\emph{The two units of distance.} The correlation theorem measures distance in units of
$M_1$. For the label $\mathcal{L}$ and a contour $\Gamma$, in the sense of
\cite[(2.46)]{B-CMP96} with or without the admissibility condition there,
\begin{equation*}
\operatorname{len}_{\mathcal{L}}(\Gamma):=\sum_i \mathsf{u}_i^{-1}\,
\bigl|\Gamma\cap(\text{zone }i)\bigr|,
\end{equation*}
where $\mathsf{u}_i$ is the $M_1$-unit of level $i$ (so that a cube of side $M$ on its own
scale has side $\mathfrak{w}:=M/M_1$ in these units), and for sets $A,B$ the distance
$d^{\wedge}(A,B)$ is the infimum of $\operatorname{len}_{\mathcal{L}}(\Gamma)$ over all
contours $\Gamma$ joining $A$ to $B$. The uniqueness theorem measures distance in lattice
units: as in \cite[(2.46)]{B-CMP96},
\begin{equation*}
\operatorname{len}(\Gamma):=\sum_j (L^j\eta)^{-1}\,\bigl|\Gamma\cap\{\iota=j\}\bigr|,
\end{equation*}
where $\eta$ is the lattice spacing of \cite[(2.46)]{B-CMP96} (not the read data
$\eta^{\mathsf{m}}_v$) and $\iota$ is the level function, equal to $j$ on the zone of level
$j$; for sets $A,B$, $d(A,B)$ is the infimum of $\operatorname{len}(\Gamma)$ over the
contours $\Gamma$ joining $A$ to $B$ that are admissible in the sense of
\cite[(2.46)]{B-CMP96}. For a set $\mathfrak{S}$ refining the label,
$\operatorname{len}_{\mathfrak{S}}$ denotes the length computed with the zones of
$\mathfrak{S}$ in place of those of the label, and $d_{\mathfrak{S}}(A,B)$ the infimum of
$\operatorname{len}_{\mathfrak{S}}(\Gamma)$ over the contours joining $A$ to $B$ admissible
for $\mathfrak{S}$; $\operatorname{len}_{\mathfrak{S}}=\operatorname{len}$ and
$d_{\mathfrak{S}}=d$ when $\mathfrak{S}$ is the label. In lattice units, a cube of side
$M_1$ has $d$-width at least $M_1$ on its own scale. The $M_1$-unit of a level being
$M_1$ lattice units of that level, the two lengths of one contour differ by the factor
$M_1$ exactly: $\operatorname{len}(\Gamma)=M_1\operatorname{len}_{\mathcal{L}}(\Gamma)$.
An admissible contour is a contour, and, by the monotonicity of the length of a contour under
refinement of the label, passing to a set refining the label does not decrease the length of
a contour. We show below that, for all sets $A,B$ and every set
$\mathfrak{S}$ refining the label,
\begin{equation}\label{9.5:eq-distance-units-a}
M_1\,d^{\wedge}(A,B)\le d_{\mathfrak{S}}(A,B).
\end{equation}

\emph{Rates.} We put $\delta_1:=\delta_0/32$ and $\delta_P:=\frac18\delta_0$, so that
$\delta_P=4\delta_1$ as numbers, while $d^{\wedge}$ is measured in units of $M_1$ and $d$ in
lattice units, the two being related by \eqref{9.5:eq-distance-units-a}. We use the statement
that all kernels of Ba{\l}aban's construction used in this chapter decay at rate at least
$4\delta_1$ in the distance $d$. For a function $f$ and a field $b$,
\begin{equation*}
f^{\sim}(x):=\sum_y e^{-\delta_1 d(x,y)}f(y),\qquad \mathsf{P}_b:=|b|^{\sim}.
\end{equation*}
The dilation disc of a region $X$ is the disc $|\lambda|<\rho_X$ of the dilation parameter
$\lambda\in\mathbb{C}$ (not to be confused with the flow constants
$\lambda^{\sharp},\underline{\lambda},\lambda_U$), with
\begin{equation}\label{9.5:eq-distance-units-2}
\rho_X:=r^0\,e^{\frac12\delta_P d^{\wedge}(X,\mathsf{C}^c)},
\end{equation}
where $r^0$ is the base radius, the value of $\rho_X$ when $d^{\wedge}(X,\mathsf{C}^c)=0$
(not the depth offset $r_0$ of the glossary);
the exponent is $\frac12\delta_P=2\delta_1$. The second part of the cell bound is used
at decay rates $\beta\ge\delta_P/16=\delta_1/4$ (here $\beta$ is a rate, not the coupling
increment $\beta_{k+1}$); in the one-run bound on the kept potential it
is used at $\beta=\frac14\delta_P\ge\delta_P/16$. In the estimate of the plain list the
rate $\frac14\delta_P$ is divided as $\frac18\delta_P+\frac18\delta_P$: one half is
absorbed through the floor $\frac18\delta_P \mathfrak{w}\ge 2$, a one-sided condition on $M$
to be met after $M_1$ is fixed, the other is spent in the
second part of the cell bound at rate $\frac18\delta_P$.

\emph{The floor on $M_1$.} We impose on $M_1$, which is chosen in
Notation~\ref{2.3:not-stages-middle}, the floor, in the sense of
Notation~\ref{2.3:not-stages-first}, to be met after $L$ is fixed and before $M$; it is
one-sided:
\begin{equation}\label{9.5:eq-distance-units-b}
\delta_1 M_1\ge 4\delta_P .
\end{equation}
With $\delta_P=\frac18\delta_0$ and $\delta_1=\delta_0/32$ it reads $M_1\ge16$; we keep it in
the form \eqref{9.5:eq-distance-units-b}, which is the form used below.

\emph{Far factors and pair circles.} Let $R_\nu$ be the region of the correlation theorem
from which the far factors are measured (not the ball radius $R$). The far factor of the
correlation theorem at a cube $\square$ is
$\vartheta(\square)=\vartheta\,e^{-\frac12\delta_P d^{\wedge}(\square,R_\nu)}$, and
$r_\square$ are the circles of that theorem built from these factors. The pair factor is
\begin{equation*}
\vartheta^{\bullet}_P(\square):=\max\Bigl\{\vartheta_P(\square),\;
\vartheta^0_P\sup_{x\in\square}e^{-\delta_1 d(x,R_\nu)/8}\Bigr\},
\end{equation*}
where $\vartheta_P(\square)$ is the factor of the cube $\square$ at the pair site,
$\vartheta^0_P$ is its base constant and the supremum is over the points $x$ of $\square$;
the pair circles, of radii $r^{\bullet}_\square$, are built from
$\vartheta^{\bullet}_P(\square)$. Under
\eqref{9.5:eq-distance-units-b}, for all sets $A,B$,
\begin{equation}\label{9.5:eq-distance-units-3}
\begin{gathered}
\tfrac12\delta_P\, d^{\wedge}(A,B)\le \tfrac12\delta_P\,d(A,B)/M_1\le \delta_1\, d(A,B)/8,\\
\delta_P\, d^{\wedge}(A,B)\le\tfrac14\delta_1\, d(A,B).
\end{gathered}
\end{equation}
Consequently, for every cube $\square$,
$\vartheta(\square)\ge\vartheta e^{-\delta_1 d(\square,R_\nu)/8}$ and
$\vartheta^{\bullet}_P(\square)=\vartheta_P(\square)$; hence the pair circles
$r^{\bullet}_\square$ are the circles $r_\square$ of the correlation theorem, and the table of
the unit shares $\theta_P,\theta_T,\theta^{\circ}_L,\theta^{\dagger}_L$ of the kept terms at
the four sites holds without change on the run of the uniqueness theorem.
\end{definition}

\begin{proof}
We prove \eqref{9.5:eq-distance-units-a}. Let $A,B$ be sets, let $\mathfrak{S}$ be a set
refining the label and let $\Gamma$ be a contour joining $A$ to $B$ that is admissible for
$\mathfrak{S}$. Passing
from the label to $\mathfrak{S}$ does not decrease the length of a contour, by the
monotonicity of the length under refinement of the label, so
$\operatorname{len}_{\mathfrak{S}}(\Gamma)\ge\operatorname{len}(\Gamma)$. The two lengths of
one contour differ by the factor $M_1$ exactly, so
$\operatorname{len}(\Gamma)=M_1\operatorname{len}_{\mathcal{L}}(\Gamma)$. An admissible
contour is a contour, and $d^{\wedge}(A,B)$ is the infimum of
$\operatorname{len}_{\mathcal{L}}$ over all contours joining $A$ to $B$, so
$\operatorname{len}_{\mathcal{L}}(\Gamma)\ge d^{\wedge}(A,B)$. Together,
\begin{equation*}
\operatorname{len}_{\mathfrak{S}}(\Gamma)\ge\operatorname{len}(\Gamma)
=M_1\operatorname{len}_{\mathcal{L}}(\Gamma)\ge M_1 d^{\wedge}(A,B).
\end{equation*}
Taking the infimum over all contours joining $A$ to $B$ admissible for $\mathfrak{S}$ gives
$d_{\mathfrak{S}}(A,B)\ge M_1d^{\wedge}(A,B)$, which is \eqref{9.5:eq-distance-units-a}. Only
the geometry of contours has been used.

Assume now \eqref{9.5:eq-distance-units-b}, and let $A,B$ be sets. By
\eqref{9.5:eq-distance-units-a} with $\mathfrak{S}$ the label,
$d^{\wedge}(A,B)\le d(A,B)/M_1$; multiplying by $\frac12\delta_P$ gives the first inequality
of the chain in the first line of \eqref{9.5:eq-distance-units-3}. Dividing
\eqref{9.5:eq-distance-units-b} by $8M_1$ gives $\frac12\delta_P/M_1\le\delta_1/8$, and
multiplying by $d(A,B)\ge0$ gives the second inequality of that chain. Doubling the resulting
inequality $\frac12\delta_Pd^{\wedge}(A,B)\le\delta_1d(A,B)/8$ gives the second line of
\eqref{9.5:eq-distance-units-3}.
Applied with $A=\square$ and $B=R_\nu$, the chain in the first line of
\eqref{9.5:eq-distance-units-3} gives
\begin{equation*}
\vartheta(\square)=\vartheta\,e^{-\frac12\delta_P d^{\wedge}(\square,R_\nu)}
\ge\vartheta\,e^{-\delta_1 d(\square,R_\nu)/8}.
\end{equation*}
Since the distance of the cube $\square$ to $R_\nu$ is the infimum over the points
$x\in\square$ of $d(x,R_\nu)$ (a contour joining $\square$ to $R_\nu$ is a contour joining
some point $x\in\square$ to $R_\nu$, and conversely, its admissibility being a condition on the
contour itself), the right side equals
$\vartheta\sup_{x\in\square}e^{-\delta_1 d(x,R_\nu)/8}$.
Hence the far factor dominates the exponential entering the second member of the maximum
that defines $\vartheta^{\bullet}_P(\square)$, the second member does not exceed the first,
and the maximum is $\vartheta_P(\square)$ itself. The pair circles $r^{\bullet}_\square$,
built from $\vartheta^{\bullet}_P(\square)=\vartheta_P(\square)$, are therefore the circles
$r_\square$ of the correlation theorem, and the table of unit shares at the four sites,
which reads the circles $r_\square$, holds without change on the run of the uniqueness
theorem.
\end{proof}

\begin{definition}[Common data, charts and the two-point oscillation]
\label{9.5:not-two-point-oscillation}
Let $Y$ be a thrown-back region whose kept factors are outputs of step $m$, let $Y^{\sharp}$ be
its strip, and let $\mathfrak{K}(Y)$ be the pair of kept factors of
Definition~\ref{9.5:def-kept-factors}, display~\eqref{9.5:eq-kept-factors-a}. The two runs
compared are $\mathsf{A}$ and $\mathsf{B}$; the letter $\mathsf{m}$ stands for either of them. We
write $\mathbf{R}_m$ for Ba{\l}aban's large-field operation $\mathbf{R}$ at step $m$ and $\pi_m$ for
the paving of the lattice by the cubes of level $m$.

\emph{Storage of the kept factors.} We adopt the following convention. The two members of
$\mathfrak{K}(Y)$ are stored outputs of $\mathbf{R}_m$, with slice region
\[
(Y^{\sharp})^{+} := Y^{\sharp} \text{ together with all } \pi_m\text{-cubes }
\infty\text{-adjacent to } Y^{\sharp},
\]
and every later production of terms, in either run, reads them through the hull: for
$v\in\mathfrak{K}(Y)$,
\begin{equation}\label{9.5:eq-two-point-oscillation-a}
\begin{gathered}
\hat{v}^{\mathsf{m}}(\xi;\Phi_Y)
:= v^{\mathsf{m}}\bigl((U^{\mathsf{m}}_m,J^{\mathsf{m}}_m)((Y^{\sharp})^{+},\xi)
\upharpoonright Y^{\sharp};\Phi_Y\bigr),\\
\tilde{v}^{\mathsf{m}}(P;\Phi_Y)
:= \hat{v}^{\mathsf{m}}\bigl(\mathrm{rd}^{\mathsf{m}}_{(Y^{\sharp})^{+}}P;\Phi_Y\bigr),
\qquad
\mathrm{rd}^{\mathsf{m}}_{(Y^{\sharp})^{+}}
:= M_{\mathfrak{B}^{\mathsf{m}}_m((Y^{\sharp})^{+})}.
\end{gathered}
\end{equation}
Here $(U^{\mathsf{m}}_m,J^{\mathsf{m}}_m)(Z,\xi)$ is the configuration that run $\mathsf{m}$
determines at level $m$ from a datum $\xi$ on the set $Z$, $\mathfrak{B}^{\mathsf{m}}_m(Z)$ is the
hull of $Z$ at level $m$ in run $\mathsf{m}$, $M_{\mathfrak{B}}$ is the averaging over the hull
$\mathfrak{B}$, so that $\mathrm{rd}^{\mathsf{m}}$ is the hull-averaging read map, and $\Phi_Y$ are
the inner variables fixed next. This convention is a choice of the scaffold of the run and asserts
nothing; it leaves the site identities and the density unchanged, as shown in the proof below.

\emph{Inner variables.} $\Phi_Y$ denotes all real variables that the members of $\mathfrak{K}(Y)$
read besides the background. They are integration variables common to both runs (the core
variables of \cite{B-CMP122b}), and they are always taken at one value, lying in the intersection
$\operatorname{supp}^{\mathsf{A}}\cap\operatorname{supp}^{\mathsf{B}}$ of their supports in the two
runs.

\emph{Data domains.} For a term $F$ with data domain $\mathfrak{D}^{\mathsf{m}}_F$ and
$0\le\varsigma\le 1$, $\mathfrak{D}^{\mathsf{m}}_F[\varsigma]$ is the sub-domain of
$\mathfrak{D}^{\mathsf{m}}_F$ on which every entry uses at most the fraction $\varsigma$ of its
room $\varrho_F$. Thus
\begin{equation}\label{9.5:eq-two-point-oscillation-1}
\mathfrak{D}^{\mathsf{m}}_F[0]\subset\mathfrak{D}^{\mathsf{m}}_F[\tfrac18]
\subset\mathfrak{D}^{\mathsf{m}}_F[\tfrac14]\subset\mathfrak{D}^{\mathsf{m}}_F[\tfrac78];
\end{equation}
in the parametrisation of the data domains by a radius parameter that the two-run comparison uses,
the domain $\mathfrak{D}^{\mathsf{m}}_F[\tfrac14]$ is the one of parameter
$\mathsf{c}-\tfrac34\varrho_F$, where $\mathsf{c}$ is the reference value of the radius parameter
in that parametrisation. We write $\mathfrak{T}^{\mathsf{m}}_F[\varsigma]$ for the set of
data on the hull's set whose slice lies in $\mathfrak{D}^{\mathsf{m}}_F[\varsigma]$; the slice of a
datum $\xi$ in run $\mathsf{m}$ is written $\operatorname{slice}^{\mathsf{m}}(\xi)$. For
$v\in\mathfrak{K}(Y)$ the data domain is
$\mathfrak{D}^{\mathsf{m}}_v := \tilde{U}^{c,\mathsf{m}}_m(Y^{\sharp},\tilde{\alpha}_0,
\tilde{\alpha}_1)$, the complex domain of configurations on $Y^{\sharp}$ at level $m$ in run
$\mathsf{m}$ with radii $\tilde{\alpha}_0,\tilde{\alpha}_1$, and the rooms are the generic rooms
$\varrho_F$ with the level index set to $m$. For the kept potential this is the standard domain of
a potential $V(X^{\sharp})$ kept by the operation $\mathbf{T}'$ of \cite[(1.102)]{B-CMP122b} and
created by $\mathbf{R}$; for the kept quadratic form
$\mathsf{Q}_Y$ we define the domain and the rooms to be the same.

\emph{Common data.} For a set $X$ and a term $F$ with data on $X$,
\begin{equation}\label{9.5:eq-two-point-oscillation-2}
\mathfrak{X}_X := \bigl\{\xi=(\mathsf{v},c):\ \mathsf{v}\in\mathfrak{T}^{\mathsf{A}}_F[\tfrac14],\
\mathsf{v}\oplus c\in\mathfrak{T}^{\mathsf{B}}_F[\tfrac14]\bigr\},
\end{equation}
where $c$ is a datum on the collar $\Lambda'_0$ of run $\mathsf{B}$. The pair
$(\mathsf{v},\mathsf{v}\oplus c)$ of the data that $\xi$ gives to the runs $\mathsf{A}$ and
$\mathsf{B}$ is the \emph{seed pair} of $\xi$. The set $\mathfrak{X}_X$ is contained in the domain of
the seeds, and on it the stored terms are defined and bounded (see the proof).

\emph{Charts.} For a configuration $P$ and a chart $\mathfrak{c}\in\mathfrak{G}_{\varepsilon}
(\Sigma_Y)$, let $v^{\sharp}(P;\Phi,\mathfrak{c})$ be the kept factor $v$ read at $P$ in the chart
$\mathfrak{c}$. At slices we put
\[
\hat{v}^{\mathsf{m}\sharp}(\xi;\Phi,\mathfrak{c})
:= v^{\mathsf{m}\sharp}\bigl(\operatorname{slice}^{\mathsf{m}}(\xi)\upharpoonright
Y^{\sharp};\Phi,\mathfrak{c}\bigr).
\]
The chart $\mathfrak{c}$, like every other argument of the terms compared between the two runs, is
a parameter common to both runs. We put
$\mathfrak{X}^{\sharp}_Y := \mathfrak{X}_{Y^{\sharp}}\times\mathfrak{G}_{\varepsilon}(\Sigma_Y)$,
and we write $\hat{V}^{\mathsf{m}\sharp}$ and $\hat{\mathfrak{q}}^{\mathsf{m}\sharp}$ for the reads
in this sense of the two members $V^{\sharp}(Y^{\sharp})$ and $\mathfrak{q}_Y=\mathsf{Q}_Y$ of
$\mathfrak{K}(Y)$.

\emph{The two-point oscillation.} For a function $f$ of $(\xi,\mathfrak{c})\in
\mathfrak{X}^{\sharp}_Y$ and of $\Phi_Y$ we define $\mathfrak{O}^{+}(f)$ by the first line of the
following display; its second line is a bound proved below:
\begin{equation}\label{9.5:eq-two-point-oscillation-b}
\begin{gathered}
\mathfrak{O}^{+}(f) := \sup_{\Phi_Y\ \mathrm{real}}\ \
\sup_{(\xi,\mathfrak{c}),(\xi',\mathfrak{c}')\in\mathfrak{X}^{\sharp}_Y}
\bigl|f(\xi,\mathfrak{c};\Phi_Y)-f(\xi',\mathfrak{c}';\Phi_Y)\bigr|,\\
\mathfrak{O}^{+}\bigl(\hat{V}^{\mathsf{m}\sharp}\bigr)
+\mathfrak{O}^{+}\bigl(\hat{\mathfrak{q}}^{\mathsf{m}\sharp}\bigr)
\le C_{\Sigma}\check{R}_m^{\,p_{\Sigma}}\qquad(\mathsf{m}=\mathsf{A},\mathsf{B}).
\end{gathered}
\end{equation}
Thus $\mathfrak{O}^{+}$ compares two points at one value of $\Phi_Y$; it is a seminorm, and it
vanishes on functions that do not depend on $(\xi,\mathfrak{c})$.

\emph{Two charts at one value of the inner variables.} In each run,
\begin{equation}\label{9.5:eq-two-point-oscillation-c}
v^{\sharp}\bigl(P;\Phi,e^{Y}w\bigr)
= v\bigl((P^{\rho_{\mathsf{r}}[Y]})^{\hat{w}^{-1}};\Phi\bigr),
\end{equation}
where $e^{Y}w\in\mathfrak{G}_{\varepsilon}(\Sigma_Y)$ is a chart, written through an element $w$
at the law sites $\Sigma_Y$ (this $w$ is not the complex source of the glossary), $\hat{w}$ is the
gauge transformation of configurations on $Y^{\sharp}$ that $w$ induces, $\rho_{\mathsf{r}}[Y]$ is
the rotation attached to $Y$ through its reading mark $\mathsf{r}_Y$, and $P^{\rho}$ denotes the
configuration $P$ transformed by $\rho$. Consequently a two-run difference at a common chart equals
the two-run
difference, at the trivial chart and at the same $\Phi$, at the pair of tube points obtained from
the seed pair by one common map $\rho_{\mathsf{r}}[Y]\hat{w}^{-1}$; the two-run estimates that
follow are written at the trivial chart and are applied at the regauged seed pairs.

\emph{Inputs.} The two-run comparison that uses these data rests on the following statements of
this book:
\begin{enumerate}
\item[(a)] the theorem on the field law and on the resolved extensions of stored terms (all its
parts), with its companion statements: the load statement, the corollary on groups of terms, the
class statements, the leaf lemma, and the further lemmas that accompany that theorem;
\item[(b)] the definition of stored derived terms with their store properties, and the lemmas that
accompany it, among them the elimination of the boundary terms and the consistency of the stores;
\item[(c)] the theorem on terms over live strips and its corollary, under the strip hypothesis,
together with the one-run expansions on labels;
\item[(d)] the corollary on reads in rooms (its parts on the rooms of the reads, on the dilated
reads and on the seeds), the lemma on reads of the kept factors in rooms, and the mixing lemma,
with the one-run statements on which these depend; the proof of the mixing lemma does not use
that the terms live over strips, and that proof is applied here at the deep terms of the site
$(\mathrm{L})$ of step $m$ and at the kept factors;
\item[(e)] the classical operator bounds (the Wilson operator, the covariances and the forms,
the potentials free of the complex sources) and the proposition on their relative bounds at the
common parameter points;
\item[(f)] the proposition on the terms of an interrupted run, in which the kept Gaussian is
excluded;
\item[(g)] geometric forgetting of old regions and its corollary, applied at the components that
hold a thrown-back region together with its chart; the accounting statement for such a component,
whose bracketed term is kept through the re-creation of the region; and the statement on the
contributions entering from the potentials kept by $\mathbf{T}'$;
\item[(h)] the corollary on pairs of runs, by which the strip bounds hold for differences and for
marks.
\end{enumerate}
\end{definition}

\begin{proof}
\emph{The storage convention leaves the site identities and the density unchanged.} The site
identities are algebraic identities among the terms; by part~(a) of the lemma on terms read
through the hull, each of them therefore holds with $\tilde{v}^{\mathsf{m}}$ of
\eqref{9.5:eq-two-point-oscillation-a} in place of $v^{\mathsf{m}}$. By part~(e) of the statement
on the extension of hull reads off the critical slice, $\tilde{v}^{\mathsf{m}}(P;\Phi_Y)=
v^{\mathsf{m}}(P;\Phi_Y)$ at every critical configuration $P$. Hence the density is not changed by
the convention;
only the extension of the kept factors off the real critical slice is, and no family of terms is
altered.

\emph{Common data.} A datum using at most a quarter of the room of each entry uses at most seven
eighths of it, so $\mathfrak{T}^{\mathsf{m}}_F[\tfrac14]\subset\mathfrak{T}^{\mathsf{m}}_F
[\tfrac78]$ by \eqref{9.5:eq-two-point-oscillation-1}. By part~(d) of the corollary on reads in
rooms, the seeds of a term $\Phi$ fill $\mathfrak{T}_{\Phi}[\tfrac78]$. Therefore both entries of
the seed pair of every $\xi\in\mathfrak{X}_X$ in \eqref{9.5:eq-two-point-oscillation-2} lie in the
domain of the seeds, which is the first assertion. By the first store property of stored derived
terms (input~(b)), at every common datum each stored term is defined, bounded and holomorphic in
$\mathsf{v}$ on the tube of the term; this is the second assertion.

\emph{$\mathfrak{O}^{+}$ is a seminorm vanishing on functions free of $(\xi,\mathfrak{c})$.} Both
points in \eqref{9.5:eq-two-point-oscillation-b} are taken at the same $\Phi_Y$. For functions
$f,g$ and at one $\Phi_Y$,
\[
\bigl|(f+g)(\xi,\mathfrak{c})-(f+g)(\xi',\mathfrak{c}')\bigr|
\le\bigl|f(\xi,\mathfrak{c})-f(\xi',\mathfrak{c}')\bigr|
+\bigl|g(\xi,\mathfrak{c})-g(\xi',\mathfrak{c}')\bigr|,
\]
and taking suprema gives $\mathfrak{O}^{+}(f+g)\le\mathfrak{O}^{+}(f)+\mathfrak{O}^{+}(g)$; for a
complex number $\lambda$ the differences of $\lambda f$ are $\lambda$ times those of $f$, so
$\mathfrak{O}^{+}(\lambda f)=|\lambda|\,\mathfrak{O}^{+}(f)$. If $h(\xi,\mathfrak{c};\Phi_Y)=
h_0(\Phi_Y)$ does not depend on $(\xi,\mathfrak{c})$, every difference in the definition vanishes,
so $\mathfrak{O}^{+}(h)=0$, and by the triangle inequality applied twice $\mathfrak{O}^{+}(f+h)=
\mathfrak{O}^{+}(f)$.

\emph{The bound in \eqref{9.5:eq-two-point-oscillation-b}.} Fix a run $\mathsf{m}$. The bound
follows from parts~(v) and~(t) of the base theorem for the kept potential, applied in run
$\mathsf{m}$ under its hypotheses in that run, from the bound of the kept quadratic form, and from
the description of the kept factors in a chart, applied in run $\mathsf{m}$.

\emph{The identity \eqref{9.5:eq-two-point-oscillation-c} and its consequence.} In each run the
identity \eqref{9.5:eq-two-point-oscillation-c} is the description of the kept factors in a chart,
combined with the statement that accompanies that description. We apply it with the same pair
$(Y,w)$ in both
runs. For configurations $P^{\mathsf{A}},P^{\mathsf{B}}$ read from the seed pair of a common datum
and one value $\Phi$ of the inner variables this gives
\begin{align*}
&v^{\mathsf{B}\sharp}\bigl(P^{\mathsf{B}};\Phi,e^{Y}w\bigr)
-v^{\mathsf{A}\sharp}\bigl(P^{\mathsf{A}};\Phi,e^{Y}w\bigr)\\
&\qquad=v^{\mathsf{B}}\bigl(((P^{\mathsf{B}})^{\rho_{\mathsf{r}}[Y]})^{\hat{w}^{-1}};\Phi\bigr)
-v^{\mathsf{A}}\bigl(((P^{\mathsf{A}})^{\rho_{\mathsf{r}}[Y]})^{\hat{w}^{-1}};\Phi\bigr),
\end{align*}
that is, the two-run difference at the common chart $e^{Y}w$ is the two-run difference at the
trivial chart and the same $\Phi$, at the pair of tube points obtained from the seed pair by the one
common map $\rho_{\mathsf{r}}[Y]\hat{w}^{-1}$. The production of terms is covariant under such maps,
by part~(g) of the base theorem for the kept potential and by the covariance statements for the
production of terms: the maps act on the arguments of the terms as common
rotations and leave the invariant terms unchanged. The field law is covariant as well: its
coefficients $A_p$ are carried to $\operatorname{Ad}A_p$, the adjoint action of the common map, by
the frame rule of the field law.
Hence an estimate proved for the two-run difference at the trivial chart, at all pairs of tube
points arising from seed pairs, applies at the regauged seed pairs, and so at every common chart;
this is how the two-run estimates that follow are used.
\end{proof}

\begin{theorem}[Layer-by-layer smallness of the kept-factor difference]
\label{9.5:thm-kept-difference-layers}
Let $\mathsf{A}$ and $\mathsf{B}$ be the two runs of the uniqueness theorem, compared on one scaffold
and on aligned labels as in Notation~\ref{9.5:not-distance-units}, and let the auxiliary
parameters satisfy the one-sided lower bounds and the upper bounds on $\gamma$ displayed in this
section, among them those of \eqref{9.5:eq-kept-elimination-a}. Assume:
\begin{enumerate}
\item[(i)] the three dictionary statements of \eqref{9.5:eq-distance-units-c} (its parts (a),
(b), (d));
\item[(ii)] the one-run theory of the kept factors in each of the two runs, in the form displayed in
\eqref{9.5:eq-distance-units-c};
\item[(iii)] the four inputs of the uniqueness theorem on which the comparison of $\mathsf{A}$ with
$\mathsf{B}$ rests, each by its statement;
\item[(iv)] four statements: the reading law over strips, with the one-run expansions on labels
and the mixing statement; the reading law at terms that meet a healed core, together with its
supplementary clause and with the containment of the classical terms at complex parameters at the
last site of a step; the classical operator bounds; and the bounds on the terms
$\mathfrak{F}^{B4}(\mathfrak{a})$ and the Gaussians $\mathfrak{Q}(\mathfrak{a})$ of an interrupted
run, with the corresponding one-run bound for the kept Gaussian $\mathsf{q}_W$;
\item[(v)] geometric forgetting of old regions (the un-normalised age weight), in both of its
forms, together with the two side conditions attached to it at components that contain a
thrown-back region, and Ba{\l}aban's waiting-time condition on treated regions
\cite[condition~(ii)]{B-CMP122a};
\item[(vi)] the two-run difference of field-independent remainder terms,
\eqref{9.5:eq-remainder-difference-a};
\item[(vii)] six lemmas stated for a narrower class of terms, read with their proofs on the class
of terms used here; each is named at the place in the proof below where it is used.
\end{enumerate}
Then for every thrown-back region $Y$, with its step $m$ as in
Definition~\ref{9.5:def-kept-factors}, and every kept factor $v\in\mathfrak{K}(Y)$ of that
definition (see \eqref{9.5:eq-kept-factors-a}) there are functions $\delta v^{(n)}$, indexed by the
$n\le m$ that are zones of the history of the label in $Y^{\sharp}$ and by $n=m$, holomorphic in
the common datum $\xi=(\mathsf{v},c)$ of $(Y^{\sharp})^{+}$, such that the statements (1)--(7)
below hold. The relations referred to in them are collected in
\eqref{9.5:eq-kept-difference-layers-a}, in which $\vartheta_f:=\vartheta^{1/512}$;
$\|\cdot\|_{\natural}$ is the norm of the resolved-extension bookkeeping (the $\natural$-norm);
$K^{rd}_{\mathfrak{K}}$ denotes the constant of the reading part; $\mathrm{rd}^{\mathsf{m}}P^{\mathsf{m}}$
is the datum read by run $\mathsf{m}$ at the site, that is, the hull-averaging read map
$\mathrm{rd}^{\mathsf{m}}$ applied to the argument $P^{\mathsf{m}}$ of run $\mathsf{m}$; $z$ is a
class of marks and $n(z)$ its layer; and $\mathsf{i}$ is the summation index of the loads
$\mathsf{l}^{\mathfrak{K},(z,\mathsf{i})}$ of the class $z$ in the resolved-extension statement:
\begin{equation}\label{9.5:eq-kept-difference-layers-a}
\begin{aligned}
&\text{(1)}\quad \sum_n \delta v^{(n)}=\hat{v}^{\mathsf{B}}-\hat{v}^{\mathsf{A}},\\
&\text{(3)}\quad \mathfrak{O}^{+}\bigl(\delta V^{(n)}(Y^{\sharp})\bigr)
+\mathfrak{O}^{+}\bigl(\delta\mathsf{Q}^{(n)}_Y\bigr)\le \mathsf{k}_n\,\check{R}_m^{p_{\Sigma}},
\qquad \mathsf{k}_n\le K_{\mathfrak{K}}\,[\![\delta]\!]^{\natural+}_n,\\
&\text{(4)}\quad \bigl\|\hat{v}^{\mathsf{B}}(\mathrm{rd}^{\mathsf{B}}P^{\mathsf{B}})
-\hat{v}^{\mathsf{B}}(\mathrm{rd}^{\mathsf{A}}P^{\mathsf{A}}\oplus c^{\mathsf{B}})\bigr\|_{\natural}
\le K^{rd}_{\mathfrak{K}}\,\vartheta_f^{m/2}\,\check{R}_m^{p_{\Sigma}}\\
&\qquad\qquad\le K_{\mathfrak{K}}\,[\![\delta]\!]^{c}_m\,\check{R}_m^{p_{\Sigma}},\\
&\text{(5)}\quad \sum_{\mathsf{i}}\mathsf{l}^{\mathfrak{K},(z,\mathsf{i})}\le \mathsf{k}_{n(z)}
+\mathbf{1}[n(z)=m]\cdot(\text{the reading constant}),\\
&\text{(7)}\quad \sum_{n\le k'}(k'-n+1)\,q_a^{(k'-n-1)_+}\,\mathsf{k}_n
\le K_{\mathfrak{K}}\,C^{TC}_{\Sigma}\,C_{env}\,\mathsf{l}^2_{s(k')}\,e_{s(k')}.
\end{aligned}
\end{equation}
\begin{enumerate}
\item[(1)] Exact resolution: the identity (1) of \eqref{9.5:eq-kept-difference-layers-a} holds
at every common datum and at every common real value $\Phi_Y$ of the inner variables.
\item[(2)] Support implies age: $\delta v^{(n)}\not\equiv 0$ only if the history of the label has
zone $n$ inside $Y^{\sharp}$. Consequently, at every later level $k'>m$ the class $(n,k')$ occurs
only in labels having a live component $C\supseteq Y$ created at a step $j(C)\le n+1$; this is the
support-implies-age property of the class loads.
\item[(3)] Production, by oscillation and not by size: with $\mathfrak{O}^{+}$ the two-point
oscillation over the common data at one value $\Phi_Y$
(see \eqref{9.5:eq-two-point-oscillation-b}), the bounds (3) of
\eqref{9.5:eq-kept-difference-layers-a} hold, where $K_{\mathfrak{K}}$ is a number, fixed before
$\gamma$ in the order of choice, that depends neither on the quantities written $D$ and $N_0$ in
the one-run bounds of hypothesis~(ii), on $|\mathsf{C}|$, on the age, on $K$, $\mathfrak{n}$ or
$k$, nor on the number, the size or the past of the thrown-back regions.
\item[(4)] The reading part: at every later site the remaining member on the left of (4) in
\eqref{9.5:eq-kept-difference-layers-a} is classical, is assigned to the class $(m,k')$, and
satisfies (4) of \eqref{9.5:eq-kept-difference-layers-a}.
\item[(5)] The sites: the bounds for the kept factor at the four sites at which it is read, with
unit shares $\theta_T$, $\theta_P$, $\theta^{\circ}_L$, $\theta^{\dagger}_L$, hold for the marks of
class $\zeta$ with the one-run size $\mathfrak{O}(v)$ replaced by
$\mathsf{k}_n\check{R}_m^{p_{\Sigma}}$. In particular clause~(ii) of the resolved-extension
statement holds class by class on $\hat{\mathfrak{E}}\cup\{\hat{e}\}$, where $\hat{\mathfrak{E}}$ is
the family of $\natural$-extensions on which that statement is formulated and $\hat{e}$ is the
extension supplied by the bounds at the four sites, with the bound (5) of
\eqref{9.5:eq-kept-difference-layers-a}.
\item[(6)] The loop closes and the five unknowns are kept: under one upper bound on $\gamma$ of
supremum form, imposed after $C_{\Sigma}$ in the order of choice (the door ceiling of
\eqref{9.5:eq-kept-elimination-a}), the kept factors introduce no new unknown; the load bound at
production, the class-load lemma, the volume and filament bound in its class form, the bound on
the $\mathbf{R}$-born boundary terms $\mathbf{B}'$, the elimination of the boundary terms
$\mathbf{B}$, the a priori lemma and the consistency corollary hold with each constant replaced
by the maximum of itself and a fixed number.
\item[(7)] Current depth: consequently the bound (7) of \eqref{9.5:eq-kept-difference-layers-a}
holds at every $k'$. In words: the contribution of the kept factors of all live thrown-back
regions, of all ages, to the inequality of step $k'$ of the comparison is small like the envelope
factor $e_{s(k')}$ at the depth $s(k')$ of $k'$.
\end{enumerate}
\end{theorem}

The theorem does not assert that the quantity $\mathsf{l}^{\mathfrak{K}}(k)$, the supremum over the
live thrown-back regions in the form in which it enters the comparison's defining display, is
small: it is of the size of the oldest class present, and as a statement at the current depth it is
false. It asserts nothing about the arrow statement of $V$, and nothing
in a normalised measure.

The proof occupies the remainder of this section.

\begin{lemma}[Schwarz contraction on a disc of radius above one]\label{9.5:lem-disc-contraction}
Let $\rho>1$, and let $f$ be a holomorphic function on the open disc
$\{\lambda\in\mathbb{C}:|\lambda|<\rho\}$. Here $\rho$ and $f$ are local to this lemma. Let
$\mathfrak{B}\ge 0$ be a constant with
\begin{equation*}
\sup_{|\lambda|<\rho}\,|f(\lambda)-f(0)|\le\mathfrak{B}.
\end{equation*}
Then
\begin{equation*}
|f(1)-f(0)|\le\frac{\mathfrak{B}}{\rho}.
\end{equation*}
\end{lemma}

\begin{proof}
If $\mathfrak{B}=0$, the hypothesis gives $f(\lambda)=f(0)$ for all $|\lambda|<\rho$. Since
$\rho>1$, this holds in particular at $\lambda=1$, and both sides of the conclusion vanish.

Let $\mathfrak{B}>0$. On the open unit disc put
\begin{equation*}
h(\lambda):=\frac{f(\rho\lambda)-f(0)}{\mathfrak{B}},\qquad |\lambda|<1.
\end{equation*}
For $|\lambda|<1$ we have $|\rho\lambda|<\rho$, so $h$ is holomorphic on the open unit disc.
By the hypothesis, $|h(\lambda)|\le1$ there, and $h(0)=0$. Schwarz's lemma therefore gives
$|h(\lambda)|\le|\lambda|$ for all $|\lambda|<1$.

Since $\rho>1$, the point $\lambda=1/\rho$ lies in the open unit disc. Evaluating there gives
\begin{equation*}
\frac{|f(1)-f(0)|}{\mathfrak{B}}=\Bigl|h\Bigl(\frac{1}{\rho}\Bigr)\Bigr|\le\frac{1}{\rho},
\end{equation*}
which is the assertion.
\end{proof}

\begin{lemma}[Vertex values against a polydisc bound]\label{9.5:lem-polydisc-vertices}
Let $\mathsf{y}$ be a finite set and $R>1$. Let $G$ be a function of the variables
$\vec{\mathsf{s}}=(\mathsf{s}(\Delta))_{\Delta\in\mathsf{y}}\in\mathbb{C}^{\mathsf{y}}$, holomorphic on a
neighbourhood of the closed polydisc $\{\vec{\mathsf{s}}:\ |\mathsf{s}(\Delta)|\le R\ \text{for all }
\Delta\in\mathsf{y}\}$, and put
\begin{equation}\label{9.5:eq-polydisc-vertices-1}
\mathfrak{p}[G]:=\prod_{\Delta\in\mathsf{y}}\int_0^1 d\mathsf{s}(\Delta)\,\partial_{\mathsf{s}(\Delta)}\,G .
\end{equation}
For $\sigma'\subset\mathsf{y}$ let $\mathbf{1}_{\sigma'}\in[0,1]^{\mathsf{y}}$ be the vertex with
$\mathsf{s}(\Delta)=1$ for $\Delta\in\sigma'$ and $\mathsf{s}(\Delta)=0$ for
$\Delta\in\mathsf{y}\setminus\sigma'$, and let $\sup|G|$ denote the supremum of $|G|$ over the closed
polydisc. Then:
\begin{itemize}
\item[(i)] $\mathfrak{p}[G]=\sum_{\sigma'\subset\mathsf{y}}(-1)^{|\mathsf{y}\setminus\sigma'|}\,
G(\mathbf{1}_{\sigma'})$;
\item[(ii)] $|\mathfrak{p}[G]|\le (R-1)^{-|\mathsf{y}|}\sup|G|$;
\item[(iii)] if $|G(\mathbf{1}_{\sigma'})|\le\beta$ for all $\sigma'\subset\mathsf{y}$ and
$\sup|G|\le\mathfrak{B}$, then
\begin{equation*}
|\mathfrak{p}[G]|\le\Bigl(\frac{2}{R-1}\Bigr)^{|\mathsf{y}|/2}(\beta\mathfrak{B})^{1/2}.
\end{equation*}
\end{itemize}
\end{lemma}

\begin{proof}
Since $R>1$, the real cube $[0,1]^{\mathsf{y}}$ lies in the polydisc, and $G$ is smooth on a
neighbourhood of it; hence the partial derivatives in \eqref{9.5:eq-polydisc-vertices-1} commute, the
iterated integrals may be taken in any order, and $\mathfrak{p}[G]$ is the integral over
$[0,1]^{\mathsf{y}}$ of the mixed derivative
$\partial^{\mathsf{y}}G:=\prod_{\Delta\in\mathsf{y}}\partial_{\mathsf{s}(\Delta)}G$.

(i) For one variable $\mathsf{s}(\Delta)$ and a smooth function $H$, the fundamental theorem of calculus
gives
\begin{equation*}
\int_0^1 d\mathsf{s}(\Delta)\,\partial_{\mathsf{s}(\Delta)}H
=H\big|_{\mathsf{s}(\Delta)=1}-H\big|_{\mathsf{s}(\Delta)=0}.
\end{equation*}
Applying this in each variable in turn, \eqref{9.5:eq-polydisc-vertices-1} becomes the product over
$\Delta\in\mathsf{y}$ of these commuting difference operations applied to $G$. Expanding the product,
each term corresponds to the choice of the endpoint $1$ for $\Delta$ in a subset $\sigma'$ and of the
endpoint $0$ for $\Delta\in\mathsf{y}\setminus\sigma'$; it equals $G(\mathbf{1}_{\sigma'})$ with the
sign $(-1)^{|\mathsf{y}\setminus\sigma'|}$, one minus sign for each lower endpoint. This is (i).

(ii) Fix $\vec{\mathsf{s}}\in[0,1]^{\mathsf{y}}$. If $|\sigma(\Delta)-\mathsf{s}(\Delta)|\le R-1$ for all
$\Delta$, then $|\sigma(\Delta)|\le|\mathsf{s}(\Delta)|+R-1\le R$, so the closed polydisc of polyradius
$R-1$ about $\vec{\mathsf{s}}$ lies in the closed polydisc of radius $R$. Cauchy's formula on the
distinguished boundary of the former gives
\begin{equation*}
\partial^{\mathsf{y}}G(\vec{\mathsf{s}})
=\prod_{\Delta\in\mathsf{y}}\Bigl(\frac{1}{2\pi i}\oint_{|\sigma(\Delta)-\mathsf{s}(\Delta)|=R-1}
\frac{d\sigma(\Delta)}{(\sigma(\Delta)-\mathsf{s}(\Delta))^{2}}\Bigr)\,G(\vec{\sigma}),
\end{equation*}
and each circle contributes at most $(2\pi)^{-1}\cdot 2\pi(R-1)\cdot(R-1)^{-2}=(R-1)^{-1}$, so
$|\partial^{\mathsf{y}}G(\vec{\mathsf{s}})|\le(R-1)^{-|\mathsf{y}|}\sup|G|$. Integrating over the cube
$[0,1]^{\mathsf{y}}$, of volume one, gives (ii).

(iii) By (i) and the hypothesis on the vertex values, $|\mathfrak{p}[G]|\le 2^{|\mathsf{y}|}\beta$, since
there are $2^{|\mathsf{y}|}$ subsets $\sigma'$. By (ii), $|\mathfrak{p}[G]|\le(R-1)^{-|\mathsf{y}|}\mathfrak{B}$.
With $\min\{a,b\}\le(ab)^{1/2}$ for $a,b\ge0$,
\begin{equation*}
|\mathfrak{p}[G]|\le\bigl(2^{|\mathsf{y}|}\beta\,(R-1)^{-|\mathsf{y}|}\mathfrak{B}\bigr)^{1/2}
=\Bigl(\frac{2}{R-1}\Bigr)^{|\mathsf{y}|/2}(\beta\mathfrak{B})^{1/2}.
\end{equation*}
For $\mathsf{y}=\emptyset$ the product in \eqref{9.5:eq-polydisc-vertices-1} is empty,
$\mathfrak{p}[G]=G=G(\mathbf{1}_{\emptyset})$, and the bound reads
$|G|\le\min\{\beta,\mathfrak{B}\}\le(\beta\mathfrak{B})^{1/2}$.
\end{proof}

\begin{lemma}[A supremum recursion without compounding]\label{9.5:lem-sup-recursion}
Let $0\le\theta<1$, let $(a_m)_{m\ge1}$ be a non-decreasing sequence of real numbers, and let
$x_m\ge0$ $(m=1,2,\dots)$ satisfy
\begin{equation}\label{9.5:eq-sup-recursion-1}
x_m\le a_m+\theta\,\sup_{m'<m}x_{m'}\qquad(m=1,2,\dots),
\end{equation}
with the convention $\sup\emptyset:=0$. Then
\begin{equation}\label{9.5:eq-sup-recursion-2}
x_m\le\frac{a_m}{1-\theta}\qquad(m=1,2,\dots).
\end{equation}
\end{lemma}

\begin{proof}
For $m=1$ the supremum in \eqref{9.5:eq-sup-recursion-1} is over the empty set. Hence
$0\le x_1\le a_1$, so $a_1\ge0$. Since $(a_m)$ is non-decreasing, $a_m\ge0$ for every $m$.

We prove \eqref{9.5:eq-sup-recursion-2} by induction on $m$, assuming it for every $m'<m$. The
supremum over $m'<m$ is then the maximum of finitely many numbers, each at most
$a_{m'}/(1-\theta)$. For $m=1$ it is $0$, which is at most $a_1/(1-\theta)$ because $a_1\ge0$
and $1-\theta>0$. For $m\ge2$ the sequence $(a_m)$ is non-decreasing and $1-\theta>0$, so
\begin{equation*}
\sup_{m'<m}x_{m'}\le\frac{a_{m-1}}{1-\theta}\le\frac{a_m}{1-\theta}.
\end{equation*}
Thus in every case $\sup_{m'<m}x_{m'}\le a_m/(1-\theta)$. We multiply this inequality by
$\theta\ge0$ and insert it into \eqref{9.5:eq-sup-recursion-1}. This gives
\begin{equation*}
x_m\le a_m+\frac{\theta\,a_m}{1-\theta}=\frac{a_m}{1-\theta},
\end{equation*}
which is \eqref{9.5:eq-sup-recursion-2} at $m$.

No finiteness of $\sup_{m'<m}x_{m'}$ is assumed in advance: it follows from the induction. The
constant $1/(1-\theta)$ does not grow with $m$.
\end{proof}

\begin{definition}[Resolved terms: family marks and reading marks]\label{9.5:def-resolved-terms}
\emph{Setting.} Fix a step $m$, an aligned label, and a connected component $C$ of the
large-field region $Z$ at the large-field stage of the operation $\mathbf{R}_m$; if $C$ becomes a
region of the second class we write $C=Y$. At level $k:=m$ we use the notation of the total bound
on the terms of the step:
\begin{itemize}
\item $h:=m-N$, where $N$ is the integer fixed in the total bound on the terms of the step at
level $k=m$ (here $N$ is not the rank of the gauge group);
\item the core $\mathsf{C}$ and the strip $C^{\sharp}$;
\item the scaled contour distance $d^{\wedge}$, the weight $w$ and the room $\varrho(\Delta)$;
\item the space $\mathfrak{P}^{\sharp}$ of production points $\mathfrak{z}=(\zeta,\varpi,\vec{\mathsf{s}})$,
whose entries are the slice $\zeta$, the auxiliary parameters $\varpi$ and the decoupling
parameters $\vec{\mathsf{s}}$.
\end{itemize}
Let $\mathsf{A}$ and $\mathsf{B}$ be the two runs of the comparison carried out in this chapter, of
lengths $\mathfrak{n}$ and $\mathfrak{n}+1$, and let $\mathsf{m}\in\{\mathsf{A},\mathsf{B}\}$. In this
comparison the slice ranges over
\begin{equation}\label{9.5:eq-resolved-terms-1}
\zeta^{\mathsf{m}}(\xi):=(U^{\mathsf{m}}_m,J^{\mathsf{m}}_m)\bigl((C^{\sharp})^{+},\xi\bigr)
\upharpoonright C^{\sharp},\qquad \xi\in\mathfrak{X}_{C^{\sharp}},
\end{equation}
where $\mathfrak{X}_{C^{\sharp}}$ is the set of common data with charts of
Notation~\ref{9.5:not-two-point-oscillation}; here
$(U^{\mathsf{m}}_m,J^{\mathsf{m}}_m)\bigl((C^{\sharp})^{+},\xi\bigr)$ consists of the configuration
$U^{\mathsf{m}}_m$ and the accompanying datum $J^{\mathsf{m}}_m$ of run $\mathsf{m}$ at step $m$ on the
enlargement $(C^{\sharp})^{+}$ of $C^{\sharp}$, determined by the common datum $\xi$, and
$\upharpoonright C^{\sharp}$ is restriction to $C^{\sharp}$.

Production points of the two runs are paired by the common entries $(\xi,\varpi,\vec{\mathsf{s}})$.
We put
\[
\mathfrak{z}^{\mathsf{m}}:=(\zeta^{\mathsf{m}}(\xi),\varpi,\vec{\mathsf{s}}).
\]
For a term $v$ of the list of terms of the step fixed in the total bound on the terms of the step,
$f^{\mathsf{m}}_v(\mathfrak{z}^{\mathsf{m}})$ is the function attached to $v$ in that bound. In the
comparison a \emph{stored}
term, that is, a term that carries a store in the sense of the definition of the store of a kept
member, is read through its own hull:
\begin{equation}\label{9.5:eq-resolved-terms-2}
f^{\mathsf{m}}_v(\mathfrak{z}^{\mathsf{m}})=\hat{F}^{\mathsf{m}}_v(\eta^{\mathsf{m}}_v),\qquad
\eta^{\mathsf{m}}_v(\xi,\varpi,\vec{\mathsf{s}})
:=\mathrm{rd}^{\mathsf{m}}_{X_v^{+}}\,\mathrm{chain}^{\mathsf{m}}(\mathfrak{z}^{\mathsf{m}}),
\end{equation}
where $\mathrm{chain}^{\mathsf{m}}(\mathfrak{z}^{\mathsf{m}})$ denotes the chain of configurations
that run $\mathsf{m}$ produces from the production point $\mathfrak{z}^{\mathsf{m}}$.
We call $\eta^{\mathsf{m}}_v$ the \emph{read datum} of $v$ in run $\mathsf{m}$. For a term that is a
renormalised group or a core, $X_v^{+}$ is the box $\square_0$ of the corollary on renormalised
groups.

\emph{Resolutions.} Let $v$ be a stored term of level $j_v$ and domain $X_v$, and write the
elements of $\mathfrak{X}_{X_v}$ as pairs $(\mathsf{v},c)$, where $\mathsf{v}$ is the datum read by
$\mathsf{A}$ and $c$ the collar entries of $\mathsf{B}$. A \emph{resolution} of $v$ is a finite family
$(G^{(n)}_v)_n$ of holomorphic functions on $\mathfrak{X}_{X_v}$ (times the common slots of $v$)
with the following three properties.
\begin{enumerate}
\item[(ru1)] For every $(\mathsf{v},c)\in\mathfrak{X}_{X_v}$,
\begin{equation}\label{9.5:eq-resolved-terms-3}
\sum_n G^{(n)}_v(\mathsf{v},c)=\hat{F}^{\mathsf{B}}_v(\mathsf{v}\oplus c)-\hat{F}^{\mathsf{A}}_v(\mathsf{v}).
\end{equation}
\item[(ru2)] $G^{(n)}_v\equiv0$ unless zone $n$ meets $X_v$; inside live or arrested components
the condition is taken in the sequence form of the zone condition on class marks.
\item[(ru3)] $G^{(n)}_v$ has the format of $v$, namely its invariance, its pointing and its slots.
Consequently the size functional $\mathfrak{N}_v[\cdot]$ of $v$, taken from the total bound on the
terms of the step, is defined on $G^{(n)}_v$.
\end{enumerate}

\emph{Family marks and the reading mark.} Given a resolution, put
\[
\eta^{\times}_v:=\eta^{\mathsf{A}}_v\oplus\bigl(\text{the collar entries of }\eta^{\mathsf{B}}_v\bigr),
\]
and define
\begin{equation}\label{9.5:eq-resolved-terms-4}
\mathsf{g}^{(n)}_v(\xi,\varpi,\vec{\mathsf{s}}):=G^{(n)}_v(\eta^{\times}_v),\qquad
\mathsf{r}_v(\xi,\varpi,\vec{\mathsf{s}}):=\hat{F}^{\mathsf{B}}_v(\eta^{\mathsf{B}}_v)
-\hat{F}^{\mathsf{B}}_v(\eta^{\times}_v).
\end{equation}
The $\mathsf{g}^{(n)}_v$ are the \emph{family marks} of $v$, and $\mathsf{r}_v$ is its
\emph{reading mark}. They satisfy, identically in $(\xi,\varpi,\vec{\mathsf{s}})$,
\begin{equation}\label{9.5:eq-resolved-terms-a}
f^{\mathsf{B}}_v(\mathfrak{z}^{\mathsf{B}})-f^{\mathsf{A}}_v(\mathfrak{z}^{\mathsf{A}})
=\sum_n\mathsf{g}^{(n)}_v+\mathsf{r}_v .
\end{equation}

\emph{The class of the reading mark.} Let $n_{X_v}$ be the least level of the label whose zone
meets $X_v^{+}$. The reading mark is assigned to the class to which the comparison assigns the
classical part of a use. This is the class $n^c_v$ in which the classical part
$[\![\delta]\!]^{c}$ of the comparison's defect carries its entry, and it is given as follows.
\begin{enumerate}
\item[(i)] For boundary-type terms, namely the boundary terms $\mathbf{B}$, $\mathbf{B}'$, the chain
terms $\mathbb{C}^{(\nu)}$, the terms of interrupted runs, the terms over inner live strips and the
inner kept members,
\[
n^c_v:=j_v\wedge n_{X_v}.
\]
This equals $n_{X_v}$, which is a zone met by $X_v^{+}$.
\item[(ii)] For pointed terms of the list whose hull lies in zones $\ge j_v$, $n^c_v$ is the
anchor zone $n$. The class of a use is fixed by the label and by the position of the anchor.
The corresponding entries of $[\![\delta]\!]^{c}_n$ are the two classical entries of the
layer-$n$ load at level $j_v$,
\[
\varsigma_{1+n-j_v}E_0\hat{\vartheta}^{\natural}_{j_v}
\qquad\text{and}\qquad
w_{n,j_v}g_{j_v}^{\kappa_0}\hat{\vartheta}^{\natural}_{j_v}.
\]
\item[(iii)] For pointed terms whose hull meets a zone $n_{X_v}<j_v$, $n^c_v:=n_{X_v}$. This
contribution is carried by the last entry of $[\![\delta]\!]^{c}$.
\end{enumerate}
In every case the reading mark carries the smallness index $\vartheta^{\sharp}_v$ taken at the
level $j_v\wedge n_{X_v}$, in the sense in which the size bounds for resolved terms fix that
index at a level.

\emph{The resolutions of the stored terms.} Each stored term has the resolution listed below.
\begin{enumerate}
\item[(a)] A boundary term $\mathbf{B}^{(\ell)}(X)$ created by the operation $\mathbf{T}$:
$G^{(n)}_v=\delta\mathbf{B}^{(\ell),(n)}(X;\cdot)$ for $n\le\ell-1$. This is given by the store
identities for boundary terms.
\item[(b)] A boundary term $\mathbf{B}'^{(\ell)}(X')$ created by the operation $\mathbf{R}$:
$G^{(n)}_v=\delta\mathbf{B}'^{(\ell),(n)}$ for $n\le\ell-1$, together with the own member
$\delta\mathbf{B}'^{(\ell),\mathrm{own}}$. The own member is assigned to the class $n:=\ell$. This
is given by the store identities for boundary terms created by $\mathbf{R}$.
\item[(c)] A left-over or frozen chain term $\mathbb{C}^{(i)}_{k_a}(X)$: the family $G^{(n)}_v$ is
given by the increments $\delta\Phi^{z}(X;\cdot)$.
This is given by the store identities for left-over and frozen chain terms.
\item[(d)] Terms over inner live strips: the same stores as in (a)--(c), taken in the form
$\mathsf{S}^{\theta}$ of the expansion of terms over inner live strips. This is given by that
expansion.
\item[(e)] Singles, tails and weights of the list: one class only, the anchor zone $n$, with
$G^{(n)}_v=\hat{F}^{\mathsf{B}}_v-\hat{F}^{\mathsf{A}}_v$ for that $n$ and $G^{(n)}_v=0$ for every
other $n$. This is given by parts (i) and (iii) of the pair lemma for list terms.
\item[(f)] Renormalised groups and cores of the list: one class $n$ only, with
\[
G^{(n)}_v=\mathfrak{i}_y[\mathcal{G}^{+}]-b^{\mathsf{B}}\,\mathfrak{i}_y[\mathcal{G}^{+}_{\hat{A}}]
\]
for that $n$ and $G^{(n)}_v=0$ for every other $n$. Here $\mathcal{G}^{+}$ is an
$\mathsf{A}$-family, $\mathfrak{i}_y$ is the renormalised group functional, and
$\mathcal{G}^{+}_{\hat{A}}$ and $b^{\mathsf{B}}$ are the family and the coefficient that the
corollary on renormalised groups attaches to $\mathcal{G}^{+}$. This is given by the corollary on
renormalised groups, part (ii) of the pair lemma for list terms, and the lemma on the marks of
list terms.
\item[(g)] The terms $\mathfrak{F}^{B4}(\mathfrak{a})\setminus\{\mathsf{q}_W\}$ of an interrupted run:
the family $G^{(n)}_v$ is given by the increments $\delta\Phi^{z}$ of part (d) of the proposition
on terms of interrupted runs.
\item[(h)] Inner kept members $v'\in\mathfrak{K}(Y')$ with $Y'\subset C$:
$G^{(n)}_{v'}=\delta\hat{v}'^{(n)}$, the class-$n$ member of the store of $v'$ constructed at
$(Y',m')$, $m'<m$, in the definition of the store of a kept member. Since $m'<m$, this family is
available by induction on $m$.
\end{enumerate}

\emph{Terms derived inside the step.} The chain terms $\mathbb{C}^{(\nu)}_m(X)$ of the present
step are derived inside the step and have no store yet. Their marks are not taken from a store
but defined directly, as follows.
For a piece $u=(\mathbb{C}^{(\nu)}_m(X),\vec{\mathsf{y}})$ and an index $(z,\mathsf{i})$,
consisting of a dial class $z$, with zone $n(z)$, and of $\mathsf{i}\in\{\mathrm{u},\mathrm{c}\}$,
let $W^{(z,\mathsf{i})}_u$ be the alternating sum, in the sense of the alternating-sum lemma for
$\natural$-extensions, of the resolved $\natural$-extension
$\bigl(\delta\mathbb{C}^{(\nu)}(X)\bigr)^{\natural,(z,\mathsf{i})}_{\mathfrak{s}_u}$, where
$\mathfrak{s}_u$ is the index of the piece $u$ in that extension. These marks satisfy
\[
\sum_{(z,\mathsf{i})} W^{(z,\mathsf{i})}_u=x^{\mathsf{B}}_u-x^{\mathsf{A}}_u ,
\]
where $x^{\mathsf{m}}_u$ is the value of the piece $u$ in run $\mathsf{m}$; they vanish unless the
zone $n(z)$ meets $X$; and they are holomorphic in the seed, which is a component of $\varpi$.

\emph{Constituents without a store.} The following constituents are purely classical: they
contain no unknown, and each lies in one class only.
\begin{itemize}
\item The Gaussian terms $v_e\in\mathfrak{Q}(\mathfrak{a})$ of live arrests of interrupted runs
nested in $C$ lie in the class $j_W$, the level of the arrest.
\item The form $\mathfrak{q}_C$ lies in the class $h+1$.
\end{itemize}
The members $\mathfrak{v}^{\mathbf{C}}(Y_4)$ that do not depend on the dial variables are split
according to the lowest zone they meet. Each of them has a part that is linear in the difference
of the couplings and a classical part.
\end{definition}

\begin{proof}
\emph{The identity \eqref{9.5:eq-resolved-terms-a}.} By \eqref{9.5:eq-resolved-terms-2}, applied in
both runs, the left-hand side of \eqref{9.5:eq-resolved-terms-a} is
\[
\hat{F}^{\mathsf{B}}_v(\eta^{\mathsf{B}}_v)-\hat{F}^{\mathsf{A}}_v(\eta^{\mathsf{A}}_v).
\]
By definition $\eta^{\times}_v=\eta^{\mathsf{A}}_v\oplus c$, where $c$ denotes the collar entries of
$\eta^{\mathsf{B}}_v$. We apply \eqref{9.5:eq-resolved-terms-3} at the datum
$(\mathsf{v},c)=(\eta^{\mathsf{A}}_v,c)$ and use the definition \eqref{9.5:eq-resolved-terms-4} of the
family marks. This gives
\[
\sum_n\mathsf{g}^{(n)}_v=\sum_nG^{(n)}_v(\eta^{\times}_v)
=\hat{F}^{\mathsf{B}}_v(\eta^{\times}_v)-\hat{F}^{\mathsf{A}}_v(\eta^{\mathsf{A}}_v).
\]
Adding $\mathsf{r}_v=\hat{F}^{\mathsf{B}}_v(\eta^{\mathsf{B}}_v)-\hat{F}^{\mathsf{B}}_v(\eta^{\times}_v)$, the
terms $\pm\hat{F}^{\mathsf{B}}_v(\eta^{\times}_v)$ cancel. The sum is therefore
$\hat{F}^{\mathsf{B}}_v(\eta^{\mathsf{B}}_v)-\hat{F}^{\mathsf{A}}_v(\eta^{\mathsf{A}}_v)$, which is
\eqref{9.5:eq-resolved-terms-a}.

The identity \eqref{9.5:eq-resolved-terms-a} is the splitting
$\delta_u=\delta^{\mathrm{fam}}+\delta^{\mathrm{rd}}$ of the difference of a use into a family part
and a reading part in the pair lemma for list terms. It is also the splitting $g=g^{1}+g^{2}$ of
the corollary on $\natural$-extensions of terms, in which $g^{1}$ corresponds to the sum of the
family marks and $g^{2}$ to the reading mark.

\emph{The class in case} (i). The rule $n^c_v:=j_v\wedge n_{X_v}$ is the class assignment of the
classical part of a use for boundary-type terms. By the anchoring property of boundary-type
terms, every term $F\in\mathfrak{L}^{B}$ of level $j$ has a cube of the level-$j$ partition
$\pi_j$ meeting the enlarged region $Z_j^{\boxplus}$. Applied to $v$, with $j=j_v$, this gives
$n_{X_v}\le j_v$. Therefore $j_v\wedge n_{X_v}=n_{X_v}$, and by the definition of $n_{X_v}$
this is a zone met by $X_v^{+}$.

\emph{Property} (ru3)\emph{, case by case.} We first prove a general fact. The projection onto a
class is defined by the scaffold alone, that is, by the assignment of slots to classes that the
label and the anchors determine. A gauge map applied in common to the datum, a common source
$(y,w)$ and a common slot each act slot by slot. Hence each of them commutes with the projection
onto a class. We now verify (ru3) in the cases (a)--(h).

Cases (a) and (b), including the own members. Each member is a first variation
\[
\int_0^1d\lambda\,\frac{d}{d\varepsilon}\Big|_{\varepsilon=0}
\mathcal{P}_X\bigl[x^{\lambda}+\varepsilon W^{(z)}\bigr],
\]
where $x^{\lambda}$ is the interpolated configuration and $\mathcal{P}_X$ is a single formula in
invariant slots. Each mark is invariant, by
induction on the level. By the first store identity, each mark is holomorphic on the tube of the
term. The format of these members asks no more than this invariance and holomorphy.

Cases (c) and (g). The increments $\delta\Phi^{z}$, and the increment of part (d) of the
proposition on terms of interrupted runs, are obtained in the same way through the first
variation of the resolution lemma for chain terms. By part (a) of the proposition on terms of
interrupted runs, they are jointly invariant with the silent slot $x$.

Case (d). By the expansion of terms over inner live strips, the stores of these terms, and hence
their formats, are those of cases (a)--(c), verified above.

Case (e). Here $G^{(n)}_v=\hat{F}^{\mathsf{B}}_v-\hat{F}^{\mathsf{A}}_v$ is a difference of two
functions pointed at the same source $(y,w)$. The pointed invariant functions on a domain form a
linear space, so $G^{(n)}_v$ is again such a function.

Case (f). Both $\mathcal{G}^{+}$ and $\mathcal{G}^{+}_{\hat{A}}$ are pointed families of the kind
required by the pair lemma for list terms and by the lemma on the marks of list terms. Hence the
difference $\mathfrak{i}_y[\mathcal{G}^{+}]-b^{\mathsf{B}}\mathfrak{i}_y[\mathcal{G}^{+}_{\hat{A}}]$
has the format of the term.

Case (h). The class members of an inner kept member are invariant under common centre maps with
values in the gauge group, acting jointly on the datum and on the stored term $\Phi$. By the
general fact above, such maps commute with the projection onto a class.

\emph{The class of the own member in case} (b). The theorem on resolved $\natural$-extensions of
stored terms would assign the own member to the class $i_u\vee\ell$, where $i_u$ is the class
attached to the read $u$. Assigning it to the class
$\ell$ instead is legitimate inside the component $C$, for the following reason. Since
$\mathbf{B}'^{(\ell)}(X')$ belongs to the first group of terms in \cite[(1.98)]{B-CMP122b}, the
region $X'$ meets $Z^{\sim}_{\ell}$; moreover, $X'$ lies in $C$. Hence the component containing
$Y$ has first index at most
$\ell+1$, and this is all that the zone condition on class marks requires of the class $\ell$.

\emph{Terms derived inside the step.} By part (i) of the alternating-sum lemma for
$\natural$-extensions,
\[
\sum_{(z,\mathsf{i})} W^{(z,\mathsf{i})}_u=x^{\mathsf{B}}_u-x^{\mathsf{A}}_u ,
\]
where $x^{\mathsf{m}}_u$ is the value of the piece $u$ in run $\mathsf{m}$.
By part (iii) of the same lemma, $W^{(z,\mathsf{i})}_u=0$ unless the zone $n(z)$ meets $X$.
Finally, $W^{(z,\mathsf{i})}_u$ is holomorphic in the seed, which is a component of $\varpi$,
because a $\natural$-extension is holomorphic by definition, and $W^{(z,\mathsf{i})}_u$ is formed
from the resolved $\natural$-extension by an alternating sum.
\end{proof}

\begin{definition}[The layer-resolved difference of a kept factor]\label{9.5:def-layer-store}
Let $C$ be a region with strip $C^{\sharp}$. Let $(\xi,c)\in\mathfrak{X}^{\sharp}$ be a common datum,
$\xi=(\mathsf{v},c)$, and let $\Phi_C$ be a real field on $C$. For a layer $n\le m$, where $m$ is
the step index paired with $C$ (as the step $m'$ is paired with an inner thrown-back region $Y'$
below; here $m$ is not the torus exponent of $\mathbb{T}_m$), the
\emph{layer-$n$ difference} of the kept potential on $C^{\sharp}$ is
\begin{equation}\label{9.5:eq-layer-store-a}
\begin{split}
\delta V^{(n)}(C^{\sharp}) := {} & \sum_{v}\ \sum_{\vec{\mathsf{y}}:\ \mathrm{lab}_4\subset C}
\Bigl\{\mathfrak{p}_{v,\vec{\mathsf{y}}}\bigl[\mathsf{g}^{(n)}_v\bigr]
+\mathbf{1}[n^{c}_v=n]\,\mathfrak{p}_{v,\vec{\mathsf{y}}}\bigl[\mathsf{r}_v\bigr]\Bigr\}
+\sum_{u}\sum_{\mathsf{i}}W^{((n,m),\mathsf{i})}_u\\
&-\sum_{v\in\mathfrak{L}^{cp}(C)}\Bigl\{\mathsf{g}^{(n)}_v+\mathbf{1}[n^{c}_v=n]\,\mathsf{r}_v\Bigr\}
+\sum_{Y_4\in\mathcal{Y}_{rim}}W^{(n)}(Y_4)\\
&+\sum_{Y_4\subset C^{\sharp}}\delta\mathfrak{v}^{\mathbf{C},(n)}(Y_4)+\delta\mathfrak{y}^{(n)}.
\end{split}
\end{equation}
The terms in \eqref{9.5:eq-layer-store-a} are as follows.
\begin{itemize}
\item In the first sum, $v$ runs over the stored units of the list of $\mathfrak{v}(C)$ of the
following three kinds: the plain units, $v\in\mathfrak{L}^{pl}=\mathfrak{L}^{B}\sqcup\mathfrak{L}^{pt}$;
the wrapped units, $v\in\mathfrak{L}^{wr}$; and, for every inner thrown-back region $Y'\subset C$
with its step $m'$, the two factors of $\mathfrak{K}(Y')$ kept inside $Y'$,
see \eqref{9.5:eq-kept-factors-a}.
The inner sum runs over the toggle tuples $\vec{\mathsf{y}}$ whose labels $\mathrm{lab}_4$ lie in $C$.
The functions $\mathsf{g}^{(n)}_v$ and $\mathsf{r}_v$ are the family marks and the reading mark of
$v$ of Definition~\ref{9.5:def-resolved-terms}, and $n^{c}_v$ is the layer index attached to $v$,
so that $\mathsf{r}_v$ enters only the layer $n=n^{c}_v$.
\item $\mathfrak{p}_{v,\vec{\mathsf{y}}}$ is the piece operator of the scaffold: for a tuple
$\vec{\mathsf{y}}=(\mathsf{y}_t)_t$ of toggle sets and a function $G$ of the decoupling parameters
$\vec{\mathsf{s}}$,
\begin{equation}\label{9.5:eq-layer-store-1}
\mathfrak{p}_{v,\vec{\mathsf{y}}}[G]
:=\Bigl(\prod_{t}\prod_{\Delta\in\mathsf{y}_t}\int_0^1 d\mathsf{s}^{(t)}(\Delta)\,
\partial_{\mathsf{s}^{(t)}(\Delta)}\Bigr)G\,\Big|_{\mathsf{s}=0\ \text{off}\ \vec{\mathsf{y}}},
\end{equation}
that is, for each $t$ and each cube $\Delta\in\mathsf{y}_t$ the derivative in
$\mathsf{s}^{(t)}(\Delta)$ is integrated over $[0,1]$, and every decoupling parameter not indexed
by $\vec{\mathsf{y}}$ is set equal to zero.
\item In the second sum, $u$ runs over the derived units and $\mathsf{i}$ over the further indices
labelling their class marks, so that $((n,m),\mathsf{i})$ is the class of the mark
$W^{((n,m),\mathsf{i})}_u$ of the derived unit $u$.
\item The third sum, over the copies $v\in\mathfrak{L}^{cp}(C)$, enters with a minus sign and
carries the same marks $\mathsf{g}^{(n)}_v$ and $\mathbf{1}[n^{c}_v=n]\,\mathsf{r}_v$ as the first
sum, without the piece operator.
\item For a rim piece $Y_4\in\mathcal{Y}_{rim}$, $W^{(n)}(Y_4)$ is the class-$n$ mark of the total
$(\mathfrak{v}-\mathfrak{v}^{\mathbf{C}})(Y_4)$.
\item The fifth sum collects the class-$n$ terms $\delta\mathfrak{v}^{\mathbf{C},(n)}(Y_4)$ of the
localised members $\mathfrak{v}^{\mathbf{C}}$ of Ba{\l}aban's remainder terms $\mathbf{C}_m$, over the
pieces $Y_4\subset C^{\sharp}$.
\item $\delta\mathfrak{y}^{(n)}$ collects the class-$n$ marks of the terms that do not depend on the
slice variable $\zeta$ of the production point $\mathfrak{z}=(\zeta,\varpi,\vec{\mathsf{s}})$ in
the decomposition of the kept potential. These are fixed values: they do not depend on $(\xi,c)$,
and therefore the two-point oscillation over common data annihilates them,
$\mathfrak{O}^{+}(\delta\mathfrak{y}^{(n)})=0$. They enter only the value part of the comparison.
\end{itemize}
For the kept quadratic form of $C$ we put
\begin{equation}\label{9.5:eq-layer-store-2}
\delta\mathfrak{q}^{(n)}_C:=\mathbf{1}[n=h+1]\,
\bigl(\hat{\mathfrak{q}}^{\mathsf{B}}_C-\hat{\mathfrak{q}}^{\mathsf{A}}_C\bigr),
\end{equation}
where $\hat{\mathfrak{q}}^{\mathsf{A}}_C$ and $\hat{\mathfrak{q}}^{\mathsf{B}}_C$ are the kept
quadratic forms of the runs $\mathsf{A}$ and $\mathsf{B}$, written as functions of the datum, and
$h$ is the layer index attached to the kept quadratic form of $C$ (not a history). Thus
$\delta\mathfrak{q}^{(n)}_C$ vanishes except in the single layer $n=h+1$.

If $C=Y$ is a thrown-back region, the data kept for the factor of $Y$ are
\begin{equation}\label{9.5:eq-layer-store-3}
\bigl(\hat{v}^{\mathsf{A}},\hat{v}^{\mathsf{B}};\ (\delta\hat{v}^{(n)})_{n\le m}\bigr),
\end{equation}
namely the kept member of $Y$ in the run $\mathsf{A}$ and in the run $\mathsf{B}$, each written as a
function of its datum, and the layer-resolved differences $\delta\hat{v}^{(n)}$, $n\le m$, of that
member. Each entry of \eqref{9.5:eq-layer-store-3} is a function of $(\xi,c;\Phi_Y)$. No auxiliary
variable is kept among these data.
\end{definition}

\begin{lemma}[Exactness, support and values of the layered difference]\label{9.5:lem-layer-store-bounds}
We use the runs $\mathsf{A}$, $\mathsf{B}$, the common data $\xi=(\mathsf{v},c)$, the strip $Y^{\sharp}$
with the set of common data $\mathfrak{X}^{\sharp}_Y$, and the layer-resolved difference of
Definition~\ref{9.5:def-layer-store}; we let $C$ be the component whose strip is $Y^{\sharp}$, so
that Definition~\ref{9.5:def-layer-store} is read with $C^{\sharp}=Y^{\sharp}$. For a real datum
$\Phi_Y$ we write:
\begin{itemize}
\item $\hat{v}^{\mathsf{m}}$ for the kept potential $V^{\sharp}(Y^{\sharp})$ of run $\mathsf{m}$,
as a function of the common datum;
\item $\delta\hat{V}^{(n)}$ for the layer-resolved difference
of Definition~\ref{9.5:def-layer-store} in zone $n\le m$, as a function of $\xi=(\mathsf{v},c)$.
\end{itemize}
Assume the following:
\begin{enumerate}
\item[(a)] in each of the two runs, $V^{\sharp}(C^{\sharp})$ is written as one and the same scaffold
sum of pieces, copies, rim totals, localised remainder terms $\mathfrak{v}^{\mathbf{C}}$ and
$\zeta$-free terms;
\item[(b)] parts (i), (iii) and (iv) of the statement on class marks of derived units; part (iv)
bounds the values of the kept member on the region $Z^{\sim}$ by
$2K_{\mathbb{C}}[\![\delta]\!]^{\natural+}_{n(z)}$, with $n(z)$ the zone of the class $z$, per unit
of $\sum_j|\Gamma^0_j\cap Y|$, where $Z^{\sim}$ and the sets $\Gamma^0_j$ are those of that statement;
\item[(c)] part (c1) of the classification of class marks, $\sum_nW^{(n)}=\delta(\mathfrak{v}
-\mathfrak{v}^{\mathbf{C}})$, and its first three and sixth entries;
\item[(d)] part (a) of the lemma on stores of derived terms;
\item[(e)] part (ii) of the locality lemma for the localised remainder terms
$\mathfrak{v}^{\mathbf{C}}$;
\item[(f)] holomorphy of the read datum: $\eta^{\mathsf{m}}_v$ is holomorphic in
$(\xi,\vec{\mathsf{s}})$;
\item[(g)] part (d) of the tube corollary and part (d) of its form for kept factors: the read
datum lies in the tube data domains $\mathfrak{D}^{\mathsf{m}}_F[\varsigma]$;
\item[(h)] the mixing lemma;
\item[(i)] the boundary terms $\mathbf{B}$ created at the healing step carry only levels that lie
strictly above the level of the healing step and not above $m$;
\item[(j)] the kept potential $V^{\sharp}(Y^{\sharp})$ is a plain sum of its slots.
\end{enumerate}
Then, with
\begin{equation*}
K_{val}:=2K_{\mathbb{C}}+K(d)+C_{rim}C^{\sharp}_L ,
\end{equation*}
where $K_{\mathbb{C}}$ is the constant of part (iv) of hypothesis (b) and $K(d)$, $C_{rim}$,
$C^{\sharp}_L$ are constants of the bounds in hypotheses (c) and (d) (the letter $d$ in $K(d)$
belongs to the name of the constant and is not the dimension), the following hold:
\begin{equation}\label{9.5:eq-layer-store-bounds-a}
\begin{aligned}
&\textstyle\sum_n\delta\hat{V}^{(n)}=\hat{v}^{\mathsf{B}}-\hat{v}^{\mathsf{A}}
\ \text{ on }\mathfrak{X}^{\sharp}_Y\times\{\Phi_Y\},
\ \text{ each member holomorphic in }\xi;\\
&\delta\hat{V}^{(n)}\not\equiv0\ \Longrightarrow\ \text{the label's history has zone $n$
inside $Y^{\sharp}$ ($n<m$), or $n=m$};\\
&\sup|\delta\hat{V}^{(n)}|\le K_{val}[\![\delta]\!]^{\natural+}_n\times
\#\{\text{cells of $Y^{\sharp}$ met by class-$n$ slots}\}.
\end{aligned}
\end{equation}
The right-hand side of the last line is a volume; it is read at the top storeys only, under
un-normalised weights.
\end{lemma}

\begin{proof}
\emph{Exactness.} By hypothesis (a), the kept potential $V^{\sharp}(C^{\sharp})$ is, in each run, the
same scaffold sum of pieces, copies, rim totals, localised remainder terms
$\mathfrak{v}^{\mathbf{C}}$ and $\zeta$-free terms. We subtract the expression for run $\mathsf{A}$
from that for run $\mathsf{B}$, term by term. For a stored unit $v$ and a toggle tuple
$\vec{\mathsf{y}}$, the difference of the two pieces is
$\mathfrak{p}_{v,\vec{\mathsf{y}}}[\hat{F}^{\mathsf{B}}_v-\hat{F}^{\mathsf{A}}_v]$, because the
piece operator $\mathfrak{p}_{v,\vec{\mathsf{y}}}$ is linear. By the identities
\eqref{9.5:eq-resolved-terms-a} of Definition~\ref{9.5:def-resolved-terms}, among them the sum
rule $\sum_nG^{(n)}_v=\hat{F}^{\mathsf{B}}_v(\mathsf{v}\oplus c)-\hat{F}^{\mathsf{A}}_v(\mathsf{v})$,
together with part (i) of hypothesis (b), this difference splits into the family marks
$\mathsf{g}^{(n)}_v$ and the reading mark $\mathsf{r}_v$, the latter placed in the zone
$n=n^c_v$. The same holds for the copies in $\mathfrak{L}^{cp}(C)$, and the differences of the
derived units split into their class marks $W^{((n,m),\mathsf{i})}_u$. By part (c1) of
hypothesis (c),
$\sum_nW^{(n)}=\delta(\mathfrak{v}-\mathfrak{v}^{\mathbf{C}})$. Comparing the resulting expression
with Definition~\ref{9.5:def-layer-store} gives
$\sum_n\delta\hat{V}^{(n)}=\hat{v}^{\mathsf{B}}-\hat{v}^{\mathsf{A}}$.

\emph{Holomorphy.} By hypothesis (f), the read
datum $\eta^{\mathsf{m}}_v$ is holomorphic in $(\xi,\vec{\mathsf{s}})$. By hypotheses (g) and (h),
it takes values in the domains of $\hat{F}^{\mathsf{m}}$ and of $G^{(n)}_v$. Each member of
Definition~\ref{9.5:def-layer-store} is a finite sum of terms built from these compositions by the
piece operator, and is therefore holomorphic in $\xi$.

\emph{Support.} Suppose $\delta\hat{V}^{(n)}\not\equiv0$, so that some summand of
Definition~\ref{9.5:def-layer-store} is not identically zero. The zone $n$ is then restricted by
the following five facts.
\begin{itemize}
\item A family mark $\mathsf{g}^{(n)}_v$ vanishes identically unless zone $n$ meets $X_v$; this is
the support condition of Definition~\ref{9.5:def-resolved-terms}, contained in
\eqref{9.5:eq-resolved-terms-a}.
\item For the class marks of derived units the same conclusion is part (iii) of hypothesis (b).
\item A reading mark enters with the indicator $\mathbf{1}[n^c_v=n]$. Here $n^c_v$ is a level met
by the enlarged domain $X^+_v$ of the unit $v$, and $X^+_v\subset C^{\sharp}\cup(\text{zone }m)$.
\item By hypothesis (i), the boundary terms $\mathbf{B}$ created at the healing step carry only
levels strictly above the level of the healing step and not above $m$.
\item The terms $\mathfrak{v}^{\mathbf{C}}$ obey the locality of part (ii) of hypothesis (e).
\end{itemize}
These five facts give that the history of the label has zone $n$ inside $Y^{\sharp}$ with $n<m$,
or $n=m$.

\emph{Values.} Part (iv) of hypothesis (b) bounds the values of the kept member on the region
$Z^{\sim}$ by $2K_{\mathbb{C}}[\![\delta]\!]^{\natural+}_{n(z)}$ per unit of
$\sum_j|\Gamma^0_j\cap Y|$. Part (a) of
hypothesis (d) and the first three and sixth entries of the classification in hypothesis (c) bound
the remaining class-$n$ slots. By hypothesis (j), the kept potential is a plain sum of these slots,
so the slot bounds add. The summand $2K_{\mathbb{C}}$ of $K_{val}$ is the constant of part (iv) of
hypothesis (b); the summands $K(d)$ and $C_{rim}C^{\sharp}_L$ come from the bounds of hypotheses
(c) and (d). Adding the
bounds slot by slot, over the class-$n$ slots, gives the last line of
\eqref{9.5:eq-layer-store-bounds-a} with the constant $K_{val}$. The count there is the number of
cells of $Y^{\sharp}$ met by class-$n$ slots, so the bound is a volume.
\end{proof}

The depth decay of the one-run oscillation rests on Schwarz's lemma (Lemma~\ref{9.5:lem-disc-contraction}) applied on a disc of points of the term of a unit. In the two runs a unit is a function of its datum. In what follows the disc is placed in the data, inside $\mathfrak{X}_{X_v}$. Every function holomorphic there then inherits the decay, and not only the term of the unit.

We use the following objects of Notation~\ref{9.5:not-distance-units}:
\begin{itemize}
\item the two runs $\mathsf{A}$, $\mathsf{B}$ on one scaffold;
\item the core $\mathsf{C}$;
\item the scaled contour distance $d^{\wedge}$, counted in units of the small block size;
\item the lattice distance $d(\cdot,\cdot)$, which is always written with its arguments; bare $d$
denotes the dimension $4$ of the glossary.
\end{itemize}
From Notation~\ref{9.5:not-two-point-oscillation} we use the common data $\mathfrak{X}_X$, the tube data $\mathfrak{T}^{\mathsf{m}}_F[\varsigma]$ and the data domains $\mathfrak{D}^{\mathsf{m}}_F[\varsigma]$. When the run $\mathsf{m}$ is fixed we write $\mathfrak{D}_F[\varsigma]$ for $\mathfrak{D}^{\mathsf{m}}_F[\varsigma]$.

\emph{Deep units.} Let $v$ be a stored unit of the two runs, of one of the three kinds in which units are stored. We write $X_v$ for its domain, $X_v^{+}$ for the enlarged domain, $\hat{X}_v$ for the set of its hull and $F_v$ for its term. We call $v$ \emph{deep$^{+}$} if $X_v^{+}\subset\mathsf{C}$, and we put
\[
\hat{D}_v:=d^{\wedge}(X_v^{+},\mathsf{C}^c).
\]
A unit with $\hat{X}_v\subset\mathsf{C}$ but $X_v^{+}\not\subset\mathsf{C}$ is not deep$^{+}$. For such a unit $d^{\wedge}(\hat{X}_v,\mathsf{C}^c)\le c_+$, where $c_+$ is the $d^{\wedge}$-width of one layer of cubes of the unit's own scale; $c_+$ is a number depending on $d$, $M$ and $M_1$ only. Wherever sums over units carry the factor $e^{-\frac14\delta_P d^{\wedge}(\hat{X}_v,\mathsf{C}^c)}$ for such a unit, that factor is at least $e^{-\frac14\delta_P c_+}$. The constant $e^{\frac12\delta_P c_+}$ is included in $C_{\mathfrak{K}}$ below.

\emph{Levels.} For a site $x$ we use the following levels and units:
\begin{itemize}
\item $\iota(x)$ is the level at $x$ of $\mathbf{B}''_m$;
\item $\lambda_{\mathfrak{S}}(x)\le\iota(x)$ is the level of the set of the hull at $x$;
\item $\ell_i$ is the $M_1$-unit of level $i$, and $\ell_{\mathfrak{S}}(x):=\ell_{\lambda_{\mathfrak{S}}(x)}$.
\end{itemize}

\emph{Two production points.} Let $\mathfrak{z}=(\xi,\varpi,\vec{\mathsf{s}})$ and $\mathfrak{z}'=(\xi',\varpi,\vec{\mathsf{s}}')$ be production points with the same auxiliary parameters $\varpi$. Write $\mathrm{chain}^{\mathsf{m}}(\mathfrak{z})$ for the chain configuration of run $\mathsf{m}\in\{\mathsf{A},\mathsf{B}\}$ at $\mathfrak{z}$. The representation of two production points by one generator holds in run $\mathsf{m}$, also when the slices of the two points differ. It gives a generator $\mathcal{K}^{\mathsf{m}}$ with current $\jmath^{\mathsf{m}}$ such that
\[
\mathrm{chain}^{\mathsf{m}}(\mathfrak{z})=e^{i\eta^{\mathsf{m}}\mathcal{K}^{\mathsf{m}}}\,
\mathrm{chain}^{\mathsf{m}}(\mathfrak{z}')\qquad\text{on }D_2 .
\]
Here $D_2$ is the domain on which this representation holds, and $\eta^{\mathsf{m}}$ denotes the
multiplier of the generator in it; it is a different object from the read datum $\eta^{\mathsf{m}}_v$
of a unit.
For a cell $y$ we put
\[
N_y(\mathcal{K}^{\mathsf{m}}):=\sup_{\tilde{\Delta}(y)}
\max\{\ell|\mathcal{K}^{\mathsf{m}}|,\ \ell^2|\nabla\mathcal{K}^{\mathsf{m}}|\},
\qquad \ell=\ell_{\iota(y)} ,
\]
where the supremum is over the set $\tilde{\Delta}(y)$ attached to the cell $y$.
The bounds on $\mathcal{K}^{\mathsf{m}}$ and $\jmath^{\mathsf{m}}$ are the first two lines of \eqref{9.5:eq-dilation-disc-a} below. The last inequality there is the amplitude bound.

\emph{Transport of the read datum.} Apply the transport lemma for read data in run $\mathsf{m}$ to the configuration $\mathrm{chain}^{\mathsf{m}}(\mathfrak{z}')$ and the generator $\mathcal{K}^{\mathsf{m}}$. For a bond $\beta$ of level $l$ with endpoints $\beta_-$, $\beta_+$ it gives
\begin{equation}\label{9.5:eq-dilation-disc-1}
\begin{aligned}
\eta^{\mathsf{m}}_v(\mathfrak{z})(\beta)
&=\mathsf{w}(\beta_-)\,e^{i\mathfrak{d}^{\mathsf{m}}(\beta)}\,
\eta^{\mathsf{m}}_v(\mathfrak{z}')(\beta)\,\mathsf{w}(\beta_+)^{-1},\\
|\mathfrak{d}^{\mathsf{m}}(\beta)|&\le 2\ell_l\sup_{P(\beta)}|\mathcal{K}^{\mathsf{m}}|,
\qquad |\mathsf{w}-1|\le 16d\,\ell_l\sup|\mathcal{K}^{\mathsf{m}}| .
\end{aligned}
\end{equation}
Here $d=4$ is the dimension, and the first supremum is taken over the set $P(\beta)$ attached to
the bond $\beta$ in the transport lemma for read data. The dilated data are defined, for
$\lambda\in\mathbb{C}$, by
\begin{equation}\label{9.5:eq-dilation-disc-2}
\eta^{\mathsf{m}}_{v,\lambda}(\beta):=\mathsf{w}(\beta_-)^{\lambda}\,
e^{i\lambda\mathfrak{d}^{\mathsf{m}}(\beta)}\,\eta^{\mathsf{m}}_v(\mathfrak{z}')(\beta)\,
\mathsf{w}(\beta_+)^{-\lambda},
\qquad \mathsf{w}^{\lambda}:=e^{\lambda\log\mathsf{w}} .
\end{equation}
The mixed dilated datum is
\[
\eta^{\times}_{v,\lambda}:=\eta^{\mathsf{A}}_{v,\lambda}\oplus
(\text{the collar entries of }\eta^{\mathsf{B}}_{v,\lambda}).
\]
At $\lambda=0$ the dilated datum is $\eta^{\mathsf{m}}_v(\mathfrak{z}')$. At $\lambda=1$ it is $\eta^{\mathsf{m}}_v(\mathfrak{z})$, by \eqref{9.5:eq-dilation-disc-1}. So the two ends of the disc are the two read data.

\emph{Constants.} For a function $f$ of cells put
\[
(f)^{\approx}(x):=\sum_y e^{-2\delta_1 d(x,y)}f(y).
\]
The constant $B_\natural$ is that of the local Lipschitz lemma in the chart. That lemma says: for two complex data $B^1$, $B^2$ with $\sup|B^i|\le a'_R$, where $a'_R$ is the chart radius, in the
chart about a real critical centre, the field $E$, its gradient, and the current increment
$\delta J$ that represent the difference of the two data in the chart satisfy
\[
\ell_{\mathfrak{S}}|E|,\ \ell_{\mathfrak{S}}^2|\nabla E|,\ \ell_{\mathfrak{S}}^3|\delta J|
\le B_\natural\,(|B^2-B^1|)^{\approx} .
\]
Put $Y^a(x):=\mathsf{T}^a_N(x)$, the $a$-th entry of the scaffold variable $\mathsf{T}$ at $x$, the
lower index $N$ being the index with which the scaffold variable is defined, at points with
$\lambda_{\mathfrak{S}}(x)=\iota(x)$, and $Y^a(x):=\mathsf{u}^a_F(x)$, the $a$-th entry of the
quantity $\mathsf{u}_F$ of the data domains of the term $F$,
at points with $n_F(x)\le\lambda_{\mathfrak{S}}(x)<\iota(x)$, where $n_F(x)$ is the level attached
to $F$ at $x$.
Here $a$ indexes the entries of a datum, and $e_a$ is the power of $\ell$ attached to the $a$-th entry
in the chart. Inequalities between slices or displacements and $Y$ are meant entrywise, entry $a$
against $Y^a$.

Let $\mathsf{I}:=Z''_m\cup\mathfrak{K}$, where $\mathfrak{K}$ is the kept region and $Z''_m$ is the
region of the run hypotheses \eqref{9.5:eq-distance-units-c}. We use:
\begin{equation}\label{9.5:eq-dilation-disc-a}
\begin{aligned}
&N_y(\mathcal{K}^{\mathsf{m}})\le\mathsf{U}\,p(y),\qquad
\ell^3|\jmath^{\mathsf{m}}|\le\mathsf{U}\,p(y),\\
&p(y)\le e^{\delta_P\mathsf{w}_{\boxplus}}e^{-\delta_P d^{\wedge}(y,\mathsf{C}^c)},
\qquad \mathsf{U}\le c_{\mathsf{U}}C_{\mathsf{a}}\,\alpha_{*,m}
\qquad(\ell=\ell_{\iota(y)});\\
&\text{on }\mathsf{I}:\ \text{slices of read data}\le\tfrac38\,Y,\ \text{and a displacement}
\le\tfrac1{64}\,Y\ \text{keeps the datum in }\mathfrak{D}_F[\tfrac14];\\
&\alpha_{0,\iota(x)}\,\ell_{\mathfrak{S}}(x)^{-e_a}\,L^{\lambda_{\mathfrak{S}}(x)-\iota(x)}
\le\varphi_{\mathrm{arr}}^{-1}\,Y^a(x);\\
&e^{-\frac{5}{16}\delta_P d^{\wedge}(x,\mathsf{C}^c)}\,
\frac{\alpha_{*,m}}{\alpha_{*,\iota(x)}}\le C^{\natural}_{age}
:=2^{3/2}(1+B^{\sharp})^{1/2}\qquad(x\in\mathsf{C});\\
&r_{\mathfrak{K}}:=\min\Bigl\{r^0,\ \frac{\varphi_{\mathrm{arr}}}
{256L^2B_\natural c_1e^{\delta_P\mathsf{w}_{\boxplus}}c_{\mathsf{U}}C_{\mathsf{a}}
(C_*/C_0)C^{\natural}_{age}}\Bigr\},\\
&\rho_v:=r_{\mathfrak{K}}\,e^{\frac12\delta_P\hat{D}_v},\qquad
C_{\mathfrak{K}}:=\frac{2e^{\frac12\delta_Pc_+}}{r_{\mathfrak{K}}} .
\end{aligned}
\end{equation}

The lines of \eqref{9.5:eq-dilation-disc-a} are obtained as follows.
\begin{itemize}
\item The first two lines are the generator bound together with the amplitude bound.
\item The third line is the room condition. The set $\mathsf{I}$ contains $\mathsf{C}$ by \eqref{9.5:eq-distance-units-c}. The condition comes from the displacement, conversion and fraction bounds on the kept region. Those bounds allow a displacement budget $Y/16$. The smaller budget $Y/64$ is used here. Thus on $\mathsf{I}$ the room is a fixed fraction of $Y$, and not $\varrho_F\mathsf{u}_F$.
\item The fourth line is the conversion inequality from the same bounds; $\varphi_{\mathrm{arr}}$ is
the constant of that inequality. At points with $\lambda_{\mathfrak{S}}=\iota$ it is the comparison
of the scaffold variable with the level. At points with $n_F\le\lambda_{\mathfrak{S}}<\iota$ it rests
on the bound, from the same conversion bounds,
\[
\frac{\alpha_{0,\iota}}{\alpha_{a,n_F}}
\Bigl(\frac{\ell_{n_F}}{\ell_{\lambda_{\mathfrak{S}}}}\Bigr)^{e_a}
L^{-(\iota-\lambda_{\mathfrak{S}})}
\le\tfrac34 .
\]
\item The fifth line is the age inequality. It is derived below.
\end{itemize}

\emph{Derivation of the age inequality.} Let $x\in\mathsf{C}$. Let $\mathsf{n}_{\mathsf{C}}$ be the
depth of the core, so that its outermost layer, in which every contour from a point of $\mathsf{C}$
to $\mathsf{C}^c$ ends, has level $m-\mathsf{n}_{\mathsf{C}}$. Put
$\sigma:=m-\mathsf{n}_{\mathsf{C}}-\iota(x)\ge0$; then $m-\iota(x)=\mathsf{n}_{\mathsf{C}}+\sigma$.
\begin{itemize}
\item The ratio bound for the radii gives
\[
\alpha_{*,m}/\alpha_{*,\iota(x)}\le2(1+B^{\sharp})^{1/2}(1+\sigma)^{1/2} .
\]
\item A contour from $x$ to $\mathsf{C}^c$ ends in the layer of level $m-\mathsf{n}_{\mathsf{C}}$;
hence, by the layer
geometry of the core, $d^{\wedge}(x,\mathsf{C}^c)\ge w\,(\sigma-1)_+$. Here $w$ is the
$d^{\wedge}$-side of a cube of the own scale (Notation~\ref{9.5:not-distance-units}).
\item The floor $\frac18\delta_P w\ge2$ on the cube side, one of the conditions under which the
depth decay of the one-run oscillation is proved, gives $\frac{5}{16}\delta_P w\ge5$, and so
\[
e^{-\frac5{16}\delta_Pd^{\wedge}(x,\mathsf{C}^c)}\le e^{-5(\sigma-1)_+}.
\]
\item Finally,
\[
\sup_{\sigma\ge0}(1+\sigma)^{1/2}e^{-5(\sigma-1)_+}=2^{1/2}.
\]
The supremum is attained at $\sigma=1$. On $[0,1]$ the function increases. For $\sigma>1$ its
logarithmic derivative is $\frac1{2(1+\sigma)}-5<0$.
\end{itemize}
Multiplying the bounds gives the fifth line of \eqref{9.5:eq-dilation-disc-a}. It is the age inequality of the one-run oscillation, read at the rate $\frac5{16}\delta_P$ in place of $\frac12\delta_P$.

\emph{The constants.} The constants $r_{\mathfrak{K}}$ and $C_{\mathfrak{K}}$ are fixed before
$\gamma$. They depend only on $d$, $L$, $M$, $M_1$, $\kappa_1$, $C_0$, $C_1$ and $\delta_0$. Here $r^0$
is the radius factor of the dilation disc of the one-run oscillation.

\begin{lemma}[The dilation disc of a deep term]\label{9.5:lem-dilation-disc}
Assume the following hypotheses:
\begin{itemize}
\item the run hypotheses \eqref{9.5:eq-distance-units-c} for $\mathsf{m}=\mathsf{A}$ and for $\mathsf{m}=\mathsf{B}$, including $\mathsf{C}\subset\mathsf{I}$;
\item the scale condition \eqref{9.5:eq-distance-units-b};
\item the standing conditions on the data of the two runs, in particular the room condition (third line of \eqref{9.5:eq-dilation-disc-a}) and the flow condition on the mixed datum over the polydisc.
\end{itemize}
Assume further the following statements:
\begin{itemize}
\item the representation of two production points by one generator, with the amplitude bound (first
and second lines of \eqref{9.5:eq-dilation-disc-a}), and the ratio bound for the radii, from which
the fifth line of \eqref{9.5:eq-dilation-disc-a} is derived above;
\item the transport lemma for read data \eqref{9.5:eq-dilation-disc-1}, and centre-gauge invariance of the classes;
\item the local Lipschitz lemma in the chart;
\item the displacement, conversion and fraction bounds on the kept region;
\item the mixing lemma, with its room bound and the extension of that bound to the polydisc;
\item the kept-path lemma;
\item the riding-rotation lemma.
\end{itemize}
Let $v$ be a deep$^{+}$ unit. Let $\mathfrak{z}=(\xi,\varpi,\vec{\mathsf{s}})$ and $\mathfrak{z}'=(\xi',\varpi,\vec{\mathsf{s}}')$ be two production points with the same $\varpi$, with $\vec{\mathsf{s}},\vec{\mathsf{s}}'$ in the polydisc and $\xi,\xi'\in\mathfrak{X}_{\mathsf{C}^{\sharp}}$. Then for every $\lambda\in\mathbb{C}$ with $|\lambda|\le\rho_v$,
\[
\eta^{\mathsf{A}}_{v,\lambda}\in\mathfrak{T}^{\mathsf{A}}_{F_v}[\tfrac14],\qquad
\eta^{\mathsf{B}}_{v,\lambda}\in\mathfrak{T}^{\mathsf{B}}_{F_v}[\tfrac14],\qquad
\eta^{\times}_{v,\lambda}\in\mathfrak{X}_{X_v},
\]
and these data depend holomorphically on $\lambda$. Neither $\rho_v$ nor the conclusion depends on any
of the following: the age of any region, the depth parameter $N_0$, the widths $\check{R}_m$, $|X_v|$,
or the number or the history of the large-field regions enclosed in $\mathsf{C}$.
\end{lemma}

\begin{proof}
\emph{Step 1: the smeared increment.} Fix $\mathsf{m}\in\{\mathsf{A},\mathsf{B}\}$ and $x\in X_v$.
All cells $y$ of $\mathfrak{B}(X_v^{+})$, the set of cells making up $X_v^{+}$, lie in $\mathsf{C}$,
because $v$ is deep$^{+}$. Since $y\subset X_v^{+}$, they satisfy
$d^{\wedge}(y,\mathsf{C}^c)\ge\hat{D}_v$.

Let $y$ be such a cell, of level $l(y)\le\iota(y)$. By \eqref{9.5:eq-dilation-disc-1} and the first line of \eqref{9.5:eq-dilation-disc-a}, on the bonds of $y$,
\[
|\mathfrak{d}^{\mathsf{m}}(y)|\le2\ell_l\sup|\mathcal{K}^{\mathsf{m}}|
\le2L^{l(y)-\iota(y)}\,\mathsf{U}\,p(y).
\]
We split the exponential in the bound on $p(y)$:
\begin{equation}\label{9.5:eq-dilation-disc-3}
\begin{aligned}
e^{-\delta_Pd^{\wedge}(y,\mathsf{C}^c)}
&\le e^{-\frac12\delta_P\hat{D}_v}\,e^{-\frac5{16}\delta_Pd^{\wedge}(y,\mathsf{C}^c)}\\
&\le e^{-\frac12\delta_P\hat{D}_v}\,e^{-\frac5{16}\delta_Pd^{\wedge}(x,\mathsf{C}^c)}
\,e^{\frac5{16}\delta_Pd^{\wedge}(x,y)}\\
&\le e^{-\frac12\delta_P\hat{D}_v}\,e^{-\frac5{16}\delta_Pd^{\wedge}(x,\mathsf{C}^c)}
\,e^{\frac54\delta_1d(x,y)} .
\end{aligned}
\end{equation}
The three inequalities hold for the following reasons.
\begin{itemize}
\item First: $\frac12+\frac5{16}\le1$ and $d^{\wedge}(y,\mathsf{C}^c)\ge\hat{D}_v$.
\item Second: the triangle inequality $d^{\wedge}(x,\mathsf{C}^c)\le d^{\wedge}(x,y)+d^{\wedge}(y,\mathsf{C}^c)$.
\item Third: by \eqref{9.5:eq-distance-units-a} and \eqref{9.5:eq-distance-units-b},
\[
\delta_Pd^{\wedge}(x,y)\le\delta_Pd(x,y)/M_1\le\tfrac14\delta_1d(x,y),
\]
so $\frac5{16}\delta_Pd^{\wedge}(x,y)\le\frac5{64}\delta_1d(x,y)\le\frac54\delta_1d(x,y)$.
\end{itemize}
The level comparison, namely the first distance property of Notation~\ref{9.5:not-distance-units}
together with the floor condition used in the conversion bounds on the kept region, gives
\[
L^{l(y)-\iota(y)}\le L^2e^{\delta_1d(x,y)/8}L^{\lambda_{\mathfrak{S}}(x)-\iota(x)} .
\]
The weight of the smear and these two factors combine as
\[
e^{-2\delta_1d(x,y)+\frac54\delta_1d(x,y)+\frac18\delta_1d(x,y)}=e^{-\frac58\delta_1d(x,y)}
\le e^{-\frac12\delta_1d(x,y)}.
\]
Moreover, by the summability bound for the lattice distance, the second distance property of
Notation~\ref{9.5:not-distance-units}, with the summability constant $c_1$,
$\sum_ye^{-\delta_1d(x,y)/2}\le c_1$. Inserting all of this into the definition of
$(\,\cdot\,)^{\approx}$ gives
\begin{equation}\label{9.5:eq-dilation-disc-4}
(|\mathfrak{d}^{\mathsf{m}}|)^{\approx}(x)\le2L^2c_1\mathsf{U}e^{\delta_P\mathsf{w}_{\boxplus}}
\cdot e^{-\frac12\delta_P\hat{D}_v}\cdot e^{-\frac5{16}\delta_Pd^{\wedge}(x,\mathsf{C}^c)}
\cdot L^{\lambda_{\mathfrak{S}}(x)-\iota(x)} .
\end{equation}

\emph{Step 2: the displacement.} Put
\[
D^a:=2\ell_{\mathfrak{S}}^{-e_a}B_\natural(|\mathfrak{d}^{\mathsf{m}}|)^{\approx},
\]
as in the displacement bounds on the kept region, and let
$|\lambda|\le\rho_v=r_{\mathfrak{K}}e^{\frac12\delta_P\hat{D}_v}$. The factor
$e^{\frac12\delta_P\hat{D}_v}$ cancels the factor $e^{-\frac12\delta_P\hat{D}_v}$ of
\eqref{9.5:eq-dilation-disc-4}. We get
\begin{equation}\label{9.5:eq-dilation-disc-5}
\begin{aligned}
|\lambda|D^a(x)
&\le r_{\mathfrak{K}}\cdot4L^2B_\natural c_1e^{\delta_P\mathsf{w}_{\boxplus}}\cdot\mathsf{U}\,
e^{-\frac5{16}\delta_Pd^{\wedge}(x,\mathsf{C}^c)}
\,\ell_{\mathfrak{S}}(x)^{-e_a}L^{\lambda_{\mathfrak{S}}(x)-\iota(x)}\\
&\le r_{\mathfrak{K}}\cdot4L^2B_\natural c_1e^{\delta_P\mathsf{w}_{\boxplus}}
c_{\mathsf{U}}C_{\mathsf{a}}C^{\natural}_{age}(C_*/C_0)
\cdot\alpha_{0,\iota(x)}\,\ell_{\mathfrak{S}}(x)^{-e_a}L^{\lambda_{\mathfrak{S}}(x)-\iota(x)}\\
&\le r_{\mathfrak{K}}\cdot4L^2B_\natural c_1e^{\delta_P\mathsf{w}_{\boxplus}}
c_{\mathsf{U}}C_{\mathsf{a}}C^{\natural}_{age}(C_*/C_0)
\cdot\varphi_{\mathrm{arr}}^{-1}\,Y^a(x)\\
&\le Y^a(x)/64 .
\end{aligned}
\end{equation}
The steps are justified as follows.
\begin{itemize}
\item The first inequality is \eqref{9.5:eq-dilation-disc-4}.
\item The second uses three facts: the amplitude bound $\mathsf{U}\le c_{\mathsf{U}}C_{\mathsf{a}}\alpha_{*,m}$; the age inequality, fifth line of \eqref{9.5:eq-dilation-disc-a}, which applies because $x\in X_v\subset\mathsf{C}$; and the relation $\alpha_{*,\iota}=(C_*/C_0)\alpha_{0,\iota}$ between the two families of radii, valid with $q_0=q_1=q$.
\item The third is the conversion inequality, fourth line of \eqref{9.5:eq-dilation-disc-a}.
\item The last uses the definition of $r_{\mathfrak{K}}$: it makes the coefficient at most $4/256=1/64$.
\end{itemize}
The local Lipschitz lemma requires the chart condition $\sup|B_\lambda|\le a'_R$ for the complex datum $B_\lambda$ of the dilated datum. This condition holds because $|\lambda|\sup|\mathfrak{d}^{\mathsf{m}}|$ is bounded by the same tube-small number. That follows from the chart-approximation bound of the displacement bounds on the kept region.

\emph{Step 3: fractions of the room.} We first treat the data of one run.
\begin{itemize}
\item The base datum $\eta^{\mathsf{m}}_v(\mathfrak{z}')$ has slice at most $\frac38Y$. This is the room condition of \eqref{9.5:eq-dilation-disc-a}, which holds at the points of the polydisc by the room bound of the mixing lemma and its extension to the polydisc.
\item By the local Lipschitz lemma in run $\mathsf{m}$, the slice of the dilated datum $\eta^{\mathsf{m}}_{v,\lambda}$ is within $|\lambda|D^a$ of the slice of the base datum.
\item By \eqref{9.5:eq-dilation-disc-5} that slice is therefore at most
\[
(\tfrac38+\tfrac1{64})Y=\tfrac{25}{64}Y.
\]
\end{itemize}
Next we treat the mixed datum.
\begin{itemize}
\item The mixed base $\eta^{\times}_{0}=\mathsf{v}\oplus c^{\mathsf{B}}$ is the initial point of the
path of the kept-path lemma.
\item Its slice is within the net displacement $D^{net}$ of the slice of the datum of $\mathsf{B}$.
The displacement, conversion and fraction bounds on the kept region give
$8C^{\mathfrak{K}}_{TL}\vartheta^{\mathfrak{K}}\le1$, where $C^{\mathfrak{K}}_{TL}$ is the constant
of the kept-path lemma and $\vartheta^{\mathfrak{K}}$ the smallness index of the kept region; hence
\[
D^{net}\le\tfrac14C^{\mathfrak{K}}_{TL}\vartheta^{\mathfrak{K}}Y\le\tfrac1{32}Y.
\]
This holds at the vertices of the polydisc by part (a) of the kept-path lemma, and on the whole
polydisc by the flow condition on the mixed datum.
\item The slice of the mixed base is therefore at most $\frac{13}{32}Y$.
\item For $\eta^{\times}_{v,\lambda}$, apply the local Lipschitz lemma in run $\mathsf{B}$ to the increment $\lambda\bigl(\mathfrak{d}^{\mathsf{A}}\oplus\mathfrak{d}^{\mathsf{B}}\upharpoonright\text{collar}\bigr)$.
\item Each entry of this increment obeys \eqref{9.5:eq-dilation-disc-5}. The reason is that the bound in the first two lines of \eqref{9.5:eq-dilation-disc-a} concerns the common scaffold and holds in both runs.
\item The slice of $\eta^{\times}_{v,\lambda}$ is therefore at most
\[
(\tfrac{13}{32}+\tfrac1{64})Y=\tfrac{27}{64}Y.
\]
\end{itemize}
Both bounds are at most $\frac7{16}Y$. A slice of this size lies in $\mathfrak{D}_F[\frac14]$: by the
description of $\mathfrak{D}_F[\frac14]$ in Notation~\ref{9.5:not-two-point-oscillation}, in which
$\mathsf{c}$ is the room coefficient and $\varrho_F$ the generic room, this amounts to
\[
\tfrac7{16}\mathsf{c}\le\mathsf{c}-\tfrac34\varrho_F,
\]
which holds when $\mathsf{c}\ge\frac1{30}$; this reduction is the one made in the displacement,
conversion and fraction bounds on the kept region.

\emph{Step 4: the gauge factors.} The conjugation by $\mathsf{w}^{\lambda}$ in
\eqref{9.5:eq-dilation-disc-2} is a centre gauge with values in the complexification $G^{c}$ of the
gauge group. By \eqref{9.5:eq-dilation-disc-1} and the first line of
\eqref{9.5:eq-dilation-disc-a}, its size is
\[
|\lambda||\log\mathsf{w}|\le18d\,|\lambda|\sup L^{l-\iota}\,\mathsf{U}\,p .
\]
For invariant units it does not change the class, by centre-gauge invariance of the classes. For
units with slots and for inner kept members it is the rotation of the riders. By the riding-rotation
lemma, applied in each run, this rotation stays within one half of the one-run margin, in each run.

\emph{Step 5: holomorphy.} In \eqref{9.5:eq-dilation-disc-2}, $\lambda$ enters only through $e^{\lambda\log\mathsf{w}}$ and $e^{i\lambda\mathfrak{d}^{\mathsf{m}}}$. Both are entire exponentials in $\lambda$.

\emph{Conclusion.} Steps 2--5 give the three memberships for $|\lambda|\le\rho_v$ and their holomorphy
in $\lambda$. The only quantity attached to $v$ that enters the constants is $\hat{D}_v$.

The rates in this proof are divided as follows.
\begin{itemize}
\item Of the rate $\delta_P$ carried by $p$: one half goes to the dilation, $\frac5{16}$ pays the age ratio pointwise, and $\frac3{16}$ is not used.
\item Of the rate $2\delta_1$ of the smear: $\frac54\delta_1$ transports the pointwise factor, $\frac18\delta_1$ pays the level comparison, and the remaining $\frac58\delta_1\ge\frac12\delta_1$ pays the summation.
\end{itemize}
The lemma contains no estimate of a difference between the two runs.
\end{proof}

The pair rate of the decoupling parameters is
\begin{equation}\label{9.5:eq-family-mark-sums-1}
\kappa^{\natural}_1:=\log\Bigl(\tfrac14\Bigl(\tfrac12\bigl(e^{\kappa_1-2}-1\bigr)^{1/2}-1\Bigr)\Bigr),
\end{equation}
defined whenever the argument of the logarithm is positive. For $\kappa_1\ge4$ one has
$\kappa^{\natural}_1\ge\tfrac12\kappa_1-5$. Indeed, put $u:=e^{\kappa_1/2}\ge e^2$ and
\begin{equation*}
h(u):=\tfrac12\bigl(e^{-2}u^2-1\bigr)^{1/2}-1-4e^{-5}u .
\end{equation*}
The claimed bound is $h(u)\ge0$. One has $h(e^2)>0$, since $\tfrac12(e^2-1)^{1/2}>1.26>1+4e^{-3}$.
Moreover $h'(u)\ge(2e)^{-1}-4e^{-5}>0$ for $u\ge e^2$, because
$e^{-2}u(e^{-2}u^2-1)^{-1/2}\ge e^{-1}$.
Since $\tfrac12(e^{\kappa_1-2}-1)^{1/2}-1<\tfrac12e^{\kappa_1/2-1}$, one has
$\kappa^{\natural}_1\le\tfrac12\kappa_1-1-\log 8$.

Squaring shows that $(2/(e^{\kappa_1}-1))^{1/2}\le e^{-\frac12(\kappa_1-1-\log 2)}$ is equivalent to
$e\le(e-1)e^{\kappa_1}$, which holds for $\kappa_1\ge1$. Moreover
$\tfrac12\kappa_1-1-\log 8<\tfrac12(\kappa_1-1-\log 2)$. Together these give, whenever
$\kappa^{\natural}_1$ is defined (in particular for $\kappa_1\ge4$),
\begin{equation}\label{9.5:eq-family-mark-sums-2}
\Bigl(\frac{2}{e^{\kappa_1}-1}\Bigr)^{1/2}\le e^{-\frac12(\kappa_1-1-\log 2)}\le e^{-\kappa^{\natural}_1}.
\end{equation}

The pair-rate floor on $\kappa_1$ is the pair of lower bounds \eqref{9.5:eq-family-mark-sums-a} below.
Its constants are the following.
\begin{itemize}
\item $c_l$ and $\lambda_*$ are the constants in the condition under which the four free sums
of hypothesis (e) of Lemma~\ref{9.5:lem-family-mark-sums} hold; both are non-negative, as fixed
with that total bound over the unit lists.
\item $a_k$, $c_{\flat}$, $\hat b_t$ and $N^{\flat}$ are the constants of the floor on $\kappa_1$
required by the family sums at the production points.
\item $\kappa$ is Ba{\l}aban's auxiliary parameter, positive as chosen in the order of choice
(Notation~\ref{2.3:not-stages-middle}).
\end{itemize}
The floor reads
\begin{equation}\label{9.5:eq-family-mark-sums-a}
\begin{aligned}
\kappa^{\natural}_1&\ge c_l\max\bigl\{2\kappa+4,\ \tfrac{31}{20}\kappa+9\bigr\}+\lambda_*,\\
\kappa^{\natural}_1&\ge(a_k+1)+3^d c_{\flat}+\log\bigl(2e^2\hat b_t^{\,2}N^{\flat}\bigr).
\end{aligned}
\end{equation}
Its features are as follows.
\begin{itemize}
\item It is one-sided, with three right members.
\item Its place in the order of choice of the constants (Notation~\ref{2.3:not-stages-middle}) is
after $L$ and $\kappa$ and before $M_1$ and $M$.
\item The last member is the floor required by the family sums, taken at the pair rate.
\item It has the shape of the floor on $\kappa_1$ of the one-run bounds at the production points,
with $\kappa_1-1$ replaced by $\kappa^{\natural}_1$.
\item It holds as soon as $\kappa_1\ge 2R_{\max}+10$, where $R_{\max}$ is the largest of its three
right members. Indeed, $R_{\max}$ is at least the first right member, which is non-negative. Hence
$\kappa_1\ge10$. The bound following \eqref{9.5:eq-family-mark-sums-1} then gives
$\kappa^{\natural}_1\ge\tfrac12\kappa_1-5\ge R_{\max}$.
\end{itemize}

We use the following extension of the size bounds to all functions in a unit's format. It is
carried below as hypothesis (b) of Lemma~\ref{9.5:lem-family-mark-sums}. For every unit $v$ of the
second type in the table of size functionals of the total bound over the unit lists (this type is
one of the types of pointed units), and every holomorphic function $G$ in the format of $v$
(Definition~\ref{9.5:def-resolved-terms}), one has, on a disc of data furnished by
Lemma~\ref{9.5:lem-dilation-disc},
\begin{equation}\label{9.5:eq-family-mark-sums-b}
\sup\bigl|G-c_G\bigr|\le\mathfrak{N}_v[G],
\end{equation}
where $c_G$ is the value of $G$ at the trivial datum, in the form of the size functional of pointed
units.

Here $\mathfrak{N}_v[G]$ denotes the bound given by the table of the size lemmas for units of the
second type. In that table, the three prefactors of the row, namely the size constant
$\mathfrak{n}_v$, the constant $E_0$ and the counterterm prefactor, are the size constants of the
unit $v$ itself. They are replaced by the corresponding quantities of $G$, namely $\|G\|$,
$\|G^+\|^c$ and the $\delta$-entries of the counterterm of $G$.

The extension \eqref{9.5:eq-family-mark-sums-b} enters
Lemma~\ref{9.5:lem-family-mark-sums} as a hypothesis. It rests on the following features of the two
size lemmas.
\begin{itemize}
\item The size lemma for analytic functions is stated for every function that is analytic and invariant
on the small-field domain and bounded by $\mathfrak{n}$.
\item The size lemma for pointed families depends on the family only through a lemma stated for
pointed families.
\item Both lemmas depend on the background only through the three regularity bounds of the
background.
\item These regularity bounds hold on the disc. They depend on the background only through the
threshold and base conditions, and on $X$ only through the domain on which $X$ is defined.
\end{itemize}

Units of the first and of the fourth type need no such extension. For them:
\begin{itemize}
\item on units of the first type, $\mathfrak{N}_v[G]:=\sup|G|$, taken per strong set $\mathcal{S}$
and summed with the weight $e^{\kappa_A|\mathcal{S}|}$, where strong sets and the rate $\kappa_A$ are
those of the table of size functionals;
\item on units of the fourth type, $\mathfrak{N}_v[G]:=\mathfrak{O}^{+}(G)$, in the notation of
\eqref{9.5:eq-two-point-oscillation-b}.
\end{itemize}
In both cases Lemma~\ref{9.5:lem-dilation-disc} and Lemma~\ref{9.5:lem-disc-contraction} supply
everything used below.

\begin{lemma}[Tuple sums of the family marks]\label{9.5:lem-family-mark-sums}
Assume the following.
\begin{enumerate}
\item[(a)] The hypothesis \eqref{9.5:eq-distance-units-c} holds for the run $\mathsf{A}$.
\item[(b)] The extension \eqref{9.5:eq-family-mark-sums-b} holds.
\item[(c)] The conclusion of Lemma~\ref{9.5:lem-dilation-disc} holds (the dilation disc of a deep
term, proved as an implication from the hypotheses listed there).
\item[(d)] The floor \eqref{9.5:eq-family-mark-sums-a} holds.
\item[(e)] The four free sums of the total bound over the unit lists hold in the following form. For
every resolved stored unit $v$ and every $\mu$ with $\mu-3c_l\ge\lambda_*$,
\begin{equation*}
\sum_{\vec{\mathsf{y}}}e^{-\mu\sum_t|\mathsf{y}_t|}\le e^{4(d_m(K_v)+1)},
\end{equation*}
where the sum runs over all tuples $\vec{\mathsf{y}}=(\mathsf{y}_t)_t$ of toggle sets of $v$, the
index $t$ labelling the four free sums.
\end{enumerate}
Let $v$ be a resolved stored unit in the sense of Definition~\ref{9.5:def-resolved-terms}, and let
$n$ be a class. Put
\begin{equation}\label{9.5:eq-family-mark-sums-3}
\operatorname{tot}^{u,(n)}(v):=\sum_{\vec{\mathsf{y}}}
\mathfrak{O}^{+}\bigl(\mathfrak{p}_{v,\vec{\mathsf{y}}}[\mathsf{g}^{(n)}_v]\bigr),
\end{equation}
the sum running over all tuples $\vec{\mathsf{y}}$ of hypothesis (e), the empty tuple included; for
$\Delta\in\mathsf{y}_t$ we write $\mathsf{s}^{(t)}(\Delta)$ for the decoupling parameter of $\Delta$
in the component $\mathsf{y}_t$. Then
\begin{equation}\label{9.5:eq-family-mark-sums-4}
\operatorname{tot}^{u,(n)}(v)\le\mathfrak{n}^{u,(n)}(v)\,e^{4(d_m(K_v)+1)},
\end{equation}
where
\begin{equation}\label{9.5:eq-family-mark-sums-5}
\mathfrak{n}^{u,(n)}(v):=
\begin{cases}
2\,\mathfrak{N}_v[G^{(n)}_v] & \text{if $v$ is not deep$^{+}$},\\[2pt]
C_{\mathfrak{K}}\,e^{-\frac12\delta_P\hat D_v}\,\mathfrak{N}_v[G^{(n)}_v] & \text{if $v$ is deep$^{+}$}.
\end{cases}
\end{equation}
Here deep$^{+}$ is meant in the sense of Lemma~\ref{9.5:lem-dilation-disc}. The quantity
$d_m(K_v)$ is the distance function $d_j$ at level $j=m$, evaluated at the set $K_v$ of the unit $v$.
The quantity $\hat D_v$ is the distance of the deep$^{+}$ unit $v$ which enters, together with
$\rho_v$ and $C_{\mathfrak{K}}$, the relation \eqref{9.5:eq-dilation-disc-a}.
\end{lemma}

\begin{proof}
Fix the variable $\Phi_C$ of \eqref{9.5:eq-two-point-oscillation-b}; the bounds below are uniform
in it. Fix also the auxiliary parameters $\varpi$ and two common data $\xi,\xi'$. Put
\begin{equation}\label{9.5:eq-family-mark-sums-6}
G(\vec{\mathsf{s}}):=\mathsf{g}^{(n)}_v(\xi,\varpi,\vec{\mathsf{s}})
-\mathsf{g}^{(n)}_v(\xi',\varpi,\vec{\mathsf{s}}).
\end{equation}

\emph{Holomorphy of $G$.} The family mark $\mathsf{g}^{(n)}_v$ is $G^{(n)}_v$ evaluated at the mixed
datum $\eta^{\times}_v$. We use three facts.
\begin{itemize}
\item The mixed datum is holomorphic on the product polydisc
$\{|\mathsf{s}^{(t)}(\Delta)|\le e^{\kappa_1}\}$.
\item On that polydisc the mixed datum takes its values in $\mathfrak{X}_{X_v}$, by the statements on
the mixed datum and on its tube.
\item $G^{(n)}_v$ is holomorphic on $\mathfrak{X}_{X_v}$ (Definition~\ref{9.5:def-resolved-terms}).
\end{itemize}
Hence $G$ is holomorphic on this polydisc.

\emph{The bound on $G$: $v$ not deep$^{+}$.} The function $G$ is a difference of two values of
$G^{(n)}_v$. Subtract from each value a constant $c$ for which
$\sup|G^{(n)}_v-c|\le\mathfrak{N}_v[G^{(n)}_v]$. For units of the second type, $c$ is the value of
$G^{(n)}_v$ at the trivial datum, in the form of the pointed size functional, and the supremum bound is
\eqref{9.5:eq-family-mark-sums-b}. This gives
\begin{equation}\label{9.5:eq-family-mark-sums-7}
|G|\le 2\sup|G^{(n)}_v-c|\le 2\,\mathfrak{N}_v[G^{(n)}_v].
\end{equation}

\emph{The bound on $G$: $v$ deep$^{+}$.} Fix $\vec{\mathsf{s}}$ and apply
Lemma~\ref{9.5:lem-dilation-disc} with $\vec{\mathsf{s}}'=\vec{\mathsf{s}}$.
\begin{itemize}
\item The function $f(\lambda):=G^{(n)}_v(\eta^{\times}_{v,\lambda})$ is holomorphic on
$|\lambda|<\rho_v$.
\item On that disc, $|f(\lambda)-f(0)|\le 2\,\mathfrak{N}_v[G^{(n)}_v]$, as in
\eqref{9.5:eq-family-mark-sums-7}.
\item The values $f(1)$ and $f(0)$ are the two values whose difference is $G(\vec{\mathsf{s}})$.
\end{itemize}
By the choice of $\rho_v$ and $C_{\mathfrak{K}}$ in \eqref{9.5:eq-dilation-disc-a}, one has
$2\rho_v^{-1}\le C_{\mathfrak{K}}e^{-\frac12\delta_P\hat D_v}$ for every deep$^{+}$ unit, whether or
not $\rho_v>1$. If $\rho_v>1$, Lemma~\ref{9.5:lem-disc-contraction} and this inequality give
\begin{equation}\label{9.5:eq-family-mark-sums-8}
|G|\le\frac{2\,\mathfrak{N}_v[G^{(n)}_v]}{\rho_v}
\le C_{\mathfrak{K}}\,e^{-\frac12\delta_P\hat D_v}\,\mathfrak{N}_v[G^{(n)}_v].
\end{equation}
If $\rho_v\le 1$, the bound \eqref{9.5:eq-family-mark-sums-7} gives the same conclusion. Indeed, the
same inequality then yields
\begin{equation*}
2\le 2\rho_v^{-1}\le C_{\mathfrak{K}}e^{-\frac12\delta_P\hat D_v}.
\end{equation*}

In either case $|G|\le\mathfrak{n}^{u,(n)}(v)$ on the whole polydisc.

\emph{The pieces.} The operator $\mathfrak{p}_{v,\vec{\mathsf{y}}}$ acts only on the decoupling
parameters and is linear. Therefore
\begin{equation*}
\mathfrak{p}_{v,\vec{\mathsf{y}}}[\mathsf{g}^{(n)}_v](\xi,\varpi)
-\mathfrak{p}_{v,\vec{\mathsf{y}}}[\mathsf{g}^{(n)}_v](\xi',\varpi)
=\mathfrak{p}_{v,\vec{\mathsf{y}}}[G].
\end{equation*}
We apply Lemma~\ref{9.5:lem-polydisc-vertices}(ii) to the set of parameters
$\{(t,\Delta):\Delta\in\mathsf{y}_t\}$, with $R=e^{\kappa_1}$. One has
$e^{\kappa_1}-1\ge e^{\kappa_1-1}$ for $\kappa_1\ge1$, and $\kappa_1\ge1$ by
\eqref{9.5:eq-family-mark-sums-10} below. We therefore obtain
\begin{equation*}
|\mathfrak{p}_{v,\vec{\mathsf{y}}}[G]|
\le\bigl(e^{\kappa_1}-1\bigr)^{-\sum_t|\mathsf{y}_t|}\sup|G|
\le\mathfrak{n}^{u,(n)}(v)\,e^{-(\kappa_1-1)\sum_t|\mathsf{y}_t|}.
\end{equation*}
Take the supremum defining $\mathfrak{O}^{+}$ in \eqref{9.5:eq-two-point-oscillation-b}, at the same
$\varpi$; the bound above is uniform in $\Phi_C,\xi,\xi'$. This gives
\begin{equation}\label{9.5:eq-family-mark-sums-9}
\mathfrak{O}^{+}\bigl(\mathfrak{p}_{v,\vec{\mathsf{y}}}[\mathsf{g}^{(n)}_v]\bigr)
\le\mathfrak{n}^{u,(n)}(v)\,e^{-(\kappa_1-1)\sum_t|\mathsf{y}_t|}.
\end{equation}

\emph{The sums.} Sum \eqref{9.5:eq-family-mark-sums-9} over all tuples $\vec{\mathsf{y}}$. The
resulting sum of exponentials is handled by hypothesis (e) at $\mu=\kappa_1-1$. We check that this
rate is admissible.
\begin{itemize}
\item By the first line of \eqref{9.5:eq-family-mark-sums-a} and the non-negativity of $c_l$ and
$\kappa$, one has $\kappa^{\natural}_1\ge c_l(2\kappa+4)+\lambda_*$. The right member is at least
$3c_l+\lambda_*$, and $3c_l+\lambda_*$ is non-negative.
\item By the upper bound $\kappa^{\natural}_1\le\tfrac12\kappa_1-1-\log8$ stated after
\eqref{9.5:eq-family-mark-sums-1}, one has $\kappa_1-1\ge2\kappa^{\natural}_1+1+2\log8$.
\end{itemize}
Combining the two items,
\begin{equation}\label{9.5:eq-family-mark-sums-10}
\kappa_1-1\ge2\kappa^{\natural}_1+1+2\log8\ge\kappa^{\natural}_1\ge3c_l+\lambda_*,
\end{equation}
that is, $\kappa_1-1-3c_l\ge\lambda_*$; in particular $\kappa_1\ge1$. Hypothesis (e) therefore gives
\begin{equation*}
\sum_{\vec{\mathsf{y}}}e^{-(\kappa_1-1)\sum_t|\mathsf{y}_t|}\le e^{4(d_m(K_v)+1)}.
\end{equation*}
Inserting this into the summed form of \eqref{9.5:eq-family-mark-sums-9} gives
\eqref{9.5:eq-family-mark-sums-4}.

In this argument no two points produced by the two runs are compared on the disc. Indeed,
$G^{(n)}_v$ is a single function. Its argument is the datum of the run $\mathsf{A}$ completed by the
collar entries of the run $\mathsf{B}$, and each of these is moved by its own run.
\end{proof}

\begin{lemma}[Chain terms of the current step are not deep]\label{9.5:lem-chain-terms}
Let $v=\mathbb{C}^{(\nu)}_m(X)$ be a chain term of the current step, of stage $j$, regarded as a
derived unit, with $K_v$ as in Lemma~\ref{9.5:lem-family-mark-sums}. For a tuple
$\vec{\mathsf{y}}$ write $u=(\mathbb{C}^{(\nu)}_m(X),\vec{\mathsf{y}})$ for the corresponding piece, and
$W^{(\zeta)}_u$ for its class marks, the class $\zeta$ ranging over the labels $((n,m),\mathsf{i})$, in
which $m$ is the index of the chain term $\mathbb{C}^{(\nu)}_m(X)$.
Assume the following.
\begin{enumerate}
\item The chain term vanishes at an empty set of integration bonds: $\mathbb{C}^{(\nu)}_m(X)=0$
unless $X$ contains an integration bond of stage $j$. This is the statement that the chain terms
vanish at an empty set of integration bonds.
\item The integration bonds of stage $j$ lie off $Z''_{j+1}$: those integrated in $B_j$ lie in
$(\Omega^c_{j+1}\setminus Z''_{j+1})^{(j)*}$, where $\Omega^c_{j+1}$ is the complement of the
small-field region $\Omega_{j+1}$, and those integrated in $A_j$ lie on the rim. This is the
statement on the location of the integration bonds of the step.
\item The inclusions
\begin{equation}\label{9.5:eq-chain-terms-2}
Z''_{j+1}\supseteq Z''_{h+1}\supset D_2\supset\mathsf{C}
\end{equation}
hold, where $h$ and $D_2$ are the index and the domain fixed with the core $\mathsf{C}$ of the
component at the healing site. This is the statement on the inclusions of the domains at the
healing site.
\item For every interior piece $u$ of the chain term and every class $\zeta$, at every point of the
common data $\mathfrak{X}_X$,
\begin{equation*}
|W^{(\zeta)}_u|\le e^{-\kappa^{\natural}_1|\vec{\mathsf{y}}|}\,
[\![(\delta\mathbb{C}^{(\nu)}(X))^{\natural,\zeta}]\!],
\end{equation*}
uniformly in the auxiliary parameters $\varpi$, of which the seed is one component. Here
$(\delta\mathbb{C}^{(\nu)}(X))^{\natural,\zeta}$ is the resolved $\natural$-extension, in class
$\zeta$, of the two-run difference $\delta\mathbb{C}^{(\nu)}(X)$ of the chain term, and
$[\![\cdot]\!]$ is its size functional. This is the alternating-sum bound for interior pieces.
\item The condition \eqref{9.5:eq-family-mark-sums-a} holds.
\end{enumerate}
Then:
\begin{itemize}
\item[(a)] $\mathbb{C}^{(\nu)}_m(X)$ is never deep: it vanishes unless $X$ contains an integration
bond of its stage, and these bonds lie off $Z''_{j+1}$, hence, by \eqref{9.5:eq-chain-terms-2},
off $Z''_{h+1}$, off $D_2$ and off the core $\mathsf{C}$.
\item[(b)] For every $n$, with $m$ the index of the chain term,
\begin{equation}\label{9.5:eq-chain-terms-1}
\sum_{\vec{\mathsf{y}}}2\sup\Bigl|\sum_{\mathsf{i}}W^{((n,m),\mathsf{i})}_u\Bigr|
\le 2\sum_{\mathsf{i}}
[\![(\delta\mathbb{C}^{(\nu)}(X))^{\natural,((n,m),\mathsf{i})}]\!]\,
e^{4(d_m(K_v)+1)} ,
\end{equation}
where the supremum is taken over the common data $\mathfrak{X}_X$ at one value of the inner
variables; twice this supremum bounds the two-point oscillation $\mathfrak{O}^{+}$ over the same
data that enters Lemma~\ref{9.5:lem-family-mark-sums}, and $d_m(K_v)$ carries the same index $m$.
\end{itemize}
\end{lemma}

\begin{proof}
(a) By the first hypothesis, $\mathbb{C}^{(\nu)}_m(X)$ is zero unless $X$ contains an integration bond
of stage $j$. By the second hypothesis every such bond lies off $Z''_{j+1}$: a bond integrated in
$B_j$ lies in $(\Omega^c_{j+1}\setminus Z''_{j+1})^{(j)*}$, and a bond integrated in $A_j$ lies on
the rim. By the inclusions \eqref{9.5:eq-chain-terms-2} of the third hypothesis, $Z''_{j+1}$ contains
$Z''_{h+1}$, which contains $D_2$, which contains the core $\mathsf{C}$; a bond lying off
$Z''_{j+1}$ therefore lies off each of $Z''_{h+1}$, $D_2$ and $\mathsf{C}$. Hence a non-vanishing
chain term $\mathbb{C}^{(\nu)}_m(X)$ always holds an integration bond of its stage outside
$Z''_{j+1}$, and therefore outside $Z''_{h+1}$, $D_2$ and $\mathsf{C}$; in particular $X$ is
contained neither in $D_2$ nor in the core $\mathsf{C}$. These are the facts asserted in (a).

(b) Fix $n$ and a tuple $\vec{\mathsf{y}}$, and let $u=(\mathbb{C}^{(\nu)}_m(X),\vec{\mathsf{y}})$.
By the statement that the pieces of chain terms of the current step carry toggles only, the entries
of $\vec{\mathsf{y}}$ are toggles only; so $u$ is an interior piece, and the bound
of the fourth hypothesis applies to each class $\zeta=((n,m),\mathsf{i})$. Summing it over
$\mathsf{i}$ with the triangle inequality gives
\begin{equation*}
\Bigl|\sum_{\mathsf{i}}W^{((n,m),\mathsf{i})}_u\Bigr|
\le e^{-\kappa^{\natural}_1|\vec{\mathsf{y}}|}\sum_{\mathsf{i}}
[\![(\delta\mathbb{C}^{(\nu)}(X))^{\natural,((n,m),\mathsf{i})}]\!] .
\end{equation*}
This bound holds at every point of $\mathfrak{X}_X$ and is uniform in $\varpi$; it therefore holds
for the supremum over $\mathfrak{X}_X$ on the left of
\eqref{9.5:eq-chain-terms-1}. Multiplying by $2$ and summing over $\vec{\mathsf{y}}$, the left-hand side
of \eqref{9.5:eq-chain-terms-1} is at most
\begin{equation*}
2\sum_{\mathsf{i}}[\![(\delta\mathbb{C}^{(\nu)}(X))^{\natural,((n,m),\mathsf{i})}]\!]
\sum_{\vec{\mathsf{y}}}e^{-\kappa^{\natural}_1|\vec{\mathsf{y}}|} .
\end{equation*}
The constants $c_l$ and $\lambda_*$ entering \eqref{9.5:eq-family-mark-sums-a} satisfy
$\kappa^{\natural}_1-3c_l\ge\lambda_*$ under that condition. Writing
\begin{equation*}
e^{-\kappa^{\natural}_1|\vec{\mathsf{y}}|}
= e^{-3c_l|\vec{\mathsf{y}}|}\,e^{-(\kappa^{\natural}_1-3c_l)|\vec{\mathsf{y}}|}
\end{equation*}
and using that the first factor is at most $1$, the sum over $\vec{\mathsf{y}}$ is bounded through
four free sums taken at the rate $\kappa^{\natural}_1-3c_l\ge\lambda_*$, and these are bounded by
$e^{4(d_m(K_v)+1)}$ exactly as in the
proof of the tuple sums of the family marks (Lemma~\ref{9.5:lem-family-mark-sums}).
This is \eqref{9.5:eq-chain-terms-1}.
\end{proof}

\emph{Vertex bounds.} Let $v$ be a stored unit, of level $j_v$ and domain $X_v$, in the sense of
Definition~\ref{9.5:def-resolved-terms}, and let $\mathsf{r}_v$ be its reading mark as defined there
(see \eqref{9.5:eq-resolved-terms-a}). Write $C^{\sharp}_{TL}$, $C^{\mathfrak{K}}_{TL}$ and
$C^{\mathbb{Y}}_{TL}$ for the constants of, respectively, the corollary on reads in rooms, the lemma
on reads at sites meeting $Z$, and the one-run expansions on labels. We put
\begin{equation}\label{9.5:eq-reading-mark-sums-1}
C^{\star}_{TL}:=\max\bigl\{C^{\sharp}_{TL},\,C^{\mathfrak{K}}_{TL},\,C^{\mathbb{Y}}_{TL},\,
C_WC_*c_3\bigr\},
\qquad
\vartheta^{\sharp}_v:=\vartheta_f^{\,j_v\wedge n_{X_v}} .
\end{equation}
Here $n_{X_v}$ is the number attached to the domain $X_v$ in the corollary on reads in rooms, whose
indices it enters through $j_v\wedge n_{X_v}$.
The last member of the maximum is the product of the three constants $C_W$, $C_*$ and $c_3$ of the
terminal estimate, whose index at level $i$ is $16C_WC_*c_3\vartheta_f^{\,i}$; here $C_W$ denotes
that estimate's constant and is distinct from the constant $C_W(s)$ of the two-point witness
theorem. The number $\vartheta^{\sharp}_v$ is the index that the corollary on
reads in rooms assigns to terms for which the first condition of
\eqref{9.5:eq-dilation-disc-a} fails; it satisfies
$\vartheta^{\sharp}_v\ge\vartheta^{n_{X_v}/2}$, the right-hand side being the index used in the
lemma on reads at sites meeting $Z$.

Consider a vertex production point, that is, one at which $\mathsf{s}(\Delta)\in\{0,1\}$ for every
$\Delta$, the common datum $\xi$ being complex. There the pair
$(\mathrm{chain}^{\mathsf{B}}(\mathfrak{z}^{\mathsf{B}}),\mathrm{chain}^{\mathsf{A}}(\mathfrak{z}^{\mathsf{A}}))$
is a vertex pair over the root seed, i.e.\ it belongs to $\mathfrak{F}^V$; here
$\mathrm{chain}^{\mathsf{m}}(\mathfrak{z}^{\mathsf{m}})$ is the chain of run $\mathsf{m}$ evaluated at
its production point $\mathfrak{z}^{\mathsf{m}}$, and the chains, the root seed, the vertex pairs,
the entry parameters $b$, $t$ and the boxes are those of the corollary on vertex pairs. This follows
from the
vertex-pair property collected in \eqref{9.5:eq-distance-units-c} and from the corollary on vertex
pairs, according to which the complex lines move only the entry data (the parameters $b$, $t$) and
leave $\vec{\mathsf{s}}=\mathbf{1}$ and the boxes fixed. We now bound the reading mark at such a
point. Terms read at sites missing the core are estimated by part~(b) of the corollary on reads in
rooms; terms read at sites meeting $Z$ are estimated by part~(b) of the lemma on reads at sites
meeting $Z$;
terms over inner strips are estimated by the one-run expansions on labels. In each of the three
cases the estimate is an application of Lemma~\ref{9.5:lem-disc-contraction} to the function
$s\mapsto\hat{F}^{\mathsf{B}}(\mathsf{v}_s)$ along the term's own path of data $\mathsf{v}_s$, on
the disc
\begin{equation*}
|s|<\bigl(8C^{\star}_{TL}\vartheta^{\sharp}_v\bigr)^{-1}.
\end{equation*}
The three cases together give, at every vertex production point, the vertex bound
\begin{equation}\label{9.5:eq-reading-mark-sums-a}
|\mathsf{r}_v|\le\beta_v:=16\,C^{\star}_{TL}\,\vartheta^{\sharp}_v\,\mathfrak{N}^{mar}_v .
\end{equation}
Here $\mathfrak{N}^{mar}_v$ is defined according to the kind of $v$.
\begin{itemize}
\item For units in the first and in the fourth rows of the table of size functionals to which
Definition~\ref{9.5:def-resolved-terms} refers, and for single
units, tails and weights, $\mathfrak{N}^{mar}_v:=\mathfrak{N}(v)$. For short single units the factor
$S^{-2}$ is retained in $\mathfrak{N}(v)$, with $S$ as in the marginal-size lemma. This rests on
part~(i) of the marginal-size lemma, namely $\Psi(0)=\Psi'(0)=0$, together with the maximum
principle applied to $\Psi/\sigma^2$, with $\Psi$ and $\sigma$ as in that lemma.
\item For cores, $\mathfrak{N}^{mar}_v$ is the marginal size $\mathfrak{N}^{mar}(v)$ of the
marginal-size lemma. By part~(ii)(B) of that lemma and by bound~(B) of the resolved growth lemma it
satisfies
\begin{equation*}
\mathfrak{N}^{mar}(v)\le C^{mar}\mathsf{W}_v^{2}\bigl[E_0\nu^2\bar{\rho}_\nu^{\,2}+b_{+}+r\bigr]t^4 ,
\end{equation*}
where $C^{mar}$, $\mathsf{W}_v$, $E_0$, $\nu$, $\bar{\rho}_\nu$, $b_{+}$, $r$ and $t$ have the
meaning they have in the marginal-size lemma and in the resolved growth lemma.
\end{itemize}
We further put
\begin{equation}\label{9.5:eq-reading-mark-sums-2}
\mathfrak{N}^{1/2}_v:=\bigl(\mathfrak{N}^{mar}_v\,\mathfrak{N}(v)\bigr)^{1/2}.
\end{equation}
Except at cores, $\mathfrak{N}^{1/2}_v=\mathfrak{N}(v)$, by the first item above. At a core,
$\mathfrak{N}^{1/2}_v$ is the smooth-core entry of the size table of the marginal-size lemma, with
its exponent $\hat{\alpha}$ halved, up to the factor $(C^{mar}C^{sm})^{1/2}$, where $C^{mar}$ is the
constant of the marginal-size lemma and $C^{sm}$ the constant of the smooth-core entry of that
table. The
exponent $\hat{\alpha}/2$ belongs to a classical quantity, of which a square root may be taken
without loss, and the square root taken in the marginal-size lemma and the one taken in
the resolved growth lemma are one and the same.

\begin{lemma}[Tuple sums of the reading marks]\label{9.5:lem-reading-mark-sums}
Let $v$ be a stored unit. Assume the following:
\begin{itemize}
\item the hypotheses on the two runs $\mathsf{A}$ and $\mathsf{B}$ collected in
\eqref{9.5:eq-distance-units-c};
\item condition \eqref{9.5:eq-family-mark-sums-b};
\item the conclusion of the dilation-disc lemma, Lemma~\ref{9.5:lem-dilation-disc};
\item the vertex bound \eqref{9.5:eq-reading-mark-sums-a};
\item the vertex-pair property in \eqref{9.5:eq-distance-units-c};
\item condition \eqref{9.5:eq-family-mark-sums-a}.
\end{itemize}
Put $C^{+}_{\mathfrak{K}}:=\max\{2,C_{\mathfrak{K}},C^{+}\}$, with a constant $C^{+}$, and
$\mathsf{c}_{rd}:=8(C^{\star}_{TL}C^{+}_{\mathfrak{K}})^{1/2}$, and let
\begin{equation*}
\mathfrak{n}^{c}(v):=\mathsf{c}_{rd}\,(\vartheta^{\sharp}_v)^{1/2}\,\mathfrak{N}^{1/2}_v\times
\begin{cases}
1 & \text{if $v$ is not deep$^{+}$},\\
e^{-\frac14\delta_P\hat{D}_v} & \text{if $v$ is deep$^{+}$}.
\end{cases}
\end{equation*}
Here deep$^{+}$ and $\hat{D}_v$ are as in Lemma~\ref{9.5:lem-dilation-disc} and
Lemma~\ref{9.5:lem-family-mark-sums}, and $K_v$ and $d_m(K_v)$ are as in
Lemma~\ref{9.5:lem-family-mark-sums}.
Then
\begin{equation}\label{9.5:eq-reading-mark-sums-3}
\mathrm{tot}^{c}(v):=\sum_{\vec{\mathsf{y}}}\mathfrak{O}^{+}\bigl(\mathfrak{p}_{v,\vec{\mathsf{y}}}
[\mathsf{r}_v]\bigr)\le\mathfrak{n}^{c}(v)\,e^{4(d_m(K_v)+1)} .
\end{equation}
\end{lemma}

\begin{proof}
The two-point oscillation $\mathfrak{O}^{+}$ is the one of \eqref{9.5:eq-two-point-oscillation-b}.
Fix a value $\Phi_C$ of the inner variables of that definition and two common data $\xi,\xi'$, and
keep the auxiliary
parameters $\varpi$ fixed. Put
\begin{equation*}
\Xi(\vec{\mathsf{s}}):=\mathsf{r}_v(\xi,\varpi,\vec{\mathsf{s}})
-\mathsf{r}_v(\xi',\varpi,\vec{\mathsf{s}}).
\end{equation*}
It equals
\begin{equation*}
\Xi(\vec{\mathsf{s}})=
\bigl[\hat{F}^{\mathsf{B}}_v(\eta^{\mathsf{B}}_v(\xi))-\hat{F}^{\mathsf{B}}_v(\eta^{\mathsf{B}}_v(\xi'))\bigr]
-\bigl[\hat{F}^{\mathsf{B}}_v(\eta^{\times}_v(\xi))-\hat{F}^{\mathsf{B}}_v(\eta^{\times}_v(\xi'))\bigr],
\end{equation*}
and it is holomorphic in $\vec{\mathsf{s}}$ on the polydisc of decoupling parameters over which the
production points of Lemma~\ref{9.5:lem-dilation-disc} range.

\emph{Bound on the polydisc.} Each of the two brackets involves the one function
$\hat{F}^{\mathsf{B}}_v$ only. Each is bounded by the argument of
Lemma~\ref{9.5:lem-family-mark-sums}, applied to $\hat{F}^{\mathsf{B}}_v$ along the dilated data
$\eta^{\mathsf{B}}_{v,\lambda}$ for the first bracket and along $\eta^{\times}_{v,\lambda}$ for the
second. Both families lie in the domain of $\hat{F}^{\mathsf{B}}_v$, by
Lemma~\ref{9.5:lem-dilation-disc}: if $v$ is deep$^{+}$, for $|\lambda|\le\rho_v$ one has
$\eta^{\mathsf{B}}_{v,\lambda}\in\mathfrak{T}^{\mathsf{B}}_{F_v}[\tfrac14]$ and
$\eta^{\times}_{v,\lambda}\in\mathfrak{X}_{X_v}$, holomorphically in $\lambda$. This gives, for each
bracket, the bound $C^{+}_{\mathfrak{K}}\mathfrak{N}(v)$ if $v$ is not deep$^{+}$ and
$C^{+}_{\mathfrak{K}}e^{-\frac12\delta_P\hat{D}_v}\mathfrak{N}(v)$ if $v$ is deep$^{+}$. The first
bracket is also bounded by the one-run total bound of run $\mathsf{B}$. Adding the two brackets,
\begin{equation*}
\sup|\Xi|\le\mathfrak{B}:=2C^{+}_{\mathfrak{K}}\,\mathfrak{N}(v)\times
\begin{cases}
1 & \text{($v$ not deep$^{+}$)},\\
e^{-\frac12\delta_P\hat{D}_v} & \text{($v$ deep$^{+}$)},
\end{cases}
\end{equation*}
the supremum being taken over that polydisc.

\emph{Bound at the vertices.} Let $\sigma'$ be a subset of the toggle set. The vertex bound
\eqref{9.5:eq-reading-mark-sums-a} applies at the datum $\xi$ and at the datum $\xi'$, and the
triangle inequality gives
\begin{equation*}
|\Xi(\mathbf{1}_{\sigma'})|\le2\beta_v .
\end{equation*}

\emph{Combination.} The operator $\mathfrak{p}_{v,\vec{\mathsf{y}}}$ acts on the variables
$\vec{\mathsf{s}}$ only and is linear. Hence $\mathfrak{p}_{v,\vec{\mathsf{y}}}[\Xi]$ is the value of
$\mathfrak{p}_{v,\vec{\mathsf{y}}}[\mathsf{r}_v]$ at $\xi$ minus its value at $\xi'$. We apply part~(iii)
of Lemma~\ref{9.5:lem-polydisc-vertices} with $G=\Xi$, $R=e^{\kappa_1}$, $\beta=2\beta_v$ and the
bound $\mathfrak{B}$ above. This gives
\begin{equation*}
|\mathfrak{p}_{v,\vec{\mathsf{y}}}[\Xi]|\le\Bigl(\frac{2}{e^{\kappa_1}-1}\Bigr)^{|\vec{\mathsf{y}}|/2}
(2\beta_v\mathfrak{B})^{1/2}.
\end{equation*}
Since $(2/(e^{\kappa_1}-1))^{1/2}\le e^{-\kappa^{\natural}_1}$ by condition
\eqref{9.5:eq-family-mark-sums-a}, which fixes the pair decoupling rate $\kappa^{\natural}_1$, we
obtain
\begin{equation*}
|\mathfrak{p}_{v,\vec{\mathsf{y}}}[\Xi]|\le e^{-\kappa^{\natural}_1|\vec{\mathsf{y}}|}
(2\beta_v\mathfrak{B})^{1/2}.
\end{equation*}
By the definitions of $\beta_v$ and $\mathfrak{B}$,
\begin{align*}
2\beta_v\mathfrak{B}
&=2\cdot16C^{\star}_{TL}\vartheta^{\sharp}_v\mathfrak{N}^{mar}_v\cdot
2C^{+}_{\mathfrak{K}}\mathfrak{N}(v)\times\{1;\,e^{-\frac12\delta_P\hat{D}_v}\}\\
&=64\,C^{\star}_{TL}C^{+}_{\mathfrak{K}}\,\vartheta^{\sharp}_v\,\mathfrak{N}^{mar}_v\,
\mathfrak{N}(v)\times\{1;\,e^{-\frac12\delta_P\hat{D}_v}\},
\end{align*}
where the first entry of the brace is taken if $v$ is not deep$^{+}$ and the second if $v$ is
deep$^{+}$. Taking the square root and using \eqref{9.5:eq-reading-mark-sums-2},
\begin{align*}
(2\beta_v\mathfrak{B})^{1/2}
&=8\bigl(C^{\star}_{TL}C^{+}_{\mathfrak{K}}\bigr)^{1/2}(\vartheta^{\sharp}_v)^{1/2}
\bigl(\mathfrak{N}^{mar}_v\mathfrak{N}(v)\bigr)^{1/2}\times\{1;\,e^{-\frac14\delta_P\hat{D}_v}\}\\
&=\mathfrak{n}^{c}(v).
\end{align*}
Taking, at the fixed value $\Phi_C$, the supremum over the common data $\xi,\xi'$ in the definition
\eqref{9.5:eq-two-point-oscillation-b} of $\mathfrak{O}^{+}$, we obtain for every tuple
\begin{equation*}
\mathfrak{O}^{+}\bigl(\mathfrak{p}_{v,\vec{\mathsf{y}}}[\mathsf{r}_v]\bigr)\le
e^{-\kappa^{\natural}_1|\vec{\mathsf{y}}|}\,\mathfrak{n}^{c}(v).
\end{equation*}
It remains to sum over the tuples $\vec{\mathsf{y}}$. The sum is performed as in the proof of
Lemma~\ref{9.5:lem-family-mark-sums}: it is carried out as four free sums at
the rate $\kappa^{\natural}_1-3c_l$, which satisfies $\kappa^{\natural}_1-3c_l\ge\lambda_*$, with
$c_l$ and $\lambda_*$ as there. These
sums produce the factor $e^{4(d_m(K_v)+1)}$, and \eqref{9.5:eq-reading-mark-sums-3} follows.
\end{proof}

The estimate takes one square root, at the production point, jointly in all toggle parameters and
in the depth. Under it stand only sizes read in one run and the classical index
$\vartheta^{\sharp}_v$; no unknown quantity of the comparison between the two runs enters.

\emph{Purely classical constituents.} Three further kinds of constituent are handled by the same
argument: the units $v_e$ among the Gaussians $\mathfrak{Q}(\mathfrak{a})$ of an interrupted run,
the kept quadratic form $\mathfrak{q}_{\mathsf{C}}$ at the core $\mathsf{C}$, and the classical part
of the localised members $\mathfrak{v}^{\mathbf{C}}$ of the remainder terms. For each of them, the
whole difference between its values in the two runs takes the place of the reading mark
$\mathsf{r}_v$ in the argument above.

Two inputs are needed, supplied by two different statements:
\begin{itemize}
\item the pointwise smallness of that difference is supplied by the pointwise smallness statement
for the classical terms;
\item its decay is supplied by the bound that each run provides for its own terms of the
corresponding kind, namely the one-run bounds on the Gaussians of the interrupted run, the
bound on the remainder terms, and the bound on the kept quadratic forms.
\end{itemize}
The estimate is carried out in the lemma on the purely classical constituents.

Throughout, $m$ is a step, $C$ is a component of $Z$ at $\mathbf{R}_m$ for an aligned label, and $n\le m$ is a layer. The stored units of the lists (m1)--(m3), the derived units $\mathbb{C}^{(\nu)}_m(X)$, the copies $\mathfrak{L}^{cp}(C)$, the rim set $\mathcal{Y}_{rim}$ and the members $\delta\mathfrak{v}^{\mathbf{C},(n)}(Y_4)$, $\delta\mathfrak{q}^{(n)}_C$ are those of Definition~\ref{9.5:def-layer-store}. The plain units are $\mathfrak{L}^{pl}=\mathfrak{L}^{B}\sqcup\mathfrak{L}^{pt}$ (list (m1)) and the wrapped units are $\mathfrak{L}^{wr}$ (list (m2)).

\emph{Weights.} For a stored unit $v$ of (m1) or (m2) we put
\begin{equation*}
W^{(n)}(v):=\bigl\{\mathfrak{N}_v[G^{(n)}_v]
+\mathbf{1}[n^c_v=n]\,\mathsf{c}_{rd}(\vartheta^{\sharp}_v)^{1/2}\mathfrak{N}^{1/2}_v\bigr\}
e^{4d_m(K_v)} .
\end{equation*}
For a derived unit $v=\mathbb{C}^{(\nu)}_m(X)$ we put
\begin{equation*}
W^{(n)}(v):=\sum_{\mathsf{i}}
[\![(\delta\mathbb{C}^{(\nu)}(X))^{\natural,((n,m),\mathsf{i})}]\!]\,e^{4d_m(K_v)} .
\end{equation*}

\emph{The seven lines.} Let $\mathsf{K}^{\flat,1/2}$ be the constant $\mathsf{K}^{\flat}$ of the per-cube bound for pointed units of the one-run analysis (the constant of line $(\ell3)$ of the homogeneity lemma for the counting bounds), computed with $\hat{\alpha}$ replaced by $\hat{\alpha}/2$ in the series of the smooth cores; here $\hat{\alpha}$ is the radius exponent entering that series and $\bar{\rho}_b$, $b\ge0$, are the quantities of the smooth cores that enter its terms. It is finite, because
\begin{equation*}
\sum_b(b+1)^3\bar{\rho}_b^{\,2}(1+\bar{\rho}_b)L^{-\hat{\alpha}b/2}<\infty .
\end{equation*}
For a copy $v\in\mathfrak{L}^{cp}(C)$ we write $\mathring{\mathsf{g}}^{(n)}_v$ for $\mathsf{g}^{(n)}_v+\mathbf{1}[n^c_v=n]\mathsf{r}_v$ minus its value at the trivial datum.
We let $\mathsf{d}^{(n)}_m$ be the least number $\psi\ge0$ such that, for every aligned label and every component $C$ of $Z$ at $\mathbf{R}_m$, the following seven inequalities hold:
\begin{equation}\label{9.5:eq-production-bound-a}
\begin{aligned}
&(\ell1)^{(n)}\quad \forall j\le m,\ \forall\square'\in\pi_j:\\
&\qquad\sum_{F\in\mathfrak{L}^{wr}(C):\,j_F=j,\,A_F\ni\square'}
W^{(n)}(F)\,e^{\varsigma_0 d_j(A_F)\mathbf{1}[j<m]}
\le\psi\,\mathsf{b}_m(m-j);\\
&(\ell2)^{(n)}\quad \forall j\le m,\ \forall\square'\in\pi_j:\\
&\qquad\sum_{F\in\mathfrak{L}^{B}(C):\,j_F=j,\,X_F\ni\square'}
W^{(n)}(F)\,e^{\varsigma_0 d_j(X_F)\mathbf{1}[j<m]}
\le\psi\,\mathsf{b}(m-j);\\
&(\ell3)^{(n)}\quad \forall\text{ cells }\Delta'\subset C:\\
&\qquad\sum_{v\in\mathfrak{L}^{pt}(C):\,y_v\in\Delta'}W^{(n)}(v)\,e^{\varsigma_1\mathsf{r}(v)}
\le\psi\,\mathsf{K}^{\flat,1/2}\,\mathsf{w}^{\vee}(\Delta')^3;\\
&(\ell4)^{(n)}\quad \forall\,\pi_m\text{-cubes }\Delta'\subset C:\\
&\qquad\sum_{v\in\mathfrak{L}^{cp}(C):\,y_v\in\Delta'}
\sup\bigl|\mathring{\mathsf{g}}^{(n)}_v\bigr|\le\psi\,\mathsf{K}^{cp};\\
&(\ell5)^{(n)}\quad \forall\,Y_4\in\mathcal{Y}_{rim}:\\
&\qquad\sup|W^{(n)}(Y_4)|\le\psi\,[3+\mathsf{K}^{0}+N(g_m)C^{\sharp}_0]\,
e^{-a\,d_{m,Z}(Y_4)},\quad a=\tfrac{31}{20}\kappa;\\
&(\ell6)^{(n)}\quad \sum_{Y_4\subset C^{\sharp}}\sup|\delta\mathfrak{v}^{\mathbf{C},(n)}(Y_4)|
\le\psi\,\mathsf{c}_{\mathbf{C}};\\
&(\ell7)^{(n)}\quad \sup|\delta\mathfrak{q}^{(n)}_C|\le\psi\,\mathfrak{q}_{\max}(m).
\end{aligned}
\end{equation}
Here $W^{(n)}(Y_4)$ in $(\ell5)^{(n)}$ is the layer-$n$ mark of the rim piece $Y_4$ of Definition~\ref{9.5:def-layer-store}. The remaining symbols of these lines are those of the one-run analysis of the kept potential: $\pi_j$ is the partition of level $j$ into cubes $\square'$; for a unit $F$, $j_F$ is its level and $A_F$, $X_F$ are the sets attached to it in the counting bounds; for a unit $v$ of $\mathfrak{L}^{pt}(C)$ or $\mathfrak{L}^{cp}(C)$, $y_v$ is its anchor point; $\mathsf{r}(v)$ is the quantity attached to a pointed unit in the per-cube bound for pointed units, distinct from the reading mark $\mathsf{r}_v$; $\mathsf{b}_m(\cdot)$ and $\mathsf{b}(\cdot)$ are the per-cube budget functions of the counting bounds for wrapped and for plain units; $N(g_m)$ and $C^{\sharp}_0$ are the numbers entering the rim bound, $N(g_m)$ being determined by the coupling $g_m$ and distinct from the rank $N$ of $\mathrm{SU}(N)$; $\mathsf{c}_{\mathbf{C}}$ is the constant of the bound for the remainder terms $\mathbf{C}_m$; and $\mathfrak{q}_{\max}(m)$ is the bound for the kept form at step $m$.

Lines $(\ell1)$--$(\ell3)$ are the hypotheses of the homogeneity lemma for the counting bounds of plain and wrapped units. Lines $(\ell4)$--$(\ell7)$ are the accompanying inequalities for the copies, the rim, the remainder terms $\mathbf{C}_m$ and the form. All seven are written here for the layer-$n$ weights.

\emph{Constants.} Let $\mathsf{A}(k):=\Sigma^{V}(k)-\mathrm{nest}^{*}(k)+2\mathfrak{q}_{\max}(k)$ be the production budget of the one-run analysis of the kept potential, in which $\Sigma^{V}(k)$ and $\mathrm{nest}^{*}(k)$ are the summed bound of the members of the kept potential and its nested part, and whose constants include $C^{+}$, $C_{3F}$ and $\mathsf{K}^{\flat}$. We put
\begin{equation*}
\mathsf{A}^{\natural}(m):=\Sigma^{V}(m)-\mathrm{nest}^{*}(m)+2\mathfrak{q}_{\max}(m),
\end{equation*}
computed with the constants $(C^{+},C_{3F},\mathsf{K}^{\flat})$ replaced by $(C^{+}_{\mathfrak{K}},C^{+}_{\mathfrak{K}},\mathsf{K}^{\flat,1/2})$.
Let $C^{\natural}_{\Sigma}$ be twice the sum of the eight coefficients of the table of that budget, changed in the same way. Then $C^{\natural}_{\Sigma}\ge C_{\Sigma}$. By item (i) of the order-of-choice statement of the one-run analysis, $C^{\natural}_{\Sigma}$ is a number depending on the same parameters as $C_{\Sigma}$ and on $r_{\mathfrak{K}}$ and $\hat{\alpha}$. The degrees in $\check{R}_m$ of the table do not change, so that
\begin{equation}\label{9.5:eq-production-bound-2}
\mathsf{A}^{\natural}(m)\le\tfrac12 C^{\natural}_{\Sigma}\check{R}_m^{p_{\Sigma}} .
\end{equation}

\emph{The nesting condition.} Put
\begin{equation*}
\varepsilon_\gamma:=\sup_{\sigma\ge\frac12R_\gamma}(2\sigma)^{p_{\Sigma}}e^{-2(\sigma-1)} .
\end{equation*}
The condition
\begin{equation}\label{9.5:eq-production-bound-b}
C_{\mathfrak{K}}\,e^{4}\cdot C_{cell}\cdot 2d\,(104Lc_R+2)^{d-1}\cdot\varepsilon_\gamma\le1
\end{equation}
is a ceiling on $\gamma$ of degree $-\infty$, in supremum form. It is the nesting condition of the one-run analysis with the constant $C_{3F}$ replaced by $C_{\mathfrak{K}}$. Here $C_{cell}$ is the constant of item (b) of the cell-counting lemma, $c_R$ is a constant of the geometry of the one-run analysis, and $R_\gamma$ is the radius threshold attached to $\gamma$.

\emph{Kept quotients.} We put
\begin{equation}\label{9.5:eq-production-bound-c}
\begin{gathered}
x_n(m):=\sup\frac{\mathfrak{O}^{+}(\delta\hat{V}^{(n)}(Y^{\sharp}))
+\mathfrak{O}^{+}(\delta\hat{\mathfrak{q}}^{(n)}_Y)}{\check{R}_m^{p_{\Sigma}}},\\
\mathsf{k}^{(<m)}_n:=\sup_{m'<m}x_n(m'),\qquad
\mathsf{k}^{(\le k')}_n:=\sup_{m'\le k'}x_n(m'),\qquad \mathsf{k}_n:=\sup_m x_n(m).
\end{gathered}
\end{equation}
The supremum defining $x_n(m)$ runs over aligned labels and over thrown-back regions $(Y,m)$ created at $\mathbf{R}_m$.

\begin{theorem}[The production bound, layer by layer]\label{9.5:thm-production-bound}
Assume the following:
\begin{enumerate}
\item[(a)] the standing conditions \textup{(F0)}--\textup{(F5)} on the common data, the charts and the two-point oscillation (Notation~\ref{9.5:not-two-point-oscillation});
\item[(b)] the condition \eqref{9.5:eq-distance-units-c} on the two units of distance, for both runs $\mathsf{m}=\mathsf{A}$ and $\mathsf{m}=\mathsf{B}$;
\item[(c)] the condition \eqref{9.5:eq-family-mark-sums-b};
\item[(d)] the chart condition \eqref{9.5:eq-two-point-oscillation-c};
\item[(e)] the homogeneity lemma for the counting bounds: for every weight $W\ge0$ and every $\psi\ge0$ such that lines $(\ell1)$--$(\ell3)$ of \eqref{9.5:eq-production-bound-a} hold with $W$ in place of $W^{(n)}$, the four sums of the counting bounds for plain and wrapped units of the one-run analysis are at most $\psi$ times their right members, computed with the constant of $(\ell3)$;
\item[(f)] the pair decoupling bound \eqref{9.5:eq-family-mark-sums-a} at rate $\kappa^{\natural}_1$;
\item[(g)] the nesting condition \eqref{9.5:eq-production-bound-b};
\item[(h)] from the one-run analysis of the kept potential:
\begin{itemize}
\item the nesting statement: to every thrown-back region $(Y',m')$ with $Y'\subset C$, created at a step $m'$ not later than the healing step of $C$, which precedes $m$, it assigns a cube $\Delta_{Y'}$; the map $Y'\mapsto\Delta_{Y'}$ is injective, the radius of $\Delta_{Y'}$ is at least $\frac12\check{R}_{m'}+1\ge\frac12R_\gamma$, and
\begin{equation*}
e^{-\frac18\delta_P\,d^{\wedge}((Y'^{\sharp})^{+},\mathsf{C}^c)}\,\check{R}_{m'}^{p_{\Sigma}}
\le\varepsilon_\gamma ,
\end{equation*}
where $\mathsf{C}$ is the core of $C$;
\item item (b) of the cell-counting lemma;
\item item (ii) of the copy bound;
\item the rim bound;
\item the production budget together with its table of coefficients.
\end{itemize}
\end{enumerate}
Then, for every step $m$, every aligned label, every component $C$ of $Z$ at $\mathbf{R}_m$ and every layer $n\le m$,
\begin{equation}\label{9.5:eq-production-bound-1}
\begin{split}
\mathfrak{O}^{+}(\delta\hat{V}^{(n)}(C^{\sharp}))+\mathfrak{O}^{+}(\delta\hat{\mathfrak{q}}^{(n)}_C)
&\le 2\mathsf{d}^{(n)}_m\mathsf{A}^{\natural}(m)
+\bigl[\tfrac12\mathsf{k}^{(<m)}_n+\tfrac12\mathsf{c}'_{rd}C_{\Sigma}\vartheta_f^{n/2}\bigr]
\check{R}_m^{d-1}\\
&\le\bigl[C^{\natural}_{\Sigma}\mathsf{d}^{(n)}_m+\tfrac12\mathsf{k}^{(<m)}_n
+\tfrac12\mathsf{c}'_{rd}C_{\Sigma}\vartheta_f^{n/2}\bigr]\check{R}_m^{p_{\Sigma}},
\end{split}
\end{equation}
where $\mathsf{c}'_{rd}:=2\mathsf{c}_{rd}/C_{\mathfrak{K}}$. No constant in the right members depends on any of the following:
\begin{itemize}
\item the count $\sum_j|\Gamma^0_j\cap C|$ and the parameter $D$ of the one-run analysis of the kept potential, or the volume $|\mathsf{C}|$ of the core $\mathsf{C}$ of $C$;
\item an age;
\item $K$, $\mathfrak{n}$, $m$ or the parameter $N_0$ of the one-run analysis of the kept potential;
\item the number, the sizes or the pasts of thrown-back regions and of arrests.
\end{itemize}
\end{theorem}

\begin{proof}
\emph{Reduction to a sum of members.} The two-point oscillation $\mathfrak{O}^{+}$ is a seminorm. By the chart condition \eqref{9.5:eq-two-point-oscillation-c} it is computed at the trivial chart over regauged seed pairs. We apply Definition~\ref{9.5:def-layer-store} and the subadditivity of $\mathfrak{O}^{+}$, and we bound the oscillation of a member by twice its supremum wherever no tuple sum is available. Together these give
\begin{align*}
\mathfrak{O}^{+}(\delta\hat{V}^{(n)}(C^{\sharp}))
&\le\sum_{v\ \text{stored}}\bigl[\mathrm{tot}^{u,(n)}(v)+\mathbf{1}[n^c_v=n]\,\mathrm{tot}^{c}(v)\bigr]
+\sum_{u\ \text{derived}}\sum_{\vec{\mathsf{y}}}
2\sup\Bigl|\sum_{\mathsf{i}}W^{((n,m),\mathsf{i})}_u\Bigr|\\
&\quad+\sum_{v\in\mathfrak{L}^{cp}(C)}2\sup|\cdot|
+\sum_{Y_4\in\mathcal{Y}_{rim}}2\sup|W^{(n)}(Y_4)|
+2\sum_{Y_4\subset C^{\sharp}}\sup|\delta\mathfrak{v}^{\mathbf{C},(n)}(Y_4)| .
\end{align*}
In this display:
\begin{itemize}
\item the stored units run through the lists (m1)--(m3);
\item $\mathrm{tot}^{u,(n)}(v)$ is the tuple sum of Lemma~\ref{9.5:lem-family-mark-sums};
\item $\mathrm{tot}^{c}(v)$ is the tuple sum of Lemma~\ref{9.5:lem-reading-mark-sums};
\item the derived term is the left member of part (b) of Lemma~\ref{9.5:lem-chain-terms};
\item the supremum in the copy term is the one estimated in line $(\ell4)^{(n)}$ of \eqref{9.5:eq-production-bound-a}.
\end{itemize}
In the same way $\mathfrak{O}^{+}(\delta\hat{\mathfrak{q}}^{(n)}_C)\le2\sup|\delta\mathfrak{q}^{(n)}_C|$. We bound the terms group by group.

\emph{Plain and wrapped units.} Let $v$ be a stored unit of (m1) or (m2).

Lemma~\ref{9.5:lem-family-mark-sums} gives $\mathrm{tot}^{u,(n)}(v)\le\mathfrak{n}^{u,(n)}(v)e^{4(d_m(K_v)+1)}$. Here $\mathfrak{n}^{u,(n)}(v)$ equals $2\mathfrak{N}_v[G^{(n)}_v]$ if $v$ is not deep$^{+}$, and equals $C_{\mathfrak{K}}e^{-\frac12\delta_P\hat{D}_v}\mathfrak{N}_v[G^{(n)}_v]$ if $v$ is deep$^{+}$.

Lemma~\ref{9.5:lem-reading-mark-sums} gives $\mathrm{tot}^{c}(v)\le\mathfrak{n}^{c}(v)e^{4(d_m(K_v)+1)}$, with
\begin{equation*}
\mathfrak{n}^{c}(v)=\mathsf{c}_{rd}(\vartheta^{\sharp}_v)^{1/2}\mathfrak{N}^{1/2}_v
\end{equation*}
if $v$ is not deep$^{+}$, and the same quantity multiplied by $e^{-\frac14\delta_P\hat{D}_v}$ if $v$ is deep$^{+}$.

The constant $\mathsf{c}_{rd}$ appears inside $W^{(n)}(v)$. If $v$ is not deep$^{+}$, then, since twice a nonnegative number plus a nonnegative number is at most twice their sum, the contribution of $v$ is at most $2e^{4}W^{(n)}(v)$. If $v$ is deep$^{+}$, we use two facts:
\begin{itemize}
\item $C^{+}_{\mathfrak{K}}=\max\{2,C_{\mathfrak{K}},C^{+}\}\ge\max\{1,C_{\mathfrak{K}}\}$;
\item $e^{-\frac12\delta_P\hat{D}_v}\le e^{-\frac14\delta_P\hat{D}_v}$.
\end{itemize}
Together they give that the contribution of $v$ is at most $e^{4}C^{+}_{\mathfrak{K}}W^{(n)}(v)e^{-\frac14\delta_P\hat{D}_v}$.

For a derived unit $u=\mathbb{C}^{(\nu)}_m(X)$, part (a) of Lemma~\ref{9.5:lem-chain-terms} shows that $u$ is never deep. Part (b), under \eqref{9.5:eq-family-mark-sums-a}, bounds its contribution by
\begin{equation*}
2\sum_{\mathsf{i}}[\![(\delta\mathbb{C}^{(\nu)}(X))^{\natural,((n,m),\mathsf{i})}]\!]\,
e^{4(d_m(K_u)+1)}=2e^{4}W^{(n)}(u).
\end{equation*}

We now apply the homogeneity lemma for the counting bounds, hypothesis (e), with $W:=W^{(n)}$ and $\psi:=\mathsf{d}^{(n)}_m$. Its hypotheses are lines $(\ell1)^{(n)}$--$(\ell3)^{(n)}$ of \eqref{9.5:eq-production-bound-a}, and these hold with $\psi=\mathsf{d}^{(n)}_m$ by the definition of $\mathsf{d}^{(n)}_m$.

Hypothesis (e) is used in exactly the generality in which it is stated. We record that the proof of the homogeneity lemma depends on the weight only through the following:
\begin{itemize}
\item lines $(\ell1)$--$(\ell3)$, the condition on the budget functions $\mathsf{b}_m(\cdot)$, $\mathsf{b}(\cdot)$, and the geometry;
\item the bound for deep units, which involves no field variable; there the weight enters only through the anchored sums, whose setting is an arbitrary weight $W(F)\ge0$;
\item the shell bound for wrapped units, taken at the value $4$ of its parameter, together with the counting of units;
\item the pile lemma: its estimate for plain units depends on the weight only through $(\ell2)$ and the spare rate $\varsigma_0$, and its estimates for pointed units only through $(\ell3)$, linearly in the constant of $(\ell3)$;
\item the two summations that follow the pile lemma, which use the pile lemma, the distance bounds with their window, item (b) of the cell-counting lemma and the zone counts of the one-run geometry.
\end{itemize}
The anchored sums require the budget functions $\mathsf{b}_m(\cdot)$, $\mathsf{b}(\cdot)$ to be at least $1$; in the proof of the homogeneity lemma this is met by passing, by homogeneity, to the weight $W/\psi$. Applied here, the lemma gives that the four sums of the counting bounds, formed with the contributions just bounded, are at most $\mathsf{d}^{(n)}_m$ times the right members of those bounds, with $\mathsf{K}^{\flat,1/2}$ in place of $\mathsf{K}^{\flat}$; these right members are the members of $\mathsf{A}^{\natural}(m)$ corresponding to the plain and wrapped units.

\emph{Copies.} Every $y_v$ with $v\in\mathfrak{L}^{cp}(C)$ lies in a $\pi_m$-cube $\Delta'\subset C$. Let $N_m(C)$ denote the number of $\pi_m$-cubes of $C$. Summing line $(\ell4)^{(n)}$ over these cubes gives
\begin{equation*}
\sum_{v\in\mathfrak{L}^{cp}(C)}2\sup|\cdot|\le2\,\mathsf{d}^{(n)}_m\mathsf{K}^{cp}N_m(C),
\qquad N_m(C)\le(100\check{R}_m)^d,
\end{equation*}
where the bound on $N_m(C)$ is item (ii) of the copy bound.

\emph{Rim.} Line $(\ell5)^{(n)}$ gives
\begin{equation*}
\sum_{Y_4\in\mathcal{Y}_{rim}}2\sup|W^{(n)}(Y_4)|
\le2\,\mathsf{d}^{(n)}_m[3+\mathsf{K}^{0}+N(g_m)C^{\sharp}_0]
\sum_{Y_4\in\mathcal{Y}_{rim}}e^{-a\,d_{m,Z}(Y_4)} .
\end{equation*}
The last sum is bounded as in the rim bound of the one-run analysis. That bound uses item (i) of the splitting of the kept potential at the rim, the summation bound for the weights $e^{-a\,d_{m,Z}}$, and the count $\mathsf{n}_\sharp(m)$ of rim cubes.

\emph{The remainder terms and the form.} Lines $(\ell6)^{(n)}$ and $(\ell7)^{(n)}$ give
\begin{equation*}
2\sum_{Y_4\subset C^{\sharp}}\sup|\delta\mathfrak{v}^{\mathbf{C},(n)}(Y_4)|
\le2\,\mathsf{d}^{(n)}_m\mathsf{c}_{\mathbf{C}},
\qquad
\mathfrak{O}^{+}(\delta\hat{\mathfrak{q}}^{(n)}_C)\le2\,\mathsf{d}^{(n)}_m\mathfrak{q}_{\max}(m).
\end{equation*}

\emph{Comparison with $\mathsf{A}^{\natural}(m)$.} Each member bounded so far is at most $\mathsf{d}^{(n)}_m$ times the corresponding member of $\mathsf{A}^{\natural}(m)$. The rim, copy, remainder and form members of the one-run budget already contain the factor $2$ coming from the bound of an oscillation by twice a supremum. We state the bound with a spare factor $2$, which covers the units that are deep but not deep$^{+}$. The members of the lists (m1), (m2), the derived units, the copies, the rim, the remainder terms and the form together contribute at most $2\mathsf{d}^{(n)}_m\mathsf{A}^{\natural}(m)$.

\emph{Inner thrown-back regions (the list (m3)).} Let $(Y',m')$ be a thrown-back region with $Y'\subset C$, created at a step $m'$ not later than the healing step of $C$, which precedes $m$, and let $v'\in\mathfrak{K}(Y')$. Then $X_{v'}=Y'^{\sharp}\subset\mathsf{C}$, and $K_{v'}$ is one cube. Put
\begin{equation*}
D':=d^{\wedge}((Y'^{\sharp})^{+},\mathsf{C}^c).
\end{equation*}

\emph{Family marks.} We apply Lemma~\ref{9.5:lem-family-mark-sums} in the case of the kept factors of an inner thrown-back region, in its deep$^{+}$ case, together with the following:
\begin{itemize}
\item the size functional $\mathfrak{N}_{v'}[G^{(n)}_{v'}]$ is given by the two-point oscillation divided by $\rho_{v'}$, by the statement expressing the size functional of a kept factor through its two-point oscillation, with $\rho_{v'}$ and $C_{\mathfrak{K}}$ as in \eqref{9.5:eq-dilation-disc-a};
\item the constant is $\tfrac12C_{\mathfrak{K}}$: it is the constant $\tfrac12C_{3F}$ of item (iv) of the bound for deep kept factors, taken with $C_{\mathfrak{K}}$ in place of $C_{3F}$;
\item the factor $e^{4(d_m(K_{v'})+1)}$ is $e^{4}$, since $K_{v'}$ is one cube.
\end{itemize}
This gives
\begin{align*}
\sum_{v'\in\mathfrak{K}(Y')}\mathrm{tot}^{u,(n)}(v')
&\le\tfrac12C_{\mathfrak{K}}e^{4}e^{-\frac12\delta_PD'}
\bigl[\mathfrak{O}^{+}(\delta\hat{V}^{(n)}(Y'^{\sharp}))+\mathfrak{O}^{+}(\delta\hat{\mathfrak{q}}^{(n)}_{Y'})\bigr]\\
&\le\tfrac12C_{\mathfrak{K}}e^{4}e^{-\frac12\delta_PD'}\,x_n(m')\,\check{R}_{m'}^{p_{\Sigma}} .
\end{align*}
The second inequality is the definition \eqref{9.5:eq-production-bound-c} of $x_n(m')$, since $(Y',m')$ is created at $\mathbf{R}_{m'}$.

By the nesting statement of hypothesis (h), the map $Y'\mapsto\Delta_{Y'}$ is injective, the radius of $\Delta_{Y'}$ is at least $\tfrac12\check{R}_{m'}+1\ge\tfrac12R_\gamma$, and
\begin{equation*}
e^{-\frac18\delta_PD'}\check{R}_{m'}^{p_{\Sigma}}\le\varepsilon_\gamma .
\end{equation*}
Since $D'\ge0$,
\begin{equation*}
e^{-\frac12\delta_PD'}\le e^{-\frac18\delta_PD'}e^{-\frac14\delta_PD'} .
\end{equation*}
Since $m'<m$, we have $x_n(m')\le\mathsf{k}^{(<m)}_n$. We sum $e^{-\frac14\delta_PD'}$ over the inner thrown-back regions by item (b) of the cell-counting lemma at rate $\frac14\delta_P$. In all, the family marks of the inner thrown-back regions contribute at most
\begin{equation*}
\tfrac12C_{\mathfrak{K}}e^{4}\varepsilon_\gamma C_{cell}\,\mathsf{n}_\partial(m)\,\mathsf{k}^{(<m)}_n
\le\tfrac12\check{R}_m^{d-1}\mathsf{k}^{(<m)}_n ,
\end{equation*}
where $\mathsf{n}_\partial(m)$ is the boundary count at step $m$ entering item (b) of the cell-counting lemma. The last inequality follows from $\mathsf{n}_\partial(m)\le2d\,(104Lc_R+2)^{d-1}\check{R}_m^{d-1}$ and the nesting condition \eqref{9.5:eq-production-bound-b}.

\emph{Reading marks.} By the rule by which reading marks are filed, $n^c_{v'}=m'$ for $v'\in\mathfrak{K}(Y')$: the reading marks of the kept factors of $(Y',m')$ lie in layer $m'$. Hence the indicator $\mathbf{1}[n^c_{v'}=n]$ retains only the inner regions with $m'=n$.

We apply Lemma~\ref{9.5:lem-reading-mark-sums} with the size functional taken to be $\mathfrak{O}^{+}(\hat{v}'^{\mathsf{B}})$, together with the bound
\begin{equation*}
\sum_{v'\in\mathfrak{K}(Y')}\mathfrak{O}^{+}(\hat{v}'^{\mathsf{B}})\le C_{\Sigma}\check{R}_{m'}^{p_{\Sigma}} ,
\end{equation*}
which the one-run analysis of the kept potential supplies for the kept factors of the region $(Y',m')$. This gives
\begin{equation*}
\sum_{v'\in\mathfrak{K}(Y')}\mathrm{tot}^{c}(v')
\le\mathsf{c}_{rd}\,\vartheta_f^{m'/2}e^{4}e^{-\frac14\delta_PD'}C_{\Sigma}\check{R}_{m'}^{p_{\Sigma}} .
\end{equation*}
We write $\frac14=\frac18+\frac18$. One factor $e^{-\frac18\delta_PD'}$ is absorbed, together with $\check{R}_{m'}^{p_{\Sigma}}$, into $\varepsilon_\gamma$ as above. The other is summed by item (b) of the cell-counting lemma at rate $\frac18\delta_P$, which is at least $\delta_P/16$, the lower bound on the rate under which that item is applied here. Since $m'=n$, the reading marks of the inner thrown-back regions contribute at most
\begin{equation*}
\mathsf{c}_{rd}\,\vartheta_f^{n/2}C_{\Sigma}e^{4}\varepsilon_\gamma C_{cell}\,\mathsf{n}_\partial(m)
\le\frac{\mathsf{c}_{rd}}{C_{\mathfrak{K}}}C_{\Sigma}\vartheta_f^{n/2}\check{R}_m^{d-1}
=\tfrac12\mathsf{c}'_{rd}C_{\Sigma}\vartheta_f^{n/2}\check{R}_m^{d-1} .
\end{equation*}
The inequality follows from the bound on $\mathsf{n}_\partial(m)$ and from \eqref{9.5:eq-production-bound-b}, which gives $e^{4}\varepsilon_\gamma C_{cell}\,2d\,(104Lc_R+2)^{d-1}\le1/C_{\mathfrak{K}}$. The equality is the definition $\mathsf{c}'_{rd}=2\mathsf{c}_{rd}/C_{\mathfrak{K}}$.

\emph{Conclusion.} We add three contributions:
\begin{itemize}
\item $2\mathsf{d}^{(n)}_m\mathsf{A}^{\natural}(m)$ from all members other than those of (m3);
\item $\frac12\check{R}_m^{d-1}\mathsf{k}^{(<m)}_n$ from the family marks of the inner thrown-back regions;
\item $\frac12\mathsf{c}'_{rd}C_{\Sigma}\vartheta_f^{n/2}\check{R}_m^{d-1}$ from their reading marks.
\end{itemize}
The sum is the first inequality of \eqref{9.5:eq-production-bound-1}. The second inequality follows from \eqref{9.5:eq-production-bound-2}, which gives $2\mathsf{d}^{(n)}_m\mathsf{A}^{\natural}(m)\le C^{\natural}_{\Sigma}\mathsf{d}^{(n)}_m\check{R}_m^{p_{\Sigma}}$, together with $\check{R}_m^{d-1}\le\check{R}_m^{p_{\Sigma}}$, which holds because $\check{R}_m\ge1$ and $p_{\Sigma}\ge d-1$ for the width $\check{R}_m$ and the exponent $p_{\Sigma}$ of the one-run analysis of the kept potential.

No constant used above involves $\sum_j|\Gamma^0_j\cap C|$, $|\mathsf{C}|$, $D$, an age, $K$, $\mathfrak{n}$, $m$, $N_0$, or the number, the sizes or the pasts of thrown-back regions and arrests.
\end{proof}

\begin{definition}[Rates and unit shares for the constituent bounds]\label{9.5:not-constituent-rates}
Fix a step $m$ and a zone $n\le m$. We write
\begin{equation*}
[\![\delta]\!]_n:=[\![\delta]\!]^{\natural+,(m)}_n
\end{equation*}
for the class load of layer $n$ read at level $m$, in the sense of the definition of the class
loads. It is formed with the boundary terms produced before the comparison's step at level $m$:
the boundary terms $\mathbf{B}$ born at the operations $\mathbf{T}$ of levels at most $m$, and the
boundary terms $\mathbf{B}'$ born at the operations $\mathbf{R}$ of levels below $m$. Its classical
part is $[\![\delta]\!]^{c}_n$; we write $E_0$ for the constant appearing in its entries and
$C^{\natural}_{cl}$ for the constant of its logarithmic entry. We put
$K^{fr}_{\mathbb{C}}:=\tfrac98K_{\mathbb{C}}$. The smallness index of layer $n$ is
\begin{equation*}
\hat{\vartheta}_n=\min\bigl\{2,2(16C_WC_*c_3\vartheta_f^n)^{1/2}\bigr\},\qquad
\hat{\vartheta}^{\natural}_n=\hat{\vartheta}_n^{1/2},
\end{equation*}
where $C_W$, $C_*$ and $c_3$ are the positive constants of the smallness index
$\hat{\vartheta}_n$, and we put
\begin{equation*}
C_{\hat{\vartheta}}:=\max\bigl\{1,(16C_WC_*c_3)^{-1/4}\bigr\}.
\end{equation*}

The classical rate of the weight bounds is $\vartheta_w:=L^{-\theta_A/8}$, where $\theta_A$ is the
exponent fixed with those bounds; the weight bounds are used under the condition
$\vartheta\ge\vartheta_w$. The rate $\vartheta_w$ is not $\vartheta_2$; the two are related by the
identity
\begin{equation}\label{9.5:eq-constituent-rates-1}
\vartheta_w^{\,n}=\vartheta_2^{\,n/4}\,\vartheta_\times^{\,n},
\end{equation}
in which the rate $\vartheta_\times$ satisfies $\vartheta_\times\le1$. The rate $\vartheta_{dl}$
fixed for the comparison satisfies $\vartheta_{dl}^{\,n}=\vartheta_f^{\,n/4}$. This is the
rate at which the classical marks of the comparison are carried.

\emph{Unit shares.} The comparison uses the bounds for the kept terms at four sites, labelled $P$,
$T$, $L^{\circ}$ and $L^{\dagger}$; the subscript of each share $\theta$ below records its site. The
subscript $T$ of $\theta_T$ is a site label, unrelated to the scaffold variable $\mathsf{T}$. The
\emph{unit share} of such a term is the coefficient read from its
bound by three changes:
\begin{itemize}
\item the size constant $\mathfrak{n}^{kept}=2^dC_{\Sigma}\mathfrak{s}_{\Sigma}$ is replaced by its
value $2^d\mathfrak{s}_{\Sigma}$ at unit sizes, where
$\mathfrak{s}_{\Sigma}:=\sup_{u\ge1}(2u+2)^{p_{\Sigma}}e^{-u}$;
\item the functions of the scaffold variable $\mathsf{T}$ are evaluated at their upper functions;
\item the rates are those of the run, with $\kappa_1-1$ replaced by $\kappa^{\natural}_1$.
\end{itemize}
In what follows $c_J$, $K^{\flat}_2$, $c_{loc}$, $c''_{\star}$, $C^{sum}_{\Sigma}$, $C_e$,
$C_{cell}$, $C^{\sharp}_0$, $\delta$ and $r_{pr}$ are the constants of the bounds for the kept
terms at the four sites; $\vartheta(\mathsf{T})$, $\bar{V}'(\mathsf{T})$ and
$\mathsf{R}^{\uparrow}(\mathsf{T})$ are the upper functions of $\mathsf{T}$ at which those bounds
are evaluated; and $N(\mathsf{T})$ is a function of the scaffold variable $\mathsf{T}$ fixed with
those bounds, not the rank $N$ of $\mathrm{SU}(N)$.
With
\begin{align*}
\mathsf{a}^{kept}_{\mu}(\mathsf{T})
&:=4e^{c_J+3/2}K^{\flat}_2\,\vartheta(\mathsf{T})\,c_{loc}\,2^d\mathfrak{s}_{\Sigma},\\
\mathsf{a}^{kept}_{\mathsf{M}}(\mathsf{T})
&:=2e^{c_J+3/2}K^{\flat}_2\,\vartheta(\mathsf{T})\,c''_{\star}\,\bar{V}'(\mathsf{T})\,
2^d\mathfrak{s}_{\Sigma},
\end{align*}
the rates, the unit shares and the derived quantities are the following. The first and third
lines hold for every zone $n$; the second holds on aligned zones only, in the sense of the
alignment condition of the comparison.
\begin{equation}\label{9.5:eq-constituent-rates-a}
\begin{aligned}
&\vartheta_f^{\,n/2}\le\vartheta_f^{\,n/4}\le C_{\hat{\vartheta}}\,\hat{\vartheta}^{\natural}_n
\le 2C_{\hat{\vartheta}}\,[\![\delta]\!]^{c}_n/E_0,\\
&(\log x_n)^{c_2}\,\vartheta_2^{\,n/4}\le[\![\delta]\!]^{c}_n/(E_0C^{\natural}_{cl}),\\
&\vartheta_w^{\,n/4}=\vartheta_2^{\,n/16}\vartheta_\times^{\,n/4}\le\vartheta_2^{\,n/16}
=\vartheta_f^{\,32n}\le\vartheta_f^{\,n/2}\le 2C_{\hat{\vartheta}}\,[\![\delta]\!]^{c}_n/E_0,\\
&\theta_P(\mathsf{T}):=N(\mathsf{T})\bigl[8\mathsf{a}^{kept}_{\mathsf{M}}(\mathsf{T})
+4\mathsf{a}^{kept}_{\mu}(\mathsf{T})\bigr],\\
&\theta_T(\mathsf{T}):=C^{sum}_{\Sigma}\,2^d\mathfrak{s}_{\Sigma}\,
e^{-\frac14\delta\kappa(7\mathsf{T}^{r_{pr}}-2)},\\
&\theta^{\circ}_L(\mathsf{T}):=4(1+C^{sum}_{\Sigma}C_e)^3C^{sum}_{\Sigma}\,2^d\mathfrak{s}_{\Sigma}\,
e^{-(\mathsf{T}^{r_{pr}}-2)},\\
&\theta^{\dagger}_L(\mathsf{T}):=(1+C^{sum}_{\Sigma}C_e)^4\,C_{\mathfrak{K}}\,e\,
2^d\mathfrak{s}_{\Sigma}\,C_{cell}\\
&\qquad\qquad\times 2d\,\bigl(104L\mathsf{R}^{\uparrow}(\mathsf{T})+2\bigr)^{d-1}
e^{-(\mathsf{T}^{r_{pr}}-2)},\\
&\bar{\Theta}:=\sup_{\mathsf{T}\ge\mathsf{T}_{\gamma}}\Bigl[\frac{4\,\theta_P(\mathsf{T})}
{3N(\mathsf{T})C^{\sharp}_0}+\frac13\,\theta^{\dagger}_L(\mathsf{T})\Bigr],\\
&\mathsf{k}^{rd}_n:=(32C^{\star}_{TL})^{1/2}C_{\Sigma}\,\vartheta_f^{\,n/2},\qquad
\mathsf{k}^{+,(<m)}_n:=\mathsf{k}^{(<m)}_n+\mathsf{k}^{rd}_n,\\
&\varsigma^{1/2}_{\nu}:=\bigl[\nu^3(1+\rho)L^{-\hat{\alpha}(\nu-1)/2}
+L^{-(\nu-1)/2}\bigr]\rho^{\,2}.
\end{aligned}
\end{equation}
Here $\mathsf{k}^{(<m)}_n$ is the quotient of \eqref{9.5:eq-production-bound-c} in the production
bound, layer by layer (Theorem~\ref{9.5:thm-production-bound}).

The share $\theta_P$ has degree $r+4r+0^{+}-(7r+1)$ in the scaffold variable $\mathsf{T}$, in the
sense of degrees used in the bounds for the kept terms at the four sites. The shares $\theta_T$,
$\theta^{\circ}_L$ and
$\theta^{\dagger}_L$ each have degree $-\infty$. The supremand of $\bar{\Theta}$ has degree
$4r+0^{+}-(7r+1)$.

The share $\theta^{\dagger}_L$ is the bound for the kept terms at the site $L^{\dagger}$ with the
constant $\tfrac12C_{3F}$ of that bound replaced by $C_{\mathfrak{K}}$. Of this constant,
$\frac12C_{\mathfrak{K}}$
serves the family marks and $C_{\mathfrak{K}}^{1/2}$ serves the geometric mean taken in the reading
mark; here $C_{\mathfrak{K}}\ge2$, and $C_{\mathfrak{K}}^{1/2}\le2^{-1/2}C_{\mathfrak{K}}$.

\emph{The weight read at $\hat{\alpha}/2$.} The definition of the class loads writes the weight
\begin{equation*}
\varsigma_{\nu}=\bigl[\nu^3(1+\rho)L^{-\hat{\alpha}(\nu-1)}+L^{-(\nu-1)/2}\bigr]\rho^{\,2}
\end{equation*}
at the exponent $\hat{\alpha}$ fixed with these weights. The classical part of a near group, that
is, of a group of terms carrying the group weight $\varsigma_{\nu}$, is bounded by a geometric
mean of
two bounds, so it is available at the exponent $\hat{\alpha}/2$ only. Accordingly, in
$[\![\delta]\!]^{c}_n$ the entry
\begin{equation*}
\sum_{j\le n}\varsigma_{1+n-j}E_0\,\hat{\vartheta}^{\natural}_j
\end{equation*}
is read with $\varsigma^{1/2}_{\nu}$ in place of $\varsigma_{\nu}$. The reading at
$\hat{\alpha}/2$ concerns this entry only; the other entries stay at $\hat{\alpha}$.
\end{definition}

\begin{proof}
\emph{The first line of \eqref{9.5:eq-constituent-rates-a}.} Since $0<\vartheta_f\le1$, we have
$\vartheta_f^{\,n/2}\le\vartheta_f^{\,n/4}$. Put $a:=16C_WC_*c_3$. Taking the square root of
$\hat{\vartheta}_n$ gives
\begin{equation*}
\hat{\vartheta}^{\natural}_n=\min\bigl\{\sqrt2,\ \sqrt2\,a^{1/4}\vartheta_f^{\,n/4}\bigr\}.
\end{equation*}
Suppose the minimum is $\sqrt2$. Then $C_{\hat{\vartheta}}\ge1$ gives
\begin{equation*}
C_{\hat{\vartheta}}\hat{\vartheta}^{\natural}_n\ge\sqrt2\ge1\ge\vartheta_f^{\,n/4}.
\end{equation*}
Otherwise $C_{\hat{\vartheta}}\ge a^{-1/4}$ gives
\begin{equation*}
C_{\hat{\vartheta}}\hat{\vartheta}^{\natural}_n=\sqrt2\,C_{\hat{\vartheta}}\,a^{1/4}\vartheta_f^{\,n/4}
\ge\sqrt2\,\vartheta_f^{\,n/4}.
\end{equation*}
In both cases $\vartheta_f^{\,n/4}\le C_{\hat{\vartheta}}\hat{\vartheta}^{\natural}_n$.

For the last inequality we use the entry $\sum_{j\le n}\varsigma_{1+n-j}E_0\hat{\vartheta}^{\natural}_j$
of $[\![\delta]\!]^{c}_n$. All entries of the classical part are nonnegative, and the term $j=n$ of
this entry is $\varsigma_1E_0\hat{\vartheta}^{\natural}_n$. Hence
$[\![\delta]\!]^{c}_n\ge\varsigma_1E_0\hat{\vartheta}^{\natural}_n$. This holds also under the reading
at $\hat{\alpha}/2$, because at $\nu=1$ the exponents of $L$ in $\varsigma_{\nu}$ vanish, so
$\varsigma^{1/2}_1=\varsigma_1$. By the definition of the class loads, $\varsigma_1=(2+\rho)\rho^2$.
Since $\rho\ge\frac12$ by the conditions under which the contraction letter $\rho$ of the comparison
is chosen,
\begin{equation*}
\varsigma_1\ge\tfrac52\cdot\tfrac14=\tfrac58\ge\tfrac12 .
\end{equation*}
Therefore
\begin{equation*}
C_{\hat{\vartheta}}\hat{\vartheta}^{\natural}_n\le C_{\hat{\vartheta}}[\![\delta]\!]^{c}_n/(\varsigma_1E_0)
\le2C_{\hat{\vartheta}}[\![\delta]\!]^{c}_n/E_0 .
\end{equation*}

\emph{The second line.} The last entry of $[\![\delta]\!]^{c}_n$ is
\begin{equation*}
E_0\min\bigl\{1,C^{\natural}_{cl}(\log x_n)^{c_2}\vartheta_2^{\,n/4}\bigr\}.
\end{equation*}
Since the entries are nonnegative, $[\![\delta]\!]^{c}_n$ is at least this entry. On an aligned
zone, that is, on a zone satisfying the alignment condition of the comparison, this minimum is
attained at its second member. Dividing by $E_0C^{\natural}_{cl}$ gives the second line.

\emph{The third line.} The rate $\vartheta_w^{\,n}$ enters the constituent bounds after two
geometric means, that is, as $\vartheta_w^{\,n/4}$. Taking fourth roots in
\eqref{9.5:eq-constituent-rates-1} gives the equality
$\vartheta_w^{\,n/4}=\vartheta_2^{\,n/16}\vartheta_\times^{\,n/4}$. The factor
$\vartheta_\times^{\,n/4}$ is at most one and is dropped. The equality
$\vartheta_2^{\,n/16}=\vartheta_f^{\,32n}$ holds by the definitions of $\vartheta_2$ and
$\vartheta_f$ as powers of $\vartheta$.
Since $\vartheta_f\le1$ and $32n\ge n/2$, we get $\vartheta_f^{\,32n}\le\vartheta_f^{\,n/2}$. The
last inequality is the first line.

\emph{No new loss of rate.} By the first and third lines, $\vartheta_f^{\,n/2}$ and
$\vartheta_w^{\,n/4}$ are at most $\vartheta_f^{\,n/4}=\vartheta_{dl}^{\,n}$. The quotient
$\mathsf{k}^{rd}_n$ is a constant times $\vartheta_f^{\,n/2}$, so it is also at most a constant
times $\vartheta_{dl}^{\,n}$. Thus each of the classical marks $\vartheta_f^{\,n/2}$,
$\vartheta_w^{\,n/4}$ and $\mathsf{k}^{rd}_n$ entering the constituent bounds through
\eqref{9.5:eq-constituent-rates-a} is at most a constant times $\vartheta_{dl}^{\,n}$. No rate
slower than the one fixed for the comparison is used, and the range in which the contraction
letter $\rho$ is chosen is not changed.

\emph{The constant of $\theta^{\dagger}_L$.} Since $C_{\mathfrak{K}}\ge2$, we have
$C_{\mathfrak{K}}^{1/2}\ge2^{1/2}$. Multiplying by $2^{-1/2}C_{\mathfrak{K}}^{1/2}$ gives
$C_{\mathfrak{K}}^{1/2}\le2^{-1/2}C_{\mathfrak{K}}$.

\emph{Convergence of the envelope at $\hat{\alpha}/2$.} Read the class loads with
$\varsigma^{1/2}_{\nu}$ in the classical entry. The part of the envelope bound for the class loads
that is proportional to $E_0$ then reads, at depth $s$,
\begin{equation*}
E_0C_{\vartheta}\rho^{\,n-s}\cdot6\sum_{i\ge0}(i+1)^3\Bigl[\bigl(L^{-\hat{\alpha}/2}/\rho\bigr)^i
+\bigl(L^{-1/2}/\rho\bigr)^i\Bigr],
\end{equation*}
where $C_{\vartheta}$ denotes the constant standing in front of this part of the envelope bound,
distinct from
$C_{\hat{\vartheta}}$. This series is finite because $\rho\ge q^{1/2}>L^{-\hat{\alpha}/2}$. The
second inequality is strict by the condition $q>L^{-\hat{\alpha}}$ imposed on $q$. Hence the ratio
$L^{-\hat{\alpha}/2}/\rho$ is strictly less than one.
For $0\le y<1$ the series $\sum_{i\ge0}(i+1)^3y^i$ is finite and depends on $y$ only. The
second geometric ratio $L^{-1/2}/\rho$ does not involve $\hat{\alpha}$ and is not changed by the
substitution $\hat{\alpha}\mapsto\hat{\alpha}/2$. Consequently the constant of the envelope bound
that absorbs the value of this series, which is not $C_{\vartheta}$, is replaced by its maximum
with a number depending only on $(q,L,\hat{\alpha})$. No range in the order of
choice is changed, and the other entries remain at $\hat{\alpha}$.
\end{proof}

\begin{lemma}[Pointed cell sums linear in size constants]\label{9.5:lem-pointed-sums}
Assume the bounds, for each kind of pointed unit, established in the proof of the collection bound for
pointed units, and the table of sizes of the pointed units on which they rest. Let $\varsigma_{\nu}$ $(\nu\ge1)$, $\rho$,
$\bar\rho_{\nu}$ and $w_{n,j}$ be the rates and weights of
Notation~\ref{9.5:not-constituent-rates}, and let $\hat\alpha\in\,]0,1[$ be the exponent that
appears as $1-\beta$ in \cite[(3.67)]{B-CMP119}. Assume further:
\begin{enumerate}
\item[(a)] $\bar\rho_{\nu}^2\le 3\le 12\rho^2$ for every index $\nu$ at which $\bar\rho_{\nu}$ is
defined; in particular $\rho^2\ge\frac14$;
\item[(b)] $\varsigma_{\nu}\ge\rho^2L^{-(\nu-1)/2}$ and $\varsigma_{\nu}\ge\nu^3\rho^2L^{-\hat\alpha(\nu-1)}$
for all $\nu\ge1$, and the same two inequalities with $\varsigma^{1/2}_{\nu}$ and $\hat\alpha/2$ in
place of $\varsigma_{\nu}$ and $\hat\alpha$, this second pair being the form of (b) used for the
second assertion below;
\item[(c)] $w_{n,j}\ge\rho^2$ for all $j\le n$;
\item[(d)] the cell weight $\mathsf{w}^{\vee}(\Delta')$ is at least $1$.
\end{enumerate}
Let $k\ge n$ be the level at which the weights are read: $d_k(K_v)$ is the distance function of
level $k$ evaluated on the region $K_v$ of the unit $v$, $r$ is the decay rate with which the table
of sizes weights it, and $\omega_{k,j}$ below is the weight of level $j$ read at level $k$. The
numbers $E_0$, $g_j^{\kappa_0}$ and $b_{+}+r$ are the size constants that the table of sizes
assigns to its entries at level $j$, and $\mathfrak{g}_n$ is the coupling factor of level $n$
carried in that table by the counterterm pieces and by the fifth entry.
In the proof of the collection bound for pointed units replace the size constants of the entries
of the table by arbitrary numbers $\ge0$, level by level, as follows: $E_0\mapsto e_j$ in the first
entry, in the boundary entry, in the tails, and inside the prefactor $C_BE_0\,\omega_{k,j}$ of the
second entry, the weight $\omega_{k,j}=L^{-(k-j)/2}$ being kept (the first, second and boundary
entries are affine in one and the same list $E^{(j)}$ of size constants of level $j$);
$g_j^{\kappa_0}\mapsto r_j$ in the third entry; $b_{+}+r\mapsto\beta_j$ in the counterterm pieces
and in the fifth entry. Write $\mathrm{sz}(v)$ for the size of a pointed unit $v$ after this
replacement. The pointed units fall into the kinds distinguished in that proof: short and long
singles, tails, smooth and rough cores, rim units, Wilson pieces and the unit of the fifth entry;
$\square^{\sim\nu}$ denotes the enlarged cube of that proof. Then for every cell $\Delta'$ of level $n$
\begin{equation}\label{9.5:eq-pointed-sums-a}
\begin{split}
\sum_{v:\,y_v\in\Delta'}\mathrm{sz}(v)\,e^{r\,d_k(K_v)}e^{\varsigma_1\mathsf{r}_v}
\le K_{pt}\,\mathsf{w}^{\vee}(\Delta')^3\Bigl[&\sum_{j\le n}\varsigma_{1+n-j}\,e_j
+\sum_{j\le n}w_{n,j}\,r_j\\
&+\mathfrak{g}_n\sum_{j\le n}L^{-(n-j)/2}\beta_j\Bigr],
\end{split}
\end{equation}
and the same bound holds with $(\hat\alpha,\varsigma_{\nu})$ replaced by
$(\hat\alpha/2,\varsigma^{1/2}_{\nu})$. Here $K_{pt}$ is the constant defined in
\eqref{9.5:eq-pointed-sums-1} below; it depends on $d$, $L$, $M$ and $\kappa$ only.
\end{lemma}

\begin{proof}
Fix the cell $\Delta'$ of level $n$ and a level $j\le n$. We put
\begin{equation*}
\nu:=1+n-j,\qquad t:=L^{-(\nu-1)},\qquad \mathsf{W}:=\mathsf{w}^{\vee}(\Delta'),
\end{equation*}
so that $\nu\ge1$, $0<t\le1$ and $\mathsf{W}\ge1$ by (d). All bounds below are bounds for the
contribution of the units of level $j$ to the left side of \eqref{9.5:eq-pointed-sums-a}; summing
them over $j\le n$ gives the right side.

\emph{The weights.} The part of the proof of the collection bound for pointed units that absorbs
the weights involves no size constant, and therefore holds unchanged after the replacement: the weights
$e^{r\,d_k(K_v)}e^{\varsigma_1\mathsf{r}_v}$ are absorbed by the decay rates carried by the entries
themselves, the part $\kappa/20$ of the decay rate pays for the sum over the domains, and for cores
and Wilson pieces the weights cost at most the factor
\begin{equation*}
\mathsf{c}_0:=e^{(8/5)\kappa 6^d+5\varsigma_1}.
\end{equation*}
The cell $\Delta'$ contains $M^4t^{-4}$ points of the lattice $\mathbb{T}^{(j)}$ of level $j$.

Besides (a)--(c) we use the elementary fact that
\begin{equation*}
\mathsf{c}_L:=\sup_{\nu\ge1}(\nu+4)\nu^2L^{-(\nu-1)/2}<\infty ,
\end{equation*}
which holds because $L>1$, so that the exponential factor dominates the polynomial one. From (a)
and (b) we also have $L^{-(\nu-1)/2}\le 4\rho^2L^{-(\nu-1)/2}\le4\varsigma_{\nu}$, and from (a)
and (c) we have $1\le4\rho^2\le4w_{n,j}$.

\emph{Short singles.} For a short single of level $j$ with prefactor $\mathfrak{n}_j$, the proof of
the collection bound gives the bound
$2C(\kappa)(4M\mathsf{W}c_{\vee}C_M)^2\bar\rho_{\nu}^2t^4\mathfrak{n}_j$ per point. Multiplying by
the $M^4t^{-4}$ points of $\Delta'$, the contribution of level $j$ is at most
$\mathsf{c}_{sh}\mathsf{W}^2\bar\rho_{\nu}^2\mathfrak{n}_j$, where
\begin{equation*}
\mathsf{c}_{sh}:=2C(\kappa)(4Mc_{\vee}C_M)^2M^4 .
\end{equation*}
In the second entry $\mathfrak{n}_j=C_B\omega_{k,j}e_j$; since $k\ge n$ we have
$\omega_{k,j}=L^{-(k-j)/2}\le L^{-(n-j)/2}=L^{-(\nu-1)/2}$, hence
$\mathfrak{n}_j\le C_BL^{-(\nu-1)/2}e_j$. By (a), $\bar\rho_{\nu}^2\le12\rho^2$, and by (b),
$\rho^2L^{-(\nu-1)/2}\le\varsigma_{\nu}$; the contribution is therefore at most
$12\mathsf{c}_{sh}C_B\mathsf{W}^2\varsigma_{\nu}e_j$. In the third entry $\mathfrak{n}_j=r_j$; by (a)
and (c), $\bar\rho_{\nu}^2\le12\rho^2\le12w_{n,j}$, and the contribution is at most
$12\mathsf{c}_{sh}\mathsf{W}^2w_{n,j}r_j$. In the boundary entry, with $\mathfrak{n}_j=e_j$, only
$27dM(\nu+2)t^{-3}$ points of level $j$ near the boundary carry a unit, instead of $M^4t^{-4}$;
the contribution is thus at most
\begin{equation*}
\frac{27d}{M^3}\,\mathsf{c}_{sh}\mathsf{W}^2(\nu+2)\bar\rho_{\nu}^2\,t\,e_j
\le\frac{324d}{M^3}\,\mathsf{c}_{sh}\mathsf{c}_L\mathsf{W}^2\varsigma_{\nu}e_j ,
\end{equation*}
where we used $\bar\rho_{\nu}^2\le12\rho^2$ from (a),
\begin{equation*}
(\nu+2)t\le(\nu+4)\nu^2L^{-(\nu-1)/2}L^{-(\nu-1)/2}\le\mathsf{c}_LL^{-(\nu-1)/2},
\end{equation*}
and $\rho^2L^{-(\nu-1)/2}\le\varsigma_{\nu}$ from (b).

\emph{Long singles.} The proof of the collection bound gives the bound $C_Lt^5\mathfrak{n}_j$ per
point, hence $C_LM^4t\,\mathfrak{n}_j$ per level. For the entries carrying $e_j$ we have
$\mathfrak{n}_j\le(1+C_B)e_j$ (as $\omega_{k,j}\le1$) and $t\le L^{-(\nu-1)/2}\le4\varsigma_{\nu}$;
for the third entry $\mathfrak{n}_j=r_j$ and $t\le1\le4w_{n,j}$. The contribution is therefore at
most $4C_LM^4(1+C_B)\varsigma_{\nu}e_j$ for the part carrying $e_j$, and at most
$4C_LM^4w_{n,j}r_j$ for the part carrying $r_j$.

\emph{Tails.} The tails are bounded by $2K(d)e_jt^5$ per point, hence by $2K(d)M^4t\,e_j$ per
level, and, since $t\le4\varsigma_{\nu}$, by $8K(d)M^4\varsigma_{\nu}e_j$.

\emph{Smooth cores.} Put $a=a_0:=2+s(\mathsf{W})$ and $b=\nu-a_0$, with $s(\cdot)$ the function of
the proof of the collection bound for pointed units. That proof bounds the contribution of the
smooth cores of level $j$ by
\begin{equation*}
C^{sm}M^4\mathsf{W}^2(2a_0+11)^2\bar\rho^2_{a_0}(b+1)^3\bar\rho_b^2(1+\bar\rho_b)\,
L^{-\hat\alpha b}e_j .
\end{equation*}
We have $L^{-\hat\alpha b}=L^{\hat\alpha(a_0-1)}L^{-\hat\alpha(\nu-1)}$, and, by the inequality
$s(\mathsf{W})\le s(1)+\log_L\mathsf{W}$ established in that proof,
\begin{equation*}
L^{\hat\alpha(a_0-1)}=L^{\hat\alpha(1+s(\mathsf{W}))}\le L^{\hat\alpha(1+s(1))}\mathsf{W}^{\hat\alpha}.
\end{equation*}
Further $(b+1)^3\le\nu^3$. By (a), $\bar\rho_b^2\le12\rho^2$ and $\bar\rho_b\le\sqrt3\le2$, so
$1+\bar\rho_b\le3(1+\rho)$ and
\begin{equation*}
\bar\rho_b^2(1+\bar\rho_b)\le36\rho^2(1+\rho).
\end{equation*}
Finally, since $\hat\alpha\in\,]0,1[$, $\bar\rho^2_{a_0}\le3$ by (a), and
$a_0\le2+s(1)+\log_L\mathsf{W}$, the factor $(2a_0+11)^2\bar\rho^2_{a_0}$ is at most a constant times
the square of a logarithm of $\mathsf{W}$, which is bounded by a constant times
$\mathsf{W}^{1-\hat\alpha}$; thus $(2a_0+11)^2\bar\rho^2_{a_0}\mathsf{W}^{\hat\alpha}\le c_{\log}\,\mathsf{W}$
with a constant $c_{\log}$ depending only on $\hat\alpha$, $L$ and $s(1)$. Collecting these bounds and using
$\nu^3\rho^2L^{-\hat\alpha(\nu-1)}\le\varsigma_{\nu}$ from (b), the contribution is at most
$\mathsf{c}_{sm}\mathsf{W}^3\varsigma_{\nu}e_j$, where
\begin{equation*}
\mathsf{c}_{sm}:=36(1+\rho)\,c_{\log}\,C^{sm}M^4L^{\hat\alpha(1+s(1))} .
\end{equation*}

\emph{The rim $a>a_0$, and rough cores with $\square^{\sim\nu+3}\not\subset Z$.} These units sit in
thin layers of $\Delta'$. The numbers of points carrying them are
\begin{equation*}
2dM^4t^{-4}\bigl[(\nu+2)L^{-(\nu-1)}+3L^{-(a-1)}\bigr]\quad\text{and}\quad 2d(\nu+4)M^4t^{-3},
\end{equation*}
respectively; compared with the full count $M^4t^{-4}$ they put a factor at most
$\mathsf{c}_LL^{-(\nu-1)/2}$. With a constant $\mathsf{c}_{rm}$ so obtained, the contribution is at
most $\mathsf{c}_{rm}\mathsf{W}^2\varsigma_{\nu}e_j$ for the part carrying $e_j$, and at most
$\mathsf{c}_{rm}\mathsf{W}^2\mathfrak{g}_nL^{-(\nu-1)/2}\beta_j$ for the part carrying $\beta_j$.

\emph{Rough cores with $\nu\le2+s(\mathsf{W})$.} By the first form of the table of sizes, the size
of such a unit $v$ is the sum, over the regions $X\ni y_v$ with $X\subset\square^{\sim\nu}$, of the
sizes of its constituents, plus the counterterm piece $6L^4\beta_j\mathfrak{a}'^{\,2}_vt^4$; it is
therefore at most
\begin{equation*}
C^{mar}\mathsf{W}^2\nu^2\bar\rho_{\nu}^2e_jt^4+6L^4\beta_j\mathfrak{a}'^{\,2}_vt^4 ,
\end{equation*}
and per level this is multiplied by $M^4t^{-4}$. Here, in the proof of the collection bound,
$\mathfrak{a}'^{\,2}_v=c_{reg}^2\mathsf{W}^2\alpha_{*,n}^2$, which equals
$c_{reg}^2C_*^2\mathsf{W}^2\mathfrak{g}_n$ up to the parameter $q_{\flat}$ of that proof. The
counterterm piece occurs for rough cores only: in a smooth core the counterterm is contained in the
renormalised group functional $\mathfrak{i}_y$. The values of $\nu$ in question satisfy
$2+s(\mathsf{W})-\nu\ge0$, hence, by $s(\mathsf{W})\le s(1)+\log_L\mathsf{W}$,
\begin{equation*}
1\le L^{(1+s(\mathsf{W}))/2}L^{-(\nu-1)/2}\le c\,\mathsf{W}^{1/2}L^{-(\nu-1)/2},
\qquad c=L^{(1+s(1))/2}.
\end{equation*}
Multiplying the per-level bound by this factor, using $\bar\rho_{\nu}^2\le12\rho^2$ from (a),
$\rho^2L^{-(\nu-1)/2}\le\varsigma_{\nu}$ from (b), bounding $\nu^2\le(2+s(1)+\log_L\mathsf{W})^2$ by
$c''\mathsf{W}^{1/2}$ with a constant $c''$ depending only on $L$ and $s(1)$, and
$\mathsf{W}^{5/2}\le\mathsf{W}^3$, one obtains, with a constant
$\mathsf{c}_{rg}$, the bound $\mathsf{c}_{rg}\mathsf{W}^3\varsigma_{\nu}e_j$ for the part carrying
$e_j$, and $\mathsf{c}_{rg}\mathsf{W}^3\mathfrak{g}_nL^{-(\nu-1)/2}\beta_j$ for the part carrying
$\beta_j$.

\emph{The fifth entry.} Its size is $c_5\mathsf{W}^2\mathfrak{N}(v)$, with $c_5$ the constant that
the table of sizes attaches to the fifth entry; there is one such unit per
cell, and $\mathfrak{N}(v)$ is proportional to $\mathfrak{g}_n\beta_n$, with no dependence on $K$.
This is the term $j=n$ of the third sum in \eqref{9.5:eq-pointed-sums-a}.

\emph{Conclusion.} Since $\mathsf{W}\ge1$, each of the bounds above is at most the corresponding
term of the bracket in \eqref{9.5:eq-pointed-sums-a} times $\mathsf{W}^3$ times its constant, and
$L^{-(\nu-1)/2}=L^{-(n-j)/2}$. Multiplying by the factor $\mathsf{c}_0$ for the weights and summing
over the kinds of units and over $j\le n$ gives \eqref{9.5:eq-pointed-sums-a} with
\begin{equation}\label{9.5:eq-pointed-sums-1}
\begin{split}
K_{pt}:=\mathsf{c}_0\Bigl[&12\mathsf{c}_{sh}(C_B+1)+\frac{324d}{M^3}\,\mathsf{c}_{sh}\mathsf{c}_L
+4C_L(C_B+2)M^4\\
&+8K(d)M^4+\mathsf{c}_{sm}+\mathsf{c}_{rm}+\mathsf{c}_{rg}+c_5\Bigr].
\end{split}
\end{equation}
Here $12\mathsf{c}_{sh}(C_B+1)$ collects the short singles of the second and third entries, and
$4C_L(C_B+2)M^4$ collects the two parts of the long singles.

For the pair $(\hat\alpha/2,\varsigma^{1/2}_{\nu})$ the same argument applies, with (b) used in the
form stated there for that pair; among the constants in \eqref{9.5:eq-pointed-sums-1}, only
$\mathsf{c}_{sm}$ and $\mathsf{c}_{rg}$ change.

Thus every scale sum in the collection bound for pointed units is a point count, times the
prefactor of level $j$, times a weight dominated by the matching weight on the right of
\eqref{9.5:eq-pointed-sums-a}; the bound is linear in the prefactors, which is the form of the
third line of \eqref{9.5:eq-production-bound-a}.
\end{proof}

\begin{proposition}[Two-run difference of field-independent remainder terms. Proof outstanding]
\label{9.5:prop-remainder-difference}
Let $\mathsf{A}$ and $\mathsf{B}$ be the two runs compared, let $n\le m$ be a zone, $m$ the step
index of Definition~\ref{9.5:def-layer-store}, and let
$C^{\sharp}$ and the terms $\delta\mathfrak{v}^{\mathbf{C},(n)}(Y_4)$ be as in
Definition~\ref{9.5:def-layer-store}. The supremum below is taken over the common data and the
real fields of that definition. Let $\mathfrak{g}_n\mathsf{z}^{\sharp}_n$ denote the entry so written
of the five-unknown part $[\![\delta]\!]^{(5)}_n$ of the defect of layer $n$, and
let $E_0$ and $\vartheta_w$ be as in Notation~\ref{9.5:not-constituent-rates}. The $z$-free
members of Ba{\l}aban's field-independent remainder terms $\mathbf{C}_m$, that is, those that do not
depend on the source coefficients $z_i$, localised in $C^{\sharp}$ and split according to the lowest
zone they meet, obey, with a constant $C^{\times}_{\mathbf{C}}$,
\begin{equation}\label{9.5:eq-remainder-difference-a}
\sum_{Y_4\subset C^{\sharp}}\sup\bigl|\delta\mathfrak{v}^{\mathbf{C},(n)}(Y_4)\bigr|
\le C^{\times}_{\mathbf{C}}\bigl[\mathfrak{g}_n\mathsf{z}^{\sharp}_n
+E_0(\log x_n)^{c_2}\vartheta_w^{n}\bigr].
\end{equation}
\end{proposition}

\smallskip
\noindent\textit{Route.}
A proof would establish \eqref{9.5:eq-remainder-difference-a} by splitting the $z$-free members
of $\mathbf{C}_m$ localised in $C^{\sharp}$ according to the lowest zone they meet, and by
treating them by kind rather than member by member: \cite{B-CMP122a} introduces the family
$\mathbf{C}_k$ and states that its details are omitted, and no list of its members is given there.
By the representation used in the proof of the correlation theorem, the $z$-free members of
$\mathbf{C}_m$ are weight-linear Wilson differences and functions of orbits. Their bound within
one run is an input of the one-run analysis and is used as such in each of $\mathsf{A}$ and
$\mathsf{B}$.

The plaquette weights $c(p)$ of these members are the couplings of the run in which they are
read, with $|c(p)|$ bounded by the position-dependent inverse coupling at $p$
(Definition~\ref{1.5:def-domains}), and the couplings of $\mathsf{A}$ and $\mathsf{B}$ differ.
Write $A^{\mathsf{m}}(c,\cdot)$
for the weighted Wilson sum of run $\mathsf{m}$ with weights $c$, which is linear in $c$, and put
$\delta c:=c^{\mathsf{B}}-c^{\mathsf{A}}$. Adding and subtracting $A^{\mathsf{A}}(c^{\mathsf{B}},\cdot)$
and using this linearity gives
\begin{equation*}
A^{\mathsf{B}}(c^{\mathsf{B}},\cdot)-A^{\mathsf{A}}(c^{\mathsf{A}},\cdot)
=A^{\mathsf{A}}(\delta c,\cdot)
+\bigl[A^{\mathsf{B}}(c^{\mathsf{B}},\cdot)-A^{\mathsf{A}}(c^{\mathsf{B}},\cdot)\bigr].
\end{equation*}
The first term $A^{\mathsf{A}}(\delta c,\cdot)$, whose argument in run $\mathsf{A}$ is the point of
$\mathsf{A}$ itself, is a function on the lattice of $\mathsf{A}$; by a bound of the correlation
theorem it is to be bounded by a constant times
$\sup|\delta c|\,\mathfrak{a}^{2}$, and it is linear in $(\mathsf{z},\mathsf{b})$. A proof must show
that this
part is controlled by the entry $\mathfrak{g}_n\mathsf{z}^{\sharp}_n$ of
$[\![\delta]\!]^{(5)}_n$, linearly and under no root; this is the first term on the right of
\eqref{9.5:eq-remainder-difference-a}.

The second term is the two-run difference of a weighted Wilson sum on slices with the common
weights $c^{\mathsf{B}}$; it is the operator part and is classical, and a proof would bound it by
the classical operator bound on the two-run difference of a weighted Wilson sum on slices. For the
members that are functions of orbits, a proof would use the classical operator bound on the
two-run difference of a $z$-free potential, which states that the activities of such differences
are at most $C_{cl}\vartheta^{n}(\log x_n)^{c_2}\bar{\varepsilon}^{pr}$, in the notation of that
bound. Both bounds are to be applied zone by
zone, in the same way as the boundary terms created by the operation $\mathbf{T}$ are split by
zones, and the resulting contributions summed over $Y_4\subset C^{\sharp}$. What a proof must
show is that the sum is bounded by
$C^{\times}_{\mathbf{C}}E_0(\log x_n)^{c_2}\vartheta_w^{n}$, the second term on the right of
\eqref{9.5:eq-remainder-difference-a}.

\begin{proposition}[Bounds on the constituents of the production defect]\label{9.5:prop-constituent-bounds}
Write $\mathcal{B}^{\sharp}_1$ and $\mathcal{B}^{\sharp}_2$ for the two balls of the sharpened rim
conditions on the dial polydisc $\mathbb{D}^{\star}$ (they enter the proof of line $(\ell 5)$ below).
Let $\mathsf{R}^{\sharp}$ be the radius of $\mathcal{B}^{\sharp}_1$ in the norm $\|\cdot\|_a$ on
$\mathbb{D}^{\star}$, and put
\begin{equation}\label{9.5:eq-constituent-bounds-1}
K_{rim}:=\frac{2\mathsf{R}^{\sharp}}{E^z_1}.
\end{equation}
Assume the following.
\begin{enumerate}
\item[(1)] The hypotheses \textup{(F0)}--\textup{(F5)} of Theorem~\ref{9.5:thm-production-bound}; the
inputs (a)--(h) of the uniqueness theorem; and the
comparison of the two units of distance in \eqref{9.5:eq-distance-units-c}.
\item[(2)] As induction hypothesis: the kept-member bound on the production defect at all sites before
the kept-member site of step $m$.
\item[(3)] Proposition~\ref{9.5:prop-remainder-difference}, whose proof is outstanding.
\item[(4)] Lemma~\ref{9.5:lem-pointed-sums}.
\item[(5)] The one-run bound on the boundary-term stores and its plaquette form; part (b) of the
zone-crossing lemma; the chain-term comparison theorem; part (ii) of the lemma on resolved terms;
the kept-member load statement; the bound on the terms of interrupted runs; the lemma on the
counterterm and weight entries of differences; the group corollary for resolved extensions; the
rim-total lemma; the table of direct uses on the rim; the corollary on strip bounds for differences
and marks; the proposition on the classical operator rows; the complexification corollary for
operator differences; and part (ii) of the copy lemma.
\end{enumerate}
Then for every aligned zone $n\le m$
\begin{equation}\label{9.5:eq-constituent-bounds-2}
\mathsf{d}^{(n)}_m\le K_D\,[\![\delta]\!]^{\natural+,(m)}_n+\bar{\Theta}\,\mathsf{k}^{+,(<m)}_n,
\qquad
K_D:=\max\{K^{B}_D,K^{pt}_D,K^{cp}_D,K^{rim}_D,K^{\mathbf{C}}_D,K^{\mathfrak{q}}_D\},
\end{equation}
with the constants $K^{B}_D,\dots,K^{\mathfrak{q}}_D$ fixed in the proof below. In detail, writing
$[\![\delta]\!]_n$ for $[\![\delta]\!]^{\natural+,(m)}_n$, the constituents of the seven lines of
\eqref{9.5:eq-production-bound-a} contribute to $\mathsf{d}^{(n)}_m$ at most the following amounts.
\begin{enumerate}
\item[(i)] Lines $(\ell 1)$, $(\ell 2)$, family marks of the boundary terms $\mathbf{B}^{(j)}$,
$\mathbf{B}'^{(j)}$ (resolved and own parts) and of the wrapped terms:
$\tfrac{1}{2B_0}(\mathsf{y}_n+\mathsf{y}'_n)$.
\item[(ii)] Line $(\ell 2)$, family marks of the frozen and left-over chain terms, of the chain terms
$\mathbb{C}^{(\nu)}_m$ of this step, and of the terms of $\mathfrak{F}^{B4}$ other than
$\mathsf{q}_W$:
\begin{equation*}
K^{fr}_{\mathbb{C}}[\![\delta]\!]_n\Bigl(\frac{1}{C^{\sharp}_0}+\frac{1}{2K(d)B_0}\Bigr)
+\sup_{\mathsf{b}\le m,\ \mathsf{T}}\frac{4\theta_P(\mathsf{T})}{3N(g_{\mathsf{b}})C^{\sharp}_0}\,
\mathsf{k}^{+,(<m)}_n ,
\end{equation*}
where the supremum runs over the scaffold variable $\mathsf{T}$ and over the steps $\mathsf{b}\le m$
at which chain terms are run, and $N(g_{\mathsf{b}})\ge1$ is the number of stages
(Definition~\ref{9.1:not-levelwise-rate}) at the step $\mathsf{b}$.
\item[(iii)] Lines $(\ell 1)$, $(\ell 2)$, reading marks of all stored units:
\begin{equation*}
\mathsf{c}_{rd}\vartheta_f^{n/2}\le 2\mathsf{c}_{rd}C_{\hat{\vartheta}}[\![\delta]\!]^{c}_n/E_0 .
\end{equation*}
\item[(iv)] Line $(\ell 2)$, the Gaussians $\mathfrak{Q}(\mathfrak{a})$:
\begin{equation*}
\mathsf{c}_G\vartheta_w^{n/4}\le 2\mathsf{c}_GC_{\hat{\vartheta}}[\![\delta]\!]^{c}_n/E_0 .
\end{equation*}
\item[(v)] Line $(\ell 3)$, the entries numbered $1$--$3$, $5$ and $\partial$ of the flat column sums:
$K^{pt}_D[\![\delta]\!]_n$.
\item[(vi)] Line $(\ell 4)$, the copies: $K^{cp}_D[\![\delta]\!]_n$, present only for $n=m$.
\item[(vii)] Line $(\ell 5)$, the rim totals:
\begin{equation*}
\tfrac13K_{rim}[\![\delta]\!]_n+\tfrac13\sup\theta^{\dagger}_L\cdot\mathsf{k}^{+,(<m)}_n .
\end{equation*}
\item[(viii)] Line $(\ell 6)$, the members $\mathfrak{v}^{\mathbf{C}}$ of the field-independent remainder
terms: $K^{\mathbf{C}}_D[\![\delta]\!]_n$.
\item[(ix)] Line $(\ell 7)$, the kept quadratic form $\mathfrak{q}_C$:
$K^{\mathfrak{q}}_D[\![\delta]\!]^{c}_n$, present only for $n=h+1$, where $h$ is the index in the
bound $\|\delta\mathsf{K}\|_{\delta}\le C_{cl}\vartheta_w^{h+1}$ for the operator difference
$\delta\mathsf{K}$ supplied by the proposition on the classical operator rows (see the proof of line
$(\ell 7)$).
\end{enumerate}
\end{proposition}

\begin{proof}
We treat the seven lines of \eqref{9.5:eq-production-bound-a} in turn.

\emph{Lines $(\ell 1)$, $(\ell 2)$: family marks.} The proof of the one-run bound on the
boundary-term stores, and the plaquette form of that bound, consist, type by type, of a size constant
multiplied by a tree count, or by the fact that the sum over labels equals the norm. The size
constants are: $K(d)B_0$ for each of $\mathbf{B}^{(j)}$ and $\mathbf{B}'^{(j)}$; $C^{\sharp}_0$ per
stage and per chain for the chain terms, with at most $m-j+1$ chains; and $b_\delta$ for the stores
of type $(\delta)$. These are per-cube activity sums of the kind that part (b) of the zone-crossing
lemma bounds for flat marks, so that lemma applies to the class-$n$ members as stated. We repeat the
estimates in order to display the constants.

(a) Boundary terms born at the operation $\mathbf{T}$, class $n\le j-1$. By the crossing
estimate of part (b) of the zone-crossing lemma, with $\mathsf{w}_{j,n}$ the crossing weight of that
lemma,
\begin{equation*}
\mathfrak{N}[G^{(n)}]\le \mathsf{S}^{+,(n)}_j[\mathbf{B}]\,e^{-\kappa_{\mathbf{B}}d_j(X)}
\le \mathsf{w}_{j,n}\,\mathsf{S}^{+,(n)}_j[\mathbf{B}]\,e^{-(\kappa_{\mathbf{B}}-1)d_j(X)} .
\end{equation*}
For $j<m$ the weight $e^{4d_m(K_F)}e^{\varsigma_0d_j}$ is at most
$e^{((8/15)+(1/5))\kappa d_j}$, as shown in the one-run bound on the boundary-term stores. Since
$1-\tfrac{8}{15}-\tfrac15=\tfrac{4}{15}$, the rate left is $\tfrac{4}{15}\kappa-1$, and
$\tfrac{4}{15}\kappa-1\ge c_d+1$ because $\kappa\ge\kappa_d\ge 20(c_d+2)$, which is the floor on
$\kappa$ used in that one-run bound; indeed $20(c_d+2)\ge\tfrac{15}{4}(c_d+2)$. No new floor on
$\kappa$ is needed. By \cite[(1.26)]{B-CMP116} the sum of these members over the regions $X$ is then
at most $K(d)\mathsf{w}_{j,n}\mathsf{S}^{+,(n)}_j[\mathbf{B}]$. For $j=m$ the members carry the decay
rate $\kappa_{out}\ge 2c^{\sharp}\kappa$ against the weight rate $\tfrac85\kappa$, by the third
property of the zone split of the boundary terms.

(b) Boundary terms born at the operation $\mathbf{R}$. In the same way their contribution is
at most
\begin{equation*}
K(d)\bigl\{\mathsf{S}^{+,(n)}_j[\mathbf{B}'^{res}]
+\mathbf{1}[j=n]\,\mathsf{S}^{+}_j[\mathbf{B}'^{own}]\bigr\}.
\end{equation*}
Each of the quantities bounded in (a) and (b) is one summand of
$\mathsf{y}^{(\le m)}_n$, respectively of $\mathsf{y}'_n$ (the weight $w^{\mathbf{B}}_{n,n}$ equals
$1$). Hence their total is at most
\begin{equation*}
K(d)(\mathsf{y}_n+\mathsf{y}'_n)=\frac{\mathsf{y}_n+\mathsf{y}'_n}{2B_0}\cdot 2K(d)B_0 ,
\end{equation*}
while the one-run per-cube budget satisfies $\mathsf{b}_m(s)\ge 2K(d)B_0$ by the one-run bound on the
boundary-term stores. This gives item (i). On inner strips the same bounds hold with $\mathsf{S}^{\theta}$
and $d^{\theta}$ in place of $\mathsf{S}$ and $d_j$, by \eqref{9.5:eq-distance-units-c} and the
corollary for inner strips contained in input (c) of the uniqueness theorem.

(c) Chain terms of a chain run at a step $\mathsf{b}\le m$. By the chain-term comparison theorem
and the value bound of the theorem on resolved extensions of stores, in the norm for which the sum over
labels equals the norm,
\begin{equation*}
\sum_{\nu}\sum_{\mathsf{i}}\mathsf{s}^{\natural,(z,\mathsf{i})}_{\nu}(\mathsf{b})
\le K^{fr}_{\mathbb{C}}[\![\delta]\!]^{\natural+,(\mathsf{b})}_n
+\tfrac43\,\theta_P(\check{\mathsf{T}}_{\mathsf{b}})\,\mathsf{k}^{+,(<\mathsf{b})}_n ;
\end{equation*}
here the inequality uses part (ii) of the lemma on resolved terms and the kept-member load statement,
together with the line of the kept members in part (e) of the kept-member bound at the site of the
stores (induction hypothesis), and, for the terms of $\mathfrak{F}^{B4}$ other than $\mathsf{q}_W$,
part (c) of the bound on the terms of interrupted runs. The one-run bound charges
$N(g_{\mathsf{b}})C^{\sharp}_0$ to the same sum, where $N(g_{\mathsf{b}})$ is the number of stages of
Definition~\ref{9.1:not-levelwise-rate}, since it allows $C^{\sharp}_0$ per stage. As
$N(g_{\mathsf{b}})\ge1$, the ratio of the two is at most
\begin{equation*}
\frac{K^{fr}_{\mathbb{C}}[\![\delta]\!]_n}{C^{\sharp}_0}
+\sup_{\mathsf{b}\le m,\ \mathsf{T}}\frac{4\theta_P(\mathsf{T})}{3N(g_{\mathsf{b}})C^{\sharp}_0}\,
\mathsf{k}^{+,(<m)}_n .
\end{equation*}

(d) Stores of type $(\delta)$, which are case $(\beta)$ of the classification of stores in the
one-run bound on the boundary-term stores. These stores satisfy the value bound of the theorem on
resolved extensions of stores at the constant $K^{fr}_{\mathbb{C}}$, by part (d) of the bound on the
terms of interrupted runs. Summed by the tree count against $2K(d)B_0$, they contribute
$K^{fr}_{\mathbb{C}}[\![\delta]\!]_n/(2K(d)B_0)$. Together with (c) this is item (ii).

\emph{Reading marks.} The units $v$ with $n^c_v=n$ form a subfamily of the units summed in the one-run
lines. Those lines hold with $\mathsf{d}=1$ for $\mathfrak{N}(v)e^{4d_m(K_v)}$, by the one-run bound on
the boundary-term stores and its plaquette form. By Lemma~\ref{9.5:lem-reading-mark-sums} the reading
marks of these units carry the common factor
$\mathsf{c}_{rd}(\vartheta^{\sharp}_v)^{1/2}=\mathsf{c}_{rd}\vartheta_f^{n/2}$, and by the bound
$(\hat{\vartheta})$ of \eqref{9.5:eq-constituent-rates-a} this is at most
$2\mathsf{c}_{rd}C_{\hat{\vartheta}}[\![\delta]\!]^{c}_n/E_0$. This is item (iii).

\emph{Gaussians.} For a vertex pair family $e$, the unit $v_e$ is a localised piece of
$-\tfrac12\langle x,\mathcal{A}_ex\rangle$, where $\mathcal{A}_e$ (a letter used only in this proof) is
a compression of the operator
$C^{*}\Delta^{(j)}C$, by the bound on the terms of interrupted runs. At common data,
$|\delta v_e|$ is at most $K_1\|\delta\mathcal{A}_e\|_{\delta}$ times the same supremum of $|x|^2$
and the same bond count that $\mathfrak{N}^G(e)$ contains. Let $c_{\mathcal{A}}$ be the lower constant
by which $\mathfrak{N}^G$ is normalised. By part $(\beta)$ of the complexification corollary for
operator differences, on $\mathfrak{F}^V$
\begin{equation*}
|\delta v_e|\le K_1\|\delta\mathcal{A}_e\|_{\delta}\cdot\frac{\mathfrak{N}^G(e)}{c_{\mathcal{A}}}
\le \frac{K_1}{c_{\mathcal{A}}}\bigl(2C_{cl}(1+\varpi_0)\vartheta_w^{j_W}\bigr)^{1/2}\mathfrak{N}^G(e).
\end{equation*}
We now use each run's own bound for the Gaussian units and the inequality between the minimum and
the mean of two numbers used in the proof of Lemma~\ref{9.5:lem-reading-mark-sums}. This gives the
weight
\begin{equation*}
\mathsf{c}_G\vartheta_w^{j_W/4}\mathfrak{N}^G(e)\times
\begin{cases}1,\\ e^{-\frac14\delta_P\hat{D}},\end{cases}
\qquad
\mathsf{c}_G:=2\Bigl[\frac{K_1}{c_{\mathcal{A}}}\bigl(2C_{cl}(1+\varpi_0)\bigr)^{1/2}
C^{+}_{\mathfrak{K}}\Bigr]^{1/2},
\end{equation*}
with the two cases as in Lemma~\ref{9.5:lem-reading-mark-sums}. These weights enter the line of case
$(\beta)$ of the classification of stores at the size constant $b_\delta$. By $(\vartheta_w)$ in
\eqref{9.5:eq-constituent-rates-a}, $\mathsf{c}_G\vartheta_w^{n/4}\le
2\mathsf{c}_GC_{\hat{\vartheta}}[\![\delta]\!]^{c}_n/E_0$. This is item (iv). We put
\begin{equation}\label{9.5:eq-constituent-bounds-3}
K^{B}_D:=\frac{1}{2B_0}+K^{fr}_{\mathbb{C}}\Bigl(\frac{1}{C^{\sharp}_0}+\frac{1}{2K(d)B_0}\Bigr)
+\frac{2(\mathsf{c}_{rd}+\mathsf{c}_G)C_{\hat{\vartheta}}}{E_0}.
\end{equation}

\emph{Line $(\ell 3)$.} A pointed unit anchored in a cell $\Delta'$ of level $n$ is assigned to class
$n$.

For the family marks we distinguish three kinds of entries. For singles, tails and weights, the
difference $[\hat{F}^{\mathsf{B}}-\hat{F}^{\mathsf{A}}]$ has norm $\mathsf{S}^{+}_j[F]$ and the same
form as $F$. For groups, the difference is
$\mathfrak{i}_y[G^{+}]-b^{\mathsf{B}}\mathfrak{i}_y[G^{+}_{\hat{A}}]$. By the group corollary for
resolved extensions, with a constant $C$ of that corollary,
\begin{equation*}
\|G^{+}\|^c\le C^{+}_X\bigl(\mathsf{S}^{+}_j+C\hat{\vartheta}_j\mathfrak{n}\bigr),\qquad
\|G^{+}_{\hat{A}}\|^c\le CC^{+}_X\hat{\vartheta}_j .
\end{equation*}
For the counterterm and weight entries, the lemma on the counterterm and weight entries of differences
replaces $(b_{+}+\tfrac23r,\tfrac83r)$ by
$(\mathsf{b}_j+\hat{\vartheta}_jb_{+},\mathsf{z}+\hat{\vartheta}_jr)$. That lemma sets
\begin{equation*}
\mathfrak{N}(\delta_u):=8C_X\bigl(\mathsf{S}_j[F]/\mathfrak{n}+\hat{\vartheta}_j\bigr)\mathfrak{N}_J(v).
\end{equation*}
Consequently $\mathfrak{N}_v[G^{(n)}_v]$ obeys the flat column sums of
Lemma~\ref{9.5:lem-pointed-sums} with the size constants of the rows replaced by the numbers
$e_j$, $r_j$, $\beta_j$ below:
\begin{equation*}
\begin{gathered}
e_j:=8C^{+}_X\bigl(\mathsf{S}^{+}_j[E]+E_0\hat{\vartheta}_j\bigr)
\quad\text{(entries $1$, $2$, $\partial$ and the tails)},\\
r_j:=8C^{+}_X\bigl(\mathsf{S}^{+}_j[R]+g_j^{\kappa_0}\hat{\vartheta}_j\bigr),\qquad
\beta_j:=\mathsf{b}_j+\mathsf{z}^{\sharp}_n+\hat{\vartheta}_j(b_{+}+r).
\end{gathered}
\end{equation*}

The reading marks are assigned to classes by the rule recorded in \eqref{9.5:eq-resolved-terms-a}.
They carry the weight
\begin{equation*}
\mathsf{c}_{rd}\vartheta_f^{(j\wedge n_{X_v})/2}\mathfrak{N}^{1/2}_v
\le\mathsf{c}_{rd}C_{\hat{\vartheta}}\hat{\vartheta}^{\natural}_{j\wedge n_{X_v}}\mathfrak{N}^{1/2}_v .
\end{equation*}
In the second case of that rule, the sum is the same table at $\hat{\alpha}/2$ with the prefactors
$E_0\hat{\vartheta}^{\natural}_j$ and $g_j^{\kappa_0}\hat{\vartheta}^{\natural}_j$. These are the
entries of $[\![\delta]\!]^{c}_n$, in the reading of $(\ell\text{-}\varsigma)$ fixed in
\eqref{9.5:eq-constituent-rates-a}. In the third case, the units form a subfamily of the one-run line
with the common factor
\begin{equation*}
\mathsf{c}_{rd}C_{\hat{\vartheta}}\hat{\vartheta}^{\natural}_{n_X}
\le 2\mathsf{c}_{rd}C_{\hat{\vartheta}}[\![\delta]\!]^{c}_{n_X}/E_0 .
\end{equation*}
By the walk that defines $K_{pt}$ in Lemma~\ref{9.5:lem-pointed-sums}, the sum of all these marks is
at most
\begin{equation*}
K_{pt}(\mathsf{w}^{\vee})^3\,C_{br}(1+\mathsf{c}_{rd}C_{\hat{\vartheta}})[\![\delta]\!]_n
+2\mathsf{c}_{rd}C_{\hat{\vartheta}}\mathsf{K}^{\flat,1/2}(\mathsf{w}^{\vee})^3
[\![\delta]\!]^{c}_n/E_0 ,
\end{equation*}
where
\begin{equation}\label{9.5:eq-constituent-bounds-4}
C_{br}:=\max\bigl\{12C^{+}_X,\ S_{\mathbf{B}},\ 2(b_{+}+r)/B_0\bigr\}.
\end{equation}

To justify the use of $C_{br}$, we check that the bracket of the claim of
Lemma~\ref{9.5:lem-pointed-sums}, at the prefactors above, is at most
$C_{br}([\![\delta]\!]^{(5)}_n+[\![\delta]\!]^{c}_n)$. We compare entry by entry with the entries that
define $[\![\delta]\!]^{(5)}_n$ and $[\![\delta]\!]^{c}_n$; below $\varsigma$ and $w$ denote the
weights that these entries carry.
\begin{itemize}
\item The entries $8C^{+}_X\varsigma\mathsf{S}^{+}_j[E]$ and $8C^{+}_Xw\mathsf{S}^{+}_j[R]$ are
$8C^{+}_X$ times entries of $[\![\delta]\!]^{(5)}_n$.
\item The entries $8C^{+}_X\varsigma E_0\hat{\vartheta}_j$ and
$8C^{+}_Xwg_j^{\kappa_0}\hat{\vartheta}_j$ are handled with $\hat{\vartheta}\le\sqrt2\,
\hat{\vartheta}^{\natural}$. They become at most $8\sqrt2\,C^{+}_X\le 12C^{+}_X$ times entries of
$[\![\delta]\!]^{c}_n$.
\item The sum over $j$ of the counterterm entries obeys
$\mathfrak{g}_n\sum_jL^{-(n-j)/2}\mathsf{z}^{\sharp}_n\le S_{\mathbf{B}}\mathfrak{g}_n\mathsf{z}^{\sharp}_n$.
\item The entry $\mathfrak{g}_nL^{-(n-j)/2}\mathsf{b}_j$ appears in $[\![\delta]\!]^{(5)}_n$ as it
stands.
\item The last entry obeys, since $\mathfrak{g}_n\le1$,
\begin{equation*}
\mathfrak{g}_n(b_{+}+r)L^{-(n-j)/2}\hat{\vartheta}_j
\le\frac{\sqrt2(b_{+}+r)}{B_0}\,w^{\mathbf{B}}_{n,j}B_0\hat{\vartheta}^{\natural}_j .
\end{equation*}
\end{itemize}
Dividing by the one-run budget, we obtain item (v) with
\begin{equation}\label{9.5:eq-constituent-bounds-5}
K^{pt}_D:=\frac{C_{br}(1+\mathsf{c}_{rd}C_{\hat{\vartheta}})K_{pt}}{\mathsf{K}^{\flat,1/2}}
+\frac{2\mathsf{c}_{rd}C_{\hat{\vartheta}}}{E_0}.
\end{equation}
This is the per-cell bound of the table of direct uses on the rim for the dialled parts in zone $n$,
namely at most $C_{rim}C^{\sharp}_L|\lambda_z|[\![\delta]\!]^{\natural+}_n$. Here it holds with
the additional weights $e^{4d_m(K_v)}e^{\varsigma_1\mathsf{r}(v)}$ of the flat column sums, which are
paid by the units' own decay rates inside $K_{pt}$.

\emph{Line $(\ell 4)$.} By the proof of part (ii) of the copy lemma, the copies are controlled by the
flat column sums at a level-$m$ cell with $\mathsf{w}^{\vee}=1$, together with the level $j=m$. All
copies are anchored at level-$m$ cells of the new set, so they belong to class $m$. The argument of
line $(\ell 3)$ therefore gives item (vi) with
\begin{equation}\label{9.5:eq-constituent-bounds-6}
K^{cp}_D:=\frac{C_{br}(1+\mathsf{c}_{rd}C_{\hat{\vartheta}})K_{pt}+16M^dK(d)C^{+}_X}{\mathsf{K}^{cp}}
+\frac{2\mathsf{c}_{rd}C_{\hat{\vartheta}}}{E_0}.
\end{equation}

\emph{Line $(\ell 5)$.} For $Y_4\in\mathcal{Y}_{rim}$ the rim total is a finite sum of pieces assigned
to classes and of values, by the rim decomposition lemma. Its class-$z$ mark is the sum of the marks of
the slots, by the rim decomposition lemma and by the corollary on strip bounds for differences and
marks.

For the chain-term part of the rim total, part (iv) of the rim-total lemma bounds the contribution by
$C_A\mathsf{w}_L^{-1}\cdot 8K_{\mathbb{C}}[\![\delta]\!]_n$ times the geometric factors of the rim
estimate, with $\mathsf{w}_L$ the rim weight of that lemma, and the values by
$2K_{\mathbb{C}}[\![\delta]\!]_n$ per cell. For the direct uses, entries
$1$--$3$ and $6$ of the table of direct uses on the rim give at most
$C_{rim}C^{\sharp}_L[\![\delta]\!]_n$ per cell, with brackets that carry the factor
$C_A\mathsf{w}_L^{-1}$.

The bound in the norm $\|\cdot\|_a$ follows directly from the following three statements. By part (c2)
of the table of direct uses on the rim, entries
$5$, $7$ and $8$, the rim potential $V^{\langle\lambda\rangle}(Y)$, which is affine in the dial
variable $\lambda$, satisfies on the dial the weighted form of the first sharpened rim condition, with
radius $(C_{\sharp}+C''_L/E_0)\alpha$. By part (i) of the corollary on strip bounds for differences
and marks, the strip bounds hold for differences and marks measured in $\|\delta\mathfrak{v}\|_a$ and
$\|W\|_a$. By part (ii) of that corollary, $\mathcal{B}^{\sharp}_1$ and $\mathcal{B}^{\sharp}_2$ are
convex balls for the supremum norm.

Now let $\varepsilon\mapsto\mathfrak{v}^{\lambda}+\varepsilon W^{(z)}$ be an affine function that
stays in a ball of radius $\mathsf{R}$ for $|\varepsilon|\le\Lambda_z$. Both
$\mathfrak{v}^{\lambda}$ and $\mathfrak{v}^{\lambda}+\Lambda_zW^{(z)}$ then lie in that ball, so
$\Lambda_z\|W^{(z)}\|_a\le 2\mathsf{R}$, and
\begin{equation*}
\|W^{(z)}\|_a\le\frac{2\mathsf{R}}{\Lambda_z}=\frac{2\mathsf{R}\,[\![\delta]\!]_{n(z)}}{E^z_1}.
\end{equation*}
With $\mathsf{R}=\mathsf{R}^{\sharp}$ this is $K_{rim}[\![\delta]\!]_{n(z)}$, by
\eqref{9.5:eq-constituent-bounds-1}. The kept part is bounded by part $(\delta)$ of the kept-member
bound at the rim sites, which is the induction hypothesis (2). The budget against which this line is
normalised is the sum of $3$, $\mathsf{K}^{0}$ and the charge $C^{\sharp}_0$ per stage for the chain
terms; it is at least $3$. This gives item (vii), and we put
\begin{equation}\label{9.5:eq-constituent-bounds-7}
K^{rim}_D:=K_{rim}/3 .
\end{equation}

\emph{Line $(\ell 6)$.} The members $\mathfrak{v}^{\mathbf{C}}$ of the field-independent remainder terms
localised in $C^{\sharp}$, split by the lowest zone they meet, are bounded by
Proposition~\ref{9.5:prop-remainder-difference}, whose proof is outstanding: it gives
\begin{equation*}
\sum_{Y_4\subset C^{\sharp}}\sup\bigl|\delta\mathfrak{v}^{\mathbf{C},(n)}(Y_4)\bigr|
\le C^{\times}_{\mathbf{C}}\bigl[\mathfrak{g}_n\mathsf{z}^{\sharp}_n
+E_0(\log x_n)^{c_2}\vartheta_w^n\bigr].
\end{equation*}
By $(\mathrm{cl})$ in \eqref{9.5:eq-constituent-rates-a} and the identity
$\vartheta_w^n=\vartheta_2^{n/4}\vartheta_{\times}^n$, this right member is at most
$C^{\times}_{\mathbf{C}}(1+1/C^{\natural}_{cl})[\![\delta]\!]^{\natural+}_n$. This gives item (viii)
with
\begin{equation}\label{9.5:eq-constituent-bounds-9}
K^{\mathbf{C}}_D:=\frac{C^{\times}_{\mathbf{C}}(1+1/C^{\natural}_{cl})}{\mathsf{c}_{\mathbf{C}}},
\end{equation}
where $\mathsf{c}_{\mathbf{C}}$ is the normalising constant of line $(\ell 6)$.

\emph{Line $(\ell 7)$.} By the proposition on the classical operator rows,
$\|\delta\mathsf{K}\|_{\delta}\le C_{cl}\vartheta_w^{h+1}$ at real common parameter points. By the
complexification corollary for operator differences it follows that
$\|\delta\mathsf{K}\|_{\delta}\le C_{\mathsf{K}}\vartheta_w^{(h+1)/4}$ on $\mathfrak{F}^V$. By the
Schur bound and the shell count,
\begin{equation*}
|\delta\mathfrak{q}_C|
\le\tfrac12K_1\|\delta\mathsf{K}\|_{\delta}\sup|B|^2\,\#(\mathbf{B}_0\cap C)
\le\tfrac12K_1\|\delta\mathsf{K}\|_{\delta}\,\mathfrak{b}'^{\,2}_m\cdot dM^d\mathsf{n}_{sh}(m).
\end{equation*}
This is to be compared with
$\mathfrak{q}_{max}=15c_{tr}B_{\sharp}^2\mathfrak{b}'^{\,2}_mM^d\mathsf{n}_{sh}(m)$. The ratio is at
most $dK_1C_{\mathsf{K}}\vartheta_w^{(h+1)/4}/(30c_{tr}B_{\sharp}^2)$. By $(\vartheta_w)$ in
\eqref{9.5:eq-constituent-rates-a}, $\vartheta_w^{(h+1)/4}\le
2C_{\hat{\vartheta}}[\![\delta]\!]^{c}_{h+1}/E_0$. This gives item (ix), for $n=h+1$, with
\begin{equation}\label{9.5:eq-constituent-bounds-8}
K^{\mathfrak{q}}_D:=\frac{dK_1C_{\mathsf{K}}C_{\hat{\vartheta}}}{15c_{tr}B_{\sharp}^2E_0}.
\end{equation}

The six constants of the seven lines are $K^{B}_D$, $K^{pt}_D$, $K^{cp}_D$, $K^{rim}_D$,
$K^{\mathbf{C}}_D$ and $K^{\mathfrak{q}}_D$, with $K^{B}_D$ serving lines $(\ell 1)$ and $(\ell 2)$
together; they are defined in \eqref{9.5:eq-constituent-bounds-3},
\eqref{9.5:eq-constituent-bounds-5}, \eqref{9.5:eq-constituent-bounds-6},
\eqref{9.5:eq-constituent-bounds-7}, \eqref{9.5:eq-constituent-bounds-9} and
\eqref{9.5:eq-constituent-bounds-8}, and $K_D$ is their maximum, as in
\eqref{9.5:eq-constituent-bounds-2}. The coefficients of $\mathsf{k}^{+,(<m)}_n$ in items (ii) and
(vii) are at most $\bar{\Theta}$, by the definition of $\bar{\Theta}$ in
\eqref{9.5:eq-constituent-rates-a}.

In summary, every two-run size that enters the comparison is bounded, class by class, by the
hypotheses (1)--(5) as follows. The bounds for the boundary terms, the values, the
copies and the kept quadratic form come from inputs already on the list of the uniqueness theorem. The
Gaussian and operator parts come from the proposition on the classical operator rows, under the
operator-difference bounds of input (e) of the uniqueness theorem. The
field-independent remainder terms are bounded only type by type, by
Proposition~\ref{9.5:prop-remainder-difference}, whose proof is outstanding.
\end{proof}

\begin{definition}[Kept factors as leaves and the jointly invariant path]\label{9.5:not-kept-leaves}
Let $Y$ be a large-field region created at step $m$ and thrown back, with strip $Y^{\sharp}$, enlarged strip
$(Y^{\sharp})^{+}$ and width $\check{R}_m$. Let $\mathfrak{K}(Y)$ be its pair of kept factors of
\eqref{9.5:eq-kept-factors-a}: the kept potential $V^{\sharp}(Y^{\sharp})$ and the kept quadratic form
$\mathsf{Q}_Y$. We write $\Lambda_k\subset\Omega_k$ for Ba{\l}aban's small-field regions of level $k$
(Definition~\ref{1.5:def-densities}) and
$\mathrm{dist}_m$ for the distance measured in blocks of level $m$. Let $\xi=(\mathsf{v},c)$ be the common
datum, with collar entries $c^{\mathsf{A}}$, $c^{\mathsf{B}}$ of the two runs. By part (a) of the
exact-representation statement, the read of run $\mathsf{B}$ has the representation
\begin{equation}\label{9.5:eq-kept-leaves-3}
\mathrm{rd}^{\mathsf{B}}\mathbf{U}'=\mathsf{w}\,e^{i\mathfrak{d}}\,\mathsf{v}\,\mathsf{w}^{-1},
\end{equation}
where $\mathbf{U}'$ is the configuration read by run $\mathsf{B}$ and $\mathfrak{d}$ is the phase increment;
we write $\mathsf{h}:=\mathsf{w}$ for its centre map. We use the following statements.
\begin{enumerate}
\item[(1)] On $Y^{\sharp}$ the new determining set has level $m$, as stated in \cite[(1.33)]{B-CMP122b}.
\item[(2)] (The strip-geometry statement, part (c).) At every read at which $Y$ is live, with
$\Omega_{m+1}$ the current region,
\begin{equation}\label{9.5:eq-kept-leaves-1}
\mathrm{dist}_m\bigl(Y^{\sharp},\Omega_{m+1}\bigr)\ \ge\ \tfrac52\,\check{R}_m .
\end{equation}
Moreover every cell at $\mathrm{dist}_m$-distance less than $\tfrac52\check{R}_m$ from $Y^{\sharp}$ has
level at most $m$, and this neighbourhood lies in the component of $\Lambda_i^{c}$ that contains $Y$, for
every $i>m$.
\item[(3)] (The determining-set lemma, part (i).) No step of the operation $\mathbf{T}$ re-blocks a
large-field region.
\item[(4)] (The exact-representation statement, part (d), and the centre-map lemma, part (iii).) Two bounds
on $|\log\mathsf{h}|$ for the centre map $\mathsf{h}=\mathsf{w}$ of \eqref{9.5:eq-kept-leaves-3}: the
bound $|\log\mathsf{w}|\le 9\,d\,y\,\rho^{l}\mathfrak{s}_u$ of part (d) of the exact-representation
statement, in the notation of that statement, and part (iii) of the centre-map lemma.
\item[(5)] (The mixed-rider lemma.) In the notation of that lemma, with its field letter written $\tilde{X}$
here to avoid a clash with the region $Y$, at one site and at the point $z$,
\begin{equation}\label{9.5:eq-kept-leaves-5}
\Bigl|\log\Bigl(\mathfrak{u}[e^{i\eta'X'}\mathbf{K}';\mathbf{K}';\mathfrak{S}'](z)\cdot
\mathfrak{u}[e^{i\eta\tilde{X}}\mathbb{K};\mathbb{K};\mathfrak{S}](z)^{-1}\Bigr)\Bigr|
\ \le\ C^{\times}_u\,L^{-\lambda(z)}\sup_{B^{top}(z)}\ell\,|X'| .
\end{equation}
The bound is relative: it controls the logarithm of the quotient of the two riders, and it carries its own
smallness index.
\item[(6)] (The rider Lipschitz lemma.) The Lipschitz bound for the passage
$(\tilde{X},\mathbb{K})\mapsto(\mathcal{K}^{\mathsf{A}},\mathbf{K}^{\mathsf{A}})$ from the pair of
arguments of the second rider in \eqref{9.5:eq-kept-leaves-5} to the pair from which the rider of run
$\mathsf{A}$ is built, whose increments are those of the field law.
\item[(7)] (The rider lemma, part (i).) The rider $\mathfrak{u}$ is a local analytic functional of the pair
$(\mathcal{K},\mathbf{K})$ of arguments from which it is built, on the top ball of the site, the ball
denoted $B^{top}(\cdot)$ as in \eqref{9.5:eq-kept-leaves-5}.
\end{enumerate}

Under these statements the following hold, and we fix the following objects.

\emph{Level.} The least level of the determining set met by $(Y^{\sharp})^{+}$ at any read at which $Y$ is
live is $m$. Consequently, for the terminal unit $(v,Y^{\sharp},m)$ one has $j\wedge n_X=m$, where $j$ is
the level of the unit, $n_X$ is the level index that the corollary on reads of terminal units in rooms
attaches to the domain $X$ of the unit, and $j\wedge n_X$ is their minimum; we put
\begin{equation}\label{9.5:eq-kept-leaves-a}
\vartheta^{\sharp}_Y:=\vartheta^{m/2}\ \text{in the runs},\qquad
\vartheta^{\sharp}_Y:=\vartheta_f^{m}\ \text{in the second case},
\end{equation}
the runs and the second case being the first and the second case of the corollary on reads of terminal
units in rooms.
In both cases $(\vartheta^{\sharp}_Y)^{1/2}\le\vartheta_f^{m/2}$; in the second case this is an equality,
and in the runs it reads $\vartheta^{m/4}\le\vartheta_f^{m/2}$.

\emph{Declaration.} The class of terminal units of the comparison (its leaves) is enlarged by the two
factors of $\mathfrak{K}(Y)$; this completes the data fixed in
\eqref{9.5:eq-two-point-oscillation-a}. For $v\in\mathfrak{K}(Y)$ we take the data domain
$\mathfrak{D}_v$ and the
room $\varrho_v$ of Notation~\ref{9.5:not-two-point-oscillation}, and we put
$\|v\|:=\mathfrak{O}^{+}(\hat{v}^{\sharp})$. Here $\hat{v}^{\sharp}$ denotes $v$ read as a function of its
datum on $Y^{\sharp}$.

For $v=V^{\sharp}(Y^{\sharp})$ the containment statements for the reads apply as stated, and so do parts
(a), (b) and (d) of the corollary on reads of terminal units in rooms. Part (d) includes the terms $\Phi$
of the comparison born in the operation $\mathbf{R}$ and the left-over terms. At a site other than its own,
$Y$ is not contained in the treated region $Z$. At its own site the lemma on reads of terminal units meeting
$Z$ applies.

\emph{Path.} The path along which a jointly invariant factor is compared keeps $\mathsf{h}$ and the
mismatch $\mathfrak{u}^{\mathsf{B}}(\tau)\,\mathfrak{u}^{\mathsf{A}}(\tau)^{-1}$ of the riders of the site,
$\tau$ being the argument at which both riders are evaluated;
with $\mathfrak{d}^{net}$ the net phase increment, it is
\begin{equation}\label{9.5:eq-kept-leaves-2}
\begin{split}
(\mathsf{v}_s,c_s)&:=\Bigl(\mathsf{h}^{s}e^{is\mathfrak{d}^{net}}(\mathsf{v}\oplus c^{\mathsf{B}})
\mathsf{h}^{-s},\ e^{sZ_{\delta}}c^{\mathsf{A}}\Bigr),\\
e^{Z_{\delta}}&:=\mathfrak{u}^{\mathsf{B}}\,(\mathfrak{u}^{\mathsf{A}})^{-1}.
\end{split}
\end{equation}

\emph{What is read rather than applied as stated.}
\begin{enumerate}
\item[(i)] No previously stated containment statement covers $\mathsf{Q}_Y$; its reads are checked
directly. The
form $\mathsf{Q}_Y$ depends on the slice $\zeta$ only through $\mathsf{V}$, the $m$-fold block average of
$\mathbf{U}$ (Definition~\ref{1.5:def-block-average}) restricted to $\mathsf{C}\setminus\Lambda$, where
$\mathsf{C}$ is the core of the component at the healing site and $\Lambda$ the small-field region of the
read; so it is a function of the hull datum itself.
\item[(ii)] In the norm, the weight $\mathfrak{n}e^{-\kappa_T d_j(X)}$ of the corollary on reads of
terminal units in rooms, with $\kappa_T$ the corollary's decay rate, is replaced by $\mathfrak{O}^{+}$.
\end{enumerate}

\emph{Smallness of the path.} At one site, and for that site's single mismatch, $|\log\mathsf{h}|$ and
$|Z_{\delta}|$ are small in the smallness index, measured in units of the margin of the datum. The room
used is the margin of the free-room statement, part (i), and of the small-field-region theorem.
\end{definition}

\begin{proof}
\emph{Level.} By statement (1), on $Y^{\sharp}$ the new determining set has level $m$. At the birth of $Y$
one has
\begin{equation*}
(Y^{\sharp})^{+}\setminus Y^{\sharp}\ \subset\ \Lambda_m\ \subset\ \Omega_m .
\end{equation*}
For later reads we use statement (2). By \eqref{9.5:eq-kept-leaves-1}, the strip keeps distance at least
$\tfrac52\check{R}_m$ from the current $\Omega_{m+1}$. Every cell at $\mathrm{dist}_m$-distance less than
$\tfrac52\check{R}_m$ from $Y^{\sharp}$ has level at most $m$, and this neighbourhood lies in the
component of $\Lambda_i^{c}$ containing $Y$ for every $i>m$. That neighbourhood is therefore part of the
large-field region, and by statement (3) no step of $\mathbf{T}$ re-blocks it. Hence, at every read at
which $Y$ is live, the least level of the determining set met by $(Y^{\sharp})^{+}$ is $m$.

For the terminal unit $(v,Y^{\sharp},m)$ this gives $j\wedge n_X=m$. The choice
\eqref{9.5:eq-kept-leaves-a} is therefore the index of the corollary on reads of terminal units in rooms,
evaluated at $j\wedge n_X=m$.

\emph{Item (ii).} By the lemma on reads of terminal units meeting $Z$, part (b) of that lemma holds at the
own site in the form
\begin{equation}\label{9.5:eq-kept-leaves-4}
\begin{split}
&\bigl|\hat{F}^{\mathsf{B}}(\mathrm{rd}^{\mathsf{B}}\mathbf{U}')
-\hat{F}^{\mathsf{B}}(\mathsf{v}\oplus c^{\mathsf{B}})\bigr|\\
&\qquad\le\ 16\,C^{\mathfrak{K}}_{TL}\,\vartheta^{\mathfrak{K}}\sup_{\mathfrak{T}^{\mathsf{B}}[1/4]}
\bigl|\hat{F}^{\mathsf{B}}\bigr| ,
\end{split}
\end{equation}
with $\vartheta^{\mathfrak{K}}:=\vartheta^{n_X/2}$ and $C^{\mathfrak{K}}_{TL}$ the constant of that lemma.
This bound comes from the Schwarz lemma applied to
$s\mapsto\hat{F}^{\mathsf{B}}(\mathsf{v}_s)-\hat{F}^{\mathsf{B}}(\mathsf{v}_0)$. Because the right-hand
side of \eqref{9.5:eq-kept-leaves-4} is a supremum over the tube $\mathfrak{T}^{\mathsf{B}}[1/4]$,
replacing the weight $\mathfrak{n}e^{-\kappa_T d_j(X)}$ by $\mathfrak{O}^{+}$ leaves it unchanged in form.
That this factor belongs to the class of holomorphic functionals of the tube data used for all terminal
units is part (a) of the holomorphy statement.

\emph{Path.} For a factor invariant jointly in the two runs, the reduction used for the other terminal
units is not available. For this reason \eqref{9.5:eq-kept-leaves-3} is used with its centre map, and the
mismatch of the riders is kept, which gives \eqref{9.5:eq-kept-leaves-2}.

\emph{Smallness.} The quantity $|\log\mathsf{h}|$, with $\mathsf{h}=\mathsf{w}$, is bounded by the two
bounds of statement (4).

By the definition in \eqref{9.5:eq-kept-leaves-2}, $Z_{\delta}$ is a logarithm of the quotient
$\mathfrak{u}^{\mathsf{B}}(\mathfrak{u}^{\mathsf{A}})^{-1}$ of the two riders of the site. By statement (7)
each rider is a local analytic functional of its pair of arguments on the top ball of the site. Statement
(5) bounds the logarithm of a quotient of two riders directly, relative and with its own smallness index,
and statement (6) controls the passage to the pair $(\mathcal{K}^{\mathsf{A}},\mathbf{K}^{\mathsf{A}})$ of
run $\mathsf{A}$, whose increments are those of the field law; these two statements together control
$|Z_{\delta}|$.

Both quantities are thus small in the index, in units of the margin of the datum. Since there is one
mismatch per site, the statement is made per site.
\end{proof}

\begin{lemma}[Site-independent bound of a piece by the oscillation]\label{9.5:lem-site-bound}
Let $f$ be a function on the set $\mathfrak{X}^{\sharp}_Y$ of common data, and let $f$ be read through a map
$\tau\mapsto\eta(\tau)\in\mathfrak{X}^{\sharp}_Y$, where $\tau=(\tau_i)_{i\in J}$ and $J$ is a finite index
set. Assume that the reading $\tau\mapsto f(\eta(\tau))$ is holomorphic on a neighbourhood of the product
polydisc
\begin{equation*}
\prod_{i\in J}\{\,|\tau_i|\le R_i\,\},\qquad R_i>1 .
\end{equation*}
For $\emptyset\neq I\subset J$ put
\begin{equation}\label{9.5:eq-site-bound-1}
\mathfrak{p}_I[G]:=\prod_{i\in I}\int_0^1 d\tau_i\,\partial_{\tau_i}\,G\,\Big|_{\tau_j=0\ (j\notin I)} .
\end{equation}
Then
\begin{equation}\label{9.5:eq-site-bound-2}
\bigl|\mathfrak{p}_I[f\circ\eta]\bigr|\le\mathfrak{O}^{+}(f)\prod_{i\in I}(R_i-1)^{-1},
\end{equation}
where $\mathfrak{O}^{+}(f)$ is the two-point oscillation of
Notation~\ref{9.5:not-two-point-oscillation}, display~\eqref{9.5:eq-two-point-oscillation-b}.
\end{lemma}

\begin{proof}
Write $\partial_{\tau_I}:=\prod_{i\in I}\partial_{\tau_i}$ and
$G(\tau):=f(\eta(\tau))-f(\eta(0))$. Since $I\neq\emptyset$, the operator $\partial_{\tau_I}$ contains at
least one derivative and therefore annihilates the constant $f(\eta(0))$; hence
$\partial_{\tau_I}(f\circ\eta)=\partial_{\tau_I}G$, and by \eqref{9.5:eq-site-bound-1}
$\mathfrak{p}_I[f\circ\eta]=\mathfrak{p}_I[G]$.

For every $\tau$ in the polydisc the two points $\eta(\tau)$ and $\eta(0)$ lie in
$\mathfrak{X}^{\sharp}_Y$, so the definition of the two-point oscillation in
\eqref{9.5:eq-two-point-oscillation-b}, applied to this pair of data, gives
\begin{equation*}
|G(\tau)|=|f(\eta(\tau))-f(\eta(0))|\le\mathfrak{O}^{+}(f)
\qquad\text{for all }\tau\in\prod_{i\in J}\{|\tau_i|\le R_i\}.
\end{equation*}

Now apply Cauchy's estimate. Fix $s\in[0,1]^I$ and let $\hat{s}\in\mathbb{C}^J$ have coordinates
$\hat{s}_i=s_i$ for $i\in I$ and $\hat{s}_j=0$ for $j\notin I$. If $|\sigma_i-s_i|=R_i-1$ for $i\in I$, then
$|\sigma_i|\le|s_i|+R_i-1\le R_i$; so the closed polydisc in the variables $(\tau_i)_{i\in I}$ with
centre $(s_i)_{i\in I}$ and radii $R_i-1$, with $\tau_j=0$ for $j\notin I$, lies in the product
polydisc, on a neighbourhood of which $G$ is holomorphic. Cauchy's integral formula in each variable
$\tau_i$, $i\in I$, on the circle of radius $R_i-1$ about $s_i$ gives
\begin{equation*}
\partial_{\tau_I}G(\hat{s})=\prod_{i\in I}\frac{1}{2\pi\mathrm{i}}\oint_{|\sigma_i-s_i|=R_i-1}
\frac{d\sigma_i}{(\sigma_i-s_i)^{2}}\;G(\sigma,0),
\end{equation*}
where $(\sigma,0)$ denotes the point with coordinates $\sigma_i$ for $i\in I$ and $0$ off $I$. Each circle has
length $2\pi(R_i-1)$ and on it $|\sigma_i-s_i|^{-2}=(R_i-1)^{-2}$; hence
\begin{equation*}
|\partial_{\tau_I}G(\hat{s})|\le\mathfrak{O}^{+}(f)\prod_{i\in I}(R_i-1)^{-1}.
\end{equation*}
Integrating this bound over $s\in[0,1]^I$, a set of volume one, gives
\begin{equation*}
|\mathfrak{p}_I[f\circ\eta]|=|\mathfrak{p}_I[G]|
\le\int_{[0,1]^I}|\partial_{\tau_I}G(\hat{s})|\,ds
\le\mathfrak{O}^{+}(f)\prod_{i\in I}(R_i-1)^{-1},
\end{equation*}
which is \eqref{9.5:eq-site-bound-2}.
\end{proof}

\begin{lemma}[The reading mark of a kept factor at later sites]\label{9.5:lem-kept-reading-mark}
Let $Y$ be a region created at step $m$, with the pair of kept factors $\mathfrak{K}(Y)$ and the
quantities fixed in Notation~\ref{9.5:not-two-point-oscillation} and
Notation~\ref{9.5:not-kept-leaves}. Assume the following:
\begin{enumerate}
\item[(a)] the three conditions of \eqref{9.5:eq-kept-leaves-a} on the kept factors
$\mathfrak{K}(Y)$, on the tube reads and on the phase increment $\mathsf{w}$;
\item[(b)] the base theorem on the kept potential in the tube, for run $\mathsf{B}$;
\item[(c)] the lemma on $\natural$-extensions;
\item[(d)] part (a) of the holomorphy statement for the read data on the one-machine polydisc;
\item[(e)] the mixing lemma for the mixed read data $\eta^{\times}$.
\end{enumerate}
Fix a vertex pair of a site that comes after the law site of step $m$ on an aligned label. Let
$\mathbf{U}$ and $\mathbf{U}'$ be the configurations of the runs $\mathsf{A}$ and $\mathsf{B}$
there, and put $\mathsf{v} := \mathrm{rd}^{\mathsf{A}}\mathbf{U}$. For a member $v$ of
$\mathfrak{K}(Y)$, with $\hat{v}^{\mathsf{B}\sharp}$ the member $v$ in run $\mathsf{B}$ written as
a function of its datum, $\Phi_Y$ the inner
variables of the kept factor, $c^{\mathsf{B}}_{col}$ the collar entries of run $\mathsf{B}$,
and $c^{\mathsf{A}}$, $c^{\mathsf{B}}$ the entries of the last argument of
$\hat{v}^{\mathsf{B}\sharp}$ in the runs $\mathsf{A}$ and $\mathsf{B}$, define the reading mark
$\mathsf{r}_Y$ by the first line of
\eqref{9.5:eq-kept-reading-mark-a}. Then $\mathsf{r}_Y$ is counted at the pair $(m,k')$, with
$m$ the step at which $Y$ is created and $k'$ the index of the site at which the read is made,
among
the classical terms (those carrying the superscript $c$ in
Lemma~\ref{9.5:lem-reading-mark-sums}), it obeys the second line of
\eqref{9.5:eq-kept-reading-mark-a}, and the two members of $\mathfrak{K}(Y)$ together obey the
third line:
\begin{equation}\label{9.5:eq-kept-reading-mark-a}
\begin{gathered}
\mathsf{r}_Y := \hat{v}^{\mathsf{B}\sharp}\bigl(\mathrm{rd}^{\mathsf{B}}\mathbf{U}';
\Phi_Y,c^{\mathsf{B}}\bigr)
-\hat{v}^{\mathsf{B}\sharp}\bigl(\mathsf{v}\oplus c^{\mathsf{B}}_{col};
\Phi_Y,c^{\mathsf{A}}\bigr),\\
|\mathsf{r}_Y| \le 8C^{\star}_{TL}\,\vartheta^{\sharp}_Y\,
\mathfrak{O}^{+}\bigl(\hat{v}^{\mathsf{B}\sharp}\bigr),\\
\sum_{v\in\mathfrak{K}(Y)}[\![\mathsf{r}_Y^{\natural}]\!]\ \le
\mathsf{k}^{rd}_m\,\check{R}_m^{p_{\Sigma}},\qquad
\mathsf{k}^{rd}_m := \bigl(32C^{\star}_{TL}\bigr)^{1/2}C_{\Sigma}\,\vartheta_f^{m/2},
\end{gathered}
\end{equation}
where the sum runs over the two members $v$ of $\mathfrak{K}(Y)$, each
summand being formed with $\hat{v}^{\mathsf{B}\sharp}$ for the member summed over, and
$[\![\mathsf{r}_Y^{\natural}]\!]$ is the $\natural$-norm, in the sense of the definition of
$\natural$-extensions, of the $\natural$-extension of $\mathsf{r}_Y$ given by the lemma on
$\natural$-extensions. Regarding $\mathsf{r}_Y$, as in the proof below, as a function of the
point $\tau$ of the one-machine polydisc of item (d), the
third line holds as well, with the same constant, when $[\![\mathsf{r}_Y^{\natural}]\!]$ is
replaced by the factor $(\beta\mathfrak{B})^{1/2}$ in the bound of
Lemma~\ref{9.5:lem-polydisc-vertices}\,(iii) applied to
$\tau\mapsto\mathsf{r}_Y(\tau)-\mathsf{r}_Y(0)$ (which, for non-empty $\mathsf{y}$, has the same
image under $\mathfrak{p}$ as $\mathsf{r}_Y$, by
Lemma~\ref{9.5:lem-polydisc-vertices}\,(i)), the prefactor
$(2/(R-1))^{|\mathsf{y}|/2}$ being kept aside.
\end{lemma}

\begin{proof}
\emph{The first bound.} Let $s\mapsto(\mathsf{v}_s,c_s)$ be the jointly invariant path of
Notation~\ref{9.5:not-kept-leaves}, which equals
$(\mathsf{v}\oplus c^{\mathsf{B}}_{col},c^{\mathsf{A}})$ at $s=0$ and
$(\mathrm{rd}^{\mathsf{B}}\mathbf{U}',c^{\mathsf{B}})$ at $s=1$, and put
$f(s):=\hat{v}^{\mathsf{B}\sharp}(\mathsf{v}_s;\Phi_Y,c_s)$, so that
$\mathsf{r}_Y=f(1)-f(0)$. By item (b), for run $\mathsf{B}$, $\hat{v}^{\mathsf{B}\sharp}$ is
holomorphic on the datum domain of run $\mathsf{B}$, and by item (a), that is by
\eqref{9.5:eq-kept-leaves-a}, the path lies in that domain for
$|s|<(8C^{\star}_{TL}\vartheta^{\sharp}_Y)^{-1}$; hence $f$ is
holomorphic on the disc $|s|<(8C^{\star}_{TL}\vartheta^{\sharp}_Y)^{-1}$, and for every $s$ in
this disc its value is taken at a point of the datum domain of run $\mathsf{B}$, at one and the
same value $\Phi_Y$ of the inner variables. Hence $f(s)-f(0)$ is a difference of values of
$\hat{v}^{\mathsf{B}\sharp}$ at two such points at one value of the inner
variables, and by the definition of the two-point oscillation
\eqref{9.5:eq-two-point-oscillation-b},
\begin{equation*}
\sup_{|s|<(8C^{\star}_{TL}\vartheta^{\sharp}_Y)^{-1}}|f(s)-f(0)|
\le \mathfrak{O}^{+}\bigl(\hat{v}^{\mathsf{B}\sharp}\bigr).
\end{equation*}
Lemma~\ref{9.5:lem-disc-contraction}, whose hypothesis $\rho>1$ reads here
$8C^{\star}_{TL}\vartheta^{\sharp}_Y<1$, applied to $f$ with
$\rho=(8C^{\star}_{TL}\vartheta^{\sharp}_Y)^{-1}$ and
$\mathfrak{B}=\mathfrak{O}^{+}(\hat{v}^{\mathsf{B}\sharp})$, gives
$|f(1)-f(0)|\le\mathfrak{B}/\rho$. This is the second line of
\eqref{9.5:eq-kept-reading-mark-a}.

\emph{The second bound.} Regard $\mathsf{r}_Y$ as a function of the point $\tau$ of the
one-machine polydisc, through the read data $\eta^{\mathsf{B}}(\tau)$ of run $\mathsf{B}$ and
the mixed read data $\eta^{\times}(\tau)$, the arguments $\Phi_Y$ and the entries
$c^{\mathsf{A}}$, $c^{\mathsf{B}}$ of the last argument being
suppressed from the notation. Thus
$\mathsf{r}_Y(\tau)=\hat{v}^{\mathsf{B}\sharp}(\eta^{\mathsf{B}}(\tau))
-\hat{v}^{\mathsf{B}\sharp}(\eta^{\times}(\tau))$. Put
\begin{equation*}
\begin{aligned}
H^{\mathsf{B}}&:=\hat{v}^{\mathsf{B}\sharp}(\eta^{\mathsf{B}}(\tau))
-\hat{v}^{\mathsf{B}\sharp}(\eta^{\mathsf{B}}(0)),\\
H^{\mathsf{A}}&:=\hat{v}^{\mathsf{B}\sharp}(\eta^{\times}(\tau))
-\hat{v}^{\mathsf{B}\sharp}(\eta^{\times}(0)).
\end{aligned}
\end{equation*}
With $g^2$ in the sense of the lemma on $\natural$-extensions (the
superscript is an index, not a power), regrouping the four terms gives
\begin{equation}\label{9.5:eq-kept-reading-mark-2}
g^2:=\mathsf{r}_Y(\tau)-\mathsf{r}_Y(0)=H^{\mathsf{B}}-H^{\mathsf{A}}.
\end{equation}
By part (a) of the holomorphy statement, $\eta^{\mathsf{B}}(\tau)$, and by the mixing lemma,
$\eta^{\times}(\tau)$, are points of the datum domain of run $\mathsf{B}$ for every $\tau$ in
the one-machine polydisc, complex $\tau$ included. Each $H^{\mathsf{m}}$ is therefore a
difference of values of $\hat{v}^{\mathsf{B}\sharp}$ at two points of the datum domain of run
$\mathsf{B}$ at one value of the inner variables, and on the whole polydisc, by the definition of
the two-point oscillation \eqref{9.5:eq-two-point-oscillation-b},
\begin{equation*}
|H^{\mathsf{m}}|\le\mathfrak{O}^{+}\bigl(\hat{v}^{\mathsf{B}\sharp}\bigr)=:\tfrac12 G,
\qquad \mathsf{m}\in\{\mathsf{A},\mathsf{B}\};
\end{equation*}
in particular $\sup|g^2|\le G$. At the vertices, the second line of
\eqref{9.5:eq-kept-reading-mark-a}, which holds at every vertex pair, applied at $\tau$ and at
$0$, gives
\begin{equation*}
|g^2|\le|\mathsf{r}_Y(\tau)|+|\mathsf{r}_Y(0)|\le\beta,\qquad
\beta:=16C^{\star}_{TL}\,\vartheta^{\sharp}_Y\,\mathfrak{O}^{+}\bigl(\hat{v}^{\mathsf{B}\sharp}\bigr).
\end{equation*}
These are the inputs of the lemma on $\natural$-extensions, with the part $g^2$ of the form
\eqref{9.5:eq-kept-reading-mark-2}, and of Lemma~\ref{9.5:lem-polydisc-vertices}\,(iii), with
vertex bound $\beta$ and supremum bound $\mathfrak{B}=G$. Both yield the factor
\begin{equation*}
(\beta G)^{1/2}
=\bigl(16C^{\star}_{TL}\vartheta^{\sharp}_Y\mathfrak{O}^{+}(\hat{v}^{\mathsf{B}\sharp})
\cdot 2\,\mathfrak{O}^{+}(\hat{v}^{\mathsf{B}\sharp})\bigr)^{1/2}
=\bigl(32C^{\star}_{TL}\vartheta^{\sharp}_Y\bigr)^{1/2}
\mathfrak{O}^{+}\bigl(\hat{v}^{\mathsf{B}\sharp}\bigr).
\end{equation*}
Thus the lemma on $\natural$-extensions, applied with the part $g^1$ equal to zero (so that the
bound on that part is zero), gives for the $\natural$-extension of $\mathsf{r}_Y$ which it
provides the bound $[\![\mathsf{r}_Y^{\natural}]\!]\le(\beta G)^{1/2}$, and
Lemma~\ref{9.5:lem-polydisc-vertices}\,(iii) gives the same factor $(\beta G)^{1/2}$ in its bound.
Summing over the two members of $\mathfrak{K}(Y)$, and using
\begin{equation*}
\sum_{v\in\mathfrak{K}(Y)}\mathfrak{O}^{+}(\hat{v}^{\mathsf{B}\sharp})
\le C_{\Sigma}\check{R}_m^{p_{\Sigma}},
\end{equation*}
which is the one-machine bound
\eqref{9.5:eq-two-point-oscillation-b}, gives
\begin{equation*}
\sum_{v\in\mathfrak{K}(Y)}[\![\mathsf{r}_Y^{\natural}]\!]
\le\bigl(32C^{\star}_{TL}\vartheta^{\sharp}_Y\bigr)^{1/2}C_{\Sigma}\,\check{R}_m^{p_{\Sigma}}.
\end{equation*}
\end{proof}

\begin{proposition}[The kept factors at later sites, layer by layer]\label{9.5:prop-kept-later-sites}
Let $\mathsf{A}$ and $\mathsf{B}$ be the two runs compared in this section. Assume the following:
\begin{enumerate}
\item[(1)] the conditions \textup{(F0)}--\textup{(F5)} and the hypotheses
$(H^{\mathsf{A}})$, $(H^{\mathsf{B}})$ on the runs $\mathsf{A}$ and $\mathsf{B}$;
\item[(2)] the conditions collected in \eqref{9.5:eq-kept-leaves-a} of
Notation~\ref{9.5:not-kept-leaves}, and the mixing condition on the data of the two runs;
\item[(3)] the bound \eqref{9.5:eq-family-mark-sums-a} and the condition
\eqref{9.5:eq-distance-units-b};
\item[(4)] Lemma~\ref{9.5:lem-layer-store-bounds};
\item[(5)] at sites of type \textup{(L$^{\dagger}$)}, Lemma~\ref{9.5:lem-dilation-disc};
\item[(6)] the one-run bounds for a kept factor at the sites of types \textup{(T)},
\textup{(P)}, \textup{(L$^{\circ}$)}, \textup{(L$^{\dagger}$)}, their corollary,
and the anchoring bound for kept factors;
\item[(7)] from the construction of resolved $\natural$-extensions: the one-M\"obius
identity, part~(d) of the M\"obius lemma, the $\natural$-extension lemma, the
$\natural$-norm lemma and the vertex-pair corollary.
\end{enumerate}
Let $(Y,m)$ be live at a site of step $k'$, so that the reading level is $k'$, and let
$v\in\mathfrak{K}(Y)$ be one of the two kept factors of \eqref{9.5:eq-kept-factors-a}.
Then, identically at the site's paired points $P^{\mathsf{A}}$, $P^{\mathsf{B}}$,
\begin{equation}\label{9.5:eq-kept-later-sites-a}
\begin{gathered}
\tilde{v}^{\mathsf{B}}(P^{\mathsf{B}})-\tilde{v}^{\mathsf{A}}(P^{\mathsf{A}})
=\sum_{n\le m}\delta\hat{v}^{(n)}\bigl(\mathrm{rd}^{\mathsf{A}}P^{\mathsf{A}}
\oplus c^{\mathsf{B}}_{col},\,c^{\mathsf{A}};\Phi_Y\bigr)+\mathsf{r}_Y,\\
\mathsf{k}^{+}_n:=\mathsf{k}^{(\le k')}_n+\mathsf{k}^{rd}_n,
\end{gathered}
\end{equation}
where the family marks $G^{1,(n)}$, formed from the summands $\delta\hat{v}^{(n)}$ of the
first line, lie in the classes $(n,k')$ and the reading mark
$\mathsf{r}_Y$ lies in the class $(m,k')$; the summand $\mathsf{k}^{rd}_n$ in the second
line comes from the regions $(Y,m)$ with $m=n$. Moreover, the following hold.
\begin{itemize}
\item[\textup{(T)}] At the step from $i$ to $i+1$, $i\ge m$,
\begin{equation}\label{9.5:eq-kept-later-sites-1}
\delta f_v(\mathbf{1})-\delta f_v(0)=\sum_{\mathsf{y}}\delta v_{\mathsf{y}};
\end{equation}
the class-$n$ pieces of the two members of $(Y,m)$ are together bounded by
\begin{equation}\label{9.5:eq-kept-later-sites-2}
\mathsf{k}^{+}_n\,\check{R}_m^{p_{\Sigma}}\,e^{-\kappa^{\natural}_1|\mathsf{y}|},
\qquad |\mathsf{y}|\ge\tfrac52\check{R}_m;
\end{equation}
and
\begin{equation}\label{9.5:eq-kept-later-sites-3}
\bigl|\delta V^{kept,(n)}(Y')\bigr|\le\theta_T(\check{\mathsf{T}}_i)\,
\mathsf{k}^{+}_n\,e^{-\kappa_V d_i(Y')},
\end{equation}
with anchored total at most $K(d)\,\theta_T\,\mathsf{k}^{+}_n$.
\item[\textup{(P)}] At stage $\nu+1$ of a chain of $\mathbf{R}_{k'}$, own or foreign:
\begin{itemize}
\item[(a)] $\delta(F_v-F_v(0))$ has a resolved $\natural$-extension with
$[\![\,\cdot^{((n,k'),\mathrm{u})}]\!]\le\mathfrak{O}^{+}(\delta\hat{v}^{(n)})$ and with
$[\![\,\cdot^{((m,k'),\mathrm{c})}]\!]$ at most the share assigned to the reading mark by
\eqref{9.5:eq-kept-reading-mark-a}; the $\mathrm{c}$-parts stored with
$\delta\hat{v}^{(n)}$ are carried with the family mark, and no further root is taken.
Consequently
\begin{equation}\label{9.5:eq-kept-later-sites-4}
2^{-n(\mathfrak{s})}\bigl|W^{(\zeta)}_{\hat{e}}\bigr|\le\nu^{m,\zeta}_{\natural}(\hat{e})
:=[\![\,\cdot^{\zeta}]\!]\prod_{\square\in A}K_A\,\vartheta(\square)\,
e^{-\kappa^{\natural}_1|\hat{\mathsf{y}}|};
\end{equation}
\item[(b)--(d)] the geometry of parts (b), (c) and (d) of the one-run bound at a
site of type \textup{(P)} holds without change, and
\begin{equation}\label{9.5:eq-kept-later-sites-5}
\sum\bigl\{[\![\,\cdot^{(n,k')}]\!]:\ \square_0\in\mathcal{A}_v\bigr\}
\le e\,2^d\,\mathfrak{s}_{\Sigma}\,\mathsf{k}^{+}_n\,
e^{\frac1{16}\delta_P d^{\wedge}(\square_0,R_{\nu})};
\end{equation}
\item[(e)] with $\mathsf{a}^{kept}_{\mu}$ and $\mathsf{a}^{kept}_{\mathsf{M}}$ the two
group constants of the contribution of the kept factors to $\mu[\,\cdot\,]$ and to
$\mathsf{M}_{*}[\,\cdot\,]$ in the one-run bound at a site of type \textup{(P)},
the kept groups add at most $\mathsf{a}^{kept}_{\mu}\mathsf{k}^{+}_n$ to
$\mu[\nu^{m,\zeta}_{\natural}]$ and at most
$\mathsf{a}^{kept}_{\mathsf{M}}\mathsf{k}^{+}_n$ to
$\mathsf{M}_{*}[\nu^{m,\zeta}_{\natural}]$. Hence part (ii) of the resolution lemma holds
on $\hat{\mathfrak{E}}\cup\{\hat{e}\}$, with $\hat{e}$ the index of the kept factor,
class by class, with $f^{\natural,\zeta}$ replaced by
$f^{\natural,\zeta}+(\theta_P/N)\,\mathsf{k}^{+}_{n(z)}$; here $N$ is the number of the
normalisation in part (ii) of the resolution lemma, not the rank of the gauge group.
\end{itemize}
\item[\textup{(L$^{\circ}$)}] The pieces are bounded by
$\mathsf{k}^{+}_n\check{R}_m^{p_{\Sigma}}e^{-\kappa^{\natural}_1\sum_t|\mathsf{y}_t|}$, and
\begin{equation}\label{9.5:eq-kept-later-sites-6}
\bigl\|\delta\mathfrak{v}^{kept,(n)}\bigr\|_{a,1}\le\theta^{\circ}_L(\check{\mathsf{T}}_{k'})
\,\mathsf{k}^{+}_n.
\end{equation}
\item[\textup{(L$^{\dagger}$)}] If $Y\subset C$ with $C$ treated at $\mathbf{R}_{k'}$, then
for the exterior pieces
\begin{equation}\label{9.5:eq-kept-later-sites-7}
\bigl\|\delta\mathfrak{v}^{kept,(n)}\bigr\|_{a,1}\le\theta^{\dagger}_L(\check{\mathsf{T}}_{k'})
\,\mathsf{k}^{+}_n,
\end{equation}
and the value and interior pieces are bounded by part (m3) of
Theorem~\ref{9.5:thm-production-bound}.
\end{itemize}
After each site the kept factor survives at the site's base, as in the one-run bounds at
sites of types \textup{(T)} and \textup{(P)}. Every output that holds a kept piece is
stored at a common datum together with the members of its class, and satisfies the
analogues of \eqref{9.5:eq-layer-store-bounds-a} and \eqref{9.5:eq-kept-later-sites-a}
for its own kind of output. No product of the riders $\mathfrak{u}$ and no sum of the
riders' mismatches over the lifetime of $(Y,m)$ is formed.
\end{proposition}

\begin{proof}
\emph{The identity \eqref{9.5:eq-kept-later-sites-a}.} We add and subtract
\[
\hat{v}^{\mathsf{B}\sharp}\bigl(\mathrm{rd}^{\mathsf{A}}P^{\mathsf{A}}\oplus
c^{\mathsf{B}}_{col};\Phi_Y,c^{\mathsf{A}}\bigr)
\]
in the left member of \eqref{9.5:eq-kept-later-sites-a}. The difference of
$\tilde{v}^{\mathsf{B}}(P^{\mathsf{B}})$ and this term is the reading mark
$\mathsf{r}_Y$ of Lemma~\ref{9.5:lem-kept-reading-mark}, defined in
\eqref{9.5:eq-kept-reading-mark-a} and taken here at the site's paired points, with
$\mathsf{v}=\mathrm{rd}^{\mathsf{A}}P^{\mathsf{A}}$. The remaining difference is a
difference of the runs $\mathsf{B}$ and $\mathsf{A}$ at the common datum
$(\mathrm{rd}^{\mathsf{A}}P^{\mathsf{A}}\oplus c^{\mathsf{B}}_{col},c^{\mathsf{A}})$ with
$\Phi=\Phi_Y$. By the exactness statement in \eqref{9.5:eq-layer-store-bounds-a}
(Lemma~\ref{9.5:lem-layer-store-bounds}) it equals the sum over the layers $n$ of
$\delta\hat{v}^{(n)}$ at that datum, and by the support statement in the same display
only layers $n\le m$ contribute. This is \eqref{9.5:eq-kept-later-sites-a}; the second
line there is a definition.

\emph{Family marks.} For the site's parameters $\tau$ put
\[
\eta^{\times}(\tau):=\Bigl(\mathrm{rd}^{\mathsf{A}}P^{\mathsf{A}}(\tau)\oplus
c^{\mathsf{B}}_{col}\bigl(P^{\mathsf{B}}(\tau)\bigr),\ \mathfrak{u}^{\mathsf{A}}(\tau)\,g'\Bigr).
\]
On the site's whole polydisc, $\eta^{\times}(\tau)\in\mathfrak{X}^{\sharp}_Y$: part~(a)
of the holomorphy lemma, applied in each of the two runs, gives the read datum of
$\mathsf{A}$ and the collar entries of $\mathsf{B}$ on that polydisc, and the mixing
condition of hypothesis~(2) places their combination in $\mathfrak{X}^{\sharp}_Y$. We
apply Lemma~\ref{9.5:lem-site-bound} with $f:=\delta\hat{v}^{(n)}$ and
$\eta:=\eta^{\times}$. It bounds each piece $\mathfrak{p}_I[f\circ\eta^{\times}]$ by
$\mathfrak{O}^{+}(\delta\hat{v}^{(n)})\prod_{i\in I}(R_i-1)^{-1}$. This is the bound used
in the one-run bounds at the four kinds of site with $\mathfrak{O}(v)$ in place of
$\mathfrak{O}^{+}(\delta\hat{v}^{(n)})$. Hence the four displayed bounds of those
statements hold with $\mathfrak{O}(v)$ replaced by $\mathfrak{O}^{+}(\delta\hat{v}^{(n)})$.
Summed over the two members of $\mathfrak{K}(Y)$, the quantity
$\mathfrak{O}^{+}(\delta\hat{v}^{(n)})$ is at most $x_n(m)\check{R}_m^{p_{\Sigma}}$, by
\eqref{9.5:eq-production-bound-c}, and $x_n(m)\le\mathsf{k}^{(\le k')}_n$ for $m\le k'$ by
\eqref{9.5:eq-production-bound-c}. These bounds carry the full rates $\kappa_1-1$, and
$\kappa_1-1\ge\kappa^{\natural}_1$.

At a site of type \textup{(P)} we work in the format of the resolved $\natural$-extensions
and put
\[
G^{1}_{\mathfrak{s}}:=\delta\hat{v}^{(n)}(\eta^{\times}_{\mathfrak{s}})
-\delta\hat{v}^{(n)}\bigl(\eta^{\times}_{\mathfrak{s}}\upharpoonright_{t=0}\bigr).
\]
By the one-M\"obius identity and part~(d) of the M\"obius lemma, $G^{1}_{\mathfrak{s}}$
is a DD-family on $(\mathsf{Y}(\mathfrak{s}),\rho_1)$. The argument is that which shows
the first family of the one-run $\natural$-extension to be a DD-family on $\rho_1$. That
family is likewise built from the data of $\mathsf{A}$ combined with the collar entries
of $\mathsf{B}$.

\emph{Reading mark.} At a site of type \textup{(P)} we use part \textup{(rd2)} of
Lemma~\ref{9.5:lem-kept-reading-mark}, displayed in
\eqref{9.5:eq-kept-reading-mark-a}, together with the $\natural$-norm lemma. The root is
taken in the toggle variables only. The extension is holomorphic in $\varpi$, and
$\varpi$ contains the pair $(b,t)$. Part \textup{(rd1)} of
Lemma~\ref{9.5:lem-kept-reading-mark} holds at complex $(b,t)$, since the family
$\mathfrak{F}^V$ contains the vertex pairs concerned, by the vertex-pair corollary. The
Cauchy estimate in $t$ is taken at the radii of the run. By
\eqref{9.5:eq-distance-units-b} these are the radii $r_{\square}$ of the correlation
theorem, in the units fixed by \eqref{9.5:eq-distance-units-a} of
Notation~\ref{9.5:not-distance-units}. Therefore the allocation in the one-run bound at a
site of type \textup{(P)} of the half $\frac12\delta_P$ of the decay carried by
$\vartheta(\square)$ applies to both kinds of class, $\mathrm{u}$ and $\mathrm{c}$. That
allocation is: one quarter to the corridors, by the lower bound on $M$ used there; one
sixteenth to part~(c) of the cell lemma; one sixteenth to the anchor, by the condition on
the weights used there; and one sixteenth to the share $\tau_d$ of that allocation. At
sites of types \textup{(T)} and \textup{(L$^{\circ}$)},
\textup{(L$^{\dagger}$)}, Lemma~\ref{9.5:lem-polydisc-vertices}(iii), applied on the
polydisc of the cells, bounds the reading mark with rate $\kappa^{\natural}_1$ per cell.

The family marks lie in the classes $(n,k')$ and are bounded through
$x_n(m)\check{R}_m^{p_{\Sigma}}$. The reading mark lies in the class $(m,k')$ and is
bounded by $\mathsf{k}^{rd}_m\check{R}_m^{p_{\Sigma}}$, by \textup{(rd2)} in
\eqref{9.5:eq-kept-reading-mark-a}. This is the origin of the two summands of
$\mathsf{k}^{+}_n$ in \eqref{9.5:eq-kept-later-sites-a}: for the region $(Y,m)$ itself the
reading mark falls in the class $m$ only, so that in class $n$ the summand
$\mathsf{k}^{rd}_n$ accounts for the regions created at step $n$. With the family marks
at rates $\kappa_1-1\ge\kappa^{\natural}_1$ and the reading
mark at rate $\kappa^{\natural}_1$, this gives the bounds on pieces in
\eqref{9.5:eq-kept-later-sites-2} and in \textup{(L$^{\circ}$)}. At a site of type
\textup{(P)} it gives the bounds on $[\![\,\cdot^{((n,k'),\mathrm{u})}]\!]$ and
$[\![\,\cdot^{((m,k'),\mathrm{c})}]\!]$ of part~(a). The estimate
\eqref{9.5:eq-kept-later-sites-4} then follows from the $\natural$-extension lemma.

\emph{Geometry, anchors, sums.} The proofs of parts (c)--(e) of the one-run bound at
sites of type \textup{(T)}, of parts (b)--(e) of that at sites of type \textup{(P)}, and
of the one-run bounds at sites of types \textup{(L$^{\circ}$)} and
\textup{(L$^{\dagger}$)} (the latter in its deep$^{+}$ case) read the kept factor $v$
only through $\mathfrak{O}(v)$, and linearly. Otherwise they read only the scaffold, as
the corollary of the one-run bounds states: these bounds enter through $\mathfrak{O}(v)$
alone, are linear in the terms, and are anchored by the scaffold only. We replace
$\mathfrak{O}(v)$ in those proofs by the bounds obtained above. We also replace the rate
$\kappa_1-1$ by $\kappa^{\natural}_1$. Then the proofs require
\eqref{9.5:eq-family-mark-sums-a} in place of the corresponding condition at rate
$\kappa_1-1$: at sites of types \textup{(T)} and \textup{(L$^{\circ}$)} the slices of the
allocation are taken per cell of $e^{-(\kappa_1-1)|\mathsf{y}|}$. At sites of types
\textup{(P)} and \textup{(L$^{\dagger}$)} the far factors $\vartheta(\square)$ and
$e^{-\frac12\delta_P D}$ are those of the run. For the reading mark at
\textup{(L$^{\dagger}$)} the far factor is $e^{-\frac14\delta_P D}$, split as
$\frac18+\frac18$ as in Theorem~\ref{9.5:thm-production-bound}. At sites of type
\textup{(L$^{\dagger}$)}, Lemma~\ref{9.5:lem-dilation-disc}, assumed there, gives for
$|\lambda|\le\rho_v$ the dilated data $\eta^{\mathsf{A}}_{v,\lambda}$,
$\eta^{\mathsf{B}}_{v,\lambda}$ in the tube data $\mathfrak{T}^{\mathsf{A}}_{F_v}[\frac14]$,
$\mathfrak{T}^{\mathsf{B}}_{F_v}[\frac14]$ and the mixed datum $\eta^{\times}_{v,\lambda}$
in $\mathfrak{X}_{X_v}$; these are the data at which the one-run bound at sites of type
\textup{(L$^{\dagger}$)} is read in its deep$^{+}$ case.

Finally, we read the anchoring bound for kept factors with $C_{\Sigma}$ replaced by
$\mathsf{k}^{+}_n$. At most $2^d$ strips meet at a probe; with $\varrho\ge1$ the variable
of the anchoring bound (not the auxiliary radius of the glossary),
$\check{R}_m\le2\varrho+2$. Moreover $(2\varrho+2)^{p_{\Sigma}}\le\mathfrak{s}_{\Sigma}
e^{\varrho}$, since $\mathfrak{s}_{\Sigma}$ is the supremum over $\varrho\ge1$ of
$(2\varrho+2)^{p_{\Sigma}}e^{-\varrho}$. These give the factor
$e\,2^d\mathfrak{s}_{\Sigma}\mathsf{k}^{+}_n$ in
\eqref{9.5:eq-kept-later-sites-5}, together with the anchored total in \textup{(T)}. Read
in the same way, the one-run contributions of the kept groups to $\mu[\,\cdot\,]$ and
$\mathsf{M}_{*}[\,\cdot\,]$ give part~(e), from which part~(ii) of the resolution lemma
follows class by class. The same reading gives \eqref{9.5:eq-kept-later-sites-3},
\eqref{9.5:eq-kept-later-sites-6} and \eqref{9.5:eq-kept-later-sites-7}.
\end{proof}

\begin{lemma}[Support implies age for the kept marks]\label{9.5:lem-kept-support-age}
Let $(Y,m)$ be a thrown-back region, created at step $m$, which is live at the reading level $k'$.
Write $\Lambda_{k'}^{c}$ for the complement of the small-field region $\Lambda_{k'}$ at level
$k'$. For a component $C$ of $\Lambda_{k'}^{c}$
write $j(C)$ for its first index, the least step at which a large-field region contained in $C$
was created, so that $\operatorname{age}(C)=k'-j(C)$ is the age of $C$ at level $k'$. Put
$h:=m-N(g_m)$, where $N(g_m)$ is the number of renormalisation steps in condition~(ii)
of \cite{B-CMP122a} at the coupling $g_m$ (a function of the coupling, always written with its
argument, and distinct from the rank $N$ of $\mathrm{SU}(N)$), and write $(x)_+:=\max\{x,0\}$.
Assume:
\begin{enumerate}
\item[(a)] the support clause of \eqref{9.5:eq-layer-store-bounds-a} in
Lemma~\ref{9.5:lem-layer-store-bounds}: a layered difference $\delta\hat{v}^{(n)}$ is not
identically zero only if the history of the label has zone $n$ inside $Y^{\sharp}$ with $n<m$,
or $n=m$;
\item[(b)] the first-index statement for zones of a history: a component of $\Lambda_{k}^{c}$,
at any level $k$, which holds a piece of a zone $n$ of a history, the zone lying in the
large-field region $Z_{n+1}$ of step $n+1$, has first index at most $n+1$;
\item[(c)] condition~(ii) of \cite{B-CMP122a}: a component to which the
operation $\mathbf{R}$ is applied at step $k$ is such that in the $N(g_k)$ renormalisation
steps preceding step $k$ no new large-field regions were created inside it;
\item[(d)] the geometric forgetting of old regions, which is stated for every component of
$\Lambda_{k}^{c}$, at any level $k$, whatever its history, together with the provision under
which it is applied to a component that contains a thrown-back region;
\item[(e)] the domain property of the outputs: every output holding a kept piece of
$\mathfrak{K}(Y)$ has $(Y,m)$ in its set of thrown-back regions and a true domain containing
$Y^{\sharp}$.
\end{enumerate}
Then a kept mark of $\mathfrak{K}(Y)$ of class $(n,k')$ --- a family mark, of filing zone $n$
and reading level $k'$, of the layered difference $\delta\hat{v}^{(n)}$, or, for $n=m$, the
reading mark $\mathsf{r}_Y$ (in the sense of
Proposition~\ref{9.5:prop-kept-later-sites}) --- is not
identically zero only in labels in which the history of $(Y,m)$ has zone $n$, or $n=m$. In such
a label the component $C$ of $\Lambda_{k'}^{c}$ holding $Y$ is live at $k'$, meets the domain of
the output, and
\begin{equation}\label{9.5:eq-kept-support-age-1}
j(C)\le n+1 .
\end{equation}
Consequently
\begin{equation}\label{9.5:eq-kept-support-age-2}
\operatorname{age}(C)\ \ge\ k'-n-1 ,
\end{equation}
and the geometric forgetting of old regions, applied to $C$, yields the age weight
$(k'-n+1)\,q_a^{(k'-n-1)_+}$, where $q_a$ is the age gain per step.
\end{lemma}

\begin{proof}
By hypothesis~(a), a family mark of class $(n,k')$ vanishes identically unless the history
of the label has zone $n$ inside $Y^{\sharp}$ with $n<m$, or $n=m$; the reading mark
$\mathsf{r}_Y$ has class $(m,k')$, so for it $n=m$. This is the first
assertion, and it shows that in every label where the mark is not identically zero we have
$n\le m$. We fix such a label and prove \eqref{9.5:eq-kept-support-age-1} in two cases,
according to the position of $n$ relative to $h=m-N(g_m)$.

\emph{Case $n\le h$.} Then zone $n$ is a zone of the history of $(Y,m)$ created before
step $m$, and it lies in the large-field region of the following step:
\begin{equation*}
\text{zone } n\ \subset\ \Omega_{n+1}^{c}\ \subset\ Z_{n+1},
\end{equation*}
where $\Omega_{n+1}$ is Ba{\l}aban's small-field domain at step $n+1$ and $Z_{n+1}$ is the
large-field region of that step (Definition~\ref{1.5:def-domains}).
The component $C$ holds a piece of zone $n$ (as is shown after the second case), and the
display above is the premise of hypothesis~(b) for this zone, so hypothesis~(b) gives
$j(C)\le n+1$.

\emph{Case $h<n\le m$.} For $n$ in this range, zone $n$ is a zone of the boundary terms
produced at step $m$ by the operation $\mathbf{R}_m$ itself;
these zones are created inside the large-field region $Z$ to which $\mathbf{R}_m$ is applied: the
boundary terms of $\mathbf{R}_m$ carry the levels $h+1,\dots,m$, the kept quadratic form
$\mathsf{Q}_Y$ carries the class $h+1$,
and the remaining contributions are the shallow pointed terms, the copies and the strip at
step $m$. These zones are not zones of the history of $(Y,m)$ created before step $m$, so the
argument of the first case does not apply to them. However, $\mathbf{R}_m$ is applied to the
component holding $Y$, and by hypothesis~(c), that is, condition~(ii) of
\cite{B-CMP122a}, no
new large-field regions were created inside this component in the $N(g_m)$ renormalisation
steps preceding step $m$. Hence every large-field region of that component created before
step $m$ has index at most $m-N(g_m)=h$, and
\begin{equation*}
j(C)\ \le\ h\ \le\ n\ \le\ n+1 .
\end{equation*}
This is the place where condition~(ii) enters: without it, a component created recently could
carry a class up to $N(g_m)$ steps older than itself, with no corresponding age to supply the
factor below.

In both cases $Y$ is contained in $Z_m^{\mathrm{new}}$, the part of the large-field region
created at step $m$, and by \cite[(1.102)]{B-CMP122b} a region so created stays in the
large-field region until the operation $\mathbf{R}$ is applied to its component. Since $(Y,m)$ is
live at $k'$, the component $C$ of $\Lambda_{k'}^{c}$ holding $Y$ is live at $k'$.
This is the case of the support-implies-age statement for the marks in which the region lies in
the part $Z_m^{\mathrm{new}}$ of the large-field region created at step $m$. From that statement
we take only the following: a class met in the strip of $Y$ has only the live component of $Y$
at hand, namely the component $C$, so that $C$ holds the piece of zone $n$, as was used in the
first case.
By hypothesis~(e), every output that holds a kept piece carries $(Y,m)$ among its thrown-back
regions and has a true domain containing $Y^{\sharp}$, and the component $C$ meets the domain of
the output. This proves that $C$ is live at $k'$ and meets the
domain of the output; the two cases prove \eqref{9.5:eq-kept-support-age-1}.

From $\operatorname{age}(C)=k'-j(C)$ and \eqref{9.5:eq-kept-support-age-1},
\begin{equation*}
\operatorname{age}(C)\ =\ k'-j(C)\ \ge\ k'-(n+1)\ =\ k'-n-1 ,
\end{equation*}
which is \eqref{9.5:eq-kept-support-age-2}. Finally, the geometric forgetting of old regions is
stated for every component of $\Lambda_{k}^{c}$, at every level $k$, whatever its history. The
component $C$
contains the thrown-back region $Y$ by construction, so the statement is used under the
provision of hypothesis~(d); applied to $C$, whose age at level $k'$ is at least $k'-n-1$, it
yields the age weight $(k'-n+1)\,q_a^{(k'-n-1)_+}$.
\end{proof}

We fix below the constants, the ceiling on $\gamma$ and the enlarged ceilings on which the elimination of the kept factors rests. Throughout, $\mathsf{T}$ is the scaffold variable of the coupling and $\mathsf{T}_{\gamma}$ its threshold. The unit shares of the kept factors at the four sites, $\theta_P$, $\theta_T$, $\theta^{\circ}_L$, $\theta^{\dagger}_L$, and the combined share $\bar{\Theta}$ are fixed in \eqref{9.5:eq-constituent-rates-a}; the kept quotients $x_n(m)$, $\mathsf{k}^{(<m)}_n$, $\mathsf{k}_n$ are fixed in \eqref{9.5:eq-production-bound-c}. The kept quotient $\mathsf{k}^{rd}_n$ is fixed in \eqref{9.5:eq-constituent-rates-a}, and $\mathsf{k}^{+}_n$ in \eqref{9.5:eq-kept-later-sites-a}.

We write $\mathrm{T}_i$ for the half-step that takes level $i$ to level $i+1$, and $\mathbf{R}_m$ for the half-step of step $m$ at which Ba{\l}aban's large-field operation $\mathbf{R}$ acts. The half-steps are ordered $\mathrm{T}_1,\mathbf{R}_1,\mathrm{T}_2,\mathbf{R}_2,\dots$. The upright letter $\mathrm{T}$ is not the degree variable $T$ of Definition~\ref{9.1:not-levelwise-rate}, nor the scaffold variable $\mathsf{T}$, nor Ba{\l}aban's large-field operations $\mathbf{T}$, $\mathbf{T}'$.

We write $\theta(\mathsf{T})$ for the function of $\mathsf{T}$ for which the feed-back coefficient of the stored boundary terms born at the half-step $\mathrm{T}_i$ reads
\begin{equation*}
\varkappa^{\mathbf{B}}_i=K_{CE}\,C'''_L\,\theta_i,\qquad \theta_i=\theta(\mathsf{T}_i).
\end{equation*}
Here the coefficient $\varkappa^{\mathbf{B}}_i$ of the volume and filament bound factors as a constant, which we write $K_{CE}C'''_L$, times $\theta_i$. The function $\theta$ is a negative power of $\mathsf{T}$.

The display below is arranged as follows.
\begin{itemize}
\item The rows defining $K_{\mathfrak{K}}$, $\Theta^{\star}$ and $K_{door}$ fix numbers and functions before $\gamma$.
\item The row that follows them is a hypothesis on $\gamma$.
\item The next row defines $\epsilon_{\mathfrak{K}}$, which depends on $\gamma$ through $\mathsf{T}_{\gamma}$.
\item The next row lists two inequalities that follow from the hypothesis. They are proved at the beginning of the proof of Proposition~\ref{9.5:prop-kept-elimination}.
\item The last row records the enlarged ceilings on $\gamma$, described after the display.
\end{itemize}
\begin{equation}\label{9.5:eq-kept-elimination-a}
\begin{aligned}
K_{\mathfrak{K}}&:=4C^{\natural}_{\Sigma}K_D
 +\frac{4C_{\Sigma}C_{\hat{\vartheta}}}{E_0}\Bigl[(32C^{\star}_{TL})^{1/2}+\mathsf{c}'_{rd}\Bigr],\\
\Theta^{\star}(\mathsf{T})&:=\frac{\theta_P(\mathsf{T})}{N(\mathsf{T})}+\theta_P(\mathsf{T})
 +\frac{\theta_T(\mathsf{T})}{\theta(\mathsf{T})}+\theta^{\circ}_L(\mathsf{T})
 +\theta^{\dagger}_L(\mathsf{T}),\\
K_{door}&:=\max\Bigl\{4C^{\natural}_{\Sigma}\Bigl(\frac{4}{3C^{\sharp}_0}+\frac13\Bigr),\ 
 64K_{\mathfrak{K}}\max\Bigl\{\frac{4}{3K_{\mathbb{C}}},\ \frac{K(d)}{C'''_L},\
 \frac{1}{K_{rim}}\Bigr\}\Bigr\},\\
&\sup_{\mathsf{T}\ge\mathsf{T}_{\gamma}}K_{door}\,\Theta^{\star}(\mathsf{T})\le 1,\\
\epsilon_{\mathfrak{K}}&:=K_{\mathfrak{K}}\sup_{\mathsf{T}\ge\mathsf{T}_{\gamma}}
 \max\Bigl\{\frac{4\theta_P}{3K_{\mathbb{C}}},\ \frac{K(d)\,\theta_T}{C'''_L\,\theta},\
 \frac{\theta^{\circ}_L+\theta^{\dagger}_L}{K_{rim}}\Bigr\},\\
&C^{\natural}_{\Sigma}\bar{\Theta}\le\frac14,\qquad \epsilon_{\mathfrak{K}}\le\frac{1}{64},\\
&\bigl(\bar{\varkappa},\bar{\varkappa}',K_{\mathbb{C}},C'''_L\bigr)\longmapsto
 \tfrac{65}{64}\bigl(\bar{\varkappa},\bar{\varkappa}',K_{\mathbb{C}},C'''_L\bigr)\\
&\qquad\text{in the third, fourth and fifth ceilings on }\gamma.
\end{aligned}
\end{equation}
The symbols in this display that are not fixed above are the following.
\begin{itemize}
\item $E_0$ is the constant of the bounds on $\mathsf{k}^{rd}_n$ and $\vartheta_f^{n/2}$ in \eqref{9.5:eq-constituent-rates-a}.
\item $C^{\sharp}_0$ is the constant in the definition of $\bar{\Theta}$ there.
\item $K(d)$ is the constant of the bound at the half-steps $\mathrm{T}_i$ in \eqref{9.5:eq-kept-later-sites-a}.
\item $N(\mathsf{T})$ is the function of $\mathsf{T}$ by which $\theta_P$ is divided in the first term of $\Theta^{\star}$, taken with the unit shares of \eqref{9.5:eq-constituent-rates-a}.
\end{itemize}

The hypothesis on $\gamma$ in \eqref{9.5:eq-kept-elimination-a} is the \emph{door ceiling}. It is a single ceiling on $\gamma$, in supremum form, imposed after $C^{\natural}_{\Sigma}$ has been fixed. It is one-sided, and it is not a cap in the sense of Definition~\ref{2.3:def-cap}.

By the degree of a function of $\mathsf{T}$ we mean the exponent of a power of $\mathsf{T}$ that bounds it, up to a constant factor, for large $\mathsf{T}$. The worst degree of $\Theta^{\star}$ comes from $\theta_P$. We use the condition on the auxiliary parameters $\mathfrak{p}$ that gives the unit shares their degrees. Under it, for every $\eta>0$ the degree of $\theta_P$ is at most $5r+\eta-(7r+1)$, which is negative for $\eta$ small. Under the same condition, the terms $\theta_T/\theta$, $\theta^{\circ}_L$ and $\theta^{\dagger}_L$ decay faster than every power of $\mathsf{T}$. Hence $\Theta^{\star}(\mathsf{T})\to0$ as $\mathsf{T}\to\infty$. A supremum over $\{\mathsf{T}\ge\mathsf{T}_{\gamma}\}$ does not increase as $\gamma$ decreases. Therefore the $\gamma$ satisfying the door ceiling form a non-empty interval $(0,\gamma_{door}]$, and we write $\gamma_{door}$ for its right end point.

The enlarged ceilings are the third, fourth and fifth of the five ceilings on $\gamma$ under which the stored boundary terms are eliminated. They are taken with the constants $\bar{\varkappa}$, $\bar{\varkappa}'$, $K_{\mathbb{C}}$, $C'''_L$ multiplied by $65/64$, as in the last row of \eqref{9.5:eq-kept-elimination-a}; the third is the ceiling stated for the pair $(2E^z_1,C'''_L)$, in which $C'''_L$ is likewise multiplied by $65/64$. We call them the enlarged third, fourth and fifth ceilings.

\begin{proposition}[Elimination of the kept factors under one coupling ceiling]
\label{9.5:prop-kept-elimination}
Assume the following, each with its hypotheses:
\begin{itemize}
\item[(a)] the production bound, layer by layer (Theorem~\ref{9.5:thm-production-bound});
\item[(b)] the bounds on the constituents of the production defect (Proposition~\ref{9.5:prop-constituent-bounds});
\item[(c)] the bounds on the kept factors at later sites, layer by layer (Proposition~\ref{9.5:prop-kept-later-sites}).
\end{itemize}
Assume further:
\begin{itemize}
\item[(d)] the door ceiling in \eqref{9.5:eq-kept-elimination-a};
\item[(e)] the five ceilings on $\gamma$ of the elimination of the stored boundary terms, the third to fifth in the enlarged form of \eqref{9.5:eq-kept-elimination-a};
\item[(f)] the following statements, in their form without kept factors:
\begin{itemize}
\item the load bound at the production sites and the residual lemma;
\item the bounds on the localised pieces of the kept potential;
\item the wrapping lemma for the boundary terms;
\item the cluster bound for the marked potential;
\item the volume and filament bound for the boundary terms born at the half-steps $\mathrm{T}_i$;
\item the bounds for the boundary terms born at the half-steps $\mathbf{R}_m$, with the bounds on their residual and on their own part;
\item the elimination of the stored boundary terms, with its a-priori bound on the class loads, the Gronwall-constant lemma and the corollary on its consequences;
\item the dial bound;
\item the first-variation lemma along the dial;
\item the corollary of the correlation theorem used in the volume and filament bound.
\end{itemize}
\end{itemize}
Then the two inequalities $C^{\natural}_{\Sigma}\bar{\Theta}\le\frac14$ and $\epsilon_{\mathfrak{K}}\le\frac1{64}$ of \eqref{9.5:eq-kept-elimination-a} hold. Moreover, by one induction along the half-steps $\mathrm{T}_1,\mathbf{R}_1,\mathrm{T}_2,\mathbf{R}_2,\dots$, the following hold for every aligned zone $n$ and every level $k'$.
\begin{itemize}
\item[(i)] The kept factors are bounded by the load of their layer:
\begin{equation*}
\mathsf{k}^{(\le k')}_n+\mathsf{k}^{rd}_n\le K_{\mathfrak{K}}\,[\![\delta]\!]^{\natural+,(k')}_n .
\end{equation*}
\item[(ii)] The bounds listed below hold on the family that includes the kept factors. In their right members $[\![\delta]\!]^{\natural+}_n$ is replaced by $(1+\epsilon_{\mathfrak{K}})[\![\delta]\!]^{\natural+}_n$; equivalently, $K_{\mathbb{C}}$, $\varkappa^{\mathbf{B}}_i$ and $K_{rim}$ are multiplied by $1+\epsilon_{\mathfrak{K}}$. The bounds are:
\begin{itemize}
\item the load bound at the production sites;
\item parts (ii) and (iv) of the bounds on the localised pieces of the kept potential;
\item the wrapping lemma for the boundary terms;
\item the first clause and the starred second clause of the cluster bound for the marked potential;
\item the exactness, support and value bounds of the volume and filament bound for the boundary terms born at the half-steps $\mathrm{T}_i$;
\item the exactness and support bounds, and the bounds on the residual and on the own part, for the boundary terms born at the half-steps $\mathbf{R}_m$.
\end{itemize}
The support statement for the marks is replaced by its form for the kept marks, Lemma~\ref{9.5:lem-kept-support-age}.
\item[(iii)] The following statements hold verbatim with the constants $\bar{a}^{\mathfrak{K}}$, $\epsilon^{\mathfrak{K}}_0$, $\epsilon^{\mathfrak{K}}_1$, $\epsilon^{\mathfrak{K}}_{\mathbf{B}}$:
\begin{itemize}
\item parts (i)--(iv) of the elimination of the stored boundary terms;
\item the a-priori bound $[\![\delta]\!]^{\natural+,(k')}_n\le E^z_1$;
\item parts (a) and (b) of the corollary on the consequences of that elimination.
\end{itemize}
The constants are
\begin{equation*}
\bar{a}^{\mathfrak{K}}:=(1+\epsilon_{\mathfrak{K}})\bar{a},\qquad
\epsilon^{\mathfrak{K}}_0:=e^{\bar{a}^{\mathfrak{K}}}\bar{a}^{\mathfrak{K}},\qquad
\epsilon^{\mathfrak{K}}_{\mathbf{B}}:=\epsilon^{\mathfrak{K}}_0+2\epsilon^{\mathfrak{K}}_1 .
\end{equation*}
Here $\bar{a}$ is the bound, in the elimination of the stored boundary terms, on the sum of the coefficients of its Gronwall recursion. The constant $\epsilon^{\mathfrak{K}}_1$ takes the place of the constant $\epsilon_1$ of that elimination. The unknowns of the comparison remain five; the kept factors add no sixth. The constants $\Lambda_r$, $C_E$, $C^z_{\sigma}$ keep their form.
\end{itemize}
\end{proposition}

\begin{proof}
\emph{The two consequences of the door ceiling.}

Each of the five terms of $\Theta^{\star}(\mathsf{T})$ is non-negative.

For the first consequence we use the bound of the combined share $\bar{\Theta}$ of \eqref{9.5:eq-constituent-rates-a} by the supremum of $\Theta^{\star}$:
\begin{equation*}
C^{\natural}_{\Sigma}\bar{\Theta}\le C^{\natural}_{\Sigma}\Bigl(\frac{4}{3C^{\sharp}_0}+\frac13\Bigr)
\sup_{\mathsf{T}\ge\mathsf{T}_{\gamma}}\Theta^{\star}(\mathsf{T}).
\end{equation*}
The first entry of the maximum defining $K_{door}$ equals $4C^{\natural}_{\Sigma}(4/(3C^{\sharp}_0)+1/3)$. Hence the right member is at most
\begin{equation*}
\frac14\,K_{door}\sup_{\mathsf{T}\ge\mathsf{T}_{\gamma}}\Theta^{\star}(\mathsf{T}),
\end{equation*}
and this is at most $1/4$ by the door ceiling.

For the second consequence, each of $\theta_P$, $\theta_T/\theta$ and $\theta^{\circ}_L+\theta^{\dagger}_L$ is bounded by $\Theta^{\star}$. Hence, at every $\mathsf{T}$,
\begin{equation*}
\max\Bigl\{\frac{4\theta_P}{3K_{\mathbb{C}}},\frac{K(d)\theta_T}{C'''_L\theta},
\frac{\theta^{\circ}_L+\theta^{\dagger}_L}{K_{rim}}\Bigr\}
\le\max\Bigl\{\frac{4}{3K_{\mathbb{C}}},\frac{K(d)}{C'''_L},\frac1{K_{rim}}\Bigr\}\,
\Theta^{\star}(\mathsf{T}).
\end{equation*}
Multiply by $K_{\mathfrak{K}}$ and use the second entry of the maximum defining $K_{door}$. This gives
\begin{equation*}
\epsilon_{\mathfrak{K}}\le\frac{1}{64}\,K_{door}\sup_{\mathsf{T}\ge\mathsf{T}_{\gamma}}
\Theta^{\star}(\mathsf{T})\le\frac1{64}.
\end{equation*}
In particular $1+\epsilon_{\mathfrak{K}}\le 65/64$.

\emph{The induction.} We suppose (i)--(iii) for all half-steps before the current one, and prove them at the current one.

For the first aligned half-steps there is no thrown-back region. There all kept quotients $\mathsf{k}^{(\le k')}_n$, $\mathsf{k}^{rd}_n$ vanish, so (i) holds, and (ii) and (iii) are the statements of hypothesis (f) in their form without kept factors.

At the current level we abbreviate
\begin{equation*}
[\![\delta]\!]_n:=[\![\delta]\!]^{\natural+,(i)}_n \text{ at the half-step } \mathrm{T}_i,
\qquad
[\![\delta]\!]_n:=[\![\delta]\!]^{\natural+,(m)}_n \text{ at the half-step } \mathbf{R}_m .
\end{equation*}

\emph{A half-step $\mathrm{T}_i$, from level $i$ to level $i+1$.}

The production array of the half-step gains the slots $V^{kept}(Y')$. This is part (e) of the construction of the kept factors at the half-steps $\mathrm{T}_i$. Consider the class-$n$ marks of these slots. By the bound at the half-steps $\mathrm{T}_i$ of Proposition~\ref{9.5:prop-kept-later-sites} (display \eqref{9.5:eq-kept-later-sites-a}), their anchored activity is at most $K(d)\,\theta_T(\check{\mathsf{T}}_i)\,\mathsf{k}^{+}_n$. By (i) at the preceding half-step, $\mathsf{k}^{+}_n\le K_{\mathfrak{K}}[\![\delta]\!]_n$, so this activity is at most $K(d)\,\theta_T\,K_{\mathfrak{K}}[\![\delta]\!]_n$. By the definition of $\epsilon_{\mathfrak{K}}$ in \eqref{9.5:eq-kept-elimination-a}, this is at most $\epsilon_{\mathfrak{K}}\,\theta_i\,C'''_L\,[\![\delta]\!]_n$.

In the first step of the volume and filament bound for the boundary terms born at $\mathrm{T}_i$, the marked potential $V^{m,z}$ is linear in the marks. Its activity therefore becomes at most
\begin{equation*}
\theta_i\,C'''_L\,(1+\epsilon_{\mathfrak{K}})\,[\![\delta]\!]^{\natural+,(i)}_n .
\end{equation*}
The corollary of the correlation theorem applied in that step is homogeneous of degree one in the marked activity. Hence the value bound in the volume and filament bound for these boundary terms holds with $\varkappa^{\mathbf{B}}_i$ replaced by $(1+\epsilon_{\mathfrak{K}})\varkappa^{\mathbf{B}}_i$.

The exactness bound follows from part (c) of the first-variation lemma along the dial. It is applied with the enlarged sum of dial directions $\sum_z W^{(z)}=x^{\mathsf{B}}-x^{\mathsf{A}}$. This sum is exact by the identity in \eqref{9.5:eq-kept-later-sites-a} expressing the difference of the kept factors of the two runs at paired points.

The support bound is Lemma~\ref{9.5:lem-kept-support-age}.

The bounds at the half-step $\mathrm{T}_i$ for the units with support $X\subset\Lambda_{i+1}$ read no kept piece. Indeed, a slot $V^{kept}(Y')$ carries the cube $Q_Y(i)$ in its label; this cube meets $Y^{\sharp}$, and $Y\subset Z$. So, by the argument of part (b)(i) of the corollary on the consequences of the elimination, a unit reading it has $X\supseteq Y'$ with $Y'\not\subset\Lambda_{i+1}$, hence $X\not\subset\Lambda_{i+1}$. Consequently the feed-back constants of these bounds are unchanged.

\emph{A half-step $\mathbf{R}_m$: the production site.}

We use part (e) of the production-site bound of Proposition~\ref{9.5:prop-kept-later-sites} (display \eqref{9.5:eq-kept-later-sites-a}) together with part (ii) of the residual lemma. They give the first inequality below. The load bound at the production sites without kept factors bounds $\frac43$ times the bracket by $K_{\mathbb{C}}[\![\delta]\!]_n$, and the induction hypothesis (i), at the half-steps before $\mathbf{R}_m$, bounds $\mathsf{k}^{+}_n$ by $K_{\mathfrak{K}}[\![\delta]\!]_n$; this gives the second inequality:
\begin{align*}
\sum_{\mathsf{i}}\sum_{\nu}\mathsf{s}^{\natural,(z,\mathsf{i})}_{\nu}
&\le\frac43\Bigl[\varepsilon^{\natural}\max\mathsf{l}^{\natural,\zeta}
 +\sum\mathsf{z}^{\natural,\zeta}\Bigr]+\frac43\,\theta_P\,\mathsf{k}^{+}_n\\
&\le K_{\mathbb{C}}[\![\delta]\!]_n+\frac43\,\theta_P\,K_{\mathfrak{K}}[\![\delta]\!]_n
 \le(1+\epsilon_{\mathfrak{K}})\,K_{\mathbb{C}}[\![\delta]\!]_n .
\end{align*}
The last inequality holds because $\frac43\theta_P K_{\mathfrak{K}}\le\epsilon_{\mathfrak{K}}K_{\mathbb{C}}$ by the definition of $\epsilon_{\mathfrak{K}}$. This is the load bound at the production sites with $K_{\mathbb{C}}$ multiplied by $1+\epsilon_{\mathfrak{K}}$.

The proof of the residual lemma reads the direction $(W,\mathfrak{w})$ through three properties: the bound $|W_e|\le2^{n(\mathfrak{s})}\nu^m_{\natural}(e)$, holomorphy, and locality. For the kept entries $\hat{e}$ these three properties are part (a) of the production-site bound of Proposition~\ref{9.5:prop-kept-later-sites}. The base point
\begin{equation*}
v^{\lambda}_{\hat{e}}:=(1-\lambda)v^{\mathsf{A}}_{\hat{e}}+\lambda v^{\mathsf{B}}_{\hat{e}}
\end{equation*}
has the one-run majorant $\nu(\hat{e})$ by convexity. This uses part (a) of the construction of the kept factors at the production site, in each of the two runs. The first variation is homogeneous of degree one. No smallness of $\nu^m$ is used: nothing is asked of the size of the marked factor.

\emph{A half-step $\mathbf{R}_m$: the two later sites.}

Consider the marks of $V(Y_4)$ with $Y_4\setminus Z^{\sharp}\neq\emptyset$. They gain the kept pieces $\delta\mathfrak{v}^{kept,(n)}$, namely the foreign pieces and the own exterior pieces. By the two bounds at the later sites of Proposition~\ref{9.5:prop-kept-later-sites} (display \eqref{9.5:eq-kept-later-sites-a}), by the induction hypothesis (i) at the half-steps before $\mathbf{R}_m$, and by the definition of $\epsilon_{\mathfrak{K}}$, these are bounded in $\|\cdot\|_a$ by
\begin{equation*}
(\theta^{\circ}_L+\theta^{\dagger}_L)\,\mathsf{k}^{+}_n\le\epsilon_{\mathfrak{K}}\,K_{rim}\,
[\![\delta]\!]_n .
\end{equation*}
Hence part (iv) of the bounds on the localised pieces of the kept potential, and the fourth bound of the cluster bound for the marked potential, hold with the factor $1+\epsilon_{\mathfrak{K}}$.

\emph{A half-step $\mathbf{R}_m$: the top storey.}

At the top storey the slots are numbers at seeds. The kept factors themselves enter $\mathbf{T}'$ and $\mathbf{T}$ by their values. The class marks of these values obey the value bound of Lemma~\ref{9.5:lem-layer-store-bounds} (display \eqref{9.5:eq-layer-store-bounds-a}). This is the per-cell bound of part (iv) of the bounds on the localised pieces of the kept potential, at $K_{\mathbb{C}}(1+\epsilon_{\mathfrak{K}})$.

Apply part (a) of the first-variation lemma along the dial at the dial radius $\Lambda_z=E^z_1/[\![\delta]\!]_n$. Under the enlarged third ceiling of \eqref{9.5:eq-kept-elimination-a} (recall $1+\epsilon_{\mathfrak{K}}\le65/64$), it gives verbatim the following bounds:
\begin{itemize}
\item the bound on the residual part of the boundary terms born at $\mathbf{R}_m$;
\item the bound on their own part;
\item parts (d) and (e) of the dial bound.
\end{itemize}
No kept factor is ever varied along the dial as a function of its datum; the kept factors enter the dial only through their values.

\emph{A half-step $\mathbf{R}_m$: the new thrown-back regions.}

Apply Theorem~\ref{9.5:thm-production-bound} and Proposition~\ref{9.5:prop-constituent-bounds} at $(m,n)$. With $\mathsf{k}^{(<m)}_n=\sup_{m'<m}x_n(m')$ from \eqref{9.5:eq-production-bound-c}, they give
\begin{align*}
x_n(m)&\le C^{\natural}_{\Sigma}\Bigl[K_D[\![\delta]\!]^{\natural+,(m)}_n
 +\bar{\Theta}\bigl(\mathsf{k}^{(<m)}_n+\mathsf{k}^{rd}_n\bigr)\Bigr]
 +\frac12\mathsf{k}^{(<m)}_n+\frac12\mathsf{c}'_{rd}C_{\Sigma}\vartheta_f^{n/2}\\
&=a_m+\Bigl(\frac12+C^{\natural}_{\Sigma}\bar{\Theta}\Bigr)\sup_{m'<m}x_n(m'),
\end{align*}
where
\begin{equation}\label{9.5:eq-kept-elimination-1}
a_m:=C^{\natural}_{\Sigma}K_D[\![\delta]\!]^{\natural+,(m)}_n
+C^{\natural}_{\Sigma}\bar{\Theta}\,\mathsf{k}^{rd}_n
+\frac12\mathsf{c}'_{rd}C_{\Sigma}\vartheta_f^{n/2}.
\end{equation}
The number $a_m$ does not decrease in $m$. Indeed, $[\![\delta]\!]^{\natural+,(m)}_n$ is a partial sum of non-negative terms plus terms that do not depend on $m$, and the remaining terms of \eqref{9.5:eq-kept-elimination-1} do not depend on $m$.

By the first consequence $C^{\natural}_{\Sigma}\bar{\Theta}\le\frac14$, the coefficient $\frac12+C^{\natural}_{\Sigma}\bar{\Theta}$ is at most $\frac34$. We apply the supremum recursion without compounding (Lemma~\ref{9.5:lem-sup-recursion}) to the sequence $x_n(m)$, with this coefficient in the role of $\theta$ and the non-decreasing sequence $a_m$. It gives
\begin{equation*}
x_n(m)\le4a_m\le4C^{\natural}_{\Sigma}K_D[\![\delta]\!]^{\natural+,(m)}_n+\mathsf{k}^{rd}_n
+2\mathsf{c}'_{rd}C_{\Sigma}\vartheta_f^{n/2},
\end{equation*}
where $4C^{\natural}_{\Sigma}\bar{\Theta}\le1$ was used once more.

The bounds in \eqref{9.5:eq-constituent-rates-a} give
\begin{equation*}
\mathsf{k}^{rd}_n\le\frac{2(32C^{\star}_{TL})^{1/2}C_{\Sigma}C_{\hat{\vartheta}}}{E_0}\,
[\![\delta]\!]^{c}_n,\qquad
\vartheta_f^{n/2}\le\frac{2C_{\hat{\vartheta}}}{E_0}\,[\![\delta]\!]^{c}_n .
\end{equation*}
Inserting them, we obtain
\begin{equation*}
x_n(m)+\mathsf{k}^{rd}_n\le4C^{\natural}_{\Sigma}K_D[\![\delta]\!]^{\natural+,(m)}_n
+\frac{4C_{\Sigma}C_{\hat{\vartheta}}}{E_0}\Bigl[(32C^{\star}_{TL})^{1/2}+\mathsf{c}'_{rd}\Bigr]
[\![\delta]\!]^{c}_n .
\end{equation*}
Since $[\![\delta]\!]^{c}_n\le[\![\delta]\!]^{\natural+,(m)}_n$ (by the definition of the class loads, of which $[\![\delta]\!]^{c}_n$ is the classical part), the right member is at most $K_{\mathfrak{K}}[\![\delta]\!]^{\natural+,(m)}_n$, by the definition of $K_{\mathfrak{K}}$ in \eqref{9.5:eq-kept-elimination-a}.

Take the supremum over the steps $m\le k'$ and use that $[\![\delta]\!]^{\natural+,(m)}_n$ does not decrease in $m$. This gives (i) at the current half-step.

\emph{Proof of (iii).}

The recursion of part (i) of the elimination of the stored boundary terms is indexed by the half-steps $\mu$. We write $u_\mu$ for the quantity it estimates and $\alpha_\mu$, $\beta_\mu$ for its coefficients. The Greek letters distinguish these coefficients from the number $a_m$ of \eqref{9.5:eq-kept-elimination-1}. The recursion reads
\begin{equation*}
u_{\mu+1}\le(1+\alpha_\mu)u_\mu+\alpha_\mu\mathfrak{a}_n+\beta_\mu .
\end{equation*}
By (ii) it holds with $\alpha_\mu$ replaced by $(1+\epsilon_{\mathfrak{K}})\alpha_\mu$. The new coefficients satisfy
\begin{equation*}
\sum_\mu(1+\epsilon_{\mathfrak{K}})\alpha_\mu\le(1+\epsilon_{\mathfrak{K}})\bar{a}=\bar{a}^{\mathfrak{K}}.
\end{equation*}
The following statements read the recursion only through numbers:
\begin{itemize}
\item parts (ii)--(iv) of the elimination;
\item the a-priori bound on the class loads;
\item the Gronwall-constant lemma;
\item the corollary on the consequences of the elimination.
\end{itemize}
They therefore hold verbatim at $\bar{a}^{\mathfrak{K}}$, $\epsilon^{\mathfrak{K}}_0$, $\epsilon^{\mathfrak{K}}_1$, $\epsilon^{\mathfrak{K}}_{\mathbf{B}}$. The arithmetic of the a-priori bound is the same as without kept factors, under the enlarged fourth ceiling of \eqref{9.5:eq-kept-elimination-a}.
\end{proof}

\begin{remark}
We record three features of the argument.

First, nesting does not compound. A kept factor feeds back through three routes.
\begin{itemize}
\item The inner kept factor itself contracts by $\frac12$, under the nesting condition of \eqref{9.5:eq-production-bound-b}.
\item Its exterior pieces contract by $C^{\natural}_{\Sigma}\bar{\Theta}\le\frac14$, under the door ceiling.
\item The chain terms that read it contract by the same factor $C^{\natural}_{\Sigma}\bar{\Theta}\le\frac14$, under the door ceiling.
\end{itemize}
No other route carries a term of the construction. Lemma~\ref{9.5:lem-sup-recursion} carries no factor per generation of nested regions.

Second, consider the stored terms over the strip created from kept pieces. They are ordinary members of $\mathsf{y}_n$. Their feed-back passes through the Gronwall inequality of the elimination of the stored boundary terms, with summable coefficients $\mathsf{w}_{\ell,n}\varkappa^{\mathbf{B}}$ and $\varkappa^{\mathbf{B}'}$. The kept quotient itself is a supremum over thrown-back regions, never a sum over steps. Thus $\theta_P(\mathsf{T}_k)$, which is not summable in $k$, is never summed.

Third, $K_{\mathfrak{K}}$ is not small and need not be. What enters every bound is $K_{\mathfrak{K}}$ times a unit share of the kept factors, and this product is small by the door ceiling.
\end{remark}

\begin{corollary}[Smallness of the kept factors at the current depth]\label{9.5:cor-current-depth}
Let $\mathsf{A}$ and $\mathsf{B}$ be the two runs compared, and let $\mathfrak{n}$ be the length of
$\mathsf{A}$. For a level $j$ its depth is $s=\mathfrak{n}-j$, and for a zone $n$ we put
$s(n):=\mathfrak{n}-n$. A class is a pair $z=(n,k')$ of a zone $n=n(z)$ and a reading level $k'$.
A resolved class is a pair $\zeta=(z,\mathsf{i})$, with $\mathsf{i}\in\{\mathrm{u},\mathrm{c}\}$
labelling the two parts into which the class is resolved. We write
$[\![\delta]\!]^{\natural+}_n$ for the load (the defect of the comparison) of layer $n$, and
$[\![\delta]\!]^{\natural+,(k')}_n$ for its part read up to the level $k'$. This part is at most
$[\![\delta]\!]^{\natural+}_n$, since the summands read beyond $k'$ are nonnegative. We write
$\mathsf{l}_s$, $e_s$ for the two factors of the envelope bound at depth $s$. The letter $\rho$ is the
contraction letter, and $q_a$ is the age gain per step. Assume the following.
\begin{itemize}
\item The hypotheses and the conclusions of the elimination of the kept factors under one coupling
ceiling, Proposition~\ref{9.5:prop-kept-elimination}, hold. This includes its part (i): for every
aligned zone $n$ and every $k'$,
\begin{equation*}
\mathsf{k}^{(\le k')}_n+\mathsf{k}^{rd}_n\le K_{\mathfrak{K}}[\![\delta]\!]^{\natural+,(k')}_n .
\end{equation*}
It also includes the bounds of its part (ii) on the marks of every class, and the constants
$\epsilon^{\mathfrak{K}}_0,\epsilon^{\mathfrak{K}}_1,\bar{a}^{\mathfrak{K}}$ of its part (iii).
\item The kept factors at later sites, layer by layer, Proposition~\ref{9.5:prop-kept-later-sites},
hold with their hypotheses. As there, $\mathsf{k}^+_n:=\mathsf{k}^{(\le k')}_n+\mathsf{k}^{rd}_n$.
\item There are numbers $\hat{\mathsf{l}}_b$, $b\ge0$, the growth factors of the envelope, such that
the envelope factors are monotone in the depth: for all depths $s$ and all $b\ge0$,
$e_{s+b}\le\rho^{-b}e_s$ and $\mathsf{l}_{s+b}\le\hat{\mathsf{l}}_b\,\mathsf{l}_s$.
\item The standing hypotheses of the envelope bound for the loads hold, and so does that bound: for
every admissible triple $(\epsilon_0,\epsilon_1,\bar{a})$ of constants of the Gronwall elimination,
and for every level $j$, with $s=\mathfrak{n}-j$,
\begin{equation*}
[\![\delta]\!]^{\natural+}_j\le C^z_{env}(\epsilon_0,\epsilon_1,\bar{a})\,\mathsf{l}_s^2e_s ,
\end{equation*}
where $C^z_{env}$ is the envelope constant. In the summation of that envelope bound the age gain
satisfies $q_a\le\rho^4$, and the summed age constant
\begin{equation*}
C^{TC}_{\Sigma}:=\sum_{b\ge0}(b+1)\,\hat{\mathsf{l}}_b^{\,2}\,q_a^{(b-1)_+}\rho^{-b}
\end{equation*}
is finite.
\item Two age inputs hold. The first is Lemma~\ref{9.5:lem-kept-support-age}: support implies age
for the kept marks. The second is geometric forgetting of old regions, as it enters
the first-variation formula along the dial.
\item For part (d), the following hold. First, the age-weighted bound on the two-cutoff difference
of the remainder holds in its parts (d) and (e). These parts bound the supremum $\mathsf{S}^+_k$, over
common data, of the two-cutoff difference of the remainder at level $k$ by a sum of terms, one of
which is a prefactor depending on the level times
\begin{equation*}
\sum_{n\le k}(k-n+1)\,q^{(k-n-1)_+}\,[\![\delta]\!]^{\natural+}_n ,
\end{equation*}
where $q$ is the factor per step of geometric forgetting of old regions. Second, the four statements
of this chapter through which that display is carried into the two hypotheses of the uniqueness
theorem named in part (d) hold: a corollary, a proposition, a remark and a theorem, which take the
age-weighted sums of the loads in that display to the depth hypothesis of the uniqueness theorem on
the remainder norms and to its smallness hypothesis on the parameters $\varepsilon_n$. We refer to
them below as the four carrying statements. Third, the proofs of parts (d) and (e) of the
age-weighted remainder bound, of the remark and of the proposition among the four carrying
statements, and of the target estimate of the comparison use a two-cutoff difference only through
the following five properties of its marks:
\begin{itemize}
\item the exactness of the marks;
\item the support of the marks;
\item the per-cell and per-block bounds by the load of the class;
\item the support--age property of class marks;
\item the holomorphy of the resolved extensions of the marks.
\end{itemize}
\end{itemize}
Put $C^{z\mathfrak{K}}_{env}:=C^z_{env}$ evaluated at
$(\epsilon^{\mathfrak{K}}_0,\epsilon^{\mathfrak{K}}_1,\bar{a}^{\mathfrak{K}})$. Then the following
hold.

\smallskip\noindent
(a) The recursion for the chain terms holds class by class. Let $\hat{\mathfrak{E}}$ be the family of
units of the chain, and let $\hat{e}$ be the unit of the kept factors of
Proposition~\ref{9.5:prop-kept-later-sites}. Here $\nu$ indexes the stages of the preparatory
integrations of one step of the large-field operation, $N$ is the number of these stages (in this
part the letter does not denote the rank of the gauge group), and $\varepsilon^{\natural}$ is the
feed-back ratio of the resolved chain recursion. On $\hat{\mathfrak{E}}\cup\{\hat{e}\}$, for every
stage $\nu$ and every resolved class $\zeta$,
\begin{equation}\label{9.5:eq-current-depth-1}
\begin{aligned}
\mathsf{s}^{\natural,\zeta}_{\nu+1}
&\le\frac{\varepsilon^{\natural}}{N}\sum_{i\le\nu}\mathsf{s}^{\natural,\zeta}_i
+\frac{\varepsilon^{\natural}}{N}\mathsf{l}^{\natural,\zeta}_{\nu+1}
+\frac{\theta_P}{N}\mathsf{l}^{\mathfrak{K},\zeta}+\mathsf{z}^{\natural,\zeta}_{\nu+1},\\
\sum_{\mathsf{i}}\mathsf{l}^{\mathfrak{K},(z,\mathsf{i})}
&\le\mathsf{k}^+_{n(z)}\le K_{\mathfrak{K}}[\![\delta]\!]^{\natural+}_{n(z)} .
\end{aligned}
\end{equation}
In particular, the recursion for the chain terms holds on the family that contains the kept factors.

\smallskip\noindent
(b) The envelope bound holds with the kept constants. For every level $j$, with $s=\mathfrak{n}-j$,
\begin{equation}\label{9.5:eq-current-depth-2}
[\![\delta]\!]^{\natural+}_j\le C^{z\mathfrak{K}}_{env}\,\mathsf{l}_s^2e_s,
\qquad\text{hence}\qquad
\mathsf{k}^+_n\le K_{\mathfrak{K}}C^{z\mathfrak{K}}_{env}\,\mathsf{l}^2_{s(n)}e_{s(n)} .
\end{equation}
Thus each class is small like $e$ at the depth of its own zone.

\smallskip\noindent
(c) The current depth. For every reading level $k'$, with $s':=\mathfrak{n}-k'$,
\begin{equation}\label{9.5:eq-current-depth-3}
\sum_{n\le k'}(k'-n+1)\,q_a^{(k'-n-1)_+}\,\mathsf{k}^+_n
\le K_{\mathfrak{K}}C^{TC}_{\Sigma}C^{z\mathfrak{K}}_{env}\,\mathsf{l}^2_{s'}e_{s'} .
\end{equation}

\smallskip\noindent
(d) Two hypotheses of the uniqueness theorem, its depth hypothesis on the remainder norms and its
smallness hypothesis on the parameters $\varepsilon_n$, follow from the display of the age-weighted
remainder bound and from the four carrying statements, which take the age-weighted sums of the loads
in that display to these two hypotheses, also when the kept members are
counted inside the family, because these statements depend on a two-cutoff difference only through
the loads $[\![\delta]\!]^{\natural+}_n$.
\end{corollary}

\begin{proof}
(a) Consider first the recursion in \eqref{9.5:eq-current-depth-1}. It is part (e) of the bound of
Proposition~\ref{9.5:prop-kept-later-sites} at the sites whose unit share is $\theta_P$. The same
part gives the first inequality of the second line. The second inequality of that line follows from
part (i) of Proposition~\ref{9.5:prop-kept-elimination}, because
$\mathsf{k}^+_n=\mathsf{k}^{(\le k')}_n+\mathsf{k}^{rd}_n$ and
$[\![\delta]\!]^{\natural+,(k')}_n\le[\![\delta]\!]^{\natural+}_n$. Together these give the
recursion, class by class, on $\hat{\mathfrak{E}}\cup\{\hat{e}\}$.

(b) We apply the envelope bound under the hypotheses of part (iii) of
Proposition~\ref{9.5:prop-kept-elimination}. Its constant is then $C^z_{env}$ evaluated at
$(\epsilon^{\mathfrak{K}}_0,\epsilon^{\mathfrak{K}}_1,\bar{a}^{\mathfrak{K}})$, which is
$C^{z\mathfrak{K}}_{env}$. This gives the first inequality of \eqref{9.5:eq-current-depth-2}. Insert
it, with $j=n$ and $s=s(n)$, into the second line of \eqref{9.5:eq-current-depth-1}; this gives the
second inequality.

(c) Fix $k'$ and put $s'=\mathfrak{n}-k'$. For a zone $n\le k'$ put $b:=k'-n\ge0$. Then
$s(n)=s'+b$, $k'-n+1=b+1$ and $(k'-n-1)_+=(b-1)_+$. By \eqref{9.5:eq-current-depth-2} and the
monotonicity of the envelope factors,
\begin{equation*}
\mathsf{k}^+_n
\le K_{\mathfrak{K}}C^{z\mathfrak{K}}_{env}\,\mathsf{l}^2_{s'+b}e_{s'+b}
\le K_{\mathfrak{K}}C^{z\mathfrak{K}}_{env}\,\hat{\mathsf{l}}_b^{\,2}\rho^{-b}\,
\mathsf{l}^2_{s'}e_{s'} .
\end{equation*}
Multiply by $(b+1)q_a^{(b-1)_+}$ and sum over $n\le k'$. Distinct zones $n$ give distinct
$b\ge0$, and every term is nonnegative. The sum is therefore at most
\begin{equation*}
K_{\mathfrak{K}}C^{z\mathfrak{K}}_{env}\,\mathsf{l}^2_{s'}e_{s'}
\sum_{b\ge0}(b+1)\,\hat{\mathsf{l}}_b^{\,2}\,q_a^{(b-1)_+}\rho^{-b}
=K_{\mathfrak{K}}C^{TC}_{\Sigma}C^{z\mathfrak{K}}_{env}\,\mathsf{l}^2_{s'}e_{s'} .
\end{equation*}
This is \eqref{9.5:eq-current-depth-3}. The last series is $C^{TC}_{\Sigma}$, finite by hypothesis.
The same summation appears in the envelope bound at a level $k$, and here it
is carried out at $k'$ in place of $k$.

(d) Take a mark of class $(n,k')$, kept or not. By Lemma~\ref{9.5:lem-kept-support-age} and the
support--age property of class marks, it is not identically zero only in labels containing a
component live at $k'$ of age at least $k'-n-1$.

Geometric forgetting of old regions bounds the large-field operation on every component, whatever
its history, with an un-normalised weight. In the first-variation formula along the dial it
therefore gives a dial bound $\mathfrak{B}_z$ that carries the factor $q_a^{(k'-n-1)_+}$. The kept
marks are among the marks of class $z$ bounded in part (ii) of
Proposition~\ref{9.5:prop-kept-elimination}.

The kept marks have the five properties listed in the last hypothesis. Their exactness and their
support are given by \eqref{9.5:eq-layer-store-bounds-a}. Being among the marks of class $z$ bounded
in part (ii) of Proposition~\ref{9.5:prop-kept-elimination}, they satisfy the per-cell and
per-block bounds by the load of the class, and they have the holomorphy of the resolved extensions
as members of these class-$z$ marks. Their support--age property is the one derived in the first
paragraph of this part. By the last hypothesis, the proofs of the statements listed there therefore
apply with the kept marks among the class marks, and the derivation supplies both hypotheses of the
uniqueness theorem with the kept members inside the family.
\end{proof}

\begin{remark}[The kept factors must be resolved by classes]
Let $\mathsf{l}^{\mathfrak{K}}(k')$ be the supremum, over all regions $(Y,m)$ live at $k'$, of the
kept quotient before its resolution by classes. This single number is not small at the current
depth.

By \eqref{9.5:eq-layer-store-bounds-a} it is at most $\sum_n\mathsf{k}^+_n$. In general it is of the
size of the oldest class present. If $n_c$ is the zone of that class, it is comparable with
$e_{s(n_c)}$, which is not small at the depth of $k'$.

Suppose it were entered as one number into the load $\mathsf{l}^{\natural,\zeta}_{\nu+1}$ of the
recursion \eqref{9.5:eq-current-depth-1}. It
would then feed $e_{s(n_c)}$ into the remainder norm $N^{R'}_r$ at every depth $r$. The depth
hypothesis of the uniqueness theorem on the remainder norms would fail as stated.

What yields \eqref{9.5:eq-current-depth-3} is the resolution by classes, with the age paid class by
class. The age is paid with the un-normalised weight of geometric forgetting of old regions.
\end{remark}

\begin{corollary}[The comparison on the enlarged class of paired reads]\label{9.5:cor-enlarged-reads}
Assume the following statements, each with its hypotheses:
\begin{itemize}
\item the membership statements for the mixed datum in this section, namely the holomorphy
statements in their parts (a) and (c) and the mixing statement, together with
Lemma~\ref{9.5:lem-dilation-disc};
\item Theorem~\ref{9.5:thm-production-bound} and part (i) of
Proposition~\ref{9.5:prop-kept-elimination};
\item the base bound of the kept factors in the run $\mathsf{B}$, and
Lemma~\ref{9.5:lem-kept-reading-mark};
\item Lemma~\ref{9.5:lem-polydisc-vertices}, the $\natural$-norm lemma, and part (b) of the
proposition on vertex values and one-run bounds (bounds holding within a single run).
\end{itemize}
Let $v\in\mathfrak{K}(Y)$ be a kept factor of the thrown-back region $Y$, and let
$\mathcal{K}^{\sharp}$ be any system of paired reads of $v$, that is, of pairs
$((P^{\mathsf{B}},c^{\mathsf{B}}),(P^{\mathsf{A}},c^{\mathsf{A}}))$, with $P^{\mathsf{A}}$,
$P^{\mathsf{B}}$ tube points of the two lattices and $c^{\mathsf{A}}$, $c^{\mathsf{B}}$ their
collar entries, obtained from base pairs by
the moves of the sites at one parameter $\tau$ (complex values of the move parameters $b$, $t$
of the sites and of the decoupling
parameters $\mathsf{s}$ being allowed), by the riders (the charts), by a common unfolding, and by
the arity-two disc in a dilation parameter $\lambda$, formed from two reads with the same
$\varpi$ (the disc on which Lemma~\ref{9.5:lem-dilation-disc} is stated), at one $\lambda$.
Then, in the
comparison of the runs $\mathsf{A}$ and $\mathsf{B}$ used for the uniqueness theorem, at every
point of $\mathcal{K}^{\sharp}$
\begin{equation}\label{9.5:eq-enlarged-reads-1}
\delta v=\sum_{n}\delta\hat{v}^{(n)}(\eta^{\times})+\mathsf{r}_Y ,
\end{equation}
where we call the sum the family part and $\mathsf{r}_Y$ the reading part, and the following hold.
\begin{enumerate}
\item[(a)] The family part is small on all of $\mathcal{K}^{\sharp}$, class by class:
$\eta^{\times}\in\mathfrak{X}^{\sharp}_Y$ at every point of $\mathcal{K}^{\sharp}$, the
$\lambda$-discs included. Consequently, for any two points of $\mathcal{K}^{\sharp}$ with the
same value of $\Phi_Y$, the inner variable at which the two-point oscillation
\eqref{9.5:eq-two-point-oscillation-b} is taken, the mixed second
difference of the family part between them is at most
$\sum_n\mathfrak{O}^{+}(\delta\hat{v}^{(n)})$, where $v$ is one of the two kept factors, the
kept potential or the kept quadratic form, so that each summand is at most the left member of
\eqref{9.5:eq-enlarged-reads-2}; and for every $n$
\begin{equation}\label{9.5:eq-enlarged-reads-2}
\mathfrak{O}^{+}(\delta\hat{V}^{(n)})+\mathfrak{O}^{+}(\delta\hat{\mathfrak{q}}^{(n)})
\le K_{\mathfrak{K}}\,[\![\delta]\!]^{\natural+}_n\,\check{R}_m^{p_{\Sigma}} .
\end{equation}
The condition on systems of paired reads that prices the arity-two disc, that is, that controls
what admitting the dilated points costs in a bound, therefore holds
for the family part at no cost, because $\mathfrak{X}^{\sharp}_Y$ is a fixed domain.
\item[(b)] The reading part is small at the vertices of $\mathcal{K}^{\sharp}$: at every point
of $\mathcal{K}^{\sharp}$ whose $\mathsf{s}$-parameters are at natural vertices ($b$, $t$ and
the seeds of the moves being allowed complex),
\begin{equation}\label{9.5:eq-enlarged-reads-3}
|\mathsf{r}_Y|\le 8C^{\star}_{TL}\,\vartheta^{\sharp}_Y\,C_{\Sigma}\,\check{R}_m^{p_{\Sigma}} ;
\end{equation}
the field-law estimate at presented pairs is read at these points, on the vertex pair family
$\mathfrak{F}^V$.
\item[(c)] Off the vertices the reading part is bounded, but not small, and its smallness there
is not needed: at complex $\mathsf{s}$,
\begin{equation}\label{9.5:eq-enlarged-reads-4}
|\mathsf{r}_Y(\tau)-\mathsf{r}_Y(0)|\le 2C_{\Sigma}\,\check{R}_m^{p_{\Sigma}} .
\end{equation}
No smallness off the vertices follows from (b) and one-run bounds: by part (b) of the
proposition on vertex values and one-run bounds, vertex bounds and one-run bounds give
only $\min\{2^{|\mathsf{y}|}\beta,\,q^{|\mathsf{y}|}\mathfrak{B}\}$, where $|\mathsf{y}|$ is the
number of decoupling parameters of the piece, $\beta$ bounds the values at the vertices,
$\mathfrak{B}$ bounds the supremum on the polydisc and $q$ is the constant of the Cauchy
estimate of that proposition (the letters $\beta$, $\mathfrak{B}$ and $q$ are local to this
clause; $q$ is not the auxiliary parameter of $\mathfrak{p}$); and this is sharp, the extremal case
being $f=c\prod_i(2\mathsf{s}_i-1)$. But every piece, the alternating difference
$\mathfrak{p}[\,\cdot\,]$ of Lemma~\ref{9.5:lem-polydisc-vertices} over a set of decoupling
parameters, is at most $\mathsf{k}^{rd}_m\check{R}_m^{p_{\Sigma}}$ times the far factor of the
corresponding field-law estimate at the pair rate, and pieces are all that those estimates read.
\end{enumerate}
Smallness of the supremum of $|\delta v|$ over $\mathcal{K}^{\sharp}$ is therefore more than these
means give, and more than the argument for the uniqueness theorem uses; what takes its place is the
conjunction of (a), (b) and (c).
\end{corollary}

\begin{proof}
\emph{Part (a).} Each operation generating $\mathcal{K}^{\sharp}$ keeps the mixed datum in the
common data: for the moves of the sites at one parameter, the riders and the common unfolding this is
the content of the holomorphy statements, parts (a) and (c), and of the mixing statement; for the
arity-two disc at one $\lambda$ it is Lemma~\ref{9.5:lem-dilation-disc}, which places the dilated
mixed datum $\eta^{\times}_{v,\lambda}$ in the common data, holomorphically in $\lambda$, for
every two production points with the same $\varpi$ and $|\lambda|\le\rho_v$. Starting from base
pairs and applying these operations, $\eta^{\times}\in\mathfrak{X}^{\sharp}_Y$ at every point of
$\mathcal{K}^{\sharp}$, the $\lambda$-discs included.

Let two points of $\mathcal{K}^{\sharp}$ with the same $\Phi_Y$ be given. By
\eqref{9.5:eq-enlarged-reads-1} the family part at each of them is
$\sum_n\delta\hat{v}^{(n)}$ evaluated at a mixed datum lying in $\mathfrak{X}^{\sharp}_Y$.
Since $\mathfrak{O}^{+}$ of \eqref{9.5:eq-two-point-oscillation-b} is the two-point
oscillation over common data at one value of the inner variables, the mixed second difference of
the $n$-th term between the two points is at most $\mathfrak{O}^{+}(\delta\hat{v}^{(n)})$;
summing over $n$ gives the first bound of (a). For \eqref{9.5:eq-enlarged-reads-2},
Theorem~\ref{9.5:thm-production-bound} bounds
$\mathfrak{O}^{+}(\delta\hat{V}^{(n)})+\mathfrak{O}^{+}(\delta\hat{\mathfrak{q}}^{(n)})$ by
the production budget of layer $n$,
\begin{equation*}
\bigl[C^{\natural}_{\Sigma}\mathsf{d}^{(n)}_m+\tfrac12\mathsf{k}^{(<m)}_n
+\tfrac12\mathsf{c}'_{rd}C_{\Sigma}\vartheta_f^{n/2}\bigr]\check{R}_m^{p_{\Sigma}},
\end{equation*}
and part (i) of Proposition~\ref{9.5:prop-kept-elimination} bounds, for every half-step index
$k'$ of the induction in that proposition, the kept quotients of layer $n$ by the load,
$\mathsf{k}^{(\le k')}_n+\mathsf{k}^{rd}_n\le K_{\mathfrak{K}}[\![\delta]\!]^{\natural+,(k')}_n$.
The bound \eqref{9.5:eq-enlarged-reads-2} is obtained from the conjunction of these two
statements. Neither bound depends on which operations produced the two points,
since both are suprema over the one set $\mathfrak{X}^{\sharp}_Y$; in particular the points added
by the dilation disc cost nothing, which is the last assertion of (a).

\emph{Part (b).} A point of $\mathcal{K}^{\sharp}$ whose $\mathsf{s}$-parameters are at natural
vertices is a vertex pair, in the family $\mathfrak{F}^V$, of a site at which
Lemma~\ref{9.5:lem-kept-reading-mark} applies. The first estimate of
\eqref{9.5:eq-kept-reading-mark-a} in Lemma~\ref{9.5:lem-kept-reading-mark} gives there
$|\mathsf{r}_Y|\le 8C^{\star}_{TL}\vartheta^{\sharp}_Y\,\mathfrak{O}^{+}(\hat{v}^{\mathsf{B}\sharp})$,
and the base bound of the kept factors in the run $\mathsf{B}$, which is among the hypotheses of
Lemma~\ref{9.5:lem-kept-reading-mark}, bounds
$\mathfrak{O}^{+}(\hat{v}^{\mathsf{B}\sharp})$ by $C_{\Sigma}\check{R}_m^{p_{\Sigma}}$. This is
\eqref{9.5:eq-enlarged-reads-3}.

\emph{Part (c).} The base bound of the kept factors in the run $\mathsf{B}$, applied twice, each
application contributing $C_{\Sigma}\check{R}_m^{p_{\Sigma}}$, gives
\eqref{9.5:eq-enlarged-reads-4} at complex $\mathsf{s}$. That (b) and one-run bounds give no
smallness off the vertices is part (b) of the proposition on vertex values and one-run bounds,
with the extremal case $f=c\prod_i(2\mathsf{s}_i-1)$ stated there. For the pieces, apply
Lemma~\ref{9.5:lem-polydisc-vertices} to $G=\mathsf{r}_Y$ as a function of the decoupling
parameters $\vec{\mathsf{s}}$: by its part (i) a piece is the alternating sum of the
values of $\mathsf{r}_Y$ at the vertices $\mathbf{1}_{\sigma'}$; by its part (iii), with $\beta$
the vertex bound \eqref{9.5:eq-enlarged-reads-3}, with $\mathfrak{B}$ the bound on the polydisc
furnished by the base bound in the run $\mathsf{B}$, and with $R>1$ the polydisc radius of
Lemma~\ref{9.5:lem-polydisc-vertices}, the piece is bounded by
$(2/(R-1))^{|\mathsf{y}|/2}(\beta\mathfrak{B})^{1/2}$. Together with the $\natural$-norm lemma this
is the second estimate of \eqref{9.5:eq-kept-reading-mark-a} in
Lemma~\ref{9.5:lem-kept-reading-mark}: every piece is at most
$\mathsf{k}^{rd}_m\check{R}_m^{p_{\Sigma}}$, with
$\mathsf{k}^{rd}_m=(32C^{\star}_{TL})^{1/2}C_{\Sigma}\vartheta_f^{m/2}$, times the far factor of
the corresponding field-law estimate at the pair rate.
\end{proof}

\begin{lemma}[Two-run bounds are read only at vertex pairs]\label{9.5:lem-vertex-pairs-only}
On the route through resolved terms (Definition~\ref{9.5:def-resolved-terms}), no statement of
this section reads a bound on the difference of the two runs $\mathsf{A}$ and $\mathsf{B}$ at a
pair of points other than
\begin{itemize}
\item[($\alpha$)] a vertex pair of $\mathfrak{F}^V$ presented over a seed inside the production of
one output, or
\item[($\beta$)] for a pair $(Y,m)$ as in Proposition~\ref{9.5:prop-kept-later-sites}, with kept
factors $\mathfrak{K}(Y)$, a vertex pair of a site after the site $(\mathrm{L})_m$, restricted to
the enlarged strip $(Y^{\sharp})^{+}$ of $Y$.
\end{itemize}
In particular:
\begin{itemize}
\item[(a)] no pair at which a two-run bound is read lies on a disc of jointly dilated pairs of
points of the two runs, formed from two reads at the same auxiliary parameters $\varpi$: the discs
of Lemma~\ref{9.5:lem-dilation-disc} are discs of the data of one
run, with the collar entries of the other run moved by that one run;
\item[(b)] the family $\mathfrak{F}^V$ is not enlarged: the pairs ($\alpha$) are the typed chains
of the output $V(X^{\sharp})$ whose production is meant in ($\alpha$), which is an output by the
proposition on hull reads and suprema
over common data, and the pairs ($\beta$) are restrictions to $(Y^{\sharp})^{+}$ of vertex pairs of
$\mathfrak{F}^V$ at later sites;
\item[(c)] the sets over which the suprema $\mathsf{S}^{+,(n)}_j[F]$ are taken are the fixed
domains $\mathfrak{X}_{X_F}$, which do not depend on later steps; consequently, for any step $m$,
no class of paired reads at the site \textup{(P)} of \eqref{9.5:eq-kept-later-sites-a} of the
operation $\mathbf{R}_m$ of step $m$ is enlarged by the discs of later steps, and no
bound on such an enlarged class is required;
\item[(d)] no product of charts along a chain of reads is formed (see the last sentence of
Proposition~\ref{9.5:prop-kept-later-sites}): the two-run counterpart of the one-run bound on
such products, namely a bound on the sum of the mismatches of the charts along a chain, is not
needed, and the law for the phase increment $\mathsf{w}$ in \eqref{9.5:eq-kept-leaves-a} is a
statement at one site at a time;
\item[(e)] no single class of paired reads closed at once under the paired moves of the sites at
one parameter, the charts, the unfolding map (a rough group-valued change of gauge at
$\Sigma_Y$) and the two-point dilation disc is used on this
route; the requirement of such a joint closure has no object.
\end{itemize}
\end{lemma}

\begin{proof}
We go through the statements of this section and list every place at which a bound on the
difference of the two runs is read, and record what each statement says about the points at
which the bound is read.

First, two-run bounds enter through the hypothesis \eqref{9.5:eq-reading-mark-sums-a} of
Lemma~\ref{9.5:lem-reading-mark-sums}.

Second, they enter through the reading mark of a kept factor, the first estimate
of \eqref{9.5:eq-kept-reading-mark-a} in Lemma~\ref{9.5:lem-kept-reading-mark}. That estimate is
stated at every vertex pair of every site after the site $(\mathrm{L})_m$ of step $m$
(Proposition~\ref{9.5:prop-kept-later-sites}), on an aligned label, with the kept factor read
on $(Y^{\sharp})^{+}$; these are pairs of type ($\beta$).

Third, they enter through the sizes of the stores, which are the suprema
$\mathsf{S}^{+,(n)}_j[F]$ over common data. By the proposition on hull reads and suprema over
common data these suprema are taken over the domains $\mathfrak{X}_{X_F}$. Each such domain is
fixed by the term $F$ itself, so no later step enters it.

Fourth, they enter through the resolved extensions of the terms. By the first of their defining
conditions these are read at vertex pairs.

Fifth, they enter through the standing hypothesis that compares the two runs at common parameter
points, which is one of the hypotheses of Proposition~\ref{9.5:prop-constituent-bounds}.

No other statement of this section reads a two-run bound. The dilation disc
(Lemma~\ref{9.5:lem-dilation-disc}), the tuple sums of the family marks
(Lemma~\ref{9.5:lem-family-mark-sums}) and the site-independent bound of a piece by the
oscillation (Lemma~\ref{9.5:lem-site-bound}) are one-run statements or scaffold statements. The
same holds for the two further statements of this section used on this route besides these three.
Each of these five statements is applied to functions on the sets
$\mathfrak{X}$ of common data, and none of them compares the two runs at a pair of points. This
proves the main assertion.

The particular assertions follow from the list.
\begin{itemize}
\item[(a)] The only discs that occur are those of Lemma~\ref{9.5:lem-dilation-disc}. That lemma is
a one-run statement: its disc is formed from the data of one run, with the collar entries of the
other run moved by that one run. Since no two-run bound is read on it, no pair lies on a disc of
jointly dilated pairs of points of the two runs formed from two reads at the same $\varpi$.
\item[(b)] The pairs of type ($\alpha$) are the typed chains of $V(X^{\sharp})$, which is an
output, and the pairs of type ($\beta$) are restrictions to $(Y^{\sharp})^{+}$ of vertex pairs of
$\mathfrak{F}^V$ at later sites. Hence no pair outside
$\mathfrak{F}^V$, or outside restrictions of its pairs, is read.
\item[(c)] This is the third entry of the list.
\item[(d)] By (c) and the list, no bound is read on a class enlarged by charts of later sites.
Hence no product of charts along a chain of reads is formed, and no bound on the accumulated
mismatch of the charts along a chain is used; the law for the phase increment $\mathsf{w}$ in
\eqref{9.5:eq-kept-leaves-a} enters the second entry only as a statement at one site at a time.
\item[(e)] By the main assertion every two-run bound is read at a vertex pair of type ($\alpha$)
or ($\beta$). By (a)--(d), neither the discs nor the products of charts along chains produce
further pairs. A class of paired reads closed at once under the four operations listed in (e) is
therefore never used, and the requirement of such a joint closure has no object. \qedhere
\end{itemize}
\end{proof}

\begin{proposition}[One-run bound on the kept weight at later sites]\label{9.5:prop-one-run-weight}
Fix one cutoff and consider one run $\mathsf{m}$ in the setting of the uniqueness theorem. Let $Y$
be a region thrown back at
step $m$, with the two kept factors $V^{\sharp}(Y^{\sharp})$ and $\mathfrak{q}_Y$ of
\eqref{9.5:eq-kept-factors-a}. We put
\begin{equation}\label{9.5:eq-one-run-weight-1}
v:=V^{\sharp}(Y^{\sharp})+\mathfrak{q}_Y,\qquad I_Y:=\exp v ,
\end{equation}
and write $\hat{v}^{\sharp}$ for $v$ as a function of the common datum in
$\mathfrak{X}^{\sharp}_Y$ and $\tilde{v}$ for $v$ as a function of the point at which a site reads it.
Assume the following.
\begin{itemize}
\item[(a)] The common data, the charts and the two-point oscillation of
\eqref{9.5:eq-two-point-oscillation-a}, on this run, together with the one-machine bound
\eqref{9.5:eq-two-point-oscillation-b} on the two-point oscillation $\mathfrak{O}^{+}$.
\item[(b)] The description of the kept factors and the jointly invariant path
\eqref{9.5:eq-kept-leaves-a}.
\item[(c)] The run $\mathsf{m}$ is one run of \eqref{9.5:eq-distance-units-c}, in one machine,
under the hypotheses displayed there.
\item[(d)] Clauses (v), (t) and (q1) of the tube theorem for the kept potential, which holds
under the hypotheses of (c).
\item[(e)] Clause (o) of the statement of the two kept factors, \eqref{9.5:eq-kept-factors-a}.
\item[(f)] Clause (a) of the holomorphy lemma for the read maps.
\end{itemize}
Then the following hold.
\begin{itemize}
\item[(i)] At every site later than the law site of level $m$ the weight $I_Y$ is read through
$\tilde{v}$. The alternating differences in which it enters there resolve into pieces, each
bounded as in Lemma~\ref{9.5:lem-site-bound} with the oscillation
$\mathfrak{O}^{+}(\hat{v}^{\sharp})$, and $\Sigma\le C_{\Sigma}\check{R}_m^{p_{\Sigma}}$, where
$\Sigma$, $C_{\Sigma}$ and $p_{\Sigma}$ are as in clause (v) of the tube theorem of
hypothesis~(d).
\item[(ii)] The bounds on the kept factors at the four later sites hold, namely those whose unit
shares are $\theta_T$, $\theta_P$, $\theta^{\circ}_L$ and $\theta^{\dagger}_L$, together with
their corollary.
\item[(iii)] Parametrise the strip of $Y$ by strip variables $t_{\square}$, one for each cube
$\square$ of the strip, with $|t_{\square}|\le r_{\square}$, the strip radii satisfying
$r_{\square}>1$. For a finite
non-empty set $D$ of pairs $(\square,r)$ put
\begin{equation}\label{9.5:eq-one-run-weight-2}
C_D[\tilde{v}]:=\sum_{D'\subset D}(-1)^{|D\setminus D'|}\,\tilde{v}(\mathbf{1}_{D'}),
\end{equation}
where $\tilde{v}(\mathbf{1}_{D'})$ is $\tilde{v}$ evaluated at $t_{\square}=1$ for the cubes of
$D'$ and $t_{\square}=0$ for all other cubes. Then
\begin{equation}\label{9.5:eq-one-run-weight-3}
\bigl|C_D[\tilde{v}]\bigr|\le \mathfrak{O}^{+}(\hat{v}^{\sharp})
\prod_{(\square,r)\in D}(r_{\square}-1)^{-1},\qquad D\neq\emptyset .
\end{equation}
\end{itemize}
In particular, under the tube theorem of hypothesis~(d), the pieces of $\log I_Y$ at the later
sites are bounded.
\end{proposition}

The proposition bounds pieces only. It does not assert the counting of these pieces in the
bookkeeping by strip variables of (iii). That counting needs in addition a weighted sum over levels,
used in place of a supremum over levels, and a ceiling condition in that bookkeeping. In the
bookkeeping of
hypothesis~(c), the bounds of (ii) need neither.

\begin{proof}
\emph{The read at a site.} Let $P$ be a site later than the law site of level $m$. The value
$\tilde{v}(P)$ is the value of $v$ at a slice point of the tube, so the read of $I_Y$ at $P$ is
$\exp\tilde{v}(P)$. This is the sense in which $I_Y$ is read through $\tilde{v}$.

\emph{Two reads at one fluctuation field.} Let $\Phi_Y$ be the fluctuation field on $Y$ read by
$I_Y$. Two reads of $v$ at the same $\Phi_Y$ differ by at most the oscillation of $v$, which here is
the two-point oscillation $\mathfrak{O}^{+}(\hat{v}^{\sharp})$ of
\eqref{9.5:eq-two-point-oscillation-a}: the two reads are two values of $\hat{v}^{\sharp}$ on common
data at one fluctuation field. This follows from clauses (v) and (q1) of the tube
theorem, hypothesis~(d), together with clause (o) of the statement of the two kept factors,
hypothesis~(e).

Clause (v) of the tube theorem also gives the bound $\Sigma\le C_{\Sigma}\check{R}_m^{p_{\Sigma}}$
stated in~(i). Clause (t) of the same theorem gives the oscillation as a
function of the scaffold variable $\mathsf{T}_m$ of the coupling.

\emph{Holomorphy.} By clause (a) of the holomorphy lemma, hypothesis~(f), the map
$\tau\mapsto\eta(\tau)\in\mathfrak{X}^{\sharp}_Y$ through which a site reads $\hat{v}^{\sharp}$ is
holomorphic on a neighbourhood of a product polydisc $\prod_i\{|\tau_i|\le R_i\}$ with $R_i>1$.

\emph{Bound on the pieces.} We apply Lemma~\ref{9.5:lem-site-bound} with $f=\hat{v}^{\sharp}$ and
the read map $\eta$ just described. Its hypothesis, the holomorphy of $\eta$, is the previous step;
the inequality
\begin{equation*}
\bigl|\hat{v}^{\sharp}\circ\eta-\hat{v}^{\sharp}\circ\eta(0)\bigr|
\le\mathfrak{O}^{+}(\hat{v}^{\sharp})
\end{equation*}
used in its proof is here the statement on two reads at one fluctuation field.
Lemma~\ref{9.5:lem-site-bound} then gives,
for every non-empty $I$,
\begin{equation}\label{9.5:eq-one-run-weight-4}
\bigl|\mathfrak{p}_I[\hat{v}^{\sharp}\circ\eta]\bigr|\le\mathfrak{O}^{+}(\hat{v}^{\sharp})
\prod_{i\in I}(R_i-1)^{-1}.
\end{equation}
Since $I_Y$ is read through $\tilde{v}$, its brackets at the site resolve into pieces of
$\tilde{v}=\hat{v}^{\sharp}\circ\eta$ of this form. This proves
(i).

Assertion (ii) is obtained from hypotheses (a), (b) and (c) together with the piece bound
\eqref{9.5:eq-one-run-weight-4}, which is the estimate of the kept weight that the four bounds of
(ii) and their corollary use.

\emph{The strip variables.} Take $\mathsf{y}=D$ and $\mathsf{s}(\Delta)=t_{\square}$ for
$\Delta=(\square,r)\in D$ in Lemma~\ref{9.5:lem-polydisc-vertices}. Clause (i) of that lemma
identifies the sum \eqref{9.5:eq-one-run-weight-2} with the piece over these variables:
\begin{equation*}
C_D[\tilde{v}]=\prod_{(\square,r)\in D}\int_0^1 dt_{\square}\,\partial_{t_{\square}}\tilde{v}
\,\Big|_{t=0\ \text{off}\ D}
=\mathfrak{p}_D[\hat{v}^{\sharp}\circ\eta].
\end{equation*}
The strip variables range over the polydisc $|t_{\square}|\le r_{\square}$, with $r_{\square}>1$.
Applying \eqref{9.5:eq-one-run-weight-4} with $I=D$ and $R_i=r_{\square}$ therefore gives
\eqref{9.5:eq-one-run-weight-3}.
\end{proof}

\begin{corollary}[The Gaussian factor kept by a live arrested region]\label{9.5:cor-kept-gaussian}
Let $\mathsf{A}$ and $\mathsf{B}$ be the two runs compared, let $W$ be a live arrested region with class
index $j_W$, and let
\begin{equation}\label{9.5:eq-kept-gaussian-1}
\mathsf{q}_W(x)=-\tfrac12\langle x,\mathcal{A}_{WW}\,x\rangle
\end{equation}
be the Gaussian factor that is part of the weight of the large-field operation $\mathbf{T}'_k(W)$;
here $x$ is the fluctuation field on $W$ and $\mathcal{A}_{WW}$ is the operator of this quadratic form
on $W$.
Assume:
\begin{enumerate}
\item[(a)] part (a) of the proposition on terms of interrupted runs, in which $\mathsf{q}_W$ is excluded
from the family $\mathfrak{F}^{B4}(\mathfrak{a})$ treated there;
\item[(b)] the classical operator rows $(\alpha)$ and $(\gamma)$, on the vertex pair family
$\mathfrak{F}^V$;
\item[(c)] the hypotheses of Lemma~\ref{9.5:lem-kept-reading-mark};
\item[(d)] in each of the runs $\mathsf{A}$ and $\mathsf{B}$, every site reads $\mathsf{q}_W$ through a
map satisfying the hypotheses of Lemma~\ref{9.5:lem-site-bound}.
\end{enumerate}
Then $\mathsf{q}_W$ is a kept member containing none of the unknowns of the comparison. Its entry in
the layer-resolved difference of Definition~\ref{9.5:def-layer-store} is the triple
$(\hat{\mathsf{q}}^{\mathsf{A}},\hat{\mathsf{q}}^{\mathsf{B}};\delta\hat{\mathsf{q}}^{(j_W)})$, in the
single class $j_W$ and in the classical part of that class. At a site where a unit $v$ reads the
datum $\eta^{\mathsf{A}}_v$ in run $\mathsf{A}$ and the datum $\eta^{\mathsf{B}}_v$ in run
$\mathsf{B}$, the difference of the two forms is the sum of a family part, the difference of
$\hat{\mathsf{q}}^{\mathsf{B}}$ and $\hat{\mathsf{q}}^{\mathsf{A}}$ at $\eta^{\mathsf{A}}_v$, and a
reading part, the difference of $\hat{\mathsf{q}}^{\mathsf{B}}$ at $\eta^{\mathsf{B}}_v$ and at
$\eta^{\mathsf{A}}_v$. The family part is bounded by
\begin{equation}\label{9.5:eq-kept-gaussian-3}
\frac{K_1}{c_{\mathcal{A}}}\Bigl(2C_{cl}(1+\varpi_0)\,\vartheta^{j_W}\Bigr)^{1/2}\,
\mathfrak{O}(\mathsf{q}_W),
\end{equation}
where $\mathfrak{O}(\mathsf{q}_W)$ is the one-run oscillation of $\mathsf{q}_W$, and $K_1$,
$c_{\mathcal{A}}$, $C_{cl}$ and $\varpi_0$ are the constants of the classical operator rows $(\alpha)$
and $(\gamma)$. The reading part obeys the first of the bounds \eqref{9.5:eq-kept-reading-mark-a},
with $\hat{\mathsf{q}}^{\mathsf{B}}$ in place of the kept potential of run $\mathsf{B}$, by the proof
of Lemma~\ref{9.5:lem-kept-reading-mark}. Every piece of either part at a site is bounded by
Lemma~\ref{9.5:lem-site-bound}.
\end{corollary}

\begin{proof}
By hypothesis (a), $\mathsf{q}_W$ is excluded from part (a) of the proposition on terms of interrupted
runs, so no bound on it can be taken from there. It is part of the weight of $\mathbf{T}'_k(W)$, and it
is therefore carried along as a kept member. By \eqref{9.5:eq-kept-gaussian-1}, $\mathsf{q}_W$ is an
explicit quadratic form in the fluctuation field with a fixed operator; none of the unknowns of the
comparison enters it. Its entry is accordingly the triple
$(\hat{\mathsf{q}}^{\mathsf{A}},\hat{\mathsf{q}}^{\mathsf{B}};\delta\hat{\mathsf{q}}^{(j_W)})$. It enters
one class only, the class $j_W$, and enters that class in its classical part.

At a site where the unit $v$ reads $\eta^{\mathsf{A}}_v$ and $\eta^{\mathsf{B}}_v$, adding and
subtracting $\hat{\mathsf{q}}^{\mathsf{B}}(\eta^{\mathsf{A}}_v)$ gives
\begin{equation}\label{9.5:eq-kept-gaussian-2}
\begin{split}
\hat{\mathsf{q}}^{\mathsf{B}}(\eta^{\mathsf{B}}_v)-\hat{\mathsf{q}}^{\mathsf{A}}(\eta^{\mathsf{A}}_v)
&=\bigl[\hat{\mathsf{q}}^{\mathsf{B}}(\eta^{\mathsf{A}}_v)
-\hat{\mathsf{q}}^{\mathsf{A}}(\eta^{\mathsf{A}}_v)\bigr]\\
&\quad+\bigl[\hat{\mathsf{q}}^{\mathsf{B}}(\eta^{\mathsf{B}}_v)
-\hat{\mathsf{q}}^{\mathsf{B}}(\eta^{\mathsf{A}}_v)\bigr];
\end{split}
\end{equation}
the first bracket is the family part and the second the reading part. This split is necessary. The
classical operator rows compare the operators of the two runs at common parameter points, so they
apply only when both forms are evaluated at one datum. In the first bracket both forms are evaluated
at the datum $\eta^{\mathsf{A}}_v$ of run $\mathsf{A}$. The rows $(\alpha)$ and $(\gamma)$ of
hypothesis (b), applied on $\mathfrak{F}^V$ to this difference, give the bound
\eqref{9.5:eq-kept-gaussian-3} by a multiple of the one-run oscillation of $\mathsf{q}_W$.

The second bracket is the difference of the single form $\hat{\mathsf{q}}^{\mathsf{B}}$ at the datum
read by run $\mathsf{B}$ and at the datum read by run $\mathsf{A}$. This is the situation of the reading
mark in Lemma~\ref{9.5:lem-kept-reading-mark}, with $\hat{\mathsf{q}}^{\mathsf{B}}$ in place of the kept
potential of run $\mathsf{B}$. Under hypothesis (c), the proof of that lemma applies word for word and
yields the first of the bounds \eqref{9.5:eq-kept-reading-mark-a} for the reading part.

Finally, fix one of the runs $\mathsf{A}$, $\mathsf{B}$ and a site reading $\mathsf{q}_W$. Hypothesis (d)
provides, in that run, the hypotheses of Lemma~\ref{9.5:lem-site-bound}. That lemma bounds every piece
$\mathfrak{p}_I[\cdot]$ of $\mathsf{q}_W$ read at the site by
$\mathfrak{O}^{+}(\mathsf{q}_W)\prod_{i\in I}(R_i-1)^{-1}$. This bound is independent of the site. Doing
this in each of the two runs bounds the pieces of both brackets of \eqref{9.5:eq-kept-gaussian-2}.
\end{proof}

\begin{remark}[Order of choice and classification of the new conditions]
\label{9.5:rem-new-constants-order}
The words \emph{floor}, \emph{ceiling}, \emph{free upward} and \emph{cap}, and the \emph{sup form}
and the \emph{degree} of a ceiling, are used in the sense of
Definition~\ref{2.3:def-cap} and Notation~\ref{2.3:not-stages-first}. The constants and conditions
used in this section enter the order of choice of Notations~\ref{2.3:not-stages-first},
\ref{2.3:not-stages-middle} and~\ref{2.3:not-stages-last} as follows.

\begin{enumerate}
\item[(a)] \emph{Order.} After $\kappa$, the constants are chosen in the order
\begin{equation}\label{9.5:eq-new-constants-order-1}
\begin{aligned}
&\kappa_1\to\kappa^{\natural}_1\to M_1\to M\to\dots\to r^{0}
 \to\{C^{\natural}_{age},\,r_{\mathfrak{K}},\,C_{\mathfrak{K}},\,C^{+}_{\mathfrak{K}}\}
 \to\{C^{\sharp}_{TL},\,C^{\mathfrak{K}}_{TL},\,C^{\mathbb{Y}}_{TL}\}\\
&\to\{C^{\star}_{TL},\,\mathsf{c}_{rd},\,\mathsf{c}'_{rd},\,\mathsf{c}_G,\,C_{\hat{\vartheta}}\}
 \to\{E_0,\,B_0,\,C^{\sharp}_0,\,b_{\delta},\,\mathsf{K}^{0},\,\mathsf{K}^{\flat}\}\\
&\to\{\mathsf{K}^{\flat,1/2},\,K_{pt},\,C_{br}\}\to\{C_{\Sigma},\,p_{\Sigma}\}
 \to C^{\natural}_{\Sigma}\to\{\mathfrak{n}^{kept},\,\mathfrak{s}_{\Sigma}\}\\
&\to\{K_{\mathbb{C}},\,K^{fr}_{\mathbb{C}},\,C'''_L,\,C_{rim},\,C^{+}_X,\,
 C^{\times}_{\mathbf{C}},\,C_{\mathsf{K}}\}\\
&\to\{K_{val},\,K_{rim},\,K_D\}\to K_{\mathfrak{K}}\to K_{door}\to\gamma_U\to g,
\end{aligned}
\end{equation}
followed by the scaffold functions. An arrow means that what stands on its right is chosen after
what stands on its left, and a brace lists constants chosen at one place. Here $r_{\mathfrak{K}}$ and
$C_{\mathfrak{K}}$ are the constants of the dilation disc \eqref{9.5:eq-dilation-disc-a};
$C^{\star}_{TL}$, $\mathsf{c}_{rd}$, $\mathsf{c}'_{rd}$, $\mathsf{c}_G$, $C_{\hat{\vartheta}}$ are the
constants of the reading marks; $\mathsf{K}^{0}$, $\mathsf{K}^{\flat}$, $\mathsf{K}^{\flat,1/2}$ are the
one-run per-cube budget functions; $C_{\Sigma}$, $C^{\natural}_{\Sigma}$ and the width exponent
$p_{\Sigma}$ are those of the production budget in Theorem~\ref{9.5:thm-production-bound};
$K_{pt}$, $K_{rim}$, $K_{\mathbb{C}}$ and $K^{fr}_{\mathbb{C}}$ are the constants of the pointed walk,
of the rim and of the chain terms; $K_D$ is the constant of
Proposition~\ref{9.5:prop-constituent-bounds}, $K_{\mathfrak{K}}$ that of
Proposition~\ref{9.5:prop-kept-elimination}, and $K_{door}$ that of the door ceiling in
\eqref{9.5:eq-kept-elimination-a}. Further, $C^{\natural}_{age}$ and $C^{+}_{\mathfrak{K}}$ belong,
with $r_{\mathfrak{K}}$ and $C_{\mathfrak{K}}$, to the constants of the reading marks and of the
dilation disc; $\mathfrak{n}^{kept}$ is the size constant of the kept factors; $K_{val}$ is the
constant of the value row; and $C_{br}$ is the constant chosen together with
$\mathsf{K}^{\flat,1/2}$ and $K_{pt}$. The remaining constants in
\eqref{9.5:eq-new-constants-order-1}, namely $r^{0}$; $C^{\sharp}_{TL}$, $C^{\mathfrak{K}}_{TL}$,
$C^{\mathbb{Y}}_{TL}$; $E_0$, $B_0$, $C^{\sharp}_0$, $b_{\delta}$; $\mathfrak{s}_{\Sigma}$; and
$C'''_L$, $C_{rim}$, $C^{+}_X$, $C^{\times}_{\mathbf{C}}$, $C_{\mathsf{K}}$, are not new in this
section: each is fixed at its place in the order of choice outside this section, and of each this
remark uses only its place in \eqref{9.5:eq-new-constants-order-1}, that is, which constants stand
before it. The constants from $r^{0}$ to $K_{door}$ are chosen after $M$ and before $\gamma_U$, among
the remaining parameters of Notation~\ref{2.3:not-stages-last}.
Each constant reads only
constants standing before it. The floor \eqref{9.5:eq-family-mark-sums-a} is imposed with $\kappa_1$
and the floor \eqref{9.5:eq-distance-units-b} with $M_1$, at the stage of
Notation~\ref{2.3:not-stages-middle} where these are chosen; $M$ is chosen under the conditions on $M$
already fixed in the order of choice, unchanged. The constants new in this section are
$C^{\natural}_{age}$, $r_{\mathfrak{K}}$, $C_{\mathfrak{K}}$, $C^{+}_{\mathfrak{K}}$; $C^{\star}_{TL}$,
$\mathsf{c}_{rd}$, $\mathsf{c}'_{rd}$, $\mathsf{c}_G$, $C_{\hat{\vartheta}}$; $\mathsf{K}^{\flat,1/2}$,
$K_{pt}$, $C_{br}$; $C^{\natural}_{\Sigma}$; $K_{val}$, $K_{rim}$, $K_D$; $K_{\mathfrak{K}}$; $K_{door}$.
The four constants $C^{\natural}_{age}$, $r_{\mathfrak{K}}$, $C_{\mathfrak{K}}$, $C^{+}_{\mathfrak{K}}$ read
only
\begin{equation*}
(d,L,M,M_1,\kappa_1,C_0,C_1,B^{\sharp};\,B_{\natural},c_1,\varphi_{\mathrm{arrest}},
c_{\mathsf{U}},C_{\mathsf{a}},\mathsf{w}_{\boxplus},c_{+}).
\end{equation*}
Here $d=4$ is the dimension, $C_0$, $C_1$ are among Ba{\l}aban's auxiliary parameters and $B^{\sharp}$
is the coupling-flow constant of Notation~\ref{0.2:not-glossary}; $B_{\natural}$ is the Lipschitz
constant of the local chart, $\mathsf{w}_{\boxplus}$ the profile, and $c_{\mathsf{U}}$ the constant
indexed by the amplitude $\mathsf{U}$; $c_1$, $\varphi_{\mathrm{arrest}}$, $C_{\mathsf{a}}$ and $c_{+}$
are constants fixed in the order of choice before $C^{\natural}_{age}$.
The constants $C^{z}_{\sigma}$, $C^{\star}_X$, $C^{\natural}_E$, $C_E$ are unchanged and keep their
places in the order of choice; none of them reads $K_{\mathfrak{K}}$. The coupling
ceiling $\gamma_U$ is chosen, at the stage of Notation~\ref{2.3:not-stages-last} where it is fixed,
under the conditions already imposed on it in Notation~\ref{2.3:not-stages-last}, together with the
ceilings
\eqref{9.5:eq-production-bound-b} and \eqref{9.5:eq-kept-elimination-a} (the latter comprising the
door ceiling and the three ceilings listed in (c) whose constants are enlarged by the factor $65/64$).

\item[(b)] \emph{No arrow returns.} $C^{\natural}_{\Sigma}$ reads none of $\mathfrak{n}^{kept}$,
$K_{\mathfrak{K}}$, $\gamma$; $K_D$ reads no $C^{\natural}_{\Sigma}$; $K_{\mathfrak{K}}$ reads
$C^{\natural}_{\Sigma}$ and $K_D$; the door function $\Theta^{\star}$ is taken at unit sizes and
contains no $C_{\Sigma}$.

\item[(c)] \emph{Classification of the new displayed conditions.}
\begin{itemize}
\item[--] \eqref{9.5:eq-family-mark-sums-a}: a floor on $\kappa_1$ (free upward), imposed after
$(L,\kappa)$ and before $M_1$, $M$; it consists of the floors on $\kappa_1$ already imposed in the
order of choice, read at the pair decoupling rate $\kappa^{\natural}_1$.
\item[--] \eqref{9.5:eq-distance-units-b}, namely $\delta_1M_1\ge4\delta_P$: a floor on $M_1$ (free
upward), imposed after $L$ and before $M$; with $\delta_P=\delta_0/8$ and $\delta_1=\delta_0/32$ it
reads $M_1\ge16$.
\item[--] \eqref{9.5:eq-production-bound-b}: a ceiling on $\gamma$ in sup form, of degree
$-\infty$; imposed with $\gamma_U$.
\item[--] The door ceiling of \eqref{9.5:eq-kept-elimination-a},
\begin{equation*}
\sup_{\mathsf{T}\ge\mathsf{T}_{\gamma}}K_{door}\,\Theta^{\star}(\mathsf{T})\le1,
\end{equation*}
is a ceiling on $\gamma$ in sup form, imposed with $\gamma_U$ and after $C^{\natural}_{\Sigma}$; under
the conditions of the order of choice on the exponent $r$ (the exponent of the condition
$p_0\ge5r$ on the degree function $p_0(\cdot)$) its worst degree is $5r-(7r+1)$ increased by
an arbitrarily small positive amount, and this is negative.
\item[--] The three ceilings on $\gamma$ of \eqref{9.5:eq-kept-elimination-a} that replace the
corresponding ceilings of the order of choice: these are those ceilings with their constants
multiplied by $65/64$; imposed with $\gamma_U$.
\end{itemize}
All other conditions read in this section, eight in number, one of them a condition accompanying
the choice of $M$, are conditions of the order of choice, read unchanged; none of them is new.

\item[(d)] \emph{Definitions and readings.} The display \eqref{9.5:eq-two-point-oscillation-a}, the
definition of the kept factors in \eqref{9.5:eq-kept-leaves-a} and the assignment of a unit's own
members (Proposition~\ref{9.5:prop-constituent-bounds}) at $\ell$ inside a thrown-back region are
definitions and carry no truth value. The
following are readings of notation, each stated at its place, and impose no condition on the
constants: the three readings of the case with kept factors that accompany
\eqref{9.5:eq-distance-units-a}, of which one is assumed in Lemma~\ref{9.5:lem-dilation-disc} and
another in Proposition~\ref{9.5:prop-constituent-bounds};
\eqref{9.5:eq-family-mark-sums-b}, for the second constituent line of
Proposition~\ref{9.5:prop-constituent-bounds} only; the reading of the pointed walk; the reading of
the weights assumed in Theorem~\ref{9.5:thm-production-bound}; the reading of the mixed datum assumed
in Lemma~\ref{9.5:lem-dilation-disc}; \eqref{9.5:eq-two-point-oscillation-c}; and the two readings
contained in \eqref{9.5:eq-kept-leaves-a}, of which one fixes which of the two symbols for the kept
quadratic form, $\mathsf{Q}_Y$ or $\mathfrak{q}_Y$, is meant and is a remark on norms, the other
concerns the phase increment $\mathsf{w}$. The two-run difference of the
field-independent remainder terms, \eqref{9.5:eq-remainder-difference-a}, is a statement of this
section in its own right; geometric forgetting of old regions is applied across a throw-back, at two
places of this section.

\item[(e)] \emph{No caps.} No inequality introduced in this section bounds from above any of $L$, $M$,
$M/M_1$, $\kappa$, $\kappa_1$, $K$, $\mathfrak{n}$, $k$, $k'$, $m$, an age $k'-m$ or $k'-n$, a depth,
$D$, $N$, $N_0$ (the quantities of these two names in Theorem~\ref{9.5:thm-production-bound}),
$\check{R}_m$, the number, size, nesting or past of the thrown-back regions and of the
arrests, or the volume; nor does any bound $g$ or $\gamma$ from below. $L$ enters only through
constants determined by it; the existential choice of $L_0$ is not affected, and no constant is given
a value. Ages occur only in the factor $q_a^{(k'-n-1)_{+}}$, as a gain. None of the following is used:
the lemma on thrown-back regions, the old-age tail statement, rarity of young large-field blocks, any
clustering property, any gap, uniqueness in the volume, injectivity. Every weight used is
un-normalised, as in geometric forgetting of old regions.
\end{enumerate}
\end{remark}

\begin{proof}
(a) For each constant in \eqref{9.5:eq-new-constants-order-1}, the constants it reads are those
entering its definition at its place in this section, and these stand before it in
\eqref{9.5:eq-new-constants-order-1}; for $C^{\natural}_{age}$, $r_{\mathfrak{K}}$, $C_{\mathfrak{K}}$,
$C^{+}_{\mathfrak{K}}$ this list is the one displayed in (a). The constants $C^{z}_{\sigma}$,
$C^{\star}_X$, $C^{\natural}_E$, $C_E$ are unchanged because none of them reads $K_{\mathfrak{K}}$: the
kept members enter the five-unknown system only through the kept feed-back ratio, which satisfies
$\epsilon_{\mathfrak{K}}\le1/64$.

(b) $C^{\natural}_{\Sigma}$ is determined by constants chosen before it, among them $r_{\mathfrak{K}}$
and $\mathsf{K}^{\flat,1/2}$; by \eqref{9.5:eq-new-constants-order-1}, $\mathfrak{n}^{kept}$,
$K_{\mathfrak{K}}$ and $\gamma$ are all chosen after $C^{\natural}_{\Sigma}$, so none of them is read.
The constant $K_D$ is the maximum of the constants listed in
Proposition~\ref{9.5:prop-constituent-bounds}; each of these is a ratio of a constant of one of the
constituent bounds of the comparison to a constant of the size functionals $\mathfrak{N}_v[\cdot]$,
and the contribution of the
kept factors is displayed apart, through the unit share $\bar{\Theta}$, so it does not enter $K_D$.
Among them $K_{rim}=2\mathsf{R}^{\sharp}/E^{z}_1$ is the constant of the order of choice taken on the
family without kept factors, $\mathsf{R}^{\sharp}$ being the rim constant of that family; hence
$K_{rim}$ reads no $C^{\natural}_{\Sigma}$ either.
Therefore $K_D$ reads no $C^{\natural}_{\Sigma}$. $K_{\mathfrak{K}}$ is chosen after both
$C^{\natural}_{\Sigma}$ and $K_D$, which it reads. Finally $\Theta^{\star}$ is evaluated at unit sizes,
so no $C_{\Sigma}$ occurs in it, and the door ceiling of \eqref{9.5:eq-kept-elimination-a}, imposed
after $C^{\natural}_{\Sigma}$, reads only constants already chosen.

(c) \eqref{9.5:eq-family-mark-sums-a} consists of lower bounds on $\kappa^{\natural}_1$, read after
$(L,\kappa)$; it is classified as a floor on $\kappa_1$, free upward, in the sense of
Notation~\ref{2.3:not-stages-first}. For \eqref{9.5:eq-distance-units-b}, inserting $\delta_1=\delta_0/32$
and $\delta_P=\delta_0/8$ gives
\begin{equation*}
\delta_1M_1\ge4\delta_P\iff\frac{\delta_0}{32}M_1\ge\frac{\delta_0}{2}\iff M_1\ge16,
\end{equation*}
a lower bound on $M_1$ alone, hence a floor, free upward. The conditions
\eqref{9.5:eq-production-bound-b}, the door ceiling of \eqref{9.5:eq-kept-elimination-a} and the
three ceilings of \eqref{9.5:eq-kept-elimination-a} with constants enlarged by the factor $65/64$
are upper bounds on quantities depending on $\gamma$, taken in sup form, and are imposed with
$\gamma_U$, the last constant before $g$; their degrees are those stated in (c), and for the door
ceiling $5r-(7r+1)$, increased by an arbitrarily small positive amount, is negative under the
conditions of the order of choice on $r$.

(d) The displays listed there define objects or fix the sense in which a notation is read; none of
them bounds a constant, so none enters the classification (c).

(e) By (c), every new condition is either a floor on $\kappa_1$ or on $M_1$, hence a lower bound, or a
ceiling on $\gamma$, hence an upper bound on $\gamma$; none of them is a cap in the sense of
Definition~\ref{2.3:def-cap}, and none bounds $g$ or $\gamma$ from below. In the floors the
quantities $\kappa_1$ and $M_1$ are bounded from below only; the ceilings bound $\gamma$ only, and
$L$ enters them only through constants it determines; no value is assigned to any constant, so the
existential choice of $L_0$ is untouched. The only place where an age enters is the
factor $q_a^{(k'-n-1)_{+}}$, which is used as a gain. The inputs of the statements of this section are
those named in their hypotheses, and the statements listed as unused in (e) are not among them;
every weight used in this section is un-normalised, as in geometric forgetting of old regions, so no
normalised measure, clustering property, gap or uniqueness statement is called upon.
\end{proof}

\begin{remark}[Why the inner variables are kept at one common value]\label{9.5:rem-inner-variables}
Let $Y$ be a thrown-back region created at step $m$, of width $\check{R}_m$, with kept factors
$\mathfrak{K}(Y)$ as in Definition~\ref{9.5:def-kept-factors}, and let $\Phi$ denote the inner
variables over which the operation $\mathbf{T}'_m(Y)$ integrates, with inner measure $d\nu(\Phi)$.
In this chapter the two runs $\mathsf{A}$ and $\mathsf{B}$ are compared with the inner variables
kept, as in the representation \cite[(1.102)]{B-CMP122b}, as common parameters held at one value.
The alternative is the representation \cite[(1.104)]{B-CMP122b}, in which the inner variables are
integrated and one gauge-invariant scalar per region, the logarithm $\log\mathbf{T}'_m(Y)1$, is
kept. For a single run this alternative is harmless on a reduced tube, that is, on a sub-tube
$\mathfrak{T}^{\mathsf{m}}_v[\varsigma]$ of data of the kept factors $v\in\mathfrak{K}(Y)$ using at
most a fraction $\varsigma<1$ of the room
(Notation~\ref{9.5:not-two-point-oscillation}). For the comparison of the two runs it produces,
besides a difference of integrands at a common value of the inner variables, of the kind measured
by the two-point oscillation of Notation~\ref{9.5:not-two-point-oscillation}, the difference of the
expectations of one function under
two normalised inner measures; the bound available for the value difference of the kept factors
controls the total-variation distance of these measures only through the exponential of a sum over
the cells of $Y$, that is, through the number of old cells in $Y$, and the per-cell weights of the
age bounds act on un-normalised integrals only. Keeping the inner variables common at one value is
what lets the comparison avoid every normalised inner measure.
\end{remark}

\begin{proof}
\emph{The alternative representation.} In \cite[(1.104)]{B-CMP122b} the product
$\prod_i \mathbf{T}'_k(Y_i)\exp\sum\mathbf{B}'$ is written as
\begin{align*}
\prod_i \mathbf{T}'_k(Y_i)\exp\sum\mathbf{B}'
&=\prod_i\bigl(\mathbf{T}'_k(Y_i)1\bigr)\\
&\quad\cdot\Bigl[\prod_i\bigl(\mathbf{T}'_k(Y_i)1\bigr)^{-1}\mathbf{T}'_k(Y_i)\exp\sum\mathbf{B}'\Bigr],
\end{align*}
so that only characteristic
functions, $\delta$-functions and the functions $\mathbf{T}'_k(Y_i)1$ multiply the action density.
Applied to $Y$, the kept object becomes
\begin{equation*}
I_Y(\zeta):=\mathbf{T}'_m(Y;\zeta)1=\int d\nu(\Phi)\,e^{v_{\mathrm{tot}}(\zeta;\Phi)},\qquad
v_{\mathrm{tot}}:=V^{\sharp}(Y^{\sharp})+\mathsf{Q}_Y ,
\end{equation*}
a function of the slice $\zeta$ alone: it has no inner variable, and no chart and no law site in
the sense of Notation~\ref{9.5:not-two-point-oscillation}. The scalar kept in this representation
is $\log I_Y(\zeta)=\log\mathbf{T}'_m(Y;\zeta)1$. A later site reads the ratio
$J_Y(\zeta):=I_Y(\zeta)/I_Y(\zeta_0)$, where $\zeta_0$ is the real base point.

\emph{One run.} Put $\Delta v:=v_{\mathrm{tot}}(\zeta;\Phi)-v_{\mathrm{tot}}(\zeta_0;\Phi)$ and let
$\mu_{\zeta_0}$
be the measure proportional to $e^{v_{\mathrm{tot}}(\zeta_0;\cdot)}d\nu$; since $\zeta_0$ is real,
it is a probability measure \cite{B-CMP122b}. Then
$J_Y(\zeta)=\langle e^{\Delta v}\rangle_{\mu_{\zeta_0}}$. Choose $\zeta_1$ in the reduced tube, on
the complex line through $\zeta_0$ in the direction of $\zeta$, such that
$\zeta_0+\tau(\zeta_1-\zeta_0)$ lies in the reduced tube for all $\tau\in\mathbb{C}$ with
$|\tau|\le1$, and write $\zeta=\zeta_0+\tau(\zeta_1-\zeta_0)$; $|\tau|$ is then the relative size
of the displacement of $\zeta$ from $\zeta_0$ along this complex line. For fixed $\Phi$, $\Delta v$
is holomorphic in $\tau$ on this disc, vanishes at $\tau=0$, and satisfies
$|\Delta v|\le C_{\Sigma}\check{R}_m^{p_{\Sigma}}$ on the disc, by the one-run bound on the kept
potential; the Schwarz
lemma therefore gives
$\sup_\Phi|\Delta v|\le|\tau|\,C_{\Sigma}\check{R}_m^{p_{\Sigma}}$. By the statement that bounds
the increments at which later sites read the kept factors at $Y^{\sharp}$, these sites have
$|\tau|\le e^{-c_Y\check{R}_m}$, with a constant $c_Y>0$ fixed in that statement. Since
$\mu_{\zeta_0}$ is a probability measure,
\begin{align*}
|J_Y(\zeta)-1|&=\bigl|\langle e^{\Delta v}-1\rangle_{\mu_{\zeta_0}}\bigr|
\le\sup_\Phi|e^{\Delta v}-1|\le e^{\sup_\Phi|\Delta v|}-1\\
&\le\exp\bigl(C_{\Sigma}\check{R}_m^{p_{\Sigma}}e^{-c_Y\check{R}_m}\bigr)-1,
\end{align*}
which is small for $\check{R}_m$ large. This is confined to a reduced tube: on the whole tube of
data the oscillation of the kept
potential grows as $g\downarrow0$, by the same one-run statement on the kept potential, so $I_Y$
may vanish there and the kept scalar $\log I_Y$ is not
defined there.

\emph{Two runs.} Write $J^{\mathsf{m}}_Y$, $\Delta v^{\mathsf{m}}$ and
$\mu^{\mathsf{m}}:=\mu^{\mathsf{m}}_{\zeta_0}$ for the objects
of run $\mathsf{m}\in\{\mathsf{A},\mathsf{B}\}$, put $\delta J:=J^{\mathsf{B}}_Y-J^{\mathsf{A}}_Y$
and $f:=e^{\Delta v^{\mathsf{A}}}-1$. Adding and subtracting
$\langle e^{\Delta v^{\mathsf{A}}}\rangle_{\mu^{\mathsf{B}}}$, and using that both measures have
total mass one,
\begin{equation}\label{9.5:eq-inner-variables-1}
\delta J=\bigl\langle e^{\Delta v^{\mathsf{B}}}-e^{\Delta v^{\mathsf{A}}}\bigr\rangle_{\mu^{\mathsf{B}}}
+\bigl[\langle f\rangle_{\mu^{\mathsf{B}}}-\langle f\rangle_{\mu^{\mathsf{A}}}\bigr].
\end{equation}
The first member of \eqref{9.5:eq-inner-variables-1} is a two-run difference of integrands at a
common value of $\Phi$, of the kind measured by the two-point oscillation over common data at one
value of the inner variables (Notation~\ref{9.5:not-two-point-oscillation}). For the second,
for every constant $a\in\mathbb{C}$,
\begin{equation*}
\bigl|\langle f\rangle_{\mu^{\mathsf{B}}}-\langle f\rangle_{\mu^{\mathsf{A}}}\bigr|
=\bigl|\langle f-a\rangle_{\mu^{\mathsf{B}}}-\langle f-a\rangle_{\mu^{\mathsf{A}}}\bigr|
\le\sup_\Phi|f-a|\;\|\mu^{\mathsf{B}}-\mu^{\mathsf{A}}\|_{\mathrm{TV}}.
\end{equation*}
Since $\mu^{\mathsf{m}}$ is proportional to $e^{v^{\mathsf{m}}_{\mathrm{tot}}(\zeta_0;\cdot)}d\nu$,
the total-variation distance of the two normalised inner measures is controlled only by the
oscillation in $\Phi$ of the value difference
\begin{equation*}
\delta v_{\mathrm{tot}}(\zeta_0;\Phi):=v^{\mathsf{B}}_{\mathrm{tot}}(\zeta_0;\Phi)
-v^{\mathsf{A}}_{\mathrm{tot}}(\zeta_0;\Phi).
\end{equation*}
The bound available for this value difference is a volume: writing $|\Gamma^0_n\cap Y|$ for the
number of cells of layer $n$ inside $Y$, clause~(iv) of the bound on the resolved
extension of the layered defects of the comparison gives
\begin{equation*}
\sup_\Phi|\delta v_{\mathrm{tot}}(\zeta_0;\Phi)|
\le\sum_n 2K_{\mathbb{C}}[\![\delta]\!]^{\natural+}_n\,|\Gamma^0_n\cap Y| .
\end{equation*}
Accordingly $\|\mu^{\mathsf{B}}-\mu^{\mathsf{A}}\|_{\mathrm{TV}}$ is bounded by a constant multiple
of
\begin{equation*}
\exp\Bigl(\sum_n 2K_{\mathbb{C}}[\![\delta]\!]^{\natural+}_n\,|\Gamma^0_n\cap Y|\Bigr)-1 .
\end{equation*}
This reads the number of old cells of $Y$, hence the depth to which older regions are nested inside
$Y$, whereas $\sup_\Phi|f|$, which by the one-run bound is at most a constant multiple of
$C_{\Sigma}\check{R}_m^{p_{\Sigma}}e^{-c_Y\check{R}_m}$, does not. The per-cell factors supplied by
geometric forgetting of old regions, and by the other per-cell weight used with it in the age
analysis, enter only un-normalised integrals $\mathbf{T}'(X)G$ of an integrand $G$ over the inner
variables of a region $X$; neither bounds $|\Gamma^0_n\cap Y|$ inside a normalised expectation.
To close the estimate of
the second member of \eqref{9.5:eq-inner-variables-1} along this route one would need a bound
\begin{equation*}
\|\mu^{\mathsf{B}}-\mu^{\mathsf{A}}\|_{\mathrm{TV}}\le C[\![\delta]\!],
\end{equation*}
with a constant $C$ and $[\![\delta]\!]:=\sum_n[\![\delta]\!]^{\natural+}_n$ the total defect of
the comparison, carrying no factor that counts cells of $Y$,
which is a statement on normalised inner measures and is not among the bounds established for the
two runs. Moreover the supports of the two inner integrations differ (what is available is one
configuration $\Phi$ lying in both), and the difference of the corresponding characteristic
functions is bounded only in un-normalised form.

\emph{Conclusion.} With the inner variables kept, as in \cite[(1.102)]{B-CMP122b}, as common
parameters at one value, every two-run difference is a difference of integrands at a common value
of $\Phi$ inside an un-normalised integral, and no normalised inner measure enters the comparison.
\end{proof}

\begin{definition}[Nested grids, cells and the scaled length]\label{9.5:def-scaled-length}
\emph{Space and grids.} Let $d\ge2$ and $L\ge2$ be integers (in the application to
Theorem~\ref{3.1:thm-main}, $d=4$ and $L$ is the odd blocking factor, subject to the lower bounds
of that theorem), let $\eta>0$ and $k\ge0$, and let $M$, $M_1$, $\mathfrak{m}$ be positive
integers with $M$ a multiple of $M_1$; in the application $M$, $M_1$, $\mathfrak{m}$ are powers of
$L$. Let $\ell_{\mathbb{T}}>0$ and let $\mathbb{T}:=(\mathbb{R}/\ell_{\mathbb{T}}\mathbb{Z})^d$, the
torus of side $\ell_{\mathbb{T}}$. We write $|x-y|$ for the sup-norm distance
on $\mathbb{T}$, $\operatorname{dist}(A,B)$ for the corresponding distance of sets, and
$B(x,r):=\{y:|x-y|<r\}$. For $0\le j\le k$ put
\begin{equation*}
M_j:=ML^j\eta,\qquad s_j:=M_1L^j\eta,\qquad u_j:=\mathfrak{m}L^j\eta .
\end{equation*}
For $0\le j\le k$, $\pi_j$ is a tiling of $\mathbb{T}$ by closed cubes of side $M_j$ of product form:
for each axis $\mu$ there is a finite set $\mathsf{g}_j^{\mu}\subset\mathbb{R}/\ell_{\mathbb{T}}\mathbb{Z}$
of equally spaced points, with spacing $M_j$, and the cubes of $\pi_j$ are the products of the closed
arcs between consecutive points.

\emph{The ambient sequence.} $\boldsymbol{\Omega}=(\Omega_j)_{0\le j\le k}$ is a sequence of closed
subsets of $\mathbb{T}$ with $\mathbb{T}=\Omega_0\supseteq\Omega_1\supseteq\dots\supseteq\Omega_k$;
we put $\Omega_{k+1}:=\emptyset$ and, for $A\subset\mathbb{T}$,
$A^{\complement}:=\operatorname{cl}(\mathbb{T}\setminus A)$. The grids and the sequence are assumed
to satisfy
\begin{equation}\label{9.5:eq-scaled-length-a}
\begin{gathered}
\mathsf{g}_{j+1}^{\mu}\subset\mathsf{g}_j^{\mu}\qquad(1\le\mu\le d,\ 0\le j<k),\\
\Omega_j\ \text{is a union of cubes of}\ \pi_j\qquad(0\le j\le k),\\
\Omega_j^{\complement}\cap\Omega_{j+1}=\emptyset\qquad(0\le j<k).
\end{gathered}
\end{equation}
We refer to these as the nesting, tiling and separation conditions. By the nesting condition every
cube of $\pi_{j+1}$ is the union of $L^d$ cubes of $\pi_j$ (the partitions $\pi_j$ into closed cubes
of size $M$ of \cite{B-CMP109} are here nested $L$-adically). The tiling condition holds when each
$\Omega_j$ is a localisation domain of scale $j$ as in \cite{B-CMP119}, and for the sequences
completed in \cite{B-CMP122a}; the separation condition is implied by the positive distance between
$\mathbb{T}\setminus\Omega_j$ and $\Omega_{j+1}$ required in \cite[(2.2)]{B-CMP96}. For the
sequences of large-field labels used in this book, the admissibility conditions imposed on those
labels give more than the separation condition: $\Omega_j^{\complement}$ and $\Omega_{j+1}$ are
separated by at least one layer of cubes of $\pi_{j+1}$ ($0\le j<k$). If the lowest set of a given
sequence is a proper subset of $\mathbb{T}$, we put $\Omega_0:=\mathbb{T}$: this adds cells of level
$0$ and admits more paths, so every count of cells made with the definitions below only grows and
every distance only shrinks, and statements for the original objects follow a fortiori.

\emph{Level and cells.} The \emph{ambient level} of $x\in\mathbb{T}$ is
$m(x):=\max\{j:x\in\Omega_j\}$. A \emph{cell of level $j$} is a cube $\Delta\in\pi_j$ with
$\Delta\subset\Omega_j$ and $\Delta\not\subset\Omega_{j+1}$; we write $\operatorname{lev}(\Delta):=j$,
and $\pi^{\mathrm{mix}}$ denotes the set of all cells (these are the cubes in proper scales of
\cite[p.~276]{B-CMP119}). Two cells \emph{touch} if they intersect; they \emph{share a wall} if
their intersection has dimension $d-1$.

\emph{Scaled length.} A \emph{path} is a Lipschitz map $\gamma:[a,b]\to\mathbb{T}$, and
\begin{equation*}
\mathsf{L}(\gamma):=\int_a^b|\dot{\gamma}(t)|\,\bigl(L^{m(\gamma(t))}\eta\bigr)^{-1}\,dt ,
\end{equation*}
where $|\dot{\gamma}|$ is the sup-norm of the derivative of a lift of $\gamma$ to $\mathbb{R}^d$.
A \emph{graph} is a finite non-empty family $\Gamma=(\gamma_e)_{e\in E}$ of paths (constant paths
allowed) whose union $|\Gamma|\subset\mathbb{T}$ is connected, and
$\mathsf{L}(\Gamma):=\sum_{e\in E}\mathsf{L}(\gamma_e)$. The tree graphs in the continuous space of
\cite{B-CMP109}, the contours of \cite[(2.46)]{B-CMP96} and their concatenations are graphs. We put
$\mathsf{d}(x,y):=\inf\{\mathsf{L}(\gamma):\gamma\ \text{a path from}\ x\ \text{to}\ y\}$,
$\mathsf{d}(A,B):=\inf\{\mathsf{d}(x,y):(x,y)\in A\times B\}$ and
$\mathsf{B}(z,r):=\{x:\mathsf{d}(z,x)<r\}$.

\emph{Seen cells and loose tree length.} For a graph $\Gamma$ and $\rho>0$ put
\begin{equation*}
\mathcal{S}_{\rho}(\Gamma):=\{\Delta\in\pi^{\mathrm{mix}}:\mathsf{d}(\Delta,|\Gamma|)<\rho\},
\qquad
\mathcal{S}_{0}(\Gamma):=\{\Delta\in\pi^{\mathrm{mix}}:\Delta\cap|\Gamma|\neq\emptyset\}.
\end{equation*}
For a finite $\mathsf{F}\subset\pi^{\mathrm{mix}}$ and $\rho\ge0$ put
$\mathsf{t}_{\rho}(\mathsf{F}):=\inf\{\mathsf{L}(\Gamma):\mathsf{F}\subset\mathcal{S}_{\rho}(\Gamma)\}$
and $d^{\circ}_{\mathrm{mix}}(\mathsf{F}):=\mathsf{t}_0(\mathsf{F})/M$. When $k=0$, $m$ vanishes
identically and $\mathsf{L}$ is the sup-norm length in units of $\eta$, so
$d^{\circ}_{\mathrm{mix}}(\mathsf{F})$ is the infimum of the sup-norm lengths, in units of the cube
side $M\eta$, of graphs meeting every cube of $\mathsf{F}$; up to the choice of norm, this is the
one-scale tree length used in this book for a single level.

\emph{Numbers.} In this section $r_0$ is a local letter, unrelated to the depth offset of clause
(iv). We put
\begin{equation*}
\begin{gathered}
r_0:=\frac{M}{4L},\qquad
\nu(\Gamma):=\max\Bigl\{1,\Bigl\lfloor\frac{2\mathsf{L}(\Gamma)}{r_0}\Bigr\rfloor\Bigr\},\qquad
D:=(L+2)^d-L^d,\\
C_d^{\mathrm{mix}}:=2^{d+3}L,\qquad c_d^{\mathrm{mix}}:=(2D^2)^{2^{d+3}L}.
\end{gathered}
\end{equation*}
Then:
\begin{itemize}
\item[(a)] $m$ is upper semicontinuous, with $\{m\ge j\}=\Omega_j$, and $m\circ\gamma$ is Borel for
every path $\gamma$;
\item[(b)] $\mathcal{S}_0(\Gamma)\subset\mathcal{S}_{\rho}(\Gamma)$ for every $\rho>0$, and
$\nu(\Gamma)=\max\{1,\lfloor 8L\mathsf{L}(\Gamma)/M\rfloor\}$;
\item[(c)] for all $x,y\in\mathbb{T}$,
$|x-y|/(L^k\eta)\le\mathsf{d}(x,y)\le|x-y|/\eta$; hence $\mathsf{d}$ is a metric inducing the
topology of $\mathbb{T}$, and $\mathsf{B}(z,r)$ is open in $\mathbb{T}$;
\item[(d)] if $x\in\mathsf{B}(z,r)$, some path $\gamma$ from $z$ to $x$ has $\mathsf{L}(\gamma)<r$,
and every point of $\gamma$ lies in $\mathsf{B}(z,r)$.
\end{itemize}
\end{definition}

\begin{proof}
We first verify the consequence of the nesting condition stated after
\eqref{9.5:eq-scaled-length-a}. Since $M_{j+1}=LM_j$ and $\mathsf{g}_j^{\mu}$ is equally
spaced with spacing $M_j$ and contains $\mathsf{g}_{j+1}^{\mu}$, each closed arc between consecutive
points of $\mathsf{g}_{j+1}^{\mu}$ is the union of exactly $L$ closed arcs between consecutive
points of $\mathsf{g}_j^{\mu}$; taking products over the $d$ axes, each cube of $\pi_{j+1}$ is the
union of $L^d$ cubes of $\pi_j$.

(a) Since $\Omega_0=\mathbb{T}$, the maximum defining $m(x)$ is taken over a non-empty set, so $m$
is defined on $\mathbb{T}$ with $0\le m\le k$. Because the sequence is decreasing, $m(x)\ge j$ holds
if and only if $x\in\Omega_j$; thus $\{m\ge j\}=\Omega_j$, which is closed, and $m$ is upper
semicontinuous. For a path $\gamma$, the sets $\{m\circ\gamma\ge j\}=\gamma^{-1}(\Omega_j)$ are
closed because $\gamma$ is continuous, so $m\circ\gamma$ is upper semicontinuous, hence Borel, and
the integrand defining $\mathsf{L}(\gamma)$ is measurable; since $0\le m\le k$ it lies between
$|\dot{\gamma}|/(L^k\eta)$ and $|\dot{\gamma}|/\eta$, so $\mathsf{L}(\gamma)$ is finite.

(b) If $\Delta\cap|\Gamma|\neq\emptyset$, then $\mathsf{d}(\Delta,|\Gamma|)\le\mathsf{d}(x,x)=0<\rho$
for a point $x$ of the intersection (the constant path at $x$ has $\mathsf{L}=0$). The identity for
$\nu$ holds because $2/r_0=8L/M$.

(c) Let $\gamma:[a,b]\to\mathbb{T}$ be a path from $x$ to $y$, with lift $\tilde{\gamma}$. As
$m\le k$,
\begin{equation*}
\mathsf{L}(\gamma)\ge\frac{1}{L^k\eta}\int_a^b|\dot{\tilde{\gamma}}(t)|\,dt
\ge\frac{|\tilde{\gamma}(b)-\tilde{\gamma}(a)|}{L^k\eta}\ge\frac{|x-y|}{L^k\eta},
\end{equation*}
the last step because the sup-norm distance on $\mathbb{T}$ is the minimum over lifts. Taking the
infimum over $\gamma$ gives the lower bound. For the upper bound, choose lifts $\tilde{x},\tilde{y}$
with $|\tilde{x}-\tilde{y}|=|x-y|$ and let $\gamma$ be the projection of the segment
$t\mapsto\tilde{x}+t(\tilde{y}-\tilde{x})$, $t\in[0,1]$; as $m\ge0$ and $L\ge2$,
$\mathsf{L}(\gamma)\le\eta^{-1}\int_0^1|\tilde{y}-\tilde{x}|\,dt=|x-y|/\eta$. Consequently
$\mathsf{d}$ is finite, and $\mathsf{d}(x,y)=0$ only if $x=y$. Reversing a path,
$t\mapsto\gamma(a+b-t)$, does not change $\mathsf{L}$, so $\mathsf{d}$ is symmetric; the
concatenation of a path from $x$ to $y$ and a path from $y$ to $w$, reparametrised on consecutive
intervals, is a path from $x$ to $w$ whose scaled length is the sum, so $\mathsf{d}$ satisfies the
triangle inequality. Thus $\mathsf{d}$ is a metric, and by the two-sided bound it is bi-Lipschitz
equivalent to the sup-norm distance, which induces the topology of $\mathbb{T}$; hence so does
$\mathsf{d}$, and the ball $\mathsf{B}(z,r)$ of $\mathsf{d}$ is open in $\mathbb{T}$.

(d) If $\mathsf{d}(z,x)<r$, by the definition of the infimum there is a path $\gamma:[a,b]\to
\mathbb{T}$ from $z$ to $x$ with $\mathsf{L}(\gamma)<r$. For $t\in[a,b]$ the restriction of $\gamma$
to $[a,t]$ is a path from $z$ to $\gamma(t)$, and since the integrand is non-negative its scaled
length is at most $\mathsf{L}(\gamma)<r$; hence $\mathsf{d}(z,\gamma(t))<r$, that is,
$\gamma(t)\in\mathsf{B}(z,r)$.
\end{proof}

\begin{lemma}[The gap between zones is automatic]\label{9.5:lem-zone-gap}
Let $A,A'\subset\mathbb{T}$ be disjoint closed sets, each a union of cubes of $\pi_j$. Then
\begin{equation*}
\operatorname{dist}(A,A')\ge M_j .
\end{equation*}
Hence, under the conditions \textup{(H1)}--\textup{(H3)} of
Definition~\ref{9.5:def-scaled-length}, displayed in \eqref{9.5:eq-scaled-length-a},
\begin{equation}\label{9.5:eq-zone-gap-1}
\operatorname{dist}\bigl(\Omega_j^{\complement},\Omega_{j+1}\bigr)\ge M_j
\qquad(0\le j<k),
\end{equation}
that is, every zone is at least one cell of its own scale thick.
\end{lemma}

\begin{proof}
We first prove the general claim. Let $x\in A$ and $x'\in A'$. Since $A$ and $A'$ are unions of
cubes of $\pi_j$, there are cubes $Q\subset A$ and $Q'\subset A'$ of $\pi_j$ with $x\in Q$ and
$x'\in Q'$. Since $\pi_j$ is of product form, we may write $Q=\prod_{\mu}I_{\mu}$ and
$Q'=\prod_{\mu}I'_{\mu}$, where for each axis $\mu$ the sets $I_{\mu}$ and $I'_{\mu}$ are closed
arcs of length $M_j$ of the circle $\mathbb{R}/\ell_{\mathbb{T}}\mathbb{Z}$, each joining two
consecutive points of the grid $\mathsf{g}_j^{\mu}$.

Suppose that for every $\mu$ the arcs $I_{\mu}$ and $I'_{\mu}$ were either equal or consecutive.
Two equal closed arcs intersect, and two consecutive closed arcs share an end point; so for each
$\mu$ we could pick a point $y_{\mu}\in I_{\mu}\cap I'_{\mu}$, and then $y=(y_{\mu})_{\mu}$ would lie
in $Q\cap Q'\subset A\cap A'$, contradicting the disjointness of $A$ and $A'$. Hence there is an
axis $\mu$ for which $I_{\mu}$ and $I'_{\mu}$ are neither equal nor consecutive. For this $\mu$,
each of the two arcs of the circle that join $I_{\mu}$ to $I'_{\mu}$ contains a whole arc of the
grid $\mathsf{g}_j^{\mu}$, of length $M_j$, lying between them. Since $x_{\mu}\in I_{\mu}$ and
$x'_{\mu}\in I'_{\mu}$, each of the two arcs of the circle joining $x_{\mu}$ to $x'_{\mu}$ has length
at least $M_j$, so the distance on the circle satisfies $|x_{\mu}-x'_{\mu}|\ge M_j$. As $|x-x'|$
is the sup-norm distance on $\mathbb{T}$,
\begin{equation*}
|x-x'|\ge|x_{\mu}-x'_{\mu}|\ge M_j .
\end{equation*}
Since $x\in A$ and $x'\in A'$ were arbitrary, $\operatorname{dist}(A,A')\ge M_j$.

For \eqref{9.5:eq-zone-gap-1} we apply the general claim with $A=\Omega_j^{\complement}$ and
$A'=\Omega_{j+1}$. By \textup{(H2)}, $\Omega_j^{\complement}$, the closure of the complement of
$\Omega_j$ in $\mathbb{T}$, is the union of the cubes of $\pi_j$ not contained in $\Omega_j$; it
is closed by definition. By \textup{(H2)} again, $\Omega_{j+1}$ is a union of cubes of
$\pi_{j+1}$, and by \textup{(H1)} every cube of $\pi_{j+1}$ is a union of cubes of $\pi_j$;
hence $\Omega_{j+1}$ is a union of cubes of $\pi_j$, and it is closed, being a finite union of
closed cubes. Finally, by \textup{(H3)} the sets $\Omega_j^{\complement}$ and $\Omega_{j+1}$ are
disjoint. The general claim therefore gives
$\operatorname{dist}(\Omega_j^{\complement},\Omega_{j+1})\ge M_j$ for $0\le j<k$.
\end{proof}

\begin{lemma}[Cells tile the torus, graded by level]\label{9.5:lem-cells-tiling}
Let the torus $\mathbb{T}$, the grids $\pi_j$, the sequence $\boldsymbol{\Omega}=(\Omega_j)_{0\le j\le k}$,
the set of cells $\pi^{\mathrm{mix}}$, the level $\operatorname{lev}(\Delta)$ of a cell $\Delta$, the
ambient level $m(x)=\max\{j : x\in\Omega_j\}$ and the constant $D=(L+2)^d-L^d$ be as in
Definition~\ref{9.5:def-scaled-length}. Then the following hold.
\begin{itemize}
\item[(a)] $\pi^{\mathrm{mix}}$ is a tiling of $\mathbb{T}$: the cells cover $\mathbb{T}$ and have pairwise
disjoint interiors. A cell of level $j$ is a union of cubes of $\pi_i$ for every $i\le j$. Every cube of
$\pi_i$ is contained in at most one cell of level $\ge i$.
\item[(b)] $m=\operatorname{lev}(\Delta)$ on $\operatorname{int}\Delta$, and
$m(x)\in\{\operatorname{lev}\Delta,\operatorname{lev}\Delta+1\}$ for $x\in\Delta$.
\item[(c)] Touching cells differ in level by at most $1$.
\item[(d)] A cell touches at most $D$ other cells.
\end{itemize}
\end{lemma}

\begin{proof}
We write $\operatorname{int}A$ for the interior of $A\subset\mathbb{T}$ and, as in
Definition~\ref{9.5:def-scaled-length}, $A^{\complement}$ for the closure of the complement of $A$ in
$\mathbb{T}$. We treat the four clauses in order.

(a) Let $x\in\mathbb{T}$ and put $j:=m(x)$, so that $x\in\Omega_j$ and $x\notin\Omega_{j+1}$. By the
second of the conditions \eqref{9.5:eq-scaled-length-a}, $x$ lies in a cube $Q\in\pi_j$ with
$Q\subset\Omega_j$. Since $x\in Q$ and $x\notin\Omega_{j+1}$, we have $Q\not\subset\Omega_{j+1}$. Hence
$Q$ is a cell of level $j$, and the cells cover $\mathbb{T}$.

Next let $\Delta$ be a cell of level $j$. We claim that
\begin{equation}\label{9.5:eq-cells-tiling-1}
\operatorname{int}\Delta\cap\Omega_{j+1}=\emptyset .
\end{equation}
Indeed, $\Omega_{j+1}$ is a union of cubes of $\pi_j$ by \eqref{9.5:eq-scaled-length-a}. None of these
cubes is $\Delta$, because $\Delta\not\subset\Omega_{j+1}$. Two distinct cubes of the tiling $\pi_j$ have
disjoint interiors, and both are closed cubes of the same side, so each cube of $\pi_j$ other than
$\Delta$ misses $\operatorname{int}\Delta$; this gives \eqref{9.5:eq-cells-tiling-1}.

Now let $\Delta\neq\Delta'$ be two cells. If they have the same level, they are distinct cubes of one
tiling and so have disjoint interiors. If $\operatorname{lev}\Delta=i<j=\operatorname{lev}\Delta'$, then
$\Delta'\subset\Omega_j\subset\Omega_{i+1}$, because the sequence $\boldsymbol{\Omega}$ is decreasing and
$j\ge i+1$. By \eqref{9.5:eq-cells-tiling-1} applied to $\Delta$, the set $\Delta'$ misses
$\operatorname{int}\Delta$; in particular the interiors are disjoint.

That a cell of level $j$, which is a cube of $\pi_j$, is a union of cubes of $\pi_i$ for every $i\le j$
is the first of the conditions \eqref{9.5:eq-scaled-length-a}.

For the last clause, suppose a cube $Q\in\pi_i$ were contained in two distinct cells of level $\ge i$.
Every interior point of $Q$ would then be a common interior point of the two cells. This contradicts the
disjointness of interiors just proved.

(b) Let $\Delta$ be a cell of level $j$. Then $\operatorname{int}\Delta\subset\Delta\subset\Omega_j$, and
$\operatorname{int}\Delta\cap\Omega_{j+1}=\emptyset$ by \eqref{9.5:eq-cells-tiling-1}. Thus
$\operatorname{int}\Delta\subset\Omega_j\setminus\Omega_{j+1}$, and $m=j$ on $\operatorname{int}\Delta$.

Let now $x\in\Delta$. First, $x\in\Omega_j$, so $m(x)\ge j$. Second, $\Delta$ is a closed cube and hence
the closure of its interior. Since $\operatorname{int}\Delta$ lies in the complement of $\Omega_{j+1}$,
its closure lies in $\Omega_{j+1}^{\complement}$, so $x\in\Omega_{j+1}^{\complement}$. By the third of the
conditions \eqref{9.5:eq-scaled-length-a}, $\Omega_{j+1}^{\complement}$ misses $\Omega_{j+2}$. Hence
$x\notin\Omega_{j+2}$ and $m(x)\le j+1$. Together, $m(x)\in\{j,j+1\}$.

(c) Let $\Delta$ and $\Delta'$ be touching cells and let $x\in\Delta\cap\Delta'$. By (b), applied to
$\Delta$ and to $\Delta'$ at the point $x$, both $\operatorname{lev}\Delta$ and $\operatorname{lev}\Delta'$
lie in $\{m(x)-1,m(x)\}$. Hence they differ by at most $1$.

(d) Let first $\operatorname{lev}\Delta=j\ge1$. Let $\Delta'\neq\Delta$ be a cell touching $\Delta$ at a
point $x\in\Delta\cap\Delta'$.
\begin{itemize}
\item By (c), $\operatorname{lev}\Delta'\ge j-1$.
\item By (a), $\Delta'$ is therefore a union of cubes of $\pi_{j-1}$, and one of them, call it $Q'$,
contains $x$.
\item $Q'\not\subset\Delta$: otherwise
$\operatorname{int}Q'\subset\Delta\cap\Delta'$, and $\Delta,\Delta'$ would share an interior point,
contrary to (a).
\end{itemize}
Thus $Q'$ is a cube of $\pi_{j-1}$ that intersects $\Delta$ (it contains $x$) and is not contained in
$\Delta$.

We count such cubes. By (a), $\Delta$ is a union of cubes of $\pi_{j-1}$. Its side is
$M_j=L\,M_{j-1}$, so it is the union of exactly $L^d$ of them. The grid $\pi_{j-1}$ is of product form
with intervals of length $M_{j-1}$ on each axis. On each axis the projection of $\Delta$ is the union of
$L$ consecutive grid intervals, and the grid intervals meeting it are these $L$ and at most the two
adjacent ones. Hence at most $(L+2)^d$ cubes of $\pi_{j-1}$ intersect $\Delta$, and $L^d$ of them lie in
$\Delta$. At most $(L+2)^d-L^d=D$ cubes can therefore serve as $Q'$.

Distinct cells $\Delta'$ contain distinct cubes $Q'$. Indeed, by the last clause of (a) a cube of
$\pi_{j-1}$ is contained in at most one cell of level $\ge j-1$. Hence $\Delta$ touches at most $D$ other
cells.

Let now $\operatorname{lev}\Delta=0$. Every cell is a union of cubes of $\pi_0$ by (a). So each cell
$\Delta'\neq\Delta$ touching $\Delta$ at $x$ contains a cube $Q'\in\pi_0$ with $x\in Q'$. Moreover
$Q'\neq\Delta$: otherwise $\operatorname{int}\Delta\subset\Delta\cap\Delta'$, contrary to (a). The cubes
of $\pi_0$ other than $\Delta$ that intersect the cube $\Delta\in\pi_0$ number at most $3^d-1$, and
distinct $\Delta'$ contain distinct $Q'$ by the last clause of (a). Finally, $3^d-1\le D$: the quantity
$(L+2)^d-L^d$ is nondecreasing in $L$ and equals $3^d-1$ at $L=1$. This proves (d).
\end{proof}

\begin{lemma}[Crossing a zone costs one cube size]\label{9.5:lem-zone-crossing}
Let $\boldsymbol{\Omega}=(\Omega_j)_{0\le j\le k}$ and the scaled length $\mathsf{L}$ be as in
Definition~\ref{9.5:def-scaled-length}. Let $0\le j<k$ and let $\gamma$ be a path from a point of
$\Omega_j^{\complement}$ to a point of $\Omega_{j+1}$. Then $\mathsf{L}(\gamma)\ge M$.
\end{lemma}

\begin{proof}
Write $\gamma\colon[a,b]\to\mathbb{T}$, with $\gamma(a)\in\Omega_j^{\complement}$ and
$\gamma(b)\in\Omega_{j+1}$. The set $\Omega_{j+1}$ is closed, being a union of closed cubes. Hence
$\{t\in[a,b]:\gamma(t)\in\Omega_{j+1}\}$ is closed, and it contains $b$. We put
\begin{equation*}
t_2:=\min\{t\in[a,b]:\gamma(t)\in\Omega_{j+1}\}.
\end{equation*}
The set $\Omega_j^{\complement}$ is closed, since it is by definition the closure of the
complement of $\Omega_j$. Hence $\{t\le t_2:\gamma(t)\in\Omega_j^{\complement}\}$ is closed, and it
contains $a$. We put
\begin{equation*}
t_1:=\max\{t\in[a,t_2]:\gamma(t)\in\Omega_j^{\complement}\}.
\end{equation*}
By Lemma~\ref{9.5:lem-zone-gap}, $\operatorname{dist}(\Omega_j^{\complement},\Omega_{j+1})\ge M_j>0$.
Since $\gamma(t_1)\in\Omega_j^{\complement}$ and $\gamma(t_2)\in\Omega_{j+1}$, we have
$\gamma(t_1)\ne\gamma(t_2)$, and therefore $t_1<t_2$.

Let $t\in\,]t_1,t_2[$. By the maximality of $t_1$, $\gamma(t)\notin\Omega_j^{\complement}$. Every point
outside $\Omega_j$ lies in the closure of the complement of $\Omega_j$, so $\gamma(t)\in\Omega_j$. By
the minimality of $t_2$, $\gamma(t)\notin\Omega_{j+1}$. The sequence $\boldsymbol{\Omega}$ is
decreasing, so $\gamma(t)\notin\Omega_i$ for every $i\ge j+1$. Hence the ambient level of $\gamma(t)$
is $m(\gamma(t))=j$.

By Definition~\ref{9.5:def-scaled-length}, $\mathsf{L}(\gamma)$ integrates the sup-norm speed
$|\dot\gamma|$ divided by the ambient scale $L^{m(\gamma(t))}\eta$. The integrand is non-negative,
and on $]t_1,t_2[$ the ambient scale equals $L^j\eta$. Restricting the integral to $[t_1,t_2]$
therefore gives
\begin{align*}
\mathsf{L}(\gamma)
&\ge (L^j\eta)^{-1}\int_{t_1}^{t_2}|\dot\gamma(t)|\,dt
\ge (L^j\eta)^{-1}\,|\gamma(t_2)-\gamma(t_1)|\\
&\ge (L^j\eta)^{-1}\operatorname{dist}(\Omega_j^{\complement},\Omega_{j+1})
\ge \frac{M_j}{L^j\eta}=M .
\end{align*}
The second inequality holds because the length of a path bounds the sup-norm distance of its
endpoints on $\mathbb{T}$. The third holds because $\gamma(t_1)\in\Omega_j^{\complement}$ and
$\gamma(t_2)\in\Omega_{j+1}$. The fourth is Lemma~\ref{9.5:lem-zone-gap}. The final equality is
$M_j=ML^j\eta$.
\end{proof}

\begin{lemma}[Cells meeting a small scaled ball]\label{9.5:lem-small-balls}
Work in the setting of Definition~\ref{9.5:def-scaled-length}: the label
$\boldsymbol{\Omega}=(\Omega_j)_{0\le j\le k}$, the ambient level $m(x)=\max\{j:\ x\in\Omega_j\}$,
the scaled metric $\mathsf{d}$ with its open balls $\mathsf{B}(z,r)$, the sup-norm balls $B(x,r)$,
the tilings $\pi_j$ and the cells of $\pi^{\mathrm{mix}}$. Let $z\in\mathbb{T}$ and
$0<r\le M/(2L)$, and put
\begin{equation*}
\mathsf{B}:=\mathsf{B}(z,r),\qquad j_*:=\min\{m(x):\ x\in\mathsf{B}\}.
\end{equation*}
Then:
\begin{enumerate}
\item[(i)] $m\le j_*+1$ on $\mathsf{B}$;
\item[(ii)] $\mathsf{B}\subset B\bigl(z,\tfrac12 M_{j_*}\bigr)$;
\item[(iii)] every cell meeting $\mathsf{B}$ has level at least $j_*$;
\item[(iv)] at most $2^d$ cells meet $\mathsf{B}$, and they have a common point.
\end{enumerate}
\end{lemma}

\begin{proof}
The ball $\mathsf{B}$ contains $z$ and $m$ takes values in $\{0,\dots,k\}$, so the minimum
defining $j_*$ is attained.

(i) Let $x_*\in\mathsf{B}$ with $m(x_*)=j_*$, and suppose that some $x^*\in\mathsf{B}$ has
$m(x^*)\ge j_*+2$. Since $x_*,x^*\in\mathsf{B}$, there are paths from $z$ to $x_*$ and from
$z$ to $x^*$, each of scaled length less than $r$. Running the first backwards and then the
second gives a path $\gamma$ from $x_*$ to $x^*$ through $z$ with
\begin{equation*}
\mathsf{L}(\gamma)<2r\le \frac{M}{L}\le M .
\end{equation*}
By the definition of $m$, $m(x_*)=j_*$ gives $x_*\notin\Omega_{j_*+1}$, so
$x_*\in\Omega_{j_*+1}^{\complement}$; and $m(x^*)\ge j_*+2$ gives
$x^*\in\Omega_{m(x^*)}\subset\Omega_{j_*+2}$, the sequence being decreasing. Moreover
$j_*+2\le m(x^*)\le k$, so $0\le j_*+1<k$. Lemma~\ref{9.5:lem-zone-crossing}, applied with
$j=j_*+1$ to $\gamma$, gives $\mathsf{L}(\gamma)\ge M$, a contradiction. Hence
$m\le j_*+1$ on $\mathsf{B}$.

(ii) Let $x\in\mathsf{B}$. Since $\mathsf{d}(z,x)<r$, there is a path $\gamma$ from $z$ to $x$
with $\mathsf{L}(\gamma)<r$; every point of $\gamma$ is joined to $z$ by an initial segment of
$\gamma$, of scaled length at most $\mathsf{L}(\gamma)<r$, so $\gamma$ lies inside
$\mathsf{B}$. By (i), $m\le j_*+1$ along $\gamma$, so the weight
$(L^{m}\eta)^{-1}$ by which the scaled length multiplies the sup-norm speed is at least
$(L^{j_*+1}\eta)^{-1}$ on $\gamma$. Therefore
\begin{equation*}
|x-z|\le\int|\dot\gamma|\le L^{j_*+1}\eta\,\mathsf{L}(\gamma)
<L^{j_*+1}\eta\,\frac{M}{2L}=\tfrac12 M_{j_*},
\end{equation*}
that is, $x\in B(z,\tfrac12 M_{j_*})$.

(iii) Let $\Delta$ be a cell and $x\in\Delta\cap\mathsf{B}$. By
Lemma~\ref{9.5:lem-cells-tiling}(a), $\Delta$ is a finite union of closed cubes, hence the
closure of its interior, so $x$ lies in the closure of $\operatorname{int}\Delta$. The ball
$\mathsf{B}$ is open in $\mathbb{T}$, so some $x'\in\operatorname{int}\Delta$ lies in
$\mathsf{B}$. By Lemma~\ref{9.5:lem-cells-tiling}(b), $m(x')=\operatorname{lev}(\Delta)$, and
therefore $j_*\le m(x')=\operatorname{lev}(\Delta)$.

(iv) Fix an axis $\mu$ and consider the open arc
\begin{equation*}
J_\mu:=\bigl]z_\mu-\tfrac12 M_{j_*},\,z_\mu+\tfrac12 M_{j_*}\bigr[ ,
\end{equation*}
of length $M_{j_*}$, which is the spacing of the grid points $\mathsf{g}_{j_*}^{\mu}$. Hence
$J_\mu$ contains at most one point of $\mathsf{g}_{j_*}^{\mu}$. If it contains a point
$p_\mu\in\mathsf{g}_{j_*}^{\mu}$, then $J_\mu\subset\,]p_\mu-M_{j_*},p_\mu+M_{j_*}[$, and the
closed grid arcs meeting $J_\mu$ are the two that end at $p_\mu$. If it contains none, then
$J_\mu$, being connected and free of grid points, lies in the interior of one closed grid arc,
which is the only grid arc it meets; in this case we put $p_\mu:=z_\mu$. In either case at
most two closed grid arcs of $\pi_{j_*}$ on the axis $\mu$ meet $J_\mu$, and each of them
contains $p_\mu$.

The sup-norm ball $B(z,\tfrac12 M_{j_*})$ is the product of the arcs $J_\mu$, and the cubes of
$\pi_{j_*}$ are products of closed grid arcs, one on each axis; such a cube meets the ball if
and only if each of its factors meets the corresponding $J_\mu$. Hence at most $2^d$ cubes of
$\pi_{j_*}$ meet $B(z,\tfrac12 M_{j_*})$, and all of them contain the point
$p:=(p_\mu)_{\mu}$.

Now let $\Delta$ be a cell meeting $\mathsf{B}$ at a point $x$. By (iii),
$\operatorname{lev}(\Delta)\ge j_*$, so by Lemma~\ref{9.5:lem-cells-tiling}(a) $\Delta$ is a
union of cubes of $\pi_{j_*}$; choose one of them containing $x$. Since $x\in\mathsf{B}$, (ii)
shows that this cube meets $B(z,\tfrac12 M_{j_*})$, so it is one of the at most $2^d$ cubes
just described, and it contains $p$. By Lemma~\ref{9.5:lem-cells-tiling}(a), a cube of
$\pi_{j_*}$ is contained in at most one cell of level at least $j_*$, so distinct cells meeting
$\mathsf{B}$ contain distinct such cubes. Therefore at most $2^d$ cells meet $\mathsf{B}$, and
each of them contains $p$.
\end{proof}

\begin{lemma}[Diameter and nets of a graph]\label{9.5:lem-graph-nets}
Let $\Gamma$ be a graph, with the scaled length $\mathsf{L}(\Gamma)$, the scaled path metric
$\mathsf{d}$ and its open balls $\mathsf{B}(z,r)$ of Definition~\ref{9.5:def-scaled-length}.
\begin{itemize}
\item[(a)] $\mathsf{d}(p,q)\le\mathsf{L}(\Gamma)$ for all $p,q\in|\Gamma|$.
\item[(b)] For every $r>0$ there are points $z_1,\dots,z_N\in|\Gamma|$ with
\begin{equation*}
N\le\max\bigl\{1,\lfloor 2\mathsf{L}(\Gamma)/r\rfloor\bigr\},
\qquad
|\Gamma|\subset\bigcup_{a=1}^{N}\mathsf{B}(z_a,r).
\end{equation*}
\end{itemize}
\end{lemma}

\begin{proof}
We write the graph as the finite family of Lipschitz paths
$\gamma_e\colon[a_e,b_e]\to\mathbb{T}$, $e\in E$, whose union $|\Gamma|$ is connected, as in
Definition~\ref{9.5:def-scaled-length}, and we write $\operatorname{Leb}$ for Lebesgue measure
on $\mathbb{R}$. For $p\in\mathbb{T}$ put $f:=\mathsf{d}(p,\cdot)$ and, for $e\in E$,
\begin{equation*}
g_e:=f\circ\gamma_e,
\qquad
w_e(t):=|\dot\gamma_e(t)|\,\bigl(L^{m(\gamma_e(t))}\eta\bigr)^{-1},
\end{equation*}
so that the scaled length of $\gamma_e$ is $\int_{a_e}^{b_e}w_e$ and
$\mathsf{L}(\Gamma)=\sum_{e\in E}\int_{a_e}^{b_e}w_e$.

For $a_e\le t\le t'\le b_e$ the triangle inequality for $\mathsf{d}$ gives
$|g_e(t)-g_e(t')|\le\mathsf{d}(\gamma_e(t),\gamma_e(t'))$, and the restriction of $\gamma_e$ to
$[t,t']$ is a path from $\gamma_e(t)$ to $\gamma_e(t')$ of scaled length $\int_t^{t'}w_e$, so
\begin{equation}\label{9.5:eq-graph-nets-1}
|g_e(t)-g_e(t')|\le\mathsf{d}(\gamma_e(t),\gamma_e(t'))\le\int_t^{t'}w_e .
\end{equation}
In particular $g_e$ is continuous. Hence, for an interval $I\subset[a_e,b_e]$, the set
$g_e(I)$ is an interval, and by \eqref{9.5:eq-graph-nets-1} its length, the supremum of
$|g_e(t)-g_e(t')|$ over $t,t'\in I$, is at most $\int_I w_e$. A relatively open set
$U\subset[a_e,b_e]$ is a countable disjoint union of intervals $I_i$; then
$g_e(U)=\bigcup_i g_e(I_i)$, and by countable subadditivity
\begin{equation}\label{9.5:eq-graph-nets-2}
\operatorname{Leb}\bigl(g_e(U)\bigr)\le\sum_i\operatorname{Leb}\bigl(g_e(I_i)\bigr)
\le\sum_i\int_{I_i}w_e=\int_U w_e .
\end{equation}

(a) Let $p,q\in|\Gamma|$ and $f=\mathsf{d}(p,\cdot)$. The set $|\Gamma|$ is compact, being a
finite union of continuous images of compact intervals, and it is connected. The restriction of
$f$ to $|\Gamma|$ is continuous: for a closed set $C\subset\mathbb{R}$ one has
$f^{-1}(C)\cap|\Gamma|=\bigcup_e\gamma_e\bigl(g_e^{-1}(C)\bigr)$, a finite union of compact
sets, since each $g_e$ is continuous. Therefore $f(|\Gamma|)$ is a compact interval; it
contains $f(p)=0$ and $f\ge0$, so $f(|\Gamma|)=[0,F]$ for some $F\ge f(q)$. Moreover
$[0,F]=f(|\Gamma|)=\bigcup_e g_e([a_e,b_e])$, and \eqref{9.5:eq-graph-nets-2} with
$U=[a_e,b_e]$ yields
\begin{equation*}
\mathsf{d}(p,q)=f(q)\le F=\operatorname{Leb}([0,F])
\le\sum_{e\in E}\int_{a_e}^{b_e}w_e=\mathsf{L}(\Gamma).
\end{equation*}

(b) Let $\{z_a\}_{a\le N}\subset|\Gamma|$ be a maximal set with $\mathsf{d}(z_a,z_b)\ge r$ for
$a\ne b$; it is finite, $|\Gamma|$ being compact. By maximality every point of $|\Gamma|$ lies
at $\mathsf{d}$-distance less than $r$ from some $z_a$, since otherwise it could be adjoined
to the set; thus $|\Gamma|\subset\bigcup_a\mathsf{B}(z_a,r)$. If $N=1$ the bound on $N$
holds. Let $N\ge2$.

The balls $\mathsf{B}(z_a,r/2)$ are pairwise disjoint: a point $y$ in
$\mathsf{B}(z_a,r/2)\cap\mathsf{B}(z_b,r/2)$ would give
$\mathsf{d}(z_a,z_b)\le\mathsf{d}(z_a,y)+\mathsf{d}(y,z_b)<r$. Fix $a$ and put
$f:=\mathsf{d}(z_a,\cdot)$, with $g_e=f\circ\gamma_e$ as above. As in (a), $f(|\Gamma|)$ is an
interval containing $f(z_a)=0$; since $N\ge2$ there is $b\ne a$, and $f(z_b)\ge r$. Hence
$[0,r/2[\,\subset f(|\Gamma|)$, and a point $y\in|\Gamma|$ with $f(y)<r/2$ lies in
$\mathsf{B}(z_a,r/2)$. Therefore
\begin{equation*}
[0,r/2[\;\subset f\bigl(|\Gamma|\cap\mathsf{B}(z_a,r/2)\bigr)=\bigcup_{e\in E}g_e(U_e^a),
\qquad
U_e^a:=\{t\in[a_e,b_e]:\gamma_e(t)\in\mathsf{B}(z_a,r/2)\}.
\end{equation*}
Each $U_e^a=g_e^{-1}([0,r/2[)$ is relatively open in $[a_e,b_e]$, because $g_e$ is continuous.
By \eqref{9.5:eq-graph-nets-2},
\begin{equation*}
r/2=\operatorname{Leb}([0,r/2[)\le\sum_{e\in E}\operatorname{Leb}\bigl(g_e(U_e^a)\bigr)
\le\sum_{e\in E}\int_{U_e^a}w_e .
\end{equation*}
For fixed $e$ the sets $U_e^a$, $a=1,\dots,N$, are pairwise disjoint, because the balls
$\mathsf{B}(z_a,r/2)$ are. Summing over $a$ therefore gives
\begin{equation*}
\frac{Nr}{2}\le\sum_{e\in E}\sum_{a=1}^{N}\int_{U_e^a}w_e
\le\sum_{e\in E}\int_{a_e}^{b_e}w_e=\mathsf{L}(\Gamma).
\end{equation*}
Thus $N\le 2\mathsf{L}(\Gamma)/r$, and since $N$ is an integer,
$N\le\lfloor2\mathsf{L}(\Gamma)/r\rfloor$. Together with the case $N=1$ this gives
$N\le\max\{1,\lfloor2\mathsf{L}(\Gamma)/r\rfloor\}$.
\end{proof}

\begin{proposition}[Cells near a graph against its scaled length]\label{9.5:prop-cells-near-graph}
Assume the conditions displayed in \eqref{9.5:eq-scaled-length-a}, and let $\rho$ satisfy
$0<\rho\le r_0$, where in this section we write $r_0:=M/(4L)$. For a graph $\Gamma$ let
$\mathcal{S}_{\rho}(\Gamma)$ denote the set of cells $\Delta$ for which there are $x\in\Delta$ and
$y\in|\Gamma|$ with $\mathsf{d}(x,y)<\rho$, and put
\begin{equation*}
  \nu(\Gamma):=\max\bigl\{1,\lfloor 8L\,\mathsf{L}(\Gamma)/M\rfloor\bigr\},
  \qquad C_d^{\mathrm{mix}}:=2^{d+3}L .
\end{equation*}
Then for every graph $\Gamma$
\begin{equation}\label{9.5:eq-cells-near-graph-1}
  \#\mathcal{S}_{\rho}(\Gamma)\le 2^d\nu(\Gamma)
  =2^d\max\bigl\{1,\lfloor 8L\,\mathsf{L}(\Gamma)/M\rfloor\bigr\}
  \le C_d^{\mathrm{mix}}\bigl(\mathsf{L}(\Gamma)/M+1\bigr).
\end{equation}
Moreover:
\begin{enumerate}
\item[(a)] if $\mathsf{L}(\Gamma)<r_0$, then $\#\mathcal{S}_{\rho}(\Gamma)\le 2^d$, and the cells of
$\mathcal{S}_{\rho}(\Gamma)$ have a common point;
\item[(b)] if $\#\mathcal{S}_{\rho}(\Gamma)>2^d$, then
\begin{equation}\label{9.5:eq-cells-near-graph-a}
  \mathsf{L}(\Gamma)\ge \#\mathcal{S}_{\rho}(\Gamma)\cdot\frac{M}{2^{d+3}L}.
\end{equation}
\end{enumerate}
\end{proposition}

\begin{proof}
We apply part (b) of the lemma on diameter and nets of a graph
(Lemma~\ref{9.5:lem-graph-nets}) with $r:=r_0$. Since $2/r_0=8L/M$, it gives
\begin{equation*}
  N\le\max\bigl\{1,\lfloor 2\mathsf{L}(\Gamma)/r_0\rfloor\bigr\}
  =\max\bigl\{1,\lfloor 8L\,\mathsf{L}(\Gamma)/M\rfloor\bigr\}=\nu(\Gamma)
\end{equation*}
points $z_1,\dots,z_N\in|\Gamma|$ such that $|\Gamma|\subset\bigcup_{a}\mathsf{B}(z_a,r_0)$.

Let $\Delta\in\mathcal{S}_{\rho}(\Gamma)$, and choose $x\in\Delta$ and $y\in|\Gamma|$ with
$\mathsf{d}(x,y)<\rho$. Since the balls $\mathsf{B}(z_a,r_0)$ cover $|\Gamma|$, there is an index $a$ with
$y\in\mathsf{B}(z_a,r_0)$. By the triangle inequality for the metric $\mathsf{d}$ and by
$\rho\le r_0$,
\begin{equation*}
  \mathsf{d}(x,z_a)\le\mathsf{d}(x,y)+\mathsf{d}(y,z_a)<\rho+r_0\le 2r_0=\frac{M}{2L}.
\end{equation*}
Thus $x\in\mathsf{B}(z_a,r_0+\rho)\subset\mathsf{B}(z_a,M/(2L))$, so $\Delta$ meets
$\mathsf{B}(z_a,M/(2L))$. Every cell of $\mathcal{S}_{\rho}(\Gamma)$ therefore meets one of the $N$ balls
$\mathsf{B}(z_a,M/(2L))$, $1\le a\le N$. The lemma on cells meeting a small scaled ball
(Lemma~\ref{9.5:lem-small-balls}) applies to each of them with $z=z_a$ and $r=M/(2L)$, since its
hypothesis $0<r\le M/(2L)$ holds; by part (iv) of that lemma each of these balls meets at most $2^d$
cells. Hence $\#\mathcal{S}_{\rho}(\Gamma)\le 2^dN\le 2^d\nu(\Gamma)$, which is the first inequality
in \eqref{9.5:eq-cells-near-graph-1}; the equality there is the definition of $\nu(\Gamma)$. For the
last inequality we use $\max\{1,t\}\le 1+t$ for $t\ge 0$ and $L\ge 1$:
\begin{equation*}
  2^d\nu(\Gamma)\le 2^d\Bigl(1+\frac{8L\,\mathsf{L}(\Gamma)}{M}\Bigr)
  \le 2^{d+3}L\Bigl(\frac{\mathsf{L}(\Gamma)}{M}+1\Bigr)
  =C_d^{\mathrm{mix}}\Bigl(\frac{\mathsf{L}(\Gamma)}{M}+1\Bigr).
\end{equation*}

(a) If $\mathsf{L}(\Gamma)<r_0$, then $2\mathsf{L}(\Gamma)/r_0<2$, so
$\lfloor 2\mathsf{L}(\Gamma)/r_0\rfloor\le 1$ and $N=1$. All cells of $\mathcal{S}_{\rho}(\Gamma)$
then meet the single ball $\mathsf{B}(z_1,M/(2L))$, and part (iv) of
Lemma~\ref{9.5:lem-small-balls} gives that there are at most $2^d$ of them and that they have a
common point.

(b) If $\#\mathcal{S}_{\rho}(\Gamma)>2^d$, then the first inequality in
\eqref{9.5:eq-cells-near-graph-1} forces $\nu(\Gamma)\ge 2$. Hence the maximum defining
$\nu(\Gamma)$ is attained by the floor, and
$\nu(\Gamma)=\lfloor 8L\,\mathsf{L}(\Gamma)/M\rfloor\le 8L\,\mathsf{L}(\Gamma)/M$. Consequently
$\#\mathcal{S}_{\rho}(\Gamma)\le 2^d\nu(\Gamma)\le 2^{d+3}L\,\mathsf{L}(\Gamma)/M$, which is
\eqref{9.5:eq-cells-near-graph-a}.
\end{proof}

\begin{remark}[Sharpness of the constant in $L$]
Let $e_1,\dots,e_d$ denote the coordinate unit vectors. Suppose that $\Omega_{j+1}$ coincides near
the origin with the half-space $\{x_1\ge 0\}$. The segment $\Gamma$ from $0$ to $nM_je_2$ lies in
$\Omega_{j+1}$, and $\mathsf{L}(\Gamma)=nM/L$. This segment touches at least $n$ cells of level $j$,
each of which belongs to $\mathcal{S}_{\rho}(\Gamma)$. A bound
$\#\mathcal{S}_{\rho}(\Gamma)\le C(\mathsf{L}(\Gamma)/M+1)$ would therefore require
$n\le C(n/L+1)$ for every $n$. As $n\to\infty$ this forces $C\ge L$, since
$n/(n/L+1)\to L$. So no bound of this form holds with $C<L$: the constant $C_d^{\mathrm{mix}}$
in \eqref{9.5:eq-cells-near-graph-1} must grow linearly in $L$, and $C_d^{\mathrm{mix}}=2^{d+3}L$
does.
\end{remark}

\begin{lemma}[Touch-connected families of cells and their count]\label{9.5:lem-connected-families}
Let the cells, the scaled length $\mathsf{L}$, the scaled metric $\mathsf{d}$, the graphs $\Gamma$
with supports $|\Gamma|$ and the families $\mathcal{S}_{\rho}(\Gamma)$, $\mathcal{S}_0(\Gamma)$ of
seen cells be as in Definition~\ref{9.5:def-scaled-length}, let $D=(L+2)^d-L^d$ be the bound of
Lemma~\ref{9.5:lem-cells-tiling}\textup{(d)} on the number of cells touching a given cell, and fix
a cell $\Delta_0$.
\begin{itemize}
\item[(a)] For every graph $\Gamma$ and every $\rho>0$ the family $\mathcal{S}_{\rho}(\Gamma)$, and
also the family $\mathcal{S}_0(\Gamma)$, is connected for the relation ``touch''.
\item[(b)] The touch-connected families $\mathcal{S}\ni\Delta_0$ with $n$ members number at most
$D^{2(n-1)}$.
\item[(c)] The families $\mathsf{F}$ contained in some touch-connected family
$\mathcal{S}\ni\Delta_0$ with $\#\mathcal{S}\le n$ number at most $(2D^2)^n$.
\end{itemize}
\end{lemma}

\begin{proof}
(a) We first show that $\mathcal{S}_0(\Gamma')$ is connected for every graph $\Gamma'$. Suppose it
were the disjoint union of two non-empty families $\mathcal{A}$ and $\mathcal{B}$ such that no
member of $\mathcal{A}$ touches a member of $\mathcal{B}$. Then no cell of $\mathcal{A}$ has a
common point with a cell of $\mathcal{B}$, so $\bigcup\mathcal{A}$ and $\bigcup\mathcal{B}$ are
disjoint; they are closed, being finite unions of closed cubes; both meet $|\Gamma'|$, since every
member of $\mathcal{S}_0(\Gamma')$ does; and their union contains $|\Gamma'|$, because the cells
cover $\mathbb{T}$ (Lemma~\ref{9.5:lem-cells-tiling}(a)) and every cell containing a point of
$|\Gamma'|$ belongs to $\mathcal{S}_0(\Gamma')$. The set $|\Gamma'|$, connected because $\Gamma'$ is
a graph in the sense of Definition~\ref{9.5:def-scaled-length}, would thus be split
into two disjoint non-empty relatively closed pieces, which is impossible.

Now let $\rho>0$ and $\Delta\in\mathcal{S}_{\rho}(\Gamma)$. There is a path $\sigma$ with
$\mathsf{L}(\sigma)<\rho$ from a point of $|\Gamma|$ to a point of $\Delta$, and
$\Gamma':=\Gamma\cup\{\sigma\}$ is a graph. If a cell $\Delta'$ meets $\sigma$ at a point $p$,
then the part of $\sigma$ from its initial point to $p$ is a path from $|\Gamma|$ to $p$ of scaled
length at most $\mathsf{L}(\sigma)$, and $\mathsf{d}$ is the path metric of the scaled length;
hence the scaled distance of $\Delta'$ from $|\Gamma|$, an infimum over its points, satisfies
\begin{equation*}
\inf_{y\in\Delta'}\mathsf{d}(y,|\Gamma|)\le\mathsf{d}(p,|\Gamma|)\le\mathsf{L}(\sigma)<\rho ,
\end{equation*}
so $\Delta'$ belongs to $\mathcal{S}_{\rho}(\Gamma)$; a cell meeting $|\Gamma|$ is at scaled
distance $0$ from it, and $0<\rho$. Hence
$\mathcal{S}_0(\Gamma')\subset\mathcal{S}_{\rho}(\Gamma)$. Moreover
$\Delta\in\mathcal{S}_0(\Gamma')$, because $\sigma$ ends in $\Delta$, and
$\mathcal{S}_0(\Gamma)\subset\mathcal{S}_0(\Gamma')$, because $|\Gamma|\subset|\Gamma'|$. So
$\mathcal{S}_{\rho}(\Gamma)$ is a union of families $\mathcal{S}_0(\Gamma')$, each connected for
touch by the first step, all containing the non-empty family $\mathcal{S}_0(\Gamma)$; such a union
is connected for touch.

(b) Let $\mathcal{S}\ni\Delta_0$ be touch-connected with $n$ members. Choose a spanning tree of the
graph on $\mathcal{S}$ whose edges join touching cells, and walk around it depth first, starting
at $\Delta_0$. Each of the $n-1$ edges is traversed twice, so the walk is a sequence
\begin{equation*}
\Delta_0=v_0,\ v_1,\ \dots,\ v_{2(n-1)}
\end{equation*}
of cells in which consecutive cells touch, and the set of its values is $\mathcal{S}$. The family
$\mathcal{S}$ is therefore determined by such a sequence. Given $v_i$, there are at most $D$
choices of $v_{i+1}$ by Lemma~\ref{9.5:lem-cells-tiling}(d), so there are at most $D^{2(n-1)}$
such sequences, and at most $D^{2(n-1)}$ families $\mathcal{S}$.

(c) A family $\mathsf{F}$ as in (c) is a subset of a touch-connected family $\mathcal{S}\ni\Delta_0$
with $n'$ members, $1\le n'\le n$. By (b) there are at most $D^{2(n'-1)}$ such $\mathcal{S}$, and
each has $2^{n'}$ subsets. Since
$2^{n'}D^{2(n'-1)}=2^nD^{2(n-1)}(2D^2)^{-(n-n')}$, the number of families $\mathsf{F}$ is at most
\begin{align*}
\sum_{n'=1}^{n}2^{n'}D^{2(n'-1)}
&\le 2^nD^{2(n-1)}\sum_{i\ge0}(2D^2)^{-i}\\
&\le 2^{n+1}D^{2(n-1)}\le(2D^2)^n .
\end{align*}
Here the geometric series is at most $2$ and the last inequality is $2\le D^2$; both hold because
$D=(L+2)^d-L^d\ge 3^d-1\ge 8$, the function $L\mapsto(L+2)^d-L^d$ being increasing and $d\ge2$.
\end{proof}

\begin{proposition}[Families of cells within scaled reach]\label{9.5:prop-families-reach}
Assume the standing conditions \eqref{9.5:eq-scaled-length-a} of
Definition~\ref{9.5:def-scaled-length}, and let $0<\rho\le r_0$, where $r_0=M/(4L)$ as in
Proposition~\ref{9.5:prop-cells-near-graph}. Let $D=(L+2)^d-L^d$ be the constant of
Lemma~\ref{9.5:lem-connected-families}, and put $c_d^{\mathrm{mix}}=(2D^2)^{2^{d+3}L}$.
For every cell $\Delta_0$ and every real $\tau\ge 0$,
\begin{equation}\label{9.5:eq-families-reach-1}
\#\bigl\{\mathsf{F}\subset\pi^{\mathrm{mix}}\ \text{finite}:\ \Delta_0\in\mathsf{F},\
\mathsf{t}_{\rho}(\mathsf{F})\le\tau\bigr\}
\le(2D^2)^{2^d\max\{1,\lfloor 8L\tau/M\rfloor\}};
\end{equation}
in particular, for $\tau=\mathsf{m}M$ with $\mathsf{m}\in\mathbb{N}$, this number is at most
$(c_d^{\mathrm{mix}})^{\mathsf{m}+1}$. If moreover $\mathsf{t}_{\rho}(\mathsf{F})<(a+1)r_0/2$
with an integer $a\ge 0$, then $\#\mathsf{F}\le 2^d\max\{1,a\}$, and the finite families
$\mathsf{F}\ni\Delta_0$ with this property number at most $(2D^2)^{2^d\max\{1,a\}}$.
\end{proposition}

\begin{proof}
We first treat the last assertion. Let $a\ge 0$ be an integer and let $\mathsf{F}$ be a finite
family of cells with $\mathsf{t}_{\rho}(\mathsf{F})<(a+1)r_0/2$. By the definition of the loose
tree length $\mathsf{t}_{\rho}$, there is a graph $\Gamma$ with
$\mathsf{F}\subset\mathcal{S}_{\rho}(\Gamma)$ and $\mathsf{L}(\Gamma)<(a+1)r_0/2$, that is,
$2\mathsf{L}(\Gamma)/r_0<a+1$. Since $r_0=M/(4L)$ we have
$2\mathsf{L}(\Gamma)/r_0=8L\mathsf{L}(\Gamma)/M$, and the integer part of a real number smaller
than the integer $a+1$ is at most $a$; hence
\begin{equation*}
\nu(\Gamma)=\max\bigl\{1,\lfloor 8L\mathsf{L}(\Gamma)/M\rfloor\bigr\}\le\max\{1,a\}.
\end{equation*}
As $0<\rho\le r_0$, Proposition~\ref{9.5:prop-cells-near-graph} applies and gives
$\#\mathcal{S}_{\rho}(\Gamma)\le 2^d\nu(\Gamma)\le 2^d\max\{1,a\}$, so that
$\#\mathsf{F}\le 2^d\max\{1,a\}$. If in addition $\Delta_0\in\mathsf{F}$, then
$\Delta_0\in\mathcal{S}_{\rho}(\Gamma)$, and $\mathcal{S}_{\rho}(\Gamma)$ is connected for the
relation ``touch'' by Lemma~\ref{9.5:lem-connected-families}(a). Thus every such $\mathsf{F}$
is contained in a touch-connected family $\mathcal{S}\ni\Delta_0$ with
$\#\mathcal{S}\le n:=2^d\max\{1,a\}$, and by Lemma~\ref{9.5:lem-connected-families}(c) the
families with this property number at most $(2D^2)^{n}$. This proves the last assertion.

For \eqref{9.5:eq-families-reach-1}, let $\tau\ge 0$ and put
$a:=\lfloor 2\tau/r_0\rfloor=\lfloor 8L\tau/M\rfloor$, an integer $\ge 0$. Then $2\tau/r_0<a+1$,
i.e.\ $\tau<(a+1)r_0/2$, so every $\mathsf{F}$ with $\mathsf{t}_{\rho}(\mathsf{F})\le\tau$
satisfies $\mathsf{t}_{\rho}(\mathsf{F})<(a+1)r_0/2$. The last assertion, applied with this
$a$, bounds the left side of \eqref{9.5:eq-families-reach-1} by
$(2D^2)^{2^d\max\{1,\lfloor 8L\tau/M\rfloor\}}$.

Finally let $\tau=\mathsf{m}M$ with $\mathsf{m}\in\mathbb{N}$. Then
$\lfloor 8L\tau/M\rfloor=8L\mathsf{m}$. Since $1\le 8L\le 8L(\mathsf{m}+1)$ and
$8L\mathsf{m}\le 8L(\mathsf{m}+1)$,
\begin{equation*}
2^d\max\{1,8L\mathsf{m}\}\le 2^{d+3}L(\mathsf{m}+1).
\end{equation*}
Since $D\ge 1$, we have $2D^2\ge 1$, and the bound \eqref{9.5:eq-families-reach-1} gives at
most
\begin{equation*}
(2D^2)^{2^{d+3}L(\mathsf{m}+1)}=\bigl((2D^2)^{2^{d+3}L}\bigr)^{\mathsf{m}+1}
=(c_d^{\mathrm{mix}})^{\mathsf{m}+1}.
\end{equation*}
\end{proof}

\begin{lemma}[Common point, separation and tree length]\label{9.5:lem-common-point}
Assume the conditions \eqref{9.5:eq-scaled-length-a} on the label $\boldsymbol{\Omega}$ and let
$0<\rho\le r_0$, where $r_0=M/(4L)$ is the radius of
Proposition~\ref{9.5:prop-cells-near-graph}. Recall that
$d^{\circ}_{\mathrm{mix}}(\mathsf{F})=\mathsf{t}_0(\mathsf{F})/M$ for a finite family $\mathsf{F}$
of cells.
\begin{enumerate}
\item[(a)] If $\mathsf{t}_{\rho}(\mathsf{F})<r_0$, then the cells of $\mathsf{F}$ have a common point
(and $\#\mathsf{F}\le 2^d$); if the cells of $\mathsf{F}$ have a common point, then
$\mathsf{t}_0(\mathsf{F})=0$; if $\mathsf{t}_0(\mathsf{F})=0$, then $\mathsf{t}_{\rho}(\mathsf{F})=0$.
The conditions in this chain are equivalent.
\item[(b)] Two cells $\Delta,\Delta'$ that do not touch satisfy
$\mathsf{d}(\Delta,\Delta')\ge M/L$. Moreover, for any two cells $\Delta,\Delta'$ and every finite
family $\mathsf{F}$ of cells containing both,
$\mathsf{t}_{\rho}(\mathsf{F})\ge\mathsf{d}(\Delta,\Delta')-2\rho$.
\item[(c)] If
\begin{equation}\label{9.5:eq-common-point-a}
M\ge 2^{d+3}L\rho ,
\end{equation}
then for every finite family $\mathsf{F}$ of cells
\begin{equation}\label{9.5:eq-common-point-1}
\mathsf{t}_0(\mathsf{F})\le 2\,\mathsf{t}_{\rho}(\mathsf{F})+2^d\rho ,
\qquad\text{that is,}\qquad
\mathsf{t}_{\rho}(\mathsf{F})\ge \tfrac12 M\,d^{\circ}_{\mathrm{mix}}(\mathsf{F})-2^{d-1}\rho .
\end{equation}
\item[(d)] Let $\mathsf{X}$ be a wall-connected finite family of cells, $X:=\bigcup\mathsf{X}$, and
\begin{equation*}
d_{\mathrm{mix}}(X):=M^{-1}\inf\bigl\{\mathsf{L}(\Gamma):\ |\Gamma|\subset X,\
\Gamma\text{ meets every cell of }\mathsf{X}\bigr\}.
\end{equation*}
Then $d_{\mathrm{mix}}(X)\le\#\mathsf{X}-1$ and
\begin{equation}\label{9.5:eq-common-point-2}
d_{\mathrm{mix}}(X)\le 2^{d+3}L\,d^{\circ}_{\mathrm{mix}}(\mathsf{X}).
\end{equation}
\end{enumerate}
\end{lemma}

\begin{proof}
Two facts on the quantities $\mathsf{t}$ are used throughout. First, $\mathsf{t}_0(\mathsf{F})$ is
an infimum of $\mathsf{L}(\Gamma)$ over graphs $\Gamma$ meeting every cell of $\mathsf{F}$, and
$\mathsf{t}_{\rho}(\mathsf{F})$ the infimum of $\mathsf{L}(\Gamma)$ over graphs $\Gamma$ with
$\mathsf{F}\subset\mathcal{S}_{\rho}(\Gamma)$. Second, a cell met by $|\Gamma|$ is at scaled distance
$0<\rho$ from $\Gamma$, so $\mathcal{S}_0(\Gamma)\subset\mathcal{S}_{\rho}(\Gamma)$ and hence
$\mathsf{t}_{\rho}(\mathsf{F})\le\mathsf{t}_0(\mathsf{F})$.

(a) Suppose $\mathsf{t}_{\rho}(\mathsf{F})<r_0$. Since $\mathsf{t}_{\rho}(\mathsf{F})$ is an infimum,
there is a graph $\Gamma$ with $\mathsf{L}(\Gamma)<r_0$ and $\mathsf{F}\subset\mathcal{S}_{\rho}(\Gamma)$.
By the short-class part of Proposition~\ref{9.5:prop-cells-near-graph}, whose hypotheses
\eqref{9.5:eq-scaled-length-a} and $0<\rho\le r_0$ are those of the present lemma,
$\#\mathcal{S}_{\rho}(\Gamma)\le 2^d$ and the cells of $\mathcal{S}_{\rho}(\Gamma)$ have a common
point. The same then holds for the subfamily $\mathsf{F}$, and $\#\mathsf{F}\le 2^d$. If the cells of
$\mathsf{F}$ have a common point $p$, the one-point graph $\{p\}$ has scaled length $0$ and meets every
cell of $\mathsf{F}$, so $\mathsf{t}_0(\mathsf{F})=0$. If $\mathsf{t}_0(\mathsf{F})=0$, then
$0\le\mathsf{t}_{\rho}(\mathsf{F})\le\mathsf{t}_0(\mathsf{F})=0$. Finally,
$\mathsf{t}_{\rho}(\mathsf{F})=0$ implies $\mathsf{t}_{\rho}(\mathsf{F})<r_0$ since $r_0>0$, which closes
the chain of implications.

(b) Let $\Delta,\Delta'$ be cells that do not touch. Since $\mathsf{d}$ is the path metric of
$\mathsf{L}$, it suffices to show that every path $\gamma$ joining $\Delta$ to $\Delta'$ with
$\mathsf{L}(\gamma)<M$ satisfies $\mathsf{L}(\gamma)\ge M/L$ (paths with $\mathsf{L}(\gamma)\ge M$
satisfy this because $L\ge 1$). Let $\gamma$ be such a path and let $i$ be the least level of a cell
meeting $\gamma$.

We claim that no point $p$ of $\gamma$ has $m(p)\ge i+2$. Let $\Delta_0$ be a cell of level $i$
meeting $\gamma$. By Lemma~\ref{9.5:lem-cells-tiling}(b), $m=i$ on $\operatorname{int}\Delta_0$; as
$m(x)=\max\{j: x\in\Omega_j\}$ and the sequence $\boldsymbol{\Omega}$ is decreasing,
$\operatorname{int}\Delta_0\subset\mathbb{T}\setminus\Omega_{i+1}$, and since the cube $\Delta_0$ is
the closure of its interior, $\Delta_0\subset\Omega_{i+1}^{\complement}$. If a point $p$ of $\gamma$
had $m(p)\ge i+2$, then $p\in\Omega_{i+2}$ (and $i+1<k$), so the part of $\gamma$ between a point of
$\gamma\cap\Delta_0$ and $p$ would be a path from a point of $\Omega_{i+1}^{\complement}$ to a point
of $\Omega_{i+2}$; by Lemma~\ref{9.5:lem-zone-crossing} its scaled length, and hence
$\mathsf{L}(\gamma)$, would be at least $M$, contrary to the choice of $\gamma$. This proves the claim.

Consequently, at every point $p$ of $\gamma$ the weight
$\ell(p)^{-1}=(L^{m(p)}\eta)^{-1}$ with which $\mathsf{L}$ counts sup-length is at least
$(L^{i+1}\eta)^{-1}$. The cells $\Delta,\Delta'$ meet $\gamma$ at its end points, so their levels are
at least $i$; by Lemma~\ref{9.5:lem-cells-tiling}(a) each is a union of cubes of $\pi_i$, and as they
do not touch they are disjoint closed sets. Lemma~\ref{9.5:lem-zone-gap} gives
$\operatorname{dist}(\Delta,\Delta')\ge M_i$, so the sup-length of $\gamma$ is at least $M_i$ and
\begin{equation*}
\mathsf{L}(\gamma)\ge\frac{M_i}{L^{i+1}\eta}=\frac{ML^i\eta}{L^{i+1}\eta}=\frac{M}{L}.
\end{equation*}
Hence $\mathsf{d}(\Delta,\Delta')\ge M/L$.

For the second assertion, let $\Delta,\Delta'$ be any two cells, let $\mathsf{F}\ni\Delta,\Delta'$,
and let $\Gamma$ be a graph with $\mathsf{F}\subset\mathcal{S}_{\rho}(\Gamma)$. Since
$\Delta,\Delta'\in\mathcal{S}_{\rho}(\Gamma)$, there are $p,q\in|\Gamma|$ with
$\mathsf{d}(\Delta,p)<\rho$ and $\mathsf{d}(q,\Delta')<\rho$, and by
Lemma~\ref{9.5:lem-graph-nets}(a), $\mathsf{d}(p,q)\le\mathsf{L}(\Gamma)$. The triangle inequality gives
\begin{equation*}
\mathsf{d}(\Delta,\Delta')\le\mathsf{d}(\Delta,p)+\mathsf{d}(p,q)+\mathsf{d}(q,\Delta')
<2\rho+\mathsf{L}(\Gamma).
\end{equation*}
Taking the infimum over all such $\Gamma$ gives
$\mathsf{t}_{\rho}(\mathsf{F})\ge\mathsf{d}(\Delta,\Delta')-2\rho$.

(c) Assume \eqref{9.5:eq-common-point-a} and let $\Gamma$ be a graph with
$\mathsf{F}\subset\mathcal{S}_{\rho}(\Gamma)$. For each $\Delta\in\mathsf{F}$ the scaled distance from
$|\Gamma|$ to $\Delta$ is less than $\rho$, so there is a path from a point of $|\Gamma|$ to a point
of $\Delta$ of scaled length less than $\rho$. Adding these $\#\mathsf{F}$ paths to $\Gamma$ gives a
graph meeting every cell of $\mathsf{F}$, of scaled length at most
$\mathsf{L}(\Gamma)+\#\mathsf{F}\cdot\rho$. By Proposition~\ref{9.5:prop-cells-near-graph},
\begin{equation*}
\#\mathsf{F}\le\#\mathcal{S}_{\rho}(\Gamma)\le 2^d\max\{1,\lfloor 8L\mathsf{L}(\Gamma)/M\rfloor\}
\le 2^d+2^{d+3}L\,\mathsf{L}(\Gamma)/M ,
\end{equation*}
and therefore, using $2^{d+3}L\rho/M\le 1$, which is \eqref{9.5:eq-common-point-a},
\begin{align*}
\mathsf{t}_0(\mathsf{F})
&\le\mathsf{L}(\Gamma)+2^d\rho+2^{d+3}L\rho\,\mathsf{L}(\Gamma)/M\\
&\le 2\,\mathsf{L}(\Gamma)+2^d\rho .
\end{align*}
Taking the infimum over all such $\Gamma$ gives the first inequality of
\eqref{9.5:eq-common-point-1}. Dividing by $2$ and writing
$\mathsf{t}_0(\mathsf{F})=M\,d^{\circ}_{\mathrm{mix}}(\mathsf{F})$ gives the second.

(d) Since $\mathsf{X}$ is wall-connected, there is a spanning tree of $\mathsf{X}$ for the relation
of sharing a wall; it has $\#\mathsf{X}-1$ edges. For each cell $\Delta$ let $c_{\Delta}$ be its
centre. If $\Delta,\Delta'$ share a wall, let $p$ be the centre of $\Delta\cap\Delta'$, which is a
whole face of the smaller of the two cells, so $p\in\Delta\cap\Delta'$. The segment $[c_{\Delta},p]$
lies in $\Delta$ and has sup-length at most $\tfrac12 M_{\operatorname{lev}(\Delta)}$; by
Lemma~\ref{9.5:lem-cells-tiling}(b) its points have ambient level at least $\operatorname{lev}(\Delta)$,
so the weight of $\mathsf{L}$ on it is at most $(L^{\operatorname{lev}(\Delta)}\eta)^{-1}$ and
\begin{equation*}
\mathsf{L}([c_{\Delta},p])\le
\frac{M L^{\operatorname{lev}(\Delta)}\eta}{2L^{\operatorname{lev}(\Delta)}\eta}=\frac M2 ;
\end{equation*}
in the same way $[p,c_{\Delta'}]\subset\Delta'$ has $\mathsf{L}\le M/2$. The union of these two
segments over all edges of the tree is a graph contained in $X$, passing through all centres and so
meeting every cell of $\mathsf{X}$, of scaled length at most $(\#\mathsf{X}-1)M$. Hence
$d_{\mathrm{mix}}(X)\le\#\mathsf{X}-1$.

For \eqref{9.5:eq-common-point-2}: if the cells of $\mathsf{X}$ have a common point $p$, the graph
$\{p\}\subset X$ meets every cell of $\mathsf{X}$ and has scaled length $0$, so
$d_{\mathrm{mix}}(X)=0$ and \eqref{9.5:eq-common-point-2} holds. Otherwise, by (a) applied with the
given $\rho\le r_0$, $\mathsf{t}_{\rho}(\mathsf{X})\ge r_0$, and since
$\mathsf{t}_{\rho}\le\mathsf{t}_0$, also $\mathsf{t}_0(\mathsf{X})\ge r_0$. Let $\Gamma$ be any graph
meeting every cell of $\mathsf{X}$. Then $\mathsf{L}(\Gamma)\ge r_0=M/(4L)$, so
$8L\mathsf{L}(\Gamma)/M\ge 2$ and $\nu(\Gamma)=\max\{1,\lfloor 8L\mathsf{L}(\Gamma)/M\rfloor\}
\le 8L\mathsf{L}(\Gamma)/M$. As every cell of $\mathsf{X}$ is met by $\Gamma$, it lies in
$\mathcal{S}_{\rho}(\Gamma)$, and Proposition~\ref{9.5:prop-cells-near-graph} gives
\begin{equation*}
\#\mathsf{X}\le\#\mathcal{S}_{\rho}(\Gamma)\le 2^d\nu(\Gamma)\le 2^{d+3}L\,\mathsf{L}(\Gamma)/M .
\end{equation*}
Taking the infimum over all such $\Gamma$ yields
$\#\mathsf{X}\le 2^{d+3}L\,\mathsf{t}_0(\mathsf{X})/M=2^{d+3}L\,d^{\circ}_{\mathrm{mix}}(\mathsf{X})$,
and with $d_{\mathrm{mix}}(X)\le\#\mathsf{X}-1<\#\mathsf{X}$ this gives \eqref{9.5:eq-common-point-2}.
\end{proof}

\begin{lemma}[The entropy sum over mixed-size families]\label{9.5:lem-entropy-sum}
Assume the conditions \eqref{9.5:eq-scaled-length-a} of
Definition~\ref{9.5:def-scaled-length}, and let $0<\rho\le r_0$. Let $z\ge 1$ and $c>0$, put
\begin{equation*}
X_z := (2D^2 z)^{2^d},
\end{equation*}
and assume the floor
\begin{equation}\label{9.5:eq-entropy-sum-a}
cM \ge 16L\,\log(2X_z).
\end{equation}
Then for every cell $\Delta_0$ the following hold, all sums running over finite families
$\mathsf{F}\subset\pi^{\mathrm{mix}}$ of cells:
\begin{enumerate}
\item[(a)] $\displaystyle\sum_{\mathsf{F}\ni\Delta_0} z^{\#\mathsf{F}}
e^{-c\,\mathsf{t}_{\rho}(\mathsf{F})} \le 2X_z$;
\item[(b)] for every $\tau\ge r_0/2$,
\begin{equation*}
\sum_{\mathsf{F}\ni\Delta_0,\ \mathsf{t}_{\rho}(\mathsf{F})\ge\tau} z^{\#\mathsf{F}}
e^{-c\,\mathsf{t}_{\rho}(\mathsf{F})} \le e^{cr_0/4}\,e^{-c\tau/2};
\end{equation*}
\item[(c)] for all cells $\Delta_0,\Delta_1$,
\begin{equation*}
\sum_{\mathsf{F}\ni\Delta_0,\Delta_1} z^{\#\mathsf{F}}
e^{-2c\,\mathsf{t}_{\rho}(\mathsf{F})} \le 2X_z\,e^{2c\rho}\,e^{-c\,\mathsf{d}(\Delta_0,\Delta_1)}.
\end{equation*}
\end{enumerate}
For $z=2e^{\kappa_1}$ one has $X_z=(4D^2)^{2^d}e^{2^d\kappa_1}$: the exponent of $e^{\kappa_1}$
is $2^d$, that is $16$ at $d=4$, as in \cite[(1.11)]{B-CMP116}.
\end{lemma}

\begin{proof}
Here $r_0=M/(4L)$ is the radius of Definition~\ref{9.5:def-scaled-length}, and
$D=(L+2)^d-L^d\ge 1$, so that $X_z\ge 1$ because $2D^2z\ge 2$. Put $q:=e^{-cr_0/2}\le 1$.
Dividing \eqref{9.5:eq-entropy-sum-a} by $16L$ gives $cr_0/4\ge\log(2X_z)$, that is
$q^{1/2}=e^{-cr_0/4}\le 1/(2X_z)$, or
\begin{equation}\label{9.5:eq-entropy-sum-1}
X_z\,q^{1/2}\le \tfrac12,
\qquad\text{hence}\qquad
X_z\,q = (X_z q^{1/2})\,q^{1/2}\le \tfrac12\,q^{1/2}\le\tfrac12 .
\end{equation}

Since $\mathsf{t}_{\rho}(\mathsf{F})\ge 0$, every finite family $\mathsf{F}\ni\Delta_0$ lies in
exactly one of the classes
\begin{equation*}
\mathcal{F}_a(\Delta_0) := \bigl\{\mathsf{F}\ni\Delta_0 :
a r_0/2\le \mathsf{t}_{\rho}(\mathsf{F})<(a+1)r_0/2\bigr\},
\qquad a=0,1,2,\dots
\end{equation*}
Every $\mathsf{F}\in\mathcal{F}_a(\Delta_0)$ satisfies
$\mathsf{t}_{\rho}(\mathsf{F})<(a+1)r_0/2$, so by the last assertion of
Proposition~\ref{9.5:prop-families-reach} (families of cells within scaled reach), whose
hypotheses are those assumed here, $\#\mathsf{F}\le 2^d\max\{1,a\}$, and the class
$\mathcal{F}_a(\Delta_0)$ has at most $(2D^2)^{2^d\max\{1,a\}}$ members. For each member,
$z\ge 1$ gives $z^{\#\mathsf{F}}\le z^{2^d\max\{1,a\}}$, and
$\mathsf{t}_{\rho}(\mathsf{F})\ge ar_0/2$ gives $e^{-c\,\mathsf{t}_{\rho}(\mathsf{F})}\le q^a$.
Hence
\begin{equation}\label{9.5:eq-entropy-sum-2}
\sum_{\mathsf{F}\in\mathcal{F}_a(\Delta_0)} z^{\#\mathsf{F}}
e^{-c\,\mathsf{t}_{\rho}(\mathsf{F})}
\le (2D^2)^{2^d\max\{1,a\}}\, z^{2^d\max\{1,a\}}\, q^a
= X_z^{\max\{1,a\}}\,q^a .
\end{equation}
All terms of the sums in (a)--(c) are non-negative, so these sums may be regrouped by classes.

(a) By \eqref{9.5:eq-entropy-sum-2}, the class $a=0$ contributes at most $X_z$, and the class
$a\ge1$ at most $(X_zq)^a$. By \eqref{9.5:eq-entropy-sum-1},
$\sum_{a\ge1}(X_zq)^a\le\sum_{a\ge1}2^{-a}=1$. The sum in (a) is therefore at most
$X_z+1\le 2X_z$, since $X_z\ge1$.

(b) Let $\tau\ge r_0/2$ and put $a_0:=\lfloor 2\tau/r_0\rfloor$; then $a_0\ge 1$. If
$\mathsf{F}\in\mathcal{F}_a(\Delta_0)$ and $\mathsf{t}_{\rho}(\mathsf{F})\ge\tau$, then
$(a+1)r_0/2>\tau$, so the integer $a$ exceeds $2\tau/r_0-1$ and therefore $a\ge a_0$. Only the
classes $a\ge a_0\ge1$ contribute, each at most $(X_zq)^a$ by \eqref{9.5:eq-entropy-sum-2}. Summing
the geometric series with ratio $X_zq\le\frac12$ and then using \eqref{9.5:eq-entropy-sum-1},
\begin{equation*}
\sum_{a\ge a_0}(X_zq)^a \le 2\,(X_zq)^{a_0}
\le 2\,\bigl(\tfrac12\,q^{1/2}\bigr)^{a_0}
\le q^{a_0/2},
\end{equation*}
where the last step uses $a_0\ge1$. Finally $a_0\ge 2\tau/r_0-1$ and $q\le1$ give
\begin{equation*}
q^{a_0/2}\le q^{(2\tau/r_0-1)/2}
= \exp\Bigl(-\frac{cr_0}{2}\cdot\frac{2\tau/r_0-1}{2}\Bigr)
= e^{cr_0/4}\,e^{-c\tau/2}.
\end{equation*}

(c) Let $\mathsf{F}$ contain both $\Delta_0$ and $\Delta_1$. By part (b) of
Lemma~\ref{9.5:lem-common-point} (common point, separation and tree length), whose hypotheses are
those assumed here, $\mathsf{t}_{\rho}(\mathsf{F})\ge\mathsf{d}(\Delta_0,\Delta_1)-2\rho$. Hence
\begin{equation*}
e^{-2c\,\mathsf{t}_{\rho}(\mathsf{F})}
= e^{-c\,\mathsf{t}_{\rho}(\mathsf{F})}\,e^{-c\,\mathsf{t}_{\rho}(\mathsf{F})}
\le e^{-c(\mathsf{d}(\Delta_0,\Delta_1)-2\rho)}\,e^{-c\,\mathsf{t}_{\rho}(\mathsf{F})}.
\end{equation*}
Multiplying by $z^{\#\mathsf{F}}$, summing over $\mathsf{F}\ni\Delta_0,\Delta_1$, enlarging the sum
of non-negative terms to all $\mathsf{F}\ni\Delta_0$, and applying (a) gives
\begin{equation*}
\sum_{\mathsf{F}\ni\Delta_0,\Delta_1} z^{\#\mathsf{F}}
e^{-2c\,\mathsf{t}_{\rho}(\mathsf{F})}
\le e^{2c\rho}\,e^{-c\,\mathsf{d}(\Delta_0,\Delta_1)}
\sum_{\mathsf{F}\ni\Delta_0} z^{\#\mathsf{F}} e^{-c\,\mathsf{t}_{\rho}(\mathsf{F})}
\le 2X_z\,e^{2c\rho}\,e^{-c\,\mathsf{d}(\Delta_0,\Delta_1)}.
\end{equation*}

For the final assertion, with $z=2e^{\kappa_1}$ one has $2D^2z=4D^2e^{\kappa_1}$, and raising to the
power $2^d$ gives $X_z=(4D^2)^{2^d}e^{2^d\kappa_1}$.
\end{proof}

\begin{remark}[Coarser tilings by unions of cells]\label{9.5:rem-coarser-tilings}
Assume the hypotheses of Proposition~\ref{9.5:prop-cells-near-graph}, namely (H1)--(H3) and
$0<\rho\le r_0=M/4L$. Let $\pi^{\bullet}$ be a tiling of $\mathbb{T}$ by closed cubes each of which
is a union of cells. An example is $\pi_k$ itself, taken also in the lower zones: every cube of
$\pi_k$ is a union of cells by (H1) and (H2). For a graph $\Gamma$ with support $|\Gamma|$, call a
member $P$ of $\pi^{\bullet}$ \emph{seen} from $\Gamma$ if $\mathsf{d}(P,|\Gamma|)<\rho$. Then
\begin{equation}\label{9.5:eq-coarser-tilings-1}
\#\{P\in\pi^{\bullet}:\ P \text{ seen from } \Gamma\}\le\#\mathcal{S}_{\rho}(\Gamma)\le 2^d\nu(\Gamma).
\end{equation}
Moreover, Proposition~\ref{9.5:prop-families-reach} and Lemma~\ref{9.5:lem-entropy-sum} hold for
$\pi^{\bullet}$, anchored at a point instead of a cell, with one more factor $2^d$. No later
argument uses that a set toggled by a single decoupling parameter is a single cell.
\end{remark}

\begin{proof}
Let $P$ be a member of $\pi^{\bullet}$ seen from $\Gamma$. Since $P$ is a finite union of cells
$\Delta$, the distance from $P$ to $|\Gamma|$ is the minimum of the distances
$\mathsf{d}(\Delta,|\Gamma|)$ over these cells. Hence some cell $\Delta\subset P$ satisfies
$\mathsf{d}(\Delta,|\Gamma|)<\rho$, that is, $\Delta\in\mathcal{S}_{\rho}(\Gamma)$.

Two distinct members of the tiling have disjoint interiors. A cell has non-empty interior, so it
cannot be contained in two distinct members. Thus distinct members contain distinct cells.

Choose one such seen cell in each seen member. The resulting map from seen members to
$\mathcal{S}_{\rho}(\Gamma)$ is injective. This gives the first inequality of
\eqref{9.5:eq-coarser-tilings-1}. The second inequality is Proposition~\ref{9.5:prop-cells-near-graph}.

Now let $\Gamma$ pass through a point $x_0$. The same choice of one seen cell in each member
assigns to each family of members seen from $\Gamma$ a sub-family of $\mathcal{S}_{\rho}(\Gamma)$.
Distinct members receive distinct cells, so this assignment is injective. The family
$\mathcal{S}_{\rho}(\Gamma)$ is touch-connected. It also contains one of the at most $2^d$ cells
containing $x_0$, since such a cell lies at scaled distance $0<\rho$ from $|\Gamma|$.

The counting bounds of Proposition~\ref{9.5:prop-families-reach} and the sums of
Lemma~\ref{9.5:lem-entropy-sum} are anchored at a cell $\Delta_0$. They therefore transfer to
$\pi^{\bullet}$, anchored at $x_0$, by summing over the at most $2^d$ choices of the anchoring cell
containing $x_0$. This accounts for the one additional factor $2^d$.
\end{proof}

\begin{definition}[Shell grids, boxes and induced sequences]\label{9.5:def-shell-grids}
We use the torus $\mathbb{T}$ of side $\ell_{\mathbb{T}}$, the sup-norm distance $\operatorname{dist}$ on it, the
lattice unit $\eta$, the integers $L$, $k$, $M$, $M_1$, the shell-cube sides $s_j=M_1L^j\eta$, the
tilings $\pi_j$ with the product form of their grids, the set $\pi^{\mathrm{mix}}$ of cells and the
decreasing sequence $\boldsymbol{\Omega}=(\Omega_j)_{0\le j\le k}$, all as in
Definition~\ref{9.5:def-scaled-length}. We extend the sides by one level and put
\begin{equation*}
s_{k+1}:=Ls_k=M_1L^{k+1}\eta .
\end{equation*}

\emph{Shell grids.} For $0\le n\le k+1$ the \emph{shell grid} $\mathcal{G}_n$ is a tiling of $\mathbb{T}$ by
closed cubes of side $s_n$, of product form in the sense of
Definition~\ref{9.5:def-scaled-length}. A tiling is said to refine another if every cube of the
second is a union of cubes of the first. The shell grids are required to satisfy
\begin{equation}\label{9.5:eq-shell-grids-a}
\begin{gathered}
\mathcal{G}_n \text{ refines } \mathcal{G}_{n+1}\ \text{ and }\ \mathcal{G}_n \text{ refines } \pi_n
\qquad (0\le n\le k),\\
5s_{k+1}<\ell_{\mathbb{T}} .
\end{gathered}
\end{equation}
The side $M_n=ML^n\eta$ of the cells of $\pi_n$ is an integer multiple of $s_n$, because $M_1$
divides $M$; the grids of $\pi_n$ and $\mathcal{G}_n$ are taken aligned, so that the partitions
are compatible with one another, as for the partitions of the lattice used in
\cite{B-CMP119}.

\emph{Neighbourhoods and shell interiors.} For $Q\in\mathcal{G}_n$ and $r>0$ we write
\begin{equation*}
N_r(Q):=\{x\in\mathbb{T}:\ \operatorname{dist}(x,Q)<r\}
\end{equation*}
for the open $r$-neighbourhood of $Q$ in the sup-norm distance. For a closed set
$W\subset\mathbb{T}$ and $0\le n\le k+1$, the \emph{$s_n$-shell interior} of $W$ is
\begin{equation}\label{9.5:eq-shell-grids-1}
W^{-[s_n]}:=\bigcup\bigl\{Q\in\mathcal{G}_n:\ N_{2s_n}(Q)\subset W\bigr\}.
\end{equation}

\emph{Boxes.} A \emph{box} is a set of the form $Z=\bigcup\mathsf{Z}$, where
$\mathsf{Z}\subset\pi^{\mathrm{mix}}$ is an arbitrary family of cells, whose members may be of
any levels. Being a finite union of closed cubes, a box is closed, so that
\eqref{9.5:eq-shell-grids-1} applies to $W=Z$.

\emph{The sequence induced on a box.} For a box $Z$ and $0\le n\le k$ we put
\begin{equation}\label{9.5:eq-shell-grids-2}
\Omega_n(Z):=\bigcup\bigl\{Q\in\mathcal{G}_n:\ Q\subset\Omega_n,\ N_{2s_n}(Q)\subset Z\bigr\},
\end{equation}
that is, $\Omega_n(Z)=\Omega_n\cap Z^{-[s_n]}$, where the intersection of two unions of cubes of
$\mathcal{G}_n$ is read cube-wise, i.e.\ on sites. The \emph{sequence induced on $Z$} is
\begin{equation*}
\boldsymbol{\Omega}(Z):=\bigl(\Omega_n(Z)\bigr)_{0\le n\le k}.
\end{equation*}

\emph{The unit-lattice margin set.} Fix once and for all a region
$\Omega_{\star}\subset\Omega_k$, the same for all boxes: the region on which the unit-lattice
operators of the step act. For a box $Z$ we put
\begin{equation}\label{9.5:eq-shell-grids-3}
\Lambda(Z):=\Omega_{\star}\cap Z^{-[s_{k+1}]},
\end{equation}
the intersection being read cube-wise, i.e.\ on sites, as in \eqref{9.5:eq-shell-grids-2}. In
what follows the only property of $\Omega_{\star}$ that is used is that it does not depend
on $Z$.
\end{definition}

\begin{lemma}[The shell lemma for mixed-size boxes]\label{9.5:lem-shell-lemma}
Let the sequence $\boldsymbol{\Omega}=(\Omega_j)_{0\le j\le k}$ and the tilings $\pi_j$ be as in
Definition~\ref{9.5:def-scaled-length}, so that the conditions \eqref{9.5:eq-scaled-length-a}
hold. Let the shell grids $\mathcal{G}_n$ $(0\le n\le k+1)$ be as in
Definition~\ref{9.5:def-shell-grids}, so that \eqref{9.5:eq-shell-grids-a} holds. Let $Z$ be a
box, the union of a finite family $\mathsf{Z}$ of cells. Let
$\boldsymbol{\Omega}(Z)=(\Omega_n(Z))_{0\le n\le k}$
be its induced sequence and $\Lambda(Z)$ its set formed with the fixed region $\Omega_{\star}$, both
as defined in Definition~\ref{9.5:def-shell-grids}. Then the following hold.
\begin{enumerate}
\item[(s1)] For every $0\le n\le k$, $\Omega_n(Z)\subset\Omega_n$ and $\Omega_n(Z)$ is a union of
cubes of $\mathcal{G}_n$. For $0\le n<k$, moreover, $\Omega_{n+1}(Z)\subset\Omega_n(Z)$, the set
$\Omega_{n+1}(Z)$ is a union of cubes of $\mathcal{G}_n$, and
\begin{equation}\label{9.5:eq-shell-lemma-a}
\operatorname{dist}\bigl(\mathbb{T}\setminus\Omega_n(Z),\,\Omega_{n+1}(Z)\bigr)
\ge\min\{(2L-3)s_n,\,M_n\}.
\end{equation}
\item[(s2)] For $x\in\mathbb{T}$, whether $x\in\Omega_n(Z)$ is decided by $\Omega_n$ and by the set
$Z\cap B(x,3s_n)$. Whether $x\in\Lambda(Z)$ is decided by $\Omega_{\star}$ and by the set
$Z\cap B(x,3s_{k+1})$.
\item[(s3)] $\boldsymbol{\Omega}(\mathbb{T})=\boldsymbol{\Omega}$.
\item[(s4)] Let $\boldsymbol{\Omega}'=(\Omega'_j)_j$ be a sequence of the same kind, formed for the
lattice unit $\eta/L$ in place of $\eta$. Suppose $\boldsymbol{\Omega}'$ is paired with
$\boldsymbol{\Omega}$, that is, $\Omega'_{n+1}=\Omega_n$. Then $\Omega'_{n+1}(Z)=\Omega_n(Z)$.
\item[(s5)] Let $\mathsf{Z}=\mathsf{Z}_1\sqcup\mathsf{Z}_2$, where no cell of $\mathsf{Z}_1$ shares a
wall with a cell of $\mathsf{Z}_2$, and let $Z_i$ be the union of $\mathsf{Z}_i$. Then for every $n$
we have $Z^{-[s_n]}=Z_1^{-[s_n]}\sqcup Z_2^{-[s_n]}$ and
$\operatorname{dist}(Z_1^{-[s_n]},Z_2^{-[s_n]})\ge 4s_n$. Consequently
$\Omega_n(Z)=\Omega_n(Z_1)\sqcup\Omega_n(Z_2)$.
\end{enumerate}
\end{lemma}

\begin{proof}
Recall from Definition~\ref{9.5:def-shell-grids} how $\Omega_n(Z)$ is built. A point $p$ lies in
$\Omega_n(Z)$ if and only if some $Q\in\mathcal{G}_n$ satisfies both $p\in Q\subset\Omega_n$ and
$N_{2s_n}(Q)\subset Z$.

We prove (s1). Let $0\le n<k$ and let $Q'\in\mathcal{G}_{n+1}$ satisfy $Q'\subset\Omega_{n+1}$ and
$N_{2s_{n+1}}(Q')\subset Z$. Let $x\in Q'$. If $|y-x|<2s_{n+1}$, then
$\operatorname{dist}(y,Q')\le|y-x|<2s_{n+1}$. Since $s_{n+1}=Ls_n$, this gives
\begin{equation*}
B(x,2Ls_n)\subset N_{2s_{n+1}}(Q')\subset Z .
\end{equation*}
Now let $p\in\mathbb{T}$ with $|p-x|<\min\{(2L-3)s_n,M_n\}$. By the gap lemma
(Lemma~\ref{9.5:lem-zone-gap}), $\operatorname{dist}(\Omega_n^{\complement},\Omega_{n+1})\ge M_n$.
Since $x\in\Omega_{n+1}$ and $|p-x|<M_n$, it follows that $p\notin\Omega_n^{\complement}$. Hence $p$
is interior to $\Omega_n$.

Let $Q\in\mathcal{G}_n$ contain $p$. By \eqref{9.5:eq-shell-grids-a}, $\mathcal{G}_n$ refines $\pi_n$,
so $Q$ lies in a cube $\Delta$ of $\pi_n$. By \eqref{9.5:eq-scaled-length-a}, $\Omega_n$ is a union of
cubes of $\pi_n$. So if $\Delta$ were not one of them, $\Delta$ would lie in $\Omega_n^{\complement}$,
and then $p\in\Delta$ would lie in $\Omega_n^{\complement}$, which is false. Thus $Q\subset\Omega_n$.

Next, if $y\in N_{2s_n}(Q)$, there is $q\in Q$ with $|y-q|<2s_n$. As $|q-p|\le s_n$, this gives
$|y-p|<3s_n$. If $|y-p|<3s_n$, then $|y-x|<3s_n+(2L-3)s_n=2Ls_n$. Therefore
\begin{equation*}
N_{2s_n}(Q)\subset B(p,3s_n)\subset B(x,2Ls_n)\subset Z ,
\end{equation*}
and so $p\in\Omega_n(Z)$. Taking $p=x$ gives $\Omega_{n+1}(Z)\subset\Omega_n(Z)$. For general $p$,
every point within sup-distance $\min\{(2L-3)s_n,M_n\}$ of $\Omega_{n+1}(Z)$ lies in $\Omega_n(Z)$;
this is \eqref{9.5:eq-shell-lemma-a}. For every $0\le n\le k$ the inclusion
$\Omega_n(Z)\subset\Omega_n$ holds by definition, and also by definition $\Omega_n(Z)$ is a union of
cubes of $\mathcal{G}_n$. For $0\le n<k$, the set $\Omega_{n+1}(Z)$ is
a union of cubes of $\mathcal{G}_{n+1}$, and each of these is a union of cubes of $\mathcal{G}_n$,
because $\mathcal{G}_n$ refines $\mathcal{G}_{n+1}$ by \eqref{9.5:eq-shell-grids-a}.

We prove (s2). By the characterisation recalled above, $x\in\Omega_n(Z)$ if and only if some
$Q\in\mathcal{G}_n$ with $x\in Q\subset\Omega_n$ has $N_{2s_n}(Q)\subset Z$. For such $Q$ we have
$N_{2s_n}(Q)\subset B(x,3s_n)$, as shown in the proof of (s1). Hence
$N_{2s_n}(Q)\subset Z$ if and only if $N_{2s_n}(Q)\subset Z\cap B(x,3s_n)$. The grid
$\mathcal{G}_n$ is fixed, so the cubes $Q$ with $x\in Q\subset\Omega_n$ are determined by
$\Omega_n$. The same argument applies at $n=k+1$ with $\Omega_{\star}$ in place of $\Omega_n$, which
is how $\Lambda(Z)$ is defined in Definition~\ref{9.5:def-shell-grids}.

We prove (s3). For every $Q\in\mathcal{G}_n$ we have $N_r(Q)\subset\mathbb{T}$. Hence
$\mathbb{T}^{-[s_n]}$ is the union of all cubes of $\mathcal{G}_n$, and $\Omega_n(\mathbb{T})$ is
the union of the cubes of $\mathcal{G}_n$ contained in $\Omega_n$. That union is $\Omega_n$, since
$\Omega_n$ is a union of cubes of $\pi_n$ and $\mathcal{G}_n$ refines $\pi_n$.

We prove (s4). In the definition of $\Omega'_{n+1}(Z)$ the shell cubes have side
\begin{equation*}
s'_{n+1}=M_1L^{n+1}\eta/L=s_n ,
\end{equation*}
and the set read is $\Omega'_{n+1}=\Omega_n$. Thus the two definitions coincide word for word. Neither
definition reads the cell decomposition of $Z$ or the tilings of the two sequences; only the set $Z$,
the shell cubes of side $s_n$ and the set $\Omega_n=\Omega'_{n+1}$ enter.

We prove (s5). Since $Z_i\subset Z$, every $Q$ with $N_{2s_n}(Q)\subset Z_i$ also has
$N_{2s_n}(Q)\subset Z$. This gives the inclusion $Z_i^{-[s_n]}\subset Z^{-[s_n]}$.

For the reverse inclusion, let $Q\in\mathcal{G}_n$ with $O:=N_{2s_n}(Q)\subset Z$. The set $O$ is an
open box of side $5s_n$, and, since $5s_n\le 5s_{k+1}<\ell_{\mathbb{T}}$, a chart of $\mathbb{T}$
identifies $O$ with an open box in $\mathbb{R}^d$; we argue in this chart. Put $E:=Z_1\cap Z_2$. This
is a finite union of sets
$\Delta\cap\Delta'$ with $\Delta\in\mathsf{Z}_1$ and $\Delta'\in\mathsf{Z}_2$. Each such set is a
product of arcs of dimension at most $d-2$, because $\Delta$ and $\Delta'$ share no wall. Hence, in
the chart $O$, each is contained in an affine subspace of dimension at most $d-2$. In particular $E$
has empty interior.

Next, $O\setminus E$ is path-connected. Let $p,q\in O\setminus E$. Consider the set of $r\in O$ such
that the segment $[p,r]$ or the segment $[r,q]$ meets $E$. It lies in the finite union of the affine
spans of $p$, respectively of $q$, with the subspaces above. These spans have dimension at most
$d-1$, so that set has measure zero. Choose $r\in O$ outside it. By convexity of $O$ in the chart,
$[p,r]\cup[r,q]$ is a path from $p$ to $q$ in $O\setminus E$.

Since $O\subset Z_1\cup Z_2$, the set $O\setminus E$ is the disjoint union of
$(O\cap Z_1)\setminus Z_2$ and $(O\cap Z_2)\setminus Z_1$. These equal $(O\setminus E)\cap Z_1$ and
$(O\setminus E)\cap Z_2$ respectively, so both are relatively closed in $O\setminus E$. By
connectedness one of them is empty; say the second. Then $O\setminus E\subset Z_1$. Since $E$ has empty
interior and $Z_1$ is closed, $O\subset\operatorname{cl}(O\setminus E)\subset Z_1$. Hence
$Q\subset Z_1^{-[s_n]}$.

Distance. Let $x\in Z_1^{-[s_n]}$ and $y\in Z_2^{-[s_n]}$. Then $x$ lies in a cube $Q$ with
$N_{2s_n}(Q)\subset Z_1$, so $B(x,2s_n)\subset Z_1$; likewise $B(y,2s_n)\subset Z_2$. If
$|x-y|<4s_n$, these two open balls would meet in a non-empty open subset of $Z_1\cap Z_2=E$. That is
impossible, since $E$ has empty interior. Hence $\operatorname{dist}(Z_1^{-[s_n]},Z_2^{-[s_n]})\ge 4s_n$,
and in particular the union is disjoint.

Finally, the argument above shows that every $Q\in\mathcal{G}_n$ with $N_{2s_n}(Q)\subset Z$ satisfies
$N_{2s_n}(Q)\subset Z_1$ or $N_{2s_n}(Q)\subset Z_2$. Applying this to the cubes $Q\subset\Omega_n$
gives $\Omega_n(Z)=\Omega_n(Z_1)\cup\Omega_n(Z_2)$. The union is disjoint because
$\Omega_n(Z_i)\subset Z_i^{-[s_n]}$.
\end{proof}

\begin{remark}[Admissibility of the induced sequences]
In this remark only, $R>0$ denotes the constant of the admissibility condition
\cite[(2.2)]{B-CMP96}, read with the block size $\mathfrak{m}$. Assume the conditions
$(2L-3)M_1>R\mathfrak{m}$ and $M>R\mathfrak{m}$, which are floor conditions in force throughout this
chapter. Then, for
$0\le n<k$, the bound \eqref{9.5:eq-shell-lemma-a} gives
\begin{equation*}
\operatorname{dist}\bigl(\mathbb{T}\setminus\Omega_n(Z),\Omega_{n+1}(Z)\bigr)
\ge\min\{(2L-3)M_1,M\}\,L^n\eta>R\,\mathfrak{m}L^n\eta=R\,u_n .
\end{equation*}
This is the admissibility of the induced sequence $\boldsymbol{\Omega}(Z)$ in the sense of
\cite[(2.2)]{B-CMP96}, with the parameters $R$ and $\mathfrak{m}$.
\end{remark}

\begin{definition}[Refinements and the contour length]\label{9.5:def-refinements}
Let $\boldsymbol{\Omega}=(\Omega_j)_{0\le j\le k}$, the unit $\eta$, the ambient level
$m(x)=\max\{j:x\in\Omega_j\}$ and the ambient scale function $\ell(x):=L^{m(x)}\eta$ be as in
Definition~\ref{9.5:def-scaled-length}.

A \emph{refinement} of $\boldsymbol{\Omega}$ is a decreasing sequence
$\hat{\boldsymbol{\Omega}}=(\hat{\Omega}_n)_{n\le\hat{k}}$ of closed subsets of $\mathbb{T}$,
together with a unit $\hat{\eta}>0$, whose scale function satisfies
\begin{equation}\label{9.5:eq-refinements-1}
\hat{\ell}(x):=L^{\hat{m}(x)}\hat{\eta}\;\le\;\ell(x)=L^{m(x)}\eta
\qquad\text{for every }x\in\mathbb{T}.
\end{equation}
Here $\hat{m}(x):=\max\{n:x\in\hat{\Omega}_n\}$ for $x\in\hat{\Omega}_1$, and $\hat{m}(x):=0$
for $x\notin\hat{\Omega}_1$.

\emph{Examples.} (a) $\boldsymbol{\Omega}$ itself, with $\hat{\eta}=\eta$, since then
$\hat{\ell}=\ell$.

(b) Every induced sequence $\boldsymbol{\Omega}(Z)=(\Omega_n(Z))_n$ of
Definition~\ref{9.5:def-shell-grids}, with $\hat{\eta}=\eta$. Indeed
$\Omega_n(Z)\subset\Omega_n$ for every $n$, so the level $m_Z(x)$ of $x$ in
$\boldsymbol{\Omega}(Z)$ is at most $m(x)$. Hence
$\ell_{\boldsymbol{\Omega}(Z)}(x):=L^{m_Z(x)}\eta\le\ell(x)$.

(c) The sequence $\boldsymbol{\Omega}'(Z)$ of the finer of the two cutoffs, which has unit
$\hat{\eta}=\eta/L$ and members $\Omega'_{n+1}(Z)=\Omega_n(Z)$. For $x\in\Omega_0(Z)$ the
level of $x$ in $\boldsymbol{\Omega}'(Z)$ is $m_Z(x)+1$, so
\begin{equation*}
\hat{\ell}(x)=L^{m_Z(x)+1}\,\frac{\eta}{L}=\ell_{\boldsymbol{\Omega}(Z)}(x).
\end{equation*}
For $x\notin\Omega_0(Z)$ we have $x\notin\Omega'_1(Z)$, hence $\hat{m}(x)=0$ and
$\hat{\ell}(x)=\eta/L$. In both cases $\hat{\ell}(x)\le\ell(x)$, by (b) and because
$\eta/L\le\eta\le\ell(x)$.

\emph{Contour length.} For a contour (graph) $\Gamma$ in $\mathbb{T}$ we write $|dx|_1$ for the
$\ell^1$ line element along $\Gamma$, and $|\Gamma\cap S|$ for the $\ell^1$ length of the part
of $\Gamma$ in a set $S$. The \emph{contour length} of $\Gamma$ with respect to the
refinement $\hat{\boldsymbol{\Omega}}$ is
\begin{equation}\label{9.5:eq-refinements-2}
\operatorname{len}_{\hat{\boldsymbol{\Omega}}}(\Gamma)
:=\int_{\Gamma}\frac{|dx|_1}{\hat{\ell}(x)}
=\sum_{j}\bigl(L^{j}\hat{\eta}\bigr)^{-1}\bigl|\Gamma\cap\{\hat{m}=j\}\bigr|.
\end{equation}
The second expression is obtained by splitting $\Gamma$ according to the value of $\hat{m}$,
on which $\hat{\ell}$ equals $L^{j}\hat{\eta}$. It is the sum
$\sum_j(L^j\hat{\eta})^{-1}|\Gamma\cap B^j(\hat{\Lambda}_j)|$ of
\cite[(2.46)]{B-CMP96}, written for $\hat{\boldsymbol{\Omega}}$ and $\hat{\eta}$. In that
sum $B^j(\hat{\Lambda}_j)$ is the subdomain on which the natural scale is $L^j\hat{\eta}$.

Finally, for points $y,y'$ we set
\begin{equation*}
\hat{d}(y,y'):=\inf_{\Gamma}\operatorname{len}_{\hat{\boldsymbol{\Omega}}}(\Gamma).
\end{equation*}
The infimum is taken over the contours joining $y$ to $y'$ that are admissible in the sense of
\cite[(2.46)]{B-CMP96}.
\end{definition}

\begin{lemma}[Refinements only lengthen scaled distances]\label{9.5:lem-refinement-lengthens}
Let $\hat{\boldsymbol{\Omega}}$ be a refinement of $\boldsymbol{\Omega}$ in the sense of
Definition~\ref{9.5:def-refinements}, and let $\Gamma$ be a graph in $\mathbb{T}$. Then
\begin{equation*}
\operatorname{len}_{\hat{\boldsymbol{\Omega}}}(\Gamma)\ \ge\ \mathsf{L}(\Gamma),
\end{equation*}
where $\mathsf{L}$ is the scaled length of Definition~\ref{9.5:def-scaled-length}. Consequently
$\hat{d}(y,y')\ge\mathsf{d}(y,y')$ for all $y,y'\in\mathbb{T}$. Moreover, let
$X_0,\dots,X_n\subset\mathbb{T}$, and let $\Gamma'$ be a tree graph that meets all the $X_i$,
or a contour $y,y_1,\dots,y_n,y'$ in the sense of \cite[(3.93)]{B-CMP99b} that meets all the
$X_i$. Then
\begin{equation*}
\operatorname{len}_{\hat{\boldsymbol{\Omega}}}(\Gamma')\ \ge\
\inf\bigl\{\mathsf{L}(\Gamma''):\ \Gamma'' \text{ a graph meeting all } X_i\bigr\}.
\end{equation*}
\end{lemma}

\begin{proof}
The contour length $\operatorname{len}_{\hat{\boldsymbol{\Omega}}}(\Gamma)$ integrates the
$1$-norm line element $|dx|_1$ along $\Gamma$ against a weight that equals $1/\hat{\ell}$ off
the boundaries of the zones and, by the attribution convention recalled below, is at least
$1/\hat{\ell}$ on them. The scaled
length $\mathsf{L}(\Gamma)$ integrates the sup-norm line element $|dx|_\infty$ against the
weight $1/\ell$. We compare the two integrands pointwise. The first comparison is
$|dx|_1\ge|dx|_\infty$. The second is $1/\hat{\ell}(x)\ge 1/\ell(x)$, which is the defining
inequality $\hat{\ell}\le\ell$ of a refinement.

At a point $x$ on the common boundary of two zones, that is, of two regions on which
$\hat{m}$ is constant, the definition of
$\operatorname{len}_{\hat{\boldsymbol{\Omega}}}$ attributes each piece of contour to one of
the two zones. Where the ambient level also changes at $x$, the level
$m(x)=\max\{j:\ x\in\Omega_j\}$ is the larger of the two ambient levels, so $\ell(x)$ is the
larger scale and $\mathsf{L}$ already uses the smaller weight $1/\ell(x)$ there. Likewise
$\hat{m}(x)=\max\{n:\ x\in\hat{\Omega}_n\}$ takes the larger refined level. Hence, whichever zone
the piece is attributed to, its weight in $\operatorname{len}_{\hat{\boldsymbol{\Omega}}}$
is at least $1/\hat{\ell}(x)$, and therefore at least $1/\ell(x)$. Integrating the two
pointwise comparisons along $\Gamma$ gives
$\operatorname{len}_{\hat{\boldsymbol{\Omega}}}(\Gamma)\ge\mathsf{L}(\Gamma)$.

Admissible contours joining $y$ to $y'$ (the contours over which $\hat{d}$ is the infimum,
Definition~\ref{9.5:def-refinements}) are graphs joining $y$ to $y'$. Taking the infimum
of the first inequality over them therefore gives $\hat{d}(y,y')\ge\mathsf{d}(y,y')$.

Finally, the tree graph or contour $\Gamma'$ of the statement, meeting all the $X_i$, is
itself a graph meeting all the $X_i$. By the first inequality its contour length is at least
$\mathsf{L}(\Gamma')$. This is at least the infimum displayed in the statement.
\end{proof}

\begin{remark}
The lemma settles whether the shells of a box (Definition~\ref{9.5:def-refinements}) can
increase the number of cells of $\pi^{\mathrm{mix}}$ lying at scaled distance less than
$\rho$ from the graph $\Gamma$ of a walk, that is, the number of cells in
$\mathcal{S}_{\rho}(\Gamma)$ (Definition~\ref{9.5:def-scaled-length}). They cannot: that
number is bounded through $\nu(\Gamma)$, a function of the ambient scaled length
$\mathsf{L}(\Gamma)$, while the shells of a box, whose gaps $(2L-3)s_n$ are far below $M_n$,
only make walks longer than their ambient scaled length, since by the lemma the contour length
of a refinement is at least $\mathsf{L}(\Gamma)$.
\end{remark}

\begin{lemma}[Short physical segments have short scaled length]\label{9.5:lem-short-segments}
Let $\boldsymbol{\Omega}=(\Omega_j)_{0\le j\le k}$ be the decreasing sequence of zones of
Definition~\ref{9.5:def-scaled-length}, with ambient level $m(\cdot)$, ambient scale
$\ell(\cdot)$ and scaled path metric $\mathsf{d}$. Let $c\in\mathbb{T}$ and $x\in\mathbb{T}$, and let
$\varrho$ be a real number with $0<\varrho<M/L$ such that $|x-c|\le\varrho\,\ell(c)$. Then
$\mathsf{d}(c,x)\le L\varrho$.
\end{lemma}

\begin{proof}
In this lemma $\varrho$ is a real number local to the statement; it is not the parameter
$\varrho$ of the glossary. Let $[c,x]$ denote a shortest segment from $c$ to $x$ in the
sup-norm distance of $\mathbb{T}$. Every point $p$ of $[c,x]$ satisfies $|p-c|\le|x-c|$.
Since $\mathsf{d}(c,x)$ is an infimum of scaled lengths over paths from $c$ to $x$, we have
$\mathsf{d}(c,x)\le\mathsf{L}([c,x])$. By Definition~\ref{9.5:def-scaled-length}, the scaled
length weighs sup-norm length by $1/\ell$. Hence $\mathsf{L}([c,x])$ is at most $|x-c|$
divided by the least value of $\ell$ on $[c,x]$. Put $j:=m(c)$, so that $c\in\Omega_j$ and
$\ell(c)=L^j\eta$.

Case $j\ge1$. Let $p\in[c,x]$. Then
\begin{equation}\label{9.5:eq-short-segments-1}
|p-c|\le|x-c|\le\varrho L^j\eta<\frac{M}{L}\,L^j\eta=ML^{j-1}\eta=M_{j-1}.
\end{equation}
Since $0\le j-1<j\le k$, Lemma~\ref{9.5:lem-zone-gap} (the gap between zones is automatic)
applies at level $j-1$ and gives
\begin{equation*}
M_{j-1}\le\operatorname{dist}\bigl(\Omega_{j-1}^{\complement},\Omega_j\bigr).
\end{equation*}
Because $c\in\Omega_j$, \eqref{9.5:eq-short-segments-1} gives
\begin{equation*}
\operatorname{dist}(p,\Omega_j)\le|p-c|
<\operatorname{dist}\bigl(\Omega_{j-1}^{\complement},\Omega_j\bigr).
\end{equation*}
Hence $p\notin\Omega_{j-1}^{\complement}$. Therefore $p\in\Omega_{j-1}$, and so
$m(p)\ge j-1$ and $\ell(p)\ge L^{j-1}\eta$. This holds for every $p\in[c,x]$, so
\begin{equation*}
\mathsf{d}(c,x)\le\mathsf{L}([c,x])\le\frac{\varrho L^j\eta}{L^{j-1}\eta}=L\varrho .
\end{equation*}

Case $j=0$. For every $p\in[c,x]$ we have $m(p)\ge0$, hence $\ell(p)\ge\eta$. Since
$|x-c|\le\varrho\,\ell(c)=\varrho\eta$, we get
$\mathsf{d}(c,x)\le\mathsf{L}([c,x])\le\varrho\le L\varrho$.
\end{proof}

The hypothesis $\varrho<M/L$ is what confines the segment. A physical radius
$\varrho\,\ell(c)$ with $\varrho\ge M/L$ may cross the whole stack of thinner and thinner zones
below $\Omega_j$. With $\varrho<M/L$ the segment $[c,x]$ enters at most the next finer zone.

\begin{definition}[Walk expansions and their influence sets]\label{9.5:def-walk-influence}
Let $\boldsymbol{\Omega}$, the tilings $\pi_j$ and their cells, the set $\pi^{\mathrm{mix}}$ of all cells,
the sides $M_j$, $s_j$, $u_j$ and the sup-norm distance $\operatorname{dist}$ on $\mathbb{T}$ be as in
Definition~\ref{9.5:def-scaled-length}. Let $\hat{\boldsymbol{\Omega}}=(\hat{\Omega}_n)_{n\le\hat{k}}$,
with unit $\hat{\eta}$, be a refinement of $\boldsymbol{\Omega}$ in the sense of
Definition~\ref{9.5:def-refinements}, admissible at $(R,\mathfrak{m})$ with $R\ge L\ge 7$ in the sense of
\cite[(2.1)--(2.2)]{B-CMP96}. We put $\hat{u}_s:=\mathfrak{m}L^{s}\hat{\eta}$. Let $\mathbb{O}\in\mathfrak{O}$,
the list of the fifteen operators of the walk expansion, and let
$\mathbb{O}=\sum_{\omega\in\mathcal{W}_{\mathbb{O}}(\hat{\boldsymbol{\Omega}})}\mathbb{O}(\omega)$ be its
expansion into walk terms, built on the anchored covers, the associated partition of unity and the
localisation of the walk terms; each walk
$\omega=((\alpha_i,X_i))_{0\le i\le n}$ is a chain or a tree, $\alpha_i$ a factor type of a fixed finite list
of types and $X_i$ its domain.

A \emph{cube of $\omega$} is a cube $X$ of one of the covers $\mathcal{D}$, $\mathcal{D}'(\square)$,
$\mathcal{D}''(\square)$ of \cite{B-CMP99b}, built on $\hat{\boldsymbol{\Omega}}$, occurring in a factor of
$\omega$, or a base cube of $\omega$; the \emph{base cubes} of $\omega$ are the cubes $\square\in\mathcal{D}$
at which a term of the form \cite[(3.95)]{B-CMP99b} or \cite[(3.105)]{B-CMP99b} of $\omega$ was generated,
and each lies in some $X_i$. For a cube $X$ with scale $s(X)$, radius $r(X)$ and centre $c(X)$, and
$n\ge 0$, $X^{(n)}$ is the
cube with centre $c(X)$ and radius $r(X)+n\hat{u}_{s(X)}$; $X^{2L+1}$ denotes $X$ enlarged about $c(X)$ by
$2L+1$ block units $\hat{u}_{s(X)}$ of its own scale, a cube of radius at most
$r(X)+(2L+1)\hat{u}_{s(X)}$, and for a base cube $\square$ it is the
cube $\square^{2L+1}$ of radius $(2L+2)\hat{u}_{s(\square)}$. We write $X_i^{\natural}$ for the set
obtained from $X_i$ by enlarging each constituent grid cube by five units of its own scale.

Throughout this section we use the following properties of the walk expansion as standing hypotheses.
First, the geometry of one walk: every cube $X$ of $\omega$ has a scale $s(X)$, a radius $r(X)$ and an
anchored centre $c(X)$, that is, a point of $\hat{u}_{s(X)}\mathbb{Z}^d$ which is a corner of an
$s(X)$-block of the block lattice of the refined region at level $s(X)$, with
\begin{equation}\label{9.5:eq-walk-influence-b}
\begin{gathered}
\text{(i)}\quad r(X)\le 5\hat{u}_{s(X)},\qquad
c(X)\in\hat{\Omega}_{s(X)},\qquad
X\subset X_i\ \text{for some } i;\\
\text{(ii)}\quad X_i=Y_1\cup\dots\cup Y_a,\quad a\le\mathsf{q},\quad
Y_1,\dots,Y_a\ \text{cover cubes},\quad Y_b\cap Y_{b+1}\neq\emptyset\ (1\le b<a);\\
\text{(iii)}\quad \operatorname{loc}(\omega):=\bigcup_i X_i^{\natural}
\subset\bigcup\{X^{(5)}: X\ \text{a cover cube of } \omega\},\\
\operatorname{loc}^{+}(\omega):=\operatorname{loc}(\omega)\cup
\bigcup\{\square^{2L+1}: \square\ \text{a base cube of } \omega \text{ of scale}\le\hat{k}-1\}.
\end{gathered}
\end{equation}
Here $\mathsf{q}$ is the largest number of cover cubes in one factor type of the fixed finite list of
types, an integer independent of every parameter (in the families displayed in \cite{B-CMP99b}, at most
four); by \eqref{9.5:eq-walk-influence-b}(i), $X^{(5)}$ has radius at most $10\hat{u}_{s(X)}$, and
$\square^{2L+1}$ has radius $(2L+2)\hat{u}_{s(\square)}$. Property \eqref{9.5:eq-walk-influence-b} is the
geometry of one walk given by the anchored covers of the walk expansion and by the families of
\cite{B-CMP99b}. Second, the decay, reading and locality of walk terms: put
\[
\hat{d}(\omega):=\inf\{\operatorname{len}_{\hat{\boldsymbol{\Omega}}}(\Gamma):
\Gamma\ \text{a tree graph meeting}\ X_0,\dots,X_n\},
\]
with $\operatorname{len}_{\hat{\boldsymbol{\Omega}}}$ the contour length of
Definition~\ref{9.5:def-refinements}.
Then, in each norm in which the operators of $\mathfrak{O}$ are estimated in
\cite[(3.94), (3.99), (3.108)]{B-CMP99b} and \cite[(1.7)]{B-CMP116}, for the terms between the
localisations at points $y,y'$ and for every real $\lambda\ge 0$,
\begin{equation}\label{9.5:eq-walk-influence-c}
\begin{gathered}
\text{(i)}\quad \sum_{\omega:\ \hat{d}(\omega)\ge\lambda}\|\mathbb{O}(\omega)\|
\le B_0\,e^{-\frac12\delta_0\lambda}\,\Phi;\\
\text{(ii)}\quad \mathbb{O}(\omega)\ \text{has its kernel in}\ X_0\times X_n
\ \text{and depends on}\ (\mathbf{U},\mathbf{J})\ \text{only through their restrictions to}\
\operatorname{loc}(\omega);\\
\text{(iii)}\quad \omega\in\mathcal{W}_{\mathbb{O}}(\hat{\boldsymbol{\Omega}})\ \text{and}\
\mathbb{O}(\omega)\ \text{depend on}\ \hat{\boldsymbol{\Omega}}\ \text{only through}\
\{\hat{\Omega}_j\cap\operatorname{loc}^{+}(\omega)\}_j,
\end{gathered}
\end{equation}
where $\Phi$ stands for the scale factors of the corresponding estimate among
\cite[(3.94), (3.99), (3.108)]{B-CMP99b}, \cite[(1.7)]{B-CMP116}, and with $B_0$, $\delta_0$ independent
of $\hat{\boldsymbol{\Omega}}$; \eqref{9.5:eq-walk-influence-c}(i) and this independence are as stated in
\cite{B-CMP99b}. Property \eqref{9.5:eq-walk-influence-c}(ii) is the reading property of the walk terms,
and \eqref{9.5:eq-walk-influence-c}(iii) is the locality lemma of the walk expansion.
In \eqref{9.5:eq-walk-influence-c}(ii),
for a composite with box-independent local end factors, namely a composite whose end factors are the
almost local operator $C$ of \cite{B-CMP109} and its adjoint $C^{*}$, of range at most $2L$ units,
the end points lie within
$\mathsf{r}_e\hat{\ell}$, $\mathsf{r}_e:=2L$, of $X_0$ and $X_n$, with $\hat{\ell}$ the refined scale
function of Definition~\ref{9.5:def-refinements} at the end point concerned, hence in
$\operatorname{loc}(\omega)$, since $5\mathfrak{m}\ge 2L$ ($\mathfrak{m}$ being a power of $L$).

We put $\mathsf{u}:=\max\{\mathfrak{m},M_1,\mathsf{r}_e\}=\max\{\mathfrak{m},M_1,2L\}$ and, for a cube $X$
of $\omega$, $X^{+}:=X^{(5)}\cup X^{2L+1}$. With $\mathsf{t}_{\rho}$ the loose tree length of
Definition~\ref{9.5:def-scaled-length}, with $\ell(x)=L^{m(x)}\eta$ the ambient scale, and for a walk
$\omega$ of $\hat{\boldsymbol{\Omega}}$, we define
\begin{equation}\label{9.5:eq-walk-influence-d}
\begin{gathered}
\rho_{\times}:=L\mathsf{u}(20\mathsf{q}+5L+7),\qquad
\mathsf{t}_{\times}:=\mathsf{t}_{\rho_{\times}},\\
S^{\sharp}(\omega):=\{\Delta\in\pi^{\mathrm{mix}}:\ \operatorname{dist}(x,\Delta)<3LM_1\ell(x)
\ \text{for some}\ x\in\operatorname{loc}^{+}(\omega)\},\\
S^{\mathrm{mix}}(\omega):=\{\Delta\in\pi^{\mathrm{mix}}:\ \operatorname{dist}(\Delta,X^{+})
\le 3L\hat{u}_{s(X)}\ \text{for some cube}\ X\ \text{of}\ \omega\}.
\end{gathered}
\end{equation}
The set $S^{\sharp}(\omega)$ is scaled by the ambient label and $S^{\mathrm{mix}}(\omega)$ by the cubes of
the walk. On the constants we impose the condition
\begin{equation}\label{9.5:eq-walk-influence-a}
M\ \ge\ 2^{d+3}L\rho_{\times}=2^{d+3}L^{2}(20\mathsf{q}+5L+7)\cdot\max\{\mathfrak{m},M_1,2L\}.
\end{equation}
Condition \eqref{9.5:eq-walk-influence-a} is the condition \eqref{9.5:eq-common-point-a} of
Lemma~\ref{9.5:lem-common-point} at $\rho=\rho_{\times}$; it gives, with $r_0=M/(4L)$ the radius of that
lemma,
\begin{equation}\label{9.5:eq-walk-influence-1}
\rho_{\times}\le\frac{r_0}{2^{d+1}},\qquad (10\mathsf{q}+5L+6)\,\mathsf{u}<\frac{M}{L}.
\end{equation}
Moreover $3LM_1\ell(x)=3s_{m(x)+1}$, and for every cube $X$ of $\omega$,
\begin{equation}\label{9.5:eq-walk-influence-2}
X^{+}\subset\{p\in\mathbb{T}:\ |p-c(X)|\le(2L+6)\hat{u}_{s(X)}\}.
\end{equation}
\end{definition}

\begin{proof}
The ingredients of $\rho_{\times}$ are positive, so $\rho_{\times}>0$. Substituting
$\rho_{\times}=L\mathsf{u}(20\mathsf{q}+5L+7)$ and $\mathsf{u}=\max\{\mathfrak{m},M_1,2L\}$ into
$2^{d+3}L\rho_{\times}$ gives the right-hand side of \eqref{9.5:eq-walk-influence-a}, and with
$\rho=\rho_{\times}$ the inequality $M\ge 2^{d+3}L\rho$ of \eqref{9.5:eq-common-point-a} is
\eqref{9.5:eq-walk-influence-a}.

Dividing \eqref{9.5:eq-walk-influence-a} by $2^{d+3}L$,
\[
\rho_{\times}\le\frac{M}{2^{d+3}L}=\frac{1}{2^{d+1}}\cdot\frac{M}{4L}=\frac{r_0}{2^{d+1}},
\]
the first inequality of \eqref{9.5:eq-walk-influence-1}; in particular $0<\rho_{\times}\le r_0$, the range
of radii admitted in Lemma~\ref{9.5:lem-common-point}. Dividing \eqref{9.5:eq-walk-influence-a}
by $L$ instead,
\[
\frac{M}{L}\ \ge\ 2^{d+3}\rho_{\times}=2^{d+3}L\,\mathsf{u}\,(20\mathsf{q}+5L+7)
\ \ge\ (20\mathsf{q}+5L+7)\,\mathsf{u}\ >\ (10\mathsf{q}+5L+6)\,\mathsf{u},
\]
since $2^{d+3}L\ge 1$, $\mathsf{u}>0$ and $20\mathsf{q}+5L+7>10\mathsf{q}+5L+6$ (as $\mathsf{q}\ge 0$). This
is the second inequality of \eqref{9.5:eq-walk-influence-1}.

By definition $s_{m(x)+1}=M_1L^{m(x)+1}\eta$, so $3LM_1\ell(x)=3M_1L^{m(x)+1}\eta=3s_{m(x)+1}$.

For \eqref{9.5:eq-walk-influence-2}, recall that $|\cdot|$ is the sup-norm distance, so a cube of centre
$c$ and radius $r$ is $\{p:|p-c|\le r\}$. The cube $X^{(5)}$ has centre $c(X)$ and radius
$r(X)+5\hat{u}_{s(X)}\le 10\hat{u}_{s(X)}$ by \eqref{9.5:eq-walk-influence-b}(i), and
$10\le 2L+6$ because $L\ge 2$. The set $X^{2L+1}$ is $X$ enlarged about $c(X)$ by $2L+1$ block units, of
radius at most
$r(X)+(2L+1)\hat{u}_{s(X)}\le(2L+6)\hat{u}_{s(X)}$, again by $r(X)\le 5\hat{u}_{s(X)}$; for a base cube
$\square$ its radius is $(2L+2)\hat{u}_{s(\square)}\le(2L+6)\hat{u}_{s(\square)}$. Hence both members of
$X^{+}=X^{(5)}\cup X^{2L+1}$ lie in the closed sup-norm ball of centre $c(X)$ and radius
$(2L+6)\hat{u}_{s(X)}$.
\end{proof}

\begin{lemma}[Influence sets lie near every connecting graph]\label{9.5:lem-tether}
Assume the conditions collected in \eqref{9.5:eq-scaled-length-a} and
\eqref{9.5:eq-shell-grids-a}, the walk-geometry conditions \eqref{9.5:eq-walk-influence-b}
and the condition \eqref{9.5:eq-walk-influence-a}. Let
$\omega=((\alpha_i,X_i))_{0\le i\le n}$ be a walk of a refinement $\hat{\boldsymbol{\Omega}}$,
and let $\Gamma$ be any graph that meets every domain $X_i$. Then
\begin{equation}\label{9.5:eq-tether-1}
\begin{split}
&S^{\sharp}(\omega)\cup S^{\mathrm{mix}}(\omega)\\
&\quad\cup\Bigl\{\Delta:\ \operatorname{dist}(p,\Delta)\le \mathsf{r}_e\,\ell(p)
\text{ for some } p\in{\textstyle\bigcup_i X_i}\Bigr\}
\subset \mathcal{S}_{\rho_{\times}}(\Gamma),
\end{split}
\end{equation}
with $S^{\sharp}(\omega)$, $S^{\mathrm{mix}}(\omega)$ and $\rho_{\times}$
as in \eqref{9.5:eq-walk-influence-d}, and $\mathsf{r}_e$ as in
Definition~\ref{9.5:def-walk-influence}.
\end{lemma}

\begin{proof}
We write $|\Gamma|$ for the set of points of $\Gamma$, and, for a point $y$ and a cell
$\Delta$, $\mathsf{d}(y,\Delta):=\inf_{z\in\Delta}\mathsf{d}(y,z)$. Since
$\mathcal{S}_{\rho_{\times}}(\Gamma)$ consists of the cells at scaled distance less than
$\rho_{\times}$ from $\Gamma$, it suffices to find, for each cell $\Delta$ of the left side
of \eqref{9.5:eq-tether-1}, a point $y\in|\Gamma|$ with $\mathsf{d}(y,\Delta)<\rho_{\times}$.
In each application of Lemma~\ref{9.5:lem-short-segments} below the radius $\varrho$
satisfies $\varrho<M/L$, as that lemma requires; this is supplied by
\eqref{9.5:eq-walk-influence-a}.

For every $i$ fix a point $y_i\in|\Gamma|\cap X_i$, which exists because $\Gamma$ meets $X_i$.
Fix $i$, let
$X^{*}$ be a cover cube of $X_i$ of the largest scale $s^{*}$ among the cover cubes of $X_i$,
and put $c^{*}:=c(X^{*})$. The anchoring and Step~1 below hold for every $i$; Steps~2 and~3
apply them to the index $i$ with $X\subset X_i$.

\emph{Anchoring.} By the anchoring of walk cubes in
Definition~\ref{9.5:def-walk-influence}, $c^{*}\in\hat{\Omega}_{s^{*}}$, so that
$\hat{m}(c^{*})\ge s^{*}$ and
\begin{equation*}
\ell(c^{*})\ge \hat{\ell}(c^{*})\ge L^{s^{*}}\hat{\eta}.
\end{equation*}
Since $\hat{u}_{s}=\mathfrak{m}L^{s}\hat{\eta}$, multiplying by $\mathfrak{m}$ gives
$\hat{u}_{s^{*}}\le \mathfrak{m}\,\ell(c^{*})$. In the
same way $\hat{u}_{s(X)}\le \mathfrak{m}\,\ell(c(X))$ for every cube $X$ of $\omega$.

\emph{Step 1: points of $X_i$.} Let $p\in X_i$. By \eqref{9.5:eq-walk-influence-b} there is a
chain $X^{*}=X^{0},X^{1},\dots,X^{a}\ni p$ of distinct cover cubes of $X_i$, with
$a\le\mathsf{q}-1$, consecutive cubes intersecting, all of radius at most
$5\hat{u}_{s^{*}}$ (the scale $s^{*}$ being the largest). Going from $c^{*}$ to a point of
$X^{0}$ costs at most one radius, and passing through each further cube of the chain at most
one diameter; hence
\begin{equation*}
|p-c^{*}|\le 5\hat{u}_{s^{*}}(1+2a)<10\,\mathsf{q}\,\hat{u}_{s^{*}}
\le 10\,\mathsf{q}\,\mathfrak{m}\,\ell(c^{*}),
\end{equation*}
where $1+2a\le 2\mathsf{q}-1<2\mathsf{q}$ and the anchoring were used. Lemma~\ref{9.5:lem-short-segments},
applied with $c=c^{*}$, $x=p$ and $\varrho=10\,\mathsf{q}\,\mathfrak{m}<M/L$, gives, since
$\mathfrak{m}\le\mathsf{u}$, where $\mathsf{u}$ is the constant of
Definition~\ref{9.5:def-walk-influence},
\begin{equation}\label{9.5:eq-tether-2}
\mathsf{d}(c^{*},p)\le 10\,\mathsf{q}\,L\,\mathsf{u}\qquad(p\in X_i).
\end{equation}
In particular \eqref{9.5:eq-tether-2} holds for $p=y_i$, and for $p=c(X)$ whenever $X$ is a
cube of $\omega$ contained in $X_i$.

Now let $p\in X_i$ and let $\Delta$ be a cell with
$\operatorname{dist}(p,\Delta)\le\mathsf{r}_e\,\ell(p)$. The cell $\Delta$ is a closed cube,
so there is $x\in\Delta$ with $|x-p|=\operatorname{dist}(p,\Delta)$.
Lemma~\ref{9.5:lem-short-segments} at the centre $p$, with $\varrho=\mathsf{r}_e$, gives
\begin{equation*}
\mathsf{d}(p,\Delta)\le\mathsf{d}(p,x)\le L\mathsf{r}_e\le L\mathsf{u}.
\end{equation*}
By the triangle
inequality for $\mathsf{d}$ and \eqref{9.5:eq-tether-2} applied to $y_i$ and to $p$,
\begin{equation*}
\mathsf{d}(y_i,\Delta)\le \mathsf{d}(y_i,c^{*})+\mathsf{d}(c^{*},p)+\mathsf{d}(p,\Delta)
\le L\mathsf{u}(20\mathsf{q}+1)<\rho_{\times}.
\end{equation*}
Hence the third set on the left of \eqref{9.5:eq-tether-1} lies in
$\mathcal{S}_{\rho_{\times}}(\Gamma)$.

\emph{Step 2: the set $S^{\mathrm{mix}}(\omega)$.} Let $\Delta\in S^{\mathrm{mix}}(\omega)$; by
\eqref{9.5:eq-walk-influence-d}, $\operatorname{dist}(\Delta,X^{+})\le 3L\hat{u}_{s(X)}$ for a
cube $X$ of $\omega$ with $X\subset X_i$ for some $i$. By the definition of the enlargement
$X^{+}$, every point of $X^{+}$ lies within $(2L+6)\hat{u}_{s(X)}$ of $c(X)$. Therefore, for
every $z\in X^{+}$, $\operatorname{dist}(c(X),\Delta)\le|c(X)-z|+\operatorname{dist}(z,\Delta)$,
and taking the infimum over $z$ and using the anchoring,
\begin{equation*}
\operatorname{dist}(c(X),\Delta)\le (5L+6)\hat{u}_{s(X)}\le (5L+6)\,\mathfrak{m}\,\ell(c(X)).
\end{equation*}
Lemma~\ref{9.5:lem-short-segments} at the centre $c(X)$, with $\varrho=(5L+6)\mathfrak{m}$ and
$x$ a point of $\Delta$ nearest to $c(X)$, gives $\mathsf{d}(c(X),\Delta)\le L(5L+6)\mathsf{u}$.
With Step~1, applied to $y_i$ and to $c(X)$,
\begin{equation*}
\begin{split}
\mathsf{d}(y_i,\Delta)&\le\mathsf{d}(y_i,c^{*})+\mathsf{d}(c^{*},c(X))+\mathsf{d}(c(X),\Delta)\\
&\le L\mathsf{u}(20\mathsf{q}+5L+6)<\rho_{\times}.
\end{split}
\end{equation*}

\emph{Step 3: the set $S^{\sharp}(\omega)$.} Let $\Delta\in S^{\sharp}(\omega)$; by
\eqref{9.5:eq-walk-influence-d} there is $x\in\operatorname{loc}^{+}(\omega)$ with
$\operatorname{dist}(x,\Delta)<3LM_1\ell(x)$. By the third of the walk-geometry conditions
\eqref{9.5:eq-walk-influence-b}, $x\in X^{+}$ for a cube $X$ of $\omega$ with $X\subset X_i$ for
some $i$. As in Step~2, $|x-c(X)|\le(2L+6)\hat{u}_{s(X)}$, which by the anchoring is at most
$(2L+6)\,\mathfrak{m}\,\ell(c(X))$,
so Lemma~\ref{9.5:lem-short-segments} at $c(X)$, with $\varrho=(2L+6)\mathfrak{m}$, gives
$\mathsf{d}(c(X),x)\le L(2L+6)\mathsf{u}$. Lemma~\ref{9.5:lem-short-segments} at the centre $x$,
with $\varrho=3LM_1<M/L$ and a point of $\Delta$ nearest to $x$, gives
$\mathsf{d}(x,\Delta)\le 3L^{2}M_1\le 3L^{2}\mathsf{u}$. Adding these to Step~1,
\begin{equation*}
\begin{split}
\mathsf{d}(y_i,\Delta)&\le\mathsf{d}(y_i,c^{*})+\mathsf{d}(c^{*},c(X))
+\mathsf{d}(c(X),x)+\mathsf{d}(x,\Delta)\\
&\le L\mathsf{u}(20\mathsf{q}+2L+6+3L)=L\mathsf{u}(20\mathsf{q}+5L+6)<\rho_{\times}.
\end{split}
\end{equation*}
In each of the three steps $y_i\in|\Gamma|$, so every cell on the left of
\eqref{9.5:eq-tether-1} lies in $\mathcal{S}_{\rho_{\times}}(\Gamma)$.
\end{proof}

\begin{proposition}[The mixed-size counting bound]\label{9.5:prop-mixed-counting}
Assume the hypotheses \textup{(H1)}--\textup{(H3)} of \eqref{9.5:eq-scaled-length-a}, the
hypothesis \textup{(H4)} of \eqref{9.5:eq-shell-grids-a}, the geometry of one walk
\eqref{9.5:eq-walk-influence-b}, $L\ge 7$, and the hypothesis \eqref{9.5:eq-walk-influence-a};
for \eqref{9.5:eq-mixed-counting-3} and \eqref{9.5:eq-mixed-counting-4} assume in addition the
end-factor clause of the second walk bound \eqref{9.5:eq-walk-influence-c}.
Let $Z$ be a box (any family of cells of any levels), let $\mathbb{O}\in\mathfrak{O}$, and let
$\omega=((\alpha_i,X_i))_{0\le i\le n}\in\mathcal{W}_{\mathbb{O}}(\boldsymbol{\Omega}(Z))$.
More generally, $\omega$ may be any walk of any admissible refinement
$\hat{\boldsymbol{\Omega}}$ of $\boldsymbol{\Omega}$; parts (a) and (a$'$) below hold for such
$\omega$ with $\operatorname{len}_{\hat{\boldsymbol{\Omega}}}$ in place of
$\operatorname{len}_{\boldsymbol{\Omega}(Z)}$ in the definition of $d_Z(\omega)$. Put
\begin{equation*}
d_Z(\omega):=\inf\bigl\{\operatorname{len}_{\boldsymbol{\Omega}(Z)}(\Gamma):\ \Gamma \text{ a tree
graph meeting } X_0,\dots,X_n\bigr\},
\end{equation*}
the least $\operatorname{len}_{\boldsymbol{\Omega}(Z)}$-length of a tree graph meeting all
localisation domains $X_0,\dots,X_n$ of $\omega$; for all points $y,y'$ it does not exceed the
quantity $d_Z(\omega,y,y')$ of \cite[(3.93)]{B-CMP99b}, the infimum of
$\operatorname{len}_{\boldsymbol{\Omega}(Z)}$ over contours $y,y_1,\dots,y_n,y'$ meeting
$X_0,\dots,X_n$.
Let $S$ stand for either of the influence sets $S^{\mathrm{mix}}(\omega)$, $S^{\sharp}(\omega)$ of
\eqref{9.5:eq-walk-influence-d}. Then the following hold.
\begin{itemize}
\item[(a)] With $C_d^{\mathrm{mix}}=2^{d+3}L$,
\begin{equation}\label{9.5:eq-mixed-counting-1}
\#S^{\mathrm{mix}}(\omega),\ \#S^{\sharp}(\omega)
\le 2^d\max\Bigl\{1,\Bigl\lfloor\frac{8L\,d_Z(\omega)}{M}\Bigr\rfloor\Bigr\}
\le C_d^{\mathrm{mix}}\Bigl(\frac{d_Z(\omega)}{M}+1\Bigr).
\end{equation}
\item[(a$'$)] Either $\#S\le 2^d$, and then the cells of $S$ have a common point whenever
$d_Z(\omega)<M/(4L)$; or
\begin{equation}\label{9.5:eq-mixed-counting-2}
d_Z(\omega)\ge \frac{\#S\cdot M}{2^{d+3}L}.
\end{equation}
\item[(b)] Let $D=(L+2)^d-L^d$ and $c_d^{\mathrm{mix}}=(2D^2)^{2^{d+3}L}$. For every cell
$\Delta_x$ and every $\mathsf{m}\in\mathbb{N}$, the families $\mathsf{Y}$ for which there exist
a box $Z$ and a walk $\omega$ of $\boldsymbol{\Omega}(Z)$ having a localisation domain that
meets $\Delta_x$, with $S^{\mathrm{mix}}(\omega)\supset\mathsf{Y}$ (or
$S^{\sharp}(\omega)\supset\mathsf{Y}$) and $d_Z(\omega)\le\mathsf{m}M$, number at most
$(c_d^{\mathrm{mix}})^{\mathsf{m}+1}$.
\end{itemize}
Moreover, let $\mathsf{m}\in\mathbb{N}$, let $x$ be a point with cell $\Delta_x\ni x$, and for a
point $y'$ with cell $\Delta_{y'}\ni y'$ put
\begin{equation*}
\mathsf{d}_Z(\mathsf{Y};x,y'):=\inf\bigl\{d_Z(\omega,x,y'):\ \omega \text{ a walk of }
\boldsymbol{\Omega}(Z) \text{ with } S^{\mathrm{mix}}(\omega)\supset\mathsf{Y}\bigr\}.
\end{equation*}
Then
\begin{gather}
\label{9.5:eq-mixed-counting-3}
\#\bigl\{\mathsf{Y}:\ \mathsf{d}_Z(\mathsf{Y};x,y')\le\mathsf{m}M \text{ for some } y'\bigr\}
\le (c_d^{\mathrm{mix}})^{\mathsf{m}+1},\\
\label{9.5:eq-mixed-counting-4}
\mathsf{d}_Z(\mathsf{Y};x,y')\ge
\mathsf{t}_{\times}\bigl(\mathsf{Y}\cup\{\Delta_x,\Delta_{y'}\}\bigr).
\end{gather}
\end{proposition}

\begin{proof}
Throughout, $\rho_{\times}$ and $\mathsf{t}_{\times}=\mathsf{t}_{\rho_{\times}}$ are as in
\eqref{9.5:eq-walk-influence-d}, where $\rho_{\times}$ is fixed with
$0<\rho_{\times}\le M/(4L)$; hence
Propositions~\ref{9.5:prop-cells-near-graph} and~\ref{9.5:prop-families-reach} apply with
$\rho=\rho_{\times}$.

\emph{Part (a).} Let $\varepsilon>0$. By the definition of $d_Z(\omega)$ as an infimum there is
a tree graph $\Gamma$ meeting all of $X_0,\dots,X_n$ with
$\operatorname{len}_{\boldsymbol{\Omega}(Z)}(\Gamma)\le d_Z(\omega)+\varepsilon$. By
Lemma~\ref{9.5:lem-refinement-lengthens}, refinements only lengthen scaled distances, so
\begin{equation*}
\mathsf{L}(\Gamma)\le\operatorname{len}_{\boldsymbol{\Omega}(Z)}(\Gamma)\le
d_Z(\omega)+\varepsilon .
\end{equation*}
Since $\Gamma$ meets every $X_i$, Lemma~\ref{9.5:lem-tether}, whose hypotheses are among those
assumed here, gives
$S^{\sharp}(\omega)\cup S^{\mathrm{mix}}(\omega)\subset\mathcal{S}_{\rho_{\times}}(\Gamma)$;
in particular $S\subset\mathcal{S}_{\rho_{\times}}(\Gamma)$. Proposition~\ref{9.5:prop-cells-near-graph}
with $\rho=\rho_{\times}$, together with the monotonicity of $t\mapsto\lfloor t\rfloor$, gives
\begin{equation*}
\#S\le\#\mathcal{S}_{\rho_{\times}}(\Gamma)\le
2^d\max\Bigl\{1,\Bigl\lfloor\frac{8L\,\mathsf{L}(\Gamma)}{M}\Bigr\rfloor\Bigr\}
\le 2^d\max\Bigl\{1,\Bigl\lfloor\frac{8L(d_Z(\omega)+\varepsilon)}{M}\Bigr\rfloor\Bigr\}.
\end{equation*}
The function $t\mapsto\lfloor t\rfloor$ is right-continuous, so for $\varepsilon>0$ small enough
$\lfloor 8L(d_Z(\omega)+\varepsilon)/M\rfloor=\lfloor 8L\,d_Z(\omega)/M\rfloor$; this is the first
inequality of \eqref{9.5:eq-mixed-counting-1}. For the second, $\lfloor t\rfloor\le t$ and
$\max\{1,t\}\le 1+t$ for $t\ge0$, and $2^d\le 2^{d+3}L$, whence
\begin{equation*}
2^d\max\Bigl\{1,\Bigl\lfloor\frac{8L\,d_Z(\omega)}{M}\Bigr\rfloor\Bigr\}
\le 2^d+2^{d+3}L\,\frac{d_Z(\omega)}{M}
\le 2^{d+3}L\Bigl(\frac{d_Z(\omega)}{M}+1\Bigr).
\end{equation*}

\emph{Part (a$'$).} This follows from the last two clauses of
Proposition~\ref{9.5:prop-cells-near-graph}, applied to the graphs $\Gamma$ of part~(a). Suppose
first $d_Z(\omega)<M/(4L)$, and choose $\varepsilon>0$ with $d_Z(\omega)+\varepsilon<M/(4L)$;
then $\mathsf{L}(\Gamma)<M/(4L)$, so by the short-class clause $\#\mathcal{S}_{\rho_{\times}}(\Gamma)\le
2^d$ and the cells of $\mathcal{S}_{\rho_{\times}}(\Gamma)$ have a common point; since
$S\subset\mathcal{S}_{\rho_{\times}}(\Gamma)$, also $\#S\le 2^d$ and the cells of $S$ have a
common point. Suppose next $\#S>2^d$. Then $\#\mathcal{S}_{\rho_{\times}}(\Gamma)\ge\#S>2^d$ for
every $\varepsilon>0$, and the last clause of Proposition~\ref{9.5:prop-cells-near-graph} gives
\begin{equation*}
d_Z(\omega)+\varepsilon\ge\mathsf{L}(\Gamma)\ge
\frac{\#\mathcal{S}_{\rho_{\times}}(\Gamma)\cdot M}{2^{d+3}L}\ge\frac{\#S\cdot M}{2^{d+3}L};
\end{equation*}
letting $\varepsilon\to0$ gives \eqref{9.5:eq-mixed-counting-2}.

\emph{Part (b).} Let $Z$, $\omega$ and $\mathsf{Y}$ be as in (b), let $\varepsilon>0$ and let
$\Gamma$ be as in part~(a). Some localisation domain $X_i$ meets $\Delta_x$, so
$\operatorname{dist}(p,\Delta_x)=0$ for a point $p\in X_i$, and $\Delta_x$ belongs to the third set
on the left of the inclusion of Lemma~\ref{9.5:lem-tether}; hence
$\mathsf{F}:=\mathsf{Y}\cup\{\Delta_x\}\subset\mathcal{S}_{\rho_{\times}}(\Gamma)$. Recall that
$\mathsf{t}_{\times}=\mathsf{t}_{\rho_{\times}}$ by \eqref{9.5:eq-walk-influence-d}, and that, by
the definition of the loose tree length used in Proposition~\ref{9.5:prop-families-reach},
$\mathsf{t}_{\rho}(\mathsf{F})\le\mathsf{L}(\Gamma)$ for every graph $\Gamma$ with
$\mathsf{F}\subset\mathcal{S}_{\rho}(\Gamma)$. Hence
$\mathsf{t}_{\times}(\mathsf{F})\le d_Z(\omega)+\varepsilon$ for every $\varepsilon>0$, and so
$\mathsf{t}_{\times}(\mathsf{F})\le d_Z(\omega)\le\mathsf{m}M$. The family $\mathsf{F}$ is finite by
part~(a). Proposition~\ref{9.5:prop-families-reach} with $\Delta_0=\Delta_x$ and
$\tau=\mathsf{m}M$ shows that there are at most
$(2D^2)^{2^d\max\{1,8L\mathsf{m}\}}$ such families $\mathsf{F}$. The map
$\mathsf{Y}\mapsto\mathsf{Y}\cup\{\Delta_x\}$ is at most two-to-one (the preimages of
$\mathsf{F}$ are among $\mathsf{F}$ and $\mathsf{F}\setminus\{\Delta_x\}$), and
$2\le(2D^2)^{2^d}$ because $D\ge1$. Therefore the number of families $\mathsf{Y}$ is at most
$(2D^2)^{2^d(\max\{1,8L\mathsf{m}\}+1)}$. Finally
$2^d(\max\{1,8L\mathsf{m}\}+1)\le 2^{d+3}L(\mathsf{m}+1)$: for $\mathsf{m}\ge1$ the left side is
$2^{d+3}L\mathsf{m}+2^d$, and for $\mathsf{m}=0$ it is $2^{d+1}$; in both cases it is at most
$2^{d+3}L(\mathsf{m}+1)$. Since $2D^2\ge1$, the number of families $\mathsf{Y}$ is at most
$(c_d^{\mathrm{mix}})^{\mathsf{m}+1}$.

\emph{The bounds \eqref{9.5:eq-mixed-counting-3} and \eqref{9.5:eq-mixed-counting-4}.} Let
$\omega$ be a walk of $\boldsymbol{\Omega}(Z)$ with $S^{\mathrm{mix}}(\omega)\supset\mathsf{Y}$,
and let $\varepsilon>0$. Since $d_Z(\omega,x,y')$ is the infimum of
$\operatorname{len}_{\boldsymbol{\Omega}(Z)}$ over contours $x,y_1,\dots,y_n,y'$ meeting all
$X_i$, there is such a contour whose graph $\Gamma$ satisfies
$\operatorname{len}_{\boldsymbol{\Omega}(Z)}(\Gamma)\le d_Z(\omega,x,y')+\varepsilon$. By
Lemma~\ref{9.5:lem-refinement-lengthens}, $\mathsf{L}(\Gamma)\le d_Z(\omega,x,y')+\varepsilon$.
By the end-factor clause of the second walk bound in \eqref{9.5:eq-walk-influence-c}, the block
of $x$ fixed in that clause meets $X_0$ and the block of $y'$ fixed there meets $X_n$; these
blocks lie in $\Delta_x$ and
$\Delta_{y'}$, so $\Delta_x$ meets $X_0$ and $\Delta_{y'}$ meets $X_n$. As in part~(b), both cells
belong to the third set of Lemma~\ref{9.5:lem-tether}, and together with
$\mathsf{Y}\subset S^{\mathrm{mix}}(\omega)$ that lemma gives
$\mathsf{Y}\cup\{\Delta_x,\Delta_{y'}\}\subset\mathcal{S}_{\rho_{\times}}(\Gamma)$. The definition
of $\mathsf{t}_{\times}$ then gives
$\mathsf{t}_{\times}(\mathsf{Y}\cup\{\Delta_x,\Delta_{y'}\})\le d_Z(\omega,x,y')+\varepsilon$;
taking the infimum over $\varepsilon$ and over such $\omega$ proves \eqref{9.5:eq-mixed-counting-4}.
The same argument with $\Delta_{y'}$ dropped shows that
$\mathsf{d}_Z(\mathsf{Y};x,y')\le\mathsf{m}M$ for some $y'$ implies
$\mathsf{t}_{\times}(\mathsf{Y}\cup\{\Delta_x\})\le\mathsf{m}M+\varepsilon$ for every
$\varepsilon>0$, hence $\mathsf{t}_{\times}(\mathsf{Y}\cup\{\Delta_x\})\le\mathsf{m}M$; the
count of part~(b), from Proposition~\ref{9.5:prop-families-reach} and the two-to-one map onward,
then gives \eqref{9.5:eq-mixed-counting-3}.
\end{proof}

\begin{definition}[Cores, box operators and M\"obius pieces]\label{9.5:not-moebius-pieces}
Boxes, the shell grids and the sequence $\boldsymbol{\Omega}(Z)$ induced on a box $Z$ are as in
Definition~\ref{9.5:def-shell-grids}. The operators $\mathbb{O}$ of the fixed list $\mathfrak{O}$
and their walk expansions are as in Definition~\ref{9.5:def-walk-influence}.

\emph{Core.} A \emph{core} is a family $\mathsf{K}\subset\pi^{\mathrm{mix}}$ of cells. We put
\begin{equation*}
K:=\bigcup\mathsf{K},\qquad \sigma_0:=\pi^{\mathrm{mix}}\setminus\mathsf{K},
\end{equation*}
so that $\sigma_0$ is the set of cells off the core.

\emph{Box operators.} Let $Z$ be a box and let $\mathbb{O}\in\mathfrak{O}$. The \emph{box operator}
$\mathbb{O}_Z$ is the operator $\mathbb{O}$ of the induced sequence $\boldsymbol{\Omega}(Z)$,
evaluated at the restriction $(\mathbf{U},\mathbf{J})|_{Z}$ of the gauge field and the source to $Z$.
Here $(\mathbf{U},\mathbf{J})$ ranges over the class of configurations on which Ba{\l}aban's operators
at the sequence $\boldsymbol{\Omega}$ are defined. The conventions are these:
\begin{itemize}
\item[(a)] the kernel of $\mathbb{O}_Z$ vanishes off $Z\times Z$;
\item[(b)] an operator of the list acting on the unit lattice is taken on the set $\tilde{\Lambda}(Z)$ of
free bonds of the unit-lattice margin set $\Lambda(Z)$;
\item[(c)] for the composite operator $T=C^{*}\Delta_kC$ of the list (the operators $C$ and
$\Delta_k$ being those of the list $\mathfrak{O}$ of Definition~\ref{9.5:def-walk-influence}),
with $\Delta_{k,Z}$ the box version of $\Delta_k$ at $Z$, we put
\begin{equation}\label{9.5:eq-moebius-pieces-1}
T_Z(b,b'):=\mathbf{1}_{\tilde{\Lambda}(Z)}(b)\,\bigl(C^{*}\Delta_{k,Z}C\bigr)(b,b')\,
\mathbf{1}_{\tilde{\Lambda}(Z)}(b').
\end{equation}
\end{itemize}

\emph{M\"obius pieces.} Let $\mathsf{Y}\subset\sigma_0$ be a finite family of cells, with union
$Y:=\bigcup\mathsf{Y}$. For each subfamily $\mathsf{Y}'\subset\mathsf{Y}$ write
$Y':=\bigcup\mathsf{Y}'$. The \emph{M\"obius piece} of $\mathbb{O}$ at the core $\mathsf{K}$ and the
family $\mathsf{Y}$ is
\begin{equation}\label{9.5:eq-moebius-pieces-2}
\mathfrak{o}_{\mathsf{K}}(\mathsf{Y})
:=\sum_{\mathsf{Y}'\subset\mathsf{Y}}(-1)^{\#(\mathsf{Y}\setminus\mathsf{Y}')}\,
\mathbb{O}_{K\cup Y'} .
\end{equation}

\emph{The cell of a point.} For $y\in\mathbb{T}$ we write $\Delta_y$ for a cell containing $y$.
When $y$ lies on the boundary of several cells, one of them is selected by a fixed half-open convention.

\emph{Standing assumptions.} In what follows we assume:
\begin{itemize}
\item[(i)] the conditions \eqref{9.5:eq-scaled-length-a} on the nested grids and the compatibility
condition \eqref{9.5:eq-shell-grids-a} on the shell grids;
\item[(ii)] $L\ge 7$ and $R\ge L$, where $R$ is the admissibility parameter of
Definition~\ref{9.5:def-walk-influence};
\item[(iii)] the lower bounds on the auxiliary parameters under which every induced sequence
$\boldsymbol{\Omega}(Z)$ is admissible, as stated in the remark accompanying the shell lemma for
mixed-size boxes (Lemma~\ref{9.5:lem-shell-lemma});
\item[(iv)] the geometric properties \eqref{9.5:eq-walk-influence-b} of one walk and the walk bounds
\eqref{9.5:eq-walk-influence-c}, for the walk expansion of every operator of the list at the induced
sequence $\boldsymbol{\Omega}(Z)$ of every box $Z$;
\item[(v)] the influence-set property \eqref{9.5:eq-walk-influence-a}.
\end{itemize}
\end{definition}

\begin{theorem}[M\"obius decoupling for mixed-size cells]\label{9.5:thm-moebius-decoupling}
Let the core $\mathsf{K}\subset\pi^{\mathrm{mix}}$, its union $K:=\bigcup\mathsf{K}$, the set
$\sigma_0=\pi^{\mathrm{mix}}\setminus\mathsf{K}$, the box operators $\mathbb{O}_Z$, the M\"obius pieces
$\mathfrak{o}_{\mathsf{K}}(\mathsf{Y})$ and the cells $\Delta_y$ be as in
Notation~\ref{9.5:not-moebius-pieces}, under the standing assumptions made there. Assume the floor
\eqref{9.5:eq-walk-influence-a} and the geometry of walks of
Definition~\ref{9.5:def-walk-influence}.

Let $\mathsf{Y}\subset\sigma_0$ be finite. For $\mathsf{Y}'\subset\mathsf{Y}$ put
$Y':=\bigcup\mathsf{Y}'$. For a box $Z$ we write
\begin{equation*}
\mathbb{O}_Z=\sum_{\omega\in\mathcal{W}(\boldsymbol{\Omega}(Z))}\mathbb{O}_Z(\omega)
\end{equation*}
for the walk expansion of Definition~\ref{9.5:def-walk-influence} in the geometry
$\boldsymbol{\Omega}(Z)$. Here the subscript $\mathbb{O}$ of $\mathcal{W}_{\mathbb{O}}$ is suppressed, and
$d_Z(\omega)$ denotes the tree length of $\omega$ in the geometry $\boldsymbol{\Omega}(Z)$ fixed in
that definition.

Assume that, for each of the $2^{\#\mathsf{Y}}$ boxes $Z=K\cup Y'$ with $\mathsf{Y}'\subset\mathsf{Y}$,
the following three hypotheses hold:
\begin{itemize}
\item[(i)] the decay bound for walk terms displayed in \eqref{9.5:eq-walk-influence-c}: for every
$\ell\ge0$ the norms of the terms $\mathbb{O}_Z(\omega)$ with $d_Z(\omega)\ge\ell$ sum to at most
$B_0e^{-\frac12\delta_0\ell}\cdot(\text{scale factors})$;
\item[(ii)] the localisation of the end blocks displayed in \eqref{9.5:eq-walk-influence-c};
\item[(iii)] the locality property of walks displayed in \eqref{9.5:eq-walk-influence-c}. By this
property, the following two things depend on the sequence $\boldsymbol{\Omega}(Z)$ only through the
points of $\operatorname{loc}^{+}(\omega)$ lying in its sets: whether $\omega$ belongs to
$\mathcal{W}(\boldsymbol{\Omega}(Z))$, and the operator built from $\omega$, as a function of
$(\mathbf{U},\mathbf{J})$ restricted to $\operatorname{loc}(\omega)$.
\end{itemize}
Throughout, $(\text{scale factors})$ stands for the scale factors of the bound in~(i).

Then, for every pair of kernel arguments $y,y'$, the following three assertions hold. They are
collected in \eqref{9.5:eq-moebius-decoupling-a}.
\begin{equation}\label{9.5:eq-moebius-decoupling-a}
\begin{aligned}
&\text{(a)}\quad
\mathfrak{o}_{\mathsf{K}}(\mathsf{Y})
=\sum_{\mathsf{Y}'\subset\mathsf{Y}}(-1)^{\#(\mathsf{Y}\setminus\mathsf{Y}')}
\sum_{\substack{\omega\in\mathcal{W}(\boldsymbol{\Omega}(K\cup Y'))\\
S^{\sharp}(\omega)\supset\mathsf{Y}}}\mathbb{O}_{K\cup Y'}(\omega);\\
&\text{(b)}\quad
\|\mathfrak{o}_{\mathsf{K}}(\mathsf{Y})\|
\le 2^{\#\mathsf{Y}}B_0\,
e^{-\frac12\delta_0\mathsf{t}_{\times}(\mathsf{Y}\cup\{\Delta_y,\Delta_{y'}\})}
\cdot(\text{scale factors})\\
&\qquad
\le 2^{\#\mathsf{Y}}C_0^{\mathrm{mix}}B_0\,
e^{-\frac14\delta_0 M d^{\circ}_{\mathrm{mix}}(\mathsf{Y}\cup\{\Delta_y,\Delta_{y'}\})}
\cdot(\text{scale factors}),\\
&\qquad
C_0^{\mathrm{mix}}:=e^{2^{d-2}\delta_0\rho_{\times}};\\
&\text{(c)}\quad
\mathfrak{o}_{\mathsf{K}}(\mathsf{Y})(y,y')=0\ \text{unless }\Delta_y,\ \Delta_{y'}\text{ and all
cells of }\mathsf{Y}\\
&\qquad\qquad\text{lie in one wall-component of }\mathsf{K}\cup\mathsf{Y}.
\end{aligned}
\end{equation}
The three assertions are to be read as follows.
\begin{itemize}
\item[(a)] The identity holds, with the sum over $\omega$ absolutely convergent.
\item[(b)] The bound holds in every norm, between the localisations at $y$ and $y'$, in which~(i)
holds for the $2^{\#\mathsf{Y}}$ geometries $\boldsymbol{\Omega}(K\cup Y')$.
Definition~\ref{9.5:def-walk-influence} states hypothesis~(i) in the first-order norms at unit scale
everywhere, in all norms for outputs in the core, and pointwise for the kernels of the unit-lattice
operators; (b) is used in those norms as stated there. The constant $C_0^{\mathrm{mix}}$ is formed from
$\delta_0$ and $\rho_{\times}$ alone and does not depend on $M$.
\item[(c)] The $(y,y')$-kernel of $\mathfrak{o}_{\mathsf{K}}(\mathsf{Y})$ vanishes unless the stated
cells lie in one wall-component. A wall-component is a maximal subfamily connected through cells that
share a wall.
\end{itemize}
\end{theorem}

\begin{proof}
\emph{Proof of~(a).} Substitute the walk expansion of each $\mathbb{O}_{K\cup Y'}$ into the
definition of $\mathfrak{o}_{\mathsf{K}}(\mathsf{Y})$ in Notation~\ref{9.5:not-moebius-pieces}. This
gives
\begin{equation}\label{9.5:eq-moebius-decoupling-1}
\mathfrak{o}_{\mathsf{K}}(\mathsf{Y})
=\sum_{\mathsf{Y}'\subset\mathsf{Y}}(-1)^{\#(\mathsf{Y}\setminus\mathsf{Y}')}
\sum_{\omega\in\mathcal{W}(\boldsymbol{\Omega}(K\cup Y'))}\mathbb{O}_{K\cup Y'}(\omega).
\end{equation}
The double sum is finite in $\mathsf{Y}'$. It is absolutely convergent in $\omega$ by hypothesis~(i),
applied with $\ell=0$. It may therefore be rearranged.

The right side of (a) is the part of \eqref{9.5:eq-moebius-decoupling-1} formed by the pairs
$(\mathsf{Y}',\omega)$ with $S^{\sharp}(\omega)\supset\mathsf{Y}$. It remains to show that the pairs
$(\mathsf{Y}',\omega)$ with $\mathsf{Y}'\subset\mathsf{Y}$, $\omega\in\mathcal{W}(\boldsymbol{\Omega}(K\cup Y'))$
and $S^{\sharp}(\omega)\not\supset\mathsf{Y}$ contribute zero.

Fix a total order of $\pi^{\mathrm{mix}}$. For such a pair let $\Delta_*$ be the first cell of
$\mathsf{Y}\setminus S^{\sharp}(\omega)$, which is non-empty. Put
\begin{equation*}
\mathsf{Y}^{\times}:=\mathsf{Y}'\,\triangle\,\{\Delta_*\},\qquad Y^{\times}:=\textstyle\bigcup\mathsf{Y}^{\times},
\qquad Z:=K\cup Y',\qquad Z^{\times}:=K\cup Y^{\times}.
\end{equation*}

\emph{Claim.} For every $n$, the sets $\Omega_n(Z)$ and $\Omega_n(Z^{\times})$ contain the same points
of $\operatorname{loc}^{+}(\omega)$. The same holds for $\Lambda(Z)$ and $\Lambda(Z^{\times})$.

\emph{Proof of the claim.} Let $x\in\operatorname{loc}^{+}(\omega)$, and recall that $m(x)$ is its
ambient level.

Let $n>m(x)$. Then $x\notin\Omega_n$, by the definition of $m(x)$. By
Lemma~\ref{9.5:lem-shell-lemma}\,(s1), $\Omega_n\supset\Omega_n(Z)$ and
$\Omega_n\supset\Omega_n(Z^{\times})$. Hence $x$ lies in neither set.

Let $n\le m(x)+1$. Since $\Delta_*\notin S^{\sharp}(\omega)$, the definition of $S^{\sharp}(\omega)$ in
\eqref{9.5:eq-walk-influence-d} gives
\begin{equation*}
\operatorname{dist}(x,\Delta_*)\ge 3s_{m(x)+1}\ge 3s_n .
\end{equation*}
The second inequality holds because $s_n=M_1L^n\eta$ increases with $n$. Consequently the open ball
$O:=B(x,3s_n)$ does not meet $\Delta_*$. Both $Z\cup\Delta_*$ and $Z^{\times}\cup\Delta_*$ equal
$K\cup Y'\cup\Delta_*$, so $Z\cup\Delta_*=Z^{\times}\cup\Delta_*$. Intersecting with $O$, which misses
$\Delta_*$, gives $Z\cap O=Z^{\times}\cap O$.

By Lemma~\ref{9.5:lem-shell-lemma}\,(s2), whether $x\in\Omega_n(Z)$ is decided by $\Omega_n$ and by
$Z\cap B(x,3s_n)$. Hence, for $n\le m(x)$, $x\in\Omega_n(Z)$ if and only if $x\in\Omega_n(Z^{\times})$.

Now let $x\in\Omega_{\star}$. Since $\Omega_{\star}\subset\Omega_k$, we have $m(x)=k$, and the case
$n=k+1$ applies. By the same item of Lemma~\ref{9.5:lem-shell-lemma}, whether $x\in\Lambda(Z)$ is
decided by $\Omega_{\star}$ and $Z\cap B(x,3s_{k+1})$. Hence $x\in\Lambda(Z)$ if and only if
$x\in\Lambda(Z^{\times})$. If $x\notin\Omega_{\star}$, neither $\Lambda(Z)$ nor $\Lambda(Z^{\times})$
contains $x$. This proves the claim.

\emph{The walk term is unchanged.} We apply the locality property~(iii). For the unit-lattice
operators it is read with the sets $\Lambda(\cdot)$ among the sets through which the operator depends
on the geometry. For these operators the building blocks by which they differ from the other
operators of the list are local operators, and whether a bond $b$ belongs to the free bonds
$\tilde{\Lambda}(Z)$, respectively $\tilde{\Lambda}(Z^{\times})$, is decided inside the influence set
$S^{\sharp}(\omega)$ of \eqref{9.5:eq-walk-influence-d}. By the claim
and~(iii), $\omega\in\mathcal{W}(\boldsymbol{\Omega}(Z^{\times}))$. Moreover, the operator built from
$\omega$ on $\boldsymbol{\Omega}(Z^{\times})$ equals the one built on $\boldsymbol{\Omega}(Z)$, as a
function of $(\mathbf{U},\mathbf{J})$ restricted to $\operatorname{loc}(\omega)$.

The local operators read $(\mathbf{U},\mathbf{J})$ inside $\Omega_0(Z)\subset Z$, respectively
$\Omega_0(Z^{\times})\subset Z^{\times}$. We have
$Z\cap\operatorname{loc}(\omega)=Z^{\times}\cap\operatorname{loc}(\omega)$. Indeed,
$\operatorname{loc}(\omega)\subset\operatorname{loc}^{+}(\omega)$ misses $\Delta_*$ by the inequality
above, and $Z$, $Z^{\times}$ differ only inside $\Delta_*$, since $Z\cup\Delta_*=Z^{\times}\cup\Delta_*$.
Therefore
\begin{equation*}
\mathbb{O}_Z(\omega)=\mathbb{O}_{Z^{\times}}(\omega).
\end{equation*}

\emph{The pairs cancel.} The set $S^{\sharp}(\omega)$ is a function of the point set
$\operatorname{loc}^{+}(\omega)$ and of the ambient sequence $\boldsymbol{\Omega}$ only. Hence it is
the same for the pair $(\mathsf{Y}^{\times},\omega)$. In particular
$S^{\sharp}(\omega)\not\supset\mathsf{Y}$ still holds. Also $\mathsf{Y}^{\times}\subset\mathsf{Y}$,
because $\Delta_*\in\mathsf{Y}$. The cell $\Delta_*$ computed from $(\mathsf{Y}^{\times},\omega)$ is
again $\Delta_*$.

Thus $(\mathsf{Y}',\omega)\mapsto(\mathsf{Y}^{\times},\omega)$ maps the class of pairs under
consideration into itself, and applying it twice returns $\mathsf{Y}'$. It is therefore an
involution, and it has no fixed point, since $\mathsf{Y}^{\times}\neq\mathsf{Y}'$. It reverses the sign
$(-1)^{\#(\mathsf{Y}\setminus\mathsf{Y}')}$, because $\#(\mathsf{Y}\setminus\mathsf{Y}^{\times})$ and
$\#(\mathsf{Y}\setminus\mathsf{Y}')$ differ by one. It also preserves the term, by the preceding
paragraph. By absolute convergence the contributions of these pairs cancel in pairs, and (a) follows.

\emph{Proof of~(b).} Fix $\mathsf{Y}'\subset\mathsf{Y}$ and $Z=K\cup Y'$. Let
$\omega\in\mathcal{W}(\boldsymbol{\Omega}(Z))$, with domains $X_0,\dots,X_n$, be a walk that survives
in~(a), that is, $S^{\sharp}(\omega)\supset\mathsf{Y}$, and whose $(y,y')$-block is non-zero.

For a point $y$ we write $\hat{j}(y)$ for the level of the localisation block of $y$ in the geometry
$\boldsymbol{\Omega}(Z)$, and $\hat{\Delta}(y)$ for the localisation block of $y$, the block of
level $\hat{j}(y)$ containing $y$. By the localisation of the end
blocks~(ii), $\hat{\Delta}(y)$ meets $X_0$ or lies within $\mathsf{r}_e\hat{\ell}$ of it, and
$\hat{\Delta}(y')$ meets $X_n$ or lies within $\mathsf{r}_e\hat{\ell}$ of it.

By Lemma~\ref{9.5:lem-shell-lemma}\,(s1), $\Omega_n(Z)\subset\Omega_n$ for every $n$. So the refined
level of each point does not exceed its ambient level, and $\hat{\ell}\le\ell$. Hence the cell
$\Delta_y$ has level $m(y)\ge\hat{j}(y)$. Since $M=L^{m'}$ is an integer power of $L$
(Notation~\ref{0.2:not-glossary}), a cell of level $m(y)$ is a union of blocks of every level
$\le m(y)$. Hence $\hat{\Delta}(y)\subset\Delta_y$, and likewise $\hat{\Delta}(y')\subset\Delta_{y'}$.
Consequently there is a point $p\in\bigcup_iX_i$ with
\begin{equation*}
\operatorname{dist}(p,\Delta_y)\le\operatorname{dist}(p,\hat{\Delta}(y))
\le\mathsf{r}_e\hat{\ell}(p)\le\mathsf{r}_e\ell(p),
\end{equation*}
and the same holds for $y'$. So $\Delta_y$ and $\Delta_{y'}$ belong to the third set on the left side
of the inclusion of Lemma~\ref{9.5:lem-tether}.

Let $\Gamma$ be any graph meeting all $X_i$. Lemma~\ref{9.5:lem-tether}, applied to $\omega$ as a
walk of the refinement $\boldsymbol{\Omega}(Z)$, gives
$S^{\sharp}(\omega)\cup\{\Delta_y,\Delta_{y'}\}\subset\mathcal{S}_{\rho_{\times}}(\Gamma)$. Hence
$\mathsf{Y}\cup\{\Delta_y,\Delta_{y'}\}\subset\mathcal{S}_{\rho_{\times}}(\Gamma)$. In particular this
holds when $\Gamma$ is a tree graph meeting all $X_i$; for such $\Gamma$, which sees every cell of this
family at scaled distance $<\rho_{\times}$, the definition of the loose tree length
$\mathsf{t}_{\times}=\mathsf{t}_{\rho_{\times}}$ gives
\begin{equation*}
\mathsf{L}(\Gamma)\ge\mathsf{t}_{\times}(\mathsf{Y}\cup\{\Delta_y,\Delta_{y'}\})=:\mathsf{t}.
\end{equation*}
By Lemma~\ref{9.5:lem-refinement-lengthens}, every such tree has
$\operatorname{len}_{\boldsymbol{\Omega}(Z)}(\Gamma)\ge\mathsf{L}(\Gamma)\ge\mathsf{t}$. Since
$d_Z(\omega)$ is the least $\operatorname{len}_{\boldsymbol{\Omega}(Z)}$-length of a tree meeting all
$X_i$, we obtain $d_Z(\omega)\ge\mathsf{t}$.

By~(a) and the triangle inequality, $\|\mathfrak{o}_{\mathsf{K}}(\mathsf{Y})\|$ is at most the sum over
$\mathsf{Y}'\subset\mathsf{Y}$ of the sums of $\|\mathbb{O}_{K\cup Y'}(\omega)\|$ over such walks. Each
such walk has $d_{K\cup Y'}(\omega)\ge\mathsf{t}$. We apply hypothesis~(i) with $\ell:=\mathsf{t}$ in
each of the $2^{\#\mathsf{Y}}$ geometries $\boldsymbol{\Omega}(K\cup Y')$, with one pair
$B_0,\delta_0$ for all of them \cite{B-CMP99b}. This gives the first inequality of~(b).

For the second inequality, apply Lemma~\ref{9.5:lem-common-point}\,(c) with $\rho:=\rho_{\times}$ and
$\mathsf{F}:=\mathsf{Y}\cup\{\Delta_y,\Delta_{y'}\}$, under the floor
\eqref{9.5:eq-walk-influence-a} assumed in the statement. It gives
$\mathsf{t}\ge\frac12Md^{\circ}_{\mathrm{mix}}(\mathsf{F})-2^{d-1}\rho_{\times}$, and therefore
\begin{equation*}
e^{-\frac12\delta_0\mathsf{t}}
\le e^{2^{d-2}\delta_0\rho_{\times}}\,e^{-\frac14\delta_0Md^{\circ}_{\mathrm{mix}}(\mathsf{F})}
=C_0^{\mathrm{mix}}\,e^{-\frac14\delta_0Md^{\circ}_{\mathrm{mix}}(\mathsf{F})}.
\end{equation*}

\emph{Proof of~(c).} Let $\mathsf{C}_y$ be the wall-component of $\Delta_y$ in
$\mathsf{K}\cup\mathsf{Y}$. By Lemma~\ref{9.5:lem-shell-lemma}\,(s5), the sequence
$\boldsymbol{\Omega}(K\cup Y')$ is the disjoint union of the sequences of the wall-components of
$\mathsf{K}\cup\mathsf{Y}'$. At every level $n$ these sequences lie at mutual distance $\ge 4s_n$.
Free bonds, averages and Dirichlet operators couple nothing across this distance (the last sentence of
Lemma~\ref{9.5:lem-shell-lemma}\,(s5), and \cite{B-CMP99b}). Hence $\mathbb{O}_{K\cup Y'}$ is a
direct sum over these wall-components. Its $(y,y')$-kernel is therefore that of the summand belonging
to the wall-component of $\Delta_y$ in $\mathsf{K}\cup\mathsf{Y}'$, and it vanishes if $\Delta_{y'}$
lies outside that component.

Suppose first that $\Delta_{y'}\notin\mathsf{C}_y$. The wall-component of $\Delta_y$ in
$\mathsf{K}\cup\mathsf{Y}'$ is contained in $\mathsf{C}_y$, so it does not contain $\Delta_{y'}$.
Hence every summand $\mathbb{O}_{K\cup Y'}$ of $\mathfrak{o}_{\mathsf{K}}(\mathsf{Y})$ vanishes at
$(y,y')$.

Suppose next that some $\Delta_1\in\mathsf{Y}\setminus\mathsf{C}_y$ exists. Pair each
$\mathsf{Y}'\not\ni\Delta_1$ with $\mathsf{Y}'\cup\{\Delta_1\}$. The wall-component of $\Delta_y$ is
the same in $\mathsf{K}\cup\mathsf{Y}'$ and in $\mathsf{K}\cup\mathsf{Y}'\cup\{\Delta_1\}$; otherwise
$\Delta_1$ would be wall-connected to $\Delta_y$ inside $\mathsf{K}\cup\mathsf{Y}$, contrary to
$\Delta_1\notin\mathsf{C}_y$. So the two summands have the same $(y,y')$-kernel. They enter with
opposite signs $(-1)^{\#(\mathsf{Y}\setminus\mathsf{Y}')}$ and cancel. Hence
$\mathfrak{o}_{\mathsf{K}}(\mathsf{Y})(y,y')=0$ unless $\Delta_{y'}\in\mathsf{C}_y$ and
$\mathsf{Y}\subset\mathsf{C}_y$, which is~(c).
\end{proof}

\begin{remark}
(1) In the proof of~(a), the influence set $S^{\mathrm{mix}}(\omega)$ of
\eqref{9.5:eq-walk-influence-d} may be used in place of $S^{\sharp}(\omega)$. The argument then rests
on the following fact, which holds by admissibility: for a walk cube $X$, levels $\ge s(X)+2$ are
absent from $X^{+}$ for both sequences $\boldsymbol{\Omega}(Z)$ and $\boldsymbol{\Omega}(Z^{\times})$.
Part~(b) holds verbatim for $S^{\mathrm{mix}}(\omega)$, since Lemma~\ref{9.5:lem-tether} covers both
sets. The set $S^{\sharp}(\omega)$ needs no admissibility argument and does not depend on the box.

(2) Consider one zone, that is, $\Omega_j=\mathbb{T}$ for every $j$ and all cells equal. There the
inequality
\begin{equation*}
\mathsf{t}_{\times}\ge\tfrac12Md^{\circ}_{\mathrm{mix}}-2^{d-1}\rho_{\times}
\end{equation*}
holds, and (b) becomes the M\"obius decoupling bound for cells of one size, with tree lengths
measured in the sup norm.
\end{remark}

\begin{theorem}[Interpolated box operators for mixed-size cells]
\label{9.5:thm-interpolated-operators}
Work in the setting of Notation~\ref{9.5:not-moebius-pieces}: $\mathsf{K}\subset\pi^{\mathrm{mix}}$ is a
core, $K$ its union, $\sigma_0=\pi^{\mathrm{mix}}\setminus\mathsf{K}$, $\mathbb{O}_Z$ the box operator of a
box $Z$, and $\mathfrak{o}_{\mathsf{K}}(\mathsf{Y})$ the M\"obius piece of a finite family
$\mathsf{Y}\subset\sigma_0$. Assume that the walk expansions of the box operators have the three
properties displayed in \eqref{9.5:eq-walk-influence-c}: the first is the bound on walk terms with
constant $B_0$ and rate $\delta_0$, the second is taken as stated there, and the third is taken in
its strengthened form. For $\mathsf{s}\colon\sigma_0\to\mathbb{C}$ put
\begin{equation*}
\mathbb{O}\langle\mathsf{s}\rangle:=\sum_{\mathsf{Y}\subset\sigma_0\ \text{finite}}
\Bigl(\prod_{\Delta\in\mathsf{Y}}\mathsf{s}(\Delta)\Bigr)\,\mathfrak{o}_{\mathsf{K}}(\mathsf{Y}),
\end{equation*}
and put $X_{\kappa}:=(4D^2)^{2^d}e^{2^d\kappa_1}$. For $c>0$ consider the floor
\begin{equation}\label{9.5:eq-interpolated-operators-a}
cM\ \ge\ 16L\bigl[2^d(\kappa_1+\log 4D^2)+\log 2\bigr],
\end{equation}
and assume \eqref{9.5:eq-interpolated-operators-a} at $c=\tfrac12\delta_0$. For points $y,y'$ write
$\Delta(y)\,A\,\Delta(y')$ for the block of an operator $A$ between the localisations at $y$ and
$y'$ in which Theorem~\ref{9.5:thm-moebius-decoupling}(b) is stated; the localisation $\Delta(y)$
is not to be confused with the cell $\Delta_y$ of $y$ fixed in
Notation~\ref{9.5:not-moebius-pieces}. Then the following hold.
\begin{itemize}
\item[(n1)] $\mathbb{O}\langle\mathsf{s}\rangle$ is multi-affine in $\mathsf{s}$. Moreover
$\mathbb{O}\langle\mathbf{1}_{\mathsf{Y}}\rangle=\mathbb{O}_{K\cup Y}$ for every finite
$\mathsf{Y}\subset\sigma_0$, and $\mathbb{O}\langle\mathbf{1}\rangle=\mathbb{O}$.
\item[(n2)] If $\mathsf{s}=0$ off a family $\mathsf{Y}_0\subset\sigma_0$, then the kernel of
$\mathbb{O}\langle\mathsf{s}\rangle$ lives on $(K\cup Y_0)^2$, is block-diagonal over the
wall-components of $\mathsf{K}\cup\mathsf{Y}_0$, and reads $(\mathsf{s},\mathbf{U},\mathbf{J})$ only
there.
\item[(n3)] On the polydisc $|\mathsf{s}(\Delta)|\le e^{\kappa_1}$ for all cells $\Delta\in\sigma_0$,
including cells of every level, however thin the zone they lie in, in every norm in which
Theorem~\ref{9.5:thm-moebius-decoupling}(b) holds,
\begin{equation}\label{9.5:eq-interpolated-operators-b}
\bigl\|\Delta(y)\,\mathbb{O}\langle\mathsf{s}\rangle\,\Delta(y')\bigr\|
\ \le\ C_N^{\mathrm{mix}}\,e^{2^d\kappa_1}\,B_0,
\qquad C_N^{\mathrm{mix}}:=8(4D^2)^{2^d}.
\end{equation}
The constant $C_N^{\mathrm{mix}}$ depends on $d$ and $L$ only; at $d=4$ the factor
$e^{2^d\kappa_1}$ is $e^{16\kappa_1}$, the exponent of \cite[(1.11)]{B-CMP116}.
\item[(n3$'$)] If \eqref{9.5:eq-interpolated-operators-a} holds at $c=\tfrac14\delta_0$, then
on the same polydisc and in the same norms
\begin{equation*}
\bigl\|\Delta(y)\,\mathbb{O}\langle\mathsf{s}\rangle\,\Delta(y')\bigr\|
\le C_N^{\mathrm{mix}}e^{2^d\kappa_1}B_0\,
e^{\frac12\delta_0\rho_{\times}}\,e^{-\frac14\delta_0\mathsf{d}(\Delta_y,\Delta_{y'})}.
\end{equation*}
\item[(n4)] If $\Delta_y$ touches no cell of $\sigma_0$, then on the same polydisc and in the
same norms
\begin{equation*}
\bigl\|\Delta(y)\bigl(\mathbb{O}\langle\mathsf{s}\rangle-\mathbb{O}\langle0\rangle\bigr)\Delta(y')\bigr\|
\ \le\ 4B_0\,e^{-\delta_0M/8L}.
\end{equation*}
\end{itemize}
\end{theorem}

\begin{proof}
The torus $\mathbb{T}$ contains finitely many cells, so $\sigma_0$ is finite and every sum over
$\mathsf{Y}\subset\sigma_0$ below is a finite sum. Throughout this proof we write, as in
Lemma~\ref{9.5:lem-common-point}, $r_0:=M/4L$; this abbreviation is used in this section only and
is unrelated to the depth offset $r_0$ of the glossary. Further, $\mathsf{t}_{\times}=\mathsf{t}_{\rho_{\times}}$ is the loose tree
length at the radius $\rho_{\times}$ of \eqref{9.5:eq-walk-influence-d}. The floor
\eqref{9.5:eq-walk-influence-a} gives
\begin{equation}\label{9.5:eq-interpolated-operators-1}
2\rho_{\times}\ \le\ \frac{M}{2^{d+2}L}\ \le\ \frac{M}{16L},
\end{equation}
so that $0<\rho_{\times}\le r_0$, and Lemmas~\ref{9.5:lem-common-point} and~\ref{9.5:lem-entropy-sum}
apply with $\rho:=\rho_{\times}$.

(n1) Each term of $\mathbb{O}\langle\mathsf{s}\rangle$ is a product of the values $\mathsf{s}(\Delta)$
at distinct cells $\Delta$ times an operator independent of $\mathsf{s}$, and the sum is finite;
hence $\mathbb{O}\langle\mathsf{s}\rangle$ is multi-affine. At $\mathsf{s}=\mathbf{1}_{\mathsf{Y}}$
only the families $\mathsf{Y}'\subset\mathsf{Y}$ survive, each with coefficient $1$, and by the
definition of the M\"obius pieces in Notation~\ref{9.5:not-moebius-pieces},
\begin{align*}
\mathbb{O}\langle\mathbf{1}_{\mathsf{Y}}\rangle
&=\sum_{\mathsf{Y}'\subset\mathsf{Y}}\ \sum_{\mathsf{Y}''\subset\mathsf{Y}'}
(-1)^{\#(\mathsf{Y}'\setminus\mathsf{Y}'')}\,\mathbb{O}_{K\cup Y''}\\
&=\sum_{\mathsf{Y}''\subset\mathsf{Y}}\mathbb{O}_{K\cup Y''}
\sum_{\mathsf{Y}''\subset\mathsf{Y}'\subset\mathsf{Y}}(-1)^{\#(\mathsf{Y}'\setminus\mathsf{Y}'')}.
\end{align*}
The inner sum equals $(1-1)^{\#(\mathsf{Y}\setminus\mathsf{Y}'')}$, which is $1$ for
$\mathsf{Y}''=\mathsf{Y}$ and $0$ otherwise. This is M\"obius inversion on the finite Boolean
lattice of subfamilies of $\mathsf{Y}$, and it gives
$\mathbb{O}\langle\mathbf{1}_{\mathsf{Y}}\rangle=\mathbb{O}_{K\cup Y}$. Taking
$\mathsf{Y}=\sigma_0$, the union $K\cup\bigcup\sigma_0$ is the union of all cells of
$\pi^{\mathrm{mix}}$, which is $\mathbb{T}$. Hence
$\mathbb{O}\langle\mathbf{1}\rangle=\mathbb{O}_{\mathbb{T}}$, the operator of
$\boldsymbol{\Omega}(\mathbb{T})$. By statement (s3) of Lemma~\ref{9.5:lem-shell-lemma},
$\boldsymbol{\Omega}(\mathbb{T})=\boldsymbol{\Omega}$, so
$\mathbb{O}\langle\mathbf{1}\rangle=\mathbb{O}$.

(n2) If $\mathsf{s}=0$ off $\mathsf{Y}_0$, then the product $\prod_{\Delta\in\mathsf{Y}}\mathsf{s}(\Delta)$
vanishes unless $\mathsf{Y}\subset\mathsf{Y}_0$, so only the pieces $\mathfrak{o}_{\mathsf{K}}(\mathsf{Y})$
with $\mathsf{Y}\subset\mathsf{Y}_0$ survive, and the surviving coefficients read $\mathsf{s}$ on
$\mathsf{Y}_0$ only. Each surviving piece is a signed sum of the box operators
$\mathbb{O}_{K\cup Y'}$ with $Y'\subset Y_0$. The kernel of such a box operator vanishes off
$(K\cup Y')^2\subset(K\cup Y_0)^2$ by Notation~\ref{9.5:not-moebius-pieces}. Block-diagonality over the
wall-components of $\mathsf{K}\cup\mathsf{Y}_0$, and the fact that only $(\mathbf{U},\mathbf{J})$ on
these components is read, follow for each surviving piece from
Theorem~\ref{9.5:thm-moebius-decoupling}(c) together with the second property in
\eqref{9.5:eq-walk-influence-c}. A finite sum of operators with these properties has them.

(n3) Let $\mathsf{s}$ lie in the polydisc. Then
$\bigl|\prod_{\Delta\in\mathsf{Y}}\mathsf{s}(\Delta)\bigr|\le e^{\kappa_1\#\mathsf{Y}}$. The first bound of
Theorem~\ref{9.5:thm-moebius-decoupling}(b), displayed in \eqref{9.5:eq-moebius-decoupling-a}, gives
\begin{equation}\label{9.5:eq-interpolated-operators-2}
\bigl\|\Delta(y)\,\mathbb{O}\langle\mathsf{s}\rangle\,\Delta(y')\bigr\|
\le B_0\sum_{\mathsf{Y}\subset\sigma_0}\bigl(2e^{\kappa_1}\bigr)^{\#\mathsf{Y}}
e^{-\frac12\delta_0\,\mathsf{t}_{\times}(\mathsf{Y}\cup\{\Delta_y,\Delta_{y'}\})}.
\end{equation}
Consider the map $\mathsf{Y}\mapsto\mathsf{F}:=\mathsf{Y}\cup\{\Delta_y,\Delta_{y'}\}$. For a given
$\mathsf{F}$, every preimage $\mathsf{Y}$ satisfies
$\mathsf{F}\setminus\{\Delta_y,\Delta_{y'}\}\subset\mathsf{Y}\subset\mathsf{F}$. Hence $\mathsf{Y}$ is
fixed by whether it contains $\Delta_y$ and whether it contains $\Delta_{y'}$, and the map is at
most $4$-to-$1$. Moreover $\#\mathsf{Y}\le\#\mathsf{F}$, and every image contains $\Delta_y$. Put
$z:=2e^{\kappa_1}\ge1$ and $c:=\frac12\delta_0$. Then the right side of
\eqref{9.5:eq-interpolated-operators-2} is at most
\begin{equation*}
4B_0\sum_{\mathsf{F}\ni\Delta_y}z^{\#\mathsf{F}}\,e^{-c\,\mathsf{t}_{\times}(\mathsf{F})},
\end{equation*}
where the sum runs over finite families of cells. For this $z$,
\begin{equation*}
X_z=(2D^2z)^{2^d}=(4D^2)^{2^d}e^{2^d\kappa_1}=X_{\kappa},
\qquad
\log(2X_z)=2^d(\kappa_1+\log4D^2)+\log2 .
\end{equation*}
Hence the floor \eqref{9.5:eq-entropy-sum-a} of Lemma~\ref{9.5:lem-entropy-sum} at $(c,z)$ is
\eqref{9.5:eq-interpolated-operators-a} at $c$, and the latter is assumed at $c=\tfrac12\delta_0$.
Lemma~\ref{9.5:lem-entropy-sum}(a) with $\rho=\rho_{\times}$ and anchor $\Delta_0=\Delta_y$ bounds
the sum by $2X_{\kappa}$. Therefore
\begin{equation*}
\bigl\|\Delta(y)\,\mathbb{O}\langle\mathsf{s}\rangle\,\Delta(y')\bigr\|
\le 8X_{\kappa}B_0=8(4D^2)^{2^d}e^{2^d\kappa_1}B_0=C_N^{\mathrm{mix}}e^{2^d\kappa_1}B_0 .
\end{equation*}

(n3$'$) Proceed as in (n3) up to the $4$-to-$1$ reduction. Every family
$\mathsf{F}=\mathsf{Y}\cup\{\Delta_y,\Delta_{y'}\}$ contains both $\Delta_y$ and $\Delta_{y'}$. Now
take $c:=\tfrac14\delta_0$, so that $\tfrac12\delta_0=2c$ in \eqref{9.5:eq-interpolated-operators-2},
and keep $z=2e^{\kappa_1}$. Under \eqref{9.5:eq-interpolated-operators-a} at $c=\frac14\delta_0$,
which is the floor \eqref{9.5:eq-entropy-sum-a} at $(c,z)$ by the computation above,
Lemma~\ref{9.5:lem-entropy-sum}(c) with $\Delta_0=\Delta_y$ and $\Delta_1=\Delta_{y'}$ gives
\begin{equation*}
4B_0\sum_{\mathsf{F}\ni\Delta_y,\Delta_{y'}}z^{\#\mathsf{F}}e^{-2c\,\mathsf{t}_{\times}(\mathsf{F})}
\le 8X_{\kappa}B_0\,e^{\frac12\delta_0\rho_{\times}}\,e^{-\frac14\delta_0\mathsf{d}(\Delta_y,\Delta_{y'})},
\end{equation*}
and $8X_{\kappa}B_0=C_N^{\mathrm{mix}}e^{2^d\kappa_1}B_0$.

(n4) The difference $\mathbb{O}\langle\mathsf{s}\rangle-\mathbb{O}\langle0\rangle$ is the sum over the
non-empty families $\mathsf{Y}\subset\sigma_0$. Since $\Delta_y$ touches no cell of $\sigma_0$, each
such $\mathsf{Y}$ contains a cell $\Delta$ that does not touch $\Delta_y$. The family
$\mathsf{F}=\mathsf{Y}\cup\{\Delta_y,\Delta_{y'}\}$ contains both $\Delta$ and $\Delta_y$, so
Lemma~\ref{9.5:lem-common-point}(b) gives $\mathsf{d}(\Delta,\Delta_y)\ge M/L$ and
$\mathsf{t}_{\times}(\mathsf{F})\ge\mathsf{d}(\Delta,\Delta_y)-2\rho_{\times}$. With
\eqref{9.5:eq-interpolated-operators-1} this yields
\begin{equation*}
\mathsf{t}_{\times}(\mathsf{F})\ \ge\ \frac{M}{L}-2\rho_{\times}\ \ge\ \frac{M}{L}-\frac{M}{16L}
\ \ge\ \frac{3M}{4L}=3r_0 .
\end{equation*}
The reduction of (n3), with $z=2e^{\kappa_1}$ and $c=\frac12\delta_0$, now runs over families with
$\mathsf{t}_{\times}(\mathsf{F})\ge3r_0$ only. Since $\tau:=3r_0\ge r_0/2$,
Lemma~\ref{9.5:lem-entropy-sum}(b) applies under the floor assumed at $c=\frac12\delta_0$, and
\begin{align*}
\bigl\|\Delta(y)\bigl(\mathbb{O}\langle\mathsf{s}\rangle-\mathbb{O}\langle0\rangle\bigr)\Delta(y')\bigr\|
&\le 4B_0\sum_{\substack{\mathsf{F}\ni\Delta_y\\ \mathsf{t}_{\times}(\mathsf{F})\ge3r_0}}
z^{\#\mathsf{F}}e^{-c\,\mathsf{t}_{\times}(\mathsf{F})}
\le 4B_0\,e^{cr_0/4}\,e^{-3cr_0/2}\\
&=4B_0\,e^{-5cr_0/4}=4B_0\,e^{-5\delta_0M/32L}\le 4B_0\,e^{-\delta_0M/8L}.
\end{align*}
\end{proof}

\begin{remark}
\leavevmode
\begin{itemize}
\item[(1)] Let the core contain a neighbourhood of a block of the block's own physical size.
Suppose it crosses a sequence of thin zones, and its far face touches a large number
$\mathcal{N}$ of fine cells of $\sigma_0$. In (n3) these cells enter only through families
$\mathsf{F}\ni\Delta_y$ for which $\mathsf{t}_{\times}(\mathsf{F})$ is at least $M$ times the
number of zones crossed, minus $2\rho_{\times}$.
This follows from the comparison of scaled distance with the number of zones crossed, established
earlier in this chapter, and from Lemma~\ref{9.5:lem-common-point}(b). Furthermore, the count of
families of cells within scaled reach of a cell, also established earlier in this chapter,
bounds the number of all families within scaled reach $\tau$ of $\Delta_y$, whatever zones lie
between, by $(2D^2)^{2^d\cdot8L\tau/M}$. Large physical neighbourhoods crossing thin zones are thus
paid for by scaled distance. Bounds obtained after Cauchy's estimate in the parameters
$\mathsf{s}$ on the polydisc of (n3), which weigh each of the $\mathcal{N}$ cells by
$e^{-(\kappa_1-1)}$ with no distance factor, are not part of (n3).
\item[(2)] Consider the equal-size analogue of the present theorem in its strengthened form,
part (0), in which Theorem~\ref{9.5:thm-moebius-decoupling}(b) and (n3) are proved again with the
walk-term bounds of \eqref{9.5:eq-walk-influence-c} replaced by their converted form: the
conversion, in parts (a)--(c) of the corresponding statement for equal-size cells, of the
walk-term bounds into bounds organised by bands of distance to the complement of the box.
That conversion spends $\tfrac14\delta_0$ of the decay.
For mixed sizes this strengthened form reads as follows:
Theorem~\ref{9.5:thm-moebius-decoupling}(b) holds at the rate $\frac14\delta_0$ in
$\mathsf{t}_{\times}$; the tree bound of Lemma~\ref{9.5:lem-common-point}(c) then yields the factor
$e^{-\frac18\delta_0Md^{\circ}_{\mathrm{mix}}}$; and Lemma~\ref{9.5:lem-entropy-sum} is applied under
\eqref{9.5:eq-interpolated-operators-a} at $c=\frac14\delta_0$. In the sup norms, where a further
$\frac18\delta_0$ is spent, it is applied at $c=\frac18\delta_0$. That conversion reads a box
$Z$ as a set, through bands of $\operatorname{dist}_{\infty}(\cdot,Z^{\complement})$, so no cell
size occurs in it. Hence the strengthened statement, part (0), and everything derived from it,
hold for mixed sizes. The constant is $C_N^{\mathrm{mix}}$ in place of the constant of the
equal-size case, and the floor is \eqref{9.5:eq-interpolated-operators-a} at $c=\delta_0/8$.
\item[(3)] The entropy sums in the proof are anchored at the single cell $\Delta_y$ and do not
involve the core, so the bounds (n3)--(n4) hold verbatim when the core is allowed to depend on
the pair of bonds at which the kernel is evaluated.
\end{itemize}
\end{remark}

\begin{lemma}[One partition and one metric for two cutoffs]\label{9.5:lem-two-cutoff-metric}
Let $\mathsf{A}$ and $\mathsf{B}$ be two cutoffs on the torus $\mathbb{T}$, with lattice units $\eta$
and $\eta/L$ respectively. A superscript $\mathsf{A}$ or $\mathsf{B}$ on $\boldsymbol{\Omega}$,
$\Omega_j$, $k$, $m$, $\ell$ indicates the cutoff, $k$ being the top index of the sequence
$(\Omega_j)_{0\le j\le k}$; $\ell_{\hat{\boldsymbol{\Omega}}}$ denotes the scale function of a
sequence $\hat{\boldsymbol{\Omega}}$ as in Definition~\ref{9.5:def-refinements}. Let
$\boldsymbol{\Omega}^{\mathsf{B}}$ be paired with $\boldsymbol{\Omega}^{\mathsf{A}}$, that is,
\begin{equation}\label{9.5:eq-two-cutoff-metric-1}
\Omega^{\mathsf{B}}_{j+1}=\Omega^{\mathsf{A}}_{j}\qquad\text{for } 0\le j\le k^{\mathsf{A}} .
\end{equation}
Then:
\begin{enumerate}
\item[(a)] the cells of $\mathsf{B}$ of level $j+1$ are the cells of $\mathsf{A}$ of level $j$;
\item[(b)] $\ell^{\mathsf{B}}=\ell^{\mathsf{A}}$ on $\Omega^{\mathsf{A}}_0$, and
$\ell^{\mathsf{B}}\le\ell^{\mathsf{A}}$ off $\Omega^{\mathsf{A}}_0$; for the metric one puts
$\Omega^{\mathsf{A}}_0:=\mathbb{T}$, and then $\ell^{\mathsf{B}}=\ell^{\mathsf{A}}$ on all of
$\mathbb{T}$;
\item[(c)] for every box $Z$, the sequence $\boldsymbol{\Omega}^{\mathsf{B}}(Z)$ of $\mathsf{B}$
induced on $Z$ (Lemma~\ref{9.5:lem-shell-lemma}) is a refinement of
$\boldsymbol{\Omega}^{\mathsf{A}}$ in the sense of Definition~\ref{9.5:def-refinements}.
\end{enumerate}
Consequently $\pi^{\mathrm{mix}}$, $\mathsf{L}$, $\mathsf{d}$, $\mathcal{S}_{\rho}$,
$\mathsf{t}_{\rho}$, $d^{\circ}_{\mathrm{mix}}$ and $S^{\sharp}$ are common objects of the two
cutoffs, and Proposition~\ref{9.5:prop-mixed-counting}, Theorem~\ref{9.5:thm-moebius-decoupling}
and Theorem~\ref{9.5:thm-interpolated-operators} hold for both cutoffs with the same
$\mathsf{t}_{\times}$ and the same constants.
\end{lemma}

\begin{proof}
(a) The tiling $\pi_j$ of a cutoff of unit $\eta$ consists of closed cubes of side
$M_j=ML^j\eta$. For $\mathsf{B}$ the side at level $j+1$ is
\begin{equation*}
ML^{j+1}\cdot\frac{\eta}{L}=ML^{j}\eta=M_j ,
\end{equation*}
the side of the cubes of $\pi^{\mathsf{A}}_j$. Thus $\pi^{\mathsf{B}}_{j+1}$ and
$\pi^{\mathsf{A}}_j$ are the same tiling of $\mathbb{T}$ by closed cubes of side $M_j$. By
\eqref{9.5:eq-two-cutoff-metric-1}, $m^{\mathsf{B}}=m^{\mathsf{A}}+1$ on
$\Omega^{\mathsf{A}}_0$ (see (b)), so a cube of this tiling is in proper scale $j+1$ for
$\mathsf{B}$ exactly when it is in proper scale $j$ for $\mathsf{A}$; that is, the cells of
$\mathsf{B}$ of level $j+1$ are the cells of $\mathsf{A}$ of level $j$.

(b) Let $x\in\Omega^{\mathsf{A}}_0=\Omega^{\mathsf{B}}_1$. By
\eqref{9.5:eq-two-cutoff-metric-1}, $x\in\Omega^{\mathsf{B}}_{n+1}$ if and only if
$x\in\Omega^{\mathsf{A}}_n$, so $m^{\mathsf{B}}(x)=m^{\mathsf{A}}(x)+1$ and
\begin{equation*}
\ell^{\mathsf{B}}(x)=L^{m^{\mathsf{A}}(x)+1}\,\frac{\eta}{L}=L^{m^{\mathsf{A}}(x)}\eta
=\ell^{\mathsf{A}}(x).
\end{equation*}
If $x\notin\Omega^{\mathsf{A}}_0=\Omega^{\mathsf{B}}_1$, then $m^{\mathsf{B}}(x)=0$ by the
convention of Definition~\ref{9.5:def-refinements}, and since $m^{\mathsf{A}}\ge0$,
\begin{equation*}
\ell^{\mathsf{B}}(x)=\frac{\eta}{L}<\eta\le L^{m^{\mathsf{A}}(x)}\eta=\ell^{\mathsf{A}}(x).
\end{equation*}
With $\Omega^{\mathsf{A}}_0:=\mathbb{T}$, as for the metric, the first case covers all of $\mathbb{T}$.

(c) By (s4) of the shell lemma for mixed-size boxes (Lemma~\ref{9.5:lem-shell-lemma}, see
\eqref{9.5:eq-shell-lemma-a}), $\Omega^{\mathsf{B}}_{n+1}(Z)=\Omega^{\mathsf{A}}_n(Z)$ for every
$n$. The computation of (b), applied to this pair, gives
$\ell_{\boldsymbol{\Omega}^{\mathsf{B}}(Z)}=\ell_{\boldsymbol{\Omega}^{\mathsf{A}}(Z)}$ on
$\Omega^{\mathsf{A}}_0(Z)$ and $\ell_{\boldsymbol{\Omega}^{\mathsf{B}}(Z)}\le\eta/L$ elsewhere. By
(s1) of the same lemma, $\Omega^{\mathsf{A}}_n(Z)\subset\Omega^{\mathsf{A}}_n$, so
$\ell_{\boldsymbol{\Omega}^{\mathsf{A}}(Z)}\le\ell^{\mathsf{A}}$ pointwise; hence
$\ell_{\boldsymbol{\Omega}^{\mathsf{B}}(Z)}\le\ell^{\mathsf{A}}$ pointwise. By (s1) applied to
$\mathsf{B}$, $\boldsymbol{\Omega}^{\mathsf{B}}(Z)$ is a decreasing sequence of unions of closed
cubes. It is therefore a refinement of $\boldsymbol{\Omega}^{\mathsf{A}}$ with unit $\eta/L$.

The definitions of $\pi^{\mathrm{mix}}$, $\mathsf{L}$, $\mathsf{d}$, $\mathcal{S}_{\rho}$,
$\mathsf{t}_{\rho}$, $d^{\circ}_{\mathrm{mix}}$, $S^{\sharp}$ and the proofs of
Proposition~\ref{9.5:prop-mixed-counting}, Theorem~\ref{9.5:thm-moebius-decoupling} and
Theorem~\ref{9.5:thm-interpolated-operators} do not read the unit $\eta$, the top level $k$ of
the sequence, or the value of any level. They read only the cells and the ratios of the scales at
neighbouring points. By (a) and (b) these data agree for $\mathsf{A}$ and $\mathsf{B}$, and by (c)
the sequences of $\mathsf{B}$ are refinements of $\boldsymbol{\Omega}^{\mathsf{A}}$. The listed
objects, the quantity $\mathsf{t}_{\times}$ and the constants are therefore the same for the two
cutoffs.
\end{proof}

\begin{corollary}[Two-cutoff difference of M\"obius pieces]\label{9.5:cor-two-cutoff-difference}
Let $\mathsf{A}$ and $\mathsf{B}$ be the two cutoffs paired as in
Lemma~\ref{9.5:lem-two-cutoff-metric}, with lattice units $\eta$ and $\eta/L$. A superscript
$\mathsf{A}$ or $\mathsf{B}$ marks an object of the corresponding cutoff, and $\delta$ in front of an
object denotes the difference between its $\mathsf{B}$-version and its $\mathsf{A}$-version. Let
$\delta>0$ be the decay rate of the weighted norm $\|\cdot\|_{\delta/4}$ in~\eqref{9.5:eq-two-cutoff-difference-3}
below. The prefix $\delta$ (difference of the two cutoffs) always stands in front of an operator
symbol, and the rate $\delta$ occurs only in exponents, so the two uses do not conflict. Assume:
\begin{enumerate}
\item[(1)] the bound (b) of the M\"obius decoupling theorem for mixed-size cells
(Theorem~\ref{9.5:thm-moebius-decoupling}, display~\eqref{9.5:eq-moebius-decoupling-a}) holds in
each of the two cutoffs $\mathsf{A}$ and $\mathsf{B}$;
\item[(2)] the two-cutoff comparison bound for natural functionals holds for the functionals
$\mathfrak{q}=e^{\delta|b-b'|/2}\mathbb{O}_Z(b,b')$, $\mathbb{O}$ in the list $\mathfrak{O}$, with
$\theta_q=\theta/4$:
\begin{equation*}
\bigl|\mathfrak{q}[\mathfrak{B}_{Z'};M\mathbf{U}']-\mathfrak{q}[\mathfrak{B}_Z;M\mathbf{U}]\bigr|
\le C_V\mathsf{N}\,L^{-\min\{\theta_q,\theta/2\}\,j_{\min}(Z)} ;
\end{equation*}
here $\mathfrak{q}[\mathfrak{B}_Z;M\mathbf{U}]$ denotes the functional $\mathfrak{q}$ evaluated on
the determining set $\mathfrak{B}_Z$ and the field argument $M\mathbf{U}$ of the box $Z$ in the
cutoff $\mathsf{A}$, and $\mathfrak{q}[\mathfrak{B}_{Z'};M\mathbf{U}']$ the same functional
evaluated on the corresponding determining set and field argument in the cutoff $\mathsf{B}$, so
that the left-hand side is $e^{\delta|b-b'|/2}|\delta\mathbb{O}_Z(b,b')|$;
\item[(3)] the floor~\eqref{9.5:eq-walk-influence-a};
\item[(4)] the floor~\eqref{9.5:eq-interpolated-operators-a} at $c=\tfrac14\delta_0$.
\end{enumerate}
Recall the M\"obius piece of Theorem~\ref{9.5:thm-moebius-decoupling}, formed in each of the two
cutoffs: for $\mathbb{O}$ in $\mathfrak{O}$, a core $\mathsf{K}$ and a finite family $\mathsf{W}$,
with the sum over subfamilies,
\begin{equation*}
\mathfrak{o}_{\mathsf{K}}(\mathsf{W})
=\sum_{\mathsf{W}'\subset\mathsf{W}}(-1)^{\#(\mathsf{W}\setminus\mathsf{W}')}\,\mathbb{O}_{K\cup W'} .
\end{equation*}
Then for every operator $\mathbb{O}$ in the list of the family bound for the interpolated operators
on cells of one size, every bond pair $(b,b')$,
the core $\mathsf{K}=\mathsf{K}(b,b')$, which contains $\Delta_b$ and $\Delta_{b'}$, and every finite
family $\mathsf{W}\subset\pi^{\mathrm{mix}}\setminus\mathsf{K}$,
\begin{equation}\label{9.5:eq-two-cutoff-difference-1}
\begin{split}
|\delta\mathfrak{o}_{\mathsf{K}}(\mathsf{W})(b,b')|
\le{}&(2C_V\mathsf{N}B_0)^{1/2}\,2^{\#\mathsf{W}}\,\vartheta_w^{\,j_{\min}(K\cup W)}\\
&\times e^{-\delta|b-b'|/4}\,
e^{-\frac14\delta_0\mathsf{t}_{\times}(\mathsf{W}\cup\{\Delta_b,\Delta_{b'}\})},
\end{split}
\end{equation}
where $\vartheta_w=L^{-\theta/8}$; moreover
\begin{equation*}
e^{-\frac14\delta_0\mathsf{t}_{\times}(\mathsf{F})}
\le (C_0^{\mathrm{mix}})^{1/2}\,e^{-\frac18\delta_0 M d^{\circ}_{\mathrm{mix}}(\mathsf{F})} ,
\qquad \mathsf{F}=\mathsf{W}\cup\{\Delta_b,\Delta_{b'}\}.
\end{equation*}
Put, for decoupling parameters $\mathsf{s}$ and every $n$, with the sum (here and
in~\eqref{9.5:eq-two-cutoff-difference-2}) over finite families
$\mathsf{W}\subset\pi^{\mathrm{mix}}\setminus\mathsf{K}(b,b')$,
\begin{equation*}
\delta\mathbb{O}^{(n)}\langle\mathsf{s}\rangle(b,b')
:=\sum_{\mathsf{W}:\,j_{\min}(K\cup W)=n}\ \prod_{\Delta\in\mathsf{W}}\mathsf{s}(\Delta)\,
\delta\mathfrak{o}_{\mathsf{K}(b,b')}(\mathsf{W})(b,b') .
\end{equation*}
Then on the polydisc $|\mathsf{s}(\Delta)|\le e^{\kappa_1}$, for every $n$ and every $(b,b')$,
\begin{equation}\label{9.5:eq-two-cutoff-difference-2}
\begin{split}
e^{\delta|b-b'|/4}\,|\delta\mathbb{O}^{(n)}\langle\mathsf{s}\rangle(b,b')|
&\le\sum_{\mathsf{W}:\,j_{\min}(K\cup W)=n}e^{\kappa_1\#\mathsf{W}}\,e^{\delta|b-b'|/4}\,
|\delta\mathfrak{o}_{\mathsf{K}(b,b')}(\mathsf{W})(b,b')|\\
&\le C_F^{\mathrm{mix}}\,e^{2^d\kappa_1}\,\vartheta_w^{\,n},
\end{split}
\end{equation}
that is, writing $\|\cdot\|_{\delta/4}$ for the weighted norm of the two-cutoff family bound named
below, namely the supremum over bond pairs of
$e^{\delta|b-b'|/4}$ times the modulus of the kernel,
\begin{equation}\label{9.5:eq-two-cutoff-difference-3}
\|\delta\mathbb{O}^{(n)}\langle\mathsf{s}\rangle\|_{\delta/4}\le C_F^{\mathrm{mix}}\,e^{2^d\kappa_1}\,
\vartheta_w^{\,n},
\qquad
C_F^{\mathrm{mix}}:=2(4D^2)^{2^d}(2C_V\mathsf{N}B_0)^{1/2}.
\end{equation}
Consequently part (iii) of the two-cutoff family bound for the interpolated operators used in
the uniqueness argument of this chapter holds for families of cells of several sizes: its display
is~\eqref{9.5:eq-two-cutoff-difference-1}, with its length factor $e^{-\frac18\delta_0Md^{\circ}}$,
where $d^{\circ}$ is the single-size loose tree length used there, read as
$e^{-\frac14\delta_0\mathsf{t}_{\times}}$ and hence bounded by
$(C_0^{\mathrm{mix}})^{1/2}e^{-\frac18\delta_0Md^{\circ}_{\mathrm{mix}}}$; the constant $C_F$ of
that bound, which enters its overall constant, is replaced by $C_F^{\mathrm{mix}}$
of~\eqref{9.5:eq-two-cutoff-difference-3};
the factor $e^{16\kappa_1}$ is unchanged ($2^d=16$); and the floors on $M$ of the family bound for
the interpolated operators on cells of one size, which it invokes, are replaced by
\eqref{9.5:eq-walk-influence-a} together with~\eqref{9.5:eq-interpolated-operators-a} at
$c=\tfrac14\delta_0$.
\end{corollary}

\begin{proof}
By Lemma~\ref{9.5:lem-two-cutoff-metric} the cells $\pi^{\mathrm{mix}}$, the levels $j_{\min}$, the
length $\mathsf{t}_{\times}$ and the constants of Theorem~\ref{9.5:thm-moebius-decoupling} are the same
in the two cutoffs, so $\delta\mathfrak{o}_{\mathsf{K}}(\mathsf{W})$ is a difference of objects
indexed by the same families.

\emph{First bound.} Since $\mathfrak{q}=e^{\delta|b-b'|/2}\mathbb{O}_Z(b,b')$ and
$\min\{\theta/4,\theta/2\}=\theta/4$, hypothesis~(2) gives, for every box $Z$,
\begin{equation*}
|\delta\mathbb{O}_Z(b,b')|\le C_V\mathsf{N}\,e^{-\delta|b-b'|/2}\,L^{-\theta j_{\min}(Z)/4}.
\end{equation*}
The M\"obius piece is the alternating sum over the $2^{\#\mathsf{W}}$ subfamilies
$\mathsf{W}'\subset\mathsf{W}$, and the difference is taken termwise, so
\begin{equation*}
|\delta\mathfrak{o}_{\mathsf{K}}(\mathsf{W})(b,b')|
\le\sum_{\mathsf{W}'\subset\mathsf{W}}|\delta\mathbb{O}_{K\cup W'}(b,b')|
\le 2^{\#\mathsf{W}}C_V\mathsf{N}\,e^{-\delta|b-b'|/2}\,L^{-\theta j_{\min}(K\cup W)/4},
\end{equation*}
because $K\cup W'\subset K\cup W$ gives $j_{\min}(K\cup W')\ge j_{\min}(K\cup W)$, while $L>1$ and
$\theta>0$.

\emph{Second bound.} By the triangle inequality and hypothesis~(1) in each cutoff, with the common
$\mathsf{t}_{\times}$ and constants of Lemma~\ref{9.5:lem-two-cutoff-metric},
\begin{equation*}
|\delta\mathfrak{o}_{\mathsf{K}}(\mathsf{W})(b,b')|
\le|\mathfrak{o}^{\mathsf{A}}_{\mathsf{K}}(\mathsf{W})(b,b')|
+|\mathfrak{o}^{\mathsf{B}}_{\mathsf{K}}(\mathsf{W})(b,b')|
\le 2\cdot2^{\#\mathsf{W}}B_0\,e^{-\frac12\delta_0\mathsf{t}_{\times}(\mathsf{F})},
\end{equation*}
with $\mathsf{F}=\mathsf{W}\cup\{\Delta_b,\Delta_{b'}\}$.

\emph{Combination.} For $x,y\ge0$ one has $\min\{x,y\}\le(xy)^{1/2}$. The product of the two
right-hand sides is
\begin{equation*}
2\,C_V\mathsf{N}B_0\,2^{2\#\mathsf{W}}\,e^{-\delta|b-b'|/2}\,L^{-\theta j_{\min}(K\cup W)/4}\,
e^{-\frac12\delta_0\mathsf{t}_{\times}(\mathsf{F})},
\end{equation*}
and its square root is the right-hand side of~\eqref{9.5:eq-two-cutoff-difference-1}, since
$L^{-\theta j/8}=\vartheta_w^{\,j}$. This proves~\eqref{9.5:eq-two-cutoff-difference-1}. The second
step of~\eqref{9.5:eq-moebius-decoupling-a} compares $e^{-\frac12\delta_0\mathsf{t}_{\times}}$ with
$C_0^{\mathrm{mix}}e^{-\frac14\delta_0Md^{\circ}_{\mathrm{mix}}}$; its square root is the displayed
comparison following~\eqref{9.5:eq-two-cutoff-difference-1}.

\emph{The sum.} On the polydisc, $|\prod_{\Delta\in\mathsf{W}}\mathsf{s}(\Delta)|\le
e^{\kappa_1\#\mathsf{W}}$, which gives the first inequality
of~\eqref{9.5:eq-two-cutoff-difference-2} by the triangle inequality. Inserting
\eqref{9.5:eq-two-cutoff-difference-1}, the factors $e^{\pm\delta|b-b'|/4}$ cancel, and on
$j_{\min}(K\cup W)=n$ the factor $\vartheta_w^{\,j_{\min}(K\cup W)}$ equals $\vartheta_w^{\,n}$.
Thus the sum is at most
\begin{equation*}
(2C_V\mathsf{N}B_0)^{1/2}\,\vartheta_w^{\,n}
\sum_{\mathsf{W}:\,j_{\min}(K\cup W)=n}(2e^{\kappa_1})^{\#\mathsf{W}}\,
e^{-\frac14\delta_0\mathsf{t}_{\times}(\mathsf{W}\cup\{\Delta_b,\Delta_{b'}\})}.
\end{equation*}
The map $\mathsf{W}\mapsto\mathsf{F}:=\mathsf{W}\cup\{\Delta_b,\Delta_{b'}\}$ is injective, because
$\mathsf{W}\cap\mathsf{K}=\emptyset$ and $\Delta_b,\Delta_{b'}\in\mathsf{K}$. Its image consists of
families containing $\Delta_b$. Moreover $\#\mathsf{F}\ge\#\mathsf{W}$ and $z:=2e^{\kappa_1}\ge1$, so
$z^{\#\mathsf{W}}\le z^{\#\mathsf{F}}$. The sum over $\{j_{\min}=n\}$ is a sub-sum of a sum of
non-negative terms; hence it is at most
\begin{equation*}
\sum_{\mathsf{F}\ni\Delta_b}z^{\#\mathsf{F}}\,e^{-c\,\mathsf{t}_{\rho_{\times}}(\mathsf{F})},
\qquad z=2e^{\kappa_1},\ \ c=\tfrac14\delta_0,
\end{equation*}
recalling $\mathsf{t}_{\times}=\mathsf{t}_{\rho_{\times}}$.

We bound this by the entropy sum over mixed-size families
(Lemma~\ref{9.5:lem-entropy-sum}(a)), applied with $\rho=\rho_{\times}$. For $z=2e^{\kappa_1}$ one
has $X_z=(4D^2e^{\kappa_1})^{2^d}=(4D^2)^{2^d}e^{2^d\kappa_1}$, so that
$\log(2X_z)=2^d(\kappa_1+\log4D^2)+\log2$. The floor of that lemma at $c=\tfrac14\delta_0$ is
therefore~\eqref{9.5:eq-interpolated-operators-a} at $c=\tfrac14\delta_0$, which is
hypothesis~(4). The lemma bounds the
sum by
\begin{equation*}
2X_z=2(4D^2)^{2^d}e^{2^d\kappa_1}.
\end{equation*}
Multiplying by $(2C_V\mathsf{N}B_0)^{1/2}\vartheta_w^{\,n}$ gives
$C_F^{\mathrm{mix}}e^{2^d\kappa_1}\vartheta_w^{\,n}$. This is the second inequality
of~\eqref{9.5:eq-two-cutoff-difference-2}. Since it holds for every $(b,b')$,
\eqref{9.5:eq-two-cutoff-difference-3} follows.

\emph{Consequence.} In part (iii) of the two-cutoff family bound for the interpolated operators
(the last assertion of the statement), the bound on the M\"obius pieces is
\eqref{9.5:eq-two-cutoff-difference-1}, with the factor
$e^{-\frac18\delta_0Md^{\circ}}$ read as $e^{-\frac14\delta_0\mathsf{t}_{\times}}$, together with
the displayed comparison with $d^{\circ}_{\mathrm{mix}}$. For the level-$n$ sums on the polydisc that
part uses the bound~\eqref{9.5:eq-two-cutoff-difference-3}, where $e^{2^d\kappa_1}=e^{16\kappa_1}$. For
the bilinear kernels of that bound (written $\Gamma$ there, not to be confused with the graphs
$\Gamma$ of this section), built from two factors $A$ and $B$, with level-$n$
differences $\delta A^{(n)}$, $\delta B^{(n)}$ formed as
$\delta\mathbb{O}^{(n)}\langle\mathsf{s}\rangle$ above and with $A^{\mathsf{A}}$, $B^{\mathsf{B}}$
their versions in the cutoffs $\mathsf{A}$ and $\mathsf{B}$, the level-$n$ difference is
\begin{equation*}
\delta\Gamma^{(n)}:=\delta A^{(n)}B^{\mathsf{B}}+A^{\mathsf{A}}\delta B^{(n)} .
\end{equation*}
The differenced factor is bounded by~\eqref{9.5:eq-two-cutoff-difference-3}. The undifferenced factor
is bounded by (n3) of the interpolated box operators for mixed-size cells
(Theorem~\ref{9.5:thm-interpolated-operators}, display~\eqref{9.5:eq-interpolated-operators-b}).
\end{proof}

\begin{corollary}[The two-cutoff bound on joint bond sets]\label{9.5:cor-joint-bond-sets}
Work at stage $n+1$ and level $j$, at the site of the proposition on the joint bond sets of
$Z_0$ and $Z$, and let $\sigma=(\sigma(\Delta))_{\Delta}$ denote the decoupling parameters at
that site. Let $\boldsymbol{\Omega}$ be the determining sequence $\mathbb{B}^{(n)}_k$ of the
stage, as in \cite[(1.14)--(1.17)]{B-CMP122a}. Let $\pi^{\mathrm{mix}}$ be its cells in proper
scales, in the sense of Definition~\ref{9.5:def-scaled-length}. Let $I$ and $E$ be the joint
$(A_j,B_j)$ bond sets of $Z_0$ and of $Z$. The three conditions on grids, cells and the scaled
length displayed in \eqref{9.5:eq-scaled-length-a} hold for $\boldsymbol{\Omega}$ and
$\pi^{\mathrm{mix}}$, by \cite{B-CMP122a} and the admissibility conditions of the stage.
Assume in addition:
\begin{itemize}
\item[(a)] the walk-expansion bounds \eqref{9.5:eq-walk-influence-c} hold for the operators
considered below, with cells in $\pi^{\mathrm{mix}}$;
\item[(b)] the two-cutoff comparison of natural functionals at the least level $j_{\min}$
holds, in the form in which it is an input of Corollary~\ref{9.5:cor-two-cutoff-difference}.
\end{itemize}
Then the conclusion of Corollary~\ref{9.5:cor-two-cutoff-difference} holds for each
\[
\mathbb{O}\in\bigl\{\,T,\ (T_{II})^{-1}=\mathcal{C}(Z_0),\ \mathcal{C}^{1/2}\,\bigr\},
\]
and it holds for
\[
\Gamma:=(T\langle\sigma\rangle)_{II^{c}}\,\mathcal{C}^{1/2}\langle\sigma\rangle ,
\]
the sets $I$ and $E$ being read as sets of bonds, and the subscripts $II$ and $II^{c}$ having
the meaning fixed in the proposition on the joint bond sets. Namely,
\begin{equation}\label{9.5:eq-joint-bond-sets-1}
\|\delta\mathsf{D}\|_{\delta/4},\ \|\delta\Gamma\|_{\delta/4}
\;\le\; C\,C_F^{\mathrm{mix}}\,e^{2^d\kappa_1}\sum_{i\ge j_{\min}}\vartheta_w^{\,i}
\end{equation}
on the whole polydisc $|\sigma(\Delta)|\le e^{\kappa_1}$, the thin windows included. The bound
\eqref{9.5:eq-joint-bond-sets-1} is the hypothesis of the proposition on joint bond sets, in the
form that proposition requires for the joint bond set with cells in proper scales.
The naturality of kernels between bonds of different strata is not part of this statement.
\end{corollary}

\begin{proof}
The proof of Corollary~\ref{9.5:cor-two-cutoff-difference} uses two kinds of input.

The cells enter that proof only through the M\"obius decoupling estimates
\eqref{9.5:eq-moebius-decoupling-a} and through the lemma on mixed-size cells that the proof
quotes. Both are statements about a family of cells satisfying the three conditions displayed in
\eqref{9.5:eq-scaled-length-a}. By the statement of the corollary, these conditions hold for the
cells $\pi^{\mathrm{mix}}$ of $\boldsymbol{\Omega}=\mathbb{B}^{(n)}_k$. Both inputs are therefore
available for $\pi^{\mathrm{mix}}$.

The operators enter that proof only through the walk-expansion bounds
\eqref{9.5:eq-walk-influence-c} and through the two-cutoff comparison of natural functionals.
These are hypotheses (a) and (b).

By the statement on the bounds for the operators of the fixed list $\mathfrak{O}$, on which the
proof of Corollary~\ref{9.5:cor-two-cutoff-difference} rests, the proof of those bounds depends
on $Z_0$ only through its set of bonds. Replacing $Z_0$ and $Z$ by their joint $(A_j,B_j)$ bond
sets $I$ and $E$ is therefore admissible in each step.

Hence the argument of Corollary~\ref{9.5:cor-two-cutoff-difference} applies without change to
the operators
\[
T,\qquad (T_{II})^{-1}=\mathcal{C}(Z_0),\qquad \mathcal{C}^{1/2},
\]
and to the kernel $\Gamma=(T\langle\sigma\rangle)_{II^{c}}\mathcal{C}^{1/2}\langle\sigma\rangle$
built from them. This holds on the whole polydisc $|\sigma(\Delta)|\le e^{\kappa_1}$, which gives
\eqref{9.5:eq-joint-bond-sets-1}.
\end{proof}

\begin{corollary}[Ba{\l}aban's decoupling dichotomy on multi-scale sets]
\label{9.5:cor-decoupling-dichotomy}
Let $\hat{\boldsymbol{\Omega}}$ be an admissible refinement of a label $\boldsymbol{\Omega}$
(the label $\boldsymbol{\Omega}$ itself included), let $\sigma_0\subset\pi^{\mathrm{mix}}$, let
$\omega$ be a walk of $\hat{\boldsymbol{\Omega}}$ whose random-walk domain satisfies
$\operatorname{rd}(\omega)\subset\operatorname{loc}(\omega)$, and put
\begin{equation*}
\mathsf{m}(\omega):=\#\{\Delta\in\sigma_0:\ \Delta\cap\operatorname{rd}(\omega)\neq\emptyset\}
\end{equation*}
(the random-walk domain $\operatorname{rd}(\omega)$ of a walk is that of \cite{B-CMP116} and
\cite[p.~276]{B-CMP119}, and $\mathsf{m}(\omega)$ counts the parameters $s$ connected with
$\omega$ in the sense of \cite[(1.11)]{B-CMP116}). Assume \eqref{9.5:eq-walk-influence-a}. Then
\begin{equation}\label{9.5:eq-decoupling-dichotomy-1}
\text{either}\quad \mathsf{m}(\omega)\le 2^{d}
\quad\text{or}\quad
\hat{d}(\omega)\ge \frac{\mathsf{m}(\omega)\,M}{2^{d+3}L}.
\end{equation}
Consequently the dichotomy of \cite[(1.11)]{B-CMP116}, namely that $m>2^{4}$ implies
$\delta_0 d(\omega)\ge\delta_1 mM$, where $m$ is the number of decoupling parameters attached
to $\omega$ and $d(\omega)$ is the length of a shortest tree meeting all localisation domains
of $\omega$, holds for the walks of every admissible refinement, in particular on every
determining set $\mathfrak{B}_Z$, with $m=\mathsf{m}(\omega)$, with $\hat{d}(\omega)$ in the
role of the tree length $d(\omega)$ and with the constant
\begin{equation}\label{9.5:eq-decoupling-dichotomy-2}
\delta_1:=\frac{\delta_0}{2^{d+3}L},
\end{equation}
which depends on $L$ as well as on $\delta_0$. Moreover, for every $\theta\in(0,1]$ (a free
parameter here, unrelated to the exponent $\theta$ of the two-cutoff comparison) with
$\theta\delta_0M\ge 2^{d+3}L\kappa_1$,
\begin{equation}\label{9.5:eq-decoupling-dichotomy-3}
e^{\kappa_1\mathsf{m}(\omega)}\,e^{-\delta_0\hat{d}(\omega)}
\le e^{2^{d}\kappa_1}\,e^{-(1-\theta)\delta_0\hat{d}(\omega)}.
\end{equation}
\end{corollary}

\begin{proof}
Let $\Delta\in\sigma_0$ meet $\operatorname{rd}(\omega)$. Since
$\operatorname{rd}(\omega)\subset\operatorname{loc}(\omega)$, the cell $\Delta$ meets
$\operatorname{loc}(\omega)$, and therefore, by \eqref{9.5:eq-walk-influence-b}, it meets
$X^{(5)}$ for some cube $X$ of $\omega$.
As $X^{(5)}\subset X^{+}$, the cell $\Delta$, which lies in $\pi^{\mathrm{mix}}$ because
$\sigma_0\subset\pi^{\mathrm{mix}}$, meets $X^{+}$, and by the definition of the influence set
in \eqref{9.5:eq-walk-influence-d} it belongs to $S^{\mathrm{mix}}(\omega)$. Distinct cells
counted in $\mathsf{m}(\omega)$ are distinct elements of $S^{\mathrm{mix}}(\omega)$, so
\begin{equation}\label{9.5:eq-decoupling-dichotomy-4}
\mathsf{m}(\omega)\le \#S^{\mathrm{mix}}(\omega).
\end{equation}
Now apply part (a$'$) of the mixed-size counting bound
(Proposition~\ref{9.5:prop-mixed-counting}) to $\omega$; that proposition applies to every walk
of every admissible refinement of $\boldsymbol{\Omega}$, under the hypotheses (H1)--(H4) of that
proposition, which are in force throughout this section, the walk geometry
\eqref{9.5:eq-walk-influence-b} and the condition $L\ge7$ of the setting of
Definition~\ref{9.5:def-walk-influence}, together with \eqref{9.5:eq-walk-influence-a}. For a
walk of the refinement $\hat{\boldsymbol{\Omega}}$, the clause of
Proposition~\ref{9.5:prop-mixed-counting} on walks of admissible refinements bounds the least
length of a tree meeting all localisation domains of $\omega$ measured by the contour length
$\operatorname{len}_{\hat{\boldsymbol{\Omega}}}$ of the refinement, that is, $\hat{d}(\omega)$.
It gives: either $\#S^{\mathrm{mix}}(\omega)\le 2^{d}$, or $\hat{d}(\omega)$ is at least
$\#S^{\mathrm{mix}}(\omega)M/(2^{d+3}L)$. Suppose $\mathsf{m}(\omega)>2^{d}$. Then
$\#S^{\mathrm{mix}}(\omega)>2^{d}$ by \eqref{9.5:eq-decoupling-dichotomy-4}, so the second
alternative holds, and by \eqref{9.5:eq-decoupling-dichotomy-4} again
\begin{equation*}
\hat{d}(\omega)\ge \frac{\#S^{\mathrm{mix}}(\omega)\,M}{2^{d+3}L}
\ge \frac{\mathsf{m}(\omega)\,M}{2^{d+3}L}.
\end{equation*}
This is \eqref{9.5:eq-decoupling-dichotomy-1}.

Since $d=4$, the condition $m>2^{4}$ of \cite[(1.11)]{B-CMP116}, read with
$m=\mathsf{m}(\omega)$, is $\mathsf{m}(\omega)>2^{d}$, and then
\eqref{9.5:eq-decoupling-dichotomy-1} multiplied by
$\delta_0$ gives $\delta_0\hat{d}(\omega)\ge\delta_0\mathsf{m}(\omega)M/(2^{d+3}L)
=\delta_1\mathsf{m}(\omega)M$ with $\delta_1$ as in \eqref{9.5:eq-decoupling-dichotomy-2}.

For \eqref{9.5:eq-decoupling-dichotomy-3} we use the two alternatives of
\eqref{9.5:eq-decoupling-dichotomy-1}. If $\mathsf{m}(\omega)\le 2^{d}$, then
$e^{\kappa_1\mathsf{m}(\omega)}\le e^{2^{d}\kappa_1}$, and
$e^{-\delta_0\hat{d}(\omega)}\le e^{-(1-\theta)\delta_0\hat{d}(\omega)}$ because
$\theta\delta_0\hat{d}(\omega)\ge0$. If $\mathsf{m}(\omega)>2^{d}$, then by
\eqref{9.5:eq-decoupling-dichotomy-1} and the assumption $2^{d+3}L\kappa_1\le\theta\delta_0M$,
\begin{equation*}
\kappa_1\mathsf{m}(\omega)\le \kappa_1\,\frac{2^{d+3}L\,\hat{d}(\omega)}{M}
\le \theta\delta_0\hat{d}(\omega),
\end{equation*}
hence
$e^{\kappa_1\mathsf{m}(\omega)}e^{-\delta_0\hat{d}(\omega)}\le e^{-(1-\theta)\delta_0\hat{d}(\omega)}$,
which is at most the right-hand side of \eqref{9.5:eq-decoupling-dichotomy-3} since
$e^{2^{d}\kappa_1}\ge1$.
\end{proof}

\begin{remark}
(i) With $\theta=1/16$, the condition under which \eqref{9.5:eq-decoupling-dichotomy-3} holds
reads $\delta_1M\ge16\kappa_1$ with $\delta_1$ from \eqref{9.5:eq-decoupling-dichotomy-2},
that is,
\begin{equation*}
\delta_0M\ge 2^{d+7}L\kappa_1 ,
\end{equation*}
a lower bound on $M$, chosen after $L$ and $\kappa_1$. In this form the floor
$\delta_1M\ge16\kappa_1$ of the localisation lemma for walk terms is used on every admissible
refinement: for $\mathsf{m}(\omega)>2^{d}$ the chain
\begin{equation*}
e^{\kappa_1\mathsf{m}(\omega)}\le e^{\delta_1\mathsf{m}(\omega)M/16}
\le e^{\delta_0\hat{d}(\omega)/16}
\end{equation*}
holds, its first step by $\delta_1M\ge16\kappa_1$ and its second by
Corollary~\ref{9.5:cor-decoupling-dichotomy}.

(ii) On multi-scale sets the constant $\delta_1$ cannot exceed $\delta_0/L$. Consider a flat
interface between the regions of levels $j$ and $j+1$ of the label, and the walk
$\omega=(\square_0,\dots,\square_N)$ of cubes of the cover $\mathcal{D}_{j+1}$ centred at the
points $i\,u_{j+1}e_2$, $0\le i\le N$, where $e_2$ is the unit vector of a coordinate axis
lying in the interface. Then $\hat{d}(\omega)\le N\mathfrak{m}$, while
$\omega$ meets at least $N\mathfrak{m}L/M$ cells of level $j$, each of side $M_j$; when these
cells are the cells of $\sigma_0$ counted in $\mathsf{m}(\omega)$, this gives
\begin{equation*}
\hat{d}(\omega)\le \frac{\mathsf{m}(\omega)\,M}{L}.
\end{equation*}
An inequality $\delta_0\hat{d}(\omega)\ge\delta_1\mathsf{m}(\omega)M$ valid for all such walks
therefore forces $\delta_1\le\delta_0/L$; the constant \eqref{9.5:eq-decoupling-dichotomy-2}
is smaller than this necessary bound only by the factor $2^{d+3}$. Thus the dependence of
$\delta_1$ on $L$ cannot be removed; the same factor $M/(2^{d+3}L)$ already appears in the
second alternative of Proposition~\ref{9.5:prop-cells-near-graph}.
\end{remark}

\begin{definition}[Constants and the two floors on $M$]\label{9.5:not-floors-on-m}
In this section $M$ is the integer $M=L^{m'}$ of Notation~\ref{0.2:not-glossary}, with $m'$
free, and the following numbers are used. Each is given by an
explicit expression in the parameters listed after it. None of them depends on the lattice unit
$\eta$, on a level $k$ or on $K$, on the age of a large-field region, on a number of zones,
windows or steps, on the size of the torus, on $N$, or on the coupling $g$.
\begin{itemize}
\item $r_0:=M/(4L)$; it depends on $(L,M)$. In this section $r_0$ denotes this number only; it is
unrelated to the depth offset $r_0$ of Notation~\ref{0.2:not-glossary}.
\item $D:=(L+2)^d-L^d$, a bound on the number of cells that one cell touches; it depends on
$(d,L)$.
\item $C_d^{\mathrm{mix}}:=2^{d+3}L$ and $c_d^{\mathrm{mix}}:=(2D^2)^{2^{d+3}L}$; they depend on
$(d,L)$.
\item $\mathsf{q}$ is the largest number of cover cubes in one factor type of the walk expansions
of the operators of the list $\mathfrak{O}$ (Definition~\ref{9.5:def-walk-influence}); among the
expansions displayed in Ba{\l}aban's papers the largest such number is $4$. It depends on no
parameter.
\item $\mathsf{u}:=\max\{\mathfrak{m},M_1,\mathsf{r}_e\}$, where $\mathsf{r}_e:=2L$ is the range of
the local end factor of the walk expansions; it depends on $(L,\mathfrak{m},M_1)$.
\item $\rho_{\times}:=L\mathsf{u}(20\mathsf{q}+5L+7)$, the tether radius, in scaled units; it
depends on $(L,\mathsf{q},\mathfrak{m},M_1)$.
\item $C_0^{\mathrm{mix}}:=\exp(2^{d-2}\delta_0\rho_{\times})$, with $\delta_0$ the decay rate of
the bound on walk terms (Definition~\ref{9.5:def-walk-influence}); it depends on
$(d,L,\mathsf{q},\mathfrak{m},M_1,\delta_0)$.
\item $X_{\kappa}:=(4D^2)^{2^d}e^{2^d\kappa_1}$; it depends on $(d,L,\kappa_1)$.
\item $C_N^{\mathrm{mix}}:=8(4D^2)^{2^d}$; it depends on $(d,L)$.
\item $C_F^{\mathrm{mix}}:=2(4D^2)^{2^d}(2C_V\mathsf{N}B_0)^{1/2}$, where $B_0$ is the constant of
the bound on walk terms and $C_V$, $\mathsf{N}$ are the constants of the two-cutoff comparison of
natural functionals used in Corollary~\ref{9.5:cor-two-cutoff-difference}; it depends on $(d,L)$
and on $C_V$, $\mathsf{N}$, $B_0$.
\item $\delta_1:=\delta_0/(2^{d+3}L)$, the rate with which the dichotomy of
\cite[(1.11)]{B-CMP116} is used on multi-scale sets
(Corollary~\ref{9.5:cor-decoupling-dichotomy}); it depends on $(d,L,\delta_0)$.
\end{itemize}

\emph{The two floors.} Two conditions are imposed on $M$. Both are floors in the sense of
Notation~\ref{2.3:not-stages-first}; neither is a ceiling, a window or a cap
(Definition~\ref{2.3:def-cap}). The first is the floor \eqref{9.5:eq-walk-influence-a},
\begin{equation*}
M\ \ge\ 2^{d+3}L^2(20\mathsf{q}+5L+7)\cdot\max\{\mathfrak{m},M_1,2L\},
\end{equation*}
placed after $(d,L,\mathfrak{m},M_1)$. Since $\mathsf{u}=\max\{\mathfrak{m},M_1,2L\}$, its right
member equals $2^{d+3}L\cdot L\mathsf{u}(20\mathsf{q}+5L+7)=C_d^{\mathrm{mix}}\rho_{\times}$; thus
the first floor reads $M\ge C_d^{\mathrm{mix}}\rho_{\times}$. It replaces the floor
$M\ge 2C_d\sqrt{d}\cdot 30LM_1$ of the single-size counting bounds, $C_d$ being the constant of
those bounds. At $d=4$ that floor reads $M\ge 120\,C_dLM_1$, while the right member of the first
floor is at least $2^7L^2(5L+7)M_1$; at $d=4$ the first floor implies that floor for every
$L\ge2$. The second is the floor
\eqref{9.5:eq-interpolated-operators-a},
\begin{equation*}
cM\ \ge\ 16L\bigl[2^d(\kappa_1+\log 4D^2)+\log 2\bigr],
\end{equation*}
which is used at three values of $c$: at $c=\tfrac12\delta_0$ in
Theorem~\ref{9.5:thm-interpolated-operators}; at $c=\tfrac14\delta_0$ in
Corollary~\ref{9.5:cor-two-cutoff-difference}, in the corollary that follows it, and in the
one-cutoff form of the extension of Theorem~\ref{9.5:thm-interpolated-operators}; and at
$c=\tfrac18\delta_0$ in the sup-norm bounds of the two-cutoff comparison and in the two-cutoff
form of that extension.
All three are served by the single floor
\begin{equation}\label{9.5:eq-floors-on-m-1}
\delta_0M\ \ge\ 2^7L\bigl[2^d(\kappa_1+\log 4D^2)+\log 2\bigr],
\end{equation}
placed after $(d,L,\delta_0,\kappa_1)$. Where an estimate of this book applies the second floor
with $\kappa_1$ replaced by $2\kappa_1+2$ or by $\kappa_1+1$, the floor \eqref{9.5:eq-floors-on-m-1}
is read at that value of $\kappa_1$. The floor \eqref{9.5:eq-floors-on-m-1} replaces the floor
$\tfrac18\delta_0M\ge 3C_d(\kappa_1+\log 8c_d)+1$ of the single-size interpolation theorem,
$C_d$ and $c_d$ being the constants of that theorem, and the doubled form of that floor.

\emph{Standing conditions.} These conditions are in force throughout:
\begin{itemize}
\item $L\ge 7$ is odd, and $R\ge L$, where $R$ is the admissibility parameter of the refinements
of Definition~\ref{9.5:def-walk-influence} (not the ball radius $R$ of
Notation~\ref{0.2:not-glossary});
\item $(2L-3)M_1>R\mathfrak{m}$ and $M>R\mathfrak{m}$, for the admissibility of shells;
\item $LM_1$ divides $M$;
\item $5s_{k+1}<\ell_{\mathbb{T}}$.
\end{itemize}

\emph{Two readings of conditions imposed elsewhere.} The exponential rate in clause (n4) of
\eqref{9.5:eq-interpolated-operators-b} is $\delta_0M/(8L)$, and the bound of that clause is
$4B_0e^{-\delta_0M/(8L)}$. The proposition of this book that uses the bound of clause (n4) imposes
the smallness condition $\tfrac32K_1\varepsilon_{\mathrm{pr}}\le\mathfrak{b}$, in which
$\varepsilon_{\mathrm{pr}}$ is the smallness parameter supplied by clause (n4) and $\mathfrak{b}$
is the bound with which it is compared; with $\varepsilon_{\mathrm{pr}}=4B_0e^{-\delta_0M/(8L)}$,
this condition reads
\begin{equation}\label{9.5:eq-floors-on-m-2}
6K_1B_0e^{-\delta_0M/(8L)}\ \le\ \frac{\mathsf{b}_{\mathtt{f}}}{8d\,\mathsf{k}_*(LM)^4},
\end{equation}
where $K_1$ is the constant multiplying $\varepsilon_{\mathrm{pr}}$ in that condition, and
$\mathsf{k}_*$ and $\mathsf{b}_{\mathtt{f}}$ are the constants of that proposition which enter its right
member $\mathfrak{b}=\mathsf{b}_{\mathtt{f}}/(8d\,\mathsf{k}_*(LM)^4)$; all three are positive, are fixed
before $M$ and do not depend on $M$. This is a condition of that proposition, not a further floor
of this notation; by (c) below it holds for every $M$ beyond a threshold, so it is a floor on $M$
placed after $L$. Wherever the dichotomy of \cite[(1.11)]{B-CMP116} is applied on a
multi-scale set, it is applied with $\delta_1$, by (d) below. The floor on $M$ relative to
$\kappa_1$ of the decoupling lemma that applies this dichotomy is read with this $\delta_1$ as its
rate; it remains a floor on $M$, placed after $(L,\kappa_1)$.

\emph{Order.} In the order of choice fixed in Notations~\ref{2.3:not-stages-middle}
and~\ref{2.3:not-stages-last}, the parameters $L$, $\mathfrak{m}$, $M_1$, $\kappa$ and $\kappa_1$
are fixed before $M$, and every quantity that depends on $M$ is fixed after $M$; both floors are
imposed at the place of $M$, and the two conditions read in the preceding paragraph remain floors
on $M$ at that place. By (a) below, parameters fixed after $\kappa_1$ and before $M$ may depend on
$C_N^{\mathrm{mix}}$, $C_0^{\mathrm{mix}}$, $X_{\kappa}$, $D$ and $\rho_{\times}$; in particular
the ceilings on $\varepsilon_1$, $\varepsilon_2$, $\varepsilon_3$, placed after $\kappa_1$ and
before $M$, may depend on them. The constant of the lemma that uses
Corollary~\ref{9.5:cor-two-cutoff-difference}, namely the constant that collects
$C_F^{\mathrm{mix}}e^{2^d\kappa_1}$, is fixed after $M$ and before $\gamma$.

With these conventions the following hold.
\begin{enumerate}
\item[(a)] $D$, $\rho_{\times}$, $C_0^{\mathrm{mix}}$, $X_{\kappa}$ and $C_N^{\mathrm{mix}}$ are
given by expressions in which $M$ does not occur. The same holds for $C_d^{\mathrm{mix}}$,
$c_d^{\mathrm{mix}}$, $\mathsf{u}$ and $\delta_1$; $C_F^{\mathrm{mix}}$ depends on $M$ at most
through $C_V$, $\mathsf{N}$ and $B_0$.
\item[(b)] If $\delta_0>0$, $M>0$ and \eqref{9.5:eq-floors-on-m-1} holds, then the second floor
holds for every $c\ge\tfrac18\delta_0$, in particular at $c=\tfrac12\delta_0$,
$\tfrac14\delta_0$ and $\tfrac18\delta_0$.
\item[(c)] If $\delta_0>0$ and $K_1$, $B_0$, $\mathsf{k}_*$, $\mathsf{b}_{\mathtt{f}}$ are positive, there
is a threshold, depending only on $(d,L,\delta_0,K_1,B_0,\mathsf{k}_*)$, the
constant $\mathsf{b}_{\mathtt{f}}$ being fixed, such that \eqref{9.5:eq-floors-on-m-2} holds for every $M$
at or beyond it.
\item[(d)] Under the first floor, for every walk $\omega$ as in
Corollary~\ref{9.5:cor-decoupling-dichotomy}, either $\mathsf{m}(\omega)\le 2^d$ or
$\delta_0\hat{d}(\omega)\ge\delta_1\mathsf{m}(\omega)M$.
\end{enumerate}
\end{definition}

\begin{proof}
(a) The list above gives each quantity by an expression in the parameters named in its line. For
$D$, $C_d^{\mathrm{mix}}$, $c_d^{\mathrm{mix}}$ and $C_N^{\mathrm{mix}}$ these parameters are
$d$ and $L$. For $\mathsf{u}$ and $\rho_{\times}$ they are $L$, $\mathsf{q}$, $\mathfrak{m}$ and
$M_1$. $C_0^{\mathrm{mix}}$ is an expression in $d$, $\delta_0$ and $\rho_{\times}$, and
$X_{\kappa}$ is an expression in $d$, $L$ and $\kappa_1$. $C_F^{\mathrm{mix}}$ is an expression in
$d$, $L$, $C_V$, $\mathsf{N}$ and $B_0$, and $\delta_1$ is an expression in $d$, $L$ and
$\delta_0$. The letter $M$ occurs in none of these expressions; among the numbers of the list it
occurs only in $r_0$. In particular $C_F^{\mathrm{mix}}$ can depend on $M$ only through $C_V$,
$\mathsf{N}$ and $B_0$.

(b) Let $c\ge\tfrac18\delta_0$. Since $M>0$, $cM\ge\tfrac18\delta_0M$. By
\eqref{9.5:eq-floors-on-m-1},
\begin{equation*}
\tfrac18\delta_0M\ \ge\ \tfrac18\cdot 2^7L\bigl[2^d(\kappa_1+\log 4D^2)+\log 2\bigr]
=16L\bigl[2^d(\kappa_1+\log 4D^2)+\log 2\bigr].
\end{equation*}
The two inequalities together give the second floor at $c$. The values $\tfrac12\delta_0$,
$\tfrac14\delta_0$ and $\tfrac18\delta_0$ are all at least $\tfrac18\delta_0$.

(c) Multiplying \eqref{9.5:eq-floors-on-m-2} by $8d\,\mathsf{k}_*(LM)^4>0$ shows that it is
equivalent to
\begin{equation*}
48\,d\,K_1B_0\mathsf{k}_*L^4M^4e^{-\delta_0M/(8L)}\ \le\ \mathsf{b}_{\mathtt{f}}.
\end{equation*}
Put $t:=\delta_0M/(8L)>0$. The power series of $e^t$ gives $e^t\ge t^5/5!$, that is,
$e^{-t}\le 5!\,t^{-5}=5!\,(8L)^5\delta_0^{-5}M^{-5}$. Hence the left member of the last display is
at most $C/M$, where
\begin{equation*}
C:=48\cdot 5!\cdot 8^5\,d\,K_1B_0\mathsf{k}_*L^9\delta_0^{-5}.
\end{equation*}
$C$ depends only on $(d,L,\delta_0,K_1,B_0,\mathsf{k}_*)$. Therefore \eqref{9.5:eq-floors-on-m-2}
holds for every $M\ge C/\mathsf{b}_{\mathtt{f}}$.

(d) Corollary~\ref{9.5:cor-decoupling-dichotomy} applies under the first floor. It gives, for
every such walk $\omega$, either $\mathsf{m}(\omega)\le 2^d$ or
$\hat{d}(\omega)\ge\mathsf{m}(\omega)M/(2^{d+3}L)$. In the second case we multiply by
$\delta_0>0$ and use $\delta_1=\delta_0/(2^{d+3}L)$:
\begin{equation*}
\delta_0\hat{d}(\omega)\ \ge\ \frac{\delta_0}{2^{d+3}L}\,\mathsf{m}(\omega)M
=\delta_1\mathsf{m}(\omega)M. \qedhere
\end{equation*}
\end{proof}

\begin{definition}[Dates, dated sets, weights and matrix norms]\label{9.5:not-dates-weights-norms}
Throughout this section $\mathsf{m}\in\{\mathsf{A},\mathsf{B}\}$ denotes one of the two runs
compared in the uniqueness theorem, and one label, in the sense of the one-run expansions on
labels, is fixed. We use the dates $d$, the dated sets
$\mathfrak{B}_d$ and the order $\preccurlyeq$ on dated sets in the form in which they are fixed
together with the arrest theorem. For a date $d$ and a point $y$ we write $\lambda_d(y)$ for the
level of $y$ in $\mathfrak{B}_d$,
\begin{equation*}
\ell_d(y) := L^{\lambda_d(y)}\eta
\end{equation*}
for the unit of length at $y$ (more generally, $\ell_n := L^{n}\eta$ for a level $n$, so that
$\ell_d(y)=\ell_{\lambda_d(y)}$), $w_d(y)$ for the actual weight attached to $y$ by $\mathfrak{B}_d$,
and $k(d)$ for the density of $d$. With the seven radii $r_a(n)$, functions of a level $n$, the
\emph{weighted thresholds} are
\begin{equation}\label{9.5:eq-dates-weights-norms-1}
\theta^{d}_{a}(y) := w_d(y)\,r_a\bigl(\lambda_d(y)\bigr).
\end{equation}

\emph{Birth dates.} Consider a triple $(F,X,j)$ consisting of a stored left-over term
$F=\mathbb{C}^{(i)}_{k_a}(X)$ of stage $i$ at step $k_a$, its region $X$ and its level $j=k_a$.
It is \emph{born} at the date $\tau := (k_a;i)$, and its
dated set is $\mathfrak{B}_\tau := \mathbb{B}^{(i)}_{k_a}$, frozen on the pieces live at $\tau$. Thus
a stored left-over term $\mathbb{C}^{(i)}_k$ is read with the dated set of its birth date $(k;i)$.

\emph{Cube side.} We put $u := M\ell_j$; this is the side of one cube of the partition $\pi_j$.

\emph{Matrix norms.} $|\cdot|$ and $\|\cdot\|$ denote the pair of norms on $N\times N$
matrices fixed together with the arrest theorem. For all such matrices $P$, $A$, $Q$ and every
$v\in\mathrm{SU}(N)$ they satisfy
\begin{equation}\label{9.5:eq-dates-weights-norms-2}
|PAQ| \le \|P\|\,|A|\,\|Q\|, \qquad \|A\| \le |A|, \qquad \|v\| = 1, \qquad |vAv^{-1}| = |A|.
\end{equation}
Consequently, for all $N\times N$ matrices $A$ and $B$,
\begin{equation}\label{9.5:eq-dates-weights-norms-3}
|AB| \le |A|\,|B|, \qquad |e^{A}-\mathbf{1}| \le e^{|A|}-1, \qquad \|e^{A}\| \le e^{|A|}.
\end{equation}
\end{definition}

\begin{proof}
We derive \eqref{9.5:eq-dates-weights-norms-3} from \eqref{9.5:eq-dates-weights-norms-2}.

First inequality. The identity $\mathbf{1}$ lies in $\mathrm{SU}(N)$, so $\|\mathbf{1}\|=1$. Apply the
first relation of \eqref{9.5:eq-dates-weights-norms-2} with $P=\mathbf{1}$ and $Q=B$, and then the
second relation to $B$:
\begin{equation*}
|AB| = |\mathbf{1}AB| \le \|\mathbf{1}\|\,|A|\,\|B\| = |A|\,\|B\| \le |A|\,|B|.
\end{equation*}

Second inequality. By the first inequality and induction on $n$, $|A^n|\le|A|^n$ for every
$n\ge1$. The exponential series converges absolutely in the finite-dimensional space of
$N\times N$ matrices, so $e^{A}-\mathbf{1}=\sum_{n\ge1}A^n/n!$. The triangle inequality applied to
the partial sums, together with the continuity of the norm $|\cdot|$, gives
\begin{equation*}
|e^{A}-\mathbf{1}| \le \sum_{n\ge1}\frac{|A^n|}{n!} \le \sum_{n\ge1}\frac{|A|^n}{n!}
= e^{|A|}-1.
\end{equation*}

Third inequality. By the triangle inequality for $\|\cdot\|$, the relation $\|\mathbf{1}\|=1$, the
second relation of \eqref{9.5:eq-dates-weights-norms-2} applied to $e^{A}-\mathbf{1}$, and the
second inequality,
\begin{equation*}
\|e^{A}\| \le \|\mathbf{1}\| + \|e^{A}-\mathbf{1}\| \le 1 + |e^{A}-\mathbf{1}| \le e^{|A|}.
\end{equation*}
\end{proof}

\begin{definition}[Scope, fattened sets and the ball collar]\label{9.5:def-fattened-sets}
We use the dates $d$, the levels $\lambda_d(y)$ and the
units $\ell_d(y)=L^{\lambda_d(y)}\eta$ of
Notation~\ref{9.5:not-dates-weights-norms}. Fix one of the two compared constructions,
$\mathsf{m}\in\{\mathsf{A},\mathsf{B}\}$, one label, and a stored term $(F,X,j)$ born at the
date $\tau$.

\emph{Fattened sets.} For a level $m$ write $\ell_m:=L^{m}\eta$, so that
$\ell_d(y)=\ell_{\lambda_d(y)}$. Let $u:=M\ell_j$, the side of one $\pi_j$-cube. Write
$\mathbb{T}_\eta$ for the lattice of spacing $\eta$ on the torus. For $0<s<1$ put
\begin{equation}\label{9.5:eq-fattened-sets-c}
X^{(s)}:=\bigl\{y\in\mathbb{T}_\eta:\ \operatorname{dist}_\infty(y,X)\le su\bigr\}.
\end{equation}
Recall the hull $\widehat{X}$ of $X$, which is a union of $\pi_j$-cubes: it is the union of $X$
with all $\pi_j$-cubes that are $\infty$-adjacent to a $\pi_j$-cube of $X$. A bond or a plaquette
is \emph{of $S$} if all its corners lie in $S$.

\emph{Scope.} Let $\Omega_1$ be the domain of the first renormalisation step of
\cite{B-CMP119}, and let $\Omega_1^{\sim}\supset\Omega_1$
be its enlargement by one layer of cubes, as fixed with the dated level-zero clause; put
$\Lambda_0:=\Omega_1^{\sim}\setminus\Omega_1$, the whole never-averaged collar of the label.
At every date $d$ the dated collar $\operatorname{Col}(d)$, on which the dated level-zero clause
is imposed, is contained in $\Lambda_0$, and older dates have the larger dated collar. The stored
term is said to satisfy the scope condition if
\begin{equation}\label{9.5:eq-fattened-sets-a}
\operatorname{dist}_\infty\bigl(X^{(3/4)},\Lambda_0\bigr)>2\eta .
\end{equation}
Condition \eqref{9.5:eq-fattened-sets-a} holds for every stored term of a label with
$\Lambda_0=\emptyset$, and therefore for every stored term of an aligned label, in both
constructions; this is shown later in this chapter, where the use of the level-zero clause is
discussed. Nothing is asserted in this chapter for a stored term whose hull $\widehat{X}$ touches
$\Lambda_0$. Under \eqref{9.5:eq-fattened-sets-a} no level-zero clause is used at a point of
$X^{+}$ or of the collar $S^{+}$ of a sub-domain $S\subset X$ (these collars are defined in
\eqref{9.5:eq-fattened-sets-b}).

\emph{The ball collar.} Let $x$ be a point and consider a frame at level $m$, in the sense of the
construction of local charts. The certificate ball of that construction,
\begin{equation*}
B(x)=\bigl\{z:\ |z-y_0|_\infty\le\varrho_0\ell_m\bigr\},
\end{equation*}
is centred at the base point $y_0$ of the block of $x$ in the frame, so that
$|x-y_0|_\infty\le\ell_m$; it is fixed together with the size condition
$M/M_1\ge3L(\varrho_0+2L)$. For a date $d$, a point $x$ and a set $Y$ put
\begin{equation}\label{9.5:eq-fattened-sets-b}
B^{\times}(x;d):=\bigl\{z:\ |z-x|_\infty\le(\varrho_0+1)\ell_d(x)\bigr\},\qquad
Y^{+}_{d}:=Y\cup\bigcup_{x\in Y}B^{\times}(x;d).
\end{equation}
Wherever, in the bounds used below for a term of date $d$, a collar is written $Y^{+}$ or $X^{+}$,
it stands for $Y^{+}_{d}$, respectively $X^{+}_{d}$, at that term's own date $d$. For the stored
term fixed above we put $X^{+}:=X^{+}_{\tau}$. The ball collar $X^{+}$ is not to be confused with
the hull $\widehat{X}$.

These sets have the following properties.
\begin{itemize}
\item[(a)] $X^{(s)}\subset\widehat{X}$ for $0<s<1$.
\item[(b)] $B^{\times}(x;d)$ contains the ball $B(x)$ of every frame whose unit set has level
at most $\lambda_d(x)$ at $x$, and its radius is smaller than the side $M\ell_d(x)$ of one
$M$-cube of the local scale.
\item[(c)] If $d''\le d$ and $S\subset Y$, then $S^{+}_{d''}\subset Y^{+}_{d}$.
\end{itemize}
\end{definition}

\begin{proof}
(a) Let $y\in X^{(s)}$. Since $s<1$, $\operatorname{dist}_\infty(y,X)\le su<u$. A point at
sup-distance less than $u$ from $X$ lies in a $\pi_j$-cube that is equal or $\infty$-adjacent to
a $\pi_j$-cube of $X$: the $\pi_j$-cubes tile the lattice by translates of one cube of side $u$,
so two points at sup-distance less than $u$ lie in cubes whose positions differ by at most one
cube in each coordinate direction. By the definition of the hull, $y\in\widehat{X}$.

(b) Consider a frame at level $m\le\lambda_d(x)$ at $x$, and let $z\in B(x)$. By the triangle
inequality,
\begin{equation*}
|z-x|_\infty\le|z-y_0|_\infty+|y_0-x|_\infty\le\varrho_0\ell_m+\ell_m=(\varrho_0+1)\ell_m .
\end{equation*}
Since $m\le\lambda_d(x)$ and $L>1$, we have $\ell_m=L^{m}\eta\le L^{\lambda_d(x)}\eta=\ell_d(x)$,
hence $|z-x|_\infty\le(\varrho_0+1)\ell_d(x)$, that is, $z\in B^{\times}(x;d)$. For the radius,
the size condition $M/M_1\ge3L(\varrho_0+2L)$ that accompanies the certificate ball gives, since
$1\le2L$,
\begin{equation*}
\varrho_0+1\le\varrho_0+2L\le\frac{M}{3LM_1};
\end{equation*}
as $3LM_1>1$,
\begin{equation*}
(\varrho_0+1)\ell_d(x)\le\frac{M\ell_d(x)}{3LM_1}<M\ell_d(x).
\end{equation*}

(c) Let $d''\le d$ and $S\subset Y$. By the monotonicity of levels in the date, clause (e1) of
the lemma on dated sets used with the arrest theorem, the levels satisfy
$\lambda_{d''}(x)\le\lambda_d(x)$ at every point $x$; as $L>1$, this gives
$\ell_{d''}(x)\le\ell_d(x)$, and hence, for every $x$,
\begin{align*}
B^{\times}(x;d'')&=\bigl\{z:|z-x|_\infty\le(\varrho_0+1)\ell_{d''}(x)\bigr\}\\
&\subset\bigl\{z:|z-x|_\infty\le(\varrho_0+1)\ell_d(x)\bigr\}=B^{\times}(x;d).
\end{align*}
Therefore, using $S\subset Y$ twice,
\begin{equation*}
S^{+}_{d''}=S\cup\bigcup_{x\in S}B^{\times}(x;d'')\subset Y\cup\bigcup_{x\in Y}B^{\times}(x;d)
=Y^{+}_{d}. \qedhere
\end{equation*}
\end{proof}

Under \eqref{9.5:eq-fattened-sets-a} the dated level-zero clause would be needed only for stored
terms outside the scope condition; the steps where it would enter are indicated in brackets in
what follows.

\begin{definition}[The relaxed continuation class]\label{9.5:def-relaxed-class}
We use the dates $d$, the dated sets $\mathfrak{B}_d$, the density $k(d)$ and the weighted
thresholds $\theta^{d}_{a}(y)$ of Notation~\ref{9.5:not-dates-weights-norms}, and the ball
collar $Y^{+}_{d}$ and the dated collar $\operatorname{Col}(d)$ of
Definition~\ref{9.5:def-fattened-sets}. For a set $\mathfrak{B}$, $\Omega_0(\mathfrak{B})$ is
the domain of \cite[p.~396]{B-CMP99b} formed for $\mathfrak{B}$, and $\mathfrak{g}^{c}$ is the
complexification of the Lie algebra $\mathfrak{g}$ of $\mathrm{SU}(N)$.

Let $d$ be a date of density $k=k(d)$, and let $Y$ be a union of cubes. The
\emph{relaxed continuation class} $\mathfrak{G}^{\natural}(\mathfrak{B}_d;Y)$ is the set of
pairs
\begin{equation*}
\tilde p=\bigl(e^{i\eta\tilde A}\tilde U,\ \tilde{\mathbb{J}}\bigr)\quad\text{on }
\Omega_0(\mathfrak{B}_d),
\end{equation*}
with $\tilde U$ an $\mathrm{SU}(N)$-valued bond field and $\tilde A$, $\tilde{\mathbb{J}}$
$\mathfrak{g}^{c}$-valued, which satisfy the following three conditions.

First, for every date $d''\le d$,
\begin{equation}\label{9.5:eq-relaxed-class-a}
\begin{gathered}
\text{on }\Omega_0(\mathfrak{B}_{d''}),\ \tilde U\ \text{satisfies \cite[(3.35)]{B-CMP99b}
and}\ \tilde A\ \text{satisfies \cite[(3.37)]{B-CMP99b}},\\
\text{both formed over }\mathfrak{B}_{d''}\text{, with }\nabla=\nabla^{\eta}_{\tilde U}
\text{ and at the radii }\alpha_0^{(99)},\ \alpha_1^{(99)};
\end{gathered}
\end{equation}
here $\alpha_0^{(99)}$ and $\alpha_1^{(99)}$ are the radii at which the coarse regularity
class $\mathfrak{K}_{109}$ of the background-response theorem is formed.

Second, at every $y\in Y^{+}_{d}$ the pointwise clauses of the continuation class
$\mathfrak{G}(\mathfrak{B}_d)$ hold:
\begin{equation}\label{9.5:eq-relaxed-class-b}
Q_a(\tilde p)(y)<5\,\theta^{d}_{a}(y),\qquad a\in\{1,3,5,6,7\},
\end{equation}
where at the points $y$ of $\operatorname{Col}(d)$ these clauses take the form of the dated
collar clauses
\begin{equation*}
\mathsf{q}_a(\tilde p)(y)<\mathsf{f}(d)\,\mathsf{a}_a .
\end{equation*}

Third, the current slot vanishes off the ball collar:
\begin{equation}\label{9.5:eq-relaxed-class-c}
\tilde{\mathbb{J}}(b)=0\qquad\text{for every bond } b \text{ that is not a bond of } Y^{+}_{d}.
\end{equation}

\emph{The relaxed readings.} The relaxed ambient convention is the refined ambient convention
(a sharpened form of the ambient convention; both fix what it means for a pair over $Y$ to be
a member of a space, a source domain or a stored-term domain over $Y$), with parts (a) and
(b), in whose part (a) the words ``a pair of $\mathfrak{G}(\mathfrak{B}_n)$ on the whole
domain'' are replaced by ``a pair of $\mathfrak{G}^{\natural}(\mathfrak{B}_n;Y)$ on the whole
domain''. Thus, with $\mathfrak{B}_n$ the set of part (a) of the refined ambient convention, a
pair over $Y$ is a member of a space, a source domain or a stored-term domain over $Y$ exactly
when it is
\begin{equation}\label{9.5:eq-relaxed-class-d}
\begin{gathered}
\text{the restriction to } Y \text{ of a pair of }
\mathfrak{G}^{\natural}(\mathfrak{B}_n;Y)\text{ on the whole domain,}\\
\text{obeying the fat direct clauses at every location of } Y .
\end{gathered}
\end{equation}
Part (b) of the refined ambient convention is then, word for word, condition
\eqref{9.5:eq-relaxed-class-b} on $Y^{+}\setminus Y$; here and below $Y^{+}$ is the ball collar
of Definition~\ref{9.5:def-fattened-sets} at the date of the set over which the convention is
read, so that $Y^{+}=Y^{+}_{d}$ for $\mathfrak{B}_d$. The requirement of part (a) that the
fat direct clauses hold at every location of $Y$ is kept unchanged. The relaxed ambient rider
is the ambient rider in which ``restrictions of pairs of $\mathfrak{G}(\mathfrak{B})$'' is read
as ``restrictions, for a space, a source domain or a stored-term domain over $Y$, of pairs of
$\mathfrak{G}^{\natural}(\mathfrak{B};Y)$''. For a stage $n$, the domain
$\mathfrak{D}^{\flat,\natural}_n(Y)$ is the inherited domain of stored terms
$\mathfrak{D}^{\flat}_n(Y)$ read with the relaxed ambient convention
\eqref{9.5:eq-relaxed-class-d} in place of the refined one. The \emph{$\natural$-reading} is
the relaxed ambient rider in force in every space, every source domain and every stored-term
domain.

\emph{The collar $Y^{+}\setminus Y$.} Condition \eqref{9.5:eq-relaxed-class-b} contains no
cube clause on $Y^{+}\setminus Y$. The only quantity that depends on part (b) of the refined
ambient convention at those points is the dimensionless size $a_{\flat}$ of the bump lemma,
which is formed from the entries $5$, $1$, $6$, $7$ alone. Clause $8$ at the locations of $Y$
is the stored-term domain's own local clause, a gauge condition on the part $\square\cap Y$ of
each cube $\square$ of the family at that location, and it is
not altered. The condition \cite[(3.35)]{B-CMP99b} on whole cubes is contained in
\eqref{9.5:eq-relaxed-class-a}.

\emph{Comparison with the continuation class and with the ambient convention.} At the
locations of $Y^{+}$ nothing is changed with respect to $\mathfrak{G}(\mathfrak{B}_d)$:
condition \eqref{9.5:eq-relaxed-class-b} is its pointwise condition there. Beyond $Y^{+}$, the
five level-attached clauses of
$\mathfrak{G}(\mathfrak{B}_d)$, whose radii are small in the coupling, are replaced by the two
clauses with coupling-free radii that the background-response theorem requires, imposed at
every set $\mathfrak{B}_{d''}$ with $d''\le d$, as in \eqref{9.5:eq-relaxed-class-a}; and the
second slot is cut off, as in \eqref{9.5:eq-relaxed-class-c}. Compared with the ambient
convention, its first half, ``a pair of $\mathfrak{K}_{109}(\mathfrak{B};\gamma_0)$ on the
whole domain of $\mathfrak{B}$'', is kept, at every earlier set and at the birth density; its
second half is kept on $Y$ and, on the collar, in the form of part (b) of the refined ambient
convention.
\end{definition}

\begin{lemma}[Sandwich, heredity and stability of the relaxed class]\label{9.5:lem-relaxed-heredity}
We use the dates, the dated sets $\mathfrak{B}_d$, the levels $\lambda_d$, the units $\ell_d$, the
weights $w_d$ and the weighted thresholds $\theta^{d}_{a}(y)$ of
Notation~\ref{9.5:not-dates-weights-norms}, the scope condition \eqref{9.5:eq-fattened-sets-a} and
the collars \eqref{9.5:eq-fattened-sets-b} of Definition~\ref{9.5:def-fattened-sets}, and the
relaxed continuation class $\mathfrak{G}^{\natural}(\mathfrak{B}_d;Y)$ of
Definition~\ref{9.5:def-relaxed-class}, with its clauses \eqref{9.5:eq-relaxed-class-a},
\eqref{9.5:eq-relaxed-class-b}, \eqref{9.5:eq-relaxed-class-c}. Assume:
\begin{itemize}
\item the order lemma for dated sets of the arrest theorem;
\item clauses (a), (d), (e1), (e2) and (e3) of the continuation-class lemma of the arrest theorem,
together with its nesting statement;
\item if the dated collar $\operatorname{Col}$ is non-empty, clauses (i) and (ii) of the collar
lemma;
\item for part (W3) only: clause (e4) of the continuation-class lemma, and, taken as stated, the
property of the maps on which the proof of (e4) rests at bonds astride the boundary of the region
$Y'$ of part (W3) below.
\end{itemize}
Let $d''\le d$ be dates of densities $k''\le k$, and let $S\subset Y$ be unions of cubes.

\emph{(W1) Sandwich.} For a pair $P=(\mathbb{U},\mathbb{J})$ put
$P^{Y}:=(\mathbb{U},\mathbf{1}_{Y^{+}_{d}}\mathbb{J})$. Then
\begin{equation}\label{9.5:eq-relaxed-heredity-a}
\begin{aligned}
&\text{(i)}\quad P\in\mathfrak{G}(\mathfrak{B}_d)\ \Longrightarrow\
P^{Y}\in\mathfrak{G}^{\natural}(\mathfrak{B}_d;Y),\\
&\text{(ii)}\quad \tilde p\in\mathfrak{G}^{\natural}(\mathfrak{B}_d;Y)\ \Longrightarrow\
\tilde p\restriction\Omega_0(\mathfrak{B}_{d''})\in
\mathfrak{K}_{109}\bigl(\mathfrak{B}_{d''};\gamma_0^{(k'')}\bigr),\\
&\text{(iii)}\quad \tilde p\in\mathfrak{G}^{\natural}(\mathfrak{B}_d;Y)\ \Longrightarrow\
a_{\flat}(\tilde p)\le\gamma_0^{(k'')} .
\end{aligned}
\end{equation}
In (i), moreover, $P^{Y}=P$ on the bonds of $Y^{+}_{d}$. In (ii) the membership holds at every
date $d''\le d$, the date $d$ itself included, each time at the density $k''$ of $d''$. In (iii)
the inequality holds in every frame of the bump lemma whose unit set is $\mathfrak{B}_{d''}$ and
whose sets $X'$ and $B$ satisfy $X'\cup B\subset Y^{+}_{d}$; $a_{\flat}$ is the dimensionless size
of the bump lemma.

\emph{(W2) Heredity.} For $\tilde p=(e^{i\eta\tilde A}\tilde U,\tilde{\mathbb{J}})$ put
\begin{equation*}
\tilde p^{S}:=\bigl(e^{i\eta\tilde A}\tilde U\restriction\Omega_0(\mathfrak{B}_{d''}),\
\mathbf{1}_{S^{+}_{d''}}\tilde{\mathbb{J}}\bigr).
\end{equation*}
Then, provided that the scope condition \eqref{9.5:eq-fattened-sets-a} holds for $S$ when
$\Lambda_0\neq\emptyset$,
\begin{equation}\label{9.5:eq-relaxed-heredity-b}
\tilde p\in\mathfrak{G}^{\natural}(\mathfrak{B}_d;Y)\ \Longrightarrow\
\tilde p^{S}\in\mathfrak{G}^{\natural}(\mathfrak{B}_{d''};S),
\end{equation}
and $\tilde p^{S}=\tilde p$ on the bonds of $S^{+}_{d''}$.

\emph{(W3) Maps.} Let $\mathfrak{O}$ be one of the six maps of the construction of the arrest
theorem, with source set $\mathfrak{B}'=\mathfrak{B}_d$ and target set
$\mathfrak{B}^{t}\preccurlyeq\mathfrak{B}'$. Let $p$ be a member of the source space of
$\mathfrak{O}$ over a union of cubes $Y'$, with a continuation
$\tilde p\in\mathfrak{G}^{\natural}(\mathfrak{B}';Y')$. Let $\tilde q$ be the continuation of the
image of $p$ under $\mathfrak{O}$ that is built in the proof of (e4): on $Y'$ it is the image; off
$Y'$ its first slot is that of the input, and its second slot is taken to be that of the input.
Then, under the hypotheses under which clause (e4) is asserted,
\begin{equation}\label{9.5:eq-relaxed-heredity-c}
\tilde q\in\mathfrak{G}^{\natural}(\mathfrak{B}^{t};Y').
\end{equation}
\end{lemma}

\begin{proof}
Throughout, clauses (e2) and (e3) are applied to the dated sets $\mathfrak{B}_{d''}$ and
$\mathfrak{B}_d$, $d''\le d$, which the order lemma makes comparable. In the form used below
they read, at every point $y$ and for every entry index $a$,
\begin{equation}\label{9.5:eq-relaxed-heredity-1}
\theta^{d}_{a}(y)\le\theta^{d''}_{a}(y),\qquad
\mathfrak{G}(\mathfrak{B}_d)\subset\mathfrak{G}(\mathfrak{B}_{d''})\subset
\mathfrak{K}_{109}\bigl(\mathfrak{B}_{d''};\gamma_0^{(k'')}\bigr).
\end{equation}
The first two clauses of $\mathfrak{K}_{109}$ are \cite[(3.35)]{B-CMP99b} and
\cite[(3.37)]{B-CMP99b}, formed at the radii $\alpha_0^{(99)}$, $\alpha_1^{(99)}$; neither of them
involves the second slot of a pair.

\emph{(W1)(i).} Clause \eqref{9.5:eq-relaxed-class-a}: let $d''\le d$. By
\eqref{9.5:eq-relaxed-heredity-1}, $P\in\mathfrak{G}(\mathfrak{B}_d)$ gives
$P\in\mathfrak{K}_{109}(\mathfrak{B}_{d''};\gamma_0^{(k'')})$, so on $\Omega_0(\mathfrak{B}_{d''})$
the first slot of $P$ obeys \cite[(3.35), (3.37)]{B-CMP99b} over $\mathfrak{B}_{d''}$. These
two clauses do not involve the second slot, and $P^{Y}$ has the first slot of $P$; hence $P^{Y}$
obeys them as well. Clause \eqref{9.5:eq-relaxed-class-b}: for $a\neq3$ the pointwise clauses at
$y\in Y^{+}_{d}$ are part of the membership $P\in\mathfrak{G}(\mathfrak{B}_d)$, and they involve
only the first slot, which is unchanged. For $a=3$, multiplying the second slot by
$\mathbf{1}_{Y^{+}_{d}}$ can only lower $Q_3$, so the clause for $P$ gives the clause for $P^{Y}$.
Clause \eqref{9.5:eq-relaxed-class-c} holds because the second slot of $P^{Y}$ vanishes off
$Y^{+}_{d}$ by construction, and $P^{Y}=P$ on the bonds of $Y^{+}_{d}$ for the same reason.

\emph{(W1)(ii).} The clauses \cite[(3.35), (3.37)]{B-CMP99b} over $\mathfrak{B}_{d''}$ are
clause \eqref{9.5:eq-relaxed-class-a} at the date $d''$. It remains to check the clause of
$\mathfrak{K}_{109}(\mathfrak{B}_{d''};\gamma_0^{(k'')})$ on the second slot, the bound of
$\ell_{d''}^{3}|\tilde{\mathbb{J}}|$ by $\gamma_0^{(k'')}$. Off $Y^{+}_{d}$ we have
$\tilde{\mathbb{J}}=0$ by \eqref{9.5:eq-relaxed-class-c}. At a bond at a point
$y\in Y^{+}_{d}\cap\Omega_1(\mathfrak{B}_{d''})$, clause \eqref{9.5:eq-relaxed-class-b} with
$a=3$, whose entry bounds $|\tilde{\mathbb{J}}|$ at the bonds at $y$, and the first inequality of
\eqref{9.5:eq-relaxed-heredity-1} give
\begin{equation*}
\ell_{d''}(y)^{3}|\tilde{\mathbb{J}}|
<\ell_{d''}(y)^{3}\cdot5\theta^{d}_{3}(y)
\le\ell_{d''}(y)^{3}\cdot5\theta^{d''}_{3}(y)
=5w_{d''}(y)\,\alpha_{0,\lambda_{d''}(y)} .
\end{equation*}
Finally, for the radius indices $i\in\{0,1\}$ used below,
\begin{equation}\label{9.5:eq-relaxed-heredity-2}
5w_{d''}(y)\,\alpha_{i,\lambda_{d''}(y)}\le\gamma_0^{(k'')}
\end{equation}
is the pointwise arithmetic carried out in the proof of the second inclusion of (e3), read at
the set $\mathfrak{B}_{d''}$ of density $k''$: if $w_{d''}(y)=1$, then
$5\alpha_{i,\lambda_{d''}(y)}\le5\bar\alpha_{k''}$, which is bounded by the weight allowed at
density $k''$ times $5\bar\alpha_{k''}$; if $w_{d''}(y)>1$, the structure at $y$ was written at a
step $k_c\le k''$, the weight at $y$ is at most the weight allowed at step $k_c$, and by clause (a)
the level at $y$ does not exceed $k_c$, so that the left side is bounded by that weight times
$5\bar\alpha_{k_c}$. In both cases the bound is at most $\gamma_0^{(k'')}$. This is a statement at
the single point $y$, and it is therefore available here, where only the pointwise clause
\eqref{9.5:eq-relaxed-class-b} is known.

If the collar is non-empty, at a bond at a point $y$ of $\operatorname{Col}(d'')$ the entry is
controlled by the dated clause $\mathsf{q}_a(\tilde p)(y)<\mathsf{f}(d)\mathsf{a}_a$ of
\eqref{9.5:eq-relaxed-class-b}. By clause (i) of the collar lemma the thresholds do not rise with
the date, so $\mathsf{f}(d)\le\mathsf{f}(d'')$ and $\mathsf{q}_a(\tilde p)(y)<\mathsf{f}(d'')
\mathsf{a}_a$; the inequality of that lemma at the date $d''$,
\begin{equation*}
\mathsf{f}(d'')\mathsf{a}_a<5\bar\alpha_{k''}\le\gamma_0^{(k'')},
\end{equation*}
gives the bound by $\gamma_0^{(k'')}$. Since $d''\le d$ was arbitrary, the membership holds at
every older date and at $d$ itself.

\emph{(W1)(iii).} Each of the four quantities whose supremum defines $a_{\flat}$ is, at a point
$z\in X'\cup B\subset Y^{+}_{d}$, a power $\ell_{d''}(z)^{e_a}$ of the unit times one of the
entries with indices $5$, $1$, $6$, $7$, divided by a power of $\eta$. By
\eqref{9.5:eq-relaxed-class-b} and the first inequality of \eqref{9.5:eq-relaxed-heredity-1},
each of them is below $5w_{d''}(z)\,\alpha_{\cdot,\lambda_{d''}(z)}$, with the radius attached to
its entry, exactly as the entry with index $3$ was treated in (ii). The arithmetic
\eqref{9.5:eq-relaxed-heredity-2} bounds each of these by $\gamma_0^{(k'')}$, and the supremum over
$z$ gives $a_{\flat}(\tilde p)\le\gamma_0^{(k'')}$.

\emph{(W2).} By \eqref{9.5:eq-fattened-sets-b}, $S^{+}_{d''}\subset Y^{+}_{d}$. Clause
\eqref{9.5:eq-relaxed-class-a} for $\tilde p^{S}$ at a date $d'''\le d''$: the dates $\le d''$ are
among the dates $\le d$, and the domain $\Omega_0$ of the dated sets never shrinks as the date
increases (a fact on dated sets established with the order lemma and the nesting statement), so
$\Omega_0(\mathfrak{B}_{d'''})\subset\Omega_0(\mathfrak{B}_{d''})$; on it the first slot of
$\tilde p^{S}$ is that of $\tilde p$, and \eqref{9.5:eq-relaxed-class-a} for $\tilde p$ at $d'''$
applies. Clause \eqref{9.5:eq-relaxed-class-b} at $y\in S^{+}_{d''}$: the entries $Q_a(\cdot)(y)$
are quantities at $y$ that do not depend on the dated set, and $y\in Y^{+}_{d}$, so
\eqref{9.5:eq-relaxed-class-b} for $\tilde p$ and the first inequality of
\eqref{9.5:eq-relaxed-heredity-1} give, for $a\neq3$,
\begin{equation*}
Q_a(\tilde p^{S})(y)=Q_a(\tilde p)(y)<5\theta^{d}_{a}(y)\le5\theta^{d''}_{a}(y);
\end{equation*}
for $a=3$ the cut $\mathbf{1}_{S^{+}_{d''}}$ only lowers $Q_3$. Clause
\eqref{9.5:eq-relaxed-class-c} is the cut, and the cut also gives $\tilde p^{S}=\tilde p$ on the
bonds of $S^{+}_{d''}$.

\emph{(W3).} Write $d^{t}$ for the date of $\mathfrak{B}^{t}$; since
$\mathfrak{B}^{t}\preccurlyeq\mathfrak{B}'=\mathfrak{B}_d$, the order lemma gives $d^{t}\le d$,
and \eqref{9.5:eq-relaxed-heredity-1} applies with $d''=d^{t}$.

Clause \eqref{9.5:eq-relaxed-class-b} at $y\in Y'$. The proof of (e4) bounds at such $y$ each
direct entry of $p$ by
\begin{equation*}
1.4\,\theta^{\mathfrak{B}'}_{a}(y)\le1.4\,\theta^{\mathfrak{B}^{t}}_{a}(y)
\end{equation*}
(by (e2)), and bounds the amount by which $\mathfrak{O}$ moves it by
$0.1\,\theta^{\mathfrak{B}^{t}}_{a}(y)$; at the site at which that proof works with a presented
pair, the entries are those of the presented pair. The entries of the image at $y$ are thus
below $1.5\,\theta^{\mathfrak{B}^{t}}_{a}(y)<5\theta^{\mathfrak{B}^{t}}_{a}(y)$. These estimates use
only the clauses of $p$ at $y$ and the allowances spent in the estimates on which the proof of (e4)
rests; what those estimates ask of the continuation is membership in
$\mathfrak{K}_{109}(\mathfrak{B}^{t};\gamma_0^{(k(d^{t}))})$, the hypothesis of the
background-response theorem. This membership is (W1)(ii) with $d'':=d^{t}$.

Clause \eqref{9.5:eq-relaxed-class-b} at a point $y\in Y'^{+}_{d^{t}}\setminus Y'$, which lies in
$Y'^{+}_{d}$ by \eqref{9.5:eq-fattened-sets-b}, as in (W2). There
$\tilde q$ equals $\tilde p$ in the first slot, and in the second slot $\tilde q$ is the input's,
so its entry with index $3$ is at most that of $\tilde p$. Hence
\eqref{9.5:eq-relaxed-class-b} for $\tilde p$ and the first inequality of
\eqref{9.5:eq-relaxed-heredity-1} give the clause for $\tilde q$ at the thresholds of
$\mathfrak{B}^{t}$.

Clause \eqref{9.5:eq-relaxed-class-a} at a date $d''\le d^{t}$. Where the first slot of
$\tilde q$ equals that of $\tilde p$, the clause is that of $\tilde p$, which holds at $d''$ since
$d''\le d$. On $Y'$ we treat the two parts of the first slot. For the $\tilde A$-part: by
\eqref{9.5:eq-relaxed-class-b} for $\tilde q$, just shown, and the first inequality of
\eqref{9.5:eq-relaxed-heredity-1} between $\mathfrak{B}_{d''}$ and $\mathfrak{B}^{t}$,
\begin{equation*}
\ell_{d''}|\tilde A|,\ \ \ell_{d''}^{2}|\nabla^{\eta}_{\tilde U}\tilde A|
<5w_{d''}\,\alpha_{1,\lambda_{d''}}\le\gamma_0^{(k'')},
\end{equation*}
the last inequality by \eqref{9.5:eq-relaxed-heredity-2}; entries with indices $6$ and $7$
bounded in this way give \cite[(3.37)]{B-CMP99b} over $\mathfrak{B}_{d''}$ at $y$, which is the
implication the proof of (e4) uses at the same place. For the $G$-valued part: the map
$\mathfrak{O}$ is a layer into the $\tilde A$-part, moving neither the $G$-valued part $\tilde U$
nor the cube gauges (as shown in the proof of (e4)), so \cite[(3.35)]{B-CMP99b} for $\tilde q$ is
that of $\tilde p$; at the site of the presented pair, on the region and the core treated there,
it is given by the cube clauses established in the proof of (e4), through the implication that
entries with indices $5$ and $8$ give \cite[(3.35)]{B-CMP99b}.

Clause \eqref{9.5:eq-relaxed-class-c}: off $Y'$ the second slot of $\tilde q$ is that of the input
$\tilde p$, which vanishes off $Y'^{+}_{d}$ by \eqref{9.5:eq-relaxed-class-c} for $\tilde p$.

It remains to consider a stencil astride $\partial Y'$. There clause (e4) rests on the statement
that off $Y'$ the image of the map is its input, which is part of the construction of the maps,
together with clauses (a) and (f) of the background-response theorem, under the hypotheses of
clause (e4); a bond astride $\partial Y'$ is treated in that statement in the same way as here.
That statement is used by statement, under the hypotheses of clause (e4); part (W3) uses it once,
for the same $\tilde q$, and adds nothing to it. This proves \eqref{9.5:eq-relaxed-heredity-c}.
\end{proof}

\medskip
\noindent\emph{Which continuation.} The requirement of a continuation is existential: one
continuation of the image suffices. Part (W3) takes the continuation that the proof of (e4)
builds for the levels $\ge1$, with the second slot equal to the input's off $Y'$, hence vanishing
off $Y'^{+}_{d}$. The uncut global image at the
operation $\mathbf{T}$ is itself a continuation at level $0$; it is an alternative used for the
level-$0$ collar clauses only, and under the scope condition \eqref{9.5:eq-fattened-sets-a} it is
not used. The cut of the second slot is made off $Y'^{+}_{d}$; it is the cut of the second
slot used in the proof of (e4) for the levels $\ge1$.

\medskip
\noindent\emph{No room is needed beyond $Y'^{+}_{d}$.} A clause with a single radius, the
same for source and target, as the radius $\alpha_1$ of \cite[(3.37)]{B-CMP99b} is, leaves no
room at points which the maps move. Beyond $Y'^{+}_{d}$ the continuation $\tilde q$ does
not move, and there equal radii suffice.

\begin{lemma}[The relaxed domains are open and gauge-invariant]\label{9.5:lem-relaxed-domains-open}
For every $n$ and every union of cubes $Y$, the set $\mathfrak{D}^{\flat,\natural}_n(Y)$ defined with the
ambient clause \eqref{9.5:eq-relaxed-class-d} contains $\mathfrak{D}^{\flat}_n(Y)$. The set
$\mathfrak{D}^{\flat,\natural}_n(Y)$ is open in the
space of pairs on $Y$, invariant under the gauge group $\mathcal{G}$, and decreasing in $n$, that is,
$\mathfrak{D}^{\flat,\natural}_n(Y)\subset\mathfrak{D}^{\flat,\natural}_{n-1}(Y)$. It is stable under
restriction: for every union of cubes $S\subset Y$, the restriction to $S$ of a pair of
$\mathfrak{D}^{\flat,\natural}_n(Y)$ lies in $\mathfrak{D}^{\flat,\natural}_n(S)$.
\end{lemma}

\begin{proof}
We write $d=(k;n)$ for the date of density $k$ at which the domains of stage $n$ are formed; the second
entry is the stage index, so that $(k;n-1)$ is the date of the same density at stage $n-1$.

\emph{Containment.} By the definition of the inherited domains, a pair of $\mathfrak{D}^{\flat}_n(Y)$
is the restriction to $Y$ of an ambient pair $P=(\mathbb{U},\mathbb{J})\in\mathfrak{G}(\mathfrak{B}_d)$
and satisfies the local clauses, which depend only on the pair on $Y$. Let $P$ be such a witness. By the
sandwich clause
\eqref{9.5:eq-relaxed-heredity-a}, part (i), of Lemma~\ref{9.5:lem-relaxed-heredity}, the pair
$P^{Y}=(\mathbb{U},\mathbf{1}_{Y^{+}_{d}}\mathbb{J})$ belongs to $\mathfrak{G}^{\natural}(\mathfrak{B}_d;Y)$,
and $P^{Y}=P$ on the bonds of $Y^{+}_{d}$. Since $Y\subset Y^{+}_{d}$, the restrictions of $P^{Y}$ and
of $P$ to $Y$ coincide. The pair on $Y$ therefore satisfies the ambient clause
\eqref{9.5:eq-relaxed-class-d} with the witness $P^{Y}$, and it satisfies the local clauses, which depend
only on its restriction to $Y$. Hence $\mathfrak{D}^{\flat}_n(Y)\subset\mathfrak{D}^{\flat,\natural}_n(Y)$.

\emph{Openness.} Let $\mathcal{E}_d$ denote the set of bonds of $\Omega_0(\mathfrak{B}_d)$. Let $\Xi_{\natural}$
denote the linear subspace of $(\mathfrak{g}^{c})^{\mathcal{E}_d}$ consisting of the currents that obey
\eqref{9.5:eq-relaxed-class-c}. Consider the set $\mathcal{W}$ of triples
$(\tilde U,\tilde A,\tilde{\mathbb{J}})$ in
\[
G^{\mathcal{E}_d}\times(\mathfrak{g}^{c})^{\mathcal{E}_d}\times \Xi_{\natural}
\]
which obey \eqref{9.5:eq-relaxed-class-a}, \eqref{9.5:eq-relaxed-class-b} and the strict local clauses.

The set $\mathcal{W}$ is open. Each of these conditions is a finite family of strict inequalities between
continuous functions of finitely many lattice variables, since the lattice is finite. Two of the
conditions are not of this form as they stand: the condition \cite[(3.35)]{B-CMP99b} and one further
clause among the conditions defining $\mathcal{W}$ are each formulated through a gauge transformation $u$.
Each of them is the union, over the gauge transformations $u$, of sets of the kind just described, hence
open, since a union of open sets is open.

Next consider the map
\[
\mathcal{W}\ni(\tilde U,\tilde A,\tilde{\mathbb{J}})\longmapsto
\bigl(e^{i\eta\tilde A}\tilde U,\tilde{\mathbb{J}}\bigr)\restriction Y .
\]
This map is open. On each bond, the map $(U,A)\mapsto e^{i\eta A}U$ from $G\times\mathfrak{g}^{c}$ to
$G^{c}$ is a submersion where $\eta|A|<1$. Already for fixed $U$, the map $A\mapsto e^{i\eta A}U$ is a
local diffeomorphism of $\mathfrak{g}^{c}$ onto an open set of $G^{c}$. The restriction to $Y$ is a
coordinate projection. In the second slot the current is free on the bonds of $Y$: by
\eqref{9.5:eq-relaxed-class-c}, elements of $\Xi_{\natural}$ are unconstrained on the bonds of
$Y\subset Y^{+}_{d}$.

The image of $\mathcal{W}$ under this map is $\mathfrak{D}^{\flat,\natural}_n(Y)$. This set is therefore
open in the space of pairs on $Y$.

\emph{Gauge invariance.} A gauge transformation $v$ acts on triples bondwise: on a bond
$\langle x,y\rangle$ it sends $\tilde U$ to $v(x)\tilde U v(y)^{-1}$ and $\tilde A$ to
$\operatorname{Ad}_{v(x)}\tilde A$, with the endpoint convention of \cite[(3.31)--(3.34)]{B-CMP99b}, and
it acts on the current slot $\tilde{\mathbb{J}}$ as on the current slot of pairs, bondwise and linearly.
The functionals $Q_a$ and $\mathsf{q}_a$ entering \eqref{9.5:eq-relaxed-class-b} are those of the
pointwise clauses of the continuation class $\mathfrak{G}(\mathfrak{B}_d)$ of the arrest theorem, and
they are $\mathcal{G}$-invariant, so \eqref{9.5:eq-relaxed-class-b} is invariant whatever the
thresholds. For the other conditions:
\begin{itemize}
\item The condition \cite[(3.35)]{B-CMP99b} is invariant because its gauge $u$ may be replaced by $uv^{-1}$.
\item The condition \cite[(3.37)]{B-CMP99b} is invariant for two reasons: the norm satisfies
$|\operatorname{Ad}_v\cdot|=|\cdot|$, and the covariant difference $\nabla_{U}$ is covariant under gauge
transformations \cite[(3.31)--(3.34)]{B-CMP99b}.
\item The condition \eqref{9.5:eq-relaxed-class-c} is invariant as well, since the action on the current
slot is bondwise and linear and therefore preserves the vanishing of $\tilde{\mathbb{J}}$ on a bond.
\end{itemize}
Hence \eqref{9.5:eq-relaxed-class-a} is invariant at every date $d''\le d$, and so are
\eqref{9.5:eq-relaxed-class-b} and \eqref{9.5:eq-relaxed-class-c}. The local clauses are the local
clause families of the inherited domains $\mathfrak{D}^{\flat}_n$, and each of these families is
$\mathcal{G}$-invariant by the definition of the inherited domains. Thus every condition
defining $\mathcal{W}$ is $\mathcal{G}$-invariant. The map of the openness step commutes with gauge
transformations, since on a bond $\langle x,y\rangle$
\[
v(x)\,e^{i\eta A}U\,v(y)^{-1}=e^{i\eta\operatorname{Ad}_{v(x)}A}\,v(x)Uv(y)^{-1}.
\]
Therefore its image $\mathfrak{D}^{\flat,\natural}_n(Y)$ is $\mathcal{G}$-invariant.

\emph{Monotonicity in $n$.} The ambient clause \eqref{9.5:eq-relaxed-class-d} is decreasing in $n$. This
follows from clause \eqref{9.5:eq-relaxed-heredity-b} of
Lemma~\ref{9.5:lem-relaxed-heredity}, applied with
\[
(d'',d,S)=\bigl((k;n-1),(k;n),Y\bigr).
\]
The local clauses are decreasing in $n$ by the monotonicity and restriction property of the local clause
families of the inherited domains $\mathfrak{D}^{\flat}_n$.

\emph{Stability under restriction.} For the ambient clause this is again
\eqref{9.5:eq-relaxed-heredity-b} of Lemma~\ref{9.5:lem-relaxed-heredity}. For the local clauses it is
the same restriction property of the local clause families of $\mathfrak{D}^{\flat}_n$.
\end{proof}

Let $\eta$ be the lattice spacing of Notation~\ref{9.5:not-dates-weights-norms}. Write $e_\mu$ for
the unit vector in the direction $\mu$,
$b=\langle b_-,b_+\rangle$ for a bond with lower endpoint $b_-$ and upper endpoint $b_+$, and
$b+\eta e_\mu:=\langle b_-+\eta e_\mu,\,b_++\eta e_\mu\rangle$ for its translate by $\eta e_\mu$. The
bond $b$ runs over the positively oriented bonds; a reversed bond carries the inverse of the bond
variable. Accordingly no symmetry of the forms below under $b\mapsto -b$ is required.

Let $U$ be a configuration with values in $\mathrm{SU}(N)$, and let $A$ be a function on bonds with
values in the complexification $\mathfrak{g}^{c}$ of the Lie algebra of $\mathrm{SU}(N)$. The
covariant lattice difference is
\begin{equation*}
(\nabla^{\eta}_{U,\mu}A)(b)
:=\eta^{-1}\bigl[\operatorname{Ad}_{U(\langle b_-,\,b_-+\eta e_\mu\rangle)}A(b+\eta e_\mu)-A(b)\bigr].
\end{equation*}
For a function $\chi$ on sites we put
\begin{equation*}
(\partial_\mu\chi)(x):=\eta^{-1}\bigl[\chi(x+\eta e_\mu)-\chi(x)\bigr],
\qquad
(\chi A)(b):=\chi(b_-)A(b).
\end{equation*}
Let $|\cdot|$ be the norm on $\mathfrak{g}^{c}$ of
Notation~\ref{9.5:not-dates-weights-norms}.

\begin{lemma}[Covariant difference of a cut-off form]\label{9.5:lem-cutoff-difference}
Let $U$, $A$ and $\chi$ be as just fixed, with $0\le\chi\le1$, and let $\nabla^{\eta}_{U,\mu}$,
$\partial_\mu\chi$, $\chi A$ and $|\cdot|$ be as defined above. Then
$|(\chi A)(b)|\le|A(b)|$. Moreover
\begin{equation}\label{9.5:eq-cutoff-difference-1}
\begin{split}
(\nabla^{\eta}_{U,\mu}(\chi A))(b)
&=\chi(b_-)\,(\nabla^{\eta}_{U,\mu}A)(b)\\
&\quad+(\partial_\mu\chi)(b_-)\,
\operatorname{Ad}_{U(\langle b_-,\,b_-+\eta e_\mu\rangle)}A(b+\eta e_\mu),
\end{split}
\end{equation}
and consequently
\begin{equation}\label{9.5:eq-cutoff-difference-2}
|(\nabla^{\eta}_{U,\mu}(\chi A))(b)|
\le|(\nabla^{\eta}_{U,\mu}A)(b)|+|(\partial_\mu\chi)(b_-)|\,|A(b+\eta e_\mu)|.
\end{equation}
\end{lemma}

\begin{proof}
Since $0\le\chi(b_-)\le1$, we have
\begin{equation*}
|(\chi A)(b)|=\chi(b_-)\,|A(b)|\le|A(b)|.
\end{equation*}

For the identity, write $\operatorname{Ad}_U$ for
$\operatorname{Ad}_{U(\langle b_-,b_-+\eta e_\mu\rangle)}$ in this proof. The bond $b+\eta e_\mu$ is
again positively oriented, and its lower endpoint is
$(b+\eta e_\mu)_-=b_-+\eta e_\mu$. Hence, by the definition of $\chi A$,
\begin{equation*}
(\chi A)(b+\eta e_\mu)=\chi(b_-+\eta e_\mu)\,A(b+\eta e_\mu).
\end{equation*}
Inserting this into the definition of $\nabla^{\eta}_{U,\mu}$ and using the linearity of
$\operatorname{Ad}_U$ to take the scalar $\chi(b_-+\eta e_\mu)$ outside gives
\begin{equation*}
\eta\,(\nabla^{\eta}_{U,\mu}(\chi A))(b)
=\chi(b_-+\eta e_\mu)\operatorname{Ad}_UA(b+\eta e_\mu)-\chi(b_-)A(b).
\end{equation*}
Now add and subtract $\chi(b_-)\operatorname{Ad}_UA(b+\eta e_\mu)$. The right-hand side becomes
\begin{equation*}
\chi(b_-)\bigl[\operatorname{Ad}_UA(b+\eta e_\mu)-A(b)\bigr]
+\bigl[\chi(b_-+\eta e_\mu)-\chi(b_-)\bigr]\operatorname{Ad}_UA(b+\eta e_\mu).
\end{equation*}
Multiply by $\eta^{-1}$. The first bracket becomes $(\nabla^{\eta}_{U,\mu}A)(b)$, and the second
becomes $(\partial_\mu\chi)(b_-)$. This is \eqref{9.5:eq-cutoff-difference-1}.

For the bound, apply the triangle inequality to \eqref{9.5:eq-cutoff-difference-1} and use
$0\le\chi(b_-)\le1$ in the first term. For the second term, use that $\operatorname{Ad}$ by an
element of $\mathrm{SU}(N)$ is an isometry of $(\mathfrak{g}^{c},|\cdot|)$. This gives
\begin{equation*}
|(\nabla^{\eta}_{U,\mu}(\chi A))(b)|
\le\chi(b_-)\,|(\nabla^{\eta}_{U,\mu}A)(b)|
+|(\partial_\mu\chi)(b_-)|\,|A(b+\eta e_\mu)|,
\end{equation*}
and \eqref{9.5:eq-cutoff-difference-2} follows from $\chi(b_-)\le1$.
\end{proof}

If in the definition of $\chi A$ one replaces $\chi(b_-)$ by
$\tfrac12\bigl(\chi(b_-)+\chi(b_+)\bigr)$, the constant bounding the differences of $\chi$ is to be
doubled; either convention serves, and the bounds of this section leave room for this factor $2$.

\begin{lemma}[Curvature of a perturbed plaquette]\label{9.5:lem-perturbed-plaquette}
Let $U$ be an $\mathrm{SU}(N)$-valued bond field and $B$ a field on bonds with values in the
complexification $\mathfrak{g}^{c}$ of the Lie algebra of $\mathrm{SU}(N)$, both on the lattice of
spacing $\eta$. Let $p$ be a plaquette with corners
$x$, $x+\eta e_\mu$, $x+\eta e_\mu+\eta e_\nu$, $x+\eta e_\nu$, and bonds
\begin{equation*}
b_1=\langle x,x+\eta e_\mu\rangle,\qquad b_2=b_4+\eta e_\mu,\qquad
b_3=b_1+\eta e_\nu,\qquad b_4=\langle x,x+\eta e_\nu\rangle .
\end{equation*}
Put $\mathbb{U}(b):=e^{i\eta B(b)}U(b)$ and
\begin{equation*}
\mathbb{U}(\partial p):=\mathbb{U}(b_1)\,\mathbb{U}(b_2)\,\mathbb{U}(b_3)^{-1}\,\mathbb{U}(b_4)^{-1},
\qquad P:=U(b_1)\,U(b_2)\,U(b_3)^{-1}\,U(b_4)^{-1},
\end{equation*}
so that $P$ is the plaquette variable $U(\partial p)$ in this ordering, and
\begin{equation*}
\beta_6:=\max_{1\le i\le 4}|B(b_i)|,\qquad
\beta_7:=\max\bigl\{|\nabla_{U,\mu}B(b_4)|,\ |\nabla_{U,\nu}B(b_1)|\bigr\}.
\end{equation*}
Then
\begin{equation}\label{9.5:eq-perturbed-plaquette-1}
|\mathbb{U}(\partial p)-\mathbf{1}|
\le |P-\mathbf{1}|\,(1+4\eta\beta_6)+2\eta^2\beta_7+8\eta^2\beta_6^2e^{4\eta\beta_6}.
\end{equation}
\end{lemma}

\begin{proof}
We write $\operatorname{Ad}_V Y:=VYV^{-1}$, and $U_i:=U(b_i)$, $B_i:=B(b_i)$ for $i=1,\dots,4$.
The covariant lattice difference $\nabla_{U,\mu}$ in the statement is, explicitly,
\begin{equation*}
(\nabla_{U,\mu}A)(b)
=\eta^{-1}\bigl[\operatorname{Ad}_{U(\langle b_-,b_-+\eta e_\mu\rangle)}A(b+\eta e_\mu)-A(b)\bigr],
\end{equation*}
where $b_-$ is the initial point of $b$.
Of the norms $|\cdot|$ and $\|\cdot\|$ of Notation~\ref{9.5:not-dates-weights-norms} we use that
$|YZ|\le|Y||Z|$ and $|YZ|\le|Y|\,\|Z\|$ for all matrices, and that $|VYW|=|Y|$ and
$\|V\|=1$ for $V,W\in\mathrm{SU}(N)$; in particular $|\operatorname{Ad}_VY|=|Y|$.

\emph{Moving the unitary factors to the right.} By definition,
\begin{equation*}
\mathbb{U}(\partial p)=e^{i\eta B_1}U_1\,e^{i\eta B_2}U_2\,U_3^{-1}e^{-i\eta B_3}\,U_4^{-1}e^{-i\eta B_4}.
\end{equation*}
For $V\in\mathrm{SU}(N)$ one has $Ve^{Y}=e^{\operatorname{Ad}_VY}V$, since
$VY^nV^{-1}=(\operatorname{Ad}_VY)^n$ for every $n$. Hence
$U_1e^{i\eta B_2}=e^{i\eta\operatorname{Ad}_{U_1}B_2}U_1$. Next
$U_1U_2U_3^{-1}=PU_4$ by the definition of $P$, and
$PU_4e^{-i\eta B_3}=e^{-i\eta\operatorname{Ad}_{PU_4}B_3}PU_4$ with
$\operatorname{Ad}_{PU_4}=\operatorname{Ad}_P\operatorname{Ad}_{U_4}$. Finally
$PU_4U_4^{-1}e^{-i\eta B_4}=Pe^{-i\eta B_4}=e^{-i\eta\operatorname{Ad}_PB_4}P$. Therefore
\begin{equation}\label{9.5:eq-perturbed-plaquette-2}
\mathbb{U}(\partial p)=e^{i\eta X_1}e^{i\eta X_2}e^{-i\eta X_3}e^{-i\eta X_4}\,P,
\end{equation}
with
\begin{equation*}
X_1:=B_1,\quad X_2:=\operatorname{Ad}_{U_1}B_2,\quad
X_3:=\operatorname{Ad}_P\operatorname{Ad}_{U_4}B_3,\quad X_4:=\operatorname{Ad}_PB_4 .
\end{equation*}

\emph{The linear term.} Writing $\operatorname{Ad}_PY=Y-(1-\operatorname{Ad}_P)Y$ in $X_3$ and $X_4$,
\begin{equation*}
\begin{split}
X_1+X_2-X_3-X_4
&=\bigl[\operatorname{Ad}_{U_1}B_2-B_4\bigr]-\bigl[\operatorname{Ad}_{U_4}B_3-B_1\bigr]
+(1-\operatorname{Ad}_P)\bigl[\operatorname{Ad}_{U_4}B_3+B_4\bigr]\\
&=\eta\nabla_{U,\mu}B(b_4)-\eta\nabla_{U,\nu}B(b_1)
+(1-\operatorname{Ad}_P)\bigl[\operatorname{Ad}_{U_4}B_3+B_4\bigr].
\end{split}
\end{equation*}
In the second line we used that $b_4$ has initial point $x$, that
$\langle x,x+\eta e_\mu\rangle=b_1$ and $b_4+\eta e_\mu=b_2$; and that $b_1$ has initial point $x$,
$\langle x,x+\eta e_\nu\rangle=b_4$ and $b_1+\eta e_\nu=b_3$.
For any matrix $Y$ one has
\begin{equation*}
-(1-\operatorname{Ad}_P)Y=PYP^{-1}-Y=(P-\mathbf{1})YP^{-1}+Y(P^{-1}-\mathbf{1}),
\end{equation*}
hence
\begin{equation*}
|(1-\operatorname{Ad}_P)Y|\le|P-\mathbf{1}|\,|Y|\,\|P^{-1}\|+|Y|\,|P^{-1}(\mathbf{1}-P)|
\le2|P-\mathbf{1}|\,|Y|.
\end{equation*}
Applied to $Y=\operatorname{Ad}_{U_4}B_3+B_4$, for which $|Y|\le|B_3|+|B_4|\le2\beta_6$, this gives
\begin{equation}\label{9.5:eq-perturbed-plaquette-3}
|X_1+X_2-X_3-X_4|\le2\eta\beta_7+4|P-\mathbf{1}|\beta_6 .
\end{equation}

\emph{The exponential product.} Put $Z_1:=i\eta X_1$, $Z_2:=i\eta X_2$, $Z_3:=-i\eta X_3$,
$Z_4:=-i\eta X_4$ and $s:=\sum_{i}|Z_i|$. Since $|X_i|=|B_i|$ by the invariance of $|\cdot|$
under $\operatorname{Ad}$, we have $s\le4\eta\beta_6$. Expanding the four exponential series,
$e^{Z_1}e^{Z_2}e^{Z_3}e^{Z_4}$ is the sum over $(n_1,\dots,n_4)$ of
$Z_1^{n_1}\cdots Z_4^{n_4}/(n_1!\cdots n_4!)$; the term of total degree $0$ is $\mathbf{1}$, those of
total degree $1$ sum to $\sum_iZ_i$, and by submultiplicativity the remaining ones have norms
summing to at most $\sum_{n\ge2}s^n/n!$. As $n!\ge2(n-2)!$ for $n\ge2$,
\begin{equation*}
\Bigl|\prod_{i=1}^4e^{Z_i}-\mathbf{1}-\sum_{i=1}^4Z_i\Bigr|
\le e^s-1-s\le\tfrac12s^2e^s\le8\eta^2\beta_6^2e^{4\eta\beta_6},
\end{equation*}
the last step because $s\mapsto\frac12s^2e^s$ is increasing and $s\le4\eta\beta_6$. Since
$\sum_iZ_i=i\eta(X_1+X_2-X_3-X_4)$, the bound \eqref{9.5:eq-perturbed-plaquette-3} yields
\begin{equation}\label{9.5:eq-perturbed-plaquette-4}
\Bigl|\prod_{i=1}^4e^{Z_i}-\mathbf{1}\Bigr|
\le\eta\bigl(2\eta\beta_7+4|P-\mathbf{1}|\beta_6\bigr)+8\eta^2\beta_6^2e^{4\eta\beta_6}.
\end{equation}

\emph{Conclusion.} By \eqref{9.5:eq-perturbed-plaquette-2},
$\mathbb{U}(\partial p)-\mathbf{1}=\bigl(\prod_ie^{Z_i}-\mathbf{1}\bigr)P+(P-\mathbf{1})$, so
\begin{equation*}
|\mathbb{U}(\partial p)-\mathbf{1}|\le\Bigl|\prod_{i=1}^4e^{Z_i}-\mathbf{1}\Bigr|\,\|P\|+|P-\mathbf{1}|,
\end{equation*}
with $\|P\|=1$ because $P\in\mathrm{SU}(N)$. Inserting \eqref{9.5:eq-perturbed-plaquette-4} and
collecting the terms in $|P-\mathbf{1}|$ gives \eqref{9.5:eq-perturbed-plaquette-1}.
\end{proof}

\begin{lemma}[Perturbing a real field inside a cube gauge]\label{9.5:lem-cube-gauge-perturbation}
Let $\square$ be a cube of the lattice of spacing $\eta$, and let $\ell\ge\eta$. Let $a\ge0$ and $e\ge0$ be numbers (the letter $e$ standing alone denotes this number; in $e^{Z}$ it denotes the exponential, and $e_\mu$ is the unit lattice vector in direction $\mu$) with
\begin{equation*}
a+e\le \tfrac{1}{30}.
\end{equation*}
Let $G=\mathrm{SU}(N)$, let $\mathfrak{g}$ be the real Lie algebra of Hermitian traceless $N\times N$ matrices, so that $e^{iY}\in G$ for $Y\in\mathfrak{g}$, and let $|\cdot|$ denote the operator norm.
Let $U$ be a $G$-valued field on the bonds of $\square$, and let $u$ be a $G$-valued gauge transformation on the sites of $\square$ with $U^{u}=e^{i\eta A}$, that is,
\begin{equation*}
u(b_-)\,U(b)\,u(b_+)^{-1}=e^{i\eta A(b)}\qquad\text{for every bond } b=\langle b_-,b_+\rangle \text{ of } \square,
\end{equation*}
where $\ell|A|\le a$ and $\ell^{2}|\nabla^{\eta}A|\le a$ on $\square$. Let $E$ be a real $\mathfrak{g}$-valued field on the bonds of $\square$ with $\ell|E|\le e$ and $\ell^{2}|\nabla_{U}E|\le e$ on $\square$, and let $\chi$ be a function on the sites of $\square$ with $0\le\chi\le1$ and $\ell|\partial\chi|\le1$. Put $(\chi E)(b):=\chi(b_-)E(b)$. Then
\begin{equation*}
U':=e^{i\eta\chi E}\,U
\end{equation*}
is $G$-valued, and $U'^{u}=e^{i\eta A''}$ with
\begin{equation}\label{9.5:eq-cube-gauge-perturbation-1}
\ell|A''|\le 1.1a+2.4e,\qquad \ell^{2}|\nabla^{\eta}A''|\le 1.1a+2.4e .
\end{equation}
\end{lemma}

\begin{proof}
Since $E$ is real, $e^{i\eta\chi(b_-)E(b)}\in G$ for every bond $b$, so $U'$ is $G$-valued. For a bond $b$ of $\square$ we have
\begin{align*}
U'^{u}(b)&=u(b_-)\,e^{i\eta(\chi E)(b)}\,u(b_-)^{-1}\cdot u(b_-)\,U(b)\,u(b_+)^{-1}
=e^{i\eta X(b)}\,e^{i\eta A(b)},
\end{align*}
where $X(b):=\operatorname{Ad}_{u(b_-)}(\chi E)(b)$. Since $\operatorname{Ad}_{u(b_-)}$ preserves the norm and $0\le\chi\le1$, $|X(b)|\le|E(b)|\le e/\ell$.

We use the following bound, the Lipschitz bound for the logarithm of a product proved together with the arrest theorem: the map $\Psi(Z_1,Z_2):=\log(e^{Z_1}e^{Z_2})$ satisfies
\begin{equation*}
\Psi(0,0)=0,\qquad |\Psi(Z_1,Z_2)-\Psi(Z_1',Z_2')|\le1.1\bigl(|Z_1-Z_1'|+|Z_2-Z_2'|\bigr)
\end{equation*}
for any two points $(Z_1,Z_2)$, $(Z_1',Z_2')$ of the convex set $\{|Z_1|+|Z_2|\le 0.04\}$. At every bond,
\begin{equation*}
|i\eta X(b)|+|i\eta A(b)|\le \frac{\eta}{\ell}(a+e)\le a+e\le\frac{1}{30}<0.04,
\end{equation*}
so we may define $A''$ by $i\eta A''(b):=\Psi(i\eta X(b),i\eta A(b))$; then $U'^{u}=e^{i\eta A''}$. Comparing with $\Psi(0,0)=0$ gives $\eta|A''(b)|\le1.1\,\eta(|X(b)|+|A(b)|)$, hence
\begin{equation*}
\ell|A''|\le1.1(a+e)\le1.1a+2.4e,
\end{equation*}
which is the first bound in \eqref{9.5:eq-cube-gauge-perturbation-1}.

For the second bound let $b$ and $b':=b+\eta e_\mu$ be bonds of $\square$. Both argument pairs $(i\eta X(b),i\eta A(b))$ and $(i\eta X(b'),i\eta A(b'))$ lie in the convex set above, so the Lipschitz bound applies to them; dividing by $\eta^{2}$ gives, for the plain differences,
\begin{equation}\label{9.5:eq-cube-gauge-perturbation-2}
|\nabla^{\eta}_{\mu}A''(b)|\le1.1\bigl(|\nabla^{\eta}_{\mu}X(b)|+|\nabla^{\eta}_{\mu}A(b)|\bigr).
\end{equation}
It remains to bound $\nabla^{\eta}_{\mu}X$. Write $x:=b_-$ and $x':=x+\eta e_\mu$, so that $b'_-=x'$. The gauge relation on the bond $\langle x,x'\rangle$ reads $u(x)U(x,x')u(x')^{-1}=e^{i\eta A_\mu(x)}$, that is,
\begin{equation*}
u(x')=e^{-i\eta A_\mu(x)}\,u(x)\,U(x,x').
\end{equation*}
Hence
\begin{equation*}
X(b')=\operatorname{Ad}_{e^{-i\eta A_\mu(x)}}\operatorname{Ad}_{u(x)}\operatorname{Ad}_{U(x,x')}(\chi E)(b').
\end{equation*}
The covariant difference satisfies $\eta\nabla_{U,\mu}(\chi E)(b)=\operatorname{Ad}_{U(x,x')}(\chi E)(b')-(\chi E)(b)$; adding and subtracting $\operatorname{Ad}_{e^{-i\eta A_\mu(x)}}\operatorname{Ad}_{u(x)}(\chi E)(b)$ we obtain
\begin{equation*}
X(b')-X(b)=\operatorname{Ad}_{e^{-i\eta A_\mu(x)}}\operatorname{Ad}_{u(x)}
\bigl[\eta\nabla_{U,\mu}(\chi E)(b)\bigr]
+\bigl(\operatorname{Ad}_{e^{-i\eta A_\mu(x)}}-1\bigr)X(b).
\end{equation*}
Since $A$ is real, $e^{-i\eta A_\mu(x)}\in G$, so both adjoint actions in the first term preserve the norm. By Lemma~\ref{9.5:lem-cutoff-difference},
\begin{equation*}
\ell^{2}|\nabla_{U,\mu}(\chi E)(b)|\le\ell^{2}|\nabla_{U,\mu}E(b)|
+\ell\,\bigl(\ell|\partial_\mu\chi(x)|\bigr)\,|E(b')|\le e+e=2e .
\end{equation*}
For the second term let $v\in G$ and $Y\in\mathfrak{g}$. The identity
\begin{equation*}
\operatorname{Ad}_vY-Y=(v-1)Yv^{-1}-Yv^{-1}(v-1)
\end{equation*}
gives $|(\operatorname{Ad}_v-1)Y|\le2|v-1|\,|Y|$. Now take $v=e^{-i\eta A_\mu(x)}$. With
\begin{gather*}
s:=\eta|A_\mu(x)|\le a\le\tfrac{1}{30},\\
|v-1|\le e^{s}-1\le s\,e^{s}\le 1.05\,s,
\end{gather*}
where the first inequality of the second line follows from the exponential series, we obtain
\begin{equation*}
|(\operatorname{Ad}_v-1)X(b)|\le2\bigl(e^{\eta|A_\mu(x)|}-1\bigr)|X(b)|
\le2.1\,\eta|A_\mu(x)|\,|X(b)|\le2.1\,\eta\,\frac{a}{\ell}\cdot\frac{e}{\ell}.
\end{equation*}
Dividing $X(b')-X(b)$ by $\eta$ and collecting the two terms,
\begin{equation*}
\ell^{2}|\nabla^{\eta}_{\mu}X(b)|\le2e+2.1\,ae\le2.1\,e,
\end{equation*}
the last step because $2.1\,a\le 2.1/30\le0.1$. Inserting this and $\ell^{2}|\nabla^{\eta}_{\mu}A(b)|\le a$ into \eqref{9.5:eq-cube-gauge-perturbation-2} gives
\begin{equation*}
\ell^{2}|\nabla^{\eta}_{\mu}A''(b)|\le1.1(a+2.1e)\le1.1a+2.4e,
\end{equation*}
which is the second bound in \eqref{9.5:eq-cube-gauge-perturbation-1}.
\end{proof}

\begin{lemma}[Geometry of the hull]\label{9.5:lem-hull-geometry}
Work in the scope of Definition~\ref{9.5:def-fattened-sets}: one run $\mathsf{m}$, one label, and a
term $F$ with its region $X$ and level $j$, born at the date $\tau$, with the notation of
Notation~\ref{9.5:not-dates-weights-norms}. In particular $u:=M\ell_j$ is the side of one
$\pi_j$-cube, and $X^{(s)}$ is the $s\cdot u$ sup-distance fattening of $X$ of
\eqref{9.5:eq-fattened-sets-c}. Write $X^{+}:=X^{+}_{\tau}$ for the ball collar of $X$ at the date
$\tau$ and $\widehat{X}$ for the hull, that is, $X$ together with all $\pi_j$-cubes that are
$\infty$-adjacent to it. Write
$\mathfrak{S}=\boldsymbol{\Omega}(\widehat{X})$ for the box tower over the hull, formed over the
label's determining set current at the read and truncated at $j$, and
$\lambda_{\mathfrak{S}}$ for its levels. Let $\varrho_0$ be the radius constant of the ball collar of
Definition~\ref{9.5:def-fattened-sets}. Assume:
\begin{itemize}
\item[(1)] $M/M_1\ge 3L(\varrho_0+2L)$;
\item[(2)] $30LM_1<M/256$;
\item[(3)] $L\ge 8$;
\item[(4)] the comparison of the levels of dated sets at different dates furnished by the arrest
theorem. It is used below in the form $\lambda_{d''}\le\lambda_\tau\le\lambda_d$ for dates
$d''\le\tau\le d$;
\item[(5)] the stratum geometry of the dated sets that accompanies the arrest theorem: the strata of
$\mathfrak{B}_\tau$, that is the sets $\{\lambda_\tau=l\}$, carry consecutive levels across adjacent
strata, and each stratum is at least one $M$-cube of its own scale, $M\ell_l$, wide.
\end{itemize}
Let $d$ be the date of the read, with $d\ge\tau$, and put $\iota:=\lambda_d$. Then the following
hold.
\begin{itemize}
\item[(a)] The collar lies in a thin fattening of $X$, and so does its $L^2\ell_j$-neighbourhood:
\begin{equation}\label{9.5:eq-hull-geometry-a}
X^{+}\subset X^{(1/24)},\qquad
\{y:\operatorname{dist}_\infty(y,X^{+})\le L^2\ell_j\}\subset X^{(1/4)}.
\end{equation}
\item[(b)] Define $\chi$ by the first line of \eqref{9.5:eq-hull-geometry-b}; then $\chi$ satisfies
its second line:
\begin{equation}\label{9.5:eq-hull-geometry-b}
\begin{gathered}
\chi(y):=\min\Bigl\{1,\max\Bigl\{0,\,2-\frac{4\operatorname{dist}_\infty(y,X)}{u}\Bigr\}\Bigr\},\\
\chi=1\ \text{on}\ X^{(1/4)},\qquad \chi=0\ \text{off}\ X^{(1/2)},
\qquad |\partial_\mu\chi|\le\frac{4}{u}.
\end{gathered}
\end{equation}
\item[(c)] Let $y$ be a cell of the determining set of $\mathfrak{S}$, and write $\tilde\Delta(y)$
for $y$ together with its neighbouring cells. If $\tilde\Delta(y)\cap X^{(3/4)}\neq\emptyset$,
then $y$ is live, that is,
\begin{equation}\label{9.5:eq-hull-geometry-c}
\lambda_{\mathfrak{S}}(y)=j\wedge\iota(y).
\end{equation}
\item[(d)] On $X^{(3/4)}$ the following holds for every date $d''\le\tau$:
\begin{equation}\label{9.5:eq-hull-geometry-d}
\lambda_{d''}\le\lambda_\tau\le\lambda_{\mathfrak{S}}+1.
\end{equation}
Moreover $\lambda_\tau\le\lambda_{\mathfrak{S}}$ there if $j=k(\tau)$. This applies when
$F=\mathbb{C}^{(i)}_{k_a}(X)$: then $\tau=(k_a;i)$, $j=k_a$ and $k(\tau)=k_a$.
\end{itemize}
\end{lemma}

\begin{proof}
Throughout, $\ell_j=L^j\eta$ and $\ell_\tau(x)=L^{\lambda_\tau(x)}\eta$.

\emph{Proof of \eqref{9.5:eq-hull-geometry-a}.} Let $n_F$ be the level of the units of $F$. The
units of $F$ are those of the stored term, read at its birth set $\mathfrak{B}_\tau$. Hence on $X$
the level of $\mathfrak{B}_\tau$ is the level of these units, and by the definition of $n_F$
\begin{equation*}
\lambda_\tau=n_F\le j\quad\text{on }X.
\end{equation*}
Let $x\in X$. By Definition~\ref{9.5:def-fattened-sets}, $X^{+}_\tau$ is the union over $x\in X$ of
sup-balls of radius $(\varrho_0+1)\ell_\tau(x)$ about $x$.
Since $\lambda_\tau(x)\le j$ we have $\ell_\tau(x)\le\ell_j$. Also $\varrho_0+1\le\varrho_0+2L$, and
hypothesis~(1) gives $\varrho_0+2L\le M/(3LM_1)$. Hypothesis~(3) and $M_1\ge1$ give $3LM_1\ge 24$.
Therefore
\begin{equation*}
(\varrho_0+1)\ell_\tau(x)\le(\varrho_0+2L)\ell_j\le\frac{M}{3LM_1}\,\ell_j\le\frac{u}{24}.
\end{equation*}
Every point of $X^{+}$ is thus within sup-distance $u/24$ of $X$. This proves the first inclusion.

For the second inclusion, note that $\varrho_0\ge0$, so hypothesis~(1) gives
$M\ge 3LM_1\cdot 2L\ge 6L^2$. Hence
\begin{equation*}
L^2\ell_j\le\frac{M\ell_j}{6}=\frac{u}{6}.
\end{equation*}
A point within $L^2\ell_j$ of $X^{+}$ is therefore within $u/24+u/6$ of $X$. Since
$1/24+1/6=5/24<1/4$, such a point lies in $X^{(1/4)}$.

\emph{Proof of \eqref{9.5:eq-hull-geometry-b}.} On $X^{(1/4)}$ we have
$\operatorname{dist}_\infty(y,X)\le u/4$, so $2-4\operatorname{dist}_\infty(y,X)/u\ge1$ and
$\chi(y)=1$. Off $X^{(1/2)}$ we have $\operatorname{dist}_\infty(y,X)\ge u/2$, so
$2-4\operatorname{dist}_\infty(y,X)/u\le0$ and $\chi(y)=0$.

For the derivative bound, recall that $\operatorname{dist}_\infty(\cdot,X)$ is $1$-Lipschitz in the
sup-norm. Hence $y\mapsto 2-4\operatorname{dist}_\infty(y,X)/u$ is $(4/u)$-Lipschitz in the sup-norm.
The map $t\mapsto\min\{1,\max\{0,t\}\}$ is $1$-Lipschitz, so $\chi$ is $(4/u)$-Lipschitz in the
sup-norm. The lattice difference $\partial_\mu\chi$ is a difference quotient of $\chi$ over one
bond, whose sup-length equals the step of the quotient. Therefore $|\partial_\mu\chi|\le4/u$.

\emph{Proof of \eqref{9.5:eq-hull-geometry-c}.} The tower $\mathfrak{S}$ is truncated at $j$ and
formed over the label's determining set current at the read, whose level at $y$ is $\iota(y)$ (see
the proof of \eqref{9.5:eq-hull-geometry-d} below). Hence a cell $y$ of its determining set has
level at most $j$ and at most $\iota(y)$, that is, at most $j\wedge\iota(y)$. It is live if
equality holds, and a distance cell if the level is strictly smaller. It suffices to show that no
distance cell $y$ has $\tilde\Delta(y)\cap X^{(3/4)}\neq\emptyset$.

By the distance-cell property of the box tower over the hull, every distance cell of $\mathfrak{S}$
lies within $3M_1\ell_j$ of $(\widehat{X})^c$.

Next, $\widehat{X}$ contains every $\pi_j$-cube $\infty$-adjacent to $X$, and these cubes have
side $u$. Hence $\operatorname{dist}_\infty(X,(\widehat{X})^c)\ge u$. By the triangle inequality,
\begin{equation*}
\operatorname{dist}_\infty\bigl(X^{(3/4)},(\widehat{X})^c\bigr)\ge u-\tfrac34u=\tfrac{u}{4}.
\end{equation*}
A cell of level at most $j$, taken together with its neighbours, has sup-diameter at most
$3\ell_j$.

Suppose $y$ is a distance cell and $q\in\tilde\Delta(y)\cap X^{(3/4)}$. Some point $p$ of $y$ lies
within $3M_1\ell_j$ of $(\widehat{X})^c$, and $|p-q|_\infty\le3\ell_j$. Hence $q$ lies within
$(3M_1+3)\ell_j$ of $(\widehat{X})^c$, so
\begin{equation*}
\tfrac{M}{4}\,\ell_j=\tfrac{u}{4}\le(3M_1+3)\ell_j .
\end{equation*}
On the other hand, hypotheses~(2) and~(3) and $M_1\ge1$ give
\begin{equation*}
\frac{M}{4}>64\cdot30LM_1\ge3M_1+3 .
\end{equation*}
This contradiction shows that $y$ is live.

\emph{Proof of \eqref{9.5:eq-hull-geometry-d}.} First, the tower $\boldsymbol{\Omega}(\widehat{X})$
is formed over the label's determining set current at the read. On a live piece this is the
piece's frozen structure (Notation~\ref{9.5:not-dates-weights-norms}), that is $\mathfrak{B}_d$, so
the level entering the live condition of \eqref{9.5:eq-hull-geometry-c} is $\lambda_d=\iota$.

By hypothesis~(4), applied to $d''\le\tau\le d$,
\begin{equation*}
\lambda_{d''}\le\lambda_\tau\le\lambda_d=\iota .
\end{equation*}

Next we show $\lambda_\tau\le j+1$ on $X^{(3/4)}$. Let $y\in X^{(3/4)}$ and choose $x\in X$ with
$|y-x|_\infty\le\frac34u$, so in particular $|y-x|_\infty\le u$. By the first step of the proof of
\eqref{9.5:eq-hull-geometry-a}, $\lambda_\tau(x)\le j$. Suppose $\lambda_\tau(y)\ge j+2$. By
hypothesis~(5) the levels of $\mathfrak{B}_\tau$ on adjacent strata are consecutive. Hence the
segment $[x,y]$ would cross the whole stratum of level $j+1$. That stratum has width at least
\begin{equation*}
M\ell_{j+1}=Lu>u ,
\end{equation*}
which is impossible since $|y-x|_\infty\le\frac34u<u<Lu$. Hence $\lambda_\tau(y)\le j+1$.

Only the inequality that the width exceeds $\frac34u=\frac34M\ell_j$ is used in this argument, so
the conclusion holds whichever of the two scales, $j$ or $j+1$, is used to measure the outermost
layer of the stratum: both widths, $u$ and $Lu$, exceed $\frac34u$.

Combining the two bounds with \eqref{9.5:eq-hull-geometry-c}, on $X^{(3/4)}$ we get
\begin{equation*}
\lambda_\tau\le\min\{\iota,j+1\}\le(j\wedge\iota)+1=\lambda_{\mathfrak{S}}+1 .
\end{equation*}
Here \eqref{9.5:eq-hull-geometry-c} applies because every cell of $\mathfrak{S}$ containing a
point of $X^{(3/4)}$ meets $X^{(3/4)}$.

Finally, let $j=k(\tau)$. Every level of $\mathfrak{B}_\tau$ is at most $k(\tau)=j$, so
$\lambda_\tau\le j$. Together with $\lambda_\tau\le\iota$ this gives
$\lambda_\tau\le j\wedge\iota=\lambda_{\mathfrak{S}}$ on $X^{(3/4)}$.
\end{proof}

\begin{lemma}[The chart on the hull]\label{9.5:lem-hull-chart}
In the standing notation of Notation~\ref{9.5:not-dates-weights-norms}, with $\eta$ as there, let
$\mathfrak{S}=\boldsymbol{\Omega}(\widehat{X})$ be the box tower over the hull $\widehat{X}$, with
levels $\lambda_{\mathfrak{S}}$, units $\ell_{\mathfrak{S}}(y)=L^{\lambda_{\mathfrak{S}}(y)}\eta$, scaled
distance $d_{\mathfrak{S}}$ and determining set $\mathfrak{B}$. For a cell $y$ of $\mathfrak{S}$ let
$\tilde\Delta(y)$ be the neighbourhood of $y$ over which the cell norm
$|A|_y:=\ell_{\mathfrak{S}}(y)\sup_{\tilde\Delta(y)}|A|$ of a field $A$ is taken. Let $\delta_1>0$ be
the decay rate of the exponential smearing: for a non-negative function $f$ on $\mathfrak{B}$ put
\begin{equation*}
f^{\approx}(x):=\sum_{y\in\mathfrak{B}}e^{-2\delta_1 d_{\mathfrak{S}}(x,y)}f(y).
\end{equation*}
Let $\mathfrak{S}$ be admissible, let $U_c$ be a real configuration satisfying
\cite[(3.35)]{B-CMP99b} and critical for $\mathfrak{S}$, and let $\mathsf{m}_c$ be the complex
parameter of the chart (a letter distinct from the machine label $\mathsf{m}$ of
Notation~\ref{9.5:not-dates-weights-norms}), with $|\mathsf{m}_c|\le\mathfrak{r}_0$, where
$\mathfrak{r}_0$ is its admissible radius. For a complex field $B$ on $\mathfrak{B}$ with
$\sup|B|\le a_R'$ write
$\mathfrak{X}(B):=\mathfrak{X}_{\mathfrak{S}}(B;U_c,\mathsf{m}_c)$ for the chart and
$\mathbf{V}(B):=e^{i\eta\mathfrak{X}(B)}U_c$ for the chart configuration, and $\mathcal{J}(\cdot)$ for
the current of a configuration. Assume the following three statements.
\begin{itemize}
\item[(a)] \emph{The local Lipschitz bound for the chart.} For complex $B^1,B^2$ on $\mathfrak{B}$ with
$\sup|B^i|\le a_R'$, with $\mathsf{b}:=|B^2-B^1|$ and $E:=\mathfrak{X}(B^2)-\mathfrak{X}(B^1)$, at
every cell $y$ of $\mathfrak{S}$ each of the three numbers
\begin{equation*}
\begin{gathered}
\ell_{\mathfrak{S}}(y)\sup_{\tilde\Delta(y)}|E|,\qquad
\ell_{\mathfrak{S}}(y)^2\sup_{\tilde\Delta(y)}|\nabla_{U_c}E|,\\
\ell_{\mathfrak{S}}(y)^3\sup_{\tilde\Delta(y)}\bigl|\mathcal{J}(\mathbf{V}(B^2))-\mathcal{J}(\mathbf{V}(B^1))\bigr|
\end{gathered}
\end{equation*}
is at most $B_{\natural}\sup_{\tilde\Delta(y)}\mathsf{b}^{\approx}$, where the constant
$B_{\natural}$ depends only on the dimension $d$, on $L$, on $M_1$ and on two further fixed
parameters $\theta$ and $R$.
\item[(b)] \emph{The regularity corollary for the chart.} For complex $B$ on $\mathfrak{B}$ with
$\sup|B|\le a_R'$, the chart $\mathfrak{X}(B)$ and the current $\mathcal{J}(\mathbf{V}(B))$ of the
chart configuration are analytic in $(B,\mathsf{m}_c)$; and $\mathfrak{X}(0)=0$.
\item[(c)] \emph{The kernel-sum bound.} For every point $x$,
$\sum_{y\in\mathfrak{B}}e^{-\delta_1 d_{\mathfrak{S}}(x,y)/2}\le c_1$.
\end{itemize}
Then for every complex $B$ on $\mathfrak{B}$ with $\sup|B|\le a_R'$ and every cell $y$ of
$\mathfrak{S}$ each of the three numbers
\begin{equation}\label{9.5:eq-hull-chart-1}
\begin{gathered}
\ell_{\mathfrak{S}}(y)\sup_{\tilde\Delta(y)}|\mathfrak{X}(B)|,\qquad
\ell_{\mathfrak{S}}(y)^2\sup_{\tilde\Delta(y)}|\nabla_{U_c}\mathfrak{X}(B)|,\\
\ell_{\mathfrak{S}}(y)^3\sup_{\tilde\Delta(y)}\bigl|\mathcal{J}(\mathbf{V}(B))-\mathcal{J}(U_c)\bigr|
\end{gathered}
\end{equation}
is at most $B_{\natural}c_1\sup|B|$;
moreover $\mathfrak{X}(0)=0$, and $B\mapsto(\mathfrak{X}(B),\mathcal{J}(\mathbf{V}(B)))$ is
holomorphic.
\end{lemma}

\begin{proof}
We apply (a) with $(B^1,B^2)=(0,B)$. This choice is allowed, since $\sup|0|=0\le a_R'$ and
$\sup|B|\le a_R'$ by assumption. Then $\mathsf{b}=|B|$.

Next, by (b), $\mathfrak{X}(0)=0$.
Consequently $E=\mathfrak{X}(B)-\mathfrak{X}(0)=\mathfrak{X}(B)$, and
$\mathbf{V}(0)=e^{0}U_c=U_c$, so that $\mathcal{J}(\mathbf{V}(0))=\mathcal{J}(U_c)$. The three numbers
of (a) are therefore exactly the three numbers in \eqref{9.5:eq-hull-chart-1}, and each
is at most $B_{\natural}\sup_{\tilde\Delta(y)}|B|^{\approx}$.

For every point $x$ and every $y$ we have $d_{\mathfrak{S}}(x,y)\ge0$ and $\delta_1>0$. Hence
$e^{-2\delta_1 d_{\mathfrak{S}}(x,y)}\le e^{-\delta_1 d_{\mathfrak{S}}(x,y)/2}$, and (c) gives
\begin{equation*}
|B|^{\approx}(x)=\sum_{y\in\mathfrak{B}}e^{-2\delta_1 d_{\mathfrak{S}}(x,y)}|B(y)|
\le\sup|B|\sum_{y\in\mathfrak{B}}e^{-\delta_1 d_{\mathfrak{S}}(x,y)/2}\le c_1\sup|B|.
\end{equation*}
Taking the supremum over $x\in\tilde\Delta(y)$ and inserting this bound into the previous paragraph
gives the bound $B_{\natural}c_1\sup|B|$ on each of the three numbers in
\eqref{9.5:eq-hull-chart-1}.

Finally, by (b) the pair $(\mathfrak{X}(B),\mathcal{J}(\mathbf{V}(B)))$ is analytic in
$(B,\mathsf{m}_c)$; holding the parameter $\mathsf{m}_c$ fixed, the map
$B\mapsto(\mathfrak{X}(B),\mathcal{J}(\mathbf{V}(B)))$ is holomorphic on $\{\sup|B|\le a_R'\}$.
\end{proof}

\begin{definition}[The ambient-centre clause and the compared data]\label{9.5:def-ambient-centre}
Fix one of the two compared constructions $\mathsf{m}\in\{\mathsf{A},\mathsf{B}\}$ and a term
$(F,X,j)$ born at the date $\tau$, in the sense of
Notation~\ref{9.5:not-dates-weights-norms}, within the scope of
Definition~\ref{9.5:def-fattened-sets}. Let $\mathfrak{T}^{\mathsf{m}}[\varsigma](X)$ be the class of
chart data about $(F,X,j)$, and let $r_c$ be the coupling-free radius that enters its definition
through the condition $|B|\le r_c$. We add to that definition one clause of membership, read
with the fattening $X^{(3/4)}$ of~\eqref{9.5:eq-fattened-sets-c}, and one ceiling on $r_c$.

\emph{The ambient-centre clause.} The real centre $U_c$ of a datum of
$\mathfrak{T}^{\mathsf{m}}[\varsigma](X)$ agrees on the bonds of $X^{(3/4)}$ with an
$\mathrm{SU}(N)$-valued configuration $U^{gl}$ on $\Omega_0(\mathfrak{B}_\tau)$ which obeys the
regularity condition \cite[(3.35)]{B-CMP99b}, at the radius $\alpha_0^{(99)}$ at which the class
$\mathfrak{K}_{109}$ is formed (the radius used in clause $(\natural 1)$ of
Definition~\ref{9.5:def-relaxed-class}), over $\mathfrak{B}_{d''}$ for every date $d''\le\tau$.
That is, there exists such a $U^{gl}$ with
\begin{equation}\label{9.5:eq-ambient-centre-a}
\begin{gathered}
U_c(b)=U^{gl}(b)\qquad\text{for every bond } b \text{ of } X^{(3/4)},\\
U^{gl}\ \text{obeys \cite[(3.35)]{B-CMP99b} at the radius } \alpha_0^{(99)}
\text{ over } \mathfrak{B}_{d''},\\
\text{for every date } d''\le\tau .
\end{gathered}
\end{equation}

\emph{The ceiling on the chart radius.} With $a_R'$ the radius and $B_{\natural}$, $c_1$ the
constants of Lemma~\ref{9.5:lem-hull-chart}, and $\alpha_1^{(99)}$ the second radius at which
$\mathfrak{K}_{109}$ is formed, put
\begin{equation}\label{9.5:eq-ambient-centre-b}
r_c^{\natural}:=\min\Bigl\{\,r_c,\ a_R',\ \frac{\alpha_1^{(99)}}{4L^2B_{\natural}c_1}\Bigr\},
\end{equation}
and read the condition $|B|\le r_c$ of the definition of $\mathfrak{T}^{\mathsf{m}}[\varsigma](X)$
as $|B|\le r_c^{\natural}$. The reading $|B|\le r_c^{\natural}$, with $r_c^{\natural}$ as
in~\eqref{9.5:eq-ambient-centre-b}, is a coupling-free ceiling on an auxiliary radius. In the
order of choice of the conditions on which the class $\mathfrak{T}^{\mathsf{m}}[\varsigma](X)$
depends, $r_c^{\natural}$ is placed after $B_{\natural}$ and $c_1$, which depend on the dimension
$d=4$, on $L$, on a parameter $\theta$ fixed before them in that order, on $M_1$ and on $R$, and
before $\gamma$. It bounds no quantity from below, and $L$
enters it only through a constant already determined at that place.

\emph{The conditions that read $r_c$.} Every condition of that order of choice whose right member
is $r_c$, or one of the radii $\mathfrak{r}^{cl}$, $\varepsilon_\tau$ whose definitions contain
$r_c$, is read at $r_c^{\natural}$ (for $\mathfrak{r}^{cl}$ and $\varepsilon_\tau$: formed with
$r_c^{\natural}$ in place of $r_c$). These conditions are
\begin{equation}\label{9.5:eq-ambient-centre-c}
\begin{gathered}
c_3\sup_{\mathsf{t}}\alpha^{s}(\mathsf{t})\le r_c^{\natural},\\
4\sup_{\widehat{X}}|A'|_c\le r_c^{\natural}\qquad(\widehat{X}\subset Y),
\end{gathered}
\end{equation}
together with the condition that the data stay in the ball of radius $r_c^{\natural}$. In the
first, $\alpha^{s}(\mathsf{t})$ is the coupling-small radius at flow time $\mathsf{t}$; this is the
condition used by the inclusion of the coupling-small classes in the chart ball, and $c_3$ is the
constant of that inclusion. In the second,
$\widehat{X}$ is the hull of Lemma~\ref{9.5:lem-hull-chart}, $Y$ is the region of the parent term
on which its slice is given, $|\cdot|_c$ is the norm in which the parent's table measures that
slice, and $A'$ is the parent's slice, small
in the sense of the parent's table, plus the moves; this is the condition used for the conclusion,
for the first two operations of the construction, that the datum read out of a point
$\mathbf{U}$ of a shape lies in the class at $\varsigma=0$, through the bound
\begin{equation*}
|B[\mathbf{U}]-B[U_e]|\le 4|\mathcal{Y}|_c
\end{equation*}
on the datum of a point of a shape, $U_e$ being the entry of the shape and $\mathcal{Y}$ the
moves. Each left member tends to $0$ with $\gamma$; each of these conditions is therefore a
ceiling on $\gamma$, in the supremum form displayed, chosen at the last place of the order. No
condition of a new kind is introduced.

\emph{The compared class.} Let $X^{+}$ be the ball collar of~\eqref{9.5:eq-fattened-sets-b}. Put
\begin{equation}\label{9.5:eq-ambient-centre-e}
\begin{aligned}
\mathfrak{T}^{\mathsf{m},\natural}_F[\varsigma](X)
:={}&\mathfrak{T}^{\mathsf{m}}[\varsigma](X)\ \text{with the clause \eqref{9.5:eq-ambient-centre-a},
with } |B|\le r_c^{\natural}\\
&\text{of \eqref{9.5:eq-ambient-centre-b} in place of } |B|\le r_c,\ \text{and with}\\
&(\mathbf{V}(B),J(\mathbf{V}(B)))\restriction X\in\mathfrak{D}^{\mathsf{m}}_F[\varsigma]
\ \text{read with the table on } X^{+},
\end{aligned}
\end{equation}
that is, with the membership in the domain $\mathfrak{D}^{\mathsf{m}}_F[\varsigma]$ read in the
convention used for the domains of the stored terms: the pair is given on $X^{+}$ and obeys the
table there.

\emph{What is read of $F$'s table.} On $X^{+}$ the table of $F$ satisfies
\begin{equation}\label{9.5:eq-ambient-centre-d}
\begin{gathered}
\mathsf{u}^{a}_F\le 1.4\,\theta^{\tau}_a\quad\text{and}\quad
\mathsf{c}_F^{+}\,\mathsf{u}^{a}_F\le 1.4\,\theta^{\tau}_a\qquad (a\in\{1,3,6,7\}),\\
E_5<1.4\,\theta^{\tau}_5,
\end{gathered}
\end{equation}
and, on $X^{+}$, the entry $E_5$ and the cube clause of $U_c$ are read inside $F$'s own table.
\end{definition}

\begin{proof}[Proof of the table bounds \eqref{9.5:eq-ambient-centre-d}]
Let $y\in X^{+}$. The table coefficient of the term satisfies the bound
$\mathsf{c}_F^{+}\le 1+\beta$ of the class of chart data, $\beta$ being the tolerance of the
tables (not the coupling increment $\beta_{k+1}$). The tolerance $\beta$ is chosen with
$\beta\le 2/5$, so that $1+\beta\le 1.4$; we use this below without further mention. The units of
the term are
\begin{equation*}
\mathsf{u}^{a}_F=\alpha_{a,n_F}\,\ell_{n_F}^{-e_a},
\end{equation*}
with $n_F$ the level of $F$'s own units and $e_a$ the exponent attached to the entry $a$. On live
pieces these units are the weighted units $w\,\mathsf{u}_\ell$, the weight $w=w_\tau(y)$ times the
unit $\mathsf{u}_\ell=r_a(\lambda_\tau(y))$, and $w\,\mathsf{u}_\ell=\theta^{\tau}_a$ by the
definition of the weighted thresholds $\theta^{d}_a(y)=w_d(y)\,r_a(\lambda_d(y))$ of
Notation~\ref{9.5:not-dates-weights-norms} at $d=\tau$. Hence on live pieces
\begin{equation*}
\mathsf{u}^{a}_F=\theta^{\tau}_a,\qquad
\mathsf{c}_F^{+}\,\mathsf{u}^{a}_F\le(1+\beta)\,\theta^{\tau}_a .
\end{equation*}
Since $\theta^\tau_a>0$ and $1+\beta\le 1.4$, this gives the first line
of~\eqref{9.5:eq-ambient-centre-d} on live pieces, both for $\mathsf{u}^{a}_F$ and for
$\mathsf{c}_F^{+}\,\mathsf{u}^{a}_F$.

Let now $F=\mathbb{C}^{(i)}_{k_a}(X)$ be a stored term of stage $i$ at step $k_a$, so that
$j=k_a$ and its domain $\mathfrak{D}^{\mathsf{m}}_F[\varsigma]$ is defined by the flat table
$\mathsf{T}^{\flat}_i$ at step $k_a$. The bounds on the flat tables of the stored terms give the
following, entry by entry: the entry $a$ of $\mathsf{T}^{\flat}_i$ is compared with
$\theta^{\tau}_a$, and we write $\theta^{\tau}$ for the family $(\theta^{\tau}_a)_a$. On the part
$Z^{\mathrm{on}}$ of the stored term's region,
$\mathsf{T}^{\flat}_i\le\theta^{\tau}$, because the coefficients of the flat table there are at
most $1$.
On the part $\mathfrak{O}$ of the stored term's region,
\begin{equation*}
\mathsf{T}^{\flat}_i\le\Bigl(f_{k_a,n}+\tfrac{9}{16}\,\theta_S\,\beta\,2^{-\nu}\Bigr)\mathsf{u}_n
\le 1.1\,\theta^{\tau},
\end{equation*}
where $n$ is the level of the units $\mathsf{u}_n$ in which the flat table is written there,
$f_{k_a,n}$ is the coefficient of the flat table at step $k_a$ and level $n$, $\theta_S$ is the
constant of its $\beta$-term and $2^{-\nu}$ the geometric factor of that term. The same bound
$\mathsf{T}^{\flat}_i\le1.1\,\theta^{\tau}$ holds on the part $W$ of the region, where the unit is
$\mathsf{F}^{(k_a)}w\,\mathsf{u}_\ell$, $\mathsf{F}^{(k_a)}$ being the step factor of the flat
table. No table coefficient occurs in the domain of a stored term, which is defined by the flat
table alone; since $1.1<1.4$, these bounds give
$\mathsf{T}^{\flat}_i\le 1.1\,\theta^{\tau}\le 1.4\,\theta^{\tau}$, the first line
of~\eqref{9.5:eq-ambient-centre-d} for a stored term, whose table is $\mathsf{T}^{\flat}_i$.
Finally, the clause on $U$ in the table of the stored term is
bounded by $(1+2\beta)\,r_a$; this is the threshold of the cube clause of $U_c$ that is read
inside $F$'s table in the cube-clause reading of Definition~\ref{9.5:def-ambient-centre}.

The entry $E_5$ is read inside $F$'s table: a datum of $\mathfrak{D}^{\mathsf{m}}_F[\varsigma]$
satisfies $E_5$ strictly below the threshold of $F$'s table at the entry $a=5$. By the bounds
above, that threshold is at most $(1+\beta)\,\theta^{\tau}_5$ on live pieces and at most
$1.1\,\theta^{\tau}_5$ for a stored term, hence at most $1.4\,\theta^{\tau}_5$ in both cases. This
gives $E_5<1.4\,\theta^{\tau}_5$, the second line of~\eqref{9.5:eq-ambient-centre-d}.
\end{proof}

\begin{theorem}[The cut-off continuation of every datum]\label{9.5:thm-cutoff-continuation}
Let $(F,X,j)$, born at the date $\tau$, be as in Definition~\ref{9.5:def-fattened-sets}. Assume:
\begin{itemize}
\item[(i)] the scope condition \eqref{9.5:eq-fattened-sets-a};
\item[(ii)] the floors on $L$ and on the auxiliary parameters under which
Lemma~\ref{9.5:lem-hull-geometry} (geometry of the hull) is stated;
\item[(iii)] the local Lipschitz bound for the chart, in the form in which it follows from the
regularity corollary for the chart, for the chart over $\mathfrak{S}$ attached to $(F,X,j)$;
\item[(iv)] clauses (e1) and (e2) of the lemma on the continuation class stated with the arrest
theorem (a hypothesis of Lemma~\ref{9.5:lem-relaxed-heredity}), and the pointwise arithmetic of
the second inclusion of its clause (e3);
\item[(v)] the table condition \eqref{9.5:eq-ambient-centre-d}.
\end{itemize}
Let $\mathsf{v}\in\mathfrak{T}^{\mathsf{m},\natural}_F[\varsigma](X)$, $\varsigma\in[0,1]$, be a datum
of the class \eqref{9.5:eq-ambient-centre-e}, with chart datum $B$, real centre $U_c$, ambient
field $U^{gl}$ of the ambient-centre clause \eqref{9.5:eq-ambient-centre-a}, and slice
$S=(\mathbf{V}(B),J(\mathbf{V}(B)))$, where $\mathbf{V}(B)=e^{i\eta\mathfrak{X}(B)}U_c$; let
$r_c^{\natural}$ be the radius fixed by the ceiling \eqref{9.5:eq-ambient-centre-b}. Let $\chi$ be
the cut-off of \eqref{9.5:eq-hull-geometry-b}, and put $\chi\mathfrak{X}(B):=0$ off $X^{(1/2)}$.
Define, on $\Omega_0(\mathfrak{B}_\tau)$,
\begin{equation}\label{9.5:eq-cutoff-continuation-a}
\tilde p(B):=\bigl(e^{i\eta\chi\mathfrak{X}(B)}U^{gl},\ \mathbf{1}_{X^{+}}\cdot J(\mathbf{V}(B))\bigr),
\end{equation}
read in the decomposition $\tilde U:=U^{gl}$, $\tilde A:=\chi\mathfrak{X}(B)$. Then the following
hold.
\begin{itemize}
\item[(E1), (E2)] The pair $\tilde p(B)$ agrees with the slice $S$, and its entries with the
slice's own entries, as follows:
\begin{equation}\label{9.5:eq-cutoff-continuation-b}
\begin{aligned}
&\text{(E1)}\quad \tilde p(B)=S,\ \text{decomposition included, on the bonds of } X^{(1/4)}\\
&\qquad\qquad \text{in the first slot and on the bonds of } X^{+}\ \text{in the second slot;}\\
&\qquad\qquad \text{in particular } \tilde p(B)\restriction X=S\restriction X,\ \text{and for every }
 y\in X^{+}:\\
&\qquad\qquad Q_a(\tilde p(B))(y)=E_a(S)(y)\quad(a\neq3),\qquad
 Q_3(\tilde p(B))(y)\le E_3(S)(y);\\
&\text{(E2)}\quad \tilde p(B)=(U^{gl},0)\quad\text{off } X^{(1/2)}.
\end{aligned}
\end{equation}
\item[(E3)] The pair lies in the relaxed continuation class of
Definition~\ref{9.5:def-relaxed-class}:
\begin{equation}\label{9.5:eq-cutoff-continuation-c}
\tilde p(B)\in\mathfrak{G}^{\natural}(\mathfrak{B}_\tau;X).
\end{equation}
\item[(E4)] At fixed $(U_c,U^{gl})$,
\begin{equation}\label{9.5:eq-cutoff-continuation-d}
\{B:\ |B|<r_c^{\natural}\}\ni B\longmapsto\tilde p(B)\quad\text{is holomorphic,}
\end{equation}
and takes its values in the open set defined by \eqref{9.5:eq-relaxed-class-a} and
\eqref{9.5:eq-relaxed-class-b} inside the affine space \eqref{9.5:eq-relaxed-class-c}.
\item[(E5)] For every date $d''\le\tau$ and every union of cubes $S'\subset X$, writing $X'$ and
$B'$ for the two sets of a frame (the sets denoted $X'$ and $B$ in part (iii) of
\eqref{9.5:eq-relaxed-heredity-a}; the prime distinguishes the frame's set from the chart
datum $B$):
\begin{equation}\label{9.5:eq-cutoff-continuation-e}
\begin{aligned}
&(\beta1)\quad \tilde p(B)\restriction\Omega_0(\mathfrak{B}_{d''})
 \in\mathfrak{K}_{109}\bigl(\mathfrak{B}_{d''};\gamma_0^{(k(d''))}\bigr);\\
&(\beta2)\quad a_{\flat}\le\gamma_0^{(k(d''))}\ \text{on } X'\cup B',\
 \text{for every frame with unit set } \mathfrak{B}_{d''}\\
&\qquad\qquad \text{and } X'\cup B'\subset X^{+};\\
&(\beta3)\quad \text{if } \tau=(k_a;n+1):\ \tilde p(B)\ \text{lies in }
 \mathfrak{K}_{109}\bigl(\mathfrak{B}^{\circ};\gamma_0^{(k_a)}\bigr)\ \text{on } X_{\mathfrak{a}},\\
&\qquad\qquad \mathfrak{B}^{\circ}:=\mathbb{B}^{(n+1)}_{k_a}\restriction X_{\mathfrak{a}};\\
&(\beta4)\quad \tilde p(B)^{S'}\in\mathfrak{G}^{\natural}(\mathfrak{B}_{d''};S').
\end{aligned}
\end{equation}
Here $a_{\flat}$ in $(\beta2)$ is evaluated at the pair $\tilde p(B)$, and $\tilde p(B)^{S'}$ is
formed as in part (i) of \eqref{9.5:eq-relaxed-heredity-a}. For $d''=\tau$, $(\beta1)$ is the
ambient membership at the birth set and density of $F$; for $d''<\tau$ it is the same membership
at every older date. $(\beta3)$ is the corresponding ambient membership at the stage-$(n+1)$
structure $\mathfrak{B}^{\circ}$, where $X_{\mathfrak{a}}$ denotes the region on which this
membership is required of the stored terms of stage $n+1$ at step $k_a$.
\end{itemize}
No constant, floor or ceiling in this statement depends on $g$, $k$, $K$, $N$, $N_0$,
$\check{R}$, an age, a depth, $|X|$ or the volume. The one new condition is the ceiling
\eqref{9.5:eq-ambient-centre-b}.
\end{theorem}

\begin{proof}
Throughout, $(\chi\mathfrak{X}(B))(b)=\chi(b_-)\,\mathfrak{X}(B)(b)$ for a bond $b$ with initial
point $b_-$.

\emph{(E1) and (E2).} By \eqref{9.5:eq-hull-geometry-b}, $\chi=1$ on $X^{(1/4)}$ and $\chi=0$ off
$X^{(1/2)}$. By the ambient-centre clause \eqref{9.5:eq-ambient-centre-a}, $U^{gl}=U_c$ on the
bonds of $X^{(3/4)}$.

Let $b$ be a bond of $X^{(1/4)}$. Then $\chi(b_-)=1$, so $\tilde A(b)=\mathfrak{X}(B)(b)$, and
$\tilde U(b)=U^{gl}(b)=U_c(b)$. The first slot of $\tilde p(B)$ at $b$ is therefore
$e^{i\eta\mathfrak{X}(B)(b)}U_c(b)=\mathbf{V}(B)(b)$, in the decomposition $(U_c,\mathfrak{X}(B))$
in which the slice is read. On the bonds of $X^{+}$ the second slot is $J(\mathbf{V}(B))$ by the
definition \eqref{9.5:eq-cutoff-continuation-a}. The bonds of $X$ are bonds of $X^{(1/4)}$ and
of $X^{+}$, so $\tilde p(B)\restriction X=S\restriction X$.

Off $X^{(1/2)}$ we have $\chi=0$, hence $\tilde A=0$ and the first slot is $U^{gl}$. By
\eqref{9.5:eq-hull-geometry-a}, $X^{+}\subset X^{(1/24)}\subset X^{(1/2)}$, so
$\mathbf{1}_{X^{+}}=0$ there and the second slot vanishes. This is (E2).

Now let $y\in X^{+}$. By hypothesis (iv), the functionals $Q_a(\cdot)(y)$ are read on the cell of
$y$ and its neighbours in $\mathfrak{B}_\tau$. These lie within distance
\begin{equation*}
3\ell_\tau(y)\le 3L\ell_j\le L^2\ell_j
\end{equation*}
of $y$; the last inequality holds under the floors on $L$ of
Lemma~\ref{9.5:lem-hull-geometry}, hypothesis (ii). By \eqref{9.5:eq-hull-geometry-a}, every point
within $L^2\ell_j$
of $X^{+}$ lies in $X^{(1/4)}$. Hence the bonds read, their shifts $b+\eta e_\mu$ and the
plaquettes on them are of $X^{(1/4)}$, where the first slot agrees with that of $S$, decomposition
included. This gives $Q_a(\tilde p(B))(y)=E_a(S)(y)$ for $a\neq3$.

For the second slot, a bond at $y$ either is of $X^{+}$, where
$\tilde{\mathbb{J}}=J(\mathbf{V}(B))$, or is not, where $\tilde{\mathbb{J}}=0$. This gives
$Q_3(\tilde p(B))(y)\le E_3(S)(y)$, and proves \eqref{9.5:eq-cutoff-continuation-b}.

\emph{(E3), clause \eqref{9.5:eq-relaxed-class-a}.} The pair $\tilde p(B)$ has the form required in
Definition~\ref{9.5:def-relaxed-class}: $\tilde U=U^{gl}$ takes values in the gauge group by
\eqref{9.5:eq-ambient-centre-a}. Fix a date $d''\le\tau$.

The bound \cite[(3.35)]{B-CMP99b} for $\tilde U=U^{gl}$ over $\mathfrak{B}_{d''}$, at the radius
$\alpha_0^{(99)}$, is part of \eqref{9.5:eq-ambient-centre-a}.

For \cite[(3.37)]{B-CMP99b} for $\tilde A=\chi\mathfrak{X}(B)$, let $b$ be a bond, whether of
$X^{+}$ or not, $x:=b_-$ and
$i:=\lambda_{d''}(x)$. The clauses of \cite[(3.37)]{B-CMP99b} are imposed on the domains
$\Omega_{i'}$, and at $x$ the strongest of the clauses on $\Omega_{i'}$, $i'\le i$, is the one
with $i'=i$; we verify that one. Put $r:=r_c^{\natural}$. By the definition
\eqref{9.5:eq-ambient-centre-e} of the class, the chart datum satisfies
$\sup|B|\le r$, and $r\le a_R'$ by \eqref{9.5:eq-ambient-centre-b}. So
Lemma~\ref{9.5:lem-hull-chart} applies; that lemma is proved as an implication from the local
Lipschitz bound for the chart, which is hypothesis (iii).

\emph{Case one:} neither $x$ nor any $x+\eta e_\mu$ lies in $X^{(1/2)}$. Then $\chi$ vanishes at
these points, so $\tilde A(b)=0$ and $\nabla_{\tilde U,\mu}\tilde A(b)=0$.

\emph{Case two:} otherwise. Since $\eta<u/4$, the point $x$ and the bonds $b$, $b+\eta e_\mu$ and
$\langle x,x+\eta e_\mu\rangle$ are of $X^{(3/4)}$, where $\tilde U=U_c$; in particular
$\nabla_{\tilde U,\mu}=\nabla_{U_c,\mu}$ on them. Let $y$ be the cell of $\mathfrak{S}$ holding
$x$. It is live by \eqref{9.5:eq-hull-geometry-c}. By \eqref{9.5:eq-hull-geometry-d},
$\ell_i\le L\ell_{\mathfrak{S}}(y)$, and $b$ and $b+\eta e_\mu$ lie in the set $\tilde\Delta(y)$
over which the bounds of Lemma~\ref{9.5:lem-hull-chart} are taken.

By the first bound of Lemma~\ref{9.5:lem-cutoff-difference} (covariant difference of a cut-off
form), $|\tilde A(b)|\le|\mathfrak{X}(B)(b)|$. Lemma~\ref{9.5:lem-hull-chart} then gives
\begin{equation*}
\ell_i|\tilde A(b)|\le L\ell_{\mathfrak{S}}(y)|\mathfrak{X}(B)(b)|\le LB_{\natural}c_1r .
\end{equation*}

By the Leibniz bound of Lemma~\ref{9.5:lem-cutoff-difference} and $|\partial_\mu\chi|\le4/u$
from \eqref{9.5:eq-hull-geometry-b},
\begin{equation*}
\begin{aligned}
\ell_i^2|\nabla_{\tilde U,\mu}\tilde A(b)|
&\le L^2\ell_{\mathfrak{S}}(y)^2|\nabla_{U_c,\mu}\mathfrak{X}(B)(b)|\\
&\qquad+\Bigl(\ell_i\cdot\frac{4}{u}\Bigr)\bigl(L\ell_{\mathfrak{S}}(y)|\mathfrak{X}(B)(b+\eta e_\mu)|\bigr)\\
&\le L^2B_{\natural}c_1r+\frac{4L}{M}\,LB_{\natural}c_1r\le 2L^2B_{\natural}c_1r .
\end{aligned}
\end{equation*}
The second step uses Lemma~\ref{9.5:lem-hull-chart} for both terms, together with
$\ell_i\le\ell_{j+1}=Lu/M$. The third uses $M\ge4$, which holds under the floors of
hypothesis (ii).

By \eqref{9.5:eq-ambient-centre-b}, $r\le\alpha_1^{(99)}/(4L^2B_{\natural}c_1)$. Hence
$2L^2B_{\natural}c_1r\le\tfrac12\alpha_1^{(99)}$ and
$LB_{\natural}c_1r\le\alpha_1^{(99)}/(4L)\le\tfrac12\alpha_1^{(99)}$. Both quantities are
therefore at most $\tfrac12\alpha_1^{(99)}<\alpha_1^{(99)}$. This is \cite[(3.37)]{B-CMP99b} at
the radius $\alpha_1^{(99)}$ on $\Omega_i$ at $b$.

This proves \eqref{9.5:eq-relaxed-class-a}.

\emph{Clause \eqref{9.5:eq-relaxed-class-b}.} Let $y\in X^{+}$. By the scope condition
\eqref{9.5:eq-fattened-sets-a}, $y\in\Omega_1(\mathfrak{B}_\tau)$. By
\eqref{9.5:eq-cutoff-continuation-b}, the table on $X^{+}$ carried by the class
\eqref{9.5:eq-ambient-centre-e}, and \eqref{9.5:eq-ambient-centre-d},
\begin{equation*}
Q_a(\tilde p(B))(y)\le E_a(S)(y)<1.4\,\theta^{\tau}_a(y)<5\,\theta^{\tau}_a(y),
\qquad a\in\{1,3,5,6,7\}.
\end{equation*}
The clause $a=1$ holds at points of $X^{+}$ in the unit $\alpha_0$ of the class
$\mathfrak{G}$ itself, because $\chi\equiv1$ on the stencils read there. The term of unit
$\alpha_1$ carrying the differences $\partial_\mu\chi$ lives on
$X^{(1/2)}\setminus X^{(1/4)}$, which is disjoint from $X^{+}$; so none of the pointwise clauses
\eqref{9.5:eq-relaxed-class-b}, in particular the clause $a=1$, is imposed there. Only
\eqref{9.5:eq-relaxed-class-a} is imposed there, and it was verified above.

\emph{Clause \eqref{9.5:eq-relaxed-class-c}} holds by the definition
\eqref{9.5:eq-cutoff-continuation-a}: the second slot $\mathbf{1}_{X^{+}}J(\mathbf{V}(B))$
vanishes on every bond not of $X^{+}$. This completes the proof of
\eqref{9.5:eq-cutoff-continuation-c}.

\emph{(E4).} By Lemma~\ref{9.5:lem-hull-chart}, $B\mapsto(\mathfrak{X}(B),J(\mathbf{V}(B)))$ is
holomorphic for $|B|<r_c^{\natural}\le a_R'$. The functions $\chi$, $U^{gl}$ and
$\mathbf{1}_{X^{+}}$ do not depend on $B$. Hence both slots of \eqref{9.5:eq-cutoff-continuation-a}
are holomorphic in $B$ at fixed $(U_c,U^{gl})$. All inequalities established above for
\eqref{9.5:eq-relaxed-class-a} and \eqref{9.5:eq-relaxed-class-b} are strict, and the lattice is
finite. So $\tilde p(B)$ lies in the open set that these two clauses define inside the affine
space \eqref{9.5:eq-relaxed-class-c}. This proves \eqref{9.5:eq-cutoff-continuation-d}.

\emph{(E5).} The following are the parts of \eqref{9.5:eq-cutoff-continuation-e}.
\begin{itemize}
\item[$(\beta1)$] This follows from \eqref{9.5:eq-cutoff-continuation-c} and part (ii) of
\eqref{9.5:eq-relaxed-heredity-a}, applied with $Y=X$ and $d=\tau$. That display belongs to
Lemma~\ref{9.5:lem-relaxed-heredity}, which is proved as an implication from that lemma's clauses
listed there.
\item[$(\beta2)$] This is part (iii) of \eqref{9.5:eq-relaxed-heredity-a}, with the same choices;
there $Y^{+}_\tau=X^{+}$.
\item[$(\beta3)$] This is $(\beta1)$ at $d''=\tau$, restricted to $X_{\mathfrak{a}}$. For a birth
date $\tau=(k_a;n+1)$ one has $\mathfrak{B}_\tau=\mathbb{B}^{(n+1)}_{k_a}$, by the statement
fixing the determining sets of the dates of the stored terms of stage $n+1$, and hence
$\mathfrak{B}_\tau\restriction X_{\mathfrak{a}}=\mathfrak{B}^{\circ}$; the density of this date is
$k_a$. Thus $(\beta3)$ is the ambient membership required of these stored terms on
$X_{\mathfrak{a}}$.
\item[$(\beta4)$] This follows from \eqref{9.5:eq-relaxed-heredity-b}, applied to
\eqref{9.5:eq-cutoff-continuation-c} with $Y=X$, $d=\tau$ and with the union of cubes of that
lemma taken to be $S'$.
\end{itemize}

Finally, the constants entering the argument are the following:
\begin{itemize}
\item $L$, $M$ and the floors of Lemma~\ref{9.5:lem-hull-geometry};
\item $B_{\natural}$ and $c_1$;
\item the radii $\alpha_0^{(99)}$, $\alpha_1^{(99)}$, $r_c$, $a_R'$;
\item the coefficient $1.4$ of \eqref{9.5:eq-ambient-centre-d} and the thresholds
$\theta^{\tau}_a$.
\end{itemize}
None of them depends on the quantities excluded in the statement. The only condition
introduced here is \eqref{9.5:eq-ambient-centre-b}.
\end{proof}

\begin{remark}
(1) Of the slice on the layer $X^{(1/2)}\setminus X^{+}$, the proof uses only the chart bound of
Lemma~\ref{9.5:lem-hull-chart}, whose constant does not depend on $g$. This is why the theorem
covers every datum of $\mathfrak{T}^{\mathsf{m},\natural}_F[\varsigma](X)$, the Schwarz-lemma data
$\mathsf{v}_s$ of the uniqueness comparison and all data over which the class supremum
$\mathsf{S}^{+}_j[F]$ is taken included. It is
also why none of the membership arguments made on $X^{+}$ has to be repeated on a larger set.

(2) Of the ambient field $U^{gl}$ beyond $X^{(3/4)}$, the proof uses only
\cite[(3.35)]{B-CMP99b}.

(3) The device used here has two parts. One multiplies the exponential of a cut-off added field
onto a global configuration; one cuts off the added current by an indicator. Together they produce
from a locally given pair a pair defined on all of $\Omega_0(\mathfrak{B}_\tau)$, in a class whose
radii do not depend on the coupling.
\end{remark}

\begin{remark}\label{9.5:prop-relaxed-estimates}
The statement of this proposition is not written out here; parts of its proof are printed below.
\end{remark}

\begin{proof}[Proof of Proposition~\ref{9.5:prop-relaxed-estimates}, continued]
\label{9.5:prf-domain-reading}
Clause (i) of Proposition~\ref{9.5:prop-relaxed-estimates} rests on two lemmas. The first is the
stage lemma on stored terms over a system, which shows that the stored terms of stage $l$ have
data on the domains of a given system and bounds them. The second is the local holomorphy lemma
for stage data, items (i)--(v), which establishes the holomorphy of the stage data on
$\mathfrak{K}_P$. We list what their proofs use of the domains to which they are applied. The
notion of a system is recalled in the third part of this argument below. We take the entries of
the dictionary in the proof of the stage lemma in their order, numbered $(0)$--$(9)$, and state
for each what it uses of the system $\mathfrak{D}_l$.
\begin{itemize}
\item[(0)] The definitions. These are finite sums at the vertices of the first family. Neither of
the two kinds of radius, the $t$-radii and the $u$-radii, occurs in a definition; radii enter
only the bounds.
\item[(1)] The holomorphy entry. It is used only in the form that $\Phi_t\restriction S$, the
restriction to $S$ of the stage datum $\Phi_t$, lies in the table-defined domain
$\mathfrak{U}_c(S)$, holomorphically, on the polydisc. This is clause (a) of the containment
hypothesis at stage $l$, together with the holomorphy asserted in item (i) of the local holomorphy
lemma.
\item[(2)] The identities. They involve neither radii nor domains. The entry uses item (iv) of
the local holomorphy lemma at $\mathfrak{U}:=\operatorname{dom}(c)$, the domain of definition of
$c$. It also uses that the function $v_e$ attached to $e$ is holomorphic on
$\mathfrak{D}_l(S_e)$, where $S_e$ is the union of cubes attached to $e$, hence, by stability of
the system under restriction, on $\mathfrak{D}_l(X)$.
\item[(3)] Counting statements. They use no property of the domains.
\item[(4)] This entry uses item (v) of the local holomorphy lemma together with clause (b) of the
containment hypothesis.
\begin{itemize}
\item[--] For the near groups it uses clause (c) of that hypothesis, which supplies on the box
$\square_0\subset Y$ the hypotheses of \cite[Lemma~4]{B-CMP109}. It also uses the
proposition that re-derives the estimates \cite[(3.37)--(3.52)]{B-CMP109} on the box
$\square_0\subset Y$.
\item[--] The $\mathbb{P}$- and $W$-pieces, the pieces of these two kinds in the stage lemma,
carry no radius. They use the field $\phi$ on bounded supports, the input of the stage lemma
through kernels holomorphic on $\mathfrak{K}_P$, or the Wilson plaquettes of a
configuration $\Phi\in\mathfrak{N}^{(l-1)}$.
\end{itemize}
\item[(5)] This entry takes a supremum over $\mathfrak{U}_T\times\mathfrak{D}_{\mathcal{S}}$, the
product of the table-defined domain $\mathfrak{U}_T$ and the domain $\mathfrak{D}_{\mathcal{S}}$
of the second factor of the stored data; the relaxation leaves this supremum unchanged.
\item[(6)] This entry uses three things.
\begin{itemize}
\item[--] Identities at each pair $(\mathbb{U},\mathbb{J})$.
\item[--] The tolerance bounds on $\mathfrak{N}^{(l)}\subset\mathfrak{K}_P$, supplied by the
statement that bounds the tolerances on $\mathfrak{K}_P$.
\item[--] Uniform convergence on $\mathfrak{D}_l(X)$ under majorants independent of
$(\mathbb{U},\mathbb{J})$.
\end{itemize}
\item[(7)] This entry uses majorants on $\mathfrak{D}_l$ independent of the parameter $\lambda$
of the stage lemma (not the gauge transformation $\lambda$ of
Lemma~\ref{1.1:lem-invariance}) and of $(\mathbb{U},\mathbb{J})$.
\item[(8)] Random-walk expansions on $\mathfrak{N}^{(l)}$. They rest on
\cite[Theorem~3.10]{B-CMP99b} and on the same statement that supplies the tolerance bounds of
entry (6).
\item[(9)] This entry forms totals.
\end{itemize}
Finally, the proof uses that $\mathfrak{D}_l(X)$ is $\mathcal{G}$-invariant.

The items of the local holomorphy lemma use the following.
\begin{itemize}
\item[(i)] Membership in $\mathfrak{K}_P$.
\begin{itemize}
\item[--] Its first two steps use four statements established earlier; the three of them that
involve the determining set hold for an arbitrary admissible determining set.
\item[--] Its third step treats table entries $a=5$ and $a=8$ through a decomposition that keeps
$U$ and the cube gauge transformations, for every shift. Entries $a=1,6,7$ are read through
first-order sizes only.
\item[--] Its base condition is $\mathbb{U}\in\mathfrak{K}_P$, with the ambient convention
imposed off $Y$.
\item[--] The second slot, the current, enters locally.
\end{itemize}
\item[(ii)] The sentence quoted for this item in the list of uses earlier in the proof of
Proposition~\ref{9.5:prop-relaxed-estimates}, together with one observation: item (i) holds at
every determining set, because it reads $\mathfrak{B}$ only through hypotheses that hold for an
arbitrary admissible $\mathfrak{B}$.
\item[(iii), (iv)] These use three things:
\begin{itemize}
\item[--] an identity of Ba{\l}aban's construction;
\item[--] the Cauchy estimates on polydiscs;
\item[--] the fact that $\operatorname{dom}(c)$ is a product.
\end{itemize}
\item[(v)] The fact that the kernel display used in this item holds on $\mathfrak{K}_P$.
\end{itemize}

Hence the stage lemma uses exactly the following properties of a system $\mathfrak{D}_l$.
\begin{itemize}
\item[--] That it is open, $\mathcal{G}$-invariant and stable under restriction. For the relaxed
domains this is Lemma~\ref{9.5:lem-relaxed-domains-open}.
\item[--] The inclusion $\mathfrak{D}_l\subset\mathfrak{N}^{(l)}$, in three forms.
\begin{itemize}
\item[--] As $\mathfrak{N}^{(l)}\subset\mathfrak{K}_P$: for the relaxed domains this is part (ii)
of \eqref{9.5:eq-relaxed-heredity-a}.
\item[--] As the fat local clauses at locations of $Y$: these are not changed by the relaxation.
\item[--] As the base condition of item (i) on $Y^{+}_{d}$: for the relaxed domains this is part
(iii) of \eqref{9.5:eq-relaxed-heredity-a}.
\end{itemize}
\item[--] The containment hypothesis at stage $l$, which is established in the second part of
this argument below.
\end{itemize}

Two sizes small in the coupling occur in these entries. One is the radius $\alpha_5$ of entry
(6). The other is the bound on the field dependence in entry (8), a constant multiple of the sum
of a term proportional to the size of the local clauses and of the term $K_\Delta\gamma_0$. They
carry the factor $5w_*$ of the local clauses, with $w_*$ the weight of those clauses, and the
factor $K_\Delta\gamma_0$ of the clause in $\gamma_0$, with $K_\Delta$ the constant of that clause.
Under the ambient convention with which these entries are proved, the field away from
$Y^{+}_{d}$ is subject to no clause small in the coupling. The domain
$\mathfrak{D}^{\flat,\natural}$ is contained in the domain defined under that convention; this
containment is established earlier in the proof of
Proposition~\ref{9.5:prop-relaxed-estimates}. Hence both sizes are available on the relaxed
domains, and nothing is lost.

\emph{Containments.} We turn to the containments listed in clause (ii) of
Proposition~\ref{9.5:prop-relaxed-estimates}. We take them together with two further inputs: the
statement that the family of stored terms respects the ambient convention, and the four steps of
the proof of the inheritance theorem, which establishes by induction on the stage the domains
and the bounds of the stored terms.

For a stored term $\tau'$ with region $X_{\tau'}$ and determining set $\mathfrak{B}_{\tau'}$, we
call $\mathfrak{G}^{\natural}(\mathfrak{B}_{\tau'};X_{\tau'})$ (Definition~\ref{9.5:def-relaxed-class})
its \emph{relaxed continuation class}. Work under the relaxed ambient convention
\eqref{9.5:eq-relaxed-class-d}. Let $\tau'$ be a stored term composed at an operation
$\mathfrak{O}$ whose target determining set is $\mathfrak{B}^{t}$ and target region is $Y'$, and
let $\tilde q$ be the pair produced there.
\begin{itemize}
\item[--] By item (5) of the lemma on the order $\preccurlyeq$ of dated sets,
$X_{\tau'}\subset Y'$ and $d_{\tau'}\le\operatorname{date}(\mathfrak{B}^{t})$, where $d_{\tau'}$
is the date of $\tau'$.
\item[--] By \eqref{9.5:eq-relaxed-heredity-c},
$\tilde q\in\mathfrak{G}^{\natural}(\mathfrak{B}^{t};Y')$.
\item[--] By \eqref{9.5:eq-relaxed-heredity-b}, applied with the two facts just stated,
\begin{equation*}
\tilde q^{X_{\tau'}}\in\mathfrak{G}^{\natural}(\mathfrak{B}_{\tau'};X_{\tau'}),
\end{equation*}
the relaxed continuation class of $\tau'$.
\end{itemize}
This is the relaxed form of the continuation clause of the arrest theorem. Combined with part (ii)
of \eqref{9.5:eq-relaxed-heredity-a}, applied to $\tilde q^{X_{\tau'}}$ at its own density, it
gives
\begin{equation*}
\tilde q^{X_{\tau'}}\restriction\Omega_0(\mathfrak{B}_{\tau'})
\in\mathfrak{K}_{109}(\mathfrak{B}_{\tau'};\gamma_0^{(k_{\tau'})}),
\end{equation*}
where $k_{\tau'}$ is the density of $\tau'$ and $\gamma_0^{(k)}$ is the radius of the current slot
at density $k$.
With this come the remaining containments listed in clause (ii) of
Proposition~\ref{9.5:prop-relaxed-estimates}, including the ambient containment at
$\mathfrak{B}_{\tau'}=\mathbb{B}^{(n+1)}_{k_a}$, the stage-$(n+1)$ structure at step $k_a$.

The pointwise clauses on $X_{\tau'}$ --- those of clause (A) of the arrest theorem and of the
collar table, verified in steps (1)--(3) of the proof of the inheritance theorem --- do not involve
the continuation class, so they are unmoved by the relaxation.

The closure argument for the domains of the stored terms, which is used here by statement, has
six items. We take them in order.
\begin{itemize}
\item[(1)] This item is unmoved by the relaxation.
\item[(2)] It asks for a continuation in $\mathfrak{K}_P$ at the density of birth. This is the
relaxed continuation clause just proved, combined with part (ii) of
\eqref{9.5:eq-relaxed-heredity-a}.
\item[(3)] It concerns the collar $X^{+}_{\tau}\setminus X_{\tau}$. There it is the second defining
clause of $\mathfrak{G}^{\natural}(\mathfrak{B}_\tau;X_\tau)$ in
Definition~\ref{9.5:def-relaxed-class}.
\item[(4)] It is the ambient containment, proved above.
\item[(5)] Its clause-free conditions are given by part (ii) of
\eqref{9.5:eq-relaxed-heredity-a}.
\item[(6)] It concerns the real domain. This is part (i) of \eqref{9.5:eq-relaxed-heredity-a},
applied at the real point supplied by the real-domain lemma.
\end{itemize}

\emph{The three statements of clause (iii).} These are the inheritance theorem, its corollary
and the claim that follows it. The stage lemma is stated for an arbitrary system, so applying it to
the relaxed domains requires only its hypotheses. Recall the definitions. A \emph{stage-$l$
system} $\mathfrak{D}_l$ assigns to every union $Y$ of cubes an open, $\mathcal{G}$-invariant set
\begin{equation*}
\mathfrak{D}_l(Y)\subset\mathfrak{N}^{(l)}(Y)
\end{equation*}
of pairs $(\mathbb{U},\mathbb{J})$ on $Y$, stable under restriction. The containment hypothesis at
stage $l$ consists of three containment clauses (a)--(c). The stage lemma reads as follows.
\begin{itemize}
\item[--] Hypotheses: a stage index $l\ge1$ not exceeding the number of stages fixed in
Definition~\ref{9.1:not-levelwise-rate}; systems $\mathfrak{D}_i$, $i\le l$, satisfying the
containment hypothesis at stage $l$; the first inductive hypothesis of the inheritance theorem;
and, for $i<l$, each left-over term $\mathbb{C}^{(i)}(S,\mathcal{S}')$ of stage $i$ a stored term
with datum on $\mathfrak{D}_i(S)\times\mathfrak{D}_{\mathcal{S}'}(S)$.
\item[--] Conclusion (b): each term it produces is a stored term with datum on
$\mathfrak{D}_l(X)\times\mathfrak{D}_{\mathcal{S}}(X)$, holomorphic there.
\item[--] Conclusion (c): the terms satisfy the stage-norm bound
$\|\mathbb{C}^{(l)}_k\|^{A,\mathfrak{D}}_{(l)}\le C_0^{\sharp}$, with $C_0^{\sharp}$ the constant
of that lemma (not the auxiliary parameter $C_0$).
\end{itemize}
The proof of clause (a) of the inheritance theorem establishes the containment hypothesis at
stage $n$ for $\mathfrak{D}_i:=\mathfrak{D}^{\flat}_i$ and then applies the stage lemma.

Accordingly, clause (iii) of Proposition~\ref{9.5:prop-relaxed-estimates} is the stage lemma
applied at the system $\mathfrak{D}_i:=\mathfrak{D}^{\flat,\natural}_i$. We check its hypotheses.
\begin{itemize}
\item[--] $\mathfrak{D}^{\flat,\natural}_i$ is a system by
Lemma~\ref{9.5:lem-relaxed-domains-open}.
\item[--] It lies in $\mathfrak{N}^{(i)}$ read with the relaxed ambient convention
\eqref{9.5:eq-relaxed-class-d}. Of this inclusion the dictionary uses, besides the local clauses
at locations of $Y$ which the relaxation does not touch, two forms. The first is
$\mathfrak{N}^{(i)}\subset\mathfrak{K}_P$, the form used in rows 1--8 of the list of uses earlier
in the proof of Proposition~\ref{9.5:prop-relaxed-estimates}; it is the first assertion of
the corresponding display. The second is the base condition, the form used in row 4 of
that list; it is the second assertion of the corresponding display.
\item[--] The containment hypothesis at stage $n$ holds by steps (1)--(4) of the proof of the
inheritance theorem. Their pointwise lines do not involve the continuation class. Their lines on
the continuation class are \eqref{9.5:eq-relaxed-heredity-b} and
\eqref{9.5:eq-relaxed-heredity-c}.
\end{itemize}

In short, the inheritance theorem is an induction on the stage $n$. Its inductive step establishes
the containment hypothesis at stage $n$, which is the containment part above, and applies the
stage lemma, which is the estimate part above.

The proof of clause (C) of the arrest theorem shows what the creation of a term uses. A term
created by $\mathfrak{O}$ is a cluster functional of two kinds of ingredient.
\begin{itemize}
\item[--] Old terms composed with the map of $\mathfrak{O}$. These are holomorphic on the source
domain of $\mathfrak{O}$: by clause (A) and the continuation clause of the arrest theorem, and,
for other terms, by the statements already proved for them.
\item[--] $\mathbb{H}$-operators. These are holomorphic on $\mathfrak{K}_P$.
\end{itemize}
The bounds are suprema over the domain. This uses the third and the first assertions of
the corresponding display, and nothing else.

The first and the third inductive hypotheses of the inheritance theorem are statements on the
domains of the constituents. They are read in relaxed form by the estimate part and the containment
part above, along the time-ordered induction of the arrest theorem. No circularity arises, since
\eqref{9.5:eq-relaxed-heredity-a}--\eqref{9.5:eq-relaxed-heredity-c} use no estimate.

The concluding inclusion of the corollary of the inheritance theorem becomes the following.
$\mathfrak{D}^{\flat,\natural}_n(X)$ contains every pair that satisfies three conditions:
\begin{itemize}
\item[--] it is the restriction of a pair of $\mathfrak{G}^{\natural}(\mathfrak{B}_n;X)$;
\item[--] it is fat on $X$;
\item[--] it obeys the remaining local clauses as stated.
\end{itemize}
The tube used in the claim that follows that corollary consists of global fields of
$\mathfrak{G}$, to which part (i) of \eqref{9.5:eq-relaxed-heredity-a} applies. The inclusion
used for the boundary factor is part (ii) of \eqref{9.5:eq-relaxed-heredity-a} at
$\mathfrak{B}_{n+1}$, restricted to the region $Z_W$ of the boundary factor, together with
stability under restriction.
\end{proof}

\begin{remark}[The one kind of step the relaxation could affect]
The argument above would not apply to a proof that used, on a continuation pair, a bound small in
the coupling at a point outside $Y^{+}_{d}$. The relaxed class of
Definition~\ref{9.5:def-relaxed-class} imposes no such bound there. Among the entries listed in
the proof above, none uses one.
\end{remark}

\begin{lemma}[Closure of the ambient-centre clause along the run]\label{9.5:lem-ambient-closure}
Fix the machine $\mathsf{m}\in\{\mathsf{A},\mathsf{B}\}$. We say that the ambient-centre clause
\eqref{9.5:eq-ambient-centre-a} \emph{holds for} $(U_c;X,\tau)$ \emph{with} $U^{gl}$ when $U^{gl}$ is
the ambient configuration of that clause for a datum about $X$ with real centre $U_c$, born at
$\tau$. A real centre $U_c$ \emph{has the clause for} a
stored term $(F,X,j)$ born at $\tau$ if the clause holds for $(U_c;X,\tau)$ with some $U^{gl}$.
The clause is a condition of membership: it restricts every class
$\mathfrak{T}^{\mathsf{m}}_F[\varsigma](X)$, and hence every supremum $\mathsf{S}^{+}_j[F]$, on both
sides of every inequality of the comparison in which these occur. The restriction is harmless
provided
\begin{itemize}
\item[(A)] every passage from a datum of the class of an output to the data at which the
production of that output reads a stored term preserves the clause, and
\item[(B)] every evaluation of a class at a particular datum is made at a datum with the clause.
\end{itemize}
Assume the following.
\begin{itemize}
\item[(i)] (Regularity of the run.) For every date $d$ and every real $V$ in the support of the
characteristic functions of the effective density, the pair
$(U(\mathfrak{B}_d,V),\mathcal{J}(U(\mathfrak{B}_d,V)))$ lies in the continuation class
$\mathfrak{G}(\mathfrak{B}_d)$; this is the regularity lemma for the points of the run together
with the first clause of the membership property of the arrest theorem. By a further clause of
that membership property, which concerns earlier dates, the same pair lies in
\begin{equation*}
\mathfrak{G}(\mathfrak{B}_{d''})
\subset\mathfrak{K}_{109}(\mathfrak{B}_{d''};\gamma_0^{(k'')})
\qquad\text{for every date } d''\le d,\quad k'':=k(d''),
\end{equation*}
and the first clause of $\mathfrak{K}_{109}(\mathfrak{B}_{d''};\gamma_0^{(k'')})$ is
\cite[(3.35)]{B-CMP99b} at the radius $\alpha_0^{(99)}$ over $\mathfrak{B}_{d''}$. When
$\Lambda_0\neq\emptyset$, both properties hold on the collar $\Lambda_0$ as well.
\item[(ii)] (Structure of chains.) Every point $\mathbf{U}$ of a shape has the form
$\mathbf{U}=e^{i\eta A'}U_c$ over one real field $U_c$, the real centre of the shape's entry;
this centre is the real minimiser of the seed and does not change along a chain, the field $A'$
carrying all moves, real parts included. In the case of the chain construction read here, the
entry is the root $q_0$, a
seed whose real centre is never changed, and every link of a chain is an increment. The slice
$S(\mathbf{U})$ is read in the decomposition $(U,A')=(U_c,\mathfrak{X}(B))$, with centre
$U_c\restriction\widehat{X}$. A stored term $(F,X,j)$, born at $\tau$, that is read by the
production of an output $\Phi$ on $Y$ born at $\tau_\Phi$ satisfies
$\widehat{X}\subset Y$, and $\tau\le\tau_\Phi$, since a stored term precedes in the
order of production every production that reads it. The Schwarz-lemma data
$\mathsf{v}\oplus c^{\mathsf{B}}$ and
$\mathsf{v}_s=e^{is\mathfrak{d}^{net}}(\mathsf{v}\oplus c^{\mathsf{B}})$, arising in clauses (a)
and (b) of each of the two Lipschitz lemmas of the comparison, are chart data of the
machine $\mathsf{B}$ about the centre of $\mathsf{B}$; here $\mathfrak{d}^{net}$ is the generator
of the one-parameter family $s\mapsto\mathsf{v}_s$, and it is real and fixed at the
centre. The local Lipschitz lemma for the chart is read in $\mathsf{B}$.
\item[(iii)] (The instantiation table.) The evaluations of a class at a particular datum made in
the comparison are of three kinds: ($\alpha$) evaluations at data of the polydisc
$\mathbb{D}_j(\widehat{X};\alpha_j^s)$ about $V_\circ=1$ (all block variables equal to $1$);
($\beta$) statements whose proofs fix
a seed as a parameter; ($\gamma$) one evaluation, read at both real centres and at the zero
source. Apart from these, a class enters only through the supremum $\mathsf{S}^{+}_j[F]$, and
then only as a bound for differences at reads over seeds.
\item[(iv)] (The moved centre.) The moved real centre $\mathsf{U}$ at healed cores obeys the two
bounds (a) and (b) of the lemma on the moved real centre, and the coupling ceiling holds.
\end{itemize}
Then (A) and (B) hold, in the following form.
\begin{itemize}
\item[(G1)] (Restriction in space and time.) If the clause holds
for $(U_c;X,\tau)$ with $U^{gl}$, then for every stored term $(F',X',j')$ born at $\tau'\le\tau$
with $X'^{(3/4)}\subset X^{(3/4)}$ the clause holds for
\begin{equation}\tag{G1}\label{9.5:eq-ambient-closure-a}
(U_c\restriction\widehat{X'};\,X',\,\tau')\quad\text{with}\quad
U^{gl}\restriction\Omega_0(\mathfrak{B}_{\tau'}).
\end{equation}
\item[(G2)] (Chains.) Let $\Phi$ be an output on $Y$, born at $\tau_\Phi$, and let $q_0$ be a root
or a seed in the class $\mathfrak{T}^{\natural}_\Phi[7/8](Y)$ defined in
\eqref{9.5:eq-ambient-centre-e}. Every
slice $S(\mathbf{U})$, in the sense of hypothesis \textup{(ii)}, of a point $\mathbf{U}$ of a
shape whose entry is $q_0$, read by a stored term $(F,X,j)$ of the production of $\Phi$, has
the clause; so has every chart datum about the same centre, in particular the Schwarz-lemma
data $\mathsf{v}\oplus c^{\mathsf{B}}$ and $\mathsf{v}_s$.
\item[(G3)] (The trivial centre.) $U_c=\mathbf{1}$ has the clause, with $U^{gl}:=\mathbf{1}$.
Hence every datum of the polydisc $\mathbb{D}_j(\widehat{X};\alpha_j^s)$ about $V_\circ=1$ has
it.
\item[(G4)] (The real points of the run.) For real $V$ in the support of the characteristic
functions of the effective density and $d$ the current date,
$U_c:=U(\mathfrak{B}_d,V)\restriction\widehat{X}$ has the clause for every stored term born at
$\tau\le d$, with $U^{gl}:=U(\mathfrak{B}_d,V)\restriction\Omega_0(\mathfrak{B}_\tau)$.
\item[(G5)] (Real gauges.) If $U_c$ has the clause and $v$ is an $\mathrm{SU}(N)$-valued gauge
transformation on $\widehat{X}$, then $U_c^{v}$ has the clause.
\item[(G6)] (The one re-centring.) Let $\Phi$ be born in the
operation $\mathbf{R}$ at step $k$ on $Y$, and let $q_0=e^{i\eta A_e}U_c$ be its root on
$\boldsymbol{\Omega}(R)$ for a region $R\supseteq\widehat{Y}$, in
$\mathfrak{T}^{\natural}_\Phi[7/8](Y)$; here $A_e$ is the field in this decomposition of the root,
$\mathfrak{S}_e$ and $D_c$ are the box tower and the block data of the construction of the moved
real centre at healed cores, and $\bar d_k$ denotes the date of the last stage of step $k$. Put
\begin{equation}\tag{G6}\label{9.5:eq-ambient-closure-b}
\mathsf{U}:=U_{\mathfrak{S}_e}(D_c).
\end{equation}
Under hypothesis \textup{(iv)}, $\mathsf{U}$ has the
clause for every stored term $(F,X,j)$ with $\widehat{X}\subset Y$, born at $\tau\le\bar d_k$ at
the density $k(\tau)=j$; thus, for $F=\mathbb{C}^{(i')}_{k_a}(X)$, with $j:=k_a$.
\item[(G6$'$)] (A stored term inside the core.) In the setting of \textup{(G6)}, let
$\widehat{X}\subset\mathfrak{K}$ and let $j\le k(\tau)$ be arbitrary (this covers the stored
left-over terms attached at a level below their top level), and let $\mathcal{V}$ be the core
component of a supported point. Then $\mathsf{U}$ has the clause for $(F,X,j)$.
\item[(G7)] (Instantiations.) Every evaluation of a class at a particular datum is made at a datum
with the clause: the evaluations of kind ($\alpha$) by \textup{(G3)}, those of kind ($\beta$)
by \textup{(A)}, and the evaluation of kind ($\gamma$) by \textup{(G4)}. The one evaluation that
takes a norm over the domain of a current slot, namely the exact competitor bound of the
comparison, is of kind ($\alpha$):
for every point $\mathbf{U}$ of the complex neighbourhood $\mathcal{U}^{c}_j(Y)$ it evaluates the
stored term at data of the polydisc $\mathbb{D}_j$ about $V_\circ=1$, the competitor living on one
lattice and never being a seed. Its centre is therefore $\mathbf{1}$, whatever domain the
current slot carries, and no compatibility of the clause with the continuation class of that
domain is needed there.
\end{itemize}
\end{lemma}

\begin{proof}
\emph{Proof of \textup{(G1)}.} Let the clause hold for $(U_c;X,\tau)$ with $U^{gl}$, and let
$(F',X',j')$ be born at $\tau'\le\tau$ with $X'^{(3/4)}\subset X^{(3/4)}$. The domains $\Omega_0$
of the determining sets are nested in the date,
$\Omega_0(\mathfrak{B}_{\tau'})\subset\Omega_0(\mathfrak{B}_\tau)$, so the restriction
$U^{gl}\restriction\Omega_0(\mathfrak{B}_{\tau'})$ is an $\mathrm{SU}(N)$-valued configuration
on $\Omega_0(\mathfrak{B}_{\tau'})$. The dates $d''\le\tau'$ are among the dates $d''\le\tau$,
so the conditions \cite[(3.35)]{B-CMP99b} at radius $\alpha_0^{(99)}$ over $\mathfrak{B}_{d''}$
required for $\tau'$ are among those that $U^{gl}$ satisfies by assumption. Finally $U_c$ and
$U^{gl}$ agree on the bonds of $X^{(3/4)}$, hence on the bonds of $X'^{(3/4)}\subset
X^{(3/4)}$, and there $U_c\restriction\widehat{X'}$ and
$U^{gl}\restriction\Omega_0(\mathfrak{B}_{\tau'})$ are the restrictions of these two
configurations.

\emph{Proof of \textup{(G2)}.} By hypothesis \textup{(ii)}, every point of a shape is
$\mathbf{U}=e^{i\eta A'}U_c$ over one real field, the real centre of the shape's entry, and this
centre does not change along a chain, all moves (real parts included) being carried by $A'$; in
the case of the chain construction read here the entry is the root $q_0$, whose real centre is
never changed, and every
link is an increment. Write $U_c^{\Phi}$ for the real centre of $q_0$. Since the slice
$S(\mathbf{U})$ is read in the decomposition $(U,A')=(U_c,\mathfrak{X}(B))$ with centre
$U_c\restriction\widehat{X}$, the centre of each slice is $U_c^{\Phi}\restriction\widehat{X}$.
Because $q_0\in\mathfrak{T}^{\natural}_\Phi[7/8](Y)$, the definition
\eqref{9.5:eq-ambient-centre-e} of this class gives the clause for $(U_c^{\Phi};Y,\tau_\Phi)$
with some $U^{gl}$. As to geometry, hypothesis \textup{(ii)} gives
$\widehat{X}\subset Y$, so that $X^{(3/4)}\subset Y\subset Y^{(3/4)}$. As to time,
the stored term is stored before the production that reads it, so $\tau\le\tau_\Phi$.
Clause \textup{(G1)}, applied with $(U_c^{\Phi};Y,\tau_\Phi)$ in place of $(U_c;X,\tau)$ and
with $(F,X,j)$ in place of $(F',X',j')$, now gives the clause for
$(U_c^{\Phi}\restriction\widehat{X};X,\tau)$, which is the assertion for every slice. The same
argument applies to every chart datum about the same centre. The data
$\mathsf{v}\oplus c^{\mathsf{B}}$ and
$\mathsf{v}_s=e^{is\mathfrak{d}^{net}}(\mathsf{v}\oplus c^{\mathsf{B}})$ are, by hypothesis
\textup{(ii)}, chart data of the machine $\mathsf{B}$ about the centre of $\mathsf{B}$, the
generator $\mathfrak{d}^{net}$ being real and fixed at the centre and the local Lipschitz lemma
for the chart being
read in $\mathsf{B}$; the argument just given, carried out in $\mathsf{B}$, gives the clause for
them.

\emph{Proof of \textup{(G3)}.} The constant configuration $\mathbf{1}$ obeys
\cite[(3.35)]{B-CMP99b} with $u=1$ and $A=0$ over every set, in particular over
$\mathfrak{B}_{d''}$ for every date $d''\le\tau$; and $U_c=\mathbf{1}$ agrees with
$U^{gl}:=\mathbf{1}$ on every bond. For the data of the polydisc
$\mathbb{D}_j(\widehat{X};\alpha_j^s)$ about $V_\circ=1$: the variational problem in which all
block values equal $1$ has minimiser $\mathbf{1}$, so the real centre of every such datum is
$\mathbf{1}$, and these data are data of the class by the inclusion of the polydisc into the
class of chart data. The first part of \textup{(G3)} therefore applies to them.

\emph{Proof of \textup{(G4)}.} By hypothesis \textup{(i)},
$(U(\mathfrak{B}_d,V),\mathcal{J}(U(\mathfrak{B}_d,V)))\in\mathfrak{G}(\mathfrak{B}_d)$, and by
the further clause of the membership property in hypothesis \textup{(i)} this pair lies in
\begin{equation*}
\mathfrak{G}(\mathfrak{B}_{d''})\subset\mathfrak{K}_{109}(\mathfrak{B}_{d''};\gamma_0^{(k'')})
\qquad\text{for every date } d''\le d,
\end{equation*}
whose first clause is \cite[(3.35)]{B-CMP99b} at the radius $\alpha_0^{(99)}$ over
$\mathfrak{B}_{d''}$. When $\Lambda_0\neq\emptyset$, the two memberships on the
collar are part of hypothesis \textup{(i)} as well. Let the stored term be born at $\tau\le d$.
Every date $d''\le\tau$ satisfies $d''\le d$, so $U(\mathfrak{B}_d,V)$, and hence its
restriction $U^{gl}:=U(\mathfrak{B}_d,V)\restriction\Omega_0(\mathfrak{B}_\tau)$, obeys
\cite[(3.35)]{B-CMP99b} at radius $\alpha_0^{(99)}$ over $\mathfrak{B}_{d''}$; it is
$\mathrm{SU}(N)$-valued because $V$ is real. The centre
$U_c=U(\mathfrak{B}_d,V)\restriction\widehat{X}$ and $U^{gl}$ are restrictions of one
configuration, and so agree on the bonds of $X^{(3/4)}$.

\emph{Proof of \textup{(G5)}.} Let the clause hold for $U_c$ with $U^{gl}$, and let $v$ be an
$\mathrm{SU}(N)$-valued gauge transformation on $\widehat{X}$. Define $\tilde v:=v$ on the sites
of $X^{(3/4)}$ and $\tilde v:=1$ elsewhere. Then $(U^{gl})^{\tilde v}$ is
$\mathrm{SU}(N)$-valued and $(U^{gl})^{\tilde v}=U_c^{v}$ on the bonds of $X^{(3/4)}$, since
there $U^{gl}=U_c$ and $\tilde v=v$. The condition \cite[(3.35)]{B-CMP99b} is gauge invariant:
if it holds for $U^{gl}$ with the gauge $u$, it holds for $(U^{gl})^{\tilde v}$ with
$u\tilde v^{-1}$. Hence $U_c^{v}$ has the clause, with $(U^{gl})^{\tilde v}$.

\emph{Proof of \textup{(G7)}.} By hypothesis \textup{(iii)}, the evaluations of a class at a
particular datum are of the three kinds ($\alpha$), ($\beta$), ($\gamma$). Those of kind
($\alpha$) are made at data of the polydisc $\mathbb{D}_j$ about $V_\circ=1$, which have the
clause by \textup{(G3)}. Those of kind ($\beta$) fix a seed as a parameter, so the data at which
they evaluate are reached from a seed by the passages of \textup{(A)}. The evaluation of kind
($\gamma$) is read at both real centres and at the zero source, and the real centres have the
clause by \textup{(G4)}. Elsewhere a class enters only through the supremum
$\mathsf{S}^{+}_j[F]$, as a bound for differences at reads over seeds, which is again covered
by \textup{(A)}. The exact competitor bound named in \textup{(G7)} evaluates the stored term, for
every point $\mathbf{U}$ of $\mathcal{U}^{c}_j(Y)$, at data of the polydisc $\mathbb{D}_j$
about $V_\circ=1$; the competitor lives on one lattice and is never a seed, so this evaluation
is of kind ($\alpha$), its centre is $\mathbf{1}$ whatever domain the current slot carries, and
\textup{(G3)} applies to it.

The proofs of \textup{(G6)} and \textup{(G6$'$)} are given separately, immediately after this
proof.
\end{proof}

\begin{lemma}[The ambient-centre clause at re-gauged centres]\label{9.5:lem-regauged-centres}
Fix one of the two compared runs $\mathsf{m}\in\{\mathsf{A},\mathsf{B}\}$
(Definition~\ref{9.6:def-comparison-chain}); we suppress the superscript $\mathsf{m}$ except in the
conclusion. Write $\mathrm{AL}$ for the frame used on the side of $U$ in the comparison, and
$\natural$ for the re-gauged frame; their data are $D^{\mathrm{AL}}$ and $D^{\natural}$. A datum is a
function on the bonds of $\mathfrak{B}_e$; we write $D_c^{\mathrm{AL}}$, $D_c^{\natural}$ for the data
of the two frames at the real centre seed $U_c$. For a datum $D$, write $U_{\mathfrak{S}_e}(D)$ for the
configuration with datum $D$ critical for the box tower $\mathfrak{S}_e$; it is real when $D$ is the
datum at a real seed. Assume the following.
\begin{itemize}
\item[(a)] The re-gauged frame satisfies the frame conditions on which the construction of the moved real centre $\mathsf{U}$ depends, with the frame size $\mathsf{s}_L$ replaced by
\begin{equation}\label{9.5:eq-regauged-centres-1}
\mathsf{s}_L^{\natural}:=1.1\,\mathsf{s}_L+80\,d\,K_0^{\natural}\,\tilde\alpha ,
\end{equation}
where $d=4$ is the dimension and $K_0^{\natural}$, $\tilde\alpha$ are the two constants with which
the re-gauged frame is defined.
\item[(b)] There is a site function $\tilde\sigma$ with values in the complexification
$\mathrm{SL}(N,\mathbb{C})$ of $\mathrm{SU}(N)$, depending holomorphically on the seed, such that
\begin{equation}\label{9.5:eq-regauged-centres-2}
U_{\mathfrak{S}_e}(D^{\natural})=U_{\mathfrak{S}_e}(D^{\mathrm{AL}})^{\tilde\sigma}.
\end{equation}
\item[(c)] The gauge-covariance statement (G5) of Lemma~\ref{9.5:lem-ambient-closure}, displayed
in \eqref{9.5:eq-ambient-closure-a}, and the two instantiation statements (G6) and (G6$'$) of the
same lemma, displayed in \eqref{9.5:eq-ambient-closure-b}, hold.
\end{itemize}
Hypotheses (a) and (b) are the two assertions of the statement on the re-gauged frame (that it
obeys the frame conditions and that it re-gauges the critical configurations), taken here by
statement.
Then the centres
\begin{equation*}
\mathsf{U}^{\mathsf{m}}:=U_{\mathfrak{S}_e^{\mathsf{m}}}\bigl(D_c^{\natural,\mathsf{m}}\bigr)
\end{equation*}
of the two runs satisfy the ambient-centre clause \eqref{9.5:eq-ambient-centre-a} wherever (G6) and (G6$'$) give that clause to the moved real centre $\mathsf{U}$ of the healed core.
\end{lemma}

\begin{proof}
We first recall the definition of the re-gauged data. Let $u_0^{\mathsf{m}}$ be the gauge
transformation that fixes the re-gauged frame in its definition. It is holomorphic in the seed and takes values in $\mathrm{SU}(N)$ at real seeds. Put $\sigma^{\mathsf{m}}(y):=u_0^{\mathsf{m}}(y)^{-1}$ for $y\in\Lambda^{(k)}\setminus Z''_k$, and $\sigma^{\mathsf{m}}(y):=\mathbf{1}$ at every other end-point of a bond of $\mathfrak{B}_e$; we write $\sigma$ for $\sigma^{\mathsf{m}}$. For a bond $b=\langle b_-,b_+\rangle$ of $\mathfrak{B}_e$,
\begin{equation*}
D^{\natural}(b):=\sigma(b_-)\,D^{\mathrm{AL}}(b)\,\sigma(b_+)^{-1}.
\end{equation*}
At the real centre seed $U_c$, at which the data $D_c^{\mathrm{AL}}$ and $D_c^{\natural}$ are taken,
$\sigma$ takes values in $\mathrm{SU}(N)$.

Put $\mathsf{U}^{\natural}:=U_{\mathfrak{S}_e}(D_c^{\natural})=\mathsf{U}^{\mathsf{m}}$ and
$\mathsf{U}^{\mathrm{AL}}:=U_{\mathfrak{S}_e}(D_c^{\mathrm{AL}})$. Since $D_c$ is the datum at the real
seed $U_c$, both sides of \eqref{9.5:eq-regauged-centres-2} at $D_c$ are real critical
configurations; by (b) they lie in one gauge orbit. Hence at this seed $\tilde\sigma$ takes values
in $\mathrm{SU}(N)$ and is a lattice gauge transformation, and
\begin{equation}\label{9.5:eq-regauged-centres-3}
\mathsf{U}^{\natural}=\bigl(\mathsf{U}^{\mathrm{AL}}\bigr)^{\tilde\sigma}.
\end{equation}
The configuration $\mathsf{U}^{\mathrm{AL}}$ is the moved real centre $\mathsf{U}$ of the healed core. It is formed in the frame $\mathrm{AL}$, which is the frame the construction of $\mathsf{U}$ itself uses on the side of $U$. We now give two derivations of the clause.

\emph{First derivation.} Wherever (G6) and (G6$'$) give the ambient-centre clause \eqref{9.5:eq-ambient-centre-a} to $\mathsf{U}$, they give it to $\mathsf{U}^{\mathrm{AL}}$, since the two coincide. Apply the gauge-covariance statement (G5) of Lemma~\ref{9.5:lem-ambient-closure}, displayed in \eqref{9.5:eq-ambient-closure-a}, with the gauge transformation
\begin{equation*}
v:=\tilde\sigma\restriction\widehat{X}.
\end{equation*}
By \eqref{9.5:eq-regauged-centres-3}, this yields the clause for $\mathsf{U}^{\natural}$, that is, for $\mathsf{U}^{\mathsf{m}}$.

\emph{Second derivation.} As stated with the construction of the moved real centre $\mathsf{U}$,
that construction uses the frame only through the four frame conditions. By hypothesis (a), these conditions hold for the re-gauged frame with $\mathsf{s}_L$ replaced by $\mathsf{s}_L^{\natural}$ of \eqref{9.5:eq-regauged-centres-1}. Consequently both assertions of the lemma on the moved real centre hold in the frame $\natural$, with $\mathsf{s}_L$ replaced by $\mathsf{s}_L^{\natural}$.

Statements (G6) and (G6$'$) therefore apply verbatim in the frame $\natural$, provided the gluing-size constant $e_{\mathsf{U}}$ is re-maximised at $\mathsf{s}_L^{\natural}$ inside the corresponding ceiling. The re-maximised constant is still of the form $C(\alpha_{0,k}+\alpha_{1,k})$, with $C$ independent of the coupling. Hence the left side of the corresponding display keeps the same form, and the ceiling is read as before. This gives \eqref{9.5:eq-ambient-centre-a} for $\mathsf{U}^{\mathsf{m}}$ directly.

The two derivations give the same clause, in each run $\mathsf{m}\in\{\mathsf{A},\mathsf{B}\}$.
\end{proof}

\begin{remark}[The ambient-centre clause carries no margin]\label{9.5:rem-no-margin}
Let $U^{gl}$ be the ambient real field and $\mathsf{U}^{gl}$ the corresponding extension of the moved
real centre, as in the ambient-centre clause~\eqref{9.5:eq-ambient-centre-a}.
\begin{itemize}
\item[(a)] Outside the region in which a quantitative local clause of the parent's class is
available, \eqref{9.5:eq-ambient-centre-a} requires only that $\mathsf{U}^{gl}$ be a gauge copy of
$U^{gl}$, unperturbed. The clause~\eqref{9.5:eq-ambient-centre-a} carries no margin and no
radius of its own, and no sequence of successively smaller margins or radii is generated
when the clause is passed to re-centred data. The following two classes carry the same
clause~\eqref{9.5:eq-ambient-centre-a}:
\begin{itemize}
\item the class of data about the terms that a production reads;
\item the class of data at which, in Lemma~\ref{9.5:lem-ambient-closure}, a class is
instantiated.
\end{itemize}
A re-centred datum belongs to the first class and never to the second.
\item[(b)] In the strong form of the clause, in which $U^{gl}$ is required to obey the direct
clauses (the clauses of the class, imposed on $U^{gl}$ itself),
\eqref{9.5:eq-ambient-centre-a} is a genuine restriction; this is part~(c) of the
counterexample to the strong form of the ambient-centre clause. The clause is used in this book
in the weak form, in which $U^{gl}$ is not required to obey the direct clauses.

Whether, in the weak form, \eqref{9.5:eq-ambient-centre-a} holds automatically is not decided
here. A heuristic suggests that it might, at the price of a lower bound on $M$. A hole of the hull
$\widehat{X}$, that is, a bounded component of its complement, has side at least $u=M\ell_j$,
where $\ell_j$ is the unit of level $j$. A holonomy of order one spread over an area
$M^{2}\ell_j^{2}$ costs, on each cube on which the class is formed, potentials of order $M^{-2}$
times the size of that cube. These potentials lie below $\alpha_0^{(99)}$ once
$M^{2}\alpha_0^{(99)}$ is large, which would impose a floor on $M$. A proof along these lines,
however, would require an extension through all four skeleta of the cube complex. No such
extension is constructed here. It is not asserted that the weak form holds automatically, and no
argument of this book uses such a statement.
\end{itemize}
\end{remark}

\begin{proof}
(a) The requirement that $\mathsf{U}^{gl}$ be a gauge copy of $U^{gl}$ involves no numerical
parameter, so no radius or margin enters it that could shrink when the clause is passed from a
datum to a re-centred datum. The clause imposed on the two classes is therefore the same, and a
re-centred datum, satisfying exactly the clause of the first class, is a datum of that class and
not a datum of the second class.

(b) The assertion that the strong form is a genuine restriction is the content of part~(c) of the
counterexample to the strong form of the ambient-centre clause.
\end{proof}

\begin{remark}[Emptiness of the never-averaged collar]\label{9.5:rem-empty-collar}
Let a label be \emph{aligned}, in the sense of the definition of common and aligned labels: its
first index, that is, the level
of its first large-field region, exceeds the threshold $n_c$, and its history satisfies the
further condition imposed there. Then, for $\gamma$ small, in both runs $\mathsf{A}$ and
$\mathsf{B}$ of Definition~\ref{9.6:def-comparison-chain} the
never-averaged collar $\Lambda_0$ of Definition~\ref{9.5:def-fattened-sets} is empty. Consequently:
\begin{itemize}
\item[(a)] every leaf of an aligned label, a leaf being as in
Definition~\ref{9.5:def-fattened-sets}, satisfies the scope condition of
Definition~\ref{9.5:def-fattened-sets};
\item[(b)] the level-zero collar lemma of the arrest theorem, which asserts the dated level-$0$
clause at the points of the dated collar $\operatorname{Col}(d)$, is not used on aligned labels.
Its parts (i) and (ii) enter the sandwich clause and the heredity clause of
Lemma~\ref{9.5:lem-relaxed-heredity}, and its part (iv) enters the clause (G4) of
\eqref{9.5:eq-ambient-closure-a}, in each case only at points of the collar. It is therefore used
only on labels with $\Lambda_0\neq\emptyset$, and, under the scope condition, only for the part of
the real centre $U_c$ at points of the collar, which by \eqref{9.5:eq-fattened-sets-a} lie at
sup-distance more than $2\eta$ from the fattened set $X^{(3/4)}$ of every leaf $(F,X,j)$
satisfying that condition.
\end{itemize}
In this section the arrest theorem enters only through: the lemma on the order
$\preccurlyeq$ of dated sets; parts (a), (d), (e1), (e2) and (e3) of the lemma of the arrest
theorem with parts (a)--(e4) assumed in Lemma~\ref{9.5:lem-relaxed-heredity}, part (e3) together
with the nesting statement; the first clause of part (e4) of that lemma, together with the lemma of
the arrest theorem that is read with that clause;
the second clause of part (e4), only inside \eqref{9.5:eq-relaxed-heredity-c} and inside the
proposition of this section that uses that display; and the level-zero collar lemma, only as
described in (b).
The uniqueness theorem compares the two runs on aligned labels only. Non-aligned labels
are not compared. Hence, wherever that theorem uses the ambient-centre clause of
Lemma~\ref{9.5:lem-ambient-closure}, the dated level-$0$ clause on the collar is vacuous.
Thus, of the arrest theorem, the ambient-centre clause as used by the uniqueness theorem rests only
on parts (e1)--(e3) of the lemma above with the nesting statement and on the lemma read with the
first clause of part (e4), together with one further statement of the chapter of the arrest
theorem, read in its
relaxed ($\natural$) form; the level-zero collar lemma is vacuous there, no other statement of that
chapter is used, and the steps with two windows remain outside the scope declared in the arrest
theorem.
\end{remark}

\begin{proof}
\emph{The first index of an aligned label is at least $2$.} By the definition of aligned
labels, every zone present in an aligned label has index at least $n_c$ (zones and their indices
being as in that definition). The threshold $n_c$ is fixed there so that the alignment condition
\begin{equation}\label{9.5:eq-empty-collar-1}
C^{\natural}_{\mathrm{cl}}\,(\log x_n)^{c_2}\,\vartheta_2^{\,n/4}\le 1
\qquad\text{for every integer } n\ge n_c
\end{equation}
holds, in which $C^{\natural}_{\mathrm{cl}}$ and $\vartheta_2$ are constants fixed in that
definition, $c_2$ is an exponent fixed in that definition, and $x_n=g_n^{-2}$ is the inverse
coupling at level $n$.
At $n=0$ the left member of \eqref{9.5:eq-empty-collar-1} equals
$C^{\natural}_{\mathrm{cl}}(\log x_0)^{c_2}$, and this exceeds $1$ for $\gamma$ small. So
\eqref{9.5:eq-empty-collar-1} fails at $n=0$, and therefore $n_c\ge 1$. An aligned label has
first index $>n_c\ge 1$, so its first index is at least $2$.

\emph{The run $\mathsf{A}$.} In $\mathsf{A}$, $\Omega_1$ is the domain of the first step,
\cite[(1.10)]{B-CMP119}. The functions on $\Omega_1^c$ are those defined in the first step,
\cite[p.~255]{B-CMP119}. The large-field region created by the first step is
$Z_1=\Lambda_1^c$, the complement of the small-field region $\Lambda_1$ of the first step, and it
satisfies $Z_1\supseteq\Omega_1^c$.

The index count in the definition of aligned labels starts at $1$ in $\mathsf{A}$, because
index $0$ is reserved for the extra step of $\mathsf{B}$. So the regions made by the first
step of $\mathsf{A}$ are exactly those of level $1$, that is, of index $1$. The two index
conventions therefore agree.

By the first paragraph, an aligned label has no region of index $1$. Hence $Z_1=\emptyset$,
and from $\Omega_1^c\subseteq Z_1$ we get $\Omega_1^c=\emptyset$. Thus $\Omega_1$ is the
whole torus. The set $\widetilde{\Omega}_1$ of Definition~\ref{9.5:def-fattened-sets} is a
subset of the torus, so
\begin{equation*}
\Lambda_0=\widetilde{\Omega}_1\setminus\Omega_1=\emptyset .
\end{equation*}

\emph{The run $\mathsf{B}$.} The finest level of $\mathsf{B}$ is $-1$, and its first step
is the step of index $0$. An aligned label is in particular common, that is, its first index is
at least $1$: indeed its first index is at least $2$. So it has no region of index $0$. The
first-step domain of
$\mathsf{B}$ is then its whole torus, and, exactly as for $\mathsf{A}$, $\Lambda_0=\emptyset$.

\emph{Proof of (a).} The scope condition of Definition~\ref{9.5:def-fattened-sets} requires
that the sup-distance from $X^{(3/4)}$ to $\Lambda_0$ exceed $2\eta$. The distance to the empty
set is $+\infty$, so this holds for every leaf $(F,X,j)$ of an aligned label, in either
run.

\emph{Proof of (b).} By Definition~\ref{9.5:def-fattened-sets}, the dated collar satisfies
$\operatorname{Col}(d)=\Lambda_0\setminus Z^{\le d}\subset\Lambda_0$ at every date $d$. On an
aligned label $\Lambda_0=\emptyset$, so $\operatorname{Col}(d)=\emptyset$ at every date.

The level-zero collar lemma asserts the dated level-$0$ clause at points of
$\operatorname{Col}(d)$. Its parts (i) and (ii) are used in the sandwich clause and the heredity
clause of Lemma~\ref{9.5:lem-relaxed-heredity} only at collar points, and its part (iv) is
used in the clause (G4) of \eqref{9.5:eq-ambient-closure-a} likewise only at collar points. There
are no collar points on an aligned label, so on such labels these uses are empty and the
level-zero collar lemma is an input neither of Lemma~\ref{9.5:lem-relaxed-heredity} nor of
Lemma~\ref{9.5:lem-ambient-closure}. On a label with $\Lambda_0\neq\emptyset$, every collar point
lies in $\Lambda_0$, since $\operatorname{Col}(d)\subset\Lambda_0$; under the scope condition
\eqref{9.5:eq-fattened-sets-a} it is therefore at sup-distance more than $2\eta$ from the fattened
set $X^{(3/4)}$ of every leaf $(F,X,j)$ satisfying it, and only the part of the real centre
$U_c$ at such points is involved. Since the uniqueness theorem reads
Lemma~\ref{9.5:lem-ambient-closure} on aligned labels only, the last assertions of the remark
follow.
\end{proof}

\begin{remark}\label{9.5:prop-flat-ring}
The statement of this proposition is not written out here; parts of its proof are printed below.
\end{remark}

\begin{proof}[Proof of Proposition~\ref{9.5:prop-flat-ring}, continued: the holonomy round the ring
against the plaquette product]
\label{9.5:prf-ring-holonomy}
We keep the setting of Proposition~\ref{9.5:prop-flat-ring}: $G=\mathrm{SU}(N)$, $h\in\mathfrak{g}$
hermitian, traceless and non-zero, with largest eigenvalue $\varepsilon_h>0$; $(F,X,j)$ is the term
born at the date $\tau$ of Notation~\ref{9.5:not-dates-weights-norms}; $u=M\ell_j$; and the ring
$\mathsf{R}$ and the box $\mathsf{N}$ are as in the corresponding display. The norms $|\cdot|$
and $\|\cdot\|$ are those of Notation~\ref{9.5:not-dates-weights-norms}, $\|\cdot\|$ being the
operator norm.

We write paths as words in oriented bonds, concatenated from left to right, and $b^{-1}$ for the
bond $b$ with reversed orientation. A configuration $\mathbb{W}$ assigns an invertible matrix
$\mathbb{W}(b)$ to each oriented bond, with $\mathbb{W}(b^{-1})=\mathbb{W}(b)^{-1}$. For a path
$c=b_1b_2\cdots b_r$ we put $\mathbb{W}(c):=\mathbb{W}(b_1)\mathbb{W}(b_2)\cdots\mathbb{W}(b_r)$,
with $\mathbb{W}(c)=\mathbf{1}$ for the empty path. Then $\mathbb{W}(cc')=\mathbb{W}(c)\mathbb{W}(c')$
and $\mathbb{W}(c^{-1})=\mathbb{W}(c)^{-1}$. In particular a path of the form $cc^{-1}$ has
$\mathbb{W}(cc^{-1})=\mathbf{1}$, and two words that differ by insertion or deletion of such a
factor have the same value. We say that two paths are equal \emph{as products of bond variables}
if they are equal after such cancellations.

\medskip
\noindent\emph{Step 1: the holonomy of the datum round the ring.}
Fix a lattice value $x^0$ with $0<x^0<u$. Let
\begin{equation*}
\mathsf{Q}:=\{x:\ \max(|x_1|,|x_2|)\le 6u,\ x_3=x_4=x^0\},
\end{equation*}
a planar square. Let $\gamma$ be the lattice square $\{\max(|x_1|,|x_2|)=5.5u\}$ in $\mathsf{Q}$,
oriented so that its upper side runs in the direction $-e_1$. Every point of $\gamma$ satisfies
$5u\le\max(|x_1|,|x_2|)\le 6u$ and $0<x_3=x_4=x^0<u$, so $\gamma\subset\mathsf{R}$, and
$\mathsf{R}\subset X$ by the corresponding display. Hence $\gamma\subset X$.

Along $\gamma$ all the factors of $\mathbf{V}(B_{\times})(\gamma)$ commute. The ordered product is
therefore the exponential of the signed sum of the exponents, the sign being $+$ or $-$ according
as $\gamma$ traverses a bond $b$ in its own orientation or against it. The forms $X_{\times}$ and
$A_{\times}$ of Proposition~\ref{9.5:prop-flat-ring} differ on $\gamma$ by the lattice gradient of
$\lambda_{\times}$. Its signed sum round the closed loop $\gamma$ vanishes, because it telescopes.
Inserting the form $A_{\times}$ and the profile $\mathsf{g}$ of
Proposition~\ref{9.5:prop-flat-ring}, with $\tau_{\times}$ as in the corresponding display, we
obtain
\begin{equation}\label{9.5:eq-ring-holonomy-1}
\begin{aligned}
\mathbf{V}(B_{\times})(\gamma)
&=\exp\Bigl(i\eta\sum_{b\in\gamma}\pm X_{\times}(b)\Bigr)
=\exp\Bigl(i\eta\sum_{b\in\gamma}\pm A_{\times}(b)\Bigr)\\
&=\exp\Bigl(-i\cdot i\,\tau_{\times}\,\ell_j^{-1}\,h\int\mathsf{g}\Bigr)
=e^{Ph},\qquad P:=6M\tau_{\times}.
\end{aligned}
\end{equation}
Here the third equality inserts the expression of $A_{\times}$ in terms of $\tau_{\times}$, $h$ and
$\mathsf{g}$ from Proposition~\ref{9.5:prop-flat-ring}, and $-i\cdot i=1$; the last equality uses
$\int\mathsf{g}=6M\ell_j=6u$ for the profile $\mathsf{g}$, so that
$\tau_{\times}\ell_j^{-1}\int\mathsf{g}=6M\tau_{\times}=P$.
Since $h$ is hermitian with eigenvalue $\varepsilon_h$, the matrix $e^{Ph}$ has the eigenvalue
$e^{P\varepsilon_h}$. Hence $\mathbf{V}(B_{\times})(\gamma)-\mathbf{1}$ has the eigenvalue
$e^{P\varepsilon_h}-1$. The norm $|\cdot|$ dominates the operator norm $\|\cdot\|$
(Notation~\ref{9.5:not-dates-weights-norms}), an operator norm dominates the spectral radius, and
$e^{x}-1\ge x$ for real $x$. Therefore
\begin{equation}\label{9.5:eq-ring-holonomy-2}
|\mathbf{V}(B_{\times})(\gamma)-\mathbf{1}|
\ge\|\mathbf{V}(B_{\times})(\gamma)-\mathbf{1}\|
\ge e^{P\varepsilon_h}-1\ge P\varepsilon_h .
\end{equation}
All factors along $\gamma$ commute, so this holonomy does not depend on the starting point of
$\gamma$. We read $\gamma$ from its lower-left corner.

\medskip
\noindent\emph{Step 2: the lattice non-abelian Stokes formula on the square.}
Let $D$ be the grid of plaquettes of $\mathsf{Q}$ that fills $\gamma$. It is an $n\times n$ grid
with $n\eta=11u$, since $\gamma$ has side $11u$. Let $o$ be its lower-left corner. We give the
vertices of $D$ the coordinates $(s,t)\in\{0,\dots,n\}^2$ of the point $o+\eta(se_1+te_2)$. Let
\begin{equation*}
\mathbb{W}=e^{i\eta\tilde A}\tilde U
\end{equation*}
be any configuration on the bonds of $D$. We use the following notation.
\begin{itemize}
\item $c(s,t)$, for $0\le s,t\le n-1$, is the plaquette with lower-left corner $(s,t)$.
\item $\partial c(s,t)$ is its boundary, read counter-clockwise from $(s,t)$.
\item $\mathsf{b}_s$ is the path from $o$ along the bottom side to $(s,0)$.
\item $\mathsf{a}_{s,t}$ is the path from $(s,0)$ straight up to $(s,t)$.
\end{itemize}
We put
\begin{equation*}
\mathsf{L}(s,t):=\mathbb{W}(\mathsf{b}_s\mathsf{a}_{s,t})\,\mathbb{W}(\partial c(s,t))\,
\mathbb{W}(\mathsf{b}_s\mathsf{a}_{s,t})^{-1}.
\end{equation*}

(i) \emph{One column.} Fix $s$. For $-1\le t\le n-1$ let $R_t:=[s,s+1]\times[0,t+1]$. We use the
following pieces of path.
\begin{itemize}
\item $\mathrm{bot}$ is the bond from $(s,0)$ to $(s+1,0)$.
\item $\mathrm{right}(t'\to t'')$ is the path from $(s+1,t')$ up to $(s+1,t'')$.
\item $\mathrm{left}(t\to t+1)$ is the bond from $(s,t)$ to $(s,t+1)$.
\item $\mathrm{top}_t$ is the bond from $(s,t)$ to $(s+1,t)$.
\end{itemize}
Reading the boundaries counter-clockwise, $\partial R_{t-1}$ from $(s,0)$ and $\partial c(s,t)$
from $(s,t)$, we have
\begin{align*}
\partial R_{t-1}&=\mathrm{bot}\cdot\mathrm{right}(0\to t)\cdot\mathrm{top}_t^{-1}\cdot
\mathsf{a}_{s,t}^{-1},\\
\partial c(s,t)&=\mathrm{top}_t\cdot\mathrm{right}(t\to t+1)\cdot\mathrm{top}_{t+1}^{-1}\cdot
\mathrm{left}(t\to t+1)^{-1}.
\end{align*}
For $t=0$ the first line reads $\mathrm{bot}\cdot\mathrm{bot}^{-1}$, since $\mathrm{top}_0=\mathrm{bot}$
and $\mathrm{right}(0\to0)$ and $\mathsf{a}_{s,0}$ are empty. Thus
$\mathbb{W}(\partial R_{-1})=\mathbf{1}$.

Multiplying, the factor $\mathrm{top}_t^{-1}\mathsf{a}_{s,t}^{-1}\mathsf{a}_{s,t}\mathrm{top}_t$
cancels. We then use
\begin{equation*}
\mathrm{right}(0\to t)\,\mathrm{right}(t\to t+1)=\mathrm{right}(0\to t+1),\qquad
\mathrm{left}(t\to t+1)^{-1}\mathsf{a}_{s,t}^{-1}=\mathsf{a}_{s,t+1}^{-1}.
\end{equation*}
This gives
\begin{equation*}
\begin{aligned}
\mathbb{W}(\partial R_{t-1})\,\mathbb{W}(\mathsf{a}_{s,t})\mathbb{W}(\partial c(s,t))
\mathbb{W}(\mathsf{a}_{s,t})^{-1}
&=\mathbb{W}\bigl(\mathrm{bot}\cdot\mathrm{right}(0\to t+1)\cdot\mathrm{top}_{t+1}^{-1}\cdot
\mathsf{a}_{s,t+1}^{-1}\bigr)\\
&=\mathbb{W}(\partial R_t).
\end{aligned}
\end{equation*}
Induction on $t$ from $t=0$, starting from $\mathbb{W}(\partial R_{-1})=\mathbf{1}$, gives
\begin{equation}\label{9.5:eq-ring-holonomy-3}
\mathbb{W}(\partial R_{n-1})=\prod_{t=0}^{n-1}\mathbb{W}(\mathsf{a}_{s,t})\,
\mathbb{W}(\partial c(s,t))\,\mathbb{W}(\mathsf{a}_{s,t})^{-1},
\end{equation}
the factors ordered with $t$ increasing from left to right.

(ii) \emph{Columns.} For $-1\le s\le n-1$ let $D_s:=[0,s+1]\times[0,n]$, so that $D_{n-1}=D$. We
use the following pieces of path.
\begin{itemize}
\item $v_s$ is the vertical path from $(s,0)$ to $(s,n)$.
\item $\mathrm{top}_n$ is, as in (i), the bond from $(s,n)$ to $(s+1,n)$.
\item $\mathrm{top}(0\to s)$ is the top side from $(0,n)$ to $(s,n)$.
\end{itemize}
Read from $(s,0)$, the boundary of the column $[s,s+1]\times[0,n]$ is
$\mathrm{bot}\cdot v_{s+1}\cdot\mathrm{top}_n^{-1}\cdot v_s^{-1}$. Moreover
$\mathsf{b}_s\,\mathrm{bot}=\mathsf{b}_{s+1}$. Read from $o$, the boundary of $D_{s-1}$ is
$\mathsf{b}_sv_s\,\mathrm{top}(0\to s)^{-1}v_0^{-1}$. For $s=0$ this is $v_0v_0^{-1}$, with value
$\mathbf{1}$. Hence, as products of bond variables,
\begin{equation*}
\begin{aligned}
\bigl[\mathsf{b}_s\cdot\partial([s,s+1]\times[0,n])\cdot\mathsf{b}_s^{-1}\bigr]\cdot\partial D_{s-1}
&=\mathsf{b}_{s+1}v_{s+1}\mathrm{top}_n^{-1}v_s^{-1}\mathsf{b}_s^{-1}\cdot
\mathsf{b}_sv_s\,\mathrm{top}(0\to s)^{-1}v_0^{-1}\\
&=\partial D_s .
\end{aligned}
\end{equation*}
For the last equality, the middle factor $v_s^{-1}\mathsf{b}_s^{-1}\mathsf{b}_sv_s$ cancels and
$\mathrm{top}_n^{-1}\mathrm{top}(0\to s)^{-1}=\mathrm{top}(0\to s+1)^{-1}$.

The column boundary is $\partial R_{n-1}$ of (i). Conjugation by $\mathbb{W}(\mathsf{b}_s)$ is
multiplicative, and $\mathbb{W}(\mathsf{b}_s)\mathbb{W}(\mathsf{a}_{s,t})=
\mathbb{W}(\mathsf{b}_s\mathsf{a}_{s,t})$. So \eqref{9.5:eq-ring-holonomy-3} gives
\begin{equation*}
\mathbb{W}(\partial D_s)=\Bigl(\prod_{t=0}^{n-1}\mathsf{L}(s,t)\Bigr)\,\mathbb{W}(\partial D_{s-1}).
\end{equation*}
Induction on $s$ from $s=0$, starting from $\mathbb{W}(\partial D_{-1})=\mathbf{1}$, gives
\begin{equation}\label{9.5:eq-ring-holonomy-4}
\mathbb{W}(\partial D)=\prod_{s=n-1}^{0}\ \prod_{t=0}^{n-1}\mathsf{L}(s,t),
\end{equation}
with $s$ decreasing and $t$ increasing from left to right. The loop $\partial D$, read from $o$, is
the square $\gamma$ with the orientation fixed in Step~1, read from its lower-left corner.

(iii) \emph{The product bound.} For matrices $A,B$,
\begin{equation*}
AB-\mathbf{1}=(A-\mathbf{1})(B-\mathbf{1})+(A-\mathbf{1})+(B-\mathbf{1}).
\end{equation*}
Hence, by submultiplicativity of $|\cdot|$,
\begin{equation*}
|AB-\mathbf{1}|\le(1+|A-\mathbf{1}|)(1+|B-\mathbf{1}|)-1 .
\end{equation*}
Applying this repeatedly to \eqref{9.5:eq-ring-holonomy-4}, and using $1+x\le e^{x}$, we get
\begin{equation}\label{9.5:eq-ring-holonomy-5}
|\mathbb{W}(\partial D)-\mathbf{1}|\le\exp\Bigl(\sum_{s,t}|\mathsf{L}(s,t)-\mathbf{1}|\Bigr)-1 .
\end{equation}
Write $W_{s,t}:=\mathbb{W}(\mathsf{b}_s\mathsf{a}_{s,t})$ for the transporter from $o$ to $(s,t)$.
Then $\mathsf{L}(s,t)-\mathbf{1}=W_{s,t}(\mathbb{W}(\partial c(s,t))-\mathbf{1})W_{s,t}^{-1}$. Since
$|W_{s,t}YW_{s,t}^{-1}|\le\|W_{s,t}\|\,|Y|\,\|W_{s,t}^{-1}\|$ for every matrix $Y$, we obtain
\begin{equation*}
|\mathsf{L}(s,t)-\mathbf{1}|\le\|W_{s,t}\|\,\|W_{s,t}^{-1}\|\,|\mathbb{W}(\partial c(s,t))-\mathbf{1}|.
\end{equation*}
The path $\mathsf{b}_s\mathsf{a}_{s,t}$ has $s+t\le 2n$ bonds. So $W_{s,t}$ is a product of at most
$2n$ bond variables, and $W_{s,t}^{-1}$ is the product of their inverses. For each bond $b$,
$\tilde U(b)$ is unitary, so
\begin{equation*}
\|\mathbb{W}(b)^{\pm1}\|\le\|e^{\pm i\eta\tilde A(b)}\|\le e^{\eta|\tilde A(b)|}.
\end{equation*}

\medskip
\noindent\emph{Step 3: part (b).}
Suppose that a pair $(\mathbb{W},\cdot)$ lies in $\mathfrak{G}(\mathfrak{B}_\tau)$ and that
$\mathbb{W}=\mathbf{V}(B_{\times})$ on the bonds of $X$. The square $\mathsf{Q}$ lies at level $j$,
with weight $1$.

The sixth clause of $\mathfrak{G}(\mathfrak{B}_\tau)$ gives
$\eta|\tilde A(b)|<5\alpha_{1,j}\eta/\ell_j$ on the bonds of $\mathsf{Q}$. By Step~2(iii), and since
$n=11M\ell_j/\eta$,
\begin{equation*}
\|W_{s,t}\|\,\|W_{s,t}^{-1}\|\le\exp\Bigl(2\cdot 2n\cdot\frac{5\alpha_{1,j}\eta}{\ell_j}\Bigr)
=e^{220M\alpha_{1,j}}\le 2 ,
\end{equation*}
the last inequality being $220M\alpha_{1,j}\le\log 2$.
The first clause of $\mathfrak{G}(\mathfrak{B}_\tau)$ gives, for every plaquette of $D$,
\begin{equation*}
|\mathbb{W}(\partial c(s,t))-\mathbf{1}|<\frac{5\alpha_{0,j}\eta^2}{\ell_j^2}.
\end{equation*}
There are
\begin{equation*}
n^2=\frac{121M^2\ell_j^2}{\eta^2}
\end{equation*}
plaquettes. Hence $\sum_{s,t}|\mathsf{L}(s,t)-\mathbf{1}|\le 1210M^2\alpha_{0,j}$.

Since $\gamma\subset X$, we have $\mathbb{W}(\gamma)=\mathbf{V}(B_{\times})(\gamma)$. Combining
\eqref{9.5:eq-ring-holonomy-2} with \eqref{9.5:eq-ring-holonomy-5} therefore yields
\begin{equation*}
P\varepsilon_h\le|\mathbb{W}(\gamma)-\mathbf{1}|\le e^{1210M^2\alpha_{0,j}}-1
\le 2420M^2\alpha_{0,j},
\end{equation*}
the last step being the inequality $e^{x}-1\le 2x$, valid for $0\le x\le 1$, at
$x=1210M^2\alpha_{0,j}$.
With $P=6M\tau_{\times}$ and $\tau_{\times}$ as in the corresponding display, this inequality
reads
\begin{equation*}
\frac{6M\varepsilon_h\,\alpha_{1,j}}{8C_{\times}}\le 2420M^2\alpha_{0,j}.
\end{equation*}
By the definition of the radii $\alpha_{i,j}$, the ratio $\alpha_{1,j}/\alpha_{0,j}$ is
$\mathsf{r}_1$; and $8\cdot2420/6=9680/3$. Hence, with $c_{\times}$ as in
the corresponding display,
\begin{equation*}
\mathsf{r}_1\le\frac{9680}{3}\,\frac{C_{\times}}{\varepsilon_h}\,M<c_{\times}M ,
\end{equation*}
which the hypothesis of part (b) of Proposition~\ref{9.5:prop-flat-ring} excludes. This proves
part (b).

\medskip
\noindent\emph{Step 4: part (c).}
The form $A_t$ of Proposition~\ref{9.5:prop-flat-ring} is closed, because $A_{\times}$ is; hence
$e^{i\eta A_t}$ is flat. Moreover $e^{i\eta A_t}$ is $G$-valued. A flat $G$-valued configuration is
therefore a minimum of the action, and it is critical for every set. On each cube it is a pure
gauge: the cube is simply connected, and its side is at most $2u$, smaller than the side $10u$ of
the hole, so that no cube contains a loop round the hole.

Consider now a real continuation, that is, one with $\mathbb{W}=\tilde U$ $G$-valued. Then every
transporter $W_{s,t}$ is unitary, so $\|W_{s,t}\|\,\|W_{s,t}^{-1}\|=1$ and
$|\mathsf{L}(s,t)-\mathbf{1}|\le|\tilde U(\partial c(s,t))-\mathbf{1}|$. For unitary $A$,
\begin{equation*}
|AB-\mathbf{1}|=|A(B-\mathbf{1})+(A-\mathbf{1})|\le|B-\mathbf{1}|+|A-\mathbf{1}|.
\end{equation*}
Applied to \eqref{9.5:eq-ring-holonomy-4}, this and the first clause, as in Step~3, give
\begin{equation*}
|\tilde U(\gamma)-\mathbf{1}|\le\sum_{s,t}|\tilde U(\partial c(s,t))-\mathbf{1}|
<n^2\cdot\frac{5\alpha_{0,j}\eta^2}{\ell_j^2}=605M^2\alpha_{0,j}.
\end{equation*}
On the bonds of $X\supseteq\gamma$ the continuation coincides with the datum $e^{i\eta A_t}$, whose
factors along $\gamma$ commute; its holonomy round $\gamma$, computed as in Step~1, is $e^{i6Mth}$.
Hence $\tilde U(\gamma)=e^{i6Mth}$. The matrix $e^{i6Mth}-\mathbf{1}$
has the eigenvalue $e^{i6Mt\varepsilon_h}-1$, and an operator norm dominates the spectral radius.
Hence, using $|e^{i\theta}-1|=2\sin(\theta/2)$ for $0\le\theta\le\pi$ and $2\sin x\ge x$ for
$0\le x\le1$,
\begin{equation*}
|e^{i6Mth}-\mathbf{1}|\ge\|e^{i6Mth}-\mathbf{1}\|\ge|e^{i6Mt\varepsilon_h}-1|
=2\sin(3Mt\varepsilon_h)\ge 3Mt\varepsilon_h
\qquad\text{for } 6Mt\varepsilon_h\le2 .
\end{equation*}
So, for $6Mt\varepsilon_h\le2$, we get $3Mt\varepsilon_h<605M^2\alpha_{0,j}$. Since
$605/3<202$, this gives $t\le 202M\alpha_{0,j}/\varepsilon_h$.
\end{proof}

For a determining set $\mathfrak{B}$ of date $d$ the continuation class $\mathfrak{G}(\mathfrak{B})$ consists of pairs $(e^{i\eta A'}U,\mathbb{J})$ on $\Omega_0(\mathfrak{B})$. A pair belongs to it if, at every point $y\in\Omega_1(\mathfrak{B})$, the clause functionals satisfy
\[
Q_a(y)<5\,\theta^{\mathfrak{B}}_a(y)\quad\text{for } a\in\{1,3,5,6,7\},
\]
and the cube clause (entry $8$) holds at five times the weight. On the dated collar $\operatorname{Col}(d):=\Omega_0(\mathfrak{B})\setminus\Omega_1(\mathfrak{B})$ the dated clauses $\mathsf{q}_a(y)<\mathsf{f}(d)\,\mathsf{a}_a$ hold instead, with table entries $\mathsf{a}_a$, where $\mathsf{a}_1=\alpha_{0,0}$. The entries are as follows:
\begin{itemize}
\item entry $1$ is the complex plaquette;
\item entry $3$ is the current;
\item entry $5$ is the real plaquette;
\item entry $6$ is the field $A'$;
\item entry $7$ is its covariant gradient.
\end{itemize}

The thresholds are $\theta^{\mathfrak{B}}_a(y)=w_{\mathfrak{B}}(y)\,r_a(\lambda_{\mathfrak{B}}(y))$, with $\ell_m:=L^m\eta$ and
\begin{gather*}
r_1(m)=r_5(m)=\alpha_{0,m}\eta^2\ell_m^{-2},\qquad r_3(m)=\alpha_{0,m}\ell_m^{-3},\\
r_6(m)=\alpha_{1,m}\ell_m^{-1},\qquad r_7(m)=\alpha_{1,m}\ell_m^{-2}.
\end{gather*}
The radii satisfy $\alpha_{0,m}\le\alpha_{1,m}$, and their quotient $\mathsf{r}_1=\alpha_{1,\cdot}/\alpha_{0,\cdot}$ does not depend on $m$.

\begin{corollary}[The unit-changed class under the wider-collar table]
\label{9.5:cor-unit-changed-class}
For a determining set $\mathfrak{B}$ let $\mathfrak{G}^{\vee}(\mathfrak{B})$ be defined exactly as
$\mathfrak{G}(\mathfrak{B})$, with two changes. First, in the clause of entry $1$ the radius $r_1(m)$ is replaced by
\begin{equation}\label{9.5:eq-unit-changed-class-1}
r_1^{\vee}(m):=\alpha_{1,m}\,\eta^{2}\ell_m^{-2}.
\end{equation}
Second, on the dated collar $\operatorname{Col}(d)$ the table entry $\mathsf{a}_1$ is replaced by $\alpha_{1,0}$.
Thus entry $1$ is measured in the unit $\alpha_1$ in place of $\alpha_0$.
\begin{enumerate}
\item[(i)] Assume the lemma on the continuation class accompanying the arrest theorem, with its properties
(e1)--(e4), and the level-zero lemma, both for $\mathfrak{G}$. Then their proofs, read with $r_1^{\vee}$ in place
of $r_1$ and with $\alpha_{1,0}$ in place of $\mathsf{a}_1$ on $\operatorname{Col}(d)$, show that (e1)--(e4) and the
level-zero lemma hold with $\mathfrak{G}^{\vee}$ in place of $\mathfrak{G}$.
\item[(ii)] Let $\mathsf{v}$ be a datum of $\mathfrak{T}^{\mathsf{m},\natural}_F[\varsigma](X)$
(see \eqref{9.5:eq-ambient-centre-e}). It has chart datum $B$, real centre $U_c$, chart $\mathfrak{X}(B)$, chart
configuration $\mathbf{V}(B)=e^{i\eta\mathfrak{X}(B)}U_c$ and slice $S=(\mathbf{V}(B),J(\mathbf{V}(B)))$. Let
$\tau$ be the date at which the triple $(F,X,j)$ of Definition~\ref{9.5:def-ambient-centre} is born and
$\mathfrak{B}_\tau$ its birth set; we write $\theta^{\tau}_a$, $w_\tau$, $\ell_\tau$, $\lambda_\tau$ for
$\theta^{\mathfrak{B}_\tau}_a$, $w_{\mathfrak{B}_\tau}$, $\ell_{\mathfrak{B}_\tau}$, $\lambda_{\mathfrak{B}_\tau}$.
The ambient field $U^{gl}$
is that of Definition~\ref{9.5:def-ambient-centre}.
Let $\chi$ be the cut-off of \eqref{9.5:eq-hull-geometry-b}, and let $X^{(1/2)}$ be as in
\eqref{9.5:eq-fattened-sets-c}. Assume:
\begin{enumerate}
\item[(a)] the slice obeys the first line of \eqref{9.5:eq-unit-changed-class-a};
\item[(b)] the strong form of the ambient-centre clause holds:
$(U^{gl},0)\in\mathfrak{G}(\mathfrak{B}_\tau)$, and $Q_5\bigl((U^{gl},0)\bigr)(y)\le 1.4\,\theta^{\tau}_5(y)$ at
every point $y$;
\item[(c)] the floor on $M$ and the ceiling on $\bar\gamma_0(\mathsf{t}_\gamma)$ in the second line of
\eqref{9.5:eq-unit-changed-class-a} hold:
\end{enumerate}
\begin{equation}\label{9.5:eq-unit-changed-class-a}
\begin{gathered}
E_a(S)(y)\le 1.1\,\theta^{\tau}_a(y)\qquad (y\in X^{(1/2)},\ a\in\{1,3,6,7\}),\\
M\ge 20L^{3},\qquad 4L^{4}\,\bar\gamma_0(\mathsf{t}_\gamma)\le 1.
\end{gathered}
\end{equation}
Then the pair on $\Omega_0(\mathfrak{B}_\tau)$ satisfies
\[
\bigl(e^{i\eta\chi\mathfrak{X}(B)}U^{gl},\ \chi J(\mathbf{V}(B))\bigr)\in\mathfrak{G}^{\vee}(\mathfrak{B}_\tau).
\]
Here $\chi\mathfrak{X}(B):=0$ off $X^{(1/2)}$, as in \eqref{9.5:eq-cutoff-continuation-a}.
\end{enumerate}
\end{corollary}

\begin{proof}
\emph{Part (i).} We go through the proofs of (e1)--(e4) and of the level-zero lemma, and mark where the radius of
entry $1$ enters.

\emph{Property (e1).} It compares the units and weights of a determining set and a later one. It reads no radius,
so its proof applies unchanged.

\emph{Property (e2).} Let $\mathfrak{B}'$ be later than $\mathfrak{B}$, and let the level rise from $\lambda$ by
$n\ge1$. Write $\mathsf{L}_0$ for the constant through which (e1) bounds the weights: $w_{\mathfrak{B}'}\le
\mathsf{L}_0^{2n}w_{\mathfrak{B}}$ where the level rises by $n$. Write $\beta_0$ for the constant of the ratio
bound \eqref{9.5:eq-unit-changed-class-2} below. The proof of (e2) for $\mathfrak{G}$ is the chain
\begin{align*}
\theta^{\mathfrak{B}'}_a
&\le \mathsf{L}_0^{2n}\,w_{\mathfrak{B}}\,r_a(\lambda+n)
\le w_{\mathfrak{B}}\,(L/3)^{n}(1+\beta_0)^{2}n^{1/2}L^{-e_an}\,r_a(\lambda)\\
&\le 3^{-n}(1+\beta_0)^{2}n^{1/2}\,\theta^{\mathfrak{B}}_a .
\end{align*}
The first inequality uses (e1). The second uses $\mathsf{L}_0^2\le L/3$ together with the ratio bound
\begin{equation}\label{9.5:eq-unit-changed-class-2}
\frac{r_a(m')}{r_a(m)}\ \ge\ (1+\beta_0)^{-2}n^{-1/2}L^{e_an},\qquad n:=m-m'\ge1 .
\end{equation}
Here $e_a$ is the power of $\ell_m^{-1}$ in $r_a(m)$, so $e_1=2$. The ratio bound
\eqref{9.5:eq-unit-changed-class-2} is the only property of the radii that the chain uses. For $r_1^{\vee}$,
since $\ell_{m'}^{-2}/\ell_m^{-2}=L^{2n}$, we have
\begin{equation*}
\frac{r_1^{\vee}(m')}{r_1^{\vee}(m)}=\frac{\alpha_{1,m'}}{\alpha_{1,m}}\,L^{2n}
=\frac{\alpha_{0,m'}}{\alpha_{0,m}}\,L^{2n}=\frac{r_1(m')}{r_1(m)} .
\end{equation*}
The middle equality holds because $\alpha_{0,\cdot}$ and $\alpha_{1,\cdot}$ have one and the same profile in the flow
time (their exponents of $\mathsf{t}$ coincide), so their quotient does not depend on the level. Hence
\eqref{9.5:eq-unit-changed-class-2} holds for $r_1^{\vee}$ with $e_1=2$, and (e2) holds for $\mathfrak{G}^{\vee}$.

\emph{Property (e3).} Property (e3) asserts
\begin{equation*}
\mathfrak{G}(\mathfrak{B}')\subset\mathfrak{G}(\mathfrak{B})\subset\mathfrak{K}_{109}(\mathfrak{B};\gamma_0^{(k)}),
\end{equation*}
and that $a_\flat\le\gamma_0^{(k)}$ on $\mathfrak{G}(\mathfrak{B})$.
\begin{itemize}
\item \emph{First inclusion.} Its proof observes that entries $1,3,5,6,7$ are quantities at $y$ not depending on
$\mathfrak{B}$, and then applies (e2). For $\mathfrak{G}^{\vee}$ it therefore follows from (e2) for
$\mathfrak{G}^{\vee}$, just proved.
\item \emph{Second inclusion.} Its proof does not read entry $1$. Entry $3$ gives the $\gamma_0$-clause, entries
$5$ and $8$ give the condition \cite[(3.35)]{B-CMP99b}, and entries $6$ and $7$ give the condition
\cite[(3.37)]{B-CMP99b}. So it holds for $\mathfrak{G}^{\vee}$ as it stands.
\item \emph{The bound on $a_\flat$.} Entry $1$ enters only through the member $\ell^2\eta^{-2}|\partial\mathbb{W}-1|$
of $a_\flat$, where $\mathbb{W}$ is the local complex field. On $\mathfrak{G}^{\vee}$ this member satisfies
\begin{equation*}
\ell^{2}\eta^{-2}|\partial\mathbb{W}-1|<5\,w\,\alpha_{1,\lambda}\le\gamma_0^{(k)} .
\end{equation*}
The last step holds because $\bar\alpha_j=\max_{m\le j}\max_i\alpha_i(\mathsf{t}_m)$ is a maximum over all $i$,
$i=1$ included.
\end{itemize}

\emph{Property (e4).} Since $\alpha_0\le\alpha_1$ we have $r_1\le r_1^{\vee}$ and $\alpha_{1,0}$ is at least the
replaced collar entry. Hence $\mathfrak{G}\subset\mathfrak{G}^{\vee}$, and the first clause of (e4) for
$\mathfrak{G}^{\vee}$ follows from the first clause for $\mathfrak{G}$. In the second clause, a map of the
construction with source set $\mathfrak{B}'$ and target set $\mathfrak{B}^t\preccurlyeq\mathfrak{B}'$ is applied to a
member $p$ over $Y'$ of its source, and $p$ has a continuation $\tilde p$. At $y\in Y'$ the clauses are those of $p$
itself. Off $Y'$ the clause of $\tilde p$ passes to $\mathfrak{B}^t$ by (e2), here (e2) for $\mathfrak{G}^{\vee}$.

\emph{The level-zero lemma.} Its four bounds estimate every quantity they use by a fraction of
$D(d)\,\alpha_{\cdot,0}$, where $D(d)$ is the factor depending on the date $d$ in the bounds of that lemma and
$\alpha_{\cdot,0}:=\min_i\alpha_{i,0}\le\alpha_{1,0}$; only the factor $\alpha_{\cdot,0}$ matters here. They
therefore remain valid
when $\mathsf{a}_1$ is replaced by $\alpha_{1,0}$ on $\operatorname{Col}(d)$.

\emph{Part (ii).} We verify the clauses of $\mathfrak{G}^{\vee}(\mathfrak{B}_\tau)$ entry by entry. The real part
carries entries $5$ and $8$ everywhere. On the support of $\chi$ we check entry $3$ and show that the cut-off
chart satisfies $|\chi\mathfrak{X}(B)|\le 1.2\,\theta^\tau_6$ and
$|\nabla(\chi\mathfrak{X}(B))|\le 1.2\,\theta^\tau_7$. Finally, entry $1$ follows from these by
Lemma~\ref{9.5:lem-perturbed-plaquette}. We carry this out.

\emph{Off $X^{(1/2)}$.} Here $\chi=0$ by \eqref{9.5:eq-hull-geometry-b}, so the pair is $(U^{gl},0)$. By (b) and
$\mathfrak{G}\subset\mathfrak{G}^{\vee}$, this pair is in $\mathfrak{G}^{\vee}(\mathfrak{B}_\tau)$.

\emph{On $X^{(1/2)}$: notation.} Fix $y$ and write $\ell:=\ell_\tau(y)$, $w:=w_\tau(y)$ and
$\alpha_i:=\alpha_{i,\lambda_\tau(y)}$. Then
\[
\theta^\tau_5(y)=w\alpha_0\eta^2\ell^{-2},\qquad \theta^\tau_6(y)=w\alpha_1\ell^{-1},\qquad
\theta^\tau_7(y)=w\alpha_1\ell^{-2}.
\]
By Definition~\ref{9.5:def-ambient-centre}, $U^{gl}$ agrees with $U_c$ on the bonds of
$X^{(3/4)}\supset X^{(1/2)}$.

\emph{Entries $5$, $8$, $3$, $6$.} Entries $5$ and $8$ are those of $U^{gl}$ and hold by (b). For entry $3$,
$|\chi J|\le|J|$. For entry $6$, $|\chi\mathfrak{X}|\le|\mathfrak{X}|$ by
Lemma~\ref{9.5:lem-cutoff-difference}. By (a), therefore, $|\chi J|\le 1.1\,\theta^\tau_3$ and
$|\chi\mathfrak{X}|\le 1.1\,\theta^\tau_6$.

\emph{Comparison at adjacent points.} For adjacent points $y,y'$,
\begin{equation}\label{9.5:eq-unit-changed-class-3}
\theta^\tau_6(y')\le L(1+\beta_0)^{2}(L/3)\,\theta^\tau_6(y)\le 0.44\,L^{2}\,\theta^\tau_6(y).
\end{equation}
Here $L$ is the cost of one change of level; $(1+\beta_0)^2$, with the constant $\beta_0$ of
\eqref{9.5:eq-unit-changed-class-2}, bounds the ratio of the radii at adjacent levels,
$\alpha_{1,m\pm1}\le(1+\beta_0)^{2}\alpha_{1,m}$; and $L/3$ is one step of the weight,
$\mathsf{L}_0^2\le L/3$.

\emph{Entry $7$.} By Lemma~\ref{9.5:lem-cutoff-difference},
\begin{equation*}
|\nabla_{U^{gl},\mu}(\chi\mathfrak{X})(b)|\le|\nabla_{U^{gl},\mu}\mathfrak{X}(b)|
+|\partial_\mu\chi(b_-)|\,|\mathfrak{X}(b+\eta e_\mu)| .
\end{equation*}
By \eqref{9.5:eq-hull-geometry-b} and \eqref{9.5:eq-hull-geometry-d}, $\ell_\tau|\partial_\mu\chi|\le 4L/M$.
We bound the two terms as follows:
\begin{itemize}
\item the first term by (a), giving $1.1\,w\alpha_1$ after multiplication by $\ell^2$;
\item in the second, $\ell|\mathfrak{X}(b+\eta e_\mu)|\le 1.1\,\ell\theta^\tau_6(y')$, which by
\eqref{9.5:eq-unit-changed-class-3} is at most $1.1\cdot0.44\,L^2w\alpha_1$.
\end{itemize}
Since $(4L/M)(1.1\cdot0.44\,L^2)=1.936\,L^3/M\le 0.0968$ by the floor $M\ge20L^3$ of
\eqref{9.5:eq-unit-changed-class-a}, we obtain
\begin{equation*}
\ell^{2}|\nabla(\chi\mathfrak{X})|\le 1.1\,w\alpha_1+\frac{4L}{M}\,(1.1\cdot0.44\,L^{2})\,w\alpha_1
\le 1.2\,w\alpha_1 .
\end{equation*}

\emph{Entry $1$: the three inputs.} Apply Lemma~\ref{9.5:lem-perturbed-plaquette} at each plaquette $p$ at $y$,
with $U^{gl}$ in place of $U$ and $\chi\mathfrak{X}(B)$ as the perturbing field. Write
$\mathbb{U}=e^{i\eta\chi\mathfrak{X}(B)}U^{gl}$ and $P=U^{gl}(\partial p)$, so that entry $1$ at $y$ is the
largest of the numbers $|\mathbb{U}(\partial p)-1|$ over the plaquettes $p$ at $y$. The inputs are:
\begin{itemize}
\item the bonds of $p$ lie at $y$ or at adjacent points, so the bound just used gives
$\ell\beta_6\le 0.49\,L^2w\alpha_1$;
\item entry $7$ gives $\ell^2\beta_7\le 1.2\,w\alpha_1$;
\item (b) gives $|P-1|\le1.4\,w\alpha_0\eta^2\ell^{-2}$, with $\alpha_0=\mathsf{r}_1^{-1}\alpha_1$.
\end{itemize}

\emph{Entry $1$: smallness.} We use $w\alpha_1\le\bar\gamma_0(\mathsf{t}_\gamma)/5$. With the ceiling in
\eqref{9.5:eq-unit-changed-class-a} this gives $L^4w\alpha_1\le1/20$ and $2L^2w\alpha_1\le 1/(10L^2)$.
Since $\eta\le\ell$, we get
\begin{equation*}
4\eta\beta_6\le 2L^2w\alpha_1\le 1/(10L^2)\le 1/10 .
\end{equation*}
Moreover $8\cdot0.49^2\,e^{1/10}\le2.2$.

\emph{Entry $1$: the bound.} Lemma~\ref{9.5:lem-perturbed-plaquette} now yields
\begin{equation*}
Q_1\le \eta^{2}\ell^{-2}w\alpha_1\Bigl[1.4\,\mathsf{r}_1^{-1}(1+2L^{2}w\alpha_1)+2.4+2.2\,L^{4}w\alpha_1\Bigr].
\end{equation*}
The term $2.4$ is $2\cdot1.2$, from $2\eta^2\beta_7$. The blocking factor $L$ is an odd integer larger than $1$, so
$L\ge3$ and $2L^2w\alpha_1\le 1/(10L^2)\le 1/90$. Since $\mathsf{r}_1\ge1$, the first member of the bracket is at
most $1.4(1+1/90)<1.5$, and the bracket is at most $1.5+2.4+0.11<5$. Hence $Q_1<5\,w\alpha_1\eta^2\ell^{-2}$, which
is the entry-$1$ clause of $\mathfrak{G}^{\vee}(\mathfrak{B}_\tau)$.
\end{proof}

\begin{corollary}[The two-cutoff domain inclusion]\label{9.5:cor-domain-inclusion}
Let $F=\mathbb{C}^{(i)}_{k_a}(X)$ be a stored left-over term of stage $i$ at step $k_a$, and read all
domains in the relaxed sense, that is, with \eqref{9.5:eq-relaxed-class-d} in force
(Definition~\ref{9.5:def-relaxed-class}). Assume:
\begin{itemize}
\item[(1)] clauses (a)--(f) of the domain-inclusion lemma for the flat table $\mathsf{T}^{\flat}_i$;
\item[(2)] clauses (d1)--(d5) of the domain-inclusion lemma for the five later layers;
\item[(3)] the local clauses of the increment statement on the set $W$ of that statement;
\item[(4)] the cut-off continuation of every datum (Theorem~\ref{9.5:thm-cutoff-continuation}), in the
form \eqref{9.5:eq-cutoff-continuation-c};
\item[(5)] part (iii) of the inherited estimates on the relaxed domains.
\end{itemize}
Then clause (d6) of the domain-inclusion lemma for the five later layers holds with no proviso on the
ambient continuation, in the following form. For the table-defined domain
$\mathfrak{D}_F=\mathfrak{U}[\mathsf{T}^{\flat}_i](X)$ of $F$, read in the relaxed sense,
\begin{equation}\label{9.5:eq-domain-inclusion-1}
\mathfrak{U}[\mathsf{T}^{\flat}_i](X)\cap
\bigl\{\tilde p:\ \tilde p \text{ satisfies } \eqref{9.5:eq-relaxed-class-d}\bigr\}
\subset \mathfrak{D}^{\flat,\natural}_i(X),
\end{equation}
and, for either machine $\mathsf{m}$ and every $\varsigma\in[0,1]$, every slice of a datum of
$\mathfrak{T}^{\mathsf{m},\natural}_F[\varsigma](X)$ is a pair belonging to the left-hand side of
\eqref{9.5:eq-domain-inclusion-1}. On this set part (a) of the inheritance theorem for stored terms
bounds $F$ by $C_0^{\sharp}e^{-a_m d_m(X)}$, where $C_0^{\sharp}$, $a_m$ and the index $m$ are the
constant, the decay rate and the level index of that part, and $F$ is holomorphic there.
\end{corollary}

\begin{proof}
Membership of a pair in $\mathfrak{D}^{\flat,\natural}_i(X)$ requires the local clauses on $X$, the
direct clauses in their fattened form at every location of the union of cubes $Y$ entering the
definition of $\mathfrak{D}^{\flat,\natural}_i(X)$, and the ambient condition
\eqref{9.5:eq-relaxed-class-d}, whose global half asks that the pair be the restriction of a pair
given on the whole domain. We supply these in turn, both
for a pair $\tilde p$ in the intersection on the left of \eqref{9.5:eq-domain-inclusion-1} and for
the slice of a datum of $\mathfrak{T}^{\mathsf{m},\natural}_F[\varsigma](X)$.

\emph{The local clauses.} These are the conclusions of hypotheses (1), (2) and (3): clauses (a)--(f)
of the domain-inclusion lemma for the flat table, clauses (d1)--(d5) of the domain-inclusion lemma for
the five later layers, and the local clauses of the increment statement on the set $W$ of that
statement. Each is used in the form in which it is assumed; none is derived anew here.

\emph{The fattened direct clauses.} The direct clauses, in the fattened form in which clause (d) of
the domain-inclusion lemma for the flat table states them, are required at every location of $Y$.
They are supplied by that clause (d), hypothesis (1), which
holds with the margin $(5/12)\beta$, where $\beta$ is the parameter entering the entries of the flat
table $\mathsf{T}^{\flat}_i$.

\emph{The ambient condition.} For a pair $\tilde p$ in the intersection on the left of
\eqref{9.5:eq-domain-inclusion-1}, the ambient condition \eqref{9.5:eq-relaxed-class-d}, and in
particular its global half, holds by membership in that intersection. For the slice of a datum of
$\mathfrak{T}^{\mathsf{m},\natural}_F[\varsigma](X)$, the ambient condition
\eqref{9.5:eq-relaxed-class-d}, and in particular its global half, is supplied by
\eqref{9.5:eq-cutoff-continuation-c} of Theorem~\ref{9.5:thm-cutoff-continuation}, hypothesis (4),
applied to that datum. Together with the two preceding paragraphs this gives
\eqref{9.5:eq-domain-inclusion-1}.

\emph{The slices.} That the slice of a datum of $\mathfrak{T}^{\mathsf{m},\natural}_F[\varsigma](X)$
lies in the table-defined domain $\mathfrak{U}[\mathsf{T}^{\flat}_i](X)=\mathfrak{D}_F$ is part of the
definition of that class, \eqref{9.5:eq-ambient-centre-e}; that it satisfies
\eqref{9.5:eq-relaxed-class-d} was shown in the preceding paragraph. Hence every such slice lies in
the left-hand side of \eqref{9.5:eq-domain-inclusion-1}.

\emph{The bound.} By Proposition~\ref{9.5:prop-relaxed-estimates}, part (iii), hypothesis (5), parts
(a)--(c) of the inheritance theorem for stored terms hold on the relaxed domains. Applying part (a) on
$\mathfrak{D}^{\flat,\natural}_i(X)$, which by \eqref{9.5:eq-domain-inclusion-1} contains the set in
question, gives that $F$ is holomorphic there and bounded by $C_0^{\sharp}e^{-a_m d_m(X)}$.
\end{proof}

\begin{corollary}[Tube points over the frozen region]\label{9.5:cor-frozen-tube-points}
Fix a step $k$ and a stage $n$, let $Z_W$ be the frozen region and $Z_W^{+}$ its ball collar, and
put $\mathfrak{B}^{\circ}:=\mathbb{B}^{(n+1)}_k\restriction Z_W$, the restriction to $Z_W$ of the
stage-$(n+1)$ structure at step $k$. By the definition of the hull domain, for a set $Y$ meeting
$Z_W^{+}$ a pair $(\mathbb{U},\mathbb{J})$ on $Y$ lies in the hull domain $\mathfrak{D}_H(Y)$ when
\begin{enumerate}
\item[(i)] it obeys on $Y$ the clauses of that definition, namely the local clause family
$\mathfrak{L}_{n+1}$ and the fat direct clauses of the regular-pair space
$\mathfrak{N}^{(n+1)}(Y)$, each at the locations of $Y$ at which that definition imposes it; and
\item[(ii)] it is the restriction to $Y$ of a continuation lying in the coarse regularity class
$\mathfrak{K}_{109}(\mathfrak{B}^{\circ};\gamma_0^{(k_a)})$, where $\gamma_0^{(k_a)}$ is the
current-slot radius at the density $k_a$.
\end{enumerate}
Assume the following.
\begin{itemize}
\item[(a)] Clause $(\beta3)$ of the cut-off continuation of every datum,
Theorem~\ref{9.5:thm-cutoff-continuation}, displayed in
\eqref{9.5:eq-cutoff-continuation-e}, holds at stage $n+1$.
\item[(b)] For every set $Y$ meeting $Z_W^{+}$ and every tube point $\tilde p$ of the uniqueness
comparison over $Z_W^{+}$, the restriction of $\tilde p$ to $Y$ satisfies clauses
\textup{(d1)}--\textup{(d5)} of the increment statement.
\end{itemize}
Then, for every set $Y$ meeting $Z_W^{+}$, the restriction to $Y$ of every tube point of the
uniqueness comparison over $Z_W^{+}$, that is of every complex pair
$\tilde p=(e^{i\eta\tilde A}\tilde U,\tilde{\mathbb{J}})$ of that comparison over $Z_W^{+}$, lies
in $\mathfrak{D}_H(Y)$.
\end{corollary}

\begin{proof}
Let $Y$ be a set meeting $Z_W^{+}$, and let $\tilde p$ be a tube point of the uniqueness
comparison over $Z_W^{+}$ whose restriction to $Y$ we test for membership in
$\mathfrak{D}_H(Y)$. This membership consists of requirements (i) and (ii) of the statement. We
verify them in turn.

Requirement (ii) is a continuation of $\tilde p$ lying in the coarse regularity class
$\mathfrak{K}_{109}(\mathfrak{B}^{\circ};\gamma_0^{(k_a)})$. This continuation is supplied by
clause $(\beta3)$ of \eqref{9.5:eq-cutoff-continuation-e}, applied at stage $n+1$; this is
hypothesis~(a). Theorem~\ref{9.5:thm-cutoff-continuation} requires the ambient field $U^{gl}$ of
the ambient-centre clause \eqref{9.5:eq-ambient-centre-a}. This is the membership clause of
Definition~\ref{9.5:def-ambient-centre}, which asks that the real centre $U_c$ of a datum about a
leaf $(F,X,j)$ born at date $\tau$ agree on the bonds of $X^{(3/4)}$ with a gauge-group-valued
configuration $U^{gl}$ on $\Omega_0(\mathfrak{B}_\tau)$ that obeys the first of the two regularity
bounds of the class $\mathfrak{K}_{109}$, the one formed at the radius $\alpha_0^{(99)}$, over
$\mathfrak{B}_{d''}$ for every date $d''\le\tau$.
Lemma~\ref{9.5:lem-ambient-closure} shows that this clause is preserved along the run, so the
data at which Theorem~\ref{9.5:thm-cutoff-continuation} is applied satisfy it.

Requirement (i), that is the clauses $\mathfrak{L}_{n+1}$ and the fat direct clauses of
$\mathfrak{N}^{(n+1)}(Y)$ at their respective locations, is the content of clauses
\textup{(d1)}--\textup{(d5)} of the increment statement for the restriction of $\tilde p$ to $Y$,
and these hold by hypothesis~(b).

Both requirements (i) and (ii) therefore hold, so the restriction of $\tilde p$ to $Y$ lies in
$\mathfrak{D}_H(Y)$. Since $Y$ was an arbitrary set meeting $Z_W^{+}$ and $\tilde p$ an arbitrary
tube point of the uniqueness comparison over $Z_W^{+}$, the corollary follows.
\end{proof}

\begin{lemma}[A maximum principle for coupled scale-indexed inequalities]
\label{9.6:lem-maximum-principle}
Let $I$ be a finite set and $I_0\subset I$. Let $D=(D_i)_{i\in I}$ and
$E=(E_i)_{i\in I}$ be real families with $D_i\ge 0$ and $E_i>0$ for all
$i\in I$, let $\sigma=(\sigma_i)_{i\in I_0}$ satisfy $\sigma_i\ge 0$, and let
$\mu=(\mu_{ij})_{i\in I_0,\,j\in I}$ satisfy $\mu_{ij}\ge 0$. Assume that for
every $i\in I_0$
\begin{equation}\label{9.6:eq-maximum-principle-1}
D_i\le\sum_{j\in I}\mu_{ij}D_j+\sigma_i,\qquad
\sum_{j\in I}\mu_{ij}E_j\le\tfrac12 E_i,\qquad
\sigma_i\le\tfrac12 E_i,
\end{equation}
and that $D_i\le E_i$ for every $i\in I\setminus I_0$. Then $D_i\le E_i$ for
every $i\in I$.
\end{lemma}

No triangularity of the coefficients $(\mu_{ij})$ with respect to any ordering
of $I$ is assumed.

Proof outstanding.

\paragraph{Route.}
A proof would have to establish, from the hypotheses displayed in
\eqref{9.6:eq-maximum-principle-1} together with $D_i\le E_i$ off $I_0$, that
the maximum of the ratios $D_i/E_i$ over the finite set $I$ does not exceed
$1$. If this maximum exceeded $1$, it would be attained at an index of $I_0$,
since all ratios with index outside $I_0$ are at most $1$. At such an index
the first inequality of \eqref{9.6:eq-maximum-principle-1}, the bound
$D_j\le(\max_{k\in I}D_k/E_k)\,E_j$ for every $j\in I$ combined with
$\mu_{ij}\ge 0$, and the second and third inequalities of
\eqref{9.6:eq-maximum-principle-1} would give that this maximum is at most
one half of itself plus one half, which is incompatible with its exceeding $1$.

\begin{definition}[The comparison chain between the two runs]\label{9.6:def-comparison-chain}
Let $\mathsf{A}$ and $\mathsf{B}$ be the two runs compared in
Definition~\ref{9.1:def-comparison-families}. Let $n$ and the lower depth $r_0$ be as there. Let
$\prec$ be the order on the block letters fixed in Definition~\ref{9.1:def-hybrid-lists}; below,
$\bullet$ denotes a block letter and $c$ a component of $\bullet$. The \emph{comparison chain} is the
family of competitors $G_s^c$, $r_0 \le s \le n$, with $c$ running over the components of the blocks,
together with the numbers $g_r^c$ and $g_r^{\bullet}$, $r_0 \le r < n$, defined by the following
downward recursion.

At the top depth, for every component $c$,
\begin{equation*}
G_n^c := \bigl(\mathfrak{T}F_n^{\mathsf{B}}\bigr)^c ,
\end{equation*}
where $F_s^{\mathsf{A}}$ and $F_s^{\mathsf{B}}$ denote the families of the runs $\mathsf{A}$ and
$\mathsf{B}$ at depth $s$, and $\mathfrak{T}F_s^{\mathsf{B}}$ is the transplant of the latter.

Let $r = n-1, n-2, \dots, r_0$ be taken in decreasing order. For fixed $r$, let the block letters
$\bullet$ be taken along $\prec$. Let $N_r^{\bullet}$ be the comparison family of
Definition~\ref{9.1:def-comparison-families} at the hybrid list $\mathfrak{s}^{\natural}_{r,\bullet}$
of Definition~\ref{9.1:def-hybrid-lists}. This list is formed from the competitors $G_s^{c'}$, $s>r$,
and $G_r^{c'}$ for the components $c'$ of the blocks $\circ \prec \bullet$, all of which have already
been defined. Let $M_r^{\bullet}$ be
the production difference at depth $r$ in the block $\bullet$. Then, for each component $c$ of
$\bullet$,
\begin{equation}\label{9.6:eq-comparison-chain-1}
G_r^c :=
\begin{cases}
\bigl(N_r^{\bullet}+M_r^{\bullet}\bigr)^c
  & \text{if } \bigl\|\bigl(N_r^{\bullet}+M_r^{\bullet}\bigr)^c\bigr\|^c
    \le (C_V+1)\,\mathfrak{u}_r^c ,\\[2pt]
\bigl(\mathfrak{T}F_r^{\mathsf{B}}-F_r^{\mathsf{A}}\bigr)^c & \text{otherwise,}
\end{cases}
\end{equation}
and
\begin{equation*}
g_r^c := \bigl\|G_r^c\bigr\|^c , \qquad g_r^{\bullet} := \max_{c\in\bullet} g_r^c .
\end{equation*}
Here $\|\cdot\|^c$ is the norm on component $c$. The quantity $\mathfrak{u}_r^c$ is the reading unit of
component $c$ at depth $r$, and $C_V$ is the constant in the competitor bound of
Definition~\ref{9.1:def-hybrid-lists}.
\end{definition}

Proof outstanding.

\noindent\emph{Route.}
The recursion is well posed only if, at every depth $r$ with $r_0 \le r < n$ and every block letter
$\bullet$, the list $\mathfrak{s}^{\natural}_{r,\bullet}$ is a hybrid list in the sense of
Definition~\ref{9.1:def-hybrid-lists}, that is, every competitor entering this list satisfies
$\|G_s^c\|^c \le (C_V+1)\,\mathfrak{u}_s^c$, so that the comparison family $N_r^{\bullet}$ of
Definition~\ref{9.1:def-comparison-families} exists there. For the entries given by the first case
of \eqref{9.6:eq-comparison-chain-1} this bound is the defining condition of that case. What has to
be established is the same bound for the top entries $(\mathfrak{T}F_n^{\mathsf{B}})^c$ and for the
second-case entries $(\mathfrak{T}F_s^{\mathsf{B}}-F_s^{\mathsf{A}})^c$ at each depth $s$, from bounds
on the transplanted $\mathsf{B}$-family and the $\mathsf{A}$-family at that depth.

\begin{lemma}[Well-definedness and structure of the comparison chain]\label{9.6:lem-chain-structure}
Assume the inputs of the transplanted one-lattice response bounds
(Definition~\ref{9.1:def-transplant-response}) and the existence of the two-run comparison families
with their bounds (Definition~\ref{9.1:def-comparison-families}). Then the comparison chain between
the two runs (Definition~\ref{9.6:def-comparison-chain}) is well defined. Moreover, the following
hold for every depth $r$ of the chain, every block $\bullet$ and every component $c$ of $\bullet$.
\begin{enumerate}
\item[(c1)] Each $G_r^c$ is a list-structured $\mathsf{A}$-family on $c$. It is holomorphic on
$\mathfrak{P}$, and its local members read $W{\upharpoonright}Y^{+}$. Its totals satisfy the first
line of \eqref{9.6:eq-chain-structure-a}, in the sense of
Definition~\ref{9.1:def-totals-units-distance}.
\item[(c2)] $G_r^c$ obeys the second line of \eqref{9.6:eq-chain-structure-a}, where $C_V$ is the
constant in the selection threshold of Definition~\ref{9.6:def-comparison-chain}. Consequently every
hybrid list $\mathfrak{s}^{\natural}_{r,\bullet}$ belongs to $C_{\mathfrak{K}}\mathfrak{K}^c$.
\item[(c3)] For $r<n$, the component distance between the runs is bounded by $g_r^c$, and $g_r^c$ and
$g_r^{\bullet}$ are bounded by the sum of the norms of the production difference $M_r$ and of the
comparison family $N_r$. At $r=n$ one has the bound by $C_V\,\mathfrak{u}_n^c$. These are the last
lines of \eqref{9.6:eq-chain-structure-a}.
\end{enumerate}
\begin{equation}\label{9.6:eq-chain-structure-a}
\begin{aligned}
&\text{(c1)}\qquad G_r^c \approx \varphi_r^{\mathsf{B},c}-\varphi_r^{\mathsf{A},c};\\
&\text{(c2)}\qquad g_r^c \le (C_V+1)\,\mathfrak{u}_r^c;\\
&\text{(c3)}\qquad \mathsf{D}_r^c \le g_r^c \le \|M_r^c\|+\|N_r^c\|,\\
&\phantom{\text{(c3)}}\qquad\text{hence}\quad
g_r^{\bullet} \le \|M_r^{\bullet}\|+\|N_r^{\bullet}\| \qquad (r<n);\\
&\phantom{\text{(c3)}}\qquad g_n^c \le C_V\,\mathfrak{u}_n^c.
\end{aligned}
\end{equation}
\end{lemma}

\noindent Proof outstanding.

\noindent\emph{Route.}
The proof would be a downward induction along $r$ and the order $\prec$ of the sub-steps: given (c1)
and (c2) at the earlier entries, $\mathfrak{s}^{\natural}_{r,\bullet}$ is a hybrid list as in
\eqref{9.1:eq-hybrid-lists-a}, so $M_r^{\bullet}$ and $N_r^{\bullet}$ exist, and each of the two
alternatives in the definition of $G_r^c$ has totals $\varphi_r^{\mathsf{B},c}-\varphi_r^{\mathsf{A},c}$,
the first by component-wise additivity of totals together with the comparison family, the second by
\eqref{9.1:eq-run-pair-transplant-a}, with norm at most $(C_V+1)\mathfrak{u}_r^c$ by
\eqref{9.1:eq-totals-units-distance-a}. What has to be established is that any mixture of the two
alternatives, selected by one norm dichotomy for each pair of runs, each $r$ and each component $c$
(hence independently of the parameters), is again a list-structured $\mathsf{A}$-family, holomorphic
on $\mathfrak{P}$ and with the stated read sets. Then (c2) would be the selection rule, and (c3) would
follow from the triangle inequality in the first case and from the failed threshold in the second
case, after which one takes the maximum over $c\in\bullet$.

\begin{proposition}[The two-run estimate with the power-law floor]\label{9.6:prop-two-run-estimate}
Assume hypothesis (J1) of \eqref{9.1:eq-tuned-family-a}, the transplanted one-lattice response
bounds and the exact shift identity of Definition~\ref{9.1:def-transplant-response}, the two-run
comparison families of Definition~\ref{9.1:def-comparison-families} with constant $C_\sigma$, and the
ceilings on the coupling of Definition~\ref{9.1:def-coupling-ceilings}. Consider a pair of lists,
frozen or local, alive to depth $r_0$, the two runs being indexed $\mathsf{A}$ (length $n$) and
$\mathsf{B}$ (length $n+1$), with
\begin{equation*}
\zeta^{\mathsf{B}}_{(r_0)}=\zeta^{\mathsf{A}}_{(r_0)}\quad\text{on }\mathfrak{P},
\end{equation*}
complex deviations on the tube being allowed. Then for $r_0\le r\le n$
\begin{equation}\label{9.6:eq-two-run-estimate-1}
g_r^P\le C_E e_r,\qquad g_r^A\le 4C_E e_r,\qquad g_r^R\le 4C_E\hat{\mathfrak{r}}_r e_r,
\end{equation}
and for every component $c$ of the block $\bullet\in\{P,A,R\}$ the component distance
$\mathsf{D}_r^c$ obeys the bound stated for $g_r^{\bullet}$. For $r_0\le r<n$
\begin{align}
d_{r+1}&\le \mathsf{A}C_E\Bigl[\frac{\rho^{n-r}}{1-\rho}+(r-r_0+1)\Pi_n\Bigr]
\le \mathsf{A}C_E\mathsf{n}_r e_r,\label{9.6:eq-two-run-estimate-2}\\
d^{\sharp}_{r+1}&\le c_\rho\mathsf{A}C_E\mathsf{n}_r e_r.\label{9.6:eq-two-run-estimate-3}
\end{align}
Here $g_r^{\bullet}$, $\bullet\in\{P,A,R\}$, are the bounds of the blocks $G_r^{\bullet}$ of the
comparison chain of the pair, $g_r^c$ those of their components $G_r^c$, and $\mathsf{D}_r^c$ the
component distances between the two runs, as in
Lemma~\ref{9.6:lem-chain-structure}; $d_s$ is the coupling distance and $d^{\sharp}_{r+1}$ its
weighted tail, as they enter Definition~\ref{9.1:def-transplant-response}; $\mathsf{A}$ is the
constant of hypothesis (J1) in \eqref{9.1:eq-tuned-family-a}; and $C_E$, $e_r=\rho^{n-r}+\Pi_n$,
$\hat{\mathfrak{r}}_r$, $\mathsf{n}_r$, $c_\rho$
and the power-law floor $\Pi_n$ are the constants fixed in \eqref{9.1:eq-rate-constants-a}, $C_E$
being read according to the rule stated with \eqref{9.1:eq-rate-constants-a}.
\end{proposition}

\begin{proof}
\emph{Step 1: the system of inequalities.} Fix $r$ with $r_0\le r<n$ and a block
$\bullet\in\{P,A,R\}$. By (c3) of Lemma~\ref{9.6:lem-chain-structure} (see
\eqref{9.6:eq-chain-structure-a}), $g_r^{\bullet}\le\|M_r^{\bullet}\|+\|N_r^{\bullet}\|$. We bound
$\|M_r^{\bullet}\|$ by the rows
(P), (A), (R) of item (a) of \eqref{9.1:eq-transplant-response-a}, and $\|N_r^{\bullet}\|$ by the
bounds of
Definition~\ref{9.1:def-comparison-families}, including the bound on $N_r^{R'}$ displayed in
\eqref{9.1:eq-comparison-families-a}; the norms $\|G_s^c\|$ of the members of the chain occurring on
the right-hand sides are expressed through the quantities $g_s^c$ of
Lemma~\ref{9.6:lem-chain-structure}, and we write $g=(g_s^c)$ for the array of these bounds, so that
the bracket $B_r$ of \eqref{9.1:eq-comparison-families-a} is
accordingly written $B_r(g,d)$; moreover we use that the square of the reference coupling is at
most one. This gives
\begin{align}
g_r^P&\le \alpha_r d^{\sharp}_{r+1}+\varkappa\sum_{s>r}\bigl(q^{s-r}g_s^P+g_s^R\bigr)
+\varkappa''_r\sum_{s>r}q^{s-r}g_s^A+C_\sigma e_r,\label{9.6:eq-two-run-estimate-4}\\
g_r^A&\le \alpha_r^A d^{\sharp}_{r+1}
+\varkappa_A\Bigl[g_r^P+\sum_{s>r}\bigl(q^{s-r}(g_s^P+g_s^A)+g_s^R\bigr)\Bigr]\notag\\
&\qquad+g_r^R+C_\sigma e_r,\label{9.6:eq-two-run-estimate-5}\\
g_r^R&\le 2\mathfrak{r}_r\Lambda_r B_r(g,d)+2\varkappa'_r\sum_{r<s<r+\ell_r}g_s^R
+\mathfrak{r}_r\Lambda_r(2C_\sigma+1)e_r.\label{9.6:eq-two-run-estimate-6}
\end{align}
In \eqref{9.6:eq-two-run-estimate-6} the factor $2$ arises because the row (R) of item (a) of
\eqref{9.1:eq-transplant-response-a} and the bound on $N_r^{R'}$ in
\eqref{9.1:eq-comparison-families-a} each carry one bracket $\mathfrak{r}_r\Lambda_rB_r(g,d)$ and one
band $\varkappa'_r\sum_{r<s<r+\ell_r}g_s^R$.

It remains to express the coupling distances through the $g$'s. Both runs subtract the same
constant $c_0$ at every step, so that the difference
$\delta\zeta_{(u)}:=\zeta^{\mathsf{B}}_{(u)}-\zeta^{\mathsf{A}}_{(u)}$ of the running shifts obeys
\begin{equation*}
\delta\zeta_{(u+1)}=\delta\zeta_{(u)}-\bigl(b^{\mathsf{B}}_{(u)}-b^{\mathsf{A}}_{(u)}\bigr),
\end{equation*}
and $d_{r_0}=0$ because the two running shifts agree on $\mathfrak{P}$ at depth $r_0$. Summing this
recursion from $u=r_0$ to $u=s$ and bounding each difference $b^{\mathsf{B}}_{(u)}-b^{\mathsf{A}}_{(u)}$
by item (b) of \eqref{9.1:eq-transplant-response-a}, with the constant $\mathsf{A}$ of hypothesis
(J1) in \eqref{9.1:eq-tuned-family-a}, we obtain
\begin{equation}\label{9.6:eq-two-run-estimate-7}
d_{s+1}\le\mathsf{A}\sum_{r_0\le u\le s}g_u^P\qquad(r_0\le s<n).
\end{equation}
The restriction $s<n$ is due to the fact that the shift $\zeta^{\mathsf{B}}_{(n+1)}$ of the longer
run is contained in the last member $G_n$ of the chain. We substitute
\eqref{9.6:eq-two-run-estimate-7} for every coupling distance occurring in
\eqref{9.6:eq-two-run-estimate-4}--\eqref{9.6:eq-two-run-estimate-6}, in the weighted tails
$d^{\sharp}_{r+1}$ as well as in the brackets. All entries of the resulting system are
nonnegative. For the envelope $E$ of Step~2 we write
$d_{s+1}(E):=\mathsf{A}\sum_{r_0\le u\le s}E_u^P$, and $d^{\sharp}_{r+1}(E)$, $B_r(E)$ for the
expressions obtained from $d^{\sharp}_{r+1}$, $B_r(g,d)$ after this substitution and with $g$
replaced by $E$.

\emph{Step 2: the maximum principle and the sums.} We apply the maximum principle for coupled
scale-indexed inequalities (Lemma~\ref{9.6:lem-maximum-principle}) with
\begin{equation*}
I:=\{(r,\bullet):r_0\le r\le n,\ \bullet\in\{P,A,R\}\},\qquad
I_0:=\{(r,\bullet)\in I:r<n\},
\end{equation*}
$D_{(r,\bullet)}:=g_r^\bullet$, the coefficient array of Lemma~\ref{9.6:lem-maximum-principle} being
that of the nonnegative coefficients of the substituted
system \eqref{9.6:eq-two-run-estimate-4}--\eqref{9.6:eq-two-run-estimate-6}, $\sigma$ its
constant terms, namely $\sigma_{(r,P)}=\sigma_{(r,A)}=C_\sigma e_r$ and
$\sigma_{(r,R)}=\mathfrak{r}_r\Lambda_r(2C_\sigma+1)e_r$, and with the envelope
\begin{equation*}
E_r^P:=C_E e_r,\qquad E_r^A:=4C_E e_r,\qquad E_r^R:=4C_E\hat{\mathfrak{r}}_r e_r .
\end{equation*}
The set $I$ is finite, $D\ge0$, $\sigma\ge0$ and the coefficients are nonnegative. For an index of
$I_0$, the sum of $E$ weighted by the coefficients of its row, which
Lemma~\ref{9.6:lem-maximum-principle} compares with half the envelope, is
the right-hand side of that row with $g$ replaced by $E$ and the constant term removed.

We use $q\le\rho^2$, which holds by the choice of $\rho$ in \eqref{9.1:eq-rate-constants-a}, and
$e_{r+m}\le\rho^{-m}e_r$ for $m\ge0$, which follows from $0<\rho<1$:
\begin{equation*}
e_{r+m}=\rho^{n-r-m}+\Pi_n\le\rho^{-m}\bigl(\rho^{n-r}+\Pi_n\bigr)=\rho^{-m}e_r .
\end{equation*}
All sums below run over indices not exceeding $n$. We claim
\begin{equation}\label{9.6:eq-two-run-estimate-a}
\begin{aligned}
&\text{(o1)}\quad \sum_{s\ge r}q^{s-r}e_s\le\frac{e_r}{1-\rho},\qquad
\sum_{s>r}q^{s-r}e_s\le\frac{\rho\,e_r}{1-\rho};\\
&\text{(o2)}\quad \sum_{r_0\le u\le r}e_u\le\frac{\rho^{n-r}}{1-\rho}+(r-r_0+1)\Pi_n
\le\mathsf{n}_r e_r;\\
&\text{(o3)}\quad \sum_{s>r}\hat{\mathfrak{r}}_s e_s\le 2\Pi_n\le 2e_r;\\
&\text{(o4)}\quad \sum_{m\ge0}q^m\mathsf{n}_{r+m}e_{r+m}
\le e_r\Bigl[\frac{\mathsf{n}_r}{1-\rho}+\frac{\rho}{(1-\rho)^2}\Bigr]
\le\mathsf{n}_r e_r\,\frac{1+\rho}{1-\rho};\\
&\text{(o5)}\quad d^{\sharp}_{r+1}(E)\le c_\rho\mathsf{A}C_E\mathsf{n}_r e_r .
\end{aligned}
\end{equation}

For (o1): by the two inequalities just stated,
$q^{s-r}e_s\le\rho^{2(s-r)}\rho^{-(s-r)}e_r=\rho^{s-r}e_r$,
and summing the geometric series over $s\ge r$, respectively over $s>r$, gives the two bounds.

For (o2): since $e_u=\rho^{n-u}+\Pi_n$,
\begin{equation*}
\sum_{r_0\le u\le r}e_u=\sum_{r_0\le u\le r}\rho^{n-u}+(r-r_0+1)\Pi_n
\le\frac{\rho^{n-r}}{1-\rho}+(r-r_0+1)\Pi_n ,
\end{equation*}
and the last expression is at most $\mathsf{n}_r(\rho^{n-r}+\Pi_n)=\mathsf{n}_re_r$ by the
definition of $\mathsf{n}_r$ in \eqref{9.1:eq-rate-constants-a}.

For (o3): write $e_s=\rho^{n-s}+\Pi_n$. By the first inequality of item (f) of
\eqref{9.1:eq-coupling-ceilings-a}, with the constant $C_T$ of that item,
$\hat{\mathfrak{r}}_s\le C_T\rho^{-(n-s)/2}\hat{\mathfrak{r}}_n$ for
$s\le n$, so that, with $m=n-s$,
\begin{equation*}
\sum_{s>r}\hat{\mathfrak{r}}_s e_s\le C_T\hat{\mathfrak{r}}_n\sum_{m\ge0}\rho^{m/2}
+\Pi_n\sum_{s}\hat{\mathfrak{r}}_s\le 2\Pi_n\le 2e_r ;
\end{equation*}
here the first term is at most $\Pi_n$ by the definition of the floor $\Pi_n$ in
\eqref{9.1:eq-rate-constants-a}, the second term is at most $\Pi_n$ by
$\sum_s\hat{\mathfrak{r}}_s\le1$ in item (e) of \eqref{9.1:eq-coupling-ceilings-a}, and
$\Pi_n\le e_r$ holds because $e_r=\rho^{n-r}+\Pi_n$.

For (o4): as in (o1), $q^me_{r+m}\le\rho^me_r$. By the definition of $\mathsf{n}_r$ in
\eqref{9.1:eq-rate-constants-a}, $\mathsf{n}_{r+m}\le\mathsf{n}_r+m$ for $m\ge0$, so that
\begin{equation*}
\sum_{m\ge0}\rho^m\mathsf{n}_{r+m}\le\frac{\mathsf{n}_r}{1-\rho}+\sum_{m\ge0}m\rho^m ,
\end{equation*}
and $\sum_{m\ge0}m\rho^m=\rho/(1-\rho)^2$; this is the first inequality. The second is equivalent to
$\rho/(1-\rho)^2\le\mathsf{n}_r\rho/(1-\rho)$, that is, to $\mathsf{n}_r(1-\rho)\ge1$, which holds by
the same definition.

For (o5): on the envelope, (o2) gives, for $r_0\le s<n$,
\begin{equation*}
d_{s+1}(E)=\mathsf{A}C_E\sum_{r_0\le u\le s}e_u\le\mathsf{A}C_E\mathsf{n}_se_s .
\end{equation*}
We insert these bounds into $d^{\sharp}_{r+1}(E)$ and
use $L^{-1/2}<q$, which holds by the choice of $q$ in \eqref{9.1:eq-rate-constants-a}, together with
(o4) at $r$ and at $r+1$. The application at $r$ contributes
$\mathsf{A}C_E\mathsf{n}_re_r(1+\rho)/(1-\rho)$; the application at $r+1$ contributes
$\mathsf{A}C_E\mathsf{n}_{r+1}e_{r+1}(1+\rho)/(1-\rho)$, and by $\mathsf{n}_{r+1}\le\frac32\mathsf{n}_r$
and $e_{r+1}\le e_r/\rho$ this is at most $\frac{3}{2\rho}\mathsf{A}C_E\mathsf{n}_re_r(1+\rho)/(1-\rho)$.
Hence
\begin{equation*}
d^{\sharp}_{r+1}(E)\le\mathsf{A}C_E\mathsf{n}_re_r\,
\frac{(1+\rho)\bigl(1+\frac{3}{2\rho}\bigr)}{1-\rho}\le c_\rho\mathsf{A}C_E\mathsf{n}_re_r ,
\end{equation*}
the last inequality by the definition of $c_\rho$ in \eqref{9.1:eq-rate-constants-a}.

\emph{Step 3: the three rows.} Let $r_0\le r<n$.

\emph{Row $P$.} In \eqref{9.6:eq-two-run-estimate-4} with $g=E$ the terms other than the constant
term are bounded as follows:
\begin{align*}
\alpha_rd^{\sharp}_{r+1}(E)&\le c_\rho\alpha_r\mathsf{A}C_E\mathsf{n}_re_r
\qquad\text{by (o5)},\\
\varkappa\sum_{s>r}q^{s-r}E_s^P&\le\varkappa C_E\,\frac{\rho e_r}{1-\rho}
\qquad\text{by the second bound of (o1)},\\
\varkappa\sum_{s>r}E_s^R=4\varkappa C_E\sum_{s>r}\hat{\mathfrak{r}}_se_s&\le8\varkappa C_Ee_r
\qquad\text{by (o3)},\\
\varkappa''_r\sum_{s>r}q^{s-r}E_s^A&\le4\varkappa''_rC_E\,\frac{\rho e_r}{1-\rho}
\qquad\text{by the second bound of (o1)}.
\end{align*}
Items (a), (b) and (b$'$) of \eqref{9.1:eq-coupling-ceilings-a} give
$c_\rho\alpha_r\mathsf{A}\mathsf{n}_r\le\frac18$, $\varkappa\rho/(1-\rho)\le\frac1{32}$,
$8\varkappa\le\frac14$ and $4\varkappa''_r\rho/(1-\rho)\le\frac3{32}$, so that the right-hand side
of \eqref{9.6:eq-two-run-estimate-4} at $g=E$, less its constant term, is at most
\begin{equation*}
C_Ee_r\Bigl[c_\rho\alpha_r\mathsf{A}\mathsf{n}_r+\frac{\varkappa\rho}{1-\rho}
+8\varkappa+\frac{4\varkappa''_r\rho}{1-\rho}\Bigr]
\le C_Ee_r\Bigl[\frac18+\frac1{32}+\frac14+\frac3{32}\Bigr]=\frac12E_r^P .
\end{equation*}
Moreover $\sigma_{(r,P)}=C_\sigma e_r\le\frac12E_r^P$ by the definition of
$C_E$ in \eqref{9.1:eq-rate-constants-a}.

\emph{Row $A$.} In \eqref{9.6:eq-two-run-estimate-5} with $g=E$: the first term is at most
$c_\rho\alpha_r^A\mathsf{A}C_E\mathsf{n}_re_r$ by (o5); the square bracket is evaluated with
$E_r^P=C_Ee_r$, the second bound of (o1) and (o3):
\begin{align*}
E_r^P+\sum_{s>r}\bigl(q^{s-r}(E_s^P+E_s^A)+E_s^R\bigr)
&\le C_Ee_r\Bigl[1+\frac{5\rho}{1-\rho}+8\Bigr]\\
&\le C_Ee_r\Bigl[9+\frac5{1-\rho}\Bigr];
\end{align*}
finally $E_r^R=4C_E\hat{\mathfrak{r}}_re_r$. By items (c), (d) and
(e) of \eqref{9.1:eq-coupling-ceilings-a}, each of $c_\rho\alpha_r^A\mathsf{A}\mathsf{n}_r$,
$\varkappa_A(9+5/(1-\rho))$ and $4\hat{\mathfrak{r}}_r$ is at most $\frac14$, whence the right-hand
side of \eqref{9.6:eq-two-run-estimate-5} at $g=E$, less its constant term, is at most
\begin{equation*}
C_Ee_r\Bigl[c_\rho\alpha_r^A\mathsf{A}\mathsf{n}_r
+\varkappa_A\Bigl(9+\frac5{1-\rho}\Bigr)+4\hat{\mathfrak{r}}_r\Bigr]\le\frac34C_Ee_r
\le\frac12E_r^A .
\end{equation*}
Also $\sigma_{(r,A)}=C_\sigma e_r\le\frac12C_Ee_r\le\frac12E_r^A$.

\emph{Row $R$.} Evaluating the bracket at the envelope with (o1)--(o4), the coupling distances
through (o2) and (o4), and the entries $s=r$ of both the $P$ and the $A$ components being included
in the term $5/(1-\rho)$, we obtain
\begin{equation*}
\begin{aligned}
2\mathfrak{r}_r\Lambda_rB_r(E)&\le2\mathfrak{r}_r\Lambda_rC_Ee_r\Bigl[
\mathsf{A}\Bigl(\frac{\mathsf{n}_r}{1-\rho}+\frac{\rho}{(1-\rho)^2}\Bigr)+\frac5{1-\rho}+8\Bigr]\\
&\le C_E\hat{\mathfrak{r}}_re_r=\frac14E_r^R,
\end{aligned}
\end{equation*}
the second inequality by the definition of $\hat{\mathfrak{r}}_r$ in
\eqref{9.1:eq-rate-constants-a}. For the band, let $r<s<r+\ell_r$ and $m=s-r$, so that
$1\le m<\ell_r$. The second inequality of item (f) of \eqref{9.1:eq-coupling-ceilings-a} gives
$\hat{\mathfrak{r}}_s\le2(1+m)\hat{\mathfrak{r}}_r$, and $e_s\le\rho^{-m}e_r\le\rho^{-\ell_r}e_r$;
moreover
\begin{equation*}
\sum_{1\le m<\ell_r}2(1+m)=(\ell_r-1)(\ell_r+2)\le2\ell_r^2 .
\end{equation*}
Hence, using item (g$'$) of \eqref{9.1:eq-coupling-ceilings-a} in the last step,
\begin{equation*}
2\varkappa'_r\sum_{r<s<r+\ell_r}E_s^R\le4\varkappa'_r\ell_r^2\rho^{-\ell_r}E_r^R\le\frac14E_r^R ,
\end{equation*}
and altogether the right-hand side of \eqref{9.6:eq-two-run-estimate-6} at $g=E$, less its
constant term, is at most $\frac12E_r^R$. For the constant term, the definitions of
$C_E$ and $\hat{\mathfrak{r}}_r$ in \eqref{9.1:eq-rate-constants-a} give the first inequality in
\begin{equation*}
\sigma_{(r,R)}=\mathfrak{r}_r\Lambda_r(2C_\sigma+1)e_r\le\frac{\hat{\mathfrak{r}}_r}{16}\,2C_Ee_r
\le\frac12E_r^R .
\end{equation*}

\emph{Step 4: the boundary and the conclusion.} For $r=n$, that is, for the indices of $I$ not in
$I_0$, the last assertion of (c3) of Lemma~\ref{9.6:lem-chain-structure} gives, with the constant
$C_V$ of Lemma~\ref{9.6:lem-chain-structure}, $g_n^c\le C_V\mathfrak{u}_n^c$ for every component
$c$; summing over the components of each block, with the reading units at depth $n$, gives
\begin{equation*}
g_n^P\le C_V\le C_Ee_n,\qquad g_n^A\le2C_V\mathsf{l}_n^{-1}\le4C_Ee_n,\qquad
g_n^R\le C_V\mathfrak{r}_n\le4C_E\hat{\mathfrak{r}}_ne_n ;
\end{equation*}
here $e_n=1+\Pi_n\ge1$ and $\hat{\mathfrak{r}}_n\ge16\mathfrak{r}_n$, and with these the three
right-hand inequalities follow from the definition of $C_E$ in terms of $C_V$ in
\eqref{9.1:eq-rate-constants-a}. The entries of $E$ are
positive. Thus all hypotheses of Lemma~\ref{9.6:lem-maximum-principle} hold, and it yields $g\le E$,
which is \eqref{9.6:eq-two-run-estimate-1}. Since $\mathsf{D}_r^c\le g_r^c$ by (c3) of
Lemma~\ref{9.6:lem-chain-structure} and $g_r^c\le g_r^{\bullet}$ for every component $c$ of the
block $\bullet$, the same bounds hold for the component distances.

Finally, for $r_0\le r<n$, \eqref{9.6:eq-two-run-estimate-7}, the bound $g_u^P\le C_Ee_u$ and (o2)
give
\begin{equation*}
d_{r+1}\le\mathsf{A}C_E\sum_{r_0\le u\le r}e_u
\le\mathsf{A}C_E\Bigl[\frac{\rho^{n-r}}{1-\rho}+(r-r_0+1)\Pi_n\Bigr]\le\mathsf{A}C_E\mathsf{n}_re_r ,
\end{equation*}
which is \eqref{9.6:eq-two-run-estimate-2}. Since the substituted expression for $d^{\sharp}_{r+1}$
has nonnegative entries and $g\le E$, we have $d^{\sharp}_{r+1}\le d^{\sharp}_{r+1}(E)$, and (o5)
gives \eqref{9.6:eq-two-run-estimate-3}.
\end{proof}

\begin{remark}[Closing the estimate with pair rows slaved to an exact competitor]
\label{9.6:rem-slaved-pair-rows}
In the proof of the two-run estimate with the power-law floor
(Proposition~\ref{9.6:prop-two-run-estimate}), the bounds for the $A$-, $Q$- and $R$-rows may
be supplied as pair rows, that is, as bounds for the pair of lists expressed through a component
$\mathsf{K}$ read on each constituent of the pair. Two such components occur in this
chapter: the two-cutoff seminorm $\mathsf{S}$ read on the slices of $X^{+}$, whose competitor is
built on $X$, and the seminorm $\mathsf{S}^{\mathcal{K}}$, whose competitor is built on $X^{+}$.
They are supplied, together with their pair rows and their competitors, respectively by the
statement of this chapter that furnishes $\mathsf{S}$ on the slices of $X^{+}$ with its
competitor built on $X$, and by the statement of this chapter that furnishes
$\mathsf{S}^{\mathcal{K}}$ with its competitor built on $X^{+}$. Each such family of pair rows
is slaved to an exact competitor $G$, that is, the bound
\begin{equation}\label{9.6:eq-slaved-pair-rows-1}
\|G\|^{c}\le C_X\bigl(\mathsf{K}+\hat{\vartheta}\,\|F\|\bigr)
\end{equation}
holds. Here $F$ is the family being compared and $\|F\|$ its norm, $\|\cdot\|^{c}$ is the norm
of the class of list bounds on the set $c$, and $C_X$ and $\hat{\vartheta}$ are the constant and
the rate entering this bound for the competitor built on $X$ or on $X^{+}$. In this situation the
bounds of Proposition~\ref{9.6:prop-two-run-estimate} hold with each of the quantities
$g_r^{P}$, $g_r^{A}$, $g_r^{R}$ that is fed by one of these rows replaced by
\begin{equation}\label{9.6:eq-slaved-pair-rows-2}
\max\{g_r^{\bullet},\,C_X\mathsf{K}\},
\end{equation}
and with the constants adjusted as follows:
\begin{itemize}
\item the constant $\mathsf{A}$ of the bound $|\beta_r[\cdot]|\le\mathsf{A}\|\cdot\|$ on the
step map assumed in Proposition~\ref{9.6:prop-two-run-estimate} becomes $2\mathsf{A}$;
\item every constant $\varkappa$ of the response rows, with its indices, is multiplied by
$C_X^2$;
\item the weight $\Lambda_r$ becomes $C_X\Lambda_r$;
\item the constant $C_E$ is chosen so as to include $2C_X$ among the quantities it is built
from.
\end{itemize}
This is the non-circular form of the comparison families \eqref{9.1:eq-comparison-families-a}.
The resulting bounds are linear in $\mathsf{K}$.
\end{remark}

\begin{proof}
The bounds on the $A$-, $Q$- and $R$-rows enter the proof of
Proposition~\ref{9.6:prop-two-run-estimate} only at the point at which the row bounds are
substituted. From that point on, the proof appeals only to the following inputs:
\begin{itemize}
\item the three bounds for $g_r^{P}$, $g_r^{A}$, $g_r^{R}$ in the statement of
Proposition~\ref{9.6:prop-two-run-estimate};
\item the bound on the coupling distance $d_{r+1}$ in the same statement;
\item the boundary condition imposed at the initial depth in
Proposition~\ref{9.6:prop-two-run-estimate};
\item the ceilings \eqref{9.1:eq-coupling-ceilings-a}.
\end{itemize}
Substituting the pair rows together with the slaving inequality
\eqref{9.6:eq-slaved-pair-rows-1} in place of the row bounds yields the same inequalities with
each quantity fed by these rows replaced by \eqref{9.6:eq-slaved-pair-rows-2} and the constants
changed as listed in the statement; since none of the inputs read after that point involves the
row bounds, the argument from there on is unchanged. The linearity of the resulting bounds in
$\mathsf{K}$ is the content of the statement of this chapter asserting that the pair-row bounds
depend linearly on $\mathsf{K}$.
\end{proof}

\begin{proposition}[Matching the two runs at depth zero]\label{9.6:prop-depth-zero-match}
Work in the setting of Definition~\ref{9.1:def-tuned-family}, under its standing hypotheses
(J1)--(J5) (see \eqref{9.1:eq-tuned-family-a}) and under the hypotheses of the two-run estimate
with the power-law floor, Proposition~\ref{9.6:prop-two-run-estimate}; assume moreover
$\kappa_0\ge10$, where $\kappa_0$ is the exponent in the bound on $a_u$ in (J2). For $r>0$ write
$D(0,r):=\{v\in\mathbb{C}:|v|<r\}$. For a run of length $n$ with $\zeta_{(n)}=\zeta$, let
$\zeta^{n}_{(0)}(\zeta)$ denote its running shift at depth zero. Let $\zeta_*^n$ be the real alive
run of length $n$ given by (J3), and put
\begin{equation}\label{9.6:eq-depth-zero-match-1}
h_n(v):=\zeta^{n}_{(0)}\bigl(\zeta_*^n+v\bigr),\qquad v\in D(0,R).
\end{equation}
Let $r_J$ be the radius of (J3), let $a_u$ be as in \eqref{9.1:eq-tuned-family-a}, and let
$\mathsf{A}$, $C_E$, $\rho$, $e_u$ be as in (J1) and \eqref{9.1:eq-rate-constants-a}.
Then for every $n$ the following hold.
\begin{enumerate}
\item[(a)] For every $\hat w\in\mathbb{C}$ with $|\hat w|\le r_J$ there is exactly one
$\mathfrak{n}_n(\hat w)\in D(0,\tfrac83 r_J)$ with
\begin{equation}\label{9.6:eq-depth-zero-match-2}
h_{n+1}\bigl(\mathfrak{n}_n(\hat w)\bigr)=h_n(\hat w),
\end{equation}
and $|\mathfrak{n}_n(\hat w)|\le 2|\hat w|$.
\item[(b)] With
\begin{equation}\label{9.6:eq-depth-zero-match-3}
\Upsilon_n:=2a_n+\sum_{u<n}\min\Bigl\{3a_u,\ \frac{2\mathsf{A}C_E e_u}{r_J}\Bigr\}
\end{equation}
one has, for $|\hat w|\le r_J$,
\begin{equation}\label{9.6:eq-depth-zero-match-4}
|\mathfrak{n}_n(\hat w)-\hat w|\le \Upsilon_n|\hat w|,
\qquad
\Upsilon_n\le C\,n^{-3/2}(\log n)^{p_0+1}\quad(n\ge2),
\end{equation}
with a constant $C$ independent of $n$; the second bound is asserted for $n\ge2$ only, because
its right-hand side vanishes at $n=1$ ($\log 1=0$).
\item[(c)] With $N_n:=h_n'(0)$, one has $N_{n+1}\,\mathfrak{n}_n'(0)=N_n$.
\end{enumerate}
\end{proposition}

\begin{proof}
Throughout, the run of length $n$ is denoted by the superscript $\mathsf{A}$ and the run of
length $n+1$ by the superscript $\mathsf{B}$; $b^{\mathsf{A}}_{(u)}(\zeta)$ and
$b^{\mathsf{B}}_{(u)}(\zeta)$ denote the step shifts at depth $u$ of these runs started at
$\zeta_{(n)}=\zeta$, respectively $\zeta_{(n+1)}=\zeta$.

\emph{Step 1: existence, uniqueness and holomorphy of $\mathfrak{n}_n$.} Fix $\hat w$ with
$|\hat w|\le r_J$. The functions $h_n$ and $h_{n+1}$ are holomorphic on $D(0,R)$, and the closed
disc $|v|\le\tfrac83 r_J$ is contained in $D(0,R)$; these are the facts on which the application
of Rouch\'e's theorem below rests. We apply Rouch\'e's theorem to the function
$v\mapsto h_{n+1}(v)-h_n(\hat w)$ on the single circle $|v|=\tfrac83 r_J$; the comparison on
this circle uses only (J2) and (J3) of Definition~\ref{9.1:def-tuned-family}. It gives that this
function has exactly one zero in $D(0,\tfrac83 r_J)$, which we call $\mathfrak{n}_n(\hat w)$;
this is \eqref{9.6:eq-depth-zero-match-2} together with its uniqueness. The same comparison gives
$|\mathfrak{n}_n(\hat w)|\le 2|\hat w|$. In particular $\mathfrak{n}_n(0)=0$. Since the zero is
simple and $h_{n+1}(v)-h_n(\hat w)\neq0$ on the circle, the argument principle gives
\begin{equation*}
\mathfrak{n}_n(\hat w)=\frac{1}{2\pi i}\oint_{|v|=\frac83 r_J}
\frac{v\,h_{n+1}'(v)}{h_{n+1}(v)-h_n(\hat w)}\,dv ,
\end{equation*}
and the right-hand side is holomorphic in $\hat w\in D(0,r_J)$, because $h_n$ is holomorphic and
the denominator does not vanish on the contour. This completes (a).

\emph{Step 2: the identity (c).} Since $\mathfrak{n}_n(0)=0$ and $\mathfrak{n}_n$ is
holomorphic, differentiating \eqref{9.6:eq-depth-zero-match-2} at $\hat w=0$ gives
\begin{equation*}
h_{n+1}'(0)\,\mathfrak{n}_n'(0)=h_n'(0),
\end{equation*}
which is the identity $N_{n+1}\mathfrak{n}_n'(0)=N_n$ of (c), since $N_n=h_n'(0)$ and
$N_{n+1}=h_{n+1}'(0)$.

\emph{Step 3: the input from the two-run estimate.} For $u<n$ and $|\hat w|\le r_J$ put
\begin{equation}\label{9.6:eq-depth-zero-match-5}
\delta b_u(\hat w):=b^{\mathsf{B}}_{(u)}\bigl(\zeta_*^{n+1}+\mathfrak{n}_n(\hat w)\bigr)
-b^{\mathsf{A}}_{(u)}\bigl(\zeta_*^{n}+\hat w\bigr).
\end{equation}
By Step 1, $\hat w\mapsto\zeta_*^{n+1}+\mathfrak{n}_n(\hat w)$ is holomorphic on $D(0,r_J)$, so
$\delta b_u$ is holomorphic on $D(0,r_J)$. Consider the frozen pair of lists, in the sense of
Proposition~\ref{9.6:prop-two-run-estimate}, formed by the run
$\mathsf{A}$ started at $\zeta_*^n+\hat w$ and the run $\mathsf{B}$ started at
$\zeta_*^{n+1}+\mathfrak{n}_n(\hat w)$. Their running shifts at depth zero are $h_n(\hat w)$ and
$h_{n+1}(\mathfrak{n}_n(\hat w))$, which coincide by \eqref{9.6:eq-depth-zero-match-2}. Hence
Proposition~\ref{9.6:prop-two-run-estimate} applies to this frozen pair with $r_0=0$ and gives
$g_u^P\le C_E e_u$ for $0\le u\le n$, where $g_u^P$ is the quantity bounded at depth $u$ in that
proposition, taken for this pair. With the constant $\mathsf{A}$ of (J1), which bounds
$|\beta_r[\cdot]|$ by $\mathsf{A}\|\cdot\|$, we obtain
\begin{equation}\label{9.6:eq-depth-zero-match-6}
|\delta b_u(\hat w)|\le \mathsf{A}\,g_u^P\le \mathsf{A}C_E e_u ,\qquad u<n .
\end{equation}

\emph{Step 4: the bound on $\mathfrak{n}_n(\hat w)-\hat w$.} By the recursion
$\zeta_{(r)}=\zeta_{(r+1)}+b_{(r)}-c_0$ of Definition~\ref{9.1:def-tuned-family}, summed from
$r=n-1$ down to $r=0$,
\begin{equation*}
\zeta^{n}_{(0)}(\zeta)=\zeta+\sum_{u=0}^{n-1}\bigl(b^{\mathsf{A}}_{(u)}(\zeta)-c_0\bigr),
\end{equation*}
and likewise for the run $\mathsf{B}$ with $n$ replaced by $n+1$. Inserting these two formulas
into \eqref{9.6:eq-depth-zero-match-2}, with $\mathfrak{n}:=\mathfrak{n}_n(\hat w)$, gives
\begin{equation*}
\zeta_*^{n+1}-\zeta_*^{n}+\mathfrak{n}-\hat w
+b^{\mathsf{B}}_{(n)}\bigl(\zeta_*^{n+1}+\mathfrak{n}\bigr)-c_0
+\sum_{u<n}\delta b_u(\hat w)=0 .
\end{equation*}
At $\hat w=0$, where $\mathfrak{n}_n(0)=0$, the same identity reads
\begin{equation*}
\zeta_*^{n+1}-\zeta_*^{n}+b^{\mathsf{B}}_{(n)}\bigl(\zeta_*^{n+1}\bigr)-c_0
+\sum_{u<n}\delta b_u(0)=0 .
\end{equation*}
Subtracting the second identity from the first,
\begin{equation}\label{9.6:eq-depth-zero-match-7}
\begin{split}
\mathfrak{n}_n(\hat w)-\hat w
={}&-\Bigl[b^{\mathsf{B}}_{(n)}\bigl(\zeta_*^{n+1}+\mathfrak{n}_n(\hat w)\bigr)
-b^{\mathsf{B}}_{(n)}\bigl(\zeta_*^{n+1}\bigr)\Bigr]\\
&-\sum_{u<n}\bigl[\delta b_u(\hat w)-\delta b_u(0)\bigr].
\end{split}
\end{equation}
For each $u<n$ the function $\delta b_u-\delta b_u(0)$ is holomorphic on $D(0,r_J)$, vanishes at
$0$, and by \eqref{9.6:eq-depth-zero-match-6} has modulus at most $2\mathsf{A}C_E e_u$ there.
By Schwarz's lemma,
\begin{equation*}
|\delta b_u(\hat w)-\delta b_u(0)|\le \frac{2\mathsf{A}C_E e_u}{r_J}\,|\hat w| ,
\qquad |\hat w|<r_J ,
\end{equation*}
and this extends to $|\hat w|=r_J$ by continuity. The alternative bounds come from (J2), by which
the step shift at depth $u$ changes by at most $a_u$ times the change of the starting shift.
Since $\mathfrak{n}_n(0)=0$ and $|\mathfrak{n}_n(\hat w)|\le2|\hat w|$ by Step~1, the triangle
inequality gives
\begin{equation*}
\begin{split}
|\delta b_u(\hat w)-\delta b_u(0)|
&\le\bigl|b^{\mathsf{B}}_{(u)}\bigl(\zeta_*^{n+1}+\mathfrak{n}_n(\hat w)\bigr)
-b^{\mathsf{B}}_{(u)}\bigl(\zeta_*^{n+1}\bigr)\bigr|
+\bigl|b^{\mathsf{A}}_{(u)}\bigl(\zeta_*^{n}+\hat w\bigr)-b^{\mathsf{A}}_{(u)}\bigl(\zeta_*^{n}\bigr)\bigr|\\
&\le 2a_u|\hat w|+a_u|\hat w|=3a_u|\hat w| ,
\end{split}
\end{equation*}
and in the same way the bracket in the first line of \eqref{9.6:eq-depth-zero-match-7} is at most
$a_n\cdot2|\hat w|=2a_n|\hat w|$. Taking for each $u<n$ the
smaller of the two bounds and summing gives the first inequality in
\eqref{9.6:eq-depth-zero-match-4} with $\Upsilon_n$ as in \eqref{9.6:eq-depth-zero-match-3}.

\emph{Step 5: the size of $\Upsilon_n$.} By (J2), $\sum_u a_u\le\tfrac13$, so
$\Upsilon_n\le 2a_n+3\sum_{u<n}a_u\le\tfrac53$ for every $n$; hence for any fixed $n_1$ the
bound in \eqref{9.6:eq-depth-zero-match-4} holds for $2\le n\le n_1$ after enlarging $C$. Now
split the sum at
\begin{equation*}
u_0:=n-\Bigl\lceil\frac{2\log n}{|\log\rho|}\Bigr\rceil .
\end{equation*}
For $u<u_0$ use the second entry of the minimum, with $e_u=\rho^{n-u}+\Pi_n$. Since
$\rho^{n-u_0}\le\rho^{2\log n/|\log\rho|}=n^{-2}$,
\begin{equation*}
\sum_{u<u_0}\frac{2\mathsf{A}C_E e_u}{r_J}
\le\frac{2\mathsf{A}C_E}{r_J}\Bigl(\frac{n^{-2}}{1-\rho}+n\,\Pi_n\Bigr),
\end{equation*}
and $n\,\Pi_n\le C\,n^{-3/2}(\log n)^{p_0+1}$ by the definition of the power-law floor $\Pi_n$
in \eqref{9.1:eq-rate-constants-a}. For $u_0\le u<n$ use the first entry $3a_u$. Take $n_1$ so
large that $u_0\ge n/2$ and $c_0n/2\ge e^{2p_0/3}$ for $n>n_1$. By (J2), $a_u$ is at most a
constant independent of $u$ times
$\theta_u^{-3/2}(\log\theta_u)^{p_0}+\theta_u^{-(\kappa_0-4)/2}/R$. Since
$\theta_u=X+c_0u\ge c_0n/2$ for $u\ge u_0$, since $\theta\mapsto\theta^{-3/2}(\log\theta)^{p_0}$
is decreasing for $\theta\ge e^{2p_0/3}$, and since $\kappa_0\ge10$ gives
$(\kappa_0-4)/2\ge3$, each such $a_u$, and also $a_n$, is at most a constant times
$n^{-3/2}(\log n)^{p_0}$. There are at most $2\log n/|\log\rho|+1$ indices with $u_0\le u<n$, so
\begin{equation*}
2a_n+\sum_{u_0\le u<n}3a_u\le C\,n^{-3/2}(\log n)^{p_0+1}.
\end{equation*}
Adding the two parts gives the second inequality in \eqref{9.6:eq-depth-zero-match-4}.

\emph{Order of the constructions.} The argument is not circular. The map $\mathfrak{n}_n$ is
constructed in Step 1 from (J2) and (J3) alone; the two-run estimate,
Proposition~\ref{9.6:prop-two-run-estimate}, is then applied in Step 3 only to frozen pairs; the
bound on $\Upsilon_n$ follows in Steps 4 and 5. Local pairs, the other kind of pair admitted in
Proposition~\ref{9.6:prop-two-run-estimate}, are formed only afterwards, under the
condition $\Upsilon_n\le\tfrac14$, and the two-run estimate for local pairs is applied to them
only then; no step uses a conclusion that depends on it.
\end{proof}

\begin{lemma}[Pointwise comparison of the running shifts under a common source]
\label{9.6:lem-pointwise-shift}
Let $\mathsf{A}$, $\mathsf{B}$ be a pair of runs of lengths $n$ and $n+1$ with one common source
and the transplant, as in Definition~\ref{9.1:def-run-pair-transplant}, so that
\eqref{9.1:eq-run-pair-transplant-b} holds, and assume the hypotheses of
Proposition~\ref{9.6:prop-two-run-estimate} with $r_0=0$; in particular the pair is alive to
depth $0$ and $\zeta^{\mathsf{B}}_{(0)}=\zeta^{\mathsf{A}}_{(0)}$ on $\mathfrak{P}$. In the
estimate below, $\mathsf{A}$ also denotes the constant of the bound
$|\beta_r[\,\cdot\,]|\le\mathsf{A}\|\cdot\|$ assumed in
Proposition~\ref{9.6:prop-two-run-estimate}, which carries the same letter as the run
$\mathsf{A}$; $C_E$, $\rho$, $\Pi_n$ are the constants, and $e_r=\rho^{\,n-r}+\Pi_n$ the
envelope, of that proposition. For $x$ on the continuum
torus let $\omega(x):=u+W^{\mathsf{A}}(x)$ be the local source at $x$, where $W^{\mathsf{A}}$ is
the residual source of the run $\mathsf{A}$ and $u\in\mathbb{C}$ is the $x$-independent part of
the local source. Let $r_J$ be the source radius, and assume
$\omega(x)\in\bar D(0,r_J/2)$, the closed disc of radius $r_J/2$ about $0$.
Let $\varphi$ be the localising function of the pair, let $r\ge\mathsf{w}$, and let $x$ be a
point at which $\varphi\equiv1$ down to depth $r$. Then
\begin{equation}\label{9.6:eq-pointwise-shift-1}
\bigl|\zeta^{\mathsf{B}}_{(r)}(x)-\zeta^{\mathsf{A}}_{(r)}(x)\bigr|
\le d_r+\Bigl[\frac{\mathsf{A}C_E}{2(1-\rho)}+r_J\Bigr]e_r
\le C_{PM}(1+r)\,e_r ,
\end{equation}
where the constant $C_{PM}$ is
\begin{equation*}
C_{PM}:=\frac{3\,\mathsf{A}C_E}{2(1-\rho)}+r_J .
\end{equation*}
\end{lemma}

\begin{proof}
The first inequality of \eqref{9.6:eq-pointwise-shift-1} is the local-source bound for the
running shift, applied at $x$: it splits the difference of the running
shifts at $x$ into the coupling distance $d_r$ of the two runs at depth $r$ and the contribution
of the local source $\omega(x)\in\bar D(0,r_J/2)$, and bounds that contribution by
$\bigl[\mathsf{A}C_E/(2(1-\rho))+r_J\bigr]e_r$. We prove the second
inequality.

We bound $d_r$ by Proposition~\ref{9.6:prop-two-run-estimate}, which
applies with $r_0=0$ by hypothesis. A depth $r$ of the run $\mathsf{A}$ satisfies $r\le n$,
since $\mathsf{A}$ has length $n$. If $r=0$, then $d_0=0$, since the pair is matched at
depth $0$ in the sense of Definition~\ref{9.1:def-run-pair-transplant}, so that its couplings
coincide at depth $0$; the bound on $d_r$ displayed
below then holds, its right-hand side being non-negative. If $r\ge1$, then $0\le r-1<n$, and the
proposition, applied with $r-1$ in place of $r$, gives
\begin{equation*}
d_r\le \mathsf{A}C_E\Bigl[\frac{\rho^{\,n-r+1}}{1-\rho}+r\,\Pi_n\Bigr].
\end{equation*}
Since $0<\rho<1$ we have $\rho^{\,n-r+1}\le\rho^{\,n-r}$ and $1\le 1/(1-\rho)$. With
$e_r=\rho^{\,n-r}+\Pi_n$ and $\Pi_n\ge0$ this yields
\begin{equation*}
d_r\le \mathsf{A}C_E\Bigl[\frac{\rho^{\,n-r}}{1-\rho}+r\,\Pi_n\Bigr]
\le \frac{\mathsf{A}C_E}{1-\rho}\,(1+r)\,e_r .
\end{equation*}
Moreover $1\le 1+r$, so
\begin{equation*}
\Bigl[\frac{\mathsf{A}C_E}{2(1-\rho)}+r_J\Bigr]e_r
\le\Bigl[\frac{\mathsf{A}C_E}{2(1-\rho)}+r_J\Bigr](1+r)\,e_r .
\end{equation*}
Adding the last two bounds gives the second inequality of
\eqref{9.6:eq-pointwise-shift-1}, since
\begin{equation*}
\frac{\mathsf{A}C_E}{1-\rho}+\frac{\mathsf{A}C_E}{2(1-\rho)}+r_J
=\frac{3\,\mathsf{A}C_E}{2(1-\rho)}+r_J=C_{PM} .
\end{equation*}
At points of the layers, where $\varphi$ is not identically $1$ down to depth $r$, the
comparison of the two running shifts is made through the residual sources
$W^{\mathsf{c}}$, $\mathsf{c}\in\{\mathsf{A},\mathsf{B}\}$, of the two runs and is not asserted
here.
\end{proof}

\begin{lemma}[The match at depth zero and torus-free offsets]\label{9.6:lem-torus-free-offset}
Let $\mathsf{A}$ and $\mathsf{B}$ be the two runs, of lengths $n$ and $n+1$, compared in
Proposition~\ref{9.6:prop-two-run-estimate}. Let $\mathsf{w}=\mathsf{w}(g)$ be the frozen depth, let
$d_r$ be the coupling distance between the two runs at depth $r$, and let
$e_r=\rho^{n-r}+\Pi_n$ be as in that proposition.
\begin{itemize}
\item[(i)] The pinning and the matching of the two runs take place at depth $0$ of the class
$\mathfrak{S}$ of Definition~\ref{9.1:def-tuned-family}. Consequently
Proposition~\ref{9.6:prop-two-run-estimate} applies with $r_0=0$.
\item[(ii)] At the frozen depth,
\begin{equation}\label{9.6:eq-torus-free-offset-1}
d_{\mathsf{w}}\le \mathsf{A}\,C_E\,\max\bigl\{(1-\rho)^{-1},\mathsf{w}(g)\bigr\}\,e_{\mathsf{w}} .
\end{equation}
Here $\mathsf{A}$ denotes the constant of the bound $|\beta_r[\cdot]|\le\mathsf{A}\|\cdot\|$ in
\textup{(J1)} of Definition~\ref{9.1:def-tuned-family}, not the run $\mathsf{A}$, and $C_E$ is the
constant of Proposition~\ref{9.6:prop-two-run-estimate}.
The factor $\max\{(1-\rho)^{-1},\mathsf{w}(g)\}$ depends on $g$ only. The coupling distance
$d_{\mathsf{w}}$ is in general not zero; this distance, through the bound
\eqref{9.6:eq-torus-free-offset-1}, is carried by the constants $C_{PM}$, $C_\eta$ and $C_{TC}$
of the increment bound used in the proof of the uniqueness theorem.
\item[(iii)] The offset $a_n$ of \eqref{9.1:eq-tuned-family-a} does not depend on the torus.
\end{itemize}
\end{lemma}

\begin{proof}
(i) The two runs are pinned at depth $0$ of the class $\mathfrak{S}$ of
Definition~\ref{9.1:def-tuned-family} and matched there. In the notation
of Proposition~\ref{9.6:prop-two-run-estimate}, this means that their running shifts coincide at
depth $0$ on the parameter set, $\zeta^{\mathsf{B}}_{(0)}=\zeta^{\mathsf{A}}_{(0)}$ on
$\mathfrak{P}$. This is the matching hypothesis of that proposition with $r_0=0$. The
proposition therefore applies with $r_0=0$: its bounds on $g^P_r$, $g^A_r$, $g^R_r$ hold for
$0\le r\le n$, and its bounds on $d_{r+1}$, $d^{\sharp}_{r+1}$ hold for $0\le r<n$.

(ii) We use the bound on $d_{r+1}$ of Proposition~\ref{9.6:prop-two-run-estimate}, with $r_0=0$
by part (i) and with $r=\mathsf{w}-1$, admissible since $1\le\mathsf{w}\le n$. Then
$r-r_0+1=\mathsf{w}$ and $n-r=n-\mathsf{w}+1$, so
\begin{equation}\label{9.6:eq-torus-free-offset-2}
d_{\mathsf{w}}\le \mathsf{A}\,C_E\Bigl[\frac{\rho^{\,n-\mathsf{w}+1}}{1-\rho}
+\mathsf{w}\,\Pi_n\Bigr].
\end{equation}
Since $0\le\rho<1$, we have $\rho^{\,n-\mathsf{w}+1}\le\rho^{\,n-\mathsf{w}}$. Bounding each of
the two coefficients $(1-\rho)^{-1}$ and $\mathsf{w}$ by their maximum gives
\begin{align*}
\frac{\rho^{\,n-\mathsf{w}+1}}{1-\rho}+\mathsf{w}\,\Pi_n
&\le \max\bigl\{(1-\rho)^{-1},\mathsf{w}\bigr\}\bigl(\rho^{\,n-\mathsf{w}}+\Pi_n\bigr)\\
&= \max\bigl\{(1-\rho)^{-1},\mathsf{w}\bigr\}\,e_{\mathsf{w}} .
\end{align*}
Inserting this into \eqref{9.6:eq-torus-free-offset-2} gives
\eqref{9.6:eq-torus-free-offset-1}. The rate $\rho$ is a fixed constant, and
$\mathsf{w}=\mathsf{w}(g)$ is a function of $g$. Hence the factor
$\max\{(1-\rho)^{-1},\mathsf{w}(g)\}$ depends on $g$ only.

(iii) By \eqref{9.1:eq-tuned-family-a}, $a_n=\theta_n-\zeta^n_*(0)$, where $\zeta^n_*$ is the
unique real alive run of length $n$. This run is built by the recursion
$\zeta_{(r)}=\zeta_{(r+1)}+b_{(r)}-c_0$ of Definition~\ref{9.1:def-tuned-family} from the step
shifts $b_{(r)}$. At the depths $r<\mathsf{w}$, the steps of the run read frozen step shifts
$b_{(r)}$, and these do not depend on the torus. This is the
reason for which the offset $a_n$ does not depend on the torus.
\end{proof}

\begin{remark}[Three notes on the threshold, the class constant and the radii]
\label{9.6:rem-threshold-notes}
We use the reading units, components and distances of
Definition~\ref{9.1:def-totals-units-distance} and the hybrid lists of
Definition~\ref{9.1:def-hybrid-lists}.
\begin{itemize}
\item[(i)] \emph{Threshold.} The choice $\mathfrak{u}_s^{\mathsf{P}}:=1$ of the reading unit of the
component $\mathsf{P}$ is a valid threshold for the transplant $\mathfrak{T}F$ of a
$\mathsf{B}$-family $F$.
\item[(ii)] \emph{Class constant.} The class constant of \eqref{9.1:eq-hybrid-lists-a} is
\begin{equation*}
C_{\mathfrak{K}}=3C_V+5 .
\end{equation*}
The per-component dichotomy of Lemma~\ref{9.6:lem-chain-structure} and the numerical fractions
in the proof of Proposition~\ref{9.6:prop-two-run-estimate} (see
\eqref{9.6:eq-two-run-estimate-a}) are used as stated there. The input bounds formulated for
lists in the class $(C_V+3)\mathfrak{K}^c$ hold on the class $C_{\mathfrak{K}}\mathfrak{K}^c$;
these inputs are hypothesis (a) of the transplanted one-lattice response bounds (the statement
containing \eqref{9.1:eq-transplant-response-a}), the bounds on
the blocks $P$, $A$ and on the component $R''$ assumed in
Proposition~\ref{9.6:prop-two-run-estimate}, the localisation statement for local pair lists
(its class of local members, in which a frozen member is a local member at constant background,
and its constants $E^{\sharp}$, $K_2'$, $\varrho_\mu$), and the bounds for $P^{\flat}$ on local
pairs. This reading imposes no cap in the sense of Definition~\ref{2.3:def-cap}. The
localisation statement for local pair lists and parts (c) and (d) of the uniqueness theorem,
\eqref{9.8:eq-uniqueness-theorem-a}, depend on it.
\item[(iii)] \emph{Radii.} In this item $r$ and $r_j$ denote radii, not the depth index $r$ of
Proposition~\ref{9.6:prop-two-run-estimate} and of the proof of (ii). The radius bound for pair
lists, $r_j\le 1.6\,r$ for the radii $r_j$ of the correlation theorem, rests on the two facts
\begin{equation*}
r_j\in[r,2r),\qquad \frac{1/4}{1-L^{-1/2}}\le 0.6\quad\text{for } L\ge 3 .
\end{equation*}
Definition~\ref{9.1:def-hybrid-lists} uses
the inequality $(C_V+2)\mathfrak{m}\le E_0$, which is one of the conditions defining the
coupling ceiling $\gamma_U$.
\end{itemize}
\end{remark}

\begin{proof}
(i) Both members of the transplant $\mathfrak{T}F$ are runs. The $\mathsf{P}$-slot of a run at
level $k$ obeys the run bound of item (a) of the uniform run bounds,
\begin{equation*}
\|H^{(k)}\|_k\le \tfrac{9}{16}E_0 ,
\end{equation*}
where, in the notation of that statement, $H^{(k)}$ is the $\mathsf{P}$-slot at level $k$ and
$\|\cdot\|_k$ the norm at level $k$.
The class $\mathfrak{K}$ contains no clause bounding an individual $\mathsf{P}$-slot. Its
cumulative bound $|\mathsf{P}^{(\le j)}|_j\le 2\mathsf{E}$, where $\mathsf{E}$ is the constant of
the uniform run bounds, is the one applied to the frozen slots.
Hence the threshold $\mathfrak{u}_s^{\mathsf{P}}=1$ is valid for both members of
$\mathfrak{T}F$.

(ii) The passage from the class radius $C_V+3$ to $C_{\mathfrak{K}}=3C_V+5$ costs the factor
\begin{equation*}
\frac{C_{\mathfrak{K}}}{C_V+3}=\frac{3C_V+5}{C_V+3}\le 3,
\end{equation*}
because $C_V\ge 0$ and $3C_V+5\le 3C_V+9$. This factor depends neither on $g$ nor on $n$. The
class radius enters only the constants $K_2'$ and $\varrho_\mu$, the ceiling on the unmarked
activities, the smallness condition formulated with the constant $\alpha^{\sharp}$ of the value
family, the companion smallness condition with weights, the constant $\mathsf{E}$ of the uniform
run bounds, the constant $E_1$, and one factor $C_\sigma$. Each of these enters only through a
ceiling on $\gamma$ or on the smallness parameter $\varepsilon_1$, and those ceilings are fixed
after $C_V$. No condition bounding
$L$, $M$, $M_1$ or any other quantity of Definition~\ref{2.3:def-cap} from above therefore
arises. That the inputs listed in (ii) hold on the class $C_{\mathfrak{K}}\mathfrak{K}^c$ is
the content of item (v) of the setting statement accompanying the localisation statement for
local pair lists, where these inputs are derived again on that class. With this choice the
arguments of
Lemma~\ref{9.6:lem-chain-structure} and Proposition~\ref{9.6:prop-two-run-estimate} are
unchanged.

An alternative to the class $C_{\mathfrak{K}}\mathfrak{K}^c$, not used in this book, replaces
the block $A$ by a single local family. In the notation of
Definition~\ref{9.1:def-totals-units-distance} this family is
\begin{equation*}
\mathcal{P}:=E(\cdot\,;W)-Q_{W(p_y)}=P^{\flat}_{W(p_y)}+D(\cdot\,;W),
\end{equation*}
with unit $\mathfrak{u}_s^{\mathcal{P}}:=\mathsf{l}_s^{-1}(1+\mathfrak{m}/E_0)$. A frozen member
is the local member at constant background, and $D:=\mathcal{P}-\mathcal{P}|_{\mathrm{const}}$,
so that $\|G^D\|\le 2g^{\mathcal{P}}$. A hybrid local member is then bounded by
$E_0+(C_V+1)(E_0+2\mathfrak{m})$. Under the ceiling $2(C_V+1)\mathfrak{m}\le E_0$ this is at
most
\begin{equation*}
E_0+(C_V+1)E_0+E_0=(C_V+3)E_0 ,
\end{equation*}
so the class is $(C_V+3)\mathfrak{K}^c$.

The price of this alternative is a factor at most $2$ on the $A$-entries of the rows (P), (A),
(R) of \eqref{9.1:eq-transplant-response-a} and on $N_r^{A}$. The reason is that those rows
bound the maximum over $\{P^{\flat},D\}$, whereas the merged family requires the sum. Meeting
this factor would require halving the ceilings (b$'$) and (d) of
\eqref{9.1:eq-coupling-ceilings-a} and replacing $5$ by $9$ in $\hat{\mathfrak{r}}_r$, and
then establishing Proposition~\ref{9.6:prop-two-run-estimate} again with these changes.

(iii) By the radius statement of the correlation theorem, $r_j\in[r,2r)$. The pair contributes
the factor
$\tfrac14/(1-L^{-1/2})$. This factor is at most $0.6$ whenever $L\ge 3$: in that case
$L^{-1/2}\le 3^{-1/2}<0.58$, so $1-L^{-1/2}>0.42$ and $\tfrac14/0.42<0.6$.
\end{proof}

\begin{lemma}[Maximum principle for coupled inequalities with a margin]
\label{9.6:lem-maximum-principle-slack}
Let $I$ be a finite set and $I_0\subset I$. Let $D=(D_i)_{i\in I}$ and
$\sigma=(\sigma_i)_{i\in I_0}$ be families of non-negative real numbers, let
$E=(E_i)_{i\in I}$ be a family of strictly positive real numbers, let
$\mu=(\mu_{ij})_{i\in I_0,\,j\in I}$ be a family of non-negative real numbers, and let
$0\le\epsilon<\tfrac12$. Suppose that for every $i\in I_0$
\begin{equation}\label{9.6:eq-maximum-principle-slack-1}
D_i\le\sum_{j\in I}\mu_{ij}D_j+\sigma_i,\qquad
\sum_{j\in I}\mu_{ij}E_j\le\Bigl(\frac12+\epsilon\Bigr)E_i,\qquad
\sigma_i\le\frac12\,E_i,
\end{equation}
and that $D_i\le E_i$ for every $i\in I\setminus I_0$. Then
\begin{equation}\label{9.6:eq-maximum-principle-slack-2}
D_i\le\frac{E_i}{1-2\epsilon}\qquad\text{for every } i\in I.
\end{equation}
\end{lemma}

\begin{proof}
Since $I$ is finite and $E_i>0$ for every $i$, the number
\begin{equation*}
t:=\max_{i\in I}\frac{D_i}{E_i}
\end{equation*}
is well defined, and $t\ge0$ because $D\ge0$. By definition $D_j\le tE_j$ for every $j\in I$.
Since $0\le\epsilon<\tfrac12$, we have $0<1-2\epsilon\le1$.

If $t\le1$, then for every $j\in I$ we get $D_j\le tE_j\le E_j\le E_j/(1-2\epsilon)$, which
is \eqref{9.6:eq-maximum-principle-slack-2}.

Suppose $t>1$, and let $i\in I$ be an index at which the maximum is attained, so that
$D_i=tE_i$. If $i\notin I_0$, the hypothesis $D_i\le E_i$ would give $t\le1$; hence
$i\in I_0$, and the three inequalities \eqref{9.6:eq-maximum-principle-slack-1} hold at $i$.
Using first the inequality for $D_i$, then $\mu_{ij}\ge0$ together with $D_j\le tE_j$, and
finally the bound on $\sum_j\mu_{ij}E_j$ together with $\sigma_i\le\frac12E_i$, we obtain
\begin{equation*}
tE_i=D_i\le\sum_{j\in I}\mu_{ij}D_j+\sigma_i
\le t\sum_{j\in I}\mu_{ij}E_j+\sigma_i
\le t\Bigl(\frac12+\epsilon\Bigr)E_i+\frac12\,E_i .
\end{equation*}
Dividing by $E_i>0$ gives $t\le t(\frac12+\epsilon)+\frac12$, that is,
$t(\frac12-\epsilon)\le\frac12$. As $\frac12-\epsilon>0$, this yields $t\le1/(1-2\epsilon)$,
and therefore $D_j\le tE_j\le E_j/(1-2\epsilon)$ for every $j\in I$, which is
\eqref{9.6:eq-maximum-principle-slack-2}.
\end{proof}

\begin{theorem}[Uniqueness up to one real parameter, the renormalised coupling at the reference length (clause (v) of Theorem~0, as this chapter proves it)]
\label{9.7:thm-uniqueness}
Let $N=2$; the quantifier prefix is that of Notation~\ref{2.1:not-prefix}, and hypothesis (H7) of
Notation~\ref{2.2:not-hypotheses} is assumed. The pure numbers on which this statement depends are
among the Stage-0 data of the order of choice (Notation~\ref{2.3:not-stages-first}), and we put
$\lambda_U:=c_0/2$, where $c_0$ is the one-loop constant of Definition~\ref{1.2:def-flow-constants}.
For every $\varrho>0$ there is $\gamma_U>0$ such that for every
$g\le\gamma_U$, with $X:=g^{-2}$ and $\varrho_2=\varrho/2$, the following hold. Throughout, $m_0\ge0$
is an integer, $\ell_g$ is the reference length, and $\mathbb{T}$ is a torus of side
$2L^{m_0}\ell_g$. In this statement the depth $n$ is prescribed and the bare inverse coupling is
tuned, instead of being fixed by first passage; first passage returns only in (v-c)(2) and in
(v-d)(iii)($\beta'$), which rests on it.
\begin{itemize}
\item[(v-a)] \emph{Exact pin; asymptotic freedom of the tuned bare coupling.} For every depth
$n\ge0$ there is exactly one real number $a_n$ with the following property: the $n$-step recursion
of the running inverse couplings of Definition~\ref{1.2:def-couplings}, with the coefficients of
the infinite lattice and with the coupling-dependent choices fixed in terms of $X$, started at
$x_0=a_n$, is admissible and ends exactly at $x_n=X$. The number $a_n$ depends only on $L$,
$\mathfrak{p}$, $g$ and $n$, and not on the torus. Moreover
\begin{align*}
\lambda_U n-2c_2&\le a_n-X\le 3\lambda_U n+2c_2,\\
\lambda_U(n-l)-c'&\le a_n-a_l\le 3\lambda_U(n-l)+c'\qquad(0\le l\le n),
\end{align*}
where $c_2$ is the flow constant of Definition~\ref{1.2:def-flow-constants} and $c'$ is a constant
independent of $n$, $l$ and the torus.
No monotonicity at a single step is asserted. For a number $\hat w$, the run at depth $n$ from
$a_n-\hat w$ is called the \emph{tuned run with offset $\hat w$}, and for $\hat w=0$ simply the
\emph{tuned run}.

\item[(v-b)] \emph{The one-parameter limit family.} Fix $\mathbb{T}$, an integer
$\mathsf{p}\ge2$, and Lipschitz test functions $\vec f=(f_1,\dots,f_{\mathsf{p}})$ on $\mathbb{T}$
whose supports are pairwise at distance at least $\mathsf{d}>0$. For real or complex $x$ write
$S^n(x;\vec f)$ for the connected $\mathsf{p}$-point function of the smeared curvature observables
of Definition~\ref{1.3:def-curvature-field} on $\mathbb{T}$, computed with $n$ renormalisation
steps between the bare lattice and the reference scale from the bare inverse coupling $x$. The
limits
\begin{equation*}
S^{(\hat w)}(\vec f):=\lim_{n\to\infty}S^n(a_n-\hat w;\vec f)
\end{equation*}
of Theorem~\ref{7.3:thm-tuned-family} exist for $|\operatorname{Re}\hat w|\le\varrho_2/2$,
$|\operatorname{Im}\hat w|<\upsilon_w$, where $\upsilon_w>0$ depends on the torus. The map
$\hat w\mapsto S^{(\hat w)}(\vec f)$ is real-analytic on $(-\varrho_2/2,\varrho_2/2)$ and real for
real $\hat w$; it is holomorphic on the strip $|\operatorname{Re}\hat w|<\varrho_2/2$,
$|\operatorname{Im}\hat w|<\upsilon_w$; nothing is
asserted for larger imaginary parts.

Write $D(0,\varrho)$ for the open disc of radius $\varrho$ about $0$ in $\mathbb{C}$. For each $n$
there is a
map $\mathfrak{n}_n$, holomorphic on $D(0,\varrho)$ and real for real $\hat w$, with
$\mathfrak{n}_n(0)=0$, such that the run at
depth $n+1$ from $a_{n+1}-\mathfrak{n}_n(\hat w)$ and the run at depth $n$ from $a_n-\hat w$ have
equal coupling at the reference scale; it satisfies
\begin{equation*}
|\mathfrak{n}_n(\hat w)-\hat w|\le\Upsilon_n|\hat w|,\qquad
\Upsilon_n\le C\,n^{-3/2}(\log n)^{p_0+1}\quad(n\ge2),
\end{equation*}
for a constant $C$ independent of $n$ and of the torus,
so that $\sum_n\Upsilon_n<\infty$. Let $N_n$ be the derivative with respect to the bare shift, at
shift $0$, of the running shift at the reference scale of clause (iii-c) of
Theorem~\ref{3.1:thm-main}, taken along the tuned run at depth $n$. Then $N_n\in[\frac23,\frac43]$,
$N_{n+1}\mathfrak{n}_n'(0)=N_n$, and the
sequence $(N_n)$ converges. The quantities $a_n$, $\mathfrak{n}_n$, $N_n$, $\Upsilon_n$ and every
constant in (v-a) and (v-b) are independent of the torus; the torus enters only through
$\upsilon_w$, the rate constant of Theorem~\ref{7.3:thm-tuned-family}, and a lower bound $n_0$
on $n$.

\item[(v-c)] \emph{Exhaustion.}
\begin{itemize}
\item[(1)] Let $(a'_n)$ be any real sequence with $|a'_n-a_n|\le\varrho_2/2$. Then
\begin{equation*}
S^n(a'_n;\vec f)-S^{(a_n-a'_n)}(\vec f)\to0\qquad(n\to\infty).
\end{equation*}
Hence two admissible bare families whose offsets $a_n-a'_n$ converge to the same $\hat w$ have
equal limits $S^{(\hat w)}$, and an admissible family converges, without passage to subsequences,
if and only if $S^{(a_n-a'_n)}$ converges; in particular it converges whenever its offset
converges.
\item[(2)] Suppose moreover
\begin{equation*}
\varrho\ge6(B_\sharp+1),
\end{equation*}
and that $g$ lies below a further threshold not exceeding $\gamma_U$. Then every bare inverse
coupling $x$ whose first-passage run stops at depth $K$ satisfies $|x-a_K|\le\varrho_2/2$, and every
subsequential limit of clause (ii-$\mathbb{T}$-b) of Theorem~\ref{3.1:thm-main} is $S^{(\hat w)}$
for some $\hat w$ with $|\hat w|\le\varrho_2/2$.
\item[(3)] The set of continuum limits reachable on $\mathbb{T}$ from the tube, that is, from bare
inverse couplings within $\varrho_2/2$ of the tuned ones, is exactly
$\{S^{(\hat w)}:\ |\hat w|\le\varrho_2/2\}$.
\end{itemize}
It is not asserted that $\hat w\mapsto S^{(\hat w)}$ is injective or non-constant; neither this
non-degeneracy property nor its negation is proved. So the statement ``equal limits imply equal
parameters'' is not made. A calibration by a physical quantity is available only conditionally on
this non-degeneracy and is not adopted.

\item[(v-d)] \emph{The marginal coordinate is the renormalised coupling.}
\begin{itemize}
\item[(i)] Re-file each healed large-field remnant $R^{(i)}$ into the regular family at the level
\begin{equation*}
i'':=i+\log_L(M\check R_i),
\end{equation*}
at which its partition cube is a single cell, where $\check R_i$ is the number attached by the
construction to the remnant $R^{(i)}$ such that $M\check R_i$ is the side of that partition cube
measured in lattice units of level $i$; this is done by the re-filing lemmas. Let $\hat x$
be the inverse
coupling, at the reference scale, of the run so re-filed, label the family by $t:=X-\hat x$, and
normalise the $\mathsf{p}$-point functions by $\widehat{\mathsf{N}}_n(t)^{\mathsf{p}}$. Then the
$t$-labelled normalised family converges as $n\to\infty$ at the exponential rate
$n^c\rho^{n/2}$, for some constant exponent $c$ independent of $n$, where $\rho<1$ is the flow
contraction constant, bounded below in terms of the
Stage-0 data of Notation~\ref{2.3:not-stages-first} and the flow constants of
Definition~\ref{1.2:def-flow-constants} only; the differences between the re-filed runs are bounded
at level
$u\le n$ above the reference scale by $C_E\rho^{\,n-u}$, with no additional term decaying
only polynomially in $n$. No exponential
rate is asserted for the family pinned on the first-passage coupling $x_K$ of
Theorem~\ref{3.1:thm-main}, for the family read through the bare operator, or for the family
labelled by a bare shift $\hat w\ne0$: the rates available for these are $n^{-1/2}$ times a
power of $\log n$, respectively rates below every power, and in no case of the form
$L^{-\mathrm{const}\cdot n}$.
The bare offset $\hat w$ and the renormalised coordinate $t$ parametrise the same family; at each
depth the change of coordinates is real-analytic, with derivative in $[\frac23,\frac43]$ near $0$,
in view of clause (iii-c) of Theorem~\ref{3.1:thm-main}.
\item[(ii)] Two admissible families on $\mathbb{T}$ whose inverse couplings at the reference scale
converge to the same value have the same limit. This follows from (v-c)(1) and the following
bi-Lipschitz property of the reference-scale coupling, the bi-Lipschitz lemma: writing
$x_{\mathrm{ref},n}(\hat w)$ for the inverse coupling at the reference scale of the
run at depth $n$ from $a_n-\hat w$, the map $\hat w\mapsto x_{\mathrm{ref},n}(\hat w)$ is
bi-Lipschitz on $|\hat w|\le\varrho_2/2$ uniformly in $n$, with lower constant
\begin{equation*}
\prod_{k<n}\bigl(1-\operatorname{Lip}b_{k+1}\bigr)\ \ge\ 1-\sum_{k<n}\operatorname{Lip}b_{k+1}
\ \ge\ \tfrac13
\end{equation*}
and upper constant at most $\frac53$, both obtained from the bound of the running-shift statement,
for runs of $n$ steps,
\begin{equation*}
\sum_{k<n}\operatorname{Lip}_{\overline{D}(0,\varrho)}b_{k+1}\le\tfrac23,
\end{equation*}
in which $\operatorname{Lip}_{\overline{D}(0,\varrho)}$ is the Lipschitz constant on the closed disc
of radius $\varrho$; and $x_{\mathrm{ref},n}$ converges uniformly as $n\to\infty$, because
$\sum_n\Upsilon_n<\infty$ by (v-b).
\item[(iii)] \emph{Up to scale, exactly.} The data of the construction for a limit are
$(m_0,\hat x)$, equivalently $(m_0,g,\hat w)$; by ($\gamma$) below, $\hat x$ is also the limit as
$n\to\infty$ of $x_{\mathrm{ref},n}(\hat w)$ of (ii). The \emph{canonical scaffold} of reference
$X$ is the sequence $(X+c_0u)_{u\ge0}$, with $c_0=2\lambda_U$ as above, indexed by the level $u$
above the reference scale. For $s\ge1$ put
$X_s:=X+c_0s$ and $g_s:=X_s^{-1/2}$, so that $g_s<g$, and
\begin{equation*}
D_s(g):=2c_2+\sum_{u<s}\iota_u(X),\qquad
j_U:=1+\bigl\lceil(\varrho_2+c')/\lambda_U\bigr\rceil,
\end{equation*}
where $\iota_u(X)$ is the per-level drift bound, at level $u$, of the running inverse couplings on
the canonical scaffold of reference $X$ (recall $X=g^{-2}$). Write $a_n(g)$ for the number $a_n$ of
(v-a) at reference coupling $g$;
write $S^n_{m_0}[x](\vec f)$ for the function $S^n(x;\vec f)$ of (v-b) on the torus of side
$2L^{m_0}$ reference lengths, whose bare lattice has $2L^{m_0+n}$ sites a side, that is,
$S^n_{m_0}[x](\vec f)=S^T_{\mathsf{p};n,m_0}[x](\vec f)$ in the notation of
Definition~\ref{1.3:def-curvature-field}, the bracket recording the bare inverse coupling, and
$S^{(\hat w)}_{m_0}[g](\vec f)$ for its limit of (v-b) at reference coupling $g$.
\begin{itemize}
\item[($\alpha$)] \emph{Finite cutoff.} For $0\le s\le K$, every bare inverse coupling $x_0$ and
every tuple of test functions,
\begin{equation*}
S^T_{\mathsf{p};K,m}[x_0](f_1,\dots,f_{\mathsf{p}})
=S^T_{\mathsf{p};K-s,m+s}[x_0]\bigl(f_1(L^{-s}\,\cdot\,),\dots,f_{\mathsf{p}}(L^{-s}\,\cdot\,)\bigr),
\end{equation*}
in which $K$ is the number of renormalisation steps, $m$ is the torus index of
Definition~\ref{1.1:def-lattice}, and both members are
one lattice integral, the plaquette barycentres being named with the reference length moved $s$
levels toward the cutoff; in the notation above,
$S^n_{m_0}[x](\vec f)=S^{n-s}_{m_0+s}[x]\bigl(\vec f(L^{-s}\,\cdot\,)\bigr)$ for $0\le s\le n$.
Moreover, for $1\le s\le n$, since the canonical scaffold of reference $X$, read from level $s$ on,
is the canonical
scaffold of reference $X_s$, the tuned bare datum carries two labels exactly:
\begin{equation*}
a_n(g)-\hat w=a_{n-s}(g_s)-\bigl(\hat w+\delta_{n,s}(g)\bigr),
\end{equation*}
with $\delta_{n,s}(g)$ real. If moreover $D_s(g)\le\frac23\varrho$, which holds in particular under
the condition \eqref{9.7:eq-uniqueness-1} of ($\beta$), then $|\delta_{n,s}(g)|\le\frac32D_s(g)$,
and $\delta_{n,s}(g)\to\delta_s(g)$ as $n\to\infty$, the convergence being summable in $n$.
\item[($\beta$)] \emph{Continuum limits, prescribed exponent.} Let $s\ge1$ satisfy, together with
$\varrho$ and $g$, the one-sided condition
\begin{equation}\label{9.7:eq-uniqueness-1}
\varrho\ge48c_2,\qquad \sum_{u<s}\iota_u(X)\le\varrho/24 .
\end{equation}
The first member is a floor on $\varrho$; for each $s$ the second is a ceiling on $g$, chosen last;
equivalently, for each $g$ the second restricts $s$ to a range $s\le s_{\mathrm{dil}}(g;\varrho)$,
where $s_{\mathrm{dil}}(g;\varrho)$ increases to infinity as $g$ decreases to zero, so that at fixed
$\varrho$ and $g$ the condition bounds $s$. Then for every real $\hat w$ with
$|\hat w|<\varrho_2/4$ and all real Lipschitz $\vec f$ with pairwise separated supports on the torus
of side $2L^{m_0}$ reference lengths,
\begin{equation*}
S^{(\hat w)}_{m_0}[g](\vec f)
=S^{(\hat w+\delta_s(g))}_{m_0+s}[g_s]\bigl(\vec f(L^{-s}\,\cdot\,)\bigr),
\end{equation*}
both members being full limits, without passage to subsequences, of the families of (v-b) at
$(m_0,g)$ and at $(m_0+s,g_s)$. Thus the continuum limit of the member $(m_0,\hat x)$ is the
$s$-fold change of units of the continuum limit of the member $(m_0+s,\hat x_s)$, where $\hat x_s$
is the limit, which exists by ($\gamma$), of the inverse couplings $s$ levels above the reference
scale of the first member's own tuned runs, and
\begin{equation*}
\hat x_s-\hat x\in[\lambda_\sharp s-c_5,\ B_\sharp s].
\end{equation*}
The identifications compose in $s$.
\item[($\beta'$)] \emph{Continuum limits at every reference coupling; the exponent confined to a
window.} Write
$\gamma_{U'}\le\gamma_U$ for the further threshold of (v-c)(2), and assume the floor
$\varrho\ge6(B_\sharp+1)$ of (v-c)(2) and $g\le\gamma_{U'}$. For every $X'$ whose coupling
$g':=X'^{-1/2}$ satisfies $g'\le\gamma_{U'}$, with $X'$ on either side of $X$, and every real
$\hat w$ with $|\hat w|\le\varrho_2/2$, there are an integer $\sigma$ and a real $\hat w'$ with
$|\hat w'|\le\varrho_2/2$, depending only on $L$, $\mathfrak{p}$, $g$, $g'$ and $\hat w$, and
neither on the torus nor on the test functions, such that $\sigma=j+s'$ with integers $j$, $s'$
satisfying
\begin{equation*}
|j|\le j_U-1,\qquad s'(X'-X)\ge0,\qquad
\frac{|X'-X|-B_\sharp}{B_\sharp}<|s'|<\frac{|X'-X|+B_\sharp+c_5}{\lambda_\sharp},
\end{equation*}
and such that for every $m_0$ above a floor depending only on $g$ and $g'$ and all real Lipschitz
$\vec f$ with pairwise separated supports,
\begin{equation*}
S^{(\hat w)}_{m_0}[g](\vec f)
=S^{(\hat w')}_{m_0+\sigma}[g']\bigl(\vec f(L^{-\sigma}\,\cdot\,)\bigr).
\end{equation*}
Thus the member $(m_0;X,\hat w)$ is the exact $L^\sigma$-dilate of the member
$(m_0+\sigma;X',\hat w')$, the full limit of its own tuned family; the common bare data have one
first-passage run, which first passes at $X$ and, $|s'|$ levels toward the cutoff, at $X'$.
Moreover, for every $s\ge s_\flat(g)$, where $s_\flat(g)$ is a floor depending only on $g$, there is
such an $X'>X$ with $\sigma\in[s,\,s+\lfloor c_5/\lambda_\sharp\rfloor]$, and $\sigma=s$ when
$c_5<\lambda_\sharp$.
\item[($\gamma$)] \emph{The intermediate couplings.} For every fixed $s\ge0$ the inverse coupling
$s$ levels above the reference scale of the tuned runs at depth $n$ from $a_n(g)-\hat w$ converges
as $n\to\infty$, uniformly in $|\hat w|\le\varrho_2/2$, to a real-analytic function
$\hat x_s(\hat w)$; $\hat x_0=\hat x$, and $\hat x_s-\hat x\in[\lambda_\sharp s-c_5,\ B_\sharp s]$.
In particular the corollary to clause (0) of Theorem~\ref{3.1:thm-main} holds along the tuned
family without passage to subsequences.
\item[($\delta$)] Nothing relates distinct values of $\hat x$ on one torus and at one reference
length. It is not asserted that the tuned bare data $a_n(g)$ and $a_{n-s}(g')$ coincide for any
$g'$: they are pinned at the reference scale on two different scaffolds, and what is exact is the
two-label identity of ($\alpha$). Nor are the following asserted: a prescribed exponent across an
arbitrary reference coupling, beyond the window $[s,\,s+\lfloor c_5/\lambda_\sharp\rfloor]$ of
($\beta'$); the identity of ($\beta$) outside the window \eqref{9.7:eq-uniqueness-1} (outside it,
the identity follows from the present clause read at reference coupling $g_s$ with its own radius
\begin{equation*}
\varrho_s:=48c_2+24\sum_{u<s}\iota_u(X)+8|\hat w|
\end{equation*}
whenever $g_s\le\gamma_U(\varrho_s)$, where $\gamma_U(\cdot)$ is the threshold $\gamma_U$ as a
function of the radius; this reading is not adopted); convergence, rather than clustering, of the
parameter $\hat w'$ of ($\beta'$), which would follow only from the non-degeneracy property of
(v-c); the chart maps between the two labellings, $(X,\hat w)\mapsto(X',\hat w')$ and
$\hat x_s\mapsto\hat x'$, $\hat x'$ being the reference-scale label of the member at $g'$, which
none of ($\alpha$)--($\gamma$) uses; a continuous dilation; dimensional transmutation; injectivity
of $\hat w\mapsto S^{(\hat w)}$.
\end{itemize}
\end{itemize}

\item[(v-e)] \emph{The same-cutoff Lipschitz step.} Fix one lattice and one scaffold, and take two
input lists in the class of the correlation theorem, complex sources allowed. The outputs of
Ba{\l}aban's complete renormalisation step, the large-field operations $\mathbf{R}$, $\mathbf{B}$,
$\mathbf{T}$ included, differ by at most cutoff-free multiples of the difference of the inputs,
measured in the weighted norm of the correlation theorem: the energy and small-field components
with a weight proportional to $\theta(g_k)$, the pure component at second order with the weight
$\theta_k^2\mathfrak{m}_k$, and the boundary components and those of the operations
$\mathbf{T}'_k$ by sums weighted with $e^{-p_0(g_k)}$, $p_0(\cdot)$ being the degree function of
the coupling, where $\theta(\cdot)$, $\theta_k$ and
$\mathfrak{m}_k$ are the weights fixed with the correlation theorem. No constant depends on $K$,
$k$, the torus or position. Old large-field regions are forgotten geometrically from their first
step: by the lemma on geometric forgetting of old regions they lose a fixed factor less than one
per step of their
life, paid from Ba{\l}aban's entropy constant under the lower bound $p_0\ge6r+1$, $r$ being
Ba{\l}aban's parameter in the conditions on the degree exponent $p_0$. This is the step that
Theorem~\ref{7.3:thm-tuned-family} and (v-b) iterate.
\end{itemize}
\end{theorem}

In words, item (v-d)(iii) of Theorem~\ref{9.7:thm-uniqueness} says the following. At every finite
cutoff the $L$-adic change of the reference length is an exact identity of one lattice integral
under two namings; in the continuum limit the one-parameter family is carried into itself exactly
by that change: the limit theory with parameter $\hat w$ on the torus of $2L^{m_0}$ reference
lengths at reference coupling $g$ is the $s$-fold change of units of the limit theory on the torus
of $2L^{m_0+s}$ reference lengths at the weaker reference coupling $(g^{-2}+c_0s)^{-1/2}$, with the
definite parameter $\hat w+\delta_s(g)$ --- the member whose reference-scale coupling is the limit
$\hat x_s$ of the first member's own running couplings $s$ levels up, $\hat x_s-\hat x$ between
$\lambda_\sharp s-c_5$ and $B_\sharp s$ --- for every $s$ inside a displayed window (a floor on the
window radius; for each $s$ a ceiling on $g$), and, through Ba{\l}aban's first-passage
bookkeeping, at every weaker reference coupling with the exponent confined to a two-sided window,
within $\lfloor c_5/\lambda_\sharp\rfloor$ of any prescribed large $s$;
the running couplings $s$ levels up converge along the tuned family; coincidence of the two
families' bare data, prescribed exponents across arbitrary references, the charts between the
labellings, continuous dilations, dimensional transmutation and injectivity are not asserted.

Clause (v) is proved in the remaining sections of this chapter as an implication: items (v-a),
(v-b), (v-c) with its corollary on first-passage runs, and (v-d) from the uniqueness theorem and
its corollary together with the further statements named there, and item (v-e) along the two
routes described next; these statements rest in turn on statements named there, some of which are
stated with their proof incomplete at a marked step, and the standing of each is printed with it.
The Lipschitz step (v-e) is obtained there along two routes: one through the
one-lattice Lipschitz step together with the carry lemma, and one through a second form of that
step on one lattice together with its own auxiliary estimate.

\begin{definition}[The two runs at the last lattice and their source polydiscs]
\label{9.8:def-source-polydiscs}
The window $\mathcal{N}'_w$ of complex centres and its radii $\upsilon_w$, $R_0$, $R_1$ are those
fixed with the last lattice $\mathbb{T}_w$; the convention \eqref{9.1:eq-run-pair-transplant-b} of
Definition~\ref{9.1:def-run-pair-transplant} is in force. Let
$\mathsf{p}\ge2$ be the number of test functions $f_1,\dots,f_{\mathsf{p}}$, with the norm
$\|\cdot\|_\sharp$ of Definition~\ref{1.3:def-curvature-field}. Let $n$ be a length with
$\Upsilon_n\le\frac14$, where $\Upsilon_n$ and the matching map $\mathfrak{n}_n$ are those of
Proposition~\ref{9.6:prop-depth-zero-match}. Let $\hat w\in\mathcal{N}'_w$, the window of complex
centres, contained in the disc $|\hat w|<\upsilon_w$; $R_0$ and $R_1$ are the two source radii
attached to it. With $r_J$ the radius of hypothesis (J3) and $r_s^{pr}$ the source radius of
hypothesis (J4) of Definition~\ref{9.1:def-tuned-family}, and $\varrho_2=\varrho/2$, we assume
throughout this definition that the window and the radii satisfy
\begin{equation}\label{9.8:eq-source-polydiscs-5}
\begin{gathered}
\upsilon_w<R_0/2,\qquad R_0<r_s^{pr},\qquad \upsilon_w+R_0/2\le19r_J/64,\\
(3/2)R_0\le r_J,\qquad \upsilon_w\le\tfrac12\varrho_2 .
\end{gathered}
\end{equation}
Put
$u:=\operatorname{Re}\hat w$.
The two runs are the following.
\begin{itemize}
\item[$\mathsf{A}$:] the run of length $n$ with real centre $\zeta_*^n+u$ and continuum source
\begin{equation}\label{9.8:eq-source-polydiscs-1}
W^{\mathsf{A}}(z):=i\operatorname{Im}\hat w+\sum_{j=1}^{\mathsf{p}}z_jf_j .
\end{equation}
\item[$\mathsf{B}$:] the run of length $n+1$ with centre $\zeta_*^{n+1}+\mathfrak{n}_n(u)$ and source
\begin{equation}\label{9.8:eq-source-polydiscs-2}
W^{\mathsf{B}}(z):=\mathfrak{n}_n\circ\bigl(u+W^{\mathsf{A}}(z)\bigr)-\mathfrak{n}_n(u),
\end{equation}
the composition being taken pointwise.
\end{itemize}
Here $\zeta_*^n$ is the real alive run of length $n$ of Definition~\ref{9.1:def-tuned-family}. For
$\xi\in\mathbb{C}$ and coefficients $z_j$ we put
\begin{equation*}
\Bigl\|\xi+\sum_{j}z_jf_j\Bigr\|_{\flat}:=|\xi|+\sum_{j}|z_j|\,\|f_j\|_\sharp ,
\end{equation*}
and we define the two
polydiscs
\begin{equation}\label{9.8:eq-source-polydiscs-3}
\begin{aligned}
P_{UV}&:=\Bigl\{z\in\mathbb{C}^{\mathsf{p}}:\ |z_j|\le\frac{R_0}{2\mathsf{p}\|f_j\|_\sharp},\
1\le j\le\mathsf{p}\Bigr\},\\
P_{IR}&:=\Bigl\{z\in\mathbb{C}^{\mathsf{p}}:\ |z_j|\le\frac{R_1}{2\mathsf{p}\|f_j\|_\sharp},\
1\le j\le\mathsf{p}\Bigr\}.
\end{aligned}
\end{equation}
For $z\in P_{UV}$ the following hold.
\begin{itemize}
\item[(a)] One has
\begin{equation*}
\|W^{\mathsf{A}}(z)\|_\flat<\upsilon_w+R_0/2<R_0<r_s^{pr}.
\end{equation*}
\item[(b)] The values of $\omega:=u+W^{\mathsf{A}}(z)$ lie in
$\bar D(0,19r_J/64)\subset\bar D(0,r_J/2)$, and on $\bar D(0,r_J/2)$ one has
\begin{equation*}
|\mathfrak{n}_n'-1|\le2\Upsilon_n\le\tfrac12 .
\end{equation*}
\item[(c)] $W^{\mathsf{B}}(z)\in\mathcal{W}[(3/2)R_0]\subset\mathcal{W}[r_J]$, where
$\mathcal{W}[r]$ is the class of complex sources of radius $r$ of hypothesis (J4) of
Definition~\ref{9.1:def-tuned-family}, membership in which is the pointwise bound
$|W(x)|<r$ at every point $x$, and
$|\mathfrak{n}_n(u)|\le\frac58\varrho_2$.
\item[(d)] Both runs belong to the class of runs admitted by hypothesis (J4) of
Definition~\ref{9.1:def-tuned-family}, and the pair $(\mathsf{A},\mathsf{B})$ to the class of pairs
of Definition~\ref{9.1:def-run-pair-transplant}.
\end{itemize}
For $\mathsf{c}\in\{\mathsf{A},\mathsf{B}\}$, with $\tilde Z^{\mathsf{c}}$ the normalised sourced
integral of the run $\mathsf{c}$, we put
\begin{equation}\label{9.8:eq-source-polydiscs-4}
\mathsf{h}^{\mathsf{c}}:=\frac{\tilde Z^{\mathsf{c}}(W^{\mathsf{c}})}{\tilde Z^{\mathsf{c}}(0)} .
\end{equation}
\end{definition}

\begin{proof}
(a) The source $W^{\mathsf{A}}(z)$ of \eqref{9.8:eq-source-polydiscs-1} has constant part
$\xi=i\operatorname{Im}\hat w$ and coefficients $z_j$, so by the definition of $\|\cdot\|_\flat$
\begin{equation*}
\|W^{\mathsf{A}}(z)\|_\flat=|\operatorname{Im}\hat w|+\sum_{j=1}^{\mathsf{p}}|z_j|\,\|f_j\|_\sharp .
\end{equation*}
For $z\in P_{UV}$ each summand is at most $R_0/(2\mathsf{p})$, so the sum is at most
$\mathsf{p}\cdot R_0/(2\mathsf{p})=R_0/2$. Since $|\operatorname{Im}\hat w|\le|\hat w|<\upsilon_w$ for
$\hat w\in\mathcal{N}'_w$, the first inequality of (a) follows. The remaining two inequalities of (a)
are the first two relations in \eqref{9.8:eq-source-polydiscs-5}.

(b) At every point $x$ one has $|f_j(x)|\le\|f_j\|_\infty\le\|f_j\|_\sharp$, and
\begin{equation*}
\omega(x)=u+i\operatorname{Im}\hat w+\sum_{j=1}^{\mathsf{p}}z_jf_j(x)
=\hat w+\sum_{j=1}^{\mathsf{p}}z_jf_j(x).
\end{equation*}
Hence, with $|\hat w|<\upsilon_w$ and the bound $\sum_j|z_j|\,\|f_j\|_\sharp\le R_0/2$ from the proof
of (a),
\begin{equation*}
|\omega(x)|\le|\hat w|+\sum_{j=1}^{\mathsf{p}}|z_j|\,\|f_j\|_\sharp<\upsilon_w+R_0/2\le19r_J/64
\end{equation*}
by the third relation in \eqref{9.8:eq-source-polydiscs-5}, so the values of $\omega$ lie in
$\bar D(0,19r_J/64)$. Further $19r_J/64<32r_J/64=r_J/2$ gives the inclusion of discs. By
Proposition~\ref{9.6:prop-depth-zero-match}, the map $\mathfrak{n}_n$ is holomorphic on the open
disc $D(0,r_J)$ and
\begin{equation*}
|\mathfrak{n}_n(v)-v|\le\Upsilon_n|v|\le\Upsilon_nr_J\qquad(|v|<r_J).
\end{equation*}
Thus the function $v\mapsto\mathfrak{n}_n(v)-v$ is holomorphic on $D(0,r_J)$ and bounded there in
modulus by $\Upsilon_nr_J$. Let $v_0$ be a point with $|v_0|\le r_J/2$. The open disc
$D(v_0,r_J/2)$ lies in $D(0,r_J)$, and Cauchy's inequality for the derivative at the centre of
this disc gives
\begin{equation*}
|\mathfrak{n}_n'(v_0)-1|\le\frac{\Upsilon_nr_J}{r_J/2}=2\Upsilon_n ,
\end{equation*}
and $2\Upsilon_n\le\frac12$ because $\Upsilon_n\le\frac14$.

(c) By (b), $|\mathfrak{n}_n'|\le1+\frac12=\frac32$ on $\bar D(0,r_J/2)$. The point $z=0$ lies in
$P_{UV}$, and there $\omega=u+i\operatorname{Im}\hat w=\hat w$, so by (b)
$|u|\le|\hat w|\le19r_J/64$; thus $u\in\bar D(0,r_J/2)$. At every point $x$ both $u$ and
$\omega(x)$ lie in the convex set $\bar D(0,r_J/2)\subset D(0,r_J)$, hence so does the segment
$u+\mathsf{s}\,W^{\mathsf{A}}(z)(x)$, $0\le\mathsf{s}\le1$, and \eqref{9.8:eq-source-polydiscs-2}
gives
\begin{equation*}
W^{\mathsf{B}}(z)(x)=\int_0^1\mathfrak{n}_n'\bigl(u+\mathsf{s}\,W^{\mathsf{A}}(z)(x)\bigr)\,
d\mathsf{s}\;W^{\mathsf{A}}(z)(x).
\end{equation*}
Since $|f_j(x)|\le\|f_j\|_\sharp$, one has $|W^{\mathsf{A}}(z)(x)|\le\|W^{\mathsf{A}}(z)\|_\flat<R_0$
by (a); with $|\mathfrak{n}_n'|\le\frac32$ on the segment this yields
\begin{equation*}
|W^{\mathsf{B}}(z)(x)|\le\tfrac32|W^{\mathsf{A}}(z)(x)|\le\tfrac32\|W^{\mathsf{A}}(z)\|_\flat
<\tfrac32R_0 ,
\end{equation*}
so $W^{\mathsf{B}}(z)\in\mathcal{W}[(3/2)R_0]$. The inequality $(3/2)R_0\le r_J$, which gives the
inclusion $\mathcal{W}[(3/2)R_0]\subset\mathcal{W}[r_J]$, is the fourth relation in
\eqref{9.8:eq-source-polydiscs-5}. Finally, since $|u|\le r_J/2<r_J$,
Proposition~\ref{9.6:prop-depth-zero-match} gives $|\mathfrak{n}_n(u)-u|\le\Upsilon_n|u|$, and with
$\Upsilon_n\le\frac14$, $|u|<\upsilon_w$ and the last relation in
\eqref{9.8:eq-source-polydiscs-5} we obtain
\begin{equation*}
|\mathfrak{n}_n(u)|\le(1+\Upsilon_n)|u|\le\tfrac54|u|<\tfrac54\upsilon_w\le\tfrac58\varrho_2 ,
\end{equation*}
which is the bound on the centre shift of $\mathsf{B}$ in (c).

(d) By (a) the source of $\mathsf{A}$ has $\|\cdot\|_\flat$-size below $r_s^{pr}$, by (c) the source
of $\mathsf{B}$ lies in $\mathcal{W}[r_J]$ and its centre shift is at most $\frac58\varrho_2$; these
place both runs in the class of hypothesis (J4) of Definition~\ref{9.1:def-tuned-family}. Adding
centre and source, the run $\mathsf{A}$ is read at $\zeta_*^n+v$ with $v:=\omega=u+W^{\mathsf{A}}(z)$,
whose values lie in $D(0,r_J)$ by (b), and by \eqref{9.8:eq-source-polydiscs-2} the run
$\mathsf{B}$ is read at
\begin{equation*}
\zeta_*^{n+1}+\mathfrak{n}_n(u)+W^{\mathsf{B}}(z)=\zeta_*^{n+1}+\mathfrak{n}_n(v).
\end{equation*}
Thus $\mathsf{A}$ is read at $v$ and $\mathsf{B}$ at $\mathfrak{n}_n(v)$, which is the
parametrisation of pairs through $\mathfrak{n}_n$ in Definition~\ref{9.1:def-run-pair-transplant}.
\end{proof}

\begin{lemma}[Localisation of the difference of boundary exponents]
\label{9.8:lem-exponent-localisation}
Let the runs $\mathsf{A}$ (of length $n$) and $\mathsf{B}$ (of length $n+1$), their sources
$W^{\mathsf{A}}(z)$, $W^{\mathsf{B}}(z)$, the centre $u=\operatorname{Re}\hat{w}$ and the
polydisc $P_{UV}$ be as in Definition~\ref{9.8:def-source-polydiscs}, and put
$\omega:=u+W^{\mathsf{A}}$. For $\mathsf{c}\in\{\mathsf{A},\mathsf{B}\}$ write
$e^{(r),\mathsf{c}}(X;w)$ for the boundary exponent of the run $\mathsf{c}$ at depth $r$ on
the localisation domain $X$ at the source $w$, that is, the summand on $X$ contributed by the
depth-$r$ step to the sum $\mathcal{E}$ of boundary exponents of the run $\mathsf{c}$ in the
factorisation $Z=e^{\mathcal{E}}\tilde{Z}$ of its sourced integral.
Assume:
\begin{itemize}
\item[(i)] the standing hypotheses \textup{(J1)}--\textup{(J4)} of \eqref{9.1:eq-tuned-family-a};
\item[(ii)] the two hypotheses of the two-run estimate, Proposition~\ref{9.6:prop-two-run-estimate},
other than \textup{(J1)}, which is contained in (i), and other than the bounds on the two-run
comparison families, which form item (iii);
\item[(iii)] the bounds on the two-run comparison families for pairs of local lists, in all
components of the production blocks, the component $R'$ with its bound
\eqref{9.1:eq-comparison-families-a} included.
\end{itemize}
Then for every depth $r$ with $\mathsf{w}\le r<n$ there are a constant $c_r$ independent of
$z$, and functions $G_r(Y;z)$, with $Y$ ranging over all localisation domains of the torus at
depth $r$ (domains that wrap around the torus included), each holomorphic in a neighbourhood
of $P_{UV}$ and depending on $z$ only through the restriction
$W^{\mathsf{A}}{\upharpoonright}Y^{+}$, such that
\begin{equation}\label{9.8:eq-exponent-localisation-1}
\sum_{X}e^{(r),\mathsf{B}}\bigl(X;W^{\mathsf{B}}(z)\bigr)
-\sum_{X}e^{(r),\mathsf{A}}\bigl(X;W^{\mathsf{A}}(z)\bigr)
=\sum_{Y}G_r(Y;z)+c_r
\end{equation}
and
\begin{equation}\label{9.8:eq-exponent-localisation-2}
\sup_{z\in P_{UV}}|G_r(Y;z)|\le C_\eta\,(1+r)\,e_r\,e^{-(1-\delta)\kappa^{c}d(Y)},
\end{equation}
where
\begin{equation}\label{9.8:eq-exponent-localisation-3}
\begin{gathered}
C_\eta:=M^4C_dC_\delta C_E\Bigl[8E_0\mathsf{l}_*+\frac{2}{C_t\,g^{-2}}\Bigr]+2M^4C_{PM},\\
\mathsf{l}_*:=\sup_{r\ge0}\frac{\mathsf{l}_r}{1+r},\qquad
C_\delta:=\sup_{x\ge0}(1+x)e^{-\delta\kappa^{c}x}.
\end{gathered}
\end{equation}
Here $d(Y)$ is the tree length of $Y$; $\kappa^{c}$ is the fixed decay rate, and $\delta$
the fixed margin, of the classes of list bounds entering Lemma~\ref{9.6:lem-chain-structure}
(the superscript in $\kappa^{c}$ is part of the symbol and is not a component index);
$C_d$ is the constant in the bound $\#\{y\in Y\}\le M^4C_d(d(Y)+1)$ on the number of points of
a localisation domain; and $E_0\mathsf{l}_r$ (for the components $P^{\flat}$ and $D$) and
$(C_t\theta_r)^{-1}$ (for the components $Q$ and $R$) are the weights with which the components
of the comparison families enter the boundary exponents. The constant $C_\eta$ does not depend
on $n$, on $r$ or on the torus; it depends on $g$.
\end{lemma}

\begin{proof}
\emph{Matching at depth zero.} The pair of local lists is matched at depth zero for every frozen
value of the source. Indeed, let $\omega_0$ be a frozen value of $\omega=u+W^{\mathsf{A}}$; by
Definition~\ref{9.8:def-source-polydiscs} it lies in $\bar D(0,r_J/2)\subset D(0,r_J)$, and by
the definition of $W^{\mathsf{B}}$ the corresponding value of $\mathfrak{n}_n(u)+W^{\mathsf{B}}$
is $\mathfrak{n}_n(\omega_0)$. The depth-zero matching proposition
(Proposition~\ref{9.6:prop-depth-zero-match}) gives
$h_{n+1}(\mathfrak{n}_n(\omega_0))=h_n(\omega_0)$, where, as in that proposition,
\[
h_n(v)=\zeta^n_{(0)}(\zeta_*^n+v),\qquad h_{n+1}(v)=\zeta^{n+1}_{(0)}(\zeta_*^{n+1}+v),
\]
so the depth-zero running shifts of the two runs
coincide: $\zeta^{\mathsf{B}}_{(0)}=\zeta^{\mathsf{A}}_{(0)}$ on $\mathfrak{P}$. This is the
hypothesis of the two-run estimate (Proposition~\ref{9.6:prop-two-run-estimate}) with $r_0=0$.
Together with hypotheses (i)--(iii), that estimate therefore holds for the comparison chain
(Definition~\ref{9.6:def-comparison-chain}) of the local pair, for $z\in\mathfrak{P}$ and
$0\le r\le n$:
\begin{equation}\label{9.8:eq-exponent-localisation-4}
g_r^{P}\le C_Ee_r,\qquad g_r^{A}\le 4C_Ee_r,\qquad g_r^{R}\le 4C_E\hat{\mathfrak{r}}_re_r .
\end{equation}
Of the bound on the component $R'$ assumed in (iii), only its part at the unit configuration
$V=1$ is used below: it contains the bound on the component $R$, which enters the term with
$G_r^{R}(Y,1,0;z)$.

\emph{Representation of the boundary exponents.} By the torus bound of the correlation theorem,
\[
e^{(r)}(X;w)=\bar e^{(r)}(X;w)-\bar e^{(r)}(X;0)+(\text{one-point pieces}),
\]
where
\[
\bar e^{(r)}(X;w):=\sum_{y\in X}E^{(r)}(X,1,y;w)+R^{(r)}(X,1;w),
\]
$E^{(r)}(X,1,y;w)$ being the term of this representation attached to the point $y\in X$,
evaluated at the point $p_y$ attached to $y$, and $R^{(r)}(X,1;w)$ its remainder on $X$; the
argument $1$ is the unit configuration $V=1$ at which the boundary exponents are read.
The one-point pieces are carried by single $M$-cubes and enter through the weights
$\log\mathfrak{z}\bigl(\theta_r-\zeta_{(r)}(p_y)\bigr)$. The piece $E$ splits into three
components,
\[
E(\cdot\,;W)=P^{\flat}_{W(p_y)}+Q_{W(p_y)}+D(\cdot\,;W),
\]
the first two of which read the source only through its value $W(p_y)$ at the point $p_y$.
With the weights of the statement, $P^{\flat}$ and $D$ enter the boundary exponent as
$E_0\mathsf{l}_r$ times, and $Q$ and $R^{(r)}$ as $(C_t\theta_r)^{-1}$ times, the corresponding
members of the depth-$r$ families.
These formulas hold for each of the two runs.

\emph{The identity.} In a pair of local lists the relation $\approx$ in (c1) of
\eqref{9.6:eq-chain-structure-a} is equality up to an additive constant independent of $z$, in
every member. The unit configuration $V=1$ is a regular datum: the level-$j$ fields
$U_j^{\mathsf{m}}$ built from it take the value $U_j^{\mathsf{m}}(1)=1$, and the current
vanishes, $\mathbf{J}=0$. Hence every member may be evaluated at $(V,\mathbf{J})=(1,0)$.
Statement (c1) of the chain-structure lemma
(Lemma~\ref{9.6:lem-chain-structure}, \eqref{9.6:eq-chain-structure-a}) gives
$G_r^{c}\approx\varphi_r^{\mathsf{B},c}-\varphi_r^{\mathsf{A},c}$ for each component $c$, where
$\varphi_r^{\mathsf{A},c}$ and $\varphi_r^{\mathsf{B},c}$ are the depth-$r$ families of the two
runs on the component $c$. Applied
to the pieces $E$ for each $y$, and to the pieces $R^{(r)}(X,1;w)$, which carry no point $y$, it
shows that the sums over $X$ of the values of the two runs at $(1,0)$ differ by the sum over $Y$
of the values of the chain at $(1,0)$. We put, with the components $G_r^{P^\flat}$ and
$G_r^{Q}$, which read the source through a frozen value $v$, taken at $v=\omega(p_y)$,
\begin{equation}\label{9.8:eq-exponent-localisation-5}
\begin{split}
G_r(Y;z):={}&\sum_{y\in Y}\Bigl[E_0\mathsf{l}_r\bigl(G_r^{P^\flat}\big|_{v=\omega(p_y)}
+G_r^{D}\bigr)+(C_t\theta_r)^{-1}G_r^{Q}\big|_{v=\omega(p_y)}\Bigr](Y,1,0,y;z)\\
&+(C_t\theta_r)^{-1}G_r^{R}(Y,1,0;z).
\end{split}
\end{equation}
When $Y$ is a single $M$-cube, we add to this
\begin{equation}\label{9.8:eq-exponent-localisation-6}
\sum_{y\in Y}\Bigl[\log\mathfrak{z}\bigl(\theta_r-\zeta^{\mathsf{B}}_{(r)}(p_y)\bigr)
-\log\mathfrak{z}\bigl(\theta_r-\zeta^{\mathsf{A}}_{(r)}(p_y)\bigr)\Bigr],
\end{equation}
which is the difference of the one-point pieces of the two runs on $Y$. We collect in $c_r$ the
terms $\bar e^{(r),\mathsf{c}}(X;0)$ of both runs, summed over $X$, and the additive constants of
the relation $\approx$; none of these depends on $z$, so $c_r$ does not
depend on $z$. With these choices the identity obtained from (c1) is
\eqref{9.8:eq-exponent-localisation-1}.

\emph{Holomorphy and locality.} By (c1), each $G_r^{c}$ is holomorphic on $\mathfrak{P}$, and
its local members on $Y$ read the source only through its restriction to $Y^{+}$. This gives the
claims for \eqref{9.8:eq-exponent-localisation-5}. For the terms
\eqref{9.8:eq-exponent-localisation-6}, the running shift $\zeta_{(r)}(p_y)$ reads the source
only within two depth-$r$ units of $p_y$, and hence only through
$W^{\mathsf{A}}{\upharpoonright}Y^{+}$.

\emph{Size.} The number of points of a localisation domain is bounded by the admissibility of
the prepared large-field data, printed without proof in \cite[pp.~178--179]{B-CMP122a}; cited
as printed. By this bound, and then by the definition of $C_\delta$ in
\eqref{9.8:eq-exponent-localisation-3} with $x=d(Y)$,
\begin{equation}\label{9.8:eq-exponent-localisation-7}
\#\{y\in Y\}\le M^4C_d\bigl(d(Y)+1\bigr)\le M^4C_dC_\delta\,e^{\delta\kappa^{c}d(Y)}.
\end{equation}

Each norm $g_r^{\bullet}$ bounds a member on $Y$ with the factor $e^{-\kappa^{c}d(Y)}$, and
\eqref{9.8:eq-exponent-localisation-4} bounds these norms as follows. The two components
$P^{\flat}$ and $D$ belong to the block $A$ and contribute at most $2\cdot4C_Ee_r$. The component
$Q$ belongs to the block $P$ and contributes at most $C_Ee_r$. Using also
$\theta_r\ge g^{-2}$, the bracket in \eqref{9.8:eq-exponent-localisation-5} is, for each $y$ and
$z\in P_{UV}$, bounded by
\[
\Bigl[8E_0\mathsf{l}_r+\frac{1}{C_t\,g^{-2}}\Bigr]C_Ee_r\,e^{-\kappa^{c}d(Y)}.
\]
Summing over $y$ with \eqref{9.8:eq-exponent-localisation-7}, and using
$\mathsf{l}_r\le\mathsf{l}_*(1+r)$ and $1\le1+r$, bounds the first line of
\eqref{9.8:eq-exponent-localisation-5} by
\[
M^4C_dC_\delta C_E\Bigl[8E_0\mathsf{l}_*+\frac{1}{C_t\,g^{-2}}\Bigr](1+r)\,e_r\,
e^{-(1-\delta)\kappa^{c}d(Y)}.
\]
For the term with $G_r^{R}$ we use $4\hat{\mathfrak{r}}_r\le\frac14$ and
$\theta_r\ge g^{-2}$:
\[
(C_t\theta_r)^{-1}\,|G_r^{R}(Y,1,0;z)|\le\frac{C_E}{4C_t\,g^{-2}}\,e_r\,e^{-\kappa^{c}d(Y)}.
\]
We have $C_\delta\ge1$ (take $x=0$) and $C_d\ge1$ (apply the count
\eqref{9.8:eq-exponent-localisation-7} to a single $M$-cube, which carries $M^4$ points and has
$d(Y)=0$), so $M^4C_dC_\delta\ge1$; moreover $1\le1+r$ and
$e^{-\kappa^{c}d(Y)}\le e^{-(1-\delta)\kappa^{c}d(Y)}$. Hence this term is at most
\[
M^4C_dC_\delta C_E\,\frac{1}{C_t\,g^{-2}}\,(1+r)\,e_r\,e^{-(1-\delta)\kappa^{c}d(Y)},
\]
which is the part of the summand $2/(C_t\,g^{-2})$ of the bracket in $C_\eta$ not used above.

For the terms \eqref{9.8:eq-exponent-localisation-6} we use the bound
$|(\log\mathfrak{z})'|\le2$ on the derivative of the one-point gauge weight on the matched
weights, used on the segment joining $\theta_r-\zeta^{\mathsf{A}}_{(r)}(p_y)$ and
$\theta_r-\zeta^{\mathsf{B}}_{(r)}(p_y)$. We combine it with the pointwise
comparison of the running shifts (Lemma~\ref{9.6:lem-pointwise-shift}), which we apply at depth
$r\ge\mathsf{w}$ at the points $p_y$ of a single $M$-cube carrying a one-point piece; there
$\omega(p_y)\in\bar D(0,r_J/2)$ for $z\in P_{UV}$ by Definition~\ref{9.8:def-source-polydiscs}.
Each summand of
\eqref{9.8:eq-exponent-localisation-6} is then in absolute value at most
\[
2\,\bigl|\zeta^{\mathsf{B}}_{(r)}(p_y)-\zeta^{\mathsf{A}}_{(r)}(p_y)\bigr|\le 2C_{PM}(1+r)e_r .
\]
A single $M$-cube carries $M^4$ points $y$ and has tree length $d(Y)=0$, so that the factor
$e^{-(1-\delta)\kappa^{c}d(Y)}$ equals $1$; these terms therefore contribute at most
$2M^4C_{PM}(1+r)e_r$ to \eqref{9.8:eq-exponent-localisation-2}.

Adding the three contributions gives \eqref{9.8:eq-exponent-localisation-2} with the constant
$C_\eta$ of \eqref{9.8:eq-exponent-localisation-3}. None of the constants entering $C_\eta$
depends on $n$, on $r$ or on the torus; $C_\eta$ depends on $g$.
\end{proof}

\begin{definition}[Closeness of the normalised partition functions at two cutoffs]
\label{9.8:def-partition-closeness}
Let $n$ be as in Definition~\ref{9.8:def-source-polydiscs}. For $\hat{w}\in\mathcal{N}'_w$ let
$u:=\operatorname{Re}\hat{w}$, the two runs
$\mathsf{A}$ (of length $n$) and $\mathsf{B}$ (of length $n+1$), their residual sources
$W^{\mathsf{A}}(z)$, $W^{\mathsf{B}}(z)$ and the polydisc $P_{IR}$ be as in
Definition~\ref{9.8:def-source-polydiscs}. For $\mathsf{c}\in\{\mathsf{A},\mathsf{B}\}$ let
$\tilde{Z}^{\mathsf{c}}$ be the normalised sourced integral (normalised partition function) of the
run $\mathsf{c}$. For each run and
each centre one normalisation is fixed, and it is used for all sources $W^{\mathsf{c}}(z)$,
$z\in P_{IR}$. Let $A$, $B$ be the volume constants of the setting of clause~(v) of
Theorem~\ref{3.1:thm-main}, and
$\mathsf{Q}_w=AB$. The normalisation fixed for each run and each centre is the one in which the
value of $\tilde{Z}^{\mathsf{c}}$ at the zero source is $\tilde{Z}^{\mathsf{c}}(0)=1/B$. For a
depth $r$ let $e_r$ denote the level-wise rate of Definition~\ref{9.1:not-levelwise-rate}. Let
$\mathsf{w}=\mathsf{w}(g)$ be the stop depth. Let $\varepsilon_n\ge0$.

The two runs satisfy \emph{closeness of the normalised partition functions at rate
$\varepsilon_n$} if, for every $\hat{w}\in\mathcal{N}'_w$ and every $z\in P_{IR}$,
\begin{equation}\label{9.8:eq-partition-closeness-a}
\begin{gathered}
\bigl|\tilde{Z}^{\mathsf{B}}\bigl(W^{\mathsf{B}}(z)\bigr)-\tilde{Z}^{\mathsf{A}}\bigl(W^{\mathsf{A}}(z)\bigr)\bigr|
\le\varepsilon_n A ,\\
\varepsilon_n\le C_{TC}(N_w;g)\,e_{\mathsf{w}} ,
\end{gathered}
\end{equation}
with one constant $C_{TC}(N_w;g)$, depending on $N_w$ and $g$, the same for every length $n$ with
$\Upsilon_n\le\tfrac14$.

In particular, the first line of \eqref{9.8:eq-partition-closeness-a} holds at the zero source,
that is at the parameter $\hat{w}=u$ and $z=0$: the real part $u$ of $\hat{w}$ lies in
$\mathcal{N}'_w$ as well, and $0\in P_{IR}$. The two runs then have the same
centres $\zeta_*^n+u$ and $\zeta_*^{n+1}+\mathfrak{n}_n(u)$ as at $\hat{w}$, hence the same
normalisations. Since
$\operatorname{Im}u=0$, Definition~\ref{9.8:def-source-polydiscs} gives
$W^{\mathsf{A}}(0)=0$ and
\begin{equation*}
W^{\mathsf{B}}(0)=\mathfrak{n}_n(u)-\mathfrak{n}_n(u)=0 .
\end{equation*}

Suppose the first line of \eqref{9.8:eq-partition-closeness-a} holds with an $n$-free majorant
$A^{\sharp}\ge A$ in
place of $A$. Then it holds with $A$ and the rate $\varepsilon_nA^{\sharp}/A$, since
\begin{equation*}
\varepsilon_nA^{\sharp}=(\varepsilon_nA^{\sharp}/A)\,A .
\end{equation*}
Since $A^{\sharp}$ does not depend on $n$ and $A$, a volume constant of the setting of
clause~(v), does not depend on $n$, the factor $A^{\sharp}/A$ does not depend on $n$ and is
absorbed into $C_{TC}$.

In this normalisation write
$\mathsf{h}^{\mathsf{c}}:=\tilde{Z}^{\mathsf{c}}/\tilde{Z}^{\mathsf{c}}(0)=B\,\tilde{Z}^{\mathsf{c}}$,
each evaluated at the source $W^{\mathsf{c}}(z)$ of its run. Multiplying the first line of
\eqref{9.8:eq-partition-closeness-a} by $B$ gives
\begin{equation*}
\sup\bigl|\mathsf{h}^{\mathsf{B}}-\mathsf{h}^{\mathsf{A}}\bigr|\le\varepsilon_n\mathsf{Q}_w ,
\end{equation*}
the supremum being taken over the same $\hat{w}$ and $z$.
\end{definition}

\begin{remark}[Content of the condition in charts]
The content of \eqref{9.8:eq-partition-closeness-a} is a bound obtained in charts, as follows. On
each label $h$ common
to the two runs, the same chart is used for both. This chart consists of the variables $V'$ on the
small-field region $\Lambda_i$ of the label, $i$ being the index of the label $h$, and the raw bond
variables elsewhere. The configuration
is written $V=V'V_\Lambda$, where $V_\Lambda$ is written in the gauge fixed by the label's
gauge-fixing condition $G_i$ on $\Lambda_i$ (the condition whose $\delta$-function $\delta_{G_i}$
enters below), so that $V_\Lambda$ depends on
the index $i$. In these coordinates the test functions, the gauge-fixing factor $\delta_{G_i}$ and
the Haar measure are common to $\mathsf{A}$ and $\mathsf{B}$. The left side of the first line of
\eqref{9.8:eq-partition-closeness-a} is then bounded by
\begin{equation*}
\sum_h\int\bigl|F_h^{\mathsf{B}}-F_h^{\mathsf{A}}\bigr|\,d\mu_h ,
\end{equation*}
plus the masses of the unpaired terms. Here $F_h^{\mathsf{c}}$ is the integrand of the $h$-th term
of the run $\mathsf{c}$ in the common chart and $\mu_h$ is the common reference measure there. A
label enters the pairing as common when its first index is at least one, and as aligned, in the
aligned label set $H_{al}$, when its first index exceeds the alignment threshold $n_c$; the terms of
either run whose labels are not paired in this way contribute their masses without pairing. No
bound on the total-variation distance between the measures of the runs $\mathsf{A}$ and
$\mathsf{B}$ on the configuration space $\mathcal{V}_w$ of the last lattice is part of the
condition.
\end{remark}

\begin{lemma}[From partition functions to mixed source derivatives]\label{9.8:lem-derivative-closeness}
Let $n$, $\hat{w}\in\mathcal{N}'_w$, $u:=\operatorname{Re}\hat{w}$, the two runs $\mathsf{A}$, $\mathsf{B}$, their
sources $W^{\mathsf{A}}(z)$, $W^{\mathsf{B}}(z)$, the norm $\|\cdot\|_{\flat}$ and the polydiscs
$P_{UV}\supset P_{IR}$ be as in Definition~\ref{9.8:def-source-polydiscs}. For a source
$v=c_0+\sum_{i=1}^{\mathsf{p}}z_if_i$ with $c_0\in\mathbb{C}$ the constant part of the source (not the
flow constant), let $W^{\mathsf{A}}$, $W^{\mathsf{B}}$ be the sources of
Definition~\ref{9.8:def-source-polydiscs} with $i\operatorname{Im}\hat{w}$ replaced by $c_0$, and put
\begin{equation*}
P:=\tilde{Z}^{\mathsf{B}}(W^{\mathsf{B}}),\qquad Q:=\tilde{Z}^{\mathsf{A}}(W^{\mathsf{A}}),
\end{equation*}
with the normalisations of \eqref{9.8:eq-partition-closeness-a}; for $c_0=i\operatorname{Im}\hat{w}$ these are
the two numbers compared in \eqref{9.8:eq-partition-closeness-a}.
Let $P_0$, $Q_0$ be their values at the zero source ($\hat{w}=u$, $z=0$), and let
$\mathsf{h}^{\mathsf{B}}:=P/P_0$ and $\mathsf{h}^{\mathsf{A}}:=Q/Q_0$, regarded as functions of $(c_0,z)$. Assume:
\begin{itemize}
\item[(a)] for $\mathsf{c}\in\{\mathsf{A},\mathsf{B}\}$ the function $\mathsf{h}^{\mathsf{c}}$ is holomorphic in
$(c_0,z)$ on $\{\|v\|_{\flat}<R_0\}$, satisfies $|\mathsf{h}^{\mathsf{c}}|\le\mathsf{Q}_w$ there, and
$\mathsf{h}^{\mathsf{c}}(0)=1$;
\item[(b)] $P_0\ge 1/B$ and $Q_0\ge 1/B$, each being the value of its own run at that run's real
centre, and $|Q|\le A$ for $c_0=i\operatorname{Im}\hat{w}$ and $z\in P_{IR}$, where $\mathsf{Q}_w=AB$;
\item[(c)] the closeness of the normalised partition functions at two cutoffs,
Definition~\ref{9.8:def-partition-closeness}, holds: the bound \eqref{9.8:eq-partition-closeness-a} with
$\varepsilon_n\le C_{TC}(N_w;g)\,e_{\mathsf{w}}$, for every $z\in P_{IR}$ and at the zero source;
\item[(d)] $|\operatorname{Im}\hat{w}|\le\upsilon_w$, $\upsilon_w+R_1/2=R_1$, and
$(\mathsf{Q}_w+1)R_1\le R_0/2$.
\end{itemize}
Then
\begin{equation}\label{9.8:eq-derivative-closeness-a}
\Bigl|\partial_{z_1}\cdots\partial_{z_{\mathsf{p}}}\big|_{0}
\bigl[\operatorname{Log}\mathsf{h}^{\mathsf{B}}-\operatorname{Log}\mathsf{h}^{\mathsf{A}}\bigr]\Bigr|
\le C_*(N_w)\,e_{\mathsf{w}}\Bigl(\frac{2\mathsf{p}}{R_1}\Bigr)^{\mathsf{p}}\prod_{i=1}^{\mathsf{p}}\|f_i\|_{\sharp},
\end{equation}
where $C_*(N_w):=2C_{TC}(N_w;g)\,\mathsf{Q}_w(1+\mathsf{Q}_w)$ and $\operatorname{Log}$ is the principal branch of
the logarithm, which is defined at $\mathsf{h}^{\mathsf{A}}$ and $\mathsf{h}^{\mathsf{B}}$ on $P_{IR}$ because both
take values in $\bar{D}(1,\tfrac12)$ there (Step~1 of the proof).

The bound \eqref{9.8:eq-derivative-closeness-a} is applied at run lengths $n\ge n_0(N_w)$, defined as follows.
Here $\ell_r$ and $\mathsf{N}(\theta_r)$ are two length thresholds at depth $r$, the second depending on the
depth only through the inverse coupling $\theta_r$. Let $n_2(g)$ be the least integer $n$ such that half of
every run length $l\ge n$ exceeds both, that is, $l/2>\max\{\ell_r,\mathsf{N}(\theta_r)\}$ at
$r=\lfloor l/2\rfloor$ for every $l\ge n$; let $n_\Upsilon$ be the least integer $n$ such that
$\Upsilon_l\le\min\{\tfrac14,\upsilon_w/(2\varrho_2)\}$ for every $l\ge n$, and put
\begin{equation}\label{9.8:eq-derivative-closeness-2}
n_0(N_w):=\max\Bigl\{2\mathsf{w}+2,\ n_2(g),\ n_\Upsilon\Bigr\}.
\end{equation}
These are floors on $n$ only.
\end{lemma}

\begin{proof}
Conditions (a) and (b) hold for both runs by the normalisation of the sourced integrals at the two
centres and by the finiteness bound on the last lattice $\mathbb{T}_w$; condition (d) records the relations
between the radii of the parameter window $\mathcal{N}'_w$. We derive \eqref{9.8:eq-derivative-closeness-a}
from (a)--(d).

\emph{Step 1: both functions stay near $1$ on $P_{IR}$.} Fix $\mathsf{c}\in\{\mathsf{A},\mathsf{B}\}$ and a source
$v\neq0$ with $\|v\|_{\flat}<R_0$. Since $\|\eta v\|_{\flat}=|\eta|\,\|v\|_{\flat}$ for $\eta\in\mathbb{C}$, the
complex line through the zero source in the direction $v$ meets the domain of (a) in the disc
$\{|\eta|<R_0/\|v\|_{\flat}\}$, of radius larger than $1$. On this disc put
\begin{equation*}
\phi(\eta):=\frac{\mathsf{h}^{\mathsf{c}}(\eta v)-1}{\mathsf{Q}_w+1}.
\end{equation*}
By (a), $\phi$ is holomorphic, $|\phi|\le(\mathsf{Q}_w+1)/(\mathsf{Q}_w+1)=1$, since
$|\mathsf{h}^{\mathsf{c}}-1|\le|\mathsf{h}^{\mathsf{c}}|+1$, and $\phi(0)=0$. The Schwarz lemma on
this disc gives $|\phi(\eta)|\le|\eta|\,\|v\|_{\flat}/R_0$; at $\eta=1$ this reads
\begin{equation}\label{9.8:eq-derivative-closeness-1}
|\mathsf{h}^{\mathsf{c}}(v)-1|\le(\mathsf{Q}_w+1)\,\frac{\|v\|_{\flat}}{R_0}.
\end{equation}
For $v=0$ the bound \eqref{9.8:eq-derivative-closeness-1} holds trivially, since $\mathsf{h}^{\mathsf{c}}(0)=1$.
On $P_{IR}$ each $|z_i|\,\|f_i\|_{\sharp}$ is at most $R_1/(2\mathsf{p})$, so the $\mathsf{p}$ source terms
contribute at most $R_1/2$ to $\|v\|_{\flat}$. In the case of interest the constant part of the source is
$c_0=i\operatorname{Im}\hat{w}$, so $|c_0|\le\upsilon_w$ by (d), and, again by (d),
$\|v\|_{\flat}\le\upsilon_w+R_1/2=R_1$. Hence on $P_{IR}$ the right side of
\eqref{9.8:eq-derivative-closeness-1} is at most $(\mathsf{Q}_w+1)R_1/R_0\le\tfrac12$ by (d), that is,
$\mathsf{h}^{\mathsf{A}}$ and $\mathsf{h}^{\mathsf{B}}$ take values in the closed disc $\bar{D}(1,\tfrac12)$.

\emph{Step 2: closeness of the normalised functions.} Dividing and regrouping,
\begin{equation*}
\mathsf{h}^{\mathsf{B}}-\mathsf{h}^{\mathsf{A}}=\frac{P}{P_0}-\frac{Q}{Q_0}
=\frac{P-Q}{P_0}+\frac{Q(Q_0-P_0)}{P_0Q_0}.
\end{equation*}
By (c) at $z$, $|P-Q|\le\varepsilon_nA$, and by (c) at the zero source, $|Q_0-P_0|\le\varepsilon_nA$.
By (b), $1/P_0\le B$, $1/(P_0Q_0)\le B^2$ and $|Q|\le A$. Therefore, on $P_{IR}$,
\begin{equation*}
|\mathsf{h}^{\mathsf{B}}-\mathsf{h}^{\mathsf{A}}|\le\varepsilon_nAB+\varepsilon_nA^2B^2
=\varepsilon_n\mathsf{Q}_w(1+\mathsf{Q}_w).
\end{equation*}

\emph{Step 3: passage to logarithms.} The disc $\bar{D}(1,\tfrac12)$ is convex and contained in
$\{|\xi|\ge\tfrac12\}\cap\{\operatorname{Re}\xi>0\}$, where the principal branch $\operatorname{Log}$ is
holomorphic with $|\operatorname{Log}'\xi|=|\xi|^{-1}\le2$. Integrating along the segment joining two points
of the disc shows that $\operatorname{Log}$ is $2$-Lipschitz on $\bar{D}(1,\tfrac12)$. By Steps 1 and 2, on $P_{IR}$,
\begin{equation*}
|\operatorname{Log}\mathsf{h}^{\mathsf{B}}-\operatorname{Log}\mathsf{h}^{\mathsf{A}}|
\le2\varepsilon_n\mathsf{Q}_w(1+\mathsf{Q}_w)\le2C_{TC}(N_w;g)\,\mathsf{Q}_w(1+\mathsf{Q}_w)\,e_{\mathsf{w}}
=C_*(N_w)\,e_{\mathsf{w}},
\end{equation*}
the second inequality by the rate $\varepsilon_n\le C_{TC}(N_w;g)\,e_{\mathsf{w}}$ of (c).

\emph{Step 4: Cauchy estimate.} With $c_0=i\operatorname{Im}\hat{w}$ fixed, the function
$\mathcal{L}(z):=\operatorname{Log}\mathsf{h}^{\mathsf{B}}-\operatorname{Log}\mathsf{h}^{\mathsf{A}}$ is holomorphic on an
open set containing the closed polydisc $P_{IR}$, namely the part of $\{\|v\|_{\flat}<R_0\}$ where both
$\mathsf{h}^{\mathsf{c}}$ have positive real part; it contains $P_{IR}$ by Step 1. The polydisc $P_{IR}$ has radii
$R_1/(2\mathsf{p}\|f_i\|_{\sharp})$, $i=1,\dots,\mathsf{p}$, and the Cauchy estimate for the mixed derivative of
first order in each variable gives
\begin{equation*}
\bigl|\partial_{z_1}\cdots\partial_{z_{\mathsf{p}}}\mathcal{L}(0)\bigr|
\le\sup_{P_{IR}}|\mathcal{L}|\prod_{i=1}^{\mathsf{p}}\frac{2\mathsf{p}\|f_i\|_{\sharp}}{R_1}
=\sup_{P_{IR}}|\mathcal{L}|\Bigl(\frac{2\mathsf{p}}{R_1}\Bigr)^{\mathsf{p}}\prod_{i=1}^{\mathsf{p}}\|f_i\|_{\sharp}.
\end{equation*}
Inserting the bound of Step 3 yields \eqref{9.8:eq-derivative-closeness-a}.
\end{proof}

\begin{lemma}[Increments of the Schwinger functions between consecutive cutoffs]
\label{9.8:lem-schwinger-increments}
Assume conditions (J1)--(J5) of Definition~\ref{9.1:def-tuned-family}, displayed in
\eqref{9.1:eq-tuned-family-a}, and the three further standing hypotheses of
Definition~\ref{9.1:def-tuned-family} under which Lemma~\ref{9.8:lem-exponent-localisation} is
stated, the second of them in its frozen and local form. Assume also the closeness
\eqref{9.8:eq-partition-closeness-a} of the normalised partition functions at two cutoffs, with
rate $\varepsilon_n$ (Definition~\ref{9.8:def-partition-closeness}).
Let $m_0\ge 0$; the run of length $n$ has its bare datum on the torus $\mathbb{T}_{m_0+n}$, in the
setting of clause (v) of Theorem~\ref{3.1:thm-main}. Let
$n\ge n_0(N_w)$, with $n_0(N_w)$ as in
\eqref{9.8:eq-derivative-closeness-a}. Let $f=(f_1,\dots,f_{\mathsf{p}})$, $\mathsf{p}\ge 2$, be
test functions, normed by $\|\cdot\|_{\sharp}$ of Definition~\ref{1.3:def-curvature-field}, whose
supports are pairwise separated by at least $\mathsf{d}$, and let $D$ be the volume
constant of the setting of clause (v) of Theorem~\ref{3.1:thm-main}. For $\hat w\in\mathcal{N}'_w$
put
\begin{equation*}
G_n(\hat w):=S^n(a_n-\hat w;f),
\end{equation*}
with $a_n$ as in \eqref{9.1:eq-tuned-family-a}. Here $S^n(\,\cdot\,;f)$ denotes the connected
Schwinger function of $f_1,\dots,f_{\mathsf{p}}$ for the run of length $n$ of
Definition~\ref{9.1:def-tuned-family}, as a function of its first argument. Put
\begin{align*}
C_1^w&:=C_{1.26}\sum_{t\ge0}(1+\mathsf{w}+t)\Bigl(\frac{DL^t}{M}+4\Bigr)^4
 e^{-\frac14\kappa^{c}\left[\mathsf{d}L^t/4M-3\right]_+},\\
C_2^w&:=2C_eC_{1.26}\,e^{\frac34\kappa}\sum_{t\ge0}\Bigl(\frac{DL^t}{M}+4\Bigr)^4
 e^{-\kappa\mathsf{d}L^t/32M},\qquad c:=\frac{\kappa}{32}.
\end{align*}
Here $\kappa^{c}$ is the constant in the exponential decay factor of the bound of
Lemma~\ref{9.8:lem-exponent-localisation} (the superscript is part of its name and is not the
constant $c$), $C_{1.26}$ is the constant of the summation bound over localisation domains
containing a fixed cube recalled in the proof below, and $C_e$ is the constant of the bound on
single boundary exponents recalled there. Let $\mathfrak{f}_n$ be the increment bound and $M_G$
the constant displayed with the conclusion of the uniqueness theorem, the three summands of
$\mathfrak{f}_n$ being formed with
the constants $C_1^w$, $C_2^w$ and $c$ above. Then, with $\mathfrak{n}_n$, $\Upsilon_n$ and
$N_n$ as in Proposition~\ref{9.6:prop-depth-zero-match}:
\begin{itemize}
\item[(a)] for every $\hat w\in\mathcal{N}'_w$,
\begin{equation*}
\bigl|\mathfrak{n}_n'(\hat w)^{\mathsf{p}}\,G_{n+1}(\mathfrak{n}_n(\hat w))-G_n(\hat w)\bigr|
\le\mathfrak{f}_n\prod_i\|f_i\|_{\sharp};
\end{equation*}
\item[(b)] for every $\hat w\in\mathcal{N}'_w$,
\begin{equation*}
|G_{n+1}(\hat w)-G_n(\hat w)|\le M_G\Bigl[\frac{2\varrho_2}{\upsilon_w}
+\mathsf{p}\,2^{\mathsf{p}}\Bigr]\Upsilon_n+\mathfrak{f}_n\prod_i\|f_i\|_{\sharp};
\end{equation*}
\item[(c)]
\begin{equation*}
\bigl|N_{n+1}^{-\mathsf{p}}G_{n+1}(0)-N_n^{-\mathsf{p}}G_n(0)\bigr|
\le\Bigl(\frac32\Bigr)^{\mathsf{p}}\mathfrak{f}_n\prod_i\|f_i\|_{\sharp};
\end{equation*}
\item[(d)] $G_n\to G_\infty$ on $\mathcal{N}'_w$ at rate $n^{-1/2}(\log n)^{p_0+1}$; the limit
$G_\infty$ is holomorphic in the interior of $\mathcal{N}'_w$ and real at real $\hat w$.
\end{itemize}
\end{lemma}

\begin{proof}
Throughout, $\mathsf{A}$ and $\mathsf{B}$ are the runs of lengths $n$ and $n+1$ of
Definition~\ref{9.1:def-tuned-family}, $W^{\mathsf{A}}(z)$ and $W^{\mathsf{B}}(z)$ their residual
sources, $z=(z_1,\dots,z_{\mathsf{p}})$, and $\partial_z^{\mathsf{p}}:=\partial_{z_1}\cdots
\partial_{z_{\mathsf{p}}}$. For a run $\mathsf{c}\in\{\mathsf{A},\mathsf{B}\}$,
$\mathcal{E}^{\mathsf{c}}$ denotes the sum over depths $r$ and domains $X$ of the boundary
exponents $e^{(r),\mathsf{c}}(X;W^{\mathsf{c}}(z))$, from which the $z$-independent vacuum
responses at the run's centre are subtracted, $Z^{\mathsf{c}}(z)$ the sourced integral of run
$\mathsf{c}$, $\tilde Z^{\mathsf{c}}$ the normalised sourced
integral and $\mathsf{h}^{\mathsf{c}}$ its normalisation by the value at the zero source, as in
Lemma~\ref{9.8:lem-derivative-closeness}; $\operatorname{Log}$ is the principal branch.

\emph{Proof of (a): reduction.} The supports of $f_1,\dots,f_{\mathsf{p}}$ are pairwise
disjoint and $\mathfrak{n}_n$ acts pointwise in the centre. Hence the representation of the
Schwinger functions as mixed source derivatives at the matched centres applies, and the left
member of (a) equals
\begin{equation}\label{9.8:eq-schwinger-increments-1}
\Bigl|\partial_z^{\mathsf{p}}\Bigl\{\bigl[\mathcal{E}^{\mathsf{B}}(W^{\mathsf{B}})
+\operatorname{Log}\mathsf{h}^{\mathsf{B}}\bigr]-\bigl[\mathcal{E}^{\mathsf{A}}(W^{\mathsf{A}})
+\operatorname{Log}\mathsf{h}^{\mathsf{A}}\bigr]\Bigr\}\Big|_{z=0}\Bigr|.
\end{equation}
Indeed, $Z^{\mathsf{c}}=e^{\mathcal{E}^{\mathsf{c}}}\tilde Z^{\mathsf{c}}(0)\,\mathsf{h}^{\mathsf{c}}$,
where $\tilde Z^{\mathsf{c}}(0)$ is the value at the zero source taken at the run's real centre.
This middle factor does not depend on $z$, so its logarithm is
annihilated by $\partial_z^{\mathsf{p}}$. The depths $r$ with $\mathsf{w}\le r\le n-1$ occur in
both runs; the depth $r=n$ occurs in $\mathsf{B}$ only. We split
\eqref{9.8:eq-schwinger-increments-1} into three parts: the depths
$\mathsf{w}\le r\le\lfloor n/2\rfloor$ (compared between the runs), the depths
$r>\lfloor n/2\rfloor$ (estimated run by run), and the $\operatorname{Log}\mathsf{h}$-part.

\emph{Depths $\mathsf{w}\le r\le\lfloor n/2\rfloor$.} Here $r<n$. By the localisation of the
difference of boundary exponents (Lemma~\ref{9.8:lem-exponent-localisation}), the difference of
the depth-$r$ sums of the two runs equals $\sum_Y G_r(Y;z)+c_r$, up to the $z$-independent vacuum
responses subtracted in $\mathcal{E}^{\mathsf{A}}$ and $\mathcal{E}^{\mathsf{B}}$. The constant
$c_r$ and these vacuum responses are $z$-free, so $\partial_z^{\mathsf{p}}$ annihilates them.
Each $G_r(Y;\cdot)$ reads $z$ only through the source
restricted to $Y^+$, so it depends on $z_i$ only if $Y^+\cap\operatorname{supp}f_i\neq\emptyset$.
The mixed derivative of $G_r(Y;\cdot)$ is therefore zero unless $Y^+$ meets all
$\mathsf{p}\ge2$ supports. Put $t:=r-\mathsf{w}$. In that case the support separation gives
\begin{equation*}
4ML^{-t}\bigl(d(Y)+3\bigr)\ge\mathsf{d};
\end{equation*}
this follows from the separation estimate in the proof of the correlation theorem, together with
the property of localisation domains used there to pass from the supports met by $Y^+$ to the
tree length $d(Y)$. Moreover, $Y$
contains a depth-$r$ $M$-cube $\square$ at distance at most one such cube from
$\operatorname{supp}f_1$, and
there are at most $(DL^t/M+4)^4$ such anchors $\square$.

Lemma~\ref{9.8:lem-exponent-localisation} bounds $G_r(Y;z)$ on $P_{UV}$, with the constants
$C_\eta$, $\delta$ and $\kappa^{c}$ of its bound, and $G_r(Y;\cdot)$ is
holomorphic near $P_{UV}$. Cauchy's inequality on $P_{UV}$ therefore gives
\begin{equation*}
\bigl|\partial_z^{\mathsf{p}}G_r(Y;z)|_{z=0}\bigr|\le C_\eta(1+r)e_r\,
e^{-(1-\delta)\kappa^{c}d(Y)}\Bigl(\frac{2\mathsf{p}}{R_0}\Bigr)^{\mathsf{p}}
\prod_i\|f_i\|_{\sharp}.
\end{equation*}
We split the exponential as
\begin{equation*}
e^{-(1-\delta)\kappa^{c}d}=e^{-\frac14\kappa^{c}d}\,e^{-(\frac34-\delta)\kappa^{c}d}.
\end{equation*}
The displayed geometric estimate is equivalent to $d(Y)\ge\mathsf{d}L^t/4M-3$, and also
$d(Y)\ge0$. Hence the first factor is at most
$e^{-\frac14\kappa^{c}[\mathsf{d}L^t/4M-3]_+}$ for every contributing $Y$. For the second factor
we use the summation bound over domains containing a fixed cube, from the proof of the
correlation theorem, at rate $(\frac34-\delta)\kappa^{c}$:
\begin{equation*}
\sum_{Y\supseteq\square}e^{-(\frac34-\delta)\kappa^{c}d(Y)}\le C_{1.26},
\end{equation*}
the term with $d(Y)=0$ included. Summing first over $Y\supseteq\square$ and then over the at most
$(DL^t/M+4)^4$ anchors bounds the contribution of depth $r$. We then sum over $r$, using
$1+r=1+\mathsf{w}+t$. By \eqref{9.1:eq-rate-constants-a}, $e_r=\rho^{n-r}+\Pi_n$; since
$n-r\ge n/2$ and $0<\rho<1$, this gives $e_r\le\rho^{n/2}+\Pi_n$. Altogether
\begin{equation}\label{9.8:eq-schwinger-increments-2}
\sum_{r=\mathsf{w}}^{\lfloor n/2\rfloor}\sum_Y\bigl|\partial_z^{\mathsf{p}}G_r(Y;z)|_{z=0}\bigr|
\le C_\eta\Bigl(\frac{2\mathsf{p}}{R_0}\Bigr)^{\mathsf{p}}\bigl(\rho^{n/2}+\Pi_n\bigr)
C_1^w\prod_i\|f_i\|_{\sharp}.
\end{equation}
This is the first summand of $\mathfrak{f}_n$.

\emph{Depths $r>\lfloor n/2\rfloor$, run by run.} These depths are not compared. Each run
$\mathsf{c}\in\{\mathsf{A},\mathsf{B}\}$ is estimated separately, and this includes the depth
$r=n$, which only $\mathsf{B}$ carries. For $z\in P_{UV}$ the boundary exponents are
holomorphic, read $z$ only through the source restricted to $X^+$, and satisfy
\begin{equation*}
|e^{(r),\mathsf{c}}(X;W^{\mathsf{c}}(z))|\le C_e\,e^{-(1-\delta)\kappa d(X)}.
\end{equation*}
Cauchy's inequality on $P_{UV}$ therefore gives
\begin{equation*}
\bigl|\partial_z^{\mathsf{p}}e^{(r),\mathsf{c}}(X;W^{\mathsf{c}}(z))|_{z=0}\bigr|
\le C_e\,e^{-(1-\delta)\kappa d(X)}\Bigl(\frac{2\mathsf{p}}{R_0}\Bigr)^{\mathsf{p}}
\prod_i\|f_i\|_{\sharp}.
\end{equation*}
The geometry is the same as for the first part: only $X$ with $X^+$ meeting all supports
contribute, and for these $d(X)\ge\mathsf{d}L^t/4M-3$ with $t=r-\mathsf{w}$. The number of anchors
is at most $(DL^t/M+4)^4$, and the factor $e^{-(\frac34-\delta)\kappa d}$ is summed to
$C_{1.26}$ as before. Put $x:=\mathsf{d}L^t/4M$. Since $[x-3]_+\ge x-3$,
\begin{equation*}
e^{-\frac14\kappa[x-3]_+}\le e^{\frac34\kappa}e^{-\kappa x/4}
=e^{\frac34\kappa}e^{-\kappa\mathsf{d}L^t/16M}.
\end{equation*}
Since $r>\lfloor n/2\rfloor$ we have $t\ge n/2-\mathsf{w}$, so $L^t\ge L^{n/2-\mathsf{w}}$ and
\begin{equation*}
e^{-\kappa\mathsf{d}L^t/16M}\le e^{-\kappa\mathsf{d}L^{n/2-\mathsf{w}}/32M}\,
e^{-\kappa\mathsf{d}L^t/32M}.
\end{equation*}
Summing over $t\ge0$ and over the two runs, and using this Cauchy estimate, we find that
these depths contribute at most
\begin{equation}\label{9.8:eq-schwinger-increments-3}
\Bigl(\frac{2\mathsf{p}}{R_0}\Bigr)^{\mathsf{p}}e^{-c\,\mathsf{d}L^{n/2-\mathsf{w}}/M}\,
C_2^w\prod_i\|f_i\|_{\sharp},\qquad c=\frac{\kappa}{32}.
\end{equation}
This is the second summand of $\mathfrak{f}_n$.

\emph{The $\operatorname{Log}\mathsf{h}$-part.} The closeness
\eqref{9.8:eq-partition-closeness-a} holds with rate $\varepsilon_n$, and
$n\ge n_0(N_w)$. By the passage from partition functions to mixed source derivatives
(Lemma~\ref{9.8:lem-derivative-closeness}, display \eqref{9.8:eq-derivative-closeness-a}), the
quantity
\begin{equation*}
\bigl|\partial_z^{\mathsf{p}}\bigl(\operatorname{Log}\mathsf{h}^{\mathsf{B}}
-\operatorname{Log}\mathsf{h}^{\mathsf{A}}\bigr)|_{z=0}\bigr|
\end{equation*}
is bounded by the third summand of $\mathfrak{f}_n$, times $\prod_i\|f_i\|_{\sharp}$.

Apply the triangle inequality to \eqref{9.8:eq-schwinger-increments-1}. Together with
\eqref{9.8:eq-schwinger-increments-2}, \eqref{9.8:eq-schwinger-increments-3} and the last bound,
this proves (a).

\emph{Proof of (b)--(d).} Parts (b)--(d) follow from (a) word for word as in the proof of
parts (b)--(d) of the increment lemma for the Schwinger functions of the tuned family on the
disc $|\hat w|\le r_J$, on which Proposition~\ref{9.6:prop-depth-zero-match} defines
$\mathfrak{n}_n$, that proof being run on $\mathcal{N}'_w$ with
the following four inputs in place of the ones used there. First, for
$n\ge n_0(N_w)$ the map $\hat w\mapsto\mathfrak{n}_n(\hat w)$ takes $\mathcal{N}'_w$ into
the set $\{\hat w+v:\hat w\in\mathcal{N}'_w,\ |v|<\upsilon_w/2\}$. Second, on that set
$|G_n'|,|G_{n+1}'|\le 2M_G/\upsilon_w$. Third,
Proposition~\ref{9.6:prop-depth-zero-match} gives the bound
$|\mathfrak{n}_n(\hat w)-\hat w|\le\Upsilon_n|\hat w|$ and the identity
$N_{n+1}\mathfrak{n}_n'(0)=N_n$. Fourth, the exponent $\kappa_0$ appearing in condition (J2) of
\eqref{9.1:eq-tuned-family-a} satisfies $\kappa_0\ge10$.
\end{proof}

\begin{theorem}[The uniqueness theorem]\label{9.8:thm-uniqueness-theorem}
The gauge group is $\mathrm{SU}(2)$. Fix the numbers $\theta_A\in\,]3/5,1[$, $\beta_I=1/5$,
$\alpha=1/32$, the H\"older exponent $\beta_H$ and its decrement $\varepsilon_H$ of
Notation~\ref{9.1:not-rho-window}, and $\beta^{\flat}=\beta_H-\varepsilon_H/2$ with
\begin{equation*}
\theta_A<\beta^{\flat}<\beta_H\le\beta_0^H:=\frac{1+\theta_A}{2}<1,
\end{equation*}
and the index letters $\theta_1,\theta_2,\theta',\mathfrak{i}$ of
Notation~\ref{9.1:not-rho-window}. Then let $L\ge L_0$ be odd. Let the auxiliary parameters
$\mathfrak{p}$ of Notation~\ref{2.1:not-prefix} be chosen in the order of
Notation~\ref{9.1:not-constant-order} and satisfy the conditions listed there, among them the
two floors
$\kappa\ge 10\log L$ and $\kappa^c\ge10\log L$, where $\kappa^c$ is the decay constant that
enters the constant $C_1^w$
of Lemma~\ref{9.8:lem-schwinger-increments}. Let the H\"older exponents satisfy
$\beta_1<\beta_2=\beta_H-\varepsilon_H<\beta_H$, where $\beta_H$ is the H\"older exponent of
Notation~\ref{9.1:not-rho-window}, $\beta_1:=\min\{\hat\alpha,\beta_1^{\circ}\}$, $\hat\alpha$ is
the exponent in the window for the kernel rate $q$ of Notation~\ref{9.1:not-rate-constants}, and
$\beta_1^{\circ}$ is a further H\"older exponent (written $\beta_1$ there), fixed in the order of
choice of Notation~\ref{9.1:not-constant-order}. Take the analysis constants of the correlation
theorem; the two sets $\flat$ and $c$ of the comparison table, fixed in
one table in the order of Notation~\ref{9.1:not-constant-order}, subject to the two conditions
imposed on them there and to the ceilings on the constants $\alpha^c$ listed there; the constants
\begin{equation*}
C_V:=4^{C_d}e^{\kappa_*},\qquad C_{\mathfrak{K}}:=3C_V+5,
\end{equation*}
where $C_d$ and $\kappa_*$ are constants fixed earlier in the order of choice of
Notation~\ref{9.1:not-constant-order};
the constants $\varepsilon_1,c_\varrho,k_\sharp,C_\sigma$ and $q,q_a,\rho,\Lambda,C_T,C_E$ of
Notation~\ref{9.1:not-rate-constants}, with $\rho$ in the window of
Notation~\ref{9.1:not-rho-window}; any $r_J>0$, and
\begin{equation*}
R:=4r_J,\qquad \varrho_2:=r_J/2,\qquad t_\circ:=c_2+r_J;
\end{equation*}
and the coupling ceiling $\gamma_U$ of Notation~\ref{9.1:not-constant-order}. Assume the
following hypotheses.
\begin{itemize}
\item[(i)] The standing hypothesis of the correlation theorem, and the statements
(J1)--(J5) of
Definition~\ref{9.1:def-tuned-family}, displayed in \eqref{9.1:eq-tuned-family-a}, on the
class $\mathfrak{S}$ of tuned runs.
\item[(ii)] The transplanted one-lattice response bounds of
Definition~\ref{9.1:def-transplant-response}.
\item[(iii)] The bounds of Definition~\ref{9.1:def-comparison-families} on the components
$P$, $A$ and $R''$ of the comparison families, for frozen and for local lists.
\item[(iv)] The conjunction of the bound on the component $R'$ of
the comparison families displayed in \eqref{9.1:eq-comparison-families-a} and of the closeness
of the normalised partition functions at two cutoffs,
Definition~\ref{9.8:def-partition-closeness}, in the form
\eqref{9.8:eq-partition-closeness-a}.
\end{itemize}
Then for every $g\le\gamma_U$, with $X:=g^{-2}$ and $\mathsf{w}:=\mathsf{w}(g)$ the stop depth
of the runs, depending on $g$ only, that enters (c) and (e), there are real
numbers $a_n$, one for every $n\ge0$, depending on $(L,\mathfrak{p},g,n)$ and on no torus, such
that the following hold. Here $\lambda:=c_0/2$, which is the constant $\lambda_U$ of clause (v-a)
of Theorem~\ref{3.1:thm-main},
and a run of length $n$ started from a number
$a$ is the run with scaffold $\Theta_X$ of Definition~\ref{9.1:def-tuned-family} whose initial
coupling is $x_{(n)}=a$.
\begin{itemize}
\item[(a)] \emph{Exact pin at depth zero; asymptotic freedom.} The run of length $n$ started
from $a_n$ is alive in the tube and has $x_{(0)}=X$ exactly, and it is the only real alive run
of length $n$ with this property. For $0\le l\le n$,
\begin{equation}\label{9.8:eq-uniqueness-theorem-a}
\begin{gathered}
\lambda n-2c_2\le a_n-X\le3\lambda n+2c_2,\\
\lambda(n-l)-c'\le a_n-a_l\le3\lambda(n-l)+c',
\end{gathered}
\end{equation}
where $\mathsf{A}$ is the constant of (J1),
\begin{equation*}
c':=2c_2+\mathsf{A}C_E\bigl[(1-\rho)^{-2}+T_\Pi\bigr],\qquad
T_\Pi:=\sup_{l}\,l\sum_{i\ge l}\Pi_i<\infty,
\end{equation*}
and $\Pi_i$ is the power-law floor of \eqref{9.1:eq-rate-constants-a}.
\item[(b)] \emph{The extra ultraviolet factor.} The matching map $\mathfrak{n}_n$ of
Proposition~\ref{9.6:prop-depth-zero-match} is holomorphic on the disc
$\{|\hat w|<2\varrho_2\}$, that is $\{|\hat w|<r_J\}$, real on the reals, and
$\mathfrak{n}_n(0)=0$; for every
$\hat w$ with $|\hat w|<r_J$ the run of length $n+1$ started from $a_{n+1}-\mathfrak{n}_n(\hat w)$
and the run of length $n$ started from $a_n-\hat w$ have equal coupling at depth zero; and
\begin{equation*}
|\mathfrak{n}_n(\hat w)-\hat w|\le\Upsilon_n|\hat w|,\qquad
\Upsilon_n\le C\,n^{-3/2}(\log n)^{p_0+1}.
\end{equation*}
The number $N_n:=\partial\zeta_{(0)}/\partial\zeta$, evaluated for the run of length $n$ at the
tuned point $\zeta=\zeta_*^n(0)$, lies in $[2/3,4/3]$; moreover
$N_{n+1}\mathfrak{n}_n'(0)=N_n$, and the sequence $(N_n)$ converges.
\item[(c)] \emph{Increments, on every torus.} Let $m_0\ge0$, $\mathsf{p}\ge2$, and let
$f_1,\dots,f_{\mathsf{p}}$ be Lipschitz functions on the continuum torus of side
$2L^{m_0}\ell_g$ whose supports are at pairwise distance at least $\mathsf{d}>0$; let $D$ be the
quantity of the setting of clause (v) of Theorem~\ref{3.1:thm-main} that enters the constants
$C_1^w$ and $C_2^w$ of Lemma~\ref{9.8:lem-schwinger-increments}. Put
\begin{equation*}
S^n(x;f):=\partial_{z_1}\cdots\partial_{z_{\mathsf{p}}}\big|_{z=0}\log Z,
\end{equation*}
where $Z$ is the partition function on the bare datum $(m_0+n,x)$ at the source
$w=\sum_iz_if_i$. Put $R_0:=3r_J/64$, $R_1:=R_0/(2\mathsf{Q}_w+2)$, $\upsilon_w:=R_1/2$ and
\begin{equation*}
\mathcal{N}'_w:=\bigl\{\hat w\in\mathbb{C}:|\operatorname{Re}\hat w|\le\varrho_2/2,\
|\operatorname{Im}\hat w|<\upsilon_w\bigr\}.
\end{equation*}
Then for every $n\ge n_0(N_w)$, with $n_0(N_w)$ as in \eqref{9.8:eq-derivative-closeness-a},
and every $\hat w\in\mathcal{N}'_w$,
\begin{align}
\bigl|S^{n+1}(a_{n+1}-\hat w;f)-S^n(a_n-\hat w;f)\bigr|
&\le M_G\bigl[2\varrho_2/\upsilon_w+\mathsf{p}2^{\mathsf{p}}\bigr]\Upsilon_n
+\mathfrak{f}_n\prod_i\|f_i\|_\sharp,\label{9.8:eq-uniqueness-theorem-1}\\
\bigl|N_{n+1}^{-\mathsf{p}}S^{n+1}(a_{n+1};f)-N_n^{-\mathsf{p}}S^n(a_n;f)\bigr|
&\le(3/2)^{\mathsf{p}}\,\mathfrak{f}_n\prod_i\|f_i\|_\sharp,\label{9.8:eq-uniqueness-theorem-2}
\end{align}
where
\begin{align*}
\mathfrak{f}_n:={}&\bigl[C_1^wC_\eta(\rho^{n/2}+\Pi_n)
+C_2^we^{-\kappa\mathsf{d}L^{n/2-\mathsf{w}}/(32M)}\bigr](2\mathsf{p}/R_0)^{\mathsf{p}}\\
&+C_*(N_w)(\rho^{n-\mathsf{w}}+\Pi_n)(2\mathsf{p}/R_1)^{\mathsf{p}},
\end{align*}
the power-law floor satisfies
\begin{equation*}
\Pi_n\le C\,n^{-(\kappa_0-4)/2}(\log n)^{c_3},
\end{equation*}
where $\kappa_0$ is the exponent of (J2) in Definition~\ref{9.1:def-tuned-family} and $c_3$ the
exponent of the logarithm in the power-law floor of \eqref{9.1:eq-rate-constants-a},
and
\begin{gather*}
C_*(N_w):=2C_{TC}\mathsf{Q}_w(1+\mathsf{Q}_w),\\
M_G:=\Bigl[C^w_{UV}(2\mathsf{p}/r_J)^{\mathsf{p}}
+\log2\cdot\bigl((4\mathsf{Q}_w+4)\mathsf{p}/r_J\bigr)^{\mathsf{p}}\Bigr]
\prod_i\|f_i\|_\sharp,
\end{gather*}
with $C_{TC}$ the constant of \eqref{9.8:eq-partition-closeness-a}, $C_1^w$, $C_2^w$,
$C_\eta$ the constants of Lemma~\ref{9.8:lem-schwinger-increments}, and $C^w_{UV}$ the constant
entering $M_G$ in part (b) of that lemma. The numbers
$C_1^w$, $C_2^w$, $C_\eta$, $\Upsilon_n$, $\Pi_n$ and $\rho$ do not depend on $m_0$.
\item[(d)] \emph{The limit family.} As $n\to\infty$, without passing to subsequences,
$S^n(a_n-\hat w;f)\to S^{(\hat w)}(f)$ uniformly on $\mathcal{N}'_w$, and
\begin{equation*}
\bigl|S^n(a_n-\hat w;f)-S^{(\hat w)}(f)\bigr|
\le C(\mathsf{p},\mathsf{d},D,N_w)\prod_i\|f_i\|_\sharp\cdot n^{-1/2}(\log n)^{p_0+1}.
\end{equation*}
The map $\hat w\mapsto S^{(\hat w)}$ is real-analytic on $]{-\varrho_2/2},\varrho_2/2[$,
holomorphic on the strip $|\operatorname{Im}\hat w|<\upsilon_w$ (holomorphy is asserted on this
strip only), and real on the reals. Every sequence of real numbers $a'_n$ with
$|a'_n-a_n|\le\varrho_2/2$ satisfies
\begin{equation*}
S^n(a'_n;f)-S^{(a_n-a'_n)}(f)\to0.
\end{equation*}
For real $\hat w$ the functionals $S^{(\hat w)}$ are symmetric, covariant under the lattice
symmetries and reflection positive, and they are translation invariant for
$f_i\in C^2$.
\item[(e)] \emph{Independence of the torus.} The numbers $a_n$, $\mathfrak{n}_n$, $N_n$,
$\Upsilon_n$, $\Pi_n$, $\rho$, and every constant of
Definitions~\ref{9.1:def-tuned-family}, \ref{9.1:def-transplant-response},
\ref{9.1:def-comparison-families} and~\ref{9.1:def-coupling-ceilings}, of
Notation~\ref{9.1:not-rate-constants}, of Propositions~\ref{9.6:prop-two-run-estimate}
and~\ref{9.6:prop-depth-zero-match} and of Lemma~\ref{9.6:lem-torus-free-offset}, are
independent of the torus. The exponent $m_0$ enters
only through $N_w$, and only in $\mathsf{Q}_w$, $C_*(N_w)$, $\upsilon_w$ and $n_0(N_w)$. The
parameters $N$ and $N_0$ of the construction of the runs (here $N$ does not denote the rank of
the gauge group, which is fixed to $2$) and the scaffold radii $\check{R}_k$ at level $k$
(written $R_k$ in
Ba{\l}aban's papers) enter (a)--(d) only through $\gamma_U$ and the numbers
$\mathsf{w}(g)$, $n_2(g)$ and $\ell_r$, which depend on $g$ only (floors on $n$ and the bands
of own levels), and, inside hypothesis (iv), only in the way permitted by the
condition of the order of choice of Notation~\ref{9.1:not-constant-order} that governs the
dependence of those hypotheses on these quantities.
\end{itemize}
\end{theorem}

\begin{proof}
Throughout, for $i\ge0$ we write $\zeta^{[i]}_{(r)}$, $x^{[i]}_{(r)}$ and $b^{[i]}_{(r)}$ for
the running shift, the coupling and the step shift at depth $r$ of the run of length $i$
started from $a_i$. By Definition~\ref{9.1:def-tuned-family} a run of length $i$ with initial
shift $\zeta$ is given by
\begin{equation}\label{9.8:eq-uniqueness-theorem-3}
\zeta_{(i)}:=\zeta,\qquad \zeta_{(r)}:=\zeta_{(r+1)}+b_{(r)}-c_0\quad(0\le r<i),\qquad
x_{(r)}:=\theta_r-\zeta_{(r)},
\end{equation}
with the ramp $\theta_r=X+c_0r$; in particular $\theta_0=X$. Starting from $a$ means
$x_{(i)}=a$, that is $\zeta=\theta_i-a$.

\emph{Clause (a).} Statement (J3) of Definition~\ref{9.1:def-tuned-family}, applied with
$t=0$, gives for every $n$ exactly one real run of length $n$ that is alive in the tube and has
depth-zero coupling $x_{(0)}=X$; its initial shift is $\zeta_*^n(0)$, and we define
$a_n:=\theta_n-\zeta_*^n(0)$, as in \eqref{9.1:eq-tuned-family-a}. Then $a_n$ is real, the run
of length $n$ started from $a_n$ is this run, and its uniqueness is the uniqueness in (J3).
Summing the recursion \eqref{9.8:eq-uniqueness-theorem-3} over $0\le r<n$ gives
\begin{equation*}
\zeta^{[n]}_{(0)}=\zeta_*^n(0)+\sum_{u<n}b^{[n]}_{(u)}-c_0n .
\end{equation*}
Since $x^{[n]}_{(0)}=X=\theta_0$, we have $\zeta^{[n]}_{(0)}=0$, hence
$\zeta_*^n(0)=c_0n-\sum_{u<n}b^{[n]}_{(u)}$, and therefore
\begin{equation}\label{9.8:eq-uniqueness-theorem-4}
a_n-X=\theta_n-\zeta_*^n(0)-X=\sum_{u<n}b^{[n]}_{(u)} .
\end{equation}
The step shifts are controlled by the asymptotic-freedom bound for the step shifts, which
concerns the runs of $\mathfrak{S}$ alive in the tube. Its input is the bound
$|b_{(u)}-c_0|\le\lambda$ at every
depth $u$, where $c_0=2\lambda$ is the one-loop constant of
Definition~\ref{1.2:def-flow-constants} (written $\beta^0$ in \cite[(0.31)]{B-CMP109}); this
input is printed without proof in \cite[Theorem~2, (0.31)]{B-CMP109}; cited as printed. In the
form in which it enters here, with its defect constant $2c_2$, the asymptotic-freedom bound for
the step shifts
states that the step shifts of the run of length $n$ started from $a_n$ satisfy, for
$0\le l\le n$,
\begin{equation}\label{9.8:eq-uniqueness-theorem-6}
\lambda(n-l)-2c_2\le\sum_{l\le u<n}b^{[n]}_{(u)}\le3\lambda(n-l)+2c_2 .
\end{equation}
With $l=0$, \eqref{9.8:eq-uniqueness-theorem-6} and \eqref{9.8:eq-uniqueness-theorem-4} give
the first line of \eqref{9.8:eq-uniqueness-theorem-a}.

For the second line let $0\le l\le n$. The same summation as in
\eqref{9.8:eq-uniqueness-theorem-4}, carried only over the depths $u<l$, gives
$x^{[n]}_{(l)}=X+\sum_{u<l}b^{[n]}_{(u)}$; and $x^{[l]}_{(l)}=a_l$ because the run of length
$l$ starts from $a_l$. Together with \eqref{9.8:eq-uniqueness-theorem-4} this yields the
decomposition
\begin{equation}\label{9.8:eq-uniqueness-theorem-5}
\begin{gathered}
a_n-a_l=\sum_{l\le u<n}b^{[n]}_{(u)}+\bigl(x^{[n]}_{(l)}-x^{[l]}_{(l)}\bigr),\\
x^{[n]}_{(l)}-x^{[l]}_{(l)}=\sum_{i=l}^{n-1}\bigl(x^{[i+1]}_{(l)}-x^{[i]}_{(l)}\bigr).
\end{gathered}
\end{equation}
The first sum in \eqref{9.8:eq-uniqueness-theorem-5} lies between $\lambda(n-l)-2c_2$ and
$3\lambda(n-l)+2c_2$ by \eqref{9.8:eq-uniqueness-theorem-6}. The second term is bounded in
absolute value by $\sum_{i\ge l}d^{(i,i+1)}_l$, where
$d^{(i,i+1)}_l$ is the coupling distance at depth $l$, in the sense of the two-run estimate
(Proposition~\ref{9.6:prop-two-run-estimate}), between the run of length $i$ started from
$a_i$ and the run of length $i+1$ started from $a_{i+1}$. For $l=0$ this term vanishes, since
$x^{[i]}_{(0)}=X$ for every $i$. For $l\ge1$ we apply the two-run estimate to this pair with
$r_0=0$ and $r=l-1$. Its standing assumptions are (J1), the response bounds of
Definition~\ref{9.1:def-transplant-response}, the bounds of
Definition~\ref{9.1:def-comparison-families} and the ceilings of
Definition~\ref{9.1:def-coupling-ceilings}; the first three are supplied by hypotheses (i),
(ii), and (iii) together with the first conjunct of (iv), and the ceilings are conditions on
$g$ alone, in force for $g\le\gamma_U$. Its hypotheses on the pair hold: both runs are alive by
the first part of (a), and both
have $\zeta_{(0)}=0$, so they are matched at depth zero; by part (i) of
Lemma~\ref{9.6:lem-torus-free-offset} the pin and the match sit at depth zero of
$\mathfrak{S}$, so the estimate applies with $r_0=0$. With the shorter run of length $i$, it
gives
\begin{equation*}
d^{(i,i+1)}_l\le\mathsf{A}C_E\Bigl[\frac{\rho^{\,i-l+1}}{1-\rho}+l\,\Pi_i\Bigr].
\end{equation*}
Summing over $i\ge l$ and using $\sum_{i\ge l}\rho^{\,i-l+1}=\rho/(1-\rho)$ we obtain
\begin{equation*}
\sum_{i\ge l}d^{(i,i+1)}_l
\le\mathsf{A}C_E\Bigl[\frac{\rho}{(1-\rho)^2}+l\sum_{i\ge l}\Pi_i\Bigr]
\le\mathsf{A}C_E\bigl[(1-\rho)^{-2}+T_\Pi\bigr],
\end{equation*}
the last step by $\rho<1$ and the definition of $T_\Pi$. Inserting this and the bounds on the
first sum into \eqref{9.8:eq-uniqueness-theorem-5} gives the second line of
\eqref{9.8:eq-uniqueness-theorem-a} with $c'=2c_2+\mathsf{A}C_E[(1-\rho)^{-2}+T_\Pi]$.

\emph{Clause (b).} For $|v|<R$ let $h_n(v):=\zeta^n_{(0)}(\zeta_*^n(0)+v)$ be the
depth-zero running shift of the run of length $n$ with initial shift $\zeta_*^n(0)+v$. The run
of length $n$ started from $a_n-\hat w$ has initial shift $\theta_n-a_n+\hat w=\zeta_*^n(0)+\hat w$,
so its depth-zero running shift is $h_n(\hat w)$; likewise the run of length $n+1$ started from
$a_{n+1}-\mathfrak{n}_n(\hat w)$ has depth-zero running shift $h_{n+1}(\mathfrak{n}_n(\hat w))$.
Matching the two runs at depth zero (Proposition~\ref{9.6:prop-depth-zero-match}) gives, for
$|\hat w|\le r_J$, exactly one $\mathfrak{n}_n(\hat w)$ with $|\mathfrak{n}_n(\hat w)|<(8/3)r_J$
and $h_{n+1}(\mathfrak{n}_n(\hat w))=h_n(\hat w)$. Since
$x_{(0)}=\theta_0-\zeta_{(0)}=X-\zeta_{(0)}$,
equality of the depth-zero running shifts is equality of the depth-zero couplings, which is the
second assertion of (b). Since $h_n(0)=\zeta^{[n]}_{(0)}=0$ and
$h_{n+1}(0)=\zeta^{[n+1]}_{(0)}=0$ by clause (a), the point $v=0$ solves
$h_{n+1}(v)=h_n(0)$ with $|v|<(8/3)r_J$, and the uniqueness gives $\mathfrak{n}_n(0)=0$. The
same proposition gives the holomorphy of $\mathfrak{n}_n$ on the disc $\{|\hat w|<r_J\}$ and its
reality on the
reals, the bound $|\mathfrak{n}_n(\hat w)-\hat w|\le\Upsilon_n|\hat w|$ with
$\Upsilon_n\le C n^{-3/2}(\log n)^{p_0+1}$, and the identity $N_{n+1}\mathfrak{n}_n'(0)=N_n$;
it uses the two-run estimate at depth $r_0=0$, which applies by part (i) of
Lemma~\ref{9.6:lem-torus-free-offset}. The location $N_n\in[2/3,4/3]$ is part of the same two
statements, Proposition~\ref{9.6:prop-depth-zero-match} and
Lemma~\ref{9.6:lem-torus-free-offset}. For the convergence of $(N_n)$: since
$\mathfrak{n}_n(0)=0$, dividing the bound $|\mathfrak{n}_n(\hat w)-\hat w|\le\Upsilon_n|\hat w|$
by $|\hat w|$ and letting $\hat w\to0$ gives $|\mathfrak{n}_n'(0)-1|\le\Upsilon_n$, and
$\mathfrak{n}_n'(0)$ is real because $\mathfrak{n}_n$ is real on the reals. Since
$\sum_n n^{-3/2}(\log n)^{p_0+1}<\infty$, we have $\sum_n\Upsilon_n<\infty$, and there is $n_1$
with $\Upsilon_n\le1/2$ for all $n\ge n_1$. For $n\ge n_1$ the identity
$N_{n+1}\mathfrak{n}_n'(0)=N_n$ and $N_n>0$ give
$\log N_{n+1}-\log N_n=-\log\mathfrak{n}_n'(0)$, and the inequality $|\log x|\le2|x-1|$ for
$|x-1|\le1/2$ gives $|\log N_{n+1}-\log N_n|\le2\Upsilon_n$. Hence $(\log N_n)$ is a Cauchy
sequence, and $(N_n)$ converges.

\emph{Clauses (c) and (d).} Fix $m_0$, $\mathsf{p}$ and $f_1,\dots,f_{\mathsf{p}}$ as in (c),
and put $G_n(\hat w):=S^n(a_n-\hat w;f)$ for $\hat w\in\mathcal{N}'_w$. The increments of the
Schwinger functions between consecutive cutoffs (Lemma~\ref{9.8:lem-schwinger-increments})
assume (J1)--(J5), the response bounds of Definition~\ref{9.1:def-transplant-response}, the
bounds of
Definition~\ref{9.1:def-comparison-families} for frozen and for local lists, the ceilings of
Definition~\ref{9.1:def-coupling-ceilings}, and the closeness bound
\eqref{9.8:eq-partition-closeness-a}. The first, second
and fifth of these are hypotheses (i), (ii) and the second conjunct of (iv) of the present
theorem; the bounds of
Definition~\ref{9.1:def-comparison-families} are hypothesis (iii) for the components $P$, $A$,
$R''$ together with the first conjunct of hypothesis (iv) for the component $R'$; the ceilings
of Definition~\ref{9.1:def-coupling-ceilings} are conditions on $g$ alone, in force for
$g\le\gamma_U$. Hence the lemma applies for $n\ge n_0(N_w)$ and $\hat w\in\mathcal{N}'_w$. Its
part (a) holds with the quantity $\mathfrak{f}_n$ of (c), whose constants $C_1^w$, $C_2^w$,
$C_\eta$ and decay constant $\kappa/32$ are those of the lemma. Its part (b) is
\eqref{9.8:eq-uniqueness-theorem-1}, because
$G_{n+1}(\hat w)=S^{n+1}(a_{n+1}-\hat w;f)$. Its part (c) is
\eqref{9.8:eq-uniqueness-theorem-2}, because $G_n(0)=S^n(a_n;f)$. The independence of $C_1^w$,
$C_2^w$, $C_\eta$, $\Upsilon_n$, $\Pi_n$ and $\rho$ from $m_0$ is part of the same lemma. This
proves (c).

Part (d) of the same lemma gives the convergence $G_n\to G_\infty$ on $\mathcal{N}'_w$, for the
whole sequence and uniformly in $\hat w$, at the rate displayed in (d) for $n\ge n_0(N_w)$,
together with the stated regularity of the limit in $\hat w$. We put
$S^{(\hat w)}(f):=G_\infty(\hat w)$. For the finitely many $n<n_0(N_w)$ the bound of (d) is
obtained by enlarging the constant $C(\mathsf{p},\mathsf{d},D,N_w)$. Let $a'_n$ be real with
$|a'_n-a_n|\le\varrho_2/2$, and put $\hat w_n:=a_n-a'_n$. Then $\hat w_n$ is real with
$|\hat w_n|\le\varrho_2/2$, so $\hat w_n\in\mathcal{N}'_w$, and
\begin{equation*}
S^n(a'_n;f)-S^{(a_n-a'_n)}(f)=G_n(\hat w_n)-G_\infty(\hat w_n)\to0
\end{equation*}
by the uniformity of the convergence on $\mathcal{N}'_w$. For real $\hat w$ and each $n$, the
functionals $S^n(a_n-\hat w;\cdot)$ are symmetric, covariant under the lattice symmetries and
reflection positive, by the symmetry statements of the correlation theorem and by (J5). Each of
these properties is a linear identity among, or the non-negativity of a finite linear combination
of, values of the functionals. Such conditions are preserved under pointwise limits, so they
hold for $S^{(\hat w)}$.

It remains to prove translation invariance for $f_i\in C^2$; let $\hat w$ be real. For a
vector $v$ write $f(\cdot-v)$ for the tuple $(f_i(\cdot-v))_i$. The bare lattice of the datum
$(m_0+n,\cdot)$ has $2L^{m_0+n}$ sites a side on the torus of side $2L^{m_0}\ell_g$, hence
spacing $\ell_gL^{-n}$; let $\Gamma_n$ be the group of its translation vectors. Since the
spacing at cutoff $n$ is $L$ times the spacing at cutoff $n+1$, $\Gamma_n\subset\Gamma_{n+1}$.
For $n'\ge0$, $u\in\Gamma_{n'}$ and every $n\ge n'$, covariance under the lattice symmetries gives
$S^n(a_n-\hat w;f(\cdot-u))=S^n(a_n-\hat w;f)$; letting $n\to\infty$ gives
$S^{(\hat w)}(f(\cdot-u))=S^{(\hat w)}(f)$ for every $u\in\bigcup_n\Gamma_n$. Now let $v$ be
arbitrary and choose $v_n\in\Gamma_n$ with $|v-v_n|\le\ell_gL^{-n}$. The functionals
$S^n(a_n-\hat w;\cdot)$ are multilinear in $(f_1,\dots,f_{\mathsf{p}})$, since the source
$w=\sum_iz_if_i$ enters $\log Z$ through observables linear in each $f_i$; hence so is their
pointwise limit $S^{(\hat w)}$. Writing $\psi_{n,i}:=f_i(\cdot-v)-f_i(\cdot-v_n)$, we obtain
\begin{equation*}
S^{(\hat w)}(f(\cdot-v))-S^{(\hat w)}(f(\cdot-v_n))=\sum_{i=1}^{\mathsf{p}}
S^{(\hat w)}\bigl(\varphi^{(n,i)}\bigr),
\end{equation*}
where the tuple $\varphi^{(n,i)}$ has the entries $f_j(\cdot-v_n)$ for $j<i$, the entry
$\psi_{n,i}$ in
place $i$, and the entries $f_j(\cdot-v)$ for $j>i$. The support of $\psi_{n,i}$ lies in the union
of the supports of $f_i(\cdot-v)$ and $f_i(\cdot-v_n)$, so the supports of the entries of
$\varphi^{(n,i)}$ are at pairwise distance at least $\mathsf{d}-|v-v_n|\ge\mathsf{d}/2$ once
$\ell_gL^{-n}\le\mathsf{d}/2$, and $S^{(\hat w)}(\varphi^{(n,i)})$ is defined. Since $f_i\in C^2$,
$\|\psi_{n,i}\|_\infty\le|v-v_n|\,\|\nabla f_i\|_\infty$ and
$\operatorname{Lip}\psi_{n,i}\le|v-v_n|\,\|\nabla^2f_i\|_\infty$, so $\|\psi_{n,i}\|_\sharp\to0$.
The torus bound of the correlation theorem bounds $|S^n(a_n-\hat w;\varphi)|$ by a constant
independent of $n$ times $\prod_i\|\varphi_i\|_\sharp$, for tuples
$\varphi=(\varphi_1,\dots,\varphi_{\mathsf{p}})$ of Lipschitz functions whose
supports are at pairwise distance at least $\mathsf{d}/2$; passing to the limit, the same bound
holds for $S^{(\hat w)}$. Since translation does not change $\|\cdot\|_\sharp$, this gives
$|S^{(\hat w)}(\varphi^{(n,i)})|\le C\|\psi_{n,i}\|_\sharp\prod_{j\ne i}\|f_j\|_\sharp\to0$, hence
\begin{equation*}
S^{(\hat w)}(f(\cdot-v))=\lim_{n\to\infty}S^{(\hat w)}(f(\cdot-v_n))=S^{(\hat w)}(f).
\end{equation*}
This proves (d).

\emph{Clause (e).} By (J1) the data of the runs are holomorphic and independent of the torus.
The definitions and bounds of Definitions~\ref{9.1:def-tuned-family},
\ref{9.1:def-transplant-response}, \ref{9.1:def-comparison-families}
and~\ref{9.1:def-coupling-ceilings}, of Notation~\ref{9.1:not-rate-constants}, of
Propositions~\ref{9.6:prop-two-run-estimate} and~\ref{9.6:prop-depth-zero-match} and of
Lemma~\ref{9.6:lem-torus-free-offset} contain no $N_w$. By part (iii) of
Lemma~\ref{9.6:lem-torus-free-offset} the numbers $a_n$ are independent of the torus, since at
depths below $\mathsf{w}$ the step shifts are the frozen ones, which do not depend on the torus.
Hence $\mathfrak{n}_n$, $N_n$ and
$\Upsilon_n$, which are built in the proof of (b) from the runs started from the $a_n$, are
independent of the torus as well, and in (c) and (d) the exponent $m_0$ enters through $N_w$ in
the places listed in (e).
\end{proof}

\begin{corollary}[First-passage offsets and the identification of subsequential limits. Proof outstanding]
\label{9.8:cor-subsequential-limits}
Assume the hypotheses of the uniqueness theorem (Theorem~\ref{9.8:thm-uniqueness-theorem}). Assume in
addition the supplementary run hypothesis. Let $r_J$ denote the depth parameter of the correlation
theorem. Let $\varpi(r_J,g)$ be the correction term, depending on $r_J$ and on the reference
coupling $g$, that enters the first-passage offset bound \eqref{9.8:eq-subsequential-limits-2}.
Assume that $r_J$ and $\varpi(r_J,g)$ satisfy
\begin{equation}\label{9.8:eq-subsequential-limits-1}
r_J\ \ge\ 6\,(B^{\sharp}+1),\qquad \varpi(r_J,g)\ \le\ 1 .
\end{equation}
Then the following hold.
\begin{enumerate}
\item[(a)] For every bare inverse coupling $x$ whose run, read against its own threshold, has its
first passage at the level $K$,
\begin{equation}\label{9.8:eq-subsequential-limits-2}
|x-a_K|\ \le\ \frac{3}{2}\,\bigl(B^{\sharp}+\varpi(r_J,g)\bigr)\ \le\ \frac{1}{2}\,\varrho_2 ,
\end{equation}
where $a_K$ is the point $a_n$ of Theorem~\ref{9.8:thm-uniqueness-theorem} at $n=K$ and
$\varrho_2=\varrho/2$.
\item[(b)] Every subsequential limit of the sequence of clause~(ii) of
Theorem~\ref{3.1:thm-main} at the complex source $w$, and of the sequence of sourced Schwinger
functions at $w$, is equal to $S^{(\hat w)}$ for some $\hat w$.
\end{enumerate}
Not asserted: the case $\mathsf{p}=1$ of a single test function, which is excluded here by
the power-law floor $\Pi_n$; any statement about $\log Z$ or about $\tilde Z(0)$; Schwinger functions
at coinciding points; loop observables; the groups $\mathrm{SU}(N)$ with $N\ge3$; uniformity in the
side $N_w$ of the last lattice; exponential rates other than those written with a hat in
Theorem~\ref{9.8:thm-uniqueness-theorem}; the third assertion of the uniqueness statement
underlying Theorem~\ref{9.8:thm-uniqueness-theorem}; the two-cutoff response
statement beyond zeroth order.
\end{corollary}

Proof outstanding.

\emph{Route.} Part~(a) is to be derived from the description of the points $a_n$ in
\eqref{9.8:eq-uniqueness-theorem-a}: $a_n$ is the reference inverse coupling $X=g^{-2}$ plus the
accumulated step shifts $b_{(u)}$, $u<n$, each satisfying $|b_{(u)}-\beta^0|\le\lambda$ for one
fixed number $\beta^0$. This description is to be
compared with the first-passage condition on $x$ at $K$, and the comparison
$\tfrac{3}{2}(B^{\sharp}+\varpi(r_J,g))\le\tfrac12\varrho_2$ is to be derived from
\eqref{9.8:eq-subsequential-limits-1}. Part~(b) is to be derived from the identification of the
limits as $S^{(\hat w)}$ in \eqref{9.8:eq-uniqueness-theorem-a}, applied along the given
subsequence.

\begin{definition}[Two independent blocking factors, dilates and germ equivalence]
\label{9.8:def-blocking-factors}
Throughout, $N=2$, and $L_0(2)$ is the constant $L_0(N)$ of hypothesis (H2) of
Notation~\ref{2.2:not-hypotheses} at $N=2$.

\emph{Frames.} Let $L$ and $L'$ be two odd integers with $L\ge L_0(2)$ and
$L'\ge L_0(2)$ which are \emph{multiplicatively independent}:
\begin{equation}\label{9.8:eq-blocking-factors-a}
L^{a}=L'^{\,a'}\ \text{ with } a,a'\in\mathbb{N}\quad\Longrightarrow\quad a=a'=0 .
\end{equation}
$L$ and $L'$ are otherwise arbitrary. Each of $L$, $L'$ is called a \emph{frame}, and
$\mathsf{F}$ denotes either of them. We put
\begin{equation*}
\omega:=\frac{\log L}{\log L'} ,
\end{equation*}
which is irrational by \eqref{9.8:eq-blocking-factors-a} (see the proof).
Every symbol of Theorem~\ref{3.1:thm-main} is read at blocking factor $\mathsf{F}$,
and carries the index $\mathsf{F}$ only where both frames occur in one formula.

\emph{Common window.} A number $\varrho>0$ is fixed, common to both frames; we put
$\varrho_2:=\varrho/2$, $W:=[-\varrho_2/2,\varrho_2/2]$, and write $W^{\circ}$ for its
interior.

\emph{Pins and increments.} For each frame, $a_n=a_n^{\mathsf{F}}$ denotes the exact pin
of clause (v-a) of Theorem~\ref{3.1:thm-main}, and we put
\begin{equation*}
\alpha_n:=a_{n+1}-a_n,\qquad
c'^{\sharp}:=\sup_{g\le\gamma_U}c'(g),\qquad
\bar{\alpha}:=3\lambda_U+c'^{\sharp},
\end{equation*}
where $c'(g)$ is the constant of clause (v-a) of Theorem~\ref{3.1:thm-main}; the
supremum is finite, $c'$ being bounded on $]0,\gamma_U]$ by its value at $\gamma_U$, by
the boundedness statement for $c'$ on $]0,\gamma_U]$ stated with the short-distance
corollary. For every reference coupling $g_{\mathsf{F}}\le\gamma_U$ one has
\begin{equation*}
\sup_n|\alpha_n|\le\bar{\alpha};
\end{equation*}
this is proved below from clause (v-a) of Theorem~\ref{3.1:thm-main} with $n-l=1$.

\emph{Floor on $\varrho$.} For both frames,
\begin{equation}\label{9.8:eq-blocking-factors-b}
\varrho_2\ge 4\,\bar{\alpha}_{\mathsf{F}}
\qquad\text{and}\qquad
\varrho\ge 6\bigl(B^{\sharp}_{\mathsf{F}}+1\bigr).
\end{equation}

\emph{Ceiling on the coupling.} For each frame the reference coupling $g_{\mathsf{F}}$
satisfies
\begin{equation}\label{9.8:eq-blocking-factors-c}
g_{\mathsf{F}}\le\gamma_\Lambda^{\mathsf{F}}
:=\min\bigl\{\gamma_U,\ \gamma_U',\ \gamma_W^{\mathbb{T}},\ \gamma_{\natural}\bigr\},
\end{equation}
where $\gamma_U$ is the threshold of the uniqueness theorem, $\gamma_U'$ the further
threshold of its corollary, $\gamma_W^{\mathbb{T}}$ the threshold of clause (iv) of
Theorem~\ref{3.1:thm-main}, and $\gamma_{\natural}$ the threshold of the short-distance
corollary; we put $X:=g_{\mathsf{F}}^{-2}$.

\emph{Floor on the torus exponent.} The torus exponents $m_0$ considered satisfy
\begin{equation}\label{9.8:eq-blocking-factors-d}
m_0\ge m_{\natural}(g_{\mathsf{F}}),
\end{equation}
where $m_{\natural}$ is the torus-exponent floor of the short-distance corollary; it
incorporates the native-depth discrepancy $j_U$, given by
$j_U-1=\lceil(\varrho_2+c'(g_{\mathsf{F}}))/\lambda_U\rceil$, and hypothesis (H5).

\emph{Objects.} Let $\mathsf{p}\in\{2,3\}$. For a $\mathsf{p}$-tuple
$\vec f=(f_1,\dots,f_{\mathsf{p}})$ we write $S^n_{m_0}(x;\vec f)$ for the connected
$\mathsf{p}$-point function on the bare lattice of $\mathbb{T}_{m_0}$, of side
$2\mathsf{F}^{m_0}\ell_g$ in reference units, with $\ell_g$ as in clauses
(ii-$\mathbb{T}$-a) and (v-b) of Theorem~\ref{3.1:thm-main}, lying $n$ levels below the
reference scale,
at bare inverse coupling $x$. By clause (ii-$\mathbb{T}$-a) of
Theorem~\ref{3.1:thm-main} it is a function of the lattice torus, of $x$ and of the
lattice weights only. For $\hat w\in W$ its \emph{tuned limit} is
\begin{equation*}
S^{(\hat w)}_{m_0}(\vec f):=\lim_{n\to\infty}S^n_{m_0}\bigl(a_n-\hat w;\vec f\bigr),
\end{equation*}
which exists by clauses (ii-$\mathbb{T}$-a) and (v-b) of Theorem~\ref{3.1:thm-main}.

\emph{Dilates.} For $\rho>0$ and a tuple $\vec\varphi$ we put
$\vec{\varphi}^{\,\rho}:=\vec\varphi(\cdot/\rho)$, componentwise and with no prefactor,
so that sup norms are preserved (the sup-normalised dilate), and for a functional $S$
on tuples
\begin{equation*}
(D_\rho S)(\vec\varphi):=S\bigl(\vec{\varphi}^{\,\rho}\bigr).
\end{equation*}

\emph{Test classes.} $\mathfrak{P}_{\mathsf{p}}(r;\vartheta)$ is the test class of the
uniqueness theorem, and $\mathfrak{T}(r;\vartheta)$ that of
clause (iv) of Theorem~\ref{3.1:thm-main}. Both are dilation covariant:
\begin{equation}\label{9.8:eq-blocking-factors-1}
\vec\varphi\in\mathfrak{P}_{\mathsf{p}}(r;\vartheta)
\iff
\vec{\varphi}^{\,\rho}\in\mathfrak{P}_{\mathsf{p}}(\rho r;\vartheta),
\qquad \|\varphi_i^{\rho}\|_\infty=\|\varphi_i\|_\infty ,
\end{equation}
and likewise for $\mathfrak{T}$. We fix $\vartheta\in\,]0,1]$ and a pair
$\Psi=(\Psi_1,\Psi_2)\in\mathfrak{T}(1;\vartheta)$ of class $C^2$, and put
\begin{equation*}
\Psi_r:=\Psi^{r},\qquad \|\Psi\|^2:=\|\Psi_1\|_\infty\|\Psi_2\|_\infty .
\end{equation*}

\emph{Germ equivalence.} For two functionals $S,S'$ on $\mathsf{p}$-tuples of test
functions, $\mathsf{p}\in\{2,3\}$, we write $S\approx S'$ if there exist $C$ and
$r_{\approx}>0$
such that, for every $0<r\le r_{\approx}$ and every tuple
$\vec\varphi\in\mathfrak{P}_{\mathsf{p}}(r;\vartheta)$ of class $C^2$,
\begin{equation}\label{9.8:eq-blocking-factors-e}
\bigl|S(\vec\varphi)-S'(\vec\varphi)\bigr|
\le C\,r^{4-\mathsf{p}}\prod_{i=1}^{\mathsf{p}}\|\varphi_i\|_\infty .
\end{equation}
\end{definition}

\begin{proof}
We justify the assertions made in the definition.

Theorem~\ref{3.1:thm-main} holds at each frame. Its quantifier prefix ranges over all
odd $L\ge L_0(N)$; here $N=2$ and each frame is an odd integer $\ge L_0(2)$.

Under \eqref{9.8:eq-blocking-factors-a}, $\omega\notin\mathbb{Q}$. Indeed, if
$\omega=u/v$ with $u,v$ integers, then $u,v\ge1$ may be taken, since $\omega>0$. Then
$v\log L=u\log L'$, that is $L^{v}=L'^{\,u}$, which contradicts
\eqref{9.8:eq-blocking-factors-a}. Any two coprime odd integers $L,L'\ge L_0(2)$
satisfy \eqref{9.8:eq-blocking-factors-a}. Indeed, let
$L^a=L'^{\,a'}$ with $a,a'\in\mathbb{N}$. If $a=0$, then $L'^{\,a'}=1$ with $L'>1$, so
$a'=0$. If $a\ge1$, then $L'^{\,a'}=L^a>1$, so $a'\ge1$. Since $L>1$, it has a prime
factor $q$, and $q$ divides $L'^{\,a'}$, hence $L'$, contrary to coprimality. Pairs of
this kind exist. For instance, any odd
$L\ge L_0(2)$ together with $L'=L+2$ will do: $\gcd(L,L+2)$ divides $2$, and both numbers
are odd.

Next, $\sup_n|\alpha_n|\le\bar{\alpha}$ at every reference coupling allowed by
\eqref{9.8:eq-blocking-factors-c}. Clause (v-a) of Theorem~\ref{3.1:thm-main} states
\begin{equation}\label{9.8:eq-blocking-factors-2}
\lambda_U(n-l)-c'\le a_n-a_l\le 3\lambda_U(n-l)+c'
\end{equation}
with $c'=c'(g)$ independent of $n$ and $l$. At $n=l$, \eqref{9.8:eq-blocking-factors-2}
gives $c'\ge0$. Comparing its two outer members gives, for $n-l\ge1$,
\begin{equation*}
\lambda_U\ge -\frac{c'}{n-l},
\end{equation*}
and letting $n-l\to\infty$ yields $\lambda_U\ge0$. With $n-l=1$,
\eqref{9.8:eq-blocking-factors-2} reads
\begin{equation*}
\lambda_U-c'\le\alpha_n\le 3\lambda_U+c' .
\end{equation*}
Since $\lambda_U\ge0$, we have $c'-\lambda_U\le 3\lambda_U+c'$. Hence
$|\alpha_n|\le 3\lambda_U+c'(g_{\mathsf{F}})$ for every $n$. By
\eqref{9.8:eq-blocking-factors-c}, $g_{\mathsf{F}}\le\gamma_U$, so
$c'(g_{\mathsf{F}})\le c'^{\sharp}$, and $|\alpha_n|\le\bar{\alpha}$.

Condition \eqref{9.8:eq-blocking-factors-b} can be met. The threshold $\gamma_U$ of the
uniqueness theorem depends on $\varrho$; we write $\gamma_U(\varrho)$ for it, and
$c'^{\sharp}(\varrho)$ for the supremum of $c'(g)$ over $g\le\gamma_U(\varrho)$. The
constant $\lambda_U$ depends on $(N,\mathsf{F})$ only. Since $c'$ is bounded on
$]0,\gamma_U]$ by its value at $\gamma_U$, by the boundedness statement for $c'$ on
$]0,\gamma_U]$ stated with the short-distance corollary, the supremum $c'^{\sharp}$ does
not increase when
$\gamma_U$ is lowered. Fix $\varrho_1>0$, and then take $\varrho$ such that, for both
frames,
\begin{equation*}
\varrho\ge\max\Bigl\{\varrho_1,\ 8\bigl(3\lambda_U+c'^{\sharp}(\varrho_1)\bigr),\
6\bigl(B^{\sharp}_{\mathsf{F}}+1\bigr)\Bigr\},
\end{equation*}
and $\gamma_U(\varrho)\le\gamma_U(\varrho_1)$. Then
$c'^{\sharp}(\varrho)\le c'^{\sharp}(\varrho_1)$, and therefore
\begin{equation*}
\varrho_2=\frac{\varrho}{2}\ge 4\bigl(3\lambda_U+c'^{\sharp}(\varrho_1)\bigr)
\ge 4\bigl(3\lambda_U+c'^{\sharp}(\varrho)\bigr)=4\,\bar{\alpha}_{\mathsf{F}},
\end{equation*}
while $\varrho\ge6(B^{\sharp}_{\mathsf{F}}+1)$ holds by the choice of $\varrho$. This
argument uses that $c'$, as a function of $g$, does not depend on $\varrho$, only the
range $]0,\gamma_U(\varrho)]$ of the supremum depending on $\varrho$; this point is
avoided in the lemma of this chapter that bounds the increments $\alpha_n$ without
recourse to $c'$.

In \eqref{9.8:eq-blocking-factors-1}, the equality of sup norms holds because
$\varphi_i(\cdot/\rho)$ takes exactly the values of $\varphi_i$. The equivalence of
memberships is the dilation covariance of the classes $\mathfrak{P}_{\mathsf{p}}$ and
$\mathfrak{T}$ as defined with the uniqueness theorem and with clause (iv) of
Theorem~\ref{3.1:thm-main}, respectively; it is used here as stated there. Applied to
$\mathfrak{T}$ with
$\rho=r$, it gives $\Psi_r\in\mathfrak{T}(r;\vartheta)$ and
$\|\Psi_{r,1}\|_\infty\|\Psi_{r,2}\|_\infty=\|\Psi\|^2$ for every $r>0$. Moreover,
$\Psi_r$ is of class $C^2$ together with $\Psi$.

Finally, $\approx$ is an equivalence relation. Reflexivity holds with $C=0$, and
symmetry is clear from \eqref{9.8:eq-blocking-factors-e}. For transitivity, suppose
$S\approx S'$ with constants $(C,r_{\approx})$ and $S'\approx S''$ with constants
$(C',r_{\approx}')$. By the triangle inequality, \eqref{9.8:eq-blocking-factors-e} holds
for the pair $(S,S'')$ with constant $C+C'$ and radius
$\min\{r_{\approx},r_{\approx}'\}$.
\end{proof}

\begin{lemma}[Two-sided two-point bounds at tuned coupling]\label{9.8:lem-tuned-two-sided}
Let $\mathsf{F}\in\{L,L'\}$ be one of the two blocking factors of
Definition~\ref{9.8:def-blocking-factors}, at which every constant and function below is taken. Let
$W$ be the offset
window and $a_n$ the pins of that definition, and let $g$ be the reference coupling at which the pins
are read. For $\mathsf{d}>0$ put
\begin{equation*}
s(\mathsf{d}):=\min\{s:\ C_W(s)\,\mathsf{F}^{-s}\le\mathsf{d}\},
\end{equation*}
where $C_W\ge1$ is non-decreasing and polylogarithmic in $\bar x_s$. Assume the following two
statements.
\begin{itemize}
\item[(a)] \emph{Clause~(iv) of Theorem~\ref{3.1:thm-main} on every torus $\mathbb{T}_m$ with
$m\ge m_{\mathbb{T}}(g_-(g))$.} The reference coupling satisfies $g\le\gamma_W^{\mathbb{T}}$, the
coupling ceiling of that clause, and there is a separation ceiling
$\mathsf{d}_T^{\star}(g)>0$,
depending on $g$ only, with the following property. Let $g_0$ be a bare coupling with first-passage
depth $K=K(g_0,g)$ in the sense of clause~(0) of Theorem~\ref{3.1:thm-main}, let
$m\ge m_{\mathbb{T}}(g_-(g))$, and
let $\mathsf{d}>0$ satisfy the depth condition
$\max\{s_W(g),s_{\mathrm{vol}}^{\star}(g),s_T(g)\}\le s(\mathsf{d})$, where $s_W(g)$ is the depth
floor of the two-point witness and $s_{\mathrm{vol}}^{\star}(g)$, $s_T(g)$ are the depth floors of
the volume transfer, all depending on $g$ only; this condition holds whenever
$\mathsf{d}<\mathsf{d}_T^{\star}(g)$. Let further $\mathsf{d}$ satisfy the ultraviolet floor
$K-s(\mathsf{d})\ge k_W$. Then for
every $(f_1,f_2)\in\mathfrak{T}(\mathsf{d};\vartheta)$, with $s=s(\mathsf{d})$,
\begin{equation*}
\begin{split}
\tfrac29\,b_0\,c_{\mathcal K}\,\|f_1\|_\infty\|f_2\|_\infty\,\bigl(g^{-2}+B^{\sharp}(s+1)\bigr)^{-2}
&\le S^T_{2;K,m}(f_1,f_2)\\
&\le\tfrac83\,b_0\,C_{\mathcal K}\,\|f_1\|_\infty\|f_2\|_\infty\,
\bigl(g^{-2}+\lambda^{\sharp}s-c_5\bigr)^{-2},
\end{split}
\end{equation*}
where $c_{\mathcal K},C_{\mathcal K}$ depend on $\vartheta$ only.
\item[(b)] \emph{The first-passage depth of tuned data.} For every $\hat w\in W$ and every $n$, the
bare datum $x_0=a_n-\hat w$ is a first-passage datum of clause~(0) of
Theorem~\ref{3.1:thm-main} at a depth $K_0$
with
\begin{equation*}
|K_0-n|\le j_U-1,\qquad j_U-1=\bigl\lceil(\varrho_2+c')/\lambda_U\bigr\rceil,
\end{equation*}
where $c'=c'(g)$ is the constant of Definition~\ref{9.8:def-blocking-factors}.
\end{itemize}
Write $X=g^{-2}$, let $\Psi$ be a pair of class $C^2$ in $\mathfrak{T}(1;\vartheta)$, $\Psi_r$ its
dilate of scale $r$, and $\|\Psi\|^2$ the product of the sup norms of its two members. Put
\begin{align*}
\mathfrak{m}(r)&:=\tfrac29\,b_0\,c_{\mathcal K}
\bigl(X+B^{\sharp}\bigl(s(r\mathsf{F}^{-(j_U-1)})+1\bigr)\bigr)^{-2},\\
\mathfrak{M}(r)&:=\tfrac83\,b_0\,C_{\mathcal K}
\bigl(X+\lambda^{\sharp}s(r\mathsf{F}^{j_U-1})-c_5\bigr)^{-2},\\
r^{*}&:=\mathsf{F}^{-(j_U-1)}\,\mathsf{d}_T^{\star}(g),
\end{align*}
the last being a ceiling on the separation. Then for all $\hat w\in W$, all $m_0\ge m_{\natural}$,
all $r\le r^{*}$, and all $n$ beyond a floor $n^{*}(r)$,
\begin{equation}\label{9.8:eq-tuned-two-sided-a}
\mathfrak{m}(r)\,\|\Psi\|^2\le S^n_{m_0}(a_n-\hat w;\Psi_r)\le\mathfrak{M}(r)\,\|\Psi\|^2,
\end{equation}
and hence the same bounds hold for $S^{(\hat w)}_{m_0}(\Psi_r)$. Moreover $\mathfrak{M}(r)\downarrow0$
and $r^2/\mathfrak{m}(r)\to0$ as $r\downarrow0$.
\end{lemma}

\begin{proof}
Fix $\hat w\in W$, $m_0\ge m_{\natural}$, $r\le r^{*}$ and an integer $n$; put
$x_0:=a_n-\hat w$. By
hypothesis~(b), $x_0$ is a first-passage datum of clause~(0) of Theorem~\ref{3.1:thm-main} at a
depth $K_0$ with
$|K_0-n|\le j_U-1$.

The function $S^n_{m_0}(x_0;\cdot)$ is read in units in which the bare lattice lies $n$ levels below
the reference scale. In these units $\Psi_r$ has scale $r$, so its scale in bare lattice spacings is
$r\mathsf{F}^n$. Read in the units in which the bare lattice lies $K_0$ levels below the reference
scale, which are the units of hypothesis~(a) for the datum $x_0$, the same pair has a scale
$\mathsf{d}$ determined by $r\mathsf{F}^n=\mathsf{d}\,\mathsf{F}^{K_0}$, that is
$\mathsf{d}=r\mathsf{F}^{\,n-K_0}$. Call $(f_1,f_2)$ this re-read pair. Since $\Psi_r$ lies in
$\mathfrak{T}(r;\vartheta)$ for every $r$, the pair $(f_1,f_2)$ lies in
$\mathfrak{T}(\mathsf{d};\vartheta)$. A change of units does not change sup norms, so
$\|f_1\|_\infty\|f_2\|_\infty=\|\Psi\|^2$, and $S^n_{m_0}(x_0;\Psi_r)=S^T_{2;K_0,m}(f_1,f_2)$ for
the torus exponent $m=m_0+n-K_0$ of $\mathbb{T}_{m_0}$ in the units of depth $K_0$, obtained by the
same change of units. From $|K_0-n|\le j_U-1$,
\begin{equation}\label{9.8:eq-tuned-two-sided-1}
r\mathsf{F}^{-(j_U-1)}\le\mathsf{d}\le r\mathsf{F}^{\,j_U-1}.
\end{equation}

The function $s(\cdot)$ is non-increasing. Indeed, if $\mathsf{d}\le\mathsf{d}'$, every $s$ with
$C_W(s)\mathsf{F}^{-s}\le\mathsf{d}$ also satisfies $C_W(s)\mathsf{F}^{-s}\le\mathsf{d}'$, so the
minimum defining $s(\mathsf{d}')$ is taken over a larger set. Applied to
\eqref{9.8:eq-tuned-two-sided-1}, this gives
\begin{equation}\label{9.8:eq-tuned-two-sided-2}
s(r\mathsf{F}^{\,j_U-1})\le s(\mathsf{d})\le s(r\mathsf{F}^{-(j_U-1)}).
\end{equation}

We now check the hypotheses of~(a) for $g_0$ with $g_0^{-2}=x_0$, $K=K_0$, the scale $\mathsf{d}$
and the pair $(f_1,f_2)$.

Hypothesis~(a) holds for every torus exponent $m\ge m_{\mathbb{T}}(g_-(g))$, with bounds that do not
depend on $m$, so the re-read exponent $m=m_0+n-K_0\ge m_0-(j_U-1)$, for $m_0\ge m_{\natural}$,
imposes no further condition.

The depth condition of~(a) is implied by $\mathsf{d}<\mathsf{d}_T^{\star}(g)$. This in turn is
implied by $r\le r^{*}$, because \eqref{9.8:eq-tuned-two-sided-1} and the definition of $r^{*}$ give
\begin{equation*}
\mathsf{d}\le r\mathsf{F}^{\,j_U-1}\le r^{*}\mathsf{F}^{\,j_U-1}=\mathsf{d}_T^{\star}(g).
\end{equation*}
The ceiling $\mathsf{d}_T^{\star}(g)$ depends on $g$ only. Hence $r^{*}$ does not depend on $n$,
$m_0$, $\hat w$ or the position of the pair.

For the ultraviolet floor we put
\begin{equation*}
n^{*}(r):=s(r\mathsf{F}^{-(j_U-1)})+k_W+j_U.
\end{equation*}
If $n\ge n^{*}(r)$, then $K_0\ge n-(j_U-1)$ and \eqref{9.8:eq-tuned-two-sided-2} give
\begin{equation*}
K_0-s(\mathsf{d})\ \ge\ n-(j_U-1)-s(r\mathsf{F}^{-(j_U-1)})\ \ge\ k_W+1\ \ge\ k_W.
\end{equation*}

Hypothesis~(a) therefore applies. It yields the displayed two-sided bound with
$\|f_1\|_\infty\|f_2\|_\infty=\|\Psi\|^2$, $g^{-2}=X$ and $s=s(\mathsf{d})$. The base
$X+B^{\sharp}(s+1)$ of the lower bound is non-decreasing in $s$, the coefficient $B^{\sharp}$ being
non-negative. By the right inequality in \eqref{9.8:eq-tuned-two-sided-2}, replacing $s(\mathsf{d})$
by $s(r\mathsf{F}^{-(j_U-1)})$ decreases the lower bound to $\mathfrak{m}(r)\|\Psi\|^2$. The base
$X+\lambda^{\sharp}s-c_5$ of the upper bound is increasing in $s$, since $\lambda^{\sharp}>0$. By the
left inequality in \eqref{9.8:eq-tuned-two-sided-2}, replacing $s(\mathsf{d})$ by
$s(r\mathsf{F}^{j_U-1})$ increases the upper bound to $\mathfrak{M}(r)\|\Psi\|^2$. This proves
\eqref{9.8:eq-tuned-two-sided-a} for every $n\ge n^{*}(r)$. The quantities $\mathfrak{m}(r)$ and
$\mathfrak{M}(r)$ do not depend on $n$, since $c_{\mathcal K}$ and $C_{\mathcal K}$ depend on
$\vartheta$ only and $X$, $B^{\sharp}$, $\lambda^{\sharp}$, $c_5$, $j_U$ are fixed. Both
inequalities in \eqref{9.8:eq-tuned-two-sided-a} are closed and
have bounds independent of $n$. They therefore pass to the limit $n\to\infty$ defining
$S^{(\hat w)}_{m_0}(\Psi_r)$.

We turn to the behaviour as $r\downarrow0$. Since $C_W\ge1$, every $s$ with
$C_W(s)\mathsf{F}^{-s}\le\mathsf{d}$ satisfies $\mathsf{F}^{-s}\le\mathsf{d}$, hence
\begin{equation*}
s(\mathsf{d})\ge\log_{\mathsf{F}}(1/\mathsf{d}),
\qquad
s(r\mathsf{F}^{\,j_U-1})\ge\log_{\mathsf{F}}(1/r)-(j_U-1).
\end{equation*}
So $s(r\mathsf{F}^{j_U-1})$, which is non-decreasing as $r$ decreases because $s(\cdot)$ is
non-increasing, tends to $+\infty$ as $r\downarrow0$. Since $\lambda^{\sharp}>0$, the base of
$\mathfrak{M}(r)$ grows without bound, and $\mathfrak{M}(r)\downarrow0$.

Since $C_W$ is polylogarithmic in $\bar x_s$, the minimum defining $s(\mathsf{d})$ satisfies
\begin{equation*}
s(\mathsf{d})\le\log_{\mathsf{F}}(1/\mathsf{d})+O\bigl(\log\log(1/\mathsf{d})\bigr).
\end{equation*}
With $\mathsf{d}=r\mathsf{F}^{-(j_U-1)}$, the base $X+B^{\sharp}(s(r\mathsf{F}^{-(j_U-1)})+1)$ of
$\mathfrak{m}(r)$ is at most a constant multiple of $\log(1/r)$ for small $r$. Hence
$\mathfrak{m}(r)\ge c\,(\log 1/r)^{-2}$ with $c>0$ independent of $r$, and
\begin{equation*}
\frac{r^2}{\mathfrak{m}(r)}\le c^{-1}\,r^2\,(\log 1/r)^2\longrightarrow0
\qquad(r\downarrow0).
\end{equation*}
\end{proof}

\begin{lemma}[One-step dilation covariance at a cluster value of the increments]
\label{9.8:lem-dilation-covariance}
Fix a frame $\mathsf{F}\in\{L,L'\}$ as in Definition~\ref{9.8:def-blocking-factors}, and take every
quantity of Theorem~\ref{3.1:thm-main} at blocking factor $\mathsf{F}$. Let $a_n$ be the pins and
$\alpha_n=a_{n+1}-a_n$ the tuned increments of Definition~\ref{9.8:def-blocking-factors}, and let
$W^{\circ}$ be the interior of the offset window $W=[-\varrho_2/2,\varrho_2/2]$. Assume the following
statements on every torus $\mathbb{T}_{m}$ with $m\ge m_{\natural}$:
\begin{itemize}
\item[(a)] the change-of-units identity, clause (v-d)(iii) of Theorem~\ref{3.1:thm-main}, with $s=1$;
\item[(b)] clauses (ii-$\mathbb{T}$-a) and (v-b) of Theorem~\ref{3.1:thm-main};
\item[(c)] clause (v-c)(1) of Theorem~\ref{3.1:thm-main}.
\end{itemize}
Let $\alpha$ be a cluster value of the sequence $(\alpha_n)$. Then for all $m_0\ge m_{\natural}$,
all $\mathsf{p}\ge2$, all separated Lipschitz $\vec{\varphi}=(\varphi_1,\dots,\varphi_{\mathsf{p}})$ on
$\mathbb{T}_{m_0}$, and all $\hat w$ with $\hat w\in W^{\circ}$ and $\hat w-\alpha\in W^{\circ}$,
\begin{equation}\label{9.8:eq-dilation-covariance-a}
S^{(\hat w)}_{m_0}(\vec{\varphi})=S^{(\hat w-\alpha)}_{m_0+1}\bigl(\vec{\varphi}^{\,\mathsf{F}}\bigr),
\end{equation}
where $\vec{\varphi}^{\,\mathsf{F}}$ is the dilate, with components
$\varphi_i^{\mathsf{F}}(y)=\varphi_i(y/\mathsf{F})$, a family of functions on $\mathbb{T}_{m_0+1}$;
it is again separated and Lipschitz, with separation multiplied by $\mathsf{F}$ and Lipschitz
constants divided by $\mathsf{F}$.
\end{lemma}

\begin{proof}
Fix $m_0$, $\mathsf{p}$, $\vec{\varphi}$ and $\hat w$ as in the statement.

\emph{An identity at every bare inverse coupling.} The torus $\mathbb{T}_{m_0}$ has side
$2\mathsf{F}^{m_0}\ell_g$, where $\ell_g$ denotes the side multiplier, depending on the reference
coupling $g$, with which the tori of clauses (ii-$\mathbb{T}$-a) and (v-b) of
Theorem~\ref{3.1:thm-main} are taken. So its bare lattice lying $n+1$ levels below the reference
scale has
$2\mathsf{F}^{m_0}\ell_g\mathsf{F}^{n+1}$ sites per side; this is also the number of sites per side of
the bare lattice of $\mathbb{T}_{m_0+1}$ lying $n$ levels below the reference scale. The two namings
therefore describe one lattice torus, one Wilson action at bare inverse coupling $x$, and one raw
integral. The observable $\mathcal{O}(\varphi_i)$ of Definition~\ref{1.3:def-curvature-field}
carries no power of
the lattice spacing, and in lattice units the weight $y\mapsto\varphi_i(y/\mathsf{F}^{n+1})$ that
$\mathcal{O}(\varphi_i)$ gives on $\mathbb{T}_{m_0}$ at depth $n+1$ equals
$\varphi_i^{\mathsf{F}}(y/\mathsf{F}^{n})$, the weight that $\mathcal{O}(\varphi_i^{\mathsf{F}})$ gives on
$\mathbb{T}_{m_0+1}$ at depth $n$. Hence the change-of-units identity, hypothesis (a) with $s=1$, reads
\begin{equation}\label{9.8:eq-dilation-covariance-1}
S^{n+1}_{m_0}(x;\vec{\varphi})=S^{n}_{m_0+1}\bigl(x;\vec{\varphi}^{\,\mathsf{F}}\bigr)
\qquad\text{for every } x \text{ and every } n.
\end{equation}

\emph{Choice of $x$.} Put $x:=a_{n+1}-\hat w$. Since $a_{n+1}=a_n+\alpha_n$, also
$x=a_n-(\hat w-\alpha_n)$, and \eqref{9.8:eq-dilation-covariance-1} becomes
\begin{equation}\label{9.8:eq-dilation-covariance-2}
S^{n+1}_{m_0}(a_{n+1}-\hat w;\vec{\varphi})
=S^{n}_{m_0+1}\bigl(a_n-(\hat w-\alpha_n);\vec{\varphi}^{\,\mathsf{F}}\bigr)
\qquad\text{for every } n.
\end{equation}
Since $\alpha$ is a cluster value of $(\alpha_n)$, there is a strictly increasing sequence $(n_j)$ with
$\alpha_{n_j}\to\alpha$.

\emph{The left side.} Since $\hat w\in W^{\circ}$, clause (ii-$\mathbb{T}$-a) of
Theorem~\ref{3.1:thm-main} on $\mathbb{T}_{m_0}$ (hypothesis (b)) gives
\begin{equation*}
S^{n}_{m_0}(a_n-\hat w;\vec{\varphi})\longrightarrow S^{(\hat w)}_{m_0}(\vec{\varphi})
\qquad(n\to\infty).
\end{equation*}
In
particular the left side of \eqref{9.8:eq-dilation-covariance-2}, taken at $n=n_j$, tends to
$S^{(\hat w)}_{m_0}(\vec{\varphi})$ as $j\to\infty$.

\emph{The right side.} Since $\hat w-\alpha\in W^{\circ}$, that is $|\hat w-\alpha|<\varrho_2/2$, and
$\alpha_{n_j}\to\alpha$, there is $j_0$ with $|\hat w-\alpha_{n_j}|\le\varrho_2/2$ for all $j\ge j_0$.
Define a real sequence by
\begin{equation*}
a'_n:=
\begin{cases}
a_n-(\hat w-\alpha_n), & n\in\{n_j : j\ge j_0\},\\
a_n, & \text{otherwise}.
\end{cases}
\end{equation*}
Then $|a'_n-a_n|\le\varrho_2/2$ for every $n$, so clause (v-c)(1) of Theorem~\ref{3.1:thm-main} on
$\mathbb{T}_{m_0+1}$ (hypothesis (c)) applies to the full sequence $(a'_n)$ and gives
\begin{equation*}
S^{n}_{m_0+1}\bigl(a'_n;\vec{\varphi}^{\,\mathsf{F}}\bigr)
-S^{(a_n-a'_n)}_{m_0+1}\bigl(\vec{\varphi}^{\,\mathsf{F}}\bigr)\longrightarrow0
\qquad(n\to\infty).
\end{equation*}
A sequence tending to zero does so along every subsequence; along $n=n_j$, $j\ge j_0$, we have
$a'_{n_j}=a_{n_j}-(\hat w-\alpha_{n_j})$ and $a_{n_j}-a'_{n_j}=\hat w-\alpha_{n_j}$, so
\begin{equation*}
S^{n_j}_{m_0+1}\bigl(a_{n_j}-(\hat w-\alpha_{n_j});\vec{\varphi}^{\,\mathsf{F}}\bigr)
-S^{(\hat w-\alpha_{n_j})}_{m_0+1}\bigl(\vec{\varphi}^{\,\mathsf{F}}\bigr)\longrightarrow0
\qquad(j\to\infty).
\end{equation*}
By clause (v-b) of Theorem~\ref{3.1:thm-main} on $\mathbb{T}_{m_0+1}$ (hypothesis (b)), the map
$\hat w'\mapsto S^{(\hat w')}_{m_0+1}(\vec{\varphi}^{\,\mathsf{F}})$ is continuous on $W$, and
$\hat w-\alpha_{n_j}\to\hat w-\alpha\in W^{\circ}$; hence
\begin{equation*}
S^{(\hat w-\alpha_{n_j})}_{m_0+1}\bigl(\vec{\varphi}^{\,\mathsf{F}}\bigr)\longrightarrow
S^{(\hat w-\alpha)}_{m_0+1}\bigl(\vec{\varphi}^{\,\mathsf{F}}\bigr)
\qquad(j\to\infty).
\end{equation*}
Together with the previous display, the
right side of \eqref{9.8:eq-dilation-covariance-2} at $n=n_j$ tends to
$S^{(\hat w-\alpha)}_{m_0+1}(\vec{\varphi}^{\,\mathsf{F}})$.

\emph{Conclusion.} The two sides of \eqref{9.8:eq-dilation-covariance-2} are equal for every $n$,
hence along $n=n_j$, and their limits along this subsequence are
$S^{(\hat w)}_{m_0}(\vec{\varphi})$ and $S^{(\hat w-\alpha)}_{m_0+1}(\vec{\varphi}^{\,\mathsf{F}})$
respectively. These limits therefore coincide, which is \eqref{9.8:eq-dilation-covariance-a}.
\end{proof}

\begin{lemma}[Convergence of the tuned increments]\label{9.8:lem-increment-convergence}
Fix one frame $\mathsf{F}\in\{L,L'\}$ of Definition~\ref{9.8:def-blocking-factors}, and let
$\alpha_n=a_{n+1}-a_n$ be the tuned increments of that frame. Assume the conditions
\eqref{9.8:eq-blocking-factors-b} and \eqref{9.8:eq-blocking-factors-c}, and assume the
following named statements:
\begin{itemize}
\item[(a)] on every torus $\mathbb{T}_{m_0}$ with $m_0\ge m_{\natural}$, clauses (v-a),
(ii-$\mathbb{T}$-a), (v-b) and (v-c)(1) of Theorem~\ref{3.1:thm-main}, read at blocking
factor $\mathsf{F}$;
\item[(b)] the change-of-units identity;
\item[(c)] clause (iv) of Theorem~\ref{3.1:thm-main}, whose
constants do not depend on the torus exponent $m_0\ge m_{\natural}$.
\end{itemize}
Then the sequence $(\alpha_n)$ converges, and
\begin{equation}\label{9.8:eq-increment-convergence-1}
\alpha_\infty:=\lim_{n\to\infty}\alpha_n\ \ge\ \lambda_U\ >\ 0 .
\end{equation}
The lemma gives no rate of convergence, and nothing on the intermediate couplings $g_{n-s}$;
it contains no infrared statement: the volume enters only through the uniformity of clause (iv) of
Theorem~\ref{3.1:thm-main} in the torus exponent.
\end{lemma}

\begin{proof}
Throughout, $W=[-\varrho_2/2,\varrho_2/2]$ and $W^{\circ}$ is its interior, as in
Definition~\ref{9.8:def-blocking-factors}; for a tuple $\vec\varphi$ and $\rho>0$,
$\vec{\varphi}^{\,\rho}$ denotes the sup-normalised dilate $\vec\varphi(\cdot/\rho)$, so
that $(\vec{\varphi}^{\,\rho})^{\rho'}=\vec{\varphi}^{\,\rho\rho'}$. The sequence
$(\alpha_n)$ is bounded: $|\alpha_n|\le\bar{\alpha}$ for all $n$, where
$\bar{\alpha}:=3\lambda_U+c'^{\sharp}$ is the bound on the tuned increments entering
\eqref{9.8:eq-blocking-factors-b}, and $c'^{\sharp}$ is the supremum over $g\le\gamma_U$ of the
constant $c'=c'(g)$ of clause (v-a) of Theorem~\ref{3.1:thm-main}; under
\eqref{9.8:eq-blocking-factors-b} we have $\bar{\alpha}\le\varrho_2/4$. Hence every cluster
value $\alpha$ of $(\alpha_n)$ satisfies $|\alpha|\le\bar{\alpha}\le\varrho_2/4$.

Suppose, to reach a contradiction, that $(\alpha_n)$ has two cluster values
$\alpha^1<\alpha^2$, and put $\eta:=\alpha^2-\alpha^1$. Then $0<\eta\le2\bar{\alpha}$;
in particular $\bar{\alpha}\ge\eta/2>0$.

\emph{Step 1: exact periodicity in the offset.} Let $\hat w$ satisfy
$\hat w,\ \hat w-\alpha^1,\ \hat w-\alpha^2\in W^{\circ}$. The set of such $\hat w$ is
$W^{\circ}\cap(W^{\circ}+\alpha^1)\cap(W^{\circ}+\alpha^2)$, an intersection of three open
intervals of length $\varrho_2$ whose centres lie at distance at most $\bar{\alpha}\le\varrho_2/4$
from $0$; it is therefore an open interval containing $(-\varrho_2/4,\varrho_2/4)$, and in
particular it is non-empty. Fix $m\ge m_{\natural}+1$ and a separated Lipschitz tuple
$\vec\psi$ on $\mathbb{T}_m$. The map $\vec\psi\mapsto\vec{\psi}^{\,1/\mathsf{F}}$ is a
bijection of tuples on $\mathbb{T}_m$ onto tuples on $\mathbb{T}_{m-1}$, with inverse
$\vec\varphi\mapsto\vec{\varphi}^{\,\mathsf{F}}$. We apply the one-step dilation covariance,
Lemma~\ref{9.8:lem-dilation-covariance}, on $\mathbb{T}_{m-1}$ (admissible since
$m-1\ge m_{\natural}$) to the tuple
$\vec{\psi}^{\,1/\mathsf{F}}$, once with the cluster value $\alpha^1$ and once with $\alpha^2$;
its hypotheses $\hat w,\hat w-\alpha^i\in W^{\circ}$ hold by the choice of $\hat w$.
Since $(\vec{\psi}^{\,1/\mathsf{F}})^{\mathsf{F}}=\vec\psi$, this gives
\begin{equation}\label{9.8:eq-increment-convergence-2}
S^{(\hat w-\alpha^1)}_{m}(\vec\psi)
=S^{(\hat w)}_{m-1}\bigl(\vec{\psi}^{\,1/\mathsf{F}}\bigr)
=S^{(\hat w-\alpha^2)}_{m}(\vec\psi).
\end{equation}
The maps $u\mapsto S^{(u)}_m(\vec\psi)$ and $u\mapsto S^{(u+\eta)}_m(\vec\psi)$ are
real-analytic on the open interval $(-\varrho_2/2,\ \varrho_2/2-\eta)$, which is non-empty
because $\eta\le2\bar{\alpha}\le\varrho_2/2$; this interval is exactly the set of $u$ with
$u,u+\eta\in W^{\circ}$, on which the real-analyticity of $u\mapsto S^{(u)}_m$ on $W^{\circ}$
provided by the statements (a) applies. Writing $u=\hat w-\alpha^2$, so that
$u+\eta=\hat w-\alpha^1$, \eqref{9.8:eq-increment-convergence-2} says that the two maps agree
for $u$ in the non-empty open interval
$\{\hat w-\alpha^2:\ \hat w,\hat w-\alpha^1,\hat w-\alpha^2\in W^{\circ}\}$, which is contained
in $(-\varrho_2/2,\varrho_2/2-\eta)$. By the identity theorem for real-analytic functions on
an interval,
\begin{equation}\label{9.8:eq-increment-convergence-3}
S^{(u)}_m(\vec\psi)=S^{(u+\eta)}_m(\vec\psi)\qquad\text{whenever } u,\ u+\eta\in W^{\circ},
\end{equation}
for every $m\ge m_{\natural}+1$ and every separated Lipschitz tuple $\vec\psi$ on
$\mathbb{T}_m$.

\emph{Step 2: a chain that never leaves the window.} Put
\begin{equation*}
\mathsf{T}:=\{\hat w\in W^{\circ}:\ \hat w-\alpha^1\in W^{\circ}\}
=W^{\circ}\cap(W^{\circ}+\alpha^1),
\end{equation*}
whatever the sign of $\alpha^1$. It is an open interval
of length
\begin{equation*}
\varrho_2-|\alpha^1|\ \ge\ \varrho_2-\bar{\alpha}\ \ge\ 3\bar{\alpha}\ >\ 2\bar{\alpha}\ \ge\ \eta ,
\end{equation*}
where we used $\varrho_2\ge4\bar{\alpha}$ and $\bar{\alpha}>0$. Fix $m_0\ge m_{\natural}$, a
separated Lipschitz tuple $\vec\varphi$ on $\mathbb{T}_{m_0}$, and $\hat w_0\in\mathsf{T}$.
We define inductively $\hat w_1,\hat w_2,\ldots\in\mathsf{T}$ with
\begin{equation}\label{9.8:eq-increment-convergence-4}
S^{(\hat w_i)}_{m_0+i}\bigl(\vec{\varphi}^{\,\mathsf{F}^i}\bigr)
=S^{(\hat w_{i+1})}_{m_0+i+1}\bigl(\vec{\varphi}^{\,\mathsf{F}^{i+1}}\bigr).
\end{equation}
Given $\hat w_i\in\mathsf{T}$, both $\hat w_i$ and $\hat w_i-\alpha^1$ lie in $W^{\circ}$, and
$m_0+i\ge m_{\natural}$; Lemma~\ref{9.8:lem-dilation-covariance} with the cluster value
$\alpha^1$, applied to the tuple $\vec{\varphi}^{\,\mathsf{F}^i}$ on $\mathbb{T}_{m_0+i}$, gives
\begin{equation*}
S^{(\hat w_i)}_{m_0+i}\bigl(\vec{\varphi}^{\,\mathsf{F}^i}\bigr)
=S^{(\hat w_i-\alpha^1)}_{m_0+i+1}\bigl(\vec{\varphi}^{\,\mathsf{F}^{i+1}}\bigr),
\qquad \hat w_i-\alpha^1\in W^{\circ}.
\end{equation*}
Write $\mathsf{T}=(\tau_-,\tau_+)$ with $\tau_+-\tau_->\eta$. This interval contains a point of
every coset $u+\eta\mathbb{Z}$: the least point of $u+\eta\mathbb{Z}$ that exceeds $\tau_-$ is
at most $\tau_-+\eta<\tau_+$. Let $v$ be such a point for $u:=\hat w_i-\alpha^1$, and pass from
$u$ to $v$ by steps of $+\eta$ or of $-\eta$, one at a time. Every intermediate point lies
between $u$ and $v$, which both belong to the interval $W^{\circ}$ (recall
$\mathsf{T}\subset W^{\circ}$); hence every intermediate point lies in $W^{\circ}$, and each step
joins two points $u',u'+\eta$ of $W^{\circ}$. Since $m_0+i+1\ge m_{\natural}+1$, each step is an
instance of \eqref{9.8:eq-increment-convergence-3} on $\mathbb{T}_{m_0+i+1}$ for the tuple
$\vec{\varphi}^{\,\mathsf{F}^{i+1}}$, and therefore
\begin{equation*}
S^{(u)}_{m_0+i+1}\bigl(\vec{\varphi}^{\,\mathsf{F}^{i+1}}\bigr)
=S^{(v)}_{m_0+i+1}\bigl(\vec{\varphi}^{\,\mathsf{F}^{i+1}}\bigr).
\end{equation*}
No starting point has to be excluded.
Setting $\hat w_{i+1}:=v\in\mathsf{T}$ gives \eqref{9.8:eq-increment-convergence-4}.
Composing \eqref{9.8:eq-increment-convergence-4} for $i=0,\ldots,k-1$ we obtain, exactly,
for every $k\ge1$,
\begin{equation}\label{9.8:eq-increment-convergence-5}
S^{(\hat w_0)}_{m_0}(\vec\varphi)=S^{(\hat w_k)}_{m_0+k}\bigl(\vec{\varphi}^{\,\mathsf{F}^k}\bigr),
\qquad \hat w_k\in\mathsf{T}\subset W .
\end{equation}

\emph{Step 3: clause (iv), uniformly in the volume, excludes this.} Let $\Psi$, $\Psi_r$,
$\|\Psi\|^2$, $r^{*}$, $\mathfrak{m}(r)$ and $\mathfrak{M}(r)$ be as in
Lemma~\ref{9.8:lem-tuned-two-sided}, the two-sided two-point bounds at tuned coupling. Take
$r_1\le r^{*}$ and, for
$k\ge1$, apply \eqref{9.8:eq-increment-convergence-5} to the pair
$\vec\varphi:=\Psi_{r_1\mathsf{F}^{-k}}$, for which
$\vec{\varphi}^{\,\mathsf{F}^k}=\Psi_{r_1}$. Since $r_1\mathsf{F}^{-k}\le r_1\le r^{*}$,
$\hat w_0\in W$ and $m_0\ge m_{\natural}$, the upper bound of
Lemma~\ref{9.8:lem-tuned-two-sided} gives for the left member of
\eqref{9.8:eq-increment-convergence-5}
\begin{equation*}
S^{(\hat w_0)}_{m_0}\bigl(\Psi_{r_1\mathsf{F}^{-k}}\bigr)\le
\mathfrak{M}\bigl(r_1\mathsf{F}^{-k}\bigr)\|\Psi\|^2\ \longrightarrow\ 0
\qquad(k\to\infty),
\end{equation*}
because $\mathfrak{M}(r)\downarrow0$ as $r\downarrow0$. For the right member,
$m_0+k\ge m_{\natural}$ and $\hat w_k\in W$, and no constant of
Lemma~\ref{9.8:lem-tuned-two-sided} depends on the torus exponent or on the point of $W$; its
lower bound gives, for every $k$,
\begin{equation*}
S^{(\hat w_k)}_{m_0+k}(\Psi_{r_1})\ \ge\ \mathfrak{m}(r_1)\|\Psi\|^2\ >\ 0 .
\end{equation*}
Choosing $k$ so large that $\mathfrak{M}(r_1\mathsf{F}^{-k})<\mathfrak{m}(r_1)$ contradicts
the equality \eqref{9.8:eq-increment-convergence-5}. Hence $(\alpha_n)$ has exactly one
cluster value.

\emph{Conclusion.} A bounded real sequence with exactly one cluster value converges; let
$\alpha_\infty$ be the limit of $(\alpha_n)$. By clause (v-a) of
Theorem~\ref{3.1:thm-main}, among the statements (a),
\begin{equation*}
\frac{a_n-a_0}{n}\ \ge\ \lambda_U-\frac{c'}{n}\qquad(n\ge1),
\end{equation*}
where $c'$ is the constant of clause (v-a) introduced above.
The left side equals $\frac1n\sum_{j=0}^{n-1}\alpha_j$, the Ces\`aro mean of $(\alpha_j)$,
which converges to $\alpha_\infty$ since $\alpha_j\to\alpha_\infty$; the right side tends to
$\lambda_U$. Therefore $\alpha_\infty\ge\lambda_U>0$, which is
\eqref{9.8:eq-increment-convergence-1}.
\end{proof}

\begin{corollary}[The limit increment is the one-loop slope]\label{9.8:cor-one-loop-slope}
Let $\mathsf{F}$ be one of the two blocking factors $L$, $L'$ with which Ba{\l}aban's step is run in
this chapter, and write
\begin{equation}\label{9.8:eq-one-loop-slope-1}
c_0(\mathsf{F})=\frac{11N^2}{12\pi^2}\,\log\mathsf{F}.
\end{equation}
In the frame $\mathsf{F}$ let $x_k$ denote the running inverse couplings of Ba{\l}aban's step with
blocking factor $\mathsf{F}$, with $x_0$ the bare inverse coupling, let $a_n$ be the tuned data and
$\alpha_n=a_{n+1}-a_n$ the tuned increments. Assume:
\begin{itemize}
\item[(a)] the hypotheses of Lemma~\ref{9.8:lem-increment-convergence}, so that $\alpha_n$
converges to $\alpha_\infty^{\mathsf{F}}$;
\item[(b)] clause \textup{(iii-d)} of Theorem~0 for the blocking factor $\mathsf{F}$, in its
Ces\`aro form: for the runs of Ba{\l}aban's step with blocking factor $\mathsf{F}$, with $K$ the
level of first passage, $(x_{K-s}-x_K)/s\to c_0(\mathsf{F})$ as $s\to\infty$, $s\le K$, read at
$s=K$ along these runs, that is, $(x_0-x_K)/K\to c_0(\mathsf{F})$ as $K\to\infty$;
\item[(c)] the first-passage depth statement for tuned data: for every $n$, the run started at the
bare inverse coupling $x_0=a_n$ has first passage at a level $K_0=K_0(n)$ with $|K_0-n|\le j_U-1$,
where $j_U-1=\lceil(\varrho_2+c')/\lambda_U\rceil$ is a fixed integer independent of $n$, and
$x_{K_0}\in(X,X+B^{\sharp}]$.
\end{itemize}
Then $\alpha_\infty^{\mathsf{F}}=c_0(\mathsf{F})$.
\end{corollary}

\begin{proof}
Fix $n$ and put $x_0=a_n$. By hypothesis~(c), $x_0$ is a first-passage datum of depth $K_0$ with
$|K_0-n|\le j_U-1$ and $x_{K_0}\in(X,X+B^{\sharp}]$. Hypothesis~(b), applied at $s=K_0$ with
$K=K_0$, gives
\begin{equation}\label{9.8:eq-one-loop-slope-2}
\frac{x_0-x_{K_0}}{K_0}\longrightarrow c_0(\mathsf{F})\qquad(K_0\to\infty).
\end{equation}
Since $|K_0-n|\le j_U-1$, we have $K_0\to\infty$ and $K_0/n\to1$ as $n\to\infty$. Since
$X<x_{K_0}\le X+B^{\sharp}$, the numbers $x_{K_0}$ are bounded uniformly in $n$, so
$x_{K_0}/n\to0$. Writing
\begin{equation*}
\frac{a_n}{n}=\frac{K_0}{n}\cdot\frac{x_0-x_{K_0}}{K_0}+\frac{x_{K_0}}{n},
\end{equation*}
we obtain from \eqref{9.8:eq-one-loop-slope-2} that $a_n/n\to c_0(\mathsf{F})$.

On the other hand $a_n-a_0=\sum_{j=0}^{n-1}\alpha_j$, and by hypothesis~(a) the increments
$\alpha_j$ converge to $\alpha_\infty^{\mathsf{F}}$. A convergent sequence has its Ces\`aro means
converging to the same limit, so $(a_n-a_0)/n\to\alpha_\infty^{\mathsf{F}}$; since $a_0/n\to0$,
also $a_n/n\to\alpha_\infty^{\mathsf{F}}$. Comparing the two limits of $a_n/n$ gives
$\alpha_\infty^{\mathsf{F}}=c_0(\mathsf{F})$, which with \eqref{9.8:eq-one-loop-slope-1} is the
assertion.
\end{proof}

\begin{remark}
This corollary is not used in part~(b) of the universality theorem of this chapter, which gives a
separate derivation of the universality of $\alpha_\infty^{\mathsf{F}}/\log\mathsf{F}$.
\end{remark}

\begin{corollary}[Exact covariance of the tuned family under blocking]\label{9.8:cor-exact-covariance}
Assume the hypotheses of Lemmas~\ref{9.8:lem-increment-convergence}
and~\ref{9.8:lem-dilation-covariance}, and let
$\alpha_\infty=\lim_{n\to\infty}\alpha_n$ be the limit of the tuned increments given by
Lemma~\ref{9.8:lem-increment-convergence}.
For $\rho>0$ write $D_\rho$ for the dilation acting on functionals of test functions,
$(D_\rho S)(\vec{\varphi}):=S(\vec{\varphi}^{\,\rho})$, where $\vec{\varphi}^{\,\rho}$ is the
sup-normalised dilate $\vec{\varphi}(\cdot/\rho)$. Let $m_0$, $\mathsf{p}$ and the separated
Lipschitz families $\vec{\varphi}$ on $\mathbb{T}_{m_0}$ range as in
Lemma~\ref{9.8:lem-dilation-covariance}. Then for every $\hat w$ with
$\hat w\in W^{\circ}$ and $\hat w-\alpha_\infty\in W^{\circ}$,
\begin{equation}\label{9.8:eq-exact-covariance-1}
S^{(\hat w)}_{m_0}=D_{\mathsf{F}}\,S^{(\hat w-\alpha_\infty)}_{m_0+1}.
\end{equation}
\end{corollary}

\begin{proof}
By Lemma~\ref{9.8:lem-increment-convergence}, the sequence $(\alpha_n)$ of tuned increments
converges, with limit $\alpha_\infty$. The limit of a convergent sequence is a cluster value of
it: every subsequence, in particular the full sequence $n_j=j$, satisfies
$\alpha_{n_j}\to\alpha_\infty$. Hence $\alpha=\alpha_\infty$ is an admissible choice in
Lemma~\ref{9.8:lem-dilation-covariance}, the one-step dilation covariance at a cluster value of
the increments.

Fix $m_0$, $\mathsf{p}$ and a separated Lipschitz family $\vec{\varphi}$ on $\mathbb{T}_{m_0}$ as
in that lemma, and fix $\hat w$ with $\hat w\in W^{\circ}$ and $\hat w-\alpha_\infty\in W^{\circ}$.
These are exactly the hypotheses of Lemma~\ref{9.8:lem-dilation-covariance} with
$\alpha=\alpha_\infty$, and its conclusion, the identity \eqref{9.8:eq-dilation-covariance-a},
reads
\begin{equation*}
S^{(\hat w)}_{m_0}(\vec{\varphi})=S^{(\hat w-\alpha_\infty)}_{m_0+1}(\vec{\varphi}^{\,\mathsf{F}}).
\end{equation*}
By the definition of $D_{\mathsf{F}}$, the right-hand side equals
$\bigl(D_{\mathsf{F}}\,S^{(\hat w-\alpha_\infty)}_{m_0+1}\bigr)(\vec{\varphi})$. Since
$\vec{\varphi}$ was an arbitrary family in the range of
Lemma~\ref{9.8:lem-dilation-covariance}, the two functionals
$S^{(\hat w)}_{m_0}$ and $D_{\mathsf{F}}\,S^{(\hat w-\alpha_\infty)}_{m_0+1}$ agree on that
whole range, which is \eqref{9.8:eq-exact-covariance-1}.

For comparison, the change-of-units identity, as used in clause (v-d)(iii) of
Theorem~\ref{3.1:thm-main}, identifies $(m_0,\hat x)$ with $(m_0+s,\hat x_s)$, where
$\hat x_s$ denotes the comparison point that this clause pairs with $\hat x$ at torus exponent
$m_0+s$, and
\begin{equation*}
\hat x_s-\hat x\in\bigl[\lambda^{\sharp}s-c_5,\;B^{\sharp}s\bigr].
\end{equation*}
The identity \eqref{9.8:eq-exact-covariance-1} sharpens this for the tuned family: indexed by
the bare offset $\hat w$, one step of blocking shifts the index by the single number
$\alpha_\infty$, in place of an interval of admissible shifts.
\end{proof}

\begin{lemma}[Comparison of two blocking factors on one lattice]\label{9.8:lem-two-factor-comparison}
Let $L$ and $L'$ be the two blocking factors of Definition~\ref{9.8:def-blocking-factors}, and
let $W$ be the window fixed there. As in that definition, every constant and exponent is read at
the blocking factor $\mathsf{F}\in\{L,L'\}$ and carries the index $\mathsf{F}$ only where both
occur. In particular $g_{\mathsf{F}}$ is the reference coupling read at $\mathsf{F}$, and
$\underline{s}'$, $\bar{s}'$, $m_{\natural}$, $\ell_{\mathsf{p}}$ and
$C^{\star\mathrm{cvx}}_{T,\mathsf{p}}$ are the depth exponents, the torus-exponent floor, the side
multipliers and the constants of the short-distance corollary, read at $\mathsf{F}$. We write
$m_{\natural}(\mathsf{F}):=m_{\natural}(g_{\mathsf{F}})$. For a lattice
torus $\mathbb{T}$ and lattice weights $\vec f=(f_1,\dots,f_{\mathsf{p}})$ on its bare sites we
write $S[x_0;\mathbb{T}](\vec f)$ for the connected $\mathsf{p}$-point function of the curvature
observables with weights $f_i$ under the bare Wilson measure on $\mathbb{T}$ at the bare datum
$x_0$.

Put
\begin{equation*}
C^{\flat}_{\mathsf{F},\mathsf{p}}
:=C^{\star\mathrm{cvx}}_{T,\mathsf{p}}(\mathsf{F})\,
\mathsf{F}^{-(4-\mathsf{p})\,\underline{s}'(g_{\mathsf{F}})} .
\end{equation*}
Assume the following, for each $\mathsf{F}\in\{L,L'\}$.
\begin{itemize}
\item[(g)] Clause (g) of the short-distance corollary, in its finite-cutoff form. Let
$\mathsf{p}\in\{2,3\}$. Call a rectangular lattice torus \emph{admissible for $\mathsf{F}$ and
$\mathsf{p}$} if its sides, in $\mathsf{F}$-reference units, lie in
$2\ell_{\mathsf{p}}(\mathsf{F})\,\mathsf{F}^{\bar{s}'(g_{\mathsf{F}})}\mathbb{N}$ and are
$\ge 2\mathsf{F}^{m_{\natural}(\mathsf{F})}$. Let a bare datum $x_0$ have a representation
$x_0=a_{\nu}^{\mathsf{F}}-\hat v$ with $\nu\in\mathbb{N}$ and $\hat v\in W$; let $r>0$ and
$\vec\psi\in\mathfrak{P}_{\mathsf{p}}(r;\vartheta)$ with
\begin{equation*}
\mathsf{F}^{-\nu}\le r\le \mathsf{F}^{\underline{s}'(g_{\mathsf{F}})}/40 ,
\end{equation*}
and put $f_i(x):=\psi_i(x/\mathsf{F}^{\nu})$ on bare sites $x$. Then for every torus
$\mathbb{T}$ admissible for $\mathsf{F}$ and $\mathsf{p}$
\begin{equation}\label{9.8:eq-two-factor-comparison-4}
\begin{split}
\bigl|S[x_0;\mathbb{T}](\vec f)-S[x_0;\mathbb{T}^{\mathsf{F}}_{m_{\natural}(\mathsf{F})}](\vec f)\bigr|
&\le C^{\star\mathrm{cvx}}_{T,\mathsf{p}}(\mathsf{F})\,
\bigl(r\,\mathsf{F}^{-\underline{s}'(g_{\mathsf{F}})}\bigr)^{4-\mathsf{p}}
\prod_{i=1}^{\mathsf{p}}\|\psi_i\|_\infty\\
&=C^{\flat}_{\mathsf{F},\mathsf{p}}\,r^{4-\mathsf{p}}\prod_{i=1}^{\mathsf{p}}\|\psi_i\|_\infty ,
\end{split}
\end{equation}
the equality being the definition of $C^{\flat}_{\mathsf{F},\mathsf{p}}$.
\end{itemize}
Let a bare datum $x_0$ have two representations
\begin{equation*}
x_0=a_n^{L}-\hat w=a_{n'}^{L'}-\hat w',\qquad \hat w,\hat w'\in W,
\end{equation*}
and put $\lambda:=L^{n}/L'^{\,n'}$. Let $\mathsf{p}\in\{2,3\}$ and
$\vec\varphi\in\mathfrak{P}_{\mathsf{p}}(r;\vartheta)$, and assume the ceilings
\begin{equation*}
r\le L^{\underline{s}'(g_L)}/40,\qquad \lambda r\le L'^{\,\underline{s}'(g_{L'})}/40,
\end{equation*}
and the floors
\begin{equation*}
r\ge L^{-n},\qquad \lambda r\ge L'^{\,-n'} .
\end{equation*}
Then
\begin{equation}\label{9.8:eq-two-factor-comparison-a}
\begin{split}
&\bigl|S^{n}_{L,m_{\natural}(L)}(x_0;\vec\varphi)
-S^{n'}_{L',m_{\natural}(L')}(x_0;\vec\varphi^{\,\lambda})\bigr|\\
&\qquad\le\Bigl[C^{\flat}_{L,\mathsf{p}}\,r^{4-\mathsf{p}}
+C^{\flat}_{L',\mathsf{p}}\,(\lambda r)^{4-\mathsf{p}}\Bigr]
\prod_{i=1}^{\mathsf{p}}\|\varphi_i\|_\infty .
\end{split}
\end{equation}
\end{lemma}

\begin{proof}
\emph{One set of lattice weights.} For $i=1,\dots,\mathsf{p}$ define the lattice weight
$f_i(x):=\varphi_i(x/L^{n})$ on bare sites $x$. Since $\varphi_i^{\lambda}=\varphi_i(\cdot/\lambda)$
and $\lambda L'^{\,n'}=L^{n}$,
\begin{equation}\label{9.8:eq-two-factor-comparison-1}
\varphi_i^{\lambda}\bigl(x/L'^{\,n'}\bigr)=\varphi_i\bigl(x/(\lambda L'^{\,n'})\bigr)
=\varphi_i\bigl(x/L^{n}\bigr)=f_i(x).
\end{equation}
The function $S^{n}_{L,m_{\natural}(L)}(x_0;\vec\varphi)$ is the connected $\mathsf{p}$-point
function, at the bare datum $x_0$, on the torus $\mathbb{T}^{L}_{m_{\natural}(L)}$, of the
curvature observables whose lattice weights are $\varphi_i(x/L^{n})$. Likewise
$S^{n'}_{L',m_{\natural}(L')}(x_0;\vec\varphi^{\,\lambda})$ is the same function on
$\mathbb{T}^{L'}_{m_{\natural}(L')}$ with weights $\varphi_i^{\lambda}(x/L'^{\,n'})$. By
\eqref{9.8:eq-two-factor-comparison-1} the two sets of weights coincide:
\begin{align*}
S^{n}_{L,m_{\natural}(L)}(x_0;\vec\varphi)
&=S[x_0;\mathbb{T}^{L}_{m_{\natural}(L)}](\vec f),\\
S^{n'}_{L',m_{\natural}(L')}(x_0;\vec\varphi^{\,\lambda})
&=S[x_0;\mathbb{T}^{L'}_{m_{\natural}(L')}](\vec f).
\end{align*}
Moreover $\|\varphi_i^{\lambda}\|_\infty=\|\varphi_i\|_\infty$, because a dilate of the argument
does not change the supremum.

\emph{A common torus.} Let $\mathbb{T}^{\natural}$ be the cubic lattice torus with
$N^{\natural}$ sites on each side, where
\begin{equation*}
N^{\natural}:=2\cdot\bigl[\ell_3(L)\,L^{m_{\natural}(L)+n}\bigr]\cdot
\bigl[\ell_3(L')\,L'^{\,m_{\natural}(L')+n'}\bigr].
\end{equation*}
In $L$-reference units the bare spacing is $L^{-n}$, so the side of $\mathbb{T}^{\natural}$
measures
\begin{equation*}
N^{\natural}/L^{n}=2\,\ell_3(L)\,L^{m_{\natural}(L)}\cdot
\bigl[\ell_3(L')\,L'^{\,m_{\natural}(L')+n'}\bigr].
\end{equation*}
The second bracket is a positive integer. The side multiplier $\ell_2(L)$ divides $\ell_3(L)$, so
$\ell_{\mathsf{p}}(L)$ divides $\ell_3(L)$ for $\mathsf{p}\in\{2,3\}$. Also
$m_{\natural}(L)\ge\bar{s}'(g_L)$, so $L^{\bar{s}'(g_L)}$ divides $L^{m_{\natural}(L)}$. The
side therefore lies in $2\ell_{\mathsf{p}}(L)L^{\bar{s}'(g_L)}\mathbb{N}$. Since $\ell_3(L)\ge1$ and
the second bracket is $\ge1$, the side is $\ge 2L^{m_{\natural}(L)}$. Exchanging the roles of the
two brackets gives the same two properties in $L'$-reference units, with $L'$, $n'$ and the
constants and exponents read at $L'$. Thus $\mathbb{T}^{\natural}$ is admissible for
$\mathsf{F}$ and $\mathsf{p}$ in hypothesis (g), both for $\mathsf{F}=L$ and for $\mathsf{F}=L'$.

\emph{Two applications of (g).} Apply hypothesis (g) at $\mathsf{F}=L$, in its finite-cutoff
form, with the representation $x_0=a_n^{L}-\hat w$ (so $\nu=n$ and $\hat v=\hat w$), the radius
$r$, the test tuple $\vec\varphi$ and the torus $\mathbb{T}^{\natural}$. The floor $r\ge L^{-n}$
and the ceiling $r\le L^{\underline{s}'(g_L)}/40$ are among the hypotheses of the lemma, and the
weights $\varphi_i(x/L^{n})$ are the weights $f_i$ of the first step. The finite-cutoff form of
(g) is the inequality at finite cutoff from which clause (g) is proved through the volume-transfer
corollary. On $\mathbb{T}^{\natural}$ only link
reflection positivity, the chessboard estimate, the comparison of tori and convexity of the bare
Wilson measure are used. Ba{\l}aban's construction runs on two cubic tori of the $L$-family. By
\eqref{9.8:eq-two-factor-comparison-4} this gives
\begin{equation}\label{9.8:eq-two-factor-comparison-2}
\bigl|S[x_0;\mathbb{T}^{\natural}](\vec f)-S[x_0;\mathbb{T}^{L}_{m_{\natural}(L)}](\vec f)\bigr|
\le C^{\flat}_{L,\mathsf{p}}\,r^{4-\mathsf{p}}\prod_{i=1}^{\mathsf{p}}\|\varphi_i\|_\infty .
\end{equation}
Apply hypothesis (g) at $\mathsf{F}=L'$ in the same way, with the representation
$x_0=a_{n'}^{L'}-\hat w'$ (so $\nu=n'$ and $\hat v=\hat w'$). The torus is again
$\mathbb{T}^{\natural}$, and the test tuple is $\vec\varphi^{\,\lambda}$ at radius $\lambda r$, for
which $\lambda r\le L'^{\,\underline{s}'(g_{L'})}/40$ and $\lambda r\ge L'^{\,-n'}$ are among the
hypotheses. By \eqref{9.8:eq-two-factor-comparison-1} its lattice weights
$\varphi_i^{\lambda}(x/L'^{\,n'})$ are again $\vec f$, and
$\|\varphi_i^{\lambda}\|_\infty=\|\varphi_i\|_\infty$. By
\eqref{9.8:eq-two-factor-comparison-4} this gives
\begin{equation}\label{9.8:eq-two-factor-comparison-3}
\bigl|S[x_0;\mathbb{T}^{\natural}](\vec f)-S[x_0;\mathbb{T}^{L'}_{m_{\natural}(L')}](\vec f)\bigr|
\le C^{\flat}_{L',\mathsf{p}}\,(\lambda r)^{4-\mathsf{p}}
\prod_{i=1}^{\mathsf{p}}\|\varphi_i\|_\infty .
\end{equation}

\emph{Conclusion.} Insert $S[x_0;\mathbb{T}^{\natural}](\vec f)$ between the two terms on the
left of \eqref{9.8:eq-two-factor-comparison-a}. Then apply the triangle inequality together with
the identifications of the first step and the bounds \eqref{9.8:eq-two-factor-comparison-2} and
\eqref{9.8:eq-two-factor-comparison-3}. This gives \eqref{9.8:eq-two-factor-comparison-a}.
\end{proof}

\begin{lemma}[Boundedness of the ratio of lattice scales]\label{9.8:lem-scale-ratio}
Let $L$ and $L'$ be the two blocking factors of
Lemma~\ref{9.8:lem-two-factor-comparison}. For a bare $x_0$, a pair of representations of
$x_0$ is a pair of identities $x_0=a^{L}_{n}-\hat w=a^{L'}_{n'}-\hat w'$ with
$\hat w,\hat w'\in W$. For such a pair put $\lambda:=L^{n}/L'^{\,n'}$. There are numbers
$0<\rho_-\le 1\le\rho_+<\infty$ and a floor $x_0^{*}$ such that
\begin{equation*}
\rho_-\le\lambda\le\rho_+
\end{equation*}
for all pairs of representations of all $x_0\ge x_0^{*}$.
\end{lemma}

\begin{proof}
A subscript $\mathsf{F}\in\{L,L'\}$ on $\mathfrak{m}$, $\mathfrak{M}$, $r^{*}$, $n^{*}(\cdot)$
and on the connected functions indicates that the object is formed in the frame
$\mathsf{F}$, that is, with blocking factor $\mathsf{F}$, its tuned sequence
$(a^{\mathsf{F}}_n)$ and its torus exponent $m_{\natural}$ (respectively $m_{\natural}'$).
We write $S^n_L(x_0;\cdot)$ for $S^n_{L,m_{\natural}}(x_0;\cdot)$ and
$S^{n'}_{L'}(x_0;\cdot)$ for $S^{n'}_{L',m_{\natural}'}(x_0;\cdot)$. Let $\Psi$, its dilates
$\Psi_r$ and $\|\Psi\|^2$ be as in Lemma~\ref{9.8:lem-tuned-two-sided}. We use
Lemma~\ref{9.8:lem-tuned-two-sided} (two-sided two-point bounds at tuned coupling) and
Lemma~\ref{9.8:lem-two-factor-comparison} (comparison of two blocking factors on one
lattice), each with torus exponent $m_0=m_{\natural}$ in the frame concerned.

\emph{The case $\lambda\ge1$.} Since $r^2/\mathfrak{m}_{L'}(r)\to0$ as $r\downarrow0$ by
Lemma~\ref{9.8:lem-tuned-two-sided}, applied in the frame $L'$, we
may choose $r_1'>0$ with
\begin{equation}\label{9.8:eq-scale-ratio-1}
r_1'\le\min\Bigl\{r^{*}_{L'},\,r^{*}_{L},\,\frac{L^{\underline{s}'}}{40},\,
\frac{L'^{\,\underline{s}'}}{40}\Bigr\},
\qquad
\mathfrak{m}_{L'}(r_1')\ge 2\bigl(C^{\flat}_{L,2}+C^{\flat}_{L',2}\bigr)r_1'^{\,2}.
\end{equation}
This choice does not depend on $x_0$ or on the pair of representations. Given a pair of
representations with
$\lambda\ge1$, put $r:=r_1'/\lambda\le r_1'$, so that $\lambda r=r_1'$. We assume for the
moment the floors
\begin{equation}\label{9.8:eq-scale-ratio-2}
r\ge L^{-n},\qquad \lambda r\ge L'^{\,-n'},\qquad n'\ge n^{*}_{L'}(r_1'),\qquad
n\ge n^{*}_{L}(r).
\end{equation}

Apply Lemma~\ref{9.8:lem-two-factor-comparison} with $\mathsf{p}=2$ and
$\vec\varphi=\Psi_r$. Its ceilings hold: $r\le r_1'\le L^{\underline{s}'}/40$ and
$\lambda r=r_1'\le L'^{\,\underline{s}'}/40$ by \eqref{9.8:eq-scale-ratio-1}. Its floors are
the first two inequalities of \eqref{9.8:eq-scale-ratio-2}. By the definition of the dilate,
$\Psi_r^{\,\lambda}=\Psi_{\lambda r}=\Psi_{r_1'}$. The product of the sup norms of the
members of $\Psi_r$ is $\|\Psi\|^2$. Inequality \eqref{9.8:eq-two-factor-comparison-a}
therefore gives
\begin{align*}
\bigl|S^n_L(x_0;\Psi_r)-S^{n'}_{L'}(x_0;\Psi_{r_1'})\bigr|
&\le\bigl(C^{\flat}_{L,2}r^2+C^{\flat}_{L',2}r_1'^{\,2}\bigr)\|\Psi\|^2\\
&\le\bigl(C^{\flat}_{L,2}+C^{\flat}_{L',2}\bigr)r_1'^{\,2}\|\Psi\|^2,
\end{align*}
where the second inequality uses $r\le r_1'$.

In the frame $L'$ the datum is $x_0=a^{L'}_{n'}-\hat w'$ with $\hat w'\in W$. We have
$r_1'\le r^{*}_{L'}$ by \eqref{9.8:eq-scale-ratio-1}, and $n'\ge n^{*}_{L'}(r_1')$ by
\eqref{9.8:eq-scale-ratio-2}. So the
lower half of \eqref{9.8:eq-tuned-two-sided-a} gives
$S^{n'}_{L'}(x_0;\Psi_{r_1'})\ge\mathfrak{m}_{L'}(r_1')\|\Psi\|^2$. Together with the last
display and the second condition in \eqref{9.8:eq-scale-ratio-1}, this yields
\begin{equation}\label{9.8:eq-scale-ratio-3}
S^n_L(x_0;\Psi_r)
\ge\Bigl(\mathfrak{m}_{L'}(r_1')-\bigl(C^{\flat}_{L,2}+C^{\flat}_{L',2}\bigr)r_1'^{\,2}\Bigr)
\|\Psi\|^2
\ge\tfrac12\,\mathfrak{m}_{L'}(r_1')\,\|\Psi\|^2 .
\end{equation}

In the frame $L$ the datum is $x_0=a^{L}_{n}-\hat w$ with $\hat w\in W$. We have
$r\le r_1'$, and $r_1'\le r^{*}_L$ by \eqref{9.8:eq-scale-ratio-1}; and
$n\ge n^{*}_L(r)$ by \eqref{9.8:eq-scale-ratio-2}.
So the upper half of \eqref{9.8:eq-tuned-two-sided-a} gives
$S^n_L(x_0;\Psi_r)\le\mathfrak{M}_L(r_1'/\lambda)\|\Psi\|^2$. Comparing with
\eqref{9.8:eq-scale-ratio-3} and dividing by $\|\Psi\|^2\neq0$, we obtain
\begin{equation}\label{9.8:eq-scale-ratio-4}
\mathfrak{M}_L(r_1'/\lambda)\ge\tfrac12\,\mathfrak{m}_{L'}(r_1').
\end{equation}
The number $\tfrac12\mathfrak{m}_{L'}(r_1')$ is positive. By
Lemma~\ref{9.8:lem-tuned-two-sided}, $\mathfrak{M}_L(r)\downarrow0$ as $r\downarrow0$, so
\begin{equation*}
\underline{r}:=\inf\bigl\{r>0:\ \mathfrak{M}_L(r)\ge\tfrac12\,\mathfrak{m}_{L'}(r_1')\bigr\}>0 .
\end{equation*}
The number $\underline{r}$ does not depend on $x_0$ or on the pair of representations. By
\eqref{9.8:eq-scale-ratio-4},
$r_1'/\lambda\ge\underline{r}$, that is, $\lambda\le r_1'/\underline{r}$. Put
\begin{equation*}
\rho_+:=\max\bigl\{1,\,r_1'/\underline{r}\bigr\}.
\end{equation*}
Then every pair of representations with $\lambda\ge1$ for which the floors
\eqref{9.8:eq-scale-ratio-2} hold satisfies $\lambda\le r_1'/\underline{r}\le\rho_+$.

\emph{The floors.} We show that \eqref{9.8:eq-scale-ratio-2} holds for all $x_0$ beyond a
floor. From $r=r_1'/\lambda$ and $\lambda=L^{n}/L'^{\,n'}$ we get
\begin{equation*}
rL^{n}=\lambda r\,L'^{\,n'}=r_1'L'^{\,n'} .
\end{equation*}
Hence the first two inequalities of \eqref{9.8:eq-scale-ratio-2} both read
$r_1'L'^{\,n'}\ge1$. This holds as soon as $n'\ge\log_{L'}(1/r_1')$. Since $r_1'$ is
fixed, the third inequality is a lower bound on $n'$ alone. For the fourth inequality,
$\lambda\ge1$ gives $n\ge n'\log_L L'$, and evaluating the depth floor $n^{*}_L(\cdot)$ of
Lemma~\ref{9.8:lem-tuned-two-sided} at $r=r_1'/\lambda$ one finds that $n-n^{*}_L(r)$ is
at least $n'\log_L L'$ minus a constant multiple of $\log n'$ and minus a further
constant, neither constant depending on $x_0$ or on the pair of representations.
Since $\hat w'$ ranges over the bounded window $W$ and $x_0=a^{L'}_{n'}-\hat w'$, for each
$n'_0$ the representations with $n'\le n'_0$ have $x_0$ in a bounded set; hence $n'\to\infty$
as $x_0\to\infty$. Therefore the lower bound just obtained for $n-n^{*}_L(r)$
tends to infinity as $x_0\to\infty$. So there is a floor $x_0^{*}$ such that all four
inequalities of \eqref{9.8:eq-scale-ratio-2} hold for every pair of representations, with
$\lambda\ge1$, of every $x_0\ge x_0^{*}$.

\emph{The case $\lambda\le1$.} Then $1/\lambda=L'^{\,n'}/L^{n}\ge1$ is the ratio formed with
the roles of the frames $L$ and $L'$ exchanged. The argument above, with $L$ and $L'$
exchanged throughout, gives a number $\rho_+'\ge1$ and a floor $x_0^{*\prime}$ such that
$1/\lambda\le\rho_+'$ for all pairs of representations of all $x_0\ge x_0^{*\prime}$ with
$\lambda\le1$. Put $\rho_-:=1/\rho_+'\in(0,1]$ and replace $x_0^{*}$ by
$\max\{x_0^{*},x_0^{*\prime}\}$. Then $\rho_-\le\lambda\le\rho_+$ for all pairs of
representations of all $x_0\ge x_0^{*}$.
\end{proof}

\begin{lemma}[No germ is invariant under a dilation]\label{9.8:lem-no-invariant-germ}
Let $S=S^{(u)}_{\mathsf{F},m}$ with $u\in W$, and let $\rho>1$, $C_A<\infty$ and $r_A>0$; here
$\rho$ is a dilation factor, and $\rho$, $C_A$, $r_A$ are local to this lemma. Let $\Psi$ be the fixed
pair of test functions of Lemma~\ref{9.8:lem-tuned-two-sided}, with dilates $\Psi_r$ and
$\|\Psi\|^2$ the product of the sup norms of its two members. It is impossible that
\begin{equation}\label{9.8:eq-no-invariant-germ-1}
\bigl|S(\Psi_r)-S(\Psi_{\rho r})\bigr|\le C_A\,r^2\,\|\Psi\|^2
\qquad\text{for all } r\le r_A .
\end{equation}
\end{lemma}

\begin{proof}
Suppose that \eqref{9.8:eq-no-invariant-germ-1} holds. We use Lemma~\ref{9.8:lem-tuned-two-sided},
which is proved as an implication from the statements named there. It gives the two-sided bound
\eqref{9.8:eq-tuned-two-sided-a} for the tuned limit. For $S$, with offset $\hat w=u\in W$, this
bound reads
\begin{equation}\label{9.8:eq-no-invariant-germ-2}
\mathfrak{m}(r)\,\|\Psi\|^2\le S(\Psi_r)\le \mathfrak{M}(r)\,\|\Psi\|^2
\qquad\text{for every } r\le r^{*}.
\end{equation}
The same lemma also gives that $\mathfrak{M}(r)\downarrow 0$ and $r^2/\mathfrak{m}(r)\to 0$ as
$r\downarrow 0$. Here $\mathfrak{m}(r)>0$ by its definition in that lemma.

\emph{Choice of the starting radius.} The number $C_A/(\rho^2-1)$ is finite and non-negative.
Hence
\begin{equation*}
\frac{C_A}{\rho^2-1}\cdot\frac{r^2}{\mathfrak{m}(r)}\to 0
\qquad\text{as } r\downarrow 0 .
\end{equation*}
We may therefore choose $r_0$ with $0<r_0\le\min\{r_A,r^{*}\}$ and
\begin{equation}\label{9.8:eq-no-invariant-germ-3}
\mathfrak{m}(r_0)>\frac{C_A\,r_0^2}{\rho^2-1}.
\end{equation}

\emph{The geometric chain.} Put $r_k:=\rho^{-k}r_0$ for $k\ge 0$. For $i\ge 1$ we have
$r_i\le r_0\le r_A$ and $\rho\,r_i=r_{i-1}$. Hence \eqref{9.8:eq-no-invariant-germ-1}, applied with
$r=r_i$, gives
\begin{equation*}
\bigl|S(\Psi_{r_i})-S(\Psi_{r_{i-1}})\bigr|\le C_A\,r_i^2\,\|\Psi\|^2 .
\end{equation*}
Write $S(\Psi_{r_k})-S(\Psi_{r_0})$ as the telescoping sum of these differences for
$i=1,\dots,k$ and apply the triangle inequality. For every $k\ge 1$ this yields
\begin{equation*}
\bigl|S(\Psi_{r_k})-S(\Psi_{r_0})\bigr|
\le C_A\,\|\Psi\|^2\sum_{i=1}^{k} r_i^2 .
\end{equation*}
The finite geometric sum satisfies
\begin{equation*}
\sum_{i=1}^{k} r_i^2
=r_0^2\sum_{i=1}^{k}\rho^{-2i}
<r_0^2\,\frac{\rho^{-2}}{1-\rho^{-2}}
=\frac{r_0^2}{\rho^2-1}.
\end{equation*}
Combining the last two displays, for every $k\ge 1$,
\begin{equation}\label{9.8:eq-no-invariant-germ-4}
\bigl|S(\Psi_{r_k})-S(\Psi_{r_0})\bigr|\le \frac{C_A\,r_0^2\,\|\Psi\|^2}{\rho^2-1}.
\end{equation}

\emph{The contradiction.} Every $r_k$ satisfies $r_k\le r_0\le r^{*}$, so
\eqref{9.8:eq-no-invariant-germ-2} applies at $r_k$ and at $r_0$. The upper bound at $r_k$ gives
$S(\Psi_{r_k})\le\mathfrak{M}(r_k)\|\Psi\|^2$, and the lower bound at $r_0$ gives
$S(\Psi_{r_0})\ge\mathfrak{m}(r_0)\|\Psi\|^2$. Together with \eqref{9.8:eq-no-invariant-germ-4}
this gives, for every $k\ge 1$,
\begin{equation*}
\bigl(\mathfrak{m}(r_0)-\mathfrak{M}(r_k)\bigr)\|\Psi\|^2
\le S(\Psi_{r_0})-S(\Psi_{r_k})
\le \frac{C_A\,r_0^2\,\|\Psi\|^2}{\rho^2-1}.
\end{equation*}
As $k\to\infty$ we have $r_k\downarrow 0$, hence $\mathfrak{M}(r_k)\to 0$. In the limit,
\begin{equation*}
\mathfrak{m}(r_0)\,\|\Psi\|^2\le \frac{C_A\,r_0^2\,\|\Psi\|^2}{\rho^2-1}.
\end{equation*}
Since $\|\Psi\|^2>0$, this says $\mathfrak{m}(r_0)\le C_A r_0^2/(\rho^2-1)$, which contradicts
\eqref{9.8:eq-no-invariant-germ-3}. Hence \eqref{9.8:eq-no-invariant-germ-1} cannot hold.
\end{proof}

\begin{definition}[Canonical representation and the scale mismatch]\label{9.8:not-scale-mismatch}
Let $L$ and $L'$ be the two blocking factors, and let $\varrho_2$ and the offset window
$W=[-\varrho_2/2,\varrho_2/2]$ be as in Definition~\ref{9.8:def-blocking-factors}; let
$a_n^{L}$, $a_{n'}^{L'}$ be the pinned values of Definition~\ref{9.8:def-blocking-factors} in the
two frames. Throughout, the conditions of Lemma~\ref{9.8:lem-increment-convergence} are assumed
in both frames $L$ and $L'$.
In the frame $L'$ we write $a'_{n'}:=a_{n'}^{L'}$ and $\alpha'_{n'}:=a'_{n'+1}-a'_{n'}$ for the
tuned increments, and $\bar{\alpha}'$ for the bound $\bar{\alpha}=3\lambda_U+c'^{\sharp}$ on the
tuned increments read in the frame $L'$ (every constant in it taken at blocking factor $L'$), so
that $\alpha'_{n'}\le\bar{\alpha}'$ for all $n'$.

\emph{The clock of the frame $L'$.} Let $\mathsf{A}'$ be the piecewise-linear interpolation of
$n'\mapsto a'_{n'}$:
\begin{equation*}
  \mathsf{A}'(t):=a'_{n'}+(t-n')\,\alpha'_{n'}\qquad (n'\le t\le n'+1,\ n'\in\mathbb{N}).
\end{equation*}
By Lemma~\ref{9.8:lem-increment-convergence}, read in the frame $L'$, the increments
$\alpha'_{n'}$ converge to a limit $\alpha'_\infty>0$; let $\underline{n}'$ be an integer such that
$\alpha'_{n'}\ge\alpha'_\infty/2$ for all $n'\ge\underline{n}'$. On $[\underline{n}',\infty)$ the map
$\mathsf{A}'$ is strictly increasing and continuous. For $x_0>a'_{\underline{n}'}$ we let
$\tau'(x_0)$ be the unique $t\ge\underline{n}'$ with $\mathsf{A}'(t)=x_0$, that is, the value at
$x_0$ of the inverse of the restriction of $\mathsf{A}'$ to $[\underline{n}',\infty)$.

\emph{Canonical representation in $L'$.} For such $x_0$ the \emph{canonical representation} of
$x_0$ in the frame $L'$ is the pair $(n',\hat w')$ given by
\begin{equation}\label{9.8:eq-scale-mismatch-1}
  n':=\lceil \tau'(x_0)\rceil,\qquad
  \hat w':=a'_{n'}-x_0=\bigl(n'-\tau'(x_0)\bigr)\,\alpha'_{n'-1}\in[0,\bar{\alpha}']\subset W .
\end{equation}

\emph{The scale mismatch.} For $\hat w\in W$ and $K$ with $a_K^{L}-\hat w>a'_{\underline{n}'}$ put
\begin{equation*}
  x_0^{(K)}:=a_K^{L}-\hat w,\qquad \theta_K:=\tau'\bigl(x_0^{(K)}\bigr).
\end{equation*}
Let $(n'_K,\hat w'_K)$ be the canonical representation of $x_0^{(K)}$ in $L'$. Set
\begin{equation*}
  \lambda_K:=\frac{L^K}{L'^{\,n'_K}} .
\end{equation*}
Define the scale mismatch $\delta_K(\hat w)$ by the first line of the following display; then
the identity on its second line holds:
\begin{equation}\label{9.8:eq-scale-mismatch-a}
\begin{aligned}
  \delta_K(\hat w)&:=K\log L-\theta_K\log L',\\
  \log\lambda_K&=\delta_K-(n'_K-\theta_K)\log L'
   =\delta_K-\hat w'_K\,\frac{\log L'}{\alpha'_{n'_K-1}} .
\end{aligned}
\end{equation}
\end{definition}

\begin{proof}
By Lemma~\ref{9.8:lem-increment-convergence}, read in the frame $L'$, the increments
$\alpha'_{n'}$ converge to a limit $\alpha'_\infty>0$. Hence there is an integer $\underline{n}'$ with
$\alpha'_{n'}\ge\alpha'_\infty/2$ for all $n'\ge\underline{n}'$.

On each interval $[n',n'+1]$ with $n'\ge\underline{n}'$, the function $\mathsf{A}'$ is affine with slope
$\alpha'_{n'}>0$. Consecutive pieces agree at the integers. Therefore $\mathsf{A}'$ is continuous
and strictly increasing on $[\underline{n}',\infty)$. Moreover, for $n'\ge\underline{n}'$,
\begin{equation*}
  \mathsf{A}'(n')\ge a'_{\underline{n}'}+(n'-\underline{n}')\,\alpha'_\infty/2,
\end{equation*}
so $\mathsf{A}'(t)\to\infty$ as $t\to\infty$. Consequently the restriction of $\mathsf{A}'$ to
$[\underline{n}',\infty)$ is a bijection onto $[a'_{\underline{n}'},\infty)$, and $\tau'$ is defined on
$(a'_{\underline{n}'},\infty)$ and takes values in $(\underline{n}',\infty)$.

By the same lemma read in the frame $L$, the increments $a^L_{n+1}-a^L_n$ converge to a positive
limit, so $a_K^L\to\infty$. Hence, for a fixed $\hat w\in W$, the condition
$a_K^{L}-\hat w>a'_{\underline{n}'}$ holds for all sufficiently large $K$.

Now let $x_0>a'_{\underline{n}'}$, write $\tau'=\tau'(x_0)$ and $n'=\lceil\tau'\rceil$. Then
$n'-1<\tau'\le n'$ and $n'\ge\underline{n}'+1$. On $[n'-1,n']$ the map $\mathsf{A}'$ is affine with slope
$\alpha'_{n'-1}$. Therefore
\begin{equation*}
  a'_{n'}-x_0=\mathsf{A}'(n')-\mathsf{A}'(\tau')=(n'-\tau')\,\alpha'_{n'-1}.
\end{equation*}
This is the equality in \eqref{9.8:eq-scale-mismatch-1}. Since $0\le n'-\tau'<1$ and
$0<\alpha'_{n'-1}\le\bar{\alpha}'$, we get $\hat w'\in[0,\bar{\alpha}']$. The inclusion
$[0,\bar{\alpha}']\subset W$ is the inequality $\bar{\alpha}'\le\varrho_2/2$ between the bound on the
increments and the half-width of $W$.

For the identity \eqref{9.8:eq-scale-mismatch-a}, the definitions of $\lambda_K$ and $\delta_K$ give
\begin{equation*}
  \log\lambda_K=K\log L-n'_K\log L'
  =\delta_K+\theta_K\log L'-n'_K\log L'
  =\delta_K-(n'_K-\theta_K)\log L'.
\end{equation*}
Apply \eqref{9.8:eq-scale-mismatch-1} to $x_0=x_0^{(K)}$, for which $\tau'=\theta_K$ and
$n'=n'_K$. This gives $\hat w'_K=(n'_K-\theta_K)\,\alpha'_{n'_K-1}$. Since $\alpha'_{n'_K-1}>0$, we
have $n'_K-\theta_K=\hat w'_K/\alpha'_{n'_K-1}$. Substituting this gives the last equality in
\eqref{9.8:eq-scale-mismatch-a}.
\end{proof}

\begin{remark}\label{9.8:thm-two-blockings}
The statement of this theorem is not written out here; parts of its proof are printed below.
\end{remark}

\begin{proof}[Proof of Theorem~\ref{9.8:thm-two-blockings}, parts (c)--(f), continued]
\label{9.8:prf-dilation-injectivity-proof}
Throughout, $\mathsf{F}\in\{L,L'\}$ denotes a frame. As in Lemma~\ref{9.8:lem-no-invariant-germ}, we
write $S^{(u)}_{\mathsf{F},m}$ for the tuned limit at offset $u$ in the frame $\mathsf{F}$ on the torus
$\mathbb{T}_m$, and we drop an index when the sentence fixes it. Thus $S^{(\hat w)}_{L}$ is the tuned
limit in the frame $L$ with the torus not indicated, and $S^{(u)}$ is the tuned limit in the frame $L'$
at offset $u$. We write
$\alpha'_\infty$ for the limit $\alpha_\infty$ of the tuned increments in the frame $L'$. The relation
$\approx$ is the germ equivalence of \eqref{9.8:eq-blocking-factors-e}. It compares two families
$S_1,S_2$ at the dilates $\vec\varphi^{\,r}$ of a test family: $|S_1(\vec\varphi^{\,r})
-S_2(\vec\varphi^{\,r})|$ is at most $C\,r^{4-\mathsf{p}}$ times the product of the sup norms of the
components of $\vec\varphi$, for all $r\le\bar r$, with a constant $C$ and a radius $\bar r>0$
depending on the pair. We call $C$ the constant and $\bar r$ the radius of the relation.

We use three elementary properties of the dilations $D_\rho$ and of $\approx$.

First, for $\rho_1,\rho_2>0$ we have $D_{\rho_1}D_{\rho_2}=D_{\rho_1\rho_2}$. A dilation does not
change sup norms, so $\vec\varphi^{\,\rho}$ is again sup-normalised. By the definition
$(D_\rho S)(\vec\varphi):=S(\vec\varphi^{\,\rho})$ and since
$\vec\varphi^{\,\rho_1}(\cdot/\rho_2)=\vec\varphi(\cdot/(\rho_1\rho_2))$,
\begin{equation*}
  \bigl(D_{\rho_1}(D_{\rho_2}S)\bigr)(\vec\varphi)
  =(D_{\rho_2}S)(\vec\varphi^{\,\rho_1})
  =S\bigl((\vec\varphi^{\,\rho_1})^{\rho_2}\bigr)
  =S(\vec\varphi^{\,\rho_1\rho_2}).
\end{equation*}
In particular $D_{1/\rho}D_\rho$ is the identity.

Second, $\approx$ is stable under a fixed dilation. If $S_1\approx S_2$ and $\mu>0$ is fixed, then
$D_\mu S_1\approx D_\mu S_2$. The reason is that the relation for $D_\mu S_1$ and $D_\mu S_2$ at scale
$r$ is the relation for $S_1$ and $S_2$ at scale $\mu r$. The constant of the relation is therefore
multiplied by $\mu^{4-\mathsf{p}}$, and the radius is divided by $\mu$. For $\mu=1/\lambda$ the factor is
$\lambda^{-(4-\mathsf{p})}$.

Third, if $S_1\approx S_2$ and $S_2\approx S_3$, then $S_1\approx S_3$ by the triangle inequality. The
constant is the sum of the two constants, and the radius is the smaller of the two radii. Below only
finitely many such relations are composed, each involving a dilation by a fixed factor. Hence no
uniformity of the constants over an infinite family is needed.

\emph{Part (c).} We run the argument of the proof of part~(b) that leads to
the corresponding display, together with the step that follows it. In it we take the two mismatch
values $\delta^a$ and $\delta^b$ of those steps both equal to $\delta_\infty(\hat w)$. The argument
then shows that every $u$ in the open interval
$(0,\alpha'_\infty)$ is a cluster value of the sequence $\hat w'_K$, along a subsequence on which
$\lambda_K\to\lambda$. We write $\lambda(u)$ for this limit. The two endpoints $u=0$ and
$u=\alpha'_\infty$ are obtained from the interior points by continuity in $u$. Along such a
subsequence the argument gives $D_{\lambda(u)}S^{(u)}\approx S^{(\hat w)}_{L}$; this is the form in
which part~(c) is used below.

\emph{Part (d).} We first work in the frame $L'$. Let $u_1,u_2\in[0,\alpha'_\infty]$ and fix one
$\hat w$. Part~(c) is applied twice with this same $\hat w$, once at $u_1$ and once at $u_2$. It gives
\begin{equation}\label{9.8:eq-dilation-injectivity-proof-1}
  D_{\lambda(u_1)}S^{(u_1)}\approx S^{(\hat w)}_{L}\approx D_{\lambda(u_2)}S^{(u_2)} .
\end{equation}
By the third property above, $D_{\lambda(u_1)}S^{(u_1)}\approx D_{\lambda(u_2)}S^{(u_2)}$.

We apply the fixed dilation $D_{1/\lambda(u_1)}$ to both sides. By the second property the relation
persists, with its constant multiplied by $\lambda(u_1)^{-(4-\mathsf{p})}$. By the first property,
$D_{1/\lambda(u_1)}D_{\lambda(u_1)}$ is the identity and
$D_{1/\lambda(u_1)}D_{\lambda(u_2)}=D_{\lambda(u_2)/\lambda(u_1)}$. Hence
\begin{equation}\label{9.8:eq-dilation-injectivity-proof-2}
  S^{(u_1)}\approx D_\rho S^{(u_2)},\qquad
  \rho=\frac{\lambda(u_2)}{\lambda(u_1)}=e^{(u_1-u_2)/b},
\end{equation}
with $b$ the universal slope of part~(a).

The other points of $W$ are reached through the exact covariance of the tuned family under blocking,
Corollary~\ref{9.8:cor-exact-covariance}. It states
\begin{equation*}
  S^{(\hat w)}_{\mathsf{F},m_0}=D_{\mathsf{F}}S^{(\hat w-\alpha_\infty)}_{\mathsf{F},m_0+1}
  \qquad\text{for } \hat w,\ \hat w-\alpha_\infty\in W^{\circ}.
\end{equation*}
Under this identity a shift of the offset by $\alpha_\infty$ corresponds to the dilation factor
$\rho=\mathsf{F}$. This is consistent with \eqref{9.8:eq-dilation-injectivity-proof-2}. Indeed, by
part~(a), $\alpha_\infty=b\log\mathsf{F}$, so a shift of $u_1-u_2$ by $\alpha_\infty$ multiplies
$e^{(u_1-u_2)/b}$ by $e^{\alpha_\infty/b}=\mathsf{F}$. For $u_1,u_2\in W$, each offset is moved into
$[0,\alpha'_\infty]$ by finitely many shifts by $\alpha'_\infty$, each between two points of
$W^{\circ}$; by Corollary~\ref{9.8:cor-exact-covariance} in the frame $L'$, each shift replaces the
tuned limit by its dilate by $L'$ or by $L'^{-1}$ on the neighbouring torus, and these finitely many
relations compose by the first and third properties. The identity also changes the torus from
$\mathbb{T}_{m_0+1}$ to $\mathbb{T}_{m_0}$. This change of torus is controlled by part~(a) of the
short-distance corollary, in its finite-cutoff form.

In the frame $L$ the same argument applies with the roles of the two frames exchanged. This proves
the corresponding display.

\emph{Part (e).} Suppose that $S^{(u_1)}_{L',m_1}\approx S^{(u_2)}_{L',m_2}$ for some $u_1>u_2$ in
$W$ and some torus exponents $m_1,m_2$. Part~(d)
gives $S^{(u_1)}\approx D_\rho S^{(u_2)}$, including the change of torus. We combine this with the
assumed relation, using the symmetry and transitivity of $\approx$. This gives
\begin{equation}\label{9.8:eq-dilation-injectivity-proof-3}
  S^{(u_2)}\approx D_\rho S^{(u_2)},\qquad \rho=e^{(u_1-u_2)/b}>1,
\end{equation}
where $\rho>1$ because $u_1>u_2$ and $b>0$ by part~(a).

We evaluate \eqref{9.8:eq-dilation-injectivity-proof-3} on the fixed $C^2$ pair
$\Psi\in\mathfrak{T}(1;\vartheta)$ and its dilates $\Psi_r$. Here $\mathsf{p}=2$, so
$r^{4-\mathsf{p}}=r^2$. Moreover $(D_\rho S^{(u_2)})(\Psi_r)=S^{(u_2)}(\Psi_{\rho r})$ by the first
property. Relation \eqref{9.8:eq-dilation-injectivity-proof-3} therefore yields constants
$C_A<\infty$ and $r_A>0$ such that
\begin{equation*}
  \bigl|S^{(u_2)}(\Psi_r)-S^{(u_2)}(\Psi_{\rho r})\bigr|\le C_A r^2\|\Psi\|^2
  \qquad\text{for all } r\le r_A .
\end{equation*}
Since $u_2\in W$ and $\rho>1$, this is excluded by Lemma~\ref{9.8:lem-no-invariant-germ}. Hence
$S^{(u_1)}_{L',m_1}\not\approx S^{(u_2)}_{L',m_2}$ whenever $u_1\neq u_2$.

Let now $u_1\neq u_2$ in $W$, and suppose that $\operatorname{Lim}_\infty[u_1]$ and
$\operatorname{Lim}_\infty[u_2]$ have a common $\mathbb{R}^4$ limit point $S_\infty$. By part~(b) of
the short-distance corollary, $S_\infty\approx S^{(u_1)}$ and $S_\infty\approx S^{(u_2)}$. By symmetry
and transitivity, $S^{(u_1)}\approx S^{(u_2)}$, which was excluded above. This proves
the corresponding display.

\emph{Part (f).} Part~(f) is part~(d) evaluated on the pair $\Psi$ and its dilates.
\end{proof}

\begin{corollary}[One germ and injective labelling of limit sets]\label{9.8:cor-one-germ}
Assume the hypotheses of Theorem~\ref{9.8:thm-two-blockings}. Let $\mathsf{F}$ be one of the two
blocking factors. For a bare inverse coupling $x_0=g_0^{-2}$, written in a representation with $n$
levels below the
reference scale and offset $\hat w\in W$ as in part~(b) of that theorem, let
$\ell_{\mathsf{F}}(x_0)=\mathsf{F}^n e^{-\hat w/b}$ be the physical length, in lattice units,
defined in that part. For $\mathsf{p}\in\{2,3\}$ write $\approx$
for the germ equivalence \eqref{9.8:eq-blocking-factors-e}, that is, agreement up to
$O(r^{4-\mathsf{p}})$ as $r\to0$ on test functions of the class
$\mathfrak{P}_{\mathsf{p}}(r;\vartheta)$.
\begin{itemize}
\item[(a)] \emph{One germ.} Along every bare sequence $x_0\to\infty$, on the reference tori
$\mathbb{T}_{m_0}$ and in every $\mathbb{R}^4$ limit, the two- and three-point functions, measured
in units of $\ell_{\mathsf{F}}(x_0)$, converge modulo $\approx$ to one germ $\mathfrak{g}_*$
(for the blocking factor $\mathsf{F}$ fixed above).
Tuning is therefore needed only to fix the unit of length.
\item[(b)] \emph{Injective labelling.} The renormalised coupling labels the $\mathbb{R}^4$ limit
sets injectively: the map $u\mapsto\operatorname{Lim}_\infty[u]$, $u\in W$, which assigns to a
renormalised offset $u$ the set of $\mathbb{R}^4$ limit points at that offset, is injective.
On the tori, consider the family of tuned torus limits $S^{(\hat w)}_{m_0}$, $\hat w\in W$; we call
$S^{(\hat w)}_{m_0}$ the member at offset $\hat w$ and read it in the unit of length
$e^{-\hat w/b}$, the factor of $\ell_{\mathsf{F}}(x_0)=\mathsf{F}^ne^{-\hat w/b}$ that depends on
the offset. Fix a $C^2$ pair $\Psi\in\mathfrak{T}(1;\vartheta)$, its dilates $\Psi_r$, and a
prescribed small value $v>0$. The family is calibrated by the length at which the connected
two-point function $S_2(\Psi_r)$ crosses $v$, in the following sense: for every member and every
$r>0$ at which its $S_2(\Psi_r)$ equals $v$, the length $r$, read in the member's unit, is a length
at which the two-point function of $\mathfrak{g}_*$ takes the value $v$ up to the tolerance of
$\approx$, and this description is the same for all members.
\end{itemize}
\end{corollary}

\begin{proof}
\emph{Part~(a).} We use three inputs. First, clause (ii-$\mathbb{T}$-a) of
Theorem~\ref{3.1:thm-main} gives the
convergence of the tuned functions $S^n_{m_0}(x_0;\vec f)$, for test functions $\vec f$ of the
class $\mathfrak{P}_{\mathsf{p}}(r;\vartheta)$, to their tuned limits
$S^{(\hat w)}_{m_0}$ without passing to subsequences and uniformly in the offset $\hat w\in W$.
This uniformity in $\hat w$ is what makes the conclusion hold along every bare sequence
$x_0\to\infty$, and not only along those on which the offset converges.

Second, part~(b) of Theorem~\ref{9.8:thm-two-blockings} states that $\delta_K(\hat w)$ converges,
uniformly on $W$, to $\delta_*+\hat w/b$. Part~(b) of Theorem~\ref{9.8:thm-two-blockings} also
states that $\ell_{\mathsf{F}}(x_0)$ may be computed in
any representation of the bare lattice: for two representations the ratio of the
two lengths tends to $1$. Measuring in units of $\ell_{\mathsf{F}}(x_0)$ therefore does not depend
on the representation chosen.

Third, part~(d) of Theorem~\ref{9.8:thm-two-blockings}, the corresponding display,
relates the tuned limits at any two offsets $u_1,u_2\in W$. It gives
\begin{equation*}
S^{(u_1)}_{m_0}\approx D_\rho S^{(u_2)}_{m_0},\qquad \rho=e^{(u_1-u_2)/b}.
\end{equation*}
At a fixed number $n$ of levels, the lengths at the two offsets are
$\mathsf{F}^ne^{-u_1/b}$ and $\mathsf{F}^ne^{-u_2/b}$. Their ratio is
\begin{equation*}
\frac{\mathsf{F}^ne^{-u_2/b}}{\mathsf{F}^ne^{-u_1/b}}=e^{(u_1-u_2)/b}=\rho .
\end{equation*}
Thus the dilation $D_\rho$ in part~(d) is exactly the change of the unit of length from
$\ell_{\mathsf{F}}$ at offset $u_2$ to $\ell_{\mathsf{F}}$ at offset $u_1$. Once distances are
measured in units of $\ell_{\mathsf{F}}(x_0)$, the tuned limits at all offsets lie in one class
modulo $\approx$. We call this class $\mathfrak{g}_*$.

Combining the three inputs, along every bare sequence $x_0\to\infty$ the two- and three-point
functions measured in units of $\ell_{\mathsf{F}}(x_0)$ converge modulo $\approx$ to
$\mathfrak{g}_*$. This holds on the reference tori by part~(d), and in every $\mathbb{R}^4$
limit by part~(e) of Theorem~\ref{9.8:thm-two-blockings}, the corresponding display, which concerns the sets $\operatorname{Lim}_\infty[u]$ of
$\mathbb{R}^4$ limit points. The offset $\hat w$ enters only through $\ell_{\mathsf{F}}(x_0)$, that
is, only through the unit of length.

\emph{Part~(b).} Injectivity of $u\mapsto\operatorname{Lim}_\infty[u]$ is the injectivity statement
of Theorem~\ref{9.8:thm-two-blockings}, the corresponding display.

For the calibration on tori, fix the pair $\Psi\in\mathfrak{T}(1;\vartheta)$, its dilates
$\Psi_r$ and the value $v$. By the argument of part~(a) with $\mathsf{p}=2$, every member
$S^{(\hat w)}_{m_0}$, $\hat w\in W$, read in its unit $e^{-\hat w/b}$, lies in the class
$\mathfrak{g}_*$: its two-point function $S_2(\Psi_r)$ agrees, up to the tolerance of $\approx$,
with the two-point function of $\mathfrak{g}_*$ evaluated on the dilate of $\Psi$ whose length,
read in the member's unit, is $r$. Hence, if $S_2(\Psi_r)=v$ for the member at offset $\hat w$,
then the two-point function of $\mathfrak{g}_*$ takes the value $v$, up to the tolerance of
$\approx$, at the length $r$ read in the member's unit. The function on which the crossing is
read, the two-point function of $\mathfrak{g}_*$, is the same for all members; only the unit
$e^{-\hat w/b}$ varies with the member. This is the calibration asserted in part~(b).
\end{proof}

The following are not asserted.
\begin{itemize}
\item A continuous dilation relating whole torus theories, or whole $\mathbb{R}^4$ limit sets,
rather than their germs modulo $\approx$. The families of tori built with the two blocking
factors have no common member. Their comparison
passes through a torus on which only reflection positivity is used, so that only the relation
$\approx$ survives.
\item A rate for the convergence in part~(b) of Theorem~\ref{9.8:thm-two-blockings}.
\item Any statement on the two-loop coefficient.
\item Analyticity in $u$ on $\mathbb{R}^4$.
\item Anything in the infrared.
\end{itemize}

\begin{remark}[Bounding the increments by the flow constants only. Stated; proof incomplete at the step marked $(\ast)$]
\label{9.8:rem-increment-bound}
Stated; the proof below is incomplete at the step marked $(\ast)$.

Fix one of the two frames $\mathsf{F}\in\{L,L'\}$ of Definition~\ref{9.8:def-blocking-factors};
every letter below is read at $\mathsf{F}$, and the conclusion holds in each frame. We use the
common $\varrho>0$, $\varrho_2=\varrho/2$, the window
$W=[-\varrho_2/2,\varrho_2/2]$, the pins $a_n$ and the tuned increments
$\alpha_n=a_{n+1}-a_n$ of Definition~\ref{9.8:def-blocking-factors}. Let $\varpi\le 1$ be the
tolerance of the scaffold comparison, and let $B^{\sharp}$ and $c_5$ be the flow constants of
clause~(0) of Theorem~\ref{3.1:thm-main}. Assume
\begin{equation}\label{9.8:eq-increment-bound-1}
\varrho_2\ \ge\ 24\,(B^{\sharp}+2\varpi)
\qquad\text{and}\qquad
\varrho_2\ \ge\ 24\,(c_5+2\varpi).
\end{equation}
Both conditions are lower bounds on $\varrho$ in terms of constants fixed before $\varrho$.
Then
\begin{equation}\label{9.8:eq-increment-bound-2}
|\alpha_n|\ \le\ \varrho_2/8\qquad\text{for all } n .
\end{equation}
Consequently, in the floor \eqref{9.8:eq-blocking-factors-b} the bound $\bar{\alpha}$ may be
read as $\varrho_2/8$; the inequality $\varrho_2\ge 4\bar{\alpha}$ then holds by itself, and
\eqref{9.8:eq-blocking-factors-b} becomes a lower bound on $\varrho$ depending only on
$B^{\sharp}$, $c_5$ and $\varpi$. In particular it does not involve $c'^{\sharp}$, the supremum
over $g\le\gamma_U$ of the function $c'(g)$ that enters $\bar{\alpha}=3\lambda_U+c'^{\sharp}$.
\end{remark}

\begin{proof}
Fix $n$. For $\hat w\in W$ let $H(\hat w)$ denote the inverse coupling reached after $n+1$ steps
of the depth-$(n+1)$ recursion started at the bare datum $a_n-\hat w$. After $n$ steps the
depth-$n$ recursion started at the same bare datum $a_n-\hat w$ reaches the value
$x_{\mathrm{ref},n}(\hat w)$, the reference coupling of the bi-Lipschitz lemma.

$(\ast)$ For every $\hat w\in W$ the $n+1$ steps of the depth-$(n+1)$ recursion at the datum
$a_n-\hat w$ are admissible, so that $H$ is defined and continuous on $W$. After $n$ steps, the
depth-$(n+1)$ recursion and the depth-$n$ recursion started at
$a_n-\hat w$ (two scaffolds, shifted against each other by one level) both remain within
$\varpi$ of one and the same native run, the run with which the scaffold comparison compares
them. Hence, after $n$ steps, the two recursions differ by at most $2\varpi$.

By the lower bound, with constant $1/3$, of the bi-Lipschitz lemma, the value
$x_{\mathrm{ref},n}(\hat w)$ of the depth-$n$ recursion satisfies
\begin{equation*}
|x_{\mathrm{ref},n}(\hat w)-X|\ \ge\ |\hat w|/3,
\end{equation*}
and $x_{\mathrm{ref},n}(\hat w)-X$ has the sign of $-\hat w$. By clause~(0) of
Theorem~\ref{3.1:thm-main}, applied to a single step, the last step of the depth-$(n+1)$
recursion lowers the inverse coupling by an amount $\beta$ with $-c_5\le\beta\le B^{\sharp}$.

Both points $\pm\varrho_2/8$ lie in $W$. At $\hat w=-\varrho_2/8$ we have
$x_{\mathrm{ref},n}(\hat w)\ge X+\varrho_2/24$; combining this with the two facts just stated
and with the first inequality of \eqref{9.8:eq-increment-bound-1}, which says
$\varrho_2/24\ge B^{\sharp}+2\varpi$, gives
\begin{equation*}
H(-\varrho_2/8)\ \ge\ X+\frac{\varrho_2}{24}-2\varpi-B^{\sharp}\ \ge\ X .
\end{equation*}
At $\hat w=\varrho_2/8$ we have $x_{\mathrm{ref},n}(\hat w)\le X-\varrho_2/24$, and the second
inequality of \eqref{9.8:eq-increment-bound-1}, which says $\varrho_2/24\ge c_5+2\varpi$, gives
\begin{equation*}
H(\varrho_2/8)\ \le\ X-\frac{\varrho_2}{24}+2\varpi+c_5\ \le\ X .
\end{equation*}
Since $H$ is continuous on $[-\varrho_2/8,\varrho_2/8]\subset W$, the intermediate value theorem
yields a point $\hat w^{\dagger}$ with $|\hat w^{\dagger}|\le\varrho_2/8$ and
$H(\hat w^{\dagger})=X$. By clause~(v-a) of Theorem~\ref{3.1:thm-main}, there is exactly one real
pin at depth $n+1$; the bare datum $a_n-\hat w^{\dagger}$ has the defining property of that pin,
and hence
\begin{equation*}
a_{n+1}=a_n-\hat w^{\dagger}.
\end{equation*}
Thus $\alpha_n=a_{n+1}-a_n=-\hat w^{\dagger}$ and $|\alpha_n|\le\varrho_2/8$, which is
\eqref{9.8:eq-increment-bound-2}.

The consequence follows. By \eqref{9.8:eq-increment-bound-2}, the number $\varrho_2/8$ bounds
every $|\alpha_n|$, so it may stand for $\bar{\alpha}$ in \eqref{9.8:eq-blocking-factors-b}. With
this reading, $4\bar{\alpha}=\varrho_2/2\le\varrho_2$, so that inequality holds for every
$\varrho>0$. What remains of \eqref{9.8:eq-blocking-factors-b}, together with
\eqref{9.8:eq-increment-bound-1}, involves only $B^{\sharp}$, $c_5$ and $\varpi$.

Step $(\ast)$ uses two facts that are not proved. First, the scaffold comparison
has to apply to the depth-$(n+1)$ scaffold at the bare data $a_n-\hat w$, $\hat w\in W$. Its
hypothesis is stated for data in the window of $a_{n+1}$, so a continuity or bootstrap argument
along $\hat w$ is needed to reach these data. Second, the $(n+1)$-st step has to be admissible
at $\hat w>0$; this is a ceiling on $g$.
\end{proof}

\begin{definition}[Weighted action, sources and constant-source correlations]
\label{9.9:not-weighted-action}
The gauge group is $\mathrm{SU}(2)$. The formulas are written for $\mathrm{SU}(N)$, with the trace
normalised by $\operatorname{tr}\mathbf{1}=1$; $N$ is the rank letter of the group, and we put
$n_{\mathfrak{g}}:=N^2-1$.

\emph{Lattice and weighted action.} Let $\mathbb{F}$ be the bare torus lattice of
Definition~\ref{1.1:def-lattice}, and let $dU$ be the product of the normalised Haar measures
over its bonds. For a bare plaquette $p$ put $a_p(U):=1-\operatorname{Re}\operatorname{tr}U(\partial p)$.
For a complex-valued function $\eta$ on the plaquettes of $\mathbb{F}$, the \emph{Wilson action
with pointwise weight} $\eta$ is
\begin{equation*}
A(\eta,U):=\sum_p \eta(p)\,a_p(U).
\end{equation*}
A function $f$ on the torus is read as the function $p\mapsto f(x(p))$ on plaquettes, and a
number as a constant function. For Lipschitz $f$ on the torus, $\mathcal{O}(f)=\sum_p f(x(p))\,a_p(U)$
is the smeared curvature observable of Definition~\ref{1.3:def-curvature-field}. Thus
$\mathcal{O}(f)=A(f,U)$.

\emph{Sources.} Let $x_0>0$ be a bare inverse coupling, and let $z=(z_1,\dots,z_{\mathsf{p}})$ be
sources with test functions $f_1,\dots,f_{\mathsf{p}}$. For a function $\eta$ on plaquettes put
$Z(x_0;\eta):=\int\exp[-A(x_0-\eta,U)]\,dU$, and set
\begin{equation}\label{9.9:eq-weighted-action-1}
Z(x_0;z):=\int\exp\Bigl[-x_0A(1,U)+\sum_i z_i\,\mathcal{O}(f_i)\Bigr]dU
=Z\Bigl(x_0;\sum_i z_if_i\Bigr).
\end{equation}

\emph{Connected functions.} Let $K$ be the number of renormalisation steps from the bare lattice to
the reference scale. For the runs of the uniqueness theorem on the bare datum $(m_0+n,x)$, that
is, with torus exponent $m_0+n$ and bare inverse coupling $x$, one has $K=n$. Put
\begin{equation*}
S^T_{\mathsf{p};K}(f_1,\dots,f_{\mathsf{p}})
:=\partial_{z_1}\cdots\partial_{z_{\mathsf{p}}}\big|_{z=0}\log Z(x_0;z).
\end{equation*}
Let $\mu_K$ be the probability measure $Z(x_0;0)^{-1}\exp[-x_0A(1,U)]\,dU$ on configurations of
$\mathbb{F}$, and let $\operatorname{Cov}_{\mu_K}$ be the covariance under it. The connected
three-point function with one constant source is
\begin{equation}\label{9.9:eq-weighted-action-2}
S^T_{3;K}(1,f_1,f_2):=\partial_{z_0}\partial_{z_1}\partial_{z_2}\big|_{0}
\log Z(x_0;\,z_0\cdot1+z_1f_1+z_2f_2).
\end{equation}

These objects have the following properties.
\begin{itemize}
\item[(a)] A constant source $z_0\cdot1$ is exactly an additive shift of the bare inverse coupling,
with no cutoff-dependent normalisation:
\begin{equation}\label{9.9:eq-weighted-action-3}
Z(x_0-z_0;z)=Z\Bigl(x_0;\,z_0\cdot1+\sum_i z_if_i\Bigr).
\end{equation}
\item[(b)] $S^T_{2;K}(f_1,f_2)=\operatorname{Cov}_{\mu_K}(\mathcal{O}(f_1),\mathcal{O}(f_2))$.
\item[(c)] $Z$ is an entire function of $(x_0,z)$, and $Z>0$ at real arguments. Consequently
$\log Z$ and all $S^T$ are real-analytic in $x_0$ on the real axis and holomorphic near it.
\end{itemize}
\end{definition}

\begin{proof}
The action $A(\eta,U)$ is linear in $\eta$, and $\mathcal{O}(f)=A(f,U)$. Hence
\begin{equation*}
-x_0A(1,U)+\sum_i z_i\,\mathcal{O}(f_i)=-A\Bigl(x_0-\sum_i z_if_i,\,U\Bigr).
\end{equation*}
This gives the second equality in \eqref{9.9:eq-weighted-action-1}.

(a) By \eqref{9.9:eq-weighted-action-1}, the left side of \eqref{9.9:eq-weighted-action-3} is the
integral of $\exp[-A(x_0-z_0-\sum_i z_if_i,U)]$. The weight in this integrand equals
$x_0-(z_0\cdot1+\sum_iz_if_i)$. So the integral is $Z(x_0;z_0\cdot1+\sum_iz_if_i)$, which is the
right side.

(c) Each $a_p$ is a continuous function on the compact configuration space, with values in $[0,2]$.
Hence $A(1,U)$ and each $\mathcal{O}(f_i)$ are bounded functions of $U$. For each $U$ the
integrand of \eqref{9.9:eq-weighted-action-1} is entire in $(x_0,z)$. On every compact set of
parameters it is bounded uniformly in $U$. $Z$ is therefore an integral, against the finite Haar
measure over a compact group, of an entire, locally bounded function of $(x_0,z)$, hence entire.
The same bound justifies differentiating under the integral sign in $(x_0,z)$.

At real $(x_0,z)$ the integrand is positive, so $Z>0$. Since $Z$ is continuous and has no zero at
real arguments, it has no zero in a complex neighbourhood of the real arguments. On such a
neighbourhood, $\log Z$, taken real at real arguments, is holomorphic, and it is real-analytic in
$x_0$ on the real axis. Every $S^T$ is a derivative of $\log Z$ in the sources at $z=0$, so each
has the same two properties in $x_0$.

(b) Write $\mathcal{O}_i=\mathcal{O}(f_i)$ and let $e_z(U)$ be the integrand of
\eqref{9.9:eq-weighted-action-1}. Differentiating under the integral sign gives
$\partial_{z_i}Z=\int\mathcal{O}_i\,e_z\,dU$. Hence
\begin{equation*}
\partial_{z_1}\partial_{z_2}\log Z(x_0;z)
=\frac{\int\mathcal{O}_1\mathcal{O}_2\,e_z\,dU}{Z(x_0;z)}
-\frac{\int\mathcal{O}_1\,e_z\,dU\,\int\mathcal{O}_2\,e_z\,dU}{Z(x_0;z)^2}.
\end{equation*}
At $z=0$ we have $e_0(U)\,dU/Z(x_0;0)=d\mu_K$. The right side then becomes
\begin{equation*}
\int\mathcal{O}_1\mathcal{O}_2\,d\mu_K-\int\mathcal{O}_1\,d\mu_K\int\mathcal{O}_2\,d\mu_K
=\operatorname{Cov}_{\mu_K}(\mathcal{O}_1,\mathcal{O}_2).\qedhere
\end{equation*}
\end{proof}

\begin{definition}[The running shift and its increments]\label{9.9:not-running-shift}
Throughout, $D(0,\varrho)\subset\mathbb{C}$ denotes the open disc of radius $\varrho$ about $0$.

\emph{Running couplings.} For a run of Ba{\l}aban's renormalisation recursion started from the bare
inverse coupling $x_0$, we write $x_k:=g_k^{-2}$ for the running inverse coupling after $k$ steps
(Definition~\ref{1.2:def-couplings}). We write $X:=g^{-2}$ for the target value.

\emph{The running shift.} Let $\zeta$ be a constant bare shift. The running shift is defined by
\begin{equation}\label{9.9:eq-running-shift-1}
\zeta_0:=\zeta,\qquad \zeta_{k+1}:=\zeta_k+b_{k+1}(\zeta)-b_{k+1}(0).
\end{equation}
Here $b_{k+1}(\zeta)$ is the number defined in \cite[(3.61)]{B-CMP109}, computed for the
whole-lattice family of the correlation theorem, that is, the family of effective densities on the
whole lattice considered there, at bare weight $x_0-\zeta$. In computing it, every threshold, every
cube size, every decomposition, every background and every reference law is frozen at its value in
the run with $\zeta=0$. No statement is made here about the run of the recursion started at
$x_0-\zeta$.

The increments of the running shift and the derivative at the origin are
\begin{equation}\label{9.9:eq-running-shift-2}
h_k(\zeta):=b_{k+1}(\zeta)-b_{k+1}(0),\qquad N_k:=\zeta_k'(0).
\end{equation}

\emph{Functions of the coupling and constants of the correlation theorem.} Let
$A,\sigma_0,C_t,C_{\mathfrak{m}},c_0,c_2,E^{*},\gamma_2,\varepsilon_1$ be the constants of the
correlation theorem (the constant $A$ is a number, not the action), and let
$\kappa_0$ be the exponent of the correlation theorem that enters its remnant bound through $v_k$
below. The letters $c_0,c_2$ here denote the constants of the correlation theorem; no
identification with the flow constants of Definition~\ref{1.2:def-flow-constants} is used in what
follows. Let $\varrho_J$ be the source radius of the correlation theorem, which is fixed
before $\gamma_J$. With $T=\log g^{-2}$ the degree variable of
Definition~\ref{9.1:not-levelwise-rate}, we put $\lambda:=c_0/2$ and
\begin{align*}
\mathcal{T}(g)&:=A_1T^{p_0}, &\mathsf{r}&:=\frac{\varepsilon_1}{2\mathcal{T}(g)},
&\tau&:=\frac{g}{\mathsf{r}},\\
\vartheta_e&:=E^{*}e^{-\frac14\gamma_2\mathcal{T}(g)^2}, &e(g)&:=2E^{*}\tau+\vartheta_e,
&\psi(g)&:=e(g)\,g^2,
\end{align*}
and $e_k:=e(g_k)$. The function $e$ is nondecreasing in $g$ on $(0,\gamma_J]$. We further put
\begin{align*}
y_k&:=\min_{i\le k}x_i-2c_2,
&v_k&:=2AC_tC_{\mathfrak{m}}\,y_k^{-(\kappa_0-4)/2},\\
\theta'_k&:=\frac{16A}{9\sigma_0}\,\psi\bigl(y_k^{-1/2}\bigr),
&\theta''_k&:=2v_k,
\end{align*}
and, for $\varrho>0$,
\begin{equation}\label{9.9:eq-running-shift-3}
d_k(\varrho):=\theta'_k+\frac{\theta''_k}{\varrho}.
\end{equation}
The degree function $p_0(\cdot)$, $p_0(g):=A_0(\log g^{-2})^{\mathfrak{p}_0}$ (a polylogarithm of
$g^{-2}$ in the sense fixed below), is the large-field threshold function of \cite{B-CMP119} in the
form in which the two-point witness theorem uses it; its exponent is written $\mathfrak{p}_0$ in
this chapter, and no identification of $\mathfrak{p}_0$ with the exponent $p_0$ of $\mathcal{T}$ is
used in what follows.

\emph{What the correlation theorem provides.} We use the following facts, all taken from the
correlation theorem.
\begin{itemize}
\item[(a)] Each $\zeta_k$ is analytic on $D(0,2\varrho_J)$ and real on real arguments.
\item[(b)] For every $\varrho\in[\varrho_J,2\varrho_J]$, on $D(0,\varrho)$ we have
$|h_k(\zeta)|\le d_k(\varrho)\,|\zeta|$, and
\begin{equation*}
\sum_{k<K}d_k(\varrho)\le\tfrac13 .
\end{equation*}
\item[(c)] Along $\zeta$, the effective action at level $k$ is minus the weighted action
$A(x_k(\cdot)-\zeta_k(\zeta),U_k)$ of Definition~\ref{9.9:not-weighted-action}, plus
further terms. Here $x_k(\cdot)$ is the position-dependent inverse coupling of
Definition~\ref{1.5:def-domains} and $U_k$ is the block field. Accordingly, we call
$x_k-\zeta_k(\zeta)$ the effective inverse coupling at level $k$ along $\zeta$.
\item[(d)] The running-shift statement comes with the following decomposition, which we call its
fine bookkeeping:
\begin{equation}\label{9.9:eq-running-shift-4}
b_{k+1}=\beta^{0}_{k+1}+f_{k+1}(\zeta_0,\dots,\zeta_k)+\mathsf{q}_{k+1}(\zeta).
\end{equation}
In this decomposition:
\begin{itemize}
\item $\beta^{0}_{k+1}$ is the term of the coupling increment $b_{k+1}$ that, in this
decomposition, carries no argument $\zeta$; it is the number so written in the correlation
theorem.
\item $f_{k+1}$ is holomorphic on the polydisc $\mathfrak{P}_k:=\prod_{j\le k}D(0,\sigma_0x_j)$,
with coordinates $\xi_0,\dots,\xi_k$, and satisfies $|f_{k+1}|\le Ae_k$ there.
\item $f_{k+1}$ obeys the memory bound
\begin{equation}\label{9.9:eq-running-shift-5}
\sup_{\frac12\mathfrak{P}_k}\ \sum_{j\le k}\sigma_0x_j
\Bigl|\frac{\partial f_{k+1}}{\partial\xi_j}\Bigr|\le\tfrac43\,Ae_k ,
\end{equation}
where $\tfrac12\mathfrak{P}_k$ is the polydisc with halved radii.
\item $\mathsf{q}_{k+1}(\zeta)$ is the remnant of the coupling increment $b_{k+1}$ at step $k+1$;
in the correlation theorem it is written $\beta[Q^{(k+1)}_{\zeta}]$, the coupling-increment
functional $\beta$ of the correlation theorem evaluated at the term $Q^{(k+1)}_{\zeta}$ so denoted
there. It is holomorphic on
$D(0,2\varrho_J)$ and satisfies $|\mathsf{q}_{k+1}|\le v_k$ there;
this is the remnant value bound of the correlation theorem.
\end{itemize}
\end{itemize}

\emph{Consequences.} For $0\le k\le K$ the following hold. Fact (e) is immediate from
\eqref{9.9:eq-running-shift-1} and \eqref{9.9:eq-running-shift-2}; facts (f) and (g) are part of
the statement of the correlation theorem, and we include their derivation from (a) and (b).
\begin{itemize}
\item[(e)] The running shift is the bare shift plus the sum of its increments:
\begin{equation}\label{9.9:eq-running-shift-6}
\zeta_k(\zeta)=\zeta+\sum_{j<k}h_j(\zeta).
\end{equation}
\item[(f)] On $D(0,2\varrho_J)$ we have $|\zeta_k(\zeta)-\zeta|\le|\zeta|/3$.
\item[(g)] The derivative at the origin is a real number in a fixed interval:
\begin{equation}\label{9.9:eq-running-shift-7}
N_k=1+\sum_{1\le j\le k}b_j'(0)\in\bigl[\tfrac23,\tfrac43\bigr].
\end{equation}
\end{itemize}

\emph{Words.} A \emph{polylogarithm of $X$} is a constant times a power of $\log X$. A bound is
\emph{free of $g$} when it is made small by a constant chosen before the coupling, rather than by a
power of $g$. The relevant constants are those of the two-point witness theorem that set the
separation of the supports in cells of the witness lattice and the length of the Gaussian segment.
The \emph{currency} of a statement is the class of observables it is about; here it is the
connected correlations of $\mathsf{p}\ge2$ smeared curvature fields of clause (iii)(b) of
Theorem~\ref{3.1:thm-main}.
\end{definition}

\begin{proof}
We prove (e), (f) and (g) from (a) and (b).

\emph{Proof of (e).} We argue by induction on $k$. For $k=0$ the sum in
\eqref{9.9:eq-running-shift-6} is empty, and $\zeta_0=\zeta$ by \eqref{9.9:eq-running-shift-1}.
Assume \eqref{9.9:eq-running-shift-6} holds for some $k<K$. By \eqref{9.9:eq-running-shift-1} and
the definition of $h_k$ in \eqref{9.9:eq-running-shift-2},
\begin{equation*}
\zeta_{k+1}(\zeta)=\zeta_k(\zeta)+h_k(\zeta)=\zeta+\sum_{j<k+1}h_j(\zeta).
\end{equation*}
This is \eqref{9.9:eq-running-shift-6} at $k+1$.

\emph{Proof of (f).} Take $\varrho=2\varrho_J\in[\varrho_J,2\varrho_J]$ in (b), and let
$\zeta\in D(0,2\varrho_J)$.

First, each $d_j(2\varrho_J)$ is nonnegative. Indeed, for $\zeta\ne0$ in the disc the inequality
$0\le|h_j(\zeta)|\le d_j(2\varrho_J)|\zeta|$ forces $d_j(2\varrho_J)\ge0$.

Now apply (e), the triangle inequality and (b). Since $k\le K$ and every term is nonnegative, the
sum over $j<k$ is at most the sum over $j<K$:
\begin{equation*}
|\zeta_k(\zeta)-\zeta|
\le\sum_{j<k}|h_j(\zeta)|
\le|\zeta|\sum_{j<k}d_j(2\varrho_J)
\le|\zeta|\sum_{j<K}d_j(2\varrho_J)
\le\tfrac13|\zeta|.
\end{equation*}

\emph{Proof of (g).} By (a), $\zeta_j$ and $\zeta_{j+1}$ are analytic on $D(0,2\varrho_J)$. By
\eqref{9.9:eq-running-shift-1}, $h_j=\zeta_{j+1}-\zeta_j$, so $h_j$ is analytic there as well.

Differentiate \eqref{9.9:eq-running-shift-6} at $0$. Since $h_j'(0)=b_{j+1}'(0)$, and after
shifting the summation index by one,
\begin{equation*}
N_k=\zeta_k'(0)=1+\sum_{j<k}h_j'(0)=1+\sum_{1\le j\le k}b_j'(0).
\end{equation*}

Next we bound each $h_j'(0)$. By \eqref{9.9:eq-running-shift-2}, $h_j(0)=b_{j+1}(0)-b_{j+1}(0)=0$.
Hence
\begin{equation*}
h_j'(0)=\lim_{\zeta\to0}\frac{h_j(\zeta)}{\zeta}.
\end{equation*}
By (b) with $\varrho=2\varrho_J$, the quotient under the limit has modulus at most
$d_j(2\varrho_J)$. Therefore $|h_j'(0)|\le d_j(2\varrho_J)$.

Combining this with the nonnegativity of the $d_j(2\varrho_J)$ shown in the proof of (f), and with
(b), we obtain
\begin{equation*}
|N_k-1|\le\sum_{j<k}d_j(2\varrho_J)\le\sum_{j<K}d_j(2\varrho_J)\le\tfrac13 .
\end{equation*}

Finally, $N_k$ is real. By (a), $\zeta_k$ is real on real arguments, so for real $t\ne0$ the
difference quotient $(\zeta_k(t)-\zeta_k(0))/t$ is real. Letting $t\to0$ along the real axis shows
that $\zeta_k'(0)$ is real. Together with $|N_k-1|\le\tfrac13$ this gives
$N_k\in[\tfrac23,\tfrac43]$, which is \eqref{9.9:eq-running-shift-7}.
\end{proof}

\begin{definition}[Tuned family, calibrating function and witness data]
\label{9.9:not-tuned-family}
Throughout, $g\le\gamma_U$.

\emph{The tuned family.} Let $m_0$ be the torus exponent at the reference scale. For each integer
$n\ge0$, $a_n$ is the real bare inverse coupling for which the run of length $n$ of the class
$\Theta_X$ on the bare datum $(m_0+n,a_n)$ has inverse coupling exactly $X$ at the reference scale,
$x_{(0)}=X$. The bare datum $(m_0+n,a_n)$ has torus exponent $m_0+n$ and bare inverse coupling
$a_n$. The run is written in the depth-indexed notation of the uniqueness theorem:
\begin{equation*}
\zeta_{(n)}:=\zeta,\qquad
\zeta_{(r)}:=\zeta_{(r+1)}+b_{(r)}-c_0,\qquad
x_{(r)}:=\theta_r-\zeta_{(r)} .
\end{equation*}
Depth $r=0$ is the reference scale, depth $r=n$ is the bare scale, and $\theta_r$ is the reference
ladder of the class $\Theta_X$, indexed by the depth only.

\emph{The correlation functions of the family.} Let $\mathsf{p}\ge2$ be an integer and let
$f=(f_1,\dots,f_{\mathsf{p}})$ be Lipschitz functions with pairwise separated supports on the
continuum torus of exponent $m_0$. Then $S^n(x;f)$ denotes the connected $\mathsf{p}$-point function
of $\mathcal{O}(f_1),\dots,\mathcal{O}(f_{\mathsf{p}})$ on the bare datum $(m_0+n,x)$. As
$n\to\infty$, $S^n(a_n-\hat w;f)$ converges to a limit
$S^{(\hat w)}(f)$ without passage to subsequences. The convergence is uniform for $\hat w$ in
\begin{equation*}
\mathcal{N}'_w:=\bigl\{\hat w\in\mathbb{C}:\ |\operatorname{Re}\hat w|\le\varrho_2/2,\
|\operatorname{Im}\hat w|<\upsilon_w\bigr\}.
\end{equation*}
The map $\hat w\mapsto S^{(\hat w)}$ is real-analytic on $]-\varrho_2/2,\varrho_2/2[$ and holomorphic
on the strip $\mathcal{N}'_w$.

\emph{The calibrating function.} We fix $\mathsf{p}=2$ and the reference pair
$f^{\circ}=(f_1^{\circ},f_2^{\circ})$ chosen in the next paragraph, and put
\begin{equation}\label{9.9:eq-tuned-family-1}
\varphi_n(\hat w):=S^n(a_n-\hat w;f^{\circ}),\qquad
\varphi(\hat w):=S^{(\hat w)}(f^{\circ})=\lim_{n\to\infty}\varphi_n(\hat w).
\end{equation}

\emph{The witness objects.} Fix a shape constant $\vartheta\in\,]0,1]$. The depth function of the
two-point witness theorem is
\begin{equation*}
s(\mathsf{d}):=\min\{s:\ L^s\mathsf{d}\ge C_W(s)\},\qquad
C_W(s):=C_0\,p_0\bigl(\bar{x}_s^{-1/2}\bigr)^2 ,
\end{equation*}
where $p_0(\cdot)$ is the degree function. The depth floor $s_W(g)$ of that theorem is a one-sided
condition depending on $g$. We fix one separation $\mathsf{d}^{\circ}$ with
$s^{\circ}:=s(\mathsf{d}^{\circ})\ge s_W(g)$.

We also fix one reference pair $f^{\circ}=(f_1^{\circ},f_2^{\circ})$ in the continuum class
$\mathfrak{T}(\mathsf{d}^{\circ};\vartheta)$. Its members are pairs of non-negative Lipschitz bumps
on the continuum torus with the following properties:
\begin{itemize}
\item each bump is supported in a ball of radius $\mathsf{d}^{\circ}$;
\item each bump has integral over the continuum torus at least
$\vartheta\|f_i\|_\infty(\mathsf{d}^{\circ})^4$;
\item each bump has Lipschitz constant at most $\|f_i\|_\infty/(\vartheta\mathsf{d}^{\circ})$;
\item the two supports lie at mutual distance in $[\mathsf{d}^{\circ},6\mathsf{d}^{\circ}]$.
\end{itemize}
The last condition is two-sided; it is a condition on the test functions.

The samples of $f^{\circ}$ on the witness lattice of every bare datum lie, up to Riemann sums, in
$\mathfrak{F}^{+}(\mathsf{d}^{\circ},\vartheta)^2$. Here $\mathfrak{F}^{+}(\mathsf{d}^{\circ},\vartheta)$
is the lattice class of the two-point witness theorem. Put $\delta^{\circ}:=L^{s^{\circ}}\mathsf{d}^{\circ}$.
The members of this class are bumps with the following properties:
\begin{itemize}
\item each bump is supported in a ball whose radius is $\delta^{\circ}$ cells of the witness lattice;
\item each bump has mass $\mathfrak{n}(f_i)\ge\vartheta(\delta^{\circ})^4\|f_i\|_\infty$;
\item each bump has slope at most $\|f_i\|_\infty/(\vartheta\delta^{\circ})$;
\item the two supports lie at distance in $[\mathsf{d}^{\circ},6\mathsf{d}^{\circ}]$.
\end{itemize}

\emph{Steps and level.} On the bare datum $(m_0+n,a_n)$ the number of renormalisation steps to the
reference scale is $K:=n$. The witness level is $k_0:=n-s^{\circ}$, which is depth $s^{\circ}$ of the
run of the class $\Theta_X$. Thus $x_{k_0}=x_{(s^{\circ})}$ is the inverse coupling of that run at
depth $s^{\circ}$; it satisfies the window
\begin{equation}\label{9.9:eq-tuned-family-a}
\underline{x}^{\Theta}_{s^{\circ}}:=X+\lambda_U s^{\circ}-c'
\;\le\; x_{k_0}\;\le\;
X+3\lambda_U s^{\circ}+c'=:\bar{x}^{\Theta}_{s^{\circ}} .
\end{equation}
Here $c'$ is the constant of clause (a) of the uniqueness theorem,
\begin{equation*}
c':=2c_2+\mathsf{A}\,C_E\bigl[(1-\rho)^{-2}+T_{\Pi}\bigr],
\end{equation*}
where $\mathsf{A}$ is the extraction constant, $\rho\in\,]0,1[$ the contraction rate and $T_{\Pi}$
the large-field floor of the uniqueness theorem, and $C_E$ is the constant of that theorem which also
enters the bound \eqref{9.9:eq-coupling-convergence-a}.
Both bounds in \eqref{9.9:eq-tuned-family-a} are independent of $n$. The pairwise increments of the run,
together with the pin $x_{(0)}=X$, also give the lower bound $x_{(r)}\ge X+\lambda r-2c_2$ at every
depth $r$, where $\lambda$ is the constant in that bound on the pairwise increments of the run.
This lower bound is free of $g$ and is used among the conditions defining
$\gamma_{\mathrm{cal}}$.

In the proofs of this section that apply the two-point witness theorem on the data above, the window
\eqref{9.9:eq-tuned-family-a} is used wherever that theorem uses its first-passage envelope
$[\underline{x}_s,\bar{x}_s]$. This includes the quantity $C_W(s)$ and the floors $s(\mathsf{d})$ and
$s_W(g)$, which are read with $\bar{x}^{\Theta}_s$ in place of $\bar{x}_s$, where, for every depth
$s$,
\begin{equation*}
\underline{x}^{\Theta}_s:=X+\lambda_U s-c',\qquad
\bar{x}^{\Theta}_s:=X+3\lambda_U s+c'.
\end{equation*}
In that reading of the two-point witness theorem on the bare data of the uniqueness theorem, given
below, the reference exponent satisfies $m_0\ge m_{\Theta}(g)$, where the torus floor
$m_{\Theta}(g)$ of the reading is the least $m_0$ for which the torus hypothesis of the two-point
witness theorem holds on the bare tori of exponent $m_0+n$ for every $n\ge0$.

\emph{The kernel.} Let $C^{\mathrm{bg}}_{K,s^{\circ}}$ be the covariance of the background part of the
unit lattice Maxwell curvature. The kernel is
\begin{equation}\label{9.9:eq-tuned-family-2}
\mathcal{K}:=\mathcal{K}_{K,s^{\circ}}(f_1^{\circ},f_2^{\circ})
:=\sum_{p,q}f_1^{\circ}(x(p))\,f_2^{\circ}(x(q))\,
\bigl[C^{\mathrm{bg}}_{K,s^{\circ}}(p,q)\bigr]^2 .
\end{equation}
It depends on the lattice ($K$, $s^{\circ}$, $L$, the torus) and on the test functions, but not on
the coupling. Under the reading of the two-point witness theorem on the data above, with
$m_0\ge m_{\Theta}(g)$, it satisfies
\begin{equation}\label{9.9:eq-tuned-family-3}
c_{\mathcal{K}}\,\vartheta^2\,\|f_1^{\circ}\|\,\|f_2^{\circ}\|
\;\le\;\mathcal{K}\;\le\;
100\,C_{\mathcal{K}}\,\|f_1^{\circ}\|\,\|f_2^{\circ}\| .
\end{equation}

\emph{The witness representation.} Put $b_0:=n_{\mathfrak{g}}/2=(N^2-1)/2$. Under the same reading,
with $m_0\ge m_{\Theta}(g)$, the two-point witness theorem reads
\begin{equation}\label{9.9:eq-tuned-family-4}
S^T_{2;K}(f_1^{\circ},f_2^{\circ})=N_{k_0}^2\,b_0\,g_{k_0}^4\,\mathcal{K}\,(1+R_K),
\qquad |R_K|\le\tfrac14 .
\end{equation}
Here $R_K$ is the relative error against the main term
$\mathfrak{M}:=N_{k_0}^2\,b_0\,g_{k_0}^4\,\mathcal{K}$. The difference
$S^T_{2;K}(f_1^{\circ},f_2^{\circ})-\mathfrak{M}$ is the sum of six members
\begin{equation*}
R^{\mathrm{prof}},\ R^{\mathrm{gx}},\ R^{\mathrm{irr}},\ R^{\mathrm{sh}},\ R^{\mathrm{bad}},\
R^{\mathrm{loc}} .
\end{equation*}
Each member is at most $\mathfrak{M}/32$ in absolute value, and $R^{\mathrm{sh}}=0$ for the
construction used in that theorem. Hence
$|S^T_{2;K}(f_1^{\circ},f_2^{\circ})-\mathfrak{M}|\le\tfrac{6}{32}\mathfrak{M}$. Only the
decomposition into the six members and the bound on each, not the bound $|R_K|\le\tfrac14$ itself,
is used in what follows. Finally, under the same reading, the main
term satisfies
\begin{equation}\label{9.9:eq-tuned-family-5}
\mathfrak{M}\ge c_M\,\vartheta^2\,g_{k_0}^4\,\|f_1^{\circ}\|\,\|f_2^{\circ}\|,
\qquad c_M:=\tfrac49\,b_0\,c_{\mathcal{K}} .
\end{equation}
\end{definition}

\begin{proof}
We give, in the order of their appearance, the statements from which the assertions above are taken.

\emph{The tuned family.} Let $g\le\gamma_U$. Clause (a) of the uniqueness theorem provides, for each
$n\ge0$, the real bare inverse coupling $a_n$ for which the run of length $n$ of the class $\Theta_X$
on $(m_0+n,a_n)$ has $x_{(0)}=X$. The depth-indexed notation $\zeta_{(r)}$, $x_{(r)}$ and the ladder
$\theta_r$ are those of that theorem.

\emph{The correlation functions.} The connected functions $S^n(x;f)$ are those of clause (c) of the
uniqueness theorem. The convergence of $S^n(a_n-\hat w;f)$ is its clause (d). This clause also gives
that the convergence holds without subsequences and uniformly on $\mathcal{N}'_w$, that the limit
$\hat w\mapsto S^{(\hat w)}$ is real-analytic on $]-\varrho_2/2,\varrho_2/2[$, and that it is
holomorphic on the strip.

\emph{The calibrating function.} The reference pair $f^{\circ}$ meets the hypotheses of clause (c)
with $\mathsf{p}=2$: both components are Lipschitz functions on the continuum torus of exponent
$m_0$, and their supports lie at mutual distance at least $\mathsf{d}^{\circ}>0$, hence are
separated. Hence $S^n(x;f^{\circ})$ is defined by clause (c), and clause (d) applies to
$f=f^{\circ}$. Clause (d), applied to
$f=f^{\circ}$, shows that $\varphi_n(\hat w)\to S^{(\hat w)}(f^{\circ})=\varphi(\hat w)$ uniformly on
$\mathcal{N}'_w$, without passage to subsequences. This is the limit written in
\eqref{9.9:eq-tuned-family-1}.

\emph{Steps and level.} On the bare datum $(m_0+n,a_n)$, $n$ renormalisation steps lead from the
bare lattice to the reference torus of exponent $m_0$. So the step count of the two-point witness theorem
is $K=n$, and depth $s^{\circ}$ of the run is the level $k_0=n-s^{\circ}$; see below, where the
two-point witness theorem is read on the bare data of the uniqueness theorem.

\emph{The window.} The window \eqref{9.9:eq-tuned-family-a} follows from the bound
\eqref{9.9:eq-coupling-convergence-a}, which expresses the convergence of the running coupling at each
depth along the tuned family, together with clause (a) of the uniqueness theorem. This derivation is
carried out below, where the two-point witness theorem is read on the bare data of the uniqueness
theorem. The pairwise increments of the run are bounded in the statement on the increments of the run
class, given below. Combined with the pin $x_{(0)}=X$, these bounds give
$x_{(r)}\ge X+\lambda r-2c_2$, with $\lambda$ the constant of that increment bound and with no
dependence on $g$.

\emph{The kernel.} By \eqref{9.9:eq-tuned-family-2}, $\mathcal{K}$ is formed from the values of
$f_1^{\circ}$ and $f_2^{\circ}$ at barycentres of plaquettes and from the covariance
$C^{\mathrm{bg}}_{K,s^{\circ}}$. That covariance is determined by the lattice ($K$, $s^{\circ}$, $L$,
the torus) and does not contain the coupling. The bounds \eqref{9.9:eq-tuned-family-3} are clause (e)
of the kernel bounds accompanying the two-point witness theorem, taken on the bare data
$(m_0+n,a_n)$ under the reading of that theorem on these data given below; the torus hypothesis of
the two-point witness theorem holds on the bare tori of exponent $m_0+n$ for every $n\ge0$ because
$m_0\ge m_{\Theta}(g)$, by the definition of $m_{\Theta}(g)$.

\emph{The witness representation.} Identity \eqref{9.9:eq-tuned-family-4} is the conclusion of the
two-point witness theorem under the reading, given below, of that theorem on the bare data
$(m_0+n,a_n)$ of the uniqueness theorem: its torus hypothesis holds on these data for every $n$
because $m_0\ge m_{\Theta}(g)$, and the window \eqref{9.9:eq-tuned-family-a} is used in place of the
first-passage envelope. Identity \eqref{9.9:eq-tuned-family-4} is asserted here only under that
reading. In that theorem, $R_K$ is defined by
$S^T_{2;K}(f_1^{\circ},f_2^{\circ})=\mathfrak{M}(1+R_K)$, so that
$\mathfrak{M}R_K=S^T_{2;K}(f_1^{\circ},f_2^{\circ})-\mathfrak{M}$. The decomposition of this
difference into the six members is the one of that theorem, with each member bounded by
$\mathfrak{M}/32$. It gives $|R_K|\le 6/32$, and $6/32<1/4$. The lower bound
\eqref{9.9:eq-tuned-family-5} on the main term, with $c_M=\tfrac49 b_0c_{\mathcal{K}}$, is the
main-term bound of the same theorem, under the same reading.
\end{proof}

\begin{definition}[The shifted main term and its remainder]\label{9.9:def-main-term}
Consider the tuned run of length $n$, with its scaffold
frozen, its level $k_0$, its inverse coupling $x_{k_0}=g_{k_0}^{-2}$ at that level, the witness
kernel $\mathcal{K}$, the calibrating function $\varphi_n$, the reference pair
$(f_1^{\circ},f_2^{\circ})\in\mathfrak{T}(\mathsf{d}^{\circ};\vartheta)$ and the witness main term
$\mathfrak{M}=N_{k_0}^2b_0g_{k_0}^4\mathcal{K}$, with $b_0=(N^2-1)/2$, all as in
Definition~\ref{9.9:not-tuned-family}. Let
$\zeta_{k_0}$ be the running shift at level $k_0$ and $\zeta'_{k_0}$ its derivative in the bare
offset $\hat w$, both as functions of $\hat w$ (Notation~\ref{9.9:not-running-shift}).
Let $r$ be the fixed value $\varrho_J^{0}$ of the source radius of the
correlation theorem, and write $D(0,2r)=\{\hat w\in\mathbb{C}:|\hat w|<2r\}$. For
$\hat w\in D(0,2r)$ put
\begin{equation}\label{9.9:eq-main-term-1}
  x_{k_0}(\hat w):=x_{k_0}-\zeta_{k_0}(\hat w),\qquad
  g_{k_0}(\hat w):=x_{k_0}(\hat w)^{-1/2},\qquad
  N_{k_0}(\hat w):=\zeta'_{k_0}(\hat w),
\end{equation}
and define the \emph{shifted main term} $\mathfrak{M}_n$ and its \emph{remainder} $E_n$ by
\begin{equation}\label{9.9:eq-main-term-a}
  \mathfrak{M}_n(\hat w):=N_{k_0}(\hat w)^2\, b_0\, x_{k_0}(\hat w)^{-2}\,\mathcal{K},
  \qquad
  E_n(\hat w):=\varphi_n(\hat w)-\mathfrak{M}_n(\hat w).
\end{equation}
With $g$ the coupling to which the family is tuned and $X=g^{-2}$
(Definition~\ref{9.9:not-tuned-family}), put $X^{\flat}:=X-2c_2$, and let $\sigma_0$ be the
constant of the correlation
theorem that fixes the radii of its polydiscs, proportional to $(LM)^{-4}$. Assume that
\begin{equation}\label{9.9:eq-main-term-2}
  \tfrac{16}{3}\,r\le \min\{\sigma_0,1\}\,X^{\flat},
\end{equation}
which is one of the conditions on $g$ in the theorem on non-degenerate calibration by the
curvature two-point witness, stated later in this section. Then:
\begin{enumerate}
\item[(a)] $x_{k_0}(\hat w)\neq0$ on $D(0,2r)$, and $\mathfrak{M}_n$ is holomorphic on $D(0,2r)$;
\item[(b)] $\varphi_n$ is holomorphic on a neighbourhood of $0$ that may depend on $n$, and,
uniformly in $n\ge n_0(N_w)$, the threshold of the holomorphy statement of the uniqueness
theorem, on the interior of the strip neighbourhood $\mathcal{N}'_w$; hence $E_n$ is holomorphic
on the intersection of each of these neighbourhoods of $0$ with $D(0,2r)$, on the second
uniformly in $n\ge n_0(N_w)$; no holomorphy of $\varphi_n$ on $D(0,2r)$ is asserted;
\item[(c)] at $\hat w=0$ one has $\mathfrak{M}_n(0)=\mathfrak{M}$, the witness main term
$N_{k_0}^2 b_0 g_{k_0}^4\mathcal{K}$, and
$E_n(0)=S^T_{2;n}(f_1^{\circ},f_2^{\circ})-\mathfrak{M}$.
\end{enumerate}
\end{definition}

\begin{proof}
(a) By the correlation theorem applied on the class $\mathfrak{S}$ of runs with frozen scaffold,
$\zeta_{k_0}$ is holomorphic on $D(0,2r)$, and on $D(0,2r)$ it satisfies
$|\zeta_{k_0}(\hat w)|\le \tfrac43\cdot 2r$. By the same theorem, since the inverse coupling of
the tuned run at depth $0$ is $X$, every inverse coupling of the tuned run at positive depth is
at least $X-2c_2$, while the one at depth $0$ is $X\ge X^{\flat}$; in particular
$X^{\flat}\le x_{k_0}$. Halving
\eqref{9.9:eq-main-term-2} gives $\tfrac83 r\le\tfrac12\min\{\sigma_0,1\}X^{\flat}$, so
\begin{equation*}
  |\zeta_{k_0}(\hat w)|\le \tfrac43\cdot 2r\le \tfrac12\min\{\sigma_0,1\}\,X^{\flat}< x_{k_0}.
\end{equation*}
For the last inequality: since $r>0$, \eqref{9.9:eq-main-term-2} forces $X^{\flat}>0$, whence
$\tfrac12\min\{\sigma_0,1\}X^{\flat}<X^{\flat}\le x_{k_0}$. Hence, for $\hat w\in D(0,2r)$,
\begin{equation*}
  \operatorname{Re} x_{k_0}(\hat w)\ge x_{k_0}-|\zeta_{k_0}(\hat w)|>0 .
\end{equation*}
In particular $x_{k_0}(\hat w)\neq0$. The principal branch of $x_{k_0}(\hat w)^{-1/2}$ is defined
on the right half-plane, so $g_{k_0}(\hat w)$ is well defined. Since $x_{k_0}(\hat w)$ is
holomorphic and does not vanish on $D(0,2r)$, the function $x_{k_0}(\hat w)^{-2}$ is holomorphic
there. The function $N_{k_0}=\zeta'_{k_0}$ is the derivative of the function $\zeta_{k_0}$,
holomorphic on $D(0,2r)$, and is therefore holomorphic there. With the factor
$N_{k_0}(\hat w)^2$ and the constants $b_0$ and $\mathcal{K}$, this gives
the holomorphy of $\mathfrak{M}_n$ on $D(0,2r)$ in \eqref{9.9:eq-main-term-a}.

(b) By Definition~\ref{9.9:not-tuned-family}, $\varphi_n$ is holomorphic on a
neighbourhood of $0$ that may depend on $n$. By the holomorphy statement of the uniqueness
theorem, it is holomorphic, uniformly in $n\ge n_0(N_w)$, on the interior of the strip
neighbourhood $\mathcal{N}'_w$. Since $\mathfrak{M}_n$ is holomorphic on $D(0,2r)\ni0$ by (a),
the difference $E_n=\varphi_n-\mathfrak{M}_n$ is holomorphic on the intersection of each of these
neighbourhoods of $0$ with $D(0,2r)$, on the second uniformly in $n\ge n_0(N_w)$. Nothing in
this argument concerns $\varphi_n$ on the whole of $D(0,2r)$.

(c) By Notation~\ref{9.9:not-running-shift}, $\zeta_{k_0}(0)=0$ and $\zeta'_{k_0}(0)=N_{k_0}$.
Hence the definitions \eqref{9.9:eq-main-term-1} and \eqref{9.9:eq-main-term-a}, together with
$x_{k_0}^{-2}=g_{k_0}^{4}$, give
\begin{equation*}
  \mathfrak{M}_n(0)=N_{k_0}^2\,b_0\,x_{k_0}^{-2}\,\mathcal{K}
  =N_{k_0}^2\,b_0\,g_{k_0}^{4}\,\mathcal{K}=\mathfrak{M}.
\end{equation*}
The calibrating function at $0$ is the connected
two-point function $S^T_{2;n}(f_1^{\circ},f_2^{\circ})$ of the reference pair along the run of
$n$ steps; the torus index is suppressed in this notation. Hence
$E_n(0)=S^T_{2;n}(f_1^{\circ},f_2^{\circ})-\mathfrak{M}$.
\end{proof}

In its own setting, the first-passage cutoff and the scaffold under which it is formulated, the
two-point witness theorem bounds the corresponding remainder by $\tfrac{6}{32}\mathfrak{M}$. The
tuned runs of Definition~\ref{9.9:not-tuned-family} are related to the runs of that setting only
through the corollary of the uniqueness theorem; for the tuned runs that bound would require a
transfer of the two-point witness theorem to the tuned family, which is not carried out here,
and the bound on $E_n(0)$ is not used in what follows. What is used
from the two-point witness theorem are its constants, its bounds on the kernel $\mathcal{K}$ and
its lower bound on the main term.

\begin{definition}[Homonymous letters and letters used by reference]\label{9.9:not-homonyms}
\emph{Letters with two meanings.} The following letters carry more than one meaning in this chapter. Each meaning is confined to the statements named beside it, and context decides which one is in force.
\begin{itemize}
\item[(a)] The letter $\rho$ has three meanings.
\begin{itemize}
\item In the witness data of Notation~\ref{9.9:not-tuned-family} and in the estimates of the two-point witness theorem, $\rho$ is the support radius, counted in cells of the witness lattice.
\item In the extraction theorem of the uniqueness argument and in the comparison of a pair of runs, $\rho\in\,]0,1[$ is the contraction rate.
\item Inside quotations from the correlation theorem, $\rho_{n,j}$ denotes a ratio; it occurs nowhere else.
\end{itemize}
\item[(b)] The letter $e$ has two meanings.
\begin{itemize}
\item $e_k$ is the function of the coupling used in the correlation theorem, as recalled in Notation~\ref{9.9:not-running-shift}.
\item In the comparison of a pair of runs (Definition~\ref{9.6:def-comparison-chain}), and only there, $e_r$ denotes
\begin{equation*}
\rho^{\,n-r}+\Pi_n ,
\end{equation*}
with $n$, $r$ and $\Pi_n$ as fixed there; $\Pi_n$ is the sum of the large-field floors $T_{\Pi}$ of the uniqueness theorem.
\end{itemize}
\item[(c)] The letter $p_0$ has two meanings.
\begin{itemize}
\item Without an argument, $p_0$ is the exponent that enters the function $T(g)$ of the correlation theorem.
\item With an argument, $p_0(\cdot)$ is the degree function, used as the threshold function of the two-point witness theorem.
\end{itemize}
\item[(d)] The letter $\mathsf{A}$ has two meanings.
\begin{itemize}
\item In the display of the extraction theorem of the uniqueness argument, $\mathsf{A}$ is the extraction constant.
\item In the comparison of a pair of runs, $\mathsf{A}$ is the shorter run of the pair; the longer run is $\mathsf{B}$, as in Definition~\ref{9.6:def-comparison-chain}.
\end{itemize}
\end{itemize}

\emph{Letters used by reference.} The following letters and statements are used in this chapter with the meaning fixed where they are introduced, and are not redefined here.
\begin{itemize}
\item[(e)] The coupling ceilings $\gamma_H,\gamma_J,\gamma_W,\gamma_U$ (Notation~\ref{0.2:not-glossary}).
\item[(f)] The floor $g_-(g)$ on the bare coupling and the floor $m_{\mathbb{T}}$ on the torus exponent
(Notation~\ref{0.2:not-glossary}).
\item[(g)] Letters of the uniqueness theorem:
\begin{itemize}
\item $\varrho_2:=r_J/2$, where $r_J$ is the parameter of the source domains $\mathcal{W}[r_J]$ of the correlation theorem;
\item the strip neighbourhood $\upsilon_w$ of real offsets;
\item the number $N_w$ of sources;
\item $n_0(N_w)$, as fixed in the uniqueness theorem.
\end{itemize}
\item[(h)] The floors and constants $k_W$, $s_*$, $C_0$, $C_D$, $r_0$ of the two-point witness theorem.
\item[(i)] The sourced representation at level $K-s$ given by the correlation theorem.
\item[(j)] The Gaussian extraction with two insertions.
\end{itemize}

\emph{Tower, bubble and triangle.} Here $C^{\mathrm{bg}}$ is the background covariance whose square occurs in the witness kernel $\mathcal{K}$ of the two-point witness theorem; it is used with the meaning fixed there. In the estimates of the two-point witness theorem we use three names.
\begin{itemize}
\item The \emph{tower} is the unit background Gaussian of the tower estimate with two insertions.
\item The \emph{bubble} is the square $[C^{\mathrm{bg}}]^2$ that occurs in $\mathcal{K}$.
\item The \emph{triangle} is the cyclic product of three factors $C^{\mathrm{bg}}$.
\end{itemize}
\end{definition}

\begin{definition}[Order of choice of the constants]\label{9.9:not-choice-of-constants}
The constants of this chapter are chosen in the following order. The outer quantifier prefix is
\begin{equation}\label{9.9:eq-choice-of-constants-1}
\exists L_0\;\;\forall\,\text{odd } L\ge L_0\;\;\exists\,\mathfrak{p}\;\;\exists\,\gamma\;\;
\exists\,\gamma_H ,
\end{equation}
where the auxiliary parameters $\mathfrak{p}$ are chosen in the order of
Notation~\ref{2.3:not-stages-first}, Notation~\ref{2.3:not-stages-middle} and
Notation~\ref{2.3:not-stages-last} (cf.\ Notation~\ref{2.1:not-prefix}).

Inside clauses (iii)(b)--(iii)(c) of Theorem~\ref{3.1:thm-main}, a radius
$\varrho_J^{0}=r_J^{0}>0$ is fixed first: this is the fixed value of the source radius
$\varrho_J$ of the correlation theorem. Then
\begin{equation}\label{9.9:eq-choice-of-constants-2}
\gamma_W(\dots;\varrho_J^{0})\le\gamma_J(\dots;\varrho_J^{0})\le\gamma_H .
\end{equation}
Here $\gamma_W$ is the coupling ceiling of the two-point witness theorem, and $\gamma_J$ is
that of the correlation theorem. Inside clause (iv) of Theorem~\ref{3.1:thm-main}, the ceiling
$\gamma_U$ of the uniqueness theorem satisfies $\gamma_U\le\gamma_H$.

In full, and in the order of the two-point witness theorem, the choices are made as follows.
\begin{enumerate}
\item The shape $\vartheta\in(0,1]$ is fixed at the same time as $\varrho_J^{0}$, hence before
the coupling ceilings displayed above. The witness theorem's $\gamma_W$, $C_0$, $C_W(s)$ and
$s_W(g)$ depend on $\vartheta$.
\item Then one further ceiling on $g$, the calibration ceiling $\gamma_{\mathrm{cal}}$, is
imposed:
\begin{equation}\label{9.9:eq-choice-of-constants-3}
\gamma_{\mathrm{cal}}=\gamma_{\mathrm{cal}}(L,\mathfrak{p},\gamma,\vartheta;\varrho_J^{0})
\le\min\{\gamma_W,\gamma_U\}.
\end{equation}
Each member of the chain of ceilings on $g$ is a one-sided condition on $g$ alone. Each holds
for every sufficiently small $g$.
\item Then the torus exponent is chosen.
\item Then the separation $\mathsf{d}^{\circ}$ is chosen, after $g$, subject to
$s(\mathsf{d}^{\circ})\ge s_W(g)$.
\item Then the pair of test functions of Notation~\ref{9.9:not-tuned-family} is chosen.
\item Finally the depth $n$ is chosen, from below.
\end{enumerate}
No cap is placed on $L$, $n$, $m_0$ or $s^{\circ}$ (the torus-exponent and depth letters of the
witness data of Notation~\ref{9.9:not-tuned-family}), and no floor is placed on $g$.
\end{definition}

\begin{proposition}[Bare-offset derivative of the witness remainder. Proof outstanding]
\label{9.9:prop-remainder-derivative}
All notation is that of Notation~\ref{9.9:not-tuned-family} and of
Definition~\ref{9.9:def-main-term}. In particular, the frozen-shifted main term
$\mathfrak{M}_n$ and the remainder $E_n=\varphi_n-\mathfrak{M}_n$ of the two-point witness
theorem along the tuned run of length $n$ are as in~\eqref{9.9:eq-main-term-a}.
There exist a coupling ceiling $\gamma_E$ and a constant $c_E$ with
\begin{equation*}
  \gamma_E\in\bigl(0,\min\{\gamma_W,\gamma_U\}\bigr], \qquad c_E<1,
\end{equation*}
both depending only on $L$, $\mathfrak{p}$, $\gamma$, $\vartheta$ and the fixed source radius
$\varrho_J^{0}$, with the following property. Let $g\le\gamma_E$, let $m_0\ge m_{\Theta}(g)$, let
$\mathsf{d}^{\circ}$ be a separation with $s(\mathsf{d}^{\circ})\ge s_W(g)$, let
$f^{\circ}=(f_1^{\circ},f_2^{\circ})\in\mathfrak{T}(\mathsf{d}^{\circ};\vartheta)$ be a reference
pair, and let $n\ge n_1$, where $n_1$ is the run-length threshold fixed for the tuned family
in this section. Then, along the frozen-shifted tuned family, the derivative of $E_n$ in the
bare offset $\hat{w}$ at $\hat{w}=0$ satisfies
\begin{equation}\label{9.9:eq-remainder-derivative-a}
  \bigl|E_n'(0)\bigr|\le c_E\,g_{k_0}^{2}\,\mathfrak{M}_n(0).
\end{equation}
The constant $c_E$ is one and the same for all depths, all reference pairs, all tori and
all $n\ge n_1$. Equivalently, with $S^T_{3;n}(1,f_1^{\circ},f_2^{\circ})$ the connected
three-point function of the run of $n$ steps with one constant source,
\begin{equation}\label{9.9:eq-remainder-derivative-1}
  \bigl|S^T_{3;n}(1,f_1^{\circ},f_2^{\circ})-\mathfrak{M}_n'(0)\bigr|
  \le c_E\,g_{k_0}^{2}\,\mathfrak{M}_n(0).
\end{equation}
\end{proposition}

Proof outstanding.

\smallskip
\noindent\textit{Route.}
A proof has to bound $E_n'(0)$ relative to $g_{k_0}^{2}\mathfrak{M}_n(0)$ with a constant that
is free of the depth $s^{\circ}=s(\mathsf{d}^{\circ})$; this uniformity in the depth is the point
at which every route except the sharp form~\eqref{9.9:eq-sufficient-conditions-c} strains: of the
displayed sufficient conditions only the sharp form is shown here to give a depth-free constant,
and a depth-free constant from the general condition~\eqref{9.9:eq-sufficient-conditions-a} needs
either a radius growing like $x_{k_0}$ or a relative precision decaying like $x_{k_0}^{-1}$ times
the radius. The obstruction is that the covariance comparison of the two-point
witness theorem bounds the remainder in absolute form, with no decay in the distance and no
validity at a complex base point. Either of the conditions
\eqref{9.9:eq-sufficient-conditions-a}, the two-point witness theorem at complex offsets with a
relation between its relative precision and its radius, or
\eqref{9.9:eq-sufficient-conditions-c}, its sharp form, which is the depth-uniform instance of
the first, implies~\eqref{9.9:eq-remainder-derivative-a}, and the proof of this implication is
given with them; \eqref{9.9:eq-sufficient-conditions-b} describes what a route inside the
real-measure covariance bookkeeping of the two-point witness theorem would need, and its
sufficiency is not proved.

As to the size of $c_E$: the coupling ceiling of the theorem that fixes the threshold $n_1$
gives $2c_N\le\tfrac13$, where $c_N$ is the constant of that theorem, and for the use
made of~\eqref{9.9:eq-remainder-derivative-a} in this section any constant
$c_E<\tfrac43-2c_N$ suffices; the proposition takes $c_E<1$.

\begin{definition}[Four hypotheses on the tuned family]\label{9.9:def-tuned-hypotheses}
Let $g\le\gamma_U$ and $X:=g^{-2}$. Let the ladder $(\theta_r)_{r\ge0}$, the class $\Theta_X$ of
runs tuned to $X$, the depth-indexed run variables $\zeta_{(r)}$, $b_{(r)}$, $x_{(r)}$, the tuned
bare inverse couplings $a_n$, the connected functions $S^n(x;f)$, the set $\mathcal{N}'_w$, the
reference pair $f^{\circ}=(f_1^{\circ},f_2^{\circ})\in\mathfrak{T}(\mathsf{d}^{\circ};\vartheta)$
and the depth $s^{\circ}=s(\mathsf{d}^{\circ})$ be as in Notation~\ref{9.9:not-tuned-family}. Let
the running shift, its increments and the constants $\lambda$, $c_2$, $\sigma_0$, $A$, $\kappa_0$,
$e_k$, $v_k$, $y_k$ and the source radius $\varrho_J$ of the correlation theorem be as in
Notation~\ref{9.9:not-running-shift}.

Write $\mathfrak{S}$ for the class of runs of the uniqueness theorem whose scaffold is frozen in
terms of $X$; the scaffold of a run is the collection of its thresholds, cube sizes,
decompositions, backgrounds and reference laws. For a run of this class,
$x_{(r)}=\theta_r-\zeta_{(r)}$ is its inverse coupling at depth $r$. The slot values of a run are
its deviations $\zeta_{(r)}$ from the ladder, and the ladder's slot origin is the point at which
all of them vanish. Write $\zeta_*^{n}$ for the bare deviation $\zeta_{(n)}$ from the ladder of
the tuned run of length $n$, that is, of the $\Theta_X$-run on the bare datum $(m_0+n,a_n)$, so
that $\zeta_{(r)}(\zeta_*^{n})$, $0\le r\le n$, are its slot values. Write $x^{(n)}_{(s)}$ for
the inverse coupling at depth $s$ of that run.

We say that the tuned family \emph{satisfies the four hypotheses on the tuned family} if
\textup{(a)}--\textup{(d)} below hold.

\begin{enumerate}
\item[(a)] \emph{Flow on the tuned runs}. This hypothesis concerns the $\Theta_X$-runs of the
class $\mathfrak{S}$, read about the tuned run's own slot values. The running shift in question
is the map
\begin{equation}\label{9.9:eq-tuned-hypotheses-1}
\zeta\longmapsto\zeta_{(\cdot)}(\zeta_*^{n}+\zeta)-\zeta_{(\cdot)}(\zeta_*^{n}),
\end{equation}
which is the response to a bare offset $\zeta$ at the tuned point. Let $x_j$ be the inverse
couplings of the run and $b_{k+1}$ the increments of the shift
\eqref{9.9:eq-tuned-hypotheses-1}, with levels numbered as in the correlation theorem. Then two
bounds hold.
\begin{enumerate}
\item[(a1)] The pairwise increment bound holds: for all levels $i<j$,
\begin{equation}\label{9.9:eq-tuned-hypotheses-2}
\lambda(j-i)-2c_2\;\le\;x_i-x_j\;\le\;3\lambda(j-i)+2c_2 .
\end{equation}
\item[(a2)] The running-shift statement holds with its fine bookkeeping. The increments
decompose as
\begin{equation}\label{9.9:eq-tuned-hypotheses-3}
b_{k+1}=\beta^0_{k+1}+f_{k+1}(\zeta_0,\dots,\zeta_k)+\mathsf{q}_{k+1}(\zeta),
\qquad
\mathsf{q}_{k+1}(\zeta):=\beta\bigl[Q^{(k+1)}_{\zeta}\bigr].
\end{equation}
The function $f_{k+1}$ is holomorphic on the polydisc
$\mathfrak{P}_k:=\prod_{j\le k}D(0,\sigma_0x_j)$, centred at the ladder's slot origin; we write
$\xi_0,\dots,\xi_k$ for the coordinates of $\mathfrak{P}_k$. On it
$|f_{k+1}|\le Ae_k$, and $f_{k+1}$ satisfies the derivative bound
\begin{equation}\label{9.9:eq-tuned-hypotheses-4}
\sup_{\frac12\mathfrak{P}_k}\;\sum_{j\le k}\sigma_0x_j
\Bigl|\frac{\partial f_{k+1}}{\partial\xi_j}\Bigr|\;\le\;\frac43\,Ae_k .
\end{equation}
The remnant $\mathsf{q}_{k+1}$ is holomorphic on $D(0,2\varrho_J)$ and satisfies
$|\mathsf{q}_{k+1}|\le v_k$ there.
\end{enumerate}
The derivative at $\zeta=0$ of \eqref{9.9:eq-tuned-hypotheses-1} has two readings. Read at
depth $0$, it is the quantity $N_n=\partial\zeta_{(0)}/\partial\zeta$ at the tuned point of the
uniqueness theorem. Read at depth $s^{\circ}$, it is $N_{k_0}(0)$, the factor of the main term of
Definition~\ref{9.9:def-main-term} at $\hat{w}=0$.

On $\mathfrak{S}$ the constants $\lambda,c_2,\sigma_0,A,\dots$ of the correlation theorem take
the values fixed for the uniqueness theorem. There $\mathsf{r}_k$ and $e_k$ are evaluated at the
lowered value of $\varepsilon_1$ used by the uniqueness theorem, for both runs of a pair, and
we write $\lambda^U$ for the value of $\lambda$ fixed for the uniqueness theorem.

\item[(b)] \emph{Holomorphy}. For $n\ge n_0(N_w)$, where the threshold $n_0(N_w)$ and
the index $N_w$ are those of the uniqueness theorem, the map
$\hat{w}\mapsto S^n(a_n-\hat{w};f)$ is holomorphic on the interior of $\mathcal{N}'_w$. The
family of these maps, $n\ge n_0(N_w)$, is uniformly bounded on that set.

\item[(c)] \emph{Common ladder}. In the run variables $x_{(r)}=\theta_r-\zeta_{(r)}$, the ladder
is a single sequence for the whole class $\Theta_X$, given by
\begin{equation}\label{9.9:eq-tuned-hypotheses-5}
\theta_0=X,\qquad \theta_{r+1}=\theta_r+c_0\quad(r\ge0).
\end{equation}
It is the same for the runs of lengths $n$ and $n+1$.

\item[(d)] \emph{Witness data on the tuned runs}. The notation of the two-point witness theorem, its kernel bounds
\textup{(a)} and \textup{(e)}, and its main-term lower bound
\begin{equation}\label{9.9:eq-tuned-hypotheses-6}
\mathfrak{M}\;\ge\;c_M\,\vartheta^2g^4\,\|f_1^{\circ}\|\,\|f_2^{\circ}\|
\end{equation}
are read on the bare data $(m_0+n,a_n)$ of the uniqueness theorem, as follows.
\begin{itemize}
\item The torus exponent is $m_0+n$ in bare units, with $n$ renormalisation steps above the
reference torus of exponent $m_0$. Hence the step count is $K:=n$ and the witness level is
$k_0:=n-s^{\circ}$.
\item In place of the first-passage envelope of the two-point witness theorem one uses the
window of the $\Theta_X$-run,
\begin{equation}\label{9.9:eq-tuned-hypotheses-7}
\underline{x}^{\Theta}_s:=X+\lambda^Us-c'\;\le\;x^{(n)}_{(s)}\;\le\;X+3\lambda^Us+c'
=:\bar{x}^{\Theta}_s .
\end{equation}
Here $c':=2c_2+\mathsf{A}C_E\bigl[(1-\rho)^{-2}+T_{\Pi}\bigr]$ is the constant of the tuning
clause of the uniqueness theorem. At $s=s^{\circ}$ the window
\eqref{9.9:eq-tuned-hypotheses-7} is \eqref{9.9:eq-tuned-family-a}.
\item The torus floor is $m_0\ge m_{\Theta}(g)$. Here $m_{\Theta}(g)$ is the least $m_0$ for which
the torus hypothesis of the two-point witness theorem holds on the bare tori of the uniqueness
theorem for all $n$.
\item The room $C_W(s)=C_0\,p_0(\bar{x}_s^{-1/2})^2$, the floors $s(\mathsf{d})$ and $s_W(g)$,
and the ratio $\bigl[p_0(g_{k_0})/p_0(\bar{x}_s^{-1/2})\bigr]^4$ in the remainder member
$R^{\mathrm{irr}}$ of the two-point witness theorem are taken with $\bar{x}^{\Theta}_s$ in place of
$\bar{x}_s=X+B^{\sharp}(s+1)$. Here $p_0(\cdot)$ is the degree function.
\end{itemize}
\end{enumerate}

\emph{On hypothesis \textup{(a)}.} For the runs of the correlation theorem itself, whose scaffold
is frozen at the unshifted run, the bounds \eqref{9.9:eq-tuned-hypotheses-2}--\eqref{9.9:eq-tuned-hypotheses-4}
and the bound on $\mathsf{q}_{k+1}$ are part of that theorem.

For the class $\mathfrak{S}$, the hypotheses of the uniqueness theorem on the flow display two
things. First, they display the running-shift statement in Lipschitz form, with Lipschitz
constants $\mathsf{a}_u$ satisfying
\begin{equation*}
\mathsf{a}_u\lesssim\theta_u^{-3/2}(\log\theta_u)^{p_0}+\theta_u^{-(\kappa_0-4)/2}/R,
\qquad \sum_u \mathsf{a}_u\le\tfrac13 ,
\end{equation*}
where $R$ is the constant of that name in those hypotheses. Second, they display the tube
\begin{equation*}
|x_{(r)}-\theta_r|\le\varrho_r\lesssim\theta_r^{1/2}(\log\theta_r)^{p_0}.
\end{equation*}
They display neither the polydisc radii $\sigma_0x_j$ with \eqref{9.9:eq-tuned-hypotheses-4}
nor a pairwise bound of the type \eqref{9.9:eq-tuned-hypotheses-2}; hypothesis (a) asks for
both. The tube is weaker than the constant $2c_2$ of \eqref{9.9:eq-tuned-hypotheses-2}. At the
cost of constants only, it would suffice for the anchoring step of the bound on the curvature of the running shift, and for the
lower bound $x_j\ge X^{\flat}$ with, say, $X^{\flat}:=X/2$.

By the tube, the tuned run's slot values lie at distance at most $\varrho_s$ from the ladder's
slot origin, where $\varrho_s$ is the tube radius at depth $s$. By the fifth of the conditions
defining the coupling ceiling $\gamma_{\mathrm{cal}}$, $\varrho_s\le\frac14\sigma_0x_s$. The disc
on which \eqref{9.9:eq-curvature-summation-a} is evaluated is taken with half its radius, namely
$t_k=\tfrac{3}{16}\sigma_0y_k$. Consequently the points used in the
curvature estimates of this section stay inside $\frac12\mathfrak{P}_k$, where
\eqref{9.9:eq-tuned-hypotheses-4} is stated (see the proof below). Halving the radius of the
disc of \eqref{9.9:eq-curvature-summation-a} doubles the level-$k$ constant of that display. Hence
it doubles the first term of the constant $c_N$, and this factor is absorbed in the ceiling imposed
on $c_N$.

The bound on the curvature of the running shift is written in the constants of the correlation
theorem and is read with the values fixed above. Suppose that on $\mathfrak{S}$ only the
Lipschitz form is assumed. Then that bound is replaced by the cruder curvature bound given in the
remark on the range and centring of the curvature bound, the remark containing
\eqref{9.9:eq-curvature-scope-a}, and a radius form of the running-shift statement, discussed in
the same remark, becomes necessary.

\emph{On hypothesis \textup{(b)}.} The display of the uniqueness theorem for $S^n$ at complex
offsets presupposes the holomorphy in (b). The mechanism behind (b) is the factorisation
$Z=e^{\mathcal{E}}\tilde{Z}$ in the hypotheses of the uniqueness theorem on the flow, with
$|\tilde{Z}|\le A$ and $\tilde{Z}(0)\ge 1/B$ on the source domain $\mathcal{W}[r_J]$ of the
correlation theorem, where $r_J=\varrho_J$, combined with Schwarz's lemma as in the proof of the
correlation theorem.

\emph{On hypothesis \textup{(c)}.} Two features of Notation~\ref{9.9:not-tuned-family} are
consistent with (c). The recursion $\zeta_{(r)}=\zeta_{(r+1)}+b_{(r)}-c_0$ subtracts the same
$c_0$ in the runs of lengths $n$ and $n+1$, and both runs have depth-zero inverse coupling
$x_{(0)}=X$. Hypothesis (c) asserts in addition that the ladder itself is common to the class,
with the form \eqref{9.9:eq-tuned-hypotheses-5}. The first equality in the convergence of the
running coupling at each depth, \eqref{9.9:eq-coupling-convergence-a}, rests on this hypothesis.

\emph{On hypothesis \textup{(d)}.} In the proof of the two-point witness theorem, the run enters
only through two properties. First, its couplings up to level $k_0$ lie in the regime of the
correlation theorem. Second, $x_{k_0}$ lies in a window of the affine shape
\eqref{9.9:eq-tuned-hypotheses-7} that does not depend on $K$. For a function $\Phi$ of the
inverse coupling and a constant $c$, each coupling ceiling in that proof has the form
$\sup_{y\ge X-c}\Phi(y)\le\mathrm{const}$. What that proof needs of $\bar{x}_s$ is only that it
bounds $x_{K-s}$ from above, independently of $K$; the room is taken at the envelope of the
witness scale, and never at $g$ in place of $g_{K-s}$. By \eqref{9.9:eq-tuned-hypotheses-7},
$\bar{x}^{\Theta}_s$ has this property on the data of the uniqueness theorem. The window of
clause (0) of Theorem~\ref{3.1:thm-main} concerns another run; it is not the window
\eqref{9.9:eq-tuned-hypotheses-7}.
\end{definition}

\begin{proof}[Proof of the inequalities stated with hypotheses \textup{(a)} and \textup{(d)}]
We prove three facts in turn: the inclusion into the half polydisc used in (a), the window
\eqref{9.9:eq-tuned-hypotheses-7} of (d), and the finiteness of the torus floor of (d).

\emph{The inclusion into $\frac12\mathfrak{P}_k$.} Fix $k$ and $j\le k$. The points at which
\eqref{9.9:eq-curvature-summation-a} evaluates $f_{k+1}$ have $j$-th coordinate $c_j+tN_j$. Here
$c_j$ is the $j$-th slot value of the tuned run, $t$ ranges over the disc of radius
$t_k=\frac{3}{16}\sigma_0y_k$, and $N_j$ is the derivative of the $j$-th component of the running
shift. We estimate the two terms separately.
\begin{itemize}
\item By the tube and the fifth condition defining $\gamma_{\mathrm{cal}}$,
$|c_j|\le\varrho_j\le\frac14\sigma_0x_j$, where $\varrho_j$ denotes the tube radius at the depth
of level $j$.
\item The derivative bound of Notation~\ref{9.9:not-running-shift}, read for the shift
\eqref{9.9:eq-tuned-hypotheses-1}, gives $|N_j|\le\frac43$. Therefore
$|tN_j|\le\frac{3}{16}\cdot\frac43\,\sigma_0y_k=\frac14\sigma_0y_k$.
\end{itemize}
Since $y_k=\min_{i\le k}x_i-2c_2$ with $c_2\ge0$, we have $y_k\le x_j$ for $j\le k$. Combining,
\begin{equation*}
|c_j|+|tN_j|\;\le\;\tfrac14\sigma_0x_j+\tfrac14\sigma_0y_k\;\le\;\tfrac12\sigma_0x_j .
\end{equation*}
This holds for every $j\le k$, so the point lies in
$\frac12\mathfrak{P}_k=\prod_{j\le k}D(0,\frac12\sigma_0x_j)$. This is the set on which
\eqref{9.9:eq-tuned-hypotheses-4} is stated.

\emph{The window \eqref{9.9:eq-tuned-hypotheses-7}.} Fix $s\le n$ and consider the
$\Theta_X$-run of length $s$. Its depth $s$ is its bare scale, so its inverse coupling at depth
$s$ equals its bare inverse coupling, namely $x^{(s)}_{(s)}=a_s$.

We compare the depth-$s$ couplings of runs of consecutive lengths $m$ and $m+1$, for
$m=s,\dots,n-1$. The convergence of the running coupling at each depth,
\eqref{9.9:eq-coupling-convergence-a}, bounds each such change by $\tau_m(s)$. Writing
$x^{(n)}_{(s)}-a_s$ as the telescoping sum of these changes, we obtain
\begin{equation}\label{9.9:eq-tuned-hypotheses-8}
\bigl|x^{(n)}_{(s)}-a_s\bigr|\;\le\;\sum_{m\ge s}\tau_m(s)
\;\le\;\mathsf{A}C_E\bigl[\rho(1-\rho)^{-2}+T_{\Pi}\bigr]\;\le\;c'-2c_2 .
\end{equation}
The second inequality in \eqref{9.9:eq-tuned-hypotheses-8} is the summed form of the same
convergence statement. The third holds because $0\le\rho<1$ gives
$\rho(1-\rho)^{-2}\le(1-\rho)^{-2}$, because $\mathsf{A}C_E\ge0$, and because
$c'-2c_2=\mathsf{A}C_E[(1-\rho)^{-2}+T_{\Pi}]$ by the definition of $c'$.

The tuning clause of the uniqueness theorem gives
\begin{equation}\label{9.9:eq-tuned-hypotheses-9}
\lambda^Us-2c_2\;\le\;a_s-X\;\le\;3\lambda^Us+2c_2 .
\end{equation}
Now write
\begin{equation*}
x^{(n)}_{(s)}-X=\bigl(x^{(n)}_{(s)}-a_s\bigr)+\bigl(a_s-X\bigr).
\end{equation*}
By \eqref{9.9:eq-tuned-hypotheses-8} the first bracket lies in $[-(c'-2c_2),\,c'-2c_2]$. By
\eqref{9.9:eq-tuned-hypotheses-9} the second lies in $[\lambda^Us-2c_2,\,3\lambda^Us+2c_2]$.
Adding the two ranges gives
\begin{equation*}
\lambda^Us-c'\;\le\;x^{(n)}_{(s)}-X\;\le\;3\lambda^Us+c' ,
\end{equation*}
which is \eqref{9.9:eq-tuned-hypotheses-7}. The bounds depend on $s$ and $X$ but not on $n$.

\emph{The torus floor.} The torus hypothesis of the two-point witness theorem is the condition
$m\ge m_{\mathbb{T}}(g_-(g))$ on the exponent $m$ of the reference-scale torus. On the bare data
$(m_0+n,a_n)$ of the uniqueness theorem the reference torus has exponent $m_0$ for every $n$.
The condition therefore reads $m_0\ge m_{\mathbb{T}}(g_-(g))$, which does not involve $n$.

Hence every $m_0\ge m_{\mathbb{T}}(g_-(g))$ satisfies the torus hypothesis on all the bare tori
$(m_0+n)$, $n\ge0$. The set of integers $m_0$ for which the hypothesis holds for all $n$ is
therefore non-empty, and it has a least element. That least element is $m_{\Theta}(g)$, which is
finite and satisfies $m_{\Theta}(g)\le m_{\mathbb{T}}(g_-(g))$.
\end{proof}

\begin{theorem}[The running coupling converges at each depth]\label{9.9:thm-coupling-convergence}
Let $g\le\gamma_U$, and let $X$, the class $\mathfrak{S}$ of runs with frozen scaffold, the class
$\Theta_X$ of runs tuned to $X$ with its reference ladder, and the tuned bare inverse couplings $a_n$
be as in Notation~\ref{9.9:not-tuned-family}. We write a run of $\Theta_X$ in the depth-indexed
notation of the uniqueness theorem. A run of length $n$ with bare deviation $\zeta$ has
\begin{equation}\label{9.9:eq-coupling-convergence-1}
\zeta_{(n)}:=\zeta,\qquad \zeta_{(r)}:=\zeta_{(r+1)}+b_{(r)}-c_0,\qquad x_{(r)}:=\theta_r-\zeta_{(r)},
\end{equation}
where $b_{(r)}$ is the coupling increment of the run at depth $r$ and $(\theta_r)_{r\ge0}$ is the
reference ladder of the class $\Theta_X$. Depth $r=0$ is the reference scale and depth $r=n$ is the
bare scale. For $n\ge 0$ and $0\le t\le n$ let $x^{(n)}_{(t)}$ denote the inverse coupling
$x_{(t)}$, at depth $t$, of the run of $\Theta_X$ of length $n$ whose bare inverse coupling is
$a_n$; thus $x^{(n)}_{(0)}=X$ for every $n$, by the tuning of $a_n$.

Assume the following.
\begin{enumerate}
\item[(a)] \emph{The comparison proposition for two matched runs of this chapter}, under the
hypotheses under which it is stated there: the input it takes by statement from the correlation
theorem, the transplanted one-lattice response bounds with their inputs, the contraction hypothesis,
and the ceilings on the coupling. Consider a pair of $\Theta_X$-runs of the class $\mathfrak{S}$
(scaffold frozen in terms of $X$), compared as in
Definition~\ref{9.6:def-comparison-chain}, with the shorter run of length $n$ and the longer run of
length $n+1$. The pair is
parametrised through the matching map $\mathfrak{n}_n$ of the uniqueness theorem: the shorter run is
taken at a parameter $v$ and the longer run at $\mathfrak{n}_n(v)$, and the two runs are matched at
depth $0$ for every $v\in D(0,r_J)$, where $r_J>0$ is the radius of the parameter disc of the pair.
Label the two runs by the superscripts $\mathsf{A}$ (shorter)
and $\mathsf{B}$ (longer), and put
\begin{equation}\label{9.9:eq-coupling-convergence-2}
d_s:=\sup_{v\in\overline{D(0,r_J)}}\bigl|\zeta^{\mathsf{B}}_{(s)}-\zeta^{\mathsf{A}}_{(s)}\bigr|,
\qquad s\le n .
\end{equation}
If the pair is alive (in the sense of the uniqueness theorem) to a depth $r_0$ and matched there,
then for $r_0\le r<n$
\begin{equation}\label{9.9:eq-coupling-convergence-3}
d_{r+1}\le \mathsf{A}\,C_E\Bigl[\frac{\rho^{\,n-r}}{1-\rho}+(r-r_0+1)\,\Pi_n\Bigr],
\end{equation}
where $\mathsf{A}$, $C_E$ and $\rho$ are the constants of the uniqueness theorem, with
$\rho\in\left]0,1\right[$, and $\Pi_n$ is its large-field sum. (The factor $\mathsf{A}$ in front of
$C_E$ is the extraction constant; the superscript $\mathsf{A}$ only labels the shorter run.)
\item[(b)] \emph{Tuning and matching at depth zero}: for the pairs of \textup{(a)}, the tuning
condition $x_{(0)}=X$ and the matching of the two runs are both imposed at depth $0$ of
$\mathfrak{S}$, so that the comparison in \textup{(a)} applies with $r_0=0$.
\item[(c)] The statement of the uniqueness theorem that the matching map fixes the origin:
$\mathfrak{n}_n(0)=0$.
\item[(d)] The bound on $\Pi_n$ displayed in the uniqueness theorem,
\begin{equation}\label{9.9:eq-coupling-convergence-4}
\Pi_n\le C\,n^{-(\kappa_0-4)/2}\,(\log n)^{c_3},
\end{equation}
with constants $C$ and $c_3$, where $\kappa_0$ is the constant of the correlation theorem denoted by
that letter there, which that theorem takes with $\kappa_0\ge7$.
\item[(e)] The statement of the uniqueness theorem on its constants: the quantities $\mathsf{A}$,
$C_E$, $\rho$, $\Pi_n$, $\Upsilon_n$, and every constant entering their bounds, are free of the torus
exponent; moreover the constants $\mathsf{A}$, $C_E$, $\rho$ and the constants $C$, $c_3$,
$\kappa_0$ of \eqref{9.9:eq-coupling-convergence-4} do not depend on the cutoff $K$, on the length
$n$ of the runs, or on the depth $r$.
\item[(f)] The ladder hypothesis of Definition~\ref{9.9:def-tuned-hypotheses}: the ladder
$(\theta_r)_{r\ge0}$ is one sequence for the whole
class $\Theta_X$, the same for the runs of lengths $n$ and $n+1$, with
\begin{equation}\label{9.9:eq-coupling-convergence-5}
\theta_0=X,\qquad \theta_{r+1}=\theta_r+c_0 .
\end{equation}
\end{enumerate}
Then for every fixed depth $t\ge1$
\begin{equation}\label{9.9:eq-coupling-convergence-a}
\bigl|x^{(n+1)}_{(t)}-x^{(n)}_{(t)}\bigr|\le \tau_n(t)
:=\mathsf{A}\,C_E\Bigl[\frac{\rho^{\,n-t+1}}{1-\rho}+t\,\Pi_n\Bigr]
\qquad\text{for } n\ge t\ge1 .
\end{equation}
Since $(\kappa_0-4)/2>1$ as soon as $\kappa_0>6$ (the correlation theorem takes $\kappa_0\ge7$, and the
proofs of the uniqueness theorem use $\kappa_0\ge10$), one has $\sum_{n\ge t}\tau_n(t)<\infty$, the
limit
\begin{equation*}
x_{\infty,(t)}:=\lim_{n\to\infty}x^{(n)}_{(t)}
\end{equation*}
exists for every $t\ge1$, and
\begin{equation}\label{9.9:eq-coupling-convergence-6}
\bigl|x^{(n)}_{(t)}-x_{\infty,(t)}\bigr|\le\sum_{n''\ge n}\tau_{n''}(t)\qquad(n\ge t).
\end{equation}
At depth $t=0$ the sequence is constant, $x^{(n)}_{(0)}=X$, by the tuning of $a_n$.
\end{theorem}

\begin{proof}
\emph{The tuned run in the depth-indexed notation.} The depth-indexed quantities of a run are those of
\eqref{9.9:eq-coupling-convergence-1}. By hypothesis \textup{(f)} the ladder $(\theta_r)$ in the last
of these definitions depends on the depth $r$ alone. It is therefore the same for every run of
$\Theta_X$, and in particular for both members of a pair as in hypothesis \textup{(a)}, and both
runs subtract the same $c_0$ at each step of \eqref{9.9:eq-coupling-convergence-1}.

By the tuning clause that the uniqueness theorem takes by statement from the correlation theorem,
for each $n$ exactly one real bare deviation whose run of length $n$ is alive, written
$\zeta_*^{n}(0)$, produces a run of length $n$ with $\zeta_{(0)}=0$, and the tuned bare inverse
coupling of Notation~\ref{9.9:not-tuned-family} is
\begin{equation*}
a_n=\theta_n-\zeta_*^{n}(0).
\end{equation*}
Hence the tuned run of length $n$ is the run with $\zeta_{(0)}=0$. By
\eqref{9.9:eq-coupling-convergence-1} its bare inverse coupling is
\begin{equation*}
x_{(n)}=\theta_n-\zeta_{(n)}=\theta_n-\zeta_*^{n}(0)=a_n .
\end{equation*}
Its reference-scale inverse coupling is
$x_{(0)}=\theta_0-\zeta_{(0)}=\theta_0$, and $\theta_0=X$ by \eqref{9.9:eq-coupling-convergence-5};
this is the tuning condition $x^{(n)}_{(0)}=X$ of the tuned family, which gives the last sentence of
the theorem.

\emph{The tuned pair.} Take the pair of hypothesis \textup{(a)}: the shorter run $\mathsf{A}$ of length
$n$ at the parameter $v$ and the longer run $\mathsf{B}$ of length $n+1$ at $\mathfrak{n}_n(v)$,
matched at depth $0$ for every $v\in D(0,r_J)$, with $d_s$ given by
\eqref{9.9:eq-coupling-convergence-2} for $s\le n$. At $v=0$ the shorter run $\mathsf{A}$ is the run
from the bare deviation $\zeta_*^{n}(0)$, whose bare inverse coupling is
$\theta_n-\zeta_*^{n}(0)=a_n$; by hypothesis \textup{(c)}, $\mathfrak{n}_n(0)=0$, so the longer run
$\mathsf{B}$ is taken at the parameter $0$ as well and is the
run from $a_{n+1}$. At $v=0$ the pair is therefore the pair of tuned runs of lengths $n$
and $n+1$, and for $0\le s\le n$
\begin{equation*}
x^{(n)}_{(s)}=\theta_s-\zeta^{\mathsf{A}}_{(s)}\big|_{v=0},\qquad
x^{(n+1)}_{(s)}=\theta_s-\zeta^{\mathsf{B}}_{(s)}\big|_{v=0}.
\end{equation*}
The same $\theta_s$ appears in both identities by hypothesis \textup{(f)}.

\emph{The increment bound.} By hypothesis \textup{(b)}, the tuning condition and the matching are
imposed at depth $0$ of $\mathfrak{S}$, so the comparison of hypothesis \textup{(a)} applies with
$r_0=0$. For $0\le r<n$,
\eqref{9.9:eq-coupling-convergence-3} with $r_0=0$ reads
\begin{equation*}
d_{r+1}\le\mathsf{A}\,C_E\Bigl[\frac{\rho^{\,n-r}}{1-\rho}+(r+1)\,\Pi_n\Bigr].
\end{equation*}
Put $t:=r+1$. As $r$ ranges over $0\le r<n$, $t$ ranges over $1\le t\le n$, and $n-r=n-t+1$.
Let $n\ge t\ge1$. The point $v=0$ lies in $\overline{D(0,r_J)}$, so the value at $v=0$ of
$|\zeta^{\mathsf{B}}_{(t)}-\zeta^{\mathsf{A}}_{(t)}|$ is at most the supremum $d_t$. With the
identities of the previous paragraph we obtain
\begin{equation}\label{9.9:eq-coupling-convergence-7}
\begin{aligned}
\bigl|x^{(n+1)}_{(t)}-x^{(n)}_{(t)}\bigr|
&=\Bigl|\bigl(\theta_t-\zeta^{\mathsf{B}}_{(t)}\bigr)
   -\bigl(\theta_t-\zeta^{\mathsf{A}}_{(t)}\bigr)\Bigr|_{v=0}
 \le d_t\\
&\le\mathsf{A}\,C_E\Bigl[\frac{\rho^{\,n-t+1}}{1-\rho}+t\,\Pi_n\Bigr]=\tau_n(t).
\end{aligned}
\end{equation}
The first equality is the cancellation of the common $\theta_t$ supplied by hypothesis \textup{(f)}.
This is \eqref{9.9:eq-coupling-convergence-a}.

\emph{Summability.} Fix $t\ge1$. By hypothesis \textup{(e)}, $\mathsf{A}$, $C_E$, $\rho$ and the
constants $C$, $c_3$, $\kappa_0$ of \eqref{9.9:eq-coupling-convergence-4} do not depend on $n$, so
$\tau_n(t)$ is, for fixed $t$, a sequence in $n$ with $n$-independent coefficients. We sum its two
parts separately.

For the first part, since $0<\rho<1$ by hypothesis \textup{(a)},
\begin{equation*}
\sum_{n\ge t}\frac{\rho^{\,n-t+1}}{1-\rho}=\frac{1}{1-\rho}\sum_{j\ge1}\rho^{\,j}
=\frac{\rho}{(1-\rho)^2}<\infty .
\end{equation*}

For the second part, with $\kappa_0\ge10$, as used in the proofs of the uniqueness theorem, the
exponent $(\kappa_0-4)/2$ is at least $3$, hence greater than $1$; more generally
$(\kappa_0-4)/2>1$ whenever $\kappa_0>6$, that is, $(\kappa_0-6)/4>0$. Since
$(\log n)^{c_3}\le n^{(\kappa_0-6)/4}$ for all sufficiently large $n$, the bound
\eqref{9.9:eq-coupling-convergence-4} gives, for those $n$,
\begin{equation*}
\Pi_n\le C\,n^{-(\kappa_0-4)/2+(\kappa_0-6)/4}=C\,n^{-(\kappa_0-2)/4}.
\end{equation*}
Here $(\kappa_0-2)/4>1$, so $\sum_n\Pi_n<\infty$, and hence $\sum_{n\ge t}t\,\Pi_n<\infty$.
Together, $\sum_{n\ge t}\tau_n(t)<\infty$.

\emph{Convergence and the tail bound.} Let $n'>n\ge t$. Every index $n''$ with $n\le n''\le n'-1$
satisfies $n''\ge t\ge1$, so \eqref{9.9:eq-coupling-convergence-a} applies to each increment, and the
triangle inequality gives
\begin{equation*}
\bigl|x^{(n')}_{(t)}-x^{(n)}_{(t)}\bigr|
\le\sum_{n''=n}^{n'-1}\bigl|x^{(n''+1)}_{(t)}-x^{(n'')}_{(t)}\bigr|
\le\sum_{n''=n}^{n'-1}\tau_{n''}(t)\le\sum_{n''\ge n}\tau_{n''}(t).
\end{equation*}
The right-hand side is the tail of a convergent series, so it tends to $0$ as $n\to\infty$. Hence
$(x^{(n)}_{(t)})_{n\ge t}$ is a Cauchy sequence of real numbers, and its limit $x_{\infty,(t)}$
exists. Letting $n'\to\infty$ in the last display, with $n$ fixed, gives
\eqref{9.9:eq-coupling-convergence-6}.
\end{proof}

\begin{remark}[The two-run estimate as a boundary-value problem]\label{9.9:rem-two-run-boundary}
Consider two runs of consecutive lengths, $\mathsf{A}$ of length $n$ and $\mathsf{B}$ of length $n+1$, compared as in
Definition~\ref{9.6:def-comparison-chain} and matched at a depth $r^{\circ}$. This remark describes
the structure of the proof of the two-run estimate of the uniqueness theorem.

That proof is a two-point boundary-value argument for a pair of families of difference sequences, both indexed
by the depth $r$. The first family consists of the irrelevant differences $g^P_r$, $g^A_r$, $g^R_r$. The second
consists of the marginal difference $d_r$. The sequences are coupled in both directions.

\emph{From the deeper and marginal differences to the irrelevant ones.} Clause~(a) of the contraction
hypothesis on the irrelevant output differences bounds the irrelevant output difference of the step at depth $r$
by the sum of three contributions.
\begin{itemize}
\item The first is a cross term $\alpha_r d^{\sharp}_{r+1}$, which comes from the marginal difference.
\item The second is a fixed constant times the sum, over depths $s>r$, of $q^{\,s-r}$ times the irrelevant
differences at depth $s$. The deeper differences are thus weighted by powers of the kernel rate $q<1$; this term
is the contraction.
\item The third consists of large-field terms, carried with the allowances $\mathfrak{r}_r\Lambda_r$. These
allowances are a power law in the running coupling and are not contracted.
\end{itemize}
The hypothesis on the newly born terms bounds their two-run difference at depth $r$ by
$C_{\sigma}e_r$, with $C_{\sigma}$ the constant of that hypothesis. Here
\begin{equation}\label{9.9:eq-two-run-boundary-1}
e_r:=\rho^{\,n-r}+\Pi_n ,
\end{equation}
with $\rho\in\,]0,1[$ and the large-field floor $\Pi_n$ of the uniqueness theorem. The letter $e_r$ is local to
this remark and to the two-run estimate; it is not the quantity $e_k$ of the correlation theorem.

\emph{From the irrelevant differences to the marginal one.} Clause~(b) of the contraction hypothesis on the
irrelevant output differences (the hypothesis whose clause~(a) was used above) is an exact identity for the
increments of the two runs at depth $s$:
\begin{equation}\label{9.9:eq-two-run-boundary-2}
b^{\mathsf{B}}_s-b^{\mathsf{A}}_s=\beta^{\mathsf{A}}[G^P_s]+\beta^{\mathsf{A}}[G^Q_s],
\qquad
\bigl|b^{\mathsf{B}}_s-b^{\mathsf{A}}_s\bigr|\le\mathsf{A}\,\|G^P_s\| .
\end{equation}
Here $\beta^{\mathsf{A}}[\cdot]$ is the increment functional of the shorter run, applied to the two terms
$G^P_s$ and $G^Q_s$ of that clause. In \eqref{9.9:eq-two-run-boundary-2}, and in
\eqref{9.9:eq-two-run-boundary-3} and \eqref{9.9:eq-two-run-boundary-5} below, the prefactor $\mathsf{A}$ is
the extraction constant of the uniqueness theorem (the constant $\mathsf{A}$ of
Theorem~\ref{9.9:thm-coupling-convergence}), not the run $\mathsf{A}$. Since the irrelevant difference
$g^P_u$ bounds $\|G^P_u\|$, summing the inequality in \eqref{9.9:eq-two-run-boundary-2} over depths gives the
marginal line
\begin{equation}\label{9.9:eq-two-run-boundary-3}
d_{s+1}\le\mathsf{A}\sum_{r^{\circ}\le u\le s}g^P_u ,
\qquad d_{r^{\circ}}=0 .
\end{equation}
In \eqref{9.9:eq-two-run-boundary-3}, $r^{\circ}$ is the depth at which the two runs are matched, and the
condition $d_{r^{\circ}}=0$ is the pin. This line is the cross term from the irrelevant differences to the
marginal one.

The increment also depends directly on the marginal coordinate, as described by the running-shift statement of
the correlation theorem. We read this dependence as entering through the slots of the step's data that carry
the running shift $\zeta$, and hence as contained in the term $\alpha_r d^{\sharp}_{r+1}$ above.

\emph{Boundary data.} The boundary data are given at the two ends.
\begin{itemize}
\item At depth $r^{\circ}$ the datum is the pin $d_{r^{\circ}}=0$.
\item At the bare end the datum is
\begin{equation*}
g^P_n\le C_V\le C_E\,e_n ,
\end{equation*}
which uses $e_n=1+\Pi_n\ge1$; here $C_V$ is the bound on the bare-end irrelevant difference and $C_E$ is the
constant of the two-run estimate, the constant $C_E$ of Theorem~\ref{9.9:thm-coupling-convergence}. This is
an irrelevant difference of order one, not a small one.
\end{itemize}

\emph{Closing the system.} The maximum principle for coupled scale-indexed inequalities
(Lemma~\ref{9.6:lem-maximum-principle}) closes the coupled system. That lemma does not require the system to be
triangular. The hypothesis of the two-run estimate that places ceilings on the coupling $g$ makes the linear
operator of the system a half-contraction, that is, a contraction of norm at most one half.

\emph{Content of the estimate.} Take two consecutive cutoffs tuned to equal reference coupling. For such a pair
the pin and the match sit at depth $0$, so
$r^{\circ}=0$, and the cone condition of clause~\textup{(v)} of Theorem~\ref{3.1:thm-main} is the content of
the contraction hypothesis on the irrelevant output differences and the hypothesis on the newly born terms
taken together. The conclusion of the two-run estimate then consists of two bounds.

The irrelevant differences at depth $r$ satisfy
\begin{equation}\label{9.9:eq-two-run-boundary-4}
g^P_r,\ g^A_r,\ g^R_r\le C_E\bigl(\rho^{\,n-r}+\Pi_n\bigr).
\end{equation}
This bound has two parts. The first is a geometric gain from the bare end, at the kernel rate $q\le\rho^2$. The
second is the power-law large-field floor $\Pi_n$.

The marginal differences satisfy, for $0\le r<n$, with the pin at depth $r^{\circ}=0$,
\begin{equation}\label{9.9:eq-two-run-boundary-5}
d_{r+1}\le\mathsf{A}\,C_E\Bigl[\frac{\rho^{\,n-r}}{1-\rho}+(r+1)\,\Pi_n\Bigr].
\end{equation}
The bound \eqref{9.9:eq-two-run-boundary-5} follows from the marginal line \eqref{9.9:eq-two-run-boundary-3}
with $s=r$ and $r^{\circ}=0$ once each $g^P_u$ is bounded by \eqref{9.9:eq-two-run-boundary-4}: indeed,
since $0<\rho<1$,
\begin{align*}
\sum_{u=0}^{r}\bigl(\rho^{\,n-u}+\Pi_n\bigr)
&=\rho^{\,n-r}\sum_{j=0}^{r}\rho^{\,j}+(r+1)\,\Pi_n\\
&\le\frac{\rho^{\,n-r}}{1-\rho}+(r+1)\,\Pi_n .
\end{align*}

\emph{Scope.} The estimate compares runs of consecutive lengths $n$ and $n+1$. It is stated for two runs of
arbitrary lengths only by telescoping in $n$. It says nothing about two continuum actions with equal running
couplings $g_s$ that are not members of the tuned family.
\end{remark}

\begin{remark}\label{9.9:prop-shift-curvature}
The statement of this proposition is not written out here; parts of its proof are printed below.
\end{remark}

\begin{proof}[Proof of Proposition~\ref{9.9:prop-shift-curvature}, continued]
\label{9.9:prf-curvature-summation}
We keep $k<k_0$ fixed as above and estimate, one after the other, the three terms on the right of
the corresponding display which are not $\zeta_k''(0)$; we then sum over $k$. Recall that
$\tfrac12\mathfrak{P}_k=\prod_{j\le k}D(0,\tfrac12\sigma_0x_j)$, that $f_{k+1}$ is holomorphic on
$\mathfrak{P}_k\supset\tfrac12\mathfrak{P}_k$ with $|f_{k+1}|\le Ae_k$, and that the refined form
of the running-shift statement of the correlation theorem, which splits $b_{k+1}$ into, among other
terms, its pure-channel part $f_{k+1}$ and the remnant $\mathsf{q}_{k+1}$, contains the derivative
bound
\begin{equation*}
  \sup_{\frac12\mathfrak{P}_k}\ \sum_{j\le k}\sigma_0x_j\,
  \Bigl|\frac{\partial f_{k+1}}{\partial\xi_j}\Bigr|\le\tfrac43\,Ae_k
\end{equation*}
and the remnant value bound
$\sup_{D(0,2\varrho_J)}|\mathsf{q}_{k+1}|\le v_k$, where
$v_k=2AC_tC_{\mathfrak{m}}\,y_k^{-(\kappa_0-4)/2}$. We write $\vec N=(N_0,\dots,N_k)$ with
$N_j=\zeta_j'(0)\in[\tfrac23,\tfrac43]$.

\smallskip\noindent$(\alpha)$ The quadratic form. Put
\begin{equation*}
  \Phi(t):=\frac{d}{dt}f_{k+1}(t\vec N)=\sum_{j\le k}(\partial_jf_{k+1})(t\vec N)\,N_j ,
  \qquad t_k:=\tfrac38\,\sigma_0y_k .
\end{equation*}
For $|t|<t_k$ and $j\le k$ we have, since $N_j\le\tfrac43$ and, by the corresponding display,
$y_k<\min_{i\le k}x_i\le x_j$,
\begin{equation*}
  |tN_j|<\tfrac43\cdot\tfrac38\,\sigma_0y_k=\tfrac12\sigma_0y_k<\tfrac12\sigma_0x_j ,
\end{equation*}
so $t\vec N\in\tfrac12\mathfrak{P}_k$; in particular $\Phi$ is holomorphic on $D(0,t_k)$. Using
$N_j\le\tfrac43$, then $\sigma_0x_j\ge\sigma_0\min_{i\le k}x_i$ together with the derivative bound
above, and finally $\min_{i\le k}x_i>y_k$ (from the corresponding display) and
$e_k\le e(y_k^{-1/2})$, we obtain for $|t|<t_k$
\begin{equation*}
\begin{aligned}
  |\Phi(t)|
  &\le\tfrac43\sum_{j\le k}\bigl|(\partial_jf_{k+1})(t\vec N)\bigr|
   \le\tfrac43\,\bigl(\sigma_0\min_{j\le k}x_j\bigr)^{-1}\cdot\tfrac43\,Ae_k\\
  &\le\frac{16}{9}\,\frac{A\,e(y_k^{-1/2})}{\sigma_0y_k}.
\end{aligned}
\end{equation*}
By the chain rule $\Phi'(0)=\sum_{i,j\le k}(\partial_i\partial_jf_{k+1})(0)N_iN_j$, and the Cauchy
inequality for the first derivative on the disc $D(0,t_k)$, for which $1/t_k=8/(3\sigma_0y_k)$, gives
\begin{equation}\label{9.9:eq-curvature-summation-a}
\begin{aligned}
  t_k&=\tfrac38\,\sigma_0y_k,\\
  \Bigl|\sum_{i,j\le k}(\partial_i\partial_jf_{k+1})(0)N_iN_j\Bigr|
  =|\Phi'(0)|
  &\le\frac{16}{9}\,\frac{A\,e(y_k^{-1/2})}{\sigma_0y_k}\cdot\frac{8}{3\sigma_0y_k}\\
  &=\frac{128}{27}\,\frac{A\,e(y_k^{-1/2})}{\sigma_0^2y_k^2}=:a_k .
\end{aligned}
\end{equation}

\smallskip\noindent$(\beta)$ The memory term. Since $0\in\tfrac12\mathfrak{P}_k$, the derivative
bound at the origin, with $\min_{j\le k}x_j=y_k+2c_2$ from the corresponding display and
$e_k\le e(y_k^{-1/2})$, gives
\begin{equation*}
  \sum_{j\le k}\bigl|(\partial_jf_{k+1})(0)\bigr|
  \le\frac{4Ae_k}{3\sigma_0(y_k+2c_2)}
  \le\ell_k:=\frac{4A}{3\sigma_0}\,\frac{e(y_k^{-1/2})}{y_k}.
\end{equation*}
Hence, with $M_k:=\max_{j\le k}|\zeta_j''(0)|$,
\begin{equation*}
  \Bigl|\sum_{j\le k}(\partial_jf_{k+1})(0)\,\zeta_j''(0)\Bigr|\le\ell_kM_k .
\end{equation*}

\smallskip\noindent$(\gamma)$ The remnant. The function $\mathsf{q}_{k+1}$ is holomorphic on
$D(0,2\varrho_J)$ and bounded there by $v_k$. The Cauchy inequality for the second derivative on this
disc gives
\begin{equation*}
  |\mathsf{q}_{k+1}''(0)|\le\frac{2v_k}{(2\varrho_J)^2}=\frac{v_k}{2\varrho_J^2}=:q_k .
\end{equation*}

\smallskip\noindent$(\delta)$ Summation. Inserting $(\alpha)$, $(\beta)$ and $(\gamma)$ into
the corresponding display yields
$|\zeta_{k+1}''(0)|\le M_k+a_k+\ell_kM_k+q_k$; since also $M_k\le M_k(1+\ell_k)$, we get
\begin{equation*}
  M_{k+1}\le M_k(1+\ell_k)+a_k+q_k,\qquad M_0=0 ,
\end{equation*}
where $M_0=|\zeta_0''(0)|=0$ because $\zeta_0=\zeta$ is the identity. By induction on $k$,
\begin{equation*}
  M_k\le\Bigl[\prod_{i<k}(1+\ell_i)\Bigr]\sum_{i<k}(a_i+q_i) ;
\end{equation*}
indeed, if this holds at $k$, then multiplying by $1+\ell_k\ge1$ and adding $a_k+q_k$, which is at
most $\prod_{i<k+1}(1+\ell_i)\,(a_k+q_k)$, gives it at $k+1$.

In the correlation theorem the function $\psi:=e\,g^2$ of the coupling and the numbers
$\theta_i'=\tfrac{16A}{9\sigma_0}\,\psi(y_i^{-1/2})$ are fixed, and $d_i(\varrho)\ge\theta_i'$, in
particular $d_i(\varrho_J)\ge\theta_i'$. Since
$\psi(y_i^{-1/2})=e(y_i^{-1/2})/y_i$,
\begin{equation*}
  \theta_i'=\frac{16A}{9\sigma_0}\,\frac{e(y_i^{-1/2})}{y_i}=\tfrac43\,\ell_i .
\end{equation*}
Hence, by the bound $\sum_{k<K}d_k(\varrho_J)\le\tfrac13$ of the correlation theorem, assumed in
Proposition~\ref{9.9:prop-shift-curvature},
\begin{equation*}
  \sum_{i<K}\ell_i\le\tfrac34\sum_{i<K}d_i(\varrho_J)\le\tfrac14,
  \qquad
  \prod_{i<K}(1+\ell_i)\le\exp\Bigl(\sum_{i<K}\ell_i\Bigr)\le e^{1/4} .
\end{equation*}
Since $k_0\le K$, the product over $i<k_0$ is bounded by the same $e^{1/4}$, and therefore
\begin{equation*}
  |\zeta_{k_0}''(0)|\le M_{k_0}\le e^{1/4}\sum_{k<k_0}(a_k+q_k).
\end{equation*}
We now use the corresponding display, which gives $y_k\ge\hat{g}^{-2}+\lambda\nu$ with
$\nu=k_0-k\in[1,k_0]$, and the monotonicity of $e$, which gives $e(y_k^{-1/2})\le e(\hat{g})$. Both
summands below are decreasing functions of $\nu$, so their sum over $\nu\ge1$ is at most their
integral over $[0,\infty)$. For the quadratic form,
\begin{equation*}
\begin{aligned}
  \sum_{k<k_0}a_k
  &\le\frac{128}{27}\,\frac{A}{\sigma_0^2}\,e(\hat{g})\sum_{\nu\ge1}(\hat{g}^{-2}+\lambda\nu)^{-2}
   \le\frac{128}{27}\,\frac{A}{\sigma_0^2}\,e(\hat{g})
     \int_0^\infty(\hat{g}^{-2}+\lambda t)^{-2}\,dt\\
  &=\frac{128}{27}\,\frac{A}{\lambda\sigma_0^2}\,e(\hat{g})\,\hat{g}^2 .
\end{aligned}
\end{equation*}
For the remnant,
\begin{equation*}
  q_k=\frac{v_k}{2\varrho_J^2}=\frac{AC_tC_{\mathfrak{m}}}{\varrho_J^2}\,y_k^{-(\kappa_0-4)/2},
\end{equation*}
and since $(\kappa_0-4)/2>1$ for $\kappa_0>6$,
\begin{equation*}
\begin{aligned}
  \sum_{k<k_0}q_k
  &\le\frac{AC_tC_{\mathfrak{m}}}{\varrho_J^2}
     \sum_{\nu\ge1}(\hat{g}^{-2}+\lambda\nu)^{-(\kappa_0-4)/2}
   \le\frac{AC_tC_{\mathfrak{m}}}{\varrho_J^2}
     \int_0^\infty(\hat{g}^{-2}+\lambda t)^{-(\kappa_0-4)/2}\,dt\\
  &=\frac{2AC_tC_{\mathfrak{m}}}{\lambda(\kappa_0-6)\varrho_J^2}\,\hat{g}^{\kappa_0-6} .
\end{aligned}
\end{equation*}
Multiplying the sum of these two bounds by $e^{1/4}$ gives exactly
$|\zeta_{k_0}''(0)|\le\mathfrak{z}(x_{k_0})$, the first bound in
the corresponding display.

Since $g_{k_0}^2=x_{k_0}^{-1}$ and $N_{k_0}\ge\tfrac23$,
\begin{equation*}
  |\zeta_{k_0}''(0)|
  \le\Bigl[\frac{x_{k_0}\mathfrak{z}(x_{k_0})}{N_{k_0}}\Bigr]N_{k_0}g_{k_0}^2
  \le\tfrac32\,x_{k_0}\mathfrak{z}(x_{k_0})\cdot N_{k_0}g_{k_0}^2
  =c_N(x_{k_0})\,N_{k_0}g_{k_0}^2 .
\end{equation*}

Monotonicity and the limit. Writing $\hat{g}(x)=(x-4c_2)^{-1/2}$, we have
\begin{equation*}
  x\,\mathfrak{z}(x)=e^{1/4}\Bigl[\frac{128}{27}\,\frac{A}{\lambda\sigma_0^2}\,
    e(\hat{g}(x))\,x\hat{g}(x)^2
   +\frac{2AC_tC_{\mathfrak{m}}}{\lambda(\kappa_0-6)\varrho_J^2}\,x\hat{g}(x)^{\kappa_0-6}\Bigr].
\end{equation*}
Here $x\hat{g}(x)^2=x/(x-4c_2)$ is decreasing and at most $2$ for $x\ge8c_2$; $e(\hat{g}(x))$ is
non-increasing in $x$, because $\hat{g}(x)$ decreases in $x$ and $e$ is nondecreasing on
$(0,\gamma_e]$; and for $\kappa_0\ge8$
\begin{equation*}
  \frac{d}{dx}\Bigl[x(x-4c_2)^{-(\kappa_0-6)/2}\Bigr]
  =(x-4c_2)^{-(\kappa_0-4)/2}\Bigl[(x-4c_2)-x\,\frac{\kappa_0-6}{2}\Bigr]\le0 .
\end{equation*}
So $c_N(x)=\tfrac32\,x\,\mathfrak{z}(x)$ is non-increasing on $x\ge8c_2+4$ with
$(x-4c_2)^{-1/2}\le\gamma_e$, the range of the run's inverse couplings. With
\begin{equation*}
  x\hat{g}(x)^{\kappa_0-6}=x\hat{g}(x)^2\,\hat{g}(x)^{\kappa_0-8}\le2\hat{g}(x)^{\kappa_0-8}
\end{equation*}
and $x\hat{g}(x)^2\le2$ in the first term, the displayed bound on $c_N(x)$ in
the corresponding display follows. For $\kappa_0>8$ both of its terms tend to $0$ as
$x\to\infty$: $\hat{g}(x)\to0$, hence $\hat{g}(x)^{\kappa_0-8}\to0$, and
$e(\hat{g}(x))\to0$ because, by its definition in the correlation theorem,
$e(g)=2E^{*}\tau(g)+\vartheta_e(g)\to0$ as $g\downarrow0$.
\end{proof}

\begin{remark}[Range and centring of the curvature bound]\label{9.9:rem-curvature-scope}
Throughout this remark, $c_N$, $\mathfrak{z}$ and
$\hat{g}(x)=(x-4c_2)^{-1/2}$ are as in Proposition~\ref{9.9:prop-shift-curvature}, as are
the exponent $\kappa_0$ and the constants $A$, $\sigma_0$, $\lambda$, $C_t$,
$C_{\mathfrak{m}}$ of the correlation theorem, and
$\varrho_J$ is the source radius of the correlation theorem. The \emph{remnant term} is
the second summand in the bound on $c_N$ of Proposition~\ref{9.9:prop-shift-curvature},
the summand that carries $\varrho_J^{-2}\,\hat{g}(x)^{\kappa_0-8}$. We write
$\operatorname{polylog}(x)$ for a factor bounded by a fixed power of $\log x$.
\begin{enumerate}
\item[(1)] For $\kappa_0\in\{7,8\}$ the remnant term does not tend to $0$ as
$x\to\infty$ at fixed $\varrho_J$. For $\kappa_0=8$ it can be made small by fixing
$\varrho_J$ large before $\gamma_J$; this is the order in which these two constants
are fixed in the proof of the uniqueness theorem. In either case it can be made small
by reading
the correlation theorem at a source radius proportional to $X/\operatorname{polylog}(X)$,
where $X=g^{-2}$, provided that theorem admits such a radius. Among the conditions of
the correlation theorem that couple its source radius $r$ to $g$ are
\begin{equation*}
\tfrac{16}{3}\,r\le\sigma_0\,g^{-2},\qquad
\tfrac{8}{3}\,r\,\varepsilon_k^{2}\le 1,
\end{equation*}
where $\varepsilon_k$ is the level-$k$ parameter of that name fixed in the correlation
theorem,
together with finitely many conditions of the form
$r\,g^{2}\operatorname{polylog}(g^{-2})\le C$ with $C$ a constant. Whether the remaining
conditions of that theorem which involve the source radius restrict $r$ further is not
verified here. Neither device is needed in the regime of the uniqueness theorem.
\item[(2)] A cruder route estimates the coupling increment directly by Cauchy's
estimate on the disc of the correlation theorem. It gives for $\zeta''_{k_0}(0)$ only a
bound of order $x_{k_0}^{-1/2}\operatorname{polylog}(x_{k_0})/\varrho_J$. This bound is
smaller than $g_{k_0}^{2}=x_{k_0}^{-1}$ only if $\varrho_J$ is taken at least of order
$x_{k_0}^{1/2}\operatorname{polylog}(x_{k_0})$. The fine bookkeeping of the running-shift
statement avoids this.
\item[(3)] \emph{Which run, and about which point.}
Proposition~\ref{9.9:prop-shift-curvature} is a statement about a run for which the
pairwise increment bounds and the running-shift statement of the correlation theorem,
with its fine bookkeeping, hold about that run. For the first-passage runs of the
correlation theorem they hold by that theorem, about the origin $\zeta=0$ of the
level variables, which is the run itself. For the tuned run of the class $\Theta_X$ of
the uniqueness theorem they hold by the hypothesis of
Definition~\ref{9.9:def-tuned-hypotheses} that transfers these two statements of the
correlation theorem to the runs of the class $\mathfrak{S}$, read about the values
$\zeta_*^{n}$ of the
level variables at the tuned run. The family of running shifts to which
Proposition~\ref{9.9:prop-shift-curvature} is then applied is the recentred family
\begin{equation}\label{9.9:eq-curvature-scope-a}
\zeta\longmapsto \zeta_{(\cdot)}\bigl(\zeta_*^{n}+\zeta\bigr)-\zeta_{(\cdot)}\bigl(\zeta_*^{n}\bigr),
\end{equation}
with the polydisc recentred accordingly. Where it is applied to the tuned run, which
has $n$ steps, it is
applied at the level $k_0=n-s^{\circ}$, with $s^{\circ}$ the depth at which the
uniqueness theorem uses it. There $x_{k_0}$ is the tuned run's inverse
coupling at depth $s^{\circ}$, which lies in the window for that coupling with ends
$\underline{x}^{\Theta}_{s^{\circ}}$ and $\bar{x}^{\Theta}_{s^{\circ}}$, and
$N_{k_0}=\zeta'_{k_0}(0)$ is the derivative at depth $s^{\circ}$ of the recentred
running shift~\eqref{9.9:eq-curvature-scope-a}. The quantity $N_n$ of the uniqueness
theorem, the derivative $\partial\zeta_{(0)}/\partial\zeta$ at the tuned point, is the
instance of this derivative at depth $0$.
\end{enumerate}
\end{remark}

\begin{proof}
\emph{Item} (1). The remnant term of the bound on $c_N(x)$ in
Proposition~\ref{9.9:prop-shift-curvature} is the expression
\begin{equation*}
\frac{3}{2}\,e^{1/4}\,\frac{4AC_tC_{\mathfrak{m}}}{\lambda(\kappa_0-6)\varrho_J^{2}}\,
\hat{g}(x)^{\kappa_0-8}
=\frac{6\,e^{1/4}AC_tC_{\mathfrak{m}}}{\lambda(\kappa_0-6)\varrho_J^{2}}\,
\hat{g}(x)^{\kappa_0-8},
\end{equation*}
with the constants $A$, $\lambda$, $C_t$, $C_{\mathfrak{m}}$ of the correlation theorem.
This expression is meaningful for every $\kappa_0>6$; we evaluate it below also for
$\kappa_0\in\{7,8\}$, although Proposition~\ref{9.9:prop-shift-curvature} is stated only
for $\kappa_0>8$.
As $x\to\infty$ we have $x-4c_2\to\infty$, and hence $\hat{g}(x)\to0$.

For $\kappa_0=8$ we have $\kappa_0-6=2$ and $\hat{g}(x)^{0}=1$, so the remnant term equals
\begin{equation*}
\frac{3\,e^{1/4}AC_tC_{\mathfrak{m}}}{\lambda\,\varrho_J^{2}}.
\end{equation*}
This quantity does not depend on $x$, so it does not tend to $0$ as $x\to\infty$. It is
this $x$-independent multiple of $\varrho_J^{-2}$ that fixing $\varrho_J$ large, before
$\gamma_J$, makes small.

For $\kappa_0=7$ the factor $\hat{g}(x)^{\kappa_0-8}$ equals
$\hat{g}(x)^{-1}=(x-4c_2)^{1/2}$. This factor tends to infinity as $x\to\infty$, so at
fixed $\varrho_J$ the remnant term does not tend to $0$ either.

Now let $\varrho_J=cX/\operatorname{polylog}(X)$ with a constant $c>0$, so that
$\varrho_J^{-2}=c^{-2}X^{-2}\operatorname{polylog}(X)^{2}$. For $\kappa_0=8$ the remnant
term is then, for every $x$,
\begin{equation*}
\frac{3\,e^{1/4}AC_tC_{\mathfrak{m}}}{\lambda\,c^{2}}\,X^{-2}\operatorname{polylog}(X)^{2},
\end{equation*}
which tends to $0$ as $X\to\infty$. For $\kappa_0=7$ we have $\kappa_0-6=1$, and the
remnant term is
\begin{equation*}
\frac{6\,e^{1/4}AC_tC_{\mathfrak{m}}}{\lambda\,c^{2}}\,(x-4c_2)^{1/2}\,
X^{-2}\operatorname{polylog}(X)^{2};
\end{equation*}
at levels whose inverse coupling $x$ is comparable to $X$ it is of order
$X^{-3/2}\operatorname{polylog}(X)^{2}$, and it tends to $0$ as $X\to\infty$.

For $\kappa_0>8$, which is the hypothesis of Proposition~\ref{9.9:prop-shift-curvature},
the exponent $\kappa_0-8$ is positive. Hence $\hat{g}(x)^{\kappa_0-8}\to0$, and the
remnant term tends to $0$ at fixed $\varrho_J$. This is the case in which that
proposition is applied to the runs of the uniqueness theorem, and there neither device
of item~(1) is used.

\emph{Item} (2). Put $h_k:=b_{k+1}-b_{k+1}(0)$, where the coupling increment $b_{k+1}$
is regarded as a function of the source variable. Cauchy's estimate on the disc of the
correlation theorem gives
\begin{equation*}
|h_k''(0)|\le \frac{d_k(2\varrho_J)}{\varrho_J},
\end{equation*}
with $d_k(\cdot)$ the level-$k$ quantity of the correlation theorem that enters its
Cauchy estimate for $h_k$ on its source disc, and which satisfies
$\sum_{k<K}d_k(\varrho_J)\le\frac13$. Hence
\begin{equation*}
|\zeta''_{k_0}(0)|\le \frac{\mathfrak{d}_W(x_{k_0})}{\varrho_J},
\end{equation*}
where $\mathfrak{d}_W(x_{k_0})$ is the level-$k_0$ Lipschitz sum of the lemma on Lipschitz
sums in the proof of the two-point witness theorem. This sum is of order
$x_{k_0}^{-1/2}\operatorname{polylog}(x_{k_0})$. The
resulting bound is therefore of order
$x_{k_0}^{-1/2}\operatorname{polylog}(x_{k_0})/\varrho_J$. Its ratio to
$g_{k_0}^{2}=x_{k_0}^{-1}$ is of order
$x_{k_0}^{1/2}\operatorname{polylog}(x_{k_0})/\varrho_J$. This ratio is small only if
$\varrho_J$ is of order $x_{k_0}^{1/2}\operatorname{polylog}(x_{k_0})$ or larger.

The fine bookkeeping of the running-shift statement splits the source-dependent part
$h_k$ of the coupling increment into
two parts. The first is the pure channel, holomorphic on the polydisc $\mathfrak{P}_k$
of radii $\sigma_0x_j$, $j\le k$. The second is the remnant $\mathsf{q}_{k+1}$, bounded
on $D(0,2\varrho_J)$ and of size $g^{\kappa_0-4}$. The radii $\sigma_0x_j$ grow with the
inverse couplings, so the pure channel's contribution does not involve $\varrho_J$. Only
the remnant sees $\varrho_J$, and it does so through the factor $\varrho_J^{-2}$ in the
second summand of $\mathfrak{z}$, which carries $\hat{g}^{\kappa_0-6}$.

\emph{Item} (3). The proof of Proposition~\ref{9.9:prop-shift-curvature} is written for a
family $\zeta_j(\zeta)$ with $\zeta_j(0)=0$. The family~\eqref{9.9:eq-curvature-scope-a}
takes at $\zeta=0$ the value
$\zeta_{(\cdot)}(\zeta_*^{n})-\zeta_{(\cdot)}(\zeta_*^{n})=0$, so it is of that
form. The hypotheses of Proposition~\ref{9.9:prop-shift-curvature} are read on the
polydisc recentred at $\zeta_*^{n}$; this is the content of the hypothesis on the tuned
family in Definition~\ref{9.9:def-tuned-hypotheses}. None of the steps
$(\alpha)$--$(\delta)$ of the proof of Proposition~\ref{9.9:prop-shift-curvature} uses
the position of the centre. So that proof applies verbatim to the recentred family.

By the chain rule, the derivative of~\eqref{9.9:eq-curvature-scope-a} at $\zeta=0$ is
the derivative of $\zeta_{(\cdot)}$ at $\zeta_*^{n}$. In a run of $n$ steps, depth $s$
is level $n-s$. Read at depth $s^{\circ}$, that is at level $k_0=n-s^{\circ}$, this
derivative is $N_{k_0}=\zeta'_{k_0}(0)$. Read at depth $0$, it is
$\partial\zeta_{(0)}/\partial\zeta$ at the tuned point, which is $N_n$.
\end{proof}

\begin{theorem}[Non-degenerate calibration by the two-point witness]
\label{9.9:thm-nondegenerate-calibration}
Let $\kappa_0>8$, where $\kappa_0$ is the exponent of the uniqueness theorem that enters
Proposition~\ref{9.9:prop-shift-curvature}; in the regime of the uniqueness theorem one has
$\kappa_0\ge 10$. Assume the following.
\begin{itemize}
\item The four hypotheses on the tuned family of Definition~\ref{9.9:def-tuned-hypotheses}. The
first concerns the $\Theta_X$-runs of the class $\mathfrak{S}$ of the uniqueness theorem. On these
runs it asks for the pairwise increment bounds of the correlation theorem,
\begin{equation*}
\lambda(j-i)-2c_2\;\le\; x_i-x_j\;\le\; 3\lambda(j-i)+2c_2\qquad(i<j),
\end{equation*}
and for its running-shift statement with its fine bookkeeping; here $\lambda$ and $c_2$ are the
flow constants of the correlation theorem. The second asserts that, for $n$ at least the index
$n_0(N_w)$ of clause~(c) of the uniqueness theorem, the Schwinger functions of the tuned family
are holomorphic and uniformly bounded in the bare offset $\hat{w}$ on the interior of the strip
neighbourhood $\mathcal{N}'_w$ of the uniqueness theorem. The third asserts that the ladder
$\theta_r$ is one sequence for the whole class $\Theta_X$, the same for the runs of lengths $n$
and $n+1$. The fourth reads the two-point witness theorem on the bare data of the uniqueness
theorem, with the window \eqref{9.9:eq-tuned-family-a} and the torus floor $m_{\Theta}(g)$.
\item The bound of Proposition~\ref{9.9:prop-remainder-derivative}, whose proof is outstanding,
with its ceiling $\gamma_E$ and its constant $c_E<1$.
\item Clause~(a) of the uniqueness theorem, together with the proposition of the uniqueness
theorem on which, with that clause, the upper end $\bar{x}^{\Theta}_{s^{\circ}}$ of the window
\eqref{9.9:eq-tuned-family-a} rests; both are used for that upper end. In the fourth hypothesis of
Definition~\ref{9.9:def-tuned-hypotheses} this upper end is obtained by summing the bounds of
Theorem~\ref{9.9:thm-coupling-convergence} over the run lengths from $s^{\circ}$ on and applying
clause~(a) of the uniqueness theorem to the run of length $s^{\circ}$.
\end{itemize}
The conclusions also use Proposition~\ref{9.9:prop-shift-curvature}, with its non-increasing
function $c_N$ of the corresponding display, and clauses (c) and (d) of the uniqueness
theorem.
Fix $\vartheta\in\,]0,1]$. Then there is a ceiling
$\gamma_{\mathrm{cal}}=\gamma_{\mathrm{cal}}(L,\mathfrak{p},\gamma,\vartheta;\varrho_J^{0})$ with
\begin{equation*}
0<\gamma_{\mathrm{cal}}\le\min\{\gamma_W,\gamma_U,\gamma_E\}.
\end{equation*}
Here $\gamma_W=\gamma_W(L,\mathfrak{p},\gamma,\vartheta;\varrho_J^{0})$ is the ceiling of the
two-point witness theorem and $\gamma_E$ that of Proposition~\ref{9.9:prop-remainder-derivative}.
The ceiling is such that every $g\le\gamma_{\mathrm{cal}}$ satisfies the following five conditions,
where $X:=g^{-2}$ and $X^{\flat}:=g^{-2}-2c_2$:
\begin{enumerate}
\item $X^{\flat}\ge 8c_2+4$;
\item $\tfrac{16}{3}\varrho_J^{0}\le\min\{\sigma_0,1\}\,X^{\flat}$, where $\sigma_0$ is the
constant of the correlation theorem that fixes the radii of its polydisc;
\item $(X^{\flat}-4c_2)^{-1/2}\le\gamma_e$, where $\gamma_e$ is the largest coupling below which the
function $e(\cdot)$ of the correlation theorem is nondecreasing. Since $e(\cdot)$ is nondecreasing
on couplings $g'$ with $\log g'^{-2}>2p_0$, this condition holds as soon as
$\log(X^{\flat}-4c_2)>2p_0$.
\item $c_N(X^{\flat})\le\tfrac16$;
\item $\varrho_s\le\tfrac14\sigma_0(\theta_s-2c_2)$ for all $s\ge0$. Here $\varrho_s$ is the tube
radius of the uniqueness theorem at depth $s$ and $\theta_s$ its ladder.
\end{enumerate}
These are conditions on $g$ alone. They read no depth, cutoff or torus, and each holds for all
sufficiently small $g$.

The number $X^{\flat}$ is a lower bound for every inverse coupling of the $\Theta_X$-run at depth at
least one. Indeed, the pairwise increment bounds, read against the pin $x_{(0)}=X$, give
$x_{(r)}\ge X+\lambda r-2c_2$ for every depth $r\ge1$. The constant $2c_2$ is that of the
correlation theorem and is fixed before $g$.

Condition~5 serves the recentred application of Proposition~\ref{9.9:prop-shift-curvature}
described in Remark~\ref{9.9:rem-curvature-scope}, \eqref{9.9:eq-curvature-scope-a}. It states that
the tube radius at every depth is at most a quarter of the polydisc radius there. It holds for small
$g$ because $\varrho_s\lesssim\theta_s^{1/2}(\log\theta_s)^{p_0}$ is among the inputs of the
uniqueness theorem and $\theta_s\ge X$.

Now take any $g\le\gamma_{\mathrm{cal}}$ and any torus exponent $m_0\ge m_{\Theta}(g)$. Then take any
separation $\mathsf{d}^{\circ}$ with $s^{\circ}:=s(\mathsf{d}^{\circ})\ge s_W(g)$; there is no upper
restriction on $s^{\circ}$. Take any reference pair
$f^{\circ}=(f_1^{\circ},f_2^{\circ})\in\mathfrak{T}(\mathsf{d}^{\circ};\vartheta)$. Let $n_1$ be the
least integer $n$ such that every $n'\ge n$ satisfies
\begin{equation*}
n'\ge n_0(N_w),\qquad n'-s^{\circ}\ge k_W,\qquad a_{n'}^{-1/2}\le g_{-}(g),
\end{equation*}
with $n_0(N_w)$ the index of clause~(c) of the uniqueness theorem. The integer $n_1$ is finite,
since clause~(a) of the uniqueness theorem gives the lower window
$a_{n'}\ge X+\lambda^{U}n'-2c_2$, which increases in $n'$; here $\lambda^{U}$ is the rate $\lambda$
of the correlation theorem as re-valued for the runs of the class $\mathfrak{S}$.

For every $n\ge n_1$ put $k_0:=n-s^{\circ}$, let $\varphi_n$ be as in
Notation~\ref{9.9:not-tuned-family}, and let $N_{k_0}=N_{k_0}(0)$, $g_{k_0}=g_{k_0}(0)$ and
$\mathfrak{M}_n(0)$ be as in Definition~\ref{9.9:def-main-term}. Then:
\begin{equation}\label{9.9:eq-nondegenerate-calibration-a}
\varphi_n'(0)=S^T_{3;n}(1,f_1^{\circ},f_2^{\circ}),
\end{equation}
the connected three-point function of $n$ steps at bare inverse coupling $a_n$;
\begin{equation}\label{9.9:eq-nondegenerate-calibration-b}
\begin{split}
\varphi_n'(0)&\ge\bigl(\tfrac43-2c_N(X^{\flat})-c_E\bigr)\,g_{k_0}^{2}\,\mathfrak{M}_n(0)\\
&\ge c_{\mathrm{cal}}\,\vartheta^{2}\,\|f_1^{\circ}\|_{\infty}\|f_2^{\circ}\|_{\infty}\,
\bigl(\bar{x}^{\Theta}_{s^{\circ}}\bigr)^{-3}>0,
\end{split}
\end{equation}
with
\begin{equation*}
c_{\mathrm{cal}}:=\bigl(\tfrac43-2c_N(X^{\flat})-c_E\bigr)\cdot\tfrac49\,b_0c_{\mathcal{K}}
\;\ge\;(1-c_E)\,\tfrac49\,b_0c_{\mathcal{K}}.
\end{equation*}
Here $c_N(x_{k_0})\le c_N(X^{\flat})\le\frac16$ by Proposition~\ref{9.9:prop-shift-curvature}. This
is because $x_{k_0}\ge X^{\flat}$ by the pairwise increment bounds and the pin, and this constant
does not read the depth $s^{\circ}$. The number $c_E<1$ is the constant of
Proposition~\ref{9.9:prop-remainder-derivative}, whose proof is outstanding. It may read
$(L,\mathfrak{p},\gamma,\vartheta;\varrho_J^{0})$ but not $n$, $s^{\circ}$, $m_0$, the pair or its
position.

Consequently, with $m_0$, $\mathsf{d}^{\circ}$ and $f^{\circ}$ fixed,
\begin{equation}\label{9.9:eq-nondegenerate-calibration-c}
\varphi'(0)=\lim_{n\to\infty}\varphi_n'(0)\;\ge\;c_{\mathrm{cal}}\,\vartheta^{2}\,
\|f_1^{\circ}\|_{\infty}\|f_2^{\circ}\|_{\infty}\,\bigl(\bar{x}^{\Theta}_{s^{\circ}}\bigr)^{-3}>0 .
\end{equation}
Thus the limiting calibrating function is non-degenerate at the calibration point. With
$q:=\varphi(0)$, the point $\hat{w}_q:=0$, at which $\varphi$ takes the value $q$, satisfies
$\varphi'(\hat{w}_q)\neq0$, in the local form given by \eqref{9.9:eq-nondegenerate-calibration-c}.
\end{theorem}

\begin{proof}[Proof of \eqref{9.9:eq-nondegenerate-calibration-a}]
Fix $n\ge n_1$, and write $f_1=f_1^{\circ}$, $f_2=f_2^{\circ}$.

We first compare two partition functions. For all complex $(\hat{w},z_1,z_2)$ we claim
\begin{equation}\label{9.9:eq-nondegenerate-calibration-1}
Z(a_n-\hat{w};z_1,z_2)=Z\bigl(a_n;\hat{w}\cdot 1+z_1f_1+z_2f_2\bigr).
\end{equation}
Both sides are defined in Notation~\ref{9.9:not-weighted-action}. On the left the inverse coupling
is shifted to $a_n-\hat{w}$ and the sources are $z_1f_1+z_2f_2$. On the right the inverse coupling
is $a_n$ and the source is $\hat{w}\cdot1+z_1f_1+z_2f_2$, where $1$ is the constant test function.
By that definition the two integrands coincide: both equal
\begin{equation*}
\exp\bigl[-A\bigl(a_n-\hat{w}-z_1f_1-z_2f_2,\,U\bigr)\bigr],
\end{equation*}
with $A(\cdot\,,U)$ the weighted action of Notation~\ref{9.9:not-weighted-action}. This proves
\eqref{9.9:eq-nondegenerate-calibration-1}.

Next we take logarithms. Both sides of \eqref{9.9:eq-nondegenerate-calibration-1} are entire
functions of $(\hat{w},z_1,z_2)$, and they are positive at real arguments. In particular the common
value at $(0,0,0)$ is positive. By continuity the function $Z(a_n-\hat{w};z_1,z_2)$ has no zero on a
neighbourhood of $(0,0,0)$ in $\mathbb{C}^3$. On such a neighbourhood $\log Z(a_n-\hat{w};z_1,z_2)$,
the branch real at real arguments, is holomorphic in $(\hat{w},z_1,z_2)$. Hence its mixed partial
derivatives exist to all orders and commute.

Now we compute $\varphi_n'(0)$. By the definition of the calibrating function in
Notation~\ref{9.9:not-tuned-family}, $\varphi_n(\hat{w})$ is the connected two-point function
$\partial_{z_1}\partial_{z_2}|_{z=0}\log Z(a_n-\hat{w};z_1,z_2)$ at bare inverse coupling
$a_n-\hat{w}$. Differentiating in $\hat{w}$ at $0$ and using the commutation of mixed partial
derivatives gives the first equality below. The second equality is
\eqref{9.9:eq-nondegenerate-calibration-1} with $\hat{w}$ renamed $z_0$:
\begin{align*}
\varphi_n'(0)
&=\partial_{\hat{w}}\big|_{0}\,\partial_{z_1}\partial_{z_2}\big|_{0}\log Z(a_n-\hat{w};z_1,z_2)\\
&=\partial_{z_0}\partial_{z_1}\partial_{z_2}\big|_{0}\,
\log Z\bigl(a_n;z_0\cdot1+z_1f_1+z_2f_2\bigr).
\end{align*}
The last expression is, by Notation~\ref{9.9:not-weighted-action}, the connected three-point
function $S^T_{3;n}(1,f_1,f_2)$ with one constant source. It is taken at $n$ steps (the step count
$K:=n$) and bare inverse coupling $a_n$. This is
\eqref{9.9:eq-nondegenerate-calibration-a}.

Finally, a remark on the constant source. The constant test function $1$ has the whole torus as its
support. Therefore the bound of clause~(iii)(b) of Theorem~\ref{3.1:thm-main} for smeared Schwinger
functions with pairwise separated supports does not apply to
$S^T_{3;n}(1,f_1^{\circ},f_2^{\circ})$, and no such bound is used for it in this section.
\end{proof}

Inequalities \eqref{9.9:eq-nondegenerate-calibration-b} and
\eqref{9.9:eq-nondegenerate-calibration-c} are proved below in this section from the hypotheses of
Theorem~\ref{9.9:thm-nondegenerate-calibration}. The identity
\eqref{9.9:eq-nondegenerate-calibration-a} uses only the definitions.

The bound \eqref{9.9:eq-nondegenerate-calibration-b} uses the following inputs:
\begin{itemize}
\item Proposition~\ref{9.9:prop-shift-curvature};
\item the first hypothesis of Definition~\ref{9.9:def-tuned-hypotheses}, which runs
Proposition~\ref{9.9:prop-shift-curvature} on the $\Theta_X$-run and gives
$N_{k_0}\in[\frac23,\frac43]$ there;
\item Proposition~\ref{9.9:prop-remainder-derivative}, whose proof is outstanding;
\item the reading of the two-point witness theorem in the fourth hypothesis;
\item for $x_{k_0}\le\bar{x}^{\Theta}_{s^{\circ}}$, the proposition of the uniqueness theorem on
which the upper end of the window rests, the ladder hypothesis and clause~(a) of the uniqueness
theorem.
\end{itemize}
Of the two-point witness theorem only its notation is used, including the main term
$\mathfrak{M}$, together with its lower bound
$\mathfrak{M}\ge c_M\vartheta^{2}g^{4}\|f_1\|_{\infty}\|f_2\|_{\infty}$. Its remainder bound
$|R_K|\le\frac14$ is not used. The bound \eqref{9.9:eq-nondegenerate-calibration-c} uses
\eqref{9.9:eq-nondegenerate-calibration-b}, clause~(d) of the uniqueness theorem, and the second
hypothesis of Definition~\ref{9.9:def-tuned-hypotheses}.

The lower bound in \eqref{9.9:eq-nondegenerate-calibration-b} decays in the depth like
\begin{equation*}
\bigl(\bar{x}^{\Theta}_{s^{\circ}}\bigr)^{-3}\asymp\bigl(X+3\lambda^{U}s^{\circ}\bigr)^{-3},
\end{equation*}
as the witness itself does: it is of order $g^{6}$ at the witness scale. The constant
$c_{\mathrm{cal}}$ is free of the depth; the bound is not. The bound is positive at every admissible
depth, and nothing bounds $s^{\circ}$ from above.

The mechanism is this. The derivative of the witness in the bare inverse coupling is the connected
correlation with the action density. In Definition~\ref{9.9:def-main-term},
$x_{k_0}(\hat{w})=x_{k_0}-\zeta_{k_0}(\hat{w})$ and $g_{k_0}(\hat{w})=x_{k_0}(\hat{w})^{-1/2}$, so
\begin{equation*}
\frac{d g_{k_0}}{d\hat{w}}\Big|_{0}=\tfrac12 N_{k_0}g_{k_0}^{3}.
\end{equation*}
The derivative of $b_0g_{k_0}^{4}\mathcal{K}$ per unit of $g_{k_0}$ is $4b_0g_{k_0}^{3}\mathcal{K}$.
Hence, with $N_{k_0}$ held fixed,
\begin{equation*}
\frac{d}{d\hat{w}}\Big|_{\hat{w}=0}\,N_{k_0}^{2}\,b_0\,g_{k_0}(\hat{w})^{4}\,\mathcal{K}
=2N_{k_0}g_{k_0}^{2}\,\mathfrak{M}_n(0).
\end{equation*}
This is the leading term, and $2N_{k_0}\ge\frac43$. The derivative stays bounded away from zero
exactly when two quantities are small against this leading term:
\begin{itemize}
\item the curvature of the running shift, which by Proposition~\ref{9.9:prop-shift-curvature} is of
relative size $c_N(X)\to0$ and comes from the correlation theorem alone;
\item the derivative of the witness remainder, bounded in
Proposition~\ref{9.9:prop-remainder-derivative}, whose proof is outstanding.
\end{itemize}

\begin{proof}[Proof of \eqref{9.9:eq-nondegenerate-calibration-b} and \eqref{9.9:eq-nondegenerate-calibration-c} in Theorem~\ref{9.9:thm-nondegenerate-calibration}]
Throughout, the data are those of Theorem~\ref{9.9:thm-nondegenerate-calibration}: $g\le\gamma_{\mathrm{cal}}$,
$X=g^{-2}$, $X^{\flat}=X-2c_2$, $m_0\ge m_{\Theta}(g)$, $s^{\circ}=s(\mathsf{d}^{\circ})\ge s_W(g)$,
$f^{\circ}\in\mathfrak{T}(\mathsf{d}^{\circ};\vartheta)$ and $n\ge n_1$.

\emph{The derivative of the main term.} By Definition~\ref{9.9:def-main-term}, see
\eqref{9.9:eq-main-term-a}, we have $\varphi_n=\mathfrak{M}_n+E_n$ on a neighbourhood of $\hat{w}=0$,
and all three functions are holomorphic there. Hence
$\varphi_n'(0)=\mathfrak{M}_n'(0)+E_n'(0)$. In
$\mathfrak{M}_n(\hat{w})=N_{k_0}(\hat{w})^2\,b_0\,x_{k_0}(\hat{w})^{-2}\,\mathcal{K}$ the factors
$b_0=(N^2-1)/2$ and $\mathcal{K}$ do not depend on $\hat{w}$, by \eqref{9.9:eq-tuned-family-a}.
The factor $x_{k_0}(\hat{w})$ does not vanish on the disc $D(0,2r)$ of Definition~\ref{9.9:def-main-term}.
The factor $N_{k_0}(\hat{w})$ does not vanish near $0$, since $N_{k_0}(0)=N_{k_0}\in[\tfrac23,\tfrac43]$,
as shown below. So we may differentiate
\begin{equation*}
\log\mathfrak{M}_n(\hat{w})=2\log N_{k_0}(\hat{w})-2\log x_{k_0}(\hat{w})+\log(b_0\mathcal{K}),
\end{equation*}
or equivalently apply the product rule. By Definition~\ref{9.9:def-main-term},
$N_{k_0}'(\hat{w})=\zeta_{k_0}''(\hat{w})$ and $x_{k_0}'(\hat{w})=-\zeta_{k_0}'(\hat{w})=-N_{k_0}(\hat{w})$.
At $\hat{w}=0$ the shifted quantities reduce to those of the run: $N_{k_0}(0)=N_{k_0}$, $x_{k_0}(0)=x_{k_0}$ and
$\mathfrak{M}_n(0)=\mathfrak{M}$. Since $g_{k_0}^2=1/x_{k_0}$, this gives
\begin{equation*}
\frac{\mathfrak{M}_n'(0)}{\mathfrak{M}_n(0)}
=\frac{2\zeta_{k_0}''(0)}{N_{k_0}}+\frac{2\zeta_{k_0}'(0)}{x_{k_0}}
=\frac{2\zeta_{k_0}''(0)}{N_{k_0}}+2N_{k_0}g_{k_0}^2 .
\end{equation*}
The number $\mathfrak{M}_n(0)=\mathfrak{M}$ is positive, by the lower bound recalled below. We bound the
curvature term and $E_n'(0)$ from below by minus their moduli and obtain
\begin{equation}\label{9.9:eq-calibration-lower-bound-a}
\begin{aligned}
\varphi_n'(0)
&\ge\mathfrak{M}_n(0)\Bigl[2N_{k_0}g_{k_0}^2-\frac{2|\zeta_{k_0}''(0)|}{N_{k_0}}\Bigr]-|E_n'(0)|\\
&\ge\mathfrak{M}_n(0)\,g_{k_0}^2\,\bigl[2N_{k_0}-2c_N-c_E\bigr],
\end{aligned}
\end{equation}
where $c_N=c_N(x_{k_0})$. The second inequality rests on three inputs, which we now supply.

\emph{The curvature of the running shift.} We apply Proposition~\ref{9.9:prop-shift-curvature}
(curvature of the running shift) to the tuned run of the class $\Theta_X$, with $X^{\flat}=X-2c_2$ and
$\kappa_0>8$.

Its hypothesis $x_j\ge X^{\flat}$ for all levels $j\le k_0$ comes from the first of the four hypotheses of
Definition~\ref{9.9:def-tuned-hypotheses}, the transfer of the pairwise coupling bounds and of the
running-shift statement of the correlation theorem to the class $\mathfrak{S}$ of the uniqueness theorem.
That hypothesis gives, for the tuned run, the pairwise
bounds of the correlation theorem, whose lower half is
\begin{equation*}
x_i-x_j\ge\lambda(j-i)-2c_2\qquad(i<j).
\end{equation*}
Combined with the normalisation $x_{(0)}=X$ at depth $0$, these bounds give
$x_{(d)}\ge X+\lambda d-2c_2\ge X^{\flat}$ at every
depth $d\ge1$. The remaining conditions that Proposition~\ref{9.9:prop-shift-curvature} places on
$X^{\flat}$ are among the conditions on $g$ that define $\gamma_{\mathrm{cal}}$.

Proposition~\ref{9.9:prop-shift-curvature}, through its bound,
controls the curvature term:
\begin{equation*}
\frac{2|\zeta_{k_0}''(0)|}{N_{k_0}}\le 2c_N(x_{k_0})\,g_{k_0}^2 .
\end{equation*}
The function $c_N$ is non-increasing and $x_{k_0}\ge X^{\flat}$. Hence
$c_N(x_{k_0})\le c_N(X^{\flat})$.

\emph{The remainder.} The bare-offset derivative of the witness remainder,
Proposition~\ref{9.9:prop-remainder-derivative}, whose proof is outstanding, gives
\begin{equation*}
|E_n'(0)|\le c_E\,g_{k_0}^2\,\mathfrak{M}_n(0),\qquad c_E<1 .
\end{equation*}
Its quantifiers are met: $g\le\gamma_{\mathrm{cal}}\le\gamma_E$, $m_0\ge m_{\Theta}(g)$,
$s(\mathsf{d}^{\circ})\ge s_W(g)$, $f^{\circ}\in\mathfrak{T}(\mathsf{d}^{\circ};\vartheta)$ and $n\ge n_1$.

\emph{The first derivative of the running shift.} We claim that the first derivative at $0$ of the
running shift lies in $[\tfrac23,\tfrac43]$ at every depth; in particular $N_{k_0}\in[\tfrac23,\tfrac43]$.
For the runs of the correlation theorem this is part of its running-shift statement. For the
tuned run of the class $\Theta_X$, the first hypothesis of Definition~\ref{9.9:def-tuned-hypotheses}
carries the running-shift statement over to the class $\mathfrak{S}$ of the uniqueness theorem. It does
so with level-wise Lipschitz constants whose sum over the levels is at most $\tfrac13$. For the running
shift recentred at the tuned point, $\zeta\mapsto\zeta_{(s)}(\zeta_*^{n}+\zeta)-\zeta_{(s)}(\zeta_*^{n})$,
this bound gives at every depth $s$
\begin{equation*}
\bigl|\eta_s(\zeta)\bigr|\le\tfrac13|\zeta|,\qquad
\eta_s(\zeta):=\zeta_{(s)}(\zeta_*^{n}+\zeta)-\zeta_{(s)}(\zeta_*^{n})-\zeta .
\end{equation*}
The function $\eta_s$ is holomorphic, vanishes at $0$, and is bounded by $\tfrac13|\zeta|$ on the disc of
offsets on which that bound holds; let $\delta$ be its radius. So $\zeta\mapsto3\eta_s(\delta\zeta)/\delta$
maps the unit disc into the closed unit disc and fixes $0$. Schwarz's
lemma then gives $|\eta_s'(0)|\le\tfrac13$. At $s=s^{\circ}$, where $\eta_s'(0)=N_{k_0}-1$, this is
$|N_{k_0}-1|\le\tfrac13$. In particular $2N_{k_0}\ge\tfrac43$.

\emph{Conclusion for \eqref{9.9:eq-nondegenerate-calibration-b}.} One of the conditions defining
$\gamma_{\mathrm{cal}}$, a condition on $g$ alone, is $c_N(X^{\flat})\le\tfrac16$. With $c_E<1$, the
bracket in \eqref{9.9:eq-calibration-lower-bound-a} satisfies
\begin{equation*}
2N_{k_0}-2c_N(x_{k_0})-c_E\ \ge\ \tfrac43-2c_N(X^{\flat})-c_E\ \ge\ \tfrac43-\tfrac13-c_E\ =\ 1-c_E\ >\ 0 .
\end{equation*}
This yields the first inequality of \eqref{9.9:eq-nondegenerate-calibration-b}.

Further, $\mathfrak{M}_n(0)=\mathfrak{M}$, and the main-term lower bound of the two-point witness
theorem gives
\begin{equation*}
\mathfrak{M}\ge c_M\,\vartheta^2\,g_{k_0}^4\,\|f_1^{\circ}\|_\infty\|f_2^{\circ}\|_\infty .
\end{equation*}
That bound uses the lower bound on the mass $\mathfrak{n}(f)$ of a bump in terms of $\vartheta$ and
$\|f\|_\infty$ printed in the two-point witness theorem, and the lower kernel bound stated there.
The fourth hypothesis of Definition~\ref{9.9:def-tuned-hypotheses}, which reads the two-point witness
theorem on the data of the uniqueness theorem, places that bound on those data. Finally,
$g_{k_0}^2=1/x_{k_0}\ge1/\bar{x}^{\Theta}_{s^{\circ}}$ by the upper window of
\eqref{9.9:eq-tuned-family-a}. Theorem~\ref{9.9:thm-nondegenerate-calibration} assumes that window
through its hypotheses: the third hypothesis of Definition~\ref{9.9:def-tuned-hypotheses} (one ladder
for the whole class $\Theta_X$), clause~(a) of the uniqueness theorem, and the proposition on the window
of the tuned runs that Theorem~\ref{9.9:thm-nondegenerate-calibration} assumes with them. Combining
these bounds with the positivity of the bracket shown above,
\begin{align*}
\varphi_n'(0)
&\ge\bigl(\tfrac43-2c_N(X^{\flat})-c_E\bigr)\,c_M\,\vartheta^2\,g_{k_0}^6\,
\|f_1^{\circ}\|_\infty\|f_2^{\circ}\|_\infty\\
&\ge\bigl(\tfrac43-2c_N(X^{\flat})-c_E\bigr)\,c_M\,\vartheta^2\,
\|f_1^{\circ}\|_\infty\|f_2^{\circ}\|_\infty\,\bigl(\bar{x}^{\Theta}_{s^{\circ}}\bigr)^{-3},
\end{align*}
which is \eqref{9.9:eq-nondegenerate-calibration-b}. Its right member does not depend on $n$, and it uses no
condition on $s^{\circ}$ beyond the floor $s^{\circ}\ge s_W(g)$ of the two-point witness theorem.

\emph{Passage to the limit, \eqref{9.9:eq-nondegenerate-calibration-c}.} Put
$\upsilon:=\tfrac12\min\{\varrho_2/2,\upsilon_w\}$. By the second hypothesis of
Definition~\ref{9.9:def-tuned-hypotheses}, the holomorphy of the calibrating functions, the functions
$\varphi_n$ with $n\ge n_0(N_w)$, the threshold of clause~(c) of the uniqueness theorem, are holomorphic on
the interior of $\mathcal{N}'_w$, and this interior contains $\bar{D}(0,\upsilon)$. By clause~(d) of the
uniqueness theorem they converge to $\varphi$ uniformly on $\bar{D}(0,\upsilon)$. Weierstrass' theorem
then gives $\varphi_n'\to\varphi'$ uniformly on $\bar{D}(0,\upsilon/2)$, and in particular
$\varphi_n'(0)\to\varphi'(0)$. The lower bound \eqref{9.9:eq-nondegenerate-calibration-b} holds for all
$n\ge n_1$ with a right member free of $n$, so it passes to the limit. This is
\eqref{9.9:eq-nondegenerate-calibration-c}. In particular $\varphi'(0)\neq0$: the non-degeneracy
hypothesis of the calibration holds, in its differential clause, at the calibration point $\hat{w}=0$,
with calibration value $\varphi(0)$.
\end{proof}

\begin{corollary}[Physical calibration near the tuned point]\label{9.9:cor-physical-calibration}
Assume the following.
\begin{itemize}
\item[(a)] The hypotheses of Theorem~\ref{9.9:thm-nondegenerate-calibration} (non-degenerate calibration
by the two-point witness) hold. We fix $g$, the torus exponent $m_0$, the separation
$\mathsf{d}^{\circ}$ with $s^{\circ}:=s(\mathsf{d}^{\circ})$, the reference pair
$f^{\circ}=(f_1^{\circ},f_2^{\circ})\in\mathfrak{T}(\mathsf{d}^{\circ};\vartheta)$ and the
index $n_1$ as there, so that, by that theorem, the lower bound
\eqref{9.9:eq-nondegenerate-calibration-c} holds. Here $a_n$ are the tuned
bare inverse couplings, $S^n(x;f)$ denotes the connected $p'$-point function of a tuple $f$
of $p'\ge2$ test functions on the bare datum $(m_0+n,x)$, and, as in that theorem,
\begin{equation*}
\varphi_n(\hat{w}):=S^n(a_n-\hat{w};f^{\circ}),\qquad \varphi(\hat{w}):=S^{(\hat{w})}(f^{\circ}).
\end{equation*}
\item[(b)] The holomorphy hypothesis among the assumptions of
Theorem~\ref{9.9:thm-nondegenerate-calibration} holds.
\item[(c)] The convergence clause of the uniqueness theorem holds. That is, for every
$p'\ge2$ and every admissible tuple $f=(f_1,\dots,f_{p'})$ of the uniqueness theorem
(Lipschitz functions with pairwise separated supports on the continuum torus of exponent
$m_0$), the following hold.
\begin{itemize}
\item[(c1)] $S^n(a_n-\hat{w};f)\to S^{(\hat{w})}(f)$ without subsequences, uniformly on the
strip neighbourhood
\begin{equation*}
\mathcal{N}'_w=\{|\operatorname{Re}\hat{w}|\le\varrho_2/2,\ |\operatorname{Im}\hat{w}|<\upsilon_w\}.
\end{equation*}
In particular (taking $\hat{w}=a_n-a'_n$, a real point of $\mathcal{N}'_w$, in this uniform
convergence), every real sequence $a'_n$ with $|a'_n-a_n|\le\varrho_2/2$ satisfies
\begin{equation*}
S^n(a'_n;f)-S^{(a_n-a'_n)}(f)\to0 .
\end{equation*}
\item[(c2)] The map $\hat{w}\mapsto S^{(\hat{w})}(f)$ is real-analytic on
$\mathopen]-\varrho_2/2,\varrho_2/2\mathclose[$.
\end{itemize}
\end{itemize}
Put $q:=\varphi(0)$. Then there are $\delta\in\mathopen]0,\varrho_2/2\mathclose[$ and an index $n_2$
such that the following hold.
\begin{itemize}
\item[(i)] The function $\varphi$ is strictly increasing on $[-\delta,\delta]$, and there
\begin{equation}\label{9.9:eq-physical-calibration-1}
\varphi'\ \ge\ \tfrac12\,c_{\mathrm{cal}}\,\vartheta^2\,\|f_1^{\circ}\|_\infty\,\|f_2^{\circ}\|_\infty\,
\bigl(\bar{x}^{\Theta}_{s^{\circ}}\bigr)^{-3}.
\end{equation}
\item[(ii)] For every $n\ge n_2$ the equation $\varphi_n(\hat{w})=q$ has exactly one solution
$\hat{w}^{(n)}$ in $[-\delta,\delta]$, and $\hat{w}^{(n)}\to0$.
\item[(iii)] Put $b_n:=a_n-\hat{w}^{(n)}$. This is the bare inverse coupling at depth $n$
calibrated by the physical condition that the connected curvature two-point function on
$f^{\circ}$ equals $q$. Then for every $p'\ge2$ and every admissible tuple $f$ we have
$S^n(b_n;f)\to S^{(0)}(f)$, without subsequences. More generally, let $b'_n$ be any real
sequence with $|a_n-b'_n|\le\delta$ and $S^n(b'_n;f^{\circ})\to q$. Then
$S^n(b'_n;f)\to S^{(0)}(f)$ for every $p'\ge2$ and every admissible tuple $f$.
\end{itemize}
\end{corollary}

The corollary asserts uniqueness of the calibrating solution on the interval
$[-\delta,\delta]$ only. No uniqueness of a solution of $\varphi_n(\hat{w})=q$ on the whole
window $[-\varrho_2/2,\varrho_2/2]$ is asserted. The non-degeneracy of the calibration at the
point $\hat{w}=0$ enters only through the local bound \eqref{9.9:eq-nondegenerate-calibration-c}.

\begin{proof}
Put
\begin{equation*}
c_{\mathrm{low}}:=c_{\mathrm{cal}}\,\vartheta^2\,\|f_1^{\circ}\|_\infty\,\|f_2^{\circ}\|_\infty\,
\bigl(\bar{x}^{\Theta}_{s^{\circ}}\bigr)^{-3}.
\end{equation*}
By \eqref{9.9:eq-nondegenerate-calibration-c} we have $\varphi'(0)\ge c_{\mathrm{low}}>0$.

We first record the convergence of derivatives. By the calibration lower bound and its limit,
established in the proof of Theorem~\ref{9.9:thm-nondegenerate-calibration} (the argument giving
\eqref{9.9:eq-nondegenerate-calibration-c}), where the Weierstrass theorem is applied to the
holomorphy hypothesis (b) and to the uniform convergence (c1) on $\mathcal{N}'_w$ with
$f=f^{\circ}$, we have $\varphi_n'\to\varphi'$ uniformly on a real segment
$[-\upsilon/2,\upsilon/2]$; here $\upsilon>0$ is the length of the segment on which that
argument establishes the uniform convergence of the derivatives.

\emph{Choice of $\delta$ and $n_2$.} By (c2), $\varphi$ is real-analytic on
$\mathopen]-\varrho_2/2,\varrho_2/2\mathclose[$, so $\varphi'$ is continuous there. Since
$\varphi'(0)\ge c_{\mathrm{low}}$, there is $\delta$ with
\begin{equation*}
0<\delta\le\upsilon/2,\qquad \delta<\varrho_2/2,
\end{equation*}
such that $\varphi'\ge\tfrac34 c_{\mathrm{low}}$ on $[-\delta,\delta]$.

Two uniform convergences hold on $[-\delta,\delta]$. First, $\varphi_n'\to\varphi'$ uniformly
on $[-\upsilon/2,\upsilon/2]\supset[-\delta,\delta]$. Second, $\varphi_n\to\varphi$ uniformly on
$[-\delta,\delta]\subset\mathcal{N}'_w$, by (c1) with $f=f^{\circ}$. Hence there is
$n_2\ge n_1$ such that for all $n\ge n_2$
\begin{equation}\label{9.9:eq-physical-calibration-2}
|\varphi_n'-\varphi'|\le\tfrac14 c_{\mathrm{low}},\qquad
|\varphi_n-\varphi|\le\tfrac14 c_{\mathrm{low}}\,\delta\qquad\text{on }[-\delta,\delta].
\end{equation}

\emph{Clause (i).} On $[-\delta,\delta]$ we have
$\varphi'\ge\tfrac34c_{\mathrm{low}}\ge\tfrac12c_{\mathrm{low}}$. This is
\eqref{9.9:eq-physical-calibration-1}. Since $\varphi'>0$ there, $\varphi$ is strictly
increasing on $[-\delta,\delta]$.

\emph{Clause (ii).} Let $n\ge n_2$. By \eqref{9.9:eq-physical-calibration-2}, on
$[-\delta,\delta]$,
\begin{equation*}
\varphi_n'\ \ge\ \varphi'-\tfrac14c_{\mathrm{low}}\ \ge\ \tfrac34c_{\mathrm{low}}-\tfrac14c_{\mathrm{low}}
\ =\ \tfrac12c_{\mathrm{low}} .
\end{equation*}
So $\varphi_n$ is strictly increasing on $[-\delta,\delta]$ with slope at least
$\tfrac12c_{\mathrm{low}}$.

By the mean value theorem and $\varphi'\ge\tfrac34c_{\mathrm{low}}$ on $[0,\delta]$, we have
$\varphi(\delta)\ge q+\tfrac34c_{\mathrm{low}}\delta$. Together with the second bound in
\eqref{9.9:eq-physical-calibration-2} this gives
\begin{equation*}
\varphi_n(\delta)\ \ge\ \varphi(\delta)-\tfrac14c_{\mathrm{low}}\delta
\ \ge\ q+\tfrac34c_{\mathrm{low}}\delta-\tfrac14c_{\mathrm{low}}\delta\ >\ q .
\end{equation*}
Likewise $\varphi(-\delta)\le q-\tfrac34c_{\mathrm{low}}\delta$, and therefore
\begin{equation*}
\varphi_n(-\delta)\ \le\ \varphi(-\delta)+\tfrac14c_{\mathrm{low}}\delta
\ \le\ q-\tfrac12c_{\mathrm{low}}\delta\ <\ q .
\end{equation*}
The function $\varphi_n$ is continuous and strictly increasing on $[-\delta,\delta]$. By the
intermediate value theorem, the equation $\varphi_n(\hat{w})=q$ has exactly one solution
$\hat{w}^{(n)}$ in $\mathopen]-\delta,\delta\mathclose[$. Since neither endpoint is a solution, this is also the
only solution in $[-\delta,\delta]$.

We now bound $\hat{w}^{(n)}$. The slope bound gives
$|\varphi_n(\hat{w}^{(n)})-\varphi_n(0)|\ge\tfrac12c_{\mathrm{low}}|\hat{w}^{(n)}|$. Since
$\varphi_n(\hat{w}^{(n)})=q$, this yields
\begin{equation*}
|\hat{w}^{(n)}|\ \le\ \frac{2\,|\varphi_n(0)-q|}{c_{\mathrm{low}}} .
\end{equation*}
By (c1) at $\hat{w}=0$ with $f=f^{\circ}$, $\varphi_n(0)\to\varphi(0)=q$. Hence
$\hat{w}^{(n)}\to0$. This proves (ii).

\emph{Clause (iii).} Fix $p'\ge2$ and an admissible tuple $f$. We have
$|b_n-a_n|=|\hat{w}^{(n)}|\le\delta$, and $\delta<\varrho_2/2$. The consequence stated in (c1),
with $a'_n:=b_n$, gives
\begin{equation*}
S^n(b_n;f)-S^{(\hat{w}^{(n)})}(f)\to0 .
\end{equation*}
By (c2) the map $\hat{w}\mapsto S^{(\hat{w})}(f)$ is real-analytic, hence continuous, on
$\mathopen]-\varrho_2/2,\varrho_2/2\mathclose[$. This interval contains $[-\delta,\delta]$, and
$\hat{w}^{(n)}\to0$ by (ii). Therefore $S^{(\hat{w}^{(n)})}(f)\to S^{(0)}(f)$, and so
$S^n(b_n;f)\to S^{(0)}(f)$. No subsequence has been extracted at any point.

For the second assertion of (iii), let $b'_n$ be real with $|a_n-b'_n|\le\delta$ and
$S^n(b'_n;f^{\circ})\to q$. Since $\delta\le\varrho_2/2$, the consequence stated in (c1),
applied to $f=f^{\circ}$ and $a'_n:=b'_n$, gives
$S^n(b'_n;f^{\circ})-\varphi(a_n-b'_n)\to0$. Hence $\varphi(a_n-b'_n)\to q=\varphi(0)$.

The points $a_n-b'_n$ lie in $[-\delta,\delta]$. There $\varphi'\ge\tfrac12c_{\mathrm{low}}$ by
(i), so the mean value theorem gives
\begin{equation*}
|a_n-b'_n|\ \le\ \frac{2\,|\varphi(a_n-b'_n)-q|}{c_{\mathrm{low}}}\ \to\ 0 .
\end{equation*}
Now fix $p'\ge2$ and an admissible tuple $f$. The consequence stated in (c1), with
$a'_n:=b'_n$, gives $S^n(b'_n;f)-S^{(a_n-b'_n)}(f)\to0$. The continuity of
$\hat{w}\mapsto S^{(\hat{w})}(f)$ from (c2), together with $a_n-b'_n\to0$, gives
$S^{(a_n-b'_n)}(f)\to S^{(0)}(f)$. Hence $S^n(b'_n;f)\to S^{(0)}(f)$.
\end{proof}

\begin{theorem}[The witness differentiated in the bare coupling]\label{9.9:thm-differentiated-witness}
The \emph{tuned family} is the family of runs of Theorem~\ref{9.9:thm-nondegenerate-calibration}:
the $\Theta_X$-runs of the class $\mathfrak{S}$ of the uniqueness theorem, whose scaffold is frozen in
terms of $X$, the $n$-th run having $n$ steps and bare inverse coupling $a_n$; its main term
$\mathfrak{M}_n$ is that of \eqref{9.9:eq-main-term-a}. Assume the following.
\begin{itemize}
\item[(a)] $\kappa_0>8$, for the parameter $\kappa_0$ of Proposition~\ref{9.9:prop-shift-curvature},
and the remaining hypotheses of Proposition~\ref{9.9:prop-shift-curvature}
(curvature of the running shift) hold for the tuned family, with the lower bound
$X^{\flat}:=X-2c_2=g^{-2}-2c_2$ of Theorem~\ref{9.9:thm-nondegenerate-calibration} for its inverse
couplings.
\item[(b)] The hypothesis of Definition~\ref{9.9:def-tuned-hypotheses} on the class $\mathfrak{S}$
holds: the pairwise bounds on the inverse couplings and the running-shift statement of the
correlation theorem, with the bounds on its holomorphic part and on its remnant, hold for the
runs of the tuned family.
\item[(c)] The conclusion of Proposition~\ref{9.9:prop-remainder-derivative}, whose proof is
outstanding, holds; that is, there are a ceiling $\gamma_E\in\,]0,\min\{\gamma_W,\gamma_U\}]$ and a
constant $c_E<1$, both depending at most on $(L,\mathfrak{p},\gamma,\vartheta;\varrho_J^{0})$, such that
for every $g\le\gamma_E$, every $m_0\ge m_{\Theta}(g)$, every separation $\mathsf{d}^{\circ}$ with
$s(\mathsf{d}^{\circ})\ge s_W(g)$, every reference pair
$f^{\circ}\in\mathfrak{T}(\mathsf{d}^{\circ};\vartheta)$ and every $n\ge n_1$, along the tuned family,
the bound \eqref{9.9:eq-remainder-derivative-a} holds.
\end{itemize}
In the frame of Theorem~\ref{9.9:thm-nondegenerate-calibration}, with the ceiling
$\gamma_{\mathrm{cal}}$ and the integer $n_1$ of that theorem (the $n_1$ of hypothesis~(c)), let the
choices be made in the following order:
\begin{itemize}
\item fix $\vartheta\in\,]0,1]$;
\item then $g\le\gamma_{\mathrm{cal}}$;
\item then $m_0\ge m_{\Theta}(g)$;
\item then a separation $\mathsf{d}^{\circ}$ with $s^{\circ}:=s(\mathsf{d}^{\circ})\ge s_W(g)$;
\item then a reference pair $f^{\circ}=(f_1^{\circ},f_2^{\circ})\in\mathfrak{T}(\mathsf{d}^{\circ};\vartheta)$;
\item then $n\ge n_1$.
\end{itemize}
Let $k_0$, $N_{k_0}$, $g_{k_0}$ and $\mathcal{K}$ be the level, the factors and the kernel of the
main term $\mathfrak{M}_n$ in \eqref{9.9:eq-main-term-a}. Then
\begin{equation}\label{9.9:eq-differentiated-witness-a}
\begin{gathered}
S^T_{3;n}(1,f_1^{\circ},f_2^{\circ})
 =2N_{k_0}^{3}\,b_0\,g_{k_0}^{6}\,\mathcal{K}\,(1+R_3),\\
|R_3|\le 1-\delta_W,\qquad \delta_W\ge\tfrac34(1-c_E)>0,
\end{gathered}
\end{equation}
where
\begin{equation*}
\delta_W:=1-\tfrac32\,c_N(X^{\flat})-\tfrac34\,c_E,
\end{equation*}
with $X^{\flat}$ as in~(a), and $c_N$ is the non-increasing function of
Proposition~\ref{9.9:prop-shift-curvature}.
The lower bound $\tfrac34(1-c_E)$ depends at most on $(L,\mathfrak{p},\gamma,\vartheta;\varrho_J^{0})$.
Neither $\delta_W$ nor this lower bound depends on $n$, $s^{\circ}$, $m_0$, the pair $f^{\circ}$ or the
position of the supports.
\end{theorem}

\begin{proof}
Throughout, a prime on $\varphi_n$, $\mathfrak{M}_n$, $E_n$ or $\zeta_{k_0}$ denotes the derivative in
the bare offset $\hat{w}$, and the argument $0$ denotes $\hat{w}=0$.

\emph{The three-point function as a derivative.} By
\eqref{9.9:eq-nondegenerate-calibration-a} we have
$\varphi_n'(0)=S^T_{3;n}(1,f_1^{\circ},f_2^{\circ})$; this is the $n$-step function at bare inverse coupling
$a_n$. By \eqref{9.9:eq-main-term-a}, $\varphi_n=\mathfrak{M}_n+E_n$ near $0$, and all three functions
are holomorphic there. Differentiating at $0$ therefore gives
\begin{equation}\label{9.9:eq-differentiated-witness-1}
S^T_{3;n}(1,f_1^{\circ},f_2^{\circ})=\mathfrak{M}_n'(0)+E_n'(0).
\end{equation}

\emph{The derivative of the main term.} In the derivation of
\eqref{9.9:eq-calibration-lower-bound-a}, the logarithm
\begin{equation*}
\log\mathfrak{M}_n(\hat{w})=2\log N_{k_0}(\hat{w})-2\log x_{k_0}(\hat{w})+\log(b_0\mathcal{K})
\end{equation*}
is differentiated; here $b_0=(N^2-1)/2$ and $\mathcal{K}$ do not depend on $\hat{w}$, by the
definition of the tuned family and of its main term \eqref{9.9:eq-main-term-a}. The result is
\begin{equation*}
\frac{\mathfrak{M}_n'(0)}{\mathfrak{M}_n(0)}
=\frac{2\zeta_{k_0}''(0)}{N_{k_0}}+\frac{2\zeta_{k_0}'(0)}{x_{k_0}}
=\frac{2\zeta_{k_0}''(0)}{N_{k_0}}+2N_{k_0}g_{k_0}^{2}.
\end{equation*}
We multiply by $\mathfrak{M}_n(0)=N_{k_0}^{2}b_0g_{k_0}^{4}\mathcal{K}$ and obtain
\begin{equation}\label{9.9:eq-differentiated-witness-2}
\mathfrak{M}_n'(0)
=2N_{k_0}^{3}b_0g_{k_0}^{6}\mathcal{K}+2N_{k_0}b_0g_{k_0}^{4}\mathcal{K}\,\zeta_{k_0}''(0)
=2N_{k_0}^{3}b_0g_{k_0}^{6}\mathcal{K}\Bigl[1+\frac{\zeta_{k_0}''(0)}{N_{k_0}^{2}g_{k_0}^{2}}\Bigr].
\end{equation}

\emph{The remainder $R_3$.} We have
$2N_{k_0}^{3}b_0g_{k_0}^{6}\mathcal{K}=2N_{k_0}g_{k_0}^{2}\mathfrak{M}_n(0)$, and this number is non-zero:
$N_{k_0}\ge\tfrac23$ by hypothesis~(b), as recorded below, $b_0>0$, $g_{k_0}>0$, and $\mathcal{K}>0$,
since in the two-point witness theorem the kernel of a reference pair of the class
$\mathfrak{T}(\mathsf{d}^{\circ};\vartheta)$ is bounded below by
$c_{\mathcal{K}}\vartheta^{2}\|f_1^{\circ}\|_\infty\|f_2^{\circ}\|_\infty>0$. We insert
\eqref{9.9:eq-differentiated-witness-2} into \eqref{9.9:eq-differentiated-witness-1} and divide the
remainder by this number:
\begin{equation*}
\begin{split}
S^T_{3;n}(1,f_1^{\circ},f_2^{\circ})
&=2N_{k_0}^{3}b_0g_{k_0}^{6}\mathcal{K}\Bigl[1+\frac{\zeta_{k_0}''(0)}{N_{k_0}^{2}g_{k_0}^{2}}\\
&\qquad+\frac{E_n'(0)}{2N_{k_0}^{3}b_0g_{k_0}^{6}\mathcal{K}}\Bigr].
\end{split}
\end{equation*}
This is the first line of \eqref{9.9:eq-differentiated-witness-a}, with
\begin{equation}\label{9.9:eq-differentiated-witness-3}
R_3=\frac{\zeta_{k_0}''(0)}{N_{k_0}^{2}g_{k_0}^{2}}+\frac{E_n'(0)}{2N_{k_0}g_{k_0}^{2}\mathfrak{M}_n(0)}.
\end{equation}

\emph{Bound on the first term of $R_3$.} Hypothesis~(b) is used in the proof of
Theorem~\ref{9.9:thm-nondegenerate-calibration}, in its item
\eqref{9.9:eq-nondegenerate-calibration-b}, for two things: to run
Proposition~\ref{9.9:prop-shift-curvature} on the tuned family, and to give
$N_{k_0}\in[\tfrac23,\tfrac43]$ there. It is used in the same way here. Hence $1/N_{k_0}\le\tfrac32$.

The pairwise bounds in~(b), applied between the level of depth $0$, at which the inverse coupling is
$x_{(0)}=X$, and the levels at depth $r\ge1$, give, with the rate $\lambda$ of the pairwise bounds in
Definition~\ref{9.9:def-tuned-hypotheses}, $x_{(r)}\ge X+\lambda r-2c_2\ge X^{\flat}$, as in
Theorem~\ref{9.9:thm-nondegenerate-calibration}; in particular $x_{k_0}\ge X^{\flat}$. Since $c_N$ is
non-increasing, it follows that $c_N(x_{k_0})\le c_N(X^{\flat})$.

Proposition~\ref{9.9:prop-shift-curvature}, in the corresponding form, bounds the
curvature of the running shift at level $k_0$. In the form in which this bound enters
\eqref{9.9:eq-calibration-lower-bound-a}, it reads
$|\zeta_{k_0}''(0)|\le c_N(x_{k_0})\,N_{k_0}\,g_{k_0}^{2}$. Hence
\begin{equation*}
\Bigl|\frac{\zeta_{k_0}''(0)}{N_{k_0}^{2}g_{k_0}^{2}}\Bigr|\le\frac{c_N(x_{k_0})}{N_{k_0}}
\le\tfrac32\,c_N(X^{\flat}).
\end{equation*}

\emph{Bound on the second term of $R_3$.} We have $g\le\gamma_{\mathrm{cal}}\le\gamma_E$. Moreover,
$m_0\ge m_{\Theta}(g)$, $s^{\circ}\ge s_W(g)$, $f^{\circ}\in\mathfrak{T}(\mathsf{d}^{\circ};\vartheta)$
and $n\ge n_1$. These are exactly the ranges of hypothesis~(c). The bound
\eqref{9.9:eq-remainder-derivative-a} therefore gives
$|E_n'(0)|\le c_E\,g_{k_0}^{2}\,\mathfrak{M}_n(0)$, and hence
\begin{equation*}
\Bigl|\frac{E_n'(0)}{2N_{k_0}g_{k_0}^{2}\mathfrak{M}_n(0)}\Bigr|\le\frac{c_E}{2N_{k_0}}\le\tfrac34\,c_E .
\end{equation*}

\emph{Conclusion for $R_3$.} Adding the two bounds in \eqref{9.9:eq-differentiated-witness-3} gives
\begin{equation*}
|R_3|\le\frac{c_N(X^{\flat})}{N_{k_0}}+\frac{c_E}{2N_{k_0}}
\le\tfrac32\,c_N(X^{\flat})+\tfrac34\,c_E=1-\delta_W .
\end{equation*}
Here $\delta_W>0$ holds exactly when $2c_N(X^{\flat})+c_E<\tfrac43$. The ceiling
$\gamma_{\mathrm{cal}}$ of Theorem~\ref{9.9:thm-nondegenerate-calibration} makes
$c_N(X^{\flat})\le\tfrac16$, so that $\tfrac32c_N(X^{\flat})\le\tfrac14$. Therefore
\begin{equation*}
\delta_W\ge1-\tfrac14-\tfrac34\,c_E=\tfrac34(1-c_E)>0,
\end{equation*}
since $c_E<1$. This is the second line of \eqref{9.9:eq-differentiated-witness-a}.

\emph{Dependence of the constants.} The lower bound $\tfrac34(1-c_E)$ depends only on $c_E$. By~(c),
$c_E$ depends at most on $(L,\mathfrak{p},\gamma,\vartheta;\varrho_J^{0})$ and not on $n$, $s^{\circ}$,
$m_0$, the pair or the position. The number $\delta_W$ is built from $c_E$ and from the value
$c_N(X^{\flat})$. Here $X^{\flat}=g^{-2}-2c_2$ is a function of $g$ alone, and $c_N$ is the fixed
function of Proposition~\ref{9.9:prop-shift-curvature}; this value depends on $g$ alone, and on no
depth, cutoff or torus. So
neither $\delta_W$ nor its lower bound depends on $n$, $s^{\circ}$, $m_0$, the pair or the position.
\end{proof}

\begin{remark}[Scope of the statement, and the converse]
The analogous display can be written in the setting of the two-point witness theorem itself, with
first-passage cutoff $K(g_0,g)$, the scaffold of that setting and every $g_0\in\,]0,g_-(g)]$; it is
clause~(iii)(c) of Theorem~\ref{3.1:thm-main} differentiated in the bare coupling. The two families of
runs are related
only through the corollary of the uniqueness theorem. The statement above is therefore made on the
tuned family, on which the calibration of
Theorem~\ref{9.9:thm-nondegenerate-calibration} is carried out.

Conversely, the display \eqref{9.9:eq-differentiated-witness-a} alone gives the lower bound for which
Proposition~\ref{9.9:prop-remainder-derivative} (whose proof is outstanding) is used in
Theorem~\ref{9.9:thm-nondegenerate-calibration}. From
\eqref{9.9:eq-nondegenerate-calibration-a} and \eqref{9.9:eq-differentiated-witness-a}, together with
$1+R_3\ge1-|R_3|\ge\delta_W$, we get
\begin{equation*}
\varphi_n'(0)\ge\delta_W\cdot2N_{k_0}^{3}b_0g_{k_0}^{6}\mathcal{K}
\ge\delta_W\,\tfrac{16}{27}\,b_0c_{\mathcal{K}}\vartheta^{2}
\|f_1^{\circ}\|_\infty\|f_2^{\circ}\|_\infty\bigl(\bar{x}^{\Theta}_{s^{\circ}}\bigr)^{-3}.
\end{equation*}
The second inequality uses two facts. First, $N_{k_0}\ge\tfrac23$, so that
$2N_{k_0}^{3}\ge\tfrac{16}{27}$. Second, the lower bound
\begin{equation*}
b_0g_{k_0}^{6}\mathcal{K}\ge b_0c_{\mathcal{K}}\vartheta^{2}\|f_1^{\circ}\|_\infty
\|f_2^{\circ}\|_\infty\bigl(\bar{x}^{\Theta}_{s^{\circ}}\bigr)^{-3},
\end{equation*}
which is the bound that enters the second inequality of \eqref{9.9:eq-nondegenerate-calibration-b}.

This lower bound on $\varphi_n'(0)$ is all that \eqref{9.9:eq-nondegenerate-calibration-c} and the
corollary on the calibration use. The display
\eqref{9.9:eq-differentiated-witness-a} does not return the bound of
Proposition~\ref{9.9:prop-remainder-derivative} with its own constant. It returns the lower bound on
$\varphi_n'(0)$ that this proposition serves to provide.
\end{remark}

\begin{remark}[Global uniqueness of the calibrated offset]\label{9.9:rem-global-uniqueness}
Let $\varphi_n$ be the calibrating two-point function along the tuned family, regarded as a
function of the real bare offset $\hat{w}$, let $\varphi$ be its limit, and let $n_1$ be the
depth threshold in the frame of Theorem~\ref{9.9:thm-differentiated-witness}. Assume the
following, for every real base offset $\hat{w}_0\in[-\varrho_2/2,\varrho_2/2]$:
\begin{itemize}
\item[(a)] the bare-offset derivative of the witness remainder,
\eqref{9.9:eq-remainder-derivative-a}, holds along the tuned family of the uniqueness theorem
with frozen scaffold, the family on which \eqref{9.9:eq-remainder-derivative-a} is stated,
shifted to the base offset $\hat{w}_0$;
\item[(b)] the curvature bound with the function $\mathfrak{z}$ holds along the same family.
\end{itemize}
That is, Theorem~\ref{9.9:thm-differentiated-witness} applies, and its conclusion
\eqref{9.9:eq-differentiated-witness-a} holds, at every bare inverse coupling
$a_n-\hat{w}_0$ with $\hat{w}_0$ in this window. These couplings are the real $a'_n$ with
$|a'_n-a_n|\le\varrho_2/2$ that the uniqueness theorem admits. This book does not establish (a)
and (b) on the whole window; Corollary~\ref{9.9:cor-physical-calibration} is the local form of
the conclusion below.

Under (a) and (b) there is a constant $c'_{\mathrm{cal}}>0$ with the following properties.
\begin{itemize}
\item[(1)] $\varphi_n'\ge c'_{\mathrm{cal}}$ on $[-\varrho_2/2,\varrho_2/2]$ for every $n\ge n_1$.
\item[(2)] $\varphi$ is strictly increasing on $[-\varrho_2/2,\varrho_2/2]$.
\item[(3)] Let $q\in\varphi(\,]-\varrho_2/4,\varrho_2/4[\,)$. Then for all sufficiently large $n$
the equation $\varphi_n(\hat{w})=q$ has exactly one solution $\hat{w}$ in
$[-\varrho_2/2,\varrho_2/2]$.
\end{itemize}
\end{remark}

\begin{proof}
Fix $\hat{w}_0\in[-\varrho_2/2,\varrho_2/2]$. By (a) and (b), the conclusion of
Theorem~\ref{9.9:thm-differentiated-witness} holds at the bare inverse coupling
$a_n-\hat{w}_0$, and this coupling is admitted by the clause of the uniqueness theorem quoted
above. That conclusion gives the bound $\varphi_n'(\hat{w}_0)\ge c'_{\mathrm{cal}}$ for
$n\ge n_1$. The point $\hat{w}_0$ was arbitrary, so (1) holds.

For (2), let $-\varrho_2/2\le u<v\le\varrho_2/2$. For $n\ge n_1$, the mean value theorem and (1)
give
\begin{equation*}
\varphi_n(v)-\varphi_n(u)\ge c'_{\mathrm{cal}}(v-u).
\end{equation*}
Letting $n\to\infty$ yields $\varphi(v)-\varphi(u)\ge c'_{\mathrm{cal}}(v-u)>0$.

For (3), write $q=\varphi(\hat{w}_q)$ with $|\hat{w}_q|<\varrho_2/4$, and put
$u_\pm:=\hat{w}_q\pm\varrho_2/4$; both points lie in the window. By (2), with the bound just
obtained,
\begin{equation*}
\varphi(u_+)\ge q+c'_{\mathrm{cal}}\varrho_2/4
\quad\text{and}\quad
\varphi(u_-)\le q-c'_{\mathrm{cal}}\varrho_2/4 .
\end{equation*}
Since $\varphi_n(u_\pm)\to\varphi(u_\pm)$, for all sufficiently large $n$ we have
$\varphi_n(u_-)<q<\varphi_n(u_+)$. For $n\ge n_1$ the derivative $\varphi_n'$ bounded in (1)
exists at every point of the window, so $\varphi_n$ is continuous there, and the intermediate
value theorem gives a solution of
$\varphi_n(\hat{w})=q$ in $]u_-,u_+[$. By (1), $\varphi_n$ is strictly increasing on the whole
window for $n\ge n_1$. It therefore takes the value $q$ at most once there, and the solution is
unique in $[-\varrho_2/2,\varrho_2/2]$.
\end{proof}

\begin{remark}[two-sided bounds do not control derivatives]\label{9.9:rem-no-derivative-control}
We use the frozen-shifted tuned family of Notation~\ref{9.9:not-tuned-family}, and the shifted
main term $\mathfrak{M}_n(\hat{w})$ and its remainder $E_n(\hat{w})$ of
Definition~\ref{9.9:def-main-term}, see~\eqref{9.9:eq-main-term-a}. The derivative bound of
Proposition~\ref{9.9:prop-remainder-derivative}, whose proof is outstanding, is the inequality
$|E_n'(0)|\le c_E\,g_{k_0}^2\,\mathfrak{M}_n(0)$ of~\eqref{9.9:eq-remainder-derivative-a}.
This remark has three parts.

\emph{(a) The remainder of the two-point witness theorem.} We write $S^T_{2;K}$ for
$S^T_{2;K,m}(f_1^{\circ},f_2^{\circ})$, the torus exponent $m$ and the reference pair being fixed.
In the two-point witness theorem the
remainder $R_K$ is not an identified next term of an expansion. It is defined by
\begin{equation*}
S^T_{2;K}=\mathfrak{M}\,(1+R_K),\qquad \mathfrak{M}=N_{k_0}^2\,b_0\,g_{k_0}^4\,\mathcal{K},
\end{equation*}
and it is bounded through the covariance decomposition
\begin{equation}\label{9.9:eq-no-derivative-control-1}
S^T_{2;K}=\mathbb{E}_{\tilde{\nu}}\bigl[\Psi^h\bigr]
+\operatorname{Cov}_{\tilde{\nu}}\bigl(\Phi^h_1,\Phi^h_2\bigr).
\end{equation}
Here $\tilde{\nu}$ is the probability law on pairs (history $h$, block field $V$) at level $k_0$
with density $\tau_h/Z$; $\tau^w_h$ is the density of the history $h$ with the complex source
$w=z_1f_1^{\circ}+z_2f_2^{\circ}$ (sampled on the witness lattice) inserted; and
\begin{equation*}
\Phi^h_i:=\partial_{z_i}\log\tau^w_h\big|_{0},\qquad
\Psi^h:=\partial_{z_1}\partial_{z_2}\log\tau^w_h\big|_{0}.
\end{equation*}
In the two-point witness theorem the source derivative is split into six pieces,
\begin{equation}\label{9.9:eq-no-derivative-control-2}
\Phi^h_i=c^h_i+N_{k_0}\,\mathcal{O}^{(k_0)}(\tilde{f}_i)+\mathcal{O}^{(k_0)}(\varepsilon^h_i)
+\mathrm{Irr}_i+\mathrm{Sh}^h_i+\mathrm{LF}^h_i ,
\end{equation}
namely a constant, the main block observable of a level-$k_0$ profile $\tilde{f}_i$ of
$f_i^{\circ}$, a profile error, an irrelevant part, a shift part and a large-field part. The
difference $|S^T_{2;K}-\mathfrak{M}|$ is then bounded by the sum of six members,
each at most $\mathfrak{M}/32$:
\begin{itemize}
\item $R^{\mathrm{prof}}$, the profile smearing, proportional to
$\mathfrak{d}_W/(\vartheta^2L^{s^{\circ}}\mathsf{d}^{\circ})$, with $s^{\circ}=s(\mathsf{d}^{\circ})$
the witness depth of Notation~\ref{9.9:not-tuned-family};
\item $R^{\mathrm{Gauss}}$, the Gaussian extraction at level $k_0$, bounded by the error terms of
the first part of the Gaussian extraction (its extraction error): a power $g^{\sigma_1}$ of the
coupling, with the exponent $\sigma_1$ of that part, times the large-field sum
$\Pi_{r(s^{\circ})}$, a term $g^2\bigl(1+r(s^{\circ})\bigr)$, with $r(s^{\circ})$ the extraction
depth of the Gaussian extraction at the witness depth, further terms, the tail of
the tower and a window term;
\item $R^{\mathrm{irr}}$, bounded by the second part of the Gaussian extraction (its bound on the
irrelevant terms) by
\begin{equation*}
7C_b\,\bigl(L^{s^{\circ}}\mathsf{d}^{\circ}\bigr)^{-2}\,
p_0\bigl((\bar{x}^{\Theta}_{s^{\circ}})^{-1/2}\bigr)^{4},
\end{equation*}
with $C_b$ the constant of that part and $p_0(\cdot)$ the degree function; since, by the
definition of the depth $s^{\circ}=s(\mathsf{d}^{\circ})$,
\begin{equation*}
L^{s^{\circ}}\mathsf{d}^{\circ}\ge C_W(s^{\circ})
=C_0\,p_0\bigl((\bar{x}^{\Theta}_{s^{\circ}})^{-1/2}\bigr)^{2},
\end{equation*}
this bound is at most $7C_bC_0^{-2}$, which is free of the coupling and small only through
the room constant $C_0$ of the depth floor, read at the upper end
$\bar{x}^{\Theta}_{s^{\circ}}$ of the coupling window at depth $s^{\circ}$
(Notation~\ref{9.9:not-tuned-family});
\item $R^{\mathcal{B}}$, the contribution of the contact event $\mathcal{B}$ (some coarse live or
healed label meets one of the supports), proportional to $e^{-c_{\mathcal{B}}\,p_0(g_{k_0})}$,
with $c_{\mathcal{B}}>0$ the decay constant of that event;
\item $R^{\mathrm{loc}}$, the bi-sourced local terms spanning both supports, proportional to
$e^{-\kappa''L^{s^{\circ}}\mathsf{d}^{\circ}/(32M)}$, with $\kappa''$ the localisation rate of
the bi-sourced terms;
\item $R^{\mathrm{sh}}$, which vanishes in the construction of the two-point witness theorem.
\end{itemize}
All six bounds use positivity throughout: covariances are taken under a probability law; the
conditional bounds use $\mathbb{E}^{\natural}$, the supremum over conditionings of the
expectation $\mathbb{E}_{\tilde{\nu}}$; the large-field block bound, healed part, is used through its
sub-multiplicativity; the covariance bound on an analysis window is an absolute bound at the
natural scale, not a sharp relative one; and Wick's theorem is used together with the positivity
$\mathcal{K}>0$ of the kernel.

\emph{(b) A two-sided bound at real offsets does not bound the derivative.} Suppose only that,
along the frozen-shifted tuned family,
\begin{equation}\label{9.9:eq-no-derivative-control-3}
|E_n(\hat{w})|\le\tfrac14\,\mathfrak{M}_n(\hat{w})\qquad\text{for every real }\hat{w}
\text{ near }0 .
\end{equation}
This is more than the two-point witness theorem gives, which concerns only the offset
$\hat{w}=0$ of each run. Then~\eqref{9.9:eq-no-derivative-control-3} does not bound $E_n'(0)$, and in particular
does not give the bound~\eqref{9.9:eq-remainder-derivative-a} of
Proposition~\ref{9.9:prop-remainder-derivative}, whose proof is outstanding. Without further
structure (a convexity, a correlation inequality or an identity) derivative control must come
from one of two routes. We know of no correlation inequality fixing a sign that would give such
derivative control for these correlations in non-abelian gauge theories, for block plaquettes,
smeared connected correlations or physical Wilson loops; the one identity of this kind known to
us concerns a single bare plaquette. The two
routes are:
\begin{equation}\label{9.9:eq-no-derivative-control-a}
\begin{gathered}
\text{(R1)}\quad
\left.
\begin{aligned}
&F\ \text{holomorphic on the disc}\ D(0,R_{\mathrm{cal}}),\\
&|F|\le\eta\ \text{on}\ D(0,R_{\mathrm{cal}})
\end{aligned}
\right\}
\ \Longrightarrow\ |F'(0)|\le\frac{\eta}{R_{\mathrm{cal}}};\\[2pt]
\text{(R2)}\quad\text{differentiation in }\hat{w}\text{ of the proof of the two-point witness
theorem.}
\end{gathered}
\end{equation}
Route~(R1) is used with $F$ the relative remainder and $\eta$ its complex relative precision
$r_{\mathbb{C}}$, once the identification of the main term holds at complex offsets $\hat{w}$.

If the strip-holomorphy hypothesis of the uniqueness theorem, under which the remainder $E_n$ is
holomorphic in $\hat{w}$ on the strip $\mathcal{N}'_w$, is
added to~\eqref{9.9:eq-no-derivative-control-3}, then a derivative bound does follow, but only
the weak bound of part~(c). In this precise sense the bound~\eqref{9.9:eq-remainder-derivative-a}
of Proposition~\ref{9.9:prop-remainder-derivative}, whose proof is outstanding, is not one more
order of the same estimates. The estimates of the two-point witness theorem bound a number. The
bound~\eqref{9.9:eq-remainder-derivative-a} bounds the variation of that number along a family,
and the statement of the two-point witness theorem says nothing about that family.

\emph{(c) How small a real remainder would have to be.} The following classical fact is stated
for orientation. Let $E$ be holomorphic on the disc $D(0,R_{\mathrm{cal}})$, and suppose that
$|E|\le B$ on the disc and $|E|\le\varepsilon$ on its real diameter (this $\varepsilon$ is
unrelated to the profile errors $\varepsilon^h_i$ of part~(a)). Then there is an absolute
constant $C$ with
\begin{equation}\label{9.9:eq-no-derivative-control-4}
|E'(0)|\le C\,\frac{\varepsilon}{R_{\mathrm{cal}}}\,
\max\Bigl\{1,\log\frac{B}{\varepsilon}\Bigr\}.
\end{equation}
It follows from the two-constants theorem, applied by harmonic majorisation of $\log|E|$ on the
two half-discs, together with Cauchy's estimate on a disc about $0$ of radius comparable to
$R_{\mathrm{cal}}/\log(B/\varepsilon)$.

Read~\eqref{9.9:eq-no-derivative-control-4} with $R_{\mathrm{cal}}$ a radius of the kind provided
by the correlation theorem and with $B$ a crude holomorphic majorant. It then says what a
real-offset identification must achieve in order to give~\eqref{9.9:eq-remainder-derivative-a}.
\begin{itemize}
\item At the fixed radius $\varrho_J^{0}$ of the correlation theorem, the real remainder,
measured in units of the main term, must satisfy
\begin{equation*}
\varepsilon\lesssim \frac{g_{k_0}^2\,\varrho_J^{0}}{\log(B/\varepsilon)} .
\end{equation*}
That is, it must be of relative order $g^{\sigma}$ with some $\sigma>2$, which is again the sharp
sufficient condition~\eqref{9.9:eq-sufficient-conditions-c}.
\item At a radius of order $X/\Pi(X)$, with $\Pi(X)$ the large-field floor sum of the uniqueness
theorem for runs tuned to $X$, the real remainder may be any
$\varepsilon=o\bigl(1/(\Pi(X)\log(B/\varepsilon))\bigr)$.
\end{itemize}
In both cases $B$ must be free of the volume. For the bounds used in this book it is not: the
Schwarz bound $\mathsf{Q}$ of the correlation theorem depends on the volume, and the strip
$\mathcal{N}'_w$ of the uniqueness theorem has width $\upsilon_w$ proportional to
$r_J/(2\mathsf{Q}_w+2)$, with $r_J$ the radius of the source domain $\mathcal{W}[r_J]$ of the
correlation theorem and
$\mathsf{Q}_w$ its volume-dependent Schwarz bound at the source $w$. The real route therefore
does not escape the requirements of the complex
route; it only moves the required precision from a circle to a diameter.
\end{remark}

\begin{proof}
We prove the non-implication of part~(b) and the estimate~(R1)
of~\eqref{9.9:eq-no-derivative-control-a}.

\emph{The non-implication.} Assume~\eqref{9.9:eq-no-derivative-control-3}. Since the absolute
value on its left side is non-negative, $\mathfrak{M}_n(\hat{w})\ge0$ for every real $\hat{w}$
near $0$ at which~\eqref{9.9:eq-no-derivative-control-3} is assumed. For $\Omega>0$ put
\begin{equation*}
E_{\Omega}(\hat{w}):=\tfrac14\,\mathfrak{M}_n(\hat{w})\,\sin(\Omega\hat{w}).
\end{equation*}
Since $|\sin(\Omega\hat{w})|\le1$ and $\mathfrak{M}_n(\hat{w})\ge0$, we have
$|E_{\Omega}(\hat{w})|\le\frac14\mathfrak{M}_n(\hat{w})$ at every such real $\hat{w}$. Hence
$E_{\Omega}$ satisfies the hypothesis~\eqref{9.9:eq-no-derivative-control-3} in place of $E_n$.

Because $\sin 0=0$, we have $E_{\Omega}(0)=0$. Therefore, for real $\hat{w}\ne0$,
\begin{equation*}
\frac{E_{\Omega}(\hat{w})-E_{\Omega}(0)}{\hat{w}}
=\frac{\Omega}{4}\,\mathfrak{M}_n(\hat{w})\,\frac{\sin(\Omega\hat{w})}{\Omega\hat{w}}.
\end{equation*}
As $\hat{w}\to0$ the last factor tends to $1$. The function $\mathfrak{M}_n$ is continuous at
$0$, being built in Definition~\ref{9.9:def-main-term} from the running shift $\zeta_{k_0}$,
holomorphic on $|\hat{w}|<2r$ under the hypotheses of Definition~\ref{9.9:def-tuned-hypotheses},
and from its derivative. Hence the
difference quotient tends to
\begin{equation*}
E_{\Omega}'(0)=\frac{\Omega\,\mathfrak{M}_n(0)}{4}.
\end{equation*}
At the offset $\hat{w}=0$ the frozen-shifted main term $\mathfrak{M}_n(0)$ is the main term
$\mathfrak{M}$ of the two-point witness theorem for the run, and the main-term lower bound of that
theorem, which rests on the kernel bound
$\mathcal{K}\ge c_{\mathcal{K}}\,\vartheta^2\,\|f_1^{\circ}\|_\infty\,\|f_2^{\circ}\|_\infty$, reads
\begin{equation*}
\mathfrak{M}\ge c_M\,\vartheta^2\,g_{k_0}^4\,\|f_1^{\circ}\|_\infty\,\|f_2^{\circ}\|_\infty,
\qquad c_M=\tfrac49\,b_0\,c_{\mathcal{K}},
\end{equation*}
so that $\mathfrak{M}_n(0)>0$ for non-zero reference functions. Hence $E_{\Omega}'(0)$ is
unbounded in $\Omega$. In particular, for every $c>0$ one has
$E_{\Omega}'(0)>c\,g_{k_0}^2\,\mathfrak{M}_n(0)$ as soon as
$\Omega>4c\,g_{k_0}^2$; this holds in particular with $c=c_E$. So no bound on the derivative
at $0$, and in particular not the
bound~\eqref{9.9:eq-remainder-derivative-a}, follows from~\eqref{9.9:eq-no-derivative-control-3}
alone.

\emph{The estimate \textup{(R1)}.} Let $F$ be holomorphic on $D(0,R_{\mathrm{cal}})$ with
$|F|\le\eta$ there, and let $0<r'<R_{\mathrm{cal}}$. Cauchy's formula on the circle of radius
$r'$ gives
\begin{equation*}
F'(0)=\frac{1}{2\pi i}\oint_{|\hat{w}|=r'}\frac{F(\hat{w})}{\hat{w}^2}\,d\hat{w}.
\end{equation*}
The circle has length $2\pi r'$, and on it the integrand has modulus at most $\eta/r'^2$.
Hence $|F'(0)|\le\eta/r'$. Letting $r'\uparrow R_{\mathrm{cal}}$ gives
$|F'(0)|\le\eta/R_{\mathrm{cal}}$.
\end{proof}

\begin{remark}[Two routes to the derivative bound]\label{9.9:rem-derivative-routes}
Proposition~\ref{9.9:prop-remainder-derivative}, whose proof is outstanding, asks for a bound of
the form
\begin{equation*}
|E_n'(0)|\le c_E\,g_{k_0}^2\,\mathfrak{M}_n(0)
\end{equation*}
on the derivative at the bare offset $\hat w=0$ of the remainder $E_n=\varphi_n-\mathfrak{M}_n$ of
the two-point witness theorem along the frozen-shifted tuned family. Here $\varphi_n$ is the
calibrating two-point function along the tuned family, $\mathfrak{M}_n$ the frozen-shifted main
term, $k_0$ the witness level, $g_{k_0}$ the running coupling at that level and
$x_{k_0}=g_{k_0}^{-2}$. Remark~\ref{9.9:rem-no-derivative-control} shows that bounds on $E_n$ at real
offsets do not control $E_n'(0)$, and that such control must come along one of the two routes
displayed in \eqref{9.9:eq-no-derivative-control-a}. Below we say what each of the two requires.

We write $X=g^{-2}$. Throughout, $\Pi(X)$ stands for whatever polylogarithm of $X$ (a power of
$\log X$ times a constant) limits the source radius of the correlation theorem, and
$\bar{x}^{\Theta}_{s^{\circ}}$ for the upper end of the window, free of $n$, that contains the
inverse coupling $x_{k_0}$ of the tuned run at the depth $s^{\circ}=s(\mathsf{d}^{\circ})$ of the
witness. We write $\rho_W:=L^{s^{\circ}}\mathsf{d}^{\circ}$ for the radius of the bumps of the
reference pair, measured in cells of the witness lattice. In the estimates of the two-point witness
theorem quoted below, $p_0$ without argument abbreviates $p_0(g_{k_0})$, the value of the degree
function $p_0(\cdot)$ at the witness-level coupling. We write $A(1,U)=A(U)$ for the action at unit
coupling.

Route (R1): Cauchy's estimate at complex offsets. This is how the correlation theorem, including
its use inside the proof of the two-point witness theorem, obtains every source derivative: by
Cauchy's estimate on a disc, never through a covariance. Suppose that the identification
\begin{equation}\label{9.9:eq-derivative-routes-1}
\varphi_n(\hat w)=\mathfrak{M}_n(\hat w)\bigl(1+\varpi_n(\hat w)\bigr),\qquad
|\varpi_n(\hat w)|\le r_{\mathbb{C}},
\end{equation}
held for complex $\hat w$ on a disc of radius $R_{\mathrm{cal}}$ about $0$, with $\varphi_n$ and
$\mathfrak{M}_n$ holomorphic there; here $\varpi_n$ is the relative remainder at complex offsets and
$r_{\mathbb{C}}$ its relative precision. Then
\begin{equation}\label{9.9:eq-derivative-routes-2}
|E_n'(0)|\le\frac{\sup_{|\hat w|=R_{\mathrm{cal}}}|E_n(\hat w)|}{R_{\mathrm{cal}}}
\le\frac{r_{\mathbb{C}}\sup_{|\hat w|=R_{\mathrm{cal}}}|\mathfrak{M}_n(\hat w)|}{R_{\mathrm{cal}}},
\end{equation}
which has the form of the bound of Proposition~\ref{9.9:prop-remainder-derivative} (whose proof is
outstanding) with $c_E$ of the order of $r_{\mathbb{C}}\,x_{k_0}/R_{\mathrm{cal}}$, since
$g_{k_0}^2\mathfrak{M}_n=\mathfrak{M}_n/x_{k_0}$. Two things are then needed: a radius
$R_{\mathrm{cal}}$ proportional to $x_{k_0}$ up to a factor that the smallness of $r_{\mathbb{C}}$
compensates, and the two-point witness theorem at complex offsets.

On the radius. The natural large radius is that of the correlation theorem,
$R_{\mathrm{cal}}\asymp c_{\mathrm{rad}}X/\Pi(X)$, with $c_{\mathrm{rad}}$ the radius constant of
that theorem, proportional to its constant $\sigma_0$ and hence to $(LM)^{-4}$. Then
\begin{equation*}
c_E\approx\frac{r_{\mathbb{C}}\,\Pi(X)\,\bar{x}^{\Theta}_{s^{\circ}}}{c_{\mathrm{rad}}X};
\end{equation*}
note the factor $\bar{x}^{\Theta}_{s^{\circ}}/X$, which grows with the depth. The conditions on the
source radius $\varrho_J$ in the correlation theorem include $\tfrac{16}{3}\varrho_J\le\sigma_0g^{-2}$,
the condition $\tfrac{8}{3}\varrho_J\varepsilon_k^2\le1$ on the level-$k$ parameter $\varepsilon_k$
of the correlation theorem, and a family of conditions each requiring $\varrho_J g^2$ times a
polylogarithm to be at most a constant, one of them being, at each step $j$, a bound by
$C\varrho_Jg_j^2$ times the $2q$-th power of a logarithm, with $q$ one of the auxiliary parameters
$\mathfrak{p}$ and $C$ a constant of the correlation theorem; which of the other conditions of the
correlation theorem that involve the source radius is binding is not verified here. Thus
$r_{\mathbb{C}}$ must be small against $c_{\mathrm{rad}}/\Pi(X)$: a genuinely small relative
precision.

The terms of the remainder of the two-point witness theorem are not of that size as that theorem
bounds them. The irreducible term is bounded, relative to $\mathfrak{M}_n$, by
\begin{equation*}
\frac{R^{\mathrm{irr}}}{\mathfrak{M}_n}\le 7C_bC_0^{-2}
\bigl[p_0(g_{k_0})/p_0(\bar{x}_{s^{\circ}}^{-1/2})\bigr]^4,
\end{equation*}
with $\bar{x}_{s^{\circ}}$ the upper end of the first-passage envelope
$[\underline{x}_{s^{\circ}},\bar{x}_{s^{\circ}}]$ of the two-point witness theorem at depth
$s^{\circ}$, and $C_b$ the constant of the bound on irreducible terms in the Gaussian extraction;
inside the Gaussian-extraction term $R^{\mathrm{GX}}$ the locality term and the tower-tail term
satisfy
\begin{equation*}
\rho_W^8p_0^4e^{-c_aD_{\square}}\le A_0^4C_D^{-1},\qquad L^{-4(r-r_1)}\le L^{4-4r_0};
\end{equation*}
here $r$ is the extraction depth, $r_1$ a depth of the Gaussian extraction, $D_{\square}$ the side,
in cells, of its locality box, $c_a$ its locality rate, and $C_D$, like $C_0$ and $r_0$, one of the
room constants of the two-point witness theorem, that is, constants chosen before the coupling which
set the separation of the supports in cells of the witness lattice and the length of the Gaussian
segment;
these are small by the constants $C_0$, $C_D$, $r_0$, which are fixed before $g$, and not by $g$;
and the profile term $R^{\mathrm{prof}}$ is only of order $g$ times a polylogarithm divided by
$\rho_W$. All of these would need constants depending on $g$: $C_0$, $C_D$, $r_0$ enlarged with $X$,
that is, a deeper one-sided floor on the witness depth and a longer Gaussian segment, of the same
kind as the witness-depth floor $s_W(g)$ and the floor on the extraction depth $r(s)$ of the
two-point witness theorem itself. The remaining terms are powers of $g_{k_0}$ or of the form
$e^{-cp_0}$.

A third thing is needed before either: holomorphy of $\varphi_n$ itself on the large disc. The
results used here do not provide it. The uniqueness theorem gives holomorphy on the strip
$|\operatorname{Im}\hat w|<\upsilon_w$ only. The correlation theorem controls the ratio
$\tilde Z/\tilde Z(0)$, where $Z=e^{\mathcal{E}}\tilde Z$ with $\mathcal{E}$ the boundary-term part
of $\log Z$ described below, through the Schwarz lemma by the bound
$|\tilde Z/\tilde Z(0)|\le\mathsf{Q}$, where $\mathsf{Q}$ grows exponentially with the volume of
the torus.
Holomorphy on the large disc would therefore have to be part of the statement of a sharp form of
the two-point witness theorem adapted to this route.

On complexification. The identification in the two-point witness theorem, analysed in
Remark~\ref{9.9:rem-no-derivative-control}, is a statement about a real measure. At complex $\hat w$
of modulus $\asymp X/\Pi(X)$ the weight at level $k_0$ carries the phase
$\operatorname{Im}\zeta_{k_0}(\hat w)\,A(1,U_{k_0})$. This phase is of order
$\rho_W^4p_0^2/\Pi(X)$ over the witness region on the analysis window of the covariance comparison,
and it is unbounded over the torus. Hence $\tilde\nu^{\hat w}$, normalised globally, has total
variation exponentially large in the volume compared with its mass.

The correlation theorem itself encounters the same obstruction. Its torus bound controls source
derivatives of
$\log\tilde Z$ through the globally normalised ratio $\mathsf{h}:=\tilde Z/\tilde Z(0)$. It uses
the volume-dependent bound $|\mathsf{h}|\le\mathsf{Q}$ and the Schwarz lemma on a polydisc shrunk by
the factor $2\mathsf{Q}+2$. This is why the constant of the torus bound depends on the torus exponent
$m$. Correspondingly, the constant of the uniqueness theorem through which this bound enters
contains the term $\log2\cdot\bigl((4\mathsf{Q}_w+4)\mathsf{p}/r_J\bigr)^{\mathsf{p}}$, which
depends on the volume, beside a term free of the volume; here $r_J$ is the source radius of the
correlation theorem as the uniqueness theorem reads it.

Only the boundary-term part $\mathcal{E}=\sum_j\sum_{\mathsf{Y}}e_j(\mathsf{Y};z)$ of $\log Z$
is localised by the polymer structure. A mixed source derivative
of $e_j(\mathsf{Y};z)$ is non-zero only if the enlargement $\mathsf{Y}^+$ of the polymer
$\mathsf{Y}$ meets every support. The volume-free bounds are a separate theorem on $\mathbb{R}^4$,
the per-ball volume-free bound, whose constant does not depend on $m$, and they again concern
separated local sources. A constant source is neither local nor separated from anything.

So route (R1) requires the lower bound and the main-term identification of the two-point witness
theorem to be re-proved, at a complex base point, in the logarithmic and polymer-expanded form in
which the correlation theorem is stated. This is the sharp form of the two-point witness theorem,
in which the relative remainder is of order a positive power of the coupling: at depth $s$ the
coefficient of the kernel $\mathcal{K}$ in the connected two-point function $S^T_{2;K}$ is
\begin{equation*}
b_0\,g_{K-s}^4\bigl(1+O(g_{K-s}^{\sigma})\bigr)
\end{equation*}
for an exponent $\sigma>0$. A proof of this form would require the coupling increment computed to
two loops, including $N_{K-s}\to1$, and the Gaussian extraction in relative precision; the
two-point witness theorem provides neither. The
covariance comparison used in its proof is absolute at the natural scale, and its sharp relative
form is not part of that theorem.

Route (R2): real differentiation term by term. This is the literal form of one more order of the
bookkeeping of the two-point witness theorem. For real $\hat w$ the family $\tilde\nu^{\hat w}$
remains a probability law. Here $w$ denotes the constant source of strength $\hat w$, and
$\Xi^{\hat w}$ ranges over the local functionals
of the bookkeeping described in Remark~\ref{9.9:rem-no-derivative-control}: $\Psi^h$, the
summands of the splitting of $\Phi^h_i$, and products of one summand of the splitting of
$\Phi^h_1$ with one summand of the splitting of $\Phi^h_2$. For each such $\Xi^{\hat w}$,
\begin{equation}\label{9.9:eq-derivative-routes-3}
\frac{d}{d\hat w}\Big|_{0}\mathbb{E}_{\tilde\nu^{\hat w}}\bigl[\Xi^{\hat w}\bigr]
=\mathbb{E}_{\tilde\nu}\bigl[\partial_{\hat w}\Xi\bigr]+\operatorname{Cov}_{\tilde\nu}(\Phi_0,\Xi).
\end{equation}
Here $\Phi_0$ is the response to a constant source:
\begin{equation}\label{9.9:eq-derivative-routes-4}
\begin{split}
\Phi_0^h:=\partial_{\hat w}\log\tau_h^{\hat w}\big|_0
={}&N_{k_0}\,A(1,U_{k_0})
+\partial_{\hat w}\bigl[\mathbf{E}^w+\mathbf{R}^w+\mathbf{B}^w-E^w\bigr]\big|_0\\
&+\partial_{\hat w}\log\mathbf{T}^w_h\big|_0 .
\end{split}
\end{equation}
The lightface $E^w$ in \eqref{9.9:eq-derivative-routes-4} is the term of the format of the
correlation theorem so written, distinct from the sourced large-field operation $\mathbf{E}^w$.
The difference terms $\mathbf{D}^w$ in the format of the correlation theorem vanish at constant
sources, so they do not appear in \eqref{9.9:eq-derivative-routes-4}.

(R2-i) The parts $\mathbb{E}_{\tilde\nu}[\partial_{\hat w}\Xi]$. Every constituent of the local
functionals $\Xi$ is holomorphic in $\hat w$ on $D(0,2\varrho_J)$, by the holomorphy clauses of the
inductive hypothesis of the correlation theorem. These constituents are the frozen families of
effective actions, the sourced large-field factors $\mathbf{T}^w_h$, the profile
$\partial_z\zeta_{k_0}(p)$ of the running shift, and the remaining local constituents of the
decomposition. Cauchy's estimate at a radius of the large kind, $c_{\mathrm{rad}}X/\Pi(X)$,
provided the correlation theorem is run at that source radius so that these constituents are
holomorphic on the disc of that radius (the same proviso as in route (R1), including the choice of
$\Pi(X)$), costs per derivative a relative factor
\begin{equation}\label{9.9:eq-derivative-routes-5}
\frac{\Pi(X)}{c_{\mathrm{rad}}X}\le g_{k_0}^2\,
\frac{\Pi(X)\,\bar{x}^{\Theta}_{s^{\circ}}}{c_{\mathrm{rad}}X}.
\end{equation}
This is harmless for the terms that are small against $c_{\mathrm{rad}}/\Pi(X)$. It is also
harmless for $R^{\mathrm{irr}}$ and $R^{\mathrm{prof}}$ after the same $g$-dependent enlargement of
$C_0$, $C_D$, $r_0$ as in route (R1). At the fixed radius $\varrho_J^{0}$ of the correlation theorem
the cost is $1/\varrho_J^{0}$ per derivative, which does not reach $g_{k_0}^2$.

(R2-ii) The parts $\operatorname{Cov}_{\tilde\nu}(\Phi_0,\Xi)$. The functional $\Phi_0$ is a sum
$\sum_\Delta\phi_0(\Delta)$ over every cell $\Delta$ of the torus at level $k_0$, and one needs,
with a constant $c'$,
\begin{equation}\label{9.9:eq-derivative-routes-6}
\sum_{\text{cells }\Delta}\bigl|\operatorname{Cov}_{\tilde\nu}(\phi_0(\Delta),\Xi)\bigr|
\le c'\,g_{k_0}^2\cdot(\text{the bound on the term }\Xi).
\end{equation}
This is a decay statement for the connected correlation, under $\tilde\nu$, between the action
density at a cell far from the supports and a functional localised at the supports. The decay must
be summable over the four-dimensional torus, that is, faster than the inverse fourth power of the
distance, uniformly in the torus exponent.

The Gaussian part is exact and harmless. Let $C^{\mathrm{bg}}$ be the covariance of the background
part of the unit lattice Maxwell curvature (the background Gaussian of the tower model). Under this
Gaussian, the covariance of $\sum_pa_p$, with $a_p$ the action density at the plaquette $p$, with a
product of two local quadratic functionals at plaquettes $p_1$, $p_2$ is the triangle
\begin{equation}\label{9.9:eq-derivative-routes-7}
\sum_pC^{\mathrm{bg}}(p_1,p)\,C^{\mathrm{bg}}(p,p_2)\,C^{\mathrm{bg}}(p_2,p_1).
\end{equation}
The covariance bound of the tower model, with $C_u$ and $\eta$ the constants of the kernel estimate
of the two-point witness theorem,
\begin{equation*}
|C^{\mathrm{bg}}(p,q)|\le C_u\eta^4(1+|x(p)-x(q)|)^{-4},
\end{equation*}
makes the sum over $p$ converge in dimension four. This merely reproduces the statement that in the
tower model the coupling enters as the explicit prefactor $x^{-2}$, with $x$ the inverse coupling
of that model, in the sense in which $x_k$ denotes the running inverse coupling.

The corrections beyond the Gaussian part are where the route meets resistance. The covariance
comparison used
in the proof of the two-point witness theorem is absolute at the natural scale, with no decay in the
distance, so it cannot be summed over the torus. The locality item of the Gaussian extraction, with
error $p_0^4e^{-c_aD_{\square}}$, decouples the exterior of a box of $D_{\square}$ cells at cost
$e^{-c_aD_{\square}}$. Were it a bound on the influence of each single cell, it would give
\eqref{9.9:eq-derivative-routes-6} for each torus, but with a constant depending on the volume, not
uniformly in the torus exponent. Whether the proof of the Gaussian extraction yields a bound of that
kind is not verified here.

A decoupling of the action density at level $k_0$ from the local functionals $\Xi$ that is summable
cell by cell is an exponential, or faster than inverse-fourth-power, clustering statement for
$\tilde\nu$ on these functionals. It is a statement of the kind that cluster expansions at the
witness lattice prove, and it is the same object that route (R1) needs.
\end{remark}

\begin{proof}
We establish the identities and estimates used in the remark:
\eqref{9.9:eq-derivative-routes-2} with the resulting form of $c_E$,
\eqref{9.9:eq-derivative-routes-3} with the first term of \eqref{9.9:eq-derivative-routes-4},
\eqref{9.9:eq-derivative-routes-5} with the cost of a derivative, and the convergence of the sum
\eqref{9.9:eq-derivative-routes-7}.

Cauchy's estimate in route (R1). Assume $E_n$ holomorphic on an open neighbourhood of the closed
disc $|\hat w|\le R_{\mathrm{cal}}$. Cauchy's formula gives
\begin{equation*}
E_n'(0)=\frac{1}{2\pi i}\oint_{|\hat w|=R_{\mathrm{cal}}}\frac{E_n(\hat w)}{\hat w^2}\,d\hat w.
\end{equation*}
The circle has length $2\pi R_{\mathrm{cal}}$ and $|\hat w|^{-2}=R_{\mathrm{cal}}^{-2}$ on it, so
$|E_n'(0)|\le\sup_{|\hat w|=R_{\mathrm{cal}}}|E_n(\hat w)|/R_{\mathrm{cal}}$. This is the first
inequality of \eqref{9.9:eq-derivative-routes-2}.

By \eqref{9.9:eq-derivative-routes-1}, $E_n=\varphi_n-\mathfrak{M}_n=\mathfrak{M}_n\varpi_n$ on the
disc, so $|E_n|\le r_{\mathbb{C}}|\mathfrak{M}_n|$ there. This is the second inequality.

Since $x_{k_0}=g_{k_0}^{-2}$, the right-hand side of \eqref{9.9:eq-derivative-routes-2} equals
\begin{equation*}
\frac{r_{\mathbb{C}}\,x_{k_0}}{R_{\mathrm{cal}}}\cdot
g_{k_0}^2\sup_{|\hat w|=R_{\mathrm{cal}}}|\mathfrak{M}_n(\hat w)|.
\end{equation*}
This is the statement that $c_E$ is of the order of $r_{\mathbb{C}}x_{k_0}/R_{\mathrm{cal}}$, up to
the factor
\begin{equation*}
\sup_{|\hat w|=R_{\mathrm{cal}}}|\mathfrak{M}_n(\hat w)|\big/\mathfrak{M}_n(0),
\end{equation*}
which the words ``of the order of'' in the remark absorb.

With $R_{\mathrm{cal}}=c_{\mathrm{rad}}X/\Pi(X)$ and $x_{k_0}\le\bar{x}^{\Theta}_{s^{\circ}}$, we
obtain
\begin{equation*}
\frac{r_{\mathbb{C}}\,x_{k_0}}{R_{\mathrm{cal}}}
=\frac{r_{\mathbb{C}}\,\Pi(X)\,x_{k_0}}{c_{\mathrm{rad}}X}
\le\frac{r_{\mathbb{C}}\,\Pi(X)\,\bar{x}^{\Theta}_{s^{\circ}}}{c_{\mathrm{rad}}X},
\end{equation*}
the displayed form of $c_E$ in route (R1).

The derivative formula. The law $\tilde\nu^{\hat w}$ on pairs (history $h$, block field $V$) at
level $k_0$ has density $\tau_h^{\hat w}/Z^{\hat w}$, where $Z^{\hat w}$ is the sum over histories
and the integral over block fields of $\tau_h^{\hat w}$. Consider $\hat w$ real. On the finite torus
the sum over histories is finite, the block fields range over a compact product of copies of
$\mathrm{SU}(N)$, and $\tau_h^{\hat w}$ is holomorphic in $\hat w$ by the holomorphy clauses of the
inductive hypothesis of the correlation theorem; so differentiation at $\hat w=0$ may be carried
under this sum and integral. Write $\mu$ for the
reference measure (sum over $h$, integral over $V$). Then
\begin{equation*}
\mathbb{E}_{\tilde\nu^{\hat w}}[\Xi^{\hat w}]
=\frac{\int\Xi^{\hat w}\tau_h^{\hat w}\,d\mu}{Z^{\hat w}}.
\end{equation*}
By the quotient rule and $\partial_{\hat w}\tau_h^{\hat w}=\tau_h^{\hat w}\,\partial_{\hat w}\log\tau_h^{\hat w}$,
its derivative at $0$ is
\begin{equation*}
\mathbb{E}_{\tilde\nu}[\partial_{\hat w}\Xi]+\mathbb{E}_{\tilde\nu}[\Xi\,\Phi_0]
-\mathbb{E}_{\tilde\nu}[\Xi]\,\mathbb{E}_{\tilde\nu}[\Phi_0],
\end{equation*}
with $\Phi_0=\partial_{\hat w}\log\tau_h^{\hat w}|_0$. The last two terms are
$\operatorname{Cov}_{\tilde\nu}(\Phi_0,\Xi)$, which proves \eqref{9.9:eq-derivative-routes-3}.

Along the constant source the effective inverse coupling at level $k_0$ is
$x_{k_0}-\zeta_{k_0}(\hat w)$, and $\zeta_{k_0}'(0)=N_{k_0}$. The action $A(x,U)$ with inverse
coupling $x$ (constant, or position-dependent as in Definition~\ref{1.5:def-domains}) is, plaquette
by plaquette, the inverse coupling times the unit-coupling summand, hence linear in $x$, with
$A(1,U)=A(U)$. Differentiating the term $-A(x_{k_0}-\zeta_{k_0}(\hat w),U_{k_0})$ of
$\log\tau_h^{\hat w}$ at $0$ therefore gives
\begin{equation*}
\partial_{\hat w}\bigl[-A(x_{k_0}-\zeta_{k_0}(\hat w),U_{k_0})\bigr]\big|_0
=\zeta_{k_0}'(0)\,A(1,U_{k_0})=N_{k_0}\,A(1,U_{k_0}),
\end{equation*}
the first term of \eqref{9.9:eq-derivative-routes-4}. The remaining terms of
\eqref{9.9:eq-derivative-routes-4} are the derivatives at $0$ of the other $\hat w$-dependent
factors of $\tau_h^{\hat w}$. The difference terms $\mathbf{D}^w$ contribute nothing, because they
vanish at constant sources.

The cost of a derivative in (R2-i). Let $F$ be holomorphic in $\hat w$ on an open neighbourhood of
the closed disc of radius $R'$, with $|F|\le B$ there. The Cauchy estimate of the first part of this
proof gives $|F'(0)|\le B/R'$. Relative to the bound $B$, the cost of one derivative is therefore
$1/R'$.

For $R'=c_{\mathrm{rad}}X/\Pi(X)$ this cost is $\Pi(X)/(c_{\mathrm{rad}}X)$. Since
$g_{k_0}^2\bar{x}^{\Theta}_{s^{\circ}}\ge g_{k_0}^2x_{k_0}=1$, we have
\begin{equation*}
\frac{\Pi(X)}{c_{\mathrm{rad}}X}\le g_{k_0}^2\bar{x}^{\Theta}_{s^{\circ}}\cdot
\frac{\Pi(X)}{c_{\mathrm{rad}}X},
\end{equation*}
which is \eqref{9.9:eq-derivative-routes-5}. For $R'=\varrho_J^{0}$ the cost is $1/\varrho_J^{0}$,
a number fixed independently of $g$.

Convergence of the triangle sum. Here $|x(p)-x(q)|$ is the distance, in units of the spacing of the
lattice that carries the background Gaussian, between the barycentres of $p$ and $q$. By the
covariance bound of the tower model,
\begin{equation*}
\sum_p\bigl|C^{\mathrm{bg}}(p_1,p)\,C^{\mathrm{bg}}(p,p_2)\bigr|
\le C_u^2\eta^8\sum_p(1+|x(p)-x(p_1)|)^{-4}(1+|x(p)-x(p_2)|)^{-4}.
\end{equation*}
The product of two positive numbers is at most the square of the larger, and
$\max\{a^2,b^2\}\le a^2+b^2$. Hence each summand on the right is at most
\begin{equation*}
(1+|x(p)-x(p_1)|)^{-8}+(1+|x(p)-x(p_2)|)^{-8}.
\end{equation*}

Let
\begin{equation*}
c_4:=\sup_{y}\sup_{j\ge0}(1+j)^{-3}\,\#\{p:\ |x(p)-y|\in[j,j+1)\},
\end{equation*}
the supremum being over points $y$, and the count over plaquettes $p$ of a unit lattice; in
dimension four, $c_4$ is finite and depends only on the dimension.
This bound also holds on the torus, whose distance shells contain no more
plaquettes than those of the infinite lattice. Therefore, for $i=1,2$,
\begin{equation*}
\sum_p(1+|x(p)-x(p_i)|)^{-8}\le c_4\sum_{j\ge0}(1+j)^3(1+j)^{-8}
=c_4\sum_{j\ge0}(1+j)^{-5}<\infty,
\end{equation*}
uniformly in the size of the torus. The remaining factor $|C^{\mathrm{bg}}(p_2,p_1)|$ of
\eqref{9.9:eq-derivative-routes-7} is at most $C_u\eta^4$. So the sum over $p$ in
\eqref{9.9:eq-derivative-routes-7} converges in dimension four, with a bound independent of the
torus.
\end{proof}

\begin{remark}[Gaussian skeleton and in-box precision]\label{9.9:rem-gaussian-skeleton}
This remark concerns the second route of Remark~\ref{9.9:rem-derivative-routes}, real
differentiation of the remainder of the two-point witness theorem member by member, towards the
bare-offset derivative bound
\begin{equation*}
|E_n'(0)|\le c_E\,g_{k_0}^2\,\mathfrak{M}_n(0)
\end{equation*}
of Proposition~\ref{9.9:prop-remainder-derivative}, whose proof is outstanding. Each of the
clauses (i)--(iii) below carries its own standing, and none of them is used as an input of the
statements proved in this section. Let $\Phi_0$ be the response to a constant source of
Remark~\ref{9.9:rem-derivative-routes}. Split $\Phi_0=\sum_c\phi_0(c)$ over the cells $c$ of the
witness lattice, the lattice of level $k_0$ on which the law $\tilde{\nu}$ of the two-point
witness theorem lives. Let $\rho_{\mathrm{W}}$ be the radius, in cells of the witness lattice, of
the envelope of the supports at the witness scale, as in the two-point witness theorem.

\smallskip
\noindent(i) \emph{The Gaussian skeleton is exact.} Let $y$ be a centred Gaussian vector with
covariance $\mathsf{t}^{-1}\Sigma$, where $\Sigma$ is positive definite and
$\mathsf{t}=x_{k_0}-\zeta_{k_0}(\hat{w})$. Write $\gamma$ for its law and
$\mathcal{S}(y)=\tfrac12\,y^{\top}\Sigma^{-1}y$ for the Gaussian's own quadratic form, so that
$\gamma$ has density proportional to $e^{-\mathsf{t}\mathcal{S}}$. By the definition of the bare
offset, as it enters the response $\Phi_0$ of Remark~\ref{9.9:rem-derivative-routes}, the
$\hat{w}$-derivative of $\zeta_{k_0}$ at $0$ is $N_{k_0}$ and $\zeta_{k_0}(0)=0$, so that
$\mathsf{t}=x_{k_0}=g_{k_0}^{-2}$ at $\hat{w}=0$. For quadratic forms
$\mathcal{Q}_i(y)=\tfrac12\,y^{\top}\mathcal{B}_iy$ with $\mathcal{B}_i$ symmetric ($i=1,2$),
\begin{equation}\label{9.9:eq-gaussian-skeleton-1}
\frac{d}{d\hat{w}}\Big|_{0}\operatorname{Cov}_{\gamma}(\mathcal{Q}_1,\mathcal{Q}_2)
=\kappa_3^{\gamma}(N_{k_0}\mathcal{S},\mathcal{Q}_1,\mathcal{Q}_2)
=2N_{k_0}\,g_{k_0}^2\operatorname{Cov}_{\gamma}(\mathcal{Q}_1,\mathcal{Q}_2);
\end{equation}
the coupling enters as an explicit prefactor. Let $\mathcal{N}_i$ be the set of cells carrying
$\mathcal{Q}_i$, and let $\phi_0^{\gamma}(c)=\tfrac12\,y^{\top}\mathcal{B}_cy$ be the Gaussian part of
$\phi_0(c)$, the share of the cell $c$ in $N_{k_0}\mathcal{S}$, so that
$\sum_c\mathcal{B}_c=N_{k_0}\Sigma^{-1}$ by the splitting of $\Phi_0$. Then, cell by cell,
$\kappa_3^{\gamma}(\phi_0^{\gamma}(c),\mathcal{Q}_1,\mathcal{Q}_2)$ is the cyclic trace
$\operatorname{tr}(\mathcal{B}_cC\mathcal{B}_1C\mathcal{B}_2C)$, where $C$ is the covariance of
$\gamma$. By the fourth-power
decay of $C^{\mathrm{bg}}$ in the kernel estimates of the two-point witness theorem, this trace is
bounded by a constant times
$(1+\operatorname{dist}(c,\mathcal{N}_1))^{-4}(1+\operatorname{dist}(c,\mathcal{N}_2))^{-4}$,
which is summable over $c\in\mathbb{Z}^4$. This bound is proved as an implication from the
identification of the unit background Gaussian's action in the block field with the quadratic form
whose covariance is the background covariance $C^{\mathrm{bg}}$. (The kernel of the two-point
witness theorem is defined through the decomposition $P_{\mathrm{ex}}=C^{\mathrm{bg}}+C^{\mathrm{fl}}$
of the witness theorem's covariance into its background and fluctuation parts; the identification
is a statement in addition to that decomposition.)

Whether the share of the far cells in this sum is carried by the box of $D_{\square}$ cells of the
Gaussian extraction is not decided here. No tail of order
$(\rho_{\mathrm{W}}/D)^4$ is exhibited here: at that order only an upper bound is obtained, and the
coefficient of that order vanishes: the orientation trace of the short-distance
germ of the witness kernel, $\operatorname{tr}(\Lambda^2I)$ with
$I=\delta-2\,\hat{n}\otimes\hat{n}$, is zero; here $\Lambda$ is the matrix of that germ, $\delta$
the identity matrix and $\hat{n}$ the unit vector in the direction of the displacement variable of
the germ.

\smallskip
\noindent(ii) \emph{In-box precision.} Consider a box of $D$ cells about the
supports; it has $\mathcal{V}(D)$ cells, with $\mathcal{V}(D)$ of the order of
$(14\rho_{\mathrm{W}}+2D)^4$.
We recall the covariance comparison on the analysis window $\{\max_i|y_i|<\mathfrak{b}\}$,
$\mathfrak{b}=\mathfrak{b}_{k_0}$, used in the proof of the two-point witness theorem, in the
following form: for functionals $\Gamma_1,\Gamma_2$ of the block field on that window,
\begin{equation}\label{9.9:eq-gaussian-skeleton-2}
\bigl|\operatorname{Cov}_{\langle\cdot\rangle}(\Gamma_1,\Gamma_2)
-\operatorname{Cov}_{\gamma}(\Gamma_1,\Gamma_2)\bigr|
\le 112\,\bigl(\delta_{\mathrm{ng}}+\epsilon_{\mathfrak{b}}^{1/2}\bigr)
\,M_4(\Gamma_1-a_1)\,M_4(\Gamma_2-a_2).
\end{equation}
Here:
\begin{itemize}
\item $\delta_{\mathrm{ng}}\le C_{\square}\,g_{k_0}\mathfrak{b}_{k_0}^3$ is the size of the
part of the action on the window that is not Gaussian, $C_{\square}$ being the constant of the
comparison's bound on that part;
\item $\epsilon_{\mathfrak{b}}:=2n_{\mathrm{var}}e^{-\mathfrak{b}^2/(2v)}$, where
$n_{\mathrm{var}}$ is the number of variables $y_i$ and $v$ the variance parameter entering the
Gaussian tail $e^{-\mathfrak{b}^2/(2v)}$; the comparison is stated under the condition
$\epsilon_{\mathfrak{b}}\le\tfrac12$;
\item $M_4$ denotes fourth moments and $a_1,a_2$ are centrings;
\item $\operatorname{Cov}_{\langle\cdot\rangle}$ is the covariance under the full law of the block
field.
\end{itemize}
Let $G_1,G_2$ be the local functionals of the two-point witness theorem attached to $f_1,f_2$,
and let $\mathbb{E}=\mathbb{E}_{\tilde{\nu}}$. Make two assumptions, of the same kind as those
under which \eqref{9.9:eq-gaussian-skeleton-2} enters the proof of the two-point witness theorem;
neither of them is among the statements established for the two-point witness theorem:
\begin{itemize}
\item[(a)] \eqref{9.9:eq-gaussian-skeleton-2} applies pair by pair to
$(\phi_0(c),(G_1-\mathbb{E}G_1)(G_2-\mathbb{E}G_2))$ for the cells $c$ of the box, with the
fourth moment of the centred product $(G_1-\mathbb{E}G_1)(G_2-\mathbb{E}G_2)$ at the moment scale
of the local functionals under which \eqref{9.9:eq-gaussian-skeleton-2} enters the proof of the
two-point witness theorem, namely a polylogarithm of $g_{k_0}^{-2}$ times
$g_{k_0}^4\|f_1\|\,\|f_2\|$, with $\|f_i\|$ denoting the test-function norm in which the
two-point witness theorem states its bounds;
\item[(b)] one cell's action density on the analysis window has the per-cell scale
$M_4(\phi_0(c)-a_c)\lesssim g_{k_0}^2\,p_0(g_{k_0})^2$.
\end{itemize}
Then the in-box comparison error sums to
\begin{equation}\label{9.9:eq-gaussian-skeleton-3}
\mathcal{V}(D)\cdot\mathrm{polylog}\cdot
\bigl(g_{k_0}\mathfrak{b}^3+\epsilon_{\mathfrak{b}}^{1/2}\bigr)\cdot g_{k_0}^6\,\|f_1\|\,\|f_2\|,
\end{equation}
against a target share $\asymp g_{k_0}^6\|f_1\|\|f_2\|$; here and below $\mathrm{polylog}$
denotes a fixed polylogarithm of $g_{k_0}^{-2}$, collecting $p_0(g_{k_0})^2$ and constants. This
error is of the same kind as the
members of the remainder of the two-point witness theorem that are closed by a ceiling on the
coupling, and, for $D$ at most a polylogarithm of $g_{k_0}^{-2}$ (in particular for the box of the
Gaussian extraction), it is closed by such a ceiling uniformly in $s$, $K$, $m$. So, under these
assumptions, the in-box precision is not where the route resists: one power of $g_{k_0}$ is to
spare.

\smallskip
\noindent(iii) \emph{Where the available estimates fall short.} At two places the estimates
available on the inputs used in this chapter do not reach the required bound. This shows only that
the route does not prove the bound of Proposition~\ref{9.9:prop-remainder-derivative},
whose proof is outstanding; it does not show that the quantities are large.

\emph{The complement of the box.} The same per-cell absolute bound has no decay, and summed over
the witness torus it gives a quantity proportional
to the torus volume. Moreover, \eqref{9.9:eq-gaussian-skeleton-2} as stated does not apply
torus-wide: its condition $\epsilon_{\mathfrak{b}}\le\tfrac12$ reads the number of variables.

\emph{The explicit $\hat{w}$-dependence.} The second place is the explicit $\hat{w}$-dependence,
at the fixed radius $\varrho_J^{0}$ of the correlation theorem, of the sourced local constituents
of the two-point witness theorem. These are the $z$-derivatives of $\mathbf{E}^w$, $\mathbf{R}^w$,
$\mathbf{B}^w$ and $\log\mathbf{T}^w_h$ that make up the members $\mathrm{Irr}_i$ and
$\mathring{S}_i$ of the splitting of the source derivatives $\Phi^h_i$ in the two-point witness
theorem, and the spanning terms. They are analytic in $w$ only on the source domains
$\mathcal{W}_j(X)$ of the
correlation theorem. These domains are contained in the discs $\{|w|<r_j\}$, with $r_j$ of the
order of $\varrho_J^{0}$ and $r_j\in[r_J,2r_J)$ for the radius parameter $r_J$ of the correlation
theorem. One $\hat{w}$-derivative of these costs a
factor $8/\varrho_J^{0}$, not $g_{k_0}^2$. Only the pure-channel constituents sit on the polydiscs
$\mathfrak{P}_k=\prod_{j\le k}\{w_j\in\mathbb{C}:|w_j|<\sigma_0x_j\}$, with $\sigma_0$ the radius
constant of the correlation theorem, where the cost is $(16/(3\sigma_0))\,g_{k_0}^2$.

Consider the members that are small by room only: $R^{\mathrm{irr}}$, bounded by
$7C_b\rho_{\mathrm{W}}^{-2}p_0(g_{k_0})^4$, with $C_b$ the constant of the Gaussian extraction's
bound for the irrelevant part; and the tower-tail, window and locality items of the
Gaussian-extraction member $R^{\mathrm{GX}}$, bounded by $L^{4-4r_0}$ and $A_0^4C_D^{-1}$, with
$C_D$ the locality room constant of the Gaussian extraction, $c_a$ its locality decay rate (the
locality item carries the factor $e^{-c_aD_{\square}}$) and $r_1$ the reference depth of the
tower-tail item, which decays in the difference $r(s)-r_1$, $r(s)$ being the extraction depth of
clause~(iv) of Theorem~\ref{3.1:thm-main}. For
these, the explicit part of the derivative meets the share $c_Eg_{k_0}^2$ of the bound of
Proposition~\ref{9.9:prop-remainder-derivative} (proof outstanding) only in one of
two ways. The first is with room growing like a power of $x_{k_0}$:
\begin{equation}\label{9.9:eq-gaussian-skeleton-4}
\rho_{\mathrm{W}}\asymp x_{k_0}^{1/2}\cdot\mathrm{polylog},\qquad
r(s)-r_1\gtrsim\tfrac14\log_L\bar{x}_s,\qquad
D_{\square}\gtrsim c_a^{-1}\log\bar{x}_s .
\end{equation}
This is not the kind of room of the two-point witness theorem. The second is a source radius for
the local terms proportional to $x_{k_0}$, that is, the correlation theorem read at a large radius.

The members whose bounds carry an exponential in $p_0$ or $\rho_{\mathrm{W}}$, namely the member
$R^{\mathrm{BAD}}$ of coarse live or healed labels meeting the supports and the member
$R^{\mathrm{loc}}$ of bi-sourced local terms spanning both supports, whose bound is exponentially
small in $\rho_{\mathrm{W}}$, absorb the factor $8/\varrho_J^{0}$ by one more power in their
ceilings:
\begin{itemize}
\item one more factor $\log\bar{x}_s$ in one of the ceiling conditions of the two-point witness
theorem that close these members;
\item one more power of the variable of the coupling-ceiling condition of the two-point witness
theorem that closes these members.
\end{itemize}
This holds granted the inputs of those members in the form in which the two-point witness theorem
uses them.
\end{remark}

\begin{proof}
\emph{Clause (i), first identity.} We use a general fact about tilted laws. Let
$d\nu_u\propto e^{u\Xi}\,d\nu$ be a family of probability laws. Then
$\frac{d}{du}\mathbb{E}_u[Y]=\operatorname{Cov}_u(\Xi,Y)$. Applying this to $Y_1Y_2$, $Y_1$ and
$Y_2$ gives
\begin{equation*}
\begin{split}
\frac{d}{du}\operatorname{Cov}_u(Y_1,Y_2)
&=\operatorname{Cov}_u(\Xi,Y_1Y_2)-\operatorname{Cov}_u(\Xi,Y_1)\,\mathbb{E}_uY_2
-\mathbb{E}_uY_1\,\operatorname{Cov}_u(\Xi,Y_2)\\
&=\kappa_3(\Xi,Y_1,Y_2).
\end{split}
\end{equation*}
The $\hat{w}$-dependence of $\gamma$ is through the weight $e^{-\mathsf{t}(\hat{w})\mathcal{S}}$.
The $\hat{w}$-derivative of the exponent at $0$ is
$-\partial_{\hat{w}}\mathsf{t}|_0\,\mathcal{S}=\partial_{\hat{w}}\zeta_{k_0}|_0\,\mathcal{S}
=N_{k_0}\mathcal{S}$. This gives the first equality in \eqref{9.9:eq-gaussian-skeleton-1}.

\emph{Clause (i), second identity.} We compute the cumulants of quadratic forms. Write
$C=\mathsf{t}^{-1}\Sigma$, and let $\mathcal{J}_1,\mathcal{J}_2,\mathcal{J}_3$ be symmetric. For
small real $u_j$,
\begin{equation*}
\begin{split}
\log\mathbb{E}_{\gamma}\exp\Bigl(\sum_j\tfrac{u_j}2\,y^{\top}\mathcal{J}_jy\Bigr)
&=-\tfrac12\log\det\Bigl(1-C\sum_ju_j\mathcal{J}_j\Bigr)\\
&=\tfrac12\sum_{\ell\ge1}\tfrac1\ell\operatorname{tr}\Bigl(\bigl(C\textstyle\sum_ju_j\mathcal{J}_j\bigr)^\ell\Bigr).
\end{split}
\end{equation*}
The coefficient of $u_1u_2$ gives
\[
\operatorname{Cov}_{\gamma}\bigl(\tfrac12y^{\top}\mathcal{J}_1y,\tfrac12y^{\top}\mathcal{J}_2y\bigr)
=\tfrac12\operatorname{tr}(\mathcal{J}_1C\mathcal{J}_2C).
\]
The coefficient of $u_1u_2u_3$
collects three copies each of $\operatorname{tr}(C\mathcal{J}_1C\mathcal{J}_2C\mathcal{J}_3)$ and
$\operatorname{tr}(C\mathcal{J}_1C\mathcal{J}_3C\mathcal{J}_2)$. These two traces are equal, by
transposition and cyclicity.
Hence
\[
\begin{split}
\kappa_3^{\gamma}\bigl(\tfrac12y^{\top}\mathcal{J}_1y,\tfrac12y^{\top}\mathcal{J}_2y,
\tfrac12y^{\top}\mathcal{J}_3y\bigr)
&=\tfrac12\cdot\tfrac13\cdot6\operatorname{tr}(\mathcal{J}_1C\mathcal{J}_2C\mathcal{J}_3C)\\
&=\operatorname{tr}(\mathcal{J}_1C\mathcal{J}_2C\mathcal{J}_3C).
\end{split}
\]
We take $\mathcal{J}_1=\Sigma^{-1}$, the matrix of $\mathcal{S}$, and
$\mathcal{J}_2=\mathcal{B}_1$, $\mathcal{J}_3=\mathcal{B}_2$, and use $\Sigma^{-1}C=\mathsf{t}^{-1}$.
This gives
\[
\kappa_3^{\gamma}(\mathcal{S},\mathcal{Q}_1,\mathcal{Q}_2)
=\mathsf{t}^{-1}\operatorname{tr}(\mathcal{B}_1C\mathcal{B}_2C)
=2\mathsf{t}^{-1}\operatorname{Cov}_{\gamma}(\mathcal{Q}_1,\mathcal{Q}_2).
\]
At $\hat{w}=0$ we have $\mathsf{t}^{-1}=g_{k_0}^2$. Multiplying by $N_{k_0}$ gives the second
equality in \eqref{9.9:eq-gaussian-skeleton-1}.

\emph{Clause (i), cell by cell.} The same cumulant formula with $\mathcal{J}_1=\mathcal{B}_c$ gives
\begin{equation*}
\kappa_3^{\gamma}(\phi_0^{\gamma}(c),\mathcal{Q}_1,\mathcal{Q}_2)
=\operatorname{tr}(\mathcal{B}_cC\mathcal{B}_1C\mathcal{B}_2C).
\end{equation*}
Under the identification named in (i), $C$ is the background covariance $C^{\mathrm{bg}}$ up to
the explicit coupling prefactor. By the fourth-power decay of $C^{\mathrm{bg}}$ in the kernel
estimates of the two-point witness theorem, there is a constant $C_{\mathrm{dec}}$ such that
$|C(p,q)|\le C_{\mathrm{dec}}(1+|x(p)-x(q)|)^{-4}$. Expanding the trace,
\begin{equation*}
\begin{split}
\operatorname{tr}(\mathcal{B}_cC\mathcal{B}_1C\mathcal{B}_2C)
=\sum{}&\mathcal{B}_c(q_1,q_1')\,C(q_1',q_2)\,\mathcal{B}_1(q_2,q_2')\\
&\cdot C(q_2',q_3)\,\mathcal{B}_2(q_3,q_3')\,C(q_3',q_1),
\end{split}
\end{equation*}
with $q_1,q_1'$ in $c$, $q_2,q_2'$ in $\mathcal{N}_1$ and $q_3,q_3'$ in $\mathcal{N}_2$. We bound
the factors as follows:
\begin{itemize}
\item $|C(q_1',q_2)|\le C_{\mathrm{dec}}(1+\operatorname{dist}(c,\mathcal{N}_1))^{-4}$;
\item $|C(q_3',q_1)|\le C_{\mathrm{dec}}(1+\operatorname{dist}(c,\mathcal{N}_2))^{-4}$;
\item $|C(q_2',q_3)|\le C_{\mathrm{dec}}$.
\end{itemize}
The remaining sums run over finitely many indices. This gives the bound by a constant, depending
on the matrices and on $C_{\mathrm{dec}}$, times
$(1+\operatorname{dist}(c,\mathcal{N}_1))^{-4}(1+\operatorname{dist}(c,\mathcal{N}_2))^{-4}$.

\emph{Clause (i), summability.} The sets $\mathcal{N}_1,\mathcal{N}_2$ are finite. So for $|c|$
large this product is at most a constant times $(1+|c|)^{-8}$. The number of $c\in\mathbb{Z}^4$
with $\ell'\le|c|<\ell'+1$ is of order $\ell'^3$, and $\sum_{\ell'}\ell'^3(1+\ell')^{-8}<\infty$.
This proves the summability.

\emph{Clause (i), the far cells.} The tail over cells at distance larger than $D$ is bounded
above by a multiple of $(\rho_{\mathrm{W}}/D)^4$, and since $\operatorname{tr}(\Lambda^2I)=0$ no
non-zero term of that order is exhibited; nothing more is claimed.

\emph{Clause (ii).} Apply \eqref{9.9:eq-gaussian-skeleton-2}, by assumption (a), to each pair
$(\phi_0(c),(G_1-\mathbb{E}G_1)(G_2-\mathbb{E}G_2))$ with $c$ in the box. By assumption (b), the
first factor $M_4(\phi_0(c)-a_c)$ is at most a constant times $g_{k_0}^2p_0(g_{k_0})^2$. Also
$\delta_{\mathrm{ng}}\le C_{\square}g_{k_0}\mathfrak{b}^3$. By assumption (a), the second factor,
the fourth moment of the centred product $(G_1-\mathbb{E}G_1)(G_2-\mathbb{E}G_2)$, is at most a
polylogarithm of $g_{k_0}^{-2}$ times $g_{k_0}^4\|f_1\|\,\|f_2\|$; this supplies the factor
$g_{k_0}^4\|f_1\|\,\|f_2\|$ in \eqref{9.9:eq-gaussian-skeleton-3}. Summing the resulting per-cell
bound over
the $\mathcal{V}(D)$ cells of the box gives \eqref{9.9:eq-gaussian-skeleton-3}. The polylogarithmic
factor collects $p_0(g_{k_0})^2$ and the constants.

The target share is the part of the bound of Proposition~\ref{9.9:prop-remainder-derivative}
(proof outstanding) to be carried by the box. It is
of order $g_{k_0}^6\|f_1\|\|f_2\|$, namely $g_{k_0}^2$ times a main term of order $g_{k_0}^4$. The
ratio of \eqref{9.9:eq-gaussian-skeleton-3} to the target is
$\mathcal{V}(D)\cdot\mathrm{polylog}\cdot(g_{k_0}\mathfrak{b}^3+\epsilon_{\mathfrak{b}}^{1/2})$.
This carries one power of $g_{k_0}$ beyond polylogarithmic factors, together with
$\epsilon_{\mathfrak{b}}^{1/2}$. For $D$ at most a polylogarithm of $g_{k_0}^{-2}$, as for the box
of the Gaussian extraction, the number $\mathcal{V}(D)$ and the polylogarithm are bounded by
polylogarithms of $g_{k_0}^{-2}$ uniformly in $s$, $K$, $m$; so the extra power of $g_{k_0}$, and
the smallness of $\epsilon_{\mathfrak{b}}^{1/2}$ in $g_{k_0}$ that the choice of the half-width
$\mathfrak{b}_{k_0}$ in the two-point witness theorem's own use of the comparison provides, make
the ratio small under a ceiling on the coupling, uniformly in
$s$, $K$, $m$. This is the same mechanism by
which the two-point witness theorem closes its remainder members of this kind.

\emph{Clause (iii), the complement of the box.} The per-cell bound obtained in the proof of (ii)
does not depend on the distance of $c$ from the supports. Its sum over all cells of the witness
torus is therefore the number of cells times that bound, which is proportional to the volume. The
condition $\epsilon_{\mathfrak{b}}\le\tfrac12$ is equivalent to
$\mathfrak{b}^2\ge2v\log(4n_{\mathrm{var}})$. Applied to the whole torus, $n_{\mathrm{var}}$ is the
number of variables of the torus, so the condition reads that number.

\emph{Clause (iii), the explicit $\hat{w}$-dependence.} The sourced local constituents are analytic
in $w$ only on the domains $\mathcal{W}_j(X)$, which lie in discs of radius $r_j$ of the order of
$\varrho_J^{0}$. Cauchy's estimate on a disc of radius $\varrho_J^{0}/8$ inside these domains
bounds one $\hat{w}$-derivative at $0$ by $8/\varrho_J^{0}$ times the supremum of the constituent
over that disc. The pure-channel constituents are holomorphic on the
polydiscs $\mathfrak{P}_k$. Cauchy's estimate on the disc
$|\hat{w}|<3\sigma_0x_{k_0}/16$ gives the factor $16/(3\sigma_0x_{k_0})=(16/(3\sigma_0))\,g_{k_0}^2$.

For a member small by room only, the explicit part of its derivative is its bound multiplied by
$8/\varrho_J^{0}$. This must be at most a fixed multiple of $g_{k_0}^2=x_{k_0}^{-1}$.
\begin{itemize}
\item For $R^{\mathrm{irr}}$ this requires
$\rho_{\mathrm{W}}^2\gtrsim56\,C_bp_0(g_{k_0})^4x_{k_0}/\varrho_J^{0}$, that is,
$\rho_{\mathrm{W}}\asymp x_{k_0}^{1/2}\cdot\mathrm{polylog}$.
\item The tower-tail item decreases like $L^{-4(r(s)-r_1)}$ in the depth difference $r(s)-r_1$, and
the locality item like $e^{-c_aD_{\square}}$ up to polylogarithmic factors. Here the lower and upper
windows of the coupling at depth $s$ are comparable to $x_{k_0}$, so these requirements give
$r(s)-r_1\gtrsim\frac14\log_L\bar{x}_s$ and $D_{\square}\gtrsim c_a^{-1}\log\bar{x}_s$.
\end{itemize}
This is \eqref{9.9:eq-gaussian-skeleton-4}. The only other way to remove the factor
$8/\varrho_J^{0}$ is to enlarge the radius of analyticity of the local terms in proportion to
$x_{k_0}$, which is the correlation theorem read at a large radius.

For $R^{\mathrm{BAD}}$, bounded by a multiple of $e^{-c'p_0(g_{k_0})}$ with $c'>0$ the decay rate
of the two-point witness theorem's bound for that member, and for $R^{\mathrm{loc}}$,
whose bound decays exponentially in $\rho_{\mathrm{W}}$, the factor $8/\varrho_J^{0}$ is a
constant. It is absorbed by strengthening the ceilings that make these exponentials small: by one
more factor $\log\bar{x}_s$ in one of the ceiling conditions that close these members, and by one
more power of the variable of the coupling-ceiling condition that closes these members. This holds
granted the inputs of these members in the form in which the two-point witness theorem uses them.
\end{proof}

The bound \eqref{9.9:eq-remainder-derivative-a} of
Proposition~\ref{9.9:prop-remainder-derivative}, whose proof is outstanding, does not follow by
carrying the remainder bookkeeping of the two-point witness theorem one order further.

That theorem's covariance bound on its analysis window, which enters the remainder of the Gaussian
extraction, is an absolute bound. It has no decay in the distance and is stated at a real base
point. Its main-term identification is likewise obtained in absolute, not relative, precision.

By Remark~\ref{9.9:rem-gaussian-skeleton}, route~(R2) of Remark~\ref{9.9:rem-derivative-routes},
the route inside the real-measure covariance bookkeeping, needs two inputs of a kind not present in
the two-point witness theorem.
\begin{itemize}
\item The insertion of the constant source is spread over the whole torus. A comparison with summable
  decay in the distance is therefore needed. The derivative in the offset of the law
  $\tilde{\nu}^{\hat w}$ is common to every member of the remainder and to the main term. It enters
  the proof of the two-point witness theorem first at the covariance bound just mentioned.
\item The sourced local terms carry a fixed source radius. For the members of the remainder that are
  small only through constants of the witness construction fixed before $g$ (see
  Remark~\ref{9.9:rem-derivative-routes}), a source radius growing with $x_{k_0}$ is needed.
\end{itemize}

Route~(R1) of Remark~\ref{9.9:rem-derivative-routes}, the complex-offset Cauchy route, needs one
input: the identification of the two-point witness theorem at a complex base
point. The bookkeeping of that theorem is a covariance bookkeeping under a real measure.
\begin{itemize}
\item A single derivative in a global insertion would need a decay in the distance that this
  bookkeeping does not contain.
\item Its complexification needs the sharp, logarithmic form of the identification.
\end{itemize}

The proposition below gives two conditions, each of which implies the derivative bound. It also
displays a third condition, which describes what route~(R2), the route inside the real-measure
covariance bookkeeping, would have to establish. The sufficiency of that third condition is not
proved here: the members of the remainder of the two-point witness theorem are bounded through
suprema over conditionings and histories, which are bounds on sizes, not expressions that can be
differentiated member by member.

For $a\in\mathbb{C}$ and $r>0$ we write $D(a,r)$ for the open disc of radius $r$ about $a$.

\begin{proposition}[Two sufficient conditions for the derivative bound]
\label{9.9:prop-sufficient-conditions}
We use the notation of Definition~\ref{9.9:def-main-term} along the frozen-shifted tuned family.
Fix
\begin{itemize}
\item a coupling $g$ and a torus exponent $m_0\ge m_{\Theta}(g)$;
\item a separation $\mathsf{d}^{\circ}$ with witness depth $s^{\circ}=s(\mathsf{d}^{\circ})$;
\item a reference pair $f^{\circ}\in\mathfrak{T}(\mathsf{d}^{\circ};\vartheta)$;
\item a length $n$.
\end{itemize}
Let $k_0$ be the witness level and $x_{k_0}>0$ the inverse coupling of the tuned run there. Put
$g_{k_0}=x_{k_0}^{-1/2}$ and $N_{k_0}=N_{k_0}(0)$. Let $\bar{x}^{\Theta}_{s^{\circ}}$ be the upper
window of \eqref{9.9:eq-tuned-family-a} for the tuned run's inverse coupling at depth $s^{\circ}$.
Put $X=g^{-2}$ and $X^{\flat}=X-2c_2$. Consider the following conditions.

\smallskip
\noindent\textup{(S1)} (Complex offsets, with a precision-to-radius relation.) There are numbers
$R>0$ and $r_{\mathbb{C}}\ge 0$, not depending on $n$, on the position or on the pair, such that:
\begin{itemize}
\item[(i)] the running-shift statement of the correlation theorem holds on $D(0,8R)$, so that in
particular $\zeta_{k_0}$ is holomorphic on $D(0,8R)$ with the first bound below; and
$R\le x_{k_0}/16$ (for instance $R\le X^{\flat}/16$);
\item[(ii)] $\hat w\mapsto\varphi_n(\hat w)$ is holomorphic on a neighbourhood of
$\{|\hat w|\le R\}$ and satisfies the second bound below on $|\hat w|=R$;
\item[(iii)] the third bound below holds.
\end{itemize}
\begin{equation}\label{9.9:eq-sufficient-conditions-a}
\begin{gathered}
|\zeta_{k_0}(\hat w)-\hat w|\le\tfrac13|\hat w|\quad(|\hat w|<8R),\\
|\varphi_n(\hat w)-\mathfrak{M}_n(\hat w)|\le r_{\mathbb{C}}\,|\mathfrak{M}_n(\hat w)|
\quad(|\hat w|=R),\qquad
r_{\mathbb{C}}\,\bar{x}^{\Theta}_{s^{\circ}}\le\tfrac18\,R .
\end{gathered}
\end{equation}

\smallskip
\noindent\textup{(S2)} (A decaying Gaussian comparison at the witness lattice.) There are a weight
$\omega:\mathbb{N}\to[0,\infty[$ with
\begin{equation*}
\sum_{c\in\mathbb{Z}^4}\omega(|c|)<\infty
\end{equation*}
and a constant $C_{\mathrm{gl}}$, neither depending on $K$, $m$, $g_0$, $s$ or the position, with
the following property. For every cell $c$ of the witness lattice, and every local functional $Y$
among those through which the two-point witness theorem bounds the members of its remainder,
\begin{equation}\label{9.9:eq-sufficient-conditions-b}
\begin{split}
&\bigl|\operatorname{Cov}_{\tilde{\nu}}(\phi_0(c),Y)
-\operatorname{Cov}_{\gamma}(\phi_0^{\gamma}(c),Y^{\gamma})\bigr|\\
&\qquad\le C_{\mathrm{gl}}\,(\vartheta_P+t^{1/2})\,g_{k_0}^2\,p_0(g_{k_0})^2\,M(Y)\,
\omega\bigl(\operatorname{dist}(c,N_1\cup N_2)\bigr).
\end{split}
\end{equation}
Here:
\begin{itemize}
\item $\tilde{\nu}$ is the law on pairs (history, block field) at level $k_0$;
\item $\gamma$, the parameters $\vartheta_P$ and $t$, and the cell summands $\phi_0(c)$ of the
response $\Phi_0$ to a constant source are as in Remark~\ref{9.9:rem-gaussian-skeleton};
\item the superscript $\gamma$ denotes the Gaussian counterpart;
\item $p_0(\cdot)$ is the degree function;
\item $M(Y)$ is the bound used for $Y$ in the two-point witness theorem;
\item $N_1$, $N_2$ are the two regions of the witness construction attached to the reference pair.
\end{itemize}
So the comparison is required to decay faster than the fourth power of the distance, and nothing
more in $g$ is asked.

\smallskip
\noindent\textup{(S3)} (The sharp form.) There are $\sigma>2$ and a constant $C$, not depending on
$n$, $s^{\circ}$, the position or the pair, such that on a neighbourhood of
$\{|\hat w|\le\varrho_J^{0}/4\}$ one has $\varphi_n=\mathfrak{M}_n(1+h_n)$, with $h_n$ holomorphic
there and
\begin{equation}\label{9.9:eq-sufficient-conditions-c}
|h_n(\hat w)|\le C\,g_{k_0}^{\sigma}\qquad(|\hat w|\le\varrho_J^{0}/4).
\end{equation}

\smallskip
\noindent Then the following hold.
\begin{itemize}
\item[(a)] If \textup{(S1)} holds, then
\begin{equation}\label{9.9:eq-sufficient-conditions-1}
|E_n'(0)|\le\tfrac34\,g_{k_0}^2\,\mathfrak{M}_n(0),
\end{equation}
that is, the bound \eqref{9.9:eq-remainder-derivative-a} of
Proposition~\ref{9.9:prop-remainder-derivative}, whose proof is outstanding, with $c_E=\frac34$.
\item[(b)] Let $g$ satisfy $X^{\flat}>0$ and $\frac{16}{3}\varrho_J^{0}\le\min\{\sigma_0,1\}X^{\flat}$,
two of the conditions on $g$ that define the ceiling $\gamma_{\mathrm{cal}}$ of
Theorem~\ref{9.9:thm-nondegenerate-calibration}. Assume the running-shift statement in the form of
Definition~\ref{9.9:def-tuned-hypotheses}, with source disc $D(0,2\varrho_J^{0})$, and let
\textup{(S3)} hold. Let $c'$ be the additive
constant of the upper window, so that
$\bar{x}^{\Theta}_{s^{\circ}}=X+3\lambda_U s^{\circ}+c'$, and let $\lambda$ be the lower rate of the
pairwise bounds of Definition~\ref{9.9:def-tuned-hypotheses}. If
\begin{equation}\label{9.9:eq-sufficient-conditions-2}
C\,(X^{\flat})^{-(\sigma-2)/2}\,
\max\Bigl\{\frac{3\lambda_U}{\lambda},\ \frac{X+c'}{X-2c_2}\Bigr\}\,
\frac{32}{\varrho_J^{0}}\le 1,
\end{equation}
then \textup{(S1)} holds with $R=\varrho_J^{0}/4$ and $r_{\mathbb{C}}=C\,g_{k_0}^{\sigma}$. Hence
\eqref{9.9:eq-sufficient-conditions-1} holds. Condition \eqref{9.9:eq-sufficient-conditions-2}
reads no depth, cutoff or torus. It holds for all sufficiently small $g$ provided
$c'(g)\,g^2\to0$ as $g\to0$.
\item[(c)] Condition \textup{(S2)}, together with analyticity in the offset of the sourced local
constituents of the two-point witness theorem on discs of radius proportional to $x_{k_0}$, is what
route~(R2) of Remark~\ref{9.9:rem-derivative-routes}, the route inside the real-measure covariance
bookkeeping, would need, and would close that route; the sufficiency is not proved here. At the
fixed radius $\varrho_J^{0}$, the members of the remainder that are small only through constants of
the witness construction fixed before $g$ are not controlled; see
Remark~\ref{9.9:rem-gaussian-skeleton}. This proposition asserts no implication
from \textup{(S2)} to \eqref{9.9:eq-remainder-derivative-a}.
\end{itemize}
\end{proposition}

\begin{proof}
\emph{(a), the shift and its derivative.} By the first bound of
\eqref{9.9:eq-sufficient-conditions-a} at $\hat w=0$ we have $\zeta_{k_0}(0)=0$. Define
\begin{equation*}
u(\hat w):=\frac{\zeta_{k_0}(\hat w)-\hat w}{\hat w}.
\end{equation*}
Then $u$ is holomorphic on $D(0,8R)\setminus\{0\}$ and satisfies $|u|\le\frac13$ there. By
Riemann's theorem on removable singularities, $u$ extends holomorphically to $D(0,8R)$, and by
continuity still $|u|\le\frac13$. Since $\zeta_{k_0}(\hat w)=\hat w\,(1+u(\hat w))$, and
$N_{k_0}(\hat w)=\zeta_{k_0}'(\hat w)$ by \eqref{9.9:eq-main-term-a},
\begin{equation*}
N_{k_0}(\hat w)=1+u(\hat w)+\hat w\,u'(\hat w)\qquad(\hat w\in D(0,8R)).
\end{equation*}

Let $|\hat w|=R$. Every point of $D(\hat w,7R)$ has modulus less than $8R$, so
$D(\hat w,7R)\subset D(0,8R)$. Cauchy's estimate for $u$ on the discs $D(\hat w,\rho)$,
$\rho<7R$, followed by the limit $\rho\to7R$, gives $|u'(\hat w)|\le\frac13\cdot\frac1{7R}$. Hence
$|\hat w\,u'(\hat w)|\le\frac1{21}$, and
\begin{equation*}
|N_{k_0}(\hat w)|\le1+\tfrac13+\tfrac1{21}=\tfrac{29}{21}\qquad(|\hat w|=R).
\end{equation*}
At $\hat w=0$, $N_{k_0}=1+u(0)$, so $|N_{k_0}|\ge\frac23$. Since
$\frac{29}{21}=\frac{29}{14}\cdot\frac23$, this gives
\begin{equation*}
|N_{k_0}(\hat w)|\le\tfrac{29}{14}\,|N_{k_0}|\qquad(|\hat w|=R).
\end{equation*}

\emph{(a), the shifted inverse coupling.} For $|\hat w|=R$, the first bound of
\eqref{9.9:eq-sufficient-conditions-a} gives
$|\zeta_{k_0}(\hat w)|\le|\hat w|+\frac13|\hat w|=\frac43R$. Using $R\le x_{k_0}/16$,
\begin{equation*}
|x_{k_0}(\hat w)|=|x_{k_0}-\zeta_{k_0}(\hat w)|\ge x_{k_0}-\tfrac43R
\ge x_{k_0}-\tfrac1{12}x_{k_0}=\tfrac{11}{12}\,x_{k_0}.
\end{equation*}
The same bound on all of $D(0,8R)$ gives
$|\zeta_{k_0}(\hat w)|<\frac43\cdot8R\le\frac23x_{k_0}$. So $x_{k_0}(\hat w)\neq0$ on $D(0,8R)$.
The formula \eqref{9.9:eq-main-term-a} therefore defines $\mathfrak{M}_n$ holomorphically on
$D(0,8R)$, and $x_{k_0}(0)=x_{k_0}$.

\emph{(a), the main term on the circle.} Recall that $\mathfrak{M}_n(0)$ is the main term of the
two-point witness theorem, which is positive, so that $|\mathfrak{M}_n(0)|=\mathfrak{M}_n(0)$. By
\eqref{9.9:eq-main-term-a},
$|\mathfrak{M}_n(\hat w)|=|N_{k_0}(\hat w)|^2\,b_0\,|x_{k_0}(\hat w)|^{-2}\,|\mathcal{K}|$. The two
previous steps give, for $|\hat w|=R$,
\begin{equation*}
|\mathfrak{M}_n(\hat w)|
\le\Bigl(\frac{29}{14}\Bigr)^2\Bigl(\frac{12}{11}\Bigr)^2|N_{k_0}|^2\,b_0\,x_{k_0}^{-2}\,
|\mathcal{K}|
=\frac{121104}{23716}\,\mathfrak{M}_n(0)<6\,\mathfrak{M}_n(0),
\end{equation*}
because $(29\cdot12)^2=121104$, $(14\cdot11)^2=23716$ and $6\cdot23716=142296$.

\emph{(a), Cauchy at the origin.} By (ii) and the previous step, $E_n=\varphi_n-\mathfrak{M}_n$ is
holomorphic on a neighbourhood of $\{|\hat w|\le R\}$. Cauchy's estimate at $0$ on the circle
$|\hat w|=R$, the second bound of \eqref{9.9:eq-sufficient-conditions-a} and the previous step give
\begin{equation*}
|E_n'(0)|\le\frac1R\max_{|\hat w|=R}|E_n(\hat w)|
\le\frac{r_{\mathbb{C}}}{R}\max_{|\hat w|=R}|\mathfrak{M}_n(\hat w)|
\le\frac{6\,r_{\mathbb{C}}}{R}\,\mathfrak{M}_n(0).
\end{equation*}
By the third bound of \eqref{9.9:eq-sufficient-conditions-a},
$r_{\mathbb{C}}/R\le1/(8\bar{x}^{\Theta}_{s^{\circ}})$. Since $\bar{x}^{\Theta}_{s^{\circ}}$ is the
upper window \eqref{9.9:eq-tuned-family-a} for $x_{k_0}$, we have
$x_{k_0}\le\bar{x}^{\Theta}_{s^{\circ}}$, and $g_{k_0}^2=x_{k_0}^{-1}$. Therefore
\begin{equation*}
|E_n'(0)|\le\frac68\,\frac{\mathfrak{M}_n(0)}{\bar{x}^{\Theta}_{s^{\circ}}}
\le\frac34\,\frac{\mathfrak{M}_n(0)}{x_{k_0}}=\frac34\,g_{k_0}^2\,\mathfrak{M}_n(0),
\end{equation*}
which is \eqref{9.9:eq-sufficient-conditions-1}. None of the numbers used depends on $n$, the
position or the pair.

\emph{(b), the radius.} Put $R=\varrho_J^{0}/4$ and $r_{\mathbb{C}}=C\,g_{k_0}^{\sigma}$. By
hypothesis the running-shift statement holds on $D(0,2\varrho_J^{0})=D(0,8R)$. This gives the
holomorphy of $\zeta_{k_0}$ there and the first bound of \eqref{9.9:eq-sufficient-conditions-a}.

By hypothesis $X^{\flat}>0$ and
$\frac{16}3\varrho_J^{0}\le\min\{\sigma_0,1\}X^{\flat}\le X^{\flat}$.
Hence $\varrho_J^{0}\le\frac3{16}X^{\flat}$ and
\begin{equation*}
R=\tfrac14\varrho_J^{0}\le\tfrac3{64}X^{\flat}\le\tfrac1{16}X^{\flat}.
\end{equation*}
The pairwise bounds of Definition~\ref{9.9:def-tuned-hypotheses}, read against the normalisation
$x_{(0)}=X$ of the tuned run, give
\begin{equation}\label{9.9:eq-sufficient-conditions-3}
x_{k_0}\ge X-2c_2+\lambda s^{\circ}=X^{\flat}+\lambda s^{\circ}\ge X^{\flat}.
\end{equation}
So $R\le x_{k_0}/16$, and clause (i) of \textup{(S1)} holds.

\emph{(b), the precision.} By \textup{(S3)}, $\varphi_n$ is holomorphic on a neighbourhood of
$\{|\hat w|\le R\}$. On $|\hat w|=R$, \eqref{9.9:eq-sufficient-conditions-c} gives
\begin{equation*}
|\varphi_n-\mathfrak{M}_n|=|\mathfrak{M}_n|\,|h_n|\le C\,g_{k_0}^{\sigma}\,|\mathfrak{M}_n|.
\end{equation*}
This is clause (ii) of \textup{(S1)} with $r_{\mathbb{C}}=C\,g_{k_0}^{\sigma}$.

Clause (iii) reads $C\,g_{k_0}^{\sigma}\,\bar{x}^{\Theta}_{s^{\circ}}\le\varrho_J^{0}/32$, that is,
\begin{equation*}
C\,g_{k_0}^{\sigma-2}\,\bigl(\bar{x}^{\Theta}_{s^{\circ}}\,g_{k_0}^2\bigr)\,
\frac{32}{\varrho_J^{0}}\le1 .
\end{equation*}
We bound the bracket. By \eqref{9.9:eq-sufficient-conditions-3} and
$\bar{x}^{\Theta}_{s^{\circ}}=X+3\lambda_U s^{\circ}+c'$,
\begin{equation*}
\bar{x}^{\Theta}_{s^{\circ}}\,g_{k_0}^2=\frac{\bar{x}^{\Theta}_{s^{\circ}}}{x_{k_0}}
\le\frac{(X+c')+3\lambda_U s^{\circ}}{(X-2c_2)+\lambda s^{\circ}}
\le\max\Bigl\{\frac{3\lambda_U}{\lambda},\ \frac{X+c'}{X-2c_2}\Bigr\}.
\end{equation*}
The last step is the inequality $(a+b)/(c+d)\le\max\{a/c,b/d\}$ for $a,b\ge0$ and $c,d>0$. When
$s^{\circ}=0$ it is the bound by the second entry. This bound does not depend on $s^{\circ}$.

Moreover $x_{k_0}\ge X^{\flat}>0$ and $\sigma>2$, so
$g_{k_0}^{\sigma-2}=x_{k_0}^{-(\sigma-2)/2}\le(X^{\flat})^{-(\sigma-2)/2}$. Hence the left-hand side
of the last condition is at most the left-hand side of \eqref{9.9:eq-sufficient-conditions-2}, and
clause (iii) follows. Thus \textup{(S1)} holds with $R$ and $r_{\mathbb{C}}$ not depending on $n$,
the position or the pair, and part~(a) gives \eqref{9.9:eq-sufficient-conditions-1}.

\emph{(b), small coupling.} As $g\to0$, $X^{\flat}=g^{-2}-2c_2\to\infty$, so
$(X^{\flat})^{-(\sigma-2)/2}\to0$. Also $(X+c')/(X-2c_2)\to1$ when $c'(g)\,g^2\to0$. Under that
proviso, \eqref{9.9:eq-sufficient-conditions-2} holds for all sufficiently small $g$.

Part~(c) asserts no implication and requires no proof.
\end{proof}

Condition (iii) of \textup{(S1)} ties the precision to the radius. At a fixed radius
$R\le\varrho_J^{0}/4$, clause (i) is the running-shift statement of the correlation theorem itself.
Clause (iii) then asks for the relative precision
$r_{\mathbb{C}}\le\varrho_J^{0}/(32\,\bar{x}^{\Theta}_{s^{\circ}})$, of the order of $g_{k_0}^2$.
This is the sharp form \textup{(S3)}.

Now take $R$ of the size of the large source radius of the correlation theorem. That radius is
$R\asymp c_{\mathrm{rad}}X/\Pi(X)$, where $c_{\mathrm{rad}}$ is the radius constant and $\Pi(X)$ the
polylogarithm of $X$ that limits it, as in route~(R1) of Remark~\ref{9.9:rem-derivative-routes},
the complex-offset Cauchy route. Clause (iii) then asks for
\begin{equation*}
r_{\mathbb{C}}\lesssim\frac{c_{\mathrm{rad}}X}{8\,\Pi(X)\,\bar{x}^{\Theta}_{s^{\circ}}},
\end{equation*}
that is, $c_{\mathrm{rad}}/\Pi(X)$ times $X/\bar{x}^{\Theta}_{s^{\circ}}$. Since
$\bar{x}^{\Theta}_{s^{\circ}}=X+3\lambda_U s^{\circ}+c'$ grows with $s^{\circ}$, this precision must
still improve with the depth. A complexification of the two-point witness theorem at that radius,
with relative precision merely small against $1/\Pi(X)$, would give
\eqref{9.9:eq-remainder-derivative-a} only on a bounded range of depths.

A constant $c_E$ independent of the depth can be obtained from \textup{(S1)} only in one of two
ways:
\begin{itemize}
\item $R$ grows like $x_{k_0}$, which goes beyond the running-shift statement of the correlation
theorem;
\item the precision decays like $x_{k_0}^{-1}R$.
\end{itemize}
Of the conditions displayed above, only \textup{(S3)} is shown here to yield
\eqref{9.9:eq-remainder-derivative-a} with a constant independent of the depth.

A proof of \textup{(S3)} would be a cluster (polymer) expansion at the witness lattice for
connected correlations of the renormalised curvature insertions with the action density. This
expansion must be uniform in the cutoff and run through small and large fields at level $k_0$. It
carries the weights of the large-field block bound, healed part, whose proof is outstanding. It is
not harder in kind than the two-point witness theorem together with the Gaussian extraction in
relative precision.

\begin{remark}[Why perturbation theory sees no obstruction]\label{9.9:rem-perturbative-sign}
This remark explains why perturbation theory gives the sign of $\varphi_n'(0)$ at once, and why that is not a proof. Let $N_{k_0}$, $\zeta''_{k_0}$ and $\mathfrak{M}_n$ be as in the lower bound \eqref{9.9:eq-calibration-lower-bound-a}.

To every finite order in the running coupling, the relative remainder of the two-point witness theorem is $c_1g_{k_0}^2+\cdots$, where $c_1$ denotes its leading coefficient, and the remaining quantities expand as
\begin{equation*}
N_{k_0}-1=O(g_{k_0}^2),\qquad \zeta''_{k_0}(0)=O(g_{k_0}^4).
\end{equation*}
The last order is the two-loop one. The increments satisfy
\begin{equation*}
b_j'\asymp\beta_1g_j^4,\qquad b_j''\asymp\beta_1g_j^6,
\end{equation*}
where $\beta_1$ is the two-loop coefficient. Summing these over the levels $j$ below $k_0$ gives the stated orders of $N_{k_0}-1$ and of $\zeta''_{k_0}(0)$.

Insert these orders into the logarithmic derivative of $\mathfrak{M}_n$ used in the proof of Theorem~\ref{9.9:thm-nondegenerate-calibration}. The leading term is $2N_{k_0}g_{k_0}^2$. The first correction is $2\zeta''_{k_0}(0)/N_{k_0}$, which is $O(g_{k_0}^4)$. The derivative of the remainder is of higher order in $g$ as well. Hence
\begin{equation*}
\varphi_n'(0)=2N_{k_0}g_{k_0}^2\,\mathfrak{M}_n(0)\,\bigl(1+O(g^2)\bigr)>0 .
\end{equation*}
Positivity holds with room to spare: to every finite order the corrections to the leading term $2N_{k_0}g_{k_0}^2$ carry two further powers of the coupling. The sign is the physical one: lowering the bare inverse coupling strengthens the fluctuations, and with them the correlation.

The two-point witness theorem is a two-sided bound. Its remainder is small by the room provided by the auxiliary parameter $C_0$ and by the coupling $g$. It is not, however, an asymptotic expansion with a differentiable error term.

Proposition~\ref{9.9:prop-remainder-derivative}, whose proof is outstanding, adds exactly this and nothing else. Its bound \eqref{9.9:eq-remainder-derivative-a} controls the derivative at $0$ of that remainder with respect to the bare offset $\hat{w}$, along the tuned family.

At every finite order of perturbation theory, therefore, the calibration is non-degenerate. A proof of non-degeneracy along the tuned family, in the form of Theorem~\ref{9.9:thm-nondegenerate-calibration}, requires in addition the bound \eqref{9.9:eq-remainder-derivative-a}.
\end{remark}

\begin{remark}[What the calibration statements do not assert]\label{9.9:rem-calibration-not-asserted}
The statements of this section on the calibration of the bare coupling by the connected curvature
two-point function do not assert any of the following.
\begin{enumerate}
\item[(a)] The bound of Proposition~\ref{9.9:prop-remainder-derivative}, whose proof is outstanding,
is the bound on the bare-offset derivative of the remainder of the two-point witness theorem. It is
not asserted that this bound follows from the inputs of the two-point witness theorem. The point at
which the estimates of that theorem stop short of the bound is described in the discussion
accompanying the differentiated witness, display~\eqref{9.9:eq-differentiated-witness-a}. This book
gives no proof of the bound from those inputs.
\item[(b)] Non-degeneracy is not asserted for Wilson loops or for plaquette densities. These
quantities lie outside the class of observables of clause~(iv) of Theorem~\ref{3.1:thm-main}.
\item[(c)] Injectivity of the calibration map is not asserted on
the whole window $[-\varrho_2/2,\varrho_2/2]$ of real offsets.
Remark~\ref{9.9:rem-global-uniqueness} states what would give it.
\item[(d)] Nothing is asserted for $N\ge 3$.
\item[(e)] Nothing is asserted about the block algebra.
\item[(f)] No sign is asserted for the coupling increment of a single renormalisation step
(Definition~\ref{1.2:def-couplings}).
\end{enumerate}
\end{remark}

\begin{definition}[Levels, depths and the filing level of a remnant]\label{9.10:not-filing-level}
Levels are counted as in the correlation theorem. The level $k$ runs from $0$ (the bare lattice) to
$K'$ (the last level of the run, at which it stops), and $\mathbb{T}^{(k)}$ is the unit lattice after
$k$ steps (Definition~\ref{1.5:def-block-average}). Depths are counted as in the uniqueness theorem:
the depth of level $k$ is $r=n-k$, where $n$ is the integer of the uniqueness theorem from which
depths are counted. Both are read on the prescribed scaffold $\Theta_X$ at $X=g^{-2}$,
with $\theta_r:=X+c_0r$. Every threshold, radius, integer, partition and tree used at depth $r$ is
read from $\theta_r$. Each of these is therefore common to both runs of a pair compared in
Definition~\ref{9.6:def-comparison-chain}.

Let a remnant, in the sense of the re-filing lemmas, be born at level $i$, that is, at depth
$s=n-i$. Put
\[
  \mathsf{a}_i:=\log_L\bigl(M\check R_i\bigr)\in\mathbb{N}.
\]
Here $\mathsf{a}_i\ge2$, because $M=L^{m'}$ with $m'\ge2$, the numbers $M>M_2>M_1$ being powers
of $L$. Subscripts written $i$, $j$ or $k$, and subscripts that are expressions in $n$ such as
$n-s$, are levels; subscripts written $r$, $s$ or $u$ are depths; thus $\mathsf{a}_s$ and $\ell_s$,
with $s$ a depth, mean $\mathsf{a}_{n-s}$ and $\mathsf{a}_{n-s}+2$, and a parenthesised subscript is
always a depth: $\check R_{(r)}:=\check R_{n-r}$.
Let $k_1(k)$ denote the lower band index at level $k$ of the correlation theorem.
The \emph{filing level}, the \emph{filing depth}, the set $\mathfrak{B}(k)$ of levels
\emph{filed into the output of step $k$}, the \emph{band}, the \emph{filed set} and the
\emph{remnant depth-range} are
\begin{equation}\label{9.10:eq-filing-level-a}
\begin{aligned}
  i''&:=i+\mathsf{a}_i, & s''&:=s-\mathsf{a}_s,\\
  \mathfrak{B}(k)&:=\{\,i:\ i''=k+1\,\}, &
  \mathrm{band}(k)&:=\{\,i:\ k_1(k)<i\le k\,\}=\{\,i\le k:\ i''>k+1\,\},\\
  \mathrm{filed}(k)&:=\{\,i:\ i''\le k\,\}, & \ell_s&:=\mathsf{a}_s+2 .
\end{aligned}
\end{equation}
The level $i''$ is the output level at which the $M\check R_i$-cubes of $\mathbb{T}^{(i)}$ have
become the unit cells $\square'_y$, $y\in\mathbb{T}^{(i'')}$. This is the cell alignment of the
re-filing lemmas, which uses that $L$ is odd, so that the $M\check R_i$-cubes are centred at points
of $\mathbb{T}^{(i'')}$, exactly as the $M$-cubes of \cite{B-CMP109} at level $j$ are centred at the
points of the block lattice $\mathbb{T}^{(j+\log_L M)}$, which is $\log_L M$ steps coarser.
The re-filing lemmas also give $\#\mathfrak{B}(k)\le2$: for a given $k$ there are at most two
levels $i$ with $i''=k+1$. This bound rests on the dichotomy stated in \cite{B-CMP122a}, used here
in the following form: where such a level exists, there is exactly one or there are two. They also
give
\begin{equation}\label{9.10:eq-filing-level-1}
  \{\,i\le k_1(k)\,\}=\{\,i:\ i''\le k\,\}\sqcup\mathfrak{B}(k).
\end{equation}
By the correlation theorem,
\[
  \#\mathrm{band}(k)=k-k_1(k)\le\mathfrak{k}(\mathsf{T}_k),
\]
where $\mathfrak{k}(\mathsf{T}_k)$ is the correlation theorem's bound on the band length in terms of
its band parameter $\mathsf{T}_k$, and by the band-length corollary of the large-field block bounds
$\mathfrak{k}(\mathsf{T}_k)\le2+\mathsf{a}_k$. For a depth $r$, the letter $\ell_r=\mathsf{a}_r+2$
of \eqref{9.10:eq-filing-level-a} is the band letter $2+\log_L(M\check R_{(r)})$ of the uniqueness
theorem.

These sets partition the levels up to $k$:
\begin{equation}\label{9.10:eq-filing-level-2}
  \{\,i\le k\,\}=\mathrm{filed}(k)\sqcup\mathfrak{B}(k)\sqcup\mathrm{band}(k).
\end{equation}
In depths, consider the output at depth $r$, produced by the step at level $k=n-r-1$. That step
reads the remnant of depth $s$ as a remnant, that is, its level $n-s$ lies in $\mathfrak{B}(k)$ or in
$\mathrm{band}(k)$, if and only if
\begin{equation}\label{9.10:eq-filing-level-3}
  r<s\le r+\mathsf{a}_s .
\end{equation}
The remnant lies in $\mathfrak{B}(k)$ exactly when $s-\mathsf{a}_s=r$, and it is then filed into
the output at depth $r$. It lies in $\mathrm{band}(k)$ exactly when
\[
  r<s<r+\mathsf{a}_s ,
\]
and its cube, measured in $\mathbb{T}^{(k)}$, is then still larger than $L$. By the share clause
of the re-filing lemmas, as an unfiled share the remnant is read from zones of depth $u$ with
\[
  s-\mathsf{a}_s\le u\le s ,
\]
and from the output depth $s''=s-\mathsf{a}_s$ on it lies inside the hatted regular family
$\hat E$. Finally, on $\Theta_X$ the exponent $\mathsf{a}_s$ and the range $\ell_s$ are functions of
the depth alone, and $\mathrm{filed}$, $\mathfrak{B}$ and $\mathrm{band}$ are the same index sets in
both runs of a pair at every common depth.
\end{definition}

\begin{proof}
\emph{Partition.} Fix $k$ and let $i\le k$. Exactly one of $i''\le k$, $i''=k+1$ and $i''>k+1$
holds. Conversely, every $i$ with $i''\le k+1$ satisfies $i\le k$, since $i\le i''-2\le k-1$.
Hence
\[
  \{\,i\le k\,\}=\mathrm{filed}(k)\sqcup\mathfrak{B}(k)\sqcup\mathrm{band}(k).
\]
Both descriptions of $\mathrm{band}(k)$ in \eqref{9.10:eq-filing-level-a} define the complement of
$\mathrm{filed}(k)\sqcup\mathfrak{B}(k)$ in $\{i\le k\}$. For the second description this was just
shown. For the first it is \eqref{9.10:eq-filing-level-1}. This proves the second equality for
$\mathrm{band}(k)$.

\emph{Reading in depths.} The output at depth $r$ is produced by the step at level $k=n-r-1$. Its
output level is therefore $k+1=n-r$. The remnant of depth $s$ has level $i=n-s$ and filing level
$i''=n-s+\mathsf{a}_s$.

Membership in $\mathfrak{B}(k)$ reads $n-s+\mathsf{a}_s=n-r$, that is, $s-\mathsf{a}_s=r$. In this
case the remnant is filed into the output at depth $r$.

Membership in $\mathrm{band}(k)$ reads $n-s\le n-r-1$ and $n-s+\mathsf{a}_s>n-r$, that is,
$r<s<r+\mathsf{a}_s$. Measured in $\mathbb{T}^{(k)}$, its cube has side $L^{i''-k}\ge L^2>L$.

So the step producing the output at depth $r$ reads the remnant of depth $s$ as a remnant, that is,
its level $n-s$ lies in $\mathfrak{B}(k)$ or in $\mathrm{band}(k)$, if and only if
\eqref{9.10:eq-filing-level-3} holds. The reading as an unfiled share, and the position inside the
hatted regular family $\hat E$ from the output depth $s''$ on, are the share clause of the
re-filing lemmas.

\emph{Both runs.} On $\Theta_X$ the radius $\check R$ is read from $\theta$. Hence $\mathsf{a}_s$,
and with it $\ell_s$, is a function of the depth alone. Therefore $\mathrm{filed}$,
$\mathfrak{B}$ and $\mathrm{band}$ are the same index sets in both runs of a pair at every common
depth.

Under the standing condition $n\ge n_2(g)$ of the uniqueness theorem these common depths exist:
beneath every depth $r\le n/2$ both runs have more than $\ell_r$ levels and more than
$\mathsf{N}(\theta_r)$ levels, where $\mathsf{N}(\theta_r)$ is the number of levels that the
scaffold prescribes beneath depth $r$.
\end{proof}

\begin{definition}[The re-filed renormalisation run]\label{9.10:def-refiled-run}
Levels $k$, depths $r$, the prescribed scaffold $\Theta_X$ with $\theta_r=X+c_0r$, the filing
level $i''=i+\mathsf{a}_i$ of a remnant born at level $i$, and the index sets $\mathfrak{B}(k)$
and $\mathrm{band}(k)$ are those of Notation~\ref{9.10:not-filing-level} and
\eqref{9.10:eq-filing-level-a}. The source is either local, $w\in\mathcal{W}_k$, or frozen,
$\hat w\in D(0,2r)$, where $D(0,2r)\subset\mathbb{C}$ is the disc centred at $0$ of radius $2r$
in which the frozen source ranges, with $r$ the radius parameter of the frozen source in the
uniqueness theorem and not a depth; a subscript $\zeta$ on a family denotes the family at the
frozen source $\zeta$. The \emph{re-filed run} is the collection of objects (a)--(g) below. Objects
carrying a hat are its own, as are the remnants $R^{(i)}$ of item~(a) and the output families
$E^{(k+1)}$ of item~(c). An object said below to be the same as in the correlation theorem is
defined by the same rule as there, the rule being read in the re-filed run, at its own couplings,
shifts and action (items (d), (f), (g)); it coincides with the object of the correlation theorem
itself where the rule reads none of these.

\begin{itemize}
\item[(a)] \emph{Remnants.} For each level $i$, the remnant $R^{(i)}=R'^{(i)}+R''^{(i)}$, its
production, its class and its vacuum value are the same as in the correlation theorem: the class
is the marginal class of families of \cite[(2.30)--(2.32)]{B-CMP119}, with
$\|R^{(i)}_\zeta\|^{\flat}\le g_i^{\kappa_0}$, where $\kappa_0$ is the exponent of the smallness
bound of the remnants in the correlation theorem, and the vacuum value $R^{(i)}(X,1;w)$ lies in the
set $e^{(i)}(X;w)$ of vacuum pieces
of the correlation theorem. A remnant is created exactly as in the correlation theorem; only its
later filing changes.

\item[(b)] \emph{Re-filing of a remnant.} For a remnant $R=R^{(i)}$, $\mathsf{RF}[R]$ is the
pointed family at scale $i''$ constructed in the re-filing lemmas. It is built from the shares
$\mathsf{sh}(X,y)$, $y\in\mathbb{T}^{(i'')}$, with assignment domains
$Y(X,y)\supseteq\square_0(X)$, and its argument is the configuration
$U_i(\square_0(X),M^{\cdot}_i(\mathfrak{b}))$ formed from the block field $U_i$ on the cube
$\square_0(X)$ and the datum $M^{\cdot}_i(\mathfrak{b})$ fixed in the construction of the
level-wise rate. By the re-filing lemmas:
$\mathsf{RF}[R]$ is related to the vacuum term $R^{\mathrm{vac}}$ of $R$ (the term entering
item~(g)) by the relation $\simeq$ of the re-filing lemmas, and $\mathsf{RF}[R](Y,1,y)=0$;
$\mathsf{RF}[R]$ belongs to the pointed class on the radii $\hat\alpha_{i''-1}$ fixed in the
re-filing lemmas, that class requiring covariance of the pointed total under every symmetry of
$\mathbb{T}^{(i'')}$; its norm obeys
\begin{equation}\label{9.10:eq-refiled-run-1}
\|\mathsf{RF}[R]\|^{\flat}\le C_{RF}\,w_{i''-1,i}\,\|R\|^{\flat},
\end{equation}
where $w_{i''-1,i}$ is the zone weight of the correlation theorem (this letter does not denote
the source) and $C_{RF}$ is the norm constant of the re-filing lemmas;
the map $R\mapsto\mathsf{RF}[R]$ is linear; and the dependence on parameters and analyticity in
$w$ pass from $R$ to $\mathsf{RF}[R]$, the value at $(Y,\cdot,y)$ reading $w$ only through its
restriction $w\restriction_{Y^{+}}$ to the enlargement $Y^{+}$ of $Y$ fixed in the re-filing
lemmas.

\item[(c)] \emph{The output families.} For every frozen $\zeta$, and locally in $w$,
\begin{equation}\label{9.10:eq-refiled-run-a}
\begin{aligned}
\hat E^{(k+1)}&:=E^{(k+1)}+\sum_{i\in\mathfrak{B}(k)}\mathsf{RF}[R^{(i)}],\\
\hat E^{(k+1)}(X,\cdot,y;w)&:=E^{(k+1)}(X,\cdot,y;w)
+\sum_{i\in\mathfrak{B}(k)}\mathsf{RF}[R^{(i)}(\cdot;w)](X,\cdot,y),\\
\hat Q^{(k+1)}&:=\hat E^{(k+1)}-P^{\flat(k+1)},\\
\hat D^{(j)}&:=\hat E^{(j)}(\cdot;w)-\hat E^{(j)}_{w(p_y)}.
\end{aligned}
\end{equation}
Here $E^{(k+1)}$ is the output family of step $k$ of the re-filed run, produced by the
production formula of the correlation theorem from the inputs of item~(g). The pure part
$P^{\flat(k+1)}$ is defined as in the correlation theorem, from the pure
inputs $\sum(P^{\flat(i)})^{\mathrm{ren}}$ and the numbers $\hat\zeta_i$, and
$p_y$ is the plaquette attached to $y$ in the correlation theorem.
The freezing defect $\hat D^{(j)}$ vanishes at constant $w$, by the statement of the re-filing
lemmas that the local re-filed family equals the frozen one where $w$ is constant on $Y^{+}$.

\item[(d)] \emph{Couplings and running shift.} With $\beta[\cdot]$ the coupling-extraction
functional,
\begin{equation}\label{9.10:eq-refiled-run-b}
\begin{aligned}
\hat b_j(\zeta)&:=\beta[\hat E^{(j)}_\zeta],\\
\hat\zeta_0&:=\zeta,\qquad \hat\zeta_{k+1}:=\hat\zeta_k+\hat b_{k+1}(\zeta)-\hat b_{k+1}(0),\\
\hat x_0&:=x_0,\qquad \hat x_{k+1}:=\hat x_k-\hat b_{k+1}(0),\\
\hat\zeta_k(p)&:=w(p)+\sum_{j\le k}\varphi_j(p)\sum_{y\in\mathbb{T}^{(j)}}h^{(j)}_y(p)
\bigl[\hat b_j(w(p_y))-\hat b_j(0)\bigr],
\end{aligned}
\end{equation}
where $\varphi_j$, $h^{(j)}_y$ are the partition functions of the running shift of the
correlation theorem. The second line is read at a frozen source $\zeta$; at a local source $w$
the running shift is the function $p\mapsto\hat\zeta_k(p)$ of the last line, the hatted form of
the running shift of the correlation theorem. The identity of
the correlation theorem for $\beta$ of the output family, applied with $\hat Q$ in place of $Q$,
gives
\begin{equation}\label{9.10:eq-refiled-run-2}
\hat b_j(\zeta)=\beta^0_j+f_j(\hat\zeta_0,\ldots,\hat\zeta_{j-1})+\beta[\hat Q^{(j)}_\zeta],
\end{equation}
where $\beta^0_j$ is the one-loop part and $f_j$ the pure part of the increment, as in the
correlation theorem.
On the scaffold $\Theta_X$, writing the subscript $(r)$ for depth $r$, with $\hat b_{(r)}$ the
increment of the step at depth $r$, and measuring the shift from the scaffold point
$\theta_r=X+c_0r$, the recursion of the uniqueness theorem for the running shift reads
\begin{equation*}
\hat\zeta_{(r)}:=\hat\zeta_{(r+1)}+\hat b_{(r)}-c_0.
\end{equation*}
By the identity of the re-filing lemmas relating the two runs,
\begin{equation}\label{9.10:eq-refiled-run-3}
\hat x_k=x_k-\sum_{i:\,i''\le k}\beta\bigl[\mathsf{RF}[R^{(i)}]\bigr],
\end{equation}
where $x_k$ is the inverse coupling of the correlation theorem for the same bare datum.

\item[(e)] \emph{Wilson weight and gauge fixing.} The weight of the Wilson term at level $k$ is
$\hat x_k(\cdot)-\hat\zeta_k(\cdot)$, where $\hat x_k(\cdot)$ is the position-dependent inverse
coupling of Definition~\ref{1.5:def-domains} built from $\hat x_k$, and the gauge-fixing
coefficient is
\begin{equation*}
a_j(y):=\hat x_j-\hat\zeta_j(p_y).
\end{equation*}
Both are carried by $\hat\zeta$. This is legitimate because the gauge-fixing statement of the
effective-action lemma of the correlation theorem holds for every configuration $U$; hence the
gauge-fixing lemma of the correlation theorem keeps a single coefficient on the sum of the
gauge-fixing term and the action, and $P^{\flat}\simeq\mathsf{P}_\xi$ at $\xi_i=\hat\zeta_i$,
where $\mathsf{P}_\xi$ is the pure family of the correlation theorem at the shift parameters
$\xi=(\xi_i)$.

\item[(f)] \emph{Choices.} Every choice of the correlation theorem's construction --- the
thresholds $\varepsilon_k$, $\delta_k$, the radius $\check R_k$, the integer $\mathsf{N}$,
the numbers $p_i(g_k)$, the radii of
\cite[(2.28)]{B-CMP119}, the domains, labels, backgrounds, operators, trees, the graph averaging,
and the choices made in the correlation theorem's flat-class bound for the output family of a
step --- is made by the same rule as there: at the run's real inverse coupling, which is now
$\hat x_k$, and, along the depth indexing on $\Theta_X$ used by the
uniqueness theorem, at $\theta_r$. The function $\chi_k$, the measures and forms of the sourced
large-field operation $\mathbf{T}^w_k$ of the correlation theorem, the objects of
\cite[(2.19)--(2.22), (3.2), (3.3), (3.16)]{B-CMP119} and the
surrogate of \cite[(1.100)]{B-CMP122a} are the same as in the correlation theorem; none of them
depends on the source.

\item[(g)] \emph{The action read by step $k$.} The action read by step $k$ of the re-filed run
(the effective action written $E_k$ in \cite{B-CMP119}, together with the remnant terms from
which the next step's remnants are produced) is
\begin{equation}\label{9.10:eq-refiled-run-c}
\begin{aligned}
&\sum_{j\le k}(\hat E^{(j)})^{\mathrm{ren}}
+\sum_{i\in\mathfrak{B}(k)}(R^{(i)})^{\mathrm{vac}}
+\sum_{k_1(k)<i\le k}R^{(i)}\\
&\qquad+\hat{\mathbf{B}}^w_k+\hat{\mathbf{D}}^w_k+\hat{\mathbf{R}}^{w,\mathrm{unf}}_k
-\hat{\mathbf{E}}^w_k-A\bigl(\hat x_k(\cdot)-\hat\zeta_k,U_k\bigr).
\end{aligned}
\end{equation}
Here the sum over $k_1(k)<i\le k$, that is over $i\in\mathrm{band}(k)$, is the sum of earlier
remnants from which Ba{\l}aban produces $R''$ in \cite{B-CMP119}, unchanged. The renormalised group
$(\hat E^{(j)})^{\mathrm{ren}}$ carries, besides the groups of $E^{(j)}$, the groups of the
re-filed remnants with their own counterterm
\begin{equation*}
-\beta\bigl[\mathsf{RF}[R^{(i)}]\bigr]\,A\bigl(h^{(i'')}_y,\cdot\bigr),
\end{equation*}
their coupling part having passed into $\hat x$ by \eqref{9.10:eq-refiled-run-3}. The terms
$\hat{\mathbf{B}}^w_k$ and $\hat{\mathbf{D}}^w_k$ are the objects $\mathbf{B}^w_k$,
$\mathbf{D}^w_k$ of the correlation theorem formed in the re-filed run;
$\hat{\mathbf{E}}^w_k$ is likewise the corresponding term $\mathbf{E}^w_k$ of the correlation
theorem's effective action formed in the re-filed run; and $A(\cdot,\cdot)$ is the weighted
Wilson action of Definition~\ref{1.1:def-wilson-action}. The unfiled term is
\begin{equation}\label{9.10:eq-refiled-run-4}
\hat{\mathbf{R}}^{w,\mathrm{unf}}_k
:=\sum_{i:\,i''\le k}\ \sum_{(X,y)}\mathsf{sh}(X,y)\bigl[R^{(i)}(X,U_k;w)-R^{(i)}(X,1;w)\bigr],
\end{equation}
the inner sum running over the pairs $(X,y)$ of level $i$ such that $y$ is not in the set
$\Lambda^{\circledR}_{i''}$ of points of the small-field region eligible for filed groups, or
$Y(X,y)$ is not contained in the region $\Lambda_{i''}$ of level $i''$ in which Ba{\l}aban's rule
requires the domain of a filed pair to lie. These are exactly the
pairs $(X,y)$ for which the associated pair $(Y(X,y),y)$ at level $i''$ is not filed under
Ba{\l}aban's rule \cite{B-CMP119} deciding which pairs are filed into the output family of a
step.
\end{itemize}

By the identity of the re-filing lemmas relating the two runs, at every configuration the
exponent of the re-filed run equals, by an add-and-subtract identity, the exponent of the
correlation theorem computed with the re-filed run's own objects. The five
operations of the effective-action lemma of the correlation theorem do not depend on which run
computed their inputs, and the partition function $Z$ and the Schwinger functions
$S_{\mathsf{p}}$ with $\mathsf{p}$ test functions are the same whichever of the two runs is used
to compute them.
\end{definition}

\begin{lemma}[Bounded level range of unfiled remnant shares]\label{9.10:lem-unfiled-range}
Use the levels, depths and filing levels of Notation~\ref{9.10:not-filing-level}: a remnant born
at level $i$ has $\mathsf{a}_i\ge 2$ and filing level $i''=i+\mathsf{a}_i$ as in
\eqref{9.10:eq-filing-level-a}. Its shares $\mathsf{sh}(X,y)$ are those of the re-filed
renormalisation run of Definition~\ref{9.10:def-refiled-run}. Assume that the partition-cube
radii satisfy $\check R_j\ge 2$ at every level $j$. Let $k$ be a level, and let $(X,y)$ be a pair
of level $i$ whose share is unfiled at level $k$, with $i''\le k$.
\begin{itemize}
\item[(a)] If $\bar X^{\sim 1}\subset\Lambda_{i''}$, where $\Lambda_{i''}$ is the complement of
the large-field region $Z_{i''}$ at level $i''$ (a \emph{small} unfiled pair), then $y$ and
$X$ lie outside $\Omega_{i''+1}$. The share $\mathsf{sh}(X,y)$ is present only in zones $n$ with
$i\le n\le i''$.
\item[(b)] If $\bar X^{\sim 1}\not\subset\Lambda_{i''}$ (a \emph{large} unfiled pair), then
either $d_i(X)\ge\mathsf{t}^{-1}$, where $\mathsf{t}^{-1}$ is the size threshold on $d_i(X)$ in
the correlation theorem's classification of terms, or $\bar X^{\sim 1}$ meets $\Lambda_{i''}^c$.
Such pairs are tail contributions of the first and third rows of the correlation theorem's
classification of terms.
\item[(c)] Consequently, in every zone $n$, the unfiled remnant shares present (the large pairs
of \textup{(b)} being counted among the tail terms of the correlation theorem and not among the
remnant shares) have levels $i$ with
\begin{equation}\label{9.10:eq-unfiled-range-a}
n-\mathsf{a}_i\le i\le n .
\end{equation}
In every zone $n$, therefore, an unfiled remnant share of level $i$ present there lies at most
$\mathsf{a}_i$ levels below $n$. The lemma places no restriction on the zone age $k-n$.
\end{itemize}
\end{lemma}

\begin{proof}
The floor $\check R_j\ge 2$ is an upper bound on the coupling. In Ba{\l}aban's construction the
radius at level zero satisfies $\check R_0\ge(\log g_0^{-2})^{r}$ with an exponent $r\ge 2$
\cite{B-CMP119}.

\emph{Part} (a). Let $(X,y)$ be a small unfiled pair of level $i$ with $i''\le k$. Measure
distances in units of the lattice spacing at level $i''$. By the filing rule of the re-filed
renormalisation run of Definition~\ref{9.10:def-refiled-run}, the share of a small pair of level
$i$ with $i''\le k$ is left unfiled only if $y$ lies within distance $(\check R_{i''}+4)M$ of
$\Lambda_{i''}^c$. Hence $y$ lies within distance $(\check R_{i''}+4)M$ of $\Lambda_{i''}^c$.

We first bound the distance from $\Lambda_{i''}^c$ to $\Omega_{i''+1}$. At level $i''$ the
large-field region is $Z_{i''}=\Lambda_{i''}^c$. Let $\bar Z_{i''}$ be the union of all
$LM\check R_{i''+1}$-cubes that intersect $Z_{i''}$. In the construction of \cite{B-CMP119},
$\bar Z_{i''}$ is surrounded by four layers of $LM\check R_{i''+1}$-cubes. The decomposition of
unity of the next step is taken for the domain $(\bar Z_{i''}^{\sim 4})^c$. Hence
$\Omega_{i''+1}$ keeps four layers of $LM\check R_{i''+1}$-cubes from $\Lambda_{i''}^c$.

The radius $\check R$ decreases by at most one power of $L$ per step, so
$L\check R_{i''+1}\ge\check R_{i''}$. Four layers of $LM\check R_{i''+1}$-cubes therefore have
width at least $4M\check R_{i''}$. Every point of $\Omega_{i''+1}$ is thus at distance at least
$4M\check R_{i''}$ from $\Lambda_{i''}^c$.

Next we locate $X$. We have $X\subset\tilde\square(y)^{\sim 2}$, the $M$-cube at $y$ enlarged by
two layers of $M$-cubes. The enlargement contributes a further $2M$ to the distance of the
points of $X$ from $\Lambda_{i''}^c$. Now
\begin{equation*}
4\check R_{i''}-(\check R_{i''}+4)-2=3\check R_{i''}-6,
\end{equation*}
and the right side is non-negative exactly when $\check R_{i''}\ge 2$, which is the assumed
floor. Hence
\begin{equation*}
(\check R_{i''}+4)M+2M\le 4M\check R_{i''} .
\end{equation*}
So $y$ and $X$ lie within the four-layer margin around $\Lambda_{i''}^c$, outside
$\Omega_{i''+1}$. The share therefore sits in zones $n$ with $n\le i''$.

It also sits in zones with $n\ge i$, because $X\subset\Omega_i$. This proves (a).

\emph{Part} (b). Let $(X,y)$ be a large unfiled pair, so $\bar X^{\sim 1}\not\subset\Lambda_{i''}$.
Then either $d_i(X)\ge\mathsf{t}^{-1}$ or $\bar X^{\sim 1}$ reaches $\Lambda_{i''}^c$. In either
case the pair belongs to the tail terms of the first and third rows of the correlation theorem's
classification, and it is estimated there. This proves (b).

\emph{Part} (c). Let $n$ be a zone present at level $k$, and let an unfiled remnant share of
level $i$ be present in it, not coming from a large pair. If $i''\le k$, part (a) gives
$i\le n\le i''$. If $i''>k$, then $n\le k<i''$, since a zone present at level $k$ has $n\le k$,
and $n\ge i$ because $X\subset\Omega_i$; again $i\le n\le i''$. Since $i''=i+\mathsf{a}_i$, this
is equivalent to $n-\mathsf{a}_i\le i\le n$, which is \eqref{9.10:eq-unfiled-range-a}. The
statement on the length of the level range in zone $n$ is a restatement of
\eqref{9.10:eq-unfiled-range-a}.

Neither (a) nor (b) constrains $k-n$. The zone age of the unfiled shares present in zone $n$ is
therefore unrestricted by this argument.
\end{proof}

The zone age $k-n$ is controlled instead by the lemma on geometric forgetting of old regions. If
a zone $n<k$ is present at level $k$ inside $X^{\sim}$, then there is a live large-field component
of age at least $k-n-1$. This is the content of the lemma on geometric forgetting of old regions
and of the first sentence of part~(d) of the theorem called the one-lattice Lipschitz step.

\begin{remark}[Healing and the last lattice]\label{9.10:rem-healing-last-lattice}
Consider the re-filed renormalisation run of Definition~\ref{9.10:def-refiled-run}, with its
hatted families as displayed in \eqref{9.10:eq-refiled-run-c}.
\begin{enumerate}
\item[(a)] \emph{Healing.} Let a large-field region $Z$ be healed at a site of the large-field
operation $\mathbf{R}$ of level $k$. The terms added back for the newly regular region are the
following members of the hatted families:
\begin{itemize}
\item the groups $\hat{\mathfrak{i}}^{(j)}_y$ of the families $\hat E^{(j)}$, $j\le k$, at the points
$y$ that now lie in $\Lambda_j^{\circledR}$. For every level $j=i''\le k$ at which filing has
already taken place, these include the groups of $\mathsf{RF}[R^{(i)}]$ at pairs that were
unfiled only because they lay beside $Z$; thus unfiled pairs of the healed region become filed;
\item the units of the levels in $\mathfrak{B}(k)$ and in $\mathrm{band}(k)$;
\item the unfiled shares beside the part of $Z$ that survives the step.
\end{itemize}
Their copies are subtracted from the term $V(Y)$ of the new action on the domain $Y$ at the
healing site: the subtracted copies are placed in $V(Y)$ and hence in the level-$k$ slots
$R'^{(k)}$ and $\mathbf{B}'^{(k)}$. Healing is an add-and-subtract of the hatted families: it
produces no new kind of term and no new memory.
\item[(b)] \emph{Last lattice.} Let $K'=n-\mathsf{w}$ be the last level of the run, where $n$ is the
top level of the run, from which depths are counted, and $\mathsf{w}$ is the last-step offset.
Remnants whose filing level satisfies $i''>K'$ are never filed. These are the remnants of depths
$s<\mathsf{a}_s+\mathsf{w}$, that is, of the shallowest $\mathsf{a}$ levels of the run. Their size is
only $g^{\kappa_0}$-small. With $C_E$ the constant and $\check{\mathfrak{r}}_s$ the depth-dependent
factor of the pair-difference bound for remnants in the re-filing lemmas, their pair difference
$D^R_s$ satisfies
\begin{equation}\label{9.10:eq-healing-last-lattice-1}
D^R_s\le 4C_E\,\check{\mathfrak{r}}_s\,\rho^{\,n-s}
\le 4C_E\,\check{\mathfrak{r}}_s\,\rho^{-\mathsf{a}_s-\mathsf{w}}\,\rho^{\,n},
\end{equation}
and is therefore geometric in $n$. The vacuum pieces $e^{(j)}(X;w)$, the pieces
$\log[\mathfrak{z}(a_j(y))/\mathfrak{z}(\hat x_j)]$ and the vacuum energy $\mathcal{E}$ (the double
sum of the pieces $e^{(j)}$) have the same form as in the construction of the correlation theorem,
with
\begin{equation}\label{9.10:eq-healing-last-lattice-2}
a_j(y)=\hat x_j-\hat\zeta_j(p_y).
\end{equation}
\end{enumerate}
\end{remark}

\begin{proof}
(a) In \cite{B-CMP122b} healing works as follows. When a large-field region $Z$ is healed, the new
action contains all the necessary $E$-terms and all the $R$-terms whose domains meet $Z$, and the
same terms are subtracted from $V(Y)$. The correlation theorem carries this step over in the same
form: in the display of its construction that records the healing step, and in parts (b) and (c)
of its proposition on remnants.

In the re-filed run the necessary terms for the newly regular region are read off from
\eqref{9.10:eq-refiled-run-c}. First, a point $y$ that now lies in $\Lambda_j^{\circledR}$, with
$j\le k$, carries the group $\hat{\mathfrak{i}}^{(j)}_y$ of $\hat E^{(j)}$. At a level $j=i''\le k$
at which filing has already occurred, a pair of $\mathsf{RF}[R^{(i)}]$ that was unfiled only because
it lay beside $Z$ now meets the filing condition, and its group is added as a filed group. Second,
the units of $\mathfrak{B}(k)$ and $\mathrm{band}(k)$ are added. Third, the unfiled shares lying
beside the part of $Z$ that survives the step are added.

Every added term belongs to one of the hatted families of \eqref{9.10:eq-refiled-run-c}. Every
subtracted copy is placed in $V(Y)$ and hence in the level-$k$ slots $R'^{(k)}$ and
$\mathbf{B}'^{(k)}$. No new kind of term therefore arises.

As for memory, a re-added group of scale $i''$ is read at every later level $k'>k$ with the
geometric kernel $\varsigma_{k'+1-i''}$. Since $k'\ge k+1$ and $i''\le k$, the index satisfies
$k'+1-i''\ge 2$. The memory of these groups is therefore geometric a fortiori.

(b) A remnant of depth $s$ is created at level $i=n-s$ and filed at level $i''=i+\mathsf{a}_i$. In
the depth indexing, with $\mathsf{a}_s$ the value of $\mathsf{a}_i$ at $i=n-s$, this reads
$i''=n-s+\mathsf{a}_s$. Hence
\begin{equation*}
i''>K'=n-\mathsf{w}\iff s<\mathsf{a}_s+\mathsf{w}.
\end{equation*}
Such a remnant is never filed within the run. The only bound on its size is therefore the class
bound of Definition~\ref{9.10:def-refiled-run}, which is a power $g^{\kappa_0}$ of the coupling.

The first inequality of \eqref{9.10:eq-healing-last-lattice-1} is the pair-difference bound for
remnants in the re-filing lemmas. For the second, write
\begin{equation*}
\rho^{\,n-s}=\rho^{\,n}\rho^{-s}.
\end{equation*}
Since $0<\rho<1$, the factor $\rho^{-s}$ increases with $s$, and $s<\mathsf{a}_s+\mathsf{w}$ gives
$\rho^{-s}\le\rho^{-\mathsf{a}_s-\mathsf{w}}$. The resulting bound is a constant times $\rho^{\,n}$,
so it is geometric in $n$.

The vacuum pieces, the pieces $\log[\mathfrak{z}(a_j(y))/\mathfrak{z}(\hat x_j)]$ and their double
sum $\mathcal{E}$ are defined from the same data and by the same formulas as in the construction of
the correlation theorem. The gauge-fixing coefficient is \eqref{9.10:eq-healing-last-lattice-2}.
\end{proof}

\begin{definition}[Re-filed lists and the substitution rule]\label{9.10:def-substitution-rule}
We use the levels and depths, the numbers $\mathsf{a}_i$, the filing level $i''=i+\mathsf{a}_i$, the
sets $\mathfrak{B}(k)$ and $\mathrm{band}(k)$ and the range letter $\ell_s$ of
Notation~\ref{9.10:not-filing-level}, display~\eqref{9.10:eq-filing-level-a}, and the re-filed
renormalisation run of Definition~\ref{9.10:def-refiled-run}. In that run the remnant
$R^{(i)}=R'^{(i)}+R''^{(i)}$ created at level $i$ is re-filed at level $i''$ as the pointed
family $\mathsf{RF}[R^{(i)}]$. Local sources are $w\in\mathcal{W}_k$ and frozen complex sources
are $\hat w\in D(0,2r)$. In item~(b), $r$ is the parameter of the frozen-source disc $D(0,2r)$.
In item~(d) the subscript $r$ of $\theta_r$,
$\mathfrak{e}^P_r$ and $\zeta_{(r)}$ is a depth. In item~(f), $r$ denotes the output depth.

\smallskip
\noindent(a) \emph{Hatted lists.} The level-$j$ lists of the first large-field estimate,
\begin{equation*}
\mathfrak{s}_j=\bigl(\zeta_j;\ \mathsf{P}^{(j)};\ P^{\flat(j)},\,Q^{(j)},\,R^{(j)};\
E^{(j)}(\cdot\,;w),\,R^{(j)}(\cdot\,;w)\bigr),
\end{equation*}
are replaced by the \emph{hatted lists}
\begin{equation*}
\hat{\mathfrak{s}}_j=\bigl(\hat\zeta_j;\ \mathsf{P}^{(j)};\ P^{\flat(j)},\,\hat Q^{(j)},\,R^{(j)};\
\hat E^{(j)}(\cdot\,;w),\,R^{(j)}(\cdot\,;w)\bigr).
\end{equation*}
The running shift $\hat\zeta_j$, the dirty part $\hat Q^{(j)}$ and the regular family
$\hat E^{(j)}$ are the hatted ones. The slots $\mathsf{P}^{(j)}$, the pure part
$P^{\flat(j)}$ and the remnant $R^{(j)}$ are those of the unhatted list.

\smallskip
\noindent(b) \emph{The class $\hat{\mathfrak{K}}_k$.} A family of hatted lists belongs to the
class $\hat{\mathfrak{K}}_k$ if, at each of its levels $j$,
\begin{equation}\label{9.10:eq-substitution-rule-1}
\begin{gathered}
\|\hat E^{(j)}\|^{\flat}\le E_0,\qquad \|P^{\flat(j)}\|^{\flat}\le E_0,\qquad
\|R^{(j)}\|^{\flat}\le g_j^{\kappa_0},\\
\|\hat Q^{(j)}\|^{\flat}\le \hat{\mathfrak{m}}_{j-1}:=10(1+C_{RF})\,y_{j-1}^{-\kappa_0/2},\\
|\hat\zeta_j|\le \tfrac{8}{3}\,r,\qquad |\hat b_j|\le b_+ +\tfrac{2}{3}\,r,\qquad
|\mathsf{P}^{(\le j)}|_j\le 2\mathsf{E}.
\end{gathered}
\end{equation}
Here $E_0$ is the norm ceiling of the class for the regular family and the pure part,
$\kappa_0$ is the exponent of the remnant bound, $y_{j-1}$ is the value at level $j-1$ of the
variable $y$ in which the bound on the dirty part is written, and $\mathsf{E}$ is the ceiling
for the size $|\mathsf{P}^{(\le j)}|_j$ at
level $j$ of the slots $\mathsf{P}$ up to level $j$.
The bound imposed on $\hat Q^{(j)}$ is the one provided by the proposition that bounds the dirty
part $\hat Q^{(j)}$ in the re-filed run of Definition~\ref{9.10:def-refiled-run}.
Under the upper bound on the coupling $10(1+C_{RF})\le C_{\mathfrak{m}}\,y$, one has
$\hat{\mathfrak{m}}_{j-1}\le\mathfrak{m}_{j-1}$. The hatted dirty part therefore satisfies the
bound $\mathfrak{m}_{j-1}$ of the unhatted class.

\smallskip
\noindent(c) \emph{Difference norms.} For two runs compared slot by slot, $\delta$ denotes the
difference of corresponding slots. We put
\begin{equation}\label{9.10:eq-substitution-rule-2}
\begin{gathered}
\hat{\mathsf{e}}_j:=\sup_{\hat w,\,w}\|\delta\hat E^{(j)}\|^{\flat},\qquad
\hat{\mathsf{q}}_j:=\|\delta\hat Q^{(j)}\|^{\flat},\qquad
\mathsf{r}_j:=\|\delta R^{(j)}\|^{\flat},\\
\hat{\mathsf{z}}_j:=\sup|\delta\hat\zeta_j|,\qquad
\hat{\mathsf{b}}_j:=\sup|\delta\hat b_j|.
\end{gathered}
\end{equation}
In $\hat{\mathsf{e}}_j$ the supremum runs over the frozen sources $\hat w$ and over the local
sources $w$. The norm $\mathsf{r}_j$ is unchanged from the unhatted lists. So are the
difference norms $\mathsf{p}_j$ and $\mathsf{p}^{\flat}_j$ of the slots $\mathsf{P}^{(j)}$ and
$P^{\flat(j)}$.

\smallskip
\noindent(d) \emph{The reduced dictionary.} The hatted lists enter the reduced slots of the
depth-indexed comparison as follows. One has $C_t\theta_r\hat Q\in\mathfrak{e}^P_r$, and the
three reduced slots, written $P$, $A$ and $R$ in the depth-indexed comparison (the letters of
$\mathfrak{e}^P_r$, $G^P_s$, $G^A_s$ and $G^R_s$), are, together with the shift and the increment,
\begin{equation}\label{9.10:eq-substitution-rule-3}
\begin{gathered}
P:=\bigl(\mathsf{P}/E_0;\ C_t\theta_s\hat Q\bigr),\qquad
A:=\mathsf{l}_s^{-1}\bigl(P^{\flat};\hat D\bigr)/E_0,\\
R:=C_t\theta_s\bigl(R'+R''\bigr),\qquad \zeta_{(r)}:=\hat\zeta,\qquad b:=\hat b.
\end{gathered}
\end{equation}
The remnant slot $R$ here is unchanged from the unhatted dictionary.

\smallskip
\noindent(e) \emph{Stored numbers and the uses of a filed remnant.} The array $x^{\mathsf{c}}$
of stored numbers of the last-lattice comparison consists of the following entries:
\begin{itemize}
\item the slots $\hat E^{(j)}$, $P^{\flat(j)}$, $\hat Q^{(j)}$, $R^{(j)}$, $\hat D^{(j)}$,
$\hat{\mathbf{B}}^{(j)}$ and $\hat{\mathbf{B}}'^{(j)}$;
\item four further terminal entries of the array;
\item the $\hat\zeta$-entries and the $\hat b$-entries;
\item the weight entries $\hat x-\hat\zeta$;
\item the source-free pieces $\mathsf{z}$, namely the vacuum constants, the
$\mathfrak{z}$-pieces and $\mathsf{s}^0/D^0$.
\end{itemize}
A filed remnant $R^{(i)}$ is a stored slot, and it is used in exactly two ways.
\begin{itemize}
\item[($\alpha$)] It is used as a remnant at the levels $k'\in[i,i''-1]$, through
$\mathfrak{B}$ and the band. It is also used, as unfiled shares, in the zones
$n\in[i,i'']$ at any later level.
\item[($\beta$)] From level $i''$ on, it is used as part of the regular family, through the
stored slot $\hat E^{(i'')}$. The two-run difference of $\hat E^{(i'')}$ contains
$\mathsf{RF}[\delta R^{(i)}]$, which is linear in $\delta R^{(i)}$, and
\begin{equation}\label{9.10:eq-substitution-rule-4}
\mathsf{S}^+[\hat E^{(i'')}]\le \mathsf{S}^+[E^{(i'')}]
+C_{RF}\,w_{i''-1,i}\,\mathsf{S}^+[R^{(i)}],
\qquad w_{i''-1,i}\le 3\ \text{on }\Theta_X.
\end{equation}
\end{itemize}

\smallskip
\noindent(f) \emph{The substitution rule.} Consider the loads, kernels, inequalities and
envelopes of the comparison estimates of this chapter in which a remnant slot is read as a
remnant. The rule concerns the following places:
\begin{itemize}
\item the remnant term of the zone bound, at zone weight $w_{n,j}$;
\item the remnant contribution $\mu_k$ in the proposition bounding the dirty part of the
unhatted lists;
\item the sums $\sum_{j\le k_1(k)}w_{k,j}\mathsf{r}_j$ and $\sum_{j\le n}w_{n,j}\mathsf{r}_j$,
where $k_1(k)$ is the last level summed at zone level $k$, the remnant difference term
$\bar w_{i,j}\,\delta^R_j$ and the kernel
$\mathfrak{w}^R$ of the one-lattice Lipschitz step;
\item the sums over depths $s>r$ of the remnant norms $\|\cdot^R_s\|$ in the depth-indexed
bounds and in their transplanted form;
\item the sum $\sum_{s>r}\|G^R_s\|$ in the bracket $B_r(G,d)$ of the $R'$-inequality;
\item the corresponding sums in the one-step input bounds of the two-cutoff comparison;
\item the sums $\sum_{i\le j}w_{j,i}\,\mathsf{S}^+_i[R]$ and
$\sum_{i\le j}w_{j,i}\,g_i^{\kappa_0}\hat\vartheta^{\natural}_i$ of the last-lattice comparison.
\end{itemize}
In each of them the index set of all $j\le n$ (respectively of all $j\le k_1$), for a zone
letter $j$ and a zone level $n$, and the index set of all $s>r$, for a depth letter $s$ and an
output depth $r$, are replaced as follows:
\begin{equation}\label{9.10:eq-substitution-rule-a}
\begin{aligned}
\{\,j:\ j\le n\,\}\ &\longmapsto\ \{\,j:\ j\le n\le j''\,\}
=\{\,j:\ n-\mathsf{a}_j\le j\le n\,\},\\
\{\,s:\ s>r\,\}\ &\longmapsto\ \{\,s:\ r<s\le r+\ell_s\,\},\qquad
\ell_s=\mathsf{a}_s+2\ \text{as in \eqref{9.10:eq-filing-level-a}} .
\end{aligned}
\end{equation}
The summand $2$ in $\ell_s$ covers the boundary cases of a remnant filed into the output of a
step through $\mathfrak{B}$, a remnant in the band, and a remnant read through its unfiled
shares. The replacement is justified as follows. By Definition~\ref{9.10:def-refiled-run} and
Lemma~\ref{9.10:lem-unfiled-range}, no remnant whose index lies outside these sets is present as
a remnant at zone level $n$ (respectively at output depth $r$).

In addition, every reading of a regular slot $E$ in these estimates is the reading of the
hatted slot $\hat E$. This applies to the first and second rows and the row marked $\partial$
of the zone bound of item~(f), to the
$\varsigma$-kernels, to $\mathsf{S}^+[E]$, and to $\|G^P\|$ and $\|G^A\|$. The norm of $\hat E$
carries, at the level $k$ read, the additive filing term
$C_{RF}\sum_{i\in\mathfrak{B}(k)}w_{k,i}\,\|\delta R^{(i)}\|^{\flat}$ of the norm
bound for the re-filed family, and this term lies inside the same finite range.

The following fact is not part of the definition; it is stated here together with the inputs on
which it depends. The substitution rule is exhaustive: no term of the comparison estimates of
this chapter other than the places listed in item~(f) reads a remnant. This fact depends on
three inputs: the enumeration
of the places at which the correlation theorem, the two large-field estimates and the two-cutoff
comparison evaluate a level-$j$ list term; the enumeration of the estimates bounding a remnant,
which shows that each of them is restricted by the age of the remnant; and the hatted form of
the one-step bounds that enter these estimates as inputs.
\end{definition}

\begin{remark}\label{9.10:thm-refiled-induction}
The statement of this theorem is not written out here; parts of its proof are printed below.
\end{remark}

\begin{lemma}[Induction step: frozen integral and flow]\label{9.10:prf-step-frozen-flow}
Let the hypotheses of Theorem~\ref{9.10:thm-refiled-induction} hold, let $0\le k<K$, and assume:
\begin{enumerate}
\item[(1)] the hatted induction hypothesis $(\hat J_k)$ of Theorem~\ref{9.10:thm-refiled-induction},
with its clauses (a)--(d);
\item[(2)] the frozen-integral proposition for the re-filed run;
\item[(3)] the flow lemma of the correlation theorem;
\item[(4)] the two statements of the correlation theorem that carry the reference relation of
clause~(a) of $(\hat J_k)$ from the pairs up to level $k$ to the pairs up to level $k+1$.
\end{enumerate}
Here $y_k$, $v_k$, $e_k$, $\lambda$, $\sigma_0$, $\beta_{\mathrm{pr}}$, the coefficients $d_j(\varrho)$
and the inputs $\mathfrak{d}_j$, $\mathfrak{d}_*(k)$, $\omega_k$ of the flow lemma are the quantities so
denoted in the correlation theorem, and $\hat y_k$ is $y_k$ read for the re-filed run.
Then the following hold.
\begin{enumerate}
\item[(i)] The frozen family $\hat E^{(k+1)}_\zeta$ is defined on $\hat\alpha_k$ and is analytic in
$\zeta\in D(0,2r)$, with
\begin{equation*}
\|\hat E^{(k+1)}_\zeta\|^{\flat}\le \tfrac{9}{16}E_0+\hat{\mathfrak{m}}_k\le E_0,
\qquad
\hat{\mathsf{n}}_{k+1}\le \hat{\mathfrak{m}}_k ,
\end{equation*}
and $\hat{\mathsf{q}}_{k+1}:=\beta[\hat Q^{(k+1)}_\zeta]$ satisfies, on $D(0,2r)$,
\begin{equation}\label{9.10:eq-step-frozen-flow-1}
|\hat{\mathsf{q}}_{k+1}|\le \hat v_k:=C_\beta C_t\,x_k\,\hat{\mathfrak{m}}_k\le v_k .
\end{equation}
\item[(ii)] For every $\varrho\in[r,2r]$ and every $\zeta\in D(0,\varrho)$,
\begin{equation*}
|\hat b_{k+1}(\zeta)-\hat b_{k+1}(0)|\le d_k(\varrho)\,|\zeta|,
\qquad
\sum_{j\le k}d_j(\varrho)\le \tfrac13 ,
\end{equation*}
and $|\hat\zeta_{k+1}-\zeta|\le\frac13|\zeta|$ on $D(0,2r)$.
\item[(iii)] One has
\begin{equation*}
|\hat b_{k+1}(0)-\beta^0_{k+1}|\le C_\beta e_k+\hat v_k\le\lambda ;
\end{equation*}
with $\hat x_{k+1}:=\hat x_k-\hat b_{k+1}(0)$, the reference relation of clause~(a) holds for all pairs
$(\hat x_i,\hat x_j)$ with $i<j\le k+1$; the quantities $\hat g_{k+1}$, $\check R_{k+1}$,
$\varepsilon_{k+1}$, $\alpha_{\cdot,k+1}$ are defined at $\hat x_{k+1}$ and satisfy
$(1+\beta_{\mathrm{pr}})\,\alpha_{\cdot,k+1}\le\hat\alpha_k$; and the rows at level $k+1$, for both
domain classes that it distinguishes, of the correlation theorem's hypothesis on the large-field
operations are satisfied.
\end{enumerate}
\end{lemma}

\begin{proof}
The proof has two parts: the frozen whole-lattice integral at level $k$, which yields~(i), and the
flow and reference at level $k+1$, which yield~(ii) and~(iii).

\smallskip
\noindent\textit{The frozen integral.}
The frozen-integral step of the correlation theorem uses, at level $k$, the band parameter
$\mathsf{T}_k$, the frozen families $E^{(j)}_\zeta$ for $j\le k$, the remnants $R^{(j)}_\zeta$ for
$j$ up to a level strictly below $k$, and the bound $|\zeta_k|\le\frac83 r$; all of these are part of
the induction hypothesis of that theorem at level $k$. For the re-filed run the frozen-integral
proposition (hypothesis~(2)) uses the following inputs, and we check that each is available.
\begin{itemize}
\item The band parameter $\mathsf{T}_k$ is the same as for the un-hatted run.
\item The frozen families $\hat E^{(j)}_\zeta$, $j\le k$, with $\|\hat E^{(j)}_\zeta\|^{\flat}\le E_0$
and $\|\hat Q^{(j)}_\zeta\|^{\flat}\le\hat{\mathfrak{m}}_{j-1}$, are given by clause~(c) of
$(\hat J_k)$.
\item The remnants $R^{(i)}_\zeta$ with $i\in\mathfrak{B}(k)$, the set of levels filed into the output
of step $k$, defined in \eqref{9.10:eq-filing-level-a}. For such $i$ the filing level is
$i''=i+\mathsf{a}_i=k+1$, hence $i=k+1-\mathsf{a}_i\le k-1$, and clause~(d) of $(\hat J_k)$ gives
$\|R^{(i)}_\zeta\|^{\flat}\le g_i^{\kappa_0}$.
\item The units of the band $\mathrm{band}(k)$ of \eqref{9.10:eq-filing-level-a}. They are the source
terms $R''$ of Ba{\l}aban's construction, read by clause~(f) of the remnant proposition inside the
second expectation of the step of the large-field operation $\mathbf{T}$, \cite[(3.26)]{B-CMP119},
and they are the same as for the un-hatted run.
\item The bound $|\hat\zeta_k|\le\frac83 r$ on $D(0,2r)$. By clause~(b) of $(\hat J_k)$,
$|\hat\zeta_k-\zeta|\le\frac13|\zeta|$, so for $|\zeta|<2r$
\begin{equation*}
|\hat\zeta_k|\le\tfrac43|\zeta|\le\tfrac43\cdot 2r=\tfrac83 r .
\end{equation*}
\end{itemize}
The frozen-integral proposition therefore applies at level $k$. Its proof is the correlation
theorem's proof of the frozen integral, in which the unmarked groups are those of that proof and the
marked groups are the renormalised groups $\mathfrak{i}_y[\hat Q^{(j)}]$ of the dirty parts, each with
its own coupling extraction $\beta[\hat Q^{(j)}]$; the marking depends on the levels only, so that
\cite[(3.39)]{B-CMP119} holds by the transformation laws \cite[(2.29), (2.32)]{B-CMP119}, and for the
re-filed part by the corresponding clause of the re-filing lemma; every estimate in it is one of the
correlation theorem's, at that theorem's constants, and the re-filed families $\mathsf{RF}[R^{(i)}]$
enter only through their membership in the pointed class and their norm bound, both supplied by the
re-filing lemma. Its conclusion at level $k$ is: the family $\hat E^{(k+1)}_\zeta$ is defined on
$\hat\alpha_k$ and is analytic in $\zeta$; its dirty part satisfies
$\hat{\mathsf{n}}_{k+1}\le\hat{\mathfrak{m}}_k$, where
$\hat{\mathsf{n}}_{k+1}=\sup_\zeta\|\hat Q^{(k+1)}_\zeta\|^{\flat}$; the extraction bound
\eqref{9.10:eq-step-frozen-flow-1} holds on $D(0,2r)$; and
\begin{equation*}
\|\hat E^{(k+1)}_\zeta\|^{\flat}\le\tfrac{9}{16}E_0+\hat{\mathfrak{m}}_k\le E_0 ,
\end{equation*}
the last inequality being the a priori ceiling $\frac{9}{16}E_0+\hat{\mathfrak{m}}_k\le E_0$ among the
hypotheses of Theorem~\ref{9.10:thm-refiled-induction}. This is~(i).

\smallskip
\noindent\textit{The flow.}
The hatted coupling increment decomposes into its one-loop part, its pure part and the extraction
$\hat{\mathsf{q}}_{k+1}$ of the dirty part:
\begin{equation}\label{9.10:eq-step-frozen-flow-2}
\hat b_{k+1}=\beta^0_{k+1}+f_{k+1}(\hat\zeta_0,\dots,\hat\zeta_k)+\hat{\mathsf{q}}_{k+1}.
\end{equation}
We apply the flow lemma of the correlation theorem (hypothesis~(3)) to
\eqref{9.10:eq-step-frozen-flow-2} with the inputs
\begin{equation*}
\mathfrak{d}_j=\sigma_0\hat x_j,\qquad
\mathfrak{d}_*(k)\ge\sigma_0\hat y_k,\qquad
\omega_k=C_\beta e_k,
\end{equation*}
and with the value bound $\hat v_k\le v_k$ of \eqref{9.10:eq-step-frozen-flow-1}. For any
$\varrho\in[r,2r]$ it gives
$|\hat b_{k+1}(\zeta)-\hat b_{k+1}(0)|\le d_k(\varrho)|\zeta|$ on $D(0,\varrho)$. The coefficients
$d_k(\varrho)$ are those defined in the correlation theorem, read at $\hat y_k$; the only input
that differs from the un-hatted run is $\hat v_k\le v_k$, and this makes the right side smaller. The
bound $\sum_{j\le k}d_j(\varrho)\le\frac13$ follows from clause~(a) of $(\hat J_k)$ and the summability
condition on the coefficients $d_j$ in the parameter choice of the correlation theorem, applied with
anchor $\hat x_k>X$. With this sum bound and the definition of the running shift
\eqref{9.10:eq-refiled-run-b}, the flow lemma yields $|\hat\zeta_{k+1}-\zeta|\le\frac13|\zeta|$ on
$D(0,2r)$. This is~(ii).

\smallskip
\noindent\textit{The reference.}
At $\zeta=0$, with the input $\omega_k=C_\beta e_k$ of the flow lemma and the value bound
\eqref{9.10:eq-step-frozen-flow-1},
\begin{equation*}
|\hat b_{k+1}(0)-\beta^0_{k+1}|\le C_\beta e_k+\hat v_k\le\lambda .
\end{equation*}
We put $\hat x_{k+1}:=\hat x_k-\hat b_{k+1}(0)$. The two statements of hypothesis~(4), applied with this
bound and with the reference relation for the pairs up to level $k$ (clause~(a) of $(\hat J_k)$), give
the reference relation for all pairs $(\hat x_i,\hat x_j)$ with $i<j\le k+1$. The quantities
$\hat g_{k+1}$, $\check R_{k+1}$, $\varepsilon_{k+1}$ and $\alpha_{\cdot,k+1}$ are then defined at
$\hat x_{k+1}$, and on the prescribed scaffold $\Theta_X$ at the scaffold value $\theta$, by the same
prescription as in the correlation theorem. They satisfy
$(1+\beta_{\mathrm{pr}})\alpha_{\cdot,k+1}\le\hat\alpha_k$. The rows at level $k+1$ of the correlation
theorem's hypothesis on the large-field operations, for the two domain classes that it
distinguishes, read only the bounds $g_j\le\gamma$ and the relations
\cite[(2.6)--(2.9)]{B-CMP119}, and these are the content of the reference relation just established;
hence these rows are satisfied. This is~(iii).

\smallskip
\noindent\textit{Identification and order of use.}
Statements~(i)--(iii) are the parts of the corollary on the flow of the re-filed run that the present
step supplies at level $k+1$: the reference relation, the running-shift bounds and the value bound for
$(\hat x,\hat\zeta,\hat b,\hat Q)$. That corollary also records the property of $f_{k+1}$ that the
correlation theorem records with the running-shift statement, and the bound
$\hat{\mathsf{N}}_K:=\hat\zeta_{K'}(0)\in[\frac23,\frac43]$; these are not among (i)--(iii).
The argument above uses, from clause~(d) of $(\hat J_k)$, the bound
$\|R^{(i)}_\zeta\|^{\flat}\le g_i^{\kappa_0}$ at levels $i<k$, which is the remnant bound of the
induction at level $i$. This involves no circularity: $\hat b_{k+1}$ reads $R^{(i)}$ only for
$i\in\mathfrak{B}(k)$, so that $i<k$, and the thresholds for $R^{(i)}$ were fixed at levels at most
$i$. The order of use is that of the correlation theorem: the running-shift statement up to level $j$
yields $R^{(j)}_\zeta$ through the remnant bound at level $j$, which in turn yields the running-shift
statement at level $k+1$ through the frozen integral; for the re-filed run the last step is taken for
$j\in\mathfrak{B}(k)$.
\end{proof}

\begin{proof}[Proof of Theorem~\ref{9.10:thm-refiled-induction}, continued: the step with
domains, renormalisation and filing]
\label{9.10:prf-step-filing}
We continue the induction step from level $k$ to level $k+1$, under the induction hypothesis of
Theorem~\ref{9.10:thm-refiled-induction} at level $k$. The sets of levels $\mathfrak{B}(k)$ and
$\mathrm{band}(k)$ and the filing level $i''$ of a remnant created at level $i$ are those of
\eqref{9.10:eq-filing-level-a}. In particular $i''=k+1$ for $i\in\mathfrak{B}(k)$, and
$\#\mathfrak{B}(k)\le 2$. This part of the proof carries the renormalisation step with domains
over to the re-filed renormalisation run of Definition~\ref{9.10:def-refiled-run}.

We first recall the shape of this step in the proof of the correlation theorem. It consists of
the following operations.
\begin{itemize}
\item The large-field operation $\mathbf{T}$ is applied.
\item The source-free decompositions of the effective action are made, together with the flat
decomposition attached to $\mathbf{T}$.
\item The gauge is fixed with the coefficient $a_k$, by the gauge-fixing lemma.
\item At the site of $\mathbf{T}$ in the table of units of the correlation theorem, the activity
lemma, the vertex lemma and the two summation lemmas, read with that table, produce the following
outputs, in Ba{\l}aban's enlarged spaces and at the decay rate $\mathsf{r}(\delta)\kappa$ of the
correlation theorem, which satisfies $\mathsf{r}(\delta)\kappa\ge(1+4\beta_s)\kappa$ with the
exponent $\beta_s$ fixed there:
\begin{itemize}
\item the regular family $E^{(k+1)}$;
\item the remnant part $R''^{(k+1)}$;
\item the $\mathbf{B}$-type terms $B^{(k+1)}$, stored together with their datum $A^{\flat}$ under
the corollary of the correlation theorem on that datum.
\end{itemize}
\item By \cite{B-CMP119}, the terms with $X\subset\Lambda_{k+1}$ coincide with the whole-lattice
ones. They are therefore the outputs of the frozen step of that proof.
\item The vacuum values are carried into the vacuum pieces $e^{(k+1)}$.
\item Ba{\l}aban's subtraction of $\beta_{k+1}(g_k)A(\varphi_{k+1},U_{k+1})$ \cite{B-CMP119} is
performed with the coupling increment $b_{k+1}(\zeta)$ in place of $\beta_{k+1}(g_k)$. After it
the weight of the action becomes $x_{k+1}(\cdot)-\zeta_{k+1}$.
\end{itemize}
We go through these operations one by one for the re-filed renormalisation run. Only the filing,
item~(vi) below, has no counterpart in the proof of the correlation theorem.

(i) \emph{The large-field operation, the decompositions, the labels and the gauge fixing.} The
operation $\mathbf{T}$, the source-free decompositions, the flat decomposition attached to
$\mathbf{T}$ and the labels do not involve the source.
They are therefore the same objects as in the proof of the correlation theorem.

The gauge fixing is carried out with the coefficient
\begin{equation*}
a_k(y)=\hat x_k-\hat\zeta_k(p_y),
\end{equation*}
as prescribed by \eqref{9.10:eq-refiled-run-b}. The gauge-fixing lemma holds for every
configuration $U$ and every complex coefficient $a$ with $|\operatorname{Im}a|\,\varepsilon_k^2\le 1$.
It is read through the last clause of the second list of smallness conditions of the correlation
theorem at
$|\hat\zeta_k|\le\frac83 r$. That reading is available here, so the gauge fixing goes through
without change.

(ii) \emph{The units at the site of $\mathbf{T}$, zone $k$.} We read the table of units of the
correlation theorem row by row.

\emph{The hatted first row.} It is the first row with the family $E^{(j)}$ replaced by
$\hat E^{(j)}$. The groups $\hat{\mathfrak{i}}^{(j)}_y$ of $\hat E^{(j)}$ include the groups of
the re-filed remnants. The pointed-family lemma of the correlation theorem is generic in the
family: its hypothesis on the family is only that it be a pointed $E$-family at scale $j$ on a
given sequence of radii, of norm at most a given number $\mathfrak{n}$. Part~(b) of the re-filing
lemma (the one among the re-filing lemmas whose parts~(a)--(d) concern the families
$\mathsf{RF}[R]$) supplies exactly this hypothesis for $\mathsf{RF}[R^{(i)}]$, with scale, radii
and norm bound
\begin{equation*}
j=i'',\qquad \text{radii }\hat\alpha_{i''-1},\qquad
\mathfrak{n}=C_{RF}\,w_{i''-1,i}\,g_i^{\kappa_0}.
\end{equation*}
This application, together with its bookkeeping of ratios, is the content of the corollary of the
re-filing lemma.

\emph{The hatted second row.} It is the second row with the family $\hat D^{(j)}$ of
\eqref{9.10:eq-refiled-run-a}, which contains the differences
\begin{equation*}
\mathsf{RF}\bigl[R^{(i)}(\cdot;w)\bigr]-\mathsf{RF}\bigl[R^{(i)}_{\hat w}\bigr],
\qquad \hat w=w(p_y),
\end{equation*}
between the family built at the local source and the family built at the source frozen at $p_y$.
The source-dependence lemma of the correlation theorem is applied at scale $i''$. Its hypotheses
are supplied by part~(c) of the re-filing lemma, which gives analyticity on
$\mathcal{W}_{i''}(Y)$ and the clause by which the family is frozen when the source is constant.
The subtraction of the value at the unit configuration is done by the subtraction lemma of the
correlation theorem.

\emph{The hatted third row.} It is the third row restricted to the levels $i\in\mathfrak{B}(k)$.
For each such $i$ the subtraction lemma is applied at the pair of levels $(n,j)=(k,i)$, with
weight $\rho^2_{k,i}\,g_i^{\kappa_0}$; recall that $\#\mathfrak{B}(k)\le2$. To these are added the
band units contained in Ba{\l}aban's source of $R''$-terms. For them part~(f) of the remnant
proposition of the correlation theorem applies without change.

\emph{The fourth, fifth and boundary rows.} These rows enter only at the second of the two sites
of the table, not at the site of $\mathbf{T}$, and they are read without change.
\begin{itemize}
\item The fifth row reads $\hat\zeta_k$ only through the bound $|\hat\zeta_k|\le\frac83 r$ and
through clause~(b) of the correlation theorem's statement on admissible sources.
\item The fourth row reads only the sup-norms of the stored $\mathbf{B}$-terms and of the stored
chains $\mathbb{C}$.
\item The near constituents of the boundary row are constituents of the regular family taken
singly. These are now constituents of $\hat E$, and the constituents coming from $\mathsf{RF}$
lie in the same class, with the same norm. Part~(ii) of the boundary lemma of the correlation
theorem consists of parts~(i) and~(iii) of the subtraction lemma, the trivial bound and
clause~(b) of the statement on admissible sources, all of which are generic in the family.
\end{itemize}

\emph{The last column of the table.} In the first two rows the norm bound enters at
$\mathfrak{n}\le E_0$. This holds a priori, by clause~(c) of the induction hypothesis of
Theorem~\ref{9.10:thm-refiled-induction} at level $k$. In the hatted third row the sum per cell
satisfies
\begin{equation}\label{9.10:eq-step-filing-1}
C\sum_{i\in\mathfrak{B}(k)}\rho^2_{k,i}\,g_i^{\kappa_0}\;\le\;C\sigma_k\;\le\;C\hat{\mathfrak{m}}_k .
\end{equation}
In the proof of the correlation theorem the corresponding sum is bounded, by the condition of
that theorem which bounds the weighted per-cell sum of the remnant units, as
\begin{equation}\label{9.10:eq-step-filing-2}
C\sum_{j\le k_1}\rho^2_{k,j}\,g_j^{\kappa_0}\;\le\;C\mathfrak{m}_k ,
\end{equation}
where $k_1=k_1(k)$ is the level for which Ba{\l}aban's source of $R''$-terms consists of the
levels $i$ with $k_1<i\le k$. The right member of
\eqref{9.10:eq-step-filing-1} is not larger than the right member of
\eqref{9.10:eq-step-filing-2}, because $\hat{\mathfrak{m}}_k\le\mathfrak{m}_k$. Hence the
constants $S_I$, $S_\omega$, $S^{\sharp}_{\partial}$, $C_{\sharp}$ of the correlation theorem,
and every ceiling in either list of smallness conditions of the correlation theorem that reads
them, are met a fortiori.

(iii) \emph{The activity, summation and vertex lemmas.} The activity lemma is linear in the unit
and is stated in terms of the unit size $\mathfrak{N}(v)$. Its output is analytic in every
parameter in which the activity $f_v$ of the unit $v$ is analytic. The two summation lemmas act on
the activities only. The vertex lemma at complex arguments is applied as in the
correlation theorem. All of them read the units only through the rows listed in~(ii), so they
apply without change.

Their outputs are:
\begin{itemize}
\item the family $E^{(k+1)}$ of the re-filed renormalisation run;
\item the remnant part $R''^{(k+1)}$, carrying the band marks;
\item the terms $B^{(k+1)}$ with their datum $A^{\flat}$, stored under the same corollary of the
correlation theorem.
\end{itemize}
All of these lie in Ba{\l}aban's enlarged spaces, at the decay rate
$\mathsf{r}(\delta)\kappa\ge(1+4\beta_s)\kappa$, as in the proof of the correlation theorem.

(iv) \emph{Terms with $X\subset\Lambda_{k+1}$.} In \cite{B-CMP119} the terms with
$X\subset\Lambda_{k+1}$ coincide with the whole-lattice ones. This is a statement about the
production formula, given its inputs. The local inputs of the re-filed renormalisation run,
restricted to $\Lambda_{k+1}$, are its frozen inputs there, for the following reasons.
\begin{itemize}
\item For the regular part this is the clause of the corresponding display~(c) by which
the local family equals the frozen family when the source is constant on $X^{+}$.
\item For the part coming from $\mathsf{RF}$ it is the same clause, supplied by part~(c) of the
re-filing lemma.
\item The units with levels in $\mathfrak{B}(k)$ and in the band are the same families in both
readings.
\end{itemize}
Therefore these terms are the constituents $E^{(k+1)}_{\zeta}$ of the frozen family
$\hat E^{(k+1)}_{\zeta}$ produced in the part of this proof concerning the frozen integral and
the flow (part~\ref{9.10:prf-step-frozen-flow} of the proof).

(v) \emph{Vacuum values.} The vacuum values are carried into the vacuum pieces $e^{(k+1)}$, as in
the proof of the correlation theorem. The families $\mathsf{RF}[R^{(i)}]$ contribute no vacuum
value, by part~(a) of the re-filing lemma.

(vi) \emph{The filing at step $k$.} Let $i\in\mathfrak{B}(k)$, so that $i''=k+1$. Consider every
pair $(X,y)$ of level $i$ with
\begin{equation*}
y\in\Lambda_{i''}^{\circledR}\qquad\text{and}\qquad Y(X,y)\subset\Lambda_{i''}.
\end{equation*}
This is the rule that \cite{B-CMP119} applies to the terms $E^{(k+1)}(X,z)$ with
$z\in\Lambda^0_{k+1}$ and $X\subset\Lambda_{k+1}$. For each such pair we move the share
\begin{equation}\label{9.10:eq-step-filing-3}
\mathsf{sh}(X,y)\bigl[R^{(i)}(X,U_{k+1})-R^{(i)}(X,1)\bigr]
\end{equation}
from $\hat{\mathbf{R}}^w$ into the group $\mathsf{RF}[R^{(i)}](Y,\cdot,y)$ of $\hat E^{(k+1)}$.
At the same time we subtract and add its counterterm
\begin{equation}\label{9.10:eq-step-filing-4}
\beta\bigl[\mathsf{RF}R^{(i)}\bigr]\,A\bigl(h^{(k+1)}_y,U_{k+1}\bigr).
\end{equation}
By the counterterm identity of the re-filing lemmas, applied with the partition functions
$\varphi_{k+1}$, this operation is an identity in the exponent.

The unfiled shares are those with $y\notin\Lambda^{\circledR}_{k+1}$ or
$Y(X,y)\not\subset\Lambda_{k+1}$. They stay in $\hat{\mathbf{R}}^{w,\mathrm{unf}}_{k+1}$ of
\eqref{9.10:eq-refiled-run-c}, written exactly in the format of the remnant terms of the
correlation theorem. They are remnant terms of Ba{\l}aban's format. For the shares with
$y\notin\Lambda^{\circledR}_{k+1}$ this is the statement of \cite{B-CMP119} that the $R''$-terms
do not depend on $\Lambda_{k+1}$; for the shares with $Y(X,y)\not\subset\Lambda_{k+1}$ it is the
sentence of \cite{B-CMP119} that ``the remaining terms are included again into the boundary
term''.
This is part~(f) of the remnant proposition of the correlation theorem, together with clause~(d)
of its induction format.

(vii) \emph{The coupling subtraction.} Subtracting $\hat b_{k+1}A(\varphi_{k+1},U_{k+1})$, the
weight of the action becomes $\hat x_{k+1}(\cdot)-\hat\zeta_{k+1}$, by the definitions collected
in \eqref{9.10:eq-refiled-run-b}. This is an add-and-subtract.

No estimate enters this part of the proof beyond the following, used at the constants of the
correlation theorem:
\begin{itemize}
\item the pointed-family lemma;
\item the subtraction lemma;
\item the source-dependence lemma;
\item the activity lemma;
\item the two summation lemmas;
\item parts~(a)--(d) of the re-filing lemma.
\end{itemize}
\end{proof}

\begin{proof}[Proof of Theorem~\ref{9.10:thm-refiled-induction}, continued: the local and large-field steps]
\label{9.10:prf-step-local-large}
We continue the induction. Let $k<K$ and assume $(\hat J_k)$. The frozen step and the filing at
level $k+1$ are carried out in the preceding part of this proof. It remains to carry out the
local step of the operation
$\mathbf{T}$ and the large-field step at level $k+1$, and to conclude $(\hat J_{k+1})$.

\emph{The local step.} Let $k+1\le K'$ and $w\in\mathcal{W}_{k+1}$. The local step is the
local step of the proof of the correlation theorem, with the same operations, performed with
the hatted definitions \eqref{9.10:eq-refiled-run-a} and \eqref{9.10:eq-refiled-run-b} in
place of the unhatted ones. We go through the points at which its inputs enter, in order.

First, the flat summation lemma of that proof is applied for $w\in\mathcal{W}_k$ with
majorants that do not depend on $w$. For the terms coming from the re-filed families
$\mathsf{RF}[R^{(i)}]$ the majorants are independent of $w$ by clause~(c) of the re-filing
lemma, whose norm takes the supremum over $w$ inside.

Second, the local step uses $|\hat\zeta_k|\le\frac83 r$. This is the bound
$|\hat\zeta_k(p)|\le\frac43\sup|w|$ of clause~(d) of $(\hat J_k)$, for sources
$w\in\mathcal{W}_k$, whose values lie in $D(0,2r)$; hence
$|\hat\zeta_k(p)|\le\frac43\cdot 2r=\frac83 r$.

Third, the vertex $\mathfrak{S}_k$ of the local step is a function of the product of the coupling
$\hat g_k$ with the same two factors as the vertex of the correlation theorem's proof; it
vanishes when this product is $0$ and contains no factor
$\hat g_k^{-2}$ \cite{B-CMP119}. Its bounds are
therefore those of the correlation theorem's proof, without change.

Fourth, the outputs of the local step have three properties. They are analytic in $w$ on
$\mathcal{W}_{k+1}(X)$. They are local: a potential in $Y\subset X$ reads $\hat\zeta_k$
restricted to $Y$ and the old terms in $Y$, that is, it reads $w$ restricted to $X^{+}$; for
a re-filed group $\mathsf{RF}[R^{(i)}(\cdot;w)](Y,\cdot,y)$ this is clause~(c) of the
re-filing lemma, by which the group reads $w$ restricted to $Y^{+}$. They are frozen at
constant sources: when $w\equiv\hat w$ on $X^{+}$, the
output equals the frozen output at $\hat w$; for the part coming from the re-filed families
this is the last assertion of clause~(c) of the re-filing lemma.

Fifth, we subtract the counterterm
\begin{equation*}
\sum_y \hat b_{k+1}\bigl(w(p_y)\bigr)\,A\bigl(h^{(k+1)}_y\varphi_{k+1},U_{k+1}\bigr),
\qquad |w(p_y)|<2r,
\end{equation*}
and move it, together with $+A(\hat\zeta_k,U_{k+1})$, into the weight. The resulting weight is
the one of clause~(a) of the corresponding format at level $k+1$.

Sixth, the running-shift statement is used exactly once in the local run, in the bound
\begin{equation}\label{9.10:eq-step-local-large-1}
|\hat\zeta_{k+1}(p)|
\le |w(p)| + \sum_j \varphi_j(p)\sum_y h^{(j)}_y(p)\,d_{j-1}(2r)\,|w(p_y)|
\le \tfrac43 \sup|w| .
\end{equation}
The middle member bounds the definition \eqref{9.10:eq-refiled-run-b} of $\hat\zeta_{k+1}$
term by term by means of the Lipschitz bounds $|\hat b_j(\zeta)-\hat b_j(0)|\le
d_{j-1}(2r)|\zeta|$ on $D(0,2r)$: for $j\le k$ those of clause~(b) of $(\hat J_k)$, and for
$j=k+1$ the same bound for $\hat b_{k+1}$, obtained in the frozen step at level $k+1$. The last
inequality holds because the partition functions $\varphi_j$ and $h^{(j)}_y$ of the running
shift take values in $[0,1]$, with $\sum_y h^{(j)}_y(p)\le1$, and because
$\sum_j d_{j-1}(2r)\le\frac13$ by the running-shift statement. Apart from these properties of
the partition functions, \eqref{9.10:eq-step-local-large-1} uses only the Lipschitz bounds and
their summability.

Seventh, the bound $|\hat b_j|\le b_+ +\frac23 r$ is in force, since the frozen step has been
performed; here $b_+$ is the constant of the correlation theorem's proof that bounds the
coupling increments at zero source.

Eighth, the freezing defect $\hat D^{(k+1)}$ of \eqref{9.10:eq-refiled-run-a} is the local
output minus the frozen output. Its bound has two parts. For the part coming from the regular
families it follows from the a-priori bound $E_0$ by the freezing lemma of the correlation
theorem's proof, applied at each later level. For the part coming from the re-filed families
it follows from clause~(c) of the re-filing lemma together with the freezing lemma, applied at
the filing scale $i''$. In neither case does the bound involve the defects of earlier levels,
so no constant compounds from level to level. This completes the local step.

\emph{The large-field step.} We apply the large-field step proposition of the correlation
theorem's proof at level $k+1$, to the frozen run and to the local run, together with the
lemma on the moves of the large-field operation, with the definitions
\eqref{9.10:eq-refiled-run-c} and the definition of the re-filed run that follows them, and
with the healing described in Remark~\ref{9.10:rem-healing-last-lattice}.

Among its inputs, the large-field step proposition reads the following two. The first is the
format of the induction hypothesis after the operation $\mathbf{T}$ at level $k+1$; this is
now the hatted format. The second is the bound
$\sup|\zeta_j|\le\frac83 r$; this now holds for $\hat\zeta_j$: for $j\le k$ by clause~(d) of
$(\hat J_k)$, as in the local step above, and for $j=k+1$ by
\eqref{9.10:eq-step-local-large-1}, since the values of $w$ lie in $D(0,2r)$. The proposition
produces the following six
outputs, which we take in its order.

(a) The chains $\mathbb{C}^{(n+1)}_k$. The large-field step proposition distinguishes two kinds
of sites, written (P) and (L) there: the sites at which it produces and stores chains are of
kind (P), and the sites at which it bounds the potential $V(Y;w)$ are of kind (L). At a site
of kind (P) the chains are produced by the
expansion lemma of the correlation theorem's proof, its table of units and the summation
lemma, and they are stored with the datum $A^{\flat}$. Their holomorphy in $w$, under
majorants independent of $w$, is given by the holomorphy corollaries for stored chains of that
proof. The units are those of the table of units, hatted as in part~(ii) of the frozen step
above. The backgrounds of the chains at a site of kind (P) are the
minimisers, in Ba{\l}aban's construction, of their own level data; they do not depend on the
complex sources. Hence this output holds for the hatted run without change.

(b) The source-part identity for the Wilson term at a large-field site. In the correlation
theorem's proof this identity states that $A(\zeta_k,U''_k)$ equals $A(\zeta_k^{\mathrm{new}},
U_k)$, plus the further terms displayed in that identity, plus the Wilson remainder
$\mathfrak{c}_w$; here $\zeta_k^{\mathrm{new}}$ is defined by clause~(a) of the format with
the partition functions $\varphi_j$ reset to $1$ on $Z$. In the hatted run the identity reads
with $\hat\zeta_k$, with $\hat\zeta_k^{\mathrm{new}}$ (the same formula with $\hat b_j$ in
place of $b_j$) and with a Wilson remainder $\hat{\mathfrak{c}}_w$. The remainder
$\hat{\mathfrak{c}}_w$ is bounded by $\theta_W$ times the bounds of Ba{\l}aban's construction,
which do not depend on the complex sources, through clause~(c) of the Wilson-remainder lemma;
that clause reads only $\sup|\hat\zeta_k|\le\frac83 r$. Hence this output holds without change.

(c) The bounds of the type \cite[(1.68)--(1.69)]{B-CMP122a} for $V(Y;w)$ at a site of kind
(L), in the form depending on the complex source $w$ that the correlation theorem's proof
establishes. They are derived there from the expansion lemma; from rows $1$, $2$, $3$, $5$ and
$\partial$ of the table of units, which in the hatted run are rows $\hat1$, $\hat2$,
$\hat3\cup\mathrm{band}$, $5$ and $\hat\partial$, hatted as in part~(ii) of the frozen step;
from the members of the table of inputs of the large-field step, namely the six members of its
error bound, the value input, the volume input and the input for stored chain terms $\mathbf{C}$;
from the strip theorem and its corollary on values; and from the filing lemma for contained
domains.

Every one of these members is a statement about units of the rows of the table of units, about
stored $\mathbf{B}$-terms, chains $\mathbb{C}$ and stored chain terms $\mathbf{C}$, and about
labels, strips and values. The sizes of the units enter only through $\mathfrak{N}(v)$, through
$\mathfrak{N}^{\flat}(v)$, the analogue of $\mathfrak{N}(v)$ for the class $\flat$, or through
the value bounds. For rows $\hat1$, $\hat2$ and
$\hat\partial$ these are the bounds of the correlation theorem at family norm at most $E_0$,
by clause~(c) of $(\hat J_k)$. For row $\hat3$ they are the bounds of the correlation theorem
at $g_i^{\kappa_0}$, by clause~(d) of $(\hat J_k)$, and they are taken over a smaller index set.
Therefore each of the six ceiling conditions of the large-field step proposition, as listed
there, is met a fortiori.

In Ba{\l}aban's construction the first part of the potential $V(Y;w)$ bounded in
\cite[(1.68)--(1.69)]{B-CMP122a} is split at a site of kind (L) according to
$X\cap Z$, into far and near pieces, by clause~(c) of the large-field step proposition and
clause~(i) of the boundary lemma. In the hatted run this split becomes the same split of the
filed pairs according to $Y\cap Z$. The split only redistributes terms: whatever is sent to
the first part is re-added. The unfiled shares beside $Z$ are terms of type $R$ of the first
part, exactly as in Ba{\l}aban's construction.

(d)--(e) The exponentiation \cite[(1.90)--(1.99)]{B-CMP122a} and its two groups of terms.
The second group becomes $R'^{(k+1)}$, whose modulus is bounded by
$e^{-p_0(g_{k+1})}e^{-\kappa d}$; this bound is
\cite[(1.100)]{B-CMP122a}, and $d$ is the distance that appears there. By their construction
these terms are Euclidean covariant at a constant source $w\equiv\hat w$
\cite[(2.32)]{B-CMP122b}; no covariance is asserted for them when $w$ is not constant. The
first group
becomes $\mathbf{B}'^{(k+1)}$, in the
relative form of clause~(d)$'$($\alpha$) of the corresponding format.
These are statements about the output of the exponentiation, given the bounds of~(c) on
$V(Y)$ and the modulus bound \eqref{9.10:eq-step-local-large-2} established below. Hence they
hold without change.

(f) The term $R''^{(k+1)}$ is the marked sub-sum of the band units of levels $j$ with
$k_1(k+1)<j\le k$, where $k_1(k+1)$ is the lower end of the band at level $k+1$ in the
correlation theorem's proof. This also holds without change, because the band is the source
of $R''$ in Ba{\l}aban's construction and is retained by the band condition of the hatted run.
Together with (d)--(e) this gives the remnant bound at level $k+1$:
\begin{equation}\label{9.10:eq-step-local-large-3}
\|R^{(k+1)}_{\zeta}\|^{\flat}\le g_{k+1}^{\kappa_0}.
\end{equation}

\emph{The moves of the large-field operation.} Each move in the lemma on the moves of the
large-field operation is an identity at fixed $w$ and is linear in the integrand. It therefore
applies to the hatted run without change. The gauge-fixing move requires in addition
\begin{equation*}
|\operatorname{Im} a|\,\varepsilon_k^2\le 1
\qquad\text{at}\qquad
a=\hat x_k-\hat\zeta_k ,
\end{equation*}
and this is the last clause of the strengthened small-field condition; here $\varepsilon_k$ is
the level-$k$ parameter with which the gauge-fixing move is stated in the correlation
theorem's proof.

\emph{The modulus bound.} The proof of the modulus bound for the large-field operation lists
the ingredients of \cite[(1.71)]{B-CMP122a} that depend on the inputs. They are the form,
which does not depend on the complex sources; the raw term
$-A\bigl((x''_k-\zeta_k)\hat\chi_1,U''_{h+2}\bigr)$ of the modulus bound's proof, in which
$\hat\chi_1$ and $U''_{h+2}$ are as in \cite[(1.71)]{B-CMP122a}; the potential
$V(X^{\sim};w)$; and
constants. The corresponding moduli are as follows.
\begin{enumerate}
\item[(1)] $|e^{V}|\le e^{\sup|V|}$, by the bound \cite[(1.69)]{B-CMP122a} in its form
depending on $w$.
\item[(2)] The form does not depend on the complex sources.
\item[(3)] The source shift in the raw Wilson term is at most $\theta_W x$, with $x$ as in the
proof of the modulus bound. This uses only $|\hat\zeta_k|\le\frac83 r$.
\item[(4)] The vertex contributes a factor $e^{\mathfrak{S}_j}$ per cell.
\item[(5)] The constants are $O(1)|Z_j|$, together with the per-zone-cell constants
$\mathsf{K}_{\mathrm{cell}}(k,j)$ with the frozen piles $\mathsf{P}(\mathfrak{a})$.
\end{enumerate}
Of these, (1) and (5) read the bounds of~(c), which hold in the hatted run a fortiori. The
moduli (2)--(4) read nothing beyond what was established above. Hence, for every integrand $F$,
\begin{equation}\label{9.10:eq-step-local-large-2}
\bigl|\hat{\mathbf{T}}'^{\,w}_k(X)F\bigr|\le e\,\mathfrak{M}_k(X)\,\sup|F| ,
\end{equation}
with the same right member $\mathfrak{M}_k(X)$ as in Ba{\l}aban's construction; this member
does not depend on the inputs.

\emph{Healing.} The healing at a large-field site of the current level is the add-and-subtract
of the hatted families of Remark~\ref{9.10:rem-healing-last-lattice}, which follows the
healing in Ba{\l}aban's construction \cite{B-CMP122b}.

\emph{Orbit data and arrested left-overs.} Every new term of type $\mathbf{B}$ is entered into
the format together with its orbit datum, as clause~(d)$'$($\gamma$) of the corresponding format requires. This is provided by the complex corollaries of
the correlation theorem's proof for the frozen and local steps, for the intermediate stages
and for the sites of kind (L). These are statements about terms of type $\mathbf{B}$ and about
the outputs of the orbit-datum construction, so they apply without change. Clause~(a) of the
large-field step proposition and clause~(d)$'$($\beta$) of the format carry unchanged the
arrested left-over terms, namely the terms arrested at earlier levels by the arrest theorem.

\emph{Conclusion.} The local step gives clause~(a) of the format at level $k+1$, the
analyticity, locality and freezing of the new terms, the bound
\eqref{9.10:eq-step-local-large-1} and the bound on $\hat D^{(k+1)}$. The large-field step gives
the chains, the Wilson-term identity, the bounds on $V(Y;w)$, the groups $R'^{(k+1)}$ and
$\mathbf{B}'^{(k+1)}$, the remnant bound \eqref{9.10:eq-step-local-large-3} and the modulus
bound \eqref{9.10:eq-step-local-large-2}. The healing, the orbit data and the arrested
left-over terms complete the format. Together with the frozen step and the filing, these give
$(\hat J_{k+1})$. This completes the induction step.
\end{proof}

\begin{proof}[Proof of Theorem~\ref{9.10:thm-refiled-induction}, continued: the comparison, running-shift and sourced-format clauses]
\label{9.10:prf-induction-conclusion}
By the induction step proved above (the local and large-field steps of the proof of
Theorem~\ref{9.10:thm-refiled-induction}), the hatted induction
hypothesis $(\hat J_k)$ holds for every $k\le K$, in particular for $k=K$. The comparison clause
for the pairs $(\hat x_i,\hat x_j)$, the hatted running-shift statement and the hatted sourced-format clause are read off from
$(\hat J_K)$:
\begin{itemize}
\item the comparison clause from clause~(a);
\item the analyticity and the reality of $\hat b_{k+1}$ and $\hat\zeta_{k+1}$, the bound
$|\hat\zeta_k-\zeta|\le\tfrac13|\zeta|$ and the bounds on the parts of the hatted regular family
from clauses~(b) and~(c);
\item the corresponding format of the hatted densities, together with
$|\hat\zeta_k(p)|\le\tfrac43\sup|w|$, from clause~(d).
\end{itemize}
The further assertions of these clauses are not among clauses~(a)--(d) as listed in
Theorem~\ref{9.10:thm-refiled-induction}; they are established, at each level $k+1\le K$, in the
induction step above together with $(\hat J_{k+1})$, and we collect them here: the decomposition
\begin{equation*}
\hat b_{k+1}\ =\ \beta^0_{k+1}+f_{k+1}(\hat\zeta_0,\dots,\hat\zeta_k)+\beta[\hat Q^{(k+1)}_\zeta],
\end{equation*}
with $f_{k+1}$ having the same memory property as in the running-shift statement of the
correlation theorem, and the value bound
\begin{equation*}
\sup_{D(0,2r)}\bigl|\beta[\hat Q^{(k+1)}_\zeta]\bigr|\ \le\ \hat v_k\ \le\ v_k ,
\end{equation*}
with $v_k$ the value constant of the running-shift statement of the correlation theorem and
$\hat v_k$ its counterpart for the hatted run;
the bound
\begin{equation*}
|\hat b_{k+1}(\zeta)-\hat b_{k+1}(0)|\ \le\ d_k(\varrho)\,|\zeta| ,\qquad \zeta\in D(0,\varrho) ;
\end{equation*}
and, for the hatted densities $\hat\rho^w_k$, the analyticity in $w$ of their terms, each of which,
localised in a region $X$, depends on $w$ only through the restriction of $w$ to the enlargement
$X^+$ of $X$ used in the sourced format of the correlation theorem, and the conservation
$\int\hat\rho^w_k\,dV_k=\int\rho^w_0\,dU$.
The remaining quantitative assertions are obtained as follows.

The first-passage level $K$ is finite by clause~(a). The first-passage inequality
$\hat x_{K+1}\le g^{-2}$ gives $\hat x_K\le g^{-2}+b_+$, with the constant $b_+$ of the comparison
clause of the correlation theorem. The bound $\sum_{k<K}d_k(\varrho)\le\tfrac13$ is the
summation statement of the correlation theorem; its hypothesis, the spacing
\begin{equation*}
\hat x_k\ \ge\ \hat x_K+\lambda(K-k)-2c_2 ,\qquad k<K ,
\end{equation*}
holds for the inverse couplings of the hatted run by clause~(a) of $(\hat J_K)$; here
$\lambda$ is the rate of the comparison clause of the correlation theorem.

For $k<K$ put $h_k:=\hat b_{k+1}-\hat b_{k+1}(0)$. This function is analytic on $D(0,2r)$ and
vanishes at $0$. Since $k+1\le K$, clause~(b) of $(\hat J_K)$ gives, for $\zeta\in D(0,2r)$,
\begin{equation*}
|h_k(\zeta)|\ \le\ d_k(2r)\,|\zeta|\ \le\ 2r\,d_k(2r).
\end{equation*}
Let $\zeta\in\bar D(0,r)$ and $0<r'<r$. The circle of radius $r'$ about $\zeta$ lies in
$D(0,2r)$, so the Cauchy estimate gives
\begin{equation*}
|h_k'(\zeta)|\ \le\ \frac{2r\,d_k(2r)}{r'} .
\end{equation*}
Letting $r'\to r$ yields $|h_k'(\zeta)|\le 2d_k(2r)$ on $\bar D(0,r)$. The closed disc is
convex, and $\hat b_{k+1}$ and $h_k$ differ by a constant. Hence
$\operatorname{Lip}_{\bar D(0,r)}\hat b_{k+1}\le 2d_k(2r)$. Summing over $k$, we obtain
\begin{equation}\label{9.10:eq-induction-conclusion-1}
\sum_{k<K}\operatorname{Lip}_{\bar D(0,r)}\hat b_{k+1}\ \le\ 2\sum_{k<K}d_k(2r)\ \le\ \tfrac23 ,
\end{equation}
where the second inequality is the summation bound of the preceding paragraph with the radius
$2r$ in the place of $\varrho$.

Finally, $\hat{\mathsf{N}}_K:=\hat\zeta_K'(0)$. The identity
\begin{equation*}
\hat{\mathsf{N}}_K-1\ =\ \sum_{k<K}h_k'(0)
\end{equation*}
is obtained for the hatted run by the argument that proves the same identity in the proof of the
running-shift statement of the correlation theorem; that argument applies verbatim.
Each $h_k$ is real on the real axis, because $\hat b_{k+1}$ is. Hence each $h_k'(0)$ is real, and
so is $\hat{\mathsf{N}}_K$. The function $\zeta\mapsto\hat\zeta_K(\zeta)-\zeta$ is analytic on
$D(0,2r)$ and, by clause~(b) of $(\hat J_K)$, bounded in modulus by $\tfrac13|\zeta|$. Hence it
vanishes at $0$, and its derivative at $0$, which is $\hat{\mathsf{N}}_K-1$, is the limit of
$(\hat\zeta_K(\zeta)-\zeta)/\zeta$ as $\zeta\to0$. Therefore
\begin{equation*}
|\hat{\mathsf{N}}_K-1|\ \le\ \tfrac13 ,
\end{equation*}
that is, $\hat{\mathsf{N}}_K\in[\tfrac23,\tfrac43]$.
\end{proof}

\begin{proof}[Proof of Theorem~\ref{9.10:thm-refiled-induction}, continued: the hatted torus bound]
The proof of the torus bound in the correlation theorem uses four inputs. We read each of them for
the hatted run, and each holds there with the same content and the same constants.

First, the vacuum-piece bound. The vacuum pieces $e^{(j)}(X;z)$ are built from the items of the
corresponding format that give the vacuum values:
\begin{itemize}
\item those of $\hat E^{(j)}(X,\cdot,y;w)$, which are those of the $E$-families, since the
re-filed families $\mathsf{RF}[R]$ have no vacuum values;
\item those of $R^{(j)}(X,1;w)$;
\item $\log\mathfrak{z}$.
\end{itemize}
These items occupy the same slots of the format and obey the same bounds as in the correlation
theorem. The bound
$|e^{(j)}(X;z)|\le C_e\,e^{-(1-\delta)\kappa d_j(X)}$, with the constant $C_e$ and the loss
$\delta$ in the exponent of the vacuum-piece bound in the proof of the torus bound of the
correlation theorem, therefore holds unchanged; here $d_j(X)$ is the quantity $d_j(X)$ of the
region $X$ fixed in the glossary, not one of the coefficients $d_k(\cdot)$ of the running shift.

Second, the real-part bound in the presence of the source $w$. Write $\hat A^w_{K'}$ for the
hatted effective action with source $w$ at the stop level $K'$, and let $\Gamma_j$ and the per-cell
constants $\mathsf{K}_{\mathrm{cell}}(K',j)$ be as in the corresponding bound of the proof of the
torus bound of the correlation theorem. Then
\begin{align*}
\operatorname{Re}\hat A^{w}_{K'}\ \le\ &-(1-\theta_W)\,A\bigl(\hat x_{K'}(\cdot),U_{K'}\bigr)
+\text{logarithmic terms}\\
&+\sum_j\mathsf{K}_{\mathrm{cell}}(K',j)\,|\Gamma_j|+\sum\mathsf{P}(\mathfrak{a}).
\end{align*}
This is the form with source $w$ of \cite[(2.49)]{B-CMP119}. It holds for the following reasons:
\begin{itemize}
\item the last column of the table of terms in the proof of the correlation theorem; for the hatted
terms the bound holds a fortiori;
\item the inequalities \cite[(2.47)--(2.48)]{B-CMP119};
\item clause~(c) of the Wilson-remainder bound of the correlation theorem;
\item the bound with right member $\mathfrak{M}_k(X)$, used for arbitrary regions;
\item the bound on the gas of live large-field components in the flat form fixed for this
chapter, applied with the hatted large-field operation $\hat{\mathbf{T}}'$; it holds with the same
content and the same constants.
\end{itemize}

Third, the floor $\tilde Z(0)\ge e^{-C_{\mathrm{low}}N^4}$, with $N$ as in the floor of the torus
bound of the correlation theorem. Its constant $C_{\mathrm{low}}$ depends on the run only through
the inequality
\begin{equation*}
\hat x_{K'}\ \le\ g^{-2}+b_++3\lambda m+2c_2 ,
\end{equation*}
which holds for the hatted run by the comparison clause established above.

Fourth, the two derivative estimates of that proof. They rest on the Schwarz lemma, the Cauchy
estimates, the anchoring at zero source and \cite[(1.26)]{B-CMP116}, and hold unchanged.

The display of the torus bound therefore holds for the hatted run with the same constants
$C_{UV}(\mathsf{d})$, $\mathsf{Q}$, $E_+^{\sharp}$ and $C_{\mathrm{low}}$.
\end{proof}

\begin{corollary}[Uniqueness inputs on the prescribed scaffold]\label{9.10:cor-scaffold-inputs}
Let $L\ge L_0$ be odd, and let $\mathfrak{p}$, $\gamma_J$ and the ceilings on $g$ be as in
Theorem~\ref{9.10:thm-refiled-induction}; let the hatted run be the run defined there, let
$X=g^{-2}$, and let $\Theta_X$ be the prescribed scaffold at $X$. Assume:
\begin{itemize}
\item[(a)] the proposition on the dirty parts $\hat Q^{(j)}$ of the hatted regular families
holds at every level $j\le K$;
\item[(b)] the induction theorem for the re-filed run, Theorem~\ref{9.10:thm-refiled-induction},
holds, so that the hatted induction hypothesis holds for every $k\le K$.
\end{itemize}
Then on $\Theta_X$ the hatted run satisfies the input conditions \textup{(J1)}--\textup{(J5)} of
the uniqueness theorem:
\begin{itemize}
\item[(J1)] its data lie in unit balls, are holomorphic, and do not depend on the torus;
\item[(J2)] the running-shift statement for the hatted run,
the corresponding display of Theorem~\ref{9.10:thm-refiled-induction}, holds together with its finite-range form, in which the
finite-range supremum $\mathfrak{S}$ enters, and with the coefficients $a_u$ of that
finite-range form the same as for the un-hatted run;
\item[(J3)] the scaffold theorem holds for the hatted running shift $\hat\zeta$;
\item[(J4)] with $A$ and $B$ the constants of condition \textup{(J4)} of the uniqueness
theorem (the letter $A$ here denotes that constant, not the Wilson action $A(U)$), the
partition function factorises as
\begin{equation*}
Z=e^{\mathcal{E}}\,\tilde Z,\qquad |\tilde Z|\le A,\qquad \tilde Z(0)\ge 1/B,
\end{equation*}
where the vacuum energy $\mathcal{E}$ differs from that of the un-hatted run only in that
each quantity entering it is replaced by its hatted counterpart;
\item[(J5)] for each $n$, the symmetry, covariance and reflection-positivity properties
required by the uniqueness theorem hold for the hatted run.
\end{itemize}
\end{corollary}

\begin{proof}
We take the five conditions in turn; throughout, the run is the hatted run of
Theorem~\ref{9.10:thm-refiled-induction}, read on the scaffold $\Theta_X$.

For \textup{(J1)}: condition \textup{(J1)} for the hatted run, namely that its data lie in unit
balls, are holomorphic and do not depend on the torus, is the data statement for the hatted
run, which is proved from hypotheses (a) and (b).

For \textup{(J2)}: the running-shift statement for the
hatted run, together with its finite-range form, in which the finite-range supremum
$\mathfrak{S}$ enters, is part of the conclusion of Theorem~\ref{9.10:thm-refiled-induction};
it therefore holds by hypothesis (b), and \textup{(J2)} uses it in its finite-range form, with
the supremum $\mathfrak{S}$. As part of the same conclusion of
Theorem~\ref{9.10:thm-refiled-induction}, the coefficients $a_u$ of the finite-range form are
the same as for the un-hatted run.

For \textup{(J3)}: the scaffold theorem is not re-proved here. Its proof, a shooting
argument, reads two statements and nothing else: the running-shift statement in the form of
\textup{(J2)}, which holds for $\hat\zeta$ by the previous paragraph, and a second statement of
this chapter, the only further input of the shooting argument, which for the hatted run is
likewise a consequence of hypotheses (a) and (b). Both inputs being available for the hatted
run, and the shooting argument reading nothing else, it applies without change to
$\hat\zeta$, and the scaffold theorem holds for $\hat\zeta$.

For \textup{(J4)}: the factorisation $Z=e^{\mathcal{E}}\tilde Z$ for the hatted run has the
same form as for the un-hatted run, $\mathcal{E}$ being obtained from the vacuum energy of the
un-hatted run by replacing each quantity entering it by its hatted counterpart. This
factorisation, together with the bounds $|\tilde Z|\le A$ and $\tilde Z(0)\ge 1/B$, is the
conclusion of the statement of this chapter on the partition function of the hatted run,
whose hypotheses are (a) and (b).

For \textup{(J5)}: symmetry, covariance and reflection positivity of the hatted run at each
$n$ are the content of the symmetry, covariance and reflection-positivity statement for the
hatted run, which is proved under hypotheses (a) and (b).

Thus \textup{(J1)}--\textup{(J5)} hold for the hatted run on $\Theta_X$ as a consequence of
hypotheses (a) and (b).
\end{proof}

\begin{definition}[Loads of a pair of re-filed lists]\label{9.10:def-pair-loads}
\emph{Convention on letters.} In this definition, and in the statements of this section that
precede the comparison of the two runs by their lengths, the letter $n$ denotes a zone level, in
the sense of the zone weights $w_{n,j}$ of the correlation theorem, and $k$ the step. From the
first statement of that comparison on, $n$ denotes the run length and $r,s$ denote depths.

\emph{Setting.} We work on one lattice and with one scaffold $\Theta_X$, the scaffold prescribed
at $X=g^{-2}$, with $\theta_k=X+c_0k$. Let
$\hat{\mathfrak{s}},\tilde{\hat{\mathfrak{s}}}\in\hat{\mathfrak{K}}_k$ be two hatted lists of
Definition~\ref{9.10:def-substitution-rule}, and put
$\delta:=\tilde{\hat{\mathfrak{s}}}-\hat{\mathfrak{s}}$, slot by slot. The slot-by-slot difference
norms $\hat{\mathsf{e}}_j$, $\mathsf{r}_j$, $\hat{\mathsf{z}}_j$, $\hat{\mathsf{b}}_j$ of $\delta$
are those of Definition~\ref{9.10:def-substitution-rule}; $\hat{\mathsf{z}}^{\sharp}_n$ denotes
the sharp form of the running-shift difference norm $\hat{\mathsf{z}}_n$ in which it enters the
loads of the one-lattice comparison. The uses of the lists are filed by the first of the four
conventions that accompany the one-lattice Lipschitz proposition for the dialled
production. The group of a re-filed remnant is \emph{pointed}, with anchor
$y\in\mathbb{T}^{(i'')}$, where $i''$ is the filing level of
\eqref{9.10:eq-filing-level-a}. An unfiled share is an unpointed unit of the remnant slot,
whose anchor is the cube of its large-field region in the lowest zone that region meets.

\emph{Loads.} The loads are those of the one-lattice comparison. They are read with the remaining
three conventions of the same proposition and with the substitution rule
\eqref{9.10:eq-substitution-rule-a}. Explicitly, for zone levels $n$ we put
\begin{equation}\label{9.10:eq-pair-loads-a}
\begin{aligned}
{}[\![\hat\delta]\!]_n &:= \sum_{j\le n}\varsigma_{1+n-j}\,\hat{\mathsf{e}}_j
 +\sum_{n-\mathsf{a}_j\le j\le n} w_{n,j}\,\mathsf{r}_j\\
&\qquad+\mathfrak{g}_n\Bigl[\hat{\mathsf{z}}^{\sharp}_n
 +\sum_{j\le n}L^{-(n-j)/2}\,\hat{\mathsf{b}}_j\Bigr],\\
\Sigma^q_k[\hat\delta] &:= \sum_{n\le k}(k-n+1)\,q^{(k-n-1)_+}\,[\![\hat\delta]\!]_n,\\
\hat C_L &:= C_L^{\sharp}\,(1+C_{RF}).
\end{aligned}
\end{equation}
Here $\varsigma_m$ is the geometric kernel, $w_{n,j}$ are the zone weights of the correlation
theorem, $\mathfrak{g}_n$ are the source-channel weights, $\mathsf{a}_j$ is as in
\eqref{9.10:eq-filing-level-a}, $q$ is the geometric rate, and $(\cdot)_+$ is the positive part;
$C_L^{\sharp}$, $C_L'$ and $\varkappa_k$ are the constants of the one-lattice Lipschitz
proposition for the dialled production, and $C_{RF}$ is the norm constant of the re-filing lemma.

\emph{Terminal load.} The terminal load $[\![\hat\delta]\!]^T_k$ is the first expression in
\eqref{9.10:eq-pair-loads-a} taken at $n=k$, with three changes. The $\mathsf{r}$-sum runs over
$j\in\mathfrak{B}(k)$, the band of \eqref{9.10:eq-filing-level-a}. The entry
$\hat{\mathsf{z}}_k$ replaces $\hat{\mathsf{z}}^{\sharp}_k$. The
$\hat{\mathsf{b}}$-sum is omitted.

\emph{Dial domain.} The dial $\lambda=(\lambda_{n,k'})$ is that of the one-lattice Lipschitz
proposition for the dialled production, and $E_1$ is the constant that bounds the dial domain of
that proposition. The dial domain built on the hatted loads is
\begin{equation*}
\mathbb{D}_{\hat\delta}:=\bigl\{\lambda:\ |\lambda_{n,k'}|\,[\![\hat\delta]\!]_n\le E_1
\ \text{for all } n,k'\bigr\}.
\end{equation*}

\emph{Constants.} The constant $C_L^{\sharp}$ of the one-lattice Lipschitz proposition for the
dialled production is replaced by
$\hat C_L$ of \eqref{9.10:eq-pair-loads-a}. This is the only new constant. The factor
$1+C_{RF}$ is the norm factor of part (b) of the re-filing lemma that bounds the norm of a
re-filed family. It acts on the re-filed constituents $\mathsf{RF}[R]$ of the first, second and
derivative entries of the hatted production, and on the filing term of the lemma on the hatted
constants. Correspondingly we put
\begin{equation}\label{9.10:eq-pair-loads-1}
\hat C_L':=\hat C_L\,E_1,\qquad \hat{\varkappa}_k:=K_{CE}\,\hat C_L\,\theta_k ,
\end{equation}
in place of $C_L'$ and $\varkappa_k$. The \emph{hatted smallness condition} is the smallness
condition of the one-lattice comparison on the Lipschitz coefficients, read with $\hat C_L'$ in
place of $C_L'$. It is a ceiling on $\gamma$, in the sense of
Notation~\ref{2.3:not-stages-first}, of the same kind as the original condition. Finally,
$\hat{\mathsf{c}}$ denotes the constant $\mathsf{c}$ of the dialled production read at
$\hat C_L'$.

\emph{Consequences.} Assume the following three statements, each together with the stated
property of its proof:
\begin{itemize}
\item[(a)] the one-lattice Lipschitz proposition for the dialled production holds: every object
derived in the dialled production is holomorphic on the dial domain and obeys the bounds of the
correlation theorem there, with the constants $C_L^{\sharp}$, $C_L'$, $\varkappa_k$ and
$\mathsf{c}$; and its proof treats the entries of the correlation theorem one by one, reading
each unit of the production only through the three structural properties
\textup{(f1)}--\textup{(f3)} of production units and through the unit's norm, as follows:
\begin{itemize}
\item entries (M), (I), (F) and (W) are stated for an arbitrary value of the parameter
$\mathfrak{n}$ on which they depend;
\item entry $(\partial)$ is homogeneous of degree one;
\item entry (A) uses only the subadditivity of the unit size $\mathfrak{N}$;
\item entry $(\Sigma)$ reads the unit only through its activity;
\item the weighted summation display used in that proof and the first three parts of
entry (R) are linear in the weights;
\item part (3) of the modulus bound, whose right member is the printed quantity
$\mathfrak{M}_k(\cdot)$, reads the unit only through $|\lambda\,\delta\hat\zeta|$;
\item entry (E) is read through the condition attached to it in the correlation theorem;
\item the fifth part of entry (R) reads the third and the fourth domain class, and then only
combinatorics free of the variable $z$;
\end{itemize}
\item[(b)] geometric forgetting of old regions holds, and its proof reads only the following
four ingredients, and no term of the action individually:
\begin{itemize}
\item the moduli, factor by factor, of the modulus bound under the sum over sequences;
\item Ba{\l}aban's replacement of that sum by a supremum in \cite{B-CMP122b};
\item the identity
\begin{equation*}
\Pi^{[\mathsf{c}-c_a]}=\Pi^{[\mathsf{c}]}\,
\exp\Bigl(-c_a\sum M^d\check R^{\,d+1}d'\Bigr),
\end{equation*}
in which $\Pi^{[\mathsf{c}]}$ denotes the product that the proof estimates at the
constant $\mathsf{c}$, $c_a$ is the amount by which that constant is lowered, $\check R$ is the
partition-cube radius, and the sum, with its factor $d'$, is the one of that proof;
\item Ba{\l}aban's induction at the constant $\mathsf{c}$;
\end{itemize}
\item[(c)] the corollary to geometric forgetting of old regions holds, and its proof has the
property stated in (b).
\end{itemize}
Assume also the hatted smallness condition. Then the following hold.
\begin{itemize}
\item[(i)] The bound of (a) holds for the hatted dialled production on $\mathbb{D}_{\hat\delta}$,
with $\hat C_L,\hat C_L',\hat\varkappa_k,\hat{\mathsf{c}}$ in place of
$C_L^{\sharp},C_L',\varkappa_k,\mathsf{c}$.
\item[(ii)] Statement (b) holds for the hatted production at the constant $\hat{\mathsf{c}}$,
under one further ceiling on $\gamma$ of the same kind as the smallness condition.
\item[(iii)] Statement (c) holds for the hatted production at the constant $\hat{\mathsf{c}}$,
under the further ceiling on $\gamma$ of (ii).
\end{itemize}
\end{definition}

\begin{proof}[Proof of the consequences stated in Definition~\ref{9.10:def-pair-loads}]
\emph{(i).} By hypothesis (a), the proof of the one-lattice Lipschitz proposition for the
dialled production treats the entries of the correlation theorem one by one and reads each unit
only through \textup{(f1)}--\textup{(f3)} and through the unit's norm, in the manner listed in
(a); each of these readings persists under the dial. By
Definition~\ref{9.10:def-substitution-rule}, the hatted units lie in the same classes. Their
norms are those of the class $\hat{\mathfrak{K}}_k$ of that definition. Each entry therefore
reads a hatted unit exactly as it reads an un-hatted one.

The new factor appears only on the re-filed constituents $\mathsf{RF}[R]$ and on the filing term
named in the definition; on both, part (b) of the re-filing lemma supplies the norm
factor $1+C_{RF}$. Hence $C_L^{\sharp}$ becomes $\hat C_L$ of \eqref{9.10:eq-pair-loads-a}, and
$C_L'$ and $\varkappa_k$ become $\hat C_L'$ and $\hat\varkappa_k$ of
\eqref{9.10:eq-pair-loads-1}. The constant $\mathsf{c}$, read at $\hat C_L'$, becomes
$\hat{\mathsf{c}}$. The dial domain is the one built on the hatted loads, namely
$\mathbb{D}_{\hat\delta}$. The smallness condition used in the proof of (a) is now the hatted
smallness condition, which is assumed. This gives (i).

\emph{(ii).} By hypothesis (b), the proof of geometric forgetting of old regions reads only the
four ingredients listed in (b), and none of them reads a term of the action individually. Each of
the four ingredients therefore applies verbatim to the hatted production at the constant
$\hat{\mathsf{c}}$, and (b) holds at $\hat{\mathsf{c}}$. Since $\hat{\mathsf{c}}\ge\mathsf{c}$,
this requires one further ceiling on $\gamma$, of the same kind as the smallness condition. This
gives (ii).

\emph{(iii).} By hypothesis (c), the proof of (c) has the same property as the proof of (b). The
argument of (ii) applies with (c) in place of (b), at the constant $\hat{\mathsf{c}}$ and under
the same further ceiling on $\gamma$.
\end{proof}

\begin{theorem}[Re-filed one-lattice Lipschitz step. Stated; proof incomplete at the step marked $(\ast)$]
\label{9.10:thm-refiled-lipschitz}
Throughout, $n$ denotes a zone level and $k$ the step. Let $\hat{\mathfrak{s}}$ and
$\tilde{\hat{\mathfrak{s}}}$ be two hatted lists of the class $\hat{\mathfrak{K}}_k$ of
Definition~\ref{9.10:def-substitution-rule}, on one and the same lattice and with one and the
same scaffold $\Theta_X$, and put $\delta:=\tilde{\hat{\mathfrak{s}}}-\hat{\mathfrak{s}}$.
The slot differences $\hat{\mathsf{e}}_j$, $\hat{\mathsf{q}}_j$, $\mathsf{r}_j$, $\mathsf{p}_j$,
$\mathsf{p}^{\flat}_j$, $\hat{\mathsf{z}}_j$, $\hat{\mathsf{b}}_j$ are those of
Definition~\ref{9.10:def-substitution-rule}; the loads $[\![\hat\delta]\!]_n$,
$[\![\hat\delta]\!]^T_k$, $\Sigma^q_k[\hat\delta]$, the dial domain $\mathbb{D}_{\hat\delta}$ and
the constants $\hat C_L$, $\hat C_L'$ and $\hat\varkappa_k=K_{CE}\hat C_L\theta_k$ are those of
Definition~\ref{9.10:def-pair-loads} and \eqref{9.10:eq-pair-loads-a}; the set $\mathfrak{B}(k)$ of
levels filed into the output of step $k$, the filing levels $i''$ and the numbers
$\mathsf{a}_i$ are those of \eqref{9.10:eq-filing-level-a}. Let $K_2$ be the constant of the
one-lattice Lipschitz step, evaluated at $\hat C_L$. Let $\hat\gamma_L$ be the coupling ceiling
$\gamma_L$ of the one-lattice Lipschitz step, lowered so that the smallness condition on the
coupling of that step holds with $\hat C_L'$ in place of $C_L'$, and so that the coupling
ceilings of Definition~\ref{9.10:def-substitution-rule} and of the lemma in hypothesis~(iv) below
hold. Let $g\le\hat\gamma_L$.
Assume:
\begin{itemize}
\item[(i)] the holomorphy proposition for the dialled production, applied to the hatted dialled
production with its constants $\mathsf{c}$ replaced by $\hat{\mathsf{c}}$ (the constants
$\mathsf{c}$ evaluated at $\hat C_L'$): every object derived in the hatted dialled production is
holomorphic on $\mathbb{D}_{\hat\delta}$ and obeys there the bounds of the correlation theorem;
\item[(ii)] geometric forgetting of old regions, at the constant $\hat{\mathsf{c}}$;
\item[(iii)] the hypothesis of the one-lattice Lipschitz step that stands among its hypotheses
beside the holomorphy proposition of \textup{(i)} and the theorem clause of \textup{(iii$'$)}, as
stated there;
\item[(iii$'$)] the clause (c) of a theorem that the one-lattice Lipschitz step takes among its
hypotheses, as stated there;
\item[(iv)] the frozen-output lemma for the hatted step: the bound of $\hat{\mathsf{e}}_{k+1}$ by
the marked activity of the rows of groups and of remnant units together with the filing term
$C_{RF}\sum_{i\in\mathfrak{B}(k)}w_{k,i}\mathsf{r}_i$.
\end{itemize}
The levels $k_1<j<k$ form the band of clause (c) of the one-lattice Lipschitz step; re-filing
leaves this band unchanged. The band remnant term $R''^{(k)}$, the remnant $R'^{(k)}$ produced at
step $k$ and the primed large-field boundary term $\hat{\mathbf{B}}'^{(k)}$ are the objects of
clauses (c) and (e) of the one-lattice Lipschitz step, formed for the hatted step; the number
$\beta_0$, the constant $C_P$, the labels $\omega$ of clause (f) of the one-lattice Lipschitz
step, their index $h(\omega)$, the sets $\Gamma_n$ and the tube of configurations
$(\mathbf{U},\mathbf{J})$ are those of clauses (c), (d) and (f) of the one-lattice Lipschitz step,
and the constants $a_k$, $c_k$, $C_t$ and the kernel $\varpi$ of part $(\hat{\mathrm g})$ are those
of its clause (g) (these $a_k$ are coefficients of that clause, not the labels $a_n$ of the
uniqueness theorem).
Then parts $(\hat{\mathrm a})$--$(\hat{\mathrm e})$ and $(\hat{\mathrm g})$ below hold, and clause
(f) of the one-lattice Lipschitz step holds with the quantities of part $(\hat{\mathrm f})$; the
first inequality of $(\hat{\mathrm a})$ is for the frozen and local output:
\begin{align}
&(\hat{\mathrm a}) && \hat{\mathsf{e}}_{k+1}\le\hat\varkappa_k[\![\hat\delta]\!]^T_k
+C_{RF}\sum_{i\in\mathfrak{B}(k)}w_{k,i}\mathsf{r}_i,\label{9.10:eq-refiled-lipschitz-a}\\
&&&\mathsf{p}^{\flat}_{k+1}\le\hat\varkappa_k\Big[\sum_{j\le k}\varsigma_{k+1-j}\mathsf{p}^{\flat}_j
+\mathfrak{g}_k\hat{\mathsf{z}}_k\Big],\notag\\
&(\hat{\mathrm b}) && \hat{\mathsf{q}}_{k+1}\le\hat\varkappa_k\Big[\sum_{j\le k}
\varsigma_{k+1-j}\hat{\mathsf{q}}_j+\sum_{i\in\mathfrak{B}(k)}w_{k,i}\mathsf{r}_i\Big]\notag\\
&&&\qquad+C_{RF}\sum_{i\in\mathfrak{B}(k)}w_{k,i}\mathsf{r}_i\notag\\
&&&\qquad+K_2\hat C_L^2(1+S_I)\theta_k^2\hat{\mathfrak{m}}_k\Big[\sum_{j\le k}
\varsigma_{k+1-j}\mathsf{p}^{\flat}_j+\mathfrak{g}_k\hat{\mathsf{z}}_k\Big],\notag\\
&(\hat{\mathrm c}) && \|\delta R''^{(k)}\|^{\flat}\le3\hat\varkappa_k\sum_{k_1<j<k}\mathsf{r}_j\notag\\
&&&\qquad+3K_2\hat C_L^2\theta_k^2(k-k_1)(1+\beta_0)^{\kappa_0}g_k^{\kappa_0}
[\![\hat\delta]\!]^T_k,\notag\\
&(\hat{\mathrm d}) && \big|\big(\hat{\mathbf{T}}'^{\tilde{\hat{\mathfrak{s}}}}_k(X)
-\hat{\mathbf{T}}'^{\hat{\mathfrak{s}}}_k(X)\big)F\big|\notag\\
&&&\qquad\le2eE_0^{-1}e^{-2(1+\beta_0)^{-1}p_0(g_k)}\sup|F|\cdot\Sigma^q_k[\hat\delta],\notag\\
&(\hat{\mathrm e}) && \|\delta R'^{(k)}\|^{\flat}\le4E_0^{-1}e^{-p_0(g_k)}\Sigma^q_k[\hat\delta],
\notag\\
&&&\text{the same for }\hat{\mathbf{B}}'^{(k)}\text{ with the constant }c_1\quad(\ast),\notag\\
&(\hat{\mathrm f}) && \mathsf{H}_\omega:=C_P\Big[[\![\hat\delta]\!]_k|\Lambda_k|
+[\![\hat\delta]\!]_\omega\sum_n|\Gamma_n\setminus\Lambda_k|\Big],\notag\\
&&&[\![\hat\delta]\!]_\omega:=\sum_{n\ge h(\omega)-1}(k-n+1)[\![\hat\delta]\!]_n,\notag\\
&(\hat{\mathrm g}) && \hat{\mathsf{b}}_{k+1}\le C_\beta\mathsf{p}_{k+1}
+C_\beta C_t\hat x_k\hat{\mathsf{q}}_{k+1},\notag\\
&&&\mathsf{p}_{k+1}\le a_k\sum_{j\le k}\varpi_{k-j}\mathsf{p}_j+c_k\hat{\mathsf{z}}_k,\notag\\
&&&\hat{\mathsf{z}}_{k+1}\le\hat{\mathsf{z}}_k+\hat{\mathsf{b}}_{k+1}.\notag
\end{align}
In part $(\hat{\mathrm d})$ the bound holds for every component $X$, every $(\mathbf{U},\mathbf{J})$
in the tube, every $w\in\mathcal{W}_k$ and every bounded $F$. In part $(\hat{\mathrm e})$ the
constant $c_1$ is the one of \cite[(1.97)]{B-CMP122a}. Part $(\hat{\mathrm f})$ (real sources,
relative form) states that clause (f) of the one-lattice Lipschitz step holds with its quantity
$\mathsf{H}_\omega$ replaced by the displayed one, the constant $C_P$ being evaluated at
$\hat C_L$. The third inequality of part $(\hat{\mathrm g})$ holds along runs.

None of the constants depends on $K$, $k$, $m$, $N$, the torus or the position. In every part the
remnant terms $\mathsf{r}$ occur only over a range of levels bounded by $\mathsf{a}$: in
$[\![\hat\delta]\!]_n$ only for $j\ge n-\mathsf{a}_j$, in parts $(\hat{\mathrm a})$ and
$(\hat{\mathrm b})$ only for $i\in\mathfrak{B}(k)$, and in part $(\hat{\mathrm c})$ only for the
levels $k_1<j<k$ of the band. No sum over all earlier remnants occurs: the re-filed images of the
older remnants are contained in $\hat{\mathsf{e}}_{j''}$ and enter through the kernel $\varsigma$.
\end{theorem}

Stated; the proof below is incomplete at the step marked $(\ast)$.

\begin{proof}
We treat the parts in order.
Hypotheses (iii) and (iii$'$) enter only through the proofs of the clauses of the one-lattice
Lipschitz step, which we follow part by part; hypotheses (i) and (ii) are used where indicated,
and hypothesis (iv) is the bound to which the first inequality of part $(\hat{\mathrm a})$ adds the
contributions of rows $\hat2$ and $5$, as recorded after its proof.

\emph{Part $(\hat{\mathrm a})$.} We follow the proof of clause (a) of the one-lattice Lipschitz
step at the site $(\mathrm{T})$ of the activity table of the correlation theorem, the site at
which the frozen output $E^{(k+1)}$ is formed: the domains are
$X\subset\Lambda_{k+1}$, the lattice is the
whole lattice, the zone is $k$, and the boundary row $\partial$ is inactive. Let $V[\cdot]$ be the
activity fed into the cluster expansion at the site $(\mathrm{T})$, as a function of the list. By
the linearity lemma of the correlation theorem, $V$ is linear in the list, so that
\begin{equation*}
V[\tilde{\hat{\mathfrak{s}}}]=V[\hat{\mathfrak{s}}]+V^m,\qquad V^m:=V[\delta].
\end{equation*}
Let $\mathsf{W}^m$ be the marked activity of $V^m$. It satisfies
$\mathsf{W}^m\le\theta_k\hat C_L[\![\hat\delta]\!]^T_k$, the contributions being bounded row by row
as follows, in the numbering of the rows of the activity table of the correlation theorem. Row
$\hat1$ consists of the groups of $\delta\hat E^{(j)}$, each with its own
coupling extraction $\beta[\delta\hat E^{(j)}]=\delta\hat b_j$; this is a difference of pointed
$E$-families, whose re-filed part is accounted for by parts (b) and (c) of the re-filing lemma,
and it is bounded by the group lemma of the correlation theorem
with its norm parameter equal to $\hat{\mathsf{e}}_j$. Row $\hat2$ consists of $\delta\hat D$, and
is bounded by the freezing-defect lemma of the correlation theorem composed with its unit lemma.
Row $\hat3$ consists of the $\delta R^{(i)}$, $i\in\mathfrak{B}(k)$, bounded by the unit lemma at
weight $w_{k,i}$. Row $5$ consists of $\delta\hat\zeta_k$, bounded by part (b) of the lemma of the
correlation theorem on the running shift. By the cluster-expansion corollary of the correlation
theorem, the outputs of the expansion with activities $V[\hat{\mathfrak{s}}]$ and
$V[\hat{\mathfrak{s}}]+V^m$ differ in norm by at most $K_{CE}\mathsf{W}^m$. This bounds the difference
of the frozen outputs:
\begin{equation*}
\|\delta E^{(k+1)}\|^{\flat}\le K_{CE}\theta_k\hat C_L[\![\hat\delta]\!]^T_k
=\hat\varkappa_k[\![\hat\delta]\!]^T_k .
\end{equation*}
The hatted output contains in addition the re-filed families $\mathsf{RF}[R^{(i)}]$,
$i\in\mathfrak{B}(k)$. By linearity of $\mathsf{RF}$ their difference is
$\sum_{i\in\mathfrak{B}(k)}\mathsf{RF}[\delta R^{(i)}]$, and by the norm bound of part (b) of the
re-filing lemma,
\begin{equation*}
\Big\|\sum_{i\in\mathfrak{B}(k)}\mathsf{RF}[\delta R^{(i)}]\Big\|^{\flat}
\le C_{RF}\sum_{i\in\mathfrak{B}(k)}w_{k,i}\mathsf{r}_i .
\end{equation*}
Adding the two bounds gives the first inequality of $(\hat{\mathrm a})$. This estimate is the
frozen-output lemma of hypothesis (iv) with the contributions of rows $\hat2$ and $5$ added, in
the same way as clause (a) of the one-lattice Lipschitz step is the corresponding lemma for the
un-hatted step with rows $2$ and $5$ added. The pure part $P^{\flat(k+1)}$ is computed from the
pure inputs and the numbers $\hat\zeta_i$ only, so its estimate is unchanged, and the second
inequality of $(\hat{\mathrm a})$ follows as in clause (a) of the one-lattice Lipschitz step.

\emph{Part $(\hat{\mathrm b})$.} This is Lemma~\ref{9.10:lem-second-difference}, the second
difference with re-filed remnant marks, which rests on the identity
\eqref{9.10:eq-second-difference-a}.

\emph{Part $(\hat{\mathrm c})$.} The term $R''^{(k)}$ is the marked sub-sum of the band units in
part (f) of the remnant proposition of the correlation theorem. The marks $\delta R^{(j)}$ enter
at weight at most $3$; for the other slots we apply the second dial estimate of the correlation
theorem, with a parameter $\mu$ on those units. This is the proof of clause (c) of the one-lattice
Lipschitz step, applied without change, because re-filing does not change the band. The other
slots now also include the $\hat Q$-groups and the units of the levels in $\mathfrak{B}(k)$; all
of them are unmarked for $R''$, and their contributions lie inside $[\![\hat\delta]\!]^T_k$.

\emph{Part $(\hat{\mathrm d})$.} Let $\lambda=(\lambda_{n,k'})$ be the dial variables and put
$\mathsf{G}(\lambda):=\hat{\mathbf{T}}'^{\lambda}_k(X)F$ on $\mathbb{D}_{\hat\delta}$; by
hypothesis (i), $\mathsf{G}$ is holomorphic there. The variable $\lambda_{n,k'}$ occurs only in
those sequences in which zone $n$ is present inside $X^{\sim}$ at a level $k'\le k$. For re-filed
remnants this follows from the filing of their uses in Definition~\ref{9.10:def-pair-loads}: a
re-filed remnant's group is pointed, anchored at a point of the lattice of its filing level $i''$.
An unfiled share of level $i$ lies, by
\eqref{9.10:eq-unfiled-range-a}, in zones $n$ with $i\le n\le i''$, so that it carries
the variables $\lambda_{n,k'}$ with $n$ its zone. For $n<k$ this forces the presence inside
$X^{\sim}$ of a component created at a step at most $n+1$, hence of age at least $k-n-1$ at level
$k$; the sequences not containing
$\lambda_{n,k'}$ do not contribute to the dependence on $\lambda_{n,k'}$. Let $\Delta_{n,k'}$ be
the supremum over $\mathbb{D}_{\hat\delta}$ of the part of $\mathsf{G}$ formed by the sequences
containing $\lambda_{n,k'}$. By hypothesis (ii),
\begin{equation*}
\Delta_{n,k'}\le2e\cdot e^{-2(1+\beta_0)^{-1}p_0(g_k)}q^{(k-n-1)_+}\sup|F| .
\end{equation*}
We apply the third dial estimate of the correlation theorem in each variable $\lambda_{n,k'}$ at
radius $E_1/[\![\hat\delta]\!]_n$, and sum over the $k-n+1$ values of $k'$ and over $n\le k$; the
weights $(k-n+1)q^{(k-n-1)_+}[\![\hat\delta]\!]_n$ so obtained form $\Sigma^q_k[\hat\delta]$, and
part $(\hat{\mathrm d})$ follows. This is the proof of clause (d) of the one-lattice Lipschitz step,
applied without change on $\mathbb{D}_{\hat\delta}$.

\emph{Part $(\hat{\mathrm e})$.} We follow the proof of clause (e) of the one-lattice Lipschitz step
without change. By hypothesis (i), the activities of \cite[(1.91)]{B-CMP122a} obey
\cite[(1.97)]{B-CMP122a} with $2c_1$ in place of $c_1$ on $\mathbb{D}_{\hat\delta}$. The same holds
when, for an integer $a\ge0$, the sum over sequences of each component is split as
$\hat{\mathbf{T}}'_{<a}+\mu\hat{\mathbf{T}}'_{\ge a}$ with $|\mu|\le q^{-a}$, as in the proof of
clause (e) of the one-lattice Lipschitz step, because by hypothesis (ii)
\cite[(1.92)]{B-CMP122a} continues to hold. Hence
\begin{equation*}
|R'(X;\lambda,\mu)|\le\eta_X:=e^{-p_0(g_k)}e^{-\kappa d_k(X)} .
\end{equation*}
The first dial estimate of the correlation theorem, applied to the function $\mu\mapsto
R'(X;\lambda,\mu)$, holomorphic on the disc $|\mu|\le q^{-a}$ and bounded there by $\eta_X$,
gives
\begin{equation*}
|R'(\lambda,1)-R'(\lambda,0)|\le2\eta_Xq^a .
\end{equation*}
Fix a zone $n\le k$ and a level $k'\le k$, and take $a:=(k-n-1)_+$. By the analysis of part
$(\hat{\mathrm d})$, every sequence containing
$\lambda_{n,k'}$ has age at least $k-n-1$, so $R'(\cdot,0)$ does not contain $\lambda_{n,k'}$.
Therefore
\begin{equation*}
R'(\lambda,1)-\big(R'(\lambda,1)\big|_{\lambda_{n,k'}=0}\big)
=\big[R'(\lambda,1)-R'(\lambda,0)\big]
-\big[R'(\lambda,1)-R'(\lambda,0)\big]\big|_{\lambda_{n,k'}=0},
\end{equation*}
and each bracket is at most $2\eta_Xq^{(k-n-1)_+}$ in modulus, whence
\begin{equation*}
\sup_{\mathbb{D}_{\hat\delta}}\Big|R'(\lambda,1)-\big(R'(\lambda,1)\big|_{\lambda_{n,k'}=0}\big)\Big|
\le4\eta_Xq^{(k-n-1)_+} .
\end{equation*}
The third dial estimate, applied as in part $(\hat{\mathrm d})$, now gives the bound of
$\|\delta R'^{(k)}\|^{\flat}$ in part $(\hat{\mathrm e})$. The inputs enter $R'$ only through the
units at the sites $(\mathrm{P})$ and $(\mathrm{L})$ of the activity table of the correlation
theorem, namely the hatted rows $\hat1$,
$\hat2$, $\hat3$ together with the band and the unfiled shares, $4$, $5$ and $\hat\partial$; and
through the un-expanded exponent of $\hat{\mathbf{T}}'$ inside $X$, estimated factor by factor
through its moduli, and
the raw weights $\hat x_j-\hat\zeta_j$. All of these lie inside $[\![\hat\delta]\!]_n$, zone by
zone. Thus the newly created remnant $R'^{(k)}$ depends on the old remnants in two ways: unfiled
ones at
zone weight $w_{n,i}$ with $n-i\le\mathsf{a}_i$, the age of the zone being paid by $q^{k-n}$; and
filed ones through $\hat{\mathsf{e}}_j$ with the kernel $\varsigma$. Either way the dependence is
geometric.

$(\ast)$ The corresponding bound for $\hat{\mathbf{B}}'^{(k)}$, with the constant $c_1$ of
\cite[(1.97)]{B-CMP122a}, is not proved here: the argument just given for $R'^{(k)}$ is not carried
over to the large-field boundary term $\hat{\mathbf{B}}'^{(k)}$. This clause of part
$(\hat{\mathrm e})$ is the subject of the carry lemma.

\emph{Parts $(\hat{\mathrm f})$ and $(\hat{\mathrm g})$.} The proofs of clauses (f) and (g) of the
one-lattice Lipschitz step apply without change. The coupling extraction splits as
\begin{equation*}
\beta[\delta\hat E^{(k+1)}]=\beta[\delta\mathsf{P}^{(k+1)}]+\beta[\delta\hat Q^{(k+1)}],
\end{equation*}
because $P^{\flat}$ agrees with $\mathsf{P}$, and because the constituents $T'$ and $\Pi^{\flat}$ of
the proof of clause (g) of the one-lattice Lipschitz step are evaluated at $\hat x_k$, as in that
proof.
\end{proof}

\begin{theorem}[Kernel form of the re-filed Lipschitz step. Stated; proof incomplete at the step
marked $(\ast)$]\label{9.10:thm-lipschitz-kernel}
Let the coupling $g$, the level $k$, the two hatted lists
$\hat{\mathfrak{s}},\tilde{\hat{\mathfrak{s}}}\in\hat{\mathfrak{K}}_k$ on one lattice and one scaffold,
and their difference $\delta$ be as in Theorem~\ref{9.10:thm-refiled-lipschitz}, whose proof is
incomplete at one step. Assume the following statements:
\begin{enumerate}
\item[(i)] the one-lattice Lipschitz step, that is, its estimates (P), (b), (w), (Q), (R$''$),
(R$'$), ($\mathbf{B}$), (D), ($\mathbf{T}$) together with the bounds established in its proof;
\item[(ii)] the identification of the pure part of the re-filed run with the pure part of the
original run;
\item[(iii)] geometric forgetting of old regions.
\end{enumerate}
For $j\le k+1$ write $\delta^P_j$, $\delta^{\hat Q}_j$, $\delta^R_j$, $\delta^{\mathbf{B}}_j$,
$\delta^{\hat D}_j$ and $\delta^{\hat\zeta}_j$ for the norms, in the normalisation of the one-lattice
Lipschitz step, of the differences between the two lists of, respectively, the pure part
$P^{\flat(j)}$, the dirty part $\hat Q^{(j)}$, the remnant $R^{(j)}$, the boundary term of level $j$,
the freezing-defect part $\hat D^{(j)}$, and the running shift $\hat\zeta_j$. The kernels
$\mathfrak{w}^P_{k,j}$, $\mathfrak{w}^{\mathbf{B}}_{k,j}$, $\mathfrak{w}^D_{k,j}$, the weights
$\bar w_{i,j}$, the constants $K_R$, $C_{\mathbf{T}}$, $\varpi'$, the rate $\vartheta_*\in(0,1)$, the
term $\Upsilon$ and the level $k_1=k_1(k)$ are those of the one-lattice Lipschitz step; $w_{k,i}$ and
$\omega_{i,j}$ are the zone weights and $K_{CE}$, $K_{\times}$, $S_I$, $C_U$ the cluster-expansion and
room constants of the correlation theorem; $\theta_k$ is the scaffold at level $k$,
$\hat{\mathfrak{m}}_k$ the bound for the dirty part, $\mathsf{v}_k$ the vertex activity, $C_{RF}$ the
norm constant of the re-filing, $r$ the source radius, the frozen sources $\hat w$ ranging over the
disc $D(0,2r)$, and
$\mathfrak{B}(k)$, $\mathrm{band}(k)$ and the filing level $j''=j+\mathsf{a}_j$ are as in
\eqref{9.10:eq-filing-level-a}. Then the estimates of the one-lattice Lipschitz step hold for the
re-filed run of Definition~\ref{9.10:def-refiled-run} in the following form.
\begin{enumerate}
\item[] The estimates (P), (b) and (w) hold as stated in the one-lattice Lipschitz step.
\item[(Q)] The dirty part satisfies
\begin{equation}\label{9.10:eq-lipschitz-kernel-1}
\begin{split}
\delta^{\hat Q}_{k+1}\le{}& K_{CE}C_U\theta_k\Big[\sum_{j\le k}\varsigma_{k+1-j}\delta^{\hat Q}_j
+\sum_{i\in\mathfrak{B}(k)}w_{k,i}\delta^R_i\Big]
+C_{RF}\sum_{i\in\mathfrak{B}(k)}w_{k,i}\delta^R_i\\
&+K_{\times}C_U^2\theta_k^2(S_I+1)\hat{\mathfrak{m}}_k\sum_{j\le k}\varsigma_{k+1-j}\delta^P_j
+K_{\times}C_U\theta_k(S_I+1)\hat{\mathfrak{m}}_k\mathsf{v}_k\delta^{\hat\zeta}_k ,
\end{split}
\end{equation}
which is the estimate bounding $\hat{\mathsf{q}}_{k+1}$ in
Theorem~\ref{9.10:thm-refiled-lipschitz}, \eqref{9.10:eq-refiled-lipschitz-a}, written in the
constants of the one-lattice Lipschitz step.
\item[(R$''$)] The estimate (R$''$) holds with the sum $\sum_{j\le k_1}w_{k,j}\delta^R_j$ inside its
second bracket replaced by $\sum_{i\in\mathfrak{B}(k)}w_{k,i}\delta^R_i$.
\item[(R$'$)] With $R'^{(k+1)}$ and $\hat{\mathbf{B}}'^{(k+1)}$ the remnant and boundary terms of
level $k+1$ bounded in the estimate (R$'$),
\begin{equation}\label{9.10:eq-lipschitz-kernel-2}
\begin{split}
&\|\delta R'^{(k+1)}\|^{\flat}+\|\delta\hat{\mathbf{B}}'^{(k+1)}\|\\
&\quad\le e^{-p_0(g_k)}K_R\Big\{\sum_{j\le k}\Big[\mathfrak{w}^P_{k,j}\big(\delta^P_j
+\delta^{\hat Q}_j\big)+\hat{\mathfrak{w}}^R_{k,j}\delta^R_j\\
&\qquad+\mathfrak{w}^{\mathbf{B}}_{k,j}
\delta^{\mathbf{B}}_j+\mathfrak{w}^D_{k,j}\delta^{\hat D}_j
+\vartheta_*^{k-j}\delta^{\hat\zeta}_j/r\Big]+\Upsilon\Big\},
\end{split}
\end{equation}
with the hatted remnant kernel
\begin{equation}\label{9.10:eq-lipschitz-kernel-a}
\hat{\mathfrak{w}}^R_{k,j}:=\sum_{i=j}^{\min\{k,j''\}}\bar w_{i,j}\,\vartheta_*^{k-i}
\le(\mathsf{a}_j+1)\Big(\frac{4x_j}{x_k}+\varpi'\Big)\vartheta_*^{(k-j'')_+}.
\end{equation}
\item[($\mathbf{B}$)] The estimate ($\mathbf{B}$) holds with the sum
$\sum_{j\le k_1}w_{k,j}\delta^R_j$ replaced by $\sum_{i\in\mathfrak{B}(k)}w_{k,i}\delta^R_i$.
\item[(D)] The estimate (D) holds for the hatted regular family $\hat E$.
\item[($\mathbf{T}$)] The estimate ($\mathbf{T}$) holds with its exponent bound replaced, for a
large-field region $X$, by
\begin{equation}\label{9.10:eq-lipschitz-kernel-3}
\begin{split}
\hat\Xi_X:={}&C_{\mathbf{T}}\sum_{h_X\le i\le k+1}\Big[g_i^2\delta^{\hat\zeta}_i
+\sum_{j\le i}\Big(\varsigma_{i+1-j}\big(\delta^P_j+\delta^{\hat Q}_j\big)
+\mathbf{1}\{i\le j''\}\bar w_{i,j}\delta^R_j\\
&\qquad+2^{-(i-j)}\delta^{\mathbf{B}}_j+\omega_{i,j}\delta^{\hat D}_j\Big)\Big]
+C_{\mathbf{T}}\Upsilon ,
\end{split}
\end{equation}
where $h_X$ is the lowest level of the inner large-field operations acting on $X$.
\end{enumerate}
\end{theorem}

Stated; the proof below is incomplete at the step marked $(\ast)$. All other estimates of
Theorem~\ref{9.10:thm-lipschitz-kernel}, and the contribution of the first member
$\|\delta R'^{(k+1)}\|^{\flat}$ to \eqref{9.10:eq-lipschitz-kernel-2}, are proved below as an
implication from the statements (i)--(iii).

As a function of the level $i$ at which a remnant of level $j$ is read, the kernel
\eqref{9.10:eq-lipschitz-kernel-a} has range $\mathsf{a}_j$ in the level age and is geometric in
$k-j''$ beyond it; in the one-lattice Lipschitz step the remnant kernel is the same sum taken over
all $j\le i\le k$, without the restriction $i\le j''$.

\begin{proof}
We follow the proof of the one-lattice Lipschitz step estimate by estimate, for the re-filed run of
Definition~\ref{9.10:def-refiled-run}, whose action at level $k$ is \eqref{9.10:eq-refiled-run-c};
for each estimate we state what enters in place of the corresponding object of the original run.

\emph{(P).} By hypothesis (ii) the pure part of the re-filed run is that of the original run, so
(P) and its proof are unchanged.

\emph{(b).} The proof of (b) in the one-lattice Lipschitz step is repeated with the large-field
operation $\mathbf{T}'$, the quantity $\Pi^{\flat}$ entering that proof, and the first object of the
list of objects of the run, wherever they enter, taken to be those of the re-filed run of
Definition~\ref{9.10:def-refiled-run}; the estimate itself is unchanged.

\emph{(Q).} The estimate \eqref{9.10:eq-lipschitz-kernel-1} is obtained from the decomposition into
three braces of the second difference with re-filed remnant marks,
Lemma~\ref{9.10:lem-second-difference} (proved as an implication from the statements named there),
identity \eqref{9.10:eq-second-difference-a}, in which the loads of the marks are the dirty-part
loads $\hat{\mathsf{q}}$ together with the remnant loads $\mathsf{r}_i$ restricted to
$i\in\mathfrak{B}(k)$.

\emph{(R$''$).} The remnants are produced as in clause (f) of the proposition producing the
remnants in the correlation theorem, so the proof of (R$''$) in the one-lattice Lipschitz step
applies unchanged; the remnants read in its second bracket are now those of the levels in
$\mathfrak{B}(k)$, as they appear in the action \eqref{9.10:eq-refiled-run-c}, which replaces the
sum over $j\le k_1$ by the sum over $i\in\mathfrak{B}(k)$.

\emph{($\mathbf{B}$).} Near the boundary $\partial\Lambda_{k+1}$ of the region $\Lambda_{k+1}$, whose
complement is the large-field region of level $k+1$, the kernels are those of the proof of
($\mathbf{B}$) in the one-lattice Lipschitz step. The summation lemma of the correlation theorem
covers the pinned kernels of \cite[(3.47)]{B-CMP119} and the kernels of the boundary terms and of
the terms $R''$. The old boundary terms $\mathbf{B}^{(j)}$ enter as the terms of the summation of
boundary terms at finer zones. The marks are the differences $\delta\mathbf{B}$, the differences of
the groups of the hatted regular families $\hat E$, the differences of the remnants of the levels in
$\mathfrak{B}(k)$, and the vertex. With these marks the proof is unchanged, and the sum over
$j\le k_1$ becomes the sum over $i\in\mathfrak{B}(k)$.

\emph{(D), (w).} The proofs are unchanged; in (D) the freezing-defect parts are those of the hatted
regular family $\hat E$.

\emph{(R$'$) and ($\mathbf{T}$).} We proceed in three steps.

\emph{First step: the remnant terms of the estimate.} In the proof of (R$'$) in the one-lattice
Lipschitz step, remnant terms enter at two places, with the weights $w_{k,j}g^{\kappa_0}$ and
$\varpi'$ respectively, where $\kappa_0$ is the exponent of the remnant class of
Definition~\ref{9.10:def-refiled-run}; the terms at these two places are the remnant units of the
table of hatted units of the run, in its third row. That reading uses the following sentence of
\cite{B-CMP122b}:
\begin{quote}
For E-terms and R-terms this includes renormalization in the form described in the proof of
Theorem 2.
\end{quote}
For the re-filed run the remnant terms entering at these two places are the remnants of the levels
in $\mathfrak{B}(k)\cup\mathrm{band}(k)$ together with the unfiled shares
$\mathsf{sh}(\cdot,y)$ of the remnants in their zones. The re-filed remnants enter instead as groups
of E-type of the hatted regular family $\hat E$, each with its own counterterm
(Definition~\ref{9.10:def-refiled-run}). This is the treatment covered by the quoted sentence:
R-terms renormalised in the form described in \cite[proof of Theorem~2]{B-CMP122b}.

\emph{Second step: what the large-field operations read.} Inside a large-field region $X$ act the
large-field operation $\mathbf{T}'_{k+1}(X)$ and the inner operations
$\mathbf{T}_{h_X},\dots,\mathbf{T}_k$. At each level $i$ with $h_X\le i\le k+1$, the exponent they read
inside $X$ is the action \eqref{9.10:eq-refiled-run-c} at level $i$. It consists of the groups of the
hatted regular families $\hat E$; the remnants of the levels in
$\mathfrak{B}(i)\cup\mathrm{band}(i)$; the unfiled shares of remnants in zones not below their level
and not above their filing level (Lemma~\ref{9.10:lem-unfiled-range},
\eqref{9.10:eq-unfiled-range-a}); the terms $\hat{\mathbf{B}}$ and $\hat{\mathbf{D}}$; and the raw
weights $\hat x_i-\hat\zeta_i$.

Consequently a remnant of level $j$ is read by an inner operation at level $i$ only if $i\le j''$.
This happens either as a remnant of $\mathfrak{B}(i)\cup\mathrm{band}(i)$ of step $i$, or as an
unfiled share in a zone $n\in[j,j'']\subset[j,i]$. In the latter case the weight supplied by the
zone-weight statement of the correlation theorem is the zone weight $w_{n,j}$. This weight is
bounded by $\bar w_{i,j}$ up to the factor $2$ of the comparison of zone weights in the correlation
theorem; that factor is absorbed in the margin contained in $\varpi'$.

This produces the indicator $\mathbf{1}\{i\le j''\}$ in the remnant term of
\eqref{9.10:eq-lipschitz-kernel-3}, and in the kernel \eqref{9.10:eq-lipschitz-kernel-a}. The
remaining terms of $\hat\Xi_X$ are read as in the one-lattice Lipschitz step, since they involve no
index set of remnants. This proves ($\mathbf{T}$).

\emph{Third step: the kernel.} By clause (c) of geometric forgetting of old regions (hypothesis
(iii)), a remnant live at level $i$ costs at level $k+1$ the factor
$\vartheta_*^{k+1-i}\le\vartheta_*^{k-i}$. By the second step, a remnant of level $j$ is read only
at the levels $i\in[j,\min\{k,j''\}]$. Summing over these levels gives the kernel
$\hat{\mathfrak{w}}^R_{k,j}$ of \eqref{9.10:eq-lipschitz-kernel-a} in the place of the remnant
kernel of the one-lattice Lipschitz step.

For its bound, note that the number of these levels is at most $j''-j+1=\mathsf{a}_j+1$ by
\eqref{9.10:eq-filing-level-a}. For each of them
\begin{equation*}
k-i\ge k-\min\{k,j''\}=(k-j'')_+ ,
\end{equation*}
so $\vartheta_*^{k-i}\le\vartheta_*^{(k-j'')_+}$. Hence
\begin{equation*}
\hat{\mathfrak{w}}^R_{k,j}\le(\mathsf{a}_j+1)\,
\max_{j\le i\le\min\{k,j''\}}\bar w_{i,j}\;\vartheta_*^{(k-j'')_+},
\end{equation*}
which gives the bound displayed in \eqref{9.10:eq-lipschitz-kernel-a} by the inequality
\begin{equation*}
\max_{j\le i\le\min\{k,j''\}}\bar w_{i,j}\le\frac{4x_j}{x_k}+\varpi',
\end{equation*}
which is taken from the proof of the one-lattice Lipschitz step (hypothesis (i)), where its
right-hand side is the bound of the remnant kernel.

\emph{The remaining ingredients.} Several ingredients of the proof of (R$'$) in the one-lattice
Lipschitz step do not involve any index set of remnants, and they apply unchanged:
\begin{itemize}
\item the rooms furnished from the enlarged class $\hat{\mathfrak{K}}^+$ by the lemma of that step
providing the rooms;
\item the bound of that step on the raw Wilson factors at $g_j^2\delta^{\hat\zeta}_j$;
\item the difference $\delta V$ of the extensive part of the exponent (the term written $V$ in the
proof of the one-lattice Lipschitz step), with its volume structure;
\item the replacement of $+2$ by $+3$ made in that proof.
\end{itemize}

The three steps, with these remaining ingredients, give the contribution of
$\|\delta R'^{(k+1)}\|^{\flat}$ to \eqref{9.10:eq-lipschitz-kernel-2} exactly as in the one-lattice
Lipschitz step. That the readings of the first and second steps also control the contribution of
the second member $\|\delta\hat{\mathbf{B}}'^{(k+1)}\|$, so that the sum of the two members is
bounded as displayed with the same constant $K_R$, is the step $(\ast)$. This completes the proof of
(R$'$) except at the step marked $(\ast)$, and with it the proof of the theorem except at that step.
\end{proof}

\begin{lemma}[Second difference with re-filed remnant marks]\label{9.10:lem-second-difference}
Work on one lattice with one scaffold $\Theta_X$, and let
$\hat{\mathfrak{s}},\tilde{\hat{\mathfrak{s}}}\in\hat{\mathfrak{K}}_k$ be two hatted lists with
difference $\delta:=\tilde{\hat{\mathfrak{s}}}-\hat{\mathfrak{s}}$, as in
Definition~\ref{9.10:def-pair-loads}. Write $P^{\flat}$, $\hat Q$, $R$, $\hat\zeta$ for the slots of
$\hat{\mathfrak{s}}$ and $\tilde P^{\flat}$, $\tilde{\hat Q}$, $\tilde R$, $\tilde{\hat\zeta}$ for those
of $\tilde{\hat{\mathfrak{s}}}$. Use the slot sizes $\hat{\mathsf{q}}_j$, $\mathsf{r}_i$,
$\mathsf{p}^{\flat}_j$, $\hat{\mathsf{z}}_k$, $\hat{\mathsf{n}}_j$ and the bound
$\hat{\mathfrak{m}}_k$ of Definition~\ref{9.10:def-substitution-rule}, and the set
$\mathfrak{B}(k)$ of levels filed into the output of step $k$ and the re-filed family
$\mathsf{RF}[\cdot]$ of Definition~\ref{9.10:def-refiled-run}. Assume the following
statements:
\begin{itemize}
\item the proposition on the dirty part of the re-filed run, including the representation of the
dirty part used in its proof;
\item the lemma on marked differences in the re-filed run;
\item parts (b) and (c) of the re-filing lemma, and the
linearity of $R\mapsto\mathsf{RF}[R]$;
\item from the cluster expansion of the correlation theorem: the group-insertion lemma, the lemma
on remnant marks, the anchor lemma, the lemma on the holomorphy domain of the marked expansion, the
cluster corollary, and the entry of the table of activities for the vertex $\mathfrak{S}_k$;
\item the mixed-difference lemma for functions holomorphic in two variables;
\item the dichotomy in part (b) of the one-lattice Lipschitz step, together with the bound for the
third brace in the proof of that step.
\end{itemize}
For a family $\mathsf{x}$ of pure groups, a family $\mathsf{y}$ of marks and a complex number $c$, let
$\mathcal{T}(\mathsf{x};\mathsf{y};c)$ be the frozen output $E^{(k+1)}$ of the
$\hat{\mathbf{T}}$-operation of the hatted step at level $k+1$, computed with the
$P^{\flat}$-groups $\mathsf{x}$, the marks $\mathsf{y}$ and the vertex placed at $c$. The marks
used below are $(\hat Q,R_{\mathfrak{B}})$: the groups $\mathfrak{i}_y[\hat Q^{(j)}]$,
$j\le k$, each with its own functional $\beta[\cdot]$, together with the units $R^{(i)}$,
$i\in\mathfrak{B}(k)$; the symbol $0$ in the mark slot means that no mark is present, and
$(\tilde{\hat Q},\tilde R_{\mathfrak{B}})$ is defined in the same way from
$\tilde{\hat{\mathfrak{s}}}$. Then the following hold.
\begin{enumerate}
\item The dirty part at level $k+1$ is
\begin{equation}\label{9.10:eq-second-difference-1}
\hat Q^{(k+1)}
=\bigl[\mathcal{T}(P^{\flat};\hat Q,R_{\mathfrak{B}};\hat\zeta_k)
-\mathcal{T}(P^{\flat};0;\hat\zeta_k)\bigr]
+\sum_{i\in\mathfrak{B}(k)}\mathsf{RF}[R^{(i)}].
\end{equation}
\item The difference $\delta\hat Q^{(k+1)}=\tilde{\hat Q}^{(k+1)}-\hat Q^{(k+1)}$ satisfies
\begin{align}
\delta\hat Q^{(k+1)}
&=\Bigl\{\bigl[\mathcal{T}(\tilde P^{\flat};\tilde{\hat Q},\tilde R_{\mathfrak{B}};
\tilde{\hat\zeta}_k)-\mathcal{T}(\tilde P^{\flat};0;\tilde{\hat\zeta}_k)\bigr]\notag\\
&\qquad-\bigl[\mathcal{T}(\tilde P^{\flat};\hat Q,R_{\mathfrak{B}};\tilde{\hat\zeta}_k)
-\mathcal{T}(\tilde P^{\flat};0;\tilde{\hat\zeta}_k)\bigr]\Bigr\}\notag\\
&\quad+\Bigl\{\bigl[\mathcal{T}(\tilde P^{\flat};\hat Q,R_{\mathfrak{B}};\tilde{\hat\zeta}_k)
-\mathcal{T}(\tilde P^{\flat};0;\tilde{\hat\zeta}_k)\bigr]\notag\\
&\qquad-\bigl[\mathcal{T}(P^{\flat};\hat Q,R_{\mathfrak{B}};\tilde{\hat\zeta}_k)
-\mathcal{T}(P^{\flat};0;\tilde{\hat\zeta}_k)\bigr]\Bigr\}\notag\\
&\quad+\Bigl\{\bigl[\mathcal{T}(P^{\flat};\hat Q,R_{\mathfrak{B}};\tilde{\hat\zeta}_k)
-\mathcal{T}(P^{\flat};0;\tilde{\hat\zeta}_k)\bigr]\notag\\
&\qquad-\bigl[\mathcal{T}(P^{\flat};\hat Q,R_{\mathfrak{B}};\hat\zeta_k)
-\mathcal{T}(P^{\flat};0;\hat\zeta_k)\bigr]\Bigr\}
+\sum_{i\in\mathfrak{B}(k)}\mathsf{RF}[\delta R^{(i)}].
\label{9.10:eq-second-difference-a}
\end{align}
In particular $\delta\hat Q^{(k+1)}$ is a sum of differences of covariant families built on
level-indexed marks, whose marking depends on the levels only, and of the re-filed pointed families
$\mathsf{RF}[\delta R^{(i)}]$, $i\in\mathfrak{B}(k)$; this is the form required of the dirty slot of
a list in $\hat{\mathfrak{K}}_k$.
\item The bound (b) of \eqref{9.10:eq-refiled-lipschitz-a} for $\hat{\mathsf{q}}_{k+1}$ holds:
\begin{align*}
\hat{\mathsf{q}}_{k+1}
&\le\hat{\varkappa}_k\Bigl[\sum_{j\le k}\varsigma_{k+1-j}\hat{\mathsf{q}}_j
+\sum_{i\in\mathfrak{B}(k)}w_{k,i}\mathsf{r}_i\Bigr]
+C_{RF}\sum_{i\in\mathfrak{B}(k)}w_{k,i}\mathsf{r}_i\\
&\quad+K_2\hat C_L^2(1+S_I)\theta_k^2\hat{\mathfrak{m}}_k
\Bigl[\sum_{j\le k}\varsigma_{k+1-j}\mathsf{p}^{\flat}_j+\mathfrak{g}_k\hat{\mathsf{z}}_k\Bigr].
\end{align*}
\end{enumerate}
\end{lemma}

\begin{proof}
\emph{The decomposition and the identity.} In the proof of the proposition on the dirty part of the
re-filed run the output of the hatted step is written as the frozen output of the
$\hat{\mathbf{T}}$-operation with the $P^{\flat}$-groups unmarked and the groups
$\mathfrak{i}_y[\hat Q^{(j)}]$, $j\le k$, and the units $R^{(i)}$, $i\in\mathfrak{B}(k)$, as marks,
the vertex being placed at $c=\hat\zeta_k$; the dirty part $\hat Q^{(k+1)}$ is the marked part of this
output, that is the output minus the output computed without marks, together with the re-filed
families $\mathsf{RF}[R^{(i)}]$, $i\in\mathfrak{B}(k)$. This is
\eqref{9.10:eq-second-difference-1}. The marking depends on the levels only (as stated in the
proposition on the dirty part), so the same set $\mathfrak{B}(k)$ and the same marking scheme serve
for $\tilde{\hat{\mathfrak{s}}}$, and \eqref{9.10:eq-second-difference-1} holds for
$\tilde{\hat{\mathfrak{s}}}$ with every slot replaced by its tilded counterpart. On the right-hand side
of \eqref{9.10:eq-second-difference-a} the three braces telescope: the second bracket of each brace
is the first bracket of the next, so their sum is
\begin{equation*}
\bigl[\mathcal{T}(\tilde P^{\flat};\tilde{\hat Q},\tilde R_{\mathfrak{B}};\tilde{\hat\zeta}_k)
-\mathcal{T}(\tilde P^{\flat};0;\tilde{\hat\zeta}_k)\bigr]
-\bigl[\mathcal{T}(P^{\flat};\hat Q,R_{\mathfrak{B}};\hat\zeta_k)
-\mathcal{T}(P^{\flat};0;\hat\zeta_k)\bigr].
\end{equation*}
By \eqref{9.10:eq-second-difference-1} for both lists, this equals $\delta\hat Q^{(k+1)}$ minus
$\sum_{i\in\mathfrak{B}(k)}(\mathsf{RF}[\tilde R^{(i)}]-\mathsf{RF}[R^{(i)}])$, and by the
linearity of $\mathsf{RF}$ the last sum is $\sum_{i\in\mathfrak{B}(k)}\mathsf{RF}[\delta R^{(i)}]$.
This proves \eqref{9.10:eq-second-difference-a}. We bound its four terms in turn.

\emph{First brace.} Here only the marks change, from $(\hat Q,R_{\mathfrak{B}})$ to
$(\tilde{\hat Q},\tilde R_{\mathfrak{B}})$, while the unmarked data are $(\tilde P^{\flat},
\tilde{\hat\zeta}_k)$ in both brackets; the two unmarked outputs cancel, and the brace is the
difference of the two marked outputs. We apply the argument of the lemma on marked differences in
the re-filed run, which is the proof of the proposition on the dirty part with the marked family
taken to be the groups $\mathfrak{i}_y[\delta\hat Q^{(j)}]$, $j\le k$, each with its own
$\beta[\cdot]$, together with the units $\delta R^{(i)}$, $i\in\mathfrak{B}(k)$. This is legitimate
because $\delta\hat Q^{(j)}$ is a pointed $E$-family, by parts (b) and (c) of the re-filing
lemma. The group-insertion lemma of the correlation theorem is applied with its size
parameter equal to $\hat{\mathsf{q}}_j$, the lemma on remnant marks to the units $\delta R^{(i)}$,
the anchor lemma with anchor weights given by the inverse fourth power of that lemma's anchor
variable, and the cluster corollary then bounds the $\flat$-norm of the first brace by
\begin{equation*}
K_{CE}C_U\theta_k\Bigl[\sum_{j\le k}\varsigma_{k+1-j}\hat{\mathsf{q}}_j
+\sum_{i\in\mathfrak{B}(k)}w_{k,i}\mathsf{r}_i\Bigr].
\end{equation*}

\emph{Second brace.} Put
\begin{equation*}
G(\lambda,\mu):=\mathcal{T}\bigl(P^{\flat}+\lambda\,\delta P^{\flat};\,\mu\hat Q,\mu R_{\mathfrak{B}};
\,\tilde{\hat\zeta}_k\bigr).
\end{equation*}
The unmarked output $\mathcal{T}(\cdot;0;\cdot)$ does not read $\mu$, so
$G(\lambda,0)=\mathcal{T}(P^{\flat}+\lambda\,\delta P^{\flat};0;\tilde{\hat\zeta}_k)$, and since
$\tilde P^{\flat}=P^{\flat}+\delta P^{\flat}$ the second brace equals the mixed second
difference $G(1,1)-G(1,0)-G(0,1)+G(0,0)$. In the expansion of $G$ there are three kinds
of vertices. The unmarked vertices have activity at most
$\tfrac98\cdot\tfrac14\nu_0\alpha_4$. The $\lambda$-marks are the groups
$\mathfrak{i}_y[\delta P^{\flat(j)}]$, each with its own $\beta[\cdot]$; by the group-insertion lemma
with size parameter $\mathsf{p}^{\flat}_j$ and the anchor lemma, using that the inverse of the weight
attached by the anchor lemma is at most $\theta_k$, their activity is
\begin{equation*}
\mathsf{W}_1\le C_U\theta_k\sum_{j\le k}\varsigma_{k+1-j}\mathsf{p}^{\flat}_j .
\end{equation*}
The $\mu$-marks are the marks $(\hat Q,R_{\mathfrak{B}})$; their activity is
\begin{equation*}
\mathsf{W}_2\le C_U\theta_k\Bigl[S_I\max_j\hat{\mathsf{n}}_j+\sigma_k\Bigr]
\le C_U\theta_k(S_I+1)\hat{\mathfrak{m}}_k ,
\end{equation*}
where $\sigma_k:=\sum_{i\in\mathfrak{B}(k)}w_{k,i}x_i^{-\kappa_0/2}$; the second inequality is the
count made in the proposition on the dirty part, together with
$\sigma_k\le 5y_k^{-\kappa_0/2}\le\hat{\mathfrak{m}}_k$, with $y_k$ as in
Definition~\ref{9.10:def-substitution-rule}. By the lemma on the holomorphy domain of the
marked expansion, for every localisation domain $Y$ the localised family $G(\lambda,\mu)(Y)$ is
holomorphic on $|\lambda|<\varrho_1$, $|\mu|<\varrho_2$ and satisfies there the bound of that lemma,
$\|G(\lambda,\mu)(Y)\|\le E^{\sharp}e^{-\kappa d(Y)}$, where the constant $E^{\sharp}$ and the size
function $d(\cdot)$ are those of that lemma, and where
\begin{equation*}
\varrho_i:=\frac{\nu_0\alpha_4}{8C_A\mathsf{W}_i},\qquad i=1,2 .
\end{equation*}
By the bound on $\mathsf{W}_2$, $\varrho_2\ge2$ under a ceiling on $g$. We now use the dichotomy in
part (b) of the one-lattice Lipschitz step. If $\varrho_1\ge2$, the mixed-difference lemma for
functions holomorphic in two variables applies on the bidisc of radii $\varrho_1,\varrho_2$, and
bounds the mixed difference by a constant multiple of the holomorphy bound divided by
$\varrho_1\varrho_2$; by the definition of $\varrho_1$ and $\varrho_2$ this is a constant multiple
of $\mathsf{W}_1\mathsf{W}_2$. If
$\varrho_1<2$, we apply the cluster corollary at each $\lambda\in\{0,1\}$ and Schwarz's lemma in
$\mu$ to $\mu\mapsto G(\lambda,\mu)-G(\lambda,0)$, which vanishes at $\mu=0$; this bounds each of
the two $\mu$-differences by $K_{CE}\mathsf{W}_2$, hence the brace by $2K_{CE}\mathsf{W}_2$, and since
$2<4/\varrho_1$ in this case, by $4K_{CE}\mathsf{W}_2/\varrho_1$. By the definition of $\varrho_1$,
$1/\varrho_1=8C_A\mathsf{W}_1/(\nu_0\alpha_4)$, so this bound is a constant multiple of
$\mathsf{W}_1\mathsf{W}_2$. In both cases the bounds on $\mathsf{W}_1$ and $\mathsf{W}_2$ bound the
$\flat$-norm of the second brace by
\begin{equation*}
K_{\times}C_U^2\theta_k^2(S_I+1)\hat{\mathfrak{m}}_k
\sum_{j\le k}\varsigma_{k+1-j}\mathsf{p}^{\flat}_j .
\end{equation*}

\emph{Third brace.} Put $G_{\ast}(c',\mu):=\mathcal{T}(P^{\flat};\mu\hat Q,\mu R_{\mathfrak{B}};c')$.
As in the second brace, the third brace equals the $(c',\mu)$-second difference
\begin{equation*}
G_{\ast}(\tilde{\hat\zeta}_k,1)-G_{\ast}(\tilde{\hat\zeta}_k,0)-G_{\ast}(\hat\zeta_k,1)
+G_{\ast}(\hat\zeta_k,0).
\end{equation*}
By the entry of
the table of activities of the correlation theorem for the vertex, the vertex $\mathfrak{S}_k$ is
linear in $c'$ with activity at most $|c'|\mathsf{v}_k$, and the Gaussian factor requires
$|c'|g_k^2$ to stay below the fixed radius constant of the Gaussian integration. Hence $G_{\ast}$ is
holomorphic in $c'$ on the disc $|c'-\hat\zeta_k|<\varrho_{\Sigma}/(4\mathsf{v}_k)$, and in
$|\mu|<\varrho_2$ with the $\mu$-activity
$\mathsf{W}_2$ bounded as above. We apply the mixed-difference lemma in the variables
$((c'-\hat\zeta_k)/\hat{\mathsf{z}}_k,\mu)$; this is the argument for the third brace in the proof of
the one-lattice Lipschitz step, with $\mathfrak{m}_k$ replaced by $\hat{\mathfrak{m}}_k$. It bounds
the $\flat$-norm of the third brace by
\begin{equation*}
K_{\times}\mathsf{v}_k\hat{\mathsf{z}}_k\cdot C_U\theta_k(S_I+1)\hat{\mathfrak{m}}_k .
\end{equation*}

\emph{Last term.} By part (b) of the re-filing lemma,
\begin{equation*}
\Bigl\|\sum_{i\in\mathfrak{B}(k)}\mathsf{RF}[\delta R^{(i)}]\Bigr\|^{\flat}
\le C_{RF}\sum_{i\in\mathfrak{B}(k)}w_{k,i}\mathsf{r}_i .
\end{equation*}

\emph{Structure of $\delta\hat Q^{(k+1)}$.} By \eqref{9.10:eq-second-difference-a},
$\delta\hat Q^{(k+1)}$ is the sum of the three braces and the last term. Each brace is a difference
of covariant families built on level-indexed marks, whose marking depends on the levels only
(proposition on the dirty part); the last term is a sum of the re-filed pointed families
$\mathsf{RF}[\delta R^{(i)}]$, $i\in\mathfrak{B}(k)$, to which part (b) of the re-filing lemma
applies. This is the last assertion of item 2.

\emph{Summation.} The vertex activity is
$\mathsf{v}_k=C_v g_k^2(\log x_k)^{2q}\le C\,\mathfrak{g}_k$, with the vertex constant $C_v$ of the
table of activities of the correlation theorem and a constant $C$ fixed by the definition of
$\mathfrak{g}_k$ in Definition~\ref{9.10:def-pair-loads}, so that the third brace is bounded by the
$\mathfrak{g}_k\hat{\mathsf{z}}_k$ term in the normalisation of the loads of
Definition~\ref{9.10:def-pair-loads}. Adding the bounds for the three braces and the last term
inside this normalisation, with $\hat C_L$ the constant of \eqref{9.10:eq-pair-loads-a} and with
$K_2$ the constant of the same name in part (b) of the one-lattice Lipschitz step, gives the bound
(b) of \eqref{9.10:eq-refiled-lipschitz-a}.
\end{proof}

\begin{remark}
Apart from substitution into the proof of the one-lattice Lipschitz step, the argument above
contains one further step: the Cauchy second difference on the holomorphy domains of the marked
expansion in the second brace, as in part (b) of that step, with the $\mu$-marks now carrying the
units $R^{(i)}$, $i\in\mathfrak{B}(k)$, beside the groups $\mathfrak{i}_y[\hat Q^{(j)}]$. Their
activity $\sigma_k$ is already included in the count of the proposition on the dirty part, so no
new holomorphy domain is needed. In the reduced normalisation, in which the loads are multiplied by
$C_t x_{k+1}$, the coefficient of the second brace satisfies
\begin{equation*}
C_t x_{k+1}\cdot K_{\times}C_U^2\theta_k^2(S_I+1)\hat{\mathfrak{m}}_k
\le K_{\times}C_U^2(S_I+1)\,\theta_k^2\,\hat v_k\,\frac{x_{k+1}}{x_k}\,\frac{1}{C_{\beta}},
\end{equation*}
with $\hat v_k$ the letter of the reduced comparison system for the dirty slot at level $k$, and the
right-hand side is small compared with $\varkappa$.
\end{remark}

\begin{corollary}[Two-slot bounds with finite-range remnant load]\label{9.10:cor-two-slot-bounds}
Work on one lattice with one scaffold $\Theta_X$. Let $\hat{\mathfrak{s}},\tilde{\hat{\mathfrak{s}}}\in
\hat{\mathfrak{K}}_k$ be two re-filed lists in the sense of
Definition~\ref{9.10:def-substitution-rule}, and put $\delta:=\tilde{\hat{\mathfrak{s}}}-\hat{\mathfrak{s}}$.
Assume the following two hypotheses.
\begin{enumerate}
\item[(1)] The conclusions (a)--(g) of the re-filed one-lattice Lipschitz step,
Theorem~\ref{9.10:thm-refiled-lipschitz} (whose proof is incomplete at one step), hold for the
pair $\hat{\mathfrak{s}},\tilde{\hat{\mathfrak{s}}}$.
\item[(2)] The two-slot corollary of the one-lattice Lipschitz step holds, with coefficients
$\alpha_r$ and $\varkappa$; its proof is read below with every entry of an unhatted list replaced
by the corresponding entry of the re-filed list according to the substitution
rule~\eqref{9.10:eq-substitution-rule-a}.
\end{enumerate}
The reduced slots of that two-slot corollary, read at depth $r$, are
\begin{equation*}
\mathfrak{e}^P_r:=\Bigl(E_0^{-1}\mathsf{P},\ (\log\theta_r)^{-c_1}E_0^{-1}P^{\flat},\
C_t\hat x_r\hat Q\Bigr),\qquad \mathfrak{e}^R_r:=C_t\hat x_r R.
\end{equation*}
Here the entries are those of the slot read at depth $r$, and the first component is measured at
radii that do not depend on $g$. The number $E_0$ is the bound of the class
$\hat{\mathfrak{K}}_k$ in Definition~\ref{9.10:def-substitution-rule}; $C_t\hat x_r$ is the factor
that converts a slot at depth $r$ into reading units; and $c_1$ is the exponent of the
normalisation of $P^{\flat}$ in that corollary. Depth $r$ corresponds to the output level $k+1$. A
slot at a depth $s>r$ corresponds to a level $j\le k$, and then $s-r=k+1-j$.
Let $q_1$ be a geometric rate with $q\le q_1^2$, and put $\ell_r=\mathsf{a}_r+2$, so that
$\ell_r\ge k-k_1(k)$, where $k_1(\cdot)$ is the level function of that corollary. Define the
\emph{finite-range remnant load}
\begin{equation}\label{9.10:eq-two-slot-bounds-a}
\mathsf{S}^+_r:=\sum_{s>r}(\ell_s+2)\,q_1^{(s-r-\ell_s)_+}\,\|\delta\mathfrak{e}^R_s\| .
\end{equation}
For $r<s\le r+\ell_s$ the weight in~\eqref{9.10:eq-two-slot-bounds-a} is $\ell_s+2$. For
$s>r+\ell_s$ it decreases geometrically at rate $q_1$.

Then, on $\hat{\mathfrak{K}}_k$, the one-step outputs $\hat{\mathsf{F}}^P_r,\hat{\mathsf{F}}^R_r$ at
depth $r$ satisfy
\begin{align}
\|\delta\hat{\mathsf{F}}^P_r\|&\le\alpha_r|\delta\xi|
+\varkappa\sum_{s>r}q_1^{s-r}\|\delta\mathfrak{e}^P_s\|+C_F\,\mathsf{S}^+_r,
\label{9.10:eq-two-slot-bounds-1}\\
\|\delta\hat{\mathsf{F}}^R_r\|&\le\mathfrak{r}_r\Lambda_r\Bigl[\sum_{s\ge r}q_1^{s-r}
\bigl(g_s^2|\delta\hat\zeta_s|+\|\delta\mathfrak{e}^P_s\|\bigr)+\mathsf{S}^+_r\Bigr]
\notag\\
&\quad+\varkappa'_r\sum_{r<s<r+\ell_r}\|\delta\mathfrak{e}^R_s\|,
\label{9.10:eq-two-slot-bounds-2}
\end{align}
where
\begin{equation*}
C_F:=3C_{RF}+1,\qquad \varkappa'_r:=6\hat\varkappa_k,\qquad
\mathfrak{r}_r=C_t\theta_r^{1-\kappa_0/2},\qquad \Lambda_r=\Lambda(\log\theta_r)^{c_3}.
\end{equation*}
The coefficients $\alpha_r$ and $\varkappa$ are those of the two-slot corollary, with $\varkappa$
read as $\hat\varkappa_k$; $\delta\xi$ is the difference, between the two lists, of the datum on
which $\alpha_r$ acts in that corollary; $\kappa_0$ is the exponent of the class
$\hat{\mathfrak{K}}_k$ in Definition~\ref{9.10:def-substitution-rule}; and $\Lambda$, $c_3$ are the
constant and the exponent of the polylogarithmic factor $\Lambda_r$ of that corollary.
\end{corollary}

\begin{proof}
We follow the proof of the two-slot corollary of the one-lattice Lipschitz step, line by line. By
hypothesis~(2), each unhatted entry is replaced by its re-filed counterpart through the
substitution rule~\eqref{9.10:eq-substitution-rule-a}. The bounds invoked are those of
hypothesis~(1), that is, of Theorem~\ref{9.10:thm-refiled-lipschitz}, whose proof is incomplete at
one step.

\emph{The bound~\eqref{9.10:eq-two-slot-bounds-1}.} We treat the three components of
$\delta\hat{\mathsf{F}}^P_r$ in turn.

The first, pure, component is bounded by clause~(g) of Theorem~\ref{9.10:thm-refiled-lipschitz},
whose proof is incomplete at one step. For the pure component that clause gives the bound of the
two-slot corollary with equality.

The second component is bounded by clause~(a) of the same theorem, whose proof is incomplete at
one step:
\begin{equation*}
\mathsf{p}^{\flat}_{k+1}\le\hat\varkappa_k\Bigl[\sum_{j\le k}\varsigma_{k+1-j}\mathsf{p}^{\flat}_j
+\mathfrak{g}_k\hat{\mathsf{z}}_k\Bigr].
\end{equation*}

The third component is $C_t\hat x_k$ times the bound of clause~(b) of
Theorem~\ref{9.10:thm-refiled-lipschitz}, whose proof is incomplete at one step. In that bound the
sums carrying the kernel $\varsigma$ are treated verbatim as in the proof of the two-slot corollary.
The remnant members of clause~(b) are
\begin{equation*}
\hat\varkappa_k\sum_{i\in\mathfrak{B}(k)}w_{k,i}\mathsf{r}_i
+C_{RF}\sum_{i\in\mathfrak{B}(k)}w_{k,i}\mathsf{r}_i ,
\end{equation*}
with $\mathfrak{B}(k)$ the set of levels filed into the output of step~$k$,
as in~\eqref{9.10:eq-filing-level-a}. We convert them to reading units, that is, we multiply by
$C_t\hat x_{k+1}$. For a level $i$ read at depth $s$ we write
$\mathsf{r}_i=\|\delta\mathfrak{e}^R_s\|/(C_t\hat x_i)$, and we use $w_{k,i}\hat x_{k+1}/\hat x_i\le 3$.
This gives
\begin{equation*}
C_t\hat x_{k+1}\bigl(\hat\varkappa_k+C_{RF}\bigr)\sum_{i\in\mathfrak{B}(k)}w_{k,i}\mathsf{r}_i
\le\bigl(3\hat\varkappa_k+3C_{RF}\bigr)\sum_{r<s\le r+\ell_s}\|\delta\mathfrak{e}^R_s\|
\le C_F\,\mathsf{S}^+_r .
\end{equation*}
The last inequality uses $C_F=3C_{RF}+1$, the constant of the two-slot dictionary of the
transplanted one-lattice response bounds, obtained from the re-filing norm constant $C_{RF}$ in
reading units, in which $w_{k,i}\hat x_{k+1}/\hat x_i\le3$. It also uses the fact that for
$r<s\le r+\ell_s$ the exponent in~\eqref{9.10:eq-two-slot-bounds-a} vanishes and the weight there
is $\ell_s+2\ge1$.

The coefficients of the $P^{\flat}$-terms and of the $\hat\zeta$-terms are bounded by $\varkappa$,
as in the proof of the two-slot corollary. Two inputs are needed for this. First,
$\hat{\mathfrak{m}}\le\mathfrak{m}$, which is part of the definition of the class
$\hat{\mathfrak{K}}_k$ in Definition~\ref{9.10:def-substitution-rule}. Second, the hatted vertex
activity is bounded by the unhatted one, $\hat{\mathsf{v}}_k\le\mathsf{v}_k$. Collecting the three
components gives~\eqref{9.10:eq-two-slot-bounds-1}.

\emph{The bound~\eqref{9.10:eq-two-slot-bounds-2}.} Here $\delta\hat{\mathsf{F}}^R_r$ is bounded by
$C_t\hat x_k$ times the sum of the bounds of clauses~(c) and~(e) of
Theorem~\ref{9.10:thm-refiled-lipschitz}, whose proof is incomplete at one step.

The band term of clause~(c), which is the same as the band term of clause~(c) of the one-lattice
Lipschitz step, gives the last term of~\eqref{9.10:eq-two-slot-bounds-2}, with
$\varkappa'_r=6\hat\varkappa_k$.

Clause~(e) gives the bound $4E_0^{-1}e^{-p_0(g_k)}\Sigma^q_k[\hat\delta]$, where $p_0(\cdot)$ is
the degree function. By~\eqref{9.10:eq-pair-loads-a},
\begin{equation*}
\Sigma^q_k[\hat\delta]=\sum_{n\le k}(k-n+1)\,q^{(k-n-1)_+}\,[\![\hat\delta]\!]_n ,
\end{equation*}
the index $n$ running over zones.

The parts of $[\![\hat\delta]\!]_n$ carrying $\hat{\mathsf{e}}$, $\hat{\mathsf{z}}$ and
$\hat{\mathsf{b}}$ give the $q_1$-sums of the first bracket of~\eqref{9.10:eq-two-slot-bounds-2}
exactly as in the proof of the two-slot corollary. That argument rests on the inequality
\begin{equation*}
\sum_{n}(k-n+1)\,q^{(k-n-1)_+}\,\varsigma_{1+n-j}\le C_q\,q_1^{k-j},
\end{equation*}
where $C_q$ is a constant of that argument.
The polylogarithmic factors are absorbed by
$e^{-p_0(g_k)}\le g_k^{\kappa_0}(\log x_k)^{-c_1-2q_{\flat}}$, where $q_{\flat}$ is the exponent
so named in the two-slot corollary.

The remnant part of $[\![\hat\delta]\!]_n$ is $\sum_{n-\mathsf{a}_j\le j\le n}w_{n,j}\mathsf{r}_j$. The
condition $n-\mathsf{a}_j\le j\le n$ is equivalent to $j\le n\le j+\mathsf{a}_j$. Hence a slot of
level $j$, at depth $s$, is read only from the zones $n\in[j,j+\mathsf{a}_j]$. Each such zone
carries the factor $(k-n+1)q^{(k-n-1)_+}$. After multiplication by $C_t\hat x_k$, the zone-weight
bound of the correlation theorem shows that each such zone carries a weight bounded by an absolute
constant.

There are at most $\ell_s$ such zones. Summing over them, and using
$(k-n)q^{k-n}\le C_qq_1^{k-n}$, we find that the coefficient of $\|\delta\mathfrak{e}^R_s\|$ is
bounded as follows, for suitable constants $C$ and $C'$:
\begin{equation*}
\begin{split}
C(\ell_s+1)\max_{n\in[j,j+\mathsf{a}_j]}(k-n+1)\,q^{(k-n-1)_+}
&\le C'(\ell_s+2)\,q_1^{(k-j-\mathsf{a}_j-1)_+}\\
&\le C'(\ell_s+2)\,q_1^{(s-r-\ell_s)_+}.
\end{split}
\end{equation*}
The last inequality holds because depth $s$ corresponds to level $j$, so that $s-r=k+1-j$ and
$\ell_s=\mathsf{a}_j+2$, the numbers $\mathsf{a}_s$ and $\mathsf{a}_j$ being the same count for
this slot; hence $s-r-\ell_s=k-j-\mathsf{a}_j-1$. The resulting sum over $s>r$ is the load
$\mathsf{S}^+_r$ of~\eqref{9.10:eq-two-slot-bounds-a}. It stands inside the bracket multiplied by
$\mathfrak{r}_r\Lambda_r$, as the whole output of clause~(e) does.

Adding the contributions of clauses~(c) and~(e) gives~\eqref{9.10:eq-two-slot-bounds-2}.
\end{proof}

In the maximum-principle system for the envelopes that follows, the load $\mathsf{S}^+_r$ is paid by
the finite-range supremum $\mathfrak{S}$ in the corresponding bound of that system.

\begin{corollary}[Three-slot hybrid bounds for the re-filed run]\label{9.10:cor-three-slot-bounds}
Let $\mathsf{A}$, $\mathsf{B}$ be a pair of runs and let $n$ be the length of the run; from here
on $r$ denotes a depth with $r<n$, and $s$ a slot.
Assume the following statements.
\begin{itemize}
\item The re-filed one-lattice Lipschitz step, Theorem~\ref{9.10:thm-refiled-lipschitz}, whose
proof is incomplete at one step. It is taken with its constants evaluated at the constant set
$c$, on the class
$\hat{\mathfrak{K}}$ of Definition~\ref{9.10:def-substitution-rule}, on the tube, the domain on
which the class statements about the production maps are made. Here the
class statements on the tube about the production maps are read with every entry hatted as in
the substitution rule \eqref{9.10:eq-substitution-rule-a}; in the hybrid construction this means
in particular $Q\mapsto\hat Q$ and $D\mapsto\hat D$.
\item The two-slot bounds with finite-range remnant load,
Corollary~\ref{9.10:cor-two-slot-bounds}.
\item Part (c) of the comparison theorem for two lists on one lattice.
\item The transplant lemma, in its clause on the coupling row.
\end{itemize}
Form the hybrid one-step outputs $\hat M^P_r$, $\hat M^A_r$, $\hat M^R_r$ from competitors
$G^P_s$, $G^A_s$, $G^R_s$ as in the un-hatted three-slot hybrid construction, with $Q$ replaced
by $\hat Q$ and $D$ by $\hat D$. Let $\mathsf{S}^+_r$ be the finite-range remnant load
\eqref{9.10:eq-two-slot-bounds-a}, formed with $\|G^R_s\|$ in place of $\|\delta\mathfrak{e}^R_s\|$.
The constants $\alpha_r$, $\alpha^A_r$, $\varkappa$, $\varkappa'_r$, $\varkappa''_r$,
$\varkappa_A$, $C_F$, $\mathfrak{r}_r$, $\Lambda_r$, $q$ and $\ell_r$ are those of
Theorem~\ref{9.10:thm-refiled-lipschitz} (whose proof is incomplete at one step) in the above
reading, together with those of Corollary~\ref{9.10:cor-two-slot-bounds}. In the last row
below, $\mathsf{A}$ also denotes the Lipschitz constant of the coupling row; $G^{\mathsf{P}}_s$
and $G^{\hat Q}_s$ denote the two components of the $P$-slot competitor $G^P_s$ that the
coupling-extraction functional reads, corresponding respectively to the summands
$\beta[P^{\flat}]$ and $\beta[\hat Q]$ of the decomposition of the hatted coupling increment
used in the proof; and $\beta^{\mathsf{A}}$
denotes the coupling-extraction functional $\beta[\cdot]$ of the run $\mathsf{A}$.

Then, for every such pair and every $r<n$, the transplanted
three-slot hybrid bounds hold under the substitution rule \eqref{9.10:eq-substitution-rule-a}
in the following form:
\begin{equation}\label{9.10:eq-three-slot-bounds-a}
\begin{aligned}
&(\hat{\mathrm{P}})\quad \|\hat M^P_r\| \le \alpha_r d^{\sharp}_{r+1}
 + \varkappa\sum_{s>r}q^{s-r}\|G^P_s\|
 + \varkappa''_r\sum_{s>r}q^{s-r}\|G^A_s\| + C_F\mathsf{S}^+_r;\\
&(\hat{\mathrm{A}})\quad \|\hat M^A_r\| \le \alpha^A_r d^{\sharp}_{r+1}
 + \varkappa_A\Bigl[\|G^P_r\| + \sum_{s>r}q^{s-r}\bigl(\|G^P_s\|+\|G^A_s\|\bigr)
 + \mathsf{S}^+_r\Bigr]\\
&\qquad\qquad\qquad + \|G^R_r\| + C_F\mathsf{S}^+_r;\\
&(\hat{\mathrm{R}})\quad \|\hat M^R_r\| \le \mathfrak{r}_r\Lambda_r\hat B_r(G,d)
 + \varkappa'_r\sum_{r<s<r+\ell_r}\|G^R_s\|,\\
&\qquad\qquad \hat B_r(G,d) := \sum_{r\le s<n}q^{s-r}g^2_{(s+1)}d_{s+1}
 + \sum_{s\ge r}q^{s-r}\bigl(\|G^P_s\|+\|G^A_s\|\bigr) + \mathsf{S}^+_r;\\
&(\hat{\mathrm{b}})\quad \hat b^{\mathsf{B}}_s-\hat b^{\mathsf{A}}_s
 = \beta^{\mathsf{A}}[G^{\mathsf{P}}_s]+\beta^{\mathsf{A}}[G^{\hat Q}_s]\ \text{ exactly, for every
 competitor},\\
&\qquad\qquad \bigl|\hat b^{\mathsf{B}}_s-\hat b^{\mathsf{A}}_s\bigr| \le \mathsf{A}\|G^P_s\|.
\end{aligned}
\end{equation}
\end{corollary}

\begin{proof}
We follow the proof of part (a) of the transfer theorem for the transplanted lists. That proof
rests on one observation. The two lists compared there, the transplanted list and the list of
the run $\mathsf{A}$ with a zero entry adjoined, are two lists on one lattice and on one
scaffold. Three statements therefore
apply to them by statement, at the set of constants $c$: the one-lattice Lipschitz step,
together with the proposition from which it is derived; the un-hatted two-slot bounds; and
part (c) of the comparison theorem for two lists on one lattice. Nothing has to be known about
how the two lists read a term.

We now run the same argument with two substitutions. The re-filed one-lattice Lipschitz step,
Theorem~\ref{9.10:thm-refiled-lipschitz} (whose proof is incomplete at one step), takes the
place of the one-lattice Lipschitz step. Corollary~\ref{9.10:cor-two-slot-bounds}
takes the place of the un-hatted two-slot bounds. Part (c) of the comparison theorem is kept.

Under the substitution rule \eqref{9.10:eq-substitution-rule-a}, the two lists are hatted lists
of the class $\hat{\mathfrak{K}}$ of Definition~\ref{9.10:def-substitution-rule}. They still lie
on one lattice and one scaffold. The class statements on the tube are statements about the
production maps, and they read their entries hatted as above, so they hold for the hatted
production. Hence Theorem~\ref{9.10:thm-refiled-lipschitz}, whose proof is incomplete at one
step, applies at the set $c$ on the class, on the tube. Corollary~\ref{9.10:cor-two-slot-bounds}
and part (c) of the comparison theorem also apply by statement. As before, nothing has to be
known about how the hatted lists read a term. This gives the rows $(\hat{\mathrm{P}})$ and
$(\hat{\mathrm{R}})$ of \eqref{9.10:eq-three-slot-bounds-a}, with $\mathsf{S}^+_r$ formed with
$\|G^R_s\|$ and with the bracket $\hat B_r(G,d)$ as displayed.

For the row $(\hat{\mathrm{A}})$ it remains to account for its remnant content, which has two
parts.
\begin{itemize}
\item The pure part $P^{\flat}$ reads no remnant.
\item The entry $\hat D$ is the freezing defect of the local production. This production reads
the same hatted entries as the production bounded in $(\hat{\mathrm{P}})$. The remnant
contribution to $(\hat{\mathrm{A}})$ is therefore the load $\mathsf{S}^+_r$, entering with the
coefficient $\varkappa_A$.
\end{itemize}
The same-depth term $\|G^R_r\|$ is carried exactly as in the un-hatted three-slot hybrid
construction. A void entry only enlarges a right member, so keeping this term as a separate
summand leaves $(\hat{\mathrm{A}})$ valid.

For the row $(\hat{\mathrm{b}})$, recall that the coupling-extraction functional $\beta$ reads
the total family, by the statement on the total family of the hatted production. For each run
the hatted coupling increment is therefore
\begin{equation*}
\hat b = \beta[\hat E] = \beta[P^{\flat}]+\beta[\hat Q].
\end{equation*}
We write this identity for the run $\mathsf{B}$ and for the run $\mathsf{A}$ at slot $s$, and
subtract. With the transplant lemma, in its clause on the coupling row, this gives the row
$(\hat{\mathrm{b}})$ of \eqref{9.10:eq-three-slot-bounds-a}: the difference
$\hat b^{\mathsf{B}}_s-\hat b^{\mathsf{A}}_s$ equals
$\beta^{\mathsf{A}}[G^{\mathsf{P}}_s]+\beta^{\mathsf{A}}[G^{\hat Q}_s]$ exactly, for every
competitor, and it is bounded in absolute value by $\mathsf{A}\|G^P_s\|$, where $\mathsf{A}$ is
the Lipschitz constant of the coupling row.
\end{proof}

\begin{proposition}[Cross-cutoff transplant of re-filed lists]\label{9.10:prop-transplant}
Let $\mathsf{A}$ and $\mathsf{B}$ be the two runs compared by the comparison chain of
Definition~\ref{9.6:def-comparison-chain}, each a re-filed renormalisation run in the sense of
Definition~\ref{9.10:def-refiled-run}. For a level-$j$ family $F$ of $\mathsf{B}$ in the class
$\flat$, its \emph{transplant} is
\begin{equation}\label{9.10:eq-transplant-1}
\mathfrak{T}F := \sum_{(Y',Y_0)} \prod \int d\mathsf{s}\;
\partial_{\mathsf{s}} F\bigl(Y';\Theta_{\mathsf{s}}(\mathbf{U},\mathbf{J});w'\bigr),
\end{equation}
with the point $y$ carried as a further argument,
$F(Y';\Theta_{\mathsf{s}}(\mathbf{U},\mathbf{J}),y;w')$, when $F$ is pointed. Here the sum runs
over the pairs $(Y',Y_0)$ of the transplant's decomposition, the product and the integrals run
over the interpolation parameters $\mathsf{s}$, and $\Theta_{\mathsf{s}}(\mathbf{U},\mathbf{J})$
is the interpolating configuration, whose arguments $\mathbf{U}$, $\mathbf{J}$ are fixed in the
transplant lemma; $w'$ is the source argument of $F$, and $C_V$ below is the constant of the
transplant lemma. The transplant re-localises $F$ onto the fat domains
$Y=Y_0^{\sim 5}$ of the lattice of $\mathsf{A}$. We write \textup{(f1)}, \textup{(f2)},
\textup{(f3)} for the three defining conditions of a family, \textup{(f3)}$^{\mathrm{bg}}$ for
\textup{(f3)} required on backgrounds only, and
$\simeq^{\mathrm{bg}}$ for the equivalence $\simeq$ of families taken on backgrounds. Assume:
\begin{enumerate}
\item[(h1)] the transplant lemma, properties (T1)--(T3). (T1): $\mathfrak{T}F$ is of class $c$,
$\|\mathfrak{T}F\|^{c}\le C_V\|F\|^{\flat}$, $\mathfrak{T}$ is linear, and the parameters (the
point $y$ and the source) pass through $\mathfrak{T}$. (T2): the totals of $F$ and
$\mathfrak{T}F$ agree on backgrounds. (T3): (i) $\mathfrak{T}F$ satisfies
\textup{(f3)}$^{\mathrm{bg}}$; (ii) $\beta^{\mathsf{A}}[\mathfrak{T}F]=\beta^{\mathsf{B}}[F]$;
moreover $\simeq^{\mathrm{bg}}$ passes through $\mathfrak{T}$ and the vacuum sums of $F$ and
$\mathfrak{T}F$ agree;
\item[(h2)] the re-filing lemma, statements (a)--(d), for the re-filed family
$\mathsf{RF}[R]$ at scale $i''$ of a remnant $R$ of level $i$; (a) is its statement on totals and
(d) is the identity $\beta[\mathsf{RF}[R]]=\beta[R]$;
\item[(h3)] the chain lemma for the hybrid lists $\mathfrak{s}^{\natural}=\mathfrak{s}^{\mathsf{A}}+G$
of the comparison chain, in its two alternatives for the competitor $G$, in both of which the
$D$-components (for hatted lists, the $\hat D$-components) of the competitor vanish at constant
sources;
\item[(h4)] the transfer theorem for transplanted lists, clauses (a) and (b), as proved for
un-hatted lists;
\item[(h5)] the chain lemma and the proof of clause (a) of the transfer theorem use, of a hybrid
list, only: its list structure; its membership in the list class $C_{\mathfrak{K}}\mathfrak{K}^{c}$
of the transfer theorem; and that
$P^{\flat}\simeq\mathsf{P}$ and $b=\beta[E]$ are statements on totals, which both alternatives
share, while the $D$-components vanish at constant sources in both.
\end{enumerate}
Then:
\begin{enumerate}
\item[(i)] the renormalised form $(\mathfrak{T}\hat E^{\mathsf{B}})^{\mathrm{ren}}$ of
$\mathfrak{T}\hat E^{\mathsf{B}}$ is an admissible $E$-slot of a
list of $\mathsf{A}$, and $\hat b=\beta[\hat E]$ is a statement on totals in both alternatives;
the competitor for the $\hat Q$-component is, in the second alternative,
\begin{equation}\label{9.10:eq-transplant-2}
G^{\hat Q}_s := \bigl(\mathfrak{T}\hat E^{\mathsf{B}}_s-\hat E^{\mathsf{A}}_s\bigr)^{\hat Q},
\end{equation}
and, in the first alternative, the $\hat Q$-component of the competitor that the chain lemma forms
there, with $Q$ replaced by $\hat Q$;
\item[(ii)] the transplanted remnants $\mathfrak{T}R^{(i),\mathsf{B}}$ are admissible $R$-slots of
the run $\mathsf{A}$, treated as remnants at the levels below $i''$ and re-filed at step $i''-1$;
\item[(iii)] clauses (a) and (b) of the transfer theorem hold for hatted lists, with the re-filed
one-lattice Lipschitz step (Theorem~\ref{9.10:thm-refiled-lipschitz}, whose proof is incomplete at
one step) and the slot bounds of Corollaries~\ref{9.10:cor-two-slot-bounds}
and~\ref{9.10:cor-three-slot-bounds} in place of their un-hatted counterparts; clause (d) of the
transfer theorem is the same statement for hatted and un-hatted lists.
\end{enumerate}
\end{proposition}

\begin{proof}
\emph{The $E$-slots.} The hatted regular family of $\mathsf{B}$ is the sum
$\hat E^{\mathsf{B}}=E^{\mathsf{B}}+\mathsf{RF}[R^{\mathsf{B}}]$ of its un-hatted regular family
and the re-filed families of its remnants (Definition~\ref{9.10:def-refiled-run}). Since
$\mathfrak{T}$ is linear by (T1),
\begin{equation*}
\mathfrak{T}\hat E^{\mathsf{B}}=\mathfrak{T}E^{\mathsf{B}}+\mathfrak{T}\mathsf{RF}[R^{\mathsf{B}}].
\end{equation*}
The first summand is the transplanted $E$-slot of the un-hatted case. For the second,
$\mathsf{RF}[R^{\mathsf{B}}]$ is a pointed family at scale $i''$, and by (T1) the point $y$
passes through $\mathfrak{T}$ as a parameter; hence $\mathfrak{T}\mathsf{RF}[R^{\mathsf{B}}]$ is a
pointed family of $\mathsf{A}$ in the class $c$ at scale $i''$. It satisfies
\textup{(f3)}$^{\mathrm{bg}}$ by (T3)(i). By (T3)(ii) applied to $F=\mathsf{RF}[R^{\mathsf{B}}]$,
and then by statement (d) of the re-filing lemma,
\begin{equation*}
\beta^{\mathsf{A}}\bigl[\mathfrak{T}\mathsf{RF}[R^{\mathsf{B}}]\bigr]
=\beta^{\mathsf{B}}\bigl[\mathsf{RF}[R^{\mathsf{B}}]\bigr]
=\beta^{\mathsf{B}}[R^{\mathsf{B}}].
\end{equation*}
Therefore $(\mathfrak{T}\hat E^{\mathsf{B}})^{\mathrm{ren}}$ has the properties that an $E$-slot
of a list of $\mathsf{A}$ is required to have. Moreover
$\beta^{\mathsf{A}}[\mathfrak{T}E^{\mathsf{B}}]=\beta^{\mathsf{B}}[E^{\mathsf{B}}]$ by (T3)(ii);
with the last display this shows that $\hat b=\beta[\hat E]$ is, in both alternatives, a
statement on totals. The comparison chain never uses an identity
$\mathfrak{T}\circ\mathsf{RF}=\mathsf{RF}\circ\mathfrak{T}$. Such an identity fails term by term,
since both operations are re-localisations; on totals it holds, by (T2) for $\mathfrak{T}$ and by
statement (a) of the re-filing lemma for $\mathsf{RF}$. The competitor for the $\hat Q$-component is
then formed as in clause~(i), exactly as the chain lemma forms it for the
$Q$-component in each of its two alternatives. This proves~(i).

\emph{The $R$-slots.} By (T1), $\mathfrak{T}R^{(i),\mathsf{B}}$ is a family of $\mathsf{A}$ at
level $i$ on fat domains. In general these domains are not unions of $M\check R_i$-cubes. The
hatted step of $\mathsf{A}$ uses $\mathfrak{T}R^{(i),\mathsf{B}}$ as a remnant, through the
remnant terms and band terms of its remnant lemma, at the levels below $i''$ only. That lemma
uses of a remnant only \textup{(f1)}, \textup{(f2)} and a norm bound, all supplied by (T1), and
not its cube structure. At step
$i''-1$ the step of $\mathsf{A}$ files the re-filed family $\mathsf{RF}[\,\cdot\,]$ of the hybrid
$R$-slot into its
$\hat E^{(i'')}$. This is possible because the shares $\mathsf{sh}(X,y)$, the assignment domain
$Y(X,y)$ and the re-filing are defined for arbitrary level-$i$ domains $X\in\mathfrak{D}_i$. Their
definition does not depend on whether $X$ is a union of cubes. The alignment condition of the
re-filing concerns the cells $\square'_y$ of $\mathbb{T}^{(i'')}$ and the $M$-cubes
$\tilde{\square}(y)$, not $X$.

Statements (a)--(d) of the re-filing lemma hold for the hybrid $R$-slot with \textup{(f3)} in its
pointed-total form, on backgrounds. Nothing more is used downstream: the one-lattice Lipschitz
step uses \textup{(f3)} and $\simeq$ only through its kernel lemma, and there on backgrounds only.
The covariance \cite[(2.32)]{B-CMP119} of the hybrid slot under partition-preserving isometries on
backgrounds follows from (T3)(i) applied to $F=R^{(i),\mathsf{B}}$; that (T3)(i) yields it rests on
three facts. First, $R^{(i),\mathsf{B}}$ has the
covariance \cite[(2.32)]{B-CMP119} at constants. Second, $\Theta_{\mathsf{s}}$ is built from the
covariant operators of Ba{\l}aban's construction, so the interpolating configurations are gauge
covariant. Third, the sum over domains in \eqref{9.10:eq-transplant-1} is invariant under the
isometries. This proves~(ii).

\emph{Conclusion.} By (i) and (ii), the hybrid list of (h3) formed from the hatted lists,
$\hat{\mathfrak{s}}^{\natural}:=\hat{\mathfrak{s}}^{\mathsf{A}}+G$, has list structure, its $E$-slot
and its $R$-slots being admissible slots of a list of $\mathsf{A}$. Its transplanted members are
of class $c$ with the norm bound of (T1), which is the membership condition for
$C_{\mathfrak{K}}\mathfrak{K}^{c}$ verified in the proof of the transfer theorem. The statements
$P^{\flat}\simeq\mathsf{P}$ and
$\hat b=\beta[\hat E]$ are statements on totals, which both alternatives share; the coupling
statement was treated in (i). The $\hat D$-components vanish at constant
sources in both alternatives, by (h3). These are the only properties of a
hybrid list that the chain lemma and the proof of clause (a) of the transfer theorem use, by
(h5). Hence the proof of clauses (a) and (b) of the transfer theorem, hypothesis (h4),
applies verbatim to hatted lists. The one-step input is the re-filed one-lattice Lipschitz step,
Theorem~\ref{9.10:thm-refiled-lipschitz}, whose proof is incomplete at one step. The bounds used
for the slots are those of Corollaries~\ref{9.10:cor-two-slot-bounds}
and~\ref{9.10:cor-three-slot-bounds}. This proves~(iii).

The proof contains no estimate. The part of the transfer theorem that carries the comparison of
lists over to the observables of Theorem~\ref{3.1:thm-main} rests only on the transplant and the
inputs of its construction; the re-filing does not enter it, and it is used here as it stands in
the transfer theorem.
\end{proof}

\begin{proposition}[Geometric source bounds for the re-filed run. Stated; proof incomplete at the step
marked $(\ast)$]\label{9.10:prop-source-bounds}
Consider the re-filed run of length $n$, its lists re-filed by the substitution rule
\eqref{9.10:eq-substitution-rule-a} of Definition~\ref{9.10:def-substitution-rule}, and the original
run from which it is obtained by that rule; assume the hypotheses of the source-row bounds of the
original run. Let
$N^P_r$, $N^A_r$, $N^{R'}_r$, $N^{R''}_r$ be the source rows of the re-filed run at depth $r$.
Let $G^R_s$ be the
competitors, $d=(d_{s+1})_s$ the coupling differences and $\hat B_r(G,d)$ the bracket of
\eqref{9.10:eq-three-slot-bounds-a}. Let $\ell_r=\mathsf{a}_r+2$ be the remnant depth range of
\eqref{9.10:eq-filing-level-a}, and let $\mathfrak{r}_r=C_t\theta_r^{1-\kappa_0/2}$ with $\kappa_0>2$.
Let $C_\sigma$ be the source constant of the original run. Let $C_\iota$ and $\vartheta_\chi$ be the
constant and the rate of the index-stationarity errors, with $\vartheta_\chi\le\rho<1$. Assume the
coupling ceiling
\begin{equation}\label{9.10:eq-source-bounds-a}
\sup_{r\ge 0}\,6C_\iota(\ell_r+3)\,\vartheta_\chi^{-\ell_r-3}\,\mathfrak{r}_r\le 1 ,
\end{equation}
whose right member is read in the normalisation of the source constant $C_\sigma$. This is an upper
bound on $g$ alone once $L$, the auxiliary parameters $\mathfrak{p}$, the exponent $\theta_\chi$
attached to the rate $\vartheta_\chi$ and $\kappa_0>2$ are fixed; it places no condition on $\rho$ and
is not a cap in the sense of Definition~\ref{2.3:def-cap}. Let
$C_V$, $C^\star_X$ and
$C^z_\sigma$ be the further constants entering the envelope constant
\[
C_E=\tfrac87\max\{C_V+1,\,2C_\sigma,\,2C^\star_X,\,2C^z_\sigma\}
\]
of the original run. Put
\begin{equation}\label{9.10:eq-source-bounds-b}
\hat C_\sigma:=2C_\sigma,\qquad
\hat C_E:=\tfrac87\max\{C_V+1,\,2\hat C_\sigma,\,2C^\star_X,\,2C^z_\sigma\}.
\end{equation}
Then at every depth $r$ of the run
\begin{equation}\label{9.10:eq-source-bounds-1}
\begin{aligned}
\|N^P_r\|,\ \|N^A_r\| &\le \hat C_\sigma\,\rho^{n-r},\\
\|N^{R''}_r\| &\le \mathfrak{r}_r\Lambda_r\,\hat C_\sigma\,\rho^{n-r},\\
\|N^{R'}_r\| &\le \mathfrak{r}_r\Lambda_r\bigl[(\hat C_\sigma+1)\rho^{n-r}+\hat B_r(G,d)\bigr]
+\varkappa'_r\sum_{r<s<r+\ell_r}\|G^R_s\| .
\end{aligned}
\end{equation}
\end{proposition}

\begin{proof}
\emph{The rows of the original run.} In this paragraph the same letters $N^P_r$, $N^A_r$,
$N^{R'}_r$, $N^{R''}_r$ denote the source rows of the original run. Its source-row bounds read
\begin{equation}\label{9.10:eq-source-bounds-2}
\begin{aligned}
\|N^P_r\|,\ \|N^A_r\| &\le C_\sigma e_r,\qquad
\|N^{R''}_r\|\le \mathfrak{r}_r\Lambda_r C_\sigma e_r,\\
\|N^{R'}_r\| &\le \mathfrak{r}_r\Lambda_r\bigl[(C_\sigma+1)e_r+B_r(G,d)\bigr]
+\varkappa'_r\sum_{r<s<r+\ell_r}\|G^R_s\| .
\end{aligned}
\end{equation}
Here $e_r=\rho^{n-r}+\Pi_n$, and $\Pi_n$ is the polynomial floor of the original run. The rows for
$N^P_r$, $N^A_r$ and $N^{R''}_r$ come from the source-row theorems for the pure slots $P$ and
$P^\flat$, their corollary, and the localisation statement, none of which carries a floor of its own.
The row for $N^{R'}_r$ is obtained
through the chain of estimates for the $R'$-part of the remnant row.

\emph{The origin of the floor.} In the original run, the floor $\Pi_n$ inside $C_\sigma e_r$ has a
single origin. This origin is the one-step index-stationarity defect of an output at depth $r$. That
defect depends on every old remnant present, taken as a remnant, each with its own index error
\[
\iota_s\le C_\iota\vartheta_\chi^{\,n-s}.
\]
Here an object at level $j=n-s$ carries the factor $\vartheta_\chi^{\,j}$, and $\vartheta_\chi\le\rho$.
The defect reads these remnants at flat zone weight. Early remnants are therefore read with an index
error of order one, and they are remembered flatly; only the corollary on the index errors of re-filed
remnants could restore their smallness. In the original run the sum of these flat remnant
terms obeys
\[
\sum_s\hat\vartheta_{n-s}\,\mathfrak{r}_s\le\frac{2C_{\mathrm{cl}}}{\Lambda}\,\Pi_n ,
\]
where $\hat\vartheta_{n-s}$ is the index-error factor carried by an object at level $n-s$ of the
original run, and $C_{\mathrm{cl}}$ and $\Lambda$ are constants of the original run ($\Lambda$ is a
constant, not the factor $\Lambda_r$). Accordingly, in the rate
hypothesis on the source rows of the original run the geometric factor of level $j$ is replaced by
$\vartheta_1^{\,j}+\Pi_{r+j}$ in both of its components, $\vartheta_1$ being the geometric rate of that
hypothesis; this is the origin of the summand $\Pi_n$ in $e_r=\rho^{n-r}+\Pi_n$.

\emph{The finite-range sum.} Under the substitution rule \eqref{9.10:eq-substitution-rule-a}, a
remnant at depth $s$ is read as a remnant only until it is filed. Hence the same sum runs only over
$r<s\le r+\ell_s$. Three inputs bound it.
\begin{itemize}
\item The zone weights $w_{r,s}$ (indexed here by the depths $r$ and $s$ rather than by levels), with
which the defect at depth $r$ reads the remnant at depth $s$,
satisfy $w_{r,s}\le 3$ on the prescribed scaffold $\Theta_X$.
\item For $r<s\le r+\ell_s+1$ one has $\ell_s\le\ell_r+1$, by the depth-regularity of the scaffold and
of the filing ranges.
\item For the same $s$ one has $\mathfrak{r}_s\le 2\mathfrak{r}_r$, by the same two inputs.
\end{itemize}
Consequently every $s$ in the sum satisfies $s\le r+\ell_r+1$. The sum thus has at most
$\ell_r+1\le\ell_r+3$ terms. For each such $s$,
\[
n-s\ge n-r-\ell_r-1\ge n-r-\ell_r-3 ,
\]
so that $\vartheta_\chi^{\,n-s}\le\vartheta_\chi^{\,n-r-\ell_r-3}$ because $\vartheta_\chi\le\rho<1$.
Therefore
\begin{equation}\label{9.10:eq-source-bounds-3}
\begin{aligned}
\sum_{r<s\le r+\ell_s} w_{r,s}\,\mathfrak{r}_s\,\iota_s
&\le 3(\ell_r+3)\cdot 2\mathfrak{r}_r\cdot C_\iota\vartheta_\chi^{\,n-r-\ell_r-3}\\
&= \bigl[6C_\iota(\ell_r+3)\vartheta_\chi^{-\ell_r-3}\mathfrak{r}_r\bigr]\,\vartheta_\chi^{\,n-r}.
\end{aligned}
\end{equation}
The index errors of remnants that have already been filed do not enter this sum. They are carried inside
the $\hat E$-slots, at the weight of the geometric kernel $\varsigma$. This is the corollary on the index
errors of re-filed remnants, applied to differences, with the remnant taken as the difference of the
remnants of the two compared runs.

\emph{The bracket is a ceiling on $g$.} The bracket in \eqref{9.10:eq-source-bounds-3} is bounded
uniformly in $r$ by \eqref{9.10:eq-source-bounds-a}. To see that this is a condition on $g$ alone,
note first that
\[
\vartheta_\chi^{-\ell_r}=(L^2M\check R_{(r)})^{\theta_\chi},
\]
where $\check R_{(r)}$ is the partition-cube radius at the level of depth $r$, and $\theta_\chi>0$ is the
fixed exponent attached to the rate $\vartheta_\chi$ for which this identity holds. This is a fixed power
of $M\check R_{(r)}$, and $M\check R_{(r)}$ is bounded by a constant times
$(\log\theta_r)^{r_{\mathrm{pr}}}$, a fixed power of $\log\theta_r$ whose exponent $r_{\mathrm{pr}}$ is
fixed by the growth of the partition-cube radius. Hence
the left member of \eqref{9.10:eq-source-bounds-a} is a power of $\log\theta_r$ times
$C_t\theta_r^{1-\kappa_0/2}$. It is maximal near $r=0$, where $\theta_0=X=g^{-2}$, and it tends to $0$ as
$g\to0$, because $\kappa_0>2$. So \eqref{9.10:eq-source-bounds-a} is an upper bound on $g$, imposed after
$L$, the auxiliary parameters $\mathfrak{p}$, $\theta_\chi$ and $\kappa_0>2$ are fixed. It places no
condition on $\rho$. It bounds none of the quantities $L$, $M$, $M_1$, $\kappa,\dots$ of
Definition~\ref{2.3:def-cap} from above, and so is not a cap.

\emph{The rows of the re-filed run.} By \eqref{9.10:eq-source-bounds-3} and
\eqref{9.10:eq-source-bounds-a}, read in
the units of $C_\sigma$, the remnant channel contributes at most
\[
C_\sigma\vartheta_\chi^{\,n-r}\le C_\sigma\rho^{n-r}
\]
to the factor $C_\sigma e_r$ of each row of \eqref{9.10:eq-source-bounds-2}, and the non-remnant
channels contribute $C_\sigma\rho^{n-r}$ to that factor. The floor $\Pi_n$, whose
only origin was the flat remnant read, is absent. In all four rows of
\eqref{9.10:eq-source-bounds-2} we may therefore replace $e_r$ by $\hat e_r:=\rho^{n-r}$ and $C_\sigma$
by $\hat C_\sigma=2C_\sigma$ of \eqref{9.10:eq-source-bounds-b}. In the row for $N^{R'}_r$ the bracket
$B_r(G,d)$ of the original run is replaced by the bracket $\hat B_r(G,d)$ of the re-filed run, defined
in \eqref{9.10:eq-three-slot-bounds-a}. This gives \eqref{9.10:eq-source-bounds-1}. The finite-range
factor is the left member of \eqref{9.10:eq-source-bounds-a}, and it is paid by that condition.

\emph{The envelope constant.} The definition of the envelope constant $C_E$ of the original run is read
with $\hat C_\sigma$ in place of $C_\sigma$. This re-values that constant as $\hat C_E$ in
\eqref{9.10:eq-source-bounds-b}.

$(\ast)$ \emph{The reading of remnants by the statements used.} The derivation above applies, to the
objects of the re-filed run, the proofs of the statements that supply the rows of the original run
\eqref{9.10:eq-source-bounds-2}. These statements are:
\begin{itemize}
\item the source-row theorems for the pure slots $P$ and $P^\flat$;
\item the localisation statement used for the floor-free rows;
\item for the row of $N^{R'}_r$, the chain of estimates through the source-carrying form of the volume and
filament bound, the proposition on the primed $\mathbf{B}$-terms, the elimination of the
$\mathbf{B}$-terms, the corollary of the rim estimate and the remark closing that chain.
\end{itemize}
The derivation uses that each of these proofs depends on a remnant only through three quantities: the
third row (the remnant row) of the lemma fixing the reading of a remnant at zone weight; the
$R''$-band; and the groups of the regular family $E^{(j)}$.

For the correlation theorem, the two statements of the one-lattice comparison and the further comparison
statement used with them, the lemma listing the quantities through which their proofs depend on
remnants gives exactly this. For the statements of this list themselves, the statements of their
outputs name no other such quantity. The remnant row there is $C_t\theta_s(R'+R'')$, with $R'$ and
$R''$ its two parts. The stored
units are the list slots $E^{(j)}$, $P^{\flat(j)}$, $Q^{(j)}$, $R^{(j)}$ and $D^{(j)}$, together with
stored boundary terms; the source-carrying form of the bound on stored terms names no other such
quantity either.

That the proofs of the statements of this list, read on the objects of the re-filed run, depend on
remnants only through the three quantities named above is the step $(\ast)$. It is not carried out here.
\end{proof}

\begin{definition}[Envelope letters and coupling ceilings]\label{9.10:def-envelope-ceilings}
Depths $r,s\ge 0$ are counted from $n$ as in Notation~\ref{9.10:not-filing-level}: depth $r$ is
level $n-r$, and $\theta_r=X+c_0r$ with $X=g^{-2}$ is the prescribed scaffold fixed there. The
letter $\rho$, with $0<\rho<1$, is the contraction rate of the comparison of the two runs,
chosen in its window. The following letters are those of the one-step bounds of
Corollary~\ref{9.10:cor-two-slot-bounds}:
\begin{itemize}
\item the coefficients $\alpha_r$, $\alpha^A_r$, $\varkappa$, $\varkappa'_r$, $\varkappa''_r$,
$\varkappa_A$ and $C_F$;
\item the Lipschitz constant $\mathsf{A}$ of the coupling inequality;
\item the geometric rate $q$;
\item the constant $C_t$ and the number $\kappa_0$ entering the decay factor $\mathfrak{r}_r$,
and the constant $\Lambda$ and the exponent $c_3$ of the logarithmic factor $\Lambda_r$;
$\mathfrak{r}_r$ and $\Lambda_r$ are defined in the next display.
\end{itemize}
Here $\check R_{(r)}$ denotes the partition-cube radius at level $n-r$, and $M=L^{m'}$. We put
\begin{equation*}
\mathsf{n}_r:=r+1+(1-\rho)^{-1},\qquad c_{\rho}:=\frac{5}{\rho(1-\rho)},\qquad
\mathfrak{r}_r:=C_t\,\theta_r^{1-\kappa_0/2},\qquad \Lambda_r:=\Lambda(\log\theta_r)^{c_3},
\end{equation*}
and
\begin{equation*}
\ell_r:=2+\log_L\bigl(M\check R_{(r)}\bigr)=\mathsf{a}_r+2 .
\end{equation*}
This $\ell_r$ is the depth-range letter of \eqref{9.10:eq-filing-level-a}.

Let $\hat C_E$ be the constant of \eqref{9.10:eq-source-bounds-b}. It is fixed with the geometric
source bounds for the re-filed run, Proposition~\ref{9.10:prop-source-bounds} (stated; proof
incomplete). The \emph{envelope letters} and the \emph{envelopes} are
\begin{equation}\label{9.10:eq-envelope-ceilings-a}
\begin{gathered}
\hat e_r:=\rho^{\,n-r},\qquad
\hat{\mathfrak{r}}_r:=2\,\mathfrak{r}_r\Lambda_r\Bigl(\mathsf{A}\Bigl[\frac{\mathsf{n}_r}{1-\rho}
+\frac{1}{(1-\rho)^{2}}\Bigr]+\frac{5}{1-\rho}+8\Bigr),\\
E^P_r:=\hat C_E\,\hat e_r,\qquad E^A_r:=4\hat C_E\,\hat e_r,\qquad
E^R_r:=4\hat C_E\,\hat{\mathfrak{r}}_r\,\hat e_r .
\end{gathered}
\end{equation}
Thus $\hat e_r$ is the envelope letter $e_r=\rho^{n-r}+\Pi_n$ of the un-hatted comparison of the
two runs with the polynomial floor $\Pi_n\equiv0$; neither $\Pi_n$ nor the constant $C_T$ of the
un-hatted growth condition $\hat{\mathfrak{r}}_s\le C_T\rho^{-(n-s)/2}\hat{\mathfrak{r}}_n$ enters
the conditions below.

The \emph{ceiling conditions} are conditions (a)--(g$'$) and (b$''$) of the display below. They
are imposed for every depth $r$, and in (f) for every $s\ge r$:
\begin{equation}\label{9.10:eq-envelope-ceilings-b}
\begin{aligned}
&\text{(a)}\quad c_{\rho}\,\alpha_r\,\mathsf{A}\,\mathsf{n}_r\le\tfrac18;\qquad
\text{(b)}\quad \varkappa\le\tfrac{1-\rho}{32};\qquad
\text{(b$'$)}\quad 4\varkappa''_r\le\tfrac{3}{32}(1-\rho);\\
&\text{(c)}\quad c_{\rho}\,\alpha^A_r\,\mathsf{A}\,\mathsf{n}_r\le\tfrac14;\qquad
\text{(d)}\quad \varkappa_A\Bigl(9+\frac{5}{1-\rho}\Bigr)\le\tfrac14;\qquad
\text{(e)}\quad \hat{\mathfrak{r}}_r\le\tfrac1{16};\\
&\text{(f)}\quad \hat{\mathfrak{r}}_s\le 2(1+s-r)\,\hat{\mathfrak{r}}_r\quad(s\ge r);\qquad
\text{(g$'$)}\quad 8\,\varkappa'_r\,\ell_r^{2}\rho^{-\ell_r}\le\tfrac14;\\
&\text{(b$''$)}\quad 8\,(\varkappa+\varkappa_A+C_F+1)\bigl(1+(1-\rho)^{-1}\bigr)\,\mathfrak{S}
\le\tfrac1{16};\\
&\qquad\quad \mathfrak{S}:=\sup_{s\ge0}\,(\ell_s+2)^2\,\hat{\mathfrak{r}}_s\,\rho^{-\ell_s};\\
&\text{for every } r:\quad
\sum_{s>r}(\ell_s+2)\,q^{(s-r-\ell_s)_+}\,\hat{\mathfrak{r}}_s\,\rho^{-(s-r)}
\le\mathfrak{S}\bigl(1+(1-\rho)^{-1}\bigr).
\end{aligned}
\end{equation}
Condition (g$'$) is the condition $4\varkappa'_r\ell_r^2\rho^{-\ell_r}\le\frac14$ of the
envelope conditions of the un-hatted comparison, read with $2\varkappa'_r$ in place of
$\varkappa'_r$.

The ceiling conditions are imposed on $g$ in the order of choice
(Notation~\ref{2.3:not-stages-first}): they are required of $g$ after $L$, $\mathfrak{p}$, the
further constants of the order of choice that precede $\rho$, and $\rho$ itself have been fixed,
and $\rho$, a letter of its window, is not restricted by them. For (b$''$) this is part (ii) of
the lemma that follows this definition. For (a)--(f) it is the corresponding property of the
envelope conditions of the un-hatted comparison of the two runs, of which they are members; for
(g$'$) it is that property of the member with $\varkappa'_r$, which doubling the coefficient
does not alter.

The last line of \eqref{9.10:eq-envelope-ceilings-b} is not one of the ceiling conditions: it is
the inequality through which (b$''$) is used. It is proved in part (i) of the lemma that follows
this definition from three properties of the letters: the bound $q\le\rho^2$ on the geometric
rate, the increment bound $\ell_s\le\ell_r+1$ for $r<s\le r+\ell_s$, and the monotonicity
$\ell_r\le\ell_s$ for $s>r$ of the depth-range letters of \eqref{9.10:eq-filing-level-a}.
\end{definition}

\begin{lemma}[The summation inequality for \textup{(b$''$)}, and \textup{(b$''$)} as a ceiling
on $g$]
Let the letters be as in Definition~\ref{9.10:def-envelope-ceilings}.
\begin{enumerate}
\item[(i)] Assume that $q\le\rho^2$, that the depth-range letters are non-decreasing in the
depth, $\ell_r\le\ell_s$ for $s>r$, and that the depth-range letters $\ell_s$ of
\eqref{9.10:eq-filing-level-a} satisfy the increment bound $\ell_s\le\ell_r+1$ whenever
$r<s\le r+\ell_s$. Then the last line of \eqref{9.10:eq-envelope-ceilings-b} holds for every
depth $r$.
\item[(ii)] Fix $L$, $\mathfrak{p}$, $\rho$ in its window, and $\kappa_0>4$, all fixed before
$g$ in the order of choice; $c_0>0$ is the fixed constant of the scaffold of
Notation~\ref{9.10:not-filing-level}. Assume that $\varkappa$, $\varkappa_A$, $C_F$,
$\mathsf{A}$, $C_t$, $\Lambda$ and $c_3$ are bounded,
for all sufficiently small $g$, by quantities fixed before $g$ in the order of choice. Then
(b$''$) holds for all $g$ below a threshold fixed after $L$, $\mathfrak{p}$, $\rho$ and
$\kappa_0$.
\end{enumerate}
\end{lemma}

\begin{proof}[Proof of part \textup{(i)}]
Fix a depth $r$. We split the sum over $s>r$ into two parts: the depths $s$ with
$r<s\le r+\ell_s$, and the depths $s$ with $s>r+\ell_s$.

\emph{Near depths.} Let $r<s\le r+\ell_s$. Then $(s-r-\ell_s)_+=0$, so $q$ enters to the power
zero. Since $0<\rho<1$ and $s-r\le\ell_s$, we have $\rho^{-(s-r)}\le\rho^{-\ell_s}$. The term of
index $s$ is therefore at most $(\ell_s+2)\hat{\mathfrak{r}}_s\rho^{-\ell_s}$.

We now count these depths. For $r<s\le r+\ell_s$ the increment bound assumed in part (i)
gives $\ell_s\le\ell_r+1$. Hence every such $s$ satisfies $r<s\le r+\ell_r+1$, so there are at
most $\ell_r+1$ such depths, and in particular at most $\ell_r+2$. Since $\ell_r\le\ell_s$ for
$s>r$ by the monotonicity assumed in part (i), we have
$\ell_r+2\le\ell_s+2$, and the sum of these terms is at most
\begin{equation*}
(\ell_r+2)\max_{s}\,(\ell_s+2)\,\hat{\mathfrak{r}}_s\,\rho^{-\ell_s}
\le\sup_{s\ge0}\,(\ell_s+2)^2\,\hat{\mathfrak{r}}_s\,\rho^{-\ell_s}=\mathfrak{S}.
\end{equation*}
Here the maximum is over the depths $s$ with $r<s\le r+\ell_s$.

\emph{Far depths.} Let $s>r+\ell_s$. Then $(s-r-\ell_s)_+=s-r-\ell_s$. Writing
$\rho^{-(s-r)}=\rho^{-\ell_s}\rho^{-(s-r-\ell_s)}$, the term of index $s$ equals
\begin{equation*}
(\ell_s+2)\,\hat{\mathfrak{r}}_s\,\rho^{-\ell_s}\Bigl(\frac{q}{\rho}\Bigr)^{s-r-\ell_s}.
\end{equation*}
Since $q\le\rho^2$ by hypothesis, we have $q/\rho\le\rho$. The term is therefore at most
\begin{equation*}
(\ell_s+2)\,\hat{\mathfrak{r}}_s\,\rho^{-\ell_s}\,\rho^{\,s-r-\ell_s}
\le\frac{\mathfrak{S}}{\ell_s+2}\,\rho^{\,s-r-\ell_s}
\le\tfrac12\,\mathfrak{S}\,\rho^{\,s-r-\ell_s},
\end{equation*}
because $(\ell_s+2)^2\hat{\mathfrak{r}}_s\rho^{-\ell_s}\le\mathfrak{S}$ and
$\ell_s+2\ge4\ge2$, since $\ell_s=\mathsf{a}_s+2\ge2$ with $\mathsf{a}_s\in\mathbb{N}$ by
Notation~\ref{9.10:not-filing-level}.

The exponents $u=s-r-\ell_s$ are positive integers, and each value of $u$ is taken by at most
two far depths. Indeed, let $s<s'$ with $s-\ell_s=s'-\ell_{s'}$. Then
$\ell_{s'}=\ell_s+(s'-s)\ge s'-s$, so $s<s'\le s+\ell_{s'}$, and the increment bound with $r$
replaced by $s$ and $s$ replaced by $s'$ gives $\ell_{s'}\le\ell_s+1$, that is, $s'-s\le1$.
Three distinct depths with the same value of $s-\ell_s$ would contain two at distance at least
two, which is excluded. The sum over the far depths is therefore at most
\begin{equation*}
\tfrac12\,\mathfrak{S}\cdot2\sum_{u\ge1}\rho^{u}=\mathfrak{S}\,\frac{\rho}{1-\rho}
\le\mathfrak{S}(1-\rho)^{-1}.
\end{equation*}

Adding the two parts gives the bound $\mathfrak{S}+\mathfrak{S}(1-\rho)^{-1}$, which is the last
line of \eqref{9.10:eq-envelope-ceilings-b}.
\end{proof}

\begin{proof}[Proof of part \textup{(ii)}]
Fix $L$, $\mathfrak{p}$, $\rho$ and $\kappa_0$ with $\kappa_0>4$.

\emph{The factors carrying $\ell_s$.} By definition, $L^{\ell_s}=L^2M\check R_{(s)}$. Hence
\begin{equation*}
\rho^{-\ell_s}=\bigl(L^{2}M\check R_{(s)}\bigr)^{\log_L\rho^{-1}},
\end{equation*}
which is a fixed power of $M\check R_{(s)}$ at fixed $\rho<1$. The partition-cube radius, whose
exponents are fixed in Notation~\ref{9.1:not-radius-exponent-floors}, satisfies
$M\check R_{(s)}\le C\,M^{c}(\log\theta_s)^{c'}$ for a constant $C$ and exponents $c,c'$ that do
not depend on $s$ or $g$. Consequently $\rho^{-\ell_s}$ is bounded by a polynomial in
$\log\theta_s$ whose coefficients depend only on $L$, $\mathfrak{p}$ and $\rho$. The same holds
for $(\ell_s+2)^2$, since
\begin{equation*}
(\ell_s+2)^2=\bigl(4+\log_L(M\check R_{(s)})\bigr)^2 .
\end{equation*}

\emph{The factor $\hat{\mathfrak{r}}_s$.} By \eqref{9.10:eq-envelope-ceilings-a},
\begin{equation*}
\hat{\mathfrak{r}}_s=2C_t\,\theta_s^{1-\kappa_0/2}\,\Lambda(\log\theta_s)^{c_3}
\Bigl(\mathsf{A}\Bigl[\frac{\mathsf{n}_s}{1-\rho}+\frac{1}{(1-\rho)^{2}}\Bigr]
+\frac{5}{1-\rho}+8\Bigr).
\end{equation*}
The bracket is an affine function of $\mathsf{n}_s$ with coefficients fixed by $\mathsf{A}$ and
$\rho$. Since $\theta_s=X+c_0s$, we have $s=(\theta_s-X)/c_0\le\theta_s/c_0$. Therefore
\begin{equation*}
\mathsf{n}_s=s+1+\frac{1}{1-\rho}\le\frac{\theta_s}{c_0}+1+\frac{1}{1-\rho} .
\end{equation*}
For $g\le1$ we have $\theta_s\ge X\ge1$, so the right side is at most
$\bigl(c_0^{-1}+1+(1-\rho)^{-1}\bigr)\theta_s$.
Hence $\hat{\mathfrak{r}}_s$ is bounded by a polynomial in $\log\theta_s$ times an affine
function of $\mathsf{n}_s$ times $\theta_s^{1-\kappa_0/2}$. It is therefore bounded by a
polynomial in $\log\theta_s$ times $\theta_s^{2-\kappa_0/2}$.

\emph{The supremum.} Combining the three factors, there is a polynomial $\mathcal{P}$ such that
\begin{equation*}
(\ell_s+2)^2\,\hat{\mathfrak{r}}_s\,\rho^{-\ell_s}\le\mathcal{P}(\log\theta_s)\,
\theta_s^{\,2-\kappa_0/2}
\qquad\text{for every }s\ge0 .
\end{equation*}
The coefficients of $\mathcal{P}$ are fixed by $L$, $\mathfrak{p}$, $\rho$, $\kappa_0$, $c_0$ and
the bounds assumed in part (ii); $\mathcal{P}$ does not depend on $s$ or $g$. Since
$\theta_s\ge\theta_0=X$ for every $s$,
\begin{equation*}
\mathfrak{S}\le\sup_{\theta\ge X}\,\mathcal{P}(\log\theta)\,\theta^{\,2-\kappa_0/2}.
\end{equation*}
For $\kappa_0>4$ the exponent $2-\kappa_0/2$ is negative, so
$\mathcal{P}(\log\theta)\theta^{2-\kappa_0/2}\to0$ as $\theta\to\infty$. The supremum on the
right is therefore finite, is attained at a bounded value of $\theta\ge X$, and tends to $0$ as
$X=g^{-2}\to\infty$, that is, as $g\to0$.

\emph{Conclusion.} By the hypothesis of part (ii), the prefactor
$8(\varkappa+\varkappa_A+C_F+1)(1+(1-\rho)^{-1})$ of (b$''$) is bounded, for all sufficiently
small $g$, by a quantity fixed before $g$. Consequently (b$''$) holds for all $g$ below a
threshold that is fixed after
$L$, $\mathfrak{p}$, $\rho$ and $\kappa_0$, so (b$''$) is a ceiling on $g$. For every fixed
$\rho<1$ in its window it is met by taking $g$ small, so it places no condition on $\rho$ beyond
that window. It bounds none of the letters chosen before $g$, so it is not a cap in the sense of
Definition~\ref{2.3:def-cap}.
\end{proof}

\begin{proposition}[Geometric envelopes for the re-filed run]\label{9.10:prop-geometric-envelopes}
Let $n$ be the length of the re-filed run, that is, of the hatted run of
Corollary~\ref{9.10:cor-scaffold-inputs}. Let $\hat e_r=\rho^{n-r}$, $\mathsf{n}_r$, $c_{\rho}$,
$\mathfrak{r}_r$, $\Lambda_r$, $\hat{\mathfrak{r}}_r$, $\ell_r$ and the envelopes $E^P_r$, $E^A_r$, $E^R_r$
be the envelope letters of \eqref{9.10:eq-envelope-ceilings-a}, and let $\hat C_{\sigma}$, $\hat C_E$ be the
constants of \eqref{9.10:eq-source-bounds-b}. Assume the following.
\begin{itemize}
\item Clause (c3) of Lemma~\ref{9.6:lem-chain-structure}, for the comparison chain of
Definition~\ref{9.6:def-comparison-chain}, and the maximum principle for coupled scale-indexed
inequalities, Lemma~\ref{9.6:lem-maximum-principle}.
\item The input (J1) for the hatted run, as delivered by Corollary~\ref{9.10:cor-scaffold-inputs}.
\item The three-slot hybrid bounds \eqref{9.10:eq-three-slot-bounds-a} of
Corollary~\ref{9.10:cor-three-slot-bounds}.
\item The geometric source bounds of Proposition~\ref{9.10:prop-source-bounds}, whose proof is
incomplete at one step.
\item The coupling ceilings (a), (b), (b$'$), (c), (d), (e), (f), (g$'$) and (b$''$) of
\eqref{9.10:eq-envelope-ceilings-b} in Definition~\ref{9.10:def-envelope-ceilings}, together with the
inequality attached there to (b$''$).
\item The geometric rate $q$ of \eqref{9.10:eq-three-slot-bounds-a} and the rate $\rho$ of the
envelope letters satisfy $q\le\rho^2$.
\end{itemize}
Let $\underline{r}$ be an integer with $0\le\underline{r}\le n$. Consider a pair of hatted lists,
frozen or local, both alive
to depth $\underline{r}$, whose running shifts at depth $\underline{r}$ agree on $\mathfrak{P}$, the set
on which the running shifts of the two lists are compared,
\begin{equation*}
\hat\zeta^{\mathsf{B}}_{(\underline{r})}=\hat\zeta^{\mathsf{A}}_{(\underline{r})}\quad\text{on }\mathfrak{P}
\end{equation*}
(complex deviations of the running shifts within the complex tube admitted for hatted lists being
allowed). Let $g^P_r$, $g^A_r$, $g^R_r$ be the unknowns and
$d_{r+1}$, $d^{\sharp}_{r+1}$ the coupling differences attached to the pair by its comparison chain
(Definition~\ref{9.6:def-comparison-chain}). In the display below, $\mathsf{A}$ denotes the Lipschitz
constant of the coupling row of \eqref{9.10:eq-three-slot-bounds-a}.
Then
\begin{equation}\label{9.10:eq-geometric-envelopes-a}
\begin{aligned}
&g^P_r\le\hat C_E\hat e_r,\qquad g^A_r\le4\hat C_E\hat e_r,\qquad
g^R_r\le4\hat C_E\hat{\mathfrak{r}}_r\hat e_r &&(\underline{r}\le r\le n),\\
&d_{r+1}\le\frac{\mathsf{A}\hat C_E\rho^{n-r}}{1-\rho}\le\mathsf{A}\hat C_E\mathsf{n}_r\hat e_r,
\qquad d^{\sharp}_{r+1}\le c_{\rho}\mathsf{A}\hat C_E\mathsf{n}_r\hat e_r &&(\underline{r}\le r<n).
\end{aligned}
\end{equation}
Hence the bounds of the first line of \eqref{9.10:eq-geometric-envelopes-a} hold also for the
class-$c$ norms $\mathsf{D}^c_r$ of the chain differences at depth $r$.
\end{proposition}

\begin{proof}
The proof has two steps. First we write
the system of coupled inequalities satisfied by the unknowns. Then we apply the maximum principle
of Lemma~\ref{9.6:lem-maximum-principle}, row by row.

\emph{Step one: the system.} Clause (c3) of Lemma~\ref{9.6:lem-chain-structure} applies to hatted
lists, since it reads only the definition of the comparison chain and the bounds of the class of
lists, and these bounds hold for the class $\hat{\mathfrak{K}}_k$ of hatted lists. We combine it with
the three-slot hybrid bounds
\eqref{9.10:eq-three-slot-bounds-a},
with their coefficients $\alpha_r$, $\alpha^A_r$, $\varkappa$, $\varkappa'_r$, $\varkappa''_r$,
$\varkappa_A$, $C_F$, in which the competitor norms $\|G^{\bullet}_s\|$ are replaced by
the unknowns $g^{\bullet}_s$. The geometric source bounds of Proposition~\ref{9.10:prop-source-bounds}
(whose proof is incomplete at the step marked $(\ast)$ there) contribute the terms
$\hat C_{\sigma}\hat e_r$ and $\mathfrak{r}_r\Lambda_r(2\hat C_{\sigma}+1)\hat e_r$. We also use the
inequality $g^2\le1$ for the reference coupling. Together these
give, for $\underline{r}\le r<n$,
\begin{align}
g^P_r&\le\alpha_rd^{\sharp}_{r+1}+\varkappa\sum_{s>r}q^{s-r}g^P_s
+\varkappa''_r\sum_{s>r}q^{s-r}g^A_s+C_F\mathsf{S}^+_r(g^R)+\hat C_{\sigma}\hat e_r,
\label{9.10:eq-geometric-envelopes-1}\\
g^A_r&\le\alpha^A_rd^{\sharp}_{r+1}
+\varkappa_A\Bigl[g^P_r+\sum_{s>r}q^{s-r}\bigl(g^P_s+g^A_s\bigr)+\mathsf{S}^+_r(g^R)\Bigr]
\notag\\
&\qquad+g^R_r+C_F\mathsf{S}^+_r(g^R)+\hat C_{\sigma}\hat e_r,
\label{9.10:eq-geometric-envelopes-2}\\
g^R_r&\le2\mathfrak{r}_r\Lambda_r\hat B_r(g^{\bullet},d)+2\varkappa'_r\sum_{r<s<r+\ell_r}g^R_s
+\mathfrak{r}_r\Lambda_r\bigl(2\hat C_{\sigma}+1\bigr)\hat e_r.
\label{9.10:eq-geometric-envelopes-3}
\end{align}
Here $\mathsf{S}^+_r(g^R)$ is the finite-range remnant load of \eqref{9.10:eq-two-slot-bounds-a}, formed
with the $g^R_s$, and $\hat B_r(g^{\bullet},d)$ is the bracket of \eqref{9.10:eq-three-slot-bounds-a},
formed with
the $g^{\bullet}_s$ and the $d_{s+1}$. The coupling row of \eqref{9.10:eq-three-slot-bounds-a} gives,
for $\underline{r}\le s<n$,
\begin{equation}\label{9.10:eq-geometric-envelopes-4}
d_{s+1}\le\mathsf{A}\sum_{\underline{r}\le u\le s}g^P_u .
\end{equation}

\emph{Step two: the maximum principle.} We apply Lemma~\ref{9.6:lem-maximum-principle} with index set
\begin{equation*}
I:=\{(r,\bullet):\underline{r}\le r\le n,\ \bullet\in\{P,A,R\}\},
\qquad I_0:=\{(r,\bullet)\in I:r<n\}.
\end{equation*}
The envelopes are $E=(E^{\bullet}_r)$ of \eqref{9.10:eq-envelope-ceilings-a}. For $(r,\bullet)\in I_0$
we split the right member of \eqref{9.10:eq-geometric-envelopes-1}--\eqref{9.10:eq-geometric-envelopes-3}
into two parts:
\begin{itemize}
\item the part $\mathcal{L}^{\bullet}_r[E]$, in which every unknown $g^{\bullet'}_s$ is replaced by
$E^{\bullet'}_s$ and every coupling difference by the bound that \eqref{9.10:eq-geometric-envelopes-4}
yields from $E$;
\item the source term $\mathfrak{u}^{\bullet}_r$, which is the last term of each line.
\end{itemize}
At every index of $I_0$ we show that $\mathcal{L}^{\bullet}_r[E]\le\frac12E^{\bullet}_r$ and
$\mathfrak{u}^{\bullet}_r\le\frac12E^{\bullet}_r$. At the indices of $I\setminus I_0$ we show that
$g^{\bullet}_n\le E^{\bullet}_n$.

\emph{Geometric sums.} We use the hypothesis $q\le\rho^2$. Since $\hat e_r=\rho^{n-r}$, we have exactly
$\hat e_{r+m}=\rho^{-m}\hat e_r$. Hence $q^{s-r}\hat e_s=(q/\rho)^{s-r}\hat e_r\le\rho^{s-r}\hat e_r$.
Summing the geometric series gives
\begin{equation}\label{9.10:eq-geometric-envelopes-5}
\sum_{s\ge r}q^{s-r}\hat e_s\le\frac{\hat e_r}{1-\rho},\qquad
\sum_{s>r}q^{s-r}\hat e_s\le\frac{\rho\,\hat e_r}{1-\rho}.
\end{equation}
Next, $\sum_{\underline{r}\le u\le r}\rho^{n-u}\le\rho^{n-r}\sum_{j\ge0}\rho^j$. Moreover
$(1-\rho)^{-1}\le\mathsf{n}_r=r+1+(1-\rho)^{-1}$. Therefore
\begin{equation}\label{9.10:eq-geometric-envelopes-6}
\sum_{\underline{r}\le u\le r}\hat e_u\le\frac{\rho^{n-r}}{1-\rho}\le\mathsf{n}_r\hat e_r .
\end{equation}
For the remnant channel we insert $E^R_s=4\hat C_E\hat{\mathfrak{r}}_s\hat e_s$ and
$\hat e_s=\rho^{-(s-r)}\hat e_r$. We then apply the inequality attached to (b$''$) in
\eqref{9.10:eq-envelope-ceilings-b}, with $\mathfrak{S}=\sup_{s\ge0}(\ell_s+2)^2\hat{\mathfrak{r}}_s\rho^{-\ell_s}$.
This gives
\begin{equation}\label{9.10:eq-geometric-envelopes-7}
\begin{aligned}
\sum_{s>r}(\ell_s+2)q^{(s-r-\ell_s)_+}E^R_s
&=4\hat C_E\hat e_r\sum_{s>r}(\ell_s+2)q^{(s-r-\ell_s)_+}\hat{\mathfrak{r}}_s\rho^{-(s-r)}\\
&\le4\hat C_E\hat e_r\,\mathfrak{S}\bigl(1+(1-\rho)^{-1}\bigr).
\end{aligned}
\end{equation}
By the definition of the finite-range remnant load in \eqref{9.10:eq-two-slot-bounds-a}, when the
unknowns are replaced by the envelopes, $\mathsf{S}^+_r$ is bounded by the
left member of \eqref{9.10:eq-geometric-envelopes-7}. This is how \eqref{9.10:eq-geometric-envelopes-7}
enters below.

\emph{Coupling differences.} Suppose $g^P_u\le E^P_u=\hat C_E\hat e_u$ for $\underline{r}\le u\le r$. Then
\eqref{9.10:eq-geometric-envelopes-4} and \eqref{9.10:eq-geometric-envelopes-6} give
\begin{equation}\label{9.10:eq-geometric-envelopes-8}
d_{r+1}\le\mathsf{A}\hat C_E\sum_{\underline{r}\le u\le r}\hat e_u
\le\frac{\mathsf{A}\hat C_E\rho^{n-r}}{1-\rho}\le\mathsf{A}\hat C_E\mathsf{n}_r\hat e_r .
\end{equation}
The bound
\begin{equation}\label{9.10:eq-geometric-envelopes-9}
d^{\sharp}_{r+1}\le c_{\rho}\mathsf{A}\hat C_E\mathsf{n}_r\hat e_r
\end{equation}
is the corresponding geometric sum of the envelope proposition for the run without re-filing;
its derivation there uses only the following three facts, each of which holds here:
\begin{itemize}
\item inequality \eqref{9.10:eq-geometric-envelopes-6};
\item the monotonicity of $r\mapsto\hat e_r=\rho^{n-r}$, which is non-decreasing because $\rho<1$;
\item the inequality $\mathsf{n}_{r+1}=\mathsf{n}_r+1\le\frac32\mathsf{n}_r$, which holds because
$\mathsf{n}_r\ge2$.
\end{itemize}
In the rows below, $d^{\sharp}_{r+1}$ is replaced by the right member of
\eqref{9.10:eq-geometric-envelopes-9}.

For the $R$-row we also need a weighted form of \eqref{9.10:eq-geometric-envelopes-5}. As above,
$q\le\rho^2$ gives $q^{s-r}\hat e_s\le\rho^{s-r}\hat e_r$, and the definition of $\mathsf{n}_r$ gives
$\mathsf{n}_{r+j}=\mathsf{n}_r+j$. Using $\sum_{j\ge0}\rho^j=(1-\rho)^{-1}$ and
$\sum_{j\ge0}j\rho^j=\rho(1-\rho)^{-2}$, we obtain
\begin{equation}\label{9.10:eq-geometric-envelopes-10}
\sum_{s\ge r}q^{s-r}\mathsf{n}_s\hat e_s\le\hat e_r\sum_{j\ge0}\rho^j\bigl(\mathsf{n}_r+j\bigr)
=\Bigl(\frac{\mathsf{n}_r}{1-\rho}+\frac{\rho}{(1-\rho)^2}\Bigr)\hat e_r .
\end{equation}
By \eqref{9.10:eq-geometric-envelopes-8}, applied at each $s$, the coupling differences $d_{s+1}$
entering the bracket $\hat B_r$ are bounded by $\mathsf{A}\hat C_E\mathsf{n}_s\hat e_s$, and
\eqref{9.10:eq-geometric-envelopes-10} sums these bounds against the geometric weights $q^{s-r}$.

\emph{The $P$-row.} We bound each term of $\mathcal{L}^P_r[E]$ in turn.
\begin{itemize}
\item The term $\alpha_rd^{\sharp}_{r+1}$ contributes $c_{\rho}\alpha_r\mathsf{A}\mathsf{n}_r\hat C_E\hat e_r$.
\item By the second inequality of \eqref{9.10:eq-geometric-envelopes-5}, the term in $\varkappa$
contributes at most $\varkappa\rho(1-\rho)^{-1}\hat C_E\hat e_r$.
\item Since $E^A_s=4\hat C_E\hat e_s$, the term in $\varkappa''_r$ contributes at most
$4\varkappa''_r\rho(1-\rho)^{-1}\hat C_E\hat e_r$.
\item The remnant term contributes at most $C_F$ times the right member of
\eqref{9.10:eq-geometric-envelopes-7}.
\end{itemize}
Hence
\begin{equation*}
\begin{aligned}
\mathcal{L}^P_r[E]
&\le\hat C_E\hat e_r\Bigl[c_{\rho}\alpha_r\mathsf{A}\mathsf{n}_r+\frac{\varkappa\rho}{1-\rho}
+\frac{4\varkappa''_r\rho}{1-\rho}+4C_F\mathfrak{S}\bigl(1+(1-\rho)^{-1}\bigr)\Bigr]\\
&\le\hat C_E\hat e_r\Bigl[\tfrac18+\tfrac1{32}+\tfrac3{32}+\tfrac1{32}\Bigr]
=\tfrac{9}{32}\hat C_E\hat e_r\le\tfrac12E^P_r .
\end{aligned}
\end{equation*}
The four numbers in the bracket come from the ceilings as follows; ceiling (b$''$) reads
\begin{equation*}
8(\varkappa+\varkappa_A+C_F+1)\bigl(1+(1-\rho)^{-1}\bigr)\mathfrak{S}\le\tfrac1{16}.
\end{equation*}
\begin{itemize}
\item Ceiling (a) gives $c_{\rho}\alpha_r\mathsf{A}\mathsf{n}_r\le\frac18$.
\item Ceiling (b), $\varkappa\le(1-\rho)/32$, gives $\varkappa\rho/(1-\rho)\le\frac1{32}$.
\item Ceiling (b$'$), $4\varkappa''_r\le\frac3{32}(1-\rho)$, gives $4\varkappa''_r\rho/(1-\rho)\le\frac3{32}$.
\item Ceiling (b$''$) gives in particular $4C_F\mathfrak{S}(1+(1-\rho)^{-1})\le\frac1{32}$.
\end{itemize}
The source term satisfies
\begin{equation*}
\mathfrak{u}^P_r=\hat C_{\sigma}\hat e_r\le\tfrac12\hat C_E\hat e_r=\tfrac12E^P_r,
\end{equation*}
because $\hat C_E\ge2\hat C_{\sigma}$ by \eqref{9.10:eq-source-bounds-b}.

\emph{The $A$-row.} We bound each term of $\mathcal{L}^A_r[E]$ in turn.
\begin{itemize}
\item The term $\alpha^A_rd^{\sharp}_{r+1}$ contributes $c_{\rho}\alpha^A_r\mathsf{A}\mathsf{n}_r\hat C_E\hat e_r$.
\item Inside the bracket multiplied by $\varkappa_A$, the term $E^P_r$ contributes $\hat C_E\hat e_r$.
By the second inequality of \eqref{9.10:eq-geometric-envelopes-5}, the sum $\sum_{s>r}q^{s-r}(E^P_s+E^A_s)$
contributes at most $5\rho(1-\rho)^{-1}\hat C_E\hat e_r$.
\item The unknown $g^R_r$ becomes $E^R_r=4\hat{\mathfrak{r}}_r\hat C_E\hat e_r$.
\item The two remnant loads contribute at most $(\varkappa_A+C_F)$ times the right member of
\eqref{9.10:eq-geometric-envelopes-7}.
\end{itemize}
Hence
\begin{equation*}
\begin{aligned}
\mathcal{L}^A_r[E]
&\le\hat C_E\hat e_r\Bigl[c_{\rho}\alpha^A_r\mathsf{A}\mathsf{n}_r
+\varkappa_A\Bigl(1+\frac{5\rho}{1-\rho}\Bigr)+4\hat{\mathfrak{r}}_r
+4(\varkappa_A+C_F)\mathfrak{S}\bigl(1+(1-\rho)^{-1}\bigr)\Bigr]\\
&\le\hat C_E\hat e_r\Bigl[\tfrac14+\tfrac14+\tfrac14+\tfrac1{32}\Bigr]
\le2\hat C_E\hat e_r=\tfrac12E^A_r .
\end{aligned}
\end{equation*}
The four numbers in the bracket come from the ceilings as follows.
\begin{itemize}
\item Ceiling (c) gives $c_{\rho}\alpha^A_r\mathsf{A}\mathsf{n}_r\le\frac14$.
\item Ceiling (d), $\varkappa_A(9+5/(1-\rho))\le\frac14$, bounds the second term, since
$1+5\rho/(1-\rho)\le9+5/(1-\rho)$.
\item Ceiling (e), $\hat{\mathfrak{r}}_r\le\frac1{16}$, gives $4\hat{\mathfrak{r}}_r\le\frac14$.
\item Ceiling (b$''$) gives $4(\varkappa_A+C_F)\mathfrak{S}(1+(1-\rho)^{-1})\le\frac1{32}$.
\end{itemize}
The source term satisfies
\begin{equation*}
\mathfrak{u}^A_r=\hat C_{\sigma}\hat e_r\le\tfrac12\hat C_E\hat e_r\le\tfrac12E^A_r .
\end{equation*}

\emph{The $R$-row.} First consider the bracket $\hat B_r$ evaluated at the envelopes and at the bounds
\eqref{9.10:eq-geometric-envelopes-8}, \eqref{9.10:eq-geometric-envelopes-9}. We use
\eqref{9.10:eq-geometric-envelopes-5}, \eqref{9.10:eq-geometric-envelopes-6},
\eqref{9.10:eq-geometric-envelopes-7} and \eqref{9.10:eq-geometric-envelopes-10}, and we include the
entries with $s=r$ of the $P$- and $A$-slots
in the term $5/(1-\rho)$. This gives
\begin{equation*}
\begin{aligned}
2\mathfrak{r}_r\Lambda_r\hat B_r(E)
&\le2\mathfrak{r}_r\Lambda_r\hat C_E\hat e_r\Bigl[\mathsf{A}\Bigl(\frac{\mathsf{n}_r}{1-\rho}
+\frac{\rho}{(1-\rho)^2}\Bigr)+\frac5{1-\rho}+4\mathfrak{S}\bigl(1+(1-\rho)^{-1}\bigr)\Bigr]\\
&\le2\mathfrak{r}_r\Lambda_r\hat C_E\hat e_r\Bigl[\mathsf{A}\Bigl(\frac{\mathsf{n}_r}{1-\rho}
+\frac1{(1-\rho)^2}\Bigr)+\frac5{1-\rho}+8\Bigr]\\
&=\hat C_E\hat{\mathfrak{r}}_r\hat e_r=\tfrac14E^R_r .
\end{aligned}
\end{equation*}
The second inequality uses two facts. First, $\rho/(1-\rho)^2\le(1-\rho)^{-2}$. Second, ceiling (b$''$)
gives $4\mathfrak{S}(1+(1-\rho)^{-1})\le\frac1{16}\le8$. The equality is the definition of
$\hat{\mathfrak{r}}_r$ in \eqref{9.10:eq-envelope-ceilings-a}.

Next, the band term. For $r<s<r+\ell_r$, write $m=s-r$, so that $1\le m<\ell_r$. Ceiling (f) gives
$\hat{\mathfrak{r}}_s\le2(1+m)\hat{\mathfrak{r}}_r$. Together with $\hat e_s=\rho^{-m}\hat e_r$ this gives
$E^R_s\le2(1+m)\rho^{-m}E^R_r\le2(1+m)\rho^{-\ell_r}E^R_r$. For $\ell\ge1$ we have
\begin{equation*}
\sum_{1\le m<\ell}2(1+m)=(\ell-1)(\ell+2)\le2\ell^2 .
\end{equation*}
Therefore
\begin{equation*}
2\varkappa'_r\sum_{r<s<r+\ell_r}E^R_s\le4\varkappa'_r\ell_r^2\rho^{-\ell_r}E^R_r\le\tfrac14E^R_r ,
\end{equation*}
where the last inequality is ceiling (g$'$), $4\varkappa'_r\ell_r^2\rho^{-\ell_r}\le\frac14$. Adding the
two bounds gives $\mathcal{L}^R_r[E]\le\frac12E^R_r$.

For the source term, the definition of $\hat{\mathfrak{r}}_r$ gives $\hat{\mathfrak{r}}_r\ge16\mathfrak{r}_r\Lambda_r$.
Hence
\begin{equation*}
\mathfrak{u}^R_r=\mathfrak{r}_r\Lambda_r\bigl(2\hat C_{\sigma}+1\bigr)\hat e_r
\le\frac{\hat{\mathfrak{r}}_r}{16}\cdot2\hat C_E\hat e_r\le\tfrac12E^R_r .
\end{equation*}
The first inequality also uses $2\hat C_{\sigma}+1\le2\hat C_E$.

\emph{The boundary $r=n$.} Here $\hat e_n=1$. At the last depth the boundary step of the
envelope proposition for the run without re-filing is repeated without change; it bounds the
unknowns at depth $n$ through
the constant $C_V$ of that proposition and the logarithmic factor $\mathsf{l}_n$ at depth $n$, and
compares these bounds with the envelope constant:
\begin{equation*}
g^P_n\le C_V\le\hat C_E\hat e_n,\qquad
g^A_n\le2C_V\mathsf{l}_n^{-1}\le4\hat C_E,\qquad
g^R_n\le C_V\mathfrak{r}_n\le4\hat C_E\hat{\mathfrak{r}}_n .
\end{equation*}
The last inequality uses $\hat{\mathfrak{r}}_n\ge16\mathfrak{r}_n$, as in the boundary step repeated
here. Thus $g^{\bullet}_n\le E^{\bullet}_n$
for each $\bullet$.

\emph{Conclusion.} All inputs of Lemma~\ref{9.6:lem-maximum-principle} have now been checked. At each
index of $I_0$ both the linear part and the source term are at most $\frac12E^{\bullet}_r$, and
$g^{\bullet}_n\le E^{\bullet}_n$ holds on $I\setminus I_0$. The lemma therefore yields
$g^{\bullet}_r\le E^{\bullet}_r$ for all
$(r,\bullet)\in I$. This is the first line of \eqref{9.10:eq-geometric-envelopes-a}. With
$g^P_u\le E^P_u$ established for $\underline{r}\le u\le n$, the second line of
\eqref{9.10:eq-geometric-envelopes-a} is
\eqref{9.10:eq-geometric-envelopes-8} together with \eqref{9.10:eq-geometric-envelopes-9}.
\end{proof}

\begin{remark}[Downstream constants without the floor]\label{9.10:rem-downstream-constants}
Consider the re-filed run under the hypotheses of
Proposition~\ref{9.10:prop-geometric-envelopes}, with
$\hat e_r=\rho^{\,n-r}$ as there. Write $e_r$ and $C_E$ for the envelope sequence and the envelope
constant of the envelope proposition for the run without re-filing. Take as antecedents:
\begin{itemize}
\item[(A)] the following four statements attached to the uniqueness theorem: the remark that reads,
of the envelope proposition for the run without re-filing, only its three envelope bounds, its
bounds on the coupling differences, its boundary condition and its hypothesis
$(\mathrm{H}6\sharp)''$; and the proposition and the two lemmas that use that envelope proposition
only through its displayed bounds;
\item[(B)] the re-filed form of the corollary of the uniqueness theorem and the re-filed statement
accompanying it.
\end{itemize}
Then the following hold.
\begin{itemize}
\item[(i)] The remark of (A) holds verbatim with each quantity replaced by its re-filed
counterpart, in particular with the constant $\hat C_E$, whose definition in
\eqref{9.10:eq-source-bounds-b} contains the term $2C^{\star}_X$.
\item[(ii)] The proposition and the two lemmas of (A) hold with $\hat e_r$ in place of $e_r$.
The polynomial floor $\Pi_n$ is replaced by $\Pi_n\equiv 0$, and its sum $T_{\Pi}$ by
$T_{\Pi}=0$.
\item[(iii)] The bookkeeping quantity
\begin{equation}\label{9.10:eq-downstream-constants-1}
\Upsilon_n:=2a_n+\sum_{u<n}\min\Bigl\{3a_u,\;
\frac{2\mathsf{A}\hat C_E\hat e_u}{r_J}\Bigr\},
\end{equation}
where $\mathsf{A}$ is the Lipschitz constant in the bounds on the coupling differences $d_{r+1}$
of \eqref{9.10:eq-geometric-envelopes-a}, has the same order $n^{-3/2}(\log n)^{p_0+1}$ as in the
uniqueness theorem.
\item[(iv)] The constant of clause (a) of the uniqueness theorem, in its re-filed form, is
\begin{equation}\label{9.10:eq-downstream-constants-2}
c':=2c_2+\mathsf{A}\hat C_E(1-\rho)^{-2}.
\end{equation}
\item[(v)] The constants $C_{PM}$ and $C_{\zeta}$ of the statements in (A) are taken with
$\hat C_E$ in place of $C_E$.
\end{itemize}
\end{remark}

\begin{proof}
(i) The remark of (A) reads only the three envelope bounds, the bounds on the coupling differences,
the boundary condition and the hypothesis $(\mathrm{H}6\sharp)''$ of the envelope proposition for
the run without re-filing. For the re-filed run the corresponding inputs are:
\begin{itemize}
\item the three bounds on $g^P_r$, $g^A_r$ and $g^R_r$;
\item the bounds on $d_{r+1}$ and $d^{\sharp}_{r+1}$;
\item the boundary condition of Proposition~\ref{9.10:prop-geometric-envelopes}, namely the
agreement of the two hatted running shifts at the initial depth of that proposition;
\item the hypothesis $(\hat{\mathrm{H}}6\sharp)''$ of
Proposition~\ref{9.10:prop-geometric-envelopes}.
\end{itemize}
Each of these is available for the re-filed run. The bounds are those of
\eqref{9.10:eq-geometric-envelopes-a}, and the boundary condition and the hypothesis
$(\hat{\mathrm{H}}6\sharp)''$ are
hypotheses of Proposition~\ref{9.10:prop-geometric-envelopes}. The remark therefore carries over
word for word, with the constant $\hat C_E$ of \eqref{9.10:eq-source-bounds-b}, whose definition
includes the term $2C^{\star}_X$.

(ii) The proposition and the two lemmas of (A) read the envelope proposition only through its
bounds. Proposition~\ref{9.10:prop-geometric-envelopes} gives these bounds in
\eqref{9.10:eq-geometric-envelopes-a}, with $e_r$ replaced by $\hat e_r$, with $\Pi_n\equiv0$ and
with $T_{\Pi}=0$. By the re-filed corollary and the re-filed statement accompanying it in (B),
which give the vanishing of the polynomial floor, $\Pi_n\equiv0$ and $T_{\Pi}=0$, for the
re-filed run, the three statements therefore hold with these substitutions.

(iii) Inserting $\hat e_u$ into the definition of $\Upsilon_n$ gives
\eqref{9.10:eq-downstream-constants-1}. By the estimate of the bare label $a_n$ under the
re-filing, the
order $n^{-3/2}(\log n)^{p_0+1}$ of $\Upsilon_n$ is the effect of one additional ultraviolet step
on the bare label $a_n$. It is not an effect of the remnant. This order is therefore the same as
in the uniqueness theorem.

(iv) With $T_{\Pi}=0$, the constant of clause (a) of the uniqueness theorem reads
\eqref{9.10:eq-downstream-constants-2}.

(v) This is the substitution of (ii) applied to the constants $C_{PM}$ and $C_{\zeta}$: they are
taken with $\hat C_E$.
\end{proof}

\begin{remark}[Repeated large-field operations paid by geometric forgetting]
\label{9.10:rem-geometric-forgetting}
We consider the repeated applications of Ba{\l}aban's large-field operation $\mathbf{R}$ to a
large-field region across successive steps, together with the chains of large-field regions
they re-create. Assume the following two statements.
\begin{itemize}
\item[(A)] (Geometric forgetting of old regions.) Every live large-field
component of age $a$ carries a geometric factor $q^{a+1}$, where $q<1$ is the forgetting rate.
The rate $q$ comes from the
entropy constant of Ba{\l}aban's large-field expansion under the one-sided condition
$p_0\ge 6r+1$ on the degree exponent $p_0$, $r$ being the radius exponent of
Notation~\ref{9.1:not-radius-exponent-floors} (not a depth).
\item[(B)] (The second-class statement.) The statement that the proof of the lemma on
geometric forgetting of old regions invokes for a live component with an outcome of the
second class.
\end{itemize}
Then the following hold.
\begin{itemize}
\item[(a)] The repeated applications of $\mathbf{R}$ and the chains of re-created regions are
paid by the factor $q^{a+1}$ per live component of
age $a$. Statement (A) holds for the hatted lists $\hat{\mathfrak{s}},
\tilde{\hat{\mathfrak{s}}}\in\hat{\mathfrak{K}}_k$ of Definition~\ref{9.10:def-pair-loads},
with the constant $\hat{\mathsf{c}}$ in place of its un-hatted counterpart.
\item[(b)] No unfiled share, in the sense of Definition~\ref{9.10:def-pair-loads}, enters the
estimates at a level age that is unbounded inside a long-lived component.
\end{itemize}
\end{remark}

\begin{proof}
(a) By (A), with (B) used for the second-class outcome, each live component of age $a$
contributes the factor $q^{a+1}$. This factor is the quantity that pays for the repeated
applications of $\mathbf{R}$ to that component and for the chains of regions it re-creates.
Clauses (d) and (e) of the one-lattice Lipschitz step and the rows $(R')$ and $(\mathbf{T})$ of
its kernel form, both taken for un-hatted lists (their re-filed forms are
Theorems~\ref{9.10:thm-refiled-lipschitz} and~\ref{9.10:thm-lipschitz-kernel}; each is
stated; proof incomplete), are already geometric in zone age (the age of the
large-field component carrying the zone, that is, the number of steps since the step that
created it) and in operation age.

Since the proof of the lemma on geometric forgetting of old regions estimates no term of the
effective action individually, the lemma holds for the hatted lists of
Definition~\ref{9.10:def-pair-loads}, with the constant $\hat{\mathsf{c}}$ in place of its
un-hatted counterpart.

An alternative route pays for these repetitions by the margin in the strict inequality
$\nu<p_0$ between the parameter $\nu$ of the large-field estimates of \cite{B-CMP122a} and
\cite{B-CMP122b} and the degree exponent $p_0$. That route is not used in this chapter.

(b) The only question about the repeated applications of $\mathbf{R}$ and the re-created
regions that involves remnants is whether an unfiled share
can enter the estimates at an unbounded level age inside a long-lived component. By
Lemma~\ref{9.10:lem-unfiled-range}, the level range of the unfiled shares in the zone at
depth $s$ is at most $\ell_s$, as displayed in \eqref{9.10:eq-unfiled-range-a}, $\ell_s$ being
the remnant depth range fixed there. Within a zone the level age of
an unfiled share is therefore bounded. The age of the zone is paid, by (a), by the factor
$q^{a+1}$ of (A). Hence no unfiled share enters the estimates at an unbounded level age.
\end{proof}

\begin{definition}[Stored arrays and loads at the last lattice]\label{9.10:def-last-lattice-loads}
The following statements on the two runs $\mathsf{A}$ and $\mathsf{B}$ of
Definition~\ref{9.6:def-comparison-chain} are restated for the re-filed objects, that is, for the
re-filed lists $\hat{\mathfrak{s}}_j$ of class $\hat{\mathfrak{K}}_k$ and under the substitution
rule \eqref{9.10:eq-substitution-rule-a} of Definition~\ref{9.10:def-substitution-rule}, with
every un-hatted quantity replaced by its hatted counterpart: the last-lattice comparison theorem;
its corollary at the comparison rate $\hat\varepsilon_n$; the two statements that follow from
that corollary, of which the second is the last-lattice localisation lemma of this chapter; and
the three items of the lemma of the last-lattice comparison that is deduced from these statements.
Throughout, $n$ is the length of the run, $k:=n-\mathsf{w}$ is the stop level, where
$\mathsf{w}$ is the last-step offset, and $j$ denotes a zone level, of depth $s=n-j$. The
hatted objects are the following.

\begin{enumerate}
\item[(a)] \emph{Centres, energies, histories, labels and floors.} The centres are
$\hat\zeta^n_*+u$ and $\hat\zeta^{n+1}_*+\hat{\mathfrak{n}}_n(u)$, where $\hat\zeta^n_*$ is the
centre of the hatted uniqueness theorem, $\hat{\mathfrak{n}}_n$ its bookkeeping map, and $u$
the variable so denoted in the un-hatted last-lattice comparison theorem. The vacuum energy
$\mathcal{E}^{\mathsf{c}}$ is subtracted at the
run's own centre (the superscript $\mathsf{c}$ marks, here and below, the objects of the
last-lattice comparison theorem); the history integrals are $\tau^{\mathsf{c}}_h$; the labels
are common to
the two runs and aligned on the prescribed scaffold $\Theta_X$ at $X=g^{-2}$. The floors are
$n_1$ and $n_0(N_w)$, where $N_w$ is the quantity so denoted in the un-hatted floors, together
with $n_2(g)$; the floor $n_2(g)$ guarantees that
$\mathrm{filed}(\cdot)$, $\mathfrak{B}(\cdot)$ and $\mathrm{band}(\cdot)$ of
\eqref{9.10:eq-filing-level-a} are the same index sets in the two runs $\mathsf{A}$ and
$\mathsf{B}$ beneath depth $n/2$, namely through the condition below, in which $\ell_r$ is the
remnant depth-range letter of \eqref{9.10:eq-filing-level-a}, $\theta_r$ is the scaffold value,
and $l$ and $\mathsf{N}(\cdot)$ are the quantities so denoted in the floor condition of the
un-hatted statement:
\begin{equation*}
l/2>\max\{\ell_r,\mathsf{N}(\theta_r)\}.
\end{equation*}

\item[(b)] \emph{The stored array and the uses of a filed remnant.} The array
$x^{\mathsf{c}}$ of stored numbers is the array of Definition~\ref{9.10:def-substitution-rule};
the classes $z=(j,k')$ and the marks $W^{(z)}$ are as there. Let $R^{(i)}$ be a filed remnant
of level $i$, with filing level $i''$. Its uses as an $E$-slot are uses of the stored slot
$\hat E^{(i'')}$, in the classes $(i'',k')$ onward; its uses as a remnant are uses in the
classes $(n',k')$ with zone $n'\in[i,i'']$.

\item[(c)] \emph{The loads.} Under the substitution rule \eqref{9.10:eq-substitution-rule-a}
the loads at the zone level $j$ and their weighted sum up to the level $k$ are
\begin{equation}\label{9.10:eq-last-lattice-loads-a}
\begin{aligned}
{}[\![\hat\delta]\!]^u_j
&:=\sum_{i\le j}\varsigma_{1+j-i}\,\mathsf{S}^+_i[\hat E]
+\sum_{j-\mathsf{a}_i\le i\le j}w_{j,i}\,\mathsf{S}^+_i[R]
+\bigl[\mathsf{y}_j+\mathsf{y}'_j\bigr]\\
&\qquad+\mathfrak{g}_j\Bigl[\hat{\mathsf{z}}^{\sharp}_j
+\sum_{i\le j}L^{-(j-i)/2}\,\hat{\mathsf{b}}_i\Bigr],\\
{}[\![\hat\delta]\!]^c_j
&:=\sum_{i\le j}\varsigma_{1+j-i}\,E_0\,\hat\vartheta^{\natural}_i
+\sum_{j-\mathsf{a}_i\le i\le j}w_{j,i}\,g_i^{\kappa_0}\,\hat\vartheta^{\natural}_i
+\sum_{i\le j}w^{\mathbf{B}}_{j,i}\,B_0\,\hat\vartheta^{\natural}_i\\
&\qquad+E_0\min\bigl\{1,\,C^{\natural}_{\mathrm{cl}}(\log\hat x_j)^{c_2}\vartheta_2^{j/4}\bigr\},\\
{}[\![\hat\delta]\!]^{\natural+}_j
&:=[\![\hat\delta]\!]^u_j+[\![\hat\delta]\!]^c_j,\\
\Sigma^{q_a}_k[\hat\delta]
&:=\sum_{j\le k}(k-j+1)\,q_a^{(k-j-1)_+}\,[\![\hat\delta]\!]^{\natural+}_j .
\end{aligned}
\end{equation}
Here $\varsigma_m$ is the geometric kernel, $w_{j,i}$ and $w^{\mathbf{B}}_{j,i}$ are the zone
weights, $\mathsf{S}^+_i[\cdot]$ is the size of the two-run difference of the stored slot read
at level $i$, $\mathfrak{g}_j$ is the source-channel weight, $\hat{\mathsf{b}}_i$ is the norm
of the difference of the coupling increments fixed in
Definition~\ref{9.10:def-substitution-rule}, $\hat{\mathsf{z}}^{\sharp}_j$ is the running-shift
difference entry so denoted in the un-hatted loads of the last-lattice comparison, read for the
hatted running shifts $\hat\zeta_j$ of Definition~\ref{9.10:def-substitution-rule},
$\mathsf{y}_j$, $\mathsf{y}'_j$ are the entries so denoted in the un-hatted loads of the
last-lattice comparison, $\hat\vartheta^{\natural}_i$ is the index-stationarity error at level
$i$, and $\hat x_j$ is the hatted inverse coupling. In the first sum of
$[\![\hat\delta]\!]^u_j$ the size of the stored $E$-slot satisfies
\begin{equation}\label{9.10:eq-last-lattice-loads-1}
\mathsf{S}^+_i[\hat E]\le \mathsf{S}^+_i[E]+3C_{RF}\,\mathsf{S}^+_{i_*}[R],
\end{equation}
where $i_*$ is the level filed into $\hat E^{(i)}$; the second summand is the filing term. The
load $[\![\hat\delta]\!]^c_j$ collects, inside each mark, the contributions of the classical
part: the index-stationarity errors $\hat\vartheta^{\natural}_i$ of its three sums, and its last
term, whose constant $C^{\natural}_{\mathrm{cl}}$ is the one so denoted in the un-hatted loads.
For a use
of a filed remnant of level $i$ as an $E$-slot, the index-stationarity error enters the
$\varsigma$-weighted entry of $\hat E^{(i'')}$, at size
$C_{RF}\,w_{i''-1,i}\,g_i^{\kappa_0}\le E_0$; it is therefore contained in the first sum of
$[\![\hat\delta]\!]^c_j$.

\item[(d)] \emph{The a priori bound and the dial domain.} Let $\hat E_1$ denote $E_1$ with
$C_L$ replaced by $\hat C_L$, and $\hat E^z_1$ denote $E^z_1$ with the same replacement. Under
the hypothesis of the a priori lemma for the loads of the last-lattice comparison, read for the
re-filed runs $\mathsf{A}$ and $\mathsf{B}$,
\begin{equation}\label{9.10:eq-last-lattice-loads-2}
[\![\hat\delta]\!]^{\natural+}_j\le \hat E^z_1 \qquad\text{for every zone level } j.
\end{equation}
The dial domain $\hat{\mathbb{D}}^{\star}$ is $\mathbb{D}^{\star}$ taken at the loads
$[\![\hat\delta]\!]^{\natural+}$ and the bound $\hat E^z_1$; the constants
$\mathsf{c}^{\star}$ are replaced by $\hat{\mathsf{c}}^{\star}$, taken at $\hat C_L'$.

\item[(e)] \emph{Rates and envelope.} The rates $q$, $q_a$, $\vartheta_2$, $\vartheta_{dl}$,
$\vartheta_{\chi}$ are those of the un-hatted last-lattice comparison, unchanged, and so is the
window of the contraction letter,
\begin{equation*}
\rho\in\bigl[\max\{q^{1/2},L^{-\vartheta_{\sharp}},\vartheta_{dl},\vartheta_{\chi}\},\,1\bigr).
\end{equation*}
The envelope letter is $\hat e_s:=\rho^{n-s}$, as in \eqref{9.10:eq-envelope-ceilings-a}, and
the constant $C_E$ is read as $\hat C_E$ of \eqref{9.10:eq-source-bounds-b} in
Proposition~\ref{9.10:prop-source-bounds} (stated there with proof incomplete at the step
marked $(\ast)$).
\end{enumerate}
\end{definition}

\begin{proof}
Three assertions are contained in the definition: the description of the uses of a filed
remnant in item~(b), the bound \eqref{9.10:eq-last-lattice-loads-1}, and the a priori bound
\eqref{9.10:eq-last-lattice-loads-2}.

\emph{Item~(b).} By Definition~\ref{9.10:def-substitution-rule}, a filed remnant $R^{(i)}$ is a
stored slot which is read as a remnant at the levels $k'\in[i,i''-1]$, through
$\mathfrak{B}(\cdot)$ and $\mathrm{band}(\cdot)$ of \eqref{9.10:eq-filing-level-a}, and, as
unfiled shares, in the zones $n'\in[i,i'']$ at any later level; and which is read as an
$E$-slot from the level $i''$ on, through the stored slot $\hat E^{(i'')}$. By
Lemma~\ref{9.10:lem-unfiled-range}, in the form \eqref{9.10:eq-unfiled-range-a}, the unfiled
shares of $R^{(i)}$ occupy only zones in this range. Hence, by the convention of this chapter on
the classes $z=(j,k')$, the uses of $R^{(i)}$ as an $E$-slot are those of the classes
$(i'',k')$ onward, and its uses as a remnant are those of the classes $(n',k')$ with
$n'\in[i,i'']$.

\emph{The bound \eqref{9.10:eq-last-lattice-loads-1}.} Let $i_*$ be the level filed into
$\hat E^{(i)}$, so that $i$ is the filing level $i_*''$ of $i_*$. By
Definition~\ref{9.10:def-substitution-rule}, the two-run difference of $\hat E^{(i)}$ contains
the re-filed family $\mathsf{RF}[\delta R^{(i_*)}]$, which is linear in $\delta R^{(i_*)}$, and,
in the notation of \eqref{9.10:eq-last-lattice-loads-a},
\begin{equation*}
\mathsf{S}^+_i[\hat E]\le \mathsf{S}^+_i[E]+C_{RF}\,w_{i-1,i_*}\,\mathsf{S}^+_{i_*}[R].
\end{equation*}
The zone weight of the filing term satisfies $w_{i-1,i_*}\le 3$ on the prescribed scaffold
$\Theta_X$, and \eqref{9.10:eq-last-lattice-loads-1} follows.

\emph{The a priori bound \eqref{9.10:eq-last-lattice-loads-2}.} Each term of
\eqref{9.10:eq-last-lattice-loads-a} is the corresponding term of the un-hatted loads of the
last-lattice comparison, read for the re-filed runs, that is, with the hatted slot sizes in place
of the un-hatted ones; and it is bounded by that term with the constant $C_L$, wherever $C_L$
enters the coefficients of the un-hatted loads, replaced by $\hat C_L$. In
$[\![\hat\delta]\!]^c_j$ the index error of a use of a filed remnant as an $E$-slot lies inside
the first sum, as stated in item~(c). The a priori
lemma for the un-hatted loads is proved by an induction along the half-steps which uses only
the structure of the loads, under the hypothesis of that lemma, which is assumed in item~(d) for
the re-filed runs; it bounds the
un-hatted loads by $E^z_1$. Since the hatted loads have this structure with $C_L$ replaced by
$\hat C_L$, the same induction applies to them and bounds them by $E^z_1$ at $\hat C_L$, which
is $\hat E^z_1$. This is \eqref{9.10:eq-last-lattice-loads-2}.
\end{proof}

\begin{lemma}[Finite-range envelope of last-lattice loads. Stated; proof incomplete at the step marked $(\ast)$]
\label{9.10:lem-load-envelope}
Assume the following.
\begin{enumerate}
\item[(i)] The conclusions of Proposition~\ref{9.10:prop-geometric-envelopes}, with the value of
$\hat C_E$ fixed in Remark~\ref{9.10:rem-downstream-constants} through its constant $C^{\star}_X$.
\item[(ii)] Conclusion (a) of Theorem~\ref{9.10:thm-refiled-lipschitz},
display~\eqref{9.10:eq-refiled-lipschitz-a} (the proof of that theorem is incomplete at one step),
in the following form: the
$\mathsf{T}$-outputs with domain in the region $\Lambda_{k+1}$ differ by at most
\begin{equation*}
\hat\varkappa_k[\![\hat\delta]\!]^{\natural T}_k
+C_{RF}\sum_{i\in\mathfrak{B}(k)}w_{k,i}\,\delta^R_i+\sigma^T_k ,
\end{equation*}
where $[\![\hat\delta]\!]^{\natural T}_k$ is the weighted zone sum of the hatted loads that enters
the $\mathsf{T}$-rows, $\delta^R_i$ is the size of the difference of the remnant slots $R$ of the two
lists at the filed level $i$, and $\sigma^T_k$ is the term of the $\mathsf{T}$-rows formed, as
$\sigma_k$ (the sum over $i\in\mathfrak{B}(k)$ of $w_{k,i}x_i^{-\kappa_0/2}$), from the levels filed
at step $k$.
\item[(iii)] The hypothesis on the weights $w^{\mathbf{B}}_{j,i}$ that enter the load
$[\![\hat\delta]\!]^c_j$ of Definition~\ref{9.10:def-last-lattice-loads}, together
with the volume and filament bound in its class-by-class form, the $\mathbf{B}$-load proposition,
the $\mathbf{B}$-elimination proposition and the a priori load lemma, each applied with the loads
$[\![\hat\delta]\!]$ and the envelope $\hat e$ in place of the loads and envelope of the run without
re-filing: their statements involve only the loads and one stored family, class by class, according
to the substitution rule~\eqref{9.10:eq-substitution-rule-a}.
\item[(iv)] $g\le\gamma_U$, and the ceilings~\eqref{9.10:eq-envelope-ceilings-b} of
Definition~\ref{9.10:def-envelope-ceilings}.
\end{enumerate}
The notation is that of Definition~\ref{9.10:def-last-lattice-loads}, in particular the loads and
stored arrays displayed in~\eqref{9.10:eq-last-lattice-loads-a}.
Then for every zone $j\le k$ aligned in the sense of Definition~\ref{9.10:def-last-lattice-loads},
with $s:=n-j$,
\begin{equation}\label{9.10:eq-load-envelope-a}
[\![\hat\delta]\!]^{\natural+}_j\le\hat C_{\mathrm{env}}\,\mathsf{l}_s^2\,\hat e_s,
\qquad
\Sigma^{q_a}_k[\hat\delta]\le C^{TC}_{\Sigma}\,\hat C_{\mathrm{env}}\,\mathsf{l}_{\mathsf{w}}^2\,
\hat e_{\mathsf{w}} ,
\end{equation}
where $\hat C_{\mathrm{env}}$ is the constant produced by part (iv) of the $\mathbf{B}$-elimination
proposition from the constants $C^{(E)}$, $C^{(R)}$, $C^{(\zeta)}$, $C^{(c)}$ of the
proof below, which contain $\hat C_E$, and
\begin{equation*}
C^{TC}_{\Sigma}=\sum_m(m+1)\hat{\mathsf{l}}_m^2q_a^{(m-1)_+}\rho^{-m}
\end{equation*}
(step (4) of that proof); these constants contain no
contribution of a polynomial floor term ($\Pi_n\equiv0$, as in
Proposition~\ref{9.10:prop-geometric-envelopes}).
\end{lemma}

\begin{proof}
The proof has the four steps of the proof of the last-lattice envelope lemma for the loads of the
run without re-filing, taken with the envelope $\hat e$. Since $\hat e_r=\rho^{n-r}$, the
relation between envelopes at two depths holds with equality:
\begin{equation}\label{9.10:eq-load-envelope-1}
\hat e_{s'}=\rho^{-(s'-s)}\,\hat e_s\qquad(s'\ge s).
\end{equation}
We first extract three consequences of~\eqref{9.10:eq-envelope-ceilings-b}. The coefficients
$\varkappa,\varkappa_A$ are non-negative, the constant $C_F$ of the inequality (b$''$) satisfies
$C_F\ge 3C_{RF}$, and $1+(1-\rho)^{-1}\ge1$, so the
inequality (b$''$) gives $8\mathfrak{S}\le1/16$ and $24C_{RF}\mathfrak{S}\le 8C_F\mathfrak{S}\le1/16$,
where $\mathfrak{S}=\sup_{r\ge0}(\ell_r+2)^2\hat{\mathfrak{r}}_r\rho^{-\ell_r}$. Since
$(1+\ell_r)\le(\ell_r+2)^2$, $2(\ell_r+1)^2\le2(\ell_r+2)^2$ and
$\mathfrak{r}_r\le\hat{\mathfrak{r}}_r/16$ for the quantities
of~\eqref{9.10:eq-envelope-ceilings-a},
it follows that for every $r$
\begin{equation}\label{9.10:eq-load-envelope-2}
\begin{gathered}
24C_{RF}(1+\ell_r)\hat{\mathfrak{r}}_r\rho^{-\ell_r}\le1,\qquad
2(\ell_r+1)^2\hat{\mathfrak{r}}_r\rho^{-\ell_r}\le1,\\
16(\ell_r+1)\mathfrak{r}_r\rho^{-\ell_r}\le1 .
\end{gathered}
\end{equation}

\emph{Step (1): the unknowns at depth $s'\ge s$.} Let $i$ be a level of depth $s':=n-i\ge s$. By
hypothesis (i), that is by the bounds~\eqref{9.10:eq-geometric-envelopes-a}, the stored-slot
difference sizes of~\eqref{9.10:eq-two-slot-bounds-a} satisfy
\begin{equation*}
\begin{gathered}
\mathsf{S}^+_i[P^{\flat}],\ \mathsf{S}^+_i[\hat D]\le4\hat C_EE_0\mathsf{l}_{s'}\hat e_{s'},\qquad
\mathsf{S}^+_i[\hat Q]\le\frac{\hat C_E\hat e_{s'}}{C_t\theta_{s'}},\qquad
\mathsf{S}^+_i[R]\le\frac{4\hat C_E\hat{\mathfrak{r}}_{s'}\hat e_{s'}}{C_t\theta_{s'}} .
\end{gathered}
\end{equation*}
For the filed slot we use the bound
$\mathsf{S}^+_i[\hat E]\le\mathsf{S}^+_i[E]+3C_{RF}\mathsf{S}^+_{i_*}[R]$
of~\eqref{9.10:eq-last-lattice-loads-a}, where $i_*$ is the level filed into $\hat E^{(i)}$, whose
depth is $s'+\mathsf{a}$ with $\mathsf{a}:=\mathsf{a}_{i_*}$ (see~\eqref{9.10:eq-filing-level-a}).
This gives
\begin{align*}
\mathsf{S}^+_i[\hat E]
&\le\hat C_E\bigl(8E_0\mathsf{l}_{s'}+(C_t\theta_{s'})^{-1}\bigr)\hat e_{s'}
+3C_{RF}\,\frac{4\hat C_E\hat{\mathfrak{r}}_{s'+\mathsf{a}}\hat e_{s'+\mathsf{a}}}
{C_t\theta_{s'+\mathsf{a}}}\\
&\le\hat C_E\bigl(8E_0\mathsf{l}_{s'}+2(C_t\theta_{s'})^{-1}\bigr)\hat e_{s'} ,
\end{align*}
because $\hat{\mathfrak{r}}_{s'+\mathsf{a}}\le2(1+\mathsf{a})\hat{\mathfrak{r}}_{s'}$ by item (f)
of~\eqref{9.10:eq-envelope-ceilings-b},
$\hat e_{s'+\mathsf{a}}=\rho^{-\mathsf{a}}\hat e_{s'}$ by~\eqref{9.10:eq-load-envelope-1},
$\theta_{s'+\mathsf{a}}\ge\theta_{s'}$ and $\mathsf{a}\le\ell_{s'}$, and the second term is
therefore at most
\begin{equation*}
24C_{RF}(1+\ell_{s'})\hat{\mathfrak{r}}_{s'}\rho^{-\ell_{s'}}\cdot
\frac{\hat C_E\hat e_{s'}}{C_t\theta_{s'}}
\le\frac{\hat C_E\hat e_{s'}}{C_t\theta_{s'}}
\end{equation*}
by the first inequality of~\eqref{9.10:eq-load-envelope-2}. For the couplings,
\begin{equation*}
\hat{\mathsf{b}}_i\le\mathsf{A}\hat C_E\hat e_{s'},\qquad
d_{s'}\le\mathsf{A}\hat C_E\mathsf{n}_{s'}\hat e_{s'},\qquad
\hat{\mathsf{z}}_{s'}\le C_{\zeta}(1+s')\hat e_{s'},
\end{equation*}
the last by the form of the bound on the running shift recorded in
Remark~\ref{9.10:rem-downstream-constants}, $C_{\zeta}$ being the constant of that bound.

\emph{Step (2)(E).} With $m:=j-i$, so that the depth of level $i$ is $s+m$, with
$\varsigma_{1+m}\le C_{\varsigma}(m+1)^3q^m$ (where $C_{\varsigma}$ is the constant of this bound on
the geometric kernel), $\mathsf{l}_{s+m}\le\hat{\mathsf{l}}_m\mathsf{l}_s$ and
$\hat e_{s+m}=\rho^{-m}\hat e_s$, step (1) gives
\begin{align*}
\sum_{i\le j}\varsigma_{1+j-i}\mathsf{S}^+_i[\hat E]
&\le\sum_m C_{\varsigma}(m+1)^3q^m\,
\hat C_E\bigl(8E_0\hat{\mathsf{l}}_m\mathsf{l}_s+2\bigr)\rho^{-m}\hat e_s
\le C^{(E)}\mathsf{l}_s\hat e_s,\\
C^{(E)}&:=2\hat C_EC_{\varsigma}(8E_0+1)\sum_m(m+1)^3\hat{\mathsf{l}}_m(q/\rho)^m<\infty ,
\end{align*}
the ratio $q/\rho$ being at most $\rho<1$ on the window of $\rho$ of
Definition~\ref{9.10:def-last-lattice-loads}.

\emph{Step (2)(R).} For $j-\mathsf{a}_i\le i\le j$ the depth $s'=n-i$ ranges over
$s\le s'\le s+\ell_s$, at most $\ell_s+1$ values, and the zone weights obey
$w_{j,i}\le C^{\theta}_w\theta_{s'}/\theta_s$, $C^{\theta}_w$ being the constant of this bound on the
zone weights. With item (f)
of~\eqref{9.10:eq-envelope-ceilings-b}, $\hat{\mathfrak{r}}_{s'}\le2(1+\ell_s)\hat{\mathfrak{r}}_s$,
and with $\hat e_{s'}\le\rho^{-\ell_s}\hat e_s$,
\begin{align*}
\sum_{j-\mathsf{a}_i\le i\le j}w_{j,i}\mathsf{S}^+_i[R]
&\le\sum_{s\le s'\le s+\ell_s}C^{\theta}_w\frac{\theta_{s'}}{\theta_s}\cdot
\frac{4\hat C_E\hat{\mathfrak{r}}_{s'}\hat e_{s'}}{C_t\theta_{s'}}\\
&\le\frac{4C^{\theta}_w\hat C_E}{C_t\theta_s}\,(\ell_s+1)\cdot2(1+\ell_s)\hat{\mathfrak{r}}_s\,
\rho^{-\ell_s}\hat e_s
\le\frac{4C^{\theta}_w\hat C_E}{C_t\theta_s}\,\hat e_s=:C^{(R)}\hat e_s ,
\end{align*}
by the second inequality of~\eqref{9.10:eq-load-envelope-2}. No polynomial floor term enters: beyond
the range $\ell_s$ the remnant entry is absent, and within it the factor $\rho^{-\ell_s}$ is absorbed
by~\eqref{9.10:eq-load-envelope-2}.

\emph{Step (2)($\mathbf{B}$).} $(\ast)$ The $\mathbf{B}$-loads are eliminated by parts (ii)--(iv) of
the $\mathbf{B}$-elimination proposition in the envelope norm
$\mathsf{l}_s^2\hat e_s$, uniformly in the number of steps. The proof of that proposition
uses only the structure of the loads and the envelope of the five unknowns of step (1); here it is
applied with the loads $[\![\hat\delta]\!]$ and the envelope $\hat e$ in place of the loads and
envelope of the run without re-filing, as in hypothesis (iii). The volume and filament
bound, the $\mathbf{B}$-load proposition, the $\mathbf{B}$-elimination proposition and the a priori
load lemma are stated for the loads of the run without re-filing and its envelope; their validity
with $[\![\hat\delta]\!]$ and $\hat e$ in place of those, on which this step rests, is not proved here.

\emph{Step (2)($\zeta$).} The source-channel term
\begin{equation*}
\mathfrak{g}_j\Bigl[\hat{\mathsf{z}}^{\sharp}_j+\sum_{i\le j}L^{-(j-i)/2}\hat{\mathsf{b}}_i\Bigr]
\end{equation*}
is bounded, with the coupling bounds of step (1), as in step (2)($\zeta$) of the proof of the
last-lattice envelope lemma for the run without re-filing, with $\hat e$ in place of the envelope
there: the geometric sums that occur have ratios $L^{-1}/\rho$ and $L^{-1/2}/\rho$, both at most
$\rho$, against which $\hat e_{s'}=\rho^{-(s'-s)}\hat e_s$ is summable. The result is a bound
$C^{(\zeta)}\mathsf{l}_s^2\hat e_s$.

\emph{Step (3): the load $[\![\hat\delta]\!]^c_j$.} Its terms carrying $E_0$ (the first) and $B_0$
are bounded as in step (3) of the proof of the last-lattice envelope lemma for the run
without re-filing, which applies without change once the index-stationarity error is bounded by
\begin{equation*}
\hat\vartheta^{\natural}_i\le C_{\vartheta}\vartheta_{dl}^{\,n-s'}\le C_{\vartheta}\rho^{n-s'}
=C_{\vartheta}\hat e_{s'} ,
\end{equation*}
with $C_{\vartheta}$ the constant of this bound on the index-stationarity error, the middle
inequality because $\vartheta_{dl}\le\rho$ on the window of $\rho$. For the term
carrying $g_i^{\kappa_0}$, the range is again $s\le s'\le s+\ell_s$ and
\begin{align*}
\sum_{j-\mathsf{a}_i\le i\le j}w_{j,i}g_i^{\kappa_0}\hat\vartheta^{\natural}_i
&\le C^{\theta}_wC_{\vartheta}2^{\kappa_0/2}\theta_s^{-1}
\sum_{s\le s'\le s+\ell_s}\theta_{s'}^{1-\kappa_0/2}\rho^{n-s'}\\
&\le C^{\theta}_wC_{\vartheta}2^{\kappa_0/2}(C_t\theta_s)^{-1}(\ell_s+1)\mathfrak{r}_s
\rho^{-\ell_s}\cdot\rho^{n-s}
\le C^{\mathrm{b}}(C_t\theta_s)^{-1}\hat e_s ,
\end{align*}
with $C^{\mathrm{b}}:=C^{\theta}_wC_{\vartheta}2^{\kappa_0/2}$, by the third inequality
of~\eqref{9.10:eq-load-envelope-2}. This is the one entry without index smallness at old levels:
beyond the range $\ell_s$ it is absent, and within it the decay of $\mathfrak{r}$ itself absorbs the
factor $\rho^{-\ell_s}$; no polynomial floor term enters. The last term of
$[\![\hat\delta]\!]^c_j$ is bounded as in the last line of step (3) of the same proof, without
change. Together these give $[\![\hat\delta]\!]^c_j\le C^{(c)}\mathsf{l}_s\hat e_s$.

\emph{Adding.} Write $[\![\hat\delta]\!]^{(5)}_j$ for the part of $[\![\hat\delta]\!]^u_j$ other than
its $\mathbf{B}$-entries, that is, the part bounded in steps (2)(E), (2)(R) and (2)($\zeta$). Steps
(2) and (3) give
\begin{equation*}
[\![\hat\delta]\!]^{(5)}_j+[\![\hat\delta]\!]^c_j
\le\bigl(C^{(E)}+C^{(R)}+C^{(\zeta)}\bigr)\mathsf{l}_s^2\hat e_s+C^{(c)}\mathsf{l}_s\hat e_s .
\end{equation*}
Part (iv) of the $\mathbf{B}$-elimination proposition (equivalently, part (a) of its corollary),
applied with $\hat e$ as in step (2)($\mathbf{B}$), then gives
$[\![\hat\delta]\!]^{\natural+}_j\le\hat C_{\mathrm{env}}\mathsf{l}_s^2\hat e_s$, where
$\hat C_{\mathrm{env}}$ is the constant that part produces from the constants $C^{(E)}$, $C^{(R)}$,
$C^{(\zeta)}$, $C^{(c)}$ above. This is the first inequality of~\eqref{9.10:eq-load-envelope-a}.

\emph{Step (4): the weighted sum.} By~\eqref{9.10:eq-last-lattice-loads-a},
\begin{equation*}
\Sigma^{q_a}_k[\hat\delta]=\sum_{j\le k}(k-j+1)q_a^{(k-j-1)_+}[\![\hat\delta]\!]^{\natural+}_j .
\end{equation*}
With $m:=k-j$ and $\mathsf{w}=n-k$, so that $n-j=\mathsf{w}+m$, the first inequality
of~\eqref{9.10:eq-load-envelope-a}, $\mathsf{l}_{\mathsf{w}+m}\le\hat{\mathsf{l}}_m\mathsf{l}_{\mathsf{w}}$
and $\hat e_{\mathsf{w}+m}=\rho^{-m}\hat e_{\mathsf{w}}$ give
\begin{align*}
\Sigma^{q_a}_k[\hat\delta]
&\le\hat C_{\mathrm{env}}\sum_m(m+1)q_a^{(m-1)_+}\mathsf{l}_{\mathsf{w}+m}^2\hat e_{\mathsf{w}+m}
\le\hat C_{\mathrm{env}}\mathsf{l}_{\mathsf{w}}^2\hat e_{\mathsf{w}}
\sum_m(m+1)\hat{\mathsf{l}}_m^2q_a^{(m-1)_+}\rho^{-m}\\
&=C^{TC}_{\Sigma}\hat C_{\mathrm{env}}\mathsf{l}_{\mathsf{w}}^2\hat e_{\mathsf{w}},
\qquad
C^{TC}_{\Sigma}=\sum_m(m+1)\hat{\mathsf{l}}_m^2q_a^{(m-1)_+}\rho^{-m},
\end{align*}
where $q_a\le q^2\le\rho^4$. This is the second inequality of~\eqref{9.10:eq-load-envelope-a}.

Hypothesis (i) is available along the whole chain $0\le s\le n$; no comparison of vacuum energies at
depths above $n/2$ enters; and the factor $\mathsf{l}_{\mathsf{w}}$ depends on $g$.
\end{proof}

\begin{lemma}[Objective, test and remainder estimates]\label{9.10:lem-objective-test-rest}
We use the hatted arrays, classes $z=(j,k')$ and loads of
Definition~\ref{9.10:def-last-lattice-loads}, in particular the weighted zone sum
$\Sigma^{q_a}_k[\hat\delta]$ of \eqref{9.10:eq-last-lattice-loads-a}, and the two runs
$\mathsf{A}$, $\mathsf{B}$ of Definition~\ref{9.6:def-comparison-chain}; the run length $n$, the
level $k$ and the last-step offset $\mathsf{w}$ are those of
Definitions~\ref{9.10:def-last-lattice-loads} and~\ref{9.10:def-pair-loads}. Assume the following.
\begin{itemize}
\item[(H$_1$)] The objective lemma, the test lemma and the remainder lemma of the uniqueness
theorem hold for the un-hatted lists, with the proofs given for them under the uniqueness
theorem.
\item[(H$_2$)] The cluster-expansion proposition with one dial per class holds in its hatted
form, on the hatted array.
\item[(H$_3$)] Geometric forgetting of old regions holds in its hatted form on the dial domain
$\hat{\mathbb{D}}^{\star}$.
\item[(H$_4$)] The torus bound holds in its hatted form; more precisely, each input of its proof
by moduli of the upper member of \cite[(2.50)]{B-CMP119} holds for the hatted quantities. These
inputs are: the sourced forms, with the complex source $w$, of \cite[(2.43)--(2.49)]{B-CMP119}
and of \cite[(1.68)--(1.69)]{B-CMP122a} established in the proof of the correlation theorem; the
printed right member $\mathfrak{M}_k(X)$ of the modulus bound; the raw Wilson strength (the size
of the Wilson factor of the activities before any expansion); and \cite[(1.97)]{B-CMP122a}.
\item[(H$_5$)] Item~(iii) of the pointed-cluster statements (those on which the filing of uses
in Definition~\ref{9.10:def-pair-loads} rests) holds, with $n_c$ its exponent and $C_{nc}$ its
constant; the geometric-rate arithmetic for its factor $q_a^{-n_c}$ supplies the inequality
$4\,q_a^{\,k-n_c}\le C_{nc}\,\rho^{\,n-\mathsf{w}}$.
\end{itemize}
Then the following hold.
\begin{itemize}
\item[(i)] \emph{Objective.} The objective lemma holds for the hatted lists, with its objective
read on the hatted arrays $x$,
\begin{equation}\label{9.10:eq-objective-test-rest-1}
F(x):=\sum_{h\in H_{al}}\int\varphi_h\,d\tau_h[x;t^{\mathsf{A}}],
\end{equation}
where $\varphi_h$ is the weight attached to the history $h$ in the objective lemma,
and with its function $\mathsf{G}_z(\lambda,\varepsilon)$ read on the
$z$-fibre of $\hat{\mathbb{D}}^{\star}$. In particular the cluster statement of the objective
lemma holds for the hatted array with its three clauses: (c1) exactness at the ends; (c2)
holomorphy and continuity on $\hat{\mathbb{D}}^{\star}$, together with geometric forgetting of
old regions on $\hat{\mathbb{D}}^{\star}$; (c3) the dial $\lambda_z$, $z=(j,k')$, occurs only in
labels in which zone $j$ is present at level $k'$. Moreover the second claim of the objective
lemma reads
\begin{equation}\label{9.10:eq-objective-test-rest-a}
\sum_{h\in H_{al}}\bigl\|\hat\tau^{\mathsf{B}}_h[t^{\mathsf{A}}]-\hat\tau^{\mathsf{A}}_h\bigr\|
\le \hat E_1^{-1}\,\Sigma^{q_a}_k[\hat\delta]\,A^{\sharp}.
\end{equation}
\item[(ii)] \emph{Test.} The test lemma holds for the hatted lists in the form in which it holds
for the un-hatted lists.
\item[(iii)] \emph{Remainder.} The remainder lemma holds for the hatted lists in the form in
which it holds for the un-hatted lists, and
\begin{equation}\label{9.10:eq-objective-test-rest-2}
4\,q_a^{\,k-n_c}\le C_{nc}\,\rho^{\,n-\mathsf{w}}=C_{nc}\,\hat e_{\mathsf{w}}.
\end{equation}
\end{itemize}
\end{lemma}

\begin{proof}
\emph{(i).} The function $F$ of \eqref{9.10:eq-objective-test-rest-1} is defined on arrays
as the sum over aligned histories $h\in H_{al}$ of the integrals of $\varphi_h$ against the
history integrals $\tau_h[x;t^{\mathsf{A}}]$ at the aligned labels $t^{\mathsf{A}}$, and
$\mathsf{G}_z(\lambda,\varepsilon)$ is taken on the $z$-fibre of $\hat{\mathbb{D}}^{\star}$.
In the objective lemma of the uniqueness theorem the cluster statement is the
cluster-expansion proposition run with one dial per class, on the dial domain of that lemma,
for the assembly $\mathcal{P}_h[x]$ built from the stored slots. In the hatted setting the
same statement is the hatted cluster-expansion proposition (H$_2$) applied to the hatted
array. We check its three clauses.

Clause (c1), exactness at the ends, rests on a statement about one run only; it therefore
holds for the hatted array exactly as for the un-hatted one. Clause (c2): holomorphy and
continuity on $\hat{\mathbb{D}}^{\star}$ are supplied by (H$_2$), and geometric forgetting of
old regions on $\hat{\mathbb{D}}^{\star}$ is hypothesis (H$_3$). Clause (c3), that $\lambda_z$
with $z=(j,k')$ occurs only in labels in which zone $j$ is present at level $k'$, is the
filing statement of Definition~\ref{9.10:def-pair-loads}. It holds for the hatted uses by
Lemma~\ref{9.10:lem-unfiled-range}, in the form \eqref{9.10:eq-unfiled-range-a}: an unfiled
share of level $i$ lies in the zones $[i,i'']$, $i''$ its filing level, and a filed group is
pointed at a point $y\in\Lambda_{i''}^{\circledR}$.

For \eqref{9.10:eq-objective-test-rest-a}: the proof, by moduli, of the upper member of
\cite[(2.50)]{B-CMP119} in the proof of the torus bound within the correlation theorem uses
only the inputs listed in (H$_4$). By (H$_4$) each of these inputs holds a fortiori for the
hatted quantities. Reading the same proof with the hatted inputs therefore gives
\eqref{9.10:eq-objective-test-rest-a}. Since the cluster statement and
\eqref{9.10:eq-objective-test-rest-a} are the parts of the objective lemma that involve the
hatted data, and the remainder of its proof under (H$_1$) is unchanged, the objective lemma
holds for the hatted lists.

\emph{(ii).} The ramps and the tests entering the test lemma, which are scaffold data, do not
depend on $z$: they belong to the scaffold on $\Theta_X$. The statement on moduli on which the
test lemma rests is a statement about one run. Hence neither input changes in passing to the
hatted lists, and the test lemma holds for them as stated in (H$_1$).

\emph{(iii).} The remainder lemma rests on the corollary of the age estimates that bounds the
contribution of old regions, applied in each of the two runs separately. In its hatted form
geometric forgetting of old regions, hypothesis (H$_3$), coincides with the un-hatted statement.
The remainder lemma therefore holds for the hatted lists as stated in (H$_1$). The geometric-rate
arithmetic for the factor $q_a^{-n_c}$, hypothesis (H$_5$), gives the
inequality in \eqref{9.10:eq-objective-test-rest-2}. The equality in
\eqref{9.10:eq-objective-test-rest-2} is the definition $\hat e_r=\rho^{\,n-r}$ read at
$r=\mathsf{w}$.
\end{proof}

\begin{theorem}[Geometric last-lattice comparison for the re-filed run. Stated; proof incomplete
at the step marked $(\ast)$]\label{9.10:thm-last-lattice-comparison}
Let $\mathsf{A}$ and $\mathsf{B}$ be the two runs of Definition~\ref{9.6:def-comparison-chain},
carried with the hatted stored arrays and loads of
Definition~\ref{9.10:def-last-lattice-loads}. Assume the following.
\begin{itemize}
\item[$(\ast)$] The statements listed below hold when read on the hatted array. The statements
are:
\begin{itemize}
\item the statements that supply one particular hypothesis of the envelope proposition of
Definition~\ref{9.10:def-envelope-ceilings};
\item the volume and filament bound, the proposition on the $\mathbf{B}$-terms, the elimination
of the $\mathbf{B}$-loads and the assignment lemma, each read at the hatted loads;
\item the cluster-expansion proposition run with one dial per class for the assembly of the
stored slots;
\item the statement fixing the domain on which the dials run;
\item the age lemma;
\item the random-walk statements, each of which concerns one run;
\item clause (a) of Theorem~\ref{9.10:thm-refiled-lipschitz};
\item the bound on the weights $w^{\mathbf{B}}_{j,i}$ of the $\mathbf{B}$-loads;
\item the smallness statement that governs the arithmetic of $q_a^{-n_c}$;
\item the correlation theorem in its sourced form;
\item the finite-volume statement for the sources $w$.
\end{itemize}
For the first two items, whose statements read the loads and one stored family class by class,
reading on the hatted array means that they hold with the hatted loads and the hatted stored
array in place of the un-hatted ones. The readings on the hatted array of these two items, and
of the bound on the weights $w^{\mathbf{B}}_{j,i}$ with which the second is read, form the step
marked $(\ast)$.
\item[(i)] The second-order smallness condition on the hatted loads at the constant $2\hat E^z_1$,
and the smallness condition on the $\natural$-loads.
\item[(ii)] $g\le\gamma_U$, and the coupling ceilings \eqref{9.10:eq-envelope-ceilings-b} hold.
\item[(iii)] $n\ge n_1$, and $k:=n-\mathsf{w}\ge 1$.
\item[$(\iota)$] $\vartheta_{\chi}\le\rho$.
\end{itemize}
Let $\hat w\in\mathcal{N}'_w$ and $z\in P_{UV}$, and let the residual source be zero. We use the
normalisation of Definition~\ref{9.10:def-last-lattice-loads}, with hatted weights
$\hat x-\hat\zeta$ and centres $\hat\zeta_*$. Then
\begin{equation}\label{9.10:eq-last-lattice-comparison-a}
\begin{aligned}
&\sum_{h\in H_{al}}\bigl\|\hat\tau^{\mathsf{B}}_h(W^{\mathsf{B}}(z))
-\hat\tau^{\mathsf{A}}_h(W^{\mathsf{A}}(z))\bigr\|\\
&\quad+\sum_{h\in H^{\mathsf{A}}\setminus H_{al}}\|\hat\tau^{\mathsf{A}}_h\|
+\sum_{h'\in H^{\mathsf{B}}\setminus H_{al}}\|\hat\tau^{\mathsf{B}}_{h'}\|
\le\hat\varepsilon_n A^{\sharp},\\
&\hat\varepsilon_n:=4\hat E_1^{-1}\Sigma^{q_a}_k[\hat\delta]+4q_a^{k-n_c}
+C_{\mathrm{test}}\vartheta_{\chi}^{k}
\le\hat C_{TC}(g,N_w)\,\hat e_{\mathsf{w}}=\hat C_{TC}\,\rho^{n-\mathsf{w}},
\end{aligned}
\end{equation}
with
\begin{equation}\label{9.10:eq-last-lattice-comparison-1}
\hat C_{TC}:=\frac{4C^{TC}_{\Sigma}\hat C_{\mathrm{env}}\mathsf{l}_{\mathsf{w}}^{2}}{\hat E^z_1}
+C_{\mathrm{nc}}+C_{\mathrm{test}}.
\end{equation}
The constant $\hat C_{TC}$ is independent of $n$, $k$, $K$, of the parameter $m_0$ of the two runs,
and of the position. The bound is geometric in $n$, with no floor.
\end{theorem}

Here $H^{\mathsf{c}}$ is the set of histories of run $\mathsf{c}\in\{\mathsf{A},\mathsf{B}\}$, and
$H_{al}$ is the set of aligned histories. The source of run $\mathsf{c}$ is $W^{\mathsf{c}}(z)$.
The constant $A^{\sharp}$ is that of the torus bound. The parameter $N_w$ is the parameter of the
sources fixed with the finite-volume statement for the sources $w$; the floor $n_0(N_w)$ of
Definition~\ref{9.10:def-last-lattice-loads} and the constant $\hat C_{TC}=\hat C_{TC}(g,N_w)$
depend on it. The load sum $\Sigma^{q_a}_k[\hat\delta]$ is
the one in \eqref{9.10:eq-last-lattice-loads-a}, and $\hat e_r=\rho^{n-r}$ is the envelope letter
of \eqref{9.10:eq-envelope-ceilings-a}. The constants $\hat C_{\mathrm{env}}$, $C^{TC}_{\Sigma}$ and
the factor $\mathsf{l}_{\mathsf{w}}$ are those of Lemma~\ref{9.10:lem-load-envelope}. The constants
$\hat E_1$, $C_{\mathrm{test}}$, $C_{\mathrm{nc}}$ and the exponent $n_c$ are those of the objective,
test and remainder estimates of Lemma~\ref{9.10:lem-objective-test-rest}. Finally, $\hat E^z_1$ is
the constant whose double enters the second-order smallness condition of hypothesis~(i).

\begin{proof}
The argument combines the objective, test and remainder estimates of
Lemma~\ref{9.10:lem-objective-test-rest} with Lemma~\ref{9.10:lem-load-envelope}, whose proof is
incomplete at one step.

At the step of hypothesis $(\ast)$, both lemmas are applied on the hatted array. Their inputs are
the statements listed in the hypothesis marked $(\ast)$, read on the hatted array. This includes
clause (a) of Theorem~\ref{9.10:thm-refiled-lipschitz}, whose proof is incomplete at one step.

\emph{The first inequality.} The three estimates of Lemma~\ref{9.10:lem-objective-test-rest}
contribute the three terms of $\hat\varepsilon_n$, each multiplied by $A^{\sharp}$.
\begin{itemize}
\item The objective estimate \eqref{9.10:eq-objective-test-rest-a} bounds the aligned differences
$\hat\tau^{\mathsf{B}}_h[t^{\mathsf{A}}]-\hat\tau^{\mathsf{A}}_h$, taken at the aligned labels
$t^{\mathsf{A}}$, by $\hat E_1^{-1}\Sigma^{q_a}_k[\hat\delta]A^{\sharp}$. It contributes the term
$4\hat E_1^{-1}\Sigma^{q_a}_k[\hat\delta]$.
\item The remainder estimate contributes the term $4q_a^{k-n_c}$. It rests on the corollary of the
age lemma, applied in each run separately, with geometric forgetting of old regions taking the
same form on the hatted array, and on the arithmetic of $q_a^{-n_c}$.
\item The test estimate contributes the term $C_{\mathrm{test}}\vartheta_{\chi}^{k}$. The ramps and
test functions entering it belong to the prescribed scaffold $\Theta_X$ and do not depend on $z$,
and the random-walk statement it uses concerns one run only.
\end{itemize}

\emph{The second inequality.} We bound the three terms of $\hat\varepsilon_n$ in turn.

For the first term, \eqref{9.10:eq-load-envelope-a} in Lemma~\ref{9.10:lem-load-envelope}, whose
proof is incomplete at one step, gives
$\Sigma^{q_a}_k[\hat\delta]\le C^{TC}_{\Sigma}\hat C_{\mathrm{env}}\mathsf{l}_{\mathsf{w}}^2
\hat e_{\mathsf{w}}$. Hence
\begin{equation*}
4\hat E_1^{-1}\Sigma^{q_a}_k[\hat\delta]
\le\frac{4C^{TC}_{\Sigma}\hat C_{\mathrm{env}}\mathsf{l}_{\mathsf{w}}^2}{\hat E_1}\,
\hat e_{\mathsf{w}}.
\end{equation*}
The first summand of $\hat C_{TC}$ in \eqref{9.10:eq-last-lattice-comparison-1} carries
$\hat E^z_1$ in place of $\hat E_1$; the passage from this bound to that summand uses the relation
between $\hat E_1$ and $\hat E^z_1$ that is fixed together with the second-order smallness
condition of hypothesis~(i).

For the second term, the remainder estimate of Lemma~\ref{9.10:lem-objective-test-rest} gives
$4q_a^{k-n_c}\le C_{\mathrm{nc}}\rho^{n-\mathsf{w}}=C_{\mathrm{nc}}\hat e_{\mathsf{w}}$.

For the third term, we use hypothesis $(\iota)$, the identity $k=n-\mathsf{w}$ and the definition
$\hat e_{\mathsf{w}}=\rho^{n-\mathsf{w}}$ of \eqref{9.10:eq-envelope-ceilings-a}. They give
\begin{equation*}
C_{\mathrm{test}}\vartheta_{\chi}^{k}\le C_{\mathrm{test}}\rho^{k}
=C_{\mathrm{test}}\rho^{n-\mathsf{w}}
=C_{\mathrm{test}}\hat e_{\mathsf{w}}.
\end{equation*}

Adding the three bounds gives
$\hat\varepsilon_n\le\hat C_{TC}\hat e_{\mathsf{w}}$, with $\hat C_{TC}$ as in
\eqref{9.10:eq-last-lattice-comparison-1}.
\end{proof}

\begin{corollary}[Comparison of the normalised partition functions]
Under the hypotheses of Theorem~\ref{9.10:thm-last-lattice-comparison},
\begin{equation}\label{9.10:eq-last-lattice-comparison-2}
\bigl|\tilde Z^{\mathsf{B}}(W^{\mathsf{B}}(z))-\tilde Z^{\mathsf{A}}(W^{\mathsf{A}}(z))\bigr|
\le\hat\varepsilon^{\star}_n A,\qquad
\hat\varepsilon^{\star}_n\le\hat C^{\star}_{TC}(N_w;g)\,\hat e_{\mathsf{w}},
\end{equation}
where
\begin{equation*}
\hat C^{\star}_{TC}:=e^{(E_+^{\sharp\sharp}-E_+^{\sharp})N_w^4}\,\hat C_{TC}.
\end{equation*}
Here $\tilde Z^{\mathsf{c}}$ is the normalised partition function of run $\mathsf{c}$, and $A$ is
the constant of the comparison of total masses from which the corollary is derived (not the action
$A(U)$). The letters $E_+^{\sharp}$ and $E_+^{\sharp\sharp}$ are two values of the constant $E_+$
of the Glossary (Notation~\ref{0.2:not-glossary}), fixed with the finite-volume statement for the
sources $w$; the factor $e^{(E_+^{\sharp\sharp}-E_+^{\sharp})N_w^4}$ is the factor by which the
passage from the history integrals to the total masses enlarges $\hat C_{TC}$.
\end{corollary}

\begin{proof}
The first line of the finite-volume statement for the sources $w$ bounds the difference of the
normalised partition functions of the two runs by the total masses of the history integrals.
It is one of the statements of the hypothesis marked $(\ast)$.

We apply it to the three sums on the left of \eqref{9.10:eq-last-lattice-comparison-a}, using
Theorem~\ref{9.10:thm-last-lattice-comparison}, whose proof is incomplete at one step. This gives
\eqref{9.10:eq-last-lattice-comparison-2}. The factor
$e^{(E_+^{\sharp\sharp}-E_+^{\sharp})N_w^4}$ multiplies $\hat C_{TC}$ in $\hat C^{\star}_{TC}$.
\end{proof}

\begin{corollary}[Comparison of the logarithmic derivatives at the origin]
Under the hypotheses of Theorem~\ref{9.10:thm-last-lattice-comparison}, let
$w=\sum_{i=1}^{\mathsf{p}}z_i f_i$. Then
\begin{equation}\label{9.10:eq-last-lattice-comparison-3}
\Bigl|\partial_{z_1}\cdots\partial_{z_{\mathsf{p}}}\big|_{0}
\bigl[\operatorname{Log}\hat{\mathsf{h}}^{\mathsf{B}}-\operatorname{Log}\hat{\mathsf{h}}^{\mathsf{A}}\bigr]
\Bigr|
\le\hat C_*(N_w)\,\hat e_{\mathsf{w}}\Bigl(\frac{2\mathsf{p}}{R_1}\Bigr)^{\mathsf{p}}
\prod_{i=1}^{\mathsf{p}}\|f_i\|_{\sharp},
\end{equation}
where
\begin{equation*}
\hat C_*(N_w):=2\hat C^{\star}_{TC}\,\mathsf{Q}_w(1+\mathsf{Q}_w).
\end{equation*}
Here $\hat{\mathsf{h}}^{\mathsf{c}}$ is the ratio attached to run $\mathsf{c}$, $\operatorname{Log}$
is the principal logarithm, and $R_1$ is the constant that enters through Cauchy's estimates on
the polydisc $P_{IR}$.
\end{corollary}

\begin{proof}
The bound is derived from \eqref{9.10:eq-last-lattice-comparison-2} of the preceding corollary in
four steps.
\begin{enumerate}
\item The Schwarz lemma is applied on complex lines in the source variables. This carries the
bound \eqref{9.10:eq-last-lattice-comparison-2} over to the ratios
$\hat{\mathsf{h}}^{\mathsf{B}}$ and $\hat{\mathsf{h}}^{\mathsf{A}}$.
\item Each ratio is controlled by the floor of its own run.
\item On the resulting range the principal logarithm is Lipschitz with constant $2$.
\item Cauchy's estimates on the polydisc $P_{IR}$ convert the resulting bound on
$\operatorname{Log}\hat{\mathsf{h}}^{\mathsf{B}}-\operatorname{Log}\hat{\mathsf{h}}^{\mathsf{A}}$
into the bound \eqref{9.10:eq-last-lattice-comparison-3} on the mixed derivative at the origin.
\end{enumerate}
The factor
$\hat C^{\star}_{TC}\hat e_{\mathsf{w}}$ comes from the preceding corollary, which rests on
Theorem~\ref{9.10:thm-last-lattice-comparison}, whose proof is incomplete at one step.
\end{proof}

The comparison of the two normalisations involves three pieces: the vacuum-energy piece at the
centre, the piece attached to the floor $n_2$, and the pieces of the one-point factor
$\mathfrak{z}$. In this comparison the gauge-fixing coefficient $a_j(y)=\hat x_j-\hat\zeta_j(p_y)$
is treated as a weight slot of class $(j,k)$; its two-run difference is measured by the norm
$\hat{\mathsf{z}}_j$ and bounded by the statement on the running-shift differences; and the move of
the dial of this class per cube is at most a constant times $E_1$, the un-hatted form of the
constant $\hat E_1$, under the coupling ceiling with exponent $q_{\flat}$
\begin{equation*}
A_0^2(\log X)^{-2(q_{\flat}-p_0)}\le 1.
\end{equation*}

\begin{lemma}[Localisation of the vacuum-energy difference]\label{9.10:lem-vacuum-localisation}
Let $\mathsf{A}$ and $\mathsf{B}$ be the two runs compared in this chapter, each carried by local lists of
the re-filed renormalisation run of Definition~\ref{9.10:def-refiled-run}. Let $W^{\mathsf{A}}(z)$ and
$W^{\mathsf{B}}(z)$ be their sources, and let $P_{UV}$ be the ultraviolet polydisc of the sources. Assume:
\begin{itemize}
\item[(a)] the inputs \textup{(J1)}--\textup{(J4)} of Corollary~\ref{9.10:cor-scaffold-inputs} for the
hatted run;
\item[(b)] the hypotheses of Proposition~\ref{9.10:prop-geometric-envelopes}, for local lists, including
the coupling ceilings \eqref{9.10:eq-envelope-ceilings-b};
\item[(c)] the localisation lemma for the vacuum energy in the proof of the uniqueness theorem, whose
constant is given as a function $C_\eta(C_E,C_{PM})$ of the envelope constant $C_E$ and the one-point
constant $C_{PM}$ of that lemma; its proof is followed below with the substitutions indicated;
\item[(d)] the conclusions \eqref{9.10:eq-geometric-envelopes-a} of
Proposition~\ref{9.10:prop-geometric-envelopes};
\item[(e)] the bound on the one-point factor $\mathfrak{z}$ in the form given in
Remark~\ref{9.10:rem-downstream-constants}, with its constant $\hat C_{PM}$.
\end{itemize}
Put $\hat C_\eta:=C_\eta(\hat C_E,\hat C_{PM})$. Then for
$\mathsf{w}\le r<n$ there are a number $c_r$ independent of $z$ and functions $G_r(Y;z)$, indexed by
all localisation domains $Y$ of the depth-$r$ torus, holomorphic in a neighbourhood of $P_{UV}$ and
depending on $z$ only through the restriction of $W^{\mathsf{A}}(z)$ to the enlargement $Y^{+}$ of $Y$
through which the localisation lemma named in (c) reads the source, such that
\begin{equation}\label{9.10:eq-vacuum-localisation-1}
\sum_X e^{(r),\mathsf{B}}\bigl(X;W^{\mathsf{B}}(z)\bigr)-\sum_X e^{(r),\mathsf{A}}\bigl(X;W^{\mathsf{A}}(z)\bigr)
=\sum_Y G_r(Y;z)+c_r
\end{equation}
and
\begin{equation}\label{9.10:eq-vacuum-localisation-a}
\sup_{z\in P_{UV}}\lvert G_r(Y;z)\rvert\le \hat C_\eta\,(1+r)\,\hat e_r\,e^{-(1-\delta)\kappa^{c}d(Y)} .
\end{equation}
Here $d(Y)$ is the size of the domain $Y$ as it enters the family norm $\|\cdot\|^{c}$, whose decay rate
is $\kappa^{c}$, and $\delta$ is the fraction of this decay spent on the count of points in the proof
named in (c).
\end{lemma}

\begin{proof}
We follow the proof of the localisation lemma named in (c), replacing every object by its counterpart
for the re-filed run and every input bound by the one listed below.

Fix one of the two runs and a localisation domain $X$ of the depth-$r$ torus. The vacuum pieces are
decomposed as
\begin{equation}\label{9.10:eq-vacuum-localisation-2}
e^{(r)}(X;w)=\bar e^{(r)}(X;w)-\bar e^{(r)}(X;0)+\bigl[\text{one-point pieces}\bigr],
\end{equation}
where the one-point pieces are the $\mathfrak{z}$-pieces $\log[\mathfrak{z}(a_r(y))/\mathfrak{z}(\hat x_r)]$,
$y\in X$, of Remark~\ref{9.10:rem-healing-last-lattice}, and
\begin{equation}\label{9.10:eq-vacuum-localisation-3}
\bar e^{(r)}(X;w):=\sum_{y\in X}\hat E^{(r)}(X,1,y;w)+R^{(r)}(X,1;w).
\end{equation}
Two facts give this decomposition the same form as in the proof named in (c). First, the value of
$\hat E^{(r)}$ at the configuration $1$ equals the value at $1$ of the regular family $E^{(r)}$ of the
correlation theorem. The reason is that the re-filed pointed family vanishes there:
$\mathsf{RF}(Y,1,y)=0$. Second, $R^{(r)}(X,1;w)$ is the vacuum value of the remnant created at depth $r$.
By Definition~\ref{9.10:def-refiled-run}, a remnant is created exactly as in the correlation theorem,
and its vacuum value $R^{(r)}(X,1;w)$ is one of the terms of $e^{(r)}(X;w)$. By
Remark~\ref{9.10:rem-healing-last-lattice}, this holds whether or not the remnant is filed later,
including the remnants on the last lattice that are never filed. Hence the term $R^{(r)}(X,1;w)$ is
present at every depth.

Next, the hatted regular family $\hat E^{(r)}$ splits into its pure part, its dirty part and its
freezing-defect part:
\begin{equation}\label{9.10:eq-vacuum-localisation-4}
\hat E^{(r)}(\cdot\,;W)=P^{\flat(r)}_{W(p_y)}+\hat Q^{(r)}_{W(p_y)}+\hat D^{(r)}(\cdot\,;W).
\end{equation}
Here $W(p_y)$ is the value of the source at the plaquette $p_y$ attached to $y$. From these three parts
and from the remnant, the proof named in (c) builds localised pieces $G^{P^\flat}_r$, $G^{\hat D}_r$,
$G^{\hat Q}_r$ and $G^R_r$. In that proof the pure and freezing-defect pieces carry the weight
$E_0\,\mathsf{l}_r$ and the dirty and remnant pieces the weight $(C_t\theta_r)^{-1}$, where $E_0$ and
$C_t$ are constants of that proof; $v$ is the source slot of the localised pieces, and $\omega(p_y)$ is
the value of the source datum of that proof at the plaquette $p_y$, at which the slot is evaluated.
We put
\begin{align*}
G_r(Y;z):={}&\sum_{y\in Y}\Bigl[E_0\,\mathsf{l}_r\bigl(G^{P^\flat}_r\big|_{v=\omega(p_y)}+G^{\hat D}_r\bigr)\\
&\quad+(C_t\theta_r)^{-1}G^{\hat Q}_r\big|_{v=\omega(p_y)}\Bigr](Y,1,0,y;z)\\
&+(C_t\theta_r)^{-1}G^R_r(Y,1,0;z)
+\bigl[\text{difference of the }\mathfrak{z}\text{-pieces of }\mathsf{B}\text{ and }\mathsf{A}\bigr].
\end{align*}
The $\mathfrak{z}$-pieces are evaluated at $a_r(y)=\hat x_r-\hat\zeta_{(r)}(p_y)$ and bounded by the
one-point bound (e). Each term is holomorphic in a neighbourhood of $P_{UV}$ and reads $z$ through
$W^{\mathsf{A}}(z)$ restricted to $Y^{+}$, as in the proof named in (c). The terms of the difference
that do not depend on $z$ form $c_r$, as there. Summing
\eqref{9.10:eq-vacuum-localisation-2}--\eqref{9.10:eq-vacuum-localisation-3} over $X$ for both runs
therefore gives \eqref{9.10:eq-vacuum-localisation-1}.

The sizes of the pieces are supplied by \eqref{9.10:eq-geometric-envelopes-a}:
\begin{itemize}
\item the two components controlled by $g^A_r$ contribute together at most
$2\cdot 4\hat C_E\hat e_r$;
\item the dirty part $\hat Q$ contributes at most $\hat C_E\hat e_r$;
\item the remnant contributes at most $4\hat C_E\hat{\mathfrak{r}}_r\hat e_r$, where
$\hat{\mathfrak{r}}_r$ is the envelope letter of \eqref{9.10:eq-envelope-ceilings-a}, and
$4\hat{\mathfrak{r}}_r\le\tfrac14$.
\end{itemize}
Moreover, with $g^{-2}$ the inverse coupling at which the prescribed scaffold is taken,
$\theta_r=g^{-2}+c_0r\ge g^{-2}$, so that $(C_t\theta_r)^{-1}\le (C_tg^{-2})^{-1}$. The number of points
of $Y$ is counted, with the lattice-counting constants $C_d$ and $C_\delta$ of the proof named in (c), by
\begin{equation}\label{9.10:eq-vacuum-localisation-5}
\#\{y\in Y\}\le M^4C_dC_\delta\,e^{\delta\kappa^{c}d(Y)}.
\end{equation}
These sizes and the count \eqref{9.10:eq-vacuum-localisation-5} replace the corresponding inputs of the
estimate in the proof named in (c). Its constant is then $C_\eta(\hat C_E,\hat C_{PM})$, that is
$\hat C_\eta$. With this replacement the estimate gives
\eqref{9.10:eq-vacuum-localisation-a}.
\end{proof}

\begin{theorem}[Exponential rate for the renormalised-coupling family]\label{9.10:thm-exponential-rate}
Objects of the re-filed run carry a hat; $n$ denotes its length, $\mathsf{w}$ the last-step offset,
$\mathsf{p}$ the number of test functions $f=(f_1,\dots,f_{\mathsf{p}})$, taken as in the uniqueness
theorem, with supports pairwise separated by $\mathsf{d}$, $X=g^{-2}$, and
$\theta_r=X+c_0r$ the points of the prescribed scaffold $\Theta_X$. The following are the objects of
the proof of the uniqueness theorem, read for the re-filed run:
\begin{itemize}
\item $G_n$, the quantity of run length $n$ compared by the step lemma of that proof, a function of
the label. Written without argument, it is taken at the bare label of that step lemma, as opposed
to the normalised and re-centred labels of (c) and (d) below. $G_n(0)$ denotes its value at the
argument $0$;
\item $\hat{\mathsf{h}}$, the ratio entering that step lemma whose logarithm is estimated separately
there;
\item the constants of that step lemma: $M_G$ and $\upsilon_w$, which enter the coefficient
$M_G[2\varrho_2/\upsilon_w+\mathsf{p}2^{\mathsf{p}}]$ of $\Upsilon_n$ in its bare-label estimate;
$C^w_1$, whose defining sum over depths is taken here with the additional
factor $1+r$ of \eqref{9.10:eq-vacuum-localisation-a} (see the proof); $C^w_2$ and the decay
constant $\bar c$ of its run-by-run estimate, where the letter $\bar c$ distinguishes this constant
from the exponent $c$ in (d); the radii $R_0$ and $R_1$, which enter through the factors
$(2\mathsf{p}/R_0)^{\mathsf{p}}$ and $(2\mathsf{p}/R_1)^{\mathsf{p}}$; the parameter $N_w$ and the
threshold $n_0(N_w)$ depending on it; the radius $r_J$ bounding the parameter $t$; and the constant
$D$ of that step lemma, on which the constant $C(\mathsf{p},\mathsf{d},D,N_w)$ of
\eqref{9.10:eq-exponential-rate-a} depends;
\item the centres $\hat\zeta^n_*(t)$ and the bookkeeping function $\hat{\mathfrak{n}}_n$;
\item the normalisations $\hat{\mathsf{N}}_n$;
\item the Schwinger function $S^n(a;f)$ at run length $n$ and label $a$.
\end{itemize}
The letter $\hat C_{\eta}$ is the constant of Lemma~\ref{9.10:lem-vacuum-localisation}; the letter
$\hat C_*(N_w)$ is the constant multiplying $\rho^{n-\mathsf{w}}$ in the last-lattice estimate
\eqref{9.10:eq-last-lattice-comparison-a}; and $\Upsilon_n$ is as in
Remark~\ref{9.10:rem-downstream-constants}.

Assume:
\begin{enumerate}
\item[(i)] the inputs for the re-filed run of Corollary~\ref{9.10:cor-scaffold-inputs};
\item[(ii)] the three-slot hybrid bounds of Corollary~\ref{9.10:cor-three-slot-bounds};
\item[(iii)] the geometric source bounds of Proposition~\ref{9.10:prop-source-bounds};
\item[(iv)] the coupling ceilings of Definition~\ref{9.10:def-envelope-ceilings}, with the envelopes
$\hat e_r=\rho^{n-r}$ of \eqref{9.10:eq-envelope-ceilings-a};
\item[(v)] the localisation \eqref{9.10:eq-vacuum-localisation-a} of
Lemma~\ref{9.10:lem-vacuum-localisation};
\item[(vi)] the comparison \eqref{9.10:eq-last-lattice-comparison-a} of
Theorem~\ref{9.10:thm-last-lattice-comparison};
\item[(vii)] the statements of (viii), and the statements on which
Corollary~\ref{9.10:cor-three-slot-bounds}, Proposition~\ref{9.10:prop-source-bounds},
Definition~\ref{9.10:def-envelope-ceilings}, Lemma~\ref{9.10:lem-vacuum-localisation} and
Theorem~\ref{9.10:thm-last-lattice-comparison} rest, hold for the objects of the re-filed run in
place of those of the original run;
\item[(viii)] the statements on which the proof of the uniqueness theorem rests; among them are the
increment statement and the extension proposition.
\end{enumerate}
Hypothesis (iii) is the conclusion of Proposition~\ref{9.10:prop-source-bounds}, whose proof is
incomplete at one step. Hypothesis (vi) is the conclusion of
Theorem~\ref{9.10:thm-last-lattice-comparison}, whose proof is incomplete at one step.
Hypothesis (i) is the conclusion of Corollary~\ref{9.10:cor-scaffold-inputs}, which is proved as an
implication; among its antecedents is the induction theorem for the re-filed run, whose proof is
outstanding.

Put
\begin{equation}\label{9.10:eq-exponential-rate-1}
\begin{split}
\hat{\mathfrak{f}}_n := {}&\Bigl[C^w_1\hat C_{\eta}\,\rho^{n/2}
+C^w_2\,e^{-\bar c\,\mathsf{d}\,L^{n/2-\mathsf{w}}/M}\Bigr]
\Bigl(\frac{2\mathsf{p}}{R_0}\Bigr)^{\mathsf{p}}\\
&+\hat C_*(N_w)\,\rho^{n-\mathsf{w}}\Bigl(\frac{2\mathsf{p}}{R_1}\Bigr)^{\mathsf{p}}.
\end{split}
\end{equation}
Then, for every run length $n$ to which the step lemma of the uniqueness theorem applies, the
following hold.
\begin{enumerate}
\item[(a)] The estimate of that step lemma holds for the re-filed run with its step bound replaced
by $\hat{\mathfrak{f}}_n$.
\item[(b)] For the bare label,
\begin{equation}\label{9.10:eq-exponential-rate-2}
|G_{n+1}-G_n|\le M_G\Bigl[\frac{2\varrho_2}{\upsilon_w}+\mathsf{p}\,2^{\mathsf{p}}\Bigr]\Upsilon_n
+\hat{\mathfrak{f}}_n\prod_{i=1}^{\mathsf{p}}\|f_i\|_{\sharp}.
\end{equation}
\item[(c)] One has $\hat{\mathfrak{n}}_n'(0)=\hat{\mathsf{N}}_n/\hat{\mathsf{N}}_{n+1}$, and
\begin{equation}\label{9.10:eq-exponential-rate-3}
\bigl|\hat{\mathsf{N}}_{n+1}^{-\mathsf{p}}G_{n+1}(0)-\hat{\mathsf{N}}_n^{-\mathsf{p}}G_n(0)\bigr|
\le\Bigl(\frac32\Bigr)^{\mathsf{p}}\hat{\mathfrak{f}}_n\prod_{i=1}^{\mathsf{p}}\|f_i\|_{\sharp}.
\end{equation}
\item[(d)] Let $t$ be real with $|t|\le r_J/4$. Consider the run whose inverse coupling at depth
zero is exactly $\hat x_{(0)}=X-t$. Put
\[
\hat a_n(t):=\theta_n-\hat\zeta^n_*(t),
\]
and let $\hat{\mathsf{N}}_n(t)$ be the derivative $\partial\hat\zeta_{(0)}/\partial\zeta$ of the
depth-zero running shift, evaluated at the centre $\hat\zeta^n_*(t)$; it is positive. Then, on every
torus and for
every $n\ge n_0(N_w)$, with an exponent $c\ge0$ independent of $n$, fixed by the bookkeeping of
$\mathsf{w}$ and of the factors $1+r$ in the step lemma,
\begin{equation}\label{9.10:eq-exponential-rate-a}
\begin{split}
\bigl|\hat{\mathsf{N}}_{n+1}(t)^{-\mathsf{p}}S^{n+1}(\hat a_{n+1}(t);f)
&-\hat{\mathsf{N}}_n(t)^{-\mathsf{p}}S^{n}(\hat a_{n}(t);f)\bigr|\\
&\le C(\mathsf{p},\mathsf{d},D,N_w)\prod_{i=1}^{\mathsf{p}}\|f_i\|_{\sharp}\;n^{c}\rho^{n/2}.
\end{split}
\end{equation}
\item[(e)] Let $t$ be as in (d). Write $K$ for the run length and $S^K_{\mathsf{p}}(a)$ for
$S^K(a;f)$. There exist
$\hat a_K(t)$ and $\hat{\mathsf{N}}_K(t)>0$ such that
$\hat{\mathsf{N}}_K(t)^{-\mathsf{p}}S^K_{\mathsf{p}}(\hat a_K(t))$ converges, exponentially in $K$,
to a limit $S^{[t]}_{\mathsf{p}}$. Fix a pair $f',f''$ of test functions, write $S^K_2(a)$ for
$S^K(a;f',f'')$, and let $S^{[t]}_2$ denote the corresponding limit. Where $S^{[t]}_2>0$, the ratios
\[
S^K_{\mathsf{p}}(\hat a_K(t))\big/S^K_2(\hat a_K(t))^{\mathsf{p}/2}
\]
are defined for all $K$ large enough, are free of $\hat{\mathsf{N}}$ and converge exponentially
in $K$ to $S^{[t]}_{\mathsf{p}}/(S^{[t]}_2)^{\mathsf{p}/2}$.
\end{enumerate}
\end{theorem}

\begin{proof}
Hypotheses (i)--(vi) are in turn consequences of the statements of (viii) read under (vii).
We run the proof of the step lemma of the uniqueness theorem, which hypothesis (viii) makes
available, with two replacements: the localisation of Lemma~\ref{9.10:lem-vacuum-localisation} at
the depths $\mathsf{w}\le r\le\lfloor n/2\rfloor$, and the comparison
\eqref{9.10:eq-last-lattice-comparison-a} of hypothesis (vi) (the conclusion of
Theorem~\ref{9.10:thm-last-lattice-comparison}, whose proof is incomplete at one step) for the
logarithm of $\hat{\mathsf{h}}$. Every other part
of that proof is used unchanged.

\emph{The first summand of \eqref{9.10:eq-exponential-rate-1}.}
Lemma~\ref{9.10:lem-vacuum-localisation} is available at every depth $r$ with
$\mathsf{w}\le r<n$; we use it at the depths $\mathsf{w}\le r\le\lfloor n/2\rfloor$ and keep the
run-by-run estimate of the step lemma above $\lfloor n/2\rfloor$. The
hypotheses of Lemma~\ref{9.10:lem-vacuum-localisation} are supplied by (i)--(iv), where (iii) is the
conclusion of
Proposition~\ref{9.10:prop-source-bounds}, whose proof is incomplete at one step; the conclusion of
Lemma~\ref{9.10:lem-vacuum-localisation} is
hypothesis (v). By
\eqref{9.10:eq-vacuum-localisation-a}, the difference of the vacuum-energy sums of the two
machines at depth $r$ is a sum of localised pieces and a source-free constant. Each piece is
bounded by $\hat C_{\eta}(1+r)\hat e_r$ times a decay factor in its domain.

At the depths $\mathsf{w}\le r\le\lfloor n/2\rfloor$ we have
$n-r\ge n-\lfloor n/2\rfloor\ge n/2$. Since $0<\rho<1$, this gives
\[
\hat e_r=\rho^{n-r}\le\rho^{n/2}.
\]
Writing $r=\mathsf{w}+j$, the $j$-th term of the sum over these depths that defines $C^w_1$
carries the factor $1+r=1+\mathsf{w}+j$ of
\eqref{9.10:eq-vacuum-localisation-a}. With the factor $(2\mathsf{p}/R_0)^{\mathsf{p}}$ of the
step lemma, these depths contribute
$C^w_1\hat C_{\eta}\rho^{n/2}(2\mathsf{p}/R_0)^{\mathsf{p}}$.

\emph{The second summand.}
The depths above $\lfloor n/2\rfloor$ are estimated run by run, exactly as in the step lemma. This
gives $C^w_2e^{-\bar c\,\mathsf{d}L^{n/2-\mathsf{w}}/M}(2\mathsf{p}/R_0)^{\mathsf{p}}$, with
$C^w_2$ and $\bar c$ the constants of the step lemma.

\emph{The third summand.}
The part of the estimate coming from the logarithm of the ratio $\hat{\mathsf{h}}$ is bounded by
\eqref{9.10:eq-last-lattice-comparison-a}, which is hypothesis (vi). This is the conclusion of
Theorem~\ref{9.10:thm-last-lattice-comparison}, whose proof is incomplete at one step. It gives
$\hat C_*(N_w)\rho^{n-\mathsf{w}}(2\mathsf{p}/R_1)^{\mathsf{p}}$.

\emph{Proof of (a).}
The three contributions add up to \eqref{9.10:eq-exponential-rate-1}, which proves (a). Compared
with the step bound of the uniqueness theorem, the factor $\rho^{n/2}+\Pi_n$ has become
$\rho^{n/2}$ and the factor $\rho^{n-\mathsf{w}}+\Pi_n$ has become $\rho^{n-\mathsf{w}}$. The floor
$\Pi_n$ does not occur, because the envelopes $\hat e_r=\rho^{n-r}$ of
\eqref{9.10:eq-envelope-ceilings-a} carry none.

\emph{Proof of (b).}
In the step lemma the bound on $|G_{n+1}-G_n|$ for the bare label is the sum of two terms:
$M_G[2\varrho_2/\upsilon_w+\mathsf{p}2^{\mathsf{p}}]\Upsilon_n$, and the step bound times
$\prod_i\|f_i\|_{\sharp}$. With the step bound replaced by $\hat{\mathfrak{f}}_n$ through (a), this is
\eqref{9.10:eq-exponential-rate-2}. The quantity $\Upsilon_n$ is the one of
Remark~\ref{9.10:rem-downstream-constants}, and its order there is unchanged. Estimate (b) is
therefore governed by $\Upsilon_n$ and is not geometric in $n$; the geometric rate enters only
after normalisation, in (c).

\emph{Proof of (c).}
By Remark~\ref{9.10:rem-downstream-constants}, the proposition on the centres
$\hat\zeta^n_*(t)$ in the proof of the uniqueness theorem holds for the re-filed run. That
proposition reads the envelope bounds only through their values, which are now $\hat e_r$ with
$\Pi_n\equiv0$. The room in the shift variable that this proposition requires is supplied by the
input of Corollary~\ref{9.10:cor-scaffold-inputs} on the running shift $\hat\zeta$. The proposition
gives
$\hat{\mathfrak{n}}_n'(0)=\hat{\mathsf{N}}_n/\hat{\mathsf{N}}_{n+1}$. The normalised estimate of the
step lemma, with the step bound $\hat{\mathfrak{f}}_n$ of (a), is then
\eqref{9.10:eq-exponential-rate-3}.

\emph{Proof of (d).}
Fix a real $t$ with $|t|\le r_J/4$. Apply (c) to the run with $\hat x_{(0)}=X-t$, at the label
$\hat a_n(t)=\theta_n-\hat\zeta^n_*(t)$ and with the normalisation
$\hat{\mathsf{N}}_n(t)=\partial\hat\zeta_{(0)}/\partial\zeta$. The positivity of
$\hat{\mathsf{N}}_n(t)$ is taken by statement from the proof of the uniqueness theorem, in which the
normalisation at the centre is positive; hypothesis (viii), read for the re-filed run as in (vii),
makes it available here. This is the passage from the step
lemma to the uniqueness theorem, carried out with $\hat{\mathfrak{f}}_n$ in place of the step bound.
Like (c), the resulting bound holds on every torus; the data of (i) entering it are torus-free by
Corollary~\ref{9.10:cor-scaffold-inputs}.

It remains to bound $(3/2)^{\mathsf{p}}\hat{\mathfrak{f}}_n$. Its first summand carries
$\rho^{n/2}$. Its second carries $e^{-\bar c\,\mathsf{d}L^{n/2-\mathsf{w}}/M}$. Its third carries
$\rho^{n-\mathsf{w}}$. For $n\ge n_0(N_w)$ the bookkeeping of $\mathsf{w}$ in the proof of the
uniqueness theorem absorbs two dependences into a factor $n^{c}$: the dependence of these three
summands on $\mathsf{w}$, and the factors $1+r$ carried by $C^w_1$. This leaves the rate
$\rho^{n/2}$ and a constant $C(\mathsf{p},\mathsf{d},D,N_w)$, which is
\eqref{9.10:eq-exponential-rate-a}.

\emph{Proof of (e).}
Take $\hat a_K(t)$ and $\hat{\mathsf{N}}_K(t)$ as in (d); by (d), $\hat{\mathsf{N}}_K(t)>0$. Put
$u_K:=\hat{\mathsf{N}}_K(t)^{-\mathsf{p}}S^K_{\mathsf{p}}(\hat a_K(t))$. By
\eqref{9.10:eq-exponential-rate-a}, for $K_2>K_1\ge n_0(N_w)$,
\begin{align*}
|u_{K_2}-u_{K_1}|&\le C(\mathsf{p},\mathsf{d},D,N_w)\prod_{i=1}^{\mathsf{p}}\|f_i\|_{\sharp}
\sum_{n\ge K_1}n^{c}\rho^{n/2}.
\end{align*}
Writing $n=K_1+j$ and using $(K_1+j)^c\le K_1^c(1+j)^c$ for $K_1\ge1$ and $c\ge0$, we get
\begin{align*}
\sum_{n\ge K_1}n^{c}\rho^{n/2}&\le K_1^{c}\rho^{K_1/2}\sum_{j\ge0}(1+j)^{c}\rho^{j/2}.
\end{align*}
The last series converges because $\rho<1$. So $(u_K)$ is Cauchy. Letting $K_2\to\infty$ at
$K_1=K$, its limit $S^{[t]}_{\mathsf{p}}$
satisfies
\[
|u_K-S^{[t]}_{\mathsf{p}}|\le C'K^{c}\rho^{K/2},
\]
with
\[
C'=C(\mathsf{p},\mathsf{d},D,N_w)\prod_{i=1}^{\mathsf{p}}\|f_i\|_{\sharp}
\sum_{j\ge0}(1+j)^{c}\rho^{j/2}.
\]
This decays exponentially in $K$.

For the ratios, assume $S^{[t]}_2>0$. Put $v_K:=\hat{\mathsf{N}}_K(t)^{-2}S^K_2(\hat a_K(t))$, with
$S^K_2$ taken at the pair $f',f''$ fixed in (e), the limit $S^{[t]}_2$ being taken at that pair.
Since $v_K\to S^{[t]}_2$ at the rate just shown, we have $v_K\ge S^{[t]}_2/2>0$ for $K$ large; since
$\hat{\mathsf{N}}_K(t)>0$, also $S^K_2(\hat a_K(t))>0$ for these $K$, so the ratio of (e) is defined
there. Note that
\[
\hat{\mathsf{N}}_K^{-\mathsf{p}}=(\hat{\mathsf{N}}_K^{-2})^{\mathsf{p}/2},
\]
so for these $K$ the ratio of (e) equals $u_K/v_K^{\mathsf{p}/2}$ and is therefore free of
$\hat{\mathsf{N}}$. On
$[S^{[t]}_2/2,\infty)$ the map $v\mapsto v^{-\mathsf{p}/2}$ is Lipschitz. The two
convergences therefore give
\[
\bigl|u_K/v_K^{\mathsf{p}/2}-S^{[t]}_{\mathsf{p}}/(S^{[t]}_2)^{\mathsf{p}/2}\bigr|
\le C''K^{c}\rho^{K/2}
\]
for $K$ large, with a constant $C''$ independent of $K$, depending only on the constants of the two
convergences, on $S^{[t]}_2$, $S^{[t]}_{\mathsf{p}}$ and $\mathsf{p}$. This is again exponential in
$K$.
\end{proof}

\begin{corollary}[The bi-Lipschitz lemma for the re-filed run]\label{9.10:cor-refiled-bilipschitz}
Assume the following two statements.
\begin{itemize}
\item[(a)] The corresponding conclusion of the induction theorem for the re-filed run
(Theorem~\ref{9.10:thm-refiled-induction}), in the form: for the run of every member of an admissible
family in the sense of the uniqueness theorem,
\begin{equation}\label{9.10:eq-refiled-bilipschitz-1}
\sum_{k<K}\operatorname{Lip}_{\bar D(0,\varrho)}\hat b_{k+1}\le\frac23 ,
\end{equation}
where $K$ is the number of steps of the run, $\hat b_{k+1}$ is the coupling increment of step $k+1$ of
the re-filed run \eqref{9.10:eq-refiled-run-b}, and $\bar D(0,\varrho)$ is the closed disc of radius
$\varrho$ about $0$ in $\mathbb{C}$.
\item[(b)] The summability $\sum_n\Upsilon_n<\infty$ of the uniqueness theorem, with $\Upsilon_n$ as in
Remark~\ref{9.10:rem-downstream-constants}.
\end{itemize}
For the $n$-th member of an admissible family, let $\hat\zeta^n_{(0)}$ be the running shift of the
$n$-th re-filed run, a function of the frozen complex source, evaluated as in
\eqref{9.10:eq-refiled-bilipschitz-2} below; let $\hat\zeta^n_*$ be the centre of the $n$-th run fixed in
Theorem~\ref{9.10:thm-exponential-rate}; and let
$\theta_0:=\theta_r|_{r=0}=X$, where $\theta_r=X+c_0r$ and $X=g^{-2}$. Put
\begin{equation}\label{9.10:eq-refiled-bilipschitz-2}
\hat x_{\mathrm{ref},n}(\hat w):=\theta_0-\hat\zeta^n_{(0)}\bigl(\hat\zeta^n_*+\hat w\bigr).
\end{equation}
Then:
\begin{itemize}
\item[(i)] the map $\hat w\mapsto\hat x_{\mathrm{ref},n}(\hat w)$ is bi-Lipschitz on $|\hat w|\le\varrho_2/2$,
uniformly in $n$: for all $\hat w,\hat w'$ in this disc,
\begin{equation}\label{9.10:eq-refiled-bilipschitz-3}
\begin{aligned}
\frac13\,|\hat w-\hat w'|
&\le\prod_{k<K}\bigl(1-\operatorname{Lip}_{\bar D(0,\varrho)}\hat b_{k+1}\bigr)\,|\hat w-\hat w'|\\
&\le\bigl|\hat x_{\mathrm{ref},n}(\hat w)-\hat x_{\mathrm{ref},n}(\hat w')\bigr|
\le\frac53\,|\hat w-\hat w'| ;
\end{aligned}
\end{equation}
\item[(ii)] $\hat x_{\mathrm{ref},n}$ converges uniformly on $|\hat w|\le\varrho_2/2$ as $n\to\infty$.
\end{itemize}
\end{corollary}

\begin{proof}
For (i), fix $n$ and write $l_k:=\operatorname{Lip}_{\bar D(0,\varrho)}\hat b_{k+1}$ for $k<K$. That the map
$\hat w\mapsto\hat x_{\mathrm{ref},n}(\hat w)$ on $|\hat w|\le\varrho_2/2$ admits the product
$\prod_{k<K}(1-l_k)$ of the per-step factors of the run as a lower Lipschitz constant, and $5/3$ as an
upper Lipschitz constant whenever \eqref{9.10:eq-refiled-bilipschitz-1} holds, is the composition step in
the proof of the bi-Lipschitz lemma for the run without re-filing, applied word for word to the re-filed
run \eqref{9.10:eq-refiled-run-b}, with $\hat b_{k+1}$ and $\hat\zeta^n_{(0)}$ in place of their
counterparts for the run without re-filing; it remains to bound the product from below. By
\eqref{9.10:eq-refiled-bilipschitz-1} each $l_k$ lies in $[0,2/3]$, in particular $0\le l_k\le1$. For numbers
$l_k\in[0,1]$ one has
\begin{equation}\label{9.10:eq-refiled-bilipschitz-4}
\prod_{k<K}(1-l_k)\ge1-\sum_{k<K}l_k ,
\end{equation}
by induction on the number of factors: if $P$ denotes the product over the first $j$ factors and
$P\ge1-\sum_{k<j}l_k$, then, since $0\le P\le1$ and $l_j\ge0$,
\begin{equation*}
P\,(1-l_j)=P-Pl_j\ge P-l_j\ge1-\sum_{k\le j}l_k .
\end{equation*}
Inserting \eqref{9.10:eq-refiled-bilipschitz-1} into \eqref{9.10:eq-refiled-bilipschitz-4} gives
$\prod_{k<K}(1-l_k)\ge1-\frac23=\frac13$, which is the first inequality of
\eqref{9.10:eq-refiled-bilipschitz-3}. Since \eqref{9.10:eq-refiled-bilipschitz-1} holds for
the run of every member of the family, with the same right side $2/3$, neither bound depends on $n$; this
is the uniformity in (i).

For (ii), the uniform convergence of $\hat x_{\mathrm{ref},n}$ as $n\to\infty$ follows from hypothesis (b),
the summability $\sum_n\Upsilon_n<\infty$, by the convergence step in the proof of the bi-Lipschitz lemma
for the run without re-filing, which applies to the re-filed run word for word, with $\Upsilon_n$ unchanged
(Remark~\ref{9.10:rem-downstream-constants}).
\end{proof}

Together with item (1) of clause (v-c) of Theorem~\ref{3.1:thm-main}, as established by the uniqueness
theorem, (i) and (ii) above give:
two admissible families whose reference-scale (renormalised) inverse couplings along the re-filed runs
converge to the same value have the same limit.

\begin{remark}[Norm of re-filed families on the tube]\label{9.10:rem-tube-norm}
Let $R^{(i)}$ be a remnant of the re-filed renormalisation run and $\mathsf{RF}[R^{(i)}]$ its re-filed
pointed family at scale $i''$, both as in Definition~\ref{9.10:def-refiled-run}. By Ba{\l}aban's
replacement of the configuration by its box minimiser, described in the proof, the norm bound of the
re-filing lemma, part (b), for $\mathsf{RF}[R^{(i)}]$ holds on the whole tube of configurations at
scale $i''$ on which the re-filing lemma is stated, with the notation and all constants unchanged.
This is the way followed in this book, and it requires no estimate beyond those of the correlation
theorem and of the re-filing lemma. A second way, which avoids that replacement, gives in place of
the bound of part (b) the bound
\begin{equation}\label{9.10:eq-tube-norm-1}
\|\mathsf{RF}[R^{(i)}]\|
\le \bigl[L^4(\log x_i)^{4r_{\mathrm{pr}}}e^{2\cdot 7^d\kappa}+2\bigr]\,C(\kappa)\,g_i^{\kappa_0},
\end{equation}
where $r_{\mathrm{pr}}$ is a fixed exponent, independent of $i$ and $k$, $C(\kappa)$ is a constant
depending on $\kappa$, $C'$ below is a constant independent of $k$, and $y_k$ is as in the ceiling
conditions among the hypotheses of Theorem~\ref{9.10:thm-refiled-induction}. Under the second way the
bound $\hat{\mathfrak{m}}_k$ on the dirty part $\hat Q^{(j)}$ of the hatted regular family of the
re-filed run (Definition~\ref{9.10:def-refiled-run}) is replaced by
\begin{equation}\label{9.10:eq-tube-norm-2}
10\bigl(1+C'L^4(\log y_k)^{4r_{\mathrm{pr}}}\bigr)\,y_k^{-\kappa_0/2},
\end{equation}
and again no estimate beyond those of the correlation theorem and of the re-filing lemma is required.
\end{remark}

\begin{proof}
Part (i) of the lemma on box minimisers used in the proof of the correlation theorem requires its
configuration to be a box minimiser. A generic point of the tube at scale $i''$ is not a box
minimiser, so that lemma does not apply there directly.

\emph{The first way.} Ba{\l}aban replaces the configuration $U_{k+1}$ by the configuration
$U_{k+1}(\square_0,M'(U))$, extended as equal to $U$ outside $\square_0$, and states that this
configuration has approximately the same regularity properties as $U$, and hence that all the
theorems needed there remain valid \cite{B-CMP109}. This device is used here as stated there, in the
same form in which the proof of the correlation theorem already uses it. Accordingly the argument of
$\mathsf{RF}[R^{(i)}]$ in
Definition~\ref{9.10:def-refiled-run} is the replaced configuration
$U_i(\square_0(X),M^{\cdot}_i(\mathfrak{b}))$ on the cube $\square_0(X)$ of $X$. At this argument
part (i) of the lemma on box minimisers applies. With it, the norm bound of the re-filing lemma,
part (b), holds at every point of the tube, with the notation and all constants unchanged.

\emph{The second way.} Take any share $\mathsf{sh}(X,y)$ with $y\in\mathbb{T}^{(i'')}$, the block
lattice at level $i''$ of Definition~\ref{1.5:def-block-average}, over which the shares of
Definition~\ref{9.10:def-refiled-run} are indexed. Bound it by the trivial bound on a share, and use,
for each share, the containment of its dependence in the
assignment domain $Y(X,y)\supseteq\square_0(X)$ of Definition~\ref{9.10:def-refiled-run} that the
containment statement and the radius lemma of the proof of the correlation theorem supply. Combining
these
bounds over the shares gives \eqref{9.10:eq-tube-norm-1}; there $R^{(i)}$ is of the marginal
class, with $\|R^{(i)}_\zeta\|^{\flat}\le g_i^{\kappa_0}$ by Definition~\ref{9.10:def-refiled-run}.

Under the second way, the proposition bounding the dirty part $\hat Q^{(j)}$ is applied with
\eqref{9.10:eq-tube-norm-1} in place of the bound of the re-filing lemma, part (b). Its induction
closes once the bound on the dirty part, namely
$\hat{\mathfrak{m}}_k = 10(1+C_{RF})y_k^{-\kappa_0/2}$, is replaced by the quantity
\eqref{9.10:eq-tube-norm-2}. That quantity is nondecreasing in $k$.

The inequality stating that the quantity \eqref{9.10:eq-tube-norm-2} is at most the un-hatted bound
$\mathfrak{m}_k$ on the dirty part is
imposed as one further ceiling condition on $g$, of the same kind as those among the hypotheses of
the induction theorem for the re-filed run (Theorem~\ref{9.10:thm-refiled-induction}).

Each bound of this section in which the norm constant $C_{RF}$ of the re-filing lemma occurs then
carries one further polylogarithmic factor. That factor is absorbed at the place where $C_{RF}$
stands, exactly as in the finite-range supremum bound. The reason is that $C_{RF}$ always occurs
there multiplied by a power $g^{\kappa_0}$ or by the factor $\hat{\mathfrak{r}}$.

In both ways the bound on the whole tube follows from estimates of the correlation theorem and of
the re-filing lemma, and no new estimate is introduced.
\end{proof}

\begin{remark}[One-sided conditions and order of coupling ceilings]\label{9.10:rem-one-sided-conditions}
The conditions on the constants used in this section are listed below. Each is one-sided, and none assigns a value to any constant.
\begin{enumerate}
\item[(1)] $L\ge L_0$ and $L$ odd, as in hypothesis (H2) of Notation~\ref{2.2:not-hypotheses}.
\item[(2)] $M$, $M_1$, $M_2$ (Notation~\ref{0.2:not-glossary}) and the radii $\check R_k$ are powers of $L$, with $M>M_2>M_1$, as in Ba{\l}aban's construction. Consequently $\mathsf{a}_i\ge2$ for every $i$.
\item[(3)] $\check R_0\ge2$, where $\check R_0$ is the partition-cube radius at level $0$.
\item[(4)] $\kappa_0>4$, where $\kappa_0$ is the exponent that enters the correlation theorem and the uniqueness theorem.
\item[(5)] $\kappa\ge10\log L$.
\item[(6)] The parameter $\rho$ is less than $1$; the conditions imposed on it are the lower bounds that place it in the window
\begin{equation*}
\rho\in\bigl[\max\{q^{1/2},\,L^{-\vartheta_{\sharp}},\,\vartheta_{dl},\,\vartheta_{\chi}\},\,1\bigr).
\end{equation*}
The window, its non-emptiness for every $L$, and the fact that its end points do not depend on $n$
are taken from the statement of this chapter that fixes the window for $\rho$; there
$\vartheta_{\sharp}$ is the exponent fixed together with the window. Every further condition
imposed on $\rho$ is a lower bound strictly below $1$, so the window is only narrowed from below
and $\rho<1$ is kept. The further lower bound $\rho\ge L^{-1/(2r_{pr})}$ is also imposed, where
$r_{pr}>0$ is the radius exponent: the exponent for which the product of $M$ with the
partition-cube radius at depth $s$ grows at most like a constant times $(\log\theta_s)^{r_{pr}}$,
with $\theta_s:=g^{-2}+c_0s$ (a property of the radii as they are defined in Ba{\l}aban's
construction, not a condition imposed here). It is fixed before $g$ (it is among the constants listed in
item~(7)); the bound $\rho\ge L^{-1/(2r_{pr})}$ at most raises the lower end of the window, which
remains non-empty for every $L$ since $L^{-1/(2r_{pr})}<1$, and no estimate uses it.
\item[(7)] The ceilings on $g$ are chosen last, after
\begin{equation*}
(L,\mathfrak{p},M,\kappa,\kappa_0,r_{pr},r_J,\rho)
\end{equation*}
and the remaining constants of the re-filed run. They are the following:
\begin{itemize}
\item a coupling ceiling which gives, in particular, $\ell_s\le\ell_{s'}+1$ whenever $s'<s\le s'+\ell_s$;
\item the two ceilings of the re-filing lemmas, whose norm constant is $C_{RF}$;
\item the two coupling conditions $10(1+C_{RF})\le C_{\mathfrak{m}}y_k$ and
$\tfrac{9}{16}E_0+\hat{\mathfrak{m}}_k\le E_0$ among the hypotheses of
Theorem~\ref{9.10:thm-refiled-induction}, with $C_{\mathfrak{m}}$, $y_k$, $E_0$ and
$\hat{\mathfrak{m}}_k$ as in that theorem;
\item the coupling ceiling, in hatted form, taken at the dial constant $\hat C_L'$;
\item the conditions (a)--(d), the hatted forms of (e) and (f), and condition (g$'$) taken with $2\varkappa'$ in place of $\varkappa'$, all collected in~\eqref{9.10:eq-envelope-ceilings-b}, together with condition (b$''$) there and the two instances of it displayed where it is applied in this chapter;
\item the inequality~\eqref{9.10:eq-source-bounds-a} of Proposition~\ref{9.10:prop-source-bounds}, stated with proof incomplete at the step marked $(\ast)$;
\item the condition $\varrho_2\ge2$ of Lemma~\ref{9.10:lem-second-difference}, where $\varrho_2$ is the room constant of the correlation theorem.
\end{itemize}
\end{enumerate}
None of these conditions bounds from above any of $L$, $K$, $k$, $n$, the parameters $m_0$, $\mathsf{N}$ and $N_0$ of the re-filed run, the torus exponent $m$, the parameter $N_w$ attached to the complex source $w$, an age, a depth, the radii $\check R_k$ or the number of large-field regions carried by the run. None bounds the coupling from below. None bounds the volume: the volume enters these conditions only through constants. The constant $L_0=L_0(N)$ of (H2) is the one existential constant. It is not modified here and is never evaluated.
\end{remark}

\begin{proof}
Item~(2). Since $M_1\ge1$ is a power of $L$ and $M_2>M_1$ is a power of $L$, we have $M_2\ge LM_1\ge L$. In the same way $M>M_2$ gives $M\ge LM_2\ge L^2$. The radius $\check R_i$ is a power of $L$ with non-negative exponent, so $\check R_i\ge1$. Hence
\begin{equation*}
\mathsf{a}_i=\log_L(M\check R_i)\ge\log_L L^2=2 .
\end{equation*}

Item~(3) follows from the lower bound $\check R_0\ge(\log g_0^{-2})^{r}$ on the partition-cube
radius in Ba{\l}aban's construction, where $r\ge2$ is the exponent of that lower bound, at $g$
below a ceiling. Indeed, the bare coupling satisfies $g_0\le g$ in the first-passage setting of
Definition~\ref{1.2:def-flow-constants}, so
\begin{equation*}
(\log g_0^{-2})^{r}\ge(\log g^{-2})^{r}\ge2
\end{equation*}
as soon as $\log g^{-2}\ge2^{1/r}$, a ceiling on $g$ depending on $r$ only.

Item~(4) is implied by the condition $\kappa_0\ge7$ of the correlation theorem and by the condition $\kappa_0\ge10$ of the uniqueness theorem. Indeed each of these bounds is stronger than $\kappa_0>4$.

Item~(5) is one of the parameter conditions of the correlation theorem.

One-sidedness. Items~(1), (2), (4) and (5) are floors on constants fixed before $g$, together
with, in item~(2), an ordering and the requirement of being powers of $L$. Item~(3) holds for
every $g$ below a ceiling, by the lower bound on $\check R_0$ used in its proof. Item~(6) consists
of lower bounds on $\rho$, each strictly below $1$. Every entry of
item~(7) is a ceiling on $g$, that is, a condition satisfied by every $g$ below a positive
threshold. That threshold depends only on the constants chosen before $g$.
For~\eqref{9.10:eq-source-bounds-a} and the conditions collected
in~\eqref{9.10:eq-envelope-ceilings-b} this is shown where they are stated
(Proposition~\ref{9.10:prop-source-bounds}, stated with proof incomplete at the step marked
$(\ast)$, and Definition~\ref{9.10:def-envelope-ceilings}); the remaining entries are introduced
as ceilings on $g$ in the statements where they first appear, with thresholds depending on
constants fixed before $g$.

Order of choice. The entries of item~(7) are finitely many, and every constant their thresholds depend on is fixed before $g$. Therefore all of them hold at once when $g$ lies below the least of these thresholds.
In particular no threshold depends on $k$ or $K$, which are not among these constants, so the two
conditions of Theorem~\ref{9.10:thm-refiled-induction} carrying the index $k$ hold for every
$k\le K$ below one threshold in $g$.

No entry bounds from above any of $K$, $k$, $n$, $m_0$, $m$, $N_w$, an age, a depth, $\mathsf{N}$,
$N_0$ or the number of large-field regions carried by the run. The constant $L$ is chosen first;
every later condition in which $L$ appears, namely items~(2) and~(5), the lower end of the window
of item~(6) and the thresholds of the entries of item~(7), is a condition on a constant chosen
after $L$ or on $g$, and not a bound on $L$. The radii $\check R_k$ appear in items~(2) and~(3)
and, through $\ell_s$, in~\eqref{9.10:eq-source-bounds-a} and in condition~(b$''$), which are
ceilings on $g$; no condition bounds them from above. No entry involves $g$ in the direction of a
lower bound.
\end{proof}

\begin{definition}[Runs, the common ladder and reading units]\label{9.11:def-runs-ladder}
Let $\mathsf{A}$, of length $n$, and $\mathsf{B}$, of length $n+1$, be the two runs of the
uniqueness theorem, between which the comparison chain of
Definition~\ref{9.6:def-comparison-chain} is built; we fix the frame in which they are compared.

\emph{Depths, levels and the ladder.} A run of length $n$ is indexed by depths $s=0,1,\dots,n$.
Depth $s$ of $\mathsf{A}$ corresponds to the level $j=n-s$ of $\mathsf{A}$.
The run $\mathsf{B}$ has one further level, finer than all levels of
$\mathsf{A}$, which carries the index $-1$ in the numbering of $\mathsf{A}$. At each depth $s$
there is one unit lattice common to both runs. The \emph{ladder} is the sequence
\begin{equation}\label{9.11:eq-runs-ladder-1}
  \theta_0:=X,\qquad \theta_{s+1}:=\theta_s+c_0\quad(s\ge 0),
\end{equation}
with $X=g^{-2}$; it is common to both runs. The \emph{couplings} of a run of length $n$ are
defined from the value $\zeta$ assigned to the run at depth $n$ by
\begin{equation}\label{9.11:eq-runs-ladder-2}
  \zeta_{(n)}:=\zeta,\qquad \zeta_{(r)}:=\zeta_{(r+1)}+b_{(r)}-c_0,\qquad
  x_{(r)}:=\theta_r-\zeta_{(r)},
\end{equation}
where $b_{(r)}$ is the marginal coefficient at depth $r$ defined below. The running coupling of
the run at depth $s$ is written $g_{(s)}$; it satisfies $g_{(s)}^2\le 1$. The small parameter of
the correlation theorem, written $\theta(g)$ there, is written here $\theta_J(g)$, the subscript
keeping it apart from the ladder $\theta_s$ of \eqref{9.11:eq-runs-ladder-1}; $\theta_J(g)\to 0$
as $g\to 0$.

Depth $r$ corresponds to the output level $k+1$, and a slot at depth $s>r$ corresponds to a level
$j\le k$, with
\begin{equation}\label{9.11:eq-runs-ladder-3}
  s-r=k+1-j;
\end{equation}
indeed, with $k+1=n-r$ and $j=n-s$ one has $k+1-j=(n-r)-(n-s)=s-r$. Accordingly, the production
at depth $r$ (the step whose output is the list at depth $r$) reads slots at depths $s>r$, that
is, at older and finer levels.

\emph{Lists.} A \emph{list} carries at each depth $s$ the slots
\begin{equation*}
  \bigl(\zeta_{(s)};\ \mathsf{P};\ P^{\flat},\,Q,\,R;\ E(\cdot;W),\,R(\cdot;W)\bigr),
\end{equation*}
with the structure of the lists of the correlation theorem: $E:=P^{\flat}+Q$ is a pointed
$E$-family, and
\begin{equation}\label{9.11:eq-runs-ladder-4}
  b_{(s)}:=\beta[E]
\end{equation}
is its marginal coefficient, $\beta[\cdot]$ being the marginal-coefficient functional of the
correlation theorem applied to the $E$-family; $R$ denotes the marginal (large-field) families;
the local members
of the list read one continuum source $W$, and the frozen members are holomorphic in $\hat{w}$.

\emph{Blocks and components.} The blocks are $\bullet\in\{P,A,R\}$, each with components $c$,
given as follows:
\begin{equation}\label{9.11:eq-runs-ladder-5}
\begin{aligned}
  P&=\bigl(\mathsf{P}/E_0 \text{ at $g$-free radii};\ C_t\theta_s Q\bigr),\\
  A&=\mathsf{l}_s^{-1}\bigl(P^{\flat};\,D\bigr)/E_0,\\
  R&=C_t\theta_s\,(R'+R''),
\end{aligned}
\end{equation}
where
\begin{equation}\label{9.11:eq-runs-ladder-6}
  \mathsf{l}_s:=(\log\theta_s)^{c_1},\qquad c_1:=p_0+q_{\flat}+1,
\end{equation}
$p_0$ being the fixed exponent of the degree function $p_0(\cdot)$, and where
\begin{equation}\label{9.11:eq-runs-ladder-7}
  D:=E(\cdot;W)-E_{W(p_y)}
\end{equation}
is the \emph{freezing defect}, $W(p_y)$ being the value of the source $W$ at the point $p_y$;
$D$ vanishes at constant sources. In \eqref{9.11:eq-runs-ladder-5}, $R'$ and $R''$ are the two
marginal families of a depth, namely the families denoted $R'$ and $R''$ in
\cite[(1.91)]{B-CMP122b}, as they are carried by the lists of the correlation theorem; the block
$R$ is denoted by the same letter as the slot $R$, whose two families $R'$, $R''$ it combines
with the factor $C_t\theta_s$. For a
block $\bullet$, $\|\cdot^{\bullet}\|$ denotes the maximum of the norms over the components of
the block.

\emph{Reading units.} The blocks are read in the reading units of the corollary of this chapter
that bounds the blocks $P$, $A$, $R$ at each depth, one unit for each component:
\begin{equation}\label{9.11:eq-runs-ladder-8}
\begin{aligned}
  \mathfrak{u}_s^{\mathsf{P}}&:=1, &
  \mathfrak{u}_s^{Q}&:=C_t\theta_s\,\mathfrak{m}_{(s)}, &
  \mathfrak{u}_s^{P^{\flat}}&:=\mathsf{l}_s^{-1},\\
  \mathfrak{u}_s^{D}&:=2\,\mathsf{l}_s^{-1}, &
  \mathfrak{u}_s^{R}&:=\mathfrak{r}_s:=C_t\,\theta_s^{1-\kappa_0/2}, &&
\end{aligned}
\end{equation}
where $\mathfrak{m}_{(s)}$ is the bound on the slot $Q$ at depth $s$.
\end{definition}

\begin{definition}[Competitors, coupling distances, transplant and hybrids]
\label{9.11:def-competitors-distances}
We use the runs, depths, levels, the common ladder, the run couplings $\zeta_{(s)}$ and the
reading units $\mathfrak{u}_s^c$ of Definition~\ref{9.11:def-runs-ladder}.

\emph{The pair and its parameters.} The run $\mathsf{A}$ has length $n$ and the run
$\mathsf{B}$ has length $n+1$. The lists of the pair are parametrised over a set
$\mathfrak{P}$, with closure $\overline{\mathfrak{P}}$: for a frozen pair, $\mathfrak{P}$ is
the disc $D(0,r_J)$; for a local pair, $\mathfrak{P}$ is a parameter set about a real point,
which we denote $u$.

\emph{Agreement of totals.} For a family $G$ and a function $\varphi$ of the data on the unit
lattices common to the two runs, we write
\begin{equation*}
G\approx\varphi
\end{equation*}
if the total of $G$ on backgrounds equals $\varphi$, component by component and in the pointed
sense. For $\mathsf{m}\in\{\mathsf{A},\mathsf{B}\}$, a depth $s$ and a component $c$, we write
$\varphi_s^{\mathsf{m},c}$ for the total of the depth-$s$ component $c$ of the run
$\mathsf{m}$, and we put $\varphi_n^{\mathsf{A},c}:=0$.

\emph{Competitors and their distances.} A \emph{competitor} at depth $s$ and component $c$ is
a family $G_s^c$ of the run $\mathsf{A}$ with
\begin{equation*}
G_s^c\approx\varphi_s^{\mathsf{B},c}-\varphi_s^{\mathsf{A},c},
\end{equation*}
and the \emph{competitor distance} is
\begin{equation*}
\mathsf{D}_s^c:=\inf\bigl\{\|G\|^c:\ G\approx\varphi_s^{\mathsf{B},c}
-\varphi_s^{\mathsf{A},c}\bigr\},
\end{equation*}
where $\|\cdot\|^c$ is the norm of the component $c$.

\emph{Coupling distances.} The \emph{coupling distances} are
\begin{equation*}
d_s:=\sup_{\overline{\mathfrak{P}}}\bigl|\zeta^{\mathsf{B}}_{(s)}
-\zeta^{\mathsf{A}}_{(s)}\bigr|\quad (s\le n),\qquad d_s:=0\quad (s>n),
\end{equation*}
and the weighted two-level distances are
\begin{equation}\label{9.11:eq-competitors-distances-a}
d^{\sharp}_{r+1}:=\sum_{i\ge0}L^{-i/2}\bigl(d_{r+1+i}+d_{r+2+i}\bigr)\ \ge\ d_{r+1}.
\end{equation}

\emph{The transplant.} For a family $F$ of the run $\mathsf{B}$ at level $j$, let
$\mathfrak{T}F$ be its cross-cutoff transplant. This is a family of the run $\mathsf{A}$; the
map $\mathfrak{T}$ is linear, the parameters of the pair pass through it unchanged, and
\begin{equation}\label{9.11:eq-competitors-distances-b}
\begin{aligned}
&(\mathfrak{T}1)\quad \|\mathfrak{T}F\|^c\le C_V\|F\|^{\flat};\\
&(\mathfrak{T}2)\quad \text{the totals of $\mathfrak{T}F$ and of $F$ agree}\\
&\phantom{(\mathfrak{T}2)}\quad \text{on whole-torus regular backgrounds};\\
&(\mathfrak{T}3)\quad \beta^{\mathsf{A}}[\mathfrak{T}F]=\beta^{\mathsf{B}}[F];\\
&(\mathfrak{T}4)\quad \text{the first output of $\mathsf{B}$ sits in an extra slot}\\
&\phantom{(\mathfrak{T}4)}\quad \text{at level $0$ of $\mathsf{A}$.}
\end{aligned}
\end{equation}
Here $\|\cdot\|^{\flat}$ is the family norm, $\beta^{\mathsf{m}}[\,\cdot\,]$ is the
coefficient functional of the run $\mathsf{m}$, and $C_V$ is the norm constant of the
transplant.

\emph{The order of sub-steps.} Inside depth $r$ the production is a finite sequence of
sub-steps. Each sub-step reads, besides the slots at depths $s>r$, only depth-$r$ blocks
produced before it, in an order $\prec$ on the set of blocks $\{P,A,R\}$. We write $\bullet$
for a generic block and $\circ$ for a block with $\circ\prec\bullet$; for a block $\bullet$ the
symbol $G_s^{\bullet}$ stands for the collection of the $G_s^c$ with $c$ a component of
$\bullet$, and $G_s$ for the collection of all $G_s^c$ at depth $s$.

\emph{Hybrids.} Let $\underline{r}$ be a depth and let $(G_s^c)_{\underline{r}\le s\le n}$ be
competitors with
\begin{equation*}
\|G_s^c\|^c\le (C_V+1)\,\mathfrak{u}_s^c .
\end{equation*}
Write $\mathfrak{s}^{\mathsf{A}}_s$ for the tuple of slots of the list of $\mathsf{A}$ at
depth $s$, and
$\mathfrak{e}_r^{\mathsf{A},\circ}$ for the block $\circ$ produced by $\mathsf{A}$ at depth
$r$. The \emph{hybrid} $\mathfrak{s}^{\natural}_{r,\bullet}$ is the list of $\mathsf{A}$ in
which
\begin{itemize}
\item the slots at each depth $s>r$ are $\mathfrak{s}^{\mathsf{A}}_s+G_s$;
\item at depth $r$, each block $\circ\prec\bullet$ is
$\mathfrak{e}_r^{\mathsf{A},\circ}+G_r^{\circ}$;
\item the coupling slots are $\zeta^{\mathsf{B}}_{(s)}$;
\item the slot at depth $n$ is $G_n$ alone.
\end{itemize}
By the closure clause of the frame of the uniqueness theorem, the hybrid lies in the enlarged
class $C_{\mathfrak{K}}\mathfrak{K}^c$, where $\mathfrak{K}^c$ is the class of the frame at the
component $c$ and $C_{\mathfrak{K}}$ is its class constant:
\begin{equation}\label{9.11:eq-competitors-distances-c}
\mathfrak{s}^{\natural}_{r,\bullet}\in C_{\mathfrak{K}}\mathfrak{K}^c .
\end{equation}

\emph{Production differences and source families.} Let
$\mathsf{F}_r^{\mathsf{A},+,\bullet}$ be the production of the block $\bullet$ at depth $r$
by $\mathsf{A}$, with the extra slot of $(\mathfrak{T}4)$ in
\eqref{9.11:eq-competitors-distances-b}, and let $F_r^{\mathsf{A},\bullet}$ be the output
block $\bullet$ of $\mathsf{A}$ at depth $r$. The \emph{production difference} is
\begin{equation*}
M_r^{\bullet}:=\mathsf{F}_r^{\mathsf{A},+,\bullet}\bigl(\mathfrak{s}^{\natural}_{r,\bullet}
\bigr)-F_r^{\mathsf{A},\bullet}.
\end{equation*}
The \emph{source family} $N_r^{\bullet}$ is the family supplied by the source-family
hypothesis of the uniqueness theorem; its components satisfy
\begin{equation*}
N_r^c\approx\varphi_r^{\mathsf{B},c}-\varphi^c\bigl[\mathsf{F}_r^{\mathsf{A},+,\bullet}
\bigl(\mathfrak{s}^{\natural}_{r,\bullet}\bigr)\bigr],
\end{equation*}
where $\varphi^c[\,\cdot\,]$ denotes the total of the component $c$.
\end{definition}

\begin{proof}
The transplant properties \eqref{9.11:eq-competitors-distances-b} are part of the data of the
cross-cutoff transplant fixed above, and the closure clause
\eqref{9.11:eq-competitors-distances-c} is that of the frame of the uniqueness theorem; the one
assertion of Definition~\ref{9.11:def-competitors-distances} proved here is the inequality in
\eqref{9.11:eq-competitors-distances-a}; we prove it for every $r$ with $r+1\ge0$.

Every coupling distance is non-negative: for $s\le n$, $d_s$ is a supremum of absolute values,
and for $s>n$, $d_s=0$. Hence every term $L^{-i/2}(d_{r+1+i}+d_{r+2+i})$ of the series defining
$d^{\sharp}_{r+1}$ is non-negative, since $L^{-i/2}>0$.

The series has only finitely many non-zero terms. Indeed, if $i\ge n-r$, then
$r+1+i\ge n+1>n$ and $r+2+i>n$, so $d_{r+1+i}=d_{r+2+i}=0$ by the definition of the coupling
distances at depths beyond $n$. Thus $d^{\sharp}_{r+1}$ is a finite sum of non-negative terms.

A sum of non-negative terms is at least any one of its terms. The term with $i=0$ is
$L^{0}(d_{r+1}+d_{r+2})=d_{r+1}+d_{r+2}$, and $d_{r+2}\ge0$. Therefore
\begin{equation*}
d^{\sharp}_{r+1}\ \ge\ d_{r+1}+d_{r+2}\ \ge\ d_{r+1},
\end{equation*}
which is the inequality in \eqref{9.11:eq-competitors-distances-a}.
\end{proof}

\begin{definition}[Hatted lists: re-filing of healed remnants]\label{9.11:def-hatted-lists}
Let $\mathsf{A}$ and $\mathsf{B}$ be the two runs of a pair, of lengths $n$ and $n+1$. Depths $s$,
levels $j=n-s$ and the common ladder $\theta_s$ are as in Definition~\ref{9.11:def-runs-ladder}.
Families, their totals on backgrounds and the relation $\approx$ are as in
Definition~\ref{9.11:def-competitors-distances}. The hatted lists defined below change the slot
into which a healed large-field remnant is entered. They do not change the measure.

\emph{The re-filing delay.} For each depth $s$ let $\mathsf{a}_s\ge 1$ be the integer
$\log_L(M\check{R}_{(s)})$ read at depth $s$ from the scaffold, as a datum of the ladder $\theta$;
it is common to both runs of the pair. Written in levels, $\mathsf{a}$ is the increment of the
re-filing level
\begin{equation*}
  i'':=i+\log_L(M\check{R}_{(n-i)}),
\end{equation*}
that is, $\mathsf{a}_{n-j}=j''-j$ for a level $j$; the subscripts of $\mathsf{a}$ and of
$\check{R}$ are always depths. For the level $j=n-s$ of depth $s$ this reads
\begin{equation*}
  j''=n-s+\mathsf{a}_s,
\end{equation*}
which is the level of depth $s-\mathsf{a}_s$. Let $\rho$ be the contraction rate of the glossary
(Notation~\ref{0.2:not-glossary}). For the integer
$\mathsf{a}_s$ defined here,
\begin{equation*}
  \rho^{-\mathsf{a}_s}=(L^{\mathsf{a}_s})^{\log_L\rho^{-1}}
  =(M\check{R}_{(s)})^{\log_L\rho^{-1}}.
\end{equation*}
If one takes instead $\mathsf{a}_s:=\lceil\log_L(M\check{R}_{(s)})\rceil\vee 1$, this identity
becomes the inequality
\begin{equation*}
  \rho^{-\mathsf{a}_s}\le\rho^{-1}(M\check{R}_{(s)})^{\log_L\rho^{-1}}.
\end{equation*}
The definition used throughout this chapter is the first one, with the identity.

\emph{The re-filing map.} The remnant slot $R_t$ produced at depth $t$ has, $\mathsf{a}_t$ steps
later, a partition cube consisting of a single cell. The re-filing map
$\operatorname{\mathsf{RF}}$ of the re-filing map lemma re-expresses $R_t$ at depth
$t-\mathsf{a}_t$.
The result is a pointed, fully covariant $E$-family built on the run's own radii. The map
$\operatorname{\mathsf{RF}}$ is linear, the total of $\operatorname{\mathsf{RF}}[R_t]$ equals
$(R_t)^{\mathrm{vac}}$, the vacuum part of $R_t$, and
\begin{equation*}
  \|\operatorname{\mathsf{RF}}[R_t]\|^{\flat}\le C_{\mathrm{RF}}\,\omega\,\|R_t\|^{\flat}.
\end{equation*}
Here $\omega$ is the weight in the bound of the re-filing map lemma; it is indexed by the level at
which the share is entered and by the level of the remnant.

\emph{Presence as a unit.} Let $j=n-t$ be the level of the slot $R_t$, and let
$j'':=j+\mathsf{a}_{n-j}=n-t+\mathsf{a}_t$. This is the level of depth $t-\mathsf{a}_t$. Whether
$R_t$ is present as a unit is decided zone by zone. Here a \emph{use} of a slot is a read of it by
a later production, and the zone-assignment rule assigns each use to a \emph{zone}, which carries
a level. The uses that the zone-assignment rule places in zones of level at
most $j''$ read $R_t$ as a remnant unit that has not yet been re-filed. Write $\nu$ for the level
of a zone. By the presence lemma, in its form for small and large pairs, the remnant units
assigned to a zone of level $\nu$ are those of levels $j\ge\nu-\mathsf{a}_{n-j}$, equivalently
$j''\ge\nu$. By the presence clause, a statement distinct from the presence lemma,
a remnant unit of level $j$ is read at level $i$ in one of two cases:
\begin{itemize}
  \item it is young, $j\le i<j''$; or
  \item it is old and not yet re-filed, $i\ge j''$; in this case it lies in a zone of level
    $\nu\le j''$ inside the enlarged component $X^{\sim}$ at level $i$.
\end{itemize}
The share $\operatorname{\mathsf{RF}}[R_t]$ is entered whole into $\hat{E}_{t-\mathsf{a}_t}$.
A share, once entered, never returns.

\emph{Hatted lists.} A \emph{hatted list} of length $n$ carries at depth $s$ the slots displayed
in the first line of~\eqref{9.11:eq-hatted-lists-a}. The un-hatted slots $\mathsf{P}_s$,
$P^{\flat}_s$, $Q_s$, $R_s$, $R_s(\cdot;W)$, the marginal-coefficient functional $\beta[\cdot]$
and the frozen evaluation are the data of the run's un-hatted list at depth $s$, as in
Definitions~\ref{9.11:def-runs-ladder} and~\ref{9.11:def-competitors-distances}; the hatted
entries are defined by the remaining lines of~\eqref{9.11:eq-hatted-lists-a}:
\begin{equation}\label{9.11:eq-hatted-lists-a}
\begin{gathered}
  \bigl(\hat{\zeta}_{(s)};\ \mathsf{P}_s;\ P^{\flat}_s,\ \hat{Q}_s,\ R_s;\
  \hat{E}_s(\cdot;W),\ R_s(\cdot;W)\bigr),\\
  \hat{Q}_s:=Q_s+\sum_{t:\ t-\mathsf{a}_t=s}\operatorname{\mathsf{RF}}_s[R_t]_{\mathrm{frozen}},
  \qquad \hat{E}_s:=P^{\flat}_s+\hat{Q}_s,\\
  \operatorname{\mathsf{RF}}_s[R_t]:=\operatorname{\mathsf{RF}}[R_t]\ \text{read at}\
  (i'',i)=(n-s,\,n-t),\qquad s=t-\mathsf{a}_t,\\
  \hat{b}_{(s)}:=\beta[\mathsf{P}_s]+\beta[\hat{Q}_s],\qquad
  \hat{\zeta}_{(s)}:=\hat{\zeta}_{(s+1)}+\hat{b}_{(s)}-c_0,\qquad
  \hat{x}_{(s)}:=\theta_s-\hat{\zeta}_{(s)}.
\end{gathered}
\end{equation}
For the run $\mathsf{B}$ the value $t=n+1$ is admitted in the sum defining $\hat{Q}_s$. The slot
$R_t$ itself is unchanged. Which later uses still read $R_t$ as a unit is settled by the presence
lemma, as described above. Through $\beta[\hat{Q}_s]$, the hatted marginal coefficient
$\hat{b}_{(s)}$ contains the marginal coefficients of the re-filed remnants
$\operatorname{\mathsf{RF}}_s[R_t]_{\mathrm{frozen}}$, $t-\mathsf{a}_t=s$; their contribution is
the exact share of the remnants in $\hat{b}_{(s)}$. The hatted couplings $\hat{\zeta}_{(s)}$ and
$\hat{x}_{(s)}$ are defined on the same ladder $\theta_s$ as the couplings of the run.

\emph{Hatted production.} The \emph{hatted production} $\hat{\mathsf{F}}_r$ at depth $r$ has
output level $k+1=n-r$. It is the production of the inductive step of the correlation theorem,
fed with the hatted list, with its uses assigned to zones by the zone-assignment rule. In detail,
it consists of the following three parts.
\begin{itemize}
  \item[($\alpha$)] The regular productions are the hatted step production, the hatted
    one-lattice Lipschitz step acting on $\hat{Q}$, and the hatted freezing defect $\hat{D}$.
    They read as remnant units not yet re-filed the slots $R_s$ with
    \begin{equation*}
      s\in\mathfrak{B}_r:=\{s>r:\ s-r=\mathsf{a}_s\}.
    \end{equation*}
    The $R''$-production marks the band $r<s\le r+\mathsf{a}_s-1$. The shares already re-filed
    are read through $\hat{E}_{s-\mathsf{a}_s}$.
  \item[($\beta$)] The large-field productions are the following: the $R'$-production, through
    its outer read at level $i=k$ and its inner reads through $\mathbf{T}'_{k+1}(X)$ and the
    inner operations $\hat{\mathbf{T}}_h,\dots,\hat{\mathbf{T}}_k$; the presence indicator of the
    $\mathbf{T}$-production; and the dial presence indicator. They read $R_s$ as a remnant unit
    wherever it is present in the sense of the presence lemma, whatever the value of $s-r$.
    Each such read carries the factor supplied by geometric forgetting of old regions. This
    gives the hatted history kernel $\hat{\mathfrak{w}}^R$ of the remnant slot, which has range
    $\mathsf{a}$ and is geometric in the age beyond that range. It also gives the window kernel
    $\mathsf{k}_{r,s}$ of the hatted remnant bound.
  \item[($\gamma$)] Parts ($\alpha$) and ($\beta$) are followed by entering, at depth $r$, the
    shares $\operatorname{\mathsf{RF}}_r[R_s]$, $s\in\mathfrak{B}_r$, into $\hat{E}_r$.
\end{itemize}

\emph{Hatted runs.} A \emph{hatted run} is a hatted list each of whose slots is produced from the
older slots by $\hat{\mathsf{F}}$, as in the re-filing lemmas. That a hatted run represents the
same measure, with the constants of the correlation theorem, is the conjunction of two
statements: the proposition on the hatted frozen slot $\hat{Q}$, and the hatted-run statement
(the domain clauses of the correlation theorem's induction with the hatted slots in place of the
un-hatted ones). Neither statement is used in this definition.
\end{definition}

\begin{definition}[Hatted blocks and their reading units]\label{9.11:def-hatted-blocks}
Let $\mathsf{A}$ and $\mathsf{B}$ be the two runs of a pair, of lengths $n$ and $n+1$, with the
pair, the totals, the competitors and the distances of
Definition~\ref{9.11:def-competitors-distances}. Let $\mathsf{a}_s$, the re-filing map
$\operatorname{\mathsf{RF}}$ and the hatted slots $\hat{Q}$, $\hat{E}$, $\hat{\zeta}$, $\hat{b}$ be those of
Definition~\ref{9.11:def-hatted-lists} and \eqref{9.11:eq-hatted-lists-a}. Here $\mathsf{P}$,
$P^{\flat}$, $Q$, $R$, $E(\cdot;W)$ are the slots of a list at depth $s$, and $E_0$ is the bound of
the $E$-family. $D$ is the un-hatted freezing defect. $\hat{E}_{W(p_y)}$ denotes $\hat{E}$ with the
source frozen at its value $W(p_y)$ at the point $p_y$. The ladder $\theta_s$, the constant $C_t$,
the logarithmic unit $\mathsf{l}_s$, the reading units $\mathfrak{m}_{(s)}$ and $\mathfrak{r}_s$ of the
un-hatted blocks, the splitting $R'+R''$ of the remnant slot and the coefficient functional
$\beta_r$ are those of the un-hatted comparison to which
Definition~\ref{9.11:def-competitors-distances} belongs.

\emph{Blocks.} At depth $s$ the \emph{hatted blocks} are measured in the same reading units as the
un-hatted ones. They are
\begin{equation}\label{9.11:eq-hatted-blocks-a}
\begin{gathered}
\hat{P} := \bigl(\mathsf{P}/E_0;\ C_t\theta_s\hat{Q}\bigr),\qquad
\hat{A} := \mathsf{l}_s^{-1}\bigl(P^{\flat};\ \hat{D}\bigr)/E_0,\qquad
\hat{R} := C_t\theta_s\,(R'+R''),\\[2pt]
\begin{aligned}
\hat{D} &:= \hat{E}(\cdot;W)-\hat{E}_{W(p_y)}\\
&= D+\sum_{s':\ s'-\mathsf{a}_{s'}=s}
\Bigl(\operatorname{\mathsf{RF}}_s[R_{s'}](\cdot;W)-\operatorname{\mathsf{RF}}_s[R_{s'}]_{W(p_y)}\Bigr).
\end{aligned}
\end{gathered}
\end{equation}
The second expression for $\hat{D}$ writes the freezing defect of $\hat{E}$ as follows. It is the
un-hatted defect $D$ plus the freezing defects of the remnants $R_{s'}$ that are re-filed at depth
$s=s'-\mathsf{a}_{s'}$ through $\operatorname{\mathsf{RF}}_s$.

The block $\hat{P}$ has two components. The re-filed remnant, in its frozen form, is carried by its
second component $C_t\theta_s\hat{Q}$. The block $\hat{A}$ has the components
$\mathsf{l}_s^{-1}P^{\flat}/E_0$ and $\mathsf{l}_s^{-1}\hat{D}/E_0$. The freezing defect of the
re-filed remnant is carried by the second of these, through the sum in
\eqref{9.11:eq-hatted-blocks-a}.

\emph{The block $\hat{R}$.} The block $\hat{R}$ is the un-hatted $R$-block
$C_t\theta_s(R'+R'')$, unchanged as a slot. What changes is where the remnant $R_s$ of depth $s$ is
read as a unit of $\hat{R}$ by a production at depth $r<s$. This is decided region by region, by
the reading rule of the hatted comparison, as follows.
\begin{itemize}
\item The regular productions at depth $r$ read $R_s$ as a unit only in two cases: for
$s\in\mathfrak{B}_r$, or for $s$ in the band $r<s\le r+\mathsf{a}_s-1$. For every other $s$ they read
$R_s$ only through its share filed in $\hat{E}_{s-\mathsf{a}_s}$ by
$\operatorname{\mathsf{RF}}_{s-\mathsf{a}_s}$ (Definition~\ref{9.11:def-hatted-lists}).
\item The inner operations of the large-field productions read $R_s$ as a unit wherever it is
present. Inside the parts of a large-field component that are old at the level read, they also do
so for $s-r>\mathsf{a}_s+1$. There they carry the factor supplied by clause~(c) of the lemma on
geometric forgetting of old regions.
\end{itemize}
This is why the window kernel $\mathsf{k}_{r,s}$ of \eqref{9.11:eq-hatted-blocks-1}, which weights
the remnant $R_s$ in the $\hat{R}$-row of the bounds at depth $r$, equals $1$ for
$s-r\le\mathsf{a}_s+1$. Beyond that range it is geometric in
$s-r$, and not zero.

No boundary term $\mathbf{B}^{(j)}$ of the large-field operation is a component of $\hat{P}$,
$\hat{A}$ or $\hat{R}$. The boundary terms, including the kernel terms
$\mathfrak{w}^{\mathbf{B}}_{k,j}\delta^{\mathbf{B}}_j$ and the norms
$\|\delta\mathbf{B}'^{(k+1)}\|$, are not components of these blocks; how they enter the hatted
comparison is set out where the boundary terms of the large-field operation are treated.

\emph{Reading units.} Let $\hat{\mathfrak{m}}_{(s)}$ be the bound on $\|\hat{Q}\|^{\flat}$ that
enters the hatted class. It is supplied by the proposition bounding the hatted slot $\hat{Q}$. The
hatted defect obeys the bound in the first line below, and the reading units
$\hat{\mathfrak{u}}_s^c$ of the components $c$ are those of the second line:
\begin{equation}\label{9.11:eq-hatted-blocks-b}
\begin{gathered}
\|\hat{D}\|^{\flat}\le 2E_0,\\
\hat{\mathfrak{u}}_s^{\mathsf{P}}:=1,\quad
\hat{\mathfrak{u}}_s^{Q}:=C_t\theta_s\hat{\mathfrak{m}}_{(s)},\quad
\hat{\mathfrak{u}}_s^{P^{\flat}}:=\mathsf{l}_s^{-1},\quad
\hat{\mathfrak{u}}_s^{D}:=2\,\mathsf{l}_s^{-1},\quad
\hat{\mathfrak{u}}_s^{R}:=\mathfrak{r}_s .
\end{gathered}
\end{equation}
The bound in the first line is a consequence of the bound $\|\hat{E}\|^{\flat}\le E_0$ of the
hatted class. That bound holds both for the local family $\hat{E}(\cdot;W)$ and for its frozen
value $\hat{E}_{W(p_y)}$. By the definition of $\hat{D}$ in \eqref{9.11:eq-hatted-blocks-a} and the
triangle inequality for $\|\cdot\|^{\flat}$, we obtain
\[
\|\hat{D}\|^{\flat}\le\|\hat{E}(\cdot;W)\|^{\flat}+\|\hat{E}_{W(p_y)}\|^{\flat}\le 2E_0 .
\]
This argument is the same as the one giving the bound $\|D\|^{\flat}\le 2E_0$ in the un-hatted
comparison. It justifies the unit $2\,\mathsf{l}_s^{-1}$ of the $\hat{D}$-component.

\emph{The hatted class.} The hatted class $\hat{\mathfrak{K}}$ is the class $\mathfrak{K}$ of the
un-hatted comparison with $\hat{Q}$, $\hat{E}$, $\hat{\mathfrak{m}}$ in place of $Q$, $E$,
$\mathfrak{m}$. Its closure clause is the clause \eqref{9.11:eq-competitors-distances-c} read for
$\hat{\mathfrak{K}}$, that is, with $\hat{Q}$, $\hat{E}$, $\hat{\mathfrak{m}}$ in place of $Q$, $E$,
$\mathfrak{m}$. The coefficient functional
$\beta_r$ reads both components of $\hat{P}$.

\emph{Pair, competitors, distances, hybrids.} The following objects are those of
Definition~\ref{9.11:def-competitors-distances}, verbatim with hats:
\begin{itemize}
\item the pair and the totals;
\item the competitors $\hat{G}_s^c$ and the competitor distances $\hat{\mathsf{D}}_s^c$;
\item the distances $d_s$ and $d^{\sharp}$;
\item the hybrids $\hat{\mathfrak{s}}^{\natural}_{r,\bullet}$, formed with $\hat{Q}$ and $\hat{D}$ in
place of $Q$ and $D$;
\item the extended production $\hat{\mathsf{F}}_r^{\mathsf{A},+,\bullet}$ and the output
$\hat{F}_r^{\mathsf{A},\bullet}$ at depth $r$.
\end{itemize}
The coupling distance is now taken for $\hat{\zeta}$:
\[
d_s:=\sup_{\overline{\mathfrak{P}}}\bigl|\hat{\zeta}^{\mathsf{B}}_{(s)}-\hat{\zeta}^{\mathsf{A}}_{(s)}\bigr|
\quad (s\le n),\qquad d_s:=0\quad (s>n).
\]
Writing $\varphi^{c}[G]$ for the total of the component $c$ of a family $G$, as in
Definition~\ref{9.11:def-competitors-distances}, the production differences and the source
families are
\begin{align*}
\hat{M}_r^{\bullet}&:=\hat{\mathsf{F}}_r^{\mathsf{A},+,\bullet}\bigl(\hat{\mathfrak{s}}^{\natural}_{r,\bullet}\bigr)
-\hat{F}_r^{\mathsf{A},\bullet},\\
\hat{N}_r^{c}&\approx\varphi_r^{\mathsf{B},c}
-\varphi^{c}\bigl[\hat{\mathsf{F}}_r^{\mathsf{A},+,\bullet}\bigl(\hat{\mathfrak{s}}^{\natural}_{r,\bullet}\bigr)\bigr].
\end{align*}
The transplant $\mathfrak{T}$ acts on hatted families with the properties
$(\mathfrak{T}1)$--$(\mathfrak{T}3)$ of \eqref{9.11:eq-competitors-distances-b}. The re-filing map
$\operatorname{\mathsf{RF}}$ and $\mathfrak{T}$ are linear, and both preserve totals:
$\operatorname{\mathsf{RF}}$ by the lemma on the re-filing map $\operatorname{\mathsf{RF}}$, which
states that $\operatorname{\mathsf{RF}}$ preserves totals, and $\mathfrak{T}$ by $(\mathfrak{T}2)$ of
\eqref{9.11:eq-competitors-distances-b}. Therefore the
identity
\[
\beta^{\mathsf{A}}\bigl[\mathfrak{T}\hat{Q}^{\mathsf{B}}\bigr]=\beta^{\mathsf{B}}\bigl[\hat{Q}^{\mathsf{B}}\bigr]
\]
is $(\mathfrak{T}3)$ applied to $\hat{Q}^{\mathsf{B}}$.

\emph{Letters.} The letters $q$, $q_a$, $q_{\flat}$, $\rho$, $\varkappa'_r$, $d^{\sharp}$, $C_\sigma$,
$C_V$, the constant $\mathsf{A}$ of the coefficient bound $|\beta_r[F]|\le\mathsf{A}\,\|F\|$ (a
constant, not to be confused with the run $\mathsf{A}$) and the index $n_c$ of the alignment
statement of the un-hatted comparison keep their meaning in the hatted comparison. So do
the weight and the window length
\[
\Lambda_r:=\Lambda(\log\theta_r)^{c_3},\qquad
\ell_r:=2+\log_L\bigl(M\check{R}_{(r)}\bigr).
\]
The letters introduced for the hatted comparison are $\mathsf{a}_s$ and the band coefficient
$C_F:=3C_{\mathrm{RF}}+1$. They also
include a free constant $C_F^A$, subject only to the one-sided condition $C_F^A\ge C_F+1$. They
further include the constant $\hat{c}_\rho$ and the hatted $R$-rate $\hat{\mathfrak{r}}_r$, both
fixed with the rates of the hatted comparison; the latter enters the sum $\mathfrak{S}$ below.
Finally we put
\begin{equation}\label{9.11:eq-hatted-blocks-1}
\begin{gathered}
\mathfrak{S}:=\sum_{s\ge0}\hat{\mathfrak{r}}_s\,\rho^{-\mathsf{a}_s},\qquad
\mathsf{k}_{r,s}:=\rho^{2(s-r-\mathsf{a}_s-1)_+},\\
C_{**}:=(1-q_a)^{-1}\sup_{i\ge0}(i+2)\,q_a^{i/2},\qquad
\hat{\Lambda}_r:=C_{**}\Lambda_r,\qquad
\mathfrak{B}_r:=\{s>r:\ s-r=\mathsf{a}_s\}.
\end{gathered}
\end{equation}

\emph{Levels and depths.} In terms of levels, the depths correspond through $k+1=n-r$ and $j=n-s$.
We put $j'':=j+\mathsf{a}_s$, where $s=n-j$ is the depth of the level $j$. Then
\[
k-j''=(n-r-1)-(n-s)-\mathsf{a}_s=s-r-\mathsf{a}_s-1 .
\]
Hence $\mathfrak{B}_r$ is the depth image of $\mathfrak{B}(k):=\{j\le k:\ j''=k+1\}$. The band
$\operatorname{band}(k):=\{j\le k:\ j''\ge k+2\}$ corresponds to the $R''$-band
$\{r<s\le r+\mathsf{a}_s-1\}$: the condition $j\le k$ reads $s\ge r+1$, and $j''\ge k+2$ reads
$s-r\le\mathsf{a}_s-1$. We write $\Sigma^{\mathrm{band}}_{s>r}$ for the sum over $s>r$ with
$s-r\le\mathsf{a}_s$. Its range is the union of the $R''$-band and $\mathfrak{B}_r$.

Three further windows are used:
\begin{itemize}
\item the \emph{unit window} $\{r<s\le r+\mathsf{a}_s+1\}$, on which
$(s-r-\mathsf{a}_s-1)_+=0$, that is, $\mathsf{k}_{r,s}=1$;
\item the window $\{r<s<r+\ell_r\}$ of the coefficients $\varkappa'_r$;
\item the window of $\Sigma^{\mathrm{band}}_{s>r}$ described above.
\end{itemize}
\end{definition}

\begin{definition}[One-sided floors on the kernel and contraction rates]\label{9.11:def-rate-floors}
Let $L$ be the blocking factor, and let $n$ be the length of the shorter run of the pair of runs
compared in Definition~\ref{9.6:def-comparison-chain}. Let $\hat{\alpha}$ and $\beta_1$ be the
exponents of the floor on the kernel rate used for the un-hatted blocks, which is taken over
unchanged in \eqref{9.11:eq-rate-floors-a}, and let $c_a^{\mathrm{age}}$ be the age rate constant. The
\emph{kernel rate} $q$, the rate of the kernel in the corollary for the un-hatted blocks, and
the \emph{age rate} $q_a$ are numbers satisfying
\begin{equation}\label{9.11:eq-rate-floors-a}
\max\{L^{-\hat{\alpha}},\,L^{-\beta_1},\,L^{-1/2}\}<q<1,
\qquad q_a:=e^{-c_a^{\mathrm{age}}}\le q^{2}.
\end{equation}
The \emph{contraction rate} $\rho$ is a number satisfying
\begin{equation}\label{9.11:eq-rate-floors-b}
\hat{\rho}_0\le\rho<1,
\qquad
\hat{\rho}_0:=\max\{q^{1/2},\,L^{-\vartheta_{\sharp}},\,\vartheta_{\mathrm{dl}},\,\vartheta_{\chi},\,
\rho_{g'}\}<1 .
\end{equation}
The first four members of $\hat{\rho}_0$ are the floors on the contraction rate used for the
un-hatted blocks, taken over unchanged. The member $\rho_{g'}$ is the lower bound on $\rho$
assumed in the verification of the last of the ceilings on the coupling; it is a number $<1$
and is stated in two forms: the lower bound $\rho\ge L^{-(q_{\flat}-p_0)/(2r_{\mathrm{pr}})}$,
under the condition $q_{\flat}>p_0$, where $q_{\flat}$ is one of the exponents among the given
parameters, and the lower bound
$\rho\ge L^{-1/(2r_{\mathrm{pr}})}$. Here $r_{\mathrm{pr}}$ is the constant of the waiting-time
bound, and $p_0$ is the fixed exponent of the degree function $p_0(\cdot)$.

In particular,
\begin{equation}\label{9.11:eq-rate-floors-1}
\begin{gathered}
q\le\rho^{2},\qquad L^{-1/2}<q\le\rho^{2}\le\rho,\\
q_a^{1/2}\le q\le\rho^{2},\qquad q_a\le\rho^{4}<\rho .
\end{gathered}
\end{equation}

The floor \eqref{9.11:eq-rate-floors-b} is a lower bound only. Raising $\rho$ within
$\hat{\rho}_0\le\rho<1$ weakens every conclusion drawn from these rates. It also leaves every
condition imposed on the rates in the bounds for the hatted blocks below a condition of the same
kind. Nothing bounds $\rho$ away
from $1$ except $\rho<1$ itself. Nothing bounds $L$, $n$, $r_0$, the depth or the volume from above.
No floor on the kernel rate $q$ other than \eqref{9.11:eq-rate-floors-a} is stated.

The antecedents of the bound for the hatted block $\hat{R}$ carry auxiliary constants of their
own, each subject only to one-sided conditions. These constants are not evaluated here, and none
of them is a hypothesis of the theorem on the hatted envelope. They are:
\begin{itemize}
\item the constant $\vartheta_*$, subject to $\vartheta_*\le q_a\le q^{2}$ and $\vartheta_*\le 1/4$;
\item the constants of the input of that bound which furnishes $\vartheta_*\le 1/4$; that input
carries three further one-sided conditions on its own constants, which are stated with that input.
\end{itemize}

\emph{Rate constants.} Let $\theta_r$ be the common ladder of the two runs
($\theta_0=X=g^{-2}$, $\theta_{s+1}=\theta_s+c_0$). The following constants are given and kept
symbolic:
\begin{itemize}
\item the constant $C_t$ and the exponent $\kappa_0$ entering $\mathfrak{r}_r$ below, and the
constant $\Lambda$ of the logarithmic weight $\Lambda_r$ below;
\item the norm constant $C_V$ of the cross-cutoff transplant;
\item the source constant $C_\sigma$;
\item the constant $\mathsf{A}$ of the coefficient bound.
\end{itemize}
Let $q$, $q_a$, $\rho$ satisfy \eqref{9.11:eq-rate-floors-a} and \eqref{9.11:eq-rate-floors-b}.
For every depth $r$ we put
\begin{equation}\label{9.11:eq-rate-floors-c}
\begin{aligned}
&e_r:=\rho^{\,n-r},\qquad
\hat{c}_\rho:=\frac{1+\rho}{\rho(1-\rho)^{2}},\qquad
C_E\ge\max\{C_V+1,\,2C_\sigma\},\\
&\hat{\mathfrak{r}}_r:=2\,\mathfrak{r}_r\,\hat{\Lambda}_r
\bigl[\mathsf{A}(1-\rho)^{-2}+6(1-\rho)^{-1}+8\bigr],
\qquad \mathfrak{r}_r:=C_t\,\theta_r^{\,1-\kappa_0/2},\\
&\hat{\Lambda}_r:=C_{**}\Lambda_r,\qquad
\Lambda_r:=\Lambda(\log\theta_r)^{c_3},\qquad c_3:=3p_0+2,\\
&C_{**}:=(1-q_a)^{-1}\sup_{m\ge0}\,(m+2)\,q_a^{m/2}.
\end{aligned}
\end{equation}
The envelope constant $C_E$ is any constant satisfying the displayed one-sided condition.
\end{definition}

\begin{proof}
We derive \eqref{9.11:eq-rate-floors-1} from \eqref{9.11:eq-rate-floors-a} and
\eqref{9.11:eq-rate-floors-b}. We then show that the constants in
\eqref{9.11:eq-rate-floors-c} are well defined.

First, $q>0$ by \eqref{9.11:eq-rate-floors-a}. Also $\rho\ge\hat{\rho}_0\ge q^{1/2}>0$, because
$q^{1/2}$ is a member of the maximum defining $\hat{\rho}_0$. Squaring the inequality
$\rho\ge q^{1/2}$ between positive numbers gives $q\le\rho^{2}$.

The lower bound $L^{-1/2}<q$ holds because $L^{-1/2}$ is a member of the maximum in
\eqref{9.11:eq-rate-floors-a}. Since $0<\rho<1$, we have $\rho^{2}\le\rho$. Together with
$q\le\rho^{2}$ this gives the chain $L^{-1/2}<q\le\rho^{2}\le\rho$.

Next, $q_a=e^{-c_a^{\mathrm{age}}}>0$. From $q_a\le q^{2}$ and $q>0$, taking square roots gives
$q_a^{1/2}\le q$, and hence $q_a^{1/2}\le q\le\rho^{2}$. Squaring $q\le\rho^{2}$ and using
$q_a\le q^{2}$ gives
\begin{equation*}
q_a\le q^{2}\le\rho^{4}.
\end{equation*}
Finally, $0<\rho<1$ gives $\rho^{3}<1$, so $\rho^{4}<\rho$. This proves
\eqref{9.11:eq-rate-floors-1}.

It remains to check \eqref{9.11:eq-rate-floors-c}. Since $0<\rho<1$, the denominators $\rho$ and
$1-\rho$ are non-zero, so $\hat{c}_\rho$ and the bracket in $\hat{\mathfrak{r}}_r$ are finite.
Since $c_0>0$, we have $\theta_r\ge\theta_0=X=g^{-2}$, and $\log g^{-2}>0$ because the reference
coupling satisfies $g<1$; hence $\log\theta_r>0$, and, the constants $C_t$ and $\Lambda$ being
positive, $\mathfrak{r}_r$ and $\Lambda_r$ are positive and finite. By
\eqref{9.11:eq-rate-floors-1} we have $0<q_a\le\rho^{4}<1$. Hence $(1-q_a)^{-1}$ is finite, and
$0<q_a^{1/2}<1$. The sequence $(m+2)q_a^{m/2}$ therefore tends to $0$ as $m\to\infty$, so its
supremum over $m\ge0$ is finite. Thus $C_{**}<\infty$, and $\hat{\Lambda}_r$ and
$\hat{\mathfrak{r}}_r$ are finite.
\end{proof}

\begin{remark}[The three memory structures of the hatted hypotheses]\label{9.11:rem-memory-structures}
For each depth $s$ let $\mathsf{a}_s$ denote the integer $\log_L(M\check{R}_{(s)})$ fixed earlier in
this chapter, where $\check{R}_{(s)}$ is the waiting time of the large-field operation at depth $s$
and $M$ is Ba{\l}aban's parameter. For depths $s>r$ we say that
\emph{$s$ is in the band of $r$} if $0<s-r\le\mathsf{a}_s$. We put
\begin{equation}\label{9.11:eq-memory-structures-a}
\begin{gathered}
\Sigma^{\mathrm{band}}_{s>r}:=\sum_{s>r:\ s-r\le\mathsf{a}_s},\qquad
\mathfrak{B}_r:=\{s>r:\ s-r=\mathsf{a}_s\},\\
\{s:\ 0<s-r\le\mathsf{a}_s+1\}\quad\text{(the unit window of $r$)}.
\end{gathered}
\end{equation}
The unit window of $r$ contains the band of $r$, and $\mathfrak{B}_r$ lies in the band of $r$.
The hatted hypotheses for the blocks $\hat{P}$, $\hat{A}$ and $\hat{R}$ of
Definition~\ref{9.11:def-hatted-blocks} carry the differences at older depths $s>r$ in three ways.
\begin{itemize}
\item[(i)] With the geometric kernel $q^{s-r}$; this is the form of the $P$-, $A$- and
$\zeta$-members.
\item[(ii)] In the hypotheses for $\hat{P}$ and $\hat{A}$, by a term of finite range summed with
$\Sigma^{\mathrm{band}}_{s>r}$ and carrying a coefficient of order one.
For $\hat{P}$ the contributions producing this term are carried
over $\mathfrak{B}_r$, which lies in the band; for $\hat{A}$ they are carried over the units
present at the step.
\item[(iii)] In the hypothesis for $\hat{R}$, with the window kernel
\begin{equation}\label{9.11:eq-memory-structures-1}
\mathsf{k}_{r,s}:=\rho^{2(s-r-\mathsf{a}_s-1)_{+}},
\end{equation}
where $(u)_{+}:=\max\{u,0\}$ for a real number $u$. This kernel equals $1$ on the unit window and
is geometric beyond it; it is the image in depths of the hatted remnant kernel
$\hat{\mathfrak{w}}^R$ of the hatted form of the one-lattice Lipschitz step. The kernel
$\hat{\mathfrak{w}}^R$ has range given by the integers $\mathsf{a}_s$ and is geometric in the age
beyond that range.
\end{itemize}
In the hypotheses for $\hat{P}$ and $\hat{A}$ both terms, (i) and (ii), are written. This covers
either form of these hypotheses, since a right member from which one of the two terms is absent
only shrinks.
\end{remark}

\begin{proof}
Let $s\in\mathfrak{B}_r$. Then $s>r$ and $s-r=\mathsf{a}_s$, so $0<s-r\le\mathsf{a}_s$ and $s$ is in
the band of $r$; thus $\mathfrak{B}_r$ lies in the band.

Let $s$ be in the band of $r$. Then $0<s-r\le\mathsf{a}_s<\mathsf{a}_s+1$, so $s$ lies in the unit
window of $r$; thus the unit window contains the band.

Let $s$ lie in the unit window, so $s\le r+\mathsf{a}_s+1$. Then $s-r-\mathsf{a}_s-1\le0$, hence
$(s-r-\mathsf{a}_s-1)_{+}=0$ and \eqref{9.11:eq-memory-structures-1} gives
$\mathsf{k}_{r,s}=\rho^{0}=1$.

Let now $s-r>\mathsf{a}_s+1$. Then the positive part in \eqref{9.11:eq-memory-structures-1} is the
excess $e:=s-r-\mathsf{a}_s-1\ge1$ itself, and $\mathsf{k}_{r,s}=\rho^{2e}=(\rho^{2})^{e}$. This is
a geometric sequence in the excess of $s-r$ over $\mathsf{a}_s+1$, which is the second part of the
description of $\mathsf{k}_{r,s}$ in~(iii).
\end{proof}

\begin{definition}[The hatted response rows with window kernel]\label{9.11:def-response-rows}
Let $q$ and $\rho$ satisfy the one-sided floors of Definition~\ref{9.11:def-rate-floors}. Consider a pair of
hatted lists as in \eqref{9.11:eq-hatted-lists-a}, formed from the run $\mathsf{A}$ of length $n$ and the run
$\mathsf{B}$ of length $n+1$. We use the following objects:
\begin{itemize}
\item $\hat{G}_s^c$ is the hatted competitor difference at depth $s$ in the component $c$, and
$\|\hat{G}_s^{\bullet}\|$ is the maximum of the norms of its components in the block $\bullet$;
\item $\hat{P}$, $\hat{A}$, $\hat{R}$ are the hatted blocks of Definition~\ref{9.11:def-hatted-blocks};
\item $\hat{M}_r^{\bullet}$ is the production difference at depth $r$ in the block $\bullet$, evaluated at a
hatted hybrid list $\hat{\mathfrak{s}}^{\natural}_{r,\bullet}$ as in the definition of the competitors,
coupling distances, transplant and hybrids of this section;
\item $d_{s}$ are the coupling distances, $d^{\sharp}_{r+1}$ is the weighted two-level distance of
\eqref{9.11:eq-competitors-distances-a}, and $g_{(s)}$ is the running coupling at depth $s$;
\item $\mathsf{a}_s$ is the integer $\log_L(M\check{R}_{(s)})$ read from the scaffold, where $\check{R}_{(s)}$ is
the waiting time of the large-field operation;
\item $\Sigma^{\mathrm{band}}_{s>r}$ is the sum over $s>r$ with $s-r\le\mathsf{a}_s$, and $\mathfrak{B}_r$ is
the set of the band's oldest slots, both as in \eqref{9.11:eq-memory-structures-a} and
Remark~\ref{9.11:rem-memory-structures}.
\end{itemize}

\emph{Hypotheses.} We assume the following.
\begin{itemize}
\item[(I1)] The three response inequalities for the production differences in the blocks $\hat{P}$,
$\hat{A}$, $\hat{R}$ and the identity for the hatted marginal coefficients of the hatted two-run corollary
hold for the pair under consideration, for every $r<n$, every block and every hatted hybrid list.
\item[(I2)] The coefficient bound of the uniqueness theorem holds for the coefficient functionals
$\beta^{\mathsf{A}}_s$ of the run $\mathsf{A}$: for every depth $s$ and every family $F$,
\begin{equation*}
|\beta^{\mathsf{A}}_s[F]|\le\mathsf{A}\,\|F\|.
\end{equation*}
Here the constant $\mathsf{A}$ is written with the same letter as the run $\mathsf{A}$.
\item[(I3)] The coefficient $\varkappa_A$ described under \emph{Constants} below satisfies
$\varkappa_A\le 1/36$. Since $\varkappa_A$ tends to $0$ with $g$, this bound holds under the smallness
condition on $g$ among the standing conditions of this chapter.
\end{itemize}

The pair is said to satisfy the \emph{hatted response rows with window kernel} if the following hold for
every $r<n$, every block $\bullet\in\{\hat{P},\hat{A},\hat{R}\}$ (written $P$, $A$, $R$ in superscripts) and
every hatted hybrid list $\hat{\mathfrak{s}}^{\natural}_{r,\bullet}$.

\smallskip\noindent(a) The three production differences obey the following rows. The row $(\hat{P})$:
\begin{equation}\label{9.11:eq-response-rows-a}
\begin{split}
\|\hat{M}_r^P\|\le{}&\alpha_r d^{\sharp}_{r+1}
+\varkappa\sum_{s>r}q^{s-r}\bigl(\|\hat{G}_s^P\|+\|\hat{G}_s^R\|\bigr)\\
&+\varkappa''_r\sum_{s>r}q^{s-r}\|\hat{G}_s^A\|
+C_F\,\Sigma^{\mathrm{band}}_{s>r}\|\hat{G}_s^R\|,
\end{split}
\end{equation}
the row $(\hat{A})$:
\begin{equation}\label{9.11:eq-response-rows-b}
\begin{split}
\|\hat{M}_r^A\|\le{}&\alpha_r^A d^{\sharp}_{r+1}
+\varkappa_A\Bigl[\|\hat{G}_r^P\|+\sum_{s>r}q^{s-r}\bigl(\|\hat{G}_s^P\|+\|\hat{G}_s^A\|
+\|\hat{G}_s^R\|\bigr)\Bigr]\\
&+\|\hat{G}_r^R\|+C_F^A\,\Sigma^{\mathrm{band}}_{s>r}\|\hat{G}_s^R\|,
\end{split}
\end{equation}
and the row $(\hat{R})$, with the function $\hat{B}_r$ and the window kernel $\mathsf{k}_{r,s}$:
\begin{equation}\label{9.11:eq-response-rows-c}
\begin{split}
\|\hat{M}_r^R\|&\le\mathfrak{r}_r\hat{\Lambda}_r\hat{B}_r(\hat{G},d)
+\varkappa'_r\sum_{r<s<r+\ell_r}\|\hat{G}_s^R\|,\\
\hat{B}_r(\hat{G},d)&:=\sum_{r\le s<n}q^{s-r}g^2_{(s+1)}d_{s+1}
+\sum_{s\ge r}q^{s-r}\bigl(\|\hat{G}_s^P\|+\|\hat{G}_s^A\|\bigr)
+\sum_{s>r}\mathsf{k}_{r,s}\|\hat{G}_s^R\|,\\
\mathsf{k}_{r,s}&:=\rho^{2(s-r-\mathsf{a}_s-1)_+}.
\end{split}
\end{equation}
The window kernel $\mathsf{k}_{r,s}$ equals $1$ on the unit window $s\le r+\mathsf{a}_s+1$. Beyond the unit
window it decreases geometrically.

\smallskip\noindent(b) For every competitor and every depth $s$, the hatted marginal coefficients satisfy
\begin{equation}\label{9.11:eq-response-rows-d}
\begin{split}
\hat{b}^{\mathsf{B}}_{(s)}-\hat{b}^{\mathsf{A}}_{(s)}
&=\beta^{\mathsf{A}}_s\bigl[\hat{G}_s^{\mathsf{P}}\bigr]+\beta^{\mathsf{A}}_s\bigl[\hat{G}_s^{\hat{Q}}\bigr],\\
\bigl|\hat{b}^{\mathsf{B}}_{(s)}-\hat{b}^{\mathsf{A}}_{(s)}\bigr|
&\le\mathsf{A}\,\|\hat{G}_s^P\|.
\end{split}
\end{equation}

\smallskip\noindent\emph{Constants.} All constants in (a) are free of $K$, $n$, $r_0$ and of the volume.
\begin{itemize}
\item For a constant $C$ free of $K$, $n$, $r_0$ and of the volume, the coefficient $\alpha_r$ satisfies
\begin{equation}\label{9.11:eq-response-rows-1}
\alpha_r\le C\Bigl(g_{(r)}^3\bigl(\log\hat{x}_{(r)}\bigr)^{p_0}+\frac{g_{(r)}^2}{\log\hat{x}_{(r)}}\Bigr),
\end{equation}
and $\alpha_r^A$ satisfies a bound of the same form.
\item The coefficients $\varkappa$, $\varkappa''_r$ and $\varkappa_A$ are proportional to
\begin{equation*}
\sup_s\bigl(a_s+K_{CE}C_L\theta_J(g_{(s)})\bigr),
\end{equation*}
and they tend to $0$ with $g$. Here $a_s$ is the depth-$s$ coefficient and $C_L$ the constant with which the
sizes of the un-hatted response rows are stated; both are fixed with those rows, and the coefficient $a_s$ is
not the integer $\mathsf{a}_s$.
\item The window coefficient is $\varkappa'_r:=6K_{CE}C_L\theta_J(g_{(r)})$. The window length is
$\ell_r=2+\log_L(M\check{R}_{(r)})$.
\item All of $\alpha_r,\alpha_r^A,\varkappa,\varkappa_A,\varkappa''_r,\varkappa'_r$ are the functions of
depth of the un-hatted response rows. The bound \eqref{9.11:eq-response-rows-1}, the bound of the same form
for $\alpha_r^A$, the proportionality of $\varkappa$, $\varkappa''_r$, $\varkappa_A$, their convergence to
$0$ with $g$, and the formula for $\varkappa'_r$ are established with those rows.
\item The band coefficients are $C_F:=3C_{\mathrm{RF}}+1$, with $C_{\mathrm{RF}}$ the constant of the
re-filing map $\operatorname{\mathsf{RF}}$, and $C_F^A\ge C_F+1$.
\item The logarithmic weights are $\Lambda_r:=\Lambda(\log\theta_r)^{c_3}$ with $c_3:=3p_0+2$, and
$\hat{\Lambda}_r:=C_{**}\Lambda_r$. Here $C_{**}$ is the absorption constant of the hatted remnant kernel; it
is fixed in the order of choice before $C_\sigma$, $C_E$ and $\gamma$. Every bound in which $\Lambda$ enters
is monotone in $\Lambda$, so the replacement of $\Lambda_r$ by $\hat{\Lambda}_r$ is admissible there.
\end{itemize}

Under the hypotheses (I1), (I2) and (I3), every pair of hatted lists satisfies the hatted response rows with
window kernel.
\end{definition}

\begin{proof}
\emph{The row $(\hat{R})$.} The inequality \eqref{9.11:eq-response-rows-c} is the response inequality for the
block $\hat{R}$ of the hatted two-run corollary in hypothesis (I1), with the same function
$\hat{B}_r(\hat{G},d)$.

Compare this row with the un-hatted row $(R)$. The third member $\sum_{s>r}\|G_s^R\|$ of $B_r$ is replaced
by $\sum_{s>r}\mathsf{k}_{r,s}\|\hat{G}_s^R\|$, and $\Lambda_r$ is replaced by $\hat{\Lambda}_r$. The
$\varkappa'$-window, the sum over $r<s<r+\ell_r$, is unchanged.

The description of $\mathsf{k}_{r,s}$ follows from its definition. If $s\le r+\mathsf{a}_s+1$, then
$s-r-\mathsf{a}_s-1\le 0$, so the positive part vanishes and $\mathsf{k}_{r,s}=\rho^0=1$. If
$s>r+\mathsf{a}_s+1$, then $\mathsf{k}_{r,s}=\rho^{2(s-r-\mathsf{a}_s-1)}$. Since $\rho<1$ by
Definition~\ref{9.11:def-rate-floors}, this is geometric in the excess $s-r-\mathsf{a}_s-1$.

\emph{The rows $(\hat{P})$ and $(\hat{A})$.} The inequalities \eqref{9.11:eq-response-rows-a} and
\eqref{9.11:eq-response-rows-b} follow from the response inequalities for the blocks $\hat{P}$ and $\hat{A}$
of the hatted two-run corollary in hypothesis (I1). The corollary's inequalities do not contain the member
$\varkappa\sum_{s>r}q^{s-r}\|\hat{G}_s^R\|$ of
\eqref{9.11:eq-response-rows-a} nor the member $\varkappa_A\sum_{s>r}q^{s-r}\|\hat{G}_s^R\|$ of
\eqref{9.11:eq-response-rows-b}. These members are non-negative, so adding them only enlarges the right
members. In its inequality for the block $\hat{A}$ the corollary carries two band members; they are bounded by
the single band member of \eqref{9.11:eq-response-rows-b}, as shown below.

Compare them with the un-hatted rows $(P)$ and $(A)$.
\begin{itemize}
\item $(\hat{P})$ is $(P)$ with $\varkappa\sum_{s>r}\|G_s^R\|$ replaced by
\begin{equation*}
\varkappa\sum_{s>r}q^{s-r}\|\hat{G}_s^R\|+C_F\,\Sigma^{\mathrm{band}}_{s>r}\|\hat{G}_s^R\|.
\end{equation*}
\item $(\hat{A})$ is $(A)$ with $\varkappa_A\sum_{s>r}\|G_s^R\|$ replaced by
\begin{equation*}
\varkappa_A\sum_{s>r}q^{s-r}\|\hat{G}_s^R\|+C_F^A\,\Sigma^{\mathrm{band}}_{s>r}\|\hat{G}_s^R\|.
\end{equation*}
The same-depth member $\|\hat{G}_r^R\|$ is kept.
\end{itemize}

The band member $C_F\Sigma^{\mathrm{band}}$ in $(\hat{P})$ is carried by the component $\hat{Q}$ of the block
$\hat{P}$. The band member $C_F^A\Sigma^{\mathrm{band}}$ in $(\hat{A})$ is carried by the component $\hat{D}$
of the block $\hat{A}$. Both members are the re-filing at depth $r$ of the band's oldest slots
$t\in\mathfrak{B}_r$. The re-filing map is linear, so the difference between the two runs of the re-filed
remnant is the re-filing of the competitor:
\begin{equation}\label{9.11:eq-response-rows-2}
\operatorname{\mathsf{RF}}_r\bigl[R^{\mathsf{A}}_t+\hat{G}_t^R\bigr]
-\operatorname{\mathsf{RF}}_r\bigl[R^{\mathsf{A}}_t\bigr]
=\operatorname{\mathsf{RF}}_r\bigl[\hat{G}_t^R\bigr].
\end{equation}
Every $t\in\mathfrak{B}_r$ satisfies $t>r$ and $t-r=\mathsf{a}_t$, so $\mathfrak{B}_r$ is contained in the
band $\{t>r:\ t-r\le\mathsf{a}_t\}$. The corollary's band members are sums over $t\in\mathfrak{B}_r$ of
multiples of $\|\hat{G}_t^R\|$; since these summands are non-negative, each such member is bounded by the same
multiple of $\Sigma^{\mathrm{band}}_{s>r}\|\hat{G}_s^R\|$.

In the inequality for the block $\hat{A}$ of the hatted two-run corollary, the band enters through two
members, one with
coefficient $\varkappa_A$ and one with coefficient $C_F$. By hypothesis (I3),
$\varkappa_A\le 1/36$, so the sum of the two coefficients is at most $C_F+1$. Hence both
members are covered by the single band coefficient $C_F^A\ge C_F+1$ of \eqref{9.11:eq-response-rows-b}.

\emph{Clause (b).} For every competitor, the difference of the hatted marginal coefficients of the two runs
is exactly $\beta^{\mathsf{A}}_s$ evaluated on the $\mathsf{P}$-component and on the $\hat{Q}$-component of
$\hat{G}_s$. This is the first line of \eqref{9.11:eq-response-rows-d}, which is the identity for the hatted
marginal coefficients of the hatted two-run corollary in hypothesis (I1). Compared with the corresponding
identity of the un-hatted response rows, the only change is that $Q$ is replaced by $\hat{Q}$.

Both components belong to the block $\hat{P}$, and the functional $\beta^{\mathsf{A}}_s$ reads their total.
The coefficient bound of hypothesis (I2), applied at depth $s$, therefore gives the second line of
\eqref{9.11:eq-response-rows-d}.
\end{proof}

\begin{remark}[The window kernel from the hatted remnant kernel]\label{9.11:rem-window-kernel}
Fix a run of length $n$. The level $k+1$ produced by the step from level $k$ corresponds to the depth
$r$ with $k+1=n-r$, and a level $j\le k$ to the depth $s$ with $j=n-s$. We put
$j'':=j+\mathsf{a}_{n-j}$, where $\mathsf{a}_{n-j}$ is the integer $\mathsf{a}_s$ read at the depth
$s=n-j$. For a
remnant $R^{(j)}$ created at level $j$, $\bar{w}_{i,j}$ is its unit weight at level $i$, and
$\bar{w}_j:=\max_i\bar{w}_{i,j}$. For a large-field region $X$, $j(X)$ is the level at which it was
created and $X^{\sim}$ is its enlarged component. The constant $C_{**}$ is
\begin{equation}\label{9.11:eq-window-kernel-1}
C_{**}:=(1-q_a)^{-1}\sup_{m\ge0}(m+2)\,q_a^{m/2}.
\end{equation}
Assume the following.
\begin{itemize}
\item[(a$_0$)] The un-hatted response row holds. It is obtained from the band and remnant rows of
the one-lattice Lipschitz step written in reading units. Its band member is
$\varkappa'_r\sum_{r<s<r+\ell_r}\|G^R_s\|$; its remnant member is the flat sum
$\sum_{s>r}\|G^R_s\|$, which is the remnant row's history kernel $\mathfrak{w}^R$ written in reading
units; and its prefactor carries the weight $\Lambda_r$.
\item[(a)] The remnant row of the hatted one-lattice Lipschitz step holds with the hatted history
kernel
\begin{equation}\label{9.11:eq-window-kernel-2}
\hat{\mathfrak{w}}^R_{k,j}:=\sum_{\substack{i\in[j,k+1]:\ R^{(j)}\text{ is read as a remnant slot}\\
\text{by the inner operation of }X\text{ at level }i}}\bar{w}_{i,j}\cdot
(\text{factor of the read at level }i),
\end{equation}
where the inner operation of $X$ at level $i$ is the large-field operation acting inside $X$ at that
level, and each summand $\bar{w}_{i,j}$ is multiplied by the factor that part (c) of geometric
forgetting of old regions assigns to the read of $R^{(j)}$ by $X$ at level $i$. Its kernels
$\mathfrak{w}^P$ and $\mathfrak{w}^D$ of the $P$- and $D$-slots are those of the un-hatted remnant
row.
The following relations hold:
\begin{equation*}
\bar{w}_j\le 4x_j/x_k+\varpi',\qquad
q_a^{1/2}\le q\le\rho^2,\qquad
\vartheta_*\le q_a\le q^2.
\end{equation*}
\item[(b)] The band row of the hatted one-lattice Lipschitz step holds. Its linear member is the
sum over the levels $j\in\operatorname{band}(k)$, the unfiled levels (those whose remnant has not yet
been re-filed) not in $\mathfrak{B}(k)$; in depths it is the window
$\varkappa'_r\sum_{r<s<r+\ell_r}\|\hat{G}^R_s\|$.
\item[(b$'$)] The growth property of the integers $\mathsf{a}_s$ holds; in particular, by its part
(i),
$\{s:\ r<s\le r+\mathsf{a}_s+1\}\subset\{s:\ r<s\le r+\tfrac32\ell_r\}$.
\item[(c)] The operation $\mathbf{R}$ reads the older $\zeta$-slots directly, and these slots
enter the remnant row with geometric weights.
\item[(d)] Part (c) of geometric forgetting of old regions holds.
\item[(e)] Re-filing rule of the hatted step: a remnant $R^{(j)}$ is read as a remnant slot by the
inner operation at level $i$ if and only if it is present at level $i$.
\item[(f)] Presence property: if a zone of index $l$ is present inside $X^{\sim}$, then
$Z_{l+1}\cap X\neq\emptyset$, where $Z_{l+1}$ is the large-field region at level $l+1$, and hence
$j(X)\le l+1$.
\item[(g)] The weight $\bar{w}$ is absorbed into $\Lambda$ exactly as in the corollary of the
un-hatted one-lattice Lipschitz step that fixes $\Lambda$.
\item[(h)] The tube condition of the hatted one-lattice Lipschitz step holds.
\end{itemize}
Then the kernel \eqref{9.11:eq-window-kernel-2} obeys the bound
\begin{equation}\label{9.11:eq-window-kernel-3}
\hat{\mathfrak{w}}^R_{k,j}\le\bar{w}_j\,C_{**}\,q_a^{(k-j'')_+/2}
\le\bar{w}_j\,C_{**}\,\rho^{2(k-j'')_+},
\end{equation}
and the row \eqref{9.11:eq-response-rows-c} of Definition~\ref{9.11:def-response-rows} holds in
the following form.
\begin{itemize}
\item[(i)] Its band member is $\varkappa'_r\sum_{r<s<r+\ell_r}\|\hat{G}^R_s\|$, that of the
un-hatted response row of (a$_0$), unchanged.
\item[(ii)] Its remnant member is $\sum_{s>r}\mathsf{k}_{r,s}\|\hat{G}^R_s\|$ with
\begin{equation}\label{9.11:eq-window-kernel-4}
\mathsf{k}_{r,s}=\rho^{2(s-r-\mathsf{a}_s-1)_+},
\end{equation}
and the weight in the prefactor $\mathfrak{r}_r\hat{\Lambda}_r$ of
\eqref{9.11:eq-response-rows-c} is $\hat{\Lambda}_r:=C_{**}\Lambda_r$.
\item[(iii)] Its $P$-, $A$- and $\zeta$-members are those of the un-hatted response row of (a$_0$),
unchanged.
\item[(iv)] $\mathsf{k}_{r,s}=1$ on the unit window $\{s:\ r<s\le r+\mathsf{a}_s+1\}$.
\end{itemize}
\end{remark}

\begin{proof}
\emph{Depth dictionary.} We have $k=n-r-1$ and $j''=n-s+\mathsf{a}_s$, so that
\begin{equation}\label{9.11:eq-window-kernel-5}
k-j''=s-r-\mathsf{a}_s-1 .
\end{equation}
This is an identity in $\mathbb{Z}$, and it translates the level sets one by one.
\begin{itemize}
\item $\mathfrak{B}(k)=\{j:\ j''=k+1\}$ is the set where \eqref{9.11:eq-window-kernel-5} equals
$-1$. This means $s-r=\mathsf{a}_s$, so $\mathfrak{B}(k)$ corresponds to $\mathfrak{B}_r$.
\item $\operatorname{band}(k)=\{j:\ j''\ge k+2\}$ corresponds to
$\{r<s\le r+\mathsf{a}_s-1\}$. By the band window of (b) together with the growth property (b$'$)
of the integers $\mathsf{a}_s$, this set is contained in $\{r<s<r+\ell_r\}$; no equality is
asserted.
\item The unit window $\{k-j''\le0\}$ corresponds to $\{r<s\le r+\mathsf{a}_s+1\}$. By part (i) of
the growth property (b$'$), this set is contained in
$\{r<s\le r+\tfrac32\ell_r\}$.
\end{itemize}

\emph{The band member.} The band member $\varkappa'_r\sum_{r<s<r+\ell_r}\|\hat{G}^R_s\|$ is the
linear window of the band row in (b). It has the form of the band member of the un-hatted response
row of (a$_0$), and it is carried into \eqref{9.11:eq-response-rows-c}
unchanged, which gives (i).

\emph{The kernel bound.} We derive \eqref{9.11:eq-window-kernel-3} from
\eqref{9.11:eq-window-kernel-2}. By (e), the levels $i$ in the sum are exactly those at which
$R^{(j)}$ is present. We treat the young and the old reads separately.

\emph{Young reads} are those with $j\le i<j''$. They carry the factor $\vartheta_*^{k+1-i}$. Summed
geometrically, their contribution is at most
\begin{equation*}
\bar{w}_j\,\frac{\vartheta_*^{(k+2-j'')_+}}{1-\vartheta_*}.
\end{equation*}

\emph{Old unfiled reads} are those with $i\ge j''$. Each lies in a zone of index $l\le j''$ inside
$X^{\sim}$ at level $i$. By (f), $Z_{l+1}\cap X\neq\emptyset$, and so
\begin{equation*}
j(X)\le l+1\le j''+1 .
\end{equation*}
Hence the factor of (d) at such a read is $\vartheta_*^{k-j''}$. There are at most $k+2-j''$ such
levels.

Now put $m:=k-j''\ge0$. Using $\vartheta_*\le q_a$, the two contributions together are at most
\begin{align*}
\hat{\mathfrak{w}}^R_{k,j}
&\le\bar{w}_j\Bigl[\frac{q_a^2}{1-q_a}+(m+2)\Bigr]q_a^{m}\\
&\le\bar{w}_j\,(1-q_a)^{-1}(m+2)\,q_a^{m/2}\cdot q_a^{m/2}\\
&\le\bar{w}_j\,C_{**}\,q_a^{m/2}\\
&\le\bar{w}_j\,C_{**}\,\rho^{2m}.
\end{align*}
We justify each line in turn.
\begin{itemize}
\item The first line uses $\vartheta_*^{m+2}/(1-\vartheta_*)\le q_a^{m+2}/(1-q_a)$ and
$(m+2)\vartheta_*^m\le(m+2)q_a^m$.
\item The second line holds because $(m+2)/(1-q_a)-(m+2)=(m+2)q_a/(1-q_a)$, and this is at least
$q_a^2/(1-q_a)$ since $0<q_a<1$.
\item The third line is the definition \eqref{9.11:eq-window-kernel-1} of $C_{**}$, applied to the
factor $(1-q_a)^{-1}(m+2)q_a^{m/2}$.
\item The fourth line is $q_a^{1/2}\le q\le\rho^2$.
\end{itemize}
The kernel is thus flat over the range of $\mathsf{a}$ and geometric in the age beyond it.

\emph{The remnant member.} In the un-hatted response row of (a$_0$) the remnant member is the flat
sum $\sum_{s>r}\|G^R_s\|$. It is the remnant row's history kernel $\mathfrak{w}^R$ written in
reading units, that is, multiplied by the ratio $C_tx_k/(C_tx_j)$. This ratio turns
$\bar{w}_j\le4x_j/x_k+\varpi'$
into a bound by $4+\varpi'$, which is absorbed into $\Lambda$ by (g).

In the hatted row the kernel is bounded by \eqref{9.11:eq-window-kernel-3}. By
\eqref{9.11:eq-window-kernel-5}, the factor $\rho^{2(k-j'')_+}$ equals
$\rho^{2(s-r-\mathsf{a}_s-1)_+}=\mathsf{k}_{r,s}$. The member therefore becomes
$\sum_{s>r}\mathsf{k}_{r,s}\|\hat{G}^R_s\|$, once the constant $C_{**}$ is absorbed in the same way
as $\bar{w}$, namely by replacing $\Lambda$ with $\hat{\Lambda}:=C_{**}\Lambda$.

This replacement is a re-valuation of $\Lambda$. In the order in which the constants are chosen,
it takes place before $C_\sigma$, $C_E$ and $\gamma$ are chosen. Every ceiling that reads $\Lambda$
is monotone in it. Passing to the logarithmic weights gives $\hat{\Lambda}_r=C_{**}\Lambda_r$. This
proves (ii).

\emph{The other members of $\hat{B}_r(\hat{G},d)$.} The $P$- and $A$-members are those of the
un-hatted response row, unchanged, because by (a) the hatted remnant row carries the same kernels
$\mathfrak{w}^P$ and $\mathfrak{w}^D$ as the un-hatted one. The $\zeta$-member is unchanged because
the older $\zeta$-slots enter with geometric weights by (c), and the tube condition (h) places its
weights $\mathfrak{g}$ before the differences $|\delta\hat{\zeta}|$ of the hatted $\zeta$-slots.
This proves (iii).

\emph{The unit window.} On $\{r<s\le r+\mathsf{a}_s+1\}$ we have $s-r-\mathsf{a}_s-1\le0$, so
$\mathsf{k}_{r,s}=\rho^0=1$. There the hatted remnant member coincides with the un-hatted flat one.
In words: the remnant production remembers a remnant with unit weight exactly as long as that
remnant exists as a remnant. This proves (iv).
\end{proof}

\begin{proposition}[Hatted source families with geometric envelope. Stated; proof incomplete at the step marked $(\ast)$]\label{9.11:prop-source-families}
Let $C_\sigma$ be the source constant, and let $\rho$ be
the contraction rate of Definition~\ref{9.11:def-rate-floors}.
Let $(\mathsf{A},\mathsf{B})$ be a pair of runs as in Definition~\ref{9.6:def-comparison-chain},
of lengths $n$ and $n+1$, that is alive to the depth $\underline{r}$ and matched at
$\underline{r}$, and put $e_r:=\rho^{n-r}$.
Fix $r$ with $\underline{r}\le r<n$, a block $\bullet$, and a hybrid list
$\hat{\mathfrak{s}}^{\natural}_{r,\bullet}$ of the kind used in the hatted response rows,
Definition~\ref{9.11:def-response-rows}.
Then, on all domains, including those that wrap around the torus, there is a list-structured
family $\hat{N}_r^{\bullet}$ of the run $\mathsf{A}$, indexed by the components $c$ of the
block $\bullet$. For each component $c$, agreement of totals holds:
\begin{equation*}
\hat{N}_r^c\approx\varphi_r^{\mathsf{B},c}
-\varphi^c\bigl[\hat{\mathsf{F}}_r^{\mathsf{A},+,\bullet}(\hat{\mathfrak{s}}^{\natural}_{r,\bullet})\bigr].
\end{equation*}
The family satisfies the bounds
\begin{equation}\label{9.11:eq-source-families-a}
\begin{aligned}
\|\hat{N}_r^P\|,\ \|\hat{N}_r^A\|&\le C_\sigma e_r,\qquad
\|\hat{N}_r^{R''}\|\le\mathfrak{r}_r\hat{\Lambda}_r C_\sigma e_r,\\
\|\hat{N}_r^{R'}\|&\le\mathfrak{r}_r\hat{\Lambda}_r\bigl[(C_\sigma+1)e_r+\hat{B}_r(\hat{G},d)\bigr]
+\varkappa'_r\sum_{r<s<r+\ell_r}\|\hat{G}_s^R\|,
\end{aligned}
\end{equation}
in which:
\begin{itemize}
\item $\hat{N}_r^P$ and $\hat{N}_r^A$ are the families at the blocks $P$ and $A$, and
$\hat{N}_r^{R'}$ and $\hat{N}_r^{R''}$ are the two members of the family at the block $R$;
\item $\hat{B}_r(\hat{G},d)$ is the response sum of \eqref{9.11:eq-response-rows-c};
\item $\mathfrak{r}_r$ and the logarithmic weight $\hat{\Lambda}_r$ are those of
\eqref{9.11:eq-rate-floors-c};
\item $\varkappa'_r$ and the window length $\ell_r$ are those of
Definition~\ref{9.11:def-response-rows}.
\end{itemize}
\end{proposition}

Stated; the proof below is incomplete at the step marked $(\ast)$.

\begin{proof}
\emph{The members at $P$, $A$ and $R''$.} These three bounds of
\eqref{9.11:eq-source-families-a} are obtained as an implication from two inputs: the statements
from which the un-hatted source-family statement is derived, read without a floor condition,
which bound the hatted class $\hat{\mathfrak{K}}$ by $C_\sigma\rho^{n-r}=C_\sigma e_r$; and the
re-filing lemma.

\emph{Comparison with the un-hatted bounds.} The bounds \eqref{9.11:eq-source-families-a} have
the shape of the bounds of the un-hatted source-family statement for the same pair of runs, with
three changes:
\begin{itemize}
\item the envelope is $e_r=\rho^{n-r}$, without the factor $\Pi_n$ of the un-hatted envelope;
\item the response sum is the hatted one, $\hat{B}_r(\hat{G},d)$ of
\eqref{9.11:eq-response-rows-c}, in place of the un-hatted one;
\item the logarithmic weight is $\hat{\Lambda}_r$ in place of $\Lambda_r$.
\end{itemize}

\emph{The term carrying the alignment index.} The alignment index $n_c$ of the alignment
statement enters the argument at one place only, through clause (iii) of the admissibility
conditions. That clause is used in the form
\begin{equation}\label{9.11:eq-source-families-1}
q_a^{\,n-r-n_c}\,e^{-p_0(g_{(r)})}\le C\,\mathfrak{r}_r\Lambda_r\,\rho^{n-r},
\end{equation}
where $q_a$ is the age rate of Definition~\ref{9.11:def-rate-floors}, $p_0(\cdot)$ is the
degree function, and $C$ is the constant of clause (iii) of the admissibility conditions. By
\eqref{9.11:eq-rate-floors-c} we have $\hat{\Lambda}_r\ge\Lambda_r$, and
$\rho^{n-r}=e_r$. Hence the right member of \eqref{9.11:eq-source-families-1} satisfies
\begin{equation*}
C\,\mathfrak{r}_r\Lambda_r\,\rho^{n-r}\le C\,\mathfrak{r}_r\hat{\Lambda}_r\,e_r .
\end{equation*}
The bound on the right is a multiple of $\mathfrak{r}_r\hat{\Lambda}_r e_r$, and the term
\eqref{9.11:eq-source-families-1} is absorbed in the summand $+1$ of the factor
$(C_\sigma+1)e_r$ of the $R'$-member of \eqref{9.11:eq-source-families-a}, as in the un-hatted
bound.

The alignment statement involves neither $\Pi_n$ nor any hatted quantity. It is therefore used
here unchanged, with the same arithmetic in the index $n_c$ as in the un-hatted case.

$(\ast)$ \emph{The $R'$-member with the hatted quantities.} It remains to show that the
contribution to $\hat{N}_r^{R'}$ other than the one bounded by
\eqref{9.11:eq-source-families-1} has norm at most
\begin{equation*}
\mathfrak{r}_r\hat{\Lambda}_r\bigl[C_\sigma e_r+\hat{B}_r(\hat{G},d)\bigr]
+\varkappa'_r\sum_{r<s<r+\ell_r}\|\hat{G}_s^R\|.
\end{equation*}
With $e_r\Pi_n$ in place of $e_r$, and with $B_r$, $\Lambda_r$ and $G_s^R$ in place of the
hatted quantities $\hat{B}_r$, $\hat{\Lambda}_r$ and $\hat{G}_s^R$, this bound, together with
the term \eqref{9.11:eq-source-families-1} carried in the summand $+1$, is the $R'$-member of
the un-hatted source-family statement. In the hatted quantities the bound is not derived here.
\end{proof}

\begin{lemma}[Re-filed shares add nothing to the source families]
\label{9.11:lem-refiled-shares}
Let $\mathsf{A}$ and $\mathsf{B}$ be the two runs of a pair, of lengths $n$ and $n+1$, with the
hatted lists of Definition~\ref{9.11:def-hatted-lists} and the hatted blocks of
Definition~\ref{9.11:def-hatted-blocks}. For a family $F$ write $\operatorname{tot}(F)$ for its
total. For $\mathsf{m}\in\{\mathsf{A},\mathsf{B}\}$ write $\operatorname{\mathsf{RF}}^{\mathsf{m}}$
for the re-filing map $\operatorname{\mathsf{RF}}$ on the lattice of run $\mathsf{m}$,
$R^{\mathsf{m}}_{(s)}$ for the $R$-slot of run $\mathsf{m}$ at depth $s$, and
$(\cdot)^{\mathrm{vac}}$ for the vacuum part of a slot or of its total. For $r<n$ let
$\mathfrak{B}_r$ be the set of depths $s>r$ with $s-r=\mathsf{a}_s$. Assume:
\begin{itemize}
\item[(a)] the total-preservation property of the re-filing map $\operatorname{\mathsf{RF}}$ of
Definition~\ref{9.11:def-hatted-lists}, established with that map: on each run's lattice the
re-filed remnant is equal in total to the vacuum part of the remnant it re-files, that is, for
$\mathsf{m}\in\{\mathsf{A},\mathsf{B}\}$ and every $R$-slot $R$,
$\operatorname{tot}(\operatorname{\mathsf{RF}}^{\mathsf{m}}[R])=\operatorname{tot}(R^{\mathrm{vac}})$;
\item[(b)] totals are additive;
\item[(c)] for every $r<n$ and every $s\in\mathfrak{B}_r$,
$\hat{G}_s^R\approx\varphi_s^{\mathsf{B},R}-\varphi_s^{\mathsf{A},R}$,
which is clause (c1) of \eqref{9.11:eq-chain-bounds-a}.
\end{itemize}
Then for every $r<n$ the re-filed parts of the $\hat{Q}$- and $\hat{D}$-components of
\begin{equation*}
\varphi_r^{\mathsf{B},\cdot}-\varphi^{\cdot}\bigl[\hat{\mathsf{F}}_r^{\mathsf{A},+,\bullet}
(\hat{\mathfrak{s}}^{\natural}_{r,\bullet})\bigr]
\end{equation*}
have total zero (up to the vacuum constant in the frozen case). Hence the source families
$\hat{N}_r^{\hat{Q}}$ and $\hat{N}_r^{\hat{D}}$ of \eqref{9.11:eq-source-families-a} may be taken to be the
families $N_r^Q$ and $N_r^D$ supplied for the $Q$- and $D$-parts alone.
\end{lemma}

\begin{proof}
The families of \eqref{9.11:eq-source-families-a} are those of
Proposition~\ref{9.11:prop-source-families}, which is stated with proof incomplete at the step
marked $(\ast)$; only the defining agreement of totals required of them there is used below.

Fix $r<n$ and $s\in\mathfrak{B}_r$, that is, $s>r$ with $s-r=\mathsf{a}_s$. Run $\mathsf{B}$'s
contribution from depth $s$ to the slot $\hat{Q}_{(r)}$ is
$\operatorname{\mathsf{RF}}^{\mathsf{B}}[R^{\mathsf{B}}_{(s)}]$. By hypothesis~(a) the total of this
contribution is
\begin{equation*}
\operatorname{tot}\bigl(\operatorname{\mathsf{RF}}^{\mathsf{B}}[R^{\mathsf{B}}_{(s)}]\bigr)
=\operatorname{tot}\bigl((R^{\mathsf{B}}_{(s)})^{\mathrm{vac}}\bigr)
=(\varphi_s^{\mathsf{B},R})^{\mathrm{vac}} .
\end{equation*}
Run $\mathsf{A}$'s extended production on the hybrid,
$\hat{\mathsf{F}}_r^{\mathsf{A},+,\bullet}(\hat{\mathfrak{s}}^{\natural}_{r,\bullet})$, re-files the
hybrid's $R$-slot at depth $s$. That slot is $R^{\mathsf{A}}_{(s)}+\hat{G}_s^R$. By hypotheses~(b)
and~(c) its total is
\begin{equation}\label{9.11:eq-refiled-shares-1}
\varphi_s^{\mathsf{A},R}+\bigl(\varphi_s^{\mathsf{B},R}-\varphi_s^{\mathsf{A},R}\bigr)
=\varphi_s^{\mathsf{B},R}.
\end{equation}
By hypothesis~(a) applied to the slot $R^{\mathsf{A}}_{(s)}+\hat{G}_s^R$, and since taking totals
and taking vacuum parts are linear in the slot, \eqref{9.11:eq-refiled-shares-1} gives
\begin{equation*}
\operatorname{tot}\bigl(\operatorname{\mathsf{RF}}^{\mathsf{A}}[R^{\mathsf{A}}_{(s)}
+\hat{G}_s^R]\bigr)=(\varphi_s^{\mathsf{B},R})^{\mathrm{vac}} .
\end{equation*}
Thus the two contributions from depth $s$ have the same total, and their difference has total
zero.

Sum over $s\in\mathfrak{B}_r$. By hypothesis~(b), the total of the sum of these differences is the
sum of their totals, which is zero. By Definition~\ref{9.11:def-hatted-lists}, the re-filed part
of the $\hat{Q}$-difference at depth $r$ is made of exactly the remnants re-filed at depth $r$,
namely those from the depths $s$ with $s-\mathsf{a}_s=r$. Hence this re-filed part has total zero.
The zero family therefore agrees in total with it.

The $\hat{Q}$-difference is the sum of its $Q$-part and its re-filed part. The family $N_r^Q$
supplied for the $Q$-part agrees in total with the $Q$-part. By hypothesis~(b), the total of the
$\hat{Q}$-difference is the sum of the totals of the two parts, that is, $\operatorname{tot}(N_r^Q)+0$.
Hence $\hat{N}_r^{\hat{Q}}:=N_r^Q$ satisfies
\begin{equation*}
\hat{N}_r^{\hat{Q}}\approx\varphi_r^{\mathsf{B},\hat{Q}}
-\varphi^{\hat{Q}}\bigl[\hat{\mathsf{F}}_r^{\mathsf{A},+,\bullet}
(\hat{\mathfrak{s}}^{\natural}_{r,\bullet})\bigr].
\end{equation*}

The $\hat{D}$-component is the same computation, made twice and then subtracted. In the notation
of Definition~\ref{9.11:def-hatted-blocks}, with $R_t=R_{(t)}$ and $\operatorname{\mathsf{RF}}_r$ the
reading of the re-filing map at depth $r=t-\mathsf{a}_t$, its re-filed part at depth $r$ is
\begin{equation*}
\sum_{t-\mathsf{a}_t=r}\Bigl(\operatorname{\mathsf{RF}}_r[R_t](\cdot;W)
-\operatorname{\mathsf{RF}}_r[R_t]_{W(p_y)}\Bigr).
\end{equation*}
The argument above is carried out once pointwise in the source $W$ and once at the constant value
$W(p_y)$. In each case the re-filed part of the difference has total zero. Subtracting the two,
the re-filed part of the $\hat{D}$-difference has total zero. As for $\hat{Q}$, the family
$\hat{N}_r^{\hat{D}}:=N_r^D$ then agrees in total with the $\hat{D}$-component of the
difference.
\end{proof}

\begin{definition}[The ceilings on the coupling for the hatted system]\label{9.11:def-coupling-ceilings}
Let $q$, $q_a$ and $\rho$ be as in Definition~\ref{9.11:def-rate-floors}. Let
$\hat{\mathfrak{r}}_r$, $\hat{\Lambda}_r$, $\Lambda_r$, $C_{**}$ and $\hat{c}_\rho$ be the constants
of \eqref{9.11:eq-rate-floors-c}; in particular
\begin{equation*}
\hat{c}_\rho=\frac{1+\rho}{\rho(1-\rho)^2}.
\end{equation*}
Let $(\theta_s)_{s\ge0}$ be the common ladder, $\theta_0:=X=g^{-2}$ and $\theta_{s+1}=\theta_s+c_0$.
Let $\mathsf{a}_s$ be the integer that fixes the band $s-r\le\mathsf{a}_s$ and the unit window
$s-r\le\mathsf{a}_s+1$ of \eqref{9.11:eq-memory-structures-a}. We put
\begin{equation}\label{9.11:eq-coupling-ceilings-1}
\mathfrak{S}:=\sum_{s\ge0}\hat{\mathfrak{r}}_s\,\rho^{-\mathsf{a}_s}.
\end{equation}

The \emph{ceilings on the coupling for the hatted system} are the conditions displayed below. Each
condition is a ceiling on the coupling $g$, that is, a floor on $X=\theta_0$; it is either imposed
after $L$, $M$, $\kappa$, $\kappa_0$, the exponents, $q$, $q_a$, $\rho$, the coefficient constant
$\mathsf{A}$, $C_t$, $\Lambda$ and $C_{**}$ are fixed, or it is already implied by the
corresponding condition among the ceilings on the coupling for the un-hatted system, which is itself
such a ceiling. Here $\alpha_r$, $\alpha_r^A$, $\varkappa$, $\varkappa'_r$, $\varkappa''_r$,
$\varkappa_A$ are the small coefficients of the block inequalities of the hatted system, $C_F$ and
$C_F^A$ its band coefficients, $\check{R}_{(r)}$ is the waiting time of the large-field operation at
depth $r$, so that the window length is $\ell_r=2+\log_L(M\check{R}_{(r)})=\mathsf{a}_r+2$, and $C_t$,
$\hat{\mathfrak{m}}_{(s)}$, $c_1$ are the constant, the hatted $Q$-bound and the exponent that enter
the reading units of the hatted blocks of Definition~\ref{9.11:def-hatted-blocks}.
The conditions read as follows.
\begin{equation}\label{9.11:eq-coupling-ceilings-a}
\begin{gathered}
\text{(\^{a})}\quad \hat{c}_\rho\,\alpha_r\,\mathsf{A}\le\tfrac18\ \text{ for every } r;
\qquad
\text{(b)}\quad \varkappa\le\frac{1-\rho}{32};\\
\text{(b$'$)}\quad 4\varkappa''_r\le\tfrac{3}{32}(1-\rho)\ \text{ for every } r;
\qquad
\text{(\^{c})}\quad \hat{c}_\rho\,\alpha_r^A\,\mathsf{A}\le\tfrac14\ \text{ for every } r;\\
\text{(d)}\quad \varkappa_A\Bigl(9+\frac{5}{1-\rho}\Bigr)\le\tfrac14.
\end{gathered}
\end{equation}
\begin{equation}\label{9.11:eq-coupling-ceilings-b}
\text{(\^{e})}\quad \hat{\mathfrak{r}}_r\le\tfrac{1}{16}\ \text{ for every } r\ge0,
\qquad\text{and}\qquad
\hat{\mathfrak{r}}_s\le\hat{\mathfrak{r}}_r\ \text{ whenever } s\ge r\ge0.
\end{equation}
\begin{equation}\label{9.11:eq-coupling-ceilings-c}
\text{(\^{g}$'$)}\quad 4\varkappa'_r\,\ell_r^{2}\,\rho^{-\ell_r}\le\tfrac14\ \text{ for every } r.
\end{equation}
\begin{equation}\label{9.11:eq-coupling-ceilings-d}
\text{(\^{h})}\quad \max\{C_F,\,C_F^A,\,\rho^{-1}\}\cdot\mathfrak{S}\le\tfrac{1}{32}.
\end{equation}
\begin{equation}\label{9.11:eq-coupling-ceilings-e}
\begin{gathered}
\text{(u)}\quad \hat{\mathfrak{u}}_s^Q=C_t\,\theta_s\,\hat{\mathfrak{m}}_{(s)}\le1\ \text{ at every depth } s,
\qquad
\mathsf{l}_s=(\log\theta_s)^{c_1}\ge1,\\
\hat{\Lambda}_r\ge\Lambda_r\ge1\ \text{ for every } r.
\end{gathered}
\end{equation}
\end{definition}

\emph{Relation to the ceilings for the un-hatted system.} Conditions (\^{a}) and (\^{c}) are the
un-hatted conditions (a) and (c) with the product $c_\rho\mathsf{n}_r$ of the un-hatted constant
$c_\rho$ and the run-length factor $\mathsf{n}_r$ replaced by $\hat{c}_\rho$; that the run-length
factor may be dropped is the content of part~(i) of the lemma that removes it from the hatted bounds.
Conditions (b), (b$'$), (d) and (\^{g}$'$) coincide with their un-hatted counterparts, read at the
hatted constants; in particular (\^{g}$'$) in \eqref{9.11:eq-coupling-ceilings-c} is, verbatim, the
corresponding un-hatted ceiling, and the floor $\rho_{g'}$ on $\rho$ that enters $\hat{\rho}_0$ in
Definition~\ref{9.11:def-rate-floors} is the floor on $\rho$ under which that ceiling can be met.
Condition (\^{e}) takes the place of the un-hatted condition (e): the members
$\sum_r\hat{\mathfrak{r}}_r\le1$ and $\Pi_n\le1$ of that condition are not imposed here, and the
rate $\hat{\mathfrak{r}}_r$ carries the weight $\hat{\Lambda}_r$ in place of $\Lambda_r$. The
monotonicity half of (\^{e}) takes the place of the un-hatted condition (f). The first member of
that condition,
\begin{equation*}
\hat{\mathfrak{r}}_s\le C_T\,\rho^{-(n-s)/2}\,\hat{\mathfrak{r}}_n,
\end{equation*}
is not used, and neither is the constant $C_T$; its second member is kept without the factor
$2(1+s-r)$, which came from the ratio $\mathsf{n}_s/\mathsf{n}_r$. Condition (\^{h}) is condition
(c) of the ceilings for the renormalised system, written for the three blocks, with
$\hat{\Lambda}_r$ inside $\hat{\mathfrak{r}}_r$ and with $\rho^{-1}$ in the maximum in place of the
un-hatted coefficient $C_F^R$.

In (u), the bound on $\hat{\mathfrak{u}}_s^Q$ is required at every depth, so that it does not depend
on the run length $n$. The remaining inequalities of (u) rest on the following:
\begin{itemize}
\item the ceiling $\log\theta_0\ge1$ on $g$, whence $\log\theta_s\ge1$ for every $s$, the ladder
being increasing;
\item the one-sided premise $\Lambda\ge1$;
\item $C_{**}\ge2$, which holds because $C_{**}=(1-q_a)^{-1}\sup_m(m+2)q_a^{m/2}$ of
\eqref{9.11:eq-rate-floors-c} is at least its term with $m=0$, namely $2/(1-q_a)$, and
$2/(1-q_a)\ge 2$ since $0<q_a<1$.
\end{itemize}
Hence, with $\Lambda_r=\Lambda(\log\theta_r)^{c_3}$ and $\hat{\Lambda}_r=C_{**}\Lambda_r$ as in
\eqref{9.11:eq-rate-floors-c},
\begin{equation*}
\mathsf{l}_s=(\log\theta_s)^{c_1}\ge1,\qquad \Lambda_r\ge1,\qquad
\hat{\Lambda}_r=C_{**}\Lambda_r\ge2\Lambda_r\ge\Lambda_r.
\end{equation*}

\emph{Consequences of} (\^{h}). Each of $C_F$, $C_F^A$ and $\rho^{-1}$ is at most
$\max\{C_F,C_F^A,\rho^{-1}\}$. Moreover $\mathfrak{S}\ge0$, since it is a sum of non-negative terms.
Multiplying each of the three numbers by $\mathfrak{S}$ and using
\eqref{9.11:eq-coupling-ceilings-d} therefore gives
\begin{equation}\label{9.11:eq-coupling-ceilings-2}
C_F\,\mathfrak{S}\le\tfrac{1}{32},\qquad C_F^A\,\mathfrak{S}\le\tfrac{1}{32},\qquad
\rho^{-1}\mathfrak{S}\le\tfrac{1}{32}.
\end{equation}
Since $C_F\ge1$, the first of these inequalities gives in turn
$\mathfrak{S}\le C_F\,\mathfrak{S}\le\tfrac{1}{32}$.

None of the conditions \eqref{9.11:eq-coupling-ceilings-a}--\eqref{9.11:eq-coupling-ceilings-e}
involves the constants $C_V$, $C_\sigma$, $C_E$ or the run length $n$.

The factor $\rho^{-1}$ occurs in the maximum of (\^{h}) for the following reason. The window kernel
$\mathsf{k}_{r,s}$ of the hatted system equals $1$ on the unit window $s-r\le\mathsf{a}_s+1$ of
\eqref{9.11:eq-memory-structures-a}. That window is one step longer than the band
$s-r\le\mathsf{a}_s$, and this one step contributes the factor $\rho^{-1}$; see the window sum in
\eqref{9.11:eq-envelope-sums-d}.

\begin{lemma}[The coupling ceilings hold along the ladder]\label{9.11:lem-ceilings-ladder}
Let the kernel rate $q$, the age rate $q_a$ and the contraction rate $\rho$ satisfy the one-sided
floors \eqref{9.11:eq-rate-floors-a} and \eqref{9.11:eq-rate-floors-b} of
Definition~\ref{9.11:def-rate-floors}. Let $c_0>0$, and let $\theta_0:=X$ and
$\theta_{s+1}:=\theta_s+c_0$, so that $\theta_s=X+sc_0$ for every $s\ge 0$. The conditions
(\^{a}), (b), (b$'$), (\^{c}), (d), (\^{e}), (\^{g}$'$), (\^{h}) and (u) are those of
Definition~\ref{9.11:def-coupling-ceilings}, displayed in \eqref{9.11:eq-coupling-ceilings-a},
\eqref{9.11:eq-coupling-ceilings-b}, \eqref{9.11:eq-coupling-ceilings-c},
\eqref{9.11:eq-coupling-ceilings-d} and \eqref{9.11:eq-coupling-ceilings-e}. The hatted
$R$-rate $\hat{\mathfrak{r}}_r$, the weights $\Lambda_r$, $\hat{\Lambda}_r$ and the constant
$\hat{c}_\rho=(1+\rho)/(\rho(1-\rho)^2)$ are those of \eqref{9.11:eq-rate-floors-c}, and
$\mathfrak{S}=\sum_{s\ge0}\hat{\mathfrak{r}}_s\rho^{-\mathsf{a}_s}$, with $\mathsf{a}_s$ the
integer of Definition~\ref{9.11:def-hatted-lists}. Put
\begin{equation*}
c_\rho:=\frac{5}{\rho(1-\rho)},\qquad \mathsf{n}_r:=r+1+(1-\rho)^{-1}\quad(r\ge 0).
\end{equation*}
\begin{enumerate}
\item[(i)] If
\begin{equation}\label{9.11:eq-ceilings-ladder-1}
c_\rho\,\alpha_r\,\mathsf{A}\,\mathsf{n}_r\le\tfrac18\qquad\text{for every }r\ge0,
\end{equation}
then (\^{a}) holds. If
\begin{equation}\label{9.11:eq-ceilings-ladder-2}
c_\rho\,\alpha^A_r\,\mathsf{A}\,\mathsf{n}_r\le\tfrac14\qquad\text{for every }r\ge0,
\end{equation}
then (\^{c}) holds.
\item[(ii)] If $\kappa_0>2$, there is a threshold
$\hat{X}_{\hat{e}}=\hat{X}_{\hat{e}}(\rho,\mathsf{A},C_t,\Lambda,C_{**},c_3,\kappa_0)$ such
that (\^{e}) holds for every $X\ge\hat{X}_{\hat{e}}$, every $r\ge0$ and every $s\ge r$.
\item[(iii)] Assume $\kappa_0>4$, and assume that the waiting times $\check{R}_{(s)}$ of the
large-field operation satisfy
\begin{equation}\label{9.11:eq-ceilings-ladder-3}
\check{R}_{(s)}\le C_{\check{R}}\,(\log\theta_s)^{r_{\mathrm{pr}}}\qquad\text{for every }s\ge0,
\end{equation}
with constants $C_{\check{R}}\ge1$ and $r_{\mathrm{pr}}\ge0$. Then there is a threshold
\begin{equation*}
\hat{X}_{\hat{h}}=\hat{X}_{\hat{h}}\bigl(\rho,\mathsf{A},C_t,\Lambda,C_{**},c_3,\kappa_0,
M,C_{\check{R}},r_{\mathrm{pr}},C_F,C^A_F,c_0\bigr)
\end{equation*}
such that (\^{h}) holds for every $X\ge\hat{X}_{\hat{h}}$. The threshold depends on $L$ only
through the constants listed and through the inequality $\lambda_\rho<\frac14$ established in
the proof, whose value is not used.
\item[(iv)] For every $X\ge e$ the clauses $\mathsf{l}_s\ge1$ and
$\hat{\Lambda}_r\ge\Lambda_r\ge1$ of (u) hold.
\end{enumerate}
The displays \eqref{9.11:eq-ceilings-ladder-1} and \eqref{9.11:eq-ceilings-ladder-2} are the
conditions (a) and (c) of the ceilings on the coupling for the un-hatted system, read at the
hatted constants and with the hatted running coupling $\hat{x}_{(s)}$ in place of $x_{(s)}$;
they are hypotheses of this lemma. The conditions (b), (b$'$), (d), (\^{g}$'$) and the first
clause of (u) are the un-hatted system's conditions (b), (b$'$), (d), (g$'$) and the first clause
of its condition (u), read at the hatted constants $\varkappa$, $\varkappa''_r$, $\varkappa_A$,
$\varkappa'_r$ (functions of depth), with (g$'$) under the floor $\rho\ge\rho_{g'}$ contained in
\eqref{9.11:eq-rate-floors-b}, and read at the bound $\hat{\mathfrak{m}}$ of the hatted class on
the hatted $Q$-slot, of which only the inequality
\begin{equation*}
C_t\,\hat{x}_k\,\hat{\mathfrak{m}}_k\le\hat{\mathsf{v}}_k/C_\beta\le\mathsf{v}_k/C_\beta
\end{equation*}
is used; these conditions, in this reading, are hypotheses of this lemma as well. No quantity
in (i)--(iv) bounds $L$, $n$, $r_0$, $\mathsf{a}_s$ or $\ell_r$ from above, and no value of
$L_0$ is used.
\end{lemma}

\begin{proof}
(i) From the definitions of $c_\rho$ and $\mathsf{n}_r$,
\begin{equation*}
c_\rho\mathsf{n}_r=\frac{5(r+1)}{\rho(1-\rho)}+\frac{5}{\rho(1-\rho)^2},
\qquad\text{so}\qquad
c_\rho\mathsf{n}_r-\frac{5}{\rho(1-\rho)^2}=\frac{5(r+1)}{\rho(1-\rho)}\ge0 .
\end{equation*}
Since $0<\rho<1$ we have $1+\rho\le2\le5$, and therefore
\begin{equation*}
\frac{5}{\rho(1-\rho)^2}\ge\frac{1+\rho}{\rho(1-\rho)^2}=\hat{c}_\rho .
\end{equation*}
Hence $\hat{c}_\rho\le c_\rho\mathsf{n}_r$ for every $r\ge0$. Since $\alpha_r\mathsf{A}\ge0$,
multiplying this inequality by $\alpha_r\mathsf{A}$ and using
\eqref{9.11:eq-ceilings-ladder-1} gives
\begin{equation*}
\hat{c}_\rho\,\alpha_r\mathsf{A}\le c_\rho\mathsf{n}_r\,\alpha_r\mathsf{A}\le\tfrac18
\qquad\text{for every }r\ge0,
\end{equation*}
which is (\^{a}). The same argument with $\alpha^A_r$ in place of $\alpha_r$,
\eqref{9.11:eq-ceilings-ladder-2} in place of \eqref{9.11:eq-ceilings-ladder-1} and $\frac14$ in
place of $\frac18$ gives (\^{c}).

(ii) Put
\begin{equation*}
K_{\hat{e}}:=2C_tC_{**}\Lambda\bigl[\mathsf{A}(1-\rho)^{-2}+6(1-\rho)^{-1}+8\bigr],
\qquad
\psi(\theta):=\theta^{1-\kappa_0/2}(\log\theta)^{c_3}\quad(\theta>1).
\end{equation*}
By the definition of the hatted $R$-rate in \eqref{9.11:eq-rate-floors-c},
$\hat{\mathfrak{r}}_r=K_{\hat{e}}\,\psi(\theta_r)$. For $\theta>1$,
\begin{equation*}
\psi'(\theta)=\theta^{-\kappa_0/2}(\log\theta)^{c_3-1}
\Bigl[\bigl(1-\tfrac{\kappa_0}{2}\bigr)\log\theta+c_3\Bigr],
\end{equation*}
and for $\kappa_0>2$ the bracket is $\le0$ exactly when $\log\theta\ge2c_3/(\kappa_0-2)$. Put
$X_\psi:=\exp\bigl(2c_3/(\kappa_0-2)\bigr)$. Then $\psi$ is non-increasing on $[X_\psi,\infty)$,
and, since the exponent $1-\kappa_0/2$ is negative, $\psi(\theta)\to0$ as $\theta\to\infty$. The
function $X\mapsto K_{\hat{e}}\psi(X)$ is therefore continuous and non-increasing on
$[X_\psi,\infty)$ with limit $0$, so the set
\begin{equation*}
\bigl\{X\ge\max\{e,X_\psi\}:\ K_{\hat{e}}\psi(X)\le\tfrac1{16}\bigr\}
\end{equation*}
is non-empty (because of the limit), closed (by continuity) and a ray (by monotonicity); let
$\hat{X}_{\hat{e}}$ be its least element. It depends only on $K_{\hat{e}}$, $c_3$ and $\kappa_0$,
hence only on the constants listed in (ii). Let $X\ge\hat{X}_{\hat{e}}$ and $r\ge0$. Since
$c_0>0$, $\theta_r=X+rc_0\ge X\ge X_\psi$, so
\begin{equation*}
\hat{\mathfrak{r}}_r=K_{\hat{e}}\psi(\theta_r)\le K_{\hat{e}}\psi(X)
\le K_{\hat{e}}\psi(\hat{X}_{\hat{e}})\le\tfrac1{16}.
\end{equation*}
If moreover $s\ge r$, then $\theta_s\ge\theta_r\ge X_\psi$, and the monotonicity of $\psi$ gives
$\psi(\theta_s)\le\psi(\theta_r)$, that is $\hat{\mathfrak{r}}_s\le\hat{\mathfrak{r}}_r$. This is
(\^{e}).

(iii) By the floors of Definition~\ref{9.11:def-rate-floors}, $q\le\rho^2$ and $q>L^{-1/2}$, so
$\rho\ge q^{1/2}>L^{-1/4}$. Since also $\rho<1$, the number $\lambda_\rho:=\log_L(1/\rho)$
satisfies $0<\lambda_\rho<\frac14$. Under the integer reading
$\mathsf{a}_s=\log_L(M\check{R}_{(s)})$ of Definition~\ref{9.11:def-hatted-lists},
\begin{equation*}
\rho^{-\mathsf{a}_s}=L^{\lambda_\rho\mathsf{a}_s}=\bigl(M\check{R}_{(s)}\bigr)^{\lambda_\rho}.
\end{equation*}
Let $X\ge e$; then $\theta_s\ge X\ge e$, so $\log\theta_s\ge1$. By
\eqref{9.11:eq-ceilings-ladder-3}, and since $M\ge1$ and $C_{\check{R}}\ge1$, both
$MC_{\check{R}}$ and $\log\theta_s$ are numbers $\ge1$; raising the exponent of a number $\ge1$
from $\lambda_\rho$ to $\frac14$ does not decrease it, hence
\begin{equation*}
\rho^{-\mathsf{a}_s}\le(MC_{\check{R}})^{\lambda_\rho}(\log\theta_s)^{\lambda_\rho r_{\mathrm{pr}}}
\le(MC_{\check{R}})^{1/4}(\log\theta_s)^{r_{\mathrm{pr}}/4}.
\end{equation*}
Put $c_\Psi:=c_3+r_{\mathrm{pr}}/4$ and $\Psi(t):=t^{1-\kappa_0/2}(\log t)^{c_\Psi}$ for $t>1$, so
that $\psi(\theta)(\log\theta)^{r_{\mathrm{pr}}/4}=\Psi(\theta)$. Inserting
$\hat{\mathfrak{r}}_s=K_{\hat{e}}\psi(\theta_s)$ from (ii) and the last bound into the definition
of $\mathfrak{S}$, and using $\theta_s=X+sc_0$, we obtain
\begin{equation}\label{9.11:eq-ceilings-ladder-4}
\max\{C_F,C^A_F,\rho^{-1}\}\,\mathfrak{S}\le K_{\hat{h}}\sum_{s\ge0}\Psi(X+sc_0),
\qquad
K_{\hat{h}}:=\max\{C_F,C^A_F,\rho^{-1}\}\,K_{\hat{e}}\,(MC_{\check{R}})^{1/4}.
\end{equation}
As in (ii), with $c_\Psi\ge0$ in place of $c_3$, the function $\Psi$ is non-increasing on
$[X_\Psi,\infty)$, where $X_\Psi:=\exp\bigl(2c_\Psi/(\kappa_0-2)\bigr)$. Since $\kappa_0>4$, the
exponent $1-\kappa_0/2$ is $<-1$. With $\eta:=\kappa_0/4-1>0$ we have
$\Psi(t)=t^{-1-\eta}\cdot t^{-\eta}(\log t)^{c_\Psi}$, and $t^{-\eta}(\log t)^{c_\Psi}$ is
continuous on $[e,\infty)$ and tends to $0$, hence is bounded there; as $\Psi$ is continuous on
$(1,\infty)$, it follows that $\int_X^\infty\Psi(t)\,dt<\infty$ for every $X>1$, and this integral
tends to $0$ as $X\to\infty$. Let $X\ge X_\Psi$ and $s\ge1$. On the interval
$[X+(s-1)c_0,\,X+sc_0]\subset[X_\Psi,\infty)$ the function $\Psi$ is at least
$\Psi(X+sc_0)$, so
\begin{equation*}
\Psi(X+sc_0)\le c_0^{-1}\int_{X+(s-1)c_0}^{X+sc_0}\Psi(t)\,dt .
\end{equation*}
Summing over $s\ge1$ and adding the term $s=0$,
\begin{equation*}
\sum_{s\ge0}\Psi(X+sc_0)\le\Gamma(X):=\Psi(X)+c_0^{-1}\int_X^\infty\Psi(t)\,dt .
\end{equation*}
The function $\Gamma$ is finite and continuous on $[X_\Psi,\infty)$; it is non-increasing there,
because $\Psi$ is and because $\Psi\ge0$; and $\Gamma(X)\to0$ as $X\to\infty$. As in (ii), the set
of $X\ge\max\{e,X_\Psi\}$ with $K_{\hat{h}}\Gamma(X)\le\frac1{32}$ is a non-empty closed ray; let
$\hat{X}_{\hat{h}}$ be its least element. For $X\ge\hat{X}_{\hat{h}}$,
\eqref{9.11:eq-ceilings-ladder-4} and the monotonicity of $\Gamma$ give
\begin{equation*}
\max\{C_F,C^A_F,\rho^{-1}\}\,\mathfrak{S}\le K_{\hat{h}}\Gamma(X)
\le K_{\hat{h}}\Gamma(\hat{X}_{\hat{h}})\le\tfrac1{32},
\end{equation*}
which is (\^{h}). The threshold $\hat{X}_{\hat{h}}$ is determined by $K_{\hat{h}}$, $c_\Psi$,
$\kappa_0$ and $c_0$, hence by the constants listed in (iii) only. The series defining
$\mathfrak{S}$ runs over all $s\ge0$, so the run length $n$ is not read; $\mathsf{a}_s$ is bounded
only through \eqref{9.11:eq-ceilings-ladder-3}, that is, by a function of $\theta_s$ that grows
without bound along the ladder, not by a uniform constant; $\ell_r$ does not occur; and no value
of $L$ or $L_0$ enters.

If $\mathsf{a}_s$ is read instead as $\max\{\lceil\log_L(M\check{R}_{(s)})\rceil,1\}$, then, as
soon as $M\check{R}_{(s)}\ge1$, one has $\mathsf{a}_s\le\log_L(M\check{R}_{(s)})+1$, hence
$\rho^{-\mathsf{a}_s}\le\rho^{-1}(M\check{R}_{(s)})^{\lambda_\rho}$, and the argument above goes
through with $K_{\hat{h}}$ multiplied by one further factor $\rho^{-1}$.

(iv) Let $X\ge e$. Then $\theta_s\ge X\ge e$ for every $s$, so $\log\theta_s\ge1$, and since
$c_1,c_3\ge0$ and $\Lambda\ge1$,
\begin{equation*}
\mathsf{l}_s=(\log\theta_s)^{c_1}\ge1,\qquad \Lambda_r=\Lambda(\log\theta_r)^{c_3}\ge1 .
\end{equation*}
Since $C_{**}\ge2$, $\hat{\Lambda}_r=C_{**}\Lambda_r\ge2\Lambda_r\ge\Lambda_r$. The remaining
conditions named in the lemma, (b), (b$'$), (d), (\^{g}$'$) and the first clause of (u), enter as
hypotheses in the reading stated there, and nothing further is to be proved.
\end{proof}

The constants are chosen in the following order: first the parameters of the statements on
which this lemma rests; then $q$ and $q_a$; then $\rho$ subject to its floors
\eqref{9.11:eq-rate-floors-b}; then $X\ge\max\{e,\hat{X}_{\hat{e}},\hat{X}_{\hat{h}}\}$, together
with the thresholds above which the conditions (a)--(d), (g$'$) and (u) of the un-hatted system
hold. The constants $C_V$, $C_\sigma$ and $C_E$ are read by none of these choices.

\begin{lemma}[A maximum principle for non-negative inequality systems]
\label{9.11:lem-maximum-principle}
Let $I$ be a finite set and $I_0\subset I$. Let $D_i\ge 0$ and $E_i>0$ for $i\in I$, and let
$\sigma_i\ge 0$ and $M_{ij}\ge 0$ for $i\in I_0$, $j\in I$. Suppose that for every $i\in I_0$
\begin{equation}\label{9.11:eq-maximum-principle-1}
D_i\le\sum_{j\in I}M_{ij}D_j+\sigma_i,\qquad
\sum_{j\in I}M_{ij}E_j\le\tfrac12 E_i,\qquad
\sigma_i\le\tfrac12 E_i,
\end{equation}
and that $D_i\le E_i$ for every $i\in I\setminus I_0$. Then $D_i\le E_i$ for every $i\in I$.
\end{lemma}

\begin{proof}
Put
\begin{equation*}
t:=\max_{i\in I}\frac{D_i}{E_i}.
\end{equation*}
This is a maximum over finitely many real numbers, since $I$ is finite, and each quotient is
defined, since $E_i>0$ for every $i\in I$. By the definition of $t$ we have $D_j\le tE_j$ for
every $j\in I$. If $t\le 1$, then $D_i\le E_i$ for every $i\in I$, which is the claim.

Suppose, to reach a contradiction, that $t>1$, and let $i\in I$ be an index at which the
maximum is attained, so that $D_i=tE_i$. There are two cases.

If $i\notin I_0$, then the hypothesis $D_i\le E_i$ for indices in $I\setminus I_0$ gives
$tE_i\le E_i$. Dividing by $E_i>0$ gives $t\le 1$, which contradicts $t>1$.

If $i\in I_0$, we use the three inequalities \eqref{9.11:eq-maximum-principle-1} at the index
$i$. The first of them bounds $D_i$. In the sum it contains, each coefficient satisfies
$M_{ij}\ge 0$ and each unknown satisfies $D_j\le tE_j$, so each term satisfies
$M_{ij}D_j\le tM_{ij}E_j$. The second inequality bounds the resulting sum, and the third bounds
$\sigma_i$. Together,
\begin{equation*}
tE_i=D_i\le\sum_{j\in I}M_{ij}D_j+\sigma_i
\le t\sum_{j\in I}M_{ij}E_j+\sigma_i
\le\tfrac12 tE_i+\tfrac12 E_i .
\end{equation*}
Subtracting $\tfrac12 tE_i$ from both ends gives $\tfrac12 tE_i\le\tfrac12 E_i$. Dividing by
$\tfrac12 E_i>0$ gives $t\le 1$, which again contradicts $t>1$.

Both cases lead to a contradiction, hence $t\le 1$, and $D_i\le tE_i\le E_i$ for every $i\in I$.
\end{proof}

No triangularity is required of the coefficients $M_{ij}$. They are not asked to vanish for
$j$ above or below $i$ in any order on $I$, and the diagonal coefficients $M_{ii}$,
$i\in I_0$, may be positive. The proof uses only the non-negativity of the $M_{ij}$ and the
row bound $\sum_{j}M_{ij}E_j\le\tfrac12 E_i$.

This matters for the systems of inequalities to which the lemma is applied in this chapter.
There, the inequality for one block at a given depth involves the unknowns of the other blocks
at that same depth. One of the inequalities also involves its own unknown: it does so through
the coupling distance at the next depth, which is one of the terms of the weighted two-level
distance. Consequently no order on the index set makes the coefficient matrix triangular, and
the lemma does not need one.

\begin{definition}[The chain of competitors from the finest level]\label{9.11:def-competitor-chain}
Let $\mathsf{A}$ and $\mathsf{B}$ be the two runs of a pair, of lengths $n$ and $n+1$ respectively,
as in Definition~\ref{9.11:def-competitors-distances}. Write $\mathfrak{T}$ for the cross-cutoff
transplant and $C_V$ for its norm constant, both from that definition. Write $\hat{F}_r^{\mathsf{A}}$
and $\hat{F}_r^{\mathsf{B}}$ for the hatted outputs of the two runs at depth $r$, and
$\hat{M}_r^{\bullet}$ for the hatted production difference of a block $\bullet$ at depth $r$, both
read with the hatted blocks of Definition~\ref{9.11:def-hatted-blocks}. The
hatted reading units $\hat{\mathfrak{u}}_r^c$ are those of \eqref{9.11:eq-hatted-blocks-b}. For a
family $G$ and a component $c$, $G^c$ denotes its component $c$, and $\|\cdot\|^c$ denotes the family
norm read on the component $c$; on a single component we also write $\|G^c\|$ for $\|G^c\|^c$. The
construction follows that of Definition~\ref{9.6:def-comparison-chain}, with the hatted blocks, the
hatted reading units and the hatted source families in place of the un-hatted ones.

\emph{Initial value.} For every component $c$ put
\begin{equation}\label{9.11:eq-competitor-chain-1}
  \hat{G}_n^c := \bigl(\mathfrak{T}\hat{F}_n^{\mathsf{B}}\bigr)^c .
\end{equation}

\emph{Recursion.} Let $r_0$ be a depth down to which the pair is alive and matched, in the sense of
the hypotheses of Proposition~\ref{9.11:prop-source-families}. We define
the remaining entries depth by depth, for $r = n-1, n-2, \dots, r_0$ in this order. At a fixed depth
$r$ we treat the blocks $\bullet$ one after another along the order $\prec$ of sub-steps. When the
block $\bullet$ at depth $r$ is reached, the entries already defined are those at the depths $s$ with
$r < s \le n$, together with those at depth $r$ that belong to blocks preceding $\bullet$ in $\prec$.
From these entries the hatted hybrid $\hat{\mathfrak{s}}^{\natural}_{r,\bullet}$ of
Definition~\ref{9.11:def-competitors-distances} is built. Let
$\hat{N}_r^{\bullet}$ be the family that Proposition~\ref{9.11:prop-source-families} supplies at this
hybrid (Proposition~\ref{9.11:prop-source-families} is stated; its proof is incomplete at the step
marked $(\ast)$). For each component $c$
of $\bullet$ we then put
\begin{equation}\label{9.11:eq-competitor-chain-a}
\begin{gathered}
  \hat{G}_r^c :=
  \begin{cases}
    \bigl(\hat{N}_r^{\bullet} + \hat{M}_r^{\bullet}\bigr)^c
      & \text{if } \bigl\|\bigl(\hat{N}_r^{\bullet} + \hat{M}_r^{\bullet}\bigr)^c\bigr\|^c
        \le (C_V+1)\,\hat{\mathfrak{u}}_r^c,\\[2pt]
    \bigl(\mathfrak{T}\hat{F}_r^{\mathsf{B}} - \hat{F}_r^{\mathsf{A}}\bigr)^c
      & \text{otherwise;}
  \end{cases}\\[4pt]
  \hat{g}_r^c := \bigl\|\hat{G}_r^c\bigr\|^c, \qquad
  \hat{g}_r^{\bullet} := \max_{c\in\bullet}\hat{g}_r^c;\\[4pt]
  \bigl\|\hat{M}_r^{\bullet}\bigr\| := \max_{c\in\bullet}\bigl\|\hat{M}_r^c\bigr\|, \qquad
  \bigl\|\hat{N}_r^{\bullet}\bigr\| := \max_{c\in\bullet}\bigl\|\hat{N}_r^c\bigr\|.
\end{gathered}
\end{equation}

Thus each entry $\hat{G}_r^c$ is fixed, given the hatted data at depth $r$ and the family
supplied by Proposition~\ref{9.11:prop-source-families}, by the initial value
\eqref{9.11:eq-competitor-chain-1} and the entries defined earlier in the recursion. The sequence
$(\hat{G}_r^c)$, with
$r_0 \le r \le n$ and $c$ running over all components, is called the \emph{chain of competitors from
the finest level}. The numbers $\hat{g}_r^c$ and $\hat{g}_r^{\bullet}$ are its component norms and its
block norms.
\end{definition}

\begin{lemma}[Well-definedness of the chain and competitor bounds]\label{9.11:lem-chain-bounds}
Let $\mathsf{A}$, of length $n$, and $\mathsf{B}$, of length $n+1$, be the two runs of a pair as in
Definition~\ref{9.11:def-competitors-distances}, parametrised over $\mathfrak{P}$, and let
$\hat{G}_r^c$, $\hat{g}_r^c$, $\hat{g}_r^{\bullet}$ for $r_0\le r\le n$, and $\hat{M}_r^{\bullet}$,
$\hat{N}_r^{\bullet}$ and the hybrids $\hat{\mathfrak{s}}^{\natural}_{r,\bullet}$ for $r_0\le r<n$, be
the entries of the chain of competitors from the finest level,
Definition~\ref{9.11:def-competitor-chain}, with the rule \eqref{9.11:eq-competitor-chain-a}. Assume:
\begin{itemize}
\item the depth-$s$ slots of both runs lie in the hatted class $\hat{\mathfrak{K}}$ of
Definition~\ref{9.11:def-hatted-blocks}, with its reading units $\hat{\mathfrak{u}}_s^c$, the
$\hat{D}$-components obeying the defect clause \eqref{9.11:eq-hatted-blocks-b};
\item the transplant clauses $(\mathfrak{T}1)$--$(\mathfrak{T}3)$ of
\eqref{9.11:eq-competitors-distances-b}, on hatted families;
\item the closure clause for hatted classes, \eqref{9.11:eq-competitors-distances-c};
\item the conclusion \eqref{9.11:eq-source-families-a} of
Proposition~\ref{9.11:prop-source-families} (hatted source families with geometric envelope),
which is stated with its proof incomplete at the step marked $(\ast)$.
\end{itemize}
Then the chain is well defined, and for every component $c$ and every block $\bullet$
\begin{equation}\label{9.11:eq-chain-bounds-a}
\begin{aligned}
\text{(c1)}&\quad \hat{G}_r^c\approx\varphi_r^{\mathsf{B},c}-\varphi_r^{\mathsf{A},c}
&&(r_0\le r\le n),\\
\text{(c2)}&\quad \hat{g}_r^c\le (C_V+1)\,\hat{\mathfrak{u}}_r^c &&(r_0\le r\le n),\\
\text{(c3)}&\quad \hat{\mathsf{D}}_r^c\le \hat{g}_r^c\le \|\hat{M}_r^c\|+\|\hat{N}_r^c\|
&&(r_0\le r<n),\\
&\quad \hat{g}_r^{\bullet}\le \|\hat{M}_r^{\bullet}\|+\|\hat{N}_r^{\bullet}\| &&(r_0\le r<n),\\
\text{(boundary)}&\quad \hat{g}_n^c\le C_V\,\hat{\mathfrak{u}}_n^c ,
\end{aligned}
\end{equation}
in the following sense.
\begin{itemize}
\item[(c1)] For $r_0\le r\le n$, $\hat{G}_r^c$ is a list-structured $\mathsf{A}$-family on $c$,
holomorphic on $\mathfrak{P}$, whose local members read $W$ on $Y^{+}$, and the relation marked (c1)
in \eqref{9.11:eq-chain-bounds-a} holds.
\item[(c2)] For $r_0\le r\le n$ the bound marked (c2) in \eqref{9.11:eq-chain-bounds-a} holds; in
consequence every hybrid $\hat{\mathfrak{s}}^{\natural}_{r,\bullet}$ lies in
$C_{\mathfrak{K}}\hat{\mathfrak{K}}^c$.
\item[(c3)] For $r_0\le r<n$ the three bounds marked (c3) in \eqref{9.11:eq-chain-bounds-a} hold, the
last being the block form.
\end{itemize}
At the boundary depth $n$ the bound marked (boundary) in \eqref{9.11:eq-chain-bounds-a} holds.
\end{lemma}

\begin{proof}
We argue by downward induction along the entries $(r,\bullet)$, ordered by decreasing $r$ and, at
fixed $r$, by the order $\prec$ of the blocks, which is the order in which
Definition~\ref{9.11:def-competitor-chain} defines the entries.

\emph{The depth $r=n$.} By Definition~\ref{9.11:def-competitor-chain},
$\hat{G}_n^c=(\mathfrak{T}\hat{F}_n^{\mathsf{B}})^c$. By $(\mathfrak{T}1)$ this is an
$\mathsf{A}$-family, and
\begin{equation*}
\|\hat{G}_n^c\|^c\le C_V\,\|\hat{F}_n^{\mathsf{B},c}\|^{\flat}\le C_V\,\hat{\mathfrak{u}}_n^c ,
\end{equation*}
where the first inequality is $(\mathfrak{T}1)$ and the second is the bound of the class
$\hat{\mathfrak{K}}$ obeyed by the depth-$n$ slots of $\mathsf{B}$ at the common ladder value
$\theta_n$, the $\hat{D}$-component being covered by the defect clause
\eqref{9.11:eq-hatted-blocks-b}. By $(\mathfrak{T}2)$ and the convention
$\varphi_n^{\mathsf{A},c}:=0$ of Definition~\ref{9.11:def-competitors-distances},
\begin{equation*}
\hat{G}_n^c\approx\varphi_n^{\mathsf{B},c}=\varphi_n^{\mathsf{B},c}-\varphi_n^{\mathsf{A},c}.
\end{equation*}
This is the boundary clause, the total in (c1), and, the reading units being non-negative, (c2) at
$r=n$. The remaining properties in (c1) at $r=n$ are those of a transplanted family given by
$(\mathfrak{T}1)$--$(\mathfrak{T}3)$.

\emph{The induction step: well-definedness.} Let $r<n$ and suppose that (c1) and (c2) hold for every
entry earlier than $(r,\bullet)$. The hybrid $\hat{\mathfrak{s}}^{\natural}_{r,\bullet}$ is then a
hybrid in the sense of Definition~\ref{9.11:def-competitors-distances}, built from entries already
defined, whose slots obey
\begin{equation*}
\|\hat{G}_s^c\|\le (C_V+1)\,\hat{\mathfrak{u}}_s^c
\end{equation*}
by the induction hypothesis (c2). The closure clause \eqref{9.11:eq-competitors-distances-c} applies
to it, so $\hat{\mathfrak{s}}^{\natural}_{r,\bullet}$ lies in $C_{\mathfrak{K}}\hat{\mathfrak{K}}^c$.
Hence the production difference $\hat{M}_r^{\bullet}$ is defined, and the family $\hat{N}_r^{\bullet}$
at this hybrid exists by the hypothesis \eqref{9.11:eq-source-families-a}, the conclusion of
Proposition~\ref{9.11:prop-source-families}, which is stated with its proof incomplete at the step
marked $(\ast)$.

\emph{The first alternative competes.} Fed with $\mathsf{A}$'s own slots, $\mathsf{A}$'s hatted
production returns $\hat{F}_r^{\mathsf{A},\bullet}$, whose component $c$ has total
$\varphi_r^{\mathsf{A},c}$. Since totals add component by component, the production difference
satisfies
\begin{equation*}
\hat{M}_r^c\approx
\varphi^c\bigl[\hat{\mathsf{F}}_r^{\mathsf{A},+,\bullet}(\hat{\mathfrak{s}}^{\natural}_{r,\bullet})\bigr]
-\varphi_r^{\mathsf{A},c}.
\end{equation*}
By the assumed relation \eqref{9.11:eq-source-families-a},
\begin{equation*}
\hat{N}_r^c\approx\varphi_r^{\mathsf{B},c}
-\varphi^c\bigl[\hat{\mathsf{F}}_r^{\mathsf{A},+,\bullet}(\hat{\mathfrak{s}}^{\natural}_{r,\bullet})\bigr].
\end{equation*}
Adding the two relations, the middle totals cancel and
\begin{equation*}
(\hat{N}_r^{\bullet}+\hat{M}_r^{\bullet})^c\approx\varphi_r^{\mathsf{B},c}-\varphi_r^{\mathsf{A},c},
\end{equation*}
so the first alternative of \eqref{9.11:eq-competitor-chain-a} is a competitor at depth $r$,
component $c$.

\emph{The second alternative competes.} By $(\mathfrak{T}2)$,
$\mathfrak{T}\hat{F}_r^{\mathsf{B}}\approx\varphi_r^{\mathsf{B}}$, and
$\hat{F}_r^{\mathsf{A}}\approx\varphi_r^{\mathsf{A}}$, so that
$(\mathfrak{T}\hat{F}_r^{\mathsf{B}}-\hat{F}_r^{\mathsf{A}})^c\approx
\varphi_r^{\mathsf{B},c}-\varphi_r^{\mathsf{A},c}$; by $(\mathfrak{T}1)$ the difference is an
$\mathsf{A}$-family. The depth-$r$ slots of both runs obey the bounds of $\hat{\mathfrak{K}}$ at the
common value $\theta_r$ (the $\hat{D}$-components by \eqref{9.11:eq-hatted-blocks-b}). Hence, by the
triangle inequality and $(\mathfrak{T}1)$,
\begin{equation*}
\|(\mathfrak{T}\hat{F}_r^{\mathsf{B}}-\hat{F}_r^{\mathsf{A}})^c\|^c
\le C_V\,\hat{\mathfrak{u}}_r^c+\hat{\mathfrak{u}}_r^c=(C_V+1)\,\hat{\mathfrak{u}}_r^c .
\end{equation*}

\emph{List structure, holomorphy and read-sets.} These are required of either alternative and of any
mixture of the two alternatives over the components $c\in\bullet$. The relations of the list
structure that tie the slot $P^{\flat}$ to $\mathsf{P}$ and fix $\hat{b}=\beta[\hat{E}]$ are
statements about totals, and both alternatives have the same totals, so they hold for both. The
$\hat{D}$-components vanish at constant sources in both alternatives: the production freezes at
constant sources, the family $\hat{N}_r^{\bullet}$ is list-structured by the hypothesis
\eqref{9.11:eq-source-families-a}, and
$\mathfrak{T}$ is linear and passes the parameters. The frozen and the local members of the
$R$-component are taken from one alternative. Holomorphy on $\mathfrak{P}$ and the read-sets, the
local members reading $W$ on $Y^{+}$, hold inside each alternative. Finally, the choice between the
alternatives is one dichotomy for each pair, each depth $r$ and each component $c$, made on the list
norm, which is a supremum over $\overline{\mathfrak{P}}$, over sources and over tori; the choice
therefore does not depend on the parameter in $\mathfrak{P}$, and the chosen family is holomorphic on
$\mathfrak{P}$. This proves (c1) at $(r,c)$ for every $c\in\bullet$, from
$(\mathfrak{T}1)$--$(\mathfrak{T}3)$, the closure clause \eqref{9.11:eq-competitors-distances-c}, the
defect clause \eqref{9.11:eq-hatted-blocks-b} and the hypothesis
\eqref{9.11:eq-source-families-a}.

\emph{Proof of} (c2). In the first alternative of \eqref{9.11:eq-competitor-chain-a} the bound
$\hat{g}_r^c\le(C_V+1)\hat{\mathfrak{u}}_r^c$ is the condition under which that alternative is
chosen; in the second alternative it is the bound just shown. Thus (c1) and (c2) hold at
$(r,\bullet)$, which closes the induction; in particular every entry of the chain is defined, and the
statement about the hybrids in (c2) was obtained in the induction step.

\emph{Proof of} (c3). Since $\hat{G}_r^c$ is a competitor at depth $r$, component $c$, by (c1), and
$\hat{\mathsf{D}}_r^c$ is the infimum of the norms of competitors
(Definition~\ref{9.11:def-competitors-distances}), $\hat{\mathsf{D}}_r^c\le\hat{g}_r^c$. In the first
alternative, by the triangle inequality,
\begin{equation*}
\hat{g}_r^c=\|(\hat{N}_r^{\bullet}+\hat{M}_r^{\bullet})^c\|^c\le\|\hat{M}_r^c\|+\|\hat{N}_r^c\| .
\end{equation*}
In the second alternative the first was rejected, that is
$\|(\hat{N}_r^{\bullet}+\hat{M}_r^{\bullet})^c\|^c>(C_V+1)\hat{\mathfrak{u}}_r^c$, so that, with
the bound of the second alternative and the triangle inequality,
\begin{equation*}
\hat{g}_r^c\le(C_V+1)\,\hat{\mathfrak{u}}_r^c
<\|(\hat{N}_r^{\bullet}+\hat{M}_r^{\bullet})^c\|^c
\le\|\hat{M}_r^c\|+\|\hat{N}_r^c\| .
\end{equation*}
Taking the maximum over $c\in\bullet$, and using that $\hat{g}_r^{\bullet}$,
$\|\hat{M}_r^{\bullet}\|$ and $\|\hat{N}_r^{\bullet}\|$ are the maxima over $c\in\bullet$ of the
corresponding component quantities (Definition~\ref{9.11:def-competitor-chain}),
\begin{equation*}
\hat{g}_r^{\bullet}=\max_{c\in\bullet}\hat{g}_r^c
\le\max_{c\in\bullet}\bigl(\|\hat{M}_r^c\|+\|\hat{N}_r^c\|\bigr)
\le\|\hat{M}_r^{\bullet}\|+\|\hat{N}_r^{\bullet}\| ,
\end{equation*}
which is the block form of (c3).
\end{proof}

\begin{theorem}[Geometric envelope for two-run differences of hatted runs. Stated; proof incomplete at the step marked $(\ast)$]
\label{9.11:thm-geometric-envelope}
Let $L\ge L_0$, and let $M$, $\kappa$, $\kappa_0$, $p_0$, $q_{\flat}$, $r_{\mathrm{pr}}$ be given.
Let the following constants and integers also be given, all kept symbolic:
\begin{itemize}
\item the constant $\mathsf{A}$ of hypothesis (J1) of the uniqueness theorem;
\item the constants $C_t$, $\Lambda$, $c_3$ of \eqref{9.11:eq-rate-floors-c};
\item the norm constant $C_V$ of the transplant $\mathfrak{T}$;
\item the constant $C_\sigma$ of Proposition~\ref{9.11:prop-source-families};
\item the constants $\alpha_r$, $\alpha_r^A$, $\varkappa$, $\varkappa''_r$, $\varkappa_A$, $\varkappa'_r$, $C_F$,
$C_F^A$ and the integers $\mathsf{a}_s\ge 1$, $\ell_r\ge 2$ of
Definitions~\ref{9.11:def-response-rows} and~\ref{9.11:def-coupling-ceilings}.
\end{itemize}
Let $q$, $q_a$ and $\rho$ satisfy the floors \eqref{9.11:eq-rate-floors-a} and
\eqref{9.11:eq-rate-floors-b} of Definition~\ref{9.11:def-rate-floors}.

We put $e_r:=\rho^{n-r}$, where $n$ is the length of the run $\mathsf{A}$ below, and
$\hat{c}_\rho:=(1+\rho)/(\rho(1-\rho)^2)$. Let $C_E\ge\max\{C_V+1,\,2C_\sigma\}$.
Let $\theta_r$ be the points of the ladder $\theta_0:=X$, $\theta_{s+1}:=\theta_s+c_0$. As in
\eqref{9.11:eq-rate-floors-c}, put
\begin{gather*}
\hat{\mathfrak{r}}_r:=2\mathfrak{r}_r\hat{\Lambda}_r\bigl[\mathsf{A}(1-\rho)^{-2}+6(1-\rho)^{-1}+8\bigr],
\qquad \mathfrak{r}_r:=C_t\theta_r^{1-\kappa_0/2},\\
\hat{\Lambda}_r:=C_{**}\Lambda_r,\qquad \Lambda_r:=\Lambda(\log\theta_r)^{c_3},\qquad c_3:=3p_0+2,\\
C_{**}:=(1-q_a)^{-1}\sup_{m\ge 0}(m+2)q_a^{m/2}.
\end{gather*}

Assume:
\begin{itemize}
\item hypothesis (J1) of the uniqueness theorem;
\item parts (a) and (b) of Definition~\ref{9.11:def-response-rows}, that is, the rows
\eqref{9.11:eq-response-rows-a}, \eqref{9.11:eq-response-rows-b}, \eqref{9.11:eq-response-rows-c}
and the clause \eqref{9.11:eq-response-rows-d};
\item the source families of Proposition~\ref{9.11:prop-source-families}, whose proof is incomplete at
one step;
\item the hatted class $\hat{\mathfrak{K}}$ with its reading units $\hat{\mathfrak{u}}_s^c$
(Definition~\ref{9.11:def-hatted-blocks}), the first three transplant clauses of
\eqref{9.11:eq-competitors-distances-b}, the hatted closure clause of
\eqref{9.11:eq-competitors-distances-c}, and the defect clause \eqref{9.11:eq-hatted-blocks-b};
\item the conditions of Definition~\ref{9.11:def-coupling-ceilings};
\item the one-sided premises $c_0>0$, $g_{(s)}^2\le 1$ for every depth $s$, and $\Lambda\ge 1$.
\end{itemize}

Fix the following data:
\begin{itemize}
\item an integer $n\ge1$ and a depth $r_0$ with $0\le r_0\le n$;
\item a pair of hatted lists, $\mathsf{A}$ of length $n$ and $\mathsf{B}$ of length $n+1$ (frozen or
local), alive to depth $r_0$, with
$\hat{\zeta}^{\mathsf{B}}_{(r_0)}=\hat{\zeta}^{\mathsf{A}}_{(r_0)}$ on the parameter set
$\mathfrak{P}$ (complex deviations on the tube being allowed);
\item a depth $r$ with $r_0\le r\le n$.
\end{itemize}
For these data, the norms $\hat{g}_r^{\bullet}$ of the competitors of the chain of
Definition~\ref{9.11:def-competitor-chain} satisfy the first line of the following display. The
coupling distance $d_{r+1}$ and the weighted two-level distance $d^{\sharp}_{r+1}$ of the pair
satisfy its second line for $r_0\le r<n$:
\begin{equation}\label{9.11:eq-geometric-envelope-a}
\begin{gathered}
\hat{g}_r^P\le C_E\rho^{n-r},\qquad \hat{g}_r^A\le 4C_E\rho^{n-r},\qquad
\hat{g}_r^R\le 4C_E\hat{\mathfrak{r}}_r\rho^{n-r},\\
d_{r+1}\le\frac{\mathsf{A}C_E\rho^{n-r}}{1-\rho},\qquad
d^{\sharp}_{r+1}\le\hat{c}_\rho\mathsf{A}C_E\rho^{n-r}.
\end{gathered}
\end{equation}
Hence the first three bounds hold also for the distances $\hat{\mathsf{D}}_r^c$, $c$ a component of
the block $P$, $A$ or $R$ respectively. No quantity in these bounds depends on $n$ otherwise than
through the factor $\rho^{n-r}$.
\end{theorem}

Stated; the proof below is incomplete at the step marked $(\ast)$.

\begin{proof}[Proof of Theorem~\ref{9.11:thm-geometric-envelope}]
Throughout the proof $e_r=\rho^{n-r}$. Thus $e_n=1$, and for every integer $m$ we have exactly
\begin{equation*}
e_{r+m}=\rho^{n-r-m}=\rho^{-m}e_r.
\end{equation*}

By the floors \eqref{9.11:eq-rate-floors-a} and \eqref{9.11:eq-rate-floors-b} we have
$0<\rho<1$, $0<q\le\rho^2$ and $L^{-1/2}<q$. Moreover $\mathfrak{r}_r>0$ and
$\hat{\Lambda}_r\ge1$ (the premise $\Lambda\ge1$ enters here); these positivity facts are used in
the continuation of the proof.

All sums over chain depths $s$ run over $r_0\le s\le n$. For sums indexed by $s>r$ or $s\ge r$ this
means in addition $s\le n$. Likewise, the window sum $\sum_{r<s<r+\ell_r}$ and the band sum
$\Sigma^{\mathrm{band}}_{s>r}$ stop at $s=n$; the latter is the sum over $s>r$ with
$s-r\le\mathsf{a}_s$, as in Definition~\ref{9.11:def-response-rows}.

\emph{The system of inequalities.}
Let $r_0\le r<n$. By clause (c3) of Lemma~\ref{9.11:lem-chain-bounds} (see
\eqref{9.11:eq-chain-bounds-a}), for each block $\bullet\in\{P,A,R\}$,
\begin{equation*}
\hat{g}_r^{\bullet}\le\|\hat{M}_r^{\bullet}\|+\|\hat{N}_r^{\bullet}\|.
\end{equation*}

We bound the first term by the rows of part (a) of Definition~\ref{9.11:def-response-rows}, applied
at the hybrids $\hat{\mathfrak{s}}^{\natural}_{r,\bullet}$ of the chain. By
Definition~\ref{9.11:def-competitor-chain}, these hybrids are built from the entries already defined.
Their slots are therefore the competitors $\hat{G}_s$, and each norm $\|\hat{G}_s^{\bullet}\|$
occurring in the rows is $\hat{g}_s^{\bullet}$. The factors $g^2_{(s+1)}\le1$ in the kernel
$\hat{B}_r$ are kept symbolically and not replaced by $1$.

We bound the second term by Proposition~\ref{9.11:prop-source-families}, whose proof is incomplete at
one step. It is applied at the same hybrids, so that its kernel $\hat{B}_r$ and its window sum are
also read at the chain's competitors. For the $R$-block we use
\begin{equation*}
\|\hat{N}_r^R\|\le\|\hat{N}_r^{R'}\|+\|\hat{N}_r^{R''}\|.
\end{equation*}
We obtain, for $r_0\le r<n$:
\begin{equation}\label{9.11:eq-geometric-envelope-b}
\begin{split}
\hat{g}_r^P&\le\alpha_r d^{\sharp}_{r+1}
+\varkappa\sum_{s>r}q^{s-r}\bigl(\hat{g}_s^P+\hat{g}_s^R\bigr)
+\varkappa''_r\sum_{s>r}q^{s-r}\hat{g}_s^A\\
&\qquad+C_F\,\Sigma^{\mathrm{band}}_{s>r}\hat{g}_s^R+C_\sigma e_r;\\
\hat{g}_r^A&\le\alpha_r^A d^{\sharp}_{r+1}
+\varkappa_A\Bigl[\hat{g}_r^P
+\sum_{s>r}q^{s-r}\bigl(\hat{g}_s^P+\hat{g}_s^A+\hat{g}_s^R\bigr)\Bigr]\\
&\qquad+\hat{g}_r^R+C_F^A\,\Sigma^{\mathrm{band}}_{s>r}\hat{g}_s^R+C_\sigma e_r;\\
\hat{g}_r^R&\le 2\mathfrak{r}_r\hat{\Lambda}_r\hat{B}_r(\hat{g},d)
+2\varkappa'_r\sum_{r<s<r+\ell_r}\hat{g}_s^R
+\mathfrak{r}_r\hat{\Lambda}_r(2C_\sigma+1)e_r,\\
\hat{B}_r(\hat{g},d)&=\sum_{r\le s<n}q^{s-r}g^2_{(s+1)}d_{s+1}
+\sum_{r\le s\le n}q^{s-r}\bigl(\hat{g}_s^P+\hat{g}_s^A\bigr)
+\sum_{r<s\le n}\mathsf{k}_{r,s}\,\hat{g}_s^R.
\end{split}
\end{equation}
Here $\mathsf{k}_{r,s}$ is the window kernel of \eqref{9.11:eq-response-rows-c}.

\emph{The $P$- and $A$-inequalities.}
The first inequality of \eqref{9.11:eq-geometric-envelope-b} is the row
\eqref{9.11:eq-response-rows-a}, with $\|\hat{G}_s^{\bullet}\|=\hat{g}_s^{\bullet}$, plus the bound
$\|\hat{N}_r^P\|\le C_\sigma e_r$ of Proposition~\ref{9.11:prop-source-families}. The second is the
row \eqref{9.11:eq-response-rows-b} plus the bound $\|\hat{N}_r^A\|\le C_\sigma e_r$ of the same
proposition.

\emph{The $R$-inequality.}
By the row \eqref{9.11:eq-response-rows-c},
\begin{equation*}
\|\hat{M}_r^R\|\le\mathfrak{r}_r\hat{\Lambda}_r\hat{B}_r(\hat{g},d)
+\varkappa'_r\sum_{r<s<r+\ell_r}\hat{g}_s^R.
\end{equation*}
Proposition~\ref{9.11:prop-source-families} gives $\|\hat{N}_r^{R''}\|\le\mathfrak{r}_r\hat{\Lambda}_r C_\sigma e_r$
and
\begin{equation*}
\|\hat{N}_r^{R'}\|\le\mathfrak{r}_r\hat{\Lambda}_r\bigl[(C_\sigma+1)e_r+\hat{B}_r(\hat{g},d)\bigr]
+\varkappa'_r\sum_{r<s<r+\ell_r}\hat{g}_s^R.
\end{equation*}
Together with the bound on $\|\hat{N}_r^R\|$ displayed above this gives
\begin{align*}
\|\hat{N}_r^R\|
&\le\mathfrak{r}_r\hat{\Lambda}_r C_\sigma e_r
+\mathfrak{r}_r\hat{\Lambda}_r\bigl[(C_\sigma+1)e_r+\hat{B}_r(\hat{g},d)\bigr]\\
&\qquad+\varkappa'_r\sum_{r<s<r+\ell_r}\hat{g}_s^R.
\end{align*}

Adding the two bounds and using $\hat{g}_r^R\le\|\hat{M}_r^R\|+\|\hat{N}_r^R\|$ we obtain
\begin{equation*}
\hat{g}_r^R\le2\mathfrak{r}_r\hat{\Lambda}_r\hat{B}_r(\hat{g},d)
+2\varkappa'_r\sum_{r<s<r+\ell_r}\hat{g}_s^R
+\mathfrak{r}_r\hat{\Lambda}_r(2C_\sigma+1)e_r,
\end{equation*}
which is the third inequality of \eqref{9.11:eq-geometric-envelope-b}. The factor $2$ in front of
the kernel and of the window sum arises because the row \eqref{9.11:eq-response-rows-c} and the
$R'$-bound of Proposition~\ref{9.11:prop-source-families} each carry the bracket
$\mathfrak{r}_r\hat{\Lambda}_r\hat{B}_r(\hat{g},d)$ and the window sum. The coefficient
$2C_\sigma+1$ is the sum of the two source contributions: $C_\sigma$ from the $R''$-component and
$C_\sigma+1$ from the $R'$-component, each multiplied by $\mathfrak{r}_r\hat{\Lambda}_r e_r$.

$(\ast)$ The remaining step is the deduction of the bounds \eqref{9.11:eq-geometric-envelope-a}
from the system \eqref{9.11:eq-geometric-envelope-b}; the proof is incomplete at this step.
\end{proof}

\begin{proof}[Proof of Theorem~\ref{9.11:thm-geometric-envelope}, continued: reduction of the coupling-distance line to the maximum principle]\label{9.11:prf-distance-reduction}
In this part of the proof we bound the coupling distances by the norms of the competitors at the $P$-block. We then bring the three inequalities
\eqref{9.11:eq-geometric-envelope-b} into the form of the maximum principle for non-negative inequality systems (Lemma~\ref{9.11:lem-maximum-principle}).

We write $\delta(\cdot):=(\cdot)^{\mathsf{B}}-(\cdot)^{\mathsf{A}}$ for the difference between the run $\mathsf{B}$ and the run $\mathsf{A}$ of the pair. A superscript $\mathsf{A}$ denotes the run. The letter $\mathsf{A}$ on the line (not as a superscript) denotes the constant of the coefficient bound
$|\beta_r[F]|\le\mathsf{A}\|F\|$ of the uniqueness theorem.

\emph{The coupling-distance line.} Both runs subtract $c_0$ on the same ladder. For $\mathsf{m}\in\{\mathsf{A},\mathsf{B}\}$ and every depth $u$,
\begin{equation*}
\hat{\zeta}^{\mathsf{m}}_{(u)}=\hat{\zeta}^{\mathsf{m}}_{(u+1)}+\hat{b}^{\mathsf{m}}_{(u)}-c_0 .
\end{equation*}
Solve each identity for $\hat{\zeta}^{\mathsf{m}}_{(u+1)}$ and subtract the identity for $\mathsf{A}$ from the one for $\mathsf{B}$. The constant $c_0$ is the same for both runs and cancels. Hence, pointwise on $\overline{\mathfrak{P}}$,
\begin{equation*}
\delta\hat{\zeta}_{(u+1)}=\delta\hat{\zeta}_{(u)}-\delta\hat{b}_{(u)} .
\end{equation*}
By the matching of the two runs at depth $r_0$, $\delta\hat{\zeta}_{(r_0)}=0$. Let $r_0\le s\le n-1$. Summing the last identity over $r_0\le u\le s$ telescopes to
\begin{equation}\label{9.11:eq-distance-reduction-1}
\delta\hat{\zeta}_{(s+1)}=\delta\hat{\zeta}_{(r_0)}-\sum_{r_0\le u\le s}\delta\hat{b}_{(u)}
=-\sum_{r_0\le u\le s}\delta\hat{b}_{(u)} .
\end{equation}

Now take absolute values in \eqref{9.11:eq-distance-reduction-1}, apply the triangle inequality, and take the supremum over $\overline{\mathfrak{P}}$. Each term $|\delta\hat{b}_{(u)}|$ is bounded using the hatted response row \eqref{9.11:eq-response-rows-d} of Definition~\ref{9.11:def-response-rows}, together with the coefficient bound of the uniqueness theorem. This gives
\begin{equation*}
|\delta\hat{b}_{(u)}|\le\mathsf{A}\|\hat{G}_u^{P}\|=\mathsf{A}\hat{g}_u^{P} .
\end{equation*}
Here $\hat{G}_u^{P}$ is the competitor at depth $u$ of the chain of competitors from the finest level, and $\hat{g}_u^{P}$ is its norm, as in \eqref{9.11:eq-competitor-chain-a}.

The coupling distance $d_{s+1}$ between the two runs, fixed in \eqref{9.11:eq-competitors-distances-a}, is bounded by the supremum over $\overline{\mathfrak{P}}$ of $|\delta\hat{\zeta}_{(s+1)}|$; this yields the first line of the display below. Its remaining lines define the quantities used in the rest of the argument.
\begin{equation}\label{9.11:eq-distance-reduction-a}
\begin{gathered}
d_{s+1}\le\sum_{r_0\le u\le s}\sup_{\overline{\mathfrak{P}}}|\delta\hat{b}_{(u)}|
\le\mathsf{A}\sum_{r_0\le u\le s}\hat{g}_u^{P}\qquad(r_0\le s<n),\\
d(\hat{g})_t:=\mathsf{A}\sum_{r_0\le u\le t-1}\hat{g}_u^{P}\quad(r_0<t\le n),
\qquad d(\hat{g})_t:=0\quad(t>n),\\
d(E)_t:=\mathsf{A}\sum_{r_0\le u\le t-1}E_u^{P}\quad(r_0<t\le n),
\qquad d(E)_t:=0\quad(t>n),\\
I:=\{(r,\bullet):r_0\le r\le n,\ \bullet\in\{P,A,R\}\},\qquad
I_0:=\{(r,\bullet)\in I:r<n\},\\
\sigma_r^{P}=\sigma_r^{A}:=C_\sigma e_r,\qquad
\sigma_r^{R}:=\mathfrak{r}_r\hat{\Lambda}_r(2C_\sigma+1)e_r,\\
M=(M_{ij})_{i\in I_0,\ j\in I},\qquad
M_{ij}:=\text{the coefficient of }D_j\text{ in the substituted inequality of index }i .
\end{gathered}
\end{equation}
In the third line, $E=(E_j)_{j\in I}$ is a vector indexed by $I$, and we write $E_u^{\bullet}:=E_{(u,\bullet)}$. In the last line, the vector $D$ and the substituted inequalities are those introduced below, where $M_{ij}\ge0$ is established.

\emph{Comparison with the majorant.} By definition, $d_{r_0}=0$ and $d_t=0$ for $t>n$. Let $t>r_0$. If $t\le n$, put $s:=t-1$. Then $r_0\le s<n$, and the first line of \eqref{9.11:eq-distance-reduction-a} reads $d_t\le d(\hat{g})_t$. If $t>n$, both sides vanish. Hence $d_t\le d(\hat{g})_t$ for every $t>r_0$ occurring in the rows.

\emph{Substitution in the rows.} Every coupling distance occurs in the three inequalities \eqref{9.11:eq-geometric-envelope-b} with a non-negative coefficient. In the $P$- and $A$-inequalities of \eqref{9.11:eq-geometric-envelope-b} the coefficients are $\alpha_rL^{-i/2}$ and $\alpha_r^AL^{-i/2}$. These arise inside the weighted distance $d^{\sharp}_{r+1}$ of \eqref{9.11:eq-competitors-distances-a}. In the $R$-inequality of \eqref{9.11:eq-geometric-envelope-b} the coefficient is $2\mathfrak{r}_r\hat{\Lambda}_rq^{s-r}g^2_{(s+1)}$, inside the first member of $\hat{B}_r$.

Replacing each $d_t$ by the larger number $d(\hat{g})_t$ therefore does not decrease any right member. So the three inequalities remain true with $d$ replaced by $d(\hat{g})$.

Consider the finite vector
\begin{equation*}
D:=(\hat{g}_u^{\bullet})_{r_0\le u\le n,\ \bullet\in\{P,A,R\}},
\qquad D_{(u,\bullet)}:=\hat{g}_u^{\bullet}\ge0 .
\end{equation*}
After the substitution, each right member is a linear form in $D$ with non-negative coefficients, plus a source. This holds since each right member of \eqref{9.11:eq-geometric-envelope-b} is linear with non-negative coefficients in the norms $\hat{g}_s^{\bullet}$ and in the distances $d_t$, and each $d(\hat{g})_t$ is itself a linear form in $D$ with non-negative coefficients. The weighted distance $d^{\sharp}$ is a series over $i\ge0$. After the substitution, however, only the finitely many terms whose distance index is at most $n$ are non-zero, since $d(\hat{g})_t=0$ for $t>n$.

For $i=(r,\bullet)\in I_0$ and $j\in I$, the number $M_{ij}$ of the last line of \eqref{9.11:eq-distance-reduction-a} is the coefficient of $D_j$ in this linear form for the substituted inequality of \eqref{9.11:eq-geometric-envelope-b} for the block $\bullet$ at depth $r$; thus $M_{ij}\ge0$. Let $\sigma_i$ be the source of that inequality, as given in \eqref{9.11:eq-distance-reduction-a}. Then for every $i\in I_0$,
\begin{equation*}
D_i\le\sum_{j\in I}M_{ij}D_j+\sigma_i .
\end{equation*}
This is the form of the hypothesis of Lemma~\ref{9.11:lem-maximum-principle}, on the index sets $I$ and $I_0$ of \eqref{9.11:eq-distance-reduction-a}, with the matrix $M\ge0$ and the sources $\sigma$.

\emph{The homogeneous part at a vector $E$.} The map $\hat{g}\mapsto d(\hat{g})$ is linear, and its value at $\hat{g}=E$ is $d(E)$. Hence, by linearity, the following holds for $i=(r,\bullet)\in I_0$. The number $\sum_{j\in I}M_{ij}E_j$ equals the homogeneous part of the substituted right member of the inequality of \eqref{9.11:eq-geometric-envelope-b} for the block $\bullet$ at depth $r$, evaluated at $\hat{g}=E$ and $d=d(E)$, with $d(E)$ as in \eqref{9.11:eq-distance-reduction-a}.
\end{proof}

\begin{lemma}[The envelope and its geometric, band and window sums]\label{9.11:prf-envelope-sums}
Let $n$ be the length of the shorter run of the pair compared in
Definition~\ref{9.6:def-comparison-chain}, and let $r_0$ be the depth at which the two runs are
matched, as in \eqref{9.11:eq-distance-reduction-a}. Let $q$ and $\rho$ be the kernel rate and the
contraction rate of Definition~\ref{9.11:def-rate-floors}. Let $\hat{\mathfrak{r}}_s$, the sum
$\mathfrak{S}:=\sum_{s\ge0}\hat{\mathfrak{r}}_s\rho^{-\mathsf{a}_s}$ and the number
$\hat{c}_\rho:=(1+\rho)/(\rho(1-\rho)^2)$ be as in
Definition~\ref{9.11:def-coupling-ceilings}. Let $\mathsf{a}_s\ge1$ be the integers that define the band
and the unit window in Remark~\ref{9.11:rem-memory-structures}. For every depth $r$ put
$e_r:=\rho^{\,n-r}$, and define the \emph{envelope}
\begin{equation}\label{9.11:eq-envelope-sums-a}
E_r^P:=C_Ee_r,\qquad E_r^A:=4C_Ee_r,\qquad E_r^R:=4C_E\hat{\mathfrak{r}}_re_r ,
\end{equation}
where $C_E>0$ is the envelope constant and $\hat{\mathfrak{r}}_r>0$; all three numbers are positive. Let
$\mathsf{k}_{r,s}:=\rho^{2(s-r-\mathsf{a}_s-1)_+}$ be the window kernel, let $d(E)_t$ be the coupling
distance of \eqref{9.11:eq-distance-reduction-a} evaluated at the envelope, with $\mathsf{A}$ the constant
of that display, let $g_{(s)}$ be the running coupling at depth $s$, which satisfies $g_{(s)}^2\le1$,
and put
\begin{equation*}
d^{\sharp}(E)_{r+1}:=\sum_{i\ge0}L^{-i/2}\bigl(d(E)_{r+1+i}+d(E)_{r+2+i}\bigr).
\end{equation*}
Then, for every depth $r$ with $r_0\le r\le n$, the following hold.
\begin{equation}\label{9.11:eq-envelope-sums-b}
\begin{aligned}
&\sum_{r\le s\le n}q^{s-r}e_s\le\frac{e_r}{1-\rho},\qquad
\sum_{r<s\le n}q^{s-r}e_s\le\frac{\rho\,e_r}{1-\rho},\qquad
\sum_{r_0\le u\le r}e_u\le\frac{e_r}{1-\rho};\\
&d(E)_t\le\frac{\mathsf{A}C_Ee_{t-1}}{1-\rho}\quad\text{for }r_0<t\le n,
\qquad d(E)_t=0\quad\text{for }t>n .
\end{aligned}
\end{equation}
For $r<s\le n$ with $s-r\le\mathsf{a}_s$ one has $e_s\le\rho^{-\mathsf{a}_s}e_r$, and the band sum, taken over
the depths $s\le n$, satisfies
\begin{equation}\label{9.11:eq-envelope-sums-c}
\Sigma^{\mathrm{band}}_{s>r}\,\hat{\mathfrak{r}}_se_s\le\mathfrak{S}\,e_r .
\end{equation}
For every $s$ with $r<s\le n$,
\begin{equation}\label{9.11:eq-envelope-sums-d}
\mathsf{k}_{r,s}e_s\le\rho^{-1}\rho^{-\mathsf{a}_s}e_r,
\qquad\text{and hence}\qquad
\sum_{r<s\le n}\mathsf{k}_{r,s}\hat{\mathfrak{r}}_se_s\le\rho^{-1}\mathfrak{S}\,e_r .
\end{equation}
Finally,
\begin{equation}\label{9.11:eq-envelope-sums-e}
\sum_{r\le s<n}q^{s-r}g_{(s+1)}^2\,d(E)_{s+1}\le\frac{\mathsf{A}C_Ee_r}{(1-\rho)^2},
\qquad
d^{\sharp}(E)_{r+1}\le\hat{c}_\rho\,\mathsf{A}C_Ee_r .
\end{equation}
\end{lemma}

\begin{proof}
By Definition~\ref{9.11:def-rate-floors}, $L^{-1/2}<q\le\rho^2$ and $\rho<1$; in particular
$q>0$, and $\rho\ge q^{1/2}>0$, so $0<\rho<1$. For depths $r\le s$ we have
\begin{equation*}
e_s=\rho^{\,n-s}=\rho^{-(s-r)}\rho^{\,n-r}=\rho^{-(s-r)}e_r ,
\end{equation*}
and this identity is used throughout.

\emph{The geometric sums.} Substituting $m:=s-r$ and using $e_s=\rho^{-m}e_r$,
\begin{equation*}
\sum_{r\le s\le n}q^{s-r}e_s=e_r\sum_{m=0}^{n-r}q^m\rho^{-m}.
\end{equation*}
Since $0<q\le\rho^2$, we have $q^m\rho^{-m}\le\rho^{2m}\rho^{-m}=\rho^m$ for every $m\ge0$. All terms being
nonnegative, the finite sum is bounded by the full geometric series:
\begin{equation*}
e_r\sum_{m=0}^{n-r}q^m\rho^{-m}\le e_r\sum_{m\ge0}\rho^m=\frac{e_r}{1-\rho}.
\end{equation*}
If the sum is restricted to $r<s\le n$, only the terms with $m\ge1$ occur, and the same comparison gives
the bound $e_r\sum_{m\ge1}\rho^m=\rho e_r/(1-\rho)$. This proves the first two inequalities
of~\eqref{9.11:eq-envelope-sums-b}.

For the third, write $u=r-j$ with $0\le j\le r-r_0$; then $e_u=\rho^{n-r+j}$, and
\begin{equation*}
\sum_{r_0\le u\le r}e_u=\sum_{u=r_0}^{r}\rho^{n-u}\le\sum_{j\ge0}\rho^{n-r+j}=\frac{\rho^{n-r}}{1-\rho}
=\frac{e_r}{1-\rho}.
\end{equation*}
By \eqref{9.11:eq-distance-reduction-a}, for $r_0<t\le n$ the distance at the envelope is
$d(E)_t=\mathsf{A}C_E\sum_{r_0\le u\le t-1}e_u$, and $d(E)_t=0$ for $t>n$. Applying the inequality
just proved with $r:=t-1\ge r_0$ gives $d(E)_t\le\mathsf{A}C_Ee_{t-1}/(1-\rho)$ for $r_0<t\le n$. This
proves~\eqref{9.11:eq-envelope-sums-b}. The sum over $u$ is geometric, so the number of its terms does
not enter the bound; only the factor $(1-\rho)^{-1}$ does.

\emph{The band sums.} Let $r<s\le n$ with $s-r\le\mathsf{a}_s$. Since $0<\rho<1$, the map
$x\mapsto\rho^{-x}$ is increasing, so $e_s=\rho^{-(s-r)}e_r\le\rho^{-\mathsf{a}_s}e_r$. Multiplying by
$\hat{\mathfrak{r}}_s$, summing over the band and then enlarging the range to all depths $s\ge0$ (the terms
are nonnegative),
\begin{equation*}
\Sigma^{\mathrm{band}}_{s>r}\,\hat{\mathfrak{r}}_se_s
\le e_r\sum_{\substack{r<s\le n\\ s-r\le\mathsf{a}_s}}\hat{\mathfrak{r}}_s\rho^{-\mathsf{a}_s}
\le e_r\sum_{s\ge0}\hat{\mathfrak{r}}_s\rho^{-\mathsf{a}_s}=\mathfrak{S}\,e_r ,
\end{equation*}
which is~\eqref{9.11:eq-envelope-sums-c}.

\emph{The window kernel.} We use the following inequality: for integers $m\ge1$ and $\mathsf{a}\ge1$,
\begin{equation}\label{9.11:eq-envelope-sums-1}
2(m-\mathsf{a}-1)_+-m\ge-\mathsf{a}-1 .
\end{equation}
Indeed, if $m\le\mathsf{a}+1$ then $(m-\mathsf{a}-1)_+=0$, and the left side is $-m\ge-\mathsf{a}-1$. If
$m\ge\mathsf{a}+1$ then $(m-\mathsf{a}-1)_+=m-\mathsf{a}-1$, and
\begin{equation*}
2(m-\mathsf{a}-1)-m=m-2\mathsf{a}-2\ge(\mathsf{a}+1)-2\mathsf{a}-2=-\mathsf{a}-1,
\end{equation*}
with equality at $m=\mathsf{a}+1$. This proves~\eqref{9.11:eq-envelope-sums-1}.

Now let $r<s\le n$ and apply~\eqref{9.11:eq-envelope-sums-1} with $m:=s-r\ge1$ and
$\mathsf{a}:=\mathsf{a}_s\ge1$. Since $0<\rho<1$, the map $x\mapsto\rho^{x}$ is decreasing, so
\begin{equation*}
\mathsf{k}_{r,s}e_s=\rho^{2(s-r-\mathsf{a}_s-1)_+}\rho^{-(s-r)}e_r
=\rho^{\,2(m-\mathsf{a}_s-1)_+-m}\,e_r
\le\rho^{-\mathsf{a}_s-1}e_r=\rho^{-1}\rho^{-\mathsf{a}_s}e_r .
\end{equation*}
Multiplying by $\hat{\mathfrak{r}}_s\ge0$ and summing over $r<s\le n$,
\begin{equation*}
\sum_{r<s\le n}\mathsf{k}_{r,s}\hat{\mathfrak{r}}_se_s
\le\rho^{-1}e_r\sum_{r<s\le n}\hat{\mathfrak{r}}_s\rho^{-\mathsf{a}_s}
\le\rho^{-1}\mathfrak{S}\,e_r ,
\end{equation*}
where the last step again enlarges the range of nonnegative terms to all $s\ge0$. This
proves~\eqref{9.11:eq-envelope-sums-d}. No bound on $\mathsf{a}_s$ or on $s-r$ has been used. On the unit
window of Remark~\ref{9.11:rem-memory-structures}, that is for $m\le\mathsf{a}_s+1$, the kernel equals $1$
and $e_s\le\rho^{-\mathsf{a}_s-1}e_r$. Beyond the window the decay $\rho^{2(m-\mathsf{a}_s-1)}$ of the kernel
dominates the growth $\rho^{-m}$ of $e_s$. The worst case is the edge of the window,
$m=\mathsf{a}_s+1$.

\emph{The distance sums.} Let $r\le s<n$. By hypothesis, $g_{(s+1)}^2\le1$. We apply the bound on $d(E)_t$
in~\eqref{9.11:eq-envelope-sums-b} at $t=s+1$, which gives
$d(E)_{s+1}\le\mathsf{A}C_Ee_s/(1-\rho)$. Then we apply the first inequality
of~\eqref{9.11:eq-envelope-sums-b} to the sum over $s$, whose range $r\le s<n$ is contained in
$r\le s\le n$, the terms being nonnegative. Together these give
\begin{equation*}
\sum_{r\le s<n}q^{s-r}g_{(s+1)}^2\,d(E)_{s+1}
\le\sum_{r\le s<n}q^{s-r}\,\frac{\mathsf{A}C_Ee_s}{1-\rho}
\le\frac{\mathsf{A}C_E}{1-\rho}\cdot\frac{e_r}{1-\rho}=\frac{\mathsf{A}C_Ee_r}{(1-\rho)^2},
\end{equation*}
which is the first inequality of~\eqref{9.11:eq-envelope-sums-e}.

For the second inequality we bound the two terms of $d^{\sharp}(E)_{r+1}$ separately, for every $i\ge0$.
Consider first $d(E)_{r+1+i}$. By the bound on $d(E)_t$ in~\eqref{9.11:eq-envelope-sums-b} at $t=r+1+i$,
together with $e_{r+i}=\rho^{-i}e_r$,
\begin{equation*}
d(E)_{r+1+i}\le\frac{\mathsf{A}C_Ee_{r+i}}{1-\rho}=\frac{\mathsf{A}C_E\rho^{-i}e_r}{1-\rho}.
\end{equation*}
Similarly, at $t=r+2+i$, using $e_{r+1+i}=\rho^{-i-1}e_r$,
\begin{equation*}
d(E)_{r+2+i}\le\frac{\mathsf{A}C_E\rho^{-i-1}e_r}{1-\rho}.
\end{equation*}
Indices beyond $n$ contribute $0$ by~\eqref{9.11:eq-envelope-sums-b}, and the right sides are positive,
so both bounds hold for every $i\ge0$. Since $L^{-1/2}<q\le\rho^2$, we have
$L^{-i/2}=(L^{-1/2})^i\le\rho^{2i}$, hence $L^{-i/2}\rho^{-i}\le\rho^{i}$. Therefore
\begin{align*}
d^{\sharp}(E)_{r+1}
&\le\frac{\mathsf{A}C_Ee_r}{1-\rho}\,\bigl(1+\rho^{-1}\bigr)\sum_{i\ge0}\rho^{i}
=\frac{\mathsf{A}C_Ee_r\,(1+\rho^{-1})}{(1-\rho)^2}\\
&=\frac{\mathsf{A}C_Ee_r\,(1+\rho)}{\rho(1-\rho)^2}=\hat{c}_\rho\,\mathsf{A}C_Ee_r ,
\end{align*}
which is the second inequality of~\eqref{9.11:eq-envelope-sums-e}.
\end{proof}

\begin{proof}[The three row inequalities of the maximum principle]\label{9.11:prf-row-inequalities}
We use the following data:
\begin{itemize}
\item the index set $I$, the subset $I_0$, the matrix $M=(M_{ij})$ and the inhomogeneous terms
$\sigma=(\sigma_i)$ fixed in \eqref{9.11:eq-distance-reduction-a};
\item the envelope $E$ of \eqref{9.11:eq-envelope-sums-a}, namely
\begin{equation*}
E_r^P=C_Ee_r,\qquad E_r^A=4C_Ee_r,\qquad E_r^R=4C_E\hat{\mathfrak{r}}_re_r,
\end{equation*}
with $e_r=\rho^{n-r}$ and the constants of \eqref{9.11:eq-rate-floors-c}.
\end{itemize}
We prove the three inequalities
\begin{equation*}
\sum_jM_{ij}E_j\le\tfrac12E_i,
\end{equation*}
one for each kind of index, together with the bounds $\sigma_i\le\frac12E_i$ on the inhomogeneous terms.
These are the inequalities used in the application of Lemma~\ref{9.6:lem-maximum-principle}.
The indices are the pairs $(r,\bullet)$ with $\bullet\in\{P,A,R\}$, and we write
\begin{equation*}
ME_r^{\bullet}:=\sum_jM_{(r,\bullet)j}E_j .
\end{equation*}
The contributions to $ME_r^{\bullet}$ are obtained from the terms of the hatted response rows
\eqref{9.11:eq-response-rows-a}, \eqref{9.11:eq-response-rows-b} and
\eqref{9.11:eq-response-rows-c}. In each term the competitor norms are replaced by the envelope and
the distance by $d(E)$. We treat the three rows in turn. Throughout, $\rho<1$.

\emph{Row $P$.} The row \eqref{9.11:eq-response-rows-a} has four contributions.

For the first, we use the bound on $d^{\sharp}(E)_{r+1}$ in \eqref{9.11:eq-envelope-sums-e}, and
then the ceiling $\hat{c}_\rho\alpha_r\mathsf{A}\le\frac18$ of \eqref{9.11:eq-coupling-ceilings-a}:
\begin{equation*}
\alpha_rd^{\sharp}(E)_{r+1}\le\alpha_r\hat{c}_\rho\mathsf{A}C_Ee_r\le\tfrac18C_Ee_r .
\end{equation*}

For the second, we insert the envelope, so that
\begin{equation*}
E_s^P+E_s^R=C_Ee_s(1+4\hat{\mathfrak{r}}_s).
\end{equation*}
We use $\hat{\mathfrak{r}}_s\le\frac1{16}$ from \eqref{9.11:eq-coupling-ceilings-b}, so that
$1+4\hat{\mathfrak{r}}_s\le\frac54$. We then use the second of the two geometric bounds in
\eqref{9.11:eq-envelope-sums-b},
\begin{equation*}
\sum_{r<s\le n}q^{s-r}e_s\le\frac{\rho e_r}{1-\rho},
\end{equation*}
and finally the ceiling $\varkappa\le(1-\rho)/32$ of \eqref{9.11:eq-coupling-ceilings-a}. This gives
\begin{align*}
\varkappa\sum_{r<s\le n}q^{s-r}\bigl(E_s^P+E_s^R\bigr)
&=\varkappa C_E\sum_{r<s\le n}q^{s-r}e_s\bigl(1+4\hat{\mathfrak{r}}_s\bigr)
\le\varkappa C_E\Bigl(1+\tfrac14\Bigr)\frac{\rho e_r}{1-\rho}\\
&\le\frac54\cdot\frac{1-\rho}{32}\cdot\frac{\rho}{1-\rho}\,C_Ee_r
=\frac{5\rho}{128}\,C_Ee_r\le\frac5{128}\,C_Ee_r .
\end{align*}

For the third, we have $E_s^A=4C_Ee_s$. The same geometric bound of
\eqref{9.11:eq-envelope-sums-b} applies, followed by the ceiling
$4\varkappa''_r\le\frac3{32}(1-\rho)$ of \eqref{9.11:eq-coupling-ceilings-a}:
\begin{equation*}
\varkappa''_r\sum_{r<s\le n}q^{s-r}E_s^A\le4\varkappa''_rC_E\frac{\rho e_r}{1-\rho}
\le\frac3{32}(1-\rho)\frac{\rho}{1-\rho}\,C_Ee_r\le\frac3{32}\,C_Ee_r .
\end{equation*}

For the fourth, we use $E_s^R=4C_E\hat{\mathfrak{r}}_se_s$ and the band sum of
\eqref{9.11:eq-envelope-sums-c},
\begin{equation*}
\Sigma^{\mathrm{band}}_{s>r}\hat{\mathfrak{r}}_se_s\le\mathfrak{S}e_r .
\end{equation*}
We then use \eqref{9.11:eq-coupling-ceilings-d} in the form $4C_F\mathfrak{S}\le\frac18$:
\begin{equation*}
C_F\Sigma^{\mathrm{band}}_{s>r}E_s^R=4C_EC_F\Sigma^{\mathrm{band}}_{s>r}\hat{\mathfrak{r}}_se_s
\le4C_EC_F\mathfrak{S}e_r\le\tfrac18C_Ee_r .
\end{equation*}

Adding the four bounds,
\begin{equation*}
ME_r^P\le\Bigl(\frac18+\frac5{128}+\frac3{32}+\frac18\Bigr)C_Ee_r=\frac{49}{128}\,C_Ee_r
\le\frac12C_Ee_r=\frac12E_r^P .
\end{equation*}
The inhomogeneous term of this row is $\sigma_r^P=C_\sigma e_r$. By the choice $C_E\ge2C_\sigma$ in
\eqref{9.11:eq-rate-floors-c}, $\sigma_r^P\le\frac12C_Ee_r=\frac12E_r^P$.

\emph{Row $A$.} The row \eqref{9.11:eq-response-rows-b} has four contributions.

For the first, we use the bound on $d^{\sharp}(E)_{r+1}$ in \eqref{9.11:eq-envelope-sums-e}, and
then the ceiling $\hat{c}_\rho\alpha_r^A\mathsf{A}\le\frac14$ of
\eqref{9.11:eq-coupling-ceilings-a}:
\begin{equation*}
\alpha_r^Ad^{\sharp}(E)_{r+1}\le\hat{c}_\rho\alpha_r^A\mathsf{A}C_Ee_r\le\tfrac14C_Ee_r .
\end{equation*}

For the second, the envelope gives
\begin{equation*}
E_s^P+E_s^A+E_s^R=C_Ee_s(1+4+4\hat{\mathfrak{r}}_s)\le(1+4+\tfrac14)C_Ee_s ,
\end{equation*}
by $\hat{\mathfrak{r}}_s\le\frac1{16}$ from \eqref{9.11:eq-coupling-ceilings-b}. With the second
geometric bound of \eqref{9.11:eq-envelope-sums-b} this yields
\begin{equation*}
\varkappa_A\Bigl[E_r^P+\sum_{r<s\le n}q^{s-r}\bigl(E_s^P+E_s^A+E_s^R\bigr)\Bigr]
\le\varkappa_AC_Ee_r\Bigl[1+\Bigl(1+4+\tfrac14\Bigr)\frac{\rho}{1-\rho}\Bigr].
\end{equation*}
The ceiling $\varkappa_A(9+5/(1-\rho))\le\frac14$ of \eqref{9.11:eq-coupling-ceilings-a} gives
\begin{equation*}
\varkappa_A\le\frac1{4\bigl(9+5/(1-\rho)\bigr)}=\frac{1-\rho}{4\bigl(9(1-\rho)+5\bigr)} .
\end{equation*}
Dropping the positive term $5/(1-\rho)$ from the first of these expressions gives
$\varkappa_A\le\frac1{36}$. Using the second expression, $1+4+\frac14=\frac{21}4$ and $\rho<1$,
\begin{equation*}
\varkappa_A\cdot\frac{21}4\cdot\frac{\rho}{1-\rho}\le\frac{21\rho}{16\bigl(9(1-\rho)+5\bigr)}
\le\frac{21}{80}.
\end{equation*}
Together, the second contribution is at most
\begin{equation*}
\Bigl(\frac1{36}+\frac{21}{80}\Bigr)C_Ee_r=\frac{209}{720}\,C_Ee_r\le\frac5{16}\,C_Ee_r .
\end{equation*}

The third contribution is the same-depth entry. By $\hat{\mathfrak{r}}_r\le\frac1{16}$ from
\eqref{9.11:eq-coupling-ceilings-b},
\begin{equation*}
E_r^R=4C_E\hat{\mathfrak{r}}_re_r\le\tfrac14C_Ee_r .
\end{equation*}

The fourth contribution is the band sum. By \eqref{9.11:eq-envelope-sums-c} and by
\eqref{9.11:eq-coupling-ceilings-d} in the form $4C_F^A\mathfrak{S}\le\frac18$,
\begin{equation*}
C_F^A\Sigma^{\mathrm{band}}_{s>r}E_s^R\le4C_EC_F^A\mathfrak{S}e_r\le\tfrac18C_Ee_r .
\end{equation*}

Adding the four bounds,
\begin{equation*}
ME_r^A\le\Bigl(\frac14+\frac5{16}+\frac14+\frac18\Bigr)C_Ee_r=\frac{15}{16}\,C_Ee_r
\le2C_Ee_r=\frac12E_r^A .
\end{equation*}
The inhomogeneous term satisfies
\begin{equation*}
\sigma_r^A=C_\sigma e_r\le\tfrac12C_Ee_r\le\tfrac12E_r^A .
\end{equation*}

\emph{Row $R$.} The row \eqref{9.11:eq-response-rows-c} has two contributions. The first comes from
the bracket $\hat{B}_r(\hat{G},d)$ of \eqref{9.11:eq-response-rows-c}, evaluated at $\hat{G}=E$ and
$d=d(E)$. We bound its three members separately.

The first member is at most $\mathsf{A}C_Ee_r/(1-\rho)^2$, by \eqref{9.11:eq-envelope-sums-e}.

For the second member, $E_s^P+E_s^A=5C_Ee_s$. The first of the two geometric bounds in
\eqref{9.11:eq-envelope-sums-b} gives
\begin{equation*}
\sum_{r\le s\le n}q^{s-r}\bigl(E_s^P+E_s^A\bigr)=5C_E\sum_{r\le s\le n}q^{s-r}e_s
\le\frac{5C_Ee_r}{1-\rho}.
\end{equation*}
The entries with $s=r$ of $P$ and $A$ are the term $m=0$ of that sum, and so are contained in the
bound $5/(1-\rho)$.

For the third member, we use the window sum of \eqref{9.11:eq-envelope-sums-d},
\begin{equation*}
\sum_{r<s\le n}\mathsf{k}_{r,s}\hat{\mathfrak{r}}_se_s\le\rho^{-1}\mathfrak{S}e_r ,
\end{equation*}
and \eqref{9.11:eq-coupling-ceilings-d} in the form $\rho^{-1}\mathfrak{S}\le\frac1{32}$:
\begin{equation*}
\sum_{r<s\le n}\mathsf{k}_{r,s}E_s^R=4C_E\sum_{r<s\le n}\mathsf{k}_{r,s}\hat{\mathfrak{r}}_se_s
\le4C_E\rho^{-1}\mathfrak{S}e_r\le\tfrac18C_Ee_r .
\end{equation*}

Hence, using $5\le6$ and $\frac18\le8$ in the second step,
\begin{align*}
\hat{B}_r\bigl(E,d(E)\bigr)
&\le C_Ee_r\bigl[\mathsf{A}(1-\rho)^{-2}+5(1-\rho)^{-1}+\tfrac18\bigr]\\
&\le C_Ee_r\bigl[\mathsf{A}(1-\rho)^{-2}+6(1-\rho)^{-1}+8\bigr].
\end{align*}
By the definition of $\hat{\mathfrak{r}}_r$ in \eqref{9.11:eq-rate-floors-c},
\begin{equation}\label{9.11:eq-row-inequalities-a}
\begin{split}
2\mathfrak{r}_r\hat{\Lambda}_r\hat{B}_r\bigl(E,d(E)\bigr)
&\le2\mathfrak{r}_r\hat{\Lambda}_r\bigl[\mathsf{A}(1-\rho)^{-2}+6(1-\rho)^{-1}+8\bigr]C_Ee_r\\
&=\hat{\mathfrak{r}}_rC_Ee_r=\tfrac14E_r^R .
\end{split}
\end{equation}
The terms $\mathsf{A}(1-\rho)^{-2}$ and $6(1-\rho)^{-1}$ in this definition are non-negative.
Retaining only the term $8$ therefore gives
\begin{equation}\label{9.11:eq-row-inequalities-1}
\hat{\mathfrak{r}}_r\ge2\mathfrak{r}_r\hat{\Lambda}_r\cdot8=16\,\mathfrak{r}_r\hat{\Lambda}_r .
\end{equation}

The second contribution is the window term. The sum
\begin{equation*}
\sum_{r<s<r+\ell_r,\ s\le n}E_s^R
\end{equation*}
has at most $\ell_r-1$ terms, indexed by $m:=s-r\in\{1,\dots,\ell_r-1\}$. For each of them:
\begin{itemize}
\item $\hat{\mathfrak{r}}_{r+m}\le\hat{\mathfrak{r}}_r$, by the monotonicity in
\eqref{9.11:eq-coupling-ceilings-b};
\item $e_{r+m}=\rho^{-m}e_r\le\rho^{-\ell_r}e_r$, since $\rho<1$ and $m<\ell_r$.
\end{itemize}
Hence
\begin{equation*}
2\varkappa'_r\sum_{r<s<r+\ell_r,\ s\le n}E_s^R
\le8C_E\varkappa'_r\hat{\mathfrak{r}}_re_r\cdot\ell_r\rho^{-\ell_r}
=2\varkappa'_r\ell_r\rho^{-\ell_r}\cdot E_r^R .
\end{equation*}
We use \eqref{9.11:eq-coupling-ceilings-c} in the form $4\varkappa'_r\ell_r^2\rho^{-\ell_r}\le\frac14$,
together with $\ell_r\ge2$:
\begin{equation*}
2\varkappa'_r\ell_r\rho^{-\ell_r}=\frac1{2\ell_r}\cdot4\varkappa'_r\ell_r^2\rho^{-\ell_r}
\le\frac1{8\ell_r}\le\frac1{16}\le\frac14 .
\end{equation*}
Adding this to \eqref{9.11:eq-row-inequalities-a},
\begin{equation*}
ME_r^R\le\tfrac14E_r^R+\tfrac14E_r^R=\tfrac12E_r^R .
\end{equation*}

It remains to bound the inhomogeneous term of this row. The choices $C_E\ge2C_\sigma$ and
$C_E\ge C_V+1\ge1$ in \eqref{9.11:eq-rate-floors-c} give
\begin{equation*}
2C_\sigma+1\le C_E+1\le2C_E .
\end{equation*}
Combining this with \eqref{9.11:eq-row-inequalities-1},
\begin{equation*}
\sigma_r^R=\mathfrak{r}_r\hat{\Lambda}_r(2C_\sigma+1)e_r
\le\frac{\hat{\mathfrak{r}}_r}{16}\cdot2C_Ee_r=\frac1{32}E_r^R\le\frac12E_r^R .
\end{equation*}
This completes the verification of the three row inequalities and of the three bounds on the
inhomogeneous terms.
\end{proof}

\begin{proof}[Proof of Theorem~\ref{9.11:thm-geometric-envelope}, continued: the finest level and
conclusion of the envelope bound]\label{9.11:prf-finest-level}
\emph{The finest level $r=n$.} In the system of inequalities set up in
\eqref{9.11:eq-distance-reduction-a}, the indices $(r,\bullet)$ that do not lie in $I_0$ are those
with $r=n$. For them we verify $D_i\le E_i$ directly.

The boundary clause of \eqref{9.11:eq-chain-bounds-a} in Lemma~\ref{9.11:lem-chain-bounds}
gives
\begin{equation}\label{9.11:eq-finest-level-1}
  \hat{g}_n^c\le C_V\,\hat{\mathfrak{u}}_n^c\qquad\text{for every component } c .
\end{equation}
Moreover $e_n=\rho^{n-n}=1$. With the envelope of \eqref{9.11:eq-envelope-sums-a} this gives
\begin{equation*}
  E_n^P=C_E,\qquad E_n^A=4C_E,\qquad E_n^R=4C_E\,\hat{\mathfrak{r}}_n .
\end{equation*}

\emph{The block $P$.} The reading units of the components of the block $P$ at depth $n$ are those
fixed in \eqref{9.11:eq-hatted-blocks-b}, and the largest of them is $\max\{1,\hat{\mathfrak{u}}_n^Q\}$.
Taking in \eqref{9.11:eq-finest-level-1} the maximum over the components of $P$ therefore gives
\begin{equation*}
  \hat{g}_n^P\le C_V\max\{1,\hat{\mathfrak{u}}_n^Q\}\le C_V\le C_E=E_n^P .
\end{equation*}
The second inequality holds because $\hat{\mathfrak{u}}_n^Q\le 1$ by the ceiling
\eqref{9.11:eq-coupling-ceilings-e}. The third holds because $C_E\ge C_V+1$ by the choice of the
rate constants \eqref{9.11:eq-rate-floors-c}.

\emph{The block $A$.} The reading units of the components of $A$ at depth $n$ are
$\mathsf{l}_n^{-1}$ and $2\mathsf{l}_n^{-1}$. Hence \eqref{9.11:eq-finest-level-1} gives
\begin{equation*}
  \hat{g}_n^A\le C_V\max\{\mathsf{l}_n^{-1},2\mathsf{l}_n^{-1}\}\le 2C_V\le 4C_E=E_n^A .
\end{equation*}
The second inequality holds because the maximum, $2\mathsf{l}_n^{-1}$, is at most $2$ by the ceiling
\eqref{9.11:eq-coupling-ceilings-e}. The third holds because $C_V\le C_E$.

\emph{The block $R$.} The reading unit of the block $R$ at depth $n$ is $\mathfrak{r}_n$, so
\eqref{9.11:eq-finest-level-1} gives
\begin{equation*}
  \hat{g}_n^R\le C_V\,\mathfrak{r}_n\le C_E\,\hat{\mathfrak{r}}_n\le 4C_E\,\hat{\mathfrak{r}}_n=E_n^R .
\end{equation*}
The second inequality uses two facts: $C_V\le C_E$, and
\begin{equation*}
  \mathfrak{r}_n\le 16\,\mathfrak{r}_n\hat{\Lambda}_n\le\hat{\mathfrak{r}}_n .
\end{equation*}
Here the first inequality holds because $\mathfrak{r}_n>0$ and $\hat{\Lambda}_n\ge 1$ by the ceiling
\eqref{9.11:eq-coupling-ceilings-e}. The second follows from the definition of $\hat{\mathfrak{r}}_n$
among the rate constants \eqref{9.11:eq-rate-floors-c}, through the additive term $8$ it contains.
The last inequality of the block-$R$ chain holds because $\hat{\mathfrak{r}}_n\ge 0$.
This proves $D_i\le E_i$ for every index $i\notin I_0$.

\emph{Conclusion.} We apply the maximum principle, Lemma~\ref{9.11:lem-maximum-principle}, on the
finite set $I$ of \eqref{9.11:eq-distance-reduction-a}, with $I_0$ the set of indices $(r,\bullet)$
with $r<n$. Its hypotheses hold for three reasons.
\begin{itemize}
\item By the reduction \eqref{9.11:eq-distance-reduction-a}, the substituted system has the form
  required in Lemma~\ref{9.11:lem-maximum-principle}, with $D_{(r,\bullet)}=\hat{g}_r^\bullet$.
\item By the three row inequalities of the maximum principle, established in the part of this
  proof devoted to them, we have $\sum_j M_{ij}E_j\le\frac12 E_i$ and
  $\sigma_i\le\frac12 E_i$ for every $i\in I_0$.
\item By the preceding paragraphs, $D_i\le E_i$ for every $i\notin I_0$.
\end{itemize}
Lemma~\ref{9.11:lem-maximum-principle} therefore gives $D_i\le E_i$ for every $i\in I$, that is,
\begin{equation}\label{9.11:eq-finest-level-2}
  \hat{g}_r^P\le C_E\,e_r,\qquad \hat{g}_r^A\le 4C_E\,e_r,\qquad
  \hat{g}_r^R\le 4C_E\,\hat{\mathfrak{r}}_r\,e_r
  \qquad (r_0\le r\le n).
\end{equation}
In addition, $\hat{\mathsf{D}}_r^c\le\hat{g}_r^c$ by (c3) of \eqref{9.11:eq-chain-bounds-a}.

\emph{The bounds on the coupling distances.} Let $r_0\le r<n$. By the reduction
\eqref{9.11:eq-distance-reduction-a}, the first bound of \eqref{9.11:eq-finest-level-2} and the
second geometric sum of \eqref{9.11:eq-envelope-sums-b}, namely
$\sum_{r_0\le u\le r}e_u\le e_r/(1-\rho)$, we obtain
\begin{equation}\label{9.11:eq-finest-level-3}
\begin{aligned}
  d_{r+1}
  &\le\mathsf{A}\sum_{r_0\le u\le r}\hat{g}_u^P
   \le\mathsf{A}C_E\sum_{r_0\le u\le r}e_u\\
  &\le\frac{\mathsf{A}C_E\,e_r}{1-\rho}
   =\frac{\mathsf{A}C_E\,\rho^{\,n-r}}{1-\rho}.
\end{aligned}
\end{equation}
Now consider a coupling distance $d_t$ occurring in the weighted two-level distance $d^{\sharp}_{r+1}$
of \eqref{9.11:eq-competitors-distances-a}. Each such $d_t$ has either $t>n$ or $r_0<t\le n$, and
is bounded as follows.
\begin{itemize}
\item If $t>n$, then $d_t=0$.
\item If $r_0<t\le n$, then \eqref{9.11:eq-finest-level-3} applied with $t-1$ in place of $r$, which
  is admissible since $r_0\le t-1<n$, gives
  \begin{equation*}
    d_t\le\frac{\mathsf{A}C_E\,e_{t-1}}{1-\rho}.
  \end{equation*}
\end{itemize}
These are the bounds that \eqref{9.11:eq-envelope-sums-b} gives for the distances $d(E)_t$
evaluated at the envelope. The computation of the bound on $d^{\sharp}(E)_{r+1}$ in
\eqref{9.11:eq-envelope-sums-e} uses only these bounds on the distances.
Carried out with the $d_t$ in place of the $d(E)_t$, it gives
\begin{equation*}
  d^{\sharp}_{r+1}\le\hat{c}_\rho\,\mathsf{A}C_E\,\rho^{\,n-r}
  \qquad(r_0\le r<n).
\end{equation*}
Together with \eqref{9.11:eq-finest-level-2}, the bound $\hat{\mathsf{D}}_r^c\le\hat{g}_r^c$ and
\eqref{9.11:eq-finest-level-3}, these are the bounds \eqref{9.11:eq-geometric-envelope-a}, which
completes the proof of Theorem~\ref{9.11:thm-geometric-envelope}.
\end{proof}

\begin{remark}[Enlarging the envelope constant and the logarithmic weight]
\label{9.11:rem-envelope-enlargement}
The proof of the geometric envelope for two-run differences of hatted runs,
Theorem~\ref{9.11:thm-geometric-envelope}, uses the envelope constant $C_E$ and the logarithmic
weight $\hat{\Lambda}_r$ only through the following two properties.

First, it uses of the envelope constant $C_E$ only the lower bound
\begin{equation}\label{9.11:eq-envelope-enlargement-1}
  C_E \;\ge\; \max\{C_V+1,\ 2C_\sigma\}.
\end{equation}

Second, it uses of the logarithmic weight $\hat{\Lambda}_r$ only that one and the same
$\hat{\Lambda}_r$ appears in three places:
\begin{itemize}
\item it multiplies the bracket in each of the two hypotheses of
  Theorem~\ref{9.11:thm-geometric-envelope} in which $\hat{\Lambda}_r$ occurs, namely the
  hatted bound on the remnant slots and the hatted third hypothesis on the history kernels;
\item it enters the definition of the hatted rate $\hat{\mathfrak{r}}_r$ among the rate
  constants~\eqref{9.11:eq-rate-floors-c}.
\end{itemize}

Consequently the conclusion of Theorem~\ref{9.11:thm-geometric-envelope} remains valid under
either or both of the following replacements:
\begin{itemize}
\item $C_E$ is replaced by any larger constant;
\item $\hat{\Lambda}_r$ is replaced, in all three places at once, by one common larger weight,
  the ceilings on the coupling \eqref{9.11:eq-coupling-ceilings-b} and
  \eqref{9.11:eq-coupling-ceilings-d} being read with the enlarged weight.
\end{itemize}
In particular, enlarging $C_E$ by a fixed factor, and $\hat{\Lambda}_r$ by a fixed factor in
these three places, after the induction of that proof has been run leaves every inequality of
the proof intact.
\end{remark}

\begin{proof}
\emph{Enlarging $C_E$.}
The inequalities in which $C_E$ occurs are those of the part of the proof of
Theorem~\ref{9.11:thm-geometric-envelope} that establishes the three row inequalities of the
maximum principle, among them \eqref{9.11:eq-row-inequalities-a}, and those of the part that
treats the finest level and draws the conclusion. Each of them is of one of two
kinds.

The first kind is homogeneous in $C_E$: its two sides are homogeneous in $C_E$ of the same
degree. After division by the common power of $C_E$ it no longer contains $C_E$. It therefore
holds for one positive value of $C_E$ if and only if it holds for every positive value, and in
particular for every larger one.

The second kind consists of the comparisons
\begin{equation*}
  C_\sigma \le \tfrac12 C_E,
  \qquad
  C_V \le C_E,
\end{equation*}
together with $C_E \ge C_V+1$ itself. All of these follow
from~\eqref{9.11:eq-envelope-enlargement-1}. Their right-hand sides increase with $C_E$ while
their left-hand sides do not involve it, so they persist when $C_E$ is enlarged. Since
\eqref{9.11:eq-envelope-enlargement-1} is a lower bound only, it is itself preserved by any
enlargement of $C_E$.

\emph{Enlarging $\hat{\Lambda}_r$.}
Every occurrence of $\hat{\Lambda}_r$ in the comparisons of the proof is of one of two kinds.

The first kind is the definition of $\hat{\mathfrak{r}}_r$.

The second kind is absorbed into $\hat{\mathfrak{r}}_r$ by the inequality
\begin{equation}\label{9.11:eq-envelope-enlargement-2}
  \mathfrak{r}_r\hat{\Lambda}_r \;\le\; \tfrac{1}{16}\,\hat{\mathfrak{r}}_r .
\end{equation}
This is how the product $2\mathfrak{r}_r\hat{\Lambda}_r$ is absorbed in the row-$R$
inequality~\eqref{9.11:eq-row-inequalities-a}. It is also how $\mathfrak{r}_n$ is compared
with $\hat{\mathfrak{r}}_n$ at the finest level.

Under a common enlargement the absorption~\eqref{9.11:eq-envelope-enlargement-2}, the definition
of $\hat{\mathfrak{r}}_r$ and the two hypotheses named in the remark are all read with the same
enlarged weight.

Apart from these comparisons, the weight enters only the ceilings on the coupling
\eqref{9.11:eq-coupling-ceilings-b} and \eqref{9.11:eq-coupling-ceilings-d}, which are monotone
in the weights $\hat{\Lambda}_r$ and are read with the enlarged weight.

Hence every inequality of the proof of Theorem~\ref{9.11:thm-geometric-envelope} is preserved
by any enlargement of $C_E$ and by any common enlargement of $\hat{\Lambda}_r$ in the three
places. The conclusion of that theorem therefore holds with the enlarged constants.
\end{proof}

\begin{proposition}[Un-decayed rows admit no pure geometric envelope]\label{9.11:prop-no-geometric-envelope}
Consider the un-hatted three-row system of response rows, with rows $(P)$, $(A)$, $(R)$, its $d$-line
\begin{equation}\label{9.11:eq-no-geometric-envelope-1}
d_{s+1}\le \mathsf{A}\sum_{u\le s} g_u^P ,
\end{equation}
and its boundary condition. Read it as what it is: a family of inequalities with non-negative
coefficients constraining the non-negative unknowns $(g_r^c)$, $c\in\{P,A,R\}$, and $(d_s)$.
Fix any $\rho\in\,]0,1[$, any $C>0$, and the constants $\varkappa>0$, $C_\sigma\ge 0$ and
$\mathfrak{r}_s\Lambda_s>0$ of the system. Then for every $n$ there are non-negative numbers
$(g_r^c)_{0\le r\le n}$ and $(d_s)$ with the following properties. They satisfy every inequality of
the system with $e_r:=\rho^{n-r}$ in the sources. They satisfy the boundary condition. And
\begin{equation}\label{9.11:eq-no-geometric-envelope-2}
g_r^P\ \ge\ \varkappa\rho\,\mathfrak{r}_{n-1}\Lambda_{n-1}(2C_\sigma+1)
\qquad\text{for every } r\le n-2 .
\end{equation}
The ladder is $\theta_{n-1}=X+(n-1)c_0$. The product $\mathfrak{r}_{n-1}\Lambda_{n-1}$ is the power
$\theta_{n-1}^{1-\kappa_0/2}(\log\theta_{n-1})^{c_3}$ up to the constant factor $C_t\Lambda$, while
$\rho^n$ decays geometrically. Hence
\begin{equation*}
\frac{\rho^n}{\mathfrak{r}_{n-1}\Lambda_{n-1}}\longrightarrow 0\qquad (n\to\infty).
\end{equation*}
Consequently, for $n$ large the bound $g_0^P\le C\rho^n$ fails for these numbers. The three rows,
the $d$-line and the boundary condition, read alone, therefore do not imply
$g_r^P\le C\rho^{n-r}$.
\end{proposition}

\begin{proof}
Put $d_s:=0$ for all $s$. For the remaining unknowns put
\begin{equation}\label{9.11:eq-no-geometric-envelope-3}
\begin{aligned}
g_s^R &:= \mathfrak{r}_s\Lambda_s(2C_\sigma+1)\rho^{n-s}\quad (0\le s<n), &\qquad g_n^R&:=0,\\
g_r^A &:= C_\sigma\rho^{n-r}\quad (r<n), & g_n^A&:=0,\\
g_r^P &:= \varkappa\sum_{r<s\le n} g_s^R + C_\sigma\rho^{n-r}\quad (r<n), & g_n^P&:=0.
\end{aligned}
\end{equation}
All these numbers are non-negative, because $\varkappa>0$, $C_\sigma\ge 0$, $\mathfrak{r}_s\Lambda_s>0$
and $\rho>0$. With $e_r=\rho^{n-r}$ the definitions read
\[
g_r^A=C_\sigma e_r,\qquad g_r^R=\mathfrak{r}_r\Lambda_r(2C_\sigma+1)e_r\qquad (r<n).
\]

We check the system row by row.

\emph{Row $(P)$.} Its right member contains the terms $\varkappa\sum_{s>r}g_s^R$ and $C_\sigma e_r$.
All its other terms are non-negative, being non-negative coefficients times non-negative unknowns.
Since $g_n^R=0$, the sum $\sum_{s>r}g_s^R$ equals the sum over $r<s\le n$ in
\eqref{9.11:eq-no-geometric-envelope-3}. Hence $g_r^P$ is at most the right member of row $(P)$.

\emph{Row $(A)$.} Its right member contains $C_\sigma e_r$, which equals $g_r^A$. Its other terms are
non-negative, so row $(A)$ holds.

\emph{Row $(R)$.} Its right member contains $\mathfrak{r}_r\Lambda_r(2C_\sigma+1)e_r$, which equals
$g_r^R$. Its other terms are non-negative, so row $(R)$ holds.

\emph{The $d$-line.} Since $d\equiv 0$ and every $g_u^P\ge 0$, we have
$d_{s+1}=0\le\mathsf{A}\sum_{u\le s}g_u^P$. This is \eqref{9.11:eq-no-geometric-envelope-1}.

\emph{The boundary condition.} It holds, since the unknowns with index $n$ are all $0$.

\emph{The lower bound.} Let $r\le n-2$. Then $s=n-1$ lies in the range $r<s\le n$ of the sum
defining $g_r^P$. The other summands and the term $C_\sigma\rho^{n-r}$ are non-negative. Therefore
\[
g_r^P\ \ge\ \varkappa\,g_{n-1}^R\ =\ \varkappa\,\mathfrak{r}_{n-1}\Lambda_{n-1}(2C_\sigma+1)\rho ,
\]
which is \eqref{9.11:eq-no-geometric-envelope-2}.

\emph{The consequence.} Let $n\ge 2$, so that $r=0$ is admissible in
\eqref{9.11:eq-no-geometric-envelope-2}. Suppose $g_0^P\le C\rho^n$. Then
\[
\varkappa\rho(2C_\sigma+1)\ \le\ C\,\frac{\rho^n}{\mathfrak{r}_{n-1}\Lambda_{n-1}} .
\]
The left side is a fixed positive number. The right side tends to $0$ as $n\to\infty$, by the
limit displayed in the statement. Hence the inequality is false for all $n$ large enough, and for
such $n$ the numbers \eqref{9.11:eq-no-geometric-envelope-3} satisfy the whole system while
$g_0^P>C\rho^n$.
\end{proof}

The same construction applies when only some of the rows are hatted. Keep the un-decayed entry only
in row $(R)$, inside its term $B_r$, and replace rows $(P)$ and $(A)$ by the hatted rows
\eqref{9.11:eq-response-rows-a} and \eqref{9.11:eq-response-rows-b} of
Definition~\ref{9.11:def-response-rows}. The construction then gives
$g_r^R\ge 2\mathfrak{r}_r\Lambda_r\,g_{n-1}^R$ for $r\le n-2$.

Now use the band term $C_F\Sigma^{\mathrm{band}}_{s>r}\|\hat{G}_s^R\|$ of
\eqref{9.11:eq-response-rows-a}. When $\mathsf{a}_{r+1}\ge 1$, the index $s=r+1$ satisfies
$s-r\le\mathsf{a}_s$ and so lies in the band. This gives
\[
g_r^P\ \ge\ C_F\cdot 2\mathfrak{r}_{r+1}\Lambda_{r+1}\,\mathfrak{r}_{n-1}\Lambda_{n-1}(2C_\sigma+1)\rho
\qquad\text{whenever } \mathsf{a}_{r+1}\ge 1 .
\]
The right side is a product of two power laws, and it is still not bounded by $C\rho^{n-r}$.

Each of the three un-decayed entries of the un-hatted system separately obstructs the geometric
bound $g_r^P\le C\rho^{n-r}$. The entry in $(P)$ does so directly. The entry in $(A)$ does so because
it feeds the sum $\sum q^{s-r}\|\hat{G}_s^A\|$ in \eqref{9.11:eq-response-rows-c}, and thence the band
term of \eqref{9.11:eq-response-rows-a}. The entry in $(R)$ does so as just shown.

\begin{remark}[Origin of the un-decayed remnant entry]\label{9.11:rem-undecayed-entry}
Write (P), (A), (R) for the un-hatted response inequalities for the blocks $P$, $A$, $R$,
transplanted across the cutoff. They are replaced in
Definition~\ref{9.11:def-response-rows} (the hatted response rows with window kernel) by the
inequality \eqref{9.11:eq-response-rows-a}, that is $(\hat{P})$, by the corresponding row
$(\hat{A})$ of that definition, and by \eqref{9.11:eq-response-rows-c}, that is $(\hat{R})$.
The one-lattice Lipschitz step is established in two forms, which we call the first and the
second form: the first is formulated with the loads $[\![\delta]\!]_n$ and
$[\![\delta]\!]^T_k$ below, the second with history kernels of the remnant and boundary slots,
among them the kernel $\mathfrak{w}^R$ below. We use the following notation.
\begin{itemize}
\item $\delta\mathfrak{e}^P_s$ and $\delta\mathfrak{e}^R_s$ are the differences between the two
runs of the components of the produced depth-$s$ block in the blocks $P$ and $R$.
\item $\mathsf{r}_j$ is the remnant difference at level $j$.
\item $\mathsf{w}_{n,j}$ are the weights with which the load $[\![\delta]\!]_n$ of the first
form reads $\mathsf{r}_j$.
\item $[\![\delta]\!]^T_k$ is the second load of the first form, and $k_1(k)$ is the level up
to which its remnant sum runs.
\item The indices $k$, $n$, $j$ are level indices of the first form, quoted with its letters,
and $x_k=g_k^{-2}$.
\end{itemize}
\begin{enumerate}
\item[(i)] Every entry of (P), (A) and (R) without decay in the depth difference arises, in
each of the two forms, from one term. In the first form it arises from the member
$\sum_{j\le n}\mathsf{w}_{n,j}\mathsf{r}_j$ of the load $[\![\delta]\!]_n$, together with the
remnant member of $[\![\delta]\!]^T_k$. In the second form it arises from the remnant history
kernel $\mathfrak{w}^R$. Passing to $(\hat{R})$ replaces the first form's image by the hatted
kernel $\hat{\mathfrak{w}}^R$ of the second form.
\item[(ii)] The windows of $(\hat{P})$, $(\hat{A})$ and $(\hat{R})$ are not the obstruction to a
geometric envelope exhibited in Proposition~\ref{9.11:prop-no-geometric-envelope} for the
un-hatted rows: each closes under one of the ceilings on the coupling of
Definition~\ref{9.11:def-coupling-ceilings}. Here a member of a row is said to close under a
condition if, with the envelope $E^P_r$, $E^A_r$, $E^R_r$ of the hatted system inserted for the
unknowns, that condition bounds the member by a fixed fraction of the envelope at depth $r$.
\end{enumerate}
\end{remark}

\begin{proof}
\emph{(i), first form.} Consider the proof of the corollary of the first form that gives the
response inequalities for the blocks $P$ and $R$. There the member
$\varkappa\sum_{s>r}\|\delta\mathfrak{e}^R_s\|$ of (P) is $C_tx_k$ times the bound of
clause~(b) of the one-lattice Lipschitz step in its first form, and this uses
\begin{equation*}
\mathsf{w}_{k,j}\,x_k/x_j\le 5/2 .
\end{equation*}
The sum $\sum_{s>r}\|\delta\mathfrak{e}^R_s\|$ inside $\mathfrak{r}_r\Lambda_r[\,\cdot\,]$ in (R)
is $C_tx_k$ times the sum of the bounds of clauses~(c) and~(e) of that form, and this uses
\begin{equation*}
\mathsf{w}_{n,j}\,\mathsf{r}_j\le \frac{5}{2C_tx_n}\,\|\delta\mathfrak{e}^R_j\| .
\end{equation*}
Both entries are therefore images of one term, namely the member
$\sum_{j\le n}\mathsf{w}_{n,j}\mathsf{r}_j$ of $[\![\delta]\!]_n$. This member weights the
remnant differences of all older levels $j\le n$ by
\begin{equation*}
\mathsf{w}_{n,j}\le \tfrac52\,x_j/x_n ,
\end{equation*}
with no decay in $n-j$. Both entries are also images of the member of $[\![\delta]\!]^T_k$ whose
sum over $\mathsf{r}_j$ runs over $j\le k_1(k)$.

\emph{(i), second form.} Here the same entry is the remnant history kernel
\begin{equation}\label{9.11:eq-undecayed-entry-1}
\mathfrak{w}^R_{k,j}:=\sum_{i=j}^{k}\bar{w}_{i,j}\,\vartheta_*^{\,k-i}
\le 4\,x_j/x_k+\varpi' ,
\end{equation}
where $\bar{w}_{i,j}$ are the unit weights of the remnant slot. The bound
\eqref{9.11:eq-undecayed-entry-1} is flat in $k-j$. By the corollary of the second form on the
remnant inputs, this kernel enters the remnant hypothesis of that form's distance bound, and these
remnant inputs remain flat. The kernel also enters the remnant line of the response corollary of
the second form,
\begin{equation*}
\sum_{s>r}\bigl(q^{s-r}\|\delta\mathfrak{e}^P_s\|+\|\delta\mathfrak{e}^R_s\|
+q^{s-r}|\delta\xi_s|\bigr),
\end{equation*}
where $\delta\xi_s$ is the difference between the two runs of the scalar quantity $\xi_s$
attached to depth $s$ in the second form; in this line only the remnant member carries no
factor $q^{s-r}$.

\emph{The hatted replacement.} Remark~\ref{9.11:rem-window-kernel} (the window kernel from the
hatted remnant kernel) replaces $\mathfrak{w}^R$ by a kernel with
\begin{equation*}
\hat{\mathfrak{w}}^R_{k,j}\le \bar{w}_j\,C_{**}\,\rho^{2(k-j'')_+},
\end{equation*}
where $\bar{w}_j$ is the unit weight of the remnant slot at level $j$.
This kernel is at most $\bar{w}_jC_{**}$ for $k\le j''=j+\mathsf{a}_j$ and decays geometrically
in $k-j''$ beyond. From it come the third member $\sum_{s>r}\mathsf{k}_{r,s}\|\hat{G}^R_s\|$ of
$\hat{B}_r$ in \eqref{9.11:eq-response-rows-c} and the replacement of $\Lambda_r$ by
$\hat{\Lambda}_r=C_{**}\Lambda_r$, that is, of the constant $\Lambda$ by $C_{**}\Lambda$.

\emph{Which inequalities inherit the entry.} The un-hatted inequalities (P) and (R) are those of
the corollary above, transplanted across the cutoff. The inequality (A) is not among that
corollary's inequalities, since the corollary has inequalities for $P$ and $R$ only. It is
obtained instead from clause~(a) of the one-lattice Lipschitz step in its first form, read in
the logarithmic units $\mathsf{l}_r$, from clause~(D) of that step in its second form, and from
two further bounds, one of them belonging to the second form. Clause~(a) is a bound of the first
form in terms of the same load $[\![\delta]\!]_n$, so the un-decayed entry of (A) is again the
image of the member $\sum_{j\le n}\mathsf{w}_{n,j}\mathsf{r}_j$. The supply of the cross-cutoff
hypothesis on the remnant slot is read at the same load, written for the two cutoffs.

Hence every un-decayed entry of (P), (A), (R) arises from one term in each form. These are the
entries for which Proposition~\ref{9.11:prop-no-geometric-envelope} shows that no pure
geometric envelope exists. The replacement of the remnant member in $(\hat{R})$ substitutes the
second form's hatted kernel for the first form's image.

\emph{(ii).} There are three windows to close.

The window $\varkappa'_r\sum_{r<s<r+\ell_r}\|\hat{G}^R_s\|$ of $(\hat{R})$ has length $\ell_r$,
which is unbounded along the ladder $\theta_s$. By Remark~\ref{9.11:rem-window-kernel} it is the
linear band window of the second form, unchanged under the passage to the hatted rows. It closes
under \eqref{9.11:eq-coupling-ceilings-c}, which is the same condition as in the un-hatted
system; this is shown in the verification of the hatted rows at the envelope.

The band members $C_F\Sigma^{\mathrm{band}}_{s>r}\|\hat{G}^R_s\|$ of $(\hat{P})$ and
$C_F^A\Sigma^{\mathrm{band}}_{s>r}\|\hat{G}^R_s\|$ of $(\hat{A})$ close under
\eqref{9.11:eq-coupling-ceilings-d}. For $\kappa_0>4$ that condition holds for every
$X\ge\hat{X}_{\hat{h}}$ by Lemma~\ref{9.11:lem-ceilings-ladder}(iii), which derives it from the
bound on the waiting time $\check{R}_{(s)}$ of the large-field operation.

The $\mathsf{k}$-member of $(\hat{R})$ closes under \eqref{9.11:eq-coupling-ceilings-d} through
\eqref{9.11:eq-envelope-sums-d}: the bound \eqref{9.11:eq-envelope-sums-d} for the
$\mathsf{k}$-sum at the envelope carries a factor $\rho^{-1}$, which
\eqref{9.11:eq-coupling-ceilings-d} accommodates. Since
$\mathsf{k}_{r,s}=\rho^{2(s-r-\mathsf{a}_s-1)_+}$ and $\rho<1$, the set on which
$\mathsf{k}_{r,s}=1$ is $\{s:\ r<s\le r+\mathsf{a}_s+1\}$, and a growth bound on $\mathsf{a}_s$
gives
\begin{equation*}
\{s:\ r<s\le r+\mathsf{a}_s+1\}\subset\{s:\ r<s\le r+\tfrac32\ell_r\}.
\end{equation*}
So the flat part of the $\mathsf{k}$-member is a window of length at most $\tfrac32\ell_r$, of
the same order as the $\varkappa'$-window.

In each of these three cases the member is controlled by the product
\begin{equation*}
\theta_s^{1-\kappa_0/2}\,(M\check{R}_{(s)})^{\log_L(1/\rho)}(\log\theta_s)^{c_3},
\end{equation*}
whose second and third factors form a polylogarithm in $\theta_s$. Under the conditions on
$\kappa_0$ in Lemma~\ref{9.11:lem-ceilings-ladder}(ii) and~(iii) the exponent $1-\kappa_0/2$ is
negative, so the power of $\theta_s$ beats the polylogarithm; the thresholds
$\hat{X}_{\hat{e}}$ and $\hat{X}_{\hat{h}}$ of that lemma make this quantitative.

We add two observations. The alignment statement carrying the index $n_c$ involves no floor in
the form in which it is used. The quantities $\mathsf{n}_r$, $C_T$, $\Pi_n$ of the un-hatted
system do not occur in the hatted one: they drop out in the bound
\begin{equation*}
\sum_{u\le r}e_u\le e_r/(1-\rho)
\end{equation*}
for the geometric sum of the $e_u=\rho^{n-u}$, and in \eqref{9.11:eq-envelope-sums-d}.
\end{proof}

\begin{remark}[The boundary slot lies outside the three-block frame]\label{9.11:rem-boundary-slot}
Consider the frame in which the hatted response rows of Definition~\ref{9.11:def-response-rows}
and the hatted source families of Proposition~\ref{9.11:prop-source-families}, whose proof is
incomplete at one step, are written. Here
$\delta^{\mathbf{B}}_j := \|\delta\mathbf{B}^{(j)}\|$ is the norm of the two-cutoff difference of
the boundary family $\mathbf{B}^{(j)}$ on the common labels of the two runs, and $\mathbf{B}'$ is the
boundary term produced by the large-field operation $\mathbf{R}$.
\begin{itemize}
\item[(a)] The frame has exactly the three blocks $\hat{P}$, $\hat{A}$, $\hat{R}$
of \eqref{9.11:eq-hatted-blocks-a}, and no boundary slot.
\item[(b)] The row $(\hat{R})$ of \eqref{9.11:eq-response-rows-c} and the bound on
$\hat{N}_r^{R'}$ in \eqref{9.11:eq-source-families-a} contain no term in $\delta^{\mathbf{B}}_j$;
the row $(R')$ of the hatted form of the one-lattice Lipschitz step reads $\delta^{\mathbf{B}}_j$
through the kernel $\mathfrak{w}^{\mathbf{B}}_{k,j}$, and the boundary slot leaves the three-block
rows through the label-wise elimination (Remark~\ref{9.11:rem-window-kernel}), the display
\eqref{9.11:eq-response-rows-c} not exhibiting the kernel through which its image enters
$\hat{B}_r(\hat{G},d)$.
\item[(c)] The sampling members $\Upsilon$ carried by the rows of the hatted form of the
one-lattice Lipschitz step do not occur in \eqref{9.11:eq-response-rows-c}; no argument removing
them is displayed beyond the convention that the local slots are functions of one continuum
source $W$.
\item[(d)] For $0<q<1$ one has $2^{-m}\le q^m$ for every integer $m\ge0$ if and only if
$q\ge\frac12$. Definition~\ref{9.11:def-rate-floors} imposes on the kernel rate $q$ the floor
\eqref{9.11:eq-rate-floors-a} and no other lower bound.
\end{itemize}
\end{remark}

\begin{proof}
(a) The blocks of the frame are those of \eqref{9.11:eq-hatted-blocks-a}. Apart from coefficients
depending only on the depth and on the run couplings, the right members of the
rows $(\hat{P})$, $(\hat{A})$, $(\hat{R})$ of Definition~\ref{9.11:def-response-rows} involve
the competitor norms $\|\hat{G}_s^P\|$, $\|\hat{G}_s^A\|$, $\|\hat{G}_s^R\|$ and the distances
$d_{s+1}$, $d^{\sharp}_{r+1}$.
The same holds for the bounds on $\hat{N}_r^P$, $\hat{N}_r^A$, $\hat{N}_r^{R''}$, $\hat{N}_r^{R'}$ in
Proposition~\ref{9.11:prop-source-families}, whose only further members are $e_r=\rho^{n-r}$ and
$\hat{B}_r(\hat{G},d)$. No norm of a fourth block occurs in either. In particular neither the row
$(\hat{R})$ nor the bound on $\|\hat{N}_r^{R'}\|$ in \eqref{9.11:eq-source-families-a} contains a
term in $\delta^{\mathbf{B}}_j$.

In the list structure of the correlation theorem, among the derived objects of a list are, for
every label, the boundary families $\mathbf{B}^{(j)}$, the operations $\mathbf{T}_{k'}$ and
$\mathbf{T}'_{k'}(X)$, the activities of \cite[(1.91)]{B-CMP122b}, and $R'$, $\mathbf{B}'$ and
$R''$. Each of these is determined by the list, so in that structure the boundary terms carry no slot
of their own. For these derived terms the volume and filament bound has nothing to estimate; for the
term $\mathbf{B}'$ produced by $\mathbf{R}$ there is an estimate to prove. In the hatted form of the
one-lattice Lipschitz step, on the other hand, the boundary families $(\mathbf{B}^{(j)})_{j\le k}$
form a slot of the list, with two-cutoff differences
$\delta^{\mathbf{B}}=(\delta^{\mathbf{B}}_j)_{j\le k}$; item (b) describes how this slot leaves the
three-block frame. In the rows without hats the two-cutoff difference of the boundary terms is
routed into the component $R'$ of the source family; its hatted counterpart is the component
estimated in \eqref{9.11:eq-source-families-a} of Proposition~\ref{9.11:prop-source-families}, which
is stated with its proof incomplete at one step.

(b) The first half of (b) was shown in the first paragraph of (a).
In the hatted form of the one-lattice Lipschitz step the slot $\delta^{\mathbf{B}}$ is read in
three places:
\begin{itemize}
\item in the row $(R')$, through the member $\mathfrak{w}^{\mathbf{B}}_{k,j}\,\delta^{\mathbf{B}}_j$;
\item in the row of the operations $\mathbf{T}$, through its boundary entry;
\item in the boundary row itself, through the first bracket
$\sum_{j\le k}2^{-(k-j)}\delta^{\mathbf{B}}_j$.
\end{itemize}
The left member of the row $(R')$ is
\begin{equation}\label{9.11:eq-boundary-slot-1}
\|\delta R'^{(k+1)}\|^{\flat}+\|\delta\mathbf{B}'^{(k+1)}\| .
\end{equation}
Its second summand is the norm of the two-cutoff difference of the term $\mathbf{B}'$ produced by
$\mathbf{R}$. Read
$\mathbf{B}^{(k+1)}$ as produced together with $\mathbf{B}'$; then the slot difference
$\delta^{\mathbf{B}}_{k+1}$ contains this summand. It is $\delta^{\mathbf{B}}$ that the label-wise
elimination removes from the three-block frame. The kernels displayed in $\hat{B}_r(\hat{G},d)$ are
$q^{s-r}$ for the members of the $P$-, $A$- and $\zeta$-type, and
$\mathsf{k}_{r,s}=\rho^{2(s-r-\mathsf{a}_s-1)_+}$ for the $\|\hat{G}^R_s\|$-member. The display
\eqref{9.11:eq-response-rows-c} does not state which of these receives the image of the eliminated
slot.

(c) Among the rows of the hatted form of the one-lattice Lipschitz step that carry a member
proportional to $\Upsilon$ are:
\begin{itemize}
\item the row $(R')$;
\item the defect row $(\hat{D})$, through a constant multiple of $E_0\Upsilon$;
\item the row of the operations $\mathbf{T}$, through a constant multiple of $\Upsilon$.
\end{itemize}
None of these members appears in \eqref{9.11:eq-response-rows-c}. Under the single-source
convention the local slots of runs, competitors and hybrids are functions of the one continuum
source $W$, no difference of two sources is formed, and sampling is carried in the weights only; no
argument removing the members in $\Upsilon$ from the rows is displayed beyond this convention.

(d) Suppose first that $q\ge\frac12$. For every integer $m\ge0$ we then have
$q^m\ge(\frac12)^m=2^{-m}$. Conversely, suppose the inequality holds for every $m\ge0$. Taking $m=1$
gives $\frac12\le q$.

By Definition~\ref{9.11:def-rate-floors}, the floor \eqref{9.11:eq-rate-floors-a} places $q$ in the
open interval with left end $\max\{L^{-\hat{\alpha}},L^{-\beta_1},L^{-1/2}\}$ and right end $1$
and requires $q_a\le q^2$;
the floor on $\rho$ of that definition bounds $\rho$, not $q$. The condition $q\ge\frac12$ is
therefore not among the
conditions imposed on $q$. A condition of that form would be one-sided and would not involve $L$.

We now compare this with the kernels of (b). The kernel of the row $(R')$ satisfies
\begin{equation}\label{9.11:eq-boundary-slot-2}
\mathfrak{w}^{\mathbf{B}}_{k,j}=\sum_{j\le i\le k}2^{-(i-j)}\vartheta_*^{\,k-i}\ \ge\ 2^{-(k-j)} .
\end{equation}
The inequality holds because every summand is non-negative and the sum contains the summand $i=k$,
which equals $2^{-(k-j)}$. As written, this kernel has base $2$ in the age $k-j$. The volume part of
the volume and filament bound gives for this member the kernel
$C_{\mathbf{B}}^v(k-j+2)\,\vartheta_*^{\,k-j}$, which is geometric in
the age and is at most $2C_{\mathbf{B}}^v\,\mathfrak{w}^{\mathbf{B}}_{k,j}$ when
$\vartheta_*\le\frac14$: indeed, for $j\le i\le k$ the first line below holds, and summing it over
the $k-j+1$ values of $i$ gives the second,
\begin{align*}
2^{-(i-j)}\vartheta_*^{\,k-i}&=\vartheta_*^{\,k-j}\bigl(2\vartheta_*\bigr)^{-(i-j)}
\ge 2^{\,i-j}\,\vartheta_*^{\,k-j}\ge\vartheta_*^{\,k-j},\\
\mathfrak{w}^{\mathbf{B}}_{k,j}&\ge(k-j+1)\vartheta_*^{\,k-j}\ge\tfrac12(k-j+2)\vartheta_*^{\,k-j}.
\end{align*}
By the filament part of the same bound, the boundary entry of the row of the
operations $\mathbf{T}$ is of the same
kind, with age factor $\vartheta_*^{\,k-j}$. The one displayed kernel of base $2$ that is not so
sharpened is the first bracket $\sum_{j\le k}2^{-(k-j)}\delta^{\mathbf{B}}_j$ of the boundary row.

There are two ways the elimination can act, and they lead to different conclusions.
\begin{itemize}
\item Suppose the elimination iterates the boundary row. Then members with kernel $2^{-(k-j)}$ reach
$\hat{B}_r(\hat{G},d)$ as kernels $2^{-m}$ in an exponent $m$ that the displayed kernels carry as
$q^{m}$; by the equivalence of the first paragraph of (d), they are dominated by $q^{m}$ for all $m$
only if $q\ge\frac12$. Under that floor they are dominated by $\mathsf{k}_{r,s}$ as well:
$q\le\rho^2$ by Definition~\ref{9.11:def-rate-floors}, $0<\rho<1$ and
$(s-r-\mathsf{a}_s-1)_+\le s-r$, so that $q^{s-r}\le\rho^{2(s-r)}\le\mathsf{k}_{r,s}$. Domination
by either kernel thus rests on a floor of this kind on the rates.
\item Suppose instead the elimination bounds $\delta^{\mathbf{B}}$ through the right member of the
row $(R')$. The kernel that the volume part of the volume and filament bound gives for the boundary
term on the right of the row $(R')$ is geometric in the age with ratio
$\vartheta_*$. Here $\vartheta_*\le q_a$, as in the derivation of the hatted kernel of the row
$(R')$ in the hatted form of the one-lattice Lipschitz step, and $q_a\le q^2$ by
\eqref{9.11:eq-rate-floors-a}, so this kernel is dominated by the displayed kernels of
$\hat{B}_r(\hat{G},d)$ and no member of base $2$ reaches it.
\end{itemize}
Which of the two the label-wise elimination performs is not displayed beside
\eqref{9.11:eq-response-rows-c}.

For the response rows without hats, before re-filing, the question does not arise: there the
$\|G^R_s\|$-member of $B_r$ carries a flat kernel, and $2^{-(k-j)}\le1$. A statement that takes
$(\hat{R})$ as a hypothesis reads only the kernels
displayed in \eqref{9.11:eq-response-rows-c}.
\end{proof}

\begin{corollary}[Geometric closeness of consecutive tuned hatted couplings]
\label{9.11:cor-consecutive-couplings}
Assume the hypotheses of Theorem~\ref{9.11:thm-geometric-envelope},
the geometric envelope for two-run differences of hatted runs, with the constants $\mathsf{A}$,
$C_E$ and the rate $\rho$ as fixed there; let $n$ be a run length they admit, and let
$e_r=\rho^{n-r}$. Let $\mathsf{A}$ and $\mathsf{B}$ be the runs with hatted lists in the sense
of Definition~\ref{9.11:def-hatted-lists} (hatted runs), of
lengths $n$ and $n+1$ that form the \emph{tuned pair}, that is, the pair with starting depth
$r_0=0$ taken at the value $v=0$ of the pair's parameter $v$. Here $r_0$ denotes the starting
depth of a pair of hatted runs as in Theorem~\ref{9.11:thm-geometric-envelope}. The letter
$\mathsf{A}$ names the run of length $n$ when it stands for a run or appears as a superscript; as
a factor in a bound it is the constant of Theorem~\ref{9.11:thm-geometric-envelope}. Write
$\hat{x}^{(n)}_{(r)}$ and
$\hat{x}^{(n+1)}_{(r)}$ for their hatted running couplings at depth $r$,
$\hat{\zeta}^{\mathsf{A}}_{(r)}$
and $\hat{\zeta}^{\mathsf{B}}_{(r)}$ for their hatted run couplings at depth $r$, and $d_r$ for
the coupling distance of the pair at depth $r$. Assume in addition:
\begin{itemize}
\item[(a)] the tuned hatted pair exists, that is, both hatted runs are defined in the sense of
the existence statement for hatted runs, and each satisfies the three run conditions that the
theorem on the hatted scaffold imposes on hatted runs, the first of which is the coefficient
bound with the constant $\mathsf{A}$;
\item[(b)] the pinning lemma, part \textup{(i)}, holds for the pair: the pinning depth and the
matching depth of the pair are both $0$;
\item[(c)] on the common ladder $\theta_s$, for every depth $0\le r\le n$,
\begin{equation}\label{9.11:eq-consecutive-couplings-1}
\bigl|\hat{x}^{(n+1)}_{(r)}-\hat{x}^{(n)}_{(r)}\bigr|
=\bigl|\hat{\zeta}^{\mathsf{B}}_{(r)}-\hat{\zeta}^{\mathsf{A}}_{(r)}\bigr|\Big|_{v=0}\le d_r .
\end{equation}
\end{itemize}
Then for every depth $0\le r\le n$
\begin{equation}\label{9.11:eq-consecutive-couplings-2}
\bigl|\hat{x}^{(n+1)}_{(r)}-\hat{x}^{(n)}_{(r)}\bigr|\le
\frac{\mathsf{A}\,C_E\,\rho^{\,n-r+1}}{1-\rho}.
\end{equation}
\end{corollary}

\begin{proof}
By hypothesis (b), the pinning depth and the matching depth of the tuned pair are both $0$. The
pair is therefore admissible in Theorem~\ref{9.11:thm-geometric-envelope} with starting depth
$r_0=0$. Since hypothesis (a) makes both runs of the pair defined, that theorem applies to them,
and its bound \eqref{9.11:eq-geometric-envelope-a} on the coupling distance is available at every
depth at which that theorem states it, in particular at depth $r-1$ for $1\le r\le n$.

Fix a depth $r$ with $0\le r\le n$. By hypothesis (c), that is, by
\eqref{9.11:eq-consecutive-couplings-1}, the left-hand side of
\eqref{9.11:eq-consecutive-couplings-2} equals
$\bigl|\hat{\zeta}^{\mathsf{B}}_{(r)}-\hat{\zeta}^{\mathsf{A}}_{(r)}\bigr|$ at $v=0$ and is at most
$d_r$. It thus remains to bound $d_r$, and we distinguish two cases.

Let first $1\le r\le n$, so that $0\le r-1\le n-1$. The bound on the coupling
distance in \eqref{9.11:eq-geometric-envelope-a}, taken at depth $r-1$, where
$e_{r-1}=\rho^{\,n-(r-1)}$, gives
\begin{equation*}
d_r\le\frac{\mathsf{A}\,C_E\,\rho^{\,n-(r-1)}}{1-\rho}
=\frac{\mathsf{A}\,C_E\,\rho^{\,n-r+1}}{1-\rho}.
\end{equation*}
Together with \eqref{9.11:eq-consecutive-couplings-1} this gives
\eqref{9.11:eq-consecutive-couplings-2}.

Let now $r=0$. Then $d_0=0$, and the right-hand side of \eqref{9.11:eq-consecutive-couplings-2}
is non-negative, so \eqref{9.11:eq-consecutive-couplings-2} holds at $r=0$ as well.
\end{proof}

Put $\tau_n(t):=\mathsf{A}C_E\rho^{\,n-t+1}/(1-\rho)$ for $0\le t\le n$; here $t$ denotes a
depth. The corollary bounds the difference of the hatted running couplings
$\hat{x}^{(n+1)}_{(t)}-\hat{x}^{(n)}_{(t)}$ of the consecutive runs at depth $t$ by $\tau_n(t)$.
At fixed $t$ this bound decays geometrically in $n$, and
\begin{equation*}
\sum_{n\ge t}\tau_n(t)=\frac{\mathsf{A}\,C_E}{1-\rho}\sum_{n\ge t}\rho^{\,n-t+1}
=\frac{\mathsf{A}\,C_E\,\rho}{(1-\rho)^2}.
\end{equation*}

Consider the modulus of the calibration proposition,
\begin{equation*}
\Upsilon_n=2a_n+\sum_{u<n}\min\Bigl\{3a_u,\ \frac{2\mathsf{A}\,C_E\,e_u}{r_J}\Bigr\}.
\end{equation*}
With $e_u=\rho^{\,n-u}$ this modulus keeps the order $n^{-3/2}(\log n)^{p_0+1}$ it has in the
calibration proposition, since that order is set by the $a_u$ and not by the $e_u$.

\begin{definition}[Standing hypotheses: the tuned family of the uniqueness theorem]
\label{9.12:not-standing-hypotheses}
Throughout this section the setting and the hypotheses are those of the uniqueness theorem:
\begin{itemize}
\item the gauge group is $\mathrm{SU}(2)$;
\item the numbers $\theta_A$, $\beta_I$, $\alpha$, $\beta_H$, $\beta^{\flat}$, $\theta_1$, $\theta_2$,
$\theta'$ and $\mathfrak{i}$ are fixed as in the setting of that theorem;
\item $L\ge L_0$ is odd, and the auxiliary parameters $\mathfrak{p}$ are as there;
\item $r_J>0$ is arbitrary, and we put
\begin{equation*}
R:=4r_J,\qquad \varrho_2:=\frac{r_J}{2},\qquad t_{\circ}:=c_2+r_J ;
\end{equation*}
here $R$ is an abbreviation of the setting of the uniqueness theorem and is distinct from the
ball radius $R$ of the glossary;
\item $g\le\gamma_U$, and $X:=g^{-2}$;
\item the hypotheses \textup{(i)}--\textup{(iv)} and $(\mathrm{J}1)\text{--}(\mathrm{J}5)$ of the
uniqueness theorem hold.
\end{itemize}
All of these are carried unchanged; none of them is discharged in this section. The scaffold
$\Theta_X$, the run of length $n$, the tuned bare shift $\zeta^n_*$, the tuned bare inverse coupling
$a_n$, the function $\varphi_n$, the transition map $\mathfrak{n}_n$, the defect bound
$\Upsilon_n$ and the number $N_n$ are defined later in this section.
\end{definition}

\begin{definition}[Runs on the scaffold and the reference-scale coupling map]
\label{9.12:def-runs-coupling-map}
Throughout, the setting and the hypotheses are those of
Definition~\ref{9.12:not-standing-hypotheses}. In particular $X=g^{-2}$, $r_J>0$ is arbitrary,
$\varrho_2:=r_J/2$, and, for $a>0$, $D(0,a)$ and $\overline{D}(0,a)$ denote the open and the
closed disc of radius $a$ about $0$ in $\mathbb{C}$.
\begin{enumerate}
\item[(a)] \emph{The scaffold.} $\Theta_X=(\theta_r)_{r\ge0}$ is the scaffold of real numbers
of the uniqueness theorem. In the convention of that theorem, a run whose shift at depth $r$
is $\zeta_{(r)}$ has inverse coupling $\theta_r-\zeta_{(r)}$ at depth $r$. In this convention
the shift at depth $0$ equals $\hat w$ exactly when the inverse coupling at depth $0$ equals
$X-\hat w$. Hence $\theta_0=X$.
\item[(b)] \emph{The run of length $n$.} Fix $n\ge1$. Let $\zeta$ be a complex number, the
\emph{bare shift}. The \emph{run of length $n$ on $\Theta_X$ from $\zeta$} is the finite
sequence $(\zeta^{(n)}_{(r)}(\zeta))_{0\le r\le n}$ defined by
\begin{equation}\label{9.12:eq-runs-coupling-map-a}
\zeta^{(n)}_{(n)}(\zeta):=\zeta,\qquad
\zeta^{(n)}_{(r)}(\zeta):=\zeta^{(n)}_{(r+1)}(\zeta)+b^{(n)}_{(r)}(\zeta)-c_0
\quad(r=n-1,\dots,0).
\end{equation}
Here $b^{(n)}_{(r)}$ is the step function of the running-shift statement of the correlation
theorem, read on the scaffold at depth $r$ of a run of length $n$; it is a function of the
bare shift $\zeta$. In each step of \eqref{9.12:eq-runs-coupling-map-a} the increment
$\zeta^{(n)}_{(r)}(\zeta)-\zeta^{(n)}_{(r+1)}(\zeta)$ depends on $\zeta$ only through
$b^{(n)}_{(r)}$. The number $c_0$ is independent of $\zeta$; it is
local to this definition and is not the constant $c_0$ of the glossary in
Definition~\ref{0.2:not-glossary} (the constant entering $\underline{\lambda}=c_0/4$). In the
form of the recursion given by the running-shift statement of the correlation theorem, the
corresponding number is $b_{k+1}(0)$. Only the independence of $c_0$ from $\zeta$ is used
below. Two $\zeta$-independent centrings change every $\zeta^{(n)}_{(r)}$ by an amount
independent of $\zeta$, so they give the same differences between runs from two bare shifts.

The \emph{inverse coupling at depth $r$} of this run is
$x^{(n)}_{(r)}(\zeta):=\theta_r-\zeta^{(n)}_{(r)}(\zeta)$. The bare inverse coupling is
$x^{(n)}_{(n)}(\zeta)=\theta_n-\zeta$. The reference-scale (depth-$0$) inverse coupling is
$x^{(n)}_{(0)}(\zeta)=X-\zeta^{(n)}_{(0)}(\zeta)$. For $0\le r\le n$ one has
\begin{equation}\label{9.12:eq-runs-coupling-map-b}
\begin{gathered}
\zeta^{(n)}_{(r)}(\zeta)=\zeta+\sum_{s=r}^{n-1}\bigl(b^{(n)}_{(s)}(\zeta)-c_0\bigr),\\
\text{in particular}\quad
\zeta^{(n)}_{(0)}(\zeta)=\zeta+\sum_{r<n}\bigl(b^{(n)}_{(r)}(\zeta)-c_0\bigr).
\end{gathered}
\end{equation}
\item[(c)] \emph{The tuned centre and the map $\varphi_n$.} The tuning statement asserts that
for every real $\hat w$ with $|\hat w|\le t_{\circ}-c_2$ there is exactly one real bare shift
whose run of length $n$ is alive to depth $0$, in the sense of the uniqueness theorem, and
has shift $\hat w$ at depth $0$. Let
$\zeta^n_*$ be this real bare shift for $\hat w=0$, so that
$\zeta^{(n)}_{(0)}(\zeta^n_*)=0$. The setting of the uniqueness theorem puts
$a_n=\theta_n-\zeta^n_*$, the tuned bare inverse coupling. The run of length $n$ with bare
inverse coupling
$x^{(n)}_{(n)}=a_n-\hat w$ has bare shift
\begin{equation*}
\zeta=\theta_n-(a_n-\hat w)=\zeta^n_*+\hat w .
\end{equation*}
For $\hat w\in D(0,2r_J)$ define
\begin{equation}\label{9.12:eq-runs-coupling-map-c}
\varphi_n(\hat w):=x^{(n)}_{(0)}(\zeta^n_*+\hat w)=X-\zeta^{(n)}_{(0)}(\zeta^n_*+\hat w).
\end{equation}
Thus $\varphi_n(\hat w)$ is the reference-scale inverse coupling of the run of length $n$
started at bare inverse coupling $a_n-\hat w$. The bare data $(n,a_n-\hat w)$, $n\ge1$,
$\hat w\in J$ (with $J$ as in (d)), form the tuned family of the uniqueness theorem recalled
in Definition~\ref{9.12:not-standing-hypotheses}. For $0\le r<n$ put
$\beta^{(n)}_r(\eta):=b^{(n)}_{(r)}(\zeta^n_*+\eta)$. Then
\begin{equation}\label{9.12:eq-runs-coupling-map-d}
\varphi_n(\hat w)=X-\zeta^n_*-\hat w-\sum_{r<n}\bigl(\beta^{(n)}_r(\hat w)-c_0\bigr).
\end{equation}
\item[(d)] \emph{Transition maps, defects and the interval $J$.} For each $n$, the conclusion
of the uniqueness theorem on the transition maps between consecutive depths provides three
objects:
\begin{itemize}
\item[--] the transition map $\mathfrak{n}_n$ between the runs of length $n$ and $n+1$;
\item[--] the defect bound $\Upsilon_n$, with
$|\mathfrak{n}_n(\hat w)-\hat w|\le\Upsilon_n|\hat w|$ for $\hat w\in D(0,r_J)$, the domain
of $\mathfrak{n}_n$;
\item[--] the number $N_n:=\partial\zeta^{(n)}_{(0)}(\zeta)/\partial\zeta$ at
$\zeta=\zeta^n_*$; equivalently $N_n=-\varphi_n'(0)$ by
\eqref{9.12:eq-runs-coupling-map-c}.
\end{itemize}
These are the objects so denoted here. Put
\begin{equation*}
J:=[-\varrho_2/2,\varrho_2/2]=[-r_J/4,r_J/4].
\end{equation*}
Then $J\subset\overline{D}(0,r_J)\subset D(0,2r_J)$. For $m\ge1$ put
$T_m:=\sum_{n\ge m}\Upsilon_n$. This sum is finite by the summability of the transition
defects, \eqref{9.12:eq-defect-summability-a}.
\end{enumerate}
\end{definition}

\begin{proof}[Proof of \eqref{9.12:eq-runs-coupling-map-b} and \eqref{9.12:eq-runs-coupling-map-d}]
Fix $n\ge1$, a bare shift $\zeta$ and $r$ with $0\le r\le n$.

For each $s$ with $r\le s\le n-1$, the recursion \eqref{9.12:eq-runs-coupling-map-a}, taken at
index $s$, gives
\begin{equation*}
\zeta^{(n)}_{(s)}(\zeta)-\zeta^{(n)}_{(s+1)}(\zeta)=b^{(n)}_{(s)}(\zeta)-c_0 .
\end{equation*}
Summing these identities over $s=n-1$ down to $s=r$, the left-hand sides telescope to
$\zeta^{(n)}_{(r)}(\zeta)-\zeta^{(n)}_{(n)}(\zeta)$. By the first half of
\eqref{9.12:eq-runs-coupling-map-a}, $\zeta^{(n)}_{(n)}(\zeta)=\zeta$. This gives the first
identity of \eqref{9.12:eq-runs-coupling-map-b}; for $r=n$ the sum is empty and the identity
is the initial condition. Taking $r=0$ gives the second identity.

For the expressions for the inverse couplings in (b), note that, by the initial condition and
because $\theta_0=X$ by (a),
\begin{align*}
x^{(n)}_{(n)}(\zeta)&=\theta_n-\zeta^{(n)}_{(n)}(\zeta)=\theta_n-\zeta,\\
x^{(n)}_{(0)}(\zeta)&=\theta_0-\zeta^{(n)}_{(0)}(\zeta)=X-\zeta^{(n)}_{(0)}(\zeta).
\end{align*}
The second equality in \eqref{9.12:eq-runs-coupling-map-c} is this
last identity at the bare shift $\zeta^n_*+\hat w$.

For \eqref{9.12:eq-runs-coupling-map-d}, apply the second identity of
\eqref{9.12:eq-runs-coupling-map-b} at the bare shift $\zeta^n_*+\hat w$:
\begin{align*}
\zeta^{(n)}_{(0)}(\zeta^n_*+\hat w)
&=\zeta^n_*+\hat w+\sum_{r<n}\bigl(b^{(n)}_{(r)}(\zeta^n_*+\hat w)-c_0\bigr)\\
&=\zeta^n_*+\hat w+\sum_{r<n}\bigl(\beta^{(n)}_r(\hat w)-c_0\bigr),
\end{align*}
where the second equality is the definition of $\beta^{(n)}_r$. Substituting this into
\eqref{9.12:eq-runs-coupling-map-c} gives \eqref{9.12:eq-runs-coupling-map-d}.
\end{proof}

\begin{lemma}[Summability of the transition defects and the depth floor]
\label{9.12:lem-defect-summability}
Assume the defect-bound conclusion of the uniqueness theorem. In particular, for every $n\ge 1$ the transition
map $\mathfrak{n}_n$ between depths $n$ and $n+1$ satisfies
$|\mathfrak{n}_n(\hat w)-\hat w|\le \Upsilon_n|\hat w|$ for every $\hat w$ in the disc, of positive radius,
on which the uniqueness theorem states it, and
\begin{equation}\label{9.12:eq-defect-summability-1}
  \Upsilon_n\le C\,n^{-3/2}(\log n)^{p_0+1}\qquad (n\ge 1),
\end{equation}
with constants $C$ and $p_0$ independent of $n$. Then
\begin{equation}\label{9.12:eq-defect-summability-a}
\begin{gathered}
  \sum_{n\ge1}\Upsilon_n<\infty,\qquad \Upsilon_n\to 0,\qquad\text{and the integer}\\
  n_1:=\min\{\,m\ge 1:\ \Upsilon_n\le 3\ \text{for every}\ n\ge m\,\}\ \text{exists.}
\end{gathered}
\end{equation}
\end{lemma}

\begin{proof}
Each $\Upsilon_n$ is non-negative, being an upper bound for the quotient
$|\mathfrak{n}_n(\hat w)-\hat w|/|\hat w|$ at any $\hat w\neq0$ where the defect bound holds.

Next we compare with a power. Whatever the value of $p_0$, one has $(\log n)^{p_0+1}n^{-1/4}\to0$ as
$n\to\infty$.
Hence there is an integer $n'$ such that $(\log n)^{p_0+1}\le n^{1/4}$ for every $n\ge n'$. Inserting
this into \eqref{9.12:eq-defect-summability-1} gives, for every $n\ge n'$,
\begin{equation*}
  \Upsilon_n\le C\,n^{-3/2}\,n^{1/4}=C\,n^{-5/4}.
\end{equation*}

We now sum. The series $\sum_{n\ge1}n^{-5/4}$ converges, because its exponent exceeds $1$. The terms
$\Upsilon_n$ are non-negative. By the comparison test, $\sum_{n\ge n'}\Upsilon_n$ is therefore at most
$C\sum_{n\ge n'}n^{-5/4}<\infty$. Adding the finitely many terms with $n<n'$ gives
$\sum_{n\ge1}\Upsilon_n<\infty$.

The terms of a convergent series tend to zero, so $\Upsilon_n\to0$. Consequently there is $m\ge1$ with
$\Upsilon_n\le3$ for every $n\ge m$. The set of such $m$ in \eqref{9.12:eq-defect-summability-a} is
therefore a non-empty set of positive integers. It has a least element, and this least element is $n_1$.
\end{proof}

\begin{lemma}[The bi-Lipschitz lemma]\label{9.12:lem-bi-lipschitz}
Assume the standing hypotheses of Notation~\ref{9.12:not-standing-hypotheses}. Let the scaffold
$\Theta_X=(\theta_r)_{r\ge0}$, the runs of length $n$ on it with step functions $b^{(n)}_{(r)}$,
the tuned bare shift $\zeta^n_*$ (whose existence is hypothesis (B) below) and the tuned bare
inverse coupling $a_n=\theta_n-\zeta^n_*$ be as in Definition~\ref{9.12:def-runs-coupling-map}. Let
$r_J>0$ be arbitrary and put $\varrho_2:=r_J/2$.
For $\rho>0$ let $D(0,\rho)$ and $\overline{D}(0,\rho)$ be the open and the closed disc in
$\mathbb{C}$ of centre $0$ and radius $\rho$; in this lemma and its proof, $\rho$ and $\varrho$
denote variable radii, not the fixed quantities of the glossary
(Notation~\ref{0.2:not-glossary}). For a set $E\subset\mathbb{C}$ and a
function $h$ on
$E$ let $\operatorname{Lip}_E h$ be the Lipschitz constant of $h$ on $E$. For $n\ge1$ and
$0\le r<n$ put
\begin{equation*}
\beta^{(n)}_r(\eta):=b^{(n)}_{(r)}(\zeta^n_*+\eta).
\end{equation*}
Assume the following.
\begin{itemize}
\item[(A)] The running-shift statement of the correlation theorem holds when read on the scaffold
$\Theta_X$ about the tuned centre with radius $r_J$, in the following form. For every $n\ge1$:
\begin{equation}\label{9.12:eq-bi-lipschitz-a}
\begin{aligned}
&\text{(R1)}\quad \text{for each } 0\le r<n,\ \beta^{(n)}_r \text{ is holomorphic on } D(0,2r_J)\\
&\qquad\quad \text{and real on } \mathopen]-2r_J,2r_J\mathclose[;\\
&\text{(R2)}\quad \text{for every } \varrho\in[r_J,2r_J] \text{ there are numbers }
d^{(n)}_r(\varrho)\ge0,\ 0\le r<n, \text{ with}\\
&\qquad\quad |\beta^{(n)}_r(\eta)-\beta^{(n)}_r(0)|\le d^{(n)}_r(\varrho)\,|\eta|
\text{ for all } \eta\in D(0,\varrho) \text{ and all } r<n,\\
&\qquad\quad \text{and } \sum_{r<n} d^{(n)}_r(\varrho)\le\tfrac13;\\
&\text{(R3)}\quad \sum_{r<n}\operatorname{Lip}_{\overline{D}(0,r_J)}\beta^{(n)}_r\le\tfrac23 .
\end{aligned}
\end{equation}
\item[(B)] The statement on the existence of the tuned centre, on which the uniqueness theorem
rests to define $a_n$, holds: for every $n\ge1$ there is a real number $\zeta^n_*$ with
$\zeta^{(n)}_{(0)}(\zeta^n_*)=0$, and $a_n=\theta_n-\zeta^n_*$.
\item[(C)] The conclusion of the uniqueness theorem on the transition maps holds: for every
$n\ge1$ there are a map $\mathfrak{n}_n$, holomorphic on $D(0,2\varrho_2)=D(0,r_J)$, real on
$\mathopen]-r_J,r_J\mathclose[$, with $\mathfrak{n}_n(0)=0$, and a defect bound $\Upsilon_n$, such
that for $\hat w\in D(0,r_J)$ the run of length $n+1$ started at the bare inverse coupling
$a_{n+1}-\mathfrak{n}_n(\hat w)$ and the run of length $n$ started at $a_n-\hat w$ have the same
reference-scale inverse coupling, and
\begin{equation*}
|\mathfrak{n}_n(\hat w)-\hat w|\le\Upsilon_n|\hat w|,
\qquad
\Upsilon_n\le C\,n^{-3/2}(\log n)^{p_0+1},
\end{equation*}
where $C$ is a constant independent of $n$ and $p_0$ is the fixed logarithmic exponent of this
defect bound, likewise independent of $n$; moreover the number $N_n$, the derivative
$\partial\zeta^{(n)}_{(0)}/\partial\zeta$ at $\zeta=\zeta^n_*$, satisfies
$N_n\in[\frac23,\frac43]$ and $N_{n+1}\mathfrak{n}_n'(0)=N_n$, and the sequence $(N_n)$
converges.
\item[(D)] The torus-independence conclusion of the uniqueness theorem holds: $a_n$,
$\mathfrak{n}_n$, $N_n$ and $\Upsilon_n$ do not depend on the torus; and its hypothesis on the
data holds: the scaffold and the step functions of its runs do not depend on the torus.
\end{itemize}
For every $n\ge1$ let $\varphi_n:D(0,2r_J)\to\mathbb{C}$ be the reference-scale (depth-$0$) inverse
coupling of the run of length $n$ on $\Theta_X$ started at the bare inverse coupling
$a_n-\hat w$, as defined in~\eqref{9.12:eq-runs-coupling-map-c}. Then the following hold.
\begin{itemize}
\item[(i)] (Analyticity; one-point and two-point bounds, every $n\ge1$.) $\varphi_n$ is
holomorphic on $D(0,2r_J)$, real on $\mathopen]-2r_J,2r_J\mathclose[$, and $\varphi_n(0)=X$. For
every $\hat w\in D(0,2r_J)$,
\begin{equation}\label{9.12:eq-bi-lipschitz-b}
\tfrac23|\hat w|\le|\varphi_n(\hat w)-X|\le\tfrac43|\hat w|,
\qquad\text{indeed}\qquad |\varphi_n(\hat w)-X+\hat w|\le\tfrac13|\hat w| .
\end{equation}
For every $\hat w,\hat w'\in\overline{D}(0,r_J)$,
\begin{equation}\label{9.12:eq-bi-lipschitz-c}
\tfrac13|\hat w-\hat w'|\le|\varphi_n(\hat w)-\varphi_n(\hat w')|\le\tfrac53|\hat w-\hat w'|,
\end{equation}
and for real $-r_J\le\hat w'<\hat w\le r_J$,
\begin{equation}\label{9.12:eq-bi-lipschitz-d}
-\tfrac53(\hat w-\hat w')\le\varphi_n(\hat w)-\varphi_n(\hat w')\le-\tfrac13(\hat w-\hat w');
\end{equation}
in particular $\varphi_n$ is strictly decreasing on $[-r_J,r_J]$. Moreover
$-\varphi_n'(0)=N_n\in[\frac23,\frac43]$. The constants $\frac13$, $\frac23$, $\frac43$,
$\frac53$ are free of $n$, of the torus, and of $L$ and $g$ within the standing hypotheses of
Notation~\ref{9.12:not-standing-hypotheses}.
\item[(ii)] (Inverse maps.) Put $J:=[-\varrho_2/2,\varrho_2/2]$ and $I_n:=\varphi_n(J)$. Then
$I_n$ is a compact interval with
\begin{equation}\label{9.12:eq-bi-lipschitz-e}
[X-\varrho_2/3,\,X+\varrho_2/3]\subset I_n\subset[X-2\varrho_2/3,\,X+2\varrho_2/3],
\end{equation}
and the inverse $\psi_n:=(\varphi_n|_J)^{-1}:I_n\to J$ satisfies, for all $c,c'\in I_n$,
\begin{equation}\label{9.12:eq-bi-lipschitz-f}
\tfrac35|c-c'|\le|\psi_n(c)-\psi_n(c')|\le3|c-c'| .
\end{equation}
\item[(iii)] (Depth-to-depth consistency; uniform convergence.) Let $n_1$ be the integer of
Lemma~\ref{9.12:lem-defect-summability}, so that $\Upsilon_n\le3$ for $n\ge n_1$. For every
$n\ge n_1$ and every real $\hat w\in J$,
\begin{gather}
\varphi_{n+1}(\mathfrak{n}_n(\hat w))=\varphi_n(\hat w),\label{9.12:eq-bi-lipschitz-g}\\
|\varphi_{n+1}(\hat w)-\varphi_n(\hat w)|\le\tfrac53\,\Upsilon_n\,|\hat w| .
\label{9.12:eq-bi-lipschitz-h}
\end{gather}
Moreover $\sum_n\Upsilon_n<\infty$; writing $T_m:=\sum_{n\ge m}\Upsilon_n$, the sequence
$(\varphi_n)$ converges uniformly on $J$ to a continuous function $\varphi_\infty$ with
\begin{equation}\label{9.12:eq-bi-lipschitz-i}
\sup_{\hat w\in J}|\varphi_m(\hat w)-\varphi_\infty(\hat w)|\le\tfrac{5\varrho_2}{6}\,T_m
\qquad\text{for every } m\ge n_1,
\end{equation}
and $\varphi_\infty(0)=X$, $\varphi_\infty$ satisfies~\eqref{9.12:eq-bi-lipschitz-b} on $J$ and
\eqref{9.12:eq-bi-lipschitz-c}--\eqref{9.12:eq-bi-lipschitz-d} on $J$, $\varphi_\infty$ is strictly
decreasing, $I_\infty:=\varphi_\infty(J)$ satisfies~\eqref{9.12:eq-bi-lipschitz-e}, and
$\psi_\infty:=(\varphi_\infty|_J)^{-1}$ satisfies~\eqref{9.12:eq-bi-lipschitz-f}.
\end{itemize}
\end{lemma}

\begin{proof}[Proof of part \textup{(i)}]
Fix $n\ge1$. Part (i) uses hypotheses (A)--(D) only and no condition on $n$ beyond $n\ge1$;
Lemma~\ref{9.12:lem-defect-summability} does not enter.

\emph{Normalisation.} By hypothesis (B), $\zeta^{(n)}_{(0)}(\zeta^n_*)=0$. Hence the definition
\eqref{9.12:eq-runs-coupling-map-c} of $\varphi_n$, evaluated at $\hat w=0$, gives
$\varphi_n(0)=X-\zeta^{(n)}_{(0)}(\zeta^n_*)=X$. Evaluating the representation
\eqref{9.12:eq-runs-coupling-map-d} at $\hat w=0$ and inserting $\varphi_n(0)=X$, we obtain
\begin{equation*}
X=X-\zeta^n_*-\sum_{r<n}\bigl(\beta^{(n)}_r(0)-c_0\bigr),
\qquad\text{that is,}\qquad
\sum_{r<n}\bigl(\beta^{(n)}_r(0)-c_0\bigr)=-\zeta^n_* ,
\end{equation*}
where $c_0$ is the number, independent of the bare shift, that appears in
\eqref{9.12:eq-runs-coupling-map-a}. Subtracting this identity from
\eqref{9.12:eq-runs-coupling-map-d}, the terms $\zeta^n_*$ and $c_0$ cancel and we obtain, for
every $\hat w\in D(0,2r_J)$,
\begin{equation}\label{9.12:eq-bi-lipschitz-1}
\varphi_n(\hat w)-X=-\hat w-\sum_{r<n}\bigl(\beta^{(n)}_r(\hat w)-\beta^{(n)}_r(0)\bigr).
\end{equation}

\emph{Analyticity.} By (R1) in~\eqref{9.12:eq-bi-lipschitz-a}, each $\beta^{(n)}_r$ is
holomorphic on $D(0,2r_J)$ and real on the real diameter $\mathopen]-2r_J,2r_J\mathclose[$. The
right-hand side of~\eqref{9.12:eq-bi-lipschitz-1} is a finite sum of such functions and of the
function $\hat w\mapsto-\hat w$, and $X=g^{-2}$ is real. Hence $\varphi_n$ is holomorphic on
$D(0,2r_J)$ and real on $\mathopen]-2r_J,2r_J\mathclose[$. Together with the normalisation this
proves the first sentence of part (i).

\emph{One-point bounds.} These come from the summed centred bound in (R2). Let
$\hat w\in D(0,2r_J)$. Apply (R2) with $\varrho=2r_J$, which lies in $[r_J,2r_J]$: for every
$r<n$,
\begin{equation*}
|\beta^{(n)}_r(\hat w)-\beta^{(n)}_r(0)|\le d^{(n)}_r(2r_J)\,|\hat w|,
\qquad \sum_{r<n}d^{(n)}_r(2r_J)\le\tfrac13 .
\end{equation*}
Hence, by~\eqref{9.12:eq-bi-lipschitz-1} and the triangle inequality,
\begin{equation*}
|\varphi_n(\hat w)-X+\hat w|
=\Bigl|\sum_{r<n}\bigl(\beta^{(n)}_r(\hat w)-\beta^{(n)}_r(0)\bigr)\Bigr|
\le\sum_{r<n}d^{(n)}_r(2r_J)\,|\hat w|\le\tfrac13|\hat w| ,
\end{equation*}
which is the second inequality of~\eqref{9.12:eq-bi-lipschitz-b}. Writing
$\varphi_n(\hat w)-X=-\hat w+(\varphi_n(\hat w)-X+\hat w)$, the triangle inequality gives
\begin{equation*}
|\hat w|-\tfrac13|\hat w|\le|\varphi_n(\hat w)-X|\le|\hat w|+\tfrac13|\hat w| ,
\end{equation*}
which is the first pair of inequalities of~\eqref{9.12:eq-bi-lipschitz-b}. For real $\hat w$ the
number $\varphi_n(\hat w)-X+\hat w$ is real, by the analyticity step, and of absolute value at
most $\frac13|\hat w|$; therefore $\varphi_n(\hat w)-X$ lies between $-\frac43\hat w$ and
$-\frac23\hat w$.

\emph{Two-point bounds.} These come from the summed Lipschitz bound (R3). Let
$\hat w,\hat w'\in\overline{D}(0,r_J)$; both lie in $D(0,2r_J)$. Writing
\eqref{9.12:eq-bi-lipschitz-1} at $\hat w$ and at $\hat w'$ and subtracting, the terms
$\beta^{(n)}_r(0)$ cancel and
\begin{equation}\label{9.12:eq-bi-lipschitz-2}
\varphi_n(\hat w)-\varphi_n(\hat w')
=-(\hat w-\hat w')-\sum_{r<n}\bigl(\beta^{(n)}_r(\hat w)-\beta^{(n)}_r(\hat w')\bigr).
\end{equation}
By the definition of the Lipschitz constant on $\overline{D}(0,r_J)$ and by (R3),
\begin{equation*}
\Bigl|\sum_{r<n}\bigl(\beta^{(n)}_r(\hat w)-\beta^{(n)}_r(\hat w')\bigr)\Bigr|
\le\sum_{r<n}\operatorname{Lip}_{\overline{D}(0,r_J)}\beta^{(n)}_r\,|\hat w-\hat w'|
\le\tfrac23|\hat w-\hat w'| .
\end{equation*}
The triangle inequality applied to~\eqref{9.12:eq-bi-lipschitz-2} therefore gives
\begin{equation*}
\bigl(1-\tfrac23\bigr)|\hat w-\hat w'|\le|\varphi_n(\hat w)-\varphi_n(\hat w')|
\le\bigl(1+\tfrac23\bigr)|\hat w-\hat w'| ,
\end{equation*}
which is~\eqref{9.12:eq-bi-lipschitz-c}. Now let $\hat w'<\hat w$ be real numbers in
$[-r_J,r_J]$. By (R1) each $\beta^{(n)}_r$ is real at $\hat w$ and at $\hat w'$, so the sum in
\eqref{9.12:eq-bi-lipschitz-2} is a real number, and by the preceding bound its absolute value is at
most $\frac23(\hat w-\hat w')$. Consequently, by~\eqref{9.12:eq-bi-lipschitz-2},
\begin{equation*}
\varphi_n(\hat w)-\varphi_n(\hat w')\in
\bigl[-(\hat w-\hat w')-\tfrac23(\hat w-\hat w'),\;
-(\hat w-\hat w')+\tfrac23(\hat w-\hat w')\bigr],
\end{equation*}
which is~\eqref{9.12:eq-bi-lipschitz-d}. The right end point $-\frac13(\hat w-\hat w')$ of this
interval is negative, so $\varphi_n(\hat w)<\varphi_n(\hat w')$ whenever
$-r_J\le\hat w'<\hat w\le r_J$; that is, $\varphi_n$ is strictly decreasing on $[-r_J,r_J]$.

\emph{Derivative at $0$.} Since each $\beta^{(n)}_r$ is holomorphic on $D(0,2r_J)$,
differentiating~\eqref{9.12:eq-bi-lipschitz-1} at $\hat w=0$ gives
\begin{equation*}
\varphi_n'(0)=-1-\sum_{r<n}\bigl(\beta^{(n)}_r\bigr)'(0).
\end{equation*}
By (R2), applied with any $\varrho\in[r_J,2r_J]$ (for definiteness $\varrho=2r_J$),
\begin{equation*}
\bigl|\bigl(\beta^{(n)}_r\bigr)'(0)\bigr|
=\lim_{\eta\to0}\frac{|\beta^{(n)}_r(\eta)-\beta^{(n)}_r(0)|}{|\eta|}
\le d^{(n)}_r(\varrho),
\end{equation*}
so that $\bigl|\sum_{r<n}(\beta^{(n)}_r)'(0)\bigr|\le\sum_{r<n}d^{(n)}_r(\varrho)\le\frac13$. Each
$(\beta^{(n)}_r)'(0)$ is real, because $\beta^{(n)}_r$ is real on the real interval
$\mathopen]-2r_J,2r_J\mathclose[$ by (R1) and its derivative at $0$ is therefore the limit of real
difference quotients along that interval. Hence
$-\varphi_n'(0)=1+\sum_{r<n}(\beta^{(n)}_r)'(0)$ is real and lies in $[\frac23,\frac43]$. By
\eqref{9.12:eq-runs-coupling-map-c}, $\varphi_n(\hat w)=X-\zeta^{(n)}_{(0)}(\zeta^n_*+\hat w)$, so
by the chain rule $-\varphi_n'(0)$ is the derivative $\partial\zeta^{(n)}_{(0)}/\partial\zeta$ at
$\zeta=\zeta^n_*$, which is the number $N_n$ of hypothesis (C). The interval obtained here agrees
with the inclusion $N_n\in[\frac23,\frac43]$ of hypothesis (C).

\emph{Uniformity.} The constants $\frac13,\frac23,\frac43,\frac53$ in
\eqref{9.12:eq-bi-lipschitz-b}--\eqref{9.12:eq-bi-lipschitz-d}, and the end points of the interval
$[\frac23,\frac43]$ containing $-\varphi_n'(0)$, were obtained above from the absolute numbers
$\frac13$ in (R2) and $\frac23$ in (R3) by the triangle inequality alone; no other number entered.
They are therefore free of $n$, and of $L$ and $g$ within the standing hypotheses of
Notation~\ref{9.12:not-standing-hypotheses}. The objects to which these bounds apply, namely
$\zeta^n_*$, $a_n$ and the step functions of the runs, are independent of the torus by
hypothesis (D). This proves part (i); parts (ii) and (iii) are proved below.
\end{proof}

\begin{remark}[The two-point constants from the centred bounds]\label{9.12:rem-two-point-constants}
Fix $n\ge 1$. We use the centred step functions $\beta^{(n)}_r$, $0\le r<n$, of the run of length $n$, the
numbers $d^{(n)}_r(\varrho)$, and the inputs (R1)--(R3) of the bi-Lipschitz lemma
(Lemma~\ref{9.12:lem-bi-lipschitz}), as displayed in~\eqref{9.12:eq-bi-lipschitz-a}. For $E\subset\mathbb{C}$,
$\operatorname{Lip}_E h$ is the Lipschitz constant of $h$ on $E$.

The two summed inputs of the bi-Lipschitz lemma give different constants.
\begin{itemize}
\item[(a)] The inequality $\sum_{r<n}d^{(n)}_r(\varrho)\le\frac13$ of (R2) controls only increments measured
from the centre $\hat w=0$. It yields the constants $\frac23$ and $\frac43$ of the one-point bounds
\eqref{9.12:eq-bi-lipschitz-b}.
\item[(b)] The inequality
\begin{equation*}
\sum_{r<n}\operatorname{Lip}_{\overline{D}(0,r_J)}\beta^{(n)}_r\le\frac23
\end{equation*}
of (R3) controls increments between two arbitrary points of $\overline{D}(0,r_J)$. It yields the constants
$\frac13$ and $\frac53$ of the two-point bounds \eqref{9.12:eq-bi-lipschitz-c}.
\item[(c)] (R3) is a consequence of (R1) and (R2) alone.
\end{itemize}

A two-point bound with the constants $\frac23$ and $\frac43$ would require
\begin{equation*}
\sum_{r<n}\operatorname{Lip}_{\overline{D}(0,r_J)}\beta^{(n)}_r\le\frac13 .
\end{equation*}
This inequality is not among the inputs of Lemma~\ref{9.12:lem-bi-lipschitz}, and no such two-point bound is
claimed. The two-point constants used in this book are $\frac13$ and $\frac53$. The constants $\frac23$ and
$\frac43$ hold in the one-point form~\eqref{9.12:eq-bi-lipschitz-b}.
\end{remark}

\begin{proof}
We prove (c). Fix $0\le r<n$ and put
\begin{equation}\label{9.12:eq-two-point-constants-1}
h(\eta):=\beta^{(n)}_r(\eta)-\beta^{(n)}_r(0),\qquad \eta\in D(0,2r_J).
\end{equation}
By (R1), $\beta^{(n)}_r$ is holomorphic on $D(0,2r_J)$, and hence so is $h$.

The range $[r_J,2r_J]$ in (R2) contains $\varrho=2r_J$, so (R2) applies with this value. For every
$\eta\in D(0,2r_J)$ it gives
\begin{equation*}
|h(\eta)|\le d^{(n)}_r(2r_J)\,|\eta|\le 2r_J\,d^{(n)}_r(2r_J).
\end{equation*}

Let $\eta_0\in\overline{D}(0,r_J)$ and $0<s<r_J$. Every point $\eta$ with $|\eta-\eta_0|=s$ satisfies
\begin{equation*}
|\eta|\le|\eta_0|+s<r_J+r_J=2r_J .
\end{equation*}
So the circle $|\eta-\eta_0|=s$ lies in $D(0,2r_J)$. Cauchy's estimate for the derivative of $h$ on this
circle gives
\begin{equation*}
|h'(\eta_0)|\le\frac{1}{s}\max_{|\eta-\eta_0|=s}|h(\eta)|\le\frac{2r_J\,d^{(n)}_r(2r_J)}{s}.
\end{equation*}
The left-hand side does not depend on $s$. Letting $s\uparrow r_J$ gives
\begin{equation}\label{9.12:eq-two-point-constants-2}
|h'(\eta_0)|\le 2\,d^{(n)}_r(2r_J)\qquad\text{for every }\eta_0\in\overline{D}(0,r_J).
\end{equation}

The closed disc $\overline{D}(0,r_J)$ is convex. For $\eta,\eta'\in\overline{D}(0,r_J)$ the segment
$t\mapsto\eta'+t(\eta-\eta')$, $t\in[0,1]$, therefore lies in $\overline{D}(0,r_J)$. On this segment
$h'=(\beta^{(n)}_r)'$, and we obtain
\begin{align*}
|\beta^{(n)}_r(\eta)-\beta^{(n)}_r(\eta')|
&=|h(\eta)-h(\eta')|
=\Bigl|\int_0^1 h'\bigl(\eta'+t(\eta-\eta')\bigr)\,dt\,(\eta-\eta')\Bigr|\\
&\le 2\,d^{(n)}_r(2r_J)\,|\eta-\eta'| ,
\end{align*}
where the inequality is~\eqref{9.12:eq-two-point-constants-2}. Thus
\begin{equation*}
\operatorname{Lip}_{\overline{D}(0,r_J)}\beta^{(n)}_r\le 2\,d^{(n)}_r(2r_J).
\end{equation*}

Summing over $0\le r<n$ and using $\sum_{r<n}d^{(n)}_r(2r_J)\le\frac13$ from (R2) with $\varrho=2r_J$, we get
\begin{equation*}
\sum_{r<n}\operatorname{Lip}_{\overline{D}(0,r_J)}\beta^{(n)}_r
\le 2\sum_{r<n}d^{(n)}_r(2r_J)\le 2\cdot\frac13=\frac23 .
\end{equation*}
This is (R3).
\end{proof}

\begin{proof}[Proof of part~(ii) of Lemma~\ref{9.12:lem-bi-lipschitz}]\label{9.12:prf-inverse-maps}
Recall from the statement of part~(ii) that $J=[-\varrho_2/2,\varrho_2/2]$ and $I_n=\varphi_n(J)$.
Since $\varrho_2/2=r_J/4\le r_J$,
the interval $J$ lies in $\mathopen]-2r_J,2r_J\mathclose[$ and in the closed disc $\overline{D}(0,r_J)$,
so part~(i) of Lemma~\ref{9.12:lem-bi-lipschitz} applies at every point of $J$ and to every pair of
points of $J$.

By part~(i), $\varphi_n$ is holomorphic on $D(0,2r_J)$ and real on $\mathopen]-2r_J,2r_J\mathclose[$.
Hence $\varphi_n|_J$ is real-valued and continuous.

We now locate the endpoints. For real $\hat w\in J$ the number $\varphi_n(\hat w)$ is real, so the bound
$|\varphi_n(\hat w)-X+\hat w|\le\tfrac13|\hat w|$ of \eqref{9.12:eq-bi-lipschitz-b} reads
\begin{equation*}
X-\hat w-\tfrac13|\hat w|\le\varphi_n(\hat w)\le X-\hat w+\tfrac13|\hat w|.
\end{equation*}
At $\hat w=\varrho_2/2$ this gives
\begin{equation*}
X-\tfrac43\cdot\tfrac{\varrho_2}{2}\le\varphi_n(\varrho_2/2)\le X-\tfrac23\cdot\tfrac{\varrho_2}{2},
\end{equation*}
that is, $\varphi_n(\varrho_2/2)\in[X-2\varrho_2/3,\,X-\varrho_2/3]$. At $\hat w=-\varrho_2/2$ it gives
\begin{equation*}
X+\tfrac23\cdot\tfrac{\varrho_2}{2}\le\varphi_n(-\varrho_2/2)\le X+\tfrac43\cdot\tfrac{\varrho_2}{2},
\end{equation*}
that is, $\varphi_n(-\varrho_2/2)\in[X+\varrho_2/3,\,X+2\varrho_2/3]$.

For $\hat w\ne\hat w'$ in $J$, since $J\subset\overline{D}(0,r_J)$, the left-hand inequality of
\eqref{9.12:eq-bi-lipschitz-c} gives
\begin{equation*}
|\varphi_n(\hat w)-\varphi_n(\hat w')|\ge\tfrac13|\hat w-\hat w'|>0 .
\end{equation*}
Thus $\varphi_n|_J$ is injective. A continuous injective real-valued function on an interval is strictly
monotone, by the intermediate value theorem. By the endpoint bounds just obtained,
\begin{equation*}
\varphi_n(\varrho_2/2)\le X-\tfrac{\varrho_2}{3}<X+\tfrac{\varrho_2}{3}\le\varphi_n(-\varrho_2/2),
\end{equation*}
so $\varphi_n|_J$ is strictly decreasing. A continuous, strictly decreasing function on the compact
interval $J$ maps it bijectively, by the intermediate value theorem, onto the compact interval
\begin{equation*}
I_n=\bigl[\varphi_n(\varrho_2/2),\,\varphi_n(-\varrho_2/2)\bigr].
\end{equation*}
In particular the inverse $\psi_n=(\varphi_n|_J)^{-1}\colon I_n\to J$ is well defined.

The left endpoint of $I_n$ is therefore at most $X-\varrho_2/3$, and the right endpoint is at least
$X+\varrho_2/3$. This gives $[X-\varrho_2/3,X+\varrho_2/3]\subset I_n$. The left endpoint is also at
least $X-2\varrho_2/3$, and the right endpoint at most $X+2\varrho_2/3$. This gives
$I_n\subset[X-2\varrho_2/3,X+2\varrho_2/3]$. Together these two inclusions are
\eqref{9.12:eq-bi-lipschitz-e}.

Let $c,c'\in I_n$, and put $\hat w=\psi_n(c)$ and $\hat w'=\psi_n(c')$. Then $\hat w,\hat w'\in J$,
$c=\varphi_n(\hat w)$ and $c'=\varphi_n(\hat w')$. Since $J\subset\overline{D}(0,r_J)$,
\eqref{9.12:eq-bi-lipschitz-c} applies to $\hat w,\hat w'$ and reads
\begin{equation*}
\tfrac13|\hat w-\hat w'|\le|c-c'|\le\tfrac53|\hat w-\hat w'|.
\end{equation*}
The right-hand inequality gives $|\psi_n(c)-\psi_n(c')|\ge\tfrac35|c-c'|$. The left-hand inequality
gives $|\psi_n(c)-\psi_n(c')|\le3|c-c'|$. This is \eqref{9.12:eq-bi-lipschitz-f}, which completes the
proof of part~(ii).
\end{proof}

\begin{proof}[Proof of Lemma~\ref{9.12:lem-bi-lipschitz}, part (iii): consistency across depths
and the uniform limit]
\label{9.12:prf-uniform-limit}
Throughout, $n_1$ is the depth floor of Lemma~\ref{9.12:lem-defect-summability}, so that
$\Upsilon_n\le 3$ for every $n\ge n_1$. We have $\varrho_2=r_J/2$ and
$J=[-\varrho_2/2,\varrho_2/2]=[-r_J/4,r_J/4]$, so that
$J\subset\overline{D}(0,r_J)\subset D(0,2r_J)$. We also write $T_m=\sum_{n\ge m}\Upsilon_n$.

Among the antecedents of the lemma is the conclusion of the uniqueness theorem on the transition
maps $\mathfrak{n}_n$. We use it in the following form. For every $n\ge1$, $\mathfrak{n}_n$ is
holomorphic on $D(0,2\varrho_2)=D(0,r_J)$, real at real points, and satisfies
$\mathfrak{n}_n(0)=0$. For $\hat w$ in that disc, the run of length $n+1$ started at the bare
inverse coupling $a_{n+1}-\mathfrak{n}_n(\hat w)$ and the run of length $n$ started at
$a_n-\hat w$ have equal depth-$0$ coupling. Finally,
\begin{equation*}
|\mathfrak{n}_n(\hat w)-\hat w|\le\Upsilon_n|\hat w| .
\end{equation*}

\emph{The identity \eqref{9.12:eq-bi-lipschitz-g}.} Let $n\ge n_1$ and let $\hat w\in J$ be real.
Since $J\subset D(0,2\varrho_2)$, the map $\mathfrak{n}_n$ is defined and real at $\hat w$. By
the conclusion just recalled, the runs of length $n+1$ from $a_{n+1}-\mathfrak{n}_n(\hat w)$ and
of length $n$ from $a_n-\hat w$ have equal depth-$0$ coupling. The depth-$0$ coupling of a run
determines its depth-$0$ inverse coupling $x_{(0)}$, and is determined by it. By
\eqref{9.12:eq-runs-coupling-map-c}, the depth-$0$ inverse coupling of the run of length $n$ from
$a_n-\hat w$ is $\varphi_n(\hat w)$. By \eqref{9.12:eq-runs-coupling-map-c} at length $n+1$,
that of the run of length $n+1$ from $a_{n+1}-\mathfrak{n}_n(\hat w)$ is
$\varphi_{n+1}(\mathfrak{n}_n(\hat w))$.

This last expression makes sense because $\mathfrak{n}_n(\hat w)$ lies in the domain
$D(0,2r_J)$ of $\varphi_{n+1}$. Indeed, the bound
$|\mathfrak{n}_n(\hat w)-\hat w|\le\Upsilon_n|\hat w|$, the floor $\Upsilon_n\le3$ and
$|\hat w|\le\varrho_2/2$ give
\begin{equation}\label{9.12:eq-uniform-limit-1}
|\mathfrak{n}_n(\hat w)|\le(1+\Upsilon_n)|\hat w|\le 4\cdot\frac{\varrho_2}{2}=r_J .
\end{equation}
Hence $\varphi_{n+1}(\mathfrak{n}_n(\hat w))=\varphi_n(\hat w)$, which is
\eqref{9.12:eq-bi-lipschitz-g}.

\emph{The increment bound \eqref{9.12:eq-bi-lipschitz-h}.} Let $n\ge n_1$ and $\hat w\in J$. The
point $\hat w$ is real with $|\hat w|\le r_J/4$. By \eqref{9.12:eq-uniform-limit-1}, and since
$\mathfrak{n}_n$ is real on real points, $\mathfrak{n}_n(\hat w)$ is real with
$|\mathfrak{n}_n(\hat w)|\le r_J$. So both points lie in $[-r_J,r_J]\subset\overline{D}(0,r_J)$.
By \eqref{9.12:eq-bi-lipschitz-g}, and then the upper bound in \eqref{9.12:eq-bi-lipschitz-c}
for $\varphi_{n+1}$ at the two points $\hat w$ and $\mathfrak{n}_n(\hat w)$ of
$\overline{D}(0,r_J)$, we obtain
\begin{align*}
|\varphi_{n+1}(\hat w)-\varphi_n(\hat w)|
&=|\varphi_{n+1}(\hat w)-\varphi_{n+1}(\mathfrak{n}_n(\hat w))|\\
&\le\tfrac53\,|\hat w-\mathfrak{n}_n(\hat w)|
\le\tfrac53\,\Upsilon_n|\hat w| .
\end{align*}
The last step uses the bound $|\mathfrak{n}_n(\hat w)-\hat w|\le\Upsilon_n|\hat w|$ once more.
This is \eqref{9.12:eq-bi-lipschitz-h}.

\emph{Uniform convergence \eqref{9.12:eq-bi-lipschitz-i}.} Let $p>m\ge n_1$ and $\hat w\in J$.
Write $\varphi_p-\varphi_m$ as the telescoping sum of the increments
$\varphi_{n+1}-\varphi_n$, $m\le n\le p-1$. Apply \eqref{9.12:eq-bi-lipschitz-h} to each
increment; this is allowed since every such $n$ is at least $n_1$. Then use $|\hat w|\le\varrho_2/2$
and $\sum_{n=m}^{p-1}\Upsilon_n\le T_m$. This gives
\begin{align*}
|\varphi_p(\hat w)-\varphi_m(\hat w)|
&\le\sum_{n=m}^{p-1}|\varphi_{n+1}(\hat w)-\varphi_n(\hat w)|
\le\tfrac53\,|\hat w|\sum_{n=m}^{p-1}\Upsilon_n\\
&\le\tfrac53\cdot\frac{\varrho_2}{2}\,T_m=\frac{5\varrho_2}{6}\,T_m .
\end{align*}
By Lemma~\ref{9.12:lem-defect-summability}, $\sum_{n\ge1}\Upsilon_n<\infty$. The numbers $T_m$
are the tails of this convergent series, so $T_m\to0$ as $m\to\infty$. Since the right-hand side
above does not depend on $\hat w\in J$, the sequence $(\varphi_n)_{n\ge n_1}$ is uniformly Cauchy
on $J$. It therefore converges uniformly on $J$ to a function $\varphi_\infty$.

Each $\varphi_n$ is continuous on $J$, being holomorphic on $D(0,2r_J)$ by part (i). Hence
$\varphi_\infty$ is continuous on $J$, as a uniform limit of continuous functions. Letting
$p\to\infty$ in the last display gives
$|\varphi_\infty(\hat w)-\varphi_m(\hat w)|\le(5\varrho_2/6)\,T_m$ for every $\hat w\in J$ and
every $m\ge n_1$. This is \eqref{9.12:eq-bi-lipschitz-i}.

\emph{Properties of $\varphi_\infty$.} By part (i), $\varphi_n(0)=X$ for every $n$, so
$\varphi_\infty(0)=\lim_n\varphi_n(0)=X$.

Fix $\hat w,\hat w'\in J$. The inequalities \eqref{9.12:eq-bi-lipschitz-b},
\eqref{9.12:eq-bi-lipschitz-c} and \eqref{9.12:eq-bi-lipschitz-d} hold for $\varphi_n$ at these
points for every $n$, with constants independent of $n$. Each member of these inequalities
depends continuously on the values $\varphi_n(\hat w)$ and $\varphi_n(\hat w')$, and non-strict
inequalities are preserved under limits. Letting $n\to\infty$ therefore gives
\eqref{9.12:eq-bi-lipschitz-b}, \eqref{9.12:eq-bi-lipschitz-c} and
\eqref{9.12:eq-bi-lipschitz-d} for $\varphi_\infty$ on $J$. Moreover $\varphi_\infty$ is real on
$J$, as a limit of the real values of the $\varphi_n$ there.

In particular, for $\hat w'<\hat w$ in $J$ these inequalities give
\begin{equation*}
\varphi_\infty(\hat w)-\varphi_\infty(\hat w')\le-\tfrac13(\hat w-\hat w')<0 .
\end{equation*}
So $\varphi_\infty$ is strictly decreasing on $J$. Being continuous, it maps $J$ bijectively onto
the compact interval
\begin{equation*}
I_\infty=\varphi_\infty(J)=[\varphi_\infty(\varrho_2/2),\varphi_\infty(-\varrho_2/2)] .
\end{equation*}

The argument that gave \eqref{9.12:eq-bi-lipschitz-e} and \eqref{9.12:eq-bi-lipschitz-f} for
$\varphi_n$ used only \eqref{9.12:eq-bi-lipschitz-b} and \eqref{9.12:eq-bi-lipschitz-c} on $J$.
It therefore applies without change to $\varphi_\infty$, and gives
\eqref{9.12:eq-bi-lipschitz-e} for $I_\infty$ and \eqref{9.12:eq-bi-lipschitz-f} for
$\psi_\infty:=(\varphi_\infty|_J)^{-1}$.

This completes the proof of part (iii). Together with the preceding parts, it completes the proof
of Lemma~\ref{9.12:lem-bi-lipschitz} as an implication from the three antecedents listed in its
statement, the third of which is the conclusion of the uniqueness theorem on the transition maps
$\mathfrak{n}_n$ used above.
\end{proof}

\begin{corollary}[Equal reference-scale coupling gives the same continuum limit]
\label{9.12:cor-same-limit}
Let $\varphi_n$ (for $n\ge 1$), $\varphi_\infty$, $\psi_\infty$, the interval $J=[-\varrho_2/2,\varrho_2/2]$,
the compact interval $I_\infty=\varphi_\infty(J)$, the depth floor $n_1$ and the tail sums $T_m$ be as
in the bi-Lipschitz lemma (Lemma~\ref{9.12:lem-bi-lipschitz}), and let $a_n$ be the tuned bare inverse
coupling. Assume the conclusions of the bi-Lipschitz lemma, and assume the following conclusions of
the uniqueness theorem, under its standing assumptions ($L\ge L_0$, $g\le\gamma_U$) and for every
$m_0\ge 0$, every $\mathsf{p}\ge 2$ and every tuple of test functions $f=(f_1,\dots,f_{\mathsf{p}})$
admitted in its statement on the finite-depth Schwinger functions $S^n(a;f)$, with $n\ge n_0(N_w)$:
\begin{itemize}
\item[$\bullet$] every sequence of real numbers $a'_n$ with $|a'_n-a_n|\le\varrho_2/2$ satisfies
$S^n(a'_n;f)-S^{(a_n-a'_n)}(f)\to 0$ as $n\to\infty$;
\item[$\bullet$] the map $\hat w\mapsto S^{(\hat w)}$ into the limit family is real-analytic on
$\mathopen]-\varrho_2/2,\varrho_2/2\mathclose[$;
\item[$\bullet$] the torus enters these statements only through the number $N_w$, namely in the rate
of convergence and in the floor $n\ge n_0(N_w)$.
\end{itemize}
Let $(\hat w_n)_{n\ge 1}$ and $(\hat w'_n)_{n\ge 1}$ be real sequences in $J$, and let $c$ be a real
number such that $\varphi_n(\hat w_n)\to c$ and $\varphi_n(\hat w'_n)\to c$ as $n\to\infty$; that is,
the two families of bare data $(n,a_n-\hat w_n)$ and $(n,a_n-\hat w'_n)$ have reference-scale inverse
couplings converging to the same value $c$. Then:
\begin{enumerate}
\item[(a)] for every $n\ge 1$,
\begin{equation}\label{9.12:eq-same-limit-1}
|\hat w_n-\hat w'_n|\le 3\,|\varphi_n(\hat w_n)-\varphi_n(\hat w'_n)|,
\end{equation}
hence $\hat w_n-\hat w'_n\to 0$;
\item[(b)] $c\in I_\infty$, and $\hat w_n\to\hat w^{\circ}$ and $\hat w'_n\to\hat w^{\circ}$, where
$\hat w^{\circ}:=\psi_\infty(c)\in J$;
\item[(c)] if $|\hat w^{\circ}|<\varrho_2/2$, then for every $m_0\ge 0$, every $\mathsf{p}\ge 2$ and
every $f=(f_1,\dots,f_{\mathsf{p}})$ as above, both $S^n(a_n-\hat w_n;f)$ and
$S^n(a_n-\hat w'_n;f)$ converge, as $n\to\infty$, to $S^{(\hat w^{\circ})}(f)$; in particular the two
families have the same limit, namely the member $S^{(\hat w^{\circ})}$ of the limit family of the
uniqueness theorem.
\end{enumerate}
The corollary makes no assertion in the case $|\hat w^{\circ}|=\varrho_2/2$ of (c); what that case
requires is set out in the remark on the endpoint case that accompanies this corollary.
\end{corollary}

\begin{proof}
(a) Fix $n\ge 1$. Since $\varrho_2=r_J/2$, the interval $J=[-r_J/4,r_J/4]$ is contained in the
closed disc $\overline{D}(0,r_J)$; hence $\hat w_n,\hat w'_n\in\overline{D}(0,r_J)$. The lower bound in
the two-point bound \eqref{9.12:eq-bi-lipschitz-c} of the bi-Lipschitz lemma, applied to the pair
$\hat w_n,\hat w'_n$, reads $\tfrac13|\hat w_n-\hat w'_n|\le|\varphi_n(\hat w_n)-\varphi_n(\hat w'_n)|$,
which is \eqref{9.12:eq-same-limit-1}. As $n\to\infty$ the right side of
\eqref{9.12:eq-same-limit-1} tends to $3|c-c|=0$, because both $\varphi_n(\hat w_n)$ and
$\varphi_n(\hat w'_n)$ tend to $c$. Hence $\hat w_n-\hat w'_n\to 0$.

(b) Let $n\ge n_1$. By the triangle inequality and the uniform bound \eqref{9.12:eq-bi-lipschitz-i} of
the bi-Lipschitz lemma, applied at the point $\hat w_n\in J$,
\begin{align*}
|\varphi_\infty(\hat w_n)-c|
&\le|\varphi_\infty(\hat w_n)-\varphi_n(\hat w_n)|+|\varphi_n(\hat w_n)-c|\\
&\le\tfrac{5\varrho_2}{6}\,T_n+|\varphi_n(\hat w_n)-c|.
\end{align*}
The first term tends to $0$, since $T_n=\sum_{k\ge n}\Upsilon_k$ is the tail of the series
$\sum_k\Upsilon_k$, which converges by the bi-Lipschitz lemma; the second tends to $0$ by hypothesis.
Thus $\varphi_\infty(\hat w_n)\to c$. The points $\varphi_\infty(\hat w_n)$ lie in
$I_\infty=\varphi_\infty(J)$, which is compact, being the image of the compact interval $J$ under the
continuous function $\varphi_\infty$; hence its limit $c$ lies in $I_\infty$. Since $\varphi_\infty$ is
strictly decreasing on $J$, the inverse $\psi_\infty=(\varphi_\infty|_J)^{-1}\colon I_\infty\to J$ is
defined, and so is $\hat w^{\circ}:=\psi_\infty(c)\in J$. Since $\hat w_n\in J$, we have
$\psi_\infty(\varphi_\infty(\hat w_n))=\hat w_n$, and the upper bound in the Lipschitz bounds
\eqref{9.12:eq-bi-lipschitz-f}, which the bi-Lipschitz lemma asserts for $\psi_\infty$, applied to the
points $\varphi_\infty(\hat w_n),c\in I_\infty$, gives
\begin{equation*}
|\hat w_n-\hat w^{\circ}|=|\psi_\infty(\varphi_\infty(\hat w_n))-\psi_\infty(c)|
\le 3\,|\varphi_\infty(\hat w_n)-c|\longrightarrow 0 .
\end{equation*}
The same argument applies to $(\hat w'_n)$, since $\varphi_n(\hat w'_n)\to c$ as well: it gives
$\varphi_\infty(\hat w'_n)\to c$ and then $\hat w'_n\to\psi_\infty(c)=\hat w^{\circ}$.

(c) Fix $m_0\ge 0$, $\mathsf{p}\ge 2$ and $f=(f_1,\dots,f_{\mathsf{p}})$ as in the hypotheses. Put
$a'_n:=a_n-\hat w_n$. Then $a'_n$ is real, since $a_n$ and $\hat w_n$ are, and
$|a'_n-a_n|=|\hat w_n|\le\varrho_2/2$, since $\hat w_n\in J$. By the first assumed conclusion of the
uniqueness theorem, $S^n(a'_n;f)-S^{(a_n-a'_n)}(f)\to 0$; as $a_n-a'_n=\hat w_n$, this is
\begin{equation}\label{9.12:eq-same-limit-2}
S^n(a_n-\hat w_n;f)-S^{(\hat w_n)}(f)\longrightarrow 0\qquad(n\to\infty).
\end{equation}
By the second assumed conclusion, $\hat w\mapsto S^{(\hat w)}$ is real-analytic on
$\mathopen]-\varrho_2/2,\varrho_2/2\mathclose[$, in particular continuous there, so
$\hat w\mapsto S^{(\hat w)}(f)$ is continuous there. By (b), $\hat w_n\to\hat w^{\circ}$, and by the
proviso $|\hat w^{\circ}|<\varrho_2/2$ the point $\hat w^{\circ}$ lies in the open interval
$\mathopen]-\varrho_2/2,\varrho_2/2\mathclose[$; hence $\hat w_n$ lies in that open interval for all
sufficiently large $n$, and continuity at $\hat w^{\circ}$ gives
\begin{equation}\label{9.12:eq-same-limit-a}
S^{(\hat w_n)}(f)\longrightarrow S^{(\hat w^{\circ})}(f)\qquad(n\to\infty).
\end{equation}
Adding \eqref{9.12:eq-same-limit-2} and \eqref{9.12:eq-same-limit-a} yields
$S^n(a_n-\hat w_n;f)\to S^{(\hat w^{\circ})}(f)$. The same two steps with $\hat w'_n$ in place of
$\hat w_n$ (using $\hat w'_n\in J$ and $\hat w'_n\to\hat w^{\circ}$ from (b)) yield
$S^n(a_n-\hat w'_n;f)\to S^{(\hat w^{\circ})}(f)$. Hence
$S^n(a_n-\hat w_n;f)-S^n(a_n-\hat w'_n;f)\to 0$ for every such $f$: the two families have the same
limit, the member $S^{(\hat w^{\circ})}$, $\hat w^{\circ}=\psi_\infty(c)$, of the limit family of the
uniqueness theorem.

The torus enters only as the third assumed conclusion of the uniqueness theorem allows, that is,
through $N_w$ in the rate of convergence in \eqref{9.12:eq-same-limit-2} and through the floor
$n\ge n_0(N_w)$; the maps $\varphi_n$, $\psi_\infty$ and the point $\hat w^{\circ}$ do not depend on the
torus.
\end{proof}

\begin{remark}[The endpoint case of the equal-coupling corollary. Proof outstanding]
\label{9.12:rem-endpoint-case}
Let $(\hat w_n)_{n\ge1}$, $(\hat w'_n)_{n\ge1}$ and $c$ be as in
Corollary~\ref{9.12:cor-same-limit}, and put $\hat w^{\circ}:=\psi_\infty(c)$. Suppose that
\begin{equation*}
|\hat w^{\circ}|=\varrho_2/2 .
\end{equation*}
Then part~(c) of Corollary~\ref{9.12:cor-same-limit} holds without the proviso
$|\hat w^{\circ}|<\varrho_2/2$. That is, for every $m_0$, every $\mathsf{p}\ge2$ and every
$f=(f_1,\dots,f_{\mathsf{p}})$ as in the uniqueness theorem, both
$S^n(a_n-\hat w_n;f)$ and $S^n(a_n-\hat w'_n;f)$ converge to $S^{(\hat w^{\circ})}(f)$ as
$n\to\infty$.

Proof outstanding.

\emph{Route.} In the proof of Corollary~\ref{9.12:cor-same-limit}, the step
\eqref{9.12:eq-same-limit-a}, namely $S^{(\hat w_n)}(f)\to S^{(\hat w^{\circ})}(f)$, needs in
this case the continuity of $\hat w\mapsto S^{(\hat w)}(f)$ at $\pm\varrho_2/2$ along
$J=[-\varrho_2/2,\varrho_2/2]$.

The uniqueness theorem gives uniform convergence
\begin{equation*}
S^n(a_n-\hat w;f)\to S^{(\hat w)}(f)
\end{equation*}
on a neighbourhood $\mathcal{N}'_w\supset J$. A uniform limit of continuous functions is
continuous. Hence the limit is continuous on the closed segment $J$ as soon as the
following holds:
\begin{itemize}
\item[] for every $m_0$, $\mathsf{p}$ and $f$ as in the uniqueness theorem and every
$n\ge n_0(N_w)$, the map $\hat w\mapsto S^n(a_n-\hat w;f)$ is continuous on
$[-\varrho_2/2,\varrho_2/2]$.
\end{itemize}
This is a property of the finite-depth quantities $S^n$ defined in the uniqueness theorem.

A proof therefore has to establish this continuity from the construction of the
finite-depth quantities $S^n$. Given it, the step \eqref{9.12:eq-same-limit-a} holds for
$|\hat w^{\circ}|=\varrho_2/2$. The remaining steps of the proof of part~(c) of
Corollary~\ref{9.12:cor-same-limit} are then unchanged, and part~(c) holds also when
$|\psi_\infty(c)|=\varrho_2/2$.
\end{remark}

\begin{remark}[What the bi-Lipschitz lemma does not assert]\label{9.12:rem-lipschitz-not-asserted}
The bi-Lipschitz lemma (Lemma~\ref{9.12:lem-bi-lipschitz}) and the corollary that equal
reference-scale coupling gives the same continuum limit
(Corollary~\ref{9.12:cor-same-limit}) are stated for the reference-scale inverse coupling
$\varphi_n(\hat w)$ of the run of length $n$ started at the bare inverse coupling
$a_n-\hat w$. Three further statements lie outside them.
\begin{enumerate}
\item \emph{The re-filed coupling.} Let $\hat x^{(n)}_{(0)}$ denote the re-filed
reference-scale inverse coupling of the uniqueness theorem. The statement that the map
\begin{equation*}
\hat w\longmapsto \hat x^{(n)}_{(0)}(a_n-\hat w)
\end{equation*}
satisfies parts (i)--(iii) of Lemma~\ref{9.12:lem-bi-lipschitz}, that is, the same
statement with $\hat x^{(n)}_{(0)}(a_n-\hat w)$ in place of $\varphi_n(\hat w)$, is not
part of Lemma~\ref{9.12:lem-bi-lipschitz} or of Corollary~\ref{9.12:cor-same-limit}.
It is not proved: it would require three further statements on the re-filed run, none of
which is among the inputs of Lemma~\ref{9.12:lem-bi-lipschitz}.
\item \emph{Injectivity of the limit family.} Injectivity of the map
$\hat w\mapsto S^{(\hat w)}$ on $\mathopen]-\varrho_2/2,\varrho_2/2\mathclose[$ is not
asserted. This is the converse of Corollary~\ref{9.12:cor-same-limit}: it would say that
different values $c$ of the limit of the reference-scale inverse couplings give different
members $S^{(\hat w)}$ of the limit family. It is not proved, and nothing in
Lemma~\ref{9.12:lem-bi-lipschitz} or in
Corollary~\ref{9.12:cor-same-limit} bears on it.
\item \emph{Uniformity in the torus at finite $n$.} Parts (i)--(iii) of
Lemma~\ref{9.12:lem-bi-lipschitz} do not depend on the torus: their inputs are
conclusion (e) and hypothesis $(\mathrm{J}1)$ of the uniqueness theorem.
Part (c) of Corollary~\ref{9.12:cor-same-limit} is a statement for each torus separately,
the torus index $m_0$ being fixed, and the torus enters exactly as conclusion (e) states
it. No uniformity in the torus beyond what conclusion (e) states is asserted.
\end{enumerate}
\end{remark}

\begin{definition}[Tori, quotient metric and test-function norms]\label{9.13:not-tori-metric-norms}
In what follows $L$ is an integer with $L\ge 2$. In the prefix of Theorem~\ref{3.1:thm-main}
(Notation~\ref{2.1:not-prefix}) $L$ is moreover odd and satisfies $L\ge L_0$. Neither the
oddness of $L$ nor the constant $L_0$ is used here.

We write $|\cdot|$ for the Euclidean norm on $\mathbb{R}^4$ and $e_1,\dots,e_4$ for the standard
basis of $\mathbb{R}^4$. For an integer $r\ge 0$ put
\begin{equation}\label{9.13:eq-tori-metric-norms-1}
  \Gamma_r:=2L^{r}\mathbb{Z}^4\subset\mathbb{R}^4,\qquad
  \mathbb{T}_r:=\mathbb{R}^4/\Gamma_r,
\end{equation}
and let $\pi_r\colon\mathbb{R}^4\to\mathbb{T}_r$ denote the quotient map. Thus $\mathbb{T}_r$ is the
torus of side $2L^{r}$; for $r=m$ it is the torus $\mathbb{T}_m$ of
Definition~\ref{1.1:def-lattice}. It carries the quotient metric
\begin{equation}\label{9.13:eq-tori-metric-norms-2}
  \operatorname{dist}_r(x,y):=\min\bigl\{\,|\tilde x-\tilde y| :
  \pi_r(\tilde x)=x,\ \pi_r(\tilde y)=y\,\bigr\}.
\end{equation}
The minimum is attained because $\Gamma_r$ is discrete.

For subsets $E,F\subset\mathbb{T}_r$ we put
\begin{equation*}
  \operatorname{dist}_r(E,F):=\inf\{\operatorname{dist}_r(x,y) : x\in E,\ y\in F\}.
\end{equation*}
For $x\in\mathbb{T}_r$ and $\rho\in\mathbb{R}$ we put
\begin{equation*}
  B_r(x,\rho):=\{y\in\mathbb{T}_r : \operatorname{dist}_r(x,y)<\rho\}.
\end{equation*}

For $f\colon\mathbb{T}_r\to\mathbb{R}$ we put
\begin{equation}\label{9.13:eq-tori-metric-norms-3}
\begin{aligned}
  \|f\|_\infty&:=\sup|f|\in[0,\infty],\\
  \operatorname{Lip}_r f&:=\sup_{x\neq y}
    \frac{|f(x)-f(y)|}{\operatorname{dist}_r(x,y)}\in[0,\infty],\\
  \|f\|_{\sharp,r}&:=\max\bigl\{\|f\|_\infty,\ 40M\cdot\operatorname{Lip}_r f\bigr\}.
\end{aligned}
\end{equation}
Here $M>0$ is the auxiliary parameter in Ba{\l}aban's tuple $\mathfrak{p}$. It is the parameter that
enters the norm $\|\cdot\|_\sharp$ of Definition~\ref{1.3:def-curvature-field}, and
$\|\cdot\|_{\sharp,r}$ is that norm with the torus index displayed. Of $M$ only the inequality
$M\ge 0$ is used in this section.

The \emph{support} of $f\colon\mathbb{T}_r\to\mathbb{R}$ is $\operatorname{supp} f$, the closure in
$\mathbb{T}_r$ of $\{x\in\mathbb{T}_r : f(x)\neq 0\}$.
\end{definition}

\begin{definition}[The bare periodic lattice]\label{9.13:def-bare-lattice}
We use the period lattices $\Gamma_r$, the tori $\mathbb{T}_r$, the integer $L\ge 2$ and the
standard basis $e_1,\dots,e_4$ of $\mathbb{R}^4$ as fixed in
Definition~\ref{9.13:not-tori-metric-norms}. Fix an integer $\mathfrak{n}\ge 0$, the \emph{total
exponent}. The set $(\mathbb{Z}+\tfrac12)^4\subset\mathbb{R}^4$ is invariant under translation by
$\mathbb{Z}^4$, hence by $\Gamma_{\mathfrak{n}}=2L^{\mathfrak{n}}\mathbb{Z}^4\subset\mathbb{Z}^4$.
The \emph{set of sites} is
\begin{equation}\label{9.13:eq-bare-lattice-1}
  \Lambda=\Lambda_{\mathfrak{n}}
  :=\bigl\{\tilde a+\Gamma_{\mathfrak{n}} : \tilde a\in(\mathbb{Z}+\tfrac12)^4\bigr\}
  \subset\mathbb{R}^4/\Gamma_{\mathfrak{n}} .
\end{equation}
Two points of $(\mathbb{Z}+\tfrac12)^4$ define the same site exactly when their difference lies in
$2L^{\mathfrak{n}}\mathbb{Z}^4$, so every site has exactly one representative whose coordinates
all lie in $\{\tfrac12,\tfrac32,\dots,2L^{\mathfrak{n}}-\tfrac12\}$. Hence $\Lambda$ is finite, with
$2L^{\mathfrak{n}}$ sites per coordinate direction and $(2L^{\mathfrak{n}})^4$ sites in all.
The group $\mathbb{Z}^4$ acts on $\Lambda$ by $a\mapsto a+v$, where
$a+v:=(\tilde a+v)+\Gamma_{\mathfrak{n}}$ for any representative $\tilde a$ of $a$.

The \emph{set of positively oriented bonds} is $B:=\Lambda\times\{1,2,3,4\}$; the bond
$(a,\mu)$ joins $a$ to $a+e_\mu$. The \emph{set of plaquettes} is
\begin{equation}\label{9.13:eq-bare-lattice-2}
  P:=\Lambda\times\{(\mu,\nu) : 1\le\mu<\nu\le 4\};
\end{equation}
the plaquette $p=(a;\mu,\nu)$ has corners $a$, $a+e_\mu$, $a+e_\mu+e_\nu$, $a+e_\nu$. Since there are
six pairs $\mu<\nu$,
\begin{equation}\label{9.13:eq-bare-lattice-3}
  |P|=6\,(2L^{\mathfrak{n}})^4 ,
\end{equation}
so that $P$ contains all plaquettes, in all six orientations, as in
Definition~\ref{1.1:def-lattice}. The \emph{abstract barycentre} of $p=(a;\mu,\nu)$ is
\begin{equation}\label{9.13:eq-bare-lattice-4}
  c(p):=\tilde a+\tfrac12 e_\mu+\tfrac12 e_\nu+\Gamma_{\mathfrak{n}}
  \in\mathbb{R}^4/\Gamma_{\mathfrak{n}} .
\end{equation}
It does not depend on the representative $\tilde a$ of $a$: if $\tilde a'=\tilde a+\gamma$ with
$\gamma\in\Gamma_{\mathfrak{n}}$, the right-hand side of \eqref{9.13:eq-bare-lattice-4} changes by
$\gamma\in\Gamma_{\mathfrak{n}}$, that is, not at all.
\end{definition}

\begin{definition}[Gauge fields, Wilson action and bare measure]\label{9.13:def-wilson-action}
We use the set of sites $\Lambda_{\mathfrak{n}}$, the set $B$ of positively oriented bonds, the set
$P$ of plaquettes and the unit vectors $e_\mu$ of the bare periodic lattice of
Definition~\ref{9.13:def-bare-lattice}.

Let $G=\mathrm{SU}(N)$, and let $\operatorname{tr}:=N^{-1}\operatorname{Tr}$ be the normalised
trace. Then $\operatorname{tr}\mathbf{1}=1$, and $|\operatorname{tr}V|\le 1$ for every $V\in G$.

The \emph{configuration space} is $\mathcal{U}:=G^{B}$. A configuration $U\in\mathcal{U}$ assigns one
group element $U(a,\mu)$ to each positively oriented bond $(a,\mu)\in B$. The space $\mathcal{U}$ is
compact. We write $dU$ for the product over $B$ of the normalised Haar measures on $G$; it is a
probability measure on $\mathcal{U}$.

For a plaquette $p=(a;\mu,\nu)\in P$ the \emph{plaquette holonomy} is
\begin{equation}\label{9.13:eq-wilson-action-1}
U(\partial p):=U(a,\mu)\,U(a+e_\mu,\nu)\,U(a+e_\nu,\mu)^{-1}\,U(a,\nu)^{-1}\in G .
\end{equation}
We put
\begin{equation}\label{9.13:eq-wilson-action-2}
a_p(U):=1-\operatorname{Re}\operatorname{tr}U(\partial p)\in[0,2],
\qquad
A(U):=\sum_{p\in P}a_p(U)\in[0,2|P|].
\end{equation}
The ranges in \eqref{9.13:eq-wilson-action-2} follow from $|\operatorname{tr}V|\le 1$. Indeed,
$\operatorname{Re}\operatorname{tr}U(\partial p)\in[-1,1]$, so $a_p(U)\in[0,2]$. Summing over the
$|P|$ plaquettes gives $A(U)\in[0,2|P|]$.

The quantity $a_p$ does not depend on the orientation in which $\partial p$ is traversed. Reversing
the orientation replaces $U(\partial p)$ by $U(\partial p)^{-1}$. For a unitary $V$ we have
$V^{-1}=V^{*}$, hence
\begin{equation*}
\operatorname{tr}V^{-1}=\overline{\operatorname{tr}V}
\quad\text{and}\quad
\operatorname{Re}\operatorname{tr}V^{-1}=\operatorname{Re}\operatorname{tr}V .
\end{equation*}

For real $x_0$ the \emph{Wilson weight} is $e^{-x_0A(U)}$. For $x_0>0$ it is the bare density of
Definition~\ref{1.1:def-wilson-action} at bare inverse coupling $x_0=g_0^{-2}$. The corresponding
probability measure on $\mathcal{U}$ is
\begin{equation}\label{9.13:eq-wilson-action-3}
\mu_{x_0}:=\frac{e^{-x_0A}\,dU}{\int_{\mathcal{U}}e^{-x_0A}\,dU}.
\end{equation}

None of the objects $\Lambda_{\mathfrak{n}}$, $P$, $\mathcal{U}$, $dU$, $a_p$, $A$, $\mu_{x_0}$ refers
to a length.
\end{definition}

\begin{definition}[Systems of units on the bare lattice]\label{9.13:def-systems-of-units}
Let $\Lambda_{\mathfrak{n}}$ be the bare periodic lattice of Definition~\ref{9.13:def-bare-lattice},
with total exponent $\mathfrak{n}$, set of plaquettes $P$, and abstract barycentre
$c(p)\in\mathbb{R}^4/\Gamma_{\mathfrak{n}}$ of a plaquette $p\in P$. We use the tori, period lattices
and quotient maps of Notation~\ref{9.13:not-tori-metric-norms}:
$\Gamma_r=2L^r\mathbb{Z}^4$, $\mathbb{T}_r=\mathbb{R}^4/\Gamma_r$, and
$\pi_r\colon\mathbb{R}^4\to\mathbb{T}_r$ the quotient map.

A \emph{system of units} on $\Lambda_{\mathfrak{n}}$ is a pair $(K,m)$ of integers with $K\ge 0$,
$m\ge 0$ and $K+m=\mathfrak{n}$; $K$ is the \emph{cutoff exponent} and $m$ the
\emph{reference} (torus) \emph{exponent}. In the system $(K,m)$ the lattice is read inside
$\mathbb{T}_m$ through the map
\begin{equation}\label{9.13:eq-systems-of-units-1}
  \iota_{K,m}\colon \mathbb{R}^4/\Gamma_{\mathfrak{n}}\to\mathbb{T}_m,\qquad
  \xi+\Gamma_{\mathfrak{n}}\mapsto L^{-K}\xi+\Gamma_m .
\end{equation}
The barycentre of $p\in P$ in the system $(K,m)$ is
\begin{equation}\label{9.13:eq-systems-of-units-2}
  x_{K,m}(p):=\iota_{K,m}(c(p))\in\mathbb{T}_m .
\end{equation}
For an integer $s$ with $0\le s\le K$, the pair $(K-s,m+s)$ is again a system of units on the
same $\Lambda_{\mathfrak{n}}$, since $(K-s)+(m+s)=\mathfrak{n}$ and $K-s\ge 0$; in it the lattice
spacing is $L^{-(K-s)}$, the torus is $\mathbb{T}_{m+s}$, of side $2L^{m+s}$, and the barycentres
are $x_{K-s,m+s}(p)$. The two systems differ only in which length is called $1$: in $(K,m)$ the
unit of length is $L^{K}$ bare spacings, in $(K-s,m+s)$ it is $L^{K-s}$ bare spacings. Thus passing
from $(K,m)$ to $(K-s,m+s)$ moves the reference length $s$ levels toward the cutoff.
\end{definition}

\begin{proof}[Well-definedness and bijectivity of $\iota_{K,m}$]
Multiplication by $L^{-K}$ is a linear automorphism of $\mathbb{R}^4$. It maps
$\Gamma_{\mathfrak{n}}$ onto $\Gamma_m$ because $m+K=\mathfrak{n}$ gives
\begin{equation*}
  L^{-K}\Gamma_{\mathfrak{n}}=L^{-K}\cdot 2L^{m+K}\mathbb{Z}^4=2L^{m}\mathbb{Z}^4=\Gamma_m .
\end{equation*}
Hence if $\xi-\xi'\in\Gamma_{\mathfrak{n}}$ then $L^{-K}\xi-L^{-K}\xi'\in\Gamma_m$. Therefore
\eqref{9.13:eq-systems-of-units-1} does not depend on the representative $\xi$, and
$\iota_{K,m}$ is a group homomorphism of the quotient groups.

In the same way $L^{K}\Gamma_m=\Gamma_{\mathfrak{n}}$. Hence
$x+\Gamma_m\mapsto L^{K}x+\Gamma_{\mathfrak{n}}$ is a well-defined map
$\mathbb{T}_m\to\mathbb{R}^4/\Gamma_{\mathfrak{n}}$. Composing it with $\iota_{K,m}$ in either
order gives the identity on representatives. So it is the inverse of $\iota_{K,m}$, and
$\iota_{K,m}$ is a bijection, indeed a group isomorphism.

The sites are then read at
\begin{equation*}
  \iota_{K,m}(\Lambda_{\mathfrak{n}})=\pi_m\bigl(L^{-K}(\mathbb{Z}+\tfrac12)^4\bigr).
\end{equation*}
This is a lattice of spacing $L^{-K}$ on the torus of side $2L^m$. It has
$2L^m/L^{-K}=2L^{m+K}$ sites per side. It is the bare lattice of
Definition~\ref{1.1:def-lattice}: sites $\varepsilon(\mathbb{Z}+\tfrac12)^4$ modulo $2L^m$, with
$\varepsilon:=L^{-K}$, on the periodic torus $\mathbb{T}_m$ of side $2L^m$ with bare spacing
$L^{-K}$. Since $c(p)\in\mathbb{R}^4/\Gamma_{\mathfrak{n}}$, the barycentre
\eqref{9.13:eq-systems-of-units-2} is a well-defined point of $\mathbb{T}_m$.

For $0\le s\le K$ the pair $(K-s,m+s)$ satisfies the defining conditions, as noted in the
definition. So the same argument, applied with $(K-s,m+s)$ in place of $(K,m)$, shows that
$\iota_{K-s,m+s}$ is a group isomorphism onto $\mathbb{T}_{m+s}$. The lattice it yields has
spacing $L^{-(K-s)}$, and the barycentres $x_{K-s,m+s}(p)$ are well defined.
\end{proof}

\begin{lemma}[The $L$-adic dilation between tori]\label{9.13:lem-torus-dilation}
We use the tori $\mathbb{T}_r=\mathbb{R}^4/\Gamma_r$ with $\Gamma_r=2L^r\mathbb{Z}^4$, the quotient maps
$\pi_r$ and the quotient metrics $\operatorname{dist}_r$ of
Notation~\ref{9.13:not-tori-metric-norms}, and the systems of units of
Definition~\ref{9.13:def-systems-of-units}. For integers $m\ge0$ and $s\ge0$ define
\begin{equation}\label{9.13:eq-torus-dilation-1}
D_s=D_s^{(m)}\colon \mathbb{T}_{m+s}\to\mathbb{T}_m,\qquad
\tilde y+\Gamma_{m+s}\mapsto L^{-s}\tilde y+\Gamma_m .
\end{equation}
\begin{itemize}
\item[(a)] $D_s$ is well defined and is a bijection (a group isomorphism), with inverse
\begin{equation*}
D_s^{-1}\colon\mathbb{T}_m\to\mathbb{T}_{m+s},\qquad
\tilde x+\Gamma_m\mapsto L^{s}\tilde x+\Gamma_{m+s}.
\end{equation*}
\item[(b)] For every $y\in\mathbb{T}_{m+s}$ one has
$\pi_m^{-1}(D_s y)=L^{-s}\cdot\pi_{m+s}^{-1}(y)$ as subsets of $\mathbb{R}^4$.
\item[(c)] For all $y,y'\in\mathbb{T}_{m+s}$ one has
\begin{equation*}
\operatorname{dist}_m(D_s y,D_s y')=L^{-s}\operatorname{dist}_{m+s}(y,y').
\end{equation*}
In particular $D_s$ is a homeomorphism.
\item[(d)] For every system of units $(K,m)$ on $\Lambda_{\mathfrak{n}}$ and every $0\le s\le K$,
\begin{equation*}
\iota_{K,m}=D_s\circ\iota_{K-s,m+s};
\end{equation*}
hence $x_{K,m}(p)=D_s(x_{K-s,m+s}(p))$ for every $p\in P$.
\end{itemize}
\end{lemma}

\begin{proof}
(a) Let $\tilde y$ and $\tilde y'$ be two representatives of the same point $y\in\mathbb{T}_{m+s}$. Then
$\tilde y-\tilde y'=2L^{m+s}v$ for some $v\in\mathbb{Z}^4$, and the images under the linear map
$\tilde y\mapsto L^{-s}\tilde y$ differ by $L^{-s}\cdot2L^{m+s}v=2L^{m}v\in\Gamma_m$. Hence the class
$L^{-s}\tilde y+\Gamma_m$ does not depend on the representative, and $D_s$ is well defined. Since
$\tilde y\mapsto L^{-s}\tilde y$ is additive on $\mathbb{R}^4$, $D_s$ is additive, that is, a group
homomorphism. The map $\tilde x+\Gamma_m\mapsto L^{s}\tilde x+\Gamma_{m+s}$ is well defined because
\begin{equation*}
L^{s}\Gamma_m=L^{s}\cdot2L^{m}\mathbb{Z}^4=2L^{m+s}\mathbb{Z}^4=\Gamma_{m+s},
\end{equation*}
so that representatives of one class of $\mathbb{T}_m$ are sent into one class of $\mathbb{T}_{m+s}$.
The two maps are mutually inverse because $L^{s}L^{-s}\tilde y=\tilde y$ and
$L^{-s}L^{s}\tilde x=\tilde x$. Thus $D_s$ is a bijective homomorphism, a group isomorphism, with the
displayed inverse.

(b) Write $y=\tilde y+\Gamma_{m+s}$. The preimage of $y$ under $\pi_{m+s}$ is the coset itself,
$\pi_{m+s}^{-1}(y)=\tilde y+\Gamma_{m+s}\subset\mathbb{R}^4$. Using
$L^{-s}\Gamma_{m+s}=\Gamma_m$ (the identity of (a) read backwards) we obtain
\begin{align*}
L^{-s}\bigl(\tilde y+\Gamma_{m+s}\bigr)
&=L^{-s}\tilde y+L^{-s}\Gamma_{m+s}=L^{-s}\tilde y+\Gamma_m\\
&=\pi_m^{-1}\bigl(\pi_m(L^{-s}\tilde y)\bigr)=\pi_m^{-1}(D_s y),
\end{align*}
where the last equality is the definition \eqref{9.13:eq-torus-dilation-1}:
$\pi_m(L^{-s}\tilde y)=L^{-s}\tilde y+\Gamma_m=D_s y$.

(c) By the definition of the quotient metric in Notation~\ref{9.13:not-tori-metric-norms},
$\operatorname{dist}_m(x,x')$ is the minimum of $|\tilde x-\tilde x'|$ over
$\tilde x\in\pi_m^{-1}(x)$ and $\tilde x'\in\pi_m^{-1}(x')$, and likewise for
$\operatorname{dist}_{m+s}$. By (b), $\pi_m^{-1}(D_s y)=L^{-s}\pi_{m+s}^{-1}(y)$ and
$\pi_m^{-1}(D_s y')=L^{-s}\pi_{m+s}^{-1}(y')$, so the pairs $(\tilde x,\tilde x')$ admitted in the
minimum defining $\operatorname{dist}_m(D_s y,D_s y')$ are exactly the pairs
$(L^{-s}\tilde y,L^{-s}\tilde y')$ with $\tilde y\in\pi_{m+s}^{-1}(y)$ and
$\tilde y'\in\pi_{m+s}^{-1}(y')$. By homogeneity of the Euclidean norm,
$|L^{-s}\tilde y-L^{-s}\tilde y'|=L^{-s}|\tilde y-\tilde y'|$. Therefore
\begin{align*}
\operatorname{dist}_m(D_s y,D_s y')
&=\min\bigl\{|\tilde x-\tilde x'| :
\tilde x\in L^{-s}\pi_{m+s}^{-1}(y),\ \tilde x'\in L^{-s}\pi_{m+s}^{-1}(y')\bigr\}\\
&=\min\bigl\{L^{-s}|\tilde y-\tilde y'| :
\tilde y\in\pi_{m+s}^{-1}(y),\ \tilde y'\in\pi_{m+s}^{-1}(y')\bigr\}\\
&=L^{-s}\operatorname{dist}_{m+s}(y,y').
\end{align*}
Consequently $D_s$ is Lipschitz with constant $L^{-s}$, and, $D_s$ being a bijection by (a), the same
identity applied to $y=D_s^{-1}x$, $y'=D_s^{-1}x'$ shows that $D_s^{-1}$ is Lipschitz with constant
$L^{s}$. Both maps are continuous, so $D_s$ is a homeomorphism.

(d) Since $K-s\ge0$, $m+s\ge0$ and $(K-s)+(m+s)=K+m=\mathfrak{n}$, the pair $(K-s,m+s)$ is again a
system of units on $\Lambda_{\mathfrak{n}}$ in the sense of
Definition~\ref{9.13:def-systems-of-units}, and $\iota_{K-s,m+s}$ takes values in
$\mathbb{T}_{m+s}$, the domain of $D_s=D_s^{(m)}$. For $\xi+\Gamma_{\mathfrak{n}}\in
\mathbb{R}^4/\Gamma_{\mathfrak{n}}$ the definition of the reading map gives
$\iota_{K-s,m+s}(\xi+\Gamma_{\mathfrak{n}})=L^{-(K-s)}\xi+\Gamma_{m+s}$, and by
\eqref{9.13:eq-torus-dilation-1}
\begin{equation*}
D_s\bigl(L^{-(K-s)}\xi+\Gamma_{m+s}\bigr)=L^{-s}L^{-(K-s)}\xi+\Gamma_m
=L^{-K}\xi+\Gamma_m=\iota_{K,m}(\xi+\Gamma_{\mathfrak{n}}).
\end{equation*}
This is the identity $\iota_{K,m}=D_s\circ\iota_{K-s,m+s}$. Applying it to
$\xi+\Gamma_{\mathfrak{n}}=c(p)$, the barycentre of $p\in P$ in
$\mathbb{R}^4/\Gamma_{\mathfrak{n}}$, and using $x_{K,m}(p)=\iota_{K,m}(c(p))$ and
$x_{K-s,m+s}(p)=\iota_{K-s,m+s}(c(p))$ from Definition~\ref{9.13:def-systems-of-units}, we obtain
$x_{K,m}(p)=D_s(x_{K-s,m+s}(p))$.
\end{proof}

\begin{definition}[Curvature observables and the sourced partition function]
\label{9.13:def-sourced-partition}
Let $(K,m)$ be a system of units on $\Lambda_{\mathfrak{n}}$ in the sense of
Definition~\ref{9.13:def-systems-of-units}, and let $x_{K,m}(p)\in\mathbb{T}_m$ be the barycentre
of the plaquette $p\in P$ in that system, as fixed there. For a function
$f\colon\mathbb{T}_m\to\mathbb{R}$, with no regularity assumed, define the function
$\mathcal{O}_{K,m}(f)\colon\mathcal{U}\to\mathbb{R}$ by
\begin{equation}\label{9.13:eq-sourced-partition-1}
  \mathcal{O}_{K,m}(f)(U) := \sum_{p\in P} f\bigl(x_{K,m}(p)\bigr)\, a_p(U),
\end{equation}
where $a_p(U)=1-\operatorname{Re}\operatorname{tr}U(\partial p)$ is the plaquette term of
Definition~\ref{9.13:def-wilson-action}. The set $P$ of plaquettes is finite, so the sum in
\eqref{9.13:eq-sourced-partition-1} has finitely many terms. The function
$\mathcal{O}_{K,m}(f)$ is the smeared curvature field $\mathcal{O}(f)$ of
Definition~\ref{1.3:def-curvature-field}, with the system of units written as indices.
Since $a_p(U)\in[0,2]$ for every $p\in P$ and every $U\in\mathcal{U}$
(Definition~\ref{9.13:def-wilson-action}), each of the $|P|$ terms of
\eqref{9.13:eq-sourced-partition-1} has modulus at most
$2\max_{p\in P}|f(x_{K,m}(p))|$. Summing these bounds gives, for every $U\in\mathcal{U}$,
\begin{equation*}
  \bigl|\mathcal{O}_{K,m}(f)(U)\bigr|
  \le 2\,|P|\,\max_{p\in P}\bigl|f\bigl(x_{K,m}(p)\bigr)\bigr| < \infty .
\end{equation*}
For $n\ge 1$, an $n$-tuple $\vec f=(f_1,\dots,f_n)$ of functions
$f_i\colon\mathbb{T}_m\to\mathbb{R}$, and a pair $(x_0,z)\in\mathbb{C}\times\mathbb{C}^n$, the
\emph{sourced partition function} is
\begin{equation}\label{9.13:eq-sourced-partition-2}
  Z_{K,m}(x_0;z;\vec f)
  := \int_{\mathcal{U}}
     \exp\Bigl[-x_0A(U)+\sum_{i=1}^{n} z_i\,\mathcal{O}_{K,m}(f_i)(U)\Bigr]\,dU ,
\end{equation}
where $A(U)=\sum_{p\in P}a_p(U)$ is the Wilson action and $dU$ is the product over the
positively oriented bonds $B$ of the normalised Haar measures, both as in
Definition~\ref{9.13:def-wilson-action}.
\end{definition}

\begin{lemma}[Entirety and positivity of the partition function]\label{9.13:lem-partition-entire}
Let $n\ge 1$ and let $\vec f=(f_1,\dots,f_n)$ be an $n$-tuple of real functions on
$\mathbb{T}_m$. The sourced partition function $Z_{K,m}(\,\cdot\,;\,\cdot\,;\vec f)$ of
Definition~\ref{9.13:def-sourced-partition} is an entire function on
$\mathbb{C}\times\mathbb{C}^n$, and $Z_{K,m}(x_0;z;\vec f)>0$ for all real $(x_0,z)$.
\end{lemma}

\begin{proof}
Identify $\mathbb{C}\times\mathbb{C}^n$ with $\mathbb{C}^{n+1}$. For $U\in\mathcal{U}$ put
\begin{equation*}
\Phi(U):=\bigl(-A(U),\,\mathcal{O}_{K,m}(f_1)(U),\,\dots,\,\mathcal{O}_{K,m}(f_n)(U)\bigr)
\in\mathbb{R}^{n+1},
\end{equation*}
write $\Phi_0(U),\dots,\Phi_n(U)$ for its components in this order, and write
$\xi=(\xi_0,\dots,\xi_n):=(x_0,z)\in\mathbb{C}^{n+1}$. Let
$\langle\xi,\eta\rangle:=\sum_{j=0}^{n}\xi_j\eta_j$
be the bilinear pairing, without complex conjugation. Then the integrand of
Definition~\ref{9.13:def-sourced-partition} is $\exp\langle\xi,\Phi(U)\rangle$.

Each $a_p(U)=1-\operatorname{Re}\operatorname{tr}U(\partial p)$ is a continuous function of $U$
with $0\le a_p(U)\le 2$, because $|\operatorname{tr}V|\le 1$ for $V\in\mathrm{SU}(N)$. Hence
$\Phi$ is continuous on $\mathcal{U}$, and for $i=1,\dots,n$
\begin{equation*}
|A(U)|\le 2|P|,\qquad
|\mathcal{O}_{K,m}(f_i)(U)|\le 2|P|\max_{p\in P}|f_i(x_{K,m}(p))|.
\end{equation*}
The Euclidean norm $|\cdot|$ of a vector in $\mathbb{R}^{n+1}$ or $\mathbb{C}^{n+1}$ is
at most the sum of the moduli of its components, so for all $U\in\mathcal{U}$
\begin{equation}\label{9.13:eq-partition-entire-1}
|\Phi(U)|\le C_\Phi:=2|P|\Bigl(1+\sum_{i=1}^{n}\max_{p\in P}|f_i(x_{K,m}(p))|\Bigr).
\end{equation}

Fix $\sigma>0$. For $|\xi|\le\sigma$ the Cauchy--Schwarz inequality and
\eqref{9.13:eq-partition-entire-1} give
\begin{equation*}
|\langle\xi,\Phi(U)\rangle|\le\sum_{j=0}^{n}|\xi_j|\,|\Phi_j(U)|\le|\xi|\,|\Phi(U)|
\le\sigma C_\Phi .
\end{equation*}
Therefore the $k$-th
term of the exponential series
\begin{equation*}
\exp\langle\xi,\Phi(U)\rangle=\sum_{k\ge 0}\frac{\langle\xi,\Phi(U)\rangle^{k}}{k!}
\end{equation*}
is bounded by $(\sigma C_\Phi)^k/k!$. Since $\sum_k(\sigma C_\Phi)^k/k!<\infty$, the series
converges uniformly in $(U,\xi)$ on $\mathcal{U}\times\{|\xi|\le\sigma\}$.

The measure $dU$ is a probability measure. Uniform convergence therefore permits integration
term by term, and this gives
\begin{equation*}
Z_{K,m}(x_0;z;\vec f)=\sum_{k\ge0}\Pi_k(\xi),\qquad
\Pi_k(\xi):=\int_{\mathcal{U}}\frac{\langle\xi,\Phi(U)\rangle^{k}}{k!}\,dU.
\end{equation*}
Expanding the $k$-th power multinomially shows that $\Pi_k$ is a homogeneous polynomial of degree
$k$ in $\xi$. Its coefficients are integrals of products of components of the bounded continuous
function $\Phi$, so they are finite. Moreover $|\Pi_k(\xi)|\le(\sigma C_\Phi)^k/k!$ for
$|\xi|\le\sigma$. By the Weierstrass test the series $\sum_k\Pi_k$ converges uniformly on
$\{|\xi|\le\sigma\}$ for every $\sigma>0$, that is, locally uniformly on $\mathbb{C}^{n+1}$.
A locally uniform limit of polynomials is holomorphic, so $Z_{K,m}(\,\cdot\,;\,\cdot\,;\vec f)$
is entire.

Now let $\xi=(x_0,z)$ be real. Then $U\mapsto\exp\langle\xi,\Phi(U)\rangle$ is continuous and
strictly positive on $\mathcal{U}$. The space $\mathcal{U}$ is compact, being a finite product of
copies of $\mathrm{SU}(N)$. Hence this function attains a minimum, and the minimum is strictly
positive. Since $dU$ is a probability measure,
\begin{equation*}
Z_{K,m}(x_0;z;\vec f)\ \ge\ \min_{U\in\mathcal{U}}\exp\langle\xi,\Phi(U)\rangle\ >\ 0. \qedhere
\end{equation*}
\end{proof}

\begin{definition}[Truncated functions at bare inverse coupling]\label{9.13:def-truncated-functions}
Let $n\ge 1$, let $\vec f=(f_1,\dots,f_n)$ be an $n$-tuple of test functions, and let
$Z_{K,m}(x_0;z;\vec f)$ be the sourced partition function of
Definition~\ref{9.13:def-sourced-partition}.

By Lemma~\ref{9.13:lem-partition-entire}, $Z_{K,m}(\cdot\,;\cdot\,;\vec f)$ is entire on
$\mathbb{C}\times\mathbb{C}^n$ and $Z_{K,m}(x_0;z;\vec f)>0$ for all real $(x_0,z)$.
Fix a real $x_0$. The function $z\mapsto Z_{K,m}(x_0;z;\vec f)$ is then holomorphic on
$\mathbb{C}^n$, and it is continuous with $Z_{K,m}(x_0;0;\vec f)>0$. Hence it does not
vanish on some neighbourhood of $z=0$ in $\mathbb{C}^n$, which we take to be a polydisc
centred at $0$. On this polydisc, $\log Z_{K,m}(x_0;\cdot\,;\vec f)$, taken real at $z=0$, is
holomorphic. Any other branch of the logarithm differs from this one by a constant, and so
has the same derivatives.

We put
\begin{equation}\label{9.13:eq-truncated-functions-1}
S^T_{n;K,m}[x_0](\vec f)
:=\partial_{z_1}\cdots\partial_{z_n}\log Z_{K,m}(x_0;z;\vec f)\Big|_{z=0}.
\end{equation}
This is the joint cumulant of order $n$ of the curvature observables
$\mathcal{O}_{K,m}(f_1),\dots,\mathcal{O}_{K,m}(f_n)$ under the
normalised Wilson measure $\mu_{x_0}$.

The quantity \eqref{9.13:eq-truncated-functions-1} is the connected (truncated) $n$-point
function $S^T_{n;K,m}$ of the setting of Part~I (see Notation~\ref{0.2:not-glossary} and
Definition~\ref{1.3:def-curvature-field}); it is the same object, with the bare inverse
coupling $x_0$ and the system of units $(K,m)$ now displayed in the notation.
For $n\ge 2$ and test functions with pairwise separated supports, these are the
Schwinger functions of Theorem~\ref{3.1:thm-main}.
\end{definition}

\begin{definition}[Dilation and periodisation]\label{9.13:def-dilation-periodisation}
We use the tori $\mathbb{T}_r=\mathbb{R}^4/\Gamma_r$, their period lattices $\Gamma_r=2L^r\mathbb{Z}^4$
and the quotient maps $\pi_r\colon\mathbb{R}^4\to\mathbb{T}_r$ of
Notation~\ref{9.13:not-tori-metric-norms}.

\emph{Dilation.} For $s\ge0$ the \emph{$L$-adic dilation} is
\begin{equation*}
\delta_s\colon\mathbb{R}^4\to\mathbb{R}^4,\qquad \delta_s(\tilde y):=L^{-s}\tilde y .
\end{equation*}
For $f\colon\mathbb{R}^4\to\mathbb{R}$ we write $f\circ\delta_s=f(L^{-s}\cdot)$; this is the function
denoted $f\circ L^{-s}\cdot$ in Part~I.

\emph{Periodisation.} For a compactly supported $f\colon\mathbb{R}^4\to\mathbb{R}$ and $r\ge0$, the
\emph{periodisation} of $f$ to $\mathbb{T}_r$ is the function
\begin{equation}\label{9.13:eq-dilation-periodisation-1}
\tau_r f\colon\mathbb{T}_r\to\mathbb{R},\qquad
\tau_r f(x):=\sum_{\tilde x\in\pi_r^{-1}(x)} f(\tilde x).
\end{equation}
\end{definition}

\begin{proof}[Well-definedness of the periodisation]
Fix $x\in\mathbb{T}_r$ and a point $\tilde x_0\in\pi_r^{-1}(x)$. Since $\pi_r$ is the quotient map by
$\Gamma_r$, we have $\pi_r^{-1}(x)=\tilde x_0+\Gamma_r$, a translate of the lattice $\Gamma_r$. The
lattice $\Gamma_r$ is discrete: any two distinct points of it are at Euclidean distance at least
$2L^r$. The same holds for its translate $\tilde x_0+\Gamma_r$. The support $\operatorname{supp}f$ is
compact, hence bounded, and so it is contained in a closed ball of some finite radius $\rho$. A set of
points in a ball of radius $\rho$ with mutual distances at least $2L^r$ is finite, because the open
balls of radius $L^r$ about these points are pairwise disjoint and all lie in the ball of radius
$\rho+L^r$. Hence only finitely many $\tilde x\in\pi_r^{-1}(x)$ lie in $\operatorname{supp}f$. All other
terms of the sum in \eqref{9.13:eq-dilation-periodisation-1} vanish. The sum is therefore a finite sum
of real numbers for every $x$, and $\tau_r f$ is a well-defined real-valued function on $\mathbb{T}_r$.
\end{proof}

\begin{lemma}[Transport of test functions to the torus]\label{9.13:lem-transport}
In this lemma $\rho$ denotes a radius. Let $r\ge 0$, let $\rho>0$ with $\rho<L^{r}$, and let
$f\colon\mathbb{R}^4\to\mathbb{R}$ satisfy $\operatorname{supp} f\subset\bar B(0,\rho)$, where
$\bar B(0,\rho):=\{\tilde x\in\mathbb{R}^4:|\tilde x|\le\rho\}$. Let $\tau_r f$ be the
periodisation of Definition~\ref{9.13:def-dilation-periodisation}. Then:
\begin{enumerate}
\item[\textup{(1)}] for every $x\in\mathbb{T}_r$ the set $\pi_r^{-1}(x)\cap\bar B(0,\rho)$ has at
most one point;
\item[\textup{(2)}] $\tau_r f(\pi_r\tilde x)=f(\tilde x)$ for $\tilde x\in\bar B(0,\rho)$, and
$\tau_r f=0$ off $\pi_r(\bar B(0,\rho))$;
\item[\textup{(3)}] $\|\tau_r f\|_\infty=\|f\|_\infty$.
\end{enumerate}
In words, $\tau_r f$ is $f$ transported to $\mathbb{T}_r$ by extension by zero. Each of the two
support floors of Part~I in the setting of Theorem~\ref{3.1:thm-main}, $2\rho<L^{r}$ and
$L^{r}\ge 2\rho+2$, implies $\rho<L^{r}$.
\end{lemma}

\begin{proof}
Recall that $\mathbb{T}_r=\mathbb{R}^4/\Gamma_r$ with $\Gamma_r=2L^{r}\mathbb{Z}^4$, and that
$\pi_r\colon\mathbb{R}^4\to\mathbb{T}_r$ is the quotient map, so $\pi_r\tilde x=\pi_r\tilde y$ if and
only if $\tilde x-\tilde y\in\Gamma_r$.

(1) Let $\tilde x,\tilde y\in\bar B(0,\rho)$. Then
\begin{equation*}
|\tilde x-\tilde y|\le|\tilde x|+|\tilde y|\le 2\rho<2L^{r}.
\end{equation*}
A non-zero element of $\Gamma_r$ has at least one coordinate equal to a non-zero integer multiple
of $2L^{r}$, hence has length at least $2L^{r}$. So if $\pi_r\tilde x=\pi_r\tilde y$, then
$\tilde x-\tilde y$ is an element of $\Gamma_r$ of length less than $2L^r$, hence $\tilde x=\tilde y$.
Thus two distinct points of $\bar B(0,\rho)$ are never congruent modulo $\Gamma_r$, which is (1).

(2) Since $\operatorname{supp} f\subset\bar B(0,\rho)$, $f$ vanishes off $\bar B(0,\rho)$. By
Definition~\ref{9.13:def-dilation-periodisation},
$\tau_r f(x)=\sum_{\tilde y\in\pi_r^{-1}(x)}f(\tilde y)$. Let $\tilde x\in\bar B(0,\rho)$ and
$x=\pi_r\tilde x$. The fibre $\pi_r^{-1}(x)$ contains $\tilde x$, and by (1) every other point of it
lies outside $\bar B(0,\rho)$, where $f$ vanishes; so the sum reduces to $f(\tilde x)$. If
$x\notin\pi_r(\bar B(0,\rho))$, then no point of $\pi_r^{-1}(x)$ lies in $\bar B(0,\rho)$, every term
of the sum vanishes, and $\tau_r f(x)=0$.

(3) By (2), every value of $\tau_r f$ is either $0$ or a value $f(\tilde x)$ with
$\tilde x\in\bar B(0,\rho)$, so $\|\tau_r f\|_\infty\le\|f\|_\infty$. Conversely, for
$\tilde x\in\mathbb{R}^4$ either $\tilde x\in\bar B(0,\rho)$, and then
$|f(\tilde x)|=|\tau_r f(\pi_r\tilde x)|\le\|\tau_r f\|_\infty$ by (2), or $f(\tilde x)=0$; hence
$\|f\|_\infty\le\|\tau_r f\|_\infty$.

Finally, $2\rho<L^{r}$ gives $\rho<L^{r}/2<L^{r}$, and $L^{r}\ge2\rho+2$ gives
$\rho\le L^{r}/2-1<L^{r}$.
\end{proof}

\begin{lemma}[The change-of-units identity]\label{9.13:lem-change-of-units}
Let $L\ge 2$ and $\mathfrak{n}\ge 0$ be integers, and let $G=\mathrm{SU}(N)$, $N\ge 2$, with the
normalised trace $\operatorname{tr}$, $\operatorname{tr}\mathbf{1}=1$. Let $\Lambda_{\mathfrak{n}}$
be the bare periodic lattice of Definition~\ref{9.13:def-bare-lattice}, with its set of plaquettes
$P$. Let $\mathcal{U}=G^{B}$ be the configuration space with product Haar measure $dU$, let $a_p$
be the plaquette functions and $A$ the Wilson action
(Definition~\ref{9.13:def-wilson-action}). Then for every pair of integers $K\ge 0$, $m\ge 0$ with
$K+m=\mathfrak{n}$, and every integer $s$ with $0\le s\le K$, the following hold, where
$D_s\colon\mathbb{T}_{m+s}\to\mathbb{T}_m$ is the $L$-adic dilation of
Lemma~\ref{9.13:lem-torus-dilation}.
\begin{itemize}
\item[(i)] For every function $f\colon\mathbb{T}_m\to\mathbb{R}$ one has
$\mathcal{O}_{K,m}(f)=\mathcal{O}_{K-s,m+s}(f\circ D_s)$ as functions on $\mathcal{U}$.
\item[(ii)] For every $n\ge 1$, every $n$-tuple $\vec f=(f_1,\dots,f_n)$ of functions
$\mathbb{T}_m\to\mathbb{R}$ and every $(x_0,z)\in\mathbb{C}\times\mathbb{C}^n$,
\begin{equation*}
Z_{K,m}(x_0;z;\vec f)=Z_{K-s,m+s}(x_0;z;\vec f\circ D_s),
\qquad \vec f\circ D_s:=(f_1\circ D_s,\dots,f_n\circ D_s).
\end{equation*}
\item[(iii)] For every $n\ge 1$, every real $x_0$ and every $\vec f$ as in \textup{(ii)},
\begin{equation}\label{9.13:eq-change-of-units-a}
S^T_{n;K,m}[x_0](\vec f)=S^T_{n;K-s,m+s}[x_0](\vec f\circ D_s).
\end{equation}
\item[(iv)] For every $f\colon\mathbb{T}_m\to\mathbb{R}$,
\begin{align*}
&\|f\circ D_s\|_\infty=\|f\|_\infty,\qquad
\operatorname{Lip}_{m+s}(f\circ D_s)=L^{-s}\operatorname{Lip}_m f,\\
&\text{hence}\quad \|f\circ D_s\|_{\sharp,m+s}\le\|f\|_{\sharp,m};\qquad
\operatorname{supp}(f\circ D_s)=D_s^{-1}(\operatorname{supp}f);
\end{align*}
for all subsets $E,F\subset\mathbb{T}_m$,
\begin{equation*}
\operatorname{dist}_{m+s}(D_s^{-1}E,D_s^{-1}F)=L^{s}\operatorname{dist}_m(E,F);
\end{equation*}
and for every $x\in\mathbb{T}_m$ and $r>0$,
\begin{equation*}
D_s^{-1}\bigl(B_m(x,r)\bigr)=B_{m+s}(D_s^{-1}x,L^{s}r).
\end{equation*}
\item[(v)] For every $n\ge 1$ and all compactly supported $f_1,\dots,f_n\colon\mathbb{R}^4\to\mathbb{R}$,
with $\tau$ the periodisation and $\delta_s$ the dilation of
Definition~\ref{9.13:def-dilation-periodisation}, one has
$\tau_m f_i\circ D_s=\tau_{m+s}(f_i\circ\delta_s)$ for each $i$; hence, for every real $x_0$,
\begin{equation}\label{9.13:eq-change-of-units-b}
S^T_{n;K,m}[x_0](\tau_m\vec f)=S^T_{n;K-s,m+s}[x_0]\bigl(\tau_{m+s}(\vec f\circ\delta_s)\bigr).
\end{equation}
If moreover $\operatorname{supp}f_i\subset\bar B(0,\rho)$ for all $i$, with $\rho<L^m$, then
$\operatorname{supp}(f_i\circ\delta_s)\subset\bar B(0,L^{s}\rho)$ with $L^{s}\rho<L^{m+s}$;
$\tau_m f_i$ is $f_i$ transported to $\mathbb{T}_m$ and $\tau_{m+s}(f_i\circ\delta_s)$ is
$f_i\circ\delta_s$ transported to $\mathbb{T}_{m+s}$ (Lemma~\ref{9.13:lem-transport}); and pairwise
Euclidean distances of supports are multiplied by $L^{s}$.
\item[(vi)] The map $f\mapsto f\circ D_s$ is a bijection from the set of real functions on
$\mathbb{T}_m$ onto the set of real functions on $\mathbb{T}_{m+s}$, with inverse
$g\mapsto g\circ D_s^{-1}$, and it preserves boundedness, continuity and the Lipschitz class.
Consequently \textup{(iii)} may be read from either side: for integers $K'\ge 0$, $m'\ge 0$ with
$K'+m'=\mathfrak{n}$, for $0\le s\le m'$, $n\ge 1$ and real $x_0$, with
$D_s=D_s^{(m'-s)}\colon\mathbb{T}_{m'}\to\mathbb{T}_{m'-s}$,
\begin{equation*}
S^T_{n;K',m'}[x_0](\vec g)=S^T_{n;K'+s,m'-s}[x_0](\vec g\circ D_s^{-1})
\end{equation*}
for all $n$-tuples $\vec g$ of real functions on $\mathbb{T}_{m'}$.
\end{itemize}
\end{lemma}

Item (v) is the identity displayed in clause (v-d)(iii) of Theorem~0,
\begin{equation*}
S^T_{n;K,m}[x_0](\vec f)=S^T_{n;K-s,m+s}[x_0]\bigl(\vec f(L^{-s}\,\cdot\,)\bigr),
\end{equation*}
in which $f(L^{-s}\,\cdot\,)$ stands for $f\circ\delta_s$, $\delta_s(\tilde y)=L^{-s}\tilde y$, here
with the transport of the test functions to the tori written out. The identity holds for $n=1$ as
well; Theorem~0 asserts nothing at $n=1$, and the identity does not change that.

The proof of items (iv)--(vi) follows that of items (i)--(iii), which we now give. Throughout, fix
$L\ge 2$, $\mathfrak{n}\ge 0$, a system of units $(K,m)$ on $\Lambda_{\mathfrak{n}}$ and an integer
$s$ with $0\le s\le K$, so that $(K-s,m+s)$ is a system of units on the same $\Lambda_{\mathfrak{n}}$
(Definition~\ref{9.13:def-systems-of-units}).

\begin{proof}[Proof of items \textup{(i)}--\textup{(iii)}]
\emph{Item (i).} Let $f\colon\mathbb{T}_m\to\mathbb{R}$. For each $p\in P$,
Lemma~\ref{9.13:lem-torus-dilation}(d) gives $x_{K,m}(p)=D_s(x_{K-s,m+s}(p))$, hence
\begin{equation*}
(f\circ D_s)\bigl(x_{K-s,m+s}(p)\bigr)=f\bigl(D_s(x_{K-s,m+s}(p))\bigr)=f\bigl(x_{K,m}(p)\bigr).
\end{equation*}
By Definition~\ref{9.13:def-sourced-partition}, the two observables $\mathcal{O}_{K,m}(f)$ and
$\mathcal{O}_{K-s,m+s}(f\circ D_s)$ are sums over the same finite set $P$ of the same functions
$a_p\colon\mathcal{U}\to[0,2]$, and by the last display the coefficient of $a_p$ is, plaquette by
plaquette, the same real number in both sums. Therefore, for every $U\in\mathcal{U}$,
\begin{align*}
\mathcal{O}_{K-s,m+s}(f\circ D_s)(U)
&=\sum_{p\in P}(f\circ D_s)\bigl(x_{K-s,m+s}(p)\bigr)\,a_p(U)\\
&=\sum_{p\in P}f\bigl(x_{K,m}(p)\bigr)\,a_p(U)=\mathcal{O}_{K,m}(f)(U).
\end{align*}

\emph{Item (ii).} Let $n\ge 1$, $\vec f=(f_1,\dots,f_n)$ and $(x_0,z)\in\mathbb{C}\times\mathbb{C}^n$.
By Definition~\ref{9.13:def-sourced-partition},
\begin{align*}
Z_{K,m}(x_0;z;\vec f)&=\int_{\mathcal{U}}
\exp\Bigl[-x_0A(U)+\sum_{i=1}^n z_i\,\mathcal{O}_{K,m}(f_i)(U)\Bigr]\,dU,\\
Z_{K-s,m+s}(x_0;z;\vec f\circ D_s)&=\int_{\mathcal{U}}
\exp\Bigl[-x_0A(U)+\sum_{i=1}^n z_i\,\mathcal{O}_{K-s,m+s}(f_i\circ D_s)(U)\Bigr]\,dU.
\end{align*}
The term $-x_0A(U)$ is common to the two integrands, since the Wilson action $A$ is defined on
$\Lambda_{\mathfrak{n}}$ without reference to units (Definition~\ref{9.13:def-wilson-action}).
By item (i) applied to each $f_i$, for every $U\in\mathcal{U}$,
\begin{equation*}
\sum_{i=1}^n z_i\,\mathcal{O}_{K,m}(f_i)(U)=\sum_{i=1}^n z_i\,\mathcal{O}_{K-s,m+s}(f_i\circ D_s)(U).
\end{equation*}
Hence the two integrands are the same function of $U\in\mathcal{U}$, for every
$(x_0,z)\in\mathbb{C}\times\mathbb{C}^n$. Both are integrated against the same measure $dU$ on the
same space $\mathcal{U}$, so the two integrals are equal.

\emph{Item (iii).} Fix a real $x_0$ and $\vec f$. By item (ii), the functions
$z\mapsto Z_{K,m}(x_0;z;\vec f)$ and $z\mapsto Z_{K-s,m+s}(x_0;z;\vec f\circ D_s)$ are the same
function on $\mathbb{C}^n$; by Lemma~\ref{9.13:lem-partition-entire} it is holomorphic, and it is
positive at $z=0$, hence, by continuity, non-vanishing on a ball around $z=0$. By
Definition~\ref{9.13:def-truncated-functions}, each side of \eqref{9.13:eq-change-of-units-a} is
obtained from the logarithm of this function on such a ball, taken real at $z=0$. The two
logarithms are holomorphic logarithms of one and the same non-vanishing function on a ball, so their
difference is continuous with values in $2\pi i\mathbb{Z}$, hence constant, and it vanishes at
$z=0$, where both equal the real logarithm of the same positive number. Thus the two logarithms
coincide near $0$, and so do all their partial derivatives at $z=0$; in particular
\begin{equation*}
\partial_{z_1}\cdots\partial_{z_n}\log Z_{K,m}(x_0;z;\vec f)\big|_{z=0}
=\partial_{z_1}\cdots\partial_{z_n}\log Z_{K-s,m+s}(x_0;z;\vec f\circ D_s)\big|_{z=0},
\end{equation*}
which is \eqref{9.13:eq-change-of-units-a}. (More is true, and equally immediate from item (i): the
random vectors $(\mathcal{O}_{K,m}(f_i))_{i\le n}$ and $(\mathcal{O}_{K-s,m+s}(f_i\circ D_s))_{i\le n}$
are the same map $\mathcal{U}\to\mathbb{R}^n$, so they have the same joint law under $\mu_{x_0}$, the
normalised Wilson measure at bare inverse coupling $x_0$, and every functional of that law agrees:
moments, cumulants, and the centred moments $S_q$ of Notation~\ref{0.2:not-glossary}.)
\end{proof}

\begin{proof}[Proof of clause (iv) of Lemma~\ref{9.13:lem-change-of-units}]
\label{9.13:prf-units-norms-transport}
Let $f\colon \mathbb{T}_m\to\mathbb{R}$.

\emph{Sup norm.} By Lemma~\ref{9.13:lem-torus-dilation}(a), $D_s\colon\mathbb{T}_{m+s}\to\mathbb{T}_m$ is a bijection. Hence
\begin{equation*}
\{(f\circ D_s)(y): y\in\mathbb{T}_{m+s}\}=\{f(x): x\in\mathbb{T}_m\},
\end{equation*}
and therefore $\|f\circ D_s\|_\infty=\|f\|_\infty$.

\emph{Lipschitz constant.} Since $D_s$ is a bijection, $y\neq y'$ holds if and only if $D_sy\neq D_sy'$, and as $(y,y')$ runs over the pairs of distinct points of $\mathbb{T}_{m+s}$, the pair $(x,x'):=(D_sy,D_sy')$ runs over the pairs of distinct points of $\mathbb{T}_m$. By Lemma~\ref{9.13:lem-torus-dilation}(c), $\operatorname{dist}_{m+s}(y,y')=L^{s}\operatorname{dist}_m(D_sy,D_sy')$. Substituting,
\begin{align*}
\operatorname{Lip}_{m+s}(f\circ D_s)
&=\sup_{y\neq y'}\frac{|f(D_sy)-f(D_sy')|}{\operatorname{dist}_{m+s}(y,y')}\\
&=\sup_{x\neq x'}\frac{|f(x)-f(x')|}{L^{s}\operatorname{dist}_m(x,x')}
=L^{-s}\operatorname{Lip}_m f,
\end{align*}
where both sides may equal $+\infty$.

\emph{The norm $\|\cdot\|_{\sharp}$.} From the two identities just proved, together with $L^{-s}\le 1$ and $M\ge 0$,
\begin{align*}
\|f\circ D_s\|_{\sharp,m+s}
&=\max\{\|f\|_\infty,\,40M L^{-s}\operatorname{Lip}_m f\}\\
&\le\max\{\|f\|_\infty,\,40M\operatorname{Lip}_m f\}=\|f\|_{\sharp,m}.
\end{align*}

\emph{Supports.} A point $y\in\mathbb{T}_{m+s}$ satisfies $(f\circ D_s)(y)\neq0$ if and only if $D_sy$ lies in $\{x: f(x)\neq 0\}$, that is,
\begin{equation*}
\{y:(f\circ D_s)(y)\neq0\}=D_s^{-1}(\{x: f(x)\neq0\}).
\end{equation*}
By Lemma~\ref{9.13:lem-torus-dilation}(c), $D_s$ is a homeomorphism, so $D_s^{-1}$ carries the closure of a set onto the closure of its image. Taking closures gives $\operatorname{supp}(f\circ D_s)=D_s^{-1}(\operatorname{supp} f)$.

\emph{Distances of sets.} Let $E,F\subset\mathbb{T}_m$. For $y\in D_s^{-1}E$ and $y'\in D_s^{-1}F$, Lemma~\ref{9.13:lem-torus-dilation}(c) gives $\operatorname{dist}_{m+s}(y,y')=L^{s}\operatorname{dist}_m(D_sy,D_sy')$. Since $D_s$ is a bijection, letting $y$ range over $D_s^{-1}E$ and $y'$ over $D_s^{-1}F$ is the same as letting $D_sy$ range over $E$ and $D_sy'$ over $F$. Taking infima on both sides therefore yields
\begin{equation*}
\operatorname{dist}_{m+s}(D_s^{-1}E,D_s^{-1}F)=L^{s}\operatorname{dist}_m(E,F).
\end{equation*}

\emph{Balls.} Let $x\in\mathbb{T}_m$ and $r>0$. Write $x=D_s(D_s^{-1}x)$ and apply Lemma~\ref{9.13:lem-torus-dilation}(c) to the pair $(D_s^{-1}x,\,y)$. This gives $\operatorname{dist}_m(x,D_sy)=L^{-s}\operatorname{dist}_{m+s}(D_s^{-1}x,y)$, so that
\begin{align*}
D_s^{-1}(B_m(x,r))
&=\{y:\operatorname{dist}_m(x,D_sy)<r\}\\
&=\{y: L^{-s}\operatorname{dist}_{m+s}(D_s^{-1}x,y)<r\}
=B_{m+s}(D_s^{-1}x,L^{s}r).
\end{align*}
This proves clause (iv).
\end{proof}

In particular, let $f_1,\dots,f_n$ be functions on $\mathbb{T}_m$ whose supports are pairwise separated, with $\operatorname{dist}_m(\operatorname{supp}f_i,\operatorname{supp}f_j)\ge\mathsf{d}$ for $i\neq j$. Then, by the support and distance identities of clause (iv), the functions $f_i\circ D_s$ on $\mathbb{T}_{m+s}$ satisfy
\begin{equation*}
\operatorname{dist}_{m+s}(\operatorname{supp}(f_i\circ D_s),\operatorname{supp}(f_j\circ D_s))\ge L^{s}\mathsf{d}.
\end{equation*}
If the pairwise distances of the supports lie in the separation window $[\mathsf{d},6\mathsf{d}]$ (a condition on the test functions), then those of the $f_i\circ D_s$ lie in $[L^{s}\mathsf{d},6L^{s}\mathsf{d}]$. Moreover, by the ball identity, a support contained in a ball of radius $r$ is carried to a support contained in a ball of radius $L^{s}r$.

\begin{proof}[Proof of clause (v) of Lemma~\ref{9.13:lem-change-of-units}]
Let $\delta_s\colon\mathbb{R}^4\to\mathbb{R}^4$ denote the dilation $\tilde y\mapsto L^{-s}\tilde y$. Let $f\colon\mathbb{R}^4\to\mathbb{R}$ be compactly supported, and let $y\in\mathbb{T}_{m+s}$. The periodisation $\tau_r f$ at a point of $\mathbb{T}_r$ is the sum of $f$ over the fibre of $\pi_r$ above that point. This sum is finite because $f$ has compact support and each fibre is a translate of $\Gamma_r$.

By Lemma~\ref{9.13:lem-torus-dilation}(b), $\pi_m^{-1}(D_sy)=L^{-s}\,\pi_{m+s}^{-1}(y)$. Hence $\tilde y\mapsto L^{-s}\tilde y$ is a bijection of $\pi_{m+s}^{-1}(y)$ onto $\pi_m^{-1}(D_sy)$, and reindexing the sum gives
\begin{equation}\label{9.13:eq-units-norms-transport-1}
(\tau_mf)(D_sy)=\sum_{\tilde x\in\pi_m^{-1}(D_sy)}f(\tilde x)
=\sum_{\tilde y\in\pi_{m+s}^{-1}(y)}f(L^{-s}\tilde y)
=\tau_{m+s}(f\circ\delta_s)(y).
\end{equation}
Thus $\tau_mf_i\circ D_s=\tau_{m+s}(f_i\circ\delta_s)$ for each $i$.

Now apply clause (iii), display \eqref{9.13:eq-change-of-units-a}, to the $n$-tuple $\vec f':=\tau_m\vec f$ of functions on $\mathbb{T}_m$. Together with \eqref{9.13:eq-units-norms-transport-1}, this gives
\begin{align*}
S^T_{n;K,m}[x_0](\tau_m\vec f)
&=S^T_{n;K-s,m+s}[x_0](\tau_m\vec f\circ D_s)\\
&=S^T_{n;K-s,m+s}[x_0](\tau_{m+s}(\vec f\circ\delta_s)),
\end{align*}
which is display \eqref{9.13:eq-change-of-units-b}.

Suppose moreover that $\operatorname{supp}f_i\subset\bar B(0,\rho)$ with $\rho<L^{m}$. Since $(f_i\circ\delta_s)(\tilde y)\neq0$ if and only if $L^{-s}\tilde y$ lies in the set where $f_i\neq0$, and since multiplication by $L^{s}$ is a homeomorphism of $\mathbb{R}^4$, we have
\begin{equation*}
\operatorname{supp}(f_i\circ\delta_s)=L^{s}\operatorname{supp}f_i\subset\bar B(0,L^{s}\rho),
\end{equation*}
and
\begin{equation*}
L^{s}\rho<L^{s}L^{m}=L^{m+s}.
\end{equation*}
Lemma~\ref{9.13:lem-transport} therefore applies both at $(m,\rho)$ and at $(m+s,L^{s}\rho)$. Both periodisations are thus transports by extension by zero: $\tau_mf_i$ is $f_i$ transported to $\mathbb{T}_m$, and $\tau_{m+s}(f_i\circ\delta_s)$ is $f_i\circ\delta_s$ transported to $\mathbb{T}_{m+s}$.

For $E,F\subset\mathbb{R}^4$, homogeneity of the Euclidean norm gives $\operatorname{dist}(L^{s}E,L^{s}F)=L^{s}\operatorname{dist}(E,F)$. Hence pairwise Euclidean distances of the supports are multiplied by $L^{s}$.

The two floors of Lemma~\ref{9.13:lem-transport} are also inherited. If $2\rho<L^{m}$, multiplying by $L^{s}$ gives $2L^{s}\rho<L^{m+s}$. If $L^{m}\ge2\rho+2$, then, using $L^{s}\ge1$,
\begin{equation*}
L^{m+s}\ge 2L^{s}\rho+2L^{s}\ge 2L^{s}\rho+2.
\end{equation*}

Finally, let $f\in C^1_c(\mathbb{R}^4)$. Since $\delta_s$ is a bijection of $\mathbb{R}^4$ with $|\delta_s\tilde y-\delta_s\tilde y'|=L^{-s}|\tilde y-\tilde y'|$, the substitution used for the sup norm and the Lipschitz constant in the proof of clause (iv) gives
\begin{equation*}
\|f\circ\delta_s\|_\infty=\|f\|_\infty,
\qquad
\operatorname{Lip}(f\circ\delta_s)=L^{-s}\operatorname{Lip}f
\quad\text{on }\mathbb{R}^4.
\end{equation*}
This proves clause (v).
\end{proof}

\begin{proof}[Proof of clause (vi) of Lemma~\ref{9.13:lem-change-of-units}]
By Lemma~\ref{9.13:lem-torus-dilation}(a), $D_s$ is a bijection with inverse $D_s^{-1}$. Hence $g\mapsto g\circ D_s^{-1}$ inverts $f\mapsto f\circ D_s$. The three properties are preserved as follows:
\begin{itemize}
\item boundedness, by the sup-norm identity of clause (iv);
\item continuity, because $D_s$ is a homeomorphism (Lemma~\ref{9.13:lem-torus-dilation}(c));
\item the Lipschitz class, because by clause (iv) the Lipschitz constant is multiplied by the finite positive factor $L^{-s}$ under $f\mapsto f\circ D_s$, and hence by $L^{s}$ under the inverse map.
\end{itemize}

Now let integers $K'\ge0$, $m'\ge0$ and $0\le s\le m'$ be given, and put $K:=K'+s$ and $m:=m'-s$. Then $K\ge0$, $m\ge0$ and $K+m=K'+m'$. Hence $(K,m)$ and $(K-s,m+s)=(K',m')$ are both systems of units on $\Lambda_{K'+m'}$, and $0\le s\le K$. Clause (iii), display \eqref{9.13:eq-change-of-units-a}, applied with $\vec f:=\vec g\circ D_s^{-1}$, reads
\begin{equation*}
S^T_{n;K'+s,m'-s}[x_0](\vec g\circ D_s^{-1})=S^T_{n;K',m'}[x_0](\vec g),
\end{equation*}
for all $\vec g$ on $\mathbb{T}_{m'}$. This proves clause (vi).
\end{proof}

This completes the proof of Lemma~\ref{9.13:lem-change-of-units}. Its content is the following. The same plaquettes of the same bare lattice, carrying the same Haar measure and the same Wilson weight at the same bare inverse coupling $x_0$, are read with the reference length moved $s$ levels toward the cutoff. The weight $e^{-x_0A}$ does not depend on the units, and the test function is carried along by the dilation that relates the two barycentre maps.

Only $L$-adic factors occur. A reading of $\Lambda_{\mathfrak{n}}$ with spacing $\theta L^{-K}$ and torus side $\theta\cdot2L^{m}$ is again of the form $(K',m')$ with $K'+m'=\mathfrak{n}$ if and only if $\theta=L^{s}$ for an integer $s$ with $-m\le s\le K$. Indeed, $\theta L^{-K}=L^{-K'}$ forces $\theta=L^{s}$ with $s:=K-K'$. Then $\theta\cdot2L^{m}=2L^{m+s}$, so that $m'=m+s$. The conditions $K'\ge0$ and $m'\ge0$ are $s\le K$ and $s\ge-m$. Conversely, every such $s$ gives the system $(K-s,m+s)$.

\begin{corollary}[The coupling shift between two readings]\label{9.13:cor-coupling-shift}
Let $K\ge 0$ and $m\ge 0$ be integers, put $\mathfrak{n}=K+m$, and let $(\Lambda_{\mathfrak{n}},x_0)$
be a bare datum: the bare periodic lattice $\Lambda_{\mathfrak{n}}$ together with the bare density
$\mu_{x_0}$ at bare inverse coupling $x_0=g_0^{-2}$. Assume the following two hypotheses.
\begin{enumerate}
\item[(a)] Clause~(0) of Theorem~\ref{3.1:thm-main} attaches to this bare datum the running inverse
couplings of Definition~\ref{1.2:def-couplings},
\begin{equation*}
x_0=g_0^{-2},\qquad x_{k+1}=x_k-\beta_{k+1}\qquad(0\le k<K),
\end{equation*}
and these satisfy the two-sided flow bound of clause~(0) with the flow letters
$\lambda^{\sharp}>0$, $B^{\sharp}$ and $c_5=2c_2$ of that clause:
\begin{equation}\label{9.13:eq-coupling-shift-a}
\lambda^{\sharp}(j-i)-c_5\;\le\;x_i-x_j\;\le\;B^{\sharp}(j-i)
\qquad\text{for all }0\le i\le j\le K.
\end{equation}
\item[(b)] The sequence $(x_k)_{0\le k\le K}$ of hypothesis~(a) is attached to the bare datum and is
indexed by the number $k$ of renormalisation steps performed from the cutoff: $x_k$ is the effective
inverse coupling after $k$ renormalisation steps from the cutoff, a function of $(\Lambda_{\mathfrak{n}},x_0)$
and of $k$ only, and it does not depend on which step is declared to be the reference length.
In the system of units $(K,m)$ of Definition~\ref{9.13:def-systems-of-units} the reference length
is $L^{K}$ bare spacings and is reached after $K$ steps; in the system $(K-s,m+s)$, $0\le s\le K$,
it is $L^{K-s}$ bare spacings and is reached after $K-s$ steps.
\end{enumerate}
Then for every integer $s$ with $0\le s\le K$ the effective inverse couplings at the reference lengths
of the two readings $(K,m)$ and $(K-s,m+s)$ of the bare datum are $x_K$ and $x_{K-s}$ respectively, and
\begin{equation}\label{9.13:eq-coupling-shift-b}
\lambda^{\sharp}s-c_5\;\le\;x_{K-s}-x_K\;\le\;B^{\sharp}s .
\end{equation}
\end{corollary}

\begin{proof}
Fix an integer $s$ with $0\le s\le K$. Both $(K,m)$ and $(K-s,m+s)$ are systems of units on the same
bare lattice, since $(K-s)+(m+s)=K+m=\mathfrak{n}$, and both readings carry the same bare inverse
coupling $x_0$; so the two readings concern one and the same bare datum $(\Lambda_{\mathfrak{n}},x_0)$.

By hypothesis~(b), in the reading $(K,m)$ the reference length is reached after $K$ renormalisation
steps from the cutoff, and the effective inverse coupling after $K$ steps is $x_K$; in the reading
$(K-s,m+s)$ the reference length is reached after $K-s$ steps, and the effective inverse coupling after
$K-s$ steps is $x_{K-s}$. Since, again by hypothesis~(b), $x_k$ depends only on the bare datum and on
$k$, and not on which step is declared to be the reference length, the same sequence $(x_k)$ serves both
readings. Hence the effective inverse couplings at the two reference lengths are $x_K$ and $x_{K-s}$,
which is the first assertion.

For the inequality, apply \eqref{9.13:eq-coupling-shift-a} with $i:=K-s$ and $j:=K$. The condition
$0\le i\le j\le K$ holds because $0\le s\le K$, and $j-i=s$. Then \eqref{9.13:eq-coupling-shift-a}
reads $\lambda^{\sharp}s-c_5\le x_{K-s}-x_K\le B^{\sharp}s$, which is \eqref{9.13:eq-coupling-shift-b}.
\end{proof}

Thus one and the same bare datum presents, in the two systems of units $(K,m)$ and $(K-s,m+s)$: the
same bare inverse coupling $x_0$; the same bare lattice $\Lambda_{\mathfrak{n}}$; truncated functions
related exactly by the change-of-units identity, Lemma~\ref{9.13:lem-change-of-units}, in its parts
\eqref{9.13:eq-change-of-units-a} and \eqref{9.13:eq-change-of-units-b}; the tori $\mathbb{T}_m$ and
$\mathbb{T}_{m+s}$; and reference-scale effective inverse couplings whose difference lies in the interval
$[\lambda^{\sharp}s-c_5,\,B^{\sharp}s]$. This is the interval that Theorem~\ref{3.1:thm-main} carries for
$\hat x_s-\hat x$; the identification of the pins $\hat x$ and $\hat x_s$ of clause~(v) with the effective
inverse couplings after $K$, respectively $K-s$, steps is part of the statement of
Theorem~\ref{3.1:thm-main} and is not proved again here.

\begin{remark}[Re-read datum, first passage and torus floor]\label{9.13:rem-reread-datum}
Let $(\Lambda_{\mathfrak{n}},x_0)$, $\mathfrak{n}=K+m$, be a bare datum to which clause~(0) of
Theorem~0 attaches the inverse couplings $(x_k)_{0\le k\le K}$. Let $s$ be an integer
with $0\le s\le K$. The \emph{re-read datum} is the same datum read in the system of units $(K-s,m+s)$, as
in the change-of-units identity (Lemma~\ref{9.13:lem-change-of-units}). We make two observations on the
quantifier prefix of Theorem~0 (Notation~\ref{2.1:not-prefix}).

\emph{First passage.} In Theorem~0 the cutoff exponent is a first passage
(Definition~\ref{1.2:def-flow-constants}):
\begin{equation}\label{9.13:eq-reread-datum-1}
K=K(g_0,g)=\min\{k\ge 0 : x_k\le g^{-2}+B^{\sharp}\}.
\end{equation}
Suppose $K-s\ge1$. The re-read cutoff exponent $K-s$ is the first-passage index of the same flow for some
target $X'$ only if the interval
\begin{equation*}
[\,x_{K-s}-B^{\sharp},\;x_{K-s-1}-B^{\sharp}\,[
\end{equation*}
is non-empty, that is, only if $x_{K-s-1}>x_{K-s}$. The flow
bound~\eqref{9.13:eq-coupling-shift-a} of clause~(0), taken with $j-i=1$, yields this only when
$\lambda^{\sharp}>c_5$, an inequality that is not assumed here. Hence the change-of-units identity does
not by itself place the re-read datum inside the part of the prefix of Theorem~0 in which the cutoff
exponent is fixed by first passage. This is consistent with the prefix of Theorem~0, according to which,
in clause~(v), the bare inverse coupling is tuned, $g_0^{-2}=a_n-\hat w$, rather than fixed by first
passage.

\emph{Torus floor.} The torus floor in the setting of Theorem~0, $m\ge m_{\mathbb{T}}(g_-(g))$ with
$m_{\mathbb{T}}$ as in the glossary (Notation~\ref{0.2:not-glossary}),
\begin{equation}\label{9.13:eq-reread-datum-2}
m_{\mathbb{T}}(g')=11+m'+\log_L R(g'),
\end{equation}
depends on the reference coupling. The re-read datum has a different reference coupling, so the floor
relevant to it is $m+s\ge m_{\mathbb{T}}(\cdot)$, evaluated at the argument appropriate to the re-read
system. This is a one-sided floor on $m+s$, which the change-of-units identity neither needs nor supplies.
\end{remark}

\begin{proof}
For a target $X'\in\mathbb{R}$, the first-passage index of the flow $(x_k)$ is
$\min\{k\ge0 : x_k\le X'+B^{\sharp}\}$, in the form~\eqref{9.13:eq-reread-datum-1} with $g^{-2}$
replaced by $X'$. Let $K-s\ge1$. If $K-s$ is this index, then $x_{K-s}\le X'+B^{\sharp}$, and
$x_{K-s-1}>X'+B^{\sharp}$ because $K-s-1$ is smaller than the minimum. These two inequalities together say
\begin{equation*}
x_{K-s}-B^{\sharp}\le X'<x_{K-s-1}-B^{\sharp},
\end{equation*}
that is, $X'$ lies in the interval displayed in the remark, which is therefore non-empty; this is the
inequality $x_{K-s-1}>x_{K-s}$. Conversely, if the interval is non-empty and $X'$ is chosen in it, the
display holds, so $x_{K-s}\le X'+B^{\sharp}<x_{K-s-1}$; then $K-s$ is the first-passage index for $X'$
as soon as $x_k>X'+B^{\sharp}$ holds also for every $k<K-s-1$, a condition not available here; this is
why the remark states only the first of these two directions.

The flow bound~\eqref{9.13:eq-coupling-shift-a} with $i=K-s-1$ and $j=K-s$, so that $j-i=1$, gives
\begin{equation*}
x_{K-s-1}-x_{K-s}\ge\lambda^{\sharp}-c_5.
\end{equation*}
The right-hand side is positive when $\lambda^{\sharp}>c_5$. Without that inequality, the lower bound does
not exclude $x_{K-s-1}\le x_{K-s}$. In that case the interval is empty, and $K-s$ is the first-passage
index for no target. The consistency with clause~(v) rests on the prefix of Theorem~0
(Notation~\ref{2.1:not-prefix}), in which the bare inverse coupling of clause~(v) is tuned,
$g_0^{-2}=a_n-\hat w$, and is not fixed by first passage.

For the torus floor, note that the floor of Theorem~0 is a function of the reference coupling
through~\eqref{9.13:eq-reread-datum-2}. By Corollary~\ref{9.13:cor-coupling-shift}, the effective
inverse couplings at the reference lengths of the readings
$(K,m)$ and $(K-s,m+s)$ are $x_K$ and $x_{K-s}$. These are in general different. The torus exponent of
the re-read system is $m+s$, so the floor that applies to it is a lower bound on $m+s$ at the argument
appropriate there.

The hypotheses of Lemma~\ref{9.13:lem-change-of-units} are that $L\ge2$, $K,m\ge0$ and $0\le s\le K$
are integers, together with, for transported test functions, a bound on their support radius relative
to $L^{m}$. None of them is a floor of the form $m+s\ge m_{\mathbb{T}}(\cdot)$, and none of the
conclusions of that lemma is such a floor; so this floor is neither an input nor a conclusion of that
identity.
\end{proof}

\begin{remark}[Limits of the change-of-units identity]\label{9.13:rem-units-limits}
Lemma~\ref{9.13:lem-change-of-units} relates the truncated functions of one and the same bare
periodic lattice $\Lambda_{\mathfrak{n}}$ read in two systems of units. It does not give the
following.
\begin{itemize}
\item[(a)] It relates no truncated functions at two distinct bare inverse couplings $x_0\neq x_0'$.
It also relates none at two distinct reference couplings on one and the same torus
$\mathbb{T}_m$. The identity moves $K$ and $m$ together, keeping $K+m=\mathfrak{n}$ and the bare
inverse coupling $x_0$ fixed: $S^T_{n;K,m}[x_0]$ is compared only with
$S^T_{n;K-s,m+s}[x_0]$, with the test functions composed with $D_s$.
\item[(b)] It gives no continuous dilation. The depth shift $s$ is an integer with
$0\le s\le K$, and the closing observation of the proof of
Lemma~\ref{9.13:lem-change-of-units}, which we recall, shows why. For a real number
$\theta>0$, a reading of $\Lambda_{\mathfrak{n}}$ with spacing $\theta L^{-K}$ and torus side
$\theta\cdot 2L^{m}$ is again a system of units $(K',m')$ with $K'+m'=\mathfrak{n}$ if and only
if $\theta=L^{s'}$ for an integer $s'$ with $-m\le s'\le K$. The values $0\le s'\le K$ are
those of Lemma~\ref{9.13:lem-change-of-units}; for $s'<0$ the system $(K-s',m+s')$ is related to
$(K,m)$ by that lemma applied from $(K-s',m+s')$ with depth shift $-s'$.
\item[(c)] It contains no statement of dimensional transmutation.
\item[(d)] It gives no injectivity of the map that sends a reference coupling, or a pin
$\hat x$, to the corresponding truncated functions.
\end{itemize}
All of this agrees with the closing sentence of item~(iii) of clause~(v-d) of
Theorem~\ref{3.1:thm-main}, which states what is not asserted there.
\end{remark}

\begin{lemma}[The limit-member identification. Proof outstanding]\label{9.13:lem-limit-member}
Let $(m_0,\hat x)$ index a member $S^{(m_0,\hat x)}$ of the limit family of clause~(v) of
Theorem~\ref{3.1:thm-main}. This member is obtained along the tuned sequence of bare data of that clause,
\begin{equation*}
  \bigl(\Lambda_{K_\nu+m_0},\,x_0^{(\nu)}\bigr),\qquad \nu\to\infty,
\end{equation*}
on $\mathbb{T}_{m_0}$, where $\nu$ denotes the index of the sequence. For each $\nu$ let
$(x^{(\nu)}_k)_{0\le k\le K_\nu}$ be the running inverse couplings that clause~(0) of
Theorem~\ref{3.1:thm-main} attaches to the bare datum $(\Lambda_{K_\nu+m_0},x_0^{(\nu)})$
(Definition~\ref{1.2:def-couplings}). Let $s\ge 0$ be an integer such that, for all
sufficiently large $\nu$, the exponent $K_\nu-s$ is not below the floor that clause~(v) of
Theorem~\ref{3.1:thm-main} places on the cutoff exponent. By the \emph{re-read datum} we mean
the same bare lattice $\Lambda_{K_\nu+m_0}$ with the same bare inverse coupling $x_0^{(\nu)}$,
read in the system of units $(K_\nu-s,\,m_0+s)$. Then the following hold.
\begin{enumerate}
\item[(1)] The depth-$s$ reference couplings $x^{(\nu)}_{K_\nu-s}$ of the re-read data
converge, as $\nu\to\infty$, to a limit $\hat x_s$, and
\begin{equation}\label{9.13:eq-limit-member-1}
  \hat x_s-\hat x\in\bigl[\lambda^{\sharp}s-c_5,\;B^{\sharp}s\bigr].
\end{equation}
\item[(2)] The re-read sequence is admissible for the uniqueness theorem of clause~(v) of
Theorem~\ref{3.1:thm-main} at $(m_0+s,\hat x_s)$. Consequently its truncated functions converge
to the member $S^{(m_0+s,\hat x_s)}$.
\item[(3)] Consequently, by Lemma~\ref{9.13:lem-change-of-units}, display
\eqref{9.13:eq-change-of-units-b}, applied at each $\nu$, and by passage to the limit on both
sides, the following holds for every $n\ge 1$ and every $n$-tuple $\vec f$ of test functions
in the class admitted by Theorem~\ref{3.1:thm-main} on $\mathbb{T}_{m_0}$:
\begin{equation}\label{9.13:eq-limit-member-2}
  S^{(m_0,\hat x)}\bigl(\tau_{m_0}\vec f\bigr)
  =S^{(m_0+s,\hat x_s)}\bigl(\tau_{m_0+s}(\vec f\circ\delta_s)\bigr).
\end{equation}
\end{enumerate}
\end{lemma}

Proof outstanding.

\medskip
\noindent\textit{Route.}
Two statements in this chapter concern a fixed bare lattice $\Lambda_{\mathfrak{n}}$ and a fixed
bare inverse coupling $x_0$ at finite cutoff. They are
Lemma~\ref{9.13:lem-change-of-units} and Corollary~\ref{9.13:cor-coupling-shift}, and the latter
is proved as an implication from \eqref{9.13:eq-coupling-shift-a}. Apply
Lemma~\ref{9.13:lem-change-of-units} to the $\nu$-th member of the tuned sequence, with
$K=K_\nu$, $m=m_0$ and the given $s$. Display \eqref{9.13:eq-change-of-units-b} then yields, for
each $\nu$, an exact identity
\begin{equation*}
  S^T_{n;K_\nu,m_0}\bigl[x_0^{(\nu)}\bigr]\bigl(\tau_{m_0}\vec f\bigr)
  =S^T_{n;K_\nu-s,m_0+s}\bigl[x_0^{(\nu)}\bigr]\bigl(\tau_{m_0+s}(\vec f\circ\delta_s)\bigr).
\end{equation*}
It relates the truncated functions of the datum on $\mathbb{T}_{m_0}$ at cutoff exponent
$K_\nu$ to the truncated functions of the same bare datum read on $\mathbb{T}_{m_0+s}$ at cutoff
exponent $K_\nu-s$. Along the tuned sequence the left-hand side converges to
$S^{(m_0,\hat x)}(\tau_{m_0}\vec f)$.

Passing from these identities to \eqref{9.13:eq-limit-member-2}, an identity between the limit
member $(m_0,\hat x)$ on $\mathbb{T}_{m_0}$ and a limit member $(m_0+s,\hat x_s)$ on
$\mathbb{T}_{m_0+s}$, requires in addition the following four inputs.
\begin{enumerate}
\item[1.] The re-read sequence of bare data must be an admissible tuned sequence of clause~(v)
at torus exponent $m_0+s$ with a pin $\hat x_s$, or be asymptotically comparable to one. In
particular, the depth-$s$ reference coupling $x^{(\nu)}_{K_\nu-s}$ (the quantity
$x_{K-s}$ of Corollary~\ref{9.13:cor-coupling-shift}) must converge along the tuned family to a
limit $\hat x_s$. Corollary~\ref{9.13:cor-coupling-shift} places the difference of the two
reference couplings in the interval of display~\eqref{9.13:eq-coupling-shift-b} for every
$\nu$:
\begin{equation*}
  \lambda^{\sharp}s-c_5\le x^{(\nu)}_{K_\nu-s}-x^{(\nu)}_{K_\nu}\le B^{\sharp}s.
\end{equation*}
Since this interval is closed, the limit $\hat x_s$ then satisfies
\eqref{9.13:eq-limit-member-1}.
\item[2.] The uniqueness theorem of clause~(v) of Theorem~\ref{3.1:thm-main} is needed at
$(m_0+s,\hat x_s)$.
It ensures that the limit of the right-hand sides of \eqref{9.13:eq-change-of-units-b} along
the re-read sequence is the member $S^{(m_0+s,\hat x_s)}$, and not merely a subsequential limit.
\item[3.] The bi-Lipschitz lemma is needed at the shifted pin $\hat x_s$, to the extent that
clause~(v-d) of Theorem~\ref{3.1:thm-main} uses it at that pin. It enters only through
clause~(v-d)(ii).
\item[4.] If the pin $\hat x$ is written in the re-filed coupling rather than in a step-indexed
effective coupling, the statement of this chapter relating the re-filed coupling to the
step-indexed effective couplings along the tuned sequence is also needed.
\end{enumerate}
Inputs~1 and~2 are items~(1) and~(2) of the lemma. Given them, item~(3) is a two-line
consequence of Lemma~\ref{9.13:lem-change-of-units}. One writes
\eqref{9.13:eq-change-of-units-b} at each $\nu$ for the bare datum
$(\Lambda_{K_\nu+m_0},x_0^{(\nu)})$ and lets $\nu\to\infty$ on both sides. The left side
converges to $S^{(m_0,\hat x)}(\tau_{m_0}\vec f)$, and by item~(2) the right side converges to
$S^{(m_0+s,\hat x_s)}(\tau_{m_0+s}(\vec f\circ\delta_s))$.

The lemma relates no two distinct pins on one torus. It asserts neither a continuous dilation
nor the injectivity of any map between members.

\begin{definition}[The canonical scaffold, tuned runs and shifted reference couplings]
\label{9.13:not-canonical-scaffold}
The following letters are used throughout this chapter.

\emph{The canonical scaffold.} Let $g$ be the reference coupling and $X=g^{-2}$. The
\emph{canonical scaffold} $\Theta_X$ is the sequence
\begin{equation*}
\theta_r := X + c_0 r, \qquad r=0,1,2,\dots, \qquad c_0 = 2\lambda .
\end{equation*}
The index $r$ is the \emph{depth}, the number of levels below the reference length
$\ell_g$. The \emph{length} $n$ of a run is its cutoff.

\emph{Bare data and Schwinger functions.} For a tuple $\vec f=(f_1,\dots,f_{\mathsf{p}})$ of test
functions we put $w=\sum_i z_i f_i$. For integers $m_0$, $n$ and a bare inverse coupling $x$,
the \emph{bare datum} $(m_0+n,x)$ at the source $w$ is the Wilson weight $\exp[-A(x-w,U)]$, in
the weighted form of Definition~\ref{1.1:def-wilson-action}, on the periodic lattice with
$2L^{m_0+n}$ sites on a side. The torus then has side $2L^{m_0}\ell_g$. We define
\begin{equation}\label{9.13:eq-canonical-scaffold-1}
S^n_{m_0}[x](\vec f) := \partial_{z_1}\cdots\partial_{z_{\mathsf{p}}}\big|_{z=0}\log Z ,
\end{equation}
where $Z$ is the partition function of the datum $(m_0+n,x)$ at the source $w$. This is the
connected Schwinger function of clause (ii-$\mathbb{T}$-a) of Theorem~\ref{3.1:thm-main}.

\emph{Runs.} A run of length $n$ starts from a \emph{bare deviation} $\zeta_{(n)}:=\zeta$. Its
bare inverse coupling is $\theta_n-\zeta$. For $r=n-1,\dots,0$ it proceeds by
\begin{equation*}
\zeta_{(r)} := \zeta_{(r+1)} + b_{(r)}(\zeta) - c_0, \qquad x_{(r)} := \theta_r - \zeta_{(r)} .
\end{equation*}
Here $\zeta_{(r)}$ is the deviation and $x_{(r)}$ the running inverse coupling at depth $r$.
By the scaffold form of the running-shift statement of the correlation theorem, the
per-depth coefficient is
\begin{equation*}
b_{(r)} = \beta^0_{n-r} + f_{n-r}\big(\zeta_{(n)},\dots,\zeta_{(r+1)}\big) + \mathsf{q}_{(r)} .
\end{equation*}
Its one-loop part is $\beta^0_{n-r}$, its higher-order part $f_{n-r}$ and its remainder
$\mathsf{q}_{(r)}$. These do not depend on the torus, and they satisfy
\begin{equation*}
\big|f_{n-u}+\mathsf{q}_{(u)}\big| \le \iota_u := C_\beta e_{(u+1)} + v_{(u+1)} ,
\end{equation*}
where $e_{(u)}$ and $v_{(u)}$ are the two per-depth sequences from which that statement composes
this drift bound. The bound $\iota_u=\iota_u(X)$ depends on the reference inverse coupling $X$;
we write the argument where it matters.

\emph{Bounded defect.} By the one-loop Ces\`aro-defect theorem, the one-loop coefficients obey
\begin{equation*}
\sum_{l=i+1}^{j}\beta^0_l = (j-i)c_0 + s_i - s_j, \qquad |s_i|\le c_2 \ \text{for all } i .
\end{equation*}
This bounds the Ces\`aro defect of the partial sums only; it gives no bound on the individual
differences $\beta^0_l-c_0$. The defect sequence is always written with a subscript, $s_i$; an
unsubscripted $s$ below is a depth. The \emph{defect-shifted deviation} is
$\tilde\zeta_{(r)} := \zeta_{(r)} + s_{n-r}$.

\emph{Constants.} Let $r_J$ be the constant of the uniqueness theorem that bounds the offsets
$v$, $|v|\le r_J$, in the matching statement recalled below; it equals the radius $\varrho$ of
Theorem~\ref{3.1:thm-main}, $r_J=\varrho$, and $\varrho_2=\varrho/2$ as there. We put
\begin{equation*}
t_{\circ} := c_2 + r_J, \qquad R := 4r_J .
\end{equation*}
In this chapter $R$ denotes this constant, not the ball radius of the observables.

\emph{Tuned runs.} Under the hypotheses of the canonical-scaffold theorem on $g$, that
theorem, applied with these constants, gives the following four statements.
\begin{itemize}
\item[(a)] The map $\zeta\mapsto\tilde\zeta_{(0)}(\zeta)$ is a real-analytic increasing
bijection from the interval $J_0$ of that theorem onto $[-t_{\circ},t_{\circ}]$, with slope in
$[\tfrac23,\tfrac43]$.
\item[(b)] Let $\hat w$ be real with $|\hat w|\le t_{\circ}-c_2=\varrho$. Then exactly one real
$\zeta$, alive in the sense of that theorem, has $\zeta_{(0)}=\hat w$. We denote it
$\zeta_*^n(\hat w;X)$, and we put $a_n(X):=\theta_n-\zeta_*^n(0;X)$.
\item[(c)] Along the run started at $\zeta_*^n(\hat w;X)$,
\begin{equation*}
|\zeta_{(r)}|\le|\hat w|+2c_2+\sum_{u<r}\iota_u .
\end{equation*}
\item[(d)] $|\zeta_*^n(\hat w)-\zeta_*^n(0)|\le\tfrac32|\hat w|$.
\end{itemize}

\emph{Comparison of runs.} Three further statements from the proof of the uniqueness theorem
are used; we recall them.

The first is the sup-distance bound. For a pair $(\mathsf{A},\mathsf{B})$ of runs, frozen and
matched at depth $r_0$ in the sense of that bound, put
$d_s:=\sup_{\bar{\mathfrak{P}}}\big|\zeta^{\mathsf{B}}_{(s)}-\zeta^{\mathsf{A}}_{(s)}\big|$.
Then
\begin{equation*}
d_{r+1}\le C'_E\Big[\frac{\rho^{n-r}}{1-\rho}+(r-r_0+1)\Pi_n\Big].
\end{equation*}
Here $C'_E$ is the prefactor of that bound, a constant multiple of its constant $C_E$. The set
$\bar{\mathfrak{P}}$ and the sequence $\Pi_n$ are those of that bound, and $\rho$ is the
contraction letter.

The second is the matching statement. Put $h_n(v):=\zeta^n_{(0)}(\zeta_*^n+v)$, where
$\zeta^n_{(0)}(\zeta)$ denotes the depth-$0$ deviation of the run of length $n$ started at the
bare deviation $\zeta$. For
$|v|\le r_J$ there is exactly one $\mathfrak{n}_n(v)$, real for real $v$, with
$h_{n+1}(\mathfrak{n}_n(v))=h_n(v)$. It satisfies
$|\mathfrak{n}_n(v)-v|\le\Upsilon_n|v|$, and $\sum_n\Upsilon_n<\infty$.

The third is the pinning statement. The pin and the match sit at depth $0$, so the
sup-distance bound applies with $r_0=0$. Moreover, $a_n$ does not depend on the torus.

\emph{Shifted reference couplings.} For $s\ge1$ we put $X_s:=\theta_s(X)=X+c_0s$ and
$g_s:=X_s^{-1/2}$. We also put
\begin{equation}\label{9.13:eq-canonical-scaffold-a}
D_s(g) := 2c_2 + \sum_{u<s}\iota_u(X) .
\end{equation}
Then $g_s<g$. Under the hypotheses of the canonical-scaffold theorem on $g$, moreover,
\begin{equation}\label{9.13:eq-canonical-scaffold-2}
D_s(g)\le 2c_2 + C_\iota\lambda^{-1}\theta_s^{1/2}(\log\theta_s)^{p_0} .
\end{equation}
If moreover $g\le\gamma_U$, then
\begin{equation}\label{9.13:eq-canonical-scaffold-3}
D_s(g)\le 2c_2 + s\,\iota_0(X) .
\end{equation}
\end{definition}

\begin{proof}
Since $s\ge1$ and $c_0s>0$, we have $X_s=X+c_0s>X$. Hence
$g_s=X_s^{-1/2}<X^{-1/2}=g$.

For \eqref{9.13:eq-canonical-scaffold-2}, the proof of the canonical-scaffold theorem bounds
the drift sum along the canonical scaffold by
\begin{equation*}
\sum_{u<s}\iota_u(X)\le C_\iota\lambda^{-1}\theta_s^{1/2}(\log\theta_s)^{p_0} .
\end{equation*}
Inserting this bound into the definition \eqref{9.13:eq-canonical-scaffold-a} gives
\eqref{9.13:eq-canonical-scaffold-2}.

For \eqref{9.13:eq-canonical-scaffold-3}, suppose $g\le\gamma_U$. Then the sequences
$e_{(u)}$ and $v_{(u)}$ of the running-shift statement in scaffold form are non-increasing in
$u$ along the scaffold. So for every $u\ge0$,
\begin{equation*}
\iota_u(X)=C_\beta e_{(u+1)}+v_{(u+1)}\le C_\beta e_{(1)}+v_{(1)}=\iota_0(X).
\end{equation*}
The sum in \eqref{9.13:eq-canonical-scaffold-a} has $s$ terms, namely $u=0,\dots,s-1$. Each is
at most $\iota_0(X)$, so the sum is at most $s\,\iota_0(X)$. This gives
\eqref{9.13:eq-canonical-scaffold-3}.
\end{proof}

\begin{lemma}[Re-rooting on the canonical scaffold]\label{9.13:lem-re-rooting}
We use the objects of Notation~\ref{9.13:not-canonical-scaffold}:
\begin{itemize}
\item the canonical scaffold $\Theta_X=(\theta_r(X))_{r\ge0}$, with $\theta_r(X)=X+c_0r$ and $X=g^{-2}$;
\item the runs on $\Theta_X$, with deviations $\zeta_{(r)}$, coefficients $b_{(r)}$ and running inverse couplings $x_{(r)}=\theta_r-\zeta_{(r)}$;
\item tubes and aliveness;
\item the tuned deviations $\zeta_*^n(\hat w;X)$ and the numbers $a_n(X)$;
\item the shifted reference inverse coupling $X_s=X+c_0s$, with $g_s=X_s^{-1/2}$;
\item the drift size $D_s(g)$ of \eqref{9.13:eq-canonical-scaffold-a}.
\end{itemize}
Fix $g\le\gamma_U(\varrho)$, $s\ge1$ and $n\ge s$.
\begin{enumerate}
\item[(a)] \emph{Shift covariance.} For all $r\ge0$,
\begin{equation}\label{9.13:eq-re-rooting-a}
\theta_{r+s}(X)=\theta_r(X_s).
\end{equation}
Hence every datum that the scaffold prescribes at depth $r+s$ of $\Theta_X$ is the datum at depth $r$ of $\Theta_{X_s}$. These data are:
\begin{itemize}
\item the parameters of Ba{\l}aban's step that Notation~\ref{9.13:not-canonical-scaffold} prescribes at each depth as functions of the value of $\theta$ at that depth;
\item the radii of \cite[(2.28)]{B-CMP119}, $\hat\alpha$, and $\check R_{(\cdot)}:=R(\theta_{\cdot})$;
\item the trees and the classes;
\item the quantities of \cite[(1.62)]{B-CMP122a};
\item the surrogate of \cite[(1.100)]{B-CMP122a} at coefficient $\theta_{\cdot}$.
\end{itemize}
None of these depends on $n$. Moreover, let $(\varrho_r)_{r\ge0}$ be a tube for $\Theta_X$, that is, a nondecreasing sequence with
\begin{equation*}
\varrho_{r+1}\,\frac{(\log\theta_r(X))^{2q}}{\theta_r(X)}\le\varpi_{\mathfrak{S}}
\qquad\text{for all } r\ge0 .
\end{equation*}
Then $(\varrho_{r+s})_{r\ge0}$ is a tube for $\Theta_{X_s}$.
\item[(b)] \emph{Truncation is a run.} Let $\zeta$ be a bare deviation. Let $(\zeta_{(r)},b_{(r)})_{0\le r\le n}$ be the length-$n$ run on $\Theta_X$ from $\zeta$. Let $(\zeta'_{(r')},b'_{(r')})_{0\le r'\le n-s}$ be the length-$(n-s)$ run on $\Theta_{X_s}$ from the same $\zeta$. Both runs have the same bare inverse coupling $\theta_n(X)-\zeta=\theta_{n-s}(X_s)-\zeta$. The second run is the first restricted to the depths $s\le r\le n$ and re-indexed by $r'=r-s$:
\begin{equation}\label{9.13:eq-re-rooting-b}
\zeta'_{(r')}=\zeta_{(r'+s)},\qquad b'_{(r')}=b_{(r'+s)}\qquad(0\le r'\le n-s).
\end{equation}
The second run is alive in the tube $(\varrho_{r'+s})_{r'}$ if and only if the first is alive at the depths $s\le r\le n$. Its depth-$0$ coupling $x'_{(0)}$ is the depth-$s$ coupling $x_{(s)}$ of the first.
\item[(c)] \emph{The tuned datum re-rooted.} Put $\zeta^n:=\zeta_*^n(0;X)$. Let $t_{n,s}:=\zeta^n_{(s)}$ be the depth-$s$ deviation of the length-$n$ run on $\Theta_X$ from $\zeta^n$. If $|t_{n,s}|\le\varrho$, then $\zeta^n=\zeta_*^{n-s}(t_{n,s};X_s)$, and
\begin{equation}\label{9.13:eq-re-rooting-c}
\begin{gathered}
a_n(X)=a_{n-s}(X_s)-\delta_{n,s},\qquad
\delta_{n,s}:=\zeta_*^{n-s}(t_{n,s};X_s)-\zeta_*^{n-s}(0;X_s)\in\mathbb{R},\\
\tfrac34\,|t_{n,s}|\le|\delta_{n,s}|\le\tfrac32\,|t_{n,s}| .
\end{gathered}
\end{equation}
Hence, for every $\hat w$,
\begin{equation}\label{9.13:eq-re-rooting-1}
a_n(X)-\hat w=a_{n-s}(X_s)-(\hat w+\delta_{n,s}).
\end{equation}
\item[(d)] \emph{Size.} Write $\zeta^n_{(r)}$ for the run from $\zeta^n$ and evaluate $\mathsf{q}_{(r)}$ along it. Let $s_i$ be the Ces\`aro defects of the one-loop coefficients in the bounded Ces\`aro-defect statement, recalled in \eqref{9.13:eq-re-rooting-2} below. Then
\begin{equation}\label{9.13:eq-re-rooting-d}
\begin{gathered}
|t_{n,s}|\le D_s(g),\\
t_{n,s}=(s_n-s_{n-s})
-\sum_{r<s}\Bigl(f_{n-r}\bigl(\zeta^n_{(n)},\dots,\zeta^n_{(r+1)}\bigr)+\mathsf{q}_{(r)}\Bigr).
\end{gathered}
\end{equation}
\end{enumerate}
\end{lemma}

\begin{proof}
\emph{(a)} By the definition of the canonical scaffold,
\begin{equation*}
\theta_{r+s}(X)=X+c_0(r+s)=(X+c_0s)+c_0r=\theta_r(X_s),
\end{equation*}
which is \eqref{9.13:eq-re-rooting-a}.

Each datum in the list is prescribed at a given depth as a function of the value of $\theta$ at that same depth, and of nothing depending on $n$. Its value at depth $r+s$ of $\Theta_X$ is therefore its prescription at the value $\theta_{r+s}(X)=\theta_r(X_s)$. That is the datum at depth $r$ of $\Theta_{X_s}$.

The tube condition is a condition on the pairs $(\theta_r,\varrho_{r+1})$, $r\ge0$, together with monotonicity, and it is invariant under the joint shift. Indeed, for $r\ge0$, by \eqref{9.13:eq-re-rooting-a},
\begin{equation*}
\varrho_{(r+1)+s}\,\frac{(\log\theta_r(X_s))^{2q}}{\theta_r(X_s)}
=\varrho_{(r+s)+1}\,\frac{(\log\theta_{r+s}(X))^{2q}}{\theta_{r+s}(X)}\le\varpi_{\mathfrak{S}} .
\end{equation*}
The inequality is the tube condition for $\Theta_X$ at the index $r+s$. The sequence $(\varrho_{r+s})_{r\ge0}$ is nondecreasing because it is a tail of a nondecreasing sequence.

\emph{(b)} Counted from the bare lattice, the step landing at depth $r$ of a length-$n$ run is Ba{\l}aban's step $k\to k+1$ with $k=n-r-1$. Here level $k\le n$ has depth $n-k$, and the lattice at depth $r$ is the same for all $n$.

By the running-shift statement on the scaffold, the coefficient of this step is
\begin{equation*}
b_{(r)}=\beta^0_{n-r}+f_{n-r}\bigl(\zeta_{(n)},\dots,\zeta_{(r+1)}\bigr)+\mathsf{q}_{(r)} .
\end{equation*}
It is torus-free. It is a function of three things:
\begin{itemize}
\item the level index $n-r$;
\item the deeper deviations $\zeta_{(u)}$, $u>r$ (equivalently, of the bare density and the $k$ previous steps);
\item the data prescribed by the scaffold at the depths read by those steps and by this step.
\end{itemize}
By causality of the renormalisation group, all of these depths are $\ge r$. The step $k\to k+1$ depends on nothing at levels above $k+1$, and on the canonical scaffold there is no further datum for it to depend on.

Now substitute $(r,n)=(r'+s,n'+s)$ with $n'=n-s$. Then:
\begin{itemize}
\item $n-r=n'-r'$;
\item the deviations $\zeta_{(u)}$, $u>r$, are the same numbers;
\item $\theta_u(X)=\theta_{u-s}(X_s)$ by (a), so the prescribed data at depth $u$ of $\Theta_X$ are those at depth $u-s$ of $\Theta_{X_s}$.
\end{itemize}

It follows that $r'\mapsto\zeta_{(r'+s)}(\zeta)$ satisfies the defining recursion of the length-$n'$ run on $\Theta_{X_s}$ with the same coefficients. We check this by downward induction on $r'$. At $r'=n'$ both sequences start from $\zeta'_{(n')}=\zeta=\zeta_{(n)}$. Assume $\zeta'_{(u')}=\zeta_{(u'+s)}$ for all $u'>r'$. Then $b'_{(r')}$ and $b_{(r'+s)}$ are the same function evaluated at the same level index, the same deeper deviations and the same prescribed data, so $b'_{(r')}=b_{(r'+s)}$. Therefore
\begin{equation*}
\zeta'_{(r')}=\zeta'_{(r'+1)}+b'_{(r')}-c_0=\zeta_{(r'+s+1)}+b_{(r'+s)}-c_0=\zeta_{(r'+s)} .
\end{equation*}
This proves \eqref{9.13:eq-re-rooting-b}.

The bare inverse couplings agree because $\theta_n(X)=\theta_{n-s}(X_s)$ by (a). The depth-$0$ coupling of the second run is
\begin{equation*}
x'_{(0)}=\theta_0(X_s)-\zeta'_{(0)}=\theta_s(X)-\zeta_{(s)}=x_{(s)} .
\end{equation*}
Aliveness of the second run in $(\varrho_{r'+s})_{r'}$ is the tube condition at its depths $0\le r'\le n-s$. By \eqref{9.13:eq-re-rooting-b} and (a), this is the tube condition of the first run at the depths $s\le r\le n$.

\emph{The stop depth of the torus integration.} Let $m_0$ be the integer such that the length-$n$ datum lives on the torus of side $2L^{m_0}\ell_g$. In the uniqueness theorem there is a stop depth $\mathsf{w}(\cdot)$: $n-\mathsf{w}\ge1$ steps are carried out on the torus, and the depths $r<\mathsf{w}$ carry frozen whole-lattice numbers only. This stop depth concerns where the torus integration stops, not the run. Every coefficient $b$ of the run is a whole-lattice number, because the running-shift statement on the scaffold is torus-free; these are the infinite-lattice coefficients of clause (v-a) of Theorem~\ref{3.1:thm-main}. By clause (iii) of the lemma placing the pin and the match at depth~$0$, $a_n$ depends on no torus and on no stop depth.

The torus floor in the scaffold version of the correlation theorem is inherited by the shifted frame. That floor at $(m_0,g)$ is $L^{m_0}\ge C_{\mathbb{T}}\sup_{u'\ge0}L^{-u'}R(\theta_{u'}(X))$. By \eqref{9.13:eq-re-rooting-a} it implies
\begin{align*}
C_{\mathbb{T}}\sup_{u\ge0}L^{-u}R(\theta_u(X_s))
&=L^{s}\,C_{\mathbb{T}}\sup_{u\ge0}L^{-(u+s)}R(\theta_{u+s}(X))\\
&\le L^{s}\,C_{\mathbb{T}}\sup_{u'\ge0}L^{-u'}R(\theta_{u'}(X))\le L^{m_0+s},
\end{align*}
which is the floor at $(m_0+s,g_s)$. In any case the floor is internal to the two invocations of the uniqueness theorem made in the proof of the limit-member identification, where the present lemma is used. No part of that theorem's internal analysis is compared there between the two frames. Only its conclusions, clause (d) of the uniqueness theorem, are invoked there: once at $(m_0,g)$ and once at $(m_0+s,g_s)$.

\emph{(c)} By (b), the run from $\zeta^n$ restricted to the depths $s\le r\le n$ is a real length-$(n-s)$ run on $\Theta_{X_s}$, with depth-$0$ deviation $\zeta'_{(0)}=t_{n,s}$. It is alive in the shifted tube $(\varrho_{r'+s})_{r'}$, because $\zeta^n$ is alive. This shifted tube is a legitimate tube for $\Theta_{X_s}$ by (a), and the scaffold version of the correlation theorem holds for every tube.

Along this run the drift inequality from the proof of the scaffold theorem holds. Write $\tilde\zeta'_{(r')}:=\zeta'_{(r')}+s_{(n-s)-r'}$. The inequality reads
\begin{equation*}
|\tilde\zeta'_{(r')}|\le|\tilde\zeta'_{(0)}|+\sum_{u<r'}\iota_u(X_s)\qquad(0\le r'\le n-s),
\end{equation*}
with $\tilde\zeta'_{(0)}=t_{n,s}+s_{n-s}$ and $|s_{n-s}|\le c_2$.

The scaffold theorem is used with the half-width $t_{\circ}=c_2+r_J$, so that $t_{\circ}-c_2=r_J=\varrho$. Suppose $|t_{n,s}|\le\varrho$. Then $|\tilde\zeta'_{(0)}|\le|t_{n,s}|+c_2\le t_{\circ}$, and the right-hand side above is at most $t_{r'}(X_s)$, for every $r'$, where $t_{r'}(X_s)$ denotes the half-width at depth $r'$ of the deviation set $J_0(\Theta_{X_s};n-s)$ of the scaffold theorem. That is, the run lies in $J_0(\Theta_{X_s};n-s)$. On this set two facts hold:
\begin{itemize}
\item $|\zeta'_{(r')}|\le t_{r'}(X_s)+c_2=\varrho_{r'}(X_s)-8r_J$, so the run is alive in the tube $(\varrho_{r'}(X_s))_{r'}$ attached to $\Theta_{X_s}$ itself;
\item $\zeta\mapsto\tilde\zeta'_{(0)}(\zeta)$ is a bijection onto $[-t_{\circ},t_{\circ}]$, by clause (a) of the scaffold theorem.
\end{itemize}
This is exactly the uniqueness argument in the proof of clause (b) of the scaffold theorem. Every real alive deviation whose run has depth-$0$ value $\hat w$ with $|\hat w|\le\varrho$ lies in $J_0(\Theta_{X_s};n-s)$, and there the map is injective. The two deviations $\zeta^n$ and $\zeta_*^{n-s}(t_{n,s};X_s)$ both have depth-$0$ value $t_{n,s}$, so $\zeta^n=\zeta_*^{n-s}(t_{n,s};X_s)$.

Then, by (a) and the definitions of $a_n$ and $a_{n-s}$,
\begin{align*}
a_n(X)&=\theta_n(X)-\zeta^n=\theta_{n-s}(X_s)-\zeta_*^{n-s}(t_{n,s};X_s)\\
&=\bigl(\theta_{n-s}(X_s)-\zeta_*^{n-s}(0;X_s)\bigr)-\delta_{n,s}
=a_{n-s}(X_s)-\delta_{n,s} .
\end{align*}
The number $\delta_{n,s}$ is real because both tuned deviations are real.

For the two-sided size we use clause (a) of the scaffold theorem. The map $\zeta\mapsto\zeta'_{(0)}(\zeta)=\tilde\zeta'_{(0)}(\zeta)-s_{n-s}$ has slope in $[\tfrac23,\tfrac43]$ on $J_0(\Theta_{X_s};n-s)$. Its image is $[-t_{\circ}-s_{n-s},\,t_{\circ}-s_{n-s}]$, which contains $[-\varrho,\varrho]$ because $|s_{n-s}|\le c_2$. So the inverse map $\hat w\mapsto\zeta_*^{n-s}(\hat w;X_s)$ has slope in $[\tfrac34,\tfrac32]$ on $[-\varrho,\varrho]$. Integrating its derivative from $0$ to $t_{n,s}$ gives $\tfrac34|t_{n,s}|\le|\delta_{n,s}|\le\tfrac32|t_{n,s}|$, which completes \eqref{9.13:eq-re-rooting-c}.

Finally, \eqref{9.13:eq-re-rooting-1} follows by subtracting $\hat w$ from both sides of the first identity in \eqref{9.13:eq-re-rooting-c}, since $\delta_{n,s}$ does not depend on $\hat w$.

\emph{(d)} Clause (c) of the scaffold theorem states that $|\zeta_{(r)}|\le|\hat w|+2c_2+\sum_{u<r}\iota_u(X)$ along the run from $\zeta_*^n(\hat w;X)$. Take $\hat w=0$ and $r=s$. This gives $|t_{n,s}|\le 2c_2+\sum_{u<s}\iota_u(X)$, which is $D_s(g)$ by \eqref{9.13:eq-canonical-scaffold-a}.

For the exact formula, sum the defining recursion $\zeta_{(r+1)}=\zeta_{(r)}-(b_{(r)}-c_0)$ of the run from $\zeta^n$ over $r<s$; this is the summation on which the drift inequality rests. Using $\zeta^n_{(0)}=0$, we get
\begin{equation*}
t_{n,s}=\zeta^n_{(s)}=-\sum_{r<s}\bigl(\beta^0_{n-r}-c_0\bigr)
-\sum_{r<s}\Bigl(f_{n-r}\bigl(\zeta^n_{(n)},\dots,\zeta^n_{(r+1)}\bigr)+\mathsf{q}_{(r)}\Bigr).
\end{equation*}
By the bounded Ces\`aro-defect statement for the one-loop coefficients, for $i<j$,
\begin{equation}\label{9.13:eq-re-rooting-2}
\sum_{l=i+1}^{j}\beta^0_l=(j-i)c_0+s_i-s_j,\qquad|s_i|\le c_2 .
\end{equation}
Taking $i=n-r-1$ and $j=n-r$ in \eqref{9.13:eq-re-rooting-2} gives $\beta^0_{n-r}-c_0=s_{n-r-1}-s_{n-r}$. The first sum therefore telescopes:
\begin{equation*}
\sum_{r<s}\bigl(s_{n-r-1}-s_{n-r}\bigr)=s_{n-s}-s_n .
\end{equation*}
Substituting this gives the second line of \eqref{9.13:eq-re-rooting-d}.
\end{proof}

\begin{remark}[The case $n=s$, and composition]
For $n=s$ the truncated run has length $0$, so $\zeta'_{(0)}=\zeta$. Then $a_0(X_s)=X_s$ and $\zeta_*^0(\hat w;X_s)=\hat w$. Whenever $|t_{s,s}|\le\varrho$, \eqref{9.13:eq-re-rooting-c} reads
\begin{equation*}
\delta_{s,s}=t_{s,s}=X_s-a_s(X).
\end{equation*}
By \eqref{9.13:eq-re-rooting-d}, and recalling $c_0=2\lambda$,
\begin{equation*}
a_s(X)-X-c_0s\in\Bigl[-2c_2-\sum_{u<s}\iota_u(X),\;2c_2+\sum_{u<s}\iota_u(X)\Bigr].
\end{equation*}
This is finer than the window $[\lambda s-2c_2,\,3\lambda s+2c_2]$ for $a_s(X)-X$ in clause (a) of the uniqueness theorem, and is consistent with it since $\iota_u\le\lambda$.

For composition, let $s'\ge1$ with $s+s'\le n$, and write $\delta_{n,s}(X)$ to record the reference inverse coupling. The lemma applies at the reference coupling $g_s$, since $g_s<g\le\gamma_U(\varrho)$. We also have $(X_s)_{s'}=X_s+c_0s'=X_{s+s'}$. Suppose (c) applies to $(n,s)$ at $X$, to $(n-s,s')$ at $X_s$, and to $(n,s+s')$ at $X$. Then \eqref{9.13:eq-re-rooting-c} gives
\begin{equation}\label{9.13:eq-re-rooting-3}
\delta_{n,s+s'}(X)=\delta_{n,s}(X)+\delta_{n-s,s'}(X_s),
\end{equation}
because both sides equal $a_{n-s-s'}(X_{s+s'})-a_n(X)$.
\end{remark}

\begin{lemma}[Convergence of the depth-$s$ data along the tuned family]\label{9.13:lem-depth-convergence}
Let $g\le\gamma_U(\varrho)$ and $s\ge 1$, let $X=g^{-2}$, and use the canonical scaffold $\Theta_X$,
$\theta_r=X+c_0r$, the runs on it, the shifted reference inverse coupling $X_s=X+c_0s$ with
$g_s=X_s^{-1/2}<g$, and the quantity $D_s(g)$ of \eqref{9.13:eq-canonical-scaffold-a}, all as in
Notation~\ref{9.13:not-canonical-scaffold}; recall $\varrho=r_J$. For $n\ge s$ we write $t_{n,s}$ for the
depth-$s$ deviation of the $g$-tuned run of length $n$, that is, of the run $(n,a_n(X))$ of length $n$ with
bare inverse coupling $a_n(X)$, and
\begin{equation*}
\delta_{n,s}:=a_{n-s}(X_s)-a_n(X)
\end{equation*}
for the label difference. For $Y\in\{X,X_s\}$ assume the following statements on the scaffold $\Theta_Y$,
under the standing hypotheses of the uniqueness theorem.
\begin{itemize}
\item The frozen-pair comparison bound: for a frozen pair of runs $\mathsf{A}$, of length $n$, and
$\mathsf{B}$, parametrised by an offset $v$ in the closure $\bar{\mathfrak{P}}$ of the disc
$\mathfrak{P}=D(0,r_J)$ and matched at depth $r_0$ (that is, having equal depth-$r_0$ deviations for
every $v$), the distances
$d_u:=\sup_{v\in\bar{\mathfrak{P}}}|\zeta^{\mathsf{B}}_{(u)}-\zeta^{\mathsf{A}}_{(u)}|$ satisfy, for $r<n$,
\begin{equation*}
d_{r+1}\le \mathsf{A}C_E\Bigl[\frac{\rho^{\,n-r}}{1-\rho}+(r-r_0+1)\,\Pi_n\Bigr],
\end{equation*}
where $\mathsf{A}C_E$ is the constant of that bound (the letter $\mathsf{A}$ in this constant is
unrelated to the run $\mathsf{A}$), $\rho<1$, and $(\Pi_n)$ is its non-negative tail
sequence.
\item The matching proposition: with $h_n(v):=\zeta^n_{(0)}(\zeta_*^n(0;Y)+v)$, where
$\zeta^n_{(0)}(\zeta)$ denotes the depth-$0$ deviation of the run of length $n$ on $\Theta_Y$ with bare
deviation $\zeta$, for every $|v|\le r_J$
there is exactly one $\mathfrak{n}_n(v)$ with $h_{n+1}(\mathfrak{n}_n(v))=h_n(v)$; it is real for real
$v$, and $|\mathfrak{n}_n(v)-v|\le\Upsilon_n(Y)\,|v|$.
\item The first clause of the pinning lemma: pin and match sit at depth $0$, so that the frozen-pair
comparison bound applies with $r_0=0$.
\item The uniqueness theorem, in the following parts: $T_\Pi:=\sup_{l\ge1}l\sum_{i\ge l}\Pi_i<\infty$;
$\mathfrak{n}_n(0)=0$ and $\Upsilon_n(Y)\le C\,n^{-3/2}(\log n)^{p_0+1}$; and the deviation maps are
$\tfrac43$-Lipschitz on the tube.
\item The scaffold theorem, clauses (a)--(d): (a) $\zeta\mapsto\tilde\zeta_{(0)}(\zeta)$ is a
real-analytic increasing bijection $J_0\to[-t_{\circ},t_{\circ}]$ with slope in $[\tfrac23,\tfrac43]$;
(b) for every real $|\hat w|\le\varrho$ exactly one real alive bare deviation
$\zeta=:\zeta_*^n(\hat w;Y)$ has $\zeta_{(0)}=\hat w$, and $a_n(Y)=\theta_n(Y)-\zeta_*^n(0;Y)$ is the
tuned bare inverse coupling of Notation~\ref{9.13:not-canonical-scaffold};
(c) along that run $|\zeta_{(r)}|\le|\hat w|+2c_2+\sum_{u<r}\iota_u$;
(d) $|\zeta_*^n(\hat w;Y)-\zeta_*^n(0;Y)|\le\tfrac32|\hat w|$.
\end{itemize}
Then the following hold.
\begin{enumerate}
\item[(i)] There is a real number $t_{\infty,s}(g)$ such that
\begin{equation}\label{9.13:eq-depth-convergence-a}
\begin{gathered}
t_{n,s}\longrightarrow t_{\infty,s}(g)\quad(n\to\infty),\\
\sum_{n\ge s}|t_{n+1,s}-t_{n,s}|
\le\sum_{n\ge s}\mathsf{A}C_E\Bigl[\frac{\rho^{\,n-s+1}}{1-\rho}+s\,\Pi_n\Bigr]<\infty .
\end{gathered}
\end{equation}
Hence the depth-$s$ inverse coupling $x^{(n)}_{(s)}=X_s-t_{n,s}$ of the $g$-tuned run of length $n$
converges, to $\hat x_s(g):=X_s-t_{\infty,s}(g)$. More generally, for every real $\hat w$ with
$|\hat w|\le\varrho$ the depth-$s$ inverse coupling $x^{(n)}_{(s)}(\hat w)$ of the run
$(n,a_n(X)-\hat w)$ converges as $n\to\infty$.
\item[(ii)] If $D_s(g)\le\tfrac23\varrho$, there is a real number $\delta_s(g)$ such that
\begin{equation}\label{9.13:eq-depth-convergence-b}
\begin{gathered}
\delta_{n,s}\longrightarrow\delta_s(g)\quad(n\to\infty),\qquad
\sum_{n\ge s}|\delta_{n+1,s}-\delta_{n,s}|<\infty,\\
\tfrac34|t_{\infty,s}|\le|\delta_s|\le\tfrac32|t_{\infty,s}|\le\tfrac32D_s(g).
\end{gathered}
\end{equation}
\end{enumerate}
\end{lemma}

\begin{proof}
All statements assumed above are applied under the standing hypotheses of the uniqueness theorem,
which contain the hypotheses of the frozen-pair comparison bound. Since $g_s<g\le\gamma_U(\varrho)$,
the matching proposition and the scaffold theorem are available on $\Theta_{X_s}$ as well as on
$\Theta_X$. For a length $n'$ and a bare inverse coupling $y$ we write
$\zeta^{(n')}_{(s)}(y)$ for the depth-$s$ deviation of the run $(n',y)$ on $\Theta_X$, as a function of
the bare inverse coupling $y=\theta_{n'}-\zeta$; thus
$t_{n,s}=\zeta^{(n)}_{(s)}(a_n(X))$.

\emph{Proof of} (i). Fix $n\ge s$. For $v\in\bar{\mathfrak{P}}$, where $\mathfrak{P}=D(0,r_J)$, consider
on $\Theta_X$ the frozen pair
\begin{equation*}
\mathsf{A}:=\bigl(n,\;a_n(X)-v\bigr),\qquad
\mathsf{B}:=\bigl(n+1,\;a_{n+1}(X)-\mathfrak{n}_n(v)\bigr).
\end{equation*}
Since $a_n(X)-v=\theta_n-(\zeta_*^n(0;X)+v)$, the bare deviation of $\mathsf{A}$ is
$\zeta_*^n(0;X)+v$ and its depth-$0$ deviation is $h_n(v)$; likewise the depth-$0$ deviation of
$\mathsf{B}$ is $h_{n+1}(\mathfrak{n}_n(v))$. By the defining property of $\mathfrak{n}_n$ in the matching
proposition these coincide, so the pair is matched at depth $0$ for every $v$. The depths of the two runs
are aligned at the reference end, and their levels are indexed by the level index of $\mathsf{A}$, so that
depth $s$ corresponds to level $n-s$. By the first clause of the pinning lemma the frozen-pair comparison
bound applies with $r_0=0$; we take $r=s-1$, which satisfies $r<n$ because $n\ge s$. Since
$n-r=n-s+1$ and $r-r_0+1=s$, it gives
\begin{equation}\label{9.13:eq-depth-convergence-1}
d_s\le\mathsf{A}C_E\Bigl[\frac{\rho^{\,n-s+1}}{1-\rho}+s\,\Pi_n\Bigr].
\end{equation}

At $v=0$ the uniqueness theorem gives $\mathfrak{n}_n(0)=0$, so that $\mathsf{A}=(n,a_n(X))$ and
$\mathsf{B}=(n+1,a_{n+1}(X))$ are the $g$-tuned runs of lengths $n$ and $n+1$. As $0\in\bar{\mathfrak{P}}$,
\begin{equation*}
|t_{n+1,s}-t_{n,s}|=|\zeta^{\mathsf{B}}_{(s)}(0)-\zeta^{\mathsf{A}}_{(s)}(0)|\le d_s ,
\end{equation*}
and summing \eqref{9.13:eq-depth-convergence-1} over $n\ge s$ gives the inequality in
\eqref{9.13:eq-depth-convergence-a}. The right-hand side is finite. The first part is the geometric series
$\sum_{n\ge s}\rho^{\,n-s+1}/(1-\rho)=\rho/(1-\rho)^2$, finite since $\rho<1$. For the second part, the
uniqueness theorem gives $T_\Pi=\sup_{l\ge1}l\sum_{i\ge l}\Pi_i<\infty$; the term $l=1$ gives
$\sum_{i\ge1}\Pi_i\le T_\Pi$, hence $\sum_{n\ge s}s\,\Pi_n\le s\,T_\Pi$. The increments of
$(t_{n,s})_{n\ge s}$ are therefore absolutely summable. The sequence is
Cauchy and has a limit $t_{\infty,s}(g)$.

By definition of the running inverse coupling, $x_{(s)}=\theta_s-\zeta_{(s)}$, and $\theta_s(X)=X+c_0s=X_s$.
Hence $x^{(n)}_{(s)}=X_s-t_{n,s}$, which converges to $X_s-t_{\infty,s}(g)=\hat x_s(g)$.

Now let $\hat w$ be real with $|\hat w|\le\varrho=r_J$, so $\hat w\in\bar{\mathfrak{P}}$, and take $v=\hat w$
in the pair above. Then $\mathsf{A}=(n,a_n(X)-\hat w)$, while $\mathsf{B}$ sits at the offset
$\mathfrak{n}_n(\hat w)$, which is real because $\hat w$ is. By the triangle inequality,
\begin{align*}
&\bigl|\zeta^{(n+1)}_{(s)}(a_{n+1}(X)-\hat w)-\zeta^{(n)}_{(s)}(a_n(X)-\hat w)\bigr|\\
&\qquad\le\bigl|\zeta^{(n+1)}_{(s)}(a_{n+1}(X)-\mathfrak{n}_n(\hat w))
-\zeta^{(n)}_{(s)}(a_n(X)-\hat w)\bigr|\\
&\qquad\quad+\bigl|\zeta^{(n+1)}_{(s)}(a_{n+1}(X)-\mathfrak{n}_n(\hat w))
-\zeta^{(n+1)}_{(s)}(a_{n+1}(X)-\hat w)\bigr| .
\end{align*}
The first term on the right is $|\zeta^{\mathsf{B}}_{(s)}(\hat w)-\zeta^{\mathsf{A}}_{(s)}(\hat w)|\le d_s$,
and $d_s$ is bounded by \eqref{9.13:eq-depth-convergence-1}. For the second term, the deviation maps are
$\tfrac43$-Lipschitz on the tube (uniqueness theorem) and
$|\mathfrak{n}_n(\hat w)-\hat w|\le\Upsilon_n(X)|\hat w|$ (matching proposition). The second term is
therefore at most $\tfrac43\Upsilon_n(X)|\hat w|\le\tfrac43\varrho\,\Upsilon_n(X)$. The uniqueness theorem
gives $\Upsilon_n(X)\le C\,n^{-3/2}(\log n)^{p_0+1}$, so $\sum_n\Upsilon_n(X)<\infty$. Together with the
summability of the bound \eqref{9.13:eq-depth-convergence-1} shown above, the depth-$s$ deviations of the
runs $(n,a_n(X)-\hat w)$ form a sequence with absolutely summable increments, hence a convergent one.
Consequently $x^{(n)}_{(s)}(\hat w)=X_s-\zeta^{(n)}_{(s)}(a_n(X)-\hat w)$ converges.

\emph{Proof of} (ii). Assume $D_s(g)\le\tfrac23\varrho$. Work on $\Theta_{X_s}$. For a length $n'$ and real
$t$ with $|t|\le\varrho$ put
\begin{equation*}
v_{n'}(t):=\zeta_*^{n'}(t;X_s)-\zeta_*^{n'}(0;X_s),
\end{equation*}
with $\zeta_*^{n'}$ from clause (b) of the scaffold theorem on $\Theta_{X_s}$. This is the real bare offset
with $h_{n'}(v_{n'}(t))=t$, where $h_{n'}$ is the map of the matching proposition on $\Theta_{X_s}$.
Indeed, $h_{n'}(v_{n'}(t))$ is the depth-$0$ deviation of the run with bare deviation
$\zeta_*^{n'}(t;X_s)$, which is $t$ by clause (b). Clause (d) gives $|v_{n'}(t)|\le\tfrac32|t|$.

The map $t\mapsto v_{n'}(t)$ is $\tfrac32$-Lipschitz on $[-\varrho,\varrho]$. To see this, note that
$\zeta_*^{n'}(t;X_s)$ is the preimage of $t$ under $\zeta\mapsto\zeta_{(0)}(\zeta)$. This map differs from
$\zeta\mapsto\tilde\zeta_{(0)}(\zeta)$ by an additive constant independent of $\zeta$, so by clause (a) it
is increasing with slope at least $\tfrac23$, and its inverse has slope at most $\tfrac32$.

By clause (c) of the scaffold theorem along the $g$-tuned run of length $n$ on $\Theta_X$ (where
$\hat w=0$), together with \eqref{9.13:eq-canonical-scaffold-a}, we have
\begin{equation*}
|t_{n,s}|\le2c_2+\sum_{u<s}\iota_u(X)=D_s(g)\le\tfrac23\varrho\qquad(n\ge s).
\end{equation*}
In particular $v_{n-s}(t_{n,s})$ is defined.

The label difference is
\begin{equation}\label{9.13:eq-depth-convergence-2}
\delta_{n,s}=v_{n-s}(t_{n,s}).
\end{equation}
Since $\theta_n(X)=X+c_0n=\theta_{n-s}(X_s)$, and $a_n(X)=\theta_n(X)-\zeta_*^n(0;X)$,
$a_{n-s}(X_s)=\theta_{n-s}(X_s)-\zeta_*^{n-s}(0;X_s)$, the identity
\eqref{9.13:eq-depth-convergence-2} is the statement
\begin{equation*}
\zeta_*^n(0;X)=\zeta_*^{n-s}(t_{n,s};X_s).
\end{equation*}
To see it, we use the shift-covariance clause of Lemma~\ref{9.13:lem-re-rooting}: $\theta_{r+s}(X)=\theta_r(X_s)$
for all $r\ge0$, and every datum of the scaffold at depth $r+s$ of $\Theta_X$ is the datum at depth $r$ of
$\Theta_{X_s}$. The run of length $n-s$ on $\Theta_{X_s}$ with bare deviation $\zeta_*^n(0;X)$ has bare
inverse coupling $\theta_{n-s}(X_s)-\zeta_*^n(0;X)=a_n(X)$; it is therefore the $g$-tuned run of length $n$
on $\Theta_X$ read from depth $s$, it is real and alive on $\Theta_{X_s}$ because the $g$-tuned run is real
and alive on $\Theta_X$, and its depth-$0$ deviation is $t_{n,s}$. Since $|t_{n,s}|\le\varrho$, clause (b)
of the scaffold theorem on $\Theta_{X_s}$ at length $n-s$ identifies its bare deviation with
$\zeta_*^{n-s}(t_{n,s};X_s)$.

Next we compare consecutive lengths at fixed $t$. By the matching proposition on $\Theta_{X_s}$, for real
$|v|\le r_J$ the offset $\mathfrak{n}_{n'}(v)$ is real,
\begin{equation*}
h_{n'+1}(\mathfrak{n}_{n'}(v))=h_{n'}(v),\qquad
|\mathfrak{n}_{n'}(v)-v|\le\Upsilon_{n'}(X_s)\,|v| .
\end{equation*}
Let $|t|\le\tfrac23\varrho$ and $v=v_{n'}(t)$. Then $|v|\le\tfrac32|t|\le\varrho$, and
$h_{n'+1}(\mathfrak{n}_{n'}(v))=h_{n'}(v_{n'}(t))=t$. Clause (b) at length $n'+1$ says that exactly one
real alive bare deviation has depth-$0$ deviation $t$. Hence $v_{n'+1}(t)=\mathfrak{n}_{n'}(v_{n'}(t))$.
Therefore
\begin{equation}\label{9.13:eq-depth-convergence-3}
|v_{n'+1}(t)-v_{n'}(t)|\le\Upsilon_{n'}(X_s)\,|v_{n'}(t)|\le\tfrac32\,\Upsilon_{n'}(X_s)\,|t| .
\end{equation}

We split the increment of $\delta_{n,s}$ using \eqref{9.13:eq-depth-convergence-2}:
\begin{align*}
|\delta_{n+1,s}-\delta_{n,s}|
&\le|v_{n+1-s}(t_{n+1,s})-v_{n+1-s}(t_{n,s})|+|v_{n+1-s}(t_{n,s})-v_{n-s}(t_{n,s})|\\
&\le\tfrac32\,|t_{n+1,s}-t_{n,s}|+\Upsilon_{n-s}(X_s)\,\varrho .
\end{align*}
The first term was bounded by the $\tfrac32$-Lipschitz property of $v_{n+1-s}$ on $[-\varrho,\varrho]$. The
second was bounded by \eqref{9.13:eq-depth-convergence-3} with $n'=n-s$ and $t=t_{n,s}$, using
$\tfrac32|t_{n,s}|\le\varrho$.

The first term is summable over $n\ge s$ by \eqref{9.13:eq-depth-convergence-a}. The second is summable
because the uniqueness theorem on $\Theta_{X_s}$ gives
$\Upsilon_{n'}(X_s)\le C\,n'^{-3/2}(\log n')^{p_0+1}$. Hence
$\sum_{n\ge s}|\delta_{n+1,s}-\delta_{n,s}|<\infty$, and $\delta_{n,s}$ converges to a limit $\delta_s(g)$.

The bounds on the sizes in \eqref{9.13:eq-depth-convergence-b} come from \eqref{9.13:eq-re-rooting-c} of
Lemma~\ref{9.13:lem-re-rooting}, passed to the limit $n\to\infty$ along the convergent sequences
$t_{n,s}\to t_{\infty,s}(g)$ and $\delta_{n,s}\to\delta_s(g)$.
\end{proof}

\begin{remark}[The two tuned runs share one bare coupling]\label{9.13:rem-shared-bare-coupling}
Let $g$, $s\ge 1$ and $n\ge s$ be as in Lemma~\ref{9.13:lem-re-rooting}. Write $X=g^{-2}$ and
$X_s=X+c_0 s=g_s^{-2}$, and use the tuned runs, their deviations and the tuned bare data $a_n$ of
Notation~\ref{9.13:not-canonical-scaffold}. We compare the tuned bare datum of the $g_s$-tuned run
of length $n-s$ with that of the $g$-tuned run of length $n$, read in units shifted by $s$ depths.
\begin{itemize}
\item[(a)] The depth-$s$ running inverse coupling of the $g$-tuned run of length $n$ is
\begin{equation}\label{9.13:eq-shared-bare-coupling-1}
x_{(s)}=X+\sum_{r<s}b_{(r)}=X_s-t_{n,s},
\end{equation}
where
\begin{equation}\label{9.13:eq-shared-bare-coupling-2}
t_{n,s}:=\zeta_*^n(0;X)_{(s)}=(s_n-s_{n-s})-\sum_{r<s}\bigl(f_{n-r}+\mathsf{q}_{(r)}\bigr),
\end{equation}
whereas the $g_s$-tuned run of length $n-s$ has depth-$0$ inverse coupling equal to $X_s$. The
two tuned data coincide if and only if $t_{n,s}=0$, and this fails in general, for two reasons:
\begin{itemize}
\item[(D1)] the bare-end defect $s_n-s_{n-s}$ of the one-loop Ces\`aro defects $(s_i)_i$;
\item[(D2)] the higher-order drift $\sum_{r<s}(f_{n-r}+\mathsf{q}_{(r)})$, of modulus at most
$\sum_{u<s}\iota_u$.
\end{itemize}
\item[(b)] One bare inverse coupling belongs to both tuned families: with the real
number $\delta_{n,s}:=a_{n-s}(X_s)-a_n(X)$, for every offset $\hat w$ the member of the $g$-tuned
family of length $n$ at offset $\hat w$ and the member of the $g_s$-tuned family of length $n-s$ at
offset $\hat w+\delta_{n,s}$ have the same bare inverse coupling,
\begin{equation}\label{9.13:eq-shared-bare-coupling-3}
a_n(X)-\hat w=a_{n-s}(X_s)-(\hat w+\delta_{n,s}),
\qquad
\tfrac34\,|t_{n,s}|\le|\delta_{n,s}|\le\tfrac32\,|t_{n,s}|.
\end{equation}
\item[(c)] $\delta_{n,s}$ converges as $n\to\infty$.
\end{itemize}
Thus the bare inverse coupling $a_n(X)-\hat w$ is one and the same number for the two runs, while
the offsets $\hat w$ and $\hat w+\delta_{n,s}$ at which the two tuned families read it differ by an
amount that converges as $n\to\infty$; the two families do not agree offset for offset.
\end{remark}

\begin{proof}
\emph{Part (a).} By Notation~\ref{9.13:not-canonical-scaffold} the tuned bare datum is
$a_n:=\theta_n-\zeta_*^n(0;X)_{(n)}$, and it is pinned by the requirement that the run's inverse
coupling at its own depth $0$ is $x_{(0)}=X$; equivalently, its deviation at depth $0$ is
$\zeta_{(0)}=0$. Read this run from depth $s$: its new depth-$0$ inverse coupling is its old
depth-$s$ inverse coupling $x_{(s)}$, which by the definition of the per-depth coefficients
$b_{(r)}$ is $x_{(0)}+\sum_{r<s}b_{(r)}$; since $x_{(0)}=X$, this is the first
equality in \eqref{9.13:eq-shared-bare-coupling-1}. The scaffold form of the flow expansion of the
correlation theorem, applied with $\zeta_{(0)}=0$, together with the expression of the deviation
through its one-loop, higher-order and remainder parts, gives the second equality in
\eqref{9.13:eq-shared-bare-coupling-1} with $t_{n,s}$ as in
\eqref{9.13:eq-shared-bare-coupling-2}.

The $g_s$-tuned run of length $n-s$ is pinned in the same way at its own depth $0$, now at the
reference inverse coupling $X_s$, so its depth-$0$ inverse coupling is $X_s$. Comparing with
\eqref{9.13:eq-shared-bare-coupling-1}, the two tuned data coincide if and only if $t_{n,s}=0$.

By \eqref{9.13:eq-shared-bare-coupling-2}, $t_{n,s}$ is the bare-end defect (D1) minus the drift
(D2). For (D1), the bound on the one-loop Ces\`aro defect gives $|s_i|\le c_2$ for every $i$, hence
$|s_n-s_{n-s}|\le 2c_2$, and it contains no statement on the convergence of the sequence
$(s_i)_i$. For (D2), the drift is the higher-order part of the deviation by which the true run
leaves the linear scaffold $\Theta_X$; by the per-depth drift bound it has modulus at most
$\sum_{u<s}\iota_u$, while the whole deviation is at most $2c_2+\sum_u\iota_u$. Neither defect is
made to vanish by the construction, which is the sense in which $t_{n,s}=0$ fails in general.

\emph{Part (b).} The identity in \eqref{9.13:eq-shared-bare-coupling-3} is the definition of
$\delta_{n,s}$: substituting $a_{n-s}(X_s)=a_n(X)+\delta_{n,s}$ into the right-hand side returns
$a_n(X)-\hat w$. That $\delta_{n,s}$ is real and satisfies the two-sided comparison with
$|t_{n,s}|$ in \eqref{9.13:eq-shared-bare-coupling-3} is part~(c) of Lemma~\ref{9.13:lem-re-rooting},
display~\eqref{9.13:eq-re-rooting-c}. Hence the bare inverse coupling of the $g$-tuned run of
length $n$ at offset $\hat w$ is the bare inverse coupling of the $g_s$-tuned run of length $n-s$ at
offset $\hat w+\delta_{n,s}$.

\emph{Part (c).} The convergence of $\delta_{n,s}$ as $n\to\infty$ is
Lemma~\ref{9.13:lem-depth-convergence}. Its mechanism is the following. The consecutive-length
estimate in the proof of the
uniqueness theorem controls the full coefficient $b$ at equal depth between the $g$-tuned runs
$\mathsf{A}$ and $\mathsf{B}$ of consecutive lengths $n$ and $n+1$, the defects (D1) and (D2)
included, through the bound comparing the coefficients of the two runs at equal depth:
$|b^{\mathsf{B}}_{(r)}-b^{\mathsf{A}}_{(r)}|$ is at most $C_E(\rho^{\,n-r}+\Pi_n)$ times the
constant that multiplies $C_E$ in item~(i) of Lemma~\ref{9.13:lem-depth-convergence}, where $C_E$
and $\rho$ are the constants of that estimate and $(\Pi_n)_n$ is its summable tail sequence. The
offset stability of item~(1) of
clause (v-c) of Theorem~\ref{3.1:thm-main} carries the resulting convergence of the depth-$s$
deviations $t_{n,s}$ to the convergence of $\delta_{n,s}$.
\end{proof}

\begin{theorem}[The limit-member identification]\label{9.13:thm-limit-member}
Let $L\ge L_0$ be odd, and let $\mathfrak{p}$, $\gamma$ and $\gamma_H$ be as in hypotheses
(H2)--(H3) of Notation~\ref{2.2:not-hypotheses}. Let $\varrho>0$, let $\gamma_U(\varrho)$ be the
coupling threshold of the uniqueness theorem (clause~(v) of Theorem~\ref{3.1:thm-main}),
chosen after
$\varrho$, and let $g\le\gamma_U(\varrho)$; put $X=g^{-2}$ and $\varrho_2=\varrho/2$. Let
$m_0\ge 0$ be such that hypothesis (H5) holds for the torus $\mathbb{T}$ of side
$2L^{m_0}\ell_g$ on which clause (ii-$\mathbb{T}$-a) is read. The uniqueness theorem itself is
stated for every $m_0\ge 0$ and asks no lower bound on the torus exponent. Let $s\ge 1$, and
let $X_s=X+c_0s$, $g_s=X_s^{-1/2}$ and $D_s(g)$ be as in
Notation~\ref{9.13:not-canonical-scaffold} and \eqref{9.13:eq-canonical-scaffold-a}.

We use the following notation. For a bare inverse coupling $x$, integers $n,m\ge 0$ and a
tuple $\vec h=(h_1,\dots,h_{\mathsf{p}})$ of test functions on the torus of side $2L^{m}$,
$S^{n}_{m}[x](\vec h)$ is the connected Schwinger function $S^T_{\mathsf{p};n,m}$ of the
observables $\mathcal{O}(h_1),\dots,\mathcal{O}(h_{\mathsf{p}})$ under the Wilson measure at
bare inverse coupling $x$ on the bare lattice of spacing $L^{-n}$ on that torus, that is, of
$2L^{m+n}$ sites a side. The labels $a_n(\cdot)$ are those of the uniqueness theorem, so that
$a_n(X)-\hat w$ is the bare inverse coupling of the tuned run of length $n$ at reference $g$
and offset $\hat w$. For a real offset $w$ and a reference coupling $g'$,
$S^{(w)}_{m}[g'](\vec h)$ is the limit of the tuned family at reference $g'$ and offset $w$
given by the uniqueness theorem; for complex $w$ in the strip of the third statement of (II)
below, $S^{(w)}_{m}[g'](\vec h)$ is its holomorphic extension there. For $n\ge s$,
$\delta_{n,s}$ is the real label difference of Lemma~\ref{9.13:lem-re-rooting}, and
$\delta_s(g)$ is its limit as $n\to\infty$, as in Lemma~\ref{9.13:lem-depth-convergence}, when
this limit exists.
Finally, $x^{(n)}_{(r)}(\hat w)$ is the running inverse coupling at depth $r$ of the run of
length $n$ with bare inverse coupling $a_n(X)-\hat w$, and $\hat x$ is the canonical
comparison point of the member with data $(m_0,g,\hat w)$.

Assume the following statements.
\begin{itemize}
\item[(I)] The change-of-units identity. For every bare inverse coupling $x_0$, every $m\ge 0$,
every tuple $\vec h$ of test functions on the torus of side $2L^m$ and every $n\ge s$,
\begin{equation*}
S^{n}_{m}[x_0](\vec h)=S^{n-s}_{m+s}[x_0](\vec h\circ L^{-s}\cdot),
\end{equation*}
where $\vec h\circ L^{-s}\cdot=(h_1(L^{-s}\cdot),\dots,h_{\mathsf{p}}(L^{-s}\cdot))$. The
right member reads the same plaquettes with the reference length moved $s$ levels toward the
cutoff.
\item[(II)] The uniqueness theorem, at the data $(m_0,g)$ and at the data $(m_0+s,g_s)$. At
a datum $(m_1,g')$ equal to either of these, and for a tuple $\vec h$ of real Lipschitz functions
on the torus of side $2L^{m_1}$ with pairwise separated supports (an admissible tuple of the
uniqueness theorem), we use four of its statements.
\begin{itemize}
\item[--] First: for every real $w$ with $|w|<\varrho_2/2$, the sequence
$S^{n}_{m_1}[a_n(g'^{-2})-w](\vec h)$ converges to $S^{(w)}_{m_1}[g'](\vec h)$ as
$n\to\infty$. The convergence holds along the full sequence, at rate
$n^{-1/2}(\log n)^{p_0+1}$.
\item[--] Second: for every sequence of reals $a'_n$ with
$|a'_n-a_n(g'^{-2})|\le\varrho_2/2$,
\begin{equation*}
S^{n}_{m_1}[a'_n](\vec h)-S^{(a_n(g'^{-2})-a'_n)}_{m_1}[g'](\vec h)\to 0
\qquad(n\to\infty).
\end{equation*}
\item[--] Third: the map $w\mapsto S^{(w)}_{m_1}[g'](\vec h)$ is real-analytic on
$\mathopen]-\varrho_2/2,\varrho_2/2\mathclose[$. It is holomorphic on the strip
$\{|\operatorname{Re}w|<\varrho_2/2,\ |\operatorname{Im}w|<\upsilon\}$, where the strip height
is $\upsilon=\upsilon_w$ at $(m_0,g)$ and $\upsilon=\upsilon'_w$ at $(m_0+s,g_s)$.
\item[--] Fourth (from the bi-Lipschitz lemma): for every real $w$ with $|w|<\varrho_2/2$, the
depth-$0$ running inverse coupling of the run of length $n$ with bare inverse coupling
$a_n(g'^{-2})-w$ converges, as $n\to\infty$, to the canonical comparison point of the member
with data $(m_1,g',w)$.
\end{itemize}
\item[(III)] The statements of the uniqueness theorem's comparison of two runs from which
Lemma~\ref{9.13:lem-depth-convergence} is derived. These are the comparison estimate with
constants $\mathsf{A}C_E$, contraction letter $\rho$ and tail sequence $\Pi_n$; the offset
maps and the matching maps, with the summable rate $\Upsilon_n$ of the matching; and the
pinning statement for the offset maps.
\item[(IV)] The tube theorem on the canonical scaffold $\Theta_X$ of
Notation~\ref{9.13:not-canonical-scaffold}, and the correlation theorem in its scaffold form
on $\Theta_X$, read with a tube of $\Theta_X$. Of the correlation theorem in its scaffold form
we use three clauses.
\begin{itemize}
\item[(a)] For every run, the partition function and the connected Schwinger functions are
the same numbers for every scaffold, every tube and every choice of widths.
\item[(b)] The asymptotic-freedom clause for the running inverse couplings $x^{(n)}_{(r)}$ of
a run on the scaffold, in terms of the per-depth coefficients $b_{(r)}$ of the run.
\item[(c)] The bound $|b_{(r)}-\beta^0_{n-r}|\le\lambda$ of the running-shift statement,
where $\beta^0_{n-r}$ is the one-loop part of $b_{(r)}$.
\end{itemize}
\item[(V)] The one-loop Ces\`aro defect bound: the one-loop Ces\`aro defects of the run are
bounded by the constant $c_2$.
\end{itemize}
Assume finally the two one-sided conditions
\begin{equation}\label{9.13:eq-limit-member-a}
\varrho\ge 48c_2\qquad\text{and}\qquad\sum_{u<s}\iota_u(X)\le\frac{\varrho}{24}.
\end{equation}
The first condition is a lower bound on $\varrho$. It is imposed after $L$ and $\mathfrak{p}$
are fixed and before $\gamma_U(\varrho)$ is chosen.

The second condition is read in two ways. We have $\sum_{u<s}\iota_u(X)\le s\,\iota_0(X)$, and
$\iota_0(X)$, which equals $C_\beta\,e_{(1)}+v_{(1)}$ with $e_{(1)}$ and $v_{(1)}$ the
constituents of the per-depth drift bound of Notation~\ref{9.13:not-canonical-scaffold},
decreases to $0$ as $g\downarrow 0$. Hence,
for each $s$, the second condition is an upper bound $g\le\gamma_{\mathrm{dil}}(s;\varrho)$
with $\gamma_{\mathrm{dil}}(s;\varrho)\le\gamma_U(\varrho)$. Equivalently, for each $g$, it is
an upper bound $s\le s_{\mathrm{dil}}(g;\varrho)$. Here $s_{\mathrm{dil}}(g;\varrho)$ does not
decrease as $g$ decreases and tends to infinity as $g\to 0$. Neither bound is evaluated.

Then, for every $\mathsf{p}\ge 2$, every $\mathsf{d}>0$, all real Lipschitz functions
$f_1,\dots,f_{\mathsf{p}}$ on $\mathbb{T}$ whose supports are pairwise at distance at least
$\mathsf{d}$, and every real $\hat w$ with $|\hat w|<\varrho_2/4=\varrho/8$, the following
hold. Here $\vec f=(f_1,\dots,f_{\mathsf{p}})$.
\begin{itemize}
\item[($\alpha$)] Finite cutoff, exactly: for every $n\ge s$,
\begin{equation}\label{9.13:eq-limit-member-b}
S^{n}_{m_0}[a_n(X)-\hat w](\vec f)
=S^{n-s}_{m_0+s}\bigl[a_{n-s}(X_s)-(\hat w+\delta_{n,s})\bigr](\vec f\circ L^{-s}\cdot).
\end{equation}
Both members are one lattice integral: over $2L^{m_0+n}$ sites a side, at one bare inverse
coupling, with plaquette weights $f_i(x(p))$. The left member names it at reference $g$ on
the torus of exponent $m_0$. The right member names it at reference $g_s$ on the same torus,
now of exponent $m_0+s$.
\item[($\beta$)] The limit:
\begin{equation}\label{9.13:eq-limit-member-c}
S^{(\hat w)}_{m_0}[g](\vec f)=S^{(\hat w+\delta_s(g))}_{m_0+s}[g_s](\vec f\circ L^{-s}\cdot).
\end{equation}
The limit $\delta_s(g)$ exists. Both members are full limits, with no passage to
subsequences, of the tuned families at $(m_0,g)$ and at $(m_0+s,g_s)$. Moreover
$|\hat w+\delta_s(g)|<\varrho/8+\varrho/8=\varrho_2/2$.
\item[($\gamma$)] The reference couplings: the limit
\begin{equation}\label{9.13:eq-limit-member-d}
\hat x_s:=\lim_{n\to\infty}x^{(n)}_{(s)}(\hat w)\ \text{exists, and}\quad
\hat x_s-\hat x\in\bigl[\lambda^{\sharp}s-c_5,\ B^{\sharp}s\bigr].
\end{equation}
The right member of \eqref{9.13:eq-limit-member-c} is the member whose limiting
reference-scale inverse coupling is $\hat x_s$. Here $\hat x_s$ is the limit of the depth-$s$
inverse couplings of the tuned runs of the left member. It depends on $\hat w$; the runs in
question are those with bare inverse coupling $a_n(X)-\hat w$, and at $\hat w=0$, where these
are the $g$-tuned runs, $\hat x_s$ is the limit $\hat x_s(g)$ of
Lemma~\ref{9.13:lem-depth-convergence}. Thus the member
with data
$(m_0,\hat x)$ is identified exactly, as a continuum limit, with the member with data
$(m_0+s,\hat x_s)$; and $\hat x_s$ is a full limit along the tuned family.
\item[($\delta$)] Holomorphic extension: both members of \eqref{9.13:eq-limit-member-c} are
holomorphic in $\hat w$ on
\begin{equation*}
\bigl\{\,|\operatorname{Re}\hat w|<\varrho/8,\ |\operatorname{Im}\hat w|<
\min\{\upsilon_w,\upsilon'_w\}\,\bigr\},
\end{equation*}
and they agree there.
\item[($\varepsilon$)] Inheritance of (H5): hypothesis (H5) for $\mathbb{T}$ at reference $g$,
with exponent $m_0$, implies (H5) for the same torus at reference $g_s$, with exponent
$m_0+s$.
\end{itemize}
\end{theorem}

\begin{proof}
Fix $\mathsf{p}$, $\mathsf{d}$, the tuple $\vec f$ and the offset $\hat w$ as in the
statement.

\emph{Consequences of \eqref{9.13:eq-limit-member-a}.} By the first condition,
$2c_2\le\varrho/24$. With the second condition and the definition
\eqref{9.13:eq-canonical-scaffold-a} of $D_s(g)$ this gives
\begin{equation}\label{9.13:eq-limit-member-1-bis}
D_s(g)=2c_2+\sum_{u<s}\iota_u(X)\le\frac{\varrho}{24}+\frac{\varrho}{24}=\frac{\varrho}{12}.
\end{equation}
In particular $D_s(g)<\varrho$, $D_s(g)\le\frac{2}{3}\varrho$ and
$\frac{3}{2}D_s(g)\le\varrho/8$. Since $c_0s>0$ we have $X_s>X$, so $g_s<g\le\gamma_U(\varrho)$.
Thus the uniqueness theorem, with the same $\varrho$, is available at $(m_0,g)$ and at
$(m_0+s,g_s)$; these are the two invocations in (II).

\emph{The rescaled tuple.} Put $\vec f':=\vec f\circ L^{-s}\cdot$, that is,
$f'_i(y)=f_i(L^{-s}y)$. The length unit at reference $g_s$ is $\ell_{g_s}=L^{-s}\ell_g$. Hence
the torus of side $2L^{m_0}\ell_g$ is the torus of side $2L^{m_0+s}$ in $\ell_{g_s}$-units:
it is the same physical torus. Each $f'_i$ is a real Lipschitz function on it, and
\begin{equation*}
\|f'_i\|_\infty=\|f_i\|_\infty,\qquad
\operatorname{Lip}f'_i=L^{-s}\operatorname{Lip}f_i\le\operatorname{Lip}f_i .
\end{equation*}
From the definition $\|f\|_{\sharp}=\max\{\|f\|_\infty,40M\operatorname{Lip}f\}$ it follows that
$\|f'_i\|_{\sharp}\le\|f_i\|_{\sharp}$. The supports of the $f'_i$ are the supports of the
$f_i$ dilated by $L^s$. They are therefore pairwise at distance at least $L^s\mathsf{d}\ge\mathsf{d}$.
So $\vec f'$ is an admissible tuple for the uniqueness theorem at $(m_0+s,g_s)$, which is
stated for every torus exponent.

\emph{Proof of ($\alpha$).} Let $n\ge s$. We first describe how the two members of
\eqref{9.13:eq-limit-member-b} are read.
\begin{itemize}
\item[--] The left member is an integral over the bare lattice of spacing $L^{-n}$ on the
torus of side $2L^{m_0}$ in $\ell_g$-units.
\item[--] The right member is an integral over the bare lattice of spacing $L^{-(n-s)}$ on the
torus of side $2L^{m_0+s}$ in $\ell_{g_s}$-units.
\item[--] Both lattices have $2L^{m_0+n}$ sites a side, and they are the same lattice.
\end{itemize}

The observable $\mathcal{O}(f)=\sum_p f(x(p))[1-\operatorname{Re}\operatorname{tr}U(\partial p)]$
carries no power of the lattice spacing. Let $p$ be a plaquette whose barycentre has lattice
coordinates $i$. Then $x(p)=L^{-n}i$ in $\ell_g$-units and $x(p)=L^{-(n-s)}i$ in
$\ell_{g_s}$-units. Moreover
\begin{equation*}
f'_j\bigl(L^{-(n-s)}i\bigr)=f_j\bigl(L^{-s}L^{-(n-s)}i\bigr)=f_j\bigl(L^{-n}i\bigr).
\end{equation*}
Hence $\mathcal{O}(f_j)$ read at reference $g$ and $\mathcal{O}(f'_j)$ read at reference $g_s$
are the same function of the bond variables $U$, for each $j$. This is the content of the
change-of-units identity (I).

The bare measure is the Wilson measure at one bare inverse coupling. By the identity
\eqref{9.13:eq-re-rooting-c} of Lemma~\ref{9.13:lem-re-rooting}, this coupling carries the two
labels
\begin{equation}\label{9.13:eq-limit-member-2-bis}
a_n(X)-\hat w=a_{n-s}(X_s)-(\hat w+\delta_{n,s}).
\end{equation}
The re-rooting lemma gives that identity for every $n\ge s$ under the condition
$D_s(g)<\varrho$, which holds because $D_s(g)\le\varrho/12$ by
\eqref{9.13:eq-limit-member-1-bis}. Apply (I) with $x_0=a_n(X)-\hat w$, $m=m_0$ and
$\vec h=\vec f$. Then write $x_0$ in the right member by \eqref{9.13:eq-limit-member-2-bis}. This
gives \eqref{9.13:eq-limit-member-b}.

It remains to identify the right member. By clause (a) of (IV), the partition function and
the connected Schwinger functions of a run of $n-s$ steps are the same numbers for every
scaffold, every tube and every choice of widths. Therefore the right
member of \eqref{9.13:eq-limit-member-b} is the finite-cutoff member of the uniqueness theorem
at the datum $(m_0+s,g_s)$. It has $n-s$ steps, bare datum $a_{n-s}(X_s)-(\hat w+\delta_{n,s})$,
and test functions $\vec f'$, all read in the $g_s$-frame.

\emph{Proof of ($\beta$).} We compute the limit of \eqref{9.13:eq-limit-member-b} as
$n\to\infty$ from each side.

\emph{The left member.} Apply the first statement of (II) at $(m_0,g)$ with $w=\hat w$; this is
allowed because $|\hat w|<\varrho/8\le\varrho_2/2$. Hence, as $n\to\infty$,
\begin{equation*}
S^{n}_{m_0}[a_n(X)-\hat w](\vec f)\to S^{(\hat w)}_{m_0}[g](\vec f),
\end{equation*}
along the full sequence, at rate $n^{-1/2}(\log n)^{p_0+1}$.

\emph{The right member: first comparison.} Put $n':=n-s$ and
\begin{equation*}
a'_{n'}:=a_{n'}(X_s)-(\hat w+\delta_{n'+s,s})\in\mathbb{R}.
\end{equation*}
By the identity \eqref{9.13:eq-re-rooting-c} and the bound \eqref{9.13:eq-re-rooting-d} of
Lemma~\ref{9.13:lem-re-rooting}, $|\delta_{n'+s,s}|\le\frac{3}{2}D_s(g)$. Together with
\eqref{9.13:eq-limit-member-1-bis} this gives
\begin{equation}\label{9.13:eq-limit-member-3}
\bigl|a'_{n'}-a_{n'}(X_s)\bigr|=\bigl|\hat w+\delta_{n'+s,s}\bigr|
\le\frac{\varrho}{8}+\frac{3}{2}D_s(g)\le\frac{\varrho}{8}+\frac{\varrho}{8}=\frac{\varrho_2}{2}.
\end{equation}
Note that $a_{n'}(X_s)=a_{n'}(g_s^{-2})$. By \eqref{9.13:eq-limit-member-3}, the second
statement of (II) at $(m_0+s,g_s)$ applies to the sequence $a'_{n'}$ and the tuple $\vec f'$.
Since $a_{n'}(X_s)-a'_{n'}=\hat w+\delta_{n'+s,s}$, it gives
\begin{equation}\label{9.13:eq-limit-member-4}
S^{n'}_{m_0+s}[a'_{n'}](\vec f')
-S^{(\hat w+\delta_{n'+s,s})}_{m_0+s}[g_s](\vec f')\to 0\qquad(n'\to\infty).
\end{equation}

\emph{The right member: convergence of the offsets.} By \eqref{9.13:eq-limit-member-1-bis},
$D_s(g)\le\varrho/12\le\frac{2}{3}\varrho$. So the statement
\eqref{9.13:eq-depth-convergence-b} of Lemma~\ref{9.13:lem-depth-convergence} applies. (That
lemma is proved as an implication from the statements listed in (III).) It gives
$\delta_{n,s}\to\delta_s(g)$. In particular $\delta_s(g)$ exists, and
\begin{equation*}
\hat w+\delta_{n'+s,s}\to\hat w+\delta_s(g)\qquad(n'\to\infty).
\end{equation*}
The bound $|\delta_{n,s}|\le\frac{3}{2}D_s(g)\le\varrho/8$ holds for every $n\ge s$, so it
passes to the limit: $|\delta_s(g)|\le\varrho/8$. Hence all the points $\hat w+\delta_{n'+s,s}$,
and their limit $\hat w+\delta_s(g)$, lie in the compact interval
\begin{equation*}
\Bigl\{w\in\mathbb{R}:\ |w|\le|\hat w|+\frac{\varrho}{8}\Bigr\}
\subset\Bigl]-\frac{\varrho_2}{2},\frac{\varrho_2}{2}\Bigr[,
\end{equation*}
since $|\hat w|+\varrho/8<\varrho/4=\varrho_2/2$. The same inequality gives the bound
$|\hat w+\delta_s(g)|<\varrho_2/2$ stated in ($\beta$).

\emph{The right member: continuity.} By the third statement of (II) at $(m_0+s,g_s)$, the map
$w\mapsto S^{(w)}_{m_0+s}[g_s](\vec f')$ is real-analytic on
$\mathopen]-\varrho_2/2,\varrho_2/2\mathclose[$.
Hence it is continuous on the compact interval above, and therefore uniformly continuous
there. Consequently
\begin{equation*}
S^{(\hat w+\delta_{n'+s,s})}_{m_0+s}[g_s](\vec f')\to
S^{(\hat w+\delta_s(g))}_{m_0+s}[g_s](\vec f')\qquad(n'\to\infty).
\end{equation*}
Combined with \eqref{9.13:eq-limit-member-4}, the right member of
\eqref{9.13:eq-limit-member-b} converges, along the full sequence, to
$S^{(\hat w+\delta_s(g))}_{m_0+s}[g_s](\vec f')$.

\emph{Conclusion of ($\beta$).} By ($\alpha$), the two members of
\eqref{9.13:eq-limit-member-b} form one sequence indexed by $n\ge s$. We have found two limits
for it, and a sequence has at most one limit. Hence the two limits are equal, which is
\eqref{9.13:eq-limit-member-c}.

\emph{Proof of ($\gamma$).} The existence of $\hat x_s=\lim_n x^{(n)}_{(s)}(\hat w)$ is
\eqref{9.13:eq-depth-convergence-a} of Lemma~\ref{9.13:lem-depth-convergence}, applicable since
$|\hat w|\le\varrho$.

We next bound the slope of the tuned run between depths $0$ and $s$, at finite $n$. We combine
the asymptotic-freedom clause (IV)(b), the bound $|b_{(r)}-\beta^0_{n-r}|\le\lambda$ of the
running-shift statement (IV)(c), where $\lambda=c_0/2$, and
the one-loop Ces\`aro defect bound (V). Together they give, for every $n\ge s$,
\begin{equation}\label{9.13:eq-limit-member-5}
x^{(n)}_{(s)}(\hat w)-x^{(n)}_{(0)}(\hat w)\in\bigl[\lambda s-2c_2,\ 3\lambda s+2c_2\bigr].
\end{equation}
This is the two-sided bound on the running couplings given by the correlation theorem. Written
with the constants $\lambda^{\sharp}$, $B^{\sharp}$ of Definition~\ref{1.2:def-flow-constants}
and $c_5=2c_2$, as in Theorem~\ref{3.1:thm-main}, it reads
\begin{equation*}
\lambda^{\sharp}s-c_5\le x^{(n)}_{(s)}(\hat w)-x^{(n)}_{(0)}(\hat w)\le B^{\sharp}s.
\end{equation*}
Both end-points converge as $n\to\infty$:
\begin{itemize}
\item[--] $x^{(n)}_{(0)}(\hat w)\to\hat x$, by the fourth statement of (II) at $(m_0,g)$,
taken from the bi-Lipschitz lemma of clause~(v);
\item[--] $x^{(n)}_{(s)}(\hat w)\to\hat x_s$, by \eqref{9.13:eq-depth-convergence-a}.
\end{itemize}
Hence the differences converge to $\hat x_s-\hat x$. A closed interval is closed, so the limit
lies in the same interval, $[\lambda^{\sharp}s-c_5,\,B^{\sharp}s]$. This proves
\eqref{9.13:eq-limit-member-d}.

It remains to show that the right member of \eqref{9.13:eq-limit-member-c} is the member with
reference inverse coupling $\hat x_s$. Consider the runs whose finite-cutoff members converge
to it: by ($\alpha$), these are the runs with bare data $a_n(X)-\hat w$, read in the
$g_s$-frame. By the statement \eqref{9.13:eq-re-rooting-b} of Lemma~\ref{9.13:lem-re-rooting},
their depth-$0$ couplings in the $g_s$-frame are the depth-$s$ couplings in the $g$-frame of the
same bare data, namely $x^{(n)}_{(s)}(\hat w)$. The limit of these couplings is $\hat x_s$.

\emph{Proof of ($\delta$).} Consider first the left member of \eqref{9.13:eq-limit-member-c}.
By the third statement of (II) at $(m_0,g)$, it is holomorphic in $\hat w$ on
$\{|\operatorname{Re}\hat w|<\varrho_2/2,\ |\operatorname{Im}\hat w|<\upsilon_w\}$. This set
contains the set displayed in ($\delta$), since $\varrho/8<\varrho_2/2$.

Consider next the right member. If $|\operatorname{Re}\hat w|<\varrho/8$, then, as
$\delta_s(g)$ is real with $|\delta_s(g)|\le\varrho/8$,
\begin{equation*}
|\operatorname{Re}(\hat w+\delta_s(g))|<\frac{\varrho}{4}=\frac{\varrho_2}{2},\qquad
\operatorname{Im}(\hat w+\delta_s(g))=\operatorname{Im}\hat w .
\end{equation*}
So, by the third statement of (II) at $(m_0+s,g_s)$, the right member is holomorphic in
$\hat w$ wherever in addition $|\operatorname{Im}\hat w|<\upsilon'_w$.

Hence both members are holomorphic on the set displayed in ($\delta$). That set is a
rectangle, hence a connected open set. It contains the real interval
$\mathopen]-\varrho/8,\varrho/8\mathclose[$, on which the two members agree by ($\beta$). By
the identity theorem
for holomorphic functions, they agree on the whole set.

\emph{Proof of ($\varepsilon$).} By the definition $g_-(g')=(g'^{-2}+B^{\sharp})^{-1/2}$, we
have $g_-(g)^{-2}=X+B^{\sharp}$ and
\begin{equation*}
g_-(g_s)^{-2}=X_s+B^{\sharp}=X+B^{\sharp}+c_0s=\theta_s(X+B^{\sharp}),
\end{equation*}
where $\theta_r(Y)=Y+c_0r$ is the canonical scaffold of $Y$
(Notation~\ref{9.13:not-canonical-scaffold}), here taken with $Y=X+B^{\sharp}$. Thus
$g_-(g)$ and $g_-(g_s)$ are the couplings at depths $0$ and $s$ of the canonical scaffold of
$X+B^{\sharp}$.

Put $\check R'_{(r)}:=R\bigl(\theta_r(X+B^{\sharp})^{-1/2}\bigr)$ for the ball radius read at
depth $r$ of that scaffold, so that
$\check R'_{(0)}=R(g_-(g))$ and $\check R'_{(s)}=R(g_-(g_s))$. By the growth property of the
radii prescribed along a scaffold, applied to the scaffold of $X+B^{\sharp}$, the radius grows
by a factor in $\{1,L\}$ per depth:
\begin{equation*}
\check R'_{(r+1)}/\check R'_{(r)}\in\{1,L\}\quad\text{for every } r.
\end{equation*}
Hence $\check R'_{(s)}\le L^s\check R'_{(0)}$, that is,
\begin{equation*}
\log_L R(g_-(g_s))\le\log_L R(g_-(g))+s.
\end{equation*}

The torus exponent is $m_{\mathbb{T}}(g')=11+m'+\log_L R(g')$. Using this, and then (H5) for
$\mathbb{T}$ at reference $g$ (which gives $m_{\mathbb{T}}(g_-(g))\le m_0$), we obtain
\begin{equation*}
m_{\mathbb{T}}(g_-(g_s))\le m_{\mathbb{T}}(g_-(g))+s\le m_0+s .
\end{equation*}
This is (H5) for the same torus at reference $g_s$, with exponent $m_0+s$.
\end{proof}

\begin{remark}[The window for the change of units and its direction]
\label{9.13:rem-dilation-window}
Let $\varrho>0$, $g\le\gamma_U(\varrho)$, $m_0\ge0$, $s\ge 1$ and $\hat w$ be as in the limit-member
identification, Theorem~\ref{9.13:thm-limit-member}, $m_0$ being the torus exponent on the $g$-side.
Write $X=g^{-2}$, $X_s=X+c_0s$ for the
shifted reference inverse coupling and $g_s=X_s^{-1/2}$ for its coupling, and let $D_s(g)$ be the
depth-$s$ drift size of \eqref{9.13:eq-canonical-scaffold-a}. For $n>s$ let $t_{n,s}$ denote the
depth-$s$ deviation of the $g$-tuned run of length $n$ from the canonical scaffold, so that the
depth-$s$ inverse coupling of that run is $x^{(n)}_{(s)}=X_s-t_{n,s}$, and let $t_{\infty,s}$ be
its limit as $n\to\infty$. Let $s_i$ be the one-loop Ces\`aro defects (the subscript indexes the
defect sequence; no confusion with the depth $s$ arises), which measure the approach of
the one-loop coefficients to $c_0$ and obey $|s_i|\le c_2$, and let $\iota_u$ be the per-depth
drift bounds.
\begin{itemize}
\item[(a)] The proof of Theorem~\ref{9.13:thm-limit-member} uses only the change-of-units
identity, clause~(v) of Theorem~0 (the uniqueness theorem) at the two reference couplings $g$
and $g_s$, both at most $\gamma_U(\varrho)$, together with statements proved, under the hypotheses
of the uniqueness theorem, in the course of its proof, the scaffold theorem and the scaffold form
of the running-shift statement of the correlation theorem on the canonical scaffold, the bound
$|s_i|\le c_2$, and the bi-Lipschitz lemma, the last only where the canonical comparison point
$\hat x$ is named and not for \eqref{9.13:eq-limit-member-c}. It does not use the old-age tail
statement or any bound on an old-age tail, any rarity statement for large-field blocks, clustering,
a mass gap, uniqueness in the volume, an injectivity statement, the scaffold comparison, the bound
$\varpi_{\mathfrak{S}}\le1$ on the tube constant or the corollary of the uniqueness theorem (the
first-passage level never enters), or any comparison between two scaffolds, because the canonical
scaffold $\Theta_X$, shifted by $s$ depths, is the canonical scaffold $\Theta_{X_s}$.
\item[(b)] The $g$-tuned run leaves the linear scaffold: $|t_{n,s}|\le D_s(g)$ for every $n$, hence
$|t_{\infty,s}|\le D_s(g)$. The family of clause~(v) at the reference coupling $g_s$, on the other
hand, is parametrised by bare offsets $w$ with $|w|\le\varrho_2/2$ in a window that does not grow
with $s$:
the tube of the scaffold theorem grows with the depth, the window of the uniqueness theorem does
not. Condition \eqref{9.13:eq-limit-member-a} says exactly that the drifted member still lies in the
window of the $g_s$-family.
\item[(c)] Two alternatives to \eqref{9.13:eq-limit-member-a} are not adopted in this book. The first
is
\begin{equation}\label{9.13:eq-dilation-window-a}
g_s\le\gamma_U(\varrho_s),\qquad
\varrho_s:=48c_2+24\sum_{u<s}\iota_u(g^{-2})+8|\hat w| ,
\end{equation}
under which clause~(v) is invoked at $g_s$ with its own radius $\varrho_s$; this condition is
one-sided, and whether it holds for all $s$ at fixed $g$ depends on how $\gamma_U(\varrho)$
depends on $\varrho$, which is not made explicit. The second is the hypothesis that the sequence
$(s_k)$ of one-loop Ces\`aro defects converges, a statement on the approach of the one-loop
coefficients to $c_0$ beyond the bound $|s_i|\le c_2$; it is not among the statements on which the
proof of Theorem~\ref{9.13:thm-limit-member} depends. Under it every occurrence of $c_2$ in (b) is
deleted, $\limsup_{n\to\infty}|t_{n,s}|\le\sum_{u<s}\iota_u$, and the lower bound that
\eqref{9.13:eq-limit-member-a} places on $\varrho$ disappears, while the ceiling
$s\le s_{\mathrm{dil}}(g;\varrho)$ on $s$ stays.
\item[(d)] A shift by $s\ge1$ moves the reference length toward the cutoff: the reference coupling
$g_s$ is weaker than $g$ and the torus exponent $m_0+s$ is larger than $m_0$. This is the direction in
which Theorem~0 is read. The opposite direction, toward the infrared at fixed $m_0$, is available for
$s\le m_0$ steps as long as the stronger reference coupling stays at most $\gamma_U(\varrho)$, by the
same proof with the roles of the two couplings exchanged; these are finitely many integer steps, and
no continuous change of units is obtained.
\end{itemize}
\end{remark}

\begin{proof}
Item (a) describes the proof of Theorem~\ref{9.13:thm-limit-member}, read off from
that proof: the only identification of two scaffolds it needs is
$\theta_{r+s}=X+c_0(r+s)=X_s+c_0r$, which says that the canonical scaffold of $X$ read from depth
$s$ on is the canonical scaffold of $X_s$.

For item (b), in the proof of Theorem~\ref{9.13:thm-limit-member} the depth-$s$ deviation of the
$g$-tuned run of length $n$ is written, using the scaffold form of the running-shift statement with
vanishing deviation at depth $0$, as
\begin{equation*}
t_{n,s}=(s_n-s_{n-s})-\sum_{r<s}\bigl(f_{n-r}+\mathsf{q}_{(r)}\bigr),
\end{equation*}
the first term being the defect of the one-loop coefficients at the bare end and the second the
higher-order drift, with $\bigl|\sum_{r<s}(f_{n-r}+\mathsf{q}_{(r)})\bigr|\le\sum_{u<s}\iota_u$.
Since $|s_n|\le c_2$ and $|s_{n-s}|\le c_2$, the first term has modulus at most $2c_2$, and therefore
\begin{equation*}
|t_{n,s}|\le 2c_2+\sum_{u<s}\iota_u(X)=D_s(g)
\end{equation*}
by \eqref{9.13:eq-canonical-scaffold-a}; this is the bound, recorded with the canonical scaffold, by
which true runs leave $\Theta_X$: at most $2c_2+\sum_{u<s}\iota_u$. The bound is uniform in $n$ and
passes to
the limit $t_{\infty,s}$. The window of the $g_s$-family is the set of bare offsets $w$ with
$|w|\le\varrho_2/2$,
fixed by the radius $\varrho$ at which clause~(v) is invoked, and $\varrho$ does not depend on $s$,
whereas $D_s(g)$ grows with $s$ through $\sum_{u<s}\iota_u$. The member of the $g$-family read at
depth $s$ is the member of the $g_s$-family at the offset displaced by the drift, and
\eqref{9.13:eq-limit-member-a} is the condition that this displaced offset lies in the window.

For item (c), observe first that $48c_2+24\sum_{u<s}\iota_u(g^{-2})=24D_s(g)$, so that
$\varrho_s=24D_s(g)+8|\hat w|$ grows with $s$ exactly as the drift bound of (b) does. If
\eqref{9.13:eq-dilation-window-a} holds, clause~(v) applies at $g_s$ with radius $\varrho_s$, and the
proof of Theorem~\ref{9.13:thm-limit-member} runs verbatim with $(\varrho,\varrho_2)$ replaced by
$(\varrho_s,\varrho_s/2)$ on the $g_s$-side. This replacement does not change the bare labels
$a_{n'}(X_s)$ of the $g_s$-tuned runs of length $n'$, which are independent of the radius: by the
scaffold theorem the tuned run is the
unique run in its tube, and the unique run in a larger tube is the same run. The condition
\eqref{9.13:eq-dilation-window-a} bounds $g_s$ from above only; at fixed $g$ its validity for all $s$
is a question about the dependence of $\gamma_U(\varrho)$ on $\varrho$. If instead the sequence
$(s_k)$ converges, then $s_n-s_{n-s}\to0$ as $n\to\infty$ at fixed $s$, and the identity displayed
above gives
\begin{equation*}
\limsup_{n\to\infty}|t_{n,s}|\le\sum_{u<s}\iota_u ,
\end{equation*}
so every contribution $c_2$ is deleted; with it the floor on $\varrho$ in
\eqref{9.13:eq-limit-member-a} disappears, while the drift $\sum_{u<s}\iota_u$ still grows with $s$,
and the ceiling $s\le s_{\mathrm{dil}}(g;\varrho)$ stays.

For item (d), $c_0>0$ gives $X_s=X+c_0s>X$ for $s\ge1$, hence $g_s<g$; the torus exponent on the
$g_s$-side is $m_0+s$. Exchanging the roles of the two couplings, one compares the family at
$g$ on the torus exponent $m_0$ with the family at the stronger coupling $(X-c_0s)^{-1/2}$, which
requires $X>c_0s$, on the
torus exponent $m_0-s$, which requires $m_0-s\ge0$, that is $s\le m_0$, and requires clause~(v) to
apply at the stronger coupling, that is $(X-c_0s)^{-1/2}\le\gamma_U(\varrho)$. Under these
conditions the argument of Theorem~\ref{9.13:thm-limit-member} applies unchanged. Since $s$ is an
integer with $s\le m_0$, only finitely many such shifts exist, and each is a shift by a whole number
of steps.
\end{proof}

\begin{corollary}[Convergence of intermediate couplings and composition of unit changes]
\label{9.13:cor-intermediate-couplings}
Assume the setting of the limit-member identification (Theorem~\ref{9.13:thm-limit-member}); in
particular $\varrho>0$ and $g\le\gamma_U(\varrho)$ are fixed there. Put $X=g^{-2}$, $X_s=X+c_0s$ and
$g_s=X_s^{-1/2}$ for
$s\ge 0$. For $n\ge s\ge 1$ let $\delta_{n,s}(g)=a_{n-s}(X_s)-a_n(X)$ be the label difference, and
let $\delta_s(g)=\lim_{n\to\infty}\delta_{n,s}(g)$ be its limit as it enters
Theorem~\ref{9.13:thm-limit-member}.
\begin{enumerate}
\item[(a)] For every fixed $s\ge1$ and every real $\hat w$ with $|\hat w|\le\varrho$, the depth-$s$ inverse
coupling $x^{(n)}_{(s)}(\hat w)$ of the run of length $n$ with bare datum $a_n(X)-\hat w$ converges as
$n\to\infty$. For $\hat w=0$, which is the $g$-tuned run, its limit is the number $\hat x_s(g)$ of
Lemma~\ref{9.13:lem-depth-convergence}. Let $\hat x=\hat x(g)$ be the canonical comparison point at
$g$ and $\hat x_s=\hat x_s(g)$, the argument $g$ of these two numbers being dropped from now on, and let
$g^{(n)}_{n-s}=(x^{(n)}_{(s)})^{-1/2}$ be the
coupling at level $n-s$ of the $g$-tuned run of length $n$. Then
\begin{equation}\label{9.13:eq-intermediate-couplings-1}
g^{(n)}_{n-s}\longrightarrow g_{\infty,s}:=\hat x_s^{-1/2}\quad(n\to\infty),\qquad
\hat x_s-\hat x\in[\lambda^{\sharp}s-c_5,\;B^{\sharp}s].
\end{equation}
In particular, along the $g$-tuned family, the limit couplings $g_{\infty,s}$ of the corollary to
clause~(0) of Theorem~\ref{3.1:thm-main} are limits of the whole sequence and not only of
subsequences.
\item[(b)] For all $s,s'\ge1$ and $n\ge s+s'$,
\begin{equation}\label{9.13:eq-intermediate-couplings-2}
\delta_{n,s+s'}(g)=\delta_{n,s}(g)+\delta_{n-s,s'}(g_s)
\end{equation}
holds exactly. Consequently $\delta_{s+s'}(g)=\delta_s(g)+\delta_{s'}(g_s)$.
\item[(c)] The identifications \eqref{9.13:eq-limit-member-c} compose: the identification for the unit
change by $s$ at reference coupling $g$, followed by the identification for the unit change by $s'$ at
reference coupling $g_s$, equals the identification for the unit change by $s+s'$ at reference
coupling $g$, whenever all three lie inside the window \eqref{9.13:eq-limit-member-a}. Thus the additive
semigroup $\mathbb{N}$ of $L$-adic unit changes toward the cutoff acts on the family of members indexed
by $(m_0,g,\hat w)$, compatibly with the limits, inside the windows.
\end{enumerate}
The following are not asserted, in agreement with clause~(v) of Theorem~\ref{3.1:thm-main}:
\begin{itemize}
\item convergence of the effective actions along the tuned family;
\item exponential rates of convergence other than those stated at the canonical comparison point $\hat x$;
\item the non-degeneracy item that clause~(v) of Theorem~\ref{3.1:thm-main} lists as not asserted;
\item uniformity in the torus;
\item convergence when $s$ grows with $n$;
\item any relation, from the identifications \eqref{9.13:eq-limit-member-c} or their composition,
between distinct values of $\hat x$ at one torus exponent $m_0$ and one reference length;
\item a continuous dilation, or dimensional transmutation;
\item injectivity of $\hat w\mapsto S^{(\hat w)}$;
\item any identification outside the window \eqref{9.13:eq-limit-member-a}.
\end{itemize}
\end{corollary}

\begin{proof}
\emph{Part (a).} Take first $\hat w=0$.

\begin{itemize}
\item Lemma~\ref{9.13:lem-depth-convergence}(i) applies at the fixed pair $(g,s)$, since
$g\le\gamma_U(\varrho)$ and $s\ge1$.
\item It gives $t_{n,s}\to t_{\infty,s}(g)$, with the constants $\mathsf{A}C_E$, $\rho$ and the
summable sequence $(\Pi_n)_n$ of that lemma,
\begin{equation*}
\sum_n|t_{n+1,s}-t_{n,s}|\le\sum_n\mathsf{A}C_E\Bigl[\frac{\rho^{\,n-s+1}}{1-\rho}+s\Pi_n\Bigr]<\infty .
\end{equation*}
\item Hence $x^{(n)}_{(s)}=X_s-t_{n,s}$ converges to $\hat x_s=X_s-t_{\infty,s}(g)$.
\item The coupling at level $n-s$ of the $g$-tuned run is $g^{(n)}_{n-s}=(x^{(n)}_{(s)})^{-1/2}$. Its
convergence to $g_{\infty,s}=\hat x_s^{-1/2}$ follows by continuity of $y\mapsto y^{-1/2}$.
\item By \eqref{9.13:eq-limit-member-d} of the limit-member identification, the limit $\hat x_s$ is the
canonical comparison point of the member at $(m_0+s,g_s)$, the member written $(m_0+s,\hat x_s)$ in
clause~(v) of Theorem~\ref{3.1:thm-main}. The inclusion
$\hat x_s-\hat x\in[\lambda^{\sharp}s-c_5,B^{\sharp}s]$ is the one stated in
clause~(v) of Theorem~\ref{3.1:thm-main} for the members $(m_0,\hat x)$ and $(m_0+s,\hat x_s)$,
where it is obtained from the change-of-units identity together with the running-shift statement.
\end{itemize}

No subsequence was extracted at any point. Hence the limits $g_{\infty,s}$ named in the corollary to
clause~(0) are, along the $g$-tuned family, limits of the full sequence.

Now let $\hat w$ be real with $|\hat w|\le\varrho$. The proof of
Lemma~\ref{9.13:lem-depth-convergence}(i) rests on three inputs, each among the named hypotheses of
that lemma:
\begin{itemize}
\item the displayed bound on the distance $d_s$ between the deviations of the two runs
compared in Definition~\ref{9.6:def-comparison-chain}, which it takes from the
proposition bounding $d_s$ in the proof of the uniqueness theorem;
\item the named hypothesis of Lemma~\ref{9.13:lem-depth-convergence} that the proof of its part~(i)
combines with this bound;
\item the summability $\sum_n\Upsilon_n<\infty$ of the rates of the matching maps $\mathfrak{n}_n$.
\end{itemize}
We read the first input with bare offset $\hat w$ in place of $0$, and with depth offset zero. With the
second and third inputs, the
argument of Lemma~\ref{9.13:lem-depth-convergence}(i) gives a summable bound on the consecutive
differences
\begin{equation*}
\bigl|x^{(n+1)}_{(s)}(\hat w)-x^{(n)}_{(s)}(\hat w)\bigr| .
\end{equation*}
So the sequence $\bigl(x^{(n)}_{(s)}(\hat w)\bigr)_n$ is Cauchy and converges.

\emph{Part (b).} The reference inverse coupling at $g_s$ is $g_s^{-2}=X_s$. Shifting it by $s'$ steps of
the scaffold gives
\begin{equation*}
X_s+c_0s'=X+c_0(s+s')=X_{s+s'} .
\end{equation*}
Hence, by the definition of the label difference at reference coupling $g_s$ and length $n-s\ge s'$,
\begin{equation*}
\delta_{n-s,s'}(g_s)=a_{n-s-s'}(X_{s+s'})-a_{n-s}(X_s).
\end{equation*}
Adding $\delta_{n,s}(g)=a_{n-s}(X_s)-a_n(X)$, the term $a_{n-s}(X_s)$ cancels, and
\begin{equation*}
\delta_{n,s}(g)+\delta_{n-s,s'}(g_s)=a_{n-s-s'}(X_{s+s'})-a_n(X)=\delta_{n,s+s'}(g).
\end{equation*}
This is \eqref{9.13:eq-intermediate-couplings-2}.

Read in terms of the re-rooting lemma (Lemma~\ref{9.13:lem-re-rooting}), the identity says the
following. The identity \eqref{9.13:eq-re-rooting-c} may be applied at $(g,s)$ and then at $(g_s,s')$,
or once at $(g,s+s')$. Both routes yield the same result, because both name the one bare number
$a_n(X)$.

Since $c_0>0$ we have $X_s\ge X$, hence $g_s\le g\le\gamma_U(\varrho)$. So the limit $\delta_{s'}(g_s)$
is taken in the same setting as $\delta_s(g)$. As $n\to\infty$, $\delta_{n,s}(g)\to\delta_s(g)$ and,
since $n-s\to\infty$, $\delta_{n-s,s'}(g_s)\to\delta_{s'}(g_s)$; hence
\eqref{9.13:eq-intermediate-couplings-2} gives $\delta_{s+s'}(g)=\delta_s(g)+\delta_{s'}(g_s)$.

\emph{Part (c).} The identification \eqref{9.13:eq-limit-member-c} for the unit change by $s$ at $g$
relates the member at $(m_0,g)$ to a member at $(m_0+s,g_s)$. The test functions $\vec f$ are replaced
by $\vec f\circ L^{-s}\cdot$, and the offset is read through the label difference $\delta_s(g)$. Follow
it by the identification for $s'$ at $g_s$. The data compose as follows.
\begin{itemize}
\item The torus exponent becomes $m_0+s+s'$.
\item The reference coupling becomes $(g_s)_{s'}=(X_s+c_0s')^{-1/2}=g_{s+s'}$.
\item The test functions become
\begin{equation*}
(\vec f\circ L^{-s}\cdot)\circ L^{-s'}\cdot=\vec f\circ L^{-(s+s')}\cdot .
\end{equation*}
\item The label differences add to $\delta_s(g)+\delta_{s'}(g_s)=\delta_{s+s'}(g)$, by part~(b).
\end{itemize}
These are exactly the data of the identification for the unit change by $s+s'$ at $g$. Hence the
composite coincides with it whenever all three identifications lie inside the window
\eqref{9.13:eq-limit-member-a}. The two instances of \eqref{9.13:eq-limit-member-c} are taken at the
reference couplings $g$ and $g_s$. Both couplings are at most $\gamma_U(\varrho)$, as shown in the proof
of part~(b), so both instances are instances of clause~(v) of Theorem~\ref{3.1:thm-main}.
\end{proof}

\begin{definition}[Frozen scaffolds, their runs and the tuned family]\label{9.13:not-scaffold-runs}
We use the notation of Theorem~\ref{3.1:thm-main} and of the uniqueness theorem. We write $X=g^{-2}$
for the inverse square of the reference coupling. For integers $r\ge 0$ we put
\begin{equation*}
\theta_r:=X+c_0 r,
\end{equation*}
the straight reference line at depth $r$. We put $\varrho_2:=\varrho/2=r_J/2$, where $r_J>0$ is
the radius of the parameter disc of the uniqueness theorem; any $r_J>0$ is allowed. The constants
$\lambda_U>0$ and $c'$ are those of clause (v-a) of Theorem~\ref{3.1:thm-main}, and
$\lambda^{\sharp}$, $B^{\sharp}$, $c_5$ are those of clause (0).

\emph{First passage.} For a bare coupling $g_0=x^{-1/2}$, the run of clause (0) from $g_0$ (the
\emph{own-threshold run}) has coupling increments $b_j$ (Definition~\ref{1.2:def-couplings}) and
inverse couplings $x_j$, and its first-passage
level is
\begin{equation*}
K(g_0,g):=\min\{k\ge 0:\ x_k\le X+B^{\sharp}\}.
\end{equation*}
Clause (0) gives $x_K\in\,]X,X+B^{\sharp}]$ and the two-sided flow inequalities
\begin{equation*}
\lambda^{\sharp}(j-i)-c_5\le x_i-x_j\le B^{\sharp}(j-i),\qquad 0\le i\le j\le K .
\end{equation*}
These numbers do not depend on the torus: by the correlation theorem the non-wrapping terms of every
family in its expansion are independent of the torus, and hence, by the lemma of the correlation
theorem's chapter that passes torus-independence from the non-wrapping terms to the numbers of the
step, so is every
increment $b_j$, every running shift $\zeta_j$, every $x_j$, and $K$.

\emph{The frozen scaffold.} The frozen scaffold $\Theta_X$ is the assignment, fixed in the setting of
the uniqueness theorem, of every coupling-dependent choice of Ba{\l}aban's step at depth $r$ to its
value at the coupling $g_{(r)}:=\theta_r^{-1/2}$. These choices are therefore functions of the depth
alone, the same for all run lengths $n$.

\emph{Runs.} Fix $n\ge 0$ and a depth-$n$ shift $\zeta$. The $\Theta_X$-run of length $n$ with
depth-$n$ shift $\zeta$ has bare inverse coupling $\theta_n-\zeta$. It is given by the recursion
\begin{equation}\label{9.13:eq-scaffold-runs-1}
\begin{gathered}
\zeta_{(n)}:=\zeta,\qquad \zeta_{(r)}:=\zeta_{(r+1)}+b_{(r)}-c_0\quad(0\le r<n),\\
x_{(r)}:=\theta_r-\zeta_{(r)}\quad(0\le r\le n),
\end{gathered}
\end{equation}
so that $x_{(n)}=\theta_n-\zeta$ is the bare inverse coupling.
Here $b_{(r)}$ is the coupling increment of Definition~\ref{1.2:def-couplings} produced by the step
of the run at depth $r$, and the run's coupling-dependent choices at
depth $r$ are evaluated at $g_{(r)}$. For $0\le r<n$ and $0\le s\le n$,
\begin{equation}\label{9.13:eq-scaffold-runs-2}
x_{(r+1)}=x_{(r)}+b_{(r)},\qquad x_{(s)}=x_{(0)}+\sum_{r<s}b_{(r)} .
\end{equation}

\emph{The tuned family.} Let $n\ge 0$ and let $t$ be real with $|t|\le r_J$. By the uniqueness
theorem there is exactly one real $\Theta_X$-run of length $n$ that is alive (in the shift tube, in
the terminology of that theorem) and has $\zeta_{(0)}=t$. We write $\zeta_*^{n}(t)$ for this run and
for its depth-$n$ shift, and we put
\begin{equation*}
a_n=a_n(X):=\theta_n-\zeta_*^{n}(0).
\end{equation*}

\emph{Schwinger functions.} Fix a number $\mathsf{p}$ of test functions and a separation
$\mathsf{d}>0$. Let
$\vec f=(f_1,\dots,f_{\mathsf{p}})$ be Lipschitz functions on the continuum torus of side
$2L^{m_0}\ell_g$ whose supports are pairwise at distance at least $\mathsf{d}$; here $\ell_g$ is the
reference length at the coupling $g$. On the bare datum $(m_0+n,x)$ we define
\begin{equation*}
S^{n}_{m_0}(x;\vec f):=\partial_{z_1}\cdots\partial_{z_{\mathsf{p}}}\big|_{0}\log Z(z) .
\end{equation*}
The bare datum $(m_0+n,x)$ consists of the periodic lattice with $2L^{m_0+n}$ sites a side and the
weight $\exp[-A(x-w,U)]$, where $w=\sum_i z_i f_i$ and $A$ is the weighted Wilson action of
Definition~\ref{1.1:def-wilson-action}; $Z(z)$ is the integral of this weight over $U$ against
product Haar measure. The limits
\begin{equation*}
S^{(\hat w)}_{X,m_0}:=\lim_{n\to\infty}S^{n}_{m_0}\big(a_n(X)-\hat w;\,\cdot\,\big),
\end{equation*}
for every offset $\hat w$ for which the convergence statement of the uniqueness theorem supplies the
limit, are the members of the one-parameter family of the uniqueness theorem on $\mathbb{T}_{m_0}$;
they are the limits of its convergence statement.

\emph{The gap constant.} We put
\begin{equation*}
j_U:=1+\big\lceil(\varrho_2+c')/\lambda_U\big\rceil .
\end{equation*}
It depends on $\varrho_2$, $c'$ and $\lambda_U$ only.
\end{definition}

\begin{proof}
We prove \eqref{9.13:eq-scaffold-runs-2}. Let $0\le r<n$. The definition of $\theta_r$ gives
$\theta_{r+1}=\theta_r+c_0$. Using this and the recursion \eqref{9.13:eq-scaffold-runs-1} for
$\zeta_{(r)}$,
\begin{align*}
x_{(r)}&=\theta_r-\zeta_{(r)}=\theta_r-\zeta_{(r+1)}-b_{(r)}+c_0\\
&=\theta_{r+1}-\zeta_{(r+1)}-b_{(r)}=x_{(r+1)}-b_{(r)} .
\end{align*}
This is the first identity. Summing it over $r=0,\dots,s-1$ telescopes to
$x_{(s)}-x_{(0)}=\sum_{r<s}b_{(r)}$, which is the second identity.
\end{proof}

\begin{lemma}[The change-of-units identity]\label{9.13:lem-change-of-units-bis}
For integers $m_0\ge 0$ and $n\ge 0$, a bare inverse coupling $x$, and Lipschitz functions
$\vec f=(f_1,\dots,f_{\mathsf{p}})$ on the continuum torus of coordinate side $2L^{m_0}$, whose
supports
are pairwise at distance at least $\mathsf{d}>0$, let $S^{n}_{m_0}(x;\vec f)$ denote the derivative
$\partial_{z_1}\cdots\partial_{z_{\mathsf{p}}}\log Z$ at $z=0$, where $Z$ is the partition function
of the periodic lattice with $2L^{m_0+n}$ sites a side at bare inverse coupling $x$, with Boltzmann
weight $\exp[-A(x-w,U)]$, $A(\cdot,U)$ being the weighted Wilson action of
Definition~\ref{1.1:def-wilson-action}, and with plaquette weights
\begin{equation*}
w(p)=\sum_{i=1}^{\mathsf{p}} z_i f_i(x(p)),
\end{equation*}
$x(p)$ being the barycentre of $p$ in units of the reference length $\ell_g$ (the letter $x$ alone
always denotes the bare inverse coupling, $x(p)$ the barycentre). The number
$S^{n}_{m_0}(x;\vec f)$ depends only on $m_0$, $n$, $x$ and on $\vec f$ as a function of the
coordinate, not on the name of the length unit. For
$0\le s\le n$ write $(\vec f\circ L^{-s}\cdot)_i(y):=f_i(L^{-s}y)$, a function on the torus of
coordinate side $2L^{m_0+s}$.
Then for all $0\le s\le n$, all $m_0\ge 0$, all $x$ and all $\vec f$ as above,
\begin{equation}\label{9.13:eq-change-of-units-a-bis}
S^{n}_{m_0}(x;\vec f)=S^{n-s}_{m_0+s}\bigl(x;\vec f\circ L^{-s}\cdot\bigr).
\end{equation}
Let $K\ge 0$ and $m\ge 0$ be integers and $x_0$ a bare inverse coupling, and write
$S^T_{\mathsf{p};K,m}[x_0]$ for the connected Schwinger function of
Theorem~\ref{3.1:thm-main} with $\mathsf{p}$ test functions on $\mathbb{T}_m$ at bare spacing
$L^{-K}$ and bare inverse coupling $x_0$; its lattice has $2L^{m+K}$ sites a side, so that
$S^T_{\mathsf{p};K,m}[x_0]=S^{K}_{m}(x_0;\cdot)$. Then for all $0\le s\le K$,
\begin{equation}\label{9.13:eq-change-of-units-1}
S^T_{\mathsf{p};K,m}[x_0](\vec f)=S^T_{\mathsf{p};K-s,m+s}[x_0]\bigl(\vec f\circ L^{-s}\cdot\bigr).
\end{equation}
The first-passage-frame identity is the case $s=K(g_0,\gamma_\star)-K(g_0,g)$ of
\eqref{9.13:eq-change-of-units-1}, read toward the infrared; here $g_0$ is the bare coupling with
$g_0^{-2}=x_0$, $K(\cdot,\cdot)$ is the first-passage depth of
Definition~\ref{1.2:def-flow-constants}, and $\gamma_\star$ is the reference coupling of that
identity.
The map $\vec f\mapsto\vec f\circ L^{-s}\cdot$ preserves the class of real Lipschitz tuples with
pairwise separated supports: sup norms are unchanged, the separation $\mathsf{d}$ becomes
$L^s\mathsf{d}$, and each Lipschitz constant is multiplied by $L^{-s}$.
\end{lemma}

\begin{proof}
The lattice with $2L^{m_0+n}$ sites a side has spacing $L^{-n}\ell_g$ and torus side
$2L^{m_0}\ell_g$. Measure lengths instead in the unit $\ell':=L^{-s}\ell_g$. Since
$L^{-n}\ell_g=L^{-(n-s)}\ell'$ and $2L^{m_0}\ell_g=2L^{m_0+s}\ell'$, the same lattice has spacing
$L^{-(n-s)}\ell'$ and side $2L^{m_0+s}\ell'$. It therefore has
$2L^{(m_0+s)+(n-s)}=2L^{m_0+n}$ sites a side, which is the lattice on which
$S^{n-s}_{m_0+s}$ is defined. Sites, bonds, plaquettes, bond variables and the Haar measure are
unchanged, and so is the bare inverse coupling $x$. Only the coordinate of a barycentre changes:
in units of $\ell'$ it is $y(p)=L^sx(p)$. Hence the plaquette weights are
\begin{equation*}
w(p)=\sum_i z_if_i(x(p))=\sum_i z_if_i(L^{-s}y(p))=\sum_i z_i(f_i\circ L^{-s})(y(p)),
\end{equation*}
which are the weights of the lattice with $2L^{(m_0+s)+(n-s)}$ sites a side, at inverse
coupling $x$, built from the test functions $\vec f\circ L^{-s}\cdot$. The two descriptions
give the same integrand, plaquette by plaquette, for every $U$ and every complex
$z=(z_1,\dots,z_{\mathsf{p}})$. Hence $Z(z)$ is the same function of $z$ in both descriptions,
and every $\partial_z$-derivative of $\log Z$ at $z=0$ is the same number. Taking
$\partial_{z_1}\cdots\partial_{z_{\mathsf{p}}}$ at $z=0$ gives \eqref{9.13:eq-change-of-units-a-bis}.

By the identification of $S^T_{\mathsf{p};K,m}[x_0]$ with $S^{K}_{m}(x_0;\cdot)$ made in the
statement, applying \eqref{9.13:eq-change-of-units-a-bis} with $n=K$, $m_0=m$ and
$x=x_0$ gives \eqref{9.13:eq-change-of-units-1}.

Now let $\vec f$ be real, Lipschitz and pairwise separated. Each $f_i\circ L^{-s}$ is real and
takes the same values as $f_i$, so its sup norm equals that of $f_i$. Its support is
$L^s\operatorname{supp}f_i$. The dilation $y\mapsto L^{-s}y$ maps the torus of coordinate side
$2L^{m_0+s}$ onto that of coordinate side $2L^{m_0}$ and multiplies torus distances by $L^{-s}$.
Hence supports at mutual distance at least $\mathsf{d}$ become supports at mutual distance at least
$L^s\mathsf{d}$. For the Lipschitz constants,
\begin{equation*}
\bigl|f_i(L^{-s}y)-f_i(L^{-s}y')\bigr|\le\operatorname{Lip}f_i\cdot L^{-s}|y-y'|,
\end{equation*}
so that $\operatorname{Lip}(f_i\circ L^{-s})\le L^{-s}\operatorname{Lip}f_i$. Since the dilation
is a bijection multiplying torus distances by exactly $L^{-s}$, the same estimate applied to
$f_i=(f_i\circ L^{-s})\circ L^{s}$ gives the reverse inequality.
Thus $\operatorname{Lip}(f_i\circ L^{-s})=L^{-s}\operatorname{Lip}f_i$.

Neither side of \eqref{9.13:eq-change-of-units-a-bis} involves the frozen scaffolds, their runs,
the tuned family of Notation~\ref{9.13:not-scaffold-runs}, or any reference coupling: both sides
are determined by $Z$ alone, which, like the connected Schwinger functions built from it in the
uniqueness theorem, does not depend on the scaffold, and $\ell_g$ serves only as the name of the
unit in which positions are read.
\end{proof}

\begin{lemma}[Convergence of the frozen coupling at fixed depth]\label{9.13:lem-frozen-coupling}
Let $g\le\gamma_U$, where $\gamma_U$ is the threshold of the uniqueness theorem. Put $X=g^{-2}$ and
$\theta_r=X+c_0r$, and use the frozen scaffold $\Theta_X$, its runs and the tuned family of
Notation~\ref{9.13:not-scaffold-runs}. Let $n_0$ be such that $\Upsilon_l\le\frac14$ for all
$l\ge n_0$. Fix $s\ge0$. For $\hat w\in\overline{D}(0,r_J/4)$ (for real $\hat w$ this is
$|\hat w|\le\varrho_2/2$) and $n\ge\max\{s,n_0\}$ put
\begin{equation*}
x^{(n)}_{(s)}(\hat w):=\theta_s-\zeta^{n}_{(s)}\bigl(\zeta_*^{n}(0)+\hat w\bigr),
\end{equation*}
the inverse coupling at depth $s$ of the $\Theta_X$-run of length $n$ from the bare datum
$a_n(X)-\hat w$. Then the sequence $(x^{(n)}_{(s)}(\hat w))_n$ is Cauchy, uniformly in $\hat w$.
Its limit $\hat{x}_s(\hat w):=\lim_n x^{(n)}_{(s)}(\hat w)$ exists uniformly on
$\overline{D}(0,r_J/4)$, is holomorphic in $\hat w$ on $D(0,r_J/4)$ and real for real $\hat w$.
For all $n'>n\ge\max\{s,n_0\}$ and all $\hat w\in\overline{D}(0,r_J/4)$, with the constants
$\mathsf{A}$, $C_E$, the contraction letter $\rho$ and the summable sequences $(\Pi_i)$,
$(\Upsilon_i)$ of the uniqueness theorem,
\begin{equation}\label{9.13:eq-frozen-coupling-a}
\begin{gathered}
\bigl|x^{(n')}_{(s)}(\hat w)-x^{(n)}_{(s)}(\hat w)\bigr|
\le\mathsf{A}C_E\Bigl[\frac{\rho^{\,n+1-s}}{(1-\rho)^2}+s\sum_{i\ge n}\Pi_i\Bigr]
+\frac53\,|\hat w|\sum_{i\ge n}\Upsilon_i,\\
\hat{x}_s(0)-X\in\bigl[\lambda_Us-2c_2,\;3\lambda_Us+2c_2\bigr],
\end{gathered}
\end{equation}
and the right-hand side of the first line tends to $0$ as $n\to\infty$. Finally,
$X-\hat{x}_0(\hat w)=\lim_n h_n(\hat w)$, the limit of the depth-zero label maps.
\end{lemma}

\begin{proof}
Fix $n\ge\max\{s,n_0\}$ and put $\mathfrak{P}=D(0,r_J)$. For $v\in\mathfrak{P}$ consider the
frozen pair formed by the $\Theta_X$-run of length $n$ from the bare shift $\zeta_*^{n}(0)+v$ and the
$\Theta_X$-run of length $n+1$ from the bare shift $\zeta_*^{n+1}(0)+\mathfrak{n}_n(v)$, where
$\mathfrak{n}_n$ is the matching map of the frozen pair. The matching proposition for frozen pairs
(obtained from a single Rouch\'e argument, using only the bound of the uniqueness theorem on the
Lipschitz constants of the increments along $\Theta_X$ and its bound on the tube) states that for
every $v\in\mathfrak{P}$
\begin{equation}\label{9.13:eq-frozen-coupling-1}
h_{n+1}\bigl(\mathfrak{n}_n(v)\bigr)=h_n(v),\qquad |\mathfrak{n}_n(v)|\le2|v|,\qquad
|\mathfrak{n}_n(v)-v|\le\Upsilon_n|v|.
\end{equation}
Here $h_n$ is the depth-zero label map of the length-$n$ run
(Notation~\ref{9.13:not-scaffold-runs}), $h_n(v)=\zeta^{n}_{(0)}(\zeta_*^{n}(0)+v)$, the depth-zero
shift of the length-$n$ run from $\zeta_*^{n}(0)+v$, so the two runs of the pair are matched at
depth $0$ for every $v$. For
$0\le u\le n$ put
\begin{equation*}
d_u:=\sup_{v\in\overline{\mathfrak{P}}}\bigl|\zeta^{n+1}_{(u)}\bigl(\zeta_*^{n+1}(0)+\mathfrak{n}_n(v)\bigr)
-\zeta^{n}_{(u)}\bigl(\zeta_*^{n}(0)+v\bigr)\bigr|.
\end{equation*}
By the first identity in \eqref{9.13:eq-frozen-coupling-1}, $d_0=0$.

Both the pin $x_{(0)}=X$ of the tuned run and the matching of the pair sit at depth $0$. Part~(i)
of the pinning lemma therefore allows the depth-wise comparison bound for frozen pairs to be applied
with starting depth $0$. That bound holds under the hypotheses of the uniqueness theorem and states
that for every $r<n$
\begin{equation}\label{9.13:eq-frozen-coupling-2}
d_{r+1}\le\mathsf{A}C_E\Bigl[\frac{\rho^{\,n-r}}{1-\rho}+(r+1)\Pi_n\Bigr].
\end{equation}
Let $s\ge1$ and take $r=s-1$; then $r<n$ because $n\ge s$. Evaluating the supremum defining $d_s$
at $v=\hat w$, which is allowed because $\overline{D}(0,r_J/4)\subset\overline{\mathfrak{P}}$, we
obtain from \eqref{9.13:eq-frozen-coupling-2}
\begin{equation}\label{9.13:eq-frozen-coupling-3}
\bigl|\zeta^{n+1}_{(s)}\bigl(\zeta_*^{n+1}(0)+\mathfrak{n}_n(\hat w)\bigr)
-\zeta^{n}_{(s)}\bigl(\zeta_*^{n}(0)+\hat w\bigr)\bigr|
\le\mathsf{A}C_E\Bigl[\frac{\rho^{\,n+1-s}}{1-\rho}+s\Pi_n\Bigr].
\end{equation}
For $s=0$ the left-hand side of \eqref{9.13:eq-frozen-coupling-3} is at most $d_0=0$. Since the
right-hand side is non-negative, \eqref{9.13:eq-frozen-coupling-3} holds for $s=0$ as well; in that
case only the Lipschitz term below remains.

Next we compare the runs of length $n+1$ from the two bare shifts $\zeta_*^{n+1}(0)+\hat w$ and
$\zeta_*^{n+1}(0)+\mathfrak{n}_n(\hat w)$. For the run of length $n+1$ from a bare shift $\zeta$
we write $b_{(u)}(\zeta)$ for its increment at depth $u$. By Notation~\ref{9.13:not-scaffold-runs},
$\zeta_{(n+1)}=\zeta$ and $\zeta_{(u)}=\zeta_{(u+1)}+b_{(u)}-c_0$. Unwinding this recursion from
depth $n+1$ down to depth $s$ gives
\begin{equation*}
\zeta^{n+1}_{(s)}(\zeta)=\zeta+\sum_{s\le u\le n}\bigl(b_{(u)}(\zeta)-c_0\bigr).
\end{equation*}
On the tube, the Lipschitz constant of $\zeta\mapsto\zeta^{n+1}_{(s)}(\zeta)$ is therefore at most
$1+\sum_u\operatorname{Lip}b_{(u)}$. The running-shift statement of the correlation theorem gives
$\sum_{k<K}\operatorname{Lip}_{\overline{D}(0,\varrho)}b_{k+1}\le\frac23$, hence the bound $\frac53$.
Along $\Theta_X$ the uniqueness theorem bounds the sum of the Lipschitz constants of the increments
by $\frac13$, which gives the better bound $\frac43\le\frac53$.

These bounds apply at the two points in question. We have $|\hat w|\le r_J/4$ and, by
\eqref{9.13:eq-frozen-coupling-1}, $|\mathfrak{n}_n(\hat w)|\le2|\hat w|\le r_J/2$, so both
$\hat w$ and $\mathfrak{n}_n(\hat w)$ lie in $\overline{D}(0,r_J/2)$. The uniqueness theorem provides
for the tuned run, at every depth $s$, the margin $|\zeta_{(s)}|+8r_J\le\varrho_s$, where
$\varrho_s$ denotes the radius of the shift tube at depth $s$. Hence the segment joining the two
bare shifts lies inside the region
where the Lipschitz bound holds, and by the last inequality of \eqref{9.13:eq-frozen-coupling-1}
\begin{equation}\label{9.13:eq-frozen-coupling-4}
\begin{split}
&\bigl|\zeta^{n+1}_{(s)}\bigl(\zeta_*^{n+1}(0)+\hat w\bigr)
-\zeta^{n+1}_{(s)}\bigl(\zeta_*^{n+1}(0)+\mathfrak{n}_n(\hat w)\bigr)\bigr|\\
&\qquad\le\frac53\,\bigl|\hat w-\mathfrak{n}_n(\hat w)\bigr|\le\frac53\,\Upsilon_n|\hat w|.
\end{split}
\end{equation}

In the difference $x^{(n+1)}_{(s)}(\hat w)-x^{(n)}_{(s)}(\hat w)$ the term $\theta_s$ cancels, so
\begin{equation*}
x^{(n+1)}_{(s)}(\hat w)-x^{(n)}_{(s)}(\hat w)
=\zeta^{n}_{(s)}\bigl(\zeta_*^{n}(0)+\hat w\bigr)-\zeta^{n+1}_{(s)}\bigl(\zeta_*^{n+1}(0)+\hat w\bigr).
\end{equation*}
We insert $\zeta^{n+1}_{(s)}(\zeta_*^{n+1}(0)+\mathfrak{n}_n(\hat w))$ and apply the triangle
inequality together with \eqref{9.13:eq-frozen-coupling-3} and \eqref{9.13:eq-frozen-coupling-4}.
This gives
\begin{equation}\label{9.13:eq-frozen-coupling-5}
\bigl|x^{(n+1)}_{(s)}(\hat w)-x^{(n)}_{(s)}(\hat w)\bigr|
\le\mathsf{A}C_E\Bigl[\frac{\rho^{\,n+1-s}}{1-\rho}+s\Pi_n\Bigr]+\frac53\,\Upsilon_n|\hat w|.
\end{equation}

Let $n'>n$. Every $i$ with $n\le i<n'$ satisfies $i\ge\max\{s,n_0\}$, so
\eqref{9.13:eq-frozen-coupling-5} holds with $n$ replaced by $i$. Summing these inequalities over
$n\le i<n'$ and using the triangle inequality bounds
$|x^{(n')}_{(s)}(\hat w)-x^{(n)}_{(s)}(\hat w)|$ by the sum of the right-hand sides. We bound the
three resulting sums in turn.

The geometric series gives
\begin{equation*}
\sum_{i\ge n}\frac{\rho^{\,i+1-s}}{1-\rho}=\frac{\rho^{\,n+1-s}}{(1-\rho)^2}.
\end{equation*}
By the summability statement of the uniqueness theorem,
$T_{\Pi}:=\sup_l l\sum_{i\ge l}\Pi_i<\infty$; here
$\Pi_n=\theta_n^{1-\kappa_0/2}(\log\theta_n)^{c_3}$, with the exponents $\kappa_0\ge10$ and $c_3$
of that theorem. Consequently
$s\sum_{i\ge n}\Pi_i\le sT_{\Pi}/n$, which tends to $0$ as $n\to\infty$. By the same theorem,
$\Upsilon_n\le C\,n^{-3/2}(\log n)^{p_0+1}$ with a constant $C$ independent of $n$, so
$\sum_i\Upsilon_i<\infty$ and its tails tend to $0$.

This proves the first line of \eqref{9.13:eq-frozen-coupling-a}, and shows that its right-hand side
tends to $0$ as $n\to\infty$, uniformly in $\hat w$ because $|\hat w|\le r_J/4$. The sequence is
therefore uniformly Cauchy on $\overline{D}(0,r_J/4)$, so $\hat{x}_s$ exists as a uniform limit.
Holomorphy on $D(0,r_J/4)$ and reality for real $\hat w$ carry over from the functions
$x^{(n)}_{(s)}$ to their uniform limit. For $s=0$ we have $\theta_0=X$, so
$x^{(n)}_{(0)}(\hat w)=X-h_n(\hat w)$; letting $n\to\infty$ gives
$X-\hat{x}_0(\hat w)=\lim_n h_n(\hat w)$. In the bi-Lipschitz lemma this limit is the limit label;
this identification is not used below.

It remains to prove the window for $\hat{x}_s(0)$. The run from $\zeta_*^{n}(0)$ has depth-zero
shift $0$, so $x_{(0)}=\theta_0=X$. Since $x_{(r+1)}=x_{(r)}+b_{(r)}$, it follows that
$x^{(n)}_{(s)}(0)-X=\sum_{r<s}b^{(n)}_{(r)}$, where $b^{(n)}_{(r)}$ are the increments of this run.
These increments satisfy $|b_{(u)}-\beta^0_u|\le\lambda_U$, the bound used in the proof of the
uniqueness theorem. By the statement on the two-loop reference increments, these satisfy
\begin{equation*}
\sum_{r<s}\beta^0_r=c_0s+s_{-1}-s_{s-1},\qquad c_0=2\lambda_U.
\end{equation*}
Hence $x^{(n)}_{(s)}(0)-X$ lies between $\lambda_Us+s_{-1}-s_{s-1}$ and
$3\lambda_Us+s_{-1}-s_{s-1}$. By the bound $|s_l|\le c_2$ on the constants of the two-loop reference
increments, from the same statement, this interval is contained in
$[\lambda_Us-2c_2,3\lambda_Us+2c_2]$. The interval is
closed and independent of $n$, so letting $n\to\infty$ gives the second line of
\eqref{9.13:eq-frozen-coupling-a}.
\end{proof}

The argument above is the estimate by which the proof of the uniqueness theorem bounds the
differences $a_n-a_l$. There, the depth-$l$ distances of the consecutive frozen pairs of lengths
$i,i+1$, summed over $i\ge l$, are bounded by
$\mathsf{A}C_E[\rho/(1-\rho)^2+l\sum_{i\ge l}\Pi_i]$, with starting depth $0$; here the same
estimate is read at the fixed depth $s$ instead. Lemma~\ref{9.13:lem-frozen-coupling} gives, for
every fixed $s$ and for runs along the frozen scaffold $\Theta_X$, the convergence of the inverse
couplings at depth $s$ along the tuned family; this is, for fixed $s$ and for the runs along
$\Theta_X$ only, the convergence of the intermediate couplings $g_{n-s}$ along the tuned family,
which clause (v) of Theorem~\ref{3.1:thm-main} does not assert in general. It says nothing about
depths $s$ that grow with $n$.
Nor does it say anything about the own-threshold inverse couplings $x^{\mathrm{own}}_{n,K}$ beyond
the bound $|x^{\mathrm{own}}-x^{\Theta}|\le\varpi(g)$ of the scaffold comparison, where that
comparison applies; here $x^{\Theta}$ is the corresponding inverse coupling along $\Theta_X$.

\begin{remark}[Tuned bare data on distinct frozen scaffolds]\label{9.13:rem-tuned-data-distinct}
Let $g\le\gamma_U$, let $X=g^{-2}$, let $\theta_r(X)=X+c_0r$ be the straight reference line and let
$\Theta_X$ be
the frozen scaffold along it, with its runs and the tuned data $a_n(X)=\theta_n(X)-\zeta_*^{n}(0)$
as in Notation~\ref{9.13:not-scaffold-runs}. Let $s\ge 1$, $n'\ge 0$, and let $X'=g'^{-2}$ with
$g'\le\gamma_U$. For $n\ge s$
write $x^{(n)}_{(s)}$ for the depth-$s$ inverse coupling of the $\Theta_X$-run of length $n$ from
$a_n(X)$, $b^{(n)}_{(r)}$ for its increment at depth $r$, and $\zeta^{(n)}_{(r)}$ for its shift at
depth $r$; in the notation of
Lemma~\ref{9.13:lem-frozen-coupling} this coupling is $x^{(n)}_{(s)}(0)$.
\begin{enumerate}
\item[(a)] Read from depth $s$ in the sense of Lemma~\ref{9.13:lem-change-of-units-bis}, the
$\Theta_X$-run of length $n'+s$ from $a_{n'+s}(X)$ is a $\Theta_{X+c_0s}$-run of length $n'$ from
the same bare datum; it is alive in the tube of $\Theta_{X+c_0s}$, and its depth-$0$
inverse coupling is
\begin{equation*}
x^{(n'+s)}_{(s)}=X+\sum_{r<s}b^{(n'+s)}_{(r)} .
\end{equation*}
\item[(b)] If $X'\neq X+c_0s$, then $a_{n'+s}(X)$, read from depth $s$ as in (a), is the bare
datum of a run on $\Theta_{X+c_0s}$, whereas $a_{n'}(X')$ is the tuned datum on $\Theta_{X'}$, and
$\Theta_{X+c_0s}\neq\Theta_{X'}$. If $X'=X+c_0s$, then $a_{n'+s}(X)=a_{n'}(X')$
holds if and only if $x^{(n'+s)}_{(s)}=X+c_0s$.
\item[(c)] The increments of the run give only the window
\begin{equation}\label{9.13:eq-tuned-data-distinct-1}
\lambda_Us-2c_2\;\le\; x^{(n'+s)}_{(s)}-X\;\le\;3\lambda_Us+2c_2 .
\end{equation}
\end{enumerate}
In particular no identity $a_{n'+s}(X)=a_{n'}(X')$ between tuned data at two references holds by
construction.
\end{remark}

\begin{proof}
\emph{What $a_n(X)$ is.} By the uniqueness theorem, for every $n\ge 0$ and every real $t$ with
$|t|\le r_J$ exactly one real $\Theta_X$-run of length $n$ alive in the tube has depth-$0$ shift
$\zeta_{(0)}=t$; this run is $\zeta_*^{n}(t)$, and $a_n(X)=\theta_n(X)-\zeta_*^{n}(0)$. Its
depth-$0$ inverse coupling is $x_{(0)}=\theta_0(X)-0=X$. Thus $a_n(X)$ is fixed by two data: the
scaffold $\Theta_X$ --- its coupling-dependent choices at depth $r$ are evaluated at
$g_{(r)}=\theta_r(X)^{-1/2}$, so depend on the depth only and not on $n$ (in
particular the action of \cite[(1.100)]{B-CMP122a} entering the scaffold's choices at depth $r$
is the pure Wilson action with coefficient $\theta_r(X)$) --- and the normalisation $x_{(0)}=X$
at depth $0$ of that scaffold. By the uniqueness theorem, moreover, $a_n(X)$ is real and depends
on $(L,\mathfrak{p},g,n)$ only; it does not depend on the torus.

\emph{Proof of (a).} By the change-of-units identity
(Lemma~\ref{9.13:lem-change-of-units-bis}, \eqref{9.13:eq-change-of-units-a-bis}), for any torus exponent
$m_0$ the lattice with $2L^{m_0+n'+s}$ sites a side at reference $X$, read from depth $s$, is the
lattice with $2L^{(m_0+s)+n'}$ sites a side at the same bare inverse coupling. The choices seen at
the new depth $r'$ are those of
$\Theta_X$ at depth $r'+s$, evaluated at
\begin{equation*}
\theta_{r'+s}(X)=(X+c_0s)+c_0r'=\theta_{r'}(X+c_0s),
\end{equation*}
so the scaffold seen is exactly $\Theta_{X+c_0s}$. The bare datum is unchanged, and the shifts
and increments at new depth $r'$ are those of the original run at depth $r'+s$; the run is
alive because the tube of $\Theta_{X+c_0s}$ at depth $r'$ is the tube of $\Theta_X$ at depth
$r'+s$, which by the uniqueness theorem contains the original run with room equal to twice the
radius that the uniqueness theorem attaches to its tube. Since
$x_{(r+1)}=x_{(r)}+b_{(r)}$, its depth-$0$ coupling is
$x^{(n'+s)}_{(s)}=X+\sum_{r<s}b^{(n'+s)}_{(r)}$.

\emph{Proof of (b).} If $X'\neq X+c_0s$, the lines $\theta_r(X')$ and $\theta_r(X+c_0s)$ differ
by the non-zero constant $X'-X-c_0s$ at every depth, so the coefficients of the Wilson action of
\cite[(1.100)]{B-CMP122a} entering the choices at depth $r$, hence the scaffolds, differ. The
scaffold comparison bounds by $\varpi(g)$, at every level, the difference
between the couplings carried by one bare datum on the own-threshold run of clause~(0) of
Theorem~\ref{3.1:thm-main} and on a
frozen scaffold; applied on $\Theta_{X+c_0s}$ and on $\Theta_{X'}$ and combined by the triangle
inequality, it bounds by $2\varpi(g)$ the difference between the couplings carried by one bare
datum on the two frozen scaffolds, which is a bound, not an identity. Now let
$X'=X+c_0s$. If $a_{n'+s}(X)=a_{n'}(X')$, the run of (a) is the $\Theta_{X'}$-run of length $n'$
from $a_{n'}(X')$, whose depth-$0$ coupling is $X'$; hence $x^{(n'+s)}_{(s)}=X+c_0s$.
Conversely, if $x^{(n'+s)}_{(s)}=X+c_0s$, the run of (a) is a real alive
$\Theta_{X'}$-run of length $n'$ with depth-$0$ shift $0$, so by the uniqueness statement above
it is $\zeta_*^{n'}(0)$ on $\Theta_{X'}$, and its bare datum is $a_{n'}(X')$.

\emph{Proof of (c).} By the uniqueness theorem the increments satisfy
$|b_{(r)}-\beta^0_r|\le\lambda_U$, where, by the two-loop form of the reference increments,
$\beta^0_l=c_0+s_{l-1}-s_l$ with $|s_l|\le c_2$ for every term of the sequence $(s_l)_l$. The sum
$\sum_{r<s}\beta^0_r$ telescopes: it equals $c_0s$ plus the difference between the terms of the
sequence $(s_l)_l$ at the two ends of the range of summation, each of modulus at most $c_2$, so
it lies in
$[c_0s-2c_2,c_0s+2c_2]$; adding $\sum_{r<s}(b_{(r)}-\beta^0_r)\in[-\lambda_Us,\lambda_Us]$ and
using $c_0=2\lambda_U$ gives \eqref{9.13:eq-tuned-data-distinct-1}. No per-step monotonicity is
asserted in clause~(v-a) of Theorem~\ref{3.1:thm-main}. Moreover the depth-$s$ coupling
$x^{(n)}_{(s)}$ depends on $n$, and by part~(ii) of the lemma characterising the tuned data this
dependence is non-zero in general. Hence the equality $x^{(n'+s)}_{(s)}=X+c_0s$ for every $n'$,
which would say that the frozen run follows the reference line exactly for $s$ steps at every
cutoff, fails in general; it is asserted neither by the uniqueness theorem nor by
clause~(v-a).

Fix a torus exponent $m_0$. For neither of the two natural choices of $X'$ do the statements
above identify the limit along
the dilated tuned family $\bigl(a_{n'+s}(X)\bigr)_{n'}$, read on the lattices of torus exponent
$m_0+s+n'$, with a member $S^{(\hat{w}')}_{X',m_0+s}$ of the family of limits at reference $X'$.
With $X'=X+c_0s$, the run
of (a) is $\zeta_*^{n'}(t_{n'})$ on $\Theta_{X+c_0s}$ with
$t_{n'}=\zeta^{(n'+s)}_{(s)}=X+c_0s-x^{(n'+s)}_{(s)}$, and the uniqueness
theorem identifies limits only while $|t|\le\varrho_2/2=r_J/4$; but $|\zeta^{(n'+s)}_{(s)}|$ is
bounded only by the radius of the tube at depth $s$, which is of order
$\theta_s(X)^{1/2}(\log\theta_s(X))^{p_0}$, and, from the increments, by
\begin{equation*}
|\zeta^{(n'+s)}_{(s)}|\le 2c_2+\Bigl|\sum_{r<s}\bigl(b^{(n'+s)}_{(r)}-\beta^0_r\bigr)\Bigr| ;
\end{equation*}
the increment bound gives only
$\bigl|\sum_{r<s}(b^{(n'+s)}_{(r)}-\beta^0_r)\bigr|\le\lambda_Us$, which is not uniform in $s$;
no bound by $r_J/4$ uniform in $s$ is asserted. With
$X'=\hat{x}_s(0)$, the limit of the depth-$s$ frozen coupling of
Lemma~\ref{9.13:lem-frozen-coupling} (see \eqref{9.13:eq-frozen-coupling-a}), the scaffolds
$\Theta_{\hat{x}_s(0)}$ and $\Theta_{X+c_0s}$ are frozen lines
offset by $X+c_0s-\hat{x}_s(0)$ at every depth, and admissibility of the dilated data in the
tube of $\Theta_{\hat{x}_s(0)}$ is a comparison of two frozen scaffolds at one cutoff from one
bare datum, which neither the scaffold comparison nor the uniqueness theorem provides. Finally,
the first-passage-frame identity concerns one run read at two
first-passage references and no tuned family, so it does not yield the identification either.
Nor does the bi-Lipschitz lemma: its map between bare and reference couplings is that of one
scaffold and says nothing about couplings read on another. The limit statement of clause~(v) of
Theorem~\ref{3.1:thm-main}, in which the dilation exponent is determined in the course of the
proof rather than prescribed in advance, is the content of the proposition proving the dilation
identity of clause~(v) with offset labels, whose proof does not pass through
an identity between tuned data; an identification at a prescribed depth $s$ is not asserted.
\end{proof}

\begin{proposition}[The limit-member identification]\label{9.13:prop-limit-member}
Let $N=2$ and assume the hypotheses \textup{(H1)}--\textup{(H7)} of
Notation~\ref{2.2:not-hypotheses}. Let $\varrho>0$ satisfy the floor
\begin{equation*}
\varrho\ge 6\,(B^{\sharp}+1)
\end{equation*}
displayed in clause \textup{(v-c)(2)} of Theorem~\ref{3.1:thm-main}, and put $\varrho_2:=\varrho/2$. Let
$\lambda_U>0$ and $c'$ be the constants of clause \textup{(v-a)} of Theorem~0, let
$\lambda^{\sharp}$, $B^{\sharp}>0$, $c_5=2c_2$ be the flow constants of clause \textup{(0)}
(Definition~\ref{1.2:def-flow-constants}), and put
\begin{equation*}
j_U:=1+\Bigl\lceil\frac{\varrho_2+c'}{\lambda_U}\Bigr\rceil .
\end{equation*}
Let $\gamma_{U'}\le\gamma_U$ be the threshold of the corollary of the uniqueness theorem.
Let $g\le\gamma_{U'}$ and $X:=g^{-2}$; let $X'$ be a real number with
$g':=X'^{-1/2}\le\gamma_{U'}$ (when $X'\ge X$, that is $g'\le g$, this is automatic);
put $u:=X'-X$. For $\tilde g\in\{g,g'\}$ write $\tilde X:=\tilde g^{-2}$, so that
$g_-(\tilde g)=(\tilde X+B^{\sharp})^{-1/2}$, and let $a_n(\tilde X)$, $n\ge0$, be the tuned
bare data at reference $\tilde g$ of Notation~\ref{9.13:not-scaffold-runs}. Let
$S^{n}_{m}(x;\vec f)$ be the connected Schwinger function of
Lemma~\ref{9.13:lem-change-of-units-bis} (the periodic lattice with $2L^{m+n}$ sites a side at
bare inverse coupling $x$, with test functions on the continuum torus of side $2L^{m}$ in the
reference length), and for an integer $r$ write $\vec f\circ L^{r}\cdot$ for the tuple of
functions $y\mapsto f_i(L^{r}y)$. Assume (a) and (c) below, and assume (b), (d), (e) at each
reference $\tilde g\in\{g,g'\}$.
\begin{enumerate}
\item[(a)] The change-of-units identity, Lemma~\ref{9.13:lem-change-of-units-bis}: the identity
\eqref{9.13:eq-change-of-units-a-bis} holds for all $0\le s\le n$, $m_0\ge0$, $x$ and $\vec f$,
and $\vec f\mapsto\vec f\circ L^{-s}\cdot$ preserves real Lipschitz tuples with pairwise
separated supports, the separation $\mathsf{d}$ becoming $L^{s}\mathsf{d}$.
\item[(b)] Clause \textup{(0)} of Theorem~0 at reference $\tilde g$: for every bare coupling
$g_0\le g_-(\tilde g)$, admissible by \textup{(H6)}, the own-threshold run of clause
\textup{(0)} from $g_0$, whose inverse coupling at level $k\ge0$ we write
$x^{\mathrm{own}}_k$, is determined by the bare data and not by the reference; its
first-passage depth
\begin{equation*}
K(g_0,\tilde g):=\min\{k\ge0:\ x^{\mathrm{own}}_k\le\tilde X+B^{\sharp}\}
\end{equation*}
is finite and satisfies
\begin{equation*}
x^{\mathrm{own}}_{K(g_0,\tilde g)}\in\,]\tilde X,\tilde X+B^{\sharp}],
\end{equation*}
and for $0\le i\le j\le K(g_0,\tilde g)$
\begin{equation*}
\lambda^{\sharp}(j-i)-c_5\le x^{\mathrm{own}}_i-x^{\mathrm{own}}_j\le B^{\sharp}(j-i).
\end{equation*}
\item[(c)] Torus-freeness, from the correlation theorem: the numbers of the own-threshold run
from $g_0$ (its inverse couplings and its first-passage depths) are functions of
$(g_0,L,\mathfrak{p})$ and of the reference alone; they depend neither on the torus exponent
nor on any test function.
\item[(d)] The uniqueness theorem at reference $\tilde g$:
\begin{enumerate}
\item[(d1)] for all $n,l\ge0$, $a_n(\tilde X)\ge\tilde X+\lambda_U n-2c_2$ and
$\lambda_U|n-l|-c'\le|a_n(\tilde X)-a_l(\tilde X)|$;
\item[(d2)] for every torus exponent $m\ge0$ admissible at reference $\tilde g$ in the sense
of \textup{(H5)}, every $\mathsf{p}\ge2$ and every $\mathsf{p}$-tuple $\vec f$ of real
Lipschitz functions on the continuum torus of side $2L^{m}\ell_{\tilde g}$ with supports
pairwise at distance at least $\mathsf{d}>0$, there exist a neighbourhood
$\mathcal{N}'_w\subset\mathbb{C}$ of $[-\varrho_2/2,\varrho_2/2]$ and constants $C$, $n_0$,
independent of $n$ and of the offset, and depending on $\vec f$ only through $\mathsf{p}$,
$\mathsf{d}$ and a length $D$ associated with the supports of the $f_i$, such that for all
$n\ge n_0$ and all $\hat w\in\mathcal{N}'_w$
\begin{equation*}
\bigl|S^{n}_{m}(a_n(\tilde X)-\hat w;\vec f)-S^{(\hat w)}_{\tilde X,m}(\vec f)\bigr|
\le C\prod_{i=1}^{\mathsf{p}}\|f_i\|_{\sharp}\; n^{-1/2}(\log n)^{p_0+1},
\end{equation*}
where $S^{(\hat w)}_{\tilde X,m}(\vec f)$ is the member $\hat w$ of the one-parameter family
of clause \textup{(v-b)} of Theorem~0; the functions
$\hat w\mapsto S^{n}_{m}(a_n(\tilde X)-\hat w;\vec f)$ are holomorphic on
$\mathcal{N}'_w$, so that their uniform limit $\hat w\mapsto S^{(\hat w)}_{\tilde X,m}(\vec f)$
is continuous on $\mathcal{N}'_w$ and real-analytic on $]-\varrho_2/2,\varrho_2/2[$.
\end{enumerate}
\item[(e)] The first half of the corollary of the uniqueness theorem at reference $\tilde g$:
for every bare inverse coupling $x$ with $x^{-1/2}\le g_-(\tilde g)$ whose own-threshold run
first passes at depth $K=K(x^{-1/2},\tilde g)$,
\begin{equation*}
|x-a_K(\tilde X)|\le\tfrac32\bigl(B^{\sharp}+\varpi(\tilde g)\bigr),
\end{equation*}
where $\varpi(\tilde g)\le1$ is the cross-scaffold modulus of the scaffold comparison.
\end{enumerate}
Then for every real $\hat w$ with $|\hat w|\le\varrho_2/2$ there exist an integer $\sigma$ and a
real $\hat w'$ with $|\hat w'|\le\varrho_2/2$, obtained from $(L,\mathfrak{p},g,g',\hat w)$
alone by a subsequence selection on torus-free numbers (they depend neither on the torus
exponent $m_0$ nor on the test functions $\vec f$ below), such that the following hold.
\begin{enumerate}
\item[(i)] Location of the exponent: there are integers $\delta$, $s$ with
\begin{equation}\label{9.13:eq-limit-member-a-bis}
\begin{gathered}
\sigma=\delta+s,\qquad |\delta|\le j_U-1,\qquad s\,u\ge0,\\
\frac{|u|-B^{\sharp}}{B^{\sharp}}<|s|<\frac{|u|+B^{\sharp}+c_5}{\lambda^{\sharp}} .
\end{gathered}
\end{equation}
In particular, for $u\ge0$,
\begin{equation*}
\frac{u}{B^{\sharp}}-j_U<\sigma<\frac{u+B^{\sharp}+c_5}{\lambda^{\sharp}}+j_U-1,
\end{equation*}
so that $\sigma$ is bounded above and below by linear functions of $u$ which both tend to
$+\infty$ with $u$.
\item[(ii)] The identity: put $\bar s^{-}(u):=0$ if $u\ge0$ and
$\bar s^{-}(u):=\lfloor(|u|+B^{\sharp}+c_5)/\lambda^{\sharp}\rfloor$ if $u<0$, and
\begin{equation*}
m^{\Lambda}(g,g'):=\max\bigl\{m_{\mathbb{T}}(g_-(g)),\,m_{\mathbb{T}}(g_-(g')),\,0\bigr\}
+j_U+\bar s^{-}(u),
\end{equation*}
a floor on the torus exponent which depends on $g$ and $g'$ only. For every
$m_0\ge m^{\Lambda}(g,g')$, every $\mathsf{p}\ge2$ and every $\mathsf{p}$-tuple $\vec f$ of
real Lipschitz functions on the continuum torus of side $2L^{m_0}\ell_g$ with pairwise
separated supports,
\begin{equation}\label{9.13:eq-limit-member-b-bis}
S^{(\hat w)}_{X,m_0}(\vec f)=S^{(\hat w')}_{X',m_0+\sigma}\bigl(\vec f\circ L^{-\sigma}\cdot\bigr).
\end{equation}
The right member is the member $\hat w'$ of the one-parameter family of clause \textup{(v-b)}
at reference coupling $g'$ on the torus of exponent $m_0+\sigma$, of side
$2L^{m_0+\sigma}\ell_{g'}$ with $\ell_{g'}=L^{-\sigma}\ell_g$, which is the same continuum
torus; that is, it is the full limit along its own tuned family
$(a_{n'}(X')-\hat w')_{n'\ge0}$ as in \textup{(d2)} at $g'$, evaluated at the same test
functions written in the reference length $\ell_{g'}$.
\item[(iii)] The coupling bookkeeping: $\sigma$, $\delta$, $s$ and $\hat w'$ are obtained along
a strictly increasing sequence of indices $n_j$, depending on $(L,\mathfrak{p},g,g',\hat w)$
only, such that for each $j$ one
own-threshold run of clause \textup{(0)}, that from the common bare datum
$a_{n_j}(X)-\hat w$, visits the window $]X,X+B^{\sharp}]$ at level $n_j-\delta$ and the window
$]X',X'+B^{\sharp}]$ at level $n_j-\sigma$, these levels being $|s|=|\sigma-\delta|$ apart.
Hence, for $u\ge0$,
\begin{equation}\label{9.13:eq-limit-member-c-bis}
\lambda^{\sharp}s-c_5-B^{\sharp}<X'-X<B^{\sharp}(s+1),
\end{equation}
and the inverse couplings of that run at the two levels differ by an amount in
$[\lambda^{\sharp}s-c_5,\,B^{\sharp}s]$, in accordance with clause \textup{(0)} of Theorem~0.
\end{enumerate}
\end{proposition}

\begin{proof}
Fix a real $\hat w$ with $|\hat w|\le\varrho_2/2$ and put
\begin{equation*}
x_n:=a_n(X)-\hat w\qquad(n\ge0),
\end{equation*}
the bare data of the member $\hat w$ at reference $g$. Since
$\hat w\in[-\varrho_2/2,\varrho_2/2]\subset\mathcal{N}'_w$, hypothesis (d2) at reference $g$
gives, for every torus exponent $m_0$ at which it applies and for every admissible tuple
$\vec f$, the convergence of the full sequence
\begin{equation*}
S^{n}_{m_0}(x_n;\vec f)\longrightarrow S^{(\hat w)}_{X,m_0}(\vec f)=:\mathfrak{S}(\vec f)
\qquad(n\to\infty).
\end{equation*}

\textit{Step 1: both references see $x_n$ as a first-passage datum.} By (d1) at reference $g$,
$a_n(X)\ge X+\lambda_U n-2c_2$; since $\hat w\le|\hat w|\le\varrho_2/2$, this gives
\begin{equation*}
x_n\ge X+\lambda_U n-2c_2-\varrho_2/2 .
\end{equation*}
Since $\lambda_U>0$, there is a least integer $\bar n\ge0$ with
$\lambda_U\bar n-2c_2-\varrho_2/2\ge\max\{B^{\sharp},u+B^{\sharp}\}$, and the same inequality
holds for every $n\ge\bar n$. For $n\ge\bar n$ we therefore have
\begin{equation*}
x_n\ge X+\max\{B^{\sharp},u+B^{\sharp}\}=\max\{X+B^{\sharp},X'+B^{\sharp}\}>0 .
\end{equation*}
Because $g_-(g)=(X+B^{\sharp})^{-1/2}$ and $g_-(g')=(X'+B^{\sharp})^{-1/2}$, this says that
$g_0^{(n)}:=x_n^{-1/2}$ satisfies $g_0^{(n)}\le\min\{g_-(g),g_-(g')\}$. Thus $g_0^{(n)}$ is an
admissible bare coupling of clause \textup{(0)} at reference $g$ and at reference $g'$, by
\textup{(H6)} at each, and hypothesis (b) applies at both references. Write
$x^{\mathrm{own}}_{n,k}$ for the inverse coupling at level $k$ of the own-threshold run from
$g_0^{(n)}$ and put
\begin{equation*}
K_n:=K(g_0^{(n)},g),\qquad K'_n:=K(g_0^{(n)},g').
\end{equation*}
By (b), $x^{\mathrm{own}}_{n,K_n}\in\,]X,X+B^{\sharp}]$ and
$x^{\mathrm{own}}_{n,K'_n}\in\,]X',X'+B^{\sharp}]$. By (c), all these numbers are functions of
$(g_0^{(n)},L,\mathfrak{p})$ and of the two references only.

\textit{Step 2: $|K_n-n|\le j_U-1$.} Let $n\ge\bar n$. By definition of $K_n$, the own-threshold
run of the bare inverse coupling $x_n$ first passes at depth $K_n$ at reference $g$, and
$x_n^{-1/2}\le g_-(g)$ by Step~1. Hypothesis (e) at reference $g$ applies and, with
$\varpi(g)\le1$ and the floor $\varrho\ge6(B^{\sharp}+1)$, gives
\begin{equation*}
|x_n-a_{K_n}(X)|\le\tfrac32\bigl(B^{\sharp}+\varpi(g)\bigr)\le\tfrac32\,(B^{\sharp}+1)
\le\frac{\varrho}{4}=\frac{\varrho_2}{2}.
\end{equation*}
Since $a_n(X)-x_n=\hat w$, the triangle inequality gives
\begin{equation*}
|a_n(X)-a_{K_n}(X)|\le|\hat w|+|x_n-a_{K_n}(X)|\le\frac{\varrho_2}{2}+\frac{\varrho_2}{2}
=\varrho_2 .
\end{equation*}
The spacing inequality of (d1) at reference $g$, with $l=K_n$, gives
$\lambda_U|n-K_n|-c'\le\varrho_2$, hence
\begin{equation*}
|n-K_n|\le\frac{\varrho_2+c'}{\lambda_U}\le\Bigl\lceil\frac{\varrho_2+c'}{\lambda_U}\Bigr\rceil
=j_U-1 .
\end{equation*}
Put $\delta_n:=n-K_n$; it is an integer in $[1-j_U,\,j_U-1]$.

\textit{Step 3: the two first passages are $|s_n|$ levels apart.} Let $n\ge\bar n$ and put
$s_n:=K_n-K'_n$.

Suppose first $u\ge0$, that is $X'\ge X$. For every level $k<K'_n$ the minimality in the
definition of $K'_n$ gives $x^{\mathrm{own}}_{n,k}>X'+B^{\sharp}\ge X+B^{\sharp}$, so no level
below $K'_n$ satisfies the condition defining $K_n$; hence $K_n\ge K'_n$ and $s_n\ge0$. The
inequalities of (b) at reference $g$ apply to the pair $(i,j)=(K'_n,K_n)$, since
$0\le K'_n\le K_n=K(g_0^{(n)},g)$, and give
\begin{equation*}
\lambda^{\sharp}s_n-c_5\le x^{\mathrm{own}}_{n,K'_n}-x^{\mathrm{own}}_{n,K_n}\le B^{\sharp}s_n .
\end{equation*}
The two windows of Step~1 give
\begin{equation*}
X'-X-B^{\sharp}<x^{\mathrm{own}}_{n,K'_n}-x^{\mathrm{own}}_{n,K_n}<X'-X+B^{\sharp}.
\end{equation*}
Combining, $B^{\sharp}s_n>u-B^{\sharp}$ and $\lambda^{\sharp}s_n-c_5<u+B^{\sharp}$, that is
\begin{equation*}
\frac{u-B^{\sharp}}{B^{\sharp}}<s_n<\frac{u+B^{\sharp}+c_5}{\lambda^{\sharp}} .
\end{equation*}
For $0\le u<B^{\sharp}$ the left member is negative and the information $s_n\ge0$ obtained
above stands beside it.

Suppose now $u<0$, that is $X'<X$. Symmetrically, for every $k<K_n$ we have
$x^{\mathrm{own}}_{n,k}>X+B^{\sharp}>X'+B^{\sharp}$, so $K'_n\ge K_n$ and $s_n\le0$, with
$|s_n|=K'_n-K_n$. The inequalities of (b) at reference $g'$ apply to the pair
$(i,j)=(K_n,K'_n)$, since $0\le K_n\le K'_n=K(g_0^{(n)},g')$, and give
\begin{equation*}
\lambda^{\sharp}|s_n|-c_5\le x^{\mathrm{own}}_{n,K_n}-x^{\mathrm{own}}_{n,K'_n}\le B^{\sharp}|s_n| ;
\end{equation*}
the two windows give
\begin{equation*}
|u|-B^{\sharp}<x^{\mathrm{own}}_{n,K_n}-x^{\mathrm{own}}_{n,K'_n}<|u|+B^{\sharp},
\end{equation*}
and exactly as before
\begin{equation*}
\frac{|u|-B^{\sharp}}{B^{\sharp}}<|s_n|<\frac{|u|+B^{\sharp}+c_5}{\lambda^{\sharp}} .
\end{equation*}
In all cases $s_n u\ge0$, the two-sided bound on $|s_n|$ holds, and the integer $s_n$ takes
finitely many values. Hence so does
\begin{equation*}
\sigma_n:=n-K'_n=(n-K_n)+(K_n-K'_n)=\delta_n+s_n .
\end{equation*}

\textit{Step 4: the offsets at reference $g'$ lie in the window.} Let $n\ge\bar n$. By definition
of $K'_n$ the own-threshold run of $x_n$ first passes at depth $K'_n$ at reference $g'$, and
$x_n^{-1/2}\le g_-(g')$ by Step~1; moreover $g'\le\gamma_{U'}$. Hypothesis (e) at reference
$g'$, with $\varpi(g')\le1$ and the floor on $\varrho$ as in Step~2, gives
\begin{equation*}
|x_n-a_{K'_n}(X')|\le\tfrac32\bigl(B^{\sharp}+\varpi(g')\bigr)\le\frac{\varrho_2}{2}.
\end{equation*}
Put $\hat w'_n:=a_{K'_n}(X')-x_n$. Then $\hat w'_n\in[-\varrho_2/2,\varrho_2/2]$ and
\begin{equation*}
x_n=a_{K'_n}(X')-\hat w'_n .
\end{equation*}

\textit{Step 5: one subsequence, free of the torus and of the test functions.} By Steps~1--4,
the sequences $(\delta_n)$, $(s_n)$, $(\sigma_n)$ and $(\hat w'_n)$, $n\ge\bar n$, are built
from $x_n=a_n(X)-\hat w$ (which depends on $L,\mathfrak{p},g,\hat w$), from the depths $K_n$,
$K'_n$ (which by (c) depend on $g_0^{(n)},L,\mathfrak{p}$ and the references) and from
$a_{K'_n}(X')$ (which depends on $L,\mathfrak{p},g'$); thus they depend on
$(L,\mathfrak{p},g,g',\hat w)$ only, and neither on a torus exponent nor on a test function.
Since the pair $(\delta_n,s_n)$ takes finitely
many values (Steps~2 and~3), some value $(\delta,s)$ is taken for infinitely many
$n\ge\bar n$; along these indices the $\hat w'_n$ lie in the compact interval
$[-\varrho_2/2,\varrho_2/2]$, so by the Bolzano--Weierstrass theorem there are indices
$\bar n\le n_1<n_2<\cdots$ with
\begin{equation*}
\delta_{n_j}=\delta,\qquad s_{n_j}=s,\qquad \sigma_{n_j}=\sigma:=\delta+s,\qquad
\hat w'_{n_j}\to\hat w'\in[-\varrho_2/2,\varrho_2/2].
\end{equation*}
These constants are obtained from $(L,\mathfrak{p},g,g',\hat w)$ alone. Note that
$K'_{n_j}=n_j-\sigma\to\infty$.

Statement \eqref{9.13:eq-limit-member-a-bis} is Steps~2 and~3 at these constants:
$|\delta|\le j_U-1$ by Step~2, and $s u\ge0$ together with the two-sided bound on $|s|$ by
Step~3. For $u\ge0$ we have $|s|=s$, and with $1-j_U\le\delta\le j_U-1$ and
\begin{equation*}
\frac{u-B^{\sharp}}{B^{\sharp}}<s<\frac{u+B^{\sharp}+c_5}{\lambda^{\sharp}}
\end{equation*}
we obtain
\begin{align*}
\sigma&=\delta+s>(1-j_U)+\frac{u}{B^{\sharp}}-1=\frac{u}{B^{\sharp}}-j_U,\\
\sigma&=\delta+s<\frac{u+B^{\sharp}+c_5}{\lambda^{\sharp}}+j_U-1,
\end{align*}
which is the bound on $\sigma$ stated in (i). For (iii), fix $j$. The own-threshold run from
the bare inverse coupling $x_{n_j}=a_{n_j}(X)-\hat w$ satisfies, by Step~1,
$x^{\mathrm{own}}_{n_j,K_{n_j}}\in\,]X,X+B^{\sharp}]$ and
$x^{\mathrm{own}}_{n_j,K'_{n_j}}\in\,]X',X'+B^{\sharp}]$, where $K_{n_j}=n_j-\delta_{n_j}=n_j-\delta$
and $K'_{n_j}=n_j-\sigma_{n_j}=n_j-\sigma$; these two levels are
$|K_{n_j}-K'_{n_j}|=|s|=|\sigma-\delta|$ apart. For $u\ge0$, the inequalities
$B^{\sharp}s_n>u-B^{\sharp}$ and $\lambda^{\sharp}s_n-c_5<u+B^{\sharp}$ of Step~3 at $n=n_j$,
where $s_{n_j}=s$ and $u=X'-X$, read $X'-X<B^{\sharp}(s+1)$ and
$X'-X>\lambda^{\sharp}s-c_5-B^{\sharp}$, which is \eqref{9.13:eq-limit-member-c-bis}; and the first
display of Step~3 at $n=n_j$ says that the inverse couplings of that run at the levels
$K'_{n_j}$ and $K_{n_j}$ differ by an amount in $[\lambda^{\sharp}s-c_5,\,B^{\sharp}s]$.

\textit{Step 6: the identity.} Fix $m_0\ge m^{\Lambda}(g,g')$, an integer $\mathsf{p}\ge2$ and a
$\mathsf{p}$-tuple $\vec f$ of real Lipschitz functions on the continuum torus of side
$2L^{m_0}\ell_g$ with supports pairwise at distance at least $\mathsf{d}>0$. Since
$m_0\ge\max\{m_{\mathbb{T}}(g_-(g)),0\}+j_U$, the torus exponent $m_0$ is admissible at
reference $g$ in the sense of \textup{(H5)}, so the convergence
$S^{n}_{m_0}(x_n;\vec f)\to\mathfrak{S}(\vec f)$ recorded at the beginning of the proof holds.

We bound $\sigma$ from below. By Step~2, $\delta\ge1-j_U$. If $u\ge0$, then $s\ge0=-\bar s^{-}(u)$.
If $u<0$, then $s\le0$ and $|s|<(|u|+B^{\sharp}+c_5)/\lambda^{\sharp}$; since $|s|$ is an
integer, $|s|\le\lfloor(|u|+B^{\sharp}+c_5)/\lambda^{\sharp}\rfloor=\bar s^{-}(u)$, that is
$s\ge-\bar s^{-}(u)$. Hence in both cases
\begin{equation*}
m_0+\sigma\ge m_0-(j_U-1)-\bar s^{-}(u)
\ge\max\{m_{\mathbb{T}}(g_-(g')),0\}+1\ge1 ,
\end{equation*}
the middle inequality being the definition of $m^{\Lambda}(g,g')$. Thus the torus exponent
$m_0+\sigma$ is admissible for clause \textup{(v)} at reference $g'$ in the sense of
\textup{(H5)}, and it is non-negative, as (a) requires.

Next we apply the change-of-units identity (a) to the datum $(n_j,m_0,x_{n_j},\vec f)$. Since
$n_j\to\infty$, we have $\sigma\le n_j$ for all $j$ large enough, and we consider only such $j$.
If $\sigma\ge0$, we apply \eqref{9.13:eq-change-of-units-a-bis} with $(n,s):=(n_j,\sigma)$. If
$\sigma<0$, we read \eqref{9.13:eq-change-of-units-a-bis} from right to left with $s:=|\sigma|$:
we apply it with $n:=n_j+|\sigma|$, torus exponent $m_0-|\sigma|=m_0+\sigma\ge0$ and the tuple
$\vec f\circ L^{|\sigma|}\cdot$, whose composition with $L^{-|\sigma|}\cdot$ is $\vec f$. In
both cases, writing
\begin{equation*}
\vec h:=\vec f\circ L^{-\sigma}\cdot ,
\end{equation*}
and using $n_j-\sigma=K'_{n_j}$ together with $x_{n_j}=a_{K'_{n_j}}(X')-\hat w'_{n_j}$ from
Step~4, we obtain
\begin{equation*}
\begin{split}
S^{n_j}_{m_0}(x_{n_j};\vec f)
&=S^{n_j-\sigma}_{m_0+\sigma}(x_{n_j};\vec h)\\
&=S^{K'_{n_j}}_{m_0+\sigma}\bigl(a_{K'_{n_j}}(X')-\hat w'_{n_j};\vec h\bigr).
\end{split}
\end{equation*}
The last member is the object of (d2) at reference $g'$, at torus exponent $m_0+\sigma$, length
$K'_{n_j}$ and offset $\hat w'_{n_j}$, which is real with $|\hat w'_{n_j}|\le\varrho_2/2$ and
hence lies in $\mathcal{N}'_w$. By (a), $\vec h$ is a tuple of real Lipschitz functions on the
torus of side $2L^{m_0+\sigma}\ell_{g'}$ with pairwise separated supports, the separation being
$L^{\sigma}\mathsf{d}>0$.

The left members $S^{n_j}_{m_0}(x_{n_j};\vec f)$ converge to $\mathfrak{S}(\vec f)$, being a
subsequence of a convergent sequence. For the right members, (d2) at $(g',m_0+\sigma)$, applied
to the tuple $\vec h$, gives for $K'_{n_j}\ge n_0$
\begin{equation*}
\begin{split}
&\bigl|S^{K'_{n_j}}_{m_0+\sigma}\bigl(a_{K'_{n_j}}(X')-\hat w'_{n_j};\vec h\bigr)
-S^{(\hat w'_{n_j})}_{X',m_0+\sigma}(\vec h)\bigr|\\
&\qquad\le C\prod_{i=1}^{\mathsf{p}}\|h_i\|_{\sharp}\;
(K'_{n_j})^{-1/2}(\log K'_{n_j})^{p_0+1},
\end{split}
\end{equation*}
where $C$ is the constant of (d2) with $\mathsf{d}$ and $D$ replaced by $L^{\sigma}\mathsf{d}$ and
$L^{\sigma}D$. Since this bound is uniform in the offset in $\mathcal{N}'_w$, it applies at the
varying offsets $\hat w'_{n_j}$, and since $K'_{n_j}\to\infty$ its right side tends to $0$ as
$j\to\infty$. Moreover, by (d2) the function $\hat w\mapsto S^{(\hat w)}_{X',m_0+\sigma}(\vec h)$
is continuous on $\mathcal{N}'_w$, which contains the closed real segment
$[-\varrho_2/2,\varrho_2/2]$ in which $\hat w'_{n_j}$ and $\hat w'$ lie; hence
$S^{(\hat w'_{n_j})}_{X',m_0+\sigma}(\vec h)\to S^{(\hat w')}_{X',m_0+\sigma}(\vec h)$. Therefore
the right members converge to $S^{(\hat w')}_{X',m_0+\sigma}(\vec h)$, and since left and right
members are equal for every large $j$,
\begin{equation*}
\mathfrak{S}(\vec f)=S^{(\hat w')}_{X',m_0+\sigma}\bigl(\vec f\circ L^{-\sigma}\cdot\bigr),
\end{equation*}
which is \eqref{9.13:eq-limit-member-b-bis}. The subsequence $(n_j)$, the integer $\sigma$ and the
offset $\hat w'$ were fixed in Step~5 before $m_0$ and $\vec f$ were chosen, so the identity
holds with the same $\sigma$ and $\hat w'$ for all admissible $m_0$ and $\vec f$ at once.
\end{proof}

\begin{remark}[The prescribed-exponent refinement]\label{9.13:rem-prescribed-exponent}
Let $g$ satisfy the hypotheses of Lemma~\ref{9.13:lem-frozen-coupling} and of
Proposition~\ref{9.13:prop-limit-member}, and let $X=g^{-2}$. Fix a depth $s\ge 0$ and a real offset
$\hat w$ with
$|\hat w|\le\varrho_2/2$, and let $\hat x_s(\hat w)$ be the limit of the depth-$s$ frozen inverse
coupling given by Lemma~\ref{9.13:lem-frozen-coupling}. Let $X'$ be any number with
$|X'-\hat x_s(\hat w)|\le r_J/8$, for instance $\hat x_s(\hat w)$ itself.

The limit-member identification (Proposition~\ref{9.13:prop-limit-member}) gives the identity
\eqref{9.13:eq-limit-member-b-bis} with a dilation exponent $\sigma$ of which only a window is known, not
a prescribed value, and with an offset label $\hat w'$. To read the dilated tuned family directly as an
admissible family of the uniqueness theorem at the prescribed reference $X'$, one would need the
following statement:
\begin{equation}\label{9.13:eq-prescribed-exponent-1}
\begin{gathered}
\text{for every sufficiently large } n',\ \text{the } \Theta_{X'}\text{-run of length } n'\\
\text{from the bare datum } a_{n'+s}(X)-\hat w \text{ is alive, and}\\
|x_{(0)}-X'|\le \varrho_2/4 \text{ for its inverse coupling } x_{(0)} \text{ at depth } 0 .
\end{gathered}
\end{equation}

Statement \eqref{9.13:eq-prescribed-exponent-1} is a comparison made at one cutoff and from one bare
datum. The two runs compared are built on the frozen scaffolds $\Theta_{X+c_0 s}$ and $\Theta_{X'}$.
Their reference lines differ, at every depth, by the same constant $X+c_0 s-X'$. This constant is
bounded by $2c_2$ plus the contribution of the two-loop reference increments $\beta^0_l$. It therefore
lies well inside half the radius of the shift
tube at every depth $r$, since that radius is of order
$\theta_r^{1/2}(\log\theta_r)^{p_0}$.

The comparison is of the same kind as the scaffold comparison, which compares, at one cutoff, the run
carrying its own threshold with the run on a frozen scaffold. In
\eqref{9.13:eq-prescribed-exponent-1} both runs are on frozen scaffolds. This case is simpler in that
the bookkeeping of neither run depends on the complex parameter that enters the bookkeeping of the
scaffold comparison.

If \eqref{9.13:eq-prescribed-exponent-1} were available, it would add two things.
\begin{itemize}
\item[(a)] The identity would hold in a prescribed-exponent form: one dilates by any prescribed depth
$s$ and lands in the family attached to $\Theta_{\hat x_s(\hat w)}$, in place of an exponent $\sigma$
known only up to a window.
\item[(b)] The dilated member would be identified by its label. Its label in the $g'$-frame, namely
the depth-$0$ limit $\hat x'$ attached to $\hat w'$ by Lemma~\ref{9.13:lem-frozen-coupling} with the
reference $X'$ in place of $X$, would equal $\hat x_s(\hat w)$ up to $2\varpi(g)$, where
$\varpi(g)\le 1$ is the cross-scaffold modulus of the scaffold comparison.
\end{itemize}
It would add nothing to the existence of the limits or to their convergence. It would also add
nothing to the identity \eqref{9.13:eq-limit-member-b-bis}.

Clause~(v) of Theorem~\ref{3.1:thm-main} asserts the identity only in the form with an exponent
$\sigma$ confined to a window and with offset labels; in that form it is the identity
\eqref{9.13:eq-limit-member-b-bis}, which Proposition~\ref{9.13:prop-limit-member} proves.
The refinements (a) and (b) are not asserted by Theorem~\ref{3.1:thm-main}. Statement
\eqref{9.13:eq-prescribed-exponent-1} is not proved in this book and is not used anywhere in it.
\end{remark}

\begin{remark}[Located exponent, both directions, infinite volume]\label{9.13:rem-located-exponent}
Let $g,g'$ be the two reference couplings of the limit-member identification,
Proposition~\ref{9.13:prop-limit-member}, with $X=g^{-2}$, $X'=g'^{-2}$ and $u=X'-X$. For a real
$\hat w$ with $|\hat w|\le\varrho_2/2$, let $\sigma=\delta+s$ and $\hat w'$ be the integer and the
offset label that this proposition provides.
\begin{enumerate}
\item[(a)] Proposition~\ref{9.13:prop-limit-member} is proved as an implication from the following
statements:
\begin{itemize}
\item the first half of the corollary of the uniqueness theorem, taken at reference coupling $g$ and
at reference coupling $g'$;
\item clause (0) of Theorem~\ref{3.1:thm-main}, namely the first passage $K(g_0,g)$, the two-sided
flow inequalities and their corollary;
\item hypotheses (H5) and (H6) of Definition~\ref{2.2:not-hypotheses}.
\end{itemize}
The lower bound $\varrho\ge 6(B^{\sharp}+1)$ and the threshold $\gamma_{U'}$ under which it holds
are those of clause (v-c) of Theorem~\ref{3.1:thm-main}.

No comparison at one cutoff of runs on the two frozen scaffolds $\Theta_X$ and $\Theta_{X'}$
enters. The only cross-scaffold comparison used is the one contained in the first half of that
corollary, between the own-threshold run of clause (0) and the $\Theta_{X'}$-run at first passage.
The own-threshold run does not depend on a scaffold, and it links the tuned family at reference $g$
to the tuned family at reference $g'$.
\item[(b)] The exponent $\sigma$ is located in the window \eqref{9.13:eq-limit-member-a-bis}; it is not
prescribed. Whether every integer of that window occurs as $\sigma$ depends on the sign, at every
step, of the increment of the own-threshold run, and clause (0) of Theorem~\ref{3.1:thm-main} does
not assert such a sign.

The label $\hat w'$ (equivalently, the label $\hat x'$ of the member in the frame of reference $g'$)
comes from a subsequence selection in the proof of Proposition~\ref{9.13:prop-limit-member}; it is
obtained by compactness and is not given as a function of $(X,\hat w,s)$. Relating it to the
limit $\hat x_s(\hat w)$ of Lemma~\ref{9.13:lem-frozen-coupling} requires a comparison of runs on two
frozen scaffolds, which is the subject of Remark~\ref{9.13:rem-prescribed-exponent}.

Different subsequences may produce different pairs $(\sigma,\hat w')$. This is consistent with the
statements that clause (v-c) of Theorem~\ref{3.1:thm-main} does not assert.
\item[(c)] Both signs of $u$ are covered. If $u>0$ then $s\ge0$, and if $u=0$ then $s=0$. If $u<0$,
the coarsening direction $X'<X$, then $s\le0$, and the torus floor $m^{\Lambda}(g,g')$ contains the
additional term $\bar s^{-}(u)=\lfloor(|u|+B^{\sharp}+c_5)/\lambda^{\sharp}\rfloor$, which is at
least $|s|$ and so compensates the negative $s$. In every case every $m_0\ge m^{\Lambda}(g,g')$
satisfies
\begin{equation}\label{9.13:eq-located-exponent-1}
m_0+\sigma\ \ge\ \max\{m_{\mathbb{T}}(g_-(g)),\,m_{\mathbb{T}}(g_-(g')),\,0\}+1 .
\end{equation}
If $X'=X$, then $s=0$ and $\sigma=\delta$ with $|\delta|\le j_U-1$, and
\eqref{9.13:eq-limit-member-b-bis} reads
\begin{equation}\label{9.13:eq-located-exponent-2}
S^{(\hat w)}_{X,m_0}(\vec f)=S^{(\hat w')}_{X,m_0+\delta}(\vec f\circ L^{-\delta}\cdot).
\end{equation}
Thus members at $(X,m_0)$, in particular those with large $|\hat w|$, may be dilates or contractions
(according to the sign of $\delta$) of members at the neighbouring torus exponents $m_0+\delta$.
Thus the families of members realised at the neighbouring torus exponents $m_0$ and $m_0+\delta$
overlap.
\item[(d)] In particular, the per-torus dilation relation \eqref{9.13:eq-limit-member-b-bis} between
members, in the offset parametrisation $(X,\hat w)$ and with $\sigma$ located in the window
\eqref{9.13:eq-limit-member-a-bis} (the window of $s$ widened by $j_U-1$ on each side), uses neither the
bi-Lipschitz lemma nor the comparison of runs on two frozen scaffolds of
Remark~\ref{9.13:rem-prescribed-exponent}; only the identification of the dilated member by chart
maps requires them.

Let $\mathcal{F}$ be the set of $\mathsf{p}$-tuples, $\mathsf{p}\ge2$, of real Lipschitz
functions on $\mathbb{R}^4$ with compact, pairwise separated supports. We write
$\operatorname{Lim}_{\infty}[X,\hat w]$ for the set of infinite-volume limit points, as
$m_0\to\infty$, of the member $\hat w$ at reference $X$, taken pointwise on $\mathcal{F}$; precisely,
it is the set of maps $\mathfrak{S}$ on $\mathcal{F}$ for which there
is a sequence of integers $m_0^{(j)}\to\infty$ with
\begin{equation*}
\mathfrak{S}(\vec f)=\lim_{j\to\infty}S^{(\hat w)}_{X,m_0^{(j)}}(\vec f)\qquad\text{for every }
\vec f\in\mathcal{F}.
\end{equation*}
Here $\vec f$ is read, in the unit $\ell_g$, on the torus of side $2L^{m_0^{(j)}}\ell_g$ as soon as
its supports lie in $(-L^{m_0^{(j)}}/2,L^{m_0^{(j)}}/2)^4$. Put
$(D_\sigma\mathfrak{S})(\vec f):=\mathfrak{S}(\vec f\circ L^{-\sigma}\cdot)$. Then
\begin{equation}\label{9.13:eq-located-exponent-3}
\operatorname{Lim}_{\infty}[X,\hat w]=D_\sigma\operatorname{Lim}_{\infty}[X',\hat w'] ,
\end{equation}
with the same $\sigma$ and $\hat w'$ as in \eqref{9.13:eq-limit-member-b-bis}. This identity is stated
in offset labels. The form labelled by the limits $\hat x_s(\hat w)$ is not obtained here.
\end{enumerate}
\end{remark}

\begin{proof}
Items (a) and (b), and the first paragraph of (d), describe the proof of
Proposition~\ref{9.13:prop-limit-member} and the content of its conclusion. We prove (c) and the
identity \eqref{9.13:eq-located-exponent-3} of (d).

(c) Since $s\cdot u\ge0$ by \eqref{9.13:eq-limit-member-a-bis}, $u>0$ forces $s\ge0$ and $u<0$ forces
$s\le0$. The case $u=0$ is treated below.

Let $u\ge0$. Then $\bar s^{-}(u)=0$. Moreover, by \eqref{9.13:eq-limit-member-a-bis} in its form for
$u\ge0$,
\begin{equation*}
\sigma\ >\ u/B^{\sharp}-j_U\ \ge\ -j_U ,
\end{equation*}
and since $\sigma$ is an integer, $\sigma\ge-(j_U-1)$. Hence $m_0\ge m^{\Lambda}(g,g')$ gives
\begin{equation*}
m_0+\sigma\ \ge\ \max\{m_{\mathbb{T}}(g_-(g)),m_{\mathbb{T}}(g_-(g')),0\}+j_U-(j_U-1),
\end{equation*}
which is \eqref{9.13:eq-located-exponent-1}.

Let $u<0$. By \eqref{9.13:eq-limit-member-a-bis}, the integer $|s|$ is strictly smaller than
$(|u|+B^{\sharp}+c_5)/\lambda^{\sharp}$. An integer strictly smaller than a real number is at most
its integer part, so $|s|\le\bar s^{-}(u)$. Therefore
\begin{equation*}
\sigma\ \ge\ -(j_U-1)-\bar s^{-}(u),
\end{equation*}
and $m_0\ge m^{\Lambda}(g,g')$ again gives \eqref{9.13:eq-located-exponent-1}.

Now let $X'=X$. In the proof of Proposition~\ref{9.13:prop-limit-member}, whose coupling
bookkeeping \eqref{9.13:eq-limit-member-c-bis} records, $\delta$ is the gap between the index $n$ and
the first passage of one own-threshold run at reference $g$, and $s$ is the separation between the
first passages of that same run at references $g$ and $g'$; these first passages are the levels
$n-\delta$ and $n-\sigma$ at which the run visits the windows $]X,X+B^{\sharp}]$ and
$]X',X'+B^{\sharp}]$. For $X'=X$ we have $g'=g$, so the two first passages coincide. Hence
$s=\sigma-\delta=0$ and $\sigma=\delta$, with $|\delta|\le j_U-1$ by
\eqref{9.13:eq-limit-member-a-bis}. Substituting $X'=X$ and $\sigma=\delta$ into
\eqref{9.13:eq-limit-member-b-bis} gives \eqref{9.13:eq-located-exponent-2}.

(d) By Proposition~\ref{9.13:prop-limit-member}, $\sigma$ and $\hat w'$ are obtained from
$(L,\mathfrak{p},g,g',\hat w)$ alone, so they do not depend on $m_0$ or on $\vec f$. Hence
\eqref{9.13:eq-limit-member-b-bis} holds with the same $(\sigma,\hat w')$ for every
$m_0\ge m^{\Lambda}(g,g')$.

Let $\vec f\in\mathcal{F}$ have supports in $(-L^{m_0}/2,L^{m_0}/2)^4$. Two points of this cube
differ by less than $L^{m_0}$ in each coordinate, which is half the period $2L^{m_0}$. So their
distance on the torus equals their Euclidean distance, and $\vec f$ is an admissible tuple on the
torus of exponent $m_0$. The supports of $\vec f\circ L^{-\sigma}\cdot$ lie in
$(-L^{m_0+\sigma}/2,L^{m_0+\sigma}/2)^4$. By the same reasoning, this tuple is admissible, in the
unit $\ell_{g'}=L^{-\sigma}\ell_g$, on the torus of exponent $m_0+\sigma$. Finally,
$\vec h\in\mathcal{F}$ if and only if $\vec h\circ L^{\sigma}\cdot\in\mathcal{F}$.

Let $\mathfrak{S}\in\operatorname{Lim}_{\infty}[X,\hat w]$ along $m_0^{(j)}\to\infty$. Fix
$\vec h\in\mathcal{F}$ and put $\vec f=\vec h\circ L^{\sigma}\cdot\in\mathcal{F}$. For all large $j$,
\eqref{9.13:eq-limit-member-b-bis} gives
\begin{equation*}
S^{(\hat w')}_{X',m_0^{(j)}+\sigma}(\vec h)=S^{(\hat w)}_{X,m_0^{(j)}}(\vec f)
\ \longrightarrow\ \mathfrak{S}(\vec h\circ L^{\sigma}\cdot)\qquad(j\to\infty).
\end{equation*}
So $\mathfrak{S}'(\vec h):=\mathfrak{S}(\vec h\circ L^{\sigma}\cdot)$ defines an element of
$\operatorname{Lim}_{\infty}[X',\hat w']$ along $m_0^{(j)}+\sigma$, and
$\mathfrak{S}=D_\sigma\mathfrak{S}'$.

Conversely, let $\mathfrak{S}'\in\operatorname{Lim}_{\infty}[X',\hat w']$ along $m'^{(j)}\to\infty$,
and put $m_0^{(j)}:=m'^{(j)}-\sigma\to\infty$. For every $\vec f\in\mathcal{F}$ and all large $j$,
\eqref{9.13:eq-limit-member-b-bis} gives
\begin{equation*}
S^{(\hat w)}_{X,m_0^{(j)}}(\vec f)=S^{(\hat w')}_{X',m'^{(j)}}(\vec f\circ L^{-\sigma}\cdot)
\ \longrightarrow\ (D_\sigma\mathfrak{S}')(\vec f).
\end{equation*}
Hence $D_\sigma\mathfrak{S}'\in\operatorname{Lim}_{\infty}[X,\hat w]$. Together the two inclusions
prove \eqref{9.13:eq-located-exponent-3}.
\end{proof}

\begin{remark}[Assembly of the up-to-scale statement]\label{9.13:rem-up-to-scale}
Clause (v-d)(iii) of Theorem~\ref{3.1:thm-main}, the up-to-scale statement, consists of the
following three parts. Throughout, $X=g^{-2}$, $\varrho_2=\varrho/2$, $\mathsf{p}$ is the
number of test functions, and we write $S^T_{\mathsf{p};K,m}[x_0]$ for the connected
Schwinger function $S^T_{\mathsf{p};K,m}$ at bare inverse coupling $x_0$.
The letter $s$ denotes in ($\alpha$) an arbitrary shift, in ($\beta$) and in the list below
the separation $s=\sigma-\delta$ of \eqref{9.13:eq-limit-member-a-bis}, and in ($\gamma$) a fixed
depth.

\emph{($\alpha$) Finite cutoff.} Part ($\alpha$) is the change-of-units identity,
Lemma~\ref{9.13:lem-change-of-units-bis}, that is \eqref{9.13:eq-change-of-units-a-bis}, read in the
first-passage notation of Theorem~\ref{3.1:thm-main}: for every bare inverse coupling $x_0$,
all $K\ge s\ge 0$, $m\ge 0$, $\mathsf{p}\ge 2$ and all $\mathsf{p}$-tuples $\vec f$,
\begin{equation*}
S^T_{\mathsf{p};K,m}[x_0](\vec f)=S^T_{\mathsf{p};K-s,m+s}[x_0](\vec f\circ L^{-s}\cdot).
\end{equation*}

\emph{($\beta$) Continuum limits.} Part ($\beta$) is the limit-member identification,
Proposition~\ref{9.13:prop-limit-member}, with its conclusions \eqref{9.13:eq-limit-member-a-bis},
\eqref{9.13:eq-limit-member-b-bis} and \eqref{9.13:eq-limit-member-c-bis}, under the standing
hypotheses of clause (v), the floor $\varrho\ge 6(B^{\sharp}+1)$ of clause (v-c)(2) and the
ceiling $g\le\gamma_{U'}$, with $j_U:=1+\lceil(\varrho_2+c')/\lambda_U\rceil$: for every real
$\hat w$ with $|\hat w|\le\varrho_2/2$ and every $X'$ with $X'^{-1/2}\le\gamma_{U'}$ (in
particular for every weaker reference coupling $g':=X'^{-1/2}\le g$) it provides an integer
$\sigma$, located by \eqref{9.13:eq-limit-member-a-bis} as $\sigma=\delta+s$ with
$|\delta|\le j_U-1$, $s(X'-X)\ge0$ and $|s|$ in the window displayed there, and a real
$\hat w'$ with $|\hat w'|\le\varrho_2/2$, both depending only on
$(L,\mathfrak{p},g,g',\hat w)$ and neither on the torus nor on the test functions, such that
for every $m_0\ge m^{\Lambda}(g,g')$ (the torus-exponent floor of
Proposition~\ref{9.13:prop-limit-member}; it depends on $g$ and $g'$ only) the functional
$S^{(\hat w)}_{X,m_0}$ is the exact $L^{\sigma}$-dilate of $S^{(\hat w')}_{X',m_0+\sigma}$,
that is, \eqref{9.13:eq-limit-member-b-bis} holds; and, by \eqref{9.13:eq-limit-member-c-bis}, one
run of clause (0) from the common bare data visits $(X,X+B^{\sharp}]$ and, $|s|$ levels
apart, $(X',X'+B^{\sharp}]$.

\emph{($\gamma$) The frozen coupling.} Part ($\gamma$) is the convergence of the frozen
coupling at fixed depth, Lemma~\ref{9.13:lem-frozen-coupling}: under the hypothesis
$g\le\gamma_U$ of that lemma, for every
fixed depth $s\ge0$, the inverse coupling of the run frozen on the scaffold $\Theta_X$, read
$s$ levels below the reference, converges along the tuned family $(a_n(X)-\hat w)_n$,
uniformly in real $|\hat w|\le\varrho_2/2$, to a function $\hat x_s$ real-analytic on
$|\hat w|<\varrho_2/2$, and $\hat x_s(0)-X\in[\lambda_U s-2c_2,\,3\lambda_U s+2c_2]$.

The up-to-scale statement does not assert any of the following: that the tuned bare data
$a_{n'+s}(X)$ and $a_{n'}(X')$ coincide for any $n'$ and any $X'$, with $s$ as in ($\beta$)
(they are tuned on the two different frozen scaffolds $\Theta_X$ and $\Theta_{X'}$); that
$\sigma$ can be prescribed, or that every integer in the window for $\sigma$ given by
\eqref{9.13:eq-limit-member-a-bis} occurs as $\sigma$; any chart map
$(X,\hat w)\mapsto(X',\hat w')$, or any map between the frozen
couplings $\hat x_s$ at $X$ and those at $X'$, relating the two labellings (none of
($\alpha$)--($\gamma$) uses such a map); any relation between distinct frozen couplings
$\hat x_s$ on one torus; a continuous dilation; dimensional transmutation; injectivity.
\end{remark}

\begin{proof}
Part ($\alpha$) is the change-of-units identity, Lemma~\ref{9.13:lem-change-of-units-bis}, read in
the first-passage notation of Theorem~\ref{3.1:thm-main}: the Schwinger function
$S^T_{\mathsf{p};K,m}[x_0]$ lives on the lattice with $2L^{m+K}$ sites a side (torus
exponent $m$, spacing $L^{-K}$) at bare inverse
coupling $x_0$, which is the lattice of \eqref{9.13:eq-change-of-units-a-bis} with run length
$n=K$, $m_0=m$ and $x=x_0$; the condition $0\le s\le n$ of that lemma is $K\ge s\ge 0$. The same
lemma shows that $\vec f\mapsto\vec f\circ L^{-s}\cdot$ preserves real Lipschitz tuples with
pairwise separated supports.

Part ($\beta$) is the limit-member identification, Proposition~\ref{9.13:prop-limit-member},
which is proved as an implication from the change-of-units identity
(Lemma~\ref{9.13:lem-change-of-units-bis}) and from clauses (v-b), (v-c)(1), the first sentence of
(v-c)(2) read at the two references $X$ and $X'$, (v-a) and (0) of
Theorem~\ref{3.1:thm-main}, under the floor $\varrho\ge 6(B^{\sharp}+1)$ and the ceiling
$g\le\gamma_{U'}$. Its conclusion \eqref{9.13:eq-limit-member-a-bis} locates $\sigma$; its
conclusion \eqref{9.13:eq-limit-member-b-bis} is the identity of the two members for every
$\mathsf{p}\ge2$ and every $\mathsf{p}$-tuple $\vec f$ of real Lipschitz functions on
$\mathbb{T}_{m_0}$ with pairwise separated supports. The right member of
\eqref{9.13:eq-limit-member-b-bis} is the member
$\hat w'$ of the family of clause (v-b) at reference $g'$ on the torus of exponent
$m_0+\sigma$ in the $g'$-reference unit $\ell_{g'}=L^{-\sigma}\ell_g$, hence of side
$2L^{m_0+\sigma}\ell_{g'}=2L^{m_0}\ell_g$, the same continuum torus as $\mathbb{T}_{m_0}$;
this member is the full limit along its own tuned family $(a_{n'}(X')-\hat w')_{n'}$.
Its conclusion \eqref{9.13:eq-limit-member-c-bis} is the statement that one run of clause (0)
from the common bare data visits the windows $(X,X+B^{\sharp}]$ and $(X',X'+B^{\sharp}]$,
$|s|$ levels apart. For $X'\ge X$, so that $s\ge0$, the corollary to clause (0) gives that
the couplings of that run at the two levels differ by an amount in
$[\lambda^{\sharp}s-c_5,\,B^{\sharp}s]$; since each lies in one of these windows of width
$B^{\sharp}$, this gives
\begin{equation*}
\lambda^{\sharp}s-c_5-B^{\sharp}<X'-X<B^{\sharp}(s+1).
\end{equation*}
The proposition applies to every weaker reference coupling $g'\le g$, since then
$X'^{-1/2}=g'$ and $g'\le g\le\gamma_{U'}$.

Part ($\gamma$) is the convergence of the frozen coupling at fixed depth,
Lemma~\ref{9.13:lem-frozen-coupling}, whose hypothesis is $g\le\gamma_U$. Its quantity
$x^{(n)}_{(s)}(\hat w)$ is the depth-$s$ inverse coupling of the length-$n$ run on $\Theta_X$
from the bare datum $a_n(X)-\hat w$. Since $\varrho_2=r_J/2$, the real interval
$[-\varrho_2/2,\varrho_2/2]=[-r_J/4,r_J/4]$ lies in $\overline{D}(0,r_J/4)$, on which the lemma
gives the Cauchy
bound \eqref{9.13:eq-frozen-coupling-a}, uniformly in $\hat w$. The limit $\hat x_s$ is
holomorphic on $D(0,r_J/4)$ and real on reals, hence real-analytic in real $\hat w$ with
$|\hat w|<\varrho_2/2$, and the lemma gives the window
$\hat x_s(0)-X\in[\lambda_U s-2c_2,\,3\lambda_U s+2c_2]$.
\end{proof}

The following are not asserted by clause (v) of Theorem~\ref{3.1:thm-main}; the assertions of
that clause are the content of Theorem~\ref{9.7:thm-uniqueness}.
\begin{itemize}
\item the case $\mathsf{p}=1$;
\item any statement about $\log Z$;
\item Schwinger functions at coinciding points;
\item expectations of Wilson loops;
\item the case $N\ge3$: the uniqueness theorem for $\mathrm{SU}(N)$, $N\ge3$, is
Proposition~\ref{9.14:prop-sun-uniqueness}, whose proof is outstanding;
\item uniformity in the size of the torus;
\item uniqueness of the state on the block algebra (the block-law statements of
Remark~\ref{3.2:rem-block-laws}, whose proofs are outstanding);
\item uniqueness in the volume;
\item convergence of the intermediate couplings $g_{n-s}$ along the tuned family of
Theorem~\ref{9.7:thm-uniqueness};
\item convergence of the effective actions;
\item exponential rates in any labelling of the family by bare data;
\item injectivity or non-constancy of the map $\hat w\mapsto S^{(\hat w)}$; in particular it is
not claimed that equal limits force equal parameters, and no calibration of the parameter by a
physical quantity is adopted;
\item Ba{\l}aban's uniqueness statement \cite[Theorem~2, (0.31)]{B-CMP109}, which is
neither proved nor used in this book.
\end{itemize}
The convergence, at each fixed depth $s$, of the couplings $g_{n-s}$ listed above along the
tuned family is proved as an implication from the same statements as
Theorem~\ref{9.7:thm-uniqueness}; it is not adopted among the assertions of clause (v).

\begin{proposition}[The uniqueness theorem for $\mathrm{SU}(N)$, $N\ge3$. Proof outstanding]
\label{9.14:prop-sun-uniqueness}
Let $N\ge3$. Assume the hypotheses \textup{(H1)}--\textup{(H7)} of
Notation~\ref{2.2:not-hypotheses} for the gauge group $\mathrm{SU}(N)$. Here $L_0=L_0(N)$ is the one
existential constant of \textup{(H2)}, which is never given a value, and every condition on $L$ is a
lower bound. Then the statements \textup{(v-a)}--\textup{(v-e)} of Theorem~\ref{9.7:thm-uniqueness}
hold for $G=\mathrm{SU}(N)$, $N\ge3$, with constants depending on $N$.
\end{proposition}

Proof outstanding.

\paragraph{Route.}
A proof would carry out the argument of this chapter with $G=\mathrm{SU}(N)$, $N\ge3$, in place of
$\mathrm{SU}(2)$, with the dependence on the rank $N$ kept explicit at each step. The rank enters in
three ways. First, it enters through the one-loop constant $c_0$ of
Definition~\ref{1.2:def-flow-constants}, and hence through the slopes of the flow, via the lower
bound $\lambda_\sharp\ge\underline{\lambda}=c_0/4$. Second, it enters through the existential
constant $L_0(N)$. Third, it enters through the constants of the Lie algebra $\mathfrak{su}(N)$ and
the constants built from them, namely the stability constants $E_+$ and $E_-$. Through these it
enters the thresholds, among them $\gamma_U$ of \textup{(H7)}. Every step of the chapter rests on
the correlation theorem. At $N\ge3$ the input required in its place is the $\mathrm{SU}(N)$ form of
that theorem, Proposition~\ref{5.13:prop-sun-correlation}, whose proof is outstanding.

Theorem~\ref{3.1:thm-main} asserts nothing at $N\ge3$ beyond clauses \textup{(0)} and \textup{(i)}.

\begin{definition}[Couplings, first passage, limit sets and test classes]\label{9.98:not-couplings-limit-sets}
In what follows $N=2$. The hypotheses \textup{(H1)}--\textup{(H7)} are those of
Notation~\ref{2.2:not-hypotheses}, and every symbol not introduced below has the meaning fixed in
Notation~\ref{0.2:not-glossary}.

\emph{Couplings and first passage.} We write $x_k=g_k^{-2}$ for the running inverse couplings of
Definition~\ref{1.2:def-couplings}; by clause \textup{(0)} of Theorem~\ref{3.1:thm-main} these do not
depend on the torus. The flow constants are
$\lambda_{\sharp}>0$, $B_{\sharp}$ and $c_5=2c_2$ of Definition~\ref{1.2:def-flow-constants}. For $t>0$ we
put
\begin{equation*}
g_-(t):=\bigl(t^{-2}+B_{\sharp}\bigr)^{-1/2};
\end{equation*}
the function $g_-$ is increasing.
With $M=L^{m'}$ the torus exponent is
\begin{equation*}
m_{\mathbb{T}}(g')=11+m'+\log_L R(g'),
\end{equation*}
where $R(\cdot)$ is non-decreasing as the coupling decreases; consequently $t\mapsto m_{\mathbb{T}}(g_-(t))$
is non-increasing in $t$. Both facts are proved below. The first-passage depth is
\begin{equation*}
K(g_0,g):=\min\bigl\{k\ge 0:\ x_k\le g^{-2}+B_{\sharp}\bigr\}.
\end{equation*}
Clause \textup{(0)} of Theorem~\ref{3.1:thm-main} states, with $K=K(g_0,g)$, the flow bounds
\begin{equation}\label{9.98:eq-couplings-limit-sets-a}
\lambda_{\sharp}(j-i)-c_5\ \le\ x_i-x_j\ \le\ B_{\sharp}(j-i)\qquad (0\le i\le j\le K),
\end{equation}
and, as the corollary of clause \textup{(0)}, the fixed-depth windows: for every integer $s$ with
$0\le s\le K$,
\begin{equation}\label{9.98:eq-couplings-limit-sets-b}
x_{K-s}\in\bigl]\,g^{-2}+\max\{0,\lambda_{\sharp}s-c_5\},\ g^{-2}+B_{\sharp}(s+1)\,\bigr].
\end{equation}

\emph{Constants and their order.} The radius $\varrho>0$ is that of \textup{(H7)}, fixed after $\gamma_H$;
$\varrho_2:=\varrho/2$; and $\vartheta\in\,]0,1]$ is fixed. The constants are chosen in the order
\begin{gather*}
L\to\mathfrak{p}\to\gamma\to\gamma_H\to\varrho\to\gamma_J
\to\{\text{the constants of the two-point witness theorem}\}\\
\to\gamma_W\to\gamma_{*}^{\dagger}\to\gamma_W^{\mathbb{T}}.
\end{gather*}
By \textup{(H7)},
$\gamma_J\le\gamma_H$ and $\gamma_W\le\gamma_J$. The ceiling of clause \textup{(v)} satisfies
$\gamma_U\le\gamma_H$. The further threshold $\gamma_U'$ of the corollary of the uniqueness theorem satisfies
$\gamma_U'\le\gamma_U$, under that corollary's floor $\varrho\ge 6(B_{\sharp}+1)$.

\emph{Observables.} As in Definition~\ref{1.3:def-curvature-field},
\begin{equation*}
\mathcal{O}(f)=\sum_p f(x(p))\bigl[1-\operatorname{Re}\operatorname{tr}U(\partial p)\bigr],
\end{equation*}
the sum running over all plaquettes of the bare lattice. With the generating function $Z(z)$ we put
\begin{equation*}
S^T_{n;K,m}(\vec f\,):=\partial_{z_1}\cdots\partial_{z_n}\log Z(z)\big|_{z=0}.
\end{equation*}
Further, $\|f\|_{\sharp}=\max\{\|f\|_\infty,40M\operatorname{Lip}f\}$ and $b_0=(N^2-1)/2$. The running shift
of clause \textup{(iii-c)} satisfies $\zeta'_k\in[2/3,4/3]$.

\emph{Clause \textup{(iv)}.} The constant $C_W(s)\ge1$ is non-decreasing in $s$; it is a fixed
polylogarithm of $\bar{x}_s:=g^{-2}+B_{\sharp}(s+1)$. We put
\begin{equation*}
s(\mathsf{d}):=\min\{s\ge0:\ C_W(s)L^{-s}\le\mathsf{d}\}.
\end{equation*}
The depth $s_W(g)$, the level $k_W(L)$, the constants $c_{\mathcal{K}}$, $C_{\mathcal{K}}$, the kernel
$\mathcal{K}_s$ and the test class $\mathfrak{T}(\mathsf{d};\vartheta)$ are those of clause \textup{(iv)} of
Theorem~\ref{3.1:thm-main}.

The symbols $\gamma_W^{\mathbb{T}}$, $s_{\mathrm{vol}}^{\star}(g)$, $s_T(g;\vartheta,\varrho)$,
$C_{VT}(g)$ and $\mathcal{K}^{\star}_s$ are those fixed in the statement of clause \textup{(iv)} uniformly
in the volume. In that statement, the coupling ceiling is $\gamma_W^{\mathbb{T}}$. The depth
floors are $s_{\mathrm{vol}}^{\star}(g)$ and $s_T(g;\vartheta,\varrho)$, and the depth condition reads
\begin{equation*}
\max\bigl\{s_W(g),\,s_{\mathrm{vol}}^{\star}(g),\,s_T(g;\vartheta,\varrho)\bigr\}\le s(\mathsf{d}).
\end{equation*}
The transfer constant is
\begin{equation*}
C_T(g):=160^4\,(40M/\vartheta)^2\,C_{VT}(g),
\end{equation*}
and the kernel is $\mathcal{K}^{\star}_s$.

\emph{Clause \textup{(v)}.} We write $X:=g^{-2}$. The numbers $a_n=a_n(g)$ are fixed exactly in
\textup{(v-a)}, and for the indices $n,l$ of clause \textup{(v-a)} they satisfy
\begin{equation*}
\lambda_U n-2c_2\le a_n-X\le 3\lambda_U n+2c_2,\qquad \lambda_U|n-l|-c'\le|a_n-a_l|.
\end{equation*}
Here $\lambda_U>0$ is determined by $(N,L)$. The offset $\hat w$ satisfies $|\hat w|\le\varrho_2/2$. On each
torus, $S^{(\hat w)}$ denotes the limit of \textup{(v-b)}. Its moduli are $\upsilon_w$, the rate constant, and
the floor $n_0$. The floor $n_0$ depends on $N_w=2L^{m_0}$ through
\begin{equation*}
\mathsf{Q}_w=e^{(E_{+}^{\sharp}+C_{\mathrm{low}})N_w^4},
\end{equation*}
where $N_w=2L^{m_0}$ is the number of sites a side of the torus $\mathbb{T}_{m_0}$ on which the moduli
of \textup{(v-b)} are read. Clause \textup{(v-c)(1)} is the exhaustion
statement. Clause \textup{(v-c)(2)} is the corollary of the uniqueness theorem. Clauses
\textup{(v-d)(i)}--\textup{(iii)} carry the bi-Lipschitz lemma and the change-of-units identity.

\emph{Frames and limit sets.} In the tuned frame, the renormalised reference coupling $\hat{x}$ corresponds to
the bare offset $\hat w$. We put
\begin{equation*}
\mathfrak{S}_m[\hat{x}]:=S^{(\hat w)}_{\mathbb{T}_m};
\end{equation*}
this is the full continuum limit, on $\mathbb{T}_m$, of clauses \textup{(ii-$\mathbb{T}$-a)} and \textup{(v)}.
The \emph{tuned family} is the family of lattice theories of clause \textup{(v)} indexed by the integer $n$
of $a_n$, whose member with index $n$ has bare inverse coupling $a_n(g)-\hat w$. For such an index $n$ and a
bare inverse coupling $x_0$, let $S^{n}_{\mathbb{T}_m}(x_0;\vec f\,)$ be the connected correlation
functions on $\mathbb{T}_m$, evaluated on tuples $\vec f$ of test functions, of the member with index $n$
of the tuned family, taken at bare inverse coupling $x_0$. Then
$\operatorname{Lim}_{\infty}[\hat{x}]$ is the set of all limit points of
\begin{equation*}
j\mapsto S^{n_j}_{\mathbb{T}_{m_j}}(a_{n_j}-\hat w;\,\cdot\,)
\end{equation*}
along sequences with $n_j\to\infty$ and $m_j\to\infty$, in any manner. In the first-passage frame,
$\operatorname{Lim}_{\infty}(g)$ is the set of joint diagonal limits of clause \textup{(ii-$\mathbb{R}^4$)}.

\emph{Test classes.} The class $\mathfrak{P}_n(\mathsf{d};\vartheta)$ consists of the $n$-tuples
$(f_1,\dots,f_n)$ of real-valued Lipschitz functions, of any signs, that satisfy three conditions: their
supports are pairwise disjoint; all the supports lie in one closed ball $\overline{B}(y,5\mathsf{d})$ for
some point $y$; and
$\operatorname{Lip}f_i\le\|f_i\|_\infty/(\vartheta\mathsf{d})$ for each $i$. The supports of a pair in
$\mathfrak{T}(\mathsf{d};\vartheta)$ lie in two closed balls of radius $\mathsf{d}$ at distance at most
$6\mathsf{d}$, and these lie in one closed ball of radius $5\mathsf{d}$. Since the pairs of
$\mathfrak{T}(\mathsf{d};\vartheta)$ are real, have disjoint supports and satisfy the same Lipschitz bound
(clause \textup{(iv)} of Theorem~\ref{3.1:thm-main}), this gives
\begin{equation*}
\mathfrak{T}(\mathsf{d};\vartheta)\subset\mathfrak{P}_2(\mathsf{d};\vartheta).
\end{equation*}

\emph{Dilation.} Let $\gamma_{\star}$ be a second reference coupling and $\mathsf{s}$ an integer. Let
$\sigma'$ be a functional on tuples with separated supports, expressed in $\gamma_{\star}$-reference units.
We put
\begin{equation*}
(D_{\mathsf{s}}\sigma')(\vec f\,):=\sigma'\bigl(\vec f(L^{\mathsf{s}}\,\cdot\,)\bigr),
\end{equation*}
read in $g$-reference units.
\end{definition}

\begin{proof}
There are three assertions to prove: the monotonicity of $g_-$, the monotonicity of
$m_{\mathbb{T}}(g_-(\cdot))$, and the geometric fact behind the inclusion of test classes.

\emph{$g_-$ is increasing.} On $]0,\infty[$ the map $t\mapsto t^{-2}$ is strictly decreasing. Hence
$t\mapsto t^{-2}+B_{\sharp}$ is strictly decreasing, and its values are positive wherever $g_-$ is defined.
The map $u\mapsto u^{-1/2}$ is strictly decreasing on $]0,\infty[$. Composing two decreasing maps gives an
increasing one, so $g_-$ is increasing.

\emph{$m_{\mathbb{T}}(g_-(\cdot))$ is non-increasing in the coupling.} Let $t\le t'$. By the first step,
$g_-(t)\le g_-(t')$, so $g_-(t')$ is the larger coupling of the two. Since $R(\cdot)$ is non-decreasing as the
coupling decreases, $R(g_-(t'))\le R(g_-(t))$. Since $\log_L$ is increasing ($L>1$), we get
$\log_L R(g_-(t'))\le\log_L R(g_-(t))$. Adding the constant $11+m'$ gives
\begin{equation*}
m_{\mathbb{T}}(g_-(t'))\le m_{\mathbb{T}}(g_-(t)).
\end{equation*}

\emph{Two balls of radius $\mathsf{d}$ at distance at most $6\mathsf{d}$ lie in one closed ball of radius
$5\mathsf{d}$.} Let $\overline{B}(y_1,\mathsf{d})$ and $\overline{B}(y_2,\mathsf{d})$ be the two balls.
By the triangle inequality the distance between the two sets is at least $|y_1-y_2|-2\mathsf{d}$. Since this
distance is at most $6\mathsf{d}$, we get $|y_1-y_2|-2\mathsf{d}\le6\mathsf{d}$, that is
$|y_1-y_2|\le 8\mathsf{d}$. Put $y=(y_1+y_2)/2$, so that $|y_i-y|\le4\mathsf{d}$ for $i=1,2$. For
$z\in\overline{B}(y_i,\mathsf{d})$ the triangle inequality gives
\begin{equation*}
|z-y|\le|z-y_i|+|y_i-y|\le\mathsf{d}+4\mathsf{d}=5\mathsf{d}.
\end{equation*}
Hence both balls lie in $\overline{B}(y,5\mathsf{d})$.
\end{proof}

\begin{definition}[The strong reference coupling]\label{9.98:def-strong-reference-coupling}
Let $L$, $\mathfrak{p}$, $\gamma$, $\gamma_H$, $\varrho$ and $\vartheta\in\,]0,1]$ be fixed as in the order of
choice of Definitions~\ref{2.3:not-stages-first}--\ref{2.3:not-stages-last}. Clause~(iv), which holds
uniformly in the volume through the volume-transfer theorem, provides, for the fixed
$\vartheta$, a coupling ceiling
$\gamma_W^{\mathbb{T}}(L,\mathfrak{p},\gamma,\vartheta;\varrho)\le\gamma_J$, given by
$\gamma_W^{\mathbb{T}}=\min\{\gamma_W,\gamma_{*}^{\dagger}\}$; it is the ceiling of the two-point
witness on every torus, Theorem~\ref{8.18:thm-torus-witness}, and clause~(iv) carries it unchanged.
The \emph{strong reference coupling} is
\begin{equation}\label{9.98:eq-strong-reference-coupling-1}
\gamma_{\star}:=\gamma_W^{\mathbb{T}}(L,\mathfrak{p},\gamma,\vartheta;\varrho)
=\min\{\gamma_W,\gamma_{*}^{\dagger}\}.
\end{equation}
Thus $\gamma_{\star}$ depends on $(L,\mathfrak{p},\gamma,\gamma_H,\varrho,\vartheta)$ only. It is fixed after
$\varrho$ and $\vartheta$ and before every weak coupling $g$ used in this chapter. It is a ceiling on a
coupling, chosen last, one-sided, and it is never assigned a value.

The statements applied with reference coupling $\gamma_{\star}$ in this chapter are exactly three. Each
carries its own ceiling on the coupling:
\begin{itemize}
\item[($\alpha$)] the volume-transfer theorem of this chapter and its corollary, together with the
correlation theorem on the small torus, with ceiling $g\le\gamma_J$;
\item[($\beta$)] clause~(iv), uniform in the volume, with ceiling
$g\le\gamma_W^{\mathbb{T}}$;
\item[($\gamma$)] clause~(0) of Theorem~\ref{3.1:thm-main}, with the ceiling $g\le\gamma_H$ of hypothesis
\textup{(H4)} of Definition~\ref{2.2:not-hypotheses}.
\end{itemize}
Each of ($\alpha$), ($\beta$), ($\gamma$) applies with reference coupling $\gamma_{\star}$.

Clause~(v) of Theorem~\ref{3.1:thm-main} is not applied at $\gamma_{\star}$: it is applied at reference
coupling $g$ on the one torus $\mathbb{T}_{m_{\natural}(g)}$, where $m_{\natural}(g)$ is a floor on the torus
exponent depending on $g$ alone. Accordingly, the ceiling $\gamma_U$ of clause~(v) does not enter
$\gamma_{\star}$. The one statement of this chapter that reads clause~(v) in the frame of the strong
reference coupling, the half of the reduction to one window of couplings that uses tuned charts, reads it
at the lowered reference coupling
\begin{equation*}
\min\{\gamma_W^{\mathbb{T}},\gamma_U\}=\min\{\gamma_{\star},\gamma_U\},
\end{equation*}
a ceiling on a coupling chosen last, at which the ceiling $g\le\gamma_U$ of clause~(v) holds; every other
statement of this chapter that is applied at the strong reference coupling is applied at $\gamma_{\star}$
itself.

The items of the corollary to clause~(v) on short distances that are stated in absolute form, that is, all
its items other than the relative two-point discrepancy at the reference scale, do not use clause~(iv)
uniformly in the volume; they use only the statements under ($\alpha$) and ($\gamma$). For those items
alone, the argument below would allow
$\gamma_{\star}$ to be raised to $\gamma_J$, and this would lengthen the prolongation, that is, raise the
lower bounds $\underline{s}(g)$ and $\underline{s}'(g)$ on the prolongation
$K(g_0,\gamma_{\star})-K(g_0,g)$ and on its version in the tuned frame. This option is not taken: the single value
\eqref{9.98:eq-strong-reference-coupling-1} is used wherever the statements ($\alpha$)--($\gamma$) are
applied at the strong reference coupling.
\end{definition}

\begin{proof}[Proof that \textup{(}$\alpha$\textup{)}--\textup{(}$\gamma$\textup{)} apply at $\gamma_{\star}$]
We check the three ceilings one at a time.

\emph{Item ($\alpha$).} By~\eqref{9.98:eq-strong-reference-coupling-1}, $\gamma_{\star}=\gamma_W^{\mathbb{T}}$.
The existence statement of clause~(iv) uniformly in the volume gives $\gamma_W^{\mathbb{T}}\le\gamma_J$.
Hence $\gamma_{\star}\le\gamma_J$. The ceiling of ($\alpha$) therefore holds at $g=\gamma_{\star}$, and the
volume-transfer theorem, its corollary and the correlation theorem on the small torus apply with reference
coupling $\gamma_{\star}$.

\emph{Item ($\gamma$).} Hypothesis \textup{(H7)} of Definition~\ref{2.2:not-hypotheses} gives
$\gamma_J\le\gamma_H$. Combining this with item ($\alpha$),
\begin{equation*}
\gamma_{\star}\le\gamma_J\le\gamma_H .
\end{equation*}
So the ceiling of \textup{(H4)} holds at $g=\gamma_{\star}$, and clause~(0) applies with reference coupling
$\gamma_{\star}$. In particular, the following are available at reference coupling $\gamma_{\star}$:
\begin{itemize}
\item the first-passage depth $K(\cdot,\gamma_{\star})$;
\item the two-sided flow bounds of clause~(0);
\item the corollary of clause~(0) on fixed-depth windows.
\end{itemize}

\emph{Item ($\beta$).} The display of clause~(iv) uniformly in the volume is stated for every
$g\le\gamma_W^{\mathbb{T}}$. By~\eqref{9.98:eq-strong-reference-coupling-1}, the inequality
$\gamma_{\star}\le\gamma_W^{\mathbb{T}}$ holds, with equality. Hence that display applies at reference
coupling $\gamma_{\star}$. In this chapter it is used at $\gamma_{\star}$ only for the two-point statements
that rest on clause~(iv).

\emph{Dependence and order.} The ceiling $\gamma_W^{\mathbb{T}}$ reads
$(L,\mathfrak{p},\gamma,\vartheta;\varrho)$. In the order of choice of
Definitions~\ref{2.3:not-stages-first}--\ref{2.3:not-stages-last}, it is chosen after $\gamma_J$, which is
itself chosen after $\gamma_H$ and $\varrho$. Hence $\gamma_{\star}$
reads only $(L,\mathfrak{p},\gamma,\gamma_H,\varrho,\vartheta)$. It is fixed after $(\varrho,\vartheta)$
and before any weak coupling $g$. It enters only through the one-sided conditions
$g\le\gamma_J$, $g\le\gamma_W^{\mathbb{T}}$ and $g\le\gamma_H$ at $g=\gamma_{\star}$, and no value of it is
used.
\end{proof}

\begin{definition}[Prolongation depths, torus floor and pure constants]
\label{9.98:def-prolongation-depths}
Let $\varrho$ and $\vartheta$ be fixed, and let $\gamma_{\star}$ be the strong reference coupling
of Definition~\ref{9.98:def-strong-reference-coupling}. The couplings, the first-passage level
$K(g_0,\cdot)$, the tuned and first-passage frames and the test classes are those of
Definition~\ref{9.98:not-couplings-limit-sets}; $B^{\sharp}$, $\lambda^{\sharp}$, $c_5$,
$g_{-}(\cdot)$ and $m_{\mathbb{T}}(\cdot)$ are those of Definition~\ref{1.2:def-flow-constants}.

\emph{Reference data.} We put
\begin{equation}\label{9.98:eq-prolongation-depths-1}
\begin{gathered}
X_{\star}:=\gamma_{\star}^{-2},\qquad
W_{\star}:=(X_{\star},\,X_{\star}+B^{\sharp}],\qquad
\overline{W}_{\star}:=[X_{\star},\,X_{\star}+B^{\sharp}],\\
m_{\star}:=m_{\mathbb{T}}(g_{-}(\gamma_{\star})),\qquad
\mathbb{T}_{\star}:=\mathbb{T}_{m_{\star}},\qquad
\mathbb{T}_{\star}':=\mathbb{T}_{m_{\star}+1}.
\end{gathered}
\end{equation}
Thus $W_{\star}$ is the first-passage window of $x_{K(g_0,\gamma_{\star})}$ and
$\overline{W}_{\star}$ is its closure. For $n\in\{2,3\}$ the exponent $m_{\star}$ coincides with
$\max\{m_{\mathbb{T}}(g_{-}(\gamma_{\star})),\log_L\ell_n\}$, where $\ell_n$ is the side of the
chessboard cells for $n$ sources in the volume-transfer theorem; so $\mathbb{T}_{\star}$ is the
smallest torus of the family at reference $\gamma_{\star}$, and no condition on $L$ enters here,
since $\ell_2\le 32L$ and $\ell_3\le 48L<L^{11}M$.

\emph{The depth floor at the reference coupling.} With $s_W$, $s_{\mathrm{vol}}^{\star}$ and $s_T$ the
depth floors of clause~(iv) of Theorem~\ref{3.1:thm-main},
\begin{equation}\label{9.98:eq-prolongation-depths-2}
s_{\star}:=\max\bigl\{s_W(\gamma_{\star}),\ s_{\mathrm{vol}}^{\star}(\gamma_{\star}),\
s_T(\gamma_{\star};\vartheta,\varrho)\bigr\}.
\end{equation}

\emph{Pure constants.} For $n\in\{2,3\}$, with $C_{VT}$, $C_{VT,n}$, $C_T$ and $Q_{\star,n}$ the
constants of the volume-transfer theorem and of its corollary at $n=2$, and
$\mathfrak{r}\in\{1,2\}$ a factor fixed there,
\begin{equation}\label{9.98:eq-prolongation-depths-3}
\begin{aligned}
C_T^{\star}\equiv C^{\star}_{T,2}&:=C_T(\gamma_{\star})
=160^4\Bigl(\frac{40M}{\vartheta}\Bigr)^{2}C_{VT}(\gamma_{\star}),\\
C^{\star}_{T,3}&:=C_{T,3}(\gamma_{\star})
:=160^4\Bigl(\frac{40M}{\vartheta}\Bigr)^{3}C_{VT,3}(\gamma_{\star}),\\
C^{\star}_{VT,n}&:=C_{VT,n}(\gamma_{\star})
=2\log 2\,\bigl(2Q_{\star,n}(\gamma_{\star})+2\bigr)^{n}
\Bigl(\frac{n\mathfrak{r}}{\varrho}\Bigr)^{n}.
\end{aligned}
\end{equation}
The constants $C^{\star\mathrm{cvx}}_{T,n}$, $C^{\star\mathrm{cvx}}_{VT,n}$ are given by the same
expressions with $Q_{\star,n}(\gamma_{\star})$ replaced by
\begin{equation}\label{9.98:eq-prolongation-depths-4}
Q^{\mathrm{cvx}}_{\star,n}:=\exp\bigl[(1+\mathsf{c}_L)(\mathsf{B}+\bar{\delta})\ell_n^4
+2nC_{\mathrm{face}}(1)\bigr],\qquad
\mathsf{c}_L:=\Bigl(\frac{2L-1}{L-1}\Bigr)^{4}.
\end{equation}
All these constants are \emph{pure}: $C_{VT}$ and $C_{VT,n}$ read the coupling only through
$\mathsf{B}(g)=E_{+}^{\sharp}+C_{\mathrm{low}}$, $\ell_n$, $C_{\mathrm{face}}$ and $\bar{\delta}$, so
that, evaluated at the fixed coupling $\gamma_{\star}$, they are numbers depending on
$(L,\mathfrak{p},\gamma,\gamma_H,\varrho,\vartheta)$ and on nothing of $g$, $g_0$, $K$, $m$, the
tuned depth, $\mathsf{d}$ or position.

\emph{Two integers.} With $\lambda_U=\lambda_U(N,L)$ and $c'$ the spacing constants of
clause~(v-a) of Theorem~\ref{3.1:thm-main}, and $c_2$ the constant with $c_5=2c_2$ of
Definition~\ref{1.2:def-flow-constants},
\begin{equation}\label{9.98:eq-prolongation-depths-5}
j_U:=1+\Bigl\lceil\frac{\varrho_2+c'}{\lambda_U}\Bigr\rceil,\qquad
n_{\natural}:=\Bigl\lceil\frac{B^{\sharp}+2c_2+\varrho_2/2}{\lambda_U}\Bigr\rceil .
\end{equation}
The integer $j_U$, the renaming integer, bounds the renaming shifts $j$, $|j|\le j_U-1$, that
enter \eqref{9.98:eq-prolongation-depths-d} below; it is pure: it reads $\varrho$, $c'$ and
$\lambda_U=\lambda_U(N,L)$ only. The integer
$n_{\natural}$ is a floor on the tuned depth $n$, immaterial in every limit; for $n\ge n_{\natural}$
one has $a_n(g)-\varrho_2/2\ge g^{-2}+B^{\sharp}$.

\emph{Prolongation depths.} For $g\le\gamma_{\star}$ write $y:=g^{-2}$, so $y\ge X_{\star}$, and put
\begin{equation}\label{9.98:eq-prolongation-depths-a}
\begin{aligned}
\underline{s}(g)&:=\Bigl\lfloor\frac{y-X_{\star}-B^{\sharp}}{B^{\sharp}}\Bigr\rfloor\vee 0,
&\qquad
\overline{s}(g)&:=\Bigl\lceil\frac{y-X_{\star}+B^{\sharp}+c_5}{\lambda^{\sharp}}\Bigr\rceil,\\
\underline{s}'(g)&:=\bigl(\underline{s}(g)-j_U+1\bigr)\vee 0,
&\qquad
\overline{s}'(g)&:=\overline{s}(g)+j_U-1 .
\end{aligned}
\end{equation}
Both $\underline{s}$ and $\overline{s}$ are, up to integer parts, linear in $g^{-2}$; the unprimed
depths bound the first-passage prolongation $\mathsf{s}$ and the primed ones the tuned-frame
prolongation $\mathsf{s}'$ (see the classification below).

\emph{The coupling ceiling $\gamma_{\natural}$.} The coupling $\gamma_{\natural}=\gamma_{\natural}(\vartheta)$ is
defined by
\begin{equation}\label{9.98:eq-prolongation-depths-b}
\gamma_{\natural}^{-2}:=X_{\star}+B^{\sharp}\,(s_{\star}+j_U+1).
\end{equation}
Since $B^{\sharp}>0$ and $s_{\star}+j_U+1>0$, one has $\gamma_{\natural}^{-2}>X_{\star}$, that is
$\gamma_{\natural}<\gamma_{\star}$; so the depths \eqref{9.98:eq-prolongation-depths-a} are defined for
$g\le\gamma_{\natural}$, and then $\underline{s}(g)\ge s_{\star}+j_U$ and
$\underline{s}'(g)\ge s_{\star}+1$ (see the proof below).

\emph{The torus floor.} For $g\le\gamma_{\star}$,
\begin{equation}\label{9.98:eq-prolongation-depths-c}
m_{\natural}(g):=j_U+\max\bigl\{m_{\mathbb{T}}(g_{-}(g)),\ m_{\star}+\overline{s}(g)\bigr\}.
\end{equation}
It is a floor on the torus exponent $m$, depends on $g$ alone, and is of the same kind as the
floor of (H5): the torus contains the $\gamma_{\star}$-reference torus, with a margin of $j_U$
levels for the renaming between the two frames.

\emph{The window for the running coupling at a pair's level.} Let
$s_{\gamma_{\star}}(\cdot)$ be the depth function of clause~(iv) at reference $\gamma_{\star}$; for
$\mathsf{d}>0$ let $\mathsf{U}(g,\mathsf{d})$ be the maximum, over the prolongations $t$ and the
renaming shifts $j$ admitted in the display, of $s_{\gamma_{\star}}(L^{-t-j}\mathsf{d})-t$; and let
$\bar{x}^{\natural}(\mathsf{d})$, $\underline{x}^{\natural}$ be as follows:
\begin{equation}\label{9.98:eq-prolongation-depths-d}
\begin{aligned}
s_{\gamma_{\star}}(\delta)&:=\min\bigl\{t\ge 0:\
C_W\bigl(X_{\star}+B^{\sharp}(t+1)\bigr)L^{-t}\le\delta\bigr\},\\
\mathsf{U}(g,\mathsf{d})&:=\max\bigl\{s_{\gamma_{\star}}(L^{-t-j}\mathsf{d})-t:\
t\in[\underline{s}(g),\overline{s}(g)]\cap\mathbb{Z},\ |j|\le j_U-1\bigr\},\\
\bar{x}^{\natural}(\mathsf{d})&:=g^{-2}+B^{\sharp}\bigl(\mathsf{U}(g,\mathsf{d})+1\bigr)+c_5,
\qquad
\underline{x}^{\natural}:=\bigl(g^{-2}-B^{\sharp}j_U-c_5\bigr)_{+}.
\end{aligned}
\end{equation}
$\mathsf{U}(g,\mathsf{d})$ depends on $g$ and $\mathsf{d}$ only, and
$[\underline{x}^{\natural},\bar{x}^{\natural}(\mathsf{d})]$ is a window, depending on $g$ and
$\mathsf{d}$ only, for the running inverse coupling at the level of a pair at separation
$\mathsf{d}$, in both frames.

\emph{Classification of the conditions.} Every condition in which these quantities occur is
one-sided, of one of the following kinds, and none is a cap in the sense of
Definition~\ref{2.3:def-cap}.
\begin{enumerate}
\item Ceilings on couplings, chosen last: $\gamma_{\star}$, $\gamma_{\natural}$, $\gamma_U$ (the
uniqueness theorem's) and $\gamma_U'$ (its corollary's).
\item Ceilings on separations: $\mathsf{d}\le L^{\underline{s}'(g)}/40$ in the tuned frame,
$\mathsf{d}\le L^{\underline{s}(g)}/40$ in the first-passage frame, and $\mathsf{d}\le 1/40$ for the
two-point law on tori with $m\ge m_{\natural}(g)$ and on infinite-volume limit points.
\item Floors: $m\ge m_{\natural}(g)$ on the torus exponent; $n\ge n_{\natural}$ and
$L^{-n}\le\mathsf{d}$ on the tuned depth $n$ of the member (here $n$ is not the number of test
functions); $K-u(\mathsf{d};g_0,g)\ge k_W$ on the cutoff, where
$u(\mathsf{d};g_0,g):=s_{\gamma_{\star}}(L^{-\mathsf{s}}\mathsf{d})-\mathsf{s}$ with
$\mathsf{s}=K(g_0,\gamma_{\star})-K(g_0,g)$ the prolongation; and $\varrho\ge 6(B^{\sharp}+1)$,
the floor on $\varrho$ of the corollary to the uniqueness theorem.
\end{enumerate}
The only two-sided conditions are the window $[\mathsf{d},6\mathsf{d}]$ of clause~(iv), a condition
on the test functions on the conclusion side, and the parameter domain
$|\hat{w}|\le\varrho_2/2$ of clause~(v). The two-sided ranges
$\mathsf{s}\in[\underline{s}(g),\overline{s}(g)]$ and
$\mathsf{s}'\in[\underline{s}'(g),\overline{s}'(g)]$ are consequences of the two-sided flow bounds
of clause~(0) of Theorem~\ref{3.1:thm-main}, not hypotheses.
\end{definition}

\begin{proof}
We verify the elementary assertions made in Definition~\ref{9.98:def-prolongation-depths}.

\emph{The inequality $48L<L^{11}M$.} The blocking factor $L$ is an odd integer larger than one,
so $L\ge 3$ and $L^{10}\ge 3^{10}>48$; hence $48L<L^{11}\le L^{11}M$, since $M=L^{m'}\ge 1$.
With $\ell_2\le 32L$ and $\ell_3\le 48L$ this gives, for $n\in\{2,3\}$,
\begin{equation*}
\log_L\ell_n<11+m',
\end{equation*}
and $11+m'$ is the part of $m_{\mathbb{T}}(g')=11+m'+\log_L R(g')$ not depending on $g'$; no
lower bound on $L$ beyond $L\ge 3$ has been used.

\emph{The implication of $g\le\gamma_{\natural}$.} Let $g\le\gamma_{\natural}$. Then
$g^{-2}\ge\gamma_{\natural}^{-2}$, and by \eqref{9.98:eq-prolongation-depths-b}
\begin{equation*}
\frac{g^{-2}-X_{\star}-B^{\sharp}}{B^{\sharp}}
\ \ge\ \frac{B^{\sharp}(s_{\star}+j_U+1)-B^{\sharp}}{B^{\sharp}}
\ =\ s_{\star}+j_U ,
\end{equation*}
where the division preserves the inequality because $B^{\sharp}>0$. The right-hand side is an
integer, so the integer part of the left-hand side is at least $s_{\star}+j_U$, which is
non-negative; by \eqref{9.98:eq-prolongation-depths-a}, $\underline{s}(g)\ge s_{\star}+j_U$.
Consequently $\underline{s}(g)-j_U+1\ge s_{\star}+1\ge 1$, so the maximum with $0$ in the definition
of $\underline{s}'(g)$ is not active and $\underline{s}'(g)\ge s_{\star}+1$. In particular
$g\le\gamma_{\natural}$ implies $\underline{s}(g)\ge s_{\star}+j_U$ and $\underline{s}'(g)\ge s_{\star}+1$.

\emph{The second form of $m_{\natural}(g)$.} Adding the integer $j_U$ to both members of the
maximum in \eqref{9.98:eq-prolongation-depths-c}, and using
$\overline{s}(g)+j_U=\overline{s}'(g)+1$ from \eqref{9.98:eq-prolongation-depths-a}, we obtain
\begin{equation*}
m_{\natural}(g)=\max\bigl\{m_{\mathbb{T}}(g_{-}(g))+j_U,\ m_{\star}+1+\overline{s}'(g)\bigr\}.
\end{equation*}
In particular $m_{\natural}(g)\ge m_{\star}+1+\overline{s}'(g)\ge m_{\star}+1$, so every torus
$\mathbb{T}_m$ with $m\ge m_{\natural}(g)$ has exponent exceeding that of the reference torus
$\mathbb{T}_{\star}$ by at least $1+\overline{s}'(g)$. The right-hand side involves only $g$ and the
fixed data $\gamma_{\star}$, $j_U$, so $m_{\natural}(g)$ depends on $g$ alone.

\emph{The quantity $\mathsf{U}(g,\mathsf{d})$.} The maximum in
\eqref{9.98:eq-prolongation-depths-d} runs over the pairs $(t,j)$ with $t$ an integer in
$[\underline{s}(g),\overline{s}(g)]$ and $j$ an integer with $|j|\le j_U-1$; this index set is
finite, and it is determined by $g$ through $\underline{s}(g)$, $\overline{s}(g)$ and by the pure
integer $j_U$. The function $s_{\gamma_{\star}}$ reads only $X_{\star}$, $B^{\sharp}$, $C_W$ and $L$.
Hence $\mathsf{U}(g,\mathsf{d})$, and with it $\bar{x}^{\natural}(\mathsf{d})$, depend on $g$ and
$\mathsf{d}$ only, and $\underline{x}^{\natural}$ on $g$ only; none of them reads $g_0$, $K$ or $m$.

\emph{Purity of the constants.} In \eqref{9.98:eq-prolongation-depths-3} and
\eqref{9.98:eq-prolongation-depths-4} the coupling enters only through the arguments of
$C_{VT}$, $C_{VT,3}$ and $Q_{\star,n}$, which read it only through $\mathsf{B}$, $\ell_n$,
$C_{\mathrm{face}}$ and $\bar{\delta}$; all of these are evaluated at the fixed coupling
$\gamma_{\star}$. The remaining factors $160^4$, $(40M/\vartheta)^n$, $2\log 2$, $n\mathfrak{r}/\varrho$
and $\mathsf{c}_L$ read $L$, $\mathfrak{p}$, $\vartheta$, $\varrho$ and $n\in\{2,3\}$ only. The same
holds for $j_U$, which reads $\varrho$, $c'$ and $\lambda_U$.
\end{proof}

\begin{lemma}[Bounded offset between tuned and first-passage depths]\label{9.98:lem-tuned-depth-offset}
Let $g$ be a reference coupling in the window of clause~(v) of Theorem~\ref{3.1:thm-main}, with
$g\le\gamma_U'$ and $\varrho\ge 6(B^\sharp+1)$, and put $X:=g^{-2}$. Let $\lambda_U$, $c'$, $c_2$ be the
constants of clause~(v-a), and let $j_U=1+\lceil(\varrho_2+c')/\lambda_U\rceil$ be the renaming
integer of Definition~\ref{9.98:def-prolongation-depths}. Let $n\ge n_{\natural}$, the floor on the
tuned depth fixed in Definition~\ref{9.98:def-prolongation-depths}, let $|\hat w|\le\varrho_2/2$, and put
\begin{equation*}
x_0:=a_n(g)-\hat w ,
\end{equation*}
the bare inverse coupling at which the depth-$n$ member of the family of clause~(v) on
$\mathbb{T}_m$ is the Wilson measure (Definition~\ref{1.1:def-wilson-action}) on the lattice with
$2L^{m+n}$ sites a side. Then $g_0:=x_0^{-1/2}\le g_-(g)$, so that the first-passage depth
$K_0:=K(g_0,g)$ of clause~(0) is defined, and
\begin{equation}\label{9.98:eq-tuned-depth-offset-a}
|K_0-n|\le j_U-1 .
\end{equation}
\end{lemma}

Two labellings of the same measures occur here. In the \emph{tuned frame} a member is labelled by
its depth $n$ and has bare inverse coupling $x_0=a_n(g)-\hat w$. In the \emph{first-passage frame} a
member is labelled by the first-passage depth $K(g_0,g)$ of its bare coupling
(Notation~\ref{9.98:not-couplings-limit-sets}). The volume-transfer theorem at the strong reference
coupling $\gamma_{\star}$ (with $X_{\star}=\gamma_{\star}^{-2}$ as in
Definition~\ref{9.98:def-prolongation-depths}) is applied to the tuned member along the
first-passage run of clause~(0) started from $x_0$. The depth of that run at reference $g$ is $K_0$,
and \eqref{9.98:eq-tuned-depth-offset-a} compares it with the tuned depth $n$.

\begin{proof}
\emph{The datum is admissible.} Since $|\hat w|\le\varrho_2/2$, we have
$x_0=a_n(g)-\hat w\ge a_n(g)-\varrho_2/2$. Clause~(v-a) gives the lower bound
$a_n(g)\ge X+\lambda_U n-2c_2$. The floor $n\ge n_{\natural}$ gives
$\lambda_U n-2c_2-\varrho_2/2\ge B^\sharp$. Together,
\begin{equation}\label{9.98:eq-tuned-depth-offset-1}
x_0\ \ge\ a_n(g)-\frac{\varrho_2}{2}\ \ge\ X+\lambda_U n-2c_2-\frac{\varrho_2}{2}\ \ge\ X+B^\sharp .
\end{equation}
Since $g_-(g)=(g^{-2}+B^\sharp)^{-1/2}=(X+B^\sharp)^{-1/2}$ and $t\mapsto t^{-1/2}$ is decreasing
on $(0,\infty)$, \eqref{9.98:eq-tuned-depth-offset-1} gives
\begin{equation*}
g_0=x_0^{-1/2}\le (X+B^\sharp)^{-1/2}=g_-(g).
\end{equation*}
Thus $g_0$ is an admissible first-passage datum at reference $g$ in the sense of clause~(0). Its run
stops at depth $K_0=K(g_0,g)$.

\emph{Location of $x_0$.} The corollary of the uniqueness theorem (clause~(v-c)(2)), whose
hypotheses $\varrho\ge 6(B^\sharp+1)$ and $g\le\gamma_U'$ are assumed here, gives for bare data of
first passage at depth $K$ at reference $g$ the bound $|x-a_K(g)|\le\varrho_2/2$, where $x$ is the
bare inverse coupling. We apply it to
$x=x_0$, which is a first-passage datum at depth $K=K_0$ by the previous step. This gives
\begin{equation}\label{9.98:eq-tuned-depth-offset-2}
|x_0-a_{K_0}(g)|\le\frac{\varrho_2}{2}.
\end{equation}

\emph{Comparison of the two tuned values.} We have $a_n(g)-x_0=\hat w$. The triangle inequality and
\eqref{9.98:eq-tuned-depth-offset-2} give
\begin{equation*}
|a_n(g)-a_{K_0}(g)|\ \le\ |\hat w|+|x_0-a_{K_0}(g)|\ \le\ \frac{\varrho_2}{2}+\frac{\varrho_2}{2}
\ =\ \varrho_2 .
\end{equation*}

\emph{Spacing.} Clause~(v-a) gives, for the two depths $n$ and $K_0$, the spacing inequality
\begin{equation*}
\lambda_U|n-K_0|-c'\le|a_n(g)-a_{K_0}(g)| .
\end{equation*}
Combined with the previous display, it yields $\lambda_U|n-K_0|\le\varrho_2+c'$, that is,
\begin{equation*}
|n-K_0|\ \le\ \frac{\varrho_2+c'}{\lambda_U}\ \le\
\Bigl\lceil\frac{\varrho_2+c'}{\lambda_U}\Bigr\rceil\ =\ j_U-1 ,
\end{equation*}
which is \eqref{9.98:eq-tuned-depth-offset-a}.
\end{proof}

The integer $j_U-1$ enters the tuned frame twice. Let $g_0=x_0^{-1/2}$ be as in
Lemma~\ref{9.98:lem-tuned-depth-offset}. Part~(a) of the lemma comparing first-passage depths at
the two reference couplings $g$ and $\gamma_{\star}$, applied with this $g_0$, gives
\begin{equation*}
\mathsf{s}:=K(g_0,\gamma_{\star})-K_0\in[\underline{s}(g),\overline{s}(g)],
\end{equation*}
with $\underline{s}(g)$ and $\overline{s}(g)$ as in \eqref{9.98:eq-prolongation-depths-a}. The
prolongation beyond the tuned depth is
\begin{equation*}
\mathsf{s}':=K(g_0,\gamma_{\star})-n=\mathsf{s}+(K_0-n).
\end{equation*}
The bounds $\underline{s}'(g)$ and $\overline{s}'(g)$ of \eqref{9.98:eq-prolongation-depths-a} are
$\underline{s}(g)$ and $\overline{s}(g)$ shifted by the integer $j_U-1$, downward and upward
respectively. Hence \eqref{9.98:eq-tuned-depth-offset-a} gives
$\mathsf{s}'\in[\underline{s}'(g),\overline{s}'(g)]$.

The number $\mathsf{s}'$ fixes two things in the $\gamma_{\star}$ frame. The first is the exponent
$m-\mathsf{s}'$ of the torus, which yields a floor on $m$. The second is the scale
$L^{-\mathsf{s}'}\mathsf{d}$ of the tuple of test functions. This scale yields a ceiling on
$\mathsf{d}$ and the gain factor: $L^{-2\mathsf{s}'}$ for the two-point functions and
$L^{-\mathsf{s}'}$ for the three-point functions.

Accordingly, in the tuned frame the quantities
\begin{equation*}
\underline{s}(g),\quad \overline{s}(g),\quad X_{\star}+B^\sharp(s_{\star}+2),\quad
m_{\star}+1+\overline{s}(g)
\end{equation*}
of the comparison at reference $\gamma_{\star}$ are replaced, respectively, by
\begin{equation*}
\underline{s}'(g),\quad \overline{s}'(g),\quad X_{\star}+B^\sharp(s_{\star}+j_U+1),\quad
m_{\natural}(g),
\end{equation*}
where $m_{\natural}(g)$ is the floor of Definition~\ref{9.98:def-prolongation-depths}. Each
replacement is a shift by the integer $j_U-1$. The exponential factor
$e^{-(2\ln L/B^\sharp)g^{-2}}$ is unchanged.

In the statements read in the first-passage labelling $K=K(g_0,g)$ --- clause~(e) of the
short-distance corollary to clause~(v), part~(a) of the theorem on the $L$-adic orbits of the limit
sets, and the corollary carrying clause~(iv) to infinite-volume limit points --- $K_0$ is the depth
itself, the shift is $0$, and $\underline{s}(g)$, $\overline{s}(g)$, $X_{\star}+B^\sharp(s_{\star}+2)$ and
$m_{\star}+1+\overline{s}(g)$ are used without change.

\begin{remark}[The extra members of the torus and depth floors]\label{9.98:rem-floor-members}
Let $\gamma_{\star}$ be the strong reference coupling of Definition~\ref{9.98:def-strong-reference-coupling}. Let
$m_{\star}$, $s_{\star}$ and the pure constants be as in
Definition~\ref{9.98:def-prolongation-depths}, and let $m_{\natural}(g)$ be the floor on the torus
exponent in \eqref{9.98:eq-prolongation-depths-c}. The floor $m_{\natural}(g)$ contains the member
$m_{\mathbb{T}}(g_{-}(g))$ of its maximum and the summand $j_U$, and $s_{\star}$ contains the member
$s_{\mathrm{vol}}^{\star}(\gamma_{\star})$. The following hold.
\begin{itemize}
\item[(a)] Let $g\le\gamma_{\star}$ be a coupling in the window of clause (v) of
Theorem~\ref{3.1:thm-main}; let $n\ge n_{\natural}$ and $\hat{w}$ with
$|\hat{w}|\le\varrho_2/2$ be as in Lemma~\ref{9.98:lem-tuned-depth-offset}; and let
$m\ge m_{\natural}(g)$. The one torus
$\mathbb{T}_{m_{\natural}(g)}$ to which clause (ii-$\mathbb{T}$-a) and clause (v) of
Theorem~\ref{3.1:thm-main} are applied satisfies hypothesis (H5) of
Definition~\ref{2.2:not-hypotheses} for that torus. Let $x_0$ and $K_0=K(x_0,g)$ be as in
Lemma~\ref{9.98:lem-tuned-depth-offset}; write $K'$ for the first-passage depth of the same run at
reference $\gamma_{\star}$, $\mathsf{s}:=K'-K_0$ for its prolongation and $\mathsf{s}':=K'-n$. The
member of the tuned family, taken along the native first-passage run of
Lemma~\ref{9.98:lem-tuned-depth-offset} at reference $g$, is a measure on the torus
$\mathbb{T}_{m+n-K_0}$. Its member at reference $\gamma_{\star}$ is a measure on the torus
$\mathbb{T}_{m-\mathsf{s}'}$. These exponents satisfy
\begin{equation}\label{9.98:eq-floor-members-a}
\begin{aligned}
m+n-K_0 &\ge m-j_U+1 \ge m_{\mathbb{T}}(g_{-}(g))+1,\\
m-\mathsf{s}' &\ge m-\overline{s}(g)-j_U+1 \ge m_{\star}+1 \ge m_{\star}.
\end{aligned}
\end{equation}
The floor $m+n-K_0\ge m_{\mathbb{T}}(g_{-}(g))+1$ is the floor on the torus exponent under which
the passage to $\mathbb{R}^4$ of the infinite-volume theorem is applied to the tuned family.
Since $m_{\mathbb{T}}(g_{-}(\cdot))$ is non-increasing in the coupling and $g\le\gamma_{\star}$,
\begin{equation*}
m_{\star}=m_{\mathbb{T}}(g_{-}(\gamma_{\star}))\le m_{\mathbb{T}}(g_{-}(g)),
\end{equation*}
so the member $m_{\mathbb{T}}(g_{-}(g))$ is at least $m_{\star}$; for all sufficiently small $g$,
however, the member $m_{\star}+\overline{s}(g)\asymp g^{-2}/\lambda^{\sharp}$ of the maximum
exceeds $m_{\mathbb{T}}(g_{-}(g))$, so the added member does not increase $m_{\natural}(g)$ there;
in every case $m_{\natural}(g)$ depends on $g$ only.
\item[(b)] The depth floor $s_{\star}$ of Definition~\ref{9.98:def-prolongation-depths} is
\begin{equation}\label{9.98:eq-floor-members-b}
s_{\star}=\max\bigl\{s_W(\gamma_{\star}),\,s_{\mathrm{vol}}^{\star}(\gamma_{\star}),
\,s_T(\gamma_{\star};\vartheta,\varrho)\bigr\}.
\end{equation}

Clause (iv) uniformly in the volume is available under two coupling ceilings: the coupling
$\gamma_W^{\mathbb{T}}$, under which its depth condition is
$\max\{s_W,s_{\mathrm{vol}}^{\star},s_T\}\le s(\mathsf{d})$, and the coupling $\gamma_{\star}^{V}$ of
Definition~\ref{9.98:def-strong-reference-coupling}. With $\gamma_{\star}^{V}$ in place of
$\gamma_{\star}$ one has
$s_{\mathrm{vol}}^{\star}(\gamma_{\star}^{V})=0$, and the floor is
$\max\{s_W(\gamma_{\star}^{V}),\,s_T(\gamma_{\star}^{V};\vartheta,\varrho)\}$; the requirements of
Definition~\ref{9.98:def-strong-reference-coupling} on the strong reference coupling are met by
this choice as well (see there). Under either choice the depth floor is a pure constant.
\end{itemize}
\end{remark}

\begin{proof}
(a) By inspection of \eqref{9.98:eq-prolongation-depths-c},
$m_{\natural}(g)\ge m_{\mathbb{T}}(g_{-}(g))$, and (H5) holds for the torus
$\mathbb{T}_{m_{\natural}(g)}$.

For the first line of \eqref{9.98:eq-floor-members-a}, Lemma~\ref{9.98:lem-tuned-depth-offset}
gives $|K_0-n|\le j_U-1$. Therefore $n-K_0\ge -j_U+1$, which is the first inequality. The
hypothesis $m\ge m_{\natural}(g)$, read through the member $m_{\mathbb{T}}(g_{-}(g))$ and the
summand $j_U$, gives $m-j_U\ge m_{\mathbb{T}}(g_{-}(g))$. This is the second inequality.

For the second line, recall $\mathsf{s}'=K'-n$. Then
\begin{equation*}
\mathsf{s}'=K'-n=(K'-K_0)+(K_0-n).
\end{equation*}
Here $K'-K_0$ is the prolongation $\mathsf{s}$, which is at most $\overline{s}(g)$ by
Definition~\ref{9.98:def-prolongation-depths}. The offset
$K_0-n$ is at most $j_U-1$ by Lemma~\ref{9.98:lem-tuned-depth-offset}. Hence
$\mathsf{s}'\le\overline{s}(g)+j_U-1$, which is the first inequality of the second line.

The member $m_{\star}+\overline{s}(g)$ of the maximum in \eqref{9.98:eq-prolongation-depths-c},
increased by the summand $j_U$, gives $m-\overline{s}(g)-j_U\ge m_{\star}$. This is the second
inequality. The last inequality, $m_{\star}+1\ge m_{\star}$, is trivial.

(b) By Definition~\ref{9.98:def-strong-reference-coupling}, the coupling
$\gamma_{\star}=\gamma_W^{\mathbb{T}}$ is the coupling ceiling of clause (iv)
uniformly in the volume in the form whose depth condition is
$\max\{s_W,s_{\mathrm{vol}}^{\star},s_T\}\le s(\mathsf{d})$. Its depth floor is therefore
$\max\{s_W,s_{\mathrm{vol}}^{\star},s_T\}$; evaluated at the reference coupling $\gamma_{\star}$
this is \eqref{9.98:eq-floor-members-b}, and this is why $s_{\mathrm{vol}}^{\star}(\gamma_{\star})$
is a member of $s_{\star}$.

Under the ceiling $\gamma_{\star}^{V}$ the member $s_{\mathrm{vol}}^{\star}(\gamma_{\star}^{V})$
vanishes (Definition~\ref{9.98:def-strong-reference-coupling}), so the floor is the maximum of the
remaining two members. The three functions are
evaluated at the fixed coupling $\gamma_{\star}$, respectively
$\gamma_{\star}^{V}$, and at the fixed $\vartheta,\varrho$. Hence either floor is
pure. That the choice $\gamma_{\star}^{V}$ meets the requirements on the strong reference coupling
is settled in Definition~\ref{9.98:def-strong-reference-coupling}.
\end{proof}

\begin{theorem}[Short-distance uniqueness given the renormalised coupling]
\label{9.98:thm-short-distance-uniqueness}
Let $\gamma_{\star}$ be the strong reference coupling of
Definition~\ref{9.98:def-strong-reference-coupling}, and put $X_{\star}:=\gamma_{\star}^{-2}$.
Let $\underline{s}(g)$, $\overline{s}(g)$, $\underline{s}'(g)$, $\overline{s}'(g)$, $\gamma_{\natural}$,
$m_{\natural}(g)$ and $\bar{x}^{\natural}(\mathsf{d})$ be as in
\eqref{9.98:eq-prolongation-depths-a}, \eqref{9.98:eq-prolongation-depths-b},
\eqref{9.98:eq-prolongation-depths-c} and \eqref{9.98:eq-prolongation-depths-d}.
Let the pure constants $C^{\star}_{T,n}$ and $C^{\star\mathrm{cvx}}_{T,n}$, the cell sides
$\ell_n$, the reference torus exponent $m_{\star}$, the tori $\mathbb{T}_{\star}$,
$\mathbb{T}_{\star}'$ and the admissible values of the
tuned prolongation $\mathsf{s}'$ be those of Definition~\ref{9.98:def-prolongation-depths}, and
write $C^{\star}_{T}:=C^{\star}_{T,2}$. Let $j_U$ be the renaming integer of
Lemma~\ref{9.98:lem-tuned-depth-offset}, and let $\gamma_U'\le\gamma_U$ be the coupling threshold
of the corollary of the uniqueness theorem. The tuned family, its continuum limits
$\mathfrak{S}_m[\hat{x}]$ on $\mathbb{T}_m$, the limit sets
$\operatorname{Lim}_{\infty}[\hat{x}]$ and $\operatorname{Lim}_{\infty}(g)$, the connected
$n$-point functions $S_n[\sigma]$ of a limit point $\sigma$ and the test classes
$\mathfrak{P}_n(\mathsf{d};\vartheta)$ are those of
Notation~\ref{9.98:not-couplings-limit-sets}. Fix $\vartheta\in\,]0,1]$, assume the floor
$\varrho\ge6(B^{\sharp}+1)$ of the corollary of the uniqueness theorem, and let
$g\le\min\{\gamma_{\natural},\gamma_U'\}$. Lengths are measured in units of the reference
length of $g$. Assume:
\begin{itemize}
\item the correlation bounds of clause (iii) of Theorem~\ref{3.1:thm-main}, and the
volume-transfer theorem at the strong reference coupling $\gamma_{\star}$;
\item clauses (v-a), (v-b) and (v-c)(1) of the uniqueness theorem (clause (v) of
Theorem~\ref{3.1:thm-main}), applied on the one torus $\mathbb{T}_{m_{\natural}(g)}$; the corollary
of the uniqueness theorem; and clauses (v-d)(i) and (v-d)(ii) of Theorem~\ref{3.1:thm-main};
\item for clause (d) below, clause (iv) of Theorem~\ref{3.1:thm-main} uniformly in the volume.
\end{itemize}
Let $\hat{w}$ be real
with $|\hat{w}|\le\varrho_2/2$, and let $\hat{x}$ be the renormalised reference coupling that
corresponds to $\hat{w}$ by clauses (v-d)(i) and (v-d)(ii) of Theorem~\ref{3.1:thm-main}. Let
$n\in\{2,3\}$, let $0<\mathsf{d}\le L^{\underline{s}'(g)}/40$, and let
$\vec f=(f_1,\dots,f_n)\in\mathfrak{P}_n(\mathsf{d};\vartheta)$. Put
\begin{equation}\label{9.98:eq-short-distance-uniqueness-a}
\Pi:=\prod_{i=1}^{n}\|f_i\|_{\infty},\qquad
\tau_n:=C^{\star}_{T,n}\,L^{-(4-n)\underline{s}'(g)}\,\mathsf{d}^{4-n}\,\Pi ,
\end{equation}
so that
\begin{equation*}
\tau_2=C^{\star}_{T}\,L^{-2\underline{s}'(g)}\,\mathsf{d}^{2}\,\Pi,\qquad
\tau_3=C^{\star}_{T,3}\,L^{-\underline{s}'(g)}\,\mathsf{d}\,\Pi ,
\end{equation*}
and let the torus floor be
\begin{equation*}
m_{\natural}(g,n):=j_U+\max\Bigl\{m_{\mathbb{T}}(g_{-}(g)),\;
\max\{m_{\star},\lceil\log_L\ell_n\rceil\}+1+\overline{s}(g)\Bigr\}.
\end{equation*}
Comparison with \eqref{9.98:eq-prolongation-depths-c} gives $m_{\natural}(g,n)\ge m_{\natural}(g)$;
the reference torus in (b)--(e) below is $\mathbb{T}_{m_{\natural}(g)}$.
Then the following hold.
\begin{itemize}
\item[(a)] (Tori.) For all $m,m'\ge m_{\natural}(g,n)$,
\begin{equation*}
\bigl|\mathfrak{S}_m[\hat{x}](\vec f\,)-\mathfrak{S}_{m'}[\hat{x}](\vec f\,)\bigr|\le\tau_n .
\end{equation*}
\item[(b)] ($\mathbb{R}^4$, all limit points at one $\hat{x}$.) For all
\begin{equation*}
\sigma,\sigma'\in\operatorname{Lim}_{\infty}[\hat{x}]\cup
\{\mathfrak{S}_m[\hat{x}]:m\ge m_{\natural}(g,n)\}
\end{equation*}
one has
\begin{equation*}
\bigl|S_n[\sigma](\vec f\,)-S_n[\sigma'](\vec f\,)\bigr|\le 2\tau_n,\qquad
\bigl|S_n[\sigma](\vec f\,)-\mathfrak{S}_{m_{\natural}(g)}[\hat{x}](\vec f\,)\bigr|\le\tau_n .
\end{equation*}
Here the map $\hat{w}\mapsto\mathfrak{S}_{m_{\natural}(g)}[\hat{x}](\vec f\,)$, with $\hat{x}$
corresponding to $\hat{w}$ as above, is real-analytic on $]{-\varrho_2/2},\varrho_2/2[$ and has a
holomorphic extension to the strip $|\operatorname{Im}\hat{w}|<\upsilon_w(m_{\natural}(g))$,
whose width depends on $g$ only (clause (v-b)).
\item[(c)] (The limits commute at short distance.) The depth index is written $\nu$ here, the
letter $n$ being reserved for the number of test functions. For $\nu$ and $m$ in the range in
which the tuned family is defined (Notation~\ref{9.98:not-couplings-limit-sets}), write
$S^{\nu}_{\mathbb{T}_m}(a_{\nu}-\hat{w};\vec f\,)$ for the connected $n$-point function, at
$\vec f$, of the member of depth $\nu$ of the tuned family on $\mathbb{T}_m$, that is, of the
Wilson measure on the lattice with $2L^{m+\nu}$ sites a side at bare inverse coupling
$a_{\nu}-\hat{w}$ (Notation~\ref{9.98:not-couplings-limit-sets}). Every cluster value of the double
array $(\nu,m)\mapsto S^{\nu}_{\mathbb{T}_m}(a_{\nu}-\hat{w};\vec f\,)$ as $\nu,m\to\infty$
--- $\nu\to\infty$ then $m\to\infty$; $m\to\infty$ along any subsequence at fixed $\nu$, then
$\nu\to\infty$; or jointly in any manner --- lies in the one closed interval
\begin{equation*}
\bigl[\,\mathfrak{S}_{m_{\natural}(g)}[\hat{x}](\vec f\,)-\tau_n,\;
\mathfrak{S}_{m_{\natural}(g)}[\hat{x}](\vec f\,)+\tau_n\,\bigr]
\end{equation*}
of length $2\tau_n$.
\item[(d)] (Relative form, $n=2$.) Let moreover $\mathsf{d}\le1/40$. For
$(f_1,f_2)\in\mathfrak{T}(\mathsf{d};\vartheta)$ and $\sigma,\sigma'$ as in (b),
\begin{equation*}
\Bigl|\frac{S_2[\sigma](f_1,f_2)}{S_2[\sigma'](f_1,f_2)}-1\Bigr|
\le\frac{9\,C^{\star}_{T}\,L^{-2\underline{s}'(g)}\,\mathsf{d}^{2}\,
\bar{x}^{\natural}(\mathsf{d})^{2}}{b_0\,c_{\mathcal{K}}} .
\end{equation*}
Since
$\underline{s}'(g)\ge(g^{-2}-X_{\star})/B^{\sharp}-(j_U+1)$ by
\eqref{9.98:eq-prolongation-depths-a}, the right-hand side is at most
\begin{equation*}
\mathsf{A}_0\,\mathsf{d}^{2}\,\bar{x}^{\natural}(\mathsf{d})^{2}
\exp\bigl[-(2\ln L/B^{\sharp})\,g^{-2}\bigr],
\end{equation*}
where
\begin{equation}\label{9.98:eq-short-distance-uniqueness-b}
\mathsf{A}_0:=\frac{9\,C^{\star}_{T}\,
L^{2(j_U-1)+2(X_{\star}+2B^{\sharp})/B^{\sharp}}}{b_0\,c_{\mathcal{K}}}
\end{equation}
is a pure constant. Equivalently: if, for one member $\sigma$ of the set
in (b), the two-point function $S_2[\sigma]$ on the pairs of $\mathfrak{T}(\mathsf{d};\vartheta)$
has an asymptotic expansion in $1/\ln(1/\mathsf{d})$ to some order, then all members have
the same expansion to that order; the existence of such an expansion beyond tree level is not
part of this statement. Without the restriction $\mathsf{d}\le1/40$: if
$L^{\underline{s}'(g)}\ge40$, so that the reference scale $\mathsf{d}=1$ of $g$ lies in the range
$\mathsf{d}\le L^{\underline{s}'(g)}/40$, then for $(f_1,f_2)\in\mathfrak{P}_2(1;\vartheta)$ and
$\sigma,\sigma'$ as in (b), clause (b) at $\mathsf{d}=1$ and the same lower bound on
$\underline{s}'(g)$ give for the absolute discrepancy
\begin{align*}
\bigl|S_2[\sigma](f_1,f_2)-S_2[\sigma'](f_1,f_2)\bigr|
&\le 2\tau_2\\
&\le 2\,C^{\star}_{T}\,L^{2(j_U+1)+2X_{\star}/B^{\sharp}}
\exp\bigl[-(2\ln L/B^{\sharp})\,g^{-2}\bigr]\,\Pi ,
\end{align*}
which is smaller than every power of the running coupling.
\item[(e)] (First-passage frame.) In this clause let $0<\mathsf{d}\le L^{\underline{s}(g)}/40$,
$\vec f\in\mathfrak{P}_n(\mathsf{d};\vartheta)$ and $\Pi$ as in
\eqref{9.98:eq-short-distance-uniqueness-a}. Every $\sigma\in\operatorname{Lim}_{\infty}(g)$,
the set of infinite-volume limit points of clause (ii-$\mathbb{R}^4$) of
Theorem~\ref{3.1:thm-main}, satisfies
\begin{equation*}
\bigl|S_n[\sigma](\vec f\,)-\mathfrak{S}_{m_{\natural}(g)}[\hat{x}'](\vec f\,)\bigr|
\le C^{\star}_{T,n}\,L^{-(4-n)\underline{s}(g)}\,\mathsf{d}^{4-n}\,\Pi
\end{equation*}
for some renormalised coupling $\hat{x}'$, namely the one corresponding to a limit
$\hat{w}'\in[-\varrho_2/2,\varrho_2/2]$ of the offsets $a_{K_j}-(g_0^{(j)})^{-2}$ of the runs along
which $\sigma$ is obtained, $K_j$ and $g_0^{(j)}$ being their depths and bare couplings (the
corollary of the uniqueness theorem). Two such $\sigma,\sigma'$ whose runs' renormalised inverse
couplings at the reference scale converge to the same value obey
\begin{equation*}
\bigl|S_n[\sigma](\vec f\,)-S_n[\sigma'](\vec f\,)\bigr|
\le 2\,C^{\star}_{T,n}\,L^{-(4-n)\underline{s}(g)}\,\mathsf{d}^{4-n}\,\Pi
\end{equation*}
(clause (v-d)(ii) of Theorem~\ref{3.1:thm-main}).
\item[(f)] (Shapes.) Clauses (a)--(e) hold, with $C^{\star}_{T,n}$ replaced by
$2C^{\star\mathrm{cvx}}_{T,n}$ (equivalently, $\tau_n$ replaced by $2\tau^{\mathrm{cvx}}_n$, where
$\tau^{\mathrm{cvx}}_n$ is defined as $\tau_n$ in \eqref{9.98:eq-short-distance-uniqueness-a} with
$C^{\star\mathrm{cvx}}_{T,n}$ in place of $C^{\star}_{T,n}$), when $\mathbb{T}_m$, $\mathbb{T}_{m'}$
are replaced by rectangular tori $\mathfrak{B}$, $\mathfrak{B}'$ and $\mathfrak{S}_m[\hat{x}]$,
$\mathfrak{S}_{m'}[\hat{x}]$ by $\mathfrak{S}_{\mathfrak{B}}[\hat{x}]$,
$\mathfrak{S}_{\mathfrak{B}'}[\hat{x}]$, each of the tori having four side lengths $b$ that satisfy
\begin{equation*}
b\in 2\ell_n L^{\overline{s}'(g)}\mathbb{N},\qquad b\ge 2L^{m_{\natural}(g,n)}
\end{equation*}
(in the frame of $\gamma_{\star}$: even numbers of chessboard cells per side, each at least the
number for the reference torus $\mathbb{T}_{\star}$, for every admissible $\mathsf{s}'$). Thermal
boxes $\mathbb{R}^3\times S^1_{\beta}$, with $S^1_{\beta}$ the circle of circumference $\beta$,
$\beta$ of that form and the other three sides tending to infinity, are included. Here
$\mathfrak{S}_{\mathfrak{B}}[\hat{x}]$ denotes any cutoff-limit point on $\mathfrak{B}$ of the
tuned family at $\hat{w}$; such points exist along subsequences. Ba{\l}aban's construction enters
only on two cubic tori of the family, in the frame of $\gamma_{\star}$ the torus
$\mathbb{T}_{m_{\natural}(g)-\mathsf{s}'}$ and the next one up, both containing
$\mathbb{T}_{\star}'$; every other volume, $\mathfrak{B}$ included, is controlled by link
reflection positivity of its bare measure and the chessboard, transfer and convexity estimates
derived from it.
\end{itemize}
\end{theorem}

For $n=2$ and $(f_1,f_2)\in\mathfrak{T}(\mathsf{d};\vartheta)$ the map
$\hat{w}\mapsto\mathfrak{S}_{m_{\natural}(g)}[\hat{x}](f_1,f_2)$ of clause (b) is moreover
non-constant on every sub-interval, by the non-constancy theorem for the reference-torus family;
this is not used in clauses (a)--(f).

Theorem~\ref{9.98:thm-short-distance-uniqueness} asserts nothing about connected $n$-point
functions with $n\ge4$; nothing about the optimal power of $\mathsf{d}$; nothing at separations
$\mathsf{d}>L^{\underline{s}'(g)}/40$, that is, beyond the fixed length $\ell_{\star}/40$, up to
the pure factor $L^{j_U-1}$, where $\ell_{\star}$ is the reference length of $\gamma_{\star}$;
neither convergence of $\mathfrak{S}_m[\hat{x}]$ as $m\to\infty$, nor a relative form of (d) at
$n=3$, nor injectivity of the dependence on $\hat{x}$, nor anything at complex $\hat{x}$ on
$\mathbb{R}^4$; nor full convergence in the cutoff of the tuned family on a non-cubic torus
$\mathfrak{B}$, clause (f) speaking of cutoff-limit points on $\mathfrak{B}$ only. Nor does it
assert uniqueness of the infinite-volume limit, uniqueness of the vacuum, continuity,
analyticity or injectivity of the map from the renormalised coupling to the infinite-volume
theory, or a mass gap.

Each of clauses (a)--(f) is derived from the hypotheses listed in the statement in a separate
proposition later in this chapter.

\begin{lemma}[Run prolongation to the strong reference coupling]\label{9.98:lem-run-prolongation}
Let $\gamma_{\star}$ be the strong reference coupling of
Definition~\ref{9.98:def-strong-reference-coupling}, let $g\le\gamma_{\star}$ and
$g_0\in\,]0,g_-(g)]$, and put
\begin{equation*}
K:=K(g_0,g),\qquad K':=K(g_0,\gamma_{\star}),\qquad \mathsf{s}:=K'-K,
\end{equation*}
first-passage depths in the sense of Definition~\ref{1.2:def-flow-constants}. Then the
following hold.
\begin{enumerate}
\item[(a)] $K'$ is well defined, $g_0\le g_-(\gamma_{\star})$, $K'\ge K$, and
\begin{equation}\label{9.98:eq-run-prolongation-a}
\underline{s}(g)\le\mathsf{s}\le\overline{s}(g),
\end{equation}
where $\underline{s}(g)$ and $\overline{s}(g)$ are the prolongation depths
of~\eqref{9.98:eq-prolongation-depths-a}.
\item[(b)] For every $n\ge2$, every $m\ge\mathsf{s}+m_{\star}$, with $m_{\star}$ as in
Definition~\ref{9.98:def-prolongation-depths}, and every $n$-tuple
$\vec f=(f_1,\dots,f_n)$ of test functions on $\mathbb{T}_m$, written in the reference unit
of $g$,
\begin{equation}\label{9.98:eq-run-prolongation-b}
S^T_{n;K,m}[g_0](\vec f\,)
=S^T_{n;K',m-\mathsf{s}}[g_0]\bigl(\vec f\,(L^{\mathsf{s}}\,\cdot\,)\bigr).
\end{equation}
Here the bare coupling is written in square brackets, and $\vec f\,(L^{\mathsf{s}}\cdot)$ is
the tuple of functions $x'\mapsto f_i(L^{\mathsf{s}}x')$ on $\mathbb{T}_{m-\mathsf{s}}$. The
two members of~\eqref{9.98:eq-run-prolongation-b} are one lattice integral under two namings of
the reference length. The right member is the object of Theorem~\ref{3.1:thm-main} at reference
coupling $\gamma_{\star}$ on the torus $\mathbb{T}_{m-\mathsf{s}}$. It lives on the same bare
lattice, of $2L^{m+K}=2L^{(m-\mathsf{s})+K'}$ sites a side, at the same bare coupling $g_0$. Its
test functions are re-expressed in the reference unit of $\gamma_{\star}$, in which supports and
separations are $L^{-\mathsf{s}}$ times smaller, sup norms are unchanged, and Lipschitz
constants are $L^{\mathsf{s}}$ times larger. Consequently
$\vec f\in\mathfrak{P}_n(\mathsf{d};\vartheta)$ if and only if
$\vec f\,(L^{\mathsf{s}}\cdot)\in\mathfrak{P}_n(L^{-\mathsf{s}}\mathsf{d};\vartheta)$, and
likewise for $\mathfrak{T}$, the classes being those of
Notation~\ref{9.98:not-couplings-limit-sets}.
\item[(c)] Let $(g_0^{(j)})_j$ be a sequence in $]0,g_-(g)]$ with $g_0^{(j)}\downarrow0$, and
let $\mathsf{s}_j$ be the prolongation defined as above with $g_0^{(j)}$ in place of $g_0$.
Then $\mathsf{s}_j\in[\underline{s}(g),\overline{s}(g)]$ takes finitely many values; hence it
is constant along a subsequence.
\end{enumerate}
\end{lemma}

\begin{proof}
\emph{Proof of} (a). Put $y:=g^{-2}$ and $x_0:=g_0^{-2}$, and write
$X_{\star}=\gamma_{\star}^{-2}$ as in Definition~\ref{9.98:def-prolongation-depths}. Since
$g\le\gamma_{\star}$, we have $y\ge X_{\star}$.

The function $t\mapsto g_-(t)=(t^{-2}+B^{\sharp})^{-1/2}$ is increasing on $]0,\infty[$. Indeed,
$t\mapsto t^{-2}+B^{\sharp}$ is positive and decreasing, and $u\mapsto u^{-1/2}$ is decreasing on
$]0,\infty[$. Hence
\begin{equation*}
g_0\le g_-(g)\le g_-(\gamma_{\star}),
\end{equation*}
which is the assertion $g_0\le g_-(\gamma_{\star})$. The inequality $g_0\le g_-(t)$ is
equivalent to $x_0\ge t^{-2}+B^{\sharp}$. The two inequalities above therefore read
\begin{equation*}
x_0\ge y+B^{\sharp}\ge X_{\star}+B^{\sharp}.
\end{equation*}

By Definition~\ref{9.98:def-strong-reference-coupling}, $\gamma_{\star}\le\gamma_H$. Together
with $g_0\le g_-(\gamma_{\star})$, this lets clause~(0) of Theorem~\ref{3.1:thm-main} at
reference coupling $\gamma_{\star}$ apply to $g_0$. It gives two conclusions. First,
$K'=\min\{k:\ x_k\le X_{\star}+B^{\sharp}\}$ is a well-defined non-negative integer. Second,
the two-sided flow bounds~\eqref{9.98:eq-couplings-limit-sets-a} hold for
$0\le i\le j\le K'$ along the run $(x_k)_k$ of running inverse couplings started from $x_0$
(Definition~\ref{1.2:def-couplings}). This run does not depend on the torus.

Likewise $g\le\gamma_{\star}\le\gamma_H$ and $g_0\le g_-(g)$. So clause~(0) at reference
coupling $g$ applies, and $K=\min\{k:\ x_k\le y+B^{\sharp}\}$. Clause~(0) at reference $g$ and
clause~(0) at reference $\gamma_{\star}$ speak of the same run. The recursion
$x_{k+1}=x_k-\beta_{k+1}$ does not involve the reference coupling, which only decides the index
at which one stops.

\emph{$K'\ge K$.} For $k<K$, minimality of $K$ gives
$x_k>y+B^{\sharp}\ge X_{\star}+B^{\sharp}$. Hence no $k<K$ satisfies
$x_k\le X_{\star}+B^{\sharp}$, and therefore $K'\ge K$. In particular $\mathsf{s}\ge0$.

\emph{The first-passage windows.} By definition of $K$, $x_K\le y+B^{\sharp}$. We show that
$x_K>y$, distinguishing two cases.
\begin{itemize}
\item If $K\ge1$, minimality gives $x_{K-1}>y+B^{\sharp}$. The flow
bounds~\eqref{9.98:eq-couplings-limit-sets-a} with $i=K-1$, $j=K$ (so $j-i=1$; admissible since
$K\le K'$) give $x_{K-1}-x_K\le B^{\sharp}$. Hence $x_K\ge x_{K-1}-B^{\sharp}>y$.
\item If $K=0$: this cannot occur when $x_0>y+B^{\sharp}$. When $x_0=y+B^{\sharp}$, one has
$K=0$ and $x_K=y+B^{\sharp}>y$.
\end{itemize}
Thus $x_K\in\,]y,y+B^{\sharp}]$. The same argument, with $X_{\star}$ in place of $y$, $K'$ in
place of $K$, and the inequality $x_0\ge X_{\star}+B^{\sharp}$, gives
$x_{K'}\in\,]X_{\star},X_{\star}+B^{\sharp}]$.

Combine the upper end of the first window with the lower end of the second, and vice versa:
\begin{equation}\label{9.98:eq-run-prolongation-1}
y-X_{\star}-B^{\sharp}<x_K-x_{K'}<y+B^{\sharp}-X_{\star}.
\end{equation}
On the other hand, the flow bounds~\eqref{9.98:eq-couplings-limit-sets-a} with $i=K\le j=K'$,
so $j-i=\mathsf{s}$, give
\begin{equation}\label{9.98:eq-run-prolongation-2}
\lambda^{\sharp}\mathsf{s}-c_5\le x_K-x_{K'}\le B^{\sharp}\mathsf{s}.
\end{equation}

\emph{Lower bound.} From the right inequality in~\eqref{9.98:eq-run-prolongation-2} and the
left one in~\eqref{9.98:eq-run-prolongation-1},
\begin{equation*}
B^{\sharp}\mathsf{s}\ge x_K-x_{K'}>y-X_{\star}-B^{\sharp},
\end{equation*}
so $\mathsf{s}>q:=(y-X_{\star}-B^{\sharp})/B^{\sharp}$. We deduce
$\mathsf{s}\ge\lfloor q\rfloor$ in two cases.
\begin{itemize}
\item If $q$ is an integer $N$, then $\mathsf{s}>N$ gives $\mathsf{s}\ge N+1\ge N$.
\item Otherwise the integer $\mathsf{s}$ satisfies
$\mathsf{s}\ge\lceil q\rceil\ge\lfloor q\rfloor$.
\end{itemize}
Together with $\mathsf{s}\ge0$, this is $\mathsf{s}\ge\underline{s}(g)$, by the definition of
$\underline{s}(g)$ in~\eqref{9.98:eq-prolongation-depths-a}.

\emph{Upper bound.} From the left inequality in~\eqref{9.98:eq-run-prolongation-2} and the right
one in~\eqref{9.98:eq-run-prolongation-1},
\begin{equation*}
\lambda^{\sharp}\mathsf{s}-c_5\le x_K-x_{K'}<y+B^{\sharp}-X_{\star}.
\end{equation*}
Since $\lambda^{\sharp}>0$, this gives
\begin{equation*}
\mathsf{s}<\frac{y-X_{\star}+B^{\sharp}+c_5}{\lambda^{\sharp}}
\le\Bigl\lceil\frac{y-X_{\star}+B^{\sharp}+c_5}{\lambda^{\sharp}}\Bigr\rceil
=\overline{s}(g),
\end{equation*}
the last equality being the definition of $\overline{s}(g)$
in~\eqref{9.98:eq-prolongation-depths-a}. This proves~\eqref{9.98:eq-run-prolongation-a} and
completes~(a).

\emph{Proof of} (b). Put $\vec h:=\vec f\,(L^{\mathsf{s}}\cdot)$, so that
$\vec h\,(L^{-\mathsf{s}}\cdot)=\vec f$. The change-of-units identity of clause~(v) of
Theorem~\ref{3.1:thm-main} states the following. For every integer $s$ with $0\le s\le K$,
\begin{equation*}
S^T_{n;K,m}[g_0](\vec f\,)=S^T_{n;K-s,m+s}[g_0]\bigl(\vec f\,(L^{-s}\cdot)\bigr):
\end{equation*}
the same plaquettes, read with the reference length moved $s$ levels toward the cutoff. Apply
it with $s=\mathsf{s}$ (equivalently, its one-level form iterated $\mathsf{s}$ times) to the
right member of~\eqref{9.98:eq-run-prolongation-b}, at depth $K'$ on $\mathbb{T}_{m-\mathsf{s}}$;
this is admissible since $0\le\mathsf{s}=K'-K\le K'$ by~(a) and $K\ge0$.
Since $K'-\mathsf{s}=K$, this gives
\begin{align*}
S^T_{n;K',m-\mathsf{s}}[g_0](\vec h\,)
&=S^T_{n;K'-\mathsf{s},(m-\mathsf{s})+\mathsf{s}}[g_0]\bigl(\vec h\,(L^{-\mathsf{s}}\cdot)\bigr)\\
&=S^T_{n;K,m}[g_0](\vec f\,),
\end{align*}
which is~\eqref{9.98:eq-run-prolongation-b}.

This is an identity of one integral. In both members the measure is the Wilson measure at bare
coupling $g_0$ on one lattice. In the reference unit of $g$ that lattice has spacing $L^{-K}$
on $\mathbb{T}_m$, of side $2L^m$. In the reference unit of $\gamma_{\star}$ it has spacing
$L^{-K'}$ on $\mathbb{T}_{m-\mathsf{s}}$, of side $2L^{m-\mathsf{s}}$. In both descriptions it
has $2L^{m+K}=2L^{(m-\mathsf{s})+K'}$ sites a side. The unit of $\gamma_{\star}$ is therefore
$L^{K'-K}=L^{\mathsf{s}}$ units of $g$, and a barycentre with $g$-coordinate $x(p)$ has
$\gamma_{\star}$-coordinate $L^{-\mathsf{s}}x(p)$. Moreover
$h_i(L^{-\mathsf{s}}x(p))=f_i(x(p))$, so the plaquette weights of the observables
$\mathcal{O}(f_i)$ of Definition~\ref{1.3:def-curvature-field} are the same numbers in both
members; only the unit in which the barycentres are named differs.

Under the change of variable $x'=L^{-\mathsf{s}}x$ the supports, separations, sup norms and
Lipschitz constants transform as follows:
\begin{itemize}
\item $\operatorname{supp}h_i=L^{-\mathsf{s}}\operatorname{supp}f_i$, and all distances are
multiplied by $L^{-\mathsf{s}}$;
\item $\|h_i\|_\infty=\|f_i\|_\infty$;
\item $\operatorname{Lip}h_i=L^{\mathsf{s}}\operatorname{Lip}f_i$.
\end{itemize}
These are the statements on supports, separations, sup norms and Lipschitz constants in~(b). As
stated in~(b), they carry membership of $\vec f$ in the class at scale $\mathsf{d}$ to
membership of $\vec h$ in the class at scale $L^{-\mathsf{s}}\mathsf{d}$, for $\mathfrak{P}_n$
and for $\mathfrak{T}$.

It remains to see that the right member is the first-passage object of
Theorem~\ref{3.1:thm-main} at reference coupling $\gamma_{\star}$, that is, that each hypothesis
of Notation~\ref{2.2:not-hypotheses} it involves holds with $\gamma_{\star}$ as reference
coupling.
\begin{itemize}
\item Hypothesis~(H4) asks $\gamma_{\star}\le\gamma_H$, which holds by
Definition~\ref{9.98:def-strong-reference-coupling}.
\item Hypothesis~(H6) asks $g_0\le g_-(\gamma_{\star})$, proved in~(a), with depth
$K'=K(g_0,\gamma_{\star})$, which is the definition of $K'$.
\item Hypothesis~(H5) asks $m-\mathsf{s}\ge m_{\mathbb{T}}(g_-(\gamma_{\star}))$. Since
$m_{\mathbb{T}}(g_-(\gamma_{\star}))=m_{\star}$ by
Definition~\ref{9.98:def-prolongation-depths}, this is the hypothesis
$m\ge\mathsf{s}+m_{\star}$.
\item Hypothesis~(H7) at reference $\gamma_{\star}$ asks $\gamma_{\star}\le\gamma_J$ and, for
the use of clause~(iv), $\gamma_{\star}\le\gamma_W^{\mathbb{T}}$. Both hold by
Definition~\ref{9.98:def-strong-reference-coupling}.
\end{itemize}
This completes~(b).

\emph{Proof of} (c). For each $j$, part~(a) applied to $g_0^{(j)}$ gives
$\mathsf{s}_j\in[\underline{s}(g),\overline{s}(g)]\cap\mathbb{Z}$. This is a finite set that
depends on $g$ alone and not on $j$. Hence some value of this set is taken by
$\mathsf{s}_j$ for infinitely many $j$, and those indices form a subsequence along which
$\mathsf{s}_j$ is constant.
\end{proof}

\begin{remark}[The direction of the prolongation]\label{9.98:rem-prolongation-direction}
Let $g\le\gamma_{\star}$ and $g_0\in\,]0,g_-(g)]$, and let $K=K(g_0,g)$, $K'=K(g_0,\gamma_{\star})$ and
$\mathsf{s}=K'-K$ be as in the run prolongation lemma (Lemma~\ref{9.98:lem-run-prolongation}).
Write $y:=g^{-2}$ and $X_{\star}:=\gamma_{\star}^{-2}$.
\begin{itemize}
\item[(a)] Reading the reference length $\mathsf{s}$ levels coarser, as in part (b) of
Lemma~\ref{9.98:lem-run-prolongation}, is the same as prolonging the first-passage run from
$g_0$ by $\mathsf{s}$ steps toward the infrared, where the coupling grows. Along these steps the
running inverse coupling decreases from the window $]y,y+B^{\sharp}]$ to the window
$W_{\star}=\,]X_{\star},X_{\star}+B^{\sharp}]$.
\item[(b)] The prolongation is admissible exactly as long as the running coupling stays below the
ceiling up to which the clauses of Theorem~\ref{3.1:thm-main} invoked along it hold. Here this is
up to the first passage into $W_{\star}$, at the coupling $\gamma_{\star}=\gamma_W^{\mathbb{T}}$, and
$\gamma_W^{\mathbb{T}}\le\gamma_J\le\gamma_H$.
\item[(c)] The identity of part (b) of Lemma~\ref{9.98:lem-run-prolongation}, read from right to
left, is the change-of-units identity in the direction of weaker coupling.
\item[(d)] Lemma~\ref{9.98:lem-run-prolongation} reads neither the never-healed old-age tail, nor
any rarity statement for large-field blocks, nor any large-field history $h$, nor any decay of
correlations. It is the bookkeeping of clause~(0) of Theorem~\ref{3.1:thm-main} on one lattice
integral.
\end{itemize}
\end{remark}

\begin{proof}
By the definition of first passage, the running inverse coupling of the run from $g_0$ lies in
$]y,y+B^{\sharp}]$ at level $K$ and in $W_{\star}$ at level $K'$. Since $g\le\gamma_{\star}$, we
have $X_{\star}\le y$. By part (a) of Lemma~\ref{9.98:lem-run-prolongation}, $K'\ge K$. Hence the
$\mathsf{s}$ levels beyond $K$ are further steps of the same run.

By the running-coupling bounds \eqref{9.98:eq-couplings-limit-sets-a}, the inverse coupling
decreases along the run. It therefore passes from the first window to the second, that is, the
coupling grows. This proves (a).

Part (b) of Lemma~\ref{9.98:lem-run-prolongation} identifies the integral at reference coupling
$g$ on $\mathbb{T}_m$ with the object of Theorem~\ref{3.1:thm-main} at reference coupling
$\gamma_{\star}$ on $\mathbb{T}_{m-\mathsf{s}}$. The clauses of Theorem~\ref{3.1:thm-main} applied
to the right member are those at reference coupling $\gamma_{\star}$. Their use requires the running
coupling to stay below their ceiling, which holds up to the first passage into $W_{\star}$ at
$\gamma_{\star}=\gamma_W^{\mathbb{T}}$. By definition
$\gamma_W^{\mathbb{T}}=\min\{\gamma_W,\gamma_{*}^{\dagger}\}\le\gamma_W$, and
$\gamma_W\le\gamma_J\le\gamma_H$ is the order of the coupling thresholds fixed in
Notation~\ref{2.3:not-stages-last}. This gives (b).

Read from right to left, the identity of part (b) of Lemma~\ref{9.98:lem-run-prolongation} passes
from the frame of $\gamma_{\star}$ to the finer reference length at the weaker coupling
$g\le\gamma_{\star}$. This is the change-of-units identity toward weaker coupling, which gives (c).

Finally, the proof of Lemma~\ref{9.98:lem-run-prolongation} uses only two inputs: clause~(0) of
Theorem~\ref{3.1:thm-main}, taken by statement, and the change-of-units identity. None of the items
listed in (d) is among these inputs, which gives (d).
\end{proof}

\begin{lemma}[The tuned member read at $\gamma_{\star}$]\label{9.98:lem-tuned-member}
Let $\gamma_{\star}$, $m_{\star}=m_{\mathbb{T}}(g_-(\gamma_{\star}))$, $\underline{s}(g)$, $\overline{s}(g)$, $\underline{s}'(g)$,
$\overline{s}'(g)$ and $m_{\natural}(g)$ be as in Definition~\ref{9.98:def-prolongation-depths}, and let $j_U$, $n_{\natural}$ be
as in Lemma~\ref{9.98:lem-tuned-depth-offset}.

Fix the following data:
\begin{itemize}
\item a coupling $g\le\min\{\gamma_{\star},\gamma_U'\}$;
\item a shift $\hat w\in[-\varrho_2/2,\varrho_2/2]$;
\item a depth $n\ge n_{\natural}$;
\item a number $\mathsf{q}\in\{2,3\}$ of observables;
\item a torus exponent $m$.
\end{itemize}
The depth-$n$ member of the tuned family of clause (v) of Theorem~0
(Theorem~\ref{3.1:thm-main}) on $\mathbb{T}_m$ is the Wilson measure on the lattice
$\Lambda_{n,m}$ with $2L^{m+n}$ sites a side, at bare inverse coupling $x_0:=a_n(g)-\hat w$.
We write $S^{n}_{\mathbb{T}_m}(x_0;\vec f\,)$ for its connected $\mathsf{q}$-point function of
$\mathcal{O}(f_1),\dots,\mathcal{O}(f_{\mathsf{q}})$, as in clause (ii-$\mathbb{T}$-a) of Theorem~0. Here
$\vec f=(f_1,\dots,f_{\mathsf{q}})$ is read in $g$-reference units, in which the reference length is $L^{n}$
spacings.

Put
\begin{equation*}
g_0:=x_0^{-1/2},\qquad K_0:=K(g_0,g),\qquad K':=K(g_0,\gamma_{\star}),
\end{equation*}
and
\begin{equation*}
\mathsf{s}:=K'-K_0,\qquad \jmath:=K_0-n,\qquad \mathsf{s}':=K'-n .
\end{equation*}
For $\mathsf{d}>0$, let $\mathfrak{P}_{\mathsf{q}}(\mathsf{d};\vartheta)$ be the set of $\mathsf{q}$-tuples of real
Lipschitz functions of any signs such that:
\begin{itemize}
\item the supports are pairwise disjoint and lie inside one closed ball of radius $5\mathsf{d}$;
\item $\operatorname{Lip} f_i\le\|f_i\|_{\infty}/(\vartheta\mathsf{d})$ for each $i$.
\end{itemize}
Then the following hold.
\begin{itemize}
\item[(a)] $g_0\le g_-(g)$. Further, $K'$ is well defined and $K'=K_0+\mathsf{s}$ with
$\underline{s}(g)\le\mathsf{s}\le\overline{s}(g)$. Moreover $|\jmath|\le j_U-1$ and
\begin{equation*}
\mathsf{s}'=\mathsf{s}+\jmath\in[\underline{s}'(g),\overline{s}'(g)].
\end{equation*}
\item[(b)] Take the reference length to be $L^{K'}$ spacings. Then $\Lambda_{n,m}$ is the bare lattice of the
torus $\mathbb{T}_{m-\mathsf{s}'}$ at cutoff $K'=K(g_0,\gamma_{\star})$. This torus is a torus of Theorem~0's
family at reference $\gamma_{\star}$ if and only if $m-\mathsf{s}'\ge m_{\star}$. If $m\ge m_{\natural}(g)$, then
$m-\mathsf{s}'\ge m_{\star}+1$ and
\begin{equation}\label{9.98:eq-tuned-member-a}
S^{n}_{\mathbb{T}_m}(x_0;\vec f\,)=S^T_{\mathsf{q};K',m-\mathsf{s}'}[g_0]\bigl(\vec f(L^{\mathsf{s}'}\cdot)\bigr).
\end{equation}
The right member is the connected $\mathsf{q}$-point function of Theorem~0 in the first-passage labelling at
reference coupling $\gamma_{\star}$ on $\mathbb{T}_{m-\mathsf{s}'}$.
\item[(c)] Put $\mathsf{d}':=L^{-\mathsf{s}'}\mathsf{d}$. If $\vec f\in\mathfrak{P}_{\mathsf{q}}(\mathsf{d};\vartheta)$, then
$\vec f(L^{\mathsf{s}'}\cdot)\in\mathfrak{P}_{\mathsf{q}}(\mathsf{d}';\vartheta)$, and the sup norms are unchanged. If
moreover $\mathsf{d}\le L^{\underline{s}'(g)}/40$, then
\begin{equation*}
\mathsf{d}'\le L^{-\underline{s}'(g)}\mathsf{d}\le 1/40 .
\end{equation*}
The spacing $\varepsilon':=L^{-K'}$ of the $\gamma_{\star}$-frame satisfies $\varepsilon'\le\mathsf{d}'$ if and only if
$L^{-n}\le\mathsf{d}$. In particular $\varepsilon'\le\mathsf{d}'$ for all sufficiently large $n$ at fixed $\mathsf{d}$.
\item[(d)] (First-passage frame.) Consider a member $(g_0,K,\mathbb{T}_m)$ of clauses (ii-$\mathbb{T}$-b) and
(ii-$\mathbb{R}^4$) of Theorem~0 with $g_0\in\,]0,g_-(g)]$ and $K=K(g_0,g)$. Read $\vec f$ in $g$-reference
units, in which the reference length is $L^{K}$ spacings. Put $n:=K$ and $x_0:=g_0^{-2}$. Then $\Lambda_{n,m}$
is the bare lattice of $\mathbb{T}_m$ at cutoff $K$, and
\begin{equation*}
S^{n}_{\mathbb{T}_m}(x_0;\vec f\,)=S^T_{\mathsf{q};K,m}[g_0](\vec f\,).
\end{equation*}
Further $K_0=K$, $\jmath=0$ and $\mathsf{s}'=\mathsf{s}\in[\underline{s}(g),\overline{s}(g)]$. In this frame:
\begin{itemize}
\item the display~\eqref{9.98:eq-tuned-member-a} is the identity of Lemma~\ref{9.98:lem-run-prolongation}(b);
\item the torus $\mathbb{T}_{m-\mathsf{s}}$ of the right member is a torus of Theorem~0's family at reference
$\gamma_{\star}$ for $m\ge m_{\star}+\mathsf{s}$;
\item $\mathsf{d}'=L^{-\mathsf{s}}\mathsf{d}\le 1/40$ whenever $\mathsf{d}\le L^{\underline{s}(g)}/40$.
\end{itemize}
\end{itemize}
\end{lemma}

\begin{proof}
(a) Since $g_-(g)=(g^{-2}+B^{\sharp})^{-1/2}$, the inequality $g_0\le g_-(g)$ is equivalent to
$x_0\ge g^{-2}+B^{\sharp}$. For $|\hat w|\le\varrho_2/2$ and $n\ge n_{\natural}$,
\begin{align*}
x_0=a_n(g)-\hat w&\ge a_n(g)-\varrho_2/2\\
&\ge g^{-2}+\lambda_U n-2c_2-\varrho_2/2\ge g^{-2}+B^{\sharp}.
\end{align*}
The second inequality is the lower half of the spacing bound on $a_n(g)$ in clause (v) of Theorem~0; the third
holds for $n\ge n_{\natural}$ by the choice of the floor $n_{\natural}$ in
Lemma~\ref{9.98:lem-tuned-depth-offset}. The bound also gives $x_0>0$, so $g_0\in\,]0,g_-(g)]$.

We now apply Lemma~\ref{9.98:lem-run-prolongation}(a). Its hypotheses are $g\le\gamma_{\star}$ and
$g_0\in\,]0,g_-(g)]$, and both hold. It yields three facts: $K'=K(g_0,\gamma_{\star})$ is well defined,
$K'\ge K_0$, and $\underline{s}(g)\le\mathsf{s}\le\overline{s}(g)$.

Next, Lemma~\ref{9.98:lem-tuned-depth-offset} applies, since $n\ge n_{\natural}$, $|\hat w|\le\varrho_2/2$ and
$g\le\gamma_U'$. It gives $|\jmath|=|K_0-n|\le j_U-1$.

Finally,
\begin{equation*}
\mathsf{s}'=K'-n=(K'-K_0)+(K_0-n)=\mathsf{s}+\jmath .
\end{equation*}
Combining $\underline{s}(g)\le\mathsf{s}\le\overline{s}(g)$ with $|\jmath|\le j_U-1$ places $\mathsf{s}+\jmath$ in the
interval $[\underline{s}'(g),\overline{s}'(g)]$ defined in~\eqref{9.98:eq-prolongation-depths-a}.

(b) Since $\mathsf{s}'=K'-n$, we have $m+n=(m-\mathsf{s}')+K'$. Measured in the reference length $L^{K'}$
spacings, the lattice $\Lambda_{n,m}$ therefore has:
\begin{itemize}
\item side $2L^{m+n}\cdot L^{-K'}=2L^{m-\mathsf{s}'}$;
\item spacing $L^{-K'}$.
\end{itemize}
Hence, in the sense of Definition~\ref{1.1:def-lattice}, it is the bare lattice of $\mathbb{T}_{m-\mathsf{s}'}$ at
cutoff $K'$. The bare coupling is $g_0$, and $K'=K(g_0,\gamma_{\star})$ is its first-passage depth at reference
$\gamma_{\star}$. The torus $\mathbb{T}_{m-\mathsf{s}'}$ is a torus of the family at reference $\gamma_{\star}$ if and
only if $m-\mathsf{s}'\ge m_{\star}$, since $m_{\star}=m_{\mathbb{T}}(g_-(\gamma_{\star}))$
(Definition~\ref{9.98:def-prolongation-depths}).

Let $m\ge m_{\natural}(g)$. By the definition of $m_{\natural}(g)$
in~\eqref{9.98:eq-prolongation-depths-c}, $m_{\natural}(g)\ge m_{\star}+\overline{s}(g)+j_U$. Together with
$\mathsf{s}\le\overline{s}(g)$ and $\jmath\le j_U-1$ from (a), this gives
\begin{equation*}
m-\mathsf{s}'=m-\mathsf{s}-\jmath\ge m_{\star}+\overline{s}(g)+j_U-\overline{s}(g)-(j_U-1)=m_{\star}+1 .
\end{equation*}

The identity~\eqref{9.98:eq-tuned-member-a} is obtained in two renamings of one lattice integral.

First step. Apply the change-of-units identity of clause (v) of Theorem~0 by $\jmath$ levels. This passes from the
tuned reference length, $L^{n}$ spacings, to the first-passage reference length at reference $g$, namely
$L^{K_0}=L^{n+\jmath}$ spacings. The lattice
$\Lambda_{n,m}$, of side $2L^{(m-\jmath)+K_0}$, becomes the bare lattice of $\mathbb{T}_{m-\jmath}$ at cutoff $K_0$,
and the test functions become $\vec f(L^{\jmath}\cdot)$:
\begin{equation}\label{9.98:eq-tuned-member-1}
S^{n}_{\mathbb{T}_m}(x_0;\vec f\,)=S^T_{\mathsf{q};K_0,m-\jmath}[g_0]\bigl(\vec f(L^{\jmath}\cdot)\bigr).
\end{equation}
Since $K_0=K(g_0,g)$, the right member is Theorem~0's first-passage object at reference $g$ on
$\mathbb{T}_{m-\jmath}$. The torus $\mathbb{T}_{m-\jmath}$ belongs to Theorem~0's family at reference $g$ because
$m-\jmath\ge m_{\mathbb{T}}(g_-(g))$
by~\eqref{9.98:eq-floor-members-a} of Remark~\ref{9.98:rem-floor-members}.

Second step. Apply Lemma~\ref{9.98:lem-run-prolongation}(b) by $\mathsf{s}$ further levels. The data are:
\begin{itemize}
\item bare coupling $g_0$, with $K=K_0$, $K'$ and $\mathsf{s}$ as in part (a) of that lemma;
\item torus exponent $m-\jmath$;
\item the tuple $\vec f(L^{\jmath}\cdot)$.
\end{itemize}
Its hypotheses are $\mathsf{q}\ge 2$ and $m-\jmath\ge\mathsf{s}+m_{\star}$. The latter is $m-\mathsf{s}'\ge m_{\star}$,
proved above. The lemma turns the right member of~\eqref{9.98:eq-tuned-member-1} into
$S^T_{\mathsf{q};K',m-\jmath-\mathsf{s}}[g_0]$ evaluated at $\vec f(L^{\jmath}L^{\mathsf{s}}\cdot)$. Since
$m-\jmath-\mathsf{s}=m-\mathsf{s}'$ and $L^{\jmath}L^{\mathsf{s}}=L^{\mathsf{s}'}$,
this is~\eqref{9.98:eq-tuned-member-a}.

(c) Let $\vec f\in\mathfrak{P}_{\mathsf{q}}(\mathsf{d};\vartheta)$, with all supports in the closed ball
$\overline{B}(y,5\mathsf{d})$. The map $x\mapsto L^{\mathsf{s}'}x$ is a bijection. Hence
$\operatorname{supp} f_i(L^{\mathsf{s}'}\cdot)=L^{-\mathsf{s}'}\operatorname{supp} f_i$, and these sets are pairwise
disjoint and contained in $\overline{B}(L^{-\mathsf{s}'}y,5\mathsf{d}')$. The sup norms are unchanged. For the
Lipschitz constants,
\begin{equation*}
\operatorname{Lip} f_i(L^{\mathsf{s}'}\cdot)=L^{\mathsf{s}'}\operatorname{Lip} f_i
\le L^{\mathsf{s}'}\frac{\|f_i\|_{\infty}}{\vartheta\mathsf{d}}=\frac{\|f_i\|_{\infty}}{\vartheta\mathsf{d}'} .
\end{equation*}
Thus $\vec f(L^{\mathsf{s}'}\cdot)\in\mathfrak{P}_{\mathsf{q}}(\mathsf{d}';\vartheta)$.

By (a), $\mathsf{s}'\ge\underline{s}'(g)$, and $L>1$. Hence, under $\mathsf{d}\le L^{\underline{s}'(g)}/40$,
\begin{equation*}
\mathsf{d}'=L^{-\mathsf{s}'}\mathsf{d}\le L^{-\underline{s}'(g)}\mathsf{d}\le 1/40 .
\end{equation*}
Finally, $\varepsilon'\le\mathsf{d}'$ reads $L^{-K'}\le L^{-\mathsf{s}'}\mathsf{d}$, that is
$L^{\mathsf{s}'-K'}\le\mathsf{d}$. Since $\mathsf{s}'-K'=-n$, this is $L^{-n}\le\mathsf{d}$, which holds for all
sufficiently large $n$ when $\mathsf{d}$ is fixed.

(d) Here $n=K$ and $x_0=g_0^{-2}$. The lattice $\Lambda_{n,m}$ has $2L^{m+K}$ sites a side; measured in the
reference length $L^{K}$ spacings, its side is $2L^{m}$ and its spacing is $L^{-K}$, so by
Definition~\ref{1.1:def-lattice} it is the bare lattice of $\mathbb{T}_m$ at cutoff $K$. The Wilson measure on it at
bare inverse coupling $x_0=g_0^{-2}$ has bare coupling $g_0$, and its connected $\mathsf{q}$-point function of
$\mathcal{O}(f_1),\dots,\mathcal{O}(f_{\mathsf{q}})$ is therefore $S^T_{\mathsf{q};K,m}[g_0](\vec f\,)$.

Further $K_0=K(g_0,g)=K$, and the reference length is $L^{K}=L^{K_0}$ spacings. Hence no first renaming is
needed: $\jmath=K_0-n=0$ and $\mathsf{s}'=\mathsf{s}$. Lemma~\ref{9.98:lem-run-prolongation}(a) applies because
$g\le\gamma_{\star}$ and $g_0\in\,]0,g_-(g)]$, and it gives $\mathsf{s}\in[\underline{s}(g),\overline{s}(g)]$.

The display~\eqref{9.98:eq-tuned-member-a} with $\mathsf{s}'=\mathsf{s}$ is the identity of
Lemma~\ref{9.98:lem-run-prolongation}(b) as stated there. That identity holds for $m\ge\mathsf{s}+m_{\star}$, and
its right member is Theorem~0's object at reference $\gamma_{\star}$ on $\mathbb{T}_{m-\mathsf{s}}$.

The argument of (c) with $\mathsf{s}$ in place of $\mathsf{s}'$ gives $\mathsf{d}'=L^{-\mathsf{s}}\mathsf{d}$. From
$\mathsf{s}\ge\underline{s}(g)$ we get $\mathsf{d}'\le L^{-\underline{s}(g)}\mathsf{d}$, and this is at most $1/40$
under $\mathsf{d}\le L^{\underline{s}(g)}/40$.
\end{proof}

\begin{remark}[The two tori in the volume transfer]\label{9.98:rem-transfer-two-tori}
Let $m_{\star}=m_{\mathbb{T}}(g_-(\gamma_{\star}))$ be the reference torus exponent at $\gamma_{\star}$, and
write $\mathbb{T}_{\star}=\mathbb{T}_{m_{\star}}$ and $\mathbb{T}_{\star}'=\mathbb{T}_{m_{\star}+1}$.

The volume transfer theorem is a statement about the connected functions of the smeared
curvature field $\mathcal{O}$ under two bare Wilson measures on two tori $\mathbb{T}_{m_1}$ and
$\mathbb{T}_{m_2}$, where $m_1,m_2\ge m_{\star}$.
The two measures have the same bare inverse coupling $x_0$, hence the same bare coupling
$g_0=x_0^{-1/2}$, and the same number $K':=K(g_0,\gamma_{\star})$ of renormalisation steps.

Its proof compares each of these two measures with the bare Wilson measure on the reference
torus $\mathbb{T}_{\star}$ with the same $x_0$ and the same $K'$, and uses the following inputs.
\begin{itemize}
\item The correlation theorem is used on one torus only, the reference torus $\mathbb{T}_{\star}$, at
$\gamma_{\star}$. It is applied along the first-passage run of the measure on that torus: the run of
$K'$ steps that starts from the bare inverse coupling $x_0$.
\item Of a torus $\mathbb{T}_{m}$ with $m\ge m_{\star}$ compared with $\mathbb{T}_{\star}$,
it uses only the following properties of its bare measure: reflection positivity in link
reflections, the chessboard estimate, the trace inequality, and the lower stability bound at
zero source, which is the lower member of clause~(i-b) of
Theorem~\ref{3.1:thm-main}. In the convexity variant of the volume transfer theorem, that bound
is used on $\mathbb{T}_{\star}'$ only, and the convexity lemma of that variant replaces it on tori
of exponent larger than $m_{\star}+1$.
\item The proof further uses the one-sided bound, with constant $\bar{\delta}$, on the difference
of the vacuum-energy densities of $\mathbb{T}_{\star}$ and $\mathbb{T}_{m}$, and clause~(0) of
Theorem~\ref{3.1:thm-main}.
\end{itemize}

Consequently, when the measures are followed to the reference coupling $\gamma_{\star}$ as
above, one statement is not needed. This is the statement that the correlation theorem, the
lower stability bound and the per-ball volume-free bound hold along the runs used in the proof
of the uniqueness theorem, clause~(v) of Theorem~\ref{3.1:thm-main}. Establishing these three
statements along those runs would give an alternative route; it is not the one taken here.
\end{remark}

\begin{proof}
Consider first the reference torus $\mathbb{T}_{\star}$. The correlation theorem is applied there to one
bare Wilson measure, with bare inverse coupling $x_0$. For the tuned member of
Lemma~\ref{9.98:lem-tuned-member} this is $x_0=a_n(g)-\hat{w}$; for a measure given by its bare
coupling $g_0$ it is $x_0=g_0^{-2}$. The run along which the theorem is applied is the sequence
of running inverse couplings of Definition~\ref{1.2:def-couplings}, started at $x_0$ and followed
for $K'$ steps. This run exists by clause~(0) of Theorem~\ref{3.1:thm-main}, and its definition
makes no reference to the runs used in the proof of the uniqueness theorem; the correlation
theorem is applied along this run and no other.

None of the inputs listed in the statement for the torus $\mathbb{T}_{m}$, nor the bound
$\bar{\delta}$, nor clause~(0) of Theorem~\ref{3.1:thm-main}, refers to a run of renormalisation
steps other than the first-passage run of $K'$ steps from $x_0$.

Hence no input of the proof refers to the runs used in the proof of the uniqueness theorem,
which is the assertion of the remark.
\end{proof}

\begin{proposition}[The volume-transfer bound for three sources]\label{9.98:prop-transfer-three-sources}
Let $g^{\flat}\le\gamma_J$ be a reference coupling, let $g_0\in\,]0,g_-(g^{\flat})]$, let
$K=K(g_0,g^{\flat})$ be the first-passage number of steps of
Definition~\ref{1.2:def-flow-constants}, put $\varepsilon:=L^{-K}$, and let $R>0$. Let $\ell_3$ be
the least power of $L$ that is $\ge 24(R+1)$; this is the chessboard cell side used for three
sources in the per-ball volume-free bound. Put
\begin{equation*}
m_{\star}(g^{\flat}):=\max\bigl\{m_{\mathbb{T}}(g_-(g^{\flat})),\ \log_L\ell_3\bigr\},
\end{equation*}
and let $m\ge m'\ge m_{\star}(g^{\flat})$ be torus exponents (here $m'$ denotes a torus exponent,
not the exponent of $M$). When $R=1$ one has $\ell_3\le 48L<L^{11}M$, where $M$ is Ba{\l}aban's
auxiliary parameter of Notation~\ref{0.2:not-glossary}, so that
$m_{\star}(g^{\flat})=m_{\mathbb{T}}(g_-(g^{\flat}))$, without any condition on $L$.

For $m''\in\{m,m'\}$ write $N_{m''}:=2L^{m''}$, let $\mathcal{C}_{m''}$ be the set of chessboard
cells of side $\ell_3$ of $\mathbb{T}_{m''}$, and let $Z_{m''}(w)$ be the partition function of
$\mathbb{T}_{m''}$ with complex source $w$ at bare coupling $g_0$ and spacing $\varepsilon$
(Notation~\ref{0.2:not-glossary}), so that
$Z_{m''}(w)/Z_{m''}(0)=\langle e^{\mathcal{O}(w)}\rangle_{m''}$, where $\langle\cdot\rangle_{m''}$
denotes expectation in the normalised Wilson measure of Definition~\ref{1.1:def-wilson-action} on
$\mathbb{T}_{m''}$ at bare coupling $g_0$; let $E_{m''}$, $\mathcal{E}_{m''}$ and
$\tilde{Z}_{m''}$ be the vacuum energy, the source counterterm and the normalised partition function
of the correlation theorem on $\mathbb{T}_{m''}$, and let $\mathcal{W}$ be that theorem's class of
admissible weights. Let $\mathfrak{r}\in\{1,2\}$ be the factor of the two readings of the correlation
theorem used in the per-ball volume-free bound: $\mathfrak{r}=1$ in the class reading, in which the
theorem is applied to complex weights of $\mathcal{W}$, and $\mathfrak{r}=2$ in the literal reading,
in which it is applied to real systems of functions. Assume:
\begin{itemize}
\item[(a)] the correlation theorem, with all its conclusions, on the torus $\mathbb{T}_{m'}$ for the
run of $K$ steps from $g_0$ (in the reading fixed by $\mathfrak{r}$); in particular its upper
stability bound: there is a constant $E_{+}^{\sharp}$, depending on $(L,\mathfrak{p},g^{\flat},\varrho)$
only and not on $K$, $g_0$, $m$, $m'$, such that for every $W\in\mathcal{W}$
descending to $\mathbb{T}_{m'}$
\begin{equation}\label{9.98:eq-transfer-three-sources-1}
\begin{split}
|Z_{m'}(W)|&=e^{E_{m'}+\operatorname{Re}\mathcal{E}_{m'}(W)}\,|\tilde{Z}_{m'}(W)|\\
&\le\exp\bigl[E_{m'}+\operatorname{Re}\mathcal{E}_{m'}(W)+E_{+}^{\sharp}N_{m'}^4\bigr];
\end{split}
\end{equation}
\item[(b)] the following ingredients of the per-ball volume-free bound, in their instance at three
sources: reflection positivity of the lattice measure, the chessboard estimate, the placement lemma,
the lemma on mirror-periodised weights, the class and literal readings of the correlation theorem,
the face-term lemma with its constant $C_{\mathrm{face}}(R)$, depending on
$(L,\mathfrak{p},g^{\flat},\varrho)$ and $R$ only and not on $K$, $g_0$, $m$, $m'$, the
holomorphy of the re-centred counterterm and the identity for its real part; and the
Schwarz--Cauchy lemma, valid for any number of sources (each is recalled in the
form used at the place where it is used);
\item[(c)] the trace inequality between tori: every mirror-periodic weight $W$
descending to $\mathbb{T}_m$ and to $\mathbb{T}_{m'}$ satisfies
$Z_m(W)\le Z_{m'}(W)^{L^{4(m-m')}}$;
\item[(d)] the lower bound of clause (i) at the last level and zero source: there is a constant
$C_{\mathrm{low}}$, depending on $(L,\mathfrak{p},g^{\flat},\varrho)$ only and not on $K$, $g_0$,
$m$, $m'$, such that
$Z_{m''}(0)\ge\exp[E_{m''}-C_{\mathrm{low}}N_{m''}^4]$ for $m''\in\{m,m'\}$;
\item[(e)] the one-sided bound on the vacuum-energy densities: there is a constant $\bar{\delta}$,
depending on $(L,\mathfrak{p},g^{\flat},\varrho)$ only and not on $K$, $g_0$, $m$, $m'$,
such that $E_{m_2}/N_{m_2}^4-E_{m_1}/N_{m_1}^4\le\bar{\delta}$ for all integers
$m_1\ge m_2\ge m_{\star}(g^{\flat})$.
\end{itemize}
Put $\mathsf{B}:=E_{+}^{\sharp}+C_{\mathrm{low}}$,
\begin{equation*}
Q_{\star,3}:=\exp\bigl[(\mathsf{B}+\bar{\delta})\ell_3^4+6C_{\mathrm{face}}(R)\bigr],\qquad
C_{VT,3}:=2\log 2\cdot(2Q_{\star,3}+2)^3\Bigl(\frac{3\mathfrak{r}}{\varrho}\Bigr)^3 .
\end{equation*}
Then for all real Lipschitz functions $F_1,F_2,F_3$ on $\mathbb{R}^4$ supported in one ball of radius
$R$ (transported to both tori), with no condition on their signs, on the components of their supports,
or on overlaps or distances between their supports,
\begin{equation}\label{9.98:eq-transfer-three-sources-2}
\bigl|S^T_{3;K,m}(F_1,F_2,F_3)-S^T_{3;K,m'}(F_1,F_2,F_3)\bigr|
\le C_{VT,3}\prod_{i=1}^3\|F_i\|_{\sharp}.
\end{equation}
The constant $C_{VT,3}$ depends on $(L,\mathfrak{p},g^{\flat},\varrho;R)$ only; it does not depend on
$K$, $g_0$, $m$, $m'$, the position of the ball, or the supports of the $F_i$.
\end{proposition}

\begin{proof}
In the corollary to clause~(v) of Theorem~\ref{3.1:thm-main} on $\mathbb{R}^4$ the proposition is
applied with $g^{\flat}=\gamma_{\star}$.
For $m''\in\{m,m'\}$ the number $N_{m''}/\ell_3=2L^{m''}/\ell_3$ is an even integer, since $\ell_3$ is
a power of $L$ and $L^{m''}\ge\ell_3$ by $m''\ge m_{\star}(g^{\flat})\ge\log_L\ell_3$. Hence
$|\mathcal{C}_{m''}|=(N_{m''}/\ell_3)^4$, that is,
\begin{equation}\label{9.98:eq-transfer-three-sources-3}
N_{m''}^4/|\mathcal{C}_{m''}|=\ell_3^4\qquad(m''\in\{m,m'\}).
\end{equation}
For $z\in\mathbb{C}^3$ put $w_z:=z_1F_1+z_2F_2+z_3F_3$, and let
\begin{equation*}
\mathbb{D}^{(3)}_{\mathfrak{r}}:=\Bigl\{z\in\mathbb{C}^3:\ \sum_{i=1}^3|z_i|\,\|F_i\|_{\sharp}
<\varrho/\mathfrak{r}\Bigr\}.
\end{equation*}

\emph{Placement.} Since the three functions are supported in one ball ($B_1=B_2=B_3$ in the notation
of the placement lemma), that lemma gives one translation $a\in\varepsilon\mathbb{Z}^4$ which puts the
ball in the interior of a cell $c$ of side $\ell_3$, at distance at least $R+2$ from all the face
planes of the cells. Both bare Wilson measures, on $\mathbb{T}_m$ and on $\mathbb{T}_{m'}$, are
invariant under the translation by $a$ (Lemma~\ref{1.1:lem-invariance}), so neither
$S^T_{3;K,m}(F_1,F_2,F_3)$ nor $S^T_{3;K,m'}(F_1,F_2,F_3)$ changes when the $F_i$ are replaced by
their translates; we assume this replacement made. There is
then a single occupied cell, and it carries all three indices.

Let $\mathfrak{G}_{\ell_3}$ be the group generated by the reflections of $\mathbb{R}^4$ in the face
planes of the cells of side $\ell_3$; for $\eta\in\mathfrak{G}_{\ell_3}$ let $\epsilon(\eta)$ be its
parity, and let $u\mapsto u^{(\epsilon(\eta))}$ be the operation attached to that parity in the
lemma on mirror-periodised weights. For $z\in\mathbb{D}^{(3)}_{\mathfrak{r}}$ put
\begin{equation}\label{9.98:eq-transfer-three-sources-4}
W_z:=\sum_{\eta\in\mathfrak{G}_{\ell_3}}\bigl(w_z\circ\eta^{-1}\bigr)^{(\epsilon(\eta))}.
\end{equation}
This weight is mirror-periodic in the sense of the lemma on mirror-periodised weights: the condition
of invariance under $\mathfrak{G}_{\ell_3}$ together with the parity operation holds by the
construction \eqref{9.98:eq-transfer-three-sources-4}, and the condition on the distance of the
support of $w_z$ from the face planes holds because of the
margin $R+2>2\varepsilon$ between the ball and the face planes. By that lemma $W_z$ descends to every
torus $\mathbb{T}_{m''}$ whose cell count $N_{m''}/\ell_3$ is even, hence to $\mathbb{T}_m$ and to
$\mathbb{T}_{m'}$, as the periodised weight of the lemma. The reflected copies
$(w_z\circ\eta^{-1})^{(\epsilon(\eta))}$ have pairwise disjoint supports, so by the same lemma
\begin{equation*}
\|W_z\|_{\sharp}\le\sum_{i=1}^3|z_i|\,\|F_i\|_{\sharp}<\varrho/\mathfrak{r}.
\end{equation*}
In the class reading ($\mathfrak{r}=1$) this gives $W_z\in\mathcal{W}$. In the literal reading one
applies the correlation theorem to the real system of six functions attached to $(F_1,F_2,F_3)$ at
$(z,\bar{z})$, with $\mathfrak{r}=2$; the number of sources enters no constant of the correlation
theorem.

\emph{Chessboard bound on the large torus.} The observable $e^{\mathcal{O}(w_z)}$ depends only on the
bond variables of the one cell $c$. Since $Z_m(0)>0$,
$|Z_m(w_z)|/Z_m(0)=|\langle e^{\mathcal{O}(w_z)}\rangle_m|$. The chessboard estimate, which rests on
reflection positivity of the lattice measure on $\mathbb{T}_m$, bounds this by the chessboard
seminorm, and the lemma on mirror-periodised weights identifies that seminorm:
\begin{equation}\label{9.98:eq-transfer-three-sources-5}
\frac{|Z_m(w_z)|}{Z_m(0)}=|\langle e^{\mathcal{O}(w_z)}\rangle_m|
\le\|e^{\mathcal{O}(w_z)}\|_{\mathrm{ch}}
=\Bigl[\frac{Z_m(W_z)}{Z_m(0)}\Bigr]^{1/|\mathcal{C}_m|}.
\end{equation}
The correlation theorem is not used here.

\emph{Passage from the large torus to the small one.} Hypothesis (c) applied to
the mirror-periodic weight $W_z$ gives $Z_m(W_z)\le Z_{m'}(W_z)^{L^{4(m-m')}}$. Since
$N_m=L^{m-m'}N_{m'}$, \eqref{9.98:eq-transfer-three-sources-3} gives
$|\mathcal{C}_m|=L^{4(m-m')}|\mathcal{C}_{m'}|$,
and therefore
\begin{equation}\label{9.98:eq-transfer-three-sources-6}
Z_m(W_z)^{1/|\mathcal{C}_m|}\le Z_{m'}(W_z)^{1/|\mathcal{C}_{m'}|}.
\end{equation}

\emph{Stability bounds.} The lemma on mirror-periodised weights, applied on $\mathbb{T}_{m'}$, whose
cell count is even, identifies the chessboard seminorm there with
$[Z_{m'}(W_z)/Z_{m'}(0)]^{1/|\mathcal{C}_{m'}|}$; hence $Z_{m'}(W_z)\ge0$, that is,
$Z_{m'}(W_z)=|Z_{m'}(W_z)|$. By hypothesis (a) at $W=W_z\in\mathcal{W}$,
\begin{equation*}
\begin{split}
Z_{m'}(W_z)=|Z_{m'}(W_z)|&=e^{E_{m'}+\operatorname{Re}\mathcal{E}_{m'}(W_z)}|\tilde{Z}_{m'}(W_z)|\\
&\le\exp\bigl[E_{m'}+\operatorname{Re}\mathcal{E}_{m'}(W_z)+E_{+}^{\sharp}N_{m'}^4\bigr],
\end{split}
\end{equation*}
and by hypothesis (d) on $\mathbb{T}_m$, $Z_m(0)\ge\exp[E_m-C_{\mathrm{low}}N_m^4]$. Inserting both
into \eqref{9.98:eq-transfer-three-sources-5}--\eqref{9.98:eq-transfer-three-sources-6} yields
\begin{equation*}
\frac{|Z_m(w_z)|}{Z_m(0)}\le\exp\Bigl\{\frac{E_{m'}+\operatorname{Re}\mathcal{E}_{m'}(W_z)
+E_{+}^{\sharp}N_{m'}^4}{|\mathcal{C}_{m'}|}
+\frac{-E_m+C_{\mathrm{low}}N_m^4}{|\mathcal{C}_m|}\Bigr\}.
\end{equation*}
By \eqref{9.98:eq-transfer-three-sources-3} on both tori, $E_{m''}/|\mathcal{C}_{m''}|
=\ell_3^4E_{m''}/N_{m''}^4$ and $N_{m''}^4/|\mathcal{C}_{m''}|=\ell_3^4$, so the exponent equals
\begin{equation*}
|\mathcal{C}_{m'}|^{-1}\operatorname{Re}\mathcal{E}_{m'}(W_z)
+\ell_3^4\Bigl(\frac{E_{m'}}{N_{m'}^4}-\frac{E_m}{N_m^4}\Bigr)
+(E_{+}^{\sharp}+C_{\mathrm{low}})\ell_3^4 .
\end{equation*}
Hypothesis (e) with $(m_1,m_2)=(m,m')$ bounds the bracket by $\bar{\delta}$, and
$E_{+}^{\sharp}+C_{\mathrm{low}}=\mathsf{B}$. Hence, for $z\in\mathbb{D}^{(3)}_{\mathfrak{r}}$,
\begin{equation}\label{9.98:eq-transfer-three-sources-7}
\frac{|Z_m(w_z)|}{Z_m(0)}\le\exp\Bigl\{|\mathcal{C}_{m'}|^{-1}\operatorname{Re}\mathcal{E}_{m'}(W_z)
+(\mathsf{B}+\bar{\delta})\ell_3^4\Bigr\}.
\end{equation}

\emph{One counterterm for both tori.} Let $\mathcal{E}''$ be the re-centred counterterm of the per-ball
volume-free bound, formed on $\mathbb{T}_{m'}$ for the system $(F_1,F_2,F_3)$:
\begin{equation*}
\mathcal{E}''(z):=|\mathcal{C}_{m'}|^{-1}\!\!\sum_{\eta\in\mathfrak{G}_{\ell_3}(\mathbb{T}_{m'})}
\Bigl[\mathcal{E}^{\mathrm{bulk}}_{\eta c}\bigl((w_z\circ\eta^{-1})^{(\epsilon(\eta))}\bigr)
\Bigr]^{(\epsilon(\eta))}+\mathcal{E}^{\partial}(w_z),
\end{equation*}
where $\mathfrak{G}_{\ell_3}(\mathbb{T}_{m'})$ is the group generated by the reflections in the cell
faces acting on $\mathbb{T}_{m'}$, $\mathcal{E}^{\mathrm{bulk}}_{c'}$ is the bulk counterterm of the
cell $c'$, $\mathcal{E}^{\partial}$ is the face part of the counterterm, and all coefficients
entering the counterterms are those of $\mathbb{T}_{m'}$. By the holomorphy statement of the per-ball
volume-free bound, $\mathcal{E}''$ is holomorphic on $\mathbb{D}^{(3)}_{\mathfrak{r}}$ and
$\mathcal{E}''(0)=0$. By the identity for the real part established in the proof of that bound,
\begin{equation*}
\operatorname{Re}\mathcal{E}''(z)=|\mathcal{C}_{m'}|^{-1}
\operatorname{Re}\mathcal{E}^{\mathrm{bulk}}_{m'}(W_z)+\operatorname{Re}\mathcal{E}^{\partial}(w_z)
\end{equation*}
identically in $z\in\mathbb{D}^{(3)}_{\mathfrak{r}}$. The face-term lemma, applied with three
ball-supported pieces, gives
\begin{equation*}
|\mathcal{C}_{m'}|^{-1}|\mathcal{E}^{\partial}_{m'}(W_z)|\le 3C_{\mathrm{face}}(R),\qquad
|\mathcal{E}^{\partial}(w_z)|\le 3C_{\mathrm{face}}(R).
\end{equation*}
By the decomposition
$\mathcal{E}_{m'}(W_z)=\mathcal{E}^{\mathrm{bulk}}_{m'}(W_z)+\mathcal{E}^{\partial}_{m'}(W_z)$ of the
counterterm into bulk and face parts, the bulk terms cancel in the difference below, and the two face
bounds give
\begin{equation}\label{9.98:eq-transfer-three-sources-8}
|\mathcal{C}_{m'}|^{-1}\operatorname{Re}\mathcal{E}_{m'}(W_z)-\operatorname{Re}\mathcal{E}''(z)
\le 6C_{\mathrm{face}}(R).
\end{equation}
For $m''\in\{m,m'\}$ define
\begin{equation*}
\mathsf{h}_{m''}(z):=e^{-\mathcal{E}''(z)}\,\frac{Z_{m''}(w_z)}{Z_{m''}(0)},\qquad
z\in\mathbb{D}^{(3)}_{\mathfrak{r}}.
\end{equation*}
At finite $K$ the function $z\mapsto Z_{m''}(w_z)$ is entire, and $Z_{m''}(0)>0$; since
$\mathcal{E}''$ is holomorphic, $\mathsf{h}_{m''}$ is holomorphic on $\mathbb{D}^{(3)}_{\mathfrak{r}}$,
and $\mathsf{h}_{m''}(0)=1$ because $\mathcal{E}''(0)=0$. By
\eqref{9.98:eq-transfer-three-sources-7} and \eqref{9.98:eq-transfer-three-sources-8},
\begin{equation*}
|\mathsf{h}_m(z)|=e^{-\operatorname{Re}\mathcal{E}''(z)}\frac{|Z_m(w_z)|}{Z_m(0)}
\le\exp\bigl[6C_{\mathrm{face}}(R)+(\mathsf{B}+\bar{\delta})\ell_3^4\bigr]=Q_{\star,3}.
\end{equation*}
For $\mathsf{h}_{m'}$, the three preceding paragraphs read with $m$ replaced by $m'$ (then $m=m'$,
hypothesis (d) is used on $\mathbb{T}_{m'}$, and the vacuum-energy bracket vanishes) give
$|\mathsf{h}_{m'}|\le\exp[\mathsf{B}\ell_3^4+6C_{\mathrm{face}}(R)]$. Hypothesis (e) with $m_1=m_2$
gives $\bar{\delta}\ge0$, hence $|\mathsf{h}_{m'}|\le Q_{\star,3}$ on $\mathbb{D}^{(3)}_{\mathfrak{r}}$.

\emph{Conclusion: the counterterm cancels.} Since $\mathsf{h}_{m''}(0)=1$, the principal logarithm
$\operatorname{Log}\mathsf{h}_{m''}$ is holomorphic near $0$, and near $0$
\begin{equation*}
\log Z_{m''}(w_z)=\log Z_{m''}(0)+\mathcal{E}''(z)+\operatorname{Log}\mathsf{h}_{m''}(z)
\end{equation*}
is a holomorphic logarithm of $Z_{m''}(w_z)$, since the exponential of the right side is
$Z_{m''}(0)e^{\mathcal{E}''(z)}\mathsf{h}_{m''}(z)=Z_{m''}(w_z)$. The connected three-point function is
$\partial_{z_1}\partial_{z_2}\partial_{z_3}\log Z_{m''}(w_z)$ at $z=0$; the constant
$\log Z_{m''}(0)$ does not contribute, so
\begin{equation}\label{9.98:eq-transfer-three-sources-9}
S^T_{3;K,m''}(F_1,F_2,F_3)=\partial_{z_1}\partial_{z_2}\partial_{z_3}\mathcal{E}''(0)
+\partial_{z_1}\partial_{z_2}\partial_{z_3}\operatorname{Log}\mathsf{h}_{m''}(0),
\end{equation}
with the same first term for $m''=m$ and $m''=m'$, because $\mathcal{E}''$ is formed on
$\mathbb{T}_{m'}$ in both cases. The Schwarz--Cauchy lemma states: if $\mathsf{h}$ is holomorphic on
the domain $\{\sum_{i=1}^n|z_i|\,\|f_i\|_{\sharp}<r/\mathfrak{r}\}$, $\mathsf{h}(0)=1$ and
$|\mathsf{h}|\le Q$ there, then
\begin{equation*}
|\partial_{z_1}\cdots\partial_{z_n}\operatorname{Log}\mathsf{h}(0)|
\le\log 2\cdot(2Q+2)^n\mathfrak{r}^n\Bigl(\frac{n}{r}\Bigr)^n\prod_{i=1}^n\|f_i\|_{\sharp}.
\end{equation*}
Its proof applies the Schwarz lemma on complex lines through $0$ to obtain
$|\operatorname{Log}\mathsf{h}|\le\log 2$ on the polydisc $\prod_i\{|z_i|\le\xi_i\}$, here with
$\xi_i:=\varrho/(2\mathfrak{r}(Q+1)\cdot3\cdot\|F_i\|_{\sharp})$, which lies in
$\mathbb{D}^{(3)}_{\mathfrak{r}}$ since $\sum_i\xi_i\|F_i\|_{\sharp}=\varrho/(2\mathfrak{r}(Q+1))
<\varrho/\mathfrak{r}$; Cauchy's coefficient estimate then bounds the mixed derivative by
$\log 2/(\xi_1\xi_2\xi_3)$. We apply the lemma with $n=3$, $r=\varrho$, $f_i=F_i$,
$Q=Q_{\star,3}$ and $\mathsf{h}=\mathsf{h}_{m''}$, whose hypotheses were established in the paragraph
on the common counterterm. Since
$\mathfrak{r}^3(3/\varrho)^3=(3\mathfrak{r}/\varrho)^3$, for each $m''\in\{m,m'\}$
\begin{equation*}
|\partial_{z_1}\partial_{z_2}\partial_{z_3}\operatorname{Log}\mathsf{h}_{m''}(0)|
\le\log 2\cdot(2Q_{\star,3}+2)^3\Bigl(\frac{3\mathfrak{r}}{\varrho}\Bigr)^3
\prod_{i=1}^3\|F_i\|_{\sharp}.
\end{equation*}
Subtracting \eqref{9.98:eq-transfer-three-sources-9} for $m''=m'$ from the same identity for $m''=m$,
the terms $\partial_{z_1}\partial_{z_2}\partial_{z_3}\mathcal{E}''(0)$ cancel, and the triangle
inequality gives \eqref{9.98:eq-transfer-three-sources-2} with
$C_{VT,3}=2\log 2\cdot(2Q_{\star,3}+2)^3(3\mathfrak{r}/\varrho)^3$.

\emph{Dependence of the constant.} $C_{VT,3}$ is a function of $\mathfrak{r}$, $\varrho$ and
$Q_{\star,3}$, and $Q_{\star,3}$ of $\mathsf{B}$, $\bar{\delta}$, $\ell_3$ and $C_{\mathrm{face}}(R)$.
We go through the quantities excluded in the statement one by one. The number of steps $K$:
by hypotheses (a) and (d), $E_{+}^{\sharp}$ and $C_{\mathrm{low}}$ do not depend on $K$, and
$C_{\mathrm{face}}(R)$ and $\bar{\delta}$ are majorants of series which, by hypotheses (b) and (e),
do not depend on $K$; the quantities that diverge with
$K$, namely $E_m$, $E_{m'}$, $\operatorname{Re}\mathcal{E}_{m'}(W_z)$ and
$\partial_{z_1}\partial_{z_2}\partial_{z_3}\mathcal{E}''(0)$, occur only in exact cancellations (the
vacuum energies in the density difference of hypothesis (e), the bulk counterterm in
\eqref{9.98:eq-transfer-three-sources-8}, and $\partial^3\mathcal{E}''(0)$ in the final subtraction).
The bare coupling $g_0$ enters no constant. The exponent $m$ enters only through reflection
positivity, the chessboard estimate and the trace inequality, which are used without any constant
depending on $m$, and through the mismatch of the densities per cell, which is bounded per cell. The
exponent $m'$ enters only through constants per site, and $\bar{\delta}$ is uniform by hypothesis
(H5) of Notation~\ref{2.2:not-hypotheses}. The position of the ball enters only through the one
translation of the placement paragraph, which changes neither Schwinger function. The distances,
overlaps and signs of the three supports are never used: the lemma on mirror-periodised weights, the
face-term lemma and the Schwarz--Cauchy lemma use only $\|F_i\|_{\sharp}$ and $R$, and the one object
that could depend on the separations, $\partial_{z_1}\partial_{z_2}\partial_{z_3}\mathcal{E}''(0)$,
is removed by the subtraction.

\emph{A two-sided bound from one-sided inputs.} Each
$|\partial_{z_1}\partial_{z_2}\partial_{z_3}\operatorname{Log}\mathsf{h}_{m''}(0)|$ is bounded in
modulus from the bound $\sup|\mathsf{h}_{m''}|\le Q_{\star,3}$ on a complex domain together with the
exact value $\mathsf{h}_{m''}(0)=1$. All the inequalities used --- the chessboard estimate, the trace
inequality, the upper stability bound of hypothesis (a), the lower bound on $Z_m(0)$ of hypothesis
(d) (which bounds $1/Z_m(0)$ from above), the one-sided vacuum bound of hypothesis (e) and the face
bounds --- point in the direction that an upper bound on $|\mathsf{h}_{m''}|$ requires; no lower bound
on $|Z_m(w_z)|$ at $z\ne0$ is used.
\end{proof}

\begin{lemma}[The trilinear averaging lemma]\label{9.98:lem-trilinear-averaging}
Let $\varepsilon$ be the lattice spacing, let $R>0$, and let $C\ge0$ be a constant. Let $D$ be a real
trilinear form on the
compactly supported real Lipschitz functions on $\mathbb{R}^4$ whose supports lie inside a ball of
radius less than $L^{m'}/2$ (in the applications the arguments are transported to tori, which this
restriction on the supports makes possible).
Assume:
\begin{enumerate}
\item[(a)] $D(F_1(\cdot-\tau),F_2(\cdot-\tau),F_3(\cdot-\tau))=D(F_1,F_2,F_3)$ for all
$\tau\in\varepsilon\mathbb{Z}^4$;
\item[(b)] whenever $\operatorname{supp}F_1\cup\operatorname{supp}F_2\cup\operatorname{supp}F_3$ is
contained in one closed ball of radius $R$,
\begin{equation}\label{9.98:eq-trilinear-averaging-a}
|D(F_1,F_2,F_3)|\le C\prod_{i=1}^{3}\|F_i\|_{\sharp}.
\end{equation}
\end{enumerate}
Let $y\in\mathbb{R}^4$, let $r$ satisfy $\varepsilon\le r\le R/8$, and let $f_1,f_2,f_3$ be real
Lipschitz functions with $\operatorname{supp}f_i\subset\overline{B}(y,r)$ for $i=1,2,3$. Then
\begin{equation}\label{9.98:eq-trilinear-averaging-1}
|D(f_1,f_2,f_3)|\le C\cdot\min\Bigl\{1,\Bigl(\frac{32r}{R}\Bigr)^{4}\Bigr\}\cdot
\prod_{i=1}^{3}\|f_i\|_{\sharp}.
\end{equation}
\end{lemma}

Here $\|\cdot\|_{\sharp}$ is the test-function norm of Definition~\ref{1.3:def-curvature-field},
$\|f\|_{\sharp}=\max\{\|f\|_\infty,40M\operatorname{Lip}f\}$, where
$\operatorname{Lip}f=\sup_{x\neq x'}|f(x)-f(x')|/|x-x'|$; $|\cdot|$ is the Euclidean norm on
$\mathbb{R}^4$, $|\cdot|_\infty$ the maximum norm, and $\overline{B}(z,\varsigma)$ the closed Euclidean
ball of centre $z$ and radius $\varsigma$.

\begin{proof}
The geometry and the gluing below are those of the bilinear averaging lemma; its segment argument
reads only that each summand vanishes on the boundary sphere of its ball, so it applies to
coefficients of either sign.
We fix the geometry first. Put $r':=\varepsilon\lceil r/\varepsilon\rceil$. Since
$r/\varepsilon\le\lceil r/\varepsilon\rceil<r/\varepsilon+1$ and $\varepsilon\le r$, we have
\begin{equation*}
r\le r'<r+\varepsilon\le 2r ,
\end{equation*}
and $r'$ is a positive integer multiple of $\varepsilon$. Put
\begin{equation*}
\Gamma:=\{\tau\in 4r'\mathbb{Z}^4:\ |\tau|_\infty\le R/4\}.
\end{equation*}
Since $4r'$ is an integer multiple of $\varepsilon$, $\Gamma\subset\varepsilon\mathbb{Z}^4$. Writing
$\tau=4r'a$ with $a\in\mathbb{Z}^4$, the condition $|\tau|_\infty\le R/4$ reads
$|a_\mu|\le R/(16r')$ for each $\mu$, and there are $2\lfloor R/(16r')\rfloor+1$ integers $a_\mu$ with
this property. Hence
\begin{equation*}
\mathcal{N}:=|\Gamma|=\bigl(2\lfloor R/(16r')\rfloor+1\bigr)^4 .
\end{equation*}
For every real $\xi\ge0$ one has $2\lfloor\xi\rfloor+1\ge\xi$: if $\xi<1$ the left side is at least
$1>\xi$, and if $\xi\ge1$ then
\begin{equation*}
2\lfloor\xi\rfloor+1>2(\xi-1)+1=2\xi-1\ge\xi .
\end{equation*}
Applied to $\xi=R/(16r')$, and using $r'<2r$,
this gives $2\lfloor R/(16r')\rfloor+1\ge R/(16r')>R/(32r)$; the same number is also at least $1$.
Therefore
\begin{equation}\label{9.98:eq-trilinear-averaging-2}
\mathcal{N}\ge\max\Bigl\{1,\Bigl(\frac{R}{32r}\Bigr)^{4}\Bigr\}.
\end{equation}
For $\tau\in\Gamma$ put $B_\tau:=\overline{B}(y+\tau,r)$. If $\tau\neq\tau'$ are in $\Gamma$, then
$|\tau-\tau'|\ge|\tau-\tau'|_\infty\ge4r'\ge4r$, so the distance between $B_\tau$ and $B_{\tau'}$ is
at least $|\tau-\tau'|-2r\ge2r>0$; the balls $B_\tau$, $\tau\in\Gamma$, are pairwise disjoint. For
$\tau\in\Gamma$ we have $|\tau|\le2|\tau|_\infty\le R/2$, so every $x\in B_\tau$ satisfies
$|x-y|\le R/2+r\le R/2+R/8<R$; thus all $B_\tau$ lie in $\overline{B}(y,R)$.

Next we glue translated copies with two independent families of signs. For
$(\sigma^{(1)},\sigma^{(2)})\in\{\pm1\}^{\Gamma}\times\{\pm1\}^{\Gamma}$ put
\begin{equation}\label{9.98:eq-trilinear-averaging-3}
\begin{aligned}
F_1&:=\sum_{\tau\in\Gamma}\sigma^{(1)}_\tau f_1(\cdot-\tau),\qquad
F_2:=\sum_{\tau\in\Gamma}\sigma^{(2)}_\tau f_2(\cdot-\tau),\\
F_3&:=\sum_{\tau\in\Gamma}\sigma^{(1)}_\tau\sigma^{(2)}_\tau f_3(\cdot-\tau).
\end{aligned}
\end{equation}
Each $F_i$ has the form $F=\sum_{\tau\in\Gamma}c_\tau f(\cdot-\tau)$ with $f=f_i$ and coefficients
$c_\tau\in\{\pm1\}$, and the summand $c_\tau f(\cdot-\tau)$ is supported in $B_\tau$. We show
$\|F\|_\infty=\|f\|_\infty$ and $\operatorname{Lip}F\le\operatorname{Lip}f$; only
$c_\tau\in\{\pm1\}$ and the disjointness of the $B_\tau$ are used. Since the $B_\tau$ are pairwise
disjoint, $F(x)=c_\tau f(x-\tau)$ for $x\in B_\tau$ and $F(x)=0$ for
$x\notin\bigcup_{\tau}B_\tau$. Hence $|F|$ takes on $B_\tau$ exactly the values of $|f|$ on
$\overline{B}(y,r)$, which contains $\operatorname{supp}f$, and $\|F\|_\infty=\|f\|_\infty$. Further,
$f$ vanishes on the sphere $\{|z-y|=r\}$: it is continuous and vanishes outside $\overline{B}(y,r)$, and
every point of the sphere is a limit of points outside the ball. Consequently $F$ vanishes on the
boundary sphere of every $B_\tau$. Let $x,x'\in\mathbb{R}^4$, and consider the segment
$x+t(x'-x)$, $t\in[0,1]$. There are four cases.
\begin{itemize}
\item[--] Neither point lies in $\bigcup_\tau B_\tau$. Then $F(x)=F(x')=0$.
\item[--] Both lie in the same $B_\tau$. Then
$|F(x)-F(x')|=|f(x-\tau)-f(x'-\tau)|\le\operatorname{Lip}f\,|x-x'|$.
\item[--] One point, say $x$, lies in $B_\tau$ and $x'$ lies in no ball. The function
$t\mapsto|x+t(x'-x)-y-\tau|$ is continuous, at most $r$ at $t=0$ and larger than $r$ at $t=1$, so some
point $z$ of the segment lies on the boundary sphere of $B_\tau$. Then $F(z)=F(x')=0$ and
\begin{equation*}
\begin{aligned}
|F(x)-F(x')|&=|F(x)-F(z)|=|f(x-\tau)-f(z-\tau)|\\
&\le\operatorname{Lip}f\,|x-z|\le\operatorname{Lip}f\,|x-x'|.
\end{aligned}
\end{equation*}
\item[--] $x\in B_\tau$ and $x'\in B_{\tau'}$ with $\tau\neq\tau'$. Let $t_1$ be the largest
$t\in[0,1]$ with $x+t(x'-x)\in B_\tau$ and $t_2$ the smallest with $x+t(x'-x)\in B_{\tau'}$; they exist
since the balls are closed, and $z:=x+t_1(x'-x)$, $z':=x+t_2(x'-x)$ lie on the boundary spheres of
$B_\tau$ and $B_{\tau'}$ respectively, because $x'\notin B_\tau$ and $x\notin B_{\tau'}$. If
$t_2<t_1$, then $z'$ would lie between $x$ and $z$, both in the convex set $B_\tau$, hence
$z'\in B_\tau\cap B_{\tau'}=\emptyset$; so $t_1\le t_2$. Using $F(z)=F(z')=0$,
\begin{equation*}
\begin{aligned}
|F(x)-F(x')|&\le|F(x)-F(z)|+|F(z')-F(x')|\\
&\le\operatorname{Lip}f\,(|x-z|+|z'-x'|)
=\operatorname{Lip}f\,(t_1+1-t_2)|x-x'|\le\operatorname{Lip}f\,|x-x'|.
\end{aligned}
\end{equation*}
\end{itemize}
Thus $\operatorname{Lip}F\le\operatorname{Lip}f$, and
\begin{equation*}
\begin{aligned}
\|F_i\|_{\sharp}&=\max\{\|F_i\|_\infty,40M\operatorname{Lip}F_i\}\\
&\le\max\{\|f_i\|_\infty,40M\operatorname{Lip}f_i\}=\|f_i\|_{\sharp},\qquad i=1,2,3,
\end{aligned}
\end{equation*}
for every choice of signs and whatever the number $\mathcal{N}$ of copies: the glued function has at
most the norm of one copy. This norm is not normalised to be of order one, and no normalisation is
needed, since
\eqref{9.98:eq-trilinear-averaging-a} is homogeneous of degree one in each argument. Moreover each $F_i$ is
a compactly supported real Lipschitz function with support in
$\bigcup_\tau B_\tau\subset\overline{B}(y,R)$.

We now average over the signs. For a function $\Phi$ of $(\sigma^{(1)},\sigma^{(2)})$ write
\begin{equation*}
\langle\Phi\rangle:=2^{-2\mathcal{N}}
\sum_{(\sigma^{(1)},\sigma^{(2)})\in\{\pm1\}^{\Gamma}\times\{\pm1\}^{\Gamma}}
\Phi(\sigma^{(1)},\sigma^{(2)}),
\end{equation*}
a finite average over the $2^{2\mathcal{N}}$ sign choices. By
trilinearity of $D$ and the definitions of $F_1,F_2,F_3$ in
\eqref{9.98:eq-trilinear-averaging-3},
\begin{equation*}
D(F_1,F_2,F_3)=\sum_{\tau,\tau',\tau''\in\Gamma}
\sigma^{(1)}_\tau\,\sigma^{(2)}_{\tau'}\,\sigma^{(1)}_{\tau''}\sigma^{(2)}_{\tau''}\,
D\bigl(f_1(\cdot-\tau),f_2(\cdot-\tau'),f_3(\cdot-\tau'')\bigr).
\end{equation*}
The coefficient is the product of $\sigma^{(1)}_\tau\sigma^{(1)}_{\tau''}$, which depends on
$\sigma^{(1)}$ only, and $\sigma^{(2)}_{\tau'}\sigma^{(2)}_{\tau''}$, which depends on $\sigma^{(2)}$
only; since the sign choices range over the product set
$\{\pm1\}^\Gamma\times\{\pm1\}^\Gamma$, its average is the product of the two averages over a single
family. If $\tau=\tau''$, then $\sigma^{(1)}_\tau\sigma^{(1)}_{\tau''}=1$. If $\tau\neq\tau''$, the
involution of $\{\pm1\}^\Gamma$ that reverses $\sigma^{(1)}_\tau$ alone is a bijection that reverses
the sign of $\sigma^{(1)}_\tau\sigma^{(1)}_{\tau''}$ and leaves the $\sigma^{(2)}$-factor unchanged, so
the average of $\sigma^{(1)}_\tau\sigma^{(1)}_{\tau''}$ equals its own negative and vanishes. In the
same way, reversing $\sigma^{(2)}_{\tau'}$ alone shows that the average of
$\sigma^{(2)}_{\tau'}\sigma^{(2)}_{\tau''}$ is $\delta_{\tau'\tau''}$, where $\delta$ denotes the
Kronecker symbol on $\Gamma$, equal to $1$ when its two indices coincide and to $0$ otherwise. Hence
\begin{equation*}
\bigl\langle\sigma^{(1)}_\tau\sigma^{(2)}_{\tau'}\sigma^{(1)}_{\tau''}\sigma^{(2)}_{\tau''}\bigr\rangle
=\delta_{\tau\tau''}\,\delta_{\tau'\tau''},
\end{equation*}
and, by linearity of the average,
\begin{equation}\label{9.98:eq-trilinear-averaging-4}
\langle D(F_1,F_2,F_3)\rangle
=\sum_{\tau\in\Gamma}D\bigl(f_1(\cdot-\tau),f_2(\cdot-\tau),f_3(\cdot-\tau)\bigr)
=\mathcal{N}\,D(f_1,f_2,f_3).
\end{equation}
The last equality is hypothesis (a) applied to $(f_1,f_2,f_3)$ with $\tau\in\Gamma$, which is
legitimate because $\Gamma\subset\varepsilon\mathbb{Z}^4$. Thus the $\mathcal{N}^3-\mathcal{N}$ cross
terms, those with $(\tau,\tau',\tau'')$ not all equal, cancel exactly, and translation invariance makes
each of the $\mathcal{N}$ diagonal terms equal to the value on the given triple.

Finally, for every sign choice the triple $(F_1,F_2,F_3)$ consists of real Lipschitz functions with
supports in $\overline{B}(y,R)$, a closed ball of radius $R$, and $\|F_i\|_{\sharp}\le\|f_i\|_{\sharp}$.
Hypothesis (b) therefore gives, by \eqref{9.98:eq-trilinear-averaging-a},
\begin{equation*}
|D(F_1,F_2,F_3)|\le C\prod_{i=1}^{3}\|F_i\|_{\sharp}\le C\prod_{i=1}^{3}\|f_i\|_{\sharp}.
\end{equation*}
Averaging over the signs and using \eqref{9.98:eq-trilinear-averaging-4},
\begin{equation*}
\begin{aligned}
\mathcal{N}\,|D(f_1,f_2,f_3)|&=|\langle D(F_1,F_2,F_3)\rangle|\\
&\le\langle|D(F_1,F_2,F_3)|\rangle\le C\prod_{i=1}^{3}\|f_i\|_{\sharp}.
\end{aligned}
\end{equation*}
By \eqref{9.98:eq-trilinear-averaging-2}, $\mathcal{N}^{-1}\le\min\{1,(32r/R)^4\}$, and
\eqref{9.98:eq-trilinear-averaging-1} follows.
\end{proof}

\begin{remark}
The construction extends to an $n$-linear form with the analogous hypotheses (a) and (b), with $n-1$
independent sign families $\sigma^{(1)},\dots,\sigma^{(n-1)}\in\{\pm1\}^\Gamma$: one puts
$F_i:=\sum_{\tau\in\Gamma}\sigma^{(i)}_\tau f_i(\cdot-\tau)$ for $i<n$ and
$F_n:=\sum_{\tau\in\Gamma}\bigl(\prod_{i<n}\sigma^{(i)}_\tau\bigr)f_n(\cdot-\tau)$. By the same
factorisation and the same involutions, the average of
$\prod_{i<n}\sigma^{(i)}_{\tau_i}\sigma^{(i)}_{\tau_n}$ is $\prod_{i<n}\delta_{\tau_i\tau_n}$. The
cases $n=2$, which is the bilinear averaging lemma, and $n=3$ are the only ones used.

No bound is applied in the proof to an individual cross triple
$(f_1(\cdot-\tau),f_2(\cdot-\tau'),f_3(\cdot-\tau''))$. Such triples enter only as summands in the
expansion of the single value $D(F_1,F_2,F_3)$, which hypothesis (b) controls without any condition on
the separation of the supports, and they are then removed exactly by the sign average. The glued
functions $F_i$ are signed; when $r$ is comparable to a separation $\mathsf{d}$ and $R$ is fixed, each is
a sum of $\mathcal{N}$ translated copies, a number of order $\mathsf{d}^{-4}$. This is why (b) is
required for all real Lipschitz triples supported in one closed ball of radius $R$, and why it matters
that $\|\cdot\|_{\sharp}$ involves only the supremum and the Lipschitz constant, for which gluing
costs nothing. The same universality of (b), for all real Lipschitz arguments in one closed ball with
no separation condition, is the form that enters the per-ball volume-free bound for $n$ balls.
\end{remark}

\begin{corollary}[The separation gain for three sources]\label{9.98:cor-separation-gain-three}
Let $g^{\flat}$, $g_0$, $K$ and the torus exponents $m\ge m'$ be as in the volume-transfer bound for
three sources (Proposition~\ref{9.98:prop-transfer-three-sources}), and take there $R:=1$. Write
$C_{VT,3}(g^{\flat})$ for the constant $C_{VT,3}$ of that proposition at $R=1$, and put
\begin{equation}\label{9.98:eq-separation-gain-three-1}
C_{T,3}(g^{\flat}) := 160^4\Bigl(\frac{40M}{\vartheta}\Bigr)^{3} C_{VT,3}(g^{\flat}).
\end{equation}
Let $0<\mathsf{d}\le 1/40$ and assume that the lattice spacing satisfies $\varepsilon=L^{-K}\le\mathsf{d}$.
Let $\mathfrak{P}_3(\mathsf{d};\vartheta)$ be the set of triples $(f_1,f_2,f_3)$ of real Lipschitz
functions on $\mathbb{R}^4$, of any signs, whose supports lie in one closed ball
$\overline{B}(y,5\mathsf{d})$ and which satisfy $\operatorname{Lip} f_i\le\|f_i\|_\infty/(\vartheta\mathsf{d})$
for $i=1,2,3$. Pairwise disjointness of the supports, which the observables of
Theorem~\ref{3.1:thm-main} carry, is not required and is not used by the bound. Then for every
$(f_1,f_2,f_3)\in\mathfrak{P}_3(\mathsf{d};\vartheta)$, the $f_i$ being transported to $\mathbb{T}_m$
and to $\mathbb{T}_{m'}$ as in Proposition~\ref{9.98:prop-transfer-three-sources},
\begin{equation}\label{9.98:eq-separation-gain-three-2}
\bigl|S^T_{3;K,m}(f_1,f_2,f_3)-S^T_{3;K,m'}(f_1,f_2,f_3)\bigr|
\le C_{T,3}(g^{\flat})\,\mathsf{d}\prod_{i=1}^{3}\|f_i\|_\infty ,
\end{equation}
uniformly in $K$, $g_0$, $y$ and in all $m\ge m'\ge m_{\star}(g^{\flat})$.
\end{corollary}

\begin{proof}
Put $D:=S^T_{3;K,m}-S^T_{3;K,m'}$. We regard $D$ as a form on triples of real Lipschitz functions
supported in balls of radius at most $1$, each such function being transported to both tori
$\mathbb{T}_m$ and $\mathbb{T}_{m'}$ as in Proposition~\ref{9.98:prop-transfer-three-sources}. We
check the hypotheses of the trilinear averaging lemma (Lemma~\ref{9.98:lem-trilinear-averaging})
for $D$ with $R=1$.

\emph{Trilinearity.} For each torus exponent $m''\in\{m,m'\}$, the connected Schwinger function
$S^T_{3;K,m''}(F_1,F_2,F_3)$ is the third joint cumulant of the random variables
$\mathcal{O}(F_1),\mathcal{O}(F_2),\mathcal{O}(F_3)$ under the lattice measure of
Definition~\ref{1.1:def-wilson-action} on $\mathbb{T}_{m''}$.
A joint cumulant of three random variables is linear in each of them, and
$\mathcal{O}(F)=\sum_p F(x(p))[1-\operatorname{Re}\operatorname{tr}U(\partial p)]$
(Definition~\ref{1.3:def-curvature-field}) is linear in $F$.
Hence each $S^T_{3;K,m''}$ is trilinear in $(F_1,F_2,F_3)$; it is real for real $F_i$, since then each
$\mathcal{O}(F_i)$ is real. The difference $D$ of two real trilinear forms is a real trilinear form.

\emph{Invariance under $\varepsilon\mathbb{Z}^4$.} Let $\tau\in\varepsilon\mathbb{Z}^4$. Transport to a
torus commutes with translation by $\tau$. The Haar measure on $\mathbb{T}_{m''}$ and the Wilson
weight are each invariant under the lattice isometry $p\mapsto p+\tau$ of the bare lattice
(Lemma~\ref{1.1:lem-invariance}). Therefore, for $m''\in\{m,m'\}$,
\begin{equation*}
S^T_{3;K,m''}\bigl(F_1(\cdot-\tau),F_2(\cdot-\tau),F_3(\cdot-\tau)\bigr)=S^T_{3;K,m''}(F_1,F_2,F_3).
\end{equation*}
Hence $D$ has property~(a) of Lemma~\ref{9.98:lem-trilinear-averaging}.

\emph{The bound on one ball of radius one.} By
Proposition~\ref{9.98:prop-transfer-three-sources} with $R=1$, the inequality
\begin{equation*}
|D(F_1,F_2,F_3)|\le C_{VT,3}(g^{\flat})\prod_{i=1}^{3}\|F_i\|_{\sharp}
\end{equation*}
holds for all real Lipschitz triples whose supports lie in one closed ball of radius $1$. This is
\eqref{9.98:eq-trilinear-averaging-a} with $C=C_{VT,3}(g^{\flat})$ and $R=1$.

\emph{Averaging.} We apply Lemma~\ref{9.98:lem-trilinear-averaging} to $D$ and to the triple
$(f_1,f_2,f_3)$. Since $(f_1,f_2,f_3)\in\mathfrak{P}_3(\mathsf{d};\vartheta)$, the supports of the
$f_i$ lie in $\overline{B}(y,r)$ with $r:=5\mathsf{d}$; moreover $r\le 1/8=R/8$ because
$\mathsf{d}\le 1/40$, and $\varepsilon\le\mathsf{d}\le r$. The
lemma gives
\begin{equation*}
|D(f_1,f_2,f_3)|\le C_{VT,3}(g^{\flat})\,(32\cdot 5\mathsf{d})^4\prod_{i=1}^{3}\|f_i\|_{\sharp}
= C_{VT,3}(g^{\flat})\,(160\mathsf{d})^4\prod_{i=1}^{3}\|f_i\|_{\sharp}.
\end{equation*}
In the range $1/160<\mathsf{d}\le 1/40$ the family of translates of that lemma may consist of a
single element, and the factor $\min\{1,\cdot\}$ of that lemma may then equal $1$. Since
$160\mathsf{d}>1$ in this range, that factor is still at most $(160\mathsf{d})^4$, so the displayed
bound holds in this range as well.

\emph{From $\|\cdot\|_{\sharp}$ to $\|\cdot\|_\infty$.} We have $M\ge 1$, since $M$ is a power of $L$,
while $\vartheta\le 1$ and $\mathsf{d}\le 1/40$. Hence $40M\ge 1\ge\vartheta\mathsf{d}$, that is,
$40M/(\vartheta\mathsf{d})\ge 1$. Using $\operatorname{Lip} f_i\le\|f_i\|_\infty/(\vartheta\mathsf{d})$, we
obtain for $i=1,2,3$
\begin{equation}\label{9.98:eq-separation-gain-three-3}
\begin{aligned}
\|f_i\|_{\sharp}
&=\max\{\|f_i\|_\infty,\,40M\operatorname{Lip} f_i\}
\le\max\Bigl\{1,\frac{40M}{\vartheta\mathsf{d}}\Bigr\}\|f_i\|_\infty\\
&=\frac{40M}{\vartheta\mathsf{d}}\,\|f_i\|_\infty .
\end{aligned}
\end{equation}

\emph{Conclusion.} We insert \eqref{9.98:eq-separation-gain-three-3} for $i=1,2,3$ into the averaged
bound and use \eqref{9.98:eq-separation-gain-three-1}:
\begin{equation*}
|D(f_1,f_2,f_3)|
\le 160^4\mathsf{d}^4\,(40M)^3(\vartheta\mathsf{d})^{-3}\,C_{VT,3}(g^{\flat})\prod_{i=1}^{3}\|f_i\|_\infty
= C_{T,3}(g^{\flat})\,\mathsf{d}\prod_{i=1}^{3}\|f_i\|_\infty .
\end{equation*}
This is \eqref{9.98:eq-separation-gain-three-2}. None of the constants used depends on $K$, $g_0$,
$y$, $m$ or $m'$ within the stated ranges.
\end{proof}

The powers of $\mathsf{d}$ balance as follows. The three $\|\cdot\|_{\sharp}$-norms cost
$\mathsf{d}^{-3}$. The averaging over at least $\max\{(160\mathsf{d})^{-4},1\}$ translates gains
$\mathsf{d}^4$. The net result is the factor $\mathsf{d}^1$, which is the gain of the averaging
argument for three sources. The same count gives $\mathsf{d}^2$ for two sources and the exponent
zero for four sources, which is the marginal case. For five or more sources it gives a loss. For
the averaging argument the exponent $4-n$ for $n$ sources cannot be improved: there are explicit
extremal $n$-linear forms for which it is attained. This is not argued here and is not used.

The bound \eqref{9.98:eq-separation-gain-three-2} is absolute. No relative form of it is asserted.
No lower bound on the connected three-point function is proved in this book. A comparison of the
discrepancy with the size of that function would therefore be relative to its tree-level size, of
order the sixth power of the coupling, for which no lower bound is proved.

\begin{proposition}[The transfer inequality at $\gamma_{\star}$]\label{9.98:prop-transfer-inequality}
Let $g$, $\hat w$ and the tuned family be as in Lemma~\ref{9.98:lem-tuned-member}; this is the
\emph{tuned frame}. Let $n\in\{2,3\}$ be the number of sources and let $\mathsf{d}>0$. Let $\nu\ge
n_{\natural}$ be a tuned depth with $L^{-\nu}\le\mathsf{d}$, and let $x_0:=a_\nu(g)-\hat w$ be the
bare inverse coupling of the depth-$\nu$ member, $g_0:=x_0^{-1/2}$. Put
\begin{equation*}
K':=K(g_0,\gamma_{\star}),\qquad \mathsf{s}':=K'-\nu,\qquad \mathsf{d}':=L^{-\mathsf{s}'}\mathsf{d}.
\end{equation*}
Then for all $m_1,m_2\ge m_{\natural}(g)$, with $m_{\natural}(g)$ as in
\eqref{9.98:eq-prolongation-depths-c}, and every $\vec f\in\mathfrak{P}_n(\mathsf{d};\vartheta)$ with
$\mathsf{d}\le L^{\underline{s}'(g)}/40$, writing $\Pi:=\prod_{i}\|f_i\|_{\infty}$,
\begin{equation}\label{9.98:eq-transfer-inequality-a}
\begin{aligned}
\bigl|S^{\nu}_{\mathbb{T}_{m_1}}(x_0;\vec f\,)-S^{\nu}_{\mathbb{T}_{m_2}}(x_0;\vec f\,)\bigr|
&=\bigl|S^T_{n;K',m_1-\mathsf{s}'}[g_0]\bigl(\vec f(L^{\mathsf{s}'}\cdot)\bigr)
-S^T_{n;K',m_2-\mathsf{s}'}[g_0]\bigl(\vec f(L^{\mathsf{s}'}\cdot)\bigr)\bigr|\\
&\le C^{\star}_{T,n}\,(\mathsf{d}')^{4-n}\,\Pi
=C^{\star}_{T,n}\,L^{-(4-n)\mathsf{s}'}\,\mathsf{d}^{4-n}\,\Pi\le\tau_n ,
\end{aligned}
\end{equation}
where $\tau_n$ is the quantity defined in \eqref{9.98:eq-short-distance-uniqueness-a},
$C^{\star}_{T,2}:=C_T^{\star}=C_T(\gamma_{\star})$ and $C^{\star}_{T,3}:=C_{T,3}(\gamma_{\star})$. The
same holds in the
\emph{first-passage frame}: for $g\le\gamma_{\star}$ and $g_0\in\,]0,g_-(g)]$ as in
Lemma~\ref{9.98:lem-run-prolongation}, with $K:=K(g_0,g)$ in place of $\nu$ and the hypothesis
$L^{-K}\le\mathsf{d}$ in place of $L^{-\nu}\le\mathsf{d}$, $S^T_{n;K,m_i}[g_0](\vec f\,)$ in place of
$S^{\nu}_{\mathbb{T}_{m_i}}(x_0;\vec f\,)$, $\mathsf{s}:=K'-K$ in place of $\mathsf{s}'$, and the
hypothesis $m_1,m_2\ge m_{\star}+\overline{s}(g)$ in place of $m_1,m_2\ge m_{\natural}(g)$.
\end{proposition}

\begin{proof}
\emph{The first equality.} In the tuned frame, Lemma~\ref{9.98:lem-tuned-member}, display
\eqref{9.98:eq-tuned-member-a}, reads the depth-$\nu$ member on $\mathbb{T}_{m_i}$ under the
naming of the reference length at $\gamma_{\star}$: it is the bare Wilson measure of
Theorem~\ref{3.1:thm-main} at reference coupling $\gamma_{\star}$, on the torus
$\mathbb{T}_{m_i-\mathsf{s}'}$, at depth $K'$, the test functions being dilated to
$\vec f(L^{\mathsf{s}'}\cdot)$; for $m_i\ge m_{\natural}(g)$ the same lemma gives
$m_i-\mathsf{s}'\ge m_{\star}$. In the first-passage frame, Lemma~\ref{9.98:lem-run-prolongation}(a)
gives $\mathsf{s}\le\overline{s}(g)$, so that $m_i\ge m_{\star}+\overline{s}(g)\ge
m_{\star}+\mathsf{s}$; this is the hypothesis of Lemma~\ref{9.98:lem-run-prolongation}(b), which
yields $S^T_{n;K,m_i}[g_0](\vec f\,)=S^T_{n;K',m_i-\mathsf{s}}[g_0](\vec f(L^{\mathsf{s}}\cdot))$,
and $m_i-\mathsf{s}\ge m_{\star}$. Only the upper bound on the prolongation is used here.

\emph{The hypotheses of the separation gains at reference coupling $\gamma_{\star}$.} For
$n=2$ we apply the separation gain for two sources, for $n=3$ the separation gain for three
sources (Corollary~\ref{9.98:cor-separation-gain-three}), each with reference coupling
$g^{\flat}:=\gamma_{\star}$ and $R:=1$. We check their hypotheses in turn; we write the tuned
frame, the first-passage frame being the same with $\mathsf{s}$ for $\mathsf{s}'$.
\begin{itemize}
\item The reference coupling satisfies $\gamma_{\star}\le\gamma_J$, by the choice of
$\gamma_{\star}:=\gamma_W^{\mathbb{T}}$ in the display that fixes the strong reference coupling.
\item $g_0\le g_-(\gamma_{\star})$, by Lemma~\ref{9.98:lem-run-prolongation}(a).
\item The depth is $K'=K(g_0,\gamma_{\star})$, the first-passage depth at reference coupling
$\gamma_{\star}$, by definition.
\item Both torus exponents satisfy $m_i-\mathsf{s}'\ge m_{\star}$, as shown above, and $m_{\star}$ is
the smallest torus exponent of the family at reference coupling $\gamma_{\star}$
(Definition~\ref{9.98:def-prolongation-depths}). The bound is stated for an ordered pair of torus
exponents, the first at least the second; we apply it with the larger of $m_1-\mathsf{s}'$,
$m_2-\mathsf{s}'$ as the first and the smaller as the second, which is
legitimate because the absolute value of the difference is symmetric in the two tori.
\item $0<\mathsf{d}'\le1/40$: from $\mathsf{d}\le L^{\underline{s}'(g)}/40$ and
$\mathsf{s}'\ge\underline{s}'(g)$ (the lower bound on the prolongation in the tuned frame,
Lemma~\ref{9.98:lem-tuned-member}),
\begin{equation*}
\mathsf{d}'=L^{-\mathsf{s}'}\mathsf{d}\le L^{\underline{s}'(g)-\mathsf{s}'}/40\le 1/40 .
\end{equation*}
In the first-passage frame Lemma~\ref{9.98:lem-run-prolongation}(a) gives $\mathsf{s}\ge
\underline{s}(g)$, and $\underline{s}(g)\ge\underline{s}'(g)$ since the two floors
\eqref{9.98:eq-prolongation-depths-a} differ by the non-negative integer $j_U-1$.
\item The spacing $\varepsilon'=L^{-K'}$ of the bare lattice in $\gamma_{\star}$-reference units
satisfies $\varepsilon'\le\mathsf{d}'$: since $K'=\nu+\mathsf{s}'$ and $L^{-\nu}\le\mathsf{d}$,
\begin{equation*}
\varepsilon'=L^{-\mathsf{s}'}L^{-\nu}\le L^{-\mathsf{s}'}\mathsf{d}=\mathsf{d}' .
\end{equation*}
\item The dilated tuple lies in $\mathfrak{P}_n(\mathsf{d}';\vartheta)$. Let $\vec f\in
\mathfrak{P}_n(\mathsf{d};\vartheta)$ have its supports in $\overline{B}(y,5\mathsf{d})$. Each
$f_i(L^{\mathsf{s}'}\cdot)$ is real and Lipschitz, with the sign of $f_i$, has its support in the
one closed ball $\overline{B}(L^{-\mathsf{s}'}y,5\mathsf{d}')$, has the same sup norm
$\|f_i\|_{\infty}$, and
\begin{equation*}
\operatorname{Lip}f_i(L^{\mathsf{s}'}\cdot)=L^{\mathsf{s}'}\operatorname{Lip}f_i
\le L^{\mathsf{s}'}\frac{\|f_i\|_{\infty}}{\vartheta\mathsf{d}}=\frac{\|f_i\|_{\infty}}{\vartheta
\mathsf{d}'} .
\end{equation*}
Neither bound requires the supports to be separated or disjoint, nor the $f_i$ to have a sign.
\end{itemize}

\emph{The bound.} The separation gain for two sources (for $n=2$), respectively
Corollary~\ref{9.98:cor-separation-gain-three} (for $n=3$), at reference coupling $\gamma_{\star}$,
applied to the
tori $\mathbb{T}_{m_1-\mathsf{s}'}$, $\mathbb{T}_{m_2-\mathsf{s}'}$ at depth $K'$ and to the tuple
$\vec f(L^{\mathsf{s}'}\cdot)$ of scale $\mathsf{d}'$, gives the first inequality of
\eqref{9.98:eq-transfer-inequality-a}, the product of the sup norms of the dilated tuple being
$\Pi$. The next equality is $(\mathsf{d}')^{4-n}=L^{-(4-n)\mathsf{s}'}\mathsf{d}^{4-n}$. The last
inequality holds because $4-n\ge1$, $L>1$ and $\mathsf{s}'\ge\underline{s}'(g)$, so that
$L^{-(4-n)\mathsf{s}'}\le L^{-(4-n)\underline{s}'(g)}$, and hence the third member of
\eqref{9.98:eq-transfer-inequality-a} is at most $\tau_n$ by
\eqref{9.98:eq-short-distance-uniqueness-a}. Only the lower bound on the
prolongation enters this gain, and only the upper bound enters the floor; no subsequence along
which $\mathsf{s}'$ is constant is used.

\emph{The constant.} $C^{\star}_{T,2}=C_T(\gamma_{\star})$ and $C^{\star}_{T,3}=C_{T,3}(\gamma_{\star})$
depend on $(L,\mathfrak{p},\gamma_{\star},\varrho,\vartheta)$ only; they depend neither on $K'$ nor
on $g_0$, $m_1$, $m_2$, $\mathsf{d}'$, the position of the ball, the signs of the $f_i$ or the
overlaps of their supports. The one quantity in the proof of the volume-transfer theorem that
depends on the separation, the derivative $\partial^n\mathcal{E}''(0)$ of the re-centred source
counterterm, is the same number on both tori and cancels in the difference. The volume of the
larger torus enters no constant; the smaller torus need only contain $\mathbb{T}_{\star}$ in the
$\gamma_{\star}$-frame, and this, and nothing else, is the content of the floor $m\ge
m_{\natural}(g)$.
\end{proof}

Clause~(v) of Theorem~\ref{3.1:thm-main} is not used in Proposition~\ref{9.98:prop-transfer-inequality}, which is
a statement at finite cutoff about bare measures; it is used, on one torus only, in
Theorem~\ref{9.98:thm-short-distance-uniqueness}.

\begin{lemma}[The cutoff limit on the reference torus]\label{9.98:lem-reference-torus-limit}
Let $g\le\gamma_U$ be a reference coupling, and put $m_{\natural}:=m_{\natural}(g)$ as in
\eqref{9.98:eq-prolongation-depths-c} of Definition~\ref{9.98:def-prolongation-depths}.
As in clause~(v) of Theorem~\ref{3.1:thm-main}, $a_n=a_n(g)$ denotes the tuned bare inverse
coupling at depth $n$,
and $S^{n}_{\mathbb{T}_m}(x_0;\vec f\,)$ denotes the connected Schwinger function, smeared with the
tuple $\vec f$, of the depth-$n$ member of the construction on $\mathbb{T}_m$ at bare inverse
coupling $x_0$. Let $\vec f$ be a separated Lipschitz pair or triple of test functions, as in
clause~(ii-$\mathbb{T}$-a) of Theorem~\ref{3.1:thm-main}. Then the following hold.
\begin{enumerate}
\item[(a)] For every real $\hat w$ with $|\hat w|\le\varrho_2/2$ the limit
\begin{equation}\label{9.98:eq-reference-torus-limit-a}
\lim_{n\to\infty} S^{n}_{\mathbb{T}_{m_{\natural}}}(a_n-\hat w;\vec f\,)
=S^{(\hat w)}_{\mathbb{T}_{m_{\natural}}}(\vec f\,)
=:\mathfrak{S}_{m_{\natural}}[\hat x](\vec f\,)
\end{equation}
exists along the full sequence $n\to\infty$, without passing to subsequences; the name
$\mathfrak{S}_{m_{\natural}}[\hat x](\vec f\,)$ denotes this limit at the given value of
$\hat w$, and is used with the same meaning in (b). The map
$\hat w\mapsto S^{(\hat w)}_{\mathbb{T}_{m_{\natural}}}(\vec f\,)$ is real-analytic on
$\mathopen{]}-\varrho_2/2,\varrho_2/2\mathclose{[}$ and extends holomorphically to the strip
$|\operatorname{Im}\hat w|<\upsilon_w(m_{\natural})$.
\item[(b)] For every real sequence $(x^{(n)})_n$ with $|x^{(n)}-a_n|\le\varrho_2/2$ for all $n$,
\begin{equation}\label{9.98:eq-reference-torus-limit-1}
S^{n}_{\mathbb{T}_{m_{\natural}}}(x^{(n)};\vec f\,)
-S^{(a_n-x^{(n)})}_{\mathbb{T}_{m_{\natural}}}(\vec f\,)\longrightarrow 0
\qquad(n\to\infty);
\end{equation}
and if moreover $a_n-x^{(n)}\to\hat w$, then
$S^{n}_{\mathbb{T}_{m_{\natural}}}(x^{(n)};\vec f\,)\to\mathfrak{S}_{m_{\natural}}[\hat x](\vec f\,)$.
\item[(c)] The moduli of the convergences in (a) and (b), namely the rate constant, the strip
width $\upsilon_w(m_{\natural})$ and the floor $n_0(N_w)$, depend on the torus only through
$N_w=2L^{m_{\natural}}$, by way of
\begin{equation}\label{9.98:eq-reference-torus-limit-2}
\mathsf{Q}_w=\exp\bigl[(E_{+}^{\sharp}+C_{\mathrm{low}})N_w^4\bigr];
\end{equation}
they are numbers depending on $g$ alone, and in (a) and (b) they multiply a term which the
limit $n\to\infty$ removes.
\end{enumerate}
\end{lemma}

\begin{proof}
\emph{Admissibility of the torus.} By Remark~\ref{9.98:rem-floor-members},
display~\eqref{9.98:eq-floor-members-a}, the floor of
\eqref{9.98:eq-prolongation-depths-c} satisfies
\begin{equation*}
m_{\natural}(g)\ge m_{\mathbb{T}}(g_-(g)).
\end{equation*}
Clauses (ii-$\mathbb{T}$-a), (v-b) and (v-c)(1) of Theorem~\ref{3.1:thm-main} hold for every
reference coupling $g\le\gamma_U$ on every torus $\mathbb{T}_m$ with $m\ge m_{\mathbb{T}}(g_-(g))$.
Since $g\le\gamma_U$ by hypothesis and $m_{\natural}$ satisfies the last display, all three clauses
apply on the single torus $\mathbb{T}_{m_{\natural}}$ at reference coupling $g$. Clause~(v) of
Theorem~\ref{3.1:thm-main} is the uniqueness theorem; clause~(ii-$\mathbb{T}$-a) is used
for pairs and triples only.

\emph{Proof of (a).} Fix a real $\hat w$ with $|\hat w|\le\varrho_2/2$ and a pair or triple
$\vec f$ as in the statement. Clause~(v-b), read on $\mathbb{T}_{m_{\natural}}$, asserts that
$S^{n}_{\mathbb{T}_{m_{\natural}}}(a_n-\hat w;\vec f\,)$ converges as $n\to\infty$, along the
whole sequence, to a limit $S^{(\hat w)}_{\mathbb{T}_{m_{\natural}}}(\vec f\,)$; this is
\eqref{9.98:eq-reference-torus-limit-a}, in which the symbol
$\mathfrak{S}_{m_{\natural}}[\hat x](\vec f\,)$ is introduced as a name for the limit. The same
clause asserts that $\hat w\mapsto S^{(\hat w)}_{\mathbb{T}_{m_{\natural}}}(\vec f\,)$ is
real-analytic on $\mathopen{]}-\varrho_2/2,\varrho_2/2\mathclose{[}$ and holomorphic on the strip
$|\operatorname{Im}\hat w|<\upsilon_w(m_{\natural})$.

\emph{Proof of (b).} Let $(x^{(n)})_n$ be real with $|x^{(n)}-a_n|\le\varrho_2/2$. Then
$a_n-x^{(n)}$ lies in the closed window $[-\varrho_2/2,\varrho_2/2]$ for every $n$, and clause
(v-c)(1), read on $\mathbb{T}_{m_{\natural}}$, is exactly the convergence
\eqref{9.98:eq-reference-torus-limit-1}. Suppose in addition that $a_n-x^{(n)}\to\hat w$; since
the window is closed, $|\hat w|\le\varrho_2/2$. Write
\begin{align*}
S^{n}_{\mathbb{T}_{m_{\natural}}}(x^{(n)};\vec f\,)
-\mathfrak{S}_{m_{\natural}}[\hat x](\vec f\,)
&=\Bigl[S^{n}_{\mathbb{T}_{m_{\natural}}}(x^{(n)};\vec f\,)
-S^{(a_n-x^{(n)})}_{\mathbb{T}_{m_{\natural}}}(\vec f\,)\Bigr]\\
&\quad+\Bigl[S^{(a_n-x^{(n)})}_{\mathbb{T}_{m_{\natural}}}(\vec f\,)
-S^{(\hat w)}_{\mathbb{T}_{m_{\natural}}}(\vec f\,)\Bigr],
\end{align*}
where we used $S^{(\hat w)}_{\mathbb{T}_{m_{\natural}}}(\vec f\,)
=\mathfrak{S}_{m_{\natural}}[\hat x](\vec f\,)$ from
\eqref{9.98:eq-reference-torus-limit-a}. The first bracket tends to zero by
\eqref{9.98:eq-reference-torus-limit-1}. By clause~(v-b), read on $\mathbb{T}_{m_{\natural}}$,
the map $\hat w\mapsto S^{(\hat w)}_{\mathbb{T}_{m_{\natural}}}(\vec f\,)$ is holomorphic on the
strip $|\operatorname{Im}\hat w|<\upsilon_w(m_{\natural})$, which contains the closed real
window $[-\varrho_2/2,\varrho_2/2]$; in particular the map is continuous on this window.
Since $a_n-x^{(n)}$ lies in this window and tends to
$\hat w$, the second bracket tends to zero. Hence
$S^{n}_{\mathbb{T}_{m_{\natural}}}(x^{(n)};\vec f\,)\to\mathfrak{S}_{m_{\natural}}[\hat x](\vec f\,)$.

\emph{Proof of (c).} Clause~(v-b) of Theorem~\ref{3.1:thm-main} states that the torus enters
its conclusions only through $\upsilon_w$, the rate constant and a floor $n_0$, and that these
read the side length $N_w=2L^{m}$ of the torus through
$\mathsf{Q}_w$ of \eqref{9.98:eq-reference-torus-limit-2}. In the uniqueness theorem the rate
of convergence does not depend on the volume; its constant does. On the torus
$\mathbb{T}_{m_{\natural}}$ one has
$N_w=2L^{m_{\natural}(g)}$, and $m_{\natural}(g)$, defined in
\eqref{9.98:eq-prolongation-depths-c}, is a function of $g$ alone. Hence the rate constant,
$\upsilon_w(m_{\natural})$ and $n_0(N_w)$ are numbers depending on $g$ alone. In (a) and (b)
they multiply a term which the limit $n\to\infty$ removes.
\end{proof}

\begin{remark}
The moduli of Lemma~\ref{9.98:lem-reference-torus-limit}(c) enter none of the
constants of the short-distance comparison of infinite-volume limit points at a given
renormalised coupling that is built on the present lemma. In that comparison clause~(v) of
Theorem~\ref{3.1:thm-main} is used only here; its factor $\exp[\mathsf{B}N_w^4]$, with
$\mathsf{B}=E_{+}^{\sharp}+C_{\mathrm{low}}$, is incurred once, at the value
$N_w=2L^{m_{\natural}(g)}$ fixed by $g$, inside a term that vanishes as $n\to\infty$.
\end{remark}

\begin{proof}[Proof of \textup{(a)} in Theorem~\ref{9.98:thm-short-distance-uniqueness}]
\label{9.98:prf-proof-two-tori}
The hypotheses of Theorem~\ref{9.98:thm-short-distance-uniqueness} are in force. Thus
$g\le\min\{\gamma_{\natural},\gamma_U'\}$, $\hat w$ is real with $|\hat w|\le\varrho_2/2$ and $n\in\{2,3\}$.
Moreover $0<\mathsf{d}\le L^{\underline{s}'(g)}/40$, and $\vec f\in\mathfrak{P}_n(\mathsf{d};\vartheta)$.
In this proof the index $n$ of $\tau_n$ in \eqref{9.98:eq-short-distance-uniqueness-a} is always
the number of sources. The tuned depth is written $\nu$, so the tuned member of depth $\nu$ has
bare inverse coupling $a_\nu-\hat w$.

Let $m,m'\ge m_{\natural}(g)$. Let $\nu$ be a tuned depth with $\nu\ge n_{\natural}$ and
$L^{-\nu}\le\mathsf{d}$. We apply the transfer inequality, Proposition~\ref{9.98:prop-transfer-inequality},
with $(m_1,m_2)=(m,m')$. Its hypotheses hold:
\begin{itemize}
\item $m_1,m_2\ge m_{\natural}(g)$;
\item $n\in\{2,3\}$;
\item the depth $\nu$ is at least $n_{\natural}$ and satisfies $L^{-\nu}\le\mathsf{d}$;
\item $\vec f\in\mathfrak{P}_n(\mathsf{d};\vartheta)$ with $\mathsf{d}\le L^{\underline{s}'(g)}/40$.
\end{itemize}
The tuned-frame prolongation $\mathsf{s}'=\mathsf{s}'(\nu)$ of the depth-$\nu$ member satisfies
$\mathsf{s}'(\nu)\ge\underline{s}'(g)$ for every $\nu$. Hence
\eqref{9.98:eq-transfer-inequality-a} gives, with $\tau_n$ as in
\eqref{9.98:eq-short-distance-uniqueness-a},
\begin{equation}\label{9.98:eq-proof-two-tori-1}
\bigl|S^{\nu}_{\mathbb{T}_m}(a_\nu-\hat w;\vec f\,)
-S^{\nu}_{\mathbb{T}_{m'}}(a_\nu-\hat w;\vec f\,)\bigr|\le\tau_n .
\end{equation}
The right-hand side $\tau_n$ does not depend on $\nu$. The bound holds for every sufficiently
large $\nu$, namely every $\nu\ge n_{\natural}$ with $L^{-\nu}\le\mathsf{d}$.

Next we let $\nu\to\infty$. By Lemma~\ref{9.98:lem-reference-torus-limit} we have
$m_{\natural}(g)\ge m_{\mathbb{T}}(g_-(g))$, so both exponents $m$ and $m'$ are at least
$m_{\mathbb{T}}(g_-(g))$. Clause (ii-$\mathbb{T}$-a) of Theorem~\ref{3.1:thm-main}, for
$n\in\{2,3\}$, therefore applies on $\mathbb{T}_m$ and on $\mathbb{T}_{m'}$. It shows that both
members on the left of \eqref{9.98:eq-proof-two-tori-1} converge as $\nu\to\infty$. Their limits
are the continuum limits $\mathfrak{S}_m[\hat x](\vec f\,)$ and
$\mathfrak{S}_{m'}[\hat x](\vec f\,)$ on the two tori at the renormalised coupling $\hat x$.

The set of pairs of reals $(a,b)$ with $|a-b|\le\tau_n$ is closed. We pass to the limit
$\nu\to\infty$ in \eqref{9.98:eq-proof-two-tori-1} and obtain
\begin{equation}\label{9.98:eq-proof-two-tori-2}
\bigl|\mathfrak{S}_m[\hat x](\vec f\,)-\mathfrak{S}_{m'}[\hat x](\vec f\,)\bigr|\le\tau_n ,
\end{equation}
for all $m,m'\ge m_{\natural}(g)$. This is \textup{(a)}.
\end{proof}

\begin{proof}[Proof of (b): triangle through the reference torus]
\label{9.98:prf-proof-triangle}
Fix $g$, $\hat w$, $n\in\{2,3\}$, $\mathsf{d}$ and $\vec f\in\mathfrak{P}_n(\mathsf{d};\vartheta)$ as in
Theorem~\ref{9.98:thm-short-distance-uniqueness}, and let $\tau_n$ be the quantity defined in
\eqref{9.98:eq-short-distance-uniqueness-a}. We abbreviate $m_{\natural}:=m_{\natural}(g)$. As in the
transfer inequality at $\gamma_{\star}$ (Proposition~\ref{9.98:prop-transfer-inequality}), a
superscript on $S$ is the depth of the tuned member, while the subscript $n$ of $\tau_n$ is the
number of test functions.

\emph{Infinite-volume limit points.} Let $\sigma\in\operatorname{Lim}_{\infty}[\hat{x}]$. By the
definition of this set there is a sequence $(n_j,m_j)_{j\ge1}$ with $n_j\to\infty$ and
$m_j\to\infty$, the two tending to infinity in any manner, such that
\begin{equation}\label{9.98:eq-proof-triangle-1}
\sigma(\vec f\,)=\lim_{j\to\infty}S^{n_j}_{\mathbb{T}_{m_j}}(a_{n_j}-\hat w;\vec f\,).
\end{equation}
Here $S_n[\sigma](\vec f\,)$ denotes the connected $n$-point function of $\sigma$ on the tuple
$\vec f$; thus $S_n[\sigma](\vec f\,)=\sigma(\vec f\,)$.

Since $n_j\to\infty$, $m_j\to\infty$ and $\mathsf{d}>0$, there is $j_0$ such that for every
$j\ge j_0$ we have $n_j\ge n_{\natural}$, $L^{-n_j}\le\mathsf{d}$ and $m_j\ge m_{\natural}$. For such
$j$ we apply Proposition~\ref{9.98:prop-transfer-inequality} with $(m_1,m_2)=(m_j,m_{\natural})$,
tuned depth $n_j$ and bare inverse coupling $x_0=a_{n_j}-\hat w$. Its hypotheses hold: $m_1=m_j\ge
m_{\natural}(g)$ by the choice of $j_0$ and $m_2=m_{\natural}(g)$; $n\in\{2,3\}$; the depth satisfies
$n_j\ge n_{\natural}$ and $L^{-n_j}\le\mathsf{d}$; and $\vec f\in\mathfrak{P}_n(\mathsf{d};\vartheta)$
with $\mathsf{d}\le L^{\underline{s}'(g)}/40$ by the hypotheses of
Theorem~\ref{9.98:thm-short-distance-uniqueness}. Display \eqref{9.98:eq-transfer-inequality-a}
therefore gives, for every $j\ge j_0$,
\begin{equation}\label{9.98:eq-proof-triangle-2}
\bigl|S^{n_j}_{\mathbb{T}_{m_j}}(a_{n_j}-\hat w;\vec f\,)
-S^{n_j}_{\mathbb{T}_{m_{\natural}}}(a_{n_j}-\hat w;\vec f\,)\bigr|\le\tau_n .
\end{equation}

We let $j\to\infty$ in \eqref{9.98:eq-proof-triangle-2}. Since $|\hat w|\le\varrho_2/2$, the cutoff
limit on the reference torus (Lemma~\ref{9.98:lem-reference-torus-limit}, display
\eqref{9.98:eq-reference-torus-limit-a}) states that the tuned members on
$\mathbb{T}_{m_{\natural}}$, evaluated on $\vec f$, converge to
$\mathfrak{S}_{m_{\natural}}[\hat{x}](\vec f\,)$ as the depth tends to infinity through all
integers, not only along a subsequence; in particular, since $n_j\to\infty$, the
second member of \eqref{9.98:eq-proof-triangle-2} tends to
$\mathfrak{S}_{m_{\natural}}[\hat{x}](\vec f\,)$. The first member tends to $\sigma(\vec f\,)$ by
\eqref{9.98:eq-proof-triangle-1}. The right side $\tau_n$ does not depend on $j$, and a non-strict
inequality is preserved under limits.
Hence
\begin{equation}\label{9.98:eq-proof-triangle-3}
\bigl|S_n[\sigma](\vec f\,)-\mathfrak{S}_{m_{\natural}}[\hat{x}](\vec f\,)\bigr|\le\tau_n .
\end{equation}
No limits are exchanged here: the limit $j\to\infty$ is taken along the same sequence
$(n_j,m_j)$ that defines $\sigma$, whether the cutoff and the volume tend to infinity jointly or in
iterated form, and the order in which $\sigma$ takes them is never altered.

\emph{Torus theories.} Let $\sigma=\mathfrak{S}_m[\hat{x}]$ with $m\ge m_{\natural}$, so that
$S_n[\sigma](\vec f\,)=\mathfrak{S}_m[\hat{x}](\vec f\,)$. Part (a) of
Theorem~\ref{9.98:thm-short-distance-uniqueness}, applied with $m'=m_{\natural}$ (both exponents are
at least $m_{\natural}(g)$), gives
$|\mathfrak{S}_m[\hat{x}](\vec f\,)-\mathfrak{S}_{m_{\natural}}[\hat{x}](\vec f\,)|\le\tau_n$; thus
\eqref{9.98:eq-proof-triangle-3} holds for these $\sigma$ as well.

\emph{The triangle.} Let $\sigma,\sigma'\in\operatorname{Lim}_{\infty}[\hat{x}]\cup
\{\mathfrak{S}_m[\hat{x}]:m\ge m_{\natural}\}$. By the triangle inequality through
$\mathfrak{S}_{m_{\natural}}[\hat{x}](\vec f\,)$ and \eqref{9.98:eq-proof-triangle-3} for $\sigma$ and
for $\sigma'$,
\begin{equation}\label{9.98:eq-proof-triangle-4}
\begin{aligned}
\bigl|S_n[\sigma](\vec f\,)-S_n[\sigma'](\vec f\,)\bigr|
&\le\bigl|S_n[\sigma](\vec f\,)-\mathfrak{S}_{m_{\natural}}[\hat{x}](\vec f\,)\bigr|
+\bigl|\mathfrak{S}_{m_{\natural}}[\hat{x}](\vec f\,)-S_n[\sigma'](\vec f\,)\bigr|\\
&\le2\tau_n .
\end{aligned}
\end{equation}

\emph{Real-analyticity and the strip.} The comparison object
$\mathfrak{S}_{m_{\natural}}[\hat{x}](\vec f\,)$ is the continuum limit on the one torus
$\mathbb{T}_{m_{\natural}(g)}$. By Lemma~\ref{9.98:lem-reference-torus-limit}, which reads clause
(v-b) of Theorem~\ref{3.1:thm-main} on this torus without change, the map
$\hat w\mapsto\mathfrak{S}_{m_{\natural}}[\hat{x}](\vec f\,)$ is real-analytic on
$\mathopen{]}-\varrho_2/2,\varrho_2/2\mathclose{[}$ and holomorphic on the strip
$|\operatorname{Im}\hat w|<\upsilon_w(m_{\natural}(g))$, whose width depends on $g$ alone.

\emph{The name $\hat{x}$.} The parameter $\hat w$ is identified with the renormalised reference
coupling $\hat{x}$ through clause (v-d)(i) of Theorem~\ref{3.1:thm-main} at each depth and through
the bi-Lipschitz lemma for the limit chart. This identification serves only to call $\hat w$ the
renormalised coupling; it enters none of the inequalities
\eqref{9.98:eq-proof-triangle-2}--\eqref{9.98:eq-proof-triangle-4}.
\end{proof}

\begin{proof}[Proof of (c): the two limits commute]\label{9.98:prf-proof-limits-commute}
We work under the hypotheses of Theorem~\ref{9.98:thm-short-distance-uniqueness}:
$g\le\min\{\gamma_{\natural},\gamma_U'\}$, $\hat w$ real with $|\hat w|\le\varrho_2/2$,
$0<\mathsf{d}\le L^{\underline{s}'(g)}/40$ and $\vec f\in\mathfrak{P}_n(\mathsf{d};\vartheta)$ with
$n\in\{2,3\}$. We abbreviate $m_{\natural}:=m_{\natural}(g)$.

The letter $n$ is used in two senses, as in the transfer inequality
(Proposition~\ref{9.98:prop-transfer-inequality}). In
$S^{n}_{\mathbb{T}_m}(a_n-\hat w;\vec f\,)$, in $a_n$ and in every limit $n\to\infty$ below, $n$ is the
depth of the tuned member. In $\tau_n$ of \eqref{9.98:eq-short-distance-uniqueness-a}, $n$ is the
number of test functions. That number is fixed throughout, so $\tau_n$ does not depend on the
depth.

\emph{The transfer bound.} Let $n\ge n_{\natural}$ be a depth with $L^{-n}\le\mathsf{d}$; we call such
a depth admissible. Let $m\ge m_{\natural}$. We apply the transfer inequality
\eqref{9.98:eq-transfer-inequality-a} with
$m_1=m$ and $m_2=m_{\natural}$. Its hypotheses hold:
\begin{itemize}
\item both torus exponents are at least $m_{\natural}(g)$;
\item the depth satisfies $n\ge n_{\natural}$ and $L^{-n}\le\mathsf{d}$;
\item $\vec f\in\mathfrak{P}_n(\mathsf{d};\vartheta)$ with $\mathsf{d}\le L^{\underline{s}'(g)}/40$.
\end{itemize}
It gives
\begin{equation}\label{9.98:eq-proof-limits-commute-1}
\bigl|S^{n}_{\mathbb{T}_m}(a_n-\hat w;\vec f\,)
-S^{n}_{\mathbb{T}_{m_{\natural}}}(a_n-\hat w;\vec f\,)\bigr|\le\tau_n .
\end{equation}

\emph{The reference torus.} Since $|\hat w|\le\varrho_2/2$ and $\vec f\in\mathfrak{P}_n(\mathsf{d};\vartheta)$
is admissible for Lemma~\ref{9.98:lem-reference-torus-limit}, the cutoff limit on the reference
torus, \eqref{9.98:eq-reference-torus-limit-a} of
Lemma~\ref{9.98:lem-reference-torus-limit}, applies on $\mathbb{T}_{m_{\natural}}$. It states that
$S^{n}_{\mathbb{T}_{m_{\natural}}}(a_n-\hat w;\vec f\,)\to\mathfrak{S}_{m_{\natural}}[\hat x](\vec f\,)$ as
$n\to\infty$, without passing to subsequences. For each admissible depth $n$ put
\begin{equation*}
\eta_n:=\bigl|S^{n}_{\mathbb{T}_{m_{\natural}}}(a_n-\hat w;\vec f\,)
-\mathfrak{S}_{m_{\natural}}[\hat x](\vec f\,)\bigr| .
\end{equation*}
Then $\eta_n\to0$ as $n\to\infty$. Moreover $\eta_n$ does not depend on $m$, since it is a quantity
on the one torus $\mathbb{T}_{m_{\natural}}$. The triangle inequality and
\eqref{9.98:eq-proof-limits-commute-1} give, for every depth $n\ge n_{\natural}$ with
$L^{-n}\le\mathsf{d}$ and every $m\ge m_{\natural}$,
\begin{equation}\label{9.98:eq-proof-limits-commute-2}
\bigl|S^{n}_{\mathbb{T}_m}(a_n-\hat w;\vec f\,)-\mathfrak{S}_{m_{\natural}}[\hat x](\vec f\,)\bigr|
\le\tau_n+\eta_n .
\end{equation}

We show that every cluster value of the double array
$(n,m)\mapsto S^{n}_{\mathbb{T}_m}(a_n-\hat w;\vec f\,)$, as $n,m\to\infty$, lies in the closed
interval
\begin{equation}\label{9.98:eq-proof-limits-commute-3}
\bigl[\,\mathfrak{S}_{m_{\natural}}[\hat x](\vec f\,)-\tau_n,\;
\mathfrak{S}_{m_{\natural}}[\hat x](\vec f\,)+\tau_n\,\bigr],
\end{equation}
which has length $2\tau_n$. This holds in each of the following three orders of passage.

\emph{First $n\to\infty$, then $m\to\infty$.} Fix $m\ge m_{\natural}$. In this order the inner limits
are the values $\mathfrak{S}_m[\hat x](\vec f\,)$, which exist by clause (ii-$\mathbb{T}$-a) of
Theorem~\ref{3.1:thm-main}. By clause~(a) of Theorem~\ref{9.98:thm-short-distance-uniqueness},
proved above, applied with $m'=m_{\natural}$,
\begin{equation*}
\bigl|\mathfrak{S}_m[\hat x](\vec f\,)-\mathfrak{S}_{m_{\natural}}[\hat x](\vec f\,)\bigr|\le\tau_n .
\end{equation*}
Hence the sequence $m\mapsto\mathfrak{S}_m[\hat x](\vec f\,)$ is confined to the interval
\eqref{9.98:eq-proof-limits-commute-3}. Being bounded, it has cluster values as $m\to\infty$. Each of
them lies in \eqref{9.98:eq-proof-limits-commute-3}, because that interval is closed.

\emph{First $m\to\infty$ along a subsequence, then $n\to\infty$.} Fix an admissible depth $n$. By
\eqref{9.98:eq-proof-limits-commute-2}, for all $m\ge m_{\natural}$,
\begin{equation*}
\bigl|S^{n}_{\mathbb{T}_m}(a_n-\hat w;\vec f\,)\bigr|
\le\bigl|\mathfrak{S}_{m_{\natural}}[\hat x](\vec f\,)\bigr|+\tau_n+\eta_n ,
\end{equation*}
so these numbers are bounded uniformly in $m$. Uniform boundedness in $m$ also follows,
alternatively, from the per-ball volume-free bound, clause (iii-b) of Theorem~\ref{3.1:thm-main},
applied at the reference coupling $\gamma_{\star}$ to the same tuned member read in the frame of
$\gamma_{\star}$, once $L^{m-\mathsf{s}'}\ge\ell(n,1,\gamma_{\star})$; here $\mathsf{s}'=K'-n$, with
$n$ the depth, is the tuned-frame prolongation, and $n$ in $\ell(n,1,\gamma_{\star})$ is the number
of test functions. The
argument does not need this. By the Bolzano--Weierstrass theorem, cluster values as $m\to\infty$
along any subsequence exist. By \eqref{9.98:eq-proof-limits-commute-2} each such value $c_n$
satisfies $|c_n-\mathfrak{S}_{m_{\natural}}[\hat x](\vec f\,)|\le\tau_n+\eta_n$. Now let
$n\to\infty$. Since $\eta_n\to0$, every cluster value $c$ of the sequence $(c_n)$ satisfies
$|c-\mathfrak{S}_{m_{\natural}}[\hat x](\vec f\,)|\le\tau_n$, so $c$ lies in
\eqref{9.98:eq-proof-limits-commute-3}.

\emph{Jointly, in any manner.} Let $(n_j,m_j)$ be admissible pairs, that is, $n_j\ge n_{\natural}$,
$L^{-n_j}\le\mathsf{d}$ and $m_j\ge m_{\natural}$. Suppose $n_j,m_j\to\infty$ and
$S^{n_j}_{\mathbb{T}_{m_j}}(a_{n_j}-\hat w;\vec f\,)\to c$. By
\eqref{9.98:eq-proof-limits-commute-2}, the $j$-th term lies within $\tau_n+\eta_{n_j}$ of
$\mathfrak{S}_{m_{\natural}}[\hat x](\vec f\,)$. Letting $j\to\infty$ gives
$|c-\mathfrak{S}_{m_{\natural}}[\hat x](\vec f\,)|\le\tau_n$.

In all three orders, therefore, every cluster value lies in the one closed interval
\eqref{9.98:eq-proof-limits-commute-3}. This is a Cauchy property in the cutoff, uniform in the
volume, at short distance. Its floor is $\tau_n$, of order $\mathsf{d}^{4-n}$ in the number $n$ of
test functions. No iterated limit is exchanged with another, no derivative in a coupling is
taken, and no sum over far pairs is formed.
\end{proof}

\begin{remark}[What a tuned limit point presupposes]\label{9.98:rem-tuned-limit-points}
Let $g$, $\hat{w}$ and $\hat{x}$ be as in
Theorem~\ref{9.98:thm-short-distance-uniqueness}. Let $n\ge n_{\natural}$ be a depth, with
$n_{\natural}$ as in Lemma~\ref{9.98:lem-tuned-depth-offset}, and $m$ a torus exponent. Let
$x_0=a_n(g)-\hat{w}$ and $K_0=K(x_0,g)$ be as in Lemma~\ref{9.98:lem-tuned-depth-offset}.
\begin{itemize}
\item[(a)] The inequalities of Theorem~\ref{9.98:thm-short-distance-uniqueness} use nothing about
the tuned-frame limit set $\operatorname{Lim}_{\infty}[\hat{x}]$ except that it is the set of
cluster values of the double array through which it is defined.
\item[(b)] To say of a limit point of the sequence $(\mathfrak{S}_m[\hat{x}])_{m\ge m_{\natural}(g)}$,
with $m_{\natural}(g)$ as in \eqref{9.98:eq-prolongation-depths-c}, along
subsequences of volumes what clause (ii-$\mathbb{R}^4$) of Theorem~\ref{3.1:thm-main} says of a
first-passage limit point, namely invariance, reflection positivity and
Osterwalder--Schrader reconstruction,
presupposes that the tuned member of depth $n$ on $\mathbb{T}_m$ can be read as the native
first-passage member $(x_0,K_0,\mathbb{T}_{m+n-K_0})$ at reference $g$ (the member of the
first-passage run of clause (0) of Theorem~\ref{3.1:thm-main} started at bare inverse coupling
$x_0$, with $K_0$
steps, on $\mathbb{T}_{m+n-K_0}$), and that this member satisfies the one-sided floor
\begin{equation}\label{9.98:eq-tuned-limit-points-1}
m+n-K_0\ \ge\ m_{\mathbb{T}}(g_{-}(g)).
\end{equation}
For $m\ge m_{\natural}(g)$, with $m_{\natural}(g)$ as in \eqref{9.98:eq-prolongation-depths-c}, the
floor \eqref{9.98:eq-tuned-limit-points-1} holds. Under it, the passage to $\mathbb{R}^4$ of the
infinite-volume theorem applies to the sequence $(\mathfrak{S}_m[\hat{x}])_{m\ge m_{\natural}(g)}$.
\end{itemize}
\end{remark}

\begin{proof}
(a) Each inequality of Theorem~\ref{9.98:thm-short-distance-uniqueness} concerns the elements of
$\operatorname{Lim}_{\infty}[\hat{x}]$. These elements are cluster values of the double array of
tuned-frame Schwinger functions. The derivation of the inequalities uses no further property of
these cluster values.

(b) The per-ball volume-free bound, clause (iii-b) of Theorem~\ref{3.1:thm-main}, is a statement
about the
derivatives of $\log Z(z)$ with respect to the sources for a Wilson measure. It therefore holds
for any description of the same Wilson measure, in particular for the tuned member once it is
identified with a first-passage member.

The tuned member of depth $n$ on $\mathbb{T}_m$ is the Wilson measure on the lattice with
$2L^{m+n}$ sites a side, at bare inverse coupling $x_0$. Along the native first-passage run from
$x_0$ at reference $g$, its number of steps at reference $g$ is $K_0=K(x_0,g)$ in the sense of
clause (0) of Theorem~\ref{3.1:thm-main}. At bare
spacing $L^{-K_0}$ in $g$-reference units, a lattice with $2L^{m+n}$ sites a side is the bare
lattice of the torus of side $2L^{m+n-K_0}$, that is, of $\mathbb{T}_{m+n-K_0}$. Hence the tuned
member is the same measure as the native first-passage member $(x_0,K_0,\mathbb{T}_{m+n-K_0})$ at
reference $g$. The per-ball volume-free bound applies to this member once the floor
\eqref{9.98:eq-tuned-limit-points-1} holds.

We now verify the floor. Since $n\ge n_{\natural}$ and $|\hat{w}|\le\varrho_2/2$ (the latter by the
choice of $\hat{w}$ as in Theorem~\ref{9.98:thm-short-distance-uniqueness}),
Lemma~\ref{9.98:lem-tuned-depth-offset} gives
$|K_0-n|\le j_U-1$, as displayed in \eqref{9.98:eq-tuned-depth-offset-a}. In particular
$K_0\le n+j_U-1$, hence $m+n-K_0\ge m-j_U+1$. By \eqref{9.98:eq-floor-members-a} of
Remark~\ref{9.98:rem-floor-members}, the hypothesis $m\ge m_{\natural}(g)$ gives
$m-j_U+1\ge m_{\mathbb{T}}(g_{-}(g))+1$, so that
\begin{equation*}
m+n-K_0\ \ge\ m-j_U+1\ \ge\ m_{\mathbb{T}}(g_{-}(g))+1 .
\end{equation*}
This implies \eqref{9.98:eq-tuned-limit-points-1}.

The proof of the infinite-volume theorem passes to $\mathbb{R}^4$ in five steps. Apply
them to the sequence $(\mathfrak{S}_m[\hat{x}])_{m\ge m_{\natural}(g)}$. These steps use the bounds
of clause (iii-b) of Theorem~\ref{3.1:thm-main} and invariance and reflection positivity on each
torus. By the
preceding paragraphs the bounds of clause (iii-b) hold for every lattice member
$(x_0,K_0,\mathbb{T}_{m+n-K_0})$ with $m\ge m_{\natural}(g)$, and the continuum torus limits
$\mathfrak{S}_m[\hat{x}]$ carry these bounds, together with invariance and reflection positivity
on each torus, by clause (ii-$\mathbb{T}$-a) of Theorem~\ref{3.1:thm-main}. With these inputs
the steps apply
without change. They yield the invariance, the reflection positivity and the
Osterwalder--Schrader reconstruction of every limit point of
$(\mathfrak{S}_m[\hat{x}])_{m\ge m_{\natural}(g)}$ along subsequences of volumes.
\end{proof}

\begin{proof}[Proof of (g): rectangular tori]
\label{9.98:prf-proof-rectangular-tori}
Fix $g$, $\hat w$, $n\in\{2,3\}$, $\mathsf{d}$ and $\vec f\in\mathfrak{P}_n(\mathsf{d};\vartheta)$ as in
Theorem~\ref{9.98:thm-short-distance-uniqueness}. Let $\hat x$ be the renormalised reference coupling
attached to $\hat w$, and put $\Pi=\prod_i\|f_i\|_\infty$. As in the transfer inequality
\eqref{9.98:eq-transfer-inequality-a}, the depth of a tuned member is written as a dot in the
superscript, $S^{\,\cdot\,}_{\mathbb{T}}(x_0;\vec f\,)$ for a torus $\mathbb{T}$. Here $x_0$ is its
bare inverse coupling, namely the
value at that depth of the tuned sequence of Lemma~\ref{9.98:lem-tuned-member}, minus $\hat w$. In
formulas the letter $n$ denotes only the number of sources; it is the subscript of $\ell_n$, $\tau_n$
and $C^{\star}_{T,n}$.

Let $\mathfrak{B}$ be a rectangular torus whose four sides $l_1,\dots,l_4$, measured in
$g$-reference units, satisfy
\begin{equation}\label{9.98:eq-proof-rectangular-tori-1}
l_\mu\in 2\ell_n L^{\overline{s}'(g)}\mathbb{N},\qquad
l_\mu\ge 2L^{m_{\natural}(g)}\qquad(\mu=1,\dots,4),
\end{equation}
where $m_{\natural}(g)$ is the floor \eqref{9.98:eq-prolongation-depths-c}. The tuned members on
$\mathfrak{B}$ are defined as in Lemma~\ref{9.98:lem-tuned-member}, with $\mathbb{T}_m$ replaced by
$\mathfrak{B}$. We write $S^{\,\cdot\,}_{\mathfrak{B}}(x_0;\vec f\,)$ for their connected $n$-point
functions.

\emph{Cell counts.} We pass to the $\gamma_{\star}$-frame of Lemma~\ref{9.98:lem-tuned-member}. Its
reference length is $L^{\mathsf{s}'}$ $g$-reference units, where $\mathsf{s}'\le\overline{s}'(g)$. In
this frame the chessboard cell for $n$ sources has side $\ell_n$. Hence $\mathfrak{B}$ has
$\mathsf{n}_\mu=l_\mu L^{-\mathsf{s}'}/\ell_n$ cells in the direction $\mu$.

By the first condition in \eqref{9.98:eq-proof-rectangular-tori-1},
$\mathsf{n}_\mu\in 2L^{\overline{s}'(g)-\mathsf{s}'}\mathbb{N}$. Since $\mathsf{s}'\le\overline{s}'(g)$,
the factor $L^{\overline{s}'(g)-\mathsf{s}'}$ is an integer, so every $\mathsf{n}_\mu$ is even.

Put $m_{\mathrm{sm}}:=m_{\natural}(g)-\mathsf{s}'$. The cubic $\gamma_{\star}$-frame torus
$\mathbb{T}_{m_{\mathrm{sm}}}$ has side $2L^{m_{\mathrm{sm}}}$ in $\gamma_{\star}$-reference units, that
is $2L^{m_{\natural}(g)}$ in $g$-reference units. It is therefore the torus $\mathbb{T}_{m_{\natural}(g)}$
of the $g$-frame. Its $\gamma_{\star}$-exponent $m_{\mathrm{sm}}$ is at least $m_{\star}+1$ by
\eqref{9.98:eq-prolongation-depths-c} and $\mathsf{s}'\le\overline{s}'(g)$, with
$m_{\star}$ as in Definition~\ref{9.98:def-prolongation-depths}. The second condition in
\eqref{9.98:eq-proof-rectangular-tori-1} gives
\begin{equation}\label{9.98:eq-proof-rectangular-tori-2}
\mathsf{n}_\mu\;\ge\;\frac{2L^{m_{\mathrm{sm}}}}{\ell_n}\;\ge\;\frac{2L^{m_{\star}+1}}{\ell_n}
\qquad(\mu=1,\dots,4).
\end{equation}
The middle member is the number of cells per direction of $\mathbb{T}_{m_{\mathrm{sm}}}$. The right
member is that of $\mathbb{T}_{\star}'$, and $\mathbb{T}_{m_{\mathrm{sm}}}\supseteq
\mathbb{T}_{\star}'\supseteq\mathbb{T}_{\star}$. Thus the cell counts of $\mathfrak{B}$ are even and
dominate, direction by direction, those of the cubic tori $\mathbb{T}_{m_{\mathrm{sm}}}$,
$\mathbb{T}_{\star}'$ and $\mathbb{T}_{\star}$.

\emph{The volume transfer with a rectangular large volume.} We run the proof of the volume-transfer
bound at reference $\gamma_{\star}$, with large volume $\mathfrak{B}$ and small volume the cubic torus
$\mathbb{T}_{m_{\mathrm{sm}}}$. For $n=3$ this is Proposition~\ref{9.98:prop-transfer-three-sources}; for
$n=2$ it is the two-source volume-transfer theorem. Its steps go through as follows.

The step that places the supports of $f_1,\dots,f_n$ in one ball with a margin is used without change.

The chessboard step is used without change. The chessboard estimate holds on every torus whose cell
count is even in each direction, and $\mathfrak{B}$ is such a torus.

The comparison step of the large with the small volume uses the monotonicity of the normalised partition
functions in the cell counts. For a torus $\mathfrak{X}$ of the $\gamma_{\star}$-frame whose cell counts
are all even, write $\mathsf{n}(\mathfrak{X})\in(2\mathbb{N})^4$ for its cell-count vector; for
$\mathsf{n}\in(2\mathbb{N})^4$ write $|\mathsf{n}|=\prod_\mu\mathsf{n}_\mu$ for the number of cells and,
for a weight $W$ (below, a mirror-periodised source weight $W_z$, or $W=0$), $Z_{\mathsf{n}}(W)$ for the
partition function of the bare Wilson measure on the torus with cell-count vector $\mathsf{n}$ with that
weight. For $\mathsf{n}\le\mathsf{n}'$ componentwise in $(2\mathbb{N})^4$ the monotonicity states
\begin{equation*}
Z_{\mathsf{n}'}(W)^{1/|\mathsf{n}'|}\le Z_{\mathsf{n}}(W)^{1/|\mathsf{n}|}.
\end{equation*}
In particular every pair of cubic or rectangular tori with
even cell counts $\mathsf{n}\le\mathsf{n}'$ can be compared. We apply it with
$\mathsf{n}=\mathsf{n}(\mathbb{T}_{m_{\mathrm{sm}}})$ and $\mathsf{n}'=\mathsf{n}(\mathfrak{B})$. The
componentwise inequality is \eqref{9.98:eq-proof-rectangular-tori-2}. No divisibility of $\mathsf{n}'$ by
$\mathsf{n}$ is needed.

The stability step is taken in the convexity form of the volume transfer. On the small cubic torus,
Ba{\l}aban's construction at reference $\gamma_{\star}$ is available. Indeed $\mathbb{T}_{m_{\mathrm{sm}}}$
is a member of the first-passage family at reference $\gamma_{\star}$, since $m_{\mathrm{sm}}\ge
m_{\star}+1$. Two bounds are used there: the stability bound of the correlation theorem, and the lower
stability bound, namely the lower member of clause (i-b) of Theorem~\ref{3.1:thm-main} at the last level
and zero source. Together they give
\begin{equation}\label{9.98:eq-proof-rectangular-tori-3}
\Bigl[\frac{Z_{m_{\mathrm{sm}}}(W_z)}{Z_{m_{\mathrm{sm}}}(0)}\Bigr]^{1/|\mathcal{C}_{m_{\mathrm{sm}}}|}
\le\exp\Bigl\{|\mathcal{C}_{m_{\mathrm{sm}}}|^{-1}\operatorname{Re}\mathcal{E}(W_z)
+\mathsf{B}\,\ell_n^4\Bigr\}.
\end{equation}
Here $Z_{m_{\mathrm{sm}}}:=Z_{\mathsf{n}(\mathbb{T}_{m_{\mathrm{sm}}})}$ is the partition function of
$\mathbb{T}_{m_{\mathrm{sm}}}$, so that $|\mathcal{C}_{m_{\mathrm{sm}}}|=|\mathsf{n}(\mathbb{T}_{m_{\mathrm{sm}}})|$;
$z$ is a complex source vector as in the volume-transfer proof, $W_z$ is its mirror-periodised source
weight, and $\mathcal{E}(W_z)$ is its source counterterm on the small torus. The constant
$\mathsf{B}=E_{+}^{\sharp}+C_{\mathrm{low}}$ is the free-energy-density mismatch.

It remains to control the drop of the zero-source pressure from the small torus to $\mathfrak{B}$. For a
torus $\mathfrak{X}$ with even cell counts, write $\pi(\mathfrak{X})$ for the zero-source pressure of its
bare Wilson measure per unit $\gamma_{\star}$-reference volume (so that the pressure per chessboard cell
is $\ell_n^4\pi(\mathfrak{X})$). By the monotonicity just quoted at $W=0$, $\pi$ is non-increasing in
each cell count, the cell counts being even.

Choose an integer $e\ge0$ such that the cubic $\gamma_{\star}$-frame torus $\mathbb{T}_{m_{\mathrm{sm}}+e}$,
of side $L^e$ times the side of $\mathbb{T}_{m_{\mathrm{sm}}}$, has side at least every side of
$\mathfrak{B}$. Then the cell-count vectors satisfy
\begin{equation*}
\mathsf{n}(\mathbb{T}_{m_{\mathrm{sm}}})\le\mathsf{n}(\mathfrak{B})\le
\mathsf{n}(\mathbb{T}_{m_{\mathrm{sm}}+e})
\end{equation*}
componentwise, and therefore
\begin{equation*}
\pi(\mathbb{T}_{m_{\mathrm{sm}}})\ge\pi(\mathfrak{B})\ge\pi(\mathbb{T}_{m_{\mathrm{sm}}+e}).
\end{equation*}
The cubic drop $\pi(\mathbb{T}_{m_{\mathrm{sm}}})-\pi(\mathbb{T}_{m_{\mathrm{sm}}+e})$ is at most
$\mathsf{c}_L(\mathsf{B}+\bar{\delta})$, where $\mathsf{c}_L=\bigl((2L-1)/(L-1)\bigr)^4$ is the factor
of the convexity lemma. This follows from clause (c) of the convexity lemma together with
one further bound: the one-sided bound $\bar{\delta}$ on the mismatch of vacuum-energy densities. That
bound is used for the single cubic pair $(\mathbb{T}_{m_{\mathrm{sm}}},\mathbb{T}_{m_{\mathrm{sm}}+1})$
only. Consequently
\begin{equation*}
0\le\pi(\mathbb{T}_{m_{\mathrm{sm}}})-\pi(\mathfrak{B})\le\mathsf{c}_L(\mathsf{B}+\bar{\delta}),
\end{equation*}
and so
\begin{equation*}
\exp\bigl\{\ell_n^4\,[\pi(\mathbb{T}_{m_{\mathrm{sm}}})-\pi(\mathfrak{B})]\bigr\}
\le\exp\bigl\{\mathsf{c}_L(\mathsf{B}+\bar{\delta})\,\ell_n^4\bigr\}.
\end{equation*}
This factor, the term $\mathsf{B}\ell_n^4$ of \eqref{9.98:eq-proof-rectangular-tori-3}, and the term
$2nC_{\mathrm{face}}$ of the face step give the exponent of
\begin{equation*}
Q^{\mathrm{cvx}}_{\star,n}
=\exp\bigl[(1+\mathsf{c}_L)(\mathsf{B}+\bar{\delta})\,\ell_n^4+2nC_{\mathrm{face}}\bigr].
\end{equation*}

The face step and the concluding step are used without change.

Hence the transfer inequality \eqref{9.98:eq-transfer-inequality-a} holds between $\mathfrak{B}$ and
$\mathbb{T}_{m_{\natural}(g)}$ with $C^{\star\mathrm{cvx}}_{T,n}$ in place of $C^{\star}_{T,n}$, and the
volume-transfer bound holds between them with $C^{\star\mathrm{cvx}}_{VT,n}$ in place of
$C^{\star}_{VT,n}$. Between two such tori $\mathfrak{B}$, $\mathfrak{B}'$ the transfer inequality
follows through $\mathbb{T}_{m_{\natural}(g)}$ by the triangle inequality, with the constant
$2C^{\star\mathrm{cvx}}_{T,n}$ in place of $C^{\star}_{T,n}$.

On $\mathfrak{B}$ nothing of Ba{\l}aban's construction is used: no clause (i), no floor, and no vacuum
energy $E_{\mathfrak{B}}$ of $\mathfrak{B}$. Only the following properties of its bare Wilson measure
enter: link reflection positivity, the chessboard estimate, the monotonicity in the cell counts, and the
convexity lemma.

\emph{Limit points on $\mathfrak{B}$.} Put
\begin{equation*}
\tau^{\mathrm{cvx}}_n:=C^{\star\mathrm{cvx}}_{T,n}\,L^{-(4-n)\underline{s}'(g)}\,
\mathsf{d}^{4-n}\,\Pi,
\end{equation*}
that is, $\tau_n$ of \eqref{9.98:eq-short-distance-uniqueness-a} with $C^{\star\mathrm{cvx}}_{T,n}$ in
place of $C^{\star}_{T,n}$. At every tuned depth admitted in \eqref{9.98:eq-transfer-inequality-a}, the
bound just proved gives, both members being taken at the same depth,
\begin{equation*}
\bigl|S^{\,\cdot\,}_{\mathfrak{B}}(x_0;\vec f\,)
-S^{\,\cdot\,}_{\mathbb{T}_{m_{\natural}(g)}}(x_0;\vec f\,)\bigr|\le\tau^{\mathrm{cvx}}_n .
\end{equation*}
Thus $S^{\,\cdot\,}_{\mathfrak{B}}(x_0;\vec f\,)$ lies in the closed interval of half-length
$\tau^{\mathrm{cvx}}_n$ around $S^{\,\cdot\,}_{\mathbb{T}_{m_{\natural}(g)}}(x_0;\vec f\,)$. By
Lemma~\ref{9.98:lem-reference-torus-limit} these centres converge, as the depth tends to infinity, to
$\mathfrak{S}_{m_{\natural}(g)}[\hat x](\vec f\,)$ (display \eqref{9.98:eq-reference-torus-limit-a}).
The numbers $S^{\,\cdot\,}_{\mathfrak{B}}(x_0;\vec f\,)$, indexed by the depth, therefore form a bounded
sequence, and cutoff-limit points on $\mathfrak{B}$ of the tuned family at $\hat w$ exist along
subsequences.

Let $\sigma_{\mathfrak{B}}(\vec f\,)$ be the limit along such a subsequence. Every term of the
subsequence lies in its interval
\begin{equation*}
\bigl[S^{\,\cdot\,}_{\mathbb{T}_{m_{\natural}(g)}}(x_0;\vec f\,)-\tau^{\mathrm{cvx}}_n,\;
S^{\,\cdot\,}_{\mathbb{T}_{m_{\natural}(g)}}(x_0;\vec f\,)+\tau^{\mathrm{cvx}}_n\bigr].
\end{equation*}
The end points converge, and the intervals are closed, so in the limit
\begin{equation}\label{9.98:eq-proof-rectangular-tori-4}
\bigl|\sigma_{\mathfrak{B}}(\vec f\,)-\mathfrak{S}_{m_{\natural}(g)}[\hat x](\vec f\,)\bigr|
\le\tau^{\mathrm{cvx}}_n .
\end{equation}
For two such limit points $\sigma_{\mathfrak{B}}$, $\sigma_{\mathfrak{B}'}$ on tori $\mathfrak{B}$,
$\mathfrak{B}'$ as above, the triangle inequality applied to \eqref{9.98:eq-proof-rectangular-tori-4}
gives
\begin{equation*}
\bigl|\sigma_{\mathfrak{B}}(\vec f\,)-\sigma_{\mathfrak{B}'}(\vec f\,)\bigr|\le 2\tau^{\mathrm{cvx}}_n .
\end{equation*}
With \eqref{9.98:eq-proof-rectangular-tori-4} in place of the corresponding bound with $\tau_n$, the
arguments for clauses (b)--(e) of Theorem~\ref{9.98:thm-short-distance-uniqueness} apply to
$\sigma_{\mathfrak{B}}$ word for word, with $\tau^{\mathrm{cvx}}_n$ in place of $\tau_n$.

\emph{Thermal boxes.} Consider $\mathbb{R}^3$ times a circle of circumference $\beta$, where
$\beta\in 2\ell_n L^{\overline{s}'(g)}\mathbb{N}$ and $\beta\ge 2L^{m_{\natural}(g)}$. Take joint limits
along tori $\mathfrak{B}_j$ satisfying \eqref{9.98:eq-proof-rectangular-tori-1}, with three sides tending
to infinity and the fourth side equal to $\beta$. The constant $\tau^{\mathrm{cvx}}_n$ does not depend on
the sides of $\mathfrak{B}$, so the argument above applies to every $\mathfrak{B}_j$: joint limit points
exist along subsequences, and each of them satisfies the bound \eqref{9.98:eq-proof-rectangular-tori-4}
and the conclusions of clauses (b)--(e) of Theorem~\ref{9.98:thm-short-distance-uniqueness}.

Convergence of the full tuned family as the cutoff is removed on a non-cubic $\mathfrak{B}$ is not
asserted; the uniqueness theorem concerns one cubic torus at a time.

The constants $C^{\star\mathrm{cvx}}_{T,n}$ and $C^{\star\mathrm{cvx}}_{VT,n}$ exceed $C^{\star}_{T,n}$
and $C^{\star}_{VT,n}$ by a pure factor. This factor comes from $\mathsf{c}_L$ in the exponent of
$Q^{\mathrm{cvx}}_{\star,n}$. Nothing in what follows depends on which of the two sets of constants is
used.
\end{proof}

\begin{proof}[Proof of (e): the first-passage frame]
We prove the two assertions of clause~(e) of
Theorem~\ref{9.98:thm-short-distance-uniqueness}, in the first-passage labelling of
clause~(ii) of Theorem~\ref{3.1:thm-main} on $\mathbb{R}^4$. Fix $n\in\{2,3\}$, a separation
$\mathsf{d}$ with $0<\mathsf{d}\le L^{\underline{s}(g)}/40$, and a tuple
$\vec f\in\mathfrak{P}_n(\mathsf{d};\vartheta)$ in $g$-reference units, where $\underline{s}(g)$
is the lower bound on the prolongation fixed in \eqref{9.98:eq-prolongation-depths-a}. Put
$\Pi:=\prod_i\|f_i\|_\infty$. For $\sigma\in\operatorname{Lim}_{\infty}(g)$ we write
$S_n[\sigma](\vec f\,):=\sigma(\vec f\,)$ for the $n$-point function of the limit point $\sigma$
evaluated at $\vec f$.

\emph{First assertion.} Let $\sigma\in\operatorname{Lim}_{\infty}(g)$. By the definition of
$\operatorname{Lim}_{\infty}(g)$ through clause~(ii) of Theorem~\ref{3.1:thm-main}, there are
bare couplings $g_0^{(j)}\downarrow 0$, depths
$K_j:=K(g_0^{(j)},g)\to\infty$ (the first-passage depth of
Definition~\ref{1.2:def-flow-constants}) and torus exponents $m_j\to\infty$, all indexed by one
common index $j$ (a joint diagonal), such that
\begin{equation}\label{9.98:eq-proof-first-passage-1}
\sigma(\vec f\,)=\lim_{j\to\infty}S^T_{n;K_j,m_j}[g_0^{(j)}](\vec f\,).
\end{equation}
Passing to a subsequence does not change the limit in \eqref{9.98:eq-proof-first-passage-1};
we shall do so finitely many times, without renaming. Since $K_j\to\infty$, after discarding finitely many
indices and thinning we may assume that the $K_j$ are strictly increasing. Since
$g_0^{(j)}\to0$, we may also assume $g_0^{(j)}\in\,]0,g_{-}(g)]$ for all $j$.

\emph{Step (i): the window.} Put $x^{(j)}:=(g_0^{(j)})^{-2}$. This is a bare inverse coupling
whose first-passage run stops at depth $K_j$. Clause~(v-c)(2) of Theorem~\ref{3.1:thm-main}
(the corollary of the uniqueness theorem) applies, because its two conditions, the floor
$\varrho\ge 6(B^\sharp+1)$ and the ceiling $g\le\gamma_U'$, are among the hypotheses of
Theorem~\ref{9.98:thm-short-distance-uniqueness}. It gives $|x^{(j)}-a_{K_j}|\le\varrho_2/2$,
that is,
\begin{equation*}
\hat w_j:=a_{K_j}-x^{(j)}\in[-\varrho_2/2,\,\varrho_2/2].
\end{equation*}
The interval $[-\varrho_2/2,\varrho_2/2]$ is compact, so after passing to a subsequence
$\hat w_j\to\hat w$ for some $\hat w\in[-\varrho_2/2,\varrho_2/2]$. Let $\hat x$ be the
renormalised reference coupling that corresponds to $\hat w$, as in the statement of
Theorem~\ref{9.98:thm-short-distance-uniqueness}.

\emph{Step (ii): the transfer.} We read the lattice integrals in the first-passage frame of
Lemma~\ref{9.98:lem-tuned-member} (see \eqref{9.98:eq-tuned-member-a}), in which the native
first-passage depth is $K_0=K_j$ and $\jmath=0$. We have $g\le\gamma_{\star}$, since
$g\le\gamma_{\natural}$ and
$\gamma_{\natural}^{-2}=\gamma_{\star}^{-2}+B^\sharp(s_{\star}+j_U+1)\ge\gamma_{\star}^{-2}$;
and $g_0^{(j)}\in\,]0,g_{-}(g)]$. Hence Lemma~\ref{9.98:lem-run-prolongation} applies to
$(g_0^{(j)},g)$, and its part~(a), display \eqref{9.98:eq-run-prolongation-a}, gives for the
prolongation
\begin{equation*}
\mathsf{s}_j:=K(g_0^{(j)},\gamma_{\star})-K_j
\end{equation*}
the bounds $\underline{s}(g)\le\mathsf{s}_j\le\overline{s}(g)$. Since $K_j\to\infty$ and
$m_j\to\infty$, for all large $j$ we have $L^{-K_j}\le\mathsf{d}$, $K_j\ge n_{\natural}$ and
$m_j\ge m_{\star}+\overline{s}(g)$. For such $j$ we apply the transfer inequality,
Proposition~\ref{9.98:prop-transfer-inequality}, display \eqref{9.98:eq-transfer-inequality-a},
in the first-passage frame ($K_0=K_j$, $\jmath=0$), with $(m_1,m_2)=(m_j,m_{\natural}(g))$ and
depth $K_j$. We check its conditions. The exponents of the
two tori in the $\gamma_{\star}$-reference frame are $m_j-\mathsf{s}_j$ and
$m_{\natural}(g)-\mathsf{s}_j$. The first satisfies
$m_j-\mathsf{s}_j\ge m_j-\overline{s}(g)\ge m_{\star}$. For the second, the definition of
$m_{\natural}(g)$ in \eqref{9.98:eq-prolongation-depths-c} together with
$\mathsf{s}_j\le\overline{s}(g)$ gives
\begin{equation*}
m_{\natural}(g)-\mathsf{s}_j\ge m_{\star}+j_U\ge m_{\star}.
\end{equation*}
The depth satisfies $L^{-K_j}\le\mathsf{d}$ and $K_j\ge n_{\natural}$. The scale of the tuple in
the $\gamma_{\star}$ frame, $\mathsf{d}'=L^{-\mathsf{s}_j}\mathsf{d}$, satisfies
\begin{equation*}
\mathsf{d}'=L^{-\mathsf{s}_j}\mathsf{d}\le L^{-\underline{s}(g)}\mathsf{d}\le\tfrac{1}{40},
\end{equation*}
by $\mathsf{s}_j\ge\underline{s}(g)$ and $\mathsf{d}\le L^{\underline{s}(g)}/40$. Finally
$\vec f\in\mathfrak{P}_n(\mathsf{d};\vartheta)$. The transfer inequality therefore gives
\begin{equation}\label{9.98:eq-proof-first-passage-2}
\begin{split}
\bigl|S^T_{n;K_j,m_j}[g_0^{(j)}](\vec f\,)-S^T_{n;K_j,m_{\natural}(g)}[g_0^{(j)}](\vec f\,)\bigr|
&\le C^{\star}_{T,n}L^{-(4-n)\mathsf{s}_j}\mathsf{d}^{4-n}\Pi\\
&\le C^{\star}_{T,n}L^{-(4-n)\underline{s}(g)}\mathsf{d}^{4-n}\Pi .
\end{split}
\end{equation}
The second inequality holds because $4-n\ge1$ for $n\in\{2,3\}$, $L>1$ and
$\mathsf{s}_j\ge\underline{s}(g)$.

\emph{Step (iii): the one torus.} Write $m_{\natural}:=m_{\natural}(g)$. The lattice integral with
data $(g_0^{(j)},K_j,\mathbb{T}_{m_{\natural}})$ is, in the tuned labelling of clause~(ii-$\mathbb{T}$-a)
of Theorem~\ref{3.1:thm-main} on $\mathbb{T}_{m_{\natural}}$, the depth-$K_j$ member of the tuned
family at bare inverse coupling $x^{(j)}=a_{K_j}-\hat w_j$. Indeed, the lattice is the same, with
$2L^{m_{\natural}+K_j}$ sites a side. The bare coupling is the same. The tuned reference length at
depth $K_j$ and the first-passage reference length at depth $K_j$ coincide, so the test functions
are the same. Hence
\begin{equation}\label{9.98:eq-proof-first-passage-3}
S^T_{n;K_j,m_{\natural}}[g_0^{(j)}](\vec f\,)
=S^{K_j}_{\mathbb{T}_{m_{\natural}}}(a_{K_j}-\hat w_j;\vec f\,).
\end{equation}
We now apply the statement of clause~(v-c)(1) on $\mathbb{T}_{m_{\natural}}$ restated in
Lemma~\ref{9.98:lem-reference-torus-limit}, display \eqref{9.98:eq-reference-torus-limit-a}. We
apply it along the strictly increasing sequence of depths $K_j$, to the real numbers $x^{(j)}$. Its
hypotheses hold: $|x^{(j)}-a_{K_j}|\le\varrho_2/2$ by step~(i);
$a_{K_j}-x^{(j)}=\hat w_j\to\hat w$; and $K_j\to\infty$. It yields
\begin{equation}\label{9.98:eq-proof-first-passage-4}
S^{K_j}_{\mathbb{T}_{m_{\natural}}}(a_{K_j}-\hat w_j;\vec f\,)\longrightarrow
\mathfrak{S}_{m_{\natural}}[\hat x](\vec f\,)\qquad(j\to\infty).
\end{equation}

\emph{Step (iv): the closed inequality.} The right side of \eqref{9.98:eq-proof-first-passage-2}
does not depend on $j$. We let $j\to\infty$ in \eqref{9.98:eq-proof-first-passage-2}. By
\eqref{9.98:eq-proof-first-passage-1}, the first term inside the absolute value tends to
$S_n[\sigma](\vec f\,)$. By \eqref{9.98:eq-proof-first-passage-3} and
\eqref{9.98:eq-proof-first-passage-4}, the second term tends to
$\mathfrak{S}_{m_{\natural}(g)}[\hat x](\vec f\,)$. Therefore
\begin{equation}\label{9.98:eq-proof-first-passage-5}
\bigl|S_n[\sigma](\vec f\,)-\mathfrak{S}_{m_{\natural}(g)}[\hat x](\vec f\,)\bigr|
\le C^{\star}_{T,n}L^{-(4-n)\underline{s}(g)}\mathsf{d}^{4-n}\Pi
\end{equation}
for the value $\hat x$ obtained in step~(i), which lies in the window of clause~(v) since
$|\hat w|\le\varrho_2/2$. The exponent in
\eqref{9.98:eq-proof-first-passage-5} is the unprimed $\underline{s}(g)$ of
\eqref{9.98:eq-prolongation-depths-a}: in the first-passage frame $\jmath=0$, and no renaming
between tuned and native depths occurs. This is the first assertion of clause~(e).

\emph{Second assertion.} Let $\sigma,\sigma'\in\operatorname{Lim}_{\infty}(g)$ be taken along joint
diagonals whose runs have reference-scale (renormalised) inverse couplings at the stopping depth
that converge to the same value. Carry out steps (i)--(iii) for each of them. This gives
renormalised couplings $\hat x$ and $\hat x'$, and two families of connected $n$-point functions
of tuned members on the one torus $\mathbb{T}_{m_{\natural}}$, namely the right sides of
\eqref{9.98:eq-proof-first-passage-3} for $\sigma$ and for $\sigma'$. These two families converge
to $\mathfrak{S}_{m_{\natural}}[\hat x](\vec f\,)$ and to
$\mathfrak{S}_{m_{\natural}}[\hat x'](\vec f\,)$ respectively. By step~(iii), the members of each
family on $\mathbb{T}_{m_{\natural}}$ have the same bare inverse couplings $x^{(j)}$ and the same
depths $K_j$ as the corresponding runs; the running couplings do not depend on the torus
(clause~(0) of Theorem~\ref{3.1:thm-main}), so the reference-scale inverse couplings of the two
families at the stopping depth are those of the two runs, and along the subsequences chosen in
steps (i)--(iii) they still converge to the same value. Clause~(v-d)(ii) of
Theorem~\ref{3.1:thm-main} states that two admissible families on a torus whose reference-scale
(renormalised) inverse couplings converge to the same value have the same limit. We apply it on
the one torus $\mathbb{T}_{m_{\natural}}$ to these two families. It gives that the two
reference-torus limits of step~(iii) coincide:
\begin{equation*}
\mathfrak{S}_{m_{\natural}}[\hat x]=\mathfrak{S}_{m_{\natural}}[\hat x'].
\end{equation*}
We apply \eqref{9.98:eq-proof-first-passage-5} to $\sigma$ with $\hat x$ and to $\sigma'$ with
$\hat x'$, and combine the two bounds by the triangle inequality through the common value
$\mathfrak{S}_{m_{\natural}}[\hat x](\vec f\,)$. This gives
\begin{equation*}
\bigl|S_n[\sigma](\vec f\,)-S_n[\sigma'](\vec f\,)\bigr|
\le 2C^{\star}_{T,n}L^{-(4-n)\underline{s}(g)}\mathsf{d}^{4-n}\Pi,
\end{equation*}
which is the second assertion of clause~(e).
\end{proof}

\begin{remark}[First-passage limit sets span a coupling window]
\label{9.98:rem-limit-set-window}
Let $g$ be a coupling in the window of clause~(v) of Theorem~\ref{3.1:thm-main} with
$g\le\min\{\gamma_{\natural},\gamma_U'\}$, as in the short-distance uniqueness corollary to
clause~(v) of Theorem~\ref{3.1:thm-main}, and let
$\operatorname{Lim}_{\infty}(g)$ be the first-passage infinite-volume limit set at $g$. Then
$\operatorname{Lim}_{\infty}(g)$ is a union over the window of renormalised couplings that the
corollary of the uniqueness theorem supplies for first-passage data at $g$: its
elements are infinite-volume limit points along first-passage data whose renormalised couplings
$\hat{x}$ range over this window, and the short-distance uniqueness corollary does not assert
that any two of its elements coincide. Consequently the
short-distance uniqueness corollary asserts no uniqueness statement for
$\operatorname{Lim}_{\infty}(g)$ at the coupling $g$ beyond clause~(e): the elements of
$\operatorname{Lim}_{\infty}(g)$ are not labelled by a single renormalised coupling.
Clause~(e), in which the two- and three-point functions of the smeared action density of every
first-passage limit point agree, in the sense and to the accuracy displayed there, with those
of the reference-torus theory $\mathfrak{S}_{m_{\natural}(g)}[\hat{x}]$ at some renormalised
coupling $\hat{x}$ in the window of clause~(v), is the exact first-passage form of clause~(b)
of the same corollary,
in which the renormalised coupling $\hat{x}$ is fixed.
\end{remark}

\begin{proof}
Let $(\hat{w}_j)_j$ be the first-passage offsets at the coupling $g$: for each index
$j$ of the sequence of cutoffs along which the first-passage data at $g$ are taken,
$\hat{w}_j$ is the first-passage offset of the first-passage datum at $g$ taken at the $j$-th
cutoff of that sequence. By the lemma on non-convergence of the
first-passage offsets, which is derived from clause~(0) of Theorem~\ref{3.1:thm-main}, from
clause~(v-a) of Theorem~\ref{3.1:thm-main} and from the corollary of the uniqueness theorem,
the sequence $(\hat{w}_j)_j$ does not converge.

By the corollary of the uniqueness theorem, each first-passage datum at $g$ is a tuned datum at
a renormalised coupling lying in the window of renormalised couplings that this corollary
supplies for first-passage data at $g$; since the offsets do
not converge, the first-passage labelling does not attach one renormalised coupling
$\hat{x}$ to the coupling $g$. The set $\operatorname{Lim}_{\infty}(g)$ is accordingly a union
over this window, and the short-distance uniqueness corollary does not assert that any two of
its elements coincide. Agreement of the
two- and three-point functions of the smeared action density of every element of
$\operatorname{Lim}_{\infty}(g)$ with those of one theory $\mathfrak{S}_{m_{\natural}(g)}[\hat{x}]$,
in the sense and to the accuracy of clause~(e), would be a statement at one fixed renormalised
coupling $\hat{x}$, which is the form of clause~(b); since the renormalised couplings of the
first-passage data range over this window, only clause~(e), which asserts this agreement for each
element at some renormalised coupling in the window of clause~(v), is available for
$\operatorname{Lim}_{\infty}(g)$. Thus clause~(e) is the exact first-passage form of clause~(b).
\end{proof}

\begin{proof}[The three-point case]\label{9.98:prf-three-point-case}
Let $n=3$. Let $g$, $\hat w$, $\mathsf{d}$ and $\vec f=(f_1,f_2,f_3)\in\mathfrak{P}_3(\mathsf{d};\vartheta)$ be as in
Theorem~\ref{9.98:thm-short-distance-uniqueness}, and put
$\Pi=\|f_1\|_\infty\|f_2\|_\infty\|f_3\|_\infty$.
Since $n$ now denotes the number of sources, we write $n'$ for the depth of the tuned member, which
Proposition~\ref{9.98:prop-transfer-inequality} denotes by the same letter as the number of sources.
Fix a tuned depth $n'\ge n_{\natural}$ with $L^{-n'}\le\mathsf{d}$.

We use the transfer inequality at $\gamma_{\star}$ (Proposition~\ref{9.98:prop-transfer-inequality}),
display~\eqref{9.98:eq-transfer-inequality-a}, at three sources. At three sources that inequality is
the separation gain for three sources (Corollary~\ref{9.98:cor-separation-gain-three}), read with
reference coupling $g^{\flat}=\gamma_{\star}$ and $R=1$. The tuple enters at the scale
\begin{equation*}
\mathsf{d}'=L^{-\mathsf{s}'}\mathsf{d},
\end{equation*}
the scale of $\vec f$ in the units of $\gamma_{\star}$.

The tuple meets the hypotheses of Corollary~\ref{9.98:cor-separation-gain-three} at this scale. Put
$f_i'=f_i(L^{\mathsf{s}'}\cdot)$.

Sup norms are unchanged by the dilation, so $\|f_i'\|_\infty=\|f_i\|_\infty$.

The condition defining $\mathfrak{P}_3$ is covariant under the change of scale. If
$\operatorname{supp}f_i\subset\overline{B}(y,5\mathsf{d})$, then
$\operatorname{supp}f_i'\subset\overline{B}(L^{-\mathsf{s}'}y,5\mathsf{d}')$. Moreover
\begin{equation*}
\operatorname{Lip}f_i'=L^{\mathsf{s}'}\operatorname{Lip}f_i
\le\frac{L^{\mathsf{s}'}\|f_i\|_\infty}{\vartheta\mathsf{d}}
=\frac{\|f_i'\|_\infty}{\vartheta\mathsf{d}'} .
\end{equation*}

The lattice spacing in the units of $\gamma_{\star}$ is $\varepsilon'=L^{-K'}$, and
$\mathsf{s}'=K'-n'$. Hence $\mathsf{d}'=L^{n'-K'}\mathsf{d}$, and the condition
$L^{-n'}\le\mathsf{d}$ on the depth is equivalent to $\varepsilon'\le\mathsf{d}'$.

Finally, $\mathsf{d}\le L^{\underline{s}'(g)}/40$ and $\mathsf{s}'\ge\underline{s}'(g)$ give
\begin{equation*}
\mathsf{d}'\le L^{\underline{s}'(g)-\mathsf{s}'}/40\le 1/40 .
\end{equation*}
Here $\mathsf{s}'\ge\underline{s}'(g)$ is the lower bound on the prolongation in the tuned frame; this
bound is the defining property of $\underline{s}'(g)$ in the notation of
Theorem~\ref{9.98:thm-short-distance-uniqueness}.

Hence, for all $m_1,m_2\ge m_{\natural}(g)$ and the tuned depth $n'$ fixed above,
\begin{equation}\label{9.98:eq-three-point-case-1}
\begin{split}
\bigl|S^{n'}_{\mathbb{T}_{m_1}}(x_0;\vec f\,)-S^{n'}_{\mathbb{T}_{m_2}}(x_0;\vec f\,)\bigr|
&\le C^{\star}_{T,3}\,L^{-\mathsf{s}'}\mathsf{d}\,\|f_1\|_\infty\|f_2\|_\infty\|f_3\|_\infty\\
&\le C^{\star}_{T,3}\,L^{-\underline{s}'(g)}\mathsf{d}\,\Pi=\tau_3 .
\end{split}
\end{equation}
The first inequality is the transfer inequality at three sources. The second uses the same lower bound
$\mathsf{s}'\ge\underline{s}'(g)$ on the prolongation together with $L>1$. The final equality is the
definition \eqref{9.98:eq-short-distance-uniqueness-a} of $\tau_n$ at $n=3$.

With \eqref{9.98:eq-three-point-case-1} in place of its two-source counterpart, items (a), (b), (c),
(e) and (g) of Theorem~\ref{9.98:thm-short-distance-uniqueness} at $n=3$ follow word for word by the
arguments of their proofs:
\begin{itemize}
\item item (a) (any two tori) follows by the argument of the proof of (a)
(proof~\ref{9.98:prf-proof-two-tori});
\item item (b) (triangle through the reference torus) follows by the argument of the proof of (b)
(proof~\ref{9.98:prf-proof-triangle});
\item item (c) (the two limits commute) follows by the argument of the proof of (c)
(proof~\ref{9.98:prf-proof-limits-commute});
\item item (e) (the first-passage frame) follows by the argument of the proof of (e);
\item item (g) (rectangular tori) follows by the argument of the proof of (g)
(proof~\ref{9.98:prf-proof-rectangular-tori}).
\end{itemize}
In (a), (b), (c) and (e), $\tau_3$ replaces $\tau_2$; in (g) the constant of the proof of (g), taken at
three sources, replaces its two-source value. On $\mathbb{R}^4$ the passage to infinite-volume limit
points goes
through the closed interval of the proof of (b) (proof~\ref{9.98:prf-proof-triangle}), here applied to
the difference of the three-point functions on two volumes.
\end{proof}

\begin{proposition}[The two-point law at every short separation]\label{9.98:prop-two-point-law}
Assume the following three statements.
\begin{enumerate}
\item[(I)] The conclusions \eqref{9.98:eq-run-prolongation-a} and
\eqref{9.98:eq-run-prolongation-b} of Lemma~\ref{9.98:lem-run-prolongation} (run prolongation to
the strong reference coupling). For $g\le\gamma_{\star}$, $g_0\in\,]0,g_-(g)]$, $K=K(g_0,g)$,
$K'=K(g_0,\gamma_{\star})$ and $\mathsf{s}=K'-K$, these say that $g_0\le g_-(\gamma_{\star})$, that
$K'\ge K$, and that $\underline{s}(g)\le\mathsf{s}\le\overline{s}(g)$. They also say that for $n\ge2$,
$m\ge\mathsf{s}+m_{\star}$ and all $\vec f$ on $\mathbb{T}_m$ one has
\[
S^T_{n;K,m}[g_0](\vec f\,)=S^T_{n;K',m-\mathsf{s}}[g_0](\vec f(L^{\mathsf{s}}\cdot)).
\]
\item[(II)] Clause~(iv) of Theorem~\ref{3.1:thm-main}, uniform in the volume,
applied once, at the
strong reference coupling $\gamma_{\star}$ of Definition~\ref{9.98:def-strong-reference-coupling}.
This includes its two-point display, its remainder bound $|R|\le\frac12$, its kernel read on the
reference torus $\mathbb{T}_{\star}$ with the bounds $c_{\mathcal{K}}$, $C_{\mathcal{K}}$, and its
uniformity statement. It is used together with part~(iii-c) of clause~(iii) of
Theorem~\ref{3.1:thm-main}, which
places the running shift in $[\frac23,\frac43]$.
\item[(III)] Clause~(0) of Theorem~\ref{3.1:thm-main}. This consists of the first-passage depth
$K(g_0,\cdot)$, the
one-step bounds \eqref{9.98:eq-couplings-limit-sets-a}, their corollary
\eqref{9.98:eq-couplings-limit-sets-b}, and the independence of the running couplings from the
torus.
\end{enumerate}
Then (a)--(d) below hold.

\smallskip\noindent\emph{(a) First-passage frame.}
Let $\vartheta\in\,]0,1]$. Let $g\le\gamma_{\natural}(\vartheta)$, with $\gamma_{\natural}$ as in
\eqref{9.98:eq-prolongation-depths-b}. Let $g_0\in\,]0,g_-(g)]$, $K:=K(g_0,g)$ and
$\mathsf{s}:=K(g_0,\gamma_{\star})-K$, so that
$\mathsf{s}\in[\underline{s}(g),\overline{s}(g)]$ by (I), which applies since
$\gamma_{\natural}(\vartheta)\le\gamma_{\star}$ by \eqref{9.98:eq-prolongation-depths-b}. Let
$m\ge m_{\star}+\overline{s}(g)$; in
particular any $m\ge m_{\natural}(g)$, with $m_{\natural}(g)$ as in
\eqref{9.98:eq-prolongation-depths-c}, is admitted. Let $\mathsf{d}$ satisfy
$0<\mathsf{d}\le 1/40$, a condition in which $g$ does not enter. Put
\[
u=u(\mathsf{d};g_0,g):=s_{\gamma_{\star}}(L^{-\mathsf{s}}\mathsf{d})-\mathsf{s},
\]
and assume
\begin{equation}\label{9.98:eq-two-point-law-b}
K-u\ge k_W .
\end{equation}
Then $u\ge1$, and every $(f_1,f_2)\in\mathfrak{T}(\mathsf{d};\vartheta)$ on $\mathbb{T}_m$ satisfies
\begin{equation}\label{9.98:eq-two-point-law-1}
S^T_{2;K,m}(f_1,f_2)=\zeta'^{\,2}_{K-u}\,b_0\,g_{K-u}^4\,
\mathcal{K}^{\star}_u(f_1,f_2)\,(1+R),\qquad |R|\le\tfrac12 .
\end{equation}
The symbols in \eqref{9.98:eq-two-point-law-1} are as follows.
\begin{itemize}
\item $g_{K-u}$ and $\zeta'_{K-u}$ are the running coupling and the running shift at level $K-u$
of the same run, namely the one started at $g_0$. One has $\zeta'_{K-u}\in[\frac23,\frac43]$ by
part~(iii-c) of clause~(iii).
\item $R$ is the remainder of clause~(iv).
\item $\mathcal{K}^{\star}_u(f_1,f_2)$ is the kernel of clause~(iv) at
reference coupling $\gamma_{\star}$ and depth $u+\mathsf{s}$ (counted at reference $\gamma_{\star}$).
It is read on the reference torus $\mathbb{T}_{\star}$ and evaluated on
$(f_1(L^{\mathsf{s}}\cdot),f_2(L^{\mathsf{s}}\cdot))$, in reference units of $\gamma_{\star}$.
Concretely, one takes the free lattice gauge field at unit coupling on the bare lattice of
$\mathbb{T}_{\star}$. Its side is $2L^{m_{\star}}$ reference lengths of $\gamma_{\star}$, with
$K+\mathsf{s}=K(g_0,\gamma_{\star})$ levels, equivalently $2L^{m_{\star}+\mathsf{s}}$ reference
lengths of $g$, with $K$ levels. This field is block-averaged to level
$K-u=(K+\mathsf{s})-(u+\mathsf{s})$; the kernel is
$(\mathbb{E}[F_BF_B])^2$, with $\mathbb{E}$ the expectation of the free field, paired with
$(f_1,f_2)$ in reference units of $g$.
\item The kernel satisfies
\[
c_{\mathcal{K}}\|f_1\|_\infty\|f_2\|_\infty\le\mathcal{K}^{\star}_u(f_1,f_2)\le
C_{\mathcal{K}}\|f_1\|_\infty\|f_2\|_\infty ,
\]
with the constants of clause~(iv), which do not depend on the torus.
\end{itemize}
Hence, uniformly in $K$, $g_0$, $m\ge m_{\star}+\overline{s}(g)$ and the position of the supports,
\begin{equation}\label{9.98:eq-two-point-law-2}
\begin{split}
\tfrac29\,b_0c_{\mathcal{K}}\,\bigl(g^{-2}+B^{\sharp}(u+1)\bigr)^{-2}\,
\|f_1\|_\infty\|f_2\|_\infty
&\le S^T_{2;K,m}(f_1,f_2)\\
&\le \tfrac83\,b_0C_{\mathcal{K}}\,\bigl(g^{-2}+\lambda^{\sharp}u-c_5\bigr)^{-2}\,
\|f_1\|_\infty\|f_2\|_\infty .
\end{split}
\end{equation}
Since $u\le\mathsf{U}(g,\mathsf{d})$, where $\mathsf{U}$ depends on $g$ and $\mathsf{d}$ only,
see \eqref{9.98:eq-prolongation-depths-d}, the following bound holds and is free of $g_0$:
\begin{equation}\label{9.98:eq-two-point-law-3}
\begin{split}
\tfrac29\,b_0c_{\mathcal{K}}\,\bar{x}^{\natural}(\mathsf{d})^{-2}\,\|f_1\|_\infty\|f_2\|_\infty
&\le S^T_{2;K,m}(f_1,f_2)\\
&\le \tfrac83\,b_0C_{\mathcal{K}}\,\bigl(g^{-2}+\lambda^{\sharp}-c_5\bigr)_+^{-2}\,
\|f_1\|_\infty\|f_2\|_\infty .
\end{split}
\end{equation}
Here $u\ge1$ is used, the right member is taken in its crudest form free of $g_0$, and
$(\eta)_+^{-2}$ is read as $+\infty$ when $\eta\le0$.

\smallskip\noindent\emph{(b) Consequences.} Under hypothesis~(II), and together with clauses~(0)
and~(ii) of Theorem~\ref{3.1:thm-main}, reflection positivity of the lattice measure and the
continuum kernel, the
displays \eqref{9.98:eq-two-point-law-1}--\eqref{9.98:eq-two-point-law-3} yield, by the derivation
by which clause~(iv) of Theorem~\ref{3.1:thm-main} draws them from its own two-point display, the
four consequences
of that clause --- that the curvature Hilbert space $\mathcal{H}^{\mathrm{curv}}$ has dimension at
least two, and the three further consequences stated in that clause --- in the following limits of
clause~(ii) of Theorem~\ref{3.1:thm-main}:
\begin{itemize}
\item in every such limit on a torus $\mathbb{T}_m$ with $m\ge m_{\star}+\overline{s}(g)$;
\item in every limit of clause~(ii) on $\mathbb{R}^4$.
\end{itemize}
They hold at every separation $\mathsf{d}\le1/40$ of the reference length of $g$. The separation
ceiling depending on $g$, which clause~(iv) carries at reference $g$
through its depth condition --- an upper bound on $\mathsf{d}$ depending on $g$, whose removal for
the clause read at reference $g$ is not asserted --- is replaced here by two other conditions.
One is the coupling ceiling
$g\le\gamma_{\natural}$. The other is the volume floor $m\ge m_{\star}+\overline{s}(g)$, which is no
restriction for limits on $\mathbb{R}^4$. On tori with
$m_{\mathbb{T}}(g_-(g))\le m<m_{\star}+\overline{s}(g)$ nothing beyond clause~(iv) as stated in
Theorem~\ref{3.1:thm-main} is asserted here, and there is no conflict with that clause.

\smallskip\noindent\emph{(c) Tuned frame.}
Let $\vartheta$ and $g$ be as in~(a), with moreover $g\le\gamma_U'$, let
$\hat{w}\in[-\varrho_2/2,\varrho_2/2]$, let $0<\mathsf{d}\le1/40$ and let $m\ge m_{\natural}(g)$.
Consider the member of depth $n\ge n_{\natural}$ of the tuned family of clause~(v) of
Theorem~\ref{3.1:thm-main} on
$\mathbb{T}_m$, with $L^{-n}\le\mathsf{d}$, in the sense of Lemma~\ref{9.98:lem-tuned-member}. Its
bare inverse coupling is $x_0:=a_n(g)-\hat{w}$; put $g_0:=x_0^{-1/2}$ and $K_0:=K(g_0,g)$. Put
\[
\mathsf{s}:=K(g_0,\gamma_{\star})-K_0,\qquad \mathsf{s}':=K(g_0,\gamma_{\star})-n,
\]
and let $\tilde s:=s_{\gamma_{\star}}(L^{-\mathsf{s}'}\mathsf{d})$ be the depth of the pair at
reference $\gamma_{\star}$. Put $u:=\tilde s-\mathsf{s}$ and assume $K_0-u\ge k_W$, the analogue of
\eqref{9.98:eq-two-point-law-b}; at fixed $\mathsf{d}$ this holds for all $n$ large enough (see
the proof). Then \eqref{9.98:eq-two-point-law-1} holds for every
$(f_1,f_2)\in\mathfrak{T}(\mathsf{d};\vartheta)$ with the following changes:
\begin{itemize}
\item $S^n_{\mathbb{T}_m}(x_0;f_1,f_2)$ stands in place of $S^T_{2;K,m}(f_1,f_2)$, and $K_0$ in
place of $K$;
\item the running coupling $g_{K_0-u}$ is that of the run from $x_0$ itself, at level $K_0-u$, and
\[
u\in\bigl[\,2-j_U,\ \mathsf{U}(g,\mathsf{d})\,\bigr];
\]
\item the kernel is that of clause~(iv) at reference $\gamma_{\star}$ and depth $\tilde s$, read on
$\mathbb{T}_{\star}$ and evaluated on $(f_1(L^{\mathsf{s}'}\cdot),f_2(L^{\mathsf{s}'}\cdot))$.
\end{itemize}
Moreover, the following bound holds and is free of $g_0$, $n$ and $m$:
\begin{equation}\label{9.98:eq-two-point-law-4}
\begin{split}
\tfrac29\,b_0c_{\mathcal{K}}\,\bar{x}^{\natural}(\mathsf{d})^{-2}\,\|f_1\|_\infty\|f_2\|_\infty
&\le S^n_{\mathbb{T}_m}(a_n-\hat{w};f_1,f_2)\\
&\le \tfrac83\,b_0C_{\mathcal{K}}\,(\underline{x}^{\natural})^{-2}\,
\|f_1\|_\infty\|f_2\|_\infty ,
\end{split}
\end{equation}
where, as in \eqref{9.98:eq-prolongation-depths-d},
$\underline{x}^{\natural}=(g^{-2}-B^{\sharp}j_U-c_5)_+$. If $X_{\star}>B^{\sharp}j_U+c_5$ then
$\underline{x}^{\natural}>0$, because $g^{-2}\ge X_{\star}$. If $X_{\star}\le B^{\sharp}j_U+c_5$, the
right member is read as $+\infty$ and only the left member is asserted.

\smallskip\noindent\emph{(d) Limits.}
For a limit $\sigma$ write $S_2[\sigma](f_1,f_2)$ for its connected two-point function of
$\mathcal{O}(f_1)$, $\mathcal{O}(f_2)$. Let $\hat{x}$ be the renormalised coupling of the tuned
family of~(c) at $(g,\hat{w})$, so that the limit set $\operatorname{Lim}_{\infty}[\hat{x}]$ and
the limits $\mathfrak{S}_m[\hat{x}]$ of Notation~\ref{9.98:not-couplings-limit-sets} are limits of
members of that family. The following holds for every
$\sigma\in\operatorname{Lim}_{\infty}[\hat{x}]$, every $\sigma=\mathfrak{S}_m[\hat{x}]$ with
$m\ge m_{\natural}(g)$, every $\mathsf{d}\le1/40$ and
every pair $(f_1,f_2)\in\mathfrak{T}(\mathsf{d};\vartheta)$ of continuously differentiable functions:
\begin{equation}\label{9.98:eq-two-point-law-a}
\begin{split}
\tfrac29\,b_0c_{\mathcal{K}}\,\bar{x}^{\natural}(\mathsf{d})^{-2}\,\|f_1\|_\infty\|f_2\|_\infty
&\le S_2[\sigma](f_1,f_2)\\
&\le \tfrac83\,b_0C_{\mathcal{K}}\,(\underline{x}^{\natural})^{-2}\,
\|f_1\|_\infty\|f_2\|_\infty .
\end{split}
\end{equation}
The pairs are taken continuously differentiable because clause~(ii) of Theorem~\ref{3.1:thm-main}
on $\mathbb{R}^4$ is
stated for compactly supported continuously differentiable test functions. Consequently, under
hypothesis~(II) and together with clauses~(0) and~(ii) of Theorem~\ref{3.1:thm-main}, reflection
positivity of the
lattice measure and the continuum kernel, the four consequences listed in~(b) hold in every limit
point on $\mathbb{R}^4$ in the tuned frame as well, at every $\mathsf{d}\le1/40$.
\end{proposition}

\begin{proof}
\emph{Proof of (a). Step 1: two namings of one integral.}
Fix the data of~(a). By (I), $\mathsf{s}\le\overline{s}(g)$, hence
$m\ge m_{\star}+\overline{s}(g)\ge m_{\star}+\mathsf{s}$. The identity
\eqref{9.98:eq-run-prolongation-b} therefore applies with $n=2$ and gives
\begin{equation}\label{9.98:eq-two-point-law-5}
S^T_{2;K,m}[g_0](f_1,f_2)=S^T_{2;K',m-\mathsf{s}}[g_0](\tilde f_1,\tilde f_2),
\qquad \tilde f_i:=f_i(L^{\mathsf{s}}\cdot),\quad K'=K+\mathsf{s}.
\end{equation}

\emph{Step 2: the rescaled pair is admissible.}
We claim $(\tilde f_1,\tilde f_2)\in\mathfrak{T}(L^{-\mathsf{s}}\mathsf{d};\vartheta)$. Every
condition defining $\mathfrak{T}$ is covariant under the dilation $x\mapsto L^{\mathsf{s}}x$, as
follows.
\begin{itemize}
\item If $f_i$ is supported in a ball of radius $\mathsf{d}$, then $\tilde f_i$ is supported in a ball
of radius $L^{-\mathsf{s}}\mathsf{d}$.
\item The masses scale as $\int\tilde f_i=L^{-4\mathsf{s}}\int f_i$. Since
$\int f_i\ge\vartheta\|f_i\|_\infty\mathsf{d}^4$ and $\|\tilde f_i\|_\infty=\|f_i\|_\infty$,
\[
\int\tilde f_i\ \ge\ \vartheta\,\|\tilde f_i\|_\infty\,(L^{-\mathsf{s}}\mathsf{d})^4 .
\]
\item The Lipschitz constants scale as $\operatorname{Lip}\tilde f_i=L^{\mathsf{s}}\operatorname{Lip}f_i$.
Since $\operatorname{Lip}f_i\le\|f_i\|_\infty/(\vartheta\mathsf{d})$,
\[
\operatorname{Lip}\tilde f_i\le\|\tilde f_i\|_\infty/(\vartheta L^{-\mathsf{s}}\mathsf{d}) .
\]
\item The separation window $\mathsf{d}\le\operatorname{dist}\le6\mathsf{d}$ becomes
$L^{-\mathsf{s}}\mathsf{d}\le\operatorname{dist}\le6L^{-\mathsf{s}}\mathsf{d}$.
\item Non-negativity and $f_i\not\equiv0$ are unchanged.
\end{itemize}

\emph{Step 3: a legitimate instance of clause~(iv) at $\gamma_{\star}$.}
We apply hypothesis~(II) at reference coupling $\gamma_{\star}$ to the datum
$(g_0,K'=K(g_0,\gamma_{\star}),\mathbb{T}_{m-\mathsf{s}})$ and to the pair $(\tilde f_1,\tilde f_2)$
of scale $\tilde{\mathsf{d}}:=L^{-\mathsf{s}}\mathsf{d}$. Each condition of the clause is met.
\begin{itemize}
\item The clause asks that the reference coupling not exceed $\gamma_W^{\mathbb{T}}$. This holds with
equality, since $\gamma_{\star}=\gamma_W^{\mathbb{T}}$ by
Definition~\ref{9.98:def-strong-reference-coupling}.
\item One has $g_0\in\,]0,g_-(\gamma_{\star})]$ by \eqref{9.98:eq-run-prolongation-a}.
\item One has $K'=K(g_0,\gamma_{\star})$ by the definition of $\mathsf{s}$.
\item The clause holds on every torus whose exponent is at least
$m_{\mathbb{T}}(g_-(\gamma_{\star}))=m_{\star}$ (Definition~\ref{9.98:def-prolongation-depths}).
By Step~1, $m-\mathsf{s}\ge m_{\star}$.
\item The depth condition at $\gamma_{\star}$ asks $\tilde s:=s_{\gamma_{\star}}(\tilde{\mathsf{d}})\ge s_{\star}$,
where
$s_{\star}=\max\{s_W(\gamma_{\star}),s_{\mathrm{vol}}^{\star}(\gamma_{\star}),s_T(\gamma_{\star})\}$ by
\eqref{9.98:eq-floor-members-b}. By the definition of the depth function of clause~(iv),
$C_W(\cdot)L^{-\tilde s}\le\tilde{\mathsf{d}}$. Since $C_W\ge1$,
\[
L^{-\tilde s}\le C_W(\cdot)L^{-\tilde s}\le\tilde{\mathsf{d}}=L^{-\mathsf{s}}\mathsf{d}
\le L^{-\mathsf{s}}/40<L^{-\mathsf{s}} ,
\]
so $\tilde s>\mathsf{s}$, that is, $\tilde s\ge\mathsf{s}+1$. For $g\le\gamma_{\natural}$, the
definitions \eqref{9.98:eq-prolongation-depths-a} and \eqref{9.98:eq-prolongation-depths-b} give
$\underline{s}(g)\ge s_{\star}+j_U$. Hence
\[
\tilde s\ \ge\ \mathsf{s}+1\ \ge\ \underline{s}(g)+1\ \ge\ s_{\star}+j_U+1\ \ge\ s_{\star},
\]
whatever $\mathsf{d}\le1/40$ is. At this point the separation ceiling depending on $g$ disappears:
the prolongation $\mathsf{s}$ alone supplies all the depth that the clause requires at
$\gamma_{\star}$.
\item The ultraviolet condition of the clause at $\gamma_{\star}$ asks $K'-\tilde s\ge k_W$. With
$u=\tilde s-\mathsf{s}$ this reads $K-u\ge k_W$, which is exactly \eqref{9.98:eq-two-point-law-b}.
Moreover $u\ge1$, because $\tilde s\ge\mathsf{s}+1$.
\end{itemize}

\emph{Step 4: the display.}
The conclusion of the clause at reference $\gamma_{\star}$, depth $\tilde s$, on
$\mathbb{T}_{m-\mathsf{s}}$ is the following. The kernel of the clause at depth $\tilde s$ is read on
the reference torus $\mathbb{T}_{m_{\mathbb{T}}(g_-(\gamma_{\star}))}=\mathbb{T}_{\star}$ and
evaluated on $(\tilde f_1,\tilde f_2)$. By the definition in~(a) this value is
$\mathcal{K}^{\star}_u(f_1,f_2)$, since $\tilde s=u+\mathsf{s}$. The clause then gives
\begin{equation}\label{9.98:eq-two-point-law-6}
S^T_{2;K',m-\mathsf{s}}(\tilde f_1,\tilde f_2)=\zeta'^{\,2}_{K'-\tilde s}\,b_0\,
g_{K'-\tilde s}^4\,\mathcal{K}^{\star}_u(f_1,f_2)\,\bigl(1+R_{K',m-\mathsf{s}}\bigr),
\qquad|R_{K',m-\mathsf{s}}|\le\tfrac12 .
\end{equation}
It also gives
\[
c_{\mathcal{K}}\|\tilde f_1\|_\infty\|\tilde f_2\|_\infty\le\mathcal{K}^{\star}_u(f_1,f_2)\le
C_{\mathcal{K}}\|\tilde f_1\|_\infty\|\tilde f_2\|_\infty .
\]

Next, $K'-\tilde s=K+\mathsf{s}-\tilde s=K-u$. By clause~(0), the running inverse couplings form one
flow $(x_k)_k$, determined by $g_0$ and independent of the torus. The reference couplings $g$ and
$\gamma_{\star}$ only decide at which level one stops counting. Hence the running coupling
$g_{K'-\tilde s}$ and the running shift $\zeta'_{K'-\tilde s}$ in
\eqref{9.98:eq-two-point-law-6} are $g_{K-u}$ and $\zeta'_{K-u}$ of the same run. Further,
$\|\tilde f_i\|_\infty=\|f_i\|_\infty$. Combining \eqref{9.98:eq-two-point-law-5} with
\eqref{9.98:eq-two-point-law-6} and putting $R:=R_{K',m-\mathsf{s}}$ gives
\eqref{9.98:eq-two-point-law-1}, together with the stated bounds on $\mathcal{K}^{\star}_u$.

\emph{Step 5: the sandwich \eqref{9.98:eq-two-point-law-2}.}
All factors on the right of \eqref{9.98:eq-two-point-law-1} are positive, with the following ranges:
\begin{itemize}
\item $\zeta'^{\,2}_{K-u}\in[\frac49,\frac{16}9]$, because $\zeta'_{K-u}\in[\frac23,\frac43]$;
\item $1+R\in[\frac12,\frac32]$;
\item the kernel value $\mathcal{K}^{\star}_u(f_1,f_2)$ lies between
$c_{\mathcal{K}}\|f_1\|_\infty\|f_2\|_\infty$ and $C_{\mathcal{K}}\|f_1\|_\infty\|f_2\|_\infty$.
\end{itemize}
Since $\frac12\cdot\frac49=\frac29$ and $\frac32\cdot\frac{16}9=\frac83$, it follows that
\begin{equation}\label{9.98:eq-two-point-law-7}
\tfrac29\,b_0c_{\mathcal{K}}\,g^4_{K-u}\,\|f_1\|_\infty\|f_2\|_\infty\le S^T_{2;K,m}(f_1,f_2)
\le\tfrac83\,b_0C_{\mathcal{K}}\,g^4_{K-u}\,\|f_1\|_\infty\|f_2\|_\infty .
\end{equation}
Now apply the corollary \eqref{9.98:eq-couplings-limit-sets-b} of clause~(0) at reference $g$, with
first passage $K=K(g_0,g)$. It gives
\[
x_{K-u}\in\bigl]\,g^{-2}+\max\{0,\lambda^{\sharp}u-c_5\},\ g^{-2}+B^{\sharp}(u+1)\,\bigr] .
\]
Since $g^4_{K-u}=x_{K-u}^{-2}$, this means
\[
\bigl(g^{-2}+B^{\sharp}(u+1)\bigr)^{-2}\le g^4_{K-u}
<\bigl(g^{-2}+\lambda^{\sharp}u-c_5\bigr)_+^{-2} .
\]
Inserting this into \eqref{9.98:eq-two-point-law-7} gives \eqref{9.98:eq-two-point-law-2}. This is
the same derivation by which clause~(iv) of Theorem~\ref{3.1:thm-main} obtains its last display from
clause~(0).

\emph{Step 6: the form \eqref{9.98:eq-two-point-law-3} free of $g_0$.}
Since $\mathsf{s}\in[\underline{s}(g),\overline{s}(g)]$, the definition of $\mathsf{U}$ in
\eqref{9.98:eq-prolongation-depths-d} (with $j=0$) gives $u\le\mathsf{U}(g,\mathsf{d})$. Hence, by
the definition of $\bar{x}^{\natural}(\mathsf{d})$ in the same display,
$g^{-2}+B^{\sharp}(u+1)\le\bar{x}^{\natural}(\mathsf{d})$. For the right member, $u\ge1$ gives
$\lambda^{\sharp}u-c_5\ge\lambda^{\sharp}-c_5$. The two bounds turn
\eqref{9.98:eq-two-point-law-2} into \eqref{9.98:eq-two-point-law-3}.

\emph{Step 7: uniformity.}
The constants are $C_W$, $s_{\star}$, $k_W$, $c_{\mathcal{K}}$, $C_{\mathcal{K}}$ and the window of
$\zeta'$. Each is either the one belonging to $\gamma_{\star}$ or pure. The exponent $m$ enters only
through the condition $m-\mathsf{s}\ge m_{\star}$. By the uniformity statement of the clause, $K$,
$g_0$ and the position of the supports enter no constant. Finally, $u$ depends on $g_0$ only through
$\mathsf{s}\in[\underline{s}(g),\overline{s}(g)]$, and this dependence is removed by $\mathsf{U}$.

\smallskip
\emph{Proof of (b).} The consequences are obtained from the displays
\eqref{9.98:eq-two-point-law-1}--\eqref{9.98:eq-two-point-law-3}, together with clauses~(0)
and~(ii) of Theorem~\ref{3.1:thm-main}, reflection positivity of the lattice measure and the
continuum kernel. The
derivation is the same as that by which clause~(iv) of Theorem~\ref{3.1:thm-main}
obtains them from its own display. The only difference is that
\eqref{9.98:eq-two-point-law-1}--\eqref{9.98:eq-two-point-law-3} hold for every
$\mathsf{d}\le1/40$ on every torus with $m\ge m_{\star}+\overline{s}(g)$.

\smallskip
\emph{Proof of (c). Step 1: the tuned member in the $\gamma_{\star}$ frame.}
As shown in the proof of Lemma~\ref{9.98:lem-tuned-depth-offset}, $x_0\ge g^{-2}+B^{\sharp}$ for
$n\ge n_{\natural}$ and $|\hat{w}|\le\varrho_2/2$. Since $g_-(g)=(g^{-2}+B^{\sharp})^{-1/2}$, this is
$g_0\le g_-(g)$. As $g\le\gamma_{\natural}(\vartheta)\le\gamma_{\star}$, hypothesis~(I) applies to
$g_0$: one has $g_0\le g_-(\gamma_{\star})$ and $\mathsf{s}\in[\underline{s}(g),\overline{s}(g)]$.
By \eqref{9.98:eq-tuned-member-a},
\[
S^n_{\mathbb{T}_m}(x_0;f_1,f_2)=S^T_{2;K',m-\mathsf{s}'}[g_0](\tilde f_1,\tilde f_2),
\]
with $x_0=a_n-\hat{w}$, $g_0=x_0^{-1/2}$, $\tilde f_i:=f_i(L^{\mathsf{s}'}\cdot)$ and
\[
K'=K(g_0,\gamma_{\star})=K_0+\mathsf{s}=n+\mathsf{s}'.
\]
The pair scale is $\tilde{\mathsf{d}}:=L^{-\mathsf{s}'}\mathsf{d}\le L^{-\mathsf{s}'}/40$. Put
$\jmath:=K_0-n$; then $|\jmath|\le j_U-1$ by Lemma~\ref{9.98:lem-tuned-depth-offset}, which applies
because $n\ge n_{\natural}$ and $|\hat{w}|\le\varrho_2/2$. Step~2 of the proof of~(a), with
$\mathsf{s}'$ in place of
$\mathsf{s}$, shows that $(\tilde f_1,\tilde f_2)\in\mathfrak{T}(\tilde{\mathsf{d}};\vartheta)$.

\emph{Step 2: legitimacy of the instance.}
The instance of hypothesis~(II) at $\gamma_{\star}$ is legitimate, exactly as in Step~3 of the proof
of~(a).
\begin{itemize}
\item $\gamma_{\star}=\gamma_W^{\mathbb{T}}$.
\item $g_0\le g_-(\gamma_{\star})$, by Step~1.
\item $m-\mathsf{s}'\ge m_{\star}$, by \eqref{9.98:eq-floor-members-a}.
\item The argument of Step~3 of the proof of~(a) gives $\tilde s\ge\mathsf{s}'+1$. By the
definitions \eqref{9.98:eq-prolongation-depths-a} and \eqref{9.98:eq-prolongation-depths-b}, for
$g\le\gamma_{\natural}$,
\[
\tilde s\ \ge\ \mathsf{s}'+1\ \ge\ \underline{s}'(g)+1\ \ge\ s_{\star}+2 .
\]
\item The ultraviolet condition asks $K'-\tilde s\ge k_W$. Since $K'=K_0+\mathsf{s}$ and
$\tilde s=u+\mathsf{s}$, one has $K'-\tilde s=K_0-u$, so this is the assumption $K_0-u\ge k_W$
of~(c). At fixed $\mathsf{d}$ it holds for all large $n$.
Indeed $K'-\tilde s=n-(\tilde s-\mathsf{s}')$. Since $\mathsf{s}'-\mathsf{s}=K_0-n=\jmath$, the
difference $\tilde s-\mathsf{s}'=s_{\gamma_{\star}}(L^{-\mathsf{s}'}\mathsf{d})-\mathsf{s}'$ equals
$u-\jmath$; by $u\le\mathsf{U}(g,\mathsf{d})$ (Step~3 below) and $|\jmath|\le j_U-1$ it satisfies
$\tilde s-\mathsf{s}'\le\mathsf{U}(g,\mathsf{d})+j_U-1$, a bound that does not depend on $n$. Hence
\[
K'-\tilde s\ \ge\ n-\mathsf{U}(g,\mathsf{d})-j_U+1 ,
\]
which is at least $k_W$ for all $n$ large enough.
\end{itemize}

\emph{Step 3: the level and the range of $u$.}
The display of the clause at $\gamma_{\star}$ carries the running coupling at level
$K'-\tilde s=K_0-u$, where $u:=\tilde s-\mathsf{s}$. As in Step~4 of the proof of~(a), by clause~(0)
this is the coupling of the one run from $x_0$ itself. The definition of $\mathsf{U}$ in
\eqref{9.98:eq-prolongation-depths-d}, with $t=\mathsf{s}\in[\underline{s}(g),\overline{s}(g)]$
(Step~1),
$j=\jmath$ and $|\jmath|\le j_U-1$, gives $u\le\mathsf{U}(g,\mathsf{d})$. On the other side,
\[
u\ \ge\ \mathsf{s}'+1-\mathsf{s}\ =\ \jmath+1\ \ge\ 2-j_U .
\]

\emph{Step 4: the window for $x_{K_0-u}$.}
There are two cases.
\begin{itemize}
\item If $u\ge0$, the corollary \eqref{9.98:eq-couplings-limit-sets-b} of clause~(0) at reference
$g$, with first passage $K_0$, gives
\[
x_{K_0-u}\in\bigl]\,g^{-2}+(\lambda^{\sharp}u-c_5)_+,\ g^{-2}+B^{\sharp}(u+1)\,\bigr]
\subset\bigl]\,\underline{x}^{\natural},\ \bar{x}^{\natural}(\mathsf{d})\,\bigr] .
\]
\item If $2-j_U\le u<0$, the level of the pair lies at most $j_U-2$ steps past $K_0$, before $K'$.
Apply the one-step bounds \eqref{9.98:eq-couplings-limit-sets-a} between $K_0$ and $K_0+|u|$,
together with $x_{K_0}\in\,]g^{-2},g^{-2}+B^{\sharp}]$. This gives
\[
x_{K_0+|u|}\in\bigl]\,g^{-2}-B^{\sharp}(j_U-2),\ g^{-2}+B^{\sharp}+c_5\,\bigr]
\subset\bigl]\,\underline{x}^{\natural},\ \bar{x}^{\natural}(\mathsf{d})\,\bigr] .
\]
For the inclusion, $\mathsf{U}\ge0$ is used, and the summand $c_5$ in the definition of
$\bar{x}^{\natural}(\mathsf{d})$ in \eqref{9.98:eq-prolongation-depths-d} serves here.
\end{itemize}

\emph{Step 5: conclusion.}
As in Step~5 of the proof of~(a), one has $\zeta'^{\,2}\in[\frac49,\frac{16}9]$ and
$1+R\in[\frac12,\frac32]$, and the kernel of~(c) lies between
$c_{\mathcal{K}}\|f_1\|_\infty\|f_2\|_\infty$ and $C_{\mathcal{K}}\|f_1\|_\infty\|f_2\|_\infty$.
By Step~4,
\[
g^4_{K_0-u}=x_{K_0-u}^{-2}\in\bigl[\,\bar{x}^{\natural}(\mathsf{d})^{-2},
\ (\underline{x}^{\natural})^{-2}\,\bigr[ .
\]
Together these give \eqref{9.98:eq-two-point-law-4}.

\smallskip
\emph{Proof of (d).}
The endpoints of \eqref{9.98:eq-two-point-law-4} depend on $(g,\mathsf{d},\vartheta)$ and on pure
constants only. They do not depend on $n$, $m$, $g_0$, $\hat{w}$ or the position of the supports.
Consider $S_2[\sigma](f_1,f_2)$ with $\sigma\in\operatorname{Lim}_{\infty}[\hat{x}]$ or
$\sigma=\mathfrak{S}_m[\hat{x}]$, $m\ge m_{\natural}(g)$. It is a limit of numbers
$S^{n_j}_{\mathbb{T}_{m_j}}(a_{n_j}-\hat{w};f_1,f_2)$ along a sequence with $n_j\to\infty$. For all
$j$ large enough one has $n_j\ge n_{\natural}$, $L^{-n_j}\le\mathsf{d}$, and the assumption
$K_0-u\ge k_W$ of~(c) holds at depth $n_j$, by Step~2 of the proof of~(c); moreover
$m_j\ge m_{\natural}(g)$, since $m_j=m$ for $\sigma=\mathfrak{S}_m[\hat{x}]$, and since the torus
exponents $m_j$ tend to infinity for $\sigma\in\operatorname{Lim}_{\infty}[\hat{x}]$. Each of these
numbers then lies in the closed interval
bounded by these endpoints. The limit therefore lies in the same closed interval, and no exchange of
limits is involved. This is \eqref{9.98:eq-two-point-law-a}. The consequences in the tuned frame then
follow from \eqref{9.98:eq-two-point-law-a} as in the proof of~(b).

\smallskip
The two frames differ only in the range of $u$. In the first-passage frame $u\ge1$ always, and the
finer right member $\frac83 b_0C_{\mathcal{K}}(g^{-2}+\lambda^{\sharp}u-c_5)^{-2}$ of
\eqref{9.98:eq-two-point-law-2} stands. In the tuned frame the pair may sit up to $j_U-2$ levels on
the far side of the stopping depth $K_0$ of the run from $x_0$. This is harmless: both frames give
a two-sided window, depending on $g$ and $\mathsf{d}$ only, for the running coupling at the level
of the pair. That window is all that the consequences in~(b) and~(d) use.
\end{proof}

\begin{lemma}[Comparison of the two frames' depths]\label{9.98:lem-frame-depth-comparison}
Let $g\le\gamma_{\star}$, let $g_0\in\,]0,g_{-}(g)]$, put $K:=K(g_0,g)$, and let
$\mathsf{s}:=K(g_0,\gamma_{\star})-K$ be the prolongation of
Lemma~\ref{9.98:lem-run-prolongation}. Let $C_W$ be the constant of clause \textup{(iv)} of
Theorem~\ref{3.1:thm-main}, written as a function of the running inverse coupling $\bar{x}$ (a fixed
non-decreasing polylogarithm of $\bar{x}$); at depth $t$ in the frame of reference coupling $g$
it is read at $\bar{x}=\bar{x}^{g}_t$ below. For integers $t$ write
\begin{equation}\label{9.98:eq-frame-depth-comparison-1}
\bar{x}^{g}_t:=g^{-2}+B^{\sharp}(t+1),\qquad \bar{x}^{\star}_t:=X_{\star}+B^{\sharp}(t+1).
\end{equation}
For $\mathsf{d}>0$ write $s(\mathsf{d})$ for the least integer $t\ge0$ with
$C_W(\bar{x}^{g}_t)L^{-t}\le\mathsf{d}$, the depth at reference coupling $g$. Let
$s_{\gamma_{\star}}(\cdot)$ be the same depth at reference coupling $\gamma_{\star}$, that is,
with $\bar{x}^{\star}_t$ in place of $\bar{x}^{g}_t$, and let
$u(\mathsf{d};g_0,g):=s_{\gamma_{\star}}(L^{-\mathsf{s}}\mathsf{d})-\mathsf{s}$ be the depth of
Proposition~\ref{9.98:prop-two-point-law}. Assume that there is a pure constant
$c^{\sharp}_W<L$ such that
\begin{equation}\label{9.98:eq-frame-depth-comparison-a}
C_W(\bar{x}')\le c^{\sharp}_W\,C_W(\bar{x})
\quad\text{whenever}\quad \bar{x}\le\bar{x}'\le\bar{x}+2B^{\sharp},\ \ \bar{x}\ge X_{\star},
\end{equation}
let $r_{\mathrm{fr}}$ be the least integer $r\ge0$ with $L^{r}\ge(c^{\sharp}_W)^{r+1}$, and let
$r'_{\mathrm{fr}}$ be the least integer $r\ge0$ with $L^{r}\ge(c^{\sharp}_W)^{j_U-1+r}$.
Then the following hold.
\begin{enumerate}
\item[(a)] For every $\mathsf{d}>0$ and every integer $t\ge0$ with
$B^{\sharp}t\le g^{-2}-X_{\star}+B^{\sharp}$,
\begin{equation}\label{9.98:eq-frame-depth-comparison-7}
s_{\gamma_{\star}}(L^{-t}\mathsf{d})-t\le s(\mathsf{d})+r_{\mathrm{fr}};
\end{equation}
the prolongation $t=\mathsf{s}$ is such an integer, so that
$u(\mathsf{d};g_0,g)\le s(\mathsf{d})+r_{\mathrm{fr}}$ for every $\mathsf{d}>0$. If moreover
$u(\mathsf{d};g_0,g)\ge0$ \textup{(}in particular for
the pairs of Proposition~\ref{9.98:prop-two-point-law}, for which $0<\mathsf{d}\le1/40$ and
$u\ge1$\textup{)}, then $s(\mathsf{d})\le u(\mathsf{d};g_0,g)+r_{\mathrm{fr}}$, and hence
\begin{equation}\label{9.98:eq-frame-depth-comparison-2}
\bigl|u(\mathsf{d};g_0,g)-s(\mathsf{d})\bigr|\le r_{\mathrm{fr}}.
\end{equation}
\item[(b)] For every $\mathsf{d}>0$, the quantities
$\mathsf{U}(g,\mathsf{d})$ and $\bar{x}^{\natural}(\mathsf{d})$ of
\eqref{9.98:eq-prolongation-depths-d} satisfy
\begin{equation}\label{9.98:eq-frame-depth-comparison-3}
\begin{gathered}
\mathsf{U}(g,\mathsf{d})\le s(\mathsf{d})+r_{\mathrm{fr}}+j_U-1+r'_{\mathrm{fr}},\\
\bar{x}^{\natural}(\mathsf{d})\le g^{-2}
+B^{\sharp}\bigl(s(\mathsf{d})+r_{\mathrm{fr}}+j_U+r'_{\mathrm{fr}}\bigr)+c_5 .
\end{gathered}
\end{equation}
At $\mathsf{d}=1$, $s(1)$ is the least $t\ge0$ with $C_W(\bar{x}^{g}_t)\le L^{t}$, and if
$s(1)\ge1$ then $s(1)<1+\log_L C_W(\bar{x}^{g}_{s(1)})$. There is a coupling ceiling, depending
only on $L$, $B^{\sharp}$, $c_5$, $j_U$, $c^{\sharp}_W$ and the function $C_W$, such that for
every $g$ at or below it
\begin{equation}\label{9.98:eq-frame-depth-comparison-9}
B^{\sharp}\bigl(s(1)+r_{\mathrm{fr}}+j_U+r'_{\mathrm{fr}}\bigr)+c_5\le g^{-2},\qquad
\bar{x}^{\natural}(1)\le 2g^{-2};
\end{equation}
and $\limsup_{g\downarrow0}\bar{x}^{\natural}(1)/g^{-2}\le1$.
\end{enumerate}
\end{lemma}

\begin{proof}
The integers $r_{\mathrm{fr}}$ and $r'_{\mathrm{fr}}$ exist: $L^{r}\ge(c^{\sharp}_W)^{r+1}$ is
equivalent to $(L/c^{\sharp}_W)^{r}\ge c^{\sharp}_W$, the inequality
$L^{r}\ge(c^{\sharp}_W)^{j_U-1+r}$ is equivalent to
$(L/c^{\sharp}_W)^{r}\ge(c^{\sharp}_W)^{j_U-1}$, and $L/c^{\sharp}_W>1$.

Write $C^{g}(t):=C_W(\bar{x}^{g}_t)$ and $C^{\star}(t):=C_W(\bar{x}^{\star}_t)$. Since
$g\le\gamma_{\star}$ we have $g^{-2}\ge X_{\star}$, so $\bar{x}^{g}_t\ge X_{\star}$ for
$t\ge-1$; two consecutive arguments differ by $B^{\sharp}\le2B^{\sharp}$, and
\eqref{9.98:eq-frame-depth-comparison-a} gives
\begin{equation}\label{9.98:eq-frame-depth-comparison-4}
C^{g}(t+1)\le c^{\sharp}_W\,C^{g}(t),\qquad\text{hence}\qquad
C^{g}(t+r)\le(c^{\sharp}_W)^{r}C^{g}(t)\quad(r\ge0),
\end{equation}
for every integer $t$ with $\bar{x}^{g}_t\ge X_{\star}$, in particular every $t\ge-1$. The same
holds for $C^{\star}$. Consequently
$C^{g}(t)L^{-t}\le C^{g}(0)(c^{\sharp}_W/L)^{t}\to0$ as $t\to\infty$, and likewise for
$C^{\star}$, so the least integers defining $s(\mathsf{d})$ and $s_{\gamma_{\star}}(\cdot)$
exist.

Put $s:=s(\mathsf{d})$, the least $t\ge0$ with $C^{g}(t)L^{-t}\le\mathsf{d}$. For an integer
$t\ge0$, the integer $s_{\gamma_{\star}}(L^{-t}\mathsf{d})$ is the least $\tau\ge0$ with
$C^{\star}(\tau)L^{-\tau}\le L^{-t}\mathsf{d}$. Writing $\tau=t'+t$ and multiplying by
$L^{t}$, the condition reads $C^{\star}(t'+t)L^{-t'}\le\mathsf{d}$; hence
\begin{equation}\label{9.98:eq-frame-depth-comparison-5}
s_{\gamma_{\star}}(L^{-t}\mathsf{d})-t
=\min\bigl\{t'\ge-t:\ C^{\star}(t'+t)L^{-t'}\le\mathsf{d}\bigr\}.
\end{equation}
For $t=\mathsf{s}$ the left member is $u$ by its definition; here $\mathsf{s}\ge0$, since
$K(g_0,\gamma_{\star})\ge K$ by part (a) of Lemma~\ref{9.98:lem-run-prolongation}.

Let $t\ge0$ be an integer with $B^{\sharp}t\le g^{-2}-X_{\star}+B^{\sharp}$. Adding
$X_{\star}+B^{\sharp}(t'+1)$ to both members gives
$X_{\star}+B^{\sharp}(t'+t+1)\le g^{-2}+B^{\sharp}(t'+2)$, that is,
$\bar{x}^{\star}_{t'+t}\le\bar{x}^{g}_{t'+1}$. For every integer $t'\ge0$, the monotonicity of
$C_W$ and \eqref{9.98:eq-frame-depth-comparison-4} at $t'$ give
\begin{equation}\label{9.98:eq-frame-depth-comparison-8}
C^{\star}(t'+t)\le C^{g}(t'+1)\le c^{\sharp}_W\,C^{g}(t').
\end{equation}

By the proof of part (a) of Lemma~\ref{9.98:lem-run-prolongation},
\begin{equation*}
g^{-2}-X_{\star}-B^{\sharp}<B^{\sharp}\mathsf{s}<g^{-2}-X_{\star}+B^{\sharp}.
\end{equation*}
In particular $t=\mathsf{s}$ is an integer as in the preceding paragraph. Adding
$X_{\star}+B^{\sharp}(t'+1)$ to each member gives
\begin{equation*}
g^{-2}+B^{\sharp}t'<X_{\star}+B^{\sharp}(t'+\mathsf{s}+1)<g^{-2}+B^{\sharp}(t'+2),
\end{equation*}
that is, $\bar{x}^{g}_{t'-1}<\bar{x}^{\star}_{t'+\mathsf{s}}<\bar{x}^{g}_{t'+1}$. Let $t'\ge0$ be
an integer, so that $t'-1\ge-1$. Since $C_W$ is non-decreasing, and by
\eqref{9.98:eq-frame-depth-comparison-4} applied at $t'-1$ and at $t'$,
\begin{equation}\label{9.98:eq-frame-depth-comparison-6}
(c^{\sharp}_W)^{-1}C^{g}(t')\le C^{g}(t'-1)\le C^{\star}(t'+\mathsf{s})\le C^{g}(t'+1)
\le c^{\sharp}_W\,C^{g}(t').
\end{equation}

Upper bound. Let $t\ge0$ be an integer with $B^{\sharp}t\le g^{-2}-X_{\star}+B^{\sharp}$, let
$r:=r_{\mathrm{fr}}$ and $t':=s+r\ge0\ge-t$. By
\eqref{9.98:eq-frame-depth-comparison-8}, then \eqref{9.98:eq-frame-depth-comparison-4}, then the
definition of $s$, and finally the definition of $r_{\mathrm{fr}}$,
\begin{align*}
C^{\star}(t'+t)L^{-t'}
&\le c^{\sharp}_W\,C^{g}(s+r)L^{-s-r}
\le(c^{\sharp}_W)^{r+1}L^{-r}\cdot C^{g}(s)L^{-s}\\
&\le(c^{\sharp}_W)^{r+1}L^{-r}\,\mathsf{d}\le\mathsf{d}.
\end{align*}
By \eqref{9.98:eq-frame-depth-comparison-5},
$s_{\gamma_{\star}}(L^{-t}\mathsf{d})-t\le s+r_{\mathrm{fr}}$, which is
\eqref{9.98:eq-frame-depth-comparison-7}. At $t=\mathsf{s}$, which is admissible by the interval
above, this gives $u\le s+r_{\mathrm{fr}}$ for every $\mathsf{d}>0$.

Lower bound. Let again $r:=r_{\mathrm{fr}}$, and assume $u\ge0$. Then
\eqref{9.98:eq-frame-depth-comparison-6} applies at $t'=u$, and $u+r\ge0$ is an admissible value
of $t$ in the definition of $s$. By \eqref{9.98:eq-frame-depth-comparison-5} at
$t=\mathsf{s}$ and $t'=u$, the left half of \eqref{9.98:eq-frame-depth-comparison-6} at $t'=u$,
and \eqref{9.98:eq-frame-depth-comparison-4} in the form
$C^{g}(u)\ge(c^{\sharp}_W)^{-r}C^{g}(u+r)$,
\begin{align*}
\mathsf{d}&\ge C^{\star}(u+\mathsf{s})L^{-u}\ge(c^{\sharp}_W)^{-1}C^{g}(u)L^{-u}
\ge(c^{\sharp}_W)^{-(r+1)}C^{g}(u+r)L^{-u}\\
&=(c^{\sharp}_W)^{-(r+1)}L^{r}\cdot C^{g}(u+r)L^{-(u+r)}.
\end{align*}
Hence
$C^{g}(u+r)L^{-(u+r)}\le(c^{\sharp}_W)^{r+1}L^{-r}\mathsf{d}\le\mathsf{d}$ by the definition of
$r_{\mathrm{fr}}$, and by the definition of $s$ as a least integer, $s\le u+r_{\mathrm{fr}}$. When
$u\ge0$ the two bounds give \eqref{9.98:eq-frame-depth-comparison-2}, which completes part (a).

Part (b). By \eqref{9.98:eq-frame-depth-comparison-7}, for every $\mathsf{d}>0$ and every
integer $t\ge0$ with $B^{\sharp}t\le g^{-2}-X_{\star}+B^{\sharp}$ one has
$s_{\gamma_{\star}}(L^{-t}\mathsf{d})-t\le s(\mathsf{d})+r_{\mathrm{fr}}$, with no condition on
the sign of the left member; by the proof of part (a) of
Lemma~\ref{9.98:lem-run-prolongation}, applied to each bare coupling in $]0,g_{-}(g)]$, these
integers $t$ include every prolongation of that lemma at reference coupling $g$.
In \eqref{9.98:eq-prolongation-depths-d} the depth entering $\mathsf{U}(g,\mathsf{d})$
is read at the arguments $L^{-t-j}\mathsf{d}$ displayed there, with $t$ such a prolongation and
$|j|\le j_U-1$, so that $j_U-1\ge0$. We compare the least depth at $L^{-t-j}\mathsf{d}$ with
the least depth at $L^{-t}\mathsf{d}$. Put $\mathsf{d}':=L^{-t}\mathsf{d}$,
$\tau_0:=s_{\gamma_{\star}}(\mathsf{d}')\ge0$ and $r':=r'_{\mathrm{fr}}$, and let $j\le j_U-1$
be an integer. By \eqref{9.98:eq-frame-depth-comparison-4} for $C^{\star}$ at $\tau_0$ with
$r=j_U-1+r'\ge0$, then the definition of $\tau_0$, then the definition of $r'_{\mathrm{fr}}$,
and finally $j\le j_U-1$ and $L>1$,
\begin{align*}
C^{\star}(\tau_0+j_U-1+r')L^{-(\tau_0+j_U-1+r')}
&\le(c^{\sharp}_W)^{j_U-1+r'}L^{-r'}\cdot L^{-(j_U-1)}\,C^{\star}(\tau_0)L^{-\tau_0}\\
&\le(c^{\sharp}_W)^{j_U-1+r'}L^{-r'}\cdot L^{-(j_U-1)}\mathsf{d}'\\
&\le L^{-(j_U-1)}\mathsf{d}'\le L^{-j}\mathsf{d}'.
\end{align*}
Hence $s_{\gamma_{\star}}(L^{-t-j}\mathsf{d})\le s_{\gamma_{\star}}(L^{-t}\mathsf{d})
+j_U-1+r'_{\mathrm{fr}}$: the change of the argument raises the least depth by at most
$j_U-1+r'_{\mathrm{fr}}$. Together with part (a), in the form
\eqref{9.98:eq-frame-depth-comparison-7}, this
gives the first inequality
of \eqref{9.98:eq-frame-depth-comparison-3}, and by the definition of
$\bar{x}^{\natural}(\mathsf{d})$ in \eqref{9.98:eq-prolongation-depths-d} the second follows.

At $\mathsf{d}=1$ the condition $C^{g}(t)L^{-t}\le1$ reads $C_W(\bar{x}^{g}_t)\le L^{t}$, so
$s(1)$ is the least $t\ge0$ with this property. If $s(1)\ge1$, minimality at $t=s(1)-1$ and the
monotonicity of $C_W$ give
\begin{equation*}
L^{s(1)-1}<C_W(\bar{x}^{g}_{s(1)-1})\le C_W(\bar{x}^{g}_{s(1)}),
\end{equation*}
that is, $s(1)<1+\log_L C_W(\bar{x}^{g}_{s(1)})$. Moreover, by the bound
$C^{g}(t)L^{-t}\le C^{g}(0)(c^{\sharp}_W/L)^{t}$ of the second paragraph, valid for $t\ge0$, and
since $C^{g}(0)=C_W(g^{-2}+B^{\sharp})$ and $L/c^{\sharp}_W>1$, every integer $t\ge0$ with
$t\log(L/c^{\sharp}_W)\ge\log C_W(g^{-2}+B^{\sharp})$ satisfies $C^{g}(t)L^{-t}\le1$; hence
\begin{equation*}
s(1)\le1+\max\Bigl\{0,\ \frac{\log C_W(g^{-2}+B^{\sharp})}{\log(L/c^{\sharp}_W)}\Bigr\}.
\end{equation*}
As $C_W$ is a polylogarithm, $s(1)$ is thus bounded by a quantity logarithmic in a
polylogarithm of $g^{-2}$, while $r_{\mathrm{fr}}$, $r'_{\mathrm{fr}}$, $j_U$, $B^{\sharp}$ and
$c_5$ do not depend on $g$; therefore
$B^{\sharp}(s(1)+r_{\mathrm{fr}}+j_U+r'_{\mathrm{fr}})+c_5$ is $o(g^{-2})$ as
$g\downarrow0$, at a rate determined by $L$, $B^{\sharp}$, $c_5$, $j_U$, $c^{\sharp}_W$ and the
function $C_W$ alone. This gives the coupling ceiling of part (b): for every $g$ at or below it,
the first inequality of \eqref{9.98:eq-frame-depth-comparison-9} holds, and the second
inequality of \eqref{9.98:eq-frame-depth-comparison-3} at $\mathsf{d}=1$ then gives
$\bar{x}^{\natural}(1)\le2g^{-2}$, the second inequality of
\eqref{9.98:eq-frame-depth-comparison-9}. Finally, for every $g$ the second
inequality of \eqref{9.98:eq-frame-depth-comparison-3} at $\mathsf{d}=1$ gives
\begin{equation*}
\frac{\bar{x}^{\natural}(1)}{g^{-2}}
\le1+\frac{B^{\sharp}(s(1)+r_{\mathrm{fr}}+j_U+r'_{\mathrm{fr}})+c_5}{g^{-2}},
\end{equation*}
and the last quotient tends to $0$ as $g\downarrow0$; hence
$\limsup_{g\downarrow0}\bar{x}^{\natural}(1)/g^{-2}\le1$.
\end{proof}

\begin{remark}[The growth hypothesis on $C_W$ and the role of the comparison]
The inequality
\eqref{9.98:eq-frame-depth-comparison-a} is a hypothesis of
Lemma~\ref{9.98:lem-frame-depth-comparison}. Whether clause
\textup{(iv)}'s $C_W$ satisfies it with $c^{\sharp}_W<L$ at the given $X_{\star}$ is not
asserted, and the inequality $c^{\sharp}_W<L$ is not a condition on $L$ adopted by any statement
of this book. If $C_W(\bar{x})=P(\log\bar{x})$ with $P$ a fixed non-decreasing polynomial of
degree $\nu$ with positive leading coefficient, then for $\bar{x}\ge2B^{\sharp}$ one has
$\bar{x}'\le\bar{x}+2B^{\sharp}\le2\bar{x}$ in \eqref{9.98:eq-frame-depth-comparison-a}, and
\begin{equation*}
\frac{C_W(\bar{x}')}{C_W(\bar{x})}\le\Bigl(1+\frac{\log2}{\log\bar{x}}\Bigr)^{\nu}\bigl(1+o(1)\bigr)
\longrightarrow1\qquad(\bar{x}\to\infty),
\end{equation*}
so that a constant $c^{\sharp}_W$ in \eqref{9.98:eq-frame-depth-comparison-a} can be taken as
close to $1$ as $X_{\star}$ is large.

The displays of Proposition~\ref{9.98:prop-two-point-law} hold with $u$ as defined there and
with the $g$-only quantity $\mathsf{U}(g,\mathsf{d})$; they do not use
Lemma~\ref{9.98:lem-frame-depth-comparison}. Part (a) of that lemma
identifies the running coupling displayed there as
the one at the depth $s(\mathsf{d})$ of the pair in the frame of reference coupling $g$, up to
$r_{\mathrm{fr}}$ levels, and part (b) gives the size of $\bar{x}^{\natural}(\mathsf{d})$ in terms
of $\bar{x}^{g}_{s(\mathsf{d})}$.
\end{remark}

\begin{proof}[Proof of Theorem~\ref{9.98:thm-short-distance-uniqueness}\,(d): the relative two-point bound]
\label{9.98:prf-proof-relative-two-point}
Let $g$, $\hat{x}$ and $\vartheta$ be as in Theorem~\ref{9.98:thm-short-distance-uniqueness}. We consider the set
\begin{equation*}
\operatorname{Lim}_{\infty}[\hat{x}]\cup\{\mathfrak{S}_m[\hat{x}] : m\ge m_{\natural}(g)\}.
\end{equation*}
It consists of the infinite-volume limit points in the tuned frame at the renormalised coupling
$\hat{x}$, together with the continuum limits on the tori $\mathbb{T}_m$ at $\hat{x}$ for
$m\ge m_{\natural}(g)$. For $\sigma$ in this set, $S_2[\sigma](f_1,f_2)$ denotes the connected
two-point function of $\sigma$ evaluated at the pair $(f_1,f_2)$.

Let the separation $\mathsf{d}$ satisfy
\begin{equation*}
0<\mathsf{d}\le\min\bigl\{L^{\underline{s}'(g)}/40,\ 1/40\bigr\},
\end{equation*}
and let $(f_1,f_2)\in\mathfrak{T}(\mathsf{d};\vartheta)$, the class of pairs of clause~(iv) of
Theorem~\ref{3.1:thm-main}. By the inclusion
$\mathfrak{T}(\mathsf{d};\vartheta)\subset\mathfrak{P}_2(\mathsf{d};\vartheta)$ of the pairs of
clause~(iv) in the admissible pairs of scale $\mathsf{d}$, recorded with the definition of the
admissible class $\mathfrak{P}_2(\mathsf{d};\vartheta)$, the pair $(f_1,f_2)$ is
admissible for part~(b) of Theorem~\ref{9.98:thm-short-distance-uniqueness}.

Fix $\sigma,\sigma'$ in the set above.

\emph{Numerator.} Part~(b) of Theorem~\ref{9.98:thm-short-distance-uniqueness} applies at $n=2$,
with $\tau_2$ given by \eqref{9.98:eq-short-distance-uniqueness-a}. It gives
\begin{equation}\label{9.98:eq-proof-relative-two-point-1}
\bigl|S_2[\sigma](f_1,f_2)-S_2[\sigma'](f_1,f_2)\bigr|\le 2\tau_2
=2C_T^{\star}L^{-2\underline{s}'(g)}\mathsf{d}^2\,\|f_1\|_\infty\|f_2\|_\infty .
\end{equation}

\emph{Denominator.} We use the tuned-frame limit sandwich \eqref{9.98:eq-two-point-law-a} of
Proposition~\ref{9.98:prop-two-point-law}, in its form written with the tuned-frame prolongation
$\mathsf{s}'$ of that proposition. In this form it requires of
$\mathsf{d}$ nothing beyond the range above; the form written with the prolongation $\mathsf{s}$
would require in addition $\mathsf{d}\le L^{1-j_U}/40$. It applies to $\sigma'$ at the pair
$(f_1,f_2)$ and gives
\begin{equation}\label{9.98:eq-proof-relative-two-point-2}
S_2[\sigma'](f_1,f_2)\ge\tfrac{2}{9}\,b_0c_{\mathcal{K}}\,
\bar{x}^{\natural}(\mathsf{d})^{-2}\,\|f_1\|_\infty\|f_2\|_\infty>0 .
\end{equation}
Here $\bar{x}^{\natural}(\mathsf{d})$ is the upper end of the window
\eqref{9.98:eq-prolongation-depths-d} of Definition~\ref{9.98:def-prolongation-depths}.

\emph{Ratio.} By \eqref{9.98:eq-proof-relative-two-point-2} the quotient
$S_2[\sigma]/S_2[\sigma']$ at $(f_1,f_2)$ is defined. Dividing
\eqref{9.98:eq-proof-relative-two-point-1} by \eqref{9.98:eq-proof-relative-two-point-2}, the
product $\|f_1\|_\infty\|f_2\|_\infty$ cancels and $2/(2/9)=9$. We obtain
\begin{equation}\label{9.98:eq-proof-relative-two-point-3}
\begin{split}
\Bigl|\frac{S_2[\sigma](f_1,f_2)}{S_2[\sigma'](f_1,f_2)}-1\Bigr|
&=\frac{\bigl|S_2[\sigma](f_1,f_2)-S_2[\sigma'](f_1,f_2)\bigr|}{S_2[\sigma'](f_1,f_2)}
\le\frac{2\tau_2}{S_2[\sigma'](f_1,f_2)}\\
&\le\frac{9C_T^{\star}L^{-2\underline{s}'(g)}\,\mathsf{d}^2\,\bar{x}^{\natural}(\mathsf{d})^2}
{b_0c_{\mathcal{K}}}.
\end{split}
\end{equation}
As $\mathsf{d}\downarrow0$, the right-hand side tends to zero like $\mathsf{d}^2$ times the
factor $\bar{x}^{\natural}(\mathsf{d})^2$, which by the form of the window
\eqref{9.98:eq-prolongation-depths-d} grows only polylogarithmically in $1/\mathsf{d}$.

\emph{In reference units.} By \eqref{9.98:eq-prolongation-depths-a}, the tuned-frame depth is
the shift $\underline{s}'(g)=\underline{s}(g)-(j_U-1)$, so
$L^{-2\underline{s}'(g)}=L^{2(j_U-1)}L^{-2\underline{s}(g)}$. Display
\eqref{9.98:eq-prolongation-depths-a} also gives the lower bound
\begin{equation*}
\underline{s}(g)\ge\frac{g^{-2}-X_{\star}-2B^{\sharp}}{B^{\sharp}} .
\end{equation*}
Since $L>1$, this yields
\begin{equation*}
L^{-2\underline{s}(g)}\le L^{2(X_{\star}+2B^{\sharp})/B^{\sharp}}\,L^{-2g^{-2}/B^{\sharp}} .
\end{equation*}
For every real $y$ we have $L^{-2y/B^{\sharp}}=\exp\bigl(-(2\ln L/B^{\sharp})\,y\bigr)$; we apply
this with $y=g^{-2}$. Altogether
\begin{equation}\label{9.98:eq-proof-relative-two-point-4}
L^{-2\underline{s}'(g)}\le L^{2(j_U-1)}\,L^{2(X_{\star}+2B^{\sharp})/B^{\sharp}}\,
\exp\Bigl[-\frac{2\ln L}{B^{\sharp}}\,g^{-2}\Bigr].
\end{equation}
We insert \eqref{9.98:eq-proof-relative-two-point-4} into
\eqref{9.98:eq-proof-relative-two-point-3}, for every $\mathsf{d}$ in the range above. This range
is $0<\mathsf{d}\le1/40$. Indeed, the two facts just used give
\begin{equation*}
\underline{s}'(g)=\underline{s}(g)-(j_U-1)\ge\frac{g^{-2}-X_{\star}}{B^{\sharp}}-(j_U+1),
\end{equation*}
and the hypothesis $g\le\gamma_{\natural}$, where
$\gamma_{\natural}^{-2}=X_{\star}+B^{\sharp}(s_{\star}+j_U+1)$ by
Definition~\ref{9.98:def-prolongation-depths}, gives
$(g^{-2}-X_{\star})/B^{\sharp}\ge s_{\star}+j_U+1$; hence $\underline{s}'(g)\ge s_{\star}\ge0$,
the depth $s_{\star}$ of Definition~\ref{9.98:def-prolongation-depths} being a maximum of depth
floors, which are depths and hence non-negative, and so $L^{\underline{s}'(g)}/40\ge1/40$. The
factors independent of $g$
and of $\mathsf{d}$ are $9C_T^{\star}$, $(b_0c_{\mathcal{K}})^{-1}$, $L^{2(j_U-1)}$ and
$L^{2(X_{\star}+2B^{\sharp})/B^{\sharp}}$. Their product is bounded by the pure constant
$\mathsf{A}_0$ of \eqref{9.98:eq-short-distance-uniqueness-b}. The relative discrepancy is
therefore at most
\begin{equation}\label{9.98:eq-proof-relative-two-point-5}
\mathsf{A}_0\,\mathsf{d}^2\,\bar{x}^{\natural}(\mathsf{d})^2
\exp\Bigl[-\frac{2\ln L}{B^{\sharp}}\,g^{-2}\Bigr].
\end{equation}
This is the reference-scale form of the relative bound, read at the separations
$0<\mathsf{d}\le1/40$, that is, within a fixed fraction of the reference length. Let $g$ lie,
in addition, below the coupling ceiling of Lemma~\ref{9.98:lem-frame-depth-comparison}. Under
that ceiling the $\mathsf{d}$-dependent factor $\mathsf{d}^2\,\bar{x}^{\natural}(\mathsf{d})^2$ of
\eqref{9.98:eq-proof-relative-two-point-5} is at most $4g^{-4}$ at every separation of this
range. Putting $\mathsf{A}:=4\mathsf{A}_0$, a pure constant, the relative discrepancy is then,
for every $0<\mathsf{d}\le1/40$, at most
\begin{equation*}
\mathsf{A}\,g^{-4}\exp\Bigl[-\frac{2\ln L}{B^{\sharp}}\,g^{-2}\Bigr].
\end{equation*}
Inserting \eqref{9.98:eq-proof-relative-two-point-4} into
\eqref{9.98:eq-proof-relative-two-point-1} instead, the discrepancy itself satisfies
\begin{equation*}
\begin{split}
&\bigl|S_2[\sigma](f_1,f_2)-S_2[\sigma'](f_1,f_2)\bigr|\\
&\qquad\le 2C_T^{\star}L^{2(j_U-1)}L^{2(X_{\star}+2B^{\sharp})/B^{\sharp}}\,
\mathsf{d}^2\,\|f_1\|_\infty\|f_2\|_\infty\\
&\qquad\quad\times\exp\Bigl[-\frac{2\ln L}{B^{\sharp}}\,g^{-2}\Bigr].
\end{split}
\end{equation*}
In reference units this discrepancy is exponentially small in the inverse coupling squared,
and hence smaller than every power of the running coupling. Indeed, along any run with first
passage at $g$, the running inverse coupling $x_{K-s}$ at a fixed depth $s$ below the cutoff is
comparable with $g^{-2}$, so every fixed power $g_{K-s}^{2p}=x_{K-s}^{-p}$ of the running coupling
decays only polynomially in $g^{-2}$, whereas the prefactor of the exponential above does not
depend on $g$.

\emph{The perturbative reading.} Suppose that one member $\sigma'$ of the set admits an asymptotic
expansion on a family of pairs $(f_1^{\mathsf{d}},f_2^{\mathsf{d}})\in\mathfrak{T}(\mathsf{d};\vartheta)$,
where $\mathsf{d}$ ranges over the values above and the family is normalised by
$\|f_1^{\mathsf{d}}\|_\infty\|f_2^{\mathsf{d}}\|_\infty=1$; the leading size below is
measured in this normalisation. The expansion reads
\begin{equation}\label{9.98:eq-proof-relative-two-point-6}
S_2[\sigma'](f_1^{\mathsf{d}},f_2^{\mathsf{d}})
=\mathfrak{M}(\mathsf{d})\Bigl(1+\sum_{k\le p}c_k\bigl(\ln\tfrac{1}{\mathsf{d}}\bigr)^{-k}
+o\bigl((\ln\tfrac{1}{\mathsf{d}})^{-p}\bigr)\Bigr),
\end{equation}
with a leading size $\mathfrak{M}(\mathsf{d})$ comparable with $(\ln\frac{1}{\mathsf{d}})^{-2}$.
Here $k$ is a summation index and $p\ge0$ an integer; neither denotes a level or a plaquette.
Then every $\sigma$ in the set has the same expansion
\eqref{9.98:eq-proof-relative-two-point-6} to order $p$. To see this, note that by
\eqref{9.98:eq-proof-relative-two-point-1}, and since
$1/\mathfrak{M}(\mathsf{d})=O((\ln\frac{1}{\mathsf{d}})^2)$,
\begin{equation*}
\begin{split}
\bigl|S_2[\sigma](f_1^{\mathsf{d}},f_2^{\mathsf{d}})-S_2[\sigma'](f_1^{\mathsf{d}},f_2^{\mathsf{d}})\bigr|
&\le2C_T^{\star}L^{-2\underline{s}'(g)}\mathsf{d}^2\\
&=\mathfrak{M}(\mathsf{d})\,O\bigl(\mathsf{d}^2(\ln\tfrac{1}{\mathsf{d}})^2\bigr).
\end{split}
\end{equation*}
Because $\mathsf{d}^2(\ln\frac{1}{\mathsf{d}})^{2+p}\to0$ as $\mathsf{d}\downarrow0$, the right-hand
side is $\mathfrak{M}(\mathsf{d})\,o((\ln\frac{1}{\mathsf{d}})^{-p})$, and this holds for every $p$.
Adding it to \eqref{9.98:eq-proof-relative-two-point-6} leaves the coefficients $c_k$, $k\le p$,
unchanged. Thus, if one member has an asymptotic expansion to some order, all have the same.

Clause~(iv) of Theorem~\ref{3.1:thm-main}, and with it
Proposition~\ref{9.98:prop-two-point-law}, determine the two-point function only up to the factor
$1+R$, with the bound $|R|\le\frac12$ on the remainder. The existence of an asymptotic expansion
\eqref{9.98:eq-proof-relative-two-point-6} is not asserted.
\end{proof}

\begin{theorem}[Weak-coupling limit sets as exact dilates]\label{9.98:thm-exact-dilates}
Let $\gamma_{\star}$ be the strong reference coupling of
Definition~\ref{9.98:def-strong-reference-coupling}, let $g\le\gamma_{\star}$, and let
$\underline{s}(g)$, $\overline{s}(g)$ be the prolongation depths of
\eqref{9.98:eq-prolongation-depths-a}. For an integer $\mathsf{s}$, a tuple
$\vec f=(f_1,\dots,f_n)$ and a family $\sigma'$ of functionals on tuples, write
$\vec f(L^{\mathsf{s}}\cdot):=(f_1(L^{\mathsf{s}}\cdot),\dots,f_n(L^{\mathsf{s}}\cdot))$ and
define the exact dilation $D_{\mathsf{s}}$ by
$(D_{\mathsf{s}}\sigma')(\vec f\,):=\sigma'(\vec f(L^{\mathsf{s}}\cdot))$. Then, with the
limit sets of Notation~\ref{9.98:not-couplings-limit-sets},
\begin{equation}\label{9.98:eq-exact-dilates-1}
\operatorname{Lim}_{\infty}(g)\subset
\bigcup_{\mathsf{s}=\underline{s}(g)}^{\overline{s}(g)}
D_{\mathsf{s}}\operatorname{Lim}_{\infty}(\gamma_{\star}).
\end{equation}
Precisely: let $\sigma\in\operatorname{Lim}_{\infty}(g)$ be the limit
\begin{equation*}
\sigma=\lim_{j\to\infty}S^T_{n;K_j,m_j}[g_0^{(j)}]
\end{equation*}
along a joint diagonal, that is, $g_0^{(j)}\downarrow0$, $K_j=K(g_0^{(j)},g)$ and
$m_j\to\infty$, the limit holding for every $n\ge2$ and every separated real tuple of
$C^1_c$ functions. Then, along a further subsequence, the integer
$\mathsf{s}_j:=K(g_0^{(j)},\gamma_{\star})-K_j$ is constant, equal to some
$\mathsf{s}\in[\underline{s}(g),\overline{s}(g)]$; along it the limit
\begin{equation*}
\sigma':=\lim_{j\to\infty}S^T_{n;K_j+\mathsf{s},\,m_j-\mathsf{s}}[g_0^{(j)}]
\end{equation*}
exists for every $n\ge2$ and every such tuple, lies in
$\operatorname{Lim}_{\infty}(\gamma_{\star})$, and $\sigma=D_{\mathsf{s}}\sigma'$, that is,
\begin{equation}\label{9.98:eq-exact-dilates-2}
\sigma(\vec f\,)=\sigma'(\vec f(L^{\mathsf{s}}\cdot))\qquad\text{for every tuple }\vec f.
\end{equation}
Thus every limit point on $\mathbb{R}^4$ at reference coupling $g$ is an exact
$L^{\mathsf{s}}$-dilate of a limit point at reference coupling $\gamma_{\star}$, obtained
along the same sequence of bare couplings.
\end{theorem}

\begin{proof}
Since $g_0^{(j)}\downarrow0$ and $g_-(g)=(g^{-2}+B^{\sharp})^{-1/2}>0$, we have
$g_0^{(j)}\in\,]0,g_-(g)]$ for all $j$ large enough; we discard the finitely many other
indices, which changes neither the limit $\sigma$ nor the properties
$g_0^{(j)}\downarrow0$, $m_j\to\infty$.

\emph{The prolongation takes finitely many values.} For each $j$ we apply the run
prolongation to the strong reference coupling,
Lemma~\ref{9.98:lem-run-prolongation}, with $g\le\gamma_{\star}$ and
$g_0=g_0^{(j)}\in\,]0,g_-(g)]$, so that its $K$ is $K_j=K(g_0^{(j)},g)$. By parts (a)
and (c) of that lemma, $K(g_0^{(j)},\gamma_{\star})$ is well defined,
$g_0^{(j)}\le g_-(\gamma_{\star})$, and the integer
$\mathsf{s}_j=K(g_0^{(j)},\gamma_{\star})-K_j$ satisfies
\eqref{9.98:eq-run-prolongation-a}, that is,
$\underline{s}(g)\le\mathsf{s}_j\le\overline{s}(g)$. The bounds $\underline{s}(g)$ and
$\overline{s}(g)$ depend on $g$ alone and not on $j$, so $\mathsf{s}_j$ ranges over the
finitely many integers of the interval $[\underline{s}(g),\overline{s}(g)]$. Hence some
integer $\mathsf{s}$ of this interval equals $\mathsf{s}_j$ for infinitely many $j$, and we
pass to the subsequence of those indices. Along it
$\mathsf{s}_j\equiv\mathsf{s}$, still $g_0^{(j)}\downarrow0$ and $m_j\to\infty$, and, being a
subsequence of a convergent sequence, $S^T_{n;K_j,m_j}[g_0^{(j)}](\vec f\,)$ still converges
to $\sigma(\vec f\,)$ for every $n\ge2$ and every tuple $\vec f$.

\emph{One lattice integral under two names.} Since $m_j\to\infty$, we have
$m_j\ge\mathsf{s}+m_{\star}$ for all $j$ large, with $m_{\star}$ the torus exponent of
Definition~\ref{9.98:def-prolongation-depths}; again we keep only these indices. For each
such $j$, part (b) of Lemma~\ref{9.98:lem-run-prolongation}, that is, the identity
\eqref{9.98:eq-run-prolongation-b} applied with $g_0=g_0^{(j)}$, $K=K_j$,
$K'=K(g_0^{(j)},\gamma_{\star})=K_j+\mathsf{s}$ and $m=m_j\ge\mathsf{s}+m_{\star}$, gives
\begin{equation}\label{9.98:eq-exact-dilates-3}
S^T_{n;K_j,m_j}[g_0^{(j)}](\vec f\,)
=S^T_{n;K_j+\mathsf{s},\,m_j-\mathsf{s}}[g_0^{(j)}](\vec f(L^{\mathsf{s}}\cdot))
\end{equation}
for every $n\ge2$ and every tuple $\vec f$.

\emph{The dilated tuples exhaust the test class.} The map
$\vec f\mapsto\vec f(L^{\mathsf{s}}\cdot)$ is a bijection of the test class of
Notation~\ref{9.98:not-couplings-limit-sets}, the separated real tuples of $C^1_c$
functions, onto itself. Indeed, each
$f_i(L^{\mathsf{s}}\cdot)$ is real and of class $C^1$, its support is
$L^{-\mathsf{s}}\operatorname{supp}f_i$, which is compact, and two such supports are
disjoint (and at distance $L^{-\mathsf{s}}$ times the original distance, hence still
separated) exactly when the original supports are; the inverse map is
$\vec h\mapsto\vec h(L^{-\mathsf{s}}\cdot)$, which has the same properties.

\emph{Definition of $\sigma'$.} Let $\vec h$ be any tuple of the test class, and put
$\vec f:=\vec h(L^{-\mathsf{s}}\cdot)$, so that $\vec f$ belongs to the test class and
$\vec h=\vec f(L^{\mathsf{s}}\cdot)$. By \eqref{9.98:eq-exact-dilates-3}, the sequence
$S^T_{n;K_j+\mathsf{s},\,m_j-\mathsf{s}}[g_0^{(j)}](\vec h\,)$ coincides term by term with
$S^T_{n;K_j,m_j}[g_0^{(j)}](\vec f\,)$, which converges to $\sigma(\vec f\,)$ by the
definition of $\sigma$. Hence the limit
\begin{equation*}
\sigma'(\vec h\,):=\lim_{j\to\infty}S^T_{n;K_j+\mathsf{s},\,m_j-\mathsf{s}}[g_0^{(j)}](\vec h\,)
\end{equation*}
exists for every $n\ge2$ and every tuple $\vec h$ of the test class, and
$\sigma'(\vec f(L^{\mathsf{s}}\cdot))=\sigma(\vec f\,)$ for every $\vec f$; this is
\eqref{9.98:eq-exact-dilates-2}, i.e. $\sigma=D_{\mathsf{s}}\sigma'$.

\emph{Membership in $\operatorname{Lim}_{\infty}(\gamma_{\star})$.} The sequence
$(g_0^{(j)},K_j+\mathsf{s},m_j-\mathsf{s})$ has $g_0^{(j)}\downarrow0$,
$K_j+\mathsf{s}=K(g_0^{(j)},\gamma_{\star})$ and $m_j-\mathsf{s}\to\infty$; it is therefore a
first-passage joint diagonal at reference coupling $\gamma_{\star}$. Its bare couplings
satisfy $g_0^{(j)}\le g_-(\gamma_{\star})$ by part (a) of
Lemma~\ref{9.98:lem-run-prolongation}, and its tori belong to the family at reference
$\gamma_{\star}$, since $m_j-\mathsf{s}\ge m_{\star}$. The limit $\sigma'$ along it exists
for every $n\ge2$ and every tuple of the test class. By the definition of the limit sets in
Notation~\ref{9.98:not-couplings-limit-sets}, whose objects are those of clause
(ii-$\mathbb{R}^4$) of Theorem~\ref{3.1:thm-main} read at reference coupling
$\gamma_{\star}$, $\sigma'$ is therefore an element of
$\operatorname{Lim}_{\infty}(\gamma_{\star})$. That
clause may be read at reference $\gamma_{\star}$ because
$\gamma_{\star}\le\gamma_J\le\gamma_H$ by the choice of $\gamma_{\star}$ in
Definition~\ref{9.98:def-strong-reference-coupling}, so that hypothesis (H7) of
Notation~\ref{2.2:not-hypotheses} holds at that reference coupling.

Since $\sigma\in\operatorname{Lim}_{\infty}(g)$ was arbitrary and
$\mathsf{s}\in[\underline{s}(g),\overline{s}(g)]$, this proves
\eqref{9.98:eq-exact-dilates-1}.
\end{proof}

We add four observations on Theorem~\ref{9.98:thm-exact-dilates}.

(i) The proof uses Lemma~\ref{9.98:lem-run-prolongation} and the definition of the limit
sets, and nothing else: no estimate on the never-healed old-age tail, no rarity of
large-field regions, no age or labelling of such regions, and no clustering property
enter. The theorem is the identity of one
lattice integral under two namings of the reference length, together with the
first-passage bookkeeping of clause (0) of Theorem~\ref{3.1:thm-main}.

(ii) Clause (v-d)(iii) of Theorem~\ref{3.1:thm-main} asserts, on one torus, no relation
between the continuum limits at distinct canonical comparison points $\hat{x}$, neither a
continuous dilation nor a dimensional transmutation. This stays so
torus by torus: the change-of-units identity sends a torus exponent $m_0$ and a
comparison point $\hat{x}$ to the torus exponent $m_0+\mathsf{s}$ and a comparison point
depending on $\hat{x}$ and $\mathsf{s}$, so that it changes the torus.
Theorem~\ref{9.98:thm-exact-dilates} is a statement on $\mathbb{R}^4$, where the torus
exponent shifted by the change-of-units identity has tended to infinity on both sides.

(iii) Taken alone, the inclusion \eqref{9.98:eq-exact-dilates-1} says nothing about
whether any limit point is non-trivial in the sense of
Definition~\ref{3.1:def-non-trivial} (hierarchical models satisfy the same inclusion),
and it asserts no uniqueness. What it provides is a frame: the dependence on the coupling
of the limit sets on $\mathbb{R}^4$, at all weak couplings, is carried entirely by exact
dilations and by the dependence inside one compact window, namely the closure
$\overline{W}_{\star}$ of the first-passage window of
Definition~\ref{9.98:def-prolongation-depths} at the one reference coupling
$\gamma_{\star}$. To name the members of this window by renormalised couplings one applies,
at reference coupling $\gamma_{\star}$, the corollary of the uniqueness theorem, clause
(v-c)(2) of Theorem~\ref{3.1:thm-main}, which places the offset
$a_{K'}(\gamma_{\star})-x$ of every first-passage datum $x=x_{K'}$, with
$K'=K(g_0,\gamma_{\star})$, in $[-\varrho_2/2,\varrho_2/2]$; that application requires
$\gamma_{\star}\le\gamma_U'$, where $\gamma_U'\le\gamma_U$ is the coupling threshold under
which that clause is stated. Neither that clause nor this condition is used
in Theorem~\ref{9.98:thm-exact-dilates}.

(iv) Read in the direction of weaker coupling, the change-of-units identity places the
$L^{\mathsf{s}}$-dilates of the limit set at reference coupling $g$ inside a limit set at a
weaker reference coupling. Theorem~\ref{9.98:thm-exact-dilates} is the opposite,
coarsening direction: it carries the limit sets at every weak coupling $g\le\gamma_{\star}$
back to the one reference coupling $\gamma_{\star}$, the ceiling fixed in
Definition~\ref{9.98:def-strong-reference-coupling}.

\begin{corollary}[Reduction to one compact window of couplings]\label{9.98:cor-reduction-compact-window}
Let $\gamma_{\star}$, $X_{\star}=\gamma_{\star}^{-2}$ and the closed window $\overline{W}_{\star}$ be as in
Definition~\ref{9.98:def-prolongation-depths}. For a reference coupling $g$ let
$\operatorname{Lim}_{\infty}(g)$ be the first-passage limit set of Theorem~\ref{9.98:thm-exact-dilates}. For
$\mathsf{s}\in\mathbb{N}$ let $D_{\mathsf{s}}$ be the exact dilation of that theorem,
$(D_{\mathsf{s}}\sigma')(\vec f\,)=\sigma'(\vec f(L^{\mathsf{s}}\cdot))$.
\begin{enumerate}
\item[(a)] Let $\mathcal{P}$ be a property of a single family
$\sigma=(S_n[\sigma])_{n\ge2}$ of infinite-volume Schwinger functions which is invariant under the
exact dilations: for every $\mathsf{s}\in\mathbb{N}$, $\sigma$ has $\mathcal{P}$ if and only if
$D_{\mathsf{s}}\sigma$ has $\mathcal{P}$. Examples of such properties are: uniqueness of the vacuum in
the Osterwalder--Schrader reconstruction of $\sigma$, that is $\ker H=\mathbb{C}\Omega$; clustering
along the coordinate axes; any qualitative Osterwalder--Schrader property; and existence of an
asymptotic expansion at short distance. Then for every $g\le\gamma_{\star}$ the following holds: if
every $\sigma'\in\operatorname{Lim}_{\infty}(\gamma_{\star})$ has $\mathcal{P}$, then every
$\sigma\in\operatorname{Lim}_{\infty}(g)$ has $\mathcal{P}$.
\item[(b)] In this part and in its proof the reference coupling $\gamma_{\star}$ is replaced by
$\min\{\gamma_{\star},\gamma_U'\}$, where $\gamma_U'\le\gamma_U$ is the further coupling threshold of
clause (v-c)(2) of Theorem~\ref{3.1:thm-main}, and $X_{\star}$, $\overline{W}_{\star}$ and $s_{\star}$ of
Definition~\ref{9.98:def-prolongation-depths}, as well as $\gamma_{\natural}$ below, are formed with it.
Let $\operatorname{Lim}_{\infty}[\hat{x}]$ denote the set of infinite-volume limit points, at
renormalised inverse coupling $\hat{x}$, of the continuum limits $S^{(\hat w)}$ of clause (v-b) of
Theorem~\ref{3.1:thm-main} (the tuned frame). Let $\gamma_{\natural}$ be the
coupling ceiling given by
\[
\gamma_{\natural}^{-2}=X_{\star}+B^{\sharp}(s_{\star}+j_U+1),
\]
where $j_U$ is the renaming integer fixed with the pure constants in
Definition~\ref{9.98:def-prolongation-depths}.
Assume the tuned-frame counterpart of Theorem~\ref{9.98:thm-exact-dilates}: for every $\hat{x}$ with
$\hat{x}^{-1/2}\le\gamma_{\natural}$ there are an integer $\mathsf{s}\in\mathbb{N}$, a set
$J_{\mathsf{s}}\subset\mathbb{R}$ containing $\hat{x}$, and a map
$\Phi_{\mathsf{s}}\colon J_{\mathsf{s}}\to\overline{W}_{\star}$, continuous, increasing, bi-Lipschitz and
real-analytic, with $\Phi_{\mathsf{s}}(J_{\mathsf{s}})=\overline{W}_{\star}$, such that
\begin{equation}\label{9.98:eq-reduction-compact-window-1}
\operatorname{Lim}_{\infty}[\hat{x}']=D_{\mathsf{s}}\operatorname{Lim}_{\infty}[\Phi_{\mathsf{s}}(\hat{x}')]
\qquad\text{for every }\hat{x}'\in J_{\mathsf{s}}.
\end{equation}
Then, for every such $\hat{x}$ and with $x^{*}:=\Phi_{\mathsf{s}}(\hat{x})\in\overline{W}_{\star}$,
each of the following statements holds at $\hat{x}$ if and only if the same statement holds at
$x^{*}$:
\begin{enumerate}
\item[(b1)] the set $\operatorname{Lim}_{\infty}[\hat{x}]$ consists of one point;
\item[(b2)] uniqueness of the vacuum, $\ker H=\mathbb{C}\Omega$, holds in the
Osterwalder--Schrader reconstruction of every member of $\operatorname{Lim}_{\infty}[\hat{x}]$;
\item[(b3)] the map $\hat{x}'\mapsto\operatorname{Lim}_{\infty}[\hat{x}']$ on $J_{\mathsf{s}}$ is
continuous at $\hat{x}$ (respectively locally injective, respectively real-analytic), the same
statement at $x^{*}$ being read for the map $x\mapsto\operatorname{Lim}_{\infty}[x]$ on
$\overline{W}_{\star}$.
\end{enumerate}
In (b3), for a map $x\mapsto\operatorname{Lim}_{\infty}[x]$ on a set $I\subset\mathbb{R}$ and a point
$p\in I$:
\emph{local injectivity at $p$} means that some neighbourhood of $p$ in $I$ contains no two distinct
points with the same limit set; \emph{continuity} (respectively \emph{real-analyticity}) \emph{at $p$}
means that on some neighbourhood of $p$ in $I$ every limit set consists of a single family
$\sigma_x$, and that for every $n\ge2$ and every tuple $\vec f$ of real $C^1_c$ functions with
pairwise separated supports the function $x\mapsto S_n[\sigma_x](\vec f\,)$ is continuous at $p$
(respectively is given, on the intersection with $I$ of some open interval around $p$, by a power
series in $x-p$ converging on that open interval).
\end{enumerate}
\end{corollary}

\begin{proof}
\emph{Part (a).} Fix $g\le\gamma_{\star}$ and $\sigma\in\operatorname{Lim}_{\infty}(g)$. By
Theorem~\ref{9.98:thm-exact-dilates}, applied at this $g$, there are an integer $\mathsf{s}$ with
$\underline{s}(g)\le\mathsf{s}\le\overline{s}(g)$ and a family
$\sigma'\in\operatorname{Lim}_{\infty}(\gamma_{\star})$ such that $\sigma=D_{\mathsf{s}}\sigma'$. The
family $\sigma'$ has $\mathcal{P}$ by hypothesis. Since $\mathcal{P}$ is invariant under
$D_{\mathsf{s}}$, the family $D_{\mathsf{s}}\sigma'=\sigma$ has $\mathcal{P}$.

\emph{Part (b): an observation on $D_{\mathsf{s}}$.} Fix $\mathsf{s}$. The map
$\vec f\mapsto\vec f(L^{\mathsf{s}}\cdot)$ is a bijection of the real $C^1_c$ tuples with pairwise
separated supports onto themselves. Its inverse is $\vec h\mapsto\vec h(L^{-\mathsf{s}}\cdot)$. A
rescaling of the argument preserves compactness of supports, the $C^1$ class and positivity of mutual
distances of supports, the distances being multiplied by $L^{\mp\mathsf{s}}$. Hence
\[
\sigma'(\vec h\,)=(D_{\mathsf{s}}\sigma')(\vec h(L^{-\mathsf{s}}\cdot))
\]
for every such tuple $\vec h$. Thus $D_{\mathsf{s}}\sigma'$ determines $\sigma'$, and $D_{\mathsf{s}}$
is injective on families. Since the inverse is itself a relabelling of the test functions,
$D_{\mathsf{s}}$ is a bijection of families whose inverse is the dilation by $L^{-\mathsf{s}}$. In
particular $D_{\mathsf{s}}$ maps any set of families bijectively onto its image.

\emph{Part (b1).} By \eqref{9.98:eq-reduction-compact-window-1}, we have
$\operatorname{Lim}_{\infty}[\hat{x}]=D_{\mathsf{s}}\operatorname{Lim}_{\infty}[x^{*}]$. By the
observation above, the two sets have the same cardinality. So one consists of exactly one point if
and only if the other does.

\emph{Part (b2).} By the same identity, the members of $\operatorname{Lim}_{\infty}[\hat{x}]$ are
exactly the families $D_{\mathsf{s}}\sigma'$ with $\sigma'\in\operatorname{Lim}_{\infty}[x^{*}]$.
Uniqueness of the vacuum is one of the dilation-invariant properties listed in part~(a): the
Osterwalder--Schrader data of $D_{\mathsf{s}}\sigma'$ are those of $\sigma'$ after a rescaling of all
Euclidean coordinates by a power of $L$, so the two reconstructed Hamiltonians are unitarily
equivalent up to a positive constant factor, the unitary carrying the vacuum vector to the vacuum
vector, and their kernels correspond under it. Hence
$D_{\mathsf{s}}\sigma'$ has it if and only if $\sigma'$ has it. Therefore it holds for every member of
$\operatorname{Lim}_{\infty}[\hat{x}]$ if and only if it holds for every member of
$\operatorname{Lim}_{\infty}[x^{*}]$.

\emph{Part (b3).} Write $J=J_{\mathsf{s}}$, $\Phi=\Phi_{\mathsf{s}}$, and
$\Lambda(x)=\operatorname{Lim}_{\infty}[x]$ for $x\in\overline{W}_{\star}$. By
\eqref{9.98:eq-reduction-compact-window-1}, for every $\hat{x}'\in J$ we have
\[
\operatorname{Lim}_{\infty}[\hat{x}']=D_{\mathsf{s}}\bigl(\Lambda(\Phi(\hat{x}'))\bigr).
\]
So on $J$ the map $\hat{x}'\mapsto\operatorname{Lim}_{\infty}[\hat{x}']$ is the composite of three
maps: the chart $\Phi$, then $\Lambda$, then $D_{\mathsf{s}}$ applied to the values.

We first describe the chart. Since $\Phi$ is bi-Lipschitz, there is a constant $c>0$ with
\[
|\Phi(u)-\Phi(v)|\ge c\,|u-v| ,
\]
so $\Phi$ is injective, hence strictly increasing, and its inverse
$\Phi^{-1}\colon\overline{W}_{\star}\to J$ is Lipschitz, hence continuous. Thus
$J=\Phi^{-1}(\overline{W}_{\star})$ is the continuous image of a compact interval, a compact interval,
and $\Phi$ is a homeomorphism of $J$ onto $\overline{W}_{\star}$; in particular it carries the
neighbourhoods of $\hat{x}$ in $J$ onto the neighbourhoods of $x^{*}$ in $\overline{W}_{\star}$. At
every point of $J$, $\Phi$ agrees, on the intersection of $J$ with an open interval around the point,
with a power series converging on that open interval;
since $\Phi$ is increasing with difference quotients at least $c$ on $J$, and the series agrees with
$\Phi$ on the intersection of $J$ with that interval, the series has derivative at least $c$ at
the point. By the real-analytic inverse function theorem applied to this series, $\Phi^{-1}$ agrees,
near every point of $\overline{W}_{\star}$, on the intersection of $\overline{W}_{\star}$ with an open
interval around the point, with a power series converging on that open interval.

We next describe the values. By the observation above, $D_{\mathsf{s}}$ is a bijection of families,
and it maps distinct sets of families to distinct sets. Moreover, by the definition of
$D_{\mathsf{s}}$, for every family $\sigma$, every $n\ge2$ and every tuple $\vec f$,
\[
S_n[D_{\mathsf{s}}\sigma](\vec f\,)=S_n[\sigma](\vec f(L^{\mathsf{s}}\cdot)).
\]

Local injectivity. For $\hat{x}',\hat{x}''\in J$ the sets $\operatorname{Lim}_{\infty}[\hat{x}']$
and $\operatorname{Lim}_{\infty}[\hat{x}'']$ coincide if and only if
$\Lambda(\Phi(\hat{x}'))=\Lambda(\Phi(\hat{x}''))$, since $D_{\mathsf{s}}$ distinguishes distinct sets.
As $\Phi$ is a bijection carrying neighbourhoods of $\hat{x}$ onto neighbourhoods of $x^{*}$, the map
on $J$ is locally injective at $\hat{x}$ if and only if $\Lambda$ is locally injective at $x^{*}$.

Continuity and real-analyticity. By the argument of part~(b1), which uses only
\eqref{9.98:eq-reduction-compact-window-1} and therefore applies at every point of $J$, and since
$\Phi$ matches
neighbourhoods, the limit sets are single families on a neighbourhood of $\hat{x}$ in $J$ if and only
if they are on a neighbourhood of $x^{*}$ in $\overline{W}_{\star}$. In that case, writing $\sigma_x$
for the member of $\Lambda(x)$, the member of $\operatorname{Lim}_{\infty}[\hat{x}']$ is
$D_{\mathsf{s}}\sigma_{\Phi(\hat{x}')}$, and for every $n$ and every tuple $\vec f$
\[
S_n[D_{\mathsf{s}}\sigma_{\Phi(\hat{x}')}](\vec f\,)
=S_n[\sigma_{\Phi(\hat{x}')}](\vec h\,),\qquad \vec h=\vec f(L^{\mathsf{s}}\cdot).
\]
Thus each coordinate function on $J$ at $\vec f$ is the coordinate function on $\overline{W}_{\star}$
at $\vec h$ composed with $\Phi$; conversely the coordinate function on $\overline{W}_{\star}$ at
$\vec h$ is the coordinate function on $J$ at $\vec h(L^{-\mathsf{s}}\cdot)$ composed with
$\Phi^{-1}$, and as $\vec f$ runs through all tuples so does $\vec h$. Composition with the
homeomorphism $\Phi$ or $\Phi^{-1}$ preserves continuity at corresponding points, and composition of
a function real-analytic at a point with a map real-analytic at the corresponding point is
real-analytic there. Hence the map on $J$ is continuous (respectively real-analytic) at $\hat{x}$ if
and only if $\Lambda$ is continuous (respectively real-analytic) at $x^{*}=\Phi(\hat{x})$.
\end{proof}

\begin{remark}
Part~(a) rests on Theorem~\ref{9.98:thm-exact-dilates} alone. Part~(b) is proved as an implication
from the tuned-frame counterpart of Theorem~\ref{9.98:thm-exact-dilates}, that is, from
\eqref{9.98:eq-reduction-compact-window-1} with its chart maps $\Phi_{\mathsf{s}}$. That statement is
itself proved as an implication from the identity on which Theorem~\ref{9.98:thm-exact-dilates} rests,
expressing the bare correlations at reference $g$ through those at the reference coupling
$\gamma_{\star}$, here in the sense of part~(b); from the torus-freeness part of clause (v-b) of
Theorem~\ref{3.1:thm-main}; from rarity of young large-field blocks; from clause (v-c)(1) of
Theorem~\ref{3.1:thm-main}; from the bi-Lipschitz lemma (the lemma on which clause (v-d)(ii) of
Theorem~\ref{3.1:thm-main} depends), which is proved as an implication from the conclusion of the
induction theorem for the re-filed run, whose proof is incomplete at one step, and from the
summability $\sum_n\Upsilon_n<\infty$; from clause (iii-c) of
Theorem~\ref{3.1:thm-main}; and from the statement that the chart maps $\Phi_{\mathsf{s}}$ fit
together into one atlas, whose proof is outstanding.

By part~(a), every dilation-invariant question about an individual limit point at any weak coupling
$g\le\gamma_{\star}$ is the same question for the first-passage limit set at the one reference
coupling $\gamma_{\star}$ of part~(a).
That coupling is moderately weak, and the corresponding bare data form a compact window. By
part~(b), the statements labelled by the renormalised coupling reduce to the window
$\overline{W}_{\star}$ of part~(b) in the same way. In (b3) the map is read on the chart domain
$J_{\mathsf{s}}$; at interior points of $J_{\mathsf{s}}$ this is the unrestricted statement.

Part~(b) rests on an identity of limit sets; part~(a) on the inclusion of
Theorem~\ref{9.98:thm-exact-dilates}. Neither asserts that any of the properties in (a), (b1),
(b2) or (b3) holds. The corollary reduces each of them to one place: an argument establishing any of
them must
control the thermodynamic limit at one moderately weak coupling, in the window
$\overline{W}_{\star}$ of part~(a) or of part~(b), and not along $g\downarrow0$. Possible such
arguments are a last-scale
cluster expansion in infinite volume, a gap uniform in the volume, or absence of zeros of the
partition function.
\end{remark}

\begin{remark}[What the reduction uses]\label{9.98:rem-reduction-uses}
We note three facts about statements surrounding Theorem~\ref{9.98:thm-exact-dilates} (weak-coupling
limit sets as exact dilates) and Corollary~\ref{9.98:cor-reduction-compact-window}
(reduction to one compact window of couplings). None of them enters the proof of
Theorem~\ref{9.98:thm-short-distance-uniqueness} (short-distance uniqueness given the
renormalised coupling), of the lemmas of this section on which that proof rests, or of the
form of clause (iv) for infinite-volume limit points at every
separation below $1/40$ of the reference length. Throughout, $X_{\star}=\gamma_{\star}^{-2}$, and
$\overline{W}_{\star}$ is the closure of the first-passage window
$\mathopen{]}X_{\star},\,X_{\star}+B^{\sharp}]$ at $\gamma_{\star}$.
\begin{enumerate}
\item[(i)] The tuned charts $\Phi_{\mathsf{s}}$ onto $\overline{W}_{\star}$ rest, among other
inputs, on the bi-Lipschitz lemma, whose proof is incomplete at one step, and on a gluing
statement for the charts, whose proof is outstanding.
The charts are needed only for two
purposes. The first is the structural identity of the second assertion of
Corollary~\ref{9.98:cor-reduction-compact-window}: for every weak renormalised coupling
$\hat{x}$, with the prolongation $\mathsf{s}$ and the chart $\Phi_{\mathsf{s}}$ of that
corollary,
\begin{equation*}
\operatorname{Lim}_{\infty}[\hat{x}]
=D_{\mathsf{s}}\operatorname{Lim}_{\infty}[\Phi_{\mathsf{s}}(\hat{x})].
\end{equation*}
The second is the transfer, from one chart to another, of statements on continuity,
analyticity or local injectivity of $\hat{x}\mapsto\operatorname{Lim}_{\infty}[\hat{x}]$;
no such statement is asserted. The charts are not needed for any of the parts (a)--(g) of
Theorem~\ref{9.98:thm-short-distance-uniqueness}.
\item[(ii)] The tuned charts, combined with a further asymptotic-freedom statement that is
not used in this chapter, would give the asymptotic dependence on $\hat{x}$ of the physical
length attached to the renormalised coupling $\hat{x}$. That relation is not asserted.
\item[(iii)] To name the window of $\operatorname{Lim}_{\infty}(\gamma_{\star})$ by
renormalised couplings, one applies the corollary of the uniqueness theorem at reference
coupling $\gamma_{\star}$. This requires $\gamma_{\star}\le\gamma_U'$, which compares two
coupling ceilings; that comparison is not asserted here. The window can instead be named by
bare first-passage data, namely the condition $x_{K'}\in\overline{W}_{\star}$ with
$K'=K(g_0,\gamma_{\star})$. This requires no comparison of $\gamma_{\star}$ with
$\gamma_U'$, and it is all that the first assertion of
Corollary~\ref{9.98:cor-reduction-compact-window} uses.
\end{enumerate}
\end{remark}

\begin{remark}[What does not follow: five counter-states]\label{9.98:rem-five-counter-states}
Let $g$ and $\hat{x}$ be as in Theorem~\ref{9.98:thm-short-distance-uniqueness}
(short-distance uniqueness given the renormalised coupling). Let $\sigma$ be an infinite-volume
limit point at renormalised coupling $\hat{x}$, that is, an element of
$\operatorname{Lim}_{\infty}[\hat{x}]$.

\emph{Not asserted.} Theorem~\ref{9.98:thm-short-distance-uniqueness},
Theorem~\ref{9.98:thm-exact-dilates} and Corollary~\ref{9.98:cor-reduction-compact-window} do not
assert any of the following.
\begin{itemize}
\item Uniqueness of the infinite-volume limit, that is, independence of the subsequence of
volumes, jointly with the cutoff.
\item Uniqueness of the vacuum. This is listed separately from the previous item: a degenerate
vacuum on periodic tori would show as absence of clustering of a limit that may itself be
unique.
\item Clustering; a mass gap.
\item Continuity, analyticity or local injectivity of the map from the renormalised coupling to
the infinite-volume theory at fixed distances.
\item Independence of the shape of the torus.
\item Agreement at short distance of the $n$-point functions for $n\ge4$ beyond the boundedness
stated in clause (f) of Theorem~\ref{9.98:thm-short-distance-uniqueness}.
\item Any improvement of the exponent $4-n$, $n\in\{2,3\}$, of $\mathsf{d}$ in the discrepancy.
\item Anything at separations $\mathsf{d}>L^{\underline{s}'(g)}/40$, that is, up to a pure
factor, beyond the fixed physical length $\ell_{\star}/40$, where $\ell_{\star}$ is the reference
length of $\gamma_{\star}$.
\item Convergence in the torus exponent $m$; a relative form of the bound at $n=3$; anything at
complex $\hat{x}$ on $\mathbb{R}^4$.
\item That the two-point function of clause (iv) of Theorem~\ref{3.1:thm-main} takes at short
distance the free-field form, beyond the two-sided bounds stated there; the values of its two
universal numbers; any determination of the one-loop coefficient from the limit.
\item Any uniqueness statement through summation of a perturbation series. The only statement of
perturbative type made is that of clause (d) of
Theorem~\ref{9.98:thm-short-distance-uniqueness}: if one limit point at a given renormalised
coupling admits an asymptotic expansion to some order, then all limit points at that coupling
admit the same expansion; the existence of such an expansion is not asserted.
\item Independence of the regularisation or of the blocking factor $L$.
\item Uniqueness in the cutoff on a torus without the words ``given the renormalised coupling''.
The conditioning cannot be dropped: this follows from the non-convergence of the first-passage
offsets, together with the non-constancy of
$\hat{x}\mapsto\mathfrak{S}_m[\hat{x}](f_1,f_2)$ on every subinterval of the range of $\hat{x}$,
for the pairs $(f_1,f_2)$ to which the non-constancy statement on tori applies.
\item Convergence in the cutoff of the tuned family on non-cubic tori.
\item Any characterisation of non-uniqueness in the volume through long-range order, and any
uniqueness statement conditional on a hypothesis excluding long-range order. No such statement is
part of Theorem~\ref{9.98:thm-short-distance-uniqueness}.
\item Anything about Wilson loops, block observables, or local observables other than the
smeared action density.
\item The old-age tail statement; the two-point witness on every torus with a depth floor growing
with the torus exponent (Theorem~\ref{8.18:thm-torus-witness}), which the proof of
Theorem~\ref{9.98:thm-short-distance-uniqueness} does not use; volume-uniform rarity of
large-field regions.
\item Anything for $N\ge3$.
\end{itemize}

\emph{Five counter-states.} Each of the following is elementary, and the five are mutually
independent. Each satisfies the following statements:
\begin{itemize}
\item clause (0) of Theorem~\ref{3.1:thm-main};
\item the invariance, reflection positivity and Osterwalder--Schrader reconstruction of clause
(ii) on $\mathbb{R}^4$;
\item bounds of the type of clause (iii-b);
\item the two-sided law of clause (iv) of Theorem~\ref{3.1:thm-main} in every volume, with its
depth floor possibly enlarged;
\item the change-of-units identity together with clause (0);
\item the bound of order $\mathsf{d}^2$ of Theorem~\ref{9.98:thm-short-distance-uniqueness}.
\end{itemize}
Each nevertheless differs from $\sigma$ on pairs of test functions at separation one, measured in
the reference units in which $\mathsf{d}$ is measured in
Theorem~\ref{9.98:thm-short-distance-uniqueness}; we say ``at distance one'' for this. The
properties listed under
(c1)--(c5), and the two facts just stated, are taken by statement from the statements of this
book that construct these five states.
\begin{itemize}
\item[(c1)] \emph{The frozen random constant.} Replace $\mathcal{O}(f)$ by
$\mathcal{O}(f)+\xi\int f$, where $\xi$ is a real random variable independent of the field. Let
$\operatorname{cum}_j(\xi)$ denote the $j$-th cumulant of $\xi$. Connected functions of a sum of
independent fields add, so the connected $j$-point function shifts by
\begin{equation}\label{9.98:eq-five-counter-states-1}
\operatorname{cum}_j(\xi)\prod_{i=1}^j\int f_i ,
\end{equation}
which, for test functions supported in balls of radius of order $\mathsf{d}$, is bounded in
absolute value by $|\operatorname{cum}_j(\xi)|\,(c\,\mathsf{d}^4)^j\,\Pi$, where
$\Pi=\prod_{i=1}^j\|f_i\|_{\infty}$ and $c$ depends only on the ratio of the radius to
$\mathsf{d}$. If $\operatorname{Var}\xi>0$ there are two or more translation-invariant vectors;
the vacuum is degenerate if and only if $\operatorname{Var}\xi>0$. Reflection positivity,
invariance, the bounds of the type of clause (iii-b), the display of clause (iv) below a larger
depth floor, and the bound of order $\mathsf{d}^2$ all survive. This state arises from the tori
if and only if $\operatorname{Var}_m\bigl(|\mathbb{T}_m|^{-1}\mathcal{O}(1)\bigr)$ does not tend
to $0$ as $m\to\infty$, where $\operatorname{Var}_m$ is the variance under the torus measure and
$\mathcal{O}(1)$ is the observable with $f\equiv1$; equivalently, the infinite-volume pressure
has a kink at the coupling considered. It also satisfies the conclusion of the lemma on
reflection positivity of the bare Wilson measure in operator form that is used in the proof of
Theorem~\ref{9.98:thm-short-distance-uniqueness}.
\item[(c2)] \emph{A decoupled Gaussian field.} Add to $\sigma$ an independent Gaussian field with
prescribed covariance, of scaling dimension
$\Delta\in[1,3]$, with the $L$-adic amplitude law of the proposition that constructs this state.
The result clusters if
$\sigma$ does, and the discrepancy is of order $\mathsf{d}^{8-2\Delta}\le\mathsf{d}^2$.
\item[(c3)] \emph{The Wick square of a decoupled free scalar.} Take the tensor product of
$\sigma$ with an independent free massless scalar field $\psi$; replace the observable
$\mathcal{O}(f)$ by $\mathcal{O}(f)+\eta\,{:}\psi^2{:}(f)$, where ${:}\psi^2{:}(f)$ is the Wick
square of $\psi$ smeared with $f$, and the stress-tensor density $\Theta_{\mu\nu}$ by
$\Theta_{\mu\nu}+T^{\psi}_{\mu\nu}$, where $\eta>0$ is a parameter and $T^{\psi}_{\mu\nu}$ is the
stress tensor of $\psi$. The result is reflection positive. Whenever $\sigma$ is
$\mathrm{O}(4)$-invariant with the local conserved traceless stress tensor $\Theta_{\mu\nu}$, the
result is $\mathrm{O}(4)$-invariant with the local conserved traceless stress tensor
$\Theta_{\mu\nu}+T^{\psi}_{\mu\nu}$; it clusters whenever $\sigma$ does. Clause (vi) of
Theorem~\ref{3.1:thm-main} displays, beside the far-sphere flux hypothesis
$(\mathrm{IR}\text{-}\mathrm{fl})_{\omega_0}$, a sufficient condition for it: one far-away
stress-tensor insertion decorrelates from the observables faster than the inverse fourth power of
its distance. If the connected functions of $\sigma$ with one insertion of $\Theta_{\mu\nu}$ have
this decay, so do those of the result with one insertion of $\Theta_{\mu\nu}+T^{\psi}_{\mu\nu}$,
the part coming from $\psi$ decaying as the inverse sixth power of the distance. With
$\sigma'_q$ for the connected $q$-point functions of the result and $\sigma_q$ for those of
$\sigma$,
\begin{equation*}
\sigma'_2-\sigma_2\asymp\eta^2\mathsf{d}^4\le C_T^{\star}\mathsf{d}^2,\qquad
\sigma'_3-\sigma_3=O(\eta^3\mathsf{d}^6),
\end{equation*}
and $\sigma'_2-\sigma_2$ at distance one is $O(\eta^2)$. The parameter $\eta$ is
scaled like
$L^{-2\mathsf{s}}$ along the dilation orbits of Theorem~\ref{9.98:thm-exact-dilates}. This state
shows that adding to the statements listed the requirement that every limit point be reflection
positive, $\mathrm{O}(4)$-invariant with a local conserved traceless stress tensor, clustering,
and subject to that sufficient condition still does not give uniqueness of the infinite-volume
limit; its
variant with a massive free scalar in place of $\psi$ shows that a mass gap on every limit point
would not give it either.
\item[(c4)] \emph{A positive-type kernel.} Add to $\sigma_2$ a smooth bounded kernel of positive
type.
\item[(c5)] \emph{A one-parameter family.} A one-parameter family of reflection-positive,
Euclidean-invariant hierarchies of Schwinger functions, in the sense of the statement of this
book that constructs the family, with one and the same ultraviolet germ in the sense defined
there and different mass gaps.
\end{itemize}

\emph{Consequence.} Read the statements listed before (c1)--(c5), including the bound of
Theorem~\ref{9.98:thm-short-distance-uniqueness}, as axioms on limit points. No deduction from
these statements gives uniqueness of the infinite-volume limit, uniqueness of the vacuum,
clustering, or continuity or analyticity of the dependence on the renormalised coupling on
$\mathbb{R}^4$. This does not say that any of these
properties fails for Yang--Mills theory.

\emph{Where the volume enters.} In the proofs of
Theorem~\ref{9.98:thm-short-distance-uniqueness} and of clause (v) of
Theorem~\ref{3.1:thm-main}, the large volume enters the measures only through the
free-energy-density mismatch $\mathsf{B}=E_{+}^{\sharp}+C_{\mathrm{low}}$ per cell at the
stopping scale. It appears in three places:
\begin{itemize}
\item globally, as $e^{\mathsf{B}(2L^m)^4}$;
\item per chessboard cell, as the factor $e^{\mathsf{B}\ell_n^4}$ in $Q_{\star,n}$;
\item in clause (v), as
\begin{equation*}
\mathsf{Q}_w=e^{(E_{+}^{\sharp}+C_{\mathrm{low}})N_w^4}.
\end{equation*}
\end{itemize}
At fixed separation this is a constant; it is compensated only by the factor $\mathsf{d}^4$
gained by averaging over translates in the proof of
Theorem~\ref{9.98:thm-short-distance-uniqueness}, hence only as $\mathsf{d}\downarrow0$ below
the fixed length $\ell_{\star}$. None of the bounds used in these proofs decays in the separation
$\mathsf{d}$; the per-ball volume-free bound asserts no decay in $\mathsf{d}$. \cite{OSe78}
states uniqueness of the infinite-volume lattice Yang--Mills measure at strong bare coupling; no
uniqueness statement for the infinite-volume limit at weak coupling is used or proved here. By
Corollary~\ref{9.98:cor-reduction-compact-window} (reduction to one compact window of couplings),
each of the following properties --- uniqueness of the infinite-volume limit; uniqueness of the
vacuum; continuity, analyticity or local injectivity in the coupling --- holds at every weak
renormalised coupling as soon as it holds at every coupling of the one compact window
$\overline{W}_{\star}$.
\end{remark}

\begin{proof}
Let $\xi$ be as in (c1) with $\operatorname{Var}\xi>0$. By the properties stated before and
under (c1), the resulting state satisfies the statements listed before (c1)--(c5), the law of
clause (iv) being read below the enlarged depth floor. The second
cumulant of $\xi$ is its variance, so by \eqref{9.98:eq-five-counter-states-1} with $j=2$ the
connected two-point function of this state differs from that of $\sigma$ by
\begin{equation*}
\operatorname{Var}\xi\int f_1\int f_2 ,
\end{equation*}
which is non-zero for every pair of test functions with $\int f_1\ne0$ and $\int f_2\ne0$, and
such pairs exist at distance one. Thus $\sigma$ and this state are two distinct objects satisfying
these
statements, read as axioms on limit points; hence these statements do not give uniqueness of the
infinite-volume limit. By the properties stated under (c1), this state has two or more
translation-invariant vectors, so the statements do not give uniqueness of the vacuum. The
difference displayed above does not change when $f_2$ is replaced by a translate, since the
integral of a translate is the integral of the function. Hence, if $\sigma$ clusters, the
connected two-point function of this state tends to $\operatorname{Var}\xi\int f_1\int f_2\ne0$
as the support of $f_2$ is moved away from that of $f_1$, so this state does not cluster; if
$\sigma$ does not cluster, $\sigma$ itself satisfies the statements without clustering. In either
case the statements do not give clustering.

Finally, fix an interior point $\hat{x}_1$ of the range of $\hat{x}$ in
Theorem~\ref{9.98:thm-short-distance-uniqueness}, a variance $v_0>0$ admissible in (c1), and, for
each $\hat{x}$ in that range, an element $\sigma_{\hat{x}}$ of
$\operatorname{Lim}_{\infty}[\hat{x}]$. Fix a pair $(f_1,f_2)$ at distance one with
$\int f_1\ne0$ and $\int f_2\ne0$. Put $v(\hat{x})=0$ for $\hat{x}\le\hat{x}_1$ and
$v(\hat{x})=v_0$ for $\hat{x}>\hat{x}_1$, and at each $\hat{x}$ apply (c1) to $\sigma_{\hat{x}}$
with a random variable of variance $v(\hat{x})$; at each coupling the resulting state satisfies
the statements listed before (c1)--(c5), as shown above. By
\eqref{9.98:eq-five-counter-states-1} with $j=2$, its connected two-point function at
$(f_1,f_2)$ is the value of $\sigma_{\hat{x}}$ there plus $v(\hat{x})\int f_1\int f_2$, and the
second term jumps by $v_0\int f_1\int f_2\ne0$ at $\hat{x}_1$. If the map $\hat{x}$ to the value
of $\sigma_{\hat{x}}$ at $(f_1,f_2)$ is discontinuous at $\hat{x}_1$, the family
$(\sigma_{\hat{x}})$ itself satisfies the statements with a dependence on the coupling that is
neither continuous nor real-analytic; if it is continuous at $\hat{x}_1$, the modified family has
a two-point function that is discontinuous at $\hat{x}_1$, hence neither continuous nor
real-analytic in $\hat{x}$. In either case the statements do not give continuity or analyticity
of the dependence on the renormalised coupling.
\end{proof}

\paragraph{Inputs used by statement.}
This section lists, for each clause of Theorem~\ref{3.1:thm-main} in its printed order, the
published statements of Ba{\l}aban used by statement in this chapter's proofs and the statements of
this book used there whose proofs are incomplete or outstanding.

\begin{description}
\item[(0)] No part of clause~(0) is proved in this chapter.
\item[(i)] No part of clause~(i) is proved in this chapter.
\item[(iii)] No part of clause~(iii) is proved in this chapter.
\item[(ii)] The torus statement (ii-$\mathbb{T}$-a) is the convergence statement of clause~(v) and
is proved with it. Its inputs are those listed under~(v). No other part of clause~(ii) is proved
here.
\item[(iv)] No part of clause~(iv) is proved in this chapter.
\item[(v)] Published statements cited by statement:
\begin{itemize}
\item \cite[Theorem~1, (181)]{B-CMP102b}
\item \cite[Proposition~9]{B-CMP102b}
\item \cite[Proposition~3]{B-CMP102b}
\item \cite[Proposition~5]{B-CMP102b}
\item \cite[Proposition~6]{B-CMP102b}
\item \cite[(189)--(190), p.~308]{B-CMP102b} (printed without proof; cited as printed)
\item \cite[Lemma~3]{B-CMP116}
\item \cite[Theorem~3.1]{B-CMP99b}
\item \cite[Theorem~3.4]{B-CMP99b}
\item \cite[Theorem~3.12]{B-CMP99b}
\item \cite[Theorem~3.13]{B-CMP99b}
\item \cite[(3.35)]{B-CMP99b}
\item \cite[(3.37)--(3.38)]{B-CMP99b}
\item \cite[(3.187)]{B-CMP99b}
\item \cite[Lemma~4]{B-CMP109}
\item \cite[(3.15)]{B-CMP109}
\item \cite[Proposition~1]{B-CMP98}
\item \cite[Proposition~2]{B-CMP98}
\item \cite[Proposition~4]{B-CMP98}
\item \cite[Proposition~5]{B-CMP98}
\item \cite[Lemma~1]{B-CMP99a}
\item \cite[Proposition~3]{B-CMP99a}
\item \cite[Proposition~5]{B-CMP99a}
\item \cite[Lemma~2.1]{B-CMP96}
\item \cite[(2.60)--(2.61)]{B-CMP96}
\item \cite[pp.~178--179]{B-CMP122a} (printed without proof; cited as printed)
\item \cite[(1.70), p.~192]{B-CMP122a} (printed without proof; cited as printed)
\item \cite[(1.68)--(1.69)]{B-CMP122a}
\item \cite[(1.90)--(1.99)]{B-CMP122a}
\item \cite[(1.100)]{B-CMP122a}
\item \cite[(1.33)]{B-CMP122b}
\item \cite[(1.61)--(1.63)]{B-CMP122b}
\item \cite[(1.68)--(1.69)]{B-CMP122b}
\item \cite[(1.71)]{B-CMP122b}
\item \cite[(1.72)]{B-CMP122b}
\item \cite[(1.89)]{B-CMP122b}
\item \cite[(1.90)--(1.91)]{B-CMP122b}
\item \cite[(1.92)--(1.100)]{B-CMP122b}
\item \cite[p.~378]{B-CMP122b}
\item \cite[p.~390]{B-CMP122b}
\item \cite[(2.18)--(2.22)]{B-CMP119}
\item \cite[(2.27)(i)--(iv)]{B-CMP119}
\item \cite[(2.30)--(2.32)]{B-CMP119}
\item \cite[(2.40)--(2.42)]{B-CMP119}
\item \cite[(2.43)--(2.49)]{B-CMP119}
\item \cite[(2.50)]{B-CMP119}
\item \cite[(3.47)]{B-CMP119}
\item \cite[(3.61)]{B-CMP119}
\end{itemize}
Statements of this book whose proofs are incomplete or outstanding:
\begin{itemize}
\item The assertions on the radius exponents, the floors on $\kappa_1$ and the fat core made in
Definition~\ref{9.1:not-radius-exponent-floors} --- proof outstanding
\item The assertions on the floor chain on shells and the ratio $A_0/A_1$ made in
Definition~\ref{9.1:not-shell-floor-chain} --- proof outstanding
\item The assertions on the level-wise rate, the number of stages and the strip width made in
Definition~\ref{9.1:not-levelwise-rate} --- proof outstanding
\item Proposition~\ref{9.1:prop-blind-site-expansions} (site expansions with island-blind inputs)
--- proof outstanding
\item The carry lemma, Theorem~\ref{9.2:thm-carry} --- stated; proof incomplete
\item The extension proposition, on which the carry lemma depends --- stated; proof incomplete
\item Zone-resolved variation of the first-group boundary terms --- stated; proof incomplete
\item The tail of the rest terms above a threshold zone --- stated; proof incomplete
\item Aligned parts of nested leaves and the unmatched term --- stated; proof incomplete
\item The first index over the full history --- stated; proof incomplete
\item The per-unit response factor at the step --- stated; proof incomplete
\item Lemma~\ref{9.2:lem-read-set} (read sets of cells meeting the domain) --- stated; proof
incomplete
\item The boundary-family entry of the Lipschitz step --- stated; proof incomplete
\item Linear size at least one for constructed regions --- stated; proof incomplete
\item Proposition~\ref{9.2:prop-modulus-majorant} (the history-wise modulus majorant) --- proof
outstanding
\item The nesting clause of the majorant hypothesis used in the proof of
Theorem~\ref{9.2:thm-geometric-forgetting} --- proof outstanding
\item Proposition~\ref{9.3:prop-propagator-stationarity} (weak stationarity of the fine-lattice
propagators) --- proof outstanding
\item The two-cutoff response statement at zeroth order, carried over to the shells used by
clauses~(v) and~(ii-$\mathbb{T}$-a), which are not the shells on which it is proved --- proof
outstanding
\item The one-run expansions on labels --- stated; proof incomplete
\item The residual clauses of the increment statement, on which the proof of clause~(v) depends
--- proof outstanding
\item The bound on the difference, between the two runs, of the remainder terms of case~(ii) of
\cite{B-CMP122a}, whose bounds that paper states without giving details --- proof outstanding
\item The bound on the difference, between the two runs, of the nested Gaussian integrals
interrupted by large-field regions --- proof outstanding
\item The bound at the current depth for the two factors kept un-integrated inside a large-field
region, which depends on the two preceding items --- proof outstanding
\item Lemma~\ref{9.1:lem-glued-minimisers} (ambient regularity of the two runs), in the
form in which the stronger regularity conditions are imposed only where a proof reads them ---
stated; proof incomplete
\item Lemma~\ref{9.6:lem-maximum-principle} (a maximum principle for coupled scale-indexed
inequalities) --- proof outstanding
\item The assertions made in Definition~\ref{9.6:def-comparison-chain} --- proof outstanding
\item Lemma~\ref{9.6:lem-chain-structure} (well-definedness and structure of the comparison chain
of Definition~\ref{9.6:def-comparison-chain}) --- proof outstanding
\item Corollary~\ref{9.8:cor-subsequential-limits} (first-passage offsets and the identification of
subsequential limits) --- proof outstanding
\item Proposition~\ref{9.9:prop-remainder-derivative} (bare-offset derivative of the witness
remainder) --- proof outstanding
\item Gaussian skeleton and in-box precision --- proof outstanding
\item Theorem~\ref{9.10:thm-refiled-induction} (the induction theorem for the re-filed run) ---
stated; proof incomplete
\item The statement bounding the regrouped terms in the two-run identity at the real centres for
the re-filed boundary terms, by following the walk of the absent run inside a tube of prescribed
radius in the depth --- proof outstanding
\item The bound on the own-history part of the boundary terms produced at a large-field step (the
terms filed at the producing block's own level, including those inside the regions healed at that
step), taken from the corresponding clause for the history part of the boundary terms --- proof
outstanding
\item Theorem~\ref{9.10:thm-refiled-lipschitz} (re-filed one-lattice Lipschitz step) --- stated;
proof incomplete
\item Kernel form of the re-filed Lipschitz step --- stated; proof incomplete
\item Geometric source bounds for the re-filed run --- stated; proof incomplete
\item Finite-range envelope of last-lattice loads --- stated; proof incomplete
\item Geometric last-lattice comparison for the re-filed run --- stated; proof incomplete
\item Hatted source families with geometric envelope --- stated; proof incomplete
\item The endpoint case of Corollary~\ref{9.12:cor-same-limit} --- proof outstanding
\item Lemma~\ref{9.13:lem-limit-member} (the limit-member identification) --- proof outstanding
\item Proposition~\ref{9.14:prop-sun-uniqueness} (the uniqueness theorem for
$\mathrm{SU}(N)$, $N\ge 3$), used only for $N\ge 3$ --- proof outstanding
\end{itemize}
The corollary after clause~(v), Theorem~\ref{9.98:thm-short-distance-uniqueness}, is proved as an
implication from clause~(v), from Corollary~\ref{9.8:cor-subsequential-limits}, from the
correlation bounds of clause~(iii) and from clause~(iv)
(the two-point witness theorem and the volume-transfer bound). Within this chapter its proof cites
no published statement by statement; of the statements listed above it uses only
Corollary~\ref{9.8:cor-subsequential-limits}.
\item[(vi)] No part of clause~(vi) is proved in this chapter.
\end{description}

None of the statements listed under~(v) is
proved with the help of clause~(v), of the corollary after it, or of any statement of this chapter
whose proof uses either of them. The constants entering the proofs of this chapter are chosen in
the order of Definitions~\ref{2.3:not-stages-first}, \ref{2.3:not-stages-middle}
and~\ref{2.3:not-stages-last}, extended by Definitions~\ref{9.1:not-radius-exponent-floors},
\ref{9.1:not-shell-floor-chain} and~\ref{9.1:not-levelwise-rate}; by Lemma~\ref{2.4:lem-no-cycle}
that order, before this extension, has no cycle, and no constant in it is bounded in terms of a
constant chosen after it. The statements contained in the three extending definitions are listed
under~(v) above.

%% ---- end inlined file: chapters/c09 ----

\clearpage
\setcounter{section}{9}
\refstepcounter{section}
\section*{Chapter~\thesection\quad $\mathrm{O}(4)$ invariance under an infrared decay hypothesis: clause (vi)\addcontentsline{toc}{section}{\protect\numberline{\thesection}$\mathrm{O}(4)$ invariance under an infrared decay hypothesis: clause (vi)}\label{ch:10}}
%% ---- begin inlined file: chapters/c10 ----
\begin{definition}[Bare lattice, sourced expectations and rotations]\label{10.1:not-bare-lattice}
Throughout this chapter the gauge group is $\mathrm{SU}(2)$, and $\operatorname{tr}$ denotes the
trace normalised by $\operatorname{tr}\mathbf{1}=1$. The objects are those of the setting of
Theorem~0, specialised to $N=2$.
\begin{itemize}
\item \emph{Bare lattice.} Put $a:=L^{-K}$ and
\begin{equation}\label{10.1:eq-bare-lattice-1}
\mathfrak{L}:=a\bigl(\mathbb{Z}^4+\tfrac12(1,1,1,1)\bigr)\big/\,2L^m\mathbb{Z}^4 ,
\qquad \hat\mu:=a e_\mu ,
\end{equation}
where $e_1,\dots,e_4$ is the standard basis of $\mathbb{R}^4$; thus $\mathfrak{L}$ is the lattice
$a(\mathbb{Z}^4+\tfrac12(1,1,1,1))$ periodised by $2L^m\mathbb{Z}^4$.
\item \emph{Links and plaquettes.} A positively oriented link $\ell=(x,\mu)$, $x\in\mathfrak{L}$, has
midpoint $x_\ell:=x+\tfrac12\hat\mu$. A plaquette $p=\{x;\mu,\nu\}$, with $\mu\neq\nu$ an unordered
pair, has barycentre $x_p:=x+\tfrac12\hat\mu+\tfrac12\hat\nu$, also written $x(p)$;
$p\pm\hat\rho:=\{x\pm\hat\rho;\mu,\nu\}$ are its translates.
\item \emph{Action density and observable.} For bond variables $U$ put
$a_p(U):=1-\operatorname{Re}\operatorname{tr}U(\partial p)$ and
$\mathcal{O}(f):=\sum_p f(x_p)\,a_p(U)$, the sum running over the plaquettes of all six
orientations.
\item \emph{Sourced expectations.} For $w$ a function from the plaquettes to $\mathbb{C}$ put
\begin{equation}\label{10.1:eq-bare-lattice-2}
S_w:=\sum_p c(p)\,a_p(U),\qquad c(p):=g_0^{-2}-w(p),\qquad Z[w]:=\int e^{-S_w}\,dU ,
\end{equation}
with $dU$ the Haar measure on the bond variables of the setting of Theorem~0, and let
$\langle\cdot\rangle_w$ be the complex normalised expectation
$\langle Y\rangle_w:=Z[w]^{-1}\int Y\,e^{-S_w}\,dU$, defined wherever $Z[w]\neq0$.
\item \emph{Schwinger functions.} With $w(p)=\sum_{i=1}^n z_i f_i(x_p)$, $z_i\in\mathbb{C}$,
\begin{equation}\label{10.1:eq-bare-lattice-3}
S^T_{n;K,m}(f_1,\dots,f_n)
=\partial_{z_1}\cdots\partial_{z_n}\log Z[w]\big|_{z=0}
=\langle\mathcal{O}(f_1);\dots;\mathcal{O}(f_n)\rangle ,
\end{equation}
the right-hand side being the joint cumulant under $\langle\cdot\rangle_0$. For $Y$ not depending
on $z$ we write
\begin{equation*}
\langle Y;\mathcal{O}(f_1);\dots;\mathcal{O}(f_n)\rangle
:=\partial_{z_1}\cdots\partial_{z_n}\langle Y\rangle_w\big|_{z=0}.
\end{equation*}
\item \emph{Norm.} $\|f\|_\sharp:=\max\{\|f\|_\infty,\ 40M\cdot\operatorname{Lip}f\}$.
\item \emph{Symmetries.} $W_4$ is the group of the $384$ signed permutation matrices;
$V:=\mathbb{R}^4$ and $\Lambda^2:=\Lambda^2\mathbb{R}^4$ are regarded as $W_4$-modules. For
$R\in \mathrm{O}(4)$ put $(R\cdot f)(x):=f(R^{-1}x)$. For $\omega$ a real antisymmetric $4\times4$
matrix put
$\xi(x):=\omega x$, that is $\xi^\rho=\omega_{\rho\sigma}x_\sigma$, and, for $f\in C^1$,
\begin{equation}\label{10.1:eq-bare-lattice-4}
\delta_\omega f:=\frac{d}{dt}\bigl(e^{t\omega}\cdot f\bigr)\Big|_{t=0}=-\,\xi\cdot\nabla f .
\end{equation}
\end{itemize}
\end{definition}

\begin{proof}
The notation contains three assertions: that the objects in \eqref{10.1:eq-bare-lattice-3} are
defined, the second equality in \eqref{10.1:eq-bare-lattice-3}, and the second equality in
\eqref{10.1:eq-bare-lattice-4}.

Since $\mathfrak{L}$ is finite, the bond variables range over a finite product of
copies of the compact group $\mathrm{SU}(2)$. The integrand $e^{-S_w}$ is continuous in $U$,
entire in $(z_1,\dots,z_n)$ and bounded uniformly in $U$ for $z$ in compact sets, so $Z[w]$ is
entire in $z$. At $z=0$ we have $S_0=g_0^{-2}\sum_p a_p(U)=g_0^{-2}A(U)$, so
$Z[0]=\int e^{-g_0^{-2}A(U)}\,dU>0$. Hence $Z[w]\neq0$ on a neighbourhood of $z=0$, on which
$\log Z[w]$ and $\langle\cdot\rangle_w$ are holomorphic in $z$, and the derivatives at $z=0$ exist.

For the cumulant identity,
\begin{equation*}
\sum_p w(p)a_p(U)=\sum_i z_i\sum_p f_i(x_p)a_p(U)=\sum_i z_i\,\mathcal{O}(f_i),
\end{equation*}
whence
\begin{align*}
e^{-S_w}&=e^{-g_0^{-2}A(U)}\exp\Bigl(\sum_i z_i\mathcal{O}(f_i)\Bigr),\\
Z[w]&=Z[0]\,\Bigl\langle\exp\Bigl(\sum_i z_i\mathcal{O}(f_i)\Bigr)\Bigr\rangle_0 .
\end{align*}
Thus
$\log Z[w]-\log Z[0]$ is the cumulant generating function of the vector
$(\mathcal{O}(f_1),\dots,\mathcal{O}(f_n))$ under $\langle\cdot\rangle_0$, and its mixed derivative
$\partial_{z_1}\cdots\partial_{z_n}$ at $z=0$ is, by the definition of joint cumulants,
$\langle\mathcal{O}(f_1);\dots;\mathcal{O}(f_n)\rangle$.

For \eqref{10.1:eq-bare-lattice-4}: since $\omega$ is antisymmetric, the transpose of $e^{t\omega}$
is $e^{-t\omega}$, so $e^{t\omega}$ is orthogonal; moreover
$\det e^{t\omega}=e^{t\operatorname{Tr}\omega}=1$, where $\operatorname{Tr}$ is the ordinary matrix
trace on $4\times4$ matrices, which vanishes on antisymmetric matrices; so
$e^{t\omega}\in \mathrm{SO}(4)\subset\mathrm{O}(4)$ and
$(e^{t\omega}\cdot f)(x)=f(e^{-t\omega}x)$. By the chain rule,
\begin{equation*}
\frac{d}{dt}f(e^{-t\omega}x)\Big|_{t=0}=\nabla f(x)\cdot(-\omega x)=-\xi(x)\cdot\nabla f(x).
\end{equation*}
\end{proof}

\begin{definition}[Quaternion components and link derivations]\label{10.1:not-link-derivations}
We use the bare lattice $\mathfrak{L}$, its links $\ell$, its plaquettes $p$ and the plaquette action
density $a_p(U) = 1 - \operatorname{Re}\operatorname{tr} U(\partial p)$ as in
Definition~\ref{10.1:not-bare-lattice}. The gauge group is $SU(2)$ and $\operatorname{tr}\mathbf{1} = 1$.
Let $\vec\sigma = (\sigma^1,\sigma^2,\sigma^3)$ be the Pauli matrices. Write $\delta^{bd}$ for the Kronecker
symbol and $\varepsilon^{bcd}$ for the totally antisymmetric symbol with $\varepsilon^{123} = 1$. Repeated
colour indices $b,c,d \in \{1,2,3\}$ are summed unless the sum is written out.

\emph{Quaternion components.} Every $W \in SU(2)$ is written as
\[
W = w^0 + i\,\vec w\cdot\vec\sigma, \qquad w^0 \in \mathbb{R},\ \vec w \in \mathbb{R}^3,
\qquad (w^0)^2 + |\vec w|^2 = 1,
\]
so that $\operatorname{tr} W = w^0$.

\emph{Holonomies through a link.} For a link $\ell \subset \partial p$, let $W_{\ell,p}$ be the holonomy of
$\partial p$ that starts at the initial point of $\ell$ and runs through $\ell$ first, in the positive
direction of $\ell$. Holonomies are written as ordered products in which the link traversed first stands on
the left, so that $U_\ell$ is the leftmost factor of $W_{\ell,p}$. Its components are denoted $w^0_{\ell,p}$
and $\vec w_{\ell,p}$. Let $\mu(\ell)$ be the direction of $\ell$ and $\hat\rho := a e_\rho$. For
$\rho \neq \mu(\ell)$, let $p(\ell,\pm\rho)$ be the plaquette spanned by $\ell$ and $\pm\hat\rho$, and set
$W_{\ell,\pm\rho} := W_{\ell,p(\ell,\pm\rho)}$.

\emph{Link derivations.} For a function $f$ of the configuration $U$, a link $\ell$ and $b = 1,2,3$, define
\[
(\partial^b_\ell f)(U) := \frac{d}{ds} f(\dots, e^{is\sigma^b} U_\ell, \dots)\Big|_{s=0}.
\]

Then for every plaquette $p$ and every link $\ell \subset \partial p$,
\begin{equation}\label{10.1:eq-link-derivations-a}
a_p = 1 - w^0_{\ell,p}, \qquad |\vec w_{\ell,p}|^2 = a_p(2 - a_p) \le 2a_p,
\end{equation}
and
\begin{equation}\label{10.1:eq-link-derivations-b}
\begin{split}
&\partial^b_\ell w^0_{\ell,p} = -w^b_{\ell,p}, \qquad
\partial^b_\ell w^d_{\ell,p} = w^0_{\ell,p}\,\delta^{bd} - \varepsilon^{bcd} w^c_{\ell,p};\\
&\text{hence}\quad \partial^b_\ell a_p = w^b_{\ell,p}, \qquad
\sum_{b=1}^{3} \partial^b_\ell w^b_{\ell,p} = 3w^0_{\ell,p} = 3(1 - a_p).
\end{split}
\end{equation}
\end{definition}

\begin{proof}
We first prove \eqref{10.1:eq-link-derivations-a}. The holonomy $U(\partial p)$ in the definition of $a_p$ and
the holonomy $W_{\ell,p}$ are products of the same four link variables around $\partial p$. They can differ
only by a cyclic change of the starting point and possibly by a reversal of orientation. Hence $W_{\ell,p}$ is
conjugate in $SU(2)$ either to $U(\partial p)$ or to $U(\partial p)^{-1}$. The trace is invariant under
conjugation. For $W = w^0 + i\vec w\cdot\vec\sigma \in SU(2)$ one has
$W^{-1} = W^* = w^0 - i\vec w\cdot\vec\sigma$, because the $\sigma^b$ are hermitian; so
$\operatorname{tr} W^{-1} = w^0 = \operatorname{tr} W$, and this number is real. Therefore
\[
\operatorname{Re}\operatorname{tr} U(\partial p) = \operatorname{tr} W_{\ell,p} = w^0_{\ell,p},
\qquad\text{that is,}\qquad a_p = 1 - w^0_{\ell,p}.
\]
From $(w^0_{\ell,p})^2 + |\vec w_{\ell,p}|^2 = 1$ we get
\[
|\vec w_{\ell,p}|^2 = 1 - (1 - a_p)^2 = a_p(2 - a_p).
\]
Since $|w^0_{\ell,p}| \le 1$, we have $a_p \ge 0$, and so $a_p(2 - a_p) \le 2a_p$.

We now prove \eqref{10.1:eq-link-derivations-b}. The link $\ell$ occurs exactly once in $\partial p$. Its
variable $U_\ell$ is the leftmost factor of $W_{\ell,p}$, and the other three factors do not depend on
$U_\ell$. Replacing $U_\ell$ by $e^{is\sigma^b}U_\ell$ therefore replaces $W_{\ell,p}$ by
$e^{is\sigma^b}W_{\ell,p}$. Differentiating at $s = 0$ gives
\[
\frac{d}{ds}\, e^{is\sigma^b} W_{\ell,p}\Big|_{s=0} = i\sigma^b W_{\ell,p}.
\]
Using $\sigma^b\sigma^c = \delta^{bc}\mathbf{1} + i\varepsilon^{bcd}\sigma^d$, we compute
\begin{align*}
(i\sigma^b)(w^0 + i w^c\sigma^c)
&= i w^0\sigma^b - w^c\sigma^b\sigma^c
 = i w^0\sigma^b - w^b - i\varepsilon^{bcd} w^c \sigma^d\\
&= -w^b + i\bigl(w^0\delta^{bd} - \varepsilon^{bcd} w^c\bigr)\sigma^d .
\end{align*}

The curve $s \mapsto e^{is\sigma^b}W_{\ell,p}$ lies in $SU(2)$. Its components are therefore real, and its
derivative is $\partial^b_\ell w^0_{\ell,p} + i(\partial^b_\ell w^d_{\ell,p})\sigma^d$. The matrices
$\mathbf{1}, i\sigma^1, i\sigma^2, i\sigma^3$ are linearly independent over $\mathbb{R}$. Comparing
coefficients with the expression just computed (with $w = w_{\ell,p}$) gives the first line of
\eqref{10.1:eq-link-derivations-b}.

By \eqref{10.1:eq-link-derivations-a}, $a_p = 1 - w^0_{\ell,p}$, so
$\partial^b_\ell a_p = -\partial^b_\ell w^0_{\ell,p} = w^b_{\ell,p}$. Finally, set $d = b$ and sum over $b$.
Since $\sum_b \delta^{bb} = 3$ and $\varepsilon^{bcb} = 0$,
\[
\sum_{b=1}^{3} \partial^b_\ell w^b_{\ell,p} = 3w^0_{\ell,p} = 3(1 - a_p),
\]
where the last equality is \eqref{10.1:eq-link-derivations-a}.
\end{proof}

\begin{lemma}[Integration by parts for the sourced measure]\label{10.1:lem-integration-by-parts}
Let
\[
V=\sum_{\ell,b} v^b_\ell(U)\,\partial^b_\ell ,
\]
where the sum runs over the links $\ell$ of $\mathfrak{L}$ and $b=1,2,3$, and the coefficients
$v^b_\ell$ are smooth functions of $U$. Put $\operatorname{div}V:=\sum_{\ell,b}\partial^b_\ell v^b_\ell$.
Then for every complex $w$ with $Z[w]\neq 0$,
\begin{equation}\label{10.1:eq-integration-by-parts-1}
\langle V S_w\rangle_w=\langle \operatorname{div}V\rangle_w .
\end{equation}
\end{lemma}

\begin{proof}
We use the bare lattice, the sourced action $S_w$, the partition function $Z[w]$ and the
expectation $\langle\cdot\rangle_w=Z[w]^{-1}\int(\cdot)\,e^{-S_w}\,dU$ as fixed
in~\ref{10.1:not-bare-lattice}, and the link derivations $\partial^b_\ell$ as fixed
in~\ref{10.1:not-link-derivations}. Here $dU$ is the product of the Haar measures of the links.

Fix a link $\ell$ and $b\in\{1,2,3\}$, and put $F:=v^b_\ell\,e^{-S_w}$. The function $F$ is
smooth, because $v^b_\ell$ is smooth and $S_w=\sum_p c(p)a_p$ is a finite sum of the smooth
functions $a_p(U)=1-\operatorname{Re}\operatorname{tr}U(\partial p)$. The matrix $e^{is\sigma^b}$
lies in the group for every real $s$. Left multiplication of the variable $U_\ell$ by
$e^{is\sigma^b}$ therefore leaves the Haar measure of that factor invariant. It follows that
\[
\int F(\dots,e^{is\sigma^b}U_\ell,\dots)\,dU
\]
does not depend on $s$. The integrand is smooth in $(s,U)$ and the domain of integration is
compact, so we may differentiate under the integral sign at $s=0$. This gives
\[
\int \partial^b_\ell\bigl(v^b_\ell\,e^{-S_w}\bigr)\,dU=0 .
\]
The operator $\partial^b_\ell$ is a derivative along a curve, so it obeys the product rule:
\[
\partial^b_\ell\bigl(v^b_\ell\,e^{-S_w}\bigr)
=\bigl(\partial^b_\ell v^b_\ell\bigr)e^{-S_w}-v^b_\ell\,\bigl(\partial^b_\ell S_w\bigr)e^{-S_w}.
\]
Summing over the finitely many pairs $(\ell,b)$ and using
$VS_w=\sum_{\ell,b}v^b_\ell\,\partial^b_\ell S_w$, we obtain
\[
\int(\operatorname{div}V)\,e^{-S_w}\,dU=\int(VS_w)\,e^{-S_w}\,dU .
\]
Since $Z[w]\neq0$, we may divide by $Z[w]$, and this is \eqref{10.1:eq-integration-by-parts-1}.
\end{proof}

\begin{definition}[The derivation of an infinitesimal displacement]
\label{10.1:def-displacement-derivation}
We use the notation fixed in Notation~\ref{10.1:not-link-derivations} (quaternion components and
link derivations): $W=w^0+i\,\vec w\cdot\vec\sigma$ for the quaternion
components of $W\in SU(2)$, the plaquette holonomies $W_{\ell,p}$ and $W_{\ell,\pm\rho}$, and the
left-invariant derivations $\partial^b_\ell$, $b=1,2,3$.

Let $\eta\colon\mathbb{R}^4\to\mathbb{R}^4$ be a vector field with components $\eta^\rho$; it enters
only through its values
$\eta^\rho(x_\ell)$ at the link midpoints $x_\ell$. For a link $\ell=(x,\mu)$ of $\mathfrak{L}$ put
\begin{equation}\label{10.1:eq-displacement-derivation-a}
\begin{gathered}
\vec u_{\ell,\rho}:=-\tfrac12\bigl(\vec w_{\ell,+\rho}-\vec w_{\ell,-\rho}\bigr)\quad(\rho\neq\mu),
\qquad \vec u_{\ell,\mu}:=0,\\
\vec v_\ell:=a^{-1}\sum_{\rho}\eta^\rho(x_\ell)\,\vec u_{\ell,\rho},
\qquad
V_\eta:=\sum_{\ell,\,b}v^b_\ell\,\partial^b_\ell .
\end{gathered}
\end{equation}
Here $\vec w_{\ell,\pm\rho}$ is the vector part of $W_{\ell,\pm\rho}=W_{\ell,p(\ell,\pm\rho)}$, where
$p(\ell,\pm\rho)$ is the plaquette spanned by $\ell$ and $\pm\hat\rho$. The sum over $\rho$ runs over the four
lattice directions. The sum defining $V_\eta$ runs over all links $\ell$ of $\mathfrak{L}$ and over
$b=1,2,3$. We call $V_\eta$ the \emph{derivation of the displacement} $\eta$.

For a plaquette $p$, a link $\ell$ and a direction $\rho$, put
\begin{equation}\label{10.1:eq-displacement-derivation-b}
\mathfrak{d}_{\ell,p,\rho}:=\vec u_{\ell,\rho}\cdot\vec w_{\ell,p}\qquad(\ell\subset\partial p),
\end{equation}
where the dot is the Euclidean inner product on $\mathbb{R}^3$. We call $\mathfrak{d}_{\ell,p,\rho}$ the
\emph{elementary density}. By \eqref{10.1:eq-displacement-derivation-a}, for $\rho\neq\mu$ it is
$-\tfrac12$ times the difference of the products $\vec w_{\ell,+\rho}\cdot\vec w_{\ell,p}$ and
$\vec w_{\ell,-\rho}\cdot\vec w_{\ell,p}$, each a product of the imaginary parts, that is the vector
components $\vec w$, of the holonomies of two plaquettes that share the link
$\ell$. For $\rho=\mu$, the direction of $\ell$, it vanishes, since $\vec u_{\ell,\mu}=0$.
\end{definition}

In the classical limit the definition reads as follows.
Suppose the configuration is obtained by parallel transport from a smooth $\mathfrak{su}(2)$-valued
connection $A=A_\mu\,dx^\mu$ on $\mathbb{R}^4$ (not the action $A(U)$),
that is, $U_\ell$ is the path-ordered exponential $\operatorname{P}\exp\int_\ell A$ along $\ell$. Both
holonomies $W_{\ell,\pm\rho}$ are based at the initial point $x$ of $\ell$, and
\[
i\,\vec u_{\ell,\rho}\cdot\vec\sigma=a^2F_{\rho\mu}(x_\ell)+O(a^4),
\]
where $F_{\rho\mu}(x_\ell)$ is understood as parallel-transported along $\ell$ from $x_\ell$ to $x$, so
that both sides are compared at $x$.
Consequently $V_\eta$ transcribes the variation $\delta A_\mu=\eta^\rho F_{\rho\mu}$. This variation is the
Lie derivative of $A$ along $\eta$ up to a gauge transformation.

The link-centred average over $\pm\rho$ in \eqref{10.1:eq-displacement-derivation-a} is what makes every
remainder in what follows even in $a$.

\begin{lemma}[The displacement derivation: covariance and divergence]
\label{10.1:lem-displacement-divergence}
Let $\eta\colon\mathbb{R}^4\to\mathbb{R}^4$, and let $V_\eta=\sum_{\ell,b}v^b_\ell\,\partial^b_\ell$ be the
derivation of Definition~\ref{10.1:def-displacement-derivation}, the sum running over the links
$\ell=(x,\mu)$ of $\mathfrak{L}$; recall that $\vec u_{\ell,\mu(\ell)}=0$, so that in every sum
$\sum_\rho$ paired with a link $\ell$ only $\rho\neq\mu(\ell)$ contributes.
\begin{enumerate}
\item[(i)] $V_\eta$ commutes with gauge transformations: for every gauge transformation $g$ and every
smooth function $G$ of the configuration, $V_\eta$ maps the function $U\mapsto G(U^g)$ to the function
$U\mapsto(V_\eta G)(U^g)$, where
$(U^g)_\ell:=g(x)U_\ell\,g(x+\hat\mu)^{-1}$ for $\ell=(x,\mu)$. Moreover $V_\eta$ does not depend on the
orientation chosen for the links, and it is covariant under every symmetry $r$ of the periodic
lattice:
\begin{equation}\label{10.1:eq-displacement-divergence-1}
V_{r\cdot\eta}(G\circ r^{-1})=(V_\eta G)\circ r^{-1},\qquad (r\cdot\eta)(x):=r_*\eta(r^{-1}x),
\end{equation}
where $r_*$ is the linear part of $r$, $r$ acts on configurations by $(rU)_{r\ell}:=U_\ell$ for every
oriented link $\ell$, and $(G\circ r^{-1})(U):=G(r^{-1}U)$.
\item[(ii)] For every plaquette $p$,
\begin{equation}\label{10.1:eq-displacement-divergence-2}
V_\eta a_p=a^{-1}\sum_{\ell\subset\partial p}\sum_\rho\eta^\rho(x_\ell)\,\mathfrak{d}_{\ell,p,\rho}.
\end{equation}
\item[(iii)] The divergence $\operatorname{div}V_\eta:=\sum_{\ell,b}\partial^b_\ell v^b_\ell$ satisfies
\begin{equation}\label{10.1:eq-displacement-divergence-3}
\begin{split}
\operatorname{div}V_\eta
&=\frac{3}{2a}\sum_\ell\sum_{\rho\neq\mu(\ell)}\eta^\rho(x_\ell)
\bigl[a_{p(\ell,\rho)}-a_{p(\ell,-\rho)}\bigr]\\
&=-\frac32\sum_{p=\{x;\mu,\rho\}}a_p\bigl[(\nabla^c_\rho\eta^\rho)(x_p)
+(\nabla^c_\mu\eta^\mu)(x_p)\bigr],
\end{split}
\end{equation}
with no summation over $\mu,\rho$ inside the bracket; here $\{x;\mu,\rho\}$ denotes the plaquette
spanned at its corner $x$ by $\hat\mu$ and $\hat\rho$, the last sum counts each plaquette once, and
\begin{equation*}
(\nabla^c_\rho\varphi)(y):=a^{-1}\bigl[\varphi(y+\tfrac12\hat\rho)-\varphi(y-\tfrac12\hat\rho)\bigr].
\end{equation*}
\item[(iv)] On $SU(2)$-valued fields, for $\ell\subset\partial p$ and $\rho\neq\mu(\ell)$,
\begin{equation}\label{10.1:eq-displacement-divergence-4}
|\mathfrak{d}_{\ell,p,\rho}|\le\tfrac12\bigl(a_{p(\ell,\rho)}+a_{p(\ell,-\rho)}\bigr)+a_p.
\end{equation}
\end{enumerate}
\end{lemma}

\begin{proof}
(i) Gauge transformations. For $h\in SU(2)$ let $\operatorname{Ad}h\in SO(3)$ be the matrix defined by
$h\sigma^bh^{-1}=\sum_d(\operatorname{Ad}h)_{db}\sigma^d$. For $\ell=(x,\mu)$ every holonomy
$W_{\ell,\cdot}$ is the holonomy of a closed loop based at $x$, so under $U\mapsto U^g$ it becomes
$g(x)W_{\ell,\cdot}\,g(x)^{-1}$.
Since $g(x)(w^0+i\vec w\cdot\vec\sigma)g(x)^{-1}=w^0+i(\operatorname{Ad}g(x)\vec w)\cdot\vec\sigma$, all
$\vec w_{\ell,\cdot}$ rotate by $\operatorname{Ad}g(x)$; as $\vec u_{\ell,\rho}$ and $\vec v_\ell$ are
linear in them with coefficients independent of $U$, $\vec v_\ell(U^g)=\operatorname{Ad}g(x)\,
\vec v_\ell(U)$. The frame rotates in the same way. Put $G^g(U):=G(U^g)$. The map $U\mapsto U^g$
sends $e^{is\sigma^b}U_\ell$ to
\[
g(x)e^{is\sigma^b}U_\ell\,g(x+\hat\mu)^{-1}
=\exp\Bigl(is\sum_d(\operatorname{Ad}g(x))_{db}\sigma^d\Bigr)(U^g)_\ell,
\]
and differentiating at $s=0$ gives
\[
(\partial^b_\ell G^g)(U)=\sum_d(\operatorname{Ad}g(x))_{db}(\partial^d_\ell G)(U^g).
\]
Hence
\[
(V_\eta G^g)(U)
=\sum_{\ell,d}\Bigl(\sum_b(\operatorname{Ad}g(x))_{db}v^b_\ell(U)\Bigr)(\partial^d_\ell G)(U^g)
=\sum_{\ell,d}v^d_\ell(U^g)(\partial^d_\ell G)(U^g)=(V_\eta G)(U^g).
\]

Orientation. The reversed link $-\ell$ has initial point $x+\hat\mu$ and $U_{-\ell}=U_\ell^{-1}$, and it
spans with $\pm\hat\rho$ the same plaquettes $p(\ell,\pm\rho)$. The product $U_\ell^{-1}W_{\ell,\rho}$ is
the holonomy from $x+\hat\mu$ back to $x$ along the other three links of $p(\ell,\rho)$, so the loop
starting at $x+\hat\mu$ and running through $-\ell$ first has holonomy
$U_\ell^{-1}\bigl(U_\ell^{-1}W_{\ell,\rho}\bigr)^{-1}=U_\ell^{-1}W_{\ell,\rho}^{-1}U_\ell$; likewise for
$-\rho$. As
$W^{-1}=w^0-i\vec w\cdot\vec\sigma$, this gives $\vec w_{-\ell,\pm\rho}=-\operatorname{Ad}(U_\ell^{-1})
\vec w_{\ell,\pm\rho}$, hence $\vec u_{-\ell,\rho}=-\operatorname{Ad}(U_\ell^{-1})\vec u_{\ell,\rho}$ and,
since $x_{-\ell}=x_\ell$, $\vec v_{-\ell}=-\operatorname{Ad}(U_\ell^{-1})\vec v_\ell$. This is how the
coefficient of a vector field transforms under $U_\ell\mapsto U_\ell^{-1}$: the curve
$e^{is\sigma^b}U_{-\ell}$ corresponds to
$U_\ell e^{-is\sigma^b}=\exp\bigl(-is\sum_d(\operatorname{Ad}U_\ell)_{db}\sigma^d\bigr)U_\ell$, so
$\partial^b_{-\ell}=-\sum_d(\operatorname{Ad}U_\ell)_{db}\partial^d_\ell$, and therefore
\[
\sum_bv^b_{-\ell}\partial^b_{-\ell}
=\sum_d\bigl(\operatorname{Ad}U_\ell\operatorname{Ad}U_\ell^{-1}\vec v_\ell\bigr)^d\partial^d_\ell
=\sum_dv^d_\ell\partial^d_\ell.
\]

Lattice symmetries. By its definition $\vec u_{\ell,\rho}$ is odd under $\hat\rho\mapsto-\hat\rho$
(this exchanges $\vec w_{\ell,+\rho}$ and $\vec w_{\ell,-\rho}$), i.e.\ $\rho$ is a vector index:
setting $\vec u_{\ell,-\rho}:=-\vec u_{\ell,\rho}$, the linear map $e\mapsto\vec u_\ell(e)$ on
$\mathbb{R}^4$ with $\vec u_\ell(\pm e_\rho)=\vec u_{\ell,\pm\rho}$ satisfies
$\vec v_\ell=a^{-1}\vec u_\ell(\eta(x_\ell))$.
The rest is relabelling. Since $W_{\ell,p}$ is determined by the oriented link $\ell$ and the plaquette
$p$, $W_{r\ell,r p}(rU)=W_{\ell,p}(U)$, so $\vec u_{r\ell}(rU)(r_*e)=\vec u_\ell(U)(e)$; with
$x_{r\ell}=r x_\ell$ and $(r\cdot\eta)(rx_\ell)=r_*\eta(x_\ell)$ this gives, writing $\vec v_\ell[\eta]$
(with components $v^b_\ell[\eta]$) for the coefficient of
Definition~\ref{10.1:def-displacement-derivation} formed with the displacement $\eta$,
$\vec v_{r\ell}[r\cdot\eta](rU)=\vec v_\ell[\eta](U)$. Moreover
$\partial^b_{r\ell}(G\circ r^{-1})(rU)=(\partial^b_\ell G)(U)$ for every oriented link $\ell$, since
$(rU)_{r\ell}=U_\ell$. Summing over $\ell$, i.e.\ over the oriented links $r\ell$, and using the
orientation independence just proved for those $r\ell$ that are negatively oriented,
\[
V_{r\cdot\eta}(G\circ r^{-1})(rU)=\sum_{\ell,b}v^b_\ell[\eta](U)(\partial^b_\ell G)(U)
=\bigl((V_\eta G)\circ r^{-1}\bigr)(rU),
\]
which is \eqref{10.1:eq-displacement-divergence-1}.

(ii) The function $a_p$ depends only on the links of $\partial p$, so $\partial^b_\ell a_p=0$ for
$\ell\not\subset\partial p$, while for $\ell\subset\partial p$ the identity
\eqref{10.1:eq-link-derivations-b} gives
$\partial^b_\ell a_p=w^b_{\ell,p}$. Hence
\[
V_\eta a_p=\sum_{\ell\subset\partial p}\vec v_\ell\cdot\vec w_{\ell,p}
=a^{-1}\sum_{\ell\subset\partial p}\sum_\rho\eta^\rho(x_\ell)\,\vec u_{\ell,\rho}\cdot\vec w_{\ell,p},
\]
which is \eqref{10.1:eq-displacement-divergence-2} by \eqref{10.1:eq-displacement-derivation-b}.

(iii) The coefficient $\vec v_\ell$ depends on $U_\ell$ only through $\vec w_{\ell,\pm\rho}$, its
coefficients $a^{-1}\eta^\rho(x_\ell)$ being independent of $U$. By
\eqref{10.1:eq-link-derivations-b},
\[
\sum_b\partial^b_\ell w^b_{\ell,\pm\rho}=3w^0_{\ell,\pm\rho}=3(1-a_{p(\ell,\pm\rho)}),
\]
so for $\rho\neq\mu(\ell)$
\[
\sum_b\partial^b_\ell u^b_{\ell,\rho}=-\tfrac32\bigl(w^0_{\ell,\rho}-w^0_{\ell,-\rho}\bigr)
=\tfrac32\bigl(a_{p(\ell,\rho)}-a_{p(\ell,-\rho)}\bigr).
\]
Multiplying by $a^{-1}\eta^\rho(x_\ell)$ and summing over $\rho$ and $\ell$ gives the first equality
in \eqref{10.1:eq-displacement-divergence-3}. For the second, with $\ell=(x,\mu)$ one has
$x_\ell=x+\frac12\hat\mu$, $p(\ell,\rho)=\{x;\mu,\rho\}$ and $p(\ell,-\rho)=\{x-\hat\rho;\mu,\rho\}$.
In the sum of the terms with $a_{p(\ell,-\rho)}$ we shift $x\mapsto x+\hat\rho$, which permutes the
points of the periodic lattice; this sum becomes
\[
\sum_x\sum_{\mu\neq\rho}\eta^\rho(x+\tfrac12\hat\mu+\hat\rho)\,a_{\{x;\mu,\rho\}}.
\]
Since
$x_p=x+\frac12\hat\mu+\frac12\hat\rho$ for $p=\{x;\mu,\rho\}$, the first line of
\eqref{10.1:eq-displacement-divergence-3} equals
\[
\frac{3}{2a}\sum_x\sum_{\mu\neq\rho}a_{\{x;\mu,\rho\}}\bigl[\eta^\rho(x_p-\tfrac12\hat\rho)
-\eta^\rho(x_p+\tfrac12\hat\rho)\bigr]
=-\frac32\sum_x\sum_{\mu\neq\rho}a_{\{x;\mu,\rho\}}(\nabla^c_\rho\eta^\rho)(x_p).
\]
Each plaquette $\{x;\mu,\rho\}$ occurs here for the ordered pairs $(\mu,\rho)$ and $(\rho,\mu)$,
contributing $(\nabla^c_\rho\eta^\rho)(x_p)$ and $(\nabla^c_\mu\eta^\mu)(x_p)$; collecting the two gives
the second line of \eqref{10.1:eq-displacement-divergence-3}.

(iv) Write $\vec w_\pm:=\vec w_{\ell,\pm\rho}$ and $\vec w_p:=\vec w_{\ell,p}$. By the Cauchy--Schwarz and
triangle inequalities, and by the inequality between the geometric and arithmetic means of two
non-negative numbers,
\[
|\vec u_{\ell,\rho}\cdot\vec w_p|\le\tfrac12\bigl(|\vec w_+|+|\vec w_-|\bigr)|\vec w_p|
\le\tfrac14\bigl(|\vec w_+|^2+|\vec w_-|^2\bigr)+\tfrac12|\vec w_p|^2.
\]
By \eqref{10.1:eq-link-derivations-a}, $|\vec w_\pm|^2\le2a_{p(\ell,\pm\rho)}$ and
$|\vec w_p|^2\le2a_p$, which gives \eqref{10.1:eq-displacement-divergence-4}.
\end{proof}

\begin{definition}[Orbital, spin and remainder densities]\label{10.1:def-orbital-spin-densities}
Let $p$ be a plaquette of $\mathfrak{L}$, let $\rho,\sigma\in\{1,\dots,4\}$, and let $\chi$ be a function.
With the elementary densities $\mathfrak{d}_{\ell,p,\rho}$ of
Definition~\ref{10.1:def-displacement-derivation}, and with the convention that in a sum over
$\ell\subset\partial p$ only the links of $\partial p$ orthogonal to $\hat\rho$ contribute (for $\ell$
parallel to $\hat\rho$ the plaquettes $p(\ell,\pm\rho)$ spanned by $\ell$ and $\pm\hat\rho$, on which
$\mathfrak{d}_{\ell,p,\rho}$ depends, do not exist, and the term is absent), set
\begin{equation}\label{10.1:eq-orbital-spin-densities-a}
\begin{split}
\mathfrak{t}_\rho(p) &:= \sum_{\ell\subset\partial p}\mathfrak{d}_{\ell,p,\rho}, \qquad
\mathfrak{s}_{\sigma\rho}(p) := \sum_{\ell\subset\partial p}a^{-1}(x_\ell-x_p)_\sigma\,
\mathfrak{d}_{\ell,p,\rho},\\
\mathfrak{r}_\rho(p) &:= \mathfrak{t}_\rho(p)-\tfrac12\bigl(a_{p+\hat\rho}-a_{p-\hat\rho}\bigr),
\end{split}
\end{equation}
where $p\pm\hat\rho$ is the translate of $p$ by $\pm\hat\rho$. With $\xi(x)=\omega x$ the rotation field,
set further
\begin{equation}\label{10.1:eq-orbital-spin-densities-b}
\begin{split}
\mathfrak{h}_\chi(p) &:= a^{-1}\sum_{\ell\subset\partial p}\sum_\rho
\bigl(\chi(x_\ell)-\chi(x_p)\bigr)\,\xi^\rho(x_\ell)\,\mathfrak{d}_{\ell,p,\rho},\\
(\nabla^s_\rho\varphi)(y) &:= (2a)^{-1}\bigl[\varphi(y+\hat\rho)-\varphi(y-\hat\rho)\bigr], \qquad
\mathsf{L}^a_\xi\varphi := \sum_\rho\xi^\rho\,\nabla^s_\rho\varphi .
\end{split}
\end{equation}
We call $\mathfrak{t}_\rho$ the orbital density, $\mathfrak{s}_{\sigma\rho}$ the spin density,
$\mathfrak{r}_\rho$ the orbital remainder density and $\mathfrak{h}_\chi$ the cutoff-variation density.
For $p=\{x;\mu,\nu\}$ one has $a^{-1}(x_\ell-x_p)=\mp\frac12 e_\nu$ on the two $\mu$-links of $p$ (the sign
$-$ for the link through $x$) and $\mp\frac12 e_\mu$ on the two $\nu$-links. By
Lemma~\ref{10.1:lem-displacement-divergence}(ii) with the constant field $\eta=e_\rho$,
$\mathfrak{t}_\rho(p)=a\,V_{e_\rho}a_p$ is the lattice's infinitesimal translation of $a_p$, and
$\mathfrak{r}_\rho(p)$ compares it with the symmetric lattice difference. Let $N(p)$ denote the set of
plaquettes sharing a link with $p$, together with $p$ and $p\pm\hat\rho$. The densities
$\mathfrak{t}_\rho(p)$, $\mathfrak{s}_{\sigma\rho}(p)$, $\mathfrak{r}_\rho(p)$ and $\mathfrak{h}_\chi(p)$
are polynomials in the links within sup-distance $2a$ of $x_p$ (with coefficients depending on $\chi$
and $\xi$ for $\mathfrak{h}_\chi$), are gauge invariant, and:
\begin{enumerate}
\item[(a)] on $SU(2)$-valued fields,
\begin{equation}\label{10.1:eq-orbital-spin-densities-c}
|\mathfrak{s}_{\sigma\rho}(p)|,\ |\mathfrak{t}_\rho(p)|,\ |\mathfrak{r}_\rho(p)|
\le 4\sum_{p'\in N(p)}a_{p'} ;
\end{equation}
\item[(b)] for every symmetry $r$ of the periodic lattice with linear part $r_*$, writing $G\circ r$ for a
function $G$ of the field evaluated on the transported field,
\begin{equation}\label{10.1:eq-orbital-spin-densities-1}
\mathfrak{s}_{\sigma\rho}(rp)\circ r=\sum_{\sigma',\rho'}(r_*)_{\sigma\sigma'}(r_*)_{\rho\rho'}\,
\mathfrak{s}_{\sigma'\rho'}(p), \qquad
\mathfrak{r}_\rho(rp)\circ r=\sum_{\rho'}(r_*)_{\rho\rho'}\,\mathfrak{r}_{\rho'}(p);
\end{equation}
that is, $\mathfrak{s}$ is a rank-two tensor density and $\mathfrak{r}$ a vector density.
\end{enumerate}
\end{definition}

\begin{proof}
Each $\mathfrak{d}_{\ell,p,\rho}$ is a polynomial in the holonomies of $p$ and of $p(\ell,\pm\rho)$; these
plaquettes and $p\pm\hat\rho$ lie within sup-distance $\frac32 a\le 2a$ of $x_p$. Since
$1-\operatorname{Re}\operatorname{tr}U\ge 0$ for $U\in SU(2)$, the trace being normalised, every
$a_{p'}$ is nonnegative on $SU(2)$-valued fields, and $a_p$ does not depend on the orientation of $p$
because $\operatorname{Re}\operatorname{tr}U^{-1}=\operatorname{Re}\operatorname{tr}U$.
By Lemma~\ref{10.1:lem-displacement-divergence}(ii), $V_\eta a_p$ is linear in $\eta$ and depends on
$\eta$ only through its values at the four midpoints $x_\ell$, $\ell\subset\partial p$. Let
$\eta^{(\sigma\rho)}_p$ be any displacement field (in the sense of
Definition~\ref{10.1:def-displacement-derivation}) whose values at these four midpoints are
\[
\eta^{(\sigma\rho)}_p(x_\ell)=a^{-1}(x_\ell-x_p)_\sigma e_\rho\qquad(\ell\subset\partial p);
\]
then that formula gives $V_{\eta^{(\sigma\rho)}_p}a_p=a^{-1}\mathfrak{s}_{\sigma\rho}(p)$, and likewise
$\mathfrak{h}_\chi(p)=V_\eta a_p$ for any displacement field $\eta$ whose values at the midpoints are
$\eta(x_\ell)=(\chi(x_\ell)-\chi(x_p))\,\xi(x_\ell)$ for $\ell\subset\partial p$. Gauge invariance of
$\mathfrak{t}_\rho$, $\mathfrak{s}_{\sigma\rho}$, $\mathfrak{h}_\chi$ and hence $\mathfrak{r}_\rho$ follows
from Lemma~\ref{10.1:lem-displacement-divergence}(i), since $a_p$ and $a_{p\pm\hat\rho}$ are gauge
invariant.

(a) Let $p=\{x;\mu,\nu\}$, and recall that for a link $\ell=(y,\alpha)$ and $\rho\ne\alpha$ one has
$p(\ell,\rho)=\{y;\alpha,\rho\}$ and $p(\ell,-\rho)=\{y-\hat\rho;\alpha,\rho\}$. For a link
$\ell\subset\partial p$
orthogonal to $\hat\rho$, Lemma~\ref{10.1:lem-displacement-divergence}(iv) gives
$|\mathfrak{d}_{\ell,p,\rho}|\le a_p+\frac12\bigl(a_{p(\ell,\rho)}+a_{p(\ell,-\rho)}\bigr)$; links parallel
to $\hat\rho$ contribute no term. By the triangle inequality,
\begin{equation}\label{10.1:eq-orbital-spin-densities-2}
|\mathfrak{t}_\rho(p)|\le\sum_{\ell\subset\partial p}|\mathfrak{d}_{\ell,p,\rho}|
\le\sum_{\substack{\ell\subset\partial p\\ \ell\perp\hat\rho}}
\Bigl[a_p+\tfrac12\bigl(a_{p(\ell,\rho)}+a_{p(\ell,-\rho)}\bigr)\Bigr].
\end{equation}
Each $p(\ell,\pm\rho)$ contains the link $\ell$ of $\partial p$ and hence lies in $N(p)$. If
$\rho\notin\{\mu,\nu\}$, all four links of $\partial p$ are orthogonal to $\hat\rho$, and the right-hand
side of \eqref{10.1:eq-orbital-spin-densities-2} is
$4a_p+\frac12\sum_{\ell\subset\partial p}\bigl(a_{p(\ell,\rho)}+a_{p(\ell,-\rho)}\bigr)$; the eight
plaquettes $p(\ell,\pm\rho)$ lie in planes containing $\hat\rho$, so they differ from $p$ and $p\pm\hat\rho$,
and their coefficients add up to $4$. Hence every $a_{p'}$, $p'\in N(p)$, carries coefficient at most $4$,
and $a_{p\pm\hat\rho}$ does not occur. If $\rho\in\{\mu,\nu\}$, say $\rho=\nu$ (the case $\rho=\mu$ is the
same with $\mu$ and $\nu$ exchanged), only the two $\mu$-links $\ell_1=(x,\mu)$ and
$\ell_2=(x+\hat\nu,\mu)$ contribute, with $p(\ell_1,\nu)=p$, $p(\ell_1,-\nu)=p-\hat\nu$,
$p(\ell_2,\nu)=p+\hat\nu$ and $p(\ell_2,-\nu)=p$, so that
\eqref{10.1:eq-orbital-spin-densities-2} reads
\begin{equation}\label{10.1:eq-orbital-spin-densities-3}
|\mathfrak{t}_\rho(p)|\le 3a_p+\tfrac12\bigl(a_{p+\hat\rho}+a_{p-\hat\rho}\bigr).
\end{equation}
Since all terms are nonnegative, in both cases $|\mathfrak{t}_\rho(p)|\le 4\sum_{p'\in N(p)}a_{p'}$. Since
$|a^{-1}(x_\ell-x_p)_\sigma|\le\frac12$ for every $\ell\subset\partial p$, we get
$|\mathfrak{s}_{\sigma\rho}(p)|\le\frac12\sum_{\ell\subset\partial p}|\mathfrak{d}_{\ell,p,\rho}|$, which is
bounded in the same way. Finally
$|\mathfrak{r}_\rho(p)|\le|\mathfrak{t}_\rho(p)|+\frac12 a_{p+\hat\rho}+\frac12 a_{p-\hat\rho}$. We insert
\eqref{10.1:eq-orbital-spin-densities-2}, not \eqref{10.1:eq-orbital-spin-densities-c}, for
$|\mathfrak{t}_\rho(p)|$: if $\rho\notin\{\mu,\nu\}$, the plaquettes $p\pm\hat\rho$ acquire coefficient
$\frac12$ and every other $p'\in N(p)$ keeps coefficient at most $4$; if $\rho\in\{\mu,\nu\}$,
\eqref{10.1:eq-orbital-spin-densities-3} gives coefficient $3$ to $a_p$ and coefficient $1$ to
$a_{p\pm\hat\rho}$. Since $p\pm\hat\rho\in N(p)$, this proves
\eqref{10.1:eq-orbital-spin-densities-c} for $\mathfrak{r}_\rho$.

(b) Write $ry=r_*y+b$, with $r_*$ a signed permutation matrix and $b$ a fixed vector. Then
$x_{rp}=r x_p$, and $r$ maps the midpoints of the links of $\partial p$ onto those of $\partial(rp)$. For
$y=x_{\ell'}$ with $\ell'\subset\partial(rp)$ we have $r^{-1}y-x_p=r_*^{-1}(y-x_{rp})$, and orthogonality
of $r_*$ gives
\[
(r\cdot\eta^{(\sigma\rho)}_p)(y)=r_*\eta^{(\sigma\rho)}_p(r^{-1}y)
=\sum_{\sigma',\rho'}(r_*)_{\sigma'\sigma}(r_*)_{\rho'\rho}\,\eta^{(\sigma'\rho')}_{rp}(y).
\]
Since $a_p\circ r^{-1}=a_{rp}$, Lemma~\ref{10.1:lem-displacement-divergence}(i) gives the first
equality below; the second follows because $V_\eta a_{rp}$ is linear in $\eta$ and depends on $\eta$ only
through its values at the midpoints of the links of $\partial(rp)$
(Lemma~\ref{10.1:lem-displacement-divergence}(ii)):
\[
\mathfrak{s}_{\sigma\rho}(p)\circ r^{-1}=a\,V_{r\cdot\eta^{(\sigma\rho)}_p}a_{rp}
=\sum_{\sigma',\rho'}(r_*)_{\sigma'\sigma}(r_*)_{\rho'\rho}\,\mathfrak{s}_{\sigma'\rho'}(rp).
\]
Multiplying by $(r_*)_{\bar\sigma\sigma}(r_*)_{\bar\rho\rho}$, summing over $\sigma,\rho$, using the
orthogonality of $r_*$ and composing with $r$ gives the first identity of
\eqref{10.1:eq-orbital-spin-densities-1}. In the same way, $r\cdot e_\rho=\sum_{\rho'}(r_*)_{\rho'\rho}
e_{\rho'}$ gives $\mathfrak{t}_\rho(p)\circ r^{-1}=\sum_{\rho'}(r_*)_{\rho'\rho}\mathfrak{t}_{\rho'}(rp)$.
Moreover $(a_{p+\hat\rho}-a_{p-\hat\rho})\circ r^{-1}=a_{rp+r_*\hat\rho}-a_{rp-r_*\hat\rho}$. If
$r_*e_\rho=s\,e_{\rho_0}$ with $s=\pm1$, this gives
\[
(a_{p+\hat\rho}-a_{p-\hat\rho})\circ r^{-1}
=s\,(a_{rp+\hat\rho_0}-a_{rp-\hat\rho_0})
=\sum_{\rho'}(r_*)_{\rho'\rho}\,(a_{rp+\hat\rho'}-a_{rp-\hat\rho'}).
\]
Hence
$\mathfrak{r}_\rho(p)\circ r^{-1}=\sum_{\rho'}(r_*)_{\rho'\rho}\mathfrak{r}_{\rho'}(rp)$, and inverting as
before gives the second identity of \eqref{10.1:eq-orbital-spin-densities-1}.
\end{proof}

We use the bare lattice $\mathfrak{L}$, the spacing $a=L^{-K}$, the plaquette densities
$a_p(U)=1-\operatorname{Re}\operatorname{tr}U(\partial p)$, the sourced coefficient
$c(p)=g_0^{-2}-w(p)$, the sourced action $S_w$, the partition function $Z[w]$, the normalised
expectation $\langle\cdot\rangle_w$ and the cumulants $\langle Y;\mathcal{O}(f_1);\dots;\mathcal{O}(f_n)\rangle$
as fixed in Notation~\ref{10.1:not-bare-lattice}.
We also use the densities $\mathfrak{t}_\rho$, $\mathfrak{s}_{\sigma\rho}$, $\mathfrak{r}_\rho$, $\mathfrak{h}_\chi$
and the operators $\nabla^s_\rho$, $\mathsf{L}^a_\xi$ of Definition~\ref{10.1:def-orbital-spin-densities},
and the half-step difference $\nabla^c_\rho$ of Lemma~\ref{10.1:lem-displacement-divergence}.

\begin{theorem}[The rotational Ward identity]\label{10.1:thm-rotational-identity}
We call the following data \emph{admissible}:
\begin{enumerate}
\item[(a)] a real antisymmetric matrix $\omega=(\omega_{\rho\sigma})$, with $\xi(x)=\omega x$, that is,
$\xi^\rho(x)=\sum_\sigma\omega_{\rho\sigma}x_\sigma$;
\item[(b)] a function $\chi\colon\mathbb{R}^4\to[0,1]$ with
$\operatorname{supp}\chi\subset\,]-L^m+2a,\,L^m-2a[^4$, so that $\chi\xi$ is a function on the torus;
\item[(c)] a complex function $w$ on the plaquettes with $Z[w]\neq0$, such that $\chi\equiv1$ on the closed
$2a$-neighbourhood of $\{x_p: w(p)\neq0\}$.
\end{enumerate}

A function $f$ on the torus is regarded as the function $p\mapsto f(x_p)$ on plaquettes. For a function $h$
on plaquettes, $\mathcal{O}(h):=\sum_p h(p)a_p$, and $\mathsf{L}^a_\xi h$ is given by
\eqref{10.1:eq-orbital-spin-densities-b} with $\varphi(x_p)=h(p)$, that is,
\[
(\mathsf{L}^a_\xi h)(p)=\sum_\rho\xi^\rho(x_p)\,(2a)^{-1}\big[h(p+\hat\rho)-h(p-\hat\rho)\big].
\]
For such $h$ put
\begin{equation}\label{10.1:eq-rotational-identity-1}
\begin{aligned}
\mathfrak{S}_\chi[h]&:=\sum_p h(p)\chi(x_p)\sum_{\rho,\sigma}\omega_{\rho\sigma}\,\mathfrak{s}_{\sigma\rho}(p),\\
\mathfrak{K}_\chi[h]&:=\sum_p h(p)\chi(x_p)\sum_{\rho}a^{-1}\xi^\rho(x_p)\,\mathfrak{r}_{\rho}(p).
\end{aligned}
\end{equation}
Let $\mathfrak{S}[h]$ and $\mathfrak{K}[h]$ denote the same sums with the factor $\chi(x_p)$ omitted.
Further let
\begin{equation}\label{10.1:eq-rotational-identity-2}
\mathfrak{B}^{\mathrm{sh}}_\chi:=g_0^{-2}\sum_p a_p\,(\mathsf{L}^a_\xi\chi)(x_p)
-g_0^{-2}\sum_p\mathfrak{h}_\chi(p)+\operatorname{div}V_{\chi\xi},
\end{equation}
where $V_{\chi\xi}$ is the derivation of the displacement $\eta=\chi\xi$ of
Lemma~\ref{10.1:lem-displacement-divergence}.
The spin term $\mathfrak{S}_\chi[c]$ contracts the rank-two tensor density $\mathfrak{s}$ with $\omega$
against the bounded weight $c\chi$. The orbital remainder $\mathfrak{K}_\chi[c]$ contracts the vector
density $\mathfrak{r}$ against the weight $c\chi a^{-1}\xi$. The shell term $\mathfrak{B}^{\mathrm{sh}}_\chi$
involves $\chi$ only through differences of its values.

For admissible data,
\begin{equation}\label{10.1:eq-rotational-identity-a}
\big\langle\mathcal{O}(\mathsf{L}^a_\xi w)\big\rangle_w
=-\big\langle\mathfrak{S}_\chi[c]\big\rangle_w-\big\langle\mathfrak{K}_\chi[c]\big\rangle_w
+\big\langle\mathfrak{B}^{\mathrm{sh}}_\chi\big\rangle_w ,
\end{equation}
and
\begin{equation}\label{10.1:eq-rotational-identity-3}
\operatorname{div}V_{\chi\xi}=-\tfrac32\sum_{p=\{x;\mu,\rho\}}a_p
\big[\xi^\rho(x_p)(\nabla^c_\rho\chi)(x_p)+\xi^\mu(x_p)(\nabla^c_\mu\chi)(x_p)\big],
\end{equation}
where $\mu,\rho$ are the two directions of $p$ (no summation over them).

Consequently, let $n\ge1$, let $\omega$ and $\chi$ be as in (a), (b), and let $f_1,\dots,f_n$ be real (or
complex) functions such that $\chi\equiv1$ on the $2a$-neighbourhood of $\bigcup_i\operatorname{supp}f_i$.
Then
\begin{equation}\label{10.1:eq-rotational-identity-b}
\begin{split}
\sum_{i=1}^n S^T_{n;K,m}(f_1,\dots,\mathsf{L}^a_\xi f_i,\dots,f_n)
&=-g_0^{-2}\big\langle(\mathfrak{S}_\chi+\mathfrak{K}_\chi)[1];\mathcal{O}(f_1);\dots;\mathcal{O}(f_n)\big\rangle\\
&\quad+\sum_{i=1}^n\big\langle(\mathfrak{S}+\mathfrak{K})[f_i];(\mathcal{O}(f_l))_{l\neq i}\big\rangle\\
&\quad+\big\langle\mathfrak{B}^{\mathrm{sh}}_\chi;\mathcal{O}(f_1);\dots;\mathcal{O}(f_n)\big\rangle ,
\end{split}
\end{equation}
where $\langle Y;(\mathcal{O}(f_l))_{l\neq i}\rangle$ denotes the bracket with the $n-1$ insertions
$\mathcal{O}(f_l)$, $l\neq i$.
\end{theorem}

\begin{proof}
Every value of $\xi$ that occurs below is multiplied by a value of $\chi$, by a difference of values of
$\chi$, or by a value of $w$, at points within sup-distance $a$ of the point where $\xi$ is evaluated.
Where such a factor is non-zero, hypotheses (b) and (c) place all the points involved in the open cube
$]-L^m+a,L^m-a[^4$. We evaluate $\xi$ at these representatives, so that $\xi$ is the linear function
$x\mapsto\omega x$ there.

\emph{The Killing field.} The field $\xi$ is linear and $\omega_{\rho\rho}=0$. Hence, for every plaquette
$p$ and every link $\ell\subset\partial p$,
\begin{equation}\label{10.1:eq-rotational-identity-4}
\xi^\rho(x_\ell)=\xi^\rho(x_p)+\sum_\sigma\omega_{\rho\sigma}(x_\ell-x_p)_\sigma,
\qquad
\xi^\rho(x_p\pm\hat\rho)=\xi^\rho(x_p)\pm a\,\omega_{\rho\rho}=\xi^\rho(x_p).
\end{equation}
For the same reason $\xi^\rho(y\pm\tfrac12\hat\rho)=\xi^\rho(y)$ for every point $y$.

\emph{The derivation on one plaquette.} Apply part (ii) of Lemma~\ref{10.1:lem-displacement-divergence}
with $\eta=\chi\xi$, and write $\chi(x_\ell)=\chi(x_p)+(\chi(x_\ell)-\chi(x_p))$. The second summand gives
$\mathfrak{h}_\chi(p)$ by \eqref{10.1:eq-orbital-spin-densities-b}. In the first summand insert
\eqref{10.1:eq-rotational-identity-4}. The sum over $\ell\subset\partial p$ of $\mathfrak{d}_{\ell,p,\rho}$ is
$\mathfrak{t}_\rho(p)$, and the sum of $a^{-1}(x_\ell-x_p)_\sigma\mathfrak{d}_{\ell,p,\rho}$ is
$\mathfrak{s}_{\sigma\rho}(p)$, both by \eqref{10.1:eq-orbital-spin-densities-a}. Therefore
\begin{equation}\label{10.1:eq-rotational-identity-5}
V_{\chi\xi}a_p=\chi(x_p)\Big[\sum_\rho a^{-1}\xi^\rho(x_p)\,\mathfrak{t}_\rho(p)
+\sum_{\rho,\sigma}\omega_{\rho\sigma}\,\mathfrak{s}_{\sigma\rho}(p)\Big]+\mathfrak{h}_\chi(p).
\end{equation}

\emph{Summation over plaquettes.} By \eqref{10.1:eq-orbital-spin-densities-a},
$\mathfrak{t}_\rho=\tfrac12(a_{p+\hat\rho}-a_{p-\hat\rho})+\mathfrak{r}_\rho$. Insert this in
\eqref{10.1:eq-rotational-identity-5}, multiply by $c(p)$ and sum over $p$. The coefficient $c(p)$ does not
depend on $U$ and $V_{\chi\xi}$ is a derivation, so $V_{\chi\xi}S_w=\sum_pc(p)V_{\chi\xi}a_p$. With
\eqref{10.1:eq-rotational-identity-1} this gives
\begin{equation}\label{10.1:eq-rotational-identity-6}
\begin{split}
V_{\chi\xi}S_w&=\sum_\rho\sum_pc(p)\chi(x_p)\xi^\rho(x_p)(2a)^{-1}(a_{p+\hat\rho}-a_{p-\hat\rho})\\
&\quad+\mathfrak{K}_\chi[c]+\mathfrak{S}_\chi[c]+\sum_pc(p)\mathfrak{h}_\chi(p).
\end{split}
\end{equation}

For fixed $\rho$, the translations $p\mapsto p\pm\hat\rho$ are bijections of the set of plaquettes of the
periodic lattice, and $x_{p\pm\hat\rho}=x_p\pm\hat\rho$. Substitute $p\mapsto p-\hat\rho$ in the term with
$a_{p+\hat\rho}$ and $p\mapsto p+\hat\rho$ in the term with $a_{p-\hat\rho}$. Then use
$\xi^\rho(x_p\mp\hat\rho)=\xi^\rho(x_p)$ from \eqref{10.1:eq-rotational-identity-4}. This is summation by
parts on the torus in the direction $\rho$:
\begin{equation}\label{10.1:eq-rotational-identity-7}
\sum_pc(p)\chi(x_p)\xi^\rho(x_p)(2a)^{-1}(a_{p+\hat\rho}-a_{p-\hat\rho})
=-\sum_pa_p\,\xi^\rho(x_p)\,\nabla^s_\rho(c\chi)(x_p).
\end{equation}
Here $(c\chi)(p)=c(p)\chi(x_p)$.

At every plaquette, $c\chi=g_0^{-2}\chi-w$. Indeed, where $w(p)\neq0$ we have $\chi(x_p)=1$ by (c), and
where $w(p)=0$ both sides equal $g_0^{-2}\chi(x_p)$. Hence
$\nabla^s_\rho(c\chi)=g_0^{-2}\nabla^s_\rho\chi-\nabla^s_\rho w$. Summing \eqref{10.1:eq-rotational-identity-7}
over $\rho$ and using the definition of $\mathsf{L}^a_\xi$ in \eqref{10.1:eq-orbital-spin-densities-b}, the
first line of \eqref{10.1:eq-rotational-identity-6} equals
\[
-g_0^{-2}\sum_pa_p(\mathsf{L}^a_\xi\chi)(x_p)+\sum_pa_p(\mathsf{L}^a_\xi w)(p)
=\mathcal{O}(\mathsf{L}^a_\xi w)-g_0^{-2}\sum_pa_p(\mathsf{L}^a_\xi\chi)(x_p).
\]

On the support of $\mathfrak{h}_\chi$ we have $c=g_0^{-2}$. Indeed, by \eqref{10.1:eq-orbital-spin-densities-b},
$\mathfrak{h}_\chi(p)\neq0$ requires $\chi(x_\ell)\neq\chi(x_p)$ for some $\ell\subset\partial p$, with
$x_\ell$ at sup-distance $\tfrac12a$ from $x_p$. Then $\chi$ is not identically one on the closed
$2a$-neighbourhood of $x_p$, so $w(p)=0$ by (c). Hence
$\sum_pc(p)\mathfrak{h}_\chi(p)=g_0^{-2}\sum_p\mathfrak{h}_\chi(p)$, and
\eqref{10.1:eq-rotational-identity-6} becomes
\begin{equation}\label{10.1:eq-rotational-identity-8}
V_{\chi\xi}S_w=\mathcal{O}(\mathsf{L}^a_\xi w)-g_0^{-2}\sum_pa_p(\mathsf{L}^a_\xi\chi)(x_p)
+\mathfrak{S}_\chi[c]+\mathfrak{K}_\chi[c]+g_0^{-2}\sum_p\mathfrak{h}_\chi(p).
\end{equation}

\emph{Integration by parts.} Since $Z[w]\neq0$, integration by parts for the sourced measure
(Lemma~\ref{10.1:lem-integration-by-parts}) applies to $V=V_{\chi\xi}$ and gives
$\langle V_{\chi\xi}S_w\rangle_w=\langle\operatorname{div}V_{\chi\xi}\rangle_w$. Take the expectation of
\eqref{10.1:eq-rotational-identity-8}, solve for $\langle\mathcal{O}(\mathsf{L}^a_\xi w)\rangle_w$, and
collect the terms $g_0^{-2}\sum_pa_p(\mathsf{L}^a_\xi\chi)(x_p)$, $-g_0^{-2}\sum_p\mathfrak{h}_\chi(p)$ and
$\operatorname{div}V_{\chi\xi}$ into $\mathfrak{B}^{\mathrm{sh}}_\chi$ as in
\eqref{10.1:eq-rotational-identity-2}. This is \eqref{10.1:eq-rotational-identity-a}.

\emph{The divergence.} Part (iii) of Lemma~\ref{10.1:lem-displacement-divergence} with $\eta=\chi\xi$ expresses
$\operatorname{div}V_{\chi\xi}$ through $(\nabla^c_\rho(\chi\xi^\rho))(x_p)$ and
$(\nabla^c_\mu(\chi\xi^\mu))(x_p)$. Since $\xi^\rho(y\pm\tfrac12\hat\rho)=\xi^\rho(y)$,
\[
(\nabla^c_\rho(\chi\xi^\rho))(y)=a^{-1}\big[\chi(y+\tfrac12\hat\rho)-\chi(y-\tfrac12\hat\rho)\big]\xi^\rho(y)
=\xi^\rho(y)(\nabla^c_\rho\chi)(y).
\]
The same holds with $\mu$ in place of $\rho$. This gives \eqref{10.1:eq-rotational-identity-3}.

\emph{The cumulant identity.} Put $w=\sum_iz_if_i$ with $z=(z_1,\dots,z_n)\in\mathbb{C}^n$, so that
$c=g_0^{-2}-\sum_iz_if_i$. The set $\{x_p:w(p)\neq0\}$ lies in $\bigcup_i\operatorname{supp}f_i$, so by the
hypothesis on $\chi$ the requirement on $\chi$ in (c) is met for every $z$; it remains to secure
$Z[w]\neq0$. The function $Z[w]$ is
the integral over a compact group of an integrand continuous in $U$ and entire in $z$, and $Z[0]>0$. Hence
$Z[w]\neq0$ for $z$ in a neighbourhood of $0$. There the data are admissible,
\eqref{10.1:eq-rotational-identity-a} holds identically, and $z\mapsto\langle Y\rangle_w$ is smooth for every
$U$-measurable bounded $Y$ independent of $z$.

Apply $\partial_{z_1}\cdots\partial_{z_n}$ at $z=0$ to \eqref{10.1:eq-rotational-identity-a}, using the
definition of the truncated bracket for $z$-free $Y$,
\[
\langle Y;\mathcal{O}(f_1);\dots;\mathcal{O}(f_n)\rangle
=\partial_{z_1}\cdots\partial_{z_n}\langle Y\rangle_w\big|_{z=0}.
\]
For $z$-free $Y$ and a source $w$ with $Z[w]\neq0$ we put
\[
\langle Y;\mathcal{O}(f)\rangle_w:=\partial_t\langle Y\rangle_{w+tf}\big|_{t=0};
\]
with $w=\sum_iz_if_i$ this gives $\partial_{z_i}\langle Y\rangle_w=\langle Y;\mathcal{O}(f_i)\rangle_w$. We
treat the four terms in turn.

(1) The left side. Since $\mathsf{L}^a_\xi$ is linear, $\mathcal{O}(\mathsf{L}^a_\xi w)=\sum_iz_i\mathcal{O}(\mathsf{L}^a_\xi f_i)$.
By the Leibniz rule and evaluation at $z_i=0$, the derivative of $z_i\langle\mathcal{O}(\mathsf{L}^a_\xi f_i)\rangle_w$
is $\langle\mathcal{O}(\mathsf{L}^a_\xi f_i);(\mathcal{O}(f_l))_{l\neq i}\rangle$. Let $w'$ be the source
obtained from $w$ by replacing $z_if_i$ with $z_i\mathsf{L}^a_\xi f_i$. Then
$\partial_{z_i}\log Z[w']=\langle\mathcal{O}(\mathsf{L}^a_\xi f_i)\rangle_{w'}$, and $w'=w$ at $z_i=0$.
The mixed partial derivatives commute, so we may set
$z_i=0$ first. Therefore this cumulant is $S^T_{n;K,m}(f_1,\dots,\mathsf{L}^a_\xi f_i,\dots,f_n)$.

(2) The spin and orbital terms. We have
$\mathfrak{S}_\chi[c]=g_0^{-2}\mathfrak{S}_\chi[1]-\sum_iz_i\mathfrak{S}_\chi[f_i]$, and likewise for
$\mathfrak{K}_\chi$. By the same Leibniz argument, $-\langle\mathfrak{S}_\chi[c]\rangle_w-\langle\mathfrak{K}_\chi[c]\rangle_w$
contributes
\[
-g_0^{-2}\langle(\mathfrak{S}_\chi+\mathfrak{K}_\chi)[1];\mathcal{O}(f_1);\dots;\mathcal{O}(f_n)\rangle
+\sum_i\langle(\mathfrak{S}_\chi+\mathfrak{K}_\chi)[f_i];(\mathcal{O}(f_l))_{l\neq i}\rangle .
\]

(3) Removal of $\chi$. Since $\chi(x_p)=1$ wherever $f_i(x_p)\neq0$, we have
$\mathfrak{S}_\chi[f_i]=\mathfrak{S}[f_i]$ and $\mathfrak{K}_\chi[f_i]=\mathfrak{K}[f_i]$.

(4) The shell term. $\mathfrak{B}^{\mathrm{sh}}_\chi$ does not depend on $z$, so it contributes
$\langle\mathfrak{B}^{\mathrm{sh}}_\chi;\mathcal{O}(f_1);\dots;\mathcal{O}(f_n)\rangle$.

Collecting (1)--(4) gives \eqref{10.1:eq-rotational-identity-b}.
\end{proof}

\begin{remark}[Domination, the abelian case and the torus form. Stated; proof incomplete at the step marked $(\ast)$]\label{10.1:rem-domination-abelian-torus}
Let $\omega$, $\chi$, $w$ be admissible data in the sense of the rotational Ward identity
(Theorem~\ref{10.1:thm-rotational-identity}), with $\xi(x)=\omega x$, and let
$\mathfrak{S}_\chi[c]$, $\mathfrak{K}_\chi[c]$, $\mathfrak{B}^{\mathrm{sh}}_\chi$ be the spin term,
the orbital remainder and the shell term of \eqref{10.1:eq-rotational-identity-a}.
\begin{enumerate}
\item[(1)] \emph{Nothing is chosen.} No lattice energy--momentum tensor (no lattice counterpart of
the stress tensor $T_{\sigma\rho}$ of item~(5)) enters the
identity \eqref{10.1:eq-rotational-identity-a}: its three functionals are given explicitly by the
displacement derivation \eqref{10.1:eq-displacement-derivation-a},
\eqref{10.1:eq-displacement-derivation-b} and the densities of
Definition~\ref{10.1:def-orbital-spin-densities}. The bulk contribution of the divergence
$\operatorname{div} V_{\chi\xi}$ vanishes wherever $\chi$ is constant: the divergence identity of
part~(iii) of Lemma~\ref{10.1:lem-displacement-divergence} holds exactly for every link, and one
summation by parts (a shift of the second sum by $\hat\rho$) brings it into plaquette form.
\item[(2)] \emph{Domination.} On $SU(2)$-valued fields, pointwise,
\begin{equation*}
|\mathfrak{S}_\chi[1]| \le C|\omega|\sum_{p'\in\mathcal{N}_\chi} a_{p'},
\end{equation*}
where
$|\omega|:=\max_{\rho,\sigma}|\omega_{\rho\sigma}|$, $C$ is a constant depending on the dimension
only, and $\mathcal{N}_\chi$ is the union of the sets $N(p)$ of
\eqref{10.1:eq-orbital-spin-densities-c} over all directions $\rho$ and all plaquettes $p$ with
$\chi(x_p)\neq0$; thus the spin term,
with its prefactor $g_0^{-2}$, is dominated by a fixed multiple of the action. The orbital
remainder $\mathfrak{K}_\chi[1]$ is dominated in the same way only with the constant
$\sup|\xi|/a$.
\item[(3)] \emph{The abelian case.} Consider the non-compact Gaussian $U(1)$ model, in which the bond
variables are real numbers and $F_p$ is the oriented sum of the bond variables over the links of
$\partial p$; there $a_p=\tfrac12 F_p^2$. Let $p\pm\hat\rho$ denote the plaquette $p$ translated
by $\pm\hat\rho$, let $F_{\rho\mu}(y)$ be the field strength of the plaquette $\{y;\rho,\mu\}$,
and, for $\ell=(x,\mu)$, let $u_{\ell,\rho}:=\tfrac12\bigl(F_{\rho\mu}(x)+F_{\rho\mu}(x-\hat\rho)\bigr)$
be the mean of the field strengths of the two plaquettes adjacent to $\ell$ in the direction
$\rho$. In this model $V_{e_\rho}F_p=\nabla^s_\rho F_p$ exactly and
\begin{equation}\label{10.1:eq-domination-abelian-torus-1}
\mathfrak{r}_\rho(p)
=-\tfrac14\Bigl[\bigl(F_{p+\hat\rho}-F_p\bigr)^2-\bigl(F_p-F_{p-\hat\rho}\bigr)^2\Bigr],
\end{equation}
an exact $\rho$-difference. Consequently $\mathfrak{K}_\chi[g_0^{-2}]$ is a shell term. For $f$
supported where $\chi\equiv1$ write
\begin{equation*}
\mathfrak{K}[f]:=\mathfrak{K}_\chi[f]=\sum_pf(x_p)\sum_\rho a^{-1}\xi^\rho(x_p)\mathfrak{r}_\rho(p);
\end{equation*}
then
\begin{equation}\label{10.1:eq-domination-abelian-torus-2}
\mathfrak{K}[f]=\tfrac14\sum_p\sum_\rho \xi^\rho(x_p)\,a^{-1}\bigl(f(x_p+\hat\rho)-f(x_p)\bigr)
\bigl(F_{p+\hat\rho}-F_p\bigr)^2
\end{equation}
is a contact term with bounded weight $\xi\cdot\nabla f$: in the free theory the orbital remainder
has no bulk part. For $SU(2)$ it has one: $V_{e_\rho}A\neq 0$ as a function.
\item[(4)] \emph{The case $N\ge 3$.} The same holds with $\vec w$ replaced by the anti-Hermitian
traceless part of $W$ and the number $3$ replaced by a Casimir number. This is not used in what
follows.
\item[(5)] \emph{Torus form; the step marked $(\ast)$.} Let $\chi$ be the
indicator of the cube $\mathopen{]}-L^m+2a,\,L^m-2a\mathclose{[}^4$. Then the
shell term sits on the cut locus of the torus $\mathbb{T}_m$ and is, classically,
\begin{equation}\label{10.1:eq-domination-abelian-torus-3}
\sum_{\sigma,\rho}2L^m\,\omega_{\rho\sigma}\int_{x_\sigma=L^m}T_{\sigma\rho}\,dS,
\end{equation}
that is, (side of the torus) $\times$ (momentum in the direction $\rho$ flowing through the face
orthogonal to $e_\sigma$): the response of the expectation $\langle\cdot\rangle$ to an
infinitesimal tilt of the period lattice. In this form the far-sphere flux hypothesis
$(\mathrm{IR}\text{-}\mathrm{fl})_{\omega_0}$ of clause~(vi) of Theorem~0 (the flux of the lattice
rotation current through far spheres vanishes along one sequence of radii) and the
statement on the
response of the expectations to a finite change of shape of the period lattice are the
infinitesimal and the finite version of one statement.
\end{enumerate}
\end{remark}

\begin{proof}
\emph{Item (1).} The functionals $\mathfrak{S}_\chi[c]$, $\mathfrak{K}_\chi[c]$,
$\mathfrak{B}^{\mathrm{sh}}_\chi$ are the expressions displayed in
Theorem~\ref{10.1:thm-rotational-identity}; they are built from the densities
$\mathfrak{s}_{\sigma\rho}(p)$, $\mathfrak{r}_\rho(p)$ of
Definition~\ref{10.1:def-orbital-spin-densities}, from $\mathfrak{h}_\chi$, from
$\mathsf{L}^a_\xi\chi$ and from $\operatorname{div} V_{\chi\xi}$, and no further object is
chosen. For the divergence, apply part~(iii) of the displacement derivation lemma
(Lemma~\ref{10.1:lem-displacement-divergence}) to $\eta=\chi\xi$. Its first form is a sum over
links; at each link $\ell$ the summand comes from the exact identity
$\sum_b\partial^b_\ell u^b_{\ell,\rho}=\tfrac32\bigl(a_{p(\ell,\rho)}-a_{p(\ell,-\rho)}\bigr)$
established in the proof of Lemma~\ref{10.1:lem-displacement-divergence}. One summation by parts,
the shift of the second sum by $\hat\rho$ (recall $p(\ell,-\rho)=\{x-\hat\rho;\mu,\rho\}$ for
$\ell=(x,\mu)$), gives the plaquette form:
\begin{align*}
\operatorname{div}V_{\chi\xi}
&=\frac{3}{2a}\sum_\ell\sum_\rho\eta^\rho(x_\ell)\bigl[a_{p(\ell,\rho)}-a_{p(\ell,-\rho)}\bigr]\\
&=-\frac32\sum_{p=\{x;\mu,\rho\}}a_p\bigl[(\nabla^c_\rho\eta^\rho)(x_p)
+(\nabla^c_\mu\eta^\mu)(x_p)\bigr].
\end{align*}
Since $\hat\rho=ae_\rho$ and $\omega$ is antisymmetric,
\begin{equation}\label{10.1:eq-domination-abelian-torus-4}
\xi^\rho\bigl(y\pm\tfrac12\hat\rho\bigr)=\sum_\nu\omega_{\rho\nu}\bigl(y_\nu\pm\tfrac12 a\,
\delta_{\nu\rho}\bigr)=\xi^\rho(y)\pm\tfrac12 a\,\omega_{\rho\rho}=\xi^\rho(y),
\end{equation}
so that $(\nabla^c_\rho(\chi\xi^\rho))(y)=\xi^\rho(y)(\nabla^c_\rho\chi)(y)$ (no sum over $\rho$).
Hence
\begin{equation*}
\operatorname{div}V_{\chi\xi}=-\frac32\sum_{p=\{x;\mu,\rho\}}a_p\bigl[\xi^\rho\nabla^c_\rho\chi
+\xi^\mu\nabla^c_\mu\chi\bigr](x_p),
\end{equation*}
which is the expression in
Theorem~\ref{10.1:thm-rotational-identity}, and every summand vanishes at plaquettes $p$ for which
$\chi$ takes equal values at $x_p\pm\frac12\hat\rho$ and at $x_p\pm\frac12\hat\mu$.

\emph{Item (2).} Let $|\omega|$ and $\mathcal{N}_\chi$ be as in the statement, $N(p)$ being as in
\eqref{10.1:eq-orbital-spin-densities-c}. By definition
\begin{equation*}
\mathfrak{S}_\chi[1]=\sum_p\chi(x_p)\sum_{\rho,\sigma}\omega_{\rho\sigma}\mathfrak{s}_{\sigma\rho}(p).
\end{equation*}
Using $0\le\chi\le1$ and the bound $|\mathfrak{s}_{\sigma\rho}(p)|\le4\sum_{p'\in N(p)}a_{p'}$ of
\eqref{10.1:eq-orbital-spin-densities-c}, which holds on $SU(2)$-valued fields, we get
\begin{equation*}
|\mathfrak{S}_\chi[1]|\le 4|\omega|\sum_{p:\chi(x_p)\neq0}\ \sum_{\rho,\sigma}\
\sum_{p'\in N(p)}a_{p'}.
\end{equation*}
Since $a_{p'}\ge0$ and a given plaquette $p'$ lies in $N(p)$ for a number of pairs $(p,\rho)$
bounded in terms of the dimension only, the right side is at most
$C|\omega|\sum_{p'\in\mathcal{N}_\chi}a_{p'}$. Multiplying by $g_0^{-2}$ gives the domination of
the spin term by a fixed multiple of the action $g_0^{-2}\sum_{p'}a_{p'}$. For
\begin{equation*}
\mathfrak{K}_\chi[1]=\sum_p\chi(x_p)\sum_\rho a^{-1}\xi^\rho(x_p)\mathfrak{r}_\rho(p)
\end{equation*}
the same
argument, with the bound on $|\mathfrak{r}_\rho(p)|$ from
\eqref{10.1:eq-orbital-spin-densities-c}, gives
\begin{equation*}
|\mathfrak{K}_\chi[1]|\le C\,a^{-1}\sup_{x_p\in\operatorname{supp}\chi}|\xi(x_p)|
\sum_{p'\in\mathcal{N}_\chi}a_{p'};
\end{equation*}
the weight $a^{-1}\xi^\rho(x_p)$ is not bounded uniformly in the lattice spacing, and the constant
the argument gives is $\sup|\xi|/a$.

\emph{Item (3).} In the non-compact Gaussian $U(1)$ model, with $a_p=\frac12F_p^2$ and
$u_{\ell,\rho}$ as stated, the lattice Bianchi identity gives $V_{e_\rho}F_p=\nabla^s_\rho F_p$
exactly and the formula \eqref{10.1:eq-domination-abelian-torus-1}. Write
$\Delta_\rho(p):=F_{p+\hat\rho}-F_p$, so that
$\mathfrak{r}_\rho(p)=-\frac14[\Delta_\rho(p)^2-\Delta_\rho(p-\hat\rho)^2]$. For a function
$\varphi$ on the torus, shifting $p\mapsto p+\hat\rho$ in the second sum gives
\[
\sum_p\varphi(x_p)\,\mathfrak{r}_\rho(p)=\tfrac14\sum_p \Delta_\rho(p)^2\bigl(\varphi(x_p+\hat\rho)
-\varphi(x_p)\bigr).
\]
Take $\varphi=f\xi^\rho$ with $f$ supported where $\chi\equiv1$ (so that $f\xi^\rho$ is a function
on the torus), and use $\xi^\rho(x_p+\hat\rho)=\xi^\rho(x_p)$, which follows from
$\omega_{\rho\rho}=0$ as in \eqref{10.1:eq-domination-abelian-torus-4}; multiplying by $a^{-1}$
and summing over $\rho$ yields \eqref{10.1:eq-domination-abelian-torus-2}; here $\chi(x_p)=1$
wherever $f(x_p)\neq0$, which is why $\mathfrak{K}_\chi[f]$ has the form used in the definition of
$\mathfrak{K}[f]$. Its weight
$\xi^\rho(x_p)a^{-1}(f(x_p+\hat\rho)-f(x_p))$ is bounded by $|\xi(x_p)|$ times the Lipschitz
constant of $f$. Taking instead $\varphi=g_0^{-2}\chi\xi^\rho$, which is a function on the torus
by admissibility, gives
\begin{equation*}
\mathfrak{K}_\chi[g_0^{-2}]=\frac14g_0^{-2}\sum_p\sum_\rho\xi^\rho(x_p)a^{-1}
\bigl(\chi(x_p+\hat\rho)-\chi(x_p)\bigr)\Delta_\rho(p)^2,
\end{equation*}
whose summand for the pair $(p,\rho)$ vanishes wherever
$\chi(x_p+\hat\rho)=\chi(x_p)$: a shell term. For the final assertion of item~(3), recall that
$A=\sum_pa_p$ and that $V_{e_\rho}$ acts as a derivation, so that part~(ii) of
Lemma~\ref{10.1:lem-displacement-divergence}, applied with the constant displacement
$\eta=e_\rho$, gives
\begin{equation*}
V_{e_\rho}A=a^{-1}\sum_p\sum_{\ell\subset\partial p}\mathfrak{d}_{\ell,p,\rho};
\end{equation*}
on $SU(2)$-valued fields this function is not identically zero.

\emph{Item (4).} The construction of the displacement derivation and of the densities uses
$\vec w$ and the number $3$, wherever it enters (for instance in the coefficient $3/2$ of
part~(iii) of Lemma~\ref{10.1:lem-displacement-divergence}); replacing $\vec w$ by the
anti-Hermitian traceless part of $W$ and $3$ by a Casimir number gives the same structure for
$N\ge3$. It is not used.

\emph{Item (5)}, the step marked $(\ast)$. The identification of the shell term, for $\chi$ the indicator of
$\mathopen{]}-L^m+2a,\,L^m-2a\mathclose{[}^4$, with the boundary expression
\eqref{10.1:eq-domination-abelian-torus-3} and
with the response to an infinitesimal tilt of the period lattice is stated and is not proved here.
\end{proof}

\begin{definition}[Transported links and the projected central difference]
\label{10.1:def-projected-difference}
Let $\mathbb{H}$ be the real span of the identity matrix $1$ and of $i\sigma^1, i\sigma^2, i\sigma^3$
in the complex
$2\times 2$ matrices, the quaternions, identified with $\mathbb{R}^4$ through
$q = q^0 + i\vec q\cdot\vec\sigma \mapsto (q^0,\vec q)$. Let $\bar q := q^0 - i\vec q\cdot\vec\sigma$ be
the conjugate, and put
\[
\langle q, q'\rangle := \operatorname{tr}(q\bar q'), \qquad |q|^2 := \langle q, q\rangle, \qquad
\Pi_q X := X - \langle X, q\rangle q \qquad (q, q', X \in \mathbb{H}),
\]
with $\operatorname{tr} 1 = 1$. Then $SU(2)$ is the unit sphere of $\mathbb{H}$, and for
$|q| = 1$ the map $\Pi_q$ is the orthogonal projection onto the tangent space of the sphere at $q$.

Fix a direction $\rho$ and a link $\ell = (x,\mu)$ with $\mu \neq \rho$. Put
\[
q_\ell := U_\ell, \qquad q_\ell^{\pm} := W_{\ell,\pm\rho}^{-1}\, U_\ell .
\]
Thus $q_\ell^{\pm}$ is the link $\ell \pm \hat\rho$ transported to $\ell$ by the two $\rho$-links joining
their endpoints. Equivalently, it is the image of $U_\ell$ under the exact lattice symmetry consisting
of the translation by $\pm\hat\rho$ followed by the gauge transformation by the $\rho$-link variables.
These elements satisfy
\begin{equation}\label{10.1:eq-projected-difference-1}
\langle q_\ell, q_\ell^{\pm}\rangle = 1 - a_{p(\ell,\pm\rho)}(U),
\end{equation}
and, with $\vec u_{\ell,\rho}$ as in \eqref{10.1:eq-displacement-derivation-a},
\begin{equation}\label{10.1:eq-projected-difference-a}
(i\,\vec u_{\ell,\rho}\cdot\vec\sigma)\,U_\ell
= \tfrac12\bigl[\Pi_{q_\ell} q_\ell^{+} - \Pi_{q_\ell} q_\ell^{-}\bigr].
\end{equation}
Consequently $a V_{e_\rho}$, where $V_{e_\rho}$ is the derivation of the constant displacement
$e_\rho$ (Definition~\ref{10.1:def-displacement-derivation}), is the central difference of the field
along $\rho$, projected on the tangent space of the sphere.

Finally, define
\[
T(q_1,q_2,q_3,q_4) := \operatorname{tr}(q_1 q_2 \bar q_3 \bar q_4), \qquad q_1,\dots,q_4 \in \mathbb{H}.
\]
This is a real $4$-linear form on $\mathbb{H}^4$ with $|T(q_1,q_2,q_3,q_4)| \le \prod_{i=1}^4 |q_i|$.
The letters $q$, $q_i$ and the form $T$ are used in this sense only within the present section.
\end{definition}

\begin{proof}
\emph{The inner product.} Let $q = q^0 + i\vec q\cdot\vec\sigma$ and
$q' = q'^0 + i\vec q\,'\cdot\vec\sigma$. Since
$(\vec q\cdot\vec\sigma)(\vec q\,'\cdot\vec\sigma) = (\vec q\cdot\vec q\,')\cdot 1
+ i(\vec q\times\vec q\,')\cdot\vec\sigma$ and the $\sigma^b$ are traceless, we get
$\langle q,q'\rangle = q^0 q'^0 + \vec q\cdot\vec q\,'$. This is the Euclidean inner product of
$\mathbb{R}^4$. The same identity gives $q\bar q = |q|^2\cdot 1$, so $\bar q = q^{-1}$ when
$|q| = 1$. Moreover $\overline{q q'} = \bar q'\bar q$, and
$|qq'|^2\cdot 1 = q q'\bar q'\bar q = |q|^2|q'|^2\cdot 1$.

\emph{Proof of \eqref{10.1:eq-projected-difference-1}.} Write $W := W_{\ell,\pm\rho} \in SU(2)$. Then
$\overline{W^{-1}U_\ell} = \bar U_\ell\,\overline{W^{-1}} = U_\ell^{-1}W$, hence
$q_\ell\,\overline{q_\ell^{\pm}} = W$ and $\langle q_\ell, q_\ell^{\pm}\rangle = \operatorname{tr} W = w^0$.
By the identity $a_p = 1 - w^0_{\ell,p}$ for $\ell\subset\partial p$, recorded in
Notation~\ref{10.1:not-link-derivations}, applied to the plaquette
$p = p(\ell,\pm\rho)$, which contains $\ell$, this equals $1 - a_{p(\ell,\pm\rho)}$.

\emph{Proof of \eqref{10.1:eq-projected-difference-a}.} Let $W = w^0 + i\vec w\cdot\vec\sigma\in SU(2)$
and $U\in SU(2)$. Then $i\vec w\cdot\vec\sigma = \frac12(W - \bar W)$ and $W + \bar W = 2w^0$, so
$WU = 2w^0 U - \bar W U$. Hence
\begin{equation}\label{10.1:eq-projected-difference-2}
-(i\vec w\cdot\vec\sigma)\,U = \tfrac12(\bar W U - W U) = \bar W U - w^0 U = \Pi_U(\bar W U).
\end{equation}
The last step uses $\langle \bar W U, U\rangle = \operatorname{tr}(\bar W U\bar U) = \operatorname{tr}\bar W = w^0$.

Take $U = U_\ell$ and $W = W_{\ell,\pm\rho}$, so that $\vec w = \vec w_{\ell,\pm\rho}$. Since
$\bar W U_\ell = W^{-1}U_\ell = q_\ell^{\pm}$, \eqref{10.1:eq-projected-difference-2} reads
$-(i\vec w_{\ell,\pm\rho}\cdot\vec\sigma)U_\ell = \Pi_{q_\ell} q_\ell^{\pm}$. By
\eqref{10.1:eq-displacement-derivation-a}, $\vec u_{\ell,\rho} = -\frac12(\vec w_{\ell,+\rho} - \vec
w_{\ell,-\rho})$. Subtracting the two cases and halving gives
\eqref{10.1:eq-projected-difference-a}.

\emph{The reading of $aV_{e_\rho}$.} For $\eta = e_\rho$, the definition
\eqref{10.1:eq-displacement-derivation-a} gives $a\vec v_\ell = \vec u_{\ell,\rho}$. Hence
$aV_{e_\rho} = \sum_{\ell,b} u^b_{\ell,\rho}\,\partial^b_\ell$. The derivation $\partial^b_\ell$
differentiates along the curve $s\mapsto e^{is\sigma^b}U_\ell$ in the variable $U_\ell$, whose velocity
at $s=0$ is $i\sigma^b U_\ell$. Therefore at each link $\ell = (x,\mu)$ with $\mu\ne\rho$ the
derivation $aV_{e_\rho}$ differentiates along the tangent vector $(i\vec u_{\ell,\rho}\cdot\vec\sigma)
U_\ell$. By \eqref{10.1:eq-projected-difference-a}, this vector is half the difference of the tangential
projections at $q_\ell$ of the two transported neighbours $q_\ell^{\pm}$. At links with $\mu = \rho$
the component vanishes, since $\vec u_{\ell,\rho} = 0$ there.

\emph{The form $T$.} Conjugation is $\mathbb{R}$-linear and the product is $\mathbb{R}$-bilinear, so
$T$ is $\mathbb{R}$-linear in each argument. The product $q_1 q_2\bar q_3\bar q_4$ lies in
$\mathbb{H}$; its trace is its real component, which is real and bounded in absolute value by the
norm of the product. By the multiplicativity of the norm, that norm equals
$|q_1||q_2||\bar q_3||\bar q_4| = \prod_i |q_i|$, since $|\bar q| = |q|$.
\end{proof}

\begin{lemma}[Orbital remainder: exact difference plus cubic term]\label{10.1:lem-orbital-remainder}
The gauge group is $SU(2)$, regarded as the unit sphere of the quaternions
$\mathbb{H}\cong\mathbb{R}^4$, with the inner product $\langle q,q'\rangle=\operatorname{tr}(q\bar q')$,
$|q|^2=\langle q,q\rangle$, the tangent projections $\Pi_q$ and the transported links $q_\ell^{\pm}$
of Definition~\ref{10.1:def-projected-difference}, and the real four-linear form
$T(q_1,q_2,q_3,q_4)=\operatorname{tr}(q_1q_2\bar q_3\bar q_4)$ on $\mathbb{H}^4$; $\mathfrak{r}_\rho(p)$ is the
orbital remainder density of \eqref{10.1:eq-orbital-spin-densities-a}. For $x\in\mathfrak{L}$ we write
$(x,\mu)$ for the link from $x$ to $x+\hat\mu$, $p=\{x;\mu,\nu\}$ for the plaquette with corners $x$,
$x+\hat\mu$, $x+\hat\nu$, $x+\hat\mu+\hat\nu$, and $p(\ell,\pm\rho)$ for the plaquette spanned by the link
$\ell$ and $\pm\hat\rho$. Fix a direction $\rho$.
\begin{itemize}
\item[(i)] (Plaquettes containing the direction $\rho$.) For $p=\{x;\rho,\nu\}$ let
$R(p):=p\cup(p+\hat\rho)$ be the $1\times2$ rectangle, and for a rectangle $R$ let
$a_R:=1-\operatorname{Re}\operatorname{tr}U(\partial R)$. Then
\begin{equation}\label{10.1:eq-orbital-remainder-1}
\mathfrak{r}_\rho(p)=\Psi_\rho(p)-\Psi_\rho(p-\hat\rho),\qquad
\Psi_\rho(p):=\tfrac12 a_{R(p)}-a_p-a_{p+\hat\rho}+\tfrac12 a_pa_{p+\hat\rho}:
\end{equation}
for such plaquettes the orbital remainder density is an exact difference in the direction $\rho$;
there is no further term.
\item[(ii)] (Plaquettes orthogonal to $\rho$.) For $p=\{x;\mu,\nu\}$ with $\rho\notin\{\mu,\nu\}$ let
$\ell_1=(x,\mu)$, $\ell_2=(x+\hat\mu,\nu)$, $\ell_3=(x+\hat\nu,\mu)$, $\ell_4=(x,\nu)$, and put
$q_i:=q_{\ell_i}$, $\delta_i:=q_{\ell_i}^+-q_i$ (so that $|\delta_i|^2=2a_{p(\ell_i,\rho)}$),
$q:=(q_1,\dots,q_4)$ and $q+\delta:=(q_1+\delta_1,\dots,q_4+\delta_4)$. For
$I\subset\{1,2,3,4\}$ let $q^{(I)}$ be $q$ with $q_i$ replaced by $\delta_i$ for every $i\in I$, and put
\[
T_k(p):=\sum_{|I|=k}T\bigl(q^{(I)}\bigr)\quad(0\le k\le4),\qquad
\sigma(p):=\sum_{i=1}^4a_{p(\ell_i,\rho)},
\]
the sum of the densities of the four side plaquettes of the cube between $p$ and $p+\hat\rho$. Then
$T(q)=1-a_p$, $T(q+\delta)=1-a_{p+\hat\rho}$, and
\begin{equation}\label{10.1:eq-orbital-remainder-2}
\begin{aligned}
\mathfrak{r}_\rho(p)&=\Psi_\rho(p)-\Psi_\rho(p-\hat\rho)+\mathfrak{r}^{\sharp}_\rho(p),\\
\Psi_\rho(p)&:=\tfrac12T_2+\tfrac34T_3+T_4
-\tfrac12\sigma\bigl(1-\tfrac12(a_p+a_{p+\hat\rho})\bigr),\\
\mathfrak{r}^{\sharp}_\rho(p)&:=\tfrac12\bigl[E^a(p)+E^a(p-\hat\rho)\bigr],\\
E^a(p)&:=\tfrac12\bigl[\sigma(p)(a_p-a_{p+\hat\rho})-T_3(p)-2T_4(p)\bigr].
\end{aligned}
\end{equation}
\item[(iii)] On $SU(2)$ fields $|E^a(p)|\le C\sigma(p)^{3/2}$ for every plaquette $p$ as in (ii), with
an absolute constant $C$: the term $\mathfrak{r}^{\sharp}_\rho$ is at least cubic in the plaquette
variables of the plaquettes containing the direction $\rho$; it has no quadratic part. For
plaquettes $p$ containing the direction $\rho$ we put $\mathfrak{r}^{\sharp}_\rho(p):=0$.
\item[(iv)] Classically, that is for the lattice field of a smooth continuum gauge field as $a\to0$,
$\Psi_\rho=O(a^6)$ and $\mathfrak{r}^{\sharp}_\rho=O(a^7)$. On constant abelian fields $\Psi_\rho$ is
quartic; for plaquettes containing the direction $\rho$,
$\Psi_\rho=-\tfrac12(1-\cos\theta)^2$, where $\theta$ is the angle of the plaquette holonomy.
\end{itemize}
\end{lemma}

\begin{proof}
All quantities in the lemma are gauge invariant. We therefore fix the gauge so that $U_\ell=1$ on
each of the finitely many links $\ell$ in direction $\rho$ that enter the quantities below. For a
link $\ell$ in a direction other than $\rho$ and for $t\in\mathbb{Z}$ we write $q(t)$, or $q_i(t)$
when $\ell=\ell_i$, for the variable of the link $\ell+t\hat\rho$. In this gauge the transport along
the links in direction $\rho$ is the identity, so $q_i^{\pm}=q_i(\pm1)$. By
\eqref{10.1:eq-projected-difference-a}, which expresses the displacement in direction $\rho$ on each
link $\ell$ in a direction other than $\rho$ as the projected central difference
$\tfrac12[\Pi_{q_\ell}q_\ell^+-\Pi_{q_\ell}q_\ell^-]$, the orbital density of
\eqref{10.1:eq-orbital-spin-densities-a} is the derivative
\begin{equation}\label{10.1:eq-orbital-remainder-3}
\mathfrak{t}_\rho(p)=dG(q)\bigl[\tfrac12\Pi(q^+-q^-)\bigr],\qquad G:=a_p,
\end{equation}
of $a_p$, regarded as a function of the variables of the links of $\partial p$, each on the unit
sphere, in the direction of the tangent vector whose component at the link $\ell$ is
$\tfrac12\Pi_{q_\ell}(q_\ell^+-q_\ell^-)$. We prove (ii) first, then (i), (iii) and (iv).

\emph{(ii).} Let $p=\{x;\mu,\nu\}$ with $\rho\notin\{\mu,\nu\}$, and write
$q(t):=(q_1(t),\dots,q_4(t))$ for the quadruple of variables of the links $\ell_i+t\hat\rho$, so that
$q(0)=q$ and, since $q_i(1)=q_i^+=q_i+\delta_i$, $q(1)=q+\delta$. On $SU(2)$ the inverse is the
quaternion conjugate and the trace is real, so with the orientation of $\partial p$ for which
$U(\partial p)=U_{\ell_1}U_{\ell_2}U_{\ell_3}^{-1}U_{\ell_4}^{-1}$ we have
$\operatorname{Re}\operatorname{tr}U(\partial p)=T(q)$, that is $G=1-T$ on the product of the four
unit spheres, with $T$ four-linear; this is the identity $T(q)=1-a_p$. The links of $p+\hat\rho$ are
the $\ell_i+\hat\rho$, with variables $q_i(1)=q_i+\delta_i$, so likewise $T(q+\delta)=1-a_{p+\hat\rho}$,
and $T(q(-1))=1-a_{p-\hat\rho}$. The boundary of the side plaquette $p(\ell_i,\rho)$ consists of
$\ell_i$, $\ell_i+\hat\rho$ and two links in direction $\rho$, whose variables are $1$; hence
$\operatorname{Re}\operatorname{tr}U(\partial p(\ell_i,\rho))=\operatorname{tr}(q_i\bar q_i(1))$, which is
$\langle q_i,q_i^+\rangle$, so $\langle q_i,q_i^+\rangle=1-a_{p(\ell_i,\rho)}$, and since
$\langle\cdot,\cdot\rangle$ is the Euclidean inner product of $\mathbb{H}\cong\mathbb{R}^4$ and
$|q_i|=|q_i^+|=1$,
$|\delta_i|^2=2-2\langle q_i^+,q_i\rangle=2a_{p(\ell_i,\rho)}$.

We work with two general quadruples $b=(b_1,\dots,b_4)$, $b'=(b'_1,\dots,b'_4)$ of unit
quaternions. Write $T(b;i\mapsto Y)$ for the value of $T$ at $b$ with its $i$-th argument replaced by
$Y\in\mathbb{H}$, let $T_k(b;b')$ be the sum over $|I|=k$ of $T$ evaluated at $b$ with $b_i$ replaced
by $b'_i-b_i$ for every $i\in I$ (so $T_0(b;b')=T(b)$), and put
\begin{align*}
\sigma(b;b')&:=\sum_{i=1}^4\bigl(1-\langle b_i,b'_i\rangle\bigr),\\
E(b;b')&:=T_2(b;b')+T_3(b;b')+T_4(b;b')-\sigma(b;b')\,T(b).
\end{align*}
By the preceding paragraph, $T_k(p)=T_k(q(0);q(1))$ and $\sigma(p)=\sigma(q(0);q(1))$.

Since $G$ is the restriction of the polynomial $1-T$ to the product of spheres, its differential at
$b$ on a tangent vector $Y=(Y_i)$ is $-\sum_iT(b;i\mapsto Y_i)$, by multilinearity. For
$X\in\mathbb{H}^4$ let $\Pi X$ denote the tangent vector with components
$\Pi_{b_i}X_i=X_i-\langle X_i,b_i\rangle b_i$. Linearity of $T$ in its $i$-th argument and
$T(b;i\mapsto b_i)=T(b)$ give
\begin{equation}\label{10.1:eq-orbital-remainder-4}
dG(b)[\Pi X]=-\sum_{i=1}^4T(b;i\mapsto X_i)+\sum_{i=1}^4\langle X_i,b_i\rangle\,T(b).
\end{equation}
Take $X=b'$. Writing $b'_i=b_i+(b'_i-b_i)$ in the $i$-th slot gives
$T(b;i\mapsto b'_i)=T(b)+T(b;i\mapsto b'_i-b_i)$, whence $\sum_iT(b;i\mapsto b'_i)=4T(b)+T_1(b;b')$;
and $\sum_i\langle b'_i,b_i\rangle=4-\sigma(b;b')$. Thus
$dG(b)[\Pi b']=-T_1(b;b')-\sigma(b;b')T(b)$. On the other hand, writing every argument of $T(b')$ as
$b_i+(b'_i-b_i)$, expanding by multilinearity and recalling $T_0(b;b')=T(b)$,
\[
T(b')=\sum_{k=0}^{4}T_k(b;b'),\qquad G(b')-G(b)=-(T_1+T_2+T_3+T_4)(b;b').
\]
Comparing the two expressions,
\begin{equation}\label{10.1:eq-orbital-remainder-5}
dG(b)[\Pi b']=G(b')-G(b)+E(b;b').
\end{equation}
Now $\tfrac12\Pi(q^+-q^-)=\tfrac12\Pi q(1)-\tfrac12\Pi q(-1)$ at the base point $q=q(0)$, so
\eqref{10.1:eq-orbital-remainder-3}, applied with \eqref{10.1:eq-orbital-remainder-5} once for
$b'=q(1)$ and once for $b'=q(-1)$, together with $G(q(\pm1))=a_{p\pm\hat\rho}$, gives
\[
\mathfrak{t}_\rho(p)=\tfrac12\bigl[a_{p+\hat\rho}-a_p+E(q(0);q(1))\bigr]
-\tfrac12\bigl[a_{p-\hat\rho}-a_p+E(q(0);q(-1))\bigr].
\]
Subtracting $\tfrac12(a_{p+\hat\rho}-a_{p-\hat\rho})$, as in the definition
\eqref{10.1:eq-orbital-spin-densities-a} of $\mathfrak{r}_\rho$,
\begin{equation}\label{10.1:eq-orbital-remainder-6}
\mathfrak{r}_\rho(p)=\tfrac12\bigl[E(q(0);q(1))-E(q(0);q(-1))\bigr].
\end{equation}

We next exchange the arguments: $E(b';b)$ has base $b'$ and increments $-(b'_i-b_i)$. Fix $I$ with
$|I|=k$ and consider $T$ at $b'$ with $b'_i$ replaced by $-(b'_i-b_i)$ for $i\in I$. In each slot
$j\notin I$ write $b'_j=b_j+(b'_j-b_j)$ and expand by multilinearity: the result is $(-1)^k$ times the
sum, over all $J\supset I$, of $T$ at $b$ with $b_i$ replaced by $b'_i-b_i$ for $i\in J$. Summing over
$|I|=k$, a set $J$ with $|J|=j$ is counted once for each of its $\binom jk$ subsets of size $k$, so
$T_k(b';b)=(-1)^k\sum_{j=k}^4\binom jk T_j(b;b')$. In particular, with $T_j=T_j(b;b')$, the pair
terms become $T_2(b';b)=T_2+3T_3+6T_4$, the triple terms $T_3(b';b)=-(T_3+4T_4)$, and the quadruple
term $T_4(b';b)=T_4$. Since $\langle\cdot,\cdot\rangle$ is symmetric, $\sigma(b';b)=\sigma(b;b')=:\sigma$.
Adding,
\begin{equation}\label{10.1:eq-orbital-remainder-7}
E(b';b)=T_2+2T_3+3T_4-\sigma\,T(b').
\end{equation}
Put $E^{s}:=\tfrac12[E(b;b')+E(b';b)]$ and $E^{a}:=\tfrac12[E(b;b')-E(b';b)]$. By the definition of
$E(b;b')$ and by \eqref{10.1:eq-orbital-remainder-7},
\[
E^{s}=T_2+\tfrac32T_3+2T_4-\tfrac12\sigma\bigl(T(b)+T(b')\bigr),\qquad
E^{a}=-\tfrac12T_3-T_4-\tfrac12\sigma\bigl(T(b)-T(b')\bigr).
\]
For $(b,b')=(q(0),q(1))$ we have $T(b)=1-a_p$ and $T(b')=1-a_{p+\hat\rho}$, hence
$T(b)+T(b')=2(1-\tfrac12(a_p+a_{p+\hat\rho}))$ and $T(b)-T(b')=a_{p+\hat\rho}-a_p$; therefore $E^s$
and $E^a$ are functions of the cube between $p$ and $p+\hat\rho$, namely $E^{s}=2\Psi_\rho(p)$ and
$E^{a}=E^{a}(p)$ as in \eqref{10.1:eq-orbital-remainder-2}. The plaquette $p-\hat\rho$ is also
orthogonal to $\rho$; its links are the $\ell_i-\hat\rho$, with variables $q_i(-1)$ and transported
links $q_i(0)$. So the pair $(b,b')=(q(-1),q(0))$ gives $E^{s}=2\Psi_\rho(p-\hat\rho)$ and
$E^{a}=E^{a}(p-\hat\rho)$. Consequently, since $E(q(0);q(-1))$ is $E(b';b)=E^s-E^a$ evaluated at
the pair $(b,b')=(q(-1),q(0))$,
\begin{align*}
E(q(0);q(1))&=2\Psi_\rho(p)+E^a(p),\\
E(q(0);q(-1))&=2\Psi_\rho(p-\hat\rho)-E^a(p-\hat\rho).
\end{align*}
Inserting these into \eqref{10.1:eq-orbital-remainder-6} yields the first line of
\eqref{10.1:eq-orbital-remainder-2}.

\emph{(i).} Let $p=\{x;\rho,\nu\}$ and let $q(t)$ be the variable of the $\nu$-link
$(x+t\hat\rho,\nu)$. The boundary of $p$ consists of the two links $(x,\rho)$, $(x+\hat\nu,\rho)$,
whose variables are $1$, and of the two $\nu$-links with variables $q(0)$ and $q(1)$; hence
$\operatorname{Re}\operatorname{tr}U(\partial p)=\operatorname{tr}(q(1)\bar q(0))$, which is
$\langle q(0),q(1)\rangle$, and
$G=1-\langle q(0),q(1)\rangle$. Only the two $\nu$-links move, since $\vec u_{\ell,\rho}=0$ on links
$\ell$ in direction $\rho$: in \eqref{10.1:eq-orbital-remainder-3} the link with variable $q(0)$ moves
in the direction $\tfrac12\Pi_{q(0)}(q(1)-q(-1))$ and the link with variable $q(1)$ in the direction
$\tfrac12\Pi_{q(1)}(q(2)-q(0))$. By bilinearity of $\langle\cdot,\cdot\rangle$,
\[
\mathfrak{t}_\rho(p)=-\bigl\langle\tfrac12\Pi_{q(0)}(q(1)-q(-1)),q(1)\bigr\rangle
-\bigl\langle q(0),\tfrac12\Pi_{q(1)}(q(2)-q(0))\bigr\rangle.
\]
Put $c_t:=\langle q(t),q(t+1)\rangle$ and $e_t:=\langle q(t-1),q(t+1)\rangle$. Since
$\Pi_{q(0)}(q(1)-q(-1))=q(1)-q(-1)-(c_0-c_{-1})q(0)$ and $|q(1)|=1$,
\[
\bigl\langle\tfrac12\Pi_{q(0)}(q(1)-q(-1)),q(1)\bigr\rangle=\tfrac12\bigl[1-e_0-c_0^2+c_{-1}c_0\bigr];
\]
since $\Pi_{q(1)}(q(2)-q(0))=q(2)-q(0)-(c_1-c_0)q(1)$ and $|q(0)|=1$,
\[
\bigl\langle q(0),\tfrac12\Pi_{q(1)}(q(2)-q(0))\bigr\rangle=\tfrac12\bigl[e_1-1-c_1c_0+c_0^2\bigr].
\]
Adding and changing sign, $\mathfrak{t}_\rho(p)=-\tfrac12(e_1-e_0)+\tfrac12c_0(c_1-c_{-1})$. The same
computation of the boundary holonomy for the plaquettes $p\pm\hat\rho$ gives $c_{\pm1}=1-a_{p\pm\hat\rho}$,
so $-\tfrac12(a_{p+\hat\rho}-a_{p-\hat\rho})=\tfrac12(c_1-c_{-1})$, and
\[
\mathfrak{r}_\rho(p)=\mathfrak{t}_\rho(p)+\tfrac12(c_1-c_{-1})
=-\tfrac12(e_1-e_0)+(c_1-c_{-1})-\tfrac12a_p(c_1-c_{-1}),
\]
where $c_0=1-a_p$ was used in $\tfrac12(c_0+1)=1-\tfrac12a_p$. The boundary of $R(p)$ consists of the
$\nu$-links with variables $q(0)$, $q(2)$ and of four links in direction $\rho$, so
$e_1=\langle q(0),q(2)\rangle=1-a_{R(p)}$; likewise $e_0=1-a_{R(p-\hat\rho)}$. Moreover
$c_1-c_{-1}=a_{p-\hat\rho}-a_{p+\hat\rho}$ and
$a_p(a_{p+\hat\rho}-a_{p-\hat\rho})=a_pa_{p+\hat\rho}-a_{p-\hat\rho}a_p$. Substituting,
\[
\mathfrak{r}_\rho(p)=\tfrac12a_{R(p)}-\tfrac12a_{R(p-\hat\rho)}-a_{p+\hat\rho}+a_{p-\hat\rho}
+\tfrac12a_pa_{p+\hat\rho}-\tfrac12a_{p-\hat\rho}a_p,
\]
and this is $\Psi_\rho(p)-\Psi_\rho(p-\hat\rho)$ by the definition in
\eqref{10.1:eq-orbital-remainder-1}, the two terms $-a_p$ and $+a_p$ cancelling.

\emph{(iii).} For quaternions, $|\operatorname{tr}z|\le|z|$ (the trace is the real part) and the norm
is multiplicative, so $|T(q_1,\dots,q_4)|\le\prod_i|q_i|$. Each $a_{p'}$ lies in $[0,2]$, so
$|\delta_i|=(2a_{p(\ell_i,\rho)})^{1/2}\le(2\sigma)^{1/2}$ and $\sigma\le8$. Each term of $T_k(p)$ has
$k$ arguments $\delta_i$ and $4-k$ arguments of norm one, hence is bounded by $(2\sigma)^{k/2}$; thus
$|T_3|\le4(2\sigma)^{3/2}$ and $|T_4|\le(2\sigma)^2$. Further, by (ii) and the expansion of $T(q+\delta)$,
\begin{align*}
|a_p-a_{p+\hat\rho}|&=|T(q+\delta)-T(q)|=|T_1+T_2+T_3+T_4|\\
&\le\sum_{k=1}^4\binom4k(2\sigma)^{k/2}\le C(2\sigma)^{1/2},
\end{align*}
using $2\sigma\le16$. Hence
\[
|E^a(p)|\le\tfrac12\bigl[C\sigma(2\sigma)^{1/2}+4(2\sigma)^{3/2}+2(2\sigma)^2\bigr]\le C\sigma(p)^{3/2},
\]
where $(2\sigma)^2\le4(2\sigma)^{3/2}$ was used and $C$ denotes absolute constants. Consequently
$|\mathfrak{r}^{\sharp}_\rho(p)|\le\tfrac12C\bigl(\sigma(p)^{3/2}+\sigma(p-\hat\rho)^{3/2}\bigr)$, and
$\sigma(p)$, $\sigma(p-\hat\rho)$ are sums of densities of plaquettes containing the direction $\rho$.

\emph{(iv).} Consider $\Psi_\rho$ on $p$ orthogonal to $\rho$. Since
\[
1-\tfrac12(a_p+a_{p+\hat\rho})=\tfrac12\bigl(T(q)+T(q+\delta)\bigr)=T(q)+\tfrac12(T_1+T_2+T_3+T_4),
\]
we have
\[
\Psi_\rho(p)=\tfrac12\bigl(T_2-\sigma T(q)\bigr)+\tfrac34T_3+T_4-\tfrac14\sigma(T_1+T_2+T_3+T_4),
\]
and, as $\sigma=\tfrac12\sum_i|\delta_i|^2$, each term after the first contains at least three
increments $\delta_i$; for a smooth field each $\delta_i$ is of order $a^2$, because
$|\delta_i|^2=2a_{p(\ell_i,\rho)}$ and the plaquette densities are of order $a^4$. For smooth fields
the increments of opposite links nearly cancel: up to terms of order $a^6$ or higher,
\[
T_2-\sigma T(q)=-\tfrac12\bigl|\delta_1+\delta_2-\delta_3-\delta_4\bigr|^2
=-\tfrac12\bigl|(\delta_1-\delta_3)+(\delta_2-\delta_4)\bigr|^2,
\]
the signs being those with which $\ell_1,\ell_2$ and $\ell_3,\ell_4$ enter $U(\partial p)$; here
$\delta_1-\delta_3$ and $\delta_2-\delta_4$ are differences of the increments of the opposite, parallel
links $\ell_1,\ell_3$ and $\ell_2,\ell_4$, each of order $a^3$, so $T_2-\sigma T(q)=O(a^6)$. Hence
$\Psi_\rho=O(a^6)$.

For $\mathfrak{r}^{\sharp}_\rho$: it is a vector density, so its classical expansion has odd dimension;
by (iii) it is at least cubic in increments of order $a^2$, and its term of order $a^6$, of even
dimension, is absent, so $\mathfrak{r}^{\sharp}_\rho=O(a^7)$. At constant $F$ this is seen on $T_3$:
with $F_\mu:=F_{\rho\mu}$ and $F_\nu:=F_{\rho\nu}$, the $a^6$-term of $T_3$ is
\[
a^6\bigl[-\operatorname{tr}F_\mu F_\nu F_\mu-\operatorname{tr}F_\mu F_\nu F_\nu
+\operatorname{tr}F_\mu F_\mu F_\nu+\operatorname{tr}F_\nu F_\mu F_\nu\bigr]=0,
\]
since $\operatorname{tr}F_\mu F_\nu F_\mu=\operatorname{tr}F_\mu F_\mu F_\nu$ and
$\operatorname{tr}F_\mu F_\nu F_\nu=\operatorname{tr}F_\nu F_\mu F_\nu$ by cyclicity of the trace.

Finally let the field be constant and abelian and let $p$ contain the direction $\rho$. Then $p$ and
$p+\hat\rho$ have holonomies with the same angle $\theta$, so $a_p=a_{p+\hat\rho}=1-\cos\theta$, and the
holonomy of $R(p)$ is the product of the two, with angle $2\theta$, so
$a_{R(p)}=1-\cos2\theta=2(1-\cos\theta)(1+\cos\theta)$. By \eqref{10.1:eq-orbital-remainder-1},
\[
\Psi_\rho(p)=(1-\cos\theta)\bigl[(1+\cos\theta)-2+\tfrac12(1-\cos\theta)\bigr]
=-\tfrac12(1-\cos\theta)^2,
\]
which is quadratic in $a_p=1-\cos\theta$ and vanishes to fourth order in $\theta$.
\end{proof}

\begin{corollary}[The sharp form of the rotational identity]\label{10.1:cor-sharp-identity}
Let the data be admissible in the sense of Theorem~\ref{10.1:thm-rotational-identity}. Define the
forward lattice difference
\[
(\nabla^+_\rho\varphi)(y) := a^{-1}\bigl[\varphi(y+\hat\rho)-\varphi(y)\bigr].
\]
Write $\Psi_\rho(p)$ and $\mathfrak{r}^{\sharp}_\rho(p)$ for the functions of
Lemma~\ref{10.1:lem-orbital-remainder}. There $\Psi_\rho(p)$ is given by part~(i) when $p$
contains the direction $\rho$ and by part~(ii) otherwise, and $\mathfrak{r}^{\sharp}_\rho(p)=0$
when $p$ contains $\rho$. Then the orbital remainder $\mathfrak{K}_\chi[c]$ of
\eqref{10.1:eq-rotational-identity-a} satisfies
\begin{equation}\label{10.1:eq-sharp-identity-a}
\begin{aligned}
\mathfrak{K}_\chi[c] &= \mathfrak{K}^{\sharp}_\chi[c]+\mathfrak{K}^{\partial}_\chi[c],\\
\mathfrak{K}^{\sharp}_\chi[c] &:= \sum_p c(p)\chi(x_p)\sum_\rho a^{-1}\xi^\rho(x_p)\,
\mathfrak{r}^{\sharp}_\rho(p),\\
\mathfrak{K}^{\partial}_\chi[c] &:= -\sum_p\sum_\rho \xi^\rho(x_p)\,
\nabla^+_\rho(c\chi)(x_p)\,\Psi_\rho(p)\\
&\phantom{:}= \sum_p\sum_\rho \xi^\rho(x_p)\,(\nabla^+_\rho w)(x_p)\,\Psi_\rho(p)\\
&\quad -g_0^{-2}\sum_p\sum_\rho \xi^\rho(x_p)\,(\nabla^+_\rho\chi)(x_p)\,\Psi_\rho(p).
\end{aligned}
\end{equation}
Here a function of plaquettes is read as a function on their barycentres, and
$x_{p+\hat\rho}=x_p+\hat\rho$.

Thus in the bulk, where $c$ and $\chi$ are constant, the orbital remainder is the cubic
functional $\mathfrak{K}^{\sharp}$. The quadratic part of the orbital remainder
$\mathfrak{K}_\chi[c]$ consists of two terms:
\begin{itemize}
\item a contact term with the bounded weight $\xi\cdot\nabla w$;
\item a term supported in the shell, classically $O(a^2)$ relative to
$\mathfrak{B}^{\mathrm{sh}}_\chi$.
\end{itemize}

Now consider the insertions of \eqref{10.1:eq-rotational-identity-b}. Let $f$ be one of the
$f_i$, let $\mathfrak{K}[f]$ be as in \eqref{10.1:eq-rotational-identity-b}, and let
$\mathfrak{K}^{\sharp}[f]$ be obtained from $\mathfrak{K}^{\sharp}_\chi[c]$ by replacing
$c(p)\chi(x_p)$ with $f(x_p)$. Since $\chi\equiv1$ on the $2a$-neighbourhood of
$\operatorname{supp}f$, keeping or dropping a factor $\chi(x_p)$ beside $f(x_p)$ changes nothing.
For them,
\[
\mathfrak{K}[f] = \mathfrak{K}^{\sharp}[f]
-\sum_p\sum_\rho \xi^\rho(x_p)\,(\nabla^+_\rho f)(x_p)\,\Psi_\rho(p).
\]
\end{corollary}

\begin{proof}
By parts (i), (ii) and (iii) of Lemma~\ref{10.1:lem-orbital-remainder}, every plaquette $p$
and every direction $\rho$ satisfy
\[
\mathfrak{r}_\rho(p)=\Psi_\rho(p)-\Psi_\rho(p-\hat\rho)+\mathfrak{r}^{\sharp}_\rho(p).
\]
When $p$ contains $\rho$ this is part~(i) together with the convention
$\mathfrak{r}^{\sharp}_\rho(p)=0$; otherwise it is part~(ii).

Put $\varphi_\rho(p):=a^{-1}c(p)\chi(x_p)\xi^\rho(x_p)$. This is a function on the plaquettes
of the torus, because $\chi\xi$ is a function on the torus by admissibility. Inserting the
identity above into the definition of $\mathfrak{K}_\chi[c]$ in
\eqref{10.1:eq-rotational-identity-a} gives
\[
\mathfrak{K}_\chi[c]=\mathfrak{K}^{\sharp}_\chi[c]
+\sum_\rho\sum_p\varphi_\rho(p)\bigl[\Psi_\rho(p)-\Psi_\rho(p-\hat\rho)\bigr].
\]

The translation $p\mapsto p+\hat\rho$ is a bijection of the plaquettes of the periodic lattice.
Hence $\sum_p\varphi_\rho(p)\Psi_\rho(p-\hat\rho)=\sum_p\varphi_\rho(p+\hat\rho)\Psi_\rho(p)$,
and so
\[
\sum_p\varphi_\rho(p)\bigl[\Psi_\rho(p)-\Psi_\rho(p-\hat\rho)\bigr]
=-a\sum_p(\nabla^+_\rho\varphi_\rho)(x_p)\,\Psi_\rho(p).
\]

Since $\xi^\rho(x)=\sum_\sigma\omega_{\rho\sigma}x_\sigma$ and $\omega_{\rho\rho}=0$, the
component $\xi^\rho$ does not depend on $x_\rho$. On the torus this gives
$\xi^\rho(x_p+\hat\rho)=\xi^\rho(x_p)$ whenever $\chi(x_p)$ or $\chi(x_p+\hat\rho)$ is non-zero,
since both points then lie within distance $a$ of
$\operatorname{supp}\chi\subset\,]-L^m+2a,L^m-2a[^4$, and so inside the fundamental cube; when
both values of $\chi$ vanish, both sides of the next identity are zero. Therefore
\[
(\nabla^+_\rho\varphi_\rho)(x_p)=a^{-1}\xi^\rho(x_p)\,\nabla^+_\rho(c\chi)(x_p).
\]
Multiplying by $-a$ gives the first expression for $\mathfrak{K}^{\partial}_\chi[c]$.

For the second expression, write $c\chi=g_0^{-2}\chi-w\chi$. By admissibility, $\chi(x_q)=1$
whenever $w(q)\neq0$. Hence $(w\chi)(x_q)=w(q)$ for every plaquette $q$, and
$\nabla^+_\rho(w\chi)(x_p)=(\nabla^+_\rho w)(x_p)$. Consequently
\[
-\nabla^+_\rho(c\chi)=\nabla^+_\rho w-g_0^{-2}\nabla^+_\rho\chi
\quad\text{at every } x_p.
\]

The same summation by parts with $\varphi_\rho(p)=a^{-1}f(x_p)\chi(x_p)\xi^\rho(x_p)$ gives the
identity for $\mathfrak{K}[f]$ stated above, which is the form entering
\eqref{10.1:eq-rotational-identity-b}. There $\chi\equiv1$ on the
$2a$-neighbourhood of $\operatorname{supp}f$, so $(f\chi)(x_q)=f(x_q)$ for every $q$, and
$\nabla^+_\rho(f\chi)(x_p)=(\nabla^+_\rho f)(x_p)$.

In the non-compact Gaussian $U(1)$ model the cubic piece $E^a(p)$ of
Lemma~\ref{10.1:lem-orbital-remainder} vanishes identically.
Identity~\eqref{10.1:eq-sharp-identity-a} is the non-abelian form of the abelian check of the
structure in the remarks following Theorem~\ref{10.1:thm-rotational-identity}.
\end{proof}

In what follows, the continuum connection is $A=(A_\mu)_{\mu=1}^{4}$, with every $A_\mu$ anti-Hermitian. We write
\begin{equation}\label{10.1:eq-classical-limit-1}
D_\mu:=\partial_\mu+[A_\mu,\cdot\,],\qquad
F_{\mu\nu}:=\partial_\mu A_\nu-\partial_\nu A_\mu+[A_\mu,A_\nu],\qquad
J_\nu:=\sum_\mu D_\mu F_{\mu\nu}.
\end{equation}
Sums over repeated indices are taken only where a summation sign is written, or where no index
occurs more than twice. The same letter also denotes the action
$A(U)=\sum_p a_p(U)$ of a lattice configuration $U$; the argument distinguishes the two uses. We put
\begin{equation}\label{10.1:eq-classical-limit-2}
O_1:=\sum_{\mu,\nu}\operatorname{tr}\bigl(D_\mu F_{\mu\nu}\,D_\mu F_{\mu\nu}\bigr),\qquad
B_{\mu\rho}:=2\sum_\nu\operatorname{tr}\bigl[D_\mu F_{\mu\nu}\bigl(D_\rho F_{\mu\nu}
+D_\mu F_{\rho\nu}\bigr)\bigr],
\end{equation}
with no summation over $\mu,\rho$ in $B_{\mu\rho}$. A link $\ell=(x,\mu)$ is the link from $x$ to
$x+\hat\mu$, with midpoint $x_\ell=x+\tfrac12\hat\mu$. The plaquette with base point $x$ spanned by
$\hat\mu$ and $\hat\nu$ is written $p=\{x;\mu,\nu\}$. The matrix $\omega$ is real antisymmetric and
$\xi(x)=\omega x$, as in the rotational Ward identity (Theorem~\ref{10.1:thm-rotational-identity}),
so that $\partial_\mu\xi^\rho=\omega_{\rho\mu}$ and $\partial_\mu\partial_\nu\xi^\rho=0$.

\begin{proposition}[The classical limit of the breaking term]\label{10.1:prop-classical-limit}
Let $A\in C^\infty_c(\mathbb{R}^4)$ be a connection as above. On $\mathfrak{L}$, with $m$ large, let
$U$ be the configuration $U_\ell:=\mathrm{P}\exp\int_\ell A$. Then, as $a\to0$, the following hold.
\begin{enumerate}
\item[(i)] The action of $U$ satisfies
\begin{equation}\label{10.1:eq-classical-limit-3}
A(U)=\sum_p a_p=\int\Bigl[-\frac14\sum_{\mu,\nu}\operatorname{tr}F_{\mu\nu}F_{\mu\nu}\Bigr]
+\frac{a^2}{24}\int O_1+O(a^4).
\end{equation}
\item[(ii)] Let $\ell=(x,\mu)$ and let $\vec v_\ell$ be the vector of
Definition~\ref{10.1:def-displacement-derivation} for the displacement $\eta=\xi$. Read both sides
in the frame at $x_\ell$. Then
\begin{equation}\label{10.1:eq-classical-limit-4}
\begin{split}
i\,\vec v_\ell\cdot\vec\sigma
={}&\bigl(\text{the left-logarithmic derivative of }U_\ell\text{ along }
\delta A_\mu=\textstyle\sum_\rho\xi^\rho F_{\rho\mu}\bigr)\\
&+a^3\Bigl[\frac16\sum_\rho\xi^\rho D_\rho^2F_{\rho\mu}
-\frac1{12}\sum_\rho\omega_{\rho\mu}D_\mu F_{\rho\mu}\Bigr](x_\ell)+O(a^5).
\end{split}
\end{equation}
\item[(iii)] With the densities of \eqref{10.1:eq-orbital-spin-densities-a},
\begin{equation}\label{10.1:eq-classical-limit-a}
\begin{split}
\sum_p\Bigl[a^{-1}\sum_\rho\xi^\rho(x_p)\mathfrak{r}_\rho(p)
+\sum_{\rho,\sigma}\omega_{\rho\sigma}\mathfrak{s}_{\sigma\rho}(p)\Bigr]
={}&V_\xi A(U)\\
={}&a^2\int\Bigl\{\frac1{24}\sum_{\mu,\rho}\omega_{\rho\mu}B_{\mu\rho}\\
&\qquad+\sum_\mu\operatorname{tr}\Bigl(J_\mu\Bigl[\frac16\sum_\rho\xi^\rho D_\rho^2F_{\rho\mu}
-\frac1{12}\sum_\rho\omega_{\rho\mu}D_\mu F_{\rho\mu}\Bigr]\Bigr)\Bigr\}\\
&+O(a^4).
\end{split}
\end{equation}
Let $\delta_\xi$ denote the variation of a local expression in $A$ under the rotation generated
by $\omega$, namely the first-order change under $A_\mu\mapsto A_\mu+t\sum_\rho\xi^\rho F_{\rho\mu}$.
Then $\sum_{\mu,\rho}\omega_{\rho\mu}B_{\mu\rho}=\delta_\xi O_1-\xi\cdot\partial O_1$. Consequently
the classical breaking term, the left-hand side of \eqref{10.1:eq-classical-limit-a}, is
$\frac{a^2}{24}\,\delta_\xi\int O_1$ plus two multiples of the field equation $J=0$.
\end{enumerate}
\end{proposition}

\begin{proof}
\emph{(i).} Fix a plaquette $p=\{x;\mu,\nu\}$. Since $a_p(U)$ is gauge invariant, we may compute
in the Fock--Schwinger gauge at $x_p$, that is, $\sum_\alpha y_\alpha A_\alpha(x_p+y)=0$. In this
gauge
\[
A_\mu(x_p+y)=\tfrac12\sum_\nu y_\nu F_{\nu\mu}(x_p)+O(|y|^2),\qquad
F(x_p+y)=\sum_{n\ge0}(n!)^{-1}(y\cdot D)^nF(x_p).
\]
The four links of $\partial p$ lie at distance $a/2$ from $x_p$ and have length $a$. Their
exponents are therefore $\frac{a^2}{4}F_{\mu\nu}(x_p)+O(a^3)$, each taken with the orientation of
$\partial p$. In particular they commute to leading order.

Commutators of the four exponents enter $\log U(\partial p)$ at order $a^5$ and at order $a^6$.
The terms of order $a^5$ are odd, and they cancel in the trace by the inversion $y\mapsto-y$
about $x_p$. In the trace, commutators enter either through
$\operatorname{tr}\bigl(F_{\mu\nu}(x_p)[\cdot,\cdot]\bigr)$ multiplied by $a^8$, or through the trace
of a commutator, which is zero.

The sum of the four exponents is, by Stokes' formula, the integral over the square of
$\partial_\mu A_\nu-\partial_\nu A_\mu=F_{\mu\nu}-[A_\mu,A_\nu]$. In the Fock--Schwinger gauge the
leading order of $[A_\mu,A_\nu]$ is
$\frac14 y_\nu y_\mu[F_{\nu\mu}(x_p),F_{\mu\nu}(x_p)]=0$, since $F_{\nu\mu}=-F_{\mu\nu}$; its next
terms are odd in $y$. The integral of $F_{\mu\nu}(x_p+y)$ over the square
$x_p+[-a/2,a/2]^2$ in the $(\mu\nu)$-plane is $a^2$ times the average of the expansion above:
\begin{itemize}
\item the term of first order in $y$ averages to zero;
\item the average of $\frac12(y\cdot D)^2F_{\mu\nu}$ is
$\frac{a^2}{24}(D_\mu^2+D_\nu^2)F_{\mu\nu}(x_p)$, because $y_\mu^2$ and $y_\nu^2$ have mean
$a^2/12$ and $y_\mu y_\nu$ has mean zero.
\end{itemize}
For an anti-Hermitian $X$ one has $\operatorname{Re}\operatorname{tr}X=0$ and
$\operatorname{Re}\operatorname{tr}X^3=0$. Hence $\operatorname{Re}\operatorname{tr}e^X
=1+\frac12\operatorname{tr}X^2+O(|X|^4)$. With $\log U(\partial p)=O(a^2)$ this gives
\begin{equation}\label{10.1:eq-classical-limit-5}
\operatorname{Re}\operatorname{tr}U(\partial p)
=1+\frac12\operatorname{tr}\Bigl(a^2F_{\mu\nu}+\frac{a^4}{24}(D_\mu^2+D_\nu^2)F_{\mu\nu}\Bigr)^2
+O(a^8),
\end{equation}
all fields being evaluated at $x_p$. Therefore
\[
a_p=-\frac{a^4}{2}\operatorname{tr}F_{\mu\nu}F_{\mu\nu}
-\frac{a^6}{24}\operatorname{tr}F_{\mu\nu}(D_\mu^2+D_\nu^2)F_{\mu\nu}+O(a^8),
\]
with no summation over $\mu,\nu$. The remainder is uniform in $p$. Moreover $a_p=0$ unless $p$
meets the support of $A$, so only $O(a^{-4})$ plaquettes contribute, and the remainders add up to
$O(a^4)$.

For fixed $\mu<\nu$, the barycentres of the plaquettes $\{x;\mu,\nu\}$ form a translate of
$a\mathbb{Z}^4$; for $m$ large the support of $A$ lies inside one period of the torus. The sum,
over such a translate, of a smooth compactly supported function multiplied by $a^4$ equals its
integral up to $O(a^\infty)$. Thus
\[
\sum_p a_p=\sum_{\mu<\nu}\int\Bigl[-\frac12\operatorname{tr}F_{\mu\nu}F_{\mu\nu}
-\frac{a^2}{24}\operatorname{tr}F_{\mu\nu}(D_\mu^2+D_\nu^2)F_{\mu\nu}\Bigr]+O(a^4).
\]
Since $\partial_\mu\operatorname{tr}(XY)=\operatorname{tr}(D_\mu X\,Y)+\operatorname{tr}(X\,D_\mu Y)$,
we have, modulo a total derivative whose integral vanishes,
\[
-\operatorname{tr}F_{\mu\nu}(D_\mu^2+D_\nu^2)F_{\mu\nu}
\simeq\operatorname{tr}(D_\mu F_{\mu\nu})^2+\operatorname{tr}(D_\nu F_{\mu\nu})^2 .
\]
Summing over $\mu<\nu$ and using $F_{\mu\nu}=-F_{\nu\mu}$ gives
$\sum_{\mu\ne\nu}\operatorname{tr}(D_\mu F_{\mu\nu})^2=O_1$; the terms with $\mu=\nu$ vanish.
Also $\sum_{\mu<\nu}\bigl(-\frac12\operatorname{tr}F_{\mu\nu}F_{\mu\nu}\bigr)
=-\frac14\sum_{\mu,\nu}\operatorname{tr}F_{\mu\nu}F_{\mu\nu}$. This is
\eqref{10.1:eq-classical-limit-3}.

\emph{(ii).} We first compute the continuum variation. Put
$\psi:=\sum_\rho\xi^\rho F_{\rho\mu}$. The left-logarithmic derivative of $U_\ell$ along
$\delta A_\mu=\psi$, read in the frame at $x_\ell$, is $\int_{-a/2}^{a/2}ds$ of $\psi$ at
$x_\ell+se_\mu$, parallel transported to $x_\ell$ along $\ell$. By the covariant Taylor expansion
the transported field is $\sum_n(s^n/n!)D_\mu^n\psi(x_\ell)$. The odd powers of $s$ integrate to
zero, and $\int_{-a/2}^{a/2}\frac{s^2}{2}\,ds=\frac{a^3}{24}$. The variation is therefore
\begin{gather*}
a\Bigl[\psi+\frac{a^2}{24}D_\mu^2\psi\Bigr](x_\ell)+O(a^5),\\
D_\mu^2\psi=\sum_\rho\xi^\rho D_\mu^2F_{\rho\mu}+2\sum_\rho\omega_{\rho\mu}D_\mu F_{\rho\mu},
\end{gather*}
where the second identity uses $\partial_\mu\xi^\rho=\omega_{\rho\mu}$ and $\partial_\mu^2\xi^\rho=0$.

Next we compute the lattice field. The plaquette $p(\ell,\pm\rho)$ spanned by $\ell$ and
$\pm\hat\rho$ has barycentre $x_\ell\pm\frac12\hat\rho$, and its holonomy $W_{\ell,\pm\rho}$ has
orientation given by $\hat\mu$ followed by $\pm\hat\rho$. Put
\[
\Phi:=F_{\mu\rho}+\frac{a^2}{24}(D_\mu^2+D_\rho^2)F_{\mu\rho}.
\]
By the computation of (i), the anti-Hermitian part $i\vec w_{\ell,\pm\rho}\cdot\vec\sigma$ of
$W_{\ell,\pm\rho}$ equals, to the order kept, $\pm a^2\Phi$ evaluated at the barycentre. Here the
anti-Hermitian part
of $W=w^0+i\vec w\cdot\vec\sigma$ is $i\vec w\cdot\vec\sigma$. We transport this to $x_\ell$. The
triangle between the two transport paths lies in the $(\mu\rho)$-plane, so its holonomy commutes
with $F_{\mu\rho}$ to the order kept. The covariant Taylor expansion in the direction $\rho$ then
gives, to the order kept,
\[
i\vec w_{\ell,\pm\rho}\cdot\vec\sigma
=\pm a^2\Bigl[\Phi\pm\frac a2D_\rho\Phi+\frac{a^2}{8}D_\rho^2\Phi\Bigr](x_\ell).
\]
By Definition~\ref{10.1:def-displacement-derivation},
$\vec u_{\ell,\rho}=-\frac12(\vec w_{\ell,+\rho}-\vec w_{\ell,-\rho})$ for $\rho\ne\mu$. The terms
with $D_\rho\Phi$ cancel, and
\[
\begin{split}
i\vec u_{\ell,\rho}\cdot\vec\sigma
&=-a^2\Bigl[\Phi+\frac{a^2}{8}D_\rho^2\Phi\Bigr]\\
&=a^2\Bigl[F_{\rho\mu}+\frac{a^2}{24}(D_\mu^2+D_\rho^2)F_{\rho\mu}
+\frac{a^2}{8}D_\rho^2F_{\rho\mu}\Bigr]+O(a^6)\\
&=a^2\Bigl[F_{\rho\mu}+\frac{a^2}{24}D_\mu^2F_{\rho\mu}+\frac{a^2}{6}D_\rho^2F_{\rho\mu}\Bigr](x_\ell)
+O(a^6),
\end{split}
\]
since $\frac1{24}+\frac18=\frac16$. The same formula holds for $\rho=\mu$, because
$\vec u_{\ell,\mu}=0$ and $F_{\mu\mu}=0$; likewise $\omega_{\mu\mu}=0$. The vector of
Definition~\ref{10.1:def-displacement-derivation} is
$\vec v_\ell=a^{-1}\sum_\rho\xi^\rho(x_\ell)\vec u_{\ell,\rho}$. Hence
\[
i\vec v_\ell\cdot\vec\sigma
=a\sum_\rho\xi^\rho\Bigl[F_{\rho\mu}+\frac{a^2}{24}D_\mu^2F_{\rho\mu}
+\frac{a^2}{6}D_\rho^2F_{\rho\mu}\Bigr](x_\ell)+O(a^5).
\]
Subtracting the continuum variation computed above and using $\frac{2}{24}=\frac1{12}$ gives
\eqref{10.1:eq-classical-limit-4}.

\emph{(iii), first equality.} This is the computation in the proof of the rotational Ward
identity (Theorem~\ref{10.1:thm-rotational-identity}), specialised to $\chi\equiv1$ and
$c\equiv1$. With these choices that computation gives $V_\xi A(U)$ through the elementary
densities $\mathfrak{d}_{\ell,p,\rho}$, in the form
\[
V_\xi A(U)=a^{-1}\sum_p\sum_{\ell\subset\partial p}\sum_\rho\xi^\rho(x_\ell)\,
\mathfrak{d}_{\ell,p,\rho}.
\]
We write
\[
\xi^\rho(x_\ell)=\xi^\rho(x_p)+\sum_\sigma\omega_{\rho\sigma}(x_\ell-x_p)_\sigma .
\]
By the definitions \eqref{10.1:eq-orbital-spin-densities-a}, this splits $V_\xi A(U)$ into
\[
\sum_p\Bigl[a^{-1}\sum_\rho\xi^\rho(x_p)\mathfrak{t}_\rho(p)
+\sum_{\rho,\sigma}\omega_{\rho\sigma}\mathfrak{s}_{\sigma\rho}(p)\Bigr].
\]
Then we insert $\mathfrak{t}_\rho(p)=\mathfrak{r}_\rho(p)+\frac12(a_{p+\hat\rho}-a_{p-\hat\rho})$.

The orbital main part telescopes to zero. Indeed, $U$ is trivial far out, so only finitely many
plaquettes carry non-zero $a_p$, and no boundary terms arise. Shifting the summation variable
gives
\[
\begin{aligned}
\sum_p a^{-1}\sum_\rho\xi^\rho(x_p)\tfrac12\bigl(a_{p+\hat\rho}-a_{p-\hat\rho}\bigr)
&=\tfrac12a^{-1}\sum_p a_p\sum_\rho\bigl(\xi^\rho(x_p-\hat\rho)-\xi^\rho(x_p+\hat\rho)\bigr)\\
&=-\sum_p a_p\sum_\rho\omega_{\rho\rho}\\
&=0,
\end{aligned}
\]
because $\omega$ is antisymmetric. This is the first equality in
\eqref{10.1:eq-classical-limit-a}.

\emph{(iii), second equality.} Write
\[
\delta^{(2)}A_\mu:=\frac16\sum_\rho\xi^\rho D_\rho^2F_{\rho\mu}
-\frac1{12}\sum_\rho\omega_{\rho\mu}D_\mu F_{\rho\mu}.
\]
The derivation $V_\xi$ acts through the link vectors $i\vec v_\ell\cdot\vec\sigma$. The
derivatives of the gauge-invariant function $A(\cdot)$ along gauge directions vanish. By (ii)
we therefore have
\[
V_\xi A(U[A])=\frac{d}{dt}\Big|_{t=0}A\bigl(U[A+t\,\delta A]\bigr),\qquad
\delta A_\mu=\sum_\rho\xi^\rho F_{\rho\mu}+a^2\delta^{(2)}A_\mu+O(a^4).
\]
We apply (i) to the family $A+t\,\delta A$ and differentiate at $t=0$.

\emph{The part $\sum_\rho\xi^\rho F_{\rho\mu}$ of $\delta A$.} This part gives
$\delta_\xi\int\bigl(-\frac14\sum\operatorname{tr}F^2\bigr)+\frac{a^2}{24}\delta_\xi\int O_1$.
The Lie derivative of $A$ along $\xi$ is
\[
\sum_\rho\bigl(\xi^\rho\partial_\rho A_\mu+\omega_{\rho\mu}A_\rho\bigr)
=\sum_\rho\xi^\rho F_{\rho\mu}+D_\mu\Bigl(\sum_\rho\xi^\rho A_\rho\Bigr).
\]
Indeed, expanding $F_{\rho\mu}$ and using $\partial_\mu\xi^\rho=\omega_{\rho\mu}$, the right-hand
side equals
$\sum_\rho\bigl(\xi^\rho\partial_\rho A_\mu+\omega_{\rho\mu}A_\rho\bigr)$ plus the terms
$-\xi^\rho\partial_\mu A_\rho+\xi^\rho\partial_\mu A_\rho$ and
$\xi^\rho[A_\rho,A_\mu]+\xi^\rho[A_\mu,A_\rho]$, which cancel in pairs.
So $\delta_\xi$ differs from
the Lie derivative by an infinitesimal gauge transformation. That transformation acts on
covariant fields by commutators, whose contributions to traces of products vanish by cyclicity of
the trace. For $X_{\alpha\beta\gamma}:=D_\alpha F_{\beta\gamma}$ we therefore have
\[
(\delta_\xi X)_{\alpha\beta\gamma}=\xi\cdot\partial X_{\alpha\beta\gamma}
+\sum_\rho\bigl(\omega_{\rho\alpha}X_{\rho\beta\gamma}+\omega_{\rho\beta}X_{\alpha\rho\gamma}
+\omega_{\rho\gamma}X_{\alpha\beta\rho}\bigr),
\]
modulo such commutators. In $\delta_\xi\operatorname{tr}(X_{\mu\mu\nu}X_{\mu\mu\nu})$, summed over
$\mu,\nu$, the term in which $\omega$ acts on the last index is
$2\sum_{\mu,\nu,\rho}\omega_{\rho\nu}\operatorname{tr}(X_{\mu\mu\nu}X_{\mu\mu\rho})$. It vanishes,
since it pairs an expression symmetric in $(\nu,\rho)$ with the antisymmetric $\omega$. The
remaining terms give
\[
\begin{split}
\delta_\xi O_1&=\xi\cdot\partial O_1+2\sum_{\mu,\rho}\omega_{\rho\mu}\sum_\nu
\operatorname{tr}\bigl[D_\mu F_{\mu\nu}\bigl(D_\rho F_{\mu\nu}+D_\mu F_{\rho\nu}\bigr)\bigr]\\
&=\xi\cdot\partial O_1+\sum_{\mu,\rho}\omega_{\rho\mu}B_{\mu\rho},
\end{split}
\]
which is the asserted identity. Since $\sum_\rho\partial_\rho\xi^\rho=\sum_\rho\omega_{\rho\rho}=0$,
we get $\int\xi\cdot\partial O_1=0$, and so
$\delta_\xi\int O_1=\int\sum_{\mu,\rho}\omega_{\rho\mu}B_{\mu\rho}$.

The same computation applies to
$\sum_{\mu,\nu}\operatorname{tr}F_{\mu\nu}F_{\mu\nu}$. Its variation is $\xi\cdot\partial$ of it
plus $2\sum_{\mu,\rho}\omega_{\rho\mu}\sum_\nu\operatorname{tr}F_{\mu\nu}F_{\rho\nu}$ plus the
analogous term in the second index. Both $\omega$-terms vanish by the symmetry in $(\mu,\rho)$,
and the integral of $\xi\cdot\partial$ vanishes as before. Hence
$\delta_\xi\int\bigl(-\frac14\sum\operatorname{tr}F^2\bigr)=0$, and this part contributes
$0+\frac{a^2}{24}\int\sum_{\mu,\rho}\omega_{\rho\mu}B_{\mu\rho}$.

\emph{The part $a^2\delta^{(2)}A$ of $\delta A$.} For any variation $\delta A$,
\[
\delta\int\Bigl(-\frac14\sum_{\mu,\nu}\operatorname{tr}F_{\mu\nu}F_{\mu\nu}\Bigr)
=-\int\sum_{\mu,\nu}\operatorname{tr}\bigl(F_{\mu\nu}D_\mu\delta A_\nu\bigr)
=\int\sum_\nu\operatorname{tr}\bigl(J_\nu\,\delta A_\nu\bigr).
\]
The first equality uses $\delta F_{\mu\nu}=D_\mu\delta A_\nu-D_\nu\delta A_\mu$ and the
antisymmetry of $F$; the second integrates by parts. The part $a^2\delta^{(2)}A$ of $\delta A$
therefore gives $a^2\int\sum_\mu\operatorname{tr}(J_\mu\,\delta^{(2)}A_\mu)$. Its contribution
through the term $\frac{a^2}{24}\int O_1$ of (i) is of order $a^4$.

Adding the two parts gives the second equality in \eqref{10.1:eq-classical-limit-a}. The
coefficients of $J$ there are the two multiples $\frac16\sum_\rho\xi^\rho D_\rho^2F_{\rho\mu}$ and
$-\frac1{12}\sum_\rho\omega_{\rho\mu}D_\mu F_{\rho\mu}$, and the first term is
$\frac{a^2}{24}\delta_\xi\int O_1$.
\end{proof}

\begin{remark}[Leading orders of the pieces]
Let $U$ be as in Proposition~\ref{10.1:prop-classical-limit}, and write
$\vec F\cdot\vec G:=-\operatorname{tr}FG$. For $p=\{x;\mu,\nu\}$,
\[
\mathfrak{s}_{\nu\rho}(p)=a^4\vec F_{\rho\mu}\cdot\vec F_{\nu\mu}+O(a^6),\qquad
\mathfrak{s}_{\mu\rho}(p)=a^4\vec F_{\rho\nu}\cdot\vec F_{\mu\nu}+O(a^6),
\]
with no summation. Consequently, the sum of $\mathfrak{s}_{\sigma\rho}$ over the plaquettes based
at $x$, namely $\sum_{\mu<\nu}\mathfrak{s}_{\sigma\rho}(\{x;\mu,\nu\})$, is, to leading order,
$a^4\sum_\mu\vec F_{\rho\mu}\cdot\vec F_{\sigma\mu}$. This is symmetric
in $(\sigma,\rho)$, so $\sum_{\rho,\sigma}\omega_{\rho\sigma}\mathfrak{s}_{\sigma\rho}$ starts at
order $a^6$.

Moreover $\mathfrak{r}_\rho(p)=O(a^7)$. Finally, the shell term $\mathfrak{B}^{\mathrm{sh}}_\chi$ of
Theorem~\ref{10.1:thm-rotational-identity} tends to
$-\int\sum_{\rho,\sigma}\xi^\rho\,\partial_\sigma\chi\,\Theta_{\sigma\rho}$, where
\[
\Theta_{\sigma\rho}:=g_0^{-2}\Bigl(\sum_\mu\vec F_{\sigma\mu}\cdot\vec F_{\rho\mu}
-\frac14\delta_{\sigma\rho}\sum_{\alpha,\beta}\vec F_{\alpha\beta}\cdot\vec F_{\alpha\beta}\Bigr),
\]
plus the Jacobian term $\operatorname{div}V_{\chi\xi}$, which carries no factor $g_0^{-2}$.
\end{remark}

\begin{definition}[Gauge-invariant local operators as hyperoctahedral modules]
\label{10.1:def-local-operators}
Let $A_\mu$ be a smooth anti-Hermitian connection with values in $\mathfrak{su}(2)$, let
$D_\mu = \partial_\mu + [A_\mu,\cdot\,]$ and
$F_{\mu\nu} = \partial_\mu A_\nu - \partial_\nu A_\mu + [A_\mu,A_\nu]$, and let
$J_\nu := \sum_\mu D_\mu F_{\mu\nu}$. The trace $\operatorname{tr}$ is normalised by
$\operatorname{tr}1 = 1$, where $1$ is the identity matrix. Assign the dimensions
$[F_{\mu\nu}] = 2$ and $[D_\alpha] = 1$.
\begin{enumerate}
\item The \emph{jet coordinates} are the symmetrised covariant derivatives
\begin{equation}\label{10.1:eq-local-operators-1}
D_{(\alpha_1}\cdots D_{\alpha_k)}F_{\mu\nu}, \qquad k \ge 0,
\end{equation}
of dimension $k+2$. Polynomials in the components of $D_{\alpha_1}\cdots D_{\alpha_k}F_{\mu\nu}$ are
reduced to polynomials in \eqref{10.1:eq-local-operators-1} by the antisymmetry
$F_{\mu\nu} = -F_{\nu\mu}$, the Bianchi identity $D_{[\alpha}F_{\mu\nu]} = 0$ and the commutator
identity $[D_\alpha,D_\beta] = \operatorname{ad}F_{\alpha\beta}$.
\item For an integer $\delta$, $\mathcal{V}_\delta$ is the space of gauge-invariant polynomials of
dimension $\delta$ in the jet coordinates. The hyperoctahedral group $W_4$ acts on
$\mathcal{V}_\delta$ through the space indices, each index transforming as in $V := \mathbb{R}^4$;
$\Lambda^2 := \Lambda^2\mathbb{R}^4$.
\item The \emph{jet module} $\mathsf{J}_k := [D^kF]$ is the $W_4$-module of the components of the
jet coordinates \eqref{10.1:eq-local-operators-1} of order $k$:
\begin{equation}\label{10.1:eq-local-operators-2}
\mathsf{J}_k = \operatorname{Sym}^{k+1}V \otimes V \ominus \operatorname{Sym}^{k+2}V, \qquad
\dim \mathsf{J}_0 = 6,\quad \dim \mathsf{J}_1 = 20,\quad \dim \mathsf{J}_2 = 45 .
\end{equation}
\item Write $\tau_a := -\tfrac{i}{2}\sigma_a$ ($a = 1,2,3$, with $\sigma_a$ the Pauli matrices),
$X = \sum_a X^a\tau_a$ for $X \in \mathfrak{su}(2)$, and $\epsilon^{abc}$ for the totally
antisymmetric symbol with $\epsilon^{123} = 1$. For $\delta \le 7$ every element of
$\mathcal{V}_\delta$ is a linear combination of bilinears $\operatorname{tr}(X_1X_2)$ and of
trilinears
$\operatorname{tr}(X_1X_2X_3) = (\mathrm{const})\sum_{a,b,c}\epsilon^{abc}X_1^aX_2^bX_3^c$, with
$X_1, X_2, X_3$ jet coordinates and a fixed numerical constant. As $W_4$-modules,
\begin{equation}\label{10.1:eq-local-operators-3}
\begin{split}
&\mathcal{V}_4 = \operatorname{Sym}^2\Lambda^2 \ (21), \qquad
\mathcal{V}_5 = \Lambda^2 \otimes \mathsf{J}_1 \ (120), \\
&\mathcal{V}_6 = \Lambda^2 \otimes \mathsf{J}_2 \oplus \operatorname{Sym}^2\mathsf{J}_1
\oplus \Lambda^3(\Lambda^2) \ (500),
\end{split}
\end{equation}
the dimensions in parentheses, and $\mathcal{V}_\delta = 0$ for $\delta = 1,2,3$.
\item For gauge-invariant local polynomials, $P \simeq P'$ means that $P - P'$ is a total
derivative $\sum_\alpha \partial_\alpha j_\alpha$ with each $j_\alpha$ gauge invariant.
\item With all indices summed,
\begin{align}
O_2 &:= \operatorname{tr}(D_\alpha F_{\mu\nu} D_\alpha F_{\mu\nu}), \qquad
O_3 := \operatorname{tr}(J_\nu J_\nu), \qquad
O_F := \operatorname{tr}(F_{\mu\nu}F_{\nu\rho}F_{\rho\mu}), \label{10.1:eq-local-operators-4}\\
Q &:= \sum_\alpha \partial_\alpha^2 \sum_\mu \operatorname{tr}(F_{\alpha\mu}F_{\alpha\mu}),
\label{10.1:eq-local-operators-5}
\end{align}
and, for $\alpha,\beta \in \{1,2,3,4\}$,
\begin{equation}\label{10.1:eq-local-operators-6}
\begin{gathered}
\mathcal{T}^{(1)}_{\alpha\beta} := \sum_\kappa \operatorname{tr}F_{\alpha\kappa}F_{\beta\kappa},
\quad
\mathcal{T}^{(2)}_{\alpha\beta} := \delta_{\alpha\beta}\sum_{\kappa,\lambda}
\operatorname{tr}F_{\kappa\lambda}F_{\kappa\lambda},
\\
\mathcal{T}^{(3)}_{\alpha\beta} := \delta_{\alpha\beta}\sum_\kappa
\operatorname{tr}F_{\alpha\kappa}F_{\alpha\kappa}.
\end{gathered}
\end{equation}
\end{enumerate}
For $\mathrm{SU}(N)$ with $N \ge 3$ the trilinears built with the symmetric invariant tensor
$d^{abc}$ are odd under charge conjugation, a symmetry of $\operatorname{Re}\operatorname{tr}$,
and therefore do not occur.
\end{definition}

\begin{proof}
The fourth item and \eqref{10.1:eq-local-operators-2} contain assertions, which we verify.

The dimensions in \eqref{10.1:eq-local-operators-2} are
$4\binom{k+4}{3} - \binom{k+5}{3}$, namely $16 - 10 = 6$, $40 - 20 = 20$ and $80 - 35 = 45$
for $k = 0,1,2$.

Since $\operatorname{tr}X = 0$ for $X \in \mathfrak{su}(2)$, a trace of a single factor vanishes.
A gauge-invariant trace therefore has at least two factors. Every jet coordinate has dimension at
least $2$, so such a trace has dimension at least $4$. This gives $\mathcal{V}_\delta = 0$ for
$\delta = 1,2,3$, and shows that a product of two traces has dimension at least $8$.

For $\delta \le 7$ there remain single traces of two or three factors; four factors would already
have dimension at least $8$. From $\sigma_a\sigma_b = \delta_{ab}\,1 + i\sum_c
\epsilon^{abc}\sigma_c$ one gets
$\tau_a\tau_b = -\tfrac14\delta_{ab}\,1 + \tfrac12\sum_c\epsilon^{abc}\tau_c$. Multiply by
$\tau_c$ and take the trace. The $\delta_{ab}$ term gives $\operatorname{tr}\tau_c = 0$, and
$\operatorname{tr}(\tau_d\tau_c)$ is a multiple of $\delta_{dc}$. Hence
$\operatorname{tr}(\tau_a\tau_b\tau_c)$ is a fixed multiple of $\epsilon^{abc}$, which is the stated
form of the trilinears.

We now read off \eqref{10.1:eq-local-operators-3}; after the reduction of the first item, the
components of the jet modules serve as independent coordinates.

In dimension $4$ the only terms are $\operatorname{tr}(F_{\mu\nu}F_{\kappa\lambda})$. Each factor
ranges over $\Lambda^2$, and $\operatorname{tr}(XY) = \operatorname{tr}(YX)$, so the monomial
depends only on the unordered pair of factors. This gives $\operatorname{Sym}^2\Lambda^2$, of
dimension $6\cdot 7/2 = 21$.

In dimension $5$ a trilinear would have dimension at least $6$. The only terms are therefore
$\operatorname{tr}(F\,DF)$, with one factor in $\Lambda^2$ and one in $\mathsf{J}_1$. This gives
$\Lambda^2 \otimes \mathsf{J}_1$, of dimension $6\cdot 20 = 120$.

In dimension $6$ there are three kinds of term.
\begin{itemize}
\item The bilinears $\operatorname{tr}(F\,D^2F)$ give $\Lambda^2 \otimes \mathsf{J}_2$, of dimension
$270$.
\item The bilinears $\operatorname{tr}(DF\,DF)$ are symmetric in their two factors and give
$\operatorname{Sym}^2\mathsf{J}_1$, of dimension $20\cdot 21/2 = 210$.
\item The trilinears $\operatorname{tr}(FFF)$ are totally antisymmetric in the three factors, by the
$\epsilon^{abc}$ form, and give $\Lambda^3(\Lambda^2)$, of dimension $\binom{6}{3} = 20$.
\end{itemize}
The total is $270 + 210 + 20 = 500$.
\end{proof}

\begin{lemma}[The parity rule for hyperoctahedral tensors]\label{10.1:lem-parity-rule}
Let $W_4$ act on $V:=\mathbb{R}^4$ through its defining representation by signed permutation
matrices, on $V^{\otimes n}$ ($n\ge 1$) through the $n$-th tensor power of that representation,
and on $\Lambda^2:=\Lambda^2\mathbb{R}^4$ by
$R\cdot(e_\rho\wedge e_\sigma):=Re_\rho\wedge Re_\sigma$ for $R\in W_4$, where $e_1,\dots,e_4$
is the standard basis. For a multi-index $\mu=(\mu_1,\dots,\mu_n)\in\{1,2,3,4\}^n$ and
$\alpha\in\{1,2,3,4\}$ put $n_\alpha(\mu):=\#\{i:\mu_i=\alpha\}$, and let $\delta_{\alpha\gamma}$
denote the Kronecker symbol. Let $\mathsf{t}\in V^{\otimes n}$ be
$W_4$-invariant with components
$\mathsf{t}=\sum_\mu\mathsf{t}_{\mu_1\cdots\mu_n}\,e_{\mu_1}\otimes\cdots\otimes e_{\mu_n}$. Let
$\mathsf{f}\colon\Lambda^2\to V^{\otimes n}$ and $\mathsf{g}\colon V\to V^{\otimes n}$ be
$W_4$-equivariant linear maps with components defined, for $\rho\neq\sigma$, by
$\mathsf{f}(e_\rho\wedge e_\sigma)=\sum_\mu\mathsf{f}_{\mu_1\cdots\mu_n;\rho\sigma}\,
e_{\mu_1}\otimes\cdots\otimes e_{\mu_n}$ and
$\mathsf{g}(e_\rho)=\sum_\mu\mathsf{g}_{\mu_1\cdots\mu_n;\rho}\,
e_{\mu_1}\otimes\cdots\otimes e_{\mu_n}$. Then, for every $\alpha\in\{1,2,3,4\}$,
\begin{equation}\label{10.1:eq-parity-rule-a}
\begin{aligned}
\mathsf{t}_{\mu_1\cdots\mu_n}\neq 0\;&\Longrightarrow\; n_\alpha(\mu)\equiv 0 \pmod 2,\\
\mathsf{f}_{\mu_1\cdots\mu_n;\rho\sigma}\neq 0\;&\Longrightarrow\;
n_\alpha(\mu)\equiv \delta_{\alpha\rho}+\delta_{\alpha\sigma} \pmod 2,\\
\mathsf{g}_{\mu_1\cdots\mu_n;\rho}\neq 0\;&\Longrightarrow\;
n_\alpha(\mu)\equiv \delta_{\alpha\rho} \pmod 2 .
\end{aligned}
\end{equation}
In words: a component of a $W_4$-invariant tensor vanishes unless every index value occurs an
even number of times; a component of an equivariant map from $\Lambda^2$ (from $V$) vanishes
unless the two free indices $\rho\neq\sigma$ (the one free index $\rho$) occur an odd number of
times and all other values an even number of times.
\end{lemma}

\begin{proof}
Fix $\alpha$ and let $\varepsilon_\alpha\in W_4$ be the reflection of the $\alpha$-th coordinate,
the diagonal matrix with $\varepsilon_\alpha e_\gamma=(-1)^{\delta_{\alpha\gamma}}e_\gamma$. Then
\[
\varepsilon_\alpha^{\otimes n}\bigl(e_{\mu_1}\otimes\cdots\otimes e_{\mu_n}\bigr)
=\prod_{i=1}^n(-1)^{\delta_{\alpha\mu_i}}\,e_{\mu_1}\otimes\cdots\otimes e_{\mu_n}
=(-1)^{n_\alpha(\mu)}\,e_{\mu_1}\otimes\cdots\otimes e_{\mu_n},
\]
so $\varepsilon_\alpha$ acts diagonally in the product basis. Invariance
$\varepsilon_\alpha^{\otimes n}\mathsf{t}=\mathsf{t}$ gives, on comparing coefficients,
$(-1)^{n_\alpha(\mu)}\mathsf{t}_{\mu_1\cdots\mu_n}=\mathsf{t}_{\mu_1\cdots\mu_n}$; if
$n_\alpha(\mu)$ is odd this reads $\mathsf{t}_{\mu_1\cdots\mu_n}=-\mathsf{t}_{\mu_1\cdots\mu_n}$,
hence $\mathsf{t}_{\mu_1\cdots\mu_n}=0$. This is the first line of \eqref{10.1:eq-parity-rule-a}.

For $\mathsf{f}$, the definition of the action on $\Lambda^2$ gives
$\varepsilon_\alpha\cdot(e_\rho\wedge e_\sigma)=(-1)^{\delta_{\alpha\rho}+\delta_{\alpha\sigma}}
e_\rho\wedge e_\sigma$. Equivariance
$\mathsf{f}(\varepsilon_\alpha\cdot(e_\rho\wedge e_\sigma))
=\varepsilon_\alpha^{\otimes n}\mathsf{f}(e_\rho\wedge e_\sigma)$,
expanded in the product basis by the identity displayed above, gives
\[
(-1)^{\delta_{\alpha\rho}+\delta_{\alpha\sigma}}\,\mathsf{f}_{\mu_1\cdots\mu_n;\rho\sigma}
=(-1)^{n_\alpha(\mu)}\,\mathsf{f}_{\mu_1\cdots\mu_n;\rho\sigma},
\]
so $\mathsf{f}_{\mu_1\cdots\mu_n;\rho\sigma}=0$ unless
$n_\alpha(\mu)\equiv\delta_{\alpha\rho}+\delta_{\alpha\sigma}\pmod 2$. Since $\rho\neq\sigma$, the
right-hand side is $1$ for $\alpha\in\{\rho,\sigma\}$ and $0$ otherwise, which is the second line
and its verbal form. For $\mathsf{g}$, equivariance applied to $e_\rho$ with
$\varepsilon_\alpha e_\rho=(-1)^{\delta_{\alpha\rho}}e_\rho$ gives in the same way
$(-1)^{\delta_{\alpha\rho}}\mathsf{g}_{\mu_1\cdots\mu_n;\rho}
=(-1)^{n_\alpha(\mu)}\mathsf{g}_{\mu_1\cdots\mu_n;\rho}$, which is the third line. Since
$\alpha$ was arbitrary, all three conclusions hold for every $\alpha\in\{1,2,3,4\}$.
\end{proof}

\begin{proposition}[The dimension-six classification]\label{10.1:prop-dimension-six}
Let $\mathcal{V}_\delta$ be the $W_4$-module of gauge-invariant local polynomials of dimension $\delta$,
let $\simeq$ denote equality modulo total derivatives $\partial_\alpha j_\alpha$ with $j$ gauge
invariant, and let $O_2$, $O_3$, $O_F$, $Q$ and $\mathcal{T}^{(1)},\mathcal{T}^{(2)},\mathcal{T}^{(3)}$
be as in Definition~\ref{10.1:def-local-operators}, and put
\begin{equation*}
O_1 := \sum_{\mu,\nu}\operatorname{tr}(D_\mu F_{\mu\nu}D_\mu F_{\mu\nu}).
\end{equation*}
An index is summed only where $\sum$ is written or
where it is repeated in an expression in which no index occurs more than twice.
Then the following hold.
\begin{enumerate}
\item[(1)] (Dimension at most five, scalars.) $\mathcal{V}_4^{W_4} = \mathbb{R}\cdot\operatorname{tr}F_{\mu\nu}F_{\mu\nu}$
and $\mathcal{V}_5^{W_4} = 0$. Hence every $W_4$-scalar of dimension at most five is an $O(4)$-scalar.
\item[(2)] (The $\Lambda^2$ channel is empty below six.)
\begin{equation*}
\operatorname{Hom}_{W_4}(\Lambda^2,\mathbb{R}) = \operatorname{Hom}_{W_4}(\Lambda^2,\mathcal{V}_4)
= \operatorname{Hom}_{W_4}(\Lambda^2,\mathcal{V}_5) = 0,
\end{equation*}
and $\Lambda^2$ is an irreducible $W_4$-module.
\item[(3)] (The vector channel.) $\operatorname{Hom}_{W_4}(V,\mathcal{V}_\delta) = 0$ for even $\delta$;
$\operatorname{Hom}_{W_4}(V,\mathcal{V}_5)$ is three-dimensional, spanned by
\begin{equation*}
v^1_\rho := \partial_\rho\operatorname{tr}F_{\alpha\beta}F_{\alpha\beta},\qquad
v^2_\rho := \operatorname{tr}(F_{\rho\nu}J_\nu),\qquad
v^3_\rho := \partial_\rho\sum_\kappa\operatorname{tr}F_{\rho\kappa}F_{\rho\kappa}
\end{equation*}
(no sum on $\rho$), and each of them is the divergence of a symmetric tensor:
\begin{equation}\label{10.1:eq-dimension-six-1}
v^1_\rho = \partial_\alpha\mathcal{T}^{(2)}_{\alpha\rho},\qquad
v^3_\rho = \partial_\alpha\mathcal{T}^{(3)}_{\alpha\rho},\qquad
v^2_\rho = \partial_\alpha\bigl(\mathcal{T}^{(1)}_{\alpha\rho}-\tfrac14\mathcal{T}^{(2)}_{\alpha\rho}\bigr).
\end{equation}
\item[(4)] (Symmetric tensors.) $\operatorname{Hom}_{W_4}(\operatorname{Sym}^2V,\mathcal{V}_4)$ is
three-dimensional, with basis $\mathcal{T}^{(1)},\mathcal{T}^{(2)},\mathcal{T}^{(3)}$; among them the
$O(4)$-equivariant ones are $\mathcal{T}^{(1)}$ and $\mathcal{T}^{(2)}$.
Moreover, by \textup{(2)}, $\operatorname{Hom}_{W_4}(V\otimes V,\mathcal{V}_4)
= \operatorname{Hom}_{W_4}(\operatorname{Sym}^2V,\mathcal{V}_4)$.
\item[(5)] (Dimension six, scalars.) Modulo total derivatives, $\mathcal{V}_6^{W_4}$ is
three-dimensional with basis $O_1,O_2,O_3$; the operator $O_F$ is dependent:
\begin{equation}\label{10.1:eq-dimension-six-2}
O_2 \simeq 2O_3 - 4O_F .
\end{equation}
The operators $O_2,O_3,O_F$ are $O(4)$-scalars and $O_1$ is not. Hence the quotient of
$\mathcal{V}_6^{W_4}$ by the span of the $O(4)$-scalars and the total derivatives is one-dimensional,
spanned by the class of $O_1$; no operator cubic in $F$ breaks $O(4)$. Modulo the field equation
$J = 0$ as well, $O_3$ drops, $O_1$ and $O_2 \simeq -4O_F$ remain independent, and $O_1$ still breaks
$O(4)$. As densities, that is modulo nothing, $\mathcal{V}_6^{W_4}$ is six-dimensional, with basis
$O_1,O_2,O_3$ and the three total derivatives $\partial_\rho v^i_\rho$, $i=1,2,3$; of these
$\partial_\rho v^1_\rho = \sum_\alpha\partial_\alpha^2\operatorname{tr}F_{\kappa\lambda}F_{\kappa\lambda}$
and $\partial_\rho v^2_\rho$ are $O(4)$-scalars and
$\partial_\rho v^3_\rho = Q$ is not. Thus there are two $O(4)$-breaking densities, $O_1$ and $Q$, and for
every test function $h$
\begin{equation}\label{10.1:eq-dimension-six-3}
\int hQ = \sum_\alpha\int(\partial_\alpha^2h)\,\mathcal{T}^{(3)}_{\alpha\alpha}.
\end{equation}
\end{enumerate}
\end{proposition}

\begin{proof}
Two facts are used throughout. First, the parity rule for hyperoctahedral tensors
(Lemma~\ref{10.1:lem-parity-rule}, display~\eqref{10.1:eq-parity-rule-a}), applied to the coefficients
with which an invariant element, or an equivariant map, is written in the monomials of
$\mathcal{V}_\delta$. Second, $F$ carries two indices and dimension two and $D$ carries one index and
dimension one, so a monomial of $\mathcal{V}_\delta$ carries exactly $\delta$ indices; the element
$-\mathrm{id}=\varepsilon_1\varepsilon_2\varepsilon_3\varepsilon_4$ of $W_4$ therefore acts on
$\mathcal{V}_\delta$ by $(-1)^\delta$, on $V$ by $-1$, and on $\Lambda^2$ and
$\operatorname{Sym}^2V$ by $+1$.

\emph{Part \textup{(1)}.} An element of $\mathcal{V}_4$ is
$\sum c_{\mu\nu,\kappa\lambda}\operatorname{tr}(F_{\mu\nu}F_{\kappa\lambda})$ with $\mu\neq\nu$,
$\kappa\neq\lambda$, the coefficients antisymmetric in each pair and symmetric under exchange of the
two pairs. If it is $W_4$-invariant, the parity rule says that $c_{\mu\nu,\kappa\lambda}$ vanishes unless
every value occurs an even number of times among $\mu,\nu,\kappa,\lambda$. Four distinct values would
each occur once, an odd number of times; since $\mu\neq\nu$ and $\kappa\neq\lambda$, the only
remaining possibility is $\{\mu,\nu\}=\{\kappa,\lambda\}$. Under a reflection $\varepsilon_\alpha$ the
monomial $\operatorname{tr}(F_{\mu\nu}F_{\mu\nu})$ acquires the square of a sign, so the reflections
impose no further condition, and the permutation subgroup $S_4\subset W_4$ acts transitively on
unordered pairs $\{\mu,\nu\}$. Hence there is one coefficient, and the invariant is a multiple of
$\operatorname{tr}F_{\mu\nu}F_{\mu\nu}$, which is invariant. For $\mathcal{V}_5$, the element
$-\mathrm{id}\in W_4$ acts by $-1$ (odd rank), so an invariant element is zero. Since
$\operatorname{tr}F_{\mu\nu}F_{\mu\nu}$ is an $O(4)$-scalar and $\mathcal{V}_\delta=0$ for
$\delta=1,2,3$ (Definition~\ref{10.1:def-local-operators}), every $W_4$-scalar of dimension at most
five is an $O(4)$-scalar.

\emph{Part \textup{(2)}.} An element of $\operatorname{Hom}_{W_4}(\Lambda^2,\mathbb{R})$ has
components $c_{\rho\sigma}$ with $\rho\neq\sigma$; each of the values $\rho,\sigma$ occurs once, which
the parity rule forbids, so $c=0$. An element of $\operatorname{Hom}_{W_4}(\Lambda^2,\mathcal{V}_4)$
has components $c_{\rho\sigma;\mu\nu,\kappa\lambda}$, antisymmetric in $\rho\sigma$, antisymmetric in
each of the pairs $\mu\nu$ and $\kappa\lambda$, and symmetric under
$(\mu\nu)\leftrightarrow(\kappa\lambda)$. By the parity rule the free values $\rho$ and $\sigma$ occur an
odd number of times among $\mu,\nu,\kappa,\lambda$ and every other value an even number of times. A
value occurs at most once in each pair, so $\rho$ and $\sigma$ occur once each, and the two remaining
slots carry one further value $\tau\notin\{\rho,\sigma\}$ twice, once in each pair. The only possible
components are therefore $c_{\rho\sigma;\rho\tau,\sigma\tau}$ with $\tau\notin\{\rho,\sigma\}$, up to
the symmetries. The transposition $(\rho\sigma)\in W_4$ gives
\begin{equation*}
c_{\sigma\rho;\sigma\tau,\rho\tau} = c_{\rho\sigma;\rho\tau,\sigma\tau},
\end{equation*}
while antisymmetry in $\rho\sigma$ and symmetry under exchange of the two pairs give
\begin{equation*}
c_{\sigma\rho;\sigma\tau,\rho\tau} = -c_{\rho\sigma;\sigma\tau,\rho\tau} = -c_{\rho\sigma;\rho\tau,\sigma\tau}.
\end{equation*}
So $c=0$. For $\mathcal{V}_5$, $-\mathrm{id}$ acts by $+1$ on $\Lambda^2$ (even rank) and by $-1$ on
$\mathcal{V}_5$ (odd rank), so an equivariant map vanishes. Irreducibility: the sign subgroup of
$W_4$ acts on the basis vector $E_{\mu\nu}=e_\mu\wedge e_\nu$ by the character
$\varepsilon\mapsto\varepsilon_\mu\varepsilon_\nu$, and the six characters so obtained are distinct;
hence a $W_4$-invariant subspace is spanned by the basis vectors $E_{\mu\nu}$ it contains, and since the
permutations act transitively on these, it is $0$ or $\Lambda^2$.

\emph{Part \textup{(3)}.} For even $\delta$, $-\mathrm{id}$ acts by $-1$ on $V$ and by $+1$ on
$\mathcal{V}_\delta$, so every equivariant map vanishes. For $\delta=5$, the $\rho$-component $j_\rho$
of $j\in\operatorname{Hom}_{W_4}(V,\mathcal{V}_5)$ is a combination of the monomials
$\operatorname{tr}(F_{\alpha\beta}D_\gamma F_{\sigma\tau})$ in which, by the parity rule, $\rho$
occurs an odd number of times and every other value an even number of times. Each $F$ carries a value
at most once, so a value occurs at most three times; this leaves two index patterns, $\{\rho\rho\rho\kappa\kappa\}$
and $\{\rho\kappa\kappa\lambda\lambda\}$ with $\kappa,\lambda\neq\rho$ distinct.

In the pattern $\{\rho\rho\rho\kappa\kappa\}$ the value $\rho$ sits in the first $F$, in $D$ and in the
second $F$, and the monomial is
$\operatorname{tr}(F_{\rho\kappa}D_\rho F_{\rho\kappa})
=\tfrac12\partial_\rho\operatorname{tr}F_{\rho\kappa}F_{\rho\kappa}$,
because $\partial_\rho\operatorname{tr}(\Phi\Psi)=\operatorname{tr}(D_\rho\Phi\,\Psi)
+\operatorname{tr}(\Phi\,D_\rho\Psi)$
for covariant $\Phi,\Psi$ and the trace is symmetric. In the pattern $\{\rho\kappa\kappa\lambda\lambda\}$
there are three places for $\rho$. If $D$ carries $\rho$, the monomial is
$\operatorname{tr}(F_{\kappa\lambda}D_\rho F_{\kappa\lambda})
=\tfrac12\partial_\rho\operatorname{tr}F_{\kappa\lambda}F_{\kappa\lambda}$.
If the first $F$ carries $\rho$, the monomial is
$\operatorname{tr}(F_{\rho\kappa}D_\lambda F_{\lambda\kappa})$. If the second $F$ carries $\rho$, the
monomial is $\operatorname{tr}(F_{\kappa\lambda}D_\kappa F_{\rho\lambda})$, and antisymmetry of
$F_{\kappa\lambda}$ together with the Bianchi identity
$D_\kappa F_{\rho\lambda}+D_\rho F_{\lambda\kappa}+D_\lambda F_{\kappa\rho}=0$ gives
\begin{equation}\label{10.1:eq-dimension-six-4}
\operatorname{tr}(F_{\kappa\lambda}D_\kappa F_{\rho\lambda})
= \tfrac12\operatorname{tr}F_{\kappa\lambda}(D_\kappa F_{\rho\lambda}-D_\lambda F_{\rho\kappa})
= \tfrac12\operatorname{tr}F_{\kappa\lambda}D_\rho F_{\kappa\lambda},
\end{equation}
summed over $\kappa,\lambda$. An equivariant $j$ is a combination of the sums of these monomials over
the values $\kappa,\lambda$ admitted by $W_4$, that is, distinct from $\rho$ and from each other; such a
restricted sum differs from the corresponding full sum over all $\kappa,\lambda$ by terms of the first
pattern. The full sums are $\tfrac12 v^1_\rho$ (for $D$ carrying $\rho$), $v^2_\rho$ (for the first $F$
carrying $\rho$, since $D_\lambda F_{\lambda\kappa}=J_\kappa$), and $\tfrac14 v^1_\rho$ by
\eqref{10.1:eq-dimension-six-4} (for the second $F$ carrying $\rho$), while the first pattern summed
over $\kappa$ is $\tfrac12 v^3_\rho$. So the span is $\{v^1,v^2,v^3\}$. Each $v^i$ is $W_4$-equivariant:
a permutation relabels its indices, and $\varepsilon_\alpha$ changes its sign exactly when
$\rho=\alpha$, the other indices entering in pairs.

The three are linearly independent, as their evaluation on abelian fields shows. Let $T$ be a fixed
non-zero element of the Lie algebra, so that $\operatorname{tr}T^2\neq0$. On the field
$A_\mu=a_\mu e^{ip\cdot x}T$ with $p,a\in\mathbb{C}^4$ one has $D=\partial$ and
$F_{\alpha\beta}=if_{\alpha\beta}e^{ip\cdot x}T$ with $f_{\alpha\beta}:=p_\alpha a_\beta-p_\beta a_\alpha$,
hence $J_\nu=-\sum_\alpha p_\alpha f_{\alpha\nu}\,e^{ip\cdot x}T$; with $p^2:=\sum_\mu p_\mu^2$,
$a^2:=\sum_\mu a_\mu^2$ and $p\cdot a:=\sum_\mu p_\mu a_\mu$,
\begin{equation}\label{10.1:eq-dimension-six-5}
\sum_{\alpha,\beta}f_{\alpha\beta}^2=2\bigl(p^2a^2-(p\cdot a)^2\bigr),\qquad
\sum_\alpha p_\alpha f_{\alpha\nu}=p^2a_\nu-(p\cdot a)p_\nu .
\end{equation}
Suppose $c_1v^1+c_2v^2+c_3v^3=0$. At $p=(1,i,0,0)$, $a=e_3$ both $p^2$ and $p\cdot a$ vanish, so by
\eqref{10.1:eq-dimension-six-5} the trace $\operatorname{tr}F_{\alpha\beta}F_{\alpha\beta}$ and $J$
vanish, and with them $v^1$ and $v^2$; since $f_{1\kappa}=\delta_{\kappa3}$, the component $v^3_1$ is
$\partial_1(-e^{2ip\cdot x}\operatorname{tr}T^2)=-2ie^{2ip\cdot x}\operatorname{tr}T^2$, which is not
zero, so $c_3=0$. Next let $p'=(1,-i,0,0)$, for which $p'^2$ and $p'\cdot a$ also vanish, let $f'$ be
formed from $p'$ and $a$ as $f$ is from $p$ and $a$, and take $A_\mu=a_\mu(e^{ip\cdot x}+e^{ip'\cdot x})T$.
Its current $J$ is the sum of the currents of the two waves, which vanish, so $v^2=0$; since
$p\cdot p'=2$, $a^2=1$ and the squares $\sum f_{\alpha\beta}^2$, $\sum f'^2_{\alpha\beta}$ vanish,
\begin{gather*}
\sum_{\alpha,\beta}f_{\alpha\beta}f'_{\alpha\beta}
=2\bigl((p\cdot p')a^2-(p\cdot a)(p'\cdot a)\bigr)=4,\\
\operatorname{tr}F_{\alpha\beta}F_{\alpha\beta}=-8e^{2ix_1}\operatorname{tr}T^2,\qquad
v^1_1=-16ie^{2ix_1}\operatorname{tr}T^2\neq0,
\end{gather*}
so $c_1=0$. Finally, at $p=e_1$, $a=e_2$ the only non-zero components of $F$ are
$F_{12}=-F_{21}=ie^{ix_1}T$, the only non-zero component of $J$ is $J_2=-e^{ix_1}T$, and
\begin{equation*}
v^2_1=\operatorname{tr}(F_{12}J_2)=-ie^{2ix_1}\operatorname{tr}T^2\neq0,
\end{equation*}
so $c_2=0$.

For the divergences, by the factor $\delta_{\alpha\rho}$,
\begin{align*}
\partial_\alpha\mathcal{T}^{(2)}_{\alpha\rho}
&=\partial_\rho\operatorname{tr}F_{\kappa\lambda}F_{\kappa\lambda}=v^1_\rho,\\
\partial_\alpha\mathcal{T}^{(3)}_{\alpha\rho}
&=\partial_\rho\sum_\kappa\operatorname{tr}F_{\rho\kappa}F_{\rho\kappa}=v^3_\rho.
\end{align*}
Finally
\begin{equation*}
v^2_\rho = \operatorname{tr}(F_{\rho\nu}D_\alpha F_{\alpha\nu})
= \partial_\alpha\operatorname{tr}(F_{\rho\nu}F_{\alpha\nu})
- \operatorname{tr}(F_{\alpha\nu}D_\alpha F_{\rho\nu}),
\end{equation*}
where $\operatorname{tr}(F_{\rho\nu}F_{\alpha\nu})=\mathcal{T}^{(1)}_{\alpha\rho}$, and the last trace is
$\tfrac14\partial_\rho\operatorname{tr}F_{\kappa\lambda}F_{\kappa\lambda}
=\tfrac14\partial_\alpha\mathcal{T}^{(2)}_{\alpha\rho}$ by \eqref{10.1:eq-dimension-six-4}. Hence
$v^2_\rho=\partial_\alpha(\mathcal{T}^{(1)}-\tfrac14\mathcal{T}^{(2)})_{\alpha\rho}$, which completes
\eqref{10.1:eq-dimension-six-1}.

\emph{Part \textup{(4)}.} An element of $\operatorname{Hom}_{W_4}(\operatorname{Sym}^2V,\mathcal{V}_4)$
has components $c_{\alpha\beta;\mu\nu,\kappa\lambda}$, symmetric in $\alpha\beta$ and otherwise as in
part~(1); read as a $W_4$-invariant tensor in all six indices, the parity rule applies to it. If
$\alpha=\beta$, every value occurs an even number of times among $\mu,\nu,\kappa,\lambda$, so the pairs
coincide as sets, exactly as in part~(1); the signs occur squared, and the configurations
$(\alpha,\{\mu,\nu\})$ form two $S_4$-orbits, according as $\alpha$ lies in the pair or not. If
$\alpha\neq\beta$, the values $\alpha,\beta$ occur once each among $\mu,\nu,\kappa,\lambda$ and a
further value $\tau\notin\{\alpha,\beta\}$ twice, as in part~(2), so the pairs are
$\{\alpha,\tau\},\{\beta,\tau\}$;
there is one coefficient, since the signs occur squared and $S_4$ is transitive on ordered triples of
distinct values. So the dimension is at most three. The three tensors are realised: $\mathcal{T}^{(3)}$
has only diagonal components, and these contain only pairs containing $\alpha$;
$\mathcal{T}^{(2)}_{\alpha\alpha}=\operatorname{tr}F_{\kappa\lambda}F_{\kappa\lambda}$ contains also the
pairs not containing $\alpha$; $\mathcal{T}^{(1)}$ has the non-zero off-diagonal components
$\sum_\kappa\operatorname{tr}F_{\alpha\kappa}F_{\beta\kappa}$. They are therefore independent, and a
basis. The tensors $\mathcal{T}^{(1)}$ and $\mathcal{T}^{(2)}$ are formed by contraction of indices and
with $\delta_{\alpha\beta}$, hence are $O(4)$-equivariant. A combination with a non-zero coefficient of
$\mathcal{T}^{(3)}$ is $O(4)$-equivariant only if $\mathcal{T}^{(3)}$ is; but then its components in
every orthonormal frame would be diagonal, so it would be a multiple of $\delta_{\alpha\beta}$, whereas at
a point where all components of $F$ vanish except $F_{12}=-F_{21}=T$, with $T$ as in part~(3),
one has $\mathcal{T}^{(3)}_{11}=\operatorname{tr}T^2\neq0$ whereas $\mathcal{T}^{(3)}_{33}=0$. Lastly
$V\otimes V=\operatorname{Sym}^2V\oplus\Lambda^2$ and
$\operatorname{Hom}_{W_4}(\Lambda^2,\mathcal{V}_4)=0$ by part~(2).

\emph{Part \textup{(5)}: spanning.} By the description of $\mathcal{V}_6$ in
Definition~\ref{10.1:def-local-operators} its monomials are of the forms
$\operatorname{tr}(F\cdot DDF)$, $\operatorname{tr}(DF\cdot DF)$ and $\operatorname{tr}(FFF)$. Since
\begin{equation*}
\operatorname{tr}(F_{\mu\nu}D_\alpha D_\beta F_{\rho\sigma})
= \partial_\alpha\operatorname{tr}(F_{\mu\nu}D_\beta F_{\rho\sigma})
- \operatorname{tr}(D_\alpha F_{\mu\nu}D_\beta F_{\rho\sigma}),
\end{equation*}
we have $\operatorname{tr}(F\cdot DDF)\simeq-\operatorname{tr}(DF\cdot DF)$. The space of total
derivatives is $W_4$-stable, so a $W_4$-invariant element is congruent modulo total derivatives to the
$W_4$-average of a combination of the monomials
$\operatorname{tr}(D_\alpha F_{\mu\nu}D_\beta F_{\rho\sigma})$ and
$\operatorname{tr}(FFF)$, and it suffices to find the invariant combinations of these.

For $\operatorname{tr}(D_\alpha F_{\mu\nu}D_\beta F_{\rho\sigma})$ the parity rule leaves the patterns
$(4,2)$ and $(2,2,2)$ of multiplicities (a value cannot occur six times, an $F$ carrying it at most
once). In the pattern $(4,2)$ the value $\kappa$ occurring four times sits in $D$, $F$, $D$, $F$, and
the monomial is $\operatorname{tr}(D_\kappa F_{\kappa\lambda}D_\kappa F_{\kappa\lambda})$; the signs
occur squared and $S_4$ is transitive on the pairs $(\kappa,\lambda)$, so the invariant combination is
$O_1$. In the pattern $(2,2,2)$ the six slots are paired by equal values, the pairings are those of the
$O(4)$-contractions, and a sum over distinct values differs from the full contraction by terms of the
pattern $(4,2)$. The contractions with the totally antisymmetric symbol are excluded by the reflections
$\varepsilon_\alpha$. Among the
full contractions, contracting $\alpha$ with $\beta$ gives $O_2$; contracting $\alpha$ and $\beta$ each
with an index of the other factor gives
\begin{equation*}
O_4 := \operatorname{tr}(D_\alpha F_{\beta\nu}D_\beta F_{\alpha\nu});
\end{equation*}
contracting each with an index of its own factor gives $O_3$; every other pairing contracts an $F$ with
itself and vanishes. Put $X_{\alpha\beta\nu}:=D_\alpha F_{\beta\nu}$. The Bianchi identity says
$X_{\alpha\beta\nu}+X_{\beta\nu\alpha}+X_{\nu\alpha\beta}=0$; squaring it (under the trace) and summing,
the three squares contribute $3\sum\operatorname{tr}(X_{\alpha\beta\nu}X_{\alpha\beta\nu})$ and the six
cross terms coincide after relabelling, so
\begin{equation*}
3\sum\operatorname{tr}(X_{\alpha\beta\nu}X_{\alpha\beta\nu})
+ 6\sum\operatorname{tr}(X_{\alpha\beta\nu}X_{\beta\nu\alpha}) = 0 .
\end{equation*}
Since $\sum\operatorname{tr}(X_{\alpha\beta\nu}X_{\alpha\beta\nu})=O_2$ and
$X_{\beta\alpha\nu}=-X_{\beta\nu\alpha}$, this gives
\begin{equation*}
O_4=-\sum\operatorname{tr}(X_{\alpha\beta\nu}X_{\beta\nu\alpha})=\tfrac12 O_2 .
\end{equation*}
For
$\operatorname{tr}(FFF)$ no value can occur four times, since it occurs at most once in each $F$; the
pattern $(2,2,2)$ then forces the monomials $\operatorname{tr}(F_{\kappa\lambda}F_{\lambda\tau}F_{\tau\kappa})$,
whose invariant combination is $O_F$. Hence, modulo total derivatives, $\mathcal{V}_6^{W_4}$ is spanned
by $O_1,O_2,O_3,O_F$.

\emph{Part \textup{(5)}: the relation.} By the Bianchi identity and
$[D_\alpha,D_\mu]=\operatorname{ad}F_{\alpha\mu}$,
\begin{equation*}
D^2F_{\mu\nu} = D_\mu J_\nu - D_\nu J_\mu + 2[F_{\alpha\mu},F_{\alpha\nu}] .
\end{equation*}
Moving $D_\mu$ across the trace,
\begin{equation*}
\operatorname{tr}(F_{\mu\nu}D_\mu J_\nu)\simeq -\operatorname{tr}(D_\mu F_{\mu\nu}J_\nu)=-O_3,
\end{equation*}
and by antisymmetry the term $-D_\nu J_\mu$ gives the
same. Relabelling $\mu\leftrightarrow\nu$ in the second term of the commutator,
\begin{equation*}
\operatorname{tr}(F_{\mu\nu}[F_{\alpha\mu},F_{\alpha\nu}])
=2\operatorname{tr}(F_{\mu\nu}F_{\alpha\mu}F_{\alpha\nu})=2O_F.
\end{equation*}
Hence
\begin{equation*}
O_2 \simeq -\operatorname{tr}F_{\mu\nu}D^2F_{\mu\nu} \simeq 2O_3 - 4O_F,
\end{equation*}
which is \eqref{10.1:eq-dimension-six-2}; so $O_1,O_2,O_3$ span modulo total derivatives.

\emph{Part \textup{(5)}: independence; $O_1$ breaks $O(4)$, also on shell.} The operators
$O_2,O_3,O_F$ are full contractions of indices and hence $O(4)$-scalars. On the abelian fields
$A_\mu=a_\mu e^{ip\cdot x}T$ of part~(3) we use the quadratic symbol $q_O(p,a)$ of a local polynomial
$O$: the value of the symmetric bilinear form polarising the part of $O$ quadratic in $A$, taken at the
pair of fields $a_\mu e^{ip\cdot x}T$ and $a_\mu e^{-ip\cdot x}T$, divided by $\operatorname{tr}T^2$.
It does not depend on $x$, and it vanishes for a total derivative, whose bilinear form at this pair is
the derivative of an $x$-independent expression. The quadratic symbols of $O_1,O_2,O_3$ are
\begin{equation*}
q_1 = \sum_{\mu,\nu} p_\mu^2 f_{\mu\nu}^2,\qquad q_2 = p^2\sum_{\mu,\nu} f_{\mu\nu}^2,\qquad
q_3 = \sum_\nu\Bigl(\sum_\mu p_\mu f_{\mu\nu}\Bigr)^2,
\end{equation*}
with $f$, $p^2$, $a^2$ and $p\cdot a$ as in part~(3). By \eqref{10.1:eq-dimension-six-5},
$q_3=p^2(p^2a^2-(p\cdot a)^2)$ and $q_2=2q_3$, as \eqref{10.1:eq-dimension-six-2} requires at
quadratic order, $O_F$ having no part quadratic in $A$. The polynomial $q_3$ is not identically zero
(it equals $1$ at $p=e_1$, $a=e_2$), and it vanishes at the point $p=(1,i,0,0)$, $a=e_3$ of the
complex shell considered below, where $q_1=2$; so $q_1$ and $q_3$ are linearly independent. Hence a
relation $c_1O_1+c_2O_2+c_3O_3\simeq0$ gives $c_1q_1+(2c_2+c_3)q_3=0$, so $c_1=0$ and $c_3=-2c_2$; by
\eqref{10.1:eq-dimension-six-2} the relation reduces to $-4c_2O_F\simeq0$, and $c_2=0$ because $O_F$,
cubic in the field, is not congruent to $0$ modulo total derivatives. So $O_1,O_2,O_3$ are independent
modulo total derivatives. On the complex shell
$p^2=0=p\cdot a$ one has $q_2=q_3=0$, while at $p=(1,i,0,0)$, $a=e_3$,
\begin{equation*}
q_1(p,a) = p_1^4+p_2^4 = 2,
\end{equation*}
and at the image $p'=((1-i),(1+i),0,0)/\sqrt2$ of $p$ under the rotation by $\pi/4$ in the
$(12)$-plane, which lies on
the shell with the same $a$, $q_1(p',a)=-2$. The quadratic symbol of an $O(4)$-scalar is unchanged when
$p$ and $a$ are rotated together, so $O_1$ is not congruent modulo total derivatives to an
$O(4)$-scalar, also modulo $J=0$. With
the spanning this shows that the quotient of $\mathcal{V}_6^{W_4}$ by the span of the $O(4)$-scalars and
the total derivatives is spanned by the class of $O_1$ and is one-dimensional; and since the only
$W_4$-invariant cubic in $F$ is $O_F$, no $F^3$-operator breaks $O(4)$. Modulo $J=0$, the operator
$O_3=\operatorname{tr}(J_\nu J_\nu)$ drops and \eqref{10.1:eq-dimension-six-2} becomes
$O_2\simeq-4O_F$. The quadratic symbol of $O_2$ vanishes on the shell while that of $O_1$ does not, and
$O_F$, being cubic in the field, is not congruent to $0$ modulo total derivatives and $J=0$; so $O_1$
and $O_2$ remain independent.

\emph{Part \textup{(5)}: densities.} If a total derivative $\partial_\rho j_\rho$ is $W_4$-invariant, it
is unchanged when $j$ is replaced by its $W_4$-average, since the divergence intertwines the actions;
hence the $W_4$-scalar total derivatives are the $\partial_\rho j_\rho$ with
$j\in\operatorname{Hom}_{W_4}(V,\mathcal{V}_5)$, that is, by part~(3), the combinations of
$\partial_\rho v^i_\rho$, $i=1,2,3$. The map $j\mapsto\partial\cdot j$ is injective on
$\operatorname{Hom}_{W_4}(V,\mathcal{V}_5)$: a kernel element would be $\partial_\beta k_{[\alpha\beta]}$
with $k\in\operatorname{Hom}_{W_4}(\Lambda^2,\mathcal{V}_4)=0$ by part~(2). Together with the
independence of $O_1,O_2,O_3$ modulo total derivatives this gives the six-dimensional basis. By the
definitions, $\partial_\rho v^1_\rho=\sum_\alpha\partial_\alpha^2\operatorname{tr}F_{\kappa\lambda}F_{\kappa\lambda}$,
and $v^1,v^2$ are $O(4)$-vectors, so
$\partial_\rho v^1_\rho$ and $\partial_\rho v^2_\rho$ are $O(4)$-scalars; and
\begin{equation*}
\partial_\rho v^3_\rho=\sum_\rho\partial_\rho^2\sum_\kappa\operatorname{tr}F_{\rho\kappa}F_{\rho\kappa}=Q.
\end{equation*}
On the abelian plane wave above, $F_{\alpha\mu}=if_{\alpha\mu}e^{ip\cdot x}T$ and
$\partial_\alpha^2e^{2ip\cdot x}=-4p_\alpha^2e^{2ip\cdot x}$, so the value of $Q$ at $x=0$ is
$4q_1(p,a)\operatorname{tr}T^2$, which takes the values $8\operatorname{tr}T^2$ and
$-8\operatorname{tr}T^2$ at the rotated pair $(p,a)$, $(p',a)$; hence $Q$ is not an $O(4)$-scalar.
Finally $Q=\sum_\alpha\partial_\alpha^2\mathcal{T}^{(3)}_{\alpha\alpha}$, and two integrations by parts
against the test function $h$ give \eqref{10.1:eq-dimension-six-3}.
\end{proof}

\begin{proposition}[The antisymmetric channel at dimension six. Stated; proof incomplete at the step marked $(\ast)$]\label{10.1:prop-antisymmetric-channel}
Let $\mathcal{V}_\delta$ be the space of gauge-invariant local polynomials of dimension $\delta$ in
the curvature $F_{\mu\nu}$ and its covariant derivatives, let $V := \mathbb{R}^4$ and
$\Lambda^2 := \Lambda^2V$, with their actions of the hyperoctahedral group $W_4$, and write $\simeq$
for equality modulo total derivatives. We say that an identity holds \emph{on shell} if it holds
modulo multiples of the field equation $J = 0$, where $J_\nu := D_\mu F_{\mu\nu}$. Then:
\begin{enumerate}
\item $\operatorname{Hom}_{W_4}(\Lambda^2,\mathcal{V}_6)$ has dimension $11$. Modulo total derivatives
it has dimension $2$. Modulo total derivatives and on shell it has dimension $1$.
\item A basis modulo total derivatives is formed by the two maps
\begin{equation}\label{10.1:eq-antisymmetric-channel-1}
\Phi_1 : \omega \longmapsto \sum_{\mu,\rho}\omega_{\rho\mu}B_{\mu\rho},
\qquad
\Phi_2 : \omega \longmapsto \sum_{\mu,\rho}\omega_{\rho\mu}
\operatorname{tr}\bigl(J_\mu D_\mu F_{\rho\mu}\bigr),
\end{equation}
where the only summations are the displayed ones. The value $\Phi_1(\omega)$ is the variation of
$O_1$ under the infinitesimal rotation generated by $\omega$
(Proposition~\ref{10.1:prop-classical-limit}), and the map $\Phi_1$ is non-zero modulo total
derivatives and on shell. The map
$\Phi_2$ is a multiple of $J$. These are exactly the two
$\Lambda^2$-structures that occur in display \eqref{10.1:eq-classical-limit-a} of
Proposition~\ref{10.1:prop-classical-limit}.
\end{enumerate}
\end{proposition}

\begin{proof}
The argument has two parts. The first gives the dimension counts; it is the incomplete step. The
second shows that the two maps \eqref{10.1:eq-antisymmetric-channel-1} are independent.

\emph{The counts.} They rest on two ingredients.

First, one needs the multiplicities of $\Lambda^2$ in $\mathcal{V}_6$, in $V\otimes\mathcal{V}_5$ and
in $\Lambda^2\otimes\mathcal{V}_4$. They are $11$, $16$ and $7$ respectively.
These numbers come from the character computation over the $384$ elements of $W_4$. That
computation is not displayed here $(\ast)$.

Second, one uses the horizontal complex
\[
\Lambda^2\otimes\mathcal{V}_4 \longrightarrow V\otimes\mathcal{V}_5 \longrightarrow \mathcal{V}_6,
\]
whose maps are divergences. One needs its exactness on gauge-invariant local forms of positive
dimension. This is the algebraic Poincar\'e lemma in its covariant form: in $d = 4$ the only
obstruction is the parity-odd form $\operatorname{tr} F\wedge F$. This exactness is used here
without proof $(\ast)$.

Granted both ingredients, the dimension of $\operatorname{Hom}_{W_4}(\Lambda^2,\mathcal{V}_6)$
modulo total derivatives is the Euler count
\[
11 - 16 + 7 = 2 .
\]

On shell, each jet module $\mathsf{J}_k$, the span of the components of $D^kF$, is replaced by its
quotient by the span of the components of $D^{k-1}J$.
The corresponding multiplicities, from the same undisplayed character computation, are $6$, $12$
and $7$. The exactness needed on shell, namely that there is no constant-coefficient conserved
current of dimension $5$, is likewise used without proof. With it, the Euler count is
\[
6 - 12 + 7 = 1 .
\]

\emph{Independence of the two maps.} The map $\Phi_2$ is a multiple of $J$, so it vanishes on
shell.

The map $\Phi_1$ does not vanish modulo total derivatives and on shell. Indeed, by
Proposition~\ref{10.1:prop-classical-limit},
$\sum_{\mu,\rho}\omega_{\rho\mu}B_{\mu\rho}$ is the variation of $O_1$ under the infinitesimal
rotation with Killing field $\xi(x) = \omega x$. By
clause~(5) of Proposition~\ref{10.1:prop-dimension-six}, $O_1$ still breaks $O(4)$ modulo total
derivatives and the field equation $J = 0$. Since $O(4)$ is generated by $SO(4)$ and one coordinate
reflection, and $O_1$ is invariant under $W_4$, which contains the coordinate reflections, $O_1$ is
not $SO(4)$-invariant in the finite-dimensional quotient of $\mathcal{V}_6$ modulo total derivatives
and on shell. On that quotient, $SO(4)$-invariance is equivalent to annihilation by every
infinitesimal rotation $\omega$; hence
$\sum_{\mu,\rho}\omega_{\rho\mu}B_{\mu\rho}\not\simeq 0$ on shell for some $\omega$, and $\Phi_1$ is
non-zero modulo total derivatives and on shell.

The map $\Phi_2$ is not $\simeq 0$. To see this, fix in this paragraph a non-zero traceless
anti-Hermitian matrix $\Xi$, so that $\operatorname{tr}\Xi^2 < 0$, vectors $p, a \in \mathbb{Z}^4$,
and the field $A_\mu(x) := a_\mu\cos(p\cdot x)\,\Xi$ (the letters $p$, $a$ and $f$ carry the
meaning given here only in the present paragraph). All commutators vanish on this field, so
$D_\mu = \partial_\mu$ there, and with $f_{\alpha\beta} := p_\alpha a_\beta - p_\beta a_\alpha$,
$p^2 := \sum_\alpha p_\alpha^2$ and
$\sum_\alpha p_\alpha f_{\alpha\mu} = p^2a_\mu - (p\cdot a)p_\mu$ one has
\begin{align*}
F_{\alpha\beta} &= -f_{\alpha\beta}\sin(p\cdot x)\,\Xi, \qquad
J_\mu = -\bigl(p^2a_\mu - (p\cdot a)p_\mu\bigr)\cos(p\cdot x)\,\Xi,\\
D_\mu F_{\rho\mu} &= -f_{\rho\mu}\,p_\mu\cos(p\cdot x)\,\Xi
\qquad (\text{no summation over } \mu).
\end{align*}
Hence
\[
\Phi_2(\omega) = \operatorname{tr}(\Xi^2)\cos^2(p\cdot x)
\sum_{\mu,\rho}\omega_{\rho\mu}f_{\rho\mu}\,p_\mu\bigl(p^2a_\mu - (p\cdot a)p_\mu\bigr).
\]
On this field every local polynomial in $A$ is a smooth function of $x$ that is $2\pi$-periodic in
each coordinate, so a total derivative has mean zero over $[0,2\pi]^4$. Take
$\omega_{12} = -\omega_{21} = 1$ and all other entries zero. The two non-zero terms of the sum
give, using $f_{21} = -f_{12}$,
\[
f_{12}\bigl[p^2(p_1a_1 + p_2a_2) - (p\cdot a)(p_1^2 + p_2^2)\bigr].
\]
At $p = (1,0,1,0)$ and $a = (1,1,0,0)$ one has $p^2 = 2$, $p\cdot a = 1$, $f_{12} = 1$,
$p_1a_1 + p_2a_2 = 1$ and $p_1^2 + p_2^2 = 1$, so this number is $1$, and the mean of
$\Phi_2(\omega)$ over $[0,2\pi]^4$ is $\tfrac12\operatorname{tr}\Xi^2 \neq 0$. Thus $\Phi_2(\omega)$
is not a total derivative.

Now suppose that $c_1\Phi_1 + c_2\Phi_2 \simeq 0$ with real $c_1, c_2$. Passing to the field
equation $J = 0$ as well, $\Phi_2$ drops out and $c_1\Phi_1$ vanishes modulo total derivatives and
on shell; since $\Phi_1$ is non-zero modulo total derivatives and on shell, $c_1 = 0$. Then
$c_2\Phi_2 \simeq 0$, and $c_2 = 0$ by the preceding paragraph.
Hence the two maps \eqref{10.1:eq-antisymmetric-channel-1} are linearly independent modulo total
derivatives. Given the count $2$
obtained at step $(\ast)$, they form a basis of
$\operatorname{Hom}_{W_4}(\Lambda^2,\mathcal{V}_6)$ modulo total derivatives.
\end{proof}

For each of the $W_4$-modules $\mathbb{R}$ (the trivial module), $V$, $\Lambda^2$ and the symmetric
square $\operatorname{Sym}^2V$ of $V$, the following table gives the dimension of the space of
$W_4$-equivariant linear maps from that module to $\mathcal{V}_\delta$, for $\delta = 4,5,6$.
\begin{equation}\label{10.1:eq-antisymmetric-channel-3}
\begin{array}{l|ccc|cc}
\text{type} & \delta = 4 & \delta = 5 & \delta = 6
& \begin{array}{c}\text{modulo total}\\ \text{derivatives}\end{array}
& \begin{array}{c}\text{modulo total derivatives}\\ \text{and on shell}\end{array}\\ \hline
\mathbb{R} & 1 & 0 & 6 & 3\ (\delta = 6) & 2\ (\delta = 6)\\
V & 0 & 3 & 0 & 0\ (\delta = 5) & \\
\Lambda^2 & 0 & 0 & 11 & 2\ (\delta = 6) & 1\ (\delta = 6)\\
\operatorname{Sym}^2 V & 3 & 0 & 37 & &
\end{array}
\end{equation}
All entries of the three $\delta$-columns are multiplicities from the character computation of
step $(\ast)$ of the proof above; the $\operatorname{Sym}^2V$ entry $37$ and the $\Lambda^2$
entry $11$ rest on it alone. The last two columns collect the counts modulo total derivatives and
on shell: for $\Lambda^2$ those of the proof above, for $V$ and $\mathbb{R}$ those of clauses~(3)
and~(5) of Proposition~\ref{10.1:prop-dimension-six}. Apart from $37$, $11$ and the $\Lambda^2$
entries of the last two columns, the entries follow from clauses~(1)--(5) of the dimension-six
classification, Proposition~\ref{10.1:prop-dimension-six}; the $\operatorname{Sym}^2V$ entry at
$\delta = 5$ vanishes because $-1\in W_4$, the product of the four coordinate reflections, acts on
$\mathcal{V}_\delta$ by $(-1)^\delta$ and trivially on $\operatorname{Sym}^2V$.

\begin{corollary}[The unique rotation-breaking scalar of dimension six]\label{10.1:cor-breaking-scalar}
Recall from Proposition~\ref{10.1:prop-dimension-six} the double-derivative density
\begin{equation}\label{10.1:eq-breaking-scalar-1}
Q=\sum_{\rho}\partial_\rho v^3_\rho=\sum_{\rho,\kappa}\partial_\rho^2\,
\operatorname{tr}F_{\rho\kappa}F_{\rho\kappa}.
\end{equation}
Among all local gauge-invariant
$W_4$-scalars of dimension at most six, the ones that are not $O(4)$-scalars are, modulo
$O(4)$-scalars:
\begin{itemize}
\item as densities, $O_1$ and the double derivative $Q$;
\item modulo total derivatives, $O_1$ alone;
\item modulo total derivatives and the field equation $J=0$, $O_1$ alone.
\end{itemize}
Let $A\in C^\infty_c(\mathbb{R}^4)$ and $U_\ell:=P\exp\int_\ell A$ on $\mathfrak{L}$, with $m$ large, as
in Proposition~\ref{10.1:prop-classical-limit}. As $a\to0$, the $a^2$-term of the Wilson action
$g_0^{-2}A(U)$ is
\begin{equation}\label{10.1:eq-breaking-scalar-2}
\frac{a^2}{24\,g_0^{2}}\int O_1 ,
\end{equation}
so that the leading lattice artefact of the Wilson action is the rotation-breaking operator
itself. Moreover, under the same hypothesis and as $a\to0$, the plaquette sum
$\sum_p\bigl[a^{-1}\xi^\rho(x_p)\mathfrak{r}_\rho(p)+\omega_{\rho\sigma}\mathfrak{s}_{\sigma\rho}(p)\bigr]$,
whose $c$- and $\chi$-weighted version $\mathfrak{S}_\chi[c]+\mathfrak{K}_\chi[c]$ is the negative of
the bulk term $-\mathfrak{S}_\chi[c]-\mathfrak{K}_\chi[c]$ of the rotational Ward identity
\eqref{10.1:eq-rotational-identity-a}, equals, at order $a^2$, the variation
$\frac{a^2}{24}\,\delta_\omega\int O_1$ of $\frac{a^2}{24}\int O_1$ under the rotation generated
by $\omega$, plus multiples of the field equation $J=0$.
\end{corollary}

\begin{proof}
By part (1) of the dimension-six classification (Proposition~\ref{10.1:prop-dimension-six}),
every $W_4$-scalar of dimension at most five is an $O(4)$-scalar, so the non-$O(4)$ scalars
all have dimension six. By part (5) of that proposition, as densities the space
$\mathcal{V}_6^{W_4}$ is spanned by $O_1,O_2,O_3$ and the three total derivatives
$\partial_\rho v^i_\rho$, of which only $O_1$ and $\partial_\rho v^3_\rho=Q$ fail to be
$O(4)$-scalars. Modulo total derivatives, the quotient of the $W_4$-scalars by the
$O(4)$-scalars is one-dimensional and spanned by $O_1$. Modulo the field equation as well,
$O_3$ drops while $O_1$ and the $O(4)$-scalar $O_2\simeq-4O_F$ remain independent, so $O_1$ still
breaks $O(4)$ and again spans the quotient; this proves the first assertion.

For the second, multiply clause (i) of the classical limit of the breaking term
(Proposition~\ref{10.1:prop-classical-limit}, display \eqref{10.1:eq-classical-limit-a}),
$A(U)=\int\bigl[-\tfrac14\sum_{\mu,\nu}\operatorname{tr}F_{\mu\nu}F_{\mu\nu}\bigr]
+\tfrac{a^2}{24}\int O_1+O(a^4)$, by $g_0^{-2}$; this gives \eqref{10.1:eq-breaking-scalar-2}.
The third assertion is clause (iii) of the same proposition: the plaquette sum equals
$V_\xi A(U)$, whose $a^2$-coefficient is
\begin{equation}\label{10.1:eq-breaking-scalar-3}
\begin{split}
\int\Bigl\{&\frac1{24}\sum_{\mu,\rho}\omega_{\rho\mu}B_{\mu\rho}\\
&+\sum_\mu\operatorname{tr}\Bigl(J_\mu\Bigl[\frac16\sum_\rho\xi^\rho D_\rho^2F_{\rho\mu}
-\frac1{12}\sum_\rho\omega_{\rho\mu}D_\mu F_{\rho\mu}\Bigr]\Bigr)\Bigr\}.
\end{split}
\end{equation}
There, by the same clause (iii), $\sum_{\mu,\rho}\omega_{\rho\mu}B_{\mu\rho}$ is the variation
of $O_1$ under the rotation generated by $\omega$, and the remaining two terms are multiples of $J$.
\end{proof}

\begin{proposition}[First-order covariant functionals depend only on strain]
\label{10.1:prop-strain-dependence}
Let $\eta$ be a vector field on $\mathbb{R}^4$, assigned dimension $-1$, and let
\begin{equation}\label{10.1:eq-strain-dependence-1}
\mathfrak{F}(\eta)=\int\bigl[\eta^\rho P_\rho+\partial_\nu\eta^\rho\,P_{\nu\rho}\bigr],
\end{equation}
where $P_\cdot\in\operatorname{Hom}_{W_4}(V,\mathcal{V}_5)$,
$P_{\cdot\cdot}\in\operatorname{Hom}_{W_4}(V\otimes V,\mathcal{V}_4)$, and repeated indices are
summed. This is the most general
$W_4\ltimes\mathbb{R}^4$-covariant local functional that is linear in $\eta$, has total dimension
at most $4$ and reads at most one derivative of $\eta$. Then there are numbers $c_1,c_2,c_3$
such that
\begin{equation}\label{10.1:eq-strain-dependence-2}
\mathfrak{F}(\eta)=\int\tfrac12\bigl(\partial_\nu\eta^\rho+\partial_\rho\eta^\nu\bigr)
\sum_{a=1}^{3}c_a\,\mathcal{T}^{(a)}_{\nu\rho}.
\end{equation}
In particular, for the Killing field $\xi(x)=\omega x$, with $\omega$ real antisymmetric, and for
the cutoff $\chi$, both as in the admissible data of
Theorem~\ref{10.1:thm-rotational-identity}, we have $\mathfrak{F}(\xi)=0$, and
\begin{equation}\label{10.1:eq-strain-dependence-3}
\mathfrak{F}(\chi\xi)=\int\xi^\rho\,\partial_\nu\chi\sum_{a=1}^{3}c_a\,\mathcal{T}^{(a)}_{\nu\rho}
\end{equation}
is a shell term, supported where $\chi$ is not locally constant.
\end{proposition}

\begin{proof}
We treat the two terms of \eqref{10.1:eq-strain-dependence-1} separately.

Consider first the term without derivatives of $\eta$. By part (3) of the dimension-six
classification
(Proposition~\ref{10.1:prop-dimension-six}), $\operatorname{Hom}_{W_4}(V,\mathcal{V}_5)$ is spanned
by $v^1_\rho$, $v^2_\rho$, $v^3_\rho$. Hence there are numbers $r_1,r_2,r_3$ with
$P_\rho=r_1v^1_\rho+r_2v^2_\rho+r_3v^3_\rho$. The same part of that proposition writes each basis
vector as the divergence of a symmetric tensor:
$v^1_\rho=\partial_\alpha\mathcal{T}^{(2)}_{\alpha\rho}$,
$v^3_\rho=\partial_\alpha\mathcal{T}^{(3)}_{\alpha\rho}$ and
$v^2_\rho=\partial_\alpha\bigl(\mathcal{T}^{(1)}_{\alpha\rho}-\tfrac14\mathcal{T}^{(2)}_{\alpha\rho}\bigr)$.
Consequently $P_\rho=\partial_\nu S_{\nu\rho}$, where
\begin{equation}\label{10.1:eq-strain-dependence-4}
S_{\nu\rho}:=r_2\,\mathcal{T}^{(1)}_{\nu\rho}+\bigl(r_1-\tfrac14r_2\bigr)\mathcal{T}^{(2)}_{\nu\rho}
+r_3\,\mathcal{T}^{(3)}_{\nu\rho}
\end{equation}
is symmetric in $\nu,\rho$. Integrating by parts, that is, discarding the total derivative
$\partial_\nu(\eta^\rho S_{\nu\rho})$, we obtain
$\int\eta^\rho P_\rho=-\int\partial_\nu\eta^\rho\,S_{\nu\rho}$.

Consider next the term with one derivative of $\eta$. As $W_4$-modules,
$V\otimes V=\operatorname{Sym}^2V\oplus\Lambda^2$. By part (2) of
Proposition~\ref{10.1:prop-dimension-six},
$\operatorname{Hom}_{W_4}(\Lambda^2,\mathcal{V}_4)=0$, so $P_{\cdot\cdot}$ vanishes on the antisymmetric
summand, that is, $P_{\nu\rho}=P_{\rho\nu}$. By part (4) of the same proposition,
$\operatorname{Hom}_{W_4}(\operatorname{Sym}^2V,\mathcal{V}_4)$ is spanned by
$\mathcal{T}^{(1)},\mathcal{T}^{(2)},\mathcal{T}^{(3)}$. Hence there are numbers $s_1,s_2,s_3$ with
$P_{\nu\rho}=\sum_a s_a\mathcal{T}^{(a)}_{\nu\rho}$.

Combining the two terms,
$\mathfrak{F}(\eta)=\int\partial_\nu\eta^\rho\,(P_{\nu\rho}-S_{\nu\rho})$, and $P-S$ is symmetric
with coefficients
\[
c_1=s_1-r_2,\qquad c_2=s_2-r_1+\tfrac14r_2,\qquad c_3=s_3-r_3
\]
on $\mathcal{T}^{(1)},\mathcal{T}^{(2)},\mathcal{T}^{(3)}$. Contracting a symmetric tensor with
$\partial_\nu\eta^\rho$ gives the same result as contracting it with the symmetric part
$\tfrac12(\partial_\nu\eta^\rho+\partial_\rho\eta^\nu)$. This proves
\eqref{10.1:eq-strain-dependence-2}.

For the Killing field, $\xi^\rho(x)=\omega_{\rho\sigma}x^\sigma$, so
$\partial_\nu\xi^\rho=\omega_{\rho\nu}$. This is antisymmetric in $\nu,\rho$ because $\omega$ is,
so the symmetric part in \eqref{10.1:eq-strain-dependence-2} vanishes identically and
$\mathfrak{F}(\xi)=0$.

For $\chi\xi$ we have
$\partial_\nu(\chi\xi^\rho)=\xi^\rho\partial_\nu\chi+\chi\,\omega_{\rho\nu}$. The second summand
is antisymmetric and drops out of the symmetric part, which is therefore
$\tfrac12(\xi^\rho\partial_\nu\chi+\xi^\nu\partial_\rho\chi)$. Contracting with the symmetric
tensor $\sum_ac_a\mathcal{T}^{(a)}_{\nu\rho}$ and renaming $\nu\leftrightarrow\rho$ in the second
half gives \eqref{10.1:eq-strain-dependence-3}. Its integrand contains the factor
$\partial_\nu\chi$, so it is supported where $\chi$ is not locally constant.
\end{proof}

\begin{remark}
Whatever a renormalisation step can create from the orbital remainder
$\mathfrak{K}_\chi[c]$ of the rotational Ward identity (Theorem~\ref{10.1:thm-rotational-identity})
at the marginal level is a coupling to the strain
$\tfrac12(\partial_\nu\eta^\rho+\partial_\rho\eta^\nu)$ of $\eta$.

In the bulk of a rotation this coupling vanishes. In the shell it changes the three coefficients
of the stress tensor in $\mathfrak{B}^{\mathrm{sh}}_\chi$. On a source, where the sourced
coefficient $c(p)=g_0^{-2}-w(p)$ of Theorem~\ref{10.1:thm-rotational-identity}, not to be confused
with the numbers $c_a$, depends on position, it leaves
\[
\int\eta^\rho\,\partial_\nu w\,\sum_{a=1}^{3}c_a\,\mathcal{T}^{(a)}_{\nu\rho}.
\]
The $\mathcal{T}^{(2)}$-part of this integral is a multiple of the left side of
\eqref{10.1:eq-rotational-identity-a}, and the other two parts are stress-tensor contact terms.

An infrared hypothesis must therefore be made for all components of the dimension-$4$ symmetric
tensor densities. The hypothesis of clause (vi) of Theorem~0 is the far-sphere flux hypothesis
$(\mathrm{IR}\text{-}\mathrm{fl})_{\omega_0}$: the flux of the lattice rotation current through
far spheres vanishes along one sequence of radii. A sufficient condition for it is its
stress-tensor decay form: the connected correlation of the curvature observables
$\mathcal{O}(f_1),\dots,\mathcal{O}(f_n)$ with one insertion of the stress-tensor density
$\Theta_{\nu\sigma}$ at a far-away point $x$ decays faster than $|x|^{-4}$, uniformly in $K$.
This sufficient form is stated component by component in
$\Theta_{\nu\sigma}$ and so meets this requirement.
\end{remark}

\begin{lemma}[Double subtraction for rotation-weighted vector densities]
\label{10.1:lem-double-subtraction}
Work on the unit lattice $\mathbb{Z}^4$ with unit vectors $e_1,\dots,e_4$ and forward differences
$(\nabla_\nu g)(x) := g(x+e_\nu)-g(x)$. For a site $x$ and an orientation $A=(\alpha\beta)$, $\alpha<\beta$,
let $(x;A)$ be the plaquette with corner $x$ spanned by $e_\alpha,e_\beta$; write $P\in A$ if $P$ has
orientation $A$, let $z(P)$ be the centre of $P$ and $|P|:=|z(P)|$, and let $P-y$ be the translate of $P$
by $-y$. Let $F$ be a complex $2$-cochain, $F_A(x):=F((x;A))$, with components extended to all index
pairs by $F_{\beta\alpha}=-F_{\alpha\beta}$; assume that $F$ is closed,
\begin{equation}\label{10.1:eq-double-subtraction-1}
\nabla_\rho F_{\nu\lambda}-\nabla_\nu F_{\rho\lambda}+\nabla_\lambda F_{\rho\nu}=0
\qquad\text{for all }\rho,\nu,\lambda,
\end{equation}
and that, for all orientations $A$ and sites $x,x'$,
\begin{equation*}
|F_A(x)|\le\varphi_0,\qquad |F_A(x)-F_A(x')|\le\varphi_1\,(1+|x-x'|)^2 .
\end{equation*}
Let $\delta>0$ and $\mathfrak{n}\ge0$, and for $\rho=1,\dots,4$ let $k_\rho$ be a kernel on pairs of
plaquettes, symmetric, $k_\rho(P,P')=k_\rho(P',P)$, with
$|k_\rho(P,P')|\le\mathfrak{n}\,e^{-\delta(|P|+|P'|)}$, and put
\begin{equation*}
Q_\rho(y;F):=\sum_{P,P'}k_\rho(P-y,P'-y)\,F(P)\,F(P').
\end{equation*}
Assume that $Q$ is of vector type: $Q_\rho(ry;r\cdot F)=\sum_{\rho'}(r_*)_{\rho\rho'}Q_{\rho'}(y;F)$ for
every lattice symmetry $r$ of $\mathbb{Z}^4$ and every $2$-cochain $F$; here $r x=r_*x+b$ with
$r_*\in W_4$ and $b\in\mathbb{Z}^4$, and $r$ acts on $2$-cochains as on $2$-forms,
$(r\cdot F)(rP)=\sigma_r(P)F(P)$ with $\sigma_r(P)=\pm1$ the sign by which $r$ changes the orientation
of $P$. Let $\xi(y)=\xi_0+\omega y$ with $\xi_0\in\mathbb{R}^4$ and $\omega$ an antisymmetric
$4\times4$ matrix. For a box $\Lambda=\{y\in\mathbb{Z}^4:\ m_i\le y_i\le n_i,\ i=1,\dots,4\}$, with
integers $m_i\le n_i$, let $\partial\Lambda$ be the set of sites that lie
in $\Lambda$ and have a nearest neighbour outside $\Lambda$, or lie outside $\Lambda$ and have a nearest
neighbour in $\Lambda$.

Then there are a constant $C(\delta)$, depending only on $\delta$, and a local function $H(y)$, bilinear in
$F$ (depending on $F$ only through kernels that decay exponentially, at rate $\delta/2$, away from $y$),
such that for every box $\Lambda$
\begin{equation}\label{10.1:eq-double-subtraction-2}
\sum_{y\in\Lambda}\sum_{\rho}\xi^\rho(y)\,Q_\rho(y;F)=\sum_{y\in\Lambda}H(y)+\mathsf{bd}(\Lambda),
\end{equation}
with
\begin{align*}
|H(y)|&\le C(\delta)\,\mathfrak{n}\,\bigl[|\omega|\varphi_0\varphi_1+(|\xi(y)|+|\omega|)\varphi_1^2\bigr],\\
|\mathsf{bd}(\Lambda)|&\le C(\delta)\,\mathfrak{n}\,\varphi_0^2\sum_{y\in\partial\Lambda}|\xi(y)|.
\end{align*}
\end{lemma}

\begin{proof}
Sums over $A,B$ run over the six orientations, sums over Greek indices over $1,\dots,4$; all sums are
written explicitly. Throughout, $C(\delta)$ denotes a constant depending only on $\delta$, which may change
from line to line. All series below converge absolutely by the exponential decay of the kernels and the
boundedness of $F$.

\emph{Step 1.} Fix a site $y$ and let $r$ be the inversion through $y$, $rx=2y-x$; it is a lattice symmetry
with $ry=y$, and $r_*$ is minus the identity. It reverses both edge directions of every plaquette, so
$\sigma_r(P)=(-1)^2=1$, and a constant cochain $F$ (that is, $F_A(x)=t_A$ for all $x$) satisfies
$r\cdot F=F$. The vector-type property gives $Q_\rho(y;F)=-Q_\rho(y;F)$, that is, $Q_\rho(y;F)=0$ for
constant cochains. This reads
\begin{equation*}
\sum_{A,B}t_At_B\sum_{P\in A,\,P'\in B}k_\rho(P,P')=0\qquad\text{for all }(t_A)\in\mathbb{C}^6.
\end{equation*}
Since $k_\rho$ is symmetric, the inner double sum is symmetric in $A,B$, and a symmetric bilinear form
whose quadratic form vanishes identically is zero (polarisation). Hence
\begin{equation}\label{10.1:eq-double-subtraction-3}
\sum_{P\in A,\,P'\in B}k_\rho(P,P')=0\qquad\text{for all }A,B.
\end{equation}

\emph{Step 2.} Put $\kappa_{\rho;AB}(u):=\sum_{P\in A}k_\rho\bigl(P,(u;B)\bigr)$ for $u\in\mathbb{Z}^4$.
By \eqref{10.1:eq-double-subtraction-3}, $\sum_u\kappa_{\rho;AB}(u)=0$; and since
$|(u;B)|\ge|u|-1$, $|\kappa_{\rho;AB}(u)|\le C(\delta)\,\mathfrak{n}\,e^{-\delta|u|}$. In
$Q_\rho(y;F)$ write, for $P\in A$ and $P'\in B$, $F(P)=F_A(y)+(F(P)-F_A(y))$ and
$F(P')=F_B(y)+(F(P')-F_B(y))$, and expand. The term $F_A(y)F_B(y)$ contributes zero by
\eqref{10.1:eq-double-subtraction-3}. The two mixed terms are equal by the symmetry of $k_\rho$; in their
sum put $P'=(y+u;B)$, so that the sum over $P\in A$ produces $\kappa_{\rho;AB}(u)$, and use
$\sum_u\kappa_{\rho;AB}(u)=0$ to drop $-F_B(y)$. Thus
\begin{equation}\label{10.1:eq-double-subtraction-4}
Q_\rho(y;F)=2\sum_{A,B}F_A(y)\sum_u\kappa_{\rho;AB}(u)F_B(y+u)+E^{(2)}_\rho(y),
\end{equation}
where
\begin{equation*}
\begin{split}
E^{(2)}_\rho(y):={}&\sum_{A,B}\sum_{P\in A,\,P'\in B}k_\rho(P-y,P'-y)\\
&\times\bigl(F(P)-F_A(y)\bigr)\bigl(F(P')-F_B(y)\bigr).
\end{split}
\end{equation*}
For $P=(x;A)$
one has $|F(P)-F_A(y)|\le\varphi_1(1+|x-y|)^2$ and $|x-y|\le|P-y|+1$, so
\begin{equation*}
|E^{(2)}_\rho(y)|\le\mathfrak{n}\,\varphi_1^2\sum_{P,P'}e^{-\delta(|P-y|+|P'-y|)}
(2+|P-y|)^2(2+|P'-y|)^2\le C(\delta)\,\mathfrak{n}\,\varphi_1^2 .
\end{equation*}

\emph{Step 3 (primitive).} A function $\kappa$ on $\mathbb{Z}^4$ with
$|\kappa(u)|\le C\mathfrak{n}e^{-\delta|u|}$ and zero sum is of the form
\begin{equation}\label{10.1:eq-double-subtraction-5}
\kappa(u)=\sum_\nu\bigl[\mu_\nu(u)-\mu_\nu(u-e_\nu)\bigr],\qquad |\mu_\nu(u)|\le C(\delta)\,\mathfrak{n}\,e^{-\delta|u|/2},
\end{equation}
obtained by summing along lines, coordinate by coordinate: with $\kappa^{(0)}:=\kappa$ and, for
$\nu=1,\dots,4$,
\begin{equation*}
\kappa^{(\nu)}(u):=
\begin{cases}
\sum_{t_1,\dots,t_\nu}\kappa(t_1,\dots,t_\nu,u_{\nu+1},\dots,u_4) & \text{if }u_1=\dots=u_\nu=0,\\
0 & \text{otherwise},
\end{cases}
\end{equation*}
one has $\kappa^{(4)}=0$ (zero sum), hence $\kappa=\sum_\nu(\kappa^{(\nu-1)}-\kappa^{(\nu)})$; for fixed
values of the other coordinates, $\kappa^{(\nu-1)}-\kappa^{(\nu)}$ has zero sum along the line in the
direction $e_\nu$, so it equals $\mu_\nu(u)-\mu_\nu(u-e_\nu)$ with $\mu_\nu(u)$ the sum of
$\kappa^{(\nu-1)}-\kappa^{(\nu)}$ over the points $u+te_\nu$, $t\le0$ (equivalently, minus the sum over
$t\ge1$). Since $|u|\ge|u|_1/2$, $|\kappa(u)|\le C\mathfrak{n}\prod_i e^{-\delta|u_i|/2}$, and the partial
sums keep this product bound, using the first expression for $u_\nu<0$ and the second for $u_\nu\ge0$; this
gives the bound in \eqref{10.1:eq-double-subtraction-5}.

Apply this to $\kappa_{\rho;AB}$, with primitive $\mu_{\nu\rho;AB}$. Summing by parts in $u$,
\begin{equation*}
\sum_u\kappa_{\rho;AB}(u)F_B(y+u)=\sum_\nu\sum_u\mu_{\nu\rho;AB}(u)\bigl[F_B(y+u)-F_B(y+u+e_\nu)\bigr],
\end{equation*}
so that, with $M_{\nu\rho;A}(y):=\sum_{B}\sum_u\mu_{\nu\rho;AB}(u)F_B(y+u)$,
\begin{equation}\label{10.1:eq-double-subtraction-6}
\sum_B\sum_u\kappa_{\rho;AB}(u)F_B(y+u)=-\sum_\nu(\nabla_\nu M_{\nu\rho;A})(y),
\qquad |M_{\nu\rho;A}|\le C(\delta)\,\mathfrak{n}\,\varphi_0 .
\end{equation}
Put $\hat\mu^0_{\nu\rho;AB}:=\sum_u\mu_{\nu\rho;AB}(u)$. Multiplying \eqref{10.1:eq-double-subtraction-5}
by $u_\nu$ and summing gives $\sum_u u_\nu\kappa(u)=-\sum_u\mu_\nu(u)$, since the term of index $\nu'$
contributes $-\sum_u\mu_{\nu'}(u)$ if $\nu'=\nu$ and zero otherwise; i.e.
$\hat\mu^0_{\nu\rho;AB}=-\sum_u u_\nu\kappa_{\rho;AB}(u)$ is the first moment of $\kappa$; in particular
$|\hat\mu^0|\le C(\delta)\mathfrak{n}$. It does not depend on the origins chosen for the plaquette positions:
replacing the corner $u$ of $(u;B)$ by its centre changes the moment by a constant vector times
$\sum_u\kappa_{\rho;AB}(u)=0$, and the positions of $P\in A$ are summed over in $\kappa$. Hence
\begin{equation*}
\hat\mu^0_{\nu\rho;AB}=-\sum_{P\in A,\,P'\in B}z(P')_\nu\,k_\rho(P,P').
\end{equation*}
Because $k_\rho$ is symmetric, it is determined by the quadratic form $Q_\rho(0;\cdot)$ (polarisation on
finitely supported cochains), and the vector-type property for $r\in W_4$ fixing the origin becomes
$\sigma_r(P)\sigma_r(P')k_\rho(rP,rP')=\sum_{\rho'}(r_*)_{\rho\rho'}k_{\rho'}(P,P')$. Since
$z(rP')=r_*z(P')$, the last display shows that $\hat\mu^0$ is a $W_4$-covariant constant tensor,
transforming as a vector in $\nu$ and in $\rho$ and as an element of $\Lambda^2$ in $A$ and in $B$.

\emph{Step 4 (Leibniz).} For functions $f,g$ on $\mathbb{Z}^4$,
$f\,\nabla_\nu g=\nabla_\nu(fg)-(\nabla_\nu f)\,g(\cdot+e_\nu)$. By
\eqref{10.1:eq-double-subtraction-6}, the first term of \eqref{10.1:eq-double-subtraction-4} (the cross
term) equals
\begin{equation*}
-2\sum_{\nu,A}\nabla_\nu\bigl(F_AM_{\nu\rho;A}\bigr)+2\sum_{\nu,A}(\nabla_\nu F_A)\,M_{\nu\rho;A}(\cdot+e_\nu).
\end{equation*}
Now $M_{\nu\rho;A}(y+e_\nu)=\sum_B\hat\mu^0_{\nu\rho;AB}F_B(y)+E^{M}_{\nu\rho;A}(y)$, where
\begin{equation*}
E^M_{\nu\rho;A}(y):=\sum_B\sum_u\mu_{\nu\rho;AB}(u)\bigl(F_B(y+e_\nu+u)-F_B(y)\bigr)
\end{equation*}
is bounded by
$\sum_u C(\delta)\mathfrak{n}e^{-\delta|u|/2}\varphi_1(2+|u|)^2\le C(\delta)\mathfrak{n}\varphi_1$; and
$|\nabla_\nu F_A|\le4\varphi_1$. Therefore the cross term is
\begin{equation}\label{10.1:eq-double-subtraction-7}
-2\sum_{\nu,A}\nabla_\nu\bigl(F_AM_{\nu\rho;A}\bigr)+Y_\rho+E^{(4)}_\rho,\qquad
Y_\rho:=2\sum_{\nu,A,B}\hat\mu^0_{\nu\rho;AB}(\nabla_\nu F_A)F_B,
\end{equation}
with $E^{(4)}_\rho:=2\sum_{\nu,A}(\nabla_\nu F_A)E^M_{\nu\rho;A}$ and
$|E^{(4)}_\rho|\le C(\delta)\mathfrak{n}\varphi_1^2$.

\emph{Step 5 ($Y$ is a divergence).} Split
$\hat\mu^0_{\nu\rho;AB}=\hat\mu^0_{\nu\rho;(AB)}+\hat\mu^0_{\nu\rho;[AB]}$ into its parts symmetric and
antisymmetric in $A,B$, and $Y_\rho=Y^s_\rho+Y^a_\rho$ accordingly. From
$\nabla_\nu(F_AF_B)=(\nabla_\nu F_A)F_B+F_A(\nabla_\nu F_B)+(\nabla_\nu F_A)(\nabla_\nu F_B)$,
\begin{equation*}
Y^s_\rho=\sum_{\nu,A,B}\hat\mu^0_{\nu\rho;(AB)}\bigl[\nabla_\nu(F_AF_B)-(\nabla_\nu F_A)(\nabla_\nu F_B)\bigr],
\end{equation*}
and the second part is bounded by $C(\delta)\mathfrak{n}\varphi_1^2$.

For the antisymmetric part, regard $\hat\mu^0$ as a $W_4$-invariant tensor with indices
$\nu,\rho,\alpha_1,\alpha_2,\beta_1,\beta_2$, where $A=(\alpha_1\alpha_2)$, $B=(\beta_1\beta_2)$. By the
parity rule for hyperoctahedral tensors (Lemma~\ref{10.1:lem-parity-rule}) a component vanishes unless
every index value occurs an even number of times. If $\nu=\rho$, the four values
$\alpha_1,\alpha_2,\beta_1,\beta_2$ must then each occur an even number of times; as $\alpha_1\ne\alpha_2$
and $\beta_1\ne\beta_2$, this forces $A=B$, where the antisymmetric part is zero. If $\nu\ne\rho$, the
values $\nu$ and $\rho$ must occur an odd number of times among $\alpha_1,\alpha_2,\beta_1,\beta_2$ and all
other values an even number of times; since no value occurs more than twice, the multiset is
$\{\nu,\rho,\lambda,\lambda\}$ with $\lambda\notin\{\nu,\rho\}$, and $\{A,B\}=\{(\nu\lambda),(\rho\lambda)\}$.
By the coordinate permutations in $W_4$ these entries are given by one constant $c$, $|c|\le C(\delta)\mathfrak{n}$,
so that, in terms of the antisymmetric components $F_{\nu\lambda}$,
\begin{equation*}
Y^a_\rho=c\sum_{\nu,\lambda}\bigl[(\nabla_\nu F_{\nu\lambda})F_{\rho\lambda}-(\nabla_\nu F_{\rho\lambda})F_{\nu\lambda}\bigr].
\end{equation*}
Applying the Leibniz rule of Step~4 to the first product,
\begin{align*}
Y^a_\rho={}&c\sum_{\nu,\lambda}\bigl[\nabla_\nu(F_{\nu\lambda}F_{\rho\lambda})-2F_{\nu\lambda}\nabla_\nu F_{\rho\lambda}\bigr]\\
&{}-c\sum_{\nu,\lambda}(\nabla_\nu F_{\nu\lambda})(\nabla_\nu F_{\rho\lambda}),
\end{align*}
the last sum being bounded by $C(\delta)\mathfrak{n}\varphi_1^2$. Since $F_{\nu\lambda}=-F_{\lambda\nu}$,
exchanging the names $\nu,\lambda$ in half of the sum and then using the closedness
\eqref{10.1:eq-double-subtraction-1} in the form $\nabla_\nu F_{\rho\lambda}-\nabla_\lambda F_{\rho\nu}=\nabla_\rho F_{\nu\lambda}$,
\begin{align*}
\sum_{\nu,\lambda}F_{\nu\lambda}\nabla_\nu F_{\rho\lambda}
&=\tfrac12\sum_{\nu,\lambda}F_{\nu\lambda}\bigl(\nabla_\nu F_{\rho\lambda}-\nabla_\lambda F_{\rho\nu}\bigr)
=\tfrac12\sum_{\nu,\lambda}F_{\nu\lambda}\nabla_\rho F_{\nu\lambda}\\
&=\tfrac14\nabla_\rho\sum_{\nu,\lambda}F_{\nu\lambda}^2-\tfrac14\sum_{\nu,\lambda}(\nabla_\rho F_{\nu\lambda})^2,
\end{align*}
the last step by $f\nabla_\rho f=\tfrac12\nabla_\rho(f^2)-\tfrac12(\nabla_\rho f)^2$. Hence
\begin{equation*}
\begin{split}
Y^a_\rho&=c\sum_\nu\nabla_\nu\Bigl(\sum_\lambda F_{\nu\lambda}F_{\rho\lambda}\Bigr)
-\tfrac{c}{2}\,\nabla_\rho\sum_{\alpha,\beta}F_{\alpha\beta}^2+E^{(5)}_\rho,\\
|E^{(5)}_\rho|&\le C(\delta)\,\mathfrak{n}\,\varphi_1^2 .
\end{split}
\end{equation*}
This is the lattice counterpart of part~(3) of the dimension-six classification
(Proposition~\ref{10.1:prop-dimension-six}).

\emph{Step 6.} Collecting \eqref{10.1:eq-double-subtraction-4}, \eqref{10.1:eq-double-subtraction-7} and
Step~5,
\begin{equation}\label{10.1:eq-double-subtraction-8}
Q_\rho(y;F)=\sum_\nu(\nabla_\nu\Sigma_{\nu\rho})(y)+\mathsf{r}_\rho(y),\qquad
\Sigma_{\nu\rho}:=-2\sum_AF_AM_{\nu\rho;A}+\sum_{A,B}s_{\nu\rho;AB}F_AF_B,
\end{equation}
where $s$ is defined, for $\nu\ne\rho$, by
\begin{equation*}
\sum_{A,B}s_{\nu\rho;AB}F_AF_B:=\sum_{A,B}\hat\mu^0_{\nu\rho;(AB)}F_AF_B
+c\sum_\lambda F_{\nu\lambda}F_{\rho\lambda},
\end{equation*}
and, for $\nu=\rho$, by the same expression with the further term
$-\tfrac{c}{2}\sum_{\alpha,\beta}F_{\alpha\beta}^2$ added; so $s$ is a $W_4$-covariant
constant tensor with $|s|\le C(\delta)\mathfrak{n}$, and $\mathsf{r}_\rho:=Q_\rho-\sum_\nu\nabla_\nu\Sigma_{\nu\rho}$
is the sum of the remainders of Steps~2, 4 and~5, bilinear in $F$ with $|\mathsf{r}_\rho|\le C(\delta)\mathfrak{n}\varphi_1^2$.
By \eqref{10.1:eq-double-subtraction-6}, $|\Sigma_{\nu\rho}|\le C(\delta)\mathfrak{n}\varphi_0^2$.

Since $\xi^\rho(y)=\xi_0^\rho+\sum_\beta\omega_{\rho\beta}y_\beta$, one has $\nabla_\nu\xi^\rho=\omega_{\rho\nu}$.
The Leibniz rule of Step~4 gives
\begin{equation*}
\xi^\rho\,\nabla_\nu\Sigma_{\nu\rho}=\nabla_\nu\bigl(\xi^\rho\Sigma_{\nu\rho}\bigr)
-\omega_{\rho\nu}\,\Sigma_{\nu\rho}(\cdot+e_\nu).
\end{equation*}
For the box $\Lambda$ of the statement and any $G$, the sum $\sum_{y\in\Lambda}(\nabla_\nu G)(y)$
telescopes to the sum of $G(y+e_\nu)$ over $y\in\Lambda$ with $y_\nu=n_\nu$ minus the sum of $G(y)$ over
$y\in\Lambda$ with $y_\nu=m_\nu$; the points $y+e_\nu$ and $y$ at which $G$ is evaluated there lie in
$\partial\Lambda$. Inserting
\eqref{10.1:eq-double-subtraction-8} into the left side of \eqref{10.1:eq-double-subtraction-2} we obtain
\eqref{10.1:eq-double-subtraction-2} with
\begin{equation*}
H(y):=\sum_\rho\xi^\rho(y)\mathsf{r}_\rho(y)-\sum_{\nu,\rho}\omega_{\rho\nu}\Sigma_{\nu\rho}(y+e_\nu),
\end{equation*}
and $\mathsf{bd}(\Lambda)$ the telescoped terms of $\sum_{\nu,\rho}\nabla_\nu(\xi^\rho\Sigma_{\nu\rho})$.
For each pair $\nu,\rho$ a point of $\partial\Lambda$ occurs in at most one of them, so each point of
$\partial\Lambda$ carries at most sixteen telescoped terms, each of modulus at most
$C(\delta)\mathfrak{n}\varphi_0^2|\xi|$ at that point; hence
$|\mathsf{bd}(\Lambda)|\le C(\delta)\mathfrak{n}\varphi_0^2\sum_{y\in\partial\Lambda}|\xi(y)|$.

It remains to bound $H$. The first sum is at most $C(\delta)\mathfrak{n}|\xi(y)|\varphi_1^2$. For the second,
write $d_A:=F_A(y+e_\nu)-F_A(y)$, so that $|d_A|\le4\varphi_1$, and expand
\begin{align*}
F_A(y+e_\nu)M_{\nu\rho;A}(y+e_\nu)&=F_A(y)\sum_B\hat\mu^0_{\nu\rho;AB}F_B(y)
+d_A\,M_{\nu\rho;A}(y+e_\nu)+F_A(y)\,E^M_{\nu\rho;A}(y),\\
F_A(y+e_\nu)F_B(y+e_\nu)&=F_A(y)F_B(y)+d_A\,F_B(y+e_\nu)+F_A(y)\,d_B .
\end{align*}
By $|F_A|\le\varphi_0$, the bound on $M$ in \eqref{10.1:eq-double-subtraction-6} and the bound on $E^M$
of Step~4, every term after the first on each line is at most $C(\delta)\mathfrak{n}\varphi_0\varphi_1$,
respectively $C\varphi_0\varphi_1$. Hence
\begin{equation*}
\begin{split}
\sum_{\nu,\rho}\omega_{\rho\nu}\Sigma_{\nu\rho}(y+e_\nu)
&=\sum_{\nu,\rho}\omega_{\rho\nu}\sum_{A,B}\bigl(-2\hat\mu^0_{\nu\rho;AB}+s_{\nu\rho;AB}\bigr)F_A(y)F_B(y)\\
&\quad+O\bigl(|\omega|\,\mathfrak{n}\,\varphi_0\varphi_1\bigr),
\end{split}
\end{equation*}
with the $O$-term bounded by $C(\delta)|\omega|\mathfrak{n}\varphi_0\varphi_1$. The main term depends only on
the part of $-2\hat\mu^0+s$ symmetric in $A,B$, and the linear map
$\omega\mapsto\sum_{\nu,\rho}\omega_{\rho\nu}(-2\hat\mu^0+s)_{\nu\rho;(AB)}$ is $W_4$-equivariant from
$\Lambda^2$ to $\operatorname{Sym}^2\Lambda^2$, because $\hat\mu^0$ and $s$ are $W_4$-covariant. To apply
part~(2) of the dimension-six classification (Proposition~\ref{10.1:prop-dimension-six}), let $\mathcal{F}$
denote the curvature variable of the continuum polynomials in $\mathcal{V}_4$, with values in the Lie
algebra $\mathfrak{su}(N)$ of $SU(N)$ (it is not the cochain $F$ of the lemma), and compose with the map
sending a symmetric $q=(q_{AB})$ to the dimension-four polynomial
$\sum_{A,B}q_{AB}\operatorname{tr}(\mathcal{F}_A\mathcal{F}_B)$. This map is $W_4$-equivariant and
injective: on $\mathcal{F}_A=t_A\,\mathsf{x}$ with $\mathsf{x}\in\mathfrak{su}(N)\setminus\{0\}$ it returns
$\operatorname{tr}(\mathsf{x}^2)\sum_{A,B}q_{AB}t_At_B$, and $\operatorname{tr}(\mathsf{x}^2)\ne0$. The
composition lies in $\operatorname{Hom}_{W_4}(\Lambda^2,\mathcal{V}_4)$, which is zero by that part. Hence
$\operatorname{Hom}_{W_4}(\Lambda^2,\operatorname{Sym}^2\Lambda^2)=0$, the main term vanishes, and
$|H(y)|\le C(\delta)\mathfrak{n}\bigl[|\omega|\varphi_0\varphi_1+|\xi(y)|\varphi_1^2\bigr]$, which implies the
bound claimed. Finally $H(y)$ is bilinear in $F$ and is built from $F$ through the kernels $k_\rho$ and
$\mu_{\nu\rho;AB}$, which decay exponentially away from $y$; so $H$ is local.
\end{proof}

\begin{definition}[Clustered tuples and their running coupling]\label{10.1:not-clustered-tuples}
In this section $N=2$, so that the gauge group is $SU(2)$, and we work in the regime of
clause~(iii) of Theorem~0. A number $\vartheta\in\,]0,1]$ is fixed. For a test function $f$ we put
\begin{equation*}
\|f\|_{\sharp}:=\max\{\|f\|_{\infty},\,40M\cdot\operatorname{Lip}f\}.
\end{equation*}
For $\omega\in\mathfrak{so}(4)$ we write $|\omega|^2:=\sum_{\mu<\nu}\omega_{\mu\nu}^2$.

For $n\ge 2$ and $\mathsf{d}>0$, a \emph{clustered tuple} is an element of the class
$\mathfrak{B}_n(\mathsf{d};\vartheta)$ below. Such a tuple comes with the centre $y$ of its ball.
The infinitesimal rotation $\delta_{\omega}$ and the action $R\cdot f$ of $R\in O(4)$ are taken
about this $y$; by translation invariance, any other centre may equally be used. We also fix the
depth $s(\mathsf{d})$ and the running coupling $g(\mathsf{d})$ at that depth, and the regime in
which all of this is used:
\begin{equation}\label{10.1:eq-clustered-tuples-a}
\begin{aligned}
&\mathfrak{B}_n(\mathsf{d};\vartheta):=\bigl\{\vec f=(f_1,\dots,f_n)\ \text{real}:\ 
\text{the }\operatorname{supp}f_i\text{ are pairwise disjoint},\\
&\qquad\qquad\operatorname{supp}f_i\subset\bar B(y,5\mathsf{d})\ \text{for one }y,\ 
\operatorname{Lip}f_i\le\|f_i\|_{\infty}/(\vartheta\mathsf{d})\ \text{for all }i\bigr\},\\
&\delta_{\omega}f:=-(\omega(x-y))\cdot\nabla f,\qquad R\cdot f:=f\circ R^{-1},\qquad
s(\mathsf{d}):=\lceil\log_L(1/\mathsf{d})\rceil,\\
&g(\mathsf{d})^{-2}:=\text{any }x\ \text{with}\ x\le g^{-2}+B^{\sharp}\,(s(\mathsf{d})+1),\\
&g\le\gamma_C:=\min\{\gamma_J,\gamma^{\otimes},\gamma_L\},\qquad
0<\mathsf{d}\le\tfrac{1}{40},\qquad m\ge m_{\mathbb{T}}(g_-(g)).
\end{aligned}
\end{equation}

In particular, $g(\mathsf{d})^{-2}$ may be taken to be the inverse running coupling
$x_{K-s(\mathsf{d})}$ of clause~(0) of Theorem~0 at any cutoff $K\ge s(\mathsf{d})$, or the limit
of $x_{K-s(\mathsf{d})}$ along the subsequence of cutoffs along which the continuum limit under
consideration is taken. Thus $g(\mathsf{d})$ is the running coupling at the depth of the tuple.

We put $\Pi:=\prod_{i=1}^{n}\|f_i\|_{\infty}$. The number $4-n$ is written out wherever it
occurs; the letter $p$ always denotes a plaquette.

The three conditions in the last line of \eqref{10.1:eq-clustered-tuples-a} are one-sided.
\begin{itemize}
\item The numbers $\gamma_J$, $\gamma^{\otimes}$ and $\gamma_L$ are, respectively, the ceilings on
the coupling of clause~(iii) of Theorem~0, of the stress-tensor-insertion correlation theorem
(Theorem~\ref{10.4:thm-insertion-correlation}), and of the ultraviolet half of clause~(vi) of
Theorem~0. They are chosen last.
\item The ceiling on $g$ can be removed at the price of a ceiling on $\mathsf{d}$ depending on
$g$, by the corollary on the exact change of reference scale.
\item The length $\mathsf{d}$ is measured in the units attached to the reference coupling $g$.
\item The condition on $m$, the exponent of the tori $\mathbb{T}_m$ of the family, is that of
hypothesis~(H5).
\end{itemize}

We fix the reference torus $\mathbb{T}_{\star}$, on which the ball $B(y,5)$ does not wrap around
the torus, so that it is isometric to a Euclidean ball. We fix a
radial cut-off $\chi$ built from the radial cut-off $\chi_1$ of unit radius; $\chi$ is identically
$1$ on $B(y,2)$ and is supported in $\bar B(y,4)$. We fix the rotation field $\xi$ and the lattice
shell functional $\mathfrak{B}^{\mathrm{sh}}_{\chi}$:
\begin{equation}\label{10.1:eq-clustered-tuples-b}
\begin{aligned}
&\mathbb{T}_{\star}:=\mathbb{T}_{m_{\mathbb{T}}(g_-(g))},\qquad
\chi:=\chi_1\bigl(|\cdot-y|/2\bigr),\qquad \xi:=\omega(x-y),\\
&\mathfrak{B}^{\mathrm{sh}}_{\chi}:=x_0\sum_p a_p\,(\mathsf{L}^{a}_{\xi}\chi)(x(p))
-x_0\sum_p\mathfrak{h}_{\chi}(p)+\operatorname{div}V_{\chi\xi}.
\end{aligned}
\end{equation}
The sums in \eqref{10.1:eq-clustered-tuples-b} run over the plaquettes $p$ of
$\mathbb{T}_{\star}$, and $x(p)$ is the barycentre of $p$.
Here $a=L^{-K}$ is the lattice spacing at the cutoff $K$, and $\mathsf{L}^{a}_{\xi}$ is the
lattice Lie derivative along $\xi$ at spacing $a$. The cut-off $\chi_1$ and the terms $a_p$,
$\mathfrak{h}_{\chi}$ and $V_{\chi\xi}$ are those of the shell functional in the rotational Ward
identity, used here as defined there. The last line of \eqref{10.1:eq-clustered-tuples-b} is the
exact decomposition of that functional, into a term carrying the lattice Lie derivative of the
cut-off, a plaquette term and a lattice divergence, given by the lemma on the exact decomposition
of the shell functional. The shell functional is used here in that form and is not proved again.
\end{definition}

\begin{lemma}[Unit-sphere shell bound blind to source positions]\label{10.1:lem-unit-shell-bound}
Let $y\in\mathbb{R}^4$ and let $\omega$ be a real antisymmetric $4\times 4$ matrix, with
$|\omega|:=(\sum_{\mu<\nu}\omega_{\mu\nu}^2)^{1/2}$. Let the reference
torus $\mathbb{T}_{\star}$, the radial cut-off $\chi$ and the shell functional
$\mathfrak{B}^{\mathrm{sh}}_{\chi}$ be as in \eqref{10.1:eq-clustered-tuples-b}
(Notation~\ref{10.1:not-clustered-tuples}), with $\xi:=\omega(x-y)$. Thus $\chi\equiv 1$ on $B(y,2)$
and $\operatorname{supp}\chi\subset\bar B(y,4)$. Write $m_{\star}:=m_{\mathbb{T}}(g_-(g))$, so that
$\mathbb{T}_{\star}=\mathbb{T}_{m_{\star}}$, and write $a:=L^{-K}$ for the lattice spacing at the cutoff.
The numbers $x_0$ and $x_K$ are the inverse couplings at the levels $0$ and $K$; at fixed $g$ the inverse
bare coupling $x_0=x_0(K)$ tends to infinity as $K\to\infty$. Finally, $P^{r}_K(0)$ is the
normalisation of the stress-tensor insertion $\mathfrak{T}^{r}$ in the channel
$r\in\{\mathbf{3},\mathbf{6}\}$.

Fix a positive integer $n$. In what follows $\vec F=(F_1,\dots,F_n)$ denotes an $n$-tuple of functions
in $C^{\infty}_c(B(y,1))$ with pairwise disjoint supports. For a lattice functional $A$ we abbreviate
\begin{equation*}
\langle A;\mathcal{O}(\vec F)\rangle_{K,m}
:=\langle A;\mathcal{O}(F_1);\dots;\mathcal{O}(F_n)\rangle_{K,m} .
\end{equation*}
Let $\mathsf{k}:=\mathsf{k}[\chi\xi]$ be the tensor weight of the shell. Let $\mathsf{k}^{\mathbf{3}}$ and
$\mathsf{k}^{\mathbf{6}}$ be its components in the channels $\mathbf{3}$ and $\mathbf{6}$, let
$\mathsf{k}^{\mathrm{diag}}=(\mathsf{k}_{\nu\nu})_{\nu}$ be its diagonal part, and let $\mathsf{a}'$ be its
antisymmetric partner.
Assume the following.
\begin{itemize}
\item[(a)] The exact channel decomposition of the shell functional. The weight $\mathsf{k}$ is traceless,
because $\chi$ is radial. The weights $\mathsf{k}$, $\mathsf{k}^{\mathbf{3}}$, $\mathsf{k}^{\mathbf{6}}$ and
$\mathsf{k}^{\mathrm{diag}}$ are supported in the shell $S:=\{x:2\le|x-y|\le 4\}$. There is a constant
$C_{\chi}$, independent of $K$, of $\omega$ and of $\vec F$, such that for every $K$
\begin{equation*}
\|\mathsf{k}^{\mathbf{3}}\|_{\sharp},\ \|\mathsf{k}^{\mathbf{6}}\|_{\sharp},\
\|\mathsf{k}^{\mathrm{diag}}\|_{\sharp}\le C_{\chi}\,|\omega| .
\end{equation*}
The finite identity
\begin{equation}\label{10.1:eq-unit-shell-bound-1}
\begin{split}
\mathfrak{B}^{\mathrm{sh}}_{\chi}
={}&-x_0\,\mathfrak{T}^{\mathbf{6}}(\mathsf{k}^{\mathbf{6}})
-2x_0\,\mathfrak{T}^{\mathbf{3}}(\mathsf{k}^{\mathbf{3}})
+\Bigl[x_0\sum_{\nu}\mathfrak{Q}_{\nu}(\mathsf{k}_{\nu\nu})\Bigr]
-\bigl[\mathbf{I}(\mathsf{a}'^{\mathsf{T}})\bigr]\\
&-\bigl[x_0\,\mathfrak{T}^{(2)}_{\chi\xi}[1]\bigr]
-\bigl[\tfrac14 x_0a^2\,\mathcal{O}(w_{\chi\xi})\bigr]
+\bigl[x_0\,\mathcal{O}(\mathsf{L}^{a}_{\xi}\chi)\bigr]
+\bigl[\operatorname{div}V_{\chi\xi}\bigr]
\end{split}
\end{equation}
holds. Here $\mathfrak{Q}_{\nu}$ are the diagonal unit insertions, and $\mathbf{I}(\mathsf{a}'^{\mathsf{T}})$
is the antisymmetric-channel term. The term $\mathfrak{T}^{(2)}_{\chi\xi}[1]$ is the second-order
stress-tensor insertion. The term $\mathcal{O}(w_{\chi\xi})$ is a curvature observable with a weight
$w_{\chi\xi}$ determined by $\chi$ and $\xi$. The term
$\mathcal{O}(\mathsf{L}^{a}_{\xi}\chi)$ is the curvature observable weighted by the lattice Lie derivative
of $\chi$ along $\xi$. Finally, $V_{\chi\xi}$ is the term of the shell functional whose lattice divergence
enters.
\item[(b)] The far form of the first-order insertion bound of the stress-tensor-insertion correlation
theorem, Theorem~\ref{10.4:thm-insertion-correlation}. Let $r\in\{\mathbf{3},\mathbf{6}\}$, let $K$
and $m$ be arbitrary in the range of that theorem, this range containing $m=m_{\star}$ and every $K$,
the coupling $g$ being in its regime, and let $\mathsf{k}'$ be a weight of the channel $r$. Let
$(F_1,\dots,F_n)$ be smooth test functions with pairwise disjoint supports, subject to no further
condition on their mutual positions, and put
$\mathsf{D}:=\operatorname{dist}(\operatorname{supp}\mathsf{k}',\bigcup_i\operatorname{supp}F_i)$. Then
\begin{equation}\label{10.1:eq-unit-shell-bound-2}
\bigl|\langle\mathfrak{T}^{r}(\mathsf{k}');\mathcal{O}(\vec F)\rangle_{K,m}\bigr|
\le C^{\otimes}_{n,r}(\mathsf{D})\,|P^{r}_K(0)|\,\|\mathsf{k}'\|_{\sharp}
\prod_{i=1}^{n}\|F_i\|_{\sharp},
\end{equation}
and $C^{\otimes}_{n,r}(\mathsf{D})$ is non-increasing in $\mathsf{D}$.
\item[(c)] The normalisation decay in the channels $\mathbf{3}$ and $\mathbf{6}$, in the form of
Lemma~\ref{10.4:lem-normalisation-decay}; in the channel $\mathbf{6}$ it is obtained from the one-loop
anisotropy inequality, Corollary~\ref{10.6:cor-anisotropy-inequality}. There are constants
$C_N$ and $X^{+}(g)$ such that, for $r\in\{\mathbf{3},\mathbf{6}\}$ and all $K$,
\begin{equation}\label{10.1:eq-unit-shell-bound-3}
x_0\,|P^{r}_K(0)|\le C_N\,x_K\le C_N\,X^{+}(g).
\end{equation}
\item[(d)] The ultraviolet estimate in the antisymmetric channel on the fixed torus $\mathbb{T}_{\star}$,
applied to the insertion $\mathbf{I}(\mathsf{a}'^{\mathsf{T}})$ and the tuples $\vec F$ of the statement:
for each such $\vec F$ there are numbers $\Gamma_K=\Gamma_K(\vec F)$ with
$|\langle\mathbf{I}(\mathsf{a}'^{\mathsf{T}});\mathcal{O}(\vec F)\rangle_{K,m_{\star}}|\le\Gamma_K$ and
$\Gamma_K\to 0$ as $K\to\infty$.
\item[(e)] Backward matching of the diagonal unit. One has
\begin{equation*}
x_0\sum_{\nu}\mathfrak{Q}_{\nu}(\mathsf{k}_{\nu\nu})
=C_{\mathfrak{Q}}\,\mathfrak{T}^{\mathbf{3}}(\mathsf{k}^{\mathbf{3}})+\theta .
\end{equation*}
Here $|C_{\mathfrak{Q}}|\le C'(\log x_0)^{2p_0}$, and the loads of $\theta$ at level $k$ are at
most $Cx_0\vartheta_1^{k}$, where $\vartheta_1$ is its rate. At each fixed $\vec F$ the correlation
$\langle\theta;\mathcal{O}(\vec F)\rangle_{K,m_{\star}}$ tends to zero as $K\to\infty$.
\item[(f)] Clause (iii-a) of Theorem~0. Each of the three terms $x_0\,\mathfrak{T}^{(2)}_{\chi\xi}[1]$,
$\tfrac14x_0a^2\,\mathcal{O}(w_{\chi\xi})$ and $x_0\,\mathcal{O}(\mathsf{L}^{a}_{\xi}\chi)$ carries a
factor $a$ or $a^2$. Against that factor, clause (iii-a) of Theorem~0 bounds its connected correlation
with $\mathcal{O}(\vec F)$, at fixed $\vec F$: the two terms carrying $a$, namely
$x_0\,\mathfrak{T}^{(2)}_{\chi\xi}[1]$ and $x_0\,\mathcal{O}(\mathsf{L}^{a}_{\xi}\chi)$, by a multiple of
$x_0a(\log x_0)^{c}$, with an exponent $c$, and the term carrying $a^2$ by a multiple of $x_0a^2$.
Moreover $x_0a(\log x_0)^{c}\to 0$ and $x_0a^2\to 0$ as $K\to\infty$.
\item[(g)] The divergence term. One has
$\operatorname{div}V_{\chi\xi}=-\tfrac32\,\mathfrak{T}^{\mathbf{3}}(\mathsf{k}^{\mathrm{diag}})
+a^2\,W_{\chi\xi}$, with a lattice functional $W_{\chi\xi}$ whose connected correlation with
$\mathcal{O}(\vec F)$ is bounded, at fixed $\vec F$, uniformly in $K$.
\end{itemize}
Then, for every $\vec F$ as above,
\begin{equation}\label{10.1:eq-unit-shell-bound-a}
\limsup_{K\to\infty}\,
\bigl|\langle\mathfrak{B}^{\mathrm{sh}}_{\chi};\mathcal{O}(F_1);\dots;\mathcal{O}(F_n)\rangle_{K,m_{\star}}\bigr|
\le C_{\mathrm{sh},n}(g)\,|\omega|\prod_{i=1}^{n}\|F_i\|_{\sharp},
\end{equation}
where
\begin{equation*}
C_{\mathrm{sh},n}(g):=C_{\chi}\,C_N X^{+}(g)\,
\bigl[C^{\otimes}_{n,\mathbf{6}}(1)+2\,C^{\otimes}_{n,\mathbf{3}}(1)\bigr].
\end{equation*}
The constant $C_{\mathrm{sh},n}(g)$ does not depend on the mutual positions of the supports of the $F_i$.
\end{lemma}

\begin{remark}
Hypothesis (b) is the far form of the bound of the stress-tensor-insertion correlation theorem,
Theorem~\ref{10.4:thm-insertion-correlation}. The proof of that theorem reads only the distance
$\mathsf{D}$ from the support of the weight to the union of the supports of the $F_i$; no lower bound on
the mutual distances of the supports of the $F_i$ enters it.
\end{remark}

\begin{proof}[Proof of Lemma~\ref{10.1:lem-unit-shell-bound}]
Fix $\vec F$ as in the statement.

\emph{Expansion.} The connected correlation $\langle A;\mathcal{O}(\vec F)\rangle_{K,m_{\star}}$ is linear
in its first entry $A$. We insert the identity \eqref{10.1:eq-unit-shell-bound-1} of hypothesis (a) into
this entry and define the six bracket correlations
\begin{align*}
\beta_1(K)&:=x_0\Bigl\langle\sum_{\nu}\mathfrak{Q}_{\nu}(\mathsf{k}_{\nu\nu});
\mathcal{O}(\vec F)\Bigr\rangle_{K,m_{\star}},\\
\beta_2(K)&:=-\langle\mathbf{I}(\mathsf{a}'^{\mathsf{T}});\mathcal{O}(\vec F)\rangle_{K,m_{\star}},\\
\beta_3(K)&:=-x_0\langle\mathfrak{T}^{(2)}_{\chi\xi}[1];\mathcal{O}(\vec F)\rangle_{K,m_{\star}},\\
\beta_4(K)&:=-\tfrac14x_0a^2\langle\mathcal{O}(w_{\chi\xi});\mathcal{O}(\vec F)\rangle_{K,m_{\star}},\\
\beta_5(K)&:=x_0\langle\mathcal{O}(\mathsf{L}^{a}_{\xi}\chi);\mathcal{O}(\vec F)\rangle_{K,m_{\star}},\\
\beta_6(K)&:=\langle\operatorname{div}V_{\chi\xi};\mathcal{O}(\vec F)\rangle_{K,m_{\star}}.
\end{align*}
With them we obtain, for every $K$,
\begin{equation}\label{10.1:eq-unit-shell-bound-4}
\begin{split}
\langle\mathfrak{B}^{\mathrm{sh}}_{\chi};\mathcal{O}(\vec F)\rangle_{K,m_{\star}}
={}&-x_0\langle\mathfrak{T}^{\mathbf{6}}(\mathsf{k}^{\mathbf{6}});\mathcal{O}(\vec F)\rangle_{K,m_{\star}}
-2x_0\langle\mathfrak{T}^{\mathbf{3}}(\mathsf{k}^{\mathbf{3}});\mathcal{O}(\vec F)\rangle_{K,m_{\star}}\\
&+\sum_{j=1}^{6}\beta_j(K).
\end{split}
\end{equation}

\emph{The distance from the shell to the sources.} By hypothesis (a) the weights $\mathsf{k}^{\mathbf{3}}$,
$\mathsf{k}^{\mathbf{6}}$ and $\mathsf{k}^{\mathrm{diag}}$ are supported in $S$. Each $F_i$ is supported
in $B(y,1)$. Let $x\in S$ and $z\in\bigcup_i\operatorname{supp}F_i$. The triangle inequality gives
\begin{equation*}
|x-z|\ge|x-y|-|z-y|>2-1=1 .
\end{equation*}
Hence the distance $\mathsf{D}$ from the support of each of these weights to
$\bigcup_i\operatorname{supp}F_i$ is at
least $1$. This uses only that each $\operatorname{supp}F_i$ lies in $B(y,1)$, and it holds whatever the
mutual positions of the supports of the $F_i$. So hypothesis (b) applies to these weights with
$m=m_{\star}$, and by the monotonicity in hypothesis (b),
$C^{\otimes}_{n,r}(\mathsf{D})\le C^{\otimes}_{n,r}(1)$.

\emph{The main terms.} Let $r\in\{\mathbf{3},\mathbf{6}\}$. Hypothesis (b) with $\mathsf{k}'=\mathsf{k}^{r}$,
together with $C^{\otimes}_{n,r}(\mathsf{D})\le C^{\otimes}_{n,r}(1)$, followed by the normalisation decay
\eqref{10.1:eq-unit-shell-bound-3} of hypothesis (c) and the bound on $\|\mathsf{k}^{r}\|_{\sharp}$ in
hypothesis (a), gives for every $K$
\begin{equation}\label{10.1:eq-unit-shell-bound-5}
\begin{split}
x_0\bigl|\langle\mathfrak{T}^{r}(\mathsf{k}^{r});\mathcal{O}(\vec F)\rangle_{K,m_{\star}}\bigr|
&\le C^{\otimes}_{n,r}(1)\,x_0|P^{r}_K(0)|\,\|\mathsf{k}^{r}\|_{\sharp}\prod_{i=1}^{n}\|F_i\|_{\sharp}\\
&\le C^{\otimes}_{n,r}(1)\,C_NX^{+}(g)\,C_{\chi}\,|\omega|\prod_{i=1}^{n}\|F_i\|_{\sharp}.
\end{split}
\end{equation}
We take \eqref{10.1:eq-unit-shell-bound-5} with $r=\mathbf{6}$, and twice with $r=\mathbf{3}$. The first
two terms on the right of \eqref{10.1:eq-unit-shell-bound-4} are then together bounded in absolute value
by $C_{\mathrm{sh},n}(g)\,|\omega|\prod_i\|F_i\|_{\sharp}$, for every $K$.

\emph{The six brackets.} We show that $\beta_j(K)\to0$ as $K\to\infty$ for $j=1,\dots,6$, at the fixed
$\vec F$.

The bracket $\beta_1$ contains the diagonal unit. By hypothesis (e) and the linearity of the connected
correlation in its first entry,
\begin{equation*}
\beta_1(K)=C_{\mathfrak{Q}}\,\langle\mathfrak{T}^{\mathbf{3}}(\mathsf{k}^{\mathbf{3}});
\mathcal{O}(\vec F)\rangle_{K,m_{\star}}
+\langle\theta;\mathcal{O}(\vec F)\rangle_{K,m_{\star}} .
\end{equation*}
The weight $\mathsf{k}^{\mathbf{3}}$ is supported in $S$. As for the main terms, hypothesis (b) with
$r=\mathbf{3}$ and $C^{\otimes}_{n,\mathbf{3}}(\mathsf{D})\le C^{\otimes}_{n,\mathbf{3}}(1)$, the bound
\eqref{10.1:eq-unit-shell-bound-3} divided by $x_0$, the bound on $\|\mathsf{k}^{\mathbf{3}}\|_{\sharp}$ in
hypothesis (a) and the bound on $|C_{\mathfrak{Q}}|$ in hypothesis (e) give
\begin{equation*}
\bigl|C_{\mathfrak{Q}}\,\langle\mathfrak{T}^{\mathbf{3}}(\mathsf{k}^{\mathbf{3}});
\mathcal{O}(\vec F)\rangle_{K,m_{\star}}\bigr|
\le C'(\log x_0)^{2p_0}\,x_0^{-1}\,C^{\otimes}_{n,\mathbf{3}}(1)\,C_NX^{+}(g)\,C_{\chi}\,|\omega|
\prod_{i=1}^{n}\|F_i\|_{\sharp}.
\end{equation*}
Since $x_0\to\infty$ as $K\to\infty$ and $(\log x_0)^{2p_0}x_0^{-1}\to0$ as $x_0\to\infty$, this term
tends to zero. The second term tends to zero by hypothesis (e). Hence $\beta_1(K)\to0$.

The bracket $\beta_2$ is the antisymmetric-channel term. By hypothesis (d),
$|\beta_2(K)|\le\Gamma_K$ and $\Gamma_K\to0$.

The brackets $\beta_3$, $\beta_4$ and $\beta_5$ carry the factor $a$ or $a^2$. By hypothesis (f), clause
(iii-a) of Theorem~0 bounds $\beta_3$ and $\beta_5$ by a multiple of $x_0a(\log x_0)^{c}$ and $\beta_4$ by
a multiple of $x_0a^2$. Since
$x_0a(\log x_0)^{c}\to0$ and $x_0a^2\to0$ as $K\to\infty$, the three brackets tend to zero.

The bracket $\beta_6$ is the divergence term. By hypothesis (g),
\begin{equation*}
\beta_6(K)=-\tfrac32\langle\mathfrak{T}^{\mathbf{3}}(\mathsf{k}^{\mathrm{diag}});
\mathcal{O}(\vec F)\rangle_{K,m_{\star}}
+a^2\langle W_{\chi\xi};\mathcal{O}(\vec F)\rangle_{K,m_{\star}} .
\end{equation*}
The correlation in the second term is bounded uniformly in $K$, and $a^2=L^{-2K}\to0$, so the second
term tends to zero. The weight $\mathsf{k}^{\mathrm{diag}}$ is supported in $S$. So the
distance estimate above and hypothesis (b), with $r=\mathbf{3}$ and
$C^{\otimes}_{n,\mathbf{3}}(\mathsf{D})\le C^{\otimes}_{n,\mathbf{3}}(1)$, apply to the first
term. Dividing \eqref{10.1:eq-unit-shell-bound-3} by $x_0$ and using the bound on
$\|\mathsf{k}^{\mathrm{diag}}\|_{\sharp}$ in hypothesis (a) then gives
\begin{equation}\label{10.1:eq-unit-shell-bound-6}
\begin{split}
\tfrac32\bigl|\langle\mathfrak{T}^{\mathbf{3}}(\mathsf{k}^{\mathrm{diag}});
\mathcal{O}(\vec F)\rangle_{K,m_{\star}}\bigr|
&\le\tfrac32\,C^{\otimes}_{n,\mathbf{3}}(1)\,|P^{\mathbf{3}}_K(0)|\,
\|\mathsf{k}^{\mathrm{diag}}\|_{\sharp}\prod_{i=1}^{n}\|F_i\|_{\sharp}\\
&\le\tfrac32\,C^{\otimes}_{n,\mathbf{3}}(1)\,C_NX^{+}(g)\,x_0^{-1}\,
C_{\chi}\,|\omega|\prod_{i=1}^{n}\|F_i\|_{\sharp}.
\end{split}
\end{equation}
Since $x_0\to\infty$ as $K\to\infty$ at fixed $g$, the right side tends to zero. Hence
$\beta_6(K)\to0$.

\emph{Conclusion.} We apply the triangle inequality to \eqref{10.1:eq-unit-shell-bound-4} and use the
bound on the main terms. For every $K$ this gives
\begin{equation*}
\bigl|\langle\mathfrak{B}^{\mathrm{sh}}_{\chi};\mathcal{O}(\vec F)\rangle_{K,m_{\star}}\bigr|
\le C_{\mathrm{sh},n}(g)\,|\omega|\prod_{i=1}^{n}\|F_i\|_{\sharp}+\sum_{j=1}^{6}|\beta_j(K)| .
\end{equation*}
The upper limit of a sum is at most the sum of the upper limits, and $\limsup_K|\beta_j(K)|=0$ for each
$j$. Taking $\limsup_{K\to\infty}$ therefore gives \eqref{10.1:eq-unit-shell-bound-a}.

The constant $C_{\mathrm{sh},n}(g)$ is built from $C_{\chi}$, $C_N$, $X^{+}(g)$ and the far-form
constants at distance $1$. None of these involves the mutual positions of the supports of the $F_i$.
Indeed, $C_{\chi}$ is independent of $\vec F$ by hypothesis (a), hypothesis (b) places no condition on
those positions, and the distance estimate uses only that
each support lies in $B(y,1)$.
\end{proof}

\begin{proposition}[Rotation defect of two- and three-point functions]
\label{10.1:prop-rotation-defect-two-three}
Work in the regime and with the notation fixed in Notation~\ref{10.1:not-clustered-tuples}. From
\eqref{10.1:eq-clustered-tuples-a} we use the class $\mathfrak{B}_n(\mathsf{d};\vartheta)$ of real
$n$-tuples, the infinitesimal rotations $\delta_{\omega}$ with $\omega\in\mathfrak{so}(4)$ and
$|\omega|^2=\sum_{\mu<\nu}\omega_{\mu\nu}^2$, the action $R\cdot f$ of $R\in O(4)$ about the
centre $y$ of the common ball, the depth $s(\mathsf{d})$ and the running coupling $g(\mathsf{d})$.
From \eqref{10.1:eq-clustered-tuples-b} we use the reference torus
$\mathbb{T}_{\star}=\mathbb{T}_{m_\star}$, $m_\star:=m_{\mathbb{T}}(g_-(g))$, the radial cut-off
$\chi$ and the shell functional $\mathfrak{B}^{\mathrm{sh}}_{\chi}$. For an $n$-tuple
$\vec f=(f_1,\dots,f_n)$ put $\Pi:=\prod_{i=1}^n\|f_i\|_\infty$, and for a continuum limit $\sigma$
write $S_n[\sigma]$ for its connected $n$-point Schwinger function. Assume the following.
\begin{itemize}
\item The ultraviolet half of clause (vi) of Theorem~0 on the one torus $\mathbb{T}_\star$: for
every smooth $n$-tuple $\vec F$ with pairwise disjoint supports in $B(y,1)$, the remainder term
$\mathsf{Br}_K(\vec F)$ of the rotational Ward identity on $\mathbb{T}_\star$ with the shell
functional $\mathfrak{B}^{\mathrm{sh}}_{\chi}$ tends to zero as $K\to\infty$. In the antisymmetric
(spin) channel this is supplied by the spin-channel ultraviolet estimate on every fixed torus; in the
vector channel it is
supplied by the reduction theorem, Theorem~\ref{10.5:thm-reduction}, whose proof is incomplete at
one step.
\item The first-order stress-tensor insertion bound of the stress-tensor-insertion correlation
theorem, Theorem~\ref{10.4:thm-insertion-correlation}, read in its far form, in the form in which it
yields the unit-sphere shell bound \eqref{10.1:eq-unit-shell-bound-a} of
Lemma~\ref{10.1:lem-unit-shell-bound}.
\item The volume-transfer theorem of clause (iv) of Theorem~0.
\item Clauses (0), (ii) and (iii) of Theorem~0, and the one-loop law, part (d) of clause (0).
\end{itemize}
No infrared hypothesis is assumed; in particular the far-sphere flux hypothesis
$(\mathrm{IR\text{-}fl})_{\omega_0}$ of clause (vi) is not among the hypotheses.

Let $n\in\{2,3\}$. Let $\sigma$ be any subsequential limit of clause (ii) of Theorem~0 at reference
coupling $g$: any limit on $\mathbb{R}^4$ of clause (ii-$\mathbb{R}^4$), or any torus limit of
clause (ii-$\mathbb{T}$-b) on any torus of the family. Let $\vec f\in\mathfrak{B}_n(\mathsf{d};\vartheta)$
have components of class $C^1$, and let $R\in O(4)$. Then the following hold.
\begin{enumerate}
\item[(a)] The rotation defect obeys
\begin{equation}\label{10.1:eq-rotation-defect-two-three-c}
\bigl|S_n[\sigma](R\cdot\vec f)-S_n[\sigma](\vec f)\bigr|
\le C_n(g,\vartheta)\,\mathsf{d}^{4-n}\,\Pi ,
\end{equation}
with
\begin{equation*}
C_n(g,\vartheta):=\bigl(C_{\mathrm{rot},n}+2C_{T,n}\bigr)(g)
=160^4\Bigl(\frac{40M}{\vartheta}\Bigr)^{n}
\bigl[\pi\sqrt{2}\,C_{\mathrm{sh},n}(g)+2C_{\mathrm{VT},n}(g)\bigr],
\end{equation*}
where $C_{\mathrm{sh},n}$ is the constant of \eqref{10.1:eq-unit-shell-bound-a},
$C_{\mathrm{rot},n}$ that of \eqref{10.1:eq-signed-array-b}, $C_{T,n}$ that of
\eqref{10.1:eq-transfer-coupling-form-a}, and $C_{\mathrm{VT},n}$ the constant of the
volume-transfer theorem. In native form,
\begin{equation}\label{10.1:eq-rotation-defect-two-three-1}
\bigl|S_n[\sigma](R\cdot\vec f)-S_n[\sigma](\vec f)\bigr|
\le 160^4\bigl[\pi\sqrt{2}\,C_{\mathrm{sh},n}(g)+2C_{\mathrm{VT},n}(g)\bigr]\,
\mathsf{d}^{4}\prod_{i=1}^n\|f_i\|_{\sharp};
\end{equation}
here the power $4$ of $\mathsf{d}$ does not depend on $n$, and the factor $\mathsf{d}^{-n}$ in
\eqref{10.1:eq-rotation-defect-two-three-c} is the cost of the $n$ Lipschitz norms, since
$\|f_i\|_\sharp\le(40M/(\vartheta\mathsf{d}))\|f_i\|_\infty$ for
$\vec f\in\mathfrak{B}_n(\mathsf{d};\vartheta)$.
\item[(b)] On the lattice, uniformly in the cutoff and in the volume,
\begin{equation*}
\limsup_{K\to\infty}\ \sup_{m\ge m_{\mathbb{T}}(g_-(g))}
\bigl|S^T_{n;K,m}(R\cdot\vec f)-S^T_{n;K,m}(\vec f)\bigr|
\le C_n(g,\vartheta)\,\mathsf{d}^{4-n}\,\Pi .
\end{equation*}
\item[(c)] In infinitesimal form: if the $f_i$ are of class $C^2$, then for every
$\omega\in\mathfrak{so}(4)$
\begin{equation*}
\Bigl|\sum_{i=1}^n S_n[\sigma](f_1,\dots,\delta_\omega f_i,\dots,f_n)\Bigr|
\le C^{D}_n(g,\vartheta)\,|\omega|\,\mathsf{d}^{4-n}\,\Pi,
\end{equation*}
with
\begin{equation*}
C^{D}_n(g,\vartheta):=160^4\Bigl(\frac{40M}{\vartheta}\Bigr)^{n}
\Bigl[C_{\mathrm{sh},n}(g)+\frac{6n}{\vartheta}\,C_{\mathrm{VT},n}(g)\Bigr].
\end{equation*}
\item[(d)] In terms of the running coupling (this uses clause (0) and its part (d)): for every
$\eta\in\,]0,1[$ the right-hand sides of (a)--(c) are bounded by
\begin{equation*}
A_{n,\eta}(g,\vartheta)\,
\exp\Bigl[-(1-\eta)\,(4-n)\,\frac{12\pi^2}{11N^2}\,g(\mathsf{d})^{-2}\Bigr]\,\Pi .
\end{equation*}
From clause (0) alone they are bounded by
\begin{equation*}
A'_n(g,\vartheta)\exp\Bigl[-(4-n)\,\frac{\ln L}{B^\sharp}\,g(\mathsf{d})^{-2}\Bigr]\,\Pi .
\end{equation*}
Equivalently, by the elementary optimisation inequality $e^{-cx}\le k!\,(cx)^{-k}$, valid for
$c,x>0$ and every integer $k\ge 0$, there are constants $A_n(g,\vartheta)$ and $c_n>0$, the same
for every integer $k\ge 0$, such that these right-hand sides are at most
$A_n(g,\vartheta)\,k!\,c_n^{-k}\,g(\mathsf{d})^{2k}\,\Pi$: the rotation defect is smaller than every
power of the running coupling, with factorial constants. At $N=2$ the exponent is
$(1-\eta)(4-n)\,3\pi^2/11$ times $g(\mathsf{d})^{-2}$. In the standard normalisation
$g_{\mathrm{std}}^2=2N\,g(\mathsf{d})^{2}$ of the running coupling (that is, $2N/x$ with
$x=g(\mathsf{d})^{-2}$ the inverse coupling) the exponent reads
\begin{equation*}
(1-\eta)\,(4-n)\,\frac{8\pi^2}{\beta_0\,g_{\mathrm{std}}^2},\qquad \beta_0=\frac{11N}{3},
\end{equation*}
so that the bound is the $(4-n)$-th power of the one-loop transmutation length
$\exp(-8\pi^2/(\beta_0 g_{\mathrm{std}}^2))$, up to the factor $1-\eta$ in the exponent.
\item[(e)] In relative form, for $n=2$: on the pairs of class $C^1$ in the class
$\mathfrak{T}(\mathsf{d};\vartheta)$ of clause (iv) of Theorem~0, under the three conditions displayed
in clause (iv) --- the condition on the coupling, the depth condition and the ultraviolet
condition ---
\begin{equation*}
\Bigl|\frac{S_2[\sigma](R\cdot f)}{S_2[\sigma](f)}-1\Bigr|
\le \frac{9}{2}\,(b_0c_{\mathcal{K}})^{-1}\,C_2(g,\vartheta)\,\mathsf{d}^2
\bigl(g^{-2}+B^\sharp(s(\mathsf{d})+1)\bigr)^2
=O\bigl(\mathsf{d}^2\log^2(1/\mathsf{d})\bigr).
\end{equation*}
\item[(f)] On expansions: if for one $R_0\in O(4)$ the function $s\mapsto S_n[\sigma](R_0\cdot
\vec f^{(s)})$, where $\vec f^{(s)}$ is a tuple of fixed shape dilated to depth $s$, has an asymptotic
expansion in powers of $g(L^{-s})^2$ to some order, then it has the same expansion for every
$R\in O(4)$. No expansion is asserted to exist.
\end{enumerate}
\end{proposition}

The proposition does not assert the following: exact rotation invariance at any fixed
$\mathsf{d}$, which is equivalent to the far-sphere flux hypothesis $(\mathrm{IR\text{-}fl})_{\omega_0}$
and fails for every torus limit (by the proposition on torus limits of this chapter); a relative
form for $n=3$, for which no lower bound on $|S_3[\sigma]|$ is available; anything for $n\ge 4$
(see the statement for every $n$ in this chapter); that the bounds are sharp; the existence of an
asymptotic expansion beyond leading order.

\begin{proof}[Proof of Proposition~\ref{10.1:prop-rotation-defect-two-three}: the flux form on the
torus]
Throughout the proof, $a:=L^{-K}$ denotes the lattice spacing at cutoff $K$, $\xi:=\omega(x-y)$ is
the rotation field about $y$, and $\mathsf{L}^{a}_{\xi}$ is the lattice Lie derivative along $\xi$ at
spacing $a$. For an $n$-tuple $\vec F$ we write $\mathcal{O}(\vec F)$ for the family of curvature
observables $\mathcal{O}(F_1),\dots,\mathcal{O}(F_n)$, and
$\langle\mathfrak{B}^{\mathrm{sh}}_{\chi};\mathcal{O}(\vec F)\rangle_{K,m_\star}$ for the connected
correlation, at cutoff $K$ on $\mathbb{T}_\star$, of the shell functional with these observables.

In this part of the proof we establish the flux form \eqref{10.1:eq-rotation-defect-two-three-a} of
the defect on the torus $\mathbb{T}_\star$. Let $\vec F$ be
a smooth $n$-tuple with pairwise disjoint supports in $B(y,1)$; recall from
\eqref{10.1:eq-clustered-tuples-b} that $\chi\equiv 1$ on $B(y,2)$, so the supports lie where the
cut-off equals one. The rotational Ward identity on $\mathbb{T}_\star$, which is exact at every
cutoff, states
\begin{equation}\label{10.1:eq-rotation-defect-two-three-2}
\sum_{i=1}^n S^T_{n;K,m_\star}\bigl(F_1,\dots,\mathsf{L}^{a}_{\xi}F_i,\dots,F_n\bigr)
=\bigl\langle\mathfrak{B}^{\mathrm{sh}}_{\chi};\mathcal{O}(\vec F)\bigr\rangle_{K,m_\star}
+\mathsf{Br}_K(\vec F).
\end{equation}
The lattice Lie derivative differs from the continuum derivative along $\xi$ by
\begin{equation*}
\bigl\|\mathsf{L}^{a}_{\xi}F-\xi\cdot\nabla F\bigr\|_{\sharp}\le C\,a^{2}\,|\omega|\,\|F\|_{C^4},
\end{equation*}
with a constant $C$ independent of $a$, $\omega$ and $F$,
and clause (iii-a) of Theorem~0 is used for the error term which this produces in the left-hand side
of \eqref{10.1:eq-rotation-defect-two-three-2}; since $a=L^{-K}\to 0$, this error term tends to zero as
$K\to\infty$. By the first hypothesis (the ultraviolet half of clause (vi) on the fixed torus
$\mathbb{T}_\star$), $\mathsf{Br}_K(\vec F)\to 0$ as $K\to\infty$ for each fixed $\vec F$.

Now let $\sigma_\star$ be a limit of clause (ii-$\mathbb{T}$-b) of Theorem~0 on $\mathbb{T}_\star$,
along a sequence of cutoffs $K_j\to\infty$. Along this sequence the left-hand side of
\eqref{10.1:eq-rotation-defect-two-three-2}, with $\mathsf{L}^{a}_{\xi}F_i$ replaced by
$\xi\cdot\nabla F_i$, converges to
\begin{equation*}
\sum_{i=1}^n S_n[\sigma_\star]\bigl(F_1,\dots,\xi\cdot\nabla F_i,\dots,F_n\bigr).
\end{equation*}
Since
$\delta_\omega F_i=-(\omega(x-y))\cdot\nabla F_i=-\xi\cdot\nabla F_i$, this limit equals
$-D_\omega[\sigma_\star](\vec F)$, where
\begin{equation}\label{10.1:eq-rotation-defect-two-three-b}
D_\omega(\vec F):=\sum_{i=1}^n S_n[\sigma_\star]\bigl(F_1,\dots,\delta_\omega F_i,\dots,F_n\bigr).
\end{equation}
Combining the three limits in \eqref{10.1:eq-rotation-defect-two-three-2}, the limit of the shell
term exists along the sequence and
\begin{equation}\label{10.1:eq-rotation-defect-two-three-a}
\lim_{j\to\infty}
\bigl\langle\mathfrak{B}^{\mathrm{sh}}_{\chi};\mathcal{O}(\vec F)\bigr\rangle_{K_j,m_\star}
=-D_\omega[\sigma_\star](\vec F)
\end{equation}
for every smooth $n$-tuple $\vec F$ with pairwise disjoint supports in $B(y,1)$.
\end{proof}

\begin{proof}[Proof of Proposition~\ref{10.1:prop-rotation-defect-two-three}, continued: signed arrays of translates and finite rotations]\label{10.1:prf-signed-array}
We keep the setting fixed earlier in this proof: the reference torus $\mathbb{T}_{\star}=\mathbb{T}_{m_{\star}}$ with $m_{\star}=m_{\mathbb{T}}(g_{-}(g))$, a subsequential limit $\sigma_{\star}$ on it in the sense of clause (ii-$\mathbb{T}$-b) of Theorem~0, taken along a sequence of cutoffs $(K_j)$, the radial cut-off $\chi$ and the shell functional $\mathfrak{B}^{\mathrm{sh}}_{\chi}$, an element $\omega\in\mathfrak{so}(4)$, that is a real antisymmetric $4\times4$ matrix, with $|\omega|^2=\sum_{\mu<\nu}\omega_{\mu\nu}^2$, the field $\xi=\omega(x-y)$, the infinitesimal rotation $\delta_{\omega}f=-(\omega(x-y))\cdot\nabla f$ and the defect $D_{\omega}(\vec F)$ of \eqref{10.1:eq-rotation-defect-two-three-b}. We abbreviate $S_n:=S_n[\sigma_{\star}]$. The letters $\sigma_{\tau}$, $\sigma^{1}_{\tau}$, $\sigma^{2}_{\tau}$ introduced below denote signs $\pm1$; they have nothing to do with the limit $\sigma_{\star}$.

\smallskip\noindent\textit{Translation invariance of the defect.}
Let $\vec F=(F_1,\dots,F_n)$ be a smooth $n$-tuple with pairwise disjoint compact supports in $B(y,1)$, and let $\tau\in\mathbb{R}^4$ be such that the translated tuple $\vec F(\cdot-\tau)$ also has its supports in $B(y,1)$. We show
\begin{equation}\label{10.1:eq-signed-array-1}
D_{\omega}\bigl(\vec F(\cdot-\tau)\bigr)=D_{\omega}(\vec F).
\end{equation}
Since $\xi(x)=\omega(x-y)=\omega(x-\tau-y)+\omega\tau$, the chain rule gives
\begin{equation*}
\xi\cdot\nabla\bigl[F(\cdot-\tau)\bigr]=\bigl[(\omega(\cdot-y)+\omega\tau)\cdot\nabla F\bigr](\cdot-\tau),
\end{equation*}
that is, $\delta_{\omega}[F(\cdot-\tau)]=(\delta_{\omega}F)(\cdot-\tau)-((\omega\tau)\cdot\nabla F)(\cdot-\tau)$. Inserting this in the $i$-th slot and using the linearity of $S_n$ in that slot,
\begin{align*}
D_{\omega}\bigl(\vec F(\cdot-\tau)\bigr)
&=\sum_{i}S_n\bigl(F_1(\cdot-\tau),\dots,(\delta_{\omega}F_i)(\cdot-\tau),\dots,F_n(\cdot-\tau)\bigr)\\
&\quad-\sum_{i}S_n\bigl(F_1(\cdot-\tau),\dots,((\omega\tau)\cdot\nabla F_i)(\cdot-\tau),\dots,
F_n(\cdot-\tau)\bigr).
\end{align*}
By the translation invariance of the torus limit $\sigma_{\star}$ (clause (ii-$\mathbb{T}$-b) of Theorem~0, in the form of the statement on translation invariance of the torus limits) every term may be translated back by $\tau$. The first sum becomes $D_{\omega}(\vec F)$; the second becomes, with $v:=\omega\tau$,
\begin{equation*}
\sum_{i}S_n\bigl(F_1,\dots,v\cdot\nabla F_i,\dots,F_n\bigr)
=\frac{d}{dt}\Big|_{t=0}S_n\bigl(\vec F(\cdot+tv)\bigr).
\end{equation*}
The identity holds because $t^{-1}(F_i(\cdot+tv)-F_i)\to v\cdot\nabla F_i$ in $\|\cdot\|_{\sharp}$ as $t\to0$ for smooth compactly supported $F_i$, and $S_n$ is multilinear and continuous in $\|\cdot\|_{\sharp}$ by clause (iii-a) of Theorem~0. The function $t\mapsto S_n(\vec F(\cdot+tv))$ is constant by the same translation invariance, so its derivative vanishes, and \eqref{10.1:eq-signed-array-1} follows.

\smallskip\noindent\textit{The signed array.}
Let $\vec f\in\mathfrak{B}_n(\mathsf{d};\vartheta)\cap C^{\infty}$, with common ball $\bar B(y,5\mathsf{d})$ (Notation~\ref{10.1:not-clustered-tuples}, display \eqref{10.1:eq-clustered-tuples-a}). Put $r_0:=5\mathsf{d}$. In the regime of that notation $\mathsf{d}\le 1/40$, hence $r_0\le 1/8$. Let
\begin{equation*}
\Gamma:=\bigl\{\tau\in 4r_0\mathbb{Z}^4:\ |\tau|_{\infty}\le\tfrac14\bigr\},\qquad \mathsf{N}:=|\Gamma|.
\end{equation*}
Each coordinate of $\tau\in\Gamma$ ranges over the $2\lfloor u\rfloor+1$ points of $4r_0\mathbb{Z}\cap[-\frac14,\frac14]$, where $u:=1/(16r_0)$. Since $2\lfloor u\rfloor+1\ge\lfloor u\rfloor+1>u$, we get
\begin{equation*}
\mathsf{N}>(16r_0)^{-4}\ge(32r_0)^{-4}.
\end{equation*}
For distinct $\tau,\tau'\in\Gamma$ some coordinate differs by at least $4r_0$, so $|\tau-\tau'|\ge4r_0$. Hence the balls $\bar B(y+\tau,r_0)$, $\tau\in\Gamma$, are pairwise at distance at least $2r_0$. Moreover $|\tau|\le2|\tau|_{\infty}\le\frac12$ and $r_0\le\frac18$, so all of them lie in $\bar B(y,\frac58)\subset B(y,1)$. The translate $f_i(\cdot-\tau)$ is supported in $\bar B(y+\tau,r_0)$.

For $n=2$ and a sign family $\sigma=(\sigma_{\tau})_{\tau\in\Gamma}\in\{\pm1\}^{\Gamma}$ put
\begin{equation*}
F^{\sigma}_i:=\sum_{\tau\in\Gamma}\sigma_{\tau}\,f_i(\cdot-\tau),\qquad i=1,2.
\end{equation*}
For $n=3$ and two sign families $\sigma^{1},\sigma^{2}\in\{\pm1\}^{\Gamma}$ put
\begin{equation*}
F_1:=\sum_{\tau}\sigma^{1}_{\tau}f_1(\cdot-\tau),\quad
F_2:=\sum_{\tau}\sigma^{2}_{\tau}f_2(\cdot-\tau),\quad
F_3:=\sum_{\tau}\sigma^{1}_{\tau}\sigma^{2}_{\tau}f_3(\cdot-\tau).
\end{equation*}
In either case we write $\vec F$ for the resulting signed arrayed tuple and establish three properties.

First, $\vec F$ is smooth, being a finite sum of smooth functions. Since the signs do not vanish, $\operatorname{supp}F_i$ is the union over $\tau\in\Gamma$ of $\operatorname{supp}f_i+\tau$. For $i\neq j$ the sets $\operatorname{supp}f_i+\tau$ and $\operatorname{supp}f_j+\tau$ are disjoint because $\operatorname{supp}f_i\cap\operatorname{supp}f_j=\emptyset$. For $\tau\neq\tau'$ they lie in disjoint balls. So the supports of $F_1,\dots,F_n$ are pairwise disjoint and lie in $B(y,1)$.

Second, on $\bar B(y+\tau,r_0)$ only the term with index $\tau$ is non-zero, and the signs have modulus one, so $\|F_i\|_{\infty}=\|f_i\|_{\infty}$.

Third, we show $\operatorname{Lip}F_i\le\operatorname{Lip}f_i$. A continuous function supported in a closed ball vanishes on the boundary sphere, since a non-zero value there would persist on a neighbourhood meeting the complement. Hence $F_i$ vanishes on every sphere $\partial B(y+\tau,r_0)$ and outside the union of the balls. Let $x,x'\in\mathbb{R}^4$.
\begin{itemize}
\item If $x,x'$ lie in one ball $\bar B(y+\tau,r_0)$, then $|F_i(x)-F_i(x')|=|f_i(x-\tau)-f_i(x'-\tau)|\le\operatorname{Lip}f_i\,|x-x'|$.
\item If neither lies in any ball, then $F_i(x)=F_i(x')=0$.
\item If $x\in\bar B(y+\tau,r_0)$ and $x'$ lies in no ball, let $e$ be the exit point of the segment $[x,x']$ from that ball. Then $F_i(e)=0$, and
\begin{equation*}
|F_i(x)-F_i(x')|=|F_i(x)-F_i(e)|\le\operatorname{Lip}f_i\,|x-e|\le\operatorname{Lip}f_i\,|x-x'|.
\end{equation*}
\item If $x\in\bar B(y+\tau,r_0)$ and $x'\in\bar B(y+\tau',r_0)$ with $\tau\neq\tau'$, let $e$ be the exit point of $[x,x']$ from the first ball and $e'$ its entry point into the second. Then $F_i(e)=F_i(e')=0$, and $e,e'$ lie on the segment in the order $x,e,e',x'$. Hence
\begin{equation*}
|F_i(x)-F_i(x')|\le\operatorname{Lip}f_i\bigl(|x-e|+|e'-x'|\bigr)\le\operatorname{Lip}f_i\,|x-x'|.
\end{equation*}
\end{itemize}
Together with the second property this gives $\|F_i\|_{\sharp}\le\|f_i\|_{\sharp}$.

Next we average over the signs. The defect $D_{\omega}$ is $n$-linear in $(F_1,\dots,F_n)$, because $S_n$ is multilinear and $\delta_{\omega}$ is linear. For $n=2$, expanding gives
\begin{equation*}
D_{\omega}(F^{\sigma}_1,F^{\sigma}_2)=\sum_{\tau,\tau'\in\Gamma}\sigma_{\tau}\sigma_{\tau'}\,
D_{\omega}\bigl(f_1(\cdot-\tau),f_2(\cdot-\tau')\bigr).
\end{equation*}
We have $2^{-\mathsf{N}}\sum_{\sigma}\sigma_{\tau}\sigma_{\tau'}=\delta_{\tau\tau'}$: for $\tau\neq\tau'$ the sum over $\sigma_{\tau}=\pm1$ factors out and vanishes. For $n=3$ the coefficient of $D_{\omega}(f_1(\cdot-\tau_1),f_2(\cdot-\tau_2),f_3(\cdot-\tau_3))$ is $\sigma^{1}_{\tau_1}\sigma^{2}_{\tau_2}\sigma^{1}_{\tau_3}\sigma^{2}_{\tau_3}$. Averaging over the two families separately gives
\begin{equation*}
2^{-2\mathsf{N}}\sum_{\sigma^1,\sigma^2}\sigma^{1}_{\tau_1}\sigma^{1}_{\tau_3}\,\sigma^{2}_{\tau_2}\sigma^{2}_{\tau_3}
=\Bigl(2^{-\mathsf{N}}\sum_{\sigma^1}\sigma^{1}_{\tau_1}\sigma^{1}_{\tau_3}\Bigr)
\Bigl(2^{-\mathsf{N}}\sum_{\sigma^2}\sigma^{2}_{\tau_2}\sigma^{2}_{\tau_3}\Bigr)
=\delta_{\tau_1\tau_3}\delta_{\tau_2\tau_3}.
\end{equation*}
In both cases only the diagonal $\tau_1=\tau_2=\dots=\tau$ survives, and the average of $D_{\omega}(\vec F)$ over the sign families equals $\sum_{\tau\in\Gamma}D_{\omega}(\vec f(\cdot-\tau))$. For $n=2$ this is the averaging identity used for pairs in the volume-transfer argument of clause (iv). Write $\operatorname{avg}_{\mathrm{signs}}$ for the average over the sign families, that is $2^{-\mathsf{N}}\sum_{\sigma}$ for $n=2$ and $2^{-2\mathsf{N}}\sum_{\sigma^1,\sigma^2}$ for $n=3$. Both $\vec f$ and $\vec f(\cdot-\tau)$ have supports in $\bar B(y,\frac58)\subset B(y,1)$, so \eqref{10.1:eq-signed-array-1} gives
\begin{equation}\label{10.1:eq-signed-array-2}
\operatorname{avg}_{\mathrm{signs}}D_{\omega}(\vec F)=\sum_{\tau\in\Gamma}D_{\omega}\bigl(\vec f(\cdot-\tau)\bigr)=\mathsf{N}\,D_{\omega}(\vec f).
\end{equation}

We now bound each signed tuple. Each $\vec F$ is a smooth tuple with pairwise disjoint supports in $B(y,1)$. By the flux form \eqref{10.1:eq-rotation-defect-two-three-a}, along the sequence $(K_j)$ defining $\sigma_{\star}$,
\begin{equation*}
D_{\omega}(\vec F)=-\lim_{j}\langle\mathfrak{B}^{\mathrm{sh}}_{\chi};\mathcal{O}(\vec F)\rangle_{K_j,m_{\star}}.
\end{equation*}
The unit-sphere shell bound, Lemma~\ref{10.1:lem-unit-shell-bound}, bounds the $\limsup$ of the modulus of this expectation and gives
\begin{equation*}
|D_{\omega}(\vec F)|\le C_{\mathrm{sh},n}(g)\,|\omega|\prod_i\|F_i\|_{\sharp}\le C_{\mathrm{sh},n}(g)\,|\omega|\prod_i\|f_i\|_{\sharp}.
\end{equation*}
Combining this with \eqref{10.1:eq-signed-array-2} and $\mathsf{N}^{-1}\le(32r_0)^4=160^4\mathsf{d}^4$ yields \eqref{10.1:eq-signed-array-a} below.

For the last inequality in that display we use $\operatorname{Lip}f_i\le\|f_i\|_{\infty}/(\vartheta\mathsf{d})$ and $40M/(\vartheta\mathsf{d})\ge1$, which holds since $M\ge1$, $\vartheta\le1$ and $\mathsf{d}\le1/40$. These give
\begin{equation*}
\|f_i\|_{\sharp}=\max\{\|f_i\|_{\infty},40M\operatorname{Lip}f_i\}\le\frac{40M}{\vartheta\mathsf{d}}\,\|f_i\|_{\infty},
\end{equation*}
and hence $\prod_i\|f_i\|_{\sharp}\le(40M/\vartheta)^n\mathsf{d}^{-n}\,\Pi$. Therefore
\begin{equation}\label{10.1:eq-signed-array-a}
\begin{split}
|D_{\omega}(\vec f)|&\le(32r_0)^4\,C_{\mathrm{sh},n}(g)\,|\omega|\prod_i\|f_i\|_{\sharp}
=160^4\,\mathsf{d}^4\,C_{\mathrm{sh},n}(g)\,|\omega|\prod_i\|f_i\|_{\sharp}\\
&\le160^4\Bigl(\frac{40M}{\vartheta}\Bigr)^{n}C_{\mathrm{sh},n}(g)\,|\omega|\,\mathsf{d}^{4-n}\,\Pi .
\end{split}
\end{equation}
The same bound holds for tuples $\vec f\in\mathfrak{B}_n(\mathsf{d};\vartheta)\cap C^2$ by mollification. What is passed to the limit is the first line of \eqref{10.1:eq-signed-array-a}, in which the tuple enters only through $D_{\omega}$ and the norms $\|f_i\|_{\sharp}$. Smooth tuples obtained by convolution converge to $\vec f$, and their images under $\delta_{\omega}$ converge to those of $\vec f$, in $\|\cdot\|_{\sharp}$. Both sides pass to the limit by the continuity of $S_n$ in $\|\cdot\|_{\sharp}$ furnished by clause (iii-a) of Theorem~0. The second line then follows for the $C^2$ tuple $\vec f$ itself from $\operatorname{Lip}f_i\le\|f_i\|_{\infty}/(\vartheta\mathsf{d})$, exactly as above.

\smallskip\noindent\textit{Finite rotations on $\mathbb{T}_{\star}$.}
For $t\in\mathbb{R}$ let $e^{t\omega}_{y}$ denote the rotation $x\mapsto y+e^{t\omega}(x-y)$, acting on test functions by $R\cdot f=f\circ R^{-1}$. Let $\vec f\in\mathfrak{B}_n(\mathsf{d};\vartheta)\cap C^2$. By the lemma on differentiability of the torus limits along rotations, the map $t\mapsto S_n(e^{t\omega}_{y}\cdot\vec f)$ is $C^1$ with derivative $D_{\omega}(e^{t\omega}_{y}\cdot\vec f)$.

Rotation about $y$ preserves $\bar B(y,5\mathsf{d})$, Euclidean distances and sup norms. Hence $e^{t\omega}_{y}\cdot\vec f\in\mathfrak{B}_n(\mathsf{d};\vartheta)\cap C^2$ with the same $\Pi$, and \eqref{10.1:eq-signed-array-a} applies to it for every $t$. Every $R\in SO(4)$ is of the form $e^{\omega}$ with $|\omega|\le\pi\sqrt2$. For $R=e^{\omega}$ acting about $y$, integrating the derivative over $t\in[0,1]$ gives
\begin{equation*}
|S_n(R\cdot\vec f)-S_n(\vec f)|\le\int_0^1|D_{\omega}(e^{t\omega}_{y}\cdot\vec f)|\,dt
\le\pi\sqrt2\cdot160^4\Bigl(\frac{40M}{\vartheta}\Bigr)^{n}C_{\mathrm{sh},n}(g)\,\mathsf{d}^{4-n}\,\Pi .
\end{equation*}

Next let $R\in O(4)\setminus SO(4)$. Then $O(4)=SO(4)\cup SO(4)\mathsf{r}$, where $\mathsf{r}\in W_4$ is a reflection, so $R=R'\mathsf{r}$ with $R'\in SO(4)$. The limit $\sigma_{\star}$ is $W_4$-invariant by clause (ii-$\mathbb{T}$-b) of Theorem~0, and by translation invariance the centre may be taken at $y$; thus $S_n(\mathsf{r}\cdot\vec f)=S_n(\vec f)$. Moreover $\mathsf{r}\cdot\vec f\in\mathfrak{B}_n(\mathsf{d};\vartheta)\cap C^2$ with the same sup norms. Applying the bound just proved to $R'$ and the tuple $\mathsf{r}\cdot\vec f$ gives the same estimate for $|S_n(R\cdot\vec f)-S_n(\vec f)|$.

Finally, this estimate involves only the values of $S_n$ at the tuples $R\cdot\vec f$ and $\vec f$. It therefore extends from $C^2$ tuples to $C^1$ tuples in $\mathfrak{B}_n(\mathsf{d};\vartheta)$ by the density of $C^2$ in the norm $\|\cdot\|_{\sharp}$ and the continuity supplied by clause (iii-a). What is passed to the limit is the estimate with $160^4\mathsf{d}^4\prod_i\|f_i\|_{\sharp}$ in place of $(40M/\vartheta)^n\mathsf{d}^{4-n}\Pi$, which the integration above also gives by the first line of \eqref{10.1:eq-signed-array-a}, since rotations and reflections about $y$ preserve sup norms and Lipschitz constants; the $C^1$ tuple itself then satisfies $\operatorname{Lip}f_i\le\|f_i\|_{\infty}/(\vartheta\mathsf{d})$, and the conversion is made as in \eqref{10.1:eq-signed-array-a}. For every $\vec f\in\mathfrak{B}_n(\mathsf{d};\vartheta)\cap C^1$ and every $R\in O(4)$,
\begin{equation}\label{10.1:eq-signed-array-b}
|S_n[\sigma_{\star}](R\cdot\vec f)-S_n[\sigma_{\star}](\vec f)|
\le\pi\sqrt2\cdot160^4\Bigl(\frac{40M}{\vartheta}\Bigr)^{n}C_{\mathrm{sh},n}(g)\,\mathsf{d}^{4-n}\,\Pi
=:C_{\mathrm{rot},n}\,\mathsf{d}^{4-n}\,\Pi .
\end{equation}
\end{proof}

\begin{proof}[Proof of Proposition~\ref{10.1:prop-rotation-defect-two-three}, continued: volume transfer and the running-coupling form]
\label{10.1:prf-transfer-coupling-form}
Throughout this part of the proof $n\in\{2,3\}$ and $p:=4-n$. We use the notation of Notation~\ref{10.1:not-clustered-tuples} and of \eqref{10.1:eq-clustered-tuples-a}: a tuple $\vec f\in\mathfrak{B}_n(\mathsf{d};\vartheta)$ has its supports in the closed ball $\bar B(y,5\mathsf{d})$, and $\Pi:=\prod_i\|f_i\|_\infty$. We write
\begin{equation*}
m_\star:=m_{\mathbb{T}}(g_-(g)),
\end{equation*}
and $\mathbb{T}_\star$ for the torus of side $2L^{m_\star}$, the smallest torus of the family. The letter $\sigma_\star$ denotes a torus limit on $\mathbb{T}_\star$ in the sense of clause (ii-$\mathbb{T}$-b) of Theorem~0. We write $D_\omega$ for the rotation defect of $\sigma_\star$ displayed in \eqref{10.1:eq-rotation-defect-two-three-b}. The constant $C_{\mathrm{rot},n}$ is the one of \eqref{10.1:eq-signed-array-b}.

\emph{The volume-transfer bound on a signed array.}
We use the volume-transfer theorem of clause (iv) of Theorem~0, one of the named antecedents of Proposition~\ref{10.1:prop-rotation-defect-two-three}. It bounds the difference between a connected Schwinger function on a torus of the family and the same function on $\mathbb{T}_\star$. Let $K$ be a cutoff of the sequences of clause (ii) at reference $g$, and let $m\ge m_\star$. Let $\vec F$ be an $n$-tuple of test functions with pairwise disjoint supports in $B(y,1)$; no condition is placed on the signs of the $F_i$ or on the positions of their supports in this ball. Then
\begin{equation}\label{10.1:eq-transfer-coupling-form-1}
\bigl|S^T_{n;K,m}(\vec F)-S^T_{n;K,m_\star}(\vec F)\bigr|
\le C_{\mathrm{VT},n}(g)\prod_{i=1}^n\|F_i\|_{\sharp}.
\end{equation}
Write $\mathsf{E}:=S^T_{n;K,m}-S^T_{n;K,m_\star}$. This functional is $n$-linear in the test functions.

Let $\vec f$ be an $n$-tuple with pairwise disjoint supports in $\bar B(y,5\mathsf{d})$. We form the signed array of translates used for \eqref{10.1:eq-signed-array-a}. That is, $r_0=5\mathsf{d}$; $\Gamma$ is the set of $\tau\in 4r_0\mathbb{Z}^4$ with $|\tau|_\infty\le\tfrac14$; $\mathsf{N}=|\Gamma|\ge(32r_0)^{-4}$; and $\vec F^{\sigma}$ is the signed arrayed tuple built from the fair sign families $\sigma_\tau$ (for $n=2$) or $\sigma^1_\tau,\sigma^2_\tau$ (for $n=3$).

Each $\vec F^{\sigma}$ has pairwise disjoint supports in $B(y,1)$, and $\|F^{\sigma}_i\|_\sharp\le\|f_i\|_\sharp$. Hence \eqref{10.1:eq-transfer-coupling-form-1} applies to every $\vec F^{\sigma}$.

We now average over sign families. For $n=2$ the average of $\sigma_\tau\sigma_{\tau'}$ is $\delta_{\tau\tau'}$. For $n=3$ the average of $\sigma^1_{\tau_1}\sigma^2_{\tau_2}\sigma^1_{\tau_3}\sigma^2_{\tau_3}$ is $\delta_{\tau_1\tau_3}\delta_{\tau_2\tau_3}$. By $n$-linearity of $\mathsf{E}$, the average of $\mathsf{E}(\vec F^{\sigma})$ is therefore $\sum_{\tau\in\Gamma}\mathsf{E}(\vec f(\cdot-\tau))$.

The identity that turns the corresponding sum for $D_\omega$ into $\mathsf{N}\,D_\omega(\vec f)$ in the derivation of \eqref{10.1:eq-signed-array-a} turns this sum into $\mathsf{N}\,\mathsf{E}(\vec f)$. The average of numbers each bounded by the right side of \eqref{10.1:eq-transfer-coupling-form-1} is bounded by the same quantity. Using $\mathsf{N}^{-1}\le(32r_0)^4=160^4\mathsf{d}^4$, we obtain
\begin{equation}\label{10.1:eq-transfer-coupling-form-2}
|\mathsf{E}(\vec f)|\le 160^4\,\mathsf{d}^4\,C_{\mathrm{VT},n}(g)\prod_{i=1}^n\|f_i\|_\sharp .
\end{equation}
No Lipschitz condition on $\vec f$ has been used so far; only the supports and the norms $\|\cdot\|_\sharp$ enter.

Now let $\vec f\in\mathfrak{B}_n(\mathsf{d};\vartheta)$. Then $\operatorname{Lip}f_i\le\|f_i\|_\infty/(\vartheta\mathsf{d})$. Since $M\ge1$, $\vartheta\le1$ and $\mathsf{d}\le1/40$, we have $40M/(\vartheta\mathsf{d})\ge1$. Therefore
\begin{equation}\label{10.1:eq-transfer-coupling-form-3}
\|f_i\|_\sharp=\max\{\|f_i\|_\infty,\,40M\operatorname{Lip}f_i\}\le\frac{40M}{\vartheta\mathsf{d}}\,\|f_i\|_\infty .
\end{equation}
Inserting \eqref{10.1:eq-transfer-coupling-form-3} into \eqref{10.1:eq-transfer-coupling-form-2} gives the corollary of the volume-transfer theorem used below:
\begin{equation}\label{10.1:eq-transfer-coupling-form-a}
\bigl|S^T_{n;K,m}(\vec f)-S^T_{n;K,m_\star}(\vec f)\bigr|
\le 160^4\Bigl(\frac{40M}{\vartheta}\Bigr)^{n}C_{\mathrm{VT},n}(g)\,\mathsf{d}^{4-n}\,\Pi
=:C_{T,n}(g)\,\mathsf{d}^{4-n}\,\Pi .
\end{equation}
The same bound holds with $R\cdot\vec f$ in place of $\vec f$, for every $R\in O(4)$. Indeed $R$ acts about $y$, so it maps $\bar B(y,5\mathsf{d})$ onto itself and preserves disjointness of supports, sup norms and Lipschitz constants. Hence $R\cdot\vec f\in\mathfrak{B}_n(\mathsf{d};\vartheta)$, with the same $\Pi$.

\emph{Item (a).}
Let $\sigma$ be a subsequential limit of clause (ii) at reference $g$ along a sequence $(g_0^{(j)},K_j,m_j)$ with $m_j\ge m_\star$. This covers both the $\mathbb{R}^4$ limits of (ii-$\mathbb{R}^4$) and the torus limits of (ii-$\mathbb{T}$-b) on every torus of the family. By clause (ii-$\mathbb{T}$-b) we may pass to a subsequence along which $S^T_{n;K_j,m_\star}$ converges to a torus limit $\sigma_\star$ on $\mathbb{T}_\star$.

Let $\vec f\in\mathfrak{B}_n(\mathsf{d};\vartheta)\cap C^1$ and $R\in O(4)$. Apply \eqref{10.1:eq-transfer-coupling-form-a} at $(K_j,m_j)$ to $\vec f$ and to $R\cdot\vec f$, and let $j\to\infty$ along the subsequence:
\begin{equation*}
\bigl|S_n[\sigma](\vec f)-S_n[\sigma_\star](\vec f)\bigr|\le C_{T,n}\,\mathsf{d}^{4-n}\Pi,
\qquad
\bigl|S_n[\sigma](R\cdot\vec f)-S_n[\sigma_\star](R\cdot\vec f)\bigr|\le C_{T,n}\,\mathsf{d}^{4-n}\Pi .
\end{equation*}
By \eqref{10.1:eq-signed-array-b},
\begin{equation*}
|S_n[\sigma_\star](R\cdot\vec f)-S_n[\sigma_\star](\vec f)|\le C_{\mathrm{rot},n}\mathsf{d}^{4-n}\Pi .
\end{equation*}
The triangle inequality through $S_n[\sigma_\star](R\cdot\vec f)$ and $S_n[\sigma_\star](\vec f)$ gives
\begin{equation*}
\bigl|S_n[\sigma](R\cdot\vec f)-S_n[\sigma](\vec f)\bigr|
\le\bigl(C_{\mathrm{rot},n}+2C_{T,n}\bigr)(g)\,\mathsf{d}^{4-n}\,\Pi
=C_n(g,\vartheta)\,\mathsf{d}^{4-n}\,\Pi .
\end{equation*}
With $C_{\mathrm{rot},n}=\pi\sqrt2\cdot160^4(40M/\vartheta)^nC_{\mathrm{sh},n}$ this constant equals
\begin{equation*}
C_n(g,\vartheta)=160^4\Bigl(\frac{40M}{\vartheta}\Bigr)^{n}\bigl[\pi\sqrt2\,C_{\mathrm{sh},n}(g)+2C_{\mathrm{VT},n}(g)\bigr].
\end{equation*}
This is item (a).

\emph{Item (b).}
We argue by contradiction. Suppose that for some $\vec f\in\mathfrak{B}_n(\mathsf{d};\vartheta)\cap C^1$ and some $R\in O(4)$ the $\limsup$ in (b) exceeds $C_n(g,\vartheta)\mathsf{d}^{4-n}\Pi$. Then there are $\epsilon>0$, cutoffs $K_j\to\infty$ of the sequences of clause (ii) at reference $g$, and $m_j\ge m_\star$ with
\begin{equation*}
\bigl|S^T_{n;K_j,m_j}(R\cdot\vec f)-S^T_{n;K_j,m_j}(\vec f)\bigr|\ge C_n\,\mathsf{d}^{4-n}\Pi+\epsilon
\quad\text{for all } j.
\end{equation*}
Apply \eqref{10.1:eq-transfer-coupling-form-a} twice, once to $\vec f$ and once to $R\cdot\vec f$, and subtract $2C_{T,n}\mathsf{d}^{4-n}\Pi$. Since $C_n-2C_{T,n}=C_{\mathrm{rot},n}$, this gives
\begin{equation*}
\bigl|S^T_{n;K_j,m_\star}(R\cdot\vec f)-S^T_{n;K_j,m_\star}(\vec f)\bigr|\ge C_{\mathrm{rot},n}\,\mathsf{d}^{4-n}\Pi+\epsilon .
\end{equation*}
By clause (ii-$\mathbb{T}$-b) this violating sequence has a subsequence along which $S^T_{n;K_j,m_\star}$ converges to a torus limit $\sigma_\star$ on $\mathbb{T}_\star$. In the limit,
\begin{equation*}
|S_n[\sigma_\star](R\cdot\vec f)-S_n[\sigma_\star](\vec f)|\ge C_{\mathrm{rot},n}\mathsf{d}^{4-n}\Pi+\epsilon .
\end{equation*}
This contradicts \eqref{10.1:eq-signed-array-b}, and item (b) follows.

\emph{Item (c).}
Let $\vec f\in\mathfrak{B}_n(\mathsf{d};\vartheta)$ with $f_i\in C^2$, let $\sigma$, $(g_0^{(j)},K_j,m_j)$ and $\sigma_\star$ be as in item (a), and put
\begin{equation*}
\mathsf{E}_j:=S^T_{n;K_j,m_j}-S^T_{n;K_j,m_\star}.
\end{equation*}
For each $i$, the tuple $(f_1,\dots,\delta_\omega f_i,\dots,f_n)$ has pairwise disjoint supports in $\bar B(y,5\mathsf{d})$, because $\operatorname{supp}\delta_\omega f_i\subset\operatorname{supp}f_i$. Summing over $i$ and letting $j\to\infty$ along the subsequence gives
\begin{equation*}
\sum_{i=1}^n S_n[\sigma](f_1,\dots,\delta_\omega f_i,\dots,f_n)
=D_\omega(\vec f)+\lim_{j\to\infty}\sum_{i=1}^n\mathsf{E}_j(f_1,\dots,\delta_\omega f_i,\dots,f_n).
\end{equation*}

The first term on the right is bounded by \eqref{10.1:eq-signed-array-a}:
\begin{equation*}
|D_\omega(\vec f)|\le160^4(40M/\vartheta)^nC_{\mathrm{sh},n}|\omega|\,\mathsf{d}^{4-n}\Pi .
\end{equation*}

For the second term we apply \eqref{10.1:eq-transfer-coupling-form-2} to $\mathsf{E}_j$ at each of the tuples $(f_1,\dots,\delta_\omega f_i,\dots,f_n)$; this is legitimate because \eqref{10.1:eq-transfer-coupling-form-2} uses only the supports and the $\sharp$-norms. For the slots $l\ne i$ we use \eqref{10.1:eq-transfer-coupling-form-3}. For the slot $i$ we use the lemma on the infinitesimal rotation of a member of the clustered class, one of the named antecedents of Proposition~\ref{10.1:prop-rotation-defect-two-three}, which reads
\begin{equation*}
\|\delta_\omega f_i\|_\sharp\le\frac{240M|\omega|}{\vartheta^2\mathsf{d}}\|f_i\|_\infty
=\frac{40M}{\vartheta\mathsf{d}}\cdot\frac{6|\omega|}{\vartheta}\,\|f_i\|_\infty .
\end{equation*}
Each of the $n$ terms is therefore at most
\begin{equation*}
160^4\,\mathsf{d}^4\,C_{\mathrm{VT},n}\Bigl(\frac{40M}{\vartheta\mathsf{d}}\Bigr)^n\frac{6|\omega|}{\vartheta}\,\Pi
=160^4\Bigl(\frac{40M}{\vartheta}\Bigr)^n\frac{6}{\vartheta}\,C_{\mathrm{VT},n}\,|\omega|\,\mathsf{d}^{4-n}\,\Pi ,
\end{equation*}
uniformly in $j$. Summing over $i$ gives the factor $6n/\vartheta$. Adding the two contributions,
\begin{equation*}
\Bigl|\sum_{i=1}^nS_n[\sigma](f_1,\dots,\delta_\omega f_i,\dots,f_n)\Bigr|
\le C^D_n(g,\vartheta)\,|\omega|\,\mathsf{d}^{4-n}\,\Pi ,
\end{equation*}
with
\begin{equation*}
C^D_n=160^4\Bigl(\frac{40M}{\vartheta}\Bigr)^n\Bigl[C_{\mathrm{sh},n}+\frac{6n}{\vartheta}C_{\mathrm{VT},n}\Bigr].
\end{equation*}
This is item (c).

\emph{Item (e).}
Let $n=2$, and let $(f_1,f_2)$ be a $C^1$ pair in clause (iv)'s class $\mathfrak{T}(\mathsf{d};\vartheta)$, under the conditions displayed in clause (iv). On every torus of the family and on $\mathbb{R}^4$, the lower member of clause (iv), one of the named antecedents of Proposition~\ref{10.1:prop-rotation-defect-two-three}, bounds $|S_2[\sigma](f_1,f_2)|$ from below, and we divide by this lower bound.
We write $\vec f$ for the pair. Then
\begin{equation*}
\Bigl|\frac{S_2[\sigma](R\cdot\vec f)}{S_2[\sigma](\vec f)}-1\Bigr|
=\frac{|S_2[\sigma](R\cdot\vec f)-S_2[\sigma](\vec f)|}{|S_2[\sigma](\vec f)|}.
\end{equation*}
Bounding the numerator by item (a) with $n=2$ and the denominator by the lower member of clause (iv) gives
\begin{equation*}
\Bigl|\frac{S_2[\sigma](R\cdot\vec f)}{S_2[\sigma](\vec f)}-1\Bigr|
\le\frac92\,(b_0c_{\mathcal{K}})^{-1}\,C_2(g,\vartheta)\,\mathsf{d}^2\bigl(g^{-2}+B^\sharp(s(\mathsf{d})+1)\bigr)^2 .
\end{equation*}
Since $s(\mathsf{d})\le\log_L(1/\mathsf{d})+1$, the last factor is $O(\log^2(1/\mathsf{d}))$.

\emph{Item (d).}
Put $s:=s(\mathsf{d})=\lceil\log_L(1/\mathsf{d})\rceil$, $X:=g^{-2}$ and $x:=g(\mathsf{d})^{-2}$. Since $s-1<\log_L(1/\mathsf{d})$ and $p\ge1$,
\begin{equation*}
\mathsf{d}^p\le L^{-p(s-1)}=L^p\,L^{-ps}.
\end{equation*}

We first bound $L^{-ps}$ using clause (0) alone. By clause (0), $x\le X+B^\sharp(s+1)$, that is, $s\ge(x-X)/B^\sharp-1$. Hence
\begin{equation}\label{10.1:eq-transfer-coupling-form-4}
L^{-ps}\le L^{p(X+B^\sharp)/B^\sharp}\,e^{-(p\ln L/B^\sharp)\,x}.
\end{equation}
Combined with the previous display, \eqref{10.1:eq-transfer-coupling-form-4} gives the weaker form stated with (d). The right sides of (a) and (b) are $C_n\mathsf{d}^{p}\Pi$ and that of (c) is $C^D_n|\omega|\mathsf{d}^p\Pi$, so they are at most
\begin{equation*}
A'_n\exp[-(4-n)(\ln L/B^\sharp)g(\mathsf{d})^{-2}]\,\Pi,
\qquad
A'_n:=L^{p}L^{p(X+B^\sharp)/B^\sharp}\max\{C_n,C^D_n\},
\end{equation*}
with the factor $|\omega|$ carried along in (c).

We next use the one-loop law (0)(d). It states that
\begin{equation*}
x_{K-s}-x_K=s\,c_0+o(s),\qquad c_0=\frac{11N^2}{12\pi^2}\ln L,\qquad x_K\le X+B^\sharp .
\end{equation*}
With $x=x_{K-s}$, or its limit value along the sequence, this gives
\begin{equation*}
s\,c_0\ge x-X-B^\sharp-\Delta(s),
\end{equation*}
where $\Delta(s)$ is the modulus of the remainder and $\Delta(s)/s\to0$.

Fix $\eta\in\,]0,1[$. Choose $s_\eta$ such that $\Delta(s)\le\frac{\eta}{1-\eta}\,s\,c_0$ for $s\ge s_\eta$. For $s\ge s_\eta$ we then have $s\,c_0/(1-\eta)\ge x-X-B^\sharp$, hence
\begin{equation*}
L^{-ps}=e^{-ps\ln L}\le e^{(1-\eta)p(\ln L/c_0)(X+B^\sharp)}\,e^{-(1-\eta)p(\ln L/c_0)\,x}.
\end{equation*}
For the finitely many $s<s_\eta$ we argue directly. There $\mathsf{d}^p\le1$, and $x\le X+B^\sharp(s_\eta+1)$, so
\begin{equation*}
1\le e^{(1-\eta)p(\ln L/c_0)(X+B^\sharp(s_\eta+1))}\,e^{-(1-\eta)p(\ln L/c_0)\,x}.
\end{equation*}
We put
\begin{equation*}
A_\eta:=\max\Bigl\{L^{p}e^{(1-\eta)p(\ln L/c_0)(X+B^\sharp)},\ 
e^{(1-\eta)p(\ln L/c_0)(X+B^\sharp(s_\eta+1))}\Bigr\},
\end{equation*}
and, using $\ln L/c_0=12\pi^2/(11N^2)$, we obtain
\begin{equation}\label{10.1:eq-transfer-coupling-form-5}
\mathsf{d}^p\le A_\eta\exp\Bigl[-(1-\eta)\,p\,\frac{12\pi^2}{11N^2}\,g(\mathsf{d})^{-2}\Bigr].
\end{equation}
Multiplying \eqref{10.1:eq-transfer-coupling-form-5} by $C_n\Pi$, respectively $C^D_n|\omega|\Pi$, gives the first form of (d) with $A_{n,\eta}:=A_\eta\max\{C_n,C^D_n\}$.

For the factorial form we need an elementary inequality. For $c>0$, $t\ge0$ and an integer $k\ge0$,
\begin{equation}\label{10.1:eq-transfer-coupling-form-6}
t^ke^{-ct}\le\Bigl(\frac{k}{ec}\Bigr)^k\le k!\,c^{-k},
\end{equation}
with $0^0:=1$. For $k=0$ the left side is $e^{-ct}\le1$ and the other two members equal $1$. For $k\ge1$, the function $t\mapsto t^ke^{-ct}$ on $[0,\infty[$ has derivative $t^{k-1}e^{-ct}(k-ct)$. It therefore attains its maximum at $t=k/c$, with value $(k/c)^ke^{-k}$. The second inequality is $k^k\le k!\,e^k$, which holds because $e^k=\sum_{j\ge0}k^j/j!\ge k^k/k!$.

Now fix one $\eta\in\,]0,1[$ and put
\begin{equation*}
c_n:=(1-\eta)(4-n)\frac{12\pi^2}{11N^2},\qquad A_n:=A_{n,\eta}.
\end{equation*}
Apply \eqref{10.1:eq-transfer-coupling-form-6} with $t=g(\mathsf{d})^{-2}$ and $c=c_n$. Then
\begin{equation*}
e^{-c_nt}=t^{-k}\,t^ke^{-c_nt}\le k!\,c_n^{-k}\,g(\mathsf{d})^{2k},
\end{equation*}
so the right sides of (a)–(c) are at most $A_n\,k!\,c_n^{-k}g(\mathsf{d})^{2k}\Pi$ for every integer $k\ge0$, with one pair $A_n,c_n$. At $N=2$ one has $c_n=(1-\eta)(4-n)\,3\pi^2/11$. This completes item (d).

\emph{Item (f).}
Let $0<\mathsf{d}_0\le1/40$, let $\vec f^{(0)}\in\mathfrak{B}_n(\mathsf{d}_0;\vartheta)\cap C^1$ be a fixed shape, and let $\vec f^{(s)}$ be its dilate to depth $s$, that is, $f^{(s)}_i(x):=f^{(0)}_i(y+L^{s}(x-y))$. The supports of $\vec f^{(s)}$ are pairwise disjoint and lie in $\bar B(y,5L^{-s}\mathsf{d}_0)$, the sup norms are unchanged, and
\begin{equation*}
\operatorname{Lip}f^{(s)}_i=L^{s}\operatorname{Lip}f^{(0)}_i
\le\frac{\|f^{(s)}_i\|_\infty}{\vartheta L^{-s}\mathsf{d}_0}.
\end{equation*}
Hence $\vec f^{(s)}\in\mathfrak{B}_n(\mathsf{d};\vartheta)\cap C^1$ with $\mathsf{d}:=L^{-s}\mathsf{d}_0\le1/40$, and, since $s$ is an integer and $\log_L(1/\mathsf{d})=s+\log_L(1/\mathsf{d}_0)$, the depth of the dilated tuple is $s(\mathsf{d})=s+s(\mathsf{d}_0)$, a fixed shift of $s$. Dilation does not change $\Pi$. Let $R_0,R\in O(4)$. Apply the factorial form of (d) to the pairs $(R\cdot\vec f^{(s)},\vec f^{(s)})$ and $(R_0\cdot\vec f^{(s)},\vec f^{(s)})$ and use the triangle inequality. For every integer $k\ge0$,
\begin{equation*}
\bigl|S_n[\sigma](R\cdot\vec f^{(s)})-S_n[\sigma](R_0\cdot\vec f^{(s)})\bigr|\le2A_n\,k!\,c_n^{-k}\,g(\mathsf{d})^{2k}\,\Pi .
\end{equation*}
The two functions of $s$ therefore differ by a quantity that is smaller than every power of the running coupling at depth $s(\mathsf{d})=s+s(\mathsf{d}_0)$. By the corollary on asymptotic orders, one of the named antecedents of Proposition~\ref{10.1:prop-rotation-defect-two-three}, an asymptotic expansion of $S_n[\sigma](R_0\cdot\vec f^{(s)})$ in powers of $g(L^{-s})^2$ to some order is then also an asymptotic expansion, with the same coefficients, of $S_n[\sigma](R\cdot\vec f^{(s)})$. This is item (f).

The bounds obtained in (a)–(c) are a positive power of $\mathsf{d}$, hence, by (d), smaller than every power of $g(\mathsf{d})$. They are not zero. For torus limits of the family a zero bound at fixed $\mathsf{d}$ is excluded by the proposition on torus limits, by which exactness at fixed $\mathsf{d}$ fails for every torus limit of the family.
\end{proof}

\begin{remark}[Why the cut-off sphere has unit radius]\label{10.1:rem-unit-radius-sphere}
In the proof of Proposition~\ref{10.1:prop-rotation-defect-two-three}, for a tuple
$\vec f=(f_1,\dots,f_n)\in\mathfrak{B}_n(\mathsf{d};\vartheta)$, $n\in\{2,3\}$, at scale
$\mathsf{d}$, the radial cut-off $\chi$
of the shell functional $\mathfrak{B}^{\mathrm{sh}}_{\chi}$ of \eqref{10.1:eq-clustered-tuples-b}
has unit radius, fixed independently of the scale $\mathsf{d}$ of the tuple. The factor
$\mathsf{d}^{4-n}$ comes from the signed array of translates \eqref{10.1:eq-signed-array-a}, not from
a dimension count. Suppose the cut-off is placed at a radius comparable to $\mathsf{d}$, with all
supports inside it, and the shell term is bounded by the same ultraviolet estimates at scale
$\mathsf{d}$. The bound obtained then contains no positive power of $\mathsf{d}$.

Two further features of the argument are recorded here. The argument is run
on the single reference torus $\mathbb{T}_{\star}$ and only afterwards carried to every torus of
the family and to $\mathbb{R}^4$. The stress-tensor insertion bound enters only in its far form.
\end{remark}

\begin{proof}
We justify each of these assertions in turn.

\emph{The estimates are not scale covariant.} The constant $C_{\mathrm{UV}}(\mathsf{d})$ of the
torus bound of the correlation theorem at separation $\mathsf{d}$ satisfies
\begin{equation}\label{10.1:eq-unit-radius-sphere-1}
  C_{\mathrm{UV}}(\mathsf{d})\le C\,\mathsf{d}^{-4}\bigl(1+\log^{+}\mathsf{d}^{-1}\bigr).
\end{equation}
Here the factor $\mathsf{d}^{-4}$ is the number of cubes of a last-level mesh of fixed spacing.
The insertion bound behaves in the same way. Suppose these estimates are applied at scale
$\mathsf{d}$ to a shell term at radius comparable to $\mathsf{d}$, with sources of size
$\mathsf{d}$. They then give a bound of order
\begin{equation*}
  \mathsf{d}^{-4}\bigl(1+\log^{+}\mathsf{d}^{-1}\bigr)\,\mathsf{d}^{-n}\prod_i\|f_i\|_{\infty},
\end{equation*}
which contains no positive power of $\mathsf{d}$.

\emph{The limiting shell term does not depend on the radius.} By the identity
\eqref{10.1:eq-rotation-defect-two-three-a}, the limiting shell term equals
$-D_{\omega}(\vec f)$, with $D_{\omega}$ as in \eqref{10.1:eq-rotation-defect-two-three-b}. This
holds at every admissible radius of the cut-off; it is the radius-independence of the limiting
shell term. Consequently, choosing the radius comparable to $\mathsf{d}$ cannot improve the bound.

The smallness is instead the ratio of the size of the tuple to the radius of the shell. It is
extracted by packing a number of signed translates of the tuple of order
$(\text{shell radius}/\mathsf{d})^4$ inside the ball of the cut-off. This requires the shell
radius to be much larger than $\mathsf{d}$. At unit radius this packing is carried out in the
proof of the bound \eqref{10.1:eq-signed-array-a}, by signed translates of the tuple placed
inside the ball of the cut-off.

\emph{The coarsest frame is the best.} Write $s(\mathsf{d})=\lceil\log_L 1/\mathsf{d}\rceil$, and
let $g(\mathsf{d})=g_{K-s(\mathsf{d})}$ be the running coupling at depth $s(\mathsf{d})$. Transport
the unit-radius bound to a cut-off at radius comparable to $\mathsf{d}$ by the exact change of
scale. The result is a constant $C(g(\mathsf{d}))$ with no power of $\mathsf{d}$. Among the scales
at which the cut-off may be placed, the coarsest, unit radius, is therefore the best.

\emph{The argument runs on one torus.} The constant $C_{\mathrm{sh},n}(g)$ of
Lemma~\ref{10.1:lem-unit-shell-bound}, displayed in \eqref{10.1:eq-unit-shell-bound-a}, contains the
insertion constants $C^{\otimes}_{n,\mathbf{3}}(1)$ and $C^{\otimes}_{n,\mathbf{6}}(1)$ at distance
one. These contain $\mathsf{Q}(m_\star)$, where $m_\star:=m_{\mathbb{T}}(g_-(g))$ is the exponent
of the reference torus $\mathbb{T}_{\star}=\mathbb{T}_{m_\star}$. Hence the constant depends on the
reference torus, and one
cannot pass directly to the limit along (ii-$\mathbb{R}^4$). The bound is proved for the limit
$\sigma_{\star}$ on $\mathbb{T}_{\star}$. It is then carried to every torus of the family in
(ii-$\mathbb{T}$-b) and to $\mathbb{R}^4$ by the volume-transfer theorem.

\emph{The far form of the insertion bound.} The argument uses the stress-tensor-insertion
correlation theorem (Theorem~\ref{10.4:thm-insertion-correlation}) in the following form: its
constant $C^{\otimes}_{n,r}(\mathsf{D})$ depends on the supports of the tensor weight $\mathsf{k}$
and of the $F_i$ only through the distance
\begin{equation*}
  \mathsf{D}=\operatorname{dist}\Bigl(\operatorname{supp}\mathsf{k},\
  \bigcup_i\operatorname{supp}F_i\Bigr),
\end{equation*}
and it imposes no condition on the mutual positions of the supports of the $F_i$. This is the
form that makes the constant
$C_{\mathrm{sh},n}(g)$ of Lemma~\ref{10.1:lem-unit-shell-bound} independent of $\mathsf{d}$ and of
the mutual positions of the supports of the $F_i$.

The signed array needs exactly this. Its arrayed tuples $\vec F^{\sigma}$ are admissible for
Lemma~\ref{10.1:lem-unit-shell-bound} with the same constant for every sign family. A version of
the insertion bound requiring the supports to be pairwise at least $\mathsf{d}$ apart, with a
constant depending on $\mathsf{d}$, would not serve in the step leading to
\eqref{10.1:eq-signed-array-a}.

Finally, the argument uses none of the following: the far-sphere flux hypothesis
$(\mathrm{IR\text{-}fl})_{\omega_0}$ of clause (vi) of Theorem~0, any decay of correlations, or any
expansion of the shell term in the running coupling at the shell radius.
\end{proof}

\begin{proposition}[Rotation defect of every $n$-point function]
\label{10.1:prop-rotation-defect-all}
Let $n\ge 2$ be an integer and let $\vartheta\in\,]0,1]$ be fixed. Assume the following statements:
\begin{itemize}
\item the reflection-positivity inequality, read at one fixed radius, together with the grid
form of its supporting lemma;
\item the bound on the correlation of two stress-tensor insertions with separated supports,
at that radius;
\item the exact change of reference scale;
\item the ultraviolet half of clause (vi) of Theorem~0, along every diagonal sequence of
cutoffs;
\item clauses (0), (ii) and (iii) of Theorem~0, and the one-loop law of clause (0).
\end{itemize}
No hypothesis on the infrared behaviour enters. In particular, the far-sphere flux hypothesis
$(\mathrm{IR\text{-}fl})_{\omega_0}$ is not assumed.

For $\delta>0$, let $\mathfrak{B}^{\infty}_n(\delta;\vartheta)$ be the set of real
$n$-tuples $\vec f=(f_1,\dots,f_n)$ with the following properties. There is one centre
$y\in\mathbb{R}^4$ with $f_i\in C^{\infty}_c(B(y,\delta))$ for every $i$. The supports of the
$f_i$ are pairwise at distance at least $\vartheta\delta$. Finally,
$\operatorname{Lip} f_i\le\|f_i\|_\infty/(\vartheta\delta)$ for every $i$. Tuples of class $C^1$
with the same properties are admitted by density. The symbols $\Pi$, $R\cdot f$,
$\delta_\omega f$, $|\omega|$ and $g(\delta)$ have the meanings fixed in
Notation~\ref{10.1:not-clustered-tuples} (see \eqref{10.1:eq-clustered-tuples-a}), with
$\mathsf{d}$ replaced by $\delta$.

The regime is the following; every condition in it is one-sided. The coupling and the size
satisfy
\begin{equation}\label{10.1:eq-rotation-defect-all-1}
g\le\gamma_G:=\min\{\gamma_{\Theta},\gamma_{O4}\},\qquad
0<\delta\le\delta_G(g):=\min\bigl\{\tfrac{1}{512},\,L^{-s_{\mathrm{fr}}(g)}\bigr\}.
\end{equation}
Here $\gamma_{\Theta}$ and $\gamma_{O4}$ are coupling ceilings of the antecedent statements
listed above, chosen last. The condition
$\delta\le L^{-s_{\mathrm{fr}}(g)}$ is a ceiling on the size $\delta$, equivalently a floor
$s_{\mathrm{fr}}(g)$ on the depth $\log_L(1/\delta)$.

The limit $\sigma$ is one of the following.
\begin{itemize}
\item On $\mathbb{R}^4$, $\sigma$ is any limit of clause (ii-$\mathbb{R}^4$) of Theorem~0.
\item On tori, $\sigma$ is any limit of clause (ii-$\mathbb{T}$-b) of Theorem~0 on
$\mathbb{T}_m$, with torus side $\mathsf{T}=2L^m\ge 96$ and
$m\ge\max\{m_{\mathbb{T}}(g_-(g)),\,m'(g,1,2)\}$, where $m'(g,1,2)$ is a floor on the torus
exponent $m$. In this case $\delta$ is in addition
subject to the $n$-dependent ceiling
\begin{equation}\label{10.1:eq-rotation-defect-all-2}
\delta\le\delta^{\mathbb{T}}_G(g,n):=\min\bigl\{\delta_G(g),\,L^{-s^{\mathbb{T}}(g,n)}\bigr\}.
\end{equation}
Here $s^{\mathbb{T}}(g,n)$ is the least $s_0$ such that
\begin{equation*}
48L^{s}\ge 16n\bigl(L^{1+\mathsf{o}+\hat b(\bar x_{s+\mathsf{o}})}+1\bigr)
\qquad\text{for all } s\ge s_0 .
\end{equation*}
The offset $\mathsf{o}$ of the depth count and the function $\hat b(\cdot)$ of the inverse
coupling are those of the volume floor of the exact change of reference scale.
\end{itemize}

Then there are a number
\begin{equation}\label{10.1:eq-rotation-defect-all-3}
\mathsf{G}_n(g):=2\sqrt6\cdot 12^3\,
\bigl[C_{\circledast}\,\mathsf{c}_n\,C_{\mathsf{N}}(g)\bigr]^{1/2}
\end{equation}
and a function $\mathsf{S}^{\vee}_{2n}(\delta)\ge 1$ with the properties stated below. In
\eqref{10.1:eq-rotation-defect-all-3} the two constants are
\begin{equation*}
C_{\circledast}=64\cdot 1400\cdot\bigl(\tfrac{256}{7}\bigr)^2,\qquad
\mathsf{c}_n=2\cdot 4^n\,\mathsf{B}_{2n}\,(\mathsf{F}_n\mathsf{B}_n)^2 .
\end{equation*}
Here $\mathsf{B}_{2n}$, $\mathsf{B}_n$ and $\mathsf{F}_n$ are constants indexed by $2n$ and by
$n$, which enter $\mathsf{G}_n(g)$ only through $\mathsf{c}_n$.
The constant $C_{\mathsf{N}}(g)$ is the constant of the bound on the correlation of two
stress-tensor insertions with separated supports, expressed in the units of the reference
coupling $g$. The function satisfies, for a constant $C'_n(g,\vartheta)$,
\begin{equation}\label{10.1:eq-rotation-defect-all-4}
\log\mathsf{S}^{\vee}_{2n}(\delta)\le C'_n(g,\vartheta)\,
\Bigl(1+\log\bigl(g^{-2}+B_\sharp\,(\log_L(1/\delta)+\mathsf{o}+3)\bigr)\Bigr)^{4r_{\mathrm{pr}}+1},
\end{equation}
where $r_{\mathrm{pr}}$ is an exponent not depending on $\delta$; and hence, for every
$\varepsilon>0$, there is a constant $C_{n,\varepsilon}(g,\vartheta)$ with
\begin{equation}\label{10.1:eq-rotation-defect-all-5}
\mathsf{S}^{\vee}_{2n}(\delta)^{n/2}\le C_{n,\varepsilon}(g,\vartheta)\,\delta^{-\varepsilon}.
\end{equation}
Both $\mathsf{G}_n(g)$ and $\mathsf{S}^{\vee}_{2n}$ depend only on
$(L,\mathfrak{p},g,n,\vartheta)$. They do not depend on the sequence of cutoffs, the diagonal
sequence along which the limit is taken, the volume, the position $y$ or the limit point
$\sigma$.

For every such $\sigma$, every $\vec f\in\mathfrak{B}^{\infty}_n(\delta;\vartheta)$, every
$\omega\in\mathfrak{so}(4)$ and every $R\in O(4)$ in the regime, the following hold.
\begin{enumerate}
\item[(a)] The infinitesimal defect satisfies
\begin{equation}\label{10.1:eq-rotation-defect-all-6}
\Bigl|\sum_{i=1}^{n}S_n[\sigma](f_1,\dots,\delta_\omega f_i,\dots,f_n)\Bigr|
\le\mathsf{G}_n(g)\,\mathsf{S}^{\vee}_{2n}(\delta)^{n/2}\,\delta\,|\omega|\,\Pi .
\end{equation}
\item[(b)] The finite defect satisfies
\begin{equation}\label{10.1:eq-rotation-defect-all-7}
\bigl|S_n[\sigma](R\cdot\vec f)-S_n[\sigma](\vec f)\bigr|
\le\pi\sqrt2\,\mathsf{G}_n(g)\,\mathsf{S}^{\vee}_{2n}(\delta)^{n/2}\,\delta\,\Pi .
\end{equation}
By \eqref{10.1:eq-rotation-defect-all-5}, the right-hand side of
\eqref{10.1:eq-rotation-defect-all-7} is at most
\begin{equation*}
\pi\sqrt2\,\mathsf{G}_n(g)\,C_{n,\varepsilon}(g,\vartheta)\,\delta^{1-\varepsilon}\,\Pi
\qquad\text{for every }\varepsilon>0 .
\end{equation*}
\item[(c)] In terms of the running coupling $g(\delta)$: for every $\eta\in\,]0,1[$ there is a
constant $A_{n,\eta}(g,\vartheta)$ such that both left-hand sides of
\eqref{10.1:eq-rotation-defect-all-6} and \eqref{10.1:eq-rotation-defect-all-7} are bounded by
\begin{equation}\label{10.1:eq-rotation-defect-all-8}
A_{n,\eta}(g,\vartheta)\,
\exp\Bigl[-(1-\eta)\,\frac{12\pi^2}{11N^2}\,g(\delta)^{-2}\Bigr]\,\Pi .
\end{equation}
This is a bound smaller than every power of $g(\delta)^2$.
\item[(d)] Together with the rotation defect of two- and three-point functions
(Proposition~\ref{10.1:prop-rotation-defect-two-three}), applied, where its own hypotheses
also hold, to tuples of
$\mathfrak{B}^{\infty}_n(\delta;\vartheta)\subset\mathfrak{B}_n(\delta;\vartheta)$, the
defect $\bigl|S_n[\sigma](R\cdot\vec f)-S_n[\sigma](\vec f)\bigr|$ is bounded for all $n\ge 2$ by
\begin{equation}\label{10.1:eq-rotation-defect-all-9}
C_n(g,\vartheta)\,\delta^{\max\{1,\,4-n\}}\,\mathsf{S}^{\vee}_{2n}(\delta)^{n/2}\,\Pi ,
\end{equation}
where, for $n\in\{2,3\}$, $C_n(g,\vartheta)$ is the constant of
Proposition~\ref{10.1:prop-rotation-defect-two-three}, and, for $n\ge 4$,
$C_n(g,\vartheta):=\pi\sqrt2\,\mathsf{G}_n(g)$.
For $n\in\{2,3\}$ the power of $\delta$ given by
Proposition~\ref{10.1:prop-rotation-defect-two-three} is the better one, and that bound
carries no factor $\mathsf{S}^{\vee}_{2n}$.
\end{enumerate}

The proposition does not assert the following:
\begin{itemize}
\item the equality $S_n[\sigma](R\cdot\vec f)=S_n[\sigma](\vec f)$ at any $\delta>0$;
\item a relative form of (a) or (b) for $n\ge 3$;
\item for $n\ge 4$, a bound on the lattice defect in the form
$\limsup_{K}\sup_m$ (such a form is part of
Proposition~\ref{10.1:prop-rotation-defect-two-three} for $n\in\{2,3\}$ only);
\item that $\mathsf{S}^{\vee}_{2n}(\delta)$ grows only polylogarithmically in $1/\delta$, except
under the verbatim reading of the dictionary of units.
\end{itemize}
\end{proposition}

\begin{proposition}[No local rotation invariance on the torus]\label{10.1:prop-torus-not-invariant}
Assume the hypotheses (H1)--(H7) of Theorem~0 with $N=2$. Fix a real function
$\varphi\in C^\infty_c(B(0,1))$ with $\varphi\ge 0$ and $\varphi\not\equiv 0$, and put
\begin{equation*}
\vartheta_{\varphi}:=\min\Bigl\{1,\ \frac{\int\varphi}{\|\varphi\|_\infty},\
\frac{\|\varphi\|_\infty}{\operatorname{Lip}\varphi}\Bigr\}.
\end{equation*}
Let the reference coupling satisfy $g\le\gamma^{\mathbb{T}}_W(L,\mathfrak{p},\gamma,\vartheta_{\varphi};\varrho)$,
the coupling ceiling of clause (iv) of Theorem~0, and let $m\ge m_{\mathbb{T}}(g_-(g))$, as in (H5);
the tori $\mathbb{T}_m$ with $m\ge m_{\mathbb{T}}(g_-(g))$ are the tori of the family, as in
Notation~\ref{10.1:not-clustered-tuples}.
Assume further:
\begin{enumerate}
\item[(a)] clause (ii-$\mathbb{T}$-b) of Theorem~0, which provides the subsequential continuum
limits on $\mathbb{T}_m$;
\item[(b)] clause (iii-a) of Theorem~0;
\item[(c)] the lower member of clause (iv) of Theorem~0 for the one shape $\varphi$, on every torus
of the family;
\item[(d)] positivity of the lattice transfer operator, and the calculus of slab operators built from it.
\end{enumerate}
Then for every subsequential limit $\sigma$ of clause (ii-$\mathbb{T}$-b) on $\mathbb{T}_m$, every
$\omega\in\mathfrak{so}(4)\setminus\{0\}$ and every $d_0>0$ there are real functions
$g_1,g_2\in C^\infty(\mathbb{T}_m)$ with disjoint supports contained in a common ball of diameter
less than $d_0$ such that
\begin{equation}\label{10.1:eq-torus-not-invariant-1}
S_2[\sigma](\delta_{\omega}g_1,g_2)+S_2[\sigma](g_1,\delta_{\omega}g_2)\neq 0,
\qquad
\delta_{\omega}g:=-(\omega(x-y))\cdot\nabla g,
\end{equation}
where $S_2[\sigma]$ is the two-point Schwinger function of $\sigma$, $\delta_{\omega}$ is the
infinitesimal rotation of \eqref{10.1:eq-clustered-tuples-a} (Notation~\ref{10.1:not-clustered-tuples}),
and $y$ is a centre of that ball; the value of the left-hand side does not depend on the choice of
the centre. Equivalently, writing
\begin{equation*}
S_2[\sigma](f_1,f_2)=\langle T,f_1\star\check f_2\rangle,\qquad T\in\mathcal{D}'(\mathbb{T}_m\setminus\{0\}),
\end{equation*}
with $\check f(x)=f(-x)$, and $u$ denoting the variable of $T$ (the separation): for no $\omega\neq 0$
and no $d_0>0$ does the distribution $(\omega u)\cdot\nabla_u T$ vanish on $B(0,d_0)\setminus\{0\}$.
\end{proposition}

\begin{proof}
The conclusion is derived from hypotheses (a)--(d) in the theorem on torus limits of
the two-point function, which shows that on a torus of the family positivity of the lattice transfer
operator, together with the slab-operator calculus, turns exact rotation invariance of small pairs
into a constant two-point kernel, and that such a kernel is excluded by hypothesis (c); the mechanism
of that deduction is described in the remark that follows.
\end{proof}

The proposition does not use the rotational Ward identity or its ultraviolet input, the stress tensor,
or any infrared input. It gives no lower bound on the size of the rotation defect, and it says nothing
about $n$-point functions with $n\ge 3$.

\begin{remark}[The mechanism behind Proposition~\ref{10.1:prop-torus-not-invariant}]
Under the hypotheses of Proposition~\ref{10.1:prop-torus-not-invariant}, on every periodic torus of
the family the continuum limit is not exactly rotation invariant at any
scale, however small. The mechanism is the following. Positivity of the transfer operator gives
separate real-analyticity of the two-point function in each coordinate of the separation; this would
spread exact invariance of small pairs over the whole cell and about a second centre, and would
thereby force a constant two-point kernel, which the lower bound of clause (iv) of Theorem~0
forbids. This is why clause (vi) of Theorem~0 carries a hypothesis in infinite volume, the far-sphere
flux hypothesis $(\mathrm{IR\text{-}fl})_{\omega_0}$ (the flux of the lattice rotation current through
far spheres vanishes along one sequence of radii), and why neither the rotation-defect bound of
Proposition~\ref{10.1:prop-rotation-defect-two-three} nor the corresponding bound for every $n$ can
be sharpened to exactness by any argument valid torus by torus.
\end{remark}

\begin{corollary}[The torus test]
Let $g$ and the family of tori $\mathbb{T}_m$, $m\ge m_{\mathbb{T}}(g_-(g))$, be as in
Proposition~\ref{10.1:prop-torus-not-invariant}, under its hypotheses. Say that a subsequential limit
$\sigma$ has \emph{exact rotation invariance of small clusters} if there is $d_0>0$ such that the
left-hand side of \eqref{10.1:eq-torus-not-invariant-1} vanishes for every $\omega\in\mathfrak{so}(4)$
and every pair of real test functions $g_1,g_2$ with disjoint supports contained in a common ball of
diameter less than $d_0$. Then no set of statements each of which holds for the limits on one torus
of the family implies exact rotation invariance of small clusters. Consequently a route to exact
rotation invariance of small clusters on $\mathbb{R}^4$ must contain a statement that is false or
meaningless for the limits on every torus of the family.
\end{corollary}

\begin{proof}
Fix a torus $\mathbb{T}_m$ of the family, and suppose that a set of statements, each of which holds
for every subsequential limit of clause (ii-$\mathbb{T}$-b) on $\mathbb{T}_m$, implied exact rotation
invariance of small clusters; then every such limit $\sigma$ on $\mathbb{T}_m$ would admit $d_0>0$ for
which the left-hand side of \eqref{10.1:eq-torus-not-invariant-1} vanishes for all $\omega$ and all
admissible pairs $g_1,g_2$, whereas Proposition~\ref{10.1:prop-torus-not-invariant}, applied to
$\sigma$, to any $\omega\neq 0$ and to this $d_0$, produces a pair for which it does not vanish. This
contradiction proves the first assertion. For the second, a route to exact rotation invariance of
small clusters rests on finitely many statements and on a deduction from them that is valid for the
limits on any torus on which those statements are meaningful and hold; their conjunction is a single
statement asserted by the route, and it implies exact rotation invariance of small clusters. If this
conjunction held (and were meaningful) for the limits on some torus of the family, the set consisting
of it alone would contradict the first assertion on that torus. Hence it is false or meaningless for
the limits on every torus of the family, which is the second assertion.
\end{proof}

\begin{remark}[Two-point rotation invariance under diagonal reflection positivity]
\label{10.1:rem-diagonal-positivity}
This remark concerns the two-point function only. It is an implication from a hypothesis that
Theorem~0 does not assume and that is not proved, and it asserts nothing about
the limits of Theorem~0 beyond that implication.

Let $e_0,\dots,e_3$ be the standard basis of $\mathbb{R}^4$. Put $v_{01}:=(e_0+e_1)/\sqrt{2}$
and let $\mathsf{r}_{01}$ be the reflection in the hyperplane $v_{01}^{\perp}$. For a test
function $f$ and $R\in O(4)$ put $R\cdot f:=f\circ R^{-1}$. For a limit $\sigma$ we write
$S_2[\sigma]$ for its connected curvature two-point function.

There is a pure number $\vartheta_0\in\,]0,1]$ with the following property. Let
$g\le\gamma^{\mathbb{T}}_W(\vartheta_0)$, which is the coupling ceiling of clause~(iv) of
Theorem~0 taken at $\vartheta_0$. Let $\sigma$ be any subsequential continuum limit on
$\mathbb{R}^4$, at reference coupling $g$, of the family of clause~(ii-$\mathbb{R}^4$) of
Theorem~0. Assume the following.
\begin{itemize}
\item[(a)] $\sigma$ is reflection positive on one-field vectors at the coordinate mirrors and
invariant under $\mathbb{R}^4\rtimes W_4$; this is the part of clause~(ii-$\mathbb{R}^4$) of
Theorem~0 that is used.
\item[(b)] $\sigma$ satisfies clause~(iii-b) of Theorem~0.
\item[(c)] The upper member of clause~(iv) of Theorem~0, with its logarithmic factor, holds at
reference coupling $g$ for the family of which $\sigma$ is a limit.
\item[(d)] $S_2[\sigma]$ is reflection positive on one-field vectors at the diagonal mirror
$v_{01}^{\perp}$, in the sense in which clause~(ii-$\mathbb{R}^4$) asserts reflection positivity
on one-field vectors at the coordinate mirrors, the coordinate mirror being replaced by
$v_{01}^{\perp}$ and the reflection in it by $\mathsf{r}_{01}$.
\end{itemize}
Then
\begin{equation}\label{10.1:eq-diagonal-positivity-1}
S_2[\sigma](R\cdot f_1,\,R\cdot f_2)\ =\ S_2[\sigma](f_1,f_2)
\end{equation}
for every $R\in O(4)$ and all real $f_1,f_2\in C^1_c(\mathbb{R}^4)$ with disjoint supports.
The deduction of \eqref{10.1:eq-diagonal-positivity-1} from (a)--(d) is carried out in the
theorem deducing two-point $O(4)$ invariance on $\mathbb{R}^4$ from (a)--(d); it is not
reproduced here. Hypothesis~(b) enters it only qualitatively.

Clause~(ii-$\mathbb{R}^4$) gives invariance under $\mathbb{R}^4\rtimes W_4$. By this
invariance, hypothesis~(d) is equivalent to the same
positivity at all twelve diagonal mirrors $\bigl((e_\mu\pm e_\nu)/\sqrt{2}\bigr)^{\perp}$,
$0\le\mu<\nu\le 3$.

Diagonal reflection positivity holds for the Wilson weight on $\mathbb{Z}^4$ and fails on
every straight torus (the lemma on diagonal reflection positivity of the Wilson weight).
For the periodic family of tori, hypothesis~(d) would
amount to independence of the limit from the shape of the box. It is not proved, and it is not
assumed by Theorem~0.

The statement is restricted to two points. For $n\ge 3$ the item that is missing is
positivity for configurations rotated by an element of $O(4)\setminus W_4$, and that
positivity is itself the conclusion sought.

Without any hypothesis, no bounded, non-constant, continuous, translation-invariant two-point
kernel $T$ on $\mathbb{R}^4\setminus\{0\}$ is reflection positive, in the sense of
hypothesis~(d) applied to the bilinear form
$(f_1,f_2)\mapsto\langle T,f_1\star\check f_2\rangle$, where $\check f_2(x)=f_2(-x)$ and
$\star$ is convolution, at all sixteen mirrors of the
hyperoctahedral group; this is proved in the theorem on hyperoctahedrally
reflection-positive translation-invariant kernels.
\end{remark}

\begin{definition}[The lattice stress-tensor density, the spin-channel tensor, and the orbital
(vector-channel) remainder of the rotational Ward identity]
\label{10.2:def-orbital-remainder}
Let $\varepsilon = L^{-K}$ be the lattice spacing at the cutoff, and let $x$ range over the
sites of the lattice of spacing $\varepsilon$. The following objects are as fixed earlier in
this chapter:
\begin{itemize}
\item the set $\mathbb{D}$ of signed directions and the relation $\mu \perp \nu$;
\item the bond variables $U(x,\mu)$, the plaquette holonomies $U_{\mu\nu}(x)$ and the map $\pi$;
\item the plaquette variables $P_{\mu\nu}(x) = \pi(U_{\mu\nu}(x))$, which are
$\operatorname{Ad}$-covariant at $x$;
\item the clover curvature $\hat F_{\mu\nu}$, which is $\operatorname{Ad}$-covariant at its base
point, where $\operatorname{Ad}_U \Phi = U \Phi U^{-1}$;
\item the lattice field equation $\mathcal{E}_{x,\mu}$ and the $\operatorname{Ad}$-invariant inner
product $\langle\cdot,\cdot\rangle$ on $\mathfrak{su}(N)$;
\item the lattice differences $\nabla^-_n$, $\nabla^{\mathrm{sym}}_n$ and $\Delta_\lambda$, which
act on real-valued site functions, and the covariant forward difference $\nabla^+_\lambda$, which
acts on $\operatorname{Ad}$-covariant site fields.
\end{itemize}
Explicitly, for a real-valued site function $f$, an $\operatorname{Ad}$-covariant site field
$\Phi$, $n \in \{1,\dots,4\}$ and $\lambda \in \mathbb{D}$,
\begin{align*}
\nabla^-_n f(x) &= \varepsilon^{-1}\bigl[f(x) - f(x-\varepsilon e_n)\bigr], \\
\nabla^{\mathrm{sym}}_n f(x) &= (2\varepsilon)^{-1}\bigl[f(x+\varepsilon e_n)
- f(x-\varepsilon e_n)\bigr], \\
\Delta_\lambda f(x) &= \varepsilon^{-1}\bigl[f(x+\varepsilon e_\lambda) - f(x)\bigr], \\
(\nabla^+_\lambda \Phi)(x) &= \varepsilon^{-1}\bigl[\operatorname{Ad}_{U(x,\lambda)}
\Phi(x+\varepsilon e_\lambda) - \Phi(x)\bigr].
\end{align*}
The symbol $\nabla^+$ is used below only in this covariant sense. Since $e_{-\nu} = -e_\nu$, for
$\nu \in \mathbb{D}$ we have
\begin{equation*}
\Delta_{-\nu} f(x) = \varepsilon^{-1}\bigl[f(x-\varepsilon e_\nu) - f(x)\bigr],
\end{equation*}
which for $\nu = n \in \{1,\dots,4\}$ equals $-\nabla^-_n f(x)$.
For $\nu, \rho \in \mathbb{D}$ put
\begin{equation}\label{10.2:eq-orbital-remainder-1}
\mathcal{T}_{\nu\rho}(x) := \frac{1}{2N\varepsilon^2}
\sum_{\mu \in \mathbb{D},\, \mu \perp \nu} \langle \hat F_{\rho\mu}(x), P_{\mu\nu}(x)\rangle .
\end{equation}
Since $\hat F_{-\rho,\mu} = -\hat F_{\rho\mu}$, we have $\mathcal{T}_{\nu,-\rho} = -\mathcal{T}_{\nu\rho}$.
The field-equation density is
\begin{equation*}
\mathcal{J}_\rho(x) := \frac{1}{2N\varepsilon^3}
\sum_{\mu \in \mathbb{D}} \langle \hat F_{\rho\mu}(x), \mathcal{E}_{x,\mu}\rangle .
\end{equation*}
For $n, \rho \in \mathbb{D}$ we define the following quantities.

The \emph{lattice stress-tensor density} is
\begin{equation}\label{10.2:eq-orbital-remainder-2}
\begin{split}
T^{\mathrm{latt}}_{n\rho}(x) &:= \tfrac12\bigl[\mathcal{T}_{n\rho}(x)
- \mathcal{T}_{-n,\rho}(x+\varepsilon e_n)\bigr] \\
&= \frac{1}{4N\varepsilon^2} \sum_{\mu \in \mathbb{D},\, \mu \perp n}
\bigl\langle \hat F_{\rho\mu}(x) + \operatorname{Ad}_{U(x,n)} \hat F_{\rho\mu}(x+\varepsilon e_n),
P_{\mu n}(x)\bigr\rangle .
\end{split}
\end{equation}
In words, it pairs the six plaquettes through the bond $(x,n)$, each read from $x$ as
``$\mu$ then $n$'', with the average of the clovers at the two ends of the bond, transported
to $x$. For every $n \in \mathbb{D}$ the first line gives
$T^{\mathrm{latt}}_{n\rho}(x) = -T^{\mathrm{latt}}_{-n,\rho}(x+\varepsilon e_n)$.

The \emph{spin-channel tensor} is $T^s_{[\nu\rho]} := \frac12(T^s_{\nu\rho} - T^s_{\rho\nu})$,
where
\begin{equation*}
T^s_{n\rho}(x) := \tfrac12\bigl[T^{\mathrm{latt}}_{n\rho}(x)
+ T^{\mathrm{latt}}_{n\rho}(x - \varepsilon e_n)\bigr].
\end{equation*}

The remaining quantities are
\begin{equation*}
Y_\rho(x) := \frac{1}{4N\varepsilon^2} \sum_{\mu,\nu \in \mathbb{D},\, \mu \perp \nu}
\bigl\langle (\nabla^+_\nu \hat F_{\rho\mu})(x - \varepsilon e_\nu),
P_{\mu\nu}(x - \varepsilon e_\nu)\bigr\rangle ,
\end{equation*}
\begin{equation*}
S^{\mathrm{latt}}(x) := \tfrac14 \sum_{\mu,\nu=1}^4
\operatorname{tr}\bigl(\hat F_{\mu\nu}(x) \hat F_{\mu\nu}(x)\bigr) .
\end{equation*}
Finally, the \emph{orbital remainder}, which is the vector channel of the bulk term of the
rotational Ward identity, is
\begin{equation*}
X_\rho(x) := \varepsilon^{-2}\bigl[Y_\rho(x) - \nabla^{\mathrm{sym}}_\rho S^{\mathrm{latt}}(x)\bigr].
\end{equation*}

These quantities satisfy the following identity exactly, as functions on configurations:
\begin{equation}\label{10.2:eq-orbital-remainder-3}
\mathcal{J}_\rho = \sum_{n=1}^4 \nabla^-_n T^{\mathrm{latt}}_{n\rho} - Y_\rho
= \sum_{n=1}^4 \nabla^-_n T^{\mathrm{latt}}_{n\rho} - \nabla^{\mathrm{sym}}_\rho S^{\mathrm{latt}}
- \varepsilon^2 X_\rho .
\end{equation}

From the first line of \eqref{10.2:eq-orbital-remainder-2} we obtain, for
$\nu, \rho \in \mathbb{D}$,
\begin{equation}\label{10.2:eq-orbital-remainder-4}
T^s_{\nu\rho}(x) = \tfrac14 \sum_{s = \pm} s\bigl[\mathcal{T}_{s\nu,\rho}(x)
+ \mathcal{T}_{s\nu,\rho}(x - s\varepsilon e_\nu)\bigr].
\end{equation}
This is a $\mathbb{D}$-uniform formula, odd in each displayed index. Hence, by the
$W_4$-naturality lemma for plaquette pairings earlier in this chapter (where the notions of a
$\mathbb{D}$-uniform formula, of a $W_4$-tensor family, vector or scalar at a site, and of
naturality at the midpoint of a bond are fixed), $T^s$ is a rank-two $W_4$-tensor family at $x$.
By the same lemma, $Y_\rho$ and $X_\rho$ are $W_4$-vectors at $x$ and $S^{\mathrm{latt}}$ is a
$W_4$-scalar, while $T^{\mathrm{latt}}_{n\rho}$ is natural at the midpoint of the bond $(x,n)$.
\end{definition}

\begin{proof}[Proof of the second form in \eqref{10.2:eq-orbital-remainder-2}, of
\eqref{10.2:eq-orbital-remainder-4} and of \eqref{10.2:eq-orbital-remainder-3}]
\emph{Reflection of plaquettes.} Let $\mu \perp \nu$ in $\mathbb{D}$ and put
$y = x - \varepsilon e_\nu$. Writing out the definition of $U_{\mu\nu}$ and using the reversal
rule $U(z,\lambda)^{-1} = U(z+\varepsilon e_\lambda, -\lambda)$, we find
\begin{align*}
U_{\mu\nu}(y)^{-1} &= U(y,\nu)\, U(x,\mu)\, U(x+\varepsilon e_\mu, -\nu)\,
U(y+\varepsilon e_\mu, -\mu), \\
U_{\mu,-\nu}(x) &= U(x,\mu)\, U(x+\varepsilon e_\mu, -\nu)\,
U(y+\varepsilon e_\mu, -\mu)\, U(y,\nu).
\end{align*}
Since $U(y,\nu) = U(x,-\nu)^{-1}$, it follows that
$U_{\mu,-\nu}(x) = U(x,-\nu)\, U_{\mu\nu}(y)^{-1}\, U(x,-\nu)^{-1}$.

Now $\pi(V^{-1}) = -\pi(V)$. Moreover $\pi(WVW^{-1}) = \operatorname{Ad}_W \pi(V)$, because the
trace is invariant under conjugation. Together these give
\begin{equation}\label{10.2:eq-orbital-remainder-5}
P_{\mu,-\nu}(x) = -\operatorname{Ad}_{U(x,-\nu)} P_{\mu\nu}(x - \varepsilon e_\nu).
\end{equation}

\emph{Second form of $T^{\mathrm{latt}}$.} Apply \eqref{10.2:eq-orbital-remainder-5} at the
point $x+\varepsilon e_n$ with $\nu = n$:
\begin{equation*}
P_{\mu,-n}(x+\varepsilon e_n) = -\operatorname{Ad}_{U(x+\varepsilon e_n,-n)} P_{\mu n}(x).
\end{equation*}
Note that $\mu \perp -n$ if and only if $\mu \perp n$. Insert this into
\eqref{10.2:eq-orbital-remainder-1} for $\mathcal{T}_{-n,\rho}(x+\varepsilon e_n)$, use the
$\operatorname{Ad}$-invariance of $\langle\cdot,\cdot\rangle$, and use
$U(x+\varepsilon e_n,-n)^{-1} = U(x,n)$. This gives
\begin{equation*}
-\mathcal{T}_{-n,\rho}(x+\varepsilon e_n) = \frac{1}{2N\varepsilon^2}
\sum_{\mu \perp n} \langle \operatorname{Ad}_{U(x,n)} \hat F_{\rho\mu}(x+\varepsilon e_n),
P_{\mu n}(x)\rangle .
\end{equation*}
Adding $\mathcal{T}_{n\rho}(x)$ and halving yields the second line of
\eqref{10.2:eq-orbital-remainder-2}. Formula \eqref{10.2:eq-orbital-remainder-4} follows by
inserting the first line of \eqref{10.2:eq-orbital-remainder-2} at $x$ and at
$x - \varepsilon e_\nu$: the terms with $s = +$ are
$\mathcal{T}_{\nu\rho}(x) + \mathcal{T}_{\nu\rho}(x-\varepsilon e_\nu)$, and those with $s = -$
are $-\mathcal{T}_{-\nu,\rho}(x) - \mathcal{T}_{-\nu,\rho}(x+\varepsilon e_\nu)$.

\emph{The field equation as a difference.} By the definition of the covariant difference,
\begin{equation*}
\varepsilon \nabla^+_{-\nu} P_{\mu\nu}(x)
= \operatorname{Ad}_{U(x,-\nu)} P_{\mu\nu}(x-\varepsilon e_\nu) - P_{\mu\nu}(x)
= -P_{\mu,-\nu}(x) - P_{\mu\nu}(x),
\end{equation*}
where the second equality is \eqref{10.2:eq-orbital-remainder-5}. Sum over $\nu \perp \mu$; the
reindexing $\nu \mapsto -\nu$ shows that both terms sum to $-\mathcal{E}_{x,\mu}$. Hence
\begin{equation*}
\langle \hat F_{\rho\mu}(x), \mathcal{E}_{x,\mu}\rangle
= -\tfrac12 \sum_{\nu \perp \mu} \langle \hat F_{\rho\mu}(x),
\varepsilon \nabla^+_{-\nu} P_{\mu\nu}(x)\rangle .
\end{equation*}

\emph{Leibniz rule.} We use the $\operatorname{Ad}$-invariance of $\langle\cdot,\cdot\rangle$,
the relation $U(x,-\nu)^{-1} = U(x-\varepsilon e_\nu, \nu)$, and
\begin{equation*}
\operatorname{Ad}_{U(x-\varepsilon e_\nu,\nu)} \hat F_{\rho\mu}(x)
= \hat F_{\rho\mu}(x-\varepsilon e_\nu)
+ \varepsilon (\nabla^+_\nu \hat F_{\rho\mu})(x-\varepsilon e_\nu),
\end{equation*}
which is the definition of the covariant difference $\nabla^+_\nu$ applied to the
$\operatorname{Ad}$-covariant site field $\hat F_{\rho\mu}$.
These give
\begin{equation*}
\langle \hat F_{\rho\mu}(x), \varepsilon \nabla^+_{-\nu} P_{\mu\nu}(x)\rangle
= \varepsilon \Delta_{-\nu} \langle \hat F_{\rho\mu}, P_{\mu\nu}\rangle(x)
+ \varepsilon \langle (\nabla^+_\nu \hat F_{\rho\mu})(x-\varepsilon e_\nu),
P_{\mu\nu}(x-\varepsilon e_\nu)\rangle .
\end{equation*}

\emph{Summation.} Multiply by $-\frac12 \cdot (2N\varepsilon^3)^{-1}$, sum over $\nu \perp \mu$
and over $\mu \in \mathbb{D}$, and exchange the two sums. The left-hand side becomes
$\mathcal{J}_\rho(x)$. The second part becomes $-Y_\rho(x)$. By
\eqref{10.2:eq-orbital-remainder-1}, the first part becomes
\begin{equation*}
-\tfrac12 \sum_{\nu \in \mathbb{D}} \Delta_{-\nu} \mathcal{T}_{\nu\rho}(x)
= \frac{1}{2\varepsilon} \sum_{\nu \in \mathbb{D}}
\bigl[\mathcal{T}_{\nu\rho}(x) - \mathcal{T}_{\nu\rho}(x-\varepsilon e_\nu)\bigr].
\end{equation*}
Group $\nu = n$ with $\nu = -n$, for which $x - \varepsilon e_\nu = x + \varepsilon e_n$. The
pair contributes
\begin{align*}
&\frac{1}{2\varepsilon}\bigl\{\mathcal{T}_{n\rho}(x) - \mathcal{T}_{n\rho}(x-\varepsilon e_n)
+ \mathcal{T}_{-n,\rho}(x) - \mathcal{T}_{-n,\rho}(x+\varepsilon e_n)\bigr\} \\
&\quad = \frac{1}{\varepsilon}\Bigl\{\tfrac12\bigl[\mathcal{T}_{n\rho}(x)
- \mathcal{T}_{-n,\rho}(x+\varepsilon e_n)\bigr]
- \tfrac12\bigl[\mathcal{T}_{n\rho}(x-\varepsilon e_n) - \mathcal{T}_{-n,\rho}(x)\bigr]\Bigr\} \\
&\quad = \nabla^-_n T^{\mathrm{latt}}_{n\rho}(x),
\end{align*}
by \eqref{10.2:eq-orbital-remainder-2} at $x$ and at $x - \varepsilon e_n$. This proves the
first equality in \eqref{10.2:eq-orbital-remainder-3}. The second equality is the definition of
$X_\rho$.
\end{proof}

\begin{proposition}[The vector-channel statement: the $x$-weighted orbital remainder is irrelevant in
the limit (conservation of the lattice stress tensor in the limit)]
\label{10.2:prop-vector-channel}
Let $\varepsilon := L^{-K}$ be the bare lattice spacing. Let $X_\rho$ be the orbital remainder of
Definition~\ref{10.2:def-orbital-remainder}, and let $\xi^\chi_\rho$ be the cut-off Killing field of
the rotation, where $\chi$ is the cut-off function of that field, admissible in the sense fixed
with the rotational Ward identity. Put
\begin{equation}\label{10.2:eq-vector-channel-1}
  \mathfrak{V}_\chi := g_0^{-2}\,\varepsilon^{6}\sum_{x}\sum_{\rho}\xi^\chi_\rho(x)\,X_\rho(x).
\end{equation}
This is the $x$-weighted, $\chi$-smeared $W_4$-vector insertion of engineering dimension seven.

Let $\langle\,\cdot\,\rangle_K$ denote the expectation in the lattice measure on the torus
$\mathbb{T}_m$ at cutoff $K$, the side index $m$ being held fixed throughout the statement, and
$\langle\,\cdot\,;\dots;\cdot\,\rangle^{\mathrm{conn}}_K$ the corresponding connected correlation.
Then there is an exponent $\eta\in(0,2)$, not depending on $K$, on $n$, on the tuple below or on
$\chi$, with the following property. For every $n \ge 2$ and every tuple
$\vec f = (f_1,\dots,f_n)$ of real functions $f_i \in C^1_c(\mathbb{R}^4)$ with pairwise disjoint
compact supports, hence pairwise separated supports, there is a constant $C(\vec f,\chi)$,
independent of $K$, such that
\begin{equation}\label{10.2:eq-vector-channel-2}
  \bigl|\bigl\langle\mathfrak{V}_\chi;\mathcal{O}(f_1);\dots;\mathcal{O}(f_n)
  \bigr\rangle^{\mathrm{conn}}_K\bigr|
  \le C(\vec f,\chi)\,\varepsilon^{2-\eta}.
\end{equation}
In particular the left side of \eqref{10.2:eq-vector-channel-2} tends to zero along the cutoffs
$K \to \infty$.

Moreover, the vacuum term of $\mathfrak{V}_\chi$, that is, its expectation
$\langle\mathfrak{V}_\chi\rangle_K$ without observables, vanishes exactly for every $K$.
\end{proposition}

A weaker form would also serve. In it, the left side of \eqref{10.2:eq-vector-channel-2} is only
required to tend to zero as $K \to \infty$, possibly more slowly than any power of the spacing
$\varepsilon$. This form suffices for clause~(vi) of Theorem~0 as well, since clause~(vi) admits
either form.

The proposition is proved, in the four sections of this chapter that follow this one, as an
implication from the vector package, the stress-tensor-insertion correlation theorem in its format
and transport halves, the reduction theorem, and the one-loop anisotropy inequality with exponent
one; each of these is printed there under a heading that records whether and how it is proved.

\begin{remark}[Why the vector channel is a separate input]\label{10.2:rem-design-decision}
With the canonical objects of Definition~\ref{10.2:def-orbital-remainder}, the bulk term that
breaks rotation invariance in the rotational Ward identity is exactly the sum of three pieces:
a spin-channel term built from the antisymmetrised spin-channel density $T^{s}_{[\nu\rho]}$, an
orbital (vector-channel) term in which the orbital remainder $X_\rho$ is weighted by the
position $x$, and a shell term supported in an annulus. The coefficient of the spin-channel term
at second order in the lattice spacing is a total derivative. The position-weighted vector term,
by contrast, carries a tree-level functional that is not identically zero, with value
$-\pi^2/(64N)$ on an explicit compactly supported configuration, and $X_\rho$ is not an exact
lattice divergence; hence, for the canonical decomposition, the breaking does not lie wholly in
the spin channel, and the
decomposition and these statements are proved as an implication from the exact lattice identity
of Section~10.1 and an explicit tree-level computation. A weaker exact decomposition, which would
remove the need for a separate input by writing $X_\rho$ as the lattice divergence of a local
gauge-invariant tensor, plus a spin-type density, plus a density proportional to the lattice
field equation, is not proved. Its order-by-order variant reduces to one identity at engineering
dimension seven, which expresses the dimension-seven classical coefficient of $X_\rho$ as the
divergence of a gauge-invariant local tensor plus a field-equation density. Two pieces of that
identity are missing: the pointwise dimension-seven coefficient of $X_\rho$ itself, of which only
integrated pairings are computed, and a gauge-invariant local antiderivative of that coefficient
corrected by a field-equation density, a sum known only to have vanishing integral for every
compactly supported configuration and hence to be the divergence of a local current that is not
manifestly gauge invariant. For this reason the position-weighted vector term is treated through
its own statement, Proposition~\ref{10.2:prop-vector-channel}.
\end{remark}

\begin{definition}[Units of length, lattice differences and the weighted insertion]
\label{10.3:not-units-insertion}
In what follows we work in the setting of the correlation theorem with $N=2$, in the notation
(with hats) of the re-filing lemmas. Put $K'=K-m$; it is the last level of the local run of
renormalisation steps in the setting of the correlation theorem. For
$0\le j\le K'$ let $T^{(j)}$ denote the lattice of level $j$ of that run. The spacing of
$T^{(K')}$ is the unit of length, which is the unit of length of the main theorem. We set
\begin{equation}\label{10.3:eq-units-insertion-1}
a:=L^{-K'},\qquad a_j:=L^{j}a=L^{-(K'-j)},\qquad \lambda_*:=(40M)^{-1},\qquad V:=\mathbb{R}^4 ;
\end{equation}
$a$ is the bare spacing and $a_j$ is the spacing of $T^{(j)}$. For a function $\psi$ on the torus
of the correlation theorem, with $e_\nu$ the unit vector in direction $\nu$,
\begin{equation}\label{10.3:eq-units-insertion-2}
\nabla_{j,\nu}\psi(y):=\psi(y+a_je_\nu)-\psi(y),\qquad
\langle\chi\rangle_{j,\nu}(y):=\int_0^1\chi(y+sa_je_\nu)\,ds ,
\end{equation}
and $b_\nu(y):=\langle y,y+a_je_\nu\rangle$ are the bonds of $T^{(j)}$. We set
\[
\|\psi\|_{\langle 2\rangle}:=\max\bigl\{\sup|\psi|,\ \sup|\partial\psi|/\lambda_*,\
\sup|\partial^2\psi|/\lambda_*^2\bigr\}.
\]
Distances, differences and the diameters $D_X:=\operatorname{diam}_j(X^+)+1$ of regions $X$
(with $X^+$ as in the setting of the correlation theorem) are said to be \emph{in $j$-units}
when they are measured in units of $a_j$; $\operatorname{Lip}_j$ is the Lipschitz constant with
respect to distance in $j$-units.

Let $\eta$ be a $C^2$ function on the torus of the correlation theorem with values in $V$. In the
applications $\eta=\varphi\xi$ with $\varphi=\chi$ or $\varphi=f_i$, where $\chi$ and $\xi$ are
fixed there; nothing below uses this form. The
\emph{weight function} is $\mathsf{h}:=\eta/a$, one and the same for all levels $j$. With
$c_0\in\{g_0^{-2},1\}$ the normalisation of the insertion, $\mathfrak{K}^{\sharp}_{\eta}[1]$ the
Killing-weighted bare vector unit and $\mathfrak{r}^{\sharp}_{\rho}(p)$ the bare vector unit at
the plaquette $p$, we set
\[
\begin{gathered}
Z(z,t):=\int\exp\Bigl[-g_0^{-2}A(U)+\sum_l z_l\,\mathcal{O}(f_l)
+t\,c_0\,\mathfrak{K}^{\sharp}_{\eta}[1]\Bigr]\,dU,\\
w:=\sum_l z_lf_l,\qquad f_l\in C^2,
\end{gathered}
\]
\[
\mathbb{D}_2:=\Bigl\{z:\ \sum_l|z_l|\,\|f_l\|_{\langle 2\rangle}<r/4\Bigr\},
\]
where $r$ is the radius of the setting of the correlation theorem. For $y\in T^{(j)}$ we put
$\hat w(y):=w(p_y)$, with $p_y$ the point attached to $y$ in the correlation theorem, and the
ratio $\tilde Z/\tilde Z(0)$ of the correlation theorem is written $\mathsf{h}_Z$.

The following symbols of the correlation theorem, of the one-lattice Lipschitz step and of the
re-filing lemmas keep their meaning: $t=L^{j-n}$, $\nu=1+n-j$, $\rho_{n,j}$,
$w_{k,j}\le\tfrac52\,x_j/x_k$, $\theta_k$, $\mathfrak{a}_n=C_{\mathfrak{a}}\alpha_{\cdot}(g_n)$,
$\mathfrak{r}_j$, $\mathfrak{g}_n=g_n^2(\log x_n)^{2q_{\flat}}$, $E_1$, $K_{CE}$, $K_2$,
$C_{RF}$, $\mathsf{a}_i$, $i''=i+\mathsf{a}_i$, $\mathfrak{B}(k)$, $k_1(k)$, $\epsilon_k$,
$\hat\varepsilon_k:=\max_{j\le k}\epsilon_j$,
$\mathsf{a}^*_k:=\max\{k-i+1:\ i\le k<i''\}$; here $q_{\flat}$ is the exponent of
\cite[(2.28)]{B-CMP119} and $q$ is the age rate. The symbols $t$ and $\nu$ in the formulas
$t=L^{j-n}$ and $\nu=1+n-j$ are those of the one-lattice Lipschitz step; they are distinct from
the parameter $t$ of $Z(z,t)$ and from the direction index $\nu$ of $\nabla_{j,\nu}$, and
$\rho_{n,j}$ is distinct from the component index $\rho$ of $\mathsf{h}^{\rho}$.

These objects have the following properties.
\begin{enumerate}
\item[(1)] For $\psi\in C^1$, $\nabla_{j,\nu}\psi=a_j\langle\partial_\nu\psi\rangle_{j,\nu}$; and
$\|\psi\|_{\langle 2\rangle}\le C(M)\|\psi\|_{C^2}$.
\item[(2)] The weight function satisfies, for all $j,\mu,\nu,\rho$,
\begin{equation}\label{10.3:eq-units-insertion-a}
|\mathsf{h}|\le\frac{\|\eta\|_{\langle 2\rangle}}{a},\qquad
\nabla_{j,\nu}\mathsf{h}^{\rho}=L^{j}\langle\partial_\nu\eta^{\rho}\rangle_{j,\nu},\qquad
|\nabla_{j,\mu}\nabla_{j,\nu}\mathsf{h}|\le L^{2j}a\,\sup|\partial^2\eta| .
\end{equation}
\item[(3)] The insertion is expanded over plaquettes:
\[
c_0\,\mathfrak{K}^{\sharp}_{\eta}[1]
=c_0\sum_p\sum_\rho\mathsf{h}^{\rho}(x(p))\,\mathfrak{r}^{\sharp}_{\rho}(p).
\]
\item[(4)] For $z\in\mathbb{D}_2$: $|w|<r/4$, $\operatorname{Lip}_jw\le a_j\lambda_*r/4$ and
$|\nabla_{j,\mu}\nabla_{j,\nu}w|\le a_j^2\lambda_*^2r/4$.
\end{enumerate}
\end{definition}

\begin{proof}
(1) By the fundamental theorem of calculus along the segment from $y$ to $y+a_je_\nu$,
\[
\psi(y+a_je_\nu)-\psi(y)=\int_0^1 a_j\,\partial_\nu\psi(y+sa_je_\nu)\,ds
=a_j\langle\partial_\nu\psi\rangle_{j,\nu}(y).
\]
Since $\lambda_*^{-1}=40M$ and $M=L^{m'}\ge1$, each of the three
quantities in $\|\psi\|_{\langle 2\rangle}$ is at most $(40M)^2$ times the
corresponding supremum of $\psi$, $\partial\psi$, $\partial^2\psi$. With
$\|\psi\|_{C^2}:=\max\{\sup|\psi|,\sup|\partial\psi|,\sup|\partial^2\psi|\}$ this gives the bound
with $C(M)=(40M)^2$.

(2) First, $|\mathsf{h}|=|\eta|/a\le\sup|\eta|/a\le\|\eta\|_{\langle 2\rangle}/a$. Next, by
linearity of $\nabla_{j,\nu}$ and by part (1) applied to $\eta^\rho$,
\[
\nabla_{j,\nu}\mathsf{h}^\rho=a^{-1}\nabla_{j,\nu}\eta^\rho
=\frac{a_j}{a}\,\langle\partial_\nu\eta^\rho\rangle_{j,\nu}
=L^j\langle\partial_\nu\eta^\rho\rangle_{j,\nu},
\]
since $a_j/a=L^j$. Finally, applying part (1) twice,
\[
\nabla_{j,\mu}\nabla_{j,\nu}\eta(y)=a_j^2\int_0^1\!\!\int_0^1
\partial_\mu\partial_\nu\eta(y+sa_je_\nu+s'a_je_\mu)\,ds\,ds',
\]
so $|\nabla_{j,\mu}\nabla_{j,\nu}\mathsf{h}|\le a^{-1}a_j^2\sup|\partial^2\eta|
=L^{2j}a\sup|\partial^2\eta|$, because $a_j^2/a=L^{2j}a$.

(3) This is the lemma expanding the Killing-weighted bare vector unit
$\mathfrak{K}^{\sharp}_{\eta}[1]$ as a sum of plaquette units, read with the weight function
$\mathsf{h}=\eta/a$ evaluated at the barycentres $x(p)$, multiplied by $c_0$.

(4) Let $z\in\mathbb{D}_2$. Then
\[
\begin{gathered}
|w|\le\sum_l|z_l|\sup|f_l|\le\sum_l|z_l|\,\|f_l\|_{\langle 2\rangle}<r/4,\\
\sup|\partial w|\le\lambda_*\sum_l|z_l|\,\|f_l\|_{\langle 2\rangle}<\lambda_*r/4,\qquad
\sup|\partial^2w|<\lambda_*^2r/4 .
\end{gathered}
\]
For $y,y'$ at distance $\delta$ in $j$-units,
$|y-y'|=a_j\delta$, hence $|w(y)-w(y')|\le\sup|\partial w|\,a_j\delta$, which gives
$\operatorname{Lip}_jw\le a_j\lambda_*r/4$. The second-difference formula displayed in part (2),
with $w$ in place of $\eta$, gives
$|\nabla_{j,\mu}\nabla_{j,\nu}w|\le a_j^2\sup|\partial^2w|\le a_j^2\lambda_*^2r/4$.
\end{proof}

\begin{definition}[The H\"older exponent, decay rates and scale sums]\label{10.3:def-exponents-rates}
The following exponents and rates are prescribed before $\gamma$, each by a one-sided condition:
in (e1) and (e2) the prescription is the lower bound, the upper bound $1$ being the range in which
the exponent, respectively the rate, is taken; in (e3) it is the upper bound on $q$.

(e1) The H\"older exponent $\hat\alpha$ of the backgrounds satisfies
\begin{equation}\label{10.3:eq-exponents-rates-a}
\tfrac12<\hat\alpha<1 .
\end{equation}
In Ba{\l}aban's notation $\hat\alpha=1-\beta$, where $\beta$ is the exponent of
\cite[(3.67)]{B-CMP119}, whose bound carries, in the scale indices $j,n$ of that paper, the factor
$O((L^jL^{-n})^{5-\beta})$ with $\beta>0$,
and of \cite[(3.31)]{B-CMP109}, which holds for $0\le\beta\le\beta_0<1$. The correlation theorem
uses the letter $\hat\alpha$ for this exponent in $]0,1[$, fixed at the place reserved for the
exponents in the order in which its parameters are chosen; here it is fixed in the
sub-interval \eqref{10.3:eq-exponents-rates-a}.

(e2) The rate $\vartheta_2$ and the rate $\vartheta_V$ satisfy
\begin{equation}\label{10.3:eq-exponents-rates-b}
L^{-(2\hat\alpha-1)}<\vartheta_2<1,\qquad
\vartheta_V:=\vartheta_2/L,\quad\text{so that}\quad L^{-2\hat\alpha}<\vartheta_V<L^{-1}.
\end{equation}
The first interval is non-empty because $2\hat\alpha-1>0$ by \eqref{10.3:eq-exponents-rates-a}
and $L>1$; the second is obtained from the first by division by $L$. The rate $\vartheta_2$ is
used together with a rate $\vartheta_1$ under the condition $\vartheta_2\le\vartheta_1$, and
$\vartheta_1$ is accordingly prescribed in
$]\max\{L^{-(2\hat\alpha-1)},L^{-1/2}\},1[$. This is a sub-interval of the window
$]\max\{L^{-\hat\alpha},L^{-1/2}\},1[$ in which $\vartheta_1$ is prescribed in the
stress-tensor-insertion correlation theorem (Theorem~\ref{10.4:thm-insertion-correlation}) and in
the vector package, both of which use $\vartheta_1$ together with the condition
$q\le\min\{\tfrac12,\vartheta_1^2\}$ on the age rate,
because $2\hat\alpha-1<\hat\alpha$ by \eqref{10.3:eq-exponents-rates-a}, whence
$L^{-(2\hat\alpha-1)}>L^{-\hat\alpha}$.

(e3) The age rate $q$ satisfies
\begin{equation}\label{10.3:eq-exponents-rates-c}
q\le\min\{\tfrac12,\vartheta_V^2\}.
\end{equation}
Here $q=e^{-c_a}$ is the rate of the lemma on geometric forgetting of old regions, in which
$c_a\ge\log 2$ is
a free number; it is now chosen, after $L$ and $\vartheta_2$ and before $\gamma$, so that in
addition $c_a\ge 2\log(L/\vartheta_2)$. Then $e^{-c_a}\le\tfrac12$ and
$e^{-c_a}\le(\vartheta_2/L)^2=\vartheta_V^2$, which is \eqref{10.3:eq-exponents-rates-c}. Since
$\vartheta_V<\vartheta_2\le\vartheta_1$, the same $q$ also satisfies
$q\le\min\{\tfrac12,\vartheta_1^2\}$, the condition under which $q$ is used together with
$\vartheta_1$ in Theorem~\ref{10.4:thm-insertion-correlation} and in the vector package.

Scale sums. Let $n$ be a level and let $\nu\ge1$ be such that $n+1-\nu$ is a level; the sum over
$\nu$ in $S^V$ below runs over these $\nu$. Let $\rho_{n,j}$ be the quantity of the correlation
theorem attached to the pair of levels $n$ and $j$, of which only the bound recalled in the proof
below is used; $\lambda$ and $c_2$ are the constants appearing in that bound. Abbreviate, in this
definition and its proof only, $\rho:=\rho_{n,n+1-\nu}$, and put
\begin{equation}\label{10.3:eq-exponents-rates-1}
\varsigma^V_{\nu}:=\nu^3\rho^2(1+\rho)\,L^{-2\hat\alpha(\nu-1)},\qquad
S^V:=\sup_n\sum_{\nu\ge1}\varsigma^V_{\nu}\,\vartheta_V^{-\nu},
\end{equation}
\begin{equation}\label{10.3:eq-exponents-rates-2}
C_q^V:=\sum_{i\ge0}(i+1)\,q^{(i-1)_+}\,\vartheta_V^{-i},
\end{equation}
\begin{equation}\label{10.3:eq-exponents-rates-3}
\mathsf{s}_n:=(n+1)\,\vartheta_V^{\,n-2},\qquad
\mathsf{s}_*:=\sup_{n\ge1}\mathsf{s}_n\bigl(1+(3\lambda n+2c_2)^{1/2}\bigr).
\end{equation}
The quantities $S^V$, $C_q^V$ and $\mathsf{s}_*$ are finite and independent of $K$.
\end{definition}

\begin{proof}
Consider first $S^V$. For every $\nu\ge1$,
\[
L^{-2\hat\alpha(\nu-1)}\vartheta_V^{-\nu}
=\vartheta_V^{-1}\bigl(L^{-2\hat\alpha}/\vartheta_V\bigr)^{\nu-1},
\]
and $0<L^{-2\hat\alpha}/\vartheta_V<1$ by \eqref{10.3:eq-exponents-rates-b}. By the correlation
theorem,
\[
\rho_{n,n+1-\nu}^2\le(1+o(1))\bigl(1+(3\lambda\nu+2c_2)g^2\bigr).
\]
The factor $1+o(1)$ being read as bounded independently of $n$ and $K$, the right-hand side is
bounded, uniformly in $n$, by a quantity growing at most linearly in $\nu$; hence $\rho^2(1+\rho)$
is bounded, uniformly in $n$, by a
quantity growing at most like $\nu^{3/2}$. Inserting both facts into
\eqref{10.3:eq-exponents-rates-1}, the general term $\varsigma^V_{\nu}\vartheta_V^{-\nu}$ is
bounded, uniformly in $n$, by $\vartheta_V^{-1}$ times a polynomial in $\nu$ times
$(L^{-2\hat\alpha}/\vartheta_V)^{\nu-1}$, a geometric factor of ratio smaller than one. The series
over $\nu$ therefore converges, with a bound independent of $n$, so $S^V<\infty$; the bound
involves only $\hat\alpha$, $L$, $\vartheta_2$, $\lambda$, $c_2$, $g$ and the bound on the factor
$1+o(1)$, and not $K$.

For $C_q^V$, the term $i=0$ equals $1$, and for $i\ge1$ the general term of
\eqref{10.3:eq-exponents-rates-2} is
$(i+1)\vartheta_V^{-1}(q/\vartheta_V)^{i-1}$. By \eqref{10.3:eq-exponents-rates-c} and
\eqref{10.3:eq-exponents-rates-b}, $q/\vartheta_V\le\vartheta_V<L^{-1}<1$, so the series converges;
its value depends only on $q$, $L$ and $\vartheta_2$.

For $\mathsf{s}_*$, the sequence
$(n+1)\vartheta_V^{\,n-2}(1+(3\lambda n+2c_2)^{1/2})$ is a geometric factor
$\vartheta_V^{\,n}$ with $\vartheta_V<1$ multiplied by $\vartheta_V^{-2}$ and by a factor growing
at most like $n^{3/2}$; it therefore tends to zero as $n\to\infty$ and its supremum over $n\ge1$
is finite. It depends only on $L$, $\vartheta_2$, $\lambda$ and $c_2$, and not on $K$.
\end{proof}

\begin{definition}[Vector-typed pointed families and the improved norm]\label{10.3:def-typed-families}
Let $j$ be a level, let $\mathsf{V}$ be a type space with complexification $\mathsf{V}_{\mathbb{C}}$, and for a
Euclidean symmetry $r$ of $T^{(j)}$ write $r_*$ for its action on $\mathsf{V}_{\mathbb{C}}$.
\begin{enumerate}
\item A \emph{$\mathsf{V}$-typed pointed $E$-family at scale $j$} is a family
$F = \{F(X,\cdot,y;h_0)\}$ that is linear in $h_0 \in \mathsf{V}_{\mathbb{C}}$, satisfies
conditions $(\mathrm{f}1)$ and $(\mathrm{f}2)$ of the correlation theorem, and satisfies the typed covariance
$(\mathrm{f}3)^{\mathsf{V}}$:
\begin{equation}\label{10.3:eq-typed-families-1}
F(rU, ry; r_* h_0) = F(U, y; h_0)
\end{equation}
for the pointed totals $F(U,y;h_0)$ of $F$ and for every Euclidean symmetry $r$ of $T^{(j)}$.
\item For a $\mathsf{V}$-typed pointed $E$-family $F$ at scale $j$, $\|F\|^{\flat}$ denotes the norm
of the correlation theorem, with rate $\kappa$. The \emph{improved norm} $\|F\|^{\mathrm{fr}}$ is defined
in the same way as $\|F\|^{\flat}$, with the weight attached to the localisation domain $X$ taken to be
\begin{equation}\label{10.3:eq-typed-families-2}
e^{(1+4\beta_s)\kappa d_j(X)} = e^{2\kappa d_j(X)}
\end{equation}
on the enlarged tube $\hat\alpha_{j-1}$ of the correlation theorem (the subscripted letter denotes this
tube, not the H\"older exponent $\hat\alpha$); here $\beta_s$ is the constant of that name in the
correlation theorem, whose value $\tfrac14$ gives the equality in \eqref{10.3:eq-typed-families-2}.
\end{enumerate}
\end{definition}

For the type space $\mathsf{V} = V = \mathbb{R}^4$ we use the following property of $V$-typed families,
taken by statement: a frozen $V$-typed family has neither a vacuum part
nor a coupling-increment part (a term of the type $\beta_{k+1}$).

The improved norm is the norm in which the outputs of the step are bounded. The output bound used in
the proof of the correlation theorem bounds the outputs at level $j = k+1$ by its constant times
$e^{-\mathsf{r}(\delta)\kappa d_{k+1}(X)}$, where $\delta$ is the parameter of that bound and
$\mathsf{r}(\delta)$ its rate factor; this rate satisfies
\[
\mathsf{r}(\delta)\kappa \ \ge\ (1+4\beta_s)\kappa = 2\kappa ,
\]
so that these outputs carry the weight \eqref{10.3:eq-typed-families-2}; in particular, the current
$\mathfrak{J}[F]$ of a frozen family is placed in the improved norm. The re-filing lemmas deliver the
re-filed remnants in the same norm.

The setting of the next lemma is the following. It is close to that of
Lemma~\ref{10.1:lem-double-subtraction}, except that there $|P|$ is measured from the centre of $P$
and here from its corner. We work on the unit lattice $\mathbb{Z}^4$ with unit vectors
$e_1,\dots,e_4$, and put $(\nabla_\nu g)(x):=g(x+e_\nu)-g(x)$.

\emph{Plaquettes and cochains.} A plaquette is a pair $P=(x;A)$ with $x\in\mathbb{Z}^4$ and
$A=\{\alpha<\beta\}\subset\{1,2,3,4\}$. We put $|P|:=|x|_2$ and $P-y:=(x-y;A)$. For a complex
$2$-cochain $f$ we write $f_A(x):=f(x;A)$, and $f_{\alpha\beta}:=f_A$ for $A=\{\alpha<\beta\}$. We also
set $f_{\beta\alpha}:=-f_{\alpha\beta}$ and $f_{\alpha\alpha}:=0$. The cochain $f$ is \emph{closed} if
\begin{equation*}
\nabla_\rho f_{\nu\lambda}-\nabla_\nu f_{\rho\lambda}+\nabla_\lambda f_{\rho\nu}=0
\qquad\text{for all }\rho,\nu,\lambda .
\end{equation*}
With the convention $f_{\alpha\alpha}=0$ this holds automatically when two of the indices coincide.

\emph{Symmetries and bonds.} The lattice symmetries $r$ are the translations and the elements of the
hyperoctahedral group $W_4$ acting about sites. They act on sites, on oriented bonds
$b_\nu(y):=\langle y,y+e_\nu\rangle$, and on cochains, $f\mapsto r\cdot f$. We write $r_*$ for the
linear part of $r$, a signed permutation matrix. The bonds $b_\nu(y)$ are called positively oriented.
Every oriented bond is either positively oriented or of the form $-b$, the positively oriented bond $b$
traversed backwards. A function $\mathfrak{F}$ of positively oriented bonds is extended to all oriented
bonds by $\mathfrak{F}(-b):=-\mathfrak{F}(b)$; this convention is used throughout. For a lattice
symmetry $r$ and a positively oriented bond $b$ we put $\operatorname{sgn}_r(b):=1$ if $rb$ is
positively oriented and $\operatorname{sgn}_r(b):=-1$ otherwise; thus $\mathfrak{F}(rb)$ is
$\operatorname{sgn}_r(b)$ times the value of $\mathfrak{F}$ at the positively oriented bond underlying
$rb$. For a family $\mathfrak{F}_\rho(b;f)$, $\rho=1,\dots,4$, of bond functions depending on cochains,
and a lattice symmetry $r$, we put
\begin{equation*}
(r\mathfrak{F})_\rho(b;f):=\sum_{\rho'}(r_*)_{\rho\rho'}\mathfrak{F}_{\rho'}(r^{-1}b;r^{-1}\cdot f),
\end{equation*}
so that the sign $\operatorname{sgn}_{r^{-1}}(b)$ enters through the convention just stated.

\emph{The data.} We are given complex numbers $k_\rho(P,P')$, $\rho=1,\dots,4$, symmetric in $(P,P')$, with
$|k_\rho(P,P')|\le\mathfrak{n}_2e^{-\delta(|P|+|P'|)}$ for some $\delta>0$. For $f$ of polynomial growth
we put
\begin{equation*}
Q_{\rho}(y;f):=\sum_{P,P'}k_{\rho}(P-y,P'-y)\,f(P)\,f(P').
\end{equation*}
The kernel is of \emph{vector type}: for every closed $f$ and every lattice symmetry $r$,
\begin{equation*}
Q_\rho(ry;r\cdot f)=\sum_{\rho'}(r_*)_{\rho\rho'}Q_{\rho'}(y;f).
\end{equation*}
Throughout, $C(\delta)$, $C(\delta',\delta'')$ and $C_Q(\delta)$ denote constants depending only on the
arguments shown.

\begin{lemma}[Divergence form of vector-typed quadratic forms]\label{10.3:lem-divergence-identity}
There is a bond current $\mathfrak{C}_{\rho}(b;f)$, quadratic in $f$ (the diagonal of a bilinear form in
$f$) and linear in $k$, of the following form: $\mathfrak{C}$ is the average of $r\mathfrak{C}^0$ over
$r\in W_4$ acting about the origin, where for every $y$ and $\nu$
\begin{equation}\label{10.3:eq-divergence-identity-1}
\begin{split}
\mathfrak{C}^0_{\rho}(b_\nu(y-e_\nu);f)={}&-2\sum_{A,B}\sum_{u\in\mathbb{Z}^4}\mu_{\nu\rho;AB}(u)\,
f_A(y)\,f_B(y+u)\\
&+\sum_{A,B}s_{\nu\rho;AB}\,f_A(y)\,f_B(y).
\end{split}
\end{equation}
Here the kernels $\mu_{\nu\rho;AB}(u)$ and the coefficients $s_{\nu\rho;AB}$ do not depend on $y$, the
kernels $\mu$ obey \eqref{10.3:eq-divergence-identity-a} below with $\delta'=\delta$,
$C_\kappa=C(\delta)\mathfrak{n}_2$ and every $0<\delta''<\delta$, and
$|s_{\nu\rho;AB}|\le C(\delta)\mathfrak{n}_2$.
Moreover, for every closed $f$ of polynomial growth and every $y\in\mathbb{Z}^4$:
\begin{enumerate}
\item[(i)] the identity
\begin{equation}\label{10.3:eq-divergence-identity-2}
Q_{\rho}(y;f)=\sum_\nu\bigl[\mathfrak{C}_{\rho}(b_\nu(y);f)-\mathfrak{C}_{\rho}(b_\nu(y-e_\nu);f)\bigr]
+\mathsf{E}_{\rho}(y;f)
\end{equation}
holds exactly;
\item[(ii)] $\mathfrak{C}$ is covariant: for every lattice symmetry $r$ and every positively oriented
bond $b$,
\begin{equation*}
\mathfrak{C}_{\rho}(rb;r\cdot f)=\sum_{\rho'}(r_*)_{\rho\rho'}\mathfrak{C}_{\rho'}(b;f);
\end{equation*}
by the convention of the setting, the left side is $\operatorname{sgn}_r(b)$ times the value of
$\mathfrak{C}_\rho(\cdot\,;r\cdot f)$ at the positively oriented bond underlying $rb$;
\item[(iii)] $\mathsf{E}_{\rho}(y;\cdot)$ is a sum of products of two differences of $f$. If
$|f_A(y+u)-f_A(y)|\le\varphi_1(1+|u|_2)^2$ for all $A$ and $u$, then
$|\mathsf{E}_{\rho}(y;f)|\le C_Q(\delta)\,\mathfrak{n}_2\,\varphi_1^2$.
\end{enumerate}
\end{lemma}

\begin{proof}
All sums below converge absolutely. The kernels decay exponentially and $f$ grows at most
polynomially, so this is immediate.

\emph{Reduction to a covariant kernel.} For a site $y$, let $r$ run over the copy of $W_4$ acting about
$y$. Replace $Q_\rho(y;f)$ by the average over $r$ of $\sum_{\rho'}((r_*)^{-1})_{\rho\rho'}Q_{\rho'}(y;r\cdot f)$.

For closed $f$ the cochain $r\cdot f$ is closed. Since $ry=y$, vector type gives
\begin{equation*}
Q_{\rho'}(y;r\cdot f)=\sum_{\rho''}(r_*)_{\rho'\rho''}Q_{\rho''}(y;f).
\end{equation*}
So each term of the average equals $Q_\rho(y;f)$, and averaging does not change $Q$ on closed
cochains.

The averaged form is again $\sum\bar k_\rho(P-y,P'-y)f(P)f(P')$, with a kernel $\bar k$ that is
symmetric and linear in $k$. It is a mean of the kernels obtained from $k$ by transporting both
plaquette arguments with an element of $W_4$ and the index $\rho$ with $(r_*)^{-1}$. An element of
$W_4$ maps a plaquette based at $x$ to one based at $r_*x+v$ with $|v|_2\le\sqrt2$. The entries of
$r_*$ are $0,\pm1$. Hence $\bar k$ satisfies the bound of the setting with $\mathfrak{n}_2$ replaced by
$e^{2\sqrt2\delta}\mathfrak{n}_2$, a factor that is absorbed into the constants $C(\delta)$ below. After
reindexing the average, $\bar k$ is covariant on all cochains. We rename $\bar k$ to $k$ and assume
from now on that $k$ is covariant on all cochains.

\emph{Step 1.} Let $r$ be the inversion through the site $y$. It fixes $y$, fixes constant cochains and
reverses vectors, so $r_*=-\mathbf{1}$. For a constant cochain $c$, covariance gives
\begin{equation*}
Q_\rho(y;c)=Q_\rho(ry;r\cdot c)=-Q_\rho(y;c),
\end{equation*}
so $Q_\rho(y;c)=0$. Now put
\begin{equation*}
K_{AB}:=\sum_{x,u\in\mathbb{Z}^4}k_\rho((x;A),(u;B)).
\end{equation*}
Then $Q_\rho(y;c)=\sum_{A,B}c_Ac_BK_{AB}$, and $K_{AB}=K_{BA}$ because $k$ is symmetric. This quadratic
form vanishes for all complex $c$, so by polarisation $K_{AB}=0$ for all $A,B$.

\emph{Step 2.} Put
\begin{equation*}
\kappa_{\rho;AB}(u):=\sum_{x}k_\rho((x;A),(u;B)).
\end{equation*}
Then $\sum_u\kappa_{\rho;AB}(u)=K_{AB}=0$, and
$|\kappa_{\rho;AB}(u)|\le C(\delta)\mathfrak{n}_2e^{-\delta|u|_2}$. Reindex $P\mapsto P+y$ and put
$\Delta_A(x):=f_A(y+x)-f_A(y)$, which is $f(P)-f_{A(P)}(y)$ for $P=(y+x;A)$. Then
\begin{equation*}
Q_\rho(y;f)=\sum_{x,u,A,B}k_\rho((x;A),(u;B))\bigl[f_A(y)+\Delta_A(x)\bigr]\bigl[f_B(y)+\Delta_B(u)\bigr].
\end{equation*}
We treat the three kinds of terms separately.
\begin{itemize}
\item The term in $f_A(y)f_B(y)$ is $\sum_{A,B}K_{AB}f_A(y)f_B(y)=0$.
\item By the symmetry of $k$ the two cross terms are equal. Their sum is
\begin{equation*}
2\sum_{A,B}f_A(y)\sum_u\kappa_{\rho;AB}(u)\bigl[f_B(y+u)-f_B(y)\bigr].
\end{equation*}
The part with $f_B(y)$ carries the factor $K_{AB}=0$.
\end{itemize}
Hence $Q_\rho(y;f)=T_1+T_2$, where
\begin{align*}
T_1&:=\sum_{x,u,A,B}k_\rho((x;A),(u;B))\,\Delta_A(x)\,\Delta_B(u),\\
T_2&:=2\sum_{A,B}f_A(y)\sum_u\kappa_{\rho;AB}(u)\,f_B(y+u).
\end{align*}

\emph{Step 3 (primitive).} Let $\kappa$ be a function on $\mathbb{Z}^4$ of exponential decay with
$\sum_u\kappa(u)=0$. We show that $\kappa(u)=\sum_\nu[\mu_\nu(u)-\mu_\nu(u-e_\nu)]$ for suitable
$\mu_\nu$.

Put $\kappa^{(0)}:=\kappa$ and, for $j=1,\dots,4$,
\begin{equation*}
\kappa^{(j)}(u_{j+1},\dots,u_4):=\sum_{t\in\mathbb{Z}}\kappa^{(j-1)}(t,u_{j+1},\dots,u_4).
\end{equation*}
In particular $\kappa^{(4)}=\sum_u\kappa(u)=0$. With $\delta_{t,0}$ the Kronecker symbol, put
\begin{align*}
\mu_j(u)&:=\delta_{u_1,0}\cdots\delta_{u_{j-1},0}\sum_{t\le u_j}
\bigl[\kappa^{(j-1)}(t,u_{j+1},\dots,u_4)-\delta_{t,0}\,\kappa^{(j)}(u_{j+1},\dots,u_4)\bigr]\\
&\phantom{:}=-\,\delta_{u_1,0}\cdots\delta_{u_{j-1},0}\sum_{t>u_j}
\bigl[\kappa^{(j-1)}(t,u_{j+1},\dots,u_4)-\delta_{t,0}\,\kappa^{(j)}(u_{j+1},\dots,u_4)\bigr].
\end{align*}
The two expressions agree because the bracket has zero sum along each line $t\in\mathbb{Z}$, by the
definition of $\kappa^{(j)}$. For $j=1$ the prefactor is empty, and $\mu_1$ is the primitive of $\kappa$
along the lines parallel to $e_1$; the construction is then continued with $\kappa^{(1)}$ on
$\{u_1=0\}$, with $\kappa^{(2)}$ on $\{u_1=u_2=0\}$, and so on.

Replacing $u$ by $u-e_j$ changes only $u_j$. Hence
\begin{equation*}
\mu_j(u)-\mu_j(u-e_j)=\delta_{u_1,0}\cdots\delta_{u_{j-1},0}
\bigl[\kappa^{(j-1)}(u_j,\dots,u_4)-\delta_{u_j,0}\,\kappa^{(j)}(u_{j+1},\dots,u_4)\bigr].
\end{equation*}
The sum over $j$ telescopes to $\kappa(u)-\delta_{u_1,0}\cdots\delta_{u_4,0}\kappa^{(4)}=\kappa(u)$.

We now bound the $\mu_j$. We claim that
\begin{equation}\label{10.3:eq-divergence-identity-a}
\begin{split}
&\text{if }|\kappa(u)|\le C_\kappa e^{-\delta'|u|_2}\text{ for all }u\in\mathbb{Z}^4,
\text{ then for every }0<\delta''<\delta':\\
&|\mu_\nu(u)|\le C(\delta',\delta'')\,C_\kappa\,e^{-\delta''|u|_2}
\qquad(u\in\mathbb{Z}^4,\ \nu=1,\dots,4).
\end{split}
\end{equation}
To see this, consider first $j=1$.
\begin{itemize}
\item For $u_1\ge0$ use the second expression. There $t>u_1\ge0$, so the term with $\delta_{t,0}$
never occurs.
\item For $u_1<0$ use the first expression. There $t\le u_1<0$, so again the term with $\delta_{t,0}$
never occurs.
\end{itemize}
In both cases only points $(t,u_2,u_3,u_4)$ with $|t|\ge|u_1|$ occur, and every such point has
Euclidean norm $\ge|u|_2$. The ray sum satisfies
\begin{equation*}
\sum_{|t|\ge|u_1|}e^{-\delta'|(t,u_2,u_3,u_4)|_2}\le C(\delta')(1+|u|_2)^{1/2}e^{-\delta'|u|_2}.
\end{equation*}
Together with the hypothesis on $\kappa$ this gives
\begin{equation*}
|\mu_1(u)|\le C(\delta')\,C_\kappa\,(1+|u|_2)^{1/2}e^{-\delta'|u|_2},
\end{equation*}
and, taking $u_1=0$, the same bound holds for $\kappa^{(1)}$ on $\mathbb{Z}^3$, which again has zero
sum. Given $0<\delta''<\delta'$, fix rates
$\delta'=\delta'_0>\delta'_1>\delta'_2>\delta'_3>\delta''$.
The factor $(1+|u|_2)^{1/2}$ is absorbed by passing from $\delta'_0$ to $\delta'_1$. Repeat the argument
with $\kappa^{(1)}$ and $\delta'_1$ in place of $\kappa$ and $\delta'$. Since $\mu_2$ vanishes unless
$u_1=0$, in which case $|u|_2=|(u_2,u_3,u_4)|_2$, this gives the bounds for $\mu_2$ and $\kappa^{(2)}$.
After four stages we obtain \eqref{10.3:eq-divergence-identity-a}.

Apply this to each $\kappa_{\rho;AB}$, with $\delta'=\delta$ and $C_\kappa=C(\delta)\mathfrak{n}_2$. This
yields kernels $\mu_{\nu\rho;AB}(u)$, independent of $y$ and linear in $k$, with
\begin{equation*}
\kappa_{\rho;AB}(u)=\sum_\nu\bigl[\mu_{\nu\rho;AB}(u)-\mu_{\nu\rho;AB}(u-e_\nu)\bigr].
\end{equation*}
Put
\begin{equation*}
M_{\nu\rho;A}(y):=\sum_{B,u}\mu_{\nu\rho;AB}(u)f_B(y+u).
\end{equation*}
Substituting this decomposition and shifting
$u\mapsto u+e_\nu$ in the terms with $\mu_{\nu\rho;AB}(u-e_\nu)$ gives
\begin{equation*}
\sum_{B,u}\kappa_{\rho;AB}(u)f_B(y+u)
=\sum_\nu\sum_{B,u}\mu_{\nu\rho;AB}(u)\bigl[f_B(y+u)-f_B(y+u+e_\nu)\bigr]
=-\sum_\nu(\nabla_\nu M_{\nu\rho;A})(y).
\end{equation*}
Hence $T_2=-2\sum_{\nu,A}f_A(y)(\nabla_\nu M_{\nu\rho;A})(y)$.

Next put $\hat\mu^0_{\nu\rho;AB}:=\sum_u\mu_{\nu\rho;AB}(u)$. By the same shift,
\begin{equation*}
\sum_uu_\nu\kappa_{\rho;AB}(u)=\sum_{\nu'}\sum_u\mu_{\nu'\rho;AB}(u)\bigl[u_\nu-(u+e_{\nu'})_\nu\bigr]
=-\hat\mu^0_{\nu\rho;AB}.
\end{equation*}
Thus $\hat\mu^0_{\nu\rho;AB}=-\sum_uu_\nu\kappa_{\rho;AB}(u)$. Replacing the position $u$ of the
plaquettes by $u+a$ changes this by $-a_\nu\sum_u\kappa_{\rho;AB}(u)=0$. So $\hat\mu^0$ does not depend
on the origins chosen for the plaquette positions. Since $k$ is covariant, $\hat\mu^0$ is therefore a
$W_4$-invariant constant tensor. Moreover $|\hat\mu^0_{\nu\rho;AB}|\le C(\delta)\mathfrak{n}_2$ by the
decay of $\kappa$.

\emph{Step 4 (Leibniz).} By the identity $g\nabla_\nu h=\nabla_\nu(gh)-(\nabla_\nu g)\,h(\cdot+e_\nu)$,
\begin{equation*}
T_2=-2\sum_{\nu,A}\nabla_\nu(f_AM_{\nu\rho;A})(y)+Y_\rho+T_3,
\end{equation*}
where
\begin{align*}
Y_\rho&:=2\sum_{\nu,A,B}\hat\mu^0_{\nu\rho;AB}(\nabla_\nu f_A)(y)f_B(y),\\
T_3&:=2\sum_{\nu,A}(\nabla_\nu f_A)(y)
\Bigl[M_{\nu\rho;A}(y+e_\nu)-\sum_B\hat\mu^0_{\nu\rho;AB}f_B(y)\Bigr].
\end{align*}
The bracket equals $\sum_{B,u}\mu_{\nu\rho;AB}(u)[f_B(y+u+e_\nu)-f_B(y)]$.

\emph{Step 5 (Y is a divergence plus differences).} Put
$\hat\mu^0_{\nu\rho;(AB)}:=\tfrac12(\hat\mu^0_{\nu\rho;AB}+\hat\mu^0_{\nu\rho;BA})$, and write
$\hat\mu^0_{\nu\rho;[AB]}$ for the antisymmetric part. Split $Y_\rho=Y^s_\rho+Y^a_\rho$ accordingly.

For the symmetric part we use $\nabla(gh)=(\nabla g)h+g\nabla h+(\nabla g)(\nabla h)$:
\begin{equation*}
Y^s_\rho=\sum_{\nu,A,B}\hat\mu^0_{\nu\rho;(AB)}
\bigl[\nabla_\nu(f_Af_B)-(\nabla_\nu f_A)(\nabla_\nu f_B)\bigr](y).
\end{equation*}

For the antisymmetric part, which is again $W_4$-invariant, we use the parity rule for hyperoctahedral
tensors (Lemma~\ref{10.1:lem-parity-rule}). It gives two cases.
\begin{itemize}
\item If $\nu=\rho$, an entry can be non-zero only if $A=B$, where the antisymmetric part vanishes.
\item If $\nu\ne\rho$, an entry can be non-zero only if $\{A,B\}=\{\{\nu,\lambda\},\{\rho,\lambda\}\}$
with $\lambda\notin\{\nu,\rho\}$.
\end{itemize}
The permutations of the axes, which lie in $W_4$, act transitively on ordered triples of distinct
directions. Hence there is one constant $c$ such that, written in the oriented components,
\begin{equation*}
Y^a_\rho=c\sum_{\nu,\lambda}\bigl[(\nabla_\nu f_{\nu\lambda})f_{\rho\lambda}
-(\nabla_\nu f_{\rho\lambda})f_{\nu\lambda}\bigr](y).
\end{equation*}
Up to sign, $c$ is twice an entry of $\hat\mu^0_{\nu\rho;[AB]}$. The terms with $\nu=\rho$ or
$\lambda\in\{\nu,\rho\}$ vanish because $f_{\alpha\alpha}=0$, so the sum may be taken unrestricted. By
the product rule $\nabla(gh)=(\nabla g)h+g\nabla h+(\nabla g)(\nabla h)$ again,
\begin{equation*}
Y^a_\rho=c\sum_{\nu,\lambda}\bigl[\nabla_\nu(f_{\nu\lambda}f_{\rho\lambda})
-2f_{\nu\lambda}\nabla_\nu f_{\rho\lambda}-(\nabla_\nu f_{\nu\lambda})(\nabla_\nu f_{\rho\lambda})\bigr].
\end{equation*}

It remains to treat the middle term. Exchange the names $\nu,\lambda$ and use
$f_{\lambda\nu}=-f_{\nu\lambda}$; this gives
$\sum f_{\nu\lambda}\nabla_\nu f_{\rho\lambda}=-\sum f_{\nu\lambda}\nabla_\lambda f_{\rho\nu}$.
Closedness gives $\nabla_\nu f_{\rho\lambda}-\nabla_\lambda f_{\rho\nu}=\nabla_\rho f_{\nu\lambda}$.
Finally $\nabla(g^2)=2g\nabla g+(\nabla g)^2$. Together,
\begin{align*}
\sum_{\nu,\lambda}f_{\nu\lambda}\nabla_\nu f_{\rho\lambda}
&=\tfrac12\sum_{\nu,\lambda}f_{\nu\lambda}\bigl(\nabla_\nu f_{\rho\lambda}
-\nabla_\lambda f_{\rho\nu}\bigr)
=\tfrac12\sum_{\nu,\lambda}f_{\nu\lambda}\nabla_\rho f_{\nu\lambda}\\
&=\tfrac14\sum_{\nu,\lambda}\bigl[\nabla_\rho(f_{\nu\lambda}^2)-(\nabla_\rho f_{\nu\lambda})^2\bigr].
\end{align*}
We write $\nabla_\rho$ of a function as $\sum_\nu\delta_{\nu\rho}\nabla_\nu$ of it.

\emph{Step 6.} We collect the results of Steps 2, 4 and 5. Put
\begin{equation*}
\mathfrak{C}^0_\rho(b_\nu(y-e_\nu);f):=-2\sum_Af_A(y)M_{\nu\rho;A}(y)
+\sum_{A,B}s_{\nu\rho;AB}f_A(y)f_B(y),
\end{equation*}
where the coefficients $s$ are defined by
\begin{equation*}
\sum_{A,B}s_{\nu\rho;AB}f_Af_B:=\sum_{A,B}\hat\mu^0_{\nu\rho;(AB)}f_Af_B
+c\Bigl[\sum_\lambda f_{\nu\lambda}f_{\rho\lambda}
-\tfrac12\delta_{\nu\rho}\sum_{\sigma,\lambda}f_{\sigma\lambda}^2\Bigr].
\end{equation*}
Put also
\begin{align*}
\mathsf{E}^0_\rho(y;f):={}&T_1+T_3-\sum_{\nu,A,B}\hat\mu^0_{\nu\rho;(AB)}(\nabla_\nu f_A)(\nabla_\nu f_B)\\
&+c\Bigl[\tfrac12\sum_{\nu,\lambda}(\nabla_\rho f_{\nu\lambda})^2
-\sum_{\nu,\lambda}(\nabla_\nu f_{\nu\lambda})(\nabla_\nu f_{\rho\lambda})\Bigr].
\end{align*}
The current $\mathfrak{C}^0$ has the form \eqref{10.3:eq-divergence-identity-1}, and
$|s|\le C(\delta)\mathfrak{n}_2$ because $|c|\le2\max|\hat\mu^0|$. Since
$b_\nu(y)=b_\nu((y+e_\nu)-e_\nu)$, with $G_\nu(x):=\mathfrak{C}^0_\rho(b_\nu(x-e_\nu);f)$ we have
\begin{equation*}
\sum_\nu\bigl[\mathfrak{C}^0_\rho(b_\nu(y);f)-\mathfrak{C}^0_\rho(b_\nu(y-e_\nu);f)\bigr]
=\sum_\nu(\nabla_\nu G_\nu)(y).
\end{equation*}
So Steps 2, 4 and 5 give
\eqref{10.3:eq-divergence-identity-2} for $\mathfrak{C}^0$ and $\mathsf{E}^0$.

Each term of $\mathsf{E}^0$ is a product of two differences, of the form $f(y+u)-f(y)$ or
$f(y+u+e_\nu)-f(y)$. Under the hypothesis of (iii) we have the following bounds:
\begin{itemize}
\item $|\nabla_\nu f_A(y)|\le4\varphi_1$;
\item $|\Delta_A(x)|\le\varphi_1(1+|x|_2)^2$;
\item $|f_B(y+u+e_\nu)-f_B(y)|\le\varphi_1(2+|u|_2)^2$.
\end{itemize}
The coefficients of these products are summable against such weights:
\begin{align*}
\sum|k_\rho|(1+|x|_2)^2(1+|u|_2)^2&\le C(\delta)\mathfrak{n}_2,\\
\sum_u|\mu_{\nu\rho;AB}(u)|(2+|u|_2)^2&\le C(\delta)\mathfrak{n}_2,
\end{align*}
the second by \eqref{10.3:eq-divergence-identity-a} with $\delta''=\delta/2$.
Also $|\hat\mu^0|,|c|\le C(\delta)\mathfrak{n}_2$.
Hence $|\mathsf{E}^0_\rho(y;f)|\le C_Q(\delta)\mathfrak{n}_2\varphi_1^2$.

\emph{Covariance.} With the convention of the setting for reversed bonds, for a bond function
$\mathfrak{F}$ the sum $\sum_\nu[\mathfrak{F}(b_\nu(y))-\mathfrak{F}(b_\nu(y-e_\nu))]$ is the sum of
$\mathfrak{F}$ over the eight
oriented bonds starting at $y$. A symmetry $r$ maps the oriented bonds starting at $r^{-1}y$ bijectively
onto those starting at $y$, so this divergence is equivariant. For a symmetry $r$ let $r\mathfrak{C}^0$
be defined as in the setting, and put
\begin{equation*}
(r\mathsf{E}^0)_\rho(y;f):=\sum_{\rho'}(r_*)_{\rho\rho'}\mathsf{E}^0_{\rho'}(r^{-1}y;r^{-1}\cdot f).
\end{equation*}
Apply \eqref{10.3:eq-divergence-identity-2} for $\mathfrak{C}^0$ at $r^{-1}y$ to the closed cochain
$r^{-1}\cdot f$, multiply by $(r_*)_{\rho\rho'}$ and sum over $\rho'$. By vector type and by the
equivariance of the divergence, this gives $Q=\operatorname{div}(r\mathfrak{C}^0)+r\mathsf{E}^0$.

The current $\mathfrak{C}^0$ is translation covariant by construction: its value at $b_\nu(y-e_\nu)$ depends
on $y$ only through $f(\cdot+y)$. Let $\mathfrak{C}$ and $\mathsf{E}$ be the averages of $r\mathfrak{C}^0$ and
$r\mathsf{E}^0$ over $r\in W_4$ acting about the origin. We check the three properties.
\begin{itemize}
\item Property (i) holds because \eqref{10.3:eq-divergence-identity-2} is linear in the pair
$(\mathfrak{C},\mathsf{E})$.
\item Property (ii) for $W_4$ follows by reindexing the average, using $(rs)_*=r_*s_*$. Translation
covariance persists because $W_4$ conjugates translations into translations. These two give
covariance under every lattice symmetry.
\item For (iii), an element of $W_4$ moves the base point of a plaquette by at most $\sqrt2$ from the
image of the original base point. So each difference of $r^{-1}\cdot f$ between the sites $r^{-1}y+u$
and $r^{-1}y$ is, up to sign, a difference $f_{A'}(y+r_*u+v)-f_{A'}(y+v)$ with $|v|_2\le\sqrt2$ and
$|r_*u|_2=|u|_2$. It equals
\begin{equation*}
\bigl[f_{A'}(y+r_*u+v)-f_{A'}(y)\bigr]-\bigl[f_{A'}(y+v)-f_{A'}(y)\bigr],
\end{equation*}
so the bound for $\mathsf{E}^0$ gives (iii) for $\mathsf{E}$ with a new $C_Q(\delta)$.
\end{itemize}
\end{proof}

\begin{remark}
Three features of the lemma matter in its later use. First, the identity is exact on closed cochains,
with an explicit remainder. Second, the current lives on bonds: a reflection maps $b_\nu(y-e_\nu)$ to
$b_\nu(y)$, and this is what makes the current covariant. Third, the hypothesis of (iii) is needed about
the one point $y$ only.

No weight function occurs in the lemma; weights enter only where the identity is applied.

On constant $f$ one has $\mathfrak{C}_{\rho}(b_\nu(y);f)=\sum_{A,B}\mathsf{T}_{\nu\rho;AB}f_Af_B$ for
every $y$, where
$\mathsf{T}_{\nu\rho;AB}:=-2\hat\mu^0_{\nu\rho;AB}+s_{\nu\rho;AB}$ is a $W_4$-invariant tensor. Its
contraction with an antisymmetric $\omega_{\nu\rho}$ vanishes, because every $W_4$-equivariant linear
map from $\Lambda^2\mathbb{R}^4$ to $\operatorname{Sym}^2\Lambda^2\mathbb{R}^4$ is zero.
This vanishing is used where the identity is applied, not in the proof above.
\end{remark}

\begin{definition}[The second-variation kernel and the comb gauge]\label{10.3:def-kernel-comb-gauge}
In this section the following standing assumptions are in force:
\begin{itemize}
\item $j\ge 1$;
\item $F$ is a $V$-typed pointed E-family at scale $j$
(Definition~\ref{10.3:def-typed-families} with $\mathsf{V}=V$) on radii $\alpha^{\circ}_j$, where
$\alpha^{\circ}_j=\hat\alpha_{j-1}$ at creation, $\hat\alpha_{j-1}$ being the radii of the
enlarged tube of Definition~\ref{10.3:def-typed-families};
\item its fresh norm satisfies $\|F\|^{\mathrm{fr}}\le\mathfrak{n}$.
\end{itemize}
In the correlation theorem the radii of a family are denoted $a_j$; here $a_j$ denotes the
spacing, so the radii are written $\alpha^{\circ}_j$.
Lengths are measured in the lattice unit of scale $j$, and $\square$ denotes the cube of
Ba{\l}aban's partition $\pi_j$ into cubes at scale $j$ that contains $y$.

\smallskip
\noindent\textit{The second-variation kernel.} For $X\ni y$ and $e\in V_{\mathbb{C}}$, let $B$
be a bond field with values in the traceless Hermitian $N\times N$ matrices, so that $\exp(iB)$
is an $SU(N)$ configuration, and write $B_{\mu}(x)=\sum_a B^a_{\mu}(x)\,t_a$ in a basis $(t_a)$ of
these matrices orthonormal for the positive-definite form $(M,M')\mapsto\operatorname{tr}(MM')$,
that is $\operatorname{tr}(t_at_b)=\delta_{ab}$, with the trace normalised by
$\operatorname{tr}\mathbf{1}=1$. Let $U_j(\cdot)$,
$J_j(\cdot)$ be the two field arguments on which a family of the correlation theorem depends,
here evaluated at the configuration $\exp(iB)$, with $U_j$ the free block field at level $j$ on
$L^{-j}\mathbb{Z}^4$ of the proof recalled below. By \cite[(4.33)]{B-CMP109} the second
derivative at $B=0$ has colour structure $\delta_{ab}$:
\begin{equation}\label{10.3:eq-kernel-comb-gauge-1}
\frac{\partial^2}{\partial B^a_{\mu}(x)\,\partial B^b_{\nu'}(x')}
F\bigl(X,U_j(\exp(iB)),J_j(\exp(iB)),y;e\bigr)\bigg|_{B=0}
=\delta_{ab}\,\Pi^X_{\mu\nu'}(x,x',y;e),
\end{equation}
and \eqref{10.3:eq-kernel-comb-gauge-1} defines the scalar kernel
$\Pi^X_{\mu\nu'}(x,x',y;e)$. Set
\begin{equation}\label{10.3:eq-kernel-comb-gauge-2}
\Pi:=\sum_{X\ni y}\Pi^X,
\end{equation}
the sum running over all $X\ni y$ of $L^{-j}\mathbb{Z}^4$.

This is the kernel that occurs in the proof of the lemma, among those on which the correlation
theorem rests, that is identified by the following step of its proof.
There, $U_j(\square_0,\cdot)$ is replaced by the free $U_j$, and the sum over $X$ is extended to
all $X\ni y$ of $L^{-j}\mathbb{Z}^4$.

\smallskip
\noindent\textit{The comb gauge.} Order the coordinate axes as $e_1,e_2,e_3,e_4$. For a lattice
point $z$, the \emph{comb path} from $y$ to $z$ runs from $y$ parallel to $e_1$ until its first
coordinate equals $z_1$, then parallel to $e_2$, $e_3$ and $e_4$ in turn, each leg ending when the
corresponding coordinate equals that of $z$. A comb path is coordinate-wise monotone, that is,
along it each coordinate varies monotonically from its value at $y$ to its value at $z$. For a
bond from $x$ to $x+e_{\mu}$, the comb paths to $x$ and to $x+e_{\mu}$ agree along their first
$\mu-1$ legs; their $\mu$-th legs lie on one line and end at two points differing by $e_{\mu}$;
after the $\mu$-th leg the comb path to $x+e_{\mu}$ is the translate by $e_{\mu}$ of the comb
path to $x$. The \emph{comb gauge centred at $y$} is the map $\mathsf{K}_y$
from plaquette cochains to bond cochains of the lattice of scale $j$, which has unit spacing in
the present units, defined by
\begin{equation}\label{10.3:eq-kernel-comb-gauge-3}
(\mathsf{K}_y f)_{\mu}(x):=\sum_{P\in\Sigma(y;x,\mu)}\pm f(P).
\end{equation}
Here $\Sigma(y;x,\mu)$ is the set of plaquettes swept between the comb path from $y$ to $x$ and
the comb path from $y$ to $x+e_{\mu}$, namely the plaquettes spanned by a bond of the comb path
to $x$ after its $\mu$-th leg and the translate by $e_{\mu}$ of that bond. The swept ribbon is
oriented so that its boundary is the comb path from $y$ to $x$, followed by the bond from $x$ to
$x+e_{\mu}$, followed by the comb path from $y$ to $x+e_{\mu}$ traversed backwards; each
plaquette is counted with the sign $+$ if its orientation agrees with that of the ribbon and
with the sign $-$ otherwise. The set
$\Sigma(y;x,\mu)$ has at most $4\,(1+|x-y|_1)$ elements. As comb paths are coordinate-wise
monotone, every $P\in\Sigma(y;x,\mu)$ satisfies
\begin{equation}\label{10.3:eq-kernel-comb-gauge-4}
|P-y|_2\le|x-y|_2+1 ,
\end{equation}
where $|P-y|_2$ denotes the largest Euclidean distance from $y$ to a vertex of $P$.
For closed $f$, that is $\mathrm{d}f=0$, one has $\mathrm{d}\mathsf{K}_y f=f$; here
$\mathrm{d}$ denotes the lattice coboundary on cochains.
\end{definition}

\begin{proof}[Proof of \eqref{10.3:eq-kernel-comb-gauge-4}]
Let $z$ be a lattice point, and let $\gamma$ be a coordinate-wise monotone lattice path from $y$
to $z$. Along $\gamma$ each coordinate moves monotonically from its value at $y$ to its value
at $z$. Hence every point $w$ of $\gamma$ has $|w_i-y_i|\le|z_i-y_i|$ for $i=1,\dots,4$, and
therefore $|w-y|_2\le|z-y|_2$.

We apply this with $z=x$, which gives $|w-y|_2\le|x-y|_2$ on the comb path to $x$. We apply it
again with $z=x+e_{\mu}$. By the triangle inequality,
$|x+e_{\mu}-y|_2\le|x-y|_2+1$, so $|w-y|_2\le|x-y|_2+1$ on the comb path to $x+e_{\mu}$.

Each $P\in\Sigma(y;x,\mu)$ is spanned by a bond of the comb path to $x$ and the translate of that
bond by $e_{\mu}$, which is a bond of the comb path to $x+e_{\mu}$. So two vertices of $P$ lie on
the comb path to $x$ and two on the comb path to $x+e_{\mu}$, and the two bounds just proved give
\eqref{10.3:eq-kernel-comb-gauge-4}.
\end{proof}

\begin{lemma}[Decay of the second-variation kernel]\label{10.3:lem-kernel-decay}
Let $j\ge 1$, let $F$ be a $\mathsf{V}$-typed pointed E-family at scale $j$ with
$\|F\|^{\mathrm{fr}}\le\mathfrak{n}$, let $\Pi^X$, $\Pi=\sum_{X\ni y}\Pi^X$ and the comb gauge
$\mathsf{K}=\mathsf{K}_y$ be as in Definition~\ref{10.3:def-kernel-comb-gauge}, and measure lengths
in units of the scale-$j$ spacing. Then the following hold.
\begin{itemize}
\item[(K1)] For every $X$ the kernel $\Pi^X$ obeys the first bound in \eqref{10.3:eq-kernel-decay-a}.
\item[(K2)] $\Pi(d\lambda,\cdot)=0$, that is, the contraction of $\Pi$ in its first argument with
the gradient $d\lambda$ of any $\lambda$ vanishes. Hence, for every closed $f$ of polynomial growth,
the quantity
\begin{equation}\label{10.3:eq-kernel-decay-1}
e^{\rho}Q[F]_{\rho}(y;f):=\frac12\sum_{x,x'}\Pi(x,x',y;e)\,B(x)\,B(x')
\end{equation}
(with $e^{\rho}$ the components of $e$ and summation over $\rho$)
does not depend on the $B$ with $dB=f$; it is a quadratic form of the kind treated in
Lemma~\ref{10.3:lem-divergence-identity}, of vector type, with kernel
$k:=\tfrac12\mathsf{K}^{T}\Pi\,\mathsf{K}$, and $k$ obeys the second bound in
\eqref{10.3:eq-kernel-decay-a}.
\item[(K3)] The kernel $\kappa_{\rho;AB}$ of Lemma~\ref{10.3:lem-divergence-identity} (its letter
is not to be confused with the rate $\kappa$) obeys the third bound in
\eqref{10.3:eq-kernel-decay-a};
hence, by \eqref{10.3:eq-divergence-identity-a}, the kernels $\mu$ and $s$ of the bond current of
$Q[F]$ obey the last two bounds in \eqref{10.3:eq-kernel-decay-a}.
\end{itemize}
\begin{equation}\label{10.3:eq-kernel-decay-a}
\begin{aligned}
&\text{(K1)}\quad |\Pi^X(x,x',y;e)|\le 4B_3^2\,\mathfrak{n}\,|e|\,\mathfrak{r}_j^{-2}\\
&\qquad\qquad\times\exp\bigl[-2\kappa d_j(X)
  -\delta_0\bigl(\operatorname{dist}(x,X)+\operatorname{dist}(x',X)\bigr)\bigr],\\
&\text{(K2)}\quad |k_{\rho}(P,P')|\le C_K\,\mathfrak{n}\,\mathfrak{r}_j^{-2}\,
  e^{-(\kappa/2M)(|P|+|P'|)},\\
&\text{(K3)}\quad |\kappa_{\rho;AB}(u)|\le C_K\,\mathfrak{n}\,\mathfrak{r}_j^{-2}\,
  e^{-(39/20)\kappa|u|_2/M},\\
&\qquad\quad\ \ |\mu_{\nu\rho;AB}(u)|\le C_K'\,\mathfrak{n}\,\mathfrak{r}_j^{-2}\,
  e^{-(77/40)\kappa|u|_2/M},\qquad |s|\le C_K'\,\mathfrak{n}\,\mathfrak{r}_j^{-2}.
\end{aligned}
\end{equation}
The constants $C_K$, $C_K'$ depend only on $d$, $M$, $\kappa$, $\delta$ and $B_3$.
\end{lemma}

\begin{proof}
Throughout we place $y$ at the origin, so that $|x|_2$, $|P|$ and $|u|_2$ are distances from $y$.

\emph{(K1).} This is the display in the proof of the bound on second derivatives of families
established for the correlation theorem, read at the fresh rate. That display states three things.
First, the second derivative at $B=0$ is a second differential of $F(X;\cdot)$, taken at the
point $(1,0)$ that corresponds to $B=0$:
$\partial^2_{B(x)B(x')}\big|_0=D^2F(X;1,0)[\mathsf{q}_x,\mathsf{q}_{x'}]$, where $\mathsf{q}_x$
denotes the variation induced by the bond variable $B(x)$. Second, $\mathsf{q}_x$ and its
lattice differences are bounded by $B_3e^{-\delta_0\operatorname{dist}(x,X)}$; we write
$\|\mathsf{q}_x\|$ for the maximum of these quantities. Third,
$|D^2F[\mathsf{q},\mathsf{q}']|\le 4\mathfrak{n}e^{-\kappa d_j(X)}\|\mathsf{q}\|\,\|\mathsf{q}'\|\,
\mathfrak{r}_j^{-2}$. The fresh norm carries the weight $e^{2\kappa d_j(X)}$, so under
$\|F\|^{\mathrm{fr}}\le\mathfrak{n}$ the factor $e^{-\kappa d_j(X)}$ is replaced by
$e^{-2\kappa d_j(X)}$. Inserting the bounds on $\mathsf{q}_x$ and $\mathsf{q}_{x'}$ then gives (K1).

\emph{(K2), gauge invariance.} Let $\Phi(B)$ be the pointed total at $e^{iB}$, namely the sum over
$X\ni y$ of $F(X,U_j(e^{iB}),J_j(e^{iB}),y;e)$, so that $\Pi$ is its second derivative at $B=0$.
By the gauge invariance of pointed E-families, property (f2) of their definition, $\Phi$ is
invariant under the gauge transformation
\[
e^{iB}\longmapsto (e^{iB})^{e^{i\varepsilon\lambda}}
=\exp i\bigl(B+\varepsilon\,d\lambda+O(\varepsilon|B|)\bigr).
\]
Moreover $D\Phi(0)=0$ by \cite[(4.14)]{B-CMP109}.

Write the first-order term in $\varepsilon$ of the exponent as $d\lambda+R_\lambda(B)$ with
$R_\lambda(B)=O(|B|)$. Differentiating the invariance once in $\varepsilon$ at $\varepsilon=0$ gives
$D\Phi(B)[d\lambda+R_\lambda(B)]=0$. Differentiating this identity once in $B$ at $B=0$, in a
direction $B'$, gives
\[
D^2\Phi(0)[d\lambda,B']+D\Phi(0)\bigl[DR_\lambda(0)B'\bigr]=0,
\]
because $R_\lambda(0)=0$. The second term vanishes since $D\Phi(0)=0$. Hence
$D^2\Phi(0)[d\lambda,B']=0$ for every $B'$, which is $\Pi(d\lambda,\cdot)=0$.

Now let $dB=dB'=f$. Then $B'-B$ is closed. Summing it along lattice paths from $y$ defines a
$\lambda$ with $B'=B+d\lambda$; the sum is independent of the path because $B'-B$ is closed. The
kernel $\Pi$ is symmetric, being a second derivative. Expanding \eqref{10.3:eq-kernel-decay-1}, the
difference between its values at $B+d\lambda$ and at $B$ is
$\Pi(d\lambda,B)+\tfrac12\Pi(d\lambda,d\lambda)$, and this vanishes. The double sums converge
absolutely for $B$ of polynomial growth, by the bound \eqref{10.3:eq-kernel-decay-2} below.

Since $d\mathsf{K}_yf=f$ for closed $f$, we may choose $B=\mathsf{K}_yf$. Then
\eqref{10.3:eq-kernel-decay-1} becomes $\sum_{P,P'}k(P,P')f(P)f(P')$ with
$k=\tfrac12\mathsf{K}^T\Pi\,\mathsf{K}$, which is the form treated in
Lemma~\ref{10.3:lem-divergence-identity}. It is of vector type by the typed covariance
$(\mathrm{f}3)^{\mathsf{V}}$ of $F$.

\emph{(K2), decay.} Let $\operatorname{St}(y,x,x')$ denote the Euclidean Steiner length of the
three points, that is, the infimum of the lengths of connected sets containing them. Since such a set
contains a curve from $y$ to $x'$, and also curves from $y$ to $x$ and from $y$ to $x'$, we have
$\operatorname{St}\ge|x'-y|_2$ and $\operatorname{St}\ge\tfrac12(|x-y|_2+|x'-y|_2)$.

Take a tree contained in $X$ and meeting all its cubes, as in the tree description of $d_j(X)$ in
\cite{B-CMP109}. Prolonged to $x$ and to $x'$, it yields a connected set joining the cube of $y$,
$x$ and $x'$. Therefore
\begin{equation}\label{10.3:eq-kernel-decay-3}
M d_j(X)+\operatorname{dist}(x,X)+\operatorname{dist}(x',X)
\ge\operatorname{St}(y,x,x')-c_dM
\end{equation}
for a constant $c_d$ depending only on $d$.

The rate $\delta_0$ chosen for the correlation theorem satisfies $\delta_0M\ge2\kappa$; this is
the lemma fixing $\delta_0$ for the correlation theorem, which rests on \cite{B-CMP116}.
The admissible window for $\delta$ in that theorem is
$\delta\le\delta_J(L)=(1-4/L)/10<1/10$, so $\tfrac14\delta\kappa<\kappa/40$ and therefore
$2\kappa\ge\tfrac{79}{40}\kappa+\tfrac14\delta\kappa$. Consequently
\[
2\kappa d_j(X)+\delta_0\bigl(\operatorname{dist}(x,X)+\operatorname{dist}(x',X)\bigr)
\ge\tfrac14\delta\kappa\,d_j(X)+\tfrac{79}{40}\tfrac{\kappa}{M}
\bigl[Md_j(X)+\operatorname{dist}(x,X)+\operatorname{dist}(x',X)\bigr].
\]
By \eqref{10.3:eq-kernel-decay-3}, the bracket on the right is at least
$\operatorname{St}-c_dM$.

Insert this into (K1) and sum over $X\ni y$. The sum of $e^{-\frac14\delta\kappa d_j(X)}$ over
$X\ni y$ is bounded by the constant $C(\kappa)$ of the correlation theorem at the rate
$\tfrac14\delta\kappa$. With $c_d'=\tfrac{79}{40}c_d$ this gives
\begin{equation}\label{10.3:eq-kernel-decay-2}
\sum_X|\Pi^X(x,x',y;e)|\le C(\kappa)\,4B_3^2\,\mathfrak{n}\,|e|\,\mathfrak{r}_j^{-2}\,
e^{c_d'\kappa}\,e^{-(79/40)(\kappa/M)\operatorname{St}(y,x,x')} .
\end{equation}

Write $k$ entrywise as
\[
k(P,P')=\tfrac12\sum_{x,\mu,x',\nu'}\mathsf{K}(x,\mu;P)\,\Pi_{\mu\nu'}(x,x',0;e)\,
\mathsf{K}(x',\nu';P').
\]
Here
$\mathsf{K}(x,\mu;P)=\pm1$ if $P\in\Sigma(0;x,\mu)$ and $0$ otherwise, and $|P|\le|x|_2+1$ on
$\Sigma(0;x,\mu)$. Hence every term obeys
$\operatorname{St}\ge\tfrac12(|x|_2+|x'|_2)\ge\tfrac12(|P|+|P'|)-1$. Splitting
$\tfrac{79}{80}=\tfrac12+\tfrac{39}{80}$, each term is bounded by \eqref{10.3:eq-kernel-decay-2}
times
\[
e^{\kappa/M}\,e^{-(\kappa/2M)(|P|+|P'|)}\,e^{-(39/80)(\kappa/M)(|x|_2+|x'|_2)} .
\]
The last factor is summable over $x,\mu,x',\nu'$ with a bound depending only on $d$, $\kappa$ and
$M$. This absorbs the polynomial multiplicities and yields the first bound on $k$.

\emph{(K3).} The sum $\kappa_{\rho;AB}(u)=\sum_{P\in A}k_{\rho}(P,(u;B))$, with $A$, $B$ the
component labels and $(u;B)$ the plaquette notation of Lemma~\ref{10.3:lem-divergence-identity},
runs over all $x$ and over
those $x'$ with $(u;B)\in\Sigma(0;x',\cdot)$; for the latter, $|x'|_2\ge|u|_2-1$. For fixed
$(x,\mu)$ the number of $P$ with $\mathsf{K}(x,\mu;P)\ne0$ is at most
$d(1+|x|_1)\le d(1+2|x|_2)$. From the two lower bounds on $\operatorname{St}$,
\[
\tfrac{79}{40}\operatorname{St}=\tfrac{39}{20}\operatorname{St}+\tfrac1{40}\operatorname{St}
\ge\tfrac{39}{20}|x'|_2+\tfrac1{80}\bigl(|x|_2+|x'|_2\bigr).
\]
Using $|x'|_2\ge|u|_2-1$ in the first term, \eqref{10.3:eq-kernel-decay-2} gives the factor
$e^{(39/20)\kappa/M}e^{-(39/20)\kappa|u|_2/M}$. The factor
$e^{-(1/80)(\kappa/M)(|x|_2+|x'|_2)}$ absorbs the multiplicity $d(1+2|x|_2)$ and the sums over
$x,\mu,x',\nu'$. This is the bound on $\kappa_{\rho;AB}$. Inserting it into
\eqref{10.3:eq-divergence-identity-a} gives the bound on $\mu$ in (K3). The bound on $s$ in (K3)
is the clause $|s|\le C(\delta)\mathfrak{n}_2$ of Lemma~\ref{10.3:lem-divergence-identity}, with
the norm bound $\mathfrak{n}_2=C_K\mathfrak{n}\mathfrak{r}_j^{-2}$ of $k$ supplied by (K2).

The constants collected here are $4B_3^2$, $C(\kappa)$, $e^{c_d'\kappa}$, $e^{\kappa/M}$,
$e^{(39/20)\kappa/M}$ and the summation constants. They depend only on $d$, $M$, $\kappa$, $\delta$
and $B_3$.
\end{proof}

\begin{definition}[The bond current and sharp part of a frozen family]\label{10.3:def-bond-current}
Throughout, $j \ge 1$, $F$ is a frozen family as fixed in Definition~\ref{10.3:def-kernel-comb-gauge},
of type $V$ as fixed there, $\pi_j$ is the partition of $T^{(j)}$ into cubes fixed there,
lengths are measured in units of the spacing $a_j$ of $T^{(j)}$, and $G^c$ denotes the
complexification of $SU(N)$.

\emph{Transported curvature.} Let $\mathsf{V}$ be a $G^c$-valued configuration on the bonds of
$T^{(j)}$, let $y'$ be a site, let $P = (x;A)$ be a plaquette (based at $x$, in the plane $A$) and let
$\gamma$ be a lattice path from $y'$ to $x$. Put
\begin{equation}\label{10.3:eq-bond-current-1}
  \mathsf{F}^{\gamma}(P) := (2i)^{-1}\bigl(W - W^{-1}\bigr),
  \qquad W := \mathsf{V}(\gamma)\,\mathsf{V}(\partial P)\,\mathsf{V}(\gamma)^{-1},
\end{equation}
and write $\mathsf{F}^{\gamma}_A(x) := \mathsf{F}^{\gamma}((x;A))$. This is a polynomial in the links of
$\mathsf{V}$ and their inverses, and, for $N = 2$, it is traceless. On matrices $X_1, X_2$ put
$\langle X_1, X_2\rangle := \operatorname{tr} X_1X_2$, with $\operatorname{tr}\mathbf{1} = 1$; for
$N = 2$ and $\vec\sigma$ the Pauli matrices this gives
$\langle \vec w\cdot\vec\sigma, \vec w'\cdot\vec\sigma\rangle = \vec w\cdot\vec w'$.

\emph{Path rule.} The set $\operatorname{Seg}(y',x)$ is the union of the cubes of $\pi_j$ that meet the
straight segment $[y',x]$. It is a localisation domain, and
\begin{equation}\label{10.3:eq-bond-current-2}
  d_j\bigl(\operatorname{Seg}(y',x)\bigr) \le |x - y'|_2/M .
\end{equation}
The path $\gamma(y',x)$ is a lattice path from $y'$ to $x$ inside $\operatorname{Seg}(y',x)$, of length
at most $c_{\gamma}(|x - y'|_2 + M)$ with $c_{\gamma} = c_{\gamma}(d)$: inside each cube, which is a box
of lattice sites, it is any monotone path between the entry and the exit site. We write
$\mathsf{F}_A(y) := \mathsf{F}^{\gamma(y,y)}_A(y)$; since $\operatorname{Seg}(y,y)$ is the cube of $y$ and
$\gamma(y,y)$ is the constant path,
\begin{equation*}
  \mathsf{F}_A(y) = (2i)^{-1}\bigl(\mathsf{V}(\partial P) - \mathsf{V}(\partial P)^{-1}\bigr),
  \qquad P = (y;A).
\end{equation*}

\emph{The current.} Let $\mu_{\nu\rho;AB}(u)$ and $s_{\nu\rho;AB}$ be the kernels of the bond current of
Lemma~\ref{10.3:lem-divergence-identity}, taken for the kernel $k = \tfrac12\mathsf{K}^T\Pi\mathsf{K}$ of
Definition~\ref{10.3:def-kernel-comb-gauge} and Lemma~\ref{10.3:lem-kernel-decay}. For the bond
$b_\nu(y - e_\nu)$ in direction $\nu$ ending at $y$, and $e \in V_{\mathbb{C}}$, define
$\mathfrak{J}[F](b_\nu(y-e_\nu);\mathsf{V};e)$ to be the $W_4$-average of
\begin{equation}\label{10.3:eq-bond-current-3}
  e^{\rho}\Bigl[-2\sum_{A,B,u}\mu_{\nu\rho;AB}(u)
    \bigl\langle \mathsf{F}^{\gamma(y,y)}_A(y), \mathsf{F}^{\gamma(y,y+u)}_B(y+u)\bigr\rangle
   + \sum_{A,B} s_{\nu\rho;AB}\bigl\langle \mathsf{F}_A(y), \mathsf{F}_B(y)\bigr\rangle\Bigr],
\end{equation}
summed over $\rho$, where each element of $W_4$ transports the kernels, the base points and the path
rule together. On fine pairs $(\mathbf{U},\mathbf{J})$, in the sense of the sourced format of the
correlation theorem, put
\begin{equation}\label{10.3:eq-bond-current-4}
  \mathfrak{J}[F](b;(\mathbf{U},\mathbf{J});e) := \mathfrak{J}[F](b;M^j(\mathbf{U});e),
\end{equation}
with $M^j$ the level-$j$ averages, the level-$j$ data of the correlation theorem.

\emph{Localisation.} The term $(y,u)$ of a bond $b$ is localised in the set $X(y,u)$ formed by
$\operatorname{Seg}(y,y+u)$ together with the at most $3\cdot 2^d$ cubes of $\pi_j$ meeting $b$ and the
two plaquettes of the term. Thus $X(y,u)$ contains both ends of $b$, and
\begin{equation}\label{10.3:eq-bond-current-5}
  d_j\bigl(X(y,u)\bigr) \le |u|_2/M + c_d ,
\end{equation}
with a constant $c_d$ depending only on $d$.
The localised current $\mathfrak{J}[F](X,\cdot\,;b;e)$ is the sum of the terms of $b$ with
$X(y,u) = X$, the dot standing for the field arguments. Put
\begin{equation}\label{10.3:eq-bond-current-6}
  (\operatorname{div}\mathfrak{J})(X,\cdot,y;e)
  := \sum_{\nu}\bigl[\mathfrak{J}(X,\cdot\,;b_\nu(y);e) - \mathfrak{J}(X,\cdot\,;b_\nu(y-e_\nu);e)\bigr],
\end{equation}
and define the \emph{sharp part} $F^{\sharp}$ and the family $\mathsf{S}[F]$ by
\begin{align}
  F^{\sharp}(X,\cdot,y;e) &:= F(X,\cdot,y;e) - (\operatorname{div}\mathfrak{J}[F])(X,\cdot,y;e),
  \label{10.3:eq-bond-current-7}\\
  \mathsf{S}[F](X,\cdot,y;\mathrm{A}) &:= -\tfrac12\sum_{\nu}\bigl[\mathfrak{J}[F](X,\cdot\,;b_\nu(y);
    \mathrm{A}_\nu) + \mathfrak{J}[F](X,\cdot\,;b_\nu(y-e_\nu);\mathrm{A}_\nu)\bigr],
  \label{10.3:eq-bond-current-8}
\end{align}
for $\mathrm{A} \in (V\otimes V)_{\mathbb{C}}$, with $\mathrm{A}_\nu := (\mathrm{A}_{\nu\rho})_{\rho}$; here
the upright $\mathrm{A}$ is a tensor, not a plane index.
\end{definition}

\begin{proof}
The assertions in the definition are justified as follows.

\emph{Polynomial form.} Since $W^{-1} = \mathsf{V}(\gamma)\mathsf{V}(\partial P)^{-1}\mathsf{V}(\gamma)^{-1}$,
both $W$ and $W^{-1}$ are products of links of $\mathsf{V}$ and their inverses, so
\eqref{10.3:eq-bond-current-1} is a polynomial in these. For $N = 2$, the Cayley--Hamilton identity for
$W \in SL(2,\mathbb{C})$ reads $W + W^{-1} = 2\operatorname{tr}(W)\,\mathbf{1}$ (recall
$\operatorname{tr}\mathbf{1} = 1$), so $W - W^{-1} = 2W - 2\operatorname{tr}(W)\,\mathbf{1}$ and
\begin{equation*}
  \operatorname{tr}\bigl(W - W^{-1}\bigr) = 2\operatorname{tr}W - 2\operatorname{tr}W = 0 .
\end{equation*}

\emph{The segment domain.} The set $\operatorname{Seg}(y',x)$ is a union of cubes of $\pi_j$ and
contains the segment $[y',x]$, which meets each of its cubes. Thus it is connected and is a
localisation domain. The segment is a tree contained in $\operatorname{Seg}(y',x)$ that meets all its
cubes. Its length is $|x-y'|_2$. By the definition of $d_j$ fixed in the glossary, under which
$d_j(X)$ is at most $M^{-1}$ times the length of any tree contained in $X$ that meets all cubes of
$X$, this gives \eqref{10.3:eq-bond-current-2}.

\emph{The path.} The walls of the cubes of $\pi_j$ lie between lattice hyperplanes, since $M$ is odd.
Hence each cube is a box of lattice sites. A monotone lattice path between two sites of a box stays in
the box. The path $\gamma(y',x)$ is made of such monotone paths, one in each cube of
$\operatorname{Seg}(y',x)$, joining its entry and exit sites. The exit site of a cube and the entry
site of the next cube met by the segment are the sites of these two cubes nearest to a common boundary
point of theirs on the segment; they differ by at most one in each coordinate, and they are joined by
a monotone path of at most $d$ bonds. Each site of this path lies within $\tfrac12$ of that boundary
point in every coordinate, hence in a cube of $\pi_j$ containing the point, and every such cube meets
the segment; so this path lies in $\operatorname{Seg}(y',x)$. Therefore $\gamma(y',x)$ lies in
$\operatorname{Seg}(y',x)$. For the
length, the cubes of $\pi_j$ meeting $[y',x]$ have side $M$ and lie in the set of points at distance at
most $\sqrt{d}\,M$ from the segment, whose volume is at most $c(d)M^{d-1}(|x-y'|_2+M)$ for a constant
$c(d)$ depending only on $d$; since each cube has volume $M^d$, there are at most
$c(d)(|x-y'|_2/M+1)$ of them. A monotone path in a box with $M$ sites on each side has at most
$d(M-1)$ bonds, and each passage to the next cube adds at most $d$ bonds, so the length of
$\gamma(y',x)$ is at most
\begin{equation*}
  c(d)\bigl(|x-y'|_2/M+1\bigr)\,dM = d\,c(d)\bigl(|x-y'|_2+M\bigr),
\end{equation*}
which is the bound of the path rule with $c_{\gamma} = d\,c(d)$.

\emph{The localisation set.} The set $X(y,u)$ contains the cubes of $\pi_j$ meeting $b$, so it
contains both ends of $b$. By \eqref{10.3:eq-bond-current-2} with $y' = y$ and $x = y+u$,
$d_j(\operatorname{Seg}(y,y+u)) \le |u|_2/M$. The cubes meeting $b$ or one of the two plaquettes
number at most $3\cdot 2^d$, and each meets or is adjacent to a cube of $\operatorname{Seg}(y,y+u)$,
since $b$ and the two plaquettes each contain $y$ or $y+u$; adjoining them therefore increases $d_j$
by at most a constant $c_d$ depending only on $d$, which gives \eqref{10.3:eq-bond-current-5}.
\end{proof}

\begin{lemma}[Non-abelian completion]\label{10.3:lem-nonabelian-completion}
Let $y$ be a site, and let $\mathsf{V}=e^{iB}$ on the bonds of a box of side $D$ about $y$, where
$B$ is a matrix-valued function on the bonds of the box with
\begin{equation}\label{10.3:eq-nonabelian-completion-1}
|B(b)|\le 2d\,(1+|b-y|)\,\varphi_0 .
\end{equation}
Let $f:=dB$, the coboundary taken component-wise in the colour. If $P$ has boundary bonds
$b_1,b_2,b_3,b_4$ in the order of $\partial P$, with $b_3,b_4$ traversed against their
orientation, then $f(P)=B(b_1)+B(b_2)-B(b_3)-B(b_4)$, and $f$ is an exactly closed cochain.
Suppose
\begin{equation}\label{10.3:eq-nonabelian-completion-2}
|f(P)|\le 2\varphi_0\,(1+|P-y|)^2,\qquad (M+D)^2\varphi_0\le c_{NA}.
\end{equation}
Let $y'$ be a site with $|y'-y|_1\le 1$, let $P=(x;A)$ be a plaquette in the box, and let
$\gamma$ be a lattice path in the box from $y'$ to $x$ of length at most
$c_{\gamma}(|P-y'|_2+M)$. Then the transported curvature $\mathsf{F}^{\gamma}(P)$ of
Definition~\ref{10.3:def-bond-current} satisfies
\begin{equation}\label{10.3:eq-nonabelian-completion-3}
|\mathsf{F}^{\gamma}(P)-f(P)|\le C_{NA}\,(M+|P-y|)^4\,\varphi_0^2 .
\end{equation}
\end{lemma}

\begin{proof}
In this proof $C$ denotes a constant whose value may change from one occurrence to the next.
Since $P$ lies in the box and $M=L^{m'}\ge 1$, the quantity $(1+|P-y|)^2\varphi_0$ is at most
$C(M+D)^2\varphi_0\le C\,c_{NA}$ by \eqref{10.3:eq-nonabelian-completion-2}; by
\eqref{10.3:eq-nonabelian-completion-1} the matrices $B(b)$ on the bonds used below are small
in the same sense, and this smallness is what the expansions below use.

By Definition~\ref{10.3:def-bond-current}, $\mathsf{F}^{\gamma}(P)=(2i)^{-1}(W-W^{-1})$ with
$W=\mathsf{V}(\gamma)\mathsf{V}(\partial P)\mathsf{V}(\gamma)^{-1}$. Since
$W^{-1}=\mathsf{V}(\gamma)\mathsf{V}(\partial P)^{-1}\mathsf{V}(\gamma)^{-1}$,
\begin{equation}\label{10.3:eq-nonabelian-completion-4}
\mathsf{F}^{\gamma}(P)=\operatorname{Ad}\mathsf{V}(\gamma)\,\bigl(F^0\bigr),\qquad
F^0:=(2i)^{-1}\bigl(\mathsf{V}(\partial P)-\mathsf{V}(\partial P)^{-1}\bigr).
\end{equation}

\emph{The holonomy of the plaquette.} We have
$\mathsf{V}(\partial P)=e^{iB(b_1)}e^{iB(b_2)}e^{-iB(b_3)}e^{-iB(b_4)}$. By the
Baker--Campbell--Hausdorff formula this product equals $\exp(if(P)+R_2)$ with
$|R_2|\le C\max_{b\in\partial P}|B(b)|^2$, the linear term being
$i(B(b_1)+B(b_2)-B(b_3)-B(b_4))=if(P)$. By \eqref{10.3:eq-nonabelian-completion-1},
$|B(b)|\le C(1+|P-y|)\varphi_0$ for $b\in\partial P$, so
$|R_2|\le C(1+|P-y|)^2\varphi_0^2$.

\emph{The sine expansion.} Write $\mathsf{V}(\partial P)=e^{iX}$ with $X:=f(P)-iR_2$. Then
$F^0=(2i)^{-1}(e^{iX}-e^{-iX})=X+O(|X|^3)$. By \eqref{10.3:eq-nonabelian-completion-2} and
the bound on $R_2$, $|X|\le 2\varphi_0(1+|P-y|)^2+C(1+|P-y|)^2\varphi_0^2$, and since
$\varphi_0\le C c_{NA}$ this gives $|F^0|\le 3\varphi_0(1+|P-y|)^2$. Moreover
$|X|^3\le C\varphi_0^3(1+|P-y|)^6\le C\,c_{NA}\,\varphi_0^2(1+|P-y|)^4$, using
$(1+|P-y|)^2\varphi_0\le Cc_{NA}$. Hence
\begin{equation}\label{10.3:eq-nonabelian-completion-5}
|F^0-f(P)|\le |R_2|+C|X|^3\le C(M+|P-y|)^4\varphi_0^2 .
\end{equation}

\emph{Adjoint transport.} Writing $\mathsf{V}(\gamma)$ as the ordered product of the factors
$e^{\pm iB(b)}$, $b\in\gamma$, and telescoping, each factor contributes at most $2|B(b)|$, so
$|\operatorname{Ad}\mathsf{V}(\gamma)-1|\le 2\sum_{b\in\gamma}|B(b)|$. Each bond $b$ of $\gamma$
lies within path distance the length of $\gamma$ of $y'$, and $|y'-y|_2\le|y'-y|_1\le 1$, so
$|b-y|\le 1+c_{\gamma}(|P-y'|_2+M)\le C(M+|P-y|)$. The number of bonds of $\gamma$ is at most
$c_{\gamma}(|P-y'|_2+M)\le C(M+|P-y|)$. By \eqref{10.3:eq-nonabelian-completion-1},
$|\operatorname{Ad}\mathsf{V}(\gamma)-1|\le C(M+|P-y|)^2\varphi_0$. This operator acts on
$F^0$, of size at most $3\varphi_0(1+|P-y|)^2$, whence
\[
\bigl|(\operatorname{Ad}\mathsf{V}(\gamma)-1)F^0\bigr|\le C(M+|P-y|)^4\varphi_0^2 .
\]

\emph{Conclusion.} By \eqref{10.3:eq-nonabelian-completion-4},
$\mathsf{F}^{\gamma}(P)-f(P)=(\operatorname{Ad}\mathsf{V}(\gamma)-1)F^0+(F^0-f(P))$, and the last
display together with \eqref{10.3:eq-nonabelian-completion-5} gives
\eqref{10.3:eq-nonabelian-completion-3}. Since the right side of
\eqref{10.3:eq-nonabelian-completion-3} does not depend on $y'$ or $\gamma$, the base point and
the path are free up to an error of this size: two admissible choices change
$\mathsf{F}^{\gamma}(P)$ by at most $2C_{NA}(M+|P-y|)^4\varphi_0^2$.
\end{proof}

\begin{lemma}[Re-localisation of a frozen vector family]\label{10.3:lem-relocalisation}
Let $j\ge 1$, let $F$ be a $\mathsf{V}$-typed pointed $E$-family at scale $j$ on its radii
$\alpha^{\circ}_j$ (the radii at which it is created at scale $j$), with $\|F\|^{\mathrm{fr}}\le\mathfrak{n}$,
let $y$ be a site of $T^{(j)}$, let
$\square\in\pi_j$ be the cube of $y$, and measure lengths in units of the spacing of $T^{(j)}$. Let
$\mathfrak{J}[F]$, $F^{\sharp}$ and $\mathsf{S}[F]$ be the current, the sharp part and the symmetrised
current of $F$ (Definition~\ref{10.3:def-bond-current}). There is a bound depending only on $L$, $M$
and $\kappa$ such that, for every $g$ below it, the following hold.
\begin{enumerate}
\item[(a)] $\mathfrak{J}[F]$, and hence $F^{\sharp}$ and $\mathsf{S}[F]$, are linear in $F$, analytic
in every parameter in which $F$ is analytic, and have on the radii of $F$ the properties (f1)
(analyticity on the spaces of the correlation theorem) and (f2) (gauge invariance) of pointed
$E$-families;
$\mathfrak{J}[F]$, and with it $\mathsf{S}[F]$, vanishes to second order at $\mathbf{1}$ ($F^{\sharp}$
does so whenever $F$ does); and
\begin{equation}\label{10.3:eq-relocalisation-1}
\sup e^{\frac{15}{8}\kappa d_j(X)}\,\bigl|\mathfrak{J}[F](X,\cdot\,;b;e)\bigr|
\le C_{\mathfrak{J}}\,\mathfrak{n}\,|e|,
\end{equation}
the supremum being taken over localisation domains $X$, bonds $b$ and configurations within the radii
of $F$.
Moreover $F^{\sharp}$ is a $\mathsf{V}$-typed and $\mathsf{S}[F]$ a
$(\mathsf{V}\otimes\mathsf{V})$-typed pointed $E$-family, and
\begin{equation}\label{10.3:eq-relocalisation-2}
\|F^{\sharp}\|\le(1+8C_{\mathfrak{J}})\,\mathfrak{n},\qquad
\|\mathsf{S}[F]\|\le 4C_{\mathfrak{J}}\,\mathfrak{n},
\end{equation}
in $\|\cdot\|^{\flat}$ and at the rate $\kappa_{\sharp}:=\tfrac{15}{8}\kappa$.
\item[(b)] For the pointed totals, and term by term in $X$,
\begin{equation}\label{10.3:eq-relocalisation-3}
F(U,y;e)=F^{\sharp}(U,y;e)
+\sum_{\sigma}\bigl[\mathfrak{J}[F](b_{\sigma}(y);U;e)-\mathfrak{J}[F](b_{\sigma}(y-e_{\sigma});U;e)\bigr].
\end{equation}
\item[(c)] Let $\mathfrak{b}$ be admissible at scale $n\ge j$ near $y$. Then the sharp core at $y$
obeys
\begin{equation}\label{10.3:eq-relocalisation-4}
\Bigl|\sum_{X\ni y,\ X\subset\square^{\sim\nu}}
\bigl[F^{\sharp}(X,\mathfrak{b},y;e)-F^{\sharp}(X,\mathbf{1},y;e)\bigr]\Bigr|
\le C_{\sharp}\,\mathfrak{n}\,|e|\,\nu^{3}\rho_{n,j}^{2}(1+\rho_{n,j})\,t^{4+2\hat\alpha}.
\end{equation}
Every constituent $F^{\sharp}(X,\cdot,y;e)$ not entering the core (a tail) obeys, on every background
$\mathfrak{b}'$ admissible at scale $n$, whether or not such a background is present at $X$,
\begin{equation}\label{10.3:eq-relocalisation-5}
\bigl|F^{\sharp}(X,\mathfrak{b}',y;e)-F^{\sharp}(X,\mathbf{1},y;e)\bigr|
\le\begin{cases}
2(1+8C_{\mathfrak{J}})\,\mathfrak{n}\,|e|\,(C_M D_X t^{2}\rho_{n,j})^{2}\,e^{-\kappa d_j(X)}
& \text{if } X\subset\tilde{\square}^{\sim 2},\\[2pt]
2(1+8C_{\mathfrak{J}})\,\mathfrak{n}\,|e|\,e^{-\kappa d_j(X)} & \text{otherwise}.
\end{cases}
\end{equation}
Here $C_M$ is the constant of part (i) of the background-difference lemma of the proof of the
correlation theorem; in the first case $d_j(X)\ge\nu-1$, in the second $d_j(X)\ge t^{-1}$. Per unit
$n$-cell, after the weights of the resummation lemma of the proof of the correlation theorem and of
\cite[(1.66)]{B-CMP122b}, the tails of level $j$ sum to at most
$C_{\sharp}\,\mathfrak{n}\,|e|\,\rho_{n,j}^{2}t^{5}$.
\item[(d)] Statement (c) holds, in the form of part (iii) of the background-difference lemma of the
proof of the correlation theorem, along holomorphic families of admissible backgrounds.
\end{enumerate}
\end{lemma}

\begin{proof}
Two conditions on $g$ are imposed below,
\begin{equation}\label{10.3:eq-relocalisation-6}
2c_{\gamma}C\alpha^{\circ}_j\le\frac{\kappa}{40M},\qquad 36M^{2}\mathfrak{a}_n\le c_{NA};
\end{equation}
here $C$ in the first condition is the constant of the holonomy bound
$|\mathsf{V}(\gamma)^{\pm1}|\le e^{C\alpha^{\circ}_j|\gamma|}$ used in the proof of
\eqref{10.3:eq-relocalisation-1}. Both conditions
are met for $g$ below a bound depending on $(L,M,\kappa)$, and together they are the ceiling on $g$
of the statement.

\emph{(a), linearity, analyticity, vanishing at $\mathbf{1}$.} By
Definition~\ref{10.3:def-bond-current}, $\mathfrak{J}[F]$ is built from the kernels $\mu$ and $s$ that
Lemma~\ref{10.3:lem-divergence-identity} attaches to the kernel
$k=\tfrac12\mathsf{K}^{T}\Pi\mathsf{K}$, and these are linear in $k$. The kernel $k$ is linear in
$\Pi$, and $\Pi$ is a second derivative at $0$ of $F$; hence $\mathfrak{J}[F]$ is linear in $F$ and
analytic in every parameter in which $F$ is analytic. Since $F^{\sharp}=F-\operatorname{div}\mathfrak{J}[F]$
and $\mathsf{S}[F]$ is a linear combination of values of $\mathfrak{J}[F]$, the same holds for them.
Each term of $\mathfrak{J}[F]$ is a product $\langle\mathsf{F}^{\gamma}(P),\mathsf{F}^{\gamma'}(P')\rangle$
of two transported curvatures, $\gamma$ and $\gamma'$ being paths given by the path rule, and each factor
vanishes at $\mathsf{V}=\mathbf{1}$, where $W=\mathbf{1}$;
so $\mathfrak{J}[F]$ and $\mathsf{S}[F]$ vanish to second order at $\mathbf{1}$. Since
$F^{\sharp}=F-\operatorname{div}\mathfrak{J}[F]$, $F^{\sharp}$ vanishes to second order at $\mathbf{1}$
whenever $F$ does.

\emph{(a), property (f1).} The level-$j$ averages $M^{j}(\mathbf{U})$ are defined and analytic on the
spaces $\mathcal{U}^{c}_j(X,\cdot)$ of the correlation theorem: they are the object of the defining
clause (iv) of these spaces in the correlation theorem, which rests on the construction in
\cite{B-CMP109} of the functions $U_n$ and of the pairs
$(U_n(M^{\cdot}(\mathbf{U})),J_n(M^{\cdot}(\mathbf{U})))$. On fine pairs
$\mathfrak{J}[F](b;(\mathbf{U},\mathbf{J});e)=\mathfrak{J}[F](b;M^{j}(\mathbf{U});e)$, and
$\mathfrak{J}(X,\cdot)$ is a polynomial in $M^{j}(\mathbf{U})^{\pm1}$ on bonds inside $X$, because each
$\mathsf{F}^{\gamma}(P)$ is a polynomial in the links of $\gamma$ and $\partial P$ and their inverses,
and for a term $(y',u)$ the path $\gamma(y',y'+u)$ lies in $\operatorname{Seg}(y',y'+u)\subset X(y',u)$.
Hence (f1) holds on the radii of $F$.

\emph{(a), property (f2).} The map $M^{j}$ is gauge covariant. Each term of $\mathfrak{J}[F]$ at a bond
with base point $y'$ is the trace of a product of two holonomies based at the one site $y'$: the matrix
$W=\mathsf{V}(\gamma)\mathsf{V}(\partial P)\mathsf{V}(\gamma)^{-1}$ is the holonomy of the loop that runs
from $y'$ along $\gamma$, around $\partial P$ and back along $\gamma$. A gauge transformation conjugates
both factors by
the same element at $y'$, and the trace is invariant under this conjugation. Hence (f2).

\emph{(a), typed covariance.} $U_j$ and $M^{j}$ commute with the Euclidean symmetries of $T^{(j)}$;
this is the fact on which the covariance (f3) itself rests, \cite[(2.29)]{B-CMP119}. The bond current is
covariant by part (ii) of Lemma~\ref{10.3:lem-divergence-identity}, and the $W_4$-average in
Definition~\ref{10.3:def-bond-current} transports kernels, base points and the path rule together, so
$\mathfrak{J}[F]$ inherits this covariance. A reflection in the direction $\sigma$ through $y$ exchanges
$b_{\sigma}(y-e_{\sigma})$ and $b_{\sigma}(y)$, with a sign. Therefore $\operatorname{div}\mathfrak{J}$ and
the symmetrised $\mathsf{S}[F]$ are covariant with the pin $y$: $F^{\sharp}$ is $\mathsf{V}$-typed and
$\mathsf{S}[F]$ is $(\mathsf{V}\otimes\mathsf{V})$-typed.

\emph{(a), the bound \eqref{10.3:eq-relocalisation-1}.} A point of the tube is an admissible background
at scale $j$, since the spaces of the correlation theorem include the $(\mathbf{U},\mathbf{J})$-extensions
on Ba{\l}aban's tubes. Hence, by the first of the conditions defining admissible backgrounds at scale $n$
in the correlation theorem, read at $n=j$,
$|\mathsf{V}(\partial P)-\mathbf{1}|\le C_{\mathfrak{a}}\alpha^{\circ}_j$ for every plaquette $P$, where
$\mathsf{V}=M^{j}(\mathbf{U})$ and $C_{\mathfrak{a}}$ is the constant of that condition. Along a path,
$|\mathsf{V}(\gamma)^{\pm1}|\le e^{C\alpha^{\circ}_j|\gamma|}$.
Conjugation by $\mathsf{V}(\gamma)$ therefore costs at most $e^{2C\alpha^{\circ}_j|\gamma|}$. For
$\gamma=\gamma(y',y'+u)$ the path rule gives $|\gamma|\le c_{\gamma}(|u|_2+M)$, and the first condition
in \eqref{10.3:eq-relocalisation-6} gives $2c_{\gamma}C\alpha^{\circ}_j\le\kappa/40M$. Therefore
\begin{equation*}
|\mathsf{F}^{\gamma}(P)|\le 2C_{\mathfrak{a}}\alpha^{\circ}_j\,e^{\kappa(|u|_2+M)/40M}.
\end{equation*}
The factor at the base point has $u=0$ and costs $e^{\kappa/40}$. By the bound (K3) of
Lemma~\ref{10.3:lem-kernel-decay}, see \eqref{10.3:eq-kernel-decay-a},
\begin{equation*}
|\mu_{\sigma\rho;AB}(u)|\le C_K'\mathfrak{n}\mathfrak{r}_j^{-2}e^{-\frac{77}{40}\kappa|u|_2/M},
\qquad |s|\le C_K'\mathfrak{n}\mathfrak{r}_j^{-2}.
\end{equation*}
Since $\tfrac{77}{40}-\tfrac{1}{40}=\tfrac{19}{10}$, the term $(y',u)$ of $\mathfrak{J}[F]$, with the
finitely many plaquette orientations $A,B$, the contraction with $e$ and the $W_4$-average included, is
at most
\begin{equation*}
C\,\mathfrak{n}\,|e|\,(\alpha^{\circ}_j/\mathfrak{r}_j)^{2}\,e^{-\frac{19}{10}\kappa|u|_2/M}.
\end{equation*}
It is localised in $X(y',u)$, and, by Definition~\ref{10.3:def-bond-current},
$d_j(X(y',u))\le|u|_2/M+c_d$ with the constant $c_d$ of that definition, so the weight it meets is
\begin{equation*}
e^{\frac{15}{8}\kappa d_j(X(y',u))}\le e^{\frac{15}{8}\kappa(|u|_2/M+c_d)}.
\end{equation*}
As $\tfrac{19}{10}-\tfrac{15}{8}=\tfrac{1}{40}$, a factor $e^{-\kappa|u|_2/40M}$ is left. At most
$(c_dM)^{4}(1+d_j(X))^{4}$ values of $u$ share one $X$, and since $d_j(X)\le|u|_2/M+c_d$ this
multiplicity is absorbed by the factor left over. Finally $\alpha^{\circ}_j/\mathfrak{r}_j$ is a quotient
of constants of \cite[(2.28)]{B-CMP119}. This gives \eqref{10.3:eq-relocalisation-1} with a constant
$C_{\mathfrak{J}}$.

\emph{(a), the bounds \eqref{10.3:eq-relocalisation-2}.} The divergence at $y$ involves the eight bonds
$b_{\sigma}(y)$ and $b_{\sigma}(y-e_{\sigma})$, $\sigma=1,\dots,4$. By \eqref{10.3:eq-relocalisation-1} it
is bounded by $8C_{\mathfrak{J}}\mathfrak{n}$ at the rate $\kappa_{\sharp}$. Adding the bound $\mathfrak{n}$
for $F$ gives the first bound. For $A\in(\mathsf{V}\otimes\mathsf{V})_{\mathbb{C}}$ with rows
$A_{\sigma}=(A_{\sigma\rho})_{\rho}$, as in Definition~\ref{10.3:def-bond-current},
$\mathsf{S}[F](\cdot\,;A)$ carries the factor $\tfrac12$ in front of the same eight terms, and
$|A_{\sigma}|\le|A|$; this gives the second bound.

\emph{(b).} By Definition~\ref{10.3:def-bond-current}, for every $X$,
\begin{equation*}
F^{\sharp}(X,\cdot,y;e)=F(X,\cdot,y;e)-(\operatorname{div}\mathfrak{J}[F])(X,\cdot,y;e).
\end{equation*}
This is \eqref{10.3:eq-relocalisation-3} term by term in $X$. Summing over $X$, with
$\mathfrak{J}[F](b;\cdot;e)=\sum_X\mathfrak{J}[F](X,\cdot;b;e)$, gives it for the pointed totals.

\emph{(c), near levels.} Put $D:=M(2\nu+3)$. Suppose first that $C_MDt^{2}\rho_{n,j}>\tfrac12$. As in the
proof of the core lemma of the correlation theorem, this happens only for $\nu\le\nu_*(L,M)$. Bound
$F^{\sharp}(X,\mathfrak{b},y;e)$ and $F^{\sharp}(X,\mathbf{1},y;e)$ separately by
\eqref{10.3:eq-relocalisation-2} and sum over $X$ with the constant $C(\kappa)$ by which the correlation
theorem bounds sums over localisation domains $X\ni y$. The resulting trivial bound
$2C(\kappa)(1+8C_{\mathfrak{J}})\mathfrak{n}|e|$ lies inside the claim, since $C_{\sharp}$ is chosen with
$C_{\sharp}\ge CL^{6\nu_*}$. From now on $C_MDt^{2}\rho_{n,j}\le\tfrac12$.

\emph{(c), tails.} By (a), the single functions $F^{\sharp}(X,\cdot,y;e)$ have (f1) and (f2), with norm
$(1+8C_{\mathfrak{J}})\mathfrak{n}$. Parts (i) and (ii) of the background-difference lemma of the proof of
the correlation theorem apply to them and give \eqref{10.3:eq-relocalisation-5}. For a tail with
$X\subset\tilde{\square}^{\sim 2}$ one has $d_j(X)\ge\nu-1$, as in the tail class of the core lemma of that
proof; in the sum per unit $n$-cell a factor $e^{-\frac12\kappa d_j(X)}\le t^{5}$ is therefore extracted
from
$e^{-\kappa d_j(X)}$ after the resummation weights, exactly as in the proof of part (i) of the derivative
lemma of the proof of the correlation theorem; this uses the standing lower bound $\kappa\ge 10\log L$ of
the correlation theorem. It comes on top of the factor $t^{4}$ carried by $(C_MD_Xt^{2}\rho_{n,j})^{2}$ in
part (i) of the background-difference lemma. For $X\not\subset\tilde{\square}^{\sim 2}$ one has
$d_j(X)\ge t^{-1}$, and
$e^{-\frac12\kappa/t}\le t^{9}$. Summed over the $t^{-4}$ points of a unit $n$-cell, the tails of level
$j$ give at most $C_{\sharp}\mathfrak{n}|e|\rho_{n,j}^{2}t^{5}$.

\emph{(c), the core: setting.} Put $\square_0:=\square^{\sim\nu+1}$ and work in the comb gauge centred at
$y$. By the first of the conditions defining admissible backgrounds in the correlation theorem, with
Ba{\l}aban's one-layer margin, and by the second and third of these conditions, the background
$\mathfrak{b}$ is, modulo gauge, $U_j(\square_0,\mathsf{V})$, where
\begin{equation*}
\mathsf{V}=e^{iB^{\mathrm{ax}}}=M^{j}(\mathfrak{b}),\qquad
B^{\mathrm{ax}}_{\mu}(x)=\sum_{\sigma<\mu}(x_{\sigma}-y_{\sigma})\mathsf{f}_{\sigma\mu}(y)
+\mathfrak{e}_{\mu}(x),
\end{equation*}
with constant coefficients $\mathsf{f}_{\sigma\mu}(y)$; we write $\mathsf{f}(y)$ for them and for the
constant cochain they define. They obey
\begin{equation*}
|\mathsf{f}(y)|\le\varphi_0:=\mathfrak{a}_nt^{2},\qquad
|\mathfrak{e}_{\mu}(x)|\le C(1+|x-y|)^{2}\mathfrak{a}_nt^{2+\hat\alpha}.
\end{equation*}
By (f2), which $F$ has and which $\mathfrak{J}[F]$ has by (a), every quantity below may be evaluated at
this representative.

\emph{Step 1.} Apply part (c) of the expansion statement of the proof of the correlation theorem; its
first-order term vanishes, the derivative at $0$ of the one-parameter function $\varphi_X$ of that
expansion being zero, $\varphi_X'(0)=0$. This gives
\begin{multline*}
\sum_{X\ni y,\ X\subset\square^{\sim\nu}}\bigl[F(X,\mathfrak{b},y;e)-F(X,\mathbf{1},y;e)\bigr]\\
=\frac12\sum_X\bigl\langle F^{(2)}(X,y;e);B^{\mathrm{ax}},B^{\mathrm{ax}}\bigr\rangle
+O\bigl(C(\kappa)\,\mathfrak{n}\,|e|\,(C_MD\rho_{n,j})^{3}t^{6}\bigr),
\end{multline*}
where $F^{(2)}(X,y;e)$ is the second derivative of $F(X,\cdot,y;e)$ at the identity.

\emph{Step 2.} Let $\mathbf{1}_{\square^{\sim\nu}}$ be the indicator of the bonds of $\square^{\sim\nu}$,
and put, for $x\in\mathbb{Z}^{4}$,
\begin{equation*}
\tilde B_{\mu}(x):=\sum_{\sigma<\mu}(x_{\sigma}-y_{\sigma})\mathsf{f}_{\sigma\mu}(y)
+(\mathbf{1}_{\square^{\sim\nu}}\mathfrak{e})_{\mu}(x),
\end{equation*}
that is, the linear part of $B^{\mathrm{ax}}$ extended to the whole lattice plus the error
$\mathfrak{e}$ kept only inside $\square^{\sim\nu}$.
Three replacements are made: $U_j(\square_0,\cdot)$ by the free $U_j$, the sum over
$X\subset\square^{\sim\nu}$ by the sum over all $X\ni y$, and $B^{\mathrm{ax}}$ by $\tilde B$. They act on
a quadratic form of size $C\mathfrak{n}|e|D^{2}\rho_{n,j}^{2}t^{4}$, and they cost the factors
$e^{-\delta_0M\nu}$ and $e^{-\frac12\kappa\nu}\le t^{5}$ of the proof of the core lemma of the
correlation theorem. The resulting errors lie inside the claim. The remaining quadratic form is
$\tfrac12\sum_{x,x'}\Pi(x,x',y;e)\tilde B(x)\tilde B(x')$. By (K2) of
Lemma~\ref{10.3:lem-kernel-decay}, it equals $e^{\rho}Q[F]_{\rho}(y;\tilde f)$ (summed over $\rho$),
where
\begin{equation*}
\tilde f:=d\tilde B=\mathsf{f}(y)+d(\mathbf{1}_{\square^{\sim\nu}}\mathfrak{e}).
\end{equation*}
Being a coboundary, $\tilde f$ is exactly closed, and it is of polynomial growth, as (K2) requires. It
obeys
\begin{equation*}
|\tilde f_A(y+u)-\tilde f_A(y)|\le\varphi_1(1+|u|)^{2},\qquad
\varphi_1:=32C\mathfrak{a}_nt^{2+\hat\alpha}.
\end{equation*}
Keeping $\tilde f$ here, rather than replacing it by its main term, avoids the loss of a factor
$t^{\hat\alpha}$ that the replacement would cost.

\emph{Step 3.} Since $\tilde f$ is closed, parts (i) and (iii) of
Lemma~\ref{10.3:lem-divergence-identity} apply to the kernel of (K2). They give
\begin{equation*}
e^{\rho}Q[F]_{\rho}(y;\tilde f)
=\sum_{\sigma}e^{\rho}\bigl[\mathfrak{C}_{\rho}(b_{\sigma}(y);\tilde f)
-\mathfrak{C}_{\rho}(b_{\sigma}(y-e_{\sigma});\tilde f)\bigr]+e^{\rho}\mathsf{E}_{\rho}(y;\tilde f).
\end{equation*}
By the kernel bound of (K2), which carries the decay factor $e^{-(\kappa/2M)(|P|+|P'|)}$, the norm bound
$\mathfrak{n}_2$ of Lemma~\ref{10.3:lem-divergence-identity} for $k=\tfrac12\mathsf{K}^{T}\Pi\mathsf{K}$ is
at most $C_K\mathfrak{n}\mathfrak{r}_j^{-2}$, and part (iii) of that lemma bounds the remainder:
\begin{equation*}
|e^{\rho}\mathsf{E}_{\rho}(y;\tilde f)|
\le C\,C_QC_K\,\mathfrak{n}\,\mathfrak{r}_j^{-2}\varphi_1^{2}\,|e|
\le C\,\mathfrak{n}\,|e|\,\rho_{n,j}^{2}t^{4+2\hat\alpha},
\end{equation*}
where $C_Q$ is the constant of part (iii) of Lemma~\ref{10.3:lem-divergence-identity}.
The second inequality uses $\mathfrak{a}_n/\mathfrak{r}_j=C\rho_{n,j}$.

\emph{Step 4.} Apply Lemma~\ref{10.3:lem-nonabelian-completion} with $B=B^{\mathrm{ax}}$ and with $D$ as
above. Since $Dt\le 5M$,
\begin{equation*}
(M+D)^{2}\varphi_0\le 36M^{2}\mathfrak{a}_n\le c_{NA},
\end{equation*}
by the second condition in \eqref{10.3:eq-relocalisation-6}; this is the smallness condition of that
lemma. The base point $y'$ of a bond $b\ni y$ satisfies $|y'-y|_1\le1$. The paths are those of the path
rule, of length $\le c_{\gamma}(|P-y'|_2+M)$. On the bonds of $\square^{\sim\nu}$, $\tilde B=B^{\mathrm{ax}}$,
so $\tilde f=dB^{\mathrm{ax}}$ on the plaquettes met by terms localised in $\square^{\sim\nu}$. The term of
$\mathfrak{J}[F](b;\mathfrak{b};e)$ and the term of $e^{\rho}\mathfrak{C}_{\rho}(b;\tilde f)$ with the same
indices carry the same kernel $\mu(u)$ or $s$; the first has the transported curvatures $\mathsf{F}$
where the second has $\tilde f$. The non-abelian completion lemma bounds
$|\mathsf{F}^{\gamma}(P)-\tilde f(P)|$ by $C_{NA}(M+|P-y|)^{4}\varphi_0^{2}$, and each factor is at most of
order $\varphi_0(1+|P-y|)^{2}$. Hence each such term with $X(\cdot,u)\subset\square^{\sim\nu}$ differs
from its counterpart by at most $C|\mu(u)|(M+|u|)^{6}\varphi_0^{3}|e|$. Summed with (K3), these
differences total at most
\begin{equation*}
C\,\mathfrak{n}\,|e|\,\mathfrak{r}_j^{-2}\mathfrak{a}_n^{3}t^{6}
\le C\,\mathfrak{n}\,|e|\,\rho_{n,j}^{2}\mathfrak{a}_nt^{6}.
\end{equation*}
The terms of $\mathfrak{C}$ with $X(\cdot,u)\not\subset\square^{\sim\nu}$ have $|u|_2\ge(\nu-1)M$. By (K3)
they sum to at most
\begin{equation*}
C\,\mathfrak{n}\,|e|\,\rho_{n,j}^{2}\nu^{4}t^{4}e^{-\frac32\kappa(\nu-1)}.
\end{equation*}

\emph{Step 5.} By (a), $\mathfrak{J}[F]$ vanishes at $\mathbf{1}$, so
$F^{\sharp}(X,\mathbf{1},y;e)=F(X,\mathbf{1},y;e)$. Hence the left side of
\eqref{10.3:eq-relocalisation-4} is
\begin{equation*}
\sum_{X\ni y,\ X\subset\square^{\sim\nu}}\bigl[F(X,\mathfrak{b},y;e)-F(X,\mathbf{1},y;e)\bigr]
-\sum_{X\ni y,\ X\subset\square^{\sim\nu}}(\operatorname{div}\mathfrak{J}[F])(X,\mathfrak{b},y;e).
\end{equation*}
The second sum collects exactly the terms of $\mathfrak{J}[F]$ at the bonds $b_{\sigma}(y)$ and
$b_{\sigma}(y-e_{\sigma})$, which contain $y$, with $X(\cdot,u)\subset\square^{\sim\nu}$. By Steps 1--3
the first sum is the divergence of $e^{\rho}\mathfrak{C}_{\rho}(\cdot\,;\tilde f)$ plus
$e^{\rho}\mathsf{E}_{\rho}$ plus the errors of Steps 1 and 2. By Step 4 the second sum is the part of
that divergence formed by the terms with $X(\cdot,u)\subset\square^{\sim\nu}$, up to an error at most
$C\mathfrak{n}|e|\rho_{n,j}^{2}\mathfrak{a}_nt^{6}$ of Step 4. Subtracting, these parts of the divergence cancel;
the terms of the divergence with
$X(\cdot,u)\not\subset\square^{\sim\nu}$ do not cancel, and they are the ones bounded at the end of
Step 4. What remains is therefore $e^{\rho}\mathsf{E}_{\rho}$ and the errors of Steps
1, 2 and 4; with $t^{6}\le t^{4+2\hat\alpha}$ they are bounded by the right side of
\eqref{10.3:eq-relocalisation-4}.

\emph{(d).} Part (a) of the expansion statement of the proof of the correlation theorem is used in place
of part (c), as in part (iii) of the background-difference lemma. The argument of (c) then gives (c) in
the form of that part along holomorphic families of admissible backgrounds.
\end{proof}

\begin{remark}
The rate $\kappa_{\sharp}$ of Lemma~\ref{10.3:lem-relocalisation} exceeds the resummation weight
$(1+3\beta_s)\kappa=\tfrac74\kappa$ of \cite[(1.66)]{B-CMP122b}, $\beta_s$ being the parameter of that
display, by $\tfrac18\kappa\ge\tfrac12\delta\kappa$, $\delta<\tfrac1{10}$ being the parameter of the
correlation theorem whose multiple $\tfrac14\delta\kappa$ is the rate of its sums over localisation
domains; this excess is what the fresh rate is used for in the proof of the correlation theorem, at the
step where $j=k$.

No datum attached to a point other than $y$, and no weight, enters the bounds of
Lemma~\ref{10.3:lem-relocalisation}. For data given point by point, such as
$(\hat w(y),\mathsf{h}(y))$, the lemma therefore applies to the family frozen at those values, as for the
core lemma of the proof of the correlation theorem. Per unit $n$-cell, which contains $t^{-4}$ points,
the core bound \eqref{10.3:eq-relocalisation-4} gives
$C_{\sharp}\mathfrak{n}|e|\varsigma^{V}_{\nu}\cdot(1+o(1))$, $\varsigma^{V}_{\nu}$ being the scale kernel
of Definition~\ref{10.3:def-exponents-rates}. The kernel $L^{-2\hat\alpha(\nu-1)}$ takes
the place of the kernel $L^{-\hat\alpha(\nu-1)}$ of the proof of the correlation theorem.
\end{remark}

\begin{lemma}[Re-localisation on extended backgrounds]\label{10.3:lem-relocalisation-extended}
Work in the setting and notation of the lemma of the correlation theorem that bounds the
core on extended backgrounds at scale $p$. Let $j \le p$, and put
\[
\nu_p := 1 + p - j, \qquad t_p := L^{\,j-p}, \qquad Q := Q^{(p)} .
\]
Let $\epsilon \le 1$. Suppose that $\mathcal{U}{\restriction}Q$ reads, in one
$G^{\mathbb{C}}$-gauge, as $e^{i\eta\mathcal{A}}$, with
\[
\max\bigl\{\ell_p|\mathcal{A}|,\ \ell_p^2|\nabla\mathcal{A}|\bigr\} \le \epsilon\, c''\mathfrak{r}_p .
\]
Here $G^{\mathbb{C}}$, the length $\ell_p$ and the constants $c''$ and $K(d)$ are those of that
lemma. Then the $\sharp$-core at $y$ obeys
\[
\bigl|\operatorname{Core}^{\sharp}_y(\operatorname{ext}_p(\mathcal{U});e)\bigr|
\le 2\epsilon^2\mathfrak{N}^V_p ,
\]
where
\begin{equation}\label{10.3:eq-relocalisation-extended-1}
\mathfrak{N}^V_p := \bigl(C_{\sharp} + 2K(d)(1+8C_{\mathfrak{J}})\bigr)\,\mathfrak{n}\,|e|\,
\nu_p^3\,\bar\rho_{p-j}^{\,2}\,(1+\bar\rho_{p-j})\,t_p^{\,4+2\hat\alpha} .
\end{equation}
\end{lemma}

\begin{proof}
We follow the proof of the lemma of the correlation theorem named in the statement, with the
function
\begin{equation}\label{10.3:eq-relocalisation-extended-2}
\Phi(\sigma) := \operatorname{Core}^{\sharp}_y\bigl(\operatorname{ext}_p(e^{i\sigma\eta\mathcal{A}});e\bigr),
\qquad |\sigma| < 1/\epsilon .
\end{equation}
Here $\eta$ is the gauge parameter and $\mathcal{A}$ the potential of that lemma.

For $|\sigma| < 1/\epsilon$ the potential $\sigma\mathcal{A}$ satisfies
\[
\max\bigl\{\ell_p|\sigma\mathcal{A}|,\ \ell_p^2|\nabla(\sigma\mathcal{A})|\bigr\}
\le |\sigma|\,\epsilon\, c''\mathfrak{r}_p < c''\mathfrak{r}_p .
\]
By the admissibility property of the extension $\operatorname{ext}_p$ stated in the setting of
that lemma, the argument $\operatorname{ext}_p(e^{i\sigma\eta\mathcal{A}})$ is therefore an
admissible background at scale $p$ on $Q^{\sim-2}$. Moreover it depends analytically on
$\sigma$. Since $Q^{\sim-2} \supseteq \square^{\sim\nu_p}$, the cube of part (c) of
Lemma~\ref{10.3:lem-relocalisation} at $\nu = \nu_p$, it is an admissible background at
scale $p$ near $y$ on the region over which the $\sharp$-core is summed.

Consequently $\Phi$ is analytic on the disc $|\sigma| < 1/\epsilon$. For each such $\sigma$,
part (c) of the re-localisation lemma for a frozen vector family
(Lemma~\ref{10.3:lem-relocalisation}) applies with $(n,\nu,t)$ replaced by
$(p,\nu_p,t_p)$. It bounds the $\sharp$-core, and gives
\[
|\Phi(\sigma)| \le \mathfrak{N}^V_p \qquad (|\sigma| < 1/\epsilon).
\]

Next, $\Phi(0) = 0$. Indeed, at $\sigma = 0$ the argument is
$\operatorname{ext}_p(\mathbf{1}) = \mathbf{1}$, and each term
$F^{\sharp}(X,\mathbf{1},y;e) - F^{\sharp}(X,\mathbf{1},y;e)$ of the $\sharp$-core vanishes.

Further, $\Phi'(0) = 0$. The derivative of $\Phi$ at $\sigma = 0$ is a first derivative, at
$\mathbf{1}$, of the terms
\[
F^{\sharp}(X,\cdot,y;e) = F(X,\cdot,y;e) - (\operatorname{div}\mathfrak{J}[F])(X,\cdot,y;e),
\]
composed with the analytic path $\sigma \mapsto \operatorname{ext}_p(e^{i\sigma\eta\mathcal{A}})$.
The first derivative of $F$ at the unit configuration vanishes, $DF(X;\mathbf{1},0) = 0$.
By part (a) of Lemma~\ref{10.3:lem-relocalisation}, $\mathfrak{J}[F]$ vanishes to second order
at $\mathbf{1}$, and hence so does its lattice divergence. By the chain rule, both
contributions to $\Phi'(0)$ vanish.

It remains to apply the second-order Schwarz bound used in the proof of the lemma named in the
statement. The function $\Phi$ is analytic and bounded by $\mathfrak{N}^V_p$ on the disc
$|\sigma| < 1/\epsilon$, and $\Phi(0) = \Phi'(0) = 0$; at $\sigma = 1$ that bound yields
$|\Phi(1)| \le 2\epsilon^2\mathfrak{N}^V_p$. Since $\mathcal{U}{\restriction}Q$ reads as
$e^{i\eta\mathcal{A}}$ in one $G^{\mathbb{C}}$-gauge, $\Phi(1)$ is the $\sharp$-core
$\operatorname{Core}^{\sharp}_y(\operatorname{ext}_p(\mathcal{U});e)$. This proves the lemma.
\end{proof}

Lemma~\ref{10.3:lem-relocalisation-extended} has the following consequence, stated in the
notation of the proof of the correlation theorem. At the cores that this proof treats as smooth,
namely those off $\Lambda^{\natural}$ and in the collar, as fixed in that proof, the
$\sharp$-core has the size obtained there, with the factor $L^{-\hat\alpha b_v}$ replaced by
$L^{-2\hat\alpha b_v}$, where $b_v = p_v - j$. The cores treated there as rough are those with
$j > p_v = n - 2 - s^{(9/8)}(\bar{\mathsf{w}}_v)$. They keep the marginal bound $\nu^2 t^4$ at
$\nu \le 3 + s^{(9/8)}(\bar{\mathsf{w}}_v)$.

\begin{lemma}[Bounds on the bare vector unit]\label{10.3:lem-bare-unit}
We use the units of length, lattice differences and the weighted insertion fixed in
Notation~\ref{10.3:not-units-insertion}. Let $c_0\in\{g_0^{-2},1\}$ and put
\begin{equation}\label{10.3:eq-bare-unit-1}
  \mathsf{v}_0(c_0) := c_0\, C\, C_0^3\, g_0^3\, (\log x_0)^{3q_{\flat}},
\end{equation}
the bound, per unit weight, on the norm $\|\cdot\|^{\flat}$ of the correlation theorem of
$c_0\mathfrak{r}^{\sharp}$. For $\rho=1,\dots,4$ let $\hat\rho$ be the
lattice unit vector in direction $\rho$, and let $\mathfrak{r}^{\sharp}_{\rho}(p)$ be the bare
vector unit at the bare plaquette $p$. For a plaquette $p'$ with plaquette variable
$W = U(\partial p')$ put
\begin{equation}\label{10.3:eq-bare-unit-6}
  a_{p'} := -\tfrac12 \operatorname{tr}\bigl[(W-\mathbf{1})^2 W^{-1}\bigr].
\end{equation}
Then:
\begin{enumerate}
\item[(i)] on the domain $\mathcal{U}^{\flat}_0$ of bare configurations with
  $|U(\partial p')-\mathbf{1}|<\alpha_{0,0}$ at every plaquette $p'$, where
  $\alpha_{0,0} = g_0 C_0 (\log x_0)^{q_{\flat}}$, one has
  $\bigl|c_0\,\mathfrak{r}^{\sharp}_{\rho}(p)\bigr| \le \mathsf{v}_0(c_0)$;
\item[(ii)] on a background $\mathfrak{b}$ admissible at scale $n$, with $t = L^{-n}$,
  \begin{equation}\label{10.3:eq-bare-unit-2}
    \bigl|c_0\,\mathfrak{r}^{\sharp}_{\rho}(p)(\mathfrak{b})\bigr|
    \le C\,\mathsf{v}_0(c_0)\,\rho_{n,0}^3\, t^6 ;
  \end{equation}
\item[(iii)] on $SU(2)$-valued fields,
  \begin{equation}\label{10.3:eq-bare-unit-3}
    \bigl|\mathfrak{r}^{\sharp}_{\rho}(p)\bigr| \le C_{\mathfrak{r}} \sum_{p'} a_{p'},
  \end{equation}
  the sum running over the plaquettes $p'$ of the two cubes lying between
  $p-\hat\rho$, $p$ and $p+\hat\rho$.
\end{enumerate}
Here $C$ denotes constants independent of the configuration, of $g_0$ and of $n$, possibly
different at different occurrences; the constant in \eqref{10.3:eq-bare-unit-1} is that of the
bound $|\mathfrak{r}^{\sharp}_{\rho}(p)|\le C\alpha^3$ in the proof. At level $0$ no
re-localisation is made: the current at level $0$ is set to $\mathfrak{J}^{(0)} := 0$.
\end{lemma}

\begin{proof}
By part (ii) of the lemma giving the polynomial form of the bare vector unit, the components of
$\mathfrak{r}^{\sharp}_{\rho}(p)$ are
\begin{equation}\label{10.3:eq-bare-unit-4}
  E^a = \tfrac12\bigl[\sigma\,(a_p - a_{p+\hat\rho}) - T_3 - 2T_4\bigr],
\end{equation}
where the superscript $a$ is the component index (not to be confused with the bare spacing $a$
or with the numbers $a_{p'}$ of \eqref{10.3:eq-bare-unit-6}), and
$\sigma$, $T_3$, $T_4$ and the variables $\delta_i$ entering them are as in that lemma.
The expression \eqref{10.3:eq-bare-unit-4} is a polynomial in the link variables. Put
\begin{equation}\label{10.3:eq-bare-unit-5}
  \alpha := \max_{p'} \bigl|U(\partial p') - \mathbf{1}\bigr|,
\end{equation}
the maximum taken over the plaquettes $p'$ of the two cubes described in (iii). We estimate
the ingredients of \eqref{10.3:eq-bare-unit-4} in turn.

First, each variable $\delta_i$ satisfies $|\delta_i|\le C\alpha$.

Second, by \eqref{10.3:eq-bare-unit-6}, $|a_{p'}|\le C\alpha^2$ for every plaquette $p'$ of the
two cubes; this bound holds also when $W$ lies in the complexification $G^c$ of the group.

Third, with $T$, $q$ and $\delta$ as in the lemma cited,
\[
  |a_p - a_{p+\hat\rho}| = |T(q+\delta) - T(q)| \le C\alpha .
\]

From these bounds, $|E^a| \le C\alpha^3$, hence $|\mathfrak{r}^{\sharp}_{\rho}(p)|\le C\alpha^3$,
the number of components being absorbed into $C$.

(i) By the defining property of $\mathcal{U}^{\flat}_0$ in (i), $\alpha < \alpha_{0,0}$, and
therefore
\[
  c_0|\mathfrak{r}^{\sharp}_{\rho}(p)| \le c_0\, C\, \alpha_{0,0}^3
  = c_0\, C\, C_0^3\, g_0^3 (\log x_0)^{3q_{\flat}}
  = \mathsf{v}_0(c_0)
\]
by \eqref{10.3:eq-bare-unit-1}.

(ii) On a background $\mathfrak{b}$ admissible at scale $n$, the admissibility condition on the
background in the correlation theorem, read at level $j=0$, gives $\alpha \le \mathfrak{a}_n t^2$
with $t = L^{-n}$. Inserting this bound into
$c_0|\mathfrak{r}^{\sharp}_{\rho}(p)| \le c_0 C\alpha^3$ gives
\[
  c_0\bigl|\mathfrak{r}^{\sharp}_{\rho}(p)(\mathfrak{b})\bigr| \le c_0\, C\, \mathfrak{a}_n^3\, t^6,
\]
and hence \eqref{10.3:eq-bare-unit-2}.

(iii) On $SU(2)$-valued fields the bound \eqref{10.3:eq-bare-unit-3} follows from part (iii)
of the same lemma together with $\sigma \le 8$.
\end{proof}

For $K''\le K'$, where $K'=K-m$, let $\mathfrak{l}^{[K'']}_i$, $r^{[K'']}_i$ and $\mathcal{W}^{[K'']}_i(X)$ denote the Lipschitz room, the radius and the source class. These are the symbols so denoted in the definition of the source classes of the correlation theorem, with $K'$ replaced by $K''$. We say that these objects are taken \emph{at the anchor} $K''$. We write $\operatorname{plaq}(X^{+})$ for the set of plaquettes of $X^{+}$, and $r$ for the radius in terms of which the correlation theorem defines its source classes (distinct from its radius constants $\mathfrak{r}_j$). In particular the \emph{own-level class} is

\begin{equation}\label{10.3:eq-reanchoring-1}
\mathcal{W}^{\circ}_j(X):=\mathcal{W}^{[j]}_j(X)
=\Bigl\{w\colon \operatorname{plaq}(X^{+})\to\mathbb{C}\ :\ |w|<\tfrac54\,r,\ \
\operatorname{Lip}_j w<r\lambda_*\Bigr\}.
\end{equation}

For a function $\psi$ on $\operatorname{plaq}(X^{+})$ we write

\begin{equation*}
\|\psi\|_{j,X}:=\max\Bigl\{\sup_{X^{+}}|\psi|,\ \operatorname{Lip}_j\psi/\lambda_*\Bigr\}.
\end{equation*}

Here $\operatorname{Lip}_j$ is the Lipschitz seminorm at level $j$ that enters the definition of the
source classes of the correlation theorem.

\begin{proposition}[Re-anchoring the source classes at each level]\label{10.3:prop-reanchoring}
Assume the following four parts of the correlation theorem.
\begin{enumerate}
\item[(a)] Its parameter sentence, together with one property of its setting. The theorem is
asserted, in particular, for every $g\le\gamma_J$, every
$x_0=g_0^{-2}>g^{-2}$ with $K=K(g_0,g)$ the first passage, every $m$ with
$m_{\mathbb{T}}(g)\le m\le K$ (that is, with the lower torus condition
$L^m\ge N_{\mathbb{T}}(g,m)$ of the setting of the correlation theorem), and every integer
$N_T\ge1$. Here $N_T$ is the
integer by which the spacing $L^{-K}$ is further divided: the bare spacing is $L^{-K}N_T^{-1}$
in the units of the correlation theorem. In that setting, moreover,
$N_{\mathbb{T}}(g,n+1)\le L\,N_{\mathbb{T}}(g,n)$ for every $n$.
\item[(b)] The sentence that $K$, $k$, $m$, the torus and the position occur in no constant of the theorem.
\item[(c)] Its clause on the sourced format. At each admissible $m$, the
sourced format $(\mathrm{IH}^w)_k$ holds for every $k\le K-m$ and every
$w\in\mathcal{W}^{[K-m]}_k$.
\item[(d)] Its clauses on the terms of level $j$. At each admissible $m$ and for every
$j\le K-m$, each term of level $j$ is defined on
$\mathcal{W}^{[K-m]}_j(X)$ by its production formula, is analytic there and depends only on
the restriction $w{\restriction}X^{+}$. Under the condition under which
Ba{\l}aban's bounds for the terms of its sourced format are stated, the bounds of the
format hold for each such term uniformly on
$\mathcal{W}^{[K-m]}_j(X)$.
\end{enumerate}
Then the following holds for every $j\le K''\le K'$. The terms of level $j$ of the sourced format $(\mathrm{IH}^w)$, namely
\[
E^{(j)}(X,\cdot,y;w),\qquad R^{(j)}(X,\cdot;w),\qquad \mathbf{B}^{(j)},
\]
are analytic on $\mathcal{W}^{[K'']}_j(X)$. They satisfy there Ba{\l}aban's bounds, under the
condition of hypothesis (d), with the constants of the correlation
theorem. In particular each of them is analytic and bounded on the own-level class $\mathcal{W}^{\circ}_j(X)$.
\end{proposition}

\begin{proof}
Fix $g$, $x_0$, $K=K(g_0,g)$, $m$ and $N_T$ as in hypothesis (a), and fix $j\le K''\le K'$. Put
\[
m'':=K-K''.
\]

First, $m''$ is an admissible torus exponent. Since $K''\le K'=K-m$, we have $m''\ge m\ge m_{\mathbb{T}}(g)$. Since $K''\ge j\ge0$, we have $m''\le K$. Hence $m_{\mathbb{T}}(g)\le m''\le K$, and we use the same $N_T\ge1$.

Second, passing from $m$ to $m''$ changes none of the data that enter the terms of level at
most $K''$. The bare lattice, of spacing $L^{-K}N_T^{-1}$ in the units of the correlation
theorem, is the same. So are the
sourced initial density $\rho_0^w$ and the sourced renormalisation operations $R^{\flat}T^{w}$
of the correlation theorem. The reason is that neither the scaffold of the construction nor the
production formulas of the levels at most $K''$ depend on $m$.

Third, the lower torus condition holds at $m''$. It holds at $m$ by hypothesis (a), that is, $L^m\ge N_{\mathbb{T}}(g,m)$. By hypothesis (a) one also has
\[
N_{\mathbb{T}}(g,n+1)\le L\,N_{\mathbb{T}}(g,n)\qquad\text{for every } n.
\]
Suppose $L^n\ge N_{\mathbb{T}}(g,n)$ for some $n$ with $m\le n<m''$. Then
\begin{equation*}
L^{n+1}=L\cdot L^n\ge L\,N_{\mathbb{T}}(g,n)\ge N_{\mathbb{T}}(g,n+1).
\end{equation*}
Starting from $n=m$, induction on $n$ gives $L^{m''}\ge N_{\mathbb{T}}(g,m'')$.

Fourth, we apply hypothesis (c) at $m''$. Since $K-m''=K''$, it gives the sourced format
$(\mathrm{IH}^w)_k$ for every $k\le K''$ and every $w\in\mathcal{W}^{[K'']}_k$. The source
classes of the run at $m''$ are, by definition, the classes at the anchor $K''$. By hypothesis
(b), the constants of the run at $m''$ are those of the run at $m$, since $m$ occurs in no
constant. Hence the format $(\mathrm{IH}^w)_k$ holds for every $k\le K''$ on the classes at
the anchor $K''$, with the constants of the correlation theorem.

Fifth, we apply hypothesis (d) at $m''$ and level $j\le K''$. Each of
$E^{(j)}(X,\cdot,y;w)$, $R^{(j)}(X,\cdot;w)$ and $\mathbf{B}^{(j)}$ is defined on
$\mathcal{W}^{[K'']}_j(X)$ by its production formula, is analytic there and depends only on
$w{\restriction}X^{+}$. Under the condition of hypothesis (d), the bounds of the format hold for it
uniformly on this class.

Sixth, we treat the own-level class. At the anchor $K''=j$ one has
\[
\mathfrak{l}^{[j]}_j=r\lambda_*,\qquad r^{[j]}_j=r\bigl(1+\tfrac14\bigr)=\tfrac54\,r,
\]
so $\mathcal{W}^{[j]}_j(X)$ is the class displayed in \eqref{10.3:eq-reanchoring-1}. The case $K''=j$ of the preceding steps therefore gives the last sentence of the proposition.

Finally, the determinations at different anchors define a single function. Let $K''_1$ and
$K''_2$ be two anchors with $j\le K''_1,K''_2\le K'$. Each class
$\mathcal{W}^{[K''_1]}_j(X)$ and $\mathcal{W}^{[K''_2]}_j(X)$ is defined by a strict bound on
$\sup_{X^{+}}|w|$ and a strict bound on $\operatorname{Lip}_j w$. Both are seminorms on
$\mathbb{C}^{\operatorname{plaq}(X^{+})}$ and the
bounds are strict, so each class is an open convex neighbourhood of $0$ in
$\mathbb{C}^{\operatorname{plaq}(X^{+})}$. Their intersection is therefore open, convex and
non-empty, hence connected. On this intersection the two determinations of a given term are produced by the same production formula, so they agree there. Consequently each term is one analytic function on the union of the classes, and on each class it satisfies the bounds stated there.
\end{proof}

\begin{remark}
In what follows, statements about the run of the correlation theorem at $(g,x_0,K,m,N_T)$ are
to be read as follows. A statement said to hold \emph{at the anchor} $K''$ means that statement
read at the torus exponent $m''=K-K''$. This includes the form of the one-lattice Lipschitz
step that concerns this run, the re-filing lemmas and the formats of this section.

The Lipschitz room $\mathfrak{l}^{[K']}_j$ of the correlation theorem is proportional to
$L^{-(K'-j)/2}$. Whether this decay is forced or a choice is immaterial for the terms of level
$j$: the exponent $\tfrac12$ measures the room that a level-$j$ term is given inside the one
induction that must reach $K'$ with summable radii, while the term itself has the room
$r\lambda_*$ at its own level.
\end{remark}

\begin{lemma}[Own-level freezing to second order]\label{10.3:lem-own-level-freezing}
Let $G(X;\cdot)$ be analytic on the own-level class $\mathcal{W}^{\circ}_j(X)$ and satisfy
$|G(X;\cdot)|\le\mathfrak{n}$ there. For $\hat w\in\mathcal{W}^{\circ}_j(X)$ and a function $\psi$ on the
plaquettes of $X^+$, write $DG(X;\hat w)[\psi]$ for the derivative of $G(X;\cdot)$ at $\hat w$ in the
direction $\psi$.
\begin{enumerate}
\item[(i)] For every constant $\hat w\in\mathbb{C}$ with $|\hat w|\le r/4$ and every function $\psi$ on
the plaquettes of $X^+$,
\begin{equation}\label{10.3:eq-own-level-freezing-1}
|DG(X;\hat w)[\psi]|\le \mathfrak{n}\,\|\psi\|_{j,X}/r .
\end{equation}
\item[(ii)] Let $\varepsilon\le 1$, and let $w$ satisfy on $X^+$
\begin{equation*}
|w|\le \frac r4,\qquad \operatorname{Lip}_j w\le \frac{\varepsilon r\lambda_*}{4},\qquad
|\nabla_j\nabla_j w|\le \frac{\varepsilon^2 r\lambda_*^2}{4}.
\end{equation*}
Let $c\in X^+$, put $\hat w:=w(c)$, and let $\ell$ be the first-order part of the Taylor polynomial
of $w$ at $c$, that is, the function on the plaquettes $p$ of $X^+$ given by
\begin{equation*}
\ell(p):=\sum_{\nu}a_j^{-1}\bigl(x(p)-x(c)\bigr)_{\nu}\,\nabla_{j,\nu}w(c),
\end{equation*}
so that $\hat w+\ell$ is the first-order lattice Taylor polynomial of $w$ at $c$. Then
\begin{equation}\label{10.3:eq-own-level-freezing-2}
|G(X;w)-G(X;\hat w)|\le \mathfrak{n}\,\varepsilon\,(1+\lambda_*D_X)
\end{equation}
and
\begin{equation}\label{10.3:eq-own-level-freezing-3}
\bigl|G(X;w)-G(X;\hat w)-DG(X;\hat w)[\ell]\bigr|\le \mathfrak{n}\,\varepsilon^2\,(1+\lambda_*D_X)^2 .
\end{equation}
\end{enumerate}
For sources $w$ in the polydisc $\mathbb{D}_2$, (ii) applies with $\varepsilon=a_j=L^{-(K'-j)}$.
\end{lemma}

\begin{proof}
Recall that $\mathcal{W}^{\circ}_j(X)=\mathcal{W}^{[j]}_j(X)$ consists of the functions $u$ on the
plaquettes of $X^+$ with $|u|<5r/4$ and $\operatorname{Lip}_j u<r\lambda_*$, $r$ being the radius of
this own-level class, and that
$\|\psi\|_{j,X}=\max\{\sup_{X^+}|\psi|,\ \operatorname{Lip}_j\psi/\lambda_*\}$.

(i) Let $\sigma\in\mathbb{C}$ with $|\sigma|\,\|\psi\|_{j,X}<r$. Since $\hat w$ is constant,
\begin{align*}
|\hat w+\sigma\psi|&\le \frac r4+|\sigma|\sup_{X^+}|\psi|<\frac r4+r=\frac{5r}4,\\
\operatorname{Lip}_j(\hat w+\sigma\psi)&=|\sigma|\operatorname{Lip}_j\psi
\le|\sigma|\lambda_*\|\psi\|_{j,X}<r\lambda_* .
\end{align*}
Hence $\hat w+\sigma\psi\in\mathcal{W}^{\circ}_j(X)$ for $|\sigma|<r/\|\psi\|_{j,X}$ (for all $\sigma$ if
$\psi=0$, when \eqref{10.3:eq-own-level-freezing-1} is trivial). The function
$\sigma\mapsto G(X;\hat w+\sigma\psi)$ is analytic and bounded by $\mathfrak{n}$ on this disc, and its
derivative at $\sigma=0$ is $DG(X;\hat w)[\psi]$. The Cauchy estimate on the disc gives
\eqref{10.3:eq-own-level-freezing-1}.

(ii) The function $w-\hat w$ vanishes at $c\in X^+$ and $\operatorname{Lip}_j(w-\hat w)=\operatorname{Lip}_jw$.
Hence $\sup_{X^+}|w-\hat w|\le(\varepsilon r\lambda_*/4)\operatorname{diam}_j(X^+)$ and
$\operatorname{Lip}_j(w-\hat w)/\lambda_*\le\varepsilon r/4$. Since $D_X=\operatorname{diam}_j(X^+)+1$,
\begin{equation*}
\|w-\hat w\|_{j,X}\le \frac{\varepsilon r}{4}\,(1+\lambda_*D_X)=:\frac rS,
\qquad S=\frac{4}{\varepsilon(1+\lambda_*D_X)} .
\end{equation*}
Put $f(\sigma):=G(X;\hat w+\sigma(w-\hat w))$. We have $|\hat w|=|w(c)|\le r/4$, so the argument of
(i) with $\psi=w-\hat w$ shows that $f$ is analytic and bounded by $\mathfrak{n}$ on $|\sigma|<S$.
Moreover $f(1)=G(X;w)$, $f(0)=G(X;\hat w)$ and $f'(0)=DG(X;\hat w)[w-\hat w]$.

Suppose first $S\ge2$. The Cauchy estimates on $|\sigma|<S$ bound the Taylor coefficients of $f$ at
$0$ by $\mathfrak{n}S^{-n}$. Summing the Taylor series at $\sigma=1$ gives
\begin{align*}
|f(1)-f(0)|&\le \mathfrak{n}\sum_{n\ge1}S^{-n}=\frac{\mathfrak{n}S^{-1}}{1-S^{-1}}\le\frac{2\mathfrak{n}}S,\\
|f(1)-f(0)-f'(0)|&\le \mathfrak{n}\sum_{n\ge2}S^{-n}=\frac{\mathfrak{n}S^{-2}}{1-S^{-1}}
\le 2\mathfrak{n}S^{-2}.
\end{align*}
Suppose next $S<2$. For $|\sigma|<2$ we have
\begin{align*}
|\hat w+\sigma(w-\hat w)|&\le \frac r4+|\sigma|\,\frac r2<\frac{5r}4,\\
\operatorname{Lip}_j(\hat w+\sigma(w-\hat w))&=|\sigma|\operatorname{Lip}_jw
\le|\sigma|\,\frac{r\lambda_*}4<r\lambda_*,
\end{align*}
because $|w|,|\hat w|\le r/4$ and $\varepsilon\le1$. So $f$ is analytic and bounded by $\mathfrak{n}$ on
$|\sigma|<2$. The Cauchy estimate then gives $|f'(0)|\le\mathfrak{n}/2$. Hence
\begin{align*}
|f(1)-f(0)|&\le 2\mathfrak{n}\le \frac{4\mathfrak{n}}S,\\
|f(1)-f(0)-f'(0)|&\le 3\mathfrak{n}\le 12\mathfrak{n}S^{-2}.
\end{align*}
In both cases $|f(1)-f(0)|\le4\mathfrak{n}/S=\mathfrak{n}\varepsilon(1+\lambda_*D_X)$, which is
\eqref{10.3:eq-own-level-freezing-2}. Also, in both cases,
\begin{equation*}
\bigl|G(X;w)-G(X;\hat w)-DG(X;\hat w)[w-\hat w]\bigr|\le12\mathfrak{n}S^{-2}
=\frac{12}{16}\,\mathfrak{n}\varepsilon^2(1+\lambda_*D_X)^2 .
\end{equation*}

By the definition of $\ell$, $\ell(c)=0$ and $\nabla_{j,\nu}\ell=\nabla_{j,\nu}w(c)$ everywhere, for
every $\nu$; hence $\nabla_j\nabla_j\ell=0$. Therefore $(w-\hat w-\ell)(c)=0$,
$\nabla_{j,\nu}(w-\hat w-\ell)(c)=0$ for every $\nu$, and
$\nabla_j\nabla_j(w-\hat w-\ell)=\nabla_j\nabla_jw$. Summing second differences along lattice paths in
$X^+$ from $c$ and using the bound on $\nabla_j\nabla_jw$ gives
\begin{align*}
\operatorname{Lip}_j(w-\hat w-\ell)&\le\frac{\varepsilon^2r\lambda_*^2}{4}\operatorname{diam}_j(X^+),\\
\sup_{X^+}|w-\hat w-\ell|&\le\frac{\varepsilon^2r\lambda_*^2}{4}\operatorname{diam}_j(X^+)^2 .
\end{align*}
Since $\lambda_*\operatorname{diam}_j(X^+)\le 1+\lambda_*D_X$ and $1+\lambda_*D_X\ge1$, this gives
\begin{equation*}
\|w-\hat w-\ell\|_{j,X}\le\frac{\varepsilon^2r}{4}\,(1+\lambda_*D_X)^2 .
\end{equation*}
By linearity of $DG(X;\hat w)$ and by (i), applied with $\psi=w-\hat w-\ell$ (admissible since
$|\hat w|\le r/4$),
\begin{equation*}
\bigl|DG(X;\hat w)[w-\hat w]-DG(X;\hat w)[\ell]\bigr|
\le\frac{\mathfrak{n}}{r}\,\|w-\hat w-\ell\|_{j,X}
\le\frac14\,\mathfrak{n}\varepsilon^2(1+\lambda_*D_X)^2 .
\end{equation*}
Adding the last two bounds and using $12/16+1/4=1$ gives \eqref{10.3:eq-own-level-freezing-3}.
\end{proof}

In what follows we use the constant
\begin{equation}\label{10.3:eq-own-level-freezing-4}
C_B':=\sup_X\,(1+\lambda_*D_X)^2\,e^{-\frac14\delta\kappa d_j(X)} ,
\end{equation}
where the supremum is over the domains $X$ at level $j$.

\begin{proposition}[Summation by parts at the level of birth]\label{10.3:prop-summation-by-parts}
Let $D(0,2r)\subset\mathbb{C}$ denote the open disc of radius $2r$ about $0$, where $r$ is the radius
fixed in the setting of this section, and let
$\{\mathfrak{V}_{\hat w}\}_{\hat w\in D(0,2r)}$ be an analytic family of $V$-typed pointed
$E$-families at scale $j$. Write $\mathfrak{J}_{\hat w}:=\mathfrak{J}[\mathfrak{V}_{\hat w}]$,
and let $\mathfrak{V}^{\sharp}_{\hat w}$ and $\mathsf{S}[\mathfrak{V}_{\hat w}]$ be the sharp part
and the symmetrised current family of $\mathfrak{V}_{\hat w}$, as in
Definition~\ref{10.3:def-bond-current} and Lemma~\ref{10.3:lem-relocalisation}. Let $\Lambda$ be a
finite set of sites of $T^{(j)}$ and $\Lambda^+\supset\Lambda$ a region. All sums over constituents
$X$ are restricted to $X\subset\Lambda^+$; since this restriction reads $X$ only, a constituent of a
current enters with both of its weights or not at all. For a site $y$ and a direction $\nu$ put
$y^{\pm}:=y\pm a_je_\nu$, and let $b_\nu(y)$ denote the bond of $T^{(j)}$ from $y$ to $y^+$. Put
$\mathsf{k}:=\partial\eta$ and $\mathsf{k}':=\eta\otimes\partial w$, that is,
$\mathsf{k}_{\nu\rho}=\partial_\nu\eta^\rho$ and $\mathsf{k}'_{\nu\rho}=\eta^\rho\partial_\nu w$;
$\mathsf{k}_\nu$, $\mathsf{k}'_\nu$ are the $V$-valued functions with these components.

\emph{(a)} For every background on which the terms are defined, identically in $(w,\mathsf{h})$,
\begin{equation}\label{10.3:eq-summation-by-parts-1}
\begin{split}
\sum_{y\in\Lambda}\mathfrak{V}_{\hat w(y)}(\cdot,y;\mathsf{h}(y))
&=\sum_{y\in\Lambda}\mathfrak{V}^{\sharp}_{\hat w(y)}(\cdot,y;\mathsf{h}(y))\\
&\quad+\sum_{y}\bigl[\mathsf{S}^{a,\partial\eta}_j+\mathsf{S}^{a,\eta\partial w}_j
+\mathsf{D}^a_j\bigr](\cdot,y)+\mathsf{rim}_j(\Lambda).
\end{split}
\end{equation}
Here, for a bond $b=b_\nu(y)$ with both ends in $\Lambda$, with
\begin{equation*}
\mathfrak{J}'_b:=\int_0^1\partial_{\hat w}\mathfrak{J}_{\hat w(y)+s(\hat w(y^+)-\hat w(y))}\,ds,
\end{equation*}
the bond terms are
\begin{equation}\label{10.3:eq-summation-by-parts-2}
\begin{aligned}
\mathsf{S}^{a,\partial\eta}_j(b)&:=-L^j\mathfrak{J}_{\hat w(y)}\bigl(b;\langle\mathsf{k}_\nu\rangle_{j,\nu}(y)\bigr),\\
\mathsf{S}^{a,\eta\partial w}_j(b)&:=-L^j\mathfrak{J}'_b\bigl(b;\langle\mathsf{k}'_\nu\rangle_{j,\nu}(y)\bigr),\\
\mathsf{D}^a_j(b)&:=-L^j\mathfrak{J}'_b\bigl(b;\langle(\eta(y^+)-\eta)\partial_\nu w\rangle_{j,\nu}(y)\bigr),
\end{aligned}
\end{equation}
and each bond term is pinned half at each end of its bond. We use the same letter for a bond term
and for the pointed family obtained by pinning it:
\begin{equation*}
\mathsf{S}^{a,\partial\eta}_j(\cdot,y):=\tfrac12\sum_{b\ni y}\mathsf{S}^{a,\partial\eta}_j(b),
\end{equation*}
the sum being over the bonds $b$ with both ends in $\Lambda$ of which $y$ is an end, and likewise for
$\mathsf{S}^{a,\eta\partial w}_j$ and $\mathsf{D}^a_j$. The rim is
\begin{equation}\label{10.3:eq-summation-by-parts-3}
\mathsf{rim}_j(\Lambda):=\sum_{b_\nu(y):\,y\in\Lambda,\,y^+\notin\Lambda}
\mathfrak{J}_{\hat w(y)}(b;\mathsf{h}(y))
-\sum_{b_\nu(y):\,y\notin\Lambda,\,y^+\in\Lambda}\mathfrak{J}_{\hat w(y^+)}(b;\mathsf{h}(y^+)).
\end{equation}

\emph{(b)} Let $\mathsf{u}:=\sup_{\hat w}\|\mathfrak{V}_{\hat w}\|^{\mathrm{fr}}$ and
$\mathsf{m}':=\sup_{|\hat w|\le r/4}\|\partial_{\hat w}\mathfrak{V}_{\hat w}\|^{\mathrm{fr}}$. Then
$\mathsf{S}^{a,\partial\eta}_j$ is linear in $\mathsf{k}{\restriction}X^+$ and analytic in $w$; it is
$(V\otimes V)$-typed, and where $(w,\mathsf{k})\equiv(\hat w,\mathsf{k}_0)$ are constant it equals
$L^j\mathsf{S}[\mathfrak{V}_{\hat w}](X,\cdot,y;\mathsf{k}_0)$. The family
$\mathsf{S}^{a,\eta\partial w}_j$ is linear in $\mathsf{k}'{\restriction}X^+$ with coefficients
analytic in $w$, and where $(w,\mathsf{k}')\equiv(\hat w,\mathsf{k}'_0)$ are constant it equals
$L^j\mathsf{S}[\partial_{\hat w}\mathfrak{V}_{\hat w}](X,\cdot,y;\mathsf{k}'_0)$. Moreover, the first
of the following bounds holds in general, and the second and third hold whenever
$|\hat w(y)|\le r/4$ at every site $y$ of $\Lambda$ and $|\partial_\nu w|\le\lambda_* r/4$
everywhere, for every $\nu$:
\begin{equation}\label{10.3:eq-summation-by-parts-4}
\begin{aligned}
\|\mathsf{S}^{a,\partial\eta}_j\|^{\flat}&\le 4C_{\mathfrak{J}}L^j\mathsf{u}\sup|\mathsf{k}|,\\
\|\mathsf{S}^{a,\eta\partial w}_j\|^{\flat}&\le 4C_{\mathfrak{J}}L^j\mathsf{m}'\sup|\mathsf{k}'|,\\
\|\mathsf{D}^a_j\|^{\flat}&\le 4C_{\mathfrak{J}}L^j\mathsf{m}'\cdot a_j\lambda_*\|\eta\|_{\langle 2\rangle}
\cdot\lambda_* r/4 .
\end{aligned}
\end{equation}
Each rim term is a constituent of $\mathfrak{J}$ with the undifferentiated weight, bounded by
$C_{\mathfrak{J}}\mathsf{u}\sup|\mathsf{h}|\,e^{-(15/8)\kappa d_j(X)}$, and pinned within one step of
$\partial\Lambda$. All the norms in (b) hold at the rate $\kappa_{\sharp}$.
\end{proposition}

\begin{proof}
Throughout, the configuration argument is suppressed, and $\mathfrak{J}_{\hat w}(b;e)$ denotes the
pointed total of the current at the bond $b$ with weight $e$, summed over the constituents
$X\subset\Lambda^+$. Every identity below holds term by term in $X$; by the remark on the restriction
$X\subset\Lambda^+$ in the statement, the two terms of a current that are paired at a bond carry the
same constituent, so the restriction is compatible with every regrouping made.

\emph{(a)} Fix $y\in\Lambda$. Part (b) of Lemma~\ref{10.3:lem-relocalisation}, applied to
$F=\mathfrak{V}_{\hat w(y)}$ with weight $e=\mathsf{h}(y)$, gives
\begin{equation}\label{10.3:eq-summation-by-parts-5}
\mathfrak{V}_{\hat w(y)}(\cdot,y;\mathsf{h}(y))=\mathfrak{V}^{\sharp}_{\hat w(y)}(\cdot,y;\mathsf{h}(y))
+\sum_\nu\bigl[\mathfrak{J}_{\hat w(y)}(b_\nu(y);\mathsf{h}(y))
-\mathfrak{J}_{\hat w(y)}(b_\nu(y^-);\mathsf{h}(y))\bigr].
\end{equation}
We sum \eqref{10.3:eq-summation-by-parts-5} over $y\in\Lambda$ and regroup the finite double sum of
currents by bonds. A bond $b=b_\nu(z)$ receives $+\mathfrak{J}_{\hat w(z)}(b;\mathsf{h}(z))$ from the
site $y=z$ if $z\in\Lambda$, and $-\mathfrak{J}_{\hat w(z^+)}(b;\mathsf{h}(z^+))$ from the site
$y=z^+$ (for which $b=b_\nu(y^-)$) if $z^+\in\Lambda$. The bonds with exactly one end in $\Lambda$
therefore contribute precisely the two sums of \eqref{10.3:eq-summation-by-parts-3}, that is,
$\mathsf{rim}_j(\Lambda)$. A bond $b=b_\nu(y)$ with both ends in $\Lambda$ carries
$\mathfrak{J}_{\hat w(y)}(b;\mathsf{h}(y))-\mathfrak{J}_{\hat w(y^+)}(b;\mathsf{h}(y^+))$. Since the
current is linear in its weight, adding and subtracting $\mathfrak{J}_{\hat w(y)}(b;\mathsf{h}(y^+))$
gives
\begin{equation}\label{10.3:eq-summation-by-parts-6}
\begin{split}
&\mathfrak{J}_{\hat w(y)}(b;\mathsf{h}(y))-\mathfrak{J}_{\hat w(y^+)}(b;\mathsf{h}(y^+))\\
&\qquad=-\mathfrak{J}_{\hat w(y)}\bigl(b;\nabla_{j,\nu}\mathsf{h}(y)\bigr)
-\bigl[\mathfrak{J}_{\hat w(y^+)}-\mathfrak{J}_{\hat w(y)}\bigr](b;\mathsf{h}(y^+)).
\end{split}
\end{equation}
In the first term, $\mathsf{h}=\eta/a$ and $a_j=L^ja$, so the identity
$\nabla_{j,\nu}\psi=a_j\langle\partial_\nu\psi\rangle_{j,\nu}$ applied to $\psi=\mathsf{h}$ gives
$\nabla_{j,\nu}\mathsf{h}=L^j\langle\partial_\nu\eta\rangle_{j,\nu}=L^j\langle\mathsf{k}_\nu\rangle_{j,\nu}$;
the first term is the bond term $\mathsf{S}^{a,\partial\eta}_j(b)$. In the second term, the family
$\hat w\mapsto\mathfrak{J}_{\hat w}$ is analytic (Lemma~\ref{10.3:lem-relocalisation}, part (a)), so
by the fundamental theorem of calculus along the segment from $\hat w(y)$ to $\hat w(y^+)$,
\begin{equation*}
\mathfrak{J}_{\hat w(y^+)}-\mathfrak{J}_{\hat w(y)}=\bigl(\hat w(y^+)-\hat w(y)\bigr)\mathfrak{J}'_b .
\end{equation*}
Now $\hat w(y^+)-\hat w(y)=a_j\langle\partial_\nu w\rangle_{j,\nu}(y)$ and
$a_j\mathsf{h}(y^+)=L^j\eta(y^+)$, so by linearity in the weight the second term of
\eqref{10.3:eq-summation-by-parts-6} equals
$-L^j\mathfrak{J}'_b\bigl(b;\eta(y^+)\langle\partial_\nu w\rangle_{j,\nu}(y)\bigr)$. As $\eta(y^+)$ is
constant in the variable of the average, and the average is linear,
\begin{equation*}
\eta(y^+)\langle\partial_\nu w\rangle_{j,\nu}(y)=\langle\eta\,\partial_\nu w\rangle_{j,\nu}(y)
+\langle(\eta(y^+)-\eta)\partial_\nu w\rangle_{j,\nu}(y),
\end{equation*}
and $\langle\eta\,\partial_\nu w\rangle_{j,\nu}=\langle\mathsf{k}'_\nu\rangle_{j,\nu}$. Hence the
second term of \eqref{10.3:eq-summation-by-parts-6} is
$\mathsf{S}^{a,\eta\partial w}_j(b)+\mathsf{D}^a_j(b)$, and a bond with both ends in $\Lambda$ carries
the sum of its three bond terms \eqref{10.3:eq-summation-by-parts-2}. Splitting each bond term into two
halves pinned at the two ends of its bond, both of which lie in $\Lambda$, the sum over these bonds is
the second sum of \eqref{10.3:eq-summation-by-parts-1}. Collecting the sharp parts from
\eqref{10.3:eq-summation-by-parts-5}, the bonds with both ends in $\Lambda$, and the rim gives
\eqref{10.3:eq-summation-by-parts-1}.

\emph{(b)} By part (a) of Lemma~\ref{10.3:lem-relocalisation}, $\mathfrak{J}[F]$ is linear in $F$ and
analytic in every parameter in which $F$ is, and the averages in \eqref{10.3:eq-summation-by-parts-2}
are linear in the functions averaged. Hence $\mathsf{S}^{a,\partial\eta}_j$ is linear in
$\mathsf{k}{\restriction}X^+$ and analytic in $w$ (through $\hat w(y)$), and
$\mathsf{S}^{a,\eta\partial w}_j$ is linear in $\mathsf{k}'{\restriction}X^+$ with coefficients
$\mathfrak{J}'_b$ analytic in $w$. Where $(w,\mathsf{k})\equiv(\hat w,\mathsf{k}_0)$ are constant, the
averages reduce to the constant weights, and by the definition of the symmetrised current family
$\mathsf{S}[F]$ (Definition~\ref{10.3:def-bond-current} and Lemma~\ref{10.3:lem-relocalisation}) the
bond terms pinned half at each end are those of the family
$L^j\mathsf{S}[\mathfrak{V}_{\hat w}](X,\cdot,y;\mathsf{k}_0)$, which is $(V\otimes V)$-typed by part
(a) of Lemma~\ref{10.3:lem-relocalisation}; likewise,
at constants, $\mathfrak{J}'_b=\partial_{\hat w}\mathfrak{J}_{\hat w}=\mathfrak{J}[\partial_{\hat w}
\mathfrak{V}_{\hat w}]$ by linearity of $\mathfrak{J}[F]$ in $F$, and $\mathsf{S}^{a,\eta\partial w}_j$
equals $L^j\mathsf{S}[\partial_{\hat w}\mathfrak{V}_{\hat w}](X,\cdot,y;\mathsf{k}'_0)$.

The bounds \eqref{10.3:eq-summation-by-parts-4} follow from the bound
$\|\mathsf{S}[F]\|^{\flat}\le 4C_{\mathfrak{J}}\mathfrak{n}$ of part (a) of
Lemma~\ref{10.3:lem-relocalisation}, at the rate $\kappa_{\sharp}$. For
$\mathsf{S}^{a,\partial\eta}_j$ it is applied with $F=\mathfrak{V}_{\hat w}$ and $\mathfrak{n}=\mathsf{u}$,
the weight being an average of $\mathsf{k}$ and so of size at most $\sup|\mathsf{k}|$. For
$\mathsf{S}^{a,\eta\partial w}_j$ it is applied with $F=\partial_{\hat w}\mathfrak{V}_{\hat w}$ at the
points of the segment from $\hat w(y)$ to $\hat w(y^+)$, which lies in $\{|\hat w|\le r/4\}$ by the
condition imposed in (b) and the convexity of the disc, so that
$\mathfrak{n}=\mathsf{m}'$, the weight being of size at most $\sup|\mathsf{k}'|$. For
$\mathsf{D}^a_j$ it is applied with the same $F$ and $\mathfrak{n}=\mathsf{m}'$; in the weight
$(\eta(y^+)-\eta)\partial_\nu w$, averaged over the segment from $y$ to $y^+$ of length $a_j$, the first
factor is at most $a_j\sup|\partial\eta|\le a_j\lambda_*\|\eta\|_{\langle 2\rangle}$ by the definition
of $\|\cdot\|_{\langle 2\rangle}$, and the second factor is bounded by $\lambda_* r/4$ by the condition
$|\partial_\nu w|\le\lambda_* r/4$ imposed in (b). Finally, each rim term in
\eqref{10.3:eq-summation-by-parts-3} is a single constituent
of $\mathfrak{J}[\mathfrak{V}_{\hat w}]$ at the weight $\mathsf{h}(y)$ or $\mathsf{h}(y^+)$, not
differentiated in $\hat w$; the bound
\begin{equation*}
\sup e^{(15/8)\kappa d_j(X)}|\mathfrak{J}[F](X,\cdot;b;e)|\le C_{\mathfrak{J}}\mathfrak{n}|e|
\end{equation*}
of part (a) of Lemma~\ref{10.3:lem-relocalisation} with $\mathfrak{n}=\mathsf{u}$ gives the stated
bound, and since its bond has exactly one end in $\Lambda$, it is pinned within one step of
$\partial\Lambda$.
\end{proof}

\begin{proposition}[Decomposition of a local output at its pin]\label{10.3:prop-output-decomposition}
In the units of length of Notation~\ref{10.3:not-units-insertion}, so that $a=L^{-K'}$,
$a_j=L^ja=L^{-(K'-j)}$ and $\lambda_*=(40M)^{-1}$,
let $y\in X^+$ be the pin and let $\mathsf{O}(X,\cdot,y;w,\mathsf{h})$ be linear in the
bounded function $\mathsf{h}{\restriction}X^+$, no regularity of $\mathsf{h}$ being required, and analytic in
$w\in\mathcal{W}^{\circ}_j(X)$, with
\begin{equation}\label{10.3:eq-output-decomposition-1}
|\mathsf{O}(X,\cdot,y;w,\mathsf{h})|\le \mathsf{o}\,\sup_{X^+}|\mathsf{h}|\,e^{-2\kappa d_j(X)},
\end{equation}
and such that, for every $w_0$ with $|w_0|\le r/4$, where $r$ is the radius of the own-level class
$\mathcal{W}^{\circ}_j(X)$ as in Lemma~\ref{10.3:lem-own-level-freezing}, every constant weight $e$ and
every function $\psi$ on $X^+$,
\begin{equation}\label{10.3:eq-output-decomposition-2}
|D_w\mathsf{O}(X,\cdot,y;w_0,e)[\psi]|\le \mathsf{m}\,|e|\,\|\psi\|_{j,X}\,e^{-2\kappa d_j(X)} .
\end{equation}
By own-level freezing, Lemma~\ref{10.3:lem-own-level-freezing}\,(i), applied to
$w\mapsto\mathsf{O}(X,\cdot,y;w,e)$, the bound \eqref{10.3:eq-output-decomposition-2} always holds with
$\mathsf{m}=\mathsf{o}/r$. Put $\hat w:=\hat w(y)$, measure $x-y$ in units of $a_j$, let
$\mathsf{h}=\eta/a$ be the weight, and let $\mathsf{k}:=\partial\eta$ and
$\mathsf{k}':=\eta\otimes\partial w$ (that is, $\mathsf{k}_{\nu\rho}=\partial_\nu\eta^\rho$ and
$\mathsf{k}'_{\nu\rho}=\eta^\rho\partial_\nu w$) as in Proposition~\ref{10.3:prop-summation-by-parts}.
Define
\begin{align*}
\mathsf{S}^{b,\partial\eta}_j(X,\cdot,y;w,\mathsf{k})
 &:=L^j\,\mathsf{O}\Bigl(X,\cdot,y;\hat w,\ x\mapsto\sum_\nu(x-y)_\nu\,\mathsf{k}_\nu(y)\Bigr),\\
\mathsf{S}^{b,\eta\partial w}_j(X,\cdot,y;w,\mathsf{k}')
 &:=L^j\sum_\nu D_w\mathsf{O}\bigl(X,\cdot,y;\hat w,\mathsf{k}'_\nu(y)\bigr)\bigl[(x-y)_\nu\bigr].
\end{align*}
Then, writing $\mathsf{O}(w,\mathsf{h})$ for $\mathsf{O}(X,\cdot,y;w,\mathsf{h})$,
\begin{equation}\label{10.3:eq-output-decomposition-3}
\mathsf{O}(w,\mathsf{h})=\mathsf{O}(\hat w,\mathsf{h}(y))+\mathsf{S}^{b,\partial\eta}_j
+\mathsf{S}^{b,\eta\partial w}_j+\mathsf{D}^b_j ,
\end{equation}
where the defect is
\begin{equation*}
\mathsf{D}^b_j:=\mathsf{O}(w,\mathsf{h})-\mathsf{O}(\hat w,\mathsf{h}(y))-\mathsf{S}^{b,\partial\eta}_j
-\mathsf{S}^{b,\eta\partial w}_j ;
\end{equation*}
and on $\mathbb{D}_2$, for each $X$,
\begin{align*}
|\mathsf{S}^{b,\partial\eta}_j(X,\cdot,y;w,\mathsf{k})|
&\le L^j\,\mathsf{o}\,D_X\sup|\mathsf{k}|\,e^{-2\kappa d_j(X)},\\
|\mathsf{S}^{b,\eta\partial w}_j(X,\cdot,y;w,\mathsf{k}')|
&\le 4L^j\,\mathsf{m}\,\lambda_*^{-1}(1+\lambda_*D_X)\sup|\mathsf{k}'|\,e^{-2\kappa d_j(X)},\\
|\mathsf{D}^b_j|&\le 3\,\mathsf{o}\,\|\eta\|_{\langle 2\rangle}\,\frac{a_j^2}{a}\,(1+\lambda_*D_X)^2\,
e^{-2\kappa d_j(X)} ;
\end{align*}
equivalently, in the fresh norm $\|\cdot\|^{\mathrm{fr}}$, which carries the weight $e^{2\kappa d_j(X)}$,
$\|\mathsf{S}^{b,\partial\eta}_j(X,\cdot)\|^{\mathrm{fr}}\le L^j\mathsf{o}D_X\sup|\mathsf{k}|$, and
similarly for $\mathsf{S}^{b,\eta\partial w}_j$.
\end{proposition}

\begin{proof}
Throughout, $x$ ranges over $X^+$; since $y\in X^+$, Lemma~\ref{10.3:lem-own-level-freezing}\,(ii)
applies with $c=y$ below. Since
$D_X=\operatorname{diam}_j(X^+)+1$, we have $|x-y|\le D_X$ in units of $a_j$.

\emph{Expansion of the weight.} On $X^+$ write $\mathsf{h}=\mathsf{h}(y)+\mathsf{l}+\mathsf{h}_2$ with
\begin{equation*}
\mathsf{l}(x):=\sum_\nu(x-y)_\nu\,a_j\,\partial_\nu\mathsf{h}(y)=L^j\sum_\nu(x-y)_\nu\,\mathsf{k}_\nu(y);
\end{equation*}
the second equality holds because $a_j\partial_\nu\mathsf{h}=(a_j/a)\,\partial_\nu\eta=L^j\mathsf{k}_\nu$.
The physical displacement from $y$ to $x$ is $a_j(x-y)$, so Taylor's formula with remainder gives
\begin{equation*}
|\mathsf{h}_2(x)|\le\tfrac12\,a_j^2|x-y|^2\sup|\partial^2\mathsf{h}|
\le\tfrac12\,\frac{a_j^2}{a}\,\sup|\partial^2\eta|\,D_X^2 .
\end{equation*}

\emph{Expansion of the source.} Write $w=\hat w+\ell+w_2$ with
$\ell(x):=\sum_\nu(x-y)_\nu\,a_j\,\partial_\nu w(y)$, the first-order Taylor term of $w$ at $y$. Since
$\mathsf{h}(y)\,a_j\,\partial_\nu w(y)=(a_j/a)\,\eta(y)\,\partial_\nu w(y)$, we have
\begin{equation*}
\mathsf{h}(y)\,a_j\,\partial_\nu w(y)=L^j\,\mathsf{k}'_\nu(y).
\end{equation*}

\emph{The identity.} We show that the defect is
\begin{equation}\label{10.3:eq-output-decomposition-4}
\begin{split}
\mathsf{D}^b_j={}&\mathsf{O}(\hat w,\mathsf{h}_2)
+\bigl\{\mathsf{O}(w,\mathsf{h}(y))-\mathsf{O}(\hat w,\mathsf{h}(y))
-D_w\mathsf{O}(\hat w,\mathsf{h}(y))[\ell]\bigr\}\\
&+\bigl\{\mathsf{O}(w,\mathsf{h}-\mathsf{h}(y))-\mathsf{O}(\hat w,\mathsf{h}-\mathsf{h}(y))\bigr\}.
\end{split}
\end{equation}
By linearity in the weight,
$\mathsf{O}(w,\mathsf{h})=\mathsf{O}(w,\mathsf{h}(y))+\mathsf{O}(w,\mathsf{h}-\mathsf{h}(y))$ and
$\mathsf{O}(\hat w,\mathsf{h}-\mathsf{h}(y))=\mathsf{O}(\hat w,\mathsf{l})+\mathsf{O}(\hat w,\mathsf{h}_2)$.
By the formula for $\mathsf{l}$ and linearity,
$\mathsf{O}(\hat w,\mathsf{l})=\mathsf{S}^{b,\partial\eta}_j(X,\cdot,y;w,\mathsf{k})$. The derivative
$D_w\mathsf{O}(\hat w,e)[\psi]$ is linear in $\psi$ and, $\mathsf{O}$ being linear in the weight, linear in
the constant $e$; hence, by the identity for $\mathsf{h}(y)a_j\partial_\nu w(y)$,
\begin{align*}
D_w\mathsf{O}(\hat w,\mathsf{h}(y))[\ell]
&=\sum_\nu D_w\mathsf{O}\bigl(\hat w,\mathsf{h}(y)\,a_j\,\partial_\nu w(y)\bigr)\bigl[(x-y)_\nu\bigr]\\
&=\mathsf{S}^{b,\eta\partial w}_j(X,\cdot,y;w,\mathsf{k}').
\end{align*}
Inserting these three relations into the right-hand side of \eqref{10.3:eq-output-decomposition-4}
turns it into
\begin{equation*}
\mathsf{O}(w,\mathsf{h})-\mathsf{O}(\hat w,\mathsf{h}(y))-\mathsf{S}^{b,\partial\eta}_j
-\mathsf{S}^{b,\eta\partial w}_j ,
\end{equation*}
which is $\mathsf{D}^b_j$; this proves \eqref{10.3:eq-output-decomposition-4}, and with it
\eqref{10.3:eq-output-decomposition-3}.

\emph{The defect.} We bound the three terms of \eqref{10.3:eq-output-decomposition-4} in turn, on
$\mathbb{D}_2$, where the hypotheses of Lemma~\ref{10.3:lem-own-level-freezing}\,(ii) hold with
$\varepsilon=a_j$, as stated at the end of that lemma. We use
\begin{equation*}
\|\eta\|_{\langle2\rangle}
=\max\bigl\{\sup|\eta|,\ \sup|\partial\eta|/\lambda_*,\ \sup|\partial^2\eta|/\lambda_*^2\bigr\}.
\end{equation*}

First term: by linearity, \eqref{10.3:eq-output-decomposition-1} applied to the weight $\mathsf{h}_2$, the
bound on $\mathsf{h}_2$, $\sup|\partial^2\eta|\le\lambda_*^2\|\eta\|_{\langle2\rangle}$ and
$\lambda_*D_X\le1+\lambda_*D_X$,
\begin{equation*}
|\mathsf{O}(\hat w,\mathsf{h}_2)|\le\tfrac12\,\mathsf{o}\,\|\eta\|_{\langle2\rangle}\,\frac{a_j^2}{a}\,
(1+\lambda_*D_X)^2e^{-2\kappa d_j(X)} .
\end{equation*}

Second term: apply Lemma~\ref{10.3:lem-own-level-freezing}\,(ii) at $\varepsilon=a_j$ and $c=y$ to
$G(X;w):=\mathsf{O}(X,\cdot,y;w,\mathsf{h}(y))$, the weight being the constant $\mathsf{h}(y)$. By
\eqref{10.3:eq-output-decomposition-1}, $|G|\le\mathfrak{n}$ with
$\mathfrak{n}=\mathsf{o}\,|\mathsf{h}(y)|\,e^{-2\kappa d_j(X)}$, and
$|\mathsf{h}(y)|=|\eta(y)|/a\le\|\eta\|_{\langle2\rangle}/a$. Since $DG(X;\hat w)[\ell]$ is
$D_w\mathsf{O}(\hat w,\mathsf{h}(y))[\ell]$, the second bound of that lemma gives
\begin{equation*}
|\{\cdots\}|\le\mathfrak{n}\,a_j^2(1+\lambda_*D_X)^2
\le\mathsf{o}\,\|\eta\|_{\langle2\rangle}\,\frac{a_j^2}{a}\,(1+\lambda_*D_X)^2e^{-2\kappa d_j(X)} .
\end{equation*}

Third term: apply the first bound of Lemma~\ref{10.3:lem-own-level-freezing}\,(ii), again at
$\varepsilon=a_j$ and $c=y$, to $G(X;w):=\mathsf{O}(X,\cdot,y;w,\mathsf{h}-\mathsf{h}(y))$. By
\eqref{10.3:eq-output-decomposition-1} one may take
$\mathfrak{n}=\mathsf{o}\sup_{X^+}|\mathsf{h}-\mathsf{h}(y)|\,e^{-2\kappa d_j(X)}$, and by the mean value
theorem
\begin{equation*}
\sup_{X^+}|\mathsf{h}-\mathsf{h}(y)|\le a_j D_X\sup|\partial\mathsf{h}|
=\frac{a_j}{a}\,D_X\sup|\partial\eta|\le\lambda_*\|\eta\|_{\langle2\rangle}\,\frac{a_j}{a}\,D_X .
\end{equation*}
Hence, using $\lambda_*D_X\le1+\lambda_*D_X$ once more,
\begin{equation*}
|\{\cdots\}|\le\mathfrak{n}\,a_j(1+\lambda_*D_X)
\le\mathsf{o}\,\|\eta\|_{\langle2\rangle}\,\frac{a_j^2}{a}\,(1+\lambda_*D_X)^2e^{-2\kappa d_j(X)} .
\end{equation*}
The three coefficients add up to $\tfrac12+1+1\le3$, which is the bound on $\mathsf{D}^b_j$.

\emph{The pin sources.} The weight $x\mapsto\sum_\nu(x-y)_\nu\mathsf{k}_\nu(y)$ is bounded on $X^+$ by
$D_X\sup|\mathsf{k}|$, because $|x-y|\le D_X$; so \eqref{10.3:eq-output-decomposition-1} gives
\begin{equation*}
|\mathsf{S}^{b,\partial\eta}_j|\le L^j\,\mathsf{o}\,D_X\sup|\mathsf{k}|\,e^{-2\kappa d_j(X)},
\end{equation*}
which is the
first bound. For the second, apply \eqref{10.3:eq-output-decomposition-2} with $w_0=\hat w$,
$e=\mathsf{k}'_\nu(y)$ and
$\psi=(x-y)_\nu$. This $\psi$ vanishes at $y$ and has $\operatorname{Lip}_j\psi=1$; the estimate for
$\|w-\hat w\|_{j,X}$ in the proof of Lemma~\ref{10.3:lem-own-level-freezing}\,(ii), in which
$\operatorname{Lip}_jw\le\varepsilon r\lambda_*/4$ yields
$\|w-\hat w\|_{j,X}\le(\varepsilon r/4)(1+\lambda_*D_X)$, gives with $\varepsilon r\lambda_*/4$ replaced by
$1$ the bound $\|\psi\|_{j,X}\le\lambda_*^{-1}(1+\lambda_*D_X)$. Summing over the four values of $\nu$
and using $|\mathsf{k}'_\nu(y)|\le\sup|\mathsf{k}'|$ gives
\begin{equation*}
|\mathsf{S}^{b,\eta\partial w}_j|\le 4L^j\,\mathsf{m}\,\lambda_*^{-1}(1+\lambda_*D_X)\sup|\mathsf{k}'|\,
e^{-2\kappa d_j(X)},
\end{equation*}
which is the second bound.
\end{proof}

\emph{Type.} If $\mathsf{O}$ is produced by a map equivariant under the lattice symmetries, then the maps
$(X,y,A)\mapsto\mathsf{O}(X,\cdot,y;\hat w,x\mapsto A(x-y))$ and
$(X,y,e\otimes e_\nu)\mapsto D_w\mathsf{O}(X,\cdot,y;\hat w,e)[(x-y)_\nu]$ are
$(V\otimes V)$-typed: a lattice symmetry $\sigma$ with linear part $\sigma_*$ sends the weight
$x\mapsto A(x-y)$ to $x\mapsto(\sigma_*A\sigma_*^{-1})(x-\sigma y)$. The factors $D_X$ and
$(1+\lambda_*D_X)^2$ are
absorbed by the spare rate $\tfrac18\kappa\ge\tfrac14\delta\kappa$: with the constant
$C_B'=\sup_X(1+\lambda_*D_X)^2e^{-\frac14\delta\kappa d_j(X)}$ of
Lemma~\ref{10.3:lem-own-level-freezing}, $\delta$ being the parameter in the exponent of $C_B'$, which
satisfies $\delta\le\tfrac12$ so that the inequality just written holds, and with
$D_X\le\lambda_*^{-1}(1+\lambda_*D_X)$,
\begin{equation*}
(1+\lambda_*D_X)^2e^{-2\kappa d_j(X)}\le C_B'\,e^{-\kappa_\sharp d_j(X)},\qquad
\kappa_\sharp=\tfrac{15}{8}\kappa,
\end{equation*}
so that $\mathsf{S}^{b,\partial\eta}_j$, $\mathsf{S}^{b,\eta\partial w}_j$ and $\mathsf{D}^b_j$ are
bounded in the norms at the rate $\kappa_\sharp$, and a fortiori in $\|\cdot\|^{\flat}$, the factor
$D_X$ costing $\lambda_*^{-1}C_B'$ and the factors $1+\lambda_*D_X$ and $(1+\lambda_*D_X)^2$ costing
$C_B'$.

\emph{Reading.} (i) The freezing in the weight is linear algebra; its first-order term is the marginal
$(V\otimes V)$-family with the bounded weight $\partial\eta$, and its second-order term carries $a_j^2/a$
where the weight itself is of size $1/a$, a gain of two powers of $L^{-(K'-j)}$. (ii) The first-order
freezing term in $w$ is the derivative of $\mathsf{O}$ at the frozen source $\hat w$ in the direction of
the linear part of $w-\hat w$, with the bounded weight $\eta\partial w$; the remainder gains
$L^{-2(K'-j)}$.

\begin{definition}[The vector part of the run and its loads]\label{10.3:def-vector-run}
Let $\hat{\mathfrak{s}}=\hat{\mathfrak{s}}(w)$ be the hatted run of the correlation theorem; it is a list.
Let $\mathsf{h}$ be an arbitrary bounded function with values in $\mathsf{V}_{\mathbb{C}}$.

\emph{Added collections.} A \emph{unit} is a triple $v=(j,S_v,f_v)$, where $f_v$ is analytic and
gauge invariant on the level-$j$ configurations on $S_v$. Let $\mathfrak{u}$ be a collection of
such units. The \emph{production with $\mathfrak{u}$ added} is the dialled production of the
hatted run in which each use $(k',n,\cdot)$ also reads $\lambda_{n,k'}f_v$ for every
$v\in\mathfrak{u}$. We write $\mathsf{T}_k[\mathfrak{u}]$ for the derivative at $0$ of the
outputs of the complete step $k\to k+1$ of this production. The derivative is taken along the
line on which all the dials $\lambda_{n,k'}$ take one common value $\lambda$. The map
$\mathfrak{u}\mapsto\mathsf{T}_k[\mathfrak{u}]$ is linear in $\mathfrak{u}$.

\emph{Level $0$.} Put
\begin{equation}\label{10.3:eq-vector-run-1}
\mathsf{x}^{\sharp}_0:=\bigl\{c_0\,\mathsf{h}^{\rho}(x(p))\,\mathfrak{r}^{\sharp}_{\rho}(p)\bigr\}_p ,
\end{equation}
summed over $\rho$. The member indexed by $p$ is pinned at the plaquette $p$, and
$\mathfrak{r}^{\sharp}_{\rho}(p)$ is the bare vector unit at $p$, bounded in
Lemma~\ref{10.3:lem-bare-unit}.
When $\mathsf{h}=\eta/a$, the sum of the members of $\mathsf{x}^{\sharp}_0$ is
$c_0\,\mathfrak{K}^{\sharp}_{\eta}[1]$. Put $\mathsf{u}_0:=\mathsf{v}_0(c_0)$.

\emph{Levels $1,\dots,k$.} Suppose the following objects are given for $1\le j\le k$.
\begin{itemize}
\item Analytic families $\hat{\mathfrak{V}}^{(j)}_{\hat w}$, $\hat w$ in the complex disc
$D(0,2r)$ of centre $0$ and radius $2r$, of
$\mathsf{V}$-typed pointed $E$-families at scale $j$. These are frozen and defined on the whole
lattice. They are universal objects: they do not depend on $\eta$, on the test functions, or on
the torus.
\item Local marginal families $\mathsf{R}^{V(i)}(X,\cdot\,;w,\mathsf{h})$, $i\le k$.
\end{itemize}
Let $\hat{\mathfrak{V}}^{\sharp(j)}$ be the sharp part of $\hat{\mathfrak{V}}^{(j)}$ in the sense
of Definition~\ref{10.3:def-bond-current}. Put
\begin{equation}\label{10.3:eq-vector-run-2}
\mathsf{x}^{\sharp}_j:=\bigl\{\hat{\mathfrak{V}}^{\sharp(j)}_{\hat w(y)}(X,\cdot,y;\mathsf{h}(y)):\
y\in\Lambda_j^{\circledR},\ X\subset\Lambda_j\bigr\}.
\end{equation}
Let $\mathsf{rim}_j(\Lambda_j^{\circledR})$ be the collection of rim terms of the summation by
parts at the level of birth (Proposition~\ref{10.3:prop-summation-by-parts}). It is applied to the
family $\hat{\mathfrak{V}}^{(j)}$ with $\Lambda=\Lambda_j^{\circledR}$ and $\Lambda^+=\Lambda_j$.
The \emph{vector part at level $k$} is
\begin{equation}\label{10.3:eq-vector-run-3}
\begin{split}
\mathfrak{u}_k[w,\mathsf{h}]:={}&\mathsf{x}^{\sharp}_0\ \cup
\bigcup_{1\le j\le k}\bigl(\mathsf{x}^{\sharp}_j\cup\mathsf{rim}_j(\Lambda_j^{\circledR})\bigr)\\
&\cup\bigl\{\text{the unfiled shares of }\mathsf{R}^{V(i)},\ i\le k\bigr\},
\end{split}
\end{equation}
where the unfiled shares are those of the re-filing lemmas.

\emph{The step.} Let $(\mathsf{O}_{k+1},\mathsf{R}^{V(k+1)})$ be the $E$-output and the
$R$-output of $\mathsf{T}_k[\mathfrak{u}_k]$. Add to the $E$-output the re-filed remnants
whose indices lie in the set $\mathfrak{B}(k)$ of the re-filing lemmas, with $\mathsf{RF}$ the
re-filing map there:
\begin{equation}\label{10.3:eq-vector-run-4}
\hat{\mathsf{O}}_{k+1}:=\mathsf{O}_{k+1}+\sum_{i\in\mathfrak{B}(k)}\mathsf{RF}\bigl[\mathsf{R}^{V(i)}\bigr].
\end{equation}
The frozen family of level $k+1$ is this output read at constant
$(w,\mathsf{h})\equiv(\hat w,e)$ on the whole lattice:
\begin{equation}\label{10.3:eq-vector-run-5}
\hat{\mathfrak{V}}^{(k+1)}_{\hat w}(X,\cdot,y;e):=\hat{\mathsf{O}}_{k+1}(X,\cdot,y;\hat w,e).
\end{equation}
On the whole lattice and at constant $(w,\mathsf{h})$ the sharp part presents the same sum as the
family itself, for two reasons:
\begin{itemize}
\item on the whole lattice no bond has one end outside $\Lambda$, so the rim of
Proposition~\ref{10.3:prop-summation-by-parts} is empty;
\item the bond sources and the defect there carry the averages
$\langle\mathsf{k}_\nu\rangle_{j,\nu}$, $\langle\mathsf{k}'_\nu\rangle_{j,\nu}$ and
$\langle(\eta(y^+)-\eta)\partial_\nu w\rangle_{j,\nu}$, which vanish when $w$ and $\mathsf{h}$
are constant.
\end{itemize}
Hence
\begin{equation}\label{10.3:eq-vector-run-6}
\sum_y e\cdot\hat{\mathfrak{V}}^{\sharp(j)}(\cdot,y)=\sum_y e\cdot\hat{\mathfrak{V}}^{(j)}(\cdot,y).
\end{equation}

\emph{Norms.} All norms are taken per unit of $\sup|\mathsf{h}|$. The
suprema are taken as follows.
\begin{itemize}
\item For the frozen objects, over $\hat w\in D(0,2r)$, that is, over constant sources, as in the
correlation theorem's hypothesis at constant sources.
\item For the local objects, over $w$ in the class of the anchor.
\item In both cases, over $|\mathsf{h}|\le1$.
\end{itemize}
Put $\mathsf{r}_0:=0$ and
\begin{equation}\label{10.3:eq-vector-run-7}
\mathsf{u}_j:=\|\hat{\mathfrak{V}}^{(j)}\|^{\mathrm{fr}}\le
\mathsf{o}_j:=\|\hat{\mathsf{O}}_j\|^{\mathrm{fr}},\qquad
\mathsf{r}_i:=\|\mathsf{R}^{V(i)}\|^{\flat}.
\end{equation}

\emph{Anchors.} The objects above are given by production formulas and do not depend on the
anchor. The norms are taken at one anchor $K''$ at a time. The rows of the typed table for the
vector part, and the lemmas that bound the vector recursion, hold at every anchor with the same
constants. The norms of the run itself are read at the anchor $K''=K'$. The decomposition of a
local output at its pin (Proposition~\ref{10.3:prop-output-decomposition}) at level $j$ is read at
the anchor $K''=j$. For $z\in\mathbb{D}_2$ the source $w$ lies in the classes of every anchor,
because
\begin{equation}\label{10.3:eq-vector-run-8}
\operatorname{Lip}_i w\le L^{-(K'-i)}\lambda_* r/4\le\mathfrak{l}^{[K'']}_i .
\end{equation}

\emph{Loads.} Let $\varsigma^V_{\nu}$ be the scale kernels of
Definition~\ref{10.3:def-exponents-rates}, and let $w_{n,i}$ and $\mathsf{a}_i$ be those of the
re-filing lemmas. Put $\mathsf{t}^{\mathrm{raw}}:=c_0C_{\mathfrak{r}}'\mathfrak{g}_0$.
The \emph{vector load at level $n$} is
\begin{equation}\label{10.3:eq-vector-run-9}
T^V_n:=\sum_{0\le j\le n}\varsigma^V_{1+n-j}\,\mathsf{u}_j
+\sum_{n-\mathsf{a}_i\le i\le n}w_{n,i}\,\mathsf{r}_i
+\mathbf{1}_{n=0}\,\mathsf{t}^{\mathrm{raw}} .
\end{equation}
Let $T^{V,T}_k$ be the same expression at $n=k$, with two changes: the sum over the
$\mathsf{r}_i$ runs over $i\in\mathfrak{B}(k)$, and the term $\mathsf{t}^{\mathrm{raw}}$ is
omitted. Finally put
\begin{equation}\label{10.3:eq-vector-run-10}
\mathsf{t}_0:=\mathsf{t}^{\mathrm{raw}}/\mathsf{v}_0(c_0)\le C_{\mathsf{t}}\,g_0^{-1}.
\end{equation}
The factor $c_0$ cancels in this ratio.
\end{definition}

\begin{lemma}[The vector units and their sizes per cell]\label{10.3:lem-unit-table}
Fix a zone $n$ and let $\tilde{\square}$ be the current-zone cube, a member of the family $\pi_n$
of cubes of zone $n$. At every site the vector units $v$ are of the following kinds:
\begin{itemize}
\item[$0$.] the bare unit at a plaquette $p$;
\item[$0^{\mathrm{raw}}$.] the same unit in the zone whose first index is $0$;
\item[$1$.] the $\sharp$-core of level $j$ at a site $y$;
\item[$1'$.] the tails of $\hat{\mathfrak{V}}^{\sharp(j)}$, taken singly;
\item[$3$.] the share of the marginal remnant $\mathsf{R}^{V(i)}$ that is not re-filed;
\item[rim.] the constituents of $\mathsf{rim}_j(\Lambda_j^{\circledR})$, the rim of
Proposition~\ref{10.3:prop-summation-by-parts} on Ba{\l}aban's region $\Lambda_j$ with one
layer removed.
\end{itemize}
All sizes are measured per unit of $\sup|\mathsf{h}|$. The supports $S_v$, the sizes
$\mathfrak{N}(v)$ and the contributions per unit cell $\tilde{\square}\in\pi_n$ of zone $n$ are
\begin{equation}\label{10.3:eq-unit-table-a}
\begin{array}{l|l|l|l}
v & S_v & \mathfrak{N}(v) & \text{per unit $n$-cell}\\
\hline
0 & \text{see below} & C\mathsf{v}_0\rho_{n,0}^3t^6
  & \le C\mathsf{v}_0\varsigma^V_{1+n}\\
0^{\mathrm{raw}} & \text{---} & \text{modulus bound \eqref{10.3:eq-unit-table-1}}
  & \text{load } \mathsf{t}^{\mathrm{raw}}\\
1 & \square^{\sim\nu+1} & C_{\sharp}\mathsf{u}_j\nu^3\rho^2(1+\rho)t^{4+2\hat\alpha}
  & C_{\sharp}\mathsf{u}_j\varsigma^V_{\nu}\\
1' & X & \text{as in clause (c) of Lemma~\ref{10.3:lem-relocalisation}}
  & C_{\sharp}\mathsf{u}_j\rho^2t^5\\
3 & X & 2\mathsf{r}_i(C_MD_Xt^2\rho)^2e^{-\kappa d_i(X)} & Cw_{n,i}\mathsf{r}_i\\
\mathrm{rim} & X & 2C_{\mathfrak{J}}\mathsf{u}_j(C_MD_X\rho)^2e^{-\frac{15}{8}\kappa d_j(X)}
  & \le C_{\mathfrak{J}}'M^3\mathsf{u}_j
\end{array}
\end{equation}
where $\mathsf{v}_0=\mathsf{v}_0(c_0)$, and where:
\begin{itemize}
\item in every row the letters $\rho_{n,j}$, $\rho$, $t$ and $\nu$ are the inherited size
factors of the correlation theorem, read at the level of the unit in zone $n$ (in row $0$,
$t=L^{-n}$); the subscript $\rho$ of the bare unit $\mathfrak{r}^{\sharp}_{\rho}(p)$ below is a
lattice direction and is unrelated to them;
\item for the kind $0$ the support $S_v$ consists of the two bare cubes at $p$;
\item for the kind $0^{\mathrm{raw}}$ the size is the modulus bound, valid on $SU(2)$,
\begin{equation}\label{10.3:eq-unit-table-1}
\Bigl|e^{\lambda\sum\mathsf{x}^{\sharp}_0}\Bigr|
\le \exp\bigl(|\lambda|\,c_0C_{\mathfrak{r}}'\,A(|\mathsf{h}|^{\sim},U)\bigr),
\end{equation}
in which $\lambda$ is the common value of the dials $\lambda_{n,k'}$ of
Definition~\ref{10.3:def-vector-run}, and
the action $A(\cdot,U)$ with a weight in its first argument and the modification
$|\mathsf{h}|^{\sim}$ of $|\mathsf{h}|$ are those of clause (c) of the modulus lemma of the
correlation theorem; its entry per cell is the load $\mathsf{t}^{\mathrm{raw}}$, that is, a raw
Wilson shift $c_0C_{\mathfrak{r}}'$ in clause (3) of the modulus recursion;
\item for the kind $3$ the entry per cell is used with $n-i\le\mathsf{a}_i$; here and in the
row rim, $C_M$ and $D_X$ are the constant and the diameter factor of the localisation bound for
single terms of the correlation theorem;
\item for the kind rim the constant $C_{\mathfrak{J}}'$ is the one produced by the localisation
bound for single terms at $t=1$, and the entry per cell occurs in the zone $n=j$ only, and is
therefore contained in $C\varsigma^V_1\mathsf{u}_j$.
\end{itemize}
The kinds are numbered as in the table of scalar units, whose kinds $2$, $5$ and $\partial$ have
no vector counterpart: there is no kind $2$, the vector units being frozen point-wise; no kind
$5$, there being no weight slot; and no boundary kind $\partial$. We define $C_L^V$ to be the
largest constant of the last column of \eqref{10.3:eq-unit-table-a}, taken with the weights that
the unit lemma of the correlation theorem attaches to units (in the same way as the table
constant $C_L$ of the scalar units and the constant $\hat C_L$ of the hatted run
$\hat{\mathfrak{s}}$ are defined), and we set
\begin{equation*}
\varkappa^V_k:=K_{CE}\,C_L^V\,\theta_k ,
\end{equation*}
where $K_{CE}$ denotes the fixed constant of the cluster expansion of the correlation theorem,
and we write $\epsilon^V_k$, $\hat\varepsilon^V_k$ for the step constants $\epsilon_k$,
$\hat\varepsilon_k$ of the statement containing the bare-level modulus estimate, with the table
constant there replaced by $C_L^V$.
\end{lemma}

\begin{proof}
We justify the rows of \eqref{10.3:eq-unit-table-a} in turn, and then the absence of further
kinds; the last sentence of the lemma is a definition.

\emph{Row $0$.} The size is clause (ii) of Lemma~\ref{10.3:lem-bare-unit}: on a background
$\mathfrak{b}$ admissible at scale $n$ one has
$|c_0\mathfrak{r}^{\sharp}_{\rho}(p)(\mathfrak{b})|\le C\mathsf{v}_0(c_0)\rho^3_{n,0}t^6$ with
$t=L^{-n}$. Read per unit cell of zone $n$, this bound is at most
$C\mathsf{v}_0\varsigma^V_{1+n}$, the entry of the last column, in terms of the scale kernels
$\varsigma^V_{\nu}$ of Definition~\ref{10.3:def-exponents-rates}.

\emph{Row $0^{\mathrm{raw}}$.} Clause (iii) of Lemma~\ref{10.3:lem-bare-unit} bounds, on
$SU(2)$-fields, $|\mathfrak{r}^{\sharp}_{\rho}(p)|$ by $C_{\mathfrak{r}}$ times the plaquette
quantities $a_{p'}$ of that clause, summed over the plaquettes $p'$ of the two cubes adjacent to
$p$ in the direction $\rho$ and in the opposite direction. Summing
over the bare units with the weight $\mathsf{h}$ and combining with clause (c) of the modulus
lemma of the correlation theorem and with the bare-level modulus estimate for the zone of first
index $0$ gives
\eqref{10.3:eq-unit-table-1}. The factor $c_0C_{\mathfrak{r}}'$ multiplying
$A(|\mathsf{h}|^{\sim},U)$ in the exponent is what enters clause (3) of the modulus recursion as
a raw Wilson shift, and this is the load $\mathsf{t}^{\mathrm{raw}}$.

\emph{Row $1$.} The size of the $\sharp$-core of level $j$ at $y$, with support
$\square^{\sim\nu+1}$, is given by clauses (c) and (d) of the re-localisation lemma
(Lemma~\ref{10.3:lem-relocalisation}). Where the correlation theorem distinguishes smooth and
rough cores, a bound of the same form is supplied by the re-localisation lemma on extended
backgrounds (Lemma~\ref{10.3:lem-relocalisation-extended}): with the parameter $\epsilon\le1$ of
that lemma, its bound
$|\operatorname{Core}^{\sharp}_y(\operatorname{ext}_p(\mathcal{U});e)|\le2\epsilon^2\mathfrak{N}^V_p$
carries the factor $\nu_p^3\bar\rho^2_{p-j}(1+\bar\rho_{p-j})t_p^{4+2\hat\alpha}$, with the
constant of that lemma in place of $C_{\sharp}$ and the additional factor $2\epsilon^2$. Per unit
cell the entry is $C_{\sharp}\mathsf{u}_j\varsigma^V_{\nu}$.

\emph{Row $1'$.} The tails of $\hat{\mathfrak{V}}^{\sharp(j)}$, taken singly and localised in $X$,
have the size given by clause (c) of Lemma~\ref{10.3:lem-relocalisation}; the localisation bound
for single terms of the correlation theorem converts it into the entry
$C_{\sharp}\mathsf{u}_j\rho^2t^5$ per unit cell.

\emph{Row $3$.} The share of $\mathsf{R}^{V(i)}$ that is not re-filed, localised in $X$, has by
the localisation bound for single terms the size
$2\mathsf{r}_i(C_MD_Xt^2\rho)^2e^{-\kappa d_i(X)}$ and the entry $Cw_{n,i}\mathsf{r}_i$ per unit
cell. The restriction $n-i\le\mathsf{a}_i$ comes from clause (Z) of the rate bound for
remnants.

\emph{Row rim.} By Proposition~\ref{10.3:prop-summation-by-parts}, each term of
$\mathsf{rim}_j(\Lambda_j^{\circledR})$ is a constituent of the current $\mathfrak{J}$ with the
undifferentiated weight $\mathsf{h}$, pinned within one step of the boundary of
$\Lambda_j^{\circledR}$. Clause (a) of Lemma~\ref{10.3:lem-relocalisation} bounds such a
constituent by $C_{\mathfrak{J}}$ times the norm of the family times the modulus $|e|$ of its
argument $e$, with the decay
$e^{-\frac{15}{8}\kappa d_j(X)}$; the localisation bound for single terms, applied at $t=1$, then
gives the size in \eqref{10.3:eq-unit-table-a} and, with the constant $C_{\mathfrak{J}}'$ of
this conversion, the entry $C_{\mathfrak{J}}'M^3\mathsf{u}_j$ per
unit cell. By the first clause of the boundary lemma for vector units these terms occur in the
zone $n=j$ only, so the entry lies inside $C\varsigma^V_1\mathsf{u}_j$.

\emph{Absent kinds.} There is no kind $2$ because the vector units are frozen point-wise. There
is no kind $5$: the exponent carried by the weight slot of the scalar units is $0$ for the
vector units, by the standing hypothesis on the vector units that sets this exponent to zero, so
that there is no weight slot and no weighted unit arises. There is no kind $\partial$: boundary
contributions are accounted for by the
boundary lemma for vector units, so there are no boundary units in the table.
\end{proof}

\begin{definition}[Production bounds for the vector part under the split rule]
\label{10.3:def-vector-production-bounds}
Let the hatted run $\hat{\mathfrak{s}}$, the added units $v=(j,S_v,f_v)$, the production with a
collection of units added, the vector part $\mathfrak{u}_k[w,\mathsf{h}]$, the dials
$\lambda_{n,k'}$ and the vector loads $T^V_n$, $T^{V,T}_k$ be as in
Definition~\ref{10.3:def-vector-run}. Let $C_L^V$, $\varkappa^V_k$ and $\epsilon^V_k$ be the
constants fixed with the table \eqref{10.3:eq-unit-table-a} of
Lemma~\ref{10.3:lem-unit-table}, let $E_1$ be the constant of that name in the hatted
production proposition for the dialled run, and put
\begin{equation*}
C_L^{V\prime} := C_L^V E_1 .
\end{equation*}
Further, $C_{RF}$ denotes the re-filing constant of the re-filing lemmas; $K_2$ is the constant
of the second-derivative statement in clause (b) of the one-lattice Lipschitz step;
$\Sigma^q_{K'}[\delta]$ is the scale sum of a direction $\delta$ in the typed form of the torus
bound; and $(\mathrm{f}1)$, $(\mathrm{f}2)$ are the first two of the properties
$(\mathrm{f}1)$--$(\mathrm{f}3)$ through which the correlation theorem's sourced format reads a
family.
We say that the \emph{production bounds for the vector part under the split rule} hold if the
following is true. There is $\gamma_V \le \hat\gamma_L$, where $\hat\gamma_L$ is the coupling
ceiling of the hatted production and $\gamma_V$ is obtained from the hatted smallness condition
read at the constant $C_L^{V\prime}$ and at the rate
$c_a$ of \eqref{10.3:eq-exponents-rates-c} (Definition~\ref{10.3:def-exponents-rates}), such
that for $g \le \gamma_V$, for every anchor $K'' \le K'$, every $k < K''$, every
$w \in \mathcal{W}^{[K'']}_k$ and every bounded $\mathsf{h}$, the clauses (T1)--(T3) below
hold.
\begin{enumerate}
\item[(T1)] (Class.) Under the split rule
\begin{equation}\label{10.3:eq-vector-production-bounds-a}
\text{at site (L), a unit } v\in\mathfrak{u}_k \text{ goes to the first part}
\iff S_v\cap Z\neq\emptyset ,
\end{equation}
every derived object of the production with $\mathfrak{u}_k[w,\mathsf{h}]$ added is holomorphic
on the dial domain
\begin{equation}\label{10.3:eq-vector-production-bounds-b}
\mathbb{D}_V := \bigl\{\, |\lambda_{n,k'}|\,\sup|\mathsf{h}|\,T^V_n \le E_1 \,\bigr\}
\end{equation}
and obeys there the bounds of the hatted production proposition for the dialled run.
\item[(T2)] (Outputs.) The outputs
\begin{equation}\label{10.3:eq-vector-production-bounds-c}
\mathsf{O}_{k+1}(X,\cdot,y;w,\mathsf{h}), \qquad \mathsf{R}^{V(k+1)}(X,\cdot;w,\mathsf{h})
\end{equation}
are linear in $\mathsf{h}$, read $(w,\mathsf{h})$ on $X^+$ only, are analytic in $w$ on
$\mathcal{W}^{[K'']}_{k+1}(X)$, have $(\mathrm{f}1)$ and $(\mathrm{f}2)$, do not depend on the
label for $X\subset\Lambda_{k+1}$, and at constant weight $\mathsf{h}$ are $\mathsf{V}$-typed:
$\mathsf{O}_{k+1}$ in the sense of $(\mathrm{f}3)^{\mathsf{V}}$, and $\mathsf{R}^{V(k+1)}$ in the
typed form of \cite[(2.32)]{B-CMP119}.
\item[(T3)] (Size bounds.) First,
\begin{equation}\label{10.3:eq-vector-production-bounds-d}
\begin{gathered}
\|\mathsf{O}_{k+1}\|^{\mathrm{fr}} \le \varkappa^V_k\,T^{V,T}_k ,
\qquad\text{hence}\\
\mathsf{o}_{k+1} \le \varkappa^V_k\,T^{V,T}_k
+ C_{RF}\sum_{i\in\mathfrak{B}(k)} w_{k,i}\,\mathsf{r}_i .
\end{gathered}
\end{equation}
Second,
\begin{equation}\label{10.3:eq-vector-production-bounds-e}
\mathsf{r}_k \le \epsilon^V_k \sum_{n\le k}(k-n+1)\,q^{(k-n-1)_+}\,T^V_n
+ 3\varkappa^V_k\sum_{k_1(k)<i<k}\mathsf{r}_i .
\end{equation}
Third, let $\delta'$ be a direction of the scalar hatted list, and let $G(\lambda,\lambda')$
denote the E-output at level $k+1$ of the production $\hat{\mathfrak{s}}+\lambda'\delta'$ with
$\lambda\,\mathfrak{u}_k[w,\mathsf{h}]$ added. Then
\begin{equation}\label{10.3:eq-vector-production-bounds-f}
\bigl\|\partial_\lambda\partial_{\lambda'}G\big|_{\lambda=\lambda'=0}\bigr\|^{\mathrm{fr}}
\le \tfrac14 K_2\,(C_L^V)^2\,\theta_k^2\,T^{V,T}_k\,[\![\hat\delta']\!]^T_k .
\end{equation}
Fourth, clauses $(\alpha)$ and $(\beta)$ of the typed form of the torus bound hold after the
replacements
\begin{equation}\label{10.3:eq-vector-production-bounds-g}
\begin{gathered}
\Sigma^q_{K'}[\delta] \;\mapsto\;
\sup|\mathsf{h}|\sum_{n\le K'}(K'-n+1)\,q^{(K'-n-1)_+}\,T^V_n ,\\
\mathsf{e}_j+\mathsf{r}_j \;\mapsto\; \sup|\mathsf{h}|\,(\mathsf{o}_j+\mathsf{r}_j).
\end{gathered}
\end{equation}
\end{enumerate}
The same clauses (T1)--(T3) hold for an added collection of marginal units that are read only
by the marginal lemma, of the two kinds in rows $2$ and $3$ of the table of marginal units, with
the flat load $T^{\mathsf{D}}_n$ of the defect families in place of the vector load $T^V_n$.
\end{definition}

The assertion that the production bounds for the vector part under the split rule hold enters,
wherever it is used, as an explicitly displayed hypothesis. Proof outstanding.

\smallskip\noindent\textit{Route.}
A proof has to show that the collection $\mathfrak{u}_k=\mathfrak{u}_k[w,\mathsf{h}]$ can be
added to the production at each place where the proof of the hatted production proposition for
the dialled run reads its inputs. It then has to derive
(T2) and (T3) from the result. The readings are those listed for the correlation theorem's
proof, in their hatted copy, and they are to be treated one by one as follows.
\begin{itemize}
\item The marginal step, the re-localisation step, the decomposition step and the Wilson-shift
step of the correlation theorem's proof read one family (a weight) through
$(\mathrm{f}1)$--$(\mathrm{f}3)$ and its norm $\mathfrak{n}$. For $\mathfrak{u}_k$ the
treatment is as follows.
\begin{itemize}
\item The marginal step reads the entries $1'$, $3$ and rim of
\eqref{10.3:eq-unit-table-a} without change; it uses no covariance and no region other
than $X$.
\item The re-localisation step is replaced by clauses (c) and (d) of
Lemma~\ref{10.3:lem-relocalisation}; no counterterm arises, and nothing is attached to another
point.
\item The decomposition step is not read, since the units are frozen point-wise; the
decomposition is made on the output, by the decomposition statement for outputs.
\item The Wilson-shift step is read for the entry $0^{\mathrm{raw}}$ only, through its
clause (c).
\end{itemize}
\item At site (L), smooth and rough cores are read by the marginal and re-localisation steps for
extended backgrounds. The latter is replaced by Lemma~\ref{10.3:lem-relocalisation-extended};
the former is read without change.
\item The reading of near points one at a time contributes nothing under the split rule
\eqref{10.3:eq-vector-production-bounds-a}, by the near-point lemma for the vector part.
\item The unit lemma reads $f_v$ linearly, together with the size $\mathfrak{N}(v)$ and
analyticity, uniformly, in every parameter in which $f_v$ is analytic. For $\mathfrak{u}_k$,
$\mathfrak{N}$ is subadditive. Let $\mathsf{W}$ be the activity bound of the hatted production
proposition for the dialled run, $\mathsf{W}_J$ the activity bound of the correlation theorem's
production, and $\mathsf{w}$ the weight with which the unit lemma sums unit sizes. Summing the
last column of \eqref{10.3:eq-unit-table-a} gives, on $\mathbb{D}_V$,
\begin{equation*}
\mathsf{W} \le \mathsf{W}_J + |\lambda_{n,k'}|\,\mathsf{w}^{-1}C_L^V\sup|\mathsf{h}|\,T^V_n
\le \mathsf{W}_J + \mathsf{w}^{-1}C_L^{V\prime}.
\end{equation*}
The second inequality holds because $|\lambda_{n,k'}|\sup|\mathsf{h}|\,T^V_n \le E_1$ on
$\mathbb{D}_V$ and $C_L^{V\prime}=C_L^VE_1$. The support of a $\sharp$-core has $n$-side at most
\begin{equation*}
(2\nu+3)ML^{-(\nu-1)} \le 5M ,
\end{equation*}
so it meets at most $5^d$ current cubes. This is the count in the second remark following the
near-point lemma of the correlation theorem's proof, with $3$ replaced by $5$.
\item The activity sums are contained in the hatted smallness condition read at
$C_L^{V\prime}$.
\item The sup-norms of derived objects are given by the induction hypothesis. The dials carried
by the stored $\mathbf{B}'$ terms are those of the carry lemma, and the bounds on
$\mathbb{D}_V$ continue to hold.
\item The weight-linear expansions do not act, since no unit carries a weight slot.
\item Brackets are read through the unit lemma, and values through the third remark following the
near-point lemma. For $\mathfrak{u}_k$, the brackets of units of the second part are taken by
the unit lemma. Units of the first part, and subtracted copies, are values at one admissible
background of whole $\sharp$-cores. They are bounded by $C_L^{V\prime}$ per unit of
$\sum|\Gamma^0_n\cap Y|$, where $Y$ is the region on which the bracket or value is read and
$\Gamma^0_n$ is the cell set of the re-filing lemmas.
\item The modulus bounds comprise the bound on $\sup|V|$, the raw Wilson shift and the two
bounds per cell. For the raw Wilson shift, the entry $0^{\mathrm{raw}}$ is read with
$|\lambda|\,c_0C_{\mathfrak{r}}' \le E_1/\mathfrak{g}_0$. The bounds per cell acquire an
additional term of at most $C_L^{V\prime}$ per cell.
\item The identities linear in the integrand, and the combinatorics that do not involve $z$,
are read without change.
\end{itemize}

(T2) follows as below.
\begin{itemize}
\item Linearity and locality are to be taken from the locality statement of the correlation
theorem's proof: potentials in $Y\subset X$ read $\zeta_k{\restriction}Y$ and the old terms in
$Y$, that is, $w{\restriction}X^+$.
\item For the bound with $\sup_{X^+}|\mathsf{h}|$, $\mathsf{h}$ is replaced by
$\mathsf{h}\mathbf{1}_{X^+}$; this is allowed because no regularity of $\mathsf{h}$ is asked.
\item For the type, the production map is equivariant under the symmetries of $T^{(k+1)}$ acting
on the pair (inputs, $\mathsf{h}$). This is the equivariance statement, together with the list of
the statements that read $(\mathrm{f}3)$.
\item The $\sharp$-presentation is covariant by clause (a) of
Lemma~\ref{10.3:lem-relocalisation}, so its differential maps constant weights to
$\mathsf{V}$-typed families.
\end{itemize}

For \eqref{10.3:eq-vector-production-bounds-d}, the proof of clause (a) of the hatted form of
the one-lattice Lipschitz step is rerun at site (T), for $X\subset\Lambda_{k+1}$, on the
whole lattice and in zone $k$, with the following data.
\begin{itemize}
\item The potential is $V^m:=V[\mathfrak{u}_k]$, the potential that the unit lemma assigns to
the collection $\mathfrak{u}_k$, defined through the linearity of the unit lemma.
\item The activity $\mathsf{W}^m$ of $V^m$ obeys $\mathsf{W}^m\le\theta_kC_L^VT^{V,T}_k$, by the
entries $0$, $1$, $1'$ and $3$ of \eqref{10.3:eq-unit-table-a}.
\item The corollary of the hatted production proposition for the dialled run is applied at the
output's own rate $\mathsf{r}(\delta)\kappa$, where $\mathsf{r}(\delta)$ is the rate factor that
the hatted one-lattice Lipschitz step attaches to the direction $\delta$ it reads; it is distinct
from the remnant norms $\mathsf{r}_i$.
\end{itemize}
In that proof, no unit is read through the induction hypothesis on sup-norms of derived
objects.

The bound \eqref{10.3:eq-vector-production-bounds-e} is clauses (c) and (e) of the hatted form of
the one-lattice Lipschitz step, read with their own letters at the load $T^V$. The proof of
clause (e) is incomplete at one step, which rests on the carry lemma.

For \eqref{10.3:eq-vector-production-bounds-f}, the route is the second-derivative statement in
clause (b) of the one-lattice Lipschitz step. It gives a function $G(\lambda,\lambda')$,
holomorphic on $|\lambda|<\varrho^{(1)}$, $|\lambda'|<\varrho^{(2)}$, with
$|G(X)|\le E^{\sharp}e^{-\mathsf{r}(\delta)\kappa\,d_{k+1}(X)}$ there. Here $E^{\sharp}$,
$\nu_0$, $\alpha_4$ and $C_A$ are the constants of that statement, $\hat C_L$ is the table
constant of the hatted production, $\mathsf{W}_1$ and $\mathsf{W}_2$ are the activities of the
two dialled directions, and the radii $\varrho^{(1)}$, $\varrho^{(2)}$ are
\begin{equation*}
\varrho^{(i)}=\frac{\nu_0\alpha_4}{8C_A\mathsf{W}_i},\qquad
\mathsf{W}_1\le\theta_kC_L^VT^{V,T}_k,\qquad
\mathsf{W}_2\le\theta_k\hat C_L[\![\hat\delta']\!]^T_k .
\end{equation*}
The Cauchy estimate at $0$ then gives
\begin{equation*}
|\partial_\lambda\partial_{\lambda'}G|
\le\frac{E^{\sharp}}{\varrho^{(1)}\varrho^{(2)}}
=E^{\sharp}\Bigl(\frac{8C_A}{\nu_0\alpha_4}\Bigr)^2\mathsf{W}_1\mathsf{W}_2
\le\tfrac14K_2\,\mathsf{W}_1\mathsf{W}_2 .
\end{equation*}
The right-hand side of \eqref{10.3:eq-vector-production-bounds-f} follows from the bounds on
$\mathsf{W}_1$ and $\mathsf{W}_2$ and from $C_L^V\ge\hat C_L$.

For \eqref{10.3:eq-vector-production-bounds-g}, the proof of the typed form of the torus bound
reads the direction only through its loads and through the age corollary for outputs.

\begin{remark}
The production bounds for the vector part under the split rule are first-order statements in
$\lambda$. Nothing is asserted at second order in $\lambda$, and no exactness is asserted off
the tangent.
\end{remark}

\begin{lemma}[The rim lemma for the vector part. Stated; proof incomplete at the step marked $(\ast)$]\label{10.3:lem-vector-rim}
Work in the setting of the production bounds for the vector part under the split rule
(Definition~\ref{10.3:def-vector-production-bounds}), whose proof is outstanding. Let
$\{\Omega_j,\Lambda_j,I_j\}$ be Ba{\l}aban's label of the large-field configuration
\cite{B-CMP119}, the letter $\Omega_j$ denoting a region of this label and not the vacuum vector,
and for each level $j$ let $\Lambda_j^{\circledR}$ be the set obtained from $\Lambda_j$ by removing
one layer of the $MR_j$-cubes of Ba{\l}aban's construction \cite{B-CMP119}. Let
$\mathsf{rim}_j(\Lambda)$ and the bond
terms $\mathsf{S}^{a,\partial\eta}_j$, $\mathsf{S}^{a,\eta\partial w}_j$, $\mathsf{D}^a_j$ be those of
the summation by parts at the level of birth, Proposition~\ref{10.3:prop-summation-by-parts}.

At the large-field operation $\mathbf{R}_k$ of step $k$ and for a large-field region $Z$, the
exponent is split into a first part, which may read the variables integrated in $Z$, and a second
part, which must not read them. The split rule
\eqref{10.3:eq-vector-production-bounds-a} sends a unit $v\in\mathfrak{u}_k$ to the first part if
and only if $S_v\cap Z\neq\emptyset$. Then the following hold.
\begin{enumerate}
\item[(i)] (Where the undifferentiated weight sits.) The rim $\mathsf{rim}_j(\Lambda_j^{\circledR})$
is pinned within one $T^{(j)}$-lattice step of $\partial\Lambda_j^{\circledR}$, and
$\partial\Lambda_j^{\circledR}\subset\Omega_j\setminus\Omega_{j+1}$. It therefore lies in zone $j$,
with $t=1$. It is attached to the large-field region of level $j$ and to no later one.
\item[(ii)] (No cut.) Assume the split rule \eqref{10.3:eq-vector-production-bounds-a}. Then at the
large-field operation of every level $k\ge j$, and for every region $Z$, each $\sharp$-core of
level $j$ is, in the part to which it is sent, a function of one background along one chain.
Parts (c) and (d) of Lemma~\ref{10.3:lem-relocalisation}, or, for a core read through
$\operatorname{ext}_p$, Lemma~\ref{10.3:lem-relocalisation-extended}, apply to it as they stand,
whatever the position of
its pin $y$ relative to $\partial Z$. No boundary term with the weight $\mathsf{h}$ arises at
$\partial Z$.
\item[(iii)] (Healing.) Suppose $\mathbf{R}_k$ heals $Z$. Let $j(Z)$ denote the step at which the
large-field region $Z$ was created, and $Z^{\sim}$ the enlarged region associated with $Z$ in
Ba{\l}aban's construction \cite{B-CMP122b}. Let $\Lambda_j^{\circledR,\mathrm{new}}$
denote the set $\Lambda_j^{\circledR}$ of the label after healing. Consider the levels $j$ with
$\Lambda_j^{\circledR,\mathrm{new}}\supsetneq\Lambda_j^{\circledR}$. For these levels the missing
terms (all missing terms with localisation domains intersecting $Z$) are added at $U_k$ and
subtracted \cite{B-CMP122b}, in the presentation of
Proposition~\ref{10.3:prop-summation-by-parts} on
$\Lambda_j^{\circledR,\mathrm{new}}\setminus\Lambda_j^{\circledR}$. Then, at the one background
$U_k$,
\begin{equation}\label{10.3:eq-vector-rim-1}
\begin{split}
&\mathsf{rim}_j(\Lambda_j^{\circledR})
+\mathsf{rim}_j\bigl(\Lambda_j^{\circledR,\mathrm{new}}\setminus\Lambda_j^{\circledR}\bigr)\\
&\qquad=\mathsf{rim}_j(\Lambda_j^{\circledR,\mathrm{new}})
+\sum_{b}\bigl[\mathsf{S}^{a,\partial\eta}_j+\mathsf{S}^{a,\eta\partial w}_j
+\mathsf{D}^a_j\bigr]_b .
\end{split}
\end{equation}
Here the sum runs over the bonds $b$ of $T^{(j)}$ joining $\Lambda_j^{\circledR}$ to
$\Lambda_j^{\circledR,\mathrm{new}}\setminus\Lambda_j^{\circledR}$, and $[\,\cdots]_b$ denotes the
bond terms of Proposition~\ref{10.3:prop-summation-by-parts} attached to $b$. Consequently the
healed rim disappears from the new exponent. The subtracted copies are values inside $Z^{\sim}$:
whole $\sharp$-cores and rim constituents of levels $j\ge j(Z)-1$.
\item[(iv)] (Count.) Per unit $n$-cell, with the weights attached to units in the correlation
theorem and the weights of \cite[(1.66)]{B-CMP122b}, the units of $\mathfrak{u}_k$ in zone $n$ sum
to at most $C_L^V\sup|\mathsf{h}|\,T^V_n$.
\end{enumerate}
\end{lemma}

Stated; the proof below is incomplete at the step marked $(\ast)$.

\begin{proof}
(i) The set $\Lambda_j^{\circledR}$ is obtained by removing one layer of the $MR_j$-cubes from
$\Lambda_j$ \cite{B-CMP119}. The region $\Omega_{j+1}$ stays away from the complement of
$\Lambda_j$ by more than that layer. Indeed, by Ba{\l}aban's construction \cite{B-CMP119}, and as
stated among the geometric properties of the label in the correlation theorem,
\begin{equation*}
\operatorname{dist}(\Lambda_j^{c},\Omega_{j+1})\;\ge\;8L\check R_{j+1}\;\ge\;8\check R_j .
\end{equation*}
The second inequality holds because $\check R_{j+1}\ge\check R_j/L$ \cite{B-CMP119}. Hence
$\partial\Lambda_j^{\circledR}\subset\Omega_j\setminus\Omega_{j+1}$.

The sum over $y$ in the definition of the level-$j$ vector terms runs over $\Lambda_j^{\circledR}$
and over nothing else \cite[(2.26)]{B-CMP119}. So the presentation of
Proposition~\ref{10.3:prop-summation-by-parts} with $\Lambda=\Lambda_j^{\circledR}$ has no boundary
other than $\partial\Lambda_j^{\circledR}$, and its only boundary term is
$\mathsf{rim}_j(\Lambda_j^{\circledR})$, whose constituents carry the undifferentiated weight
$\mathsf{h}$ (Proposition~\ref{10.3:prop-summation-by-parts}).

By the definition of $\mathsf{rim}_j$ in that proposition, every term of
$\mathsf{rim}_j(\Lambda_j^{\circledR})$ belongs to a bond $b_\nu(y)$ with exactly one end in
$\Lambda_j^{\circledR}$. The term carries the weight $\mathsf{h}$ at that end: at $y$ if
$y\in\Lambda_j^{\circledR}$, at $y^+$ if $y^+\in\Lambda_j^{\circledR}$. It is therefore pinned
within one $T^{(j)}$-lattice step of $\partial\Lambda_j^{\circledR}$. Combined with the
inclusion just proved, the rim lies in $\Omega_j\setminus\Omega_{j+1}$, that is in zone $j$, with
$t=1$. It is attached to the large-field region of level $j$, and to no later one.

(ii) Besides the large-field operation, the correlation theorem reads the exponent of the step at
two further places. At each of them the whole action sits under one integral, and no term-wise
split occurs, as stated in the correlation theorem at those places.

It remains to treat the large-field operation $\mathbf{R}_k$. There the presentation of
Proposition~\ref{10.3:prop-summation-by-parts} is an identity of the exponent at the single
configuration that the exponent reads before the split made by $\mathbf{R}_k$. After the split, the
vector part is a sum of units. Whether a term is cut by $\partial Z$ is therefore a property of the
assignment of its unit.

Compare the corresponding group of terms in the correlation theorem. There the split follows
Ba{\l}aban's split constituent by constituent, according as the localisation domain $X$ meets $Z$
or not. This is because that group contains a piece of the action, the one introduced for that group
in the correlation theorem, which Ba{\l}aban files with the Wilson term; it does not matter whether
that piece is filed whole or plaquette by plaquette. A group near $Z$ is then spread over several
backgrounds.

The vector group contains no piece of the action. Its split is an add-and-subtract in whatever
localisation is chosen, a localisation being chosen per use as in the correlation theorem. It is
subject to one constraint only: a term of the second part must not read the variables integrated
in $Z$. The split rule \eqref{10.3:eq-vector-production-bounds-a}, stated in the production bounds
for the vector part under the split rule
(Definition~\ref{10.3:def-vector-production-bounds}), whose proof is
outstanding, gives the
$\sharp$-core the one localisation domain $S_v$ and applies Ba{\l}aban's rule to it. There are two
cases.

If $S_v\subset Z^{c}$, the unit is a term of the second part. The data of the chain
\cite[(1.59)]{B-CMP122b} depend on $\tilde{\square}$ and on
$\tilde{\square}_0=\tilde{\square}^{\sim 5}$, and not on $X$. Moreover
$S_v=\square^{\sim\nu+1}$, $\nu$ here being the index $\nu_p=1+p-j$ of
Lemma~\ref{10.3:lem-relocalisation-extended} for a core read at scale $p$, not the bond direction
of $b_\nu(y)$, lies within
$(\nu+2)ML^{-(\nu-1)}\le 3M$ units of scale $n$ of $y$, hence
inside $\tilde{\square}_0$. The core is therefore one unit in the sense in which units are counted
in the correlation theorem, as are the far cores there.

If $S_v\cap Z\neq\emptyset$, the unit is a term of the first part. For $X:=S_v$ one forms the
smallest localisation domain containing $X$ and containing those components of $Z$ which intersect
$X$ \cite{B-CMP122b}. The representation \cite[(1.61)]{B-CMP122b} and its successors are then
common to the whole core.

In both cases the backgrounds are admissible on the box. The chains differ from $U_k$ only through
changes made inside the domain $\Lambda$ of Ba{\l}aban's construction \cite{B-CMP122b} (a domain of
that construction, not the set $\Lambda$ of Proposition~\ref{10.3:prop-summation-by-parts}), so a
box meeting $Z$ is admissible; Ba{\l}aban's own box
$\tilde{\square}^{\sim 5}$ at \cite[(1.59)]{B-CMP122b} does meet it. Where, in the correlation
theorem, the background is carried by a product of layers with the raw field, the core is read
through
$\operatorname{ext}_p$, and Lemma~\ref{10.3:lem-relocalisation-extended} applies. In either case the
core is a function of one background along one chain, and parts (c) and (d) of
Lemma~\ref{10.3:lem-relocalisation}, or, for a core read through $\operatorname{ext}_p$,
Lemma~\ref{10.3:lem-relocalisation-extended},
apply to it as they stand, whatever the position of $y$ relative to $\partial Z$.

Finally, the split made by $\mathbf{R}_k$ is an add-and-subtract: what is sent to the first part
is re-added. The moves applied to the exponent in the proof of the correlation theorem are linear
in
the integrand. Hence no term of the presentation is divided at $\partial Z$, and no boundary term
with the weight $\mathsf{h}$ arises there.

(iii) Suppose $\mathbf{R}_k$ heals $Z$, and let $j$ be a level with
$\Lambda_j^{\circledR,\mathrm{new}}\supsetneq\Lambda_j^{\circledR}$. $(\ast)$ Following Ba{\l}aban's
healing \cite{B-CMP122b}, all the missing terms with localisation domains intersecting the domain $Z$
are added at $U_k$ and subtracted, and the added terms are written in the presentation of
Proposition~\ref{10.3:prop-summation-by-parts} on
$\Lambda_j^{\circledR,\mathrm{new}}\setminus\Lambda_j^{\circledR}$. The proof is incomplete at this
step: the hypothesis governing this re-adding of the missing terms at healing is not stated here
and is not proved. It requires at least two inputs. The first is an admissibility input: the
background $U_k$, at which the missing terms of level $j$ are added, is admissible at the scale of
the $n$-cell of Lemma~\ref{10.3:lem-unit-table} in which these terms are counted in part (iv)
below, on $\Lambda_j^{\circledR,\mathrm{new}}\setminus\Lambda_j^{\circledR}$, so that these terms are
defined there and Proposition~\ref{10.3:prop-summation-by-parts} applies to them. The second is a
format input: the subtracted copies, read as values inside $Z^{\sim}$ cell by cell, are admissible
data for the production of Definition~\ref{10.3:def-vector-production-bounds}, whose proof is
outstanding, and obey its bounds; the count of part (iv) below relies on this cell-by-cell reading
for the subtracted copies. Granting both inputs, the exponent then contains, at
the one background $U_k$, the rim $\mathsf{rim}_j(\Lambda_j^{\circledR})$ together with the rim
$\mathsf{rim}_j(\Lambda_j^{\circledR,\mathrm{new}}\setminus\Lambda_j^{\circledR})$ of the added
presentation.

To compare these with $\mathsf{rim}_j(\Lambda_j^{\circledR,\mathrm{new}})$, let $\Lambda_1$,
$\Lambda_2$ be disjoint sets. In the definition of $\mathsf{rim}_j$, the bonds leaving
$\Lambda_1\cup\Lambda_2$ contribute the same terms to
$\mathsf{rim}_j(\Lambda_1)+\mathsf{rim}_j(\Lambda_2)$ and to $\mathsf{rim}_j(\Lambda_1\cup\Lambda_2)$.
A bond $b=b_\nu(y)$ with one end in $\Lambda_1$ and the other in $\Lambda_2$ contributes nothing to
$\mathsf{rim}_j(\Lambda_1\cup\Lambda_2)$. It contributes $+\mathfrak{J}_{\hat w(y)}(b;\mathsf{h}(y))$
to the rim of the set containing $y$, and $-\mathfrak{J}_{\hat w(y^+)}(b;\mathsf{h}(y^+))$ to the
rim of the set containing $y^+$. Therefore
\begin{equation}\label{10.3:eq-vector-rim-2}
\begin{split}
&\mathsf{rim}_j(\Lambda_1)+\mathsf{rim}_j(\Lambda_2)-\mathsf{rim}_j(\Lambda_1\cup\Lambda_2)\\
&\qquad=\sum_{b=b_\nu(y)\text{ joining }\Lambda_1,\Lambda_2}
\Bigl[\mathfrak{J}_{\hat w(y)}(b;\mathsf{h}(y))-\mathfrak{J}_{\hat w(y^+)}(b;\mathsf{h}(y^+))\Bigr].
\end{split}
\end{equation}

The proof of Proposition~\ref{10.3:prop-summation-by-parts} rewrites each bracket on the right as
\begin{equation}\label{10.3:eq-vector-rim-3}
\begin{split}
&\mathfrak{J}_{\hat w(y)}(b;\mathsf{h}(y))-\mathfrak{J}_{\hat w(y^+)}(b;\mathsf{h}(y^+))\\
&\qquad=-\mathfrak{J}_{\hat w(y)}(b;\nabla_{j,\nu}\mathsf{h}(y))
-\bigl[\mathfrak{J}_{\hat w(y^+)}-\mathfrak{J}_{\hat w(y)}\bigr](b;\mathsf{h}(y^+)).
\end{split}
\end{equation}
It then inserts $\nabla_{j,\nu}\mathsf{h}=L^j\langle\mathsf{k}_\nu\rangle_{j,\nu}$,
$\hat w(y^+)-\hat w(y)=a_j\langle\partial_\nu w\rangle_{j,\nu}$ and $a_j\mathsf{h}(y^+)=L^j\eta(y^+)$.
The result is the sum
$\mathsf{S}^{a,\partial\eta}_j+\mathsf{S}^{a,\eta\partial w}_j+\mathsf{D}^a_j$ of the bond terms
attached to $b$.

Take $\Lambda_1=\Lambda_j^{\circledR}$ and
$\Lambda_2=\Lambda_j^{\circledR,\mathrm{new}}\setminus\Lambda_j^{\circledR}$. These sets are
disjoint, and since $\Lambda_j^{\circledR}\subsetneq\Lambda_j^{\circledR,\mathrm{new}}$ their union
is $\Lambda_j^{\circledR,\mathrm{new}}$. Then \eqref{10.3:eq-vector-rim-2} and
\eqref{10.3:eq-vector-rim-3} give \eqref{10.3:eq-vector-rim-1} at the one background $U_k$.

Thus the exponent after healing carries $\mathsf{rim}_j(\Lambda_j^{\circledR,\mathrm{new}})$, and
the healed rim disappears from it. The subtracted copies are values inside $Z^{\sim}$, namely whole
$\sharp$-cores and rim constituents. A level $j<j(Z)-1$ has $Z\subset\Lambda_j^{\circledR}$, so
nothing is missing at such a level. Hence only levels $j\ge j(Z)-1$ occur.

(iv) Every entry in the last column of the table \eqref{10.3:eq-unit-table-a} of
Lemma~\ref{10.3:lem-unit-table} is a size per unit $n$-cell and per unit of $\sup|\mathsf{h}|$: it
is the size of one unit of that row multiplied by the number $t^{-4}$ of pinning points that a unit
$n$-cell contains at the value of $t$ of that row. By (i), the rim terms of level $j$ enter in zone
$j$ with $t=1$, so they contribute one pinning point per cell. Summing these entries over the units
of $\mathfrak{u}_k$ in zone $n$ gives the bound $C_L^V\sup|\mathsf{h}|\,T^V_n$. For the subtracted
copies of part (iii), which are values inside $Z^{\sim}$, their reading cell by cell in this count
is the format input missing at the step marked $(\ast)$, and for them the bound rests on that step.
\end{proof}

\begin{remark}[Single-constituent rim bounds and the age rate]\label{10.3:rem-rim-age-rate}
Let $\vartheta_2$, $\vartheta_V=\vartheta_2/L$ and $q$ be as in Definition~\ref{10.3:def-exponents-rates}
(the H\"older exponent, decay rates and scale sums), with $q$ subject to
\eqref{10.3:eq-exponents-rates-c}.
Here $q=e^{-c_a}$ is the age rate of the lemma on geometric forgetting of old regions, and not
Ba{\l}aban's auxiliary parameter denoted by the same letter. Since $q\le\vartheta_V^2$ is equivalent to
$c_a\ge2\log(L/\vartheta_2)$, this condition is a lower bound on $c_a$.
Let $\mathsf{u}_j$ be the norm of the frozen family of level $j$. Let $\mathsf{h}$ be the weight function.
\begin{enumerate}
\item[(a)] Suppose the near constituents of the rim $\mathsf{rim}_j(\Lambda_j^{\circledR})$ are bounded one
at a time, as in clauses (ii)--(iv) of the rim lemma of the correlation theorem. Then in zone $n$ each
rim cell of a constituent of level $j\le n$ carries the kernel $(\nu+2)\rho_{n,j}^2L^{-(\nu-1)}$, where
$\nu=n+1-j$. Summed over the levels $j\le n$ against the profile $\mathsf{u}_j\asymp\vartheta_V^j$,
this kernel with the factor $\rho_{n,j}^2$ removed is of order $nL^{-n}$, and this exceeds
$\vartheta_V^n$ by the factor $n\vartheta_2^{-n}$. Such bounds therefore cannot be used for the vector
part. The excess is produced
by the order ``split, then sum by parts''. In the order ``sum by parts, then split'' it does not arise
at the boundary $\partial Z$ of a large-field region $Z$ for the levels $j<j(Z)-1$, where $j(Z)$ is
the level attached to $Z$ in the setting of the rim lemma for the vector part
(Lemma~\ref{10.3:lem-vector-rim}, whose proof is incomplete at one step; only its setting is used
here).
\item[(b)] In zone $n<k$ the vector load exceeds the load at level $k$ by the factor
$\vartheta_V^{-(k-n)}$. The age factor $q^{(k-n-1)_+}$ with which zone $n$ enters, combined with
$q\le\vartheta_V^2$, leaves the geometric factor $\vartheta_V^{k-n}$, up to the fixed factor
$\vartheta_V^{-2}$. The prescription \eqref{10.3:eq-exponents-rates-c} lies within the hypotheses of
that lemma. The age lemma is neither used nor contradicted.
\end{enumerate}
\end{remark}

\begin{proof}
(a) Bounding the near constituents singly, by clauses (ii)--(iv) of the rim lemma of the correlation
theorem, gives per rim cell the bound
\begin{equation*}
27dM(\nu+2)t^{-3}\cdot C\mathsf{u}_j\rho^2t^4 .
\end{equation*}
Here $\rho=\rho_{n,n+1-\nu}=\rho_{n,j}$, and a constituent of level $j$ counted in zone $n$ has
$\nu=n+1-j$, as in the scale kernels $\varsigma^V_\nu$ of Definition~\ref{10.3:def-exponents-rates}.
The product equals $27dMC\mathsf{u}_j(\nu+2)\rho^2\,t$, and $t$, the relative spacing of level $j$ read
in zone $n$ in the notation of the correlation theorem, equals $a_j/a_n=L^{-(n-j)}=L^{-(\nu-1)}$. Hence,
up to the factor
$27dMC\mathsf{u}_j$, this is the kernel $(\nu+2)\rho^2L^{-(\nu-1)}$.

The weight requires the profile $\mathsf{u}_j\asymp\vartheta_V^j$. We sum the kernel, with the factor
$\rho^2_{n,j}$ removed, against this profile. Since
$\vartheta_V^j=\vartheta_2^jL^{-j}$, we have
$L^{-(n-j)}\vartheta_V^j=L^{-n}\vartheta_2^j$. Hence
\begin{align*}
\sum_{j\le n}(\nu+2)L^{-(\nu-1)}\vartheta_V^j
&\asymp L^{-n}\sum_{j\le n}(\nu+2)\vartheta_2^j\\
&\asymp nL^{-n}.
\end{align*}
The last comparison holds because $\nu+2=n-j+3$ and $0<\vartheta_2<1$. The term $j=0$ equals $n+3$,
and the whole sum is at most $(n+3)(1-\vartheta_2)^{-1}$. Since $\vartheta_V^n=\vartheta_2^nL^{-n}$,
\begin{equation*}
nL^{-n}=n\vartheta_2^{-n}\,\vartheta_V^n\gg\vartheta_V^n .
\end{equation*}

Consequently, for the kernel without the factor $\rho^2_{n,j}$, the bound so obtained on the
vector family that every large-field region of level $n$ would re-create is
of relative size $\hat\varepsilon_n\,n\,\vartheta_2^{-n}$. The geometric profile required of the vector
families in
the vector package in output form (Definition~\ref{10.5:def-vector-package-output}) would then fail.
So would, with it, the hypothesis $\sigma_k\le\mathsf{b}_1\vartheta_1^k$ of the backward-transport lemma
for the vector families, in which $\sigma_k$ denotes the size of the level-$k$ vector families that
lemma takes as input, $\mathsf{b}_1$ is a constant, and $\vartheta_1$ is
the rate of Definition~\ref{10.3:def-exponents-rates}.

This excess comes from the order in which two operations are performed. The first is the split into
first-part and second-part terms. The second is the summation by parts. If the split is made first,
the excess arises. If the summation by parts is made first, it is not there. Indeed, consider a
current of level $j$ and its fluxes through $\partial Z$ from the two sides of $\partial Z$. These two
fluxes belong to one telescoped sum as long as the family of level $j$ exists on both sides. They then
cancel inside that sum, and no term is left at $\partial Z$. The family exists on both sides exactly
for $j<j(Z)-1$.

(b) The load in zone $n<k$ is $\sup|\mathsf{h}|\,T^V_n\asymp\vartheta_V^n$. This is larger than the load
at level $k$ by
\begin{equation*}
\vartheta_V^{-(k-n)}=\bigl(L\,\vartheta_2^{-1}\bigr)^{k-n} .
\end{equation*}
So each level of age contributes one factor $L$ from the weight and one factor $\vartheta_2^{-1}$ from
the profile.

Zone $n$ enters clauses (d) and (e) of the recursion for the loads $[\![\hat\delta']\!]^T_k$ with the
factor $q^{(k-n-1)_+}$. Put
$i=k-n\ge1$ and use $q\le\vartheta_V^2$, which is part of \eqref{10.3:eq-exponents-rates-c}. Then
\begin{align*}
q^{(i-1)_+}\vartheta_V^{-i}
&\le\vartheta_V^{2(i-1)-i}\\
&=\vartheta_V^{-2}\,\vartheta_V^{\,i}.
\end{align*}
For $i=1$ both sides equal $\vartheta_V^{-1}$. Thus the factor $\vartheta_V^{k-n}$ remains, multiplied by
the $k$-independent constant $\vartheta_V^{-2}$.

This prescription lies within the hypotheses of the lemma on geometric forgetting of old regions. There
the rate
$q=e^{-c_a}$ is prescribed, not derived. The places at which its proof constrains the parameters are
degree inequalities on the exponent $p_0$ of the degree function $p_0(\cdot)$, followed by a ceiling on
$\gamma$, and two upper bounds on $q$ by the squares of two quantities of that lemma. Only the latter
two concern $q$; they fix $c_a$ after $L$ and impose no lower bound on $q$. A further upper bound on $q$,
such as $q\le\vartheta_V^2$, is therefore admissible. In \eqref{10.3:eq-exponents-rates-c} the number
$c_a$ is chosen in this way, after $L$ and $\vartheta_2$ and before $\gamma$, so that the ceiling on
$\gamma$ is imposed after $c_a$ is fixed.

The age lemma is a statement about the probability distribution of the large-field regions; it does
not supply a separate factor for each scale. The proofs of the modulus bound (the bound on the absolute
values of the factors of the large-field expansion) and of the activity
bound $\mathsf{W}$ bound absolute values factor by factor under the sum over sequences of large-field
regions (the histories $h$); clauses (d)
and (e) of the recursion for the loads $[\![\hat\delta']\!]^T_k$ therefore invoke the lemma on geometric
forgetting of old regions only inside bounds on absolute values, factor by factor. Nothing of the age
lemma is used here, and nothing here contradicts it.
\end{proof}

\begin{lemma}[Re-filing of the vector remnants]\label{10.3:lem-vector-refiling}
Let $i$ be a level and $i''$ its re-filing level. Let $\mathsf{R}^{V(i)}(\cdot;w,\mathsf{h})$ be the local marginal
family of level $i$ in the vector part of the run (Definition~\ref{10.3:def-vector-run}), and let
$\mathsf{r}_i=\|\mathsf{R}^{V(i)}\|^{\flat}$. Assume the following two statements.
\begin{itemize}
\item[(A)] The re-filing lemma for a marginal remnant $R$ (one of the re-filing lemmas), parts
(a)--(c), with its constant $C_{\mathrm{RF}}$ and together with its equivariance condition.
\item[(B)] The typed form of the re-filing lemmas, taken at the type $\mathsf{V}=V$.
\end{itemize}
Then parts (a)--(c) of the re-filing lemma hold for $R:=\mathsf{R}^{V(i)}(\cdot;w,\mathsf{h})$. In detail,
for every anchor $K''$, every localisation region $Y$ at level $i''$ and every pin $y$,
the re-filed family $\mathsf{RF}[\mathsf{R}^{V(i)}](Y,\cdot,y;w,\mathsf{h})$ has the following properties.
\begin{enumerate}
\item It is linear in $\mathsf{h}{\restriction}Y^+$.
\item It is analytic on $\mathcal{W}^{[K'']}_{i''}(Y)$.
\item It is equal in total to the vacuum part $(\mathsf{R}^{V(i)})^{\mathrm{vac}}$ of
$\mathsf{R}^{V(i)}$, in the sense of the re-filing lemma.
\item It is $V$-typed at constant source and constant weight, $(w,\mathsf{h})\equiv(\hat w,e)$
with $\hat w$ and $e$ constant, that is, its pointed totals satisfy $(\mathrm{f}3)^{\mathsf{V}}$
with $\mathsf{V}=V$ for every Euclidean symmetry of $T^{(i'')}$.
\item Its fresh norm obeys, with the constant $C_{\mathrm{RF}}$ of hypothesis (A),
\begin{equation}\label{10.3:eq-vector-refiling-1}
\bigl\|\mathsf{RF}[\mathsf{R}^{V(i)}]\bigr\|^{\mathrm{fr}}\le C_{\mathrm{RF}}\,w_{i''-1,i}\,\mathsf{r}_i .
\end{equation}
\end{enumerate}
\end{lemma}

\begin{proof}
Parts (a) and (c) of the re-filing lemma are statements about the family term by term, and they are
linear in the family. Hypothesis (A) therefore carries them over to $R=\mathsf{R}^{V(i)}$ term by
term.

Part (b) of the re-filing lemma contains the covariance property $(\mathrm{f}3)$. Under (A) this
property is proved from \cite[(2.32)]{B-CMP119} and from the equivariance condition of the
re-filing lemma. Both of these are
equivariance statements. By hypothesis (B), equivariance statements pass to the typed form at
$\mathsf{V}=V$.

At this type no coefficient is extracted. The reason is the property of $V$ recorded in
Definition~\ref{10.3:def-typed-families}: frozen $V$-typed families have no vacuum part and no
coefficient $\beta$. Hence $(\mathrm{f}3)$ holds in its typed form $(\mathrm{f}3)^{\mathsf{V}}$ at
constant source and constant weight, for every Euclidean symmetry of $T^{(i'')}$.

The norm in part (b) is estimated by the norm lemma of the correlation theorem's format.

For the source classes, the Lipschitz norms
of $w$ at levels $i$ and $i''$ satisfy
\begin{equation}\label{10.3:eq-vector-refiling-2}
\operatorname{Lip}_i w = L^{-\mathsf{a}_i}\operatorname{Lip}_{i''} w
< L^{-\mathsf{a}_i/2}\,\mathfrak{l}^{[K'']}_i ,
\end{equation}
which is the class condition used in part (c) of the re-filing lemma.
\end{proof}

\begin{remark}
The re-filing cannot be omitted. Suppose a remnant $\mathsf{R}^{V(i)}$ were left unfiled. It would
then be read at a later level $n$ by the norm lemma of the correlation theorem's format. That lemma
reads it with the flat kernel $w_{n,i}$, against the factor $\vartheta_V^{n-i}$.

Such a remnant can be tolerated for $n-i\le\mathsf{a}_i$. Over that range the loss is a factor
\begin{equation*}
\vartheta_V^{-\mathsf{a}_i}\le\vartheta_V^{-1}(M\check R_i)^2 ,
\end{equation*}
which is a power of $\log x_i$ and is set against $\hat\varepsilon_i$. In the flat norms of the
correlation theorem the recursion for the vector part does not close.

For this reason the remnant is carried with the whole weight function $\mathsf{h}$ up to level
$i''$. Only at level $i''$ is the decomposition of a local output at its pin
(Proposition~\ref{10.3:prop-output-decomposition}) applied to it, inside $\hat{\mathsf{O}}_{i''}$.
Consequently every tensor-typed source arises at level $i''$ as one of the pin sources
$\mathsf{S}^{b,\partial\eta}_{i''}$, $\mathsf{S}^{b,\eta\partial w}_{i''}$ of
Proposition~\ref{10.3:prop-output-decomposition}, that is, as a source term of
$\hat{\mathsf{O}}_{i''}$.
\end{remark}

We fix the constants of the vector recursion. Let
$\bar\varkappa:=\sup_k\varkappa^V_k$. For $n\ge 1$ let
\begin{equation*}
\tilde\varepsilon_n:=\max_{1\le i\le n}\hat\varepsilon^V_i\,(1+\mathsf{t}_0\mathsf{s}_i).
\end{equation*}
Let $\mathsf{C}:=2(C_q^V+1+\mathsf{s}_*)$, and
\begin{equation*}
c':=10\,C_{\mathsf{RF}}\,\mathsf{C}\,\sup_n\tilde\varepsilon_n\vartheta_V^{-\mathsf{a}^*_n}.
\end{equation*}
Here $\varkappa^V_k$, $\hat\varepsilon^V_k$ and $\epsilon^V_k$ are the step constants of the vector part, fixed
with the production bounds for the vector part under the split rule
(stated in~\ref{10.3:def-vector-production-bounds}; proof outstanding). The quantities
$T^V_n$, $T^{V,T}_k$, $\mathsf{t}_0$ are the vector loads of Definition~\ref{10.3:def-vector-run}.
The constants $\varsigma^V_\nu$, $S^V$, $C_q^V$, $\mathsf{s}_n$, $\mathsf{s}_*$ and the rate
$\vartheta_V=\vartheta_2/L$ are those of Definition~\ref{10.3:def-exponents-rates}. Finally,
$C_{\mathsf{RF}}$ is the constant of the re-filing map $\mathsf{RF}[\cdot]$ of the re-filing lemmas.
We abbreviate $\mathsf{v}_0:=\mathsf{v}_0(c_0)$.

We impose the following ceilings on the coupling $g$. Each is one-sided, and they are imposed after $L$, $\mathfrak{p}$, $\hat\alpha$, $\vartheta_2$ and $q$ have been fixed. For all $n$,
\begin{equation}\label{10.3:eq-vector-run-decay-a}
\begin{aligned}
&(\mathrm{v}1)\quad (\bar\varkappa+c')\,S^V\le\tfrac18;\\
&(\mathrm{v}2)\quad 5\,\mathsf{C}\tilde\varepsilon_n\,\mathsf{a}^*_n\vartheta_V^{-\mathsf{a}^*_n}\le\tfrac18,
\qquad \tfrac52\,\mathsf{C}\tilde\varepsilon_n\le\tfrac14,\qquad \tfrac52\,\epsilon^V_n\le\tfrac18;\\
&(\mathrm{v}3)\quad 3\,\varkappa^V_n\,\bigl(n-k_1(n)\bigr)\,\vartheta_V^{-(n-k_1(n))}\le\tfrac18;\\
&(\mathrm{v}4)\quad 4S^V(\bar\varkappa+c')\le\vartheta_V .
\end{aligned}
\end{equation}

The ceilings can be met uniformly in $K$: each left-hand side in \eqref{10.3:eq-vector-run-decay-a} tends to zero with $g$, uniformly in $K$. We check this term by term.

First, $\varkappa^V_n$ is bounded by a constant times $(\log x_n)^{p_0-q_{\flat}}$. The numbers $n-k_1(n)$ and $\mathsf{a}^*_n$ are at most $1+\log_L(M\check R_n)$. Since $\vartheta_V\in\,]L^{-2\hat\alpha},L^{-1}[$ by Definition~\ref{10.3:def-exponents-rates}, we have $\log_L(1/\vartheta_V)<2\hat\alpha<2$. Hence
\begin{equation*}
\vartheta_V^{-\mathsf{a}^*_n}\le\vartheta_V^{-1}\,\vartheta_V^{-\log_L(M\check R_n)}
=\vartheta_V^{-1}(M\check R_n)^{\log_L(1/\vartheta_V)}\le\vartheta_V^{-1}(M\check R_n)^2 .
\end{equation*}
Here $r_{\mathrm{pr}}$ is the exponent of the correlation theorem in the bound
$M\check R_n\le C(\log x_n)^{r_{\mathrm{pr}}}$, with $C$ a constant; the right-hand side is therefore
at most a constant times $(\log x_n)^{2r_{\mathrm{pr}}}$. These powers of $\log x_n$ are absorbed
under the floor $q_{\flat}\ge p_0+7r_{\mathrm{pr}}+1$ fixed in the correlation theorem, so no further
condition on $q_{\flat}$ is needed.

Second, $\hat\varepsilon^V_i$ is bounded by a constant times $x_i^{-\kappa_0/2}$ times a power of
$\log x_i$, where $\kappa_0\ge7$ is the exponent fixed in the correlation theorem that governs this
smallness in $x_i$. By Definition~\ref{10.3:def-vector-run},
$\mathsf{t}_0\le C_{\mathsf{t}}g_0^{-1}=C_{\mathsf{t}}x_0^{1/2}$. The coupling flow of the correlation
theorem, with its flow constant denoted $\lambda$ (not the dial $\lambda$ of
Definition~\ref{10.3:def-vector-run}), then gives
\begin{equation*}
\mathsf{t}_0\le C_{\mathsf{t}}\,x_0^{1/2}\le C_{\mathsf{t}}\bigl[x_i^{1/2}+(3\lambda i+2c_2)^{1/2}\bigr].
\end{equation*}
Since $\mathsf{s}_i\le\mathsf{s}_*$ and $\mathsf{s}_i(3\lambda i+2c_2)^{1/2}\le\mathsf{s}_*$ by the definition of $\mathsf{s}_*$, it follows that
\begin{equation*}
\hat\varepsilon^V_i\,\mathsf{t}_0\mathsf{s}_i\le C_{\mathsf{t}}\mathsf{s}_*\,\hat\varepsilon^V_i\,(1+x_i^{1/2}).
\end{equation*}
The right-hand side is at most a constant times $x_i^{-(\kappa_0-1)/2}$ times a power of $\log x_i$. Thus the factor $x_0$ carried by $\mathsf{t}_0$ is compensated level by level, because $\kappa_0\ge 7$ in the correlation theorem.

\begin{lemma}[Geometric decay of the frozen vector run]\label{10.3:lem-vector-run-decay}
Assume the bounds \eqref{10.3:eq-vector-production-bounds-d} and
\eqref{10.3:eq-vector-production-bounds-e} of the production bounds for the vector part under
the split rule (stated in~\ref{10.3:def-vector-production-bounds}; proof outstanding),
and the ceilings \eqref{10.3:eq-vector-run-decay-a}. Put $\mathsf{A}_V:=4S^V\mathsf{v}_0(c_0)$.
Then for all $1\le n\le K'$
\begin{equation}\label{10.3:eq-vector-run-decay-1}
\begin{gathered}
T^V_n\le\mathsf{A}_V\vartheta_V^n,\qquad
\mathsf{r}_n\le\mathsf{C}\tilde\varepsilon_n\mathsf{A}_V\vartheta_V^n,\\
\mathsf{u}_n\le\mathsf{o}_n\le(\bar\varkappa+c')\mathsf{A}_V\vartheta_V^{n-1}
\le\mathsf{v}_0(c_0)\vartheta_V^n ,
\end{gathered}
\end{equation}
and $T^V_0\le(2+\mathsf{t}_0)\mathsf{v}_0(c_0)$.
\end{lemma}

\begin{proof}
\emph{Preliminaries.} Since $\rho_{n,n}=1$, the definition of the scale kernels gives
\begin{equation*}
\varsigma^V_1=1^3\cdot1^2\cdot(1+1)\cdot L^{0}=2 .
\end{equation*}
All terms of the sum defining $S^V$ are non-negative, and its $\nu=1$ term is $\varsigma^V_1\vartheta_V^{-1}$. Hence $S^V\ge 2/\vartheta_V$, and since $\vartheta_V<1$,
\begin{equation}\label{10.3:eq-vector-run-decay-2}
\mathsf{A}_V=4S^V\mathsf{v}_0\ge\frac{8\mathsf{v}_0}{\vartheta_V}\ge 8\,\mathsf{v}_0 .
\end{equation}

\emph{The load at level $0$.} At $n=0$ the definition of the loads in
Definition~\ref{10.3:def-vector-run} has the bare term $\varsigma^V_1\mathsf{u}_0$ and the raw term
$\mathsf{t}^{\mathrm{raw}}$. Here $\mathsf{u}_0=\mathsf{v}_0$, the norm of the bare vector unit, and
$\mathsf{t}^{\mathrm{raw}}=\mathsf{t}_0\mathsf{v}_0$ by the definition of $\mathsf{t}_0$ as
$\mathsf{t}^{\mathrm{raw}}/\mathsf{v}_0(c_0)$. Since $\varsigma^V_1=2$,
\begin{equation*}
\varsigma^V_1\mathsf{u}_0+\mathsf{t}^{\mathrm{raw}}
=2\mathsf{v}_0+\mathsf{t}_0\mathsf{v}_0=(2+\mathsf{t}_0)\mathsf{v}_0 ,
\end{equation*}
and the last assertion of the lemma, $T^V_0\le(2+\mathsf{t}_0)\mathsf{v}_0$, is the bound of the
load at level $0$ by these two terms.

\emph{Induction claims.} We prove by induction on $n\ge0$ the three statements:
\begin{itemize}
\item[$(\mathrm{I}_n)$] $T^V_i\le\mathsf{A}_V\vartheta_V^i$ for $1\le i\le n$;
\item[$(\mathrm{II}_n)$] $\mathsf{r}_i\le\mathsf{C}\tilde\varepsilon_i\mathsf{A}_V\vartheta_V^i$ for $1\le i\le n$;
\item[$(\mathrm{III}_n)$] $\mathsf{o}_i\le(\bar\varkappa+c')\mathsf{A}_V\vartheta_V^{i-1}$ for
$1\le i\le n+1$.
\end{itemize}
For $n=0$ the statements $(\mathrm{I}_0)$ and $(\mathrm{II}_0)$ are empty. For $(\mathrm{III}_0)$
we use the hypothesis \eqref{10.3:eq-vector-production-bounds-d} (one of the production bounds,
whose proof is outstanding) at $k=0$. By the definition of the loads, $T^{V,T}_0$ carries no raw
term and its bare term is $\varsigma^V_1\mathsf{u}_0=2\mathsf{v}_0$; moreover
$\varkappa^V_0\le\bar\varkappa$ by the definition of $\bar\varkappa$. So
\eqref{10.3:eq-vector-production-bounds-d} at $k=0$ gives
$\mathsf{o}_1\le\bar\varkappa\,\varsigma^V_1\mathsf{v}_0$. By \eqref{10.3:eq-vector-run-decay-2},
\begin{equation*}
\bar\varkappa\,\varsigma^V_1\mathsf{v}_0=2\bar\varkappa\mathsf{v}_0\le\bar\varkappa\mathsf{A}_V
\le(\bar\varkappa+c')\mathsf{A}_V\vartheta_V^{0}.
\end{equation*}

Let $n\ge1$ and assume $(\mathrm{I}_{n-1})$, $(\mathrm{II}_{n-1})$ and $(\mathrm{III}_{n-1})$. Put $T^-:=T^V_n-w_{n,n}\mathsf{r}_n$.

\emph{Step (a): the load without its top remnant.} By the definition of the load at $n\ge1$, where there is no raw term, and with $\mathsf{u}_0=\mathsf{v}_0$, the norm of the bare vector unit,
\begin{equation*}
T^-=\varsigma^V_{1+n}\mathsf{v}_0+\sum_{1\le j\le n}\varsigma^V_{1+n-j}\mathsf{u}_j
+\sum_{n-\mathsf{a}_i\le i<n}w_{n,i}\mathsf{r}_i .
\end{equation*}
We bound the three terms in turn.

\emph{First term.} The term $\varsigma^V_{1+n}\vartheta_V^{-(n+1)}$ is one term of the sum defining $S^V$. Hence
\begin{equation*}
\varsigma^V_{1+n}\mathsf{v}_0\le S^V\vartheta_V^{n+1}\mathsf{v}_0
=\tfrac14\vartheta_V\mathsf{A}_V\vartheta_V^n .
\end{equation*}

\emph{Second term.} For $1\le j\le n$ we have $\mathsf{u}_j\le\mathsf{o}_j$, and $(\mathrm{III}_{n-1})$ gives $\mathsf{o}_j\le(\bar\varkappa+c')\mathsf{A}_V\vartheta_V^{j-1}$. Substituting $\nu=1+n-j$, so that $\vartheta_V^{j-1}=\vartheta_V^n\vartheta_V^{-\nu}$, we obtain
\begin{align*}
\sum_{1\le j\le n}\varsigma^V_{1+n-j}\mathsf{u}_j
&\le(\bar\varkappa+c')\mathsf{A}_V\sum_{1\le j\le n}\varsigma^V_{1+n-j}\vartheta_V^{j-1}
=(\bar\varkappa+c')\mathsf{A}_V\vartheta_V^n\sum_{1\le\nu\le n}\varsigma^V_\nu\vartheta_V^{-\nu}\\
&\le(\bar\varkappa+c')S^V\mathsf{A}_V\vartheta_V^n\le\tfrac18\mathsf{A}_V\vartheta_V^n ,
\end{align*}
where the last inequality is (v1).

\emph{Third term.} On the range $n-\mathsf{a}_i\le i<n$ we have $w_{n,i}\le5$ (the weights
$w_{k,j}$ are those inherited from the correlation theorem, which satisfy $w_{n,i}\le5$ on this
range and $w_{n,n}\le\frac52$). The range contains at most $\mathsf{a}^*_n$ indices, and on it $n-i\le\mathsf{a}^*_n$. By $(\mathrm{II}_{n-1})$, and because $\tilde\varepsilon$ is nondecreasing in its index by its definition as a maximum,
\begin{equation*}
\mathsf{r}_i\le\mathsf{C}\tilde\varepsilon_n\mathsf{A}_V\vartheta_V^n\vartheta_V^{-(n-i)}
\le\mathsf{C}\tilde\varepsilon_n\vartheta_V^{-\mathsf{a}^*_n}\mathsf{A}_V\vartheta_V^n .
\end{equation*}
Hence, by the first inequality of (v2),
\begin{equation*}
\sum_{n-\mathsf{a}_i\le i<n}w_{n,i}\mathsf{r}_i
\le5\,\mathsf{C}\tilde\varepsilon_n\mathsf{a}^*_n\vartheta_V^{-\mathsf{a}^*_n}\mathsf{A}_V\vartheta_V^n
\le\tfrac18\mathsf{A}_V\vartheta_V^n .
\end{equation*}

Since $\vartheta_V<1$, the three bounds together give
\begin{equation}\label{10.3:eq-vector-run-decay-3}
T^-\le\bigl(\tfrac14\vartheta_V+\tfrac18+\tfrac18\bigr)\mathsf{A}_V\vartheta_V^n
\le\tfrac12\mathsf{A}_V\vartheta_V^n .
\end{equation}

\emph{Step (b): the remnant at level $n$.} The bound \eqref{10.3:eq-vector-production-bounds-e} at $k=n$ reads
\begin{equation*}
\mathsf{r}_n\le\epsilon^V_n\sum_{0\le i\le n}(n-i+1)\,q^{(n-i-1)_+}\,T^V_i
+3\varkappa^V_n\sum_{k_1(n)<i<n}\mathsf{r}_i .
\end{equation*}
We split the first sum according to $i$.

\emph{Indices $1\le i<n$.} By $(\mathrm{I}_{n-1})$, and substituting $m=n-i$,
\begin{equation*}
\sum_{1\le i<n}(n-i+1)q^{(n-i-1)_+}T^V_i
\le\mathsf{A}_V\vartheta_V^n\sum_{1\le m\le n-1}(m+1)q^{(m-1)_+}\vartheta_V^{-m}
\le C_q^V\mathsf{A}_V\vartheta_V^n .
\end{equation*}

\emph{Index $i=n$.} The coefficient is $1$, and the term is $T^V_n=T^-+w_{n,n}\mathsf{r}_n$. With $w_{n,n}\le\frac52$ and \eqref{10.3:eq-vector-run-decay-3}, this is at most $\frac12\mathsf{A}_V\vartheta_V^n+\frac52\mathsf{r}_n$.

\emph{Index $i=0$.} The choice of $q$ in Definition~\ref{10.3:def-exponents-rates} gives
$q^{(n-1)_+}\le\vartheta_V^{2(n-1)}$. Moreover $(n+1)\vartheta_V^{2(n-1)}=\mathsf{s}_n\vartheta_V^n$,
since $\mathsf{s}_n=(n+1)\vartheta_V^{n-2}$. We use the bound $T^V_0\le(2+\mathsf{t}_0)\mathsf{v}_0$
computed above from the load at level $0$, then $\mathsf{v}_0\le\mathsf{A}_V/8$ from
\eqref{10.3:eq-vector-run-decay-2}, and finally $\mathsf{s}_n\le\mathsf{s}_*$. This gives
\begin{align*}
(n+1)q^{(n-1)_+}T^V_0
&\le(n+1)\vartheta_V^{2(n-1)}(2+\mathsf{t}_0)\mathsf{v}_0
=\mathsf{s}_n\vartheta_V^n(2+\mathsf{t}_0)\mathsf{v}_0\\
&\le\tfrac18\bigl(2\mathsf{s}_n+\mathsf{t}_0\mathsf{s}_n\bigr)\mathsf{A}_V\vartheta_V^n
\le\tfrac14\bigl(\mathsf{s}_*+\mathsf{t}_0\mathsf{s}_n\bigr)\mathsf{A}_V\vartheta_V^n .
\end{align*}

\emph{The band.} For $k_1(n)<i<n$, $(\mathrm{II}_{n-1})$ and the monotonicity of $\tilde\varepsilon$ give
\begin{equation*}
\mathsf{r}_i\le\mathsf{C}\tilde\varepsilon_n\mathsf{A}_V\vartheta_V^{i}
\le\mathsf{C}\tilde\varepsilon_n\mathsf{A}_V\vartheta_V^{n}\vartheta_V^{-(n-k_1(n))} .
\end{equation*}
There are at most $n-k_1(n)$ such indices. Therefore, by (v3),
\begin{equation*}
3\varkappa^V_n\sum_{k_1(n)<i<n}\mathsf{r}_i
\le3\varkappa^V_n\bigl(n-k_1(n)\bigr)\vartheta_V^{-(n-k_1(n))}\,
\mathsf{C}\tilde\varepsilon_n\mathsf{A}_V\vartheta_V^n
\le\tfrac18\,\mathsf{C}\tilde\varepsilon_n\mathsf{A}_V\vartheta_V^n .
\end{equation*}

\emph{Combination.} Collecting the four contributions, and moving the term $\frac52\epsilon^V_n\mathsf{r}_n$ to the left by means of the third inequality of (v2), $\frac52\epsilon^V_n\le\frac18$, we get
\begin{equation*}
\bigl(1-\tfrac18\bigr)\mathsf{r}_n\le\epsilon^V_n\Bigl[C_q^V+\tfrac12
+\tfrac14\bigl(\mathsf{s}_*+\mathsf{t}_0\mathsf{s}_n\bigr)\Bigr]\mathsf{A}_V\vartheta_V^n
+\tfrac18\,\mathsf{C}\tilde\varepsilon_n\mathsf{A}_V\vartheta_V^n .
\end{equation*}
Since $\epsilon^V_n(1+\mathsf{t}_0\mathsf{s}_n)\le\tilde\varepsilon_n$, both $\epsilon^V_n\le\tilde\varepsilon_n$ and $\epsilon^V_n\mathsf{t}_0\mathsf{s}_n\le\tilde\varepsilon_n$ hold. By $\mathsf{C}=2(C_q^V+1+\mathsf{s}_*)$ we also have $C_q^V+\frac12+\frac14(\mathsf{s}_*+1)\le\frac12\mathsf{C}$. Hence
\begin{equation*}
\tfrac78\,\mathsf{r}_n\le\tilde\varepsilon_n\mathsf{A}_V\vartheta_V^n
\Bigl[C_q^V+\tfrac12+\tfrac14(\mathsf{s}_*+1)+\tfrac18\mathsf{C}\Bigr]
\le\tfrac58\,\mathsf{C}\tilde\varepsilon_n\mathsf{A}_V\vartheta_V^n .
\end{equation*}
Thus
\begin{equation*}
\mathsf{r}_n\le\tfrac57\,\mathsf{C}\tilde\varepsilon_n\mathsf{A}_V\vartheta_V^n
\le\mathsf{C}\tilde\varepsilon_n\mathsf{A}_V\vartheta_V^n .
\end{equation*}
Together with $(\mathrm{II}_{n-1})$ this is $(\mathrm{II}_n)$.

\emph{Step (c): the load at level $n$.} By \eqref{10.3:eq-vector-run-decay-3}, $w_{n,n}\le\frac52$, step (b), and the second inequality of (v2),
\begin{equation*}
T^V_n=T^-+w_{n,n}\mathsf{r}_n\le\tfrac12\mathsf{A}_V\vartheta_V^n
+\tfrac52\,\mathsf{C}\tilde\varepsilon_n\mathsf{A}_V\vartheta_V^n
\le\bigl(\tfrac12+\tfrac14\bigr)\mathsf{A}_V\vartheta_V^n .
\end{equation*}
Together with $(\mathrm{I}_{n-1})$ this is $(\mathrm{I}_n)$.

\emph{Step (d): the output at level $n+1$.} The bound \eqref{10.3:eq-vector-production-bounds-d} at $k=n$ gives
\begin{equation*}
\mathsf{o}_{n+1}\le\varkappa^V_nT^{V,T}_n+C_{\mathsf{RF}}\sum_{i\in\mathfrak{B}(n)}w_{n,i}\mathsf{r}_i .
\end{equation*}
In the first term we use $\varkappa^V_n\le\bar\varkappa$ and $T^{V,T}_n\le\mathsf{A}_V\vartheta_V^n$.
In the second we use that at most two indices lie in $\mathfrak{B}(n)$, by the first property of the
index sets $\mathfrak{B}(n)$ in the rate statement, together with $w_{n,i}\le5$ and the bound
$\mathsf{r}_i\le\mathsf{C}\tilde\varepsilon_n\vartheta_V^{-\mathsf{a}^*_n}\mathsf{A}_V\vartheta_V^n$
obtained as in step (a) from $(\mathrm{II}_n)$. By the definition of $c'$,
\begin{equation*}
\mathsf{o}_{n+1}\le\bar\varkappa\mathsf{A}_V\vartheta_V^n
+C_{\mathsf{RF}}\cdot2\cdot5\cdot\mathsf{C}\tilde\varepsilon_n\vartheta_V^{-\mathsf{a}^*_n}
\mathsf{A}_V\vartheta_V^n\le(\bar\varkappa+c')\mathsf{A}_V\vartheta_V^n .
\end{equation*}
Together with $(\mathrm{III}_{n-1})$ this is $(\mathrm{III}_n)$, and the induction is complete.

\emph{Conclusion.} Statements $(\mathrm{I}_{K'})$, $(\mathrm{II}_{K'})$ and $(\mathrm{III}_{K'})$,
together with $\mathsf{u}_n\le\mathsf{o}_n$, give every inequality of
\eqref{10.3:eq-vector-run-decay-1} except the last; the bound on $T^V_0$ is the computation of
the load at level $0$ above. The last inequality of \eqref{10.3:eq-vector-run-decay-1} follows from
(v4) and the definition of $\mathsf{A}_V$:
\begin{equation*}
(\bar\varkappa+c')\mathsf{A}_V\vartheta_V^{n-1}
=4S^V(\bar\varkappa+c')\,\mathsf{v}_0\vartheta_V^{n-1}\le\mathsf{v}_0\vartheta_V^n .
\end{equation*}
\end{proof}

\begin{remark}
The frozen vector family of level $j$ can be read in its own-scale units, that is, with weight $\eta/a_j$ and norm $L^j\mathsf{u}_j$. Since $\vartheta_V=\vartheta_2/L$, Lemma~\ref{10.3:lem-vector-run-decay} gives
\begin{equation*}
L^j\mathsf{u}_j\le L^j\mathsf{v}_0\vartheta_V^j=\mathsf{v}_0\vartheta_2^j .
\end{equation*}
So in these units the frozen vector family of level $j$ has norm at most $\mathsf{v}_0\vartheta_2^j$.
The obstruction met by the scalar run, a bound growing like $L^{(1-\hat\alpha)K}$, comes from the
recursion for the scalar hatted run, which is the same recursion with its scale kernel
$\varsigma_\nu$ in the place of $\varsigma^V_\nu$; for that kernel no rate $\vartheta_V<L^{-1}$ is
admissible.
\end{remark}

\begin{lemma}[The load of a variation of the source]\label{10.3:lem-source-variation}
Let $K''\le K'$.
Assume the following.
\begin{itemize}
\item Parts (a), (c) and (e) of the hatted one-lattice Lipschitz step hold; part (e) is stated in
this book with its proof incomplete at one step, the step at which it rests on the carry lemma.
\item The re-anchoring of the source classes at each level,
Proposition~\ref{10.3:prop-reanchoring}, holds.
\item With the age rate $q$, the step constants $\hat\varkappa_k$, $\hat\varepsilon_k$ and the
re-filing constant $C_{\mathsf{RF}}$ of the hatted run, and with $\epsilon_k$ the step constant
with which parts (c) and (e) of the hatted one-lattice Lipschitz step are stated, put
\begin{equation}\label{10.3:eq-source-variation-1}
C_q^0:=\sum_{i\ge 0}(i+1)\,q^{(i-1)_+},\qquad \mathsf{C}'':=8C_q^0,\qquad
c_e:=3\sup_k\hat\varkappa_k+10\,C_{\mathsf{RF}}\,\mathsf{C}''\sup_i\hat\varepsilon_i ,
\end{equation}
and, with $S_{\varsigma}$ the $K$-free sum of the scale kernels of the hatted load recursion,
assume, for every $n\le K''$, the ceilings
\begin{equation}\label{10.3:eq-source-variation-2}
\begin{gathered}
S_{\varsigma}\,c_e\le\tfrac14,\qquad 5\,\mathsf{a}^*_n\,\mathsf{C}''\,\hat\varepsilon_n\le\tfrac18,\\
\tfrac52\,\mathsf{C}''\,\hat\varepsilon_n\le\tfrac18,\qquad
3\,\hat\varkappa_n\bigl(n-k_1(n)\bigr)\le\tfrac18,\\
\epsilon_n\le 2\,\hat\varepsilon_n .
\end{gathered}
\end{equation}
\end{itemize}
Let $w$ lie in the half-interior of $\mathcal{W}^{[K'']}_{K''}$, that is,
\begin{equation*}
|w|\le r/4,\qquad \operatorname{Lip}_{K''}w\le r\lambda_*/4 .
\end{equation*}
Let $\psi$ be a function with
\begin{equation*}
\|\psi\|_{K''}:=\max\{\sup|\psi|,\ \operatorname{Lip}_{K''}\psi/\lambda_*\}\le 1 ,
\end{equation*}
and let
$\delta^{\psi}:=D_w\hat{\mathfrak{s}}(w)[\psi]$ be the tangent of the hatted run in the direction
$\psi$ of the source. Then for every $n\le K''$
\begin{equation}\label{10.3:eq-source-variation-3}
[\![\hat\delta^{\psi}]\!]_n\le 3\,\mathfrak{g}^{\sharp}_n,\qquad
\mathfrak{g}^{\sharp}_n:=\max_{i\le n}\mathfrak{g}_i\le 2\,\mathfrak{g}_n .
\end{equation}
\end{lemma}

\begin{proof}
Let $\sigma\in\mathbb{C}$ with $|\sigma|<r/2$. Since $\sup|\psi|\le 1$ and
$\operatorname{Lip}_{K''}\psi\le\lambda_*$,
\begin{equation*}
|w+\sigma\psi|<\tfrac r4+\tfrac r2=\tfrac{3r}4,\qquad
\operatorname{Lip}_{K''}(w+\sigma\psi)<\tfrac{r\lambda_*}4+\tfrac{r\lambda_*}2=\tfrac{3r\lambda_*}4 .
\end{equation*}
At its own level the class $\mathcal{W}^{[K'']}_{K''}$ is $\mathcal{W}^{\circ}_{K''}$, whose members
satisfy $|\cdot|<5r/4$ and $\operatorname{Lip}_{K''}(\cdot)<r\lambda_*$; so $w+\sigma\psi$ lies in
it. By Proposition~\ref{10.3:prop-reanchoring}, applied at the anchor $K''$, the terms of every
level of the sourced format are analytic on the source classes of anchor $K''$, and the hatted run
$\hat{\mathfrak{s}}(w+\sigma\psi)$ lies, for every $k$ and every $|\sigma|<r/2$, in the class
$\hat{\mathfrak{K}}_k$ of hatted runs on which the hatted one-lattice Lipschitz step is stated.
The direction $\delta^{\psi}$ is the derivative in $\sigma$ at $\sigma=0$ of this
analytic family. The parts of the hatted one-lattice Lipschitz step are bounds on differences of two
hatted runs. The same parts hold for the derivative $\delta^{\psi}$; this follows from the lemma
that transfers difference bounds to derivatives, applied with a larger value of $\varrho$, and from
the corresponding lemma for the loads.

We list the components of $\delta^{\psi}$ that enter its loads. At level $0$ it has no component
other than the weight. Its component along the hatted running shift $\hat\zeta_k$ is
$\delta\hat\zeta_k$. Differentiating the correlation theorem's formula for the running shift,
with the hatted quantities $\hat\zeta_k$, $\hat b_j$, $\hat w$ in place of $\zeta_k$, $b_j$, $w$,
gives
\begin{equation}\label{10.3:eq-source-variation-4}
\delta\hat\zeta_k(p)=\psi(p)+\sum_j\varphi_j\sum_y h_y\,\hat b_j'\bigl(\hat w(y)\bigr)\,\psi(p_y),
\end{equation}
where $\varphi_j$ and $h_y$ are the coefficients of that formula and $p_y$ is the plaquette at
which the source is frozen at $y$.
Two results bound the size $\hat{\mathsf{z}}_k:=\sup|\delta\hat\zeta_k|$ of this component. The
first is the running-shift statement, for the hatted quantities. The second consists of the
summability bounds in the corollary on the rates of the run. Together they give
\begin{equation}\label{10.3:eq-source-variation-5}
\hat{\mathsf{z}}_k\le 1+\sum_j\operatorname{Lip}\hat b_j\le\tfrac53 .
\end{equation}
The components along the coefficients $\hat b_j$ have sizes $\hat{\mathsf{b}}_j$. By the same two
results,
\begin{equation}\label{10.3:eq-source-variation-6}
\hat{\mathsf{b}}_j\le\operatorname{Lip}\hat b_j,\qquad \sum_j\hat{\mathsf{b}}_j\le\tfrac23 .
\end{equation}

We prove by induction on $n$ the following three statements simultaneously:
\begin{equation}\label{10.3:eq-source-variation-7}
\begin{aligned}
{}[\![\hat\delta^{\psi}]\!]_i&\le 3\,\mathfrak{g}^{\sharp}_i &&(i\le n),\\
\mathsf{r}_i&\le\mathsf{C}''\,\hat\varepsilon_i\,\mathfrak{g}^{\sharp}_i &&(i\le n),\\
\hat{\mathsf{e}}_i&\le c_e\,\mathfrak{g}^{\sharp}_{i-1} &&(i\le n+1).
\end{aligned}
\end{equation}
Here $\mathsf{r}_i$ are the sizes of the re-filed remnants and $\hat{\mathsf{e}}_i$ the $E$-norms of
the direction $\delta^{\psi}$. By its definition as a running maximum, the profile
$i\mapsto\mathfrak{g}^{\sharp}_i$ is nondecreasing, and $\mathfrak{g}_i\le\mathfrak{g}^{\sharp}_i$.
The step at level $n$ consists of three estimates, taken in the following order. We first bound
$\mathsf{r}_n$; this uses the first statement of \eqref{10.3:eq-source-variation-7} only for
$i<n$. The bound on $\mathsf{r}_n$ then enters the load at level $n$. Finally we bound the $E$-norm
at level $n+1$.

The remnant at level $n$. Apply parts (c) and (e) of the hatted one-lattice Lipschitz step.
\begin{itemize}
\item For the terms $i<n$, the loads are bounded by $3\mathfrak{g}^{\sharp}_i\le
3\mathfrak{g}^{\sharp}_n$, since the profile is nondecreasing.
\item The term $i=n$ is treated by the device used for the same term in the recursion for the
vector loads. It returns a multiple $\frac18$ of $\mathsf{r}_n$ to the left-hand side.
\end{itemize}
The result is
\begin{equation}\label{10.3:eq-source-variation-8}
\bigl(1-\tfrac18\bigr)\mathsf{r}_n\le\epsilon_n\cdot 3C_q^0\,\mathfrak{g}^{\sharp}_n
+\tfrac18\,\mathsf{C}''\hat\varepsilon_n\,\mathfrak{g}^{\sharp}_n .
\end{equation}
Since $\mathsf{C}''=8C_q^0$, we have $3C_q^0=\frac38\mathsf{C}''$, and dividing
\eqref{10.3:eq-source-variation-8} by $\frac78$ gives
\begin{equation*}
\mathsf{r}_n\le\tfrac17\,\mathsf{C}''\bigl(3\epsilon_n+\hat\varepsilon_n\bigr)\mathfrak{g}^{\sharp}_n .
\end{equation*}
By the last ceiling in \eqref{10.3:eq-source-variation-2}, $\epsilon_n\le 2\hat\varepsilon_n$, so
$3\epsilon_n+\hat\varepsilon_n\le 7\hat\varepsilon_n$, and the last display gives
$\mathsf{r}_n\le\mathsf{C}''\hat\varepsilon_n\mathfrak{g}^{\sharp}_n$. This is the second statement
of \eqref{10.3:eq-source-variation-7} at $i=n$.

The load at level $n$. The load of $\delta^{\psi}$ at level $n$ collects four contributions:
\begin{itemize}
\item the $E$-norms, through the scale kernels;
\item the remnants, with the factor $\mathsf{a}^*_n$;
\item the remnant of level $n$ itself;
\item the components of $\delta^{\psi}$ along the source itself at level $n$ (the running shift
and the coefficients $\hat b_j$), bounded in \eqref{10.3:eq-source-variation-5} and
\eqref{10.3:eq-source-variation-6}.
\end{itemize}
Insert the bounds \eqref{10.3:eq-source-variation-7} for $\hat{\mathsf{e}}_i$ and $\mathsf{r}_i$,
and use that $\mathfrak{g}^{\sharp}$ is nondecreasing. This gives
\begin{equation*}
[\![\hat\delta^{\psi}]\!]_n\le S_{\varsigma}c_e\,\mathfrak{g}^{\sharp}_n
+5\,\mathsf{a}^*_n\mathsf{C}''\hat\varepsilon_n\,\mathfrak{g}^{\sharp}_n
+\tfrac52\,\mathsf{r}_n+\mathfrak{g}_n\bigl[\tfrac53+\tfrac23\bigr].
\end{equation*}
The ceilings \eqref{10.3:eq-source-variation-2} bound the first two terms by
$\frac14\mathfrak{g}^{\sharp}_n$ and $\frac18\mathfrak{g}^{\sharp}_n$. By the second statement of
\eqref{10.3:eq-source-variation-7} at $i=n$ and the third ceiling,
\begin{equation*}
\tfrac52\,\mathsf{r}_n\le\tfrac52\,\mathsf{C}''\hat\varepsilon_n\,\mathfrak{g}^{\sharp}_n
\le\tfrac18\,\mathfrak{g}^{\sharp}_n .
\end{equation*}
Finally $\mathfrak{g}_n\le\mathfrak{g}^{\sharp}_n$ and $\frac53+\frac23=\frac73$. Hence
\begin{equation*}
[\![\hat\delta^{\psi}]\!]_n\le\bigl(\tfrac14+\tfrac18+\tfrac18+\tfrac73\bigr)\mathfrak{g}^{\sharp}_n
=\tfrac{17}6\,\mathfrak{g}^{\sharp}_n\le 3\,\mathfrak{g}^{\sharp}_n .
\end{equation*}

The $E$-norm at level $n+1$. Part (a) of the hatted one-lattice Lipschitz step gives
\begin{equation*}
\hat{\mathsf{e}}_{n+1}\le 3\,\hat\varkappa_n\,\mathfrak{g}^{\sharp}_n
+10\,C_{\mathsf{RF}}\,\mathsf{C}''\hat\varepsilon_n\,\mathfrak{g}^{\sharp}_n .
\end{equation*}
Since $\hat\varkappa_n\le\sup_k\hat\varkappa_k$ and $\hat\varepsilon_n\le\sup_i\hat\varepsilon_i$,
the right-hand side is at most $c_e\,\mathfrak{g}^{\sharp}_n$, by the definition of $c_e$ in
\eqref{10.3:eq-source-variation-1}. This is the third statement of
\eqref{10.3:eq-source-variation-7} at $i=n+1$, and the induction step is complete.

It remains to show that $\mathfrak{g}^{\sharp}_n\le 2\mathfrak{g}_n$. The correlation theorem's
comparison of the couplings at different levels, as recorded among the conditions on the vector
constants, gives $\mathfrak{g}_i\le 2\mathfrak{g}_n$ for every $i\le n$. Taking the maximum over
$i\le n$ gives $\mathfrak{g}^{\sharp}_n\le 2\mathfrak{g}_n$, which completes
\eqref{10.3:eq-source-variation-3}. Where this lemma is applied, the bound just proved is the bound
on the source vertex, whose load is $\mathfrak{g}_j$ times the size of the variation of the source.
\end{proof}

\begin{lemma}[The mixed-derivative bound on the source gradient]\label{10.3:lem-mixed-derivative}
For $0\le j\le K'$ let $\mathsf{m}_j$ be the least number such that
\begin{equation*}
\bigl|D_w\hat{\mathsf{O}}_j(X,\cdot,y;w,\mathsf{h})[\psi]\bigr|
\le \mathsf{m}_j\,\sup|\mathsf{h}|\,\|\psi\|_{j,X}\,e^{-2\kappa d_j(X)}
\end{equation*}
for all $w$ with $|w|\le r/4$ and $\operatorname{Lip}_j w\le r\lambda_*/4$, that is, for $w$ in the
source class of anchor $j$. Taking $\psi\equiv 1$, this gives
\begin{equation*}
\sup_{|\hat w|\le r/4}\bigl\|\partial_{\hat w}\hat{\mathfrak{V}}^{(j)}_{\hat w}\bigr\|^{\mathrm{fr}}
\le \mathsf{m}_j .
\end{equation*}
Put $\mathsf{g}^{\sharp}_j:=\max_{i<j}\theta_i^2\mathfrak{g}^{\sharp}_i$. Assume the following.
\begin{itemize}
\item[(a)] The production bounds for the vector part under the split rule,
Definition~\ref{10.3:def-vector-production-bounds}, whose proof is outstanding, in particular the rows
\eqref{10.3:eq-vector-production-bounds-d} and \eqref{10.3:eq-vector-production-bounds-f}, with their
constants $K_2$ and $C_L^V$, the re-filing constant $C_{\mathsf{RF}}$ and the band $\mathfrak{B}(k)$.
\item[(b)] The geometric decay of the frozen vector run, Lemma~\ref{10.3:lem-vector-run-decay},
with its constants $\mathsf{A}_V$, $\bar\varkappa$, $\mathsf{C}$, $\tilde\varepsilon_i$.
\item[(c)] The load of a variation of the source, Lemma~\ref{10.3:lem-source-variation}.
\item[(d)] The first of the conditions collected in \eqref{10.3:eq-vector-run-decay-a}.
\item[(e)] The ceiling
\begin{equation}\label{10.3:eq-mixed-derivative-a}
\frac{4}{r}\,\mathsf{C}\,\tilde\varepsilon_i\cdot 10\,C_{\mathsf{RF}}\,
\vartheta_V^{-\mathsf{a}^*_i}
\le \tfrac14 K_2\,(C_L^V)^2\,\theta_i^2\,\mathfrak{g}_i
\qquad\text{for all } i .
\end{equation}
This ceiling compares a power of $x_i=g_i^{-2}$ with $x_i^{-1}$, up to logarithms.
\end{itemize}
Then, with $C_{\mathsf{m}}:=2K_2(C_L^V)^2$,
\begin{equation}\label{10.3:eq-mixed-derivative-1}
\mathsf{m}_j\le C_{\mathsf{m}}\,\mathsf{g}^{\sharp}_j\,\mathsf{A}_V\,\vartheta_V^{\,j-1}
\qquad(1\le j\le K'),
\end{equation}
and $\mathsf{m}_0=0$.
\end{lemma}

\begin{proof}
We argue by induction on $j$. Fix $0\le k<K'$ and prove \eqref{10.3:eq-mixed-derivative-1} for
$j=k+1$, assuming it for every level $n$ with $1\le n\le k$. For $k=0$ nothing is assumed.

Take the anchor $K''=k+1$. Let $w$ satisfy $|w|\le r/4$ and $\operatorname{Lip}_{k+1}w\le r\lambda_*/4$.
Let $\psi$ satisfy $\|\psi\|_{k+1,X}\le 1$, extended off $X^+$ with the same norm. Then $w$ lies in the
$\tfrac12$-interior of $\mathcal{W}^{[k+1]}_{k+1}$, and $\psi$ has global norm at most $1$ at level
$k+1$. These are the hypotheses of Lemma~\ref{10.3:lem-source-variation} at anchor $k+1$.

The output of the step is
\begin{equation*}
\mathsf{O}_{k+1}(w)=\text{the $E$-part of }
\mathsf{T}_k^{\hat{\mathfrak{s}}(w)}\bigl[\mathfrak{u}_k[w,\mathsf{h}]\bigr].
\end{equation*}
It depends on $w$ in two ways: through the added units $\mathfrak{u}_k[w,\mathsf{h}]$, and through the
hatted run $\hat{\mathfrak{s}}(w)$. We bound the two dependences separately.

\emph{The units.} For $1\le n\le k$, the derivative of the sharp units in the direction $\psi$ is the
family
\begin{equation*}
D_w\mathsf{x}^{\sharp}_n[\psi]
=\Bigl\{\partial_{\hat w}\hat{\mathfrak{V}}^{\sharp(n)}_{\hat w(y)}
\bigl(X,\cdot,y;\mathsf{h}(y)\psi(p_y)\bigr)\Bigr\}.
\end{equation*}
The map that produces sharp units from $\mathsf{V}$-typed families is linear. Hence these are the sharp
units of the $\mathsf{V}$-typed families $\partial_{\hat w}\hat{\mathfrak{V}}^{(n)}$, read with the
bounded weight $\mathsf{h}\psi$ in place of $\mathsf{h}$. Since $\sup|\mathsf{h}\psi|\le\sup|\mathsf{h}|$,
the definition of $\mathsf{m}_n$ (at $\psi\equiv1$) shows that their norm is at most $\mathsf{m}_n$.

The bare units do not read $w$.

For the re-filed remnants we use Lemma~\ref{10.3:lem-own-level-freezing}(i) at level $i$. For every
localisation region $X'$ of a remnant $\mathsf{R}^{V(i)}$ one has $\|\psi\|_{i,X'}\le\|\psi\|_{k+1}$, so
the lemma applies on the source class of the anchor. Its Cauchy estimate gives
\begin{equation*}
\bigl|D_w\mathsf{R}^{V(i)}[\psi]\bigr|\le\frac4r\,\mathsf{r}_i .
\end{equation*}

We insert these bounds into the row \eqref{10.3:eq-vector-production-bounds-d} of the production bounds
for the vector part under the split rule, whose proof is outstanding. The contribution of the units is
then at most
\begin{equation}\label{10.3:eq-mixed-derivative-2}
\bar\varkappa\sum_{1\le n\le k}\varsigma^V_{k+1-n}\,\mathsf{m}_n
+(\bar\varkappa+C_{\mathsf{RF}})\sum_{i\in\mathfrak{B}(k)}w_{k,i}\,\frac4r\,\mathsf{r}_i .
\end{equation}
Here $C_{\mathsf{RF}}$ and the band $\mathfrak{B}(k)$ are the constant and the band of the re-filing
lemmas.

\emph{The run.} The derivative in $w$ through the hatted run is the mixed derivative of the $E$-output
of the production $\hat{\mathfrak{s}}+\lambda'\delta'$ with $\lambda\mathfrak{u}_k$ added. Here
$\delta'=\delta^{\psi}=D_w\hat{\mathfrak{s}}(w)[\psi]$ is the tangent of the hatted run in the source.
We apply the row \eqref{10.3:eq-vector-production-bounds-f}, whose proof is outstanding, with this
$\delta'$. By Lemma~\ref{10.3:lem-source-variation}, the load of $\delta^{\psi}$ at level $k$ is at most
$3\mathfrak{g}^{\sharp}_k$. By Lemma~\ref{10.3:lem-vector-run-decay}, the vector load $T^{V,T}_k$
entering that row is at most $\mathsf{A}_V\vartheta_V^k$. The run contributes at most
\begin{equation}\label{10.3:eq-mixed-derivative-3}
\tfrac14K_2(C_L^V)^2\theta_k^2\,T^{V,T}_k\cdot3\mathfrak{g}^{\sharp}_k
\le\tfrac34K_2(C_L^V)^2\,\theta_k^2\mathfrak{g}^{\sharp}_k\,\mathsf{A}_V\vartheta_V^k .
\end{equation}

\emph{The remnant terms.} Lemma~\ref{10.3:lem-vector-run-decay} gives
$\mathsf{r}_i\le\mathsf{C}\tilde\varepsilon_i\mathsf{A}_V\vartheta_V^i$. Together with the ceiling
\eqref{10.3:eq-mixed-derivative-a}, this bounds the second sum in \eqref{10.3:eq-mixed-derivative-2} by
\begin{equation}\label{10.3:eq-mixed-derivative-4}
\tfrac14K_2(C_L^V)^2\,\theta_k^2\mathfrak{g}^{\sharp}_k\,\mathsf{A}_V\vartheta_V^k\cdot(1+\bar\varkappa).
\end{equation}

\emph{The first sum in \eqref{10.3:eq-mixed-derivative-2}.} We use the induction hypothesis
$\mathsf{m}_n\le C_{\mathsf{m}}\mathsf{g}^{\sharp}_n\mathsf{A}_V\vartheta_V^{n-1}$ for $1\le n\le k$. The
sequence $\mathsf{g}^{\sharp}$ is nondecreasing, since it is a maximum over an increasing range of
indices. Hence $\mathsf{g}^{\sharp}_n\le\mathsf{g}^{\sharp}_{k+1}$. Summing the kernel against
$\vartheta_V^{n-1}$ with the sum $S^V$, and using the first condition of
\eqref{10.3:eq-vector-run-decay-a}, we obtain
\begin{equation*}
\begin{aligned}
\bar\varkappa\sum_{1\le n\le k}\varsigma^V_{k+1-n}\,
C_{\mathsf{m}}\mathsf{g}^{\sharp}_n\mathsf{A}_V\vartheta_V^{n-1}
&\le\bar\varkappa\,S^V\,C_{\mathsf{m}}\,\mathsf{g}^{\sharp}_{k+1}\,\mathsf{A}_V\vartheta_V^k\\
&\le\tfrac18\,C_{\mathsf{m}}\,\mathsf{g}^{\sharp}_{k+1}\,\mathsf{A}_V\vartheta_V^k .
\end{aligned}
\end{equation*}

\emph{Summation.} By definition,
\begin{equation*}
\theta_k^2\mathfrak{g}^{\sharp}_k\le\max_{i<k+1}\theta_i^2\mathfrak{g}^{\sharp}_i
=\mathsf{g}^{\sharp}_{k+1}.
\end{equation*}
So all three contributions are multiples of
$K_2(C_L^V)^2\mathsf{g}^{\sharp}_{k+1}\mathsf{A}_V\vartheta_V^k$. Their coefficients are:
\begin{itemize}
\item $\tfrac18\cdot2$ for the units, since $C_{\mathsf{m}}=2K_2(C_L^V)^2$;
\item $\tfrac34$ for the run, by \eqref{10.3:eq-mixed-derivative-3};
\item $\tfrac38$ for the remnant terms, by \eqref{10.3:eq-mixed-derivative-4} and
$\bar\varkappa\le\tfrac12$; the latter holds because $\bar\varkappa S^V\le\tfrac18$, as used for the
first sum, and $S^V\ge2/\vartheta_V\ge2$, as in the proof of Lemma~\ref{10.3:lem-vector-run-decay}.
\end{itemize}
Adding,
\begin{equation*}
\bigl(\tfrac18\cdot2+\tfrac34+\tfrac38\bigr)K_2(C_L^V)^2=\tfrac{11}{8}K_2(C_L^V)^2
\le C_{\mathsf{m}} .
\end{equation*}
Therefore
$\mathsf{m}_{k+1}\le C_{\mathsf{m}}\mathsf{g}^{\sharp}_{k+1}\mathsf{A}_V\vartheta_V^{(k+1)-1}$. This closes
the induction.
\end{proof}

We also note the comparison $\mathsf{g}^{\sharp}_j\le4\theta_j^2\mathfrak{g}_j$. For $i<j$ one has
$\theta_i\le\theta_j$, since $\theta$ is nondecreasing in the coupling by the correlation theorem, and
$\mathfrak{g}^{\sharp}_i\le\mathfrak{g}^{\sharp}_j$, since $\mathfrak{g}^{\sharp}$ is a running maximum;
with $\mathfrak{g}^{\sharp}_j\le2\mathfrak{g}_j$ from Lemma~\ref{10.3:lem-source-variation} this gives
\begin{equation*}
\mathsf{g}^{\sharp}_j=\max_{i<j}\theta_i^2\mathfrak{g}^{\sharp}_i
\le\theta_j^2\mathfrak{g}^{\sharp}_j\le2\theta_j^2\mathfrak{g}_j\le4\theta_j^2\mathfrak{g}_j .
\end{equation*}

\begin{lemma}[Defect families with geometrically growing sources]\label{10.3:lem-defect-families}
For $1\le j\le K'$ let $\mathsf{D}_j=\mathsf{D}^a_j+\mathsf{D}^b_j$ be the local family injected at
level $j$. Here $\mathsf{D}^a_j$ is the defect of summation by parts at the level of birth
(Proposition~\ref{10.3:prop-summation-by-parts}), and $\mathsf{D}^b_j$ is the defect of the
decomposition of a local output at its pin (Proposition~\ref{10.3:prop-output-decomposition}).
Each $\mathsf{D}_j$ satisfies the conditions $(\mathrm{f}1)$ and $(\mathrm{f}2)$ on local
families, and vanishes at $\mathbf{1}$ after subtraction of its
vacuum part. Let $n^{\mathsf{D}}_j$ be its norm $\|\cdot\|^{\flat}$, taken at the rate
$(1-\tfrac14\delta)\kappa$.

The \emph{defect direction} $\mathsf{z}$ is the added collection whose slots are
\begin{equation*}
\mathsf{z}_{k+1}:=\mathsf{T}_k[\mathsf{z}_{\le k}]+\mathsf{D}_{k+1},
\end{equation*}
with $\mathsf{T}_k$ the tangent of the complete step $k\to k+1$
(Definition~\ref{10.3:def-vector-run}). The direction $\mathsf{z}$ is read only through the second
and third rows of the load system of the correlation theorem, with the weights
$\omega_{k,j}$ that those rows carry omitted; these weights are not the weights $w_{n,j}$ of
the flat load below. Let $Y_j$ and $Y^R_j$ be the two norms of $\mathsf{z}$ at
level $j$ that these rows control. Its \emph{flat load} is
\begin{equation*}
T^{\mathsf{D}}_n:=C_M''\sum_{j\le n}\rho_{n,j}^2\,Y_j+\sum_{j\le n}w_{n,j}\,Y^R_j .
\end{equation*}
Put $\mathsf{g}_V:=\vartheta_2L=\vartheta_VL^2$; then $\mathsf{g}_V>1$. Put further
\begin{equation*}
\Sigma_{\mathsf{g}}:=\sup_n\sum_{\nu\ge0}\bigl[C_M''\rho_{n,n-\nu}^2+w_{n,n-\nu}\bigr]
\mathsf{g}_V^{-\nu},
\qquad
C_q^0:=\sum_{i\ge0}(i+1)\,q^{(i-1)_+}.
\end{equation*}
The correlation theorem shows that $\Sigma_{\mathsf{g}}$ is bounded independently of the
cutoff $K'$.

Assume that the two inequalities of the last clause of the production bounds for the vector
part under the split rule (Definition~\ref{10.3:def-vector-production-bounds}; the proof of
these production bounds is outstanding) hold for $\mathsf{z}$:
\begin{align}
Y_{k+1}&\le\varkappa^V_k\,T^{\mathsf{D}}_k+n^{\mathsf{D}}_{k+1},
\label{10.3:eq-defect-families-1}\\
Y^R_k&\le\epsilon^V_k\sum_{n\le k}(k-n+1)\,q^{(k-n-1)_+}\,T^{\mathsf{D}}_n
+3\varkappa^V_k\sum_{k_1(k)<j<k}Y^R_j .
\label{10.3:eq-defect-families-2}
\end{align}
Suppose finally that the ceilings
\begin{equation}\label{10.3:eq-defect-families-a}
2\bar\varkappa\,\Sigma_{\mathsf{g}}\le1,\qquad
4C_q^0\Sigma_{\mathsf{g}}\,\hat\varepsilon^V_n+3\varkappa^V_n\bigl(n-k_1(n)\bigr)\le\tfrac12,
\qquad
\tfrac52\,\epsilon^V_n\le\tfrac18
\end{equation}
hold for every $1\le n\le K'$.

If, for some $\bar{\mathsf{n}}\ge0$, $n^{\mathsf{D}}_j\le\bar{\mathsf{n}}\,\mathsf{g}_V^j$ for
$1\le j\le K'$, then
\begin{equation*}
Y_j\le2\bar{\mathsf{n}}\,\mathsf{g}_V^j,\qquad
Y^R_j\le\bar{\mathsf{n}}\,\mathsf{g}_V^j,\qquad
T^{\mathsf{D}}_n\le2\Sigma_{\mathsf{g}}\,\bar{\mathsf{n}}\,\mathsf{g}_V^n .
\end{equation*}
\end{lemma}

\begin{proof}
The inequality $\mathsf{g}_V>1$ follows from Definition~\ref{10.3:def-exponents-rates}. That
definition gives $\vartheta_2>L^{-(2\hat\alpha-1)}$ and $\hat\alpha<1$, so
\begin{equation*}
\mathsf{g}_V=\vartheta_2L>L^{1-(2\hat\alpha-1)}=L^{2-2\hat\alpha}>1 .
\end{equation*}

We argue by induction on the level. The claim at level $k$ is that
$Y_j\le2\bar{\mathsf{n}}\,\mathsf{g}_V^j$ and $Y^R_j\le\bar{\mathsf{n}}\,\mathsf{g}_V^j$ for all
$j\le k$. All sums run over $j\ge1$, so $T^{\mathsf{D}}_0=0$.

\emph{The load.} Suppose the bounds on $Y_j$ and $Y^R_j$ hold for all $j\le n$. Insert them
into the definition of $T^{\mathsf{D}}_n$ and write $j=n-\nu$, so that $1\le j\le n$ becomes
$0\le\nu\le n-1$. Then
\begin{equation}\label{10.3:eq-defect-families-3}
\begin{split}
T^{\mathsf{D}}_n
&\le\sum_{j\le n}\bigl[2C_M''\rho_{n,j}^2+w_{n,j}\bigr]\bar{\mathsf{n}}\,\mathsf{g}_V^j\\
&=\bar{\mathsf{n}}\,\mathsf{g}_V^n\sum_{0\le\nu\le n-1}
\bigl[2C_M''\rho_{n,n-\nu}^2+w_{n,n-\nu}\bigr]\mathsf{g}_V^{-\nu}\\
&\le2\Sigma_{\mathsf{g}}\,\bar{\mathsf{n}}\,\mathsf{g}_V^n .
\end{split}
\end{equation}
The last step uses $2C_M''\rho^2+w\le2(C_M''\rho^2+w)$ termwise, the non-negativity of the
terms with $\nu\ge n$ that the sum defining $\Sigma_{\mathsf{g}}$ contains in addition, and the
definition of $\Sigma_{\mathsf{g}}$. In particular the third conclusion of the lemma follows once
the first two are proved for all levels up to $n$.

\emph{The norm $Y$.} Let $k\ge1$, and assume the claim at level $k-1$. By
\eqref{10.3:eq-defect-families-3} with $n=k-1$ (trivially if $k=1$), we have
$T^{\mathsf{D}}_{k-1}\le2\Sigma_{\mathsf{g}}\bar{\mathsf{n}}\,\mathsf{g}_V^{k-1}$. Apply
\eqref{10.3:eq-defect-families-1} at $k-1$ and use $\varkappa^V_{k-1}\le\bar\varkappa$, with
$\bar\varkappa$ the constant of the vector recursion (Lemma~\ref{10.3:lem-vector-run-decay}). The
hypothesis on $n^{\mathsf{D}}_k$ then gives
\begin{equation*}
Y_k\le\bar\varkappa\cdot2\Sigma_{\mathsf{g}}\,\bar{\mathsf{n}}\,\mathsf{g}_V^{k-1}
+\bar{\mathsf{n}}\,\mathsf{g}_V^{k}
\le\bar{\mathsf{n}}\,\mathsf{g}_V^{k-1}+\bar{\mathsf{n}}\,\mathsf{g}_V^{k}
\le2\bar{\mathsf{n}}\,\mathsf{g}_V^{k}.
\end{equation*}
The second inequality is the first ceiling in \eqref{10.3:eq-defect-families-a}, and the third
uses $\mathsf{g}_V>1$.

\emph{The norm $Y^R$.} First, a geometric sum. Write $i=k-n$ and use
$\mathsf{g}_V^{n}=\mathsf{g}_V^{k}\mathsf{g}_V^{-i}\le\mathsf{g}_V^k$, which holds because
$\mathsf{g}_V>1$. This gives
\begin{equation}\label{10.3:eq-defect-families-4}
\sum_{n\le k}(k-n+1)\,q^{(k-n-1)_+}\,\mathsf{g}_V^{n}
\le\mathsf{g}_V^k\sum_{i\ge0}(i+1)\,q^{(i-1)_+}=C_q^0\,\mathsf{g}_V^k .
\end{equation}

Next, inspect the right side of \eqref{10.3:eq-defect-families-2}. The only unknown on it is
$Y^R_k$, and it enters through the self-term $w_{k,k}Y^R_k$ of $T^{\mathsf{D}}_k$. It carries the
factor $(k-k+1)q^{0}=1$, so its coefficient is $\epsilon^V_k w_{k,k}$. Every other quantity there
is already bounded. The loads $T^{\mathsf{D}}_n$ with $n<k$ are bounded by
\eqref{10.3:eq-defect-families-3} and the claim at level $k-1$. The part of $T^{\mathsf{D}}_k$
other than $w_{k,k}Y^R_k$ involves only $Y_j$ with $j\le k$, bounded by the previous paragraph,
and $Y^R_j$ with $j<k$. The estimate \eqref{10.3:eq-defect-families-3} applies to it verbatim,
so it is at most $2\Sigma_{\mathsf{g}}\bar{\mathsf{n}}\,\mathsf{g}_V^k$.

The self-term is moved to the left as in the proof of
Lemma~\ref{10.3:lem-vector-run-decay}: with the bound $w_{k,k}\le\tfrac52$ used there, the third
ceiling in \eqref{10.3:eq-defect-families-a} bounds its coefficient $\epsilon^V_kw_{k,k}$ by
$\tfrac18$. Combining this with
\eqref{10.3:eq-defect-families-4} for the first sum, and with
$Y^R_j\le\bar{\mathsf{n}}\,\mathsf{g}_V^j\le\bar{\mathsf{n}}\,\mathsf{g}_V^k$ for the at most
$k-k_1(k)$ indices $k_1(k)<j<k$ of the second sum, gives
\begin{equation*}
\bigl(1-\tfrac18\bigr)Y^R_k
\le\Bigl[2C_q^0\Sigma_{\mathsf{g}}\,\epsilon^V_k+3\varkappa^V_k\bigl(k-k_1(k)\bigr)\Bigr]
\bar{\mathsf{n}}\,\mathsf{g}_V^k
\le\tfrac12\,\bar{\mathsf{n}}\,\mathsf{g}_V^k .
\end{equation*}
The last inequality comes from the second ceiling in \eqref{10.3:eq-defect-families-a}, read at
$n=k$. Hence
\begin{equation*}
Y^R_k\le\tfrac87\cdot\tfrac12\,\bar{\mathsf{n}}\,\mathsf{g}_V^k
=\tfrac47\,\bar{\mathsf{n}}\,\mathsf{g}_V^k\le\bar{\mathsf{n}}\,\mathsf{g}_V^k .
\end{equation*}

This proves the claim at level $k$ and completes the induction. The bound on
$T^{\mathsf{D}}_n$ then follows for every $n\le K'$ from \eqref{10.3:eq-defect-families-3}.
\end{proof}

\begin{remark}
An unrenormalised marginal family, carried forward by the tangent maps $\mathsf{T}_k$, can
compound by a factor $1+O(\varkappa^V_k)$ at each level. Over many levels this gives a cumulative
factor of the form
\begin{equation*}
\exp\Bigl(\sum_k\varkappa^V_k\Bigr).
\end{equation*}
In Lemma~\ref{10.3:lem-defect-families} this compounding does no harm. The sources grow by the
factor $\mathsf{g}_V>1$ per level, so at every level the family injected most recently dominates
the contributions of all earlier injections. The proof of Lemma~\ref{10.3:lem-defect-families}
uses no extraction, no normalisation and no bound on the tensor normalisations of the channels
$r\in\{\mathbf{3},\mathbf{6}\}$ (Lemma~\ref{10.4:lem-normalisation-decay}).
\end{remark}

\begin{remark}\label{10.3:prop-vector-package}
The statement of this proposition is not written out here; parts of its proof are printed below.
\end{remark}

\begin{proof}[Proof of Proposition~\ref{10.3:prop-vector-package}, continued]
\label{10.3:prf-package-sizes}
The hypotheses are those of Proposition~\ref{10.3:prop-vector-package}. The directions
$\delta^{\otimes}=\mathsf{y}$ and $\delta^{\mathrm{rem}}=\mathfrak{u}\cup\mathsf{z}$ and the sources
$\mathsf{S}^{\partial\eta}_j$, $\mathsf{S}^{\eta\partial w}_j$, $1\le j\le K'$, are those introduced
above. We write $\mathsf{v}_0$ for $\mathsf{v}_0(c_0)$, $\mathsf{k}=\partial\eta$ and
$\mathsf{k}'=\eta\otimes\partial w$ for the two weights, and we use $\vartheta_V=\vartheta_2/L$, that is,
$L\vartheta_V=\vartheta_2$.

\emph{Sizes of the sources.} Fix $1\le j\le K'$. The source $\mathsf{S}^{\partial\eta}_j$ is the sum of
two terms:
\begin{itemize}
\item the bond term $\mathsf{S}^{a,\partial\eta}_j$ of the summation by parts at the level of birth
(Proposition~\ref{10.3:prop-summation-by-parts}), formed with the frozen family
$\hat{\mathfrak{V}}^{(j)}_{\hat w}$, whose fresh norm is at most $\mathsf{u}_j$;
\item the pin term $\mathsf{S}^{b,\partial\eta}_j$ of the decomposition of a local output at its pin
(Proposition~\ref{10.3:prop-output-decomposition}), formed with the output $\hat{\mathsf{O}}_j$,
whose size is $\mathsf{o}_j$.
\end{itemize}
The norms in both propositions are taken at the rate $\kappa_{\sharp}=(15/8)\kappa$. The first
proposition bounds the bond term by $4C_{\mathfrak{J}}L^j\mathsf{u}_j\sup|\mathsf{k}|$. The second
bounds the pin term by $L^j\mathsf{o}_jD_X\sup|\mathsf{k}|$. We write
$D_X=\lambda_*^{-1}(\lambda_*D_X)$; the factor $\lambda_*D_X$ is absorbed into the constant $C_B'$,
which leaves the factor $C_B'\lambda_*^{-1}$. Lemma~\ref{10.3:lem-vector-run-decay} (geometric decay
of the frozen vector run) gives $\mathsf{u}_j\le\mathsf{o}_j\le\mathsf{v}_0\vartheta_V^j$. Hence
\begin{align*}
\|\mathsf{S}^{\partial\eta}_j\|^{\flat}
&\le L^j\bigl(4C_{\mathfrak{J}}\mathsf{u}_j+C_B'\lambda_*^{-1}\mathsf{o}_j\bigr)\sup|\mathsf{k}|
\le\bigl(4C_{\mathfrak{J}}+C_B'\lambda_*^{-1}\bigr)\mathsf{v}_0\,(L\vartheta_V)^j\sup|\mathsf{k}|\\
&=\bigl(4C_{\mathfrak{J}}+C_B'\lambda_*^{-1}\bigr)\mathsf{v}_0\,\vartheta_2^{\,j}\sup|\mathsf{k}|.
\end{align*}

The source $\mathsf{S}^{\eta\partial w}_j$ is treated in the same way. Its bond term is bounded by
$4C_{\mathfrak{J}}L^j\mathsf{m}'\sup|\mathsf{k}'|$, where $\mathsf{m}'$ is the supremum over
$|\hat w|\le r/4$ of the fresh norm of $\partial_{\hat w}\hat{\mathfrak{V}}^{(j)}_{\hat w}$. By the
definition of $\mathsf{m}_j$ (the choice $\psi\equiv1$ in the defining inequality), this supremum is
at most $\mathsf{m}_j$. Its pin term is bounded by
\begin{equation*}
4L^j\mathsf{m}_j\lambda_*^{-1}(1+\lambda_*D_X)\sup|\mathsf{k}'| ,
\end{equation*}
and the factor $1+\lambda_*D_X$ is again absorbed into $C_B'$. This gives the first inequality below.

For the second inequality we use the mixed-derivative bound on the source gradient
(Lemma~\ref{10.3:lem-mixed-derivative}), which states
$\mathsf{m}_j\le C_{\mathsf{m}}\mathsf{g}^{\sharp}_j\mathsf{A}_V\vartheta_V^{\,j-1}$, where
$\mathsf{A}_V=4S^V\mathsf{v}_0$. It follows that
\begin{equation*}
L^j\mathsf{m}_j\le C_{\mathsf{m}}\mathsf{g}^{\sharp}_j\cdot4S^V\vartheta_V^{-1}\mathsf{v}_0\,
(L\vartheta_V)^j=C_{\mathsf{m}}\mathsf{g}^{\sharp}_j\cdot4S^V\vartheta_V^{-1}\mathsf{v}_0\,\vartheta_2^{\,j}.
\end{equation*}
In this bound the factor $\mathsf{g}^{\sharp}_j$ is bounded by $4\theta_j^2\mathfrak{g}_j$. Altogether,
\begin{align*}
\|\mathsf{S}^{\eta\partial w}_j\|^{\flat}
&\le L^j\bigl(4C_{\mathfrak{J}}+4C_B'\lambda_*^{-1}\bigr)\mathsf{m}_j\sup|\mathsf{k}'|\\
&\le\bigl(4C_{\mathfrak{J}}+4C_B'\lambda_*^{-1}\bigr)C_{\mathsf{m}}\cdot4\theta_j^2\mathfrak{g}_j
\cdot4S^V\vartheta_V^{-1}\,\mathsf{v}_0\,\vartheta_2^{\,j}\sup|\mathsf{k}'|.
\end{align*}

The remaining properties of the sources asserted in clause~(i) are their type, their locality, their
analyticity and their universality. The first three are those established with
Propositions~\ref{10.3:prop-summation-by-parts} and~\ref{10.3:prop-output-decomposition}, together with
the display~\eqref{10.3:eq-vector-production-bounds-c} of the production bounds for the vector
part under the split rule. That display is among the hypotheses of
Proposition~\ref{10.3:prop-vector-package}, and its proof is outstanding. Universality holds because
the frozen families are built from the whole-lattice objects $\hat{\mathfrak{V}}^{(j)}_{\hat w}$ and
$\hat{\mathsf{O}}_j(\cdot\,;\hat w,\cdot)$, and these read the scaffold only.

\emph{Sizes of the loads of the remainder direction.} The remainder direction
$\delta^{\mathrm{rem}}=\mathfrak{u}\cup\mathsf{z}$ has two loads:
\begin{itemize}
\item $\sup|\mathsf{h}|\,T^V_n$, for the vector part $\mathfrak{u}$;
\item $T^{\mathsf{D}}_n$, for the defect direction $\mathsf{z}$.
\end{itemize}
Since $\mathsf{h}=\eta/a$ and $a=L^{-K'}$, we have
$\sup|\mathsf{h}|\le\|\eta\|_{\langle2\rangle}L^{K'}$. Lemma~\ref{10.3:lem-vector-run-decay} gives
$T^V_n\le\mathsf{A}_V\vartheta_V^n$ for $1\le n\le K'$. Therefore
\begin{equation}\label{10.3:eq-package-sizes-1}
\sup|\mathsf{h}|\,T^V_n\le\|\eta\|_{\langle2\rangle}\mathsf{A}_VL^{K'}\vartheta_V^{\,n}\qquad(n\ge1).
\end{equation}

The defect families are $\mathsf{D}_j=\mathsf{D}^a_j+\mathsf{D}^b_j$. Their norms $n^{\mathsf{D}}_j$ are
bounded by Propositions~\ref{10.3:prop-summation-by-parts} and~\ref{10.3:prop-output-decomposition},
with the sizes entering those bounds controlled by Lemmas~\ref{10.3:lem-vector-run-decay}
and~\ref{10.3:lem-mixed-derivative}. Here only a bound of $\mathsf{m}_j$ by a constant multiple of
$\mathsf{o}_j$ is needed. The result, with a constant $C_{\mathsf{D}}$, is
\begin{equation*}
n^{\mathsf{D}}_j\le C_{\mathsf{D}}\mathsf{v}_0\vartheta_V^{\,j}\|\eta\|_{\langle2\rangle}\,\frac{a_j^2}{a}
=C_{\mathsf{D}}\mathsf{v}_0\|\eta\|_{\langle2\rangle}L^{-K'}\mathsf{g}_V^{\,j}.
\end{equation*}
The equality holds for two reasons. First, $a_j^2/a=L^{-2(K'-j)}L^{K'}=L^{-K'}L^{2j}$. Second,
$\vartheta_V^jL^{2j}=\mathsf{g}_V^j$, because $\mathsf{g}_V=\vartheta_2L=\vartheta_VL^2$.

The hypothesis of Lemma~\ref{10.3:lem-defect-families} (defect families with geometrically growing
sources) therefore holds with
\begin{equation*}
\bar{\mathsf{n}}:=C_{\mathsf{D}}\mathsf{v}_0\|\eta\|_{\langle2\rangle}L^{-K'} .
\end{equation*}
That lemma gives $Y_j\le2\bar{\mathsf{n}}\mathsf{g}_V^j$ and $Y^R_j\le\bar{\mathsf{n}}\mathsf{g}_V^j$,
and moreover
\begin{equation}\label{10.3:eq-package-sizes-2}
T^{\mathsf{D}}_n\le2\Sigma_{\mathsf{g}}\,\bar{\mathsf{n}}\,\mathsf{g}_V^{\,n}.
\end{equation}

At $n=K'$ the relevant quantities coincide:
\begin{equation*}
L^{K'}\vartheta_V^{K'}=\vartheta_2^{K'}=L^{-K'}\mathsf{g}_V^{K'} .
\end{equation*}
By~\eqref{10.3:eq-package-sizes-1} with $\mathsf{A}_V=4S^V\mathsf{v}_0$, and
by~\eqref{10.3:eq-package-sizes-2}, both loads at $n=K'$ are at most
$C\|\eta\|_{\langle2\rangle}\mathsf{v}_0\vartheta_2^{K'}$, where
$C:=\max\{4S^V,2\Sigma_{\mathsf{g}}C_{\mathsf{D}}\}$.

\emph{Read-out.} We apply the read-out bound~\eqref{10.3:eq-vector-production-bounds-g} to the
remainder direction. This bound belongs to the production bounds for the vector part under the split
rule; it is among the hypotheses of Proposition~\ref{10.3:prop-vector-package}, and its proof is
outstanding.

\emph{Part $(\alpha)$.} Part~$(\alpha)$ is read at the loads just estimated. We split the
weighted sum of the loads into three contributions.

The vector loads with $n\ge1$ are treated first. By~\eqref{10.3:eq-package-sizes-1}, and with
$i=K'-n$, their contribution is
\begin{equation*}
\sum_{1\le n\le K'}(K'-n+1)q^{(K'-n-1)_+}\sup|\mathsf{h}|\,T^V_n
\le4S^V\|\eta\|_{\langle2\rangle}\mathsf{v}_0\vartheta_2^{K'}
\sum_{0\le i\le K'-1}(i+1)q^{(i-1)_+}\vartheta_V^{-i}.
\end{equation*}
The sum over $i$ is bounded by the constant $C_q^V$ of the vector recursion. This gives the
term $4S^VC_q^V$ in the bracket below.

The vector load with $n=0$ is estimated as in part~(b) of the proof of
Lemma~\ref{10.3:lem-vector-run-decay}. There it is shown that
$(K'+1)q^{K'-1}T^V_0\le\mathsf{s}_{K'}\vartheta_V^{K'}(2+\mathsf{t}_0)\mathsf{v}_0$. Multiplying by
$L^{K'}$ and using $L^{K'}\vartheta_V^{K'}=\vartheta_2^{K'}$, we obtain
\begin{equation*}
L^{K'}(K'+1)q^{K'-1}T^V_0\le\vartheta_2^{K'}\mathsf{s}_{K'}(2+\mathsf{t}_0)\mathsf{v}_0 .
\end{equation*}
Since $\sup|\mathsf{h}|\le\|\eta\|_{\langle2\rangle}L^{K'}$, it follows that
\begin{equation*}
(K'+1)q^{K'-1}\sup|\mathsf{h}|\,T^V_0
\le\|\eta\|_{\langle2\rangle}\mathsf{v}_0\vartheta_2^{K'}\mathsf{s}_{K'}(2+\mathsf{t}_0),
\end{equation*}
which gives the term $\mathsf{s}_{K'}(2+\mathsf{t}_0)$ in the bracket below.

The defect loads come last. By~\eqref{10.3:eq-package-sizes-2}, with $i=K'-n$, and since
$\mathsf{g}_V>1$ and $C_q^0=\sum_{i\ge0}(i+1)q^{(i-1)_+}$, we have
\begin{align*}
\sum_{n\le K'}(K'-n+1)q^{(K'-n-1)_+}T^{\mathsf{D}}_n
&\le2\Sigma_{\mathsf{g}}\bar{\mathsf{n}}\,\mathsf{g}_V^{K'}\sum_{i\ge0}(i+1)q^{(i-1)_+}\mathsf{g}_V^{-i}
\le2\Sigma_{\mathsf{g}}C_q^0\,\bar{\mathsf{n}}\,\mathsf{g}_V^{K'}\\
&=2C_{\mathsf{D}}\Sigma_{\mathsf{g}}C_q^0\,\|\eta\|_{\langle2\rangle}\mathsf{v}_0\vartheta_2^{K'} .
\end{align*}

Adding the three contributions gives
\begin{equation*}
\begin{split}
&\sum_{n\le K'}(K'-n+1)q^{(K'-n-1)_+}\bigl[\sup|\mathsf{h}|\,T^V_n+T^{\mathsf{D}}_n\bigr]\\
&\qquad\le\|\eta\|_{\langle2\rangle}\mathsf{v}_0\vartheta_2^{K'}\bigl[4S^VC_q^V
+\mathsf{s}_{K'}(2+\mathsf{t}_0)+2C_{\mathsf{D}}\Sigma_{\mathsf{g}}C_q^0\bigr].
\end{split}
\end{equation*}

The bracket is bounded uniformly in $K'$. Indeed,
$\mathsf{t}_0\mathsf{s}_{K'}\le C_{\mathsf{t}}\mathsf{s}_*(1+x_{K'}^{1/2})$. By the correlation
theorem, moreover, $x_{K'}\le X+b_++3\lambda m+2c_2$, where here $X=g^{-2}$ and $b_+$, $\lambda$
are the constants of that bound in the correlation theorem. The bracket is therefore at most a
number depending on $g$ and $m$, and not on $K'$.

The $n'$-fold derivative in $z$ at $z=0$ is now estimated by Cauchy's inequality on the polydisc with
polyradii
\begin{equation*}
\frac{(1-\upsilon)r}{4(2\mathsf{Q}+2)n'\|f_l\|_{\langle2\rangle}},\qquad l=1,\dots,n' .
\end{equation*}
This yields a constant $C_\alpha(n',\mathsf{Q})$, which is a polynomial of degree $n'+1$ in
$\mathsf{Q}$.

\emph{Part $(\beta)$.} Part~$(\beta)$ concerns $\delta e_j(X;z)$, the vacuum part of the slots
of level $j$. We treat first the slots coming from the sharp units, then the remaining slots.

For $\mathsf{x}^{\sharp}_j$ the vacuum part is
\begin{equation*}
\sum_{y\in X}\hat{\mathfrak{V}}^{\sharp(j)}_{w(p_y)}\bigl(X,\mathbf{1},y;\mathsf{h}(y)\bigr),
\end{equation*}
and each term is a function of $w$ at the single point $p_y$. Since
$w(p_y)=\sum_lz_lf_l(p_y)$, the derivative of such a term in $z_l$ and $z_{l'}$ carries the factor
$f_l(p_y)f_{l'}(p_y)$, which vanishes when $\operatorname{supp}f_l$ and $\operatorname{supp}f_{l'}$ are
disjoint. Hence the mixed derivative of this vacuum part in two sources of disjoint supports
vanishes identically. The same holds for the derivative in one source whose support misses
$\operatorname{supp}\mathsf{h}$.

For the slots coming from $\mathsf{R}^V$, from $\mathsf{rim}$ and from $\mathsf{z}$, the mixed
derivative vanishes unless $X^+$ meets every $\operatorname{supp}f_l$, and also
$\operatorname{supp}\eta$ if $n'=1$. Given the separation $\mathsf{d}$ of the supports, this requires
\begin{equation*}
4Ma_j\bigl(d_j(X)+3\bigr)\ge\mathsf{d}.
\end{equation*}

We now use the support count in clause~$(\beta)$ of the correlation theorem. We also use four
bounds established above:
\begin{itemize}
\item $\sup|\mathsf{h}|\le\|\eta\|_{\langle2\rangle}L^{K'}$;
\item the bounds on $\mathsf{u}_j$ and $\mathsf{r}_j$ of Lemma~\ref{10.3:lem-vector-run-decay};
\item the bounds on $Y_j$ and $Y^R_j$ of Lemma~\ref{10.3:lem-defect-families} with the value of
$\bar{\mathsf{n}}$ above;
\item the relations $Ma_j=ML^{-i}$ and $L^{K'}\vartheta_V^{\,j}=\vartheta_2^{K'}\vartheta_V^{-i}$,
where $i=K'-j$.
\end{itemize}
Together these give, with a constant $C'$,
\begin{equation*}
\begin{split}
&\sum_j\bigl[\sup|\mathsf{h}|\bigl(C_{\mathfrak{J}}\mathsf{u}_j+\mathsf{r}_j\bigr)+Y_j+Y^R_j\bigr]
\bigl(1+(Ma_j)^{-4}\bigr)e^{-\frac14\kappa[\mathsf{d}/(4Ma_j)-3]_+}\\
&\qquad\le C'\|\eta\|_{\langle2\rangle}\mathsf{v}_0\vartheta_2^{K'}\sum_{i\ge0}\vartheta_V^{-i}
\bigl(1+L^{4i}M^{-4}\bigr)e^{-\frac14\kappa[\mathsf{d}L^i/(4M)-3]_+}\\
&\qquad=C(\mathsf{d})\|\eta\|_{\langle2\rangle}\mathsf{v}_0\vartheta_2^{K'} .
\end{split}
\end{equation*}
The last equality defines $C(\mathsf{d})$. The series over $i$ converges. Its general term is at most
a multiple of $(L^4/\vartheta_V)^i$ times $e^{-\frac14\kappa[\mathsf{d}L^i/(4M)-3]_+}$. This
exponential decays doubly exponentially in $i$ and therefore beats the geometric growth
$(L^4/\vartheta_V)^i$.

This completes the read-out, and with it the proof of Proposition~\ref{10.3:prop-vector-package}.
\end{proof}

\begin{remark}[Remarks on the vector package]\label{10.3:rem-package-remarks}
We make five remarks on the vector package, Proposition~\ref{10.3:prop-vector-package}.
Summation over repeated indices is understood throughout, and we write $\eta_j := \eta/a_j$.

\begin{enumerate}
\item \emph{The sources named in the input of the reduction theorem.}
The reduction theorem takes as input a statement on the vector directions. In that statement the
E-sources of the typed direction $\delta^{\otimes}$ are described in words, and
Proposition~\ref{10.3:prop-vector-package} is that statement with $\vartheta_2$ prescribed as in
\eqref{10.3:eq-exponents-rates-b}. The descriptions correspond to the sources of
Propositions~\ref{10.3:prop-summation-by-parts} and~\ref{10.3:prop-output-decomposition} as follows.
\begin{itemize}
\item The first description is the term obtained by summing the divergence identity of
Lemma~\ref{10.3:lem-divergence-identity} by parts against $\eta_j$. In the notation of that
description, in which $\Sigma_{\nu\rho}$ is the quantity against which the difference of the weight is
summed, the term is
\begin{equation*}
-\sum_y (\nabla_{\nu}\eta^{\rho}_j)\,\Sigma_{\nu\rho},
\end{equation*}
in which the difference falls on the weight. It is the bond source
$\mathsf{S}^{a,\partial\eta}_j$ of Proposition~\ref{10.3:prop-summation-by-parts}.
\item The second description is the summation by parts falling on $\hat w(y)$. It is
$\mathsf{S}^{a,\eta\partial w}_j$ of Proposition~\ref{10.3:prop-summation-by-parts}.
\item The third description is the first-order freezing term. It is
$\mathsf{S}^{b,\eta\partial w}_j$ of Proposition~\ref{10.3:prop-output-decomposition}.
\item The bound on the sources of weight $\eta\otimes\partial w$ carries a factor
$\theta_j^2\mathfrak{g}_j$. The input describes this factor as a second difference with one dial on
the vector unit and one on the source vertex. It is the bound~\eqref{10.3:eq-vector-production-bounds-f}
of the production bounds for the vector part under the split rule
(stated in Definition~\ref{10.3:def-vector-production-bounds}; its proof is outstanding), combined with
Lemma~\ref{10.3:lem-source-variation}. This combination is carried out in
Lemma~\ref{10.3:lem-mixed-derivative}.
\end{itemize}
One source is not named in the input. It is of the first kind, that is, it carries the weight
$\partial\eta$: it is $\mathsf{S}^{b,\partial\eta}_j$ of
Proposition~\ref{10.3:prop-output-decomposition}, the first-order dependence of the freshly produced
family on the gradient of the weight.

The input statement of the reduction theorem rests in addition on a lemma asserting that, at the shell
and at the sources, the double subtraction leaves a marginal insertion of bounded weight whose
coefficient is a first moment of the frozen kernel. That lemma is
Proposition~\ref{10.3:prop-summation-by-parts} taken together with the evaluation of the bond current of
Lemma~\ref{10.3:lem-divergence-identity} on constant cochains,
\begin{equation*}
\mathfrak{C}_{\rho}(b_\nu;f) = \mathsf{T}_{\nu\rho;AB}f_Af_B,
\end{equation*}
with the $W_4$-invariant constant tensor $\mathsf{T} = -2\hat\mu^0 + s$, which on constant curvature
gives
\begin{equation}\label{10.3:eq-package-remarks-1}
\mathsf{S}[F](\cdot\,;A) = -A_{\nu\rho}\,\mathsf{T}_{\nu\rho;AB}\,f_Af_B.
\end{equation}

\item \emph{Own-scale weights.} Measure the level-$j$ family per unit of $\eta_j$. It is then
$L^j\hat{\mathfrak{V}}^{(j)}$, of norm at most $\mathsf{v}_0(c_0)\vartheta_2^j$. This is the
normalisation of the input of the reduction theorem, in which the weight at level $j$ is $\eta_j$ and
\begin{equation*}
\nabla_j\eta_j = \partial\eta .
\end{equation*}

\item \emph{$\eta$ is arbitrary in $C^2$.}
Proposition~\ref{10.3:prop-vector-package} does not use that $\eta$ has the form $\varphi\xi$ with a
compactly supported $\varphi$ and a Killing field $\xi$. It therefore serves all of the following:
\begin{itemize}
\item the weights $\eta = \chi\xi$ and $\eta = f_i\xi$ of the reduction theorem, with $\chi$ and the
$f_i$ compactly supported $C^2$ functions;
\item the weight $\tilde\eta$ of a further quantity read by the reduction theorem, in which the
Killing-weighted bare vector unit $\mathfrak{K}^{\sharp}_{\tilde\eta}[1]$ is replaced, at $c_0 = 1$, by
$\delta^{\otimes} + \delta^{\mathrm{rem}}$;
\item the periodic weight of the torus form of the argument, a $C^2$ function on the torus, with
its norm $\|\cdot\|_{\langle 2\rangle}$ taken on the torus.
\end{itemize}
Two facts are used by the reduction theorem and not here: the $\Lambda^2$-part of $\partial\eta$
leaves no marginal term, and the strain of a Killing field vanishes.

\item \emph{Sharpness of \eqref{10.3:eq-exponents-rates-a}.} Suppose
$\hat\alpha \le \tfrac12$. Then $2\hat\alpha \le 1$, and since $L > 1$,
\begin{equation*}
L^{-2\hat\alpha} \ge L^{-1}.
\end{equation*}
Hence the interval $]L^{-2\hat\alpha}, L^{-1}[$ is empty, and no $\vartheta_V$ as required in
\eqref{10.3:eq-exponents-rates-b} exists. Under the one regularity hypothesis on the backgrounds, their
H\"older regularity with exponent $\hat\alpha$, the bulk contribution is irrelevant if and only if
$\hat\alpha > \tfrac12$. This equivalence is the criterion for the irrelevance of the bulk contribution
on H\"older backgrounds, and it is used here in exactly the form just stated.

\item \emph{No marginal quantity is bounded.} The hypotheses of
Proposition~\ref{10.3:prop-vector-package}, listed in its statement, contain no bound on a marginal
quantity. In particular, they include neither the marginal numbers $N_r$, nor the decay bound on the
marginal numbers, nor the one-loop anisotropy inequality.
\end{enumerate}
\end{remark}

\begin{corollary}[The vector-channel statement from the vector package]\label{10.3:cor-vector-channel}
Let $\hat\alpha$ and $\vartheta_2$ be fixed as in \eqref{10.3:eq-exponents-rates-a} and
\eqref{10.3:eq-exponents-rates-b}. Let the rate $\vartheta_1$ of the reduction theorem be
prescribed as in \eqref{10.3:eq-exponents-rates-b}:
\begin{equation*}
\vartheta_1\in\bigl]\max\{L^{-(2\hat\alpha-1)},\,L^{-1/2}\},\,1\bigr[
\quad\text{with}\quad \vartheta_2\le\vartheta_1 .
\end{equation*}
Assume the following.
\begin{enumerate}
\item The vector package, Proposition~\ref{10.3:prop-vector-package}.
\item The reduction theorem, Theorem~\ref{10.5:thm-reduction}, whose proof is incomplete at one
step, and its corollary, Corollary~\ref{10.5:cor-vector-channel}, together with the format
statements on which they rest, and the stress-tensor-insertion correlation theorem,
Theorem~\ref{10.4:thm-insertion-correlation}, in the channel $\mathbf{3}$, with its format and its
transport. These format
statements are: the typed format statement, together with the typed initial format that it
contains; the typed forms of the bound statement, of the single-level statement, of the re-filing
lemmas (Lemma~\ref{10.4:lem-typed-refiling}, whose proof is incomplete at one step), of the torus
bound in its $t$-form (Proposition~\ref{10.4:prf-tensor-torus-bound}), of the cluster activity
bound and of the bare-level statement; the typed format at the $\Lambda^2$-components; and the
one-lattice Lipschitz step at the identity $\mathbf{1}$.
\item The one-loop anisotropy inequality in the channel $\mathbf{3}$, the inequality of
Corollary~\ref{10.6:cor-anisotropy-inequality}, assumed here with an exponent $\kappa$ satisfying
$\kappa>\tfrac12$.
\item Clause (iii-b) of Theorem~0.
\item The volume statement in one of its two forms: the per-ball volume-free bound in its
$t$-form, or the first volume statement.
\end{enumerate}
Then the vector-channel statement, Proposition~\ref{10.2:prop-vector-channel}, holds, with the
cutoff function $\chi_{\mathcal{R}}$ at the radius $\mathcal{R}$ taken radial.
\end{corollary}

\begin{proof}
The argument has two parts. The first supplies the hypothesis of the reduction theorem on the
vector input. The second supplies the anisotropy hypothesis of its corollary. With both in hand,
the reduction theorem and its corollary, applied under the remaining hypotheses of the statement,
give the vector-channel statement.

\emph{The vector input.} The reduction theorem takes as a hypothesis a statement about the vector
input. Its content is the following. The direction whose component at level $0$ is
$c_0\mathfrak{K}^{\sharp}_{\eta}[1]$ has the same derivative of $\log Z(z,t)$ in $t$ at $t=0$ as
the sum $\delta^{\otimes}+\delta^{\mathrm{rem}}$ of a typed direction and a remainder direction.
Here $c_0$ is a complex coefficient and $\eta$ a weight. The reduction theorem uses this
hypothesis only at the two coefficients
\begin{equation*}
c_0\in\{g_0^{-2},\,1\},
\end{equation*}
and at the two weights
\begin{equation*}
\eta\in\{\chi\xi,\;f_i\xi\}.
\end{equation*}
Here $\xi$, $\chi$ and $f_i$ are the vector field, the cutoff function and the test functions
appearing in the reduction theorem. These choices are admitted by
Proposition~\ref{10.3:prop-vector-package}. That proposition holds for every
$\eta\in C^2(\mathbb{T}_m;V)$, where $V$ is the vector space in which the weights of the vector
package take values, every $c_0\in\mathbb{C}$, every $f_l\in C^2$ and every
$z\in\mathbb{D}_2$. Its conclusion, the decomposition
\begin{equation*}
\partial_t\big|_0\log Z(z,t)=\partial_t\log Z[\delta^{\otimes}](z)
+\partial_t\log Z[\delta^{\mathrm{rem}}](z),
\end{equation*}
is therefore the required hypothesis at each of these values.

Two further points match the vector package's bounds to the form the reduction theorem uses.

First, the bounds on the sources of $\delta^{\otimes}$ are stated with the weights $\partial\eta$
and $\eta\otimes\partial w$. They involve the two quantities $\sup|\partial\eta|$ and
$\sup|\eta\otimes\partial w|$. These are exactly the supremum terms $\|\cdot\|_\infty$ in the norms
\begin{equation*}
\|\cdot\|_{\sharp}=\max\{\|\cdot\|_\infty,\,40M\operatorname{Lip}(\cdot)\}
\end{equation*}
of the weights that enter the reduction theorem. The typed part therefore
enters that theorem with the norms it requires.

Second, clause (ii) of Proposition~\ref{10.3:prop-vector-package} bounds the remainder direction
$\delta^{\mathrm{rem}}$. Each mixed $z$-derivative at $0$ of
$\partial_t\log Z[\delta^{\mathrm{rem}}]$, at separated supports, is bounded by
\begin{equation*}
C(\varphi,\xi,\mathsf{d})\,\mathsf{v}_0(c_0)\,\varepsilon_K\prod_{l}\|f_l\|_{C^2}.
\end{equation*}
Here $\varphi$ is $\chi$ or $f_i$, and $\varepsilon_K$ is the sequence tending to zero that
appears in that clause. This bound is the form of the remainder $r_K$ displayed in the reduction
theorem, with the factor $\mathsf{Q}^{n+1}$ that $r_K$ carries there.

Hence the reduction theorem applies with the vector input furnished by
Proposition~\ref{10.3:prop-vector-package}.

\emph{The anisotropy hypothesis of the corollary.} The corollary of the reduction theorem requires
an anisotropy hypothesis with an exponent $c_1>\tfrac12$. This hypothesis follows from two
statements taken together. The first is the stress-tensor-insertion correlation theorem in the
channel $r=\mathbf{3}$, which is among the hypotheses. The second is decay of the tensor
normalisation in that channel with the exponent $\kappa$ of the one-loop anisotropy inequality.
That decay follows in turn from the one-loop anisotropy inequality in the channel $\mathbf{3}$ with
exponent $\kappa$, by Lemma~\ref{10.4:lem-normalisation-decay}. The two statements together give
the anisotropy hypothesis of the corollary by Corollary~\ref{10.4:cor-tensor-correlation-decay}.
Since the one-loop anisotropy inequality is assumed with $\kappa>\tfrac12$, the anisotropy
hypothesis of the corollary holds with some $c_1>\tfrac12$.

Both inputs are now available. The reduction theorem and its corollary yield the vector-channel
statement.
\end{proof}

\begin{remark}[Decay rates of sources, sharp families and defects]
\label{10.3:rem-decay-rates}
Let $\kappa$ be the parameter of that name among Ba{\l}aban's auxiliary parameters $\mathfrak{p}$,
let $\delta \in \mathopen{]}0,\delta_J(L)\mathclose{]}$
be the fixed number of the correlation theorem, $\delta_J(L)$ being the ceiling fixed there, and
let $d_j(X)$ be as in the glossary. By the display that defines the current term of the sourced
format, that term is measured in the
fresh norm $\|\cdot\|^{\mathrm{fr}}$ of Definition~\ref{10.3:def-typed-families}, with $\beta_s$
as there, whose weight is
\begin{equation*}
  e^{(1+4\beta_s)\kappa d_j(X)} = e^{2\kappa d_j(X)},
\end{equation*}
that is, at the rate $(1+4\beta_s)\kappa = 2\kappa$. The sources and the defect of
Proposition~\ref{10.3:prop-output-decomposition}, the sources of the vector package
(Proposition~\ref{10.3:prop-vector-package}), and the sharp families $F^{\sharp}$ and
$\mathsf{S}[F]$ of Lemma~\ref{10.3:lem-relocalisation} are delivered at the smaller rate
\begin{equation}
  \label{10.3:eq-decay-rates-a}
  \kappa_{\sharp} = \tfrac{15}{8}\,\kappa .
\end{equation}
The lemma of the proof of the correlation theorem that fixes the ceiling $\delta_J(L)$ lists the
places of that proof that read the improved exponent $2\kappa$: the first part of a lemma of that
proof at $\nu = 1$ and its second part at $j = k$, a remark of that proof at $j = k$, and a place
of that proof carrying a margin $\tfrac14$. The first three of these places are served by the
exponent $\kappa_{\sharp}$ of \eqref{10.3:eq-decay-rates-a}: they spend it against
the weight $\tfrac74\kappa$ of \cite[(1.66)]{B-CMP122b}, and they retain a margin
$\tfrac12\delta\kappa < \kappa/20$.
Alternatively, if $\delta$ is restricted, inside the window $\mathopen{]}0,\delta_J(L)\mathclose{]}$,
by the one-sided condition
\begin{equation}
  \label{10.3:eq-decay-rates-1}
  \delta \le \frac{1}{10}\Bigl(1 - \frac{17}{4L}\Bigr),
\end{equation}
then every statement above holds at the rate $2\kappa$; in particular each of the four places
listed above receives the exponent $2\kappa$ that it reads.
\end{remark}

\begin{proof}
Here, to spend $a\kappa$ of a rate $c\kappa$ on a factor or on a sum means to use the part
$e^{a\kappa d_j(X)}$ of the weight $e^{c\kappa d_j(X)}$ to control that factor or to make that sum
converge, so that the weight $e^{(c-a)\kappa d_j(X)}$ remains.
The rate \eqref{10.3:eq-decay-rates-a} is obtained from the rate $2\kappa$ of the fresh norm by two
constructions, each of which starts from $2\kappa$ and spends at most $\tfrac18\kappa$. The first is
the construction of the current $\mathfrak{J}[F]$ in the re-localisation lemma
(Lemma~\ref{10.3:lem-relocalisation}), which produces $F^{\sharp}$ and $\mathsf{S}[F]$. It spends
$\tfrac14\delta\kappa$ on the sum over the regions $X$, and $\tfrac{1}{10}\kappa$ on the lattice
sums, on the relation \eqref{10.3:eq-divergence-identity-a} satisfied by the constant kernels of
the divergence lemma
(Lemma~\ref{10.3:lem-divergence-identity}), on the transport and on the count of terms. Since
$\tfrac12\delta\kappa < \kappa/20$ by the choice of the ceiling $\delta_J(L)$ in the correlation
theorem, one has
$\tfrac14\delta\kappa < \kappa/40$, and this expenditure is
\begin{equation*}
  \tfrac14\delta\kappa + \tfrac{1}{10}\kappa < \tfrac{1}{40}\kappa + \tfrac{4}{40}\kappa
  = \tfrac18\kappa .
\end{equation*}
The second is the decomposition of a local output at its pin
(Proposition~\ref{10.3:prop-output-decomposition}), which produces the pin sources
$\mathsf{S}^{b,\partial\eta}_j$, $\mathsf{S}^{b,\eta\partial w}_j$ and the defect
$\mathsf{D}^b_j$. It spends $\tfrac18\kappa$ on the factor $D_X$. In both cases what remains is at
least
\begin{equation*}
  2\kappa - \tfrac18\kappa = \tfrac{15}{8}\kappa = \kappa_{\sharp}.
\end{equation*}

At the first three of the places listed in the remark, $\kappa_{\sharp}$ is spent against
$\tfrac74\kappa$. Since
\begin{equation*}
  \kappa_{\sharp} - \tfrac74\kappa = \tfrac{15}{8}\kappa - \tfrac{14}{8}\kappa = \tfrac18\kappa,
\end{equation*}
a spare $\tfrac18\kappa$ remains; since $\tfrac12\delta\kappa < \kappa/20 < \tfrac18\kappa$, it
covers the margin $\tfrac12\delta\kappa$ that these places retain. The spare $\tfrac18\kappa$ is
spent in this way at the two parts of the lemma and at the remark, as displayed in the proof of
the correlation theorem. The fourth place, the one carrying the margin $\tfrac14$, is not covered
by this comparison; it is covered by the alternative, under which it receives the rate $2\kappa$.

For the alternative, let $\delta$ satisfy \eqref{10.3:eq-decay-rates-1}. This condition is
one-sided and is imposed inside the window $\mathopen{]}0,\delta_J(L)\mathclose{]}$. Let
$\mathsf{r}(\delta)$ be the rate function of the rate statements for the outputs of the T-site of
Definition~\ref{10.3:def-typed-families}, which are among the antecedents of the vector package
(Proposition~\ref{10.3:prop-vector-package}), so that these outputs are delivered at the rate
$\mathsf{r}(\delta)\kappa$. By those rate statements,
\eqref{10.3:eq-decay-rates-1} gives
\begin{equation*}
  \mathsf{r}(\delta)\kappa \ge \tfrac{17}{8}\kappa .
\end{equation*}
The expenditures listed above then leave the rate $2\kappa$ for every statement above, and the
rate in part (b) of the rate lemma for the re-filing map $\mathsf{RF}[\cdot]$, one of the same
rate statements, is then $\tfrac{17}{8}\kappa$. Since the four places listed in the remark read the
improved exponent $2\kappa$, each of them then receives the exponent it reads.
\end{proof}

\begin{definition}[The inner-collar cores of the vector part]\label{10.3:not-inner-collar-cores}
Let $k$ be the current step of the run, and let $\mathbf{R}_k$ be the $\mathbf{R}$ operation of that
step. The vector part of the vector package at step $k$ is $\mathfrak{u}_k$. The level $k$ of the
collar cells of the correlation theorem and the production index $k$ of the vector package are
therefore the same letter. Let $v$ be a $\sharp$-core of the vector part, at the site of
$\mathbf{R}_k$ at which the correlation theorem reads the cores of the vector part.
By this we mean a member of the family $\mathsf{x}^{\sharp}_j$ that enters the smooth-core entry of
the vector package. Such a core is counterterm-free. Its data are the following.
\begin{itemize}
\item[(a)] The level $j$, with $1\le j\le k$.
\item[(b)] The anchor $y:=y_v\in C$. Here $C$ is a connected component of the large-field region
$Z$ treated by $\mathbf{R}_k$, and by the geometry of the correlation theorem $C$ is a box.
\item[(c)] The zone $n:=\pi(y)$ of $y$, and the numbers
\begin{equation}\label{10.3:eq-inner-collar-cores-1}
t:=L^{j-n},\qquad \nu:=1+n-j .
\end{equation}
\item[(d)] The box
\begin{equation*}
\square_0:=\square^{\sim\nu+1}(y).
\end{equation*}
This is the core's box of $2\nu+3$ level-$j$ blocks. We put $X_v:=\square_0$; this is the set
that the correlation theorem attaches to a core whose anchor it calls far, and for such a core it
requires $X_v\subset Z$.
\item[(e)] For $p\ge j$, the box $Q^{(p)}$, which is the union of $[\,\square_0\,]_p$ with the
$\pi_p$-cubes lying within two layers of it. Here, for a set $A$, $[\,A\,]_p$ is the set of
$\pi_p$-cubes whose interiors meet $A$.
\item[(f)] The weight $\bar{\mathsf{w}}_v=3$ of the core.
\end{itemize}
In this definition $\nu$ and $p$ denote the integers just introduced, not the tensor index and the
plaquette of the rest of the chapter.
Let $s(\cdot)$ be the integer function of the correlation theorem. By its definition, for a weight
$\bar{\mathsf{w}}$, $s(\bar{\mathsf{w}})$ is the least integer $s\ge 1$ for which the right member of
the inequality of the correlation theorem that defines $s(\cdot)$, which depends on the weight
$\bar{\mathsf{w}}$ and on a parameter of the core, written $a_v$ for the core $v$ above, is at most
$1$ at every value $a_v\ge 2+s$. With the standing choice of the first stage of the vector package,
we put
\begin{equation}\label{10.3:eq-inner-collar-cores-2}
s^*:=s(\bar{\mathsf{w}}_v)=s(3),\qquad a^*:=2+s^*,\qquad p^*:=n-a^* .
\end{equation}
Since $s(\cdot)$ takes values in the integers $s\ge 1$ by its definition, $s^*\ge 1$, and hence
\begin{equation*}
a^*=2+s^*\ge 3 .
\end{equation*}
The \emph{inner-collar cores of the vector part} are the cores $v$ as above that satisfy
\begin{equation}\label{10.3:eq-inner-collar-cores-3}
\nu\ge a^*+1,\qquad \square_0\subset C,\qquad Q_v:=Q^{(p^*)}(y)\not\subset C .
\end{equation}
By the definition of $\nu$ in \eqref{10.3:eq-inner-collar-cores-1}, the first condition of
\eqref{10.3:eq-inner-collar-cores-3} reads
\begin{equation*}
1+n-j\ge a^*+1 .
\end{equation*}
This is equivalent to $n-a^*\ge j$, that is, $p^*\ge j$. Consequently $Q^{(p^*)}$ is defined by
item (e), and so is $Q_v$.

In words, these are the cores of weight $\bar{\mathsf{w}}_v=3$ whose set $X_v=\square_0$ lies in
the treated component $C$ but whose box $Q^{(p)}$, at $p=n-2-s(\bar{\mathsf{w}}_v)$, is not
contained in $Z$. At the weight $\bar{\mathsf{w}}_v=3$ we have
\begin{equation*}
n-2-s(\bar{\mathsf{w}}_v)=n-2-s^*=n-a^*=p^* ,
\end{equation*}
so that the exponent $p$ in this description is $p^*$, and the box named in it is $Q_v$.

Units. A $\pi_p$-cube has side $M\,L^{-(k-p)}$ current sites. The level-$n$ unit, in which
$\operatorname{dist}_n$ and $d_n$ are measured, is the side of a $\pi_n$-cube.
\end{definition}

\begin{lemma}[Geometry of an inner-collar core]\label{10.3:lem-core-geometry}
Let $v$ be an inner-collar core of the vector part in the sense of
Definition~\ref{10.3:not-inner-collar-cores}. As there, $v$ has level $j$ and anchor $y\in C$, where $C$
is a component of the treated class $Z$. Its zone is $n=\pi(y)$, and we put $\nu=1+n-j$. Its box is
$\square_0=\square^{\sim\nu+1}(y)$, and we use the integers $a^*\ge 3$ and $p^*=n-a^*$ and the set
$Q_v=Q^{(p^*)}(y)$. The defining conditions of the class are
\[
\nu\ge a^*+1,\qquad \square_0\subset C,\qquad Q_v\not\subset C .
\]
One level-$k$ unit is the side $M$ of a $\pi_k$-cube, in current sites. The level-$k$ distance is
$\operatorname{dist}_k$ and the tree length is $d_k$, both measured in these units. Further, $\Omega_k$
is the regular region, $\check{R}_k$ the separation radius, $\partial C$ the boundary of $C$, and
$\Delta_y$ the $\pi_k$-cube containing $y$. Finally, $\mathfrak{r}_v$ denotes the support set of $v$.
Let $L\ge 8$. Then the following hold.
\begin{enumerate}
\item[(G1)] The set $Q_v$ contains $\square_0$. Its side is at most $2a^*+11$ $\pi_{p^*}$-cubes, and
hence at most the side of one $\pi_{k-1}$-cube:
\begin{equation}\label{10.3:eq-core-geometry-a}
\square_0\subseteq Q_v,\qquad
(2a^*+11)\,ML^{-a^*}\le L^{a^*-1}\,ML^{-a^*}=M/L\ \text{sites}.
\end{equation}
\item[(G2)] With $K_v:=[\square_0]_k$, and with $\#\mathcal{F}$ denoting the number of cubes in a
family $\mathcal{F}$ of $\pi_k$-cubes, one has
\begin{equation}\label{10.3:eq-core-geometry-b}
\begin{gathered}
Q_v\not\subset Z,\qquad [Q_v]_k\not\subset Z,\qquad \#[Q_v]_k\le 2^d,\\
K_v\subset C,\qquad \#K_v\le 2^d,\qquad d_k(K_v)\le 2^d-1\le 6^d,\qquad K_v\subset [Q_v]_k .
\end{gathered}
\end{equation}
\item[(G3)] The anchor lies in the collar, its zone is the current level, and $j<k$:
\begin{equation}\label{10.3:eq-core-geometry-c}
\begin{gathered}
y\in C\cap\Omega_k,\qquad n=\pi(y)=k,\qquad j<k,\qquad
\mathfrak{r}_v\subset\Omega_k\cup(C\cap\Omega_k),\\
\operatorname{dist}_k(y,\Omega_k^c)\ge 2\check{R}_k-1\ge 2\check{R}_k-5\ge 71 .
\end{gathered}
\end{equation}
Conversely, every $y\in C$ with $\pi(y)<k$ and $1+\pi(y)-j\ge a^*+1$ satisfies
$Q^{(p^*)}(y)\subset C$, where $p^*=\pi(y)-a^*$.
\item[(G4)] The set $Q_v$ meets $\Delta_y$ in a set with non-empty interior, and, with
$Q_5(\Delta_y):=\{\operatorname{dist}_k(\cdot,\Delta_y)\le 2\}$,
\begin{equation}\label{10.3:eq-core-geometry-d}
Q_v\subset\{\operatorname{dist}_k(\cdot,\Delta_y)\le 1\}\subset Q_5(\Delta_y),\qquad
Q_v\subset\Omega_{k-1}.
\end{equation}
\end{enumerate}
\end{lemma}

\begin{proof}
\emph{(G1).} By definition, $Q_v=Q^{(p^*)}(y)$ is $[\square_0]_{p^*}$ with two further layers of
$\pi_{p^*}$-cubes. It therefore contains $\square_0$. The number of cubes of this set is given by the
side count of the boxes $Q^{(p)}$, $j\le p\le n-3$, established in the proof of the correlation
theorem and taken here by statement. Namely, for $j\le p\le n-3$, writing
$a:=n-p$, so that $a\ge 3$, the box $Q^{(p)}$ has at most
\[
(2\nu+3)L^{-(p-j)}+6\le 2a+11\le L^{a-1}
\]
cubes of $\pi_p$ on a side. We use this count at $p=p^*$, that is, with $a=a^*\ge 3$. In
that case $p^*-j=\nu-1-a^*$, which is non-negative because $\nu\ge a^*+1$. At $\nu=a^*+1$ the first
term equals $2a^*+5$. Each unit increase of $\nu$ multiplies it by $(2\nu+5)/((2\nu+3)L)<1$. Hence the
first inequality holds with $a=a^*$.

For the second inequality, $2a+11\le L^{a-1}$ holds for all $a\ge 3$ and $L\ge 8$. At $a=3$ it reads
$17\le 64\le L^2$. Passing from $a$ to $a+1$ adds $2$ to the left member and multiplies the right
member by $L\ge 8$, so the gap increases.

By the units convention of Definition~\ref{10.3:not-inner-collar-cores}, a $\pi_{p^*}$-cube has side
$ML^{-(k-p^*)}\le ML^{-a^*}$ sites. The side of $Q_v$ is therefore at most
$L^{a^*-1}\cdot ML^{-a^*}=M/L$ sites, which is the side of one $\pi_{k-1}$-cube. This proves
\eqref{10.3:eq-core-geometry-a}.

\emph{(G2).} The geometry of the treated class established in the proof of the correlation theorem,
which we take here by statement, provides two facts about $Z$. First, $C$ is a union
of $\pi_k$-cubes. Second, distinct components of $Z$ are at least $\check{R}_k\ge 38$ units apart.

By (G1), $Q_v$ has side at most $M/L$ sites, which is less than one $\pi_k$-cube, and $Q_v$ contains
$y\in C$.
Suppose a $\pi_{p^*}$-cube of $Q_v$ is not contained in $C$; the hypothesis $Q_v\not\subset C$ says
exactly that such a cube exists. It lies in a $\pi_k$-cube $\Delta_0$, and $\Delta_0\not\subset C$.
Since $Q_v$ lies within one unit of $y\in C$, no other component of $Z$ is reached. Hence
$\Delta_0$ is a $\pi_k$-cube disjoint from $Z$. It follows that $Q_v\not\subset Z$. Moreover, $Q_v$
contains a
$\pi_{p^*}$-cube inside $\Delta_0$, so $\Delta_0\in[Q_v]_k$, and therefore $[Q_v]_k\not\subset Z$.
Finally, a box of side less than one unit meets at most two $\pi_k$-cubes in each direction, so
$[Q_v]_k$ consists of at most $2^d$ $\pi_k$-cubes.

We turn to $K_v$. The box $\square_0\subset C$ has side $2\nu+3$ $\pi_j$-cubes. Since $n\le k$, a
$\pi_j$-cube has side $L^{-(k-j)}\le L^{-(n-j)}$ level-$k$ units, so the side of $\square_0$ is, in
level-$k$ units, at most
\[
(2\nu+3)L^{-(n-j)}=(2\nu+3)L^{-(\nu-1)}\le \tfrac{11}{512}<1 .
\]
To see this, note that $\nu\ge a^*+1\ge 4$ and $L\ge 8$. At $\nu=4$ the value is at most $11/8^3$, and
each unit increase of $\nu$ multiplies it by $(2\nu+5)/((2\nu+3)L)<1$.

Since $C$ is a union of $\pi_k$-cubes, $K_v=[\square_0]_k\subset C$. By the same count per direction
as for $[Q_v]_k$, the $\pi_k$-cubes meeting the box $\square_0$ in interior form a box of at most two
cubes per direction; hence $K_v$ is a wall-connected family of at most $2^d$ $\pi_k$-cubes. A
wall-connected family of $N$ cubes
has a spanning tree of $N-1$ unit steps, so $d_k(K_v)\le 2^d-1\le 6^d$. Lastly, $\square_0\subset Q_v$
by (G1), and hence $K_v\subset[Q_v]_k$. This proves \eqref{10.3:eq-core-geometry-b}.

\emph{(G3).} By (G1), $Q_v$ has diameter at most $1/L\le 2$ units. It contains $y\in C$
and meets $C^c$. Hence $y$ lies inside $C$, within $M/L$ sites of $\partial C$.

The collar estimate established in the proof of the correlation theorem, which we take here by
statement, states that
$\operatorname{dist}_k(\partial C,\Omega_k^c)\ge
2\check{R}_k$ on both sides of $\partial C$. Consequently the collar $C\cap\Omega_k$ of $C$ is at least
$2\check{R}_k$ units thick and its points have zone $k$. Therefore $y\in C\cap\Omega_k$ and
$n=\pi(y)=k$.
Moreover, $y$ lies within $M/L$ sites, that is within $1/L\le 1$ unit, of $\partial C$, while every
point of $\Omega_k^c$ lies at least $2\check{R}_k$ units from $\partial C$. By the triangle inequality,
\[
\operatorname{dist}_k(y,\Omega_k^c)\ge 2\check{R}_k-1\ge 2\check{R}_k-5\ge 71,
\]
where the middle inequality is immediate and the last inequality uses $\check{R}_k\ge 38$.

In the same collar estimate, for a core whose box $\square_0$
meets $C$ and has
$K_v\not\subset Z$, the
conclusion
\[
n=\pi(y)=k,\qquad \mathfrak{r}_v\subset\Omega_k\cup(C\cap\Omega_k),\qquad
\operatorname{dist}_k(y,\Omega_k^c)\ge 2\check{R}_k-5\ge 71
\]
is drawn from the fact that $y$ lies within $2$ units of $\partial C$. Here that fact is obtained from
$Q_v\not\subset C$, as shown above. The same argument then yields the inclusion
$\mathfrak{r}_v\subset\Omega_k\cup(C\cap\Omega_k)$.

With $n=k$ we get $\nu=1+k-j\ge a^*+1\ge 4$, so $j\le k-3<k$. This proves
\eqref{10.3:eq-core-geometry-c}.

Conversely, let $y\in C$ with $\pi(y)<k$ and $1+\pi(y)-j\ge a^*+1$, so that $p^*=\pi(y)-a^*\ge j$
and $Q^{(p^*)}(y)$ is defined as above. The points of $C$ within $2\check{R}_k$ units of $\partial C$
lie in the collar $C\cap\Omega_k$, whose points have zone $k$; hence $y$ lies at least $2\check{R}_k$
units inside $C$. The count of (G1) uses only $1+\pi(y)-j\ge a^*+1$ and $L\ge 8$, not the condition
$Q_v\not\subset C$; so $Q^{(p^*)}(y)$, which contains $y$, has side at most $1/L$ unit, and therefore
$Q^{(p^*)}(y)\subset C$. Consequently no such $y\in C$ with $\pi(y)<k$ satisfies
$Q^{(p^*)}(y)\not\subset C$, so the defining condition $Q_v\not\subset C$ of the class is met only by
anchors of zone $k$, and the predicate defining the class reads only the $z$-free scaffold. This
proves the last assertion of (G3).

\emph{(G4).} By (G3), $\Delta_y$ is a $\pi_k$-cube of level $k$. The set $Q_v$ contains $y$ and the
$\pi_{p^*}$-cubes of $\square_0$ around $y$. Hence $Q_v$ meets $\Delta_y$ in a set with non-empty
interior. By (G1), $Q_v$ has side at most $1/L$ of a $\pi_k$-cube. Therefore
\[
Q_v\subset\{\operatorname{dist}_k(\cdot,\Delta_y)\le 1\}\subset Q_5(\Delta_y).
\]
The cube inclusion established in the proof of a lemma used for the correlation theorem, which we
take here by statement, is the following: for a cell
$\Delta'$ of level $n'$, the cube
$\{\operatorname{dist}_{n'}(\cdot,\Delta')\le 1\}$ lies in $\Omega_{n'-1}$. Applied with $n'=k$ and
$\Delta'=\Delta_y$, this gives $Q_v\subset\Omega_{k-1}$, which proves
\eqref{10.3:eq-core-geometry-d}.
\end{proof}

\begin{definition}[Hull, halo and start label of an inner-collar core]\label{10.3:def-hull-halo-label}
Let $v$ be an inner-collar core of the vector part in the sense of
Notation~\ref{10.3:not-inner-collar-cores}, with level $j$ ($1\le j\le k$, $k$ the current step),
anchor $y\in C$, component $C$ of the treated class $Z$, zone $n=\pi(y)$, $\nu=1+n-j$, box
$\square_0=\square^{\sim\nu+1}(y)$, for $p\ge j$ the box $Q^{(p)}$ obtained from $[\square_0]_p$
by adding two layers of $\pi_p$-cubes, weight $\bar{\mathsf{w}}_v=3$, and with the integer
$s^*:=s(3)$ of the weight function $s(\cdot)$, $a^*:=2+s^*$, $p^*:=n-a^*$ fixed there. We attach
to $v$ the following data.
\begin{enumerate}
\item[(a)] \emph{Core data.} We put $X_v:=\square_0\subset C$ and
$K_v:=[X_v]_k=[\square_0]_k$.
\item[(b)] \emph{Read level and read box.} Since $\bar{\mathsf{w}}_v=3$, we put
\begin{equation*}
p_v:=n-2-s(\bar{\mathsf{w}}_v)=p^*,\qquad Q_v:=Q^{(p_v)}=Q^{(p^*)}(y).
\end{equation*}
By the hypotheses of Notation~\ref{10.3:not-inner-collar-cores}, $Q_v\not\subset C$; since $Q_v$
is connected and contains $\square_0\subset C$, an inclusion $Q_v\subset Z$ would force
$Q_v\subset C$, hence $Q_v\not\subset Z$. This definition imposes on $v$ no inclusion
$Q^{(p)}\subset Z$.
The function $f_v$ of $v$, whose argument is a point $(\zeta,\varpi,\vec{\mathsf{s}})$ of the
source polydisc $\mathfrak{P}$ of the correlation theorem, is evaluated at its extension reading
$\operatorname{ext}_v$ on this box $Q_v$, by the extension rule for plain units of the
correlation theorem.
\item[(c)] \emph{Hull, halo, start label.} The \emph{hull}, the \emph{halo}, the
\emph{start label} and the \emph{reading set} of $v$ are
\begin{equation}\label{10.3:eq-hull-halo-label-a}
\begin{gathered}
\hat{X}_v:=Q_v,\qquad
\mathfrak{h}_v:=\{\Delta\in\pi_k:\ \operatorname{dist}_k(\Delta,Q_v)\le 2\},\\
\operatorname{lab}_0:=K_v\cup\mathfrak{h}_v,\qquad
\mathcal{R}_0:=\hat{X}_v\cup X_v .
\end{gathered}
\end{equation}
For a smooth core the plain-unit clause of the correlation theorem itself takes the hull to be
$Q_v$; for $v$ the hull $\hat{X}_v:=Q_v$ is set here by definition. The start label
$\operatorname{lab}_0$ of \eqref{10.3:eq-hull-halo-label-a} is set here in
place of the start label $\operatorname{lab}_0:=K_v$ of the piece clause of the correlation
theorem. The formula for $\mathcal{R}_0$ in \eqref{10.3:eq-hull-halo-label-a} is that of the
piece clause of the correlation theorem, applied with the hull just set. The \emph{support} of
$v$ is
\begin{equation*}
S_v:=\mathfrak{r}_v:=\bigcup_{\Delta_y}\operatorname{St}(\Delta_y),
\end{equation*}
the union of the stars $\operatorname{St}(\Delta_y)$ of the cells $\Delta_y$ read by $v$ under
the reading rule of the correlation theorem; it is the set read by the exposed current bound
and by the exposed regularity bound of the correlation theorem, and it is not the hull.
\item[(d)] \emph{Domain.} We put $X^{\oplus}_v:=Y(\operatorname{lab}_0)$, where $Y(A)$
is $A$ together with the components of $Z$ meeting $A$, and
$\mathfrak{E}^{\oplus}_v:=\mathbb{C}^{(N)}_k(X^{\oplus}_v)$, where $\mathbb{C}^{(N)}_k$ is the
functional of the correlation theorem that consumes an exposed domain at level $k$, here
evaluated on the domain $X^{\oplus}_v$, $N$ being the
number of zones ($k=h+N$, with $h$ the base level). The pieces and the value of $v$ are
assigned, by definition, to the domains that the extension of the correlation theorem's data to
the inner-collar cores of the vector part prescribes for them.
\end{enumerate}
Item (a) repeats the data $X_v:=\square_0$ that the plain-unit clause of the correlation theorem
assigns to a core. Items (b)--(d) are not consequences of the correlation theorem, apart from the
formula for $\mathcal{R}_0$ and, when $v$ is a smooth core, the hull $\hat{X}_v=Q_v$, both noted
in (c): they extend
its data for plain units, whose plain-unit clause includes the inclusion $Q^{(p)}\subset Z$ and
whose piece clause takes $\operatorname{lab}_0:=K_v$, and they are adopted here by definition.
These data have the following properties.
\begin{enumerate}
\item[(i)] $K_v\subset C$; $v$ is a plain interior unit of the first part of the sourced
expansion, a member of $\mathfrak{L}^{\mathrm{pt}}$, and a member neither of
$\mathfrak{L}^{c}$ nor of $\mathfrak{L}^{\mathrm{ex}}$; the data (a)--(d) do not change this
classification.
\item[(ii)] $\mathcal{R}_0=Q_v$.
\item[(iii)] As a set, $\operatorname{lab}_0=\mathfrak{h}_v$, and $\mathfrak{h}_v$ consists of
at most $6^d\le 7^d=c_{\mathfrak{h}}$ cubes.
\item[(iv)] $X^{\oplus}_v=C\cup\mathfrak{h}_v$.
\item[(v)] \emph{Expansion.} The core $v$ is expanded by the interior link rules of the
correlation theorem: no function $\zeta_{\square}$ of the partition of unity and no transport
parameter $\mathsf{t}$ enter; link $1$ is the single bracket of Ba{\l}aban's expansion, expressed
by the fundamental theorem of calculus in the interpolation parameter $\mathsf{s}^{(1)}$ of
link $1$ alone over the cells disjoint from the enlargement $\Lambda^{\sim}$ of the link's domain
$\Lambda$, and links $2$--$4$ are expanded over the cells disjoint from $\Lambda$. The pieces are
the tuples $\mathfrak{p}=(v;\mathsf{y}_1,\dots,\mathsf{y}_4)$, $\mathsf{y}_t$ a family of
link $t$, the tuple of four empty families being the value
$\operatorname{val}_v=f_v(\operatorname{chain}(\vec{0}))$, where $\operatorname{chain}(\vec{0})$
is the carrier pair at the site $(\mathrm{L})$ of $\mathbf{R}_k$
(Notation~\ref{10.3:not-inner-collar-cores}) at the point where all link
parameters vanish. In particular $v$ has no $\mathsf{t}$-piece.
\end{enumerate}
If instead $s^*$ were taken to be $s^{(9/8)}(3)=s(27/8)$, the core $v$ would carry a transport
piece $\mathfrak{q}$,
compact in the sense of the correlation theorem's classification of pieces, under which the
transport piece of a core, smooth or rough, is compact, with first label
$\operatorname{lab}_1(\mathfrak{q}):=K_{\mathfrak{q}}\cup\mathfrak{h}_v$, $K_{\mathfrak{q}}$ the
cube set of $\mathfrak{q}$, and the same support $S_v=\mathfrak{r}_v$, which for that value of
$s^*$ would also be the set read by the summation bound of the correlation theorem; with
$s^*=s(3)$, as here, there is no piece $\mathfrak{q}$ and no label $\operatorname{lab}_1$.
\end{definition}

\begin{proof}
(i) Let $\Delta\in[\square_0]_k$. Then $\Delta$ meets $\square_0$ in its interior, hence meets
$C$ in its interior, because $\square_0\subset C$ by the hypotheses of
Notation~\ref{10.3:not-inner-collar-cores}. The partitions $\pi_p$ nest and $Z$ is a union of
$\pi_k$-cubes, so its component $C$ is a union of $\pi_k$-cubes; as distinct $\pi_k$-cubes
have disjoint interiors, $\Delta$ is one of the cubes of $C$, so $\Delta\subset C$. Thus
$K_v=[\square_0]_k\subset C$. The data $X_v$ and $K_v$ of (a) are those that the correlation
theorem itself assigns to a core, so its class criterion and its interior clause apply to $v$
as they stand; the data (b)--(d) enter neither. The class criterion of the correlation theorem
states: for a plain
$v$ whose $X_v$ is connected and lies in one component $C$, one has
$K_v=[X_v]_k\subset C$ if and only if $X_v\subset Z$, and the interior and the boundary
alternatives are complementary. Here $X_v=\square_0$ is a box contained in $C$, and $v$ is
plain, its live-piece datum $\mathcal{P}_v$ being empty by
Notation~\ref{10.3:not-inner-collar-cores}. Hence $v$ is a plain interior unit of the
first part of the sourced expansion, and the data (a)--(d) do not change this classification;
it belongs to $\mathfrak{L}^{\mathrm{pt}}$, which the plain-unit clause of the correlation
theorem defines as the set of units of the sourced expansion at the site $(\mathrm{L})$ of
$\mathbf{R}_k$ with $\mathcal{P}_v=\emptyset$ and $X_v\subset Z$. By the interior clause of the
correlation theorem, every interior first-part unit ($X_v\subset Z$), plain or wrapped, is
expanded by the interior link rules, with no $\zeta_{\square}$ and no $\mathsf{t}$; this gives
(v), the rule for link $1$ and for links $2$--$4$ being that of the first-part clause of the
sourced expansion and of the link clause of the correlation theorem, and the form of the pieces
and of the value being that of the piece clause of the
correlation theorem. Finally, the class statements of the correlation theorem say that no core
lies in $\mathfrak{L}^{c}$, since $\square^{\sim\nu+1}$ lies on one side of $\partial Z$; this
applies to $v$ because $\square_0\subset C$, so $v\notin\mathfrak{L}^{c}$. By the same class
statements of the correlation theorem, $v\notin\mathfrak{L}^{\mathrm{ex}}$.

(ii) The box $Q_v=Q^{(p_v)}$ is $[\square_0]_{p_v}$ enlarged by two layers of $\pi_{p_v}$-cubes,
so $\square_0\subset Q_v$. Therefore
\begin{equation*}
\mathcal{R}_0=\hat{X}_v\cup X_v=Q_v\cup\square_0=Q_v .
\end{equation*}

(iii) By \eqref{10.3:eq-core-geometry-b} of Lemma~\ref{10.3:lem-core-geometry},
$K_v\subset[Q_v]_k\subset\mathfrak{h}_v$; the second inclusion holds because a $\pi_k$-cube
meeting $Q_v$ in its interior is at $\operatorname{dist}_k$-distance $0$ from $Q_v$, hence
belongs to $\mathfrak{h}_v$. Hence
$\operatorname{lab}_0=K_v\cup\mathfrak{h}_v=\mathfrak{h}_v$. The halo consists of at most
$6^d\le 7^d=c_{\mathfrak{h}}$ cubes by the count of the cubes of the halo in the lemma on
the radius of an inner-collar core.

(iv) By definition,
\begin{equation*}
X^{\oplus}_v=Y(\operatorname{lab}_0)
=\operatorname{lab}_0\cup\{\text{components of } Z \text{ meeting } \operatorname{lab}_0\}.
\end{equation*}
The component $C$ meets $\operatorname{lab}_0$, since $\operatorname{lab}_0\supset K_v$, $K_v$
is non-empty ($\square_0$ has non-empty interior) and $K_v\subset C$ by (i). In the units of
$\operatorname{dist}_k$, the cubes of $\operatorname{lab}_0$ lie within distance $3+1/L<6$ of
$\partial C$, by Lemma~\ref{10.3:lem-core-geometry}, whereas every other component of $Z$ lies
at distance at least $38$ from $C$ in the same units, by the separation radius
$\check{R}_k\ge 38$ of the correlation theorem. Since $6<38$, no component of $Z$ other than $C$
meets $\operatorname{lab}_0$, and, by (iii),
\begin{equation*}
X^{\oplus}_v=\operatorname{lab}_0\cup C=C\cup\mathfrak{h}_v .
\end{equation*}
Accordingly $\mathfrak{E}^{\oplus}_v=\mathbb{C}^{(N)}_k(X^{\oplus}_v)$ is
$\mathbb{C}^{(N)}_k(C\cup\mathfrak{h}_v)$.
\end{proof}

\begin{lemma}[Inner-collar cores are never deep]\label{10.3:lem-never-deep}
Let $v$ be an inner-collar core. We use the notation of Lemma~\ref{10.3:lem-core-geometry} and
Definition~\ref{10.3:def-hull-halo-label}:
\begin{itemize}
\item $X_v=\square_0\subset C$, where $C$ is a component of the large-field set $Z$;
\item $\hat{X}_v=Q_v$ is the hull of $v$;
\item $K_v=[\square_0]_k$, and $\operatorname{lab}_0=K_v\cup\mathfrak{h}_v$ is its start label;
\item $y_v\in\square_0$ is the point of $v$ fixed in Lemma~\ref{10.3:lem-core-geometry}, and
$\Delta_{y_v}$ is the $\pi_k$-cube containing $y_v$.
\end{itemize}
Let $\mathsf{C}$ be the deep core set of the correlation theorem, where $h$ is its base level
and $k=h+N$; here $N=k-h$ is the number of zones and is unrelated to the $N$ of
$\mathrm{SU}(N)$. The deep core set satisfies $\mathsf{C}\subset Z''_{h+1}$.
Recall that a unit is called \emph{deep} if $X_v\subset\mathsf{C}$. Then $v$ is not deep:
\begin{equation}\label{10.3:eq-never-deep-a}
X_v=\square_0\not\subset\mathsf{C},\qquad \hat{X}_v=Q_v\not\subset\mathsf{C};
\end{equation}
consequently the deep term $\hat{D}_v$ that the correlation theorem attaches to $v$ vanishes,
$\hat{D}_v=0$.
Moreover, the cube $\Delta_{y_v}$ belongs to the start label and is a level-$k$ cell of the
collar $C\cap\Omega_k$ of weight three:
\begin{equation}\label{10.3:eq-never-deep-b}
\Delta_{y_v}\in[\square_0]_k=K_v\subset\operatorname{lab}_0,\qquad
\Delta_{y_v}\subset C\cap\Omega_k,\qquad \bar{\mathsf{w}}(\Delta_{y_v})=3.
\end{equation}
\end{lemma}

\begin{proof}
We first prove \eqref{10.3:eq-never-deep-a}.

The cells of $\mathsf{C}$ have levels at most $h$. This comes from the construction of the deep
core set in the correlation theorem, where every cell of $\mathsf{C}$ touching the complement of
$\mathsf{C}$ has level $h$.

The cube $\Delta_{y_v}$, by contrast, has level $k=h+N>h$. Indeed, by
Lemma~\ref{10.3:lem-core-geometry}, in the form \eqref{10.3:eq-core-geometry-c}, $\Delta_{y_v}$
is a cell of the collar $C\cap\Omega_k$, and the correlation theorem records that the cells of
this collar have level $k$. Hence $\Delta_{y_v}$ lies in the zone $k$ of $C$, which is contained
in $C\setminus Z''_{h+1}$; the cells of $C\setminus Z''_{h+1}$ have levels in $]h,k]$, and the
correlation theorem gives $\mathsf{C}\subset Z''_{h+1}$.

Hence $\Delta_{y_v}$ is disjoint from $\mathsf{C}$, so $y_v\notin\mathsf{C}$. Since
$y_v\in\square_0=X_v$, it follows that $X_v\not\subset\mathsf{C}$. This is the first half of
\eqref{10.3:eq-never-deep-a}.

The second half holds independently. The hull satisfies $\hat{X}_v=Q_v\not\subset C$, by
Definition~\ref{10.3:def-hull-halo-label}, and $C\supseteq\mathsf{C}$; hence
$\hat{X}_v\not\subset\mathsf{C}$. This is the criterion by which the correlation theorem declares
such units never deep. By the definition of deep units there, the deep term $\hat{D}_v$ that the
correlation theorem attaches to $v$ vanishes.

We now prove \eqref{10.3:eq-never-deep-b}. The cube $\Delta_{y_v}$ contains $y_v$, and
$y_v\in\square_0$. Therefore $\Delta_{y_v}$ is a $\pi_k$-cube meeting $\square_0$ in its
interior. By the definition of $[\,\cdot\,]_k$ as the set of $\pi_k$-cubes meeting a set in
interior, we get
\begin{equation*}
\Delta_{y_v}\in[\square_0]_k=K_v.
\end{equation*}
Since $\operatorname{lab}_0=K_v\cup\mathfrak{h}_v$ by Definition~\ref{10.3:def-hull-halo-label},
the inclusion $K_v\subset\operatorname{lab}_0$ is immediate.

Next, by Lemma~\ref{10.3:lem-core-geometry}, in the form \eqref{10.3:eq-core-geometry-c},
$\Delta_{y_v}$ is a cell of the collar $C\cap\Omega_k$ with $\bar{\mathsf{w}}(\Delta_{y_v})=3$.
Its level is $k$: the correlation theorem records that the cells of the collar $C\cap\Omega_k$
have level $k$. This proves \eqref{10.3:eq-never-deep-b}.
\end{proof}

\begin{definition}[Dial, norms, loads and table of the vector package]\label{10.3:not-vector-dial-table}
In this section the letters of the vector package keep the meaning they have there. We recall those
used below and fix how their indices are read. Throughout, $n$ is a zone of the run with anchor
index $k'$, $j\le n$ is a level, and $\mathsf{h}$ is the test field.

\begin{enumerate}
\item[(a)] \emph{The dial.} The production with the vector part $\mathfrak{u}$ added is the dialled
production, in which each use of the production at anchor index $k'$ and zone $n$ (a use with index
$(k',n,\cdot)$) also reads $\lambda_{n,k'}f_v$, $v\in\mathfrak{u}$. Here
$f_v$ is the function of the unit $v$ and $\lambda_{n,k'}$ is the \emph{dial}. Thus at zone $n$ of the run
with anchor index $k'$ the uses of the production read the dialled units $\lambda_{n,k'}f_v$.

\item[(b)] \emph{Norms and loads.} The norms are taken per unit of $\sup|\mathsf{h}|$. The suprema are
over $|\mathsf{h}|\le 1$ throughout and, for the frozen objects, also over $\hat{w}\in D(0,2r)$, the disc
of radius $2r$ of the vector package. In terms of the frozen vector part $\hat{\mathfrak{V}}^{(j)}$ of
level $j$, the frozen object $\hat{\mathsf{O}}_j$ of level $j$, whose norm majorises that of
$\hat{\mathfrak{V}}^{(j)}$, and the object $\mathsf{R}^{V(i)}$ of level $i$ of the
vector package, measured in its norms $\|\cdot\|^{\mathrm{fr}}$ and $\|\cdot\|^{\flat}$, they are
\begin{equation}\label{10.3:eq-vector-dial-table-1}
\mathsf{u}_j:=\|\hat{\mathfrak{V}}^{(j)}\|^{\mathrm{fr}}\le\mathsf{o}_j:=\|\hat{\mathsf{O}}_j\|^{\mathrm{fr}},
\qquad
\mathsf{r}_i:=\|\mathsf{R}^{V(i)}\|^{\flat},\qquad \mathsf{r}_0:=0 .
\end{equation}
The typed statements of the vector package and the lemmas that accompany them hold at every anchor with
the same constants. The
\emph{loads} are
\begin{equation}\label{10.3:eq-vector-dial-table-2}
T^V_n:=\sum_{0\le j\le n}\varsigma^V_{1+n-j}\,\mathsf{u}_j
+\sum_{n-\mathsf{a}_i\le i\le n}w_{n,i}\,\mathsf{r}_i+\mathbf{1}_{\{n=0\}}\,\mathsf{t}^{\mathrm{raw}},
\qquad \mathsf{t}^{\mathrm{raw}}:=c_0\,C_{\mathfrak{r}}'\,\mathfrak{g}_0 ,
\end{equation}
where $C_{\mathfrak{r}}'$ and $\mathfrak{g}_0$ are constants of the vector package, and $c_0$ is a
constant of the vector package, not the constant $c_0$ of $\underline{\lambda}=c_0/4$.
The scale weights $\varsigma^V_{\nu}$ are defined in (c) below.
For $1\le j\le k$, where $k$ is the level at which the run is read (so that the current zone is $n=k$),
the $\sharp$-cores of level $j$ are the members of the family
\begin{equation*}
\mathsf{x}^{\sharp}_j:=\bigl\{\hat{\mathfrak{V}}^{\sharp(j)}_{\hat{w}(y)}(X,\cdot,y;\mathsf{h}(y))\colon
y\in\Lambda_j^{\circledR},\ X\subset\Lambda_j\bigr\},
\end{equation*}
where $\Lambda_j$ is the level-$j$ set of the vector package over whose subsets $X$ the family runs, and
$\Lambda_j^{\circledR}$ is its set of level-$j$ anchor points $y$.
These are frozen universal objects. They do not depend on the vector field $\eta$, on the test functions,
or on the torus.

\item[(c)] \emph{The table.} The vector package bounds the units by a table. It runs over all sites, in
zone $n$, with $\tilde{\square}\in\pi_n$ the current-zone cube, and its sizes are per unit of
$\sup|\mathsf{h}|$. Its columns are: the unit $v$; its hull $S_v$; its size $\mathfrak{N}(v)$; the
statement from which the size is taken; and the entry per unit $n$-cell. The row of the $\sharp$-core
(row~1) is
\begin{equation}\label{10.3:eq-vector-dial-table-a}
\begin{array}{l|l|l|l}
\text{unit } v & S_v & \mathfrak{N}(v) & \text{per unit $n$-cell}\\ \hline
\text{$\sharp$-core, level $j$, at $y$} & \square^{\sim\nu+1}
& C_{\sharp}\,\mathsf{u}_j\,\nu^3\rho^2(1+\rho)\,t^{4+2\hat{\alpha}} & C_{\sharp}\,\mathsf{u}_j\,\varsigma^V_{\nu}
\end{array}
\end{equation}
Here $\nu$, $\rho$ and $t$ are the letters of the row; they are identified below, in
\eqref{10.3:eq-vector-dial-table-3} and the paragraph following it, as $\nu=1+n-j$,
$\rho=\rho_{n,j}$, the run ratio defined in (f), and $t=L^{j-n}$. The column recording the statement
from which the size
$\mathfrak{N}(v)$ is taken is omitted from the display; for row~1 it refers to items (c) and (d) of the
core lemma of the vector package and, where the correlation theorem reads smooth or rough cores, to the
smooth-core variant of that lemma.
The last entry of the table's row~0, whose weight $\varsigma^V_{1+n}$ is the value of
$\varsigma^V_{1+n-j}$ at $j=0$, is at most
$C\,\mathsf{v}_0\,\varsigma^V_{1+n}$, where $C$ is the constant of that row and
$\mathsf{v}_0=\mathsf{v}_0(c_0)$ is the number of the recursion of the vector package, a function of $c_0$.
We write $C_L^V$ for the
largest constant of the last column, taken with the weights $\mathsf{w}_{\mathrm{site}}$ and
$\mathsf{w}_L$ of the lemma of the correlation theorem recalled in (e).

\emph{Meaning of the last column.} The entry per unit $n$-cell is obtained as follows. Fix one unit cell
of level $n$ in zone $n$. Take the units of the row anchored in that cell, and add up their sizes
$\mathfrak{N}(v)$. The entry is this total, per unit of $\sup|\mathsf{h}|$. For row~1 the units anchored
in one unit cell of level $n$ are the $\sharp$-cores of level $j$ at the $t^{-4}=L^{4(n-j)}$ anchor points
of level $j$ that the cell contains. A cube of $\pi_n$, of side $M$, is the union of $M^4$ unit cells of
level $n$, and therefore carries $M^4$ times the entry. The vector package uses the
column in this sense when it sums the column of the table. Row~1 anchors at the points $y$ of (b), and
these are the anchor points $y_v$ of the units of this section. Hence replacing $S_v$ by $\mathfrak{r}_v$
in this section does not change the column.

\emph{The indices of row~1.} In \eqref{10.3:eq-vector-dial-table-a} the index is
\begin{equation}\label{10.3:eq-vector-dial-table-3}
\nu=1+n-j ,
\end{equation}
which is the age of the level-$j$ core in zone $n$, counted so that a core of level $n$ has age one. The
two remaining letters of row~1 are $t=L^{j-n}$ and $\rho=\rho_{n,j}=\rho_{n,n+1-\nu}$, the run ratio
of (f) at the levels $n$ and $j$. Consequently the row-1 entry per unit $n$-cell at zone $n$ and level
$j$ is
\begin{equation}\label{10.3:eq-vector-dial-table-4}
C_{\sharp}\,\varsigma^V_{1+n-j}\,\mathsf{u}_j ,
\end{equation}
that is, $C_{\sharp}$ times the $j$-th summand of the first sum in
\eqref{10.3:eq-vector-dial-table-2}.

\emph{Signs.} The quantities $\mathsf{u}_j$ and $\mathsf{r}_i$ are norms, hence nonnegative, and
$\mathsf{t}^{\mathrm{raw}}>0$. The quantities $\varsigma^V_{\nu}$ and $w_{n,i}$ are the per-cell scale
weights of rows $0$, $1$ and $3$ of the table; the first of them is
\begin{equation*}
\varsigma^V_{\nu}:=\nu^3\rho^2(1+\rho)\,L^{-2\hat{\alpha}(\nu-1)},\qquad \rho=\rho_{n,n+1-\nu},
\end{equation*}
with the run ratio of (f). Of the weights $w_{n,i}$, and of the windows $\mathsf{a}_i$, only
their indices and their nonnegativity are used below.

\item[(d)] \emph{The domain and the constant.} There is a coupling ceiling $\gamma_V\le\hat{\gamma}_L$,
where $\hat{\gamma}_L$ is the coupling ceiling of the second-order ceiling statement
read at the constant
$C_L^{V\prime}$ defined below, with the following property. For $g\le\gamma_V$, every anchor
$K''\le K'$ of the vector package, where $K'$ is the anchor of the vector package that bounds the
anchors admitted here, every level $k<K''$, every source $w$ in the admissible source class
$\mathcal{W}^{[K'']}_k$ of the vector package at level $k$, and every
bounded $\mathsf{h}$, and with the split rule under which a
unit $v\in\mathfrak{u}_k$ is assigned to the first part if and only if $S_v\cap Z\neq\emptyset$, every
derived object of the dialled production with the vector part $\mathfrak{u}_k[w,\mathsf{h}]$ added is
holomorphic on
\begin{equation}\label{10.3:eq-vector-dial-table-5}
\mathbb{D}_V:=\bigl\{\,|\lambda_{n,k'}|\,\sup|\mathsf{h}|\,T^V_n\le E_1\,\bigr\},
\end{equation}
and obeys there the bounds of the corresponding form of the extension proposition. We read
$\mathbb{D}_V$ as the polydisc
condition imposed at every index $(n,k')$ of the run. It is the standing domain of all the typed
statements of the vector package. Here $E_1:=2(S_{\varsigma}+1)E_0$ is the number fixed with the
dialled production, where $S_{\varsigma}$ is its scalar scale series and $E_0$ is a number of the
dialled production, and we put $C_L^{V\prime}:=C_L^V E_1$.

\item[(e)] \emph{Sizes and site weight.} The lemma of the correlation theorem that bounds the terms
$v_p$ into which it expands a unit $v$, indexed there by $p$ (an index of that lemma, not the scale
letter $p$ of (c) and (f)), has the following form; in it $\eta_p$, $m(p)$, $\mathsf{R}_p$ and
$\rho_A(v)$ are the prefactor, the exponent, the map and the radius attached there to the index $p$ and
the unit $v$ (the prefactor $\eta_p$ is unrelated to the vector field $\eta$ of (b)). It assumes the
size hypothesis
$|f_v\circ\mathsf{R}_p-f_v(1)|\le\mathfrak{N}(v)$ on the disc $|\mathsf{t}|<\rho_A(v)$. It concludes that
\begin{equation*}
|v_p|\le 2\,\mathsf{w}_{\mathrm{site}}^{-1}\,\eta_p\,e^{-(\kappa_1-1)m(p)}\,\mathfrak{N}(v).
\end{equation*}
A further statement accompanying this lemma carries the size $\tfrac{81}{64}\mathfrak{N}^{\flat}(v)$
and the weight $\mathsf{w}_L$.
Accordingly $\mathfrak{N}(v)$ and $\mathfrak{N}^{\flat}(v)$ denote sizes, and
$\mathsf{w}=\mathsf{w}_{\mathrm{site}}$ denotes the site weight. The vector package writes $\mathsf{W}$
and $\mathsf{W}_J$ for the per-cell total of the weighted output of this lemma. The left member of the
inner-collar estimate of this section is a sum of sizes. The inequality for $\mathsf{W}$ is therefore not
used in this section.

\item[(f)] \emph{The size assigned to the inner-collar cores.} Let $v\in\mathfrak{L}^{V,\mathrm{col}}$ be
an inner-collar core of level $j$ in zone $n$ (Notation~\ref{10.3:not-inner-collar-cores}), and let
$s^*=s(3)$, $a^*=2+s^*$ and $p^*=n-a^*$ be as there, $s(\cdot)$ being the integer function of the
correlation theorem, so that $\nu=1+n-j\ge a^*+1$ by the definition of the class, hence $j\le p^*$.
Put $\nu_{p^*}:=1+p^*-j=\nu-a^*\ge 1$ and $t_{p^*}:=L^{j-p^*}$. For levels $i\le l$
of the run let
\begin{equation*}
\rho_{l,i}:=\frac{\alpha_{\cdot}(g_l)}{\alpha_{\cdot}(g_i)}
\end{equation*}
be the \emph{run ratio}, where $\alpha_{\cdot}$ is the radius function of the correlation theorem, read in
the form $\alpha_{\cdot}(g)=c\,g(\log g^{-2})^{q}$ of the radii of that theorem, with a constant $c>0$ and
an integer $q\ge 2$ (only the ratio, from which $c$ cancels, is used), and
$g_l$ is the coupling of the run at level $l$; we write $x_l=g_l^{-2}$. The size of $v$ is
\begin{equation*}
\begin{gathered}
\mathfrak{N}^{\flat}(v):=|\lambda_{n,k'}|\cdot 4\,\epsilon_{a^*}^2\,\mathfrak{N}^{V,\mathrm{run}}_{p^*}(v),\\
\mathfrak{N}^{V,\mathrm{run}}_{p^*}(v):=\mathsf{u}_j\,\sup|\mathsf{h}|\,
\bigl[C_{\sharp}\,\nu_{p^*}^3\,\rho_{p^*,j}^2\,(1+\rho_{p^*,j})+2K(d)(1+8C_{\mathfrak{J}})\bigr]\,
t_{p^*}^{4+2\hat{\alpha}},
\end{gathered}
\end{equation*}
where
\begin{equation*}
\epsilon_{a^*}=3\,c_{\epsilon}\,(2a^*+11)\,\bar{\rho}_{a^*}\,L^{-2(a^*-1)}\le 1
\end{equation*}
is the gauge smallness of Notation~\ref{10.3:not-inner-collar-cores}, with $\bar{\rho}_{a^*}$ as below,
$K(d)$ is the constant of the correlation theorem bounding the sum over connected domains that contain a
given cube, and $C_{\mathfrak{J}}$ is the constant of the norm bound of the vector package on the cores;
these two enter the bound of the core tails, and we put $C_V:=C_{\sharp}+2K(d)(1+8C_{\mathfrak{J}})$.
This is the size of the
dialled unit $\lambda_{n,k'}f_v$ read at the one scale $p^*$, with the run ratio $\rho_{p^*,j}$ of that
scale. Let $\gamma^{\circ}=(8c_2+4)^{-1/2}$ be the first coupling ceiling of the correlation theorem, and
let
\begin{equation*}
\begin{split}
\gamma_{\rho}:=\sup\Bigl\{\gamma'\in\,\bigl]0,\min\{\gamma^{\circ},(3/4)^{1/2}e^{-q}\}\bigr]\colon
&\ (1+4c_2\gamma'^2)\\
&\times\Bigl(1+\frac{(3\lambda a^*+2c_2)\,\gamma'^2}{\log\gamma'^{-2}}\Bigr)^{2q}\le 2\Bigr\},
\end{split}
\end{equation*}
a positive number depending on $c_2$, $q$, $\lambda$ and $a^*$ only, where $c_2$ and $\lambda$ are the
numbers of the two-sided bound
\begin{equation*}
\lambda(l-i)-2c_2\le x_i-x_l\le 3\lambda(l-i)+2c_2\qquad(0\le i<l\le K)
\end{equation*}
that the run obeys in the correlation theorem. Assume the standing hypotheses $g\le\gamma^{\circ}$ and
$x_i>g^{-2}$ for the levels $i$ of the run, and $g\le\gamma\le\min\{\gamma^{\circ},\gamma_{\rho}\}$;
the condition $\gamma\le\gamma_{\rho}$ is one more one-sided smallness condition on the coupling,
chosen after $L$, $q$ and the numbers $c_2$, $\lambda$, $a^*$; it joins the coupling ceilings of the run,
lowers the admissible coupling, and adds no hypothesis to Theorem~0 and no condition on $L$. Then
\begin{equation*}
\rho_{p^*,j}\le 2^{1/2}\,\rho_{n,j},\qquad \rho_{n,j}\ge(3/4)^{1/2},
\end{equation*}
and the size satisfies
\begin{equation}\label{10.3:eq-vector-dial-table-b}
\mathfrak{N}^{\flat}(v)\le|\lambda_{n,k'}|\,\sup|\mathsf{h}|\,F'_{\mathrm{col}}\,C_{\sharp}\,\mathsf{u}_j\,
\nu^3\,\rho_{n,j}^2\,(1+\rho_{n,j})\,t^{4+2\hat{\alpha}} ,
\end{equation}
where $j$ is the level of $v$, $\nu=1+n-j$ and $t=L^{j-n}$ as in \eqref{10.3:eq-vector-dial-table-3}.
The factor $C_{\sharp}\,\mathsf{u}_j\,\nu^3\rho_{n,j}^2(1+\rho_{n,j})\,t^{4+2\hat{\alpha}}$ is the size
$\mathfrak{N}(v)$ of row~1 in \eqref{10.3:eq-vector-dial-table-a} per unit of $\sup|\mathsf{h}|$, with
$\rho=\rho_{n,j}$ as in (c). The constants are
\begin{equation}\label{10.3:eq-vector-dial-table-6}
\begin{gathered}
F_{\mathrm{col}}:=36\,c_{\epsilon}^2\,(2a^*+11)^2\,\bar{\rho}_{a^*}^2\,L^{4+2\hat{\alpha}a^*}\,
\frac{C_V}{C_{\sharp}},\qquad F'_{\mathrm{col}}:=2^{3/2}\,F_{\mathrm{col}},\\
F^{*\prime}_{\mathrm{col}}:=\max\Bigl\{F'_{\mathrm{col}},\ \frac{27\,C^{\mathrm{mar}}\,\bar{\rho}_{a^*}^2\,
L^{2\hat{\alpha}(a^*-1)}}{4\,C_{\sharp}}\Bigr\},
\end{gathered}
\end{equation}
where $\bar{\rho}_{a^*}$ is the member of index $a^*$ of the radii sequence $(\bar{\rho}_b)$ of the
correlation theorem. These constants depend on no level, zone or age, and not on
$(\lambda_{n,k'},\mathsf{h})$.
The constant $F^{*\prime}_{\mathrm{col}}$ is the larger of the two constants produced in the derivation
of the inner-collar constant. Its second
member is the constant of the young cores, namely those of age $\nu\le a^*$, which are bounded by the
marginal bound of the correlation theorem with the marginal constant $C^{\mathrm{mar}}$.
This second member is not used for the class $\mathfrak{L}^{V,\mathrm{col}}$; it enters below only
through the majorant $F^{*\prime}_{\mathrm{col}}\ge F'_{\mathrm{col}}$. The positivity of
$\gamma_{\rho}$, the ratio comparison and the
inequality \eqref{10.3:eq-vector-dial-table-b} are proved in the lemma of this section that derives the
inner-collar constant.
\end{enumerate}
\end{definition}

\begin{proof}[Proof of the index reading in (c)]
We prove \eqref{10.3:eq-vector-dial-table-3}, the values of $t$ and $\rho$, the entry
\eqref{10.3:eq-vector-dial-table-4}, and the signs.

\emph{The hull of row~1.} The hull $\square^{\sim\nu+1}$ is written in the notation of the correlation
theorem, which the table borrows. There it is a cube of $2(\nu+1)+1=2\nu+3$ level-$j$ blocks of side $M$
in each direction. Its side is therefore the level-$j$ length $(2\nu+3)M$.

\emph{Conversion to level $n$.} In units of the sites of level $n$ a level-$j$ length is multiplied by
$L^{-(n-j)}$. For the side of the hull this gives $(2\nu+3)ML^{-(n-j)}$.

\emph{Comparison with the count of the vector package.} When the vector package sums the column of the
table, it bounds the $n$-side of the row-1 hull by
\begin{equation*}
(2\nu+3)\,M\,L^{-(\nu-1)}\le 5M .
\end{equation*}
It concludes that at most $5^d$ current cubes are met. The two expressions for the $n$-side agree
exactly when their exponents of $L$ agree, that is, when $n-j=\nu-1$. This is
\eqref{10.3:eq-vector-dial-table-3}, $\nu=1+n-j$. At $j=n$ it gives $\nu=1$, which fixes the
normalisation of the age stated in (c).

\emph{Two consistency checks.} The first sum of the load \eqref{10.3:eq-vector-dial-table-2} carries the
weight $\varsigma^V_{1+n-j}$ at the norm $\mathsf{u}_j$ of level $j$. This is $\varsigma^V_{\nu}$ under
\eqref{10.3:eq-vector-dial-table-3}. Row~0 carries $\varsigma^V_{1+n}$, the value of $\varsigma^V_{\nu}$
at $j=0$.

\emph{The entry per unit $n$-cell.} Substitute \eqref{10.3:eq-vector-dial-table-3} into the last entry
$C_{\sharp}\,\mathsf{u}_j\,\varsigma^V_{\nu}$ of \eqref{10.3:eq-vector-dial-table-a}. This gives
$C_{\sharp}\,\varsigma^V_{1+n-j}\,\mathsf{u}_j$, which is \eqref{10.3:eq-vector-dial-table-4}. By the
previous paragraph it is $C_{\sharp}$ times the $j$-th summand of the first sum of $T^V_n$.

\emph{The letters $t$ and $\rho$.} The letter $t$ of row~1 is the letter $t=L^{j-n'}$ of the
inner-collar estimate of this section, taken at $n'=n$; so $t=L^{j-n}$. The letter $\rho$ of row~1 is
the run ratio $\rho_{n,j}$ of (f): the size of row~1 is taken from the core-size statement of the
vector package, whose bound for a background admissible at scale $n\ge j$ carries the factor
$\rho_{n,j}^2(1+\rho_{n,j})$. By \eqref{10.3:eq-vector-dial-table-3}, $n+1-\nu=j$, so $\rho_{n,j}$ is
also the ratio $\rho_{n,n+1-\nu}$ of the weight $\varsigma^V_{\nu}$ in (c). One unit cell of level $n$
contains $t^{-4}=L^{4(n-j)}$ anchor points of level $j$, and each carries one row-1 unit of size
$C_{\sharp}\,\mathsf{u}_j\,\nu^3\rho_{n,j}^2(1+\rho_{n,j})\,t^{4+2\hat{\alpha}}$ per unit of
$\sup|\mathsf{h}|$. Since
$t^{2\hat{\alpha}}=L^{-2\hat{\alpha}(n-j)}=L^{-2\hat{\alpha}(\nu-1)}$, the definition of
$\varsigma^V_{\nu}$ gives for the total over that cell the identity
\begin{equation*}
\begin{split}
t^{-4}\,C_{\sharp}\,\mathsf{u}_j\,\nu^3\rho_{n,j}^2(1+\rho_{n,j})\,t^{4+2\hat{\alpha}}
&=C_{\sharp}\,\mathsf{u}_j\,\nu^3\rho_{n,j}^2(1+\rho_{n,j})\,L^{-2\hat{\alpha}(\nu-1)}\\
&=C_{\sharp}\,\mathsf{u}_j\,\varsigma^V_{\nu},
\end{split}
\end{equation*}
so that this total is the entry per unit $n$-cell of row~1.

\emph{Signs.} By \eqref{10.3:eq-vector-dial-table-1}, $\mathsf{u}_j$ and $\mathsf{r}_i$ are norms, hence
nonnegative. By its definition in \eqref{10.3:eq-vector-dial-table-2}, $\mathsf{t}^{\mathrm{raw}}>0$. The
weight $\varsigma^V_{\nu}$ is nonnegative by its definition in (c), $\rho_{n,n+1-\nu}$ being a ratio of
positive radii. The weights $w_{n,i}$ enter row~$3$ as majorants of nonnegative sizes, and they are
nonnegative.
\end{proof}

\begin{lemma}[Radius, support and halo of an inner-collar core]\label{10.3:lem-core-radius}
Let $v$ be an inner-collar core, read at the level $n=k$, with base point $y=y_v$. Its hull
$\hat{X}_v:=Q_v=Q^{(p^*)}(y)$, its halo
$\mathfrak{h}_v=\{\Delta\in\pi_k:\operatorname{dist}_k(\Delta,\hat{X}_v)\le 2\}$, its start label
$\operatorname{lab}_0=K_v\cup\mathfrak{h}_v$ with $K_v=[\square_0]_k$, and its support
$\mathfrak{r}_v=\bigcup\operatorname{St}(\Delta_y)$ are as in
Definition~\ref{10.3:def-hull-halo-label}. Here $\Delta_y$ is the cell containing $y$ of the
partition of the correlation theorem, whose cells have levels in $\,]h,k]$ by
\eqref{10.3:eq-never-deep-a}, and $\operatorname{St}(\Delta_y)$ is its star.
Put
\[
\operatorname{rad}(v):=\max_{x\in\hat{X}_v}\operatorname{dist}_n(x,y),
\]
where $\operatorname{dist}_n$ is the sup-norm distance in level-$n$ units, in which a $\pi_n$-cube
has side $1$.
Let $c_{\mathfrak{h}}:=7^d$, let $c_l=1+234d$ be the constant of the correlation theorem, and let
$c_S$ be the cardinality constant of the lemma on the star of a cell.
Then
\begin{equation}\label{10.3:eq-core-radius-a}
\begin{aligned}
&\operatorname{rad}(v)\le \tfrac{1}{L}<1\le 5,\\
&\mathfrak{r}_v\subset Q_5(\Delta_y),\qquad \#\operatorname{St}(\Delta_y)\le c_S,\qquad
[\mathfrak{r}_v]_k\subset\mathfrak{h}_v,\\
&\#\mathfrak{h}_v\le 6^d\le c_{\mathfrak{h}},\qquad d_k(\mathfrak{h}_v)\le 6^d-1,\\
&d_k(\operatorname{lab}_0)=d_k(\mathfrak{h}_v)\le 6^d-1\le d_k(K_v)+c_lc_{\mathfrak{h}} .
\end{aligned}
\end{equation}
We put $c_0:=6^d+c_lc_{\mathfrak{h}}$.
Moreover $\mathfrak{h}_v$ is wall-connected. If $\varsigma_1\ge 0$ is the weight rate of the correlation
theorem, then $e^{\varsigma_1\operatorname{rad}(v)}\le e^{\varsigma_1/L}\le e^{\varsigma_1}$.
\end{lemma}

\begin{proof}
\emph{The unit of $\operatorname{dist}_n$.}
Distances $\operatorname{dist}_n$ are counted in $\pi_n$-cubes, so that a $\pi_n$-cube has side $1$.
This is the unit in which the correlation theorem counts cubes, in four places:
\begin{itemize}
\item the halo $\mathfrak{h}_v$ holds at most $7^d=c_{\mathfrak{h}}$ cubes;
\item the lemma on the star of a cell calls $\{\operatorname{dist}_{n'}(\cdot,\Delta')\le 1\}$ the cube of
$\Delta'$;
\item distinct components made of $\pi_k$-cubes are at least $\check{R}_k$ units apart;
\item for cores the theorem states the bound $\operatorname{rad}\le 5$.
\end{itemize}
The last bound is the $n$-diameter of the hull $\square^{\sim\nu+1}$ of a rough core, that is, of a
core whose hull is the box $\square^{\sim\nu+1}$ of side $2\nu+3$ level-$j$ blocks, at $\nu=1$:
\[
(2\nu+3)L^{-(\nu-1)}=5 \quad\text{at }\nu=1 ,
\]
that is, five $\pi_n$-cubes.
A $\pi_k$-cube has side $M$ sites of the current lattice and a $\pi_p$-cube side $ML^{-(k-p)}$ sites;
hence a $\pi_p$-cube has side $L^{-(k-p)}$ level-$k$ units.

\emph{The radius.}
We have $n=k$ and $Q_v=Q^{(p^*)}(y)$. By \eqref{10.3:eq-core-geometry-a}, $Q_v$ has side at most
$2a^*+11\le L^{a^*-1}$ $\pi_{p^*}$-cubes. Since $p^*=n-a^*$, a $\pi_{p^*}$-cube has side $L^{-a^*}$
level-$k$ units. Hence the side of $Q_v$ in level-$k$ units is at most
\[
L^{a^*-1}\cdot L^{-a^*}=L^{-1}.
\]
The point $y$ lies in $Q_v$. Two points of a box are at sup-norm distance at most its side, so
\[
\operatorname{rad}(v)=\max_{x\in Q_v}\operatorname{dist}_k(x,y)\le\tfrac{1}{L}
\]
in the sup-norm. If $\operatorname{dist}_n$ is read as Euclidean distance, the same argument gives
$\operatorname{rad}(v)\le\sqrt{d}/L=2/L$ at $d=4$. This is again less than $1$, since $L>2$.
In either reading $\operatorname{rad}(v)<1\le 5$, so the bound
$\operatorname{rad}\le 5$ of the correlation theorem applies to $v$.

The weight rate $\varsigma_1$ of the correlation theorem is nonnegative: the theorem's factor
$e^{5\varsigma_1}$ is at least $1$. Monotonicity of the exponential then gives
\[
e^{\varsigma_1\operatorname{rad}(v)}\le e^{\varsigma_1/L}\le e^{\varsigma_1}.
\]

\emph{The support.}
We use three properties of the star $\operatorname{St}(\Delta')$ of a cell $\Delta'$ of level $n'$,
taken from the lemma on the star of a cell.
\begin{enumerate}
\item[(i)] The members of $\operatorname{St}(\Delta')$ have levels $n'-1$, $n'$, $n'+1$.
\item[(ii)] $\#\operatorname{St}(\Delta')\le c_S$.
\item[(iii)] $\bigcup\operatorname{St}(\Delta')$ is contained in
$\{\operatorname{dist}_{n'}(\cdot,\Delta')\le L+2\}$. The number $L+2$ may be replaced by $2$ if no
member has level $n'+1$. If $y_v\in\Delta'$, then $[\operatorname{St}(\Delta')]_k\subset\mathfrak{h}_v$.
\end{enumerate}
Apply these with $\Delta'=\Delta_y$, so that $n'=k$. By \eqref{10.3:eq-never-deep-a}, the cells of the
partition have levels in $\,]h,k]$. Hence no cell has level $k+1$, and by (i) no member of
$\operatorname{St}(\Delta_y)$ has level $k+1$. The sharpened form of (iii) therefore gives
\[
\mathfrak{r}_v=\bigcup\operatorname{St}(\Delta_y)\subset\{\operatorname{dist}_k(\cdot,\Delta_y)\le 2\}
=Q_5(\Delta_y),
\]
which is a box of $k$-side at most $5$. Property (ii) gives $\#\operatorname{St}(\Delta_y)\le c_S$.
Since $y_v=y\in\Delta_y$, the last part of (iii) gives $[\mathfrak{r}_v]_k\subset\mathfrak{h}_v$.

\emph{The halo.}
By \eqref{10.3:eq-core-geometry-b}, $Q_v$ meets at most $2$ cubes of $\pi_k$ in each direction. A cube
$\Delta\in\pi_k$ with $\operatorname{dist}_k(\Delta,Q_v)\le 2$ lies, in each direction, within two cube
widths of these. Since $\operatorname{dist}_k$ is the sup-norm distance, the cubes within distance $2$
of the box $Q_v$ form a box of $\pi_k$-cubes with at most $2+2+2=6$ cubes per direction. Hence
\[
\#\mathfrak{h}_v\le 6^d\le 7^d=c_{\mathfrak{h}},
\]
and $\mathfrak{h}_v$, being a box of cubes, is wall-connected. A wall-connected family of $N$ cubes has
a spanning tree of $N-1$ unit steps, so $d_k(\mathfrak{h}_v)\le 6^d-1$.

\emph{The start label.}
By Lemma~\ref{10.3:lem-core-geometry}, $K_v\subset[Q_v]_k$. Every $\Delta\in[Q_v]_k$ meets $Q_v$ and
so satisfies $\operatorname{dist}_k(\Delta,Q_v)=0\le 2$. Hence $[Q_v]_k\subset\mathfrak{h}_v$, so
$K_v\subset\mathfrak{h}_v$ and $\operatorname{lab}_0=K_v\cup\mathfrak{h}_v=\mathfrak{h}_v$. Therefore
\[
d_k(\operatorname{lab}_0)=d_k(\mathfrak{h}_v)\le 6^d-1 .
\]
Since $d_k(K_v)\ge 0$, $c_l=1+234d\ge 1$ and $c_{\mathfrak{h}}=7^d\ge 6^d$, we get
\[
6^d-1\le c_lc_{\mathfrak{h}}\le d_k(K_v)+c_lc_{\mathfrak{h}} .
\]
This is the label-length bound in the form used by the correlation theorem: the tree length of a
start label exceeds that of its core part, here $K_v$, by at most $c_lc_{\mathfrak{h}}$, with
$c_{\mathfrak{h}}=7^d$. It enters the interior bound of the correlation theorem for pieces anchored
at the start label $K_v\cup\mathfrak{h}_v$, with the constant $c_0$ of the statement.
\end{proof}

\begin{remark}[Why one extension covers every list]\label{10.3:rem-one-extension}
Let $v$ be an inner-collar core, as in Definition~\ref{10.3:def-hull-halo-label} and
\eqref{10.3:eq-hull-halo-label-a}. Its box $X_v=\square_0=\square^{\sim\nu+1}$ lies in one
component $C$ of $Z$, with $K_v=[\square_0]_k\subset C$. Its read box $Q_v=Q^{(p_v)}$ is not
contained in $Z$. Its hull is $\hat{X}_v=Q_v$, and its halo $\mathfrak{h}_v$ consists of the
$\pi_k$-cubes $\Delta$ with $\operatorname{dist}_k(\Delta,\hat{X}_v)\le 2$. Its start label is
$\operatorname{lab}_0=K_v\cup\mathfrak{h}_v$, and its consumer is
$\mathfrak{E}^{\oplus}_v=\mathbb{C}^{(N)}_k(X^{\oplus}_v)$, where
$X^{\oplus}_v=Y(\operatorname{lab}_0)=C\cup\mathfrak{h}_v$.

The correlation theorem has lists of units indexed by exposure, by far cores and by smooth cores.
The inclusion of $v$ in these lists requires no modification of that theorem's construction beyond
the declaration of Definition~\ref{10.3:def-hull-halo-label}. That declaration fixes four
things at once: the level $p_v$, the box $Q_v$, the exposure of $v$ and its filing. Once these are
fixed, the criteria that the correlation theorem states for its lists place $v$ in each of them.
\end{remark}

\begin{proof}
\emph{The pattern for far cores.} The correlation theorem states a rule for the smooth far cores of
a plain second-part unit, and Definition~\ref{10.3:def-hull-halo-label} follows this pattern for
$v$, with the level $p_v$ fixed there. For a smooth far core $u$ of such a unit the level is
\begin{equation*}
p_u:=n-2-s^{(9/8)}(\bar{\mathsf{w}}_u).
\end{equation*}
No condition $Q^{(p)}\subset Z$ is imposed. The box $Q_u$, which lies in the enlargement of
$[X_u]_k$ specified in that rule, is kept inside the domain of the consumer, which contains
the label at which the consumer sits ($\operatorname{lab}_1\supseteq\mathfrak{h}_u$ for the far
cores of the correlation theorem, $\operatorname{lab}_0=K_v\cup\mathfrak{h}_v\supseteq\mathfrak{h}_v$
for $v$), and not inside $Z$.

In the same rule a core's box $X_u$ is what is filed, through $K_u\subset\mathfrak{h}_u$; it is not
what the unit reads. For $v$ this box is $X_v=\square_0=\square^{\sim\nu+1}$, filed through
$K_v\subset\mathfrak{h}_v$.

The reason given there is a property of the consumer $\mathfrak{E}^{\oplus}_v$: its domain
contains $\mathfrak{h}_v\supseteq[Q_v]_k$ and lies in the filing domain $Y_4(\mathfrak{p})$ of the
pieces $\mathfrak{p}=(v;\mathsf{y}_1,\dots,\mathsf{y}_4)$ of $v$, where $\mathsf{y}_t$ is the
link-$t$ family. For $v$ this property holds as follows.
\begin{itemize}
\item The domain of $\mathfrak{E}^{\oplus}_v$ is $C\cup\mathfrak{h}_v$, which contains
$\mathfrak{h}_v$.
\item Every $\pi_k$-cube meeting $Q_v$ is at $\operatorname{dist}_k$-distance $0$ from
$\hat{X}_v=Q_v$, in particular at distance at most $2$. Hence $[Q_v]_k\subset\mathfrak{h}_v$.
\item The inclusion $C\cup\mathfrak{h}_v\subset Y_4(\mathfrak{p})$ is verified for $v$ in the
lemma of this section that checks the domain of the consumer of an inner-collar core against the
filing domains of its pieces, with the data of Definition~\ref{10.3:def-hull-halo-label}.
\end{itemize}

\emph{Exposure.} The correlation theorem states its criterion for exposure in terms of size. The
size attributed to a pointed unit is larger than the unit's own supremum $\mathfrak{n}_v$ on its
own family. That supremum is all that the first of the weight conditions of the correlation
theorem gives at the unit's own family. On $C$ the surplus is supplied by four contributions stated
in the correlation theorem: the current estimate, the cube estimate, the regularity estimate, and,
as the fourth, the pair formed by the marginal and the interior estimates. Off $C$ nothing in that
construction supplies it.

For $v$, the size that the first line of the table of constants of the correlation theorem
attributes to it is $t^{4+2\hat{\alpha}}$, and this exceeds its own supremum. Its
read box $Q_v$ leaves $C$. By the stated criterion, $v$ therefore needs the exterior surplus, that
is, $v$ is exposed in the sense of the correlation theorem.

The enumeration of exposed units in the correlation theorem does not contain $v$ for one reason
only. Under that theorem's rule for interior cores one has $Q_v\subset Z$, and that rule is the one
replaced for $v$.

\emph{The value of $v$.} The statement of the correlation theorem that defines the values of
pointed units covers smooth far cores of two kinds: interior cores, and the cores of a plain
second-part unit. It contains $v$ through the first of these. Its content is the following.
Recall that $p_v$ is the level of the box, and $\mathcal{U}(m)$ is the $U$-slot of a point $m$ of
the source polydisc $\mathfrak{P}$. The function $f_v(m)$ means the value
$\mathfrak{i}^{(j)}_y$, the indices $j$ and $y$ being those of the core $v$, at the extension
reading
\begin{equation*}
\operatorname{ext}_v(m):=(U_p,J_p)\bigl(Q_v,\,M^p(\mathcal{U}(m)|_{Q_v})\bigr),\qquad p:=p_v,
\end{equation*}
where $(U_p,J_p)$ and $M^p$ are the level-$p$ maps of that statement and $\mathcal{U}(m)|_{Q_v}$ is
the restriction of $\mathcal{U}(m)$ to $Q_v$; here the index $p$ is a level, so that $M^p$ is not a
power of the constant $M=L^{m'}$, and $U_p$ is not the bond variable $U$. Of the data of $v$ this
statement therefore reads the
pair $(Q_v,p_v)$ and nothing else. The only thing that changes
for $v$ is the rule fixing $p_v$ for interior cores.

\emph{Conclusion.} The one declaration of Definition~\ref{10.3:def-hull-halo-label} fixes the
quadruple consisting of $p_v$, $Q_v$, the exposure of $v$ and its filing. The criteria stated in
the correlation theorem then place $v$ in each of its lists:
\begin{itemize}
\item the list of exposed units in the exterior current estimate and in the statement of the
correlation theorem;
\item the clause stated for far cores of a plain second-part unit, in the third part of the
enlarged-cube estimate of the correlation theorem;
\item the clause stated for smooth cores of such a unit in the interior estimate;
\item the lists of the label rule, of the filing rule, of the supply rule and of the exterior
collar estimate of the correlation theorem;
\item the corresponding part of the table of constants;
\item the anchor entry of the plaquette-interior estimate.
\end{itemize}
These criteria are stated for exposed units, for far cores or for smooth cores, and $v$ meets each
of them through the four data so fixed. Hence the extension of every list follows from the one
declaration of Definition~\ref{10.3:def-hull-halo-label}.
\end{proof}

\begin{lemma}[Admissible cubes around the read box]\label{10.3:lem-certified-cube}
Let $v$ be an inner-collar core with point $y$. Write $X_v=\square_0$ for its box of $(2\nu+3)$
level-$j$ blocks, which lies in the deep core set $\mathsf{C}$, and $K_v=[\square_0]_k$. Write
$Q_v=Q^{(p^*)}(y)$ for its read box, which is not contained in $\mathsf{C}$, and $\hat{X}_v=Q_v$
for its hull. Its halo is
$\mathfrak{h}_v$, its support is $\mathfrak{r}_v=\bigcup\operatorname{St}(\Delta_y)$, and
$X^{\oplus}_v=\mathsf{C}\cup\mathfrak{h}_v$. These objects are as in
Definition~\ref{10.3:def-hull-halo-label} and Lemma~\ref{10.3:lem-core-radius}. Assume the
following two statements of the correlation theorem, used here by statement.
\begin{enumerate}
\item[(a)] \emph{The admissible-cube statement of the correlation theorem, read at the zone index
$n$ of $y$.}

Let $S\subset\mathfrak{r}_v$ be a box with $y\in S$ whose sides are at most one cube of
$\pi_{n-1}$. Then $S$ lies in an admissible cube of index $n-1$ or $n$, and the cells of that
index of this cube are at $\operatorname{dist}_n<1$ from $\Delta_y$. Hence this cube lies inside
$\mathfrak{r}_v$, which in turn satisfies
\begin{equation*}
\mathfrak{r}_v\subset\mathfrak{h}_v\subset X^{\oplus}_v ,
\end{equation*}
the first inclusion in the sense $[\mathfrak{r}_v]_n\subset\mathfrak{h}_v$. The statement is
read at a point $y$ in the position $\operatorname{dist}_n(y,\Omega_n^c)\ge 71>1$, that is, at
the boundary of $\mathsf{C}$.

The statement closes with the following sentence. For a far core $v'$ in the sense of
clause~(d) of the correlation theorem, with read-box data $(a_{v'},p_{v'})$, the read box
$Q_{v'}=Q^{(p_{v'})}\subset\mathfrak{r}_{v'}$ is itself a box of the kind just described: by the
read-box side bound of the correlation theorem its sides are at most $L^{a_{v'}-1}$ cubes of
$\pi_{p_{v'}}$, hence at most one cube of $\pi_{n-1}$, and it lies within one level-$n$ unit of
$y$, hence in $\Omega_{n-1}$ by the global inclusion of the correlation theorem that places the
points within one level-$n$ unit of $y$ in $\Omega_{n-1}$.
\item[(b)] \emph{The star statement.} Let $n'$ be a level and $\Delta'\in\pi_{n'}$. Every box
$V\subset Q_5(\Delta')\cap\Omega_{n'-1}$ that meets $\Delta'$ lies in
$\bigcup\operatorname{St}(\Delta')$. The last clause of this statement reads: for a hull
$\hat{X}$ containing a point of $\Delta'$,
\begin{equation*}
\Bigl[\bigcup\operatorname{St}(\Delta')\Bigr]_{n'}
\subset\{\Delta\in\pi_{n'}:\operatorname{dist}_{n'}(\Delta,\hat{X})\le 2\},
\end{equation*}
the right-hand side being the halo of $\hat{X}$.
\end{enumerate}
Then the following hold.
\begin{enumerate}
\item[(1)] $Q_v$ lies in an admissible cube of index $k-1$ or $k$ whose cells of that index are
at $\operatorname{dist}_k<1$ from $\Delta_y$, and this cube lies in $\mathfrak{r}_v$. Moreover,
$Q_v\subset\mathfrak{r}_v$ and $[\mathfrak{r}_v]_k\subset\mathfrak{h}_v\subset X^{\oplus}_v$.
\item[(2)] The zone index of $y$ is $n=k$. Furthermore $\operatorname{dist}_k(y,\Omega_k^c)\ge 71$
and $\mathfrak{r}_v\subset\Omega_k\cup(\mathsf{C}\cap\Omega_k)$.
\end{enumerate}
\end{lemma}

\begin{proof}
\emph{Zone index of $y$.} By \eqref{10.3:eq-core-geometry-c} of
Lemma~\ref{10.3:lem-core-geometry}, $\operatorname{dist}_k(y,\Omega_k^c)\ge 71$, so $y$ lies in
$\Omega_k$. The zone index of $y$ is $n=k$; this is obtained from $Q_v\not\subset\mathsf{C}$ by
the same geometry of the core, \eqref{10.3:eq-core-geometry-c}. This is the first
assertion of~(2).

We check the hypotheses of statement~(a) for the box $S:=Q_v$ at the zone index $n=k$.

\emph{$S$ contains $y$.} The read box $Q_v=Q^{(p^*)}(y)$ is the hull $\hat{X}_v$, and it
contains $y$.

\emph{The sides of $S$ are at most one cube of $\pi_{k-1}$.} This is
\eqref{10.3:eq-core-geometry-a} of Lemma~\ref{10.3:lem-core-geometry}.

\emph{$S\subset\mathfrak{r}_v$.} We apply statement~(b) with $n'=k$ and $\Delta'=\Delta_y$.
Its hypotheses hold for the box $V:=Q_v$, for three reasons:
\begin{itemize}
\item $Q_v$ is a box;
\item $Q_v\subset Q_5(\Delta_y)\cap\Omega_{k-1}$, by \eqref{10.3:eq-core-geometry-d} of
Lemma~\ref{10.3:lem-core-geometry};
\item $Q_v$ meets $\Delta_y$, since $y\in Q_v\cap\Delta_y$.
\end{itemize}
Hence $Q_v\subset\bigcup\operatorname{St}(\Delta_y)=\mathfrak{r}_v$. The hull
$\hat{X}_v=Q_v$ contains $y$, so the last clause of statement~(b) gives
$[\mathfrak{r}_v]_k\subset\mathfrak{h}_v$. Finally, $\mathfrak{h}_v\subset X^{\oplus}_v$
because $X^{\oplus}_v=\mathsf{C}\cup\mathfrak{h}_v$.

\emph{Position of $y$.} Statement~(a) is used at the boundary of $\mathsf{C}$, for a point $y$
with $\operatorname{dist}_k(y,\Omega_k^c)\ge 71$. This holds, since
$\operatorname{dist}_k(y,\Omega_k^c)\ge 71>1$ by \eqref{10.3:eq-core-geometry-c} of
Lemma~\ref{10.3:lem-core-geometry}.

All hypotheses of statement~(a) on the box therefore hold for $S=Q_v$. Statement~(a) then gives
the admissible cube in~(1), together with the fact that this cube lies in $\mathfrak{r}_v$, and
the chain of inclusions in~(1) was obtained above.

We now turn to the closing sentence of statement~(a). It is stated for far cores of clause~(d)
of the correlation theorem. Its geometry uses only the data $(j,n,a_{v'})$ of the far core $v'$;
for $v$ these data are $(j,k,a^*)$, and $p^*$ takes the place of $p_{v'}$. Its requirements hold
for $v$ as follows:
\begin{itemize}
\item the side bound is \eqref{10.3:eq-core-geometry-a};
\item $Q_v\subset\Omega_{k-1}$ is \eqref{10.3:eq-core-geometry-d};
\item $Q_v\subset\mathfrak{r}_v$ was shown above.
\end{itemize}
Hence the sentence holds with the words ``far core of clause~(d)'' replaced by ``inner-collar
core'', and $Q_v$ is again a box of the kind required in statement~(a).

For~(2), the hypothesis of the first case of the correlation theorem's geometry statement for its
units is that $X_v$ meets $\mathsf{C}$; it holds for $v$, since $X_v=\square_0\subset\mathsf{C}$.
The argument given there for that case proceeds from $K_v\not\subset Z$, which does
not hold for $v$. The conclusion of that case is instead obtained from $Q_v\not\subset\mathsf{C}$,
which holds for $v$ as recorded in the setting of the lemma, together with the geometry of the
core, \eqref{10.3:eq-core-geometry-c}. Since $X_v\subset\mathsf{C}$ while
$Q_v\not\subset\mathsf{C}$, the read box of $v$ crosses the boundary of $\mathsf{C}$, which is
the situation of the first case, and there:
\begin{itemize}
\item $n=k$ was obtained at the beginning of the proof from $Q_v\not\subset\mathsf{C}$ and
\eqref{10.3:eq-core-geometry-c};
\item $\operatorname{dist}_k(y,\Omega_k^c)\ge 71$ is \eqref{10.3:eq-core-geometry-c};
\item the inclusion $\mathfrak{r}_v\subset\Omega_k\cup(\mathsf{C}\cap\Omega_k)$ is read off
the same display \eqref{10.3:eq-core-geometry-c}, at the point $y$ whose read box satisfies
$Q_v\not\subset\mathsf{C}$.
\end{itemize}
\end{proof}

\begin{lemma}[The consumer's domain lies in every piece's domain]\label{10.3:lem-consumer-domain}
Let $v$ be an inner-collar core, read at level $k$ and exposed in the sense of
Definition~\ref{10.3:def-hull-halo-label}. Write $X_v=\square_0$ for its box, $K_v=[\square_0]_k$,
$\hat{X}_v=Q_v$ for its hull, $\mathfrak{h}_v$ for its halo and
$\operatorname{lab}_0:=K_v\cup\mathfrak{h}_v$ for its start label, all as in that definition. Write
$\mathfrak{r}_v=\bigcup\operatorname{St}(\Delta_y)$ for its support, the union of the stars of the
cells $\Delta_y$ listed in Lemma~\ref{10.3:lem-core-radius}, and $\operatorname{St}_v$ for the set of
cells of these stars. Let $C$ be the component of the large-field region $Z$ that contains
$\square_0$, and put $\mathcal{R}_0:=\hat{X}_v\cup X_v=Q_v$. For a union $A$ of $\pi_k$-cubes put
\begin{equation*}
Y(A):=A\cup\bigcup\{C' : C'\text{ a component of }Z,\ C'\cap A\neq\emptyset\}.
\end{equation*}
Put $X^{\oplus}_v:=Y(\operatorname{lab}_0)$, and let $\mathfrak{E}^{\oplus}_v:=\mathbb{C}^{(N)}_k(X^{\oplus}_v)$
be the consumer at $X^{\oplus}_v$. Here and below $\mathcal{C}$ denotes the class of admissible
consumers of the correlation theorem, each $\mathfrak{E}\in\mathcal{C}$ carrying a domain $X$ and an
analyticity domain $\mathfrak{D}_{\mathfrak{E}}$ for the carrier pair.
By \eqref{10.3:eq-hull-halo-label-a}, $X^{\oplus}_v=C\cup\mathfrak{h}_v$.

Labels are defined as follows. Consider a labelled triple $F=(A_F,\mathcal{R}_F,\|F\|)$, with label
$A_F$ and reading set $\mathcal{R}_F$, and an admissible family $\mathsf{y}\in\mathfrak{Y}(F)$ with
components $\mathsf{y}_a$; the third entry $\|F\|$ plays no role in what follows and is suppressed
below. Put
\begin{equation}\label{10.3:eq-consumer-domain-1}
\operatorname{lab}(F,\mathsf{y}):=A_F\cup[\mathsf{y}]_k\cup\bigcup_a\bigl(\operatorname{Cor}_a\cup
\operatorname{Hub}_a\bigr).
\end{equation}
Here $\operatorname{Cor}_a$ is the set of intermediate cubes of a chain of cubes from the parent cube
$\operatorname{pa}(\Delta_a)$ of the source cell $\Delta_a$ of $\mathsf{y}_a$ to a nearest cube of
$A_F$, and $\operatorname{Hub}_a$ is the set of
hub corridors. Now let $\mathfrak{p}=(v;\mathsf{y}_1,\dots,\mathsf{y}_4)$ be a piece of $v$; empty
families are allowed. For $t=1,\dots,4$ put
$\operatorname{lab}_t:=\operatorname{lab}((\operatorname{lab}_{t-1},\mathcal{R}_{t-1}),\mathsf{y}_t)$
if $\mathsf{y}_t\neq\emptyset$, and $\operatorname{lab}_t:=\operatorname{lab}_{t-1}$ if
$\mathsf{y}_t=\emptyset$. Here $\mathcal{R}_1,\mathcal{R}_2,\mathcal{R}_3$ are the reading sets of the
expansion of pieces in the correlation theorem. Finally put
$Y_4(\mathfrak{p}):=Y(\operatorname{lab}_4)\in\mathbf{D}_k^Z$.
For $t\in\{1,\dots,4\}$ we say that the \emph{reading condition} holds at $t$ if every
$\Delta\in\mathsf{y}_t$ which sees $\mathcal{R}_{t-1}$, in the sense of the reader geometry of the
correlation theorem, satisfies
\begin{equation*}
\operatorname{dist}_k(\operatorname{pa}(\Delta),\operatorname{lab}_{t-1})\le c_{\rho}\rho(\Delta).
\end{equation*}

Assume the following statements of the correlation theorem and of \cite{B-CMP122a}.
\begin{itemize}
\item[(a)] The current bound for exposed units of the correlation theorem. Let $u$ be a unit declared
exposed, with support $\mathfrak{r}_u$ and consumer $\mathfrak{E}^{\oplus}_u$ at $X^{\oplus}_u$, and let
$C_u$ be the component of $Z$ met by $X^{\oplus}_u$. Suppose that
$\mathfrak{E}^{\oplus}_u\in\mathcal{C}$, that $X^{\oplus}_u\subset Y_4(\mathfrak{p})$ and
$\mathfrak{r}_u\subset Y_4(\mathfrak{p})$ for every piece $\mathfrak{p}$ of $u$, and that the cells of
$\mathfrak{r}_u$ on either side of $\partial C_u$ lie in the set $\mathfrak{T}$ of (f), off every live
piece. Then for every $m\in\mathfrak{P}$, every $y'\in\mathfrak{r}_u$ and every
$a\in\{1,3,5,6,7\}$ the current $Q_a(\operatorname{chain}(m))$ of index $a$ satisfies
\begin{equation*}
Q_a(\operatorname{chain}(m))(y')<\theta_a(y')\le 3c_{\alpha}\,w(\Delta_{y'})\,
\mathbf{r}_a(\pi(y')),
\end{equation*}
where $\theta_a(y')$ is the threshold of index $a$ at $y'$, $c_{\alpha}$ a constant of the
correlation theorem and $\mathbf{r}_a(\pi(y'))$ the radius of index $a$ read at the zone $\pi(y')$.
Among its further clauses, on every admissible cube $\square\subset\mathfrak{r}_u$ of the
correlation theorem (the cubes around the read box singled out there), the gauge-group-valued part
of the carrier configuration admits a gauge on $\square$ in
which its field $A$ satisfies
\begin{equation*}
\ell_{\mathsf{m}}|A|<3w\cdot 2BM\alpha_{0,\mathsf{m}},\qquad
\ell_{\mathsf{m}}^2|\nabla A|<3w\cdot 2BM\alpha_{0,\mathsf{m}},
\end{equation*}
with the length unit $\ell_{\mathsf{m}}$, the radius $\alpha_{0,\mathsf{m}}$ and the constant $B$ of
the correlation theorem. The bound is derived
for $u$ as follows. Statement (b) is applied at $\mathfrak{E}^{\oplus}_u$; this requires
$\mathfrak{E}^{\oplus}_u\in\mathcal{C}$ and $X^{\oplus}_u\subset Y_4(\mathfrak{p})$ for every piece
$\mathfrak{p}$ of $u$. The bound is then read in three clauses:
$(\alpha)$ on the cells of $C_u$ inside $X^{\oplus}_u$, from statement (d), with clause index
$\pi(y')$;
$(\beta)$ on $\mathfrak{r}_u\cap\mathfrak{T}$ off every live piece, from the inherited clause at the
shell index $\pi(y')$. For this the source $\zeta$ lies in the small-field class on
$Y_4(\mathfrak{p})$ and has Ba{\l}aban's clauses of level $\pi(y')$ at every cell of
$Y_4(\mathfrak{p})$ of zone $\pi(y')$, which requires $\mathfrak{r}_u\subset Y_4(\mathfrak{p})$;
moreover the operation $\mathbf{R}$ changes nothing on $Z^c$.
$(\gamma)$ on the cells of a live piece of a component of $Z$ other than $C_u$ lying inside
$X^{\oplus}_u$.
\item[(b)] The domain clause of the correlation theorem. For every consumer
$\mathfrak{E}\in\mathcal{C}$ whose domain $X$ satisfies $X\subset Y_4(\mathfrak{p})$, every
$\zeta\in\mathfrak{S}_{\mathrm{src}}$, all parameters listed there (among them $|\sigma|\le
e^{\kappa_1}$) and each of the three carrier positions, the restriction to $X$ of the
gauge-transformed carrier pair lies in $\mathfrak{D}_{\mathfrak{E}}$.
\item[(c)] The inheritance theorem of the correlation theorem. Every $X\in\mathbf{D}_k^Z$ satisfying
the clauses of \cite[(1.63)]{B-CMP122a} indexes a member of $\mathcal{C}$ carrying the stored datum
$\mathfrak{D}^{\flat}_N(X)$. An index meeting $\partial Z$ is admissible in the same way as $C$, the
class being indexed by $X$.
\item[(d)] Every component of $Z$ satisfies the clauses of (c). Every set of
$\mathbf{D}_k^Z$ containing a set that satisfies them satisfies them as well.
\item[(e)] The reading condition at the links $t\ge 2$ for the units of the correlation theorem. For
every unit $u$ of the correlation theorem, with box $X_u$, star cells $\operatorname{St}_u$ and start
label $K_u$, every piece of $u$ and every $t\ge 2$, the correlation theorem establishes the reading
condition at $t$ case by case, as follows.
\begin{itemize}
\item A member of $\mathsf{y}_t$ meeting $X_u$, a cell of $\operatorname{St}_u$ or an earlier
$\mathsf{y}_s$ has its parent in $\operatorname{lab}_{t-1}\supseteq\operatorname{lab}_1$ or sharing
a wall with a cube of it, hence at distance at most $1$.
\item For a source cell at a hub touched or met before, a cube of $\operatorname{lab}_{t-1}$ shares a
wall with that hub or lies in it, at distance at most $h_{\sharp}\le c_{\rho}$.
\item A member touching the rim of $\operatorname{lab}_{t-1}$ (distance at most $1$), of the strip of
a live piece (one unit more), or of a component of $Z$ met by $\operatorname{lab}_{t-1}$, satisfies
\begin{equation*}
\operatorname{dist}_k(\operatorname{pa}(\Delta),\operatorname{lab}_{t-1})\le 100\check{R}_k+1
\le c_{\rho}(9\check{R}_k-2)\le c_{\rho}\rho(\Delta),
\end{equation*}
since the nearest hub is at distance at least $9\check{R}_k$ from any rim of $Z$.
\end{itemize}
\item[(f)] The geometry of the correlation theorem, and the position of the read box:
\begin{itemize}
\item hubs lie at distance at least $9\check{R}_k$ inside $Z$, and the nearest hub is at distance at
least $9\check{R}_k$ from any rim of $Z$;
\item the live pieces $X_{\mathfrak{a}}$ of $C$ lie in $\Omega_k^c$ and satisfy
\begin{equation*}
X_{\mathfrak{a}}\subset Z_h\subset\mathsf{C}\subset Z''_k,\qquad
\mathsf{C}_C\subset Z''_{h+1}\subset Z''_k\cap C,
\end{equation*}
where $Z_h$ is the large-field region at the base level $h$, $Z''_{h+1}$, $Z''_k$ are the nested
large-field sets, $\mathsf{C}$ is the deep core set and $\mathsf{C}_C$ the deep core set of $C$;
\item live pieces of other components lie on those components, at distance at least $38$;
\item $w=1$ on $\mathfrak{T}:=\Omega_k\cup Z^c$;
\item the read box $Q_v$ lies within distance $M/L$ of $\partial C$.
\end{itemize}
\end{itemize}
Then the following hold.
\begin{enumerate}
\item For every piece $\mathfrak{p}$ of $v$, empty families included,
\begin{equation}\label{10.3:eq-consumer-domain-a}
\begin{gathered}
\operatorname{lab}_4\supseteq\operatorname{lab}_3\supseteq\operatorname{lab}_2\supseteq
\operatorname{lab}_1\supseteq\operatorname{lab}_0=K_v\cup\mathfrak{h}_v,\\
X^{\oplus}_v=Y(\operatorname{lab}_0)\subset Y(\operatorname{lab}_4)=Y_4(\mathfrak{p}),
\end{gathered}
\end{equation}
with $Y_4(\mathfrak{p})=X^{\oplus}_v$ for the all-empty piece. Moreover
$\mathfrak{r}_v\subset Y_4(\mathfrak{p})$.
\item The reading condition holds at $t=1,\dots,4$.
\item $\mathfrak{E}^{\oplus}_v\in\mathcal{C}$, and $v$ meets the requirements of (a); hence the
current bound of (a) holds for $v$, and clause $(\gamma)$ is void.
\end{enumerate}
\end{lemma}

\begin{proof}
Let $\mathfrak{p}=(v;\mathsf{y}_1,\dots,\mathsf{y}_4)$ be an arbitrary piece of $v$. We check the
requirements of (a) in the order in which its derivation uses them. We then prove
\eqref{10.3:eq-consumer-domain-a} and the reading condition.

\emph{Exposure and start label.} That $v$ is exposed, and that its start label is
$K_v\cup\mathfrak{h}_v$ rather than $K_v$, is fixed in Definition~\ref{10.3:def-hull-halo-label}.

\emph{The consumer lies in $\mathcal{C}$.} By \eqref{10.3:eq-hull-halo-label-a},
$X^{\oplus}_v=C\cup\mathfrak{h}_v$. As a set, $\operatorname{lab}_0=\mathfrak{h}_v$, since
$K_v\subset[Q_v]_k\subset\mathfrak{h}_v$ by Definition~\ref{10.3:def-hull-halo-label}; and
$\mathfrak{h}_v$ is connected
and meets $C$. Hence $X^{\oplus}_v$ is a connected member of $\mathbf{D}_k^Z$ containing the component
$C$. By (d), $C$ satisfies the clauses of (c), and so does its superset
$X^{\oplus}_v$. By (c), $X^{\oplus}_v$ indexes a member of $\mathcal{C}$, and the fact that it meets
$\partial Z$ does not affect admissibility. Thus $\mathfrak{E}^{\oplus}_v\in\mathcal{C}$. The domain
$X^{\oplus}_v$ has the same shape as the domains of the exposed units of the inner-collar class: a
component of $Z$ together with a connected set of cubes meeting it.

\emph{The support lies in every $Y_4(\mathfrak{p})$.} Consider
\begin{equation*}
\mathfrak{r}_v\subset\bigcup[\mathfrak{r}_v]_k\subset\bigcup\mathfrak{h}_v
\subset\bigcup\operatorname{lab}_0\subset Y(\operatorname{lab}_0)\subset Y_4(\mathfrak{p}).
\end{equation*}
The inclusions hold for the following reasons.
\begin{itemize}
\item The first holds because the partitions nest, so each cell of $\mathfrak{r}_v$ lies in a
$\pi_k$-cube meeting $\mathfrak{r}_v$ in interior.
\item The second is the inclusion $[\mathfrak{r}_v]_k\subset\mathfrak{h}_v$ of the support in the halo,
by Lemma~\ref{10.3:lem-core-radius}.
\item The third is $\mathfrak{h}_v\subset\operatorname{lab}_0$.
\item The fourth is $A\subset Y(A)$.
\item The last is \eqref{10.3:eq-consumer-domain-a}, proved below.
\end{itemize}

\emph{The cells of clause $(\beta)$.} The cells of $\mathfrak{r}_v$ on either side of $\partial C$ lie
in $\Omega_k$, by the bound $\operatorname{dist}_k(y,\Omega_k^c)\ge 71$ of
\eqref{10.3:eq-core-geometry-c}, for the points $y$ of the core to which that display refers. Hence
they lie in $\Omega_k\subset\mathfrak{T}$. They are off every
live piece. The live pieces of $C$ lie in $\Omega_k^c$ by (f), and those of other components lie on
those components, at distance at least $38$. By (f), $w=1$ on these cells.

\emph{Clause $(\gamma)$ is void.} No live piece of another component meets
$C\cup\mathfrak{h}_v=X^{\oplus}_v$, since such pieces lie on other components, at distance at least
$38$.

\emph{Clause $(\alpha)$.} On the cells of $C$ inside $X^{\oplus}_v$, the bound is read from (d) at
the clause index $\pi(y')$, as in the derivation (a).

\emph{The label recursion.} In \eqref{10.3:eq-consumer-domain-1}, $A_F$ is one member of the union.
Hence $\operatorname{lab}(F,\mathsf{y})\supseteq A_F$ for every triple $F$ and every
$\mathsf{y}\in\mathfrak{Y}(F)$: cells are only added, whatever $\mathcal{R}_F$ is. Apply this with
$F=(\operatorname{lab}_{t-1},\mathcal{R}_{t-1})$. This gives
$\operatorname{lab}_t\supseteq\operatorname{lab}_{t-1}$ when $\mathsf{y}_t\neq\emptyset$, and
$\operatorname{lab}_t=\operatorname{lab}_{t-1}$ holds by definition when $\mathsf{y}_t=\emptyset$.
Induction over $t=1,\dots,4$ gives the first line of \eqref{10.3:eq-consumer-domain-a}.

\emph{Monotonicity of $Y$.} Let $A\subset B$. Every component of $Z$ meeting $A$ meets $B$, so
$Y(A)\subset Y(B)$. With $A=\operatorname{lab}_0$ and $B=\operatorname{lab}_4$ this gives the second
line of \eqref{10.3:eq-consumer-domain-a}. For the all-empty piece,
$\operatorname{lab}_4=\operatorname{lab}_0$, so $Y_4(\mathfrak{p})=X^{\oplus}_v$. Thus (b) is applied
at $\mathfrak{E}^{\oplus}_v$ within its hypothesis $X\subset Y_4(\mathfrak{p})$ for the all-empty
piece too; this piece is the value $\operatorname{val}_v$. The filing clause of the correlation
theorem files that value at the domain
$Y_1(v):=Y(\mathfrak{K}_v)$ with $\mathfrak{K}_v:=K_v\cup\mathfrak{h}_v=\operatorname{lab}_0$, so
$Y_1(v)=X^{\oplus}_v$.

\emph{The reading condition.} For a triple $F$ the condition reads: every $\Delta\in\pi^{\mathfrak{s}}$
which sees $\mathcal{R}_F$ has
$\operatorname{dist}_k(\operatorname{pa}(\Delta),A_F)\le c_{\rho}\rho(\Delta)$. It is not needed for
the monotonicity step. It is the hypothesis under which the bounds of the correlation theorem on
labels and on sums over families apply.

At $t=1$ we have $\mathcal{R}_0=Q_v$. This set meets no hub: hubs lie at distance at least
$9\check{R}_k$ inside $Z$ by (f), while $Q_v$ lies within $M/L$ of $\partial C$, also by (f). So a
member $\Delta$ of $\mathsf{y}_1$ sees $\mathcal{R}_0$ only by meeting $Q_v$. Then
\begin{equation*}
\operatorname{pa}(\Delta)\in[Q_v]_k\subset\mathfrak{h}_v=\operatorname{lab}_0,
\end{equation*}
and the distance is $0$.

At $t\ge 2$ we read the cases of (e) for $v$, with $Q_v$ in place of $X_u$ and using
$\operatorname{lab}_0\subset\operatorname{lab}_{t-1}$ from the label recursion. Its cases are these.
\begin{itemize}
\item A member meeting $Q_v$, a cell of $\operatorname{St}_v$ or an earlier $\mathsf{y}_s$ has its
parent in $\operatorname{lab}_{t-1}\supseteq\operatorname{lab}_1$ or sharing a wall with a cube of
it, hence at distance at most $1$.
\item For a source cell at a hub touched or met before, a cube of $\operatorname{lab}_{t-1}$ shares a
wall with that hub or lies in it, at distance at most $h_{\sharp}\le c_{\rho}$.
\item A member touching the rim of $\operatorname{lab}_{t-1}$ (distance at most $1$), of the strip of a
live piece (one unit more), or of a component of $Z$ met by $\operatorname{lab}_{t-1}$, satisfies
\begin{equation*}
\operatorname{dist}_k(\operatorname{pa}(\Delta),\operatorname{lab}_{t-1})\le 100\check{R}_k+1
\le c_{\rho}(9\check{R}_k-2)\le c_{\rho}\rho(\Delta),
\end{equation*}
since by (f) the nearest hub is at distance at least $9\check{R}_k$ from any rim of $Z$.
\end{itemize}

The admissible families $\mathfrak{Y}(F)$ are defined through $\mathcal{R}_F$ and the site geometry
only. Replacing $K_v$ by $K_v\cup\mathfrak{h}_v$ as start label therefore leaves the families of $v$
unchanged; it enlarges only the labels.

\emph{Conclusion.} Every requirement of (a) is met for $v$: exposure,
$\mathfrak{E}^{\oplus}_v\in\mathcal{C}$, $X^{\oplus}_v\subset Y_4(\mathfrak{p})$ for every piece
(the all-empty one included), $\mathfrak{r}_v\subset Y_4(\mathfrak{p})$, and the cells of clause
$(\beta)$. Clause $(\gamma)$ is void. The current bound of (a) therefore holds for $v$.
\end{proof}

\begin{lemma}[Regularity of the chain on the read box]\label{10.3:lem-box-regularity}
Let $v$ be an exposed unit of the inner-collar class $\mathfrak{L}^{c}$ with core at the point $y$. Let
$n$ be the zone index of $y$. Let $k=h+N$, with $h$ the base level and $N$ the number of zones, and let
$\Delta_y$ be the cell of $y$ in the partition $\pi_k$ into level-$k$ cubes. Let
$\bar{\mathsf{w}}_v:=3w^{\vee}(\Delta_y)$ be the weight of $v$, where $w^{\vee}$ is the cell weight of
the correlation theorem. Let
$S:=Q_v=Q^{(p^*)}(y)$ be the read box of $v$, let $\mathfrak{r}_v$ be the support of $v$, and put
$\ell:=\ell_{n-1}$. Assume the following.
\begin{enumerate}
\item[(a)] $S$ is a box contained in $\mathfrak{r}_v$, of the kind admitted in the geometric
description of the support of the consumer. The extended curvature bound and the extended cube-gauge
statement of the correlation theorem hold on $\mathfrak{r}_v$. Here $u$ denotes the cube gauge of
the extended curvature bound on the admissible cube of Lemma~\ref{10.3:lem-certified-cube}.
\item[(b)] Part (i) of the regularity lemma of the correlation theorem holds, and so does the
extended regularity statement of that theorem. The latter says the following. Replace the curvature
bound and the cube-gauge statement by their extended forms. Then the proof of part (i) of the
regularity lemma applies word for word on a box $S\subset\mathfrak{r}_v$ as in \textup{(a)}, at the
zone index $n$ of $y$. Its conclusion is that, for every point $m=(\zeta,\varpi,\vec{\mathsf{s}})$ of
the source polydisc $\mathfrak{P}$, the restriction $\operatorname{chain}(m)|_S$ has the form
$(e^{i\eta\mathcal{A}^0},\mathbf{J}^0)$, where $\eta$ is the parameter of that name in the gauge
reading of the correlation theorem (distinct from the vector fields denoted $\eta$ earlier in this
chapter), and
\begin{equation}\label{10.3:eq-box-regularity-1}
\ell\,|\mathcal{A}^0|,\quad \ell^2\,|\nabla\mathcal{A}^0|,\quad \ell^3\,|\mathbf{J}^0|
\;\le\;\mathfrak{a}'_v:=c_{\mathrm{reg}}\,\bar{\mathsf{w}}_v\,\alpha_{*,n}.
\end{equation}
\item[(c)] The weight statement of the correlation theorem holds: the cell weight function $w$ of that
theorem satisfies $w(\Delta)=1$ on every cell $\Delta$ of $\Omega_k\cup Z^{c}$ that lies off the live
pieces, and the weight of every exposed unit of
$\mathfrak{L}^{c}$ is $3$. Further,
$\Delta_y\subset C\cap\Omega_k$, where $C\cap\Omega_k$ is the collar, of old level $k$, of the set
$C$ that carries the core of $v$. Finally $\operatorname{dist}_k(y,\Omega_k^{c})\ge 71$, and the
cells at $\operatorname{dist}_k\le 1$ from $\Delta_y$ lie off every live piece.
\item[(d)] The regularity ceiling on the coupling imposed in the correlation theorem holds. It is a
supremum over all $x\ge\gamma^{-2}$, stated with the weight $\bar{\mathsf{w}}_v=3$; in particular, for
every $x\ge\gamma^{-2}$,
\begin{equation*}
60\,d\,L\,M\cdot 3\,c_{\mathrm{reg}}\,\alpha_*(x)\le 1 .
\end{equation*}
Moreover $\alpha_{*,n}=\alpha_*(x)$ for a value $x\ge\gamma^{-2}$ of the inverse squared coupling,
namely the value at which the radius of zone $n$ is read.
\item[(e)] The source polydisc $\mathfrak{P}$ carries no disc for the transport parameter
$\mathsf{t}$.
\item[(f)] The statement of the correlation theorem that the $J$-slot is not read holds; thus $v$ is a
first-slot reader, that is, a unit that reads only the first slot of the carrier, such as a smooth
core read at $\operatorname{ext}_v$.
\end{enumerate}
Then $\bar{\mathsf{w}}_v=3$. For every $m\in\mathfrak{P}$, in the cube gauge $u$ restricted to $S$,
one has
\begin{equation}\label{10.3:eq-box-regularity-2}
\operatorname{chain}(m)|_S=(e^{i\eta\mathcal{A}^0},\mathbf{J}^0).
\end{equation}
The fields in \eqref{10.3:eq-box-regularity-2} satisfy
\begin{equation}\label{10.3:eq-box-regularity-3}
\ell\,|\mathcal{A}^0|,\quad \ell^2\,|\nabla\mathcal{A}^0|,\quad \ell^3\,|\mathbf{J}^0|
\;\le\;\mathfrak{a}'_v=3\,c_{\mathrm{reg}}\,\alpha_{*,n},
\end{equation}
and
\begin{equation}\label{10.3:eq-box-regularity-4}
60\,d\,L\,M\,\mathfrak{a}'_v\le 1 .
\end{equation}
In particular $\ell|\mathcal{A}^0|\le\mathfrak{a}'_v$ and $\ell^2|\nabla\mathcal{A}^0|\le\mathfrak{a}'_v$
on $\mathfrak{P}$ at $S=Q_v$.
\end{lemma}

In the application, hypothesis \textup{(a)} is supplied by Lemma~\ref{10.3:lem-certified-cube} and
Lemma~\ref{10.3:lem-consumer-domain}.

\begin{proof}
\emph{The weight of $v$.}
The value $\bar{\mathsf{w}}_v=3$ is part of \textup{(c)}; we check it directly from the definition.
By hypothesis \textup{(c)}, $\Delta_y\subset C\cap\Omega_k$ and
$\operatorname{dist}_k(y,\Omega_k^{c})\ge 71$. Every cell at $\operatorname{dist}_k\le 1$ from
$\Delta_y$ lies at level-$k$ distance from $y\in\Delta_y$ at most one plus the diameter of a unit
cube, which is less than $71\le\operatorname{dist}_k(y,\Omega_k^{c})$; hence it lies in
$\Omega_k$. By \textup{(c)} these cells also lie off every live piece. The weight statement in
\textup{(c)} gives $w=1$ on each cell of $\Omega_k\cup Z^{c}$ off the live pieces, so $w=1$ on every
cell at $\operatorname{dist}_k\le 1$ from $\Delta_y$.

By the first and third clauses of the statement of the correlation theorem that fixes
$\bar{\mathsf{w}}_v$, the cell weight
$w^{\vee}(\Delta_y)$ is determined by the values of $w$ on the cells at $\operatorname{dist}_k\le 1$
from $\Delta_y$, and equals $1$ when all of these values are $1$. Hence $w^{\vee}(\Delta_y)=1$ and
\begin{equation*}
\bar{\mathsf{w}}_v=3\,w^{\vee}(\Delta_y)=3 .
\end{equation*}
This is the value $\bar{\mathsf{w}}_v=3$ that the weight statement assigns to every exposed unit of
$\mathfrak{L}^{c}$.

\emph{The hypotheses of the extended regularity statement.}
That statement takes four inputs: the box $S$, the zone index $n$, the weight $\bar{\mathsf{w}}_v$,
and the clauses of the carrier. It uses no datum of the class of the unit, so it applies to $v$ once
these inputs are supplied. They are supplied as follows.
\begin{itemize}
\item $S=Q_v$ is a box contained in $\mathfrak{r}_v$ of the admitted kind, by \textup{(a)}.
\item The zone index is the zone index $n$ of $y$, by definition.
\item $\bar{\mathsf{w}}_v=3$, by the first paragraph.
\item The clauses of the carrier: the extended curvature bound and the extended cube-gauge statement
hold on $\mathfrak{r}_v$, by \textup{(a)}; and the clause of the undilated carrier which asks that the
coefficient of a first-slot reader be zero on receding blocks applies to $v$, since $v$ is a
first-slot reader by \textup{(f)}.
\end{itemize}

\emph{The argument on $S$.}
By \textup{(b)}, the proof of part (i) of the regularity lemma now runs word for word on $S$ at the
zone index $n$. That statement specifies four readings under which the proof transfers.
\begin{itemize}
\item The three entries of that proof measured in level-$(n-1)$ units, namely its sixth, seventh and
eighth, are multiplied by $L^2$.
\item The values $\alpha_{*,n-1}$, $\alpha_{*,n}$, $\alpha_{*,n+1}$ of the radii at the three
adjacent levels lie within a factor $2$ of one another. This follows from the comparison of radii at
neighbouring levels in the correlation theorem.
\item The Baker--Campbell--Hausdorff expansion of that proof is applied with $V=1$, the identity.
\item The gauge used is the cube gauge $u$ of the extended curvature bound on the admissible cube of
Lemma~\ref{10.3:lem-certified-cube}, restricted to $S$.
\end{itemize}

With these readings the conclusion of \textup{(b)} gives \eqref{10.3:eq-box-regularity-2}, together
with \eqref{10.3:eq-box-regularity-1} at the weight $\bar{\mathsf{w}}_v$. Since, by \textup{(e)},
$\mathfrak{P}$ carries no disc for $\mathsf{t}$, the carrier is not dilated. The bound therefore holds
at $\mathfrak{a}'_v$ itself, and substituting $\bar{\mathsf{w}}_v=3$ gives
\eqref{10.3:eq-box-regularity-3}.

\emph{The ceiling.}
The regularity ceiling in \textup{(d)} is a supremum over all $x\ge\gamma^{-2}$. By \textup{(d)},
$\alpha_{*,n}=\alpha_*(x)$ for a value $x\ge\gamma^{-2}$; the supremum covers that value of $x$, and
there the display of \textup{(d)} reads
\begin{equation*}
60\,d\,L\,M\cdot 3\,c_{\mathrm{reg}}\,\alpha_{*,n}\le 1 .
\end{equation*}
By \eqref{10.3:eq-box-regularity-3},
$3\,c_{\mathrm{reg}}\,\alpha_{*,n}=\mathfrak{a}'_v$; the ceiling is stated with the weight $3$, which is
the value of $\bar{\mathsf{w}}_v$ found above. Hence the last display is
\eqref{10.3:eq-box-regularity-4}. No property of $v$ other than $\bar{\mathsf{w}}_v=3$ enters this
step.

The final assertion is the case of \eqref{10.3:eq-box-regularity-3} for $\mathcal{A}^0$ and
$\nabla\mathcal{A}^0$, valid for every $m\in\mathfrak{P}$.

Were the carrier at $S$ dilated by the transport parameter, the same argument would give
\eqref{10.3:eq-box-regularity-3} with $\tfrac98\mathfrak{a}'_v$ in place of $\mathfrak{a}'_v$, under a
correspondingly strengthened ceiling; by \textup{(e)} this case does not arise here.
\end{proof}

\begin{lemma}[Holomorphy and size of the dialled core]\label{10.3:lem-core-size-bound}
Use the notation of Notation~\ref{10.3:not-vector-dial-table}, in particular the level $k$, the source polydisc
$\mathfrak{P}$, with points $m=(\zeta,\varpi,\vec{\mathsf{s}})$, and the carrier $m\mapsto\mathcal{U}(m)$ on
$\mathfrak{P}$.
Let $v\in\mathfrak{L}^{V,\mathrm{col}}$ be a unit of the
inner-collar class, with weight $\bar{\mathsf{w}}_v=3$, core point $y$, level $j$, zone $n=k$, dial
$\lambda_{n,k'}$, test field $\mathsf{h}$ and frozen-family norm $\mathsf{u}_j$. Let $\square\in\pi_j$ be the cube
containing $y$, put $\nu:=1+n-j$, and let $\square^{\sim\nu}$ be $\square$ enlarged by $\nu$ layers of cubes of
$\pi_j$. Let $s(\cdot)$ be the integer function of the correlation theorem, and put
\begin{equation*}
s^*:=s(3),\qquad a^*:=2+s^*,\qquad p^*:=n-a^*,\qquad Q_v:=Q^{(p^*)}(y) ,
\end{equation*}
the last being the read box of $y$ at level $p^*$.
For a level $p\ge j$ put $\nu_p:=1+p-j$ and $t_p:=L^{j-p}$, and let $\rho_{p,j}$ be the run ratio of the
correlation theorem: the quotient of its running radius at the coupling $g_p$ by the same radius at $g_j$.
For a background $\mathfrak{b}$ and a charge argument $e$, the $\sharp$-core at $y$ is
\begin{equation*}
\operatorname{Core}^{\sharp}_y(\mathfrak{b};e)=\sum_{X\ni y,\ X\subset\square^{\sim\nu}}
\bigl[F^{\sharp}(X,\mathfrak{b},y;e)-F^{\sharp}(X,1,y;e)\bigr],
\end{equation*}
the sum running over $X\in\mathbf{D}_j$, where $F^{\sharp}$ is the $\sharp$-family of the frozen family
$\hat{\mathfrak{V}}^{\sharp(j)}$ of the vector package at level $j$, frozen at $y$
(Notation~\ref{10.3:not-vector-dial-table}), with norm $\mathfrak{n}=\mathsf{u}_j$; $|e|$ denotes the modulus of
the charge argument. Put
\begin{equation}\label{10.3:eq-core-size-bound-1}
\begin{split}
\mathfrak{N}^{V,\mathrm{run}}_{p^*}(v):=\mathsf{u}_j\,\sup|\mathsf{h}|\,
\bigl[&C_{\sharp}\,\nu_{p^*}^3\,\rho_{p^*,j}^{2}\,(1+\rho_{p^*,j})\\
&+2K(d)(1+8C_{\mathfrak{J}})\bigr]\,t_{p^*}^{4+2\hat{\alpha}} ,
\end{split}
\end{equation}
where $C_{\mathfrak{J}}$ and $K(d)$ are the constants of hypotheses (B) and (C) below. Finally put
\begin{equation}\label{10.3:eq-core-size-bound-2}
\epsilon_{a^*}:=3c_{\epsilon}\,(2a^*+11)\,\bar{\rho}_{a^*}\,L^{-2(a^*-1)} .
\end{equation}
Assume the following.
\begin{enumerate}
\item[(A)] (The core bound of the vector package of Notation~\ref{10.3:not-vector-dial-table} on admissible
backgrounds.) For every level $p$ with $j\le p$
and every background $\mathfrak{b}$ admissible at scale $p$ on a box containing $\square^{\sim\nu_p}$,
\begin{equation*}
\Bigl|\sum_{X\ni y,\ X\subset\square^{\sim\nu_p}}
\bigl[F^{\sharp}(X,\mathfrak{b},y;e)-F^{\sharp}(X,1,y;e)\bigr]\Bigr|
\le C_{\sharp}\,\mathfrak{n}\,|e|\,\nu_p^3\,\rho_{p,j}^2\,(1+\rho_{p,j})\,t_p^{4+2\hat{\alpha}} .
\end{equation*}
\item[(B)] (The single-term bound of the vector package of Notation~\ref{10.3:not-vector-dial-table}.) For
every such $\mathfrak{b}$ and every
$X\in\mathbf{D}_j$ containing $y$,
\begin{equation*}
\bigl|F^{\sharp}(X,\mathfrak{b},y;e)-F^{\sharp}(X,1,y;e)\bigr|
\le 2(1+8C_{\mathfrak{J}})\,\mathfrak{n}\,|e|\,e^{-\kappa d_j(X)} ;
\end{equation*}
this is the norm bound $(1+8C_{\mathfrak{J}})\mathfrak{n}$ of $F^{\sharp}$, at a rate not smaller than
$\kappa$, applied on the two backgrounds $\mathfrak{b}$ and $1$.
\item[(C)] (Tree lengths and the polymer sum of the correlation theorem.) If $X\in\mathbf{D}_j$ contains two
cubes $\Delta,\Delta'$ of $\pi_j$, then $d_j(X)$ is at least the distance of $\Delta$ and $\Delta'$ measured in
units of the spacing of level $j$; and, for the dimensional constant $c_d$ of the tree-length sums of the
correlation theorem,
\begin{equation*}
\sum_{X\in\mathbf{D}_j,\ X\supset\square}e^{-(c_d+1)d_j(X)}\le K(d) .
\end{equation*}
\item[(D)] (Parameter floors.) $\kappa/20\ge c_d+2$, $\kappa\ge 10\log L$, and $\hat{\alpha}<1$.
\item[(E)] (The gauge-smallness formula of the correlation theorem.) Let $a_v$ be an integer and $p:=n-a_v$,
with $j\le p$, and let $Q^{(p)}(y)$ be the read box of $y$ at level $p$. For $m\in\mathfrak{P}$, a gauge of
$\mathcal{U}(m)$ on $Q^{(p)}(y)$ of
size $\ell_{k-1}|\mathcal{A}^0|,\ \ell_{k-1}^2|\nabla\mathcal{A}^0|\le\mathfrak{a}'_v$ yields on $Q^{(p)}(y)$,
in a single gauge of the complexification $G^{\mathbb{C}}$ of the gauge group, the form
$\mathcal{U}(m)=e^{i\eta_{\mathrm{exp}}\mathcal{A}^0}$, where $\eta_{\mathrm{exp}}$ is the prefactor with which
the correlation theorem writes this exponential parametrisation, with
\begin{equation*}
\max\{\ell_p|\mathcal{A}^0|,\ \ell_p^2|\nabla\mathcal{A}^0|\}\le\epsilon\,c''\,\mathfrak{r}_p,\qquad
\epsilon\le c_{\epsilon}\,\bar{\mathsf{w}}_v\,(2a_v+11)\,\bar{\rho}_{a_v}\,L^{-2(a_v-1)} ;
\end{equation*}
here $c''$ is the enlargement factor of hypothesis (F) below, $\mathfrak{r}_p$ is the scale-$p$ bound on gauge
fields in the correlation theorem, and $\ell_p$ is the lattice spacing of level $p$.
\item[(F)] (Admissibility on any box, from the correlation theorem.) For every box $Q$: if $\mathbf{U}$ lies in
the small-field ball of the correlation theorem enlarged by the factor $c''$, then the background
$U_p(Q,M^p(\mathbf{U}))$ which the correlation theorem constructs from $\mathbf{U}$ on $Q$ at scale $p$ is an
admissible background at scale $p$, in the sense of the correlation theorem, on the box $Q^{\sim-2}$ of the
correlation theorem. In the situation of (E) with $\epsilon\le 1$, this applies, for every $\sigma\in\mathbb{C}$
with $|\sigma|<1/\epsilon$, to $e^{i\sigma\eta_{\mathrm{exp}}\mathcal{A}^0}$ on $Q^{(p)}(y)$: the background
$\operatorname{ext}_p(e^{i\sigma\eta_{\mathrm{exp}}\mathcal{A}^0})$ is admissible at scale $p$ on
$(Q^{(p)}(y))^{\sim-2}$, which contains $\square^{\sim\nu}$, and it depends analytically on $\sigma$.
\item[(G)] (Second-order vanishing and the Schwarz device.) In the situation of (E) and (F), the function
\begin{equation*}
\Phi(\sigma):=\operatorname{Core}^{\sharp}_y\bigl(\operatorname{ext}_p(e^{i\sigma\eta_{\mathrm{exp}}
\mathcal{A}^0});e\bigr),\qquad |\sigma|<1/\epsilon ,
\end{equation*}
satisfies $\Phi(0)=0$ and $\Phi'(0)=0$; and if $\epsilon\le 1$ and $\Phi$ is holomorphic on $|\sigma|<1/\epsilon$
with $\Phi(0)=\Phi'(0)=0$ and $|\Phi|\le B$ there, then $|\Phi(1)|\le 2\epsilon^2B$.
\item[(H)] (The integer function $s$ of the correlation theorem.) $\epsilon_{a^*}\le 1$; this holds by the
definition of $s$ at the argument $3$, that is, the
bound of (E) at $\bar{\mathsf{w}}_v=3$ and $a_v=a^*=2+s(3)$ does not exceed $1$.
\item[(I)] (The clause of the correlation theorem on the evaluation of smooth far cores at their extension
level, with its level rule.) That clause
evaluates a smooth far core on the extension, at the level $p_v:=n-2-s^{(9/8)}(\bar{\mathsf{w}}_v)$, of the
configuration on its box, where $s^{(9/8)}(\bar{\mathsf{w}}):=s(\tfrac98\bar{\mathsf{w}})$; under the
supplementary condition on this level rule, the extension level used for $v$ is
$p^*$, and the evaluation of $v$ at level $p^*$ is the one prescribed by that clause. Accordingly, for every
$m\in\mathfrak{P}$,
\begin{equation*}
f_v(m)=\operatorname{Core}^{\sharp}_y\bigl(\operatorname{ext}_{p^*}(\mathcal{U}(m)\restriction Q_v);
\mathsf{h}(y)\bigr),
\end{equation*}
where $\mathcal{U}(m)$ is the configuration component of the carrier; at the point of $\mathfrak{P}$ denoted $1$
in Notation~\ref{10.3:not-vector-dial-table} the configuration is the identity configuration $1$, so that
$f_v(1)=\operatorname{Core}^{\sharp}_y(\operatorname{ext}_{p^*}(1);\mathsf{h}(y))$;
the carrier
$m\mapsto\mathcal{U}(m)$ is holomorphic on
$\mathfrak{P}$; for every $m\in\mathfrak{P}$ the configuration $\mathcal{U}(m)\restriction Q_v$ lies in the
small-field ball of (F) enlarged by the factor $c''$, and the extension map $\operatorname{ext}_{p^*}$ is analytic
on that ball; and the $\sharp$-core depends analytically on its configuration argument.
\end{enumerate}
Assume also $j\le p^*$. Then $m\mapsto\lambda_{n,k'}f_v(m)$ is holomorphic on $\mathfrak{P}$, $f_v(1)=0$, and
for every $m\in\mathfrak{P}$ one has
$|f_v(m)|\le 2\epsilon_{a^*}^2\,\mathfrak{N}^{V,\mathrm{run}}_{p^*}(v)$ and
\begin{equation}\label{10.3:eq-core-size-bound-a}
\bigl|\lambda_{n,k'}f_v(m)-\lambda_{n,k'}f_v(1)\bigr|\le\mathfrak{N}^{\flat}(v)
:=|\lambda_{n,k'}|\cdot 4\,\epsilon_{a^*}^2\,\mathfrak{N}^{V,\mathrm{run}}_{p^*}(v) .
\end{equation}
We put $\|v\|_{\star}:=\mathfrak{N}^{\flat}(v)$.
\end{lemma}

\begin{proof}
We apply hypotheses (A)--(I) at the level $p:=p^*$ and on the box $Q:=Q_v$.

First, $j\le p^*$ is assumed, and $\nu_{p^*}=1+p^*-j=\nu-a^*$, which is at most $\nu$. Second,
$Q_v=Q^{(p^*)}(y)$ holds by the definition of $Q_v$.

Third, the gauge representation. The regularity of the chain on the read box,
Lemma~\ref{10.3:lem-box-regularity}, taken at the set
$S:=Q_v$, gives a gauge representation of $\mathcal{U}(m)$ on $Q_v$ with
\begin{equation*}
\ell_{k-1}|\mathcal{A}^0|\le\mathfrak{a}'_v,\qquad \ell_{k-1}^2|\nabla\mathcal{A}^0|\le\mathfrak{a}'_v ,
\end{equation*}
and this holds at every point $m$ of $\mathfrak{P}$, not only at a distinguished one. Hypothesis (E), taken at
$a_v=a^*$, that is at $p=n-a^*=p^*$, converts this representation into the form
$\mathcal{U}(m)=e^{i\eta_{\mathrm{exp}}\mathcal{A}^0}$ on $Q_v$ with
\begin{equation*}
\max\{\ell_{p^*}|\mathcal{A}^0|,\ \ell_{p^*}^2|\nabla\mathcal{A}^0|\}\le\epsilon\,c''\,\mathfrak{r}_{p^*},\qquad
\epsilon\le c_{\epsilon}\,\bar{\mathsf{w}}_v\,(2a_v+11)\,\bar{\rho}_{a_v}\,L^{-2(a_v-1)} .
\end{equation*}
Substituting
$\bar{\mathsf{w}}_v=3$ and $a_v=a^*$, the right-hand side is $\epsilon_{a^*}$ of
\eqref{10.3:eq-core-size-bound-2}. Hence, at every $m\in\mathfrak{P}$, this form holds with
$\epsilon:=\epsilon_{a^*}$, and $\epsilon_{a^*}\le 1$ by hypothesis (H).

None of the hypotheses used refers to the large-field sets $Z''_{h+1}$, to the core set $\mathsf{C}$, or to the
class of the unit: (A) and (B) are stated for every admissible background, (C) and (D) are statements about
polymers and parameters, and the admissibility statement (F) holds for every box $Q$, with no condition on where
the box lies; in particular it holds at $Q_v$.

Fix $m\in\mathfrak{P}$, write $e:=\mathsf{h}(y)$, and let $\Phi_m$ be the function of (G) at $p=p^*$,
$Q^{(p)}(y)=Q_v$ and $\epsilon=\epsilon_{a^*}$, with the gauge field $\mathcal{A}^0$ of $\mathcal{U}(m)$
supplied by (E). For $|\sigma|<1/\epsilon_{a^*}$ put
$\mathfrak{b}_\sigma:=\operatorname{ext}_{p^*}(e^{i\sigma\eta_{\mathrm{exp}}\mathcal{A}^0})$. By (F),
$\mathfrak{b}_\sigma$ is admissible at scale $p^*\ge j$ on $Q_v^{\sim-2}$, which contains
$\square^{\sim\nu}$, and therefore also $\square^{\sim\nu_{p^*}}$, since $\nu_{p^*}\le\nu$; and, the
$\sharp$-core depending analytically on its configuration argument by (I), $\Phi_m$ is
holomorphic on $|\sigma|<1/\epsilon_{a^*}$. We split the sum defining
$\Phi_m(\sigma)=\operatorname{Core}^{\sharp}_y(\mathfrak{b}_\sigma;e)$ into two parts.

The constituents $X\ni y$ with $X\subset\square^{\sim\nu_{p^*}}$. By hypothesis (A) at $p=p^*$ and
$\mathfrak{b}=\mathfrak{b}_\sigma$, which is admissible at scale $p^*$ on a box containing
$\square^{\sim\nu_{p^*}}$, the modulus of their sum is at most
\begin{equation*}
C_{\sharp}\,\mathfrak{n}\,|e|\,\nu_{p^*}^3\,\rho_{p^*,j}^2\,(1+\rho_{p^*,j})\,t_{p^*}^{4+2\hat{\alpha}} ;
\end{equation*}
the ratio appearing here is the run ratio $\rho_{p^*,j}$ at the level $p^*$.

The remaining constituents $X\ni y$ with $X\subset\square^{\sim\nu}$ and $X\not\subset\square^{\sim\nu_{p^*}}$.
Such an $X\in\mathbf{D}_j$ is connected and contains $\square$ and a cube of $\pi_j$ outside
$\square^{\sim\nu_{p^*}}$, whose distance from $\square$ in units of the spacing of level $j$ is at least
$\nu_{p^*}$; by the first part of (C), $d_j(X)\ge\nu_{p^*}$. By (B), each of these terms has modulus at most
$2(1+8C_{\mathfrak{J}})\mathfrak{n}|e|e^{-\kappa d_j(X)}$. By (D),
\begin{equation*}
c_d+1\le\tfrac{\kappa}{20}-1\le\tfrac{\kappa}{20},
\end{equation*}
so
\begin{equation*}
\kappa-c_d-1\ge\tfrac{19}{20}\kappa\ge\tfrac{19}{2}\log L\ge 9\log L\ge 0 .
\end{equation*}
Writing
$e^{-\kappa d_j(X)}=e^{-(\kappa-c_d-1)d_j(X)}e^{-(c_d+1)d_j(X)}$ and using $d_j(X)\ge\nu_{p^*}$ together with
the second part of (C),
\begin{align*}
\sum_{X}\bigl|F^{\sharp}(X,\mathfrak{b}_\sigma,y;e)-F^{\sharp}(X,1,y;e)\bigr|
&\le 2(1+8C_{\mathfrak{J}})\mathfrak{n}|e|\,e^{-(\kappa-c_d-1)\nu_{p^*}}
\sum_{X\in\mathbf{D}_j,\ X\supset\square}e^{-(c_d+1)d_j(X)}\\
&\le 2K(d)(1+8C_{\mathfrak{J}})\mathfrak{n}|e|\,L^{-9\nu_{p^*}}\\
&\le 2K(d)(1+8C_{\mathfrak{J}})\mathfrak{n}|e|\,t_{p^*}^{4+2\hat{\alpha}} ,
\end{align*}
the sum on the left running over the remaining constituents. The last step holds because
$t_{p^*}=L^{-(\nu_{p^*}-1)}$, $9\nu_{p^*}\ge 6(\nu_{p^*}-1)$, $4+2\hat{\alpha}\le 6$ by $\hat{\alpha}<1$,
$\nu_{p^*}\ge 1$ since $j\le p^*$, and $L\ge 1$, so that
\begin{equation*}
L^{-9\nu_{p^*}}\le L^{-6(\nu_{p^*}-1)}\le L^{-(4+2\hat{\alpha})(\nu_{p^*}-1)} .
\end{equation*}

Adding the two parts, for every $\sigma$ with $|\sigma|<1/\epsilon_{a^*}$,
\begin{equation}\label{10.3:eq-core-size-bound-3}
\begin{split}
|\Phi_m(\sigma)|\le B(v):=\mathfrak{n}\,|e|\,
\bigl[&C_{\sharp}\,\nu_{p^*}^3\,\rho_{p^*,j}^2\,(1+\rho_{p^*,j})\\
&+2K(d)(1+8C_{\mathfrak{J}})\bigr]\,t_{p^*}^{4+2\hat{\alpha}} ,
\end{split}
\end{equation}
a bound independent of $m$. By (G), $\Phi_m(0)=\Phi_m'(0)=0$, and the Schwarz device of (G), applied with
$\epsilon=\epsilon_{a^*}\le 1$ and $B=B(v)$, gives $|\Phi_m(1)|\le 2\epsilon_{a^*}^2B(v)$. When
$\epsilon_{a^*}<1$ this is the Schwarz lemma: $\sigma\mapsto\Phi_m(\sigma)/\sigma^2$ is holomorphic on
$|\sigma|<1/\epsilon_{a^*}$, bounded by $B(v)/r^2$ on each circle $|\sigma|=r$ with $1<r<1/\epsilon_{a^*}$, so by
the maximum principle $|\Phi_m(1)|\le B(v)/r^2$, and letting $r\to 1/\epsilon_{a^*}$ gives
$|\Phi_m(1)|\le\epsilon_{a^*}^2B(v)\le 2\epsilon_{a^*}^2B(v)$.

At $\sigma=1$ the configuration $e^{i\eta_{\mathrm{exp}}\mathcal{A}^0}$ is $\mathcal{U}(m)$ on $Q_v$, so by
hypothesis (I), $\Phi_m(1)=f_v(m)$. At $\sigma=0$ it is the identity configuration $1$, so by (I),
\begin{equation*}
f_v(1)=\operatorname{Core}^{\sharp}_y(\operatorname{ext}_{p^*}(1);\mathsf{h}(y))=\Phi_m(0)=0 .
\end{equation*}
By the choice of $F^{\sharp}$, $\mathfrak{n}=\mathsf{u}_j$, the norm of the frozen family per unit of
$\sup|\mathsf{h}|$, and
$|e|=|\mathsf{h}(y)|\le\sup|\mathsf{h}|$; since the right-hand side of \eqref{10.3:eq-core-size-bound-3} is
linear in $|e|$ with a non-negative coefficient, comparison with \eqref{10.3:eq-core-size-bound-1} gives
$B(v)\le\mathfrak{N}^{V,\mathrm{run}}_{p^*}(v)$. Hence, for every $m\in\mathfrak{P}$,
\begin{equation*}
|f_v(m)|\le 2\epsilon_{a^*}^2\,\mathfrak{N}^{V,\mathrm{run}}_{p^*}(v).
\end{equation*}
Since this bound holds both at $m$ and at the point $1$ of $\mathfrak{P}$, the triangle inequality gives
\begin{equation*}
\begin{aligned}
\bigl|\lambda_{n,k'}f_v(m)-\lambda_{n,k'}f_v(1)\bigr|
&\le|\lambda_{n,k'}|\,\bigl(|f_v(m)|+|f_v(1)|\bigr)\\
&\le|\lambda_{n,k'}|\cdot 4\epsilon_{a^*}^2\,\mathfrak{N}^{V,\mathrm{run}}_{p^*}(v)=\mathfrak{N}^{\flat}(v) ,
\end{aligned}
\end{equation*}
which is \eqref{10.3:eq-core-size-bound-a}. Since moreover $f_v(1)=0$, the left member equals
$|\lambda_{n,k'}|\,|f_v(m)|$ and is in fact at most $\tfrac12\mathfrak{N}^{\flat}(v)$.

Holomorphy follows from (I). The map $m\mapsto f_v(m)$ is the composite of three maps: the holomorphic carrier
$m\mapsto\mathcal{U}(m)$ on $\mathfrak{P}$, restricted to $Q_v$, whose values lie for every $m\in\mathfrak{P}$ in
the small-field ball of (F) enlarged by the factor $c''$; the extension map $\operatorname{ext}_{p^*}$, which is
analytic on that ball; and the $\sharp$-core, which depends analytically on its configuration
argument. Multiplication by the constant $\lambda_{n,k'}$ preserves holomorphy. Hence
$m\mapsto\lambda_{n,k'}f_v(m)$ is holomorphic on $\mathfrak{P}$.
\end{proof}

In the correlation theorem, the evaluation of a smooth far core at its extension level bounds the difference
$f_v(m)-f_v(1)$ of the core's values at a point $m\in\mathfrak{P}$ and at the point $1$ by $2\epsilon^2$ times
the corresponding size. The factor $4\epsilon_{a^*}^2$ in \eqref{10.3:eq-core-size-bound-a} is what the triangle
inequality gives from the bound on $|f_v|$ at $m$ and at $1$, without use of $f_v(1)=0$; the factor $4$ is kept
so that $\mathfrak{N}^{\flat}(v)$ has one form at all its uses.

The vector package of Notation~\ref{10.3:not-vector-dial-table} states its bound on the $\sharp$-core in the
weaker form obtained from $B(v)$ by majorising
$\rho_{p^*,j}$ by $\bar{\rho}_{p^*-j}$ and the bracket in \eqref{10.3:eq-core-size-bound-3} by
\begin{equation*}
\bigl(C_{\sharp}+2K(d)(1+8C_{\mathfrak{J}})\bigr)\,\nu_{p^*}^3\,\bar{\rho}^2_{p^*-j}\,(1+\bar{\rho}_{p^*-j}) ;
\end{equation*}
in this lemma the form \eqref{10.3:eq-core-size-bound-3}, with the run ratio, is the one used.

By \eqref{10.3:eq-never-deep-a}, the unit $v$ is never deep. The
comparison of $\mathfrak{N}^{\flat}(v)$ with the entry of the first row of the table of
Notation~\ref{10.3:not-vector-dial-table}, whose ratio is the run ratio $\rho_{n,j}$, is the inequality recorded
beside the size letter in \eqref{10.3:eq-vector-dial-table-b}; it is not part of this lemma, and neither the
statement nor the proof above uses it. It requires a comparison of the run ratios at the levels $p^*$ and $n$.

\begin{lemma}[The anchor cell's star contains the read box]\label{10.3:lem-star-contains-box}
Let $v$ be an inner-collar core at level $n=k$, in the notation of
Lemma~\ref{10.3:lem-core-geometry} and Lemma~\ref{10.3:lem-core-radius}. Its ingredients are:
\begin{itemize}
\item the anchor point $y=y_v$ and the read box $Q_v=Q^{(p^*)}(y)$;
\item the hull $\hat{X}_v=Q_v$;
\item the anchor cell $\Delta_y$, which is the cell of level $k$ containing $y$ of the partition into
cells of Lemma~\ref{10.3:lem-core-geometry}; the cells of this partition have levels in $]h,k]$;
\item the support $\mathfrak{r}_v=\bigcup\operatorname{St}(\Delta_y)$;
\item the halo $\mathfrak{h}_v=\{\Delta\in\pi_k:\operatorname{dist}_k(\Delta,\hat{X}_v)\le 2\}$.
\end{itemize}
For a set $A$, $[A]_k$ denotes the set of $\pi_k$-cubes meeting $A$ in interior points. For a
cell $\Delta'$ of level $k$ put $Q_5(\Delta')=\{\operatorname{dist}_k(\cdot,\Delta')\le 2\}$.
Assume the following two statements, both from the proof of the correlation theorem.
\begin{enumerate}
\item[(a)] The following facts about the star $\operatorname{St}(\Delta')$ of a cell $\Delta'$ of
level $n'=k$ of the partition, together with the argument by which (a1) and (a3) are established.
\begin{itemize}
\item[(a0)] The union of the star satisfies
\begin{equation*}
\bigcup\operatorname{St}(\Delta')\subset\{\operatorname{dist}_{n'}(\cdot,\Delta')\le L+2\},
\end{equation*}
with $2$ in place of $L+2$ if no member of $\operatorname{St}(\Delta')$ has level $n'+1$. Since no
cell of the partition has level $k+1$, at $n'=k$ this reads
$\bigcup\operatorname{St}(\Delta')\subset Q_5(\Delta')$.
\item[(a1)] Let $V$ be a box contained in $Q_5(\Delta')\cap\Omega_{k-1}$ and meeting $\Delta'$ in
a set with non-empty interior.
Then every cell of the partition that meets $V$ belongs to $\operatorname{St}(\Delta')$.
Consequently $V\subset\bigcup\operatorname{St}(\Delta')$.
\item[(a2)] The cube $\{\operatorname{dist}_k(\cdot,\Delta')\le 1\}$ lies in $\Omega_{k-1}$.
\item[(a3)] If $y_v\in\hat{X}_v\cap\Delta'$, then $[\operatorname{St}(\Delta')]_k\subset\mathfrak{h}_v$.
\end{itemize}
\item[(b)] The reading statement, in the following form: the cube around the read box of
Lemma~\ref{10.3:lem-certified-cube} and every constituent $X$ of the core (a set contained in
$\square^{\sim\nu}$) lie in $\{\operatorname{dist}_k(\cdot,\Delta_y)\le 1\}$, and
\begin{equation*}
\{\operatorname{dist}_k(\cdot,\Delta_y)\le 1\}\subset\bigcup\operatorname{St}(\Delta_y).
\end{equation*}
\end{enumerate}
Then
\begin{equation}\label{10.3:eq-star-contains-box-1}
\begin{gathered}
\text{(i)}\ \ Q_v\subset\mathfrak{r}_v,\qquad
\text{(ii)}\ \ \{\operatorname{dist}_k(\cdot,\Delta_y)\le 1\}\subset\mathfrak{r}_v,\qquad
\text{(iii)}\ \ [\operatorname{St}(\Delta_y)]_k\subset\mathfrak{h}_v .
\end{gathered}
\end{equation}
In particular, the cube around the read box of Lemma~\ref{10.3:lem-certified-cube} and the constituents
$X\subset\square^{\sim\nu}$ of the core lie in $\mathfrak{r}_v$.
\end{lemma}

\begin{proof}
We apply hypothesis (a) with $\Delta':=\Delta_y$ and $n'=k$. We first check that the data of $v$
have the properties through which the proof of (a) reads $V$ and $y_v$.

Take $V:=Q_v$. By display \eqref{10.3:eq-core-geometry-d} of
Lemma~\ref{10.3:lem-core-geometry}, three things hold. First, $Q_v$ is a box containing $y$.
Second, $Q_v$ meets $\Delta_y$ in a set with interior points. Third,
\begin{equation}\label{10.3:eq-star-contains-box-2}
Q_v\subset\{\operatorname{dist}_k(\cdot,\Delta_y)\le 1\}\subset Q_5(\Delta_y),
\qquad Q_v\subset\Omega_{k-1}.
\end{equation}
Hence $V$ is a box contained in $Q_5(\Delta_y)\cap\Omega_{k-1}$ that meets $\Delta_y$ in a set
with non-empty interior. Moreover
$\hat{X}_v=Q_v$: this is the hull of a smooth core formed from $X_v:=\square_0$. Since
$y_v=y\in Q_v$ and $y\in\Delta_y$, we obtain $y_v\in\hat{X}_v\cap\Delta_y$. Thus $V=Q_v$
satisfies the hypotheses of (a1), and $y_v$ those of (a3).

We now recall how that proof reads $V$. The cells of the partition meeting $V$ have level at
least $k-1$, and they meet $Q_5(\Delta_y)$. The box $V$ is convex. Take a generic interior point
of $V$ lying in such a cell, and join it by a segment to an interior point of $V\cap\Delta_y$.
This segment stays in $V$ and avoids every $(d-2)$-face of the partition. Consecutive cells
along it therefore share a $(d-1)$-face (a wall), so they belong to $\operatorname{St}(\Delta_y)$.

(ii) The cube $\{\operatorname{dist}_k(\cdot,\Delta_y)\le 1\}$ is a box. It is contained in
$Q_5(\Delta_y)$ and contains $\Delta_y$, so it meets $\Delta_y$ in a set with non-empty interior.
By (a2) it lies in $\Omega_{k-1}$. It is therefore an
admissible box $V$ in (a1), and (a1) gives
\begin{equation}\label{10.3:eq-star-contains-box-3}
\{\operatorname{dist}_k(\cdot,\Delta_y)\le 1\}\subset\bigcup\operatorname{St}(\Delta_y)=\mathfrak{r}_v .
\end{equation}
(This inclusion is also the second part of hypothesis (b).)
By hypothesis (b), the cube around the read box of Lemma~\ref{10.3:lem-certified-cube} and the
constituents $X\subset\square^{\sim\nu}$ of the core lie in
$\{\operatorname{dist}_k(\cdot,\Delta_y)\le 1\}$. By \eqref{10.3:eq-star-contains-box-3} they
therefore lie in $\mathfrak{r}_v$.

(i) The first inclusion of \eqref{10.3:eq-star-contains-box-2}, combined with
\eqref{10.3:eq-star-contains-box-3}, gives $Q_v\subset\mathfrak{r}_v$. The same inclusion also
follows from (a1) applied directly to the admissible box $V=Q_v$.

(iii) We have $y_v\in\hat{X}_v\cap\Delta_y$ and $n'=k$, so the last clause (a3) applies. Its
proof runs as follows. Let $x\in\bigcup\operatorname{St}(\Delta_y)$. By (a0), $x$ lies in
$Q_5(\Delta_y)$, hence in a $\pi_k$-cube of $Q_5(\Delta_y)$ that is at $\operatorname{dist}_k\le 1$
from $\Delta_y\ni y_v$. Since $y_v\in\hat{X}_v\cap\Delta_y$ and $\Delta_y$ has $k$-side $1$, this
cube is at $\operatorname{dist}_k\le 1+1=2$ from $\hat{X}_v$, and so
it belongs to $\mathfrak{h}_v$ by the definition of the halo. Hence
$[\operatorname{St}(\Delta_y)]_k\subset\mathfrak{h}_v$.
\end{proof}

\begin{lemma}[Amplitude-one reading astride the boundary]\label{10.3:lem-amplitude-one-reading}
Let $Z$ be the large-field region and let $v$ be an exposed unit of the correlation theorem or a
unit of the inner-collar class $\mathfrak{L}^{c}$. Write $X_v:=\square_0$ for its box, $\hat{X}_v$
for its hull, $Q_v$ for its
read box, $\mathfrak{h}_v$ for its halo, $\mathfrak{r}_v$ for its support region, $K_v$ for its
core cubes, $X^{\oplus}_v$ for the set at which its consumer acts, $f_v$ for its dialled
function and $\operatorname{val}_v:=f_v(\operatorname{chain}(\vec{0}))$ for its value. Let
$\mathfrak{p}=(v;\mathsf{y}_1,\dots,\mathsf{y}_4)$ be a piece with unit $v$.
In what follows $\mathsf{s}$ denotes the link interpolation parameters; $\sigma$ denotes a site
family, and $Y(\sigma)$ the set $\sigma$ together with the components of $Z$ meeting it;
$b$ is the source variable of the link; $\mathcal{Y}(\mathsf{s};b)$, $\mathcal{J}[\mathcal{Y}(\mathsf{s};b)]$
and $\mathsf{u}(\mathsf{s})$ are the three entries of the read triple of the correlation theorem,
formed at the parameters $\mathsf{s}$ and at $b$; and $(\mathbf{U},\mathbf{J})$ is the pair of
fields on which these objects are evaluated.
We say that a functional $F$ \emph{reads at amplitude one on the closed set $\mathcal{R}$}
if $F$ depends on the $\mathsf{s}$-dependent objects only through
\begin{equation*}
\mathcal{A}(\mathsf{s}):=\bigl(\mathcal{Y}(\mathsf{s};b),\ \mathcal{J}[\mathcal{Y}(\mathsf{s};b)],
\ \mathsf{u}(\mathsf{s})\bigr)
\end{equation*}
restricted to $\mathcal{R}$, and at amplitude one: no response kernel
$(\delta/\delta b)\mathcal{Y}$ and no derivative in an amplitude is read; here, on each
component $C$ of $Y(\sigma)$, the restriction $\mathcal{A}(\mathsf{s})|_C$ is a function of
$(b|_C,\mathsf{s}|_C,(\mathbf{U},\mathbf{J})|_C)$ and equals $(0,0,1)$ when $b|_C=0$, a property
of the read triple and not of $F$.
Assume the following.
\begin{enumerate}
\item[(a)] The correlation theorem, in its reading-set clauses, takes $\mathcal{R}_0:=\hat{X}_v\cup
X_v$ for its exposed units and $\mathcal{R}_1:=\hat{X}_v\cup X_v$ for every unit of $\mathfrak{L}^{c}$,
and provides for these units functionals that read at amplitude one on these sets; for units
whose read box lies inside one component $C$ of $Y(\sigma)$ (the interior case) it takes for
$\mathcal{R}_0$ the read box of that unit, contained in $C$, the functional provided for such a
unit not reading the current $\mathbf{J}$; the formation
of composite reading functionals takes, for each link index $t$, $\mathcal{R}:=\mathcal{R}_{t-1}$ with
$\mathcal{R}_{t-1}\supseteq X_v$, both for units with $X_v\cap Z=\emptyset$ and for the others.
\item[(b)] By the definitions of the hull, the read box, the support region, the halo and the set
at which the consumer acts, $\hat{X}_v\cup X_v=Q_v$, and
\begin{equation*}
Q_v\subset\mathfrak{r}_v\subset\bigcup\mathfrak{h}_v\subset X^{\oplus}_v .
\end{equation*}
\item[(c)] The point $\operatorname{chain}(\vec{0})$ lies in the source polydisc $\mathfrak{P}$,
as established in the step of the proof of the correlation theorem that places the value at this
point.
\item[(d)] The assignment rule for exposed values places the value of a unit of
$\mathcal{V}^{\mathrm{ext}}$, the class of units whose values are exposed, among the exposed
values with a label; for $v$ the label is
$\mathfrak{K}_v:=K_v\cup\mathfrak{h}_v$. The supply rule, in its exposed clause, lets such
values enter the per-cell constant $\mathsf{K}_{\mathrm{val}}$ into which it collects the values
of units. They are summed per cell by the exposed collar summation at
$\bar{\mathsf{w}}_v=3$ and anchored as in the exposed clause of the anchoring lemma for exposed
families, that is by $\mathsf{K}^0\mapsto
\mathsf{K}^0+\mathsf{K}^{\flat,\mathrm{ex}}$, with $\mathsf{K}^{\flat}_J$ enlarged to
$\mathsf{K}^{\flat}_+$ inside $\mathsf{K}^{\flat,\mathrm{ex}}$ as in the enlargement lemma for the
collar constants. Moreover $v\in\mathcal{V}^{\mathrm{ext}}$,
because the cell $\Delta_0$ of $v$ lies in $\mathfrak{K}_v\setminus Z$.
\end{enumerate}
Then:
\begin{enumerate}
\item[(i)] the functional that the correlation theorem provides for $v$ reads at amplitude one on
$\mathcal{R}_0:=\hat{X}_v\cup X_v=Q_v$, no further condition being involved, and
\begin{equation}\label{10.3:eq-amplitude-one-reading-1}
\mathcal{R}_0=Q_v\subset\mathfrak{r}_v\subset\bigcup\mathfrak{h}_v\subset X^{\oplus}_v
\subset Y_4(\mathfrak{p});
\end{equation}
\item[(ii)] the value obeys, with $f_v(1)$ the reference value of the dialled function $f_v$,
\begin{equation}\label{10.3:eq-amplitude-one-reading-2}
\bigl|\operatorname{val}_v-f_v(1)\bigr|\le\mathfrak{N}^{\flat}(v);
\end{equation}
\item[(iii)] $\operatorname{val}_v$ is an exposed value with label
$\mathfrak{K}_v=K_v\cup\mathfrak{h}_v$, and enters $\mathsf{K}_{\mathrm{val}}$ through the exposed
clause of the supply rule, with the constants described in (d).
\end{enumerate}
Items (ii) and (iii) concern the value at $\operatorname{chain}(\vec{0})$; pieces whose link
families are non-empty are treated in the lemma on pieces with non-empty link families.
\end{lemma}

\begin{proof}
(i) The requirement that the definition of reading at amplitude one places on a functional is
its first sentence: the
functional reads $\mathcal{A}(\mathsf{s})$ restricted to a closed set $\mathcal{R}$, at amplitude
one. Here $\mathcal{R}$ is any closed set; nothing is required of its position relative to the
components $C$ of $Y(\sigma)$. The remaining clause of the definition concerns the triple
$\mathcal{A}(\mathsf{s})$ itself and involves no choice made by the functional. This is the form
in which the definition is used in hypothesis (a): there the sets $\mathcal{R}_{t-1}$ contain
$X_v$ whether or not $X_v$ meets $Z$, and the set $\mathcal{R}_1$ is taken for every unit of
$\mathfrak{L}^{c}$, which by the definition of $\mathfrak{L}^{c}$ lies astride $\partial Z$, as
does its reading set. The containment in one component $C$ is a property of the interior case,
not a part of the definition.

For the unit $v$, the reading set supplied by (a) is $\mathcal{R}_0:=\hat{X}_v\cup X_v$ if $v$
is an exposed unit and $\mathcal{R}_1:=\hat{X}_v\cup X_v$ if $v\in\mathfrak{L}^{c}$; in either
case it is $\hat{X}_v\cup X_v$, which by (b) equals $Q_v$. By (a) the correlation theorem provides
for $v$ a functional reading at amplitude one on this set. Since the definition contains no
further condition, it is
met for $v$ with $\mathcal{R}_0=Q_v$. The first three inclusions in
\eqref{10.3:eq-amplitude-one-reading-1} are hypothesis (b). The last one,
$X^{\oplus}_v\subset Y_4(\mathfrak{p})$, is Lemma~\ref{10.3:lem-consumer-domain}, applied to the
piece $\mathfrak{p}$ with unit $v$. This proves (i).

(ii) By definition $\operatorname{val}_v=f_v(\operatorname{chain}(\vec{0}))$, where $f_v$ is the
dialled function of $v$, and by hypothesis (c) the point $\operatorname{chain}(\vec{0})$ lies in the
source polydisc $\mathfrak{P}$. Lemma~\ref{10.3:lem-core-size-bound}
(holomorphy and size of the dialled core) holds at every point of $\mathfrak{P}$. At the point
$\operatorname{chain}(\vec{0})$ it gives, through the size bound
\eqref{10.3:eq-core-size-bound-a}, the inequality \eqref{10.3:eq-amplitude-one-reading-2}.

(iii) By (d), the cell $\Delta_0$ of $v$ lies in $\mathfrak{K}_v\setminus Z$, so that
$v\in\mathcal{V}^{\mathrm{ext}}$, and the assignment rule for exposed values places
$\operatorname{val}_v$ among the exposed values with label $\mathfrak{K}_v=K_v\cup\mathfrak{h}_v$.
The exposed clause of
the supply rule then takes $\operatorname{val}_v$ into $\mathsf{K}_{\mathrm{val}}$. The value is summed
per cell by the exposed collar summation at $\bar{\mathsf{w}}_v=3$ and anchored as in the exposed
clause of the anchoring lemma for exposed families, by
$\mathsf{K}^0\mapsto\mathsf{K}^0+\mathsf{K}^{\flat,\mathrm{ex}}$, where
$\mathsf{K}^{\flat}_J$ is enlarged to $\mathsf{K}^{\flat}_+$ inside $\mathsf{K}^{\flat,\mathrm{ex}}$
as in the enlargement lemma for the collar constants.
This is (iii).
\end{proof}

\begin{remark}[The dilated alternative and its filing of pieces]\label{10.3:rem-dilated-alternative}
Let $v$ be an inner-collar core. We use the notation of Definition~\ref{10.3:def-hull-halo-label}
and Lemma~\ref{10.3:lem-core-radius}:
\begin{itemize}
\item $X_v:=\square_0\subset Z$, and $C$ denotes the component of $Z$ containing $\square_0$;
$K_v:=[\square_0]_k$;
\item $Q_v$ is the read box, and $\hat{X}_v=Q_v$ is the hull;
\item $\mathfrak{h}_v$ is the halo, and $\mathfrak{r}_v=\bigcup\operatorname{St}(\Delta_y)$ is the support.
\end{itemize}

In this section $v$ is expanded as an interior first-part unit: there is no partition of unity and
no transport piece, and its pieces are estimated through $\varepsilon^{\mathrm{int}}(k)$ and
$\mathfrak{v}^{\mathrm{val}}$. Nothing in the present remark is used in that expansion.

We record how the pieces of $v$ would be filed under the dilated alternative, in which $v$ carries a
transport piece $p$ (the $\mathsf{t}$-piece); the choice between $X_v:=\square_0$ and $X_v:=Q_v$ is
immaterial below (see the step on the independence of the choice of $X_v$ near the end of the
proof). In this alternative the label of the first link is
\[
\operatorname{lab}_1:=K_v\cup\mathfrak{h}_v ,
\]
the same set as the start label $\operatorname{lab}_0$ of
Definition~\ref{10.3:def-hull-halo-label}.
It equals $\mathfrak{h}_v$, because
$K_v\subset[Q_v]_k\subset\mathfrak{h}_v$ by \eqref{10.3:eq-core-geometry-b}. Under the dilated
alternative the pieces would be filed as follows.
\begin{enumerate}
\item[(a)] The non-empty link-$2$--$4$ tuples of the first member would be taken by clause (c) of
the exterior link-interpolation statement in the proof of the correlation theorem, on the label
$\operatorname{lab}_1$.
\item[(b)] The piece $p$ would be taken by clause (e) of the same statement and by the summation of
compact pointed pieces, $p$ being compact.
\end{enumerate}
Were the dilated alternative adopted, the hypotheses of the statements named in (a) and (b) would
hold for $v$ as checked in the proof below, once the three scope conditions listed at the end of that
proof are widened by declaration so as to include $v$.
\end{remark}

\begin{proof}
\emph{Item (a).} We go through the proof of clause (c) of the exterior link-interpolation statement
step by step for $v$.

First, a non-empty family $\mathsf{y}_t$ contains a source cell at tree distance
$d_k\ge 9\check{R}_k-1$ from $Z^c$.

Next, the label $\operatorname{lab}_4$ of the fourth link contains $\operatorname{lab}_1$, and it
contains a cube not contained in $Z$. For $v$ this is a cube $\Delta_0\in[Q_v]_k$ with
$\Delta_0\not\subset Z$; one exists since $Q_v\not\subset Z$ and $Z$ is a union of $\pi_k$-cubes,
and it lies in $\mathfrak{h}_v$ by \eqref{10.3:eq-core-geometry-b}.

Since $\operatorname{lab}_4$ is connected, it follows that
\begin{equation*}
d_k(\operatorname{lab}_4)\ \ge\ 9\check{R}_k-2\ \ge\ \check{R}_k-2=:D .
\end{equation*}

The decay rates required in clause (c) are obtained by applying clauses (c) and (e) of the
rate-summation lemma in the proof of the correlation theorem three times. This is permitted provided
\begin{equation*}
\kappa_1-1\ \ge\ c_l(a+7)+\lambda_* ,
\end{equation*}
and this inequality is implied by the lower bound on $\kappa_1$ among the parameter conditions of the
correlation theorem.

The last step of clause (c) of the exterior link-interpolation statement, the transfer to a cube not
contained in $Z$, picks a cube $\Delta_v\in K_v$ with
$\Delta_v\not\subset Z$. For $v$ the cube is taken in the link-$1$ label: $\Delta_v:=\Delta_0$, which
lies in $\mathfrak{h}_v\subset\operatorname{lab}_1$ and is not contained in $Z$. The per-cell step
for the exposed pieces is carried out as for the exposed pieces in the proof of the correlation
theorem.

\emph{Item (b).} Clause (e) of the exterior link-interpolation statement estimates every pointed
piece $p$ of the classes listed there by the
estimate for pointed pieces in the proof of the correlation theorem, of which the summation of compact
pointed pieces is a corollary. It assigns to $p$ the exterior size
$\mathfrak{N}^{\mathrm{ex}}_A(p)$ of that clause, among whose factors are $2\,\mathsf{w}_L^{-1}$,
the weight $(\eta_p\mathsf{f}_v)^{1/2}$, $\tfrac{81}{64}\,\mathfrak{N}^{\flat}(v)$ and the decay
factor $e^{-(\kappa_1-1)m(p)}$. Here $\mathsf{f}_v\le 1$ is the weight attached to $v$ in clause (e),
and $\eta_p$ and $m(p)$ are taken in the notation of that clause.
The compact pieces $p$ are then summed over the units $v$, the link families $\mathsf{y}$ and the
indices $s$ and $\square$ of the summation of compact pointed pieces, by that summation. We check its
hypotheses in turn.

\emph{Compactness.} The piece $p$ is compact by definition. In the summation of compact pointed
pieces, a pointed piece is
compact if $v$ is a short single, a Wilson piece or a core, smooth or rough; and $v$ is a core.

\emph{Stars.} For a compact pointed unit the summation of compact pointed pieces sets
\begin{equation*}
S_v:=\bar{S}_v:=\mathfrak{r}_v=\bigcup\operatorname{St}(\Delta_{y}),\qquad
\operatorname{St}_v:=\operatorname{St}(\Delta_{y}) .
\end{equation*}

\emph{No source cell on the star.} The hypothesis of the lemma on source-free stars is that
$\operatorname{St}_v\cup\partial\operatorname{St}_v$ contains no source cell. It holds because
$\Delta_y$ lies within $1/L<2$ units of $\partial Z$ (Lemma~\ref{10.3:lem-core-radius}).

\emph{The datum of the theorem on sums over stars.} The datum is a triple
$\mathsf{u}=(\operatorname{St}_{\mathsf{u}},A_{\mathsf{u}},\mu_{\mathsf{u}})$ such that
\begin{itemize}
\item $\operatorname{St}_{\mathsf{u}}\subset\pi^{\mathfrak{s}}$ is finite and wall-connected, and
$\operatorname{St}_{\mathsf{u}}\cup\partial\operatorname{St}_{\mathsf{u}}$ contains no source cell;
\item $A_{\mathsf{u}}\in\mathbf{D}_k$ and $[\operatorname{St}_{\mathsf{u}}]_k\subset A_{\mathsf{u}}$;
\item $\mu_{\mathsf{u}}\ge 0$.
\end{itemize}
It is taken at $(\operatorname{St}(\Delta_y),\mathfrak{h}_v,\mu_v)$, where $\mu_v\ge 0$ is taken in
the notation of the theorem on sums over stars. The inclusion
$[\operatorname{St}(\Delta_y)]_k\subset\mathfrak{h}_v$ holds by
Lemma~\ref{10.3:lem-core-radius}.

\emph{Rim mass.} The rim-mass clause of the star lemma attached to the summation of compact pointed
pieces reads only $\Delta_y$ and the bound $\#\mathfrak{h}_v\le 7^d$. Its four bounds hold for $v$;
in them $\mathsf{b}_{\square}(\cdot)$ is the quantity, and $C^0_{\square}$, $\epsilon_{\partial}$
and $N_{\partial}$ are the constants, of that lemma, and $\operatorname{St}=\operatorname{St}_v$.
\begin{itemize}
\item $d_k(\mathfrak{h}_v)\le c_{\mathfrak{h}}-1$, the tree length of the halo being counted from
one cube.
\item $\mathsf{b}_{\square}(\operatorname{St}_v)\le C^0_{\square}\,c_S(c_S+3)$.
\item $e^{\epsilon_{\partial}\#\partial\operatorname{St}}\le e$.
\item For each cell $\Delta_1$, the cubes $\Delta'$ with $\Delta_1\in\partial\operatorname{St}(\Delta')$
number at most $N_{\partial}$.
\end{itemize}

\emph{The three factors.} The summation of compact pointed pieces assembles its bound from a collar
factor, a geometric factor and an activity factor.
\begin{itemize}
\item The collar factor is the exterior collar bound, weighted by $\mathsf{f}_v^{1/2}\le 1$, the
square root of the weight of $v$ from clause (e), at $\bar{\mathsf{w}}_v=3$.
\item The geometric factor is idle. The units of the class $\mathfrak{L}^{c}$ of the correlation
theorem, to which $v$ does not belong, and the pieces $p$ of the plain units of class (c) of the
classification of first-part units in the proof of the correlation theorem, lie along
$\partial C$, where $\bar{\mathsf{w}}=3$; for $v$ likewise $\bar{\mathsf{w}}_v=3$.
\item The activity factor is the bound on the activity of $v$ of
Lemma~\ref{10.3:lem-core-size-bound}.
\end{itemize}
The remaining auxiliary statements of the summation of compact pointed pieces use only the family of
stars $\operatorname{St}$ and the geometry of the seal of that summation.

\emph{Independence of the choice of $X_v$.} The datum $\mathsf{u}$ above carries
$A_{\mathsf{u}}=\mathfrak{h}_v$. This is the same set whether $X_v:=\square_0$ or $X_v:=Q_v$, since
$\square_0$ enters only through $K_v\subset\mathfrak{h}_v$.

\emph{Scope.} Three scope conditions would be extended under the dilated alternative by declaration,
so as to include $v$:
\begin{itemize}
\item the clause of the summation of compact pointed pieces fixing which units it counts, which as
stated covers the plain units of class (c) of the classification of first-part units;
\item the list of classes in clause (e) of the exterior link-interpolation statement;
\item the scope of the pointed half-power bound in the proof of the correlation theorem, which
as stated covers the compact pointed pieces of the plain units of classes (c) and (d) of the
classification of first-part units.
\end{itemize}
\end{proof}

\begin{definition}[Links, source polydisc and labels for inner-collar cores]
\label{10.3:not-links-polydisc-labels}
Throughout the treatment of inner-collar cores the following objects are fixed.

\emph{The class.} Let $\mathfrak{L}^{V,\mathrm{col}}$, the class of \emph{inner-collar cores}, consist
of the units $v$ that are smooth
far $\sharp$-cores of the vector part at the site $(\mathrm{L})$ of the operation $\mathbf{R}_k$,
the site labelled so in the expansion of the correlation theorem.
Each such unit has level $j$, anchor $y:=y_v$ and zone $n=k$, and the number $\nu:=1+n-j$
satisfies $\nu\ge a^*+1$. It lies
in a component $C$ of $Z$, and it satisfies
\begin{equation*}
X_v:=\square_0\subset C,\qquad K_v:=[\square_0]_k\subset C,\qquad
p_v:=p^*=n-2-s(3),\qquad Q_v:=Q^{(p^*)}(y)\not\subset C .
\end{equation*}
Such a unit is exposed in the sense of Definition~\ref{10.3:def-hull-halo-label}. Its hull is
$\hat{X}_v=Q_v$ and its halo is
\begin{equation*}
\mathfrak{h}_v=\{\Delta\in\pi_k:\ \operatorname{dist}_k(\Delta,\hat{X}_v)\le 2\}.
\end{equation*}
By Lemma~\ref{10.3:lem-core-radius}, \eqref{10.3:eq-core-radius-a}, the halo consists of at most
$6^d$ cubes, and $6^d\le c_{\mathfrak{h}}=7^d$. Its support set is $S_v:=\mathfrak{r}_v$, the
union of stars
of Lemma~\ref{10.3:lem-core-radius}, and its size is the number $\mathfrak{N}^{\flat}(v)$ of
Lemma~\ref{10.3:lem-core-size-bound}, \eqref{10.3:eq-core-size-bound-a}.

\emph{Links.} The site $(\mathrm{L})$ carries four links, $t=1,2,3,4$. Link $t$ is a site geometry
$(\pi^{\mathfrak{s}},\mathfrak{H},\mathfrak{F})$ in the sense of the definition of site
geometries and labels (whose derived notions, namely seeing, the reading condition, admissible
families and labels, are recalled below), with
\begin{equation*}
\pi^{\mathfrak{s}}:=\sigma_0^{(t)},\qquad
\mathfrak{H}:=\{\text{components of }\mathsf{K}\},\qquad
\mathsf{K}:=\begin{cases}\Lambda^{\sim}, & t=1,\\ \Lambda, & t=2,3,4.\end{cases}
\end{equation*}
Here $\Lambda$ is the domain of the link and $\Lambda^{\sim}$ its enlargement.
The members of $\mathfrak{H}$ are the \emph{hubs}. By the geometry of the hubs fixed with the
links, they are boxes with sides at most $h_{\sharp}$,
at mutual distance at least $19\check{R}_k$, and the cells sharing a wall with them have level $k$.
We put $\mathfrak{F}$ equal to the set of members of $\pi^{\mathfrak{s}}$ sharing a wall with a hub;
its members are the
\emph{source cells}, and they have level $k$.

By the expansion of the correlation theorem, in its rule for interior first-part units at the
site $(\mathrm{L})$, the site families of the four links are the following:
\begin{itemize}
\item $\sigma_0^{(1)}$ is the family of cells disjoint from $\Lambda^{\sim}$;
\item for $t\ge 2$, $\sigma_0^{(t)}$ is the family of cells disjoint from $\Lambda$.
\end{itemize}
Since $\Lambda\subset\Lambda^{\sim}$, a cell disjoint from $\Lambda^{\sim}$ or from $\Lambda$ does
not meet the interior of $\Lambda$; hence both kinds of family satisfy the condition imposed on
site families in the four-link expansion:
\begin{equation}\label{10.3:eq-links-polydisc-labels-1}
\text{no member of }\sigma_0^{(t)}\text{ meets the interior of }\Lambda,\qquad t=1,2,3,4 .
\end{equation}

\emph{Seeing, the reading condition, admissible families.} The notions of this paragraph are those
of the definition of site geometries and labels; we repeat their wording. Let $\mathcal{R}'$ be a
closed set. A
member of $\pi^{\mathfrak{s}}$ \emph{sees} $\mathcal{R}'$ if it meets $\mathcal{R}'$, or if it is a
source cell sharing a wall with a hub that meets $\mathcal{R}'$. For $\Delta\in\pi^{\mathfrak{s}}$
we let $\rho(\Delta)$ be the least cardinality of a connected subfamily of $\pi^{\mathfrak{s}}$ that
contains $\Delta$ and a source cell. Let $F=(A_F,\mathcal{R}_F)$ be a pair consisting of a label
$A_F$ (a set of the kind fixed in the definition of site geometries and labels) and a closed set
$\mathcal{R}_F$. The pair satisfies the \emph{reading condition} if every
$\Delta\in\pi^{\mathfrak{s}}$ that sees $\mathcal{R}_F$ satisfies
\begin{equation}\label{10.3:eq-links-polydisc-labels-2}
\operatorname{dist}_i(\operatorname{pa}(\Delta),A_F)\le c_{\rho}\,\rho(\Delta),
\qquad c_{\rho}:=120\ \ge\ h_{\sharp}=112 ,
\end{equation}
where $\operatorname{pa}(\Delta)$ is the parent cube of $\Delta$ and $\operatorname{dist}_i$ denotes
distance in level-$i$ units, $i$ being the level index of the definition of site geometries and
labels.
Let $\mathfrak{Y}(F)$ denote the set of finite non-empty families $\mathsf{y}\subset\pi^{\mathfrak{s}}$
such that every component of $\mathsf{y}$ contains a source cell and a member seeing
$\mathcal{R}_F$; these are the \emph{admissible families} of $F$.
For $\mathsf{y}\in\mathfrak{Y}(F)$ we put
\begin{equation*}
\operatorname{lab}(F,\mathsf{y}):=A_F\cup[\mathsf{y}]_k\cup\bigcup_{\iota}
\bigl(\operatorname{Cor}_{\iota}\cup\operatorname{Hub}_{\iota}\bigr),
\end{equation*}
where the union runs over the corridors $\operatorname{Cor}_{\iota}$ and hub corridors
$\operatorname{Hub}_{\iota}$ attached to $\mathsf{y}$ in the definition of site geometries and
labels, indexed by $\iota$.
For a set $A$ we put
\begin{equation*}
Y(A):=A\cup\bigcup\bigl\{C'\colon C'\ \text{a component of }Z,\ C'\cap A\neq\emptyset\bigr\}.
\end{equation*}

\emph{Polydisc.} The source polydisc of the four-link expansion is
\begin{equation*}
\mathfrak{P}:=\mathfrak{S}_{\mathrm{src}}\times\operatorname{supp}\Pi\times\prod_{t=1}^{4}
\bigl\{\mathsf{s}^{(t)}:\ |\mathsf{s}^{(t)}(\Delta)|\le e^{\kappa_1}\ \text{for}\
\Delta\in\sigma_0^{(t)}\bigr\}.
\end{equation*}
Its points are written $m=(\zeta,\varpi,\vec{\mathsf{s}})$: the source coordinate $\zeta$ ranges over
$\mathfrak{S}_{\mathrm{src}}$; the factor $\operatorname{supp}\Pi$ is the support of $\Pi$ as
fixed in the four-link expansion, and the coordinate $\varpi$ ranges over it; and the link
interpolation parameters $\vec{\mathsf{s}}=(\mathsf{s}^{(1)},\dots,\mathsf{s}^{(4)})$ range over the
last four factors. We let $\operatorname{chain}(m)$ be
the carrier pair of the site $(\mathrm{L})$ at the stage $\mathfrak{b}_4$ fixed for that site in the
four-link expansion, evaluated at the point $m$; when
$\zeta$ and $\varpi$ are held fixed we write $\operatorname{chain}(\vec{\mathsf{s}})$.

\emph{Labels and reading sets of $v$.} The start label and the first reading set are those of
Definition~\ref{10.3:def-hull-halo-label}:
\begin{equation*}
\operatorname{lab}_0=K_v\cup\mathfrak{h}_v,\qquad
\mathcal{R}_0=\hat{X}_v\cup X_v=Q_v .
\end{equation*}
By Definition~\ref{10.3:def-hull-halo-label}, $\operatorname{lab}_0$ equals $\mathfrak{h}_v$ as a
set. The equality $\hat{X}_v\cup X_v=Q_v$ is that of \eqref{10.3:eq-core-geometry-a}.

Given link families $\mathsf{y}_1,\dots,\mathsf{y}_4$, $\mathsf{y}_t\subset\sigma_0^{(t)}$, each
empty or admissible for its link $t$ in the sense of the pieces below, the labels
$\operatorname{lab}_t$ and reading sets $\mathcal{R}_t$ are defined recursively in $t$ as follows.
For $t\ge 2$ put $Y^{\mathrm{tr}}_{t-1}:=Y(\operatorname{lab}_{t-1})$ and
\begin{equation*}
\begin{split}
\mathcal{R}_{t-1}:=Q_v\ &\cup\ \bigcup_{t'<t}\ \bigcup\mathsf{y}_{t'}
\ \cup\ \bigl(\text{hubs touched or met by }\mathsf{y}_1,\dots,\mathsf{y}_{t-1}\bigr)\\
&\cup\ \bigl(\text{boundary layer of }Y^{\mathrm{tr}}_{t-1}\bigr),
\end{split}
\end{equation*}
the words ``touched'', ``met'' and ``boundary layer'' having the meaning of the definition of site
geometries and labels.
This is the reading set $\mathcal{R}_{t-1}$ of that definition, with $\hat{X}_v\cup X_v=Q_v$ in the
place of $X_v$. Since $v$ is plain (Definition~\ref{10.3:def-hull-halo-label}), its live-piece
datum $\mathcal{P}_v$ is empty, and the set $Y^{\mathrm{tr}}_{t}$ of the definition of site
geometries and labels reduces to $Y(\operatorname{lab}_t)$.
For $t=1,\dots,4$ we put
\begin{equation*}
\operatorname{lab}_t:=
\begin{cases}
\operatorname{lab}\bigl((\operatorname{lab}_{t-1},\mathcal{R}_{t-1}),\mathsf{y}_t\bigr),
& \mathsf{y}_t\neq\emptyset,\\
\operatorname{lab}_{t-1}, & \mathsf{y}_t=\emptyset ,
\end{cases}
\end{equation*}
as in the definition of site geometries and labels.

\emph{Pieces and value.} A piece of $v$, in the sense of the expansion of the correlation theorem
read with the start label and reading sets above, is a tuple
$\mathfrak{p}=(v;\mathsf{y}_1,\dots,\mathsf{y}_4)$
whose entries are not all empty and satisfy
\begin{equation*}
\mathsf{y}_t\in\mathfrak{Y}\bigl((\operatorname{lab}_{t-1},\mathcal{R}_{t-1})\bigr)\cup\{\emptyset\},
\qquad t=1,\dots,4 ,
\end{equation*}
where $\mathfrak{Y}$ is computed in the site geometry of link $t$.
The piece denotes the function of $(\zeta,\varpi)$ obtained by applying to
$f_v(\operatorname{chain}(\vec{\mathsf{s}}))$, with $\mathsf{s}^{(t)}=0$ off $\mathsf{y}_t$, the
commuting operators $\int_0^1 d\mathsf{s}^{(t)}(\Delta)\,\partial_{\mathsf{s}^{(t)}(\Delta)}$, one for
each $\Delta\in\mathsf{y}_t$, $t=1,\dots,4$:
\begin{equation}\label{10.3:eq-links-polydisc-labels-3}
\mathfrak{p}:=\Bigl[\prod_{t=1}^{4}\prod_{\Delta\in\mathsf{y}_t}\int_0^1 d\mathsf{s}^{(t)}(\Delta)\,
\partial_{\mathsf{s}^{(t)}(\Delta)}\Bigr]\,
f_v\bigl(\operatorname{chain}(\vec{\mathsf{s}})\bigr)
\Big|_{\mathsf{s}^{(t)}=0\ \text{off}\ \mathsf{y}_t} ,
\end{equation}
where, by abuse of notation, $f_v$ denotes the dialled function $\lambda_{n,k'}f_v$ of
Lemma~\ref{10.3:lem-core-size-bound}. The value of $v$ is
\begin{equation*}
\operatorname{val}_v:=f_v\bigl(\operatorname{chain}(\vec{0})\bigr).
\end{equation*}
The filing domain of a piece is
\begin{equation*}
Y_4(\mathfrak{p}):=Y(\operatorname{lab}_4),
\end{equation*}
which belongs to the class $\mathbf{D}_k^Z$ of filing domains of the expansion of the correlation
theorem.

\emph{Constants.} We put
\begin{equation*}
a:=\tfrac{31}{20}\,\kappa,\qquad D:=\check{R}_k-2,\qquad
c_0:=6^d+c_l\,c_{\mathfrak{h}},\qquad c_l=1+234d ,
\end{equation*}
where $\kappa$ is Ba{\l}aban's auxiliary parameter of that name.
\end{definition}

\begin{remark}\label{10.3:thm-interior-piece-bound}
The statement of this theorem is not written out here; parts of its proof are printed below.
\end{remark}

\begin{proof}[Proof of Theorem~\ref{10.3:thm-interior-piece-bound}, clause (a$'$): the pieces]
\label{10.3:prf-proof-pieces}
Throughout, $v\in\mathfrak{L}^{V,\mathrm{col}}$. We use the links, the source polydisc
$\mathfrak{P}$, the labels $\operatorname{lab}_t$, the reading sets $\mathcal{R}_{t-1}$, the
domains $Y_4(\mathfrak{p})=Y(\operatorname{lab}_4)$ and the pieces
$\mathfrak{p}=(v;\mathsf{y}_1,\dots,\mathsf{y}_4)$ of
Notation~\ref{10.3:not-links-polydisc-labels}. In particular
$\operatorname{lab}_0=K_v\cup\mathfrak{h}_v$, which as a set of $\pi_k$-cubes equals
$\mathfrak{h}_v$, and $\mathcal{R}_0=\hat{X}_v\cup X_v=Q_v$. A \emph{reader} $F$ is a label
$A_F$ together with a closed reading set $\mathcal{R}_F$ and a size letter; only the first two
enter below. We recall the three definitions of the
site geometry that are used below.
\begin{itemize}
\item A member of $\pi^{\mathfrak{s}}$ sees a closed set $\mathcal{R}$ if it meets
$\mathcal{R}$, or if it is a source cell sharing a wall with a hub that meets $\mathcal{R}$.
\item $\mathfrak{Y}(F)$ is the set of finite non-empty $\mathsf{y}\subset\pi^{\mathfrak{s}}$
such that every component of $\mathsf{y}$ contains a source cell and a member seeing
$\mathcal{R}_F$.
\item $F$ satisfies the reader condition if every $\Delta\in\pi^{\mathfrak{s}}$ seeing
$\mathcal{R}_F$ satisfies
$\operatorname{dist}_k(\operatorname{pa}(\Delta),A_F)\le c_{\rho}\,\rho(\Delta)$; here
$c_{\rho}=120\ge h_{\sharp}=112$, and $\rho(\Delta)$ is the least cardinality of a connected
subfamily containing $\Delta$ and a source cell.
\end{itemize}

In the proof of the correlation theorem, the clause on pieces rests on six items:
\begin{enumerate}
\item[(1)] holomorphy of the unit's function on $\mathfrak{P}$;
\item[(2)] the expansion in source cells, link by link, in which the components without a
source cell, or seeing nothing of $\mathcal{R}_{t-1}$, vanish;
\item[(3)] locality;
\item[(4)] the summation bound for labels, together with
$\operatorname{lab}_4\supset Y_4\setminus Z$ and $\operatorname{lab}_4\subset Y_4$, which give
$d_{k,Z}(Y_4)\le d_k(\operatorname{lab}_4)$;
\item[(5)] the reader condition;
\item[(6)] the chart.
\end{enumerate}
There the start label is $K_v$ and the reading set satisfies $\mathcal{R}_0\subset C$, so that
$\mathcal{R}_0\subset C\subset Y_4$. For our $v$ the start label is
$\operatorname{lab}_0=K_v\cup\mathfrak{h}_v$, and $\mathcal{R}_0=Q_v\not\subset C$. We verify the
six items in turn.

\medskip\noindent(1) Holomorphy.
This item does not involve the label. For $v$ it is the hypothesis of
Theorem~\ref{10.3:thm-interior-piece-bound}: $m\mapsto f_v(m)$ is holomorphic on $\mathfrak{P}$
and $|f_v(m)-f_v(1)|\le\mathfrak{N}^{\flat}(v)$ there. The bound for the exposed extension on
which this hypothesis rests is stated for every $m\in\mathfrak{P}$, where $\mathfrak{P}$ is the
product over all four links.

\medskip\noindent(2) The expansion in source cells, link by link.
We use the lemma of the correlation theorem's proof that expands a reader in source cells at one
link. Its hypotheses are three:
\begin{itemize}
\item holomorphy of the function $F$ in the link parameter;
\item membership of $F$ in the sourced class of readers;
\item the lemma on site families.
\end{itemize}
Its conclusion is
\begin{equation}\label{10.3:eq-proof-pieces-1}
F(\mathsf{R}(\mathbb{1}))-F(\mathsf{R}(0))=\sum_{\mathsf{y}\in\mathfrak{Y}(F)}F_{\mathsf{y}} ,
\end{equation}
where the sum runs over the finite non-empty $\mathsf{y}\subset\pi^{\mathfrak{s}}$ such that every
component of $\mathsf{y}$ holds a source cell and a member seeing $\mathcal{R}$; that is, over
$\mathfrak{Y}=\mathfrak{Y}(F)$.

The lemma on site families is stated for an arbitrary subfamily $\sigma_0$ of the underlying cell
partition. It therefore holds for $\sigma_0^{(1)}$ and for $\sigma_0^{(t)}$, $t\ge2$.

The requirement on a reader of the sourced class contains no clause on the position of its
reading set $\mathcal{R}$.

At link one, $F:=f_v\circ\operatorname{chain}$ reads the carrier through its $U$-slot on $Q_v$
alone. Indeed, by the definition of the value of a pointed unit, $f_v(m)$ is
$\mathfrak{i}^{(j)}_y$ evaluated at the extension reading $\operatorname{ext}_v(m)$, which, at the
scale $p:=p_v$, is formed from the box $Q_v$ and from the restriction
$\mathcal{U}(m)|_{Q_v}$ of the $U$-slot $\mathcal{U}(m)$ of $\operatorname{chain}(m)$, and from
nothing else; and the $J$-slot is not read. The $U$-slot of the link's output is the
$U$-component of the link's deformed pair $\mathsf{R}(\mathsf{s})$, namely
$e^{i\eta\mathcal{Y}(\mathsf{s};b)}\mathbf{K}$ on $Q_v$, where $\eta$ denotes the amplitude of
the link's deformation; the other component of $\mathsf{R}(\mathsf{s})$ fills the $J$-slot.
This $U$-slot is a function of $\mathcal{Y}|_{Q_v}$ at amplitude one. The gauge map of the link
is idle, because the value of a slot-free unit is a function of the gauge orbit.

Hence $f_v$ is a reader of the sourced class with reading set
$\mathcal{R}:=\mathcal{R}_0=Q_v$. This is the reading set of the correlation theorem's proof,
with the inclusion in $C$ dropped. In that proof the same expansion lemma is applied, under the
same choice of $\sigma_0^{(1)}$ (the cells disjoint with $\Lambda^{\sim}$), to the link-one reader
$(K_v,\mathcal{R}_0:=X_v)$ of the units whose $X_v$ misses $Z$; that reading set lies wholly in
$Z^{c}$. A reading set astride $\partial Z$ is therefore within the scope of the lemma, whose
requirement on readers does not involve the position of the reading set.

At the links $t\ge2$ we apply the expansion lemma to the function of $\mathsf{s}^{(t)}$ obtained
from $f_v\circ\operatorname{chain}$ by the differentiations and integrations already performed in
$\mathsf{s}^{(1)},\dots,\mathsf{s}^{(t-1)}$, the later parameters being left free. This function is
a reader of the sourced class with $\mathcal{R}:=\mathcal{R}_{t-1}$, and the lemma gives the index
set $\mathfrak{Y}$ at $\mathcal{R}_{t-1}$.

Summing the four instances of \eqref{10.3:eq-proof-pieces-1}, in the order $t=1,\dots,4$, gives
\begin{equation}\label{10.3:eq-proof-pieces-2}
f_v(\operatorname{chain}(\vec{\mathbb{1}}))
=\operatorname{val}_v+\sum_{(\mathsf{y}_1,\dots,\mathsf{y}_4)}\mathfrak{p},
\end{equation}
where each $\mathsf{y}_t$ ranges over
$\mathfrak{Y}((\operatorname{lab}_{t-1},\mathcal{R}_{t-1}))\cup\{\emptyset\}$, the $\mathsf{y}_t$
are not all empty, and the all-empty tuple contributes
$\operatorname{val}_v=f_v(\operatorname{chain}(\vec{0}))$. Each $\mathfrak{p}$ is obtained from
$f_v$ by derivatives in the parameters $\mathsf{s}^{(t)}(\Delta)$, $\Delta\in\mathsf{y}_t$,
integrals over $[0,1]$, and evaluation at $\mathsf{s}^{(t)}=0$ off $\mathsf{y}_t$. Each of these
operations is linear, so $\mathfrak{p}$ is linear in $f_v$.

It remains to check that the index sets are unaffected by the enlarged label and by the position
of $\mathcal{R}_0$. The definition of $\mathfrak{Y}(F)$ involves $\mathcal{R}_F$ and the site
geometry only. The label $A_F$ enters the new label $\operatorname{lab}(F,\mathsf{y})$ and the
reader condition, and nothing else.

Since $\mathcal{R}_F=Q_v$ meets $Z^{c}$, a member seeing $Q_v$ may be a cell of
$Z^{c}\cap\Omega_k$ beside $\partial C$. Such members belong to
$\pi^{\mathfrak{s}}=\sigma_0^{(1)}$, like all cells off $\Lambda^{\sim}$: every family runs over
all cells off $\Lambda^{\sim}$ (resp.\ $\Lambda$ for $t\ge2$), those of $X_v$ included. The
families through such members are members of $\mathfrak{Y}(F)$ by the one definition, and the
identity \eqref{10.3:eq-proof-pieces-1} holds for exactly this index set. Thus neither the
enlargement of the label nor $\mathcal{R}_0\not\subset C$ changes $\mathfrak{Y}$ in kind.

\medskip\noindent(3) Locality.
For a family $\mathsf{y}$, let $\mathsf{K}_{\mathsf{y}}$ denote the set of hubs sharing a wall
with $\mathsf{y}$ or met by $\mathcal{R}$, and let $X_F$ denote the set on which $F$ reads the
backgrounds. The third clause of the expansion lemma states that $F_{\mathsf{y}}$ reads the
sources $b$ on $\mathsf{K}_{\mathsf{y}}\cup\bigcup\mathsf{y}$ only, and reads the backgrounds
$((\mathbf{U},\mathbf{J}),\mathbf{A},w)$
on $X_F\cup\mathsf{K}_{\mathsf{y}}\cup\bigcup\mathsf{y}$ only. We apply it at each link with
$\mathcal{R}=\mathcal{R}_{t-1}$. Consequently $\mathfrak{p}$ reads the backgrounds on
\begin{equation*}
Q_v\cup\textstyle\bigcup_t\bigcup\mathsf{y}_t\cup{}\text{(the hubs touched or met)},
\end{equation*}
and reads the sources on the hubs touched or met. We show that each of these sets lies in
$Y_4(\mathfrak{p})$.

First, the box. Since $[Q_v]_k\subset\mathfrak{h}_v$ (a cube meeting $Q_v=\hat{X}_v$ is at
$\operatorname{dist}_k$ zero from it) and $\operatorname{lab}_4\supseteq\operatorname{lab}_0$
(each new label $\operatorname{lab}(F,\mathsf{y})$ contains $A_F$ by its definition in
Notation~\ref{10.3:not-links-polydisc-labels}), we have
\begin{equation}\label{10.3:eq-proof-pieces-3}
Q_v\subset\bigcup[Q_v]_k\subset\bigcup\mathfrak{h}_v\subset\bigcup\operatorname{lab}_0
\subset\bigcup\operatorname{lab}_4\subset Y_4(\mathfrak{p}).
\end{equation}

Second, the families. In the same way,
\begin{equation*}
\textstyle\bigcup\mathsf{y}_t\subset\bigcup[\mathsf{y}_t]_k\subset\bigcup\operatorname{lab}_t
\subset\bigcup\operatorname{lab}_4 .
\end{equation*}

Third, the hubs touched. A hub $H$ touched by $\mathsf{y}_t$ lies in a component $C'$ of $Z$, at
least $9\check{R}_k$ from the rim of $Z$ (the geometry of the hubs in the proof of the correlation
theorem). The touching member shares a wall with $H$, so its parent, a cube of
$\operatorname{lab}_t$, lies in $C'$. Hence $\operatorname{lab}_4$ meets $C'$, and
$H\subset C'\subset Y(\operatorname{lab}_4)$.

Fourth, the hubs met. No hub is met by $Q_v$; this is shown under (5) below.

Hence every $\mathfrak{p}$ reads inside $Y_4(\mathfrak{p})$. This is the step at which the proof
of the correlation theorem uses $\mathcal{R}_0\subset C\subset Y_4$. It uses only the consequence
$\mathcal{R}_0\subset Y_4$, and for $v$ this consequence is supplied by
\eqref{10.3:eq-proof-pieces-3}, since $\mathcal{R}_0=Q_v$. The same device occurs in the
correlation theorem for its smooth far cores.
There no inclusion of the box in $Z$ is required; the box $Q_v\subset[\square_0]_k^{(1)}$ is kept
inside the consumer's domain ($\mathfrak{h}_v\subset\operatorname{lab}_1$) and not inside $Z$.
Here $\operatorname{lab}_0$ takes the place of $\operatorname{lab}_1$.

\medskip\noindent(4) The summation bound and the tree length.
The start label $\operatorname{lab}_0=\mathfrak{h}_v$ belongs to $\mathbf{D}_k$, being connected
by \eqref{10.3:eq-core-radius-a}. The summation bound for labels in the proof of the correlation
theorem applies under the reader condition, which is verified under (5). It gives, for each $t$,
$\operatorname{lab}_t\in\mathbf{D}_k$ and
\begin{equation}\label{10.3:eq-proof-pieces-4}
d_k(\operatorname{lab}_t)\le d_k(\operatorname{lab}_{t-1})+c_l|\mathsf{y}_t| .
\end{equation}

Next, $Y_4:=Y(\operatorname{lab}_4)\in\mathbf{D}_k^Z$. By the definition
$Y(A)=A\cup\{\text{components of }Z\text{ meeting }A\}$ we have
$\operatorname{lab}_4\subset Y_4$, and $Y_4\setminus Z\subset\operatorname{lab}_4$ because the
components added lie in $Z$. Therefore an optimal tree of $\operatorname{lab}_4$ is admissible for
$d_{k,Z}(Y_4)$, and $d_{k,Z}(Y_4(\mathfrak{p}))\le d_k(\operatorname{lab}_4)$.

Finally, let $\Delta_0$ be the cube of $\operatorname{lab}_0$ not contained in $Z$ that is
provided by the geometry of inner-collar cores. For every tuple, the all-empty one included,
$\operatorname{lab}_4\supseteq\operatorname{lab}_0\ni\Delta_0$ with $\Delta_0\not\subset Z$. Since
$\Delta_0\subset\bigcup\operatorname{lab}_4\subset Y_4(\mathfrak{p})$, it follows that
$Y_4(\mathfrak{p})\setminus Z\neq\emptyset$. This is the last sentence of (a$'$).

\medskip\noindent(5) The reader condition.
We verify the reader condition afresh at the start reader
$(\operatorname{lab}_0,Q_v)$ and at each of its descendants.

Take first $t=1$, in the geometry of link one. We claim that a member $\Delta$ sees
$\mathcal{R}_0=Q_v$ if and only if it meets $Q_v$. It suffices that no hub meets $Q_v$. The hubs
lie at least $9\check{R}_k$ from the rim of $Z$: in the correlation theorem a source cell lies in
the collar of $\Lambda^{\sim}$, at $d_k\ge9\check{R}_k-1$ from $Z^{c}$. On the other hand, $Q_v$
holds a point of $\Delta_0\subset Z^{c}$ and has diameter less than one unit. Since
$\check{R}_k\ge38$, we have $9\check{R}_k\ge342>2$, and the claim follows.

Now let $\Delta$ meet $Q_v$, meeting being understood as having a common interior point. Then
$\operatorname{pa}(\Delta)$, the $\pi_k$-cube holding $\Delta$, has an interior point in $Q_v$.
Hence
\begin{align*}
&\operatorname{pa}(\Delta)\in[Q_v]_k\subset\mathfrak{h}_v=\operatorname{lab}_0,\\
&\operatorname{dist}_k(\operatorname{pa}(\Delta),\operatorname{lab}_0)=0\le c_{\rho}\,\rho(\Delta).
\end{align*}
In the proof of the correlation theorem it is also noted that such a member has level $k$. This is
not needed here, but it holds: $\Delta$ lies within one unit of the cell $\Delta_{y_v}$ at the
anchor, hence in $\Omega_k$ on either side of $\partial C$ by \eqref{10.3:eq-core-geometry-c}, and
the level there is $k$. The remaining case of the correlation theorem's verification, a source
cell at a hub met by $\mathcal{R}_0$, does not arise at $t=1$, since no hub meets $Q_v$.

Take next $t\ge2$, in the geometry of link $t$. The proof of the correlation theorem gives a list
of the members seeing $\mathcal{R}_{t-1}$. We use it with $Q_v$ in place of $X_v$ and with
$\operatorname{lab}_0\subset\operatorname{lab}_{t-1}$. A member $\Delta$ seeing
$\mathcal{R}_{t-1}$ falls into one of the following cases.
\begin{itemize}
\item $\Delta$ meets $Q_v$. The distance is $0$, as at $t=1$.
\item $\Delta$ meets an earlier $\bigcup\mathsf{y}_s$, $s<t$. Cells of the one underlying
partition meet only if they are equal, so $\Delta$ is a cell of $\mathsf{y}_s$. Its parent
lies in $[\mathsf{y}_s]_k\subset\operatorname{lab}_s\subset\operatorname{lab}_{t-1}$, and the
distance is $0$.
\item $\Delta$ is a source cell at a hub touched or met before. A cube of
$\operatorname{lab}_{t-1}$ shares a wall with that hub or lies in it. The distance is therefore at
most $h_{\sharp}\le c_{\rho}\le c_{\rho}\,\rho(\Delta)$, since $\rho(\Delta)\ge1$.
\item $\Delta$ meets the boundary layer of $Y(\operatorname{lab}_{t-1})$. Either $\Delta$ touches
the rim of $\operatorname{lab}_{t-1}$, and the distance is at most $1\le c_{\rho}\,\rho(\Delta)$.
Or $\Delta$ lies along the rim of a component $C'$ of $Z$ which $\operatorname{lab}_{t-1}$
meets, and the distance is at most $100\check{R}_k+1$. In the latter case $\Delta$ lies within
one unit of the rim of $Z$; every hub lies at least $9\check{R}_k$ from the rim of $Z$ (the
geometry of the hubs), and every source cell shares a wall with a hub. Hence
$\operatorname{dist}_k(\Delta,\bigcup\mathfrak{F})\ge9\check{R}_k-2$, and with the property
$\rho(\Delta)\ge\operatorname{dist}_k(\Delta,\bigcup\mathfrak{F})$ of $\rho$ this gives
\begin{equation}\label{10.3:eq-proof-pieces-5}
100\check{R}_k+1\le c_{\rho}(9\check{R}_k-2)\le c_{\rho}\,\rho(\Delta).
\end{equation}
The first inequality holds because, at $\check{R}_k=38$,
\begin{equation*}
c_{\rho}(9\check{R}_k-2)=120\cdot340=40800\ge3801=100\check{R}_k+1,
\end{equation*}
and the difference
\begin{equation*}
c_{\rho}(9\check{R}_k-2)-(100\check{R}_k+1)=980\check{R}_k-241
\end{equation*}
is increasing in $\check{R}_k$.
\item $\Delta$ lies in a live strip. This case is void, because $\mathcal{P}_v=\emptyset$.
\end{itemize}
So the start reader $(\operatorname{lab}_0,Q_v,\mathfrak{N}^{\flat}(v))$, in the geometry of link
one, and each of its descendants $(\operatorname{lab}_{t-1},\mathcal{R}_{t-1},\cdot)$, in the
geometry of link $t$, satisfy the reader condition. This verifies the reader condition (5); with
it the hypothesis left pending in (4) is met, and the absence of hubs met by $Q_v$ used in (3)
was shown above at $t=1$.

\medskip\noindent(6) The chart.
This item is void for $v$: the unit $v$ is slot-free, and for a slot-free unit the set of slots
is empty, the $\sharp$-form of the consumer coincides with the consumer, and every sentence on
rotations is void.

\medskip
Thus the items of the proof of the correlation theorem that involve $K_v$ involve it in two ways
only. They use it as the start label, which is replaced by $\operatorname{lab}_0$; the reader
condition for $\operatorname{lab}_0$ was verified again under (5). And they use it through
$\mathcal{R}_0\subset C\subset Y_4$, which is replaced by
$\mathcal{R}_0\subset\bigcup\mathfrak{h}_v\subset Y_4(\mathfrak{p})$ as in
\eqref{10.3:eq-proof-pieces-3}. This proves clause (a$'$) of
the corresponding display.

\medskip
It remains to verify, for every tuple, the containment behind the holomorphy hypothesis. The proof
of the bound for the exposed extension applies the consumer's estimate at
$\mathfrak{E}:=\mathfrak{E}^{\oplus}_v$, and it requires that the domain $X^{\oplus}_v$ lie in
$Y_4(\mathfrak{p})$. For $v$ we have $X^{\oplus}_v:=Y(\operatorname{lab}_0)$. For every tuple,
the all-empty one included, $\operatorname{lab}_0\subset\operatorname{lab}_4$ by the definition of
$\operatorname{lab}(F,\mathsf{y})$, and $Y$ is monotone. Hence
\begin{equation*}
X^{\oplus}_v=Y(\operatorname{lab}_0)\subset Y(\operatorname{lab}_4)=Y_4(\mathfrak{p}).
\end{equation*}
This is \eqref{10.3:eq-consumer-domain-a} of Lemma~\ref{10.3:lem-consumer-domain}, read with the
labels $\operatorname{lab}_t$ exactly as defined in
Notation~\ref{10.3:not-links-polydisc-labels}.
\end{proof}

\begin{proof}[Proof of Theorem~\ref{10.3:thm-interior-piece-bound}, continued: the piece bound and the
four-link summation]\label{10.3:prf-proof-four-link-sum}
In this part of the proof we establish clause (b$'$) and the first inequality of clause (c$'$) of
Theorem~\ref{10.3:thm-interior-piece-bound}, displayed in the corresponding display.
The links, the source polydisc $\mathfrak{P}$, its points $m=(\zeta,\varpi,\vec{\mathsf{s}})$ and the
labels $\operatorname{lab}_t$ are those of Notation~\ref{10.3:not-links-polydisc-labels}.

\medskip
\noindent\emph{The piece bound.} Let $v\in\mathfrak{L}^{V,\mathrm{col}}$ and let
$\mathfrak{p}=(v;\mathsf{y}_1,\dots,\mathsf{y}_4)$ be one of the pieces of $f_v$ in clause (a$'$).
By the definition of the pieces in clause (a$'$), the four link families of a piece are not all
empty: the all-empty tuple gives the value $\operatorname{val}_v$, which is not a piece; it is
carried by $\mathfrak{v}^{\mathrm{val}}$ with label $\mathfrak{K}_v:=K_v\cup\mathfrak{h}_v$ and is not
a summand of $\mathfrak{v}^{\mathrm{int},V}$. Hence there
is a $t\in\{1,2,3,4\}$ with $\mathsf{y}_t\neq\emptyset$.

Only the shallow case of the piece bound of the correlation theorem occurs here. Indeed, by
\eqref{10.3:eq-never-deep-a} inner-collar cores are never deep, so the two deep cases of that bound,
the one handled by the summation of boundary terms at finer zones and the deep pointed case, do not
arise. The shallow case
does not use the label at all; its only input is a bound on the oscillation of the function on the
polydisc, which here is the holomorphy hypothesis of Theorem~\ref{10.3:thm-interior-piece-bound}:
$m\mapsto f_v(m)$ is holomorphic on $\mathfrak{P}$, and
\begin{equation}\label{10.3:eq-proof-four-link-sum-1}
|f_v(m)-f_v(1)|\le\mathfrak{N}^{\flat}(v)\qquad\text{for all }m\in\mathfrak{P}.
\end{equation}

The shallow case rests on the Cauchy estimate for link pieces of the correlation theorem. Let $F$ be
holomorphic on the polydisc of a link, write $\mathsf{R}(\mathsf{s})$ for the link's deformed pair,
and let $F_{\mathsf{y}}$ be the piece of $F$ belonging to a link family $\mathsf{y}$. The estimate
states: if $c$ is a number independent of $\mathsf{s}$ and $\mathfrak{n}\ge 0$ is such that
$|F(\mathsf{R}(\mathsf{s}))-c|\le\mathfrak{n}$ on the polydisc, then
\begin{equation}\label{10.3:eq-proof-four-link-sum-2}
|F_{\mathsf{y}}|\le\mathfrak{n}\,e^{-(\kappa_1-1)|\mathsf{y}|}.
\end{equation}
In the correlation theorem this estimate is used on the product of the four polydiscs. The pieces are
Cauchy integrals over the product polydisc in the variables
$(\mathsf{s}^{(1)},\dots,\mathsf{s}^{(4)})$, and the holomorphy required is joint holomorphy in all of
them. The Cauchy integral is taken in the complexified parameters $\sigma(\Delta)$ replacing
$\mathsf{s}(\Delta)$, factor by factor over the cells $\Delta$ of
$\bigcup_t\mathsf{y}_t$, on the circles
\begin{equation*}
|\sigma(\Delta)-\mathsf{s}(\Delta)|=e^{\kappa_1}-1,\qquad \mathsf{s}(\Delta)\in[0,1].
\end{equation*}
Each factor contributes at most $(e^{\kappa_1}-1)^{-1}$. Moreover $e^{\kappa_1}-1\ge e^{\kappa_1-1}$:
this is $e^{\kappa_1}(1-e^{-1})\ge1$, which holds since $\kappa_1\ge1$ by the floor on $\kappa_1$
among the hypotheses of Theorem~\ref{10.3:thm-interior-piece-bound} (see Step~6 below). Thus each factor
contributes at most $e^{-(\kappa_1-1)}$. Since the four link families carry separate variables, there
are $\sum_t|\mathsf{y}_t|$ factors.

We apply this with $F:=f_v$, $c:=f_v(1)$ and $\mathfrak{n}:=\mathfrak{N}^{\flat}(v)$. The number $c$
does not depend on $\vec{\mathsf{s}}$. Joint holomorphy on $\mathfrak{P}$ is the holomorphy hypothesis
recalled above, and the oscillation bound is \eqref{10.3:eq-proof-four-link-sum-1}. Let
$\partial_{\mathsf{y}}$ denote the operation
that produces from a function of the link parameters its piece for the tuple
$(\mathsf{y}_1,\dots,\mathsf{y}_4)$. Because some $\mathsf{y}_t$ is non-empty, $\partial_{\mathsf{y}}$
annihilates the $\vec{\mathsf{s}}$-independent constant $c$, so
$\partial_{\mathsf{y}}f_v=\partial_{\mathsf{y}}(f_v-c)$. The estimate therefore gives
\begin{equation}\label{10.3:eq-proof-four-link-sum-3}
|\mathfrak{p}|\le\mathfrak{N}^{\flat}(v)\,e^{-(\kappa_1-1)\sum_t|\mathsf{y}_t|}.
\end{equation}
This is clause (b$'$). It shows that the size of $v$ entering the sums below is
$\|v\|_{\star}:=\mathfrak{N}^{\flat}(v)$.

\medskip
\noindent\emph{The sum: what is to be proved.} Recall from clause (c$'$) that
$\mathfrak{v}^{\mathrm{int},V}(Y_4)$ is the sum of the pieces $\mathfrak{p}$, over
$v\in\mathfrak{L}^{V,\mathrm{col}}$, with $Y_4(\mathfrak{p})=Y_4$. Here $Y_4(\mathfrak{p})$ is
$\operatorname{lab}_4(\mathfrak{p})$ together with the components of $Z$ that meet it. Put
$D:=\check{R}_k-2$ and $a:=\tfrac{31}{20}\kappa$. We prove
\begin{equation}\label{10.3:eq-proof-four-link-sum-4}
\sup_{\Delta\in\pi_k,\ \Delta\not\subset Z}\ \sum_{Y_4\ni\Delta}e^{a\,d_{k,Z}(Y_4)}
\bigl|\mathfrak{v}^{\mathrm{int},V}(Y_4)\bigr|
\le e^{-D}\,(1+C_{\Sigma}C_e)^4\,\mathsf{Q}^V_{a+9}.
\end{equation}
The left member is $\|\mathfrak{v}^{\mathrm{int},V}\|_{a,1}$. The second inequality of clause (c$'$),
the bound on $\mathsf{Q}^V_r$ for $r\le a+9$, is proved in the next part of the proof of
Theorem~\ref{10.3:thm-interior-piece-bound}.

\medskip
\noindent\emph{Step 1: a cube off $Z$ lies in the last label.} Let $\Delta\in\pi_k$ with
$\Delta\not\subset Z$, and let $\mathfrak{p}$ be a piece with $\Delta\subset Y_4(\mathfrak{p})$.
By the definition of $Y_4(\mathfrak{p})$ recalled above, a $\pi_k$-cube
$\Delta\subset Y_4(\mathfrak{p})$ with $\Delta\not\subset Z$ has interior
off $Z$. It is therefore a cube of $\operatorname{lab}_4(\mathfrak{p})$.

\medskip
\noindent\emph{Step 2: the last label is long.} As noted at the start of the piece bound, some
$\mathsf{y}_t$ is non-empty. Fix such a $t$; the argument does not need $K_v\subset Z$.

By the definition of the admissible link families of the correlation theorem, each component of
$\mathsf{y}_t$ contains a \emph{source cell} $\Delta^{\mathfrak{s}}$. This is a level-$k$ cube sharing
a wall with a hub of the link, that is, with $\Lambda^{\sim}$ or with $\Lambda$. It satisfies
$\operatorname{pa}(\Delta^{\mathfrak{s}})=\Delta^{\mathfrak{s}}$, so
\begin{equation*}
\Delta^{\mathfrak{s}}\in[\mathsf{y}_t]_k\subset\operatorname{lab}_t(\mathfrak{p})
\subset\operatorname{lab}_4(\mathfrak{p}).
\end{equation*}
By the geometric lemma of the correlation theorem on the position of source cells,
\begin{equation*}
\operatorname{dist}_k(\Delta^{\mathfrak{s}},Z^c)\ge 9\check{R}_k-1.
\end{equation*}
The cube $\Delta$ of Step~1 has its interior off $Z$, as shown there, so $\Delta$ lies in the
closure of $Z^c$. Hence
\begin{equation*}
\operatorname{dist}_k(\Delta,\Delta^{\mathfrak{s}})\ge\operatorname{dist}_k(\Delta^{\mathfrak{s}},Z^c)
\ge 9\check{R}_k-1.
\end{equation*}
The label $\operatorname{lab}_4(\mathfrak{p})$ belongs to the domain class $\mathbf{D}_k$. It is
therefore connected, by clause (a) of the summation lemma for reader families of the correlation
theorem. By Step~1 and the above it contains both $\Delta$ and $\Delta^{\mathfrak{s}}$.

We now apply the tree-length bound: a tree graph meeting two cubes $Q,Q'$ has length at least
$\operatorname{dist}(Q,Q')$. In level-$k$ units this yields
\begin{equation}\label{10.3:eq-proof-four-link-sum-5}
d_k(\operatorname{lab}_4(\mathfrak{p}))\ge\operatorname{dist}_k(\Delta,\Delta^{\mathfrak{s}})-1
\ge 9\check{R}_k-2\ge\check{R}_k-2=D .
\end{equation}

\medskip
\noindent\emph{Step 3: readers and chains.} A \emph{reader} is a triple $F=(A_F,\mathcal{R}_F,\|F\|)$
that satisfies the reading condition of the summation lemma in the site geometry in which it is used.
Its entries are a set $A_F$ of $\pi_k$-cubes, a reading set $\mathcal{R}_F$, and a size $\|F\|\ge0$.
The admissible link families of a reader $F$ form the set $\mathfrak{Y}(F)$. For
$\mathsf{y}\in\mathfrak{Y}(F)$, $\operatorname{lab}(F,\mathsf{y})$ is the link label.

Put
\begin{equation*}
\mathfrak{L}^0_V:=\bigl\{F^0_v:=(\operatorname{lab}_0(v),\,Q_v,\,\mathfrak{N}^{\flat}(v)):
\ v\in\mathfrak{L}^{V,\mathrm{col}}\bigr\}.
\end{equation*}
By the reading condition of the summation lemma, verified at $t=1$ for the link geometries in the part
of this proof that treats the pieces (the proof part~\ref{10.3:prf-proof-pieces}), the members of
$\mathfrak{L}^0_V$ are readers in the geometry
of link~$1$.

For $t=1,\dots,4$ define, in the geometry of link $t$,
\begin{equation*}
(\mathfrak{L}^{t-1}_V)':=\bigl\{(\operatorname{lab}(F,\mathsf{y}_t),\,\mathcal{R}_t,\,
\|F\|e^{-(\kappa_1-1)|\mathsf{y}_t|}):\ F\in\mathfrak{L}^{t-1}_V,\ \mathsf{y}_t\in\mathfrak{Y}(F)\bigr\},
\end{equation*}
and set $\mathfrak{L}^t_V:=\mathfrak{L}^{t-1}_V\cup(\mathfrak{L}^{t-1}_V)'$. Here
$(\mathfrak{L}^{t-1}_V)'$ is the derived family $\mathfrak{L}'$ of clause (e) of the summation lemma,
taken with the reading set $\mathcal{R}':=\mathcal{R}_t$. By the same reading condition, the members of
$\mathfrak{L}^t_V$ are readers in the geometry of link $t+1$.

To a tuple $(v;\mathsf{y}_1,\dots,\mathsf{y}_4)$ we associate a chain. Put $F^0:=F^0_v$. For
$t=1,\dots,4$ put $F^t:=F^{t-1}$ if $\mathsf{y}_t=\emptyset$, and otherwise
\begin{equation*}
F^t:=\bigl(\operatorname{lab}(F^{t-1},\mathsf{y}_t),\,\mathcal{R}_t,\,
\|F^{t-1}\|e^{-(\kappa_1-1)|\mathsf{y}_t|}\bigr).
\end{equation*}
By the definition of the labels of a piece in Notation~\ref{10.3:not-links-polydisc-labels},
$A_{F^4}=\operatorname{lab}_4(\mathfrak{p})$. By \eqref{10.3:eq-proof-four-link-sum-3},
\begin{equation}\label{10.3:eq-proof-four-link-sum-6}
\|F^4\|=\mathfrak{N}^{\flat}(v)\,e^{-(\kappa_1-1)\sum_t|\mathsf{y}_t|}\ge|\mathfrak{p}|.
\end{equation}
Distinct tuples give distinct chains.

We regard $\mathfrak{L}^4_V$ as counted along chains, a member being listed once for each chain that
produces it. So counted, the sums over $\mathfrak{L}^4_V$ below majorise the corresponding sums over
tuples. The chains range over all admissible choices, not only over those coming from pieces.
Dropping that restriction only adds non-negative terms; this is the majorant book-keeping of the
correlation theorem.

\medskip
\noindent\emph{Step 4: reduction to one reader sum.} For a family $\mathfrak{L}$ of readers and
$r\ge0$ put
\begin{equation}\label{10.3:eq-proof-four-link-sum-7}
\mathsf{Q}_r[\mathfrak{L}]:=\sup_{Q\in\pi_k}\ \sum_{F\in\mathfrak{L}:\ A_F\ni Q}\|F\|\,
e^{r\,d_k(A_F)} .
\end{equation}
This is the anchored weighted reader sum of the summation lemma.

Let $\Delta\in\pi_k$ with $\Delta\not\subset Z$. We use four facts. The triangle inequality applies
to the sum defining $\mathfrak{v}^{\mathrm{int},V}(Y_4)$. By Step~1, $Y_4(\mathfrak{p})\ni\Delta$
implies $\operatorname{lab}_4(\mathfrak{p})\ni\Delta$. Next,
$d_{k,Z}(Y_4(\mathfrak{p}))\le d_k(\operatorname{lab}_4(\mathfrak{p}))$. Finally, for $d\ge D$,
\begin{equation*}
e^{ad}=e^{-d}e^{(a+1)d}\le e^{-D}e^{(a+1)d}.
\end{equation*}
By \eqref{10.3:eq-proof-four-link-sum-5} this applies with $d=d_k(\operatorname{lab}_4(\mathfrak{p}))$.
Combining these with the chains of Step~3 and with \eqref{10.3:eq-proof-four-link-sum-6}, we get
\begin{align}
\sum_{Y_4\ni\Delta}e^{a\,d_{k,Z}(Y_4)}\bigl|\mathfrak{v}^{\mathrm{int},V}(Y_4)\bigr|
&\le\sum_{\mathfrak{p}:\ \operatorname{lab}_4(\mathfrak{p})\ni\Delta}
e^{a\,d_k(\operatorname{lab}_4(\mathfrak{p}))}\,|\mathfrak{p}| \notag\\
&\le\sum_{\text{chains}:\ A_{F^4}\ni\Delta}e^{-D}\,e^{(a+1)d_k(A_{F^4})}\,\|F^4\| \notag\\
&\le e^{-D}\,\mathsf{Q}_{a+1}[\mathfrak{L}^4_V]. \label{10.3:eq-proof-four-link-sum-8}
\end{align}
The last step holds because $\Delta$ is one of the cubes $Q\in\pi_k$ over which the supremum in
\eqref{10.3:eq-proof-four-link-sum-7} is taken.

\medskip
\noindent\emph{Step 5: four applications of the summation lemma.} Fix $t\in\{4,3,2,1\}$ and the
corresponding rate $r\in\{a+1,a+3,a+5,a+7\}$, and work in the geometry of link $t$.

Since $\mathfrak{L}^t_V$ is the union of $\mathfrak{L}^{t-1}_V$ and $(\mathfrak{L}^{t-1}_V)'$, the sum
in \eqref{10.3:eq-proof-four-link-sum-7} over $\mathfrak{L}^t_V$ is at most the sum of the two sums.
Hence
\begin{equation*}
\mathsf{Q}_r[\mathfrak{L}^t_V]\le\mathsf{Q}_r[\mathfrak{L}^{t-1}_V]
+\mathsf{Q}_r[(\mathfrak{L}^{t-1}_V)'].
\end{equation*}

Clause (c) of the summation lemma, in the form of the first assertion of its clause (e), gives
$\mathsf{Q}_r[\mathfrak{L}']\le C_{\Sigma}\,\mathsf{T}^*_{r+2}[\mathfrak{L}]$, provided
\begin{equation}\label{10.3:eq-proof-four-link-sum-9}
\kappa_1-1\ge c_l\,r+\lambda_* .
\end{equation}
It applies to $\mathfrak{L}=\mathfrak{L}^{t-1}_V$: every $F\in\mathfrak{L}^{t-1}_V$ is a reader in the
geometry of link $t$ (Step~3), and every $\mathsf{y}_t$ used belongs to $\mathfrak{Y}(F)$. Therefore
\begin{equation*}
\mathsf{Q}_r[(\mathfrak{L}^{t-1}_V)']\le C_{\Sigma}\,\mathsf{T}^*_{r+2}[\mathfrak{L}^{t-1}_V].
\end{equation*}

The second assertion of clause (e) of the summation lemma holds for any family, in any site geometry
of any level, with any reading sets obeying the reading condition. It states
$\mathsf{T}^*_{r}[\mathfrak{L}]\le C_e\,\mathsf{Q}_{r}[\mathfrak{L}]$, hence
\begin{equation*}
\mathsf{T}^*_{r+2}[\mathfrak{L}^{t-1}_V]\le C_e\,\mathsf{Q}_{r+2}[\mathfrak{L}^{t-1}_V].
\end{equation*}

Finally, $\mathsf{Q}_r\le\mathsf{Q}_{r+2}$, because the sizes are non-negative and $d_k\ge0$.
Combining the union bound, the two bounds just displayed and $\mathsf{Q}_r\le\mathsf{Q}_{r+2}$,
\begin{equation}\label{10.3:eq-proof-four-link-sum-10}
\mathsf{Q}_r[\mathfrak{L}^t_V]\le(1+C_{\Sigma}C_e)\,\mathsf{Q}_{r+2}[\mathfrak{L}^{t-1}_V].
\end{equation}
Applying \eqref{10.3:eq-proof-four-link-sum-10} successively for $(t,r)=(4,a+1)$, $(3,a+3)$,
$(2,a+5)$, $(1,a+7)$ gives
\begin{equation}\label{10.3:eq-proof-four-link-sum-11}
\mathsf{Q}_{a+1}[\mathfrak{L}^4_V]\le(1+C_{\Sigma}C_e)^4\,\mathsf{Q}_{a+9}[\mathfrak{L}^0_V].
\end{equation}

\medskip
\noindent\emph{Step 6: the parameter floor.} The largest rate at which clause (c) of the summation
lemma is used is $r=a+7$. Hence \eqref{10.3:eq-proof-four-link-sum-9} holds at all four uses once
$\kappa_1-1\ge c_l(a+9)+\lambda_*$. This in turn follows from the floor
\begin{equation*}
\kappa_1\ge 1+c_l(2\kappa+4)+\lambda_*,
\end{equation*}
which is among the parameter hypotheses of Theorem~\ref{10.3:thm-interior-piece-bound}. To see this,
note first that
\begin{equation*}
a+9=\tfrac{31}{20}\kappa+9\le\tfrac85\kappa
\quad\Longleftrightarrow\quad \tfrac{1}{20}\kappa\ge9
\quad\Longleftrightarrow\quad\kappa\ge180,
\end{equation*}
and $\kappa\ge180$ is a hypothesis of that theorem. Second, $\tfrac85\kappa\le2\kappa+4$. Since
$c_l>0$, it follows that $c_l(a+9)\le c_l(2\kappa+4)$.

The constants $C_{\Sigma}$, $C_e$ of the summation lemma depend on the reading condition, on the
admissible families $\mathfrak{Y}(F)$, on the site geometry and on the floors involving $\lambda_*$.
No constant of that lemma depends on the sizes, faces, volumes or ages of the regions involved, nor on
their number, nor on the family $\mathfrak{L}$. In particular nothing in Steps~5 and~6 uses
$A_F\subset Z$, $A_F\subset C$ or $\mathcal{R}_F\subset C$, where, when $K_v\subset Z$, $C$ denotes
the component of $Z$ containing $K_v$.

\medskip
\noindent\emph{Step 7: conclusion.} The members of $\mathfrak{L}^0_V$ have
$A_{F^0_v}=\operatorname{lab}_0(v)$ and $\|F^0_v\|=\mathfrak{N}^{\flat}(v)$. Comparing
\eqref{10.3:eq-proof-four-link-sum-7} with the definition of $\mathsf{Q}^V_r$ in clause (c$'$) shows
$\mathsf{Q}_{a+9}[\mathfrak{L}^0_V]=\mathsf{Q}^V_{a+9}$. Inserting
\eqref{10.3:eq-proof-four-link-sum-11} into \eqref{10.3:eq-proof-four-link-sum-8} and taking the
supremum over $\Delta\in\pi_k$ with $\Delta\not\subset Z$ gives
\begin{equation*}
\|\mathfrak{v}^{\mathrm{int},V}\|_{a,1}\le e^{-(\check{R}_k-2)}\,(1+C_{\Sigma}C_e)^4\,
\mathsf{Q}^V_{a+9}.
\end{equation*}
This is \eqref{10.3:eq-proof-four-link-sum-4}, the first inequality of clause (c$'$).
\end{proof}

\begin{lemma}[Anchoring at a cube of the halo]\label{10.3:prf-proof-halo-anchor}
Work in the setting of Theorem~\ref{10.3:thm-interior-piece-bound}. Recall the constants
\[
a=\tfrac{31}{20}\kappa,\qquad c_{\mathfrak{h}}=7^d,\qquad c_l=1+234d,\qquad
c_0=6^d+c_l c_{\mathfrak{h}}
\]
of clause (c$'$) of Theorem~\ref{10.3:thm-interior-piece-bound} and of
Lemma~\ref{10.3:lem-core-radius}.
For $r>0$ let, as in clause (c$'$) of Theorem~\ref{10.3:thm-interior-piece-bound},
\[
\mathsf{Q}^V_r:=\sup_{Q\in\pi_k}\ \sum_{v\in\mathfrak{L}^{V,\mathrm{col}}:\ Q\in\operatorname{lab}_0(v)}
\mathfrak{N}^{\flat}(v)\,e^{r\,d_k(\operatorname{lab}_0(v))}.
\]
Then
\begin{equation}\label{10.3:eq-proof-halo-anchor-a}
\begin{split}
\mathsf{Q}^V_{a+9}\le{}& e^{(a+9)c_l c_{\mathfrak{h}}}\,(c_0+2)(2c_0+5)^d\\
&\times\sup_{\Delta'}\ \sum_{v\in\mathfrak{L}^{V,\mathrm{col}}:\ y_v\in\Delta'}
\mathfrak{N}^{\flat}(v)\,e^{(a+9)\,d_k(K_v)},
\end{split}
\end{equation}
where the supremum runs over the cells $\Delta'$ of level $k$ of the collar of the deep core
set $\mathsf{C}$.
\end{lemma}

\begin{proof}
The bound \eqref{10.3:eq-proof-halo-anchor-a} is the first step towards the second inequality of
clause (c$'$) of Theorem~\ref{10.3:thm-interior-piece-bound}. What distinguishes it from the
anchoring inside $K_v$ is that the anchoring cube $Q$ is only required to lie in the start label
$\operatorname{lab}_0(v)$: it may lie in $\mathfrak{h}_v\setminus K_v$, it may lie off
$\mathsf{C}$, and it may or may not lie in the large-field region $Z$.

Fix $Q\in\pi_k$ and a unit $v\in\mathfrak{L}^{V,\mathrm{col}}$ with
$Q\in\operatorname{lab}_0(v)$. Write $y_v$ for the point of $v$ and $\Delta_{y_v}$ for the cell
of level $k$ that contains it.

\emph{Position of the point.} By Lemma~\ref{10.3:lem-never-deep} (inner-collar cores are never
deep), in the form \eqref{10.3:eq-never-deep-b}, the point $y_v$ lies in the cell
$\Delta':=\Delta_{y_v}$, which is a level-$k$ cell of the collar of $\mathsf{C}$, and
\[
\Delta'\in K_v\subset\operatorname{lab}_0(v).
\]

\emph{Distance from the anchor.} By Lemma~\ref{10.3:lem-core-radius} (radius, support and halo
of an inner-collar core), in the form \eqref{10.3:eq-core-radius-a}, the start label is
$\operatorname{lab}_0(v)=K_v\cup\mathfrak{h}_v=\mathfrak{h}_v$, and this set of $\pi_k$-cubes
is connected. It contains both $Q$ and $\Delta'$. For a connected set $A$ of $\pi_k$-cubes, the
$\operatorname{dist}_k$-distance of any two of its members is at most $d_k(A)+1$; this is the
connectivity inequality for tree lengths of sets of $\pi_k$-cubes. Applied to
$A=\operatorname{lab}_0(v)$ it gives
\begin{equation}\label{10.3:eq-proof-halo-anchor-1}
i:=\operatorname{dist}_k(\Delta',Q)\le d_k(\operatorname{lab}_0(v))+1 .
\end{equation}

\emph{Tree length of the start label.} By \eqref{10.3:eq-core-radius-a},
$d_k(\operatorname{lab}_0(v))\le d_k(K_v)+c_l c_{\mathfrak{h}}$. By
\eqref{10.3:eq-core-geometry-b} (geometry of an inner-collar core),
$d_k(K_v)\le 2^d-1\le 6^d$. Together,
\begin{equation}\label{10.3:eq-proof-halo-anchor-2}
d_k(\operatorname{lab}_0(v))\le d_k(K_v)+c_l c_{\mathfrak{h}}\le 6^d+c_l c_{\mathfrak{h}}=c_0 ,
\end{equation}
so that, by \eqref{10.3:eq-proof-halo-anchor-1}, $i\le c_0+1$. The first inequality in
\eqref{10.3:eq-proof-halo-anchor-2}, multiplied by $a+9>0$ and exponentiated, gives
\begin{equation}\label{10.3:eq-proof-halo-anchor-3}
e^{(a+9)\,d_k(\operatorname{lab}_0(v))}\le e^{(a+9)c_l c_{\mathfrak{h}}}\,e^{(a+9)\,d_k(K_v)} .
\end{equation}

\emph{Summation over the units.} Each unit $v\in\mathfrak{L}^{V,\mathrm{col}}$ with
$Q\in\operatorname{lab}_0(v)$ determines exactly one cell $\Delta'=\Delta_{y_v}$. By the two
preceding paragraphs, $\Delta'$ is a level-$k$ cell of the collar of $\mathsf{C}$ with
$\operatorname{dist}_k(\Delta',Q)\le c_0+1$. Group the units according to this cell, apply
\eqref{10.3:eq-proof-halo-anchor-3} to every term, and then drop the condition
$Q\in\operatorname{lab}_0(v)$ in the inner sum; the last step only enlarges the sum, since every
term is nonnegative. This gives
\begin{align*}
\sum_{v\in\mathfrak{L}^{V,\mathrm{col}}:\ Q\in\operatorname{lab}_0(v)}
\mathfrak{N}^{\flat}(v)\,e^{(a+9)d_k(\operatorname{lab}_0(v))}
&\le e^{(a+9)c_l c_{\mathfrak{h}}}
\sum_{\Delta'}
\ \sum_{v\in\mathfrak{L}^{V,\mathrm{col}}:\ y_v\in\Delta'}
\mathfrak{N}^{\flat}(v)\,e^{(a+9)d_k(K_v)},
\end{align*}
where $\Delta'$ runs over the level-$k$ cells of the collar of $\mathsf{C}$ with
$\operatorname{dist}_k(\Delta',Q)\le c_0+1$.

These cells are $\pi_k$-cubes. Those at distance $i$ from $Q$ number at most $(2i+3)^d$.
Hence the number of cells in the outer sum is at most
\[
\sum_{0\le i\le c_0+1}(2i+3)^d\le (c_0+2)(2c_0+5)^d ,
\]
because there are $c_0+2$ values of $i$ and each term is at most $(2(c_0+1)+3)^d$. Bounding
each inner sum by its supremum over the level-$k$ cells of the collar of $\mathsf{C}$, and then
taking the supremum over $Q\in\pi_k$, gives \eqref{10.3:eq-proof-halo-anchor-a}.

\emph{The same transfer in the correlation theorem.} The step just made is the transfer from the
anchoring cube to the cell of the point that the correlation theorem makes in its
last step for its exposed pointed units. There those units are summed on their enlarged label
$\operatorname{lab}_1=K_v\cup\mathfrak{h}_v$.

There the anchor is a cube $\Delta\in\pi_k$, and a unit whose enlarged label
$\operatorname{lab}_1$ contains $\Delta$ has its point $y_v$ in a level-$k$ cell $\Delta'$,
lying in zone $k$, in the collar of $\mathsf{C}$ or in $\Omega_k$ next to $\partial\mathsf{C}$,
with $\bar{\mathsf{w}}(\Delta')=3$ there. Since $\operatorname{lab}_1$ is connected,
\[
i:=\operatorname{dist}_k(\Delta',\Delta)\le d_k(\operatorname{lab}_1)+1\le d_k(K_v)+c_l c_{\mathfrak{h}}+1 .
\]
Because $d_k(K_v)\ge 0$ and $d_k(K_v)\ge i-1-c_l c_{\mathfrak{h}}$, this yields
\[
e^{(a+8)d_k(\operatorname{lab}_1)}\le e^{(a+8)c_l c_{\mathfrak{h}}}\,e^{(a+9)d_k(K_v)}\,
e^{-(i-1-c_l c_{\mathfrak{h}})_+}.
\]
The cells $\Delta'$ at distance $i$ number at most $(2i+3)^d$. For each of them, the per-cell
bound for exposed units recalled below, taken at $r=a+9\le\frac85\kappa$ (which holds because
$\kappa\ge 180$), gives
\[
\sum_{y_v\in\Delta'}\mathfrak{N}^{\flat}(v)\,e^{(a+9)d_k(K_v)}\le 27\,\mathsf{K}^{\flat}.
\]

Here the start label $\operatorname{lab}_0$ replaces $\operatorname{lab}_1$; since
$d_k(K_v)\le 2^d-1$, the distance $i$ is bounded by $c_0+1$, and no decaying factor is needed.
The cells $\Delta'$ counted in \eqref{10.3:eq-proof-halo-anchor-a} are level-$k$ cells
of zone $k$: cells of the global partition in their proper scale, which by
\eqref{10.3:eq-never-deep-b} happen to be $\pi_k$-cubes. The count made above is therefore the
one the correlation theorem makes for its exposed pointed units. It is not a count of fine cells
under a coarse cube, which the closing sentence of the per-cell bound excludes; that sentence
concerns the zones $n'<k$ of untreated components, where a $\pi_k$-cube holds $L^{d(k-n')}$
cells of proper scale.

It remains to bound the supremand in \eqref{10.3:eq-proof-halo-anchor-a}.

\emph{The per-cell bound for exposed units.} Let $\varsigma_1\ge 0$ be the rate of the radius
weight of the correlation theorem. The following statement belongs to the proof of the
correlation theorem, and is used here as stated there. Consider every cell $\Delta'$
that is either a cell of the current global partition off $\mathsf{C}$ and the treated components
(its level $n'$ being that of its zone), or a cell of the collar of $\mathsf{C}$, and that
contains the point $y_v$ of an exposed unit. For every such $\Delta'$ and every
$r\le\frac85\kappa$,
\begin{equation*}
\sum_{v\ \text{exposed}:\ y_v\in\Delta'}\mathfrak{N}^{\flat}(v)\,e^{r\,d_k(K_v)}\,
e^{\varsigma_1\operatorname{rad}(v)}\le\mathsf{K}^{\flat}\,\bar{\mathsf{w}}(\Delta')^3,
\qquad \bar{\mathsf{w}}(\Delta'):=3\,w^{\vee}(\Delta').
\end{equation*}
Here $\bar{\mathsf{w}}(\Delta')=3$ unless $\Delta'$ neighbours the window of a foreign live
arrest. Where it does, $\bar{\mathsf{w}}(\Delta')$ is bounded by the window weight of that arrest.

The statement has a second clause, for the pieces carrying the transport parameter
$\mathsf{t}$, which is not needed below. It is a bound per cell of the global partition in its
proper scale. The passage from cells to one $\pi_k$-cube is made, in the correlation theorem, by
a sum over proper-scale cells and never by a count of level-$k$ cells; the count in
\eqref{10.3:eq-proof-halo-anchor-a} is consistent with this for the reason given above.

Its constant $\mathsf{K}^{\flat}$ is, like that of the parent bound for the pointed units in
cells $\Delta'\subset\mathsf{C}$, a number independent of the coupling once $M$ is fixed; every
scale sum entering it runs over levels $j\le n'$ at fixed $n'$, so that it depends neither on
$K$ nor on $k$.

\emph{The hypotheses of the per-cell bound for the units of \eqref{10.3:eq-proof-halo-anchor-a}.}
The units in question are the $v\in\mathfrak{L}^{V,\mathrm{col}}$ with $y_v\in\Delta'$. We check
each hypothesis in turn.
\begin{itemize}
\item The cell. $\Delta'=\Delta_{y_v}$ is a level-$k$ cell of the collar of $\mathsf{C}$ by
\eqref{10.3:eq-never-deep-b}, so $n'=k$.
\item Exposure. The unit $v$ is exposed, in the sense of the list of exposed units of the
correlation theorem as extended to $\mathfrak{L}^{V,\mathrm{col}}$ by the declaration among the
hypotheses of Theorem~\ref{10.3:thm-interior-piece-bound}.
\item The rate. The condition $r\le\frac85\kappa$ at $r=a+9$ reads
$\frac{31}{20}\kappa+9\le\frac85\kappa$, which is equivalent to $\kappa\ge 180$. This is a
hypothesis of Theorem~\ref{10.3:thm-interior-piece-bound}.
\item The weight. $\bar{\mathsf{w}}(\Delta')=3$, because no window of a foreign live arrest
neighbours a cell of the collar of $\mathsf{C}$: this follows from
\eqref{10.3:eq-never-deep-b}, and it is the value $\bar{\mathsf{w}}(\Delta')=3$ used in the last
step of the correlation theorem recalled above.
\item The size. $\mathfrak{N}^{\flat}(v)$ is the size of $v$ fixed in the hypotheses of
Theorem~\ref{10.3:thm-interior-piece-bound}, where it bounds the deviation of $f_v$ on
$\mathfrak{P}$ from its value at the base point.
\item The label. $K_v:=[\square_0]_k$ is the label of $v$.
\item The radius. $\operatorname{rad}(v):=\max_{x\in\hat{X}_v}\operatorname{dist}_k(x,y_v)$ is its
radius, read on the hull $\hat{X}_v=Q_v$, by \eqref{10.3:eq-core-radius-a}.
\end{itemize}

The weighted clause, for the pieces carrying $\mathsf{t}$, is not used: it is not called for
under the declaration among the hypotheses of Theorem~\ref{10.3:thm-interior-piece-bound}.

The supremand of \eqref{10.3:eq-proof-halo-anchor-a} carries no radius factor. It is bounded
directly in the proof of the per-cell bound for vector cores. When the radius
factor is restored through $e^{\varsigma_1\operatorname{rad}(v)}\le e^{\varsigma_1}$, as in
\eqref{10.3:eq-proof-per-cell-a}, these units also satisfy the per-cell bound for exposed units
in its own form.
\end{proof}

\begin{proposition}[The per-cell bound for vector cores]\label{10.3:prf-proof-per-cell}
Let $\Delta'$ be a level-$k$ cell of the collar of $C$; it is a cube of $\pi_k$, and its zone is
$n=n'=k$. Let $r\le \tfrac85\kappa$, and work on the standing domain
\begin{equation*}
\mathbb{D}_V=\bigl\{\,|\lambda_{n,k'}|\,\sup|\mathsf{h}|\,T^V_n\le E_1\,\bigr\}
\end{equation*}
of Notation~\ref{10.3:not-vector-dial-table}, imposed at every index $(n,k')$ of the run.
Let the coupling ceiling satisfy $\gamma\le\gamma_\rho$, where $\gamma_\rho>0$ is the one-sided
smallness of the coupling under which \eqref{10.3:eq-vector-dial-table-b} holds: it depends only
on $c_2$, $q$, $a^*$ and the slope constant of the comparison inequality of the running coupling,
it is chosen after $L$, and it ensures that the run ratios at the fixed gap $n-p^*=a^*$ satisfy
$\rho_{p^*,j}\le 2^{1/2}\rho_{n,j}$, where $\rho_{i,j}$ denotes the ratio of the run's radii at
levels $i$ and $j$, the ratio carried by row~$1$ of the table of
Notation~\ref{10.3:not-vector-dial-table}.
Then the units $v$ of the inner-collar class $\mathfrak{L}^{V,\mathrm{col}}$ anchored in $\Delta'$
satisfy
\begin{equation}\label{10.3:eq-proof-per-cell-1}
\begin{split}
&\sum_{v\in\mathfrak{L}^{V,\mathrm{col}}:\ y_v\in\Delta'}\mathfrak{N}^{\flat}(v)\,
e^{r\,d_k(K_v)}\,e^{\varsigma_1\operatorname{rad}(v)}\\
&\qquad\le e^{(8/5)\kappa\,6^d+5\varsigma_1}\,M^4\,F'_{\mathrm{col}}\,\bar{C}_L^V\,E_1\,
\bar{\mathsf{w}}(\Delta')^3 ,
\end{split}
\end{equation}
where $\bar{\mathsf{w}}(\Delta')=3$ is the weight of the cell, $\varsigma_1\ge 0$ is the radius
rate in the per-cell bound of the correlation theorem, $F'_{\mathrm{col}}:=2^{3/2}F_{\mathrm{col}}$
is the constant of \eqref{10.3:eq-vector-dial-table-b}, with $F_{\mathrm{col}}$ as there,
$\bar{C}_L^V:=\max\{C_\sharp,C_L^V\}$, and $C_\sharp$,
$C_L^V$, $E_1$ are as in Notation~\ref{10.3:not-vector-dial-table}.
\end{proposition}

\begin{proof}
The per-cell bound for exposed units in the proof of the correlation theorem rests on eight
ingredients, each a property of the units summed over. We take them one by one for the units
$v\in\mathfrak{L}^{V,\mathrm{col}}$ with $y_v\in\Delta'$, and state what each requires of such a
unit. Throughout, $j$ is the level of $v$, $n=n'=k$ is the zone of $\Delta'$, $t=L^{j-n}$, and
$\nu=1+n-j$ is the age in zone $n$ of a level-$j$ core. The quantities $\mathsf{u}_j$ are norms,
hence nonnegative, and the per-cell scale weights of row~$1$ of the table of
Notation~\ref{10.3:not-vector-dial-table} are, by \eqref{10.3:eq-vector-dial-table-a},
\begin{equation*}
\varsigma^V_\nu=\nu^3\rho_{n,j}^{2}(1+\rho_{n,j})\,L^{-2\hat{\alpha}(\nu-1)},\qquad j=n+1-\nu,
\end{equation*}
which are nonnegative since the run ratio $\rho_{n,j}$ is positive.

\emph{The anchor count.} In the correlation theorem the proof of the per-cell bound is local at
the anchor point $y$ and at the zone of $y$, and uses that one cell of zone $n'$ contains
$M^4t^{-4}$ anchor points of level $j$. This ingredient depends only on $y_v\in\Delta'$ and the
level $j$ of the
unit: it is the count of anchor points of level-$j$ units in one cell of zone $n'$. Row~$1$ of
the table of Notation~\ref{10.3:not-vector-dial-table}, the row of the $\sharp$-cores of
level $j$ at $y$, anchors the vector cores at their points $y$, one core for each level $j$ and
each point $y$, and these points are the points $y_v$ of the
units of $\mathfrak{L}^{V,\mathrm{col}}$. Its column ``per unit $n$-cell'',
\eqref{10.3:eq-vector-dial-table-a}, is the total of that row's sizes, per unit of
$\sup|\mathsf{h}|$, over the units anchored in one cell of unit side at level $n$. Such a cell
holds $t^{-4}$ anchor points of level $j$, so that the level-$j$ entry of the column is
\begin{equation*}
t^{-4}\cdot C_\sharp\,\mathsf{u}_j\,\nu^3\rho_{n,j}^{2}(1+\rho_{n,j})\,t^{4+2\hat{\alpha}}
=C_\sharp\,\mathsf{u}_j\,\varsigma^V_\nu ,
\end{equation*}
where $C_\sharp\,\mathsf{u}_j\,\nu^3\rho_{n,j}^{2}(1+\rho_{n,j})\,t^{4+2\hat{\alpha}}$ is the
per-unit entry of row~$1$; the cell $\Delta'$, of side $M$, holds $M^4$ such cells. The anchor
count passes from the cell of unit side to $\Delta'$: since $t=L^{-(\nu-1)}$,
\begin{equation*}
\begin{split}
M^4t^{-4}\cdot C_\sharp\,\mathsf{u}_j\,\nu^3\rho_{n,j}^{2}(1+\rho_{n,j})\,t^{4+2\hat{\alpha}}
&=C_\sharp\,\mathsf{u}_j\,M^4\nu^3\rho_{n,j}^{2}(1+\rho_{n,j})\,L^{-2\hat{\alpha}(\nu-1)}\\
&=M^4\,C_\sharp\,\mathsf{u}_j\,\varsigma^V_\nu .
\end{split}
\end{equation*}
This identity is used in the size series below.

\emph{Row sums around the anchor.} The proof of the per-cell bound in the correlation theorem
sums, around $y$, four further rows of that theorem's table: row~$2$, row~$3$, the
boundary row and row~$5$, the last being free of $K$. This ingredient is not used for any
core. Moreover, the table of Notation~\ref{10.3:not-vector-dial-table} has no row~$2$, since
its units are frozen point by point; no row~$5$, since its units carry no weight slot, the
weight exponent being zero for them; and no
boundary row, by the boundary lemma for the vector cores, as recorded with that table in
Notation~\ref{10.3:not-vector-dial-table}. This
ingredient is therefore void for the units of $\mathfrak{L}^{V,\mathrm{col}}$.

\emph{Label and radius.} The correlation theorem bounds the label and the radius of cores and of
the units it calls Wilson pieces by $d_k(K_v)\le 6^d$ and $\operatorname{rad}\le 5$. This
ingredient depends only on the label
$K_v=[X_v]_k$ and the hull $\hat{X}_v$. For $v\in\mathfrak{L}^{V,\mathrm{col}}$ the core
$X_v=\square_0$ is kept. Its side is $(2\nu+3)L^{-(\nu-1)}<1$ times the side of a cube of
$\pi_k$, so by
Lemma~\ref{10.3:lem-core-geometry}, display \eqref{10.3:eq-core-geometry-b},
\begin{equation*}
d_k(K_v)\le 2^d-1\le 6^d .
\end{equation*}
The radius is
\begin{equation*}
\operatorname{rad}(v)=\max_{x\in\hat{X}_v}\operatorname{dist}_n(x,y_v),
\end{equation*}
computed on the hull $\hat{X}_v=Q_v$, which, by the definition of the class, is the hull the
correlation theorem assigns to a smooth core, unchanged for the vector cores. By
Lemma~\ref{10.3:lem-core-radius}, display
\eqref{10.3:eq-core-radius-a}, the box
$Q_v$ has $n$-diameter at most $1/L$. Hence
\begin{equation*}
\operatorname{rad}(v)\le \tfrac1L<1\le 5 .
\end{equation*}
The bound $\operatorname{rad}\le 5$ of the correlation theorem is the $n$-diameter of the hull
$\square^{\sim\nu+1}$ of a rough core at $\nu=1$, and it covers $v$ a fortiori. The support
$\mathfrak{r}_v$ of $v$ has $k$-diameter at most $5$ by \eqref{10.3:eq-core-radius-a}. The
radius is not computed on the support. Nor does the support column of the table of
Notation~\ref{10.3:not-vector-dial-table} enter that table's per-cell column, which totals by
anchor. It follows that for every $r\le \frac85\kappa$, in particular for $r=a+9$, the rate at
which the correlation theorem uses the bound (the inequality $a+9\le\frac85\kappa$ is one of the
parameter conditions of this section),
\begin{equation*}
e^{r\,d_k(K_v)}\le e^{(8/5)\kappa\,6^d}.
\end{equation*}
Since $\varsigma_1\ge 0$ and $L\ge 1$,
\begin{equation*}
e^{\varsigma_1\operatorname{rad}(v)}\le e^{\varsigma_1/L}\le e^{\varsigma_1}\le e^{5\varsigma_1}.
\end{equation*}

\emph{The retained rate.} The correlation theorem keeps in every decaying row a rate at least
$(\tfrac45-\tfrac12\delta)\kappa$ unspent, with $\delta$ as in that theorem. A core spends no
rate, because its label and radius
factors are bounded by the constants just obtained. This ingredient is not used.

\emph{The rim case.} For the far cores of the correlation theorem the rim case is void, because
their scale rule is $p_v=n-2-s^{(9/8)}$ always, with no clause ``$Q\subset Z$''. For
$v\in\mathfrak{L}^{V,\mathrm{col}}$ the
scale rule is
\begin{equation*}
p_v:=p^*=n-2-s(3),
\end{equation*}
taken without dilation, by the convention on the integer function $s$ fixed for the inner-collar
class, and with no clause ``$Q^{(p)}\subset Z$''. The rim alternative of the
correlation theorem's proof is therefore void for $v$, for the same reason.

\emph{The size series.} The correlation theorem uses the smooth-core series
\begin{equation*}
\sum_b (b+1)^3\,\bar{\rho}_b^{\,2}\,(1+\bar{\rho}_b)\,L^{-\hat{\alpha}b},
\end{equation*}
together with a bound for rough cores; here $\bar{\rho}_b$ is that theorem's majorant of the run
ratio, indexed by the gap $b$ and independent of the coupling, and it is a different quantity from
the run ratio $\rho_{n,j}$ carried by row~$1$ of the table of
Notation~\ref{10.3:not-vector-dial-table}. This is the class's per-cell size series: the size of a
unit, summed over the anchors of the anchor count and over the age. For the vector cores the size
of a unit is $\mathfrak{N}^{\flat}(v)$, the size of the unit $\lambda_{n,k'}f_v$, that is
$|\lambda_{n,k'}|$ times a size of the core $f_v$, defined at
\eqref{10.3:eq-core-size-bound-a}. In \eqref{10.3:eq-vector-dial-table-b} this size is read at the
scale $p^*$ of the core, where it carries the run ratio $\rho_{p^*,j}$, and is compared with the
run ratio at level $n$ through $\rho_{p^*,j}\le 2^{1/2}\rho_{n,j}$, valid for
$\gamma\le\gamma_\rho$. In this form \eqref{10.3:eq-vector-dial-table-b} gives
\begin{equation}\label{10.3:eq-proof-per-cell-2}
\begin{split}
\mathfrak{N}^{\flat}(v)
&\le |\lambda_{n,k'}|\,\sup|\mathsf{h}|\,F'_{\mathrm{col}}\,
C_\sharp\,\mathsf{u}_j\,\nu^3\,\rho^{2}_{n,j}\,(1+\rho_{n,j})\,t^{4+2\hat{\alpha}}\\
&=|\lambda_{n,k'}|\,\sup|\mathsf{h}|\,F'_{\mathrm{col}}\,\mathfrak{N}(v),
\end{split}
\end{equation}
where $\mathfrak{N}(v)$ denotes the per-unit entry of row~$1$, the product following
$F'_{\mathrm{col}}$ in the first line of \eqref{10.3:eq-proof-per-cell-2}. The units of
$\mathfrak{L}^{V,\mathrm{col}}$ of level $j$ with $y_v\in\Delta'$ are among the row-$1$ units of
level $j$ anchored at the anchor points of level $j$ in $\Delta'$, one core for each such point;
by the anchor count there are at most $M^4t^{-4}$ of them, and each obeys
\eqref{10.3:eq-proof-per-cell-2} with the same right member, which does not depend on the anchor.
Hence, by the identity of the anchor count, the total of these sizes over the units of level $j$
anchored in $\Delta'$ is
\begin{equation*}
\begin{split}
\sum_{\substack{v\in\mathfrak{L}^{V,\mathrm{col}}:\ y_v\in\Delta'\\ v\text{ of level }j}}
\mathfrak{N}^{\flat}(v)
&\le M^4t^{-4}\,|\lambda_{n,k'}|\,\sup|\mathsf{h}|\,F'_{\mathrm{col}}\,
C_\sharp\,\mathsf{u}_j\,\nu^3\rho^{2}_{n,j}(1+\rho_{n,j})\,t^{4+2\hat{\alpha}}\\
&=M^4\,|\lambda_{n,k'}|\,\sup|\mathsf{h}|\,F'_{\mathrm{col}}\,C_\sharp\,\mathsf{u}_j\,
\varsigma^V_\nu,\qquad \nu=1+n-j .
\end{split}
\end{equation*}
Rough cores and young cores do not belong to $\mathfrak{L}^{V,\mathrm{col}}$, so no bound for
them is needed. This ingredient therefore holds, with the correlation theorem's series replaced by
the per-cell column of row~$1$. Its closure over $j$ is the next ingredient.

\emph{The scale sum.} In the correlation theorem every scale sum runs over $j\le n'$ at fixed
$n'$ and is free of $K$ and of $k$; for that theorem's cores it is the geometric series in $b$.
For the vector cores we sum the per-cell column of row~$1$ over $1\le j\le n'=n=k$. By
\eqref{10.3:eq-vector-dial-table-a}, the level-$j$ term $C_\sharp\mathsf{u}_j\varsigma^V_{1+n-j}$
is $C_\sharp$ times the $j$-th summand of the first sum
$\sum_{0\le j\le n}\varsigma^V_{1+n-j}\mathsf{u}_j$ in the load $T^V_n$. All summands of that
first sum, and the remaining terms of $T^V_n$ in \eqref{10.3:eq-vector-dial-table-a}, are
nonnegative, so the partial sum over $1\le j\le n$ is at most the first sum, hence at most
$T^V_n$; multiplied by $C_\sharp$ it is at most $C_\sharp T^V_n$. Since
$\bar{C}_L^V=\max\{C_\sharp,C_L^V\}\ge C_\sharp$ and $T^V_n\ge 0$,
\begin{equation}\label{10.3:eq-proof-per-cell-3}
\sum_{1\le j\le n}C_\sharp\,\mathsf{u}_j\,\varsigma^V_{1+n-j}\le C_\sharp\,T^V_n
\le \bar{C}_L^V\,T^V_n .
\end{equation}
The remaining factor obeys
\begin{equation*}
|\lambda_{n,k'}|\,\sup|\mathsf{h}|\,T^V_n\le E_1 \quad\text{on }\mathbb{D}_V,
\end{equation*}
by the definition of that set. The right member $\bar{C}_L^VE_1$ involves none of $K$, $k$, $j$,
$n$, the age or the run. Indeed, the constant $C_L^{V\prime}:=C_L^VE_1$ is fixed in
Notation~\ref{10.3:not-vector-dial-table} before the quantification over the cutoff and over the
level $k$ below it, and the rows of the table hold at every cutoff and every level below it with
the same constants. So
the scale
sum closes for these units on $\mathbb{D}_V$ with a right member free of $K$. The closing step
is the defining inequality of $\mathbb{D}_V$, which takes the place of the factor
$L^{-\hat{\alpha}b}$ of the correlation theorem.

What this establishes is the per-cell inequality on the standing domain $\mathbb{D}_V$. That
domain is the one on which every bound of the expansion with the vector part $\mathfrak{u}_k$
added is stated, and only the definition of $\mathbb{D}_V$ is used here. The argument does not
show that $\mathbb{D}_V$ has a size uniform in $n$. That is the content of the statement
displayed at \eqref{10.3:eq-proof-collar-sum-a}.

\emph{The polynomial and the cell weight.} In the correlation theorem a polynomial in
$s^{(9/8)}\le s(1)+\log_L\bar{\mathsf{w}}+1$ is bounded by a constant multiple of
$\bar{\mathsf{w}}$, which produces
the third power of $\bar{\mathsf{w}}(\Delta')$. For the vector cores the polynomial in $s$ enters
through the factor $(2a^*+11)^2$ of $F'_{\mathrm{col}}$, at $a^*=2+s(3)$. Since $s$ is taken
without dilation, as in the rim case above, and $\bar{\mathsf{w}}(\Delta')=3$, this factor is a
constant inside $F'_{\mathrm{col}}$. Moreover
$\bar{\mathsf{w}}(\Delta')^3=27\ge 1$, so multiplying a nonnegative constant by
$\bar{\mathsf{w}}(\Delta')^3$ gives a majorant of it. The weighted clause of the correlation
theorem's per-cell bound
is not used.

\emph{Conclusion.} Collecting the label and radius bounds, the per-cell size bound and the closure
of the scale sum, we have, for every $v\in\mathfrak{L}^{V,\mathrm{col}}$ with $y_v\in\Delta'$
(first line), for every level $j$ (second line), and, in the last line, on $\mathbb{D}_V$:
\begin{equation}\label{10.3:eq-proof-per-cell-a}
\begin{gathered}
d_k(K_v)\le 6^d,\qquad \operatorname{rad}(v)\le \tfrac1L\le 5,\\
\sum_{\substack{v\in\mathfrak{L}^{V,\mathrm{col}}:\ y_v\in\Delta'\\ v\text{ of level }j}}
\mathfrak{N}^{\flat}(v)\le M^4\,|\lambda_{n,k'}|\,\sup|\mathsf{h}|\,F'_{\mathrm{col}}\,
C_\sharp\,\mathsf{u}_j\,\varsigma^V_{1+n-j},\\
\sum_{1\le j\le n}C_\sharp\,\mathsf{u}_j\,\varsigma^V_{1+n-j}\le \bar{C}_L^V\,T^V_n,\qquad
|\lambda_{n,k'}|\,\sup|\mathsf{h}|\,T^V_n\le E_1\ \text{ on }\mathbb{D}_V .
\end{gathered}
\end{equation}
By the first line of \eqref{10.3:eq-proof-per-cell-a}, every summand on the left of
\eqref{10.3:eq-proof-per-cell-1} is at most
$e^{(8/5)\kappa\,6^d+5\varsigma_1}\mathfrak{N}^{\flat}(v)$.
Summing the second line over $1\le j\le n$ and applying the third gives
\begin{equation*}
\begin{split}
\sum_{v\in\mathfrak{L}^{V,\mathrm{col}}:\ y_v\in\Delta'}\mathfrak{N}^{\flat}(v)\,
e^{r\,d_k(K_v)}\,e^{\varsigma_1\operatorname{rad}(v)}
&\le e^{(8/5)\kappa\,6^d+5\varsigma_1}\,M^4\,F'_{\mathrm{col}}\,\bar{C}_L^V\,
|\lambda_{n,k'}|\,\sup|\mathsf{h}|\,T^V_n\\
&\le e^{(8/5)\kappa\,6^d+5\varsigma_1}\,M^4\,F'_{\mathrm{col}}\,\bar{C}_L^V\,E_1\\
&\le e^{(8/5)\kappa\,6^d+5\varsigma_1}\,M^4\,F'_{\mathrm{col}}\,\bar{C}_L^V\,E_1\,
\bar{\mathsf{w}}(\Delta')^3 ,
\end{split}
\end{equation*}
which is \eqref{10.3:eq-proof-per-cell-1}. Since the radius bound also gives
$e^{\varsigma_1\operatorname{rad}(v)}\le e^{\varsigma_1}$, the same chain with $e^{\varsigma_1}$ in
place of $e^{5\varsigma_1}$ shows that the left member of
\eqref{10.3:eq-proof-per-cell-1} is at most
$M^4e^{(8/5)\kappa\,6^d+\varsigma_1}F'_{\mathrm{col}}\bar{C}_L^VE_1\,\bar{\mathsf{w}}(\Delta')^3$,
hence at most $\mathsf{K}^{\flat,V}\bar{\mathsf{w}}(\Delta')^3$, where the per-cell collar
constant of the vector class is
\begin{align*}
\mathsf{K}^{\flat,V}&:=M^4\,e^{(8/5)\kappa\,6^d+\varsigma_1}\,F^{*\prime}_{\mathrm{col}}\,
\bar{C}_L^V\,E_1,\\
F^{*\prime}_{\mathrm{col}}&:=\max\Bigl\{F'_{\mathrm{col}},\ \tfrac{27}{4}\,C^{\mathrm{mar}}\,
\bar{\rho}_{a^*}^{\,2}\,L^{2\hat{\alpha}(a^*-1)}/C_\sharp\Bigr\}\ \ge F'_{\mathrm{col}} ;
\end{align*}
here $C^{\mathrm{mar}}$ is the constant of the marginal bound for young cores, so that the second
member of the maximum is the constant of the young cores, which do not belong to
$\mathfrak{L}^{V,\mathrm{col}}$ and enter here only through $F^{*\prime}_{\mathrm{col}}\ge
F'_{\mathrm{col}}$.

It remains to check that no ingredient of the correlation theorem's proof uses a property the
vector cores lack. The anchor count, the label and radius bounds, and the rim case depend only on
$y_v\in\Delta'$ together with the level of $\Delta'$, the label and the hull, and the scale rule
for $p_v$; the vector cores have all of these in the same form. The row sums around the anchor,
the retained rate and the weighted clause are not used for any core. The size series, the scale
sum and the polynomial use the per-cell count and size series of the class. For the vector cores
these are the per-cell column of row~$1$ of the table of
Notation~\ref{10.3:not-vector-dial-table}, summed over $j$ by the definition of $T^V_n$ and
closed by $\mathbb{D}_V$.
\end{proof}

\begin{proposition}[The sum over a collar cell]\label{10.3:prf-proof-collar-sum}
Use the dial $\lambda_{n,k'}$, the norms $\mathsf{u}_j$ and $\mathsf{r}_i$, the raw term
$\mathsf{t}^{\mathrm{raw}}$, the per-cell scale weights $\varsigma^V_\nu$ and $w_{n,i}$, the rate
windows $\mathsf{a}_i$, the loads $T^V_n$, the constants $C_{\sharp}$, $C_L^V$ and
$C_L^{V\prime}=C_L^VE_1$, and the domain
\begin{equation*}
\mathbb{D}_V=\bigl\{|\lambda_{n,k'}|\,\sup|\mathsf{h}|\,T^V_n\le E_1\bigr\}
\end{equation*}
of the vector package, all as fixed in Notation~\ref{10.3:not-vector-dial-table}. Let
$\mathfrak{L}^{V,\mathrm{col}}$ be the inner-collar class of the vector package; a member $v$
has a level $\operatorname{lev}v$, an anchor point $y_v$, a size $\mathfrak{N}^{\flat}(v)$, the
set $K_v=[\square_0]_k$, and the radius $\operatorname{rad}(v)$ of
Lemma~\ref{10.3:lem-core-radius}. Put $a=\tfrac{31}{20}\kappa$, $c_l=1+234d$,
$c_{\mathfrak{h}}=7^d$, $c_0=6^d+c_lc_{\mathfrak{h}}$, $s^*=s(3)$ and $a^*=2+s^*$, where
$s(\cdot)$ is the integer function of the correlation theorem, whose definition makes the gauge
smallness of a core at cell weight $3$ and depth $a^*$ at most $1$. Let
$\varsigma_1$ and $\bar{\mathsf{w}}(\Delta')$ be the radius weight and the cell weight of the
per-cell collar bound of the correlation theorem, let $\mathsf{K}^{\flat}_J$ be that
bound's per-cell constant, and let $F_{\mathrm{col}}$ be the scale-reading constant displayed in
\eqref{10.3:eq-proof-collar-sum-2} below.

Assume the hypotheses under which Theorem~\ref{10.3:thm-interior-piece-bound} is stated, in
particular $\kappa\ge180$, and let $K$, $k\le K$ and the label be arbitrary. Assume in addition
the one-sided ceiling
\begin{equation}\label{10.3:eq-proof-collar-sum-9}
g\le\gamma\le\min\{\gamma^{\circ},\gamma_{\rho}\},\qquad \gamma^{\circ}:=(8c_2+4)^{-1/2},
\end{equation}
on the coupling, where
\begin{equation}\label{10.3:eq-proof-collar-sum-10}
\begin{gathered}
\gamma_{\rho}:=\sup\Bigl\{\gamma'\in\bigl]0,\min\{\gamma^{\circ},(3/4)^{1/2}e^{-q}\}\bigr]:\
E_{\rho}(\gamma')\le 2\Bigr\},\\
E_{\rho}(\gamma'):=(1+4c_2\gamma'^2)\Bigl(1+\frac{(3\lambda_{\mathrm{fl}}a^*+2c_2)\,\gamma'^2}
{\log\gamma'^{-2}}\Bigr)^{2q};
\end{gathered}
\end{equation}
here $q$ is the integer exponent of the logarithm in the radii of the run, and $c_2\ge0$,
$\lambda_{\mathrm{fl}}\ge0$ are the constants of the two-sided comparison
\begin{equation*}
\lambda_{\mathrm{fl}}(j-i)-2c_2\le x_i-x_j\le 3\lambda_{\mathrm{fl}}(j-i)+2c_2,
\qquad 0\le i<j\le K,
\end{equation*}
satisfied by the running couplings $x_i=g_i^{-2}$ of the correlation theorem, which also satisfy
$x_i>g^{-2}$ for $i\le k$. Let the zone be
$n=k\ge 4$, let $\Delta'\in\pi_n$ be a cell of level $k$ of the collar of the region $C$ over
whose collar cells the supremum in \eqref{10.3:eq-proof-halo-anchor-a} is taken, and let
$(\lambda,\mathsf{h})\in\mathbb{D}_V$. Such a cell carries the cell weight
$\bar{\mathsf{w}}(\Delta')=3$ of the per-cell collar bound of the correlation theorem.
Define
\begin{equation}\label{10.3:eq-proof-collar-sum-1}
\bar{C}_L^V:=\max\{C_{\sharp},C_L^V\},\qquad
\mathsf{K}^{\flat,V}:=M^4\,e^{\frac85\kappa\cdot 6^d+\varsigma_1}\,F^{*\prime}_{\mathrm{col}}\,
\bar{C}_L^V\,E_1,
\end{equation}
where
\begin{equation}\label{10.3:eq-proof-collar-sum-2}
\begin{aligned}
F'_{\mathrm{col}}&:=2^{3/2}\,\frac{36c_{\epsilon}^2(2a^*+11)^2\bar{\rho}^2_{a^*}
L^{4+2\hat{\alpha}a^*}C_V}{C_{\sharp}}=2^{3/2}F_{\mathrm{col}},\\
F^{*\prime}_{\mathrm{col}}&:=\max\Bigl\{F'_{\mathrm{col}},\
\frac{27\,C^{\mathrm{mar}}\bar{\rho}^2_{a^*}L^{2\hat{\alpha}(a^*-1)}}{4\,C_{\sharp}}\Bigr\};
\end{aligned}
\end{equation}
in particular $F^{*\prime}_{\mathrm{col}}\ge F'_{\mathrm{col}}$. Then the following hold.
\begin{enumerate}
\item[(a)] The per-cell bound
\begin{equation}\label{10.3:eq-proof-collar-sum-3}
\sum_{\substack{v\in\mathfrak{L}^{V,\mathrm{col}}\\ y_v\in\Delta'}}
\mathfrak{N}^{\flat}(v)\,e^{(a+9)d_k(K_v)}\,e^{\varsigma_1\operatorname{rad}(v)}
\le\mathsf{K}^{\flat,V}\le\mathsf{K}^{\flat,V}\,\bar{\mathsf{w}}(\Delta')^3 .
\end{equation}
\item[(b)] The number $\mathsf{K}^{\flat,V}$ is independent of $K$.
\item[(c)] With $\mathsf{K}^{\flat}_+:=\mathsf{K}^{\flat}_J+\mathsf{K}^{\flat,V}$, the anchored
weighted reader sum $\mathsf{Q}^V_{a+9}$ of \eqref{10.3:eq-proof-halo-anchor-a} satisfies
\begin{equation}\label{10.3:eq-proof-collar-sum-4}
\mathsf{Q}^V_{a+9}\le e^{(a+9)c_lc_{\mathfrak{h}}}(c_0+2)(2c_0+5)^d\cdot 27\,\mathsf{K}^{\flat}_+
=:\mathsf{K}^{\flat,\mathrm{ex}}_{\mathrm{pl}} .
\end{equation}
\end{enumerate}
\end{proposition}

\begin{proof}
Throughout, $\Delta'$, $n=k\ge 4$ and $(\lambda,\mathsf{h})\in\mathbb{D}_V$ are fixed as in the
statement. We recall from Notation~\ref{10.3:not-vector-dial-table} the definition of the load
at zone $n$:
\begin{equation}\label{10.3:eq-proof-collar-sum-5}
T^V_n=\sum_{0\le j\le n}\varsigma^V_{1+n-j}\,\mathsf{u}_j
+\sum_{n-\mathsf{a}_i\le i\le n}w_{n,i}\,\mathsf{r}_i+\mathbf{1}_{n=0}\,\mathsf{t}^{\mathrm{raw}} .
\end{equation}

\emph{Step 1 (geometry).} By the geometric bounds \eqref{10.3:eq-proof-per-cell-a} for vector
cores, for every $v\in\mathfrak{L}^{V,\mathrm{col}}$ we have
\begin{equation*}
e^{(a+9)d_k(K_v)}\le e^{(a+9)\cdot 6^d}\le e^{\frac85\kappa\cdot 6^d},
\qquad e^{\varsigma_1\operatorname{rad}(v)}\le e^{\varsigma_1}.
\end{equation*}
The last inequality of the first chain is $a+9\le\frac85\kappa$; since $a=\frac{31}{20}\kappa$ and
$\frac85-\frac{31}{20}=\frac1{20}$, it is the condition $\kappa\ge 180$, which is among the
parameter conditions under which Theorem~\ref{10.3:thm-interior-piece-bound} is stated. We note
also that $r=a+9$ is thereby an admissible exponent $r\le\frac85\kappa$ in the per-cell collar
bound of the correlation theorem.

\emph{Step 2 (one level $j$).} By \eqref{10.3:eq-proof-per-cell-a}, the members of
$\mathfrak{L}^{V,\mathrm{col}}$ of level $j$ with
$y_v\in\Delta'$ are units of the first row of the table of
Notation~\ref{10.3:not-vector-dial-table}, that is $\sharp$-cores of level $j$, anchored in the
$n$-cell $\Delta'$. The hull of such a unit is a cube of $2\nu+3$ level-$j$ blocks of side $M$ in
each direction, whose side in level-$n$ units is $(2\nu+3)ML^{-(n-j)}$; the activity row of the
table bounds this side by $(2\nu+3)ML^{-(\nu-1)}$, whence $n-j=\nu-1$, and the age of these
units in zone $n$ is $\nu=1+n-j$. The entry of the first row in the
last column of that table (the column headed ``per unit $n$-cell'') is
$C_{\sharp}\mathsf{u}_j\varsigma^V_\nu$
and nothing else. Its entry in the size column is, per unit of $\sup|\mathsf{h}|$,
$C_{\sharp}\mathsf{u}_j\nu^3\rho^2_{n,j}(1+\rho_{n,j})t^{4+2\hat{\alpha}}$, where
$t:=L^{j-n}=L^{-(\nu-1)}$ and $\rho_{n,j}$ is the ratio of the radii of the run at levels $n$
and $j$; the weight itself is
\begin{equation*}
\varsigma^V_\nu=\nu^3\rho^2_{n,j}(1+\rho_{n,j})\,L^{-2\hat{\alpha}(\nu-1)},\qquad j=n+1-\nu .
\end{equation*}
By \eqref{10.3:eq-vector-dial-table-b}, which holds under the ceiling
\eqref{10.3:eq-proof-collar-sum-9}, every member $v\in\mathfrak{L}^{V,\mathrm{col}}$ of level $j$
with $y_v\in\Delta'$ satisfies
\begin{equation*}
\mathfrak{N}^{\flat}(v)\le|\lambda_{n,k'}|\,\sup|\mathsf{h}|\,F'_{\mathrm{col}}\,
C_{\sharp}\,\mathsf{u}_j\,\nu^3\rho^2_{n,j}(1+\rho_{n,j})\,t^{4+2\hat{\alpha}},
\end{equation*}
and this right member does not depend on $y_v$. The members of
$\mathfrak{L}^{V,\mathrm{col}}$ of level $j$ with $y_v\in\Delta'$ are among the first-row units
of level $j$ anchored at the points of level $j$ in $\Delta'$, with one $\sharp$-core for each
such point; as in the proof of the per-cell collar bound of the correlation theorem, the cell
$\Delta'$ holds $M^4t^{-4}$ points of level $j$ ($d=4$). Hence, since
$t^{-4}t^{4+2\hat{\alpha}}=t^{2\hat{\alpha}}=L^{-2\hat{\alpha}(\nu-1)}$ and $\nu=1+n-j$,
\begin{equation}\label{10.3:eq-proof-collar-sum-6}
\begin{aligned}
\sum_{\substack{v\in\mathfrak{L}^{V,\mathrm{col}}:\ \operatorname{lev}v=j\\ y_v\in\Delta'}}
\mathfrak{N}^{\flat}(v)
&\le M^4t^{-4}\,|\lambda_{n,k'}|\,\sup|\mathsf{h}|\,F'_{\mathrm{col}}\,
C_{\sharp}\,\mathsf{u}_j\,\nu^3\rho^2_{n,j}(1+\rho_{n,j})\,t^{4+2\hat{\alpha}}\\
&=M^4\,|\lambda_{n,k'}|\cdot\sup|\mathsf{h}|\cdot F'_{\mathrm{col}}\,
C_{\sharp}\,\mathsf{u}_j\,\varsigma^V_{1+n-j};
\end{aligned}
\end{equation}
the last step is an identity.

\emph{Step 3 (the sum over $j$).} For each $1\le j\le n$, the number
$C_{\sharp}\mathsf{u}_j\varsigma^V_{1+n-j}$ is $C_{\sharp}$ times the $j$-th summand of the first
sum in \eqref{10.3:eq-proof-collar-sum-5}. Hence
\begin{equation}\label{10.3:eq-proof-collar-sum-7}
\begin{aligned}
\sum_{1\le j\le n}C_{\sharp}\,\mathsf{u}_j\,\varsigma^V_{1+n-j}
&\le C_{\sharp}\Bigl(\sum_{0\le j\le n}\varsigma^V_{1+n-j}\mathsf{u}_j
+\sum_{n-\mathsf{a}_i\le i\le n}w_{n,i}\mathsf{r}_i+\mathbf{1}_{n=0}\mathsf{t}^{\mathrm{raw}}\Bigr)\\
&=C_{\sharp}\,T^V_n\le\bar{C}_L^V\,T^V_n .
\end{aligned}
\end{equation}
The first inequality adds the terms $\varsigma^V_{1+n}\mathsf{u}_0$, $w_{n,i}\mathsf{r}_i$ and
$\mathbf{1}_{n=0}\mathsf{t}^{\mathrm{raw}}$, which are nonnegative by the sign conventions
\eqref{10.3:eq-vector-dial-table-a}: $\mathsf{u}_j$ and $\mathsf{r}_i$ are norms,
$\mathsf{t}^{\mathrm{raw}}>0$, and $\varsigma^V_\nu$, $w_{n,i}$ are per-cell scale weights,
nonnegative as majorants of nonnegative sizes. The equality is
\eqref{10.3:eq-proof-collar-sum-5}. The last inequality holds because
$\bar{C}_L^V=\max\{C_{\sharp},C_L^V\}\ge C_{\sharp}$. If $C_L^V$, defined in
Notation~\ref{10.3:not-vector-dial-table} as the largest constant of the last column of the table,
already dominates $C_{\sharp}$, then $\bar{C}_L^V=C_L^V$; whatever weights enter the constants of
that column, the chain \eqref{10.3:eq-proof-collar-sum-7} holds with $C_{\sharp}$, and the maximum
is taken on one side only and does not depend on $K$. No factor of the first row lies outside this
domination, since its per-cell entry is $C_{\sharp}\mathsf{u}_j\varsigma^V_\nu$ alone.

\emph{Step 4 (the domain).} By the definition of $\mathbb{D}_V$ at the index $(n,k')$ of the zone
and run in question, $|\lambda_{n,k'}|\cdot\sup|\mathsf{h}|\cdot T^V_n\le E_1$. Hence
\begin{equation}\label{10.3:eq-proof-collar-sum-8}
|\lambda_{n,k'}|\,\sup|\mathsf{h}|\,\bar{C}_L^V\,T^V_n\le\bar{C}_L^V\,E_1 .
\end{equation}

\emph{Assembly.} Splitting the sum according to the level $j$ of $v$, $1\le j\le n$, and using
Step 1 for the exponential factors, then \eqref{10.3:eq-proof-collar-sum-6} summed over $j$,
then \eqref{10.3:eq-proof-collar-sum-7}, then \eqref{10.3:eq-proof-collar-sum-8}, we obtain
\begin{align*}
&\sum_{\substack{v\in\mathfrak{L}^{V,\mathrm{col}}\\ y_v\in\Delta'}}
\mathfrak{N}^{\flat}(v)\,e^{(a+9)d_k(K_v)}e^{\varsigma_1\operatorname{rad}(v)}
\le e^{\frac85\kappa\cdot6^d+\varsigma_1}\sum_{1\le j\le n}
\sum_{\substack{\operatorname{lev}v=j\\ y_v\in\Delta'}}\mathfrak{N}^{\flat}(v)\\
&\qquad\le e^{\frac85\kappa\cdot6^d+\varsigma_1}M^4F'_{\mathrm{col}}\,|\lambda_{n,k'}|
\sup|\mathsf{h}|\cdot C_{\sharp}\sum_{1\le j\le n}\varsigma^V_{1+n-j}\mathsf{u}_j\\
&\qquad\le e^{\frac85\kappa\cdot6^d+\varsigma_1}M^4F'_{\mathrm{col}}\,|\lambda_{n,k'}|
\sup|\mathsf{h}|\cdot\bar{C}_L^V\,T^V_n\\
&\qquad\le M^4e^{\frac85\kappa\cdot6^d+\varsigma_1}F'_{\mathrm{col}}\,\bar{C}_L^V E_1
\le\mathsf{K}^{\flat,V}\le\mathsf{K}^{\flat,V}\,\bar{\mathsf{w}}(\Delta')^3 .
\end{align*}
The penultimate inequality is $F'_{\mathrm{col}}\le F^{*\prime}_{\mathrm{col}}$ together with the
definition \eqref{10.3:eq-proof-collar-sum-1}, and the last one is $\bar{\mathsf{w}}(\Delta')^3=27\ge1$.
This is \eqref{10.3:eq-proof-collar-sum-3}, and proves (a).

The units of $\mathfrak{L}^{V,\mathrm{col}}$ thus fall under the displayed inequality of the
per-cell collar bound of the correlation theorem at $\bar{\mathsf{w}}=3$ once its constant
$\mathsf{K}^{\flat}_J$ is enlarged to $\mathsf{K}^{\flat}_+=\mathsf{K}^{\flat}_J+\mathsf{K}^{\flat,V}$;
this is the single enlargement of $\mathsf{K}^{\flat}$ by the vector part's smooth-core series
declared among the hypotheses under which Theorem~\ref{10.3:thm-interior-piece-bound} is stated.
Its content is the first
row's per-cell column of the table, summed over the levels by the definition
\eqref{10.3:eq-proof-collar-sum-5} of $T^V_n$, with the constant $C_L^V$, on the domain
$\mathbb{D}_V$ of the typed format statement of the vector package, at the constant
$C_L^{V\prime}=C_L^VE_1$. Nothing is estimated anew, and neither the conclusion of the class
clause of the typed format statement nor the load recursion lemma of the vector package is used.
The other places where the correlation theorem uses $\mathsf{K}^{\flat}$ --- its conditions that read
$\mathsf{K}^{\flat}$ through $\mathsf{K}^{\flat,\mathrm{ex}}$, and the replacement of
$\mathsf{K}^0$ by $\mathsf{K}^0+\mathsf{K}^{\flat,\mathrm{ex}}$ --- read $\mathsf{K}^{\flat}_+$ in
place of $\mathsf{K}^{\flat}_J$; this is an enlargement, on one side only, of numbers fixed after
$M$ in the order of choice of the parameters.

\emph{Step 5 (independence of $K$, and the uniformity of the domain).} The right member
$\mathsf{K}^{\flat,V}$ of \eqref{10.3:eq-proof-collar-sum-3} is independent of $K$, since it
depends only on $M$, $\kappa$, $d$, $\varsigma_1$; $F^{*\prime}_{\mathrm{col}}$, fixed after $L$;
$\bar{C}_L^V$, fixed after $L$; and $E_1$, the radius number in the definition of
$\mathbb{D}_V$. This proves (b).

The domain on which \eqref{10.3:eq-proof-collar-sum-3} holds is $\mathbb{D}_V$, whose radius in
$\lambda$ at zone $n$ is $E_1/(\sup|\mathsf{h}|\,T^V_n)$; its uniformity in $n$ is the statement
that $\sup_nT^V_n$ is finite with a bound independent of $K$. This is supplied by the load
recursion lemma of the vector package, which we recall here and which is not used in
Steps~1--4. It states: under its hypotheses --- among them ceilings on $g$, each one-sided and taken
after $(L,\mathfrak{p},\hat{\alpha},\vartheta_2,q)$, the first of which is
$(\bar{\varkappa}+c')S^V\le\frac18$ and the last $4S^V(\bar{\varkappa}+c')\le\vartheta_V$, where
$\bar{\varkappa}$ and $c'$ are numbers fixed in the vector package ($\bar{\varkappa}$ is unrelated
to the decay rate $\varkappa$ of the two-point witness theorem) --- and with
$\mathsf{A}_V:=4S^V\mathsf{v}_0$, where $\mathsf{v}_0$ is a number fixed in the vector package,
one has, among other bounds, for all $1\le n\le K'$, $K'$ the top index of the
run,
\begin{equation*}
T^V_n\le\mathsf{A}_V\vartheta_V^n,\qquad
\mathsf{u}_n\le\mathsf{o}_n\le(\bar{\varkappa}+c')\mathsf{A}_V\vartheta_V^{n-1}\le
\mathsf{v}_0\vartheta_V^n,
\end{equation*}
together with a bound for $T^V_0$. Since $\vartheta_V=\vartheta_2/L<1$, this gives
\begin{equation}\label{10.3:eq-proof-collar-sum-a}
\sup_{1\le n\le K'}T^V_n\le\mathsf{A}_V\,\vartheta_V\le\mathsf{A}_V ,
\end{equation}
a number that depends only on $S^V$ and $\mathsf{v}_0$ and is independent of $K$ once
$(L,\mathfrak{p},\hat{\alpha},\vartheta_2,q)$ are fixed. Here $S^V$ is the scale series
$S^V=\sum_{\nu\ge1}\varsigma^V_\nu\vartheta_V^{-\nu}$, for which $\varsigma^V_1=2$ gives
$S^V\ge 2/\vartheta_V$. It depends only on $\vartheta_V$ and on the weights
$(\varsigma^V_\nu)_\nu$, which read $L$, $\hat{\alpha}$ and the run ratios $\rho_{n,j}$; by the
finiteness clause of the vector package for these weights, which rests on the bound
\begin{equation*}
\rho^2_{n,n+1-\nu}\le(1+o(1))\bigl(1+(3\lambda_{\mathrm{fl}}\nu+2c_2)g^2\bigr),
\end{equation*}
the series is finite for every $\vartheta_2>L^{-(2\hat{\alpha}-1)}$ and independent of $K$.
The bound for $T^V_0$ carries a factor bounded by a constant times $g_0^{-1}$, which depends on
the bare coupling $g_0$, but only $n=k\ge4$ occurs here, so that bound is never invoked. Therefore
\begin{equation*}
\mathbb{D}_V\supseteq\bigl\{(\lambda,\mathsf{h}):\ |\lambda_{n,k'}|\sup|\mathsf{h}|\le E_1/\mathsf{A}_V\bigr\}
\end{equation*}
uniformly in $1\le n\le K'$ and in $K'$, and along the run the total size of the first-row units
anchored in one collar cell satisfies
\begin{align*}
\sum_{\substack{v\in\mathfrak{L}^{V,\mathrm{col}}\\ y_v\in\Delta'}}\mathfrak{N}^{\flat}(v)
&\le M^4F^{*\prime}_{\mathrm{col}}\bar{C}_L^V\,|\lambda_{n,k'}|\sup|\mathsf{h}|\,
\mathsf{A}_V\vartheta_V^n\\
&\le M^4F^{*\prime}_{\mathrm{col}}\bar{C}_L^V\mathsf{A}_V
\end{align*}
on the set where $|\lambda|\sup|\mathsf{h}|\le1$,
a number independent of $K$; the first inequality is \eqref{10.3:eq-proof-collar-sum-6} summed
over $1\le j\le n$, followed by \eqref{10.3:eq-proof-collar-sum-7},
$F'_{\mathrm{col}}\le F^{*\prime}_{\mathrm{col}}$ and $T^V_n\le\mathsf{A}_V\vartheta_V^n$, the second
uses $\vartheta_V^n\le1$ and $|\lambda_{n,k'}|\sup|\mathsf{h}|\le1$.

The dependence runs in one direction only. The hypotheses of the load recursion lemma include rows
of the typed format statement of the vector package, and that format statement is established
entry by entry, the inequality \eqref{10.3:eq-proof-collar-sum-3} being one of its entries. The
load recursion lemma therefore depends on the present proposition and is not a premise of
\eqref{10.3:eq-proof-collar-sum-3}, which holds on $\mathbb{D}_V$ as the class clause of the typed
format statement describes it; it is stated here only to record that the window $\mathbb{D}_V$ in
$\lambda$ does not close as $n$ grows. Nothing of the load recursion lemma, and nothing of the
conclusion of the class clause, enters Steps~1--4, so there is no circularity.

\emph{The activity row of the table.} The table of Notation~\ref{10.3:not-vector-dial-table} also
carries an activity row, obtained, using the subadditivity of $\mathfrak{N}$, by summing the
table's last column into a bound for an activity-side weight. Sizes are turned into bounds on
pieces by a separate weight estimate of the vector package, not by that row. The per-cell collar
bound of the correlation theorem bounds a sum of sizes, so the size chain of the assembly is the
one required. That row is consistent with it, its middle member being the inverse of a weight
times the size total of all rows; it is not used, except for its bound on the side of the hull of
a first-row unit in
level-$n$ units, which fixes the age $\nu=1+n-j$ in Step 2, and as corroboration that the last
column of the table is a total per cell.

\emph{Insertion into the anchored sum.} By \eqref{10.3:eq-proof-halo-anchor-a},
\begin{equation*}
\mathsf{Q}^V_{a+9}\le e^{(a+9)c_lc_{\mathfrak{h}}}(c_0+2)(2c_0+5)^d
\sup_{\Delta'}\sum_{\substack{v\in\mathfrak{L}^{V,\mathrm{col}}\\ y_v\in\Delta'}}
\mathfrak{N}^{\flat}(v)\,e^{(a+9)d_k(K_v)},
\end{equation*}
the supremum being over the cells $\Delta'$ of level $k$ of the collar of $C$.
Since $\varsigma_1\operatorname{rad}(v)\ge0$, each such sum is at most the left member of
\eqref{10.3:eq-proof-collar-sum-3}, hence at most
$\mathsf{K}^{\flat,V}\bar{\mathsf{w}}(\Delta')^3=27\,\mathsf{K}^{\flat,V}$. As
$\mathsf{K}^{\flat}_J\ge0$,
\begin{align*}
\mathsf{Q}^V_{a+9}
&\le e^{(a+9)c_lc_{\mathfrak{h}}}(c_0+2)(2c_0+5)^d\cdot27\,\mathsf{K}^{\flat,V}\\
&\le e^{(a+9)c_lc_{\mathfrak{h}}}(c_0+2)(2c_0+5)^d\cdot27\,\mathsf{K}^{\flat}_+
=\mathsf{K}^{\flat,\mathrm{ex}}_{\mathrm{pl}},
\end{align*}
which is \eqref{10.3:eq-proof-collar-sum-4} and proves (c). Substituting $a=\frac{31}{20}\kappa$,
$c_{\mathfrak{h}}=7^d$ and $c_0=6^d+c_l7^d$, so that $c_0+2=6^d+c_l7^d+2$ and
$2c_0+5=2\cdot6^d+2c_l7^d+5$, and \eqref{10.3:eq-proof-collar-sum-1}, we obtain
\begin{equation}\label{10.3:eq-proof-collar-sum-b}
\begin{aligned}
\mathsf{K}^{\flat,\mathrm{ex}}_{\mathrm{pl}}
={}&e^{(\frac{31}{20}\kappa+9)c_l\cdot7^d}\,(6^d+c_l7^d+2)\,(2\cdot6^d+2c_l7^d+5)^d\cdot27\\
&\cdot\Bigl(\mathsf{K}^{\flat}_J+M^4e^{\frac85\kappa\cdot6^d+\varsigma_1}
F^{*\prime}_{\mathrm{col}}\,\bar{C}_L^V\,E_1\Bigr),
\end{aligned}
\end{equation}
with $c_l=1+234d$, $F^{*\prime}_{\mathrm{col}}$ as in \eqref{10.3:eq-proof-collar-sum-2} and
$a^*=2+s(3)$. This is a finite number, independent of $K$, of $k$ and of the argument
$\mathsf{T}$ of the interior ceiling polynomial $\mathsf{P}_{\mathrm{int}}$ (see
the corresponding display of Theorem~\ref{10.3:thm-interior-piece-bound}), depending only on quantities fixed before the coupling
ceilings $\gamma$; it is a number fixed after $M$, of the
same kind as the enlarged per-cell constant $\mathsf{K}^{\flat,\mathrm{ex}}$ of the correlation
theorem, which is $e^{(a+9)c_lc_{\mathfrak{h}}}\cdot27$ times a geometric count times
$\mathsf{K}^{\flat}$, but not equal to it. No transfer of the anchor from the read box $Q$ to a
cube off $Z$
is needed here, because $\mathsf{Q}^V_{a+9}$ is anchored at an arbitrary cube; the transfer factors
of the correlation theorem therefore do not occur, and the finite count $(c_0+2)(2c_0+5)^d$ stands
in place of its geometric count. No distinction between shallow and deep cores is made, no count of
the form $2d(104L\check{R}_{h+1}+2)^{d-1}$ and no distance weight is read: the anchor in
$\mathfrak{h}_v\setminus K_v$ is served only by the connectedness of the start label
$\operatorname{lab}_0$ and by the condition, among the hypotheses under which
Theorem~\ref{10.3:thm-interior-piece-bound} is stated, that a tree graph meeting two cubes has
length at least their distance.

\emph{The other ceilings.} The number $\mathsf{K}^{\flat,V}$ carries the factor
$M^4F^{*\prime}_{\mathrm{col}}\bar{C}_L^V$, which is $M^4L^{4+2\hat{\alpha}a^*}$ times a constant
depending on $L$.
It enters the first summand of the interior ceiling polynomial $\mathsf{P}_{\mathrm{int}}$ through
$\mathsf{K}^{\flat,\mathrm{ex}}_{\mathrm{pl}}$, as in the corresponding display of Theorem~\ref{10.3:thm-interior-piece-bound}, and is
read only by the ceiling
\begin{equation*}
\sup_{\mathsf{T}\ge\mathsf{T}_{\gamma}}\mathsf{P}_{\mathrm{int}}(\mathsf{T})\,
e^{-(\mathsf{T}^{r_{\mathrm{pr}}}-2)}\le\tfrac12 ,
\end{equation*}
a one-sided ceiling on $\gamma_J$ taken after $M$, hence after $L$. Its dependence on $L$ is
therefore no upper restriction on $L$, in the same way as the dependence of $F_{\mathrm{col}}$ on
$L$ inside the ceilings on $g$ of the load recursion lemma. The pair $(\lambda,\mathsf{h})$ is
confined to $\mathbb{D}_V$, the standing domain of the typed format statement of the vector
package, which is not a new condition; its uniformity in $n$ is
\eqref{10.3:eq-proof-collar-sum-a}.

The ceiling $\gamma\le\gamma_{\rho}$ of \eqref{10.3:eq-proof-collar-sum-9} is used only through
the per-unit bound \eqref{10.3:eq-vector-dial-table-b} read in Step 2. It is a one-sided smallness
condition on the coupling: $\gamma_{\rho}$ depends only on $c_2$, $q$, $\lambda_{\mathrm{fl}}$ and
$a^*$, which are fixed before $\gamma$ in the order of choice of the parameters, so it is chosen
after $L$, $q$ and these constants. It is positive. Indeed, $8c_2+4\ge4$ gives
$]0,\gamma^{\circ}]\subset\,]0,\frac12]$; on this interval $E_{\rho}$ is continuous and
nondecreasing, because $\gamma'^2$ increases and $\gamma'^2/\log\gamma'^{-2}$ increases, the
logarithm $\log\gamma'^{-2}$ being decreasing in $\gamma'$ and at least $\log4>0$; and
$E_{\rho}(\gamma')\to1<2$ as $\gamma'\to0$. Hence $\gamma_{\rho}>0$, $E_{\rho}(\gamma_{\rho})\le2$
by continuity, and every $\gamma\le\gamma_{\rho}$ satisfies $E_{\rho}(\gamma)\le2$ and
$\gamma\le(3/4)^{1/2}e^{-q}$. The ceiling lowers $\gamma$ and leaves every rate, degree and floor
and the order of choice unchanged; no quantity fixed before $\gamma$ depends on it, in particular
neither $\mathsf{K}^{\flat,V}$ nor $\mathsf{K}^{\flat,\mathrm{ex}}_{\mathrm{pl}}$ does, and it is
not among the hypotheses of Theorem~0.
\end{proof}

\begin{proof}[Proof of Theorem~\ref{10.3:thm-interior-piece-bound}, concluded: clauses (d$'$), (e$'$)
and the enlargement of the start label]
\label{10.3:prf-proof-ceiling}
We prove (d$'$) first. Write the interior ceiling polynomial $\mathsf{P}_{\mathrm{int}}(\mathsf{T})$ of
the correlation theorem as the sum of its two summands,
\begin{equation}\label{10.3:eq-proof-ceiling-1}
\mathsf{P}_{\mathrm{int}}(\mathsf{T})
=(1+C_{\Sigma}C_e)^4\,\mathsf{P}^{(1)}(\mathsf{T})+\mathsf{P}^{(2)}(\mathsf{T}).
\end{equation}
Here $(1+C_{\Sigma}C_e)^4\,\mathsf{P}^{(1)}(\mathsf{T})$ is the first summand, whose factor
$\mathsf{P}^{(1)}(\mathsf{T})$ is the braced expression in the display defining
$\mathsf{P}_{\mathrm{int}}$ in the proof of the correlation theorem, at the step fixing the interior
ceiling. The term $\mathsf{P}^{(2)}(\mathsf{T})$ is the
second summand, the one carrying the prefactor $(1+C_{\Sigma}C_e)^3$.

With $\mathsf{P}_{\mathrm{int}}$ so written, that step establishes, at
$\mathsf{T}:=\check{\mathsf{T}}_k$, the following two facts. The first is
\begin{equation}\label{10.3:eq-proof-ceiling-2}
\mathsf{Q}^{\star,\mathrm{pl}}_{a+9}\,(1+C_{\Sigma}C_e)^4
\le(1+C_{\Sigma}C_e)^4\,\mathsf{P}^{(1)}(\check{\mathsf{T}}_k).
\end{equation}
The second is that, with the inequalities $\check{R}_k\ge\check{\mathsf{T}}_k^{\,r_{\mathrm{pr}}}$ and
$\check{\mathsf{T}}_k\ge\mathsf{T}_{\gamma}$, which that step has at every $k\le K$,
\begin{equation}\label{10.3:eq-proof-ceiling-3}
\varepsilon^{\mathrm{int}}(k)+\varepsilon^{\mathrm{ext},\sigma}(k)
\le\mathsf{P}_{\mathrm{int}}(\check{\mathsf{T}}_k)\,
e^{-(\check{\mathsf{T}}_k^{\,r_{\mathrm{pr}}}-2)}\le\tfrac12 .
\end{equation}
Here $\varepsilon^{\mathrm{int}}(k)$ is the interior ceiling of the correlation theorem,
\begin{equation*}
\varepsilon^{\mathrm{int}}(k)
=e^{-(\check{R}_k-2)}\,(1+C_{\Sigma}C_e)^4\,\mathsf{Q}^{\star,\mathrm{pl}}_{a+9}.
\end{equation*}
Neither the definition of $\mathsf{P}_{\mathrm{int}}$ nor these two facts involve the start label.
Both therefore hold without change when
the start label $K_v$ is replaced by $\operatorname{lab}_0=K_v\cup\mathfrak{h}_v$.

Clause (d$'$) enlarges the first summand by $(1+C_{\Sigma}C_e)^4\,\mathsf{K}^{\flat,\mathrm{ex}}_{\mathrm{pl}}$.
In this proof we write the enlarged polynomial as
\begin{equation}\label{10.3:eq-proof-ceiling-4}
\mathsf{P}^{\mathrm{pl}}_{\mathrm{int}}(\mathsf{T})
:=(1+C_{\Sigma}C_e)^4\bigl(\mathsf{P}^{(1)}(\mathsf{T})
+\mathsf{K}^{\flat,\mathrm{ex}}_{\mathrm{pl}}\bigr)+\mathsf{P}^{(2)}(\mathsf{T}).
\end{equation}
By \eqref{10.3:eq-proof-ceiling-2} and the sum over a collar cell treated above, which is the
estimate through which the constant $\mathsf{K}^{\flat,\mathrm{ex}}_{\mathrm{pl}}$ of
\eqref{10.3:eq-proof-collar-sum-b} enters, the first fact takes the enlarged form
\begin{equation}\label{10.3:eq-proof-ceiling-5}
\bigl(\mathsf{Q}^{\star,\mathrm{pl}}_{a+9}+\mathsf{Q}^{V}_{a+9}\bigr)(1+C_{\Sigma}C_e)^4
\le(1+C_{\Sigma}C_e)^4\bigl(\mathsf{P}^{(1)}(\check{\mathsf{T}}_k)
+\mathsf{K}^{\flat,\mathrm{ex}}_{\mathrm{pl}}\bigr).
\end{equation}
In other words, the left side is at most the first summand of
$\mathsf{P}^{\mathrm{pl}}_{\mathrm{int}}(\check{\mathsf{T}}_k)$.
Inequality \eqref{10.3:eq-proof-ceiling-5} is the enlarged form of \eqref{10.3:eq-proof-ceiling-2};
it is in this form that the first fact enters when the bounds on interior pieces are re-run below at
the enlarged start label.

Accordingly, the interior ceiling is now read as
\begin{equation}\label{10.3:eq-proof-ceiling-6}
\varepsilon^{\mathrm{int}}(k)
:=e^{-(\check{R}_k-2)}\,(1+C_{\Sigma}C_e)^4\bigl(\mathsf{Q}^{\star,\mathrm{pl}}_{a+9}
+\mathsf{K}^{\flat,\mathrm{ex}}_{\mathrm{pl}}\bigr).
\end{equation}
In what follows, $\varepsilon^{\mathrm{int}}(k)$ denotes the right side of
\eqref{10.3:eq-proof-ceiling-6}.
We claim that, under the interior ceiling condition of the correlation theorem, which is among the
hypotheses of Theorem~\ref{10.3:thm-interior-piece-bound} and is written out below, with the
enlarged polynomial, as \eqref{10.3:eq-proof-ceiling-8}, and with the two inequalities on
$\check{R}_k$ and $\check{\mathsf{T}}_k$ recorded before \eqref{10.3:eq-proof-ceiling-3},
\begin{equation}\label{10.3:eq-proof-ceiling-7}
\varepsilon^{\mathrm{int}}(k)+\varepsilon^{\mathrm{ext},\sigma}(k)
\le\mathsf{P}^{\mathrm{pl}}_{\mathrm{int}}(\check{\mathsf{T}}_k)\,
e^{-(\check{\mathsf{T}}_k^{\,r_{\mathrm{pr}}}-2)}\le\tfrac12 .
\end{equation}

For the first inequality, consider the terms in which the enlargement enters. In
\eqref{10.3:eq-proof-ceiling-6} the new term is
$e^{-(\check{R}_k-2)}(1+C_{\Sigma}C_e)^4\mathsf{K}^{\flat,\mathrm{ex}}_{\mathrm{pl}}$. The
corresponding new term on the right of \eqref{10.3:eq-proof-ceiling-7} is
$(1+C_{\Sigma}C_e)^4\mathsf{K}^{\flat,\mathrm{ex}}_{\mathrm{pl}}\,
e^{-(\check{\mathsf{T}}_k^{\,r_{\mathrm{pr}}}-2)}$. Since
$\check{R}_k\ge\check{\mathsf{T}}_k^{\,r_{\mathrm{pr}}}$, we have
$e^{-(\check{R}_k-2)}\le e^{-(\check{\mathsf{T}}_k^{\,r_{\mathrm{pr}}}-2)}$, so the former term is at
most the latter. The remaining terms on the left,
\begin{equation*}
e^{-(\check{R}_k-2)}(1+C_{\Sigma}C_e)^4\,\mathsf{Q}^{\star,\mathrm{pl}}_{a+9}
+\varepsilon^{\mathrm{ext},\sigma}(k),
\end{equation*}
are exactly those of \eqref{10.3:eq-proof-ceiling-3}, and are bounded there, through
\eqref{10.3:eq-proof-ceiling-2}, by
$\mathsf{P}_{\mathrm{int}}(\check{\mathsf{T}}_k)\,e^{-(\check{\mathsf{T}}_k^{\,r_{\mathrm{pr}}}-2)}$.
By \eqref{10.3:eq-proof-ceiling-4} and \eqref{10.3:eq-proof-ceiling-1}, the sum of the two bounds is
$\mathsf{P}^{\mathrm{pl}}_{\mathrm{int}}(\check{\mathsf{T}}_k)\,
e^{-(\check{\mathsf{T}}_k^{\,r_{\mathrm{pr}}}-2)}$; adding them gives the first inequality of
\eqref{10.3:eq-proof-ceiling-7}.

For the second inequality, the interior ceiling condition of the correlation theorem, read with
$\mathsf{P}^{\mathrm{pl}}_{\mathrm{int}}$ in place of $\mathsf{P}_{\mathrm{int}}$, is
\begin{equation}\label{10.3:eq-proof-ceiling-8}
\sup_{\mathsf{T}\ge\mathsf{T}_{\gamma}}\mathsf{P}^{\mathrm{pl}}_{\mathrm{int}}(\mathsf{T})\,
e^{-(\mathsf{T}^{r_{\mathrm{pr}}}-2)}\le\tfrac12 .
\end{equation}
Since $\check{\mathsf{T}}_k\ge\mathsf{T}_{\gamma}$, the supremum in
\eqref{10.3:eq-proof-ceiling-8} is at least the value of the supremand at
$\mathsf{T}=\check{\mathsf{T}}_k$, which gives the second inequality of
\eqref{10.3:eq-proof-ceiling-7}. We check that \eqref{10.3:eq-proof-ceiling-8} keeps the form it
has in the correlation theorem.

The supremand involves only quantities fixed before $\gamma$, and the added
summand $(1+C_{\Sigma}C_e)^4\mathsf{K}^{\flat,\mathrm{ex}}_{\mathrm{pl}}$ does not depend on
$\mathsf{T}$: for the terms of $\mathsf{P}_{\mathrm{int}}$ the former is as in the correlation
theorem, and for the added summand both are the content of clause (e$'$), proved below
independently of (d$'$).
Its contribution $(1+C_{\Sigma}C_e)^4\mathsf{K}^{\flat,\mathrm{ex}}_{\mathrm{pl}}\,
e^{-(\mathsf{T}^{r_{\mathrm{pr}}}-2)}$ is therefore finite on $\mathsf{T}\ge\mathsf{T}_{\gamma}$ and
tends to $0$ as $\mathsf{T}\to\infty$. The rest of the supremand has these two properties as in the
correlation theorem, since $r_{\mathrm{pr}}\ge2$. Hence the supremand is finite and tends to $0$.
Condition \eqref{10.3:eq-proof-ceiling-8} therefore remains a one-sided ceiling on $\gamma$, of
degree $-\infty$, as in the correlation theorem. The added summand is of the same kind as the
summand $\mathsf{K}^{\flat,\mathrm{ex}}$ inside $\mathsf{P}^{(2)}(\mathsf{T})$, added in clause (d) of
the exterior-piece bound. Since
\eqref{10.3:eq-proof-ceiling-7} holds for every $k\le K$, clause (d$'$) is proved.

Clause (e$'$) is read off from the display \eqref{10.3:eq-proof-collar-sum-b}, which expresses
$\mathsf{K}^{\flat,\mathrm{ex}}_{\mathrm{pl}}$ in terms of the quantities listed in (e$'$), and none of
the quantities that (e$'$) excludes occurs in it.

It remains to conclude for Theorem~\ref{10.3:thm-interior-piece-bound} at the enlarged start label.
We re-run, at the reader $(K_v\cup\mathfrak{h}_v,\,Q_v)$, the two parts of the proof of the
correlation theorem that bound the interior pieces and that fix the interior ceiling. The two steps
of that proof which concern deep units are not used, because the units considered here are never
deep.

No step of that proof fails under the enlargement. Its only two clauses that mention the position of
$K_v$ are $\mathcal{R}_0\subset C\subset Y_4$ and $K_v\subset Z$, where $C$ is the set for which
$Y(\operatorname{lab}_0)=C\cup\mathfrak{h}_v$, with $Y(A)$ denoting $A$ together with the components
of $Z$ that meet $A$, and where $Y_4=Y_4(\mathfrak{p})$ is the
filing domain of the piece $\mathfrak{p}$. These are used there only to obtain
$\mathcal{R}_0\subset Y_4$ and the non-emptiness of the tuple $\mathfrak{p}$. Both conclusions hold
here: the first by the inclusions
$\operatorname{lab}_0\supseteq\mathfrak{h}_v\supseteq[\,Q_v\,]_k$, the second by the definition of
$\mathfrak{p}$. The following statements carry no clause on the start label or on the position of the
reading set:
\begin{itemize}
\item the admissible families $\mathfrak{Y}(F)$ of a reader;
\item the read-triple statement of the correlation theorem;
\item parts (a), (c) and (e) of its summation lemma.
\end{itemize}

The inputs of this proof are the following.
\begin{itemize}
\item The hypothesis of Theorem~\ref{10.3:thm-interior-piece-bound} on $f_v$ for
$v\in\mathfrak{L}^{V,\mathrm{col}}$: holomorphy of $m\mapsto f_v(m)$ on $\mathfrak{P}$, and
$|f_v(m)-f_v(1)|\le\mathfrak{N}^{\flat}(v)$ there.
\item The four geometric statements on covers and labels that the proof of the correlation theorem
uses, in the form in which it uses them.
\item The collar bound of the exterior-piece bound, with its proof, re-run in the sum over a collar
cell.
\item The following quantities of the vector package, used with their explicit defining
expressions: the constant $C_{\sharp}$, the dial $\lambda_{n,k'}$, the
norms $\mathsf{u}_j$, $\mathsf{o}_j$, $\mathsf{r}_i$, the loads $T^V_n$, the constants $C_L^V$ and
$C_L^{V\prime}$, and the domain $\mathbb{D}_V$.
\item The recursion of the vector package for the vector part, with its explicit expressions, used
as a stated input.
\item Inside $\mathsf{K}^{\flat,V}$, the following enter only as fixed constants whose explicit
values are not used:
$\varsigma_1,\ \mathsf{K}^{\flat}_J,\ C_V,\ c_{\epsilon},\ \bar{\rho}_{a^*},\ C^{\mathrm{mar}}$, a
monotonicity statement of the correlation theorem, and $E_1$.
\item Of $\varsigma^V_{\nu}$, $w_{n,i}$ and $\mathsf{a}_i$, only their indices and signs.
\end{itemize}
\end{proof}

\begin{lemma}[Curvature bounds cover the first-link deformation]\label{10.3:lem-first-link-deformation}
Let $v\in\mathfrak{L}^{V,\mathrm{col}}$ be exposed in the sense of
Notation~\ref{10.3:not-links-polydisc-labels}. For each link $t\in\{1,2,3,4\}$ of the site let
$\Lambda$, $\Lambda^{\sim}$ and $\sigma_0^{(t)}$ be as in
Notation~\ref{10.3:not-links-polydisc-labels}, with $\sigma_0^{(1)}$ the family of cells of the
partition of the first link that are disjoint with $\Lambda^{\sim}$.
Recall from Notation~\ref{10.3:not-links-polydisc-labels} that the source polydisc is defined for an
arbitrary choice, for each $t=1,\dots,4$, of a family $\sigma^{(t)}$ of cells of the partition of
link $t$ that satisfies
\begin{equation}\label{10.3:eq-first-link-deformation-1}
\text{no member of } \sigma^{(t)} \text{ meets the interior of } \Lambda ,
\end{equation}
namely, with $\Pi$ and $\mathfrak{S}_{\mathrm{src}}$ as in
Notation~\ref{10.3:not-links-polydisc-labels},
\begin{equation}\label{10.3:eq-first-link-deformation-2}
\mathfrak{P}:=\mathfrak{S}_{\mathrm{src}}\times\operatorname{supp}\Pi\times
\prod_{t=1}^{4}\bigl\{\,|\mathsf{s}^{(t)}(\Delta)|\le e^{\kappa_1}:\ \Delta\in\sigma^{(t)}\bigr\}
\ni m=(\zeta,\varpi,\vec{\mathsf{s}}),
\end{equation}
with $\operatorname{chain}(m)$ and $\mathcal{U}(m)$ as in
Notation~\ref{10.3:not-links-polydisc-labels}. The families $\sigma_0^{(t)}$ of the site are such a
choice (see the beginning of the proof), and in what follows $\mathfrak{P}$ denotes
\eqref{10.3:eq-first-link-deformation-2} with $\sigma^{(t)}=\sigma_0^{(t)}$. Assume the following.
\begin{itemize}
\item[(a)] The curvature bound for exposed units that accompanies the correlation theorem holds,
as stated there, for every $v$ exposed in the sense of
Notation~\ref{10.3:not-links-polydisc-labels} and every $m\in\mathfrak{P}$.
\item[(b)] The two source conditions of the correlation theorem at the carrier site hold for every
choice of the families $\sigma^{(t)}$ above, with $\vec{\mathsf{s}}$ free on the cells of the region
$X$ of the correlation theorem's step as well. The first of them is the condition read at the
consumer $\mathfrak{E}^{\oplus}_v$. The second is joint holomorphy in the parameters
$\mathsf{s}^{(t)}(\Delta)$, $\Delta\in\sigma^{(t)}$, $t=1,\dots,4$, each link with its own family,
on the product polydisc
\begin{equation*}
\prod_{t=1}^{4}\bigl\{\,|\mathsf{s}^{(t)}(\Delta)|\le e^{\kappa_1}:\ \Delta\in\sigma^{(t)}\bigr\}.
\end{equation*}
\end{itemize}
Let $\mathsf{y}_1\subset\sigma_0^{(1)}$ and $(\zeta,\varpi)\in\mathfrak{S}_{\mathrm{src}}\times
\operatorname{supp}\Pi$. Let $\mathsf{s}^{(1)}$ satisfy $|\mathsf{s}^{(1)}(\Delta)|\le e^{\kappa_1}$
for $\Delta\in\mathsf{y}_1$ and $\mathsf{s}^{(1)}(\Delta)=0$ for
$\Delta\in\sigma_0^{(1)}\setminus\mathsf{y}_1$. For $t=2,3,4$ let $\mathsf{s}^{(t)}$ be arbitrary in
its disc $\{|\mathsf{s}^{(t)}(\Delta)|\le e^{\kappa_1}:\Delta\in\sigma_0^{(t)}\}$. Then:
\begin{enumerate}
\item[(i)] the carrier pair of Notation~\ref{10.3:not-links-polydisc-labels}, evaluated at the
parameters just fixed (the first link interpolated on $\mathsf{y}_1$ and switched off on
$\sigma_0^{(1)}\setminus\mathsf{y}_1$, the other links anywhere in their discs), is
$\operatorname{chain}(m)$ for the point
\begin{equation}\label{10.3:eq-first-link-deformation-3}
m=(\zeta,\varpi,\vec{\mathsf{s}})\in\mathfrak{P},
\end{equation}
and the conclusion of the curvature bound for exposed units holds at this $m$;
\item[(ii)] the two source conditions of hypothesis (b) hold at this $m$.
\end{enumerate}
\end{lemma}

\begin{proof}
By the definition recalled above, the polydisc \eqref{10.3:eq-first-link-deformation-2}
carries the $\vec{\mathsf{s}}$-coordinates of all four links. Among them are the coordinates of the
first link over any family $\sigma^{(1)}$ obeying \eqref{10.3:eq-first-link-deformation-1}.

The first link's family $\sigma_0^{(1)}$, the cells disjoint with $\Lambda^{\sim}$, obeys
\eqref{10.3:eq-first-link-deformation-1}, and so do the families $\sigma_0^{(t)}$, $t=2,3,4$, the
cells disjoint with $\Lambda$; this is recorded in Notation~\ref{10.3:not-links-polydisc-labels}.
So the families of the site are an admissible choice in \eqref{10.3:eq-first-link-deformation-2}.

We now check the coordinates of $m$ one by one. First, $\zeta\in\mathfrak{S}_{\mathrm{src}}$ and
$\varpi\in\operatorname{supp}\Pi$. Next, for $\Delta\in\mathsf{y}_1$ we have
$|\mathsf{s}^{(1)}(\Delta)|\le e^{\kappa_1}$. For $\Delta\in\sigma_0^{(1)}\setminus\mathsf{y}_1$ we
have $\mathsf{s}^{(1)}(\Delta)=0$, and $0\le e^{\kappa_1}$. Finally, for $t\ge 2$ each
$\mathsf{s}^{(t)}$ lies in its disc by assumption. Hence $m\in\mathfrak{P}$.

By the definition of $\operatorname{chain}(\cdot)$ in
Notation~\ref{10.3:not-links-polydisc-labels}, the carrier pair evaluated at the parameters fixed
in the statement, the first link interpolated on $\mathsf{y}_1$ and switched off on
$\sigma_0^{(1)}\setminus\mathsf{y}_1$ and the other coordinates at these values, is
$\operatorname{chain}(m)$. The curvature bound for exposed units is asserted for every exposed $v$
and every $m\in\mathfrak{P}$. By hypothesis (a) its conclusion therefore holds at this $m$. This
proves (i).

For (ii), recall that the proof of the curvature bound for exposed units uses the first source
condition at the consumer $\mathfrak{E}^{\oplus}_v$, with $\vec{\mathsf{s}}$ free on the cells of
the region $X$ as well. By hypothesis (b), this condition and the second one hold for every choice
of the families $\sigma^{(t)}$ obeying \eqref{10.3:eq-first-link-deformation-1}, in particular for
the families $\sigma_0^{(t)}$ of the site, shown above to be such a choice. The second condition is
joint holomorphy in the parameters of all four links, each link with its own family, on the product
polydisc
\begin{equation*}
\bigl\{\,|\mathsf{s}^{(t)}(\Delta)|\le e^{\kappa_1}:\ \Delta\in\sigma_0^{(t)},\ t=1,\dots,4\bigr\},
\end{equation*}
which contains the $\vec{\mathsf{s}}$-coordinates of $m$. So both conditions hold at this $m$.
Among the consequences of the correlation theorem is the same joint holomorphy on the same product
polydisc. This proves (ii).
\end{proof}

\begin{lemma}[Real values unchanged on a box astride the boundary]\label{10.3:lem-real-values-unchanged}
Let $v$ be an inner-collar core, with core box $\square_0$, hull $\hat{X}_v$ and read level
$p^*=n-a^*$ as in Definition~\ref{10.3:def-hull-halo-label}. Recall from
Definition~\ref{10.3:def-hull-halo-label} that the hull is the box $\hat{X}_v=Q_v$ formed by
$[\square_0]_{p^*}$ together with two $\pi_{p^*}$-layers. Let
$\operatorname{ext}_v$ be the extension reading of $v$ on its hull, as defined in the correlation
theorem. We write
$\operatorname{chain}(m)$ for the configuration of the link chain of the correlation theorem at a
point $m$ of the source polydisc $\mathfrak{P}$, and call it a chain configuration. The real chain
configuration is $\operatorname{chain}(\vec{0})$, the chain configuration with link parameters
$\vec{\mathsf{s}}=\vec{0}$, at which the value $\operatorname{val}_v$ of $v$ is read; the real
values of $v$ are the values read there. Assume the
following three parts of the correlation theorem.
\begin{enumerate}
\item[(a)] The local form of the box-restriction clause of the correlation theorem, in its second
alternative. Let $j$ be a level, $X$ a domain, and $\square'$ any box that contains $X$ with two
$\pi_j$-layers of margin (the box of \cite[clause~(iv)]{B-CMP109}). Let $\mathfrak{b}$ be a
configuration such that the pair $(\mathfrak{b},\square')$ is admissible for the box-restriction
clause, in the sense of the correlation theorem. Then the restriction of $\mathfrak{b}$
to $X$ is, modulo gauge, the box minimiser on $\square'$ of the data $\mathsf{V}$, written
$U_j(\square',\mathsf{V})$. Here $\mathsf{V}$ denotes the level-$j$ data of $\mathfrak{b}$ on the
determining set of $\square'$, the set on which the level-$j$ data that determine the box
minimiser on $\square'$ are taken.
\item[(b)] The box-minimiser identity. A minimiser restricted to a box with margin is the box
minimiser of its own determining data. The construction of \cite[clause~(iv)]{B-CMP109} that
yields this is carried out there for an arbitrary domain $X$.
\item[(c)] The chain statement. The chain configurations differ from $U_k$ only
through changes made inside the domain $\Lambda$ of the link. In particular, for a chain
configuration $\mathfrak{b}$ and a box $\square'$ as in~(a) that meets $Z$, the pair
$(\mathfrak{b},\square')$ is admissible for~(a).
\end{enumerate}
Then the real values are unchanged: $\operatorname{ext}_v$ coincides with the real chain
configuration modulo gauge on $[\square_0]_{p^*}$. This holds whatever the position of $Q_v$
relative to $\partial C$, with $C$ as in Lemma~\ref{10.3:lem-core-geometry}; in particular it
holds when $Q_v$ lies astride $\partial C$ and when $Q_v$ meets $Z$.
\end{lemma}

\begin{proof}
We first check that hypothesis~(a) applies, that is, that $Q_v$ is a box of the kind it admits.
Take the level $j=p^*$, the domain $X=[\square_0]_{p^*}$ and the box $\square'=Q_v$.

By its definition (Definition~\ref{10.3:def-hull-halo-label}), $Q_v$ is the box $[\square_0]_{p^*}$
enlarged by two $\pi_{p^*}$-layers. So $Q_v$ is a box, it contains $[\square_0]_{p^*}$, and its
margin around $[\square_0]_{p^*}$ consists of two $\pi_{p^*}$-layers. These are exactly the
requirements of~(a) at
\begin{equation*}
(j,X,\square')=(p^*,[\square_0]_{p^*},Q_v).
\end{equation*}

By~(c), the chain configurations differ from $U_k$ only inside $\Lambda$, and for a chain
configuration $\mathfrak{b}$ the pair $(\mathfrak{b},\square')$ is admissible for~(a) also when the
box $\square'$ meets $Z$. Hence the pair formed by the real chain configuration and the box $Q_v$ is
admissible for~(a), also when $Q_v$ meets $Z$.

Apply~(a) with these arguments to the real chain configuration. Its restriction to
$[\square_0]_{p^*}$ is, modulo gauge, the box minimiser $U_{p^*}(Q_v,\mathsf{V})$ on $Q_v$. Here
$\mathsf{V}$ denotes its level-$p^*$ data on the determining set of $Q_v$.

By the definition of the extension reading in the correlation theorem, $\operatorname{ext}_v$ is
the box minimiser $U_{p^*}(Q_v,\mathsf{V})$ on $Q_v$ of these level-$p^*$ data $\mathsf{V}$, read
on the hull $Q_v$. With the preceding paragraph, $\operatorname{ext}_v$ and the real chain
configuration coincide modulo gauge on $[\square_0]_{p^*}$. This is the
assertion that the real values are unchanged.

No condition on the position of $Q_v$ has been used: (a) admits any box with two layers of
margin, identity~(b) holds for an arbitrary domain, and by~(c) a box meeting $Z$ is admissible.
Hence the argument is unchanged when $Q_v$ lies astride $\partial C$ or meets $Z$.
\end{proof}

\begin{remark}[The near reader]
Let the near reader of $v$ be the reader of the correlation theorem attached to points $y$ near
$C$, that is, in the collar or the ring of $C$
(Lemma~\ref{10.3:lem-core-geometry}). Two alternatives are distinguished for the core $v$,
labelled $(\alpha)$ and $(\beta)$; in alternative $(\alpha)$, $v$ carries no transport
parameter $\mathsf{t}$.

Under alternative $(\alpha)$ the near reader of $v$ reads none of the following: the bounds for
points near $C$ of the off-diagonal supplement to the correlation theorem, the pointed half-power
bound, and the half-power line of the collar statement of the correlation theorem.

Under alternative $(\beta)$ the near reader of $v$ reads the following inputs. Each enters in the
form in which the correlation theorem states it.

The first input is the half-power line of the collar statement, which is also the half-power line
of the off-diagonal supplement. In the notation of the off-diagonal supplement to the correlation
theorem, whose symbols are listed after the second display below, it reads
\begin{equation*}
\mathsf{N}_y\bigl(\mathsf{t}\,\hat{\chi}_{\square}\mathbb{H}_s\bigr)
\le |\mathsf{t}|\,c_{\chi}\,\mathsf{U}\,F_s(y)
\le \dotsb
\le c_{\chi}\,\mathsf{w}_L\,\mathsf{U}\,\min\{1,\mathsf{f}(y)^{1/2}\},
\end{equation*}
with $\mathsf{U}=\mathsf{U}_J$ for the piece of $v$ concerned, and with $s=4$.

The second input, in the same notation, is the family of bounds for points near $C$ of the
off-diagonal supplement, at a
point $y$ in the collar or the ring of $C$, at the levels $h,\dots,k$. In the present case the
level is $k$, by the geometry of an inner-collar core (Lemma~\ref{10.3:lem-core-geometry}). The
increment there is at most
\begin{equation*}
2(1+b_+)^{1/2}K_D\,c_{\chi}\,\varphi_{\flat}^{-1}\,\frac{\mathsf{w}_L\mathsf{U}_J}{\alpha_{*,k}},
\qquad
\frac{\mathsf{w}_L\mathsf{U}_J}{\alpha_{*,k}}=\frac{6B_HA_1}{C_{\theta'}C_*}.
\end{equation*}
The ratio on the right is free of the coupling. Under the parameter floor $(F_{\theta'})$ of the
correlation theorem, the increment is at most $1/16$. The entries of the supplement that concern a
location in the large-field set $Z_h$ at the base level $h$ are void for these cores, since by
Lemma~\ref{10.3:lem-core-geometry} no such location occurs for them.

All symbols in the two displays are those fixed in the off-diagonal supplement to the correlation
theorem, from which both displays are taken, namely $\mathsf{N}_y$, $\hat{\chi}_{\square}$,
$\mathbb{H}_s$ with its index $s$, $F_s$, $\mathsf{f}$, $\mathsf{U}$, $\mathsf{U}_J$, $c_{\chi}$,
$b_+$, $K_D$, $\varphi_{\flat}$, $B_H$, $C_{\theta'}$ and $C_*$, while $A_1$ is Ba{\l}aban's
auxiliary parameter; the set $S_v$ below and the set $Z_h$ above are likewise those of the
supplement. In particular $\mathsf{w}_L$ is the weight of the pointed half-power bound, and
$\alpha_{*,k}$ is the constant $\alpha_{*,n}$ of the gauge reading at $n=k$.

The hypotheses under which these bounds are read are:
\begin{itemize}
\item $y$ lies in the collar of $C$ at level $k$, by \eqref{10.3:eq-core-geometry-c};
\item $s=4$ and $\mathsf{U}=\mathsf{U}_J$;
\item the set $S_v$ of the off-diagonal supplement is the set $\mathfrak{r}_v$, the union of the
support of $v$ and the star $\operatorname{St}(\Delta_y)$ of the cell $\Delta_y$ of $y$;
\item the parameter floor $(F_{\theta'})$ holds.
\end{itemize}

The third input is the pointed half-power bound, which is stated at the link for every compact
pointed piece of the plain units of clauses~(c) and~(d) of the definition of plain units in the
correlation theorem. The piece of $v$ at hand satisfies the definition of a compact pointed piece,
so the bound covers it.

Of the data attached to $v$, the near reader under $(\beta)$ uses only the position of $y$, the
cell $\Delta_y$ and the set $S_v=\mathfrak{r}_v$; the hull $\hat{X}_v$ is not used.
\end{remark}

\begin{proposition}[Inner-collar cores under the extended correlation theorem]
\label{10.3:prop-cores-extended}
Assume the following.
\begin{enumerate}
\item[(1)] The parameter conditions of Theorem~\ref{10.3:thm-interior-piece-bound}: those of the
correlation theorem, together with the floors $(F_{\kappa_1})$, $(F_w)$ and $\kappa\ge 180$, the bound
$\gamma\le\gamma_J$, and the conditions $(c\text{-reg})$ and $(c\text{-int})$.
\item[(2)] The declared extension of the correlation theorem at its first stage, which is also a
hypothesis of Theorem~\ref{10.3:thm-interior-piece-bound}, in the following form.
For a unit $v$ of the inner-collar class $\mathfrak{L}^{V,\mathrm{col}}$ the core is
$X_v=\square_0$, of class (b) and interior in the sense of the correlation theorem's rule for
interior cores. The read box is $Q_v=Q^{(p_v)}$ with $p_v=n-2-s(\bar{\mathsf{w}}_v)$, where $n$ is
the level of the unit $v$, without the
requirement that $Q_v$ be contained in $Z$. The unit is declared exposed, with hull
$\hat{X}_v=Q_v$; its support is $\mathfrak{r}_v$, with
$Q_v\subset\mathfrak{r}_v\subset\bigcup\mathfrak{h}_v$, and its reading set is
$\mathcal{R}_0=\hat{X}_v\cup X_v=Q_v$. The units $v$ of the inner-collar
class $\mathfrak{L}^{V,\mathrm{col}}$ are entered in the correlation theorem's enumerations of
unit classes, and the interior-piece bound is applied to them at the start label
$\operatorname{lab}_0=K_v\cup\mathfrak{h}_v$. At the first stage of the extension the first two
links of a piece are treated by the two link prescriptions of that stage, the live-piece
datum $\mathcal{P}_v$ is
empty, neither the partition of unity $\zeta_{\square}$ nor the transport parameter $\mathsf{t}$
enters, and the pieces are the tuples $\mathfrak{p}=(v;\mathsf{y}_1,\dots,\mathsf{y}_4)$. In the
summation the per-cell collar constant $\mathsf{K}^{\flat}_J$ is replaced by
$\mathsf{K}^{\flat}_+:=\mathsf{K}^{\flat}_J+\mathsf{K}^{\flat,V}$, where $\mathsf{K}^{\flat,V}$ is
the per-cell collar constant of the vector part.
\item[(3)] $(\lambda,\mathsf{h})\in\mathbb{D}_V$, as in Theorem~\ref{10.3:thm-interior-piece-bound}.
\end{enumerate}
Then, for every $v\in\mathfrak{L}^{V,\mathrm{col}}$, every $K$, every $k\le K$ and every label, the
following hold.
\begin{enumerate}
\item[(i)] For every auxiliary statement of the proof of the correlation theorem that is applied
to $v$, every hypothesis of that statement is satisfied. The correlation theorem's wording has to
be extended by hypothesis~(2) only where a hypothesis restricts the units to which it applies
--- its enumerations of unit classes, its rule fixing the read box, and its exposure clauses ---
and in the anchoring of the interior-piece bound, in which the anchored weighted reader sum
$\mathsf{Q}^V_{a+9}:=\mathsf{Q}_{a+9}[\mathfrak{L}^0_V]$ is bounded by
$\mathsf{K}^{\flat,\mathrm{ex}}_{\mathrm{pl}}$. All other hypotheses are of geometric or
analytic content, including those on which the mechanism of the proof of the interior-piece bound
rests, and are satisfied as the correlation theorem states them, for the unit with
the data of hypothesis~(2).
\item[(ii)] Nothing further is required of the correlation theorem, and none of its exponents,
counts, floors or ceilings changes. In particular
\begin{equation}\label{10.3:eq-cores-extended-1}
\sum_{\mathfrak{p}}\mathfrak{p}+\operatorname{val}_v=f_v\bigl(\operatorname{chain}(\vec{\mathbb{1}})\bigr)
\end{equation}
holds exactly. The constants $\mathsf{K}^{\flat,\mathrm{ex}}$, $\mathsf{K}^{\flat}_+$ and
$\mathsf{K}^{\flat,\mathrm{ex}}_{\mathrm{pl}}$ are numbers fixed after $M$, and they enter the
conditions $(c\text{-}\sharp)$, $(c\text{-}\sharp')$, $(c\text{-br})$ and $(c\text{-int})$ of the
correlation theorem, whose form is unchanged. The correlation theorem's rule for interior cores, the
complementarity of its
classes (b) and (c), and its statement that no core lies in $\mathfrak{L}^{c}$ hold as stated there.
\item[(iii)] Each auxiliary statement of the proof of the correlation theorem that the
proofs of Lemmas~\ref{10.3:lem-certified-cube}, \ref{10.3:lem-consumer-domain},
\ref{10.3:lem-box-regularity}, \ref{10.3:lem-core-size-bound},
\ref{10.3:lem-star-contains-box}, \ref{10.3:lem-amplitude-one-reading},
\ref{10.3:lem-first-link-deformation} and~\ref{10.3:lem-real-values-unchanged}, and of
Theorem~\ref{10.3:thm-interior-piece-bound}, invoke, whose proof is not reproduced here, and whose
hypotheses are all of geometric or analytic content, of the kind described in the last sentence
of~(i), is applied to $v$ by its statement, in the
same way as in the
two steps of the proof of the correlation theorem that apply it to that theorem's own units: the
step that fixes the data of an exposed unit, namely its reading set and the holomorphy of its
function on the source polydisc $\mathfrak{P}$ together with the bound on its increment there, and
one further step of that proof. None of these
statements has a hypothesis that distinguishes $v$ from those
units.
\end{enumerate}
\end{proposition}

\begin{proof}
Fix $v\in\mathfrak{L}^{V,\mathrm{col}}$, $K$, $k\le K$, a label and $(\lambda,\mathsf{h})\in\mathbb{D}_V$.

\emph{Proof of (i).} The hypotheses of the auxiliary statements applied to $v$ fall into two groups.
The first group consists of the geometric and analytic hypotheses. They are checked statement by
statement in Lemmas~\ref{10.3:lem-certified-cube}, \ref{10.3:lem-consumer-domain},
\ref{10.3:lem-box-regularity}, \ref{10.3:lem-core-size-bound},
\ref{10.3:lem-star-contains-box}, \ref{10.3:lem-amplitude-one-reading},
\ref{10.3:lem-first-link-deformation} and~\ref{10.3:lem-real-values-unchanged}.
In each of these lemmas the hypothesis concerned is met in the form in which the correlation
theorem states it, for the unit with the data of hypothesis~(2); the wording of the hypothesis
itself is that of the correlation theorem, unamended. The same holds for the mechanism
of the interior-piece bound. Its proof, carried out in the proof of
Theorem~\ref{10.3:thm-interior-piece-bound} and in the proof that the
interior ceiling absorbs the constant, uses the auxiliary statements of the correlation theorem in
the form in which that theorem states them.

The second group consists of the remaining hypotheses, those named in the second sentence
of~(i). A unit $v$ of $\mathfrak{L}^{V,\mathrm{col}}$ is listed in those enumerations, has the read
box $Q_v=Q^{(p_v)}$ without the requirement $Q_v\subset Z$, and is exposed with hull $\hat{X}_v=Q_v$
by hypothesis~(2), and by nothing else. The interior-piece bound at the start label
$\operatorname{lab}_0=K_v\cup\mathfrak{h}_v$ is Theorem~\ref{10.3:thm-interior-piece-bound}; in its
proof the anchored weighted reader sum $\mathsf{Q}^V_{a+9}$ is bounded by
$\mathsf{K}^{\flat,\mathrm{ex}}_{\mathrm{pl}}$ through the per-cell transfer evaluated at
$\operatorname{lab}_0$. That theorem is proved under hypotheses~(1), (2) and~(3) together with the
holomorphy of $m\mapsto f_v(m)$ for $m$ in the source polydisc $\mathfrak{P}$ and the bound
$|f_v(m)-f_v(1)|\le\mathfrak{N}^{\flat}(v)$ there, which are supplied by
Lemmas~\ref{10.3:lem-certified-cube}, \ref{10.3:lem-consumer-domain},
\ref{10.3:lem-box-regularity} and~\ref{10.3:lem-core-size-bound}; so the anchoring hypothesis is
satisfied. This proves~(i).

\emph{Proof of (ii).} The correlation theorem states, as one of its consequences, that the
splitting of $f_v(\operatorname{chain}(\vec{\mathbb{1}}))$ into pieces changes the localisation of
terms and never the exponent, and that the splitting is exact. For $v\in\mathfrak{L}^{V,\mathrm{col}}$
the splitting into $\operatorname{val}_v$ and the pieces $\mathfrak{p}$ is the pieces clause of Theorem~\ref{10.3:thm-interior-piece-bound}. This is the
identity~\eqref{10.3:eq-cores-extended-1}. No exponent is therefore altered in passing to these
units.

The three constants named in~(ii) are, by their definitions, numbers fixed after $M$ and enter
the four conditions
named there in unchanged form; consequently the extension~(2) imposes no floor or
ceiling beyond those of hypothesis~(1), and no count of the correlation theorem is altered either.
The three statements of the correlation theorem named in the last sentence of~(ii) are used as
stated in that theorem. No
hypothesis is added to the correlation theorem, and none of its statements is modified beyond
hypothesis~(2). This proves~(ii).

\emph{Proof of (iii).} The proof of the correlation theorem applies each statement described
in~(iii) to its own units in the same way, in the two steps described in~(iii). The hypotheses of
these statements are among those of the first group
in~(i), and there they are met for $v$ as for the correlation theorem's own units. Hence none of
them distinguishes $v$. This proves~(iii).
\end{proof}

\begin{remark}[Clauses added to the correlation theorem's lists]\label{10.3:rem-added-clauses}
Throughout this remark $v$ denotes a smooth far $\sharp$-core of the vector part
$\hat{\mathfrak{V}}^{\sharp(j)}$, lying in the pointed class $\mathfrak{L}^{\mathrm{pt}}$, with
\begin{equation}\label{10.3:eq-added-clauses-1}
X_v=\square_0\subset C,\qquad Q_v:=Q^{(n-2-s(\bar{\mathsf{w}}_v))}\not\subset Z ,
\end{equation}
where $C$ is the core set of the correlation theorem and $\partial C$ its boundary.
Such a core is declared exposed, with hull $\hat{X}_v:=Q_v$, support $S_v:=\mathfrak{r}_v$ and start
label $\operatorname{lab}_0:=K_v\cup\mathfrak{h}_v$, as in Definition~\ref{10.3:def-hull-halo-label}.
It is read at $\operatorname{ext}_v$ by the pointed reading rule on $Q_v\subset\mathfrak{r}_v$, at
$(n,\bar{\mathsf{w}}_v)=(k,3)$.
For every such $v$ the lists of the correlation theorem are extended by the following clauses. Each
item names the clause of the correlation theorem and then states the clause added to it.
\begin{enumerate}
\item[(i)] \emph{Holomorphy at the extension reading.} The correlation theorem states, for every
pointed member of the exterior class $\mathfrak{L}^{\mathrm{ex}}$ --- a single, a Wilson piece or a
boundary constituent of $\mathfrak{L}^{c}$; a single, a Wilson piece, a boundary constituent or a
far core underlying a transported piece, that is, a $\mathsf{t}$-piece of the class
$\mathfrak{L}^{Y_0}$ of such pieces --- holomorphy on $\mathfrak{P}$ and
the tabulated size $\mathfrak{N}^{\flat}$ at $\bar{\mathsf{w}}_v:=3w^{\vee}(\Delta_{y_v})$ and at the
zone $n\le k$ of $y_v$. For compact pieces it states these on a joint polydisc in $(m,\mathsf{t})$.
Added clause: for every $v$ as above, read at $\operatorname{ext}_v$ on $Q_v\subset\mathfrak{r}_v$ at
$(n,\bar{\mathsf{w}}_v)=(k,3)$, the map $m\mapsto f_v(m)$ is holomorphic on $\mathfrak{P}$ and
\begin{equation}\label{10.3:eq-added-clauses-2}
|f_v(m)-f_v(1)|\le\mathfrak{N}^{\flat}(v)\qquad(m\in\mathfrak{P}),
\end{equation}
where $\mathfrak{N}^{\flat}(v)$ is the size of the vector part. There is no transport parameter
$\mathsf{t}$ and no joint polydisc. The conclusion that the correlation theorem draws for pointed
units on $C$ from its size bound on $C$ holds for such a $v$ by \eqref{10.3:eq-added-clauses-2}
instead of that bound.
\item[(ii)] \emph{The exposed units.} The correlation theorem enumerates the exposed units: in the
class $\mathfrak{L}^{c}$ it lists them and states that no core lies in $\mathfrak{L}^{c}$ (the box
$\square^{\sim\nu+1}$ lies on one side of $\partial Z$, and the units near $\partial Z$ are
boundary constituents), and in $\mathfrak{L}^{Y_0}$ every pointed piece is exposed. Added clause: in
$\mathfrak{L}^{\mathrm{pt}}$, the cores $v$ as above are exposed; for them $y_v$ lies within two
units of $\partial C$ and $n=k$, the first geometric clause of the exterior consumer being read from
$Q_v\not\subset C$, and the read box leaves $C$.
They take, unchanged at $\bar{\mathsf{w}}_v=3$, the support $\mathfrak{r}_v$, the halo
$\mathfrak{h}_v$, the domain $X^{\oplus}_v:=Y(K_v\cup\mathfrak{h}_v)$ and the consumer
$\mathfrak{E}^{\oplus}_v$ at it. They also take the first and third geometric clauses of the
exterior consumer, the
exterior current bound, the exterior regularity bound and the exterior small-piece integration
clause. The scope of the exterior clauses is extended to these $v$. The statement that no core lies
in $\mathfrak{L}^{c}$ is unchanged.
\item[(iii)] \emph{The interior line of the table of pieces.} For plain interior units the
correlation theorem assigns the all-empty tuples to the interior-volume values. It assigns the
non-empty tuples leaving $Z$ to $\varepsilon^{\mathrm{int}}$, and the non-empty tuples staying in
$C$ to the interior-volume class. Added cell, for the exposed cores $v$ as above (start label
$\operatorname{lab}_0\not\subset Z$, so that no tuple stays inside $C$):
\begin{itemize}
\item the all-empty tuple goes to $\mathfrak{v}^{\mathrm{val}}$, with
$\mathfrak{K}_v:=K_v\cup\mathfrak{h}_v\in\mathcal{V}^{\mathrm{ext}}$, placed at
$Y_1(v)=Y(K_v\cup\mathfrak{h}_v)=X^{\oplus}_v$, and with the constant change
$\mathsf{K}^0\mapsto\mathsf{K}^0+\mathsf{K}^{\flat,\mathrm{ex}}$, $\mathsf{K}^{\flat}_+$ standing in
place of $\mathsf{K}^{\flat}$;
\item every non-empty tuple $\mathfrak{p}=(v;\mathsf{y}_1,\dots,\mathsf{y}_4)$ goes to
$\varepsilon^{\mathrm{int}}$, with
$\mathsf{Q}^{\star,\mathrm{pl}}_{a+9}\mapsto\mathsf{Q}^{\star,\mathrm{pl}}_{a+9}
+\mathsf{K}^{\flat,\mathrm{ex}}_{\mathrm{pl}}$;
\item the interior-volume class holds none of their tuples.
\end{itemize}
\item[(iv)] \emph{The partition of pieces.} The correlation theorem states that at the carrier
every piece of every unit lies in exactly one of its classes. Among these classes are
$\mathfrak{v}^{\mathrm{val}}$, which holds the all-empty tuples, and $\varepsilon^{\mathrm{int}}$,
which holds the non-empty interior tuples whose domain leaves $Z$. Another is the interior-volume
class, which holds the non-empty interior tuples staying inside $C$. Added note: for the exposed
interior cores $v$ as above, the value lies in $\mathfrak{v}^{\mathrm{val}}$ through
$\mathcal{V}^{\mathrm{ext}}$ and all non-empty tuples lie in $\varepsilon^{\mathrm{int}}$. Nothing
lies in the interior-volume class.
\item[(v)] \emph{Placing the values of exposed units.} For an exposed pointed unit among the units
whose values the correlation theorem places (a short single or a Wilson piece astride
$\partial Z$), the correlation theorem puts
$\mathfrak{K}_v:=K_v\cup\mathfrak{h}_v=\operatorname{lab}_1$, the start label it assigns to these
units, so that $Y_1(v)=Y(\operatorname{lab}_1)=X^{\oplus}_v$. It places the value
$\operatorname{val}_v=f_v(\operatorname{chain}(\vec{0}))$ where
$|\operatorname{val}_v-f_v(1)|\le\mathfrak{N}^{\flat}(v)$ is proved, the first consumer estimate
being applied at $\mathfrak{E}^{\oplus}_v$ inside its hypothesis $X\subset Y_4$ for the all-empty
tuple as well. There $\mathfrak{K}_v$ is
connected and $d_k(\mathfrak{K}_v)\le d_k(K_v)+c_lc_{\mathfrak{h}}$, and the factor
$e^{a c_lc_{\mathfrak{h}}}$ this costs the weight lies inside $\mathsf{K}^{\flat,\mathrm{ex}}$.
Added clause: likewise for $v$ as above, with
$\mathfrak{K}_v:=K_v\cup\mathfrak{h}_v=\operatorname{lab}_0$ and
$Y_1(v)=Y(\operatorname{lab}_0)=X^{\oplus}_v$.
\item[(vi)] \emph{The supply of per-cell constants.} For the exposed pointed units of
$\mathcal{V}^{\mathrm{ext}}$ the correlation theorem makes the change
$\mathsf{K}^0\mapsto\mathsf{K}^0+\mathsf{K}^{\flat,\mathrm{ex}}$, where the constant
$\mathsf{K}^{\flat,\mathrm{ex}}$ of its supply clause is a fixed multiple of
$e^{(a+9)c_lc_{\mathfrak{h}}}\,\mathsf{K}^{\flat}$,
a number chosen after $M$, obtained by summing per cell at $\bar{\mathsf{w}}=3$ and anchored as for
the exterior pieces of the correlation theorem. Added clause:
the same holds for the exposed cores $v$ of the vector part, with
\begin{equation}\label{10.3:eq-added-clauses-3}
\mathsf{K}^{\flat}\mapsto\mathsf{K}^{\flat}_+:=\mathsf{K}^{\flat}+\mathsf{K}^{\flat,V}.
\end{equation}
\item[(vii)] \emph{The reader class and real values.} The correlation theorem states that real
values are unchanged. At real $\zeta$ and $\vec{\mathsf{s}}=\vec{\mathbb{1}}$ the chain is
Ba{\l}aban's minimiser, and the gauge reading, applied with $(j,X,\square_0)$ replaced by
$(p,[\square_0]_p,Q_v)$, where $Q_v$ is $[\square_0]_p$ with two $\pi_p$-layers, gives
$\operatorname{ext}_v=\operatorname{chain}$
modulo gauge on $[\square_0]_p$, with $p:=p_v$ of item~(viii); the $\mathbf{J}$-slot of the carrier
pair is not read, and the
reader belongs to the correlation theorem's reader class with reading set
$\mathcal{R}_0:=Q_v\subset C$. For $v$ as above the last sentence is replaced by the following.
The $\mathbf{J}$-slot of the carrier pair is not read, and the reader belongs to that reader class
with reading set $\mathcal{R}_0:=\hat{X}_v\cup X_v=Q_v$. Moreover, for every tuple $\mathfrak{p}$,
\begin{equation}\label{10.3:eq-added-clauses-4}
\begin{split}
Q_v&\subset\mathfrak{r}_v\subset\textstyle\bigcup[\mathfrak{r}_v]_k
\subset\bigcup\mathfrak{h}_v\subset\bigcup\operatorname{lab}_0\\
&\subset Y(\operatorname{lab}_0)=X^{\oplus}_v\subset Y_4(\mathfrak{p}),
\end{split}
\end{equation}
so that the read box lies inside the consumer's domain, though not inside $C$. Real values are
unchanged; the box is admitted through the alternative in which a box meeting $Z$ is admissible.
\item[(viii)] \emph{The choice of $p_v$.} For interior cores the correlation theorem puts
\begin{equation*}
p_v:=\max\{p:\ j\le p\le n-2-s(\bar{\mathsf{w}}_v),\ Q^{(p)}\subset Z\}.
\end{equation*}
For $v$ as above, which is smooth with $\nu\ge 3+s(3)$, one puts instead
\begin{equation}\label{10.3:eq-added-clauses-5}
\begin{gathered}
p_v:=n-2-s(\bar{\mathsf{w}}_v),\qquad a_v:=n-p_v=2+s(3),\\
b_v:=p_v-j,\qquad Q_v:=Q^{(p_v)}.
\end{gathered}
\end{equation}
The hull $\hat{X}_v:=Q_v$ and the radius
$\operatorname{rad}(v):=\max_{x\in Q_v}\operatorname{dist}_n(x,y_v)$
are taken unchanged.
\end{enumerate}
\end{remark}

\begin{proof}
Item~(viii) fixes $Q_v$ and is checked first; then come items~(i), (ii) and~(vii), whose chain of
inclusions the remaining items~(iii)--(vi) use.

\emph{Item~(viii).} At $(n,\bar{\mathsf{w}}_v)=(k,3)$ we have $n-p_v=2+s(\bar{\mathsf{w}}_v)=2+s(3)$,
which is the stated value of $a_v$. The age is $\nu=1+n-j$, and therefore
\begin{equation*}
b_v=p_v-j=n-2-s(3)-j=\nu-3-s(3)\ge 0
\end{equation*}
by the smoothness condition $\nu\ge 3+s(3)$. Hence $j\le p_v\le n-2-s(\bar{\mathsf{w}}_v)$, so $p_v$
lies in the range over which the interior rule maximises. Only the condition $Q^{(p)}\subset Z$ is
dropped, and \eqref{10.3:eq-added-clauses-1} records that it fails at $p=p_v$.

\emph{Item~(i).} With $Q_v$ as in \eqref{10.3:eq-added-clauses-5}, the reading of $f_v(m)$ at
$\operatorname{ext}_v$ uses the pair $(Q_v,p_v)$ only. The inclusion $Q_v\subset\mathfrak{r}_v$ is
part of \eqref{10.3:eq-core-radius-a}. Holomorphy of $m\mapsto f_v(m)$ on $\mathfrak{P}$ and the bound
\eqref{10.3:eq-added-clauses-2}, with $\mathfrak{N}^{\flat}(v)$ the size of the vector part, are
Lemma~\ref{10.3:lem-core-size-bound}, in the form \eqref{10.3:eq-core-size-bound-a}. No transport
parameter enters that lemma,
so no joint polydisc is needed. Hence the conclusion that the correlation theorem draws for pointed
units on $C$ from its size bound on $C$ follows for $v$ from \eqref{10.3:eq-added-clauses-2}.

\emph{Item~(ii).} By \eqref{10.3:eq-added-clauses-1}, $\square_0\subset C$ while the read box
$Q_v$ is not contained in $Z$; as the declaration of these cores states, the read box is not
contained in $C$ either. The
first geometric clause of the exterior consumer is read from $Q_v\not\subset C$, at $y_v$ within two
units of $\partial C$ and $n=k$. The size assigned to $v$ exceeds its own supremum on its own
reading, and its read box $Q_v$ leaves $C$; by the exposure criterion of the correlation theorem,
recalled in Remark~\ref{10.3:rem-one-extension}, $v$ therefore needs the exterior surplus and is
exposed. The objects $\mathfrak{r}_v$,
$\mathfrak{h}_v$, $X^{\oplus}_v$ and $\mathfrak{E}^{\oplus}_v$ are defined for $v$ in
Definition~\ref{10.3:def-hull-halo-label} and \eqref{10.3:eq-core-radius-a}. Each of the exterior
clauses listed in item~(ii) is taken verbatim, at the weight $\bar{\mathsf{w}}_v=3$. The core $v$
lies in $\mathfrak{L}^{\mathrm{pt}}$, not in $\mathfrak{L}^{c}$.
Hence the statement that no core lies in $\mathfrak{L}^{c}$ is not affected.

\emph{Item~(vii).} Since $\square_0\subset[\square_0]_{p_v}$ and $Q_v=Q^{(p_v)}$ is
$[\square_0]_{p_v}$ together with two $\pi_{p_v}$-layers, we have $X_v=\square_0\subset Q_v$. With
$\hat{X}_v=Q_v$ this gives $\hat{X}_v\cup X_v=Q_v$. We now take the chain
\eqref{10.3:eq-added-clauses-4} link by link.
\begin{itemize}
\item $Q_v\subset\mathfrak{r}_v$ is contained in \eqref{10.3:eq-core-radius-a}.
\item $\mathfrak{r}_v\subset\bigcup[\mathfrak{r}_v]_k$ holds by the definition of $[\,\cdot\,]_k$.
\item $\bigcup[\mathfrak{r}_v]_k\subset\bigcup\mathfrak{h}_v$ is the halo statement in
\eqref{10.3:eq-core-radius-a}.
\item $\mathfrak{h}_v\subset\operatorname{lab}_0$ holds because $\operatorname{lab}_0=K_v\cup\mathfrak{h}_v$.
\item $\bigcup\operatorname{lab}_0\subset Y(\operatorname{lab}_0)$ holds because $Y(A)$ is $A$
together with the components of $Z$ that meet $A$.
\item $Y(\operatorname{lab}_0)=X^{\oplus}_v$ is the definition of $X^{\oplus}_v$.
\item $X^{\oplus}_v\subset Y_4(\mathfrak{p})$ for every tuple is
\eqref{10.3:eq-consumer-domain-a}.
\end{itemize}
The choice $\mathcal{R}_0:=\hat{X}_v\cup X_v$ is the correlation theorem's rule for readers: at
link $t$ the label is $\operatorname{lab}_t$ and the reader has reading set $\mathcal{R}_{t-1}$,
starting from $\mathcal{R}_0:=\hat{X}\cup X$ of the unit read; for its exposed units of
$\mathfrak{L}^{c}$ astride $\partial Z$ the correlation theorem states $\mathcal{R}_1:=\hat{X}\cup X$.
That the reader of $v$ belongs to this reader class, as the correlation theorem defines it, is
shown in \eqref{10.3:eq-proof-per-cell-a}. A reader of that class with reading set
$\mathcal{R}_0=Q_v$ therefore reads only inside the consumer's domain. The claim that real values are
unchanged is Lemma~\ref{10.3:lem-real-values-unchanged}; it uses the alternative in which a box
meeting $Z$ is admissible.

\emph{Items~(iii) and~(iv).} By \eqref{10.3:eq-added-clauses-4},
$Q_v\subset\bigcup\operatorname{lab}_0$. Since $Q_v\not\subset Z$, also
$\bigcup\operatorname{lab}_0\not\subset Z$, so no tuple of $v$ stays inside $C$. The interior-volume
class therefore holds none of them.

The core $v$ is an inner-collar core in the sense of Definition~\ref{10.3:def-hull-halo-label},
that is, a member of the class $\mathfrak{L}^{V,\mathrm{col}}$. Every non-empty tuple
$\mathfrak{p}=(v;\mathsf{y}_1,\dots,\mathsf{y}_4)$ is bounded by the
interior-piece bound at the enlarged start label, Theorem~\ref{10.3:thm-interior-piece-bound}, in
the corresponding form. Its hypothesis for
$v\in\mathfrak{L}^{V,\mathrm{col}}$ is item~(i).
The tuple therefore lies in $\varepsilon^{\mathrm{int}}$, with the constant
$\mathsf{Q}^{\star,\mathrm{pl}}_{a+9}$ enlarged by $\mathsf{K}^{\flat,\mathrm{ex}}_{\mathrm{pl}}$.

The all-empty tuple carries $\operatorname{val}_v=f_v(\operatorname{chain}(\vec{0}))$. This is placed
in $\mathfrak{v}^{\mathrm{val}}$ at $\mathfrak{K}_v=\operatorname{lab}_0\in\mathcal{V}^{\mathrm{ext}}$,
as in item~(v) below, with the constant change of item~(vi). Item~(iv) records exactly this
assignment, so each piece of $v$ lies in exactly one class.

\emph{Item~(v).} The label $\mathfrak{K}_v=K_v\cup\mathfrak{h}_v$ is connected by
\eqref{10.3:eq-core-radius-a}. The domain $Y_1(v)=Y(\operatorname{lab}_0)=X^{\oplus}_v$ belongs to
the domain class $\mathbf{D}_k^{Z}$ by \eqref{10.3:eq-consumer-domain-a}. The bound for
$\operatorname{val}_v-f_v(1)$ by $\mathfrak{N}^{\flat}(v)$ is proved on this domain through item~(i).
The first consumer estimate of the correlation theorem is then used at $\mathfrak{E}^{\oplus}_v$,
inside its hypothesis $X\subset Y_4$, for the all-empty tuple as well. This is the same procedure
as for the exposed pointed units of the original clause of item~(v).

\emph{Item~(vi).} The per-cell sum for $v$ is the exterior per-cell collar sum at the cell
$\Delta'=\Delta_{y_v}$, which lies in the collar of $C$. At that cell $\bar{\mathsf{w}}(\Delta')=3$,
$d_k(K_v)\le 6^d$ and $\operatorname{rad}(v)\le 1/L$, by \eqref{10.3:eq-proof-per-cell-a}.
Hence the same exterior constant
applies, with
the per-cell constant of the vector part added, as in \eqref{10.3:eq-added-clauses-3}. The reference
value $f_v(1)$ is the volume entry of the typed bound for the vector part on the domain
$\mathbb{D}_V$; it is not an entry of the correlation theorem.
\end{proof}

\begin{remark}[Insufficiency of the unextended read-box rule]\label{10.3:rem-unextended-rule}
Let $v$ be an inner-collar core of the vector part in the sense of
Notation~\ref{10.3:not-inner-collar-cores}. It has level $j$, anchor $y\in C$ with $C$ a component
of $Z$, zone $n=\pi(y)$, $t=L^{j-n}$ and $\nu=1+n-j\ge a^*+1$. Moreover
$\square_0=\square^{\sim\nu+1}(y)\subset C$ and $Q^{(p^*)}(y)\not\subset C$, where $s^*=s(3)$,
$a^*=2+s^*$ and $p^*=n-a^*$. The weight is $\bar{\mathsf{w}}_v=3$.

Keep $v$ at the grade of an interior unit of the correlation theorem, with $X_v=\square_0$ and
$K_v=[\square_0]_k\subset C$. Read it through the read-box rule of that theorem without extension:
\begin{equation}\label{10.3:eq-unextended-rule-1}
p_v:=\max\bigl\{p:\ j\le p\le n-2-s(\bar{\mathsf{w}}_v),\ Q^{(p)}\subset Z\bigr\}.
\end{equation}
If the maximum exists, $v$ is \emph{smooth}, and we put $a_v:=n-p_v\ge 3$ and $b_v:=p_v-j$; its
smoothing box is $Q^{(p_v)}$. Otherwise $v$ is \emph{rough}. Then the following hold.
\begin{enumerate}
\item[(a)] Either $v$ is smooth with $a^*<a_v\le\nu-1$ (a \emph{rim core}), or
$Q^{(j)}=\square^{\sim\nu+3}\not\subset Z$ (a core \emph{rough by position}).
\item[(b)] Let $\Delta'$ be a $\pi_k$-cube of $C$ with a face on $\partial C$ (a \emph{collar
$k$-cell}). The bounds that this rule gives for the level-$j$ cores in $\Delta'$ are as follows.
They are stated after division by $|\lambda|\sup|\mathsf{h}|$. The rim cores contribute at most
\begin{equation}\label{10.3:eq-unextended-rule-2}
c_{\mathrm{rim}}\,M^4C_V\,\mathsf{u}_j\,\nu^3(2\nu+9)^2(\nu+4)\,\bar{\rho}^2_{\nu-1}(1+\bar{\rho}_0)\,
L^{5-(\nu-1)},
\end{equation}
where
\begin{equation*}
c_{\mathrm{rim}}:=72d\,c_{\epsilon}^2\bigl(1-L^{-(2\hat{\alpha}-1)}\bigr)^{-1}.
\end{equation*}
The cores rough by position contribute at most
\begin{equation}\label{10.3:eq-unextended-rule-3}
2d\,M^4C^{\mathrm{mar}}\,\mathsf{u}_j\,(\nu+4)\nu^2\bar{\rho}^2_\nu\,L^{-(\nu-1)}.
\end{equation}
The ratio of either bound to the per-cell scale weight $\varsigma^V_\nu\propto
L^{-2\hat{\alpha}(\nu-1)}$ is
$\asymp\operatorname{poly}(\nu)\,L^{(2\hat{\alpha}-1)(\nu-1)+5}$, which is unbounded in $\nu$.
\item[(c)] Summed over the levels $j$ against the profile
$\mathsf{u}_j\le\mathsf{v}_0(\vartheta_2/L)^j$, the bounds of (b) per collar $k$-cell exceed the
target $C_L^V S^V\mathsf{v}_0(\vartheta_2/L)^k$ of the $\sharp$-core entry of the table of
Notation~\ref{10.3:not-vector-dial-table} by a factor
$\asymp k^3\bar{\rho}_k^2\vartheta_2^{-k}$ when $\vartheta_2<1$. If $\vartheta_2=1$ were admissible,
the loss would still be polynomial in $k$.
\end{enumerate}
The loss in (c) is bounded by no constant fixed before $\gamma$ and is not a power of a logarithm;
the ratio in (b) is unbounded in $\nu$.
Neither is removed by any lower bound on $L$ or by any upper bound on the coupling through $\gamma$.
The unextended rule therefore gives no bound for the inner-collar cores of the per-cell form that the
table of Notation~\ref{10.3:not-vector-dial-table} requires.
\end{remark}

\begin{proof}
\emph{The read-box condition.} For an interior unit of the correlation theorem the live-piece datum
$\mathcal{P}_v$ is empty, and no transport parameter $\mathsf{t}$ enters. The pieces of $v$ are the
tuples $(v;\mathsf{y}_1,\dots,\mathsf{y}_4)$ with start label $\operatorname{lab}_0:=K_v$ and reading
set $\mathcal{R}_0:=\hat{X}_v\cup X_v$. Each is bounded in modulus by
$\|v\|_{\star}e^{-(\kappa_1-1)\sum_{i=1}^{4}|\mathsf{y}_i|}$, where, the unit being shallow,
$\|v\|_{\star}=\mathfrak{N}^{\flat}(v)$.

The correlation theorem states the following holomorphy property: for every
$v\in\mathfrak{L}^{\mathrm{pt}}$ the map $m\mapsto f_v(m)$ is holomorphic on $\mathfrak{P}$ and
satisfies $|f_v(m)-f_v(1)|\le\mathfrak{N}^{\flat}(v)$, where $1$ is the point of $\mathfrak{P}$ so
written in the correlation theorem. That theorem proves it on $C$, from four of its
inputs: the curvature bound, the cube-gauge bound, the regularity statement, and the marginal and
interior size bounds. Each of these controls the carrier $\operatorname{chain}(m)$ on cells of $C$
only; the curvature bound is read at the auxiliary consumer
$\mathfrak{E}^C=\mathbb{C}^{(N)}_k(C)$. Consequently the core can be read at the extension reading
$\operatorname{ext}_p$ only on a box $Q^{(p)}\subset C$. This is the content of the condition
$Q^{(p)}\subset Z$ in \eqref{10.3:eq-unextended-rule-1}.

For a smooth core, the interior size bound is replaced by the sharp reading bound of the vector
package at $\operatorname{ext}_{p_v}$. On $Q^{(p_v)}\subset C$ the hypotheses of that bound are
supplied by the curvature bound, the cube-gauge bound and the regularity statement together with the
gauge-smallness bound of the correlation theorem, exactly as for the cores of the bulk.

\emph{Proof of (a).} Since $\bar{\mathsf{w}}_v=3$, the upper end of the range in
\eqref{10.3:eq-unextended-rule-1} is $n-2-s(3)=n-a^*=p^*$. The boxes $Q^{(p)}$ increase with $p$:
each $\pi_{p+1}$-cube is a union of $\pi_p$-cubes, so $[\square_0]_{p+1}$ with two layers of
$\pi_{p+1}$-cubes contains $[\square_0]_p$ with two layers of $\pi_p$-cubes.
Hence $Q^{(p)}\supseteq Q^{(p^*)}$ for every $p\ge p^*$.

By the defining condition of the inner-collar cores (Notation~\ref{10.3:not-inner-collar-cores}),
$Q^{(p^*)}(y)\not\subset C$. Since $Q^{(p^*)}(y)$ is a box containing $\square_0\subset C$ and $C$ is a
component of $Z$, it follows that $Q^{(p^*)}(y)\not\subset Z$. So no $p\ge p^*$ satisfies
$Q^{(p)}\subset Z$. This leaves two cases.
\begin{itemize}
\item The maximum $p_v$ exists with $j\le p_v<p^*$. Then $a_v=n-p_v>n-p^*=a^*$ and
$a_v\le n-j=\nu-1$, so $v$ is a rim core.
\item No $p\ge j$ satisfies $Q^{(p)}\subset Z$. Since the boxes increase with $p$, this happens
exactly when the smallest of them, $Q^{(j)}=\square^{\sim\nu+3}$, is not contained in $Z$, so $v$ is
rough by position.
\end{itemize}

\emph{Sizes of single cores.} All sizes below are stated after division by
$|\lambda|\sup|\mathsf{h}|$.

Consider first a rim core, with $a:=a_v$ and $b:=b_v=p_v-j=\nu-1-a$. The sharp reading bound gives
\begin{equation}\label{10.3:eq-unextended-rule-4}
|f_v(m)-f_v(1)|\le 4\epsilon_a^2\,\mathfrak{N}^V_{n-a}.
\end{equation}
The gauge-smallness bound of the correlation theorem at weight $\bar{\mathsf{w}}_v=3$ gives
\begin{equation*}
\epsilon_a\le 3c_{\epsilon}(2a+11)\bar{\rho}_aL^{-2(a-1)}.
\end{equation*}
The size furnished by the sharp reading bound at level $n-a$ is
\begin{equation*}
\mathfrak{N}^V_{n-a}=C_V\,\mathsf{u}_j\,\nu_p^3\,\bar{\rho}^2_b(1+\bar{\rho}_b)\,
L^{-(4+2\hat{\alpha})b}.
\end{equation*}
Here $\nu_p$ is the age parameter of the sharp reading bound at level $p$, which satisfies
$\nu_p\le\nu$, and $C_V$ is the constant of that bound.

A core rough by position has the marginal size $C^{\mathrm{mar}}\mathsf{u}_j\nu^2\bar{\rho}^2_\nu t^4$
of the vector package. The correlation theorem counts these cores as follows. In a collar $k$-cell
(defined below) the condition $\square^{\sim\nu+3}\not\subset Z$ holds at no more than
$2d(\nu+4)M^4t^{-3}$ anchor points, and the resulting series converges:
\begin{equation*}
\sum_\nu(\nu+4)\nu^2\bar{\rho}_\nu^2L^{-(\nu-1)}<\infty.
\end{equation*}

\emph{Counts per collar $k$-cell.} Fix a collar $k$-cell $\Delta'$; it has at most $2d$ faces on
$\partial C$. Fix the level $j$, so that $t=L^{-(\nu-1)}$. Let $T^{(j)}$ be the set of level-$j$
anchor points of the correlation theorem. By that theorem, $\Delta'$ contains $M^4t^{-4}$ points of
$T^{(j)}$.
\begin{itemize}
\item[(i)] \emph{The whole sub-class.} By Lemma~\ref{10.3:lem-core-geometry} (see
\eqref{10.3:eq-core-geometry-a}), the box $Q^{(p^*)}(y)$ has side at most $2a^*+11$
$\pi_{p^*}$-cubes, and $y$ lies in its central cube $\square_0$. Hence $Q^{(p^*)}(y)\not\subset C$
forces $\operatorname{dist}(y,\partial C)\le(a^*+6)ML^{-a^*}$ sites. Along one face this is a slab of
thickness $(a^*+6)ML^{-a^*}$, and the points of $\Delta'$ in such slabs number at most
\begin{equation*}
2d\,(a^*+6)ML^{-a^*}t^{-1}\,(Mt^{-1})^3=2d(a^*+6)L^{-a^*}\cdot M^4t^{-4}.
\end{equation*}
This is a fraction of the points of the cell that does not depend on $\nu$, and it is all that the
thinness of the collar yields. Let $j_{\min}$ denote the least level of a core. The values of $\nu$
present are all $\nu$ with $a^*+1\le\nu\le k+1-j_{\min}$, and their number is unbounded in $k$.
\item[(ii)] \emph{Rim cores with $a_v=a$.} Since $p_v=n-a<p^*$, the level $n-a+1\le p^*$ lies in the
range of \eqref{10.3:eq-unextended-rule-1}. By maximality of $p_v$, therefore,
$Q^{(n-a+1)}(y)\not\subset Z$. The argument of (i), with $a^*$ replaced by $a-1$, shows that this
forces $\operatorname{dist}(y,\partial C)\le(a+5)ML^{-(a-1)}$. Such $y$ number at most
$2d(a+5)L^{-(a-1)}M^4t^{-4}$ points of $\Delta'$. This is of the same order as the correlation
theorem's own localisation: $y$ within $(\nu+2)L^{-(\nu-1)}+3L^{-(a_v-1)}$ level-$n$ units of a face
of $\partial Z$.
\item[(iii)] \emph{Rough by position.} By the count of the correlation theorem recalled above, there
are at most $2d(\nu+4)M^4t^{-3}$ such points.
\end{itemize}

\emph{Kernel per collar $k$-cell at level $j$.} For the rim cores, multiply the count (ii), in which
$t^{-4}=L^{4(\nu-1)}$, by the size \eqref{10.3:eq-unextended-rule-4}. Two substitutions enter:
\begin{equation*}
4\epsilon_a^2\le 36c_{\epsilon}^2(2a+11)^2\bar{\rho}_a^2L^{-4(a-1)},
\qquad \nu_p^3\le\nu^3.
\end{equation*}
Summing over $a$ gives at most
\begin{align*}
\sum_{a^*<a\le\nu-1}
&2d(a+5)L^{-(a-1)}M^4L^{4(\nu-1)}\cdot 36c_{\epsilon}^2(2a+11)^2\bar{\rho}_a^2L^{-4(a-1)}\\
&\cdot C_V\mathsf{u}_j\nu^3\bar{\rho}^2_{\nu-1-a}(1+\bar{\rho}_{\nu-1-a})
L^{-(4+2\hat{\alpha})(\nu-1-a)}.
\end{align*}
The exponent of $L$ in the $a$-th term is
\begin{align*}
4(\nu-1)-5(a-1)-(4+2\hat{\alpha})(\nu-1)+(4+2\hat{\alpha})a
&=-2\hat{\alpha}(\nu-1)-5a+5+(4+2\hat{\alpha})a\\
&=-2\hat{\alpha}(\nu-1)+(2\hat{\alpha}-1)a+5.
\end{align*}

Now $2\hat{\alpha}-1>0$. Indeed $\hat{\alpha}>\frac12$ is the one-sided choice of the free H\"older
exponent made in the vector package. It is the choice under which the $\sharp$-core entry of its
table carries the factor $t^{4+2\hat{\alpha}}$, and under which its profile window
$\log_L(1/\vartheta_V)<2\hat{\alpha}<2$ is written.

The powers of $L$ in the $a$-sum therefore form a geometric progression with ratio
$L^{2\hat{\alpha}-1}>1$. The sum is dominated by its top term $a=\nu-1$, times
$(1-L^{-(2\hat{\alpha}-1)})^{-1}$ and the polynomial in the radii $\bar{\rho}$. At $a=\nu-1$ the
exponent equals
\begin{equation*}
-2\hat{\alpha}(\nu-1)+(2\hat{\alpha}-1)(\nu-1)+5=5-(\nu-1).
\end{equation*}
Also at $a=\nu-1$ one has $a+5=\nu+4$, $2a+11=2\nu+9$ and $\nu-1-a=0$. With $c_{\mathrm{rim}}$ as
in (b), this gives \eqref{10.3:eq-unextended-rule-2}: a flat kernel $L^{-(\nu-1)}$.

Compare the per-cell $\sharp$-core entry of the table of Notation~\ref{10.3:not-vector-dial-table}:
\begin{equation*}
M^4t^{-4}\,C_{\sharp}\mathsf{u}_j\nu^3\bar{\rho}^2(1+\bar{\rho})\,t^{4+2\hat{\alpha}}
=M^4C_{\sharp}\mathsf{u}_j\nu^3\bar{\rho}^2(1+\bar{\rho})\,L^{-2\hat{\alpha}(\nu-1)}.
\end{equation*}

For the cores rough by position, multiply the count (iii) by the marginal size:
\begin{equation*}
2d(\nu+4)M^4t^{-3}\cdot C^{\mathrm{mar}}\mathsf{u}_j\nu^2\bar{\rho}^2_\nu t^4
=2dM^4C^{\mathrm{mar}}\mathsf{u}_j(\nu+4)\nu^2\bar{\rho}^2_\nu L^{-(\nu-1)}.
\end{equation*}
This is \eqref{10.3:eq-unextended-rule-3}, again flat.

The ratio of either bound to $\varsigma^V_\nu\propto L^{-2\hat{\alpha}(\nu-1)}$ is
$\asymp\operatorname{poly}(\nu)L^{(2\hat{\alpha}-1)(\nu-1)+5}$, which is unbounded in $\nu$. This
proves (b).

The mechanism is the following. The thinness factor $L^{-(a-1)}$ of the count fails to beat the loss
$L^{2\hat{\alpha}a}$ of the H\"older gain, and a core close to $\partial C$ cannot take that gain: its
smoothing box must stay inside $C$, so $b_v$ is small.

\emph{Sum over levels against the profile.} The profile of the vector package gives
$\mathsf{u}_j\le\mathsf{v}_0\vartheta_V^j$ with $\vartheta_V=\vartheta_2/L$, and $\bar{\rho}_\nu^2$
grows at most polynomially in $\nu$, by the growth bound for the radii sequence $\bar{\rho}$ of the
correlation theorem. Per collar $k$-cell, at $j=k+1-\nu$,
\begin{equation*}
(\vartheta_2/L)^jL^{-(\nu-1)}=(\vartheta_2/L)^k\,\vartheta_2^{-(\nu-1)}.
\end{equation*}
By (b), the sum over levels is therefore of the form
\begin{equation*}
\mathsf{v}_0(\vartheta_2/L)^kL^5\sum_{a^*<\nu\le k+1}\operatorname{poly}(\nu)\,
\bar{\rho}^2_\nu\,\vartheta_2^{-(\nu-1)}.
\end{equation*}

For $\vartheta_2<1$ the $\nu$-sum is $\asymp\operatorname{poly}(k)\bar{\rho}_k^2\vartheta_2^{-k}$.
Relative to the target $C_L^VS^V\mathsf{v}_0(\vartheta_2/L)^k$ of the $\sharp$-core entry, the loss is
$\asymp k^3\bar{\rho}_k^2\vartheta_2^{-k}$. This grows with $k$, hence is bounded only in terms of
$K$; it cannot be absorbed into any
constant fixed before $\gamma$, and it is not a power of a logarithm. If $\vartheta_2=1$ were
admissible, the loss would still be $\operatorname{poly}(k)$, a count of steps. This proves (c).

The thinness of the collar could help only if the profile bound were an average over the cells of
$C$ rather than a bound per cell. The table of the vector package
(Notation~\ref{10.3:not-vector-dial-table}) is stated per unit $n$-cell.

Neither loss is removed by a lower bound on $L$ or by an upper bound on the coupling through
$\gamma$: each is a power of $L$ or of $\vartheta_2^{-1}$ whose exponent grows with $\nu$ or with $k$.

\emph{Two further observations.}
\begin{itemize}
\item[($\alpha'$)] Suppose the interior grade is kept and only the start label is enlarged to
$K_v\cup\mathfrak{h}_v$, with the carrier controlled at the consumer $\mathfrak{E}^{\oplus}_v$. This is
the extended treatment without its read-box rule. The interior ceiling of the correlation theorem is
\begin{equation*}
\varepsilon^{\mathrm{int}}(k):=e^{-(\check{R}_k-2)}(1+C_{\Sigma}C_e)^4\,
\mathsf{Q}^{\star,\mathrm{pl}}_{a+9},
\end{equation*}
with $a:=\tfrac{31}{20}\kappa$ the constant so fixed in the correlation theorem (not the rim index
used above), and where
\begin{equation*}
\mathsf{Q}^{\star,\mathrm{pl}}_r:=\sup_{Q\in\pi_k}\sum_{K_v\ni Q}\|v\|_{\star}e^{r\,d_k(K_v)}.
\end{equation*}
This ceiling is anchored at $K_v$. The flat factor of the enlarged label, and its anchoring at the
cubes of $\mathfrak{h}_v\setminus K_v$, must therefore be accounted for separately. They can be
carried only by the declaration of the inner-collar cores as a class of their own
(Notation~\ref{10.3:not-inner-collar-cores}), with the enlarged per-cell collar constant
$\mathsf{K}^{\flat,\mathrm{ex}}_{\mathrm{pl}}$. This intermediate treatment therefore requires of the
correlation theorem the same addition as the extended treatment, and coincides with it.
\item[($\beta'$)] Without that extension, no reading of the core at $p^*$ is available at the interior
grade. On the cells of $Q^{(p^*)}(y)\setminus C$ the carrier $\operatorname{chain}(m)$ is deformed by
the link interpolation parameters $\mathsf{s}$, through the families over the cells disjoint from
the link's domain $\Lambda$ or from its enlargement $\Lambda^{\sim}$.
It is controlled only at a consumer whose domain lies in
$Y_4$ of the piece. For the tuples of an interior core without the extension, with
$\operatorname{lab}_4\subset C$, one has $Y_4=C$, and $C\not\supseteq Q^{(p^*)}(y)\setminus C$. This
excludes the H\"older gain at the interior grade under the rule \eqref{10.3:eq-unextended-rule-1}.
\qedhere
\end{itemize}
\end{proof}

\begin{lemma}[The scale-reading factor lowers only coupling ceilings]\label{10.3:lem-scale-reading-factor}
The notation is that of Notation~\ref{10.3:not-vector-dial-table} and of
Lemma~\ref{10.3:lem-core-size-bound}. Let $v\in\mathfrak{L}^{V,\mathrm{col}}$ be a core of level $j$ and zone
$n$, with site $y$. Write $\nu=1+n-j$ for its age and $t=L^{j-n}=L^{-(\nu-1)}$. Put $s^*=s(3)$, $a^*=2+s^*$,
$p^*=n-a^*$, and
\begin{equation*}
b^*=p^*-j=\nu-1-a^* ,\qquad \nu_{p^*}=1+b^*=\nu-a^* ,\qquad t_{p^*}=L^{j-p^*}=L^{-b^*} .
\end{equation*}
Here $j\le p^*$, as in the definition of the class $\mathfrak{L}^{V,\mathrm{col}}$, and $p^*<n$, that is,
$a^*\ge1$. Assume the following.
\begin{enumerate}
\item[(a)] The constants of the correlation theorem are chosen in its order of choice. In that order $L$ comes
before $\kappa$, and $\kappa$ is subject only to one-sided lower bounds, among them $\kappa\ge 10\log L$ and
$\kappa\ge\kappa_d$. After $\kappa$ come $\kappa_1$, $M_1$, $M$ and the remaining analysis constants, and the
coupling ceiling $\gamma$ is chosen last. In particular $L$, $\kappa$, $c_l$, $c_{\mathfrak{h}}$, $c_{\epsilon}$,
the radii $\bar{\rho}_b$ and $a^*$ are fixed before $\gamma$, and the ceiling $\gamma_V$ of the vector package
is a coupling ceiling chosen after them. The constants $c_{\epsilon}$, $\bar{\rho}_b$, $a^*$, $\hat{\alpha}$,
$C_V$, $C_{\sharp}$ and $C^{\mathrm{mar}}$ are fixed after $L$ and depend neither on $\kappa$ nor on $K$, and
the table constant $C_L^V$ is fixed after $L$ and does not depend on $K$. The flow number $c_2$, the number
$\lambda_0$ and the exponent $q$ of (g) are fixed before $\gamma$. The H\"older exponent $\hat{\alpha}$ of the
vector package satisfies $\hat{\alpha}>\tfrac12$.
\item[(b)] The gauge-smallness bound of the correlation theorem at weight $\bar{\mathsf{w}}_v=3$ and age
parameter $a^*$, applied to the field $\mathcal{U}(m)$ on the box $Q_v=Q^{(p^*)}$ for every $m\in\mathfrak{P}$,
gives the gauge smallness
\begin{equation*}
\epsilon_{a^*}:=3c_{\epsilon}(2a^*+11)\,\bar{\rho}_{a^*}\,L^{-2(a^*-1)}\le 1 .
\end{equation*}
If the gauge-smallness bound is taken with the dilated weight $\tfrac98\bar{\mathsf{w}}$, that is, with
$s^{(9/8)}(\bar{\mathsf{w}}):=s(\tfrac98\bar{\mathsf{w}})$ in place of $s$, the gauge smallness is
$\tfrac98\epsilon_{a^*}$; it is at most $1$ at $a^*=2+s^{(9/8)}(3)$, and its square is
$\tfrac{81}{64}\epsilon_{a^*}^2$.
\item[(c)] Lemma~\ref{10.3:lem-core-size-bound} gives, for every $m\in\mathfrak{P}$,
$|f_v(m)|\le 2\epsilon_{a^*}^2\,\mathfrak{N}^{V,\mathrm{run}}_{p^*}(v)$ and $f_v(1)=0$, where
\begin{equation}\label{10.3:eq-scale-reading-factor-1}
\mathfrak{N}^{V,\mathrm{run}}_{p^*}(v)=\mathsf{u}_j\sup|\mathsf{h}|\,
\bigl[C_{\sharp}\,\nu_{p^*}^3\rho_{p^*,j}^2(1+\rho_{p^*,j})+2K(d)(1+8C_{\mathfrak{J}})\bigr]\,
t_{p^*}^{4+2\hat{\alpha}} ,
\end{equation}
$\rho_{p^*,j}$ is the run ratio of (g), and $K(d)$, $C_{\mathfrak{J}}$ are the constants of the vector package
for which $C_V=C_{\sharp}+2K(d)(1+8C_{\mathfrak{J}})$: $K(d)$ bounds the weighted sum over connected polymers
containing a given cube, and $C_{\mathfrak{J}}$ is the constant in the norm bound $(1+8C_{\mathfrak{J}})\mathsf{u}_j$
of the $\sharp$-family. In that lemma the single-scale reading lemma of the vector package is applied at $p=p^*$
on the box $Q_v$, with $\mathfrak{n}=\mathsf{u}_j$, $e=\mathsf{h}(y)$ and $\epsilon=\epsilon_{a^*}$, the core term
carrying the run ratio $\rho_{p^*,j}$ rather than its majorant $\bar{\rho}_{b^*}$. The size bound of $v$ is
$\mathfrak{N}^{\flat}(v):=|\lambda_{n,k'}|\,4\epsilon_{a^*}^2\,\mathfrak{N}^{V,\mathrm{run}}_{p^*}(v)$; it
majorises $|\lambda_{n,k'}f_v(m)-\lambda_{n,k'}f_v(1)|$ on $\mathfrak{P}$.
\item[(d)] The radii $\bar{\rho}_b$ of the correlation theorem are positive and nondecreasing in $b$.
\item[(e)] The young cores, those of age at most $a^*$, are estimated by the marginal bound of the correlation
theorem; for a young core of level $j'$, of age $\nu'=1+n-j'\le a^*$ and with $t'=L^{j'-n}$, that bound is, per
unit of $\sup|\mathsf{h}|$, $9C^{\mathrm{mar}}\mathsf{u}_{j'}\nu'^2\bar{\rho}_{\nu'}^2t'^4$.
\item[(f)] The recursion lemma of the vector package holds, together with its four ceilings on $g$. Each ceiling
is one-sided and chosen after $L$, the auxiliary parameters, $\hat{\alpha}$, $\vartheta_2$ and $q$, and its left
member tends to $0$ with $g$, uniformly in $K$. The first and the fourth ceiling read
\begin{equation*}
(\bar{\varkappa}+c')S^V\le\tfrac18,\qquad 4S^V(\bar{\varkappa}+c')\le\vartheta_V ,
\end{equation*}
respectively, where $\vartheta_V=\vartheta_2/L$. Under its hypotheses, among which are the four ceilings and the
finiteness $S^V<\infty$, the
recursion lemma gives $T^V_n\le\mathsf{A}_V\vartheta_V^n$ for all $1\le n\le K'$, where $K'$ is the last index of
its recursion. The numbers $\mathsf{A}_V$, $S^V$, $\varsigma^V_\nu$ and $\vartheta_V$ do not depend on $C_L^V$,
and, apart from the left members of the four ceilings, no number entering the recursion lemma depends on
$C_L^V$. The numbers $\mathsf{A}_V$ and $S^V$, and the number $E_1$ that fixes $\mathbb{D}_V$, do not depend on
$K$. The vector package prescribes $\vartheta_1$
in $\mathopen{]}\max\{L^{-(2\hat{\alpha}-1)},L^{-1/2}\},1\mathclose{[}$ and $\vartheta_2$ in
$\mathopen{]}L^{-(2\hat{\alpha}-1)},\vartheta_1\mathclose{]}$; in particular $\vartheta_V=\vartheta_2/L<1$.
Finally, $S^V<\infty$ for every $\vartheta_2>L^{-(2\hat{\alpha}-1)}$, and the load satisfies
\begin{equation*}
\sum_{1\le j'\le n}\varsigma^V_{1+n-j'}\,\mathsf{u}_{j'}\le T^V_n .
\end{equation*}
\item[(g)] The coupling $g$ and the ceiling $\gamma$ satisfy $g\le\gamma\le\gamma^\circ$, where
$\gamma^\circ:=(8c_2+4)^{-1/2}$ and $c_2\ge0$ is the flow number of the correlation theorem. The couplings
$g_i$, $0\le i\le K$, of the run along which $v$ is read, with $x_i=g_i^{-2}$, satisfy $x_n>g^{-2}$ and, with a
number $\lambda_0\ge0$,
\begin{equation}\label{10.3:eq-scale-reading-factor-2}
\lambda_0(i'-i)-2c_2\le x_i-x_{i'}\le 3\lambda_0(i'-i)+2c_2,\qquad 0\le i<i'\le K .
\end{equation}
The radius function $\alpha_\cdot$ of the correlation theorem has the form
$\alpha_\cdot(g')=c\,g'(\log g'^{-2})^q$ with a constant $c>0$ and a positive integer $q$, and the run ratio is
$\rho_{i,i'}:=\alpha_\cdot(g_i)/\alpha_\cdot(g_{i'})$.
\item[(h)] For a $\sharp$-core of level $j'$ and zone $n$, of age $\nu'=1+n-j'$ and with $t'=L^{j'-n}$, the entry
of the first row of the table of the vector package, per unit of $\sup|\mathsf{h}|$, is
\begin{equation*}
\mathfrak{N}_1=C_{\sharp}\,\mathsf{u}_{j'}\,\nu'^3\rho_{n,j'}^2(1+\rho_{n,j'})\,t'^{4+2\hat{\alpha}},
\end{equation*}
and the column of that row is
\begin{equation*}
\varsigma^V_{\nu'}=\nu'^3\rho_{n,n+1-\nu'}^2(1+\rho_{n,n+1-\nu'})\,L^{-2\hat{\alpha}(\nu'-1)} .
\end{equation*}
We write $\mathfrak{N}_1(v)$ for this entry at $v$. For a collar $\pi_k$-cell $\Delta'$ of zone $n=k$ and a level
$j'\le n$, the members of $\mathfrak{L}^{V,\mathrm{col}}$ of level $j'$ with site in $\Delta'$ are at most one
per point, in $\Delta'$, of the set of level-$j'$ points at which $\sharp$-cores are anchored, and $\Delta'$
contains at most $M^4t'^{-4}$ points of that set.
\end{enumerate}
Then the following hold.
\begin{enumerate}
\item[(i)] For $0<\gamma'\le\gamma^\circ$ put
\begin{equation}\label{10.3:eq-scale-reading-factor-3}
E_\rho(\gamma')=(1+4c_2\gamma'^2)\Bigl(1+\frac{(3\lambda_0a^*+2c_2)\gamma'^2}{\log\gamma'^{-2}}\Bigr)^{2q},
\end{equation}
and
\begin{equation*}
\gamma_\rho:=\sup\bigl\{\gamma'\in\mathopen{]}0,\min\{\gamma^\circ,(3/4)^{1/2}e^{-q}\}\mathclose{]}:
E_\rho(\gamma')\le 2\bigr\}.
\end{equation*}
Then $\gamma_\rho>0$. Let $\gamma\le\gamma_\rho$, and put $C_\rho=2^{1/2}$ and $C_x=C_\rho^3=2^{3/2}$. Then
\begin{equation}\label{10.3:eq-scale-reading-factor-4}
\rho_{p^*,j}\le C_\rho\,\rho_{n,j},\qquad \rho_{n,j'}\ge(3/4)^{1/2}\quad(0\le j'\le n),
\end{equation}
and hence
\begin{equation*}
\rho_{p^*,j}^2(1+\rho_{p^*,j})\le C_x\,\rho_{n,j}^2(1+\rho_{n,j}),\qquad
1\le\tfrac34\,\rho_{n,j'}^2(1+\rho_{n,j'})\quad(0\le j'\le n).
\end{equation*}
Put
\begin{equation}\label{10.3:eq-scale-reading-factor-5}
F_{\mathrm{col}}:=36\,c_{\epsilon}^2(2a^*+11)^2\,\bar{\rho}_{a^*}^2\,\frac{C_V}{C_{\sharp}}\,
L^{4+2\hat{\alpha}a^*},\qquad F'_{\mathrm{col}}:=C_xF_{\mathrm{col}}=2^{3/2}F_{\mathrm{col}} .
\end{equation}
Then
\begin{equation}\label{10.3:eq-scale-reading-factor-6}
\mathfrak{N}^{\flat}(v)\le|\lambda_{n,k'}|\sup|\mathsf{h}|\;F'_{\mathrm{col}}\;\mathfrak{N}_1(v) .
\end{equation}
In particular the bound on $\mathfrak{N}^{\flat}(v)$ recorded in \eqref{10.3:eq-vector-dial-table-b} holds with
the factor $F'_{\mathrm{col}}$. Suppose the gauge-smallness bound
is taken with the dilated weight $\tfrac98\bar{\mathsf{w}}$, that is, with $s^{(9/8)}$ in place of
$s$. Then the same holds with $a^*=2+s^{(9/8)}(3)$ and with $F'_{\mathrm{col}}$ multiplied by $81/64$.
\item[(ii)] Let $\gamma\le\gamma_\rho$. The bound of (e) for a young core is at most
$\tfrac{27}{4}C^{\mathrm{mar}}\bar{\rho}_{a^*}^2L^{2\hat{\alpha}(a^*-1)}/C_{\sharp}$ times the entry
$\mathfrak{N}_1$ of (h) for that core. The number
\begin{equation}\label{10.3:eq-scale-reading-factor-7}
F^{*\prime}_{\mathrm{col}}:=\max\bigl\{F'_{\mathrm{col}},\;
\tfrac{27}{4}C^{\mathrm{mar}}\bar{\rho}_{a^*}^2L^{2\hat{\alpha}(a^*-1)}/C_{\sharp}\bigr\}
\end{equation}
depends on $L$, does not depend on $\kappa$, and is fixed after $L$.
\item[(iii)] Replace $C_{\sharp}$ by $C_{\sharp}F^{*\prime}_{\mathrm{col}}$ in the first row of the table of the
vector package. This changes, among the numbers of the recursion lemma, only numbers entering the left members of
its four ceilings.
Among the numbers of the correlation theorem it enters through the per-cell constant
$\mathsf{K}^{\flat,V}$ in $\mathsf{P}_{\mathrm{int}}$. The collar constants
$\mathsf{K}^{\flat,\mathrm{ex}}$, $\mathsf{K}^{\flat,\mathrm{ex}}_{\mathrm{pl}}$ and $\mathsf{K}^{\flat}_+$ are
numbers fixed after $M$. Each of these quantities enters only one-sided ceilings on the coupling, chosen after
$L$, so it can only lower $\gamma_J$ and $\gamma_V$. The comparisons of (i) and (ii) add the one-sided ceiling
$\gamma\le\gamma_\rho$ to the ceilings on the coupling, beside $\gamma\le\gamma^\circ$; the number $\gamma_\rho$
depends only on $c_2$, $q$, $\lambda_0$ and $a^*$, it is chosen after $L$, it can only lower $\gamma$, and no
number fixed before $\gamma$ depends on it. No condition on $L$ arises.
\item[(iv)] The interval $\mathopen{]}L^{-(2\hat{\alpha}-1)},\vartheta_1\mathclose{]}$ of (f), for every
$\vartheta_2$ in which the condition $S^V<\infty$ of the recursion lemma holds (a condition that no ceiling on
$g$ meets), is two-sided in the auxiliary number $\vartheta_2$ only, and it is non-empty whatever the blocking
factor $L$, because $\vartheta_2:=\vartheta_1$ belongs to it. The number $F_{\mathrm{col}}$ does not depend on
$\vartheta_2$, and the interval does not depend on $C_L^V$.
\item[(v)] The per-cell bound \eqref{10.3:eq-proof-collar-sum-a}, in the sum over a collar cell, holds on
$\mathbb{D}_V$. Moreover
\begin{equation*}
\mathbb{D}_V\supseteq\bigl\{(\lambda,\mathsf{h}):|\lambda_{n,k'}|\sup|\mathsf{h}|\le E_1/\mathsf{A}_V\bigr\}
\quad\text{uniformly in } 1\le n\le K' ,
\end{equation*}
so the $\lambda$-window does not close as $n$ grows. Let $\gamma\le\gamma_\rho$, put
$\bar{C}_L^V=\max\{C_{\sharp},C_L^V\}$, and let $\Delta'$ be a collar $\pi_k$-cell of zone $n=k$, with
$1\le n\le K'$. Assume (b) and (c) for every member $v'$ of $\mathfrak{L}^{V,\mathrm{col}}$ of zone $n$ with site
in $\Delta'$, with the level, the site and the age of $v'$ in place of those of $v$. Then, the sum running over
the members $v'$ of $\mathfrak{L}^{V,\mathrm{col}}$ with site in
$\Delta'$ and level between $1$ and $n$,
\begin{equation}\label{10.3:eq-scale-reading-factor-8}
\begin{split}
\sum_{v'}\mathfrak{N}^{\flat}(v')
&\le M^4F'_{\mathrm{col}}\bar{C}_L^V\,|\lambda_{n,k'}|\sup|\mathsf{h}|\,T^V_n\\
&\le M^4F^{*\prime}_{\mathrm{col}}\bar{C}_L^V\,|\lambda_{n,k'}|\sup|\mathsf{h}|\,\mathsf{A}_V\vartheta_V^n ,
\end{split}
\end{equation}
and the constant $M^4F^{*\prime}_{\mathrm{col}}\bar{C}_L^V\mathsf{A}_V$ of the last member does not depend on
$K$, while $\vartheta_V^n\le1$.
\end{enumerate}
\end{lemma}

\begin{proof}
We treat the order of choice first, then (i), (ii), (iii), (iv), and finally (v).

\emph{The order of choice.} By (a), $\kappa$ is fixed after $L$ by one-sided lower bounds. The constants $c_l$,
$c_{\mathfrak{h}}$, $c_{\epsilon}$, the radii $\bar{\rho}_b$ and $a^*=2+s(3)$ are, like $\kappa$, fixed before the
coupling ceiling $\gamma$. The ceiling $\gamma_V$ of the vector package is a coupling ceiling chosen after
them.

The constant $\mathsf{K}^{\flat,\mathrm{ex}}$ of the correlation theorem is
\begin{equation*}
\mathsf{K}^{\flat,\mathrm{ex}}=e^{(a+9)c_lc_{\mathfrak{h}}}\cdot 27S_d'\,\mathsf{K}^{\flat},\qquad
a=\tfrac{31}{20}\kappa,
\end{equation*}
where $\mathsf{K}^{\flat}$ is the per-cell collar constant of the correlation theorem and $27S_d'$ is a numerical
constant fixed by that theorem; $\mathsf{K}^{\flat,\mathrm{ex}}$ is a number fixed
after $M$. It enters the
smallness conditions on the coupling of the correlation theorem in two ways: through $\mathsf{K}^0$, and through
the interior ceiling polynomial
$\mathsf{P}_{\mathrm{int}}$ in the condition on interior pieces of Theorem~\ref{10.3:thm-interior-piece-bound}.
The same holds for $\mathsf{K}^{\flat,\mathrm{ex}}_{\mathrm{pl}}$, whose expanded form is
\eqref{10.3:eq-proof-collar-sum-b}, and for $\mathsf{K}^{\flat}_+$. All three are therefore admissible with
respect to the order of (a). They enter only one-sided conditions, and these can only lower $\gamma_J$ and
$\gamma_V$. This proves the part of (iii) concerning these three constants.

\emph{Proof of (i).} We first show that $\gamma_\rho>0$. Since $(\gamma^\circ)^{-2}=8c_2+4\ge4$, every
$\gamma'\in\mathopen{]}0,\gamma^\circ\mathclose{]}$ satisfies $\gamma'\le\frac12$ and
$\log\gamma'^{-2}\ge\log4>0$. On this interval $E_\rho$ is continuous and nondecreasing: the coefficients $4c_2$
and $3\lambda_0a^*+2c_2$ are nonnegative, $\gamma'^2$ increases with $\gamma'$, and so does
$\gamma'^2/\log\gamma'^{-2}$, since $\log\gamma'^{-2}$ decreases and stays positive. Moreover
$E_\rho(\gamma')\to1<2$ as $\gamma'\to0$. Hence $\gamma_\rho>0$, $E_\rho(\gamma_\rho)\le2$ by continuity, and every
$\gamma\le\gamma_\rho$ satisfies $E_\rho(\gamma)\le2$ and $\gamma\le(3/4)^{1/2}e^{-q}$. From now on
$\gamma\le\gamma_\rho$.

Put $\alpha(x)=x^{-1/2}(\log x)^q$ for $x>1$. If $x_i>1$, then $g_i=x_i^{-1/2}$ and $\log g_i^{-2}=\log x_i$, so
$\alpha_\cdot(g_i)=c\,\alpha(x_i)$ by (g). The constant $c$ cancels in the run ratio:
$\rho_{i,i'}=\alpha(x_i)/\alpha(x_{i'})$ whenever $x_i>1$ and $x_{i'}>1$.

By (g),
\begin{equation*}
x_n>g^{-2}\ge\gamma^{-2}\ge(\gamma^\circ)^{-2}=8c_2+4>8c_2 .
\end{equation*}
Hence $2c_2<x_n/4$, so
$x_n-2c_2\ge\frac34x_n$; also $x_n\ge4$ and $\log x_n\ge\log4>1$. We read
\eqref{10.3:eq-scale-reading-factor-2} in both its members at $(i,i')=(p^*,n)$, where $i'-i=a^*$:
\begin{equation*}
x_n-2c_2\le x_n+\lambda_0a^*-2c_2\le x_{p^*}\le x_n+3\lambda_0a^*+2c_2 .
\end{equation*}
Its lower member at $(i,i')=(j',n)$, for
$0\le j'<n$, gives $x_{j'}\ge x_n+\lambda_0(n-j')-2c_2\ge x_n-2c_2$, and this holds trivially for $j'=n$. In
particular $x_{p^*}\ge\frac34x_n\ge3$ and $x_{j'}\ge3$ for $0\le j'\le n$, so all the run ratios below are given
by $\alpha$.

\emph{The upper comparison.} Since $\rho_{p^*,j}=\alpha(x_{p^*})/\alpha(x_j)$ and
$\rho_{n,j}=\alpha(x_n)/\alpha(x_j)$,
\begin{equation*}
\frac{\rho_{p^*,j}}{\rho_{n,j}}=\frac{\alpha(x_{p^*})}{\alpha(x_n)}
=\Bigl(\frac{x_n}{x_{p^*}}\Bigr)^{1/2}\Bigl(\frac{\log x_{p^*}}{\log x_n}\Bigr)^q .
\end{equation*}
For the first factor we use the lower member: $x_{p^*}\ge x_n-2c_2\ge\frac34x_n$, which is positive, hence
\begin{equation*}
\frac{x_n}{x_{p^*}}\le\frac{x_n}{x_n-2c_2}=1+\frac{2c_2}{x_n-2c_2}\le1+\frac{8c_2}{3x_n}
\le1+\frac{4c_2}{x_n}\le1+4c_2\gamma^2 ,
\end{equation*}
the last step because $x_n>\gamma^{-2}$ and $c_2\ge0$. For the second factor we use the upper member and
$\log(1+u)\le u$ for $u\ge0$:
\begin{equation*}
\log x_{p^*}\le\log(x_n+3\lambda_0a^*+2c_2)=\log x_n+\log\Bigl(1+\frac{3\lambda_0a^*+2c_2}{x_n}\Bigr)
\le\log x_n+\frac{3\lambda_0a^*+2c_2}{x_n} .
\end{equation*}
Both logarithms are positive, as $x_{p^*}\ge3$ and $x_n\ge4$. Dividing by $\log x_n$, and using that
$u\mapsto u\log u$ is increasing on $[1,\infty\mathclose{[}$ while $x_n>\gamma^{-2}\ge4$, we get
\begin{equation*}
\frac{\log x_{p^*}}{\log x_n}\le1+\frac{3\lambda_0a^*+2c_2}{x_n\log x_n}
\le1+\frac{(3\lambda_0a^*+2c_2)\gamma^2}{\log\gamma^{-2}} .
\end{equation*}
Therefore $(\rho_{p^*,j}/\rho_{n,j})^2\le E_\rho(\gamma)\le2$; since $C_\rho^2=2$, this is
$\rho_{p^*,j}\le C_\rho\rho_{n,j}$.

\emph{The lower bound.} The function $\alpha$ is differentiable on $\mathopen{]}1,\infty\mathclose{[}$ with
\begin{equation*}
\alpha'(x)=x^{-3/2}(\log x)^{q-1}\bigl(q-\tfrac12\log x\bigr),
\end{equation*}
so $\alpha$ is nonincreasing on $[e^{2q},\infty\mathclose{[}$. Since $\gamma\le(3/4)^{1/2}e^{-q}$, we have
$\frac34x_n>\frac34\gamma^{-2}\ge e^{2q}$; hence $x_n-2c_2\ge\frac34x_n>e^{2q}$, and $\alpha$ is nonincreasing on
$[x_n-2c_2,\infty\mathclose{[}$. For $0\le j'\le n$ we have $x_{j'}\ge x_n-2c_2$, so
$\alpha(x_{j'})\le\alpha(x_n-2c_2)$ and
\begin{equation*}
\rho_{n,j'}=\frac{\alpha(x_n)}{\alpha(x_{j'})}\ge\frac{\alpha(x_n)}{\alpha(x_n-2c_2)}
=\Bigl(\frac{x_n-2c_2}{x_n}\Bigr)^{1/2}\Bigl(\frac{\log x_n}{\log(x_n-2c_2)}\Bigr)^q\ge(3/4)^{1/2},
\end{equation*}
because $x_n-2c_2\ge\frac34x_n$, and $1<x_n-2c_2\le x_n$ gives $0<\log(x_n-2c_2)\le\log x_n$. This proves
\eqref{10.3:eq-scale-reading-factor-4}.

\emph{Consequences.} The function $x\mapsto x^2(1+x)$ is increasing on $[0,\infty\mathclose{[}$. Since
$C_\rho\ge1$, we
have $1+C_\rho\rho\le C_\rho(1+\rho)$ for $\rho\ge0$, so
\begin{equation*}
\rho_{p^*,j}^2(1+\rho_{p^*,j})\le C_\rho^2\rho_{n,j}^2(1+C_\rho\rho_{n,j})\le C_\rho^3\rho_{n,j}^2(1+\rho_{n,j})
=C_x\rho_{n,j}^2(1+\rho_{n,j}).
\end{equation*}
By the lower bound,
\begin{equation*}
\rho_{n,j'}^2(1+\rho_{n,j'})\ge\frac34\bigl(1+(3/4)^{1/2}\bigr)\ge\frac34\bigl(1+\frac79\bigr)=\frac43 ,
\end{equation*}
where $(3/4)^{1/2}\ge\frac79$ because $\frac{49}{81}\le\frac34$, that is, $196\le243$. Hence
$1\le\frac34\rho_{n,j'}^2(1+\rho_{n,j'})$.

\emph{The majorant.} By (c), $\mathfrak{N}^{\flat}(v)=|\lambda_{n,k'}|\,4\epsilon_{a^*}^2
\mathfrak{N}^{V,\mathrm{run}}_{p^*}(v)$. Since $a^*\ge1$, $\nu_{p^*}=\nu-a^*\le\nu$; moreover $\nu\ge1$ and
$\frac34\le C_x$. By the two consequences just proved,
\begin{equation*}
\begin{split}
C_{\sharp}\nu_{p^*}^3\rho_{p^*,j}^2(1+\rho_{p^*,j})+2K(d)(1+8C_{\mathfrak{J}})
&\le\bigl[C_xC_{\sharp}\nu^3+\tfrac34\,2K(d)(1+8C_{\mathfrak{J}})\bigr]\rho_{n,j}^2(1+\rho_{n,j})\\
&\le C_x\bigl[C_{\sharp}+2K(d)(1+8C_{\mathfrak{J}})\bigr]\nu^3\rho_{n,j}^2(1+\rho_{n,j})\\
&=C_xC_V\,\nu^3\rho_{n,j}^2(1+\rho_{n,j}).
\end{split}
\end{equation*}
Since $t_{p^*}=L^{-b^*}$, \eqref{10.3:eq-scale-reading-factor-1} gives
\begin{equation*}
\mathfrak{N}^{V,\mathrm{run}}_{p^*}(v)\le C_xC_V\,\mathsf{u}_j\sup|\mathsf{h}|\,
\nu^3\rho_{n,j}^2(1+\rho_{n,j})\,L^{-(4+2\hat{\alpha})b^*} .
\end{equation*}
By (b),
$4\epsilon_{a^*}^2=36\,c_{\epsilon}^2(2a^*+11)^2\bar{\rho}_{a^*}^2L^{-4(a^*-1)}$. Hence
\begin{equation}\label{10.3:eq-scale-reading-factor-9}
\begin{split}
\mathfrak{N}^{\flat}(v)&\le|\lambda_{n,k'}|\sup|\mathsf{h}|\;36\,c_{\epsilon}^2(2a^*+11)^2
\bar{\rho}_{a^*}^2\,C_xC_V\\
&\qquad\cdot\mathsf{u}_j\,\nu^3\rho_{n,j}^2(1+\rho_{n,j})\,L^{-4(a^*-1)-(4+2\hat{\alpha})b^*} .
\end{split}
\end{equation}
On the other hand $t=L^{-(\nu-1)}=L^{-(a^*+b^*)}$, because $\nu-1=a^*+b^*$, so by (h)
\begin{equation*}
\mathfrak{N}_1(v)=C_{\sharp}\mathsf{u}_j\nu^3\rho_{n,j}^2(1+\rho_{n,j})L^{-(4+2\hat{\alpha})(a^*+b^*)} .
\end{equation*}
Thus the right
member of \eqref{10.3:eq-scale-reading-factor-9} equals $|\lambda_{n,k'}|\sup|\mathsf{h}|\,\mathfrak{N}_1(v)$
multiplied by
\begin{equation*}
36\,c_{\epsilon}^2(2a^*+11)^2\bar{\rho}_{a^*}^2\,\frac{C_xC_V}{C_{\sharp}}\;
L^{-4(a^*-1)-(4+2\hat{\alpha})b^*+(4+2\hat{\alpha})(a^*+b^*)} .
\end{equation*}
In the exponent the terms in $b^*$ cancel:
\begin{equation}\label{10.3:eq-scale-reading-factor-10}
-4(a^*-1)-(4+2\hat{\alpha})b^*+(4+2\hat{\alpha})(a^*+b^*)=-4a^*+4+(4+2\hat{\alpha})a^*=4+2\hat{\alpha}a^* .
\end{equation}
So the factor equals $C_xF_{\mathrm{col}}=F'_{\mathrm{col}}$ of \eqref{10.3:eq-scale-reading-factor-5}, and this
proves \eqref{10.3:eq-scale-reading-factor-6}.

Now take the gauge-smallness bound with the dilated weight $\tfrac98\bar{\mathsf{w}}$, that is, with
$s^{(9/8)}(\bar{\mathsf{w}}):=s(\tfrac98\bar{\mathsf{w}})$.
By the last sentence of (b), the gauge smallness is then $\tfrac98\epsilon_{a^*}$, which is at most $1$ at
$a^*=2+s^{(9/8)}(3)$, and its square is $\tfrac{81}{64}\epsilon_{a^*}^2$. The computation above, with this $a^*$
throughout and with $\tfrac{81}{64}\epsilon_{a^*}^2$ in place of $\epsilon_{a^*}^2$, is unchanged except for
this factor.
This is the last assertion of (i).

\emph{Proof of (ii).} For a young core write $j'$, $\nu'=1+n-j'\le a^*$ and $t'=L^{-(\nu'-1)}$ for its level,
its age and its number $L^{j'-n}$. By (e) and (h), the ratio of the marginal bound to the
entry $\mathfrak{N}_1$ of that core is
\begin{equation*}
\frac{9C^{\mathrm{mar}}\mathsf{u}_{j'}\nu'^2\bar{\rho}_{\nu'}^2t'^4}
{C_{\sharp}\mathsf{u}_{j'}\nu'^3\rho_{n,j'}^2(1+\rho_{n,j'})t'^{4+2\hat{\alpha}}}
=\frac{9C^{\mathrm{mar}}\bar{\rho}_{\nu'}^2\,t'^{-2\hat{\alpha}}}
{C_{\sharp}\,\nu'\,\rho_{n,j'}^2(1+\rho_{n,j'})} .
\end{equation*}
Here $1/\nu'\le1$ since $\nu'\ge1$; $\bar{\rho}_{\nu'}^2\le\bar{\rho}_{a^*}^2$ by (d), since $\nu'\le a^*$;
$t'^{-2\hat{\alpha}}=L^{2\hat{\alpha}(\nu'-1)}\le L^{2\hat{\alpha}(a^*-1)}$, since $\nu'-1\le a^*-1$; and
$1/(\rho_{n,j'}^2(1+\rho_{n,j'}))\le\frac34$ by the lower bound in (i), which holds for every level $j'\le n$. So
the ratio is at most $\frac{27}{4}C^{\mathrm{mar}}\bar{\rho}_{a^*}^2L^{2\hat{\alpha}(a^*-1)}/C_{\sharp}$.

By (a), both members of the maximum in \eqref{10.3:eq-scale-reading-factor-7} are $L$-dependent numbers fixed
after $L$, in which $\kappa$ does not enter.
Hence $F^{*\prime}_{\mathrm{col}}$ is an $L$-dependent number, independent of $\kappa$, fixed after $L$.

\emph{Proof of (iii).} Replace $C_{\sharp}$ by $C_{\sharp}F^{*\prime}_{\mathrm{col}}$ in the first row of the table
of the vector package. The factor $F^{*\prime}_{\mathrm{col}}$ then enters the following chain:
\begin{itemize}
\item the table constant $C_L^V$;
\item the per-level numbers whose supremum is $\bar{\varkappa}$, each proportional to $C_L^V$, and hence
$\bar{\varkappa}$ itself;
\item through the vector package's numbers built from $C_L^V$, the numbers $\tilde{\varepsilon}_n$ and $c'$.
\end{itemize}
These are the numbers through which $C_L^V$ enters the left members of the four ceilings. By (f), $C_L^V$
enters no other number of the recursion lemma: the numbers $\mathsf{A}_V$, $S^V$, $\varsigma^V_\nu$ and
$\vartheta_V$ do not depend on it.

By (f), each of the four ceilings is a one-sided ceiling on $g$, chosen after $L$, the auxiliary parameters,
$\hat{\alpha}$, $\vartheta_2$ and $q$,
and its left member tends to $0$ with $g$, uniformly in $K$. So a larger $C_L^V$ is met by a smaller $g$, chosen
after $L$, and no condition on $L$ arises.

On the side of the correlation theorem, the per-cell constant $\mathsf{K}^{\flat,V}$ is proportional to
$F^{*\prime}_{\mathrm{col}}\bar{C}_L^V$, where $\bar{C}_L^V=\max\{C_{\sharp},C_L^V\}$ is fixed after $L$ by (a).
It stands in $\mathsf{P}_{\mathrm{int}}$, which enters the condition on interior pieces of
Theorem~\ref{10.3:thm-interior-piece-bound}, after $M$. This is again a ceiling on a coupling, chosen after $L$.

The comparisons of (i) and (ii) use the ceiling $\gamma\le\gamma_\rho$. By its definition in (i), $\gamma_\rho$
depends only on $\gamma^\circ$, which is a function of $c_2$, and on $E_\rho$, whose coefficients are built from
$c_2$, $q$, $\lambda_0$ and $a^*$; by (a) these are fixed before $\gamma$, and $a^*$ after $L$. So
$\gamma\le\gamma_\rho$ is one more one-sided ceiling on the coupling, chosen after $L$, and it can only lower
$\gamma$. No number fixed before $\gamma$ depends on it, since $C_\rho=2^{1/2}$ and $C_x=2^{3/2}$ are pure numbers.
Together with the paragraph on the collar constants under \emph{The order of choice}, this proves (iii). None of
these dependences imposes any condition on $L$.

\emph{Proof of (iv).} By (f), $S^V<\infty$ for every $\vartheta_2>L^{-(2\hat{\alpha}-1)}$. With
$\vartheta_V=\vartheta_2/L$, this lower end is the condition $\log_L(1/\vartheta_V)<2\hat{\alpha}$, since
\begin{equation*}
\log_L(1/\vartheta_V)=1+\log_L(1/\vartheta_2)<2\hat{\alpha}
\quad\Longleftrightarrow\quad \vartheta_2>L^{-(2\hat{\alpha}-1)} .
\end{equation*}
The upper end is $\vartheta_2\le\vartheta_1<1$. Here, by (f), $\vartheta_1$ is the number that the vector
package prescribes in $\mathopen{]}\max\{L^{-(2\hat{\alpha}-1)},L^{-1/2}\},1\mathclose{[}$. Once $\vartheta_2$ is
chosen, the first and the fourth ceiling of (f) are ceilings on $g$ after $(L,\vartheta_2)$.

The lower end of the interval $\mathopen{]}L^{-(2\hat{\alpha}-1)},\vartheta_1\mathclose{]}$ is determined by
$(L,\hat{\alpha})$, and the upper end $\vartheta_1$ is an auxiliary number of the vector package chosen after
$(L,\hat{\alpha})$. So the interval is two-sided in $\vartheta_2$ only; it is not two-sided in $L$, $M$, $K$,
$g$, the parameter $\sigma$ of Theorem~\ref{10.3:thm-interior-piece-bound}, $\zeta$ or the
exponent $p_0$. It is non-empty for every value of the blocking factor. Indeed,
$L>1$ and $\hat{\alpha}>\tfrac12$, by (a), give
$L^{-(2\hat{\alpha}-1)}<1$, and $\vartheta_1$ is prescribed above $L^{-(2\hat{\alpha}-1)}$, so
$\vartheta_2:=\vartheta_1$ is always available. Hence no lower bound on $L$ is needed.

The number $F_{\mathrm{col}}$ of \eqref{10.3:eq-scale-reading-factor-5} does not contain $\vartheta_2$. The
interval depends only on $L$, $\hat{\alpha}$ and $\vartheta_1$, not on $C_L^V$.

\emph{Proof of (v).} The per-cell bound \eqref{10.3:eq-proof-collar-sum-a}, in the sum over a collar cell, holds on
$\mathbb{D}_V$ by the definition of $\mathbb{D}_V$. By Notation~\ref{10.3:not-vector-dial-table},
$\mathbb{D}_V$ is the set of $(\lambda,\mathsf{h})$ with $|\lambda_{n,k'}|\sup|\mathsf{h}|\,T^V_n\le E_1$. By the
recursion lemma of (f), $T^V_n\le\mathsf{A}_V\vartheta_V^n$ for $1\le n\le K'$, and $\vartheta_V=\vartheta_2/L<1$
by (f), because $\vartheta_2\le\vartheta_1<1$ and $L>1$; hence $T^V_n\le\mathsf{A}_V$. So every
$(\lambda,\mathsf{h})$ with $|\lambda_{n,k'}|\sup|\mathsf{h}|\le E_1/\mathsf{A}_V$ lies in $\mathbb{D}_V$, for all
$1\le n\le K'$:
\begin{equation*}
\mathbb{D}_V\supseteq\bigl\{(\lambda,\mathsf{h}):|\lambda_{n,k'}|\sup|\mathsf{h}|\le E_1/\mathsf{A}_V\bigr\}
\quad\text{uniformly in } 1\le n\le K' .
\end{equation*}
The recursion lemma of (f) enters only this inclusion; the per-cell bound \eqref{10.3:eq-proof-collar-sum-a}
itself does not use it.
By (f), the number $E_1$ that fixes the domain $\mathbb{D}_V$ does not depend on $K$.
Thus the $\lambda$-window does not close as $n$ grows.

Now let $\gamma\le\gamma_\rho$ and let $\Delta'$ be as in (v). Fix a level $1\le j'\le n$ and put
$t'=L^{j'-n}$. For every member $v'$ of $\mathfrak{L}^{V,\mathrm{col}}$ of level $j'$ with site in $\Delta'$,
hypotheses (b) and (c) hold by assumption, so \eqref{10.3:eq-scale-reading-factor-6}, applied to $v'$ in place of
$v$, bounds $\mathfrak{N}^{\flat}(v')$ by $|\lambda_{n,k'}|\sup|\mathsf{h}|\,F'_{\mathrm{col}}$ times the entry
$\mathfrak{N}_1$ of (h) at level $j'$; this bound does not depend on the site of $v'$. By (h) there are at most
$M^4t'^{-4}$ such members, and the definition of $\varsigma^V_{\nu'}$ in (h) gives the identity
$M^4t'^{-4}\mathfrak{N}_1=M^4C_{\sharp}\mathsf{u}_{j'}\varsigma^V_{1+n-j'}$. Hence, the sum running over the members
$v'$ of level $j'$ with site in $\Delta'$,
\begin{equation*}
\sum_{v'}\mathfrak{N}^{\flat}(v')\le M^4F'_{\mathrm{col}}\,|\lambda_{n,k'}|\sup|\mathsf{h}|\,
C_{\sharp}\,\mathsf{u}_{j'}\,\varsigma^V_{1+n-j'} .
\end{equation*}
Summing over $1\le j'\le n$, and using the inequality on $T^V_n$ in (f) together with $C_{\sharp}\le\bar{C}_L^V$,
gives the first inequality of \eqref{10.3:eq-scale-reading-factor-8}. The second follows from the bound
$T^V_n\le\mathsf{A}_V\vartheta_V^n$ of (f) and from $F'_{\mathrm{col}}\le F^{*\prime}_{\mathrm{col}}$, which holds
by \eqref{10.3:eq-scale-reading-factor-7}. The constant $M^4F^{*\prime}_{\mathrm{col}}\bar{C}_L^V\mathsf{A}_V$ of its
last member does not depend on $K$ by (a) and (f): $M$ is one
of the analysis constants of (a), $F^{*\prime}_{\mathrm{col}}$ and $\bar{C}_L^V$ are fixed after $L$ and do not
depend on $K$, and $\mathsf{A}_V$ does not depend on $K$; moreover $\vartheta_V^n\le1$ since $\vartheta_V<1$.
\end{proof}

\begin{definition}[Standing notation for tensor insertions]\label{10.4:not-standing-notation}
In this section the following letters keep, without
change, the meanings fixed for the lattice Ward identity earlier in this chapter, in the
setting of the correlation theorem, in the one-lattice Lipschitz step, and at the
introduction of the slot weights.
\begin{enumerate}
\item Lattice and observables: the lattice spacing $a = L^{-K}$; the plaquette action density
$a_p = 1 - \operatorname{Re}\operatorname{tr} U(\partial p)$; the observable $\mathcal{O}(f)$;
the plaquette weight
\begin{equation*}
c(p) = g_0^{-2} - w(p).
\end{equation*}
\item Ward identity: the parametrisation $W_{\ell,p} = w^0 + i\vec w\cdot\vec\sigma$ of the
plaquette variables based at a link $\ell$, valid for the gauge group $SU(2)$, whose
components $w^0$, $\vec w$ are not to be confused with $w(p)$ above or with $w_{k,j}$
below; the vectors
$\vec u_{\ell,\rho}$; the link-plaquette variation $\mathfrak{d}_{\ell,p,\rho}$; the
translation part $\mathfrak{t}_\rho$, the lattice stress tensor $\mathfrak{s}_{\sigma\rho}$
and the orbital remainder $\mathfrak{r}_\rho$; the transition term $\mathfrak{h}_\chi$; the
symmetric, centred and forward lattice differences $\nabla^{s}$, $\nabla^{c}$,
$\nabla^{+}$; the lattice Lie derivative $\mathsf{L}^a_\xi$; the vector field $V_\eta$ on
the configuration space; and the infinitesimal rotation field $\xi = \omega x$.
\item Age letters and step constants:
\begin{equation*}
t = L^{j-n}, \qquad \nu = 1 + n - j,
\end{equation*}
together with $\rho_{n,j}$, $w_{k,j}$, $\theta$, $\varsigma_\nu$, the age constant $c_a$, and
\begin{equation*}
\varkappa_k = K_C\, E\, C_L\, \theta_k, \qquad
\mathfrak{g}_n = g_n^2 (\log x_n)^{2q_\flat},
\end{equation*}
where the factors $K_C$, $E$, $C_L$ of $\varkappa_k$ are the constants fixed in the
one-lattice Lipschitz step;
the read-out letters $E_1$, $\mathbb{D}_\delta$, $\Sigma^q_k[\delta]$.
\item Slot weights: $a_k$, $c_k$, $\varpi_a$, $\omega$, $S_\omega$.
\end{enumerate}
The letter $\omega$ denotes the infinitesimal rotation in $\xi = \omega x$ and, among the
slot weights and in $S_\omega$, the geometric majorant of the slot weights and pure-kernel
constants; the context decides which is meant.
In this section $q_\flat$ denotes the exponent of \cite[(2.28)]{B-CMP119}, in the form in
which it enters the correlation theorem, and $q$ denotes the age rate
\begin{equation*}
q = e^{-c_a},
\end{equation*}
where the constant $c_a$ is recalled together with the age letters above and, like them, is
used without change. In this section the undecorated letter $q$ denotes only this age rate,
and not the auxiliary parameter of the same name listed in the glossary.
\end{definition}

For $\eta\colon\mathbb{R}^4\to\mathbb{R}^4$ and a plaquette function $c$ we use the following conventions. The variable $p$ runs over the plaquettes and $\ell$ over the links of the lattice of spacing $a=L^{-K}$ on the torus $\mathbb{T}_m$. We write $x_p=x(p)$ for the barycentre of $p$ and $x_\ell$ for the midpoint of $\ell$, and $p\pm\hat\rho$ denotes the translate of $p$ by $\pm a$ in direction $\rho$. The letters $a_p$, $\mathfrak{d}_{\ell,p,\rho}$, $\mathfrak{t}_\rho$, $\mathfrak{s}_{\sigma\rho}$, $\mathfrak{r}_\rho$, $\mathfrak{h}_\chi$, $\nabla^{s}$, $\nabla^{+}$ and $V_\eta$ are those of the standing notation for tensor insertions (\ref{10.4:not-standing-notation}).
We write $\|\partial^2\eta\|_\infty$ for the supremum over points, over $\rho$ and over unit vectors $e$ of $|\partial_e^2\eta^\rho|$. Put
\begin{equation}\label{10.4:eq-rotation-field-1}
\mathfrak{K}_\eta[c] := \sum_p c(p)\sum_\rho a^{-1}\eta^\rho(x_p)\,\mathfrak{r}_\rho(p),
\end{equation}
\begin{equation}\label{10.4:eq-rotation-field-2}
\mathfrak{T}_\eta(p) := a^{-1}\sum_{\ell\subset\partial p}\sum_\rho
\bigl(\eta^\rho(x_\ell)-\eta^\rho(x_p)\bigr)\,\mathfrak{d}_{\ell,p,\rho},
\qquad
\mathfrak{T}_\eta[c] := \sum_p c(p)\,\mathfrak{T}_\eta(p).
\end{equation}
Take $\xi=\omega x$ and a cutoff $\chi$. The orbital-remainder functional of the rotational Ward identity for $\xi$ with cutoff $\chi$, there written $\mathfrak{K}_\chi[c]$, coincides with $\mathfrak{K}_{\chi\xi}[c]$. For $\eta=\chi\xi$ one has
\begin{equation*}
\mathfrak{T}_\eta(p)=\chi(x_p)\sum_{\sigma,\rho}\omega_{\rho\sigma}\,\mathfrak{s}_{\sigma\rho}(p)+\mathfrak{h}_\chi(p).
\end{equation*}

\begin{lemma}[The rotation field acting on a weighted plaquette sum]\label{10.4:lem-rotation-field}
Let $\eta\colon\mathbb{R}^4\to\mathbb{R}^4$ have support in $]-L^m+2a,\,L^m-2a[^4$, and let $c$ be a plaquette function. Then the following hold.
\begin{enumerate}
\item[(i)] With $V_\eta S_c := \sum_p c(p)\,V_\eta a_p$,
\begin{equation}\label{10.4:eq-rotation-field-3}
V_\eta S_c = -\sum_p a_p\sum_\rho \nabla^{s}_\rho(c\eta^\rho)(p)
+\mathfrak{K}_\eta[c]+\mathfrak{T}_\eta[c].
\end{equation}
\item[(ii)] Put
\begin{equation*}
\mathfrak{K}^\sharp_\eta[c] := \sum_p c(p)\sum_\rho a^{-1}\eta^\rho(x_p)\,\mathfrak{r}^\sharp_\rho(p).
\end{equation*}
Then
\begin{equation}\label{10.4:eq-rotation-field-4}
\mathfrak{K}_\eta[c] = \mathfrak{K}^\sharp_\eta[c]
-\sum_p\sum_\rho \nabla^{+}_\rho(c\eta^\rho)(p)\,\Psi_\rho(p).
\end{equation}
\item[(iii)] If $\eta\in C^2$, then
\begin{equation}\label{10.4:eq-rotation-field-5}
\mathfrak{T}_\eta(p)=\sum_{\sigma,\rho}(\partial_\sigma\eta^\rho)(x_p)\,\mathfrak{s}_{\sigma\rho}(p)
+\mathfrak{T}^{(2)}_\eta(p),
\end{equation}
and on $SU(2)$-fields
\begin{equation}\label{10.4:eq-rotation-field-6}
\bigl|\mathfrak{T}^{(2)}_\eta(p)\bigr|\le 2a\,\|\partial^2\eta\|_\infty\sum_{p'\in N(p)}a_{p'}.
\end{equation}
Here $N(p)$ is the family of plaquettes neighbouring $p$ used in the lemma on the variation of plaquette variables along $V_\eta$; $p$ itself occurs in it.
\end{enumerate}
\end{lemma}

\begin{proof}
(i) Part (ii) of the lemma on the variation of plaquette variables along $V_\eta$ gives
\begin{equation}\label{10.4:eq-rotation-field-7}
V_\eta a_p = a^{-1}\sum_{\ell\subset\partial p}\sum_\rho \eta^\rho(x_\ell)\,\mathfrak{d}_{\ell,p,\rho}
= a^{-1}\sum_\rho\eta^\rho(x_p)\,\mathfrak{t}_\rho(p)+\mathfrak{T}_\eta(p).
\end{equation}
The second form follows from the first by adding and subtracting $\eta^\rho(x_p)$ under the sum, by the translation part $\mathfrak{t}_\rho(p)=\sum_{\ell\subset\partial p}\mathfrak{d}_{\ell,p,\rho}$, and by the definition \eqref{10.4:eq-rotation-field-2} of $\mathfrak{T}_\eta(p)$. Into the second form we insert the decomposition of the standing notation,
\begin{equation*}
\mathfrak{t}_\rho(p)=\tfrac12\bigl(a_{p+\hat\rho}-a_{p-\hat\rho}\bigr)+\mathfrak{r}_\rho(p).
\end{equation*}
We then multiply by $c(p)$ and sum over $p$. The terms containing $\mathfrak{r}_\rho$ add up to $\mathfrak{K}_\eta[c]$ by \eqref{10.4:eq-rotation-field-1}. The terms containing $\mathfrak{T}_\eta(p)$ add up to $\mathfrak{T}_\eta[c]$.

It remains to treat, for each $\rho$ separately, the function $\varphi(p):=c(p)\eta^\rho(x_p)$. On the torus the translations $p\mapsto p\pm\hat\rho$ are bijections of the set of plaquettes. Hence
\begin{equation*}
\sum_p\varphi(p)a_{p+\hat\rho}=\sum_p\varphi(p-\hat\rho)a_p,
\qquad
\sum_p\varphi(p)a_{p-\hat\rho}=\sum_p\varphi(p+\hat\rho)a_p .
\end{equation*}
Recall that $(\nabla^{s}_\rho\varphi)(p)=(2a)^{-1}\bigl(\varphi(p+\hat\rho)-\varphi(p-\hat\rho)\bigr)$. The two identities above therefore give the summation by parts
\begin{equation*}
\sum_p\varphi(p)\,(2a)^{-1}\bigl(a_{p+\hat\rho}-a_{p-\hat\rho}\bigr)
=-\sum_p a_p\,(\nabla^{s}_\rho\varphi)(p).
\end{equation*}
Summing over $\rho$ produces the first term on the right of \eqref{10.4:eq-rotation-field-3}. This proves (i).

(ii) Fix $\rho$. By the exact-difference lemma for the orbital remainder,
\begin{equation*}
\mathfrak{r}_\rho(p)=\Psi_\rho(p)-\Psi_\rho(p-\hat\rho)+\mathfrak{r}^\sharp_\rho(p).
\end{equation*}
When the plane of $p$ contains $\rho$ the difference is exact, and we set $\mathfrak{r}^\sharp_\rho(p):=0$; with this convention, which is also the one in force in the definition of $\mathfrak{K}^\sharp_\eta[c]$ in (ii), the displayed decomposition holds for every $p$. Inserting this into \eqref{10.4:eq-rotation-field-1}, the terms containing $\mathfrak{r}^\sharp_\rho$ give $\mathfrak{K}^\sharp_\eta[c]$.

For the difference terms, put $\varphi(p):=a^{-1}c(p)\eta^\rho(x_p)$. Translating the second sum by $\hat\rho$ on the torus and using $(\nabla^{+}_\rho\varphi)(p)=a^{-1}\bigl(\varphi(p+\hat\rho)-\varphi(p)\bigr)$, we obtain
\begin{equation*}
\sum_p\varphi(p)\bigl[\Psi_\rho(p)-\Psi_\rho(p-\hat\rho)\bigr]
=\sum_p\bigl(\varphi(p)-\varphi(p+\hat\rho)\bigr)\Psi_\rho(p)
=-a\sum_p(\nabla^{+}_\rho\varphi)(p)\,\Psi_\rho(p).
\end{equation*}
Since $a\,\nabla^{+}_\rho\varphi=\nabla^{+}_\rho(c\eta^\rho)$, this equals $-\sum_p\nabla^{+}_\rho(c\eta^\rho)(p)\Psi_\rho(p)$. Summing over $\rho$ gives \eqref{10.4:eq-rotation-field-4}.

The corollary giving the cubic remainder of the rotational Ward identity is the case of (ii) in which $\partial_\rho\eta^\rho=0$ for each $\rho$.

(iii) For $\ell\subset\partial p$ one has $|x_\ell-x_p|=a/2$. Taylor's formula with remainder gives
\begin{equation*}
\eta^\rho(x_\ell)-\eta^\rho(x_p)=\sum_\sigma(x_\ell-x_p)_\sigma(\partial_\sigma\eta^\rho)(x_p)+R_{\ell,\rho},
\end{equation*}
with $|R_{\ell,\rho}|\le\tfrac12|x_\ell-x_p|^2\|\partial^2\eta\|_\infty=\tfrac18a^2\|\partial^2\eta\|_\infty$. Insert this into \eqref{10.4:eq-rotation-field-2}. The first-order terms contribute
\begin{equation*}
a^{-1}\sum_{\ell\subset\partial p}\sum_{\rho,\sigma}(x_\ell-x_p)_\sigma(\partial_\sigma\eta^\rho)(x_p)\,\mathfrak{d}_{\ell,p,\rho},
\end{equation*}
which, with the lattice stress tensor $\mathfrak{s}_{\sigma\rho}$ of the standing notation, is the first term on the right of \eqref{10.4:eq-rotation-field-5}. The remainder is
\begin{equation*}
\mathfrak{T}^{(2)}_\eta(p)=a^{-1}\sum_{\ell\subset\partial p}\sum_\rho R_{\ell,\rho}\,\mathfrak{d}_{\ell,p,\rho}.
\end{equation*}

Write $\mu(\ell)$ for the direction of $\ell$. On $SU(2)$-fields, part (iv) of the lemma on the variation of plaquette variables along $V_\eta$ gives
\begin{equation*}
\sum_{\ell\subset\partial p}\sum_{\rho\ne\mu(\ell)}|\mathfrak{d}_{\ell,p,\rho}|\le 14\sum_{p'\in N(p)}a_{p'}.
\end{equation*}
The right side counts twelve terms $a_p$ and a contribution $\tfrac12 a_{p'}$ for each adjacent plaquette $p'$, with $p$ itself occurring four times among them. With the bound on $R_{\ell,\rho}$ this gives
\begin{equation*}
\bigl|\mathfrak{T}^{(2)}_\eta(p)\bigr|
\le a^{-1}\cdot\tfrac18a^2\|\partial^2\eta\|_\infty\cdot14\sum_{p'\in N(p)}a_{p'}
=\tfrac74\,a\,\|\partial^2\eta\|_\infty\sum_{p'\in N(p)}a_{p'},
\end{equation*}
which is at most the right side of \eqref{10.4:eq-rotation-field-6}.
\end{proof}

\begin{lemma}[Null directions and the strain-smeared slot terms]\label{10.4:lem-null-directions}
Let $\xi=\omega x$ be the infinitesimal rotation field, put $A:=\sum_p a_p$, and for a vector
field $\eta$ write $\nabla^{s}\cdot\eta:=\sum_\rho\nabla^{s}_\rho\eta^\rho$.
\begin{enumerate}
\item[(i)] Let $w'$ be a complex plaquette function at which the sourced partition function does
not vanish, and suppose
\begin{equation*}
\operatorname{dist}\bigl(\operatorname{supp}\eta,\{x_p:\ w'(p)\neq0\}\bigr)>a/2 .
\end{equation*}
Then $g_0^{-2}\langle V_\eta A\rangle_{w'}=\langle\operatorname{div}V_\eta\rangle_{w'}$. Consequently,
for $m\ge1$ and $f_1,\dots,f_m$ with $\operatorname{dist}(\operatorname{supp}\eta,\operatorname{supp}f_l)>a/2$
for $l=1,\dots,m$,
\begin{equation}\label{10.4:eq-null-directions-a}
\begin{gathered}
\langle\mathsf{N}(\eta);\mathcal{O}(f_1);\dots;\mathcal{O}(f_m)\rangle=0,\\
\mathsf{N}(\eta):=V_\eta A-g_0^{2}\operatorname{div}V_\eta
=\mathfrak{T}_\eta[1]+\mathfrak{K}_\eta[1]-\mathcal{O}(\nabla^{s}\cdot\eta)
-g_0^{2}\operatorname{div}V_\eta ,
\end{gathered}
\end{equation}
where, by the divergence formula for the fields $V_\eta$,
\begin{equation}\label{10.4:eq-null-directions-2}
\operatorname{div}V_\eta=-\frac32\sum_{p=\{x;\mu,\rho\}}a_p\Bigl[(\nabla^{c}_\rho\eta^\rho)(x_p)
+(\nabla^{c}_\mu\eta^\mu)(x_p)\Bigr].
\end{equation}
\item[(ii)] For every $f$, where $\mathfrak{K}[f]$ denotes the slot orbital term of the rotational
Ward identity,
\begin{equation}\label{10.4:eq-null-directions-3}
\mathfrak{K}[f]=\mathfrak{K}_{f\xi}[1],\qquad
\mathfrak{T}_{f\xi}[1]=\mathfrak{S}[f]+\sum_p\mathfrak{h}_f(p),\qquad
\nabla^{s}\cdot(f\xi)=\mathsf{L}^a_\xi f,
\end{equation}
where $\mathfrak{h}_f$ is the transition term of the rotational Ward identity with the weight
$\chi$ replaced by $f$:
\begin{equation}\label{10.4:eq-null-directions-4}
\mathfrak{h}_f(p):=a^{-1}\sum_{\ell\subset\partial p}\sum_\rho\bigl(f(x_\ell)-f(x_p)\bigr)
\xi^\rho(x_\ell)\,\mathfrak{d}_{\ell,p,\rho}.
\end{equation}
\item[(iii)] Let $n\ge2$, and let $f_1,\dots,f_n$ have supports at mutual distances at least
$\mathsf{d}>a/2$. Then for each $i$,
\begin{equation}\label{10.4:eq-null-directions-b}
\begin{split}
\bigl\langle(\mathfrak{S}+\mathfrak{K})[f_i];(\mathcal{O}(f_l))_{l\neq i}\bigr\rangle
&=S^T_n(f_1,\dots,\mathsf{L}^a_\xi f_i,\dots,f_n)
-\Bigl\langle\sum_p\mathfrak{h}_{f_i}(p);(\mathcal{O}(f_l))_{l\neq i}\Bigr\rangle\\
&\quad+g_0^{2}\bigl\langle\operatorname{div}V_{f_i\xi};(\mathcal{O}(f_l))_{l\neq i}\bigr\rangle .
\end{split}
\end{equation}
\item[(iv)] Put $\mathfrak{T}^0_f:=\sum_p\mathfrak{h}_f(p)-\mathcal{O}(\mathsf{L}^a_\xi f)$. Under the
hypotheses of (iii), and with $\chi$ as in the rotational Ward identity,
\begin{equation}\label{10.4:eq-null-directions-c}
\begin{split}
\sum_i S^T_n(\dots,\mathsf{L}^a_\xi f_i,\dots)
&=-g_0^{-2}\bigl\langle(\mathfrak{S}_\chi+\mathfrak{K}_\chi)[1];\mathcal{O}(f_1);\dots;
\mathcal{O}(f_n)\bigr\rangle\\
&\quad+\bigl\langle\mathfrak{B}^{\mathrm{sh}}_\chi;\mathcal{O}(f_1);\dots;\mathcal{O}(f_n)\bigr\rangle\\
&\quad-\sum_i\bigl\langle\mathfrak{T}^0_{f_i}-g_0^{2}\operatorname{div}V_{f_i\xi};
(\mathcal{O}(f_l))_{l\neq i}\bigr\rangle .
\end{split}
\end{equation}
\end{enumerate}
\end{lemma}

\begin{proof}
(i) Let $p$ be a plaquette with $V_\eta a_p\neq0$. Then $\eta(x_\ell)\neq0$ for some link
$\ell\subset\partial p$, and $|x_\ell-x_p|=a/2$. If moreover $w'(p)\neq0$, the hypothesis gives
$\operatorname{dist}(x_p,\operatorname{supp}\eta)>a/2$, hence $\operatorname{dist}(x_\ell,\operatorname{supp}\eta)>0$
by the triangle inequality, which contradicts $\eta(x_\ell)\neq0$. So $w'(p)V_\eta a_p=0$ for every
$p$. The action of the sourced measure at the source $w'$ is $g_0^{-2}A$ plus a linear combination
of the $a_p$ with coefficients $w'(p)$, and $V_\eta$ annihilates the latter; hence $V_\eta$ applied
to this action equals $g_0^{-2}V_\eta A$. The integration-by-parts lemma for the fields $V_\eta$,
applicable because the sourced partition function does not vanish at $w'$, gives that the
expectation at $w'$ of $V_\eta$ applied to this action equals $\langle\operatorname{div}V_\eta\rangle_{w'}$,
and the first assertion follows.

Now let $f_1,\dots,f_m$ be as in (i), and for $z\in\mathbb{C}^m$ let $w'_z(p):=\sum_l z_lf_l(x_p)$.
Then $\{x_p:w'_z(p)\neq0\}\subset\bigcup_l\operatorname{supp}f_l$, so the distance hypothesis holds for
every $z$, and the first assertion, multiplied by $g_0^2$, reads
$\langle\mathsf{N}(\eta)\rangle_{w'_z}=0$ for every $z$ at which the sourced partition function does
not vanish, in particular for all $z$ near the origin, where it does not vanish. The observable
$\mathsf{N}(\eta)$ does
not depend on $z$. Applying $\partial_{z_1}\cdots\partial_{z_m}$ at $z=0$ produces, by the definition
of the connected correlations through the sourced expectations, the left side of the first line of
\eqref{10.4:eq-null-directions-a}, which therefore vanishes.

For the decomposition of $\mathsf{N}(\eta)$ in \eqref{10.4:eq-null-directions-a}, apply
Lemma~\ref{10.4:lem-rotation-field}(i) (the rotation
field acting on a weighted plaquette sum) with $c\equiv1$, so that $S_c=\sum_pa_p=A$, and
\begin{equation*}
V_\eta A=V_\eta S_1=-\sum_pa_p\sum_\rho\nabla^{s}_\rho\eta^\rho(p)+\mathfrak{K}_\eta[1]
+\mathfrak{T}_\eta[1].
\end{equation*}
The first sum on the right is $\sum_pa_p(\nabla^{s}\cdot\eta)(p)=\mathcal{O}(\nabla^{s}\cdot\eta)$.
Subtracting $g_0^2\operatorname{div}V_\eta$ gives that decomposition.

(ii) The first identity is the definition of $\mathfrak{K}[f]$: with $c\equiv1$ and $\eta=f\xi$, the
sum defining $\mathfrak{K}_\eta[c]$ is $\sum_p\sum_\rho a^{-1}f(x_p)\xi^\rho(x_p)\mathfrak{r}_\rho(p)$,
which is $\mathfrak{K}[f]$. For the second, $\mathfrak{T}_{f\xi}[1]$ is the sum over $p$, over
$\ell\subset\partial p$ and over $\rho$ of
$a^{-1}\bigl(f(x_\ell)\xi^\rho(x_\ell)-f(x_p)\xi^\rho(x_p)\bigr)\mathfrak{d}_{\ell,p,\rho}$. Since
$\xi^\rho(x)=\omega_{\rho\sigma}x_\sigma$ is linear,
\begin{equation*}
f(x_\ell)\xi^\rho(x_\ell)-f(x_p)\xi^\rho(x_p)
=f(x_p)\,\omega_{\rho\sigma}(x_\ell-x_p)_\sigma+\bigl(f(x_\ell)-f(x_p)\bigr)\xi^\rho(x_\ell),
\end{equation*}
as one sees by adding and subtracting $f(x_p)\xi^\rho(x_\ell)$. The second term, summed, is
$\sum_p\mathfrak{h}_f(p)$ by \eqref{10.4:eq-null-directions-4}. The first term, multiplied by
$a^{-1}\mathfrak{d}_{\ell,p,\rho}$ and summed over $\ell\subset\partial p$ and over $\rho$, is
$f(x_p)\,\omega_{\rho\sigma}\mathfrak{s}_{\sigma\rho}(p)$ by the definition of the lattice stress
tensor, and $\sum_pf(x_p)\,\omega_{\rho\sigma}\mathfrak{s}_{\sigma\rho}(p)=\mathfrak{S}[f]$ by the
definition of the spin term. For the third identity, $\omega$ is antisymmetric, so
$\omega_{\rho\rho}=0$ and $\xi^\rho$ does not depend on $x_\rho$. The difference $\nabla^{s}_\rho$
combines values at points differing only in the $\rho$-coordinate, so
$\nabla^{s}_\rho(f\xi^\rho)=\xi^\rho\nabla^{s}_\rho f$ for each $\rho$, and
$\sum_\rho\xi^\rho\nabla^{s}_\rho f=\mathsf{L}^a_\xi f$ by the definition of the lattice Lie
derivative.

(iii) Fix $i$ and apply (i) with $\eta=f_i\xi$, $m=n-1\ge1$, and the functions $f_l$, $l\neq i$.
Since $\operatorname{supp}(f_i\xi)\subset\operatorname{supp}f_i$, its distance to each
$\operatorname{supp}f_l$, $l\neq i$, is at least $\mathsf{d}>a/2$. By
\eqref{10.4:eq-null-directions-a},
\begin{equation*}
\bigl\langle\mathfrak{T}_{f_i\xi}[1]+\mathfrak{K}_{f_i\xi}[1]-\mathcal{O}(\nabla^{s}\cdot(f_i\xi))
-g_0^2\operatorname{div}V_{f_i\xi};(\mathcal{O}(f_l))_{l\neq i}\bigr\rangle=0 .
\end{equation*}
Insert the three identities of (ii) with $f=f_i$, and use multilinearity of the connected
correlation:
\begin{equation*}
\begin{split}
\bigl\langle(\mathfrak{S}+\mathfrak{K})[f_i];(\mathcal{O}(f_l))_{l\neq i}\bigr\rangle
&=\bigl\langle\mathcal{O}(\mathsf{L}^a_\xi f_i);(\mathcal{O}(f_l))_{l\neq i}\bigr\rangle
-\Bigl\langle\sum_p\mathfrak{h}_{f_i}(p);(\mathcal{O}(f_l))_{l\neq i}\Bigr\rangle\\
&\quad+g_0^2\bigl\langle\operatorname{div}V_{f_i\xi};(\mathcal{O}(f_l))_{l\neq i}\bigr\rangle .
\end{split}
\end{equation*}
The first term on the right is $S^T_n(f_1,\dots,\mathsf{L}^a_\xi f_i,\dots,f_n)$, by the definition
of the connected Schwinger functions as the connected correlations of the observables
$\mathcal{O}(\cdot)$ and their symmetry in the entries. This is \eqref{10.4:eq-null-directions-b}.

(iv) By the definition of $\mathfrak{T}^0_{f_i}$ and the identification of the first term just made,
\eqref{10.4:eq-null-directions-b} is equivalent to
\begin{equation*}
\bigl\langle(\mathfrak{S}+\mathfrak{K})[f_i];(\mathcal{O}(f_l))_{l\neq i}\bigr\rangle
=-\bigl\langle\mathfrak{T}^0_{f_i}-g_0^2\operatorname{div}V_{f_i\xi};
(\mathcal{O}(f_l))_{l\neq i}\bigr\rangle .
\end{equation*}
Replacing, in the rotational Ward identity, each slot term
$\langle(\mathfrak{S}+\mathfrak{K})[f_i];(\mathcal{O}(f_l))_{l\neq i}\rangle$ by the right side of
this identity gives \eqref{10.4:eq-null-directions-c}.
\end{proof}

The identity \eqref{10.4:eq-null-directions-b} is the exact-vanishing identity for the slot terms,
written in the notation of the rotational Ward identity: the slot terms carry no vector-channel
content of their own. The functional $\mathfrak{T}^0_f$ is the lattice stress tensor smeared with
the weight $\partial_\sigma f\cdot\xi^\rho$, with its trace part removed at tree level. Since
\begin{equation*}
\partial_\sigma(f\xi^\rho)=\partial_\sigma f\cdot\xi^\rho+f\,\omega_{\rho\sigma}
\end{equation*}
and $\omega$ is antisymmetric, the symmetric part of the weight $\partial_\sigma f\cdot\xi^\rho$ is
the \emph{strain} $\tfrac12(\partial_\sigma\eta^\rho+\partial_\rho\eta^\sigma)$ of the compactly
supported field $\eta=f\xi$. The same holds in the shell with $\eta=\chi\xi$. These are the only
tensor weights that the Ward identity produces.

\begin{lemma}[The diagonal lattice stress tensor as a plaquette weight]\label{10.4:lem-diagonal-stress}
Use the standing notation for tensor insertions (Notation~\ref{10.4:not-standing-notation}),
with $SU(2)$-valued link variables, the plaquette variables being written in the Pauli
parametrisation $W_{\ell,p} = w^0 + i\vec w\cdot\vec\sigma$.
Write $e_\rho$ for the unit vector in direction $\rho$ and $\hat\rho = a e_\rho$.
For a plaquette $p$ and a link $\ell$, let $x_p = x(p)$ be the barycentre of $p$ and $x_\ell$ the
midpoint of $\ell$. For a link $\ell$ of direction $\mu$ and a direction $\nu \neq \mu$, let
$p(\ell,\pm\nu)$ be the plaquette spanned by $\ell$ and $\pm\hat\nu$, and let
$\vec w_{\ell,\pm\nu} = \vec w_{\ell,p(\ell,\pm\nu)}$. Let $|\vec w_p|$ denote the common value of
$|\vec w_{\ell,p}|$ over $\ell \in \partial p$, and let $p \pm \hat\nu$ be the translate of $p$
by $\pm\hat\nu$.

Let $p = \{x;\mu,\nu\}$, $\ell_1 := (x,\mu)$, $\ell_3 := (x+\hat\nu,\mu)$. Then
$\mathfrak{s}_{\sigma\rho}(p) = 0$ unless $\sigma \in \{\mu,\nu\}$;
$\mathfrak{s}_{\nu\mu}(p) = \mathfrak{s}_{\mu\nu}(p) = 0$; and
\begin{equation}\label{10.4:eq-diagonal-stress-1}
\begin{split}
\mathfrak{s}_{\nu\nu}(p) ={}& \tfrac34|\vec w_p|^2
 + \tfrac18\bigl(|\vec w_{p+\hat\nu}|^2 + |\vec w_{p-\hat\nu}|^2\bigr)\\
 &- \tfrac18\bigl(|\vec w_{\ell_1,+\nu}+\vec w_{\ell_1,-\nu}|^2
 + |\vec w_{\ell_3,+\nu}+\vec w_{\ell_3,-\nu}|^2\bigr),
\end{split}
\end{equation}
where $|\vec w_{p'}|^2 = a_{p'}(2-a_{p'})$ for every plaquette $p'$.

Consequently, for $\varphi\colon \mathbb{R}^4 \to \mathbb{C}$ and each direction $\nu$, writing
$\sum_{p\ni\nu}$ for the sum over the plaquettes having $\nu$ as one of their two directions,
\begin{equation}\label{10.4:eq-diagonal-stress-2}
\sum_{p\ni\nu}\varphi(x_p)\,\mathfrak{s}_{\nu\nu}(p)
 = 2\sum_{p\ni\nu}\tilde\varphi_\nu(x_p)\,a_p - \mathfrak{Q}_\nu(\varphi),
\end{equation}
where $\tilde\varphi_\nu := \varphi + \tfrac18 a^2\Delta_\nu\varphi$, with
\begin{equation*}
\Delta_\nu\varphi(y) := a^{-2}\bigl[\varphi(y+\hat\nu)+\varphi(y-\hat\nu)-2\varphi(y)\bigr],
\end{equation*}
and
\begin{equation}\label{10.4:eq-diagonal-stress-3}
\begin{split}
\mathfrak{Q}_\nu(\varphi) :={}& \sum_{p\ni\nu}\tilde\varphi_\nu(x_p)\,a_p^2\\
 &+ \tfrac18\sum_{p\ni\nu}\varphi(x_p)
 \bigl(|\vec w_{\ell_1,+\nu}+\vec w_{\ell_1,-\nu}|^2
 + |\vec w_{\ell_3,+\nu}+\vec w_{\ell_3,-\nu}|^2\bigr),
\end{split}
\end{equation}
with $\ell_1,\ell_3$ attached to each $p$ as above. The density of $\mathfrak{Q}_\nu$ has no
quadratic part on constant fields ($a_p^2$ is quartic, and
$\vec w_{\ell,+\nu}+\vec w_{\ell,-\nu}$ is a covariant $\nu$-difference); classically it is
$O(a^6)$.
\end{lemma}

\begin{proof}
In the formula defining $\mathfrak{s}_{\sigma\rho}(p)$ in the standing notation, each link
$\ell\in\partial p$ enters through the $\sigma$-component of the displacement
$a^{-1}(x_\ell - x_p)$ and through the link-plaquette variation
$\mathfrak{d}_{\ell,p,\rho}$. For $p = \{x;\mu,\nu\}$ the displacements are
\begin{equation*}
a^{-1}(x_{\ell_1}-x_p) = -\tfrac12 e_\nu,\qquad a^{-1}(x_{\ell_3}-x_p) = +\tfrac12 e_\nu,
\end{equation*}
and $\mp\tfrac12 e_\mu$ on the two links of $\partial p$ of direction $\nu$, namely $(x,\nu)$
and $(x+\hat\mu,\nu)$. Since no displacement has a component outside the directions $\mu,\nu$,
$\mathfrak{s}_{\sigma\rho}(p)=0$ unless $\sigma\in\{\mu,\nu\}$. Moreover, the variation of a link
in its own direction vanishes: $\vec u_{\ell,\mu(\ell)} = 0$, where $\mu(\ell)$ denotes the
direction of $\ell$, and accordingly $\mathfrak{d}_{\ell,p,\rho} = 0$ whenever $\rho$ is the
direction of $\ell$. The component $\sigma = \nu$ of the displacement is carried only by
$\ell_1,\ell_3$, whose direction is $\mu$, which is the second index $\rho$ of
$\mathfrak{s}_{\nu\mu}$; so $\mathfrak{s}_{\nu\mu}(p) = 0$. The component
$\sigma = \mu$ is carried only by the two links of direction $\nu$, which is the second index of
$\mathfrak{s}_{\mu\nu}$; so $\mathfrak{s}_{\mu\nu}(p) = 0$. This proves the first two claims.

For the diagonal entry, only $\ell_1$ and $\ell_3$ contribute, with the coefficients $-\tfrac12$
and $+\tfrac12$ read off above:
\begin{equation}\label{10.4:eq-diagonal-stress-4}
\mathfrak{s}_{\nu\nu}(p) = -\tfrac12\,\mathfrak{d}_{\ell_1,p,\nu}
 + \tfrac12\,\mathfrak{d}_{\ell_3,p,\nu}.
\end{equation}
Now $p = p(\ell_1,+\nu) = p(\ell_3,-\nu)$, so $\vec w_{\ell_1,p} = \vec w_{\ell_1,+\nu}$ and
$\vec w_{\ell_3,p} = \vec w_{\ell_3,-\nu}$. The explicit form of the link-plaquette variation
in the Pauli parametrisation gives
\begin{align*}
\mathfrak{d}_{\ell_1,p,\nu} &= -\tfrac12|\vec w_p|^2
 + \tfrac12\,\vec w_{\ell_1,-\nu}\cdot\vec w_{\ell_1,+\nu},\\
\mathfrak{d}_{\ell_3,p,\nu} &= \tfrac12|\vec w_p|^2
 - \tfrac12\,\vec w_{\ell_3,+\nu}\cdot\vec w_{\ell_3,-\nu}.
\end{align*}
Inserting these into \eqref{10.4:eq-diagonal-stress-4} yields
\begin{equation*}
\mathfrak{s}_{\nu\nu}(p) = \tfrac12|\vec w_p|^2
 - \tfrac14\bigl[\vec w_{\ell_1,-\nu}\cdot\vec w_{\ell_1,+\nu}
 + \vec w_{\ell_3,+\nu}\cdot\vec w_{\ell_3,-\nu}\bigr].
\end{equation*}
To each bracketed product we apply the identity
$-u\cdot v = \tfrac12(|u|^2+|v|^2) - \tfrac12|u+v|^2$. The modulus $|\vec w_{\ell,p}|$ of a
plaquette vector does not depend on the link $\ell\in\partial p$ at which it is based. Hence
$|\vec w_{\ell_1,+\nu}| = |\vec w_{\ell_3,-\nu}| = |\vec w_p|$. Moreover,
$p(\ell_1,-\nu) = p-\hat\nu$ and $p(\ell_3,+\nu) = p+\hat\nu$, so
$|\vec w_{\ell_1,-\nu}| = |\vec w_{p-\hat\nu}|$ and $|\vec w_{\ell_3,+\nu}| = |\vec w_{p+\hat\nu}|$.
This gives
\begin{equation*}
\begin{split}
\mathfrak{s}_{\nu\nu}(p) ={}& \tfrac12|\vec w_p|^2
 + \tfrac18\bigl(|\vec w_{p-\hat\nu}|^2 + |\vec w_p|^2 + |\vec w_{p+\hat\nu}|^2
 + |\vec w_p|^2\bigr)\\
 &- \tfrac18\bigl(|\vec w_{\ell_1,+\nu}+\vec w_{\ell_1,-\nu}|^2
 + |\vec w_{\ell_3,+\nu}+\vec w_{\ell_3,-\nu}|^2\bigr).
\end{split}
\end{equation*}
Since $\tfrac12+\tfrac14 = \tfrac34$, this is \eqref{10.4:eq-diagonal-stress-1}.

The relation $|\vec w_{p'}|^2 = a_{p'}(2-a_{p'})$ comes from two facts. For an element of
$SU(2)$ one has $(w^0)^2 + |\vec w|^2 = 1$. The normalised trace of $W_{\ell,p'}$ is
$w^0 = 1 - a_{p'}$. Combining them,
\begin{equation*}
|\vec w_{p'}|^2 = 1-(1-a_{p'})^2 = a_{p'}(2-a_{p'}).
\end{equation*}

For the summed identity, multiply \eqref{10.4:eq-diagonal-stress-1} by $\varphi(x_p)$ and sum
over $p\ni\nu$. In the two terms containing $|\vec w_{p\pm\hat\nu}|^2$, shift the summation
variable $p \mapsto p\mp\hat\nu$, using $x_{p\pm\hat\nu} = x_p\pm\hat\nu$:
\begin{equation*}
\sum_{p\ni\nu}\varphi(x_p)\,|\vec w_{p\pm\hat\nu}|^2
 = \sum_{p\ni\nu}\varphi(x_p\mp\hat\nu)\,|\vec w_p|^2 .
\end{equation*}
The coefficient of $|\vec w_p|^2$ then becomes
\begin{equation*}
\tfrac34\varphi(x_p) + \tfrac18\bigl(\varphi(x_p+\hat\nu)+\varphi(x_p-\hat\nu)\bigr)
 = \varphi(x_p) + \tfrac18 a^2\Delta_\nu\varphi(x_p) = \tilde\varphi_\nu(x_p).
\end{equation*}
By $|\vec w_p|^2 = 2a_p - a_p^2$, we have
\begin{equation*}
\sum_{p\ni\nu}\tilde\varphi_\nu(x_p)\,|\vec w_p|^2
 = 2\sum_{p\ni\nu}\tilde\varphi_\nu(x_p)\,a_p - \sum_{p\ni\nu}\tilde\varphi_\nu(x_p)\,a_p^2 .
\end{equation*}
The remaining terms are $-\tfrac18\sum_{p\ni\nu}\varphi(x_p)(\cdots)$, with the bracket of
\eqref{10.4:eq-diagonal-stress-1}. Collecting the terms with $a_p^2$ and the bracket terms into
$\mathfrak{Q}_\nu(\varphi)$ as in \eqref{10.4:eq-diagonal-stress-3} gives
\eqref{10.4:eq-diagonal-stress-2}.

Finally, on a constant abelian field one has $\vec w_{\ell,-\nu} = -\vec w_{\ell,+\nu}$, so the
bracket terms in the density of $\mathfrak{Q}_\nu$ vanish there; the remaining part of that
density, $a_p^2$, is quartic.
\end{proof}

\begin{remark}
In the diagonal channel, Lemma~\ref{10.4:lem-diagonal-stress} shows that the lattice stress
tensor and the orientation-weighted Wilson density differ by a bare unit with no marginal part.
So the diagonal half of the stress-tensor-insertion correlation theorem is a statement about
\begin{equation*}
A(c,U) = \sum_p c(p)\,a_p(U),
\end{equation*}
for a plaquette weight $c$ depending on the orientation of $p$ (a general weight, not the
specific weight $c(p)=g_0^{-2}-w(p)$ of the standing notation). This is the functional of the
correlation theorem: linear in the weight $c$, gauge invariant plaquette by plaquette, and entire
in the link variables on the complexification of the structure group. It lies outside the weight
class of the correlation theorem
only because the Lipschitz condition in the definition of that class fails: the Lipschitz
seminorm of $c$ is not small across orientations.
\end{remark}

\begin{lemma}[Hypercubic decomposition of symmetric tensors]
\label{10.4:lem-hypercubic-decomposition}
Let $V=\mathbb{R}^4$ with standard basis $e_1,\dots,e_4$, on which the hyperoctahedral group $W_4$
acts by signed permutations of the basis vectors. Let $D=\{\pm1\}^4\subset W_4$ be the sign
subgroup, let $S_4\subset W_4$ be the subgroup of permutations, and write $\chi_{\mu\nu}(\varepsilon)
:=\varepsilon_\mu\varepsilon_\nu$ for $\varepsilon\in D$ and $\mu\neq\nu$. Let
$\Lambda^2=\Lambda^2V$, put $E_{\mu\nu}:=e_\mu\wedge e_\nu$ for $\mu\neq\nu$ (so that
$E_{\nu\mu}=-E_{\mu\nu}$; the $E_{\mu\nu}$ with $\mu<\nu$ form a basis of $\Lambda^2$), and let
$\odot$ denote the
symmetric product; $\operatorname{Sym}^2V$ is identified with the real symmetric $4\times4$
matrices, $\delta$ is the identity matrix, and $\operatorname{diag}(h_1,\dots,h_4)$ is the diagonal
matrix with entries $h_\mu$. Put
\begin{equation*}
\mathsf{V}_{\mathbf{1}}:=\mathbb{R}\delta,\qquad
\mathsf{V}_{\mathbf{3}}:=\Bigl\{\operatorname{diag}(h_1,\dots,h_4):\ \sum_\mu h_\mu=0\Bigr\},\qquad
\mathsf{V}_{\mathbf{6}}:=\operatorname{span}\{e_\sigma\odot e_\rho:\ \sigma<\rho\}.
\end{equation*}
Then:
\begin{enumerate}
\item[(a)] $\operatorname{Sym}^2V=\mathsf{V}_{\mathbf{1}}\oplus\mathsf{V}_{\mathbf{3}}\oplus
\mathsf{V}_{\mathbf{6}}$ as $W_4$-modules; the three summands are irreducible and pairwise
inequivalent.
\item[(b)] $\operatorname{Hom}_{W_4}(\mathsf{V}_r,\mathbb{R})=0$ for $r=\mathbf{3},\mathbf{6}$.
\item[(c)] $\dim\operatorname{Hom}_{W_4}(\mathsf{V}_r,\operatorname{Sym}^2\Lambda^2)=1$ for
$r=\mathbf{1},\mathbf{3},\mathbf{6}$, with generators
\begin{equation}
\label{10.4:eq-hypercubic-decomposition-1}
\begin{aligned}
\mathsf{T}^{\mathbf{1}}(\delta)&:=\sum_{\mu<\nu}E_{\mu\nu}^{\odot2},\qquad
\mathsf{T}^{\mathbf{3}}(h):=\sum_{\mu<\nu}(h_\mu+h_\nu)\,E_{\mu\nu}^{\odot2},\\
\mathsf{T}^{\mathbf{6}}(e_\sigma\odot e_\rho)&:=\sum_{\kappa\notin\{\sigma,\rho\}}
E_{\sigma\kappa}\odot E_{\rho\kappa},
\end{aligned}
\end{equation}
where $h$ stands for $\operatorname{diag}(h_1,\dots,h_4)\in\mathsf{V}_{\mathbf{3}}$.
\item[(d)] $\operatorname{Hom}_{W_4}(V\otimes\mathsf{V}_r,\mathbb{R}\oplus\operatorname{Sym}^2
\Lambda^2)=0$ for $r=\mathbf{1},\mathbf{3},\mathbf{6}$.
\end{enumerate}
\end{lemma}

\begin{proof}
An element $\varepsilon\in D$ acts by $\varepsilon\,e_\mu=\varepsilon_\mu e_\mu$, hence by
$\varepsilon\,(e_\sigma\odot e_\rho)=\chi_{\sigma\rho}(\varepsilon)\,e_\sigma\odot e_\rho$ for
$\sigma\neq\rho$ and trivially on $e_\mu\odot e_\mu$; a permutation $\pi\in S_4$ acts by
$e_\mu\mapsto e_{\pi\mu}$. In matrix language $W_4$ acts on $\operatorname{Sym}^2V$ by conjugation
with signed permutation matrices, which preserves the diagonal matrices, the trace and the
off-diagonal matrices; hence the three subspaces are $W_4$-invariant, and their sum is
$\operatorname{Sym}^2V$ because a symmetric matrix is the sum of a multiple of $\delta$, a traceless
diagonal matrix and an off-diagonal symmetric matrix, uniquely.

(a) Since $\varepsilon_\mu^2=1$, the group $D$ acts trivially on the diagonal matrices, on which
$S_4$ acts by permuting four points. This permutation module is $\mathbf{1}\oplus$ (the standard
representation of $S_4$), the latter irreducible; the trivial summand is
$\mathsf{V}_{\mathbf{1}}$ and the standard one is the sum-zero subspace $\mathsf{V}_{\mathbf{3}}$.
The six lines $\mathbb{R}\,e_\sigma\odot e_\rho$, $\sigma<\rho$, carry the characters
$\chi_{\sigma\rho}$ of $D$, which are non-trivial and pairwise distinct (evaluate at the elements of
$D$ with exactly one entry $-1$). Hence a $D$-invariant subspace of $\mathsf{V}_{\mathbf{6}}$ is the
sum of its $\chi_{\sigma\rho}$-isotypic parts, that is, it is spanned by some of the six lines; and
$S_4$ permutes the six lines transitively. So a non-zero $W_4$-invariant subspace of
$\mathsf{V}_{\mathbf{6}}$ is all of $\mathsf{V}_{\mathbf{6}}$. The dimensions are $1$, $3$, $6$, so
the three irreducible modules are pairwise inequivalent.

(b) $\mathsf{V}_{\mathbf{3}}$ and $\mathsf{V}_{\mathbf{6}}$ are irreducible and non-trivial
($S_4$ acts non-trivially on $\mathsf{V}_{\mathbf{3}}$, $D$ acts non-trivially on
$\mathsf{V}_{\mathbf{6}}$). The kernel of an equivariant map into the trivial module $\mathbb{R}$
is an invariant subspace of codimension at most one, hence the whole module, and the map is zero.

(c) The vector $E_{\mu\nu}\odot E_{\kappa\lambda}$ transforms under $D$ with the character
$\chi_{\mu\nu}\chi_{\kappa\lambda}$, which is $\varepsilon\mapsto\prod_{\alpha}\varepsilon_\alpha$
over the symmetric difference of $\{\mu,\nu\}$ and $\{\kappa,\lambda\}$; it is trivial if and only
if $\{\mu,\nu\}=\{\kappa,\lambda\}$. Therefore
\begin{equation*}
(\operatorname{Sym}^2\Lambda^2)^D=\operatorname{span}\{E_{\mu\nu}^{\odot2}:\ \mu<\nu\},
\end{equation*}
and $S_4$ acts on it as the permutation module on the six pairs. Since $D$ acts trivially on the
diagonal matrices, every $W_4$-equivariant map from the diagonal matrices into
$\operatorname{Sym}^2\Lambda^2$ takes values in $(\operatorname{Sym}^2\Lambda^2)^D$, and is an
$S_4$-equivariant map from the permutation module on four points to the permutation module on six
pairs. Such a map is given by a function on the $S_4$-orbits of pairs (point, pair), and there are
two orbits: the point lies in the pair, or it does not. Hence
\begin{equation*}
\dim\operatorname{Hom}_{W_4}(\mathsf{V}_{\mathbf{1}}\oplus\mathsf{V}_{\mathbf{3}},
\operatorname{Sym}^2\Lambda^2)=2 .
\end{equation*}
One dimension is $\operatorname{Hom}_{W_4}(\mathsf{V}_{\mathbf{1}},\operatorname{Sym}^2\Lambda^2)
\cong(\operatorname{Sym}^2\Lambda^2)^{W_4}$, which is a line by the classification of
$W_4$-invariant quadratic forms in the curvature; directly, $(\operatorname{Sym}^2\Lambda^2)^{W_4}$
is the space of $S_4$-fixed vectors of the permutation module on the six pairs, hence a line,
$S_4$ being transitive on the pairs. It is generated by
$\mathsf{T}^{\mathbf{1}}(\delta)$ in \eqref{10.4:eq-hypercubic-decomposition-1}. The other
dimension is $\operatorname{Hom}_{W_4}(\mathsf{V}_{\mathbf{3}},\operatorname{Sym}^2\Lambda^2)$. The
map $\mathsf{T}^{\mathbf{3}}$ is the restriction to $\mathsf{V}_{\mathbf{3}}$ of the equivariant map
belonging to the orbit ``the point lies in the pair'', and $\mathsf{T}^{\mathbf{3}}\neq0$: if
$h_\mu+h_\nu=0$ for all pairs, then for three distinct indices $\mu,\nu,\kappa$ the relations
$h_\mu+h_\nu=h_\mu+h_\kappa=h_\nu+h_\kappa=0$ give $h_\mu=0$, so $h=0$.

For $\mathsf{V}_{\mathbf{6}}$ fix $\sigma<\rho$. By the computation of characters above,
$\chi_{\mu\nu}\chi_{\kappa\lambda}=\chi_{\sigma\rho}$ if and only if the two pairs share exactly one
index $\kappa$ and their other indices are $\sigma$ and $\rho$; so the $\chi_{\sigma\rho}$-isotypic
part of $\operatorname{Sym}^2\Lambda^2$ is
\begin{equation*}
\operatorname{span}\{E_{\sigma\kappa}\odot E_{\rho\kappa}:\ \kappa\notin\{\sigma,\rho\}\},
\end{equation*}
in accordance with the description of the $D$-isotypic parts of $\operatorname{Sym}^2\Lambda^2$
in the lemma on hyperoctahedral invariants of the curvature. These are two vectors, exchanged by
the transposition of the two
values of $\kappa$, which fixes $e_\sigma\odot e_\rho$. A $W_4$-equivariant map sends
$e_\sigma\odot e_\rho$ into this isotypic part, and its image is fixed by that transposition, so it
is a multiple of the sum of the two vectors. The multiple does not depend on $(\sigma,\rho)$, by
transitivity of $S_4$ on the pairs, since $\pi\in S_4$ maps $E_{\sigma\kappa}\odot E_{\rho\kappa}$
to $E_{\pi\sigma\,\pi\kappa}\odot E_{\pi\rho\,\pi\kappa}$. Hence the dimension is at most one. The
sum is fixed by the stabiliser of the line $\mathbb{R}\,e_\sigma\odot e_\rho$ in the way
$e_\sigma\odot e_\rho$ is: the transposition $(\sigma\rho)$ fixes each term, $\odot$ being
symmetric; a sign change of $\kappa$ changes both factors; and $D$ acts on both by
$\chi_{\sigma\rho}$. Hence $\mathsf{T}^{\mathbf{6}}$ in \eqref{10.4:eq-hypercubic-decomposition-1}
is $W_4$-equivariant: for $\pi\in S_4$ the formula gives
$\mathsf{T}^{\mathbf{6}}(e_{\pi\sigma}\odot e_{\pi\rho})=\pi\,\mathsf{T}^{\mathbf{6}}
(e_\sigma\odot e_\rho)$, and $\varepsilon\in D$ multiplies both sides by
$\varepsilon_\sigma\varepsilon_\rho$. It is non-zero, the vectors $E_{\sigma\kappa}\odot
E_{\rho\kappa}$ being linearly independent.

(d) The element $-\mathrm{id}_V\in W_4$ acts by $-1$ on $V$ and by $+1$ on $\operatorname{Sym}^2V$,
hence by $-1$ on $V\otimes\mathsf{V}_r$; it acts by $+1$ on $\mathbb{R}$ and on
$\operatorname{Sym}^2\Lambda^2$. An equivariant map $\Phi$ therefore satisfies
$\Phi(x)=-\Phi(x)$ for every $x$, and $\Phi=0$.
\end{proof}

\noindent\textit{Quadratic forms.} For antisymmetric families $F=(F_{\mu\nu})$ with values in the
Lie algebra of $SU(N)$ write $\vec F_{\mu\nu}\cdot\vec G_{\mu\nu}:=-\operatorname{tr}F_{\mu\nu}G_{\mu\nu}$
(the colour inner product) and $|\vec F_{\mu\nu}|^2:=\vec F_{\mu\nu}\cdot\vec F_{\mu\nu}$. The three
quadratic curvature forms are
\begin{equation}
\label{10.4:eq-hypercubic-decomposition-2}
\begin{aligned}
\mathsf{t}^{\mathbf{1}}(F)&:=\tfrac12\sum_{\mu<\nu}|\vec F_{\mu\nu}|^2,\qquad
\mathsf{t}^{\mathbf{3}}(h;F):=\tfrac12\sum_{\mu<\nu}(h_\mu+h_\nu)\,|\vec F_{\mu\nu}|^2,\\
\mathsf{t}^{\mathbf{6}}(h_0;F)&:=\sum_{\sigma\neq\rho}(h_0)_{\sigma\rho}\sum_\kappa
\vec F_{\rho\kappa}\cdot\vec F_{\sigma\kappa},
\end{aligned}
\end{equation}
for $h\in\mathsf{V}_{\mathbf{3}}$ and $h_0\in\mathsf{V}_{\mathbf{6}}$, the latter written as an
off-diagonal symmetric matrix $((h_0)_{\sigma\rho})$. In terms of the three marginal tensor vertices
$\mathcal{T}^{(1)},\mathcal{T}^{(2)},\mathcal{T}^{(3)}$, in the notation of the classification of
$W_4$-invariant quadratic forms in the curvature, the channel $\mathbf{1}$ corresponds to
$\mathcal{T}^{(2)}$, the channel $\mathbf{3}$ to the traceless part of $\mathcal{T}^{(3)}$, and the
channel $\mathbf{6}$ to the off-diagonal part of $\mathcal{T}^{(1)}$; the statement of that
classification that the space of $W_4$-equivariant maps from $\operatorname{Sym}^2V$ into
$\operatorname{Sym}^2\Lambda^2$ is three-dimensional is, by (a) and (c), the decomposition
$3=1+1+1$ into the three lines of (c). By Schur's lemma and (a), a $W_4$-equivariant linear map
on the three marginal slots is diagonal in $r$: the summands $\mathsf{V}_{\mathbf{1}},
\mathsf{V}_{\mathbf{3}},\mathsf{V}_{\mathbf{6}}$ are irreducible and pairwise inequivalent, so no
non-zero equivariant map connects two different ones. In particular a $W_4$-equivariant
normalisation of the three slots, a priori a $3\times3$ matrix, is given by three numbers.

\begin{remark}
All three pieces lie in $\mathbf{1}\oplus\mathbf{9}$ under $O(4)$: $\mathsf{V}_{\mathbf{1}}$ is the
trace $\mathbf{1}$, and $\mathsf{V}_{\mathbf{3}}\oplus\mathsf{V}_{\mathbf{6}}$ is the space
$\mathbf{9}$ of traceless symmetric tensors. In the forms
\eqref{10.4:eq-hypercubic-decomposition-2} this is visible from the identity
\begin{equation*}
\sum_{\mu<\nu}(h_\mu+h_\nu)\,|\vec F_{\mu\nu}|^2=\sum_\mu h_\mu\sum_\kappa|\vec F_{\mu\kappa}|^2,
\end{equation*}
so that $\mathsf{t}^{\mathbf{3}}$ and $\mathsf{t}^{\mathbf{6}}$ are contractions of a symmetric
tensor with $\sum_\kappa\vec F_{\mu\kappa}\cdot\vec F_{\nu\kappa}$. Consequently the Weyl-type
channel $\mathbf{10}$ of $O(4)$ is not met at first order in a tensor source transforming in
$\operatorname{Sym}^2V$.
\end{remark}

\begin{definition}[Typed families and their marginal coefficient]\label{10.4:def-typed-families}
Let $\mathsf{V}$ be a $W_4$-module.
\begin{enumerate}
\item A \emph{$\mathsf{V}$-typed pointed $E$-family at scale $j$} is a family $F$ that satisfies
conditions (f1) and (f2) of the definition of pointed $E$-families at scale $j$ together with
the $\mathsf{V}$-typed form (f3)$^{\mathsf{V}}$ of condition (f3), namely
\begin{equation}\label{10.4:eq-typed-families-1}
F(\mathsf{o}U,\mathsf{o}y;\mathsf{o}_*h_0) = F(U,y;h_0)\qquad\text{for all } \mathsf{o}\in W_4 .
\end{equation}
Here $\mathsf{o}_*$ is the action of $\mathsf{o}$ on $\mathsf{V}$. The condition (S) of that
definition is not imposed. The norm of such a family is
\begin{equation}\label{10.4:eq-typed-families-2}
\|F\| := \sup_{|h_0|\le 1}\ \sup\ e^{\kappa d_j(X)}\,\bigl|F(X,\cdot,y;h_0)\bigr| .
\end{equation}
In this formula the inner supremum is over the arguments of $F$ other than $h_0$.
\emph{$\mathsf{V}$-typed marginal families} are defined in the same way, starting from marginal
families.

\item \emph{Counterterm families.} Consider a lattice of spacing $a$ carrying the configuration
$U$. Let $\varphi$ be a scalar weight, and let $h_0\in(\mathsf{V}_r)_{\mathbb{C}}$ for
$r=\mathbf{3},\mathbf{6}$. For $r=\mathbf{3}$ write $h=(h_\mu)$ for the diagonal entries of
$h_0$, and put
\begin{equation}\label{10.4:eq-typed-families-3}
\mathsf{A}^{\mathbf{3}}(\varphi,U;h) := A(\varphi\,\mathsf{c}_h,U),\qquad
\mathsf{c}_h(p) := h_\mu+h_\nu\quad\text{for } p\parallel(\mu\nu).
\end{equation}
Here $p\parallel(\mu\nu)$ means that the plaquette $p$ lies in a $(\mu\nu)$-coordinate plane; for
a plaquette weight $c$ we write
\begin{equation*}
A(c,U) := \sum_p c(p)\,\bigl[1-\operatorname{Re}\operatorname{tr}U(\partial p)\bigr],
\end{equation*}
and $\varphi\,\mathsf{c}_h$ is the plaquette weight $p\mapsto\varphi(x(p))\,\mathsf{c}_h(p)$.
For $r=\mathbf{6}$ put
\begin{equation}\label{10.4:eq-typed-families-4}
\mathsf{A}^{\mathbf{6}}(\varphi,U;h_0)
:= \sum_p \varphi(x(p))\sum_{\sigma\neq\rho}(h_0)_{\sigma\rho}\,\mathfrak{s}_{\sigma\rho}(p).
\end{equation}
Here the lattice stress tensor $\mathfrak{s}_{\sigma\rho}(p)$ is formed with the links of that
lattice.

Let $\mathsf{k}$ be a tensor weight. Then $\mathsf{A}^{r}(\mathsf{k},U)$ is defined by the same
formulas, with $\varphi(x(p))h_0$ replaced by $\mathsf{k}(x(p))$; for $r=\mathbf{3}$ we write
$\mathsf{k}(x)_\mu$ for the diagonal entry $\mathsf{k}(x)_{\mu\mu}$:
\begin{align}
\mathsf{A}^{\mathbf{3}}(\mathsf{k},U)
&:= \sum_p \bigl(\mathsf{k}(x(p))_\mu+\mathsf{k}(x(p))_\nu\bigr)
\bigl[1-\operatorname{Re}\operatorname{tr}U(\partial p)\bigr]
\quad(p\parallel(\mu\nu)),\label{10.4:eq-typed-families-5}\\
\mathsf{A}^{\mathbf{6}}(\mathsf{k},U)
&:= \sum_p\sum_{\sigma\neq\rho}\mathsf{k}(x(p))_{\sigma\rho}\,
\mathfrak{s}_{\sigma\rho}(p).\label{10.4:eq-typed-families-6}
\end{align}
These counterterm families are polynomials in the links. They are therefore entire in each link
variable on the complexified group. For $\mathsf{A}^{\mathbf{6}}$ the lattice stress tensor is
written in the parametrisation $W=w^0+i\vec w\cdot\vec\sigma$ of $SU(2)$ (see
Lemma~\ref{10.4:lem-diagonal-stress}), so that $N=2$ there; the complexified group is then
$SL(2,\mathbb{C})\supset SU(2)$, with $\vec w$ read as the Pauli coefficients of
$\tfrac12(W-W^{-1})$. They are gauge invariant, linear
in the weight, and $\mathsf{V}_r$-typed.

\item \emph{The marginal coefficient.} Let $r\in\{\mathbf{1},\mathbf{3},\mathbf{6}\}$ and let
$F$ be a $\mathsf{V}_r$-typed family. Let $\Pi(h_0)$ be the tensor of second moments of the
Hessian, at the trivial configuration, of the total functional considered in the proof of one of
the lemmas of the correlation theorem, where this tensor is constructed; it depends linearly on
$F$ and on the marginal slot $h_0$. Its part
symmetric under exchange of the two pairs of indices
equals $\beta^{r}[F]\,\mathsf{T}^{r}(h_0)$ for a unique number $\beta^{r}[F]\in\mathbb{C}$.
This number is the \emph{marginal coefficient} of $F$, and it is linear in $F$. For
$r\in\{\mathbf{3},\mathbf{6}\}$, the values of $r$ for which $\mathsf{A}^{r}$ is defined, the
\emph{typed renormalised family} is
\begin{equation}\label{10.4:eq-typed-families-7}
F^{\mathrm{ren},r} := F - F(1) - \beta^{r}[F]\,\mathsf{A}^{r},
\end{equation}
where $F(1)$ denotes the family evaluated at the trivial configuration $U\equiv 1$.
For frozen typed families with $r=\mathbf{3},\mathbf{6}$ one has $F(1)=0$; the properties of
frozen families used for this are that their value at $U\equiv1$ does not depend on the point
$y$ and is linear in $h_0$.
\end{enumerate}
\end{definition}

\begin{proof}
Linearity of the counterterm families in the weight can be read off the formulas. The right-hand
sides of \eqref{10.4:eq-typed-families-3} and \eqref{10.4:eq-typed-families-5} are linear in
$\varphi\,\mathsf{c}_h$ and in $\mathsf{k}$ respectively. The right-hand sides of
\eqref{10.4:eq-typed-families-4} and \eqref{10.4:eq-typed-families-6} are linear in the
products $\varphi(x(p))(h_0)_{\sigma\rho}$ and in $\mathsf{k}(x(p))_{\sigma\rho}$. For
$\mathsf{A}^{\mathbf{3}}$, polynomiality and gauge invariance can also be read off
\eqref{10.4:eq-typed-families-3} and \eqref{10.4:eq-typed-families-5}: each summand is a weight
times $1-\operatorname{Re}\operatorname{tr}U(\partial p)$, where $U(\partial p)$ is the ordered
product of the link variables of $\partial p$ and their inverses; for $V\in SU(N)$, for every
$N$, one has
\begin{equation*}
\operatorname{Re}\operatorname{tr}V=\tfrac12\bigl(\operatorname{tr}V+\operatorname{tr}V^{-1}\bigr),
\end{equation*}
and $V^{-1}$ is the adjugate of $V$, a polynomial in its entries; a gauge transformation replaces
$U(\partial p)$ by a conjugate, which has the same trace. For $\mathsf{A}^{\mathbf{6}}$, where
$N=2$, the lattice stress tensor $\mathfrak{s}_{\sigma\rho}(p)$ is a polynomial in inner products
of the vectors $\vec w$ of plaquette variables based at common links, as
Lemma~\ref{10.4:lem-diagonal-stress} displays for the diagonal entries and as the definition of
the lattice stress tensor gives for all entries. Each such $\vec w$ is the vector of Pauli
coefficients of $\tfrac12(W-W^{-1})$ for a plaquette variable $W$, which is a product of link
variables and their inverses, hence a polynomial in the links. A gauge transformation conjugates
all plaquette variables based at a link by one and the same element $g\in SU(2)$; since
\begin{equation*}
g\bigl(w^0+i\vec w\cdot\vec\sigma\bigr)g^{-1}=w^0+i(R_g\vec w)\cdot\vec\sigma
\end{equation*}
with $R_g\in SO(3)$ the adjoint rotation, it rotates these vectors simultaneously by $R_g$ and
leaves their inner products unchanged. The
$\mathsf{V}_r$-typedness of $\mathsf{A}^{r}$ is the hyperoctahedral covariance of the plaquette
action and of the lattice stress tensor, that is, the identities for $W_4$ acting on
$\mathfrak{s}_{\sigma\rho}$ and on the plaquette orientations that accompany the definition of
the lattice stress tensor.

Next we show that the marginal coefficient is well defined. By item (3), $\Pi(h_0)$ is
complex-linear in $h_0\in(\mathsf{V}_r)_{\mathbb{C}}$. The trivial configuration
$U\equiv1$ is fixed by $W_4$, and the second moments defining $\Pi(h_0)$ are taken relative to
the base point of the family, which the covariance \eqref{10.4:eq-typed-families-1} moves
together with the configuration. Hence
the covariance \eqref{10.4:eq-typed-families-1} of $F$, applied at and around $U\equiv1$, gives
$\Pi(\mathsf{o}_*h_0)=\mathsf{o}\cdot\Pi(h_0)$ for $\mathsf{o}\in W_4$, with $\mathsf{o}$ acting on
the tensor through its action on the index pairs; the exchange of the two pairs commutes with
this action. So the symmetrised part of $\Pi(h_0)$ is a $W_4$-equivariant complex-linear map
$(\mathsf{V}_r)_{\mathbb{C}}\to(\operatorname{Sym}^2\Lambda^2)_{\mathbb{C}}$. By
Lemma~\ref{10.4:lem-hypercubic-decomposition}(c), the space
$\operatorname{Hom}_{W_4}(\mathsf{V}_r,\operatorname{Sym}^2\Lambda^2)$ is one-dimensional and
spanned by $\mathsf{T}^{r}$. Complexifying, the space of complex-linear $W_4$-equivariant maps
$(\mathsf{V}_r)_{\mathbb{C}}\to(\operatorname{Sym}^2\Lambda^2)_{\mathbb{C}}$ is the
complexification of that real space. It is therefore a complex line spanned by the complex-linear
extension of $\mathsf{T}^{r}$. Hence the symmetrised part of $\Pi(h_0)$ equals
$\beta^{r}[F]\,\mathsf{T}^{r}(h_0)$ for exactly one $\beta^{r}[F]\in\mathbb{C}$, since
$\mathsf{T}^{r}\neq0$. Linearity of $F\mapsto\beta^{r}[F]$ follows from uniqueness and from the
linearity of $\Pi$ in $F$.

Finally, let $F$ be a frozen typed family with $r=\mathbf{3},\mathbf{6}$. Its value at
$U\equiv1$ does not depend on the point $y$ and is linear in $h_0$ (item (3)); write it
$F(1;h_0)$. Since $U\equiv1$ is fixed by $W_4$, the covariance
\eqref{10.4:eq-typed-families-1} taken at $U\equiv1$ gives
$F(1;\mathsf{o}_*h_0)=F(1;h_0)$ for all $\mathsf{o}\in W_4$. Thus $h_0\mapsto F(1;h_0)$ is a
$W_4$-invariant complex-linear functional on $(\mathsf{V}_r)_{\mathbb{C}}$, that is, an element of
$\operatorname{Hom}_{W_4}(\mathsf{V}_r,\mathbb{R})\otimes\mathbb{C}$. By
Lemma~\ref{10.4:lem-hypercubic-decomposition}(b), this space is zero for these two values of
$r$, so $F(1)=0$.
\end{proof}

\begin{lemma}[Bounds on the typed marginal coefficient]\label{10.4:lem-coefficient-bounds}
Let $F$ be a $\mathsf{V}_r$-typed family at scale $j$ in the sense of
Definition~\ref{10.4:def-typed-families}, and let $\beta^{r}[F]\in\mathbb{C}$ be its typed marginal
coefficient, defined there. Assume the following.
\begin{enumerate}
\item[(I)] \emph{The second-moment lemmas of the correlation theorem, in the part preceding their
appeal to hypercubic symmetry.} Let $y$ be the base point. For every family $F$ at scale $j$, these
arguments bound and compare the whole tensor
\begin{equation}\label{10.4:eq-coefficient-bounds-1}
\sum_{x,x'} D^2F[\mathsf{h}_x,\mathsf{h}_{x'}]\,(x_\kappa-y_\kappa)(x'_\lambda-y_\lambda)
\end{equation}
and they use no symmetry of it. The comparison states that the tensor is the same for $F$ and $F'$
whenever $F\simeq F'$, in the sense in which the second-moment lemmas of the correlation theorem use
this relation, and that it is formed from the non-wrapping terms of $F$ only. The bounds
are $C_\beta\|F\|_j$ on $\mathcal{U}_j$ and
$C_\beta\,\mathfrak{t}_j^{-2}\,\|F\|^\flat_j$ on $\mathcal{U}^\flat_j$.
\item[(II)] \emph{Linear potentials.} A vector potential that is linear in $x$ is annihilated by
$\partial^*\partial$, where $\partial$ is the lattice exterior difference on potentials and
$\partial^*$ its adjoint. Since
$L$ is odd, it reproduces itself under the symmetric averages. It is therefore the minimiser of the
variational problem of \cite[(4.40)--(4.41), (5.37)--(5.38)]{B-CMP109}, with zero multiplier.
\item[(III)] \emph{Level independence.} Consider a typed family $F$ at scale $j$ whose typed kernel
$\hat\Pi_j$ and counterterm kernel $\hat\Delta^{r}_j$, the Fourier transforms in the
level-independence display of the correlation theorem, satisfy
\begin{equation*}
\hat\Pi_j(\mathrm{p})-\beta^{r}[F]\,\hat\Delta^{r}_j(\mathrm{p})=O(|\mathrm{p}|^3)
\end{equation*}
as the lattice momentum $\mathrm{p}$ tends to $0$.
For such a family, the level-independence display of the correlation theorem gives
$\beta^{r}[F\circ U_k]=\beta^{r}[F]$ for every $k\ge j$.
\end{enumerate}
Then the following hold.
\begin{enumerate}
\item[(i)] If $F\simeq F'$, then $\beta^{r}[F]=\beta^{r}[F']$. Moreover, $\beta^{r}$ reads the
non-wrapping terms of $F$ only.
\item[(ii)] On $\mathcal{U}_j$ one has $|\beta^{r}[F]|\le C_\beta\|F\|_j$. On
$\mathcal{U}^\flat_j$ one has
\begin{equation}\label{10.4:eq-coefficient-bounds-2}
|\beta^{r}[F]|\le C_\Pi\,\mathfrak{g}_j^{-1}\,\|F\|^\flat_j,
\qquad
C_\Pi:=\frac{C_\beta\,\mathfrak{r}_\circ^{2}}{(c_T C_*)^{2}} .
\end{equation}
\item[(iii)] For $r,r'\in\{\mathbf{3},\mathbf{6}\}$ and every $k$,
\begin{equation}\label{10.4:eq-coefficient-bounds-3}
\beta^{r}[\mathsf{A}^{r'}\circ U_k]=\delta_{rr'} .
\end{equation}
Moreover, $\beta^{r}[F\circ U_k]=\beta^{r}[F]$ for a typed family $F$ at scale $j\le k$.
\end{enumerate}
\end{lemma}

\begin{proof}
\emph{Parts (i) and (ii).} By (I), the comparison and the bounds of the second-moment lemmas are
statements about the whole tensor \eqref{10.4:eq-coefficient-bounds-1}. No symmetry of that tensor
is read in proving them. The coefficient $\beta^{r}[F]$ is read from the same tensor. It is the
coefficient of its exchange-symmetric part along the generator $\mathsf{T}^{r}(h_0)$ of
Lemma~\ref{10.4:lem-hypercubic-decomposition}(c). The comparison therefore gives
$\beta^{r}[F]=\beta^{r}[F']$ for $F\simeq F'$, and it shows that only non-wrapping terms enter.
This is (i).

The bound of (I) on $\mathcal{U}_j$ gives $|\beta^{r}[F]|\le C_\beta\|F\|_j$. On
$\mathcal{U}^\flat_j$, (I) gives the bound $C_\beta\,\mathfrak{t}_j^{-2}\|F\|^\flat_j$. It remains
to compute $\mathfrak{t}_j^{-2}$. We have $\mathfrak{t}_j^{-2}=\mathfrak{r}_\circ^2/\mathfrak{r}_j^2$
and $\mathfrak{r}_j=c_T\,g_j\,C_*\,(\log x_j)^{q_\flat}$. Since
$\mathfrak{g}_j=g_j^2(\log x_j)^{2q_\flat}$, it follows that
\begin{equation*}
\mathfrak{t}_j^{-2}
=\frac{\mathfrak{r}_\circ^2}{c_T^2C_*^2\,g_j^2(\log x_j)^{2q_\flat}}
=\mathfrak{r}_\circ^2\,(c_TC_*)^{-2}\,\mathfrak{g}_j^{-1}.
\end{equation*}
Hence $C_\beta\,\mathfrak{t}_j^{-2}=C_\Pi\,\mathfrak{g}_j^{-1}$, which is
\eqref{10.4:eq-coefficient-bounds-2}. The correlation theorem bounds this quantity by
$\mathfrak{t}_j^{-2}\le C_t\,x_j$, with the constant $C_t$ of the correlation theorem, which
discards the logarithm. Here the logarithm is kept, and the
identity is exact.

\emph{Part (iii), the case $r\neq r'$.} The family $\mathsf{A}^{r'}\circ U_k$ is
$\mathsf{V}_{r'}$-typed. The exchange-symmetric part of its second-moment tensor is a
$W_4$-equivariant linear function of $h_0\in\mathsf{V}_{r'}$. By
Lemma~\ref{10.4:lem-hypercubic-decomposition}(a), $\mathsf{V}_r$ and $\mathsf{V}_{r'}$ are
irreducible and inequivalent. Schur's lemma therefore gives that this part has no component along
$\mathsf{T}^{r}$, that is, $\beta^{r}[\mathsf{A}^{r'}\circ U_k]=0$.

\emph{Part (iii), the case $r=r'$.} Take the linear potential whose curvature is a constant
curvature $F$. In this paragraph $F$ denotes that curvature and not a family. By (II) this linear
potential is the minimiser, with zero multiplier, at every level $k$. Hence the background $U_k$ has
constant curvature $F$ to first order.

Let $h_y$ be the weight attached to $y$ in the correlation theorem, with $\sum_y h_y=1$. On each
plaquette $p$ the plaquette term $a_p:=1-\operatorname{Re}\operatorname{tr}U_k(\partial p)$ of the
action and the lattice stress tensor have the small-curvature expansions established in the
small-curvature expansion lemma for the plaquette action and the lattice stress tensor:
\begin{equation*}
a_p=\tfrac12\,\eta^4\,|\vec F_{\mu\nu}|^2+\dots,
\qquad
\mathfrak{s}_{\sigma\rho}(p)=\eta^4\,\vec F_{\rho\mu}\cdot\vec F_{\sigma\mu}+\dots .
\end{equation*}
Inserting them into the definition of the counterterm families shows that, per orientation, the
quadratic part of $\mathsf{A}^{r}(h_y,U_k;h_0)$ is
\begin{equation*}
\mathsf{t}^{r}(h_0;F)\,\sum_p h_y(x_p)\,\eta^4 .
\end{equation*}
By $\sum_y h_y=1$ this equals $\mathsf{t}^{r}(h_0;F)$. This is the quadratic curvature form of the
generator $\mathsf{T}^{r}(h_0)$. Hence the coefficient of $\mathsf{A}^{r}\circ U_k$ along
$\mathsf{T}^{r}$ is $1$, for every $k$, which is \eqref{10.4:eq-coefficient-bounds-3}.

\emph{Part (iii), the second claim.} This is (III) once its hypothesis
$\hat\Pi_j-\beta^{r}[F]\,\hat\Delta^{r}_j=O(|\mathrm{p}|^3)$ is checked. The typed kernel
$\hat\Pi_j$ is even in $\mathrm{p}$, by the element $-\mathbf{1}$ of $W_4$. It vanishes at
$\mathrm{p}=0$ and is transverse, by the
Ward--Takahashi identities. Its part of second order in $\mathrm{p}$ is the tensor
\eqref{10.4:eq-coefficient-bounds-1}, whose exchange-symmetric part is
$\beta^{r}[F]\,\mathsf{T}^{r}(h_0)$. Hence
\begin{equation*}
\hat\Pi_j-\beta^{r}[F]\,\hat\Delta^{r}_j=O(|\mathrm{p}|^3),
\end{equation*}
and (III) gives $\beta^{r}[F\circ U_k]=\beta^{r}[F]$ for $j\le k$.
\end{proof}

\begin{lemma}[Irrelevance after one typed counterterm]\label{10.4:lem-typed-irrelevance}
Let $r\in\{\mathbf{3},\mathbf{6}\}$. Let $F$ be a $\mathsf{V}_{r}$-typed pointed E-family at scale $j$ in the
sense of Definition~\ref{10.4:def-typed-families}, on radii $a_j$ (the radii at scale $j$ on which the
family is given), with $\|F\|\le\mathfrak{n}$.
Let $\mathfrak{b}$ be a background admissible at scale $n\ge j$ near the point $y$. Write
$t=L^{j-n}$ and $\nu=1+n-j$. Assume the following.
\begin{enumerate}
\item[(a)] The steps of the proof of the irrelevance lemma of the correlation theorem that precede its
appeal to reflections and permutations of the coordinate axes, namely the bounds at the near levels,
the bounds in the tail and the core of that proof, hold for the pointed family $F(\cdot\,;h_0)$ for
each fixed $h_0$.
\item[(b)] Part (c) of the hypercubic decomposition of symmetric tensors,
Lemma~\ref{10.4:lem-hypercubic-decomposition}.
\item[(c)] Parts (ii) and (iii) of the bounds on the typed marginal coefficient,
Lemma~\ref{10.4:lem-coefficient-bounds}.
\item[(d)] Part (b) of the lemma of this chapter on the counterterm family of the channel $\mathbf{6}$,
namely the bound on the counterterm family used in the proof below.
\end{enumerate}
Then there is a constant $C_I'$ such that for every $h_0\in(\mathsf{V}_{r})_{\mathbb{C}}$ the quantity
\begin{equation}\label{10.4:eq-typed-irrelevance-1}
\mathfrak{i}^{r}_y(\mathfrak{b};h_0)
:=\sum_{X\ni y}\bigl[F(X,\mathfrak{b},y;h_0)-F(X,1,y;h_0)\bigr]
-\beta^{r}[F]\,\mathsf{A}^{r}(h_y,\mathfrak{b};h_0)
\end{equation}
obeys
\begin{equation}\label{10.4:eq-typed-irrelevance-2}
\bigl|\mathfrak{i}^{r}_y(\mathfrak{b};h_0)\bigr|
\le C_I'\,\mathfrak{n}\,|h_0|\,\nu^{3}\rho_{n,j}^{2}\,(1+\rho_{n,j})\,t^{4+\hat\alpha}.
\end{equation}
Here $h_y$ is the scalar weight at $y$ of the irrelevance lemma of the correlation theorem. The
counterterm family $\mathsf{A}^{r}$ and the coefficient $\beta^{r}[F]$ are those of
Definition~\ref{10.4:def-typed-families}.
\end{lemma}

\begin{proof}
Fix $h_0$. We run the proof of the irrelevance lemma of the correlation theorem for the pointed family
$F(\cdot\,;h_0)$. Its steps are taken in their order, up to the point at which that proof appeals to
reflections and permutations of the axes, as hypothesis (a) provides.

\emph{Near levels and tail.} These parts of the proof use only norms of the family. They therefore
apply to $F(\cdot\,;h_0)$ as they stand. The norm $\|F\|$ of Definition~\ref{10.4:def-typed-families}
is the supremum over $|h_0|\le 1$ of the weighted norm used there.

The only new contribution is that of the counterterm. By part (ii) of
Lemma~\ref{10.4:lem-coefficient-bounds} and by hypothesis (d) it obeys, with a constant $C$,
\begin{equation}\label{10.4:eq-typed-irrelevance-3}
|\beta^{r}[F]|\,\bigl|\mathsf{A}^{r}(h_y,\mathfrak{b};h_0)\bigr|\le C\,\mathfrak{n}\,(1+\rho_{n,j}^{2}),
\end{equation}
and it enters these parts only through this bound.

\emph{Core.} The core of the proof consists of the following operations:
\begin{itemize}
\item the passage to the comb gauge;
\item the identity $\varphi_X'(0)=0$, where $\varphi_X$ is the function of the polymer $X$ so denoted
in that proof;
\item the appeal to part (c) of the auxiliary estimate of that proof;
\item the insertion of clause (c3) of that proof;
\item the extension of the sum to all $X\ni y$;
\item the reduction of the colour structure.
\end{itemize}
None of these operations uses a lattice symmetry other than translations, as is recorded with the
definition of pointed E-families of the correlation theorem, on which
Definition~\ref{10.4:def-typed-families} is modelled. Definition~\ref{10.4:def-typed-families}
requires of a typed family that it be translation covariant, and it does not require the symmetry
condition imposed on the families of the correlation theorem. Hence the core applies to
$F(\cdot\,;h_0)$ without change. It gives
\begin{equation}\label{10.4:eq-typed-irrelevance-4}
\sum_{X\ni y}\bigl[F(X,\mathfrak{b},y;h_0)-F(X,1,y;h_0)\bigr]
=\frac12\sum_{\kappa<\mu,\ \lambda<\nu'}\Pi_{\mu\nu',\kappa\lambda}(h_0)\,
\operatorname{tr}F_{\kappa\mu}(y)F_{\lambda\nu'}(y)+\text{(an error term)}.
\end{equation}
In \eqref{10.4:eq-typed-irrelevance-4}, $\Pi(h_0)$ is the second-moment tensor of
Definition~\ref{10.4:def-typed-families}, and $F_{\kappa\mu}(y)$ are the curvature components of
$\mathfrak{b}$ at $y$. The error term is bounded by a constant times the right-hand side
of \eqref{10.4:eq-typed-irrelevance-2}.

\emph{The quadratic part.} At this point the proof of the irrelevance lemma of the correlation
theorem uses reflections to show that only the diagonal entries $\mu=\nu'$, $\kappa=\lambda$ of the
tensor survive. It then uses permutations of the axes to show that these entries are all equal. Here
that argument is replaced by hypothesis (b).

The trace $\operatorname{tr}F_{\kappa\mu}(y)F_{\lambda\nu'}(y)$ is symmetric under the exchange of
the two curvature factors. Hence the quadratic term in \eqref{10.4:eq-typed-irrelevance-4} sees only
the part of $\Pi(h_0)$ that is symmetric under the exchange of the two pairs. This part is
$W_4$-equivariant in $h_0$. By part (c) of Lemma~\ref{10.4:lem-hypercubic-decomposition}, the space
$\operatorname{Hom}_{W_4}(\mathsf{V}_{r},\operatorname{Sym}^2\Lambda^2)$ is one-dimensional and is
spanned by $\mathsf{T}^{r}$. By the definition of $\beta^{r}[F]$ in
Definition~\ref{10.4:def-typed-families}, this part therefore equals
$\beta^{r}[F]\,\mathsf{T}^{r}(h_0)$.

Consequently the quadratic term of \eqref{10.4:eq-typed-irrelevance-4} is
$\beta^{r}[F]\,\mathsf{t}^{r}(h_0;F(y))$. Here $\mathsf{t}^{r}$ is the quadratic curvature form
belonging to $\mathsf{T}^{r}$.

\emph{The counterterm family.} The counterterm family $\mathsf{A}^{r}(h_y,\cdot\,;h_0)$ is an entire
pointed covariant family. It has radius $O(1)$ and is $\mathsf{V}_{r}$-typed. By part (iii) of
Lemma~\ref{10.4:lem-coefficient-bounds} its typed marginal coefficient equals $1$. The same analysis
as above, applied to this family, gives
\begin{equation}\label{10.4:eq-typed-irrelevance-5}
\mathsf{A}^{r}(h_y,\mathfrak{b};h_0)=\mathsf{t}^{r}(h_0;F(y))
+O\bigl(|h_0|\,\nu^{3}\mathfrak{a}_n^{2}\,t^{4+\hat\alpha}\bigr).
\end{equation}
This is \cite[(3.66)]{B-CMP119}, with the type carried along.

\emph{Cancellation.} Multiply \eqref{10.4:eq-typed-irrelevance-5} by $\beta^{r}[F]$ and subtract the
result from \eqref{10.4:eq-typed-irrelevance-4}. By \eqref{10.4:eq-typed-irrelevance-1}, the left side
becomes $\mathfrak{i}^{r}_y(\mathfrak{b};h_0)$. The marginal parts
$\beta^{r}[F]\,\mathsf{t}^{r}(h_0;F(y))$ cancel at the point $y$. What remains are the errors of the
two expansions, and they are estimated exactly as at the end of the proof of the irrelevance lemma of
the correlation theorem. Together with the near-level and tail bounds, this gives
\eqref{10.4:eq-typed-irrelevance-2}.
\end{proof}

\begin{remark}
Let $\mathsf{V}$ be a $W_4$-module with $\operatorname{Hom}_{W_4}(\mathsf{V},\operatorname{Sym}^2\Lambda^2)=0$.
For such a module the same argument applies with no counterterm: the part of $\Pi(h_0)$ seen by the
trace is then a $W_4$-equivariant map from $\mathsf{V}$ to $\operatorname{Sym}^2\Lambda^2$ and
therefore vanishes, and the bound \eqref{10.4:eq-typed-irrelevance-2} holds for
$\sum_{X\ni y}[F(X,\mathfrak{b},y;h_0)-F(X,1,y;h_0)]$ itself. This is the irrelevance lemma for
$\mathsf{V}$-typed pointed families.
The closing remarks of the irrelevance lemma of the correlation theorem hold verbatim for
Lemma~\ref{10.4:lem-typed-irrelevance}. The cancellation requires the core and the counterterm piece
to be taken on one and the same background. The tails are bounded singly.
\end{remark}

\begin{lemma}[Bounds for the off-diagonal tensor vertex]\label{10.4:lem-off-diagonal-vertex}
Throughout, $A(\varphi,U)=\sum_p\varphi\,[1-\operatorname{Re}\operatorname{tr}U(\partial p)]$ is the weighted
action of the correlation theorem. A weight on $\mathbb{R}^4$ is read at $x(p)$ and a weight on
plaquettes at $p$. We write $a_p:=1-\operatorname{Re}\operatorname{tr}U(\partial p)$.
Let $c$ be a complex scalar weight, and let $|h_0|\le 1$. In \textup{(b)} and in the last assertion
of \textup{(c)} the argument $h_0$ is suppressed from the notation.
\begin{itemize}
\item[(a)] There are polynomials $\tau'_p$, independent of $c$, each in the links within two lattice
units of the plaquette $p$, such that
\begin{equation}\label{10.4:eq-off-diagonal-vertex-1}
\mathsf{A}^{\mathbf{6}}(c,U;h_0)=\sum_p c(x(p))\,\tau'_p .
\end{equation}
\item[(b)] We work in the setting of part~(b) of the weight lemma of the correlation theorem, with
the following notation.
\begin{itemize}
\item[--] $\square$ is a unit cell of scale $n$.
\item[--] $\mathbf{U}$ is a configuration on the lattice of spacing $L^{k-n}\eta$, in the units of
scale $n$. In this part $\eta$ is the lattice spacing of the definition of the counterterm
families, not a rotation field.
\item[--] $H$ is the field of that setting by which $\mathbf{U}$ is displaced to
$e^{i\eta H}\mathbf{U}$; it is not the Hamiltonian.
\item[--] $\alpha$ and $a\le1$ are numbers; $a$ is a size parameter, not a lattice spacing.
\item[--] The norms of scale $n$ are those of that setting.
\end{itemize}
Suppose that on $\square$ every plaquette $p$ satisfies
\begin{equation*}
|\mathbf{U}(\partial p)-1|\le \alpha\,(L^{k-n}\eta)^2 ,
\end{equation*}
and that $H$ and $\nabla H$ are bounded by $a$ in the norms of scale $n$. Then, with a constant
$C_S'$ independent of $c$, $h_0$ and $\square$,
\begin{align*}
\bigl|\mathsf{A}^{\mathbf{6}}(c1_{\square},\mathbf{U})\bigr|
&\le C_S'\,\sup|c|\;\alpha^2,\\
\bigl|\mathsf{A}^{\mathbf{6}}(c1_{\square},e^{i\eta H}\mathbf{U})
-\mathsf{A}^{\mathbf{6}}(c1_{\square},\mathbf{U})\bigr|
&\le C_S'\,\sup|c|\;(\alpha+a)\,a .
\end{align*}
\item[(c)] Put $C_{\mathrm{dom}}:=4\cdot4\cdot21$ and, for a plaquette $p$,
\begin{equation*}
\widetilde{|c|}(p):=\max\bigl\{|c(x(p'))| : p'=p \text{ or } p' \text{ shares a link with } p\bigr\}.
\end{equation*}
Let $G=\mathrm{SU}(2)$, and let $G^{c}$ be its complexification, on which the counterterm families
are entire. For $U$ with values in $G$ (rather than in $G^{c}$),
\begin{equation}\label{10.4:eq-off-diagonal-vertex-2}
|\mathsf{A}^{\mathbf{6}}(c,U;h_0)|\le C_{\mathrm{dom}}\,A(\widetilde{|c|},U).
\end{equation}
Consequently, if $\widetilde{|c|}\le\vartheta\bar c$ with $\bar c\ge 0$, then
$|e^{\mathsf{A}^{\mathbf{6}}(c,U)}|\le e^{C_{\mathrm{dom}}\vartheta A(\bar c,U)}$.
\end{itemize}
\end{lemma}

\begin{proof}
(a) By the definition of the counterterm families (Definition~\ref{10.4:def-typed-families}),
\begin{equation*}
\mathsf{A}^{\mathbf{6}}(c,U;h_0)=\sum_p c(x(p))\sum_{\sigma\ne\rho}(h_0)_{\sigma\rho}\,
\mathfrak{s}_{\sigma\rho}(p).
\end{equation*}
This is \eqref{10.4:eq-off-diagonal-vertex-1} with
$\tau'_p:=\sum_{\sigma\ne\rho}(h_0)_{\sigma\rho}\mathfrak{s}_{\sigma\rho}(p)$.
These coefficients do not involve $c$. They are polynomials in the links within two lattice units
of $p$, since that definition forms $\mathfrak{s}_{\sigma\rho}(p)$ as such a polynomial.

(b) The entries $\mathfrak{s}_{\sigma\rho}(p)$ are formed from the link--plaquette variations
$\mathfrak{d}=\vec u\cdot\vec w$. Each of these is bilinear in the coefficients of
$\tfrac12(W-W^{-1})$ of two plaquettes.

Write such a term as $\Gamma(v_1,v_2)$, where $\Gamma$ is a bilinear form and $v_1,v_2$ are these
coefficient vectors for the configuration $\mathbf{U}$. Let $v_1',v_2'$ be the corresponding vectors
for $e^{i\eta H}\mathbf{U}$. Then $|\Gamma(v_1,v_2)|$ is at most a constant times $|v_1|\,|v_2|$.

Under the hypothesis on $\square$, each of $v_1,v_2$ is bounded by $\alpha(L^{k-n}\eta)^2$. Hence
$|\Gamma(v_1,v_2)|$ is bounded by a constant times $\alpha^2(L^{k-n}\eta)^4$.

Part~(b) of the weight lemma of the correlation theorem contains a computation, which applies here.
It shows that each coefficient vector moves by at most $4a(L^{k-n}\eta)^2(1+\alpha)$ when
$\mathbf{U}$ is replaced by $e^{i\eta H}\mathbf{U}$. Bilinearity then gives
\begin{equation*}
\Gamma(v_1',v_2')-\Gamma(v_1,v_2)=\Gamma(v_1'-v_1,\,v_2')+\Gamma(v_1,\,v_2'-v_2).
\end{equation*}
The size bound and the movement bound just stated estimate each term on the right.

The weight $c1_{\square}$ is bounded by $\sup|c|$. Each entry of $h_0$ is bounded by $|h_0|\le1$. The
cell contains $6(L^{k-n}\eta)^{-4}$ plaquettes. Multiplying the per-plaquette bounds by this number
gives the two inequalities of (b).

(c) Let $p=\{x;\mu,\nu\}$. By the diagonal lattice stress tensor lemma
(Lemma~\ref{10.4:lem-diagonal-stress}), $\mathfrak{s}_{\sigma\rho}(p)=0$ unless
$\sigma\in\{\mu,\nu\}$. The same lemma gives $\mathfrak{s}_{\mu\nu}(p)=\mathfrak{s}_{\nu\mu}(p)=0$.
Hence among the pairs with $\sigma\ne\rho$ only the four with $\sigma\in\{\mu,\nu\}$ and
$\rho\notin\{\mu,\nu\}$ contribute to $\tau'_p$.

Let $\mathcal{P}(p)$ denote the set consisting of $p$ and the plaquettes sharing a link with $p$.
By the entrywise bound for the lattice stress tensor at $G$-valued configurations,
\begin{equation*}
|\mathfrak{s}_{\sigma\rho}(p)|\le 4\sum_{p'\in\mathcal{P}(p)}a_{p'} .
\end{equation*}
A plaquette has four links, and each link lies in five further plaquettes. So $p$ shares a link with
$4\cdot5=20$ others, and $\mathcal{P}(p)$ has $21$ elements.

Using $|(h_0)_{\sigma\rho}|\le1$ and (a), we obtain
\begin{align*}
|\mathsf{A}^{\mathbf{6}}(c,U;h_0)|
&\le\sum_p|c(x(p))|\cdot4\cdot4\sum_{p'\in\mathcal{P}(p)}a_{p'}\\
&=4\cdot4\sum_{p'}a_{p'}\sum_{p\in\mathcal{P}(p')}|c(x(p))| .
\end{align*}
The equality uses that $p'\in\mathcal{P}(p)$ if and only if $p\in\mathcal{P}(p')$.

For $G$-valued $U$ we have $\operatorname{Re}\operatorname{tr}U(\partial p')\le1$, so $a_{p'}\ge0$.
Moreover $|c(x(p))|\le\widetilde{|c|}(p')$ for every $p\in\mathcal{P}(p')$. Hence the right-hand side
is at most
\begin{equation*}
4\cdot4\cdot21\sum_{p'}\widetilde{|c|}(p')\,a_{p'}.
\end{equation*}
This proves \eqref{10.4:eq-off-diagonal-vertex-2}.

Finally, $|e^{\mathsf{A}^{\mathbf{6}}(c,U)}|=e^{\operatorname{Re}\mathsf{A}^{\mathbf{6}}(c,U)}$. By
\eqref{10.4:eq-off-diagonal-vertex-2} this is at most $e^{C_{\mathrm{dom}}A(\widetilde{|c|},U)}$. Since
$a_p\ge0$ and $\widetilde{|c|}\le\vartheta\bar c$, we have
$A(\widetilde{|c|},U)\le\vartheta A(\bar c,U)$, which gives the last assertion.
\end{proof}

\begin{remark}
For the channel $\mathbf{3}$ no separate statement is needed. There
$\mathsf{A}^{\mathbf{3}}(\varphi,U;h)=A(\varphi\mathsf{c}_h,U)$ is a weighted action. Parts (a)--(c)
of the weight lemma of the correlation theorem apply to it as they stand with $c=\varphi\mathsf{c}_h$,
since that lemma requires only that $c$ be a complex weight.
\end{remark}

\begin{definition}[Tensor sources and the frozen tangent run]\label{10.4:def-tensor-sources}
Throughout, the parameters of the correlation theorem are in force, together with the order in
which they are chosen and the regime of that theorem. Two further standing conditions are
assumed: the standing condition in which the constant $\gamma_L$ is fixed, and the standing
condition in which the constant $\gamma_U$ is fixed. The quantities
$x_k$, $\zeta_k$, $b_k$, $E^{(k)}$, $Q^{(k)}$ of the
correlation theorem are taken throughout in their hatted form, as defined in the re-filing
lemmas that accompany the uniqueness theorem; the same quantities without hat are treated in a
separate remark of this chapter. Let
$r\in\{\mathbf{3},\mathbf{6}\}$, and let
$\mathsf{k}\colon\mathbb{T}_m\to(\mathsf{V}_r)_{\mathbb{C}}$ be Lipschitz, with
\begin{equation*}
\|\mathsf{k}\|_\sharp:=\max\bigl\{\|\mathsf{k}\|_\infty,\ \operatorname{Lip}\mathsf{k}/(\lambda_*N)\bigr\},
\end{equation*}
where $\lambda_*$ is the constant of that name among the parameters of the correlation theorem.
For tensor weights the norm $\|\cdot\|_\sharp$ is normalised in this way; this differs from the
norm $\|f\|_\sharp$ of scalar test functions.
$\mathcal{H}_j(X)$ denotes the weight classes of typed families, taken with $\mathsf{V}=\mathsf{V}_r$.
Write $a=L^{-K}$ for the bare lattice spacing (not to be confused with the plaquette quantity
$a_p$). The plaquette variables $a_p$, the lattice stress tensor
$\mathfrak{s}_{\sigma\rho}(p)$, the second difference $\Delta_\nu$, the functionals
$\mathfrak{Q}_\nu$ and the notation $p\ni\nu$ (the direction $\nu$ is one of the two directions
spanning $p$) are those of Lemma~\ref{10.4:lem-diagonal-stress}.

\emph{Channel sources.} On the bare lattice put $\mathfrak{T}^{r}(\mathsf{k}):=\mathsf{A}^{r}(\mathsf{k},U)$,
the counterterm family of Definition~\ref{10.4:def-typed-families} with tensor weight
$\mathsf{k}$. Explicitly,
\begin{align*}
\mathfrak{T}^{\mathbf{3}}(\mathsf{k})&=\sum_p\,[\mathsf{k}_{\mu\mu}+\mathsf{k}_{\nu\nu}](x(p))\,a_p
\qquad(p\parallel(\mu\nu)),\\
\mathfrak{T}^{\mathbf{6}}(\mathsf{k})&=\sum_p\sum_{\sigma\neq\rho}\mathsf{k}_{\sigma\rho}(x(p))\,
\mathfrak{s}_{\sigma\rho}(p),
\end{align*}
and
\begin{equation*}
Z(z,t):=\int\exp\Bigl[-g_0^{-2}A(U)+\sum_i z_i\,\mathcal{O}(f_i)+t\,\mathfrak{T}^{r}(\mathsf{k})\Bigr]\,dU.
\end{equation*}

\emph{The full tensor source.} For $\mathsf{k}\colon\mathbb{T}_m\to(\operatorname{Sym}^2V)_{\mathbb{C}}$
Lipschitz, where $V=\mathbb{R}^4$ with its $W_4$-action, put, with no factor $g_0^{-2}$,
\begin{equation}\label{10.4:eq-tensor-sources-1}
\mathfrak{T}^{\otimes}(\mathsf{k}):=\sum_p\sum_{\sigma,\rho}\mathsf{k}_{\sigma\rho}(x(p))
\Bigl[\tfrac12\bigl(\mathfrak{s}_{\sigma\rho}+\mathfrak{s}_{\rho\sigma}\bigr)(p)
-\delta_{\sigma\rho}\,a_p\Bigr].
\end{equation}
Write $\mathsf{k}=\tfrac14(\operatorname{tr}\mathsf{k})\,\delta+\mathsf{k}^{\mathbf{3}}+\mathsf{k}^{\mathbf{6}}$,
where $\operatorname{tr}\mathsf{k}:=\sum_\sigma\mathsf{k}_{\sigma\sigma}$ is the tensor trace (so
$\operatorname{tr}\delta=4$), not the normalised colour trace.
Here $\mathsf{k}^{\mathbf{3}}$ takes values in $(\mathsf{V}_{\mathbf{3}})_{\mathbb{C}}$, the traceless
diagonal tensors, and $\mathsf{k}^{\mathbf{6}}$ takes values in $(\mathsf{V}_{\mathbf{6}})_{\mathbb{C}}$,
the off-diagonal ones. Then the following identity holds; it is proved after the definition:
\begin{equation}\label{10.4:eq-tensor-sources-a}
\begin{split}
\mathfrak{T}^{\otimes}(\mathsf{k})
={}&\mathfrak{T}^{\mathbf{6}}(\mathsf{k}^{\mathbf{6}})+2\,\mathfrak{T}^{\mathbf{3}}(\mathsf{k}^{\mathbf{3}})
-\sum_\nu\mathfrak{Q}_\nu(\mathsf{k}_{\nu\nu})\\
&\quad+\tfrac14 a^2\sum_\nu\sum_{p\ni\nu}(\Delta_\nu\mathsf{k}_{\nu\nu})(x(p))\,a_p .
\end{split}
\end{equation}
Thus the trace cancels identically at the marginal level. The last two terms are bare units:
the first has no marginal part, and the second carries an explicit factor $a^2$.

\emph{Frozen tangent run.} Let $\zeta\in\mathbb{C}$ with $|\zeta|<2\bar r$, where $\bar r$
denotes the radius called $r$ among the parameters of the correlation theorem, renamed here
because $r$ is the channel index; and let $h_0\in(\mathsf{V}_r)_{\mathbb{C}}$ be
constant. Put
\begin{equation*}
\rho_0^{\zeta,t}:=\exp\bigl[-A(x_0-\zeta,U)+t\,\mathsf{A}^{r}(1,U;h_0)\bigr].
\end{equation*}
Here, for a plaquette weight $\varphi$,
$A(\varphi,U):=\sum_p\varphi(p)\,[1-\operatorname{Re}\operatorname{tr}U(\partial p)]$ is the action
with weight $\varphi$; for the constant weight $x_0-\zeta$ it equals $(x_0-\zeta)A(U)$.
A \emph{frozen tangent run} is the $t$-derivative at $t=0$ of the frozen whole-lattice run of the
correlation theorem started from $\rho_0^{\zeta,t}$.
\end{definition}

\begin{proof}
We prove \eqref{10.4:eq-tensor-sources-a}. Since $\mathsf{k}_{\rho\sigma}=\mathsf{k}_{\sigma\rho}$,
we may exchange the names of the two summation indices in the half of
\eqref{10.4:eq-tensor-sources-1} that contains $\mathfrak{s}_{\rho\sigma}$. This gives
\begin{equation*}
\mathfrak{T}^{\otimes}(\mathsf{k})=\sum_p\sum_{\sigma\neq\rho}\mathsf{k}_{\sigma\rho}(x(p))\,
\mathfrak{s}_{\sigma\rho}(p)
+\sum_\nu\sum_p\mathsf{k}_{\nu\nu}(x(p))\bigl[\mathfrak{s}_{\nu\nu}(p)-a_p\bigr].
\end{equation*}
For $\sigma\neq\rho$ the first two pieces of the decomposition of $\mathsf{k}$ are diagonal, so
$\mathsf{k}_{\sigma\rho}=\mathsf{k}^{\mathbf{6}}_{\sigma\rho}$. The first sum is therefore
$\mathfrak{T}^{\mathbf{6}}(\mathsf{k}^{\mathbf{6}})$.

Fix $\nu$. By Lemma~\ref{10.4:lem-diagonal-stress}, $\mathfrak{s}_{\nu\nu}(p)=0$ unless $p\ni\nu$.
The consequence stated in that lemma, applied with $\varphi=\mathsf{k}_{\nu\nu}$, so that
$\tilde\varphi_\nu=\mathsf{k}_{\nu\nu}+\tfrac18a^2\Delta_\nu\mathsf{k}_{\nu\nu}$, then gives
\begin{equation*}
\begin{split}
\sum_p\mathsf{k}_{\nu\nu}(x(p))\,\mathfrak{s}_{\nu\nu}(p)
={}&2\sum_{p\ni\nu}\mathsf{k}_{\nu\nu}(x(p))\,a_p\\
&\quad+\tfrac14a^2\sum_{p\ni\nu}(\Delta_\nu\mathsf{k}_{\nu\nu})(x(p))\,a_p-\mathfrak{Q}_\nu(\mathsf{k}_{\nu\nu}).
\end{split}
\end{equation*}
Sum over $\nu$. Insert $\mathsf{k}_{\nu\nu}=\tfrac14\operatorname{tr}\mathsf{k}+\mathsf{k}^{\mathbf{3}}_{\nu\nu}$
into the first term on the right, and no other term. Every plaquette $p\parallel(\mu\nu)$ contains
exactly the two directions $\mu,\nu$. Hence
\begin{equation*}
\sum_\nu\sum_{p\ni\nu}\mathsf{k}^{\mathbf{3}}_{\nu\nu}(x(p))\,a_p
=\mathfrak{T}^{\mathbf{3}}(\mathsf{k}^{\mathbf{3}})
\end{equation*}
and
\begin{equation*}
\sum_\nu\sum_{p\ni\nu}\tfrac14(\operatorname{tr}\mathsf{k})(x(p))\,a_p
=\tfrac12\sum_p(\operatorname{tr}\mathsf{k})(x(p))\,a_p .
\end{equation*}
On the other hand,
\begin{equation*}
\sum_\nu\sum_p\mathsf{k}_{\nu\nu}(x(p))\,a_p=\sum_p(\operatorname{tr}\mathsf{k})(x(p))\,a_p .
\end{equation*}
Collecting terms, the diagonal contribution equals
\begin{equation*}
\begin{split}
&2\,\mathfrak{T}^{\mathbf{3}}(\mathsf{k}^{\mathbf{3}})
+\Bigl(2\cdot\tfrac12-1\Bigr)\sum_p(\operatorname{tr}\mathsf{k})(x(p))\,a_p\\
&\quad+\tfrac14a^2\sum_\nu\sum_{p\ni\nu}(\Delta_\nu\mathsf{k}_{\nu\nu})(x(p))\,a_p
-\sum_\nu\mathfrak{Q}_\nu(\mathsf{k}_{\nu\nu}).
\end{split}
\end{equation*}
The trace term has coefficient zero. Adding $\mathfrak{T}^{\mathbf{6}}(\mathsf{k}^{\mathbf{6}})$
gives \eqref{10.4:eq-tensor-sources-a}.
\end{proof}

\begin{theorem}[The stress-tensor-insertion correlation theorem]
\label{10.4:thm-insertion-correlation}
Work in the setting of Definition~\ref{10.4:def-tensor-sources}, with channel
$r\in\{\mathbf{3},\mathbf{6}\}$. Assume the following statements, which enter as hypotheses:
\begin{itemize}
\item the correlation theorem, taken by statement;
\item clauses (a)--(e) and (g) of the one-lattice Lipschitz step, the extension proposition, and
geometric forgetting of old regions;
\item clause (c) of the main theorem of the complex-weight small-field analysis that fixes the
ceiling $\gamma_U$, together with the bound of that analysis for complex weights;
\item the re-filing lemmas;
\item the typed step with a marginal tensor slot \eqref{10.4:eq-typed-step-a}, in the form that
includes the lemma on irrelevance after one typed counterterm
(Lemma~\ref{10.4:lem-typed-irrelevance}), the
fluctuation kernel for typed slots \eqref{10.4:eq-typed-slot-kernel-a}, the response of one step to
the tensor vertex \eqref{10.4:eq-vertex-response-a}, and the re-filing of typed marginal families
\eqref{10.4:eq-typed-refiling-a};
\item the torus bound in its typed form, and the form of the extension proposition used with the
correlation theorem;
\item if $r=\mathbf{6}$, in addition the sourced format of the local tangent run in the form
\eqref{10.4:eq-sourced-format-b}.
\end{itemize}
Then there is $\gamma^{\otimes}\in(0,\min\{\gamma_L,\gamma_U\}]$ ($\gamma_L$ the ceiling on $g$
fixed with the one-lattice Lipschitz step), chosen subject to the ceilings
\eqref{10.4:eq-vertex-response-a} and \eqref{10.4:eq-forward-transport-a} after every constant in
the order of choice of the correlation theorem and after $(\vartheta_1,q,c_\natural)$, such that for
$g\le\gamma^{\otimes}$ and for the parameters $(x_0,K,m,N_T)$ of the correlation theorem
($x_0$ the initial inverse coupling) the following hold.
\begin{enumerate}
\item[(a)] For $k\le K$ the frozen tangent run has the following format. For $r=\mathbf{3}$ the
weight is $x_k-\zeta_k-t\tau_k\mathsf{c}_{h_0}$; for $r=\mathbf{6}$ the weight is
$x_k-\zeta_k$ and the format carries in addition the vertex
$+t\tau_k\mathsf{A}^{\mathbf{6}}(1,U_k;h_0)$. The running tensor coefficient is
$\tau_k=N_r(0,k;\zeta)$, and the typed slots are
$\delta\mathsf{P}^{(j)},\delta P^{\flat(j)},\delta Q^{(j)},\delta R^{(j)}$, $j\le k$. Here
$N_r(0,0)=1$, the function $N_r(0,k;\cdot)$ is analytic on the source disc of the correlation
theorem (the open disc about $0$ of twice its source radius) and real on reals, and
\begin{equation}\label{10.4:eq-insertion-correlation-a}
\begin{gathered}
N_r(0,k+1)=(1+\nu^{r}_k)\,N_r(0,k)+\iota_{k+1},\\
\nu^{r}_k(\zeta):=\beta^{r}\bigl[\delta^{r}\Phi_{k+1}
\bigl(g_k,s_k;\mathsf{P}^{(\le k)}_\xi\bigr)\bigr]\Big|_{\xi_i=\zeta_i(\zeta)},\qquad
|\nu^{r}_k|\le C_\nu\, g_k^2(\log x_k)^{2p_0},\\
|\iota_{k+1}|\le\sum_{j\le k}\mathsf{i}_{k+1,j}\,|N_r(0,j)|,\qquad
\sum_{k<K}\sum_{j\le k}\mathsf{i}_{k+1,j}\,c_\natural^{\,k+1-j}\le\tfrac16 ,\\
\text{hence}\quad N_r(0,k;\zeta)=\bigl(1+\epsilon_k(\zeta)\bigr)P^{r}_k(\zeta),\\
P^{r}_k:=\prod_{l<k}(1+\nu^{r}_l),\qquad |\epsilon_k|\le\tfrac13 ;
\end{gathered}
\end{equation}
here $\delta^{r}\Phi_{k+1}$ is the typed tangent of the family $\Phi_{k+1}(\mathfrak{q},s,u;h)$ of
the $(k+1)$-st step, defined beside \eqref{10.4:eq-vertex-response-a}, where also the argument
$s_k$ and the slot variables $\xi=(\xi_i)_{i\le k}$ of the pure families
$\mathsf{P}^{(\le k)}_\xi$ are fixed; the $\xi_i$ are evaluated at the shifts $\zeta_i(\zeta)$ of
the running-shift statement of the correlation theorem (the letters $\xi_i$ here are unrelated to
the rotation field of this chapter).
In particular $N_r(0,k)\neq0$, and one sets $N_r(j,k):=N_r(0,k)/N_r(0,j)$.
\item[(b)] The increment has the one-loop form
\begin{equation}\label{10.4:eq-insertion-correlation-b}
\nu^{r}_k(\zeta)=-\frac{\mathsf{m}_r(k+1)}{x_k-\zeta_k(\zeta)}+\varrho^{r}_k(\zeta),\qquad
\mathsf{m}_r(k+1):=-\beta^{r}\bigl[\mathsf{L}^{r}_{k+1}(L^{(\le k)})\bigr],
\end{equation}
where $\mathsf{m}_r(k+1)$ is a real number independent of the sources $z$, formed from the
scaffold of level $k$, the one-loop families $L^{(j)}$, $j\le k$, and the field cut $T_k$ at
level $k$ (Lemma~\ref{10.6:lem-cutoff-cost});
here $\mathsf{L}^{r}_{k+1}$ is the one-loop part of the family $\Phi_{k+1}$, defined beside
\eqref{10.4:eq-vertex-response-a}; and
\begin{equation*}
|\mathsf{m}_r(k+1)|\le C_{\mathsf{m}}(\log x_k)^{2p_0},\qquad
|\varrho^{r}_k|\le C_\varrho\, g_k^3(\log x_k)^{4p_0}.
\end{equation*}
\item[(c)] The step maps $R^\flat$, $T^{w,t}$, the initial sourced density $\rho_0^{w,t}$, the
weights $\varphi_j$, the kernels $h^{(j)}_y$, the
points $p_y$ and the region $X^{+}$ below are those of the sourced format
\eqref{10.4:eq-sourced-format-a}. For $k\le K'$, $w\in\mathcal{W}_k$ and $\mathsf{k}$ with
$\mathsf{k}/\|\mathsf{k}\|_\sharp\in\mathcal{H}_k$, the derivative
$\partial_t|_{t=0}(R^\flat T^{w,t})^k\rho_0^{w,t}$ has the first-order format
\eqref{10.4:eq-sourced-format-a}. It consists of the vertex with local weight
\begin{equation}\label{10.4:eq-insertion-correlation-c}
\begin{split}
\tau_k(p):={}&\mathsf{k}(p)+\sum_{j\le k}\varphi_j(p)\sum_y h^{(j)}_y(p)\\
&\times\bigl[N_r(0,j;w(p_y))-N_r(0,j-1;w(p_y))\bigr]\mathsf{k}(p_y),
\end{split}
\end{equation}
of typed groups $\mathfrak{i}^{r,(j)}_y$ on the frozen tangent families, of freezing defects, and
of typed $\mathbf{R}$- and $\mathbf{B}$-terms. All of these are linear in $\mathsf{k}$, analytic in
$w$, and depend on $(w,\mathsf{k})$ only through their restrictions to $X^{+}$; each of them has
flat norm at most $C_{\otimes}\bar\varkappa\,\mathfrak{g}_j\,|P^{r}_j(0)|\,\|\mathsf{k}\|_\sharp$,
with a constant $C_\otimes$, where $j$ is the scale index that the term carries in the format
\eqref{10.4:eq-sourced-format-a}.
\item[(d)] For $n\ge1$, if the supports of $\mathsf{k},f_1,\dots,f_n$ are pairwise at distance at
least $\mathsf{d}>0$, then
\begin{equation}\label{10.4:eq-insertion-correlation-d}
\bigl|\langle\mathfrak{T}^{r}(\mathsf{k});\mathcal{O}(f_1);\dots;\mathcal{O}(f_n)\rangle_{K,m}\bigr|
\le C^{\otimes}_n(\mathsf{d})\,|P^{r}_{K'}(0)|\,\|\mathsf{k}\|_\sharp\prod_{i=1}^n\|f_i\|_\sharp ,
\end{equation}
where, with $Z(z,t)$ as in Definition~\ref{10.4:def-tensor-sources}, formed at cutoff $K$ on
$\mathbb{T}_m$, the truncated correlation is
\begin{equation*}
\langle\mathfrak{T}^{r}(\mathsf{k});\mathcal{O}(f_1);\dots;\mathcal{O}(f_n)\rangle_{K,m}
:=\partial_t\partial_{z_1}\cdots\partial_{z_n}\log Z(z,t)\Big|_{z=0,\,t=0},
\end{equation*}
and the constant $C^{\otimes}_n(\mathsf{d})$ depends only on
$L,\mathfrak{p},g,m,n,\mathsf{d}$ and $\vartheta_1,q,c_\natural$, not on $K$, $g_0$ or the
positions of the supports; its dependence on the torus is through the factor $\mathsf{Q}^{n+1}$
only.
\end{enumerate}
\end{theorem}

The proof of Theorem~\ref{10.4:thm-insertion-correlation} is assembled in the course of this
section from the hypotheses listed in its statement.

\begin{corollary}[Decay of the tensor correlation under normalisation decay]
\label{10.4:cor-tensor-correlation-decay}
Let $r\in\{\mathbf{3},\mathbf{6}\}$, and let the coupling $g$ and all data of the
stress-tensor-insertion correlation theorem (Theorem~\ref{10.4:thm-insertion-correlation}) be as
in that theorem; in particular $x_0$, $K$, $m$ are its data and $K'$ is its last sourced level.
Assume:
\begin{itemize}
\item[(a)] the bound \eqref{10.4:eq-insertion-correlation-d} of that theorem, with its constant
$C^{\otimes}_n(\mathsf{d})$, for $n\ge 1$ and for test functions whose supports, together with the
support of $\mathsf{k}$, are pairwise at distance at least $\mathsf{d}>0$;
\item[(b)] the decay of the normalisation: there are $\kappa\in\mathbb{R}$ (here and below $\kappa$
denotes this exponent, not the auxiliary parameter of $\mathfrak{p}$) and a constant $C_N$
independent of $K$ (the subscript in $C_N$ refers to the normalisation and indicates no
dependence on the rank $N$) such that
\begin{equation}\label{10.4:eq-tensor-correlation-decay-a}
\Bigl|\frac{P^{r}_k(0)}{P^{r}_j(0)}\Bigr| \le C_N\Bigl(\frac{x_k}{x_j}\Bigr)^{\kappa}
\qquad\text{for all } 0\le j\le k\le K ;
\end{equation}
only an upper bound on the ratio is assumed.
\end{itemize}
Then, for such $\mathsf{k},f_1,\dots,f_n$,
\begin{equation}\label{10.4:eq-tensor-correlation-decay-1}
\bigl|\langle \mathfrak{T}^{r}(\mathsf{k});\mathcal{O}(f_1);\dots;\mathcal{O}(f_n)\rangle_{K,m}\bigr|
\le C^{\otimes}_n(\mathsf{d})\,C_N\Bigl(\frac{x_{K'}}{x_0}\Bigr)^{\kappa}
\|\mathsf{k}\|_\sharp\prod_{i=1}^n\|f_i\|_\sharp .
\end{equation}
Assume moreover $\kappa\ge 0$ and $K\ge 1$, and recall that $K=K(g_0,g)$ is the first-passage
level of the data of that theorem, for which the first-passage relation
$x_0\ge X+\lambda K-2c_2$ holds. If $\lambda>0$ and $X\ge 2c_2$, the right-hand side of
\eqref{10.4:eq-tensor-correlation-decay-1} is bounded as follows:
\begin{equation}\label{10.4:eq-tensor-correlation-decay-2}
\begin{split}
&C^{\otimes}_n(\mathsf{d})\,C_N\Bigl(\frac{x_{K'}}{x_0}\Bigr)^{\kappa}
\|\mathsf{k}\|_\sharp\prod_{i=1}^n\|f_i\|_\sharp\\
&\qquad\le C^{\otimes}_n(\mathsf{d})\,C_N\,\lambda^{-\kappa}\,x_{K'}^{\kappa}\,K^{-\kappa}\,
\|\mathsf{k}\|_\sharp\prod_{i=1}^n\|f_i\|_\sharp .
\end{split}
\end{equation}
If in addition there is a constant $C'(g,m)$, depending only on $g$ and $m$, with
$x_{K'}\le C'(g,m)$, then the right-hand side of \eqref{10.4:eq-tensor-correlation-decay-2} is at
most $C(g,m)K^{-\kappa}\|\mathsf{k}\|_\sharp\prod_{i=1}^n\|f_i\|_\sharp$, where
\begin{equation*}
C(g,m)=C^{\otimes}_n(\mathsf{d})\,C_N\,\lambda^{-\kappa}\,C'(g,m)^{\kappa},
\end{equation*}
a constant which depends, besides on $g$ and $m$, also on $n$, $\mathsf{d}$, $r$, $\kappa$, $C_N$
and on the quantities read by $C^{\otimes}_n(\mathsf{d})$.
For $\kappa<0$ only \eqref{10.4:eq-tensor-correlation-decay-1} is asserted.
\end{corollary}

\begin{proof}
Apply
\eqref{10.4:eq-tensor-correlation-decay-a} with $j=0$ and $k=K'$, which is admissible since
$0\le K'\le K$. Since $P^{r}_0(0)=1$ (at $k=0$ the product defining $P^{r}_k$ in
Theorem~\ref{10.4:thm-insertion-correlation} is empty, in accordance with $N_r(0,0)=1$), this gives
\begin{equation*}
|P^{r}_{K'}(0)|\le C_N\Bigl(\frac{x_{K'}}{x_0}\Bigr)^{\kappa}.
\end{equation*}
Inserting this into the right-hand side of \eqref{10.4:eq-insertion-correlation-d}, which is
available by hypothesis (a), gives \eqref{10.4:eq-tensor-correlation-decay-1}.

For \eqref{10.4:eq-tensor-correlation-decay-2} let $\kappa\ge0$, $K\ge1$, $\lambda>0$ and
$X\ge 2c_2$. Then $X-2c_2\ge 0$, and the first-passage relation gives
\begin{equation*}
x_0\ge X+\lambda K-2c_2\ge \lambda K>0 .
\end{equation*}
The inverse couplings $x_k=g_k^{-2}$ are positive, and $y\mapsto y^{-\kappa}$ is non-increasing on
$]0,\infty[$; hence $x_0^{-\kappa}\le(\lambda K)^{-\kappa}=\lambda^{-\kappa}K^{-\kappa}$, whence
\begin{equation*}
\Bigl(\frac{x_{K'}}{x_0}\Bigr)^{\kappa}
= x_{K'}^{\kappa}\,x_0^{-\kappa}
\le \lambda^{-\kappa}\,x_{K'}^{\kappa}\,K^{-\kappa}.
\end{equation*}
Multiplying by $C^{\otimes}_n(\mathsf{d})\,C_N\,\|\mathsf{k}\|_\sharp\prod_{i=1}^n\|f_i\|_\sharp$
gives \eqref{10.4:eq-tensor-correlation-decay-2}. If moreover $x_{K'}\le C'(g,m)$, then, since
$x_{K'}>0$ and $y\mapsto y^{\kappa}$ is non-decreasing on $]0,\infty[$, we have
$x_{K'}^{\kappa}\le C'(g,m)^{\kappa}$; inserting this into the right-hand side of
\eqref{10.4:eq-tensor-correlation-decay-2} gives the bound
$C(g,m)K^{-\kappa}\|\mathsf{k}\|_\sharp\prod_{i=1}^n\|f_i\|_\sharp$ with the constant $C(g,m)$ of
the statement.
\end{proof}

The exponents of interest in \eqref{10.4:eq-tensor-correlation-decay-a} are the target exponent
$\kappa>-\tfrac12$ and the strongest form $\kappa=1$.

\begin{remark}[Limits of the first-order theorem]\label{10.4:rem-first-order-limits}
Let $t$ denote the parameter multiplying the tensor source in the stress-tensor-insertion
correlation theorem (Theorem~\ref{10.4:thm-insertion-correlation}). That theorem does not
cover the following three cases.
\begin{enumerate}
\item Second order in $t$. At that order the representation
$\operatorname{Sym}^2(\operatorname{Sym}^2 V)$, with $V=\mathbb{R}^4$, occurs, and it meets the
channels $\mathbf{10}$ and $\mathbf{2}$ of the anisotropy analysis.
\item Correlations with two tensor insertions $\mathfrak{T}^{r}(\mathsf{k})$ and
$\mathfrak{T}^{r'}(\mathsf{k}')$, with channels $r,r'$ and tensor weights
$\mathsf{k},\mathsf{k}'$.
\item The case in which the support of $\mathsf{k}$ overlaps the support of one of the $f_i$ in
the torus bound \eqref{10.4:eq-insertion-correlation-d} of
Theorem~\ref{10.4:thm-insertion-correlation}. By contrast, the first-order format
\eqref{10.4:eq-insertion-correlation-c} carries no separation hypothesis.
\end{enumerate}
\end{remark}

\begin{proof}
Each item follows from the form of the hypotheses and conclusions of
Theorem~\ref{10.4:thm-insertion-correlation}.

For the first item: each conclusion of that theorem is either a statement about the tangent
run, that is, about the derivative $\partial_t$ at $t=0$, or a bound on a connected correlation
containing exactly one insertion $\mathfrak{T}^{r}(\mathsf{k})$; no conclusion concerns second
order in $t$. The theorem is therefore a statement at first order in $t$ only.

For the second item: the running-format and one-loop statements of that theorem do not involve
a tensor weight, while \eqref{10.4:eq-insertion-correlation-c} is linear in a single tensor
weight $\mathsf{k}$ and the correlation bounded in \eqref{10.4:eq-insertion-correlation-d}
contains exactly one insertion $\mathfrak{T}^{r}(\mathsf{k})$.

For the third item: the bound \eqref{10.4:eq-insertion-correlation-d} is asserted only under
the hypothesis that the supports of $\mathsf{k},f_1,\dots,f_n$ are pairwise at distance at
least $\mathsf{d}>0$. The statement \eqref{10.4:eq-insertion-correlation-c}, on the other hand,
is asserted for every $w\in\mathcal{W}_k$ and every $\mathsf{k}$ with
$\mathsf{k}/\|\mathsf{k}\|_\sharp\in\mathcal{H}_k$, and imposes no condition on supports.
\end{proof}

\begin{proof}[Transfer of the scalar lemmas to a tensor slot. Stated; proof incomplete at the step marked $(\ast)$]\label{10.4:prf-tensor-slot-transfer}
Let $r\in\{\mathbf{3},\mathbf{6}\}$. In the proof of the stress-tensor-insertion correlation theorem
(Theorem~\ref{10.4:thm-insertion-correlation}), the table below records, for each item of the
correlation theorem, what that item reads of the source term and how it is read in the channel $r$.
For $r=\mathbf{6}$, the weight-linear bound and the third item of the Wilson-shift bound are replaced
by two statements of the sourced format \eqref{10.4:eq-sourced-format-b}, and these two statements
are not proved here.
The list, given with the vector package (Proposition~\ref{10.3:prop-vector-package}), of the places
at which the correlation theorem reads clause (f3) of its format is gone through, place by place,
at the end of the proof.

In the table the items of the correlation theorem are named by their content. The second column
says what the item reads of the source term. An entry ``verbatim'' means that the item, with its
proof, applies without change. An entry ``same'' in the column $r=\mathbf{6}$ means that the entry
of the column $r=\mathbf{3}$ applies. Here $t$ is the parameter multiplying the tensor source, as in
Theorem~\ref{10.4:thm-insertion-correlation}. The letters $G$ (the gauge-fixing term), $a_j$ (the
gauge-fixing coefficient), $a_p$ (the plaquette density), $\theta_W$ and $x$ (the constant and the
inverse coupling of the Wilson-shift bound), $\hat\chi_\square$ (the cube partition of unity), $V$
(the potential of the step), $\mathfrak{N}(v)$ (the unit functional, evaluated on a unit $v$) and
$\mathfrak{a}$ (the radius of the family at the level read) are those of the correlation theorem;
$C_S'$ is the constant of Lemma~\ref{10.4:lem-off-diagonal-vertex}(b).

\begin{center}
\small
\begin{tabular}{|p{0.20\textwidth}|p{0.22\textwidth}|p{0.18\textwidth}|p{0.18\textwidth}|}
\hline
item of the correlation theorem & what it reads of the source term
 & $r=\mathbf{3}$ (plaquette weight $\mathsf{c}_{\mathsf{k}}$) & $r=\mathbf{6}$ \\
\hline
weighted action $A(c,U)$, initial density $\rho^w_0$, and the level-zero clause
 & the weight, linearly
 & verbatim with the weight \eqref{10.4:eq-tensor-slot-transfer-1}; no new slot
 & level-zero vertex read directly (see below) \\
\hline
the items of the correlation theorem's second section & nothing & verbatim & verbatim \\
\hline
gauge-fixing lemma & one coefficient on $G+A$; criticality of $U_{k+1}$
 & not used: the vertex is a potential, by the typed complex-factor proposition; $a_j$ stays
 free of $t$ & same \\
\hline
source-reading lemma, items (B1)--(B9) & $h$ linearly; (B1)--(B7): (f1), (f2), norm;
 (B8), (B9): (f3) & typed items (B8), (B9), whose proof is outstanding & same \\
\hline
format-and-activity lemma & $V$ through format and activity & verbatim & verbatim \\
\hline
complex-factor proposition & a scalar $s$ & the typed complex-factor proposition & same \\
\hline
conditions (c1)--(c3), ($\ell$0), the chains & nothing depending on $z$ & verbatim & verbatim \\
\hline
the coefficient lemma, the marginal-coefficient lemma and its flat form
 & coefficient lemma: nothing; marginal-coefficient lemma and its flat form: a tensor, then one
 number
 & coefficient lemma verbatim; the other two: Lemma~\ref{10.4:lem-coefficient-bounds} & same \\
\hline
covariance-free lemma & no covariance, no other region $X$ & verbatim & verbatim \\
\hline
irrelevance lemma & one sentence on $W_4$ & Lemma~\ref{10.4:lem-typed-irrelevance} & same \\
\hline
vertex lemma & $\sup|c|$, linearity, $a_p\ge 0$ & verbatim
 & Lemma~\ref{10.4:lem-off-diagonal-vertex} \\
\hline
analyticity lemma & analyticity on $\mathcal{W}_j$ & verbatim; for the slot $\mathsf{k}$, the
 analyticity statement of Proposition~\ref{10.3:prop-vector-package} & same \\
\hline
the lemma on units, the summation lemma and its flat form, the lemma on the complex potential,
 the homogeneity condition, and the flat section of the correlation theorem
 & units through the unit functional $\mathfrak{N}(v)$ of the correlation theorem; parameters
 under parameter-free majorants
 & verbatim ($t$ one more holomorphic parameter) & verbatim \\
\hline
five-line decomposition of the units; derivative lemma & lines: groups; analyticity and
 covariance-free lemmas; covariance-free lemma; derived; vertex; derivative lemma singly
 & first line by Lemma~\ref{10.4:lem-typed-irrelevance}; fifth line contributes (see below)
 & same, with Lemma~\ref{10.4:lem-off-diagonal-vertex}(b) \\
\hline
identification $P^{\flat}\simeq\mathsf{P}$ & the two integrals have one integrand
 & the $t$-derivative of that identity & same \\
\hline
identity lemma & identities linear in the integrand; gauge-fixing identity at any coefficient
 & verbatim & verbatim \\
\hline
weight-linear bound & weight-linear, bounded by weighted absolute sums & verbatim
 & the sourced format \eqref{10.4:eq-sourced-format-b} $(\ast)$ \\
\hline
Wilson-shift bound, third item & raw Wilson shift at most $\theta_W x$
 & verbatim, by the vertex lemma (c)
 & the corresponding statement of the sourced format \eqref{10.4:eq-sourced-format-b}, to rest
 on Lemma~\ref{10.4:lem-off-diagonal-vertex}(c) $(\ast)$ \\
\hline
combinatorial proposition & sup-norms, combinatorics free of $z$ & verbatim & verbatim \\
\hline
read-out & Schwarz and Cauchy estimates, connectivity
 & the torus bound of the vector package, in the parameter $t$ & same \\
\hline
\end{tabular}
\end{center}

The entries that need comment beyond a reference are as follows.

\emph{The weight.} In the channel $\mathbf{3}$ the source is the plaquette weight
$\mathsf{c}_{\mathsf{k}}$ of the tensor weight $\mathsf{k}$. The weighted action and the initial
density of the correlation theorem are therefore read with the weight
\begin{equation}\label{10.4:eq-tensor-slot-transfer-1}
c = x_0 - w - t\,\mathsf{c}_{\mathsf{k}},
\end{equation}
linearly, as in the correlation theorem. No new slot is created. In the channel $\mathbf{6}$ the
level-zero vertex is read directly: the level-zero bare argument of the vector package is applied,
with Lemma~\ref{10.4:lem-off-diagonal-vertex}(c) in place of the scalar vertex lemma.

\emph{The gauge-fixing lemma.} It reads a single coefficient, the one on $G+A$, and the
criticality of $U_{k+1}$. It is not used in either channel, because the vertex is a potential, as
established in the typed complex-factor proposition.
The gauge-fixing coefficient $a_j$ does not depend on $t$.

\emph{The fifth line of the decomposition.} The vertex line of the decomposition of the units is
present and contributes. For each cube $\square$, with localising function $\hat\chi_\square$, put
\begin{equation}\label{10.4:eq-tensor-slot-transfer-2}
f_\square := \mathsf{A}^{r}(\tau_k\hat\chi_\square,\cdot),
\qquad
\mathfrak{N}(f_\square) \le C_S'(2M)^4\,\sup|\tau_k|\,\mathfrak{a}^2 .
\end{equation}
The first line of the decomposition, the groups, is read by the typed irrelevance lemma
(Lemma~\ref{10.4:lem-typed-irrelevance}). In the derivative lemma the counterterm piece
$\beta^{r}[F]\,\mathsf{A}^{r}(h_y,\cdot)$ of the typed irrelevant group is treated singly. For
$r=\mathbf{3}$ it is read by part (b) of the vertex lemma of the correlation theorem, and for
$r=\mathbf{6}$ by Lemma~\ref{10.4:lem-off-diagonal-vertex}(b).

\emph{The identification $P^{\flat}\simeq\mathsf{P}$.} In the correlation theorem this
identification rests on the fact that the two integrals have one integrand. Here it is used through
the derivative of that identity in $t$. This is legitimate because the vertex is a potential in
both integrals.

\emph{The read-out.} In both channels the read-out of the correlation theorem, by the Schwarz and
Cauchy estimates and connectivity, is replaced by the torus bound of the vector package, in the
parameter $t$.

$(\ast)$ \emph{The weight-linear bound and the third item of the Wilson-shift bound for
$r=\mathbf{6}$.} In the correlation theorem the weight-linear bound states that the terms are linear
in the weight and bounded by weighted absolute sums. The third item of the Wilson-shift bound states
that the raw Wilson shift is at most $\theta_W x$, and it follows from part (c) of the vertex lemma.
For $r=\mathbf{3}$ both apply verbatim, because the source is then a plaquette weight. For
$r=\mathbf{6}$ both are replaced by the corresponding statements of the sourced format
\eqref{10.4:eq-sourced-format-b}. The counterpart of the Wilson-shift bound is to rest
on Lemma~\ref{10.4:lem-off-diagonal-vertex}(c). These two statements are not proved here; the proof
is incomplete at this step.

\emph{The list of places.} The places at which the correlation theorem reads clause (f3) are nine: the
items (B8), (B9) of the source-reading lemma; one place in the complex-factor proposition; the
sentence on $W_4$ of the irrelevance lemma; one place in the group formed by the flat form of the
summation lemma, the lemma on the complex potential, the homogeneity condition and the flat section
of the correlation theorem; two places in the identification $P^{\flat}\simeq\mathsf{P}$ and its
proposition; one place in the combinatorial proposition; and two further places of the correlation
theorem. Compare each with the typed format statement of the vector package. At six of them (the
place in the flat group, the two places of the identification, the place in the combinatorial
proposition and the two further places) the reading agrees with that statement. Two entries change,
and only two; the first of them covers two of the nine places.
\begin{itemize}
\item At the items (B8), (B9) of the source-reading lemma, and at the reading of the same clause in
the complex-factor proposition, the pure channel contributes, whereas in the vector package it does
not.
\item At the sentence on $W_4$ of the irrelevance lemma, the vector package has $\Pi=0$, whereas here
$\Pi=\beta^{r}\mathsf{T}^{r}$. This is the change absorbed by
Lemma~\ref{10.4:lem-typed-irrelevance}.
\end{itemize}
\end{proof}

\emph{Vertex potentials.}
We work on the whole lattice, at the step from level $k$ to level $k+1$, with the scaffold free of the
source variables $z$. For the localised counterterms
\begin{equation*}
f_\square := \mathsf{A}^{r}(\hat\chi_\square,\cdot\,;h_0)
\end{equation*}
let $\mathcal{S}^{r}(Y,\mathbf{U},\mathbf{J},\mathsf{A};h_0)$, $Y\in\mathcal{D}_k$, be the pieces that
the localisation lemma of the step produces for them; $\mathcal{D}_k$ is the set of domains $Y$ of these
pieces. These pieces are the vertex potentials. They are linear in $h_0$, analytic on the domain of the
analyticity lemma of the step, localised in $Y$, gauge invariant and zero at $\mathsf{A}=0$. They obey
\begin{equation*}
\begin{split}
\sum_{Y}\mathcal{S}^{r}(Y,U_{k+1},J_{k+1},\mathsf{A};h_0)
&= \mathsf{A}^{r}\bigl(1,U_k(e^{iB'}V^{(k)});h_0\bigr)-\mathsf{A}^{r}(1,U_{k+1};h_0),\\
B' &= C\mathsf{A}-h\tilde D(C\mathsf{A}),
\end{split}
\end{equation*}
and the bound
\begin{equation*}
|\mathcal{S}^{r}(Y)|\le C_{\mathcal{S}}\,|h_0|\,e^{-(1-2\delta)\kappa d_k(Y)}.
\end{equation*}
The bound follows from the localisation lemma, whose input bounds on the localised counterterms are the
bounds for the traceless-diagonal vertex and the bounds for the off-diagonal tensor vertex
(Lemma~\ref{10.4:lem-off-diagonal-vertex}(b)), each taken at the absolute radii
$\alpha=(1+\beta_{\mathrm{pr}})\alpha_0$ and $a\le C\varepsilon_1$.

\emph{The ceiling on $u_0$.}
The constant $\varepsilon_1$ is one of the constants of the analysis. It is fixed, in the same way as
the ceiling imposed for complex sources $w$, so that the potentials of the proposition on a complex
factor near one on the Gaussian use at most half of the constants of the format conditions
\cite[(1.36), (1.42), (1.43)]{B-CMP116} of the cluster-expansion lemma for one step. Then
\begin{equation}\label{10.4:eq-vertex-response-a}
\begin{gathered}
\text{$u_0>0$ is the largest number such that $u_0\,C_{\mathcal{S}}$ does not exceed}\\
\text{the other half of each of these constants.}
\end{gathered}
\end{equation}
Condition \eqref{10.4:eq-vertex-response-a} is one-sided: it is a ceiling on $u_0C_{\mathcal{S}}$.
Once $\varepsilon_1$ and $u_0$ have been fixed in this way, the quantities $\mathsf{r}_k$, $\tau_k$,
$e_k$ and $\gamma$ are re-evaluated.

Throughout, $D(a,\rho)$ is the open disc in $\mathbb{C}$ of centre $a$ and radius $\rho$.
Here $\tau_k$ denotes the radius ratio of the correlation theorem; the estimates below use it only
through the inequalities $g_k/\mathsf{r}_k\le\tau_k<1$. It is not
the running tensor coefficient entering \eqref{10.4:eq-insertion-correlation-a}.
The index $r$ labels the channel of the tensor vertex, $\mathsf{V}_r$ being the corresponding
$W_4$-module (Definition~\ref{10.4:def-typed-families}), and $T=T_k$ is the bound entering the
characteristic functions $\chi(|B(b)|<T)$ below.

\begin{proposition}[Response of one step to the tensor vertex; Stated; proof incomplete at the step
marked $(\ast)$]\label{10.4:prop-vertex-response}
Let $k\ge 0$ and $T=T_k$. Let
\begin{equation*}
(\mathfrak{q},s,u)\in D(0,\mathsf{r}_k)\times D(1,\sigma_0)\times D(0,u_0),
\end{equation*}
let $h\in\mathbb{B}^w_k$, and let $h_0\in(\mathsf{V}_r)_{\mathbb{C}}$ with $|h_0|\le 1$.
Let $\Phi_{k+1}(\mathfrak{q},s,u;h)$ be the family that the
cluster-expansion lemma for one step produces for the potentials
\begin{equation*}
\begin{split}
V(Y) := {}& s\cdot\tfrac12\langle B,\mathcal{Q}^\Psi(Y,\mathfrak{q}B)B\rangle
+\mathcal{V}^\Phi(Y,\mathfrak{q}B)\\
&+\mathcal{V}'(Y,\mathfrak{q}B)[h]+u\,\mathcal{S}^{r}(Y,\mathfrak{q}B;h_0).
\end{split}
\end{equation*}
Then the following hold.
\begin{enumerate}
\item[(a)] $\|\Phi_{k+1}\|_{k+1}\le E^{*}$. The family is holomorphic in $(\mathfrak{q},s,u)$ and along
holomorphic families in $\mathbb{B}^w_k$. At $\mathfrak{q}=g_k$ its total is the integral of clause (d) of
the proposition on a complex factor near one on the Gaussian, with the term
\begin{equation*}
+u\bigl[\mathsf{A}^{r}\bigl(1,U_k(e^{iB'}V^{(k)});h_0\bigr)-\mathsf{A}^{r}(1,U_{k+1};h_0)\bigr]
\end{equation*}
added in the exponent.
\item[(b)] For every $X$, and for each term of the expansion \cite[(2.1)--(2.14)]{B-CMP116} defining
$\Phi_{k+1}(X)$,
\begin{equation*}
\Phi_{k+1}(X;-\mathfrak{q},s,u;h)=\Phi_{k+1}(X;\mathfrak{q},s,u;h).
\end{equation*}
Moreover, $\Phi_{k+1}(0,s,u;h)$ does not depend on $(u,h)$.
\item[(c)] Define
\begin{equation*}
\delta^{r}\Phi_{k+1}(\mathfrak{q},s;h):=\partial_u\Phi_{k+1}\big|_{u=0}.
\end{equation*}
This is a $\mathsf{V}_r$-typed family in the sense of Definition~\ref{10.4:def-typed-families}, linear
in $h_0$. On the polydisc $D(0,\mathsf{r}_k)\times D(1,\sigma_0)$ it satisfies
$\|\delta^{r}\Phi_{k+1}\|\le 2E^{*}/u_0$; it is even in $\mathfrak{q}$; and it is zero at
$\mathfrak{q}=0$. Hence, for all $s\in D(1,\sigma_0)$ and $h\in\mathbb{B}^w_k$,
\begin{equation*}
\|\delta^{r}\Phi_{k+1}(g_k,s;h)\|_{k+1}\le c^{\otimes}_k:=\frac{2E^{*}}{u_0}\,\tau_k^2 .
\end{equation*}
\item[(d)] Define
\begin{equation*}
\mathsf{L}^{r}_{k+1}(s;h):=\tfrac12\,\partial_{\mathfrak{q}}^2\,
\delta^{r}\Phi_{k+1}(0,s;h).
\end{equation*}
This is affine in $h$, and $\|\mathsf{L}^{r}_{k+1}\|\le (2E^{*}/u_0)\,\mathsf{r}_k^{-2}$. Moreover,
\begin{align*}
\bigl\|\delta^{r}\Phi_{k+1}(g_k,s;h)-g_k^2\,\mathsf{L}^{r}_{k+1}(s;h)\bigr\|
&\le \frac{2E^{*}}{u_0}\,\frac{\tau_k^4}{1-\tau_k^2},\\
\bigl\|\mathsf{L}^{r}_{k+1}(s;h)-\mathsf{L}^{r}_{k+1}(1;h)\bigr\|
&\le \frac{4E^{*}}{u_0}\,\mathsf{r}_k^{-2}\,\frac{|s-1|}{\sigma_0},\\
\bigl\|\mathsf{L}^{r}_{k+1}(s;\tilde h)-\mathsf{L}^{r}_{k+1}(s;h)\bigr\|
&\le \frac{2E^{*}}{u_0}\,\mathsf{r}_k^{-2}\,\frac{|\tilde h-h|_k}{2\mathsf{E}}.
\end{align*}
\item[(e)] Write the expansions in $\mathfrak{q}$ as
\begin{equation*}
\mathcal{S}^{r}(\mathfrak{q}B)=\mathfrak{q}\,\mathcal{S}_1(B)+\mathfrak{q}^2\mathcal{S}_2(B)+\dots,
\qquad V(\mathfrak{q}B)\big|_{u=0}=\mathfrak{q}\,V_1(B)+\dots .
\end{equation*}
Let $\langle\cdot\rangle_0$ be the expectation with respect to the normalised measure
$d\mu_{C^{(k)}}\prod\chi(|B|<T_k)$, and let
$\langle F;G\rangle_0:=\langle FG\rangle_0-\langle F\rangle_0\langle G\rangle_0$. Then the total of
$\mathsf{L}^{r}_{k+1}(1;h)$ is
\begin{equation*}
\langle\mathcal{S}_2\rangle_0+\langle\mathcal{S}_1;V_1\rangle_0 .
\end{equation*}
This is the one-loop response.
\end{enumerate}
\end{proposition}

\begin{proof}
\emph{Clause (a).} The cluster-expansion lemma for one step uses, of the potentials it is given, nothing
beyond the format conditions \cite[(1.36), (1.42), (1.43)]{B-CMP116} and the activity bound.

By the choice of $\varepsilon_1$ in the ceiling, the potentials of the proposition on a complex factor
near one on the Gaussian use at most half of the format constants. The vertex term contributes
$u\,\mathcal{S}^{r}$. By the bound on $\mathcal{S}^{r}(Y)$, together with $|h_0|\le 1$ and $|u|<u_0$, this
term is bounded by $u_0C_{\mathcal{S}}e^{-(1-2\delta)\kappa d_k(Y)}$. By
\eqref{10.4:eq-vertex-response-a}, $u_0C_{\mathcal{S}}$ is at most the other half of the format
constants. Hence $V$ satisfies the hypotheses of the cluster-expansion lemma on the whole polydisc and
for all $h\in\mathbb{B}^w_k$. The lemma then gives $\|\Phi_{k+1}\|_{k+1}\le E^{*}$, with the stated
holomorphy.

At $\mathfrak{q}=g_k$, the summation identity for the vertex potentials displayed above shows that the
vertex term adds exactly
\begin{equation*}
u\bigl[\mathsf{A}^{r}\bigl(1,U_k(e^{iB'}V^{(k)});h_0\bigr)-\mathsf{A}^{r}(1,U_{k+1};h_0)\bigr]
\end{equation*}
to the exponent of the integral of clause (d) of that proposition.

\emph{Clause (b)} $(\ast)$. Every potential is invariant under
$(\mathfrak{q},B)\mapsto(-\mathfrak{q},-B)$. Indeed, $\mathcal{V}^\Phi$, $\mathcal{V}'[h]$ and
$\mathcal{S}^{r}$ are functions of $\mathfrak{q}B$. The remaining potential,
$\langle B,\mathcal{Q}^\Psi(\mathfrak{q}B)B\rangle$, is quadratic in $B$ and otherwise a function of
$\mathfrak{q}B$.

In the expansion \cite[(2.1)--(2.14)]{B-CMP116}, in the form in which the cluster-expansion lemma for
one step uses it, a term is an integral over two fields $(B,X')$; here $X'$ is the second integration
field, denoted $X$ in \cite{B-CMP116}, and it is not a polymer. The integrand is built from:
\begin{itemize}
\item the potentials $V$;
\item the characteristic functions $\chi(|B(b)|<T)$ and their complements $\chi^{c}$;
\item Gaussian measures whose covariances are built from $\Delta_k$ and $C^{(k)}(Z_0,\sigma)$;
\item the factor $\exp\bigl(-s^{1/2}\langle B,\Gamma_k(Z_0,\sigma)X'\rangle\bigr)$.
\end{itemize}
The letters $\Delta_k$, $C^{(k)}(Z_0,\sigma)$, $\Gamma_k(Z_0,\sigma)$ and, below, $H(Z)$ are those of
\cite[(2.1)--(2.14)]{B-CMP116}.
Each of these is invariant under $(B,X')\mapsto(-B,-X')$, once $\mathfrak{q}$ is replaced by
$-\mathfrak{q}$ in the potentials. The change of variables $(B,X')\mapsto(-B,-X')$ therefore gives the
equality of clause (b) for each term of \cite[(2.14)]{B-CMP116}. The equality then holds for the
functions $H(Z)$ of that expansion and for $\Phi_{k+1}(X)$, since these are finite sums of such terms
and convergent expansions of those sums.
$(\ast)$ The verification, on each term of \cite[(2.1)--(2.14)]{B-CMP116}, that each of its factors
has this invariance is not carried out here.

At $\mathfrak{q}=0$ all potentials vanish. Hence $V=0$ at $\mathfrak{q}=0$ whatever $(u,h)$, and
$\Phi_{k+1}(0,s,u;h)$ does not depend on $(u,h)$.

\emph{Clause (c).} Fix $X$ and all remaining arguments of the family, and let
\begin{equation*}
F(u) := \Phi_{k+1}(X;\mathfrak{q},s,u;h)-\Phi_{k+1}(X;\mathfrak{q},s,0;h).
\end{equation*}
By (a), $F$ is analytic on $D(0,u_0)$. After multiplication by the weight of the norm, $F$ is bounded
by $2E^{*}$, and $F(0)=0$. The Cauchy estimate at the centre of $D(0,u_0)$ gives
$|\partial_uF(0)|\le 2E^{*}/u_0$, again after multiplication by the weight. Taking the supremum,
$\|\delta^{r}\Phi_{k+1}\|\le 2E^{*}/u_0$ on the polydisc.

The equality in (b) holds for all $u\in D(0,u_0)$. Differentiating it in $u$ at $u=0$ shows that
$\delta^{r}\Phi_{k+1}$ is even in $\mathfrak{q}$. Since $\Phi_{k+1}(0,s,u;h)$ does not depend on $u$,
$\delta^{r}\Phi_{k+1}(0,s;h)=0$.

We now use an elementary bound. Let $f$ be even and analytic on $D(0,R)$, with $|f|\le M$ and
$f(0)=0$. Since $f$ is even, its Taylor series contains only even powers. Hence $G(v):=f(\sqrt v)$ is
well defined and analytic on $D(0,R^2)$, with $|G|\le M$ and $G(0)=0$. The Schwarz lemma gives
$|G(v)|\le M|v|/R^2$, that is,
\begin{equation*}
|f(z)|\le M\,\frac{|z|^2}{R^2}\qquad (z\in D(0,R)).
\end{equation*}
Apply this with $f=\delta^{r}\Phi_{k+1}(X;\cdot,s;h)$, pointwise and after multiplication by the weight,
$R=\mathsf{r}_k$, $M=2E^{*}/u_0$ and $z=g_k$, which lies in $D(0,\mathsf{r}_k)$ because
$g_k/\mathsf{r}_k\le\tau_k<1$. This gives
\begin{equation*}
\|\delta^{r}\Phi_{k+1}(g_k,s;h)\|_{k+1}\le\frac{2E^{*}}{u_0}\,\frac{g_k^2}{\mathsf{r}_k^2}
\le\frac{2E^{*}}{u_0}\,\tau_k^2=c^{\otimes}_k .
\end{equation*}
The family $\delta^{r}\Phi_{k+1}$ is $\mathsf{V}_r$-typed by the equivariance statement for the step:
the step's production of the family from its potentials is equivariant under $W_4$, and
$\mathcal{S}^{r}$ is typed.

\emph{Clause (d).} Keep the notation $f$, $G$, $M$, $R$ of clause (c), and write
$G(v)=\sum_{n\ge1}G_nv^n$. The Cauchy estimates on $D(0,R^2)$ give $|G_n|\le M/R^{2n}$. With
$\mathsf{u}:=|z|^2/R^2$ we therefore have
\begin{equation*}
\bigl|G(v)-G'(0)v\bigr|=\Bigl|\sum_{n\ge2}G_nv^n\Bigr|\le M\sum_{n\ge2}\mathsf{u}^n
=M\,\frac{\mathsf{u}^2}{1-\mathsf{u}},\qquad v=z^2 .
\end{equation*}
Also $G'(0)=\tfrac12 f''(0)$ and $|G'(0)|\le M/R^2$.

With $f$ as in clause (c), $\tfrac12 f''(0)$ is $\mathsf{L}^{r}_{k+1}(s;h)$ at $X$, so
$\|\mathsf{L}^{r}_{k+1}\|\le(2E^{*}/u_0)\mathsf{r}_k^{-2}$. At $z=g_k$ we have
$\mathsf{u}=g_k^2/\mathsf{r}_k^2\le\tau_k^2<1$, and $G'(0)\,v=g_k^2\,\mathsf{L}^{r}_{k+1}(s;h)$. Since
$\mathsf{u}\mapsto\mathsf{u}^2/(1-\mathsf{u})$ is increasing on $[0,1)$, this gives the second bound
of (d).

For the dependence on $s$, let $\mathcal{M}:=(2E^{*}/u_0)\mathsf{r}_k^{-2}$ and consider
$s\mapsto\mathsf{L}^{r}_{k+1}(s;h)-\mathsf{L}^{r}_{k+1}(1;h)$. This function is analytic on
$D(1,\sigma_0)$ and, after multiplication by the weight, bounded by $2\mathcal{M}$. It vanishes at
$s=1$. The Schwarz lemma on $D(1,\sigma_0)$ bounds it by $2\mathcal{M}|s-1|/\sigma_0$, which is the
third bound of (d).

For the dependence on $h$ we argue as in the Lipschitz lemma in the history variable. The bound
$\mathcal{M}$ holds for histories in the ball of radius $4\mathsf{E}$. Take $h,\tilde h$ of norm at most
$2\mathsf{E}$ and put $t\mapsto\mathsf{L}^{r}_{k+1}(s;h+t(\tilde h-h))$. This function is affine in $t$
and bounded by $\mathcal{M}$ for $|t|\le 2\mathsf{E}/|\tilde h-h|_k$. Cauchy's estimate for its
derivative gives the fourth bound of (d).

It remains to show that $\mathsf{L}^{r}_{k+1}$ is affine in $h$. The derivative
$\partial_u\partial^2_{\mathfrak{q}}$ at $\mathfrak{q}=0$ of the logarithm of the integral of the
exponential of the potentials is the sum of two terms: a mean of $\mathcal{S}_2$, and a covariance of
$\mathcal{S}_1$ with the first-order part of $V$. The history $h$ enters this first-order part linearly,
through $\mathcal{V}'(Y,\mathfrak{q}B)[h]$.

\emph{Clause (e).} Use the same expansion at $s=1$. Put
$\mathcal{S}:=\sum_Y\mathcal{S}^{r}(Y,\mathfrak{q}B;h_0)$ and $V^0:=\sum_YV(Y)|_{u=0}$. The derivative in
$u$ at $u=0$ of the logarithm of $\langle\exp\sum_YV(Y)\rangle_0$ is
\begin{equation*}
\frac{\langle\mathcal{S}\,e^{V^0}\rangle_0}{\langle e^{V^0}\rangle_0}.
\end{equation*}
Insert $\mathcal{S}=\mathfrak{q}\mathcal{S}_1+\mathfrak{q}^2\mathcal{S}_2+\dots$ and
$e^{V^0}=1+\mathfrak{q}V_1+O(\mathfrak{q}^2)$.

The coefficient of $\mathfrak{q}$ is $\langle\mathcal{S}_1\rangle_0$. It vanishes by parity: $\mathcal{S}_1$
is odd in $B$, because $\mathcal{S}^{r}$ is a function of $\mathfrak{q}B$, while
$d\mu_{C^{(k)}}\prod\chi(|B|<T_k)$ is even.

The coefficient of $\mathfrak{q}^2$ is
\begin{equation*}
\langle\mathcal{S}_2\rangle_0+\langle\mathcal{S}_1V_1\rangle_0
-\langle\mathcal{S}_1\rangle_0\langle V_1\rangle_0
=\langle\mathcal{S}_2\rangle_0+\langle\mathcal{S}_1;V_1\rangle_0 .
\end{equation*}
This coefficient is $\tfrac12\partial_{\mathfrak{q}}^2$ at $\mathfrak{q}=0$ of the derivative above. It is
therefore the total of $\mathsf{L}^{r}_{k+1}(1;h)$.
\end{proof}

\begin{remark}
The proof of the proposition on a complex factor near one on the Gaussian tolerates an anisotropic
quadratic form. By the device of the remark that follows that proposition, one carries
$\tfrac12(1-s)\langle B,\Delta^{(k)}B\rangle$ as one more form of the kind admitted in
\cite[(1.42)]{B-CMP116}. The tensor vertex is a potential of this kind.

The small bound of clause (c) of that proposition, $\|\Phi(g_k,s;h)\|\le e_k$ uniformly in $s$, does
not carry over. It rests on the fact that the one-loop term is free of $s$, and the one-loop term is
not free of a traceless deformation.

At first order, clauses (c) and (d) of Proposition~\ref{10.4:prop-vertex-response} take its place.
The response is $g_k^2$ times one family, which sees the coupling only through the cutoff
$\prod\chi(|B|<T_k)$, up to an error of order $\tau_k^4$.

For $r=\mathbf{1}$ the same computation gives
$\langle\mathcal{S}_2\rangle_0+\langle\mathcal{S}_1;V_1\rangle_0$, which by the lemma on the
$s$-independence of the one-loop term is a field-free constant up to terms produced by the cutoff
$\prod\chi(|B|<T_k)$. That is, $\mathsf{m}_{\mathbf{1}}=0$, which is one of the statements of the
correlation theorem.
\end{remark}

\begin{remark}[Evenness of the scalar step in the coupling]\label{10.4:rem-scalar-evenness}
Stated; the proof below is incomplete at the step marked $(\ast)$.
Let $k \ge 0$, and let $\Phi_{k+1}(\mathfrak{q},s,u;h)$ be the family of
Proposition~\ref{10.4:prop-vertex-response}, defined on its polydisc
$(\mathfrak{q},s,u) \in D(0,\mathsf{r}_k)\times D(1,\sigma_0)\times D(0,u_0)$, where $D(a,r)$
denotes the disc of centre $a$ and radius $r$, for
$h \in \mathbb{B}^w_k$, $|h_0| \le 1$. At $u = 0$ it is the family
$\Phi_{k+1}(\mathfrak{q},s;h)$ of the one-step proposition of the correlation theorem, and this
family is even in $\mathfrak{q}$: for every argument $X$ of the family,
\begin{equation}\label{10.4:eq-scalar-evenness-1}
  \Phi_{k+1}(X;-\mathfrak{q},s;h) = \Phi_{k+1}(X;\mathfrak{q},s;h),
  \qquad \mathfrak{q} \in D(0,\mathsf{r}_k),\ s \in D(1,\sigma_0),\ h \in \mathbb{B}^w_k .
\end{equation}
Consequently clause~(c) of the one-step proposition of the correlation theorem holds with
$E^{*}\tau_k^2+\vartheta_k$ in place of the constant $e_k = 2E^{*}\tau_k+\vartheta_k$ of that
clause, $\vartheta_k$ denoting the second summand of $e_k$, which is unchanged; and the
pure-kernel constant $a_k$ of the transport of the slot weights may be taken with $\tau_k^2$
in place of $\tau_k$. Here $\tau_k$ is the radius ratio of the correlation theorem that enters
clause~(c) of Proposition~\ref{10.4:prop-vertex-response}, not the running tensor coefficient.
\end{remark}

\begin{proof}
The disc $D(0,\mathsf{r}_k)$ is invariant under $\mathfrak{q} \mapsto -\mathfrak{q}$, so both
sides of \eqref{10.4:eq-scalar-evenness-1} are defined. The point $u = 0$ lies in $D(0,u_0)$.
$(\ast)$ By clause~(b) of
Proposition~\ref{10.4:prop-vertex-response}, whose proof is incomplete at one step, the identity
\begin{equation*}
  \Phi_{k+1}(X;-\mathfrak{q},s,u;h) = \Phi_{k+1}(X;\mathfrak{q},s,u;h)
\end{equation*}
holds term by term for all $(\mathfrak{q},s,u)$ in the polydisc and all $h \in \mathbb{B}^w_k$.
Taking $u = 0$ gives \eqref{10.4:eq-scalar-evenness-1}.
\end{proof}

\begin{lemma}[The fluctuation kernel for typed slots]\label{10.4:lem-typed-slot-kernel}
Let $k$ be a level and let $\delta H^{(j)}$, $j\le k$, be arbitrary $\mathsf{V}_r$-typed
E-families on $\mathcal{U}_j$ in the sense of Definition~\ref{10.4:def-typed-families}. Write
$\delta h$ for this collection, $|\delta h|_k$ for its norm in the weighted form of the
fluctuation-kernel lemma of the correlation theorem, and $h_0\in(\mathsf{V}_r)_{\mathbb{C}}$.
The arguments $\mathbf{U},\mathbf{J},\mathsf{A}$, the fields $B'$ and $V^{(k)}$, and the
constants $C_w,\varepsilon_1,C_1,C_2,\bar q,\delta,\kappa,\kappa_1$ are those of the
fluctuation-kernel lemma of the correlation theorem in its weighted form.
Then there are functions
\begin{equation*}
\mathcal{V}'(Y,\mathbf{U},\mathbf{J},\mathsf{A})[\delta h;h_0],
\end{equation*}
one for each region $Y$, with the following properties. They are linear in $\delta h$ and in
$h_0$. They have the analyticity, locality and invariance properties stated in the
fluctuation-kernel lemma. They vanish at $\mathsf{A}=0$. Finally, they satisfy
\begin{equation}\label{10.4:eq-typed-slot-kernel-a}
\begin{gathered}
\sum_Y\mathcal{V}'(Y)
=\delta h^{\mathrm{ren},r}\bigl(U_k(e^{iB'}V^{(k)})\bigr)-\delta h^{\mathrm{ren},r}(U_{k+1}),\\
|\mathcal{V}'(Y)|\le C_w\,\varepsilon_1\,C_1\,M^{\bar q}\,e^{C_2\kappa_1}\,|\delta h|_k\,|h_0|\,
e^{-(1-2\delta)\kappa d_k(Y)}.
\end{gathered}
\end{equation}
\end{lemma}

Stated; the proof below is incomplete at the step marked $(\ast)$.

\begin{proof}
We follow the proof of part~(ii) of the fluctuation-kernel lemma of the correlation theorem.
That proof runs through nine links, which we label (B1)--(B9) in the order in which that proof
takes them, and for each link we say what it uses of the
families and why this is available for typed families.

\emph{Links (B1)--(B7).} These links use three things only: the properties (f1) and (f2) of
the families and the norm of the families. By Definition~\ref{10.4:def-typed-families}, a
$\mathsf{V}_r$-typed E-family has properties (f1) and (f2). It has the typed invariance
$(\mathrm{f3})^{\mathsf{V}}$ in place of (f3), and it need not satisfy condition (S). The norm
used is the one fixed in that definition. Hence (B1)--(B7) hold for the typed families
$\delta H^{(j)}$ word for word. The slot weights $\varpi_{k-j}$ entering $|\delta h|_k$ are
those of the weighted form of the fluctuation-kernel lemma.

\emph{Step $(\ast)$: links (B8) and (B9).}

Link (B8) rests on \cite[(4.35)--(4.37)]{B-CMP109} and \cite[\S5]{B-CMP109}. It uses the
invariance (f3) at one place only, namely \cite[(5.12)--(5.15)]{B-CMP109}. Its purpose there
is to reduce the Taylor coefficient of order two in the momentum variable $p$ of
\cite[(5.12)--(5.15)]{B-CMP109} of the tensor $\hat\Pi$ of \cite[\S5]{B-CMP109}, as a
function of $p$, to a single number.

For typed families this reduction is supplied by part~(c) of the hypercubic decomposition of
symmetric tensors (Lemma~\ref{10.4:lem-hypercubic-decomposition}). The space
$\operatorname{Hom}_{W_4}(\mathsf{V}_r,\operatorname{Sym}^2\Lambda^2)$ is one-dimensional and
is spanned by $\mathsf{T}^r$. The part of the order-$p^2$ coefficient that is symmetric under
exchange of the two pairs is $W_4$-equivariant in $h_0$. It is therefore a multiple of
$\mathsf{T}^r(h_0)$, and the factor is the number $\beta^r[\,\cdot\,]$ of
Definition~\ref{10.4:def-typed-families}.

Terms odd in $p$ do not occur. The reason is that the element $-1$ of $W_4$ acts trivially on
$\mathsf{V}_r$, so the typed coefficient is invariant under $p\mapsto-p$. The relations
\cite[(5.30)--(5.42)]{B-CMP109} are linear equations between Taylor coefficients, and they are
used with the typed coefficient in place of the untyped one.

In link (B9) the counterterm is the untyped marginal coefficient $\beta_j$ of the
fluctuation-kernel lemma at scale $j$ times explicit expansions. In the typed case it is
$\beta^r$ times the expansions of $\mathsf{A}^r\circ U_j$. By part~(iii) of
Lemma~\ref{10.4:lem-coefficient-bounds} we have
\begin{equation*}
\beta^r[\mathsf{A}^{r'}\circ U_j]=\delta_{rr'} .
\end{equation*}
Hence the order-$p^2$ symbol of these expansions is that of $\mathsf{T}^r$.

The two groups of terms, the order-$p^2$ part coming from (B8) and the counterterm of (B9),
then cancel. This follows from the definition
\begin{equation*}
F^{\mathrm{ren},r}=F-F(1)-\beta^r[F]\,\mathsf{A}^r
\end{equation*}
in Definition~\ref{10.4:def-typed-families}, applied to $F=\delta H^{(j)}$. What is not
carried out at $(\ast)$ is the verification that each of the linear relations
\cite[(5.30)--(5.42)]{B-CMP109} between Taylor coefficients holds with the typed coefficient
$\beta^r[\delta H^{(j)}]\,\mathsf{T}^r(h_0)$ in place of the scalar one.

\emph{Pointed typed families.} For pointed typed families the bound in
\eqref{10.4:eq-typed-slot-kernel-a} is obtained by composing two statements, each applied as
stated. The first is the irrelevance after one typed counterterm
(Lemma~\ref{10.4:lem-typed-irrelevance}). The second is the pointed-family lemma of the
correlation theorem.
\end{proof}

\begin{corollary}[Tangent of the pure recursion in a typed direction]
\label{10.4:cor-pure-recursion-tangent}
Let $k\ge 0$. Assume the following.
\begin{itemize}
\item[(a)] The typed-slot kernel $\mathcal{V}'$ of \eqref{10.4:eq-typed-slot-kernel-a} exists with the
properties stated there. This is the conclusion of Lemma~\ref{10.4:lem-typed-slot-kernel} on the
fluctuation kernel for typed slots, whose proof is incomplete at one step.
\item[(b)] The pure recursion holds, together with clause~(c) of the theorem bounding the pure
recursion and the term-by-term lemma used in the proof of that clause.
\item[(c)] Clauses~(a) and~(c) of the response of one step to the tensor vertex,
Proposition~\ref{10.4:prop-vertex-response}, hold; the proof of that proposition is
incomplete at one step.
\end{itemize}
Then the tangent of the pure recursion in a typed direction $\delta\mathsf{P}^{(\le k)}$ is
\begin{equation}\label{10.4:eq-pure-recursion-tangent-1}
\begin{split}
\delta\mathsf{P}^{(k+1)}
&=\tau_k\,\delta^{r}\Phi_{k+1}\bigl(g_k,s_k;\mathsf{P}^{(\le k)}_\xi\bigr)
+D_h\Phi_{k+1}\bigl(g_k,s_k;\mathsf{P}^{(\le k)}_\xi\bigr)\bigl[\delta\mathsf{P}^{(\le k)}\bigr],\\
\tau^{\mathrm{pure}}_{k+1}&=\tau_k+\beta^{r}\bigl[\delta\mathsf{P}^{(k+1)}\bigr];
\end{split}
\end{equation}
here $D_h$ denotes the derivative of $\Phi_{k+1}$ in its history argument $h$.
The quantities $L^{(k+1)}$, $s_k$ and $\xi_k$ do not depend on the parameter along which
the typed direction is taken. Write $p_j:=\|\delta\mathsf{P}^{(j)}\|_j$. Then
\begin{equation}\label{10.4:eq-pure-recursion-tangent-a}
p_{k+1}\le a_k\sum_{j\le k}\varpi_{k-j}\,p_j+c^{\otimes}_k\,|\tau_k| .
\end{equation}
\end{corollary}

\begin{proof}
The map $\Phi_{k+1}(\mathfrak{q},s,u;h)$ is holomorphic in $(\mathfrak{q},s,u)$ and along
holomorphic families in $\mathbb{B}^w_k$. This is part~(a) of
Proposition~\ref{10.4:prop-vertex-response}, assumed in hypothesis~(c); the proof of that
proposition is incomplete at one step.

We first prove \eqref{10.4:eq-pure-recursion-tangent-1}. The pure recursion gives
$\mathsf{P}^{(k+1)}$ as the value of the step's family $\Phi_{k+1}$ at
$(\mathfrak{q},s,u;h)=(g_k,s_k,0;\mathsf{P}^{(\le k)}_\xi)$. We differentiate this along the typed
direction and apply the chain rule for holomorphic maps. The derivative has two contributions.

The first comes from the vertex slot $u$, in which the typed direction enters with coefficient
$\tau_k$. By part~(c) of Proposition~\ref{10.4:prop-vertex-response}, the derivative of the
step's family in this slot is
\begin{equation*}
\delta^{r}\Phi_{k+1}(\mathfrak{q},s;h)=\partial_u\Phi_{k+1}\big|_{u=0}.
\end{equation*}
This yields the term $\tau_k\,\delta^{r}\Phi_{k+1}$.

The second comes from the history argument $h$. It yields
$D_h\Phi_{k+1}[\delta\mathsf{P}^{(\le k)}]$.

The second line of \eqref{10.4:eq-pure-recursion-tangent-1} is the pure recursion's rule for the
running coefficient, read in the typed direction: the increment is the typed marginal coefficient
$\beta^{r}$ of the new family $\delta\mathsf{P}^{(k+1)}$.

Next, $L^{(k+1)}$ and $s_k$ do not depend on that parameter, because the typed direction leaves
the Gaussian integration of the step untouched. Moreover $\delta\xi=0$. Indeed, by parts~(b)
and~(c) of the hypercubic decomposition of symmetric tensors
(Lemma~\ref{10.4:lem-hypercubic-decomposition}),
a $\mathsf{V}_r$-typed family has no component along $\mathsf{V}_{\mathbf{1}}$. Hence $\xi_k$
does not depend on it either.

It remains to prove \eqref{10.4:eq-pure-recursion-tangent-a}. We bound the two terms of
\eqref{10.4:eq-pure-recursion-tangent-1} in $\|\cdot\|_{k+1}$ separately.

For the second term we repeat the proof of clause~(c) of the theorem bounding the pure recursion,
with the vertex potential $\mathcal{V}'[h+\lambda\,\delta h]$, $\lambda$ a complex parameter in a
small disc about $0$, in place of $\mathcal{V}'[h]$, and take the derivative at $\lambda=0$.
Here $\mathcal{V}'$ is the kernel of hypothesis~(a), \eqref{10.4:eq-typed-slot-kernel-a},
supplied by Lemma~\ref{10.4:lem-typed-slot-kernel}, whose proof is incomplete at one step. That
kernel is linear in $\delta h$, and it has the analyticity, locality and invariance required there.
In this repetition the term-by-term lemma of hypothesis~(b) is applied to each term of the
expansion, exactly as in that proof. The outcome is
\begin{equation*}
\bigl\|D_h\Phi_{k+1}\bigl(g_k,s_k;\mathsf{P}^{(\le k)}_\xi\bigr)\bigl[\delta\mathsf{P}^{(\le k)}\bigr]
\bigr\|_{k+1}
\le a_k\sum_{j\le k}\varpi_{k-j}\,p_j .
\end{equation*}

For the first term we use part~(c) of Proposition~\ref{10.4:prop-vertex-response}, whose proof is
incomplete at one step. It gives
\begin{equation*}
\|\delta^{r}\Phi_{k+1}(g_k,s;h)\|_{k+1}\le c^{\otimes}_k
\qquad\text{for all } s\in D(1,\sigma_0),\ h\in\mathbb{B}^w_k .
\end{equation*}
Applied at $s=s_k$ and $h=\mathsf{P}^{(\le k)}_\xi$, this bounds the first term by
$c^{\otimes}_k|\tau_k|$.

Adding the two bounds gives \eqref{10.4:eq-pure-recursion-tangent-a}.
\end{proof}

Throughout, $r\in\{\mathbf{3},\mathbf{6}\}$, and $\mathsf{k}$ is the tensor weight, with values in
$(\mathsf{V}_r)_{\mathbb{C}}$. The production $\mathfrak{s}$ is the production of the complete step.
For a direction $\delta$ and a parameter $\lambda$, the production in which $\delta$ is switched on
linearly is $\mathfrak{s}+\lambda\delta$; we call it the \emph{dialled production}, and $\lambda$ the
dial parameter.

A \emph{$\mathsf{V}_r$-typed direction with a marginal slot} is a sequence
$\delta=(\delta_j)_{0\le j\le k}$ with entries
\begin{equation*}
\delta_j=(\mu_j;\,\delta E^{(j)},\delta R^{(j)}).
\end{equation*}
Here $(\delta E^{(j)},\delta R^{(j)})_{0\le j\le k}$ is a typed direction, in the sense of the typed
directions with weight classes $\mathcal{H}_j(X)$, taken at $\mathsf{V}=\mathsf{V}_r$. The $\mu_j$ are
numbers which are analytic functions of $\hat w\in D(0,2r_w)$. Here $r_w$ is the radius of the disc on
which the typed directions take the sources $\hat w$; it is not the channel label $r$. Written without
an argument, $\mu_j$ always denotes this number; the plaquette weight defined below is always written
with its argument, $\mu_k(\cdot)$ or $\mu_k(p)$. The following requirements hold.
\begin{enumerate}
\item For $r=\mathbf{3}$ the Wilson weight of the dialled production is
$x_k(\cdot)-\zeta_k-\lambda\mathsf{c}_{\mu_k(\cdot)}$. For $r=\mathbf{6}$ its exponent carries the
vertex $\lambda\mathsf{A}^{\mathbf{6}}(\mu_k(\cdot),U_k)$. In both cases $\mu_k(p)$ is the
$(\mathsf{V}_r)_{\mathbb{C}}$-valued weight
\begin{equation*}
\mu_k(p):=\mu_0\,\mathsf{k}(p)
+\sum_{j\le k}\varphi_j(p)\sum_y h_y(p)\,[\mu_j-\mu_{j-1}](w(p_y))\,\mathsf{k}(p_y),
\end{equation*}
where $\varphi_j$, $h_y$, the points $p_y$ and the source $w$ are the letters of the same names in the
definition of the typed directions.
\item The groups at scale $j$ are formed with their own typed marginal coefficient
$\beta^{r}[\delta E^{(j)}]$, and
\begin{equation*}
\beta^{r}[\delta E^{(j)}]=\mu_j-\mu_{j-1}.
\end{equation*}
\item The direction carries the frozen slots $\delta\mathsf{P}^{(j)}$, $\delta P^{\flat(j)}$ and
$\delta Q^{(j)}$, with $\delta P^{\flat(j)}\simeq\delta\mathsf{P}^{(j)}$. This relation means that
the two slots are given by integrals with one integrand, the inputs being read as totals.
\end{enumerate}

We put $\mathsf{z}_j:=\sup_{\hat w\in D(0,2r_w)}|\mu_j|$. We write $\mathsf{e}_j$, $\mathsf{r}_j$,
$\mathsf{p}^\flat_j$, $\mathsf{q}_j$ for the slot sizes of $\delta E^{(j)}$, $\delta R^{(j)}$,
$\delta P^{\flat(j)}$, $\delta Q^{(j)}$, measured in the norms in which the one-lattice Lipschitz step
measures the corresponding slots. The loads of $\delta$ are the loads of the one-lattice Lipschitz step,
with the tensor slot in the place of the coupling slot; in particular,
\begin{equation}\label{10.4:eq-typed-step-1}
[\![\delta]\!]_n:=\sum_{j\le n}\varsigma_{1+n-j}\,\mathsf{e}_j+\sum_{j\le n}w_{n,j}\,\mathsf{r}_j
+\mathfrak{g}_n\Bigl[\mathsf{z}^\sharp_n+\sum_{j\le n}L^{-(n-j)/2}\,|\mu_j-\mu_{j-1}|\Bigr].
\end{equation}

Below, the base proposition of the one-lattice Lipschitz step is the proposition that the extension
proposition extends. The number $\gamma_L$ is the smallness threshold of the one-lattice Lipschitz
step. The letters $K_2$, $C_L$, $S_I$, $\mathfrak{m}_k$ and $k_1(k)$ in estimate $(\mathrm{b})$ below
are the constants and the index that appear in row $(\mathrm{b})$ of the one-lattice Lipschitz step.
The load $[\![\delta]\!]^{T}_k$ in estimate $(\mathrm{a})$ below is the load bearing the superscript
$T$ among the loads of the one-lattice Lipschitz step, formed from the slots of $\delta$ with the
tensor slot in the place of the coupling slot.

\begin{proposition}[The typed step with a marginal tensor slot. Stated; proof incomplete at the step
marked $(\ast)$]\label{10.4:prop-typed-step}
Let $\delta$ be a $\mathsf{V}_r$-typed direction with a marginal slot, $r\in\{\mathbf{3},\mathbf{6}\}$,
and let $g\le\gamma_L$. Then the following hold.
\begin{enumerate}
\item The base proposition of the one-lattice Lipschitz step holds for the dialled production
$\mathfrak{s}+\lambda\delta$ on $\mathbb{D}_\delta$.
\item The derivative in $\lambda$ at $\lambda=0$ of the outputs of the complete step is a slot
$\delta_{k+1}$ of the same kind.
\item The estimates of the one-lattice Lipschitz step hold for $\delta_{k+1}$. Namely,
\begin{equation}\label{10.4:eq-typed-step-a}
\begin{aligned}
&(\mathrm{a})\quad \mathsf{e}_{k+1}\le\varkappa_k\,[\![\delta]\!]^{T}_k,\qquad
\mathsf{p}^\flat_{k+1}\le\varkappa_k\Bigl[\sum_{j\le k}\varsigma_{k+1-j}\,\mathsf{p}^\flat_j
+\mathfrak{g}_k\mathsf{z}_k\Bigr],\\
&(\mathrm{b})\quad \mathsf{q}_{k+1}\le\varkappa_k\Bigl[\sum_{j\le k}\varsigma_{k+1-j}\,\mathsf{q}_j
+\sum_{j\le k_1(k)}w_{k,j}\,\mathsf{r}_j\Bigr]\\
&\qquad\qquad\quad +K_2C_L^2(1+S_I)\,\theta_k^2\,\mathfrak{m}_k
\Bigl[\sum_{j\le k}\varsigma_{k+1-j}\,\mathsf{p}^\flat_j+\mathfrak{g}_k\mathsf{z}_k\Bigr],\\
&(\mathrm{g})^{\otimes}\quad \mu_{k+1}-\mu_k=\beta^{r}[\delta\mathsf{P}^{(k+1)}]
+\beta^{r}[\delta Q^{(k+1)}],\\
&\qquad\qquad\quad |\beta^{r}[\delta Q^{(k+1)}]|\le C_\Pi\,\mathfrak{g}_{k+1}^{-1}\,\mathsf{q}_{k+1},
\qquad |\mu_{k+1}-\mu_k|\le C_\Pi\,\mathfrak{g}_{k+1}^{-1}\,\mathsf{e}_{k+1}.
\end{aligned}
\end{equation}
Estimates $(\mathrm{c})$, $(\mathrm{d})$ and $(\mathrm{e})$ of the one-lattice Lipschitz step hold
with the load \eqref{10.4:eq-typed-step-1}. Finally, the tangent inequality
\eqref{10.4:eq-pure-recursion-tangent-a} holds for the frozen slots $\delta\mathsf{P}^{(j)}$ of $\delta$.
\end{enumerate}
\end{proposition}

\begin{proof}
The proof runs through the inputs of the one-lattice Lipschitz step, one at a time. It uses the
replacements of the transfer of the scalar lemmas to a tensor slot
(\ref{10.4:prf-tensor-slot-transfer}, whose proof is incomplete at one step), and gives, for each
input, the reason it holds for the dialled production.

\emph{Families, groups and weights.} In the scalar case the four lemmas of the correlation theorem
that carry the family take one family, that is one weight, through the three conditions on families
of the correlation theorem, together with its norm $\mathfrak{n}=\sup|c|$; they are stated for
arbitrary $\mathfrak{n}$.
Each group is formed with its own coefficient $\beta[\delta E^{(j)}]=\delta b_j$.

In the present case the groups are typed. Each is formed with its own coefficient
$\beta^{r}[\delta E^{(j)}]=\mu_j-\mu_{j-1}$. What remains of a typed group after its counterterm is
bounded by the irrelevance after one typed counterterm, Lemma~\ref{10.4:lem-typed-irrelevance}.

For $r=\mathbf{3}$, the weight $\zeta_k+\lambda\,\delta\zeta_k$ of the coupling slot is replaced by
$\zeta_k+\lambda\mathsf{c}_{\mu_k(\cdot)}$. The correlation theorem's lemma on complex plaquette
weights reads the weight through $\sup|c|$, so it applies to this weight. For $r=\mathbf{6}$, the
vertex $\lambda\mathsf{A}^{\mathbf{6}}(\mu_k(\cdot),U_k)$ is controlled by the bounds for the
off-diagonal tensor vertex, Lemma~\ref{10.4:lem-off-diagonal-vertex}.

\emph{Expansions.} Two further inputs of the correlation theorem are its weight-linear expansion
identity and parts (a)--(c) of the correlation theorem's proposition that consumes it. Both read the
weight through weight-linear expansions.

\emph{The shift of the Wilson weight.} Entry (3) of the correlation theorem's lemma on the shift of
the Wilson weight is the bound on the raw shift of the Wilson weight. In the scalar case it is read
under the inequality $|\lambda\,\delta\zeta|\le E_1/\mathfrak{g}_n$, where, since
$\mathfrak{g}_n=g_n^2(\log x_n)^{2q_\flat}$ and $x_n=g_n^{-2}$,
\begin{equation*}
\frac{E_1}{\mathfrak{g}_n}=E_1\,x_n\,(\log x_n)^{-2q_\flat}.
\end{equation*}
Here it is read under the same inequality for $|\lambda\mu_n|$. For $r=\mathbf{3}$ the bound on the
raw shift then follows from part (c) of the correlation theorem's lemma on complex plaquette weights.
For $r=\mathbf{6}$ it follows from part (c) of Lemma~\ref{10.4:lem-off-diagonal-vertex}, with the
constant $C_{\mathrm{dom}}$ entering the ceiling.

\emph{Equivariance.} The production map is applied here to lists that are not $W_4$-invariant. Its
equivariance on such lists is supplied by the equivariance statement for the production map.

\emph{Estimates} (a)--(e). $(\ast)$ The proofs of (a)--(e) in the one-lattice Lipschitz step rest on
the corollary of its base proposition, on two auxiliary lemmas of that step, and on geometric
forgetting of old regions. None of these proofs uses the type of the direction. They therefore apply
with the load \eqref{10.4:eq-typed-step-1} in place of the load of the coupling slot, provided that
these proofs, and the base proposition they rest on, read the coupling slot only through the entries
listed above. The step marked $(\ast)$ is this proviso: the verification, input by input against the
list of inputs of the one-lattice Lipschitz step, that every input used in the proofs of (a)--(e) and
of that base proposition depends on the coupling slot only through the entries treated above
(families, groups and weights; expansions; the shift of the Wilson weight; equivariance). This
verification is not carried out here, and the same verification is required for the proof of
$(\mathrm{g})^{\otimes}$ below.

\emph{Estimate} $(\mathrm{g})^{\otimes}$. The identity rests on
$\delta P^{\flat(k+1)}\simeq\delta\mathsf{P}^{(k+1)}$: the two integrals defining these slots have one
integrand, the inputs being read as totals. By part (i) of the bounds on the typed marginal
coefficient, Lemma~\ref{10.4:lem-coefficient-bounds}, it follows that
\begin{equation*}
\beta^{r}[\delta P^{\flat(k+1)}]=\beta^{r}[\delta\mathsf{P}^{(k+1)}].
\end{equation*}
The two bounds in $(\mathrm{g})^{\otimes}$ are the flat bound of part (ii) of the same lemma,
\begin{equation*}
|\beta^{r}[F]|\le C_\Pi\,\mathfrak{g}_{k+1}^{-1}\,\|F\|^\flat_{k+1}
\qquad\text{on }\mathcal{U}^\flat_{k+1}.
\end{equation*}

\emph{Freezing defects.} At first order, the defects created by freezing the tensor weight
$\mathsf{k}$ are of type $V\otimes\mathsf{V}_r$. By part (d) of the hypercubic decomposition of
symmetric tensors, Lemma~\ref{10.4:lem-hypercubic-decomposition},
\begin{equation*}
\operatorname{Hom}_{W_4}\bigl(V\otimes\mathsf{V}_r,\ \mathbb{R}\oplus\operatorname{Sym}^2\Lambda^2\bigr)=0.
\end{equation*}
These defects therefore have no marginal part and no vacuum part.
\end{proof}

\begin{lemma}[Re-filing of typed marginal families. Stated; proof incomplete at the step marked $(\ast)$]
\label{10.4:lem-typed-refiling}
Let $r \in \{\mathbf{3},\mathbf{6}\}$. Then
\begin{equation}\label{10.4:eq-typed-refiling-a}
\begin{gathered}
\text{the first re-filing lemma, its corollary on irrelevance, and the two}\\
\text{re-filing statements on the re-filed transport letters}\\
\text{hold for every marginal family that is } \mathsf{V}_{r}\text{-typed in the sense of}\\
\text{Definition~\ref{10.3:def-typed-families},}
\end{gathered}
\end{equation}
with the untyped inputs of clause (d) of the first re-filing lemma and of its corollary on
irrelevance replaced by Lemmas~\ref{10.4:lem-coefficient-bounds}
and~\ref{10.4:lem-typed-irrelevance}.
\end{lemma}

\begin{proof}
We go through the first re-filing lemma clause by clause and then turn to its corollary on
irrelevance and to the two re-filing statements on the re-filed transport letters.

Clauses (a)--(c) of the first re-filing lemma are statements about each term of a family
separately, and each of them is linear in the family. These clauses are therefore applied to a
typed family one term at a time, and by linearity they hold for the family as a whole, except that
part (f3) of clause (b) needs two equivariance statements in addition (see below).

$(\ast)$ The following steps of the transfer are not carried out: part (f3) of clause (b) of the
first re-filing lemma, clause (d) of that lemma together with its corollary on irrelevance, and the
two re-filing statements on the re-filed transport letters. What the first two of these require is
indicated below.

Part (f3) of clause (b) is to be derived from two equivariance statements which accompany the
re-filing lemmas and which pass to the differential; this derivation for typed families is not
carried out here.

Clause (d) and the corollary on irrelevance use three inputs in the untyped case: the coefficient
bound, the flat extraction bound, and the irrelevance lemma. For typed families these inputs are to
be replaced as follows; let $F$ be a $\mathsf{V}_{r}$-typed pointed $E$-family at scale $j$ as in
Lemma~\ref{10.4:lem-typed-irrelevance}. The coefficient bound is to be replaced by
Lemma~\ref{10.4:lem-coefficient-bounds}, clauses (i)--(iii), whose clause (ii) gives, in the norms
$\|\cdot\|_j$ and $\|\cdot\|^{\flat}_j$ of that lemma, the bound
$|\beta^{r}[F]| \le C_\beta \|F\|_j$ on $\mathcal{U}_j$. The flat extraction
bound is to be replaced by the flat bound of the same lemma,
\begin{equation*}
|\beta^{r}[F]| \le C_\Pi\, \mathfrak{g}_j^{-1} \|F\|^{\flat}_j \quad \text{on } \mathcal{U}^\flat_j .
\end{equation*}
The irrelevance lemma is to be replaced by the bound on $\mathfrak{i}^{r}_y(\mathfrak{b};h_0)$ in
Lemma~\ref{10.4:lem-typed-irrelevance}. Clause (d) and the corollary on irrelevance are to be read
for typed families with these replacements; this reading is not carried out here.

Finally, the typed marginal coefficient $\beta^{r}$ of the remnant family left by the re-filing is
included in the coefficient $\mu$ of the re-filing lemmas. These are the steps which, once carried
out, give \eqref{10.4:eq-typed-refiling-a}.
\end{proof}

In what follows we work along the frozen tangent run in a $\mathsf{V}_r$-typed direction with a marginal slot, in the sense of Proposition~\ref{10.4:prop-typed-step}, whose proof is incomplete at one step. The running tensor coefficient is $\tau_k := \mu_k$, with $\tau_0 = 1$. Every other level-zero slot of the direction vanishes, so that the quantities $p_k$, $\hat e_k$, $p^\flat_k$, $\hat q_k$ and $r_k$ introduced below all vanish at $k=0$.
The remaining quantities are fixed as follows.
\begin{itemize}
\item $\hat e_k := \|\delta\hat E^{(k)}\|^{\flat}$ is the norm $\|\cdot\|^{\flat}$ of the re-filed slot $\delta\hat E^{(k)}$ of the direction. This quantity, the quantities $\hat q_k$, $r_k$ and $p^\flat_k$, the index set $\mathfrak{B}(k)$ and the constant $C_{RF}$ are those of the re-filing lemmas, in the typed form of Lemma~\ref{10.4:lem-typed-refiling}, whose proof is incomplete at one step. The quantity $p_k := \|\delta\mathsf{P}^{(k)}\|_k$ is as in Corollary~\ref{10.4:cor-pure-recursion-tangent}.
\item $\epsilon_k$, $\hat\varepsilon_k$, $C_q'$ and $S_1$ are the constants of the one-lattice Lipschitz step. By their definition there, $\epsilon_k\le\hat\varepsilon_k$, and $\hat\varepsilon_j\le\hat\varepsilon_k$ for $j\le k$.
\item $K_2'$ is the constant that multiplies the $\theta_k^2$-term in inequality (b) of Proposition~\ref{10.4:prop-typed-step}; in the notation of that proposition, $K_2' = K_2C_L^2(1+S_I)$.
\item $k_1(k)\le k$ is the level of the one-lattice Lipschitz step up to which the remnant terms are summed in inequality (b) of Proposition~\ref{10.4:prop-typed-step}.
\item The parameters $\vartheta_1$ and $q$ satisfy
\begin{equation*}
\max\{L^{-\hat\alpha},L^{-1/2}\} < \vartheta_1 < 1, \qquad q \le \min\{\tfrac12,\vartheta_1^2\}.
\end{equation*}
\item $j''$ denotes the re-filing level of $j$.
\item $\bar a$ and $\bar\varkappa$ denote bounds, uniform in $k$, for $a_k$ and for the step constants $\varkappa_k$.
\end{itemize}
The tangent families obey the bound \eqref{10.4:eq-pure-recursion-tangent-a} of Corollary~\ref{10.4:cor-pure-recursion-tangent}:
\begin{equation*}
\text{(v1)}\quad p_{k+1} \le a_k\sum_{j\le k}\varpi_{k-j}p_j + c^{\otimes}_k|\tau_k| .
\end{equation*}
The further inequalities assumed are the following:
\begin{equation}\label{10.4:eq-forward-transport-b}
\begin{aligned}
&\text{(v2)}\quad \hat e_{k+1} \le \varkappa_k T_k + C_{RF}\sum_{i\in\mathfrak{B}(k)} w_{k,i}\,r_i,
 \qquad p^\flat_{k+1} \le \varkappa_k T^\flat_k,\\
&\qquad\quad T_n := \sum_{j\le n}\varsigma_{1+n-j}\hat e_j + \sum_{j\le n<j''} w_{n,j}\,r_j
 + \mathfrak{g}_n\Bigl[|\tau_n| + \sum_{1\le j\le n}L^{-(n-j)/2}|\tau_j-\tau_{j-1}|\Bigr],\\
&\qquad\quad T^\flat_n := \sum_{j\le n}\varsigma_{1+n-j}p^\flat_j
 + \mathfrak{g}_n\Bigl[|\tau_n| + \sum_{1\le j\le n}L^{-(n-j)/2}|\tau_j-\tau_{j-1}|\Bigr],\\
&\qquad\quad |\tau_j-\tau_{j-1}| \le C_\Pi\,\mathfrak{g}_j^{-1}\hat e_j \quad (j\ge 1);\\
&\text{(v3)}\quad \hat q_{k+1} \le \varkappa_k\sum_{j\le k}\varsigma_{k+1-j}\hat q_j
 + (\varkappa_k + C_{RF})\sum_{i\in\mathfrak{B}(k)} w_{k,i}\,r_i\\
&\qquad\qquad\quad + K_2'\theta_k^2\hat{\mathfrak{m}}_k\Bigl[\sum_{j\le k}\varsigma_{k+1-j}p^\flat_j
 + \mathfrak{g}_k|\tau_k|\Bigr];\\
&\text{(v4)}\quad r_k \le \epsilon_k\sum_{n\le k}(k-n+1)\,q^{(k-n-1)_+}T_n
 + 3\varkappa_k\sum_{k_1(k)<j<k} r_j;\\
&\text{(v5)}\quad \tau_{k+1} = (1+\nu_k)\tau_k + \iota_{k+1},\\
&\qquad\quad |\iota_{k+1}| \le C_\beta a_k\sum_{j\le k}\varpi_{k-j}p_j
 + C_\Pi\,\mathfrak{g}_{k+1}^{-1}\hat q_{k+1}.
\end{aligned}
\end{equation}
Thus $T^\flat_n$ is $T_n$ with $p^\flat$ in place of $\hat e$ and without the remnant terms in $r$; the bound on $p^\flat_{k+1}$ is inequality (a) of Proposition~\ref{10.4:prop-typed-step}, whose proof is incomplete at one step, for $p^\flat$.

These inequalities have the following origins.
\begin{itemize}
\item Inequality (v1) is Corollary~\ref{10.4:cor-pure-recursion-tangent}.
\item Inequalities (v2)--(v4) are the inequalities (a), (b), (c) and (e) of the typed step with a marginal tensor slot (Proposition~\ref{10.4:prop-typed-step}, whose proof is incomplete at one step). They are written in the form given by the re-filing of typed marginal families (Lemma~\ref{10.4:lem-typed-refiling}, whose proof is incomplete at one step).
\item Inequality (v5) is the identity (g) of the same proposition, $\mu_{k+1}-\mu_k = \beta^{r}[\delta\mathsf{P}^{(k+1)}] + \beta^{r}[\delta Q^{(k+1)}]$, together with its bound $|\beta^{r}[\delta Q^{(k+1)}]|\le C_\Pi\,\mathfrak{g}_{k+1}^{-1}\hat q_{k+1}$. Here $\delta\mathsf{P}^{(k+1)}$ is split as in Corollary~\ref{10.4:cor-pure-recursion-tangent} into $\tau_k\,\delta^{r}\Phi_{k+1}$ and $D_h\Phi_{k+1}[\delta\mathsf{P}^{(\le k)}]$, and $\nu_k := \beta^{r}[\delta^{r}\Phi_{k+1}]$. By linearity of $\beta^{r}$, the identity (v5) holds with
\begin{equation*}
\iota_{k+1} := \beta^{r}\bigl[D_h\Phi_{k+1}[\delta\mathsf{P}^{(\le k)}]\bigr] + \beta^{r}[\delta Q^{(k+1)}],
\end{equation*}
and the bound in (v5) follows from the bound on $\beta^{r}[\delta Q^{(k+1)}]$ and from
\begin{equation*}
\bigl|\beta^{r}\bigl[D_h\Phi_{k+1}[\delta\mathsf{P}^{(\le k)}]\bigr]\bigr|
\le C_\beta a_k\sum_{j\le k}\varpi_{k-j}p_j .
\end{equation*}
\end{itemize}

The ceilings are one-sided conditions on $g$. They are imposed after $L$, the auxiliary parameters $\mathfrak{p}$, $\vartheta_1$, $q$, $\omega$ and $c_\natural$ have been fixed, where $c_\natural$ is a fixed number with $1 < c_\natural < \min\{\vartheta_1^{-1},\omega^{-1},q^{-1}\}$. We put
\begin{equation}\label{10.4:eq-forward-transport-1}
\begin{aligned}
&\mathsf{C} := 4C_q', \qquad
\mathsf{a}^{*}_k := \max\{k-i+1 : i\le k<i''\}, \qquad
\bar c^{\otimes}_k := \max_{j\le k}c^{\otimes}_j,\\
&C_{\mathsf{b}} := 12\,C_\Pi\vartheta_1^{-1}\sum_{i\ge0}(1+i)\bigl(L^{-1/2}/\vartheta_1\bigr)^i,\\
&\bar y_k := \max_{j\le k}\Bigl[6K_2'\theta_j^2\hat{\mathfrak{m}}_j
 + 20(1+C_{RF})\,\mathsf{C}\,\hat\varepsilon_j\,\vartheta_1^{-\mathsf{a}^{*}_j}\Bigr].
\end{aligned}
\end{equation}
The number $\mathsf{a}^{*}_k$ is the range of the unfiled remnants at level $k$; for an index $i$ at which the maximum is attained, $\mathsf{a}^{*}_k\le1+\log_L(M\check R_i)$, with $\check R_i$ as in the re-filing lemmas. The series defining $C_{\mathsf{b}}$ converges because $\vartheta_1 > L^{-1/2}$. The ceilings read:
\begin{equation}\label{10.4:eq-forward-transport-a}
\begin{aligned}
&(1)\quad \bar a\,S_\omega/\omega \le \tfrac12;\qquad
(2)\quad \bar\varkappa\,(3S_1 + C_{\mathsf{b}}) \le \tfrac14;\\
&(3)\quad 10\,\mathsf{C}(1+C_{RF})\,\hat\varepsilon_k(\mathsf{a}^{*}_k+1)\,\vartheta_1^{-\mathsf{a}^{*}_k}
 \le \min\{\tfrac14,\varkappa_k\};\\
&(4)\quad 6\varkappa_k\,(k-k_1(k))\,\vartheta_1^{-(k-k_1(k))} \le \tfrac14;\qquad
(5)\quad \sup_k C_\beta c^{\otimes}_k \le 1-c_\natural^{-1};\\
&(6)\quad \mathfrak{g}_j \le 2\mathfrak{g}_k \ \text{ and } \ \mathfrak{g}_k \le 4(1+k-j)\,\mathfrak{g}_j
 \quad (j\le k);\\
&(7)\quad \sum_{k<K}\Bigl[2C_\beta (S_\omega/\omega)\,c_\natural(1-\omega c_\natural)^{-1}a_k\bar c^{\otimes}_k
 + 2C_\Pi\, c_\natural(1-\vartheta_1 c_\natural)^{-1}\bar y_k\Bigr] \le \tfrac16 .
\end{aligned}
\end{equation}
Condition (6) is the form in which the bounds on the reference flow of the couplings are used below; like the other ceilings, it is a hypothesis of the lemma.

\begin{lemma}[Forward transport of the tangent rows]\label{10.4:lem-forward-transport}
Assume that (v1), that is \eqref{10.4:eq-pure-recursion-tangent-a}, and (v2)--(v5), that is \eqref{10.4:eq-forward-transport-b}, hold at every level $k$, and that the ceilings \eqref{10.4:eq-forward-transport-a} hold. Put
\begin{equation*}
G_k := \sum_{i\le k}\mathfrak{g}_i|\tau_i|\,\vartheta_1^{k-i}, \qquad
P_k := \prod_{l<k}(1+\nu_l).
\end{equation*}
Then, for every level $k$, the following hold.
\begin{enumerate}
\item[(i)] The tangent and re-filed quantities satisfy
\begin{equation*}
\begin{gathered}
p_{k+1} \le 2\sum_{i\le k}c^{\otimes}_i\,\omega^{k-i}|\tau_i|, \qquad T_k \le 2G_k, \qquad
\hat e_{k+1},\ p^\flat_{k+1} \le 3\varkappa_k G_k,\\
r_k \le \mathsf{C}\,\hat\varepsilon_k G_k, \qquad \hat q_{k+1} \le \bar y_k G_k .
\end{gathered}
\end{equation*}
\item[(ii)] The memory term satisfies $|\iota_{k+1}| \le \sum_{j\le k}\mathsf{i}_{k+1,j}|\tau_j|$, where $\mathbf{1}_{\{j<k\}}$ equals $1$ if $j<k$ and $0$ otherwise, and
\begin{equation*}
\mathsf{i}_{k+1,j} := 2C_\beta\,(a_kS_\omega/\omega)\,c^{\otimes}_j\,\omega^{k-j}\mathbf{1}_{\{j<k\}}
 + 2C_\Pi\,\bar y_k\,\vartheta_1^{k-j}.
\end{equation*}
Moreover, summed over all levels,
\begin{equation*}
\sum_{k<K}\sum_{j\le k}\mathsf{i}_{k+1,j}\,c_\natural^{k+1-j} \le \tfrac16 .
\end{equation*}
\item[(iii)] For $j\le k$, and with $C_G := \tfrac83(1-\vartheta_1c_\natural)^{-1}$,
\begin{equation*}
\Bigl|\frac{\tau_k}{P_k}-1\Bigr| \le \tfrac13, \qquad
c_\natural^{-(k-j)} \le \Bigl|\frac{P_k}{P_j}\Bigr| \le (2-c_\natural^{-1})^{k-j}, \qquad
G_k \le C_G\,\mathfrak{g}_k|P_k| .
\end{equation*}
\end{enumerate}
\end{lemma}

\begin{proof}
All quantities are non-negative except $\tau_k$, $\nu_k$, $\iota_{k+1}$ and $P_k$, which enter through their moduli. Sums over empty index sets vanish, and we put $G_{-1} := 0$. Recall that $\vartheta_1<1$ and, since $c_\natural<\min\{\vartheta_1^{-1},\omega^{-1}\}$, that $\omega c_\natural<1$ and $\vartheta_1c_\natural<1$.

\emph{Part (i): the tangent families.} We prove the bound on $p_{k+1}$ by induction on $k$. Assume that
\begin{equation*}
p_j \le 2\sum_{i<j}c^{\otimes}_i\,\omega^{j-1-i}|\tau_i| \qquad \text{for all } j\le k .
\end{equation*}
For $j=0$ this holds because $p_k$ vanishes at $k=0$ and the sum is empty, so the assumption holds at $k=0$. Inserting it and exchanging the order of summation gives
\begin{equation*}
a_k\sum_{j\le k}\varpi_{k-j}p_j
\le 2a_k\sum_{i<k}c^{\otimes}_i|\tau_i|\sum_{i<j\le k}\varpi_{k-j}\,\omega^{j-1-i}.
\end{equation*}
Writing $a := k-j$, we have $\varpi_{k-j}\,\omega^{j-1-i} = \varpi_a\omega^{-a}\cdot\omega^{k-1-i}$. Recall that $S_\omega = \sum_{a\ge0}\varpi_a\omega^{-a}$. The inner sum is therefore at most $\omega^{k-1-i}S_\omega = (S_\omega/\omega)\,\omega^{k-i}$. Hence
\begin{equation}\label{10.4:eq-forward-transport-2}
a_k\sum_{j\le k}\varpi_{k-j}p_j
\le 2\,(a_kS_\omega/\omega)\sum_{i<k}c^{\otimes}_i\,\omega^{k-i}|\tau_i|
\le \sum_{i<k}c^{\otimes}_i\,\omega^{k-i}|\tau_i| .
\end{equation}
The second inequality is ceiling (1), because $a_k\le\bar a$. Adding $c^{\otimes}_k|\tau_k|$ as in (v1) gives
\begin{equation*}
p_{k+1} \le \sum_{i<k}c^{\otimes}_i\,\omega^{k-i}|\tau_i| + c^{\otimes}_k|\tau_k|
\le 2\sum_{i\le k}c^{\otimes}_i\,\omega^{k-i}|\tau_i| ,
\end{equation*}
which is the claim at $k$.

\emph{Two facts about $G$.} First, $\mathfrak{g}_k|\tau_k|\le G_k$, since this is the term $i=k$ of $G_k$. Second, $G_n\le\vartheta_1^{-(k-n)}G_k$ for $n\le k$. Indeed, the terms of $G_k$ are non-negative, so
\begin{equation*}
G_k \ge \sum_{i\le n}\mathfrak{g}_i|\tau_i|\,\vartheta_1^{k-i} = \vartheta_1^{k-n}G_n .
\end{equation*}
By the second fact, $G_{j-1}\le\vartheta_1^{-(n-j+1)}G_n$ for $j\le n$, and $G_n\le\vartheta_1^{-(k-n)}G_k$ for $n\le k$. The two weighted sums that occur below are therefore bounded by $G_n$ times $\sum_{\nu\ge1}\varsigma_\nu\vartheta_1^{-\nu}$, and by $G_k$ times $\sum_{m\ge0}(m+1)q^{(m-1)_+}\vartheta_1^{-m}$. The latter series converges because $q\le\vartheta_1^2<\vartheta_1$. These are the series bounded by $S_1$ and by $C_q'$, so that
\begin{equation*}
\sum_{j\le n}\varsigma_{1+n-j}G_{j-1} \le S_1G_n, \qquad
\sum_{n\le k}(k-n+1)\,q^{(k-n-1)_+}G_n \le C_q'\,G_k .
\end{equation*}

\emph{Part (i): the re-filed quantities.} We prove by induction on $k$ that at level $k$
\begin{equation*}
T_k\le 2G_k,\quad r_k\le\mathsf{C}\hat\varepsilon_kG_k,\quad
\hat e_{k+1},\,p^\flat_{k+1}\le3\varkappa_kG_k,\quad \hat q_{k+1}\le\bar y_kG_k .
\end{equation*}
Assume these at all levels below $k$; at $k=0$ the assumption is vacuous. Since the level-zero quantities vanish and $G_{-1}=0$, the assumption gives the following bounds.
\begin{itemize}
\item For $1\le j\le k$, $\hat e_j\le3\bar\varkappa G_{j-1}$, $p^\flat_j\le3\bar\varkappa G_{j-1}$, and $\hat q_j\le\bar y_{j-1}G_{j-1}\le\bar y_kG_{j-1}$. The last inequality holds because $\bar y$ is a running maximum. For $j=0$ the three quantities vanish, so that $\hat e_0\le3\bar\varkappa G_{-1}$, $p^\flat_0\le3\bar\varkappa G_{-1}$ and $\hat q_0\le\bar y_kG_{-1}$ hold with equality, both sides being zero.
\item For all $j<k$, $r_j\le\mathsf{C}\hat\varepsilon_jG_j$, and for all $n<k$, $T_n\le2G_n$.
\end{itemize}
Put $T_k^- := T_k - w_{k,k}r_k$. The index $j=k$ always satisfies $k\le k<k''$, because $k''>k$ by the re-filing lemmas, so $T_k^-$ is $T_k$ with the remnant term $j=k$ removed.

The $\hat e$-sum in $T_k$ is at most
\begin{equation*}
3\bar\varkappa\sum_{j\le k}\varsigma_{1+k-j}G_{j-1} \le 3\bar\varkappa S_1G_k .
\end{equation*}
For the increment sum, we use in turn the increment bound in (v2), the bound on $\hat e_j$, ceiling (6), and the second fact. For $1\le j\le k$ this gives
\begin{equation*}
\begin{aligned}
\mathfrak{g}_kL^{-(k-j)/2}|\tau_j-\tau_{j-1}|
&\le L^{-(k-j)/2}\,\frac{\mathfrak{g}_k}{\mathfrak{g}_j}\,C_\Pi\cdot3\bar\varkappa\,G_{j-1}\\
&\le 12\,C_\Pi\bar\varkappa\,\vartheta_1^{-1}(1+k-j)\bigl(L^{-1/2}/\vartheta_1\bigr)^{k-j}G_k .
\end{aligned}
\end{equation*}
Summed over $j$, this is at most $C_{\mathsf{b}}\bar\varkappa G_k$. The term $\mathfrak{g}_k|\tau_k|$ is at most $G_k$ by the first fact.

The unfiled remnants with $j<k$, that is $j<k<j''$, satisfy $k-j+1\le\mathsf{a}^{*}_k$ by the definition of $\mathsf{a}^{*}_k$. There are at most $\mathsf{a}^{*}_k$ of them. On that range $w_{k,j}\le\tfrac52\,x_j/x_k\le5$. Moreover $r_j\le\mathsf{C}\hat\varepsilon_jG_j\le\mathsf{C}\hat\varepsilon_j\vartheta_1^{-(k-j)}G_k$, and $\hat\varepsilon_j\le\hat\varepsilon_k$. Together these remnants contribute at most $5\mathsf{C}\hat\varepsilon_k\mathsf{a}^{*}_k\vartheta_1^{-\mathsf{a}^{*}_k}G_k$.

Ceiling (2) gives $\bar\varkappa(3S_1+C_{\mathsf{b}})\le\frac14$. Ceiling (3), with $1+C_{RF}\ge1$ and $\mathsf{a}^{*}_k+1\ge\mathsf{a}^{*}_k$, gives $5\mathsf{C}\hat\varepsilon_k\mathsf{a}^{*}_k\vartheta_1^{-\mathsf{a}^{*}_k}\le\frac18$. Hence
\begin{equation*}
T_k^- \le \bigl(\tfrac14+\tfrac18+1\bigr)G_k = \tfrac{11}{8}G_k .
\end{equation*}

We next bound $r_k$ using (v4). Its term $n=k$ is $T_k=T_k^-+w_{k,k}r_k$, with $w_{k,k}\le\tfrac52$ and $\tfrac52\epsilon_k\le\tfrac18$. The latter holds because $\epsilon_k\le\hat\varepsilon_k$ and $10\,\mathsf{C}\hat\varepsilon_k\le\tfrac14$ by ceiling (3), all its other factors being at least $1$, while $\mathsf{C}=4C_q'\ge4$, since $C_q'$ bounds a series whose term $m=0$ equals $1$. For $n<k$ we use $T_n\le2G_n$, and at $n=k$ we use $T_k^-\le2G_k$. The second weighted-sum bound then gives
\begin{equation*}
\epsilon_k\Bigl[\sum_{n<k}(k-n+1)q^{(k-n-1)_+}T_n + T_k^-\Bigr] \le 2C_q'\epsilon_kG_k .
\end{equation*}
In the last sum of (v4), each $j$ with $k_1(k)<j<k$ has $r_j\le\mathsf{C}\hat\varepsilon_j\vartheta_1^{-(k-j)}G_k$, where $\hat\varepsilon_j\le\hat\varepsilon_k$ and $k-j\le k-k_1(k)$, and there are fewer than $k-k_1(k)$ such $j$. Writing $k_1:=k_1(k)$, we obtain
\begin{equation*}
\bigl(1-\tfrac18\bigr)r_k
\le 2C_q'\epsilon_kG_k + 3\varkappa_k(k-k_1)\,\mathsf{C}\hat\varepsilon_k\vartheta_1^{-(k-k_1)}G_k
\le \bigl(\tfrac12+\tfrac18\bigr)\mathsf{C}\hat\varepsilon_kG_k .
\end{equation*}
Here $2C_q'\epsilon_k\le2C_q'\hat\varepsilon_k=\tfrac12\mathsf{C}\hat\varepsilon_k$, and the second term is bounded by ceiling (4), which gives $3\varkappa_k(k-k_1)\vartheta_1^{-(k-k_1)}\le\tfrac18$. Dividing by $\tfrac78$ yields
\begin{equation*}
r_k\le\tfrac57\,\mathsf{C}\hat\varepsilon_kG_k\le\mathsf{C}\hat\varepsilon_kG_k .
\end{equation*}
By ceiling (3), $10\,\mathsf{C}\hat\varepsilon_k\le\tfrac14$, because $\vartheta_1^{-\mathsf{a}^{*}_k}\ge1$. Therefore
\begin{equation*}
T_k = T_k^- + w_{k,k}r_k \le \tfrac{11}{8}G_k + \tfrac52\,\mathsf{C}\hat\varepsilon_kG_k
\le \bigl(\tfrac{11}{8}+\tfrac1{16}\bigr)G_k \le 2G_k .
\end{equation*}

In (v2) the set $\mathfrak{B}(k)$ has at most two elements, by the cardinality bound of the re-filing lemmas in the typed form of Lemma~\ref{10.4:lem-typed-refiling}, whose proof is incomplete at one step. Its elements $i$ satisfy $i\le k<i''$, by the same lemmas, so that, as for the unfiled remnants, $w_{k,i}\le5$ and $r_i\le\mathsf{C}\hat\varepsilon_k\vartheta_1^{-\mathsf{a}^{*}_k}G_k$. Hence
\begin{equation*}
\hat e_{k+1} \le 2\varkappa_kG_k + C_{RF}\cdot2\cdot5\,\mathsf{C}\hat\varepsilon_k\vartheta_1^{-\mathsf{a}^{*}_k}G_k
\le 3\varkappa_kG_k ,
\end{equation*}
since $10\,\mathsf{C}C_{RF}\hat\varepsilon_k\vartheta_1^{-\mathsf{a}^{*}_k}\le\varkappa_k$ by ceiling (3). The same holds for $p^\flat$. By the bound on $p^\flat_j$, the increment bound above, the first fact and ceiling (2),
\begin{equation*}
T^\flat_k \le (3\bar\varkappa S_1 + C_{\mathsf{b}}\bar\varkappa + 1)\,G_k \le \tfrac54G_k ,
\end{equation*}
so that $p^\flat_{k+1}\le\frac54\varkappa_kG_k\le3\varkappa_kG_k$.

Finally, we bound the three terms of (v3).
\begin{itemize}
\item By $\hat q_j\le\bar y_kG_{j-1}$ and the first weighted-sum bound, the first term is at most $\bar\varkappa S_1\bar y_kG_k$.
\item As for $\hat e_{k+1}$, using $\varkappa_k\le1$, the second term is at most $10(1+C_{RF})\,\mathsf{C}\hat\varepsilon_k\vartheta_1^{-\mathsf{a}^{*}_k}G_k$.
\item By $p^\flat_j\le3\bar\varkappa G_{j-1}$ and the first fact, the third term is at most $K_2'\theta_k^2\hat{\mathfrak{m}}_k(3\bar\varkappa S_1+1)G_k$.
\end{itemize}
Ceiling (2) gives $\bar\varkappa S_1\le\frac1{12}$ and $3\bar\varkappa S_1+1\le\frac54$. The definition of $\bar y_k$ gives $10(1+C_{RF})\mathsf{C}\hat\varepsilon_k\vartheta_1^{-\mathsf{a}^{*}_k}\le\frac12\bar y_k$ and $K_2'\theta_k^2\hat{\mathfrak{m}}_k\le\frac16\bar y_k$. Therefore
\begin{equation*}
\hat q_{k+1} \le \bigl(\tfrac1{12}+\tfrac12+\tfrac5{24}\bigr)\bar y_kG_k \le \bar y_kG_k .
\end{equation*}
This completes the induction and the proof of (i).

\emph{Part (ii).} We bound the two terms of (v5). By \eqref{10.4:eq-forward-transport-2}, the first term is at most
\begin{equation*}
2C_\beta(a_kS_\omega/\omega)\sum_{j<k}c^{\otimes}_j\omega^{k-j}|\tau_j| .
\end{equation*}
By (i), the second term is at most $C_\Pi\bar y_k\,\mathfrak{g}_{k+1}^{-1}G_k$. Ceiling (6) gives $\mathfrak{g}_i\le2\mathfrak{g}_{k+1}$ for $i\le k+1$, so
\begin{equation*}
\mathfrak{g}_{k+1}^{-1}G_k = \sum_{i\le k}\frac{\mathfrak{g}_i}{\mathfrak{g}_{k+1}}\,\vartheta_1^{k-i}|\tau_i|
\le 2\sum_{i\le k}\vartheta_1^{k-i}|\tau_i| .
\end{equation*}
Adding the two bounds gives $|\iota_{k+1}|\le\sum_{j\le k}\mathsf{i}_{k+1,j}|\tau_j|$.

For the sum, fix $k$. Summing geometric series in $k-j$, and using $c^{\otimes}_j\le\bar c^{\otimes}_k$ for $j<k$, we get
\begin{equation*}
\begin{aligned}
\sum_{j<k}c^{\otimes}_j\,\omega^{k-j}c_\natural^{k+1-j}
&\le c_\natural\,\bar c^{\otimes}_k\sum_{m\ge1}(\omega c_\natural)^m
\le c_\natural(1-\omega c_\natural)^{-1}\bar c^{\otimes}_k,\\
\sum_{j\le k}\vartheta_1^{k-j}c_\natural^{k+1-j}
&\le c_\natural(1-\vartheta_1c_\natural)^{-1}.
\end{aligned}
\end{equation*}
Hence $\sum_{j\le k}\mathsf{i}_{k+1,j}c_\natural^{k+1-j}$ is at most the bracket in ceiling (7). Summing over $k<K$, ceiling (7) gives the bound $\frac16$.

\emph{Part (iii).} Since $\nu_l=\beta^{r}[\delta^{r}\Phi_{l+1}]$, the bound on $\beta^{r}$ used for (v5), together with the typed response bound $c^{\otimes}_l$ of Corollary~\ref{10.4:cor-pure-recursion-tangent}, gives $|\nu_l|\le C_\beta c^{\otimes}_l$. Ceiling (5) then gives $C_\beta c^{\otimes}_l\le1-c_\natural^{-1}$. Hence
\begin{equation*}
c_\natural^{-1}\le|1+\nu_l|\le2-c_\natural^{-1}
\end{equation*}
for every $l$. In particular $P_k\ne0$. Taking the product over $j\le l<k$, since $P_k/P_j=\prod_{j\le l<k}(1+\nu_l)$, gives the two-sided bound on $|P_k/P_j|$.

Put $\tilde\tau_k := \tau_k/P_k$. Note that $\tilde\tau_0=1$, because $P_0=1$ is an empty product and $\tau_0=1$. Dividing (v5) by $P_{k+1}=(1+\nu_k)P_k$ gives
\begin{equation}\label{10.4:eq-forward-transport-3}
\tilde\tau_{k+1}-\tilde\tau_k = \frac{\iota_{k+1}}{P_{k+1}} .
\end{equation}
By (ii), $|\tau_j|=|\tilde\tau_j||P_j|$, and the lower bound $|P_{k+1}/P_j|\ge c_\natural^{-(k+1-j)}$, we obtain
\begin{equation*}
\begin{gathered}
\Bigl|\frac{\iota_{k+1}}{P_{k+1}}\Bigr|
\le \sum_{j\le k}\mathsf{i}_{k+1,j}|\tilde\tau_j|\,\Bigl|\frac{P_j}{P_{k+1}}\Bigr|
\le \hat\iota_{k+1}\max_{j\le k}|\tilde\tau_j|,\\
\hat\iota_{k+1} := \sum_{j\le k}\mathsf{i}_{k+1,j}\,c_\natural^{k+1-j}.
\end{gathered}
\end{equation*}
By \eqref{10.4:eq-forward-transport-3}, therefore,
\begin{equation*}
\max_{j\le k+1}|\tilde\tau_j|\le(1+\hat\iota_{k+1})\max_{j\le k}|\tilde\tau_j| .
\end{equation*}
Starting from $|\tilde\tau_0|=1$ and using $\sum_k\hat\iota_{k+1}\le\frac16$ from (ii), we find
\begin{equation*}
\max_{j\le k}|\tilde\tau_j| \le \prod_{l<k}(1+\hat\iota_{l+1}) \le \exp\sum_{l<k}\hat\iota_{l+1} \le e^{1/6}.
\end{equation*}
Summing \eqref{10.4:eq-forward-transport-3} over $l<k$ then gives
\begin{equation*}
|\tilde\tau_k-1| \le \sum_{l<k}\hat\iota_{l+1}\max_{j\le l}|\tilde\tau_j| \le \tfrac16e^{1/6} \le \tfrac13 .
\end{equation*}

Finally, for $i\le k$ we have $|\tau_i|=|\tilde\tau_i||P_i|\le\frac43|P_i|$ and $|P_i|\le c_\natural^{k-i}|P_k|$, the latter from the lower bound on $|P_k/P_i|$. Ceiling (6) gives $\mathfrak{g}_i\le2\mathfrak{g}_k$. Hence
\begin{equation*}
G_k \le \tfrac43\cdot2\,\mathfrak{g}_k|P_k|\sum_{i\le k}(\vartheta_1c_\natural)^{k-i}
\le \tfrac83(1-\vartheta_1c_\natural)^{-1}\mathfrak{g}_k|P_k| = C_G\,\mathfrak{g}_k|P_k| .
\end{equation*}
\end{proof}

\begin{proof}[The running normalisation and its one-loop form: proof of
\eqref{10.4:eq-insertion-correlation-a} and \eqref{10.4:eq-insertion-correlation-b}]
\label{10.4:prf-running-normalisation}
This part of the proof of the stress-tensor-insertion correlation theorem
(Theorem~\ref{10.4:thm-insertion-correlation}) derives the format
\eqref{10.4:eq-insertion-correlation-a} and the one-loop form
\eqref{10.4:eq-insertion-correlation-b} from the following statements, which are its hypotheses:
\begin{enumerate}
\item the typed step with a marginal tensor slot, \eqref{10.4:eq-typed-step-a} of
Proposition~\ref{10.4:prop-typed-step}, whose proof is incomplete at one step;
\item parts (a), (c) and (d) of Proposition~\ref{10.4:prop-vertex-response} (response of one
step to the tensor vertex), whose proof is incomplete at one step;
\item the forward transport of the tangent rows, Lemma~\ref{10.4:lem-forward-transport},
together with the hypotheses under which that lemma is stated, for the frozen tangent run:
the first of these hypotheses, which precedes \eqref{10.4:eq-forward-transport-b}, the
transport inequalities
\eqref{10.4:eq-forward-transport-b}, and the ceilings
\eqref{10.4:eq-forward-transport-a};
\item the reference comparisons for the running shift, in the two forms used below.
\end{enumerate}

\emph{Existence.} By Proposition~\ref{10.4:prop-typed-step}, whose proof is incomplete at one
step, applied at each level $k \le K$, the tangent of a run is the run of the tangents.
Hence the frozen tangent run exists level by level in the format of
\eqref{10.4:eq-insertion-correlation-a}: the weight $x_k-\zeta_k-t\tau_k\mathsf{c}$ for
$r=\mathbf{3}$ (for $r=\mathbf{6}$, the vertex $+t\tau_k\mathsf{A}^{\mathbf{6}}(1,U_k;h_0)$),
with $\tau_k = N_r(0,k;\zeta)$ and $N_r(0,0)=1$; the recursion of
\eqref{10.4:eq-insertion-correlation-a} for $N_r(0,k+1)$, with the step increment $\nu^{r}_k$
and the memory term $\iota_{k+1}$; and the typed slots
$\delta\mathsf{P}^{(j)}, \delta P^{\flat(j)}, \delta Q^{(j)}, \delta R^{(j)}$, $j \le k$, of
\eqref{10.4:eq-insertion-correlation-a}.

\emph{Analyticity and reality.} The analyticity of $N_r(0,k;\cdot)$ on $D(0,2r)$ and its reality
on real arguments are those of the run of the correlation theorem, transmitted through the
holomorphy of $\Phi_{k+1}$ in $(\mathfrak{q},s,u)$ and along holomorphic families in
$\mathbb{B}^w_k$ asserted in Proposition~\ref{10.4:prop-vertex-response}(a), whose proof is
incomplete at one step. Here and below, the letter $r$ in $D(0,2r)$ and in the bounds
$\tfrac83 rg_k^2$ and $6rg_k^4$ is the number fixing the disc $D(0,2r)$ of
\eqref{10.4:eq-insertion-correlation-a}, not the channel label $r$ of $\nu^{r}_k$.

\emph{Transport.} Apply Lemma~\ref{10.4:lem-forward-transport} with $\tau_k = N_r(0,k;\zeta)$ and
$\nu_l = \nu^{r}_l$, so that $P_k = P^{r}_k$. Part (ii) of the lemma is the bound on the memory
term $\iota_{k+1}$ in \eqref{10.4:eq-insertion-correlation-a}, together with the bound
$\tfrac16$ on the weighted sum of its kernel $\mathsf{i}_{k+1,j}$. By part (iii) with $j = 0$
and $P^{r}_0 = 1$, $|P^{r}_k| \ge c_\natural^{-k} > 0$. Part (iii) also gives
$|\tau_k/P_k - 1| \le \tfrac13$; setting $\epsilon_k := N_r(0,k)/P^{r}_k - 1$, this is
$N_r(0,k) = (1+\epsilon_k)P^{r}_k$ with $|\epsilon_k| \le \tfrac13$, and since
$|1+\epsilon_k| \ge \tfrac23$, we have $N_r(0,k) \ne 0$.

By \eqref{10.4:eq-insertion-correlation-a},
$\nu^{r}_k = \beta^{r}[\delta^{r}\Phi_{k+1}(g_k,s_k;h)]$, with the parameters $\xi_i$ of
\eqref{10.4:eq-insertion-correlation-a} set equal to $\zeta_i(\zeta)$, and with
\begin{equation*}
s_k = 1 - \zeta_k g_k^2, \qquad |s_k - 1| \le \tfrac83\, r g_k^2, \qquad
h = \mathsf{P}^{(\le k)}_\xi = L^{(\le k)} + \Phi^{(\le k)}, \qquad
|\Phi^{(\le k)}|_k \le S_\varpi \sup_{j<k} e_j .
\end{equation*}
Here the bound on $|s_k - 1|$ is the first of the reference comparisons in the last
hypothesis listed above. The decomposition of the pure family $\mathsf{P}^{(\le k)}_\xi$ into
its one-loop part $L^{(\le k)}$ and the remainder $\Phi^{(\le k)}$, and the bound on
$|\Phi^{(\le k)}|_k$ with the constants $S_\varpi$ and $e_j$ (and the history unit
$\mathsf{E}$ used below), are those of the bounds on the pure families $\mathsf{P}^{(j)}$,
$j \le k$, of the frozen run, which fix $S_\varpi$, $e_j$ and $\mathsf{E}$; they are not
consequences of the hypotheses listed above.
By Proposition~\ref{10.4:prop-vertex-response}(c), whose proof is incomplete at one step,
$\|\delta^{r}\Phi_{k+1}(g_k,s;h)\|_{k+1} \le c^{\otimes}_k$ for all $s \in D(1,\sigma_0)$ and
$h \in \mathbb{B}^w_k$, where
$c^{\otimes}_k = (2E^{*}/u_0)\,g_k^2\mathsf{r}_k^{-2}$ is the typed response bound.
Proposition~\ref{10.4:prop-vertex-response} writes $c^{\otimes}_k$ as $2E^{*}/u_0$ times the
square of its radius ratio, which it denotes $\tau_k$ and whose square equals
$g_k^2\mathsf{r}_k^{-2}$; in this proof we write $g_k^2\mathsf{r}_k^{-2}$ for that square, to
keep $\tau_k$ for $N_r(0,k;\zeta)$, and we read the bounds of part (d) in the same way below.
With the
norm bound $|\beta^{r}[F]| \le C_\beta\|F\|$ on the typed marginal coefficient, the step
increment obeys
\begin{equation*}
|\nu^{r}_k| \le C_\beta c^{\otimes}_k
= C_\beta\,\frac{2E^{*}}{u_0}\Bigl(\frac{2T_k}{\varepsilon_1}\Bigr)^{2} g_k^{2},
\end{equation*}
where we recall that the radius of the coupling slot of
Proposition~\ref{10.4:prop-vertex-response} is $\mathsf{r}_k = \varepsilon_1/(2T_k)$, with
$T_k$ the cut $T = T_k$ of that proposition.

\emph{One-loop form.} With $\nu^{r}_k$, $s_k$ and $h$ as recalled above,
write $\mathsf{L} := \mathsf{L}^{r}_{k+1}$ and split
\begin{align*}
\delta^{r}\Phi_{k+1}(g_k,s_k;h) - g_k^2\mathsf{L}(1;L^{(\le k)})
&= \bigl[\delta^{r}\Phi_{k+1}(g_k,s_k;h) - g_k^2\mathsf{L}(s_k;h)\bigr]
 + g_k^2\bigl[\mathsf{L}(s_k;h) - \mathsf{L}(1;h)\bigr]\\
&\quad + g_k^2\bigl[\mathsf{L}(1;h) - \mathsf{L}(1;L^{(\le k)})\bigr].
\end{align*}
Proposition~\ref{10.4:prop-vertex-response}(d), whose proof is incomplete at one step, bounds
the three brackets in turn: the first by its estimate of
$\delta^{r}\Phi_{k+1}(g_k,s;h) - g_k^2\mathsf{L}^{r}_{k+1}(s;h)$; the second by its Lipschitz
estimate in $s$ together with the bound on $|s_k - 1|$; the third by its Lipschitz estimate in
$h$ together with $h - L^{(\le k)} = \Phi^{(\le k)}$ and the bound on $|\Phi^{(\le k)}|_k$:
\begin{align*}
\bigl\|\delta^{r}\Phi_{k+1}(g_k,s_k;h) - g_k^2\mathsf{L}(s_k;h)\bigr\|
&\le \frac{2E^{*}}{u_0}\,\frac{(g_k^2\mathsf{r}_k^{-2})^2}{1 - g_k^2\mathsf{r}_k^{-2}},\\
g_k^2\bigl\|\mathsf{L}(s_k;h) - \mathsf{L}(1;h)\bigr\|
&\le g_k^2\,\frac{4E^{*}}{u_0}\,\mathsf{r}_k^{-2}\,\frac{|s_k-1|}{\sigma_0}\\
&\le \frac{2E^{*}}{u_0}\, g_k^2\mathsf{r}_k^{-2}\,\frac{16r}{3\sigma_0}\, g_k^2 ,\\
g_k^2\bigl\|\mathsf{L}(1;h) - \mathsf{L}(1;L^{(\le k)})\bigr\|
&\le \frac{2E^{*}}{u_0}\, g_k^2\mathsf{r}_k^{-2}\,\frac{S_\varpi \sup_{j<k} e_j}{2\mathsf{E}}.
\end{align*}
Applying the typed marginal coefficient and the norm bound $|\beta^{r}[F]| \le C_\beta\|F\|$
recalled above,
\begin{equation}\label{10.4:eq-running-normalisation-1}
\begin{split}
\bigl|\nu^{r}_k - g_k^2\beta^{r}[\mathsf{L}^{r}_{k+1}(1;L^{(\le k)})]\bigr|
\le C_\beta\,\frac{2E^{*}}{u_0}\Bigl[
&\frac{(g_k^2\mathsf{r}_k^{-2})^2}{1 - g_k^2\mathsf{r}_k^{-2}}
+ g_k^2\mathsf{r}_k^{-2}\,\frac{16r}{3\sigma_0}\, g_k^2\\
&+ g_k^2\mathsf{r}_k^{-2}\,\frac{S_\varpi \sup_{j<k} e_j}{2\mathsf{E}}\Bigr].
\end{split}
\end{equation}
By the second of the reference comparisons in the last hypothesis listed above,
$|g_k^2 - (x_k - \zeta_k)^{-1}| \le 6 r g_k^4$. Define
\begin{equation*}
\mathsf{m}_r(k+1) := -\beta^{r}[\mathsf{L}^{r}_{k+1}(1;L^{(\le k)})],
\end{equation*}
which is the number denoted $-\beta^{r}[\mathsf{L}^{r}_{k+1}(L^{(\le k)})]$ in
\eqref{10.4:eq-insertion-correlation-b}. By the bound
$\|\mathsf{L}^{r}_{k+1}\| \le (2E^{*}/u_0)\mathsf{r}_k^{-2}$ of
Proposition~\ref{10.4:prop-vertex-response}(d), whose proof is incomplete at one step,
$|\mathsf{m}_r(k+1)| \le C_\beta (2E^{*}/u_0)\mathsf{r}_k^{-2}$. Setting
$\varrho^{r}_k := \nu^{r}_k + \mathsf{m}_r(k+1)/(x_k - \zeta_k)$, we have the identity
\begin{equation*}
\varrho^{r}_k = \bigl(\nu^{r}_k - g_k^2\beta^{r}[\mathsf{L}^{r}_{k+1}(1;L^{(\le k)})]\bigr)
- \mathsf{m}_r(k+1)\bigl(g_k^2 - (x_k-\zeta_k)^{-1}\bigr),
\end{equation*}
so that $\nu^{r}_k = -\mathsf{m}_r(k+1)/(x_k - \zeta_k) + \varrho^{r}_k$, and $|\varrho^{r}_k|$
is at most the right side of \eqref{10.4:eq-running-normalisation-1} plus
$C_\beta(2E^{*}/u_0)\mathsf{r}_k^{-2}\cdot 6rg_k^4$.
\end{proof}

\begin{remark}[The transport without re-filing and the floor it needs]
\label{10.4:rem-transport-without-refiling}
In this remark $\kappa_0$ is the exponent in the bound on the re-filed transport letters
$\hat\varepsilon_j$ displayed below; Ba{\l}aban's conditions leave it free, and the ceilings on $g$
under which Lemma~\ref{10.4:lem-forward-transport} is proved use
only the one-sided condition $\kappa_0\ge 7$ on it. It enters through the
one-sided bound
\begin{equation*}
\hat\varepsilon_j \lesssim x_j^{-\kappa_0/2}\cdot(\text{a power of } \log x_j),
\end{equation*}
and $k_1(k)$ is as in the forward transport of the tangent rows
(Lemma~\ref{10.4:lem-forward-transport}). Suppose that the transport of
the tangent rows is carried out without the re-filing lemmas. Then the proof of
part~(iii) of Lemma~\ref{10.4:lem-forward-transport} requires, in place of the factor
$c_\natural^{k+1-j}$, a bound
\begin{equation}\label{10.4:eq-transport-without-refiling-1}
\Bigl|\frac{P_j}{P_{k+1}}\Bigr| \le A\Bigl(\frac{x_j}{x_k}\Bigr)^{a_-}
\end{equation}
with a constant $A$ and an exponent $a_-\ge 0$, together with the convergence of
\begin{equation}\label{10.4:eq-transport-without-refiling-2}
\sum_k \varkappa_k\, x_k^{-a_-} \sum_{j\le k_1(k)} \hat\varepsilon_j\, x_j^{a_-}.
\end{equation}
Given the bound on $\hat\varepsilon_j$ above, this is the one-sided floor
\begin{equation*}
\kappa_0 \ge 7 + 2a_-
\end{equation*}
on $\kappa_0$, where $a_-\ge 0$ is any exponent for which
\eqref{10.4:eq-transport-without-refiling-1} holds; if
\eqref{10.4:eq-normalisation-decay-a} holds with exponent $1$ together with the matching lower
bound, that is, in two-sided form, then \eqref{10.4:eq-transport-without-refiling-1} holds with
$a_-=1$. In this variant the bound
\eqref{10.4:eq-insertion-correlation-d} of the stress-tensor-insertion correlation
theorem carries, added to $|P_{K'}|$, the floor term
\begin{equation*}
C\,\frac{x_0}{X}\,\epsilon^{\sharp}\!\Bigl(\frac{x_0}{3}\Bigr),
\end{equation*}
the quantity denoted $\Phi_{K'}$ (not the step's family $\Phi_{k+1}$) in the statement fixing
the target exponent for the decay of the tensor correlation, in which the function
$\epsilon^{\sharp}$ is introduced.
The argument of this chapter does not use this variant; it uses the re-filing lemmas.
\end{remark}

\begin{proof}
In the bound for $\hat q_{k+1}$ in \eqref{10.4:eq-forward-transport-b} the remnant
sum is $\sum_{i\in\mathfrak{B}(k)} w_{k,i}\, r_i$, with the range $\mathfrak{B}(k)$
supplied by the re-filing lemmas. Without them this sum runs over all
$j\le k_1(k)$, with the flat weight $w_{k,j}$. Consequently the kernel
$\mathsf{i}_{k+1,j}$ of the memory term in part~(ii) of
Lemma~\ref{10.4:lem-forward-transport} contains the term
\begin{equation*}
2C_\Pi\,\varkappa_k\,\mathsf{C}\,\hat\varepsilon_j ,
\end{equation*}
which has no decay in $k-j$. Part~(iii) of that lemma uses the factor
$c_\natural^{k+1-j}$ as the bound for $|P_j/P_{k+1}|$, and part~(ii) weights the
kernel by that factor. For a kernel term without decay in $k-j$ the proof of
part~(iii) needs instead the bound \eqref{10.4:eq-transport-without-refiling-1}.
Weighting the term displayed above by the right-hand side of
\eqref{10.4:eq-transport-without-refiling-1} and summing over $j\le k_1(k)$ and over
$k$ gives $2C_\Pi\,\mathsf{C}\,A$ times the double sum
\eqref{10.4:eq-transport-without-refiling-2}. The proof of part~(iii) therefore also
needs this double sum to converge. The floor $\kappa_0\ge 7+2a_-$ and the term added to
$|P_{K'}|$ in \eqref{10.4:eq-insertion-correlation-d} are then as stated above.
\end{proof}

\begin{proposition}[The sourced format of the local tangent run]\label{10.4:prf-sourced-format}
Stated; the proof below is incomplete at the step marked $(\ast)$. The step marked $(\ast)$
concerns $r=\mathbf{6}$ only; for $r=\mathbf{3}$ the proposition is an implication from the
hypotheses (i)--(vi) below.

Let $r\in\{\mathbf{3},\mathbf{6}\}$. Let $g$ and the data of the correlation theorem satisfy the
standing hypotheses of Theorem~\ref{10.4:thm-insertion-correlation}, with the ceilings fixed
there, and let $K'$ be as in \eqref{10.4:eq-insertion-correlation-c}. Assume:
\begin{enumerate}
\item[(i)] the typed step with a marginal tensor slot, \eqref{10.4:eq-typed-step-a}
(Proposition~\ref{10.4:prop-typed-step}, whose proof is incomplete at one step);
\item[(ii)] the forward transport of the tangent rows, Lemma~\ref{10.4:lem-forward-transport};
\item[(iii)] the $\mathbf{R}$-step proposition in the proof of the correlation theorem;
\item[(iv)] the splitting of the weighted Wilson action at the $\mathbf{R}$-step in the proof of
the correlation theorem, which is linear in the weight and bounded by weighted absolute sums;
\item[(v)] the lemma on the cut-off functions $\varphi_j$ in the proof of the correlation
theorem, together with its closing remark that the transition layers of the $\varphi_j$ are
disjoint;
\item[(vi)] the response of one step to the tensor vertex,
Proposition~\ref{10.4:prop-vertex-response}(c), whose proof is incomplete at one step.
\end{enumerate}
Then \eqref{10.4:eq-insertion-correlation-c} holds: for every $k\le K'$, every
$w\in\mathcal{W}_k$ and every tensor weight $\mathsf{k}$ with
$\mathsf{k}/\|\mathsf{k}\|_\sharp\in\mathcal{H}_k$, the derivative at $t=0$ of the density
$\rho^{w,t}_k$ obtained after $k$ steps has the first-order format
\begin{equation}\label{10.4:eq-sourced-format-a}
\begin{aligned}
\partial_t|_{0}\,\rho^{w,t}_k&=\sum_{\text{labels}}\chi_k\,\partial_t|_{0}
\bigl[\mathbf{T}^{w,t}_k\exp A^{w,t}_k\bigr],\\
\partial_t|_{0}\,A^{w,t}_k&=\mathsf{A}^{r}(\tau_k(\cdot),U_k)+\delta\mathbf{E}_k
+\delta\mathbf{D}_k+\delta\mathbf{R}_k+\delta\mathbf{B}_k-\delta E_k,
\end{aligned}
\end{equation}
the sum running over the labels of the expansion of the correlation theorem,
with the following ingredients, and with the $\flat$-norm bounds stated in
\eqref{10.4:eq-insertion-correlation-c}.
\begin{enumerate}
\item[(a)] $\tau_k(p)$ is the weight of \eqref{10.4:eq-insertion-correlation-c}, given by the
formula of item (a) of the sourced format in the proof of the correlation theorem, with
$h^{(j)}_y\ge 0$ and $\sum_y h^{(j)}_y=1$.
\item[(b)] The irrelevant part is
\begin{equation*}
\delta\mathbf{E}_k:=\sum_{j\le k}\ \sum_{y\in\Lambda_j^{\circledR}}
\mathfrak{i}^{r,(j)}_y\bigl(U_k;\mathsf{k}(p_y)\bigr)\Big|_{\hat w=w(p_y)}
+\bigl(\text{the unmatched counterterm pieces}\bigr),
\end{equation*}
where $\Lambda_j^{\circledR}$ is the set of points $y$ indexing the scale-$j$ typed groups, and
the $\mathfrak{i}^{r,(j)}_y$ are the typed irrelevant groups at scale $j$, in the sense of
Lemma~\ref{10.4:lem-typed-irrelevance}, built on the frozen tangent families.
\item[(c)] The local slots $\delta E^{(j)}(X,\cdot,y;w,\mathsf{k})$ are linear in
$\mathsf{k}\in\mathcal{H}_j(X)$, analytic in $w$ on $\mathcal{W}_j(X)$, and equal to the frozen
slot wherever $(w,\mathsf{k})$ is constant on $X^{+}$; the freezing defect $\delta D^{(j)}$ is
the local slot minus the frozen one.
\item[(d)] $\delta\mathbf{R}_k$ and $\delta\mathbf{B}_k$ are the $t$-derivatives of the terms of
items (d) and (d)$'$ of the sourced format in the proof of the correlation theorem. They are
derived objects in the sense of the large-field bounds for the linearised step: the terms
produced by the operation $\mathbf{B}$ carry no slot of their own in the format and are bounded
through the terms from which they derive.
\end{enumerate}
Throughout the format the test functions, measures, trees, surrogates and denominators of the
expansion of the correlation theorem do not depend on $t$.
\end{proposition}

\begin{proof}
We argue by induction on $k$. The step from $k$ to $k+1$ is the typed step
\eqref{10.4:eq-typed-step-a} of Proposition~\ref{10.4:prop-typed-step}, whose proof is
incomplete at one step; it consists of the following operations, taken in the order of the
step of the correlation theorem.

\emph{The pure part.} The pure part of the step is the tangent of the pure recursion in a
typed direction, Corollary~\ref{10.4:cor-pure-recursion-tangent}.

\emph{The frozen integral.} The frozen integral of the step yields the typed slots
$\delta E^{(k+1)}$ and $\delta P^{\flat(k+1)}$. The bridge between the frozen integral and
the pure recursion compares the slot $\delta P^{\flat(k+1)}$ with the tangent
$\delta\mathsf{P}^{(k+1)}$ of the pure part, and the remaining slot is
\begin{equation*}
\delta Q^{(k+1)}:=\delta E^{(k+1)}-\delta P^{\flat(k+1)}.
\end{equation*}
These slots enter item (b) at level $k+1$.

\emph{The flow.} The running tensor coefficient is carried from level $k$ to level $k+1$ by
the last of the relations \eqref{10.4:eq-forward-transport-b} and is controlled by the forward
transport of the tangent rows, Lemma~\ref{10.4:lem-forward-transport}. The reference inverse
couplings $x_k$, the number of steps $K$ and the thresholds do not depend on $t$ and are left
untouched.

\emph{The $\mathbf{T}$-steps.} The operations $\mathbf{T}$ of the step of the correlation
theorem are applied unchanged. The new vertex difference
\begin{equation*}
\mathsf{A}^{r}\bigl(\tau_k,U_k(e^{iB'}V^{(k)})\bigr)-\mathsf{A}^{r}(\tau_k,U_{k+1})
\end{equation*}
is a function of the quantity $g_kCA$ of the $\mathbf{T}$-steps (in the notation of the proof of
the correlation theorem) that vanishes at $0$ and carries no factor $g_k^{-2}$; this is the
statement established in the $\mathbf{T}$-steps of the proof of the correlation theorem for the
term denoted $\mathfrak{S}_k$ in that proof (not the spin term of the Ward identity), now read
for the vertex. Next the term
\begin{equation*}
\sum_y\beta^{r}\bigl[\delta E^{(k+1)}_{w(p_y)}\bigr]\,
\mathsf{A}^{r}\bigl(h^{(k+1)}_y\varphi_{k+1},U_{k+1};\mathsf{k}(p_y)\bigr)
\end{equation*}
is subtracted and moved into the vertex. This is an add-and-subtract, so the total is
unchanged.

\emph{The weight is zone-wise.} Fix $n$ and a plaquette $p$ in the zone $n$, that is the set
$\Omega_n\setminus\Omega_{n+1}$, where the $\Omega_n$ are the regions of the proof of the
correlation theorem. There $\varphi_j(p)=1$ for $j\le n$ and $\varphi_j(p)=0$ for $j>n+1$, and
the transition layers of the $\varphi_i$ are disjoint, by the lemma of hypothesis (v) and its
closing remark. Hence in the sum over $j$ defining $\tau_k(p)$ in
\eqref{10.4:eq-insertion-correlation-c} only the indices $j\le n+1$ survive, those with
$j\le n$ with coefficient one. Replace every $\mathsf{k}(p_y)$ in that sum by $\mathsf{k}(p)$.
Since $\sum_y h^{(j)}_y(p)=1$, the $j$-th term becomes
$\varphi_j(p)\,[N_r(0,j;\cdot)-N_r(0,j-1;\cdot)]\,\mathsf{k}(p)$, and since $N_r(0,0;\cdot)=1$
by \eqref{10.4:eq-insertion-correlation-a}, the sum telescopes:
\begin{align*}
&\mathsf{k}(p)+\sum_{j\le n}\bigl[N_r(0,j;\cdot)-N_r(0,j-1;\cdot)\bigr]\mathsf{k}(p)
+\varphi_{n+1}(p)\bigl[N_r(0,n+1;\cdot)-N_r(0,n;\cdot)\bigr]\mathsf{k}(p)\\
&\qquad=\bigl(1-\varphi_{n+1}(p)\bigr)N_r(0,n;\cdot)\,\mathsf{k}(p)
+\varphi_{n+1}(p)\,N_r(0,n+1;\cdot)\,\mathsf{k}(p),
\end{align*}
a convex combination of $N_r(0,n;\cdot)\mathsf{k}$ and $N_r(0,n+1;\cdot)\mathsf{k}$. Because
$h^{(j)}_y\ge0$ and $\sum_yh^{(j)}_y(p)=1$, the error made in the $j$-th term by the
replacement is at most $|N_r(0,j)-N_r(0,j-1)|$ times the supremum of
$|\mathsf{k}(p_y)-\mathsf{k}(p)|$ over the $y$ with $h^{(j)}_y(p)\neq0$. The total error is
therefore bounded by
\begin{align*}
&\sum_{j\le n+1}\bigl|N_r(0,j)-N_r(0,j-1)\bigr|\cdot
\sup\bigl\{|\mathsf{k}(p_y)-\mathsf{k}(p)|:\ h^{(j)}_y(p)\neq0\bigr\}\\
&\qquad\le C\sum_{j\le n+1}c_\natural^{\,n+1-j}\,|P_{n+1}|\,\mathfrak{l}_j\,\|\mathsf{k}\|_\sharp,
\qquad c_\natural<L^{1/2},
\end{align*}
with the constants $\mathfrak{l}_j$ of the large-field bounds for the linearised step.
Combining the convex combination with this error bound gives, on the zone $n$,
$|\tau_k(p)|\le 2|P_n(0)|\,\|\mathsf{k}\|_\sharp$. This is the
analogue of the statement, in the proof of the correlation theorem, that $x_j(p)$ lies
between two consecutive $x_i$. On the older zones $\tau_k$ is not bounded by $C|P_k|$ when
$P$ decays, and need not be: the loads of the large-field bounds for the linearised step are
zone-wise, and on the zone $n$
\begin{equation*}
\mathfrak{g}_n\sup_{\text{zone }n}|\tau_k|\le 2\,\mathfrak{g}_n\,|P_n|\,\|\mathsf{k}\|_\sharp .
\end{equation*}

\emph{The $\mathbf{R}$-step.} The $\mathbf{R}$-step proposition of hypothesis (iii) applies by
items (c)--(e) of \eqref{10.4:eq-typed-step-a}, which is part of
Proposition~\ref{10.4:prop-typed-step}, whose proof is incomplete at one step. The splitting
of hypothesis (iv) is used with the weight $\tau_k$. For $r=\mathbf{3}$ it applies verbatim.

$(\ast)$ For $r=\mathbf{6}$ this step rests on the following input, which is not derived
here. It is written in the notation of the proof of the correlation theorem. The splitting of
hypothesis (iv) reads
\begin{equation*}
A(\zeta_k,U''_k)=A(\zeta_k^{\mathrm{new}},U_k)+A(\zeta_k\hat\chi_1,U''_{h+2})+\mathfrak{c}_w,
\end{equation*}
where $U''_k$ and $U''_{h+2}$ are configurations, $\zeta_k^{\mathrm{new}}$ and
$\zeta_k\hat\chi_1$ are weights and $\mathfrak{c}_w$ is the remainder term; the subscript
$h+2$ is the index of that splitting, and the letter $h$ there is unrelated to the partition
weights $h^{(j)}_y$ and to the history variable below. Further, $\theta_W$ and $x$ are the
quantities so denoted in the third item of the modification lemma of that proof, the bound on
the raw shift of the Wilson weight. Write $\tau_k^{\mathrm{new}}$ for the weight formed from
$\tau_k$ as $\zeta_k^{\mathrm{new}}$ is formed from $\zeta_k$, and $\mathfrak{c}^{\mathbf{6}}_w$
for the corresponding remainder term. Let $\Delta^{\mathbf{6}}_k$ denote the raw shift of the
Wilson weight that the third item of the modification lemma bounds, computed for the exponent
that carries the vertex $\mathsf{A}^{\mathbf{6}}(\tau_k,\cdot)$. The input is
\begin{equation}\label{10.4:eq-sourced-format-b}
\begin{aligned}
&\mathsf{A}^{\mathbf{6}}(\tau_k,U''_k)=\mathsf{A}^{\mathbf{6}}(\tau_k^{\mathrm{new}},U_k)
+\mathsf{A}^{\mathbf{6}}(\tau_k\hat\chi_1,U''_{h+2})+\mathfrak{c}^{\mathbf{6}}_w,\\
&\qquad\text{each of the three terms on the right obeying the bound of hypothesis (iv),}\\
&\qquad\text{by weighted absolute sums of }|\tau_k|,\\
&\bigl|\Delta^{\mathbf{6}}_k\bigr|\le\theta_W\,x .
\end{aligned}
\end{equation}
For the first line, the expansions \cite[(1.16)--(1.17), (1.36), (1.46)--(1.47)]{B-CMP122b}
are applied to the two plaquette factors of each term of
$\mathsf{A}^{\mathbf{6}}(\tau_k,\cdot)$; all these expansions are small, but only their
boundedness is needed. The last line is the third item of the modification lemma in the proof
of the correlation theorem, the bound on the raw shift of the Wilson weight, read for
$\Delta^{\mathbf{6}}_k$; the input for it is Lemma~\ref{10.4:lem-off-diagonal-vertex}(c).

\emph{The bounds.} The $\flat$-norm bounds of \eqref{10.4:eq-insertion-correlation-c} follow
from Lemma~\ref{10.4:lem-forward-transport}(i) and (iii), applied to the suprema over
$(\hat w,w,\mathsf{k})$; the inequalities to which that lemma is applied are stated for these
suprema. It remains to control the dependence on $\hat w$. By the response of one step to the
tensor vertex, Proposition~\ref{10.4:prop-vertex-response}(c), whose proof is incomplete at
one step, $\delta^{r}\Phi_{l+1}(g_l,\cdot\,;\cdot)$ is bounded by $c^{\otimes}_l$ on
$D(1,\sigma_0)\times\mathbb{B}^w_l$. Item (a) of the Schwarz-lemma bounds, in the variable
$s$, and the argument of the lemma on the dependence of the pure recursion on the history, in
the variable $h$, then give
\begin{align*}
|\nu_l(\hat w)-\nu_l(0)|
&\le C_\beta\,c^{\otimes}_l\Bigl[\frac{2|s_l(\hat w)-1|}{\sigma_0}
+\frac{|\mathsf{P}_{\xi(\hat w)}-\mathsf{P}_0|_l}{2\mathsf{E}}\Bigr]\\
&\le C\,g_l^{4}\,(\log x_l)^{3p_0}.
\end{align*}
In the second inequality, where $r$ denotes the radius of the disc $D(0,2r)$ on which $\hat w$
ranges and not the channel index, one uses $|s_l-1|\le\tfrac83\,r\,g_l^{2}$, and
$|\mathsf{P}_\xi-\mathsf{P}_0|_l\le 2\mathfrak{c}\cdot\tfrac83\,r$ by item (a) of the theorem
on the pure recursion, $\mathfrak{c}$ being the constant of that theorem (this constant
$\mathfrak{c}$ is unrelated to the remainder terms $\mathfrak{c}_w$,
$\mathfrak{c}^{\mathbf{6}}_w$ above), the difference being,
slot by slot, at most $C\sum_j\varpi_{l-j}c_{j-1}$ with the constants $c_j$ of the pure
recursion. The right-hand side is summable in $l$. Since
\begin{equation*}
\frac{P_j(\hat w)}{P_j(0)}=\prod_{l<j}\frac{1+\nu_l(\hat w)}{1+\nu_l(0)},
\end{equation*}
it follows that $\tfrac12\le|P_j(\hat w)/P_j(0)|\le 2$ on $D(0,2r)$ for $g$ below a ceiling.
\end{proof}

\begin{proposition}[The torus bound for the tensor insertion]\label{10.4:prf-tensor-torus-bound}
Work in the setting of the stress-tensor-insertion correlation theorem
(Theorem~\ref{10.4:thm-insertion-correlation}), for a channel $r\in\{\mathbf{3},\mathbf{6}\}$. Let
$\delta$ be the tangent direction of \eqref{10.4:eq-insertion-correlation-c}, and let $G_l$, $P_l$,
$T_l$, $\hat e_l$ and $r_l$ be the quantities of the forward transport lemma
(Lemma~\ref{10.4:lem-forward-transport}). In this setting $qc_\natural<1$, and the flat norm unit
satisfies $\mathfrak{g}_l\le2\mathfrak{g}_{K'}$ for $l\le K'$.
Assume the following.
\begin{enumerate}
\item[(a)] The read-out of the vector package's source, applied with the direction $\delta$. The
partition function factorises as $Z=e^{\mathcal{E}}\tilde Z$, where $\mathcal{E}$ is the sum of the
read-out's local terms $e_j(X;z,t)$ over the levels $j\le K'$ and the regions $X$ of level $j$. Let
$\bar r>0$ be the source radius of the read-out. There is a constant $C_\alpha(n,\mathsf{Q})$,
depending on $n$ and on the torus constant $\mathsf{Q}$, such that, with
$\mathsf{h}:=\tilde Z/\tilde Z(0,0)$ and for $n\ge 1$,
\begin{equation}\label{10.4:eq-tensor-torus-bound-1}
\bigl|\partial_t\partial_{z_1}\cdots\partial_{z_n}\operatorname{Log}\mathsf{h}(0,0)\bigr|
\le C_\alpha(n,\mathsf{Q})\,\bar r^{-n}E_1^{-1}\,\Sigma^q_{K'}[\delta]\,
\|\mathsf{k}\|_\sharp\prod_{i=1}^n\|f_i\|_\sharp ,
\end{equation}
where
\begin{equation}\label{10.4:eq-tensor-torus-bound-2}
\Sigma^q_{K'}[\delta]
:=\sum_{l\le K'}(K'-l+1)\,q^{(K'-l-1)_+}\,[\![\delta]\!]_l ,
\end{equation}
and $[\![\delta]\!]_l$ is the load of $\delta$ at level $l$. The local terms are built on the typed
tangent slots of the frozen tangent run of \eqref{10.4:eq-insertion-correlation-c}: the
$t$-derivative $\partial_t e_j(X;z,0)$ is the sum of
$\sum_{y\in X}\delta E^{(j)}(X,1,y;w,\mathsf{k})$, of $\delta R^{(j)}(X,1;w,\mathsf{k})$, of the
contribution of the vertex $\mathsf{A}^{r}$ at the trivial configuration, and of the frozen typed
vacuum parts. There is a constant $C_e$ such that the sum of the first two of these is bounded in
absolute value by
\begin{equation*}
C_e(\hat e_j+r_j)\,\|\mathsf{k}\|_\sharp\,e^{-(1-\delta_0)\kappa d_j(X)} ,
\end{equation*}
with $\delta_0$ as in hypothesis~(d).
\item[(b)] The forward transport lemma, Lemma~\ref{10.4:lem-forward-transport}.
\item[(c)] Part~(b) of the hypercubic decomposition of symmetric tensors,
Lemma~\ref{10.4:lem-hypercubic-decomposition}: $\operatorname{Hom}_{W_4}(\mathsf{V}_r,\mathbb{R})=0$ for
$r=\mathbf{3},\mathbf{6}$.
\item[(d)] The summation of boundary terms in the proof of the correlation theorem. For a level $j$
let $\mathfrak{t}_j$ denote the radius of the correlation theorem at level $j$, and for a region $X$
let $X^+$ denote its enlargement in the correlation theorem. The summation applies to local
terms $e_j(X;z,t)$ of $\mathcal{E}$ that have two properties. First, the $t$-derivative at $t=0$
is bounded by
\begin{equation*}
\mathsf{u}_j\,\|\mathsf{k}\|_\sharp\,e^{-(1-\delta_0)\kappa d_j(X)} ,
\end{equation*}
with nonnegative numbers $\mathsf{u}_j$ and with $\delta_0$ the loss parameter in the exponent of
the correlation theorem's bound on local terms. Second, the mixed $z$-derivative of that
$t$-derivative vanishes unless $X^+$ meets
$\operatorname{supp}\mathsf{k}$ and every $\operatorname{supp}f_i$. For such terms, the contribution
of $\mathcal{E}$ to $\partial_t\partial_{z_1}\cdots\partial_{z_n}\log Z(0,0)$ is bounded by a
constant times $\|\mathsf{k}\|_\sharp\prod_i\|f_i\|_\sharp$ times
\begin{equation*}
\sum_j \mathsf{u}_j\bigl(1+(M\mathfrak{t}_j)^{-4}\bigr)
e^{-\frac14\kappa[\mathsf{d}/(4M\mathfrak{t}_j)-3]_+} .
\end{equation*}
\end{enumerate}
Then the torus bound \eqref{10.4:eq-insertion-correlation-d} holds for every $n\ge1$, every
$\mathsf{d}>0$, and all $\mathsf{k},f_1,\dots,f_n$ whose supports are pairwise at distance at least
$\mathsf{d}$.
\end{proposition}

\begin{proof}
From $Z=e^{\mathcal{E}}\tilde Z$ and the definition of $\mathsf{h}$ we have the identity
\begin{equation*}
\log Z=\mathcal{E}+\operatorname{Log}\mathsf{h}+\log\tilde Z(0,0).
\end{equation*}
Since $n\ge1$, the constant $\log\tilde Z(0,0)$ is annihilated by $\partial_{z_1}$. Hence
\begin{equation}\label{10.4:eq-tensor-torus-bound-3}
\langle\mathfrak{T}^{r}(\mathsf{k});\mathcal{O}(f_1);\dots;\mathcal{O}(f_n)\rangle
=\partial_t\partial_{z_1}\cdots\partial_{z_n}\log Z(0,0)=(\alpha)+(\beta).
\end{equation}
Here $(\alpha)$ is the mixed derivative of $\operatorname{Log}\mathsf{h}$ at $(0,0)$, and $(\beta)$
is that of $\mathcal{E}$. We bound the two terms in turn.

\emph{The bulk term $(\alpha)$.} By hypothesis~(a), $(\alpha)$ is bounded by the right side of
\eqref{10.4:eq-tensor-torus-bound-1}, so it suffices to bound $\Sigma^q_{K'}[\delta]$.

For each level $l$, part~(i) of Lemma~\ref{10.4:lem-forward-transport} gives
\begin{equation*}
[\![\delta]\!]_l\le T_l+(\text{zone-wise }\mathsf{z}^\sharp_l)\le 3G_l ,
\end{equation*}
where $\mathsf{z}^\sharp_l$ is the marginal slot size; in particular, part~(i) contains the bound
$T_l\le2G_l$ for the first summand.

By part~(iii) of the same lemma, $G_l\le C_G\,\mathfrak{g}_l|P_l|$. Inserting these two bounds into
\eqref{10.4:eq-tensor-torus-bound-2} gives
\begin{align*}
\Sigma^q_{K'}[\delta]
&\le 3\sum_{l\le K'}(K'-l+1)\,q^{(K'-l-1)_+}\,G_l\\
&\le 3C_G\sum_{l\le K'}(K'-l+1)\,q^{(K'-l-1)_+}\,\mathfrak{g}_l|P_l| .
\end{align*}

Part~(iii) of Lemma~\ref{10.4:lem-forward-transport}, read with $k=K'$ and $j=l$, gives
$c_\natural^{-(K'-l)}\le|P_{K'}/P_l|$, that is,
\begin{equation*}
|P_l|\le c_\natural^{K'-l}|P_{K'}| .
\end{equation*}
Moreover $\mathfrak{g}_l\le2\mathfrak{g}_{K'}$ for $l\le K'$, as recorded in the statement. Writing
$i:=K'-l$, we obtain
\begin{equation*}
\Sigma^q_{K'}[\delta]
\le 6C_G\,\mathfrak{g}_{K'}|P_{K'}|\sum_{i\ge0}(i+1)\,q^{(i-1)_+}c_\natural^{i} .
\end{equation*}
For $i\ge1$ the summand equals $c_\natural(i+1)(qc_\natural)^{i-1}$. Since $qc_\natural<1$, as
recorded in the statement, the series converges. Therefore
\begin{equation}\label{10.4:eq-tensor-torus-bound-4}
\Sigma^q_{K'}[\delta]\le C\,\mathfrak{g}_{K'}|P_{K'}|
\end{equation}
with a constant $C$ that depends on $C_G$, $q$ and $c_\natural$ only. Through
\eqref{10.4:eq-tensor-torus-bound-1}, this bounds $(\alpha)$ by
\begin{equation*}
C_\alpha(n,\mathsf{Q})\,\bar r^{-n}E_1^{-1}\,C\,\mathfrak{g}_{K'}\,|P_{K'}|\,
\|\mathsf{k}\|_\sharp\prod_i\|f_i\|_\sharp .
\end{equation*}

\emph{The boundary terms $(\beta)$.} By hypothesis~(a), the $t$-derivative at $t=0$ of the local
term $e_j(X;z,t)$ is the sum of the slots $\delta E^{(j)}$ and $\delta R^{(j)}$, of the contribution
of the vertex $\mathsf{A}^{r}$ at the trivial configuration, and of the frozen typed vacuum parts.
The last two vanish. First, the vertex $\mathsf{A}^{r}$ vanishes at the trivial configuration.
Second, the frozen typed vacuum parts vanish, by hypothesis~(c): for $r=\mathbf{3},\mathbf{6}$ there
is no nonzero $W_4$-equivariant map from $\mathsf{V}_r$ to $\mathbb{R}$. Hence
\begin{equation}\label{10.4:eq-tensor-torus-bound-5}
\partial_t e_j(X;z,0)
=\sum_{y\in X}\delta E^{(j)}(X,1,y;w,\mathsf{k})+\delta R^{(j)}(X,1;w,\mathsf{k}).
\end{equation}

By hypothesis~(a), the right side of \eqref{10.4:eq-tensor-torus-bound-5} is bounded by
\begin{equation*}
C_e(\hat e_j+r_j)\,\|\mathsf{k}\|_\sharp\,e^{-(1-\delta_0)\kappa d_j(X)} .
\end{equation*}

Its mixed $z$-derivative vanishes unless $X^+$ meets $\operatorname{supp}\mathsf{k}$ and every
$\operatorname{supp}f_i$. Since $n\ge1$ and the supports are pairwise at distance at least
$\mathsf{d}$, a non-vanishing term has
\begin{equation*}
4M\mathfrak{t}_j\bigl(d_j(X)+3\bigr)\ge\mathsf{d} .
\end{equation*}
Thus hypothesis~(d) applies with $\mathsf{u}_j:=C_e(\hat e_j+r_j)$. It reduces $(\beta)$ to a
constant times $\|\mathsf{k}\|_\sharp\prod_i\|f_i\|_\sharp$ times
\begin{equation*}
\sum_j(\hat e_j+r_j)\bigl(1+(M\mathfrak{t}_j)^{-4}\bigr)
e^{-\frac14\kappa[\mathsf{d}/(4M\mathfrak{t}_j)-3]_+}.
\end{equation*}

To bound this sum, we use four facts.
\begin{itemize}
\item Part~(i) of Lemma~\ref{10.4:lem-forward-transport}, read with $k=j-1$ and with $k=j$, gives
$\hat e_j\le3\varkappa_{j-1}G_{j-1}$ and $r_j\le\mathsf{C}\hat\varepsilon_jG_j$.
\item Part~(iii) of that lemma gives, for $l=j-1$ and $l=j$, $G_l\le C_G\mathfrak{g}_l|P_l|$ and
$|P_l|\le c_\natural^{K'-l}|P_{K'}|$.
\item $\mathfrak{g}_l\le2\mathfrak{g}_{K'}$ for $l\le K'$, as recorded in the statement.
\item With $j=K'-i$, one has $\mathfrak{t}_{K'-i}=L^{-i}\mathfrak{t}_{K'}$.
\end{itemize}
The first three give, with $i:=K'-j$ and $c_\natural^{K'-j+1}=c_\natural\,c_\natural^{\,i}$,
\begin{equation*}
\hat e_j+r_j\le 2C_G\,\mathfrak{g}_{K'}|P_{K'}|\,c_\natural^{\,i}
\bigl(3c_\natural\varkappa_{j-1}+\mathsf{C}\hat\varepsilon_j\bigr).
\end{equation*}
The step constants are majorised through $\bar\varkappa$ in the form
$3c_\natural\varkappa_{j-1}+\mathsf{C}\hat\varepsilon_j\le4\bar\varkappa$. Inserting this and the
fourth fact, we obtain
\begin{align*}
&\sum_j(\hat e_j+r_j)\bigl(1+(M\mathfrak{t}_j)^{-4}\bigr)
e^{-\frac14\kappa[\mathsf{d}/(4M\mathfrak{t}_j)-3]_+}\\
&\quad\le 4\bar\varkappa C_G\,\mathfrak{g}_{K'}|P_{K'}|\cdot 2\sum_{i\ge0}c_\natural^{i}
\bigl(1+L^{4i}(M\mathfrak{t}_{K'})^{-4}\bigr)
e^{-\frac14\kappa[\mathsf{d}L^{i}/(4M\mathfrak{t}_{K'})-3]_+}\\
&\quad\le C(\mathsf{d})\,|P_{K'}| .
\end{align*}
In the last step, once $\mathsf{d}L^i/(4M\mathfrak{t}_{K'})>3$ the exponential equals
$e^{3\kappa/4}e^{-cL^i}$ with $c:=\kappa\mathsf{d}/(16M\mathfrak{t}_{K'})>0$; this dominates the
growth of $c_\natural^{\,i}\bigl(1+L^{4i}(M\mathfrak{t}_{K'})^{-4}\bigr)$, so the series over $i$
converges. This bounds $(\beta)$ by a constant $C(\mathsf{d})$ times
$|P_{K'}|\,\|\mathsf{k}\|_\sharp\prod_i\|f_i\|_\sharp$.

\emph{Conclusion.} Insert the bounds on $(\alpha)$ and $(\beta)$ into
\eqref{10.4:eq-tensor-torus-bound-3}, with $|P_{K'}|$ the product $|P^{r}_{K'}(0)|$ of
\eqref{10.4:eq-insertion-correlation-d}. This gives \eqref{10.4:eq-insertion-correlation-d}. The
constant obtained is assembled from $C_\alpha(n,\mathsf{Q})$, $\bar r$, $E_1$, the constant $C$ of
\eqref{10.4:eq-tensor-torus-bound-4}, $\mathfrak{g}_{K'}$, the constant of hypothesis~(d), and
$C(\mathsf{d})$, which collects $C_e$, $\bar\varkappa$, $C_G$, $c_\natural$, $\mathfrak{g}_{K'}$,
$\kappa$, $M$, $L$, $\mathfrak{t}_{K'}$ and $\mathsf{d}$; the torus constant $\mathsf{Q}$ enters
through $C_\alpha(n,\mathsf{Q})$ alone.
\end{proof}

\begin{remark}[Passage to infinite volume]
The passage to infinite volume proceeds exactly as for the vector package's source. Along a
sequence of tori $\mathbb{T}_{m_j}$ with numbers of renormalisation steps $K_j$, the floor
$K_j\ge\mathsf{Q}(m_j)^{j}$ serves if and only if $|P_{K'}|\,\mathsf{Q}^{n+1}\to0$ as $j\to\infty$,
with $K'$ and $\mathsf{Q}$ those of the $j$-th torus. This holds under
the normalisation decay \eqref{10.4:eq-tensor-correlation-decay-a} with positive exponent.
Alternatively, one uses the per-ball volume-free bound in the form it takes for the vector
package's source.
\end{remark}

\begin{lemma}[From the one-loop inequality to normalisation decay]\label{10.4:lem-normalisation-decay}
Fix the channel $r \in \{\mathbf{3}, \mathbf{6}\}$. The coefficients $\beta^{0}_l$ and the shifts $s_l$
below are those of the correlation theorem, related to each other as in (b), and $e_l$, $v_l$ are
the error sequences of the reference bounds in (c). Assume the following.
\begin{enumerate}
\item[(a)] The one-loop anisotropy inequality holds: there are $\kappa \in \mathbb{R}$ and
$c_2^{(r)} \ge 0$, depending only on $L$ and on the gauge group (and on neither the complex sources
nor the field), such that
\begin{equation}\label{10.4:eq-normalisation-decay-a}
\sum_{l=i+1}^{j} \bigl[\mathsf{m}_r(l) - \kappa\, \beta^{0}_l\bigr] \;\ge\; -c_2^{(r)}
\qquad \text{for all } 0 \le i < j .
\end{equation}
The inequality is one-sided.
\item[(b)] The drift bound for the shifted inverse couplings holds, with the step window of the
correlation theorem. That is, for $0 \le l \le K$ the shifts satisfy $|s_l| \le c_2$, for
$0 \le l < K$ they satisfy $\beta^{0}_{l+1} = c_0 + s_l - s_{l+1}$, the numbers
$x'_l := x_l - s_l$ are positive and satisfy $|1/x_l - 1/x'_l| \le 2c_2/x_l^2$, and for
$0 \le l < K$
\begin{equation*}
x'_l - x'_{l+1} = c_0 + \bigl(b_{l+1}(0) - \beta^{0}_{l+1}\bigr) \in [\lambda, 3\lambda].
\end{equation*}
\item[(c)] The reference bounds hold: at $\zeta = 0$ the increments $\nu^{r}_l(0)$ are real and
satisfy $|\nu^{r}_l(0)| \le 1/2$, and $|b_{l+1}(0) - \beta^{0}_{l+1}| \le C_\beta e_l + v_l$, where
\begin{equation*}
\sum_{l=j}^{k-1} (C_\beta e_l + v_l)/x'_l \le C
\qquad \text{for all } 0 \le j < k \le K ,
\end{equation*}
with $C$ independent of $K$.
\item[(d)] The one-loop expansion \eqref{10.4:eq-insertion-correlation-b} of the step increment
$\nu^{r}_k$ holds. It is part of the stress-tensor-insertion correlation theorem
(Theorem~\ref{10.4:thm-insertion-correlation}).
\end{enumerate}
Then, with the same $\kappa$, there is a constant $C_N$, independent of $K$, such that
\begin{equation*}
\bigl|P^{r}_k(0)/P^{r}_j(0)\bigr| \le C_N \,(x_k/x_j)^{\kappa}
\qquad \text{for all } 0 \le j \le k \le K .
\end{equation*}
This is the normalisation decay hypothesis \eqref{10.4:eq-tensor-correlation-decay-a} of the
decay of the tensor correlation under normalisation decay
(Corollary~\ref{10.4:cor-tensor-correlation-decay}).
\end{lemma}

\begin{proof}
For $j = k$ the ratio equals $1$ and there is nothing to prove, so let $j < k$. Throughout we
work at $\zeta = 0$.

\emph{Reduction to a lower bound on a one-loop sum.}
By (c), at $\zeta = 0$ the increments $\nu_l := \nu^{r}_l(0)$ are real and satisfy
$|\nu_l| \le 1/2$; we write $u_l := 1/x_l$. The
ratio of normalisations is the product of the step factors $1 + \nu_l$ for $j \le l \le k-1$, and
each factor is positive. Since $\log(1+\nu) \le \nu$ for $\nu > -1$, we obtain
\begin{equation*}
\log \bigl|P^{r}_k(0)/P^{r}_j(0)\bigr| = \sum_{l=j}^{k-1} \log(1+\nu_l)
\le \sum_{l=j}^{k-1} \nu_l .
\end{equation*}
At $\zeta = 0$ the denominator in \eqref{10.4:eq-insertion-correlation-b} is $x_l$. Hypothesis (d)
therefore gives
\begin{equation}\label{10.4:eq-normalisation-decay-1}
\sum_{l=j}^{k-1} \nu_l = -\sum_{l=j}^{k-1} \mathsf{m}_r(l+1)\, u_l
+ \sum_{l=j}^{k-1} \varrho^{r}_l(0).
\end{equation}
Hypothesis (d) also bounds the remainder: $|\varrho^{r}_l| \le C_\varrho g_l^3 (\log x_l)^{4p_0}$.
Summed over $l$, using $g_l^3 = x_l^{-3/2}$ and the at least linear growth of $x_l$ in $k - l$
that follows from (b), this gives $\sum_l |\varrho^{r}_l(0)| \le C$.

It therefore suffices to prove
\begin{equation}\label{10.4:eq-normalisation-decay-2}
\sum_{l=j}^{k-1} \mathsf{m}_r(l+1)\, u_l \;\ge\; \kappa \log(x_j/x_k) - C ,
\end{equation}
with $C$ independent of $K$, $j$ and $k$, and for $\kappa$ of either sign.

\emph{Consequences of the drift bound.}
Write $u'_l := 1/x'_l$. By hypothesis (b) we have $x'_l - x'_{l+1} \ge \lambda > 0$, so $x'_l$ is
decreasing in $l$, and since the $x'_l$ are positive, $u'_l$ is positive and increasing. Hypothesis
(b) also gives $|s_l| \le c_2$ and $|u_l - u'_l| \le 2c_2 u_l^2$.

Introduce
\begin{equation*}
D_l := \mathsf{m}_r(l) - \kappa\,\beta^{0}_l, \qquad S(l,k) := \sum_{i=l+1}^{k} D_i .
\end{equation*}
Applying \eqref{10.4:eq-normalisation-decay-a} to the pair of levels $(l, k)$ gives
$S(l,k) \ge -c_2^{(r)}$ for $l < k$. Moreover $S(k,k) = 0$.

Since $\mathsf{m}_r(l+1) = D_{l+1} + \kappa\beta^{0}_{l+1}$, we decompose
\begin{equation}\label{10.4:eq-normalisation-decay-3}
\sum_{l=j}^{k-1} \mathsf{m}_r(l+1)\, u_l
= \sum_{l=j}^{k-1} D_{l+1}\, u'_l
+ \sum_{l=j}^{k-1} \mathsf{m}_r(l+1)\,(u_l - u'_l)
+ \kappa \sum_{l=j}^{k-1} \beta^{0}_{l+1}\, u'_l .
\end{equation}
We bound the three sums in turn.

\emph{First sum.}
Since $D_{l+1} = S(l,k) - S(l+1,k)$, summation by parts gives
\begin{equation*}
\sum_{l=j}^{k-1} D_{l+1}\, u'_l
= u'_j\, S(j,k) + \sum_{l=j+1}^{k-1} (u'_l - u'_{l-1})\, S(l,k) .
\end{equation*}
Here $u'_j > 0$, each difference $u'_l - u'_{l-1}$ is non-negative because $u'$ is increasing, and
each $S(l,k)$ is at least $-c_2^{(r)}$. The right-hand side is therefore at least
\begin{equation*}
-c_2^{(r)} \Bigl(u'_j + \sum_{l=j+1}^{k-1} (u'_l - u'_{l-1})\Bigr) = -c_2^{(r)}\, u'_{k-1} .
\end{equation*}

\emph{Second sum.}
By hypothesis (d) at level $l$ we have $|\mathsf{m}_r(l+1)| \le C_{\mathsf{m}} (\log x_l)^{2p_0}$.
Combined with $|u_l - u'_l| \le 2c_2 u_l^2$ this gives
\begin{equation*}
\Bigl| \sum_{l=j}^{k-1} \mathsf{m}_r(l+1)\,(u_l - u'_l) \Bigr|
\le 2c_2 \sum_{l=j}^{k-1} C_{\mathsf{m}} (\log x_l)^{2p_0}\, u_l^2 \le C .
\end{equation*}

\emph{Third sum.}
Use $\beta^{0}_{l+1} = c_0 + (s_l - s_{l+1})$, from hypothesis (b), to split
\begin{equation*}
\kappa \sum_{l=j}^{k-1} \beta^{0}_{l+1}\, u'_l
= \kappa c_0 \sum_{l=j}^{k-1} u'_l + \kappa \sum_{l=j}^{k-1} (s_l - s_{l+1})\, u'_l .
\end{equation*}
For the last sum, summation by parts gives
\begin{equation*}
\sum_{l=j}^{k-1} (s_l - s_{l+1})\, u'_l
= s_j u'_j + \sum_{l=j+1}^{k-1} s_l\,(u'_l - u'_{l-1}) - s_k u'_{k-1} .
\end{equation*}
Since $|s_l| \le c_2$ by (b) and $u'$ is increasing, its modulus is at most
$2c_2 u'_{k-1} \le 2c_2 u'_k$.

For the first sum, hypothesis (b) gives
\begin{equation*}
c_0 u'_l = (x'_l - x'_{l+1})\, u'_l - \bigl(b_{l+1}(0) - \beta^{0}_{l+1}\bigr)\, u'_l ,
\end{equation*}
and by (c)
\begin{equation*}
\sum_{l=j}^{k-1} \bigl|b_{l+1}(0) - \beta^{0}_{l+1}\bigr|\, u'_l
\le \sum_{l=j}^{k-1} (C_\beta e_l + v_l)\, u'_l \le C .
\end{equation*}

It remains to compare $\sum_l (x'_l - x'_{l+1})\, u'_l$ with a logarithm. Set
$t_l := (x'_l - x'_{l+1})\, u'_l$. Since $x'_{l+1} > 0$ and $x'_l - x'_{l+1} > 0$ by (b), we have
$0 < t_l < 1$ and $1 - t_l = x'_{l+1}/x'_l$, so that $\log(x'_l/x'_{l+1}) = -\log(1 - t_l)$. For
$0 \le t < 1$ the number $-\log(1-t) - t = \sum_{n \ge 2} t^n/n$ lies between $0$ and
$\sum_{n\ge 2} t^n = t^2/(1-t)$. Here
\begin{equation*}
\frac{t_l^2}{1 - t_l} = (x'_l - x'_{l+1})^2\, u'_l\, u'_{l+1}
\le (3\lambda)^2\, u'_l\, u'_{l+1} \le \frac{(3\lambda)^2}{2}\bigl(u'^{\,2}_l + u'^{\,2}_{l+1}\bigr) .
\end{equation*}
Summing over $j \le l \le k-1$ gives
\begin{equation*}
\Bigl| \sum_{l=j}^{k-1} \frac{x'_l - x'_{l+1}}{x'_l} - \log\frac{x'_j}{x'_k} \Bigr|
\le \sum_{l=j}^{k} (3\lambda)^2\, u'^{\,2}_l \le C .
\end{equation*}
Finally, $\log(x'_j/x'_k)$ and $\log(x_j/x_k)$ differ by
$\log(u'_k/u_k) - \log(u'_j/u_j)$. This is bounded because
$|u'_l/u_l - 1| \le 2c_2 u_l$.

\emph{Conclusion.}
Collecting the three bounds in \eqref{10.4:eq-normalisation-decay-3}, we obtain
\begin{equation*}
\sum_{l=j}^{k-1} \mathsf{m}_r(l+1)\, u_l
\ge \kappa \log(x_j/x_k) - c_2^{(r)} u'_{k-1} - C - |\kappa|\, C .
\end{equation*}
Since $u'$ is increasing, $u'_{k-1} \le u'_k \le u'_K$, and by (b)
$u'_K \le u_K + 2c_2 u_K^2$ with $u_K = g_K^2$; so the term $c_2^{(r)} u'_{k-1}$, like the bound
$2|\kappa| c_2 u'_k$ used in the third sum, is at most a constant multiple of
$g_K^2 + 2c_2 g_K^4$, which enters $C$.
In the three bounds, $\kappa$ enters only through the main term $\kappa \log(x_j/x_k)$ and
through terms estimated in absolute value. This is \eqref{10.4:eq-normalisation-decay-2}, for
$\kappa$ of either sign.

Inserting \eqref{10.4:eq-normalisation-decay-2} and the bound on the remainders $\varrho^{r}_l$
into \eqref{10.4:eq-normalisation-decay-1} gives
\begin{equation*}
\log \bigl|P^{r}_k(0)/P^{r}_j(0)\bigr| \le \kappa \log(x_k/x_j) + C .
\end{equation*}
Exponentiating yields the assertion with $C_N := e^{C}$. All the constants above are independent
of $K$.
\end{proof}

\begin{remark}[The cut-free one-loop numbers]\label{10.4:rem-cut-free-numbers}
Stated; proof incomplete at the step marked $(\ast)$.

Write $\chi_T$ for the cut at height $T$. This is the factor restricting the fluctuation fields to size below $T$ in the normalised Gaussian average $d\mu_{C^{(k)}}$ of the step from level $k$ to level $k+1$.

The one-loop number $\mathsf{m}_r(k+1)$ enters the stress-tensor-insertion correlation theorem through \eqref{10.4:eq-insertion-correlation-b}. It depends on the cut $T_k$ of that step.

Its cut-free value $\mathsf{m}^{\circ}_r(k+1)$ is the same number computed with $d\mu_{C^{(k)}}$ alone, without the factor $\chi_{T_k}$. This value does not depend on the coupling $g$.

There are constants $C$ and $c$ such that
\begin{equation}\label{10.4:eq-cut-free-numbers-1}
\bigl|\mathsf{m}_r(k+1)-\mathsf{m}^{\circ}_r(k+1)\bigr|
\le C\,e^{-cT_k^2}\,(\log x_k)^{2p_0}.
\end{equation}

Granted \eqref{10.4:eq-cut-free-numbers-1}, the one-sided inequality \eqref{10.4:eq-normalisation-decay-a} of Lemma~\ref{10.4:lem-normalisation-decay} may be read either for $\mathsf{m}_r$ or for $\mathsf{m}^{\circ}_r$.
\end{remark}

\begin{proof}
Let $T$ and $T'$ be two cut heights. In the expansions \cite[(2.3), (2.22)]{B-CMP116} we replace the cut-off factor $\chi_T$ by the interpolated factor
\begin{equation*}
\chi_{T'}+\lambda\,(\chi_T-\chi_{T'}),
\end{equation*}
where $\lambda$ is a complex parameter.
This factor equals $\chi_{T'}$ at $\lambda=0$ and equals $\chi_T$ at $\lambda=1$. The one-loop number computed with it is therefore a function of $\lambda$, and this function interpolates between the two cuts.

The bound \eqref{10.4:eq-cut-free-numbers-1} is to be obtained $(\ast)$ by regarding this one-loop number as a function of $\lambda$ on the disc
\begin{equation*}
|\lambda|\le e^{\frac14\gamma_2 T'^2}
\end{equation*}
and applying the Schwarz lemma in $\lambda$ there; here $\gamma_2$ is the constant so denoted in \cite{B-CMP116}.

Three things are not established here. The first is that the hypotheses of the Schwarz lemma hold on this disc for the function of $\lambda$ so defined. The second is the passage from that application to \eqref{10.4:eq-cut-free-numbers-1}. The third is which of the two heights $T$, $T'$ is the cut $T_k$, and how the value $\mathsf{m}^{\circ}_r(k+1)$, computed without any cut, is reached from the interpolation between two finite cuts.
\end{proof}

\begin{lemma}[Terminal matching for linear recursions]\label{10.4:lem-terminal-matching}
Let $N \ge 1$ be an integer, the index of the last level of the recursion (in this lemma the letter
does not denote the rank of $SU(N)$). Let
$\mathfrak{g}_0,\dots,\mathfrak{g}_N$ be positive numbers and $c_{\mathfrak{g}}$ a constant with
$\mathfrak{g}_k \le c_{\mathfrak{g}}\,\mathfrak{g}_l$ whenever $k \le l \le N$. Let
$\vartheta_1 \in \left]0,1\right[$, let $C_\Pi > 0$, and let $\mathsf{k}_{k,j} \ge 0$
($0 \le j < k \le N$) be a kernel, with
\begin{equation*}
\varkappa^\sharp := \sup_{1 \le k \le N} \sum_{j<k} \mathsf{k}_{k,j}\,\vartheta_1^{-(k-j)} .
\end{equation*}
A \emph{solution with sources} $(\sigma_k)_{0 \le k \le N}$ is a pair of sequences
$(e_k)_{0\le k\le N} \subset [0,\infty[$ and $(\mu_k)_{0 \le k \le N} \subset \mathbb{C}$ with
$e_0 \le \sigma_0$ and, for $1 \le k \le N$,
\begin{equation}\label{10.4:eq-terminal-matching-a}
e_k \le \sum_{j<k} \mathsf{k}_{k,j}\bigl(e_j + \mathfrak{g}_j|\mu_j|\bigr) + \sigma_k,
\qquad
|\mu_k - \mu_{k-1}| \le C_\Pi\,\mathfrak{g}_k^{-1} e_k .
\end{equation}
Put $c_\mu := 2C_\Pi c_{\mathfrak{g}}\vartheta_1/(1-\vartheta_1)$ and assume
$\varkappa^\sharp(1+c_\mu) \le \tfrac12$. Let $\mathsf{b}_0,\mathsf{b}_1$ be real numbers and let
$(e_k),(\mu_k)$ be a solution with sources
$(\sigma_k)$ such that $\sigma_0 \le \mathsf{b}_0$, $\sigma_k \le \mathsf{b}_1\vartheta_1^k$ for
$k \ge 1$, and $\mu_N = 0$. Then, with $\mathsf{b} := \mathsf{b}_1 + \varkappa^\sharp\mathsf{b}_0$,
\begin{equation}\label{10.4:eq-terminal-matching-1}
e_k \le 2\mathsf{b}\,\vartheta_1^k \quad (1 \le k \le N),
\qquad
\mathfrak{g}_k|\mu_k| \le c_\mu\mathsf{b}\,\vartheta_1^k \quad (0 \le k \le N).
\end{equation}
Moreover:
\begin{enumerate}
\item[(i)] a solution with sources as above (with $\mu_N$ not prescribed) such that $\mu_0 = 1$ and
$c_\mu\mathsf{b} < \mathfrak{g}_0$ has $\mu_N \neq 0$;
\item[(ii)] consider a linear system of directions: a complex vector space of directions $\delta$,
each carrying a solution $(e_k[\delta]),(\mu_k[\delta])$ of
\eqref{10.4:eq-terminal-matching-a}, with the same $N$, $\mathfrak{g}_k$, $c_{\mathfrak{g}}$,
$\vartheta_1$, $C_\Pi$ and $\mathsf{k}_{k,j}$, with sources $(\sigma_k[\delta])$, such that each
$\mu_k[\cdot]$ is linear and the sources of $\delta - C\delta'$ are at most those of $\delta$ plus
$|C|$ times those of $\delta'$. For a direction $\delta$ write $\mathsf{b}_0[\delta]$,
$\mathsf{b}_1[\delta]$ for numbers with $\sigma_0[\delta] \le \mathsf{b}_0[\delta]$ and
$\sigma_k[\delta] \le \mathsf{b}_1[\delta]\vartheta_1^k$ ($k \ge 1$), and
$\mathsf{b}[\delta] := \mathsf{b}_1[\delta] + \varkappa^\sharp\mathsf{b}_0[\delta]$. Let
$\delta^c,\delta^T$ be two directions with $\mu_0[\delta^T] = 1$,
$c_\mu\mathsf{b}[\delta^T] \le \tfrac12\mathfrak{g}_0$ and $\mu_0[\delta^c] = 0$. Then
$C := \mu_N[\delta^c]/\mu_N[\delta^T]$ is defined,
$|C| \le 2c_\mu\mathsf{b}[\delta^c]/\mathfrak{g}_0$, and $\theta := \delta^c - C\delta^T$ satisfies
\begin{equation*}
e_k[\theta] \le 4\mathsf{b}[\delta^c]\,\vartheta_1^k \quad (k \ge 1),
\qquad
\mathfrak{g}_k|\mu_k[\theta]| \le 2c_\mu\mathsf{b}[\delta^c]\,\vartheta_1^k .
\end{equation*}
\end{enumerate}
\end{lemma}

\begin{proof}
Since there are finitely many levels, the number
\begin{equation*}
D := \max_{1 \le k \le N} e_k\vartheta_1^{-k}
\end{equation*}
is finite, and $e_l \le D\vartheta_1^l$ for $1 \le l \le N$.

We first bound the marginal slot by $D$. Because $\mu_N = 0$, for every $0 \le k \le N$ the
telescoping identity
\begin{equation*}
\mu_k = \mu_k - \mu_N = -\sum_{l=k+1}^{N} (\mu_l - \mu_{l-1})
\end{equation*}
holds (the sum is empty for $k = N$). By the second row of \eqref{10.4:eq-terminal-matching-a},
$|\mu_l - \mu_{l-1}| \le C_\Pi\mathfrak{g}_l^{-1}e_l$, and for $l > k$ the hypothesis on the
$\mathfrak{g}$'s gives $\mathfrak{g}_k\mathfrak{g}_l^{-1} \le c_{\mathfrak{g}}$. Hence, using
$e_l \le D\vartheta_1^l$ for $l \ge k+1 \ge 1$ and summing the geometric series,
\begin{align*}
\mathfrak{g}_k|\mu_k|
&\le C_\Pi c_{\mathfrak{g}} \sum_{l=k+1}^{N} e_l
\le C_\Pi c_{\mathfrak{g}} D \sum_{l \ge k+1} \vartheta_1^{l}
= \frac{C_\Pi c_{\mathfrak{g}} D\,\vartheta_1^{k+1}}{1-\vartheta_1}
= \tfrac12 c_\mu D\,\vartheta_1^k
\end{align*}
for all $0 \le k \le N$.

We insert this in the first row of \eqref{10.4:eq-terminal-matching-a}. Fix $1 \le k \le N$ and
multiply that row by $\vartheta_1^{-k} = \vartheta_1^{-(k-j)}\vartheta_1^{-j}$; since
$\sigma_k\vartheta_1^{-k} \le \mathsf{b}_1$,
\begin{equation*}
e_k\vartheta_1^{-k}
\le \sum_{j<k}\mathsf{k}_{k,j}\vartheta_1^{-(k-j)}
\bigl[e_j\vartheta_1^{-j} + \mathfrak{g}_j|\mu_j|\vartheta_1^{-j}\bigr] + \mathsf{b}_1 .
\end{equation*}
In the bracket, $e_j\vartheta_1^{-j} \le D$ for $j \ge 1$, while for $j = 0$ we have
$e_0 \le \sigma_0 \le \mathsf{b}_0$; thus $e_j\vartheta_1^{-j} \le \max\{\mathsf{b}_0,D\}$ for all
$j < k$. By the preceding paragraph $\mathfrak{g}_j|\mu_j|\vartheta_1^{-j} \le \tfrac12 c_\mu D$.
The kernel sum is at most $\varkappa^\sharp$ by definition. Therefore, using
$\max\{\mathsf{b}_0,D\} \le \mathsf{b}_0 + D$ and then
$\varkappa^\sharp(1+\tfrac12 c_\mu) \le \varkappa^\sharp(1+c_\mu) \le \tfrac12$ (note
$c_\mu \ge 0$),
\begin{align*}
e_k\vartheta_1^{-k}
&\le \varkappa^\sharp\bigl[\max\{\mathsf{b}_0,D\} + \tfrac12 c_\mu D\bigr] + \mathsf{b}_1
\le \varkappa^\sharp\mathsf{b}_0 + \varkappa^\sharp\bigl(1+\tfrac12 c_\mu\bigr)D + \mathsf{b}_1\\
&\le \varkappa^\sharp\mathsf{b}_0 + \tfrac12 D + \mathsf{b}_1 = \mathsf{b} + \tfrac12 D .
\end{align*}
Taking the maximum over $1 \le k \le N$ gives $D \le \mathsf{b} + \tfrac12 D$, and since $D$ is
finite, $D \le 2\mathsf{b}$. This is the first bound in \eqref{10.4:eq-terminal-matching-1}; the
second follows from $\mathfrak{g}_k|\mu_k| \le \tfrac12 c_\mu D\vartheta_1^k$ and $D \le 2\mathsf{b}$.

(i) Suppose, to the contrary, that $\mu_N = 0$. Then all hypotheses of the first part hold, and the
second bound of \eqref{10.4:eq-terminal-matching-1} at $k = 0$, together with $\mu_0 = 1$, gives
$\mathfrak{g}_0 = \mathfrak{g}_0|\mu_0| \le c_\mu\mathsf{b}$, contradicting
$c_\mu\mathsf{b} < \mathfrak{g}_0$. Hence $\mu_N \neq 0$.

(ii) Since $\mathfrak{g}_0 > 0$, the hypothesis $c_\mu\mathsf{b}[\delta^T] \le \tfrac12\mathfrak{g}_0$
gives $c_\mu\mathsf{b}[\delta^T] < \mathfrak{g}_0$, and $\mu_0[\delta^T] = 1$; by (i),
$\mu_N[\delta^T] \neq 0$, so $C$ is defined. The element $\theta = \delta^c - C\delta^T$ is a
direction of the system. By linearity of $\mu_k[\cdot]$,
\begin{equation*}
\mu_N[\theta] = \mu_N[\delta^c] - C\mu_N[\delta^T] = 0,
\qquad
\mu_0[\theta] = \mu_0[\delta^c] - C\mu_0[\delta^T] = -C .
\end{equation*}
The sources of $\theta$ are at most those of $\delta^c$ plus $|C|$ times those of $\delta^T$, so
$\theta$ satisfies the source hypotheses with
$\mathsf{b}_0[\theta] := \mathsf{b}_0[\delta^c] + |C|\mathsf{b}_0[\delta^T]$ and
$\mathsf{b}_1[\theta] := \mathsf{b}_1[\delta^c] + |C|\mathsf{b}_1[\delta^T]$, whence
$\mathsf{b}[\theta] = \mathsf{b}[\delta^c] + |C|\mathsf{b}[\delta^T]$. Since $\mu_N[\theta] = 0$, the
first part applies to $\theta$; its second bound at $k = 0$, with
$c_\mu\mathsf{b}[\delta^T] \le \tfrac12\mathfrak{g}_0$, gives
\begin{equation*}
\mathfrak{g}_0|C| = \mathfrak{g}_0|\mu_0[\theta]|
\le c_\mu\bigl(\mathsf{b}[\delta^c] + |C|\mathsf{b}[\delta^T]\bigr)
\le c_\mu\mathsf{b}[\delta^c] + \tfrac12\mathfrak{g}_0|C| ,
\end{equation*}
that is, $|C| \le 2c_\mu\mathsf{b}[\delta^c]/\mathfrak{g}_0$. Inserting this bound and again
$c_\mu\mathsf{b}[\delta^T] \le \tfrac12\mathfrak{g}_0$,
\begin{equation*}
\mathsf{b}[\theta] = \mathsf{b}[\delta^c] + |C|\mathsf{b}[\delta^T]
\le \mathsf{b}[\delta^c]
+ \frac{2\mathsf{b}[\delta^c]}{\mathfrak{g}_0}\,c_\mu\mathsf{b}[\delta^T]
\le 2\mathsf{b}[\delta^c] .
\end{equation*}
The bounds \eqref{10.4:eq-terminal-matching-1} for $\theta$ then read
\begin{gather*}
e_k[\theta] \le 2\mathsf{b}[\theta]\vartheta_1^k \le 4\mathsf{b}[\delta^c]\vartheta_1^k
\quad (k \ge 1),\\
\mathfrak{g}_k|\mu_k[\theta]| \le c_\mu\mathsf{b}[\theta]\vartheta_1^k
\le 2c_\mu\mathsf{b}[\delta^c]\vartheta_1^k .
\end{gather*}
\end{proof}

\begin{remark}
The rows \eqref{10.4:eq-terminal-matching-a} are the form in which the typed tangent rows enter: $e_k$
is the flat norm of the raw typed family produced at level $k$, and $\mu_k - \mu_{k-1}$ is its typed
marginal coefficient $\beta^{r}$, so that the second row is the flat bound of
Lemma~\ref{10.4:lem-coefficient-bounds}; the kernel $\mathsf{k}_{k,j}$ collects the products
$\varkappa_k\varsigma_\nu$ and the age kernel acting through the remnants and their re-filing.
Written in the re-filed transport quantities $\hat e_k$, $\hat q_k$, $r_k$, $p^\flat_k$ of
Lemma~\ref{10.4:lem-forward-transport}, the typed step with a marginal tensor slot
(Proposition~\ref{10.4:prop-typed-step}, whose proof is incomplete at one step), in the form
\eqref{10.4:eq-typed-step-a}, together with the re-filing of typed marginal families
(Lemma~\ref{10.4:lem-typed-refiling}, whose proof is incomplete at one step), in the form
\eqref{10.4:eq-typed-refiling-a}, give \eqref{10.4:eq-terminal-matching-a} with
$\varkappa^\sharp$ at most a constant multiple of $\bar\varkappa$, the constant depending on $S_1$,
$C_{\mathsf{b}}$, $C_q'$ and $C_{RF}$: the slot $r_k$ is eliminated by its bound among the rows
\eqref{10.4:eq-forward-transport-b}, in the same way as in the proof of part~(i) of the forward
transport of the tangent rows (Lemma~\ref{10.4:lem-forward-transport}).

In Lemma~\ref{10.4:lem-terminal-matching} the marginal slot is solved backwards from the last level,
through the terminal condition $\mu_N = 0$, while the irrelevant slots are transported forwards from
level zero, as in part~(b) of the volume lemma for the untyped tangent rows. No bound on
$\mu_k[\delta^T]$, the normalisation, is used in either direction; only the
fact that one step moves it by at most $C_\Pi\mathfrak{g}_k^{-1}e_k$, which is the second row of
\eqref{10.4:eq-terminal-matching-a}. In particular, part~(ii) subtracts the marginal part of
$\delta^c$ by a multiple of $\delta^T$ without any assumption that the normalisation
$\mu_k[\delta^T]$ stays bounded or bounded away from zero along the run.
\end{remark}

\begin{lemma}[The trading coefficient in absolute units]\label{10.4:lem-trading-coefficient}
\begin{enumerate}
\item[(a)] Let $N\ge 1$, $\omega\in\,]0,1[$ and $C_\beta>0$. Let
$(p_k)_{0\le k\le N}\subset[0,\infty[$ and $(\mu_k)_{0\le k\le N}\subset\mathbb{C}$. Let
$\varpi_a\ge 0$ $(a\ge 0)$ and $a^{\natural}_k,c^{\otimes}_k\ge 0$ $(0\le k\le N-1)$. Let
$\pi_k,\mathsf{y}_k\ge 0$ $(1\le k\le N)$, and let $\mathsf{f},\Pi_\omega,Y_\omega$ be constants.
Let $S_\omega$ be a constant with $\sum_{a\ge 0}\varpi_a\omega^{-a}\le S_\omega$. Suppose that
\begin{equation}\label{10.4:eq-trading-coefficient-1}
\begin{gathered}
p_0\le\mathsf{f},\qquad
p_{k+1}\le a^{\natural}_k\sum_{j=0}^{k}\varpi_{k-j}\,p_j+c^{\otimes}_k|\mu_k|+\pi_{k+1},\\
|\mu_{k+1}-\mu_k|\le C_\beta\,p_{k+1}+\mathsf{y}_{k+1}\qquad(0\le k\le N-1),\qquad \mu_N=0,
\end{gathered}
\end{equation}
and that $\pi_k\le\Pi_\omega\,\omega^k$ and $\mathsf{y}_k\le Y_\omega\,\omega^k$ for $1\le k\le N$. All suprema
over $k$ below are taken over $0\le k\le N-1$. Suppose further that
\begin{equation*}
\omega^{-1}\sup_k a^{\natural}_k\,S_\omega\le\tfrac14,\qquad
\sup_k c^{\otimes}_k\,C_\beta\,(1-\omega)^{-1}\le\tfrac14 .
\end{equation*}
Then
\begin{equation}\label{10.4:eq-trading-coefficient-2}
|\mu_0|\le 4(1-\omega)^{-1}\bigl[C_\beta\bigl(a^{\natural}_{(0)}\mathsf{f}+\omega\Pi_\omega\bigr)
+\omega Y_\omega\bigr],\qquad
a^{\natural}_{(0)}:=\sup_k a^{\natural}_k\,\varpi_k\,\omega^{-k}.
\end{equation}
\item[(b)] There are constants $\underline{x}_0$ and $C'$ such that the following holds for
$x_0\ge\underline{x}_0$. Consider a bare typed unit without marginal part, of absolute norm
$\mathsf{f}$ on $\mathcal{U}_0$. Let $C=\mu_N[\delta^c]/\mu_N[\delta^T]$ be its trading coefficient, with
$\delta^c$ the direction of the unit and $\delta^T$ a direction with $\mu_0[\delta^T]=1$, as in
consequence (ii) of the terminal matching lemma (Lemma~\ref{10.4:lem-terminal-matching}). Then
\begin{equation*}
|C|\le C'g_0^2(\log x_0)^{2p_0}\,\mathsf{f},
\end{equation*}
where the exponent $p_0$ is the fixed exponent of the degree function $p_0(\cdot)$, not the initial
member of the sequence $(p_k)$ of (a).
\end{enumerate}
\end{lemma}

Stated; the proof below is incomplete at the step marked $(\ast)$.

\begin{proof}
\emph{Proof of} (a). Put
\begin{equation*}
P^{\ast}:=\max_{1\le k\le N}p_k\,\omega^{-k},\qquad M^{\ast}:=\max_{0\le k\le N}|\mu_k|\,\omega^{-k}.
\end{equation*}
Both maxima are finite, since they are taken over finitely many indices.

\emph{The bound on $M^{\ast}$.} Fix $0\le k\le N$. Since $\mu_N=0$, the telescoping identity gives
$\mu_k=-\sum_{l=k+1}^{N}(\mu_l-\mu_{l-1})$. Every index $l$ in this sum satisfies $1\le l\le N$. For
such $l$, the third inequality of \eqref{10.4:eq-trading-coefficient-1} (with $k$ replaced by $l-1$),
together with $p_l\le P^{\ast}\omega^l$ and $\mathsf{y}_l\le Y_\omega\omega^l$, gives
$|\mu_l-\mu_{l-1}|\le(C_\beta P^{\ast}+Y_\omega)\,\omega^l$. Summing the geometric series over $l>k$
yields $|\mu_k|\le(C_\beta P^{\ast}+Y_\omega)\,\omega^{k+1}(1-\omega)^{-1}$. Multiplying by
$\omega^{-k}$ and taking the maximum over $k$, we obtain
\begin{equation}\label{10.4:eq-trading-coefficient-3}
M^{\ast}\le\omega(1-\omega)^{-1}\bigl(C_\beta P^{\ast}+Y_\omega\bigr).
\end{equation}

\emph{The bound on $P^{\ast}$.} Fix $0\le k\le N-1$ and divide the second inequality of
\eqref{10.4:eq-trading-coefficient-1} by $\omega^{k+1}$. We bound the four resulting terms in turn.

The summand $j=0$ is at most
$\omega^{-1}a^{\natural}_k\varpi_k\omega^{-k}\,\mathsf{f}\le\omega^{-1}a^{\natural}_{(0)}\mathsf{f}$,
because $p_0\le\mathsf{f}$.

For the summands $1\le j\le k$ we use $p_j\le P^{\ast}\omega^j$. Substituting $a=k-j$, they contribute
at most
\begin{equation*}
\omega^{-1}a^{\natural}_k P^{\ast}\sum_{j=1}^{k}\varpi_{k-j}\,\omega^{-(k-j)}
\le\omega^{-1}a^{\natural}_k S_\omega P^{\ast}\le\tfrac14 P^{\ast} .
\end{equation*}
Here the first inequality uses $\sum_{a\ge0}\varpi_a\omega^{-a}\le S_\omega$, and the second uses the
first condition of (a).

The $\mu$-term satisfies $c^{\otimes}_k|\mu_k|\,\omega^{-(k+1)}\le\omega^{-1}c^{\otimes}_kM^{\ast}$. By
\eqref{10.4:eq-trading-coefficient-3} and the second condition of (a), which says that
$c^{\otimes}_k(1-\omega)^{-1}\le 1/(4C_\beta)$, we get
\begin{equation*}
\omega^{-1}c^{\otimes}_k M^{\ast}\le c^{\otimes}_k(1-\omega)^{-1}\bigl(C_\beta P^{\ast}+Y_\omega\bigr)
\le\tfrac14 P^{\ast}+\tfrac14 Y_\omega/C_\beta .
\end{equation*}

Finally, $\pi_{k+1}\omega^{-(k+1)}\le\Pi_\omega$.

Adding the four bounds gives, for $0\le k\le N-1$,
\begin{equation}\label{10.4:eq-trading-coefficient-4}
p_{k+1}\,\omega^{-(k+1)}\le\omega^{-1}a^{\natural}_{(0)}\mathsf{f}+\tfrac12 P^{\ast}
+\tfrac14 Y_\omega/C_\beta+\Pi_\omega .
\end{equation}
The maximum over $k$ of the left side is $P^{\ast}$. Because $P^{\ast}$ is finite, the term
$\tfrac12 P^{\ast}$ can be absorbed into the left side, and
\begin{equation*}
P^{\ast}\le 2\bigl[\omega^{-1}a^{\natural}_{(0)}\mathsf{f}+\Pi_\omega+\tfrac14 Y_\omega/C_\beta\bigr].
\end{equation*}

\emph{Conclusion.} Inserting this into \eqref{10.4:eq-trading-coefficient-3} gives
\begin{align*}
|\mu_0|\le M^{\ast}
&\le\omega(1-\omega)^{-1}\bigl(2C_\beta\omega^{-1}a^{\natural}_{(0)}\mathsf{f}+2C_\beta\Pi_\omega
+\tfrac12 Y_\omega+Y_\omega\bigr)\\
&=(1-\omega)^{-1}\bigl[2C_\beta a^{\natural}_{(0)}\mathsf{f}+2C_\beta\omega\Pi_\omega
+\tfrac32\omega Y_\omega\bigr].
\end{align*}
All terms in the bracket are nonnegative: $\Pi_\omega\ge\pi_k\omega^{-k}\ge0$ and
$Y_\omega\ge\mathsf{y}_k\omega^{-k}\ge0$. Since $2\le4$ and $\tfrac32\le4$, this implies
\eqref{10.4:eq-trading-coefficient-2}.

\emph{Proof of} (b). Step $(\ast)$. The instantiation of (a) --- the identification of $p_j$,
$\mu_k$, $\pi_k$, $\mathsf{y}_k$ with the quantities of the typed tangent run, whose $p$-row is
\eqref{10.4:eq-pure-recursion-tangent-a}, and the passage from the bound on $|\mu_0|$ to a bound on
$|C|$ --- is not given in this proof; the paragraphs below fix the constants it would use.

\emph{The constant $a^{\natural}_k$.} For $0\le j\le k$ let $D_{H^{(j)}}\Phi_{k+1}$ denote the
derivative of the step's family in the level-$j$ slot of the history, and $\|\cdot\|_{j\to k+1}$ its
operator norm from the level-$j$ norm to the level-$(k+1)$ norm, as in the bound on the history
derivatives of the step's family. We rely on Proposition~\ref{10.4:prop-vertex-response}(b), whose
proof is incomplete at one step. By it, $D_h\Phi_{k+1}(\mathfrak{q},s;h)$ is even in $\mathfrak{q}$ and
vanishes at $\mathfrak{q}=0$. In the argument for the bound on the history derivatives, let
$F(\mathfrak{q})$ denote this derivative as a holomorphic function of $\mathfrak{q}$; Schwarz's lemma
may therefore be applied to $F(\sqrt v)$ in the variable $v=\mathfrak{q}^2$. This gives
\begin{equation*}
\bigl\|D_{H^{(j)}}\Phi_{k+1}(g_k,s;h)\bigr\|_{j\to k+1}\le a^{\natural}_k\,\varpi_{k-j},\qquad
a^{\natural}_k:=\frac{E^{*}\tau_k^2}{2\mathsf{E}}\asymp g_k^2(\log x_k)^{2p_0},
\end{equation*}
with $\tau_k$ the radius ratio of the correlation theorem, as in the definition of $c^{\otimes}_k$ in
Proposition~\ref{10.4:prop-vertex-response}(c).
The constant $a_k$ of that bound carries $\tau_k$ in place of $\tau_k^2$.

\emph{The floor on $x_0$.} By the reference comparison of the running couplings,
$a^{\natural}_k/a^{\natural}_0\le2$ for $k\le x_0/(6\lambda)$, and $a^{\natural}_k/a^{\natural}_0\le x_0g^2$
for all $k$. Moreover the slot weights are $\varpi_a=L^{-\beta_1 a}$, with $\beta_1$ their decay
exponent, so that $\varpi_k\omega^{-k}=(L^{-\beta_1}/\omega)^k$. From these,
$a^{\natural}_{(0)}\le 3a^{\natural}_0$ for $x_0\ge\underline{x}_0$.

\emph{The bound on $|C|$.} Accordingly, a bare typed unit without marginal part, of absolute norm
$\mathsf{f}$ on $\mathcal{U}_0$, would have $|C|\le C'g_0^2(\log x_0)^{2p_0}\mathsf{f}$.

\emph{The $\flat$-norms.} Here the inequalities \eqref{10.4:eq-terminal-matching-a} satisfied by the
$\flat$-norms give only a contribution $2c_\mu\varkappa^\sharp\mathsf{b}_0/\mathfrak{g}_0$, where
$\mathsf{b}_0$ is the $\flat$-norm of the unit. Indeed, Lemma~\ref{10.4:lem-terminal-matching} with
$\mathsf{b}_1=0$ has $\mathsf{b}=\varkappa^\sharp\mathsf{b}_0$, and its consequence (ii) bounds the
coefficient by $2c_\mu\mathsf{b}/\mathfrak{g}_0$.

The slot $j=0$ of the pure history is the bare-slot bound at level zero, written in the letters of the
pure recursion; it is part of \eqref{10.4:eq-typed-slot-kernel-a}.

With the constant $a_k$ of the history-derivative bound in place of $a^{\natural}_k$, the bound is
$C'g_0(\log x_0)^{p_0}\mathsf{f}$, and the exponents in the table of exponents of the typed directions
therefore do not depend on the evenness.
\end{proof}

\begin{corollary}[The bound for the full symmetric-tensor source]\label{10.4:cor-symmetric-tensor-bound}
In the setting of Theorem~\ref{10.4:thm-insertion-correlation}, in particular for
$g\le\gamma^{\otimes}$ and with $x_0\ge X+\lambda K-2c_2$, let $n\ge 1$, let $\mathsf{k}$ be a
Lipschitz function on the torus of Definition~\ref{10.4:def-tensor-sources} with values in
$\operatorname{Sym}^2 V$, where $V=\mathbb{R}^4$, and let
$f_1,\dots,f_n$ be test functions such that the supports of $\mathsf{k},f_1,\dots,f_n$ are pairwise at
distance at least $\mathsf{d}>0$. Assume:
\begin{itemize}
\item the decomposition \eqref{10.4:eq-tensor-sources-a} of $\mathfrak{T}^{\otimes}(\mathsf{k})$ into
its channel pieces (Definition~\ref{10.4:def-tensor-sources});
\item the terminal-matching lemma, Lemma~\ref{10.4:lem-terminal-matching}, and the trading coefficient
lemma, Lemma~\ref{10.4:lem-trading-coefficient}, whose proof is incomplete at one step;
\item the torus bound of the correlation theorem;
\item the bound \eqref{10.4:eq-insertion-correlation-d} of the stress-tensor-insertion correlation
theorem, Theorem~\ref{10.4:thm-insertion-correlation};
\item the normalisation decay \eqref{10.4:eq-tensor-correlation-decay-a} for $r=\mathbf{3}$ and for
$r=\mathbf{6}$, with one and the same $\kappa\in\mathbb{R}$;
\item in the case $\kappa<0$, the bound $x_0\le C_*K\,x_{K'}$ with a constant $C_*$ independent of $K$.
\end{itemize}
Then
\begin{equation*}
\bigl|\langle \mathfrak{T}^{\otimes}(\mathsf{k});\mathcal{O}(f_1);\dots;\mathcal{O}(f_n)\rangle_{K,m}\bigr|
\le C_n(\mathsf{d})\,K^{\max\{-\kappa,0\}}\,\|\mathsf{k}\|_\sharp\prod_{i=1}^n\|f_i\|_\sharp
\,\bigl(1+o(1)\bigr).
\end{equation*}
Here $C_n(\mathsf{d})$ denotes a constant depending only on $L$, $\mathfrak{p}$, $g$, $m$, $n$,
$\mathsf{d}$, $\vartheta_1$, $q$, $c_\natural$, $C_N$, $\kappa$ and, for $\kappa<0$, $C_*$, and not on
$K$, $g_0$ or the position of the supports.
\end{corollary}

\begin{proof}
Write $\mathsf{k}=\tfrac14(\operatorname{tr}\mathsf{k})\,\delta+\mathsf{k}^{\mathbf{3}}+\mathsf{k}^{\mathbf{6}}$,
with $\delta$ the identity of $V$ and $\mathsf{k}^{r}$ taking values in $\mathsf{V}_r$. By
\eqref{10.4:eq-tensor-sources-a}, $\mathfrak{T}^{\otimes}(\mathsf{k})$ is the sum of four terms:
$\mathfrak{T}^{\mathbf{6}}(\mathsf{k}^{\mathbf{6}})$; $2\,\mathfrak{T}^{\mathbf{3}}(\mathsf{k}^{\mathbf{3}})$;
the bare units $-\sum_\nu\mathfrak{Q}_\nu(\mathsf{k}_{\nu\nu})$, which have no marginal part; and a last
sum of bare units carrying an explicit factor of the square of the lattice spacing. The trace part of
$\mathsf{k}$ cancels identically at the marginal level. The connected correlation is linear in its
first argument, so it suffices to estimate the four resulting correlations.

The pieces $\mathsf{k}^{\mathbf{3}},\mathsf{k}^{\mathbf{6}}$ are obtained from $\mathsf{k}$ by fixed
linear projections of $\operatorname{Sym}^2V$ applied pointwise. Hence their supports lie in that of
$\mathsf{k}$, the separation hypothesis holds for them, and
$\|\mathsf{k}^{r}\|_\sharp\le C\|\mathsf{k}\|_\sharp$ with $C$ depending only on these projections. By
\eqref{10.4:eq-insertion-correlation-d}, combined with \eqref{10.4:eq-tensor-correlation-decay-a} as in
Corollary~\ref{10.4:cor-tensor-correlation-decay}, for $r=\mathbf{3},\mathbf{6}$,
\begin{equation*}
\bigl|\langle \mathfrak{T}^{r}(\mathsf{k}^{r});\mathcal{O}(f_1);\dots;\mathcal{O}(f_n)\rangle_{K,m}\bigr|
\le C^{\otimes}_n C_N\Bigl(\frac{x_{K'}}{x_0}\Bigr)^{\kappa}\|\mathsf{k}^{r}\|_\sharp
\prod_{i=1}^n\|f_i\|_\sharp .
\end{equation*}
For $\kappa\ge 0$, that corollary, under $x_0\ge X+\lambda K-2c_2$, bounds the right-hand side by
$C(g,m)K^{-\kappa}\|\mathsf{k}^{r}\|_\sharp\prod_i\|f_i\|_\sharp$, and
$K^{-\kappa}\le 1=K^{\max\{-\kappa,0\}}$. For $\kappa<0$ the factor in the display equals
$(x_0/x_{K'})^{-\kappa}$, a positive power of the ratio of the initial inverse coupling to the inverse
coupling at the last sourced level. Since $-\kappa>0$, the hypothesis $x_0\le C_*K\,x_{K'}$, assumed
for this case, gives
\begin{equation*}
\Bigl(\frac{x_0}{x_{K'}}\Bigr)^{-\kappa}\le C_*^{-\kappa}\,K^{-\kappa}
=C_*^{-\kappa}\,K^{\max\{-\kappa,0\}},
\end{equation*}
so that in this case the right-hand side of the preceding display is at most
$C^{\otimes}_n C_N C_*^{-\kappa}K^{\max\{-\kappa,0\}}\|\mathsf{k}^{r}\|_\sharp\prod_i\|f_i\|_\sharp$.

The bare units $\mathfrak{Q}_\nu(\mathsf{k}_{\nu\nu})$ are traded by the terminal-matching lemma
(Lemma~\ref{10.4:lem-terminal-matching}) and the trading coefficient lemma
(Lemma~\ref{10.4:lem-trading-coefficient}, whose proof is incomplete at one step). Inside the
correlation they are replaced by
\begin{equation*}
C_{\mathfrak{Q}}\,\mathfrak{T}^{\mathbf{3}}(\cdot)+C'_{\mathfrak{Q}}\,\mathcal{O}(\cdot)
+O(\vartheta_1^{K'}),
\end{equation*}
where the argument in both terms is formed from the diagonal entries $\mathsf{k}_{\nu\nu}$ whose units
are traded, each of which, being an entry of $\mathsf{k}$, has $\sharp$-norm at most
$C\|\mathsf{k}\|_\sharp$, and where the coefficients $C_{\mathfrak{Q}},C'_{\mathfrak{Q}}$ so produced
satisfy
\begin{equation*}
|C_{\mathfrak{Q}}|+|C'_{\mathfrak{Q}}|\le C\,g_0^2(\log x_0)^{2p_0}.
\end{equation*}
The term $C_{\mathfrak{Q}}\,\mathfrak{T}^{\mathbf{3}}(\cdot)$ is estimated as in the preceding paragraph by
\eqref{10.4:eq-insertion-correlation-d}. The term $C'_{\mathfrak{Q}}\,\mathcal{O}(\cdot)$ is a connected
correlation of action-density observables and is estimated by the torus bound of the correlation
theorem. Since $x_0=g_0^{-2}$, the coefficient bound reads $C(\log x_0)^{2p_0}/x_0$, which tends to
zero as $x_0\to\infty$; the traded terms, which carry this coefficient, and the remainder
$O(\vartheta_1^{K'})$ enter the factor $1+o(1)$ in the statement. The last term of the decomposition,
\begin{equation*}
\tfrac14a^2\sum_\nu\sum_{p\ni\nu}(\Delta_\nu\mathsf{k}_{\nu\nu})(x(p))\,a_p,
\end{equation*}
with $a=L^{-K}$ the bare spacing and $a_p$ as in Definition~\ref{10.4:def-tensor-sources}, collects at
each plaquette $p$ the sum of the weights of the two directions spanning $p$; by the form of
$\mathfrak{T}^{\mathbf{3}}$ in Definition~\ref{10.4:def-tensor-sources} it is a diagonal tensor source
with weight $\tfrac14a^2\Delta_\nu\mathsf{k}_{\nu\nu}$. Since $a\le 1$, the function
$a^2\Delta_\nu\mathsf{k}_{\nu\nu}$ is a linear combination, with coefficients bounded independently of
$a$, of $\mathsf{k}_{\nu\nu}$ and its two translates by $\pm a$ in the direction $\nu$; translation
changes neither the supremum nor the Lipschitz constant, so the $\sharp$-norm of the weight is at most
$C\|\mathsf{k}\|_\sharp$. Its trace part is estimated by the torus bound of the correlation theorem and
its traceless-diagonal part by \eqref{10.4:eq-insertion-correlation-d}.
Adding the four estimates gives the assertion.
\end{proof}

\begin{remark}
The format statement accompanying this bound is the clause
\eqref{10.4:eq-insertion-correlation-c} of Theorem~\ref{10.4:thm-insertion-correlation}, read with the
diagonal normalisation $(N_{\mathbf{1}},N_{\mathbf{3}},N_{\mathbf{6}})$, where $N_{\mathbf{1}}$ is the
normalisation of the trace channel, whose marginal part cancels identically in
\eqref{10.4:eq-tensor-sources-a}.
\end{remark}

\begin{definition}[The vector package in output form]\label{10.5:def-vector-package-output}
Let $\xi=\omega x$ be an infinitesimal rotation field on $\mathbb{R}^4$. Let $\varphi\in C^2_c(\mathbb{R}^4)$, where $\varphi$ is either the cutoff function $\chi$ of the rotational Ward identity or one of the test functions $f_i$, and put $\eta:=\varphi\xi$. The \emph{strain} of $\eta$ is
\begin{equation}\label{10.5:eq-vector-package-output-1}
\mathsf{k}[\eta]_{\sigma\rho}:=\tfrac12\bigl(\partial_\sigma\eta^\rho+\partial_\rho\eta^\sigma\bigr).
\end{equation}
Its decomposition into the $W_4$-irreducible pieces of $\operatorname{Sym}^2V$ reads
\begin{equation}\label{10.5:eq-vector-package-output-2}
\mathsf{k}[\eta]_{\sigma\rho}
=\tfrac14(\xi\cdot\nabla\varphi)\,\delta_{\sigma\rho}
+\mathsf{k}^{\mathbf{3}}[\eta]_{\sigma\rho}+\mathsf{k}^{\mathbf{6}}[\eta]_{\sigma\rho},
\end{equation}
where $\delta_{\sigma\rho}$ is the Kronecker symbol. Let $c_0\in\{g_0^{-2},1\}$ and set
\begin{equation}\label{10.5:eq-vector-package-output-3}
\mathsf{v}_0(c_0):=c_0\,C\,C_0^{3}\,g_0^{3}\,(\log x_0)^{3q_\flat}.
\end{equation}
This is the size, per unit weight, of $c_0\mathfrak{r}^\sharp$ in the norm of the correlation theorem. Here $C$ is a constant, $C_0$ is Ba{\l}aban's auxiliary parameter, $x_0$ is the initial inverse coupling and $q_\flat$ is Ba{\l}aban's exponent.

Assume that the irrelevance exponent satisfies $\hat\alpha>\tfrac12$. Assume also that the parameter $q$ of the lemma on geometric forgetting of old regions satisfies $q\le L^{-2}$. Let $w=\sum_{l=1}^n z_lf_l$ with $n\ge1$ and $f_l\in C^2_c(\mathbb{R}^4)$. For a direction $\delta$, write $Z(z,t)[\delta]$ for the partition function with the complex source $w$, deformed by the parameter $t$ along $\delta$. Here $t$ is a local deformation parameter, and neither the letter $t=X-\hat x$ of the uniqueness theorem nor the age ratio of the one-lattice Lipschitz step. All $\flat$-norms below are read in the re-filed (hatted) quantities $\hat e_k,\hat q_k,\dots$ of the re-filing lemmas. Let $\delta_\eta$ denote the direction whose level-zero slot is $c_0\mathfrak{K}^\sharp_\eta[1]$, with $\mathfrak{K}^\sharp_\eta$ as in Lemma~\ref{10.4:lem-rotation-field}(ii).

The \emph{vector package in output form} is the following assertion. There exist $\vartheta_2\in(0,\vartheta_1]$ and constants $C_V$, $c_V$, all three independent of $K$, such that for every such $w$ there are directions $\delta^{\otimes}$ and $\delta^{\mathrm{rem}}$ with
\begin{equation}\label{10.5:eq-vector-package-output-a}
\partial_t\log Z(z,t)[\delta_\eta]
=\partial_t\log Z(z,t)\bigl[\delta^{\otimes}+\delta^{\mathrm{rem}}\bigr],
\end{equation}
and such that the following two properties hold.
\begin{enumerate}
\item[(i)] $\delta^{\otimes}$ is a $(\operatorname{Sym}^2V\oplus\Lambda^2)$-typed direction with zero level-zero slot and sources of type $E$ at the levels $j\ge1$ (in the slots $\delta E^{(j)}$). These sources are of two kinds.
\begin{itemize}
\item Sources with the weight $\partial\eta$. Their $\flat$-norm at level $j$ is at most $C_V\mathsf{v}_0(c_0)\vartheta_2^{\,j}$.
\item Sources with the weight $\eta^\rho\partial_\nu w$. Their $\flat$-norm at level $j$ is at most $C_V\mathsf{v}_0(c_0)\theta_j^2\mathfrak{g}_j\vartheta_2^{\,j}$.
\end{itemize}
\item[(ii)] If the supports of $f_1,\dots,f_n$ are separated by $\mathsf{d}$, then the mixed derivative at $z=0$ satisfies
\begin{equation}\label{10.5:eq-vector-package-output-4}
\Bigl|\,\partial_{z_1}\cdots\partial_{z_n}\,
\partial_t\log Z(z,t)[\delta^{\mathrm{rem}}]\bigr|_{z=0}\Bigr|
\le C(\varphi,\xi,\mathsf{d})\,\mathsf{v}_0(c_0)\,(\log x_0)^{c_V}\,\varepsilon_K
\prod_{l=1}^{n}\|f_l\|_{C^2},
\end{equation}
where $\varepsilon_K\to0$ as $K\to\infty$.
\end{enumerate}
\end{definition}

We now read these bounds in absolute units, that is, in the units of \eqref{10.4:eq-terminal-matching-a} after multiplication by $\mathfrak{g}_0^{-1}$. The sources induced by $g_0^{-2}\mathfrak{K}^\sharp$ then have the following coefficients.
\begin{itemize}
\item On the weight $\partial\eta$, the coefficients are at most of order $\mathsf{v}_0(g_0^{-2})/\mathfrak{g}_0$. This is comparable with $g_0^{-1}$ times a polylogarithmic factor.
\item On the weight $\eta\,\partial w$, the coefficients are at most of order $\mathsf{v}_0(g_0^{-2})\theta_0^2$. This is comparable with $g_0$ times a polylogarithmic factor.
\end{itemize}

Proof outstanding.

\medskip\noindent\textit{Route.}
A proof has to derive the assertion above from the bulk, second-order freezing and rim statements of the vector package. These statements rest on two results: the exact difference decomposition of the orbital remainder into $\Psi_\rho(p)-\Psi_\rho(p-\hat\rho)$ and the cubic part $\mathfrak{r}^\sharp_\rho$, and the double-subtraction lemma for Killing-weighted vector densities. Both are to be read in the re-filed (hatted) quantities $\hat e_k,\hat q_k,\dots$ of the re-filing lemmas, under the one-sided conditions $\hat\alpha>\tfrac12$ and $q\le L^{-2}$.

For the sources with the weight $\partial\eta$, the starting point is the step of the double-subtraction lemma that writes the vector density as a lattice divergence $\sum_\nu\nabla_\nu\Sigma_{\nu\rho}$ plus a remainder, where $\Sigma_{\nu\rho}$ is the local tensor, built from the curvature variables, that this step provides. This representation is read for the non-Killing weight $\eta$. Let $a_j$ denote the spacing of the level-$j$ lattice, put $\eta_j:=\eta/a_j$, and let $\nabla$ denote the difference operator of that lattice. At level $j$ the representation is summed by parts against $\eta_j$, which produces
\begin{equation*}
-\sum_y\bigl(\nabla_\nu\eta^\rho_j\bigr)\Sigma_{\nu\rho},
\qquad\text{with}\qquad \nabla\eta_j=\partial\eta .
\end{equation*}
The $\flat$-norm bound $C_V\mathsf{v}_0(c_0)\vartheta_2^{\,j}$ of (i) has to be established for these terms.

For the sources with the weight $\eta^\rho\partial_\nu w$, two contributions enter: the terms in which the summation by parts falls on $\hat w(y)$, and the first-order freezing term. The factor $\theta_j^2\mathfrak{g}_j$ in their bound has to be obtained from the second-difference bound, with one variation parameter on the vector unit and one on the source vertex. The load of the source vertex is $\mathfrak{g}_j|\delta w|$, where $\delta w$ is the variation of the source $w$ at that vertex.

Finally, the bound \eqref{10.5:eq-vector-package-output-4} has to be proved for the remaining direction $\delta^{\mathrm{rem}}$.

Throughout, $\xi=\omega x$ is an infinitesimal rotation of $\mathbb{R}^4$: we have $\xi^\rho(x)=\sum_\sigma\omega_{\rho\sigma}x_\sigma$ with $\omega_{\rho\sigma}=-\omega_{\sigma\rho}$. Hence $\partial_\sigma\xi^\rho=\omega_{\rho\sigma}$, and $\xi^\rho$ does not depend on $x_\rho$. For test functions $f_1,\dots,f_n$ and any functional $\Xi$ of the link variables we abbreviate
\begin{equation*}
\langle \Xi;\mathcal{O}(\vec f)\rangle:=\langle \Xi;\mathcal{O}(f_1);\dots;\mathcal{O}(f_n)\rangle .
\end{equation*}
We write $a_p:=1-\operatorname{Re}\operatorname{tr}U(\partial p)$, and $a=L^{-K}$ denotes the bare lattice spacing.

For $\varphi\in C^2_c(\mathbb{R}^4)$ and $\eta:=\varphi\xi$ let
\begin{equation*}
\mathsf{k}[\eta]_{\sigma\rho}:=\tfrac12\bigl(\partial_\sigma\eta^\rho+\partial_\rho\eta^\sigma\bigr)
\end{equation*}
be the strain of $\eta$. Since $\omega_{\rho\rho}=0$, we have $\operatorname{tr}\mathsf{k}[\eta]=\sum_\rho\partial_\rho(\varphi\xi^\rho)=\xi\cdot\nabla\varphi$. Decomposing according to Lemma~\ref{10.4:lem-hypercubic-decomposition}(a),
\begin{equation*}
\mathsf{k}[\eta]=\tfrac14(\xi\cdot\nabla\varphi)\,\delta+\mathsf{k}^{\mathbf{3}}[\eta]+\mathsf{k}^{\mathbf{6}}[\eta].
\end{equation*}

Consider the combination $-g_0^{-2}\mathfrak{K}_\chi[1]+\sum_i\mathfrak{K}[f_i]$, in which the $i$-th summand is inserted at the slot of $\mathcal{O}(f_i)$, that is, in the connected function from which $\mathcal{O}(f_i)$ has been removed. Here $\mathfrak{K}_\chi[1]$ stands for $\mathfrak{K}_{\chi\xi}[1]$ in the notation of Lemma~\ref{10.4:lem-rotation-field}, and $\mathfrak{K}[f_i]=\mathfrak{K}_{f_i\xi}[1]$ by Lemma~\ref{10.4:lem-null-directions}(ii). Two facts decompose the combination into four terms. The first is Lemma~\ref{10.4:lem-rotation-field}(ii). The second is the splitting of the orbital remainder into its cubic part $\mathfrak{K}^\sharp$ and a sum of exact $\rho$-differences of the potentials $\Psi_\rho$. Together they give
\begin{equation}\label{10.5:eq-reduction-a}
\begin{split}
-g_0^{-2}\mathfrak{K}_\chi[1]+\sum_i\mathfrak{K}[f_i]
&=-g_0^{-2}\mathfrak{K}^\sharp_{\chi\xi}[1]
+g_0^{-2}\sum_p\sum_\rho\xi^\rho(x_p)(\nabla^+_\rho\chi)(x_p)\,\Psi_\rho(p)\\
&\quad+\sum_i\Bigl(\mathfrak{K}^\sharp_{f_i\xi}[1]
-\sum_p\sum_\rho\xi^\rho(x_p)(\nabla^+_\rho f_i)(x_p)\,\Psi_\rho(p)\Bigr).
\end{split}
\end{equation}
We call the four terms on the right, in order: the first (bulk) term, the second (shell) term, and the third and fourth (slot) terms.

The second and fourth terms are diagonal $\operatorname{Sym}^2V$-typed bare terms. Their index is $\rho\rho$ and their weight is $\mathsf{k}[\eta]_{\rho\rho}=\xi^\rho\partial_\rho\varphi$, with $\eta=\varphi\xi$ and $\varphi=\chi$, respectively $\varphi=f_i$.
\begin{itemize}
\item The factor $\Psi_\rho(p)$ sits on the cube between $p$ and $p+\hat\rho$ and is invariant under the reflection of that cube. Hence the type holds after a relocalisation by half a bare lattice spacing, which costs an explicit factor $a$ on a derivative of the weight.
\item Both terms are polynomial in the link variables.
\item They have no marginal part, by clause (iv) of the lemma representing $\mathfrak{r}_\rho$ as an exact $\rho$-difference of $\Psi_\rho$ plus the cubic part $\mathfrak{r}^\sharp_\rho$.
\item Their absolute norm is $O(1)$ per unit weight.
\end{itemize}

\begin{theorem}[The reduction theorem. Stated; proof incomplete at the step marked $(\ast)$]\label{10.5:thm-reduction}
Assume the following.
\begin{itemize}
\item[(a)] Clause (iii-b) of Theorem~0, by statement.
\item[(b)] The typed step with a marginal tensor slot,
Proposition~\ref{10.4:prop-typed-step}, in the form \eqref{10.4:eq-typed-step-a} at
$r=\mathbf{3}$ and $r=\mathbf{6}$, together with the following:
the fluctuation kernel for typed slots, Lemma~\ref{10.4:lem-typed-slot-kernel} and
\eqref{10.4:eq-typed-slot-kernel-a};
the re-filing of typed marginal families, Lemma~\ref{10.4:lem-typed-refiling} and
\eqref{10.4:eq-typed-refiling-a};
and the sourced format of the local tangent run, \eqref{10.4:eq-sourced-format-b},
at $r=\mathbf{6}$.
\item[(c)] The typed step for $\Lambda^2$-typed directions, and the one-lattice Lipschitz
step at the type $\mathbf{1}$.
\item[(d)] The torus bound for the tensor insertion,
Proposition~\ref{10.4:prf-tensor-torus-bound}.
\item[(e)] The vector package in output form, \eqref{10.5:eq-vector-package-output-a}.
\end{itemize}
Then there are numbers $\gamma_0,\varepsilon,\epsilon^{\mathrm{ct}},\epsilon^{\mathrm{sh}}$ with the following property. They depend only on $K$, $m$, $g_0$ and the scaffold; they do not depend on $\vec f$, $\chi$, $\omega$, $n$ or the position. For every $n\ge2$, every $\vec f\in(C^\infty_c)^n$ with pairwise disjoint supports, and every radial admissible $\chi$ (admissibility includes that $\chi\in C^\infty_c(\mathbb{R}^4)$ and $\chi\equiv1$ near $\bigcup_i\operatorname{supp}f_i$),
\begin{equation}\label{10.5:eq-reduction-b}
\begin{split}
&-g_0^{-2}\langle\mathfrak{K}_\chi[1];\mathcal{O}(f_1);\dots;\mathcal{O}(f_n)\rangle
+\sum_i\langle\mathfrak{K}[f_i];(\mathcal{O}(f_l))_{l\neq i}\rangle\\
&\quad=\gamma_0\sum_iS^T_n(f_1,\dots,\xi\cdot\nabla f_i,\dots,f_n)
+\varepsilon\,\langle\mathfrak{B}^{\mathrm{sh}}_\chi;\mathcal{O}(f_1);\dots;\mathcal{O}(f_n)\rangle\\
&\qquad+\epsilon^{\mathrm{ct}}\sum_i\bigl\langle\mathfrak{T}^{\mathbf{3}}
(\mathsf{k}^{\mathbf{3}}[f_i\xi]);(\mathcal{O}(f_l))_{l\neq i}\bigr\rangle\\
&\qquad+\epsilon^{\mathrm{sh}}\bigl\langle\mathfrak{T}^{\mathbf{3}}(\mathsf{k}^{\mathbf{3}}[\chi\xi]);
\mathcal{O}(f_1);\dots;\mathcal{O}(f_n)\bigr\rangle+r_K,
\end{split}
\end{equation}
with
\begin{equation}\label{10.5:eq-reduction-1}
\begin{gathered}
|\gamma_0|+|\epsilon^{\mathrm{ct}}|\le C\mathsf{b}^{\mathrm{ct}},\qquad
|\epsilon^{\mathrm{sh}}|\le C\mathsf{b}^{\mathrm{sh}},\qquad
|\varepsilon|\le Cg_0^2\mathsf{b}^{\mathrm{sh}},\\
\mathsf{b}^{\mathrm{ct}}:=g_0(\log x_0)^{c},\qquad
\mathsf{b}^{\mathrm{sh}}:=g_0^{-1}(\log x_0)^{c},\\
|r_K|\le C(\vec f,\chi,\omega)\,\mathsf{Q}^{n+1}(\log x_0)^{c}
\bigl[\varepsilon_K+\vartheta_1^{K'}+a\bigr].
\end{gathered}
\end{equation}
\end{theorem}

Of the hypotheses, Proposition~\ref{10.4:prop-typed-step}, Lemmas~\ref{10.4:lem-typed-slot-kernel} and~\ref{10.4:lem-typed-refiling}, and the sourced format \eqref{10.4:eq-sourced-format-b} are stated; proof incomplete. The torus bound for the tensor insertion is proved as an implication from the statements named with it. The proof of the vector package in output form is outstanding.

In the bounds, and throughout the proof, $c$ denotes an unspecified exponent whose value may change from one occurrence to the next, and $\varepsilon_K\to0$ is the sequence of clause (ii) of the vector package in output form. Radiality of $\chi$ enters the fourth step of the proof. For radial $\chi$ the trace $\xi\cdot\nabla\chi$ of the strain $\mathsf{k}[\chi\xi]$ vanishes. Indeed, write $\chi(x)=\chi_1(|x|)$ with the radial profile $\chi_1$. At $x=0$ we have $\xi(0)=0$; for $x\neq0$,
\begin{equation*}
\xi\cdot\nabla\chi(x)=\frac{\chi_1'(|x|)}{|x|}\,\xi\cdot x=0,\qquad\text{since}\qquad
\xi\cdot x=\sum_{\rho,\sigma}\omega_{\rho\sigma}x_\rho x_\sigma=0 .
\end{equation*}

\begin{proof}
\emph{First step (types).} Apply the vector package in output form, \eqref{10.5:eq-vector-package-output-a}, whose proof is outstanding:
\begin{itemize}
\item to the first term of \eqref{10.5:eq-reduction-a}, with $c_0=g_0^{-2}$ and $\eta=\chi\xi$;
\item to the third term, with $c_0=1$ and $\eta=f_i\xi$.
\end{itemize}

The vector package leaves a remainder direction $\delta^{\mathrm{rem}}$, whose mixed $z$-derivatives at separated supports are bounded by its clause (ii) with the factor $\varepsilon_K$. Up to $\delta^{\mathrm{rem}}$, the two terms are typed directions. Their weights are
\begin{equation*}
\partial(\chi\xi)=\chi\,\omega^T+\partial\chi\otimes\xi,\qquad \partial(f_i\xi),
\end{equation*}
and, for the first term, in addition $\chi\xi^\rho\partial_\nu w$.

The $\Lambda^2$-components of these directions have no marginal slot, by the statement that $\Lambda^2$-typed families carry no marginal slot. The proof of the typed-direction bound without marginal slot then bounds their contributions; in the hatted (re-filed) notation these bounds are at most $C\mathsf{v}_0(c_0)\vartheta_1^{K'}$, with $c_0=g_0^{-2}$ for the first term and $c_0=1$ for the third.

The $\operatorname{Sym}^2V$-components have the weights $\mathsf{k}[\chi\xi]$, $\mathsf{k}[f_i\xi]$ and $\operatorname{sym}(\xi\otimes\partial w)$.
\begin{itemize}
\item The strain of the Killing field $\xi$ vanishes, since $\partial_\sigma\xi^\rho=\omega_{\rho\sigma}$ is antisymmetric. Hence
\begin{equation*}
\mathsf{k}[\chi\xi]_{\sigma\rho}=\tfrac12\bigl(\xi^\rho\partial_\sigma\chi+\xi^\sigma\partial_\rho\chi\bigr),
\end{equation*}
which is supported in the shell.
\item For the third weight, $\chi\equiv1$ on the supports of the $f_l$, so $\chi\xi^\rho\partial_\nu w=\xi^\rho\partial_\nu w$. Since $w=\sum_lz_lf_l$ and $\partial_\sigma(f_l\xi^\rho)=\xi^\rho\partial_\sigma f_l+f_l\omega_{\rho\sigma}$, where the last summand has vanishing symmetric part, we have
\begin{equation*}
\operatorname{sym}(\xi\otimes\partial w)=\sum_lz_l\,\mathsf{k}[f_l\xi].
\end{equation*}
\end{itemize}

Apply $\partial_{z_i}|_{z=0}$ to the last direction. Near $\operatorname{supp}f_i$ the remaining sources $z_lf_l$, $l\ne i$, vanish. The result is therefore a direction typed with the weight $\mathsf{k}[f_i\xi]$ in the presence of the sources $l\neq i$, with its frozen coefficients taken at $\hat w=0$. The terms of higher order in $z_i$ contribute nothing, since they have no part multilinear in $z_1,\dots,z_n$. The second and fourth terms of \eqref{10.5:eq-reduction-a} are typed as they stand, as described after that display.

\emph{Second step (channels).} By Lemma~\ref{10.4:lem-hypercubic-decomposition}(a) we have $\operatorname{Sym}^2V=\mathsf{V}_{\mathbf{1}}\oplus\mathsf{V}_{\mathbf{3}}\oplus\mathsf{V}_{\mathbf{6}}$. So each of the directions above splits into its components $r=\mathbf{1},\mathbf{3},\mathbf{6}$. The recursive inequalities \eqref{10.4:eq-terminal-matching-a} hold channel by channel. The channels communicate only through derivatives of the weights, that is, through $V\otimes\mathsf{V}_r$-typed defects. By Lemma~\ref{10.4:lem-hypercubic-decomposition}(d) these defects have no marginal part.

\emph{Third step (contact).} Fix $i$ and put $\tilde\eta:=f_i\xi$ and $\mathsf{k}:=\mathsf{k}[\tilde\eta]$. Let $\delta^c$ be the sum of the contact directions at the slot $i$ of the first, third and fourth terms. We use three counter-directions, all bare:
\begin{equation*}
\mathsf{N}(\tilde\eta),\qquad \mathfrak{T}^{\mathbf{3}}(\mathsf{k}^{\mathbf{3}}),\qquad \mathcal{O}(\xi\cdot\nabla f_i).
\end{equation*}
Here $\mathsf{N}(\tilde\eta)$ is the null direction of Lemma~\ref{10.4:lem-null-directions}, defined in \eqref{10.4:eq-null-directions-a}. By Lemma~\ref{10.4:lem-rotation-field}(iii) and Lemma~\ref{10.4:lem-diagonal-stress},
\begin{equation}\label{10.5:eq-reduction-2}
\begin{split}
\mathsf{N}(\tilde\eta)&=\mathfrak{T}^{\mathbf{6}}(\mathsf{k}^{\mathbf{6}})
+2\,\mathfrak{T}^{\mathbf{3}}(\mathsf{k}^{\mathbf{3}})
+\Bigl[\mathfrak{S}[f_i]+\Lambda^2\text{-part of }\sum_p\mathfrak{h}_{f_i}(p)\Bigr]
+\mathfrak{K}_{\tilde\eta}[1]\\
&\quad-\sum_\nu\mathfrak{Q}_\nu(\mathsf{k}_{\nu\nu})
-g_0^2\operatorname{div}V_{\tilde\eta}
+(\text{terms with an explicit factor }a).
\end{split}
\end{equation}
Here $\mathfrak{T}_{\tilde\eta}[1]=\mathfrak{S}[f_i]+\sum_p\mathfrak{h}_{f_i}(p)$ by Lemma~\ref{10.4:lem-null-directions}(ii), and Lemma~\ref{10.4:lem-rotation-field}(iii) and Lemma~\ref{10.4:lem-diagonal-stress} expand the second summand. Its marginal slots at level $0$ are $(\mu^{\mathbf{6}}_0,\mu^{\mathbf{3}}_0,\mu^{\mathbf{1}}_0)=(1,2,0)$. Its other terms are of the kinds already treated in the first step and carry no factor $g_0^{-2}$.

Choose $c_{\mathbf{6}}$, then $c_{\mathbf{3}}$, then $c_{\mathbf{1}}$; the system is triangular by Schur's lemma, Lemma~\ref{10.4:lem-hypercubic-decomposition}(a). The choice is made so that
\begin{equation*}
\theta:=\delta^c-c_{\mathbf{6}}\,\mathsf{N}(\tilde\eta)-c_{\mathbf{3}}\,
\mathfrak{T}^{\mathbf{3}}(\mathsf{k}^{\mathbf{3}})-c_{\mathbf{1}}\,\mathcal{O}(\xi\cdot\nabla f_i)
\end{equation*}
has all three frozen marginal slots equal to $0$ at the level $K'$. This is consequence (ii) of the terminal matching lemma, Lemma~\ref{10.4:lem-terminal-matching}, applied channel by channel.

Its hypothesis $c_\mu\mathsf{b}[\delta^T]\le\frac12\mathfrak{g}_0$ for $\delta^T=\mathsf{N}(\tilde\eta)$ holds within the ceilings on the constants. In $\mathsf{b}=\mathsf{b}_1+\varkappa^\sharp\mathsf{b}_0$ the contributions are as follows.
\begin{itemize}
\item The level-$0$ terms $\mathfrak{Q}$ and $\Psi$ have $\flat$-norm $\mathsf{b}_0\le C\mathfrak{g}_0$ and enter with the factor $\varkappa^\sharp$.
\item The sources of $\mathfrak{K}^\sharp_{\tilde\eta}[1]$ have $\mathsf{b}_1\le C_V\mathsf{v}_0(1)$, which is of the order of $g_0\mathfrak{g}_0(\log x_0)^{c}$.
\item The term $g_0^2\operatorname{div}V_{\tilde\eta}$ shifts $\mu^{\mathbf{3}}_0$ by $O(g_0^2)$.
\end{itemize}

Then every load of $\theta$ at level $k$ is at most $C\mathsf{b}[\delta^c]\vartheta_1^k$. The torus bound for the tensor insertion, which is proved as an implication from the statements named with it, applied to a direction whose loads at level $k$ are at most $C\mathsf{b}[\delta^c]\vartheta_1^k$, gives
\begin{equation*}
\bigl|\langle\theta;(\mathcal{O}(f_l))_{l\neq i}\rangle\bigr|\le C\mathsf{Q}^n\,\mathsf{b}[\delta^c]\,
\vartheta_1^{K'} .
\end{equation*}

Apply \eqref{10.4:eq-null-directions-a} with $\eta=\tilde\eta$ against the $n-1\ge1$ functions $f_l$, $l\neq i$. This gives exactly
\begin{equation*}
\langle\mathsf{N}(\tilde\eta);(\mathcal{O}(f_l))_{l\neq i}\rangle=0 .
\end{equation*}
By construction,
\begin{equation*}
\delta^c=\theta+c_{\mathbf{6}}\mathsf{N}(\tilde\eta)+c_{\mathbf{3}}\mathfrak{T}^{\mathbf{3}}(\mathsf{k}^{\mathbf{3}})+c_{\mathbf{1}}\mathcal{O}(\xi\cdot\nabla f_i),
\end{equation*}
and, since $\mathcal{O}$ is linear in its test function,
\begin{equation*}
\langle\mathcal{O}(\xi\cdot\nabla f_i);(\mathcal{O}(f_l))_{l\neq i}\rangle=S^T_n(f_1,\dots,\xi\cdot\nabla f_i,\dots,f_n).
\end{equation*}
Therefore
\begin{equation}\label{10.5:eq-reduction-3}
\begin{split}
\langle\delta^c;(\mathcal{O}(f_l))_{l\neq i}\rangle
&=c_{\mathbf{3}}\bigl\langle\mathfrak{T}^{\mathbf{3}}(\mathsf{k}^{\mathbf{3}});
(\mathcal{O}(f_l))_{l\neq i}\bigr\rangle\\
&\quad+c_{\mathbf{1}}S^T_n(f_1,\dots,\xi\cdot\nabla f_i,\dots,f_n)+O(\vartheta_1^{K'}).
\end{split}
\end{equation}

$(\ast)$ \emph{Sizes of the coefficients.} The contributions of the pieces of $\delta^c$ to the coefficients $c_{\mathbf{6}},c_{\mathbf{3}},c_{\mathbf{1}}$ are bounded as follows.
\begin{itemize}
\item The contact sources of the first term and the sources of the third term contribute at most $Cg_0(\log x_0)^{c}$. This follows from the vector package in output form, whose proof is outstanding, together with consequence (ii) of Lemma~\ref{10.4:lem-terminal-matching}.
\item The fourth term, and the terms $\mathfrak{Q}$ and $\Psi$ inside $c_{\mathbf{6}}\mathsf{N}(\tilde\eta)$, contribute at most $Cg_0^2(\log x_0)^{2p_0}$. This follows from the trading coefficient in absolute units, Lemma~\ref{10.4:lem-trading-coefficient}, whose proof is incomplete at one step.
\end{itemize}
The hypotheses of Lemma~\ref{10.4:lem-trading-coefficient}, namely its recursive inequalities with their constants, are not established in this proof for the directions of the fourth term and of the terms $\mathfrak{Q}$, $\Psi$.

The coefficients $c_{\mathbf{6}},c_{\mathbf{3}},c_{\mathbf{1}}$ are ratios of frozen numbers and are therefore universal: they do not depend on $\vec f$, $\chi$, $\omega$, $n$ or the position. We put $\gamma_0:=c_{\mathbf{1}}$ and $\epsilon^{\mathrm{ct}}:=c_{\mathbf{3}}$. Since $g_0^2\le g_0$, the two bounds above combine to $|\gamma_0|+|\epsilon^{\mathrm{ct}}|\le C\mathsf{b}^{\mathrm{ct}}$.

\emph{Fourth step (shell).} Put $\mathsf{k}:=\mathsf{k}[\chi\xi]$. It is supported at distance at least $\mathcal{R}-R_f$ from the sources, where $R_f$ denotes a radius such that $\bigcup_l\operatorname{supp}f_l$ lies in the ball of radius $R_f$ about the origin. Let $\delta^{\mathrm{sh}}$ be the sum of the shell directions of the first and second terms. We use two counter-directions, $\sum_p\mathfrak{h}_\chi(p)$ and $\mathfrak{T}^{\mathbf{3}}(\mathsf{k}^{\mathbf{3}})$. Here
\begin{equation*}
\begin{split}
\sum_p\mathfrak{h}_\chi(p)&=\mathfrak{T}^{\mathbf{6}}(\mathsf{k}^{\mathbf{6}})
+2\,\mathfrak{T}^{\mathbf{3}}(\mathsf{k}^{\mathbf{3}})+(\Lambda^2\text{-part})\\
&\quad-\sum_\nu\mathfrak{Q}_\nu(\mathsf{k}_{\nu\nu})
+(\text{terms with an explicit factor }a).
\end{split}
\end{equation*}
Its marginal slots at level $0$ are $(1,2,0)$. The trace channel does not occur, since $\operatorname{tr}\mathsf{k}=\xi\cdot\nabla\chi=0$ for radial $\chi$.

Proceeding as in the third step, we obtain
\begin{equation}\label{10.5:eq-reduction-4}
\begin{split}
\langle\delta^{\mathrm{sh}};\mathcal{O}(\vec f)\rangle
&=c'_{\mathbf{6}}\Bigl\langle\sum_p\mathfrak{h}_\chi(p);\mathcal{O}(\vec f)\Bigr\rangle
+c'_{\mathbf{3}}\bigl\langle\mathfrak{T}^{\mathbf{3}}(\mathsf{k}^{\mathbf{3}});
\mathcal{O}(\vec f)\bigr\rangle\\
&\quad+O\bigl(\mathsf{b}^{\mathrm{sh}}\vartheta_1^{K'}\bigr).
\end{split}
\end{equation}

$(\ast)$ Here $|c'_r|\le C\mathsf{b}^{\mathrm{sh}}$:
\begin{itemize}
\item for the first term, by the vector package in output form, whose proof is outstanding;
\item for the second term, by Lemma~\ref{10.4:lem-trading-coefficient}, whose proof is incomplete at one step, applied with $\mathsf{f}=x_0\,O(1)$, which gives at most $C(\log x_0)^{2p_0}$.
\end{itemize}

By the definition of the shell flux functional $\mathfrak{B}^{\mathrm{sh}}_\chi$,
\begin{equation}\label{10.5:eq-reduction-5}
\sum_p\mathfrak{h}_\chi(p)=-g_0^2\,\mathfrak{B}^{\mathrm{sh}}_\chi
+\sum_pa_p\,(\mathsf{L}^a_\xi\chi)(x_p)+g_0^2\operatorname{div}V_{\chi\xi}.
\end{equation}
The middle term is $\mathcal{O}(\mathsf{L}^a_\xi\chi)$. For radial $\chi$ we have $\|\mathsf{L}^a_\xi\chi\|_\sharp\le Ca^2$, and clause (iii-b) of Theorem~0 applies to its correlation with $\mathcal{O}(\vec f)$. Furthermore,
\begin{equation*}
g_0^2\operatorname{div}V_{\chi\xi}=-\tfrac32\,g_0^2\,
\mathfrak{T}^{\mathbf{3}}\bigl(\mathsf{k}^{\mathrm{diag}}[\chi\xi]\bigr)+O(a^2),
\end{equation*}
where $\mathsf{k}^{\mathrm{diag}}$ is the diagonal part of $\mathsf{k}$. Since $\operatorname{tr}\mathsf{k}[\chi\xi]=0$, this diagonal part is $\mathsf{k}^{\mathbf{3}}[\chi\xi]$. Inserting \eqref{10.5:eq-reduction-5} into \eqref{10.5:eq-reduction-4} we set
\begin{equation*}
\varepsilon:=-c'_{\mathbf{6}}\,g_0^2,\qquad
\epsilon^{\mathrm{sh}}:=c'_{\mathbf{3}}-\tfrac32\,c'_{\mathbf{6}}\,g_0^2 .
\end{equation*}
Since $|c'_r|\le C\mathsf{b}^{\mathrm{sh}}$ and $g_0^2\le1$, this gives $|\varepsilon|\le Cg_0^2\mathsf{b}^{\mathrm{sh}}$ and $|\epsilon^{\mathrm{sh}}|\le C\mathsf{b}^{\mathrm{sh}}$.

Summing \eqref{10.5:eq-reduction-3} over $i$ and adding \eqref{10.5:eq-reduction-4} yields \eqref{10.5:eq-reduction-b}. The remainder $r_K$ collects three kinds of terms:
\begin{itemize}
\item the contributions of $\delta^{\mathrm{rem}}$ (factor $\varepsilon_K$);
\item the $\Lambda^2$-components and the terms $O(\vartheta_1^{K'})$ and $O(\mathsf{b}^{\mathrm{sh}}\vartheta_1^{K'})$ of the third and fourth steps (factor $\vartheta_1^{K'}$);
\item the terms with an explicit factor $a$, the term $\mathcal{O}(\mathsf{L}^a_\xi\chi)$, and the terms $O(a^2)$ (factor $a$).
\end{itemize}
\end{proof}

\begin{corollary}[The vector-channel statement from the diagonal channel]
\label{10.5:cor-vector-channel}
Let $c_1$ be a number with $c_1>\tfrac12$. Assume the following.
\begin{enumerate}
\item The conclusions of the reduction theorem, Theorem~\ref{10.5:thm-reduction}, whose proof
is incomplete at one step. These are the identity \eqref{10.5:eq-reduction-b}, with coefficients
$\gamma_0,\varepsilon,\epsilon^{\mathrm{ct}},\epsilon^{\mathrm{sh}}$ and remainder $r_K$ obeying
the bounds stated there.
\item The \emph{diagonal-channel bound}. Let $n'\ge1$, let $\mathsf{k}$ be a Lipschitz tensor
weight with values in $\mathsf{V}_{\mathbf{3}}$, and let $f_1,\dots,f_{n'}$ be test functions
such that the supports of $\mathsf{k},f_1,\dots,f_{n'}$ are pairwise at distance at least
$\mathsf{d}$. Then, with a constant $C_{n'}(\mathsf{d})$ depending only on $n'$ and $\mathsf{d}$,
\begin{equation}\label{10.5:eq-vector-channel-a}
\bigl|\langle\mathfrak{T}^{\mathbf{3}}(\mathsf{k});\mathcal{O}(f_1);\dots;
\mathcal{O}(f_{n'})\rangle_{K,m}\bigr|
\le C_{n'}(\mathsf{d})\,K^{-c_1}\,\|\mathsf{k}\|_\sharp\prod_{i=1}^{n'}\|f_i\|_\sharp .
\end{equation}
\item The lattice Ward identity with $n$ observables.
\item The complementary ultraviolet statement.
\end{enumerate}
Then, for every $n\ge2$, every $f_1,\dots,f_n\in C^\infty_c$ with pairwise disjoint supports and
every radial admissible cutoff $\chi$ as in the reduction theorem, the left-hand side of
\eqref{10.5:eq-reduction-b} tends to zero as $K\to\infty$ on every fixed torus $\mathbb{T}_m$,
and along every diagonal sequence $(K,m)$ obeying the volume floor. This is the ultraviolet
statement for the orbital remainder, the vector-channel statement, for radial admissible
cutoffs. Together with the complementary ultraviolet statement it gives the ultraviolet
statement, which is by definition the conjunction of the two.
\end{corollary}

\begin{proof}
Fix $n\ge2$ and test functions $f_1,\dots,f_n$ with pairwise disjoint supports. Fix a radial
admissible cutoff $\chi$ as in the reduction theorem, and let $\xi=\omega x$ be the rotation
field entering \eqref{10.5:eq-reduction-b}. The limit is taken either as $K\to\infty$ on a fixed
torus or along a diagonal sequence obeying the volume floor.

Recall that $x_k=g_k^{-2}$, so that $x_0=g_0^{-2}$ and $g_0=x_0^{-1/2}$. We use that, for the
first-passage level $K=K(g_0,g)$, one has $x_0\asymp K$: $x_0$ is bounded above and below by
positive constant multiples of $K$. In the reduction theorem, whose proof is incomplete at one
step, the coefficient
scales are $\mathsf{b}^{\mathrm{ct}}=g_0(\log x_0)^{c}$ and $\mathsf{b}^{\mathrm{sh}}=g_0^{-1}(\log
x_0)^{c}$. Since $g_0^2\,\mathsf{b}^{\mathrm{sh}}=g_0(\log x_0)^c$, the coefficient bounds
stated there read
\begin{gather*}
|\gamma_0|+|\epsilon^{\mathrm{ct}}| \le C\,x_0^{-1/2}(\log x_0)^{c},\\
|\epsilon^{\mathrm{sh}}| \le C\,x_0^{1/2}(\log x_0)^{c},\qquad
|\varepsilon| \le C\,x_0^{-1/2}(\log x_0)^{c}.
\end{gather*}

\emph{The terms carrying $\epsilon^{\mathrm{ct}}$ and $\epsilon^{\mathrm{sh}}$.} Consider the
correlations
\begin{gather*}
\langle\mathfrak{T}^{\mathbf{3}}(\mathsf{k}^{\mathbf{3}}[f_i\xi]);
(\mathcal{O}(f_l))_{l\ne i}\rangle,\\
\langle\mathfrak{T}^{\mathbf{3}}(\mathsf{k}^{\mathbf{3}}[\chi\xi]);\mathcal{O}(f_1);\dots;
\mathcal{O}(f_n)\rangle
\end{gather*}
in \eqref{10.5:eq-reduction-b}. We bound
them by the diagonal-channel bound \eqref{10.5:eq-vector-channel-a}, applied first with $n'=n-1$
and weights $\mathsf{k}^{\mathbf{3}}[f_i\xi]$, then with $n'=n$ and weight
$\mathsf{k}^{\mathbf{3}}[\chi\xi]$. These weights are fixed once $f_1,\dots,f_n$, $\chi$ and
$\omega$ are fixed, so their norms $\|\cdot\|_\sharp$, the norms $\|f_l\|_\sharp$ and the
constants $C_{n'}(\mathsf{d})$ do not depend on $K$. The separation hypothesis of
\eqref{10.5:eq-vector-channel-a} is met, with $\mathsf{d}>0$ the least distance between the
supports involved. Indeed, the weight $\mathsf{k}^{\mathbf{3}}[f_i\xi]$ is supported in
$\operatorname{supp}f_i$, and the compact sets $\operatorname{supp}f_1,\dots,\operatorname{supp}f_n$
are pairwise disjoint, hence pairwise at positive distance. The weight
$\mathsf{k}^{\mathbf{3}}[\chi\xi]$ is a channel piece of the strain of $\chi\xi$, and it vanishes
on the neighbourhood of $\bigcup_i\operatorname{supp}f_i$ on which the cutoff is identically one.
There $\chi\xi=\xi=\omega x$, and $\partial_\sigma\xi^\rho+\partial_\rho\xi^\sigma
=\omega_{\rho\sigma}+\omega_{\sigma\rho}=0$ because $\omega$ is antisymmetric.
Each correlation is then at most a $K$-independent constant
times $K^{-c_1}$. The two terms of \eqref{10.5:eq-reduction-b} are therefore bounded by constants
times $|\epsilon^{\mathrm{ct}}|K^{-c_1}$ and $|\epsilon^{\mathrm{sh}}|K^{-c_1}$. By the bounds
above and $x_0\asymp K$,
\begin{equation*}
|\epsilon^{\mathrm{sh}}|\,K^{-c_1}\le C\,x_0^{1/2}(\log x_0)^{c}K^{-c_1}
\le C'\,K^{\frac12-c_1}(\log(C'K))^{c}\longrightarrow0,
\end{equation*}
because $c_1>\tfrac12$. Also $|\epsilon^{\mathrm{ct}}|K^{-c_1}\le C\,x_0^{-1/2}(\log
x_0)^cK^{-c_1}\to0$. Hence $\epsilon^{\mathrm{ct}}K^{-c_1}\to0$ and
$\epsilon^{\mathrm{sh}}K^{-c_1}\to0$, and both terms tend to zero.

\emph{The term carrying $\gamma_0$.} We have $|\gamma_0|\le C\,x_0^{-1/2}(\log x_0)^c$, so
$\gamma_0\to0$. The connected Schwinger functions $S^T_n(\dots,\xi\cdot\nabla f_i,\dots)$ in
\eqref{10.5:eq-reduction-b} are bounded: the functions $\xi\cdot\nabla f_i$ are smooth with
compact support in $\operatorname{supp}f_i$, so the supports stay pairwise disjoint, and the
bound is uniform in $K$ on a fixed torus by the torus bound of the correlation theorem, and
uniform along a diagonal sequence obeying the volume floor by the per-ball volume-free bound.
Hence $\gamma_0\sum_iS^T_n(\dots,\xi\cdot\nabla
f_i,\dots)\to0$.

\emph{The remainder.} The remaining term of \eqref{10.5:eq-reduction-b} is the remainder $r_K$.
It obeys the bound stated in Theorem~\ref{10.5:thm-reduction}, whose proof is incomplete at one
step, and that bound is used here in the form $r_K\to0$ as $K\to\infty$, on a fixed torus and
along a diagonal sequence obeying the volume floor; in particular $r_K$ is bounded.

\emph{The term carrying $\varepsilon$.} Three inputs combine here: the lattice Ward identity with
$n$ observables, the complementary ultraviolet statement, and the limits just obtained, including
that of $r_K$. By these,
the quantity
\begin{equation*}
(1+\varepsilon)\langle\mathfrak{B}^{\mathrm{sh}}_\chi;\mathcal{O}(f_1);\dots;
\mathcal{O}(f_n)\rangle
\end{equation*}
is a sum of bounded terms. Moreover $|\varepsilon|\le C\,x_0^{-1/2}(\log
x_0)^c\to0$, so $|\varepsilon|\le\tfrac12$ for all large $K$. For such $K$,
\begin{align*}
&\bigl|\varepsilon\langle\mathfrak{B}^{\mathrm{sh}}_\chi;\mathcal{O}(f_1);\dots;
\mathcal{O}(f_n)\rangle\bigr|\\
&\qquad=\frac{|\varepsilon|}{|1+\varepsilon|}\,\bigl|(1+\varepsilon)\langle
\mathfrak{B}^{\mathrm{sh}}_\chi;\mathcal{O}(f_1);\dots;\mathcal{O}(f_n)\rangle\bigr|\\
&\qquad\le2|\varepsilon|\,\bigl|(1+\varepsilon)\langle\mathfrak{B}^{\mathrm{sh}}_\chi;
\mathcal{O}(f_1);\dots;\mathcal{O}(f_n)\rangle\bigr|,
\end{align*}
and this tends to zero. Thus
\begin{equation*}
\varepsilon\langle\mathfrak{B}^{\mathrm{sh}}_\chi;\mathcal{O}(f_1);\dots;
\mathcal{O}(f_n)\rangle\longrightarrow0.
\end{equation*}

This disposes of the terms of \eqref{10.5:eq-reduction-b} carrying $\gamma_0$, $\varepsilon$,
$\epsilon^{\mathrm{ct}}$ and $\epsilon^{\mathrm{sh}}$, and of the remainder $r_K$. Every term on
the right-hand side of \eqref{10.5:eq-reduction-b} therefore tends to zero, and
\eqref{10.5:eq-reduction-b} gives that its left-hand side tends to zero, on a fixed torus
and along a diagonal sequence obeying the volume floor. This is the ultraviolet statement for the
orbital remainder, for radial admissible cutoffs; with the complementary ultraviolet statement,
which is assumed, the ultraviolet statement follows.
\end{proof}

The diagonal-channel bound \eqref{10.5:eq-vector-channel-a} with exponent $c_1$ follows from the
stress-tensor-insertion correlation theorem in the channel $r=\mathbf{3}$ together with the
normalisation decay \eqref{10.4:eq-tensor-correlation-decay-a} at $r=\mathbf{3}$ and
$\kappa=c_1$; this is Corollary~\ref{10.4:cor-tensor-correlation-decay}.

\begin{remark}[The form of the Ward identity used in the conclusion]
\label{10.5:rem-conclusion-identity}
Assume the hypotheses of the reduction theorem (Theorem~\ref{10.5:thm-reduction}), whose proof is
incomplete at one step, and let $\gamma_0,\varepsilon,\epsilon^{\mathrm{ct}},\epsilon^{\mathrm{sh}}$
be the numbers it supplies, which depend on $(K,m,g_0)$ and the scaffold only.
Let $n\ge 2$, let $f_1,\dots,f_n$ be smooth compactly supported functions with pairwise disjoint
supports, let $\chi$ be a radial admissible cutoff, and let $\xi=\omega x$. Write
\begin{align*}
Y_i &:= \bigl\langle \mathfrak{T}^{\mathbf{3}}(\mathsf{k}^{\mathbf{3}}[f_i\xi]);
(\mathcal{O}(f_l))_{l\neq i}\bigr\rangle, \qquad 1\le i\le n,\\
Y^{\mathrm{sh}}_\chi &:= \bigl\langle \mathfrak{T}^{\mathbf{3}}(\mathsf{k}^{\mathbf{3}}[\chi\xi]);
\mathcal{O}(f_1);\dots;\mathcal{O}(f_n)\bigr\rangle
\end{align*}
for the two correlations of the diagonal tensor source that occur in \eqref{10.5:eq-reduction-b},
and write $\mathsf{u}_n$ for the sum of the ultraviolet terms on the right side of the rotational
Ward identity at $n$ observables.
Without the vector-channel statement \eqref{10.5:eq-vector-channel-a}, inserting the identity
\eqref{10.5:eq-reduction-b} of Theorem~\ref{10.5:thm-reduction}, whose proof is incomplete at one
step, into the rotational Ward identity at $n$ observables gives
\begin{equation}\label{10.5:eq-conclusion-identity-1}
\begin{split}
&(1-\gamma_0)\sum_{i=1}^n S^T_{n;K,m}\bigl(f_1,\dots,\xi\cdot\nabla f_i,\dots,f_n\bigr)\\
&\quad= (1+\varepsilon)\bigl\langle \mathfrak{B}^{\mathrm{sh}}_\chi;\mathcal{O}(f_1);\dots;
\mathcal{O}(f_n)\bigr\rangle
+ \epsilon^{\mathrm{ct}}\sum_{i=1}^n Y_i + \epsilon^{\mathrm{sh}}Y^{\mathrm{sh}}_\chi\\
&\qquad + \mathsf{u}_n + o(1),
\end{split}
\end{equation}
where the argument $\xi\cdot\nabla f_i$ stands in the $i$-th place, the left side is
$(1-\gamma_0)$ times the left side of the rotational Ward identity at $n$ observables, $\mathsf{u}_n$
is carried unchanged from that identity, and the term written $o(1)$ contains the remainder $r_K$
of \eqref{10.5:eq-reduction-b}. This is the form of the relative ultraviolet identity on which the
conclusion step operates: the three shell functionals of the relative ultraviolet identity are
reduced in \eqref{10.5:eq-conclusion-identity-1} to one, built from the marginal tensor vertex
$\mathcal{T}^{(3)}$, the diagonal one in the enumeration of the three marginal tensor vertices
$\mathcal{T}^{(a)}$, and there is one contact functional.
\end{remark}

\begin{proof}
By the reduction theorem (Theorem~\ref{10.5:thm-reduction}), whose proof is incomplete at one
step, the identity \eqref{10.5:eq-reduction-b} holds for the data fixed above, with the numbers
$\gamma_0,\varepsilon,\epsilon^{\mathrm{ct}},\epsilon^{\mathrm{sh}}$, which do not depend on
$f_1,\dots,f_n$, $\chi$, $\omega$, $n$ or the position, and with remainder $r_K$.

The left side of \eqref{10.5:eq-reduction-b} is the combination of orbital-remainder
correlations
\begin{equation*}
-g_0^{-2}\bigl\langle \mathfrak{K}_\chi[1];\mathcal{O}(f_1);\dots;\mathcal{O}(f_n)\bigr\rangle
+ \sum_{i=1}^n \bigl\langle \mathfrak{K}[f_i];(\mathcal{O}(f_l))_{l\neq i}\bigr\rangle .
\end{equation*}
In the rotational Ward identity at $n$ observables we replace this combination by the right side
of \eqref{10.5:eq-reduction-b}. We then collect terms as follows.

First, the left side of the rotational Ward identity at $n$ observables is
\begin{equation*}
\sum_{i=1}^n S^T_{n;K,m}\bigl(f_1,\dots,\xi\cdot\nabla f_i,\dots,f_n\bigr),
\end{equation*}
so the term of \eqref{10.5:eq-reduction-b} carrying $\gamma_0$ is $\gamma_0$ times it. Moving it
across produces the factor $1-\gamma_0$.

Second, the right side of the rotational Ward identity at $n$ observables contains the shell flux
correlation $\langle\mathfrak{B}^{\mathrm{sh}}_\chi;\mathcal{O}(f_1);\dots;\mathcal{O}(f_n)\rangle$
with coefficient one. The term $\varepsilon\langle\mathfrak{B}^{\mathrm{sh}}_\chi;\mathcal{O}(f_1);
\dots;\mathcal{O}(f_n)\rangle$ joins it, which produces the factor $1+\varepsilon$.

Third, the terms of \eqref{10.5:eq-reduction-b} containing correlations of
$\mathfrak{T}^{\mathbf{3}}$ are, in the notation above, $\epsilon^{\mathrm{ct}}\sum_i Y_i$ and
$\epsilon^{\mathrm{sh}}Y^{\mathrm{sh}}_\chi$.

Fourth, the ultraviolet terms of the Ward identity are carried unchanged; their sum is
$\mathsf{u}_n$.

Fifth, the remainder $r_K$ of \eqref{10.5:eq-reduction-b}, with the bound stated in
Theorem~\ref{10.5:thm-reduction}, is absorbed into the term written $o(1)$.

This gives \eqref{10.5:eq-conclusion-identity-1}. The vector-channel statement
\eqref{10.5:eq-vector-channel-a} enters none of these steps.

It remains to count functionals. The only shell term besides the flux is
$Y^{\mathrm{sh}}_\chi$. It is built from $\mathfrak{T}^{\mathbf{3}}$, the diagonal channel, that
is, from the vertex $\mathcal{T}^{(3)}$, the diagonal one in the enumeration of the three marginal
tensor vertices $\mathcal{T}^{(a)}$. The only contact term is $\sum_i Y_i$.
\end{proof}

\begin{remark}[Decay exponents asked of the diagonal correlation]
\label{10.5:rem-decay-exponents}
Stated; proof incomplete at the step marked $(\ast)$.

Write
\begin{equation*}
\mathcal{N}_3 := \sup\bigl|\langle \mathfrak{T}^{\mathbf{3}};\dots\rangle\bigr|
\end{equation*}
for the supremum of the diagonal correlation, that is, of the correlations of the tensor source
in the channel $\mathbf{3}$ with the observables, which is the quantity that the hypothesis
\eqref{10.5:eq-vector-channel-a} of Corollary~\ref{10.5:cor-vector-channel} bounds. The
reduction theorem (Theorem~\ref{10.5:thm-reduction}),
whose proof is incomplete at one step, is built from the four terms of
\eqref{10.5:eq-reduction-a}: the bulk term and the shell term, carried by the cutoff $\chi$,
and the two slot terms, carried by the test function $f_i$, one of them cubic and the other
carrying the potential $\Psi_\rho(p)$. Each term has a shadow on a tensor weight, that is, it
generates a direction sourced by the tensor-source functional with the tensor weight shown, with a
coefficient measured in absolute units. The bulk term has two such shadows, a contact part and
a shell part. The following table lists, for each shadow, the weight on which it sits, the bound
on its coefficient, and the order of $\mathcal{N}_3$ that it asks for. In the table,
$\mathrm{polylog}$ denotes a power of $\log x_0$ with an exponent independent of $K$, not the
same at each occurrence. The last column is the order of $\mathcal{N}_3$, as $K\to\infty$,
under which the product of the coefficient with $\mathcal{N}_3$ tends to zero, read with
$g_0^{-2}=x_0$ comparable to $K$.
\begin{center}
\begin{tabular}{lllll}
 & term & shadow on & coefficient & asked of $\mathcal{N}_3$\\
\hline
(a) & slot term with $\Psi_\rho$ & $\mathsf{k}^{\mathbf{3}}[f_i\xi]$
 & $\le C g_0^{2}(\log x_0)^{2p_0}$ & $o(K/\mathrm{polylog})$\\
(b) & cubic slot term & $\mathsf{k}[f_i\xi]$
 & $\le C g_0\,\mathrm{polylog}$ & $o(K^{1/2}/\mathrm{polylog})$\\
(c) & bulk term, contact part & $\mathsf{k}[f_i\xi]$
 & $\le C g_0\,\mathrm{polylog}$ & $o(K^{1/2}/\mathrm{polylog})$\\
(d) & shell term with $\Psi_\rho$ & $\mathsf{k}^{\mathbf{3}}[\chi\xi]$
 & $\le C(\log x_0)^{2p_0}$ & $o(1/\mathrm{polylog})$\\
(e) & bulk term, shell part & $\mathsf{k}[\chi\xi]$
 & $\le C g_0^{-1}\,\mathrm{polylog}$ & $o(K^{-1/2}/\mathrm{polylog})$
\end{tabular}
\end{center}

The coefficients in rows (a) and (d) are those of the trading coefficient in absolute units,
Lemma~\ref{10.4:lem-trading-coefficient}, whose proof is incomplete at one step. If the
recursion of Lemma~\ref{10.4:lem-terminal-matching} is fed with the flat-norm bounds alone
(that is, with $\mathsf{b}_1=0$), the coefficients in rows (a) and (d) are only $C\theta_0$ and
$Cx_0\theta_0$, and the orders asked become $o(\theta_0^{-1})$ and $o(1/(K\theta_0))$. The
coefficients in rows (b) and (e) come from the vector package in output form
\eqref{10.5:eq-vector-package-output-a}, stated in
Definition~\ref{10.5:def-vector-package-output} and not proved (proof outstanding), together
with terminal matching for linear recursions
(Lemma~\ref{10.4:lem-terminal-matching}). The coefficient in row (c) comes from the same vector
package, combined with the second-difference bound in which one difference falls on the vector
unit and the other on the source vertex (the factor $\theta_j^2\mathfrak{g}_j$ of
\eqref{10.5:eq-vector-package-output-a}).

Row (e) asks the most. A bound $\mathcal{N}_3\le C K^{-c_1}$ as in
\eqref{10.5:eq-vector-channel-a} gives row (e) precisely when $c_1>\frac12$, which is the
hypothesis of Corollary~\ref{10.5:cor-vector-channel}.

Suppose the vector package in output form is proved with its sources measured on absolute
radii. Then the entry of row (e) in the last column becomes $o(1/\mathrm{polylog})$, and
\eqref{10.5:eq-vector-channel-a} is needed only for some $c_1>0$.
\end{remark}

\begin{proof}[Proof of the last assertion]
$(\ast)$ The absolute-radius form of \eqref{10.5:eq-vector-package-output-a} would follow from
the vector counterpart of the evenness statement for the tensor-sourced step. That counterpart
concerns the unit $\mathfrak{q}^{-2}\mathfrak{r}^\sharp(\cdot\,;\mathfrak{q}B)$, in which the
second argument $B$ of the cubic part $\mathfrak{r}^\sharp$ of the orbital remainder is scaled
by the coupling slot $\mathfrak{q}$, and asserts that its response is an even function of
$\mathfrak{q}$ with a removable singularity at $\mathfrak{q}=0$, and is bounded by its values on
the circle $|\mathfrak{q}|=\mathsf{r}_0$, which are at most $C(\log x_0)^{2p_0}$. This
statement is stated here without proof; in particular it is not contained in
Definition~\ref{10.5:def-vector-package-output}, whose sources are not measured on absolute
radii.

Under $(\ast)$ the coefficient of the shell part of the bulk term is at most
$C(\log x_0)^{2p_0}$ in place of $Cg_0^{-1}\,\mathrm{polylog}$, so its product with
$\mathcal{N}_3$ tends to zero as soon as $\mathcal{N}_3=o(1/\mathrm{polylog})$.
Granting $(\ast)$, the order asked in row (e) therefore becomes $o(1/\mathrm{polylog})$. The
other rows then ask $o(K/\mathrm{polylog})$, $o(K^{1/2}/\mathrm{polylog})$ and
$o(1/\mathrm{polylog})$.
Hence every row asks at most $\mathcal{N}_3=o(1/\mathrm{polylog})$ as $K\to\infty$.

Let $c_1>0$ and assume \eqref{10.5:eq-vector-channel-a} for this $c_1$. Then $\mathcal{N}_3$ is
at most a constant independent of $K$ times $K^{-c_1}$. Since $\mathrm{polylog}$ is a power of
$\log x_0$ and
$x_0$ is comparable to $K$, we have $K^{-c_1}\,\mathrm{polylog}\to 0$. Therefore
$\mathcal{N}_3=o(1/\mathrm{polylog})$, and every row of the table is satisfied with
\eqref{10.5:eq-vector-channel-a} for some $c_1>0$.
\end{proof}

\begin{definition}[Own-coordinate-free cutoffs of the rotation field]
\label{10.5:def-own-coordinate-cutoffs}
Let $\xi=\omega x$ be the infinitesimal rotation field of the standing notation
(Notation~\ref{10.4:not-standing-notation}), with $\omega_{\rho\sigma}=-\omega_{\sigma\rho}$, and put
$|\omega|:=\max_{\rho,\sigma}|\omega_{\rho\sigma}|$. A periodic vector field $\eta$ on
$\mathbb{R}^4$ is \emph{own-coordinate-free} if, for each $\rho$, the component $\eta^\rho$ does
not depend on $x_\rho$. The rotation field $\xi$ has this property as well.

Let $\chi_1$ be the radial cutoff profile of this chapter, let $\mathcal{R}>0$ satisfy
$4\mathcal{R}<L^m$, and let $\theta_{\mathcal{R}}$ be the $2L^m$-periodic function on
$\mathbb{R}$ equal to $\chi_1(|s|/\mathcal{R})$ for $|s|\le L^m$. Define
\begin{equation}\label{10.5:eq-own-coordinate-cutoffs-1}
(\eta^{\natural}_{\mathcal{R}})^\rho(x):=\xi^\rho(x)\prod_{\tau\neq\rho}\theta_{\mathcal{R}}(x_\tau),
\end{equation}
where $\xi^\rho(x)$ is evaluated with the coordinates $x_\tau$, $\tau\neq\rho$, taken in
$(-L^m,L^m]$. Put $B_{\mathcal{R}}:=\{x:\ |x_\mu|\le\mathcal{R}\ \text{for all}\ \mu\}$.
Then $\eta^{\natural}_{\mathcal{R}}$ is own-coordinate-free, $C^\infty$ and periodic, and it
equals $\xi$ on $B_{\mathcal{R}}$. Moreover, with a constant $C$ that depends only on $\chi_1$,
\begin{equation*}
|\partial\eta^{\natural}_{\mathcal{R}}|\le C|\omega|,\qquad
|\partial^2\eta^{\natural}_{\mathcal{R}}|\le C|\omega|/\mathcal{R}.
\end{equation*}
The field $\eta^{\natural}_{\mathcal{R}}$ is not compactly supported in $\mathbb{R}^4$: its
$\rho$-th component lives on the tube $\{|x_\tau|\le2\mathcal{R},\ \tau\neq\rho\}$ along the
$\rho$-axis.
\end{definition}

\begin{proof}
We use that $\chi_1$ is smooth, equal to one on $[0,1]$ and zero on $[2,\infty)$, and we write
$A_j:=\sup_{[0,2]}|\chi_1^{(j)}|$ for $j=0,1,2$.

Since $\omega$ is antisymmetric, $\omega_{\rho\rho}=0$, so
$\xi^\rho(x)=\sum_{\sigma\neq\rho}\omega_{\rho\sigma}x_\sigma$ does not contain $x_\rho$; this is
why $\xi$ is own-coordinate-free. The product $\Theta_\rho(x):=\prod_{\tau\neq\rho}\theta_{\mathcal{R}}(x_\tau)$
does not contain $x_\rho$ either, so
\eqref{10.5:eq-own-coordinate-cutoffs-1} is independent of $x_\rho$.

On $(-L^m,L^m]$ the function $s\mapsto\chi_1(|s|/\mathcal{R})$ is identically one for
$|s|\le\mathcal{R}$, hence smooth at $s=0$. It vanishes for $|s|\ge2\mathcal{R}$, and
$2\mathcal{R}<L^m/2$. Therefore its periodic extension $\theta_{\mathcal{R}}$ is $C^\infty$ and
vanishes in a neighbourhood of the points $L^m+2L^m\mathbb{Z}$. The representative coordinates
used for $\xi^\rho$ jump only at those points, and $\Theta_\rho$ vanishes near them. Away from them,
$\xi^\rho$ is linear. Hence \eqref{10.5:eq-own-coordinate-cutoffs-1} is $C^\infty$ and periodic.

If $x\in B_{\mathcal{R}}$, then $\theta_{\mathcal{R}}(x_\tau)=1$ for every $\tau$, so
$\eta^{\natural}_{\mathcal{R}}(x)=\xi(x)$.

On the support of $\Theta_\rho$ we have $|x_\sigma|\le2\mathcal{R}$ for $\sigma\neq\rho$, so
$|\xi^\rho(x)|\le6\mathcal{R}|\omega|$. Moreover
\begin{equation*}
|\Theta_\rho|\le A_0^3,\qquad
|\partial_\sigma\Theta_\rho|\le\frac{A_1A_0^2}{\mathcal{R}},\qquad
|\partial_\kappa\partial_\sigma\Theta_\rho|\le\frac{A^{(2)}}{\mathcal{R}^2},
\end{equation*}
where $A^{(2)}:=\max\{A_2A_0^2,A_1^2A_0\}$. By the product rule,
\begin{align*}
\partial_\sigma(\eta^{\natural}_{\mathcal{R}})^\rho
&=\omega_{\rho\sigma}\Theta_\rho+\xi^\rho\,\partial_\sigma\Theta_\rho,\\
\partial_\kappa\partial_\sigma(\eta^{\natural}_{\mathcal{R}})^\rho
&=\omega_{\rho\sigma}\partial_\kappa\Theta_\rho+\omega_{\rho\kappa}\partial_\sigma\Theta_\rho
+\xi^\rho\,\partial_\kappa\partial_\sigma\Theta_\rho .
\end{align*}
Inserting the bounds above gives the following two estimates:
\begin{align*}
|\partial_\sigma(\eta^{\natural}_{\mathcal{R}})^\rho|
&\le(A_0^3+6A_1A_0^2)\,|\omega|,\\
|\partial_\kappa\partial_\sigma(\eta^{\natural}_{\mathcal{R}})^\rho|
&\le(2A_1A_0^2+6A^{(2)})\,|\omega|/\mathcal{R}.
\end{align*}

Finally, $(\eta^{\natural}_{\mathcal{R}})^\rho$ vanishes unless $|x_\tau|\le2\mathcal{R}$ (modulo
$2L^m$) for all $\tau\neq\rho$, and it is independent of $x_\rho$. Its support is therefore
contained in the tube along the $\rho$-axis and is invariant under translations in $x_\rho$.
Consequently, no component that is not identically zero has compact support in $\mathbb{R}^4$.
\end{proof}

\begin{lemma}[Own-coordinate-free cutoffs have no diagonal strain]\label{10.5:lem-no-diagonal-strain}
Let $\eta$ be own-coordinate-free in the sense of Definition~\ref{10.5:def-own-coordinate-cutoffs}.
Then, on every lattice, the following hold.
\begin{enumerate}
\item[(i)] $\operatorname{div} V_\eta = 0$.
\item[(ii)] $\nabla^{s}\cdot\eta = 0$. Moreover, let $w$ be any plaquette function and put
$c := g_0^{-2} - w$. Then, for each $\rho$, $\nabla^{s}_\rho(c\eta^\rho) = -\eta^\rho\nabla^{s}_\rho w$.
\item[(iii)] $\mathfrak{K}_\eta[1] = \mathfrak{K}^\sharp_\eta[1]$, that is, the term containing $\Psi_\rho$ in
Lemma~\ref{10.4:lem-rotation-field}(ii) is absent.
\item[(iv)] $\mathfrak{T}_\eta(p)$ contains no term $\mathfrak{d}_{\ell,p,\rho}$ with $x_\ell - x_p$
parallel to $e_\rho$. Thus $\mathfrak{T}_\eta$ has no diagonal component, and the strain
$\mathsf{k}[\eta]$ is purely off-diagonal.
\item[(v)] Let $f_1,\dots,f_n$ be test functions, and suppose that $\eta = \xi$ on the
$2a$-neighbourhood of $\bigcup_i \operatorname{supp} f_i$. Then
\begin{equation}\label{10.5:eq-no-diagonal-strain-a}
\begin{split}
\sum_{i=1}^n S^T_n(f_1,\dots,\mathsf{L}^a_\xi f_i,\dots,f_n)
&= -g_0^{-2}\bigl\langle(\mathfrak{T}_\eta+\mathfrak{K}^\sharp_\eta)[1];
\mathcal{O}(f_1);\dots;\mathcal{O}(f_n)\bigr\rangle\\
&\quad+\sum_{i=1}^n\bigl\langle(\mathfrak{S}+\mathfrak{K})[f_i];(\mathcal{O}(f_l))_{l\neq i}\bigr\rangle .
\end{split}
\end{equation}
If moreover $\eta\in C^2$, then
\begin{equation}\label{10.5:eq-no-diagonal-strain-1}
\mathfrak{T}_\eta[1] = \mathfrak{T}^{\mathbf{6}}(\mathsf{k}[\eta])
+ \sum_p\sum_{\sigma\neq\rho}\mathsf{a}_{\sigma\rho}(x_p)\,\mathfrak{s}_{\sigma\rho}(p)
+ \mathfrak{T}^{(2)}_\eta[1],
\end{equation}
where $\mathsf{a}_{\sigma\rho} := \tfrac12(\partial_\sigma\eta^\rho - \partial_\rho\eta^\sigma)$.
Here $\mathfrak{T}^{(2)}_\eta[1] := \sum_p \mathfrak{T}^{(2)}_\eta(p)$, with
$\mathfrak{T}^{(2)}_\eta(p)$ as in Lemma~\ref{10.4:lem-rotation-field}(iii).
\end{enumerate}
\end{lemma}

\begin{proof}
All parts except the last rest on one observation. Since $\eta^\rho$ does not depend on $x_\rho$,
any two points that differ only in their $\rho$-th coordinate carry the same value of $\eta^\rho$.
Every lattice difference of $\eta^\rho$ in the direction $\rho$ (symmetric, centred or forward)
compares values of $\eta^\rho$ at such pairs of points. Each such difference therefore vanishes.

(i) By the divergence formula for $V_\eta$, which holds for every $\eta$ and is stated in
Lemma~\ref{10.4:lem-null-directions}(i),
\begin{equation*}
\operatorname{div} V_\eta = -\tfrac32\sum_{p=\{x;\mu,\rho\}} a_p
\bigl[(\nabla^{c}_\rho\eta^\rho)(x_p)+(\nabla^{c}_\mu\eta^\mu)(x_p)\bigr].
\end{equation*}
Each bracket is a sum of centred differences of $\eta^\rho$ in the direction $\rho$ and of
$\eta^\mu$ in the direction $\mu$. By the observation above, $\nabla^{c}_\rho\eta^\rho = 0$ for
every $\rho$. Hence $\operatorname{div} V_\eta = 0$.

(ii) By the observation, $\nabla^{s}_\rho\eta^\rho = 0$ for every $\rho$, and summing over $\rho$
gives $\nabla^{s}\cdot\eta = 0$. For the second identity, $\nabla^{s}_\rho(c\eta^\rho)$ compares
values of $c\eta^\rho$ at two points differing only in their $\rho$-th coordinate. The factor
$\eta^\rho$ takes the same value at both points, so
$\nabla^{s}_\rho(c\eta^\rho) = \eta^\rho\nabla^{s}_\rho c$. Since $g_0^{-2}$ is constant,
$\nabla^{s}_\rho c = -\nabla^{s}_\rho w$, and the identity follows.

(iii) By Lemma~\ref{10.4:lem-rotation-field}(ii) with $c = 1$,
\begin{equation*}
\mathfrak{K}_\eta[1] = \mathfrak{K}^\sharp_\eta[1]
- \sum_p\sum_\rho(\nabla^{+}_\rho\eta^\rho)(p)\,\Psi_\rho(p).
\end{equation*}
The forward difference $\nabla^{+}_\rho\eta^\rho$ vanishes by the observation, so the last sum is
zero.

(iv) By definition,
\begin{equation*}
\mathfrak{T}_\eta(p) = a^{-1}\sum_{\ell\subset\partial p}\sum_\rho
\bigl(\eta^\rho(x_\ell)-\eta^\rho(x_p)\bigr)\mathfrak{d}_{\ell,p,\rho}.
\end{equation*}
If $x_\ell - x_p$ is parallel to $e_\rho$, then $x_\ell$ and $x_p$ differ only in their $\rho$-th
coordinate. Hence $\eta^\rho(x_\ell) = \eta^\rho(x_p)$, and the coefficient of
$\mathfrak{d}_{\ell,p,\rho}$ is zero. Where $\eta$ is differentiable, the same independence gives
$\partial_\rho\eta^\rho = 0$, so the diagonal entries of the matrix of derivatives of $\eta$ vanish.
Consequently the strain $\mathsf{k}[\eta]$, the symmetric part of this matrix, has zero diagonal.

(v) Let $z\in\mathbb{C}^n$ lie in a neighbourhood of the origin on which $Z(z)\neq 0$, where
$w := \sum_i z_i f_i$. Put $c := g_0^{-2} - w$ and $S_c := \sum_p c(p)\,a_p$. By
Lemma~\ref{10.4:lem-rotation-field}(i),
\begin{equation}\label{10.5:eq-no-diagonal-strain-2}
V_\eta S_c = -\sum_p a_p\sum_\rho\nabla^{s}_\rho(c\eta^\rho)(p)
+ \mathfrak{K}_\eta[c] + \mathfrak{T}_\eta[c].
\end{equation}

We first treat the first term on the right. By (ii),
$\nabla^{s}_\rho(c\eta^\rho) = -\eta^\rho\nabla^{s}_\rho w$. By the argument of (ii) with $c$
replaced by $w$, $\eta^\rho\nabla^{s}_\rho w = \nabla^{s}_\rho(w\eta^\rho)$. Hence the first term
equals $\sum_p a_p\,(\nabla^{s}\cdot(w\eta))(p)$. The function $w\eta$ vanishes off
$\bigcup_i\operatorname{supp} f_i$, and on that set $\eta = \xi$; therefore $w\eta = w\xi$. By
Lemma~\ref{10.4:lem-null-directions}(ii) and linearity in the weight,
$\nabla^{s}\cdot(w\xi) = \mathsf{L}^a_\xi w$. The first term is therefore
$\mathcal{O}(\mathsf{L}^a_\xi w)$.

We next treat the remaining terms. By their definitions, $\mathfrak{K}_\eta[c]$ and
$\mathfrak{T}_\eta[c]$ are sums over $p$ of $c(p)$ times quantities independent of $c$. Hence
\begin{equation*}
\mathfrak{K}_\eta[c] + \mathfrak{T}_\eta[c]
= g_0^{-2}(\mathfrak{K}_\eta+\mathfrak{T}_\eta)[1] - (\mathfrak{K}_\eta+\mathfrak{T}_\eta)[w].
\end{equation*}
We now identify the two pieces at the weight $w$.
\begin{itemize}
\item For $\mathfrak{K}$: whenever $w(p)\neq 0$, we have $\eta(x_p) = \xi(x_p)$. Hence
$\mathfrak{K}_\eta[w] = \mathfrak{K}_{w\xi}[1]$, and this equals $\mathfrak{K}[w]$ by
Lemma~\ref{10.4:lem-null-directions}(ii).
\item For $\mathfrak{T}$: whenever $w(p)\neq 0$, the points $x_\ell$, $\ell\subset\partial p$, lie
within distance $a$ of $x_p$, hence in the $2a$-neighbourhood of the supports. There $\eta = \xi$,
so $\mathfrak{T}_\eta[w] = \mathfrak{T}_\xi[w]$.
\end{itemize}
To evaluate $\mathfrak{T}_\xi[w]$, write
$w(x_\ell)\xi^\rho(x_\ell) - w(x_p)\xi^\rho(x_p)$ as the sum of
$(w(x_\ell)-w(x_p))\xi^\rho(x_\ell)$ and $w(x_p)(\xi^\rho(x_\ell)-\xi^\rho(x_p))$. Inserting this
into the definition of $\mathfrak{T}_{w\xi}(p)$ gives, for every $p$,
\begin{equation*}
\mathfrak{T}_{w\xi}(p) = \mathfrak{h}_w(p) + w(x_p)\,\mathfrak{T}_\xi(p).
\end{equation*}
Summing over $p$ and using Lemma~\ref{10.4:lem-null-directions}(ii), we obtain
\begin{equation*}
\mathfrak{T}_\xi[w] = \mathfrak{T}_{w\xi}[1] - \sum_p\mathfrak{h}_w(p) = \mathfrak{S}[w].
\end{equation*}

By the integration-by-parts identity for the vector field $V_\eta$,
$\langle V_\eta S_c\rangle_w = \langle\operatorname{div} V_\eta\rangle_w$, and this vanishes by (i).
Taking expectations in \eqref{10.5:eq-no-diagonal-strain-2} therefore gives
\begin{equation}\label{10.5:eq-no-diagonal-strain-3}
0 = \langle\mathcal{O}(\mathsf{L}^a_\xi w)\rangle_w
+ g_0^{-2}\langle(\mathfrak{K}_\eta+\mathfrak{T}_\eta)[1]\rangle_w
- \langle(\mathfrak{K}+\mathfrak{S})[w]\rangle_w .
\end{equation}
By (iii), $\mathfrak{K}_\eta[1] = \mathfrak{K}^\sharp_\eta[1]$ in the middle term.

Now apply $\partial_{z_1}\cdots\partial_{z_n}$ to \eqref{10.5:eq-no-diagonal-strain-3} at $z = 0$.
For a function $F$ of the field that does not depend on $z$, this mixed derivative of
$\langle F\rangle_w$ is the truncated expectation
$\langle F;\mathcal{O}(f_1);\dots;\mathcal{O}(f_n)\rangle$. Since $\mathsf{L}^a_\xi$ and
$\mathcal{O}$ are linear, $\mathcal{O}(\mathsf{L}^a_\xi w) = \sum_i z_i\,\mathcal{O}(\mathsf{L}^a_\xi f_i)$.
The mixed derivative of its expectation is therefore
\begin{equation*}
\sum_i\bigl\langle\mathcal{O}(\mathsf{L}^a_\xi f_i);(\mathcal{O}(f_l))_{l\neq i}\bigr\rangle
= \sum_i S^T_n(f_1,\dots,\mathsf{L}^a_\xi f_i,\dots,f_n).
\end{equation*}
Likewise $\mathfrak{K}[w]$ and
$\mathfrak{S}[w]$ are linear in $w$, by their expressions in
Lemma~\ref{10.4:lem-null-directions}(ii). The last term of \eqref{10.5:eq-no-diagonal-strain-3}
thus yields $-\sum_i\langle(\mathfrak{S}+\mathfrak{K})[f_i];(\mathcal{O}(f_l))_{l\neq i}\rangle$.
Rearranging gives \eqref{10.5:eq-no-diagonal-strain-a}.

Finally, let $\eta\in C^2$. Summing Lemma~\ref{10.4:lem-rotation-field}(iii) over $p$ gives
\begin{equation*}
\mathfrak{T}_\eta[1] = \sum_p\sum_{\sigma,\rho}(\partial_\sigma\eta^\rho)(x_p)\,\mathfrak{s}_{\sigma\rho}(p)
+ \mathfrak{T}^{(2)}_\eta[1].
\end{equation*}
Write $\partial_\sigma\eta^\rho$ as the sum of its symmetric part, the strain $\mathsf{k}[\eta]$,
and its antisymmetric part $\mathsf{a}_{\sigma\rho}$. The diagonal of $\mathsf{a}$ vanishes
identically, so the antisymmetric part contributes the sum over $\sigma\neq\rho$ in
\eqref{10.5:eq-no-diagonal-strain-1}. By (iv) the diagonal of $\mathsf{k}[\eta]$ vanishes, so only
its off-diagonal channel contributes; this is the term $\mathfrak{T}^{\mathbf{6}}(\mathsf{k}[\eta])$.
This proves \eqref{10.5:eq-no-diagonal-strain-1}.
\end{proof}

We use the notation of the own-coordinate-free cutoff before stating the theorem. Let $\mathcal{R}$ satisfy $4\mathcal{R}<L^m$, and let $\eta^{\natural}_{\mathcal{R}}$ be the own-coordinate-free cutoff of the rotation field $\xi=\omega x$ from Definition~\ref{10.5:def-own-coordinate-cutoffs}. Put
\begin{equation}\label{10.5:eq-reduction-without-shell-1}
\mathfrak{B}^{\natural}_{\mathcal{R}}
:=-g_0^{-2}\,\mathfrak{T}^{\mathbf{6}}\bigl(\mathsf{k}[\eta^{\natural}_{\mathcal{R}}]\bigr).
\end{equation}

This functional has the following properties.
\begin{itemize}
\item Its weight $\mathsf{k}[\eta^{\natural}_{\mathcal{R}}]$ is bounded by $C|\omega|$, since $|\partial\eta^{\natural}_{\mathcal{R}}|\le C|\omega|$ by Definition~\ref{10.5:def-own-coordinate-cutoffs}.
\item Its weight has Lipschitz constant at most $C|\omega|/\mathcal{R}$.
\item Its weight vanishes on the box $B_{\mathcal{R}}$.
\item It is $-x_0=-g_0^{-2}$ times the tensor-source functional of an off-diagonal, hence traceless, weight.
\item It therefore lies in the class of sources, in the analysis of the Ward identity, that consist of the factor $x_0$ on a traceless $\operatorname{Sym}^2V$-typed weight. It has no admixture free of $x_0$.
\end{itemize}

For $\eta=\eta^{\natural}_{\mathcal{R}}$, consider the term
\begin{equation*}
\sum_p\sum_{\sigma\neq\rho}\mathsf{a}_{\sigma\rho}(x_p)\,\mathfrak{s}_{\sigma\rho}(p)
\end{equation*}
of the decomposition of $\mathfrak{T}_\eta[1]$ in Lemma~\ref{10.5:lem-no-diagonal-strain}(v). It is $\Lambda^2$-typed. Its weight is bounded and Lipschitz, and equals $\omega$ on $B_{\mathcal{R}}$. It therefore belongs to the class of $\Lambda^2$-typed directions to which the ultraviolet bound of Corollary~\ref{10.3:cor-vector-channel}, the vector-channel statement from the vector package, applies.

\begin{theorem}[Reduction without a shell term]\label{10.5:thm-reduction-without-shell}
Let $\mathcal{R}$ satisfy $4\mathcal{R}<L^m$, and write $\eta^{\natural}=\eta^{\natural}_{\mathcal{R}}$. Assume the following.
\begin{itemize}
\item[(a)] The hypotheses of the reduction theorem hold. The reduction theorem is Theorem~\ref{10.5:thm-reduction}, which is stated with its proof incomplete at the step marked $(\ast)$. In these hypotheses the vector package in output form~\eqref{10.5:eq-vector-package-output-a}, whose proof is outstanding, is read for $\eta=\eta^{\natural}_{\mathcal{R}}$.
\item[(b)] The conclusions of the first three steps of the proof of Theorem~\ref{10.5:thm-reduction} hold. These are the typing of the units~\eqref{10.5:eq-reduction-a}, their splitting into the channels $\mathbf{1},\mathbf{3},\mathbf{6}$, and the treatment of the contact directions at each slot.
\item[(c)] The conclusions of Lemma~\ref{10.5:lem-no-diagonal-strain} and of Lemma~\ref{10.4:lem-hypercubic-decomposition}(d) hold.
\end{itemize}
Let $n\ge 2$, and let $\vec f\in(C^\infty_c)^n$ have pairwise disjoint supports. Let $\gamma_0$, $\epsilon^{\mathrm{ct}}$, $\varepsilon$ be the numbers of~\eqref{10.5:eq-reduction-b}. Then
\begin{equation}\label{10.5:eq-reduction-without-shell-2}
\begin{split}
&-g_0^{-2}\bigl\langle\mathfrak{K}^\sharp_{\eta^{\natural}}[1];\mathcal{O}(f_1);\dots;\mathcal{O}(f_n)\bigr\rangle
+\sum_i\bigl\langle\mathfrak{K}[f_i];(\mathcal{O}(f_l))_{l\neq i}\bigr\rangle\\
&\quad=\gamma_0\sum_i S^T_n(\dots,\xi\cdot\nabla f_i,\dots)
+\varepsilon\bigl\langle\mathfrak{B}^{\natural}_{\mathcal{R}};\mathcal{O}(f_1);\dots;\mathcal{O}(f_n)\bigr\rangle\\
&\qquad+\epsilon^{\mathrm{ct}}\sum_i\bigl\langle\mathfrak{T}^{\mathbf{3}}(\mathsf{k}^{\mathbf{3}}[f_i\xi]);
(\mathcal{O}(f_l))_{l\neq i}\bigr\rangle+r_K .
\end{split}
\end{equation}
Here $r_K$ obeys the bound on the remainder that accompanies~\eqref{10.5:eq-reduction-b} in the reduction theorem. In that bound the constant depending on $\vec f$, $\chi$ and $\omega$ is replaced by a constant depending on $\vec f$, $\mathcal{R}$ and $\omega$. In particular there is no term $\epsilon^{\mathrm{sh}}$.
\end{theorem}

\begin{proof}
The argument uses Lemma~\ref{10.5:lem-no-diagonal-strain} and Lemma~\ref{10.4:lem-hypercubic-decomposition}(a), (d). Lemma~\ref{10.5:lem-no-diagonal-strain} applies because $\eta^{\natural}$ is own-coordinate-free by Definition~\ref{10.5:def-own-coordinate-cutoffs}. Lemma~\ref{10.4:lem-hypercubic-decomposition} is a statement about $W_4$-modules and is used as it stands.

\emph{The first three steps.} Three ingredients are taken as stated:
\begin{itemize}
\item the reduction theorem, Theorem~\ref{10.5:thm-reduction}, whose proof is incomplete at the step marked $(\ast)$; the numbers $\gamma_0$, $\epsilon^{\mathrm{ct}}$, $\varepsilon$ used here are those of~\eqref{10.5:eq-reduction-b} and carry that status;
\item hypothesis~(a), which contains the vector package in output form~\eqref{10.5:eq-vector-package-output-a}, whose proof is outstanding;
\item hypothesis~(b), the first three steps of the proof of Theorem~\ref{10.5:thm-reduction}.
\end{itemize}
Those three steps are used without change. Only the fourth step, which treats the shell, differs; we give it in full.

\emph{The shell unit without a potential.} The shell unit \textup{U2} of~\eqref{10.5:eq-reduction-a} is the term carrying the difference potential $\Psi_\rho(p)$. By Lemma~\ref{10.5:lem-no-diagonal-strain}(iii), $\mathfrak{K}_{\eta^{\natural}}[1]=\mathfrak{K}^\sharp_{\eta^{\natural}}[1]$, so no $\Psi$-term occurs. Hence $\textup{U2}=0$.

\emph{The shell sources of \textup{U1}.} These sources carry the weight $\mathsf{k}[\eta^{\natural}]$. By Lemma~\ref{10.5:lem-no-diagonal-strain}(iv), this strain is purely off-diagonal. By Lemma~\ref{10.4:lem-hypercubic-decomposition}(a), $\mathsf{V}_{\mathbf{6}}$ is the off-diagonal piece of $\operatorname{Sym}^2V$. The weight is therefore of type $\mathsf{V}_{\mathbf{6}}$.

\emph{The channel is preserved.} By Lemma~\ref{10.4:lem-hypercubic-decomposition}(a), the pieces $\mathsf{V}_{\mathbf{1}},\mathsf{V}_{\mathbf{3}},\mathsf{V}_{\mathbf{6}}$ are irreducible and pairwise inequivalent. By Schur's lemma, a $\mathsf{V}_{\mathbf{6}}$-typed direction therefore stays in the channel $\mathbf{6}$, up to defects in the derivatives of the weight. These defects are of two kinds.
\begin{itemize}
\item The defects of first order in the derivatives of the weight are $V\otimes\mathsf{V}_{\mathbf{6}}$-typed. By Lemma~\ref{10.4:lem-hypercubic-decomposition}(d), $\operatorname{Hom}_{W_4}(V\otimes\mathsf{V}_{\mathbf{6}},\mathbb{R}\oplus\operatorname{Sym}^2\Lambda^2)=0$, so they have no marginal part.
\item The defects of second order have marginal parts of all types. Each carries an explicit factor: the square of the lattice spacing at level $j$ times $|\partial^2\mathsf{k}|$. This factor is at most $CL^{-2(K-j)}$. It outweighs each of the bounds on the running normalisations $N_r(0,k;\zeta)$ established in this chapter, including the crude bound contained in~\eqref{10.4:eq-insertion-correlation-b} of the stress-tensor-insertion correlation theorem. The second-order defects therefore contribute only to the remainder $r_K$.
\end{itemize}

\emph{The counter-direction.} Let $\delta^{\mathrm{sh}}$ denote the shell directions of \textup{U1}. In the channel $\mathbf{6}$ the counter-direction is, by~\eqref{10.5:eq-reduction-without-shell-1},
\begin{equation*}
\mathfrak{T}^{\mathbf{6}}\bigl(\mathsf{k}[\eta^{\natural}]\bigr)=-g_0^2\,\mathfrak{B}^{\natural}_{\mathcal{R}} .
\end{equation*}
Up to the factor $-g_0^2$, this is the shell functional itself. The shell directions stay in the channel $\mathbf{6}$ apart from the defects just described, so this single counter-direction suffices. We expand $\delta^{\mathrm{sh}}$ against it by the procedure of the third step of the proof of Theorem~\ref{10.5:thm-reduction} (hypothesis~(b)), applied in the channel $\mathbf{6}$ alone. This gives
\begin{align*}
\bigl\langle\delta^{\mathrm{sh}};\mathcal{O}(f_1);\dots;\mathcal{O}(f_n)\bigr\rangle
&=c'_{\mathbf{6}}\bigl\langle\mathfrak{T}^{\mathbf{6}}\bigl(\mathsf{k}[\eta^{\natural}]\bigr);
\mathcal{O}(f_1);\dots;\mathcal{O}(f_n)\bigr\rangle\\
&\quad+O\bigl(\mathsf{b}^{\mathrm{sh}}\vartheta_1^{K'}\bigr).
\end{align*}
Here the coefficient $c'_{\mathbf{6}}$ is, as in the third step, a ratio of frozen numbers, and hence universal. Inserting the previous display gives
\begin{equation*}
c'_{\mathbf{6}}\,\mathfrak{T}^{\mathbf{6}}\bigl(\mathsf{k}[\eta^{\natural}]\bigr)
=-c'_{\mathbf{6}}g_0^2\,\mathfrak{B}^{\natural}_{\mathcal{R}}
=\varepsilon\,\mathfrak{B}^{\natural}_{\mathcal{R}},
\end{equation*}
with $\varepsilon=-c'_{\mathbf{6}}g_0^2$. This is the number $\varepsilon$ of~\eqref{10.5:eq-reduction-b} in the reduction theorem, whose status is recorded in hypothesis~(a).

\emph{Conclusion.} The shell directions never leave the channel $\mathbf{6}$, apart from the defects described above. No counter-direction in the channel $\mathbf{3}$ is needed for them. The error term of the expansion of $\delta^{\mathrm{sh}}$ enters $r_K$, together with the contributions of the second-order defects. Hence~\eqref{10.5:eq-reduction-without-shell-2} holds, with $\gamma_0$ and $\epsilon^{\mathrm{ct}}$ as in~\eqref{10.5:eq-reduction-b} from the first three steps (hypotheses~(a) and~(b)), and with no term $\epsilon^{\mathrm{sh}}$.
\end{proof}

\begin{corollary}[$O(4)$ invariance with own-coordinate-free cutoffs]
\label{10.5:cor-invariance-cutoffs}
Let $n\ge 2$, let $f_1,\dots,f_n$, the rotation field $\xi=\omega x$ and the coefficients
$\gamma_0$, $\varepsilon$, $\epsilon^{\mathrm{ct}}$, $r_K$ be as in the reduction theorem,
\eqref{10.5:eq-reduction-b}. Write $a$ for the lattice spacing and $S^T_n$ for $S^T_{n;K,m}$. For
a radius $\mathcal{R}$ let $\eta^{\natural}_{\mathcal{R}}$ be the own-coordinate-free cutoff of
radius $\mathcal{R}$; only radii $\mathcal{R}$ for which $\eta^{\natural}_{\mathcal{R}}$ coincides
with $\xi$ on the $2a$-neighbourhood of $\bigcup_i\operatorname{supp} f_i$ are considered. Assume:
\begin{enumerate}
\item[(1)] the hypotheses of the reduction theorem;
\item[(2)] the vector package in output form, \eqref{10.5:eq-vector-package-output-a}, taken for
$\eta=\eta^{\natural}_{\mathcal{R}}$;
\item[(3)] the bound \eqref{10.5:eq-vector-channel-a} of Corollary~\ref{10.5:cor-vector-channel}, in
the single channel $\mathbf{3}$, for $\mathsf{V}_{\mathbf{3}}$-valued plaquette weights
$\mathsf{k}$, with the exponent condition $c_1>\tfrac12$ replaced by $c_1>-\tfrac12$, that is, the
correlations $\langle\mathfrak{T}^{\mathbf{3}}(\mathsf{k});\mathcal{O}(f_1);\dots\rangle$ are bounded
by a constant times $K^{-c_1}$ with $-c_1<\tfrac12$;
\item[(4)] the ultraviolet statement on the spin and antisymmetric terms.
\end{enumerate}
Then:
\begin{enumerate}
\item[(a)] For each fixed $\mathcal{R}$ as above, as $K\to\infty$, on every fixed torus
$\mathbb{T}_m$ and along every diagonal sequence with the volume floor, as in
Corollary~\ref{10.5:cor-vector-channel}, in particular along the sequence of lattice theories along
which the limit of clause (ii-$\mathbb{R}^4$) of Theorem~0 is taken, the left side of the identity
of Theorem~\ref{10.5:thm-reduction-without-shell},
\begin{equation}\label{10.5:eq-invariance-cutoffs-1}
-g_0^{-2}\bigl\langle\mathfrak{K}^\sharp_{\eta^{\natural}_{\mathcal{R}}}[1];
\mathcal{O}(f_1);\dots;\mathcal{O}(f_n)\bigr\rangle
+\sum_{i}\bigl\langle\mathfrak{K}[f_i];(\mathcal{O}(f_l))_{l\ne i}\bigr\rangle ,
\end{equation}
equals
\begin{equation}\label{10.5:eq-invariance-cutoffs-5}
\varepsilon\bigl\langle\mathfrak{B}^{\natural}_{\mathcal{R}};\mathcal{O}(f_1);\dots;
\mathcal{O}(f_n)\bigr\rangle+o(1),
\end{equation}
and the relative factor $\varepsilon$ of this term on the flux tends to zero.
\item[(b)] Fix an infinitesimal rotation $\omega_0$, a real antisymmetric $4\times4$ matrix, and
take $\xi=\omega_0x$. Assume moreover that hypotheses (1)--(4) hold for every $n\ge 2$ and for all
$f_1,\dots,f_n$ as above; assume clauses (ii-$\mathbb{R}^4$) and (iii-b) of Theorem~0, and, for
every $n\ge 2$ and all such $f_1,\dots,f_n$, the variant, for the box-and-tube flux
$\mathfrak{B}^{\natural}_{\mathcal{R}}$, of the far-sphere flux hypothesis
$(\mathrm{IR}\text{-}\mathrm{fl})_{\omega_0}$ of clause (vi) of Theorem~0, at $\omega_0$:
\begin{equation}\label{10.5:eq-invariance-cutoffs-a}
\liminf_{\mathcal{R}\to\infty}\ \limsup_{j}\
\bigl|\bigl\langle\mathfrak{B}^{\natural}_{\mathcal{R}};
\mathcal{O}(f_1);\dots;\mathcal{O}(f_n)\bigr\rangle_j\bigr| = 0 ,
\end{equation}
where $\langle\cdot\rangle_j$ is the connected expectation in the $j$-th lattice theory of the
sequence along which the limit of clause (ii-$\mathbb{R}^4$) is taken. Then, for every $n\ge 2$,
the limit Schwinger functions $S^T_{n;\infty}$ of clause (ii-$\mathbb{R}^4$) are
$O(4)$-invariant, by the argument of the conclusion step, from clauses (ii-$\mathbb{R}^4$) and
(iii-b), the ultraviolet statement in hypothesis~(4), conclusion~(a) and
\eqref{10.5:eq-invariance-cutoffs-a}.
\end{enumerate}
\end{corollary}

\begin{proof}
Fix $\mathcal{R}$ and write $\eta=\eta^{\natural}_{\mathcal{R}}$. This field is
own-coordinate-free. Clause~(v) of Lemma~\ref{10.5:lem-no-diagonal-strain} applies to it for those
$\mathcal{R}$ for which $\eta$ coincides with $\xi$ on the $2a$-neighbourhood of
$\bigcup_i\operatorname{supp} f_i$. For such $\mathcal{R}$ it gives the identity
\eqref{10.5:eq-no-diagonal-strain-a}. The functionals $\mathfrak{T}_\eta+\mathfrak{K}^\sharp_\eta$
and $\mathfrak{S}+\mathfrak{K}$ enter the connected correlations linearly in the first entry, so
\eqref{10.5:eq-no-diagonal-strain-a} may be regrouped as
\begin{equation}\label{10.5:eq-invariance-cutoffs-2}
\begin{split}
\sum_i S^T_n(\dots,\mathsf{L}^a_\xi f_i,\dots)
={}&-g_0^{-2}\bigl\langle\mathfrak{T}_\eta[1];\mathcal{O}(f_1);\dots;\mathcal{O}(f_n)\bigr\rangle
+\sum_i\bigl\langle\mathfrak{S}[f_i];(\mathcal{O}(f_l))_{l\ne i}\bigr\rangle\\
&+\Bigl[-g_0^{-2}\bigl\langle\mathfrak{K}^\sharp_\eta[1];\mathcal{O}(f_1);\dots;
\mathcal{O}(f_n)\bigr\rangle
+\sum_i\bigl\langle\mathfrak{K}[f_i];(\mathcal{O}(f_l))_{l\ne i}\bigr\rangle\Bigr].
\end{split}
\end{equation}
The bracket is the expression \eqref{10.5:eq-invariance-cutoffs-1}.

Hypotheses~(1) and~(2) are the hypotheses of
Theorem~\ref{10.5:thm-reduction-without-shell}. That theorem therefore replaces the bracket by
\begin{equation}\label{10.5:eq-invariance-cutoffs-3}
\begin{split}
&\gamma_0\sum_i S^T_n(\dots,\xi\cdot\nabla f_i,\dots)
+\varepsilon\bigl\langle\mathfrak{B}^{\natural}_{\mathcal{R}};\mathcal{O}(f_1);\dots;
\mathcal{O}(f_n)\bigr\rangle\\
&\qquad+\epsilon^{\mathrm{ct}}\sum_i\bigl\langle\mathfrak{T}^{\mathbf{3}}
(\mathsf{k}^{\mathbf{3}}[f_i\xi]);(\mathcal{O}(f_l))_{l\ne i}\bigr\rangle+r_K .
\end{split}
\end{equation}
This expression contains no shell term. The only correlations in \eqref{10.5:eq-invariance-cutoffs-3}
of the type controlled by hypothesis~(3) are those carrying the coefficient
$\epsilon^{\mathrm{ct}}$. They are correlations of $\mathfrak{T}^{\mathbf{3}}$ in the channel
$\mathbf{3}$, with weight $\mathsf{k}^{\mathbf{3}}[f_i\xi]$, against the $n-1\ge 1$ observables
$\mathcal{O}(f_l)$, $l\ne i$. The weight $\mathsf{k}^{\mathbf{3}}[f_i\xi]$ is supported in the
support of $f_i$, which lies at distance at least $\mathsf{d}$ from the supports of the $f_l$,
$l\ne i$; so the separation condition of \eqref{10.5:eq-vector-channel-a} holds, with $n'=n-1$.
Hypothesis~(3) is used at exactly this term, with $c_1>-\tfrac12$:
it bounds the term by $\epsilon^{\mathrm{ct}}$ times a constant multiple of $K^{-c_1}$. With the
coefficient $\epsilon^{\mathrm{ct}}$ of the reduction theorem, \eqref{10.5:eq-reduction-b}, and
$x_0\asymp K$, the product $\epsilon^{\mathrm{ct}}K^{-c_1}$ tends to zero as $K\to\infty$ because
$c_1>-\tfrac12$. The remainder $r_K$ of the reduction theorem tends to zero as $K\to\infty$. The
coefficient $\gamma_0$ tends to zero and the $S^T_n$ are bounded, so the first term of
\eqref{10.5:eq-invariance-cutoffs-3} tends to zero. Hence, for fixed $\mathcal{R}$, the bracket
equals \eqref{10.5:eq-invariance-cutoffs-5} as $K\to\infty$; the relative factor $\varepsilon$ is
that of the reduction theorem, \eqref{10.5:eq-reduction-b}, and tends to zero there. This is
conclusion~(a).

On the first line of \eqref{10.5:eq-invariance-cutoffs-2}, expand $\mathfrak{T}_\eta[1]$ as in the
second part of clause~(v) of Lemma~\ref{10.5:lem-no-diagonal-strain}, which applies since
$\eta^{\natural}_{\mathcal{R}}$ is of class $C^2$. The terms of this expansion are treated as
follows.
\begin{itemize}
\item The term containing the antisymmetric part $\mathsf{a}_{\sigma\rho}$ of $\partial\eta$, and
the sum $\sum_i\mathfrak{S}[f_i]$, tend to zero as $K\to\infty$ by the ultraviolet statement of
hypothesis~(4), and so belong to the remainder $o(1)$.
\item The term $g_0^{-2}\mathfrak{T}^{(2)}_\eta[1]$ carries the explicit factor
$a\|\partial^2\eta\|_\infty$ supplied by clause~(iii) of Lemma~\ref{10.4:lem-rotation-field}, the
rotation field acting on a weighted plaquette sum. For fixed
$\mathcal{R}$, $\|\partial^2\eta^{\natural}_{\mathcal{R}}\|_\infty$ does not depend on the lattice,
the correlation multiplying $a\|\partial^2\eta\|_\infty$ is bounded uniformly in $K$ by the same
clause, and $a=L^{-K}\to0$; so this term belongs to the remainder $o(1)$.
\item The terms of the expansion not listed above are, by the definition of the box-and-tube flux
$\mathfrak{B}^{\natural}_{\mathcal{R}}$, the flux correlation appearing in
\eqref{10.5:eq-invariance-cutoffs-5}, with coefficient one; together with the term carrying
$\varepsilon$ in \eqref{10.5:eq-invariance-cutoffs-3} they give the coefficient $1+\varepsilon$.
\end{itemize}
Since $\gamma_0\to0$ and the $S^T_n$ are bounded, both
$\gamma_0\sum_i S^T_n(\dots,\xi\cdot\nabla f_i,\dots)$ and
$\gamma_0\sum_i S^T_n(\dots,\mathsf{L}^a_\xi f_i,\dots)$ are $o(1)$, so they differ by $o(1)$. We
replace the first by the second, the difference being absorbed in the remainder $o(1)$, and move
the second to the left side, which produces the factor $1-\gamma_0$.
Combining \eqref{10.5:eq-invariance-cutoffs-2} with \eqref{10.5:eq-invariance-cutoffs-3} in this
way gives, for each fixed $\mathcal{R}$,
\begin{equation}\label{10.5:eq-invariance-cutoffs-4}
(1-\gamma_0)\sum_i S^T_n(\dots,\mathsf{L}^a_\xi f_i,\dots)
=(1+\varepsilon)\bigl\langle\mathfrak{B}^{\natural}_{\mathcal{R}};\mathcal{O}(f_1);\dots;
\mathcal{O}(f_n)\bigr\rangle+o(1),
\end{equation}
where $o(1)$ tends to zero in the regime of conclusion~(a). This identity is conclusion~(a) in the
form in which it is used.

For (b), in the proof of the conclusion step the Ward identity enters only through an identity of
the shape \eqref{10.5:eq-invariance-cutoffs-4}, there with the shell flux
$\mathfrak{B}^{\mathrm{sh}}_\chi$ in the place of $\mathfrak{B}^{\natural}_{\mathcal{R}}$. The steps
of that proof after the first therefore run unchanged with clauses (ii-$\mathbb{R}^4$)
and (iii-b) of Theorem~0, the ultraviolet statement of hypothesis~(4),
\eqref{10.5:eq-invariance-cutoffs-4}, and the box-and-tube variant
\eqref{10.5:eq-invariance-cutoffs-a} of the far-sphere flux hypothesis
$(\mathrm{IR}\text{-}\mathrm{fl})_{\omega_0}$,
applied for every $n$ and $f_1,\dots,f_n$ of (b), at $\omega_0$.
They yield $O(4)$-invariance.
\end{proof}

\begin{remark}[The infrared flux hypothesis on boxes and tubes]
\label{10.5:rem-box-tube-clustering}
The infrared hypothesis of the box-and-tube variant is the condition
\eqref{10.5:eq-invariance-cutoffs-a} in Corollary~\ref{10.5:cor-invariance-cutoffs},
\[
  \liminf_{\mathcal{R}\to\infty}\ \limsup_{j}\,
  \bigl|\langle \mathfrak{B}^{\natural}_{\mathcal{R}};\mathcal{O}(f_1);\dots;
  \mathcal{O}(f_n)\rangle_j\bigr| = 0 .
\]
It follows from the stress-tensor decay condition that is sufficient for the far-sphere flux
hypothesis $(\mathrm{IR}\text{-}\mathrm{fl})_{\omega_0}$ of clause~(vi) of Theorem~0, by which one
stress-tensor insertion far away decorrelates from the observables, in its pointwise form
(stated later in this chapter), restricted to the
off-diagonal components of the tensor weight. The decay threshold it requires of the
pointwise form is the same as for the infrared hypothesis on spherical shells, namely
$|x|^{-4}$. It does not follow from the far-sphere flux hypothesis
$(\mathrm{IR}\text{-}\mathrm{fl})_{\omega_0}$ itself, which concerns one sequence of radii
and sees one spherical shell.

In exchange, two admixtures drop out of this hypothesis. When the contributions of the
rotational Ward identity (Theorem~\ref{10.1:thm-rotational-identity}) are assigned to the
hypotheses, two admixtures carrying no factor of the initial inverse coupling
$x_0=g_0^{-2}$ would otherwise enter this one: the diagonal marginal vertex
$\mathcal{T}^{(3)}$ and the diagonal tensor of the Jacobian of the displacement. Here they
are identically absent.
\end{remark}

\begin{proof}
Write $\eta^{\natural}_{\mathcal{R}}$ for the own-coordinate-free cutoff of radius
$\mathcal{R}$. Its $\rho$-th component is supported in the tube
$\{|x_\tau|\le 2\mathcal{R},\ \tau\neq\rho\}$ along the $\rho$-axis. The cross-section
of this tube orthogonal to that axis is a three-dimensional cube of side $4\mathcal{R}$.
On the box $B_{\mathcal{R}}=\{|x_\mu|\le\mathcal{R}\text{ for all }\mu\}$ one has
$\eta^{\natural}_{\mathcal{R}}=\xi$, and the strain of the rotation field $\xi=\omega x$
vanishes, since $\partial_\sigma\xi^\rho=\omega_{\rho\sigma}$ is antisymmetric in
$(\rho,\sigma)$; so $\mathsf{k}[\eta^{\natural}_{\mathcal{R}}]$ vanishes on $B_{\mathcal{R}}$,
and its support lies at distance at least $\mathcal{R}$ from the origin.
Hence the support of the tensor weight $\mathsf{k}[\eta^{\natural}_{\mathcal{R}}]$ meets
at most $C\mathcal{R}^3$ unit cubes at each distance $D\ge\mathcal{R}$ from the origin.
The stress-tensor insertion enters the box-and-tube flux functional
$\mathfrak{B}^{\natural}_{\mathcal{R}}$ of Corollary~\ref{10.5:cor-invariance-cutoffs} with
the weight $\mathsf{k}[\eta^{\natural}_{\mathcal{R}}]$, so the correlation of
$\mathfrak{B}^{\natural}_{\mathcal{R}}$ with the observables is a sum of contributions of the
unit cubes met by the support of that weight.

Own-coordinate-freeness means that $(\eta^{\natural}_{\mathcal{R}})^\rho$ does not depend
on $x_\rho$. Therefore $\partial_\rho(\eta^{\natural}_{\mathcal{R}})^\rho = 0$ for
each $\rho$. These derivatives are the diagonal components of
$\mathsf{k}[\eta^{\natural}_{\mathcal{R}}]$, so only its off-diagonal components occur.
This is why the pointwise decay condition is needed only for off-diagonal components.

The pointwise stress-tensor decay condition bounds
the contribution of the weight restricted to a unit cube at distance $D$ from the origin
by $C D^{-4-\delta}$ for some $\delta>0$, uniformly in $j$. Sum these bounds over the unit
cubes met by the
support of $\mathsf{k}[\eta^{\natural}_{\mathcal{R}}]$, grouped by their distance from
the origin. With at most $C\mathcal{R}^3$ cubes at each distance $D\ge\mathcal{R}$, this
gives
\begin{equation}\label{10.5:eq-box-tube-clustering-1}
  \limsup_j\,\bigl|\langle \mathfrak{B}^{\natural}_{\mathcal{R}};\mathcal{O}(f_1);\dots;
  \mathcal{O}(f_n)\rangle_j\bigr|
  \le C\mathcal{R}^3\int_{\mathcal{R}}^{\infty} D^{-4-\delta}\,dD
  \le C\mathcal{R}^{-\delta}.
\end{equation}
The right-hand side of \eqref{10.5:eq-box-tube-clustering-1} tends to $0$ as
$\mathcal{R}\to\infty$, so \eqref{10.5:eq-invariance-cutoffs-a} holds.

For $D\ge\mathcal{R}$ one has $\mathcal{R}^3 D^{-4-\delta}\le D^{-1-\delta}$, which is
integrable at infinity for every $\delta>0$. Hence any exponent $\delta>0$ suffices, that
is, the threshold is $|x|^{-4}$.

As for the non-implication asserted in the statement: the far-sphere flux hypothesis
$(\mathrm{IR}\text{-}\mathrm{fl})_{\omega_0}$, which holds along one sequence of radii,
controls the flux through one spherical shell only, whereas the functional
$\mathfrak{B}^{\natural}_{\mathcal{R}}$ is spread
along the tubes over all distances $D\ge\mathcal{R}$. That form therefore does not supply
the bound \eqref{10.5:eq-box-tube-clustering-1}.

Finally, since the diagonal components of $\mathsf{k}[\eta^{\natural}_{\mathcal{R}}]$ vanish
identically, as shown above, the two diagonal admixtures named in the statement do not enter
the hypothesis.
\end{proof}

\begin{remark}[The deciding inequality on one-loop numbers]\label{10.5:rem-deciding-inequality}
The one-channel bound \eqref{10.5:eq-vector-channel-a} with exponent $c_1$ is a bound on one bare
correlation. It is the weakest member of the following chain of sufficient conditions, each of
which implies the one before it.
\begin{enumerate}
\item[(a)] The bound \eqref{10.5:eq-vector-channel-a} follows from the stress-tensor-insertion
correlation theorem in the channel $r=\mathbf{3}$ together with the normalisation decay
\eqref{10.4:eq-tensor-correlation-decay-a} for $r=\mathbf{3}$ and $\kappa=c_1$. The latter is an
inequality on the products $P^{\mathbf{3}}_k(0)$ built from one sequence of numbers of the frozen
tangent run.
\item[(b)] Under the remaining hypotheses of Lemma~\ref{10.4:lem-normalisation-decay}, this
normalisation decay follows from the inequality \eqref{10.4:eq-normalisation-decay-a} for
$r=\mathbf{3}$. The latter is an inequality on partial sums of the one-loop numbers
$\mathsf{m}_{\mathbf{3}}(l)$. At $\kappa=1$ it is the one-sided half of a two-sided
representation of these partial sums, displayed in the proof below.
\item[(c)] Assume the index stationarity
\begin{equation}\label{10.5:eq-deciding-inequality-1}
\sum_{l}\bigl|\mathsf{m}_{\mathbf{3}}(l)-\mathsf{m}^{\infty}_{\mathbf{3}}\bigr|<\infty .
\end{equation}
Then \eqref{10.4:eq-normalisation-decay-a} for $r=\mathbf{3}$ and a given $\kappa$ follows from
\begin{equation}\label{10.5:eq-deciding-inequality-2}
\mathsf{m}^{\infty}_{\mathbf{3}}\ \ge\ \kappa\, c_0 ,
\end{equation}
which is one number against one number. Condition~\eqref{10.5:eq-deciding-inequality-1} is of
the same kind as the geometric stationarity $|\beta^0_{l+1}-\beta^0_l|\le C\vartheta^l$ of the
one-loop coefficients $\beta^0_l$ of \eqref{10.4:eq-normalisation-decay-a}, given by the
index-stationarity statement.
\item[(d)] Under \eqref{10.5:eq-deciding-inequality-1} and the remaining hypotheses of (a) and
(b), what the vector-channel statement
(Corollary~\ref{10.5:cor-vector-channel}) asks of this number, for radial admissible cutoffs, is
therefore $\mathsf{m}^{\infty}_{\mathbf{3}}>\tfrac12 c_0$. What the corollary on $O(4)$
invariance with own-coordinate-free cutoffs (Corollary~\ref{10.5:cor-invariance-cutoffs}) asks
is $\mathsf{m}^{\infty}_{\mathbf{3}}>-\tfrac12 c_0$. The second condition says only that the
traceless diagonal direction is not relevant by more than half the one-loop coefficient $c_0$ of
the coupling.

Suppose the exponent required in the bound \eqref{10.5:eq-vector-channel-a} is lowered by one
half, as in the improved form of the reduction theorem, Theorem~\ref{10.5:thm-reduction}; that
improved form is stated in this section without proof. Then the two thresholds become
$\mathsf{m}^{\infty}_{\mathbf{3}}>0$ and $\mathsf{m}^{\infty}_{\mathbf{3}}>-c_0$ respectively.
\item[(e)] The channel $\mathbf{6}$ may be used instead of $\mathbf{3}$ throughout. The sign
condition is therefore needed in only one of the two channels.
\end{enumerate}

The value $\kappa=1$ in \eqref{10.4:eq-normalisation-decay-a}, that is, the two-sided
representation \eqref{10.5:eq-deciding-inequality-3}, is indicated by the following
heuristic, which is not used in any proof. Let $a_p=1-\operatorname{Re}\operatorname{tr}U(\partial p)$
be the plaquette term of $A(U)$, and let $t$ be a real parameter (in this heuristic only). Then
\begin{equation*}
-\sum_p\bigl(x_k-t\tau_k\,\mathsf{c}_h(p)\bigr)a_p
\end{equation*}
is $x_k$ times the Wilson action of the constant metric $1-(t\tau_k/x_k)\,h$, where $h$ is the
metric-type tensor with orientation weight $\mathsf{c}_h(p)$.
A constant metric is removed by a linear change of coordinates, and the logarithmically divergent
part of the one-loop term is covariant under it. One step therefore multiplies the whole bracket
by $x_{k+1}/x_k$, which gives
\begin{equation*}
\tau_{k+1}=\tau_k\bigl(1-\beta^0_{k+1}/x_k\bigr).
\end{equation*}
\end{remark}

\begin{proof}
(a) Corollary~\ref{10.4:cor-tensor-correlation-decay} is applied with $r=\mathbf{3}$ and
$\kappa=c_1$. Its input is the stress-tensor-insertion correlation theorem at $r=\mathbf{3}$
together with \eqref{10.4:eq-tensor-correlation-decay-a}. Its conclusion is the bound
\begin{equation*}
\bigl|\langle\mathfrak{T}^{\mathbf{3}}(\mathsf{k});\mathcal{O}(f_1);\dots;\mathcal{O}(f_n)\rangle\bigr|
\le C^{\otimes}_nC_N\,(x_{K'}/x_0)^{c_1}\,\|\mathsf{k}\|_\sharp\prod_i\|f_i\|_\sharp ,
\end{equation*}
with the constants $C^{\otimes}_n$ and $C_N$ of that corollary. Since $x_0\asymp K$ at fixed
$(g,m)$, this is the bound \eqref{10.5:eq-vector-channel-a} with exponent $c_1$.

(b) This is Lemma~\ref{10.4:lem-normalisation-decay} in the channel $r=\mathbf{3}$.

Before proving (c) we record the case $\kappa=1$ mentioned in (b). Consider the two-sided
representation
\begin{equation}\label{10.5:eq-deciding-inequality-3}
\sum_{l=i+1}^{j}\mathsf{m}_{\mathbf{3}}(l)
=\sum_{l=i+1}^{j}\beta^0_l+s^{(\mathbf{3})}_i-s^{(\mathbf{3})}_j,
\qquad \sup_{l\ge 0}|s^{(\mathbf{3})}_l|\le c_2^{(\mathbf{3})},
\end{equation}
for all $0\le i<j$. This is the analogue, for the tensor vertex, of the two-sided representation
of the one-loop coefficients recalled below. It is also the restriction to the channel
$\mathbf{9}$ of the two-sided one-loop representation for the full tensor vertex, whose constants
in that channel are $c_{\mathbf{9}}=0$ and $\kappa_{\mathbf{9}}=1$. The inequality
\eqref{10.4:eq-normalisation-decay-a} for $r=\mathbf{3}$ and $\kappa=1$, with
$2c_2^{(\mathbf{3})}$ as its constant, is the one-sided half of
\eqref{10.5:eq-deciding-inequality-3}, since the latter gives
\begin{equation*}
\sum_{l=i+1}^{j}\bigl[\mathsf{m}_{\mathbf{3}}(l)-\beta^0_l\bigr]
=s^{(\mathbf{3})}_i-s^{(\mathbf{3})}_j\ \ge\ -2c_2^{(\mathbf{3})} .
\end{equation*}

(c) The argument uses, besides \eqref{10.5:eq-deciding-inequality-1} and
\eqref{10.5:eq-deciding-inequality-2}, only the two-sided representation of the one-loop
coefficients: there are shifts $s_l$ of the inverse couplings with $|s_l|\le c_2$ and
\begin{equation*}
\beta^0_{l+1}=c_0+s_l-s_{l+1} .
\end{equation*}
Summing over $l$ from $i$ to $j-1$ telescopes to
\begin{equation*}
\sum_{l=i+1}^{j}\bigl(\beta^0_l-c_0\bigr)=s_i-s_j .
\end{equation*}
Put
\begin{equation*}
\mathsf{D}_\infty:=\sum_l\bigl|\mathsf{m}_{\mathbf{3}}(l)-\mathsf{m}^{\infty}_{\mathbf{3}}\bigr| ,
\end{equation*}
which is finite by \eqref{10.5:eq-deciding-inequality-1}. For $0\le i<j$ we decompose
\begin{align*}
\sum_{l=i+1}^{j}\bigl[\mathsf{m}_{\mathbf{3}}(l)-\kappa\beta^0_l\bigr]
&=\sum_{l=i+1}^{j}\bigl(\mathsf{m}_{\mathbf{3}}(l)-\mathsf{m}^{\infty}_{\mathbf{3}}\bigr)
+(j-i)\bigl(\mathsf{m}^{\infty}_{\mathbf{3}}-\kappa c_0\bigr)-\kappa\,(s_i-s_j) .
\end{align*}
The first term is at least $-\mathsf{D}_\infty$. The second is non-negative by
\eqref{10.5:eq-deciding-inequality-2}. The third is at least $-2|\kappa|c_2$. The left side is
therefore bounded below by $-(\mathsf{D}_\infty+2|\kappa|c_2)$. This is
\eqref{10.4:eq-normalisation-decay-a} for $r=\mathbf{3}$, with $\mathsf{D}_\infty+2|\kappa|c_2$
as its constant.

(d) Corollary~\ref{10.5:cor-vector-channel} assumes the bound \eqref{10.5:eq-vector-channel-a}
with $c_1>\tfrac12$. By
(a)--(c), taken with $\kappa=c_1$, it is enough that \eqref{10.5:eq-deciding-inequality-2} holds for
some $\kappa>\tfrac12$. Since $c_0>0$, such a $\kappa$ exists if and only if
$\mathsf{m}^{\infty}_{\mathbf{3}}>\tfrac12c_0$: if $\mathsf{m}^{\infty}_{\mathbf{3}}>\tfrac12c_0$,
take $\kappa=\mathsf{m}^{\infty}_{\mathbf{3}}/c_0$; conversely, $\kappa>\tfrac12$ together with
\eqref{10.5:eq-deciding-inequality-2} gives $\mathsf{m}^{\infty}_{\mathbf{3}}>\tfrac12c_0$.

Corollary~\ref{10.5:cor-invariance-cutoffs} assumes the bound \eqref{10.5:eq-vector-channel-a}
with $c_1>-\tfrac12$. The
same argument, with $-\tfrac12$ in place of $\tfrac12$, gives the condition
$\mathsf{m}^{\infty}_{\mathbf{3}}>-\tfrac12c_0$.

Lowering the required exponent by $\tfrac12$ lowers each of the thresholds $\tfrac12c_0$ and
$-\tfrac12c_0$ by $\tfrac12c_0$. This gives the values $0$ and $-c_0$.

(e) Two inputs are used. The first is Lemma~\ref{10.4:lem-null-directions}. The second is the
identity \eqref{10.5:eq-reduction-b} of the reduction theorem, Theorem~\ref{10.5:thm-reduction}
(stated; proof incomplete at the step marked $(\ast)$). Together they control the correlations of
$\mathfrak{T}^{\mathbf{6}}(\mathsf{k})+2\,\mathfrak{T}^{\mathbf{3}}(\mathsf{k})$ when
$\mathsf{k}=\mathsf{k}[\eta]$ is a strain. Hence the channel $\mathbf{6}$ may be used instead of
$\mathbf{3}$ throughout: the chain (a)--(d) may be run with $\mathbf{6}$ in place of $\mathbf{3}$,
and the inequality on the limiting one-loop number is needed in one of the two channels only.
\end{proof}

\begin{remark}[Why exact identities leave one channel]\label{10.5:rem-exact-identities}
Throughout, $x_0=g_0^{-2}$ is the initial inverse coupling, and $\mathcal{O}(\vec f)$ abbreviates
the list $\mathcal{O}(f_1);\dots;\mathcal{O}(f_n)$ inside a connected correlation. For a weight $c$
on plaquettes we write $A(c):=\sum_p c(p)\,a_p$, so that $A=A(1)$ is the plaquette sum of
Lemma~\ref{10.4:lem-null-directions}. We write $\partial^b_\ell$ for the derivative in the link
variable $U_\ell$ in the colour direction $b$.

\emph{Null directions from arbitrary fields.} Any vector field $\mathsf{V}$ on the bare
configuration space gives a null direction $x_0\,\mathsf{V}A-\operatorname{div}\mathsf{V}$, in the
sense in which $\mathsf{N}(\eta)$ is one in \eqref{10.4:eq-null-directions-a}. The following fields
are of this kind.
\begin{enumerate}
\item The split-clover vector field of the lattice Ward identity, taken against the field $V_\eta$
of Lemma~\ref{10.4:lem-null-directions}, both built on the same field $\eta$. Their difference is a
null direction. Its marginal content
has two parts: a term of the type $\tfrac32\,\mathfrak{T}^{\mathbf{3}}$, which comes from the
Jacobian $\operatorname{div}V_\eta$, and $x_0$ times a field carrying an explicit factor $a^2$.
\item The gradient fields
\begin{equation*}
\mathsf{V}_c:=\sum_{\ell,b}\bigl(\partial^b_\ell A(c,\cdot)\bigr)\,\partial^b_\ell .
\end{equation*}
For arbitrary plaquette weights $c$ and observables at points separated from the support of $c$
(as in the separation hypothesis of Lemma~\ref{10.4:lem-null-directions}, the sources must not
involve the links on which $\mathsf{V}_c$ acts), they give the exact identity
\begin{equation}\label{10.5:eq-exact-identities-1}
x_0\Bigl\langle \sum_{\ell,b}\partial^b_\ell A(c)\,\partial^b_\ell A;\mathcal{O}(f_1);\dots\Bigr\rangle
=-12\,\bigl\langle A(c);\mathcal{O}(f_1);\dots\bigr\rangle .
\end{equation}
Consequently the orientation-weighted plaquette sum $A(c)$ is itself $-x_0/12$ times a product of
field equations of dimension six.
\item Direction-restricted fields, that is, fields acting on the links of one direction only.
\end{enumerate}

\emph{What the identities give.} Let $Y^{\mathbf{3}}$ and $Y^{\mathbf{6}}$ be the two bare
correlations on strains, in the channels $\mathbf{3}$ and $\mathbf{6}$. In each of the cases above,
the terminal matching lemma (Lemma~\ref{10.4:lem-terminal-matching}) turns the identity into a linear
relation
\begin{equation}\label{10.5:eq-exact-identities-2}
\alpha_{\mathbf{3}}\,Y^{\mathbf{3}}+\alpha_{\mathbf{6}}\,Y^{\mathbf{6}}=o(1)
\end{equation}
with numerical coefficients $\alpha_{\mathbf{3}},\alpha_{\mathbf{6}}$. The next assertion is
heuristic: the correlations $Y^{\mathbf{3}}$ and $Y^{\mathbf{6}}$ are expected to be of order $1/K$,
and hence not $o(\vartheta_1^{K'})$. On this heuristic ground, every relation
\eqref{10.5:eq-exact-identities-2} must be proportional to the one furnished by $\mathsf{N}(\eta)$:
the net marginal content of a null direction is the conserved combination.

\emph{The classical level.} The only gauge-covariant field of dimension one that is linear in $\eta$
is $\eta^\rho F_{\rho\mu}$. There are three gauge-covariant vector fields of dimension five, which we
denote $v^1,v^2,v^3$, the third being $v^3=\partial_\rho\mathcal{T}^{(3)}_{\rho\rho}$; of them only
$v^2$ is a multiple of the field equation, and $v^3$ is not.

\emph{Positivity.} Let $\mathcal{O}_{\mu\nu}(\varphi)$ be the part of $\mathcal{O}(\varphi)$
carried by plaquettes of orientation $\mu\nu$. Reflection positivity of the lattice measure,
together with $0\le\mathcal{O}_{\mu\nu}(\varphi)\le\mathcal{O}(\varphi)$ for $\varphi\ge0$, bounds
only the untruncated two-point functions. In these the disconnected part, of order $a^{-8}g_0^4$,
dominates.

\emph{Diagonal strains.} Let $\eta$ be a twice differentiable vector field on $\mathbb{R}^4$ whose
strain $\tfrac12(\partial_\sigma\eta^\rho+\partial_\rho\eta^\sigma)$ is purely diagonal, that is,
\begin{equation}\label{10.5:eq-exact-identities-3}
\partial_\sigma\eta^\rho+\partial_\rho\eta^\sigma=0\qquad(\sigma\ne\rho).
\end{equation}
Then the following hold.
\begin{itemize}
\item $\partial_\tau\partial_\sigma\eta^\rho=0$ for pairwise distinct $\tau,\sigma,\rho$.
\item $\eta^\rho=\sum_{\sigma\ne\rho}A^\rho_\sigma(x_\rho,x_\sigma)$ for functions
$A^\rho_\sigma$ of two variables (these functions $A^\rho_\sigma$ are unrelated to the plaquette
sum $A$).
\item In each coordinate plane there is a function $\psi=\psi_{\rho\sigma}(x_\rho,x_\sigma)$ with
$(A^\rho_\sigma,A^\sigma_\rho)=(\partial_\rho\psi,-\partial_\sigma\psi)$.
\end{itemize}
Such fields live on slabs and tubes, and the only compactly supported one is $\eta=0$.

\emph{The channel $\mathbf{3}$ alone.} Consider such an $\eta$ vanishing near the sources, for
instance $\eta=\varphi(x_1)e_1$ on the torus. Since the strain of $\eta$ is diagonal, its
off-diagonal piece $\mathsf{k}^{\mathbf{6}}[\eta]$ vanishes, and the decomposition of
$\mathsf{N}(\eta)$ into channels used in the proof of the reduction theorem
(Theorem~\ref{10.5:thm-reduction}, whose proof is incomplete at one step) leaves in the null relation
\eqref{10.4:eq-null-directions-a} the channel $\mathbf{3}$ alone:
\begin{equation}\label{10.5:eq-exact-identities-4}
2\bigl\langle\mathfrak{T}^{\mathbf{3}}(\mathsf{k}^{\mathbf{3}}[\eta]);\mathcal{O}(\vec f)\bigr\rangle
=-\bigl\langle\mathfrak{K}_\eta[1]-\mathfrak{Q}+\cdots;\mathcal{O}(\vec f)\bigr\rangle ,
\end{equation}
the dots standing for the remaining terms of $\mathsf{N}(\eta)$. Thus the vector package
(Proposition~\ref{10.3:prop-vector-package}) alone
controls the integrals of the diagonal correlation over the transverse directions. This is a
positive by-product.

With the sources inside the slab, the same identity reads
\begin{equation}\label{10.5:eq-exact-identities-5}
\begin{split}
x_0\cdot 2\bigl\langle\mathfrak{T}^{\mathbf{3}}(\mathsf{k}^{\mathbf{3}}[\eta]);
\mathcal{O}(\vec f)\bigr\rangle
&=-\sum_i S^T_n(\dots,\varphi\,\partial_1 f_i,\dots)\\
&\quad+(\text{vector, trace and contact terms}).
\end{split}
\end{equation}
Combined with \eqref{10.4:eq-insertion-correlation-d} of the stress-tensor-insertion correlation
theorem, \eqref{10.5:eq-exact-identities-5} bounds $x_0|P^{\mathbf{3}}_{K'}|$ from below whenever
its right side is not zero. The upper bound on $x_0|P^{\mathbf{3}}_{K'}|$ that the normalisation
statement for the channel $\mathbf{3}$ asks for would, by the same identity, require a lower bound on
a unit-lattice correlation. So the
normalisation of one traceless channel is dynamical information.
\end{remark}

\begin{proof}
We prove the exact identity \eqref{10.5:eq-exact-identities-1} and the assertions on diagonal
strains. The statements labelled heuristic above are not claimed.

\emph{The identity \eqref{10.5:eq-exact-identities-1}.} Apply the null direction
$x_0\,\mathsf{V}A-\operatorname{div}\mathsf{V}$ to $\mathsf{V}=\mathsf{V}_c$. By the definition of
$\mathsf{V}_c$,
\begin{equation*}
\mathsf{V}_cA=\sum_{\ell,b}\partial^b_\ell A(c)\,\partial^b_\ell A .
\end{equation*}
The vanishing of the connected correlations of the null direction with
$\mathcal{O}(f_1);\dots$, for observables at separated points, therefore gives
\begin{equation*}
x_0\Bigl\langle\sum_{\ell,b}\partial^b_\ell A(c)\,\partial^b_\ell A;\mathcal{O}(f_1);\dots\Bigr\rangle
=\bigl\langle\operatorname{div}\mathsf{V}_c;\mathcal{O}(f_1);\dots\bigr\rangle .
\end{equation*}
The divergence is taken with respect to the product Haar measure on the link variables; since each
$\partial^b_\ell$ preserves the Haar measure, the definition of $\mathsf{V}_c$ gives
\begin{equation*}
\operatorname{div}\mathsf{V}_c=\sum_{\ell,b}\partial^b_\ell\partial^b_\ell A(c).
\end{equation*}
This second-order derivative of the weighted plaquette sum is evaluated by the second-derivative
identity for weighted plaquette sums used with the integration by parts for the sourced measure
(Lemma~\ref{10.1:lem-integration-by-parts}): its connected correlation with
$\mathcal{O}(f_1);\dots$ equals $-12\,\langle A(c);\mathcal{O}(f_1);\dots\rangle$ for every weight
$c$. This is \eqref{10.5:eq-exact-identities-1}.

\emph{Second derivatives.} Let $\tau,\sigma,\rho$ be pairwise distinct and write
$s_{\alpha\beta}:=\partial_\alpha\eta^\beta+\partial_\beta\eta^\alpha$. The following is an
algebraic identity, verified by expanding the right side:
\begin{equation*}
2\,\partial_\tau\partial_\sigma\eta^\rho
=\partial_\sigma s_{\tau\rho}+\partial_\tau s_{\sigma\rho}-\partial_\rho s_{\tau\sigma}.
\end{equation*}
All three $s$ on the right have distinct indices, so they vanish by
\eqref{10.5:eq-exact-identities-3}. Hence $\partial_\tau\partial_\sigma\eta^\rho=0$; this is the
Killing computation.

\emph{Decomposition.} Fix $\rho$ and let $\sigma_1,\sigma_2,\sigma_3$ be the other indices. Since
$\partial_{\sigma_j}\partial_{\sigma_k}\eta^\rho=0$ for $j\ne k$, the derivative
$\partial_{\sigma_j}\eta^\rho$ depends only on $(x_\rho,x_{\sigma_j})$. Integrating in
$x_{\sigma_1}$, then $x_{\sigma_2}$, then $x_{\sigma_3}$ from $0$ writes $\eta^\rho$ as the sum of
three terms $\int_0^{x_{\sigma_j}}(\partial_{\sigma_j}\eta^\rho)(x_\rho,t)\,dt$, the integrand
depending on $(x_\rho,t)$ only, so that each term is a function of
$(x_\rho,x_{\sigma_j})$ only, and the value of $\eta^\rho$ where $x_{\sigma_1}$, $x_{\sigma_2}$ and
$x_{\sigma_3}$ all vanish, a function of $x_\rho$. Absorbing this last function into one of the three
terms gives $\eta^\rho=\sum_{\sigma\ne\rho}A^\rho_\sigma(x_\rho,x_\sigma)$.

\emph{Coordinate planes.} Fix $\sigma\ne\rho$ and insert the decomposition into
\eqref{10.5:eq-exact-identities-3}. Only the summand $A^\rho_\sigma$ of $\eta^\rho$ depends on
$x_\sigma$, and only the summand $A^\sigma_\rho$ of $\eta^\sigma$ depends on $x_\rho$. Hence
\begin{equation*}
\partial_\sigma A^\rho_\sigma+\partial_\rho A^\sigma_\rho=0
\end{equation*}
as functions of $(x_\rho,x_\sigma)$. This says that the $1$-form
$A^\rho_\sigma\,dx_\rho-A^\sigma_\rho\,dx_\sigma$ on the $(x_\rho,x_\sigma)$-plane is closed. By
the Poincar\'e lemma it is $d\psi_{\rho\sigma}$ for a function $\psi_{\rho\sigma}(x_\rho,x_\sigma)$,
that is,
$(A^\rho_\sigma,A^\sigma_\rho)=(\partial_\rho\psi_{\rho\sigma},-\partial_\sigma\psi_{\rho\sigma})$.
A summand depending on one coordinate is, if compactly supported in that coordinate, supported in a
slab; one depending on two coordinates is, under the same proviso, supported in a tube; neither is
compactly supported in $\mathbb{R}^4$ unless it vanishes.

\emph{Compact support.} Suppose $\eta^\rho$ vanishes as soon as $|x_\tau|>R$ for some $\tau$. Take
$|x_{\sigma_1}|>R$. Then $\sum_j A^\rho_{\sigma_j}(x_\rho,x_{\sigma_j})=0$, and varying
$x_{\sigma_2}$ with the other coordinates fixed shows that $A^\rho_{\sigma_2}(x_\rho,\cdot)$ is
constant. The same holds for $A^\rho_{\sigma_3}$, and, taking $|x_{\sigma_2}|>R$, for
$A^\rho_{\sigma_1}$. Hence $\eta^\rho$ is a function of $x_\rho$ alone. It vanishes for
$|x_{\sigma_1}|>R$ and therefore everywhere. Thus a sum of functions of two variables that vanishes
for large $x_\tau$, for every $\tau$, is zero, and the only compactly supported field with purely
diagonal strain is $\eta=0$.
\end{proof}

\begin{definition}[The cut-free one-loop numbers and their partial sums]
\label{10.6:not-cut-free-numbers}
We recall the one-loop response of the tensor-sourced step, in the traceless-diagonal channel
$\mathbf{3}$, as fixed in Proposition~\ref{10.4:prop-vertex-response}, whose proof is incomplete
at one step. For $k \ge 0$ the one-loop drift number at step $k+1$ is
\[
\mathsf{m}_{\mathbf{3}}(k+1)
:= -\beta^{\mathbf{3}}\bigl[\mathsf{L}^{\mathbf{3}}_{k+1}(1;L^{(\le k)})\bigr].
\]
Here $\mathsf{L}^{\mathbf{3}}_{k+1}(s;h)$ is the one-loop family of the tensor-sourced step, with
its parameters $s$ and $h$ as fixed there; it is affine in $h$. It is taken at $s = 1$ and
$h = L^{(\le k)}$, where $L^{(\le k)}$ denotes the collection $(L^{(j)})_{j \le k}$ of the
one-loop terms of the coupling flow, $L^{(j)}$ being the one-loop effective action produced by
the $j$-th step, in the representation fixed in the one-loop coupling-flow theorem. The
expectations entering it are normalised expectations with respect to the measure
$d\mu_{C^{(k)}}\prod \chi(|B| < T_k)$, that is, the Gaussian measure $d\mu_{C^{(k)}}$ of the
fluctuation field $B$ of the step, with covariance $C^{(k)}$, cut by $\prod \chi(|B| < T_k)$.

For $l \ge 1$ we write $\mathsf{m}^{\circ}_{\mathbf{3}}(l)$ for the cut-free number. It is
obtained from the same expression with $k = l-1$, that is, with
\[
s = 1, \qquad h = L^{(\le l-1)},
\]
and with every expectation taken with the normalised Gaussian measure $d\mu_{C^{(l-1)}}$ alone,
without the factor $\prod \chi(|B| < T_{l-1})$.

Its partial sums are
\begin{equation}\label{10.6:eq-cut-free-numbers-1}
n_{\mathbf{3}}(k) := \sum_{l=1}^{k} \mathsf{m}^{\circ}_{\mathbf{3}}(l) \quad (k \ge 1),
\qquad n_{\mathbf{3}}(0) := 0 .
\end{equation}

We also recall, from the one-loop coupling-flow theorem, the one-loop coefficients
$\beta^{0}_{l} := \beta[L^{(l)}]$ of the coupling: there is a sequence of real numbers, one for
each level and each of absolute value at most $c_2$, such that for all $0 \le i < j$ the sum
$\sum_{l=i+1}^{j} \beta^{0}_{l}$ equals $(j-i)c_0$ plus the number at level $i$ minus the
number at level $j$, where
\[
c_0 = \frac{11\,C_{\mathfrak{g}}}{24\pi^2}\log L > 0,
\]
with $C_{\mathfrak{g}}$ the Casimir constant of $\mathfrak{g}$, defined by
$\sum_a [t_a,[t_a,X]] = C_{\mathfrak{g}} X$ for every $X \in \mathfrak{g}$ and a
$\operatorname{tr}$-orthonormal basis $(t_a)$ of $\mathfrak{g}$ ($\operatorname{tr}\mathbf{1} = 1$).
\end{definition}

For $0 \le j < k$, by \eqref{10.6:eq-cut-free-numbers-1} we obtain, by cancelling the
common terms $l \le j$ of the two sums,
\[
n_{\mathbf{3}}(k) - n_{\mathbf{3}}(j) = \sum_{l=j+1}^{k} \mathsf{m}^{\circ}_{\mathbf{3}}(l).
\]
When $j = 0$ this holds directly, since $n_{\mathbf{3}}(0) = 0$.

\begin{definition}[Lattices, constant curvatures and affine fields]\label{10.6:not-affine-fields}
Fix an integer $k\ge 1$; the levels are $j=0,1,\dots,k$.

\emph{Lattices.} For each level $j$, $\mathbb{L}_j$ denotes the $\mathbb{Z}^4$-limit of the
level-$j$ lattice $T^{(j)}$ of Ba{\l}aban's construction \cite[(0.1)]{B-CMP109}, measured in the
units in which level $k$ is the unit lattice. Thus $\mathbb{L}_j$ has spacing
\begin{equation}\label{10.6:eq-affine-fields-1}
\ell_j:=L^{\,j-k},\qquad \eta:=\ell_0=L^{-k}=:N^{-1}.
\end{equation}
The lattices $\mathbb{L}_j$ are lattices of centres of blocks, and since $L$ is odd they are
nested, $\mathbb{L}_{j+1}\subset\mathbb{L}_j$ \cite[(0.1)]{B-CMP109}. For a bond $b$ we write
$x_b$ for its midpoint, for a plaquette $p$ we write $x_p$ for its centre, and $x_c$ denotes in
the same way the centre of the lattice element $c$ by which it is indexed. For
$y\in\mathbb{L}_{j+1}$, $B(y)\subset\mathbb{L}_j$ is the block of $y$. For $x\in\mathbb{L}_j$,
$y(x)$ is the centre of the block containing $x$. Finally, $\square$ denotes a unit cell of
$\mathbb{L}_k$.

\emph{Constant curvatures.} Fix $t_0\in\mathfrak{g}$ with $\operatorname{tr}t_0^2=1$. Let
$\mathsf{f}=(\mathsf{f}_{\kappa\lambda})_{\kappa,\lambda=1}^{4}$ be a real antisymmetric
$4\times 4$ matrix. Its bilinear form is
$\mathsf{f}(u,v):=\sum_{\kappa,\lambda}\mathsf{f}_{\kappa\lambda}u_\kappa v_\lambda$ for
$u,v\in\mathbb{R}^4$. Writing $\mathsf{f}_{\mu\rho}^{2}$ for the square of the entry
$\mathsf{f}_{\mu\rho}$, put
\begin{equation}\label{10.6:eq-affine-fields-2}
\phi_\mu:=\sum_{\rho}\mathsf{f}_{\mu\rho}^{2}\quad(\mu=1,\dots,4),\qquad
\Sigma_{\mathsf{f}}:=\sum_{\rho<\sigma}\mathsf{f}_{\rho\sigma}^{2}.
\end{equation}
For orientation weights $h=(h_\mu)_{\mu=1}^{4}$ with $\sum_\mu h_\mu=0$, let
$\mathsf{t}^{\mathbf{3}}(h;F)=\tfrac12\sum_{\mu<\nu}(h_\mu+h_\nu)|F_{\mu\nu}|^2$ be the reference
quadratic form of the orientation-weighted action $\mathsf{A}^{\mathbf{3}}$ of the
stress-tensor-insertion correlation theorem (Theorem~\ref{10.4:thm-insertion-correlation}). Write
$\mathsf{t}^{\mathbf{3}}(h;\mathsf{f}):=\mathsf{t}^{\mathbf{3}}(h;\mathsf{f}t_0)$ for its value at
the constant curvature $F=\mathsf{f}t_0$, whose $(\mu\nu)$-component is
$F_{\mu\nu}=\mathsf{f}_{\mu\nu}t_0$. Then
\begin{equation}\label{10.6:eq-affine-fields-3}
\mathsf{t}^{\mathbf{3}}(h;\mathsf{f})
=\tfrac12\sum_{\mu<\nu}(h_\mu+h_\nu)\,\mathsf{f}_{\mu\nu}^{2}
=\tfrac12\sum_{\mu}h_\mu\,\phi_\mu .
\end{equation}

\emph{Affine fields.} An \emph{affine field} on $\mathbb{L}_j$ is a link phase $\mathsf{a}$, that
is, a real number $\mathsf{a}(b)$ attached to each bond $b$ of $\mathbb{L}_j$, of the form
\begin{equation}\label{10.6:eq-affine-fields-4}
\mathsf{a}(b)=\ell_j\Bigl(\beta_\mu+\sum_{\kappa}\beta_{\kappa\mu}(x_b)_\kappa\Bigr)
\qquad\text{for } b\parallel e_\mu ,
\end{equation}
with real coefficients $\beta_\mu$, $\beta_{\kappa\mu}$. \emph{The} affine field of
curvature $\mathsf{f}$ and centre $y$ is
\begin{equation}\label{10.6:eq-affine-fields-5}
\mathsf{a}^{y}_{\mathsf{f}}(b):=\tfrac12\,\ell_j\sum_{\kappa}\mathsf{f}_{\kappa\mu}\,(x_b-y)_\kappa
\qquad\text{for } b\parallel e_\mu .
\end{equation}
It is of the form \eqref{10.6:eq-affine-fields-4}, with
$\beta_{\kappa\mu}=\tfrac12\mathsf{f}_{\kappa\mu}$ and
$\beta_\mu=-\tfrac12\sum_\kappa\mathsf{f}_{\kappa\mu}y_\kappa$. Finally,
$\mathcal{U}_{\mathsf{f}}$ is the configuration
$\mathcal{U}_{\mathsf{f}}(b):=\exp\bigl(i\,\mathsf{a}^{0}_{\mathsf{f}}(b)\,t_0\bigr)$, defined by
this formula on every level $j$.
\end{definition}

\begin{proof}
We verify the identity \eqref{10.6:eq-affine-fields-3}. The remaining assertions are definitions,
apart from the nesting $\mathbb{L}_{j+1}\subset\mathbb{L}_j$, which is \cite[(0.1)]{B-CMP109},
and the statement that \eqref{10.6:eq-affine-fields-5} is of the form
\eqref{10.6:eq-affine-fields-4} with the coefficients given, which follows by writing
$(x_b-y)_\kappa=(x_b)_\kappa-y_\kappa$ in \eqref{10.6:eq-affine-fields-5}.

\emph{First equality.} By definition, $\mathsf{t}^{\mathbf{3}}(h;F)$ is one half of the sum over
$\mu<\nu$ of $(h_\mu+h_\nu)$ times the squared norm of the component $F_{\mu\nu}$. For
$F=\mathsf{f}t_0$ we have $F_{\mu\nu}=\mathsf{f}_{\mu\nu}t_0$. The normalisation
$\operatorname{tr}t_0^2=1$ then gives the squared norm $\mathsf{f}_{\mu\nu}^{2}$, which is the
first equality.

\emph{Second equality.} The numbers $\mathsf{f}_{\mu\nu}^{2}$ are symmetric in $(\mu,\nu)$,
because $\mathsf{f}_{\nu\mu}=-\mathsf{f}_{\mu\nu}$. The weight $h_\mu+h_\nu$ is symmetric as
well. Hence the sum over $\mu<\nu$ is one half of the sum over ordered pairs with $\mu\ne\nu$:
\[
\sum_{\mu<\nu}(h_\mu+h_\nu)\mathsf{f}_{\mu\nu}^{2}
=\tfrac12\sum_{\mu\ne\nu}h_\mu\mathsf{f}_{\mu\nu}^{2}
+\tfrac12\sum_{\mu\ne\nu}h_\nu\mathsf{f}_{\mu\nu}^{2}.
\]
Exchanging the names $\mu$ and $\nu$ in the second sum and using
$\mathsf{f}_{\nu\mu}^{2}=\mathsf{f}_{\mu\nu}^{2}$ turns it into the first sum. The right-hand
side therefore equals $\sum_{\mu}h_\mu\sum_{\nu\ne\mu}\mathsf{f}_{\mu\nu}^{2}$.

Antisymmetry also gives $\mathsf{f}_{\mu\mu}=0$, so
$\sum_{\nu\ne\mu}\mathsf{f}_{\mu\nu}^{2}=\sum_\nu\mathsf{f}_{\mu\nu}^{2}=\phi_\mu$ by
\eqref{10.6:eq-affine-fields-2}. Multiplying by $\tfrac12$ gives the second equality in
\eqref{10.6:eq-affine-fields-3}.
\end{proof}

\begin{lemma}[Reproduction and criticality of constant abelian fields]\label{10.6:lem-constant-fields}
Fix $0\le j\le k$ and a real antisymmetric $4\times 4$ matrix $\mathsf{f}=(\mathsf{f}_{\kappa\mu})$.
Let $\mathcal{U}_{\mathsf{f}}$ be the constant abelian configuration of curvature $\mathsf{f}$ on
$\mathbb{L}_j$, and let $\mathsf{a}$ be its affine field, as fixed in
Definition~\ref{10.6:not-affine-fields}. For a plaquette $p$ of $\mathbb{L}_j$ we write
$p\parallel(\mu\nu)$ if $p$ lies in a coordinate plane of directions $\mu,\nu$ and $\partial p$
traverses first a bond of direction $\mu$ forwards, then a bond of direction $\nu$ forwards.
We put $a_p(U):=1-\operatorname{Re}\operatorname{tr}U(\partial p)$.
\begin{enumerate}
\item[(a)] For every plaquette $p\parallel(\mu\nu)$ of $\mathbb{L}_j$ one has
$\mathcal{U}_{\mathsf{f}}(\partial p)=\exp(i\ell_j^2\mathsf{f}_{\mu\nu}t_0)$.
\item[(b)] Let $j<k$ and let $\mathsf{f}$ be small. Let $M$ be the averaging map of
\cite[(0.12)]{B-CMP109}, built with \cite[(0.10)--(0.11)]{B-CMP109}. Then $M$ maps
$\mathcal{U}_{\mathsf{f}}$ on $\mathbb{L}_j$ to $\mathcal{U}_{\mathsf{f}}$ on $\mathbb{L}_{j+1}$.
\item[(c)] Let $\mathsf{c}$ be a weight $p\mapsto\mathsf{c}(p)$ on the plaquettes of $\mathbb{L}_j$
that depends only on the coordinate plane $\{\mu,\nu\}$ of $p$.
Then $\mathcal{U}_{\mathsf{f}}$ is a critical point of
$A(\mathsf{c},\cdot)=\sum_p\mathsf{c}(p)\,a_p$. In particular $J(\mathcal{U}_{\mathsf{f}})=0$,
where $J(U)$ denotes the gradient at $U$ of the Wilson action $A(1,\cdot)$, and
$\nabla\mathsf{A}^{\mathbf{3}}(\mathcal{U}_{\mathsf{f}})=0$. Thus the background is critical for the
vertex as well, and on these fields the term linear in the fluctuation field is absent from the
expansion of the action and of the vertex about the background.
\item[(d)] Let $U_j$ denote the minimiser map of the level-$j$ step and $V^{(j)}$ the critical point
of that step, both regarded as functions of the configuration, and let $D(\,\cdot\,)(\mathbf{1})$
denote the derivative at the unit configuration $\mathbf{1}$. The linearised minimiser maps
$H_j:=DU_j(\mathbf{1})$ and the derivatives
$DV^{(j)}(\mathbf{1})$ of the critical-point maps send an affine field $\mathsf{a}'$ with coefficient
matrix $(\beta_{\kappa\mu})$, in the sense of Definition~\ref{10.6:not-affine-fields}, to affine
fields with the same antisymmetric part
$\beta_{\kappa\mu}-\beta_{\mu\kappa}$ of the coefficient matrix, modulo gradients $d\lambda$;
here $d\lambda$ denotes the lattice gradient of a function $\lambda$ on the sites.
\end{enumerate}
In formulas, with $\mathsf{a}'$ as in (d):
\begin{equation}\label{10.6:eq-constant-fields-a}
\begin{gathered}
\mathcal{U}_{\mathsf{f}}(\partial p)=e^{i\ell_j^2\mathsf{f}_{\mu\nu}t_0}\ \ (p\parallel(\mu\nu)),\qquad
M\bigl(\mathcal{U}_{\mathsf{f}}\ \text{on}\ \mathbb{L}_j\bigr)
=\mathcal{U}_{\mathsf{f}}\ \text{on}\ \mathbb{L}_{j+1},\\
\nabla A(\mathsf{c},\cdot)(\mathcal{U}_{\mathsf{f}})=0,\qquad
H_j\mathsf{a}',\ DV^{(j)}(\mathbf{1})\mathsf{a}'\ \text{affine with the same}\
\beta_{\kappa\mu}-\beta_{\mu\kappa},\ \text{modulo}\ d\lambda .
\end{gathered}
\end{equation}
\end{lemma}

\begin{proof}
Write $e_\mu$ for the unit vector of direction $\mu$. The bond variable of $\mathcal{U}_{\mathsf{f}}$
on a bond $b$ of $\mathbb{L}_j$ of direction $\mu$ is $\exp(i\mathsf{a}(b)t_0)$, where
$\mathsf{a}(b)=\tfrac12\ell_j\,\mathsf{f}(x_b,e_\mu)=\tfrac12\ell_j\sum_\kappa\mathsf{f}_{\kappa\mu}(x_b)_\kappa$.
All bond variables lie in the one-parameter subgroup $\{\exp(i\theta t_0)\}$. Hence they commute,
and the phase of a product along a path is the sum of the signed phases of its bonds.

\emph{Clause (a).} Let $p\parallel(\mu\nu)$. Then $\partial p$ consists of four bonds. Two have
direction $\mu$, with midpoints $x_p\mp\tfrac12\ell_je_\nu$, traversed forwards and backwards
respectively. Two have direction $\nu$, with midpoints $x_p\pm\tfrac12\ell_je_\mu$, traversed
forwards and backwards respectively. The two $\mu$-bonds contribute
\[
\tfrac12\ell_j\sum_\kappa\mathsf{f}_{\kappa\mu}\bigl[(x_p-\tfrac12\ell_je_\nu)_\kappa
-(x_p+\tfrac12\ell_je_\nu)_\kappa\bigr]=\tfrac12\ell_j\,\mathsf{f}_{\nu\mu}(-\ell_j).
\]
The two $\nu$-bonds contribute $\tfrac12\ell_j\,\mathsf{f}_{\mu\nu}\ell_j$. Since $\mathsf{f}$ is
antisymmetric, the total phase is
$\tfrac12\ell_j(\mathsf{f}_{\nu\mu}(-\ell_j)+\mathsf{f}_{\mu\nu}\ell_j)=\ell_j^2\mathsf{f}_{\mu\nu}$,
which is clause (a).

\emph{Clause (b).} In \cite[(0.11)]{B-CMP109} the transporter $U(y,x)$ between two points $y,x$
of $\mathbb{L}_j$ is the group mean \cite[(0.10)]{B-CMP109} of the variables of a set $G(y,x)$ of
lattice contours from $y$ to $x$, indexed by the permutations of the coordinate directions; for a
bond $c$ of $\mathbb{L}_{j+1}$ from $c_-$ to $c_+$, \cite[(0.12)]{B-CMP109} defines the averaged
variable $\bar U(c)=M(U)(c)$ from $U(c)$, the product of the bond variables along $c$, and a group
mean, over $x\in B(c_-)$, of loop variables. We use three observations.

First, let the elements entering the group mean \cite[(0.10)]{B-CMP109} be commuting elements
$e^{i\theta_nt_0}$ with small $\theta_n$. Then that mean is $e^{i\bar\theta t_0}$, where $\bar\theta$
is the arithmetic mean of the $\theta_n$. Here we use that $\mathsf{f}$ is small.

Second, let $\Gamma\in G(y,x)$ be one of the contours over which $U(y,x)$ is averaged. The phase of
$\Gamma$ is the phase of the straight segment from $y$ to
$x$ plus the flux of $\mathsf{f}$ through the surface between $\Gamma$ and the segment. Here the
phase of the segment from $u$ to $w$ is $\tfrac12\mathsf{f}(u,w)$; for a bond $b$ from $u$ to
$w=u+\ell_je_\mu$ this is the value $\mathsf{a}(b)=\tfrac12\mathsf{f}(x_b,w-u)$, since
$x_b=\tfrac12(u+w)$. The flux through a planar polygon is the integral over it of the constant
two-form $\mathsf{f}$; for the parallelogram spanned by $u$ and $v$ it is $\mathsf{f}(u,v)$. Fix a pair of
directions. The contours that take $\mu$ before $\nu$ and those that take $\nu$ before $\mu$ give
opposite fluxes, and they occur equally often among the permutations. Therefore the arithmetic mean
of the fluxes vanishes, and $U(y,x)$ carries exactly the straight-segment phase.

Third, fix a bond $c$ of $\mathbb{L}_{j+1}$ from $c_-$ to $c_+$. For $x\in B(c_-)$, the loop
appearing in \cite[(0.12)]{B-CMP109} has, by the second observation, the flux of the parallelogram
spanned by $x-c_-$ and $c_+-c_-$. This flux is linear in $x-c_-$. By the first observation the
group mean of these loop variables is $e^{i\bar\theta t_0}$, with $\bar\theta$ the arithmetic mean of
their fluxes over $x\in B(c_-)$. Since $L$ is odd, the block $B(c_-)$ is symmetric about $c_-$, and
the mean over $B(c_-)$ of a function linear in $x-c_-$ vanishes; hence $\bar\theta=0$.

Hence $\bar U(c)=U(c)$, the product of the bond variables along $c$. Let $\mu$ be the direction of
$c$. Every bond $b\subset c$ has $(x_b)_\kappa=(x_c)_\kappa$ for $\kappa\ne\mu$, and
$\mathsf{f}_{\mu\mu}=0$. The phase of $U(c)$ is therefore
\[
\sum_{b\subset c}\mathsf{a}(b)=\tfrac12L\ell_j\sum_\kappa\mathsf{f}_{\kappa\mu}(x_c)_\kappa
=\tfrac12\ell_{j+1}\,\mathsf{f}(x_c,e_\mu),
\]
because $L\ell_j=\ell_{j+1}$. This is the phase of $\mathcal{U}_{\mathsf{f}}$ on the bond $c$ of
$\mathbb{L}_{j+1}$, which proves (b).

\emph{Clause (c).} Let $b$ be a bond and $X\in\mathfrak{g}$ a direction. The derivative of
$A(\mathsf{c},\cdot)$ at $\mathcal{U}_{\mathsf{f}}$, at $b$ in the direction $X$, is
\[
\sum_{p\ni b}\pm\,\mathsf{c}(p)\operatorname{tr}\bigl(X'\sin(\ell_j^2\mathsf{f}_{\mu\nu}t_0)\bigr),
\qquad p\parallel(\mu\nu),
\]
where $X'$ is the transport of $X$ along $\partial p$. This transport conjugates $X$ by a product of
bond variables of $\mathcal{U}_{\mathsf{f}}$. Each such product commutes with
$\sin(\ell_j^2\mathsf{f}_{\mu\nu}t_0)$, so by cyclicity of the trace
$\operatorname{tr}(X'\sin(\ell_j^2\mathsf{f}_{\mu\nu}t_0))
=\operatorname{tr}(X\sin(\ell_j^2\mathsf{f}_{\mu\nu}t_0))$.

The plaquettes containing $b$ come in pairs, two in each coordinate plane containing $b$. The
boundaries of the two plaquettes of a given plane traverse $b$ in opposite directions, so they enter
with opposite signs. Both lie in the same coordinate plane, so the hypothesis on $\mathsf{c}$ gives
them equal weights.
By clause (a) and the identity just proved, they also have equal values. Hence each pair cancels and
the derivative vanishes for every $b$ and every $X$.

The Wilson action is the case $\mathsf{c}\equiv1$, which gives $J(\mathcal{U}_{\mathsf{f}})=0$. The
orientation-weighted action $\mathsf{A}^{\mathbf{3}}$ is of the form $A(\mathsf{c},\cdot)$ with a
weight depending only on the coordinate plane of $p$, which gives
$\nabla\mathsf{A}^{\mathbf{3}}(\mathcal{U}_{\mathsf{f}})=0$.
Consequently, since both $J$ and $\nabla\mathsf{A}^{\mathbf{3}}$ vanish at $\mathcal{U}_{\mathsf{f}}$,
the term linear in the fluctuation field is absent from the expansion of the action and of the
vertex about $\mathcal{U}_{\mathsf{f}}$.

\emph{Clause (d).} Clause (d) is the lemma on linear potentials under the block-averaging step,
applied to the maps $H_j$ and $DV^{(j)}(\mathbf{1})$; that lemma states that a linear potential
reproduces itself under the symmetric block averages when $L$ is odd, and that it is the minimiser,
with zero multiplier.
\end{proof}

\begin{lemma}[Second moments read on constant curvature fields]\label{10.6:lem-second-moment}
Fix a level $j\in\{0,\dots,k\}$, and use the lattices, bonds and affine fields of
Notation~\ref{10.6:not-affine-fields}. For a bond $b$ of $\mathbb{L}_j$ let $\delta_b$ be the field equal
to $1$ on $b$ and to $0$ elsewhere; write $b\parallel\mu$ if $b$ is parallel to the direction $\mu$; for a
real function $\omega$ on the sites of $\mathbb{L}_j$ and a bond $b$ with initial point $x_-$ and final
point $x_+$ put $(d\omega)(b):=\omega(x_+)-\omega(x_-)$. Let $\psi=(\psi_y)_{y\in\mathbb{L}_j}$ be a
family of gauge-invariant functions of configurations on $\mathbb{L}_j$, analytic near $\mathbf{1}$, and
translation covariant: $\psi_{y+a}(U)=\psi_y(U(\cdot+a))$ for every lattice vector $a$ of $\mathbb{L}_j$.
Write $D\psi_y(\mathbf{1})$, $D^2\psi_y(\mathbf{1})$ for the derivatives at $B=0$ of
$B\mapsto\psi_y(e^{iB})$, put
\begin{equation*}
K_y(b,b'):=D^2\psi_y(\mathbf{1})[t_0\delta_b,\,t_0\delta_{b'}],
\end{equation*}
assume that for some constants $C,\delta>0$
\begin{equation}\label{10.6:eq-second-moment-1}
|K_y(b,b')|\le C\,e^{-\delta(|x_b-y|+|x_{b'}-y|)/\ell_j}\qquad\text{for all } y,b,b',
\end{equation}
and let $\Psi:=\sum_y\psi_y$. Then:
\begin{enumerate}
\item[(i)] $\sum_{b'}K_y(b,b')(d\omega)(b')=0$ for every polynomially bounded $\omega$;
\item[(ii)] the number
\begin{equation*}
\mathsf{cf}[\psi](\mathsf{f}):=\tfrac12\sum_{b,b'}K_y(b,b')\,\mathsf{a}^{y}_{\mathsf{f}}(b)\,
\mathsf{a}^{y}_{\mathsf{f}}(b')
\end{equation*}
depends neither on $y$ nor on the centre of the affine field;
\item[(iii)] with
\begin{equation*}
m_{\mu\kappa,\nu\lambda}:=\ell_j^2\sum_{b\parallel\mu,\ b'\parallel\nu}K_y(b,b')\,(x_b-y)_\kappa\,
(x_{b'}-y)_\lambda ,
\end{equation*}
the part of order $p^2$ of the $(\mu,\nu)$ component of the Fourier transform of the Hessian of $\Psi$
at $\mathbf{1}$ in the directions $t_0\delta_b$, $t_0\delta_{b'}$, as a function of the Fourier variable
$p\in\mathbb{R}^4$ dual to $x_{b'}-x_b$,
is $\ell_j^{-2}\sum_{\kappa,\lambda}p_\kappa p_\lambda m_{\mu\kappa,\nu\lambda}$, the array $m$ is
antisymmetric in $(\mu\kappa)$ and in $(\nu\lambda)$, and
$\mathsf{cf}[\psi](\mathsf{f})=\frac18\sum_{\kappa,\lambda,\mu,\nu}\mathsf{f}_{\kappa\mu}
\mathsf{f}_{\lambda\nu}m_{\mu\kappa,\nu\lambda}$;
\item[(iv)] if $\Psi$ is $\mathsf{V}_{\mathbf{3}}$-typed, with orientation weight $h=(h_\mu)$,
$\sum_\mu h_\mu=0$, in the sense of the definition of the orientation-weighted action
$\mathsf{A}^{\mathbf{3}}$ and of the typed coefficient $\beta^{\mathbf{3}}$, then
\begin{equation}\label{10.6:eq-second-moment-2}
\mathsf{cf}[\psi](\mathsf{f})=\beta^{\mathbf{3}}[\Psi]\cdot\mathsf{t}^{\mathbf{3}}(h;\ell_j^2\mathsf{f}).
\end{equation}
\end{enumerate}
\end{lemma}

\begin{proof}
Translation covariance gives $K_{y+a}(b+a,b'+a)=K_y(b,b')$, since the translate of $\delta_{b+a}$ by
$a$ is $\delta_b$. The second derivative is symmetric, so $K_y(b,b')=K_y(b',b)$. By
\eqref{10.6:eq-second-moment-1} every sum below, of $K_y$ against polynomially growing factors, converges
absolutely.

(i) By \cite[(4.14)]{B-CMP109}, $D\psi_y(\mathbf{1})=0$. Let first $\omega$ be finitely supported and
$u_\varepsilon:=\exp(i\varepsilon\omega t_0)$. Gauge invariance gives
$\psi_y\bigl((e^{iB})^{u_\varepsilon}\bigr)=\psi_y(e^{iB})$ for all small $\varepsilon$ and $B$. The chart
coordinate of $(e^{iB})^{u_\varepsilon}$ is analytic in $(\varepsilon,B)$ near $(0,0)$. At $B=0$ the
factors commute and $\mathbf{1}^{u_\varepsilon}=\exp(i\varepsilon t_0\,d\omega)$, so the
$\varepsilon$-derivative at $\varepsilon=0$ of this coordinate is $t_0\,d\omega$.

Differentiate the identity once in $\varepsilon$ at $\varepsilon=0$ and then once in $B$ at $B=0$ in the
direction $t_0\delta_b$. This gives $D^2\psi_y(\mathbf{1})[t_0\delta_b,t_0\,d\omega]$ plus a term in which
$D\psi_y(\mathbf{1})$ is applied to a mixed derivative of the chart coordinate. The sum vanishes, and the
second term vanishes because $D\psi_y(\mathbf{1})=0$. Since $d\omega$ is finitely supported, bilinearity
turns the first term into $\sum_{b'}K_y(b,b')(d\omega)(b')$, which is therefore $0$.

For polynomially bounded $\omega$, apply this to the truncations $\omega_R$ of $\omega$ to the ball of
radius $R$ about $y$. We have $(d\omega_R)(b')=(d\omega)(b')$ for $b'$ well inside that ball, and
$|d\omega_R|$ is bounded by a polynomial uniformly in $R$. Dominated convergence, which holds by
\eqref{10.6:eq-second-moment-1}, lets $R\to\infty$.

(ii) Two affine fields of curvature $\mathsf{f}$ differ by $d\omega$, with $\omega$ a polynomial of degree
at most two. Replacing $\mathsf{a}$ by $\mathsf{a}+d\omega$ in the quadratic form changes it by
$\sum_{b,b'}K_y(b,b')\mathsf{a}(b)(d\omega)(b')+\frac12\sum_{b,b'}K_y(b,b')(d\omega)(b)(d\omega)(b')$,
using the symmetry of $K_y$. Both sums vanish by (i), summed first over $b'$. Hence the value does not
depend on the centre of the affine field.

Next let $a$ be a lattice vector. Write the value at $y+a$ with $b=c+a$, $b'=c'+a$, and use
$K_{y+a}(c+a,c'+a)=K_y(c,c')$. It becomes the quadratic form of $K_y$ evaluated on the translate
$c\mapsto\mathsf{a}^{y+a}_{\mathsf{f}}(c+a)$, which is again an affine field of curvature $\mathsf{f}$. By
what was just shown, this equals $\mathsf{cf}[\psi](\mathsf{f})$ computed at $y$.

(iii) Let $e_\nu$ be the unit vector in direction $\nu$. Constant fields are of the form $d\omega$ with
$\omega$ linear: for $\omega(x)=(x-y)_\nu/\ell_j$, $d\omega$ equals $1$ on bonds parallel to $\nu$ and
$0$ on the others. By (i) and the symmetry of $K_y$,
\begin{equation}\label{10.6:eq-second-moment-3}
\sum_{b'\parallel\nu}K_y(b,b')=0,\qquad \sum_{b\parallel\mu}K_y(b,b')=0 .
\end{equation}

Let $S=(S_{\kappa\rho})$ be a real symmetric matrix and
$\omega(x):=(2\ell_j)^{-1}\sum_{\kappa,\rho}S_{\kappa\rho}(x-y)_\kappa(x-y)_\rho$. For the bond $b\parallel\mu$
from $x$ to $x+\ell_je_\mu$, with midpoint $x_b=x+\frac12\ell_je_\mu$,
\begin{equation*}
(d\omega)(b)=\sum_\kappa S_{\kappa\mu}(x-y)_\kappa+\tfrac12\ell_jS_{\mu\mu}
=\sum_\kappa S_{\kappa\mu}(x_b-y)_\kappa .
\end{equation*}
Put $T_{\nu\lambda}(b):=\sum_{b'\parallel\nu}K_y(b,b')(x_{b'}-y)_\lambda$. Then (i) gives
$\sum_{\nu,\lambda}S_{\lambda\nu}T_{\nu\lambda}(b)=0$ for every symmetric $S$, so
$T_{\nu\lambda}(b)=-T_{\lambda\nu}(b)$. Multiplying by $\ell_j^2(x_b-y)_\kappa$ and summing over
$b\parallel\mu$ gives $m_{\mu\kappa,\nu\lambda}=-m_{\mu\kappa,\lambda\nu}$. The same argument in the first
variable, using the symmetry of $K_y$, gives $m_{\mu\kappa,\nu\lambda}=-m_{\kappa\mu,\nu\lambda}$.

Let $\Pi(b,b'):=\sum_yK_y(b,b')$ be the kernel of the Hessian of $\Psi$ at $\mathbf{1}$ in the directions
$t_0\delta_b,t_0\delta_{b'}$. It is translation invariant. For $b\parallel\mu$, $b'\parallel\nu$ write
$\Pi_{\mu\nu}(b,b')$ and $z:=x_{b'}-x_b$. The Fourier transform
$\sum_{b'\parallel\nu}\Pi_{\mu\nu}(b,b')e^{ip\cdot z}$ does not depend on the choice of $b\parallel\mu$. Its
part of order $p^2$ is $-\frac12\sum_{\kappa,\lambda}p_\kappa p_\lambda\sum_{b'\parallel\nu}
\Pi_{\mu\nu}(b,b')z_\kappa z_\lambda$.

Fix $y$. Translation covariance rewrites $\sum_{y'}K_{y'}(b,b')$ as a sum of $K_y$ over the translates of
the pair $(b,b')$, and $z$ is translation invariant. As $b$ runs over the bonds parallel to $\mu$ in one
cell, these translates exhaust all pairs. Hence
\begin{equation*}
\sum_{b'\parallel\nu}\Pi_{\mu\nu}(b,b')z_\kappa z_\lambda
=\sum_{b\parallel\mu,\ b'\parallel\nu}K_y(b,b')z_\kappa z_\lambda .
\end{equation*}
Write $z=(x_{b'}-y)-(x_b-y)$ and expand. The term $(x_{b'}-y)_\kappa(x_{b'}-y)_\lambda$ vanishes on summing
over $b$, and the term $(x_b-y)_\kappa(x_b-y)_\lambda$ on summing over $b'$, both by
\eqref{10.6:eq-second-moment-3}. The two cross terms give $-(m_{\mu\kappa,\nu\lambda}+m_{\mu\lambda,\nu\kappa})
\ell_j^{-2}$. The part of order $p^2$ is therefore
$\frac12\ell_j^{-2}\sum_{\kappa,\lambda}p_\kappa p_\lambda(m_{\mu\kappa,\nu\lambda}+m_{\mu\lambda,\nu\kappa})
=\ell_j^{-2}\sum_{\kappa,\lambda}p_\kappa p_\lambda m_{\mu\kappa,\nu\lambda}$, the last step by the symmetry
of $p_\kappa p_\lambda$. Thus $m$ is the tensor of second moments of the Hessian of $\Psi$ from which the
typed coefficient $\beta^{\mathbf{3}}$ is read.

It remains to express $\mathsf{cf}[\psi](\mathsf{f})$ through $m$. For $b\parallel\mu$ the affine field of
Notation~\ref{10.6:not-affine-fields} is
$\mathsf{a}^{y}_{\mathsf{f}}(b)=\frac{\ell_j}{2}\sum_\kappa\mathsf{f}_{\kappa\mu}(x_b-y)_\kappa$. Inserting
this into the definition in (ii) and grouping the bonds by direction,
\begin{equation*}
\begin{split}
\mathsf{cf}[\psi](\mathsf{f})
&=\tfrac12\Bigl(\tfrac{\ell_j}{2}\Bigr)^2\sum_{\kappa,\lambda,\mu,\nu}\mathsf{f}_{\kappa\mu}
\mathsf{f}_{\lambda\nu}\sum_{b\parallel\mu,\ b'\parallel\nu}K_y(b,b')(x_b-y)_\kappa(x_{b'}-y)_\lambda\\
&=\tfrac18\sum_{\kappa,\lambda,\mu,\nu}\mathsf{f}_{\kappa\mu}\mathsf{f}_{\lambda\nu}
m_{\mu\kappa,\nu\lambda} .
\end{split}
\end{equation*}

(iv) By (iii), $m$ is antisymmetric in $(\mu\kappa)$ and in $(\nu\lambda)$; exchanging the names of $b$
and $b'$ in its definition and using $K_y(b,b')=K_y(b',b)$ gives
$m_{\mu\kappa,\nu\lambda}=m_{\nu\lambda,\mu\kappa}$. Hence $m$ is an element of
$\operatorname{Sym}^2\Lambda^2$. Since $\Psi$ is $\mathsf{V}_{\mathbf{3}}$-typed, the one-dimensionality of
the $W_4$-equivariant maps from the module $\mathsf{V}_{\mathbf{3}}$ into
$\operatorname{Sym}^2\Lambda^2$ and the definition of $\beta^{\mathbf{3}}$ read in the lattice units of
level $j$ give
\begin{equation*}
m_{\mu\kappa,\nu\lambda}=\beta^{\mathbf{3}}[\Psi]\,(h_\mu+h_\kappa)
(\delta_{\mu\nu}\delta_{\kappa\lambda}-\delta_{\mu\lambda}\delta_{\kappa\nu})\,\ell_j^4 .
\end{equation*}
The normalisation of this display is the one fixed by the definition of $\beta^{\mathbf{3}}$: with
$\tilde p:=\ell_jp$,
$\ell_j^{-2}\sum_{\kappa,\lambda}p_\kappa p_\lambda m_{\mu\kappa,\nu\lambda}
=\beta^{\mathbf{3}}[\Psi]\bigl(\delta_{\mu\nu}\sum_\kappa(h_\mu+h_\kappa)\tilde p_\kappa^{\,2}
-(h_\mu+h_\nu)\tilde p_\mu\tilde p_\nu\bigr)$, which is $\beta^{\mathbf{3}}[\Psi]$ times the symbol of
$\mathsf{T}^{\mathbf{3}}$, and on $\mathsf{A}^{\mathbf{3}}$ itself $\beta^{\mathbf{3}}=1$.

Insert this into (iii). Since $\mathsf{f}$ is antisymmetric,
$\sum_{\kappa,\lambda,\mu,\nu}\mathsf{f}_{\kappa\mu}\mathsf{f}_{\lambda\nu}(h_\mu+h_\kappa)
(\delta_{\mu\nu}\delta_{\kappa\lambda}-\delta_{\mu\lambda}\delta_{\kappa\nu})
=\sum_{\mu,\kappa}(h_\mu+h_\kappa)(\mathsf{f}_{\kappa\mu}^2-\mathsf{f}_{\kappa\mu}\mathsf{f}_{\mu\kappa})
=2\sum_{\mu,\kappa}(h_\mu+h_\kappa)\mathsf{f}_{\kappa\mu}^2$. Hence
\begin{equation*}
\mathsf{cf}[\psi](\mathsf{f})=\tfrac18\cdot2\sum_{\mu,\kappa}(h_\mu+h_\kappa)\mathsf{f}_{\kappa\mu}^2\,
\ell_j^4\,\beta^{\mathbf{3}}[\Psi]
=\beta^{\mathbf{3}}[\Psi]\cdot\tfrac12\sum_{\mu<\kappa}(h_\mu+h_\kappa)(\ell_j^2\mathsf{f}_{\kappa\mu})^2 .
\end{equation*}
The right-hand side is $\beta^{\mathbf{3}}[\Psi]\cdot\mathsf{t}^{\mathbf{3}}(h;\ell_j^2\mathsf{f})$, which is
\eqref{10.6:eq-second-moment-2}.
\end{proof}

\begin{remark}[Inner maps enter only through their linearisations]\label{10.6:rem-inner-maps}
Let $U_j$ be the inner map of clause (CF4) of Lemma~\ref{10.6:lem-constant-fields}, carrying
configurations on $\mathbb{L}_j$ to configurations on the fine lattice and sending $\mathbf{1}$ to
$\mathbf{1}$. Write $H_j$ for its linearisation at $\mathbf{1}$, and let
$\psi_y=\tilde{\psi}_y\circ U_j$. Here $\tilde{\psi}=(\tilde{\psi}_y)$ are gauge-invariant
functions, analytic near $\mathbf{1}$, of configurations on the fine lattice, and
$\psi=(\psi_y)$ satisfies the hypotheses of Lemma~\ref{10.6:lem-second-moment}. Then, for bond
fields $B$ on $\mathbb{L}_j$,
\begin{equation}\label{10.6:eq-inner-maps-1}
D^2\psi_y(\mathbf{1})[B,B]=D^2\tilde{\psi}_y(\mathbf{1})[H_jB,H_jB],
\end{equation}
so the second-order terms of $U_j$ do not contribute. The identity \eqref{10.6:eq-inner-maps-1}
uses only that the inner map sends $\mathbf{1}$ to $\mathbf{1}$ and that $\tilde{\psi}_y$ is gauge
invariant, so it holds for every inner map with these two properties. Consequently
$\mathsf{cf}[\tilde{\psi}\circ U_j](\mathsf{f})$ is the second Taylor coefficient of
$\tilde{\psi}_y$ along $\mathcal{U}_{\mathsf{f}}$ on the fine lattice.
\end{remark}

\begin{proof}
We prove \eqref{10.6:eq-inner-maps-1}. The
chain rule at $\mathbf{1}$ expresses $D^2(\tilde{\psi}_y\circ U_j)(\mathbf{1})[B,B]$ as the sum
of two terms. The first is the right side of \eqref{10.6:eq-inner-maps-1}. The second is the
first derivative of $\tilde{\psi}_y$ at $\mathbf{1}$ applied to $D^2U_j(\mathbf{1})[B,B]$. This
second term vanishes. Indeed, $\tilde{\psi}_y$ is gauge invariant, so $D\tilde{\psi}_y(\mathbf{1})=0$
by \cite[(4.14)]{B-CMP109}.

For the consequence, recall that
$K_y(b,b')=D^2\psi_y(\mathbf{1})[t_0\delta_b,t_0\delta_{b'}]$. The kernel $K_y$ decays
exponentially, so the defining sum of Lemma~\ref{10.6:lem-second-moment}\,(ii) converges
absolutely. This extends the bilinear form $D^2\psi_y(\mathbf{1})$ to the affine field, as in the
proof of Lemma~\ref{10.6:lem-second-moment}\,(i), and \eqref{10.6:eq-inner-maps-1} is read for the
form so extended. By bilinearity of the second derivative,
\[
\mathsf{cf}[\psi](\mathsf{f})=\tfrac12\sum_{b,b'}K_y(b,b')\,\mathsf{a}^{y}_{\mathsf{f}}(b)\,
\mathsf{a}^{y}_{\mathsf{f}}(b')
=\tfrac12\,D^2\psi_y(\mathbf{1})[t_0\mathsf{a}^{y}_{\mathsf{f}},t_0\mathsf{a}^{y}_{\mathsf{f}}].
\]
By \eqref{10.6:eq-inner-maps-1} this equals
$\tfrac12\,D^2\tilde{\psi}_y(\mathbf{1})[H_jt_0\mathsf{a}^{y}_{\mathsf{f}},
H_jt_0\mathsf{a}^{y}_{\mathsf{f}}]$. Two facts now identify this quantity. The first is clause
(CF4) of Lemma~\ref{10.6:lem-constant-fields}, applied to the image
$H_jt_0\mathsf{a}^{y}_{\mathsf{f}}$ of the affine field under the linearised inner map. The
second is the independence of $\mathsf{cf}$ from $y$ and from the centre of the affine field,
which is Lemma~\ref{10.6:lem-second-moment}\,(ii). Together they show that the quantity is the
second Taylor coefficient of $\tilde{\psi}_y$ along $\mathcal{U}_{\mathsf{f}}$ on the fine
lattice.
\end{proof}

\begin{lemma}[Terms linear in a vertex gradient drop]\label{10.6:lem-gradient-terms-drop}
Let $W$ denote a configuration on $\mathbb{L}_{j+1}$, and let
\begin{equation}\label{10.6:eq-gradient-terms-drop-1}
T(W)=\sum_{b}\bigl\langle \mathsf{s}(W)(b),\mathsf{w}(W)(b)\bigr\rangle ,
\end{equation}
the sum running over the bonds $b$ of $\mathbb{L}_j$. Suppose that $\mathsf{s}$ and $\mathsf{w}$ satisfy the following:
\begin{itemize}
\item they are analytic near $\mathbf{1}$ and translation covariant under $\mathbb{L}_{j+1}$;
\item $\mathsf{s}(\mathbf{1})=\mathsf{w}(\mathbf{1})=0$;
\item their linearised kernels
\[
S(b,c):=D\mathsf{s}(\mathbf{1})(b)[\delta_c],\qquad
\mathsf{K}_{\mathsf{w}}(b,c):=D\mathsf{w}(\mathbf{1})(b)[\delta_c]
\]
are exponentially decaying.
\end{itemize}
If
\begin{gather*}
\sum_c S(b,c)\,\mathsf{a}(c)=0\quad\text{for every affine field }\mathsf{a},\\
\sum_c \mathsf{K}_{\mathsf{w}}(b,c)\,\xi(c)=0\quad\text{for every constant field }\xi,
\end{gather*}
then every second moment of the Hessian of $T$ at $\mathbf{1}$ vanishes. Consequently $\beta[T]=0$ and $\beta^{r}[T]=0$.
\end{lemma}

Here affine fields are those introduced in~\ref{10.6:not-affine-fields}, and $\xi$ denotes a field that is constant in the index $c$ of the kernels.

\begin{proof}
We first compute the Hessian of $T$ at $\mathbf{1}$. Differentiate \eqref{10.6:eq-gradient-terms-drop-1} twice at $\mathbf{1}$ in the directions $\delta_c,\delta_{c'}$. The Leibniz rule produces four terms. Two of them pair a second derivative of one factor with the other factor evaluated at $\mathbf{1}$. These two terms vanish, because $\mathsf{s}(\mathbf{1})=\mathsf{w}(\mathbf{1})=0$. What remains is the Hessian kernel $\Pi^{T}$ of the functional $T$ of \eqref{10.6:eq-gradient-terms-drop-1},
\begin{equation}\label{10.6:eq-gradient-terms-drop-2}
\Pi^{T}(c,c')=\sum_b\Bigl[\bigl\langle S(b,c),\mathsf{K}_{\mathsf{w}}(b,c')\bigr\rangle
+\bigl\langle \mathsf{K}_{\mathsf{w}}(b,c),S(b,c')\bigr\rangle\Bigr],
\end{equation}
where the bracket is the pairing of \eqref{10.6:eq-gradient-terms-drop-1} applied to the two first derivatives.

Next we record three identities that follow from the hypotheses. Every constant field is affine. For each $\kappa$, the field $c\mapsto (x_c-x_b)_\kappa$, with $b$ held fixed, is affine. Applying the hypothesis on $S$ to these fields, and the hypothesis on $\mathsf{K}_{\mathsf{w}}$ to constant fields, gives, for every bond $b$ and every $\kappa$,
\begin{equation}\label{10.6:eq-gradient-terms-drop-3}
\sum_c S(b,c)=0,\qquad
\sum_c S(b,c)\,(x_c-x_b)_\kappa=0,\qquad
\sum_{c'}\mathsf{K}_{\mathsf{w}}(b,c')=0 .
\end{equation}

Fix $\kappa,\lambda$. The second moment of the Hessian is
\[
\sum_{c}\sum_{c'}\Pi^{T}(c,c')\,(x_{c'}-x_c)_\kappa(x_{c'}-x_c)_\lambda ,
\]
where $c$ runs over the bonds in one cell (the normalisation in which the second moments entering $\beta[\cdot]$ are taken) and $c'$ over all bonds. Insert \eqref{10.6:eq-gradient-terms-drop-2}. The kernels decay exponentially, so the triple sum over $b,c,c'$ converges absolutely and its order may be changed. The weight depends only on $x_{c'}-x_c$. Hence translation covariance under $\mathbb{L}_{j+1}$ lets us rewrite the sum so that $b$ runs over the bonds in one cell while $c,c'$ run over all bonds. It therefore suffices to show that, for each fixed $b$, the sum over all $(c,c')$ vanishes.

Set $u=u(c):=x_c-x_b$ and $u'=u'(c'):=x_{c'}-x_b$. Then $x_{c'}-x_c=u'-u$, and
\[
(x_{c'}-x_c)_\kappa(x_{c'}-x_c)_\lambda=u'_\kappa u'_\lambda-u'_\kappa u_\lambda-u_\kappa u'_\lambda+u_\kappa u_\lambda .
\]

Consider first the term $\langle S(b,c),\mathsf{K}_{\mathsf{w}}(b,c')\rangle$.
\begin{itemize}
\item The summand $u'_\kappa u'_\lambda$ does not depend on $c$. Summing first over $c$ gives
$\bigl\langle\sum_c S(b,c),\mathsf{K}_{\mathsf{w}}(b,c')\bigr\rangle\,u'_\kappa u'_\lambda=0$, by the first
identity of \eqref{10.6:eq-gradient-terms-drop-3}.
\item In the cross summands $u'_\kappa u_\lambda$ and $u_\kappa u'_\lambda$, summing first over $c$ gives
$\bigl\langle\sum_c S(b,c)\,u_\lambda,\mathsf{K}_{\mathsf{w}}(b,c')\bigr\rangle\,u'_\kappa=0$, respectively
$\bigl\langle\sum_c S(b,c)\,u_\kappa,\mathsf{K}_{\mathsf{w}}(b,c')\bigr\rangle\,u'_\lambda=0$, by the second
identity of \eqref{10.6:eq-gradient-terms-drop-3}.
\item The summand $u_\kappa u_\lambda$ does not depend on $c'$. Summing first over $c'$ gives
$\bigl\langle S(b,c),\sum_{c'}\mathsf{K}_{\mathsf{w}}(b,c')\bigr\rangle\,u_\kappa u_\lambda=0$, by the third
identity of \eqref{10.6:eq-gradient-terms-drop-3}.
\end{itemize}
Each step uses \eqref{10.6:eq-gradient-terms-drop-3}.

The term $\langle \mathsf{K}_{\mathsf{w}}(b,c),S(b,c')\rangle$ is handled in the same way with the roles of $c$ and $c'$ exchanged.
\begin{itemize}
\item The summand $u_\kappa u_\lambda$ vanishes after summation over $c'$, by $\sum_{c'}S(b,c')=0$.
\item The cross summands vanish after summation over $c'$, by $\sum_{c'}S(b,c')\,u'_\kappa=0$ and
$\sum_{c'}S(b,c')\,u'_\lambda=0$.
\item The summand $u'_\kappa u'_\lambda$ vanishes after summation over $c$, by
$\sum_c\mathsf{K}_{\mathsf{w}}(b,c)=0$.
\end{itemize}

Hence every second moment of the Hessian of $T$ at $\mathbf{1}$ vanishes. The coefficients $\beta[T]$ and $\beta^{r}[T]$ are read off from these second moments, so both vanish.
\end{proof}

\begin{lemma}[The one-loop number as a half trace]\label{10.6:lem-half-trace}
Let $0\le j<k$. Put $\mathcal{A}_j:=A\circ U_j$ and
$\mathcal{S}_j:=\mathsf{A}^{\mathbf{3}}(1,U_j(\cdot);h_0)$, with $h_0$ the vector of orientation
weights. These are functions of configurations on
$\mathbb{L}_j$, and $\mathcal{A}_0=A$, $\mathcal{S}_0=\mathsf{A}^{\mathbf{3}}$.

For a configuration $W$ on $\mathbb{L}_{j+1}$ we use the following objects.
\begin{itemize}
\item $V^{(j)}(W)$ is the critical point of \cite[(2.2)--(2.3)]{B-CMP109}.
\item $\nabla$ and $\nabla^2$ are the first and second derivatives in the chart
$B'\mapsto e^{iB'}V^{(j)}(W)$.
\item $C$ is Ba{\l}aban's linear map from the fluctuation variables $B$ to bond fields,
\cite[(2.5)]{B-CMP109}; to first order in the source-strength parameter $\mathfrak{q}$
introduced below, the chart variable is $B'=CB$.
\item $d\mu_{C^{(j)}}$ is the Gaussian measure in the variables $B$ with covariance
$C^{(j)}=(C^*\Delta^{(j)}C)^{-1}$, where $\Delta^{(j)}$ is the operator of the quadratic form in
$B'$ of \cite[(2.5)]{B-CMP109}, expanded in \cite[(1.5)]{B-CMP109}.
\item The fluctuation covariance is
\[
\mathcal{G}_P^{(j+1)}(W):=C\,(C^*\Delta^{(j)}C)^{-1}C^* ,
\]
which is the covariance of $B'=CB$ under $d\mu_{C^{(j)}}$.
\item $\langle\cdot\rangle_0$ is the expectation in the normalised measure $d\mu_{C^{(j)}}$,
Gaussian and without field cut, and $\langle\cdot\,;\cdot\rangle_0$ is the corresponding truncated
expectation.
\end{itemize}

Let $h$ be Ba{\l}aban's operator of \cite[(1.5)]{B-CMP109} (not the weight vector $h_0$), let
$\tilde{D}$ be the nonlinear map entering Ba{\l}aban's change of variables in \cite{B-CMP109}, and
let $\tilde{C}^{(2)}$ be the map, homogeneous of degree two, appearing in \cite[(1.5)]{B-CMP109}.
With $\mathfrak{q}$ the source-strength parameter of the vertex, put
$B'=\mathfrak{q}CB-h\tilde{D}(\mathfrak{q}CB)$. The vertex is
$\mathcal{S}_j\bigl(e^{iB'}V^{(j)}(W)\bigr)$ minus its value at $\mathfrak{q}=0$; it therefore has
an expansion
$\mathfrak{q}\mathcal{S}_j^{[1]}(B)+\mathfrak{q}^2\mathcal{S}_j^{[2]}(B)+O(\mathfrak{q}^3)$; thus
$\mathcal{S}_j^{[1]}$ and $\mathcal{S}_j^{[2]}$ are the Taylor coefficients in $\mathfrak{q}$ of the
vertex at level $j$. The remaining part
$V_{\mathrm{rem}}$ of the exponent of the fluctuation integrand satisfies
$V_{\mathrm{rem}}(\mathfrak{q}B)|_{u=0}=\mathfrak{q}V_{\mathrm{rem}}^{[1]}(B)+O(\mathfrak{q}^2)$,
where $u$ denotes the variables of Ba{\l}aban's exponential axial gauge fixing,
\cite[(0.14)--(0.16)]{B-CMP109}.
\begin{enumerate}
\item[(a)] One has
\[
\langle\mathcal{S}_j^{[2]}\rangle_0+\langle\mathcal{S}_j^{[1]};V_{\mathrm{rem}}^{[1]}\rangle_0
=\tfrac12\operatorname{Tr}\bigl[\mathcal{G}_P^{(j+1)}(W)\,\nabla^2\mathcal{S}_j(V^{(j)}(W))\bigr]
-\langle\nabla\mathcal{S}_j,\mathsf{w}_1\rangle+\langle\nabla\mathcal{S}_j,\mathsf{w}_2\rangle .
\]
Here $\mathsf{w}_1:=\langle h\tilde{C}^{(2)}(CB)\rangle_0$ and
$\mathsf{w}_2:=\langle CB\cdot V_{\mathrm{rem}}^{[1]}(B)\rangle_0$, and $\nabla\mathcal{S}_j$ is
taken at $V^{(j)}(W)$.
\item[(b)] The typed coefficient $\beta^{\mathbf{3}}$ of each of the last two terms in (a)
vanishes.
\item[(c)] For $y\in\mathbb{L}_{j+1}$ let $\square_{j+1}(y)$ be the cell of $y$, and put
\[
\psi^{(j+1)}_y(W):=\tfrac12\sum_{b\in\square_{j+1}(y)}
\operatorname{tr}\bigl(\mathcal{G}_P^{(j+1)}\nabla^2\mathcal{S}_j\bigr)(b,b),
\]
where $\operatorname{tr}$ here denotes the trace, over the fibre at the bond $b$, of the diagonal
block at $(b,b)$ of the kernel of $\mathcal{G}_P^{(j+1)}\nabla^2\mathcal{S}_j$ (not the normalised
trace of the fundamental representation), so that
$\sum_y\psi^{(j+1)}_y=\tfrac12\operatorname{Tr}[\mathcal{G}_P^{(j+1)}\nabla^2\mathcal{S}_j]$.
Then
\begin{equation}\label{10.6:eq-half-trace-a}
\mathsf{m}^{\circ}_{\mathbf{3}}(j+1)\cdot\mathsf{t}^{\mathbf{3}}(h_0;\ell_{j+1}^2\mathsf{f})
=-\,\mathsf{cf}[\psi^{(j+1)}](\mathsf{f}),
\end{equation}
where $\mathsf{m}^{\circ}_{\mathbf{3}}(j+1)$ is the cut-free number of
Definition~\ref{10.6:not-cut-free-numbers}.
\end{enumerate}
\end{lemma}

\begin{proof}
\emph{(a)} The homogeneous second-order part of the map $\tilde{D}$ is $\tilde{C}^{(2)}$, by the
identification of the second-order parts of these maps in \cite{B-CMP109}. Hence
\[
B'=\mathfrak{q}CB-h\tilde{D}(\mathfrak{q}CB)
=\mathfrak{q}CB-\mathfrak{q}^2h\tilde{C}^{(2)}(CB)+O(\mathfrak{q}^3).
\]

Expand $\mathcal{S}_j(e^{iB'}V^{(j)}(W))$ in the chart to second order in $B'$:
it equals $\mathcal{S}_j(V^{(j)}(W))$ plus
\[
\langle\nabla\mathcal{S}_j,B'\rangle+\tfrac12\nabla^2\mathcal{S}_j[B',B']
\]
plus terms of order three in $B'$. Insert the expansion of $B'$ and collect powers of
$\mathfrak{q}$. This gives
\[
\mathcal{S}_j^{[1]}(B)=\langle\nabla\mathcal{S}_j,CB\rangle,\qquad
\mathcal{S}_j^{[2]}(B)=\tfrac12\nabla^2\mathcal{S}_j[CB,CB]
-\langle\nabla\mathcal{S}_j,h\tilde{C}^{(2)}(CB)\rangle .
\]

Since $\mathcal{S}_j^{[1]}$ is linear in $B$ and $d\mu_{C^{(j)}}$ is centred,
$\langle\mathcal{S}_j^{[1]}\rangle_0=0$.

The covariance of $CB$ under $d\mu_{C^{(j)}}$ is $CC^{(j)}C^*=\mathcal{G}_P^{(j+1)}(W)$. Therefore
\[
\tfrac12\langle\nabla^2\mathcal{S}_j[CB,CB]\rangle_0
=\tfrac12\operatorname{Tr}\bigl[\mathcal{G}_P^{(j+1)}(W)\nabla^2\mathcal{S}_j\bigr].
\]
The gradient is evaluated at $V^{(j)}(W)$ and does not depend on $B$, so
\[
\langle\langle\nabla\mathcal{S}_j,h\tilde{C}^{(2)}(CB)\rangle\rangle_0
=\langle\nabla\mathcal{S}_j,\mathsf{w}_1\rangle .
\]

Because $\langle\mathcal{S}_j^{[1]}\rangle_0=0$, we have
$\langle\mathcal{S}_j^{[1]};V_{\mathrm{rem}}^{[1]}\rangle_0
=\langle\mathcal{S}_j^{[1]}V_{\mathrm{rem}}^{[1]}\rangle_0$. This equals
\[
\langle\langle\nabla\mathcal{S}_j,CB\rangle V_{\mathrm{rem}}^{[1]}(B)\rangle_0
=\langle\nabla\mathcal{S}_j,\mathsf{w}_2\rangle .
\]
Adding the three contributions gives (a).

\emph{(b)} We apply the lemma that terms linear in a vertex gradient drop
(Lemma~\ref{10.6:lem-gradient-terms-drop}) with
\[
\mathsf{s}(W):=\nabla\mathcal{S}_j(V^{(j)}(W))
\quad\text{and}\quad
\mathsf{w}:=\mathsf{w}_i,\ i=1,2 .
\]
Then $T=\langle\nabla\mathcal{S}_j,\mathsf{w}_i\rangle$ has the form required there.

First, $\mathsf{s}(\mathbf{1})=0$ by \cite[(4.14)]{B-CMP109}.

Second, we check that the kernel $S$ annihilates affine fields. For a field $\mathsf{a}$ the chain rule gives
\[
D\mathsf{s}(\mathbf{1})[\mathsf{a}]
=\nabla^2\mathcal{S}_j(\mathbf{1})\cdot DV^{(j)}(\mathbf{1})\mathsf{a} .
\]
By Remark~\ref{10.6:rem-inner-maps}, which accompanies the second-moment lemma
(Lemma~\ref{10.6:lem-second-moment}), with $H_j$ and the lattice coboundary $d$ as there,
\[
\nabla^2\mathcal{S}_j(\mathbf{1})=H_j^*(d^*\mathsf{c}_{h_0}d)H_j .
\]
Now let $\mathsf{a}$ be affine. By the last clause of Lemma~\ref{10.6:lem-constant-fields},
\eqref{10.6:eq-constant-fields-a},
$H_jDV^{(j)}(\mathbf{1})\mathsf{a}$ is affine on $\mathbb{L}_0$ modulo a field $d\lambda$.
\begin{itemize}
\item The field $d\lambda$ is annihilated by $d$, since $d\circ d=0$.
\item The image under $d$ of an affine field is constant on each orientation.
\end{itemize}
Hence $\mathsf{c}_{h_0}d(H_jDV^{(j)}(\mathbf{1})\mathsf{a})$ is constant on each orientation, and
$d^*$ of it vanishes. Thus $\sum_cS(b,c)\mathsf{a}(c)=0$ for every affine $\mathsf{a}$.

Third, we check that the kernel of $\mathsf{w}_i$ annihilates constant fields. The vector
$\mathsf{w}_i(\mathbf{1})$ is an $\operatorname{Ad}$-invariant vector of $\mathfrak{g}$. Since
$\mathfrak{g}$ is semisimple, $\mathsf{w}_i(\mathbf{1})=0$. The map $\mathsf{w}_i$ is gauge
covariant, so
\[
D\mathsf{w}_i(\mathbf{1})[d\lambda]=i[\lambda,\mathsf{w}_i(\mathbf{1})]=0 .
\]
This holds in particular on constant fields. Hence
$\sum_c\mathsf{K}_{\mathsf{w}}(b,c)\beta(c)=0$ for every constant field $\beta$.

Fourth, the exponential decay of the kernels $S$ and $\mathsf{K}_{\mathsf{w}}$ follows from the
theorem on the exponential locality of the step's kernels in a background field.

Lemma~\ref{10.6:lem-gradient-terms-drop} therefore gives
$\beta^{\mathbf{3}}[\langle\nabla\mathcal{S}_j,\mathsf{w}_i\rangle]=0$ for $i=1,2$.

Whatever $V_{\mathrm{rem}}^{[1]}$ contains enters through $\mathsf{w}_2$ only. This covers the cubic
term of the action, the gradients of the Haar measure and of the Jacobian, and the gradients of the
older one-loop terms $L^{(\le j)}$, written as in Definition~\ref{10.6:not-cut-free-numbers}. So all
of these are covered by the case $i=2$.

\emph{(c)} By Definition~\ref{10.6:not-cut-free-numbers}, $\mathsf{m}^{\circ}_{\mathbf{3}}(j+1)$ is the
cut-free number at step $j+1$. By Proposition~\ref{10.4:prop-vertex-response}, whose proof is
incomplete at one step and which identifies the total of the one-loop family with the one-loop
response, that number is $-\beta^{\mathbf{3}}$ of the total one-loop response,
\[
\mathsf{m}^{\circ}_{\mathbf{3}}(j+1)
=-\beta^{\mathbf{3}}\bigl[\langle\mathcal{S}_j^{[2]}\rangle_0
+\langle\mathcal{S}_j^{[1]};V_{\mathrm{rem}}^{[1]}\rangle_0\bigr],
\]
with the uncut Gaussian expectation.

The coefficient $\beta^{\mathbf{3}}$ is additive, and it is read from the total rather than from
the separate terms; this is the reading statement for $\beta^{\mathbf{3}}$ in the proposition on
the one-loop family. So by (a) and (b),
\[
\mathsf{m}^{\circ}_{\mathbf{3}}(j+1)=-\beta^{\mathbf{3}}[\Psi^{(j+1)}],
\qquad
\Psi^{(j+1)}:=\tfrac12\operatorname{Tr}\bigl[\mathcal{G}_P^{(j+1)}\nabla^2\mathcal{S}_j\bigr].
\]
Grouping the bonds $b$ of the trace according to the cell $\square_{j+1}(y)$ containing them
gives $\Psi^{(j+1)}=\sum_y\psi^{(j+1)}_y$.

Each $\psi^{(j+1)}_y$ is gauge invariant by \cite[(2.16)]{B-CMP109}. It is local, with
exponentially decaying Hessian kernel, by the theorem on the exponential locality of the step's
kernels in a background field.

Clause (iv) of Lemma~\ref{10.6:lem-second-moment}, applied at level $j+1$ to the family
$\psi^{(j+1)}$, gives
\[
\beta^{\mathbf{3}}[\Psi^{(j+1)}]\cdot\mathsf{t}^{\mathbf{3}}(h_0;\ell_{j+1}^2\mathsf{f})
=\mathsf{cf}[\psi^{(j+1)}](\mathsf{f}).
\]
Combined with the previous display, this is \eqref{10.6:eq-half-trace-a}.

Finally, by cyclicity of the trace,
\[
\tfrac12\operatorname{Tr}\bigl[\mathcal{G}_P^{(j+1)}\nabla^2\mathcal{S}_j\bigr]
=\tfrac12\operatorname{Tr}\bigl(C^{(j)}C^*\nabla^2\mathcal{S}_jC\bigr).
\]
The right-hand side is the half trace $\tfrac12\operatorname{Tr}(C^{(j)}\Delta^{\otimes,D})$ whose
coefficient defines the one-loop number in the determinant form
$\det(\Delta^{(j)}-\mathsf{k}\,g_j^2\Delta^{\otimes,D})^{-1/2}$ of the one-loop term, with
$\mathsf{k}$ the second source-strength parameter, the one multiplying the inserted operator in the
fluctuation determinant; there $\Delta^{\otimes,D}=C^*\nabla^2\mathcal{S}_jC$ is the operator
inserted with weight $\mathsf{k}$.
\end{proof}

\begin{definition}[The truncated product and the trace per cell]
\label{10.6:def-truncated-product}
We use the notation fixed in Notation~\ref{10.6:not-affine-fields}: $k\ge 1$; the lattices $\mathbb{L}_j$,
$j=0,\dots,k$, of spacing $\ell_j=L^{j-k}$; $\eta=\ell_0=L^{-k}$, also written $N^{-1}$; a unit cell
$\square$ of
$\mathbb{L}_k$; the element $t_0\in\mathfrak{g}$ with $\operatorname{tr}t_0^2=1$; the real
antisymmetric $4\times4$ matrix $\mathsf{f}=(\mathsf{f}_{\kappa\lambda})$ with
$\mathsf{f}(u,v)=\sum_{\kappa,\lambda}\mathsf{f}_{\kappa\lambda}u_\kappa v_\lambda$; and the constant
abelian configuration $\mathcal{U}_{\mathsf{f}}$. Repeated indices
$\kappa,\lambda,\kappa',\lambda'$ are summed over $1,\dots,4$.

\emph{Sectors.} The complexified algebra decomposes as $\mathfrak{g}_{\mathbb{C}}=\bigoplus_q\mathfrak{g}_q$,
where $[t_0,X]=qX$ for
$X\in\mathfrak{g}_q$; the adjoint action $\operatorname{Ad}(e^{i\theta t_0})$ is multiplication by
$e^{iq\theta}$ on
$\mathfrak{g}_q$. By the sector identity for $C_{\mathfrak{g}}$,
$\sum_q q^2\dim\mathfrak{g}_q=\operatorname{tr}_{\mathfrak{g}}(\operatorname{ad}_{t_0}^2)=C_{\mathfrak{g}}$.
Every operator considered in what follows commutes with $\operatorname{Ad}(e^{i\theta t_0})$, hence
is a direct sum
over $q$; the sector $q$ is fixed until the contributions of the sectors are summed.

\emph{Cells and fibres.} A cell $\alpha$ is a site or a bond of some $\mathbb{L}_j$; its position
$x_\alpha$ is the site itself, or the midpoint $x_b$ of the bond $b$. In Ba{\l}aban's construction
the variable of a bond $b$ is attached to its initial point $b_-$; we transport it to $x_b$ along the
half bond, which is a unitary change of variables, bond by bond. Let $\mathcal{C}_j$, $\mathcal{X}_j$
denote the $\ell^2$ spaces of the sites, respectively the bonds, of $\mathbb{L}_j$ with values in
$\mathfrak{g}_q$ and weight $\ell_j^4$.

\emph{Reduced kernels.} On $\mathcal{U}_{\mathsf{f}}$ the phase of the straight segment from
$x_\alpha$ to $x_\beta$ is $\frac12\mathsf{f}(x_\alpha,x_\beta)$. An operator $T$ on the background
$\mathcal{U}_{\mathsf{f}}$ which is gauge covariant and covariant under the translations of
$\mathbb{L}_k$ has the form $T(\alpha,\beta)=e^{(iq/2)\mathsf{f}(x_\alpha,x_\beta)}t(\alpha,\beta)$,
where $t$, \emph{the reduced kernel} of $T$, is invariant under the ordinary translations of
$\mathbb{L}_k$. The reduced kernel of a product $TS$ is
\begin{equation}\label{10.6:eq-truncated-product-1}
(\alpha,\beta)\longmapsto\sum_\gamma t(\alpha,\gamma)\,s(\gamma,\beta)\,
e^{(iq/2)\mathsf{f}(x_\alpha-x_\gamma,\;x_\gamma-x_\beta)} .
\end{equation}

\emph{The ring.} $(\delta_\kappa t)(\alpha,\beta):=(x_\alpha-x_\beta)_\kappa\,t(\alpha,\beta)$, that is
$\delta_\kappa t=[x_\kappa,t]$ with $x_\kappa$ the multiplication by $(x_\alpha)_\kappa$; it is a
derivation. For $a\in\mathbb{N}^4$ put $\delta^{a}:=\delta_1^{a_1}\cdots\delta_4^{a_4}$.
Let $\mathfrak{A}$ be the set of bounded operators between the spaces $\mathcal{C}_j$,
$\mathcal{X}_j,\dots$ which commute with the translations of $\mathbb{L}_k$ and for which
$\delta^{a}t$ is bounded for every $a$ (such operators are called \emph{tame}; operators with
exponentially decaying kernels are tame). Put $\mathfrak{A}_{\mathsf{f}}:=\mathfrak{A}\otimes
\mathbb{C}[\mathsf{f}]/(\mathsf{f})^3$, with $\mathbb{C}[\mathsf{f}]$ the polynomial ring in the
entries $\mathsf{f}_{\kappa\lambda}$, $\kappa<\lambda$; $t=t^{(0)}+t^{(1)}+t^{(2)}$ is the
decomposition by degree in $\mathsf{f}$, and $[\mathsf{f}^2]t:=t^{(2)}$. The \emph{truncated
product} is
\begin{equation}\label{10.6:eq-truncated-product-a}
t\star s:=ts+\frac{iq}{2}\,\mathsf{f}_{\kappa\lambda}(\delta_\kappa t)(\delta_\lambda s)
-\frac{q^2}{8}\,\mathsf{f}_{\kappa\lambda}\mathsf{f}_{\kappa'\lambda'}
(\delta_\kappa\delta_{\kappa'}t)(\delta_\lambda\delta_{\lambda'}s)\quad\bmod(\mathsf{f})^3 .
\end{equation}
It has the following properties: $\star$ is associative; $(t\star s)^\dagger=s^\dagger\star t^\dagger$,
the entries of $\mathsf{f}$ being real; $t$ is $\star$-invertible if and only if $t^{(0)}$ is
invertible in $\mathfrak{A}$, and the coefficients of the inverse are finite sums of words in
$(t^{(0)})^{-1}$ and the $\delta^{a}t^{(n)}$.

\emph{Trace.} For $t$ mapping one of the spaces into itself,
$\operatorname{Tr}_{\square}t:=\sum_{\alpha\in\square}\operatorname{tr}t(\alpha,\alpha)$, the sum
running over the cells of the one type on which $t$ acts whose positions lie in the unit cell
$\square$ of $\mathbb{L}_k$ (taken half-open, so that its translates by $\mathbb{L}_k$ tile
$\mathbb{R}^4$), and $\operatorname{tr}$ being the trace of endomorphisms of the fibre
$\mathfrak{g}_q$. Since the phase in~\eqref{10.6:eq-truncated-product-1} vanishes at
$\alpha=\beta$,
\begin{equation}\label{10.6:eq-truncated-product-b}
\operatorname{Tr}_{\square}(t\star s)=\operatorname{Tr}_{\square}(ts)
=\operatorname{Tr}_{\square}(s\star t).
\end{equation}

\emph{Fourier form.} For kernels invariant under the translations of $\mathbb{L}_0$, between two of
the spaces of $\mathbb{L}_0$, write $[\alpha]$ for the class of the cell $\alpha$ modulo these
translations (one class for sites, four for bonds, one for each direction). For cells
$\alpha,\beta$ of fixed classes, $t(\alpha,\beta)$ depends only on $u=x_\alpha-x_\beta$, which
ranges over a translate of $\eta\mathbb{Z}^4$ depending only on the two classes; we write
$t_{[\alpha][\beta]}(u):=t(\alpha,\beta)$ and put
\[
\hat t_{[\alpha][\beta]}(\xi):=\sum_u e^{-i\xi u}\,t_{[\alpha][\beta]}(u),
\]
the sum over that translate. Thus $\hat t(\xi)$ is a matrix of linear maps between fibres, its rows
indexed by the classes of the cells of the range and its columns by those of the domain, and the
inversion formula
\begin{gather*}
t(\alpha,\beta)=\eta^4\int_{\mathrm{BZ}_\eta}\bar{d}\xi\,e^{i\xi(x_\alpha-x_\beta)}
\hat t_{[\alpha][\beta]}(\xi),\\
\bar{d}\xi:=\frac{d^4\xi}{(2\pi)^4},\qquad \mathrm{BZ}_\eta:=[-\pi/\eta,\pi/\eta]^4,
\end{gather*}
holds. Under this correspondence $\delta_\kappa$ becomes $i\partial_\kappa$, with
$\partial_\kappa=\partial/\partial \xi_\kappa$, and the truncated product becomes
\begin{equation}\label{10.6:eq-truncated-product-2}
\begin{split}
\widehat{t\star s}&=\hat t\hat s-\frac{iq}{2}\,\mathsf{f}_{\kappa\lambda}\,\partial_\kappa\hat t\,
\partial_\lambda\hat s\\
&\quad-\frac{q^2}{8}\,\mathsf{f}_{\kappa\lambda}\mathsf{f}_{\kappa'\lambda'}\,
\partial_\kappa\partial_{\kappa'}\hat t\,\partial_\lambda\partial_{\lambda'}\hat s ,
\end{split}
\end{equation}
the products of $\hat t$ and $\hat s$ and of their derivatives being matrix products; moreover
$\operatorname{Tr}_{\square}t=\int_{\mathrm{BZ}_\eta}\bar{d}\xi\,\operatorname{tr}\hat t(\xi)$, where
$\operatorname{tr}$ runs over the fibre $\mathfrak{g}_q$ and over the classes.
\end{definition}

\begin{proof}
\emph{The kernel formula.} The reduced kernel of $TS$ at $(\alpha,\beta)$ is
\[
e^{-(iq/2)\mathsf{f}(x_\alpha,x_\beta)}\sum_\gamma
e^{(iq/2)[\mathsf{f}(x_\alpha,x_\gamma)+\mathsf{f}(x_\gamma,x_\beta)]}\,t(\alpha,\gamma)\,s(\gamma,\beta).
\]
By bilinearity and
$\mathsf{f}(x_\gamma,x_\gamma)=0$ (antisymmetry),
$\mathsf{f}(x_\alpha-x_\gamma,x_\gamma-x_\beta)=\mathsf{f}(x_\alpha,x_\gamma)
-\mathsf{f}(x_\alpha,x_\beta)+\mathsf{f}(x_\gamma,x_\beta)$, whence the exponent
in~\eqref{10.6:eq-truncated-product-1}.

\emph{The product is well defined and has a unit.} By the Leibniz rule
$\delta_\kappa(ts)=(\delta_\kappa t)s+t(\delta_\kappa s)$, every $\delta^{a}(ts)$ is a finite sum of
products $(\delta^{a'}t)(\delta^{a''}s)$; so $\mathfrak{A}$ is closed under products, and under the
$\delta_\kappa$. Since $\delta_\kappa$ acts on the $\mathfrak{A}$ factor and commutes with
multiplication by polynomials in $\mathsf{f}$, and reduction mod $(\mathsf{f})^3$ is a ring
homomorphism, \eqref{10.6:eq-truncated-product-a} defines a $\mathbb{C}[\mathsf{f}]$-bilinear product
on $\mathfrak{A}_{\mathsf{f}}$. The identity $\mathbf{1}$ has kernel supported on $\alpha=\beta$, so
$\delta_\kappa \mathbf{1}=0$ and $\mathbf{1}\star s=s$, $t\star\mathbf{1}=t$.

\emph{Associativity.} By bilinearity it suffices to take $t_1,t_2,t_3\in\mathfrak{A}$. Since
$\frac12(\frac{iq}{2})^2=-\frac{q^2}{8}$, for $t,s\in\mathfrak{A}$ the right side
of~\eqref{10.6:eq-truncated-product-a} is the sum over $\gamma$ in~\eqref{10.6:eq-truncated-product-1}
with the exponential replaced by the terms of degree at most two of its Taylor series in
$\mathsf{f}$: the factors $(x_\alpha-x_\gamma)_\kappa$, $(x_\gamma-x_\beta)_\lambda$ are absorbed as
$\delta_\kappa t$, $\delta_\lambda s$. Write
$\Xi_1:=\mathsf{f}(x_\alpha-x_\gamma,x_\gamma-x_{\gamma'})$
and $\Xi_2:=\mathsf{f}(x_\alpha-x_{\gamma'},x_{\gamma'}-x_\beta)$. In $(t_1\star t_2)\star t_3$ the
operator $\delta$ applied to $t_1\star t_2$ multiplies its kernel at $(\alpha,\gamma')$ by
$(x_\alpha-x_{\gamma'})_\kappa$; hence the kernel of $(t_1\star t_2)\star t_3$ is
$\sum_{\gamma,\gamma'}t_1(\alpha,\gamma)t_2(\gamma,\gamma')t_3(\gamma',\beta)$ times the degree-$\le2$
part of $e^{(iq/2)\Xi_1}e^{(iq/2)\Xi_2}=e^{(iq/2)(\Xi_1+\Xi_2)}$. By the identity used for the kernel
formula, applied
twice, $\Xi_1+\Xi_2=\Xi$, where
$\Xi:=\mathsf{f}(x_\alpha,x_\gamma)+\mathsf{f}(x_\gamma,x_{\gamma'})+\mathsf{f}(x_{\gamma'},x_\beta)
-\mathsf{f}(x_\alpha,x_\beta)$. In the same way the kernel of $t_1\star(t_2\star t_3)$ carries the
degree-$\le2$ part of $e^{(iq/2)(\Xi'_1+\Xi'_2)}$ with
$\Xi'_1=\mathsf{f}(x_\gamma-x_{\gamma'},x_{\gamma'}-x_\beta)$,
$\Xi'_2=\mathsf{f}(x_\alpha-x_\gamma,x_\gamma-x_\beta)$, and again $\Xi'_1+\Xi'_2=\Xi$. Thus both
$(t_1\star t_2)\star t_3$ and $t_1\star(t_2\star t_3)$ have the kernel
$\sum_{\gamma,\gamma'}t_1(\alpha,\gamma)t_2(\gamma,\gamma')t_3(\gamma',\beta)$ times the degree-$\le2$
part of $e^{(iq/2)\Xi}$. Writing
$x_\alpha-x_{\gamma'}=(x_\alpha-x_\gamma)+(x_\gamma-x_{\gamma'})$ and
$x_\gamma-x_\beta=(x_\gamma-x_{\gamma'})+(x_{\gamma'}-x_\beta)$, each coefficient of this expansion
is a finite sum of kernels of products $(\delta^{a}t_1)(\delta^{a'}t_2)(\delta^{a''}t_3)$ of bounded
operators, so the common kernel is that of one element of $\mathfrak{A}_{\mathsf{f}}$, and the two
sides agree.

\emph{Adjoint.} The charge $q$ is real, since $t_0\in\mathfrak{g}$ is hermitian in the convention
$\mathcal{U}_{\mathsf{f}}=\exp(i\mathsf{a}^{0}_{\mathsf{f}}t_0)$, so that $\operatorname{ad}_{t_0}$
is self-adjoint for $(X,Y)\mapsto\operatorname{tr}X^\dagger Y$ on
$\mathfrak{g}_{\mathbb{C}}$. On $\mathfrak{A}_{\mathsf{f}}$ the adjoint acts on the $\mathfrak{A}$
factor and conjugates scalar coefficients, fixing the real entries of $\mathsf{f}$. As $x_\kappa$ is
self-adjoint,
\[
(\delta_\kappa t)^\dagger=[x_\kappa,t]^\dagger=[t^\dagger,x_\kappa]=-\delta_\kappa(t^\dagger).
\]
Hence $(ts)^\dagger=s^\dagger t^\dagger$;
\begin{align*}
\Bigl(\tfrac{iq}{2}\mathsf{f}_{\kappa\lambda}(\delta_\kappa t)(\delta_\lambda s)\Bigr)^\dagger
&=-\tfrac{iq}{2}\mathsf{f}_{\kappa\lambda}(\delta_\lambda s^\dagger)(\delta_\kappa t^\dagger)
=\tfrac{iq}{2}\mathsf{f}_{\kappa\lambda}(\delta_\kappa s^\dagger)(\delta_\lambda t^\dagger),\\
\Bigl(-\tfrac{q^2}{8}\mathsf{f}_{\kappa\lambda}\mathsf{f}_{\kappa'\lambda'}
(\delta_\kappa\delta_{\kappa'}t)(\delta_\lambda\delta_{\lambda'}s)\Bigr)^\dagger
&=-\tfrac{q^2}{8}\mathsf{f}_{\kappa\lambda}\mathsf{f}_{\kappa'\lambda'}
(\delta_\lambda\delta_{\lambda'}s^\dagger)(\delta_\kappa\delta_{\kappa'}t^\dagger);
\end{align*}
in the first line the two signs from the derivations cancel, and the last step exchanges the names
$\kappa,\lambda$ and uses $\mathsf{f}_{\lambda\kappa}=-\mathsf{f}_{\kappa\lambda}$; in the second line
the four signs from the derivations cancel, and after exchanging $\kappa\leftrightarrow\lambda$,
$\kappa'\leftrightarrow\lambda'$ (two sign changes of $\mathsf{f}$) the right side is the third term
of $s^\dagger\star t^\dagger$.

\emph{Invertibility.} Let $P_0(t,s)=ts$, and let $P_1$, $P_2$ be the second and third terms
of~\eqref{10.6:eq-truncated-product-a}, homogeneous of degree $1$, $2$ in the explicit $\mathsf{f}$.
Then $(t\star s)^{(n)}=\sum_{n_1+n_2+n_3=n}P_{n_3}(t^{(n_1)},s^{(n_2)})$; in particular
$(t\star s)^{(0)}=t^{(0)}s^{(0)}$. If $t\star s=\mathbf{1}=s\star t$, then
$t^{(0)}s^{(0)}=\mathbf{1}=s^{(0)}t^{(0)}$ and
$t^{(0)}$ is invertible in $\mathfrak{A}$. Conversely let $v:=(t^{(0)})^{-1}\in\mathfrak{A}$. Applying
$\delta_\kappa$ to $vt^{(0)}=\mathbf{1}$ gives $(\delta_\kappa v)t^{(0)}+v\,\delta_\kappa t^{(0)}=0$,
that is
$\delta_\kappa v=-v(\delta_\kappa t^{(0)})v$; by induction on $|a|$ and the Leibniz rule every
$\delta^{a}v$ is a finite sum of words in $v$ and the $\delta^{a'}t^{(0)}$. Put $s^{(0)}:=v$ and, for
$n=1,2$,
\[
s^{(n)}:=-v\sum_{\substack{n_1+n_2+n_3=n\\ n_2<n}}P_{n_3}\bigl(t^{(n_1)},s^{(n_2)}\bigr),
\]
which makes $(t\star s)^{(n)}=0$ for $n=1,2$, so $t\star s=\mathbf{1}$. The derivatives $\delta^{a}$
entering
$P_{n_3}$ fall on $t^{(n_1)}$ and on $s^{(n_2)}$, $n_2<n$, and by the Leibniz rule and the formula for
$\delta_\kappa v$ each $s^{(n)}$ is a finite sum of words in $v$ and the $\delta^{a}t^{(m)}$. The same
construction on the other side gives $s'$ with $s'\star t=\mathbf{1}$, and by associativity
\[
s'=s'\star(t\star s)=(s'\star t)\star s=s.
\]

\emph{The trace and~\eqref{10.6:eq-truncated-product-b}.} By invariance under $\mathbb{L}_k$,
$t(\alpha+z,\alpha+z)=t(\alpha,\alpha)$ for $z\in\mathbb{L}_k$, so $\operatorname{Tr}_{\square}$ does not
depend on the choice of $\square$. The sums over $\gamma$ below converge absolutely by the
Cauchy--Schwarz inequality, the kernel of a bounded operator between the weighted $\ell^2$ spaces
being square summable in each variable for fixed value of the other. At $\alpha=\beta$ the exponent
in~\eqref{10.6:eq-truncated-product-1} is $\frac{iq}{2}\mathsf{f}(u,-u)=0$, $u=x_\alpha-x_\gamma$, by
antisymmetry; since the terms of~\eqref{10.6:eq-truncated-product-a} of degree one and two are the
Taylor terms of this exponential, their kernels vanish on the diagonal: the degree-one term carries
$\mathsf{f}(u,-u)$ and the degree-two term $\mathsf{f}(u,-u)^2$. Thus
$(t\star s)(\alpha,\alpha)=(ts)(\alpha,\alpha)$, which gives the first equality, and the same for
$s\star t$. Finally, write every cell $\gamma$ of the type indexing the range of $s$ (the domain of
$t$) uniquely as
$\gamma=\gamma'+z$, $\gamma'\in\square$, $z\in\mathbb{L}_k$. By translation invariance,
\[
\operatorname{Tr}_{\square}(ts)=\sum_{\gamma'\in\square}\sum_{z\in\mathbb{L}_k}\sum_{\alpha\in\square}
\operatorname{tr}\,t(\alpha-z,\gamma')\,s(\gamma',\alpha-z),
\]
and as $(\alpha,z)$ runs over the cells in $\square$ and $\mathbb{L}_k$, $\alpha-z$ runs once over
all cells of its type. Using the cyclicity of the trace of compositions of linear maps between the
finite-dimensional fibres, the right side is $\operatorname{Tr}_{\square}(st)
=\operatorname{Tr}_{\square}(s\star t)$.

\emph{Fourier form.} Let $t$ be invariant under $\mathbb{L}_0$. The positions of the cells of one
class form a translate of $\eta\mathbb{Z}^4$; hence for cells $\alpha,\beta$ of fixed classes,
$t(\alpha,\beta)$ depends only on $u=x_\alpha-x_\beta$, which ranges over a translate of
$\eta\mathbb{Z}^4$ depending only on the two classes, so that $t_{[\alpha][\beta]}$ and
$\hat t_{[\alpha][\beta]}$ are well defined. As $\eta^4\int_{\mathrm{BZ}_\eta}\bar{d}\xi\,e^{i\xi(u-u')}$
equals $1$ if $u=u'$ and $0$ otherwise, for $u,u'$ in the
same translate of $\eta\mathbb{Z}^4$, the inversion formula holds. The coefficients of
$i\partial_\kappa\hat t_{[\alpha][\beta]}$ are $u_\kappa t_{[\alpha][\beta]}(u)$, those of
$\delta_\kappa t$, so $\delta_\kappa\leftrightarrow i\partial_\kappa$; and since the sum over $\gamma$
in a composition runs over the classes of $\gamma$ and, within each class, over a translate of
$\eta\mathbb{Z}^4$, the composition of translation-invariant kernels is carried to the matrix product
$\hat t\hat s$. Substituting
in~\eqref{10.6:eq-truncated-product-a}, $i^2=-1$ in the first-order term and $i^4=1$ in the second
give~\eqref{10.6:eq-truncated-product-2}. Lastly
$t(\alpha,\alpha)=\eta^4\int\bar{d}\xi\,\hat t_{[\alpha][\alpha]}(\xi)$, and
$\square$, of volume one, contains $\eta^{-4}=N^4$ cells of each class, whence
$\operatorname{Tr}_{\square}t=\int_{\mathrm{BZ}_\eta}\bar{d}\xi\operatorname{tr}\hat t(\xi)$, with
$\operatorname{tr}$ over the fibre and the classes.
\end{proof}

\begin{proposition}[Second Taylor coefficients in the truncated ring]\label{10.6:prop-truncated-ring}
Let $\mathsf{f}$ be a constant curvature matrix and let $\varepsilon$ denote a real parameter.
\begin{enumerate}
\item[(a)] Let $\mathsf{T}_i(U)$ be finitely many gauge-covariant operators of finite range,
analytic in the configuration $U$. Let $E$ be an algebraic expression in them, built from sums,
products, adjoints and inverses, whose flat inverses exist in $\mathfrak{A}$; here the flat inverses
are the inverses occurring in $E$, taken on the flat configuration $\mathcal{U}_0 = \mathbf{1}$. Put
\[
\psi_y(U) := \sum_{\alpha \in \square(y)} \operatorname{tr} E(\mathsf{T}(U))(\alpha,\alpha),
\]
where $\square(y)$ is the unit cell at $y$. For each $i$ let $\mathsf{t}_i \in \mathfrak{A}_{\mathsf{f}}$ be
the Taylor polynomial of the reduced kernel of $\mathsf{T}_i(\mathcal{U}_{\mathsf{f}})$. Let
$E_{\star}(\mathsf{t})$ be the same expression evaluated in $\mathfrak{A}_{\mathsf{f}}$, with every product
replaced by $\star$. Then
\begin{equation}\label{10.6:eq-truncated-ring-a}
\mathsf{cf}[\psi](\mathsf{f}) = [\mathsf{f}^2]\operatorname{Tr}_{\square} E_{\star}(\mathsf{t}).
\end{equation}
\item[(b)] (The on-shell reading.) Let $B$ be a finitely supported bond field. The operators of
Ba{\l}aban's step at the background $e^{i\varepsilon B}$ are algebraic expressions in the finite-range
generating operators of that step listed earlier in this section together with the current $J$
(the middle term of \cite[(1.5)]{B-CMP109}) and the Lagrange multipliers of the averaging
constraints (the second fundamental forms of the level sets of the block-averaging maps).

The Taylor coefficients in $\varepsilon$ of $J$ and of the multipliers are multilinear forms in $B$
with exponentially decaying kernels and so extend to affine $B$.
Along affine fields the Taylor coefficients of orders one and two of $J$ and of the multipliers
vanish, and so do their values at the flat configuration. Consequently, to second order along affine
fields, the functionals of Ba{\l}aban's step have the same jets as the
expressions obtained from them by deleting every term that contains $J$ or a multiplier as a
factor.
\end{enumerate}
\end{proposition}

\begin{proof}
\emph{Part} (a). Let $B$ be a finitely supported bond field. Consider the analytic function
$\varepsilon \mapsto \psi_y(e^{i\varepsilon B})$.

Each $\mathsf{T}_i(e^{i\varepsilon B})$ is analytic in $\varepsilon$ and has finite range. Its Taylor
coefficients are therefore finite-range operators. At $\varepsilon = 0$ every inverse occurring in $E$
is a flat inverse, which by hypothesis exists in $\mathfrak{A}$. The inverse of a perturbation of an
invertible operator expands by the following identity: if $s_\varepsilon$ is a differentiable family
of invertible operators, then $\tfrac{d}{d\varepsilon}s_\varepsilon^{-1}
= -s_\varepsilon^{-1}\bigl(\tfrac{d}{d\varepsilon}s_\varepsilon\bigr)s_\varepsilon^{-1}$. Applying this
identity repeatedly to each inverse in $E$, and expanding sums, products and adjoints term by term, we
find the following. The Taylor coefficients of $E(\mathsf{T}(e^{i\varepsilon B}))$ are noncommutative
polynomials in the flat inverses and in the Taylor coefficients of the
$\mathsf{T}_i(e^{i\varepsilon B})$. The Taylor coefficients of $\psi_y(e^{i\varepsilon B})$ are the
sums over $\alpha \in \square(y)$ of the traces of their diagonal values at $(\alpha,\alpha)$.

The Taylor coefficient of order $n$ depends on $B$ only through the Taylor coefficients of the
$\mathsf{T}_i(e^{i\varepsilon B})$. It is therefore an $n$-linear form in $B$. Its kernel is a finite
sum of products of finite-range kernels and of kernels of flat inverses. The kernels of the flat
inverses decay exponentially by the Combes--Thomas argument \cite{CT73}. Hence each Taylor
coefficient of $\psi_y(e^{i\varepsilon B})$ is a multilinear form in $B$ with exponentially decaying
kernel. Because of this decay, the sums defining the form converge for bond fields of polynomial
growth. The Taylor coefficients therefore extend to affine fields $B$.

Take $B = \mathsf{a}^{y}_{\mathsf{f}} t_0$. By the definition of $K_y$ and of $\mathsf{cf}$ in
Lemma~\ref{10.6:lem-second-moment}, the extended coefficient of $\varepsilon^2$ is
\[
\tfrac12 \sum_{b,b'} K_y(b,b')\, \mathsf{a}^{y}_{\mathsf{f}}(b)\, \mathsf{a}^{y}_{\mathsf{f}}(b')
= \mathsf{cf}[\psi](\mathsf{f}).
\]
On this field, $\mathsf{T}_i(e^{i\varepsilon \mathsf{a}^{y}_{\mathsf{f}} t_0})$ is the operator
$\mathsf{T}_i$ on the configuration $\mathcal{U}_{\varepsilon \mathsf{f}}$. The Taylor polynomial in
$\varepsilon$ of its reduced kernel, taken modulo $\varepsilon^3$, is $\mathsf{t}_i$ with $\mathsf{f}$
replaced by $\varepsilon \mathsf{f}$. Since each Taylor coefficient of
$\mathsf{T}_i(e^{i\varepsilon B})$ at a pair of cells depends on $B$ only through its values on the
bonds within the range of $\mathsf{T}_i$ from those cells, the extended coefficients at
$B = \mathsf{a}^{y}_{\mathsf{f}} t_0$ are the same noncommutative polynomials in the flat inverses and
in the Taylor coefficients in $\varepsilon$ of the reduced kernels of the
$\mathsf{T}_i(\mathcal{U}_{\varepsilon\mathsf{f}})$, composed by the kernel formula.

On such operators the product rule is the kernel formula. By
Definition~\ref{10.6:def-truncated-product}, the truncation of the kernel formula modulo
$(\mathsf{f})^3$ is the product $\star$ of \eqref{10.6:eq-truncated-product-a}. Consequently:
\begin{itemize}
\item the reduced kernel of a product is, modulo $(\mathsf{f})^3$, the $\star$-product of the reduced
kernels;
\item the reduced kernel of an inverse is the $\star$-inverse;
\item the reduced kernel of an adjoint is the adjoint, with $(t\star s)^{\dagger} = s^{\dagger}\star
t^{\dagger}$.
\end{itemize}
Each $\star$-inverse exists by the same definition, since the degree-zero part of the sub-expression
to be inverted is the corresponding flat operator, which is invertible in $\mathfrak{A}$ by
hypothesis. Hence the jet to second order in $\varepsilon$ of the reduced kernel of
$E(\mathsf{T}(\mathcal{U}_{\varepsilon\mathsf{f}}))$, that is, the polynomials above composed by the
kernel formula and truncated modulo $\varepsilon^3$, is $E_{\star}(\mathsf{t})$ with $\mathsf{f}$
replaced by $\varepsilon\mathsf{f}$.

The phase in the kernel formula vanishes at coinciding arguments
(Definition~\ref{10.6:def-truncated-product}). So the sum over $\alpha \in \square(y)$ of
$\operatorname{tr} E(\alpha,\alpha)$ is $\operatorname{Tr}_{\square}$ of this reduced kernel. Reading
off the coefficient of $\varepsilon^2$ gives \eqref{10.6:eq-truncated-ring-a}.

\emph{Part} (b). Let $B$ be finitely supported. The operators of Ba{\l}aban's step at
$e^{i\varepsilon B}$ are algebraic expressions in the finite-range generating operators listed earlier
in this section, in the current $J$, which enters through the middle term of \cite[(1.5)]{B-CMP109},
and in the Lagrange multipliers of the constraints, which enter as the second fundamental forms of the
level sets of the block-averaging maps. The Taylor coefficients of these quantities are multilinear
forms in $B$ with exponentially decaying kernels and therefore extend to affine $B$ as in part (a).
We show that along affine fields the Taylor coefficients of orders $0$, $1$ and $2$ of $J$ and of the
multipliers vanish.

By the reproduction property and the criticality property of constant abelian fields,
\eqref{10.6:eq-constant-fields-a} of Lemma~\ref{10.6:lem-constant-fields}, the
configuration $\mathcal{U}_{\varepsilon\mathsf{f}}$ solves the constrained Euler--Lagrange equations
exactly on every level, with zero current and zero multiplier. Indeed:
\begin{itemize}
\item by the reproduction property, for $|\varepsilon|$ small, which suffices for the Taylor
coefficients at $\varepsilon = 0$, block averaging maps $\mathcal{U}_{\varepsilon\mathsf{f}}$ on
$\mathbb{L}_j$ to $\mathcal{U}_{\varepsilon\mathsf{f}}$ on $\mathbb{L}_{j+1}$;
\item by the criticality property, $\mathcal{U}_{\varepsilon\mathsf{f}}$ is a critical point of the
action for every orientation weight.
\end{itemize}

Differentiating these constrained Euler--Lagrange equations in $\varepsilon$ at $\varepsilon = 0$, we
see that the Taylor coefficients of $\mathcal{U}_{\varepsilon\mathsf{f}}$ solve, order by order, the
linearised constrained equations at the flat configuration with affine data. Now let $U_j(\cdot)$ be
the minimiser of the constrained variational problem at level $j$ and $V^{(j)}(\cdot)$ the critical
point, both as functions of the background, and consider the Taylor coefficients of
$U_j(e^{i\varepsilon B})$ and of $V^{(j)}(e^{i\varepsilon B})$, extended to affine $B$ by the decay of
their kernels. At $B = \mathsf{a}^{y}_{\mathsf{f}} t_0$ they solve the same linearised equations. The
constrained flat Hessians have exponentially decaying inverses. Therefore the polynomially bounded
solution of these linearised equations is unique modulo gauge.

Hence, along affine fields, the Taylor coefficients of the minimiser and of the critical point agree,
modulo gauge transformations, with those of the exact solution $\mathcal{U}_{\varepsilon\mathsf{f}}$.
The current and the multipliers are gauge covariant and vanish at that solution, so they vanish also
at every gauge transform of it. So the Taylor coefficients of orders $0$, $1$ and $2$ of $J$ and of
the multipliers vanish along affine fields; order $0$ is the case $\varepsilon = 0$, in which the
configuration is the flat one.

Consider a term of the expressions that contains $J$ or a multiplier as a factor. Its Taylor
coefficient of order at most two is a sum of products of Taylor coefficients of its factors whose
orders add up to at most two. Every such product contains a vanishing coefficient of that factor.
Deleting the term therefore does not change the jet to second order along affine fields.
\end{proof}

The expressions obtained by this deletion are the expressions whose second Taylor coefficients are
computed in the remainder of this section. All identities stated for them there are identities in
$\mathfrak{A}_{\mathsf{f}}$, and products there are written without $\star$.

\begin{lemma}[The generators on a constant background]\label{10.6:lem-background-generators}
Let $\mathsf{f}$ be small enough for Lemma~\ref{10.6:lem-constant-fields} to apply. For
$0\le j\le k$ let $\mathcal{U}_{\mathsf{f}}$ denote the constant abelian configuration of curvature
$\mathsf{f}$ on the lattice $\mathbb{L}_j$ of spacing $\ell_j=L^{j-k}$, so that $\ell_0=\eta$.
\begin{enumerate}
\item[(i)] \emph{The bare lattice.} Place the fibre of each bond at its midpoint. Then the two
phases in the kernel of the lattice gradient $D_j\in\mathfrak{A}(\mathcal{C}_j,\mathcal{X}_j)$ on
$\mathcal{U}_{\mathsf{f}}$ are phases of straight segments. Its reduced kernel therefore does not
depend on $\mathsf{f}$: it is
\[
(D_j\lambda)(b)=\ell_j^{-1}\bigl(\lambda(b_+)-\lambda(b_-)\bigr),
\]
where $b_\pm$ are the endpoints of the bond $b$. Its symbol at momentum $k$ (in $d_\mu(k)$ and in
$\mathsf{t}_1$ below the letter $k$ denotes a momentum, not the top level) is
$d_\mu(k)=(2i/\ell_j)\sin(\ell_jk_\mu/2)$.

For $\mu<\nu$ let $\mathsf{H}_{\mu\nu}$ be the Hessian at $\mathcal{U}_{\mathsf{f}}$ of the sum
$\sum_{p\parallel(\mu\nu)}a_p$ on $\mathbb{L}_0$ of the plaquette terms $a_p$ of the action over
the plaquettes parallel to the $(\mu\nu)$-plane. This Hessian is computed plaquette by plaquette
as in \cite[(3.10)]{B-CMP99b}. It is symmetric and of range one plaquette. Its reduced symbol has
the scaling form $\eta^{-2}\mathsf{t}_1(\eta k;\eta^2\mathsf{f})$, where $\mathsf{t}_1$ is a
trigonometric polynomial in its first argument whose coefficients are entire functions of the
second. Put
\[
\mathsf{H}_0:=\sum_{\mu<\nu}\mathsf{H}_{\mu\nu},\qquad
\mathsf{O}:=\sum_{\mu<\nu}(h_\mu+h_\nu)\,\mathsf{H}_{\mu\nu}
=\nabla^2\mathsf{A}^{\mathbf{3}}(\mathcal{U}_{\mathsf{f}}),
\]
where $h_\mu$ are the orientation weights of the orientation-weighted action
$\mathsf{A}^{\mathbf{3}}$. Then
\begin{equation}\label{10.6:eq-background-generators-a}
\mathsf{H}_{\mu\nu}D_0=0\quad(\text{each }\mu<\nu);\qquad\text{hence}\quad
\mathsf{H}_0D_0=0,\quad \mathsf{O}D_0=0,\quad D_0^{\dagger}\mathsf{O}=0 .
\end{equation}
The flat parts are $(\mathsf{H}^{(0)}_{\mu\nu}A,A)=\|(dA)_{\mu\nu}\|^2$, where $(dA)_{\mu\nu}$ is
the $(\mu\nu)$-component of the lattice exterior derivative of $A$.

\item[(ii)] \emph{One step $\mathbb{L}_j\to\mathbb{L}_{j+1}$}, $0\le j<k$. The following operators
are all taken on $\mathcal{U}_{\mathsf{f}}$.
\begin{itemize}
\item $\tilde{\mathsf{Q}}_{j+1}\in\mathfrak{A}_{\mathsf{f}}(\mathcal{X}_j,\mathcal{X}_{j+1})$ is
the linearisation of the averaging map $M$ of \cite[(0.10)--(0.12)]{B-CMP109}, that is, the
linear operator of \cite[(2.4)]{B-CMP109}, taken at $\mathcal{U}_{\mathsf{f}}$.
\item $\mathcal{C}^{\circ}_j:=\ell^2(\mathbb{L}_j\setminus\mathbb{L}_{j+1})$.
\item $\mathsf{s}_{j+1}\in\mathfrak{A}_{\mathsf{f}}(\mathcal{X}_j,\mathcal{C}^{\circ}_j)$ is the
linearisation $B'\mapsto(B'(y(x),x))_x$ of the averaged contour variables of
\cite[(2.5)]{B-CMP109}.
\item $\mathrm{ev}\colon\mathcal{C}_j\to\mathcal{C}_{j+1}$ is the restriction.
\item $\iota\colon\mathcal{C}^{\circ}_j\to\mathcal{C}_j$ is the extension by $0$.
\item $\mathsf{e}\colon\mathcal{C}_{j+1}\to\mathcal{C}_j$ is defined by letting
$(\mathsf{e}\lambda)(x)$ be the transport of $\lambda(y(x))$ to $x$.
\item $\mathsf{b}\colon\mathcal{C}_j\to\mathcal{C}_{j+1}$ is the covariant block mean
\cite[(3.19)]{B-CMP99b} for one step.
\end{itemize}
The reduced kernels of $\mathrm{ev},\iota,\mathsf{e},\mathsf{b}$ do not depend on $\mathsf{f}$, and
\[
\mathrm{ev}\,\iota=0,\qquad \mathrm{ev}\,\mathsf{e}=1,\qquad \mathsf{b}\,\mathsf{e}=1 .
\]
Moreover
\begin{equation}\label{10.6:eq-background-generators-b}
\tilde{\mathsf{Q}}_{j+1}D_j=D_{j+1}\,\mathrm{ev},\qquad
\mathsf{s}_{j+1}D_j=\iota^{\dagger}(1-\mathsf{e}\,\mathrm{ev});\qquad\text{hence}\quad
\mathsf{s}_{j+1}D_j\iota=1,\quad \mathsf{s}_{j+1}D_j\mathsf{e}=0 .
\end{equation}
With $\sigma_{j+1}:=\mathsf{b}\,\iota\,\mathsf{s}_{j+1}$ and
$\mathsf{Q}_{j+1}:=\tilde{\mathsf{Q}}_{j+1}+D_{j+1}\sigma_{j+1}$ one has
\begin{equation}\label{10.6:eq-background-generators-c}
\mathsf{Q}_{j+1}D_j=D_{j+1}\mathsf{b}.
\end{equation}

\item[(iii)] \emph{Flat parts and composites.} The flat part $\mathsf{Q}^{(0)}_{j+1}$ is the plain
average
\[
(\mathsf{Q}^{(0)}_{j+1}A)_\mu(y)=L^{-5}\sum_{x\in B(y)}\sum_{i<L}A_\mu(x+i\ell_je_\mu).
\]
It is a contraction $\mathcal{X}_j\to\mathcal{X}_{j+1}$ with non-negative kernel, and
$\mathsf{b}^{(0)}$ is the plain block mean. Put
\[
Q:=\mathsf{Q}_k\cdots\mathsf{Q}_1,\qquad
Q':=\mathsf{b}\cdots\mathsf{b}\ (k\text{ factors}),\qquad
\tilde{Q}:=\tilde{\mathsf{Q}}_k\cdots\tilde{\mathsf{Q}}_1 .
\]
Then $Q^{(0)}$ and $Q'^{(0)}$ are the operators $Q_k$ and $Q'_k$ of
\cite[(1.42)--(1.45)]{B-CMP95}. With $\mathrm{ev}_{(k)}$ the restriction
$\mathcal{C}_0\to\mathcal{C}_k$,
\begin{equation}\label{10.6:eq-background-generators-d}
QD_0=D_kQ',\qquad \tilde{Q}D_0=D_k\,\mathrm{ev}_{(k)},\qquad
\tilde{Q}=Q-D_k\sigma_{(k)}\quad\text{for some }
\sigma_{(k)}\in\mathfrak{A}_{\mathsf{f}}(\mathcal{X}_0,\mathcal{C}_k).
\end{equation}
\end{enumerate}
\end{lemma}

\begin{proof}
\emph{Part (i).}

\emph{The gradient.} Take the fibre of a bond at its midpoint. The kernel of $D_j$ on
$\mathcal{U}_{\mathsf{f}}$ carries one phase from each endpoint of the bond. Each of these is the
phase of the straight segment joining the endpoint to the midpoint. Hence the reduced kernel is the
plain difference quotient displayed in (i). Its symbol is obtained by applying it to
$e^{ik\cdot x}$ and referring the result to the midpoint of the bond. This gives
$\ell_j^{-1}(e^{i\ell_jk_\mu/2}-e^{-i\ell_jk_\mu/2})=(2i/\ell_j)\sin(\ell_jk_\mu/2)$.

\emph{Properties of $\mathsf{H}_{\mu\nu}$.} These are read off from the Hessian of a single
plaquette term in \cite[(3.10)]{B-CMP99b}, summed over the plaquettes parallel to $(\mu\nu)$.
\begin{itemize}
\item $\mathsf{H}_{\mu\nu}$ is symmetric, being a Hessian.
\item It has range one plaquette, since each $a_p$ depends only on the bonds of $p$.
\item On $\mathcal{U}_{\mathsf{f}}$ the holonomy of every plaquette $p\parallel(\mu\nu)$ of
$\mathbb{L}_0$ is $\exp(i\eta^2\mathsf{f}_{\mu\nu}t_0)$, by Lemma~\ref{10.6:lem-constant-fields}
(see \eqref{10.6:eq-constant-fields-a}) with $\ell_0=\eta$. The bond variables enter through
difference quotients of spacing $\eta$. This gives the scaling form
$\eta^{-2}\mathsf{t}_1(\eta k;\eta^2\mathsf{f})$, with the stated dependence on the two
arguments.
\end{itemize}

\emph{Proof of \eqref{10.6:eq-background-generators-a}.} Fix $\mu<\nu$ and write
$F:=\sum_{p\parallel(\mu\nu)}a_p$. The sum $F$ is gauge invariant: $F(U^u)=F(U)$ for every gauge
transformation $u$. By the criticality of $\mathcal{U}_{\mathsf{f}}$ established in
Lemma~\ref{10.6:lem-constant-fields}, $\mathcal{U}_{\mathsf{f}}$ is a critical point of $F$.

We use the principle that at a critical point the Hessian of an invariant function annihilates
the orbit directions and does not depend on the chart. In detail:
\begin{itemize}
\item Differentiating $F(U^u)=F(U)$ once in $u$ at the identity, in the direction of a gauge
function $\lambda$, shows that the first derivative of $F$ at every configuration $U$ vanishes
on the tangent to the gauge orbit through $U$.
\item Differentiating this identity once more at $\mathcal{U}_{\mathsf{f}}$, in a direction $A$,
gives $\nabla^2F(\mathcal{U}_{\mathsf{f}})[A,D_0\lambda]$ plus the first derivative of $F$ at
$\mathcal{U}_{\mathsf{f}}$ applied to the derivative of the orbit tangent. Here $D_0\lambda$ is
the orbit tangent at $\mathcal{U}_{\mathsf{f}}$.
\item The second term vanishes because $\mathcal{U}_{\mathsf{f}}$ is critical. For the same
reason the second-order term is independent of the chart used.
\end{itemize}
Hence $(A,\mathsf{H}_{\mu\nu}D_0\lambda)=0$ for all $A$ and $\lambda$, that is,
$\mathsf{H}_{\mu\nu}D_0=0$.

Summing over $\mu<\nu$, with weights $1$ and $h_\mu+h_\nu$ respectively, gives
$\mathsf{H}_0D_0=0$ and $\mathsf{O}D_0=0$. Since $\mathsf{O}$ is a real linear combination of the
symmetric operators $\mathsf{H}_{\mu\nu}$, it is symmetric. Taking adjoints in $\mathsf{O}D_0=0$
gives $D_0^{\dagger}\mathsf{O}=0$.

\emph{Flat parts.} At $\mathsf{f}=0$ every plaquette holonomy is $\mathbf{1}$. The quadratic form
of the plaquette-by-plaquette Hessian of \cite[(3.10)]{B-CMP99b} then reduces, for the plaquettes
parallel to $(\mu\nu)$, to $\|(dA)_{\mu\nu}\|^2$.

\emph{Part (ii).}

\emph{The $\mathsf{f}$-independence and the three relations.} The reduced kernels of
$\mathrm{ev},\iota,\mathsf{e},\mathsf{b}$ do not depend on $\mathsf{f}$. This follows from the
computation of the phases of contours on $\mathcal{U}_{\mathsf{f}}$ made in the proof of the
reproduction property in Lemma~\ref{10.6:lem-constant-fields}. The three relations hold for the
following reasons.
\begin{itemize}
\item $\mathrm{ev}\,\iota=0$, because $\iota\lambda$ vanishes on $\mathbb{L}_{j+1}$.
\item $\mathrm{ev}\,\mathsf{e}=1$, because for $x\in\mathbb{L}_{j+1}$ the point $y(x)$ is $x$
itself.
\item $\mathsf{b}\mathsf{e}=1$, because on $B(y)$ the function $\mathsf{e}\lambda$ consists of
the transports of the single value $\lambda(y)$, and their covariant block mean at $y$ is
$\lambda(y)$.
\end{itemize}

\emph{Proof of \eqref{10.6:eq-background-generators-b}.} The map $M$ and the parallel transports
along contours are gauge covariant:
\[
M(U^u)=M(U)^{u|_{\mathbb{L}_{j+1}}},\qquad U^u(y,x)=u(y)\,U(y,x)\,u(x)^{-1}.
\]
Take $u=e^{i\epsilon\lambda}$ with $\lambda\in\mathcal{C}_j$, and differentiate at $\epsilon=0$
at $U=\mathcal{U}_{\mathsf{f}}$. The tangent to $\epsilon\mapsto\mathcal{U}_{\mathsf{f}}^{u}$ is
$D_j\lambda$.

For the averaging map:
\begin{itemize}
\item The derivative of the left side of the first identity is
$\tilde{\mathsf{Q}}_{j+1}D_j\lambda$.
\item By the reproduction property of Lemma~\ref{10.6:lem-constant-fields}, $M$ maps
$\mathcal{U}_{\mathsf{f}}$ on $\mathbb{L}_j$ to $\mathcal{U}_{\mathsf{f}}$ on $\mathbb{L}_{j+1}$.
So the derivative of the right side is the orbit tangent on $\mathbb{L}_{j+1}$ in the direction
$\lambda|_{\mathbb{L}_{j+1}}$, namely $D_{j+1}\mathrm{ev}\,\lambda$.
\end{itemize}
This proves the first identity in \eqref{10.6:eq-background-generators-b}.

For the contour variables, take $x\in\mathbb{L}_j\setminus\mathbb{L}_{j+1}$. The second
covariance relation with $y=y(x)$ shows that the linearised contour variable at $x$ changes by
$\lambda(x)$ minus the transport of $\lambda(y(x))$ to $x$, that is, by
$((1-\mathsf{e}\,\mathrm{ev})\lambda)(x)$. Restricting to
$\mathbb{L}_j\setminus\mathbb{L}_{j+1}$ is the application of $\iota^{\dagger}$. This proves
$\mathsf{s}_{j+1}D_j=\iota^{\dagger}(1-\mathsf{e}\,\mathrm{ev})$.

For the two consequences, note that $\iota^{\dagger}\iota=1$, since $\iota$ is the extension by
zero. Hence
\[
\mathsf{s}_{j+1}D_j\iota=\iota^{\dagger}\iota-\iota^{\dagger}\mathsf{e}\,\mathrm{ev}\,\iota=1,
\qquad
\mathsf{s}_{j+1}D_j\mathsf{e}=\iota^{\dagger}(\mathsf{e}-\mathsf{e}\,\mathrm{ev}\,\mathsf{e})=0,
\]
by $\mathrm{ev}\,\iota=0$ and $\mathrm{ev}\,\mathsf{e}=1$.

\emph{Proof of \eqref{10.6:eq-background-generators-c}.} By the definition of
$\mathsf{Q}_{j+1}$ and the first identity in \eqref{10.6:eq-background-generators-b},
\[
\mathsf{Q}_{j+1}D_j=D_{j+1}\bigl(\mathrm{ev}+\sigma_{j+1}D_j\bigr).
\]
By the second identity in \eqref{10.6:eq-background-generators-b},
$\sigma_{j+1}D_j=\mathsf{b}\,\iota\iota^{\dagger}(1-\mathsf{e}\,\mathrm{ev})$. Here
$\iota\iota^{\dagger}$ is multiplication by the indicator of
$\mathbb{L}_j\setminus\mathbb{L}_{j+1}$. The function $(1-\mathsf{e}\,\mathrm{ev})\lambda$ vanishes
on $\mathbb{L}_{j+1}$, because $\mathrm{ev}(1-\mathsf{e}\,\mathrm{ev})=\mathrm{ev}-\mathrm{ev}=0$.
So $\iota\iota^{\dagger}$ acts on it as the identity. Therefore
\[
\mathrm{ev}+\sigma_{j+1}D_j=\mathrm{ev}+\mathsf{b}(1-\mathsf{e}\,\mathrm{ev})
=\mathrm{ev}+\mathsf{b}-\mathrm{ev}=\mathsf{b},
\]
using $\mathsf{b}\mathsf{e}=1$, and \eqref{10.6:eq-background-generators-c} follows.

\emph{Part (iii).}

\emph{The flat part of $\mathsf{Q}_{j+1}$.} At $\mathsf{f}=0$, the operator
$\tilde{\mathsf{Q}}^{(0)}_{j+1}$ is the plain average displayed in (iii) plus contour terms. Let
$c$ be a bond of $\mathbb{L}_{j+1}$ in direction $e_\mu$ with endpoints $c_\pm$. At $c$ the
contour terms equal the mean over $x\in B(c_-)$ of $\{A(c_-,x)-A(c_+,x')\}/(L\ell_j)$, where
$x'=x+\ell_{j+1}e_\mu$ is the corresponding point of $B(c_+)$. This mean equals
$-(D_{j+1}\sigma^{(0)}_{j+1}A)(c)$. Since
$\mathsf{Q}^{(0)}_{j+1}=\tilde{\mathsf{Q}}^{(0)}_{j+1}+D_{j+1}\sigma^{(0)}_{j+1}$, the contour
terms cancel, and $\mathsf{Q}^{(0)}_{j+1}$ is the plain average. Its kernel is non-negative, with
all weights equal to $L^{-5}$.

\emph{The contraction property.} Write $\|A\|^2=\ell_j^4\sum_b|A(b)|^2$ on $\mathcal{X}_j$.
Jensen's inequality for the average of $L^5$ terms gives
$|(\mathsf{Q}^{(0)}_{j+1}A)_\mu(y)|^2\le L^{-5}\sum_{x\in B(y)}\sum_{i<L}|A_\mu(x+i\ell_je_\mu)|^2$.
Sum over $y$ with weight $\ell_{j+1}^4=L^4\ell_j^4$. As $y$ runs over $\mathbb{L}_{j+1}$, $x$ over
$B(y)$ and $i$ over $\{0,\dots,L-1\}$, each bond of $\mathbb{L}_j$ in direction $\mu$ occurs
exactly $L$ times. Hence
\[
\|(\mathsf{Q}^{(0)}_{j+1}A)_\mu\|^2\le L^4\ell_j^4\,L^{-5}\cdot L\sum_b|A_\mu(b)|^2
=\|A_\mu\|^2 .
\]
At $\mathsf{f}=0$ the transports are trivial, so $\mathsf{b}^{(0)}\lambda(y)=L^{-4}\sum_{x\in
B(y)}\lambda(x)$, the plain block mean. The $k$-fold composites of these flat operators are the
operators $Q_k$ and $Q'_k$ of \cite[(1.42)--(1.45)]{B-CMP95}.

\emph{Proof of \eqref{10.6:eq-background-generators-d}.} For $0\le j\le k$ put
\[
Q_j:=\mathsf{Q}_j\cdots\mathsf{Q}_1,\qquad
\tilde{Q}_j:=\tilde{\mathsf{Q}}_j\cdots\tilde{\mathsf{Q}}_1,\qquad
Q'_j:=\mathsf{b}\cdots\mathsf{b}\ (j\text{ factors}),
\]
with all three equal to $1$ for $j=0$. Let $\mathrm{ev}_{(j)}$ be the restriction
$\mathcal{C}_0\to\mathcal{C}_j$. All three identities are proved by induction on $j$, and the case
$j=k$ is \eqref{10.6:eq-background-generators-d}.

\emph{First identity.} If $Q_jD_0=D_jQ'_j$, then by \eqref{10.6:eq-background-generators-c}
\[
Q_{j+1}D_0=\mathsf{Q}_{j+1}D_jQ'_j=D_{j+1}\mathsf{b}\,Q'_j=D_{j+1}Q'_{j+1}.
\]

\emph{Second identity.} If $\tilde{Q}_jD_0=D_j\mathrm{ev}_{(j)}$, then by the first identity in
\eqref{10.6:eq-background-generators-b}
\[
\tilde{Q}_{j+1}D_0=\tilde{\mathsf{Q}}_{j+1}D_j\mathrm{ev}_{(j)}=D_{j+1}\mathrm{ev}\,
\mathrm{ev}_{(j)}=D_{j+1}\mathrm{ev}_{(j+1)}.
\]

\emph{Third identity.} Start with $\sigma_{(0)}:=0$. Suppose $\tilde{Q}_j=Q_j-D_j\sigma_{(j)}$.
Then $\tilde{Q}_{j+1}=\tilde{\mathsf{Q}}_{j+1}(Q_j-D_j\sigma_{(j)})$. We expand the two terms using
$\tilde{\mathsf{Q}}_{j+1}=\mathsf{Q}_{j+1}-D_{j+1}\sigma_{j+1}$.
\begin{itemize}
\item The first term is
$\tilde{\mathsf{Q}}_{j+1}Q_j=Q_{j+1}-D_{j+1}\sigma_{j+1}Q_j$.
\item For the second term, \eqref{10.6:eq-background-generators-c} gives
$\tilde{\mathsf{Q}}_{j+1}D_j\sigma_{(j)}
=D_{j+1}\mathsf{b}\,\sigma_{(j)}-D_{j+1}\sigma_{j+1}D_j\sigma_{(j)}$.
\end{itemize}
Subtracting, and using the inductive hypothesis for $\tilde{Q}_j$,
\[
\tilde{\mathsf{Q}}_{j+1}(Q_j-D_j\sigma_{(j)})
=Q_{j+1}-D_{j+1}\bigl(\sigma_{j+1}\tilde{Q}_j+\mathsf{b}\,\sigma_{(j)}\bigr).
\]
So $\tilde{Q}_{j+1}=Q_{j+1}-D_{j+1}\sigma_{(j+1)}$ with
$\sigma_{(j+1)}:=\sigma_{j+1}\tilde{Q}_j+\mathsf{b}\,\sigma_{(j)}$. This operator is a sum of
composites of $\mathsf{b}$, $\iota$, $\mathsf{s}_{j+1}$, $\tilde{\mathsf{Q}}_1,\dots,
\tilde{\mathsf{Q}}_j$ and $\sigma_{(j)}$, all tame translation-invariant operators on
$\mathcal{U}_{\mathsf{f}}$, and composition of operators corresponds to the product $\star$ of
their reduced kernels (Definition~\ref{10.6:def-truncated-product}). Hence, by induction,
$\sigma_{(k)}\in\mathfrak{A}_{\mathsf{f}}(\mathcal{X}_0,\mathcal{C}_k)$.
\end{proof}

\begin{lemma}[Analyticity in the flux and a uniform majorant]\label{10.6:lem-uniform-majorant}
Let $Q=\mathsf{Q}_k\star\cdots\star\mathsf{Q}_1$ be the $k$-fold composite of the corrected averages
$\mathsf{Q}_{j+1}$, $0\le j\le k-1$, and let $Q'$ be the second $k$-fold composite of one-step maps
introduced with $Q$ earlier in this section, in which $\mathsf{f}$ enters only through the explicit
phases. Write their kernels as
$Q(c,b;\mathsf{f})$ and $Q'(c,b;\mathsf{f})$, where $c$ is a bond of $\mathbb{L}_k$ and $b$ a bond of
$\mathbb{L}_0$. For complex $\mathsf{f}\in\mathbb{C}^6$ with $|\mathsf{f}|\le c_L$,
\begin{equation}\label{10.6:eq-uniform-majorant-a}
|Q(c,b;\mathsf{f})|\le \bar{Q}(c,b),
\end{equation}
where $\bar{Q}$ is a non-negative kernel supported on $\{|x_b-x_c|\le 2\}$ with
$\|\bar{Q}\|_{\mathcal{X}_0\to\mathcal{X}_k}\le C_Q$ uniformly in $k$, where the constant $C_Q$
depends only on $L$ and on the gauge group. The same holds for $Q'$. For
every $n\ge 0$ and every multi-index $\alpha$ of position derivations,
\begin{equation}\label{10.6:eq-uniform-majorant-1}
\|\delta^{\alpha}Q^{(n)}\|,\ \|\delta^{\alpha}Q'^{(n)}\|
\le 2^{|\alpha|}\,C_Q\,c_L^{-n}\,|\mathsf{f}|^{n}.
\end{equation}
\end{lemma}

\begin{proof}
\emph{Analyticity of one step.} The block average of \cite[(0.12)]{B-CMP109}, the group mean
\cite[(0.10)]{B-CMP109} and the contour variables are analytic functions of configurations with
values in the complexification of the gauge group, in a neighbourhood of the regular configurations
\cite[(0.5)--(0.12)]{B-CMP109}, \cite[Proposition~3]{B-CMP98}.

Let $c$ be a bond of $\mathbb{L}_{j+1}$ and $x$ a site. Consider the matrix elements of
$\tilde{\mathsf{Q}}_{j+1}$ and of $\mathsf{s}_{j+1}$ between $c$ (respectively $x$) and the bonds of
$B(c_-)\cup B(c_+)$, written in the radial gauge at $x_c$. In this gauge the bond variables of
$\mathcal{U}_{\mathsf{f}}$ on $B(c_-)\cup B(c_+)$ are entire functions of the dimensionless flux
$\ell_j^2\mathsf{f}\in\mathbb{C}^6$, and they lie in the neighbourhood of the regular configurations just
described when $|\ell_j^2\mathsf{f}|<c_L$. These matrix elements are therefore analytic
functions of $\ell_j^2\mathsf{f}$ in the region $|\ell_j^2\mathsf{f}|<c_L$.

The reduced kernels are obtained from these matrix elements through the explicit phases. These phases
are entire functions of $\mathsf{f}$, and they are evaluated at points lying within distance
$2\ell_{j+1}$ of $x_c$. Hence the reduced kernels are analytic in the same region.

\emph{One-step majorant.} We combine this analyticity with the estimate \cite[(139)]{B-CMP98}. In the
notation of that paper, the estimate reads
\[
|Q''(V_0)A|\le C_1'L^2\alpha_0\,Q''|A| .
\]
The combination gives a constant $C_L$, depending on $L$, such that the corrected average satisfies
\begin{equation}\label{10.6:eq-uniform-majorant-2}
|\mathsf{Q}_{j+1}(c,b;\mathsf{f})|
\le e^{C_L\ell_j^2|\mathsf{f}|}\,\mathsf{Q}^{(0)}_{j+1}(c,b)
+C_L\ell_j^2|\mathsf{f}|\,\bar{E}_{j+1}(c,b).
\end{equation}
Here $\mathsf{Q}^{(0)}_{j+1}$ is the degree-zero part of $\mathsf{Q}_{j+1}$. The kernel
$\bar{E}_{j+1}\ge 0$ is supported on the same bonds as $\mathsf{Q}_{j+1}$, and
$\|\bar{E}_{j+1}\|\le C_L$.

\emph{Phases of the composite.} Write $Q_j$ for the composite of the first $j$ steps, so that $Q=Q_k$.
In the $\star$-product $\mathsf{Q}_{j+1}\star Q_j$ the phase has modulus at most
\[
\exp\bigl(2|q|\,|\mathsf{f}|\,\ell_{j+1}\ell_j\bigr).
\]

\emph{Multiplying out.} We take absolute values in the $k$-fold composite. Each one-step factor is
bounded by \eqref{10.6:eq-uniform-majorant-2}, and each phase by the bound just stated. This bounds
$|Q(c,b;\mathsf{f})|$ by the composite of the non-negative kernels
\[
e^{C_L\ell_j^2|\mathsf{f}|}\mathsf{Q}^{(0)}_{j+1}+C_L\ell_j^2|\mathsf{f}|\,\bar{E}_{j+1},
\qquad 0\le j\le k-1,
\]
multiplied by $\prod_j\exp(2|q|\,|\mathsf{f}|\,\ell_{j+1}\ell_j)$. We now use
\[
\sum_j\ell_j^2\le 2,\qquad \sum_j\ell_j\ell_{j+1}\le 2 .
\]
Multiplying out this composite and using these two inequalities yields a non-negative kernel which
dominates $|Q(c,b;\mathsf{f})|$ and whose $\mathcal{X}_0\to\mathcal{X}_k$ norm is bounded uniformly
in $k$.

The resulting bound, for complex $|\mathsf{f}|\le c_L$, is \cite[(143)]{B-CMP98} read in the present
notation. It gives
\eqref{10.6:eq-uniform-majorant-a} with a non-negative kernel $\bar{Q}$ supported on
$\{|x_b-x_c|\le 2\}$ and with $\|\bar{Q}\|_{\mathcal{X}_0\to\mathcal{X}_k}\le C_Q$ uniformly in $k$.

For $Q'$ the flux enters the one-step factors only through the explicit phases. The same argument
therefore gives the same bound.

\emph{Cauchy estimates.} Since $\ell_j=L^{j-k}\le L^{-1}$ for $j\le k-1$, we have
$|\ell_j^2\mathsf{f}|<c_L$ on a neighbourhood of the closed disc $|\mathsf{f}|\le c_L$. By the first
step each one-step reduced kernel is analytic there, and the phases are entire, so the finite
$\star$-composite $Q(c,b;\cdot)$ is analytic on a neighbourhood of $|\mathsf{f}|\le c_L$. By
\eqref{10.6:eq-uniform-majorant-a}, it is bounded by $\bar{Q}(c,b)$ on $|\mathsf{f}|\le c_L$. Fix
$\mathsf{f}\ne 0$. The degree-$n$ part satisfies $Q^{(n)}(z\mathsf{f})=z^nQ^{(n)}(\mathsf{f})$, and it is
the $n$-th Taylor coefficient of $z\mapsto Q(c,b;z\mathsf{f})$. The Cauchy estimate on the disc
$|z|\le c_L/|\mathsf{f}|$ gives
\[
|Q^{(n)}(c,b;\mathsf{f})|\le c_L^{-n}|\mathsf{f}|^n\,\bar{Q}(c,b).
\]
The position derivation $\delta_\kappa$ multiplies the kernel by the $\kappa$-th component of
$x_b-x_c$, up to sign. On the support of $\bar{Q}$ this component has modulus at most $2$, so
\[
|\delta^{\alpha}Q^{(n)}(c,b;\mathsf{f})|\le 2^{|\alpha|}c_L^{-n}|\mathsf{f}|^n\,\bar{Q}(c,b).
\]
An operator whose kernel is dominated entrywise by a non-negative kernel has norm at most the norm of
that kernel. Combined with $\|\bar{Q}\|\le C_Q$, this gives \eqref{10.6:eq-uniform-majorant-1} for $Q$.
The same argument gives \eqref{10.6:eq-uniform-majorant-1} for $Q'$.
\end{proof}

\begin{definition}[Conditional covariances for Hessian, gauge and averaging data]
\label{10.6:def-conditional-covariance}
In this definition a \emph{morphism} from one space to another is an element of the ring
$\mathfrak{A}_{\mathsf{f}}$ of tame translation-invariant operators truncated modulo
$(\mathsf{f})^3$, Definition~\ref{10.6:def-truncated-product}, acting between
these spaces. Products of morphisms are $\star$-products, and adjoints are taken in
$\mathfrak{A}_{\mathsf{f}}$.

The data are the following:
\begin{itemize}
\item a morphism $\mathsf{H}$ on a space $\mathcal{X}$ with $\mathsf{H} = \mathsf{H}^{\dagger}$;
\item a morphism $D \colon \mathcal{C} \to \mathcal{X}$ with $\mathsf{H} D = 0$;
\item a morphism $\mathsf{Q} \colon \mathcal{X} \to \mathcal{Y}$;
\item a morphism $\Xi \colon \mathcal{C}' \to \mathcal{Y}$.
\end{itemize}
A \emph{conditional covariance} for $(\mathsf{H}, D, \mathsf{Q}, \Xi)$ is a morphism
$\mathcal{G}$ on $\mathcal{X}$ with $\mathcal{G} = \mathcal{G}^{\dagger}$ that satisfies the
two conditions
\begin{equation}\label{10.6:eq-conditional-covariance-a}
\text{(c1)}\quad \mathsf{Q}\,\mathcal{G} = \Xi\,\Psi_0 ,
\qquad\qquad
\text{(c2)}\quad \mathcal{G}\,\mathsf{H} A = A + D\,\Gamma ,
\end{equation}
where in (c1) $\Psi_0 \colon \mathcal{X} \to \mathcal{C}'$ is some morphism, and where (c2)
is required in the following sense: for every morphism $A$ into $\mathcal{X}$ such that
$\mathsf{Q} A = \Xi\,\Psi$ for some morphism $\Psi$, there exists a morphism $\Gamma$, which
may depend on $A$, for which the identity (c2) holds.
In (c2), $\Psi$ takes values in $\mathcal{C}'$ and $\Gamma$ takes values in $\mathcal{C}$.
Both have the same domain as $A$.
\end{definition}

\begin{lemma}[Uniqueness of the trace on gauge-invariant forms]\label{10.6:lem-trace-unique}
Let $\mathcal{G}_1$ and $\mathcal{G}_2$ be conditional covariances, in the sense of
Definition~\ref{10.6:def-conditional-covariance}, for the same data
$(\mathsf{H}, D, \mathsf{Q}, \mathsf{D})$, and let $P$ satisfy $PD = 0$ and $D^{\dagger}P = 0$. Then
\[
\operatorname{Tr}_{\square}(\mathcal{G}_1 P) = \operatorname{Tr}_{\square}(\mathcal{G}_2 P).
\]
\end{lemma}

\begin{proof}
By the first condition in \eqref{10.6:eq-conditional-covariance-a} of
Definition~\ref{10.6:def-conditional-covariance}, applied to $\mathcal{G}_2$,
there is a morphism $X_0$ with $\mathsf{Q}\mathcal{G}_2 = \mathsf{D}X_0$. Hence $A := \mathcal{G}_2$
is a morphism into $\mathcal{X}$ with $\mathsf{Q}A = \mathsf{D}X_0$, and the second condition in
\eqref{10.6:eq-conditional-covariance-a}, applied to
$\mathcal{G}_1$, applies to it: there is $Y$ with
\[
\mathcal{G}_1\mathsf{H}\mathcal{G}_2 = \mathcal{G}_2 + DY .
\]
In the same way, the first condition in \eqref{10.6:eq-conditional-covariance-a} for
$\mathcal{G}_1$ makes $A := \mathcal{G}_1$ admissible in the second condition for
$\mathcal{G}_2$, which gives $Y'$ with $\mathcal{G}_2\mathsf{H}\mathcal{G}_1 = \mathcal{G}_1 + DY'$.
Taking adjoints and using $\mathcal{G}_1 = \mathcal{G}_1^{\dagger}$,
$\mathcal{G}_2 = \mathcal{G}_2^{\dagger}$, $\mathsf{H} = \mathsf{H}^{\dagger}$, we get
\[
\mathcal{G}_1\mathsf{H}\mathcal{G}_2 = \mathcal{G}_1 + Y'^{\dagger}D^{\dagger}.
\]
Comparing the two expressions for $\mathcal{G}_1\mathsf{H}\mathcal{G}_2$ yields
$\mathcal{G}_1 - \mathcal{G}_2 = DY - Y'^{\dagger}D^{\dagger}$, and by linearity of
$\operatorname{Tr}_{\square}$,
\[
\operatorname{Tr}_{\square}\bigl((\mathcal{G}_1 - \mathcal{G}_2)P\bigr)
= \operatorname{Tr}_{\square}(DYP) - \operatorname{Tr}_{\square}(Y'^{\dagger}D^{\dagger}P).
\]
The second term vanishes because $D^{\dagger}P = 0$. For the first we use the cyclicity of the
trace per cell, which comes from \eqref{10.6:eq-truncated-product-b}: for operators $t, s$ it reads
\[
\operatorname{Tr}_{\square}(t\star s) = \operatorname{Tr}_{\square}(ts)
= \operatorname{Tr}_{\square}(s\star t).
\]
Comparing its second and third members gives
$\operatorname{Tr}_{\square}(ts) = \operatorname{Tr}_{\square}(s\star t)$.
The same identity written for the pair $(s,t)$ gives
$\operatorname{Tr}_{\square}(s\star t) = \operatorname{Tr}_{\square}(st)$ from its first two members.
Hence $\operatorname{Tr}_{\square}(ts) = \operatorname{Tr}_{\square}(st)$.

Applied with $t = D$ and $s = YP$, this gives
$\operatorname{Tr}_{\square}(DYP) = \operatorname{Tr}_{\square}(YPD)$, and $YPD = 0$ because $PD = 0$.
Hence $\operatorname{Tr}_{\square}\bigl((\mathcal{G}_1 - \mathcal{G}_2)P\bigr) = 0$.
\end{proof}

\begin{lemma}[The one-step propagator and its algebraic identities]
\label{10.6:lem-one-step-identities}
Let $\mathsf{H}_j=\mathsf{H}_j^{\dagger}$ be an operator on $\mathcal{X}_j$ with $\mathsf{H}_jD_j=0$. For
$j=0$ such an operator is supplied by the bare-lattice case of
Lemma~\ref{10.6:lem-background-generators}. Write $\tilde{\mathsf{Q}}=\tilde{\mathsf{Q}}_{j+1}$ and
$\mathsf{s}=\mathsf{s}_{j+1}$. Let $c>0$ be the weight with which $\mathsf{s}^{\dagger}\mathsf{s}$ enters
the quadratic part $G^{(2)}$ of the contour term in \cite[(1.5)]{B-CMP109}, so that
$G^{(2)}=c\,\mathsf{s}^{\dagger}\mathsf{s}$, and put
\begin{equation}\label{10.6:eq-one-step-identities-1}
\begin{aligned}
\mathsf{F}&:=\mathsf{H}_j+c\,\mathsf{s}^{\dagger}\mathsf{s},\qquad
\mathsf{K}:=\bigl(\mathsf{F}+\tilde{\mathsf{Q}}^{\dagger}\tilde{\mathsf{Q}}\bigr)^{-1},\\
\mathsf{N}_{j+1}&:=\mathsf{K}\tilde{\mathsf{Q}}^{\dagger}
 \bigl(\tilde{\mathsf{Q}}\mathsf{K}\tilde{\mathsf{Q}}^{\dagger}\bigr)^{-1},\qquad
\mathsf{H}_{j+1}:=\bigl(\tilde{\mathsf{Q}}\mathsf{K}\tilde{\mathsf{Q}}^{\dagger}\bigr)^{-1}-1,\\
\mathcal{G}_P^{(j+1)}&:=\mathsf{K}-\mathsf{K}\tilde{\mathsf{Q}}^{\dagger}
 \bigl(\tilde{\mathsf{Q}}\mathsf{K}\tilde{\mathsf{Q}}^{\dagger}\bigr)^{-1}\tilde{\mathsf{Q}}\mathsf{K}.
\end{aligned}
\end{equation}
The identities below hold whenever the inverses in these definitions exist. For the operators
$\mathsf{H}_0,\mathsf{H}_1,\dots$ obtained by iterating \eqref{10.6:eq-one-step-identities-1} from
the bare-lattice case $j=0$ of Lemma~\ref{10.6:lem-background-generators}, these inverses exist at
each fixed $j$. If $c$ is replaced by any other positive constant, the identities (f1)--(f4) below
hold for the operators so modified, with the same proofs, and $\mathsf{N}_{j+1}$ and
$\mathsf{H}_{j+1}$ are unchanged. With
$\mathsf{N}=\mathsf{N}_{j+1}$ and $\mathcal{G}_P=\mathcal{G}_P^{(j+1)}$ one has
\begin{equation}\label{10.6:eq-one-step-identities-a}
\begin{aligned}
&\text{(f1)}\quad \tilde{\mathsf{Q}}\mathsf{N}=1,\quad \tilde{\mathsf{Q}}\mathcal{G}_P=0,\quad
 \mathcal{G}_P\mathsf{F}=1-\mathsf{N}\tilde{\mathsf{Q}},\quad
 \mathsf{F}\mathsf{N}=\tilde{\mathsf{Q}}^{\dagger}\mathsf{H}_{j+1};\\
&\text{(f2)}\quad \mathsf{s}\mathsf{N}=0,\quad\text{hence}\quad
 \mathsf{H}_j\mathsf{N}=\tilde{\mathsf{Q}}^{\dagger}\mathsf{H}_{j+1},\quad
 \mathsf{H}_{j+1}=\mathsf{N}^{\dagger}\mathsf{H}_j\mathsf{N};\\
&\text{(f3)}\quad \mathsf{N}D_{j+1}=D_j\mathsf{e},\quad \mathsf{H}_{j+1}D_{j+1}=0;\\
&\text{(f4)}\quad \mathcal{G}_P^{(j+1)}\ \text{is a conditional covariance for}\\
&\qquad\quad (\mathsf{H}_j,D_j,\tilde{\mathsf{Q}}_{j+1},D_{j+1})
 \ \text{(Definition~\ref{10.6:def-conditional-covariance})}.
\end{aligned}
\end{equation}
Moreover, $\mathcal{G}_P^{(j+1)}$ is the inverse of $\mathsf{F}$ on $\ker\tilde{\mathsf{Q}}$. If
$\mathsf{H}_j$ is the jet of $\nabla^2\mathcal{A}_j$ along the constant abelian configuration
$\mathcal{U}_{\mathsf{f}}$, as are the operators obtained by iteration (see the remark following the
proof), then $\mathcal{G}_P^{(j+1)}$ coincides with the jet along $\mathcal{U}_{\mathsf{f}}$ of
$C(C^{*}\Delta^{(j)}C)^{-1}C^{*}$, the covariance of $B'=CB$ under the Gaussian fluctuation measure
of \cite[(2.2)--(2.5)]{B-CMP109}.
\end{lemma}

\begin{proof}
Throughout, we use the following relations between $D_j$, $D_{j+1}$ and the one-step maps, which
are given by \eqref{10.6:eq-background-generators-b}:
\[
\mathsf{s}D_j\iota=1,\qquad \mathsf{s}D_j\mathsf{e}=0,\qquad
\tilde{\mathsf{Q}}D_j\iota=D_{j+1}\mathrm{ev}\,\iota=0,\qquad \tilde{\mathsf{Q}}D_j\mathsf{e}=D_{j+1}.
\]
The operator in the third relation vanishes because $\mathrm{ev}\,\iota=0$: $\iota$ extends by zero
off the coarse lattice and $\mathrm{ev}$ restricts to it. The operator $\mathsf{K}$ is self-adjoint, since
$\mathsf{H}_j$, $\mathsf{s}^{\dagger}\mathsf{s}$ and $\tilde{\mathsf{Q}}^{\dagger}\tilde{\mathsf{Q}}$ are.
Hence $\mathcal{G}_P^{\dagger}=\mathcal{G}_P$ and $\mathsf{H}_{j+1}^{\dagger}=\mathsf{H}_{j+1}$.
Note also that $\mathcal{G}_P=\mathsf{K}-\mathsf{N}\tilde{\mathsf{Q}}\mathsf{K}$ by definition.

\emph{Flat invertibility at fixed $j$.} We argue by induction on $j$. Suppose that
$\mathsf{H}_0,\dots,\mathsf{H}_j$ have been obtained by iterating
\eqref{10.6:eq-one-step-identities-1} from the case $j=0$ of
Lemma~\ref{10.6:lem-background-generators}, and that $\mathsf{H}_j$ is the jet of
$\nabla^2\mathcal{A}_j$ along $\mathcal{U}_{\mathsf{f}}$: for $j=0$ this is
Lemma~\ref{10.6:lem-background-generators}, and the step from $j$ to $j+1$ is the first item of the
remark following the proof. Let $\gamma_0>0$ be the lower bound of the operator $C^{*}\Delta^{(j)}C$
at the unit configuration, used here as stated in \cite{B-CMP99b}. The flat parts are the values at
$\mathsf{f}=0$, that is, at the unit
configuration. There, by the identification at the end of the proof taken at $\mathsf{f}=0$, the
flat part of $\Delta^{(j)}$ is $\mathsf{H}_j^{(0)}+c\,\mathsf{s}^{(0)\dagger}\mathsf{s}^{(0)}$, and
the second term vanishes on $\{\mathsf{s}^{(0)}B=0\}$. By the bound,
the flat part $\mathsf{H}_j^{(0)}$ thus satisfies
$\langle B,\mathsf{H}_j^{(0)}B\rangle\ge\gamma_0\|B\|^2$ on
$\{\tilde{\mathsf{Q}}^{(0)}B=0,\ \mathsf{s}^{(0)}B=0\}$.

Let $B\in\ker\tilde{\mathsf{Q}}^{(0)}$ and write $B=B_1+D_j\iota\mathsf{s}^{(0)}B$ with
$B_1:=(1-D_j\iota\mathsf{s}^{(0)})B$. The relations above hold for every $\mathsf{f}$, hence for the
flat parts; they give $\tilde{\mathsf{Q}}^{(0)}B_1=0$ and
$\mathsf{s}^{(0)}B_1=\mathsf{s}^{(0)}B-\mathsf{s}^{(0)}B=0$. Since $\mathsf{H}_jD_j=0$, we have
$\langle B,\mathsf{H}_j^{(0)}B\rangle=\langle B_1,\mathsf{H}_j^{(0)}B_1\rangle$. Therefore
\[
\langle B,\mathsf{F}^{(0)}B\rangle\ge\gamma_0\|B_1\|^2+c\,\|\mathsf{s}^{(0)}B\|^2,
\qquad \|B\|\le\|B_1\|+\|D_j\iota\|\,\|\mathsf{s}^{(0)}B\|,
\]
so $\mathsf{F}^{(0)}$ is bounded below on $\ker\tilde{\mathsf{Q}}^{(0)}$ by a positive constant. In
addition, $\tilde{\mathsf{Q}}^{(0)}$ has a bounded right inverse, namely the flat part of the operator
denoted $h$ in \cite[(1.5)]{B-CMP109}. By these two facts the flat parts of
$\mathsf{F}+\tilde{\mathsf{Q}}^{\dagger}\tilde{\mathsf{Q}}$ and of
$\tilde{\mathsf{Q}}\mathsf{K}\tilde{\mathsf{Q}}^{\dagger}$ are invertible. The operators
$\tilde{\mathsf{Q}}$, $\mathsf{s}$ and $\mathsf{H}_0$ lie in the truncated ring $\mathfrak{A}_{\mathsf{f}}$
by Lemma~\ref{10.6:lem-background-generators}, and so does $\mathsf{H}_j$, which is obtained from
them by the operations in \eqref{10.6:eq-one-step-identities-1}, whose inverses exist in
$\mathfrak{A}_{\mathsf{f}}$ at the lower levels by this step applied there; hence
$\mathsf{F}+\tilde{\mathsf{Q}}^{\dagger}\tilde{\mathsf{Q}}$ and
$\tilde{\mathsf{Q}}\mathsf{K}\tilde{\mathsf{Q}}^{\dagger}$ lie in $\mathfrak{A}_{\mathsf{f}}$. An element
of $\mathfrak{A}_{\mathsf{f}}$ whose flat part is invertible is itself invertible: its terms of
positive degree in $\mathsf{f}$ are nilpotent modulo $(\mathsf{f})^3$, so its inverse is a finite
geometric series. Hence the inverses in \eqref{10.6:eq-one-step-identities-1} exist at fixed $j$.

\emph{(f1).} First,
$\tilde{\mathsf{Q}}\mathsf{N}=\tilde{\mathsf{Q}}\mathsf{K}\tilde{\mathsf{Q}}^{\dagger}
(\tilde{\mathsf{Q}}\mathsf{K}\tilde{\mathsf{Q}}^{\dagger})^{-1}=1$. Consequently
$\tilde{\mathsf{Q}}\mathcal{G}_P=\tilde{\mathsf{Q}}\mathsf{K}-\tilde{\mathsf{Q}}\mathsf{N}\tilde{\mathsf{Q}}\mathsf{K}=0$.

For the third identity, start from
$\mathsf{K}\mathsf{F}=1-\mathsf{K}\tilde{\mathsf{Q}}^{\dagger}\tilde{\mathsf{Q}}$ and from
$\mathsf{N}\tilde{\mathsf{Q}}\mathsf{K}\tilde{\mathsf{Q}}^{\dagger}=\mathsf{K}\tilde{\mathsf{Q}}^{\dagger}$.
They give
\[
\mathcal{G}_P\mathsf{F}=(1-\mathsf{N}\tilde{\mathsf{Q}})\mathsf{K}\mathsf{F}
=1-\mathsf{N}\tilde{\mathsf{Q}}-\mathsf{K}\tilde{\mathsf{Q}}^{\dagger}\tilde{\mathsf{Q}}
+\mathsf{N}\tilde{\mathsf{Q}}\mathsf{K}\tilde{\mathsf{Q}}^{\dagger}\tilde{\mathsf{Q}}
=1-\mathsf{N}\tilde{\mathsf{Q}}.
\]
Finally,
$(\mathsf{F}+\tilde{\mathsf{Q}}^{\dagger}\tilde{\mathsf{Q}})\mathsf{N}
=\tilde{\mathsf{Q}}^{\dagger}(\tilde{\mathsf{Q}}\mathsf{K}\tilde{\mathsf{Q}}^{\dagger})^{-1}$.
Subtracting $\tilde{\mathsf{Q}}^{\dagger}\tilde{\mathsf{Q}}\mathsf{N}=\tilde{\mathsf{Q}}^{\dagger}$
gives $\mathsf{F}\mathsf{N}=\tilde{\mathsf{Q}}^{\dagger}\mathsf{H}_{j+1}$.

\emph{(f2).} Apply $(D_j\iota)^{\dagger}$ to $\mathsf{F}\mathsf{N}=\tilde{\mathsf{Q}}^{\dagger}\mathsf{H}_{j+1}$.

On the left side,
$(D_j\iota)^{\dagger}\mathsf{H}_j\mathsf{N}=(\mathsf{H}_jD_j\iota)^{\dagger}\mathsf{N}=0$. What
remains is $c\,(\mathsf{s}D_j\iota)^{\dagger}\mathsf{s}\mathsf{N}=c\,\mathsf{s}\mathsf{N}$.

On the right side,
\[
(D_j\iota)^{\dagger}\tilde{\mathsf{Q}}^{\dagger}\mathsf{H}_{j+1}
=(\tilde{\mathsf{Q}}D_j\iota)^{\dagger}\mathsf{H}_{j+1}
=(D_{j+1}\mathrm{ev}\,\iota)^{\dagger}\mathsf{H}_{j+1}=0.
\]

Hence $c\,\mathsf{s}\mathsf{N}=0$, and since $c>0$, $\mathsf{s}\mathsf{N}=0$. It follows that
$\mathsf{H}_j\mathsf{N}=\mathsf{F}\mathsf{N}-c\,\mathsf{s}^{\dagger}\mathsf{s}\mathsf{N}
=\tilde{\mathsf{Q}}^{\dagger}\mathsf{H}_{j+1}$. Using (f1),
$\mathsf{N}^{\dagger}\mathsf{H}_j\mathsf{N}=(\tilde{\mathsf{Q}}\mathsf{N})^{\dagger}\mathsf{H}_{j+1}
=\mathsf{H}_{j+1}$.

\emph{Independence of the weight.} Let $\mathsf{N}'$ and $\mathsf{H}'_{j+1}$ be formed as in
\eqref{10.6:eq-one-step-identities-1} with another positive constant in place of $c$. The proofs of
(f1) and (f2) apply to them and give $\tilde{\mathsf{Q}}\mathsf{N}'=1$, $\mathsf{s}\mathsf{N}'=0$ and
$\mathsf{H}_j\mathsf{N}'=\tilde{\mathsf{Q}}^{\dagger}\mathsf{H}'_{j+1}$. Hence
\[
(\mathsf{F}+\tilde{\mathsf{Q}}^{\dagger}\tilde{\mathsf{Q}})\mathsf{N}'
=\mathsf{H}_j\mathsf{N}'+c\,\mathsf{s}^{\dagger}\mathsf{s}\mathsf{N}'
+\tilde{\mathsf{Q}}^{\dagger}\tilde{\mathsf{Q}}\mathsf{N}'
=\tilde{\mathsf{Q}}^{\dagger}(\mathsf{H}'_{j+1}+1),
\]
so $\mathsf{N}'=\mathsf{K}\tilde{\mathsf{Q}}^{\dagger}(\mathsf{H}'_{j+1}+1)$. Applying
$\tilde{\mathsf{Q}}$ gives
$1=(\tilde{\mathsf{Q}}\mathsf{K}\tilde{\mathsf{Q}}^{\dagger})(\mathsf{H}'_{j+1}+1)$, hence
$\mathsf{H}'_{j+1}=(\tilde{\mathsf{Q}}\mathsf{K}\tilde{\mathsf{Q}}^{\dagger})^{-1}-1=\mathsf{H}_{j+1}$,
and then $\mathsf{N}'=\mathsf{N}$.

\emph{(f3).} Let $A:=D_j\mathsf{e}$. Then
$\mathsf{F}A=\mathsf{H}_jD_j\mathsf{e}+c\,\mathsf{s}^{\dagger}\mathsf{s}D_j\mathsf{e}=0$ and
$\tilde{\mathsf{Q}}A=D_{j+1}$. Hence
$(\mathsf{F}+\tilde{\mathsf{Q}}^{\dagger}\tilde{\mathsf{Q}})A=\tilde{\mathsf{Q}}^{\dagger}D_{j+1}$,
that is, $A=\mathsf{K}\tilde{\mathsf{Q}}^{\dagger}D_{j+1}$. Applying $\tilde{\mathsf{Q}}$ gives
$D_{j+1}=\tilde{\mathsf{Q}}A=(\tilde{\mathsf{Q}}\mathsf{K}\tilde{\mathsf{Q}}^{\dagger})D_{j+1}$, so
$(\tilde{\mathsf{Q}}\mathsf{K}\tilde{\mathsf{Q}}^{\dagger})^{-1}D_{j+1}=D_{j+1}$. Therefore
$\mathsf{N}D_{j+1}=\mathsf{K}\tilde{\mathsf{Q}}^{\dagger}D_{j+1}=D_j\mathsf{e}$ and
$\mathsf{H}_{j+1}D_{j+1}=D_{j+1}-D_{j+1}=0$.

\emph{(f4).} We verify the two conditions of
\eqref{10.6:eq-conditional-covariance-a}. We have already noted that $\mathcal{G}_P^{\dagger}=\mathcal{G}_P$.

Condition (c1) holds with $X_0=0$, by $\tilde{\mathsf{Q}}\mathcal{G}_P=0$ from (f1).

For (c2), let $\tilde{\mathsf{Q}}A=D_{j+1}X$ and put
$A_1:=(1-D_j\iota\mathsf{s})(A-D_j\mathsf{e}X)$. Then
$\tilde{\mathsf{Q}}(A-D_j\mathsf{e}X)=D_{j+1}X-D_{j+1}X=0$ and $\tilde{\mathsf{Q}}D_j\iota=0$, so
$\tilde{\mathsf{Q}}A_1=0$. Also $\mathsf{s}D_j\iota=1$, so $\mathsf{s}A_1=0$. Moreover,
\[
A-A_1=D_j\bigl(\mathsf{e}X+\iota\mathsf{s}(A-D_j\mathsf{e}X)\bigr),
\]
so that $A-A_1$ lies in the range of $D_j$.
Since $\mathsf{H}_jD_j=0$ and $\mathsf{s}A_1=0$, (f1) gives
\[
\mathcal{G}_P\mathsf{H}_jA=\mathcal{G}_P\mathsf{H}_jA_1=\mathcal{G}_P\mathsf{F}A_1
=(1-\mathsf{N}\tilde{\mathsf{Q}})A_1=A_1=A+D_jY,
\]
with $Y:=-\mathsf{e}X-\iota\mathsf{s}(A-D_j\mathsf{e}X)$.

\emph{Identification.} By (f1), $\mathcal{G}_P$ takes values in $\ker\tilde{\mathsf{Q}}$. For
$B\in\ker\tilde{\mathsf{Q}}$ we have $\mathcal{G}_P\mathsf{F}B=B$. In this sense $\mathcal{G}_P$ is the
inverse of $\mathsf{F}$ on $\ker\tilde{\mathsf{Q}}$.

Let now $\mathsf{H}_j$ be the jet of $\nabla^2\mathcal{A}_j$ along $\mathcal{U}_{\mathsf{f}}$. Let $J$
be the current entering \cite[(1.5)]{B-CMP109}, and let $\Delta^{(j)}$ be the operator of the
quadratic form of the fluctuation integral of \cite[(2.2)--(2.5)]{B-CMP109}. Since
$\mathcal{U}_{\mathsf{f}}$ is critical (Lemma~\ref{10.6:lem-constant-fields}),
$J(\mathcal{U}_{\mathsf{f}})=0$, and \cite[(1.5)]{B-CMP109} gives
$\Delta^{(j)}=\mathsf{H}_j+G^{(2)}$, with $G^{(2)}$ as in the first step of the proof. Since
$G^{(2)}=c\,\mathsf{s}^{\dagger}\mathsf{s}$, this says $\Delta^{(j)}=\mathsf{F}$. The map
$B\mapsto B'=CB$ parametrises the fluctuation variables satisfying the linearised averaging
constraint of \cite[(2.2)--(2.5)]{B-CMP109}, that is, $\ker\tilde{\mathsf{Q}}$. Hence
$C(C^{*}\Delta^{(j)}C)^{-1}C^{*}$ is the inverse of $\Delta^{(j)}$, and therefore of $\mathsf{F}$, on
$\ker\tilde{\mathsf{Q}}$.

It coincides with $\mathcal{G}_P^{(j+1)}$. Both operators are self-adjoint and vanish on
$\operatorname{ran}\tilde{\mathsf{Q}}^{\dagger}=(\ker\tilde{\mathsf{Q}})^{\perp}$: for $\mathcal{G}_P$
because $\mathcal{G}_P\tilde{\mathsf{Q}}^{\dagger}=(\tilde{\mathsf{Q}}\mathcal{G}_P)^{\dagger}=0$ by (f1),
and for the other because $C^{*}$ vanishes on $(\operatorname{ran}C)^{\perp}$ and
$\operatorname{ran}C=\ker\tilde{\mathsf{Q}}$. An operator $T$ with these two properties which inverts
$\mathsf{F}$ on $\ker\tilde{\mathsf{Q}}$ is unique. Indeed, such a $T$ takes values in
$\ker\tilde{\mathsf{Q}}$, and for every $x\in\mathcal{X}_j$ and every $B\in\ker\tilde{\mathsf{Q}}$,
\[
\langle B,\mathsf{F}Tx\rangle=\langle T\mathsf{F}B,x\rangle=\langle B,x\rangle,
\]
so $Tx$ is an element of $\ker\tilde{\mathsf{Q}}$ whose image under $\mathsf{F}$ has the same
orthogonal projection onto $\ker\tilde{\mathsf{Q}}$ as $x$; this element is unique because the
compression of $\mathsf{F}$ to $\ker\tilde{\mathsf{Q}}$ is invertible, its flat part being bounded
below by the first step. Hence $C(C^{*}\Delta^{(j)}C)^{-1}C^{*}=\mathcal{G}_P^{(j+1)}$.
\end{proof}

\begin{remark}
By (f2) and (f3), $\mathsf{H}_{j+1}$ is again self-adjoint and satisfies
$\mathsf{H}_{j+1}D_{j+1}=0$. The lemma can therefore be applied level by level, starting from the
case $j=0$ of Lemma~\ref{10.6:lem-background-generators}.

The covariance $C(C^{*}\Delta^{(j)}C)^{-1}C^{*}$ of $B'=CB$ is the operator written
$\mathcal{G}_P^{(j+1)}(W)$ in Lemma~\ref{10.6:lem-half-trace}; along $\mathcal{U}_{\mathsf{f}}$ it
coincides, by the last assertion of Lemma~\ref{10.6:lem-one-step-identities}, with the operator
defined in \eqref{10.6:eq-one-step-identities-1}, which justifies the use of one symbol for both.

The operators so defined are identified as follows. By \eqref{10.6:eq-one-step-identities-1}, (f1)
and (f2),
\[
\tfrac12\langle B,\mathsf{H}_{j+1}B\rangle
=\min\bigl\{\tfrac12\langle B',\mathsf{H}_jB'\rangle:\ \tilde{\mathsf{Q}}B'=B\bigr\}.
\]
By \eqref{10.6:eq-constant-fields-a}, $\mathcal{U}_{\mathsf{f}}$ is an unconstrained critical point
of $\mathcal{A}_j$ (Lemma~\ref{10.6:lem-constant-fields}) which satisfies the averaging constraint,
its average being $\mathcal{U}_{\mathsf{f}}$ on the next lattice. Hence at $\mathcal{U}_{\mathsf{f}}$
the first-order term of $\mathcal{A}_j$ and the multiplier of the averaging constraint vanish, and
the second-order jet of $\mathcal{A}_{j+1}$ along $\mathcal{U}_{\mathsf{f}}$ is the constrained
minimum of the second-order jet of $\mathcal{A}_j$, which is the right side of the last display.
With this, \eqref{10.6:eq-truncated-ring-a} and Remark~\ref{10.6:rem-inner-maps} (inner maps enter
only through their linearisations):
\begin{itemize}
\item the operator $\mathsf{H}_j$ is the jet of $\nabla^2\mathcal{A}_j$ along $\mathcal{U}_{\mathsf{f}}$;
\item $\mathsf{N}_{(j)}:=\mathsf{N}_1\cdots\mathsf{N}_j$ is the jet along $\mathcal{U}_{\mathsf{f}}$ of
the derivative $DU_j$ of the composite map $U_j$ through which $\mathcal{A}_j$ is defined, modulo
the range of $D_0$;
\item $\mathsf{O}_j:=\mathsf{N}_{(j)}^{\dagger}\mathsf{O}\mathsf{N}_{(j)}$ is the jet of
$\nabla^2\mathcal{S}_j$.
\end{itemize}
Finally, $\mathsf{O}_jD_j=0$. This follows from (f3), iterated over the levels, together with
\eqref{10.6:eq-background-generators-a}.
\end{remark}

\begin{lemma}[The Gaussian tower]\label{10.6:lem-gaussian-tower}
Let $\mathsf{H}_0$ be the total Hessian on $\mathcal{X}_0$ of Lemma~\ref{10.6:lem-background-generators}, and for
$i\ge 0$ let $\mathsf{N}_{i+1}$, $\mathcal{G}_P^{(i+1)}$, $\mathsf{H}_{i+1}$ be the one-step objects of
Lemma~\ref{10.6:lem-one-step-identities}, built at level $i$ from $\mathsf{H}_i$ with
$\tilde{\mathsf{Q}}_{i+1}$, $\mathsf{s}_{i+1}$, and let $\mathsf{e}_{i+1}$ be the transport-extension
$\mathsf{e}$ of the step from level $i$ to level $i+1$. We recall the $j$-fold composites
\[
\mathsf{N}_{(j)}=\mathsf{N}_1\cdots\mathsf{N}_j,\qquad
\tilde{Q}_j=\tilde{\mathsf{Q}}_j\cdots\tilde{\mathsf{Q}}_1,\qquad
\mathsf{e}_{(j)}=\mathsf{e}_1\cdots\mathsf{e}_j
\]
($\tilde{Q}_j$ and $\mathsf{e}_{(j)}$ as in Lemma~\ref{10.6:lem-background-generators}, $\mathsf{N}_{(j)}$ as
in Lemma~\ref{10.6:lem-one-step-identities}), with $\mathsf{N}_{(0)}=1$. For $j\ge 1$ put
\[
\mathcal{G}^{(\le j)}:=\sum_{i=0}^{j-1}\mathsf{N}_{(i)}\,\mathcal{G}_P^{(i+1)}\,\mathsf{N}_{(i)}^{\dagger}.
\]
Then
\begin{equation}\label{10.6:eq-gaussian-tower-a}
\begin{aligned}
&\text{(t1)}\quad \tilde{Q}_j\,\mathcal{G}^{(\le j)}=0;\\
&\text{(t2)}\quad \mathcal{G}^{(\le j)}\ \text{is a conditional covariance for }
 (\mathsf{H}_0,D_0,\tilde{Q}_j,D_j);\\
&\text{(t3)}\quad \tilde{Q}_j\mathsf{N}_{(j)}=1,\qquad
 \mathsf{H}_0\mathsf{N}_{(j)}=\tilde{Q}_j^{\dagger}\mathsf{H}_j,\qquad
 \mathsf{N}_{(j)}D_j=D_0\,\mathsf{e}_{(j)}.
\end{aligned}
\end{equation}
\end{lemma}

\begin{proof}
We first note that Lemma~\ref{10.6:lem-one-step-identities} applies at every level $i$. For $i=0$,
$\mathsf{H}_0=\mathsf{H}_0^{\dagger}$ and $\mathsf{H}_0D_0=0$ by
Lemma~\ref{10.6:lem-background-generators}. For $i\ge1$, assuming these two properties at level $i-1$,
$\mathsf{H}_i=\mathsf{N}_i^{\dagger}\mathsf{H}_{i-1}\mathsf{N}_i$ by (f2) of
\eqref{10.6:eq-one-step-identities-a} at level $i-1$, hence $\mathsf{H}_i$ is self-adjoint, and
$\mathsf{H}_iD_i=0$ by (f3) of \eqref{10.6:eq-one-step-identities-a} at level $i-1$. Hence
(f1)--(f4) of \eqref{10.6:eq-one-step-identities-a} hold at every level. We argue by induction
on $j$.

\emph{The case $j=1$.} Here $\mathcal{G}^{(\le 1)}=\mathcal{G}_P^{(1)}$, $\tilde{Q}_1=\tilde{\mathsf{Q}}_1$,
$\mathsf{N}_{(1)}=\mathsf{N}_1$, $\mathsf{e}_{(1)}=\mathsf{e}_1$. Then (t1) is the identity
$\tilde{\mathsf{Q}}\mathcal{G}_P=0$ of (f1); (t2) is (f4) at level $0$; in (t3), the first identity
is $\tilde{\mathsf{Q}}\mathsf{N}=1$ of (f1), the second is
$\mathsf{H}_0\mathsf{N}_1=\tilde{\mathsf{Q}}_1^{\dagger}\mathsf{H}_1$, the consequence of (f2) at
level $0$, and the third is $\mathsf{N}_1D_1=D_0\mathsf{e}_1$, which is (f3) at level $0$.

\emph{From $j$ to $j+1$.} Assume (t1)--(t3) at $j\ge1$. We use $\tilde{Q}_{j+1}=
\tilde{\mathsf{Q}}_{j+1}\tilde{Q}_j$, $\mathsf{N}_{(j+1)}=\mathsf{N}_{(j)}\mathsf{N}_{j+1}$,
$\mathsf{e}_{(j+1)}=\mathsf{e}_{(j)}\mathsf{e}_{j+1}$ and
\begin{equation}\label{10.6:eq-gaussian-tower-1}
\mathcal{G}^{(\le j+1)}=\mathcal{G}^{(\le j)}+\mathsf{N}_{(j)}\,\mathcal{G}_P^{(j+1)}\,
\mathsf{N}_{(j)}^{\dagger}.
\end{equation}

(t3). By (t3) at $j$ and $\tilde{\mathsf{Q}}\mathsf{N}=1$ of (f1) at level $j$,
$\tilde{Q}_{j+1}\mathsf{N}_{(j+1)}=\tilde{\mathsf{Q}}_{j+1}(\tilde{Q}_j\mathsf{N}_{(j)})\mathsf{N}_{j+1}
=\tilde{\mathsf{Q}}_{j+1}\mathsf{N}_{j+1}=1$. By (t3) at $j$ and then the consequence of (f2) at
level $j$,
\[
\mathsf{H}_0\mathsf{N}_{(j+1)}=\tilde{Q}_j^{\dagger}\mathsf{H}_j\mathsf{N}_{j+1}
=\tilde{Q}_j^{\dagger}\tilde{\mathsf{Q}}_{j+1}^{\dagger}\mathsf{H}_{j+1}
=\tilde{Q}_{j+1}^{\dagger}\mathsf{H}_{j+1}.
\]
By (f3) at level $j$ and (t3) at $j$,
$\mathsf{N}_{(j+1)}D_{j+1}=\mathsf{N}_{(j)}D_j\mathsf{e}_{j+1}=D_0\mathsf{e}_{(j)}\mathsf{e}_{j+1}
=D_0\mathsf{e}_{(j+1)}$.

(t1). By \eqref{10.6:eq-gaussian-tower-1},
\[
\tilde{Q}_{j+1}\mathcal{G}^{(\le j+1)}
=\tilde{\mathsf{Q}}_{j+1}\bigl(\tilde{Q}_j\mathcal{G}^{(\le j)}\bigr)
+\tilde{\mathsf{Q}}_{j+1}\bigl(\tilde{Q}_j\mathsf{N}_{(j)}\bigr)\mathcal{G}_P^{(j+1)}
\mathsf{N}_{(j)}^{\dagger}.
\]
The first term vanishes by (t1) at $j$. By (t3) at $j$ the second equals
$\tilde{\mathsf{Q}}_{j+1}\mathcal{G}_P^{(j+1)}\mathsf{N}_{(j)}^{\dagger}$, which vanishes by
$\tilde{\mathsf{Q}}\mathcal{G}_P=0$ of (f1) at level $j$.

(t2). We check the conditions of Definition~\ref{10.6:def-conditional-covariance} for
$(\mathsf{H}_0,D_0,\tilde{Q}_{j+1},D_{j+1})$; here $\mathsf{H}_0=\mathsf{H}_0^{\dagger}$ and
$\mathsf{H}_0D_0=0$. Each $\mathcal{G}_P^{(i+1)}$ is self-adjoint, being a conditional covariance by
(f4), so each summand $\mathsf{N}_{(i)}\mathcal{G}_P^{(i+1)}\mathsf{N}_{(i)}^{\dagger}$, and hence
$\mathcal{G}^{(\le j+1)}$, is self-adjoint. Condition (c1) of
\eqref{10.6:eq-conditional-covariance-a} holds with $X_0=0$ by (t1) at $j+1$. For (c2), let $A$ be
a morphism into $\mathcal{X}_0$ with $\tilde{Q}_{j+1}A=D_{j+1}X$, and put
\[
B:=\tilde{Q}_jA,\qquad A':=A-\mathsf{N}_{(j)}B .
\]
By (t3) at $j$, $\tilde{Q}_jA'=B-\tilde{Q}_j\mathsf{N}_{(j)}B=0$. Since $A=A'+\mathsf{N}_{(j)}B$,
by \eqref{10.6:eq-gaussian-tower-1} the operator $\mathcal{G}^{(\le j+1)}\mathsf{H}_0A$ is the sum of
$\mathcal{G}^{(\le j)}\mathsf{H}_0A'$, $\mathcal{G}^{(\le j)}\mathsf{H}_0\mathsf{N}_{(j)}B$,
$\mathsf{N}_{(j)}\mathcal{G}_P^{(j+1)}\mathsf{N}_{(j)}^{\dagger}\mathsf{H}_0A'$ and
$\mathsf{N}_{(j)}\mathcal{G}_P^{(j+1)}\mathsf{N}_{(j)}^{\dagger}\mathsf{H}_0\mathsf{N}_{(j)}B$,
which we treat in turn.

First, since $\tilde{Q}_jA'=0=D_j\cdot 0$, condition (c2) for $\mathcal{G}^{(\le j)}$, which holds by
(t2) at $j$, gives $Y$ with $\mathcal{G}^{(\le j)}\mathsf{H}_0A'=A'+D_0Y$.

Second, by (t3) at $j$ and the self-adjointness of $\mathcal{G}^{(\le j)}$,
$\mathcal{G}^{(\le j)}\mathsf{H}_0\mathsf{N}_{(j)}B=\mathcal{G}^{(\le j)}\tilde{Q}_j^{\dagger}\mathsf{H}_jB
=(\tilde{Q}_j\mathcal{G}^{(\le j)})^{\dagger}\mathsf{H}_jB=0$ by (t1) at $j$.

Third, since $\mathsf{H}_0$ and $\mathsf{H}_j$ are self-adjoint, (t3) at $j$ gives
$\mathsf{N}_{(j)}^{\dagger}\mathsf{H}_0=(\mathsf{H}_0\mathsf{N}_{(j)})^{\dagger}=\mathsf{H}_j\tilde{Q}_j$, so
$\mathsf{N}_{(j)}^{\dagger}\mathsf{H}_0A'=\mathsf{H}_j\tilde{Q}_jA'=0$ and the term
$\mathsf{N}_{(j)}\mathcal{G}_P^{(j+1)}\mathsf{N}_{(j)}^{\dagger}\mathsf{H}_0A'$ vanishes.

Fourth, by the same identity and $\tilde{Q}_j\mathsf{N}_{(j)}=1$,
$\mathsf{N}_{(j)}^{\dagger}\mathsf{H}_0\mathsf{N}_{(j)}B=\mathsf{H}_j\tilde{Q}_j\mathsf{N}_{(j)}B=
\mathsf{H}_jB$. Moreover $\tilde{\mathsf{Q}}_{j+1}B=\tilde{Q}_{j+1}A=D_{j+1}X$, so (f4) at level $j$,
that is, condition (c2) for $(\mathsf{H}_j,D_j,\tilde{\mathsf{Q}}_{j+1},D_{j+1})$, gives $Y'$ with
$\mathcal{G}_P^{(j+1)}\mathsf{H}_jB=B+D_jY'$. With (t3) at $j$,
\[
\mathsf{N}_{(j)}\mathcal{G}_P^{(j+1)}\mathsf{N}_{(j)}^{\dagger}\mathsf{H}_0\mathsf{N}_{(j)}B
=\mathsf{N}_{(j)}\bigl(B+D_jY'\bigr)=\mathsf{N}_{(j)}B+D_0\mathsf{e}_{(j)}Y'.
\]

Adding the four terms,
\[
\mathcal{G}^{(\le j+1)}\mathsf{H}_0A=A'+D_0Y+\mathsf{N}_{(j)}B+D_0\mathsf{e}_{(j)}Y'
=A+D_0\bigl(Y+\mathsf{e}_{(j)}Y'\bigr),
\]
which is (c2). This proves (t1)--(t3) at $j+1$.
\end{proof}

\begin{lemma}[The background-field propagator as a conditional covariance]
\label{10.6:lem-representation-covariance}
Work on the constant abelian configuration $\mathcal{U}_{\mathsf{f}}$, with Ba{\l}aban's parameter $a$ of
\cite[(3.26)]{B-CMP99b} equal to $1$. Put
$\Delta_s := D_0^{\dagger}D_0$ and
\begin{equation}\label{10.6:eq-representation-covariance-1}
\begin{split}
G' &:= (\Delta_s + Q'^{\dagger}Q')^{-1}, \qquad
M' := (Q'G'G'Q'^{\dagger})^{-1}, \qquad
R := 1 - G'Q'^{\dagger}M'Q'G',\\
G &:= (\mathsf{H}_0 + D_0RD_0^{\dagger} + Q^{\dagger}Q)^{-1}, \qquad
\mathsf{M} := (QGQ^{\dagger})^{-1}, \qquad
\mathcal{G}_B := G - GQ^{\dagger}\mathsf{M}QG .
\end{split}
\end{equation}
These are the operators of \cite[(3.25)--(3.27)]{B-CMP99b} on $\mathcal{U}_{\mathsf{f}}$. In the flat
case $\mathsf{f} = 0$ the operators inverted in \eqref{10.6:eq-representation-covariance-1} are
invertible by the bounds of \cite[Proposition~1.1, (1.89)--(1.90)]{B-CMP95},
\cite[(1.42)--(1.45)]{B-CMP95} and \cite[(1.99)--(1.101)]{B-CMP95}. Then $\mathcal{G}_B$ is a
conditional covariance, in the sense of
Definition~\ref{10.6:def-conditional-covariance}, for the data $(\mathsf{H}_0, D_0, \tilde{Q}, D_k)$.
\end{lemma}

\begin{proof}
We use three identities of Lemma~\ref{10.6:lem-background-generators}, namely
\eqref{10.6:eq-background-generators-a} and \eqref{10.6:eq-background-generators-d}:
\begin{equation}\label{10.6:eq-representation-covariance-2}
\mathsf{H}_0D_0 = 0, \qquad QD_0 = D_kQ', \qquad Q = \tilde{Q} + D_k\sigma_{(k)} ,
\end{equation}
the last of these being the content of \eqref{10.6:eq-background-generators-d} used below.

\emph{Properties of $R$.} Since $\Delta_s$ and $Q'^{\dagger}Q'$ are self-adjoint, so are $G'$ and $M'$.
Hence $P := G'Q'^{\dagger}M'Q'G'$ is self-adjoint. Moreover
$P^2 = G'Q'^{\dagger}M'(Q'G'G'Q'^{\dagger})M'Q'G' = G'Q'^{\dagger}M'Q'G' = P$, because
$M'^{-1} = Q'G'G'Q'^{\dagger}$. Therefore $R = 1 - P$ satisfies $R = R^{\dagger} = R^2$. Next,
$Q'G'R = Q'G' - (Q'G'G'Q'^{\dagger})M'Q'G' = Q'G' - Q'G' = 0$. Finally, since
$\Delta_s = G'^{-1} - Q'^{\dagger}Q'$,
\[
\Delta_sG'R = (G'^{-1} - Q'^{\dagger}Q')G'R = R - Q'^{\dagger}(Q'G'R) = R .
\]

\emph{The identity $GD_0R = D_0G'R$.} Expand $G^{-1}$ from
\eqref{10.6:eq-representation-covariance-1}. The first identity of
\eqref{10.6:eq-representation-covariance-2} gives $\mathsf{H}_0D_0G'R = 0$. The second gives
$Q^{\dagger}QD_0G'R = Q^{\dagger}D_kQ'G'R$. Together with $D_0^{\dagger}D_0 = \Delta_s$ this yields
\[
G^{-1}D_0G'R = D_0R\Delta_sG'R + Q^{\dagger}D_kQ'G'R = D_0R\cdot R + 0 = D_0R ,
\]
where we used $\Delta_sG'R = R$, $R^2 = R$ and $Q'G'R = 0$. Applying $G$ gives
\begin{equation}\label{10.6:eq-representation-covariance-3}
GD_0R = D_0G'R ,
\end{equation}
the analogue on $\mathcal{U}_{\mathsf{f}}$ of \cite[(1.95), (1.97)]{B-CMP95}. Using
\eqref{10.6:eq-representation-covariance-3} twice, then the second identity of
\eqref{10.6:eq-representation-covariance-2} and $Q'G'R = 0$, we get
\begin{equation}\label{10.6:eq-representation-covariance-4}
\mathcal{G}_BD_0R = (1 - GQ^{\dagger}\mathsf{M}Q)D_0G'R
= D_0G'R - GQ^{\dagger}\mathsf{M}D_kQ'G'R = D_0G'R .
\end{equation}

\emph{Self-adjointness.} The operators $\mathsf{H}_0$, $D_0RD_0^{\dagger}$ (because $R = R^{\dagger}$)
and $Q^{\dagger}Q$ are self-adjoint. Hence $G = G^{\dagger}$, and so $\mathsf{M} = \mathsf{M}^{\dagger}$.
Therefore $\mathcal{G}_B^{\dagger} = G - GQ^{\dagger}\mathsf{M}QG = \mathcal{G}_B$.

\emph{Condition \textup{(c1)}.} By definition of $\mathsf{M}$,
$Q\mathcal{G}_B = QG - (QGQ^{\dagger})\mathsf{M}QG = QG - QG = 0$. By the third identity of
\eqref{10.6:eq-representation-covariance-2},
$\tilde{Q}\mathcal{G}_B = Q\mathcal{G}_B - D_k\sigma_{(k)}\mathcal{G}_B = D_k(-\sigma_{(k)}\mathcal{G}_B)$.
This is \textup{(c1)} with $X_0 := -\sigma_{(k)}\mathcal{G}_B$.

\emph{Condition \textup{(c2)}.} Let $A$ be a morphism into $\mathcal{X}_0$ with
$\tilde{Q}A = D_kX$ for some morphism $X$. By the third identity of
\eqref{10.6:eq-representation-covariance-2},
\[
QA = \tilde{Q}A + D_k\sigma_{(k)}A = D_kX_1, \qquad X_1 := X + \sigma_{(k)}A .
\]
Put $\lambda := Q'^{\dagger}(Q'Q'^{\dagger})^{-1}X_1$ and $A_0 := A - D_0\lambda$. Then, by the second
identity of \eqref{10.6:eq-representation-covariance-2},
\[
QA_0 = D_kX_1 - D_kQ'Q'^{\dagger}(Q'Q'^{\dagger})^{-1}X_1 = 0 .
\]
By the first identity, $\mathsf{H}_0A = \mathsf{H}_0A_0 + \mathsf{H}_0D_0\lambda = \mathsf{H}_0A_0$. Now
$\mathsf{H}_0 = G^{-1} - D_0RD_0^{\dagger} - Q^{\dagger}Q$. We also have
$\mathcal{G}_BG^{-1}A_0 = A_0 - GQ^{\dagger}\mathsf{M}QA_0 = A_0$ and $Q^{\dagger}QA_0 = 0$. By
\eqref{10.6:eq-representation-covariance-4},
$\mathcal{G}_BD_0RD_0^{\dagger}A_0 = D_0G'RD_0^{\dagger}A_0$. Hence
\[
\mathcal{G}_B\mathsf{H}_0A = \mathcal{G}_B(G^{-1} - D_0RD_0^{\dagger} - Q^{\dagger}Q)A_0
= A_0 - D_0G'RD_0^{\dagger}A_0 = A + D_0Y ,
\]
with $Y := -\lambda - G'RD_0^{\dagger}A_0$. This is \textup{(c2)}, and the proof is complete.
\end{proof}

\begin{proposition}[The partial sums as one bare-lattice trace]\label{10.6:prop-bare-trace}
Let $n_{\mathbf{3}}(k)=\sum_{l=1}^{k}\mathsf{m}^{\circ}_{\mathbf{3}}(l)$ be the partial sum of the
cut-free one-loop numbers (see \ref{10.6:not-cut-free-numbers}). Let
$\mathcal{G}_B$ be Ba{\l}aban's conditional covariance,
$G$ the regulated propagator and $\mathsf{M}=(QGQ^{\dagger})^{-1}$. Let $\mathsf{O}$ be the weighted
Hessian, and let $\sum_q$ denote the sum over the charge sectors $\mathfrak{g}_q$. Then at every fixed $k$
\begin{equation}\label{10.6:eq-bare-trace-a}
\begin{split}
n_{\mathbf{3}}(k)\cdot\mathsf{t}^{\mathbf{3}}(h;\mathsf{f})
&=-\sum_q[\mathsf{f}^2]\tfrac12\operatorname{Tr}_{\square}(\mathcal{G}_B\mathsf{O})\\
&=-\sum_q[\mathsf{f}^2]\tfrac12\operatorname{Tr}_{\square}(G\mathsf{O})
+\sum_q[\mathsf{f}^2]\tfrac12\operatorname{Tr}_{\square}
\bigl(\mathsf{M}\cdot QG\mathsf{O}GQ^{\dagger}\bigr).
\end{split}
\end{equation}
\end{proposition}

\begin{proof}
Fix $j$ with $0\le j<k$. Lemma~\ref{10.6:lem-half-trace} gives \eqref{10.6:eq-half-trace-a}, which reads
\[
\mathsf{m}^{\circ}_{\mathbf{3}}(j+1)\cdot\mathsf{t}^{\mathbf{3}}(h;\ell^2_{j+1}\mathsf{f})
=-\mathsf{cf}[\psi^{(j+1)}](\mathsf{f}),
\]
where $\psi^{(j+1)}_y$ is the half trace of $\mathcal{G}_P^{(j+1)}\nabla^2\mathcal{S}_j$ over the
cell $\square_{j+1}(y)$.

We evaluate the right-hand side by Proposition~\ref{10.6:prop-truncated-ring}, that is, by
\eqref{10.6:eq-truncated-ring-a}. This turns the constant-field value $\mathsf{cf}[\psi^{(j+1)}](\mathsf{f})$
into the coefficient $[\mathsf{f}^2]$ of a trace per cell. That trace is taken of the same expression
evaluated with the truncated product $\star$, and it splits over the charge sectors $\mathfrak{g}_q$.

We then pass from the cells $\square_{j+1}$ to the unit cell $\square$ by scale invariance, which has
three ingredients:
\begin{itemize}
\item The four-dimensional Wilson action carries no scale weight, so the step $j\to j+1$, with $j$
lattices beneath it, is the same at every $k$.
\item The unit cell $\square$ contains $\ell_{j+1}^{-4}$ cells $\square_{j+1}$.
\item The symbol satisfies $\mathsf{t}^{\mathbf{3}}(h;\ell^2_{j+1}\mathsf{f})
=\ell^4_{j+1}\mathsf{t}^{\mathbf{3}}(h;\mathsf{f})$.
\end{itemize}
The factor $\ell^4_{j+1}$ on the left is thus compensated by the number $\ell^{-4}_{j+1}$ of cells
$\square_{j+1}$ in $\square$. The contribution of step $j+1$ is therefore
$\mathsf{m}^{\circ}_{\mathbf{3}}(j+1)\mathsf{t}^{\mathbf{3}}(h;\mathsf{f})
=-\sum_q[\mathsf{f}^2]\tfrac12\operatorname{Tr}_{\square}(\mathcal{G}_P^{(j+1)}\mathsf{O}_j)$.
Here $\mathsf{O}_j$ is the vertex Hessian transported to level $j$.

By the definition of the partial sums, \ref{10.6:not-cut-free-numbers},
$n_{\mathbf{3}}(k)=\sum_{l=1}^{k}\mathsf{m}^{\circ}_{\mathbf{3}}(l)
=\sum_{j<k}\mathsf{m}^{\circ}_{\mathbf{3}}(j+1)$. Summing the contributions over $j<k$ gives the first
equality of
\begin{equation}\label{10.6:eq-bare-trace-1}
-n_{\mathbf{3}}(k)\mathsf{t}^{\mathbf{3}}(h;\mathsf{f})
=\sum_q[\mathsf{f}^2]\sum_{j<k}\tfrac12\operatorname{Tr}_{\square}(\mathcal{G}_P^{(j+1)}\mathsf{O}_j)
=\sum_q[\mathsf{f}^2]\tfrac12\operatorname{Tr}_{\square}(\mathcal{G}^{(\le k)}\mathsf{O}).
\end{equation}

The second equality in \eqref{10.6:eq-bare-trace-1} comes from the Gaussian tower
(Lemma~\ref{10.6:lem-gaussian-tower}). There the tower covariance $\mathcal{G}^{(\le k)}$ is the sum over
$j<k$ of $\mathsf{N}_{(j)}\mathcal{G}_P^{(j+1)}\mathsf{N}_{(j)}^{\dagger}$. We move the factors
$\mathsf{N}_{(j)}$ and $\mathsf{N}_{(j)}^{\dagger}$ around the trace per cell by
\eqref{10.6:eq-truncated-product-b}. The $j$-th term then pairs $\mathcal{G}_P^{(j+1)}$ with the
transported Hessian $\mathsf{O}_j$.

Next we replace $\mathcal{G}^{(\le k)}$ by $\mathcal{G}_B$. By the conditional-covariance clause of
Lemma~\ref{10.6:lem-gaussian-tower} (the second item of \eqref{10.6:eq-gaussian-tower-a}), taken at
$j=k$, $\mathcal{G}^{(\le k)}$ is a conditional covariance
for the data $(\mathsf{H}_0,D_0,\tilde{Q}_k,D_k)$. By Lemma~\ref{10.6:lem-representation-covariance},
$\mathcal{G}_B$ is a conditional covariance for $(\mathsf{H}_0,D_0,\tilde{Q},D_k)$. Here
$\tilde{Q}=\tilde{Q}_k$ is the $k$-fold composite, so the two sets of data agree.

By \eqref{10.6:eq-background-generators-a} of Lemma~\ref{10.6:lem-background-generators}, $P=\mathsf{O}$
satisfies $\mathsf{O}D_0=0$ and $D_0^{\dagger}\mathsf{O}=0$; this is the hypothesis of
Lemma~\ref{10.6:lem-trace-unique}, whose gauge map is here $D=D_0$, the gauge map of the data.
Lemma~\ref{10.6:lem-trace-unique} therefore gives
$\operatorname{Tr}_{\square}(\mathcal{G}^{(\le k)}\mathsf{O})=\operatorname{Tr}_{\square}(\mathcal{G}_B\mathsf{O})$.
Inserting this into \eqref{10.6:eq-bare-trace-1} and changing sign proves the first equality of
\eqref{10.6:eq-bare-trace-a}.

For the second equality we use the expression of Ba{\l}aban's conditional covariance $\mathcal{G}_B$
given at its definition in this section.
Read through the cyclicity of the trace per cell
\eqref{10.6:eq-truncated-product-b}, it expresses $\operatorname{Tr}_{\square}(\mathcal{G}_B\mathsf{O})$ as
$\operatorname{Tr}_{\square}(G\mathsf{O})
-\operatorname{Tr}_{\square}(\mathsf{M}\cdot QG\mathsf{O}GQ^{\dagger})$. Multiplying by $-\tfrac12$,
taking $[\mathsf{f}^2]$ and summing over $q$ gives the second line of \eqref{10.6:eq-bare-trace-a}.
\end{proof}

\begin{remark}
Identity \eqref{10.6:eq-bare-trace-a} holds at fixed $k$. It uses no estimate, no determinant, no
co-area density and no ghost: the first-order response reads the covariance only.

Equivalently, let $x$ be the parameter multiplying the action in the Boltzmann weight, and let
$\mathcal{G}$ be any conditional covariance for the data of
Lemma~\ref{10.6:lem-representation-covariance} (the value below is independent of the choice, by
Lemma~\ref{10.6:lem-trace-unique}). Then $x^{-1}\cdot\tfrac12\operatorname{Tr}_{\square}(\mathcal{G}\mathsf{O})$
is the Laplace coefficient of
the conditional expectation of $\mathsf{A}^{\mathbf{3}}$ given that the $k$-fold block average of $U$
equals $V$. This is a gauge-invariant quantity. The normaliser $z$ of \cite[(0.15)]{B-CMP109} does not
depend on the field.
\end{remark}

\begin{definition}[The background operator and admissible infrared families]
\label{10.6:def-infrared-families}
Let $\mathsf{H}_0$ be the total Hessian and
$\mathsf{O}=\sum_{\mu<\nu}(h_\mu+h_\nu)\mathsf{H}_{\mu\nu}$ the weighted Hessian on the constant
abelian configuration $\mathcal{U}_{\mathsf{f}}$, $h=(h_\mu)$ being the orientation weights, and let
$D_0$ be the bare lattice gradient, all as in
Lemma~\ref{10.6:lem-background-generators}. Let $\operatorname{Tr}_{\square}$ be the trace per unit
cell of Definition~\ref{10.6:def-truncated-product}.
\begin{enumerate}
\item The \emph{background operator} is $H := \mathsf{H}_0 + D_0 D_0^{\dagger} \in
\mathfrak{A}_{\mathsf{f}}(\mathcal{X}_0,\mathcal{X}_0)$ (in this section $H$ denotes this lattice
operator). Its flat part is
$H^{(0)} = d^{\dagger} d + d\, d^{\dagger} = -\Delta^{\eta}\otimes\mathbf{1}_4$, where $d$ is the
flat lattice exterior derivative and $\Delta^{\eta}$ the scalar lattice Laplacian at spacing
$\eta$. The symbol of $-\Delta^{\eta}$ is
\begin{equation}\label{10.6:eq-infrared-families-1}
\varepsilon(k) := \sum_{\mu}\frac{4}{\eta^2}\sin^2\Bigl(\frac{\eta k_\mu}{2}\Bigr),
\qquad \frac{4}{\pi^2}|k|^2 \le \varepsilon(k) \le |k|^2 \quad (k\in\mathrm{BZ}_\eta).
\end{equation}
\item $\mathsf{J} := (1+H^{(0)})^{1/2}$, and for an operator $v$ on $\mathcal{X}_0$ we set
$\|v\|_{\natural} := \|\mathsf{J}^{-1} v\, \mathsf{J}^{-1}\|$, with $\|\cdot\|$ the operator norm.
\item With $Q$, $R$ and $G$ as in the construction of Ba{\l}aban's conditional covariance
$\mathcal{G}_B$ (see Lemma~\ref{10.6:lem-representation-covariance}), we set
$\mathsf{W}_B := Q^{\dagger} Q - D_0(1-R)D_0^{\dagger}$; then $G = (H+\mathsf{W}_B)^{-1}$.
\item The \emph{trace norm per cell} is $\|t\|_{1,\square} := \operatorname{Tr}_{\square}|t|$. It
satisfies
\begin{equation}\label{10.6:eq-infrared-families-2}
|\operatorname{Tr}_{\square}(ts)| \le \|t\|_{1,\square}\,\|s\|,
\end{equation}
and $\|t\|_{1,\square} \le r\|t\|$ whenever $t$ factors through a space whose fibre in the
Bloch--Floquet decomposition over the dual of $\mathbb{L}_k$ has dimension $r$; this dimension is
$4\dim\mathfrak{g}_q$ for $\mathcal{X}_k$ and $\dim\mathfrak{g}_q$ for $\mathcal{C}_k$.
\item There is a constant $C_H$ such that, for $n\le 2$ and $|\alpha|\le 8$, uniformly in
$\eta\le 1$,
\begin{equation}\label{10.6:eq-infrared-families-a}
\|\delta^{\alpha} H^{(n)}\|_{\natural} \le C_H|\mathsf{f}|^n, \qquad
\|\delta^{\alpha}\mathsf{O}^{(n)}\|_{\natural} \le C_H|h|\,|\mathsf{f}|^n .
\end{equation}
\item An operator $\mathsf{W} = \mathsf{W}^{\dagger} \in \mathfrak{A}_{\mathsf{f}}(\mathcal{X}_0,\mathcal{X}_0)$
is an \emph{admissible infrared family} with constants $(\gamma,C_v,C_1)$ if
\begin{equation}\label{10.6:eq-infrared-families-b}
\begin{aligned}
&\text{(a1)}\quad H^{(0)} + \mathsf{W}^{(0)} \ge \gamma\,\mathsf{J}^2;\\
&\text{(a2)}\quad \|\delta^{\alpha}\mathsf{W}^{(n)}\|_{\natural} \le C_v|\mathsf{f}|^n
\qquad (n\le 2,\ |\alpha|\le 8);\\
&\text{(a3)}\quad \|\mathsf{J}^{-1}\mathsf{W}^{(n)}\mathsf{J}^{-1}\|_{1,\square} \le C_1|\mathsf{f}|^n
\qquad (n\le 2).
\end{aligned}
\end{equation}
For such $\mathsf{W}$ the \emph{regulated half-trace} is
$\tau[\mathsf{W}] := [\mathsf{f}^2]\,\tfrac12\operatorname{Tr}_{\square}\bigl((H+\mathsf{W})^{-1}\mathsf{O}\bigr)$.
\end{enumerate}
\end{definition}

\begin{proof}
We prove the two bounds in \eqref{10.6:eq-infrared-families-1}, the identity
$G=(H+\mathsf{W}_B)^{-1}$, the two trace-norm inequalities, and
\eqref{10.6:eq-infrared-families-a}.

\emph{The bounds on $\varepsilon$.} For $k\in\mathrm{BZ}_\eta$ each $\eta k_\mu/2$ lies in
$[-\pi/2,\pi/2]$, where
$(4/\pi^2)x^2 \le \sin^2 x \le x^2$. Hence
$\frac{4}{\eta^2}\sin^2(\eta k_\mu/2)$ lies between $\frac{4}{\pi^2}k_\mu^2$ and $k_\mu^2$, and
summing over $\mu$ gives the two inequalities in \eqref{10.6:eq-infrared-families-1}.

\emph{The identity for $G$.} By definition
$G = (\mathsf{H}_0 + D_0 R D_0^{\dagger} + Q^{\dagger}Q)^{-1}$, and
$\mathsf{H}_0 + D_0 R D_0^{\dagger} + Q^{\dagger}Q
= \mathsf{H}_0 + D_0D_0^{\dagger} + Q^{\dagger}Q - D_0(1-R)D_0^{\dagger} = H + \mathsf{W}_B$.

\emph{The trace-norm inequalities.} At fixed $\eta$ the Bloch--Floquet decomposition over the
dual of $\mathbb{L}_k$
writes $t$ as a field of operators $t(\xi)$ on fibres of finite dimension, $\xi$ running over the
dual torus with normalised measure $\bar{d}\xi$; $|t|$ corresponds to $|t(\xi)|$, and
$\|s\| = \sup_\xi \|s(\xi)\|$. The trace per cell is the integral of the fibre traces,
$\operatorname{Tr}_{\square} t = \int \bar{d}\xi\,\operatorname{Tr} t(\xi)$, where
$\operatorname{Tr}$ is the (unnormalised) trace on the fibre. For matrices,
$|\operatorname{Tr}(t(\xi)s(\xi))| \le \operatorname{Tr}|t(\xi)|\,\|s(\xi)\| \le
\operatorname{Tr}|t(\xi)|\,\|s\|$; integrating over $\xi$ gives
\eqref{10.6:eq-infrared-families-2}. If $t$ factors through a space whose fibre has dimension
$r$, each $t(\xi)$ has rank at most $r$, so
$\operatorname{Tr}|t(\xi)| \le r\|t(\xi)\| \le r\|t\|$, and integration gives
$\|t\|_{1,\square}\le r\|t\|$.

\emph{The bounds \eqref{10.6:eq-infrared-families-a}.} By
Lemma~\ref{10.6:lem-background-generators}, the reduced symbols of $\mathsf{H}_{\mu\nu}$, hence of
$\mathsf{H}_0$ and $\mathsf{O}$, are of the
scaling form $\eta^{-2}\mathsf{t}_1(\eta k;\eta^2\mathsf{f})$, with $\mathsf{t}_1$ a trigonometric
polynomial in its first argument whose coefficients are entire in the second (for $\mathsf{O}$ the
Hessians $\mathsf{H}_{\mu\nu}$ enter with the coefficients $h_\mu+h_\nu$); the same holds for
$D_0D_0^{\dagger}$, whose symbol, built from $d_\mu(k) = (2i/\eta)\sin(\eta k_\mu/2)$, is
$\eta^{-2}$ times a trigonometric polynomial in $\eta k$, free of $\mathsf{f}$. Extracting the
degree-$n$ part in $\mathsf{f}$ and applying $\delta^{\alpha}$, the symbols of
$\delta^{\alpha}H^{(n)}$ and $\delta^{\alpha}\mathsf{O}^{(n)}$ are
\begin{equation}\label{10.6:eq-infrared-families-3}
i^{|\alpha|}\,\eta^{2n-2+|\alpha|}\,
(\partial^{\alpha}_{\kappa}\partial^{n}_{\phi}\mathsf{t}_1)(\eta k;0)[\mathsf{f}^n].
\end{equation}
Here $\partial_{\phi}$ denotes differentiation of $\mathsf{t}_1$ in its second argument.
Moreover $\mathsf{t}_1(\kappa;0) = O(|\kappa|^2)$ and $\partial\mathsf{t}_1(\kappa;0) = O(|\kappa|)$.
Hence for $(n,|\alpha|) = (0,0)$ and $(0,1)$, where the power of $\eta$ is $-2$ and $-1$
respectively, the
symbols \eqref{10.6:eq-infrared-families-3} are bounded by $C|k|^2$ and $C|k|$ respectively. In all
other cases with $n\le 2$, $|\alpha|\le 8$ we have $2n-2+|\alpha|\ge 0$, so
$\eta^{2n-2+|\alpha|}\le 1$ for $\eta\le 1$, and the function
$(\partial^{\alpha}_{\kappa}\partial^{n}_{\phi}\mathsf{t}_1)(\cdot\,;0)$, being a trigonometric
polynomial, is bounded; these symbols are bounded by $C|\mathsf{f}|^n$.

Since $H^{(0)} = -\Delta^{\eta}\otimes\mathbf{1}_4$ acts by the scalar symbol $\varepsilon(k)$, the
operator $\mathsf{J}^{-1}v\,\mathsf{J}^{-1}$ has symbol $(1+\varepsilon(k))^{-1}\hat v(k)$, where
$\hat v(k)$ is the symbol of $v$, and $\|v\|_{\natural}$ is the supremum over $\mathrm{BZ}_\eta$ of
the norm of this symbol. By \eqref{10.6:eq-infrared-families-1},
$(1+\varepsilon(k))^{-1} \le (1+\frac{4}{\pi^2}|k|^2)^{-1} \le 1$. Therefore the case $(0,0)$
contributes at most $C|k|^2(1+\frac{4}{\pi^2}|k|^2)^{-1} \le \frac{\pi^2}{4}C$, the case $(0,1)$
at most $C|k|(1+\frac{4}{\pi^2}|k|^2)^{-1} \le \frac{\pi}{4}C$, and every other case at most
$C|\mathsf{f}|^n$. This gives \eqref{10.6:eq-infrared-families-a}, the factor $|h|$ in the second
bound coming from the coefficients $h_\mu+h_\nu$ of $\mathsf{O}$, uniformly in $\eta\le 1$.
\end{proof}

\begin{lemma}[Change of infrared regulator at bounded cost]\label{10.6:lem-regulator-change}
Let $\mathsf{W}_0$ and $\mathsf{W}_1$ be admissible infrared families with the same constants
$(\gamma,C_v,C_1)$ in the sense of Definition~\ref{10.6:def-infrared-families}. Then
\begin{equation*}
|\tau[\mathsf{W}_1]-\tau[\mathsf{W}_0]|\le C_{\mathrm{local}}\,|h|\,|\mathsf{f}|^2 ,
\end{equation*}
where $|h|$ is the factor in the bound \eqref{10.6:eq-infrared-families-a} on the derivatives
$\delta^{\alpha}\mathsf{O}^{(n)}$, and the constant $C_{\mathrm{local}}$ depends on
$(\gamma,C_v,C_1,C_H)$ only; in particular it depends neither on $\eta$ nor on $k$.
\end{lemma}

\begin{proof}
\emph{Interpolation.} For $s\in[0,1]$ put $\mathsf{W}_s:=(1-s)\mathsf{W}_0+s\mathsf{W}_1$ and
$\dot{\mathsf{W}}:=\mathsf{W}_1-\mathsf{W}_0$. The family $\mathsf{W}_s$ is admissible with the same
constants $(\gamma,C_v,C_1)$. Indeed $\mathsf{W}_s=\mathsf{W}_s^{\dagger}$; for condition (a1) of
\eqref{10.6:eq-infrared-families-b},
\begin{equation*}
H^{(0)}+\mathsf{W}_s^{(0)}=(1-s)\bigl(H^{(0)}+\mathsf{W}_0^{(0)}\bigr)
+s\bigl(H^{(0)}+\mathsf{W}_1^{(0)}\bigr)\ge\gamma\mathsf{J}^2 ;
\end{equation*}
for conditions (a2) and (a3) of the same display, the degree components and the derivations
$\delta^{\alpha}$ are linear, so
$\delta^{\alpha}\mathsf{W}_s^{(n)}=(1-s)\delta^{\alpha}\mathsf{W}_0^{(n)}+s\,\delta^{\alpha}
\mathsf{W}_1^{(n)}$, and the triangle inequality for $\|\cdot\|_{\natural}$ and for
$\|\cdot\|_{1,\square}$ gives the bounds of \eqref{10.6:eq-infrared-families-b} with the same
constants. Let $G_s:=(H+\mathsf{W}_s)^{-1}$ be the $\star$-inverse
(Definition~\ref{10.6:def-truncated-product}) and $g:=G_s^{(0)}$. By
\eqref{10.6:eq-truncated-product-a} the degree-zero part of a truncated product is the ordinary
product of the degree-zero parts, so $g=(H^{(0)}+\mathsf{W}_s^{(0)})^{-1}$. Conjugating (a1) by
$\mathsf{J}^{-1}$ gives $\mathsf{J}^{-1}(H^{(0)}+\mathsf{W}_s^{(0)})\mathsf{J}^{-1}\ge\gamma$, whose
inverse is $\mathsf{J}g\mathsf{J}$; hence
\begin{equation*}
0\le\mathsf{J}g\mathsf{J}\le\gamma^{-1}.
\end{equation*}
Differentiating $(H+\mathsf{W}_s)\star G_s=1$ in $s$ gives
$\dot{\mathsf{W}}\star G_s+(H+\mathsf{W}_s)\star\partial_sG_s=0$; multiplying on the left by $G_s$
with the truncated product and using its associativity, we get
$\partial_sG_s=-G_s\star\dot{\mathsf{W}}\star G_s$. Since $[\mathsf{f}^2]$ and
$\operatorname{Tr}_{\square}$ are linear,
\begin{align*}
\tau[\mathsf{W}_1]-\tau[\mathsf{W}_0]
&=\int_0^1\partial_s\tau[\mathsf{W}_s]\,ds\\
&=-\int_0^1[\mathsf{f}^2]\tfrac12\operatorname{Tr}_{\square}
\bigl(G_s\star\dot{\mathsf{W}}\star G_s\star\mathsf{O}\bigr)\,ds .
\end{align*}
We apply \eqref{10.6:eq-truncated-product-b} twice. First, since by that identity the trace per
cell of a truncated product is the trace per cell of the ordinary products of the degree
components of the factors, and $\operatorname{Tr}_{\square}$ is cyclic for ordinary products,
being the integral of the traces on the finite-dimensional fibres of the Bloch--Floquet
decomposition, we have
$\operatorname{Tr}_{\square}(t\star t')=\operatorname{Tr}_{\square}(t'\star t)$ for
$t,t'\in\mathfrak{A}_{\mathsf{f}}$, and with associativity this brings the left factor $G_s$
round to the right. The second
application writes the trace of the truncated product of $\dot{\mathsf{W}}$ with
$G_s\star\mathsf{O}\star G_s$ through the degree components of the two factors. This gives
\begin{equation}\label{10.6:eq-regulator-change-1}
\begin{split}
\tau[\mathsf{W}_1]-\tau[\mathsf{W}_0]
&=-\int_0^1 ds\,[\mathsf{f}^2]\tfrac12\operatorname{Tr}_{\square}
\bigl(\dot{\mathsf{W}}\star(G_s\star\mathsf{O}\star G_s)\bigr)\\
&=-\int_0^1 ds\,\tfrac12\sum_{n+m=2}\operatorname{Tr}_{\square}
\bigl(\dot{\mathsf{W}}^{(n)}\cdot(G_s\star\mathsf{O}\star G_s)^{(m)}\bigr).
\end{split}
\end{equation}

\emph{Words.} Let $\mathcal{V}$ be the set of the operators $\delta^{\alpha}H^{(n)}$,
$\delta^{\alpha}\mathsf{W}_s^{(n)}$ and $\delta^{\alpha}\mathsf{O}^{(n)}$, where
$(n,\alpha)\neq(0,0)$ is required for the first two kinds. We claim that for $n\le2$ and
$|\alpha|\le4-2n$ every $\delta^{\alpha}G_s^{(n)}$ is a sum of at most $c_*$ terms, $c_*$ an
absolute number, each of the form
\begin{equation*}
(\text{numerical coefficient})\cdot\mathsf{f}\cdots\mathsf{f}\cdot
g\,v_1\,g\,v_2\cdots v_r\,g,\qquad v_i\in\mathcal{V},\ r\le8,
\end{equation*}
where each $v_i$ carries a derivation $\delta^{\beta}$ of order $|\beta|\le8$, the dots
$\mathsf{f}\cdots\mathsf{f}$
denote a product of entries of $\mathsf{f}$, and the total degree in $\mathsf{f}$ of each term
is $n$. The proof is by induction on $n$ and $|\alpha|$. For $n=0$, $G_s^{(0)}=g$, and the rule
$\delta g=-g\,\bigl(\delta(H^{(0)}+\mathsf{W}_s^{(0)})\bigr)\,g$, applied repeatedly together
with the Leibniz rule for the derivations $\delta_{\kappa}$, expresses $\delta^{\alpha}g$ as such
a sum with vertices $\delta^{\beta}H^{(0)}$, $\delta^{\beta}\mathsf{W}_s^{(0)}$, $\beta\neq0$,
which lie in $\mathcal{V}$. For $n\ge1$, the component of degree $n$ of the identity
$(H+\mathsf{W}_s)\star G_s=1$, written out with \eqref{10.6:eq-truncated-product-a}, expresses
$(H^{(0)}+\mathsf{W}_s^{(0)})G_s^{(n)}$ as a combination, with numerical coefficients and entries
of $\mathsf{f}$, of products of derivatives of degree components of $H+\mathsf{W}_s$ with
derivatives of $G_s^{(m)}$, $m<n$, the degrees adding up to $n$; multiplying on the left by $g$
and differentiating with the Leibniz rule and the rule for $\delta g$ gives the claim at
$(n,\alpha)$ from the induction hypothesis. Consequently, expanding
$(G_s\star\mathsf{O}\star G_s)^{(m)}$ with \eqref{10.6:eq-truncated-product-a}, which places at
most two derivations on each factor of each truncated product, one obtains a finite sum of words
$g\,v_1\,g\cdots v_r\,g$ (with numerical coefficients and entries of $\mathsf{f}$) containing
exactly one vertex of $\mathsf{O}$-type; the number of these terms, their coefficients and the
lengths $r$ are bounded by absolute numbers, since they arise from the expansions of
$\delta^{\alpha}G_s^{(n)}$ just described and from these two truncated products.

\emph{Estimate of one word.} For each such word we insert $\mathsf{J}^{-1}\mathsf{J}=1$ after
$\dot{\mathsf{W}}^{(n)}$ and after each $v_i$, and $\mathsf{J}\mathsf{J}^{-1}=1$ before each $v_i$;
we then multiply the word by $\mathsf{J}^{-1}$ on the left and by $\mathsf{J}$ on the right, which
does not change $\operatorname{Tr}_{\square}$ by its cyclicity on the finite-dimensional fibres.
Finally we use the bound $|\operatorname{Tr}_{\square}(tt')|
\le\|t\|_{1,\square}\|t'\|$ of Definition~\ref{10.6:def-infrared-families} with
$t=\mathsf{J}^{-1}\dot{\mathsf{W}}^{(n)}\mathsf{J}^{-1}$:
\begin{equation}\label{10.6:eq-regulator-change-2}
\begin{split}
\bigl|\operatorname{Tr}_{\square}\bigl(\dot{\mathsf{W}}^{(n)}g\,v_1\,g\cdots v_r\,g\bigr)\bigr|
&=\bigl|\operatorname{Tr}_{\square}\bigl((\mathsf{J}^{-1}\dot{\mathsf{W}}^{(n)}\mathsf{J}^{-1})
(\mathsf{J}g\mathsf{J})(\mathsf{J}^{-1}v_1\mathsf{J}^{-1})(\mathsf{J}g\mathsf{J})\cdots
(\mathsf{J}g\mathsf{J})\bigr)\bigr|\\
&\le2C_1|\mathsf{f}|^n\cdot\gamma^{-(r+1)}\prod_{i=1}^{r}\|v_i\|_{\natural}.
\end{split}
\end{equation}
Here $\|\mathsf{J}^{-1}\dot{\mathsf{W}}^{(n)}\mathsf{J}^{-1}\|_{1,\square}\le2C_1|\mathsf{f}|^n$
by (a3) for $\mathsf{W}_1$ and $\mathsf{W}_0$ and the triangle inequality; the operator norm of
the remaining product is at most the product of the norms, with
$\|\mathsf{J}g\mathsf{J}\|\le\gamma^{-1}$ for each of the $r+1$ factors and
$\|\mathsf{J}^{-1}v_i\mathsf{J}^{-1}\|=\|v_i\|_{\natural}$. Each vertex has derivatives of order
at most eight, so \eqref{10.6:eq-infrared-families-a} applies to the vertices of $H$- and
$\mathsf{O}$-type, uniformly in $\eta\le1$, and (a2), for the admissible family $\mathsf{W}_s$,
to those of $\mathsf{W}_s$-type. Since exactly one vertex is of $\mathsf{O}$-type, and the degrees
of the vertices together with the explicit entries of $\mathsf{f}$ make up the degree $m$, the
product of the vertex norms with the explicit factors of $\mathsf{f}$ is at most
$\max\{C_H,C_v\}^r\,|h|\,|\mathsf{f}|^m$.

\emph{Conclusion.} Inserting these bounds into \eqref{10.6:eq-regulator-change-1}, using
$|\mathsf{f}|^n\,|\mathsf{f}|^m=|\mathsf{f}|^2$ for $n+m=2$, the absolute bounds on the number of
words, their coefficients and their lengths $r$, and the factor $\tfrac12$ and the integral over
$s\in[0,1]$, we obtain
$|\tau[\mathsf{W}_1]-\tau[\mathsf{W}_0]|\le C_{\mathrm{local}}\,|h|\,|\mathsf{f}|^2$
with $C_{\mathrm{local}}$ built from $C_1$, $\gamma^{-1}$ and $\max\{C_H,C_v\}$ and absolute numbers
only; it depends neither on $\eta$ nor on $k$.

No power counting enters the proof: one factor of each word is of trace class per cell, and every
other factor is bounded after the sandwich $\mathsf{J}^{-1}\cdot\mathsf{J}^{-1}$. The logarithmic
part of $\tau[\mathsf{W}]$ itself sits in the traces of the words $g\,v_1\,g\cdots v_r\,g$ that
contain no vertex of $\mathsf{W}$-type; these cancel in the difference, in which every word
carries the factor $\dot{\mathsf{W}}$.
\end{proof}

\begin{proposition}[Admissible regulator and bounded block-mass term]
\label{10.6:prop-regulator-admissible}
\begin{enumerate}
\item[(i)] Ba{\l}aban's infrared term $\mathsf{W}_B$ is an admissible infrared family in the sense of
Definition~\ref{10.6:def-infrared-families}, with $\gamma=\gamma_0$, the constant of
\cite[Proposition~1.1, (1.90)]{B-CMP95}, and with constants $C_v$, $C_1$ depending only on $L$ and
on the gauge group.
\item[(ii)] There is a constant $C_S$, depending only on $L$ and on the gauge group, such that,
uniformly in $k$,
\begin{equation}\label{10.6:eq-regulator-admissible-1}
\Bigl|[\mathsf{f}^2]\,\tfrac12\operatorname{Tr}_{\square}\bigl(\mathsf{M}\cdot QG\mathsf{O}GQ^\dagger\bigr)\Bigr|
\le C_S\,|h|\,|\mathsf{f}|^2 .
\end{equation}
\end{enumerate}
\end{proposition}

\begin{proof}
We work with the operators of Ba{\l}aban's construction on the constant configuration
$\mathcal{U}_{\mathsf{f}}$ at $a=1$:
\begin{align*}
\Delta_s&:=D_0^\dagger D_0, & G'&:=(\Delta_s+Q'^\dagger Q')^{-1}, & M'&:=(Q'G'G'Q'^\dagger)^{-1},\\
R&:=1-G'Q'^\dagger M'Q'G', & G&:=(\mathsf{H}_0+D_0RD_0^\dagger+Q^\dagger Q)^{-1},
& \mathsf{M}&:=(QGQ^\dagger)^{-1},
\end{align*}
products and inverses being taken in $\mathfrak{A}_{\mathsf{f}}$, that is, with the truncated
product $\star$ of \eqref{10.6:eq-truncated-product-a}. Since $H=\mathsf{H}_0+D_0D_0^\dagger$,
Ba{\l}aban's infrared term is
\begin{equation}\label{10.6:eq-regulator-admissible-2}
\mathsf{W}_B=Q^\dagger\star Q-D_0\star(1-R)\star D_0^\dagger,\qquad H+\mathsf{W}_B=G^{-1}.
\end{equation}
By the rule $(t\star s)^\dagger=s^\dagger\star t^\dagger$ of
Definition~\ref{10.6:def-truncated-product}, $G'$, $M'$ and $R$ are self-adjoint, hence so is
$\mathsf{W}_B$. The letter $C$ denotes constants depending only on $L$ and the gauge group.

Two facts are used repeatedly. First, $\delta_\kappa$ obeys the Leibniz rule for composition, and
each term of \eqref{10.6:eq-truncated-product-a} is a composition of $\delta$-derivatives (at most two
on each factor) multiplied by a monomial in $\mathsf{f}$. Hence $\delta^\alpha(t\star s)^{(n)}$ is a
finite sum of constant multiples of
$\mathsf{f}$-monomials of degree $m$ times $(\delta^{\beta}t^{(n_1)})(\delta^{\beta'}s^{(n_2)})$, with
$m+n_1+n_2=n$; if the coefficients of $t$ and $s$ obey bounds $C|\mathsf{f}|^{n}$, so do those of
$t\star s$. Second, the fibres of $\mathcal{X}_k$ and $\mathcal{C}_k$ over a unit cell of
$\mathbb{L}_k$ have dimensions $4\dim\mathfrak{g}_q$ and $\dim\mathfrak{g}_q$; an operator that
factors through one of these spaces as $AB$ has trace norm per cell at most that dimension times
$\|A\|\,\|B\|$. Moreover $H^{(0)}\ge0$, so $\mathsf{J}\ge1$ and $\|\mathsf{J}^{-1}\|\le1$.

\emph{(i), condition (a1).} At $\mathsf{f}=0$ the configuration $\mathcal{U}_{\mathsf{f}}$ is the unit
configuration, and $Q^{(0)}$, $Q'^{(0)}$, $R^{(0)}$ are Ba{\l}aban's $Q_k$, $Q'_k$, $R$ of
\cite{B-CMP95}. By \eqref{10.6:eq-regulator-admissible-2},
\[
H^{(0)}+\mathsf{W}_B^{(0)}=\mathsf{H}_0^{(0)}+D_0R^{(0)}D_0^\dagger+Q^{(0)\dagger}Q^{(0)},
\]
which is Ba{\l}aban's operator $\Delta_a$ at $a=1$ and trivial background. By
\cite[Proposition~1.1, (1.89)--(1.90)]{B-CMP95}, $\Delta_a=G^{-1}\ge\gamma_0(-\Delta+I)$ with
$\gamma_0>0$ independent of $k$ and of the torus and depending only on the dimension. Here
$-\Delta=H^{(0)}$, so $H^{(0)}+\mathsf{W}_B^{(0)}\ge\gamma_0(1+H^{(0)})=\gamma_0\mathsf{J}^2$, which
is \eqref{10.6:eq-infrared-families-b}(a1) with $\gamma=\gamma_0$.

\emph{The block mass.} By the first fact, $\delta^\alpha(Q^\dagger\star Q)^{(n)}$ is a finite sum of
constant multiples of $\mathsf{f}$-monomials of degree $n-n'-n''$ times
$(\delta^{\alpha'}Q^{(n')})^\dagger(\delta^{\alpha''}Q^{(n'')})$. By
Lemma~\ref{10.6:lem-uniform-majorant}, \eqref{10.6:eq-uniform-majorant-a} and its Cauchy
consequence, $\|\delta^{\alpha'}Q^{(n')}\|\le2^{|\alpha'|}C_Qc_L^{-n'}|\mathsf{f}|^{n'}$, and likewise
for $n''$; so each term is bounded by $C|\mathsf{f}|^n$. Each term is a composite
$\mathcal{X}_0\to\mathcal{X}_k\to\mathcal{X}_0$, hence factors through $\mathcal{X}_k$. Since
$\|\mathsf{J}^{-1}\|\le1$, multiplying by $\mathsf{J}^{-1}$ on either side does not increase these
bounds, and the second fact gives the trace-norm bound. So the block mass obeys (a2) and (a3).

\emph{The gauge projection.} By Lemma~\ref{10.6:lem-background-generators}, $D_0$ has an
$\mathsf{f}$-free reduced kernel with symbol $d_\mu(k)$, so $D_0^\dagger\star D_0=\sum_\mu
\bar d_\mu\star d_\mu$. Since $d_\mu$ depends on $k_\mu$ only, $\delta_\kappa d_\mu=0$ for
$\kappa\ne\mu$, so every correction term of \eqref{10.6:eq-truncated-product-a} carries a factor
$\mathsf{f}_{\mu\mu}=0$. Hence $D_0^\dagger\star D_0=\sum_\mu|d_\mu|^2=\varepsilon(k)$ exactly, and
$\Delta_s$ is $\mathsf{f}$-free.

Next, $Q'^{(0)}$ is the $k$-fold plain block mean, so $Q'^{(0)\dagger}Q'^{(0)}\lambda$ equals, on each
unit cube, the average of $\lambda$ over that cube. The Poincar\'e inequality on the unit cubes gives
$\|\lambda-Q'^{(0)\dagger}Q'^{(0)}\lambda\|^2\le C_{\mathrm P}\langle\lambda,\Delta_s\lambda\rangle$
with $C_{\mathrm P}$ uniform in $\eta$. Hence
\[
\|\lambda\|^2\le2\|\lambda-Q'^{(0)\dagger}Q'^{(0)}\lambda\|^2+2\|Q'^{(0)}\lambda\|^2
\le\max\{2C_{\mathrm P},2\}\,\langle\lambda,(\Delta_s+Q'^{(0)\dagger}Q'^{(0)})\lambda\rangle .
\]
Adding $\Delta_s\le\Delta_s+Q'^{(0)\dagger}Q'^{(0)}$ yields
\begin{equation}\label{10.6:eq-regulator-admissible-3}
\Delta_s+Q'^{(0)\dagger}Q'^{(0)}\ge\gamma'(1+\Delta_s),\qquad
\gamma':=\tfrac12\min\{1,(\max\{2C_{\mathrm P},2\})^{-1}\}.
\end{equation}
With $\mathsf{J}_s:=(1+\Delta_s)^{1/2}$, \eqref{10.6:eq-regulator-admissible-3} gives
$\|\mathsf{J}_sG'^{(0)}\mathsf{J}_s\|\le1/\gamma'$. By the inversion rule of
Definition~\ref{10.6:def-truncated-product}, the coefficients $\delta^\alpha G'^{(n)}$ are finite sums
of words in $G'^{(0)}$, $\delta^\alpha\varepsilon$ and $\delta^\alpha(Q'^\dagger\star Q')^{(n)}$, in
which the factors other than $G'^{(0)}$ stand between factors $G'^{(0)}$. For $|\alpha|\ge1$ the
symbol of $\delta^\alpha\varepsilon$ is at most $C(1+|k|)\le C(1+\varepsilon(k))$ on
$\mathrm{BZ}_\eta$, so $\|\mathsf{J}_s^{-1}(\delta^\alpha\varepsilon)\mathsf{J}_s^{-1}\|\le C$. The
coefficients $\delta^\alpha(Q'^\dagger\star Q')^{(n)}$ are bounded by $C|\mathsf{f}|^n$, by the first
fact and the bound on $\delta^\alpha Q'^{(n)}$ in Lemma~\ref{10.6:lem-uniform-majorant}. Inserting
$\mathsf{J}_s\mathsf{J}_s^{-1}$ beside every factor $G'^{(0)}$, and using
$\|G'^{(0)}\mathsf{J}_s\|,\|\mathsf{J}_sG'^{(0)}\|\le1/\gamma'$, shows that $\delta^\alpha G'^{(n)}$
is a bounded operator on $\mathcal{C}_0$ of norm at most $C|\mathsf{f}|^n$.

By \cite[(1.42)--(1.45)]{B-CMP95}, $\gamma_0\le Q'_kG'^{2}_kQ'^{*}_k$, so $M'^{(0)}\le\gamma_0^{-1}$.
The coefficients $\delta^\alpha M'^{(n)}$ are words in $M'^{(0)}$ and in coefficients of
$Q'\star G'\star G'\star Q'^\dagger$, and are bounded likewise. Hence, by the first fact,
$\delta^\alpha(1-R)^{(n)}=\delta^\alpha(G'Q'^\dagger M'Q'G')^{(n)}$ is bounded on $\mathcal{C}_0$ by
$C|\mathsf{f}|^n$; each of its terms contains the factor $Q'$ followed later by $Q'^\dagger$, so it
factors through $\mathcal{C}_k$.

Finally, $\|\mathsf{J}^{-1}\delta^\beta D_0\|\le C$. For $\beta=0$ the symbols are
$d_\mu(k)(1+\varepsilon(k))^{-1/2}$, of modulus at most one. For $|\beta|\ge1$, $\delta^\beta$ acts
on symbols as a derivative in $k$, and
$d_\mu(k)=(2\mathrm{i}/\eta)\sin(\eta k_\mu/2)$; so $\delta^\beta d_\mu$ is $\eta^{|\beta|-1}$ times a
bounded function. By the first fact, $\delta^\alpha(D_0\star(1-R)\star D_0^\dagger)^{(n)}$ is a finite
sum of $\mathsf{f}$-monomials times $(\delta^\beta D_0)(\delta^{\alpha'}(1-R)^{(n')})
(\delta^{\beta'}D_0)^\dagger$. Sandwiched between factors $\mathsf{J}^{-1}$, each term is bounded by
$C|\mathsf{f}|^n$ and factors through $\mathcal{C}_k$. By the second fact, (a2) and (a3) follow. With
the block mass and \eqref{10.6:eq-regulator-admissible-2}, this proves (i).

\emph{(ii).} The operator $\mathsf{M}\cdot QG\mathsf{O}GQ^\dagger$ acts on $\mathcal{X}_k$, so the
trace per cell is over $\mathcal{X}_k$, of fibre dimension $4\dim\mathfrak{g}_q$. By
\cite[(1.99)--(1.101)]{B-CMP95}, $QGQ^{*}$ is bounded from below by a positive constant, so
$\mathsf{M}^{(0)}$ is bounded. By the inversion rule, the coefficients of $\mathsf{M}$ are words in
$\mathsf{M}^{(0)}$ and $\delta^\alpha(Q\star G\star Q^\dagger)^{(n)}$.

Each occurrence of $G$ between factors $Q$ or $\mathsf{O}$ is expanded in words as in the proof of
the change-of-regulator lemma (Lemma~\ref{10.6:lem-regulator-change}): $g:=G^{(0)}$ alternates
with $\delta$-derivatives of coefficients of $H$, of $\mathsf{W}_B$ and of $\mathsf{O}$. Inserting
$\mathsf{J}\mathsf{J}^{-1}$ beside every $g$, the word is bounded using the following estimates:
\begin{itemize}
\item $\|(\delta^\alpha Q^{(n)})\mathsf{J}^{-1}\|\le\|\delta^\alpha Q^{(n)}\|$, bounded by
Lemma~\ref{10.6:lem-uniform-majorant};
\item $\|\mathsf{J}g\mathsf{J}\|\le\gamma_0^{-1}$, by (a1) in (i);
\item the bounds \eqref{10.6:eq-infrared-families-a} for the coefficients of $H$ and $\mathsf{O}$;
\item the bounds (a2) of \eqref{10.6:eq-infrared-families-b} for $\mathsf{W}_B$, from (i).
\end{itemize}
The factor $\mathsf{O}$ occurs exactly once, so the degree-two coefficient is an operator on
$\mathcal{X}_k$ of norm at most $C|h||\mathsf{f}|^2$. Its trace per cell is at most
$4\dim\mathfrak{g}_q$ times this. None of the constants depends on $k$, which gives
\eqref{10.6:eq-regulator-admissible-1}.
\end{proof}

\begin{lemma}[Second-order magnetic functional calculus on the lattice]
\label{10.6:lem-functional-calculus}
Let $H\in\mathfrak{A}_{\mathsf{f}}(\mathcal{X}_0,\mathcal{X}_0)$ be the background Feynman operator of
Definition~\ref{10.6:def-infrared-families}, with parts $H^{(n)}$ of degree $n$ in $\mathsf{f}$, and
with part of degree zero the scalar $H^{(0)}=\varepsilon\mathbf{1}$. Work in the Fourier
representation, write $\partial_\kappa$ for the derivative in $k_\kappa$, and let subscripts on
$\varepsilon$ denote $k$-derivatives:
\[
\varepsilon_\kappa=\tfrac{2}{\eta}\sin(\eta k_\kappa),\qquad
\varepsilon_{\kappa\kappa'}=2\delta_{\kappa\kappa'}\cos(\eta k_\kappa),
\]
repeated indices being summed. A function of $\varepsilon$ stands for the element of
$\mathfrak{A}_{\mathsf{f}}$ of degree zero whose symbol is that function of $\varepsilon(k)$ times
$\mathbf{1}$. For $\varphi\in C^\infty([0,\infty[)$ put
\begin{equation}\label{10.6:eq-functional-calculus-a}
\begin{split}
\varphi(H)_{\star}:={}&\varphi(\varepsilon)+\varphi'(\varepsilon)\bigl(H^{(1)}+H^{(2)}\bigr)
+\tfrac12\varphi''(\varepsilon)\bigl(H^{(1)}\bigr)^2\\
&-\tfrac{q^2}{8}\,\mathsf{f}_{\kappa\lambda}\mathsf{f}_{\kappa'\lambda'}\varepsilon_{\kappa\kappa'}
\Bigl[\tfrac13\varphi'''(\varepsilon)\varepsilon_\lambda\varepsilon_{\lambda'}
+\tfrac12\varphi''(\varepsilon)\varepsilon_{\lambda\lambda'}\Bigr].
\end{split}
\end{equation}
Then:
\begin{enumerate}
\item[(i)] for $\varphi(\lambda)=e^{-\sigma\lambda}$, the element
$U(\sigma):=\varphi(H)_{\star}$ is the solution in $\mathfrak{A}_{\mathsf{f}}$ of
$\partial_\sigma U=-H\star U$, $U(0)=\mathbf{1}$; explicitly
\begin{equation}\label{10.6:eq-functional-calculus-1}
U(\sigma)=e^{-\sigma\varepsilon}\Bigl[\mathbf{1}-\sigma\bigl(H^{(1)}+H^{(2)}\bigr)
+\tfrac{\sigma^2}{2}\bigl(H^{(1)}\bigr)^2+\mathsf{a}_2(\sigma)\Bigr],
\end{equation}
\begin{equation}\label{10.6:eq-functional-calculus-2}
\mathsf{a}_2(\sigma):=\tfrac{q^2}{8}\,\mathsf{f}_{\kappa\lambda}\mathsf{f}_{\kappa'\lambda'}
\varepsilon_{\kappa\kappa'}\Bigl[\tfrac{\sigma^3}{3}\varepsilon_\lambda\varepsilon_{\lambda'}
-\tfrac{\sigma^2}{2}\varepsilon_{\lambda\lambda'}\Bigr];
\end{equation}
\item[(ii)] $\varphi_1(H)_{\star}\star\varphi_2(H)_{\star}=(\varphi_1\varphi_2)(H)_{\star}$ for
$\varphi_1,\varphi_2\in C^\infty([0,\infty[)$, and $\lambda(H)_{\star}=H$, $\lambda$ denoting here
the identity function on $[0,\infty[$.
\end{enumerate}
\end{lemma}

\begin{proof}
By Definition~\ref{10.6:def-truncated-product}, in Fourier variables (where $\delta_\kappa$ acts as
$i\partial_\kappa$) the truncated product of \eqref{10.6:eq-truncated-product-a} reads
\[
t\star s=ts-\tfrac{iq}{2}\mathsf{f}_{\kappa\lambda}\,\partial_\kappa t\,\partial_\lambda s
-\tfrac{q^2}{8}\mathsf{f}_{\kappa\lambda}\mathsf{f}_{\kappa'\lambda'}\,
\partial_\kappa\partial_{\kappa'}t\,\partial_\lambda\partial_{\lambda'}s
\pmod{(\mathsf{f})^3}.
\]
Each $\mathsf{f}$ carries one degree, so the part of degree $n$ of $t\star s$ is
$\sum_{a+b=n}t^{(a)}s^{(b)}$, plus $-\frac{iq}{2}\mathsf{f}_{\kappa\lambda}\sum_{a+b=n-1}
\partial_\kappa t^{(a)}\partial_\lambda s^{(b)}$ for $n\ge1$, plus
$-\frac{q^2}{8}\mathsf{f}_{\kappa\lambda}\mathsf{f}_{\kappa'\lambda'}\partial_\kappa\partial_{\kappa'}
t^{(0)}\partial_\lambda\partial_{\lambda'}s^{(0)}$ for $n=2$. In particular
$\mathbf{1}\star s=s=s\star\mathbf{1}$. We use that $\mathsf{f}$ is antisymmetric, that
$\partial_\kappa H^{(0)}=\varepsilon_\kappa\mathbf{1}$, and that $H^{(n)}$ does not depend on $\sigma$.

(i) Let $U$ solve $\partial_\sigma U=-H\star U$, $U(0)=\mathbf{1}$. In degree zero the equation reads
$\partial_\sigma U^{(0)}=-\varepsilon U^{(0)}$, $U^{(0)}(0)=\mathbf{1}$, so $U^{(0)}=E:=e^{-\sigma
\varepsilon}\mathbf{1}$, a scalar. Since $E$ is invertible we write $U=E(\mathbf{1}+u_1+u_2)$ with
$u_n$ of degree $n$ and $u_n(0)=0$; then $\partial_\sigma U=-\varepsilon U+E(u_1'+u_2')$, the prime
denoting $\partial_\sigma$, and $\partial_\lambda E=-\sigma\varepsilon_\lambda E$.

Degree one. The part of degree one of $H\star U$ is $\varepsilon Eu_1+H^{(1)}E
-\frac{iq}{2}\mathsf{f}_{\kappa\lambda}\varepsilon_\kappa\partial_\lambda E$. The terms
$\varepsilon Eu_1$ cancel on the two sides, and the equation becomes
\[
-Eu_1'=H^{(1)}E-\tfrac{iq}{2}\mathsf{f}_{\kappa\lambda}\varepsilon_\kappa\partial_\lambda E .
\]
Since $\partial_\lambda E=-\sigma\varepsilon_\lambda E$ and
$\mathsf{f}_{\kappa\lambda}\varepsilon_\kappa\varepsilon_\lambda=0$ by antisymmetry, the last term
vanishes; thus $u_1'=-H^{(1)}$, and with $u_1(0)=0$ we get $u_1=-\sigma H^{(1)}$.

Degree two. After the same cancellation of $\varepsilon Eu_2$, the part of degree two of the
equation is
\[
\begin{split}
-Eu_2'={}&H^{(2)}E+H^{(1)}Eu_1
-\tfrac{iq}{2}\mathsf{f}_{\kappa\lambda}\bigl[\partial_\kappa H^{(1)}\partial_\lambda E
+\varepsilon_\kappa\partial_\lambda(Eu_1)\bigr]\\
&-\tfrac{q^2}{8}\mathsf{f}_{\kappa\lambda}\mathsf{f}_{\kappa'\lambda'}\varepsilon_{\kappa\kappa'}
\partial_\lambda\partial_{\lambda'}E .
\end{split}
\]
Inserting $u_1=-\sigma H^{(1)}$ and $\partial_\lambda E=-\sigma\varepsilon_\lambda E$, the square
bracket equals
\[
E\bigl[-\sigma\,\partial_\kappa H^{(1)}\varepsilon_\lambda
+\sigma^2\varepsilon_\kappa\varepsilon_\lambda H^{(1)}
-\sigma\,\varepsilon_\kappa\partial_\lambda H^{(1)}\bigr].
\]
Contracted with $\mathsf{f}_{\kappa\lambda}$, the middle term vanishes by antisymmetry; exchanging
$\kappa\leftrightarrow\lambda$ in the first term turns it into
$+\sigma\mathsf{f}_{\kappa\lambda}\varepsilon_\kappa\partial_\lambda H^{(1)}$ (the scalar
$\varepsilon_\lambda$ commutes with $\partial_\kappa H^{(1)}$), which cancels the third. Moreover
$E^{-1}\partial_\lambda\partial_{\lambda'}E=\sigma^2\varepsilon_\lambda\varepsilon_{\lambda'}
-\sigma\varepsilon_{\lambda\lambda'}$. Dividing by $-E$,
\[
u_2'=-H^{(2)}+\sigma\bigl(H^{(1)}\bigr)^2+\tfrac{q^2}{8}\mathsf{f}_{\kappa\lambda}
\mathsf{f}_{\kappa'\lambda'}\varepsilon_{\kappa\kappa'}\bigl(\sigma^2\varepsilon_\lambda
\varepsilon_{\lambda'}-\sigma\varepsilon_{\lambda\lambda'}\bigr),
\]
and integrating from $u_2(0)=0$ gives $u_2=-\sigma H^{(2)}+\frac{\sigma^2}{2}(H^{(1)})^2
+\mathsf{a}_2(\sigma)$ with $\mathsf{a}_2$ as in \eqref{10.6:eq-functional-calculus-2}. Each
degree-wise equation is equivalent to the displayed equation for $u_n'$, so
\eqref{10.6:eq-functional-calculus-1} is a solution and the only one. Finally, with
$\varphi(\lambda)=e^{-\sigma\lambda}$ one has $\varphi^{(m)}(\varepsilon)=(-\sigma)^m
e^{-\sigma\varepsilon}$; inserting this into \eqref{10.6:eq-functional-calculus-a}, the last term
becomes $-\frac{q^2}{8}\mathsf{f}_{\kappa\lambda}\mathsf{f}_{\kappa'\lambda'}\varepsilon_{\kappa\kappa'}
[-\frac{\sigma^3}{3}\varepsilon_\lambda\varepsilon_{\lambda'}+\frac{\sigma^2}{2}
\varepsilon_{\lambda\lambda'}]e^{-\sigma\varepsilon}=e^{-\sigma\varepsilon}\mathsf{a}_2(\sigma)$, and
the other terms give those of \eqref{10.6:eq-functional-calculus-1}.

(ii) For $\varphi(\lambda)=\lambda$ one has $\varphi'=1$, $\varphi''=\varphi'''=0$, so
\eqref{10.6:eq-functional-calculus-a} gives $\varepsilon\mathbf{1}+H^{(1)}+H^{(2)}=H$.

Exponentials. Fix $\sigma_2$ and put $V_1(\sigma_1):=U(\sigma_1)\star U(\sigma_2)$,
$V_2(\sigma_1):=U(\sigma_1+\sigma_2)$. Since $\star$ is bilinear with $\sigma$-independent
coefficients and associative (Definition~\ref{10.6:def-truncated-product}), (i) gives
$\partial_{\sigma_1}V_1=(-H\star U(\sigma_1))\star U(\sigma_2)=-H\star V_1$ and
$V_1(0)=\mathbf{1}\star U(\sigma_2)=U(\sigma_2)$; by (i) also
$\partial_{\sigma_1}V_2=-H\star V_2$, $V_2(0)=U(\sigma_2)$. The difference $W=V_1-V_2$ solves
$\partial_{\sigma_1}W=-H\star W$, $W(0)=0$; in degree $n$ this reads
$\partial_{\sigma_1}W^{(n)}=-\varepsilon W^{(n)}$ plus terms built from $W^{(m)}$, $m<n$, and
their $k$-derivatives, so induction on $n$ gives $W^{(0)}=W^{(1)}=W^{(2)}=0$. Hence
$U(\sigma_1)\star U(\sigma_2)=U(\sigma_1+\sigma_2)$, which is (ii) for
$\varphi_j(\lambda)=e^{-\sigma_j\lambda}$, because $\varphi_1\varphi_2=e^{-(\sigma_1+\sigma_2)\lambda}$.

General $\varphi_1,\varphi_2$. Fix $k$. By \eqref{10.6:eq-functional-calculus-a} the symbol of
$\varphi(H)_{\star}$ near $k$ is a combination of $\varphi^{(a)}(\varepsilon(\cdot))$, $a\le3$, with
coefficients independent of $\varphi$. The product $\star$ at $k$ involves at most two
$k$-derivatives of each factor at $k$, and a $k$-derivative of $\varphi^{(a)}(\varepsilon(\cdot))$
of order at most two is a combination of $\varphi^{(a+1)}(\varepsilon(k))$,
$\varphi^{(a+2)}(\varepsilon(k))$ with coefficients given by derivatives of $\varepsilon$. Hence the
left side of (ii) at $k$ is bilinear in the jets $(\varphi_1^{(a)}(\varepsilon(k)))_{a\le5}$,
$(\varphi_2^{(b)}(\varepsilon(k)))_{b\le5}$ in $\mathbb{R}^6$; by the Leibniz rule applied to
$(\varphi_1\varphi_2)^{(m)}$, $m\le3$, so is the right side. The jet of $e^{-\sigma\lambda}$ at
$\varepsilon(k)$ is $e^{-\sigma\varepsilon(k)}(1,-\sigma,\sigma^2,\dots,-\sigma^5)$; for six distinct
values of $\sigma$ these are nonzero multiples of the rows of a Vandermonde matrix and span
$\mathbb{R}^6$. Two bilinear maps agreeing on spanning sets in each argument agree, so (ii) holds at
$k$, and $k$ was arbitrary.
\end{proof}

\begin{remark}
If $\varepsilon$ is replaced by the continuum symbol $|k|^2$ (so that $\varepsilon_\kappa=2k_\kappa$,
$\varepsilon_{\kappa\kappa'}=2\delta_{\kappa\kappa'}$), formula \eqref{10.6:eq-functional-calculus-2}
becomes $\mathsf{a}_2(\sigma)=-\frac{\sigma}{3}\,k\Theta^2k+\frac14\Theta_{\kappa\lambda}
\Theta_{\lambda\kappa}$ with $\Theta:=\sigma q\mathsf{f}$, which is the term of second order in
$\Theta$ of $e^{\sigma|k|^2}(\det\cos\Theta)^{-1/2}\exp\bigl(-\sigma k(\tan\Theta/\Theta)k\bigr)$,
where $(\det\cos\Theta)^{-1/2}\exp\bigl(-\sigma k(\tan\Theta/\Theta)k\bigr)$ is Mehler's closed form
for the symbol of the heat semigroup of the continuum covariant Laplacian in the constant
field $\mathsf{f}$.
\end{remark}

\begin{lemma}[The proper-time regulator and the heat integral]\label{10.6:lem-proper-time}
For $\lambda\ge 0$ put
\begin{equation}\label{10.6:eq-proper-time-1}
\psi(\lambda):=\frac{1-e^{-\lambda}}{\lambda}=\int_0^1 e^{-s\lambda}\,ds,\qquad
w(\lambda):=\frac{1}{\psi(\lambda)}-\lambda=\frac{\lambda}{e^{\lambda}-1},
\end{equation}
both understood at $\lambda=0$ by continuity. With the functional calculus
\eqref{10.6:eq-functional-calculus-a} and the heat semigroup $U(s)$ of
Lemma~\ref{10.6:lem-functional-calculus}\,(i), put
\begin{equation}\label{10.6:eq-proper-time-2}
\tilde{G}:=\psi(H)_{\star}=\int_0^1 U(s)\,ds,\qquad \bar{\mathsf{W}}:=w(H)_{\star}.
\end{equation}
Then the following hold.
\begin{enumerate}
\item[(i)] $(H+\bar{\mathsf{W}})\star\tilde{G}=1$ in $\mathfrak{A}_{\mathsf{f}}$.
\item[(ii)] $\bar{\mathsf{W}}$ is an admissible infrared family in the sense of
Definition~\ref{10.6:def-infrared-families}, with $\gamma=\tfrac12$ and with constants $C_v$, $C_1$
that are absolute constants times the constants of Lemma~\ref{10.6:lem-background-generators}; in
particular they do not depend on $\eta$.
\item[(iii)] The regulated half-trace of $\bar{\mathsf{W}}$ is
\begin{equation}\label{10.6:eq-proper-time-a}
\begin{split}
\tau[\bar{\mathsf{W}}]&=[\mathsf{f}^2]\tfrac12\operatorname{Tr}_{\square}(\tilde{G}\mathsf{O})
=\int_0^1 ds\int_{\mathrm{BZ}_\eta}\bar{d}k\; e^{-s\varepsilon}\cdot\tfrac12\operatorname{tr}\Big\{
\mathsf{O}^{(2)}-sH^{(1)}\mathsf{O}^{(1)}\\
&\qquad+\Big[-sH^{(2)}+\frac{s^2}{2}\big(H^{(1)}\big)^2+\mathsf{a}_2(s)\Big]\mathsf{O}^{(0)}\Big\}.
\end{split}
\end{equation}
Here $\operatorname{tr}$ is the trace of the matrix-valued symbol at fixed momentum, over its bond
and $\mathfrak{g}_q$ indices, so that $\operatorname{Tr}_{\square}$ of an element of
$\mathfrak{A}_{\mathsf{f}}$ is the integral over $\mathrm{BZ}_\eta$ with respect to $\bar{d}k$ of
$\operatorname{tr}$ of its symbol.
\item[(iv)] With $N=L^k$, so that $\eta=N^{-1}$,
\begin{equation}\label{10.6:eq-proper-time-b}
\begin{split}
\sum_q\tau[\bar{\mathsf{W}}]&=\int_0^{N^2}\mathsf{h}(\sigma)\,d\sigma,\\
\mathsf{h}(\sigma)&:=\sum_q\int_{\mathrm{BZ}_1}\bar{d}\kappa\; e^{-\sigma\varepsilon}\cdot
\tfrac12\operatorname{tr}\Big\{\mathsf{O}^{(2)}-\sigma H^{(1)}\mathsf{O}^{(1)}\\
&\qquad+\Big[-\sigma H^{(2)}+\frac{\sigma^2}{2}\big(H^{(1)}\big)^2+\mathsf{a}_2(\sigma)\Big]
\mathsf{O}^{(0)}\Big\},
\end{split}
\end{equation}
where in the definition of $\mathsf{h}$ the symbols $\varepsilon$, $H^{(n)}$, $\mathsf{O}^{(n)}$ and
$\mathsf{a}_2$ are those of the unit lattice $\eta=1$, evaluated at $\kappa$.
\end{enumerate}
\end{lemma}

\begin{proof}
The two expressions for $\psi$ in \eqref{10.6:eq-proper-time-1} agree because
$\int_0^1e^{-s\lambda}\,ds=(1-e^{-\lambda})/\lambda$, and those for $w$ agree because
$1/\psi(\lambda)-\lambda=\lambda/(1-e^{-\lambda})-\lambda=\lambda e^{-\lambda}/(1-e^{-\lambda})
=\lambda/(e^{\lambda}-1)$. The right-hand side of \eqref{10.6:eq-functional-calculus-a} depends
linearly on $\varphi$ through $\varphi,\varphi',\varphi'',\varphi'''$ evaluated at $\varepsilon$.
Since $e^{-s\lambda}$ is smooth on $[0,1]\times[0,\infty[$, the derivatives of $\psi$ are obtained by
differentiating under the integral, so $\psi(H)_{\star}=\int_0^1(e^{-s\,\cdot})(H)_{\star}\,ds$, and by
Lemma~\ref{10.6:lem-functional-calculus}\,(i) the integrand is $U(s)$. This is the second equality
in \eqref{10.6:eq-proper-time-2}.

\emph{(i).} By Lemma~\ref{10.6:lem-functional-calculus}\,(ii), $\lambda(H)_{\star}=H$, so by linearity
$H+\bar{\mathsf{W}}=(\lambda+w)(H)_{\star}=(1/\psi)(H)_{\star}$. The function
$1/\psi(\lambda)=\lambda/(1-e^{-\lambda})$ lies in $C^\infty([0,\infty[)$, since $\psi$ is smooth and
$\psi(\lambda)=\int_0^1e^{-s\lambda}ds>0$. Applying
Lemma~\ref{10.6:lem-functional-calculus}\,(ii) to $\varphi_1=1/\psi$, $\varphi_2=\psi$ gives
$(H+\bar{\mathsf{W}})\star\tilde{G}=(1)(H)_{\star}$, and \eqref{10.6:eq-functional-calculus-a} for the
constant function $1$ is $1$, all derivatives vanishing. Exchanging $\varphi_1$ and $\varphi_2$ gives
also $\tilde{G}\star(H+\bar{\mathsf{W}})=1$, so $\tilde{G}=(H+\bar{\mathsf{W}})^{-1}$ in
$\mathfrak{A}_{\mathsf{f}}$.

\emph{(ii).} We check the conditions \eqref{10.6:eq-infrared-families-b}. For (a1): for
$\lambda\ge0$,
\begin{equation}\label{10.6:eq-proper-time-3}
\lambda+w(\lambda)=\frac{\lambda}{1-e^{-\lambda}}\ge\max\{1,\lambda\}\ge\tfrac12(1+\lambda),
\end{equation}
because $0\le 1-e^{-\lambda}\le\min\{1,\lambda\}$. In \eqref{10.6:eq-functional-calculus-a} every term
except $\varphi(\varepsilon)$ carries a factor $H^{(1)}$, $H^{(2)}$ or $\mathsf{f}_{\kappa\lambda}
\mathsf{f}_{\kappa'\lambda'}$, hence has degree at least one in $\mathsf{f}$; so
$\bar{\mathsf{W}}^{(0)}=w(\varepsilon)$, and, from $\lambda(H)_{\star}=H$, $H^{(0)}=\varepsilon$.
By \eqref{10.6:eq-proper-time-3},
$H^{(0)}+\bar{\mathsf{W}}^{(0)}=\varepsilon+w(\varepsilon)\ge\tfrac12(1+\varepsilon)=\tfrac12\mathsf{J}^2$,
which is (a1) with $\gamma=\tfrac12$.

For (a2) and (a3): $w\in C^\infty([0,\infty[)$, and for every $m$ there is a constant $C_m$ with
\begin{equation}\label{10.6:eq-proper-time-4}
|w^{(m)}(\lambda)|\le C_m(1+\lambda)e^{-\lambda}\qquad(\lambda\ge0).
\end{equation}
By \eqref{10.6:eq-functional-calculus-a}, the symbols of $\delta^\alpha\bar{\mathsf{W}}^{(n)}$ are
finite sums of terms $w^{(m)}(\varepsilon(k))$ times products of factors
$\partial^\beta\varepsilon$ and $\partial^\beta H^{(n')}$, each factor bounded by
$C(1+|k|)^2$; the degree-$n$ part carries $n$ factors of $\mathsf{f}$, whence the factor
$|\mathsf{f}|^n$. Together with \eqref{10.6:eq-proper-time-4}, the suprema of these symbols over
$k$ and their integrals $\int_{\mathrm{BZ}_\eta}\bar{d}k$ are bounded uniformly in $\eta$; the
suprema bound the norms $\|\cdot\|_{\natural}$ of (a2), the integrals the trace norms per cell
$\|\mathsf{J}^{-1}\cdot\,\mathsf{J}^{-1}\|_{1,\square}$ of (a3).

\emph{(iii).} By Definition~\ref{10.6:def-infrared-families} and (i),
$\tau[\bar{\mathsf{W}}]=[\mathsf{f}^2]\tfrac12\operatorname{Tr}_{\square}((H+\bar{\mathsf{W}})^{-1}
\mathsf{O})=[\mathsf{f}^2]\tfrac12\operatorname{Tr}_{\square}(\tilde{G}\mathsf{O})$. Insert
$\tilde{G}=\int_0^1U(s)\,ds$ and the explicit form
$U(s)=e^{-s\varepsilon}[1-s(H^{(1)}+H^{(2)})+\tfrac{s^2}{2}(H^{(1)})^2+\mathsf{a}_2(s)]$ of
Lemma~\ref{10.6:lem-functional-calculus}\,(i). The part of degree zero in $\mathsf{f}$ of $U(s)$ is
$e^{-s\varepsilon}$, that of degree one is $-se^{-s\varepsilon}H^{(1)}$, and that of degree two is
$e^{-s\varepsilon}[-sH^{(2)}+\tfrac{s^2}{2}(H^{(1)})^2+\mathsf{a}_2(s)]$. Collecting the terms of total
degree two in $\mathsf{f}$, these three parts of $U(s)$ multiply $\mathsf{O}^{(2)}$, $\mathsf{O}^{(1)}$
and $\mathsf{O}^{(0)}$ respectively, and writing the trace per cell as the integral over
$\mathrm{BZ}_\eta$ of the traced symbol gives \eqref{10.6:eq-proper-time-a}.

\emph{(iv).} By the scaling form in Lemma~\ref{10.6:lem-background-generators},
$H^{(n)}(k)=\eta^{2n-2}H^{(n)}_{\eta=1}(\eta k)$, and likewise for $\mathsf{O}^{(n)}$ and, with
$n=0$, for $\varepsilon$; hence the first momentum derivatives of $\varepsilon$ satisfy
$\varepsilon_\mu(k)=\eta^{-1}\varepsilon_{\eta=1,\mu}(\eta k)$ and the second ones
$\varepsilon_{\mu\nu}(k)=\varepsilon_{\eta=1,\mu\nu}(\eta k)$. Substitute
$k=\kappa/\eta$, $s=\eta^2\sigma$. Then $e^{-s\varepsilon(k)}=e^{-\sigma\varepsilon_{\eta=1}(\kappa)}$.
In the bracket of \eqref{10.6:eq-proper-time-a}: $\mathsf{O}^{(2)}$ carries $\eta^2$;
$sH^{(1)}\mathsf{O}^{(1)}$ carries $\eta^2\cdot\eta^0\cdot\eta^0$; $sH^{(2)}\mathsf{O}^{(0)}$ carries
$\eta^2\cdot\eta^2\cdot\eta^{-2}$; $\tfrac{s^2}{2}(H^{(1)})^2\mathsf{O}^{(0)}$ carries
$\eta^4\cdot\eta^{-2}$; and in $\mathsf{a}_2(s)$, where each first momentum derivative of
$\varepsilon$ carries $\eta^{-1}$ and each second one carries $\eta^{0}$, the term with $s^3$ and two
first derivatives and the term with $s^2$ and one second derivative both carry $\eta^4$, while the
remaining second derivative in front of the bracket carries $\eta^0$, so
$\mathsf{a}_2(s)\mathsf{O}^{(0)}$ carries $\eta^4\cdot\eta^{-2}$. Thus the bracket is $\eta^2$ times
the same expression in unit-lattice symbols at $(\kappa,\sigma)$. Since $ds=\eta^2d\sigma$, and since
$\bar{d}k$ is a constant multiple of Lebesgue measure, so that the linear substitution
$k=\kappa/\eta$ maps $\mathrm{BZ}_\eta$ onto $\mathrm{BZ}_1$ and gives $\bar{d}k=\eta^{-4}\bar{d}\kappa$,
the powers of $\eta$ cancel and $\eta$ disappears from the integrand. The range $s\in[0,1]$ becomes
$\sigma\in[0,\eta^{-2}]=[0,N^2]$. Summing over $q$ gives \eqref{10.6:eq-proper-time-b}.
\end{proof}

\begin{lemma}[First-order symbols at zero momentum]\label{10.6:lem-zero-momentum-symbols}
Let $\eta=1$, so that the spacing of $\mathbb{L}_0$ is $\ell_0=1$ and the symbol of the lattice
gradient is $d_\mu(k)=2i\sin(k_\mu/2)$, $\mu=1,\dots,4$. Let $h=(h_1,\dots,h_4)$ be the
orientation weights of $\mathsf{O}=\sum_{\mu<\nu}(h_\mu+h_\nu)\mathsf{H}_{\mu\nu}$, put
$c_{\mu\nu}:=h_\mu+h_\nu$, write $|h|$ for a norm of $h$ (the constant $C$ below depending on the
choice of norm), and let
$H=\mathsf{H}_0+D_0D_0^{\dagger}$ be the background Feynman operator, the product being that of
$\mathfrak{A}_{\mathsf{f}}$. For $t\in\mathfrak{A}_{\mathsf{f}}$ write $t^{(n)}(k)$ for the
degree-$n$ part in $\mathsf{f}$ of its reduced symbol at momentum $k$, and $[\,\cdot\,]_{\nu\mu}$
for the entry in row $\nu$ and column $\mu$. Then the following hold, with a constant $C$:
\begin{equation}\label{10.6:eq-zero-momentum-symbols-a}
\begin{aligned}
\text{(i)}\quad & \bigl[\mathsf{O}^{(0)}(k)\bigr]_{\nu\nu}=\sum_{\mu\neq\nu}c_{\mu\nu}|d_\mu(k)|^2,\\
& \bigl[\mathsf{O}^{(0)}(k)\bigr]_{\nu\mu}=-c_{\mu\nu}\,d_\nu(k)\overline{d_\mu(k)}\quad(\nu\neq\mu),\\
& |d_\mu(k)|^2=4\sin^2(k_\mu/2);\\
\text{(ii)}\quad & \bigl[\mathsf{H}^{(1)}_{\mu\nu}(0)\bigr]_{\nu\mu}=\tfrac{3iq}{2}\mathsf{f}_{\mu\nu},\\
& \bigl[\mathsf{H}^{(1)}_{\mu\nu}(0)\bigr]_{\mu\nu}=\tfrac{3iq}{2}\mathsf{f}_{\nu\mu},\\
& \text{all other entries of }\mathsf{H}^{(1)}_{\mu\nu}(0)\text{ vanish},\\
& \bigl[(D_0D_0^{\dagger})^{(1)}(0)\bigr]_{\nu\mu}=\tfrac{iq}{2}\mathsf{f}_{\mu\nu},\\
& \bigl[\mathsf{O}^{(1)}(0)\bigr]_{\nu\mu}=\tfrac{3iq}{2}c_{\mu\nu}\mathsf{f}_{\mu\nu},\\
& \bigl[H^{(1)}(0)\bigr]_{\nu\mu}=2iq\,\mathsf{f}_{\mu\nu};\\
\text{(iii)}\quad & H^{(n)},\ \mathsf{O}^{(n)}\ \text{are trigonometric polynomials in }k,\\
& |H^{(n)}|\le C|\mathsf{f}|^n,\qquad |\mathsf{O}^{(n)}|\le C|h|\,|\mathsf{f}|^n .
\end{aligned}
\end{equation}
In part (ii), $\mu<\nu$ in the statement on $\mathsf{H}^{(1)}_{\mu\nu}(0)$, and each entry in part
(ii) is the stated scalar times the identity of $\mathfrak{g}_q$.
\end{lemma}

\begin{proof}
Parts (i) and (iii) come from the description of the generators in
Lemma~\ref{10.6:lem-background-generators}. For part (i): the flat part of $\mathsf{H}_{\mu\nu}$ is
given there by $(\mathsf{H}^{(0)}_{\mu\nu}A,A)=\|(dA)_{\mu\nu}\|^2$. Writing $\hat{A}_\nu(k)$ for
the Fourier transform of $A_\nu$, in momentum space $(dA)_{\mu\nu}$ becomes
$(dA)_{\mu\nu}(k)=d_\mu(k)\hat{A}_\nu(k)-d_\nu(k)\hat{A}_\mu(k)$. Expanding
\[
\begin{aligned}
|d_\mu\hat{A}_\nu-d_\nu\hat{A}_\mu|^2&=|d_\mu|^2|\hat{A}_\nu|^2+|d_\nu|^2|\hat{A}_\mu|^2\\
&\quad-\overline{\hat{A}_\nu}\,d_\nu\bar{d}_\mu\,\hat{A}_\mu
-\overline{\hat{A}_\mu}\,d_\mu\bar{d}_\nu\,\hat{A}_\nu
\end{aligned}
\]
and reading off the Hermitian symbol, $\mathsf{H}^{(0)}_{\mu\nu}$ has the entries
$(\nu,\nu)$: $|d_\mu|^2$, $(\mu,\mu)$: $|d_\nu|^2$, $(\nu,\mu)$: $-d_\nu\bar{d}_\mu$,
$(\mu,\nu)$: $-d_\mu\bar{d}_\nu$, and no other. Summing with the weights $c_{\mu\nu}=c_{\nu\mu}$,
the entry $(\nu,\nu)$ of $\mathsf{O}^{(0)}$ receives $c_{\mu\nu}|d_\mu|^2$ from each plane
$\{\mu,\nu\}$, $\mu\neq\nu$, and the entry $(\nu,\mu)$, $\nu\ne\mu$, receives only the
contribution $-c_{\mu\nu}d_\nu\bar{d}_\mu$ of the plane $\{\mu,\nu\}$. Finally
$|d_\mu|^2=|2i\sin(k_\mu/2)|^2=4\sin^2(k_\mu/2)$.

For part (iii): by Lemma~\ref{10.6:lem-background-generators}, at $\eta=1$ the reduced symbol of
$\mathsf{H}_{\mu\nu}$ is $\mathsf{t}_1(k;\mathsf{f})$, a trigonometric polynomial in $k$ with
coefficients entire in $\mathsf{f}$. Its degree-$n$ part is therefore a trigonometric polynomial
in $k$ whose coefficients are homogeneous polynomials of degree $n$ in $\mathsf{f}$. It is thus
bounded by $C|\mathsf{f}|^n$. The symbol $d_\mu(k)$ of $D_0$ is $\mathsf{f}$-free and
trigonometric. By the truncated product (Definition~\ref{10.6:def-truncated-product}), the
degree-$n$ part of $D_0D_0^{\dagger}$ is a finite sum of products of $n$ entries of
$\mathsf{f}$ with derivatives of $d$ and $\bar d$, hence again a trigonometric polynomial bounded
by $C|\mathsf{f}|^n$. The statements for $H=\mathsf{H}_0+D_0D_0^\dagger$ follow by summation.
Those for $\mathsf{O}$ follow from $|c_{\mu\nu}|\le|h_\mu|+|h_\nu|$, which is bounded by a constant
times $|h|$.

Part (ii). Write $d=(d_\sigma)_\sigma$ for the column symbol of $D_0$ and $\partial_\kappa$ for
$\partial/\partial k_\kappa$. The product in $\mathfrak{A}_{\mathsf{f}}$ is $\star$, so by
\eqref{10.6:eq-background-generators-a} we have $\mathsf{H}_{\mu\nu}\star d=0$. Since $d$ is
$\mathsf{f}$-free, $d^{(0)}=d$ and $d^{(n)}=0$ for $n\ge1$. The degree-one part of
$\mathsf{H}_{\mu\nu}\star d$ is therefore $\mathsf{H}^{(1)}_{\mu\nu}d$ plus the degree-one term of
the truncated product with $t=\mathsf{H}^{(0)}_{\mu\nu}$ and $s=d$. On reduced symbols
$\delta_\kappa$ is $\pm i\,\partial_\kappa$ according to the sign convention of the Fourier
transform; since two derivations enter, the degree-one term
$\tfrac{iq}{2}\mathsf{f}_{\kappa\lambda}(\delta_\kappa t)(\delta_\lambda s)$ carries the factor
$(\pm i)^2=-1$ in either convention, and reads
$-\tfrac{iq}{2}\mathsf{f}_{\kappa\lambda}(\partial_\kappa t)(\partial_\lambda s)$. Thus
the degree-one part of \eqref{10.6:eq-background-generators-a} becomes
\begin{equation}\label{10.6:eq-zero-momentum-symbols-1}
\mathsf{H}^{(1)}_{\mu\nu}(k)\,d(k)=\tfrac{iq}{2}\,\mathsf{f}_{\kappa\lambda}\,
\partial_\kappa\mathsf{H}^{(0)}_{\mu\nu}(k)\,\partial_\lambda d(k),
\end{equation}
summed over $\kappa,\lambda$. We compare the terms of first order in $k$. We have
$d_\sigma=ik_\sigma+O(|k|^3)$ and $\partial_\lambda d_\sigma=i\delta_{\lambda\sigma}+O(|k|^2)$.
By the entries of $\mathsf{H}^{(0)}_{\mu\nu}$ found above and
$|d_\mu|^2=4\sin^2(k_\mu/2)=k_\mu^2+O(|k|^4)$, $d_\nu\bar{d}_\mu=k_\nu k_\mu+O(|k|^4)$, up to
terms of order four $\mathsf{H}^{(0)}_{\mu\nu}$ has the entries $(\nu,\nu)$: $k_\mu^2$,
$(\mu,\mu)$: $k_\nu^2$, $(\nu,\mu)=(\mu,\nu)$: $-k_\mu k_\nu$.

Row $\nu$ of the right side of \eqref{10.6:eq-zero-momentum-symbols-1} is, to first order,
$\tfrac{iq}{2}\,i\sum_{\kappa,\lambda}\mathsf{f}_{\kappa\lambda}
\bigl[\partial_\kappa\mathsf{H}^{(0)}_{\mu\nu}\bigr]_{\nu\lambda}$. For $\lambda=\nu$ only
$\kappa=\mu$ contributes, with $2k_\mu$. For $\lambda=\mu$, $\kappa=\nu$ gives $-k_\mu$, while
$\kappa=\mu$ carries $\mathsf{f}_{\mu\mu}=0$, $\mathsf{f}$ being antisymmetric. Hence row $\nu$
equals
\[
\tfrac{iq}{2}\cdot i\cdot\bigl[\mathsf{f}_{\mu\nu}\cdot2k_\mu+\mathsf{f}_{\nu\mu}\cdot(-k_\mu)\bigr]
=-\tfrac{3q}{2}\,\mathsf{f}_{\mu\nu}\,k_\mu .
\]
Row $\nu$ of the left side is, to first order,
$i\sum_\sigma[\mathsf{H}^{(1)}_{\mu\nu}(0)]_{\nu\sigma}k_\sigma$, since
$\mathsf{H}^{(1)}_{\mu\nu}(k)=\mathsf{H}^{(1)}_{\mu\nu}(0)+O(|k|)$ and $d=O(|k|)$. As this holds
for all $k$, $[\mathsf{H}^{(1)}_{\mu\nu}(0)]_{\nu\mu}=\tfrac{3iq}{2}\mathsf{f}_{\mu\nu}$ and
$[\mathsf{H}^{(1)}_{\mu\nu}(0)]_{\nu\sigma}=0$ for $\sigma\neq\mu$.

In row $\mu$, $\lambda=\mu$, $\kappa=\nu$ gives $\mathsf{f}_{\nu\mu}\cdot2k_\nu$, and
$\lambda=\nu$, $\kappa=\mu$ gives $\mathsf{f}_{\mu\nu}\cdot(-k_\nu)$. The sum, times
$\tfrac{iq}{2}\,i$, is $-\tfrac{3q}{2}\mathsf{f}_{\nu\mu}k_\nu$. Hence
$[\mathsf{H}^{(1)}_{\mu\nu}(0)]_{\mu\nu}=\tfrac{3iq}{2}\mathsf{f}_{\nu\mu}$ and the other entries
of row $\mu$ vanish. For $\rho\notin\{\mu,\nu\}$, row $\rho$ of $\mathsf{H}^{(0)}_{\mu\nu}$ is
zero, so the right side vanishes there. Then
$\sum_\sigma[\mathsf{H}^{(1)}_{\mu\nu}(0)]_{\rho\sigma}k_\sigma=0$ for all $k$, and row $\rho$
vanishes.

For $D_0D_0^\dagger$, the entry $(\nu,\mu)$ is $d_\nu\star\bar{d}_\mu$. By the same reading of the
truncated product,
\[
d_\nu\star\bar{d}_\mu=d_\nu\bar{d}_\mu-\tfrac{iq}{2}\,\mathsf{f}_{\nu\mu}\,\partial_\nu d_\nu\,
\partial_\mu\bar{d}_\mu+\dotsb,
\]
where the dots are of degree two. Only $\kappa=\nu$,
$\lambda=\mu$ contribute, because $d_\nu$ depends on $k_\nu$ alone. With $\partial_\nu d_\nu(0)=i$
and $\partial_\mu\bar{d}_\mu(0)=-i$ this gives
$[(D_0D_0^\dagger)^{(1)}(0)]_{\nu\mu}=-\tfrac{iq}{2}\mathsf{f}_{\nu\mu}
=\tfrac{iq}{2}\mathsf{f}_{\mu\nu}$, which also vanishes for $\nu=\mu$.

Summing over planes, the entry $(\nu,\mu)$ of $\mathsf{O}^{(1)}(0)$ comes only from the plane
$\{\mu,\nu\}$, which gives
$[\mathsf{O}^{(1)}(0)]_{\nu\mu}=\tfrac{3iq}{2}c_{\mu\nu}\mathsf{f}_{\mu\nu}$
in either order of $\mu$ and $\nu$. Likewise
$[H^{(1)}(0)]_{\nu\mu}=(\tfrac{3iq}{2}+\tfrac{iq}{2})\mathsf{f}_{\mu\nu}=2iq\,\mathsf{f}_{\mu\nu}$.
For $\nu=\mu$ the entries of $\mathsf{O}^{(1)}(0)$ and $H^{(1)}(0)$ vanish, by the vanishing of
the diagonal entries of $\mathsf{H}^{(1)}_{\mu\nu}(0)$ and of $(D_0D_0^\dagger)^{(1)}(0)$ found
above, and so do the stated right-hand sides, since $\mathsf{f}_{\nu\nu}=0$; hence these two
formulas of part (ii) hold for all $\nu,\mu$.
\end{proof}

The factor $2=\tfrac32+\tfrac12$ in $[H^{(1)}(0)]_{\nu\mu}=2iq\,\mathsf{f}_{\mu\nu}$ is a
consequence of \eqref{10.6:eq-background-generators-a} together with the truncated product; of the
plaquette sums only their flat quadratic forms $\mathsf{H}^{(0)}_{\mu\nu}$ enter the proof of
part (ii) of Lemma~\ref{10.6:lem-zero-momentum-symbols}.

\begin{lemma}[Gaussian asymptotics of lattice heat integrals]\label{10.6:lem-heat-asymptotics}
Let $\varepsilon(\kappa)=\sum_{\mu}4\sin^2(\kappa_\mu/2)$ be the lattice Laplace symbol at $\eta=1$.
Let $P$ be a trigonometric polynomial on
$\mathbb{R}^4$ such that
\begin{equation}\label{10.6:eq-heat-asymptotics-1}
P(\kappa)=P^{[m]}(\kappa)+O(|\kappa|^{m+1})\qquad(\kappa\to 0),
\end{equation}
where $P^{[m]}$ is a polynomial, homogeneous of degree $m$. Then, for $\sigma\ge 1$,
\begin{equation}\label{10.6:eq-heat-asymptotics-2}
\int_{\mathrm{BZ}_1}\bar{d}\kappa\,e^{-\sigma\varepsilon(\kappa)}P(\kappa)
=\sigma^{-2-m/2}\Bigl[\int_{\mathbb{R}^4}\bar{d}u\,e^{-|u|^2}P^{[m]}(u)+O(\sigma^{-1/2})\Bigr].
\end{equation}
\end{lemma}

\begin{proof}
Substitute $\kappa=u/\sqrt{\sigma}$; then $\bar{d}\kappa=\sigma^{-2}\bar{d}u$, and the left side of
\eqref{10.6:eq-heat-asymptotics-2} equals
$\sigma^{-2}\int_{\sqrt{\sigma}\,\mathrm{BZ}_1}\bar{d}u\,e^{-\sigma\varepsilon(u/\sqrt{\sigma})}
P(u/\sqrt{\sigma})$.

Two bounds on the exponent. For $\kappa\in\mathrm{BZ}_1$ one has $|\kappa_\mu|\le\pi$ for each
$\mu$, and $(2/\pi)y\le\sin y\le y$ for $0\le y\le\pi/2$ gives
$(4/\pi^2)\kappa_\mu^2\le4\sin^2(\kappa_\mu/2)\le\kappa_\mu^2$; summing over $\mu$,
$(4/\pi^2)|\kappa|^2\le\varepsilon(\kappa)\le|\kappa|^2$ on $\mathrm{BZ}_1$; with
$\kappa=u/\sqrt{\sigma}$ and multiplication by $\sigma$ this gives
$(4/\pi^2)|u|^2\le\sigma\varepsilon(u/\sqrt{\sigma})\le|u|^2$ on $\sqrt{\sigma}\,\mathrm{BZ}_1$.
Next, $\cos x\le 1-x^2/2+x^4/24$ for real $x$ gives $4\sin^2(x/2)=2(1-\cos x)\ge x^2-x^4/12$;
summing over $\mu$ and using $\sum_\mu\kappa_\mu^4\le|\kappa|^4$ yields
$0\le|\kappa|^2-\varepsilon(\kappa)\le|\kappa|^4/12$, hence
$0\le|u|^2-\sigma\varepsilon(u/\sqrt{\sigma})\le|u|^4/(12\sigma)$. For $0\le a\le b$ the mean value
theorem gives $|e^{-a}-e^{-b}|\le(b-a)e^{-a}$; with $a=\sigma\varepsilon(u/\sqrt{\sigma})$,
$b=|u|^2$ and the lower bound on $a$,
\begin{equation}\label{10.6:eq-heat-asymptotics-3}
\bigl|e^{-\sigma\varepsilon(u/\sqrt{\sigma})}-e^{-|u|^2}\bigr|
\le\frac{|u|^4}{12\sigma}\,e^{-(4/\pi^2)|u|^2}\qquad(u\in\sqrt{\sigma}\,\mathrm{BZ}_1).
\end{equation}

The polynomial factor. Since $P$ and $P^{[m]}$ are bounded on the bounded set $\mathrm{BZ}_1$, the
bound \eqref{10.6:eq-heat-asymptotics-1} near $0$ extends to $|P(\kappa)-P^{[m]}(\kappa)|\le
C|\kappa|^{m+1}$ on all of $\mathrm{BZ}_1$, with $C$ depending on $P$. By homogeneity
$P^{[m]}(u/\sqrt{\sigma})=\sigma^{-m/2}P^{[m]}(u)$, so
\begin{equation}\label{10.6:eq-heat-asymptotics-4}
\bigl|P(u/\sqrt{\sigma})-\sigma^{-m/2}P^{[m]}(u)\bigr|\le C\sigma^{-(m+1)/2}|u|^{m+1}
\qquad(u\in\sqrt{\sigma}\,\mathrm{BZ}_1).
\end{equation}
Homogeneity also gives $|P^{[m]}(u)|\le C|u|^m$, so by \eqref{10.6:eq-heat-asymptotics-4} and
$\sigma\ge1$, $|P(u/\sqrt{\sigma})|\le C\sigma^{-m/2}(|u|^m+|u|^{m+1})$ on
$\sqrt{\sigma}\,\mathrm{BZ}_1$.

Assembly. On $\sqrt{\sigma}\,\mathrm{BZ}_1$ write
\begin{align*}
e^{-\sigma\varepsilon(u/\sqrt{\sigma})}P(u/\sqrt{\sigma})
&=\sigma^{-m/2}e^{-|u|^2}P^{[m]}(u)
+\bigl(e^{-\sigma\varepsilon(u/\sqrt{\sigma})}-e^{-|u|^2}\bigr)P(u/\sqrt{\sigma})\\
&\quad+e^{-|u|^2}\bigl(P(u/\sqrt{\sigma})-\sigma^{-m/2}P^{[m]}(u)\bigr).
\end{align*}
By \eqref{10.6:eq-heat-asymptotics-3} and the bound on $|P(u/\sqrt{\sigma})|$, the second term is
at most $C\sigma^{-1-m/2}|u|^4(|u|^m+|u|^{m+1})e^{-(4/\pi^2)|u|^2}$, whose integral over
$\mathbb{R}^4$ is $O(\sigma^{-1-m/2})$, hence $O(\sigma^{-(m+1)/2})$ because $\sigma\ge1$. By
\eqref{10.6:eq-heat-asymptotics-4} the third term integrates to at most
$C\sigma^{-(m+1)/2}\int\bar{d}u\,|u|^{m+1}e^{-|u|^2}=O(\sigma^{-(m+1)/2})$. In the first term,
replace $\sqrt{\sigma}\,\mathrm{BZ}_1$ by $\mathbb{R}^4$: if the ball of radius $r_0>0$ about $0$
lies in $\mathrm{BZ}_1$, the complement of $\sqrt{\sigma}\,\mathrm{BZ}_1$ lies in
$\{|u|\ge r_0\sqrt{\sigma}\}$, where
$|P^{[m]}(u)|e^{-|u|^2}\le C|u|^m e^{-r_0^2\sigma/2}e^{-|u|^2/2}$; since $|u|^m e^{-|u|^2/2}$ is
integrable over $\mathbb{R}^4$, this gives
\[
\sigma^{-m/2}\int_{\sqrt{\sigma}\,\mathrm{BZ}_1}\bar{d}u\,e^{-|u|^2}P^{[m]}(u)
=\sigma^{-m/2}\Bigl[\int_{\mathbb{R}^4}\bar{d}u\,e^{-|u|^2}P^{[m]}(u)+O(e^{-c\sigma})\Bigr],
\]
and $e^{-c\sigma}\le C\sigma^{-1/2}$. Collecting the three terms, the integral over
$\sqrt{\sigma}\,\mathrm{BZ}_1$ equals $\sigma^{-m/2}\bigl[\int_{\mathbb{R}^4}\bar{d}u\,
e^{-|u|^2}P^{[m]}(u)+O(\sigma^{-1/2})\bigr]$, since $O(\sigma^{-(m+1)/2})=\sigma^{-m/2}
O(\sigma^{-1/2})$. Multiplying by $\sigma^{-2}$ gives
\eqref{10.6:eq-heat-asymptotics-2}.
\end{proof}

We record the Gaussian moments used below. Put
$\mathsf{n}:=\int_{\mathbb{R}^4}\bar{d}u\,e^{-|u|^2}=1/(16\pi^2)$, which is $\pi^2/(2\pi)^4$ for
$\bar{d}u=(2\pi)^{-4}d^4u$, and $\langle F\rangle:=\mathsf{n}^{-1}\int\bar{d}u\,e^{-|u|^2}F(u)$.
The coordinates are independent centred Gaussians of variance $\tfrac12$, with
$\langle u_\alpha^4\rangle=\tfrac34$; hence, with no summation over $\mu$,
\begin{equation}\label{10.6:eq-heat-asymptotics-5}
\langle u_\alpha u_\beta\rangle=\tfrac12\delta_{\alpha\beta},\qquad
\langle u_\alpha u_\beta u_\mu^2\rangle
=\tfrac14\bigl(\delta_{\alpha\beta}+2\delta_{\alpha\mu}\delta_{\beta\mu}\bigr).
\end{equation}

\begin{lemma}[The large-time heat coefficient]\label{10.6:lem-heat-coefficient}
Let $h=(h_\mu)$ satisfy $\sum_\mu h_\mu=0$, let the orientation weights be
$c_{\mu\nu}=h_\mu+h_\nu$, and let $\mathsf{h}$ be the heat integrand
defined in \eqref{10.6:eq-proper-time-b}. Then $\mathsf{h}$ is continuous on $[0,\infty[$,
$|\mathsf{h}(\sigma)|\le C|h||\mathsf{f}|^2$ for $0\le\sigma\le 1$, and for $\sigma\ge 1$
\begin{equation}\label{10.6:eq-heat-coefficient-1}
\Bigl|\,\mathsf{h}(\sigma)+\frac{11\,C_{\mathfrak{g}}}{48\pi^2}\,\sigma^{-1}\,
\mathsf{t}^{\mathbf{3}}(h;\mathsf{f})\Bigr|
\le C_{\mathrm{HK}}\,|h|\,|\mathsf{f}|^2\,\sigma^{-3/2}.
\end{equation}
\end{lemma}

\begin{proof}
Recall from \eqref{10.6:eq-proper-time-b} that, with unit-lattice symbols,
\begin{equation*}
\begin{split}
\mathsf{h}(\sigma)=\sum_q\int_{\mathrm{BZ}_1}\bar{d}\kappa\,e^{-\sigma\varepsilon(\kappa)}\,
\tfrac12\operatorname{tr}\Bigl\{&\mathsf{O}^{(2)}-\sigma H^{(1)}\mathsf{O}^{(1)}\\
&+\Bigl[-\sigma H^{(2)}+\tfrac{\sigma^2}{2}\bigl(H^{(1)}\bigr)^2+\mathsf{a}_2(\sigma)\Bigr]
\mathsf{O}^{(0)}\Bigr\},
\end{split}
\end{equation*}
where, by Lemma~\ref{10.6:lem-functional-calculus} (summation over repeated indices),
\begin{equation*}
\mathsf{a}_2(\sigma)=\frac{q^2}{8}\,\mathsf{f}_{\rho\lambda}\mathsf{f}_{\rho'\lambda'}\,
\varepsilon_{\rho\rho'}\Bigl[\frac{\sigma^3}{3}\,\varepsilon_\lambda\varepsilon_{\lambda'}
-\frac{\sigma^2}{2}\,\varepsilon_{\lambda\lambda'}\Bigr].
\end{equation*}

\emph{Continuity and the bound on $[0,1]$.} By Lemma~\ref{10.6:lem-zero-momentum-symbols},
item (y3) of \eqref{10.6:eq-zero-momentum-symbols-a},
$H^{(n)}$ and $\mathsf{O}^{(n)}$ are trigonometric polynomials bounded by $C|\mathsf{f}|^n$ and
$C|h||\mathsf{f}|^n$ respectively; $\mathsf{a}_2(\sigma)$ is a polynomial in $\sigma$ whose
coefficients are trigonometric polynomials bounded by $C|\mathsf{f}|^2$; and
$0\le e^{-\sigma\varepsilon}\le 1$. Hence the integrand is jointly continuous in
$(\sigma,\kappa)\in[0,\infty[\times\mathrm{BZ}_1$ and, for $0\le\sigma\le1$, each of its terms is
bounded by $C|h||\mathsf{f}|^2$. Since $\mathrm{BZ}_1$ is compact and the sum over $q$ is finite,
$\mathsf{h}$ is continuous on $[0,\infty[$ and $|\mathsf{h}(\sigma)|\le C|h||\mathsf{f}|^2$
on $[0,1]$.

\emph{Conventions for $\sigma\ge1$.} Fix $\sigma\ge1$ and a sector $q$, and put
$c_\nu:=\sum_{\mu\ne\nu}c_{\mu\nu}$. We use the Gaussian asymptotics of
Lemma~\ref{10.6:lem-heat-asymptotics}, with its order $m$ and principal part $P^{[m]}$ of the
integrand at $\kappa=0$.
We write $\mathsf{n}:=\int\bar{d}u\,e^{-|u|^2}=1/(16\pi^2)$ and
$\langle P\rangle:=\mathsf{n}^{-1}\int\bar{d}u\,e^{-|u|^2}P(u)$ for a polynomial $P$, so that
\begin{equation*}
\langle u_\alpha u_\beta\rangle=\tfrac12\,\delta_{\alpha\beta},\qquad
\langle u_\alpha u_\beta u_\mu^2\rangle=\tfrac14\bigl(\delta_{\alpha\beta}
+2\delta_{\alpha\mu}\delta_{\beta\mu}\bigr).
\end{equation*}
Every integrand below is a finite sum of trigonometric polynomials in $\kappa$, independent of $h$
and $\mathsf{f}$, multiplied by monomials of degree one in the components of $h$ and two in those
of $\mathsf{f}$, and by powers of $\sigma$: indeed $\mathsf{O}$ depends linearly on the weights,
and $H^{(n)}$, $\mathsf{O}^{(n)}$, $\mathsf{a}_2$ are of degree $n$, $n$, $2$ in $\mathsf{f}$.
Applying Lemma~\ref{10.6:lem-heat-asymptotics} to each of these finitely many trigonometric
polynomials, the implied constant in every $O(\cdot)$ written below is of the form
$C|h||\mathsf{f}|^2$. The trace $\operatorname{tr}$ runs over the bond directions
and over $\mathfrak{g}_q$, on which the leading symbols displayed in
Lemma~\ref{10.6:lem-zero-momentum-symbols} act as multiples of the identity; the trace over
$\mathfrak{g}_q$ therefore contributes the factor $\dim\mathfrak{g}_q$. The lowest-order parts at
$\kappa=0$ that we use are: $\varepsilon_{\rho\rho'}\to2\delta_{\rho\rho'}$,
$\varepsilon_\lambda\to2\kappa_\lambda$, and, by item (y1) in
\eqref{10.6:eq-zero-momentum-symbols-a}, $|d_\mu|^2\to\kappa_\mu^2$ while
$d_\nu\bar d_\mu$ is of order two with principal part a multiple of $\kappa_\nu\kappa_\mu$; in
particular $\mathsf{O}^{(0)}=O(|\kappa|^2)$, with principal part $\mathsf{O}^{(0)[2]}$ satisfying
$\mathsf{O}^{(0)[2]}_{\nu\nu}(\kappa)=\sum_{\mu\ne\nu}c_{\mu\nu}\kappa_\mu^2$ and
$\langle\mathsf{O}^{(0)[2]}_{\nu\mu}\rangle=0$ for $\nu\ne\mu$.

Put $\int:=\int_{\mathrm{BZ}_1}\bar{d}\kappa$ and split the sector-$q$ integral of the brace
(without the factor $\frac12$) as $T_1+\dots+T_5$, where
\begin{gather*}
T_1:=\int e^{-\sigma\varepsilon}\operatorname{tr}\mathsf{O}^{(2)},\quad
T_2:=-\sigma\!\int e^{-\sigma\varepsilon}\operatorname{tr}H^{(1)}\mathsf{O}^{(1)},\quad
T_3:=-\sigma\!\int e^{-\sigma\varepsilon}\operatorname{tr}H^{(2)}\mathsf{O}^{(0)},\\
T_4:=\frac{\sigma^2}{2}\int e^{-\sigma\varepsilon}\operatorname{tr}\bigl(H^{(1)}\bigr)^2
\mathsf{O}^{(0)},\qquad
T_5:=\int e^{-\sigma\varepsilon}\,\mathsf{a}_2(\sigma)\operatorname{tr}\mathsf{O}^{(0)};
\end{gather*}
here $\mathsf{a}_2(\sigma)$ is a scalar, so
$\operatorname{tr}\mathsf{a}_2\mathsf{O}^{(0)}=\mathsf{a}_2\operatorname{tr}\mathsf{O}^{(0)}$.

\emph{The term $T_1$.} Lemma~\ref{10.6:lem-heat-asymptotics} with $m=0$ gives
$T_1=O(\sigma^{-2})$, which is $O(\sigma^{-3/2})$ for $\sigma\ge1$; it does not contribute.

\emph{The term $T_3$.} Since $\mathsf{O}^{(0)}=O(|\kappa|^2)$, the polynomial
$\operatorname{tr}H^{(2)}\mathsf{O}^{(0)}$ has $m=2$, and
Lemma~\ref{10.6:lem-heat-asymptotics} gives an integral $O(\sigma^{-3})$; hence
$T_3=\sigma\cdot O(\sigma^{-3})=O(\sigma^{-2})$, and it does not contribute.

\emph{The term $T_2$.} Here $m=0$ and $P^{[0]}=P(0)$. By item (y2) in
\eqref{10.6:eq-zero-momentum-symbols-a},
$H^{(1)}(0)_{\mu\nu}=2iq\mathsf{f}_{\nu\mu}$ and
$\mathsf{O}^{(1)}(0)_{\nu\mu}=\frac{3iq}{2}c_{\mu\nu}\mathsf{f}_{\mu\nu}$, both vanishing on the
diagonal. Hence, using $\mathsf{f}_{\nu\mu}=-\mathsf{f}_{\mu\nu}$ and $c_{\mu\nu}=c_{\nu\mu}$,
\begin{align*}
\operatorname{tr}H^{(1)}(0)\mathsf{O}^{(1)}(0)
&=\dim\mathfrak{g}_q\sum_{\mu\ne\nu}(2iq\mathsf{f}_{\nu\mu})\tfrac{3iq}{2}c_{\mu\nu}
\mathsf{f}_{\mu\nu}\\
&=3q^2\dim\mathfrak{g}_q\sum_{\mu\ne\nu}c_{\mu\nu}\mathsf{f}_{\mu\nu}^2\\
&=6q^2\dim\mathfrak{g}_q\sum_{\mu<\nu}c_{\mu\nu}\mathsf{f}_{\mu\nu}^2 .
\end{align*}
Therefore
\begin{align*}
T_2&=-\sigma\cdot\sigma^{-2}\bigl[\mathsf{n}\cdot6q^2\dim\mathfrak{g}_q
\sum_{\mu<\nu}c_{\mu\nu}\mathsf{f}^2_{\mu\nu}+O(\sigma^{-1/2})\bigr]\\
&=\mathsf{n}q^2\dim\mathfrak{g}_q\,\sigma^{-1}\bigl(-6\sum_{\mu<\nu}c_{\mu\nu}
\mathsf{f}^2_{\mu\nu}\bigr)+O(\sigma^{-3/2}).
\end{align*}

\emph{The term $T_4$.} Here $m=2$ and
$P^{[2]}=\operatorname{tr}\bigl(H^{(1)}(0)\bigr)^2\mathsf{O}^{(0)[2]}$. By item (y2) in
\eqref{10.6:eq-zero-momentum-symbols-a},
\begin{equation*}
\bigl(H^{(1)}(0)^2\bigr)_{\nu\nu}=\sum_\mu(2iq)^2\mathsf{f}_{\mu\nu}\mathsf{f}_{\nu\mu}
=4q^2\sum_\mu\mathsf{f}^2_{\mu\nu}=4q^2\phi_\nu .
\end{equation*}
The off-diagonal entries of
$\mathsf{O}^{(0)[2]}$ have vanishing Gaussian average, since $\langle u_\nu u_\mu\rangle=0$ for
$\nu\ne\mu$, and
$\langle\mathsf{O}^{(0)[2]}_{\nu\nu}\rangle=\sum_{\mu\ne\nu}c_{\mu\nu}\cdot\frac12=c_\nu/2$; thus
$\langle\mathsf{O}^{(0)[2]}_{\nu\mu}\rangle=\delta_{\nu\mu}c_\nu/2$, and therefore
\begin{align*}
\langle P^{[2]}\rangle&=\dim\mathfrak{g}_q\sum_\nu4q^2\phi_\nu\,c_\nu/2,\\
T_4&=\frac{\sigma^2}{2}\sigma^{-3}\bigl[\mathsf{n}\langle P^{[2]}\rangle+O(\sigma^{-1/2})\bigr]\\
&=\mathsf{n}q^2\dim\mathfrak{g}_q\,\sigma^{-1}\sum_\nu\phi_\nu c_\nu+O(\sigma^{-3/2}).
\end{align*}

\emph{The term $T_5$.} By item (y1) in \eqref{10.6:eq-zero-momentum-symbols-a},
$\operatorname{tr}\mathsf{O}^{(0)}$ has principal part
$\dim\mathfrak{g}_q\sum_\nu\sum_{\mu\ne\nu}c_{\mu\nu}\kappa_\mu^2
=\dim\mathfrak{g}_q\sum_\mu c_\mu\kappa_\mu^2$. Inserting the principal parts of
$\varepsilon_{\rho\rho'}$, $\varepsilon_\lambda$, the part of $\mathsf{a}_2(\sigma)$
proportional to $\sigma^3$ has principal part
$\frac{q^2}{8}\cdot\frac{\sigma^3}{3}\cdot2\cdot4\sum_\rho(\mathsf{f}\kappa)_\rho^2$, where
$(\mathsf{f}\kappa)_\rho=\sum_\lambda\mathsf{f}_{\rho\lambda}\kappa_\lambda$, homogeneous
of degree two, and the part proportional to $\sigma^2$ has principal part
$-\frac{q^2}{8}\cdot\frac{\sigma^2}{2}\cdot4\sum_{\rho\lambda}\mathsf{f}^2_{\rho\lambda}$,
of degree zero; that is,
\begin{equation*}
\mathsf{a}_2\to\frac{q^2}{4}\Bigl[\frac{4\sigma^3}{3}\sum_\rho(\mathsf{f}\kappa)_\rho^2
-\sigma^2\cdot2\Sigma_{\mathsf{f}}\Bigr],\qquad
\sum_\rho(\mathsf{f}\kappa)_\rho^2=\sum_{\lambda\lambda'}(\mathsf{f}^{\top}\mathsf{f})_{\lambda\lambda'}
\kappa_\lambda\kappa_{\lambda'},
\end{equation*}
where $\sum_{\rho\lambda}\mathsf{f}^2_{\rho\lambda}=2\Sigma_{\mathsf{f}}$,
$(\mathsf{f}^{\top}\mathsf{f})_{\mu\mu}=\sum_\rho\mathsf{f}^2_{\rho\mu}=\phi_\mu$ and
$\operatorname{tr}\mathsf{f}^{\top}\mathsf{f}=2\Sigma_{\mathsf{f}}$. We apply
Lemma~\ref{10.6:lem-heat-asymptotics} separately to the $\sigma^3$-part (where the product with
$\operatorname{tr}\mathsf{O}^{(0)}$ has $m=4$, giving $\sigma^3\sigma^{-4}=\sigma^{-1}$) and to the
$\sigma^2$-part ($m=2$, giving $\sigma^2\sigma^{-3}=\sigma^{-1}$). With the fourth moments,
\begin{equation*}
\langle\sum_{\lambda\lambda'}(\mathsf{f}^{\top}\mathsf{f})_{\lambda\lambda'}u_\lambda u_{\lambda'}
u_\mu^2\rangle
=\frac14\bigl(\operatorname{tr}\mathsf{f}^{\top}\mathsf{f}+2(\mathsf{f}^{\top}\mathsf{f})_{\mu\mu}\bigr)
=\frac14(2\Sigma_{\mathsf{f}}+2\phi_\mu),
\end{equation*}
and with $\langle u_\mu^2\rangle=\frac12$, we obtain
$T_5=\mathsf{n}q^2\dim\mathfrak{g}_q\,\sigma^{-1}\,[\cdot]+O(\sigma^{-3/2})$ with
\begin{equation*}
[\cdot]=\sum_\mu c_\mu\Bigl\{\frac13\cdot\frac14(2\Sigma_{\mathsf{f}}+2\phi_\mu)
-\frac12\,\Sigma_{\mathsf{f}}\cdot\frac12\Bigr\}
=\sum_\mu c_\mu\Bigl(\frac{\phi_\mu}{6}-\frac{\Sigma_{\mathsf{f}}}{12}\Bigr).
\end{equation*}

\emph{The general formula.} The computation uses the weights only as coefficients in items (y1),
(y2) of \eqref{10.6:eq-zero-momentum-symbols-a};
for arbitrary symmetric orientation weights $c_{\mu\nu}=c_{\nu\mu}$ (the symmetry is used in
$T_2$ and in the regrouping of $\operatorname{tr}\mathsf{O}^{(0)}$ in $T_5$), adding
$T_1,\dots,T_5$ gives
\begin{equation}\label{10.6:eq-heat-coefficient-a}
\begin{split}
\int_{\mathrm{BZ}_1}\bar{d}\kappa\,e^{-\sigma\varepsilon}\operatorname{tr}\{\dots\}
={}&\mathsf{n}q^2\dim\mathfrak{g}_q\,\sigma^{-1}\Bigl[-6\sum_{\mu<\nu}c_{\mu\nu}\mathsf{f}^2_{\mu\nu}
+\frac76\sum_\nu c_\nu\phi_\nu-\frac{\Sigma_{\mathsf{f}}}{12}\sum_\nu c_\nu\Bigr]\\
&+O(\sigma^{-3/2}),
\end{split}
\end{equation}
where $\{\dots\}$ is the brace in \eqref{10.6:eq-proper-time-b} and $1+\frac16=\frac76$.

\emph{Specialisation.} For $c_{\mu\nu}=h_\mu+h_\nu$ with $\sum_\mu h_\mu=0$:
$c_\nu=\sum_{\mu\ne\nu}(h_\mu+h_\nu)=-h_\nu+3h_\nu=2h_\nu$, so $\sum_\nu c_\nu=0$; and, as
$\mathsf{f}_{\mu\mu}=0$,
$\sum_{\mu<\nu}c_{\mu\nu}\mathsf{f}^2_{\mu\nu}=\sum_{\mu\ne\nu}h_\mu\mathsf{f}^2_{\mu\nu}
=\sum_\mu h_\mu\phi_\mu$. The bracket in \eqref{10.6:eq-heat-coefficient-a} becomes
\begin{equation*}
\Bigl(-6+2+\frac13\Bigr)\sum_\mu h_\mu\phi_\mu=-\frac{11}{3}\sum_\mu h_\mu\phi_\mu
=-\frac{11}{3}\cdot2\,\mathsf{t}^{\mathbf{3}}(h;\mathsf{f}),
\end{equation*}
by the identity $\mathsf{t}^{\mathbf{3}}(h;\mathsf{f})=\frac12\sum_\mu h_\mu\phi_\mu$, fixed with
the notation for lattices, constant curvatures and affine fields in
Notation~\ref{10.6:not-affine-fields}. Multiplying by
$\frac12$, summing over the finitely many
sectors $q$, and using the sector identity for $C_{\mathfrak{g}}$,
$\sum_q q^2\dim\mathfrak{g}_q
=\operatorname{tr}_{\mathfrak{g}}(\operatorname{ad}^2_{t_0})=C_{\mathfrak{g}}$, and
$\mathsf{n}=1/(16\pi^2)$,
\begin{equation*}
\begin{split}
\mathsf{h}(\sigma)&=\frac12\sum_q q^2\dim\mathfrak{g}_q\,\mathsf{n}\,\sigma^{-1}
\Bigl(-\frac{22}{3}\Bigr)\mathsf{t}^{\mathbf{3}}
+O(\sigma^{-3/2})\\
&=-\frac{C_{\mathfrak{g}}}{2}\cdot\frac{22}{3}\cdot\frac{1}{16\pi^2}\,\sigma^{-1}
\mathsf{t}^{\mathbf{3}}+O(\sigma^{-3/2}),
\end{split}
\end{equation*}
and $\frac{C_{\mathfrak{g}}}{2}\cdot\frac{22}{3}\cdot\frac{1}{16\pi^2}
=\frac{11C_{\mathfrak{g}}}{48\pi^2}$. Since every $O(\sigma^{-3/2})$ is bounded by
$C|h||\mathsf{f}|^2\sigma^{-3/2}$, this is \eqref{10.6:eq-heat-coefficient-1}.
\end{proof}

\begin{remark}[Consistency checks of the heat coefficient]\label{10.6:rem-consistency-checks}
The bracket in the large-time heat coefficient \eqref{10.6:eq-heat-coefficient-a},
\begin{equation}\label{10.6:eq-consistency-checks-1}
-6\sum_{\mu<\nu}c_{\mu\nu}\mathsf{f}^2_{\mu\nu}
+\frac{7}{6}\sum_{\nu}c_\nu\phi_\nu
-\frac{\Sigma_{\mathsf{f}}}{12}\sum_{\nu}c_\nu ,
\end{equation}
passes three consistency checks.
\begin{itemize}
\item[($\alpha$)] For the constant orientation weight $c_{\mu\nu}\equiv 1$ the bracket
\eqref{10.6:eq-consistency-checks-1} vanishes. Thus the trace channel $\mathbf{1}$ carries no
drift. This agrees with the lemma on the one-loop drift numbers of the tensor normalisations, in
the section of Lemma~\ref{10.4:lem-normalisation-decay}, which gives $\mathsf{m}_{\mathbf{1}}=0$
for the trace channel.
\item[($\beta$)] A derivation by operators for constant $\mathsf{f}$ in the continuum reproduces
the bracket \eqref{10.6:eq-consistency-checks-1} term by term.
\item[($\gamma$)] The corresponding operator computation for the coupling itself produces the
rational number $\frac{11}{3}$. This is the rational number that appears in the heat coefficient.
\end{itemize}
\end{remark}

\begin{proof}
($\alpha$) Let $c_{\mu\nu}=1$ for all $\mu<\nu$. Each index $\nu$ has three partners, so the row
sums are $c_\nu=3$, and $\sum_\nu c_\nu=12$. Moreover
$\sum_{\mu<\nu}c_{\mu\nu}\mathsf{f}^2_{\mu\nu}=\Sigma_{\mathsf{f}}$. Each unordered pair
$\{\mu,\nu\}$ occurs once in $\phi_\mu$ and once in $\phi_\nu$, so
$\sum_\nu\phi_\nu=2\Sigma_{\mathsf{f}}$. Hence \eqref{10.6:eq-consistency-checks-1} equals
\[
-6\Sigma_{\mathsf{f}}+\frac{7}{6}\cdot 3\cdot 2\Sigma_{\mathsf{f}}
-\frac{\Sigma_{\mathsf{f}}}{12}\cdot 12=(-6+7-1)\Sigma_{\mathsf{f}}=0 .
\]
This agrees with the vanishing $\mathsf{m}_{\mathbf{1}}=0$ of the drift in the trace channel.

($\beta$) Work in the continuum with constant curvature $\mathsf{f}$ on the charge-$q$ sector. Let
$D$ denote the covariant derivative there and $D^2$ its Laplacian. Let $\Theta$ denote the
antisymmetric matrix $s\,q\,\mathsf{f}$ at proper time $s$. Let
$\theta_\mu=\sum_\nu\Theta^2_{\mu\nu}$ be its row sums of squares and $\Sigma_\Theta$ its total
square, formed as $\phi_\mu$ and $\Sigma_{\mathsf{f}}$ are formed from $\mathsf{f}$. Conjugation
by the heat semigroup rotates the derivative:
\[
e^{sD^2}D_\nu e^{-sD^2}=\bigl(e^{\mp 2i\Theta}D\bigr)_\nu .
\]
From this,
\begin{equation}\label{10.6:eq-consistency-checks-2}
\begin{gathered}
\bigl(D_\mu e^{sD^2}D_\nu^{\dagger}\bigr)(x,x)
=\frac{\mathsf{n}_s}{2s}\Bigl[\frac{\Theta}{\sin\Theta}\,e^{\mp i\Theta}\Bigr]_{\mu\nu},\\
\mathsf{n}_s:=e^{sD^2}(x,x)=(4\pi s)^{-2}\Bigl(1+\frac{\operatorname{tr}\Theta^2}{12}+\dots\Bigr),
\end{gathered}
\end{equation}
where the dots denote terms of higher order in $\Theta$. Let $W$ denote the
$\mathfrak{g}_q$-valued fluctuation one-form, so that $D_\mu W_\nu-D_\nu W_\mu$ is its linearised
curvature. Combine \eqref{10.6:eq-consistency-checks-2} with the spin factor $e^{\mp 2i\Theta}$.
Then the expectation of $|D_\mu W_\nu-D_\nu W_\mu|^2$ together with its spin term has
$\Theta^2$-part
\[
(4\pi s)^{-2}s^{-1}\Bigl[\frac{7}{6}(\theta_\mu+\theta_\nu)-6\Theta^2_{\mu\nu}
-\frac{\Sigma_\Theta}{6}\Bigr].
\]
Sum this over $\mu<\nu$ against $c_{\mu\nu}$. Since $c_\nu=\sum_{\mu\ne\nu}c_{\mu\nu}$, we have
$\sum_{\mu<\nu}c_{\mu\nu}(\theta_\mu+\theta_\nu)=\sum_\nu c_\nu\theta_\nu$ and
$\sum_{\mu<\nu}c_{\mu\nu}=\frac12\sum_\nu c_\nu$. The sum therefore equals
\[
(4\pi s)^{-2}s^{-1}\Bigl[-6\sum_{\mu<\nu}c_{\mu\nu}\Theta^2_{\mu\nu}
+\frac{7}{6}\sum_\nu c_\nu\theta_\nu-\frac{\Sigma_\Theta}{12}\sum_\nu c_\nu\Bigr].
\]
This has the form of \eqref{10.6:eq-consistency-checks-1} with $\Theta$ in place of $\mathsf{f}$.
Summed against $c_{\mu\nu}$, the operator derivation thus gives the bracket of
\eqref{10.6:eq-heat-coefficient-a} term by term.

($\gamma$) Consider the logarithms of the determinants of two operators on the charge-$q$
sector. The vector operator is the fluctuation operator on $\mathfrak{g}_q$-valued one-forms. The
ghost operator is the covariant Laplacian $-D^2$ on $\mathfrak{g}_q$-valued functions. For the
coupling itself, the $\Theta^2$-part of $-\operatorname{Tr}\log(\text{vector})
+2\operatorname{Tr}\log(\text{ghost})$ is
\[
(4\pi s)^{-2}\Bigl[4\Sigma_\Theta-\frac{\Sigma_\Theta}{3}\Bigr]
=(4\pi s)^{-2}\,\frac{11}{3}\,\Sigma_\Theta .
\]
The rational number $\frac{11}{3}$ found here is the one that enters the large-time heat
coefficient of Lemma~\ref{10.6:lem-heat-coefficient} through the constant
$11C_{\mathfrak{g}}/48\pi^2$.
\end{proof}

\begin{theorem}[Linear drift of the cut-free one-loop sums]\label{10.6:thm-one-loop-drift}
Let Ba{\l}aban's renormalisation step, its fluctuation integral and its propagators in a background
field be those of \cite[(0.5)--(0.16), (1.4)--(1.5), (2.1)--(2.12)]{B-CMP109} and
\cite[(3.3)--(3.10), (3.19)--(3.26)]{B-CMP99b}. Let the flat bounds be those of
\cite[Proposition~1.1, (1.89)--(1.90)]{B-CMP95} and \cite[(1.42)--(1.45), (1.99)--(1.101)]{B-CMP95}.
Let the one-step gap at $U=\mathbf{1}$ be that of \cite{B-CMP99b}. Let
$\mathsf{m}^{\circ}_{\mathbf{3}}(l)$, $l\ge 1$, and $n_{\mathbf{3}}(k)$, $k\ge 0$, be the cut-free
one-loop numbers and their partial sums (see Notation~\ref{10.6:not-cut-free-numbers}).
Assume the following.
\begin{enumerate}
\item[(a)] The one-loop coupling-flow statement holds, namely: fix a basis $(t_a)$ of $\mathfrak{g}$
that is orthonormal for the trace form, with $\operatorname{tr}\mathbf{1}=1$; define $C_{\mathfrak{g}}$
by $\sum_a[t_a,[t_a,X]]=C_{\mathfrak{g}}X$ and put
$c_0:=(11C_{\mathfrak{g}}/24\pi^2)\log L$; there are numbers $s_i$, $i\ge0$, with $|s_i|\le c_2$
such that, for all $0\le i<j$, the one-loop coefficients $\beta^{0}_{l}$ of the coupling satisfy
\begin{equation}\label{10.6:eq-one-loop-drift-1}
\sum_{l=i+1}^{j}\beta^{0}_{l}=(j-i)\,c_0+s_i-s_j .
\end{equation}
\item[(b)] For the orientation-weighted action $\mathsf{A}^{\mathbf{3}}$, the number
$\mathsf{m}_{\mathbf{3}}(l)$ is minus the typed
second-moment coefficient $\beta^{\mathbf{3}}$ of the total one-loop functional at step $l$. The
coefficient $\beta^{\mathbf{3}}$ reads totals, and
$\dim\operatorname{Hom}_{W_4}(\mathsf{V}_{\mathbf{3}},\operatorname{Sym}^2\Lambda^2)=1$.
\item[(c)] For $L$ odd, a linear potential $A$ satisfies
$\partial^{*}\partial A=0$, where $\partial$ is the lattice derivative of potentials and
$\partial^{*}$ its adjoint, and reproduces itself under the symmetric block averages; it is
therefore the minimiser of the step's variational problem, with zero multiplier. The kernels of the
step are local in the background: their second derivatives $\partial_c\partial_{c'}$ in two
background variables $c,c'$ are bounded by
$C_{\mathrm{loc}}e^{-\delta_{\mathrm{loc}}\operatorname{diam}}$, with constants
$C_{\mathrm{loc}},\delta_{\mathrm{loc}}>0$.
\end{enumerate}
Then there is $c^{(3)}<\infty$, depending only on $L$ and on the gauge group, such that for every
$k\ge 0$
\begin{equation}\label{10.6:eq-one-loop-drift-2}
n_{\mathbf{3}}(k)=k\,c_0-s^{(3)}_k,\qquad
c_0=\frac{11C_{\mathfrak{g}}}{24\pi^2}\log L,\qquad s^{(3)}_0=0,\qquad |s^{(3)}_k|\le c^{(3)} .
\end{equation}
Hence, for all $0\le i<j$,
\begin{equation*}
\sum_{l=i+1}^{j}\mathsf{m}^{\circ}_{\mathbf{3}}(l)=(j-i)\,c_0+s^{(3)}_i-s^{(3)}_j .
\end{equation*}
This is the form of \eqref{10.6:eq-one-loop-drift-1}, with the same $c_0$. Together with
hypothesis~(a),
\begin{equation}\label{10.6:eq-one-loop-drift-a}
\Bigl|\sum_{l=i+1}^{j}\bigl[\mathsf{m}^{\circ}_{\mathbf{3}}(l)-\beta^{0}_{l}\bigr]\Bigr|
\le 2c^{(3)}+2c_2 .
\end{equation}
\end{theorem}

\begin{proof}
Hypotheses (b) and (c) and the published statements listed above are the standing hypotheses of
this section. The proof uses them only through Notation~\ref{10.6:not-cut-free-numbers},
Proposition~\ref{10.6:prop-bare-trace}, Definition~\ref{10.6:def-infrared-families},
Lemma~\ref{10.6:lem-regulator-change}, Proposition~\ref{10.6:prop-regulator-admissible},
Lemma~\ref{10.6:lem-proper-time} and Lemma~\ref{10.6:lem-heat-coefficient}. Hypothesis (a) is used
only in the last step.

Fix $k\ge1$. Let $N=L^{k}$, so that the bare spacing is $\eta=L^{-k}=N^{-1}$. Fix orientation
weights $h$ and a constant curvature matrix $\mathsf{f}$. All sums $\sum_q$ below run over the
charges $q$ of $\operatorname{ad}t_0$. The sectors $\mathfrak{g}_q$ decompose $\mathfrak{g}$, so
such a sum has at most $\dim\mathfrak{g}$ terms.

We first obtain a trace identity with a bounded remainder. The partial sum is one bare-lattice
trace, by \eqref{10.6:eq-bare-trace-a} of Proposition~\ref{10.6:prop-bare-trace}:
\begin{equation*}
n_{\mathbf{3}}(k)\,\mathsf{t}^{\mathbf{3}}(h;\mathsf{f})
=-\sum_q[\mathsf{f}^2]\tfrac12\operatorname{Tr}_{\square}(G\mathsf{O})
+\sum_q[\mathsf{f}^2]\tfrac12\operatorname{Tr}_{\square}(\mathsf{M}\,QG\mathsf{O}GQ^{\dagger}).
\end{equation*}
Here the regulated propagator is $G=(H+\mathsf{W}_B)^{-1}$, $\mathsf{W}_B$ being Ba{\l}aban's
infrared term. Hence, by the definition of the regulated half-trace $\tau$ in
Definition~\ref{10.6:def-infrared-families}, the first sum on the right is
$-\sum_q\tau[\mathsf{W}_B]$.

The regulator $\mathsf{W}_B$ is admissible by Proposition~\ref{10.6:prop-regulator-admissible}(i).
Its constants are $\gamma=\gamma_0$, the constant of \cite[Proposition~1.1, (1.90)]{B-CMP95}, and
$C_v$, $C_1$ depending only on $L$ and on the gauge group. The regulator $\bar{\mathsf{W}}$ is
admissible by Lemma~\ref{10.6:lem-proper-time}, with $\gamma=\tfrac12$ and constants depending only
on the gauge group.

Condition (a1) of Definition~\ref{10.6:def-infrared-families} is preserved when $\gamma$ is
decreased. Conditions (a2) and (a3) are preserved when $C_v$ and $C_1$ are increased. Hence both
regulators are admissible with the common constants
\[
\gamma=\min\{\gamma_0,\tfrac12\}
\]
and with the larger of the respective $C_v$ and of the respective $C_1$.

We apply Lemma~\ref{10.6:lem-regulator-change} with $\mathsf{W}_1=\mathsf{W}_B$ and
$\mathsf{W}_0=\bar{\mathsf{W}}$. It gives
$|\tau[\mathsf{W}_B]-\tau[\bar{\mathsf{W}}]|\le C_{\mathrm{reg}}|h||\mathsf{f}|^2$. Here
$C_{\mathrm{reg}}$ depends only on these common constants and on $C_H$, and not on $\eta$ or $k$.
By Proposition~\ref{10.6:prop-regulator-admissible}(ii), uniformly in $k$,
\[
|[\mathsf{f}^2]\tfrac12\operatorname{Tr}_{\square}(\mathsf{M}\,QG\mathsf{O}GQ^{\dagger})|
\le C_S|h||\mathsf{f}|^2 .
\]
By \eqref{10.6:eq-proper-time-b} of Lemma~\ref{10.6:lem-proper-time}, $\sum_q\tau[\bar{\mathsf{W}}]$
is the heat integral $\int_0^{N^2}\mathsf{h}(\sigma)\,d\sigma$. Collecting these facts,
\begin{equation}\label{10.6:eq-one-loop-drift-3}
n_{\mathbf{3}}(k)\,\mathsf{t}^{\mathbf{3}}(h;\mathsf{f})
=-\int_0^{N^2}\mathsf{h}(\sigma)\,d\sigma+\mathsf{r}_k(h,\mathsf{f}),\qquad
|\mathsf{r}_k|\le(\dim\mathfrak{g})(C_{\mathrm{reg}}+C_S)|h||\mathsf{f}|^2,
\end{equation}
where
\[
\mathsf{r}_k
=-\sum_q\bigl(\tau[\mathsf{W}_B]-\tau[\bar{\mathsf{W}}]\bigr)
+\sum_q[\mathsf{f}^2]\tfrac12\operatorname{Tr}_{\square}(\mathsf{M}\,QG\mathsf{O}GQ^{\dagger}).
\]

Next we evaluate the heat integral. By Lemma~\ref{10.6:lem-heat-coefficient}, $\mathsf{h}$ is
continuous on $[0,\infty[$, so the integral exists. Split it at $\sigma=1$, which is allowed since
$N^2\ge1$. On $[0,1]$, $|\mathsf{h}(\sigma)|\le C|h||\mathsf{f}|^2$, so
$|\int_0^1\mathsf{h}|\le C|h||\mathsf{f}|^2$. On $[1,N^2]$, write
\[
\mathsf{h}(\sigma)=-(11C_{\mathfrak{g}}/48\pi^2)\,\sigma^{-1}\mathsf{t}^{\mathbf{3}}(h;\mathsf{f})
+e(\sigma),\qquad |e(\sigma)|\le C_{\mathrm{HK}}|h||\mathsf{f}|^2\sigma^{-3/2} .
\]
We have $\int_1^{N^2}\sigma^{-1}d\sigma=\log N^2$ and
$\int_1^{N^2}\sigma^{-3/2}d\sigma\le\int_1^{\infty}\sigma^{-3/2}d\sigma=2$. Therefore
\begin{equation}\label{10.6:eq-one-loop-drift-4}
-\int_0^{N^2}\mathsf{h}(\sigma)\,d\sigma
=\frac{11C_{\mathfrak{g}}}{48\pi^2}\,\log N^2\cdot\mathsf{t}^{\mathbf{3}}(h;\mathsf{f})+\mathsf{r}'_k,
\qquad |\mathsf{r}'_k|\le(C+2C_{\mathrm{HK}})|h||\mathsf{f}|^2 .
\end{equation}

Now we combine and specialise. Since $N=L^k$, we have $\log N^2=2k\log L$, and so
\[
\frac{11C_{\mathfrak{g}}}{48\pi^2}\cdot 2k\log L=k\,c_0 .
\]
Substituting \eqref{10.6:eq-one-loop-drift-4} into \eqref{10.6:eq-one-loop-drift-3} gives
\[
\bigl(n_{\mathbf{3}}(k)-k\,c_0\bigr)\,\mathsf{t}^{\mathbf{3}}(h;\mathsf{f})
=\mathsf{r}_k+\mathsf{r}'_k .
\]
Evaluate at $h=(1,1,-1,-1)$, $\mathsf{f}_{12}=1$ and $\mathsf{f}_{\mu\nu}=0$ for all other
$\mu<\nu$. There $\mathsf{t}^{\mathbf{3}}(h;\mathsf{f})=1$, and $|h||\mathsf{f}|^2$ is a fixed
number. With $|h||\mathsf{f}|^2$ evaluated at this choice,
\[
|n_{\mathbf{3}}(k)-k\,c_0|\le c^{(3)}:=|h||\mathsf{f}|^2\,
\bigl[(\dim\mathfrak{g})(C_{\mathrm{reg}}+C_S)+C+2C_{\mathrm{HK}}\bigr],
\]
a constant depending only on $L$ and on the gauge group, and not on $k$.

Put $s^{(3)}_k:=k\,c_0-n_{\mathbf{3}}(k)$. For $k\ge1$ the bound just proved is
$|s^{(3)}_k|\le c^{(3)}$. For $k=0$, $n_{\mathbf{3}}(0)=0$ gives $s^{(3)}_0=0$. This proves
\eqref{10.6:eq-one-loop-drift-2}.

For $0\le i<j$, the definition of the partial sums gives
\[
\sum_{l=i+1}^{j}\mathsf{m}^{\circ}_{\mathbf{3}}(l)=n_{\mathbf{3}}(j)-n_{\mathbf{3}}(i)
=(j\,c_0-s^{(3)}_j)-(i\,c_0-s^{(3)}_i)=(j-i)\,c_0+s^{(3)}_i-s^{(3)}_j .
\]
Subtracting \eqref{10.6:eq-one-loop-drift-1} of hypothesis (a), in which $c_0$ is the same
constant, gives
\[
\sum_{l=i+1}^{j}\bigl[\mathsf{m}^{\circ}_{\mathbf{3}}(l)-\beta^0_l\bigr]
=(s^{(3)}_i-s_i)-(s^{(3)}_j-s_j),
\]
whose absolute value is at most $2c^{(3)}+2c_2$. This is \eqref{10.6:eq-one-loop-drift-a}.
\end{proof}

\begin{corollary}[Identification and bound of the one-loop number]\label{10.6:cor-one-loop-number}
Assume the hypotheses of Theorem~\ref{10.6:thm-one-loop-drift}, that is, the inputs listed there
from which the linear drift of the cut-free one-loop sums is derived, and let $c_0$, $c^{(3)}$ and
$s^{(3)}_k$ be the constants and the bounded sequence of that theorem. Then the following hold.
\begin{enumerate}
\item[(i)] For every $k\ge 0$,
\begin{equation}\label{10.6:eq-one-loop-number-1}
\mathsf{m}^{\circ}_{\mathbf{3}}(k+1)
=-\beta^{\mathbf{3}}\Bigl[\,W\mapsto \tfrac12\operatorname{Tr}\bigl(C^{(k)}\,
\nabla^2(\mathsf{A}^{\mathbf{3}}\circ U_k)\bigr)\Bigr].
\end{equation}
The term $\langle\mathcal{S}_1;V_1\rangle_0$ has $\beta^{\mathbf{3}}=0$. Here, as in
Lemma~\ref{10.6:lem-half-trace}, $\mathcal{S}_1$ and $\mathcal{S}_2$ are the coefficients of
$\mathfrak{q}$ and $\mathfrak{q}^2$ in the vertex at $\mathfrak{q}B$, $V_1$ is the coefficient of
$\mathfrak{q}$ in the fluctuation potential at $\mathfrak{q}B$ with the remaining source parameter
set to zero, and $\langle\cdot\rangle_0$,
$\langle\cdot\,;\cdot\rangle_0$ are the normalised Gaussian expectation of covariance $C^{(k)}$
without the cut and its truncated form; the configuration $W$, the fluctuation field $B$ and the
linear map $C$ are those of Lemma~\ref{10.6:lem-half-trace}. In particular
$\mathsf{m}^{\circ}_{\mathbf{3}}$ does not depend on the one-loop families of the levels $\le k$.
Moreover, computed without the cut, that is, with the expectation $\langle\cdot\rangle_0$, the
following three numbers are one and the same: $\mathsf{m}_{\mathbf{3}}(k+1)$; the
traceless-diagonal one-loop coefficient at step $k+1$ of the one-loop form of the tensor-sourced
coupling increment, in the sign convention of that form; and $-g_k^{-2}$ times the coefficient
$\beta^{\mathbf{3}}$ of the one-loop term of the tensor source at level $k+1$, namely of the term
linear in the source strength $\mathsf{k}$ of the logarithm of the one-loop determinant factor.
\item[(ii)] $\sup_{l\ge 1}\bigl|\mathsf{m}^{\circ}_{\mathbf{3}}(l)\bigr|\le c_0+2c^{(3)}$.
\end{enumerate}
\end{corollary}

\begin{proof}
(i) Fix $k\ge 0$ and apply the half-trace lemma (Lemma~\ref{10.6:lem-half-trace}) with $j=k$.
By the definition of the one-loop drift numbers, the number $\mathsf{m}^{\circ}_{\mathbf{3}}(k+1)$
is $-\beta^{\mathbf{3}}$ of the sum of the one-loop family
$\mathsf{L}^{\mathbf{3}}_{k+1}(1;\cdot)$ at level $k+1$ over its localisation domains, computed
with the Gaussian expectation without the cut. By the identification of this sum as the one-loop
response, it is $\langle\mathcal{S}_2\rangle_0+\langle\mathcal{S}_1;V_1\rangle_0$.

By part (a) of the lemma, this sum equals
\[
\tfrac12\operatorname{Tr}\bigl[\mathcal{G}_P^{(k+1)}(W)\,\nabla^2\mathcal{S}_k(V^{(k)}(W))\bigr]
-\langle\nabla\mathcal{S}_k,\mathsf{w}_1\rangle+\langle\nabla\mathcal{S}_k,\mathsf{w}_2\rangle .
\]
By part (b) of the lemma, the last two terms have $\beta^{\mathbf{3}}=0$. Hence $\beta^{\mathbf{3}}$
of the sum is $\beta^{\mathbf{3}}$ of the half trace. Expanding the vertex in the variable $B'$ of
Lemma~\ref{10.6:lem-half-trace}, whose term of first order in $\mathfrak{q}$ is
$\mathfrak{q}CB$, we use that $CB$ has covariance $C\,C^{(k)}C^{*}$ under $\langle\cdot\rangle_0$; by
cyclicity of the trace the half trace therefore equals
$\tfrac12\operatorname{Tr}\bigl(C^{(k)}\,C^{*}\nabla^2\mathcal{S}_k(V^{(k)}(W))\,C\bigr)$, so that,
by the form of the vertex, the half trace equals
$\tfrac12\operatorname{Tr}\bigl(C^{(k)}\nabla^2(\mathsf{A}^{\mathbf{3}}\circ U_k)\bigr)$ evaluated at
the configuration $W$, all objects being those of level $k$; hence
\eqref{10.6:eq-one-loop-number-1}.

Since $\mathcal{S}_1(B)=\langle\nabla\mathcal{S}_k,CB\rangle$ is linear in $B$, we have
$\langle\mathcal{S}_1\rangle_0=0$ and hence
$\langle\mathcal{S}_1;V_1\rangle_0=\langle\nabla\mathcal{S}_k,\mathsf{w}_2\rangle$, with
$\mathsf{w}_2=\langle CB\cdot V_1(B)\rangle_0$; this term has $\beta^{\mathbf{3}}=0$ by part (b).
Everything that $V_1$ contains, including the gradients of the older one-loop terms, those of the
levels $\le k$, enters the decomposition of part (a) through $\mathsf{w}_2$ only; hence
these terms do not contribute to \eqref{10.6:eq-one-loop-number-1}.

We turn to the three numbers. In the sign convention of the one-loop form of the tensor-sourced
coupling increment, its traceless-diagonal one-loop coefficient at step $k+1$ is
$\mathsf{m}_{\mathbf{3}}(k+1)$. The third number is defined not through the one-loop response but
through the one-loop term of the tensor source at level $k+1$: the term linear in the source
strength $\mathsf{k}$ of the logarithm of the one-loop determinant factor has as coefficient $g_k^2$
times the half trace of $C^{(k)}$ against the Hessian of the tensor source, that is, $g_k^2$ times
$\tfrac12\operatorname{Tr}$ of their product, and the third number is $-g_k^{-2}$ times
$\beta^{\mathbf{3}}$ of this term. This half trace is
$\tfrac12\operatorname{Tr}\bigl(C^{(k)}\,C^{*}\nabla^2\mathcal{S}_k\,C\bigr)$, identified
above with $\tfrac12\operatorname{Tr}\bigl(C^{(k)}\nabla^2(\mathsf{A}^{\mathbf{3}}\circ U_k)\bigr)$.
Hence $g_k^{-2}$ times $\beta^{\mathbf{3}}$ of the one-loop term of the tensor source equals
\[
\beta^{\mathbf{3}}\Bigl[\,W\mapsto \tfrac12\operatorname{Tr}\bigl(C^{(k)}\,
\nabla^2(\mathsf{A}^{\mathbf{3}}\circ U_k)\bigr)\Bigr],
\]
which by \eqref{10.6:eq-one-loop-number-1} is $-\mathsf{m}^{\circ}_{\mathbf{3}}(k+1)$; that is, the
third number is the cut-free number $\mathsf{m}^{\circ}_{\mathbf{3}}(k+1)$ of
\eqref{10.6:eq-one-loop-number-1}. Computed without the cut, $\mathsf{m}_{\mathbf{3}}(k+1)$ is by
definition the cut-free number $\mathsf{m}^{\circ}_{\mathbf{3}}(k+1)$, and hence so is the
traceless-diagonal one-loop coefficient; thus, computed without the cut, the three numbers
coincide.

(ii) Let $l\ge 1$. Apply the second assertion of Theorem~\ref{10.6:thm-one-loop-drift} with
$i=l-1$ and $j=l$. The sum on its left consists of the single term
$\mathsf{m}^{\circ}_{\mathbf{3}}(l)$, so
\[
\mathsf{m}^{\circ}_{\mathbf{3}}(l)=c_0+s^{(3)}_{l-1}-s^{(3)}_{l}.
\]
The theorem gives $c_0=(11C_{\mathfrak{g}}/24\pi^2)\log L$, which is positive since $L>1$ and
$C_{\mathfrak{g}}>0$, and $|s^{(3)}_k|\le c^{(3)}$ for every $k\ge 0$; we apply the latter with
$k=l-1$ and $k=l$. Hence
$|\mathsf{m}^{\circ}_{\mathbf{3}}(l)|\le c_0+2c^{(3)}$ for every $l\ge 1$.
\end{proof}

Before the lemma we recall the one-loop drift number with the field cut. Consider the step from level
$k$ to level $k+1$ and put $T := T_k$. Let $\langle\cdot\rangle^{T}$ denote the expectation with
respect to the normalised measure $d\mu_{C^{(k)}}\prod_b\chi(|B(b)| < T)$, and let
$\langle F;G\rangle^{T}$ denote the corresponding truncated expectation. Write
$\mathcal{S}^{\mathbf{3}}(\mathfrak{q}B) = \mathfrak{q}\,\mathcal{S}_1(B)
+ \mathfrak{q}^2\,\mathcal{S}_2(B) + \dots$ for the expansion in the source strength
$\mathfrak{q}$ of the sourced vertex of the orientation-weighted action. Let $V_1$ be the first
coefficient in $\mathfrak{q}$ of the function $V$ of
Proposition~\ref{10.4:prop-vertex-response}(d),(e), whose proof is incomplete at one step. By the
definition of $\mathsf{m}_{\mathbf{3}}$ and by Proposition~\ref{10.4:prop-vertex-response}(d),(e),
the one-loop drift number with the field cut is
\begin{equation*}
\mathsf{m}_{\mathbf{3}}(k+1) = -\beta^{\mathbf{3}}\big[\langle\mathcal{S}_2\rangle^{T}
+\langle\mathcal{S}_1;V_1\rangle^{T}\big].
\end{equation*}

\begin{lemma}[The cluster expansion with quadratic and linear sources]
\label{10.6:lem-sourced-expansion}
Assume the following three statements.
\begin{enumerate}
\item[(a)] The cluster-expansion lemma accompanying the correlation theorem. Suppose a localised
potential is given on each domain $Y$ as a function of $B$. Suppose it is analytic on
$\{|B| < T \text{ on } Y\}$ and obeys the bounds of that lemma on its size and decay. These are its
format bounds, and the constants appearing in them are called the format constants. Then there is a
family $\Phi$ with
\begin{equation*}
|\Phi(X)| \le E^{*}e^{-\kappa d_{k+1}(X)},
\end{equation*}
analytic in every parameter on which the potential depends analytically with these bounds, and
$\exp\sum\Phi = \mathsf{I}$.
\item[(b)] A further potential of the type admitted by the lemma in (a), adjoined within the
format constants to a potential satisfying the requirements of that lemma, is carried through that
lemma with the same conclusion. This is the device by which
$\tfrac12(1-s)\langle B,\Delta^{(k)}B\rangle$ is carried as one further such potential; here $s$ is
the interpolation parameter of that device and $\Delta^{(k)}$ the level-$k$ fluctuation operator.
\item[(c)] Part (a) of Proposition~\ref{10.4:prop-vertex-response}, whose proof is incomplete at
one step, together with the allocation of the format constants under which one half of them is
reserved for the vertex.
\end{enumerate}
For a localisation domain $Y$ write $[\mathfrak{q}^n]$ for the coefficient of $\mathfrak{q}^n$,
and put
\begin{equation*}
\mathcal{S}_2(Y;B) := [\mathfrak{q}^2]\,\mathcal{S}^{\mathbf{3}}(Y,\mathfrak{q}B),\qquad
V_1(Y;B) := [\mathfrak{q}^1]\,V(Y,\mathfrak{q}B),
\end{equation*}
where $V(Y,\cdot)$ is the localised form of the function $V$ recalled above. Let
$\mathsf{P}(Y;\cdot)$ be a localisation of a component of $h\tilde{C}^{(2)}(CB)(b_0)$, with $b_0$ a
fixed bond, $h$ the orientation weight, $C$ and $\tilde{C}^{(2)}$ the maps entering the proposition
named in (c), and $Y(b_0)$ the localisation domain attached to $b_0$. Let $\xi$ be a fixed vector
with $|\xi| \le 1$ in the space in which $B(b_0)$ takes its values, and let
$\langle\cdot,\cdot\rangle$ be the inner product there. For complex parameters
$\lambda,\lambda',\varepsilon,\varepsilon'$ define the sourced potential
\begin{equation}\label{10.6:eq-sourced-expansion-a}
V^{\mathrm{src}}(Y) := \lambda\,\mathcal{S}_2(Y;B) + \lambda'\,\mathsf{P}(Y;B)
+ \varepsilon\,\langle \xi,B(b_0)\rangle\,\mathbf{1}_{Y = Y(b_0)} + \varepsilon'\,V_1(Y;B).
\end{equation}
Then there are an absolute constant $a \ge 1$ and a constant $c_{\mathrm{src}} > 0$ with the
following property. Put
\begin{equation}\label{10.6:eq-sourced-expansion-1}
\varrho_T := c_{\mathrm{src}}\,T^{-a}.
\end{equation}
For $|\lambda|,|\lambda'|,|\varepsilon|,|\varepsilon'| \le \varrho_T$ the potentials
\eqref{10.6:eq-sourced-expansion-a} use at most the half of the format constants that the
allocation in (c) reserves for the vertex. Statement (a) applies to
\eqref{10.6:eq-sourced-expansion-a} and returns a family $\Phi$ with
$|\Phi(X)| \le E^{*}e^{-\kappa d_{k+1}(X)}$ and $\exp\sum\Phi = \mathsf{I}$. This family is analytic
in $(\lambda,\lambda',\varepsilon,\varepsilon')$ on the polydisc
$|\lambda|,|\lambda'|,|\varepsilon|,|\varepsilon'| < \varrho_T$.
\end{lemma}

\begin{proof}
Statement (a) asks three things of the potential: analyticity on $\{|B| < T \text{ on } Y\}$,
localisation, and the format bounds. We apply it to \eqref{10.6:eq-sourced-expansion-a}. This
potential depends linearly, hence analytically, on $\lambda,\lambda',\varepsilon,\varepsilon'$.
Consequently, once the three requirements are met, the analyticity clause of (a) yields a family
$\Phi$ analytic in these four parameters.

We first bound the two terms obtained as coefficients in $\mathfrak{q}$. Fix $B$ with $|B| < T$ on
$Y$, and let $F(\mathfrak{q})$ stand for $\mathcal{S}^{\mathbf{3}}(Y,\mathfrak{q}B)$ or for
$V(Y,\mathfrak{q}B)$. By part (a) of the proposition named in (c), $F$ is analytic in $\mathfrak{q}$
wherever $|\mathfrak{q}B| < \varepsilon_1$. It is bounded there by the format constants times
$e^{-(1-2\delta)\kappa d_k(Y)}$. Here $\varepsilon_1$ is the radius in the field variable on which
these bounds hold, and $\delta$ is the constant in the decay rate of the format bounds. Define the
radius $\mathsf{r}_k := \varepsilon_1/(2T)$ and take the circle
\begin{equation*}
|\mathfrak{q}| = 2\mathsf{r}_k = \varepsilon_1/T .
\end{equation*}
On this circle $|\mathfrak{q}B| < \varepsilon_1$ wherever $|B| < T$. Cauchy's formula gives, for
$n = 1,2$,
\begin{equation}\label{10.6:eq-sourced-expansion-2}
\big|[\mathfrak{q}^n]F\big|
= \Big|\frac{1}{2\pi i}\oint_{|\mathfrak{q}| = 2\mathsf{r}_k}
\frac{F(\mathfrak{q})}{\mathfrak{q}^{n+1}}\,d\mathfrak{q}\Big|
\le \Big(\frac{T}{\varepsilon_1}\Big)^{n}\max_{|\mathfrak{q}| = 2\mathsf{r}_k}|F(\mathfrak{q})| .
\end{equation}
With $n = 2$ and $F = \mathcal{S}^{\mathbf{3}}(Y,\mathfrak{q}B)$ this bounds $\mathcal{S}_2(Y;B)$.
With $n = 1$ and $F = V(Y,\mathfrak{q}B)$ it bounds $V_1(Y;B)$. In both cases the bound is the
format constants, times $(T/\varepsilon_1)^n$, times $e^{-(1-2\delta)\kappa d_k(Y)}$.

Next we check localisation and analyticity in $B$. Each of the four terms of
\eqref{10.6:eq-sourced-expansion-a} is attached to the domain $Y$:
\begin{itemize}
\item $\mathcal{S}_2(Y;\cdot)$ and $V_1(Y;\cdot)$ through the localised functions
$\mathcal{S}^{\mathbf{3}}(Y,\cdot)$ and $V(Y,\cdot)$;
\item $\mathsf{P}(Y;\cdot)$ by its choice as a localisation;
\item the linear term, which vanishes unless $Y = Y(b_0)$.
\end{itemize}
On $\{|B| < T \text{ on } Y\}$ the functions $\mathcal{S}_2(Y;\cdot)$ and $V_1(Y;\cdot)$ are analytic
in $B$, being Taylor coefficients in $\mathfrak{q}$ of functions analytic for
$|\mathfrak{q}B| < \varepsilon_1$. The linear term is linear in $B(b_0)$. The function
$\mathsf{P}(Y;\cdot)$, being the localisation of a component of $h\tilde{C}^{(2)}(CB)(b_0)$, is
analytic in $B$ on the same set.

The linear term is bounded directly: on $\{|B| < T\}$ it satisfies
$|\langle \xi,B(b_0)\rangle| \le T$, and it is present only on the single domain $Y(b_0)$. The
same Cauchy estimate on the circle $|\mathfrak{q}| = 2\mathsf{r}_k$ bounds $\mathsf{P}(Y;\cdot)$,
the localisation of a component of $h\tilde{C}^{(2)}(CB)(b_0)$, on $\{|B| < T\}$ by the format
constants, times a power of $T$, times $e^{-(1-2\delta)\kappa d_k(Y)}$.

Now we fix $a$ and $c_{\mathrm{src}}$. Let $T^{a_0}$ be the largest power of $T$ occurring in the
bounds of $\mathcal{S}_2(Y;\cdot)$, $V_1(Y;\cdot)$ and $\mathsf{P}(Y;\cdot)$. Each of the four
functions enters \eqref{10.6:eq-sourced-expansion-a} multiplied by one of the four parameters.
Take $a \ge \max\{a_0,1\}$, so that the factor $T^{-a}$ in $\varrho_T$ compensates the power
$T^{a_0}$ and the factor $T$ in the bound of the linear term, and take $c_{\mathrm{src}} > 0$ small
against the format constants. Then, for $|\lambda|,|\lambda'|,|\varepsilon|,|\varepsilon'| \le
\varrho_T$, the products
\begin{equation*}
\lambda\,\mathcal{S}_2(Y;B),\qquad \lambda'\,\mathsf{P}(Y;B),\qquad
\varepsilon\,\langle \xi,B(b_0)\rangle\mathbf{1}_{Y=Y(b_0)},\qquad \varepsilon'\,V_1(Y;B)
\end{equation*}
are bounded by the half of the format constants that the allocation in (c) reserves for the vertex,
times $e^{-(1-2\delta)\kappa d_k(Y)}$.

By (b), the four source terms may be adjoined to the vertex as further potentials of the type
admitted by the lemma in (a), in the same way as $\tfrac12(1-s)\langle B,\Delta^{(k)}B\rangle$.
The potential \eqref{10.6:eq-sourced-expansion-a} therefore meets
the three requirements of (a). Applied to it, statement (a) returns the family $\Phi$ with
$|\Phi(X)| \le E^{*}e^{-\kappa d_{k+1}(X)}$ and $\exp\sum\Phi = \mathsf{I}$, analytic in
$(\lambda,\lambda',\varepsilon,\varepsilon')$ on the polydisc
$|\lambda|,|\lambda'|,|\varepsilon|,|\varepsilon'| < \varrho_T$.
\end{proof}

The application of statement (a) made in Lemma~\ref{10.6:lem-sourced-expansion} is of the same kind
as part (a) of the proposition named in (c). The bounds of statement (a) make no use of gauge
invariance of the potential.

\begin{lemma}[Gradient terms drop under the cut measure]\label{10.6:lem-gradient-drop-cut}
Let $\langle\cdot\rangle^{T}$, with $T = T_k$, be the cut Gaussian expectation. Let $\mathcal{S}_1$,
$\mathcal{S}_2$, $V_1$, $C$, $B$, $h$ and $\tilde{C}^{(2)}$ be as in Lemma~\ref{10.6:lem-half-trace},
with $k$ in place of $j$. For a background configuration $W$ on $\mathbb{L}_{k+1}$, the derivatives
$\nabla\mathcal{S}_k$ and $\nabla^2\mathcal{S}_k$ are evaluated at $V^{(k)}(W)$.

Assume the cluster expansion with quadratic and linear sources for the cut measure, that is,
Lemma~\ref{10.6:lem-sourced-expansion} in the form \eqref{10.6:eq-sourced-expansion-a}. Let $\Phi(X)$
denote the family it provides. Then
\begin{equation}\label{10.6:eq-gradient-drop-cut-1}
\langle \mathcal{S}_2\rangle^{T} + \langle \mathcal{S}_1 ; V_1\rangle^{T}
= \tfrac12 \operatorname{Tr}\bigl[\mathcal{G}^{T}\,\nabla^2\mathcal{S}_k\bigr]
- \langle \nabla\mathcal{S}_k , \mathsf{w}_1^{T}\rangle
+ \langle \nabla\mathcal{S}_k , \mathsf{w}_2^{T}\rangle ,
\end{equation}
where
\begin{equation}\label{10.6:eq-gradient-drop-cut-2}
\mathcal{G}^{T} := C\,\langle B B^{\dagger}\rangle^{T} C^{*}, \qquad
\mathsf{w}_1^{T} := \bigl\langle h\tilde{C}^{(2)}(CB)\bigr\rangle^{T}, \qquad
\mathsf{w}_2^{T} := \langle CB\cdot V_1\rangle^{T}.
\end{equation}
Moreover, the last two terms of \eqref{10.6:eq-gradient-drop-cut-1}, regarded as functionals of $W$,
satisfy
\[
\beta^{\mathbf{3}}\bigl[\langle \nabla\mathcal{S}_k , \mathsf{w}_1^{T}\rangle\bigr]
= \beta^{\mathbf{3}}\bigl[\langle \nabla\mathcal{S}_k , \mathsf{w}_2^{T}\rangle\bigr] = 0 .
\]
\end{lemma}

\begin{proof}
The argument is that of Lemma~\ref{10.6:lem-half-trace}\,(a) and (b), with the cut expectation in place
of the Gaussian one. We indicate the points at which the cut enters.

\emph{The identity.} The formulas
\[
\mathcal{S}_1(B) = \langle\nabla\mathcal{S}_k, CB\rangle, \qquad
\mathcal{S}_2(B) = \tfrac12\nabla^2\mathcal{S}_k[CB,CB]
- \langle \nabla\mathcal{S}_k, h\tilde{C}^{(2)}(CB)\rangle
\]
obtained in the proof of Lemma~\ref{10.6:lem-half-trace}\,(a) are identities between functions of $B$.
They do not involve the measure, so they hold here with $k$ in place of $j$.

The cut Gaussian measure is invariant under $B\mapsto -B$. Hence $\langle B\rangle^{T} = 0$, and so
\[
\langle\mathcal{S}_1\rangle^{T} = \langle\nabla\mathcal{S}_k, C\langle B\rangle^{T}\rangle = 0 .
\]
Consequently
\[
\langle\mathcal{S}_1;V_1\rangle^{T} = \langle\mathcal{S}_1V_1\rangle^{T}
= \langle\nabla\mathcal{S}_k,\langle CB\cdot V_1\rangle^{T}\rangle
= \langle\nabla\mathcal{S}_k,\mathsf{w}_2^{T}\rangle .
\]
By linearity and cyclicity of the trace,
\[
\tfrac12\langle\nabla^2\mathcal{S}_k[CB,CB]\rangle^{T}
= \tfrac12\operatorname{Tr}\bigl[C\langle BB^{\dagger}\rangle^{T}C^{*}\nabla^2\mathcal{S}_k\bigr].
\]
Also $\langle\langle\nabla\mathcal{S}_k,h\tilde{C}^{(2)}(CB)\rangle\rangle^{T} =
\langle\nabla\mathcal{S}_k,\mathsf{w}_1^{T}\rangle$. Adding these gives \eqref{10.6:eq-gradient-drop-cut-1}.

\emph{The vanishing.} Each of the last two terms has the form
\[
F_i(W) = \sum_b \langle\mathsf{s}(W)(b),\mathsf{w}(W)(b)\rangle ,
\qquad \mathsf{s}(W) := \nabla\mathcal{S}_k(V^{(k)}(W)), \qquad \mathsf{w} = \mathsf{w}_i^{T},
\quad i = 1,2,
\]
with $b$ over the bonds of $\mathbb{L}_k$.
We apply the lemma that terms linear in a vertex gradient drop
(Lemma~\ref{10.6:lem-gradient-terms-drop}), with its kernels
$S(b,c) := D\mathsf{s}(\mathbf{1})(b)[\delta_c]$ and
$\mathsf{K}_{\mathsf{w}}(b,c) := D\mathsf{w}(\mathbf{1})(b)[\delta_c]$. The first factor $\mathsf{s}$ is
that of Lemma~\ref{10.6:lem-half-trace}\,(b) with $k$ in place of $j$, so every hypothesis of
Lemma~\ref{10.6:lem-gradient-terms-drop} bearing on $\mathsf{s}$, verified there, holds unchanged.

For the second factor, the cut is $\operatorname{Ad}$-invariant and local. Hence $\mathsf{w}_i^{T}$ is gauge
covariant, and $\mathsf{w}_i^{T}(\mathbf{1})$ is an $\operatorname{Ad}$-invariant vector of
$\mathfrak{g}$. It is therefore $0$, because $\mathfrak{g}$ is semisimple. For every $\mathfrak{g}$-valued
function $\lambda$ on the sites, gauge covariance then gives
\[
D\mathsf{w}_i^{T}(\mathbf{1})[d\lambda] = i[\lambda,\mathsf{w}_i^{T}(\mathbf{1})] = 0 ,
\]
in particular on constant fields. Thus $\sum_c\mathsf{K}_{\mathsf{w}}(b,c)\beta(c) = 0$ for every
constant field $\beta$.

It remains to establish the exponential decay of the kernel of $D\mathsf{w}_i^{T}(\mathbf{1})$. Take the
potential of \eqref{10.6:eq-sourced-expansion-a} with $\lambda = \lambda' = 0$ (the strengths of the
quadratic source and of the localised potential $\mathsf{P}$), with the linear source of strength
$\varepsilon$ at the bond $b_0$ and the coupling $\varepsilon'$ to $V_1$. Then
\[
\langle B^{a}(b_0)V_1\rangle^{T}
= \partial_{\varepsilon}\partial_{\varepsilon'}\sum_X\Phi(X)\Big|_{\varepsilon=\varepsilon'=0} .
\]
Here $\Phi(X)$ does not depend on $\varepsilon$ unless $X\ni b_0$. Each term is analytic in the background
on $X$ and satisfies
\[
\bigl|\partial_{\varepsilon}\partial_{\varepsilon'}\Phi(X)\bigr|
\le E^{*}\varrho_T^{-2}e^{-\kappa d_{k+1}(X)} ,
\]
with the constants $E^{*}$ and $\kappa$ of \eqref{10.6:eq-sourced-expansion-a}.
Since $\Phi(X)$ depends on the background only through its values on $X$, a term contributes to the
derivative of $\mathsf{w}_2^{T}(b_0)$ in the direction of a bond $c$ only if $X$ contains both $b_0$ and
$c$. This yields the exponential decay of the kernel of $D\mathsf{w}_2^{T}(\mathbf{1})$. The same argument
applies to $\mathsf{w}_1^{T}$, with the strength $\lambda'$ of the localised potential $\mathsf{P}$ in
\eqref{10.6:eq-sourced-expansion-a} in place of $(\varepsilon,\varepsilon')$.

All hypotheses of Lemma~\ref{10.6:lem-gradient-terms-drop} therefore hold for $\mathsf{w} = \mathsf{w}_1^{T}$
and for $\mathsf{w} = \mathsf{w}_2^{T}$. That lemma gives $\beta^{\mathbf{3}} = 0$ for each of the last two
terms of \eqref{10.6:eq-gradient-drop-cut-1}.
\end{proof}

\begin{lemma}[The cost of the sharp field cutoff]\label{10.6:lem-cutoff-cost}
For a level $k$ let $T_k:=A_1(\log x_k)^{p_0}$ be the field cut at level $k$. Assume the following
four statements.
\begin{enumerate}
\item[(a)] The cluster-expansion lemma of the analysis of the correlation theorem, applied to the
potential with quadratic and linear sources of Lemma~\ref{10.6:lem-sourced-expansion}, display
\eqref{10.6:eq-sourced-expansion-a}: with the absolute exponent $a\ge1$ and the constant
$c_{\mathrm{src}}$ of that lemma, for $|\lambda|,|\lambda'|,|\varepsilon|,|\varepsilon'|\le
\varrho_{T_k}=c_{\mathrm{src}}T_k^{-a}$ it returns a family $\Phi$ with
$|\Phi(X)|\le E^{*}e^{-\kappa d_{k+1}(X)}$, with $\kappa$ the decay constant of the cluster
expansion, analytic in every parameter on which the potential
depends analytically with the same format bounds, and with $\exp\sum_X\Phi(X)=\mathsf{I}$, the
integral of $e^{\sum_YV(Y)}$ against the cut Gaussian measure
$d\mu_{C^{(k)}}\prod_b\chi(|B(b)|<T_k)$.
\item[(b)] \v{S}id\'ak's inequality \cite{Sid67}: for a centred Gaussian vector $(\xi_i)_i$ and
$t_i>0$, $P\bigl(\bigcap_i\{|\xi_i|<t_i\}\bigr)\ge\prod_iP(|\xi_i|<t_i)$.
\item[(c)] The two-constants lemma: if $\varphi$ is holomorphic on a bounded plane domain with
$|\varphi|\le\mathsf{M}_{*}$ there, and $\limsup|\varphi|\le\mathsf{m}_{*}$ at every point of a
part $E$ of its boundary, then for $z$ in the domain
\begin{equation*}
\log|\varphi(z)|\le\omega(z)\log\mathsf{m}_{*}+\bigl(1-\omega(z)\bigr)\log\mathsf{M}_{*},
\end{equation*}
where $\omega(z)$ is the harmonic measure of $E$ in the domain at $z$.
\item[(d)] The lemma of the analysis of the correlation theorem by which the quantities computed on
a torus $\mathbb{T}$ of level $k$ pass to the limit $\mathbb{T}\nearrow\mathbb{Z}^4$.
\end{enumerate}
Let $\mathsf{m}_{\mathbf{3}}(k+1)$ be the one-loop drift number of the traceless-diagonal channel
$\mathbf{3}$ at step $k+1$, computed with the field cut $|B|<T_k$, and let
$\mathsf{m}^{\circ}_{\mathbf{3}}(k+1)$ be its cut-free value, as in
Definition~\ref{10.6:not-cut-free-numbers}.
Then there are constants $C_{\mathrm{cut}},c_{\mathrm{cut}}>0$, depending only on the constants of the
analysis of the correlation theorem and independent of $k$ and of $g$, such that
\begin{equation}\label{10.6:eq-cutoff-cost-1}
\bigl|\mathsf{m}_{\mathbf{3}}(k+1)-\mathsf{m}^{\circ}_{\mathbf{3}}(k+1)\bigr|
\le C_{\mathrm{cut}}\,e^{-c_{\mathrm{cut}}T_k^{2}}.
\end{equation}
\end{lemma}

\begin{proof}
Throughout, $T=T_k$, $C=C^{(k)}$, $\langle\cdot\rangle^{T}$ denotes the normalised expectation with
respect to $d\mu_{C}\prod_b\chi(|B(b)|<T)$, and $\langle\cdot\rangle_0$ the normalised expectation
with respect to $d\mu_C$ without the cut. For a background $W$ put
\begin{equation}\label{10.6:eq-cutoff-cost-2}
\mathsf{D}(W):=\langle\mathcal{S}_2\rangle^{T}-\langle\mathcal{S}_2\rangle_0 .
\end{equation}
The numbers $\mathsf{m}_{\mathbf{3}}(k+1)$ and $\mathsf{m}^{\circ}_{\mathbf{3}}(k+1)$ are
$-\beta^{\mathbf{3}}$ of the totals $\langle\mathcal{S}_2\rangle^{T}+\langle\mathcal{S}_1;V_1\rangle^{T}$
and $\langle\mathcal{S}_2\rangle_0+\langle\mathcal{S}_1;V_1\rangle_0$. By the one-loop number as a
half trace (Lemma~\ref{10.6:lem-half-trace}, parts (a) and (b)) and by Lemma~\ref{10.6:lem-gradient-drop-cut}
(gradient terms drop under the cut measure, an implication from the statements named there), each
total is a half trace plus gradient terms of the form $\langle\nabla\mathcal{S}_k,\cdot\rangle$, and
the gradient parts of both totals have $\beta^{\mathbf{3}}=0$. Hence
\begin{equation}\label{10.6:eq-cutoff-cost-3}
\mathsf{m}^{\circ}_{\mathbf{3}}(k+1)-\mathsf{m}_{\mathbf{3}}(k+1)=\beta^{\mathbf{3}}[\mathsf{D}] .
\end{equation}
The second-order term is a quadratic form in the fluctuation field, which we write
$\mathcal{S}_2(B)=:\langle B,\mathsf{S}_2B\rangle$; by locality $\|\mathsf{S}_2\|\le C_6$ for a
constant $C_6$. We work on a
torus $\mathbb{T}$ of level $k$, with $|\mathbb{T}|$ sites. Writing $d_{\mathfrak{g}}$ for the real
dimension of $\mathfrak{g}$, the field $B$ has four bonds per site and $d_{\mathfrak{g}}$ real
components per bond, so the number $n_{\mathbb{T}}$ of real variables satisfies
$n_{\mathbb{T}}\le4d_{\mathfrak{g}}|\mathbb{T}|$. It remains to bound $\beta^{\mathbf{3}}[\mathsf{D}]$;
this is done in five stages.

\emph{Real parameters.} Let $W$ be real and regular and $\lambda$ real with
$|\lambda|\le\varrho_T$, where, decreasing $c_{\mathrm{src}}$ if necessary (hypothesis (a) remains
valid on the smaller polydisc), $2|\lambda|\,\|\mathsf{S}_2\|\,\|C^{(k)}\|\le\frac12$; this is a
condition on $c_{\mathrm{src}}$. Define
\begin{equation}\label{10.6:eq-cutoff-cost-4}
\Delta(\lambda,W):=\log\int d\mu_C\prod_b\chi_T\,e^{\lambda\mathcal{S}_2}
-\log\int d\mu_C\,e^{\lambda\mathcal{S}_2},
\end{equation}
with $\chi_T=\chi(|B(b)|<T)$. Since $d\mu_C$ has density proportional to
$e^{-\frac12\langle B,C^{-1}B\rangle}$, the product $e^{\lambda\langle B,\mathsf{S}_2B\rangle}d\mu_C$ has
density proportional to $e^{-\frac12\langle B,(C^{-1}-2\lambda\mathsf{S}_2)B\rangle}$. As operators,
$C^{-1}-2\lambda\mathsf{S}_2\ge\|C\|^{-1}-2|\lambda|\,\|\mathsf{S}_2\|\ge(2\|C\|)^{-1}$, so
$C_\lambda:=(C^{-1}-2\lambda\mathsf{S}_2)^{-1}$ satisfies $0<C_\lambda\le2\|C\|$, and the ratio of
the two integrals in \eqref{10.6:eq-cutoff-cost-4} is the probability of the cut event under the
Gaussian measure of covariance $C_\lambda$:
\begin{equation}\label{10.6:eq-cutoff-cost-5}
\Delta(\lambda,W)=\log\mu_{C_\lambda}\Bigl(\bigcap_b\{|B(b)|<T\}\Bigr).
\end{equation}
In particular $\Delta\le0$. The ball $\{|B(b)|<T\}$ contains the cube
$\{|B^{a'}(b)|<T/\sqrt{d_{\mathfrak{g}}}\ \text{for all components } a'\}$. By \v{S}id\'ak's inequality,
hypothesis (b), applied to the centred Gaussian vector of all $n_{\mathbb{T}}$ components under
$\mu_{C_\lambda}$, the probability in \eqref{10.6:eq-cutoff-cost-5} is at least the product of the
$n_{\mathbb{T}}$ one-component probabilities. Each component has variance at most
$\|C_\lambda\|\le2\|C^{(k)}\|$, and $P(|\xi|\ge t)\le2e^{-t^2/(2\operatorname{Var}\xi)}$ for a
centred Gaussian variable $\xi$, so each factor is
at least $1-2e^{-c_1T^2}$ with
\begin{equation}\label{10.6:eq-cutoff-cost-6}
c_1:=\bigl(4d_{\mathfrak{g}}\sup_k\|C^{(k)}\|\bigr)^{-1}.
\end{equation}
Hence, for $T\ge P_0$, where $P_0>0$ is a constant depending only on $c_1$ such that
$2e^{-c_1P_0^2}\le\frac12$, and using
$\log(1-x)\ge-2x$ for $0\le x\le\frac12$,
\begin{equation}\label{10.6:eq-cutoff-cost-7}
0\ge\Delta(\lambda,W)\ge n_{\mathbb{T}}\log\bigl(1-2e^{-c_1T^2}\bigr)\ge-4n_{\mathbb{T}}e^{-c_1T^2}.
\end{equation}

\emph{Complex parameters.} $\Delta$ is analytic on $D(0,\varrho_T)$ times the complex tube of
configurations around the real regular ones on which hypothesis (a) applies, and there
$|\Delta|\le C_2|\mathbb{T}|$. Indeed, the first logarithm in \eqref{10.6:eq-cutoff-cost-4} is
$\sum_X\Phi(X)$, where $\Phi$ is the family that Lemma~\ref{10.6:lem-sourced-expansion} returns
under hypothesis (a) for the potential $V(Y)=\lambda\mathcal{S}_2(Y;B)$ (the other source parameters
set to zero), so it is analytic in $\lambda$ and $W$ and
$|\sum_X\Phi(X)|\le E^{*}\sum_Xe^{-\kappa d_{k+1}(X)}$. The second logarithm is the Gaussian integral
$\log\det(1-2\lambda C\mathsf{S}_2)^{-1/2}=-\frac12\log\det(1-2\lambda C\mathsf{S}_2)$; since
$\|2\lambda C\mathsf{S}_2\|\le\frac12$, expanding the logarithm in powers of
$2\lambda C\mathsf{S}_2$ and using $|\operatorname{Tr}\mathsf{R}^n|\le n_{\mathbb{T}}\|\mathsf{R}\|^n$
for $\mathsf{R}:=2\lambda C\mathsf{S}_2$ gives
$|\log\det(1-2\lambda C\mathsf{S}_2)|\le n_{\mathbb{T}}\log2$, and it is analytic in $\lambda$.

\emph{Two constants.} We first derive from hypothesis (c) its form on a disc. Let $\varphi$ be
holomorphic on $D(0,3\rho)$ with $|\varphi|\le\mathsf{M}_{*}$ there and $|\varphi|\le\mathsf{m}_{*}$
on the real diameter $(-3\rho,3\rho)$, where $\mathsf{m}_{*}\le\mathsf{M}_{*}$. The harmonic measure
of the diameter in the upper half disc of radius $3\rho$, at $z=\rho e^{i\vartheta}$, $0<\vartheta<\pi$,
equals
\begin{equation}\label{10.6:eq-cutoff-cost-8}
1-\frac{2}{\pi}\Bigl[\arctan\frac{\sin\vartheta}{3+\cos\vartheta}
+\arctan\frac{\sin\vartheta}{3-\cos\vartheta}\Bigr]\ge1-\frac{4}{\pi}\arctan\frac13>\frac12,
\end{equation}
and it is larger for $|z|<\rho$; the lower half disc is the reflection of the upper one. So by (c),
applied in the half disc containing $z$ with $E$ the diameter, $\log|\varphi|$, bounded by the
harmonic measure times $\log\mathsf{m}_{*}$ plus its complement times $\log\mathsf{M}_{*}$, is at most
$\frac12\log(\mathsf{m}_{*}\mathsf{M}_{*})$ on $\bar{D}(0,\rho)$, that is,
$|\varphi|\le(\mathsf{m}_{*}\mathsf{M}_{*})^{1/2}$ there. Fix a real $\mathfrak{g}$-valued bond field
$\dot{W}$ with $\sup|\dot{W}|\le1$, let
$r_W>0$ be such that $e^{i\zeta\dot{W}}$ lies in the tube of the preceding stage for $|\zeta|<r_W$,
and set
\begin{equation}\label{10.6:eq-cutoff-cost-9}
F(\lambda,\zeta):=\Delta(\lambda,e^{i\zeta\dot{W}})\qquad\text{on } D(0,\varrho_T)\times D(0,r_W).
\end{equation}
Apply the disc form first in $\lambda$, at real $\zeta$, with $\rho=\varrho_T/3$: then
$e^{i\zeta\dot{W}}$
is real and regular, so by \eqref{10.6:eq-cutoff-cost-7}
$|F|\le\mathsf{m}_{*}:=4n_{\mathbb{T}}e^{-c_1T^2}$ on
the real diameter, and $|F|\le\mathsf{M}_{*}:=C_2|\mathbb{T}|$ throughout (enlarging $C_2$ so that
$C_2\ge16d_{\mathfrak{g}}$, we have $\mathsf{m}_{*}\le\mathsf{M}_{*}$, since
$n_{\mathbb{T}}\le4d_{\mathfrak{g}}|\mathbb{T}|$); hence
$|F|\le(\mathsf{m}_{*}\mathsf{M}_{*})^{1/2}$ for $|\lambda|\le\varrho_T/3$ and real $\zeta$. Apply the
disc form next in $\zeta$, at such complex $\lambda$, with $\rho=r_W/3$, the bound on the real
diameter now being $(\mathsf{m}_{*}\mathsf{M}_{*})^{1/2}$ and the bound on the disc $\mathsf{M}_{*}$.
This gives, using $n_{\mathbb{T}}\le4d_{\mathfrak{g}}|\mathbb{T}|$,
\begin{equation}\label{10.6:eq-cutoff-cost-10}
|F|\le\mathsf{m}_{*}^{1/4}\mathsf{M}_{*}^{3/4}\le C_3|\mathbb{T}|\,e^{-c_1T^2/4}
\qquad\text{on } D(0,\varrho_T/3)\times D(0,r_W/3).
\end{equation}
The Cauchy estimates on the polydisc of radii $\varrho_T/3$ and $r_W/3$ give
\begin{equation}\label{10.6:eq-cutoff-cost-11}
\bigl|\partial_\lambda\partial_\zeta^2F(0,0)\bigr|
\le\frac{1!\,2!\,C_3|\mathbb{T}|e^{-c_1T^2/4}}{(\varrho_T/3)(r_W/3)^2}
=54\,C_3|\mathbb{T}|\,e^{-c_1T^2/4}\varrho_T^{-1}r_W^{-2}.
\end{equation}

\emph{Plane waves.} Differentiating \eqref{10.6:eq-cutoff-cost-4} at $\lambda=0$ gives
$\partial_\lambda\Delta|_{\lambda=0}=\langle\mathcal{S}_2\rangle^{T}-\langle\mathcal{S}_2\rangle_0
=\mathsf{D}$, so $\partial_\lambda\partial_\zeta^2F(0,0)=D^2\mathsf{D}(\mathbf{1})[\dot{W},\dot{W}]$,
the second
derivative being taken in the chart $B\mapsto e^{iB}$. Let $\Pi^{\mathsf{D}}(b,b'):=
D^2\mathsf{D}(\mathbf{1})[t_0\delta_b,t_0\delta_{b'}]$ on bonds of level $k+1$, and
$\hat{\Pi}^{\mathsf{D}}_{\mu\nu}(p)$ its Fourier transform. Take
$\dot{W}(c):=\cos(p\cdot x_c)\,t_0$ on the bonds $c\parallel e_\mu$ of level $k+1$ and $\dot{W}=0$
elsewhere,
with $p=p_\nu e_\nu$, $\nu\neq\mu$, in the dual torus and $p_\nu\notin\{0,\pi\}$. Then
\begin{equation}\label{10.6:eq-cutoff-cost-12}
D^2\mathsf{D}(\mathbf{1})[\dot{W},\dot{W}]=\tfrac12L^{-4}|\mathbb{T}|\,\hat{\Pi}^{\mathsf{D}}_{\mu\mu}(p).
\end{equation}
Combining with \eqref{10.6:eq-cutoff-cost-11} and $\varrho_T^{-1}=c_{\mathrm{src}}^{-1}T^{a}$, the
factor $|\mathbb{T}|$ cancels and
\begin{equation}\label{10.6:eq-cutoff-cost-13}
|\hat{\Pi}^{\mathsf{D}}_{\mu\mu}(p)|\le\epsilon_T:=C_4T^{a}e^{-c_1T^2/4},
\end{equation}
uniformly in $p$ and in $\mathbb{T}$.

\emph{Extraction.} By Lemma~\ref{10.6:lem-second-moment} (second moments read on constant curvature
fields), parts (iii) and (iv), applied to $\mathsf{D}$, together with the evenness of
$\Pi^{\mathsf{D}}$ and its decay, the order-zero and odd-order terms of
$\hat{\Pi}^{\mathsf{D}}_{\mu\mu}$ vanish and the order-two term is read off from the symbol of
$\mathsf{T}^{\mathbf{3}}$: on $p=p_\nu e_\nu$,
\begin{equation}\label{10.6:eq-cutoff-cost-14}
\hat{\Pi}^{\mathsf{D}}_{\mu\mu}(p_\nu e_\nu)=\beta^{\mathbf{3}}[\mathsf{D}]\,(h_\mu+h_\nu)\,p_\nu^2
+\varpi(p_\nu),\qquad|\varpi(p_\nu)|\le C_Rp_\nu^4,
\end{equation}
with $C_R\le\sum_x|x|^4|\Pi^{\mathsf{D}}(x)|\le C_5T^{a}$. For the last bound, $\mathsf{D}$ is the
difference of $\sum_X\partial_\lambda\Phi(X)|_{\lambda=0}$, whose members satisfy
\begin{equation*}
\bigl|\partial_\lambda\Phi(X)|_{\lambda=0}\bigr|\le(E^{*}/\varrho_T)\,e^{-\kappa d_{k+1}(X)}
\end{equation*}
by the Cauchy estimate in
$\lambda$ on $D(0,\varrho_T)$, and the Gaussian trace
\begin{equation*}
\operatorname{Tr}(C\mathsf{S}_2)
=\partial_\lambda\Bigl(-\tfrac12\log\det(1-2\lambda C\mathsf{S}_2)\Bigr)\Big|_{\lambda=0},
\end{equation*}
whose kernel decays by the locality bound for the kernels of the renormalisation step in a
background field and the corresponding trace formula
$\operatorname{Tr}(C\mathsf{S}_2)=\sum_b(C\mathsf{S}_2)(b,b)$, whose summands are such kernels: the
second derivative in the background
field at two bonds $c,c'$ of such a kernel is bounded by $C_7e^{-\delta\operatorname{diam}}$, where
$\delta>0$ is the decay rate of that bound and $\operatorname{diam}$ is the diameter of the set formed
by $c$, $c'$ and the arguments of the kernel.
Since \eqref{10.6:eq-cutoff-cost-13} is uniform in $\mathbb{T}$, after $\mathbb{T}\nearrow\mathbb{Z}^4$
by hypothesis (d) the momentum $p_\nu$ is arbitrary. Choose
\begin{equation}\label{10.6:eq-cutoff-cost-15}
p_\nu^2:=(\epsilon_T/C_R)^{1/2},\qquad h=(1,1,-1,-1),\qquad(\mu,\nu)=(1,2),
\end{equation}
so that $h_\mu+h_\nu=2$. Then \eqref{10.6:eq-cutoff-cost-13} and \eqref{10.6:eq-cutoff-cost-14} give
$2|\beta^{\mathbf{3}}[\mathsf{D}]|\,p_\nu^2\le\epsilon_T+C_Rp_\nu^4=2\epsilon_T$, that is,
\begin{equation}\label{10.6:eq-cutoff-cost-16}
|\beta^{\mathbf{3}}[\mathsf{D}]|\le\frac{\epsilon_T}{p_\nu^2}=(\epsilon_TC_R)^{1/2}
\le C_8\,T^{a}e^{-c_1T^2/8}.
\end{equation}
Finally $T^{a}e^{-c_1T^2/8}\le\bigl(\sup_{t\ge0}t^{a}e^{-c_1t^2/16}\bigr)e^{-c_1T^2/16}$, and with
\eqref{10.6:eq-cutoff-cost-3} this is \eqref{10.6:eq-cutoff-cost-1} with $c_{\mathrm{cut}}=c_1/16$.
The constants $c_1,C_2,\dots,C_8$, $c_{\mathrm{src}}$, $a$, $E^{*}$, $\kappa$ and $r_W$ are those of
the analysis of the correlation theorem and do not depend on $k$ or $g$.
\end{proof}

\begin{lemma}[Summability of the cutoff corrections]\label{10.6:lem-cutoff-sum}
Assume the reference regime of the correlation theorem. In it, for $X = g^{-2}$ and
$0 \le l \le K$, the running couplings $x_l = g_l^{-2}$ satisfy
\begin{equation}\label{10.6:eq-cutoff-sum-1}
x_l \ge X + \lambda (K-l) - 2c_2 ,
\end{equation}
where $\lambda > 0$ and $c_2$ are the constants of the running-coupling bounds of that regime and
do not depend on $l$, $K$ or $g$. Let $C_{\mathrm{cut}}, c_{\mathrm{cut}} > 0$ be the constants of
the cost of the sharp field cutoff (Lemma~\ref{10.6:lem-cutoff-cost}); they depend neither on $k$
nor on $g$. Let $T_k = A_1 (\log x_k)^{p_0}$ be the field cut of the correlation theorem, with
Ba{\l}aban's parameters $A_1$ and $p_0$. Then there is a threshold, depending only on
$C_{\mathrm{cut}}$, $c_{\mathrm{cut}}$, $A_1$, $p_0$, $\lambda$ and $c_2$, such that for all
$X = g^{-2}$ above it
\begin{equation}\label{10.6:eq-cutoff-sum-a}
S_{\mathrm{cut}} := \sup_{K} \sum_{l=1}^{K} C_{\mathrm{cut}}\, e^{-c_{\mathrm{cut}} T_{l-1}^{2}} \le 1 .
\end{equation}
Thus \eqref{10.6:eq-cutoff-sum-a} holds below a ceiling on $g$, and nothing is asserted above it.
\end{lemma}

\begin{proof}
Put $X' := X - 2c_2$ and write $n := K - l$. As $l$ runs over $1, \dots, K$, the index $n$
runs over $0, \dots, K-1$.

Recall that Ba{\l}aban's exponent $p_0$ satisfies $2p_0 > 1$. Set
$\varphi(t) := c_{\mathrm{cut}} A_1^{2} (\log t)^{2p_0}$ for $t > 1$. Then
\[
\varphi(t) / \log t = c_{\mathrm{cut}} A_1^{2} (\log t)^{2p_0 - 1} .
\]
Because $2p_0 > 1$, this quotient tends to $\infty$ as $t \to \infty$. Hence there is
$t_* \ge e$, depending only on $c_{\mathrm{cut}}$, $A_1$ and $p_0$, such that
$\varphi(t) \ge 2 \log t$ for all $t \ge t_*$.

Suppose $X$ is so large that $X' \ge t_*$. Let $1 \le l \le K$. Since $0 \le l-1 \le K-1$,
\eqref{10.6:eq-cutoff-sum-1} applies at the index $l-1$ and gives
\[
x_{l-1} \ge X' + \lambda (n+1) \ge X' + \lambda n \ge t_* \ge e .
\]
The function $t \mapsto (\log t)^{2p_0}$ is increasing on $t \ge 1$, so
\[
T_{l-1}^{2} = A_1^{2} (\log x_{l-1})^{2p_0} \ge A_1^{2} \bigl(\log (X' + \lambda n)\bigr)^{2p_0} .
\]
This bounds the $l$-th term of \eqref{10.6:eq-cutoff-sum-a} by
\begin{equation}\label{10.6:eq-cutoff-sum-2}
C_{\mathrm{cut}}\, e^{-c_{\mathrm{cut}} T_{l-1}^{2}}
\le C_{\mathrm{cut}} \exp\bigl(-c_{\mathrm{cut}} A_1^{2} (\log(X' + \lambda n))^{2p_0}\bigr).
\end{equation}
Consequently, for every $K$,
\begin{equation}\label{10.6:eq-cutoff-sum-3}
\sum_{l=1}^{K} C_{\mathrm{cut}}\, e^{-c_{\mathrm{cut}} T_{l-1}^{2}}
\le C_{\mathrm{cut}} \sum_{n \ge 0}
\exp\bigl(-c_{\mathrm{cut}} A_1^{2} (\log(X' + \lambda n))^{2p_0}\bigr).
\end{equation}
The right-hand side does not depend on $K$, so it also bounds $S_{\mathrm{cut}}$.

We show that the series in \eqref{10.6:eq-cutoff-sum-3} converges and tends to $0$ as
$X \to \infty$. Since $X' + \lambda n \ge t_*$ for every $n \ge 0$, the $n$-th term of the
series is $e^{-\varphi(X' + \lambda n)} \le (X' + \lambda n)^{-2}$. The function
$u \mapsto (X' + \lambda u)^{-2}$ is decreasing on $[0, \infty)$, so comparison with its
integral gives
\[
\sum_{n \ge 0} (X' + \lambda n)^{-2}
\le X'^{-2} + \int_0^{\infty} (X' + \lambda u)^{-2}\, du
= X'^{-2} + (\lambda X')^{-1} .
\]
Therefore
\[
S_{\mathrm{cut}} \le C_{\mathrm{cut}} \bigl( X'^{-2} + (\lambda X')^{-1} \bigr),
\qquad X' = X - 2c_2 .
\]
This bound is finite and tends to $0$ as $X = g^{-2} \to \infty$. In particular it is at most
$1$ for all $X$ above a threshold depending only on $C_{\mathrm{cut}}$, $c_{\mathrm{cut}}$,
$A_1$, $p_0$, $\lambda$ and $c_2$, which is \eqref{10.6:eq-cutoff-sum-a}.
\end{proof}

\begin{corollary}[The one-loop anisotropy inequality with $\kappa = 1$]
\label{10.6:cor-anisotropy-inequality}
Assume the hypotheses of Theorem~\ref{10.6:thm-one-loop-drift}, and the coupling-flow identity
for the one-loop coefficients $\beta^{0}_{l}$ of the coupling: for all $0 \le i < j$,
\begin{equation}\label{10.6:eq-anisotropy-inequality-1}
\sum_{l=i+1}^{j} \beta^{0}_{l} = (j-i)\,c_0 + s_i - s_j ,
\qquad |s_i| \le c_2 \ \text{for every } i ,
\end{equation}
with $c_0 = (11C_{\mathfrak{g}}/24\pi^2)\log L > 0$, the coefficient of
Theorem~\ref{10.6:thm-one-loop-drift}. Let $c^{(3)} = c^{(3)}(L,G)$ be the constant of that theorem.
\begin{enumerate}
\item For every real $\kappa \le 1$ the cut-free numbers satisfy the one-loop anisotropy inequality
\begin{equation}\label{10.6:eq-anisotropy-inequality-2}
\sum_{l=i+1}^{j}\bigl[\mathsf{m}^{\circ}_{\mathbf{3}}(l) - \kappa\,\beta^{0}_{l}\bigr]
\ge -c_2^{(\mathbf{3})}
\qquad \text{for all } 0 \le i < j ,
\end{equation}
with $c_2^{(\mathbf{3})} := 2c^{(3)} + 2|\kappa|\,c_2$. For no $\kappa > 1$ is there a finite
$c_2^{(\mathbf{3})}$ for which \eqref{10.6:eq-anisotropy-inequality-2} holds.
\item Assume in addition the hypotheses of Lemma~\ref{10.6:lem-cutoff-cost}, in the regime of the
correlation theorem, and the one-sided ceiling $S_{\mathrm{cut}} \le 1$ of
Lemma~\ref{10.6:lem-cutoff-sum}, display~\eqref{10.6:eq-cutoff-sum-a}. Then the same holds for the
numbers $\mathsf{m}_{\mathbf{3}}(l)$ with the cut, in place of $\mathsf{m}^{\circ}_{\mathbf{3}}(l)$,
with $c_2^{(\mathbf{3})} := 2c^{(3)} + 2|\kappa|\,c_2 + 1$.
\end{enumerate}
\end{corollary}

\begin{proof}
Fix $0 \le i < j$ and a real $\kappa$. By Theorem~\ref{10.6:thm-one-loop-drift},
$\sum_{l=i+1}^{j}\mathsf{m}^{\circ}_{\mathbf{3}}(l) = (j-i)c_0 + s^{(3)}_i - s^{(3)}_j$ with
$|s^{(3)}_k| \le c^{(3)}$ for every $k$. Subtracting $\kappa$ times
\eqref{10.6:eq-anisotropy-inequality-1} gives
\begin{equation}\label{10.6:eq-anisotropy-inequality-3}
\sum_{l=i+1}^{j}\bigl[\mathsf{m}^{\circ}_{\mathbf{3}}(l) - \kappa\,\beta^{0}_{l}\bigr]
= (1-\kappa)(j-i)\,c_0 + \bigl(s^{(3)}_i - s^{(3)}_j\bigr) - \kappa\,(s_i - s_j).
\end{equation}
The two differences on the right are bounded in absolute value:
$|s^{(3)}_i - s^{(3)}_j| \le 2c^{(3)}$ and $|\kappa(s_i - s_j)| \le 2|\kappa|\,c_2$.

Let $\kappa \le 1$. Then $1 - \kappa \ge 0$, $j - i > 0$ and $c_0 > 0$, so the first term on the
right of \eqref{10.6:eq-anisotropy-inequality-3} is non-negative, and the right-hand side is at least
$-2c^{(3)} - 2|\kappa|\,c_2$. This is \eqref{10.6:eq-anisotropy-inequality-2} with
$c_2^{(\mathbf{3})} = 2c^{(3)} + 2|\kappa|\,c_2$.

Let $\kappa > 1$. Then $(1-\kappa)c_0 < 0$, so by \eqref{10.6:eq-anisotropy-inequality-3} the sum is
at most $(1-\kappa)(j-i)c_0 + 2c^{(3)} + 2|\kappa|\,c_2$, which tends to $-\infty$ as $j - i \to \infty$.
Hence for every finite constant $c$ there are $0 \le i < j$ for which the sum is smaller than $-c$, and
\eqref{10.6:eq-anisotropy-inequality-2} fails for every finite $c_2^{(\mathbf{3})}$.

With the cut, write
\[
\sum_{l=i+1}^{j}\bigl[\mathsf{m}_{\mathbf{3}}(l) - \kappa\,\beta^{0}_{l}\bigr]
= \sum_{l=i+1}^{j}\bigl[\mathsf{m}^{\circ}_{\mathbf{3}}(l) - \kappa\,\beta^{0}_{l}\bigr]
+ \sum_{l=i+1}^{j}\bigl[\mathsf{m}_{\mathbf{3}}(l) - \mathsf{m}^{\circ}_{\mathbf{3}}(l)\bigr].
\]
By Lemma~\ref{10.6:lem-cutoff-cost}, applied at each level $l$ of the sum, the $l$-th term of the last
sum is at most $C_{\mathrm{cut}}e^{-c_{\mathrm{cut}}T_{l-1}^2}$ in absolute value. These bounds are
non-negative and the levels $l$ lie in $\{1,\dots,K\}$. Hence, by
\eqref{10.6:eq-cutoff-sum-a}, the absolute value of the last sum is at most
$\sum_{l=1}^{K}C_{\mathrm{cut}}e^{-c_{\mathrm{cut}}T_{l-1}^2} \le S_{\mathrm{cut}} \le 1$.
Thus the sum with the cut differs from the cut-free sum by at most one. For $\kappa \le 1$ it is
therefore at least $-2c^{(3)} - 2|\kappa|\,c_2 - 1$. For $\kappa > 1$ it is at most
$(1-\kappa)(j-i)c_0 + 2c^{(3)} + 2|\kappa|\,c_2 + 1$, which tends to $-\infty$ as $j-i \to \infty$, so
no finite constant works.
\end{proof}

\begin{corollary}[Two-sided decay of the tensor normalisation; stated; proof incomplete at the step
marked $(\ast)$]\label{10.6:cor-normalisation-decay}
Assume the one-loop hypothesis on the tensor normalisation flow, the coupling-flow theorem and the
reference regime of the correlation theorem. Then, for $0 \le j \le k \le K$,
\begin{equation}\label{10.6:eq-normalisation-decay-1}
C_N^{-1}\,\frac{x_k}{x_j}
\;\le\; \left|\frac{P^{\mathbf{3}}_{k}(0)}{P^{\mathbf{3}}_{j}(0)}\right|
\;\le\; C_N\,\frac{x_k}{x_j}.
\end{equation}
Consequently, in the setting of the stress-tensor-insertion correlation theorem in the channel
$\mathbf{3}$, the statement on the tensor normalisation carried there with the parameter $c_1$ holds
with $c_1 = 1$.
\end{corollary}

\begin{proof}
Throughout, $\nu_l$ and $\varrho_l$ denote the relative increments $\nu^{\mathbf{3}}_{l}$ and the
remainders $\varrho^{\mathbf{3}}_{l}$ of the normalisation flow in the channel $\mathbf{3}$, evaluated at
$\zeta = 0$. The product is $P^{\mathbf{3}}_{k} = \prod_{l<k}(1+\nu_l)$. Hence, for $0 \le j \le k \le K$,
\[
\frac{P^{\mathbf{3}}_{k}(0)}{P^{\mathbf{3}}_{j}(0)} = \prod_{l=j}^{k-1}(1+\nu_l).
\]
All sums over $l$ below run over the same window $j \le l \le k-1$.

\emph{Upper bound.} The upper bound in \eqref{10.6:eq-normalisation-decay-1} is the normalisation-decay
lemma applied with exponent $\kappa = 1$. That lemma has four hypotheses: the coupling-flow theorem, the
reference regime, the one-loop hypothesis on the tensor normalisation flow, and the one-loop anisotropy
inequality with exponent $\kappa$. The first three are assumed in the statement. The fourth is supplied
by Corollary~\ref{10.6:cor-anisotropy-inequality}, which holds for every $\kappa \le 1$ and so in
particular for $\kappa = 1$.
Since the numbers $\mathsf{m}_{\mathbf{3}}(l)$ are the numbers with the field cut, that corollary is used
in the regime of the correlation theorem and under the ceiling \eqref{10.6:eq-cutoff-sum-a},
$S_{\mathrm{cut}} \le 1$. Its constant is then
$c_2^{(\mathbf{3})} = 2c^{(3)} + 2c_2 + 1$.

\emph{Lower bound.} The argument rests on three facts.
\begin{itemize}
\item For $|\nu| \le \tfrac12$ one has $\log(1+\nu) \ge \nu - \nu^2$.
\item Let $C_\nu$ denote the constant in the bound $|\nu_l| \le C_\nu\, g_l^2(\log x_l)^{2p_0}$ on the
relative increments. Then
\[
\sum_l \nu_l^2 \;\le\; C_\nu^2 \sum_l g_l^4(\log x_l)^{4p_0} \;\le\; C ,
\]
where the last inequality holds by the reference regime.
\item The increments satisfy
\[
\sum_l \nu_l \;=\; -\sum_l \mathsf{m}_{\mathbf{3}}(l+1)\,u_l \;+\; \sum_l \varrho_l ,
\]
with $u_l$ the Abel weights.
\end{itemize}

$(\ast)$ The next step is the Abel summation in the proof of the normalisation-decay lemma; it is taken
from that proof and is not carried out in the present proof. Put
$D_l := \mathsf{m}_{\mathbf{3}}(l) - \beta^{0}_{l}$, and let $S(l,k)$ be the partial
sums of the $D_i$ that enter that summation. In the present setting these partial sums are bounded in
modulus, not only from below, by $2c^{(3)} + 2c_2 + S_{\mathrm{cut}}$; here $c^{(3)}$ is the bound on the
coboundary $s^{(3)}_l$, $c_2$ the constant of the coupling-flow theorem and $S_{\mathrm{cut}}$ the summed
cut correction of \eqref{10.6:eq-cutoff-sum-a}. With $u'_l$ the second family of Abel weights, that
summation then gives
\begin{equation}\label{10.6:eq-normalisation-decay-2}
\Bigl|\sum_l D_{l+1}\,u'_l\Bigr| \;\le\; 2\bigl(2c^{(3)} + 2c_2 + S_{\mathrm{cut}}\bigr)\,u'_{k-1}.
\end{equation}
The remaining steps of the proof of the normalisation-decay lemma are asserted to hold with both signs
as written there; they too are taken from that proof and are not carried out in the present proof.
Combined with
\eqref{10.6:eq-normalisation-decay-2}, they give
\begin{equation}\label{10.6:eq-normalisation-decay-3}
\sum_l \mathsf{m}_{\mathbf{3}}(l+1)\,u_l \;=\; \log\frac{x_j}{x_k} + O(1).
\end{equation}

The assembly of the lower bound from these pieces uses two further inputs: that $|\nu_l| \le \tfrac12$
for $j \le l \le k-1$, so that the first fact applies term by term, and that $\bigl|\sum_l \varrho_l\bigr|$
is bounded independently of $j$, $k$ and $K$. Neither is established in this proof; both belong to the
step marked $(\ast)$. Granted them, we assemble the pieces. Take logarithms in the product formula
above. Apply the inequality
$\log(1+\nu_l) \ge \nu_l - \nu_l^2$ term by term, and use the bound on $\sum_l \nu_l^2$. Insert the
identity for $\sum_l \nu_l$, and then \eqref{10.6:eq-normalisation-decay-3}. The result is
\[
\log\left|\frac{P^{\mathbf{3}}_{k}(0)}{P^{\mathbf{3}}_{j}(0)}\right|
\;\ge\; \log\frac{x_k}{x_j} + O(1),
\]
which is the lower bound in \eqref{10.6:eq-normalisation-decay-1}.
\end{proof}

\begin{definition}[Test-function moduli, separated classes and the action of rotations]
\label{10.7:not-test-function-moduli}
Test functions are functions on $\mathbb{R}^4$. Recall that for bounded random variables
$Y_1,\dots,Y_N$ the joint cumulant is
$\kappa(Y_1,\dots,Y_N) := \partial_{z_1}\cdots\partial_{z_N}\log E[e^{\sum_i z_iY_i}]|_{z=0}$,
and that for a compactly supported $h\colon\mathbb{R}^4\to\mathbb{R}$ the periodic lift
$h^{(m)}$ is the function $\sum_{v\in 2L^m\mathbb{Z}^4}h(\cdot+v)$, read as a function on
$\mathbb{T}_m$. Throughout this chapter the following notation is used.
\begin{enumerate}
\item \emph{Norms and moduli.} For $f\colon\mathbb{R}^4\to\mathbb{R}$, $\operatorname{Lip}f$ is
the Lipschitz constant for the Euclidean metric, and
$\|f\|_{\sharp} := \max\{\|f\|_\infty,\ 40M\operatorname{Lip}f\}$. For $f\in C^1_c(\mathbb{R}^4)$
we write $\|\nabla f\|_\infty := \sup|\nabla f|$; then $\operatorname{Lip}f=\|\nabla f\|_\infty$.
The support radius is $\rho_f := \max\{|x| : x\in\operatorname{supp}f\}$, with $\rho_f:=0$ if
$f\equiv 0$. For a function $u$ on $\mathbb{R}^4$ (scalar or vector valued, $|\cdot|$ the
Euclidean norm) and $\delta\ge 0$, the modulus of continuity is
$m_u(\delta) := \sup\{|u(x)-u(y)| : |x-y|\le\delta\}$; it tends to $0$ as $\delta\to 0$
whenever $u$ is continuous and compactly supported. For $f\in C^1_c(\mathbb{R}^4)$ and
$s\ge 0$ the rotation modulus is
\begin{equation}\label{10.7:eq-test-function-moduli-1}
\mu_f(s) := \max\Bigl\{\,m_f(s\rho_f),\ 40M\bigl(m_{\nabla f}(s\rho_f)
+ s\|\nabla f\|_\infty\bigr)\Bigr\};
\end{equation}
it is non-decreasing in $s$ and tends to $0$ as $s\to 0$. For $A\subset\mathbb{R}^4$,
$A^{+\theta} := \{x : \operatorname{dist}(x,A)\le\theta\}$.
\item \emph{Rotations.} $O(4)$ acts on test functions on the left by
$(R\cdot f)(x) := f(R^{-1}x)$. For $\omega\in so(4)$ (a real antisymmetric $4\times4$
matrix), $R^{\omega}_t := e^{t\omega}$, and for $f\in C^1_c(\mathbb{R}^4)$
\begin{equation}\label{10.7:eq-test-function-moduli-2}
\delta_{\omega}f := \frac{d}{dt}\Big|_{t=0} R^{\omega}_t\cdot f = -(\omega x)\cdot\nabla f .
\end{equation}
$\|\omega\|$ and $\|Q\|$ denote operator norms of matrices, and
\begin{equation}\label{10.7:eq-test-function-moduli-3}
\|e^{t\omega}-\mathbb{1}\| = \Bigl\|\int_0^t\omega e^{s\omega}\,ds\Bigr\| \le |t|\,\|\omega\| .
\end{equation}
$W_4$ is the group of signed permutation matrices, and
$w_0 := \operatorname{diag}(-1,1,1,1)\in W_4$.
\item \emph{Separated classes.} For $n\ge 2$, $\mathcal{R}>0$, $\mathsf{d}>0$ and
$r\in\{1,2,\dots,\infty\}$, $\mathcal{T}^{(r)}_n(\mathcal{R},\mathsf{d})$ is the set of tuples
$\vec h=(h_1,\dots,h_n)$ of real functions $h_i\in C^r_c(\mathbb{R}^4)$ with
$\operatorname{supp}h_i\subset B(0,\mathcal{R}/2)$, the open ball of radius $\mathcal{R}/2$
about $0$, and $\operatorname{dist}(\operatorname{supp}h_i,\operatorname{supp}h_k)\ge\mathsf{d}$
for $i\ne k$. A tuple of real $C^1_c$ functions with pairwise disjoint compact supports is
called \emph{admissible}. A cut-off $\chi$ is \emph{admissible at radius} $\mathcal{R}$ if
$\chi\in C^\infty_c(\mathbb{R}^4)$ and $\chi\equiv 1$ on $B(0,\mathcal{R})$.
\item \emph{Lattice isometries.} For a bare plaquette $p$ put
$a_p := 1-\operatorname{Re}\operatorname{tr}U(\partial p)\in[0,2]$, so that
$\mathcal{O}(\varphi)=\sum_p\varphi(x(p))\,a_p$, the sum running over all bare plaquettes of
all six orientations. An isometry $\sigma$ of the torus lattice acts on configurations by
$(\sigma U)(\langle x,y\rangle) := U(\langle\sigma^{-1}x,\sigma^{-1}y\rangle)$, and for a
function $F$ of the field $F^{\sigma}(U) := F(\sigma U)$. If $F$ depends only on the bonds
in a set $S$, then $F^{\sigma}$ depends only on the bonds in $\sigma^{-1}(S)$ (for the
involutions used below, $\sigma^{-1}(S)=\sigma(S)$). Moreover
$a_p(\sigma U)=a_{\sigma^{-1}p}(U)$, and $\mathcal{O}(\varphi)^{\sigma}=\mathcal{O}(\varphi\circ\sigma)$
for every plaquette weight $\varphi$.
\item \emph{Orientation.} In the derivation of the rotational Ward identity the rotated
observable is $\mathcal{O}(f)\circ R := \mathcal{O}(R^{-1}\cdot f)$, so the derivative
$\frac{d}{dt}|_{t=0}\langle F\circ R^{\omega}_t\rangle$ appearing there concerns the tuples
$R^{-\omega}_t\cdot\vec h$. In what follows $D^{\omega}$ denotes the derivative at $t=0$ along
$t\mapsto R^{\omega}_t\cdot\vec h$; it is the left side of the rotational Ward identity at
$-\omega$. Since every hypothesis is assumed for all $\omega\in so(4)$, this is a relabelling,
and the named terms of the rotational Ward identity are read at $-\omega$.
\end{enumerate}
\end{definition}

\begin{proof}
We justify the assertions contained in the definition, in the order in which they occur.

For $f\in C^1_c(\mathbb{R}^4)$ and $x,y\in\mathbb{R}^4$ the mean value inequality along the
segment from $y$ to $x$ gives $|f(x)-f(y)|\le\|\nabla f\|_\infty|x-y|$, so
$\operatorname{Lip}f\le\|\nabla f\|_\infty$; conversely, for every $x$ and every unit vector $e$,
$|\nabla f(x)\cdot e|=\lim_{h\to0}|f(x+he)-f(x)|/|h|\le\operatorname{Lip}f$, and taking
$e=\nabla f(x)/|\nabla f(x)|$ where $\nabla f(x)\neq0$ gives
$\|\nabla f\|_\infty\le\operatorname{Lip}f$.

A continuous compactly supported $u$ is uniformly continuous on $\mathbb{R}^4$, which is the
statement $m_u(\delta)\to0$ as $\delta\to0$. For $f\in C^1_c$ both $f$ and $\nabla f$ are
continuous and compactly supported. The set over which the supremum defining $m_u(\delta)$ is
taken grows with $\delta$, so $m_u$ is non-decreasing; since $\rho_f\ge0$, the arguments
$s\rho_f$ and the term $s\|\nabla f\|_\infty$ are non-decreasing in $s$, and the maximum of
non-decreasing functions is non-decreasing. This gives the monotonicity of $\mu_f$ in
\eqref{10.7:eq-test-function-moduli-1}; as $s\to0$ one has $s\rho_f\to0$, so each of
$m_f(s\rho_f)$, $m_{\nabla f}(s\rho_f)$ and $s\|\nabla f\|_\infty$ tends to $0$, hence
$\mu_f(s)\to0$.

The action is a left action because
$(R\cdot(R'\cdot f))(x)=f(R'^{-1}R^{-1}x)=((RR')\cdot f)(x)$. Since $\omega^{T}=-\omega$,
$(e^{s\omega})^{T}=e^{-s\omega}=(e^{s\omega})^{-1}$, so $e^{s\omega}$ is orthogonal and
$R^{\omega}_t\in O(4)$. By definition $(R^{\omega}_t\cdot f)(x)=f(e^{-t\omega}x)$, and the
chain rule gives
$\frac{d}{dt}|_{t=0}f(e^{-t\omega}x)=\nabla f(x)\cdot(-\omega x)$, which is
\eqref{10.7:eq-test-function-moduli-2}. Since $\frac{d}{ds}e^{s\omega}=\omega e^{s\omega}$,
the fundamental theorem of calculus gives the equality in
\eqref{10.7:eq-test-function-moduli-3}; as $\|e^{s\omega}\|=1$ by orthogonality,
$\|\omega e^{s\omega}\|\le\|\omega\|$ for every $s$, and integrating over the interval between
$0$ and $t$ gives the inequality.

Let $\sigma$ be a lattice isometry and let $F$ depend only on the bonds in $S$. The values of
$\sigma U$ on the bonds $b\in S$ are $U(\sigma^{-1}b)$, so $F^{\sigma}(U)=F(\sigma U)$ is
determined by the values of $U$ on $\sigma^{-1}(S)$; for an involution $\sigma^{-1}=\sigma$.
The holonomy of $\sigma U$ around $\partial p$ is the ordered product of the
$(\sigma U)(b)=U(\sigma^{-1}b)$ over the bonds $b$ of $\partial p$, that is, the holonomy of
$U$ around $\sigma^{-1}(\partial p)=\partial(\sigma^{-1}p)$, possibly with another base point
and orientation. A change of base point conjugates the holonomy, which leaves the trace
unchanged by cyclicity, and a reversal of orientation replaces the holonomy $V\in SU(N)$ by
$V^{-1}=V^{*}$, with $\operatorname{tr}V^{*}=\overline{\operatorname{tr}V}$, which leaves
$\operatorname{Re}\operatorname{tr}$ unchanged. Hence $a_p(\sigma U)=a_{\sigma^{-1}p}(U)$.
Since $p\mapsto q=\sigma^{-1}p$ is a bijection of the bare plaquettes (of all orientations)
and $x(\sigma q)=\sigma x(q)$,
\begin{equation*}
\mathcal{O}(\varphi)^{\sigma}(U)=\sum_p\varphi(x(p))\,a_{\sigma^{-1}p}(U)
=\sum_q\varphi(\sigma x(q))\,a_q(U)=\mathcal{O}(\varphi\circ\sigma)(U).
\end{equation*}

Finally, $(R^{\omega}_t)^{-1}=e^{-t\omega}=R^{-\omega}_t$, so
$\mathcal{O}(f)\circ R^{\omega}_t=\mathcal{O}(R^{-\omega}_t\cdot f)$; the derivative at $t=0$
of $\langle F\circ R^{\omega}_t\rangle$ is therefore the derivative along
$t\mapsto R^{-\omega}_t\cdot\vec h$, which is $D^{-\omega}$. Replacing $\omega$ by $-\omega$,
$D^{\omega}$ is the left side of the rotational Ward identity at $-\omega$. Since
$\omega\mapsto-\omega$ is a bijection of $so(4)$ and every hypothesis is assumed for all
$\omega\in so(4)$, nothing is lost by this relabelling.
\end{proof}

\begin{definition}[The smeared block algebra and the symbol rotation]
\label{10.7:def-smeared-block-algebra}
Let
\begin{equation*}
\mathbb{L}_s := L^{-s}\bigl(\mathbb{Z}+\tfrac12\bigr)^4 ,
\end{equation*}
the half-offset lattices; since $L$ is odd they are nested, $\mathbb{L}_s\subset\mathbb{L}_{s+1}$.
The smeared block algebra is $\mathfrak{A}^{sm}_{loc} := \mathbb{C}[\mathfrak{S}]$, the free
commutative unital $*$-algebra on the set $\mathfrak{S}$ of symbols $(F,\varphi)$ with
$F \in \mathfrak{A}^{\tau}_{loc}(\mathbb{R}^4)$ and $\varphi \in C_c(\mathbb{R}^4)$. Here the
algebra $\mathfrak{A}^{\tau}_{loc}(\mathbb{R}^4)$, the translations $\tau_v$ and the action
$\alpha$ of translations and of $W_4$ on it, the realisation $(\,\cdot\,)_K$ at cutoff $K$, the
coherent families and the functional $\omega_\infty$ on $\mathbb{C}[\mathfrak{S}]$ are those fixed
with the objects of clause (ii-e), where the smeared block algebra is introduced. At cutoff $K$ a
symbol is realised by
\begin{equation}
\label{10.7:eq-smeared-block-algebra-1}
(F,\varphi)_K
= \sum_{v \in L^{-K}\mathbb{Z}^4} L^{-4K}\,\varphi(v)\,(\alpha_{\tau_v}F)_K .
\end{equation}
The group $\mathbb{R}^4 \rtimes W_4$ acts on symbols by
\begin{equation*}
\alpha_a(F,\varphi) := \bigl(F,\varphi(\cdot - a)\bigr), \qquad
\alpha_w(F,\varphi) := \bigl(\alpha_w F,\varphi\circ w^{-1}\bigr),
\end{equation*}
for $a \in \mathbb{R}^4$ and $w \in W_4$. Each $\alpha_a$ is a bijection of $\mathfrak{S}$, with
inverse $\alpha_{-a}$, and so extends, in the same way as $\beta_R$ below, to a unital
$*$-automorphism of $\mathbb{C}[\mathfrak{S}]$, again denoted $\alpha_a$. The commensurable motion
group is
$\mathfrak{G} := \mathbb{Z}[1/L]^4 \rtimes W_4$; it acts on coherent families through its action
on bare configurations, $g\cdot V_K(U) = V_K(g\cdot U)$.

For $R \in O(4)$, the symbol rotation $\beta_R$ is the unital $*$-automorphism of
$\mathbb{C}[\mathfrak{S}]$ determined on the generators by
\begin{equation}
\label{10.7:eq-smeared-block-algebra-2}
\beta_R(F,\varphi) := \bigl(F,\varphi\circ R^{-1}\bigr).
\end{equation}
It satisfies, for $R, S \in O(4)$ and $a \in \mathbb{R}^4$,
\begin{equation}
\label{10.7:eq-smeared-block-algebra-3}
\beta_R\beta_S = \beta_{RS}, \qquad \beta_R\alpha_a = \alpha_{Ra}\beta_R .
\end{equation}
The rotation-invariance statement $\omega(F\circ R) = \omega(F)$ for a functional $\omega$ on the
algebra, read on $\mathbb{C}[\mathfrak{S}]$
with the rotation acting on the test function $\varphi$, is the statement $(S_\beta)(R)$:
\begin{equation}
\label{10.7:eq-smeared-block-algebra-a}
\omega_\infty\circ\beta_R = \omega_\infty .
\end{equation}
\end{definition}

\begin{proof}
We check the three assertions contained in the definition.

Nesting of the lattices. Since $L$ is odd, $(L-1)/2$ is an integer, so for $n \in \mathbb{Z}$
\begin{equation*}
L\bigl(n+\tfrac12\bigr) = \Bigl(Ln + \tfrac{L-1}{2}\Bigr) + \tfrac12 \in \mathbb{Z}+\tfrac12 ,
\end{equation*}
that is, $L(\mathbb{Z}+\tfrac12) \subset \mathbb{Z}+\tfrac12$. Hence
\begin{equation*}
\mathbb{L}_s = L^{-(s+1)}\,L\bigl(\mathbb{Z}+\tfrac12\bigr)^4
\subset L^{-(s+1)}\bigl(\mathbb{Z}+\tfrac12\bigr)^4 = \mathbb{L}_{s+1},
\end{equation*}
and the lattices $\mathbb{L}_s$ are nested.

The map $\beta_R$ is a unital $*$-automorphism. For $\varphi \in C_c(\mathbb{R}^4)$ the function
$\varphi\circ R^{-1}$ is continuous, and its support is $R(\operatorname{supp}\varphi)$, which is
compact; so \eqref{10.7:eq-smeared-block-algebra-2} defines a map
$\mathfrak{S}\to\mathfrak{S}$. Since $\mathbb{C}[\mathfrak{S}]$ is free as a commutative unital
$*$-algebra on $\mathfrak{S}$, this map extends uniquely to a unital $*$-homomorphism of
$\mathbb{C}[\mathfrak{S}]$, again denoted $\beta_R$. On the generators,
$(F,\varphi)\mapsto(F,\varphi\circ R)$ is a two-sided inverse of
$(F,\varphi)\mapsto(F,\varphi\circ R^{-1})$, so the map on $\mathfrak{S}$ is a bijection of the
free generating set; the unital $*$-homomorphism $\beta_{R^{-1}}$ it induces composes with
$\beta_R$, in either order, to a unital $*$-homomorphism that is the identity on $\mathfrak{S}$,
hence, by uniqueness of extensions, the identity on $\mathbb{C}[\mathfrak{S}]$. Thus $\beta_R$ is a
unital $*$-automorphism.

The relations \eqref{10.7:eq-smeared-block-algebra-3}. Both sides of each relation are unital
$*$-homomorphisms of $\mathbb{C}[\mathfrak{S}]$, so by the uniqueness just used it suffices to
compare them on a symbol $(F,\varphi)$. First,
\begin{equation*}
\beta_R\beta_S(F,\varphi) = \beta_R\bigl(F,\varphi\circ S^{-1}\bigr)
= \bigl(F,\varphi\circ S^{-1}\circ R^{-1}\bigr) = \bigl(F,\varphi\circ(R\circ S)^{-1}\bigr)
= \beta_{RS}(F,\varphi).
\end{equation*}
Second, $\beta_R\alpha_a(F,\varphi) = \beta_R(F,\varphi(\cdot-a))$ is the symbol whose test
function is $x \mapsto \varphi(R^{-1}x - a)$, while
$\alpha_{Ra}\beta_R(F,\varphi) = \alpha_{Ra}(F,\varphi\circ R^{-1})$ is the symbol whose test
function is
\begin{equation*}
x \mapsto \varphi\bigl(R^{-1}(x - Ra)\bigr) = \varphi\bigl(R^{-1}x - a\bigr),
\end{equation*}
by linearity of $R^{-1}$. The two symbols coincide, which gives
$\beta_R\alpha_a = \alpha_{Ra}\beta_R$.
\end{proof}

\begin{lemma}[The stabiliser of the half-offset lattices]\label{10.7:lem-lattice-stabiliser}
Let $W_4$ be the group of signed permutation matrices, and for $L$ and $s$ let
$\mathbb{L}_s = L^{-s}(\mathbb{Z}+\tfrac12)^4$ be the half-offset lattice.
\begin{enumerate}
\item[(1)] $\{R \in O(4) : R\mathbb{Z}^4 \subseteq \mathbb{Z}^4\} = W_4$.
\item[(2)] Let $L \ge 3$ be odd, and let $g(x) = Rx + a$ be an isometry of $\mathbb{R}^4$ with
$g(\mathbb{L}_s) = \mathbb{L}_{s'}$. Then $s = s'$, $R \in W_4$ and $a \in L^{-s}\mathbb{Z}^4$.
Conversely, every such $g$ maps $\mathbb{L}_s$ onto itself.
\item[(3)] If $R \in O(4)\setminus W_4$ and $Q$ is a closed axis-parallel cube of positive side,
then $RQ$ is not a finite union of closed axis-parallel boxes.
\end{enumerate}
\end{lemma}

\begin{proof}
(1) Let $R \in O(4)$ with $R\mathbb{Z}^4 \subseteq \mathbb{Z}^4$. For each $\mu$ the column
$Re_\mu$ lies in $\mathbb{Z}^4$ and has Euclidean norm $1$, since $R$ is orthogonal. An integer
vector of norm $1$ has exactly one non-zero entry, equal to $\pm1$; hence
$Re_\mu = \pm e_{\nu(\mu)}$ for some index $\nu(\mu)$. The columns of $R$ are mutually orthogonal,
so $\nu(\mu) \ne \nu(\mu')$ for $\mu \ne \mu'$, and $\nu$ is a permutation of $\{1,2,3,4\}$. Thus
$R$ is a signed permutation matrix, $R \in W_4$. Conversely, a signed permutation matrix is
orthogonal and has integer entries, so it maps $\mathbb{Z}^4$ into $\mathbb{Z}^4$.

(2) The difference set of $(\mathbb{Z}+\tfrac12)^4$ with itself is $\mathbb{Z}^4$: coordinatewise,
$(n+\tfrac12)-(m+\tfrac12) = n-m$, and every integer $k$ equals $(k+\tfrac12)-\tfrac12$.
Scaling by $L^{-s}$ gives
$\mathbb{L}_s - \mathbb{L}_s = L^{-s}\mathbb{Z}^4$. Since the translation part cancels in
differences, $g(x)-g(y) = R(x-y)$, and therefore
\begin{equation*}
L^{-s'}\mathbb{Z}^4 = \mathbb{L}_{s'} - \mathbb{L}_{s'} = g(\mathbb{L}_s) - g(\mathbb{L}_s)
= R(\mathbb{L}_s - \mathbb{L}_s) = L^{-s}R\mathbb{Z}^4 ,
\end{equation*}
that is, $R\mathbb{Z}^4 = L^{s-s'}\mathbb{Z}^4$. Because $R$ is an isometry, the shortest non-zero
vectors of $R\mathbb{Z}^4$ have norm $1$, while those of $L^{s-s'}\mathbb{Z}^4$ have norm
$L^{s-s'}$. Hence $L^{s-s'} = 1$, and since $L \ge 3$, $s = s'$. Then $R\mathbb{Z}^4 = \mathbb{Z}^4$,
and $R \in W_4$ by (1).

Signed permutations preserve $(\mathbb{Z}+\tfrac12)^4$: a permutation of coordinates clearly does,
and a change of sign does because $-(n+\tfrac12) = (-n-1)+\tfrac12$. Hence
$R\mathbb{L}_s \subseteq \mathbb{L}_s$. Choose any $x \in \mathbb{L}_s$; then $g(x) \in \mathbb{L}_s$
(as $s'=s$) and $Rx \in \mathbb{L}_s$, so
\begin{equation*}
a = g(x) - Rx \in \mathbb{L}_s - \mathbb{L}_s = L^{-s}\mathbb{Z}^4 .
\end{equation*}
For the converse, let $R \in W_4$ and $a \in L^{-s}\mathbb{Z}^4$. By the preceding sentence,
$R\mathbb{L}_s \subseteq \mathbb{L}_s$ and $R^{-1}\mathbb{L}_s \subseteq \mathbb{L}_s$, since
$R^{-1} \in W_4$; so $R\mathbb{L}_s = \mathbb{L}_s$. Translation by $a = L^{-s}k$ with
$k \in \mathbb{Z}^4$ maps $L^{-s}(\mathbb{Z}+\tfrac12)^4$ bijectively onto itself, since
$(n+\tfrac12)+k_\mu = (n+k_\mu)+\tfrac12$ in each coordinate. Hence $g(\mathbb{L}_s) = \mathbb{L}_s$.

(3) Suppose that $RQ = \bigcup_{i \le k} B_i$ with closed axis-parallel boxes $B_i$. The set $RQ$
is closed, so each boundary point of $RQ$ lies in some $B_j$; it does not lie in the interior of
$B_j$, because that interior is contained in the interior of $RQ$. Hence it lies in $\partial B_j$,
which is contained in the union of the coordinate hyperplanes carrying the facets of $B_j$.
Consequently $\partial(RQ)$ is contained in a finite union of coordinate hyperplanes
$P_1,\dots,P_r$, each of the form $\{x_\nu = t\}$.

Since $R \notin W_4$, some column $n = Re_\mu$ is not of the form $\pm e_\nu$: otherwise, as in (1),
orthogonality of the columns would make $R$ a signed permutation matrix. Let $F$ be a facet of $Q$
with normal $e_\mu$; it is a three-dimensional cube of positive side. As $R$ is a homeomorphism,
$\partial(RQ) = R\,\partial Q \supseteq RF$, and $RF$ is a facet of $RQ$ with normal $n$, lying in a
hyperplane $H$ with normal $n$ and having positive three-dimensional measure in $H$. A coordinate
hyperplane $P_i = \{x_\nu = t\}$ has normal $e_\nu$; as $n$ and $e_\nu$ are unit vectors with
$n \ne \pm e_\nu$, they are linearly independent, so $P_i \ne H$ and $P_i \cap H$ is an affine
two-plane, which is null for the three-dimensional measure of $H$. Therefore
\begin{equation*}
RF \subseteq H \cap \partial(RQ) \subseteq \bigcup_{i \le r} (H \cap P_i)
\end{equation*}
is contained in a finite union of null subsets of $H$, contradicting the positive measure of $RF$.
\end{proof}

\begin{corollary}[No transport of block observables under non-hypercubic rotations]
\label{10.7:cor-no-block-transport}
Let $L\ge 3$ be odd, as in Lemma~\ref{10.7:lem-lattice-stabiliser}.
\begin{enumerate}
\item[(a)] An isometry $g(x)=Rx+a$ of $\mathbb{R}^4$ maps some $\mathbb{L}_s$ onto some
$\mathbb{L}_{s'}$ if and only if $g\in\mathfrak{G}$. In that case $g$ maps $\mathbb{L}_K$
onto itself for every $K\ge t(g)$, where
\begin{equation*}
t(g):=\min\{t : a\in L^{-t}\mathbb{Z}^4\}.
\end{equation*}
\item[(b)] Let $\sigma$ be an isometry of $\mathbb{R}^4$ with linear part $R$. The prescription
\begin{equation}\label{10.7:eq-no-block-transport-1}
(\sigma U)(\langle x,y\rangle)=U(\langle \sigma^{-1}x,\sigma^{-1}y\rangle)
\end{equation}
defines a map of bare configurations at depth $K$ (configurations $U$ on the bonds
$\langle x,y\rangle$ with $x,y\in\mathbb{L}_K$) only when $\sigma$ preserves $\mathbb{L}_K$;
if $R\notin W_4$, this happens for no $K$. Consequently, for $R\notin W_4$ the transport
$F\mapsto F^{\sigma}$ of block observables has no extension to $\sigma$, and of the maps on the
symbols of the smeared block algebra of the pattern ``transport the observable $F$, move the
test function $\varphi$'' only the symbol rotation $\beta_R$ of
Definition~\ref{10.7:def-smeared-block-algebra} remains. A redefinition of the transport through
the block geometry is also unavailable: a rotated block is not a union of blocks.
\end{enumerate}
\end{corollary}

\begin{proof}
(a) If $g(\mathbb{L}_s)=\mathbb{L}_{s'}$, part~(2) of
Lemma~\ref{10.7:lem-lattice-stabiliser} gives $s=s'$, $R\in W_4$ and
$a\in L^{-s}\mathbb{Z}^4$; conversely, if $R\in W_4$ and $a\in L^{-s}\mathbb{Z}^4$ for some $s$,
the converse half of that part gives $g(\mathbb{L}_s)=\mathbb{L}_s$. These motions are the
elements of $\mathfrak{G}$. Let $K\ge t(g)$. Since $L^{K-t(g)}$ is a positive integer,
\begin{equation*}
a\in L^{-t(g)}\mathbb{Z}^4=L^{-K}\bigl(L^{K-t(g)}\mathbb{Z}^4\bigr)\subseteq L^{-K}\mathbb{Z}^4,
\end{equation*}
and the converse half of part~(2) of Lemma~\ref{10.7:lem-lattice-stabiliser}, applied with
$s=K$, gives $g(\mathbb{L}_K)=\mathbb{L}_K$.

(b) For the right-hand side of \eqref{10.7:eq-no-block-transport-1} to be defined for every
bond of the depth-$K$ lattice, $\sigma^{-1}x$ must lie in $\mathbb{L}_K$ for every
$x\in\mathbb{L}_K$, that is, $\sigma^{-1}(\mathbb{L}_K)\subseteq\mathbb{L}_K$. Since
$\mathbb{L}_K=L^{-K}(\mathbb{Z}+\tfrac12)^4$, the differences of points of $\mathbb{L}_K$ form
exactly $L^{-K}\mathbb{Z}^4$, and $\sigma^{-1}x-\sigma^{-1}y=R^{-1}(x-y)$. Hence
$R^{-1}L^{-K}\mathbb{Z}^4\subseteq L^{-K}\mathbb{Z}^4$, i.e.\
$R^{-1}\mathbb{Z}^4\subseteq\mathbb{Z}^4$. By part~(1) of
Lemma~\ref{10.7:lem-lattice-stabiliser}, $R^{-1}\in W_4$, so $R\in W_4$. Moreover, a signed
permutation matrix maps $\mathbb{Z}^4$ onto itself; fixing $x_0\in\mathbb{L}_K$ and using
$\mathbb{L}_K=x_0+L^{-K}\mathbb{Z}^4$ and $\sigma^{-1}x_0\in\mathbb{L}_K$,
\begin{equation*}
\sigma^{-1}(\mathbb{L}_K)=\sigma^{-1}x_0+R^{-1}L^{-K}\mathbb{Z}^4
=\sigma^{-1}x_0+L^{-K}\mathbb{Z}^4=\mathbb{L}_K .
\end{equation*}
So \eqref{10.7:eq-no-block-transport-1} defines a map at depth $K$ only when $\sigma$ preserves
$\mathbb{L}_K$, and this forces $R\in W_4$; for $R\notin W_4$ it happens at no depth. The
transport $F\mapsto F^{\sigma}$ evaluates a block observable $F$ on the transformed configuration
$\sigma U$, which for $R\notin W_4$ is not a bare configuration at any depth; hence this transport
has no extension to such $\sigma$. In a map of the pattern ``transport $F$, move $\varphi$'' the
first component is therefore unavailable, and what remains is the motion of the test function,
which is the symbol rotation $\beta_R$ of Definition~\ref{10.7:def-smeared-block-algebra}.
Finally, a block is a closed axis-parallel cube of positive side, and by part~(3) of
Lemma~\ref{10.7:lem-lattice-stabiliser} its image under $R\in O(4)\setminus W_4$ is not a finite
union of closed axis-parallel boxes. The blocks of a given level are closed axis-parallel cubes of
one fixed positive side centred on a lattice, so only finitely many of them meet the bounded set
$RQ$, and a union of blocks equal to $RQ$ would be such a finite union. Hence a rotated block is
not a union of blocks, and the transport cannot be redefined through the block geometry.
\end{proof}

\begin{proposition}[The single-symbol rotation identity is empty]
\label{10.7:prop-single-symbol-identity}
Let $\mathfrak{S}$ be the set of generating symbols of the smeared block algebra
(Definition~\ref{10.7:def-smeared-block-algebra}), and let $\omega_\infty$ be the limit state on the
smeared block algebra $\mathfrak{A}^{sm}_{loc}$.
For every $(F,\varphi)\in\mathfrak{S}$ and every $T\in GL_4(\mathbb{R})$ with $|\det T|=1$, in particular
for every $R\in O(4)$,
\begin{equation}\label{10.7:eq-single-symbol-identity-1}
\omega_\infty\bigl((F,\varphi\circ T^{-1})\bigr)=\omega_\infty\bigl((F,\varphi)\bigr).
\end{equation}
\end{proposition}

\begin{proof}
Every $R\in O(4)$ satisfies $|\det R|=1$, so it suffices to treat a general $T\in GL_4(\mathbb{R})$ with
$|\det T|=1$.

On a generating symbol the limit state factorises; this is clause~(e) of the statement in this chapter
that lists the properties of $\omega_\infty$ on $\mathfrak{A}^{sm}_{loc}$, and it reads: for every
$(F,\psi)\in\mathfrak{S}$,
\begin{equation}\label{10.7:eq-single-symbol-identity-2}
\omega_\infty\bigl((F,\psi)\bigr)=\Bigl(\int_{\mathbb{R}^4}\psi(x)\,dx\Bigr)\,\omega_\infty(F).
\end{equation}
We apply \eqref{10.7:eq-single-symbol-identity-2} twice: once with $\psi=\varphi\circ T^{-1}$ and once with
$\psi=\varphi$. The two right-hand sides contain the same factor $\omega_\infty(F)$. Hence
\eqref{10.7:eq-single-symbol-identity-1} follows once the integrals of $\varphi\circ T^{-1}$ and of
$\varphi$ agree.

To compare the integrals we change variables by $x=Ty$. Then $dx=|\det T|\,dy$, and
\begin{equation*}
\int_{\mathbb{R}^4}\varphi(T^{-1}x)\,dx
=|\det T|\int_{\mathbb{R}^4}\varphi(y)\,dy
=\int_{\mathbb{R}^4}\varphi(y)\,dy,
\end{equation*}
where the last equality uses $|\det T|=1$. This gives \eqref{10.7:eq-single-symbol-identity-1}.

The identity \eqref{10.7:eq-single-symbol-identity-1} holds for volume-preserving shears just as it does for
rotations. It therefore carries no information about rotational symmetry.
\end{proof}

\begin{proposition}[The product identity as isotropy of probes]\label{10.7:prop-probe-isotropy}
Let $\omega_\infty$ be the limit state on the smeared block algebra $\mathfrak{A}^{sm}_{loc}$, and let
$\alpha_h$, $h \in \mathfrak{G}$, and $\beta_w$ be the automorphisms of
Definition~\ref{10.7:def-smeared-block-algebra}. Let $w \in W_4$, and let
$X = \prod_i (F_i,\varphi_i)$ be a finite product of generating symbols $(F_i,\varphi_i) \in \mathfrak{S}$.
Then
\begin{equation}\label{10.7:eq-probe-isotropy-1}
\omega_\infty(\beta_w X) = \omega_\infty\Bigl(\prod_i (\alpha_{w^{-1}} F_i, \varphi_i)\Bigr).
\end{equation}
\end{proposition}

\begin{proof}
Three facts are used.
\begin{enumerate}
\item By the invariance of the limit state under the commensurable motion group, which is part of
the conclusion step, for every $h \in \mathfrak{G}$,
\begin{equation*}
\omega_\infty \circ \alpha_h = \omega_\infty .
\end{equation*}
\item The automorphism $\alpha_{w^{-1}}$ is multiplicative.
\item On a generating symbol it acts by
\begin{equation*}
\alpha_{w^{-1}}(F,\psi) = (\alpha_{w^{-1}}F,\ \psi \circ w).
\end{equation*}
\end{enumerate}

The element $w^{-1}$ lies in $W_4$, and $W_4 \subset \mathfrak{G}$. The invariance, applied with
$h = w^{-1}$ to $\beta_w X$, therefore gives
\begin{equation*}
\omega_\infty(\beta_w X) = \omega_\infty(\alpha_{w^{-1}} \beta_w X).
\end{equation*}
By Definition~\ref{10.7:def-smeared-block-algebra}, the symbol rotation $\beta_w$ replaces each factor
$(F_i,\varphi_i)$ of $X$ by $(F_i, \varphi_i \circ w^{-1})$.

Since $\alpha_{w^{-1}}$ is multiplicative, it acts on $\beta_w X$ factor by factor. Its action on a
generating symbol, taken with $F = F_i$ and $\psi = \varphi_i \circ w^{-1}$, then gives
\begin{equation*}
\alpha_{w^{-1}} \beta_w X = \prod_i \bigl(\alpha_{w^{-1}} F_i,\ \varphi_i \circ w^{-1} \circ w\bigr)
= \prod_i (\alpha_{w^{-1}} F_i, \varphi_i).
\end{equation*}
Inserting this into the previous identity yields \eqref{10.7:eq-probe-isotropy-1}.
\end{proof}

Read through \eqref{10.7:eq-probe-isotropy-1}, the sentence $(S_\beta)$,
\eqref{10.7:eq-smeared-block-algebra-a}, read at $w \in W_4$, asserts the following. Truncated
correlations are unchanged when every probe is rotated by $w^{-1}$ about its own position, the
positions being held fixed. This is an isotropy property of the probes. It is not the invariance of
the state, which is the identity $\omega_\infty \circ \alpha_w = \omega_\infty$ and holds exactly.
For a rotation $R \notin W_4$ no such reformulation of $(S_\beta)$ at $R$ is available. Nothing
proved here decides $(S_\beta)$ for the limit state $\omega_\infty$, and $(S_\beta)$ is not asserted.

\begin{lemma}[Positivity of the free abelian covariance]\label{10.7:lem-abelian-covariance-positive}
On $\mathbb{R}^4$ put, for $z \neq 0$ and $t > 0$,
\begin{equation}\label{10.7:eq-abelian-covariance-positive-1}
G(z) := \frac{1}{4\pi^2 |z|^2}, \qquad
p_t(z) := (4\pi t)^{-2} e^{-|z|^2/4t} .
\end{equation}
Let $f \colon \mathbb{R}^4 \to \mathbb{R}$ be bounded, measurable and compactly supported. Then the
integral
\begin{equation*}
Q(f) := \int_{\mathbb{R}^4}\int_{\mathbb{R}^4} f(x)\, G(x-y)\, f(y)\, d^4x\, d^4y
\end{equation*}
is absolutely convergent, and
\begin{equation}\label{10.7:eq-abelian-covariance-positive-2}
Q(f) = \int_0^\infty \bigl\| p_{t/2} \ast f \bigr\|_{L^2(\mathbb{R}^4)}^2 \, dt \;\ge\; 0 .
\end{equation}
\end{lemma}

\begin{proof}
We first note three properties of the kernel $p_t$. Since
$p_t(z) = \prod_{\alpha=1}^4 (4\pi t)^{-1/2} e^{-z_\alpha^2/4t}$ and
$\int_{\mathbb{R}} e^{-u^2/4t}\, du = (4\pi t)^{1/2}$, we have $\int p_t = 1$; clearly
$p_t(-z) = p_t(z)$. For the semigroup property $p_s \ast p_t = p_{s+t}$ it suffices, by the same
product structure, to treat one coordinate. Completing the square,
\begin{equation*}
\frac{(u-v)^2}{s} + \frac{v^2}{t}
= \frac{s+t}{st}\Bigl(v - \frac{tu}{s+t}\Bigr)^2 + \frac{u^2}{s+t},
\end{equation*}
so that
\begin{equation*}
\int_{\mathbb{R}} (4\pi s)^{-1/2} e^{-(u-v)^2/4s} (4\pi t)^{-1/2} e^{-v^2/4t}\, dv
= (4\pi s)^{-1/2}(4\pi t)^{-1/2}\Bigl(\frac{4\pi st}{s+t}\Bigr)^{1/2} e^{-u^2/4(s+t)} ,
\end{equation*}
which is $(4\pi(s+t))^{-1/2} e^{-u^2/4(s+t)}$. Finally, for $z \neq 0$ the substitution
$t = |z|^2/4u$, $dt = -|z|^2 (4u^2)^{-1} du$, $(4\pi t)^{-2} = u^2 \pi^{-2}|z|^{-4}$, gives
\begin{equation}\label{10.7:eq-abelian-covariance-positive-3}
\int_0^\infty p_t(z)\, dt = \int_0^\infty \frac{e^{-u}}{4\pi^2 |z|^2}\, du = G(z).
\end{equation}

Absolute convergence. Choose $\rho_f > 0$ with $f = 0$ outside the ball $\{|x| \le \rho_f\}$. In polar
coordinates, the three-sphere having area $2\pi^2$, for every $\rho > 0$
\begin{equation*}
\int_{|z| \le \rho} |z|^{-2}\, d^4z = 2\pi^2 \int_0^\rho r^{-2} r^3\, dr = \pi^2 \rho^2 .
\end{equation*}
If $|x|, |y| \le \rho_f$ then $|x - y| \le 2\rho_f$, hence, by Tonelli's theorem,
\begin{align*}
\int\!\!\int |f(x)|\, G(x-y)\, |f(y)|\, d^4x\, d^4y
&\le \|f\|_\infty^2 \int_{|x|\le \rho_f} \int_{|z| \le 2\rho_f} G(z)\, d^4z\, d^4x \\
&= \|f\|_\infty^2 \cdot \frac{\pi^2 \rho_f^4}{2} \cdot \frac{\pi^2 (2\rho_f)^2}{4\pi^2}
 < \infty .
\end{align*}

The identity. Fix $t > 0$. Since $p_{t/2}$ is integrable and $f$ is bounded with compact support,
$p_{t/2} \ast f(z) = \int p_{t/2}(z-x) f(x)\, d^4x$ is defined for every $z$, and
$\|p_{t/2} \ast f\|_{L^2} \le \|p_{t/2}\|_{L^1} \|f\|_{L^2} < \infty$. By evenness of $p_{t/2}$ and
the semigroup property,
\begin{equation*}
\int p_{t/2}(z-x)\, p_{t/2}(z-y)\, d^4z
= \int p_{t/2}(x-z)\, p_{t/2}(z-y)\, d^4z = p_t(x-y) .
\end{equation*}
Consequently
\begin{align*}
\iiint p_{t/2}(z-x)\,|f(x)|\,p_{t/2}(z-y)\,|f(y)|\, d^4z\, d^4x\, d^4y
&= \iint |f(x)|\, p_t(x-y)\, |f(y)|\, d^4x\, d^4y \\
&\le (4\pi t)^{-2}\|f\|_{L^1}^2 < \infty ;
\end{align*}
so Fubini's theorem applies to the expansion of the square and yields
\begin{equation}\label{10.7:eq-abelian-covariance-positive-4}
\bigl\| p_{t/2} \ast f \bigr\|_{L^2}^2 = \int\!\!\int f(x)\, p_t(x-y)\, f(y)\, d^4x\, d^4y .
\end{equation}
The function $(t, x, y) \mapsto f(x) p_t(x-y) f(y)$ is measurable on
$(0,\infty) \times \mathbb{R}^4 \times \mathbb{R}^4$. By Tonelli's theorem and
\eqref{10.7:eq-abelian-covariance-positive-3}, applied off the diagonal $\{x = y\}$, which is a null
set,
\begin{equation*}
\int_0^\infty \!\! \int\!\!\int |f(x)|\, p_t(x-y)\, |f(y)|\, d^4x\, d^4y\, dt
= \int\!\!\int |f(x)|\, G(x-y)\, |f(y)|\, d^4x\, d^4y ,
\end{equation*}
which is finite by the absolute convergence shown above. Hence Fubini's theorem permits the
interchange of the $t$-integration with the $(x,y)$-integration in the signed integrand:
integrating \eqref{10.7:eq-abelian-covariance-positive-4} over $t \in (0,\infty)$ and using
\eqref{10.7:eq-abelian-covariance-positive-3} again gives
$\int_0^\infty \|p_{t/2} \ast f\|_{L^2}^2\, dt = Q(f)$. The integrand on the left is non-negative
for every $t$, so $Q(f) \ge 0$, which is \eqref{10.7:eq-abelian-covariance-positive-2}.
\end{proof}

We note two further facts about the kernel $G$ of
\eqref{10.7:eq-abelian-covariance-positive-1} and the Gaussian field built on it. The form $C$
defined below is the Maxwell covariance, and the centred Gaussian field $A$ with covariance $C$
constructed below is the free Maxwell field; it serves in this section as a counter-model.

First, for $z \neq 0$ and $\alpha, \beta \in \{1,2,3,4\}$ we have
$\partial_\beta |z|^{-2} = -2 z_\beta |z|^{-4}$ and
$\partial_\alpha\bigl(z_\beta |z|^{-4}\bigr) = \delta_{\alpha\beta}|z|^{-4} - 4 z_\alpha z_\beta |z|^{-6}$.
Hence
\begin{equation}\label{10.7:eq-abelian-covariance-positive-5}
\partial_\alpha \partial_\beta G(z)
= \frac{1}{4\pi^2}\Bigl( -2\,\delta_{\alpha\beta}\, |z|^{-4} + 8\, z_\alpha z_\beta\, |z|^{-6} \Bigr) .
\end{equation}

Second, let $\mathcal{V} := C_c(\mathbb{R}^4;\mathbb{R}^4)$ be the space of continuous, compactly
supported vector-valued test functions $f = (f_1,\dots,f_4)$, and put
\begin{equation}\label{10.7:eq-abelian-covariance-positive-6}
C(f,g) := \sum_{\mu=1}^4 \int_{\mathbb{R}^4}\int_{\mathbb{R}^4}
f_\mu(x)\, G(x-y)\, g_\mu(y)\, d^4x\, d^4y ,
\qquad f, g \in \mathcal{V} .
\end{equation}
Each integral converges absolutely: with $h := |f_\mu| + |g_\mu|$, which is bounded, measurable and
compactly supported, we have $|f_\mu(x)|\,|g_\mu(y)| \le h(x) h(y)$, and the double integral of
$h(x) G(x-y) h(y)$ is finite by Lemma~\ref{10.7:lem-abelian-covariance-positive}. The form $C$ is
bilinear, and it is symmetric because $G(-z) = G(z)$. It is positive semidefinite: for
$f^{(1)}, \dots, f^{(n)} \in \mathcal{V}$ and real $c_1, \dots, c_n$, bilinearity gives
\begin{equation*}
\sum_{i,j=1}^n c_i c_j\, C(f^{(i)}, f^{(j)}) = C(f,f)
= \sum_{\mu=1}^4 Q(f_\mu) \ge 0 ,
\qquad f := \sum_{i=1}^n c_i f^{(i)} ,
\end{equation*}
where each component $f_\mu$ of $f$ is bounded, measurable and compactly supported, so that
Lemma~\ref{10.7:lem-abelian-covariance-positive} applies to it.

Consequently there is a centred Gaussian process $(A(f))_{f \in \mathcal{V}}$ with covariance $C$,
which is linear almost surely. Indeed, for $f^{(1)}, \dots, f^{(n)} \in \mathcal{V}$ the matrix
$\Gamma := (C(f^{(i)},f^{(j)}))_{i,j=1}^n$ is real, symmetric and positive semidefinite, so it has a
symmetric positive semidefinite square root $\Gamma^{1/2}$; if $\eta$ is a standard Gaussian vector
in $\mathbb{R}^n$, then $\Gamma^{1/2}\eta$ is a centred Gaussian vector with covariance
$\Gamma^{1/2}(\Gamma^{1/2})^{\mathsf{T}} = \Gamma$. These finite-dimensional laws are consistent:
permuting $f^{(1)}, \dots, f^{(n)}$ permutes the entries of $\Gamma$ accordingly, and the marginal of a
centred Gaussian vector with covariance $\Gamma$ on a subfamily of coordinates is the centred
Gaussian law whose covariance is the corresponding principal submatrix of $\Gamma$. Kolmogorov's
extension theorem therefore yields the process. For $f, g \in \mathcal{V}$ and real $a, b$, the
variable $A(af+bg) - aA(f) - bA(g)$ is a linear combination of the entries of the centred Gaussian
vector $(A(af+bg), A(f), A(g))$, hence centred Gaussian, and by bilinearity of $C$ its variance is
\begin{equation*}
C\bigl(af+bg-af-bg,\, af+bg-af-bg\bigr) = C(0,0) = 0 ;
\end{equation*}
hence $A(af+bg) = aA(f) + bA(g)$ almost surely.

Finally, the orthogonal group $O(4)$ acts on $\mathcal{V}$ by
\begin{equation*}
(R \cdot f)_\mu(x) := \sum_{\nu=1}^4 R_{\mu\nu}\, f_\nu(R^{-1}x),
\qquad R \in O(4),\ f \in \mathcal{V} .
\end{equation*}

\begin{lemma}[Euclidean invariance in law of the free abelian field]
\label{10.7:lem-abelian-invariance}
For $R\in O(4)$ and a vector test function $f\in\mathcal{V}$ put
$(R\cdot f)(x)=R\,f(R^{-1}x)$, and for $a\in\mathbb{R}^4$ put
$(\tau_a f)(x)=f(x-a)$. Then for all $f,g\in\mathcal{V}$ and all $R\in O(4)$
\begin{equation}\label{10.7:eq-abelian-invariance-1}
C(R\cdot f,R\cdot g)=C(f,g).
\end{equation}
Consequently the Gaussian field $A$ with covariance $C$ is invariant in law under
$\mathbb{R}^4\rtimes O(4)$.
\end{lemma}

\begin{proof}
Here $C(f,g)$ is the double integral of $f(x)\cdot g(y)$ against $G(x-y)$. For bounded,
compactly supported $f$ and $g$ it converges absolutely. To see this, apply
Lemma~\ref{10.7:lem-abelian-covariance-positive} to $h=|f|+|g|$. Then
$|f(x)\cdot g(y)|\,|G(x-y)|\le h(x)\,|G(x-y)|\,h(y)$, and the right-hand side is
integrable. The changes of variables below are therefore legitimate.

Since $R^{-1}=R^{\mathsf T}$, we have
\begin{equation*}
C(R\cdot f,R\cdot g)
=\int\!\!\int \bigl(R f(R^{\mathsf T}x)\bigr)\cdot\bigl(R g(R^{\mathsf T}y)\bigr)\,
G(x-y)\,d^4x\,d^4y .
\end{equation*}
Substitute $x=Rx'$ and $y=Ry'$, whose Jacobian is $|\det R|=1$. The relation
$Ru\cdot Rv=u\cdot R^{\mathsf T}R\,v$ together with $R^{\mathsf T}R=\mathbb{1}$ gives
$Rf(x')\cdot Rg(y')=f(x')\cdot g(y')$. Since $G$ is radial,
$G(R(x'-y'))=G(x'-y')$. Hence
\begin{equation*}
C(R\cdot f,R\cdot g)
=\int\!\!\int f(x')\cdot g(y')\,G(x'-y')\,d^4x'\,d^4y'
=C(f,g),
\end{equation*}
which is \eqref{10.7:eq-abelian-invariance-1}.

For translations, substitute $x=x'+a$ and $y=y'+a$. Because $G$ depends only on $x-y$,
this gives $C(\tau_a f,\tau_a g)=C(f,g)$.

Every element of $\mathbb{R}^4\rtimes O(4)$ is a composition of a rotation and a
translation, so $C$ is invariant under the whole group. The law of the Gaussian field
$A$ is determined by its covariance $C$. Hence $\bigl(A(R\cdot f)\bigr)_{f\in\mathcal{V}}$
and $\bigl(A(\tau_a f)\bigr)_{f\in\mathcal{V}}$ have the same law as
$\bigl(A(f)\bigr)_{f\in\mathcal{V}}$.
\end{proof}

\begin{lemma}[Covariance of the smeared free field strengths]
\label{10.7:lem-field-strength-covariance}
For a real $\varphi\in C^1_c(\mathbb{R}^4)$ and $\mu,\nu\in\{1,2,3,4\}$ let $f^{\mu\nu,\varphi}$ be the
vector test function with components
\begin{equation}\label{10.7:eq-field-strength-covariance-1}
(f^{\mu\nu,\varphi})_\lambda:=\delta_{\lambda\mu}\,\partial_\nu\varphi-\delta_{\lambda\nu}\,\partial_\mu\varphi ,
\qquad \lambda=1,\dots,4,
\end{equation}
and put $F_{\mu\nu}(\varphi):=A(f^{\mu\nu,\varphi})$. Formally,
$F_{\mu\nu}(\varphi)=\int\varphi\,(\partial_\mu A_\nu-\partial_\nu A_\mu)$.
Let $R\in O(4)$. On scalar test functions put $(R\cdot\varphi)(x)=\varphi(R^{-1}x)$, and on vector
test functions put $(R\cdot f)_\lambda(x)=\sum_\kappa R_{\lambda\kappa}f_\kappa(R^{-1}x)$. Then
\begin{equation}\label{10.7:eq-field-strength-covariance-2}
R\cdot f^{\mu\nu,\varphi}=\sum_{\alpha,\beta}R_{\alpha\mu}R_{\beta\nu}\,f^{\alpha\beta,R\cdot\varphi}.
\end{equation}
Consequently, with
\begin{equation*}
F^R_{\mu\nu}(\varphi):=\sum_{\alpha,\beta}R_{\alpha\mu}R_{\beta\nu}\,F_{\alpha\beta}(R\cdot\varphi),
\end{equation*}
one has, for all real $\varphi,\psi\in C^1_c(\mathbb{R}^4)$ and all indices $\mu,\nu,\rho,\sigma$,
\begin{equation}\label{10.7:eq-field-strength-covariance-3}
E\bigl[F^R_{\mu\nu}(\varphi)\,F^R_{\rho\sigma}(\psi)\bigr]=E\bigl[F_{\mu\nu}(\varphi)\,F_{\rho\sigma}(\psi)\bigr].
\end{equation}
\end{lemma}

\begin{proof}
Write $\chi:=R\cdot\varphi$, so that $\chi(x)=\varphi(R^{-1}x)$. By the chain rule, and because
$(R^{-1})_{\nu\beta}=R_{\beta\nu}$ when $R^{\mathsf T}R=\mathbb{1}$,
\begin{equation*}
\partial_\beta\chi(x)=\sum_\nu(\partial_\nu\varphi)(R^{-1}x)\,(R^{-1})_{\nu\beta}
=\sum_\nu R_{\beta\nu}\,(\partial_\nu\varphi)(R^{-1}x).
\end{equation*}
In vector form this reads $\nabla\chi(x)=R\,(\nabla\varphi)(R^{-1}x)$. Multiplying by
$R^{\mathsf T}=R^{-1}$ gives
\begin{equation}\label{10.7:eq-field-strength-covariance-4}
(\partial_\nu\varphi)(R^{-1}x)=\sum_\beta R_{\beta\nu}\,\partial_\beta(R\cdot\varphi)(x).
\end{equation}

We now compare the two sides of \eqref{10.7:eq-field-strength-covariance-2} component by
component. By \eqref{10.7:eq-field-strength-covariance-1}, the left-hand side has components
\begin{equation*}
(R\cdot f^{\mu\nu,\varphi})_\lambda(x)
=R_{\lambda\mu}(\partial_\nu\varphi)(R^{-1}x)-R_{\lambda\nu}(\partial_\mu\varphi)(R^{-1}x).
\end{equation*}
The right-hand side has components
\begin{align*}
\sum_{\alpha,\beta}R_{\alpha\mu}R_{\beta\nu}\bigl(\delta_{\lambda\alpha}\partial_\beta\chi
-\delta_{\lambda\beta}\partial_\alpha\chi\bigr)(x)
&=R_{\lambda\mu}\sum_\beta R_{\beta\nu}\partial_\beta\chi(x)
-R_{\lambda\nu}\sum_\alpha R_{\alpha\mu}\partial_\alpha\chi(x)\\
&=R_{\lambda\mu}(\partial_\nu\varphi)(R^{-1}x)-R_{\lambda\nu}(\partial_\mu\varphi)(R^{-1}x).
\end{align*}
The second line follows by applying \eqref{10.7:eq-field-strength-covariance-4} once with the
index $\nu$ and once with the index $\mu$. The two sides therefore agree, which proves
\eqref{10.7:eq-field-strength-covariance-2}.

Next, $f\mapsto A(f)$ is linear. Hence, by the definition of $F_{\alpha\beta}$ and by
\eqref{10.7:eq-field-strength-covariance-2},
\begin{equation*}
F^R_{\mu\nu}(\varphi)=A\Bigl(\sum_{\alpha,\beta}R_{\alpha\mu}R_{\beta\nu}f^{\alpha\beta,R\cdot\varphi}\Bigr)
=A\bigl(R\cdot f^{\mu\nu,\varphi}\bigr).
\end{equation*}
The same holds for $F^R_{\rho\sigma}(\psi)=A(R\cdot f^{\rho\sigma,\psi})$. The covariance of the
Gaussian field $A$ is $E[A(f)A(g)]=C(f,g)$. Using this and then the Euclidean invariance in law
of the free abelian field (Lemma~\ref{10.7:lem-abelian-invariance}), we obtain
\begin{align*}
E\bigl[F^R_{\mu\nu}(\varphi)F^R_{\rho\sigma}(\psi)\bigr]
&=C\bigl(R\cdot f^{\mu\nu,\varphi},R\cdot f^{\rho\sigma,\psi}\bigr)
=C\bigl(f^{\mu\nu,\varphi},f^{\rho\sigma,\psi}\bigr)\\
&=E\bigl[F_{\mu\nu}(\varphi)F_{\rho\sigma}(\psi)\bigr].
\end{align*}
This is \eqref{10.7:eq-field-strength-covariance-3}.
\end{proof}

\begin{lemma}[The field-strength kernel of the free abelian field]
\label{10.7:lem-field-strength-kernel}
Let $\varphi,\psi\in C^1_c(\mathbb{R}^4)$ have disjoint supports, and let
$\mu,\nu,\rho,\sigma\in\{1,2,3,4\}$. Then
\begin{equation}
\label{10.7:eq-field-strength-kernel-1}
\mathbb{E}\bigl[F_{\mu\nu}(\varphi)\,F_{\rho\sigma}(\psi)\bigr]
=\int\!\!\int \varphi(x)\,\psi(y)\,K_{\mu\nu,\rho\sigma}(x-y)\,d^4x\,d^4y ,
\end{equation}
where, for $z\in\mathbb{R}^4\setminus\{0\}$,
\begin{equation}
\label{10.7:eq-field-strength-kernel-2}
\begin{split}
K_{\mu\nu,\rho\sigma}(z)=\pi^{-2}\Bigl[
&(\delta_{\mu\rho}\delta_{\nu\sigma}-\delta_{\mu\sigma}\delta_{\nu\rho})\,|z|^{-4}\\
&-2\,(\delta_{\mu\rho}z_\nu z_\sigma-\delta_{\mu\sigma}z_\nu z_\rho
-\delta_{\nu\rho}z_\mu z_\sigma+\delta_{\nu\sigma}z_\mu z_\rho)\,|z|^{-6}\Bigr].
\end{split}
\end{equation}
In particular the function $k(z):=K_{12,12}(z)$ is
\begin{equation}
\label{10.7:eq-field-strength-kernel-3}
k(z)=\frac{z_3^2+z_4^2-z_1^2-z_2^2}{\pi^2\,|z|^6},
\end{equation}
and $\sum_{\mu<\nu}K_{\mu\nu,\mu\nu}(z)=0$ for every $z\neq0$.
\end{lemma}

\begin{proof}
Throughout, $G(z)=(4\pi^2|z|^2)^{-1}$ is the covariance kernel of the free abelian field. For
$\alpha,\beta\in\{1,2,3,4\}$ put
\begin{equation*}
I_{\alpha\beta}:=\int\!\!\int \partial_\alpha\varphi(x)\,G(x-y)\,\partial_\beta\psi(y)\,d^4x\,d^4y .
\end{equation*}
We write the smeared field strengths through the Gaussian field smeared with the derivatives of
the test function and insert the covariance $G$. This gives
\begin{equation}
\label{10.7:eq-field-strength-kernel-4}
\mathbb{E}\bigl[F_{\mu\nu}(\varphi)F_{\rho\sigma}(\psi)\bigr]
=\delta_{\mu\rho}I_{\nu\sigma}-\delta_{\mu\sigma}I_{\nu\rho}
-\delta_{\nu\rho}I_{\mu\sigma}+\delta_{\nu\sigma}I_{\mu\rho}.
\end{equation}

The integrand of $I_{\alpha\beta}$ vanishes unless $x\in\operatorname{supp}\varphi$ and
$y\in\operatorname{supp}\psi$. These supports are compact and disjoint, so they lie at positive
distance from each other. Hence $(x,y)\mapsto G(x-y)$ is smooth on a neighbourhood of
$\operatorname{supp}\varphi\times\operatorname{supp}\psi$. We integrate by parts twice.

First, for fixed $y$ we integrate by parts in $x$. This moves $\partial_\alpha$ from $\varphi$
onto $x\mapsto G(x-y)$, with a minus sign, and produces $(\partial_\alpha G)(x-y)$.

Second, for fixed $x$ we integrate by parts in $y$. This moves $\partial_\beta$ from $\psi$
onto $y\mapsto(\partial_\alpha G)(x-y)$, with a second minus sign. Since
$\partial_{y_\beta}\bigl[(\partial_\alpha G)(x-y)\bigr]=-(\partial_\beta\partial_\alpha G)(x-y)$,
a third sign appears. Altogether
\begin{equation}
\label{10.7:eq-field-strength-kernel-5}
I_{\alpha\beta}=-\int\!\!\int\varphi(x)\,\psi(y)\,(\partial_\alpha\partial_\beta G)(x-y)\,d^4x\,d^4y .
\end{equation}

For $z\neq0$ we have $\partial_\alpha G(z)=-z_\alpha/(2\pi^2|z|^4)$. Differentiating once more,
\begin{equation*}
-(\partial_\alpha\partial_\beta G)(z)
=\frac{1}{2\pi^2}\Bigl[\delta_{\alpha\beta}|z|^{-4}-4z_\alpha z_\beta|z|^{-6}\Bigr].
\end{equation*}
We insert this into \eqref{10.7:eq-field-strength-kernel-5} and then into
\eqref{10.7:eq-field-strength-kernel-4}, and collect terms.

The coefficient of $|z|^{-4}$ is
\begin{equation*}
\tfrac{1}{2\pi^2}\bigl(\delta_{\mu\rho}\delta_{\nu\sigma}-\delta_{\mu\sigma}\delta_{\nu\rho}
-\delta_{\nu\rho}\delta_{\mu\sigma}+\delta_{\nu\sigma}\delta_{\mu\rho}\bigr)
=\pi^{-2}(\delta_{\mu\rho}\delta_{\nu\sigma}-\delta_{\mu\sigma}\delta_{\nu\rho}).
\end{equation*}
The coefficient of $|z|^{-6}$ is
\begin{equation*}
-\tfrac{4}{2\pi^2}\bigl(\delta_{\mu\rho}z_\nu z_\sigma-\delta_{\mu\sigma}z_\nu z_\rho
-\delta_{\nu\rho}z_\mu z_\sigma+\delta_{\nu\sigma}z_\mu z_\rho\bigr).
\end{equation*}
This is \eqref{10.7:eq-field-strength-kernel-2}, and \eqref{10.7:eq-field-strength-kernel-1}
follows. The double integral converges absolutely, because $K_{\mu\nu,\rho\sigma}(x-y)$ is
continuous on the compact set $\operatorname{supp}\varphi\times\operatorname{supp}\psi$.

For $(\mu\nu,\rho\sigma)=(12,12)$, formula \eqref{10.7:eq-field-strength-kernel-2} gives
\begin{equation*}
K_{12,12}(z)=\pi^{-2}\bigl[|z|^{-4}-2(z_2^2+z_1^2)|z|^{-6}\bigr]
=\pi^{-2}|z|^{-6}\bigl(|z|^2-2z_1^2-2z_2^2\bigr).
\end{equation*}
This is \eqref{10.7:eq-field-strength-kernel-3}.

Finally, take $\mu<\nu$, so that $\delta_{\mu\nu}=0$. Then
\eqref{10.7:eq-field-strength-kernel-2} gives
\begin{equation*}
K_{\mu\nu,\mu\nu}(z)=\pi^{-2}\bigl[|z|^{-4}-2(z_\mu^2+z_\nu^2)|z|^{-6}\bigr].
\end{equation*}
We sum over the six pairs $\mu<\nu$. Each index lies in exactly three pairs, so each $z_\mu^2$
occurs three times, and $\sum_{\mu<\nu}(z_\mu^2+z_\nu^2)=3|z|^2$. Hence
\begin{equation*}
\sum_{\mu<\nu}K_{\mu\nu,\mu\nu}(z)=\pi^{-2}\bigl[6-6\bigr]|z|^{-4}=0
\qquad(z\neq0).
\qedhere
\end{equation*}
\end{proof}

\begin{proposition}[Failure of probe isotropy under a quarter turn]
\label{10.7:prop-isotropy-fails-hypercubic}
Let $r>0$ and $0<\delta<r/4$. Let $\varphi\in C^1_c(\mathbb{R}^4)$ be radial, with $\varphi\ge 0$,
$\operatorname{supp}\varphi\subset\bar B(0,\delta)$ and $\int\varphi>0$, and put
$\psi:=\varphi(\cdot-re_3)$. Let $R$ be the quarter turn in the $(1,3)$-plane, $Re_1=e_3$,
$Re_3=-e_1$, $Re_2=e_2$, $Re_4=e_4$; thus $R\in W_4$. Put
\begin{equation*}
c_{r,\delta}:=\frac{\bigl(\int\varphi\bigr)^2\,(r^2-4r\delta)}{\pi^2\,(r+2\delta)^6}>0 .
\end{equation*}
Then, with $R\cdot f:=f\circ R^{-1}$, so that $R\cdot\varphi=\varphi$ and
$R\cdot\psi=\varphi(\cdot+re_1)$:
\begin{enumerate}
\item[(a)] $\mathbb{E}[F_{12}(\varphi)F_{12}(\psi)]\ge c_{r,\delta}$;
\item[(b)] $\mathbb{E}[F_{12}(R\cdot\varphi)F_{12}(R\cdot\psi)]\le -c_{r,\delta}$;
\item[(c)] $\mathbb{E}[F_{32}(R\cdot\varphi)F_{32}(R\cdot\psi)]=\mathbb{E}[F_{12}(\varphi)F_{12}(\psi)]$.
\end{enumerate}
\end{proposition}

\begin{proof}
We first unwind $R\cdot\varphi$ and $R\cdot\psi$. Since $\varphi$ is radial and $R$ is orthogonal,
$R\cdot\varphi=\varphi$. Since $rRe_3=-re_1$ and $\varphi$ is radial,
\begin{equation*}
(R\cdot\psi)(x)=\varphi(R^{-1}x-re_3)=\varphi\bigl(R^{-1}(x+re_1)\bigr)=\varphi(x+re_1).
\end{equation*}
Hence $\operatorname{supp}\psi\subset\bar B(re_3,\delta)$ and
$\operatorname{supp}R\cdot\psi\subset\bar B(-re_1,\delta)$; as $r>2\delta$, each of these balls is
disjoint from $\bar B(0,\delta)$. The hypotheses of Lemma~\ref{10.7:lem-field-strength-kernel} are
therefore met for both pairs, and that lemma gives
\begin{equation*}
\mathbb{E}[F_{12}(\varphi)F_{12}(\chi)]=\iint\varphi(x)\chi(y)\,k(x-y)\,dx\,dy,
\qquad \chi\in\{\psi,\,R\cdot\psi\},
\end{equation*}
with $k(z)=K_{12,12}(z)=(z_3^2+z_4^2-z_1^2-z_2^2)/(\pi^2|z|^6)$.

(a) Let $x\in\bar B(0,\delta)$ and $y\in\bar B(re_3,\delta)$. Then $z=x-y=-re_3+w$ with
$w=x-(y-re_3)$ and $|w|\le 2\delta$. Hence $|z_3|\ge r-2\delta>0$, so $z_3^2\ge(r-2\delta)^2$,
while $z_1^2+z_2^2=w_1^2+w_2^2\le 4\delta^2$. The numerator of $k$ therefore satisfies
\begin{equation*}
z_3^2+z_4^2-z_1^2-z_2^2\ge (r-2\delta)^2-4\delta^2=r^2-4r\delta>0,
\end{equation*}
the last inequality because $\delta<r/4$. Moreover $|z|\le r+2\delta$, and so
$k(z)\ge (r^2-4r\delta)/(\pi^2(r+2\delta)^6)$ on the support of $\varphi(x)\psi(y)$. Since
$\varphi(x)\psi(y)\ge 0$ and $\iint\varphi(x)\psi(y)\,dx\,dy=(\int\varphi)(\int\psi)=(\int\varphi)^2$,
integrating this bound gives (a).

(b) Let $x\in\bar B(0,\delta)$ and $y\in\bar B(-re_1,\delta)$. Then $z=x-y=re_1+w$ with $|w|\le 2\delta$,
so $z_1^2\ge(r-2\delta)^2$ and $z_3^2+z_4^2=w_3^2+w_4^2\le 4\delta^2$. Hence
\begin{equation*}
z_3^2+z_4^2-z_1^2-z_2^2\le 4\delta^2-(r-2\delta)^2=-(r^2-4r\delta)<0.
\end{equation*}
Since $|z|\le r+2\delta$, we have $|z|^{-6}\ge (r+2\delta)^{-6}$, and because the numerator is negative,
$k(z)\le -(r^2-4r\delta)/(\pi^2(r+2\delta)^6)$ on the support of $\varphi(x)(R\cdot\psi)(y)$. This
product is non-negative and has integral $(\int\varphi)^2$, so integrating gives (b).

(c) The matrix entries of $R$ are $R_{\alpha\mu}=(Re_\mu)_\alpha$; thus $R_{\alpha1}=\delta_{\alpha3}$
and $R_{\beta2}=\delta_{\beta2}$. Hence, in the notation of
Lemma~\ref{10.7:lem-field-strength-covariance},
\begin{equation*}
F^R_{12}(\chi)=\sum_{\alpha\beta}R_{\alpha1}R_{\beta2}F_{\alpha\beta}(R\cdot\chi)=F_{32}(R\cdot\chi)
\end{equation*}
for $\chi=\varphi$ and for $\chi=\psi$. The covariance of the smeared free field strengths
(Lemma~\ref{10.7:lem-field-strength-covariance}), applied with $\mu\nu=\rho\sigma=12$, gives
\begin{equation*}
\mathbb{E}[F_{32}(R\cdot\varphi)F_{32}(R\cdot\psi)]=\mathbb{E}[F^R_{12}(\varphi)F^R_{12}(\psi)]
=\mathbb{E}[F_{12}(\varphi)F_{12}(\psi)],
\end{equation*}
which is (c).
\end{proof}

For $S\in O(4)$ put $S\cdot\varphi:=\varphi\circ S^{-1}$ and
\begin{equation}\label{10.7:eq-isotropy-fails-generic-1}
F^{S}_{12}(\varphi):=\sum_{\alpha,\beta=1}^{4}S_{\alpha1}S_{\beta2}\,F_{\alpha\beta}(S\cdot\varphi).
\end{equation}
For the quarter turn $R$ of Proposition~\ref{10.7:prop-isotropy-fails-hypercubic}, since $\varphi$ is
radial, this action gives $R\cdot\varphi=\varphi$ and $R\cdot\psi=\varphi(\cdot+re_1)$, as in that
proposition. Moreover
$R_{\alpha1}=\delta_{\alpha3}$ and $R_{\beta2}=\delta_{\beta2}$, so $F^{R}_{12}(\varphi)=F_{32}(R\cdot\varphi)$.
This is the quantity in part~(c) of Proposition~\ref{10.7:prop-isotropy-fails-hypercubic}.

\begin{proposition}[Failure of probe isotropy under a non-hypercubic rotation]
\label{10.7:prop-isotropy-fails-generic}
Let $r$, $\delta$, $\varphi$ and $\psi$ be as in
Proposition~\ref{10.7:prop-isotropy-fails-hypercubic}, and let $c_0:=c_{r,\delta}>0$ be the constant
defined there. For $t\in[0,\pi/2]$ let $R_t$ be the rotation by
$t$ in the $(1,3)$-plane:
\begin{equation*}
R_te_1=\cos t\,e_1+\sin t\,e_3,\qquad R_te_3=-\sin t\,e_1+\cos t\,e_3,\qquad R_te_2=e_2,\quad R_te_4=e_4,
\end{equation*}
so that $R_{\pi/2}$ is the quarter turn $R$ of that proposition. Put
\begin{equation*}
g(t):=E\bigl[F_{12}(R_t\cdot\varphi)\,F_{12}(R_t\cdot\psi)\bigr].
\end{equation*}
Then $g$ is continuous on $[0,\pi/2]$, $g(0)\ge c_0$ and $g(\pi/2)\le-c_0$. Consequently
$g(t)\neq g(0)$ for some $t\in(0,\pi/2)$, and for every such $t$ one has $R_t\notin W_4$. On the other
hand, the covariant quantity
$E\bigl[F^{R_t}_{12}(\varphi)\,F^{R_t}_{12}(\psi)\bigr]$
equals $g(0)$ for every $t$.
\end{proposition}

\begin{proof}
Each $R_t$ is orthogonal with determinant one. The supports of $R_t\cdot\varphi$ and $R_t\cdot\psi$ are
contained in $R_t\bar B(0,\delta)$ and $R_t\bar B(re_3,\delta)$ respectively. These balls lie at distance
at least $r-2\delta>0$, so the supports are disjoint. Hence Lemma~\ref{10.7:lem-field-strength-kernel}
applies. Changing variables
$x=R_tx'$, $y=R_ty'$, with Jacobian one, we obtain
\begin{equation}\label{10.7:eq-isotropy-fails-generic-2}
g(t)=\int\!\!\int\varphi(x')\,\psi(y')\,k\bigl(R_t(x'-y')\bigr)\,dx'\,dy'.
\end{equation}

\emph{Continuity.} On the support of the integrand we have $|x'|\le\delta$ and $|y'-re_3|\le\delta$.
Therefore $|R_t(x'-y')|=|x'-y'|\ge r-2\delta$. From the formula for $k$ in
Lemma~\ref{10.7:lem-field-strength-kernel} we get
\begin{equation*}
|k(z)|=\frac{|z_3^2+z_4^2-z_1^2-z_2^2|}{\pi^2|z|^6}\le\frac{1}{\pi^2|z|^4}
\le\frac{1}{\pi^2(r-2\delta)^4}\qquad(|z|\ge r-2\delta).
\end{equation*}
The function $k$ is continuous off the origin, so $t\mapsto k(R_t(x'-y'))$ is continuous for each
$(x',y')$ in the support. The integrand of \eqref{10.7:eq-isotropy-fails-generic-2} is dominated by
$\varphi(x')\psi(y')/(\pi^2(r-2\delta)^4)$, which is integrable because $\varphi,\psi\ge0$ are continuous with
compact support. By dominated convergence, $g$ is continuous on $[0,\pi/2]$.

\emph{Endpoints.} Since $R_0=\mathbb{1}$, part~(a) of Proposition~\ref{10.7:prop-isotropy-fails-hypercubic}
gives $g(0)=E[F_{12}(\varphi)F_{12}(\psi)]\ge c_0$. Since $R_{\pi/2}=R$, part~(b) gives
$g(\pi/2)\le-c_0$.

\emph{Conclusion.} The number $c_0$ is positive, so $g(\pi/2)<0$. By continuity at $\pi/2$ there is
$t\in(0,\pi/2)$ with $g(t)<0<c_0\le g(0)$. For every $t\in(0,\pi/2)$ the entry $(R_t)_{11}=\cos t$ lies in
$(0,1)$. The entries of a signed permutation matrix all lie in $\{0,\pm1\}$, so $R_t\notin W_4$.

\emph{Covariance.} The kernel of Lemma~\ref{10.7:lem-field-strength-kernel} is built from
Kronecker deltas, the coordinates $z_\nu$ and $|z|$. For $S\in O(4)$ these satisfy
$\delta_{\mu\rho}=\sum_{a,c}S_{\mu a}S_{\rho c}\delta_{ac}$, $(Sz)_\nu=\sum_bS_{\nu b}z_b$ and $|Sz|=|z|$.
The first and third of these follow from $SS^{\mathsf T}=\mathbb{1}$. Substituting these identities term by
term in the explicit formula for $K$ in Lemma~\ref{10.7:lem-field-strength-kernel} gives
\begin{equation*}
K_{\alpha\beta,\gamma\epsilon}(Sz)=\sum_{a,b,c,d}S_{\alpha a}S_{\beta b}S_{\gamma c}S_{\epsilon d}\,
K_{ab,cd}(z).
\end{equation*}
Now take $S=R_t$. Apply bilinearity of the expectation, the change of variables used for
\eqref{10.7:eq-isotropy-fails-generic-2}, and $\sum_\alpha S_{\alpha1}S_{\alpha a}=\delta_{1a}$ (and likewise for
the index $2$). These yield
\begin{equation*}
E\bigl[F^{R_t}_{12}(\varphi)F^{R_t}_{12}(\psi)\bigr]
=\int\!\!\int\varphi(x)\psi(y)K_{12,12}(x-y)\,dx\,dy=g(0).
\qedhere
\end{equation*}
\end{proof}

In words: the observable $F_{12}(\varphi)$ is the continuum analogue of a symbol $(F,\varphi)$ with $F$ a
plaquette variable in the $(1,2)$-plane. Parts~(a) and~(b) of
Proposition~\ref{10.7:prop-isotropy-fails-hypercubic}, together with the proposition above, show that the
identity \eqref{10.7:eq-smeared-block-algebra-a} for products fails in a theory that is manifestly invariant
under $\mathbb{R}^4\rtimes O(4)$. It fails for hypercubic and non-hypercubic $R$ alike. Part~(c), and its
extension to $R_t$ proved above, is the genuine covariance.

\begin{proposition}[Hypercubic rotations as maps of configurations]
\label{10.7:prop-hypercubic-configurations}
Let $w\in W_4$ and let $f$ be a real Lipschitz function on $\mathbb{T}_m$. Then, for every $K$
and every configuration $U$ at cutoff $K$ on $\mathbb{T}_m$,
\begin{equation}\label{10.7:eq-hypercubic-configurations-1}
\mathcal{O}(f)(wU)=\mathcal{O}(w^{-1}\cdot f)(U).
\end{equation}
Consequently, for every $n\ge 2$ and every $n$-tuple $\vec f=(f_1,\dots,f_n)$ of real Lipschitz
functions on $\mathbb{T}_m$ for which $S^T_{n;K,m}(\vec f)$ is defined, $S^T_{n;K,m}(w\cdot\vec f)$
is defined and
\begin{equation}\label{10.7:eq-hypercubic-configurations-2}
S^T_{n;K,m}(w\cdot\vec f)=S^T_{n;K,m}(\vec f),
\end{equation}
at every $(K,m)$, where $w\cdot\vec f=(w\cdot f_1,\dots,w\cdot f_n)$.
\end{proposition}

\begin{proof}
Here $R\cdot f=f\circ R^{-1}$ is the action of $O(4)$ on test functions fixed in the notation on
test-function moduli, separated classes and the action of rotations
(see~\ref{10.7:not-test-function-moduli}), so that $w^{-1}\cdot f=f\circ w$. A signed permutation
matrix $w$ maps $L^{-K}\mathbb{Z}^4$ and its translate $L^{-K}(\mathbb{Z}+\tfrac12)^4$ onto
themselves, and it maps the period lattice $2L^m\mathbb{Z}^4$ onto itself; hence it acts on the
site set of the lattice of spacing $L^{-K}$ in $\mathbb{T}_m$. It sends the oriented bond from $x$ to
$x+L^{-K}e_\mu$ to the segment from $wx$ to $wx+L^{-K}we_\mu$. Since $we_\mu=\pm e_\nu$ for
some $\nu$, this segment is a bond, possibly with reversed orientation. Likewise $w$ maps each
plaquette $p$ to a plaquette $wp$, and $x(wp)=w\,x(p)$ because $w$ is linear. We use the action of
$w$ on configurations given by $(wU)(b)=U(w^{-1}b)$ for oriented bonds $b$, with the convention
$U(b^{-1})=U(b)^{-1}$; it is under this action that the invariance of the Wilson measure is
invoked below.

Put $a_p(U)=1-\operatorname{Re}\operatorname{tr}U(\partial p)$. The holonomy $(wU)(\partial p)$
is the product of $U$ along $\partial(w^{-1}p)$. It may start at another corner, which
conjugates it, and it may run in the opposite orientation, which replaces it by its inverse.
The trace is invariant under conjugation, and for $V\in SU(N)$ one has
$\operatorname{tr}V^{-1}=\overline{\operatorname{tr}V}$, so that
$\operatorname{Re}\operatorname{tr}V^{-1}=\operatorname{Re}\operatorname{tr}V$. Hence
$a_p(wU)=a_{w^{-1}p}(U)$. Since $p'\mapsto wp'$ is a bijection of the set of plaquettes, the
substitution $p=wp'$ gives
\begin{align*}
\mathcal{O}(f)(wU)&=\sum_p f(x(p))\,a_{w^{-1}p}(U)=\sum_{p'}f(w\,x(p'))\,a_{p'}(U)\\
&=\sum_{p'}(w^{-1}\cdot f)(x(p'))\,a_{p'}(U)=\mathcal{O}(w^{-1}\cdot f)(U),
\end{align*}
which is \eqref{10.7:eq-hypercubic-configurations-1}.

Let $\vec f$ be a tuple for which $S^T_{n;K,m}(\vec f)$ is defined. Since the admissibility
conditions on tuples of test functions are $O(4)$-stable
(see~\ref{10.7:not-test-function-moduli}), $S^T_{n;K,m}(w\cdot\vec f)$ is defined as well.
Write $\langle\cdot\rangle_{K,m}$ for the expectation under the Wilson measure at cutoff $K$ on
$\mathbb{T}_m$. For real $z=(z_1,\dots,z_n)$ put
\begin{equation*}
Z_{\vec f}(z)=\Bigl\langle\exp\Bigl(\sum_{i=1}^n z_i\,\mathcal{O}(f_i)\Bigr)\Bigr\rangle_{K,m}.
\end{equation*}
Apply \eqref{10.7:eq-hypercubic-configurations-1} with $w^{-1}\in W_4$ in place of $w$. This
gives $\mathcal{O}(w\cdot f_i)(U)=\mathcal{O}(f_i)(w^{-1}U)$, so the integrand of
$Z_{w\cdot\vec f}(z)$ is the function $\exp(\sum_i z_i\mathcal{O}(f_i))$ evaluated at
$w^{-1}U$. The Wilson measure on $\mathbb{T}_m$ is invariant under $U\mapsto w^{-1}U$, by the
lemma on invariances of the Wilson measure; applying this to the function
$\exp(\sum_i z_i\mathcal{O}(f_i))$ gives $Z_{w\cdot\vec f}(z)=Z_{\vec f}(z)$ for all real $z$.

Both functions equal $1$ at $z=0$ and are smooth, since the lattice is finite and the
observables are bounded. Their logarithms are therefore defined and coincide for real $z$ near
$0$. The connected Schwinger function $S^T_{n;K,m}(\vec f)$ is the joint cumulant
$\kappa(\mathcal{O}(f_1),\dots,\mathcal{O}(f_n))$, that is, the derivative
$\partial_{z_1}\cdots\partial_{z_n}\log Z_{\vec f}$ at $z=0$. This is a derivative along real
directions, and therefore \eqref{10.7:eq-hypercubic-configurations-2} follows.
\end{proof}

For general $R\in O(4)$ the action on observables at cutoff $K$ is fixed, in the notation on
test-function moduli, separated classes and the action of rotations
(see~\ref{10.7:not-test-function-moduli}), by
\begin{equation*}
\mathcal{O}(f)\circ R=\mathcal{O}(R^{-1}\cdot f)=\sum_p f(R\,x(p))\,a_p,
\end{equation*}
where $a_p=1-\operatorname{Re}\operatorname{tr}U(\partial p)$: the lattice is not moved, only the
weights are. There $O(4)$ acts on $C^1_c(\mathbb{R}^4)$ by $R\cdot f=f\circ R^{-1}$, with
$\operatorname{supp}(R\cdot f)=R(\operatorname{supp}f)$,
$\nabla(R\cdot f)(x)=R\nabla f(R^{-1}x)$ and $\|R\cdot f\|_{\sharp}=\|f\|_{\sharp}$; rotations
preserve separations of supports and balls, mapping the centre and keeping the radius. Hence the
class of observables $\mathcal{O}(f)$ and every admissibility condition on tuples of test
functions are $O(4)$-stable, and $S^T_n(R\cdot\vec f)$ is defined whenever $S^T_n(\vec f)$ is.

On $W_4$ the action on test functions is thus implemented by a map of configurations, and it is
an exact symmetry at every $(K,m)$. For $R\in O(4)$ with $R\notin W_4$ the action is implemented
by no map of configurations, and its invariance is the content of the theorem on $O(4)$
invariance of the continuum limit.

\begin{lemma}[Classical expansion of the plaquette density]
\label{10.7:lem-plaquette-expansion}
Let $N\ge 2$, let $A=(A_1,\dots,A_4)$ be a smooth $\mathfrak{su}(N)$-valued $1$-form on $\mathbb{R}^4$,
let $e_\mu$ denote the unit vectors of $\mathbb{R}^4$, and put
$F_{\mu\nu}=\partial_\mu A_\nu-\partial_\nu A_\mu+[A_\mu,A_\nu]$. For $\varepsilon>0$ and a link $(x,\mu)$
of $\varepsilon\mathbb{Z}^4$ let $U^{\varepsilon}(x,\mu)=u(\varepsilon)$ be the ordered exponential along the
link, later points on the right: $u'(t)=u(t)A_\mu(x+te_\mu)$ on $[0,\varepsilon]$, $u(0)=\mathbb{1}$. Write
$U^{\varepsilon}[A]$ for the resulting configuration and, for $\mu\ne\nu$,
\begin{equation*}
U^{\varepsilon}_{\mu\nu}(x)=U^{\varepsilon}(x,\mu)\,U^{\varepsilon}(x+\varepsilon e_\mu,\nu)\,
U^{\varepsilon}(x+\varepsilon e_\nu,\mu)^{*}\,U^{\varepsilon}(x,\nu)^{*}
\end{equation*}
for the plaquette variable $U(\partial p)$ of the plaquette $p$ with lowest corner $x$ in the
$\mu\nu$-plane. For $c\in C_c(\mathbb{R}^4)$ let
$A(c,U)=\sum_p c(x(p))\,[1-\operatorname{Re}\operatorname{tr}U(\partial p)]$, the sum running over the
plaquettes of $\varepsilon\mathbb{Z}^4$. Let $\operatorname{Tr}=N\operatorname{tr}$ be the unnormalised trace.
Thus $A(c,\cdot)$ is the smeared observable $\mathcal{O}(c)$; the letter $A$ with the two arguments
$(c,U)$ denotes this weighted sum, and the letter $A$ alone always denotes the $1$-form.
Each $O(\varepsilon^k)$ below is bounded in matrix norm by $C\varepsilon^k$, with $C$ uniform for $x$ in
compact sets. Then:
\begin{enumerate}
\item[(i)] $U^{\varepsilon}(x,\mu)=\exp\bigl(\varepsilon A_\mu(x+\tfrac12\varepsilon e_\mu)\bigr)+O(\varepsilon^3)$;
\item[(ii)] $U^{\varepsilon}_{\mu\nu}(x)=\mathbb{1}+\varepsilon^2F_{\mu\nu}(x)+O(\varepsilon^3)$;
\item[(iii)] for every unitary $V$ one has
$1-\operatorname{Re}\operatorname{tr}V=\tfrac12\operatorname{tr}[(V-\mathbb{1})^{*}(V-\mathbb{1})]$, and hence
\begin{equation*}
1-\operatorname{Re}\operatorname{tr}U^{\varepsilon}_{\mu\nu}(x)
=-\tfrac12\varepsilon^4\operatorname{tr}(F_{\mu\nu}(x)^2)+O(\varepsilon^5),
\end{equation*}
the orders $\varepsilon^0,\dots,\varepsilon^3$ vanishing identically in $A$;
\item[(iv)] the expansion
\begin{equation*}
\sum_{\mu<\nu}[1-\operatorname{Re}\operatorname{tr}U^{\varepsilon}_{\mu\nu}(x)]
=-\tfrac14\varepsilon^4\sum_{\mu,\nu}\operatorname{tr}(F_{\mu\nu}F_{\mu\nu})(x)+O(\varepsilon^5)
\end{equation*}
holds, and $P_4[A]:=-\frac14\sum_{\mu,\nu}\operatorname{tr}(F_{\mu\nu}F_{\mu\nu})$ is an
$O(4)$-scalar density of scaling dimension $4$: for $R\in O(4)$ and the rotated $1$-form
$(A^R)_\mu(x)=\sum_\nu R_{\mu\nu}A_\nu(R^{-1}x)$,
\begin{equation*}
F^{A^R}_{\mu\nu}(x)=\sum_{\gamma,\beta}R_{\mu\gamma}R_{\nu\beta}F_{\gamma\beta}(R^{-1}x),
\qquad P_4[A^R](x)=P_4[A](R^{-1}x),
\end{equation*}
and $P_4[A^\lambda](x)=\lambda^4P_4[A](\lambda x)$ for $A^\lambda_\mu(x)=\lambda A_\mu(\lambda x)$,
$\lambda>0$. Moreover, for $c\in C_c(\mathbb{R}^4)$,
\begin{equation*}
\lim_{\varepsilon\to0}A(c,U^{\varepsilon}[A])
=-\tfrac14\int c\sum_{\mu,\nu}\operatorname{tr}(F_{\mu\nu}F_{\mu\nu}),
\end{equation*}
so that the action $g_0^{-2}A(c,U^{\varepsilon}[A])$ regularises
$-(4Ng_0^2)^{-1}\int c\sum_{\mu,\nu}\operatorname{Tr}F_{\mu\nu}^2$; that is,
$g_{\mathrm{YM}}^2=2Ng_0^2$ against
$S=-(2g_{\mathrm{YM}}^2)^{-1}\int\sum_{\mu,\nu}\operatorname{Tr}F_{\mu\nu}^2$.
\end{enumerate}
For $N=2$, $\operatorname{tr}V$ is real on $SU(2)$, so $\operatorname{Re}\operatorname{tr}=\operatorname{tr}$; for
$N\ge3$ the real part is kept and nothing else changes.
\end{lemma}

\begin{proof}
Since $A_\mu$ is anti-Hermitian, $(uu^{*})'=u\,(A_\mu+A_\mu^{*})\,u^{*}=0$, so $uu^{*}=\mathbb{1}$ and
$u(t)$ is unitary, $\|u(t)\|=1$; all link and plaquette variables are unitary.

(i) Put $a(t)=A_\mu(x+te_\mu)$. Inserting $u(s)=\mathbb{1}+\int_0^s u(r)a(r)\,dr$ twice into
$u(\varepsilon)=\mathbb{1}+\int_0^\varepsilon u(s)a(s)\,ds$ (Picard iteration to second order) gives
\begin{equation*}
u(\varepsilon)=\mathbb{1}+\int_0^\varepsilon a(s)\,ds+\int_0^\varepsilon\!\!\int_0^s a(r)a(s)\,dr\,ds
+\int_0^\varepsilon\!\!\int_0^s\!\!\int_0^r u(q)a(q)a(r)a(s)\,dq\,dr\,ds ,
\end{equation*}
and the last term has norm at most $\varepsilon^3\sup\|a\|^3/6$. Taylor expansion of $a$ about
$\varepsilon/2$ to second order gives $\int_0^\varepsilon a=\varepsilon a(\varepsilon/2)+O(\varepsilon^3)$,
the first-order term integrating to zero. The double integral differs from
$\frac12(\int_0^\varepsilon a)^2$ by $\frac12\int\!\!\int_{r<s}[a(r),a(s)]\,dr\,ds$, and
$\|[a(r),a(s)]\|=O(\varepsilon)$ on $[0,\varepsilon]^2$; so this ordering ambiguity is $O(\varepsilon^3)$
and the double integral equals $\frac12\varepsilon^2a(\varepsilon/2)^2+O(\varepsilon^3)$. Since
\begin{equation*}
\exp(\varepsilon a(\varepsilon/2))
=\mathbb{1}+\varepsilon a(\varepsilon/2)+\tfrac12\varepsilon^2a(\varepsilon/2)^2+O(\varepsilon^3),
\end{equation*}
(i) follows.

(ii) By (i) and unitarity, $U^{*}=e^{-X}+O(\varepsilon^3)$ whenever $U=e^{X}+O(\varepsilon^3)$ with $X$
anti-Hermitian; and since all factors have norm one, replacing each by its exponential costs
$O(\varepsilon^3)$. Hence $U^{\varepsilon}_{\mu\nu}(x)=e^{X_1}e^{X_2}e^{X_3}e^{X_4}+O(\varepsilon^3)$ with
\begin{align*}
X_1&=\varepsilon A_\mu(x+\tfrac12\varepsilon e_\mu), &
X_2&=\varepsilon A_\nu(x+\varepsilon e_\mu+\tfrac12\varepsilon e_\nu),\\
X_3&=-\varepsilon A_\mu(x+\varepsilon e_\nu+\tfrac12\varepsilon e_\mu), &
X_4&=-\varepsilon A_\nu(x+\tfrac12\varepsilon e_\nu).
\end{align*}
Each $X_i=O(\varepsilon)$, so $e^{X_i}=\mathbb{1}+X_i+\frac12X_i^2+O(\varepsilon^3)$ and
\begin{equation*}
e^{X_1}e^{X_2}e^{X_3}e^{X_4}=\mathbb{1}+\sum_iX_i+\tfrac12\sum_iX_i^2+\sum_{i<j}X_iX_j+O(\varepsilon^3).
\end{equation*}
First order: by Taylor expansion to first order with second-order remainder,
$X_1+X_3=-\varepsilon^2\partial_\nu A_\mu(x)+O(\varepsilon^3)$ and
$X_2+X_4=\varepsilon^2\partial_\mu A_\nu(x)+O(\varepsilon^3)$, so
$\sum_iX_i=\varepsilon^2(\partial_\mu A_\nu-\partial_\nu A_\mu)(x)+O(\varepsilon^3)$. Second order: up to
$O(\varepsilon^3)$ one may put $X_1=-X_3=\varepsilon A_\mu(x)$ and $X_2=-X_4=\varepsilon A_\nu(x)$; then
$\frac12\sum_iX_i^2=\varepsilon^2(A_\mu^2+A_\nu^2)$, while the six products $X_iX_j$, $i<j$, are
$\varepsilon^2$ times $A_\mu A_\nu$, $-A_\mu^2$, $-A_\mu A_\nu$, $-A_\nu A_\mu$, $-A_\nu^2$, $A_\mu A_\nu$.
The sum of all second-order terms is $\varepsilon^2[A_\mu,A_\nu](x)$, which proves (ii).

(iii) For unitary $V$, expanding and using $V^{*}V=\mathbb{1}$ gives
\begin{equation*}
(V-\mathbb{1})^{*}(V-\mathbb{1})=V^{*}V-V-V^{*}+\mathbb{1}=2\cdot\mathbb{1}-V-V^{*};
\end{equation*}
taking $\tfrac12\operatorname{tr}$ and using $\operatorname{tr}\mathbb{1}=1$ and
$\operatorname{tr}V^{*}=\overline{\operatorname{tr}V}$ gives the identity. For
$V=U^{\varepsilon}_{\mu\nu}(x)$, (ii) gives $V-\mathbb{1}=\varepsilon^2F_{\mu\nu}(x)+O(\varepsilon^3)$, so
$(V-\mathbb{1})^{*}(V-\mathbb{1})=\varepsilon^4F_{\mu\nu}^{*}F_{\mu\nu}+O(\varepsilon^5)$. Because
$F_{\mu\nu}$ takes values in $\mathfrak{su}(N)$, $F_{\mu\nu}^{*}=-F_{\mu\nu}$, and the claimed expansion
follows; as $V-\mathbb{1}=O(\varepsilon^2)$ for every $A$, the orders $\varepsilon^0$ to $\varepsilon^3$
vanish identically.

(iv) Summing (iii) over $\mu<\nu$ and using $F_{\nu\mu}=-F_{\mu\nu}$, $F_{\mu\mu}=0$, so that
$\sum_{\mu<\nu}\operatorname{tr}F_{\mu\nu}^2=\frac12\sum_{\mu,\nu}\operatorname{tr}(F_{\mu\nu}F_{\mu\nu})$,
gives the expansion. Since $R^{-1}=R^{\mathsf T}$,
\begin{align*}
\partial_\mu(A^R)_\nu(x)&=\sum_{\gamma,\beta}R_{\mu\gamma}R_{\nu\beta}(\partial_\gamma A_\beta)(R^{-1}x),\\
\bigl[(A^R)_\mu,(A^R)_\nu\bigr](x)&=\sum_{\gamma,\beta}R_{\mu\gamma}R_{\nu\beta}\bigl[A_\gamma,A_\beta\bigr](R^{-1}x),
\end{align*}
which is the transformation law of $F$ stated in (iv). Inserting it into $P_4[A^R]$ and using
$\sum_\mu R_{\mu\gamma}R_{\mu\gamma'}=\delta_{\gamma\gamma'}$ twice gives $P_4[A^R](x)=P_4[A](R^{-1}x)$.
Likewise $F^{A^\lambda}_{\mu\nu}(x)=\lambda^2F_{\mu\nu}(\lambda x)$, whence the factor $\lambda^4$.

Plaquettes of $\varepsilon\mathbb{Z}^4$ correspond to pairs $(x,\mu<\nu)$ with
$x(p)=x+\frac12\varepsilon(e_\mu+e_\nu)$. For $\varepsilon\le1$ only $x$ in a fixed compact neighbourhood
of $\operatorname{supp}c$ contribute, at most $C\varepsilon^{-4}$ of them, and there each bracket
$1-\operatorname{Re}\operatorname{tr}U^{\varepsilon}_{\mu\nu}(x)$ is $O(\varepsilon^4)$ by (iii). By
uniform continuity of $c$, replacing $c(x(p))$ by $c(x)$ changes $A(c,U^{\varepsilon}[A])$ by a quantity
tending to zero. By the expansion above,
\begin{equation*}
\sum_{x}c(x)\sum_{\mu<\nu}[1-\operatorname{Re}\operatorname{tr}U^{\varepsilon}_{\mu\nu}(x)]
=\sum_{x}\varepsilon^4c(x)P_4[A](x)+O(\varepsilon),
\end{equation*}
and the first term on the right is a Riemann sum of the continuous compactly supported function
$cP_4[A]$, which converges to $\int cP_4[A]$. Finally
\begin{equation*}
g_0^{-2}\cdot\bigl(-\tfrac14\bigr)\sum_{\mu,\nu}\operatorname{tr}F_{\mu\nu}^2
=-(4Ng_0^2)^{-1}\sum_{\mu,\nu}\operatorname{Tr}F_{\mu\nu}^2,
\end{equation*}
and equating $(4Ng_0^2)^{-1}$ with $(2g_{\mathrm{YM}}^2)^{-1}$ gives
$g_{\mathrm{YM}}^2=2Ng_0^2$.

For $V\in SU(2)$ the eigenvalues are $e^{i\theta}$ and $e^{-i\theta}$, so $\operatorname{tr}V=\cos\theta$
is real; for $N\ge3$ the argument above used only the real part and is unchanged.
\end{proof}

\begin{remark}[The objects on which rotations act]\label{10.7:rem-rotation-objects}
On the smeared block algebra $\mathfrak{A}^{sm}_{loc}$ the rotation sentence $(S_\beta)$ of
\eqref{10.7:eq-smeared-block-algebra-a} says the following.
\begin{enumerate}
\item For single symbols it is the empty identity of
Proposition~\ref{10.7:prop-single-symbol-identity}.
\item For products it is isotropy of fixed hypercubic probes
(Proposition~\ref{10.7:prop-probe-isotropy}); it is not $O(4)$-invariance.
\item It fails in the free Maxwell field
(Propositions~\ref{10.7:prop-isotropy-fails-hypercubic}
and~\ref{10.7:prop-isotropy-fails-generic}).
\item It cannot be repaired inside the algebras of block observables
(Lemma~\ref{10.7:lem-lattice-stabiliser} and Corollary~\ref{10.7:cor-no-block-transport}).
\end{enumerate}
The objects on which a rotation acts are the test functions $f$ of the smeared curvature
observables $\mathcal{O}(f)$ of clause (iii)(b) of the main theorem. The density
$1-\operatorname{Re}\operatorname{tr}U(\partial p)$ of these observables is the lattice
regularisation of the $O(4)$-scalar action density $\operatorname{tr}F_{\mu\nu}F_{\mu\nu}$
(compare Lemma~\ref{10.7:lem-plaquette-expansion}). Rotation invariance is to be stated for
the connected functions $S^T_{n;\infty}$ on $\mathbb{R}^4$, $n\ge 2$, of item (b2) of that
clause.

A tensor class may also be formulated. Let $W$ be a finite-dimensional inner-product space
carrying a representation $D_W$ of $O(4)$. Let $O^{\varepsilon}$ be a $W_4$-organised family
of lattice operators of dimension $\Delta$ at spacing $\varepsilon$, placed in $W$ by a map
$\iota_W$. For a $W$-valued test function $h$ set
\begin{equation*}
\mathcal{O}_W(h):=\sum_x \varepsilon^{4-\Delta}\,
\bigl\langle h(x),\iota_W\bigl(O^{\varepsilon}(x)\bigr)\bigr\rangle_W ,
\qquad
(R\cdot h):=D_W(R)\,h\circ R^{-1}.
\end{equation*}
This class is $O(4)$-stable as a class: since $D_W$ is a representation,
$h\mapsto D_W(R)\,h\circ R^{-1}$ is an action of $O(4)$ on $W$-valued test functions. It is
not among the observables of the main theorem, and on perturbative grounds, not examined in
this book, it would need $W_4$-piecewise relative normalisations. It is
stated here as a formulation only and is not used.
\end{remark}

\begin{proof}
For a single symbol, Proposition~\ref{10.7:prop-single-symbol-identity} shows that the
identity holds for every $T\in GL_4(\mathbb{R})$ with $|\det T|=1$. It therefore holds for
volume-preserving shears as much as for rotations. Hence it is empty as a statement about
rotational symmetry.

For a product of symbols and $w\in W_4$, Proposition~\ref{10.7:prop-probe-isotropy} rewrites
$(S_\beta)(w)$ as follows: truncated correlations are unchanged when every probe is rotated by
$w^{-1}$ about its own position, with the positions fixed. This is isotropy of the probes. It
is not invariance of the state under the motion.

In the free Maxwell field this isotropy fails for the quarter turn in the $(1,3)$-plane,
which lies in $W_4$. Indeed, parts (a) and (b) of
Proposition~\ref{10.7:prop-isotropy-fails-hypercubic} give a correlation that is at least
$c_0>0$ before the turn and at most $-c_0$ after it. Proposition~\ref{10.7:prop-isotropy-fails-generic}
gives a rotation $R_t\notin W_4$ for which the rotated-probe correlation differs from the
unrotated one.

Lemma~\ref{10.7:lem-lattice-stabiliser} shows that the isometries mapping a half-offset lattice
$\mathbb{L}_s$ onto a half-offset lattice have linear part in $W_4$. Corollary~\ref{10.7:cor-no-block-transport}
draws the consequence: the transport of block observables has no extension to linear parts
$R\notin W_4$. The only map on symbols of the pattern ``transport the observable, move the
test function'' that remains is $\beta_R$, to which the preceding failures apply. This proves
the four items of the remark.
\end{proof}

\begin{lemma}[Rotation stability of the hypothesis classes]\label{10.7:lem-class-stability}
Let $R\in O(4)$, let $k\in\{1,2\}$, and let $f\in C^k_c(\mathbb{R}^4)$ be real. Let $R\cdot f$ denote
the rotated function in the notation of the paragraph on test-function moduli, separated classes and
the action of rotations (\ref{10.7:not-test-function-moduli}), $(R\cdot f)(x)=f(R^{-1}x)$. Then
$R\cdot f\in C^k_c(\mathbb{R}^4)$ is real, and
\begin{equation}\label{10.7:eq-class-stability-1}
\operatorname{supp}(R\cdot f)=R(\operatorname{supp} f),\qquad \|R\cdot f\|_{\sharp}=\|f\|_{\sharp}.
\end{equation}
Passing from a tuple $\vec f=(f_1,\dots,f_n)$, $n\ge 2$, of real functions in $C^k_c(\mathbb{R}^4)$
to $R\cdot\vec f=(R\cdot f_1,\dots,R\cdot f_n)$
preserves the pairwise distances of the supports, the radii of enclosing balls, and the ball
$B(0,\mathcal{R}/2)$. Hence:
\begin{enumerate}
\item[(a)] the hypothesis classes of the per-ball volume-free bound, of the construction of the
subsequential limit families on $\mathbb{R}^4$, of the hypothesis of the rotational Ward identity
and of the hypothesis of the dimension-six coefficient bound (in either of its two forms), and
every constant these classes read, are invariant under $\vec f\mapsto R\cdot\vec f$;
\item[(b)] the per-ball volume-free bound holds for $R\cdot\vec f$ with the same constant
$C^\infty_n$ of that bound, at the same separation $\mathsf{d}$ and the same ball radius;
\item[(c)] for a real antisymmetric $4\times4$ matrix $\omega$ and $R^{\omega}_t=e^{t\omega}$,
every bound of the hypothesis of the dimension-six coefficient bound in its first form which reads
the $f_i$ only through rotation-invariant norms holds along $R^{\omega}_t\cdot\vec f$ uniformly
in $t$.
\end{enumerate}
\end{lemma}

\begin{proof}
The map $x\mapsto R^{-1}x$ is linear, hence smooth, and $R\cdot f$ is its composition with $f$, so
it is real and of class $C^k$. By the chain rule, for $x\in\mathbb{R}^4$ and indices $\mu,\lambda$,
\begin{gather*}
\partial_\mu(R\cdot f)(x)=\sum_{\nu}(R^{-1})_{\nu\mu}(\partial_\nu f)(R^{-1}x),\\
\partial_\mu\partial_\lambda(R\cdot f)(x)
=\sum_{\nu,\sigma}(R^{-1})_{\nu\mu}(R^{-1})_{\sigma\lambda}(\partial_\nu\partial_\sigma f)(R^{-1}x),
\end{gather*}
the second identity being used when $k=2$; both right sides are continuous in $x$.

Since $R$ is a bijection, $\{x:(R\cdot f)(x)\neq0\}=R\{y:f(y)\neq0\}$. Since $R$ is a
homeomorphism, it commutes with closure, which gives the first identity
of \eqref{10.7:eq-class-stability-1}. The image of the compact set $\operatorname{supp}f$ under
the continuous map $R$ is compact, so $R\cdot f\in C^k_c(\mathbb{R}^4)$.

Since $R^{-1}$ is a bijection of $\mathbb{R}^4$, we have
$\|R\cdot f\|_\infty=\sup_y|f(y)|=\|f\|_\infty$. Put $u=R^{-1}x$ and $v=R^{-1}y$. Since $R$ is an
isometry, $|x-y|=|u-v|$, and the map $(x,y)\mapsto(u,v)$ is a bijection of the pairs with
$x\neq y$ onto the pairs with $u\neq v$. Therefore
\begin{equation*}
\operatorname{Lip}(R\cdot f)=\sup_{x\neq y}\frac{|f(R^{-1}x)-f(R^{-1}y)|}{|x-y|}
=\sup_{u\neq v}\frac{|f(u)-f(v)|}{|u-v|}=\operatorname{Lip}f .
\end{equation*}
Since $\|\cdot\|_\sharp$ is the maximum of $\|\cdot\|_\infty$ and $40M\operatorname{Lip}(\cdot)$,
this gives the second identity of \eqref{10.7:eq-class-stability-1}.

For sets $A,B\subset\mathbb{R}^4$, write
$\operatorname{dist}(A,B)=\inf\{|a-b|:a\in A,\ b\in B\}$; the identity $|Ra-Rb|=|a-b|$ gives
$\operatorname{dist}(RA,RB)=\operatorname{dist}(A,B)$. By the first identity
of \eqref{10.7:eq-class-stability-1}, the pairwise distances of the supports of $R\cdot\vec f$
are therefore those of $\vec f$. Next, $A\subset B(c,r)$ holds if and only if
$RA\subset B(Rc,r)$, so enclosing balls correspond with the same radii. Finally, $R0=0$ and
$R$ is a bijective isometry, so $R\,B(0,\mathcal{R}/2)=B(0,\mathcal{R}/2)$. Hence
$\operatorname{supp}f\subset B(0,\mathcal{R}/2)$ if and only if
$\operatorname{supp}(R\cdot f)\subset B(0,\mathcal{R}/2)$.

The classes in (a) and the constants they read are expressed through the following data:
regularity $C^k_c$, reality, the separations of the supports, the radii of enclosing balls,
containment in $B(0,\mathcal{R}/2)$, and the norms $\|\cdot\|_\sharp$. Each of these data is
unchanged by $\vec f\mapsto R\cdot\vec f$, which proves (a).

For (b): by (a), $R\cdot\vec f$ satisfies the hypotheses of the per-ball volume-free bound with
the same separation $\mathsf{d}$, the same ball radius and the same norms. The constant
$C^\infty_n$ of that bound is read from these data, so it is the same.

For (c): since $\omega^{T}=-\omega$, we have
$(e^{t\omega})^{T}e^{t\omega}=e^{-t\omega}e^{t\omega}=\mathbb{1}$, so each $R^{\omega}_t$ lies in
$O(4)$. By (a), $R^{\omega}_t\cdot\vec f$ lies in the hypothesis class for every $t$, so the bound
applies at each $t$. Since it reads the $f_i$ only through norms invariant under $O(4)$, its right
side takes the same value at $R^{\omega}_t\cdot\vec f$ for every $t$. It is thus uniform in $t$.
\end{proof}

\begin{definition}[Reflection-positivity data of a subsequential limit]\label{10.7:def-reflection-data}
Let $(K_j,m_j)_{j\ge 1}$ be the diagonal sequence along which a subsequential limit is taken, and
write $S_j:=S^T_{n;K_j,m_j}$ for the connected Schwinger functions along it. The
\emph{reflection-positivity data} of the limit consist of the limit functionals $S^T_n$ defined in
(S1) below, together with
the properties (S1)--(S4) and the hypothesis (H-O4) displayed below. Properties (S1)--(S4) are
supplied by the infinite-volume theorem and the per-ball volume-free bound. Hypothesis (H-O4) is
the conclusion of the $O(4)$-invariance theorem
for this limit. Admissible tuples and admissible combinations are understood in the sense of
Notation~\ref{10.7:not-test-function-moduli}.
\begin{equation}\label{10.7:eq-reflection-data-a}
\begin{aligned}
&(\mathrm{S1})\quad S^T_n(\vec f)=\lim_{j\to\infty}S_j(\vec f)
  \ \text{exists for every admissible tuple } \vec f;\\
&(\mathrm{S2})\quad S^T_n \ \text{is symmetric, real and multilinear on admissible
  combinations};\\
&(\mathrm{S3})\quad S^T_n \ \text{is exactly invariant under } \mathbb{R}^4\rtimes W_4;\\
&(\mathrm{S4})\quad |S^T_n(\vec f)|\le C^{\infty}_n(\mathsf{d};r)\prod_{i=1}^{n}\|f_i\|_{\sharp}
  \ \text{whenever each } f_i \text{ is supported in some closed ball}\\
&\qquad\qquad \text{of radius } r \text{ and the supports are pairwise at distance at least }
  \mathsf{d};\\
&(\mathrm{H\text{-}O4})\quad \text{the conclusion of the $O(4)$-invariance theorem holds for }
  S^T_n .
\end{aligned}
\end{equation}
The constant
$C^{\infty}_n(\mathsf{d};r)$ does not depend on the positions of the balls. Only this bound for
the limit functional $S^T_n$ is used in what follows; the corresponding bounds for the $S_j$ are
not. Until the step of the argument below at which the full hypothesis (H-O4) is invoked, it is
used only for the motions $\hat R=1\oplus R$ with
$R\in SO(3)$.

The data come with the following coordinates and motions:
\begin{itemize}
\item points of $\mathbb{R}^4$ are written $x=(x^0,\vec x)$ with $\vec x\in\mathbb{R}^3$;
\item $\theta$ is the reflection of the coordinate $x^0$, that is,
  $\theta(x^0,\vec x)=(-x^0,\vec x)$, so that $\theta\in W_4$;
\item $\mathbb{R}^4_+:=\{x\in\mathbb{R}^4 : x^0>0\}$;
\item for $R\in O(3)$, $\hat R:=1\oplus R$;
\item for $t\in\mathbb{R}$, $\tau_t$ is the translation by $te_0$, where $e_0=(1,0,0,0)$ is the
  unit vector of the $x^0$-axis;
\item $\mathfrak{E}_+$ is the set of Euclidean motions of $\mathbb{R}^4$ that map
  $\mathbb{R}^4_+$ into itself.
\end{itemize}
\end{definition}

\begin{lemma}[Euclidean motions preserving the half-space]\label{10.7:lem-half-space-motions}
Write $x=(x^0,\vec x)\in\mathbb{R}^4$, let $\theta(x^0,\vec x)=(-x^0,\vec x)$ be the time
reflection, and let $\mathfrak{E}_+$ be the set of Euclidean motions $g$ of $\mathbb{R}^4$
mapping the half-space $\{x\in\mathbb{R}^4 : x^0>0\}$ into itself. Then
\begin{equation*}
\mathfrak{E}_+=\bigl\{\,x\mapsto(1\oplus R)x+b \;:\; R\in O(3),\ b\in\mathbb{R}^4,\ b^0\ge 0\,\bigr\}
\cong[0,\infty[\,\times\bigl(\mathbb{R}^3\rtimes O(3)\bigr),
\end{equation*}
and $g\mapsto g^*:=\theta g^{-1}\theta$ is an involution of $\mathfrak{E}_+$.
\end{lemma}

\begin{proof}
Let $g(x)=Ax+b$ with $A\in O(4)$. Then $g\in\mathfrak{E}_+$ exactly when
\begin{equation*}
A_{00}\,s+\sum_{i=1}^{3}A_{0i}x^i+b^0>0\qquad\text{for all } s>0,\ \vec x\in\mathbb{R}^3 .
\end{equation*}
If some $A_{0i}\neq 0$, fix $s$ and take $\vec x=-\lambda(A_{01},A_{02},A_{03})$ with
$\lambda>0$ large; the left-hand side becomes negative. Hence $A_{0i}=0$ for $i=1,2,3$. The first
row of $A$ is a unit vector, so $A_{00}=\pm1$; letting $s\to\infty$ in $A_{00}s+b^0>0$ gives
$A_{00}\ge 0$, so $A_{00}=1$. Letting $s\to 0$ gives $b^0\ge 0$. Since $A^{\mathsf T}A=\mathbb{1}$,
the first column of $A$ is also a unit vector; its entry $A_{00}$ equals $1$, so the first column is
$e_0$. Therefore $A=1\oplus R$ with $R\in O(3)$. Conversely, for such $A$ and $b^0\ge0$ the time
coordinate of $g(x)$ is $x^0+b^0>0$ whenever $x^0>0$. This proves the equality of sets. The map
$g\mapsto(b^0,\vec b,R)$ is a bijection onto $[0,\infty[\,\times\mathbb{R}^3\times O(3)$, and
composing $x\mapsto(1\oplus R)x+(t,\vec a)$ after $x\mapsto(1\oplus R')x+(t',\vec a')$ gives
$x\mapsto(1\oplus RR')x+(t+t',\vec a+R\vec a')$, which is the product of
$[0,\infty[$ under addition with the semidirect product $\mathbb{R}^3\rtimes O(3)$.

For $g(x)=(1\oplus R)x+(t,\vec a)$ with $t\ge0$ we have
$g^{-1}(y)=(1\oplus R)^{-1}(y-(t,\vec a))$, and since $\theta x-(t,\vec a)=(-x^0-t,\vec x-\vec a)$,
\begin{equation*}
g^*(x)=\theta(1\oplus R)^{-1}\bigl(\theta x-(t,\vec a)\bigr)=(1\oplus R^{-1})x+(t,-R^{-1}\vec a).
\end{equation*}
As $t\ge 0$, $g^*\in\mathfrak{E}_+$. Finally, $\theta^2=\mathbb{1}$ gives
$(g^*)^*=\theta\,\theta g\theta\,\theta=g$ and
$(gh)^*=\theta h^{-1}\theta\,\theta g^{-1}\theta=h^*g^*$.
\end{proof}

\begin{lemma}[Reflected separations grow under forward time translation]
\label{10.7:lem-reflected-separations}
Write points of $\mathbb{R}^4$ as $x = (x^0,\vec x)$ with $x^0 \in \mathbb{R}$ and
$\vec x \in \mathbb{R}^3$. Let $e_0 = (1,\vec 0)$ be the unit vector in the time direction,
and let $\theta$ be the time reflection $\theta(x^0,\vec x) = (-x^0,\vec x)$. For
$a, b \in \mathbb{R}^4$ with $a^0 \ge 0$ and $b^0 \ge 0$, and for $s \ge 0$,
\begin{equation}\label{10.7:eq-reflected-separations-1}
\bigl|\theta(a+se_0) - (b+se_0)\bigr|^2 = (a^0+b^0+2s)^2 + |\vec a - \vec b|^2
\ge |\theta a - b|^2 .
\end{equation}
\end{lemma}

\begin{proof}
The point $a + se_0$ has time component $a^0+s$ and spatial component $\vec a$, so
$\theta(a+se_0) = (-(a^0+s),\vec a)$. The point $b+se_0$ is $(b^0+s,\vec b)$. Hence
\begin{equation*}
\theta(a+se_0) - (b+se_0) = \bigl(-(a^0+b^0+2s),\ \vec a - \vec b\bigr).
\end{equation*}
The squared Euclidean norm of this vector is the square of its time component plus the
squared norm of its spatial component. This gives the equality in
\eqref{10.7:eq-reflected-separations-1}.

The equality holds for every $s \ge 0$. At $s = 0$ it reads
\begin{equation*}
|\theta a - b|^2 = (a^0+b^0)^2 + |\vec a - \vec b|^2 .
\end{equation*}
Since $a^0, b^0 \ge 0$ and $s \ge 0$, we have
\begin{equation*}
a^0+b^0+2s \ge a^0+b^0 \ge 0 .
\end{equation*}
Squaring preserves this inequality between non-negative numbers, so
$(a^0+b^0+2s)^2 \ge (a^0+b^0)^2$. Adding $|\vec a - \vec b|^2$ to both sides and using
the case $s = 0$ of the equality gives the inequality in
\eqref{10.7:eq-reflected-separations-1}.
\end{proof}

\begin{definition}[Centred correlation functions and their finite-cutoff form]
\label{10.7:def-centred-functions}
For an integer $k\ge 0$ let $\Pi^*(k)$ denote the set of partitions of $\{1,\dots,k\}$ without
singleton blocks, and for a finite set $S$ let $\Pi(S)$ denote the set of all partitions of $S$; the
empty set has exactly one partition, the empty one. For $\vec h=(h_1,\dots,h_k)$ and
$B\subset\{1,\dots,k\}$ write $\vec h_B=(h_l)_{l\in B}$. With $S^T_n$ the connected functions of the
subsequential limit of Definition~\ref{10.7:def-reflection-data}, the \emph{centred functions} are
\begin{equation}\label{10.7:eq-centred-functions-1}
S^c_k(\vec h):=\sum_{\pi\in\Pi^*(k)}\ \prod_{B\in\pi}S^T_{|B|}(\vec h_B),
\end{equation}
so that $S^c_0=1$ and $S^c_1=0$.

At finite cutoff, let $\langle\cdot\rangle_{K,m}$ be the expectation in the lattice measure at bare
spacing $L^{-K}$ on $\mathbb{T}_m$, let $h_l^{(m)}$ be the periodic lift of $h_l$ to $\mathbb{T}_m$,
and put $X_l:=\mathcal{O}(h_l^{(m)})$ and $Y_l:=X_l-\langle X_l\rangle_{K,m}$; for $B$ non-empty let
$\kappa(\vec h_B)$ be the joint cumulant of $(Y_l)_{l\in B}$ under $\langle\cdot\rangle_{K,m}$. For
$|B|\ge 2$ it equals the joint cumulant of $(X_l)_{l\in B}$, the two generating logarithms differing
by the linear term $\sum_l z_l\langle X_l\rangle_{K,m}$ (see the proof). Then:
\begin{enumerate}
\item[(a)] the centred moment has the finite-cutoff form
\begin{equation}\label{10.7:eq-centred-functions-2}
\Bigl\langle\prod_{l=1}^k Y_l\Bigr\rangle_{K,m}
=\sum_{\pi\in\Pi^*(k)}\ \prod_{B\in\pi}\kappa(\vec h_B);
\end{equation}
\item[(b)] for every admissible $\vec h$ (Definition~\ref{10.7:def-reflection-data}), with
$(K_j,m_j)$ the diagonal sequence and $\mathsf{d}$, $r$ as there, and with $Y_{j,l}$ the variable
$Y_l$ at cutoff $(K_j,m_j)$,
\begin{equation*}
S^c_k(\vec h)=\lim_{j}\Bigl\langle\prod_{l}Y_{j,l}\Bigr\rangle_{K_j,m_j};
\end{equation*}
\item[(c)] $S^c_k$ is symmetric, real and multilinear, invariant under every motion leaving all
$S^T_n$ invariant, and, for every admissible $\vec h$,
\begin{equation}\label{10.7:eq-centred-functions-3}
|S^c_k(\vec h)|\le\mathbf{C}_k(\mathsf{d};r)\prod_{l=1}^k\|h_l\|_{\sharp},
\qquad
\mathbf{C}_k(\mathsf{d};r):=\sum_{\pi\in\Pi^*(k)}\ \prod_{B\in\pi}C^\infty_{|B|}(\mathsf{d};r).
\end{equation}
\end{enumerate}
\end{definition}

\begin{proof}
(a) At fixed $(K,m)$ the lattice is finite, the bond variables range over the compact group
$SU(N)$ and $1-\operatorname{Re}\operatorname{tr}U(\partial p)$ is bounded, so each $X_l$, hence
each $Y_l$, is bounded. Therefore
\begin{equation*}
F(z):=\Bigl\langle \exp\Bigl(\sum_{l}z_lY_l\Bigr)\Bigr\rangle_{K,m},\qquad z\in\mathbb{C}^k,
\end{equation*}
is entire, with $F(0)=1$, and the principal logarithm $\mathcal{L}_Y:=\log F$ is analytic on a
neighbourhood of $0$, where $F=e^{\mathcal{L}_Y}$. By the definition of the joint cumulant,
$\kappa(\vec h_S)=\partial_{z_S}\mathcal{L}_Y(0)$ for non-empty $S$,
where $\partial_{z_S}:=\prod_{l\in S}\partial/\partial z_l$. Define $\mathcal{L}_X$ in the same way
with $X_l$ in place of $Y_l$, so that $\partial_{z_S}\mathcal{L}_X(0)$ is the joint cumulant of
$(X_l)_{l\in S}$. Since
$\sum_l z_lY_l=\sum_l z_lX_l-\sum_l z_l\langle X_l\rangle_{K,m}$, we have
\begin{equation*}
\mathcal{L}_Y(z)=\mathcal{L}_X(z)-\sum_l z_l\langle X_l\rangle_{K,m}.
\end{equation*}
Hence $\partial_{z_l}\mathcal{L}_Y(0)=\langle X_l\rangle_{K,m}-\langle X_l\rangle_{K,m}=0$, while
for $|S|\ge 2$ the linear term has vanishing derivative and $\kappa(\vec h_S)$ equals
$\partial_{z_S}\mathcal{L}_X(0)$, the joint cumulant of $(X_l)_{l\in S}$.

We show by induction on $|S|$ that, near $z=0$ and with $\mathcal{L}=\mathcal{L}_Y$,
\begin{equation*}
\partial_{z_S}e^{\mathcal{L}}
=e^{\mathcal{L}}\sum_{\pi\in\Pi(S)}\ \prod_{B\in\pi}\partial_{z_B}\mathcal{L}.
\end{equation*}
For $S=\emptyset$ both sides equal $e^{\mathcal{L}}$, the empty partition contributing the empty
product $1$. Let the identity hold for $S$ and let $t\notin S$. Applying $\partial_{z_t}$ and the
product rule (mixed partial derivatives of the analytic $\mathcal{L}$ commute),
$\partial_{z_{S\cup\{t\}}}e^{\mathcal{L}}$ equals
\begin{multline*}
e^{\mathcal{L}}\sum_{\pi\in\Pi(S)}\Bigl[\partial_{z_t}\mathcal{L}\prod_{B\in\pi}\partial_{z_B}\mathcal{L}\\
+\sum_{B\in\pi}\partial_{z_{B\cup\{t\}}}\mathcal{L}\prod_{B'\in\pi,\,B'\neq B}\partial_{z_{B'}}\mathcal{L}\Bigr].
\end{multline*}
Every partition of $S\cup\{t\}$ arises exactly once from a term of this sum: either $\{t\}$
is a block, and removing it gives $\pi\in\Pi(S)$ (first term), or $t$ lies in a block $B\cup\{t\}$
with $B\neq\emptyset$, and removing $t$ gives $\pi\in\Pi(S)$ with the distinguished block $B$
(second term). This is the identity for $S\cup\{t\}$.

Take $S=\{1,\dots,k\}$ and $z=0$. Differentiating under the finite-dimensional integral of a bounded
integrand gives $\partial_{z_S}F(0)=\langle\prod_lY_l\rangle_{K,m}$. On the right, $e^{\mathcal{L}(0)}=1$;
every partition with a singleton block $\{l\}$ contains the factor $\partial_{z_l}\mathcal{L}_Y(0)=0$,
and the remaining partitions form $\Pi^*(k)$, each block carrying $\kappa(\vec h_B)$. This is
\eqref{10.7:eq-centred-functions-2}; for $k=0,1$ it reads $1=1$ and $0=0$.

(b) For $|B|\ge 2$, $\kappa(\vec h_B)$ is the connected Schwinger function $S^T_{|B|;K,m}$ of the
lifted sub-tuple $(h_l^{(m)})_{l\in B}$. Every sub-tuple $\vec h_B$ of an admissible $\vec h$ is
admissible, so the first item of \eqref{10.7:eq-reflection-data-a} gives, along the diagonal
sequence, the convergence
of each such factor at $(K_j,m_j)$ to $S^T_{|B|}(\vec h_B)$. The right side of
\eqref{10.7:eq-centred-functions-2} is a finite sum of finite products of these factors, so it
converges to \eqref{10.7:eq-centred-functions-1}, and so does the left side.

(c) Symmetry, reality and multilinearity of $S^c_k$ follow from
\eqref{10.7:eq-centred-functions-1} and the corresponding properties of the $S^T_n$
(Definition~\ref{10.7:def-reflection-data}): a permutation of $\{1,\dots,k\}$ maps $\Pi^*(k)$
bijectively onto itself and blocks onto blocks, a product of real numbers is real, and each index
$l$ lies in exactly one block of each $\pi$, so that each product is linear in $h_l$.
If a motion $\sigma$ leaves every $S^T_n$ invariant, each product in
\eqref{10.7:eq-centred-functions-1} is unchanged, hence so is $S^c_k$. For the bound, the estimate
recorded in \eqref{10.7:eq-reflection-data-a} gives, for each admissible sub-tuple,
\begin{equation*}
|S^T_{|B|}(\vec h_B)|\le C^\infty_{|B|}(\mathsf{d};r)\prod_{l\in B}\|h_l\|_{\sharp}.
\end{equation*}
Since the blocks of each $\pi\in\Pi^*(k)$ partition $\{1,\dots,k\}$, it follows that
\begin{equation*}
|S^c_k(\vec h)|
\le\sum_{\pi\in\Pi^*(k)}\ \prod_{B\in\pi}C^\infty_{|B|}(\mathsf{d};r)\prod_{l\in B}\|h_l\|_{\sharp}
=\mathbf{C}_k(\mathsf{d};r)\prod_{l=1}^k\|h_l\|_{\sharp},
\end{equation*}
which is \eqref{10.7:eq-centred-functions-3}.
\end{proof}

\begin{lemma}[Telescoping estimate for moved entries]\label{10.7:lem-telescoping-estimate}
Let $\vec h=(h_1,\dots,h_k)$ be admissible with radius $r$ and separation at least $\mathsf{d}$, and
put $\rho:=\max_l\rho_{h_l}$. Let $g(x)=Qx+b$, with $Q\in O(4)$ and $b\in\mathbb{R}^4$, be a Euclidean
motion of $\mathbb{R}^4$ such that
\begin{equation*}
\eta_g:=\|Q-\mathbb{1}\|\,\rho+|b|\le \mathsf{d}/4 .
\end{equation*}
For $I\subset\{1,\dots,k\}$ let $\vec h^{(g;I)}$ be the tuple obtained from $\vec h$ by replacing $h_i$
with $g\cdot h_i$ for every $i\in I$. Then
\begin{equation}\label{10.7:eq-telescoping-estimate-1}
\begin{split}
\bigl|S^c_k(\vec h^{(g;I)})-S^c_k(\vec h)\bigr|
&\le \mathbf{C}_k(\mathsf{d}/2;\,r+\mathsf{d}/4)\\
&\qquad\times\sum_{i\in I}\|g\cdot h_i-h_i\|_{\sharp}\prod_{l\neq i}\|h_l\|_{\sharp}.
\end{split}
\end{equation}
\end{lemma}

\begin{proof}
Throughout, $(g\cdot f)(x)=f(g^{-1}x)$, $\rho_f$ is the support radius of $f$, so that
$\operatorname{supp}f\subset\{|x|\le\rho_f\}$, and $A^{+\theta}$ is the closed
$\theta$-neighbourhood of $A$. Of admissibility with radius $r$ and separation at least $\mathsf{d}$
we use that every entry is supported in a ball of radius $r$ and that the supports of distinct
entries are at distance at least $\mathsf{d}$.

Supports. If $|x|\le\rho$, then $|g(x)-x|=|(Q-\mathbb{1})x+b|\le\eta_g$. Hence, for every $l$,
\begin{equation*}
\operatorname{supp}(g\cdot h_l)=g(\operatorname{supp}h_l)\subset(\operatorname{supp}h_l)^{+\eta_g},
\end{equation*}
and $g\cdot h_i-h_i$ is supported in
$\operatorname{supp}h_i\cup g(\operatorname{supp}h_i)\subset(\operatorname{supp}h_i)^{+\eta_g}$.
If $\operatorname{supp}h_l$ lies in a ball of radius $r$, then $(\operatorname{supp}h_l)^{+\eta_g}$
lies in the concentric ball of radius $r+\eta_g\le r+\mathsf{d}/4$. For $l\neq l'$ the sets
$(\operatorname{supp}h_l)^{+\eta_g}$ and $(\operatorname{supp}h_{l'})^{+\eta_g}$ are at distance at
least $\mathsf{d}-2\eta_g\ge\mathsf{d}/2$.

Norms. Since $g^{-1}$ is an isometry, $\|g\cdot h_l\|_\infty=\|h_l\|_\infty$ and
$\operatorname{Lip}(g\cdot h_l)=\operatorname{Lip}h_l$, so $\|g\cdot h_l\|_{\sharp}=\|h_l\|_{\sharp}$:
moving an entry does not change its $\sharp$-norm.

Telescoping. For $I=\emptyset$ there is nothing to prove. Otherwise write
$I=\{i_1<\dots<i_p\}$ and set $\vec h^{[j]}:=\vec h^{(g;\{i_1,\dots,i_j\})}$ for $0\le j\le p$, so that
$\vec h^{[0]}=\vec h$ and $\vec h^{[p]}=\vec h^{(g;I)}$. Then
\begin{equation*}
S^c_k(\vec h^{(g;I)})-S^c_k(\vec h)=\sum_{j=1}^{p}\bigl(S^c_k(\vec h^{[j]})-S^c_k(\vec h^{[j-1]})\bigr).
\end{equation*}
The tuples $\vec h^{[j]}$ and $\vec h^{[j-1]}$ differ only in the slot $i_j$. Each summand of
$S^c_k$ is a product over the blocks of a partition in $\Pi^*(k)$, and the slot $i_j$ lies in
exactly one block of each such partition. Telescoping these products over blocks, only the factor
of the block containing $i_j$ changes; by the multilinearity stated among the
reflection-positivity data \eqref{10.7:eq-reflection-data-a} of
Definition~\ref{10.7:def-reflection-data}, that factor, and hence $S^c_k$, is linear in the slot
$i_j$. Therefore the $j$-th
difference equals $S^c_k(\vec h^{[j,\Delta]})$, where $\vec h^{[j,\Delta]}$ carries
$g\cdot h_{i_j}-h_{i_j}$ in the slot $i_j$ and agrees with $\vec h^{[j-1]}$ elsewhere. By the
paragraph on supports, the $l$-th entry of $\vec h^{[j,\Delta]}$ is supported in
$(\operatorname{supp}h_l)^{+\eta_g}$ for every $l$; so every entry lies in a ball of radius
$r+\mathsf{d}/4$, and distinct entries are at distance at least $\mathsf{d}/2$. The bound for
$S^c_k$ stated in \eqref{10.7:eq-reflection-data-a} therefore applies with separation
$\mathsf{d}/2$ and radius $r+\mathsf{d}/4$. Together with the invariance of the $\sharp$-norms of
the moved entries $g\cdot h_{i_1},\dots,g\cdot h_{i_{j-1}}$, it gives
\begin{equation*}
\bigl|S^c_k(\vec h^{[j,\Delta]})\bigr|\le\mathbf{C}_k(\mathsf{d}/2;\,r+\mathsf{d}/4)\,
\|g\cdot h_{i_j}-h_{i_j}\|_{\sharp}\prod_{l\neq i_j}\|h_l\|_{\sharp}.
\end{equation*}
Summing over $j=1,\dots,p$ and using the triangle inequality gives
\eqref{10.7:eq-telescoping-estimate-1}.
\end{proof}

We note the consequence that is used below. Let $f$ be of class $C^1$ with
$\operatorname{supp}f\subset\{|x|\le\rho_f\}$. Let $m_{\nabla f}(\delta)$ be the supremum of
$|\nabla f(x)-\nabla f(y)|$ over $|x-y|\le\delta$, and put
$\delta_g:=\|Q-\mathbb{1}\|(\rho_f+|b|)+|b|$. Then
\begin{equation}\label{10.7:eq-telescoping-estimate-2}
\|g\cdot f-f\|_{\sharp}\le\max\Bigl\{\|\nabla f\|_\infty\,\delta_g,\;
40M\bigl(\|Q-\mathbb{1}\|\,\|\nabla f\|_\infty+m_{\nabla f}(\delta_g)\bigr)\Bigr\}.
\end{equation}

Let $x$ be a point at which $f(g^{-1}x)$ and $f(x)$ are not both zero, or at which
$\nabla f(g^{-1}x)$ and $\nabla f(x)$ are not both zero. Then $x$ or $g^{-1}x$ lies in
$\operatorname{supp}f$. If $g^{-1}x=y\in\operatorname{supp}f$, then $x=Qy+b$ with $|y|\le\rho_f$;
if $x\in\operatorname{supp}f$, then $|x|\le\rho_f$. In either case
$|x|\le\rho_f+|b|$. Moreover $g^{-1}x-x=(Q^{-1}-\mathbb{1})x-Q^{-1}b$, and $Q^{-1}$ is orthogonal,
so $\|Q^{-1}-\mathbb{1}\|=\|Q^{-1}(\mathbb{1}-Q)\|=\|Q-\mathbb{1}\|$. Hence
$|g^{-1}x-x|\le\delta_g$.

By the mean value theorem this gives $|f(g^{-1}x)-f(x)|\le\|\nabla f\|_\infty\delta_g$, which is
the first entry of the maximum.

For the second entry, note that $\nabla(g\cdot f)(x)=Q\,\nabla f(g^{-1}x)$. Therefore
\begin{equation*}
\nabla(g\cdot f-f)(x)=(Q-\mathbb{1})\nabla f(g^{-1}x)+\bigl(\nabla f(g^{-1}x)-\nabla f(x)\bigr).
\end{equation*}
The first term has modulus at most $\|Q-\mathbb{1}\|\,\|\nabla f\|_\infty$, and the second at most
$m_{\nabla f}(\delta_g)$. Since $\operatorname{Lip}u\le\|\nabla u\|_\infty$ for $u\in C^1$,
\eqref{10.7:eq-telescoping-estimate-2} follows.

As $g\to\mathrm{id}$, that is as $Q\to\mathbb{1}$ and $b\to0$, both $\delta_g\to0$ and
$m_{\nabla f}(\delta_g)\to0$, the latter because $\nabla f$ is continuous with compact support and
hence uniformly continuous. Thus $\|g\cdot f-f\|_{\sharp}\to0$.

Now let $(\theta\cdot\vec u,\vec f)$ be admissible with $k$ entries, radius $r$ and separation at
least $\mathsf{d}$, and let $g\in\mathfrak{E}_+$ tend to the identity. Then $\eta_g\to0$, so
$\eta_g\le\mathsf{d}/4$ eventually. Apply Lemma~\ref{10.7:lem-telescoping-estimate} with $I$ the
set of slots occupied by $\vec f$. The resulting bound has a constant
$\mathbf{C}_k(\mathsf{d}/2;\,r+\mathsf{d}/4)$ that does not depend on $g$, and it is multiplied by
$\sum_i\|g\cdot f_i-f_i\|_{\sharp}$, which tends to $0$ by
\eqref{10.7:eq-telescoping-estimate-2}. Hence
$S^c_k(\theta\cdot\vec u,\,g\cdot\vec f)\to S^c_k(\theta\cdot\vec u,\vec f)$ as $g\to\mathrm{id}$
within $\mathfrak{E}_+$.

Let $\mathbb{R}^4_+:=\{x\in\mathbb{R}^4:\ x^0>0\}$. For a tuple $\vec u=(u_1,\dots,u_n)$ write
$\theta\cdot\vec u:=(\theta\cdot u_1,\dots,\theta\cdot u_n)$, where $\theta$ is the time reflection.

Let $\mathcal{V}_+$ be the free complex vector space on the symbols $[\vec f\,]$. Here $\vec f$ runs over
the admissible tuples all of whose supports lie in $\mathbb{R}^4_+$, and the symbol $[\emptyset]$ of the
empty tuple is included.

For two such symbols put
\begin{equation}\label{10.7:eq-reflection-form-1}
\langle[\vec u\,],[\vec f\,]\rangle:=S^c(\theta\cdot\vec u,\vec f\,).
\end{equation}
Here $(\theta\cdot\vec u,\vec f\,)$ is the concatenated tuple. It is admissible, because the reflected
supports lie in $\{x^0<0\}$. The form \eqref{10.7:eq-reflection-form-1} is extended to
$\mathcal{V}_+\times\mathcal{V}_+$ conjugate-linearly in the first and linearly in the second argument.

For $g\in\mathfrak{E}_+$ put $g_\sharp[\vec f\,]:=[g\cdot\vec f\,]$, extended linearly.

\begin{proposition}[Positivity of the reflection form and the Hilbert space]
\label{10.7:prop-reflection-form}
The form \eqref{10.7:eq-reflection-form-1} on $\mathcal{V}_+$ is Hermitian and positive semidefinite.
\end{proposition}

\begin{proof}
\emph{Hermiticity.} We use three properties of $S^c$, the centred correlation function of the
subsequential limit (Definition~\ref{10.7:def-reflection-data}):
\begin{itemize}
\item it is real;
\item it is symmetric under permutations of its entries;
\item it is invariant under $\theta$.
\end{itemize}
Since $\theta^2$ is the identity, $\theta$-invariance followed by symmetry gives, for symbols
$[\vec u\,],[\vec f\,]$,
\begin{equation*}
\langle[\vec f\,],[\vec u\,]\rangle=S^c(\theta\cdot\vec f,\vec u)
=S^c(\vec f,\theta\cdot\vec u)=S^c(\theta\cdot\vec u,\vec f\,)=\langle[\vec u\,],[\vec f\,]\rangle.
\end{equation*}
Admissible tuples have real entries, so the last quantity is real by
Definition~\ref{10.7:def-reflection-data}, and it equals its complex conjugate.

Now take $w=\sum_\alpha a_\alpha[\vec u_\alpha]$ and $v=\sum_\beta b_\beta[\vec f_\beta]$. Sesquilinearity
then gives
\begin{equation*}
\overline{\langle w,v\rangle}
=\sum_{\alpha,\beta}a_\alpha\bar b_\beta\,\overline{\langle[\vec u_\alpha],[\vec f_\beta]\rangle}
=\sum_{\alpha,\beta}\bar b_\beta a_\alpha\langle[\vec f_\beta],[\vec u_\alpha]\rangle
=\langle v,w\rangle .
\end{equation*}

\emph{Positivity: the lattice variables.} Let $v=\sum_\alpha c_\alpha[\vec f_\alpha]$ with
$\vec f_\alpha=(f_{\alpha,1},\dots,f_{\alpha,n_\alpha})$. Define
\begin{itemize}
\item $\delta>0$, the minimum of $x^0$ over all occurring supports;
\item $\rho$, the largest support radius $\rho_{f_{\alpha,i}}$ of the occurring entries.
\end{itemize}

Let $(K_j,m_j)$ be the subsequence along which $S^c$ is the limit of the finite-cutoff centred functions
(Definition~\ref{10.7:def-centred-functions}). Let $\varepsilon_j:=L^{-K_j}$ be the lattice spacing and
$\langle\cdot\rangle_{K_j,m_j}$ the expectation in the lattice measure on $\mathbb{T}_{m_j}$.

Consider $j$ with $\varepsilon_j\le\delta$ and $L^{m_j}>\rho+1$. These inequalities hold for all
sufficiently large $j$.

Write $m=m_j$ and $\varepsilon=\varepsilon_j$. For such $j$, the plaquettes $p$ entering
$\mathcal{O}(f^{(m)}_{\alpha,i})$ have barycentres $x(p)$ in the occurring supports. Their corners
therefore satisfy
\begin{equation*}
x^0\in[\delta-\varepsilon/2,\ \rho+\varepsilon/2]\subset\,]0,L^m[ .
\end{equation*}
Hence each $\mathcal{O}(f^{(m)}_{\alpha,i})$ depends only on bonds with both endpoints in the open half-torus
$H^+_{0,0}$.

Put $X_{\alpha,i}:=\mathcal{O}(f^{(m)}_{\alpha,i})$ and let
$Y_{\alpha,i}:=X_{\alpha,i}-\langle X_{\alpha,i}\rangle_{K_j,m_j}$ be the centred variables. Then
\begin{equation*}
F:=\sum_\alpha c_\alpha\prod_i Y_{\alpha,i}
\end{equation*}
is measurable with respect to the bonds in $H^+_{0,0}$.

Let $r_{0,0}$ be the lattice reflection. Reflection positivity of the lattice measure, applied with
$\nu=0$ and $c=0$, gives
\begin{equation}\label{10.7:eq-reflection-form-2}
\big\langle\,\overline{F\circ r_{0,0}}\;F\,\big\rangle_{K_j,m_j}\ge 0 .
\end{equation}

\emph{Positivity: the reflected variables.} By the definition of the transformed configuration
$U\mapsto r_{0,0}U$ under lattice isometries (the functions $F^{\sigma}$ of the isometry-transformed
configuration),
\begin{equation*}
X_{\alpha,i}(r_{0,0}U)=\mathcal{O}\big((\theta\cdot f_{\alpha,i})^{(m)}\big)(U).
\end{equation*}
This uses that the periodic lift commutes with $\theta$, which holds because
$\theta(2L^m\mathbb{Z}^4)=2L^m\mathbb{Z}^4$.

By the invariances of the lattice measure,
$\langle X\rangle_{K_j,m_j}=\langle X\circ r_{0,0}\rangle_{K_j,m_j}$. Consequently $Y_{\alpha,i}\circ r_{0,0}$
is the centred variable $Y'_{\alpha,i}$ of $\theta\cdot f_{\alpha,i}$. The $Y'_{\alpha,i}$ are real, and
therefore
\begin{equation*}
\overline{F\circ r_{0,0}}=\sum_\alpha\bar c_\alpha\prod_i Y'_{\alpha,i}.
\end{equation*}

\emph{Positivity: passage to the limit.} Substituting this into \eqref{10.7:eq-reflection-form-2} yields
\begin{equation*}
0\le\sum_{\alpha,\beta}\bar c_\alpha c_\beta
\Big\langle\prod_i Y'_{\alpha,i}\prod_l Y_{\beta,l}\Big\rangle_{K_j,m_j}.
\end{equation*}
Each expectation on the right is the finite-cutoff centred function of the admissible tuple
$(\theta\cdot\vec f_\alpha,\vec f_\beta)$. Along the subsequence it converges to
$S^c(\theta\cdot\vec f_\alpha,\vec f_\beta)$, by Definition~\ref{10.7:def-centred-functions}.

Letting $j\to\infty$ and using \eqref{10.7:eq-reflection-form-1} together with sesquilinearity, we obtain
\begin{equation*}
0\le\sum_{\alpha,\beta}\bar c_\alpha c_\beta\,S^c(\theta\cdot\vec f_\alpha,\vec f_\beta)
=\langle v,v\rangle .
\qedhere
\end{equation*}
\end{proof}

Let
\begin{equation*}
\mathcal{V}_+^{0}:=\{v\in\mathcal{V}_+:\ \langle w,v\rangle=0\ \text{for all } w\in\mathcal{V}_+\}
\end{equation*}
be the null space of the form. Let $\mathcal{H}$ be the completion of $\mathcal{V}_+/\mathcal{V}_+^{0}$, let
$\iota:\mathcal{V}_+\to\mathcal{H}$ be the quotient map, and put $\Omega:=\iota[\emptyset]$. Since
$\langle[\emptyset],[\emptyset]\rangle=S^c(\emptyset)=1$, we have $\|\Omega\|=1$.

The space $\mathcal{H}$ is separable. Indeed, consider a countable family of $C^1_c(\mathbb{R}^4_+)$
functions, each supported in a ball with rational centre and rational radius, which is
$\|\cdot\|_{\sharp}$-dense in the $C^1_c$ functions supported in such balls, and the tuples whose entries
lie in this family. By the telescoping estimate for moved entries
(Lemma~\ref{10.7:lem-telescoping-estimate}), these tuples have dense $\iota$-image.

\begin{remark}\label{10.7:thm-rotation-representation}
The statement and proof of this theorem are not written out here.
\end{remark}

\begin{remark}[Rotations in planes containing the time axis]\label{10.7:rem-time-plane-rotations}
Stated; proof incomplete at the step marked $(\ast)$.
Assume the hypotheses of the $O(4)$-invariance theorem. Let $e_0$ be the unit vector of the
time axis, fix a plane containing $e_0$, let $\omega$ be the generator of the rotations in that
plane, and let $R^{\omega}_s = e^{s\omega}$, $s \in \mathbb{R}$, be the rotation by the angle $s$
in that plane. For symbols in $\mathcal{V}_+$ whose supports are mapped into the positive
half-space by $R^{\omega}_s$ for all $|s| < s_0$, for some $s_0 > 0$ (symbols supported in a
wedge), put $P_s \xi := (R^{\omega}_s)_{\sharp}\xi$, the image of $\xi$ under the action of
$R^{\omega}_s$ on symbols as in Theorem~\ref{10.7:thm-rotation-representation}, for
$|s| < s_0$. Then $(P_s)$ is a symmetric local semigroup on these symbols:
for such $\xi, \eta$ and $|s| < s_0/2$,
\begin{equation*}
\langle P_s\xi, \eta\rangle = \langle \xi, P_s\eta\rangle,
\qquad
\|P_s\xi\|^2 = \langle \xi, P_{2s}\xi\rangle,
\end{equation*}
and the matrix elements $\langle P_s\xi, \eta\rangle$, and $s \mapsto P_s\xi$ for the seminorm
$\|\cdot\|$, are continuous in $s$. Two further assertions concern this structure:
\begin{equation}\label{10.7:eq-time-plane-rotations-a}
\begin{aligned}
\text{(M1)}\quad &\text{such a structure has a unique self-adjoint generator,}\\
&\text{with the symbols as analytic vectors;}\\
\text{(M2)}\quad &\text{there is a strongly continuous unitary representation of }
\mathbb{R}^4 \rtimes SL(2,\mathbb{C})\\
&\text{with } \Omega \text{ invariant.}
\end{aligned}
\end{equation}
Under the hypotheses of the $O(4)$-invariance theorem the limit is $O(4)$-invariant.
The spectrum condition beyond $H \ge 0$ would follow only from (M1) and (M2) of
\eqref{10.7:eq-time-plane-rotations-a}, neither of which is supplied here: (M1) is stated
with its proof incomplete, and (M2) is not proved.
\end{remark}

\begin{proof}
The semigroup properties are proved as an implication from the $O(4)$-invariance theorem.
That theorem gives invariance of the centred functions $S^c_k$ under each $R^{\omega}_s$.
Since $\theta$ reverses the $e_0$-axis of the plane and fixes its other axis, conjugating
$R^{\omega}_s$ by $\theta$ reverses the angle:
\begin{equation*}
\theta R^{\omega}_s\theta = R^{\omega}_{-s}.
\end{equation*}
Hence $\theta (R^{\omega}_s)^{-1}\theta = \theta R^{\omega}_{-s}\theta = R^{\omega}_s$. The
computation that gives symmetry of the transfer semigroup in the rotation representation theorem
(Theorem~\ref{10.7:thm-rotation-representation}) applies with $g = R^{\omega}_s$ and yields
$\langle P_s\xi, \eta\rangle = \langle \xi, P_s\eta\rangle$.

Since $R^{\omega}_sR^{\omega}_s = R^{\omega}_{2s}$, we have $P_sP_s = P_{2s}$ for
$|s| < s_0/2$, all rotations involved then keeping the supports in the positive half-space.
The symmetry then gives
\begin{equation*}
\|P_s\xi\|^2 = \langle P_s\xi, P_s\xi\rangle = \langle \xi, P_sP_s\xi\rangle
= \langle \xi, P_{2s}\xi\rangle.
\end{equation*}
Continuity in $s$ follows from the telescoping estimate for moved entries
(Lemma~\ref{10.7:lem-telescoping-estimate}). The matrix elements $\langle P_s\xi,\eta\rangle$
are built from the centred functions $S^c_k$ evaluated at tuples in which the entries of
$\xi$ are moved by $R^{\omega}_s$. For $s$ close to $t$, the lemma is applied with
$g = R^{\omega}_{s-t}$ to the tuples already moved by $R^{\omega}_t$; here $g$ has no
translation part, so
\begin{equation*}
\eta_g = \|R^{\omega}_{s-t} - \mathbb{1}\|\,\max_l \rho_{h_l} \le \mathsf{d}/4
\end{equation*}
once $|s-t|$ is small, $\rho_{h_l}$ being the support radii of the entries $h_l$. The bound
the lemma gives tends to zero as $s \to t$ for entries of class $C^1$ with compact support:
for such an entry $h$, $\|R^{\omega}_{s-t}\cdot h - h\|_\infty \to 0$ and
\begin{equation*}
\operatorname{Lip}(R^{\omega}_{s-t}\cdot h - h)
\le \|\nabla(R^{\omega}_{s-t}\cdot h) - \nabla h\|_\infty \to 0
\end{equation*}
as $s \to t$, so that $\|R^{\omega}_{s-t}\cdot h - h\|_{\sharp} \to 0$. Thus
$s \mapsto \langle P_s\xi,\eta\rangle$ is continuous. By the symmetry and
$P_sP_t = P_{s+t}$, for $|s|, |t| < s_0/2$,
\begin{equation*}
\|P_s\xi - P_t\xi\|^2 = \langle \xi, P_{2s}\xi\rangle - 2\langle \xi, P_{s+t}\xi\rangle
+ \langle \xi, P_{2t}\xi\rangle,
\end{equation*}
which tends to zero as $s \to t$; so $s \mapsto P_s\xi$ is continuous for $\|\cdot\|$.

$(\ast)$ Assertion (M1) of \eqref{10.7:eq-time-plane-rotations-a} is stated; its proof is
incomplete, the construction of the generator not being carried out here. That is the
assertion that the symmetric local semigroup $(P_s)$ has a unique self-adjoint generator with
the symbols as analytic vectors.

Assertion (M2) of \eqref{10.7:eq-time-plane-rotations-a} is not proved. It would require an
integrability theorem for representations of Lie algebras by unbounded operators, and no
such theorem is among the results this book establishes or cites.
\end{proof}

\begin{definition}[The three hypotheses in the form used]\label{10.7:def-hypotheses-used-form}
Fix a diagonal sequence $(K_j,m_j)_{j\ge 1}$ of pairs (number of renormalisation steps, torus
exponent) and an integer $r\ge 2$, and let
$\omega\in\mathfrak{so}(4)$, $n\ge 2$, $\mathcal{R}>0$ and $\mathsf{d}>0$ be arbitrary; the
cut-off functions $\chi$ that occur below range over those admissible at radius $\mathcal{R}$
in the sense of the rotational Ward identity (Theorem~\ref{10.1:thm-rotational-identity}). Test
functions, the class $\mathcal{T}^{(r)}_n(\mathcal{R},\mathsf{d})$ of separated real $C^r_c$
tuples supported in $B(0,\mathcal{R}/2)$, and the action $\vec h\mapsto R^{\omega}_t\cdot\vec h$
of the one-parameter rotation group $R^{\omega}_t=e^{t\omega}$ on such tuples are as in
Definition~\ref{10.7:not-test-function-moduli}. For a test function $h$ we write $h^{(m)}$ for
its periodic lift to $\mathbb{T}_m$, and $\kappa(Y_1,Y_2,\dots)$ for the joint cumulant of
finitely many bounded random variables $Y_1,Y_2,\dots$ under the lattice measure on
$\mathbb{T}_m$ at bare spacing $L^{-K}$. For
$\vec h=(h_1,\dots,h_n)\in\mathcal{T}^{(r)}_n(\mathcal{R},\mathsf{d})$ put
\begin{equation*}
D^{\omega}_{K,m}(\vec h):=\frac{d}{dt}\Big|_{t=0} S^T_{n;K,m}\bigl(R^{\omega}_t\cdot\vec h\bigr).
\end{equation*}
The three hypotheses in the form used are the following.

\begin{enumerate}
\item[(a)] \emph{Connected splitting (the rotational Ward identity in connected form).} For every
$\chi$ admissible at radius $\mathcal{R}$ there
are an index $j_2=j_2(n,\mathcal{R},\mathsf{d},\omega,\chi)$ (equivalently, the cutoff floor
$K\ge K_1$ and the volume floor $L^m\ge\max\{\ell_1(\chi),\mathcal{R}\}$ of
Theorem~\ref{10.1:thm-rotational-identity}) and, for $(K,m)=(K_j,m_j)$ with $j\ge j_2$, bounded
Borel functions $N^{\omega,\chi}_{K,m}$ and $A^{\omega,\chi}_{K,m}$ of the bare field and, for
each real $h\in C^r_c$ with $\operatorname{supp}h\subset B(0,\mathcal{R}/2)$, a bounded Borel
function $E^{\omega}_{h;K,m}$ of the bare field, such that for every
$\vec h\in\mathcal{T}^{(r)}_n(\mathcal{R},\mathsf{d})$
\begin{equation}\label{10.7:eq-hypotheses-used-form-a}
D^{\omega}_{K,m}(\vec h)=\mathcal{N}^{\omega,\chi}_{K,m}(\vec h)
+\mathcal{B}^{\omega,\chi}_{K,m}(\vec h),
\end{equation}
where
\begin{align*}
\mathcal{N}^{\omega,\chi}_{K,m}(\vec h)
&:=\kappa\bigl(N^{\omega,\chi}_{K,m},\mathcal{O}(h_1^{(m)}),\dots,\mathcal{O}(h_n^{(m)})\bigr)\\
&\qquad+\sum_{i=1}^{n}\kappa\bigl(\mathcal{O}(h_1^{(m)}),\dots,E^{\omega}_{h_i;K,m},\dots,
\mathcal{O}(h_n^{(m)})\bigr),\\
\mathcal{B}^{\omega,\chi}_{K,m}(\vec h)
&:=\kappa\bigl(A^{\omega,\chi}_{K,m},\mathcal{O}(h_1^{(m)}),\dots,\mathcal{O}(h_n^{(m)})\bigr),
\end{align*}
and in the $i$-th summand the entry $E^{\omega}_{h_i;K,m}$ stands in the $i$-th place, in
place of $\mathcal{O}(h_i^{(m)})$.

\item[(b)] \emph{Vanishing of the ultraviolet insertion.} For every
$\vec h\in\mathcal{T}^{(r)}_n(\mathcal{R},\mathsf{d})$ and every admissible $\chi$,
$\mathcal{N}^{\omega,\chi}_{K_j,m_j}(\vec h)\to 0$ as $j\to\infty$.

\item[(c)] \emph{Smallness of the infrared insertion along admissible cut-off functions.} For
every $\vec h\in\mathcal{T}^{(r)}_n(\mathcal{R},\mathsf{d})$,
\begin{equation*}
\inf\Bigl\{\limsup_{j\to\infty}\bigl|\mathcal{B}^{\omega,\chi}_{K_j,m_j}(\vec h)\bigr|
\;:\;\chi\ \text{admissible at radius}\ \mathcal{R}\Bigr\}=0.
\end{equation*}
\end{enumerate}

\emph{The weakest form consumed.} What is used of (a)--(c) is only this: for every
$\vec h\in\mathcal{T}^{(r)}_n(\mathcal{R},\mathsf{d})$ there are, for each $\chi$ admissible at
radius $\mathcal{R}$, an index $j(\chi)$, depending also on $n$, $\mathcal{R}$, $\mathsf{d}$ and
$\omega$, and real numbers $\mathcal{N}^\chi_j$, $\mathcal{B}^\chi_j$ for $j\ge j(\chi)$, such
that, with $D_j:=D^{\omega}_{K_j,m_j}(\vec h)$,
\begin{equation}\label{10.7:eq-hypotheses-used-form-b}
\begin{gathered}
D_j=\mathcal{N}^\chi_j+\mathcal{B}^\chi_j\quad(j\ge j(\chi)),\qquad
\lim_{j\to\infty}\mathcal{N}^\chi_j=0\ \text{for every admissible}\ \chi,\\
\inf_{\chi\ \text{admissible at radius}\ \mathcal{R}}\ \limsup_{j\to\infty}
\bigl|\mathcal{B}^\chi_j\bigr|=0.
\end{gathered}
\end{equation}
\end{definition}

In the application, $N^{\omega,\chi}_{K,m}$ collects every term of
Theorem~\ref{10.1:thm-rotational-identity}
that is classically suppressed by $\varepsilon^2$, where $\varepsilon=L^{-K}$ is the bare
spacing, namely the terms written $-\mathfrak{B}^{-\omega}_\chi$, the bulk
part of $J^{-\omega}_\chi$ and the $\varepsilon^2$-corrections of $B^{-\omega}_\chi$ in
Theorem~\ref{10.1:thm-rotational-identity}; $A^{\omega,\chi}_{K,m}$ is the total dimension-four
annulus functional
$\mathfrak{F}=B^4+J^4$, where
\begin{equation*}
B^4=\sum_x\varepsilon^4(-\omega)_{\rho\sigma}\,x_\rho\,\nabla_\nu\chi\,
\Theta^{\varepsilon}_{\nu\sigma}
\end{equation*}
and $J^4$ is the dimension-four annulus piece of the divergence term; and
$E^{\omega}_{h;K,m}$ is the observable-side remainder of
Theorem~\ref{10.1:thm-rotational-identity}. This
description is given for orientation only: the internal structure of $N^{\omega,\chi}_{K,m}$,
$A^{\omega,\chi}_{K,m}$ and $E^{\omega}_{h;K,m}$ is never used in the argument that consumes
the three hypotheses.

\emph{Passage from the product form.} Fix $(K,m)$, write $\langle\,\cdot\,\rangle$ for the
expectation under the lattice measure at that cutoff, and for a monomial
$F=\prod_{i\in I}\mathcal{O}(h_i^{(m)})$ with $I\subset\{1,\dots,n\}$ write $F\circ R_t$ for the
same monomial with each $h_i$ replaced by $R^{\omega}_t\cdot h_i$. Suppose that the rotational
Ward identity, Theorem~\ref{10.1:thm-rotational-identity}, in its product form, delivers
\begin{equation*}
\frac{d}{dt}\Big|_{t=0}\langle F\circ R_t\rangle=\langle F\cdot W\rangle+\langle E_F\rangle
\end{equation*}
for every such monomial, with one bounded Borel function $W$ of the bare field, independent of
$F$, and with the Leibniz remainder
\begin{equation*}
E_{\prod_{i\in I}\mathcal{O}(h_i^{(m)})}
=\sum_{i\in I}E_{h_i}\prod_{k\in I\setminus\{i\}}\mathcal{O}(h_k^{(m)}).
\end{equation*}
The case $F=1$, $I=\emptyset$, is included: there the left-hand side vanishes, the remainder is
$E_1=0$ (an empty sum), and the identity reduces to the Schwinger--Dyson identity with the
constant function $1$ as insertion; it gives $\langle W\rangle=0$, which is needed below. Under
these assumptions the lemma of this chapter on the passage from the product form to the
connected form of the cumulants gives \eqref{10.7:eq-hypotheses-used-form-a}
with $N^{\omega,\chi}_{K,m}+A^{\omega,\chi}_{K,m}:=W$ and $E^{\omega}_{h;K,m}:=E_h$; this mapping
is stated below as a corollary of that lemma.

\begin{lemma}[The forms used follow from the stated hypotheses]\label{10.7:lem-hypotheses-reduction}
Fix an infinitesimal rotation $\omega$ as in Definition~\ref{10.7:def-hypotheses-used-form}, a
radius $\mathcal{R}$ and a tuple $\vec h\in\mathcal{T}^{(r)}_n(\mathcal{R},\mathsf{d})$. Write
$\varepsilon=L^{-K}$ for the bare lattice spacing and $\varepsilon_j=L^{-K_j}$. Along the
diagonal $(K,m)=(K_j,m_j)$ and for $\chi$ admissible at radius $\mathcal{R}$, write
$\mathcal{N}^{\chi}_j=\mathcal{N}^{\omega,\chi}_{K_j,m_j}(\vec h)$ and
$\mathcal{B}^{\chi}_j=\mathcal{B}^{\omega,\chi}_{K_j,m_j}(\vec h)$ for the two terms of the decomposition in
Definition~\ref{10.7:def-hypotheses-used-form}.
\begin{enumerate}
\item[(a)] The dimension-six coefficient bound is stated in two forms, the first of which implies the
second. Suppose every
term of which the lattice function $N^{\omega,\chi}_{K,m}$ is a sum, and every
$E^{\omega}_{h_i;K,m}$, is an insertion of one of the three types $\mathfrak{B}$, $J$ and $E$-type
listed in the statement of the
dimension-six coefficient bound. Then the second form of the dimension-six coefficient
bound implies the ultraviolet hypothesis in the form used in
\eqref{10.7:eq-hypotheses-used-form-b}. A fortiori its first form, which implies the second,
does so.
\item[(b)] Suppose the far-sphere flux hypothesis of clause~(vi), $(\mathrm{IR}\text{-}\mathrm{fl})_{\omega_0}$,
is stated at the rotation $\omega$ fixed above and
for the functional $A^{\omega,\chi}_{K,m}$ of Definition~\ref{10.7:def-hypotheses-used-form}, in the
double-limit form
\begin{equation}\label{10.7:eq-hypotheses-reduction-2}
a_{\mathcal{R}'}:=\limsup_{j\to\infty}\bigl|\mathcal{B}^{\chi_{\mathcal{R}'}}_j\bigr|
\longrightarrow0
\qquad(\mathcal{R}'\to\infty),
\end{equation}
where, for $\mathcal{R}'\ge\mathcal{R}$, $\chi_{\mathcal{R}'}$ is the radial cut-off at radius
$\mathcal{R}'$; the
intended functional is the total annulus functional of order $\varepsilon^0$,
$\mathfrak{F}=B^4+J^4$. Then it implies the infrared hypothesis in the form used in
\eqref{10.7:eq-hypotheses-used-form-b}.
\item[(c)] The clause, in the section of this chapter stating the hypotheses, asserting ``the same
bounds with
$F\circ R_t$ in place of $F$, uniformly for $t$ in compact sets'' is not used.
\end{enumerate}
\end{lemma}

\begin{proof}
(a) For every insertion $I$ of the three types of (a), the second form of the dimension-six
coefficient bound gives
\begin{equation}\label{10.7:eq-hypotheses-reduction-1}
\bigl|\kappa\bigl(I;\mathcal{O}(\vec h)\bigr)\bigr|
\le C\,\eta(\varepsilon)\prod_{i=1}^n\|h_i\|_{\sharp},
\qquad
\eta(\varepsilon)=\varepsilon^{2-\vartheta}+e^{-c\,p_0(g_0(\varepsilon))}.
\end{equation}
Here $\mathcal{O}(\vec h)$ abbreviates the slots $\mathcal{O}(h_1^{(m)}),\dots,\mathcal{O}(h_n^{(m)})$,
$\vartheta<2$ is the exponent and $c>0$ the constant of that bound, $g_0(\varepsilon)$ is the bare
coupling at spacing $\varepsilon$, and
$C$ does not depend on $K$, $m$, $g_0$ or the positions of the supports. The cumulant $\kappa$
is linear in each of its slots. We expand $N^{\omega,\chi}_{K_j,m_j}$ into its terms and use the
hypothesis of (a) that each term, and each $E^{\omega}_{h_i;K_j,m_j}$, is an insertion of one of these
types. Applying
\eqref{10.7:eq-hypotheses-reduction-1} to every resulting cumulant gives
\begin{equation*}
|\mathcal{N}^{\chi}_j|\le n'\,C\,\eta(\varepsilon_j)\prod_{i=1}^n\|h_i\|_{\sharp},
\end{equation*}
where $n'$ is the number of cumulants so obtained.

Write $g_0^{(j)}:=g_0(\varepsilon_j)$ for the bare coupling of the $j$-th member of the diagonal.
Along the diagonal $g_0^{(j)}\downarrow 0$. Clause~(0) of Theorem~0 gives, under the standing
assumptions,
\begin{equation*}
K_j\ge \frac{(g_0^{(j)})^{-2}-g^{-2}}{B^{\sharp}}-1 .
\end{equation*}
Hence $K_j\to\infty$, and therefore $\varepsilon_j=L^{-K_j}\to0$. Since $\vartheta<2$, this gives
$\varepsilon_j^{2-\vartheta}\to0$. For the second term of $\eta$, the degree function $p_0(\cdot)$
satisfies
\begin{equation*}
p_0\bigl(g_0^{(j)}\bigr)=A_0\bigl(\log (g_0^{(j)})^{-2}\bigr)^{p_0}\longrightarrow\infty ,
\end{equation*}
so $e^{-c\,p_0(g_0^{(j)})}\to0$. Thus $\eta(\varepsilon_j)\to0$ and $\mathcal{N}^{\chi}_j\to0$ along
the diagonal, which is the form required in \eqref{10.7:eq-hypotheses-used-form-b}. The first form
implies the second, so it also yields this conclusion.

(b) Let $\mathcal{R}'\ge\mathcal{R}$. The radial cut-off $\chi_{\mathcal{R}'}$ is admissible at
radius $\mathcal{R}$. By the hypothesis of (b), that is, the far-sphere flux hypothesis for the
functional $A^{\omega,\chi}_{K,m}$ in its double-limit form \eqref{10.7:eq-hypotheses-reduction-2},
we have $a_{\mathcal{R}'}\to0$ as $\mathcal{R}'\to\infty$.
The infimum of $\limsup_j|\mathcal{B}^{\chi}_j|$ over the admissible $\chi$ is nonnegative. It is at
most $a_{\mathcal{R}'}$ for every $\mathcal{R}'\ge\mathcal{R}$, so it equals $0$. This is the form
required in \eqref{10.7:eq-hypotheses-used-form-b}.

(c) By inspection of the sections of this chapter that use the three hypotheses of
Definition~\ref{10.7:def-hypotheses-used-form}, those arguments evaluate the identity
\begin{equation*}
D^{\omega}_{K,m}(\vec h)=\mathcal{N}^{\omega,\chi}_{K,m}(\vec h)+\mathcal{B}^{\omega,\chi}_{K,m}(\vec h)
\end{equation*}
and both of its terms only at $t=0$, and only on tuples of the class
$\mathcal{T}^{(r)}_n(\mathcal{R},\mathsf{d})$, which the rotations $R^{\omega}_s$ preserve. A rotated
tuple $R^{\omega}_s\cdot\vec h$ is treated as a new member of that class, to which the
hypotheses apply directly at $t=0$.
Hence the uniform-in-$t$ clause is not used.
\end{proof}

\begin{lemma}[Joint cumulants of bounded variables]\label{10.7:lem-joint-cumulants}
Let $Y_1,\dots,Y_n$ be bounded random variables on a probability space, with expectation
$\mathbb{E}$, and let $z=(z_1,\dots,z_n)\in\mathbb{C}^n$. The joint cumulant is
\begin{equation*}
\kappa(Y_1,\dots,Y_n):=\partial_{z_1}\cdots\partial_{z_n}\log G(z)\big|_{z=0},
\qquad
G(z):=\mathbb{E}\Bigl[e^{\sum_{i=1}^n z_iY_i}\Bigr],
\end{equation*}
where $\log G$ is the branch near $z=0$ with $\log G(0)=0$.
\begin{enumerate}
\item[(a)] $G$ is entire and $G(0)=1$; in particular $\kappa$ is well defined.
\item[(b)] Let $\mathbb{A}_n=\mathbb{C}[z]/(z_1^2,\dots,z_n^2)$ and
$\bar G:=\mathbb{E}\bigl[\prod_{i=1}^n(1+z_iY_i)\bigr]\in\mathbb{A}_n$, and put
\begin{equation*}
\log\bar G:=\sum_{k=1}^{n}\frac{(-1)^{k-1}}{k}\,(\bar G-1)^k .
\end{equation*}
Writing $[z_1\cdots z_n]F$ for the coefficient of the monomial $z_1\cdots z_n$ in $F$,
\begin{equation}\label{10.7:eq-joint-cumulants-1}
\kappa(Y_1,\dots,Y_n)=[z_1\cdots z_n]\log\bar G
=\sum_{\pi}(-1)^{|\pi|-1}(|\pi|-1)!\prod_{B\in\pi}\mathbb{E}\Bigl[\prod_{i\in B}Y_i\Bigr],
\end{equation}
where $\pi$ ranges over the partitions of $\{1,\dots,n\}$ and $|\pi|$ is the number of
blocks of $\pi$.
\item[(c)] $\kappa$ is symmetric and multilinear. If $S$ is a measure-preserving map of the
probability space into itself, then
$\kappa(Y_1\circ S,\dots,Y_n\circ S)=\kappa(Y_1,\dots,Y_n)$. If $n\ge2$ and $c$ is a
constant, then $\kappa(c,Y_2,\dots,Y_n)=0$; consequently, for $n\ge2$, $\kappa$ is unchanged
under $Y_i\mapsto Y_i-\mathbb{E}[Y_i]$. If $Y_1,\dots,Y_n$ are real, then
$\kappa(Y_1,\dots,Y_n)$ is real.
\item[(d)] Let $Z(z)$ be the partition function with sources $w=\sum_i z_if_i$ on
$\mathbb{T}_m$ at bare spacing $L^{-K}$. Then
\begin{equation*}
\partial_{z_1}\cdots\partial_{z_n}\log Z(0)
=\kappa\bigl(\mathcal{O}(f_1),\dots,\mathcal{O}(f_n)\bigr)=S^T_{n;K,m}(f_1,\dots,f_n).
\end{equation*}
\end{enumerate}
\end{lemma}

\begin{proof}
(a) Let $C_i:=\sup|Y_i|$. For every $z$,
\begin{equation*}
\sum_{k\ge0}\frac{1}{k!}\Bigl|\sum_i z_iY_i\Bigr|^k\le e^{\sum_i|z_i|C_i}.
\end{equation*}
The exponential series therefore converges absolutely under $\mathbb{E}$. Expanding it
and taking expectations term by term gives a power series for $G$ that converges for all
$z\in\mathbb{C}^n$, so $G$ is entire. Moreover $G(0)=\mathbb{E}[1]=1$. Hence $|G-1|<1$ on a
neighbourhood of $0$, and there the series
$\log G=\sum_{k\ge1}(-1)^{k-1}(G-1)^k/k$ defines the holomorphic branch with
$\log G(0)=0$. Thus $\kappa$ is well defined.

(b) For $I\subset\{1,\dots,n\}$ write $z^I=\prod_{i\in I}z_i$ and
$m_I=\mathbb{E}[\prod_{i\in I}Y_i]$, so that $m_\emptyset=1$.

Let $\mathsf{T}\colon\mathbb{C}[[z]]\to\mathbb{A}_n$ keep the monomials all of whose
exponents are at most one. It is a ring homomorphism, and it preserves the coefficient of
$z_1\cdots z_n$. The coefficient of a monomial with all exponents at most one in a power
series $F$ at $0$ equals the corresponding mixed derivative of $F$ at $0$, since all
factorials are $1$. Hence $\kappa=[z_1\cdots z_n]\,\mathsf{T}(\log G)$.

Since $e^{\sum_iz_iY_i}=\prod_ie^{z_iY_i}$ and $\mathsf{T}(e^{z_iY_i})=1+z_iY_i$, taking
expectations coefficientwise gives
\begin{equation*}
\mathsf{T}(G)=\bar G=\sum_Iz^Im_I .
\end{equation*}

Next, $\mathsf{T}$ commutes with the logarithm series. Indeed, the series for $\log G$
converges locally uniformly near $0$, so its Taylor coefficients are the sums of those of
the terms. Applying $\mathsf{T}$ term by term gives $\sum_k(-1)^{k-1}(\bar G-1)^k/k$. The
element $\bar G-1$ lies in the ideal $\mathfrak{m}=(z_1,\dots,z_n)$ of $\mathbb{A}_n$, and
$\mathfrak{m}^{n+1}=0$, so the terms with $k>n$ vanish. Therefore
$\mathsf{T}(\log G)=\log\bar G$, which is the first equality in
\eqref{10.7:eq-joint-cumulants-1}.

For the second equality, write
\begin{equation*}
(\bar G-1)^k=\Bigl(\sum_{I\neq\emptyset}z^Im_I\Bigr)^k
=\sum_{I_1,\dots,I_k\neq\emptyset}z^{I_1}\cdots z^{I_k}\,m_{I_1}\cdots m_{I_k}.
\end{equation*}
If two of the $I_j$ overlap, the product $z^{I_1}\cdots z^{I_k}$ contains a square $z_i^2$
and vanishes in $\mathbb{A}_n$. Hence the coefficient of $z_1\cdots z_n$ collects exactly
the ordered $k$-tuples of pairwise disjoint non-empty sets with union $\{1,\dots,n\}$. Each
partition with $k$ blocks arises from $k!$ such tuples. So this coefficient is
$k!\sum_{|\pi|=k}\prod_{B\in\pi}m_B$. Multiplying by $(-1)^{k-1}/k$ and summing over $k$
gives the partition sum in \eqref{10.7:eq-joint-cumulants-1}.

(c) Symmetry. A permutation $\sigma$ of $\{1,\dots,n\}$ maps partitions bijectively onto
partitions with the same number of blocks. Moreover $\mathbb{E}[\prod_{i\in B}Y_{\sigma(i)}]=m_{\sigma(B)}$.
Hence the right side of \eqref{10.7:eq-joint-cumulants-1} is unchanged under $\sigma$.

Multilinearity. In each term of \eqref{10.7:eq-joint-cumulants-1}, each index $i$ lies in
exactly one block $B$. So $Y_i$ enters only through the factor $\mathbb{E}[\prod_{j\in B}Y_j]$,
which is linear in $Y_i$.

Invariance under measure-preserving maps. Since $\prod_{i\in B}(Y_i\circ S)=(\prod_{i\in B}Y_i)\circ S$
and $\mathbb{E}[F\circ S]=\mathbb{E}[F]$, every moment in \eqref{10.7:eq-joint-cumulants-1}
is unchanged when each $Y_i$ is replaced by $Y_i\circ S$.

Constant slot. We have
\begin{equation*}
\mathbb{E}\bigl[e^{z_1c+\sum_{i\ge2}z_iY_i}\bigr]=e^{z_1c}\,\mathbb{E}\bigl[e^{\sum_{i\ge2}z_iY_i}\bigr].
\end{equation*}
Its logarithm near $0$ is therefore $z_1c$ plus a function of $z_2,\dots,z_n$: both sides
are holomorphic, vanish at $0$ and have the same exponential. When $n\ge2$, the derivative
$\partial_{z_1}\partial_{z_2}$ annihilates both summands, so $\kappa(c,Y_2,\dots,Y_n)=0$.

Centring. For $n\ge2$, multilinearity gives
\begin{equation*}
\kappa(\dots,Y_i-\mathbb{E}[Y_i],\dots)=\kappa(\dots,Y_i,\dots)-\kappa(\dots,\mathbb{E}[Y_i],\dots).
\end{equation*}
By symmetry and the constant-slot rule, the last term vanishes.

Reality. If the $Y_i$ are real, every moment in \eqref{10.7:eq-joint-cumulants-1} is real,
and hence so is $\kappa$.

(d) Let $\mu_{K,m}$ denote the lattice measure on $\mathbb{T}_m$ at bare spacing $L^{-K}$.
By the definition of the partition function with sources,
\begin{equation*}
\frac{Z(z)}{Z(0)}=\mathbb{E}_{\mu_{K,m}}\Bigl[e^{\sum_iz_i\mathcal{O}(f_i)}\Bigr].
\end{equation*}
Each $\mathcal{O}(f_i)$ is bounded. It is a finite sum over the plaquettes of the finite
lattice, and $|1-\operatorname{Re}\operatorname{tr}U(\partial p)|\le2$ because
$\operatorname{tr}\mathbb{1}=1$ and the eigenvalues of $U(\partial p)$ have modulus one.
Parts (a) and (b) therefore apply with $Y_i=\mathcal{O}(f_i)$. Near $z=0$ we have
$\log Z(z)=\log Z(0)+\log G(z)$, so
$\partial_{z_1}\cdots\partial_{z_n}\log Z(0)=\kappa(\mathcal{O}(f_1),\dots,\mathcal{O}(f_n))$.
This is the connected Schwinger function $S^T_{n;K,m}(f_1,\dots,f_n)$.
\end{proof}

\begin{lemma}[The moment--cumulant derivative lemma]\label{10.7:lem-cumulant-derivative}
Let $n\ge 1$. On a common probability space, with expectation $\mathbb{E}$, let $X_1(t),\dots,X_n(t)$
(for $t$ in a neighbourhood of $0$), $W$ and $E_1,\dots,E_n$ be bounded random variables, and put
$X_i:=X_i(0)$. Suppose that for every $I\subset\{1,\dots,n\}$ the function
\begin{equation*}
m_I(t):=\mathbb{E}\Bigl[\prod_{i\in I}X_i(t)\Bigr]
\end{equation*}
is differentiable at $t=0$, with
\begin{equation}\label{10.7:eq-cumulant-derivative-1}
\dot m_I(0)=\mathbb{E}\Bigl[W\prod_{k\in I}X_k\Bigr]
+\sum_{i\in I}\mathbb{E}\Bigl[E_i\prod_{k\in I\setminus\{i\}}X_k\Bigr].
\end{equation}
For $I=\emptyset$ the empty product is $1$, so $m_\emptyset\equiv1$, and
\eqref{10.7:eq-cumulant-derivative-1} reads $0=\mathbb{E}[W]$. Then
$t\mapsto\kappa(X_1(t),\dots,X_n(t))$ is differentiable at $t=0$, and
\begin{equation}\label{10.7:eq-cumulant-derivative-2}
\frac{d}{dt}\Big|_{t=0}\kappa(X_1(t),\dots,X_n(t))
=\kappa(W,X_1,\dots,X_n)+\sum_{i=1}^{n}\kappa(X_1,\dots,X_{i-1},E_i,X_{i+1},\dots,X_n).
\end{equation}
\end{lemma}

\begin{proof}
We work in $\mathbb{A}_n=\mathbb{C}[z]/(z_1^2,\dots,z_n^2)$, whose elements are uniquely written as
$\sum_I a_I z^I$ with $z^I:=\prod_{i\in I}z_i$, $I\subset\{1,\dots,n\}$. We write $[z^I]a:=a_I$ for the
coefficient of $z^I$ in $a$, and $\mathfrak{m}$ for the maximal ideal of elements with vanishing
constant term. Every product of $n+1$ elements of $\mathfrak{m}$ is zero, so
$\mathfrak{m}^{n+1}=0$. On $1+\mathfrak{m}$ we use the logarithm of the joint cumulants lemma
(Lemma~\ref{10.7:lem-joint-cumulants}(b)), $\log M=\sum_{k=1}^{n}(-1)^{k-1}(M-1)^k/k$, a polynomial
map; adding terms with $k>n$ does not change it, since $(M-1)^k=0$ for $k\ge n+1$. An element
$M\in1+\mathfrak{m}$ is invertible with $M^{-1}=\sum_{k=0}^{n}(1-M)^k$, because
$M\sum_{k=0}^{n}(1-M)^k=1-(1-M)^{n+1}=1$.

Put $M(t):=\sum_I z^I m_I(t)=\mathbb{E}[\prod_{k}(1+z_kX_k(t))]$. Each coefficient $m_I$ is
differentiable at $0$, so $M$ is differentiable at $0$ as a map into the finite-dimensional space
$\mathbb{A}_n$; since $m_\emptyset\equiv1$, we have $M(t)\in1+\mathfrak{m}$ and
$\dot M:=\dot M(0)\in\mathfrak{m}$. By Lemma~\ref{10.7:lem-joint-cumulants}(b),
$\kappa(X_1(t),\dots,X_n(t))=[z_1\cdots z_n]\log M(t)$, a polynomial in the $m_I(t)$; hence it is
differentiable at $0$, and by the chain rule its derivative is the coefficient of $z_1\cdots z_n$ in
the directional derivative of $\log$ at $M(0)$ in the direction $\dot M$. For $M\in1+\mathfrak{m}$
and $H\in\mathfrak{m}$, differentiating the polynomial term by term in the commutative ring
$\mathbb{A}_n$ gives
\begin{align*}
\frac{d}{ds}\Big|_{s=0}\log(M+sH)
&=\sum_{k=1}^{n}(-1)^{k-1}H(M-1)^{k-1}
=H\sum_{k=0}^{n-1}(1-M)^k\\
&=HM^{-1}-H(1-M)^n=HM^{-1},
\end{align*}
where the third equality uses $M^{-1}=\sum_{k=0}^{n}(1-M)^k$, and the last one uses
$H(1-M)^n\in\mathfrak{m}^{n+1}=0$. Therefore
\begin{equation}\label{10.7:eq-cumulant-derivative-3}
\frac{d}{dt}\Big|_{t=0}\kappa(X_1(t),\dots,X_n(t))=[z_1\cdots z_n]\bigl(M(0)^{-1}\dot M\bigr).
\end{equation}

Now put
\begin{equation*}
P:=\mathbb{E}\Bigl[W\prod_{k=1}^{n}(1+z_kX_k)\Bigr],\qquad
Q_i:=\mathbb{E}\Bigl[E_i\prod_{k\ne i}(1+z_kX_k)\Bigr]\quad(1\le i\le n).
\end{equation*}
For every $I$, $[z^I]P=\mathbb{E}[W\prod_{k\in I}X_k]$. Moreover
$[z^I](z_iQ_i)=\mathbb{E}[E_i\prod_{k\in I\setminus\{i\}}X_k]$ if $i\in I$, and $=0$ if $i\notin I$.
Hence, for $I\ne\emptyset$, the hypothesis \eqref{10.7:eq-cumulant-derivative-1} says precisely that
$\dot M$ and $P+\sum_i z_iQ_i$ have the same coefficient of $z^I$. Their constant terms are
$\dot m_\emptyset(0)=0$ and $[z^\emptyset]P=\mathbb{E}[W]$ respectively, since $z_iQ_i\in\mathfrak{m}$. Thus
\begin{equation}\label{10.7:eq-cumulant-derivative-4}
\dot M=P+\sum_{i=1}^{n}z_iQ_i\quad\text{in }\mathbb{A}_n
\qquad\text{if and only if}\qquad \mathbb{E}[W]=0.
\end{equation}
This is exactly the place where the case $I=\emptyset$ of the hypothesis, $0=\mathbb{E}[W]$, is used.
Without it one would have $\dot M=P+\sum_iz_iQ_i-\mathbb{E}[W]$, and
\eqref{10.7:eq-cumulant-derivative-3} would acquire the additional term
$-\mathbb{E}[W]\,[z_1\cdots z_n]M(0)^{-1}$, which does not vanish in general. Inserting
\eqref{10.7:eq-cumulant-derivative-4} into \eqref{10.7:eq-cumulant-derivative-3} and using linearity of
$[z_1\cdots z_n]$, the derivative equals
$[z_1\cdots z_n](M(0)^{-1}P)+\sum_i[z_1\cdots z_n](M(0)^{-1}z_iQ_i)$. We identify the two kinds of
terms in turn.

First term. Adjoin a further variable $w$ with $w^2=0$. The algebra
$\mathbb{A}_n[w]/(w^2)$ is the algebra $\mathbb{A}_{n+1}$ in the variables $w,z_1,\dots,z_n$, and
$M(0)+wP$ lies in its $1+\mathfrak{m}$. Expanding each power and using $w^2=0$,
$(M(0)-1+wP)^k=(M(0)-1)^k+kwP(M(0)-1)^{k-1}$. Summing, with $\log$ taken to order $n+1$ as in
Lemma~\ref{10.7:lem-joint-cumulants}(b) for $n+1$ variables, gives
\begin{equation*}
\log\bigl(M(0)+wP\bigr)=\log M(0)+wP\sum_{k=0}^{n}(1-M(0))^k=\log M(0)+wM(0)^{-1}P.
\end{equation*}
On the other hand, expanding $(1+wW)$ gives
$M(0)+wP=\mathbb{E}[(1+wW)\prod_k(1+z_kX_k)]$. Since $\log M(0)$ contains no $w$, we obtain
\begin{equation*}
[z_1\cdots z_n]\bigl(M(0)^{-1}P\bigr)
=[wz_1\cdots z_n]\log\mathbb{E}\Bigl[(1+wW)\prod_{k=1}^{n}(1+z_kX_k)\Bigr]
=\kappa(W,X_1,\dots,X_n).
\end{equation*}
The last equality is Lemma~\ref{10.7:lem-joint-cumulants}(b) applied to the $n+1$ bounded variables
$W,X_1,\dots,X_n$.

The $i$-th term. Let $M_i:=M(0)|_{z_i=0}=\mathbb{E}[\prod_{k\ne i}(1+z_kX_k)]$. Since $z_i^2=0$, one
has $z_ia=z_i\,(a|_{z_i=0})$ for every $a\in\mathbb{A}_n$. Setting $z_i=0$ is a ring homomorphism,
so it carries $M(0)^{-1}$ to $M_i^{-1}$. Hence $z_iM(0)^{-1}Q_i=z_iM_i^{-1}Q_i$. Both $M_i$ and
$Q_i$ lie in the subalgebra generated by the $z_k$, $k\ne i$, which is the algebra
$\mathbb{A}_{n-1}$ in these variables, and therefore
\begin{equation*}
[z_1\cdots z_n]\bigl(M(0)^{-1}z_iQ_i\bigr)=\Bigl[\prod_{k\ne i}z_k\Bigr]\bigl(M_i^{-1}Q_i\bigr).
\end{equation*}
The same computation as for the first term, with $W$ replaced by $E_i$ and the family
$(X_k)_{k\ne i}$ in place of $(X_k)_{k}$, now applies. Adjoining $w$ with $w^2=0$, we have
$M_i+wQ_i=\mathbb{E}[(1+wE_i)\prod_{k\ne i}(1+z_kX_k)]$ and
$\log(M_i+wQ_i)=\log M_i+wM_i^{-1}Q_i$. Lemma~\ref{10.7:lem-joint-cumulants}(b), applied to the $n$
bounded variables $E_i,(X_k)_{k\ne i}$, then shows that this coefficient equals
$\kappa(E_i,(X_k)_{k\ne i})$. By the symmetry of $\kappa$ (Lemma~\ref{10.7:lem-joint-cumulants}(c)),
this is $\kappa(X_1,\dots,X_{i-1},E_i,X_{i+1},\dots,X_n)$.

Adding the first term and the $n$ terms just identified gives
\eqref{10.7:eq-cumulant-derivative-2}.
\end{proof}

\begin{corollary}[From the product form to the connected form of the rotational Ward identity]
\label{10.7:cor-connected-form}
Let $(K,m)$ be a cutoff with $L^m\ge\mathcal{R}/2$, let
$\vec h=(h_1,\dots,h_n)\in\mathcal{T}^{(r)}_n(\mathcal{R},\mathsf{d})$ with $r\ge1$, and write
$\langle\cdot\rangle$ for the expectation of the lattice measure at cutoff $(K,m)$ on
$\mathbb{T}_m$. Let $W$ and $E_{h_1},\dots,E_{h_n}$ be bounded Borel functions of the bare
field. For $I\subset\{1,\dots,n\}$ put
\begin{equation*}
F_I:=\prod_{i\in I}\mathcal{O}\bigl(h_i^{(m)}\bigr),\qquad
E_{F_I}:=\sum_{i\in I}E_{h_i}\prod_{k\in I\setminus i}\mathcal{O}\bigl(h_k^{(m)}\bigr),
\end{equation*}
with $F_\emptyset:=1$ and $E_{F_\emptyset}:=0$. Suppose that the rotational Ward identity
holds in the product form, that is, for every $I\subset\{1,\dots,n\}$,
\begin{equation}\label{10.7:eq-connected-form-1}
\frac{d}{dt}\Big|_{t=0}\bigl\langle F_I\circ R^{\omega}_t\bigr\rangle
=\langle F_I W\rangle+\langle E_{F_I}\rangle ,
\end{equation}
where, as in the product form of the rotational Ward identity, $F_I\circ R^{\omega}_t$ denotes
the observable $F_I$ read on the test functions rotated back,
\begin{equation*}
F_I\circ R^{\omega}_t:=\prod_{i\in I}\mathcal{O}\bigl((R^{-\omega}_t\cdot h_i)^{(m)}\bigr).
\end{equation*}
Then
\begin{equation}\label{10.7:eq-connected-form-2}
\begin{split}
D^{-\omega}_{K,m}(\vec h)
&=\kappa\bigl(W,\mathcal{O}(h_1^{(m)}),\dots,\mathcal{O}(h_n^{(m)})\bigr)\\
&\quad+\sum_{i=1}^n\kappa\bigl(\mathcal{O}(h_1^{(m)}),\dots,E_{h_i},\dots,
\mathcal{O}(h_n^{(m)})\bigr);
\end{split}
\end{equation}
that is, the identity \eqref{10.7:eq-hypotheses-used-form-a} of
Definition~\ref{10.7:def-hypotheses-used-form} holds at $-\omega$ at the cutoff $(K,m)$, for
every $\chi$ admissible at radius $\mathcal{R}$, with
$N^{-\omega,\chi}_{K,m}+A^{-\omega,\chi}_{K,m}:=W$ for any splitting of $W$ into bounded
Borel functions $N^{-\omega,\chi}_{K,m},A^{-\omega,\chi}_{K,m}$, and with
$E^{-\omega}_{h_i;K,m}:=E_{h_i}$.
\end{corollary}

\begin{proof}
For $t\in\mathbb{R}$ and $1\le i\le n$ put
$X_i(t):=\mathcal{O}\bigl((R^{-\omega}_t\cdot h_i)^{(m)}\bigr)$, so that
$X_i(0)=\mathcal{O}(h_i^{(m)})$, and $m_I(t):=\langle\prod_{i\in I}X_i(t)\rangle$; by the
definition of $F_I\circ R^{\omega}_t$ in the statement, the function
$\langle F_I\circ R^{\omega}_t\rangle$ of \eqref{10.7:eq-connected-form-1} is
$m_I(t)$. Each $X_i(t)$ is bounded, being a finite linear combination, with coefficients
$(R^{-\omega}_t\cdot h_i)^{(m)}(x(p))$, of the variables
$1-\operatorname{Re}\operatorname{tr}U(\partial p)\in[0,2]$ over the finitely many plaquettes
$p$ of the lattice on $\mathbb{T}_m$. The functions $W$, $E_{h_i}$ are
bounded by hypothesis. The computation of
Lemma~\ref{10.7:lem-finite-cutoff-structure}(c) (which applies, since
$h_i\in C^r_c\subset C^1_c$, $\operatorname{supp}h_i\subset B(0,\mathcal{R}/2)$ and
$L^m\ge\mathcal{R}/2$), carried out with the expectation in place of
the cumulant $\kappa$, shows that $m_I$ is $C^1$ on $\mathbb{R}$; in particular it is
differentiable at $0$. By \eqref{10.7:eq-connected-form-1},
\begin{equation*}
\dot m_I(0)=\Bigl\langle W\prod_{i\in I}X_i\Bigr\rangle
+\sum_{i\in I}\Bigl\langle E_{h_i}\prod_{k\in I\setminus i}X_k\Bigr\rangle,
\qquad X_i:=X_i(0),
\end{equation*}
and for $I=\emptyset$, where $m_\emptyset\equiv1$ and $E_{F_\emptyset}=0$, this reads
$0=\langle W\rangle$. These are the hypotheses of the moment--cumulant derivative lemma
(Lemma~\ref{10.7:lem-cumulant-derivative}) with $E_i:=E_{h_i}$. Since
$S^T_{n;K,m}(R^{-\omega}_t\cdot\vec h)=\kappa(X_1(t),\dots,X_n(t))$, that lemma gives
\begin{equation*}
\begin{split}
D^{-\omega}_{K,m}(\vec h)
&=\frac{d}{dt}\Big|_{t=0}\kappa\bigl(X_1(t),\dots,X_n(t)\bigr)\\
&=\kappa(W,X_1,\dots,X_n)+\sum_{i=1}^n\kappa(X_1,\dots,E_{h_i},\dots,X_n),
\end{split}
\end{equation*}
which is \eqref{10.7:eq-connected-form-2}. For any splitting $W=N+A$, multilinearity of
$\kappa$ in its first entry gives $\kappa(W,\dots)=\kappa(N,\dots)+\kappa(A,\dots)$, and
\eqref{10.7:eq-connected-form-2} becomes \eqref{10.7:eq-hypotheses-used-form-a} at $-\omega$,
for every admissible $\chi$, with $N^{-\omega,\chi}_{K,m}:=N$, $A^{-\omega,\chi}_{K,m}:=A$
and $E^{-\omega}_{h_i;K,m}:=E_{h_i}$.

The case $I=\emptyset$, namely $0=\langle W\rangle$, is, for the Schwinger--Dyson functional
$W=g_0^{-2}\partial_X A(U)-\operatorname{div}X$, the lattice Schwinger--Dyson identity with
test functional identically equal to $1$; here $X$ is the vector field on the space of bond
configurations entering that identity (not to be confused with the variables $X_i(t)$ above),
$\partial_X$ the derivative along it, $\operatorname{div}X$ its divergence, and $A(U)$ the
Wilson action. In the connected form
\eqref{10.7:eq-connected-form-2} itself additive constants in $W$ are invisible: a joint
cumulant of at least two bounded variables, one of which is constant, vanishes, by
multilinearity and the vanishing of a joint cumulant with an independent, in particular a
constant, entry.
\end{proof}

\begin{lemma}[The periodic lift preserves norms and separations]\label{10.7:lem-periodic-lift}
Let $\pi\colon\mathbb{R}^4\to\mathbb{T}_m=\mathbb{R}^4/2L^m\mathbb{Z}^4$ be the quotient map. Let $d$ be the
torus distance, $d(\pi(x),\pi(x'))=\min_{v\in 2L^m\mathbb{Z}^4}|x-x'-v|$, and let
$\operatorname{Lip}_d$ be the Lipschitz constant with respect to $d$. For a function $h$ on
$\mathbb{R}^4$ with compact support, its periodic lift is the function $h^{(m)}$ on
$\mathbb{T}_m$ given by
\begin{equation}\label{10.7:eq-periodic-lift-1}
h^{(m)}(\pi(x))=\sum_{v\in 2L^m\mathbb{Z}^4}h(x+v).
\end{equation}
Let $h$ be Lipschitz with $\operatorname{supp}h\subset\bar B(0,\rho)$, where $L^m\ge 2\rho>0$.
\begin{enumerate}
\item[(a)] For each $\xi\in\mathbb{T}_m$, at most one term of the sum defining $h^{(m)}(\xi)$ is
non-zero.
\item[(b)] $\|h^{(m)}\|_\infty=\|h\|_\infty$ and $\operatorname{Lip}_d(h^{(m)})\le\operatorname{Lip}h$.
Hence $\|h^{(m)}\|_{\sharp}\le\|h\|_{\sharp}$.
\item[(c)] $\operatorname{supp}h^{(m)}=\pi(\operatorname{supp}h)$, and this set lies in the torus
ball of radius $\rho$ about $\pi(0)$. If $h'$ satisfies the same hypotheses, then the torus
distance between the supports of $h^{(m)}$ and $h'^{(m)}$ equals the Euclidean distance between
the supports of $h$ and $h'$.
\item[(d)] For $w\in W_4$ set $\sigma_w(\pi(x)):=\pi(wx)$ and $(w\cdot h)(x):=h(w^{-1}x)$. Then
$(w\cdot h)^{(m)}=h^{(m)}\circ\sigma_w^{-1}$. The same holds for every affine isometry whose
linear part preserves $2L^m\mathbb{Z}^4$. For $R\in O(4)$, the function $R\cdot h$ satisfies
(a)--(c).
\end{enumerate}
\end{lemma}

\begin{proof}
Throughout, $v$ and $v'$ range over $2L^m\mathbb{Z}^4$.

\emph{Preliminaries.} We have $h=0$ on $\{|y|>\rho\}$ by the support hypothesis, and $h=0$ on
$\{|y|=\rho\}$ by continuity. Hence $h=0$ on $\{|y|\ge\rho\}$. In the same way, $h$ vanishes on
the topological boundary of $\operatorname{supp}h$. Indeed, every such boundary point is a limit
of points outside the support, and $h$ is zero at those points.

(a) Suppose the terms $h(x+v)$ and $h(x+v')$ were both non-zero, with $v\neq v'$. Then
$|x+v|<\rho$ and $|x+v'|<\rho$. This gives
\begin{equation*}
|v-v'|<2\rho\le L^m<2L^m .
\end{equation*}
But distinct points of $2L^m\mathbb{Z}^4$ are at distance at least $2L^m$, a contradiction.

(b) By (a), every value of $h^{(m)}$ is either a value of $h$ or $0$. Conversely, if $h(x)\neq0$
then the term with $v=0$ is non-zero, so by (a) we get $h^{(m)}(\pi(x))=h(x)$. Hence
$\|h^{(m)}\|_\infty=\|h\|_\infty$.

For the Lipschitz bound, put $H:=\sum_v h(\cdot+v)$ on $\mathbb{R}^4$, so that
$H=h^{(m)}\circ\pi$. Put $S_v:=\operatorname{supp}h-v$. These sets are closed.

\emph{Separation of the $S_v$.} Take points $y-v\in S_v$ and $y'-v'\in S_{v'}$ with $v\neq v'$.
Then
\begin{equation*}
|(y-v)-(y'-v')|\ge|v-v'|-|y|-|y'|\ge 2L^m-2\rho>0 .
\end{equation*}
So the $S_v$ are pairwise at distance at least $2L^m-2\rho>0$.

\emph{Form of $H$.} Let $x\in S_v$ and $v'\ne v$. Then
\begin{equation*}
|x+v'|\ge|v'-v|-|x+v|\ge 2L^m-\rho>\rho ,
\end{equation*}
so $h(x+v')=0$. Hence $H=h(\cdot+v)$ on $S_v$, and $H=0$ off $\bigcup_v S_v$. By the
preliminaries, each $h(\cdot+v)$ vanishes on $\partial S_v$.

\emph{A consequence.} Let $H(y)=0$. Then $h(y+v)=0$ for every $v$. If $y\in S_v$, this holds
because $h(y+v)=H(y)$. Otherwise $y+v\notin\operatorname{supp}h$.

\emph{The cases.} Take $x,y\in\mathbb{R}^4$.
\begin{itemize}
\item Let $H(x)=0$ or $H(y)=0$. By symmetry we may assume $H(y)=0$. If also $H(x)=0$ there is
nothing to prove. Otherwise $x\in S_v$ for some $v$, and by the consequence above
$|H(x)-H(y)|=|h(x+v)-h(y+v)|\le\operatorname{Lip}h\,|x-y|$.
\item Let $x,y\in S_v$. Then $|H(x)-H(y)|=|h(x+v)-h(y+v)|\le\operatorname{Lip}h\,|x-y|$.
\item Let $x\in S_v$ and $y\in S_{v'}$ with $v\neq v'$. Put $z(s):=x+s(y-x)$ for $s\in[0,1]$.
Let $s_1:=\sup\{s\in[0,1]:z(s)\in S_v\}$. Since $S_v$ is closed, $z(s_1)\in S_v$. Since
$y\notin S_v$, we have $s_1<1$, and $z(s)\notin S_v$ for $s>s_1$. So $z(s_1)\in\partial S_v$.
Let $s_2:=\inf\{s\in[s_1,1]:z(s)\in S_{v'}\}$; the set is non-empty because it contains $1$.
Then $z(s_2)\in S_{v'}$. Moreover $s_2>s_1$, because $z(s_1)\in S_v$ and $S_v\cap S_{v'}=\emptyset$.
The points $z(s)$ with $s_1\le s<s_2$ lie outside $S_{v'}$, so $z(s_2)\in\partial S_{v'}$.
Consequently $h(z(s_1)+v)=0=h(z(s_2)+v')$, and therefore
\begin{align*}
|H(x)-H(y)|&\le|h(x+v)-h(z(s_1)+v)|+|h(z(s_2)+v')-h(y+v')|\\
&\le\operatorname{Lip}h\,(s_1+1-s_2)\,|x-y|\le\operatorname{Lip}h\,|x-y| .
\end{align*}
\end{itemize}
So $\operatorname{Lip}H\le\operatorname{Lip}h$.

Now let $\xi=\pi(x)$ and $\xi'=\pi(x')$. The function $H$ is $2L^m\mathbb{Z}^4$-periodic, so for
every $v$
\begin{equation*}
|h^{(m)}(\xi)-h^{(m)}(\xi')|=|H(x)-H(x'+v)|\le\operatorname{Lip}h\,|x-x'-v| .
\end{equation*}
Minimising over $v$ gives $\operatorname{Lip}_d(h^{(m)})\le\operatorname{Lip}h$. Since
$\|f\|_{\sharp}=\max\{\|f\|_\infty,40M\operatorname{Lip}f\}$, it follows that
$\|h^{(m)}\|_{\sharp}\le\|h\|_{\sharp}$.

(c) \emph{The support.} By (a) and the form of $H$,
\begin{equation*}
\{h^{(m)}\neq0\}=\pi(\{h\neq0\}).
\end{equation*}
The set $\pi(\operatorname{supp}h)$ is compact, hence closed, and it contains this set. Conversely,
$\pi(\operatorname{supp}h)$ lies in the closure of $\pi(\{h\neq0\})$, because $\pi$ is continuous.
Hence $\operatorname{supp}h^{(m)}=\pi(\operatorname{supp}h)$. For $|y|\le\rho$ we have
$d(\pi(y),\pi(0))\le|y|\le\rho$.

\emph{Separations.} Let $S=\operatorname{supp}h$ and $S'=\operatorname{supp}h'$, both contained in
$\bar B(0,\rho)$. Then
\begin{equation*}
\operatorname{dist}_d(\pi(S),\pi(S'))=\inf_{y\in S,\,y'\in S',\,v}|y-y'+v| .
\end{equation*}
The terms with $v=0$ give the Euclidean distance $\operatorname{dist}(S,S')$, which is at most
$2\rho$. The terms with $v\neq0$ satisfy
\begin{equation*}
|y-y'+v|\ge|v|-|y|-|y'|\ge 2L^m-2\rho\ge2\rho .
\end{equation*}
Therefore the infimum equals $\operatorname{dist}(S,S')$.

(d) \emph{Signed permutations.} The matrix $w^{-1}$ maps $2L^m\mathbb{Z}^4$ bijectively onto
itself, so $\sigma_w$ is well defined. For $x\in\mathbb{R}^4$,
\begin{equation*}
(w\cdot h)^{(m)}(\pi(x))=\sum_v h(w^{-1}x+w^{-1}v)=\sum_{v'}h(w^{-1}x+v')
=h^{(m)}(\sigma_w^{-1}(\pi(x))).
\end{equation*}

\emph{Affine isometries.} Let $a(x)=Ax+b$ with $A$ orthogonal and $A(2L^m\mathbb{Z}^4)=2L^m\mathbb{Z}^4$.
Set $\sigma_a(\pi(x)):=\pi(ax)$ and $(a\cdot h)(x):=h(a^{-1}x)$. Then
$a^{-1}(x+v)=a^{-1}x+A^{-1}v$, and the same substitution $v'=A^{-1}v$ gives
$(a\cdot h)^{(m)}=h^{(m)}\circ\sigma_a^{-1}$.

\emph{Rotations.} For $R\in O(4)$, the function $R\cdot h$ is Lipschitz with
$\operatorname{supp}(R\cdot h)=R\operatorname{supp}h\subset\bar B(0,\rho)$. So it satisfies the
hypotheses of the lemma, and (a)--(c) apply to it.
\end{proof}

\begin{lemma}[Finite-cutoff correlation functions and their rotation derivative]
\label{10.7:lem-finite-cutoff-structure}
Fix $K$ and $m$, and consider the lattice of spacing $L^{-K}$ on $\mathbb{T}_m$. For a plaquette $p$ write
$a_p := 1-\operatorname{Re}\operatorname{tr} U(\partial p)$, so that
$\mathcal{O}(\varphi)=\sum_p \varphi(x(p))\,a_p$. Let $\pi\colon\mathbb{R}^4\to\mathbb{T}_m$ be the quotient map.
Let $\tilde x(p)\in\mathbb{R}^4$ be the representative of $x(p)$ with $|\tilde x(p)|_\infty\le L^m$.
For a function $h$ on $\mathbb{R}^4$ let
\begin{equation*}
h^{(m)}(\pi(x))=\sum_{v\in 2L^m\mathbb{Z}^4}h(x+v)
\end{equation*}
be its periodic lift. For $R\in O(4)$ put $R\cdot h:=h\circ R^{-1}$. Write $\kappa(a_{\vec p}):=\kappa(a_{p_1},\dots,a_{p_n})$.
\begin{enumerate}
\item[(a)] For functions $\varphi_1,\dots,\varphi_n$ on $\mathbb{T}_m$,
\begin{equation}\label{10.7:eq-finite-cutoff-structure-1}
\kappa\bigl(\mathcal{O}(\varphi_1),\dots,\mathcal{O}(\varphi_n)\bigr)
=\sum_{\vec p}\prod_{i=1}^n\varphi_i(x(p_i))\,\kappa(a_{p_1},\dots,a_{p_n}),
\end{equation}
and this is a finite sum. Consequently $S^T_{n;K,m}$ is multilinear.
\item[(b)] Let $\operatorname{supp}h\subset B(0,\mathcal{R}/2)$ and $L^m\ge\mathcal{R}/2$. Then for every $R\in O(4)$ and every plaquette $p$,
\begin{equation*}
(R\cdot h)^{(m)}(x(p))=h(R^{-1}\tilde x(p)).
\end{equation*}
\item[(c)] Let $\omega$ be a real antisymmetric $4\times4$ matrix. Let $R^{\omega}_t=e^{t\omega}$ and $\delta_{\omega}f=-(\omega x)\cdot\nabla f$.
Let $h_1,\dots,h_n\in C^1_c(\mathbb{R}^4)$ with $\operatorname{supp}h_i\subset B(0,\mathcal{R}/2)$, and let $L^m\ge\mathcal{R}/2$. Then the function
\begin{equation*}
\varphi_{K,m}(t):=S^T_{n;K,m}(R^{\omega}_t\cdot h_1,\dots,R^{\omega}_t\cdot h_n)
\end{equation*}
is $C^1$ on $\mathbb{R}$, and its derivative is
\begin{equation}\label{10.7:eq-finite-cutoff-structure-2}
\varphi_{K,m}'(t)=\sum_{i=1}^n S^T_{n;K,m}\bigl(R^{\omega}_t\cdot h_1,\dots,
R^{\omega}_t\cdot\delta_{\omega}h_i,\dots,R^{\omega}_t\cdot h_n\bigr).
\end{equation}
Here $R^{\omega}_t\cdot\delta_{\omega}h_i=\delta_{\omega}(R^{\omega}_t\cdot h_i)$. In particular the derivative along the orbit at $t=0$ is
\begin{equation*}
D^{\omega}_{K,m}(\vec h):=\varphi_{K,m}'(0)
=\sum_{i=1}^nS^T_{n;K,m}(h_1,\dots,\delta_{\omega}h_i,\dots,h_n).
\end{equation*}
\end{enumerate}
\end{lemma}

\begin{proof}
(a) Since $\operatorname{tr}\mathbb{1}=1$, we have $|\operatorname{Re}\operatorname{tr}U(\partial p)|\le1$, so $0\le a_p\le2$. The $a_p$ are therefore bounded random variables. The lattice on $\mathbb{T}_m$ has finitely many plaquettes, so each $\mathcal{O}(\varphi_i)=\sum_p\varphi_i(x(p))a_p$ is a finite linear combination of the $a_p$.

The joint cumulant is multilinear (Lemma~\ref{10.7:lem-joint-cumulants}). Expanding each argument therefore gives \eqref{10.7:eq-finite-cutoff-structure-1}, as a sum over the finitely many tuples $\vec p=(p_1,\dots,p_n)$.

The maps $h\mapsto h^{(m)}$ and $\varphi\mapsto\mathcal{O}(\varphi)$ are linear. The right side of \eqref{10.7:eq-finite-cutoff-structure-1} is linear in each $\varphi_i$ separately. Hence $S^T_{n;K,m}$, the joint cumulant of the observables $\mathcal{O}(h_i^{(m)})$, is multilinear.

(b) Since $R$ is orthogonal, $\operatorname{supp}(R\cdot h)=R\operatorname{supp}h\subset B(0,\mathcal{R}/2)$. By definition of the lift,
\begin{equation*}
(R\cdot h)^{(m)}(x(p))=\sum_{v\in2L^m\mathbb{Z}^4}(R\cdot h)(\tilde x(p)+v).
\end{equation*}
Let $v\ne0$. Some coordinate of $v$ has modulus at least $2L^m$, while every coordinate of $\tilde x(p)$ has modulus at most $L^m$. Therefore
\begin{equation*}
|\tilde x(p)+v|\ge|\tilde x(p)+v|_\infty\ge 2L^m-L^m=L^m\ge\mathcal{R}/2 .
\end{equation*}
So $\tilde x(p)+v\notin B(0,\mathcal{R}/2)$, and the term indexed by $v$ vanishes. Only the term $v=0$ remains, and it equals $(R\cdot h)(\tilde x(p))=h(R^{-1}\tilde x(p))$.

(c) Since $\omega$ is antisymmetric, $e^{t\omega}$ is orthogonal for every $t$. Hence (b) applies with $R=R^{\omega}_t$, and $(R^{\omega}_t)^{-1}=e^{-t\omega}$. By the definition of $S^T_{n;K,m}$, then (a), then (b),
\begin{equation}\label{10.7:eq-finite-cutoff-structure-3}
\varphi_{K,m}(t)=\sum_{\vec p}\kappa(a_{\vec p})\prod_{i=1}^nh_i\bigl(e^{-t\omega}\tilde x(p_i)\bigr),
\end{equation}
a finite sum.

Fix $y\in\mathbb{R}^4$ and put $z=e^{-t\omega}y$. By the chain rule,
\begin{equation*}
\frac{d}{dt}h_i(e^{-t\omega}y)=\nabla h_i(z)\cdot(-\omega z)=(\delta_{\omega}h_i)(z).
\end{equation*}
This derivative is continuous in $t$, because $h_i\in C^1$. By the product rule, \eqref{10.7:eq-finite-cutoff-structure-3} is therefore $C^1$, with
\begin{equation*}
\varphi_{K,m}'(t)=\sum_{i=1}^n\sum_{\vec p}\kappa(a_{\vec p})\,(\delta_{\omega}h_i)(e^{-t\omega}\tilde x(p_i))
\prod_{j\ne i}h_j(e^{-t\omega}\tilde x(p_j)).
\end{equation*}

Next, $\operatorname{supp}\delta_{\omega}h_i\subset\operatorname{supp}h_i\subset B(0,\mathcal{R}/2)$, because $\delta_\omega h_i$ vanishes wherever $\nabla h_i$ does. Applying (a) and (b) to the tuple in which $h_i$ is replaced by $\delta_{\omega}h_i$ identifies the inner sum with the $i$-th term of \eqref{10.7:eq-finite-cutoff-structure-2}.

It remains to check $R^{\omega}_t\cdot\delta_{\omega}h_i=\delta_{\omega}(R^{\omega}_t\cdot h_i)$. We use $\nabla(h_i\circ e^{-t\omega})(y)=e^{t\omega}\nabla h_i(z)$, which holds because $(e^{-t\omega})^{\mathsf T}=e^{t\omega}$. We also use that $\omega$ commutes with $e^{-t\omega}$. Then
\begin{equation*}
\begin{split}
\delta_{\omega}(R^{\omega}_t\cdot h_i)(y)&=-(\omega y)\cdot e^{t\omega}\nabla h_i(z)
=-(e^{-t\omega}\omega y)\cdot\nabla h_i(z)\\
&=-(\omega z)\cdot\nabla h_i(z)=(R^{\omega}_t\cdot\delta_{\omega}h_i)(y).
\end{split}
\end{equation*}

Finally, $R^{\omega}_0$ is the identity, so setting $t=0$ in \eqref{10.7:eq-finite-cutoff-structure-2} gives the formula for $D^{\omega}_{K,m}(\vec h)$.
\end{proof}

\begin{lemma}[Exact hyperoctahedral invariance at finite cutoff]
\label{10.7:lem-finite-cutoff-invariance}
For every $w\in W_4$ and every tuple $\vec h=(h_1,\dots,h_n)$ at which $S^T_{n;K,m}$ is defined,
\begin{equation*}
S^T_{n;K,m}(w\cdot\vec h)=S^T_{n;K,m}(\vec h),
\qquad w\cdot\vec h=(w\cdot h_1,\dots,w\cdot h_n).
\end{equation*}
\end{lemma}

\begin{proof}
Since $w$ is a signed permutation matrix, it maps the half-offset lattice
$\mathbb{L}_K=L^{-K}(\mathbb{Z}+\tfrac12)^4$ onto itself and maps $2L^m\mathbb{Z}^4$ onto
itself. Hence $\sigma_w(\pi(x))=\pi(wx)$ is a well-defined bijection of $\mathbb{T}_m$. It
permutes the sites and bonds of the lattice on $\mathbb{T}_m$, and therefore also its
plaquettes. Because $w$ is linear, $x(\sigma_w q)=\sigma_w x(q)$ for every plaquette $q$.

Write $a_p(U)=1-\operatorname{Re}\operatorname{tr}U(\partial p)$. Let $\sigma_w^{-1}U$ denote
the configuration $(\sigma_w^{-1}U)_b=U_{\sigma_w b}$, where the variable of a bond traversed
against its orientation is the inverse. The holonomy of $\sigma_w^{-1}U$ around $\partial q$
is the holonomy of $U$ around $\partial(\sigma_w q)$, up to conjugation and inversion. Since
$\operatorname{Re}\operatorname{tr}$ is invariant under both operations,
$a_q(\sigma_w^{-1}U)=a_{\sigma_w q}(U)$.

By Lemma~\ref{10.7:lem-periodic-lift}(d), $(w\cdot h)^{(m)}=h^{(m)}\circ\sigma_w^{-1}$. We
substitute $p=\sigma_w q$; as $q$ runs over all plaquettes, so does $p$. This gives
\begin{equation*}
\begin{split}
\mathcal{O}\bigl((w\cdot h)^{(m)}\bigr)(U)
&=\sum_p h^{(m)}\bigl(\sigma_w^{-1}x(p)\bigr)\,a_p(U)
=\sum_q h^{(m)}\bigl(x(q)\bigr)\,a_{\sigma_w q}(U)\\
&=\mathcal{O}\bigl(h^{(m)}\bigr)(\sigma_w^{-1}U).
\end{split}
\end{equation*}
Let $Y_i=\mathcal{O}(h_i^{(m)})$. These are bounded, being finite sums of functions bounded
by $2\|h_i^{(m)}\|_\infty$. The display shows that
$\mathcal{O}((w\cdot h_i)^{(m)})=Y_i\circ S$ with $S(U)=\sigma_w^{-1}U$.

By the lemma on the invariances of the Wilson measure, $S$ preserves the
Wilson measure on $\mathbb{T}_m$ at spacing $L^{-K}$. Now $S^T_{n;K,m}(\vec h)$ is the joint
cumulant $\kappa(Y_1,\dots,Y_n)$ under this measure, and likewise
$S^T_{n;K,m}(w\cdot\vec h)=\kappa(Y_1\circ S,\dots,Y_n\circ S)$. Lemma~\ref{10.7:lem-joint-cumulants}(c)
gives $\kappa(Y\circ S)=\kappa(Y)$ for measure-preserving $S$, and the claim follows.
\end{proof}

\begin{lemma}[The equicontinuity estimate]\label{10.7:lem-equicontinuity}
Let $n\ge 2$ and $\rho>0$. Along the diagonal, $j\mapsto(g_0^{(j)},K_j,m_j)$ with $g_0^{(j)}\le g_-(g)$
and $K_j=K(g_0^{(j)},g)$, write $S_j:=S^T_{n;K_j,m_j}$.
\begin{enumerate}
\item[(a)] For every admissible tuple $\vec f$ of $C^1_c$ functions the limit
$S^T_n(\vec f):=\lim_{j}S_j(\vec f)$ exists. The functional $S^T_n$ is symmetric, real, and
additive on admissible combinations.
\item[(b)] Let $j_1=j_1(n,\rho)$ be the least index such that
\begin{equation}\label{10.7:eq-equicontinuity-1}
m_j\ \ge\ \max\bigl\{m_{\mathbb{T}}(g_-(g)),\ \log_L\ell(n,\rho)\bigr\}\qquad\text{for all } j\ge j_1 .
\end{equation}
Let $j\ge j_1$ and $\mathsf{d}'>0$. Let $A_1,\dots,A_n\subset B(0,\rho)$ be sets at pairwise distance
at least $\mathsf{d}'$. Let $\vec g,\vec g'$ be tuples of real, compactly supported Lipschitz functions
with $\operatorname{supp}g_k\cup\operatorname{supp}g'_k\subset A_k$ for $k=1,\dots,n$. Then
\begin{equation}\label{10.7:eq-equicontinuity-a}
\begin{gathered}
|S_j(\vec g)|\ \le\ C^\infty_n(\mathsf{d}';\rho)\prod_{k=1}^{n}\|g_k\|_{\sharp},\\
|S_j(\vec g)-S_j(\vec g')|\ \le\ C^\infty_n(\mathsf{d}';\rho)\sum_{k=1}^{n}\|g_k-g'_k\|_{\sharp}
\prod_{l\neq k}\max\bigl\{\|g_l\|_{\sharp},\|g'_l\|_{\sharp}\bigr\}.
\end{gathered}
\end{equation}
Both inequalities of \eqref{10.7:eq-equicontinuity-a} also hold with $S^T_n$ in place of $S_j$, for
tuples $\vec g,\vec g'$ of $C^1_c$ functions as above. The index $j_1$ depends neither on
$\mathsf{d}'$ nor on the sets $A_k$.
\item[(c)] $S^T_n(w\cdot\vec g)=S^T_n(\vec g)$ for every $w\in W_4$.
\end{enumerate}
\end{lemma}

\begin{proof}
Part (a) is the statement taken from the infinite-volume theorem, read along the diagonal; its proof
is the proof of that theorem. The proof below uses (a) only to pass to the limit $j\to\infty$.

(b) Let $j\ge j_1$. By \eqref{10.7:eq-equicontinuity-1}, $m_j\ge\log_L\ell(n,\rho)$. This is the
floor in the per-ball volume-free bound, and it gives
\begin{equation}\label{10.7:eq-equicontinuity-2}
L^{m_j}\ \ge\ 8n(\rho+1)\ >\ 2\rho\ >\ 0 .
\end{equation}
Each of the functions $g_k$, $g'_k$ and $g_k-g'_k$ is Lipschitz and supported in
$A_k\subset B(0,\rho)$. Hence the hypotheses of Lemma~\ref{10.7:lem-periodic-lift} hold for each of
them with $m=m_j$. By parts (b) and (c) of that lemma:
\begin{itemize}
\item each periodic lift satisfies $\|g_k^{(m_j)}\|_{\sharp}\le\|g_k\|_{\sharp}$, and likewise for
$g'_k$ and $g_k-g'_k$;
\item its support lies in the torus ball of radius $\rho$ about $\pi(0)$;
\item torus separations between such lifts equal Euclidean separations, so the lifts of the $k$-th
and $l$-th entries, $k\ne l$, are supported at torus distance at least $\mathsf{d}'$.
\end{itemize}
The periodic lift is a sum of translates of its argument, so it is linear. Therefore
$(g_k-g'_k)^{(m_j)}=g_k^{(m_j)}-g_k'^{(m_j)}$.

The per-ball volume-free bound now applies on $\mathbb{T}_{m_j}$ to the lifted tuple. Its coupling
hypotheses are met: $g_0^{(j)}\le g_-(g)$, $K_j=K(g_0^{(j)},g)$, and, by
\eqref{10.7:eq-equicontinuity-1}, $m_j\ge m_{\mathbb{T}}(g_-(g))$. It bounds $|S_j(\vec g)|$ by
$C^\infty_n(\mathsf{d}';\rho)\prod_k\|g_k^{(m_j)}\|_{\sharp}$. With the norm inequality above this is
the first inequality of \eqref{10.7:eq-equicontinuity-a}. The same holds for every tuple whose $k$-th
entry is supported in $A_k$ for each $k$.

By part (d) of Lemma~\ref{10.7:lem-joint-cumulants}, $S_j(\vec g)$ is the mixed derivative
$\partial_{z_1}\cdots\partial_{z_n}\log Z(z)$ at $z=0$. In other words, it is the joint cumulant of the
observables $\mathcal{O}(g_1^{(m_j)}),\dots,\mathcal{O}(g_n^{(m_j)})$. By part (c) of the same lemma
the cumulant is multilinear. Since $\mathcal{O}$ and the lift are both linear in the test function,
$S_j$ is multilinear in $\vec g$.

Telescoping slot by slot gives
\begin{equation}\label{10.7:eq-equicontinuity-3}
S_j(\vec g)-S_j(\vec g')\ =\ \sum_{k=1}^{n}
S_j\bigl(g'_1,\dots,g'_{k-1},\,g_k-g'_k,\,g_{k+1},\dots,g_n\bigr).
\end{equation}
The right side is a telescoping sum. Each of its terms is the difference of two consecutive tuples in
the chain from $\vec g$ to $\vec g'$, one entry changed at a time, by multilinearity in the $k$-th
slot. In the $k$-th term, each entry is supported in the corresponding $A_l$. Applying the first
inequality of \eqref{10.7:eq-equicontinuity-a} to it gives the bound
\begin{equation*}
C^\infty_n(\mathsf{d}';\rho)\,\|g_k-g'_k\|_{\sharp}\prod_{l<k}\|g'_l\|_{\sharp}
\prod_{l>k}\|g_l\|_{\sharp}.
\end{equation*}
Each of these norms is at most $\max\{\|g_l\|_{\sharp},\|g'_l\|_{\sharp}\}$. Summing over $k$ gives the
second inequality of \eqref{10.7:eq-equicontinuity-a}.

The constant $C^\infty_n(\mathsf{d}';\rho)$ does not depend on $j$. Let $\vec g,\vec g'$ be $C^1_c$
tuples as in (b), for which (a) gives the limits of $S_j(\vec g)$ and $S_j(\vec g')$. Letting
$j\to\infty$ in both inequalities of \eqref{10.7:eq-equicontinuity-a} yields the same inequalities
for $S^T_n$.

Finally, $j_1$ is defined by \eqref{10.7:eq-equicontinuity-1}. That condition involves only $n$,
$\rho$, $g$ and the sequence $(m_j)$. Hence $j_1$ depends neither on $\mathsf{d}'$ nor on the $A_k$.

(c) Let $w\in W_4$. By Lemma~\ref{10.7:lem-finite-cutoff-invariance},
$S_j(w\cdot\vec g)=S_j(\vec g)$ for every $j$. Taking $j\to\infty$ and using (a) on both sides gives
$S^T_n(w\cdot\vec g)=S^T_n(\vec g)$.
\end{proof}

\begin{lemma}[Simultaneous subsequential convergence from the torus bound]
\label{10.7:lem-simultaneous-convergence}
Fix the reference coupling $g$. Let $m_j \to \infty$, let $g_0^{(j)} \downarrow 0$, let
$K_j = K(g_0^{(j)}, g)$, and write $S_j := S^T_{n;K_j,m_j}$. Then every subsequence of
$(m_j, g_0^{(j)}, K_j)_j$ has a further subsequence along which $S_j(\vec f)$ converges, for every
$n \ge 2$ and every admissible tuple $\vec f = (f_1,\dots,f_n)$ of functions in
$C^1_c(\mathbb{R}^4)$, admissibility being understood as in Lemma~\ref{10.7:lem-equicontinuity}.
The limit, as a function of $\vec f$, is symmetric, real, additive and $W_4$-invariant, and it obeys
the limiting form of the equicontinuity estimate \eqref{10.7:eq-equicontinuity-a}.
\end{lemma}

\begin{proof}
The inputs are the equicontinuity estimate, Lemma~\ref{10.7:lem-equicontinuity}, which rests on
the per-ball volume-free bound, and exact hyperoctahedral invariance at finite cutoff,
Lemma~\ref{10.7:lem-finite-cutoff-invariance}.

\emph{A countable skeleton.} Fix a real function $\psi \in C^\infty_c(\mathbb{R}^4)$ with
$\psi \ge 0$, $\operatorname{supp}\psi$ contained in the closed unit ball and $\int \psi = 1$. For
$r > 0$ put $\psi_r(x) = r^{-4}\psi(x/r)$. Let $\mathcal{S}$ be the set of finite sums
$\sum_i q_i\, \psi_{r_i}(\cdot - a_i)$ with $q_i \in \mathbb{Q}$, $a_i \in \mathbb{Q}^4$ and
$r_i \in \mathbb{Q}_+$. The set $\mathcal{S}$ is countable, and its elements are real, smooth and
compactly supported.

\emph{Approximation by the skeleton.} Let $f \in C^1_c(\mathbb{R}^4)$ be real, let
$V \supset \operatorname{supp} f$ be open, and let $\eta > 0$. We claim that there is
$\varphi \in \mathcal{S}$ with $\operatorname{supp} \varphi \subset V$ and
$\|f - \varphi\|_{\sharp} \le \eta$. Let
$\delta = \operatorname{dist}(\operatorname{supp} f, \mathbb{R}^4 \setminus V) > 0$. Since $f$ and
$\nabla f$ are uniformly continuous, $f * \psi_r \to f$ and
$\nabla(f * \psi_r) = (\nabla f) * \psi_r \to \nabla f$ uniformly as $r \to 0$. For a $C^1$
function the mean value theorem bounds the Lipschitz constant by the supremum of the gradient, so
$\|f - f*\psi_r\|_{\sharp} \to 0$. Choose $r \in \mathbb{Q}_+$ with $r < \delta/3$ and
$\|f - f*\psi_r\|_{\sharp} \le \eta/3$. For $h \in \mathbb{Q}_+$ set
\begin{equation*}
\Sigma_h(x) = \sum_{a \in h\mathbb{Z}^4} h^4 f(a)\, \psi_r(x - a).
\end{equation*}
The sum is finite, only the points $a \in \operatorname{supp} f$ contribute, and therefore
$\operatorname{supp} \Sigma_h$ lies in the $r$-neighbourhood of $\operatorname{supp} f$, hence in
$V$. The functions $\Sigma_h$ and $\nabla \Sigma_h$ are Riemann sums of the integrals defining
$f * \psi_r$ and $f * \nabla\psi_r = \nabla(f*\psi_r)$. The integrands
$a \mapsto f(a)\psi_r(x-a)$ and
$a \mapsto f(a)\nabla\psi_r(x-a)$ are supported in $\operatorname{supp} f$ and are uniformly
continuous in $a$, uniformly in $x$. Hence $\Sigma_h \to f * \psi_r$ in $C^1$ as $h \to 0$, and
$\|\Sigma_h - f*\psi_r\|_{\sharp} \to 0$. Fix $h$ with
$\|\Sigma_h - f*\psi_r\|_{\sharp} \le \eta/3$; then
every point $a \in h\mathbb{Z}^4$ lies in $\mathbb{Q}^4$. Finally replace each of the finitely many
coefficients $h^4 f(a)$, $a \in h\mathbb{Z}^4 \cap \operatorname{supp} f$, by a rational number
$q_a$ with
\begin{equation*}
\|\psi_r\|_{\sharp} \sum_a |q_a - h^4 f(a)| \le \eta/3 .
\end{equation*}
The function $\varphi = \sum_a q_a \psi_r(\cdot - a)$ lies in $\mathcal{S}$ and has support in
$V$. Since $\|\cdot\|_{\sharp}$ is translation invariant,
\begin{equation*}
\|\varphi - \Sigma_h\|_{\sharp} \le \|\psi_r\|_{\sharp} \sum_a |q_a - h^4 f(a)| \le \eta/3,
\end{equation*}
and the triangle inequality gives $\|f - \varphi\|_{\sharp} \le \eta$.

\emph{Diagonal extraction.} Fix a subsequence of the given sequence; it again satisfies the
hypotheses. For each $n \ge 2$ the admissible $n$-tuples with entries in $\mathcal{S}$ form a
countable set, and so does their union over $n$. Enumerate them as $\vec s_1, \vec s_2, \dots$.
The entries of $\vec s_l$ have compact supports, contained in some ball $B(0,\rho)$ and pairwise at
some positive distance. By the first estimate in \eqref{10.7:eq-equicontinuity-a}, the sequence
$S_j(\vec s_l)$ is bounded for $j \ge j_1(n,\rho)$, and the finitely many earlier indices do not
affect boundedness. Extract successively nested subsequences along which $S_j(\vec s_1)$,
$S_j(\vec s_2)$, \dots converge. By Cantor's diagonal procedure there is a single subsequence
$(j(i))_{i \ge 1}$ along which $S_{j(i)}(\vec s_l)$ converges for every $l$.

\emph{Extension to all admissible tuples.} Let $\vec f$ be an admissible tuple of real
$C^1_c$ functions and let $\varepsilon > 0$. Let $\mathsf{d}_0 > 0$ be the least mutual distance of
the supports $\operatorname{supp} f_k$, and let $\rho_0 > 0$ be such that all of them lie in the open
ball $B(0,\rho_0)$. Put
\begin{equation*}
V_k := \{x : \operatorname{dist}(x, \operatorname{supp} f_k) < \mathsf{d}_0/4\} \cap B(0,\rho_0).
\end{equation*}
Each $V_k$ is open, contains $\operatorname{supp} f_k$ and lies in $B(0,\rho_0)$. By the triangle
inequality, $V_k$ and $V_l$ are at distance at least $\mathsf{d}_0/2$ for $k \ne l$. Apply the
difference estimate in \eqref{10.7:eq-equicontinuity-a} with $A_k = V_k$,
$\mathsf{d}' = \mathsf{d}_0/2$ and $\rho = \rho_0$. It yields $\eta > 0$, independent of $j$, such
that
\begin{equation*}
|S_j(\vec f) - S_j(\vec g)| \le \varepsilon/3
\end{equation*}
for all $j \ge j_1(n,\rho_0)$ and every real compactly supported Lipschitz tuple $\vec g$ with
$\operatorname{supp} g_k \subset V_k$ and $\|f_k - g_k\|_{\sharp} \le \eta$. By the approximation
step, choose $g_k \in \mathcal{S}$ with $\operatorname{supp} g_k \subset V_k$ and
$\|f_k - g_k\|_{\sharp} \le \eta$. Then $\vec g$ is an admissible skeleton tuple, so
$S_{j(i)}(\vec g)$ converges. For $i, i'$ large enough we have
$j(i), j(i') \ge j_1(n,\rho_0)$ and
$|S_{j(i)}(\vec g) - S_{j(i')}(\vec g)| \le \varepsilon/3$, and hence
\begin{align*}
|S_{j(i)}(\vec f) - S_{j(i')}(\vec f)|
&\le |S_{j(i)}(\vec f) - S_{j(i)}(\vec g)| + |S_{j(i)}(\vec g) - S_{j(i')}(\vec g)| \\
&\quad + |S_{j(i')}(\vec g) - S_{j(i')}(\vec f)| \\
&\le \varepsilon .
\end{align*}
Thus $S_{j(i)}(\vec f)$ is Cauchy and converges, for every $n \ge 2$ and every admissible $\vec f$.

\emph{Properties of the limit.} Symmetry under permutation of the entries, reality and additivity
on admissible combinations hold for each $S_j$. They are identities between values at admissible
tuples, and so they survive the pointwise limit. For $w \in W_4$ the tuple $w\cdot\vec f$ is again
admissible, because $w$ is an orthogonal map fixing the origin and so preserves mutual distances of
supports and balls about $0$. By exact hyperoctahedral invariance at finite cutoff
(Lemma~\ref{10.7:lem-finite-cutoff-invariance}), $S_{j(i)}(w\cdot\vec f) = S_{j(i)}(\vec f)$ for
every $i$, and letting $i \to \infty$ gives $W_4$-invariance of the limit. Finally, the estimates
\eqref{10.7:eq-equicontinuity-a} hold for all $j \ge j_1(n,\rho)$ with constants independent of
$j$. Letting $j = j(i) \to \infty$ in them gives their limiting form for the limit.
\end{proof}

\begin{lemma}[Uniform continuity of the correlation functions on the orthogonal group]
\label{10.7:lem-orthogonal-continuity}
Let $n\ge 2$, and let $S$ be one of the following: the limit $S^T_n$ along the diagonal; any limit
obtained in Lemma~\ref{10.7:lem-simultaneous-convergence}; or $S_j$ for $j\ge j_1(n,\rho_0)$.
Let $\vec f=(f_1,\dots,f_n)$ be an admissible tuple whose supports have minimal separation
$\mathsf{d}_0$ and lie in $B(0,\rho_0)$. Let $R,R'\in O(4)$, write $R\cdot\vec f=(R\cdot f_1,\dots,R\cdot f_n)$,
and let $\|R'-R\|$ be the operator norm on $\mathbb{R}^4$. Put
\begin{equation*}
s:=\|R'-R\|, \qquad s\le \frac{\mathsf{d}_0}{4\rho_0}.
\end{equation*}
Then
\begin{equation}\label{10.7:eq-orthogonal-continuity-1}
\bigl|S(R'\cdot\vec f)-S(R\cdot\vec f)\bigr|
\le C^{\infty}_n(\mathsf{d}_0/2;\rho_0)\sum_{i=1}^{n}\mu_{f_i}(s)\prod_{k\ne i}\|f_k\|_{\sharp}.
\end{equation}
For $S=S_j$ this holds for every $j\ge j_1(n,\rho_0)$ with the same right-hand side. Hence the map
$R\mapsto S(R\cdot\vec f)$ is uniformly continuous on $O(4)$.
\end{lemma}

\begin{proof}
The notion of an admissible tuple, the rotation modulus $\mu_f$, the support radius $\rho_f$, the
neighbourhoods $A^{+\theta}$ and the action $(R\cdot f)(x)=f(R^{-1}x)$ are those fixed in the
notation for test-function moduli, separated classes and the action of rotations,
\ref{10.7:not-test-function-moduli}.

We first show $\|R'\cdot f-R\cdot f\|_{\sharp}\le\mu_f(s)$ for each $f_i$ of the tuple, written $f$.
If $R'\cdot f-R\cdot f$ or its gradient is non-zero at $x$, then $R'^{-1}x$ or $R^{-1}x$ lies in
$\operatorname{supp}f$. Since $R,R'$ preserve the Euclidean norm, $|x|\le\rho_f$. Moreover
$R'^{-1}-R^{-1}=(R'-R)^{T}$ has operator norm $s$, so
\begin{equation*}
|R'^{-1}x-R^{-1}x|\le s|x|\le s\rho_f .
\end{equation*}
Thus $|(R'\cdot f)(x)-(R\cdot f)(x)|=|f(R'^{-1}x)-f(R^{-1}x)|$ is a difference of values of $f$
at two points at distance at most $s\rho_f$. For the gradient, the chain rule gives
$\nabla(R\cdot f)(x)=R\,\nabla f(R^{-1}x)$ for orthogonal $R$, whence
\begin{equation}\label{10.7:eq-orthogonal-continuity-2}
\nabla(R'\cdot f)(x)-\nabla(R\cdot f)(x)
=(R'-R)\nabla f(R'^{-1}x)+R\bigl[\nabla f(R'^{-1}x)-\nabla f(R^{-1}x)\bigr].
\end{equation}
The first term has norm at most $s\sup|\nabla f|$. The second, since $\|R\|=1$, is a difference of
values of $\nabla f$ at two points at distance at most $s\rho_f$. The Lipschitz constant of the
$C^1$ compactly supported function $R'\cdot f-R\cdot f$ is the supremum of the norm of its gradient.
These two bounds are therefore the quantities entering the definition of $\mu_f(s)$, and
$\|R'\cdot f-R\cdot f\|_{\sharp}\le\mu_f(s)$.

Next, the containers. For $i=1,\dots,n$ set
\begin{equation*}
A_i:=(R\operatorname{supp}f_i)^{+\mathsf{d}_0/4}\cap B(0,\rho_0).
\end{equation*}
Since $R$ is an isometry fixing $0$, the sets $R\operatorname{supp}f_i$ lie in $B(0,\rho_0)$ and are
pairwise at distance at least $\mathsf{d}_0$. Enlarging each by $\mathsf{d}_0/4$ lowers the
distances by at most $\mathsf{d}_0/2$, so the $A_i$ are pairwise at distance at least
$\mathsf{d}_0/2$. Since a set is contained in its own $\mathsf{d}_0/4$-neighbourhood and
$R(\operatorname{supp}f_i)\subset B(0,\rho_0)$, we have $R(\operatorname{supp}f_i)\subset A_i$.
For $y\in\operatorname{supp}f_i$,
\begin{equation*}
|R'y-Ry|\le s|y|\le s\rho_0\le \mathsf{d}_0/4 ,
\end{equation*}
and $|R'y|=|y|$, so also $R'(\operatorname{supp}f_i)\subset A_i$.

Finally, apply the estimate \eqref{10.7:eq-equicontinuity-a} of
Lemma~\ref{10.7:lem-equicontinuity}. For $S=S_j$ with $j\ge j_1(n,\rho_0)$ apply it at finite $j$.
For $S^T_n$ apply the same estimate in the limit, which $S^T_n$ satisfies by
Lemma~\ref{10.7:lem-equicontinuity}; for the limits of
Lemma~\ref{10.7:lem-simultaneous-convergence} apply the same estimate in the limit, which these
functionals satisfy by Lemma~\ref{10.7:lem-simultaneous-convergence}. Take $\mathsf{d}'=\mathsf{d}_0/2$,
$\rho=\rho_0$, the sets $A_1,\dots,A_n$, and the
real compactly supported Lipschitz tuples $g_k=R\cdot f_k$, $g'_k=R'\cdot f_k$, with
$\operatorname{supp}g_k\cup\operatorname{supp}g'_k\subset A_k$. An orthogonal change of variables
preserves the supremum norm and the Lipschitz constant, so
$\|R\cdot f_k\|_{\sharp}=\|R'\cdot f_k\|_{\sharp}=\|f_k\|_{\sharp}$. Together with
$\|R'\cdot f_i-R\cdot f_i\|_{\sharp}\le\mu_{f_i}(s)$, the estimate yields
\eqref{10.7:eq-orthogonal-continuity-1}. Its constant $C^{\infty}_n(\mathsf{d}_0/2;\rho_0)$ does
not depend on $j\ge j_1(n,\rho_0)$.

The right-hand side of \eqref{10.7:eq-orthogonal-continuity-1} depends on $R,R'$ only through $s$.
It tends to zero as $s\downarrow 0$, because $\mu_{f_i}(s)\to 0$ for the continuously
differentiable compactly supported $f_i$. Hence $R\mapsto S(R\cdot\vec f)$ is uniformly continuous
on $O(4)$.
\end{proof}

\begin{lemma}[Continuity of rotated test functions]\label{10.7:lem-rotated-continuity}
Let $\omega$ be a real antisymmetric $4\times 4$ matrix, let $R^{\omega}_t = e^{t\omega}$, $t\in\mathbb{R}$, be the
one-parameter group of rotations it generates, acting on test functions by
$(R^{\omega}_t f)(x) = f(R^{\omega}_{-t}x)$, and let $f\in C^1_c(\mathbb{R}^4)$. Then
\begin{equation}\label{10.7:eq-rotated-continuity-1}
\|R^{\omega}_t f - f\|_{\sharp} \le \mu_f(|t|\,\|\omega\|) \longrightarrow 0
\qquad (t\to 0),
\end{equation}
where $\mu_f$ is the rotation modulus of $f$ and $\|\cdot\|$ is the operator norm on matrices.
\end{lemma}

\begin{proof}
In the proof of Lemma~\ref{10.7:lem-orthogonal-continuity} it is shown that
$\|R'\cdot f - R\cdot f\|_{\sharp} \le \mu_f(s)$ for $R,R'\in O(4)$ and $s = \|R'-R\|$. We apply this
with $R' = R^{\omega}_t$ and $R = \mathbb{1}$, so that $s = \|e^{t\omega} - \mathbb{1}\|$. Since $\omega$
is antisymmetric, $e^{u\omega}$ is orthogonal and $\|e^{u\omega}\| = 1$ for every real $u$. Hence, for
$t\ge 0$,
\begin{equation*}
\|e^{t\omega} - \mathbb{1}\|
= \Bigl\|\int_0^t \omega\, e^{u\omega}\,du\Bigr\|
\le \int_0^t \|\omega\|\,\|e^{u\omega}\|\,du
= t\,\|\omega\|,
\end{equation*}
and for $t<0$ the same computation over $[t,0]$ gives $|t|\,\|\omega\|$. As $\mu_f$ is nondecreasing,
$\mu_f(s)\le \mu_f(|t|\,\|\omega\|)$, which is the inequality in \eqref{10.7:eq-rotated-continuity-1}.
Finally $|t|\,\|\omega\|\to 0$ as $t\to 0$, and $\mu_f(r)\to 0$ as $r\to 0$ for $f\in C^1_c$, which gives
the limit.
\end{proof}

\begin{lemma}[Differentiability of rotated test functions]\label{10.7:lem-rotated-differentiability}
Let $f\in C^2_c$. Then $\delta_{\omega}f\in C^1_c$, $\operatorname{supp}\delta_{\omega}f\subset\operatorname{supp}f$, and
for $t\neq 0$
\begin{equation}\label{10.7:eq-rotated-differentiability-1}
\Bigl\|\frac{R^{\omega}_t f-f}{t}-\delta_{\omega}f\Bigr\|_{\sharp}
\le \mu_{\delta_{\omega}f}\bigl(|t|\,\|\omega\|\bigr)\longrightarrow 0
\qquad (t\to 0).
\end{equation}
\end{lemma}

\begin{proof}
Write $g=\delta_{\omega}f=-(\omega x)\cdot\nabla f$. The coefficient $\omega x$ is linear in $x$ and
$\nabla f\in C^1_c$, so $g\in C^1$. Since $\nabla f$ vanishes on the open complement of
$\operatorname{supp}f$, so does $g$; hence $\operatorname{supp}g\subset\operatorname{supp}f$, which is
compact, and $g\in C^1_c$.

Since $\delta_{\omega}$ is the generator of the one-parameter group $R^{\omega}_s$ on test functions,
\begin{equation*}
\frac{d}{ds}\,(R^{\omega}_s f)(x)=(R^{\omega}_s g)(x)
\end{equation*}
for every $x$ and $s$. By the fundamental theorem of calculus,
$R^{\omega}_t f-f=\int_0^t R^{\omega}_s g\,ds$, and since $g=t^{-1}\int_0^t g\,ds$,
\begin{equation}\label{10.7:eq-rotated-differentiability-2}
\frac{R^{\omega}_t f-f}{t}-g=\frac1t\int_0^t\bigl(R^{\omega}_s g-g\bigr)\,ds .
\end{equation}
Put $h(s,x)=(R^{\omega}_s g)(x)-g(x)$. This integrand is jointly continuous in $(s,x)$, of class $C^1$
in $x$, and its $x$-gradient is jointly continuous; differentiation under the integral is therefore
permitted, and the gradient of the right side of \eqref{10.7:eq-rotated-differentiability-2} is
$t^{-1}\int_0^t\nabla_x h(s,x)\,ds$. The right side of \eqref{10.7:eq-rotated-differentiability-2} is
an average over $s$ in the interval between $0$ and $t$, so its sup norm is at most
$\sup_{|s|\le|t|}\|h(s,\cdot)\|_{\infty}$ and the sup norm of its gradient is at most
$\sup_{|s|\le|t|}\sup_x|\nabla_x h(s,x)|$. For a $C^1$ function the Lipschitz constant equals the sup
norm of the gradient, so, by the definition of $\|\cdot\|_{\sharp}$,
\begin{equation*}
\Bigl\|\frac{R^{\omega}_t f-f}{t}-g\Bigr\|_{\sharp}
\le \sup_{|s|\le|t|}\bigl\|R^{\omega}_s g-g\bigr\|_{\sharp}.
\end{equation*}
Finally apply the continuity of rotated test functions (Lemma~\ref{10.7:lem-rotated-continuity}) to
$g\in C^1_c$: for each $s$ with $|s|\le|t|$ it gives
$\|R^{\omega}_s g-g\|_{\sharp}\le\mu_g(|s|\,\|\omega\|)$, and the supremum over $|s|\le|t|$ yields the
bound $\mu_g(|t|\,\|\omega\|)$ of \eqref{10.7:eq-rotated-differentiability-1}. The same lemma gives
$\mu_g(|t|\,\|\omega\|)\to 0$ as $t\to 0$.
\end{proof}

\begin{lemma}[Rotated separated classes and the rotated generator]
\label{10.7:lem-rotated-generator}
Let $R\in O(4)$, let $\omega$ be a real antisymmetric $4\times 4$ matrix, and let
$R_s=R^{\omega}_s=e^{s\omega}$ for $s\in\mathbb{R}$. Then
\begin{gather*}
R\cdot\mathcal{T}^{(r)}_n(\mathcal{R},\mathsf{d})=\mathcal{T}^{(r)}_n(\mathcal{R},\mathsf{d}),\\
\|R\cdot f\|_{\sharp}=\|f\|_{\sharp},\\
R_s\cdot\delta_{\omega}f=\delta_{\omega}(R_s\cdot f),
\end{gather*}
the second identity for every bounded Lipschitz $f$ (in particular every real $f\in C^k_c$,
$k\ge1$), the third for every real $f\in C^1_c$ and every $s\in\mathbb{R}$.
\end{lemma}

\begin{proof}
The action of rotations on test functions fixed in~\ref{10.7:not-test-function-moduli}
is $(R\cdot f)(x)=f(R^{-1}x)$, applied componentwise to tuples
$\vec f=(f_1,\dots,f_n)$. Its generator is
$\delta_{\omega}f=-(\omega x)\cdot\nabla f$. In particular $R\cdot(R'\cdot f)=(RR')\cdot f$.

\emph{The class.} Since $R^{-1}$ is linear, $R\cdot f=f\circ R^{-1}$ is of class $C^r$
whenever $f$ is, and it is real-valued when $f$ is.
Moreover $\operatorname{supp}(R\cdot f)=R(\operatorname{supp}f)$ is compact.
Since $R$ is an isometry fixing $0$, it maps $B(0,\mathcal{R}/2)$ onto itself, and it
preserves pairwise distances of supports and radii of enclosing balls. Hence every tuple in
$\mathcal{T}^{(r)}_n(\mathcal{R},\mathsf{d})$ is mapped to a tuple with the same regularity,
supports in $B(0,\mathcal{R}/2)$ and the same pairwise separations, that is
\begin{equation*}
R\cdot\mathcal{T}^{(r)}_n(\mathcal{R},\mathsf{d})\subset\mathcal{T}^{(r)}_n(\mathcal{R},\mathsf{d}).
\end{equation*}
Applying this inclusion to $R^{-1}$ and using $R\cdot(R^{-1}\cdot\vec f)=\vec f$ gives the
reverse inclusion.

\emph{The norm.} The map $R^{-1}$ is a bijection, so
$\|R\cdot f\|_\infty=\|f\|_\infty$. It is also an isometry, so for all $x,y$
\begin{equation*}
|f(R^{-1}x)-f(R^{-1}y)|\le \operatorname{Lip}f\,|R^{-1}x-R^{-1}y|
=\operatorname{Lip}f\,|x-y| .
\end{equation*}
Hence $\operatorname{Lip}(R\cdot f)\le\operatorname{Lip}f$. Applying this to $R^{-1}$ gives
equality, and therefore $\|R\cdot f\|_{\sharp}=\|f\|_{\sharp}$.

\emph{The generator.} We have $R_s^{-1}=e^{-s\omega}=R_{-s}$, so
$(R_s\cdot f)(x)=f(R_{-s}x)$. Since $\omega$ commutes with each term of the series
$e^{-s\omega}=\sum_{j\ge0}(-s)^j\omega^j/j!$, we have $R_{-s}\omega=\omega R_{-s}$; and
the transpose of $R_{-s}$ moves across the Euclidean inner product. By the
chain rule, $\nabla(f\circ R_{-s})(x)=R_{-s}^{T}\nabla f(R_{-s}x)$, and so
\begin{equation}\label{10.7:eq-rotated-generator-1}
\begin{split}
\delta_{\omega}(R_s\cdot f)(x)&=-(\omega x)\cdot R_{-s}^{T}\nabla f(R_{-s}x)
=-(R_{-s}\omega x)\cdot\nabla f(R_{-s}x)\\
&=-(\omega R_{-s}x)\cdot\nabla f(R_{-s}x)=(\delta_{\omega}f)(R_{-s}x)
=(R_s\cdot\delta_{\omega}f)(x).
\end{split}
\end{equation}
This is the third identity.
\end{proof}

\begin{lemma}[Differentiability of the limit along one-parameter groups]
\label{10.7:lem-limit-differentiability}
Let $\vec h=(h_1,\dots,h_n)\in\mathcal{T}^{(2)}_n(\mathcal{R},\mathsf{d})$ and let $\omega\neq0$. Write
$R_t:=R^{\omega}_t=e^{t\omega}$ and $R_t\cdot\vec h:=(R_t\cdot h_1,\dots,R_t\cdot h_n)$, and put
\begin{equation*}
t_0:=\frac{\mathsf{d}}{2\|\omega\|\mathcal{R}},\qquad
C_*:=C^{\infty}_n(\mathsf{d}/2;\mathcal{R}/2),\qquad
c(\vec h,\omega):=\sum_{i=1}^{n}S^T_n(h_1,\dots,\delta_{\omega}h_i,\dots,h_n),
\end{equation*}
where each tuple in the sum is an admissible $C^1_c$ tuple. Then for $0<|t|\le t_0$
\begin{equation}\label{10.7:eq-limit-differentiability-1}
\begin{split}
&\Bigl|\frac{S^T_n(R_t\cdot\vec h)-S^T_n(\vec h)}{t}-c(\vec h,\omega)\Bigr|\\
&\quad\le C_*\Bigl[\sum_{i}\mu_{\delta_{\omega}h_i}(|t|\|\omega\|)\prod_{k\neq i}\|h_k\|_{\sharp}\\
&\qquad\quad+\sum_{i}\sum_{l<i}\mu_{h_l}(|t|\|\omega\|)\,\|\delta_{\omega}h_i\|_{\sharp}
\prod_{k\neq i,l}\|h_k\|_{\sharp}\Bigr],
\end{split}
\end{equation}
and the right-hand side tends to $0$ as $t\to0$. The same holds with $S_j$ and $D_j(\vec h)$ in place
of $S^T_n$ and $c(\vec h,\omega)$, for every $j\ge j_1(n,\mathcal{R}/2)$, with the same right-hand
side, hence uniformly in $j$; here $D_j(\vec h)$ is the derivative at $t=0$ of
$t\mapsto S_j(R_t\cdot\vec h)$ given by Lemma~\ref{10.7:lem-finite-cutoff-structure}(c).
\end{lemma}

\begin{proof}
Throughout, $0<|t|\le t_0$. By Lemma~\ref{10.7:lem-rotated-differentiability}, each
$\delta_{\omega}h_i$ lies in $C^1_c$ and $\operatorname{supp}\delta_{\omega}h_i\subset\operatorname{supp}h_i$.

\emph{Containers.} Let $x\in B(0,\mathcal{R}/2)$. Since $R_s$ is a rotation, $|R_sx|=|x|$ for all
$s$, and $R_tx-x=\int_0^t\omega R_sx\,ds$, so that
\begin{equation*}
|R_tx-x|\le|t|\,\|\omega\|\,|x|\le\frac{|t|\,\|\omega\|\,\mathcal{R}}{2}
\le\frac{t_0\|\omega\|\mathcal{R}}{2}=\frac{\mathsf{d}}{4}.
\end{equation*}
Put $A_k:=(\operatorname{supp}h_k)^{+\mathsf{d}/4}\cap B(0,\mathcal{R}/2)$. Since
$\operatorname{supp}h_k\subset B(0,\mathcal{R}/2)$ and $R_t$ preserves $|x|$, the set $A_k$ contains
$\operatorname{supp}h_k$ and $R_t\operatorname{supp}h_k=\operatorname{supp}(R_t\cdot h_k)$; since
$\operatorname{supp}\delta_{\omega}h_k\subset\operatorname{supp}h_k$, it also contains
$\operatorname{supp}\delta_{\omega}h_k$ and $R_t\operatorname{supp}\delta_{\omega}h_k$. The supports of
the $h_k$ are pairwise at distance at least $\mathsf{d}$, so the $A_k$ are pairwise at distance at
least $\mathsf{d}-2\cdot\mathsf{d}/4=\mathsf{d}/2$, and $A_k\subset B(0,\mathcal{R}/2)$. Hence the
equicontinuity estimate (Lemma~\ref{10.7:lem-equicontinuity},
\eqref{10.7:eq-equicontinuity-a}) applies, with $\rho=\mathcal{R}/2$ and $\mathsf{d}'=\mathsf{d}/2$, to
every tuple of real $C^1_c$ functions whose $k$-th entry is supported in $A_k$: for $S_j$ with
$j\ge j_1(n,\mathcal{R}/2)$, and in its limiting form for $S^T_n$, with constant
$C^{\infty}_n(\mathsf{d}/2;\mathcal{R}/2)=C_*$. In the form used below it bounds the absolute value of
the correlation function of such a tuple by $C_*$ times the product of the $\|\cdot\|_{\sharp}$ norms
of its entries.

\emph{Telescoping.} Let $S$ denote either $S^T_n$ or $S_j$ with $j\ge j_1(n,\mathcal{R}/2)$. The
functions $S_j$ are multilinear (Lemma~\ref{10.7:lem-finite-cutoff-structure}(a)), and $S^T_n$ is
additive on admissible combinations (Lemma~\ref{10.7:lem-equicontinuity}(a)) and, as the limit of the
$S_j$, homogeneous in each slot; all tuples written below have $k$-th entry supported in $A_k$ and are
admissible. Set $u_k:=R_t\cdot h_k$, and write $(u_{<i},g,h_{>i})$ for the tuple
$(u_1,\dots,u_{i-1},g,h_{i+1},\dots,h_n)$, with similar notation for other mixed tuples. Replacing the
entries one at a time, rotated slots first, and using additivity in slot $i$,
\begin{equation*}
S(\vec u)-S(\vec h)=\sum_{i=1}^{n}\bigl[S(u_{\le i},h_{>i})-S(u_{<i},h_{\ge i})\bigr]
=\sum_{i=1}^{n}S(u_{<i},u_i-h_i,h_{>i}).
\end{equation*}
Put $g_{t,i}:=(u_i-h_i)/t-\delta_{\omega}h_i$, supported in $A_i$. Dividing by $t$, writing
$(u_i-h_i)/t=g_{t,i}+\delta_{\omega}h_i$ and subtracting
$c:=\sum_iS(h_{<i},\delta_{\omega}h_i,h_{>i})$ gives
\begin{equation}\label{10.7:eq-limit-differentiability-2}
\frac{S(\vec u)-S(\vec h)}{t}-c
=\sum_{i}S(u_{<i},g_{t,i},h_{>i})
+\sum_{i}\bigl[S(u_{<i},\delta_{\omega}h_i,h_{>i})-S(h_{<i},\delta_{\omega}h_i,h_{>i})\bigr].
\end{equation}
For $S=S^T_n$, $c=c(\vec h,\omega)$. For $S=S_j$, Lemma~\ref{10.7:lem-finite-cutoff-structure}(c)
gives the derivative of $t\mapsto S_j(R_t\cdot\vec h)$ at $t=0$ as
$\sum_iS_j(h_1,\dots,\delta_{\omega}h_i,\dots,h_n)$, so $c=D_j(\vec h)$. The bracket in
\eqref{10.7:eq-limit-differentiability-2} telescopes over $l<i$, replacing the entries $1,\dots,i-1$
one at a time and using additivity in slot $l$:
\begin{equation}\label{10.7:eq-limit-differentiability-3}
\begin{split}
&S(u_{<i},\delta_{\omega}h_i,h_{>i})-S(h_{<i},\delta_{\omega}h_i,h_{>i})\\
&\quad=\sum_{l<i}S(h_{<l},\,u_l-h_l,\,u_{l+1},\dots,u_{i-1},\,\delta_{\omega}h_i,\,h_{>i}).
\end{split}
\end{equation}

\emph{Bounds.} By Lemma~\ref{10.7:lem-rotated-generator}, $\|u_k\|_{\sharp}=\|h_k\|_{\sharp}$. Since
$h_i\in C^2_c$, Lemma~\ref{10.7:lem-rotated-differentiability} gives
$\|g_{t,i}\|_{\sharp}\le\mu_{\delta_{\omega}h_i}(|t|\|\omega\|)$; the equicontinuity estimate then
bounds the $i$-th term of the first sum in \eqref{10.7:eq-limit-differentiability-2} by
\begin{equation*}
C_*\,\|g_{t,i}\|_{\sharp}\prod_{k<i}\|u_k\|_{\sharp}\prod_{k>i}\|h_k\|_{\sharp}
\le C_*\,\mu_{\delta_{\omega}h_i}(|t|\|\omega\|)\prod_{k\neq i}\|h_k\|_{\sharp}.
\end{equation*}
Since $h_l\in C^1_c$, Lemma~\ref{10.7:lem-rotated-continuity} gives
$\|u_l-h_l\|_{\sharp}\le\mu_{h_l}(|t|\|\omega\|)$, and the equicontinuity estimate bounds the $l$-th
term on the right of \eqref{10.7:eq-limit-differentiability-3} by
\begin{equation*}
C_*\,\mu_{h_l}(|t|\|\omega\|)\,\|\delta_{\omega}h_i\|_{\sharp}\prod_{k\neq i,l}\|h_k\|_{\sharp}.
\end{equation*}
Summing over $i$ and $l<i$ gives \eqref{10.7:eq-limit-differentiability-1} for $S^T_n$, and the same
inequality for $S_j$ with $D_j(\vec h)$ in place of $c(\vec h,\omega)$, for every
$j\ge j_1(n,\mathcal{R}/2)$ with the same constant $C_*$, hence uniformly in $j$. The right-hand side
tends to $0$ as $t\to0$ because $\mu_{\delta_{\omega}h_i}(s)\to0$ and $\mu_{h_l}(s)\to0$ as $s\to0$
(Lemmas~\ref{10.7:lem-rotated-continuity} and~\ref{10.7:lem-rotated-differentiability}).
\end{proof}

\begin{lemma}[Convergence of the finite-cutoff rotation derivative]
\label{10.7:lem-derivative-convergence}
Let $\vec h=(h_1,\dots,h_n)\in\mathcal{T}^{(2)}_n(\mathcal{R},\mathsf{d})$, and let $\omega$, $R^{\omega}_t=e^{t\omega}$,
$D_j(\vec h)$ and $c(\vec h,\omega)$ be as in Lemma~\ref{10.7:lem-limit-differentiability}. Then
\begin{equation}\label{10.7:eq-derivative-convergence-1}
\lim_{j\to\infty} D_j(\vec h)=c(\vec h,\omega).
\end{equation}
Moreover, the function
\begin{equation}\label{10.7:eq-derivative-convergence-2}
s\longmapsto c(R^{\omega}_s\cdot\vec h,\omega)
=\sum_{i=1}^n S^T_n\bigl(R^{\omega}_s\cdot(h_1,\dots,\delta_{\omega}h_i,\dots,h_n)\bigr)
\end{equation}
is continuous on $\mathbb{R}$.
\end{lemma}

\begin{proof}
For $i=1,\dots,n$, write $\vec f^{(i)}$ for the tuple $(h_1,\dots,\delta_{\omega}h_i,\dots,h_n)$, in which
$h_i$ is replaced by $\delta_{\omega}h_i$. Since $h_i$ is real and of class $C^2_c$, the function
$\delta_{\omega}h_i=-(\omega x)\cdot\nabla h_i$ is real and of class $C^1_c$, and its support is
contained in $\operatorname{supp}h_i\subset B(0,\mathcal{R}/2)$. Hence the supports of the entries of
$\vec f^{(i)}$ are contained in those of the entries of $\vec h$. The separation of the supports is
therefore at least that of $\vec h$, and each $\vec f^{(i)}$ is an admissible $C^1_c$ tuple, as recorded in
Lemma~\ref{10.7:lem-limit-differentiability}.

First assertion. The number $D_j(\vec h)$ is the derivative at $t=0$ of
$t\mapsto S_j(R^{\omega}_t\cdot\vec h)$. For the indices $j\ge j_1(n,\mathcal{R}/2)$, for which
Lemma~\ref{10.7:lem-limit-differentiability} treats it, Lemma~\ref{10.7:lem-finite-cutoff-structure}(c)
evaluated at $t=0$ gives
\begin{equation*}
D_j(\vec h)=\sum_{i=1}^n S_j\bigl(\vec f^{(i)}\bigr).
\end{equation*}
The infinite-volume theorem gives $S_j(\vec f)\to S^T_n(\vec f)$ along the diagonal for
every admissible tuple $\vec f$. We apply it to each of the $n$ admissible tuples $\vec f^{(i)}$ and sum
over $i$. This yields
\begin{equation*}
\lim_{j\to\infty}D_j(\vec h)=\sum_{i=1}^n S^T_n\bigl(\vec f^{(i)}\bigr)=c(\vec h,\omega),
\end{equation*}
which is \eqref{10.7:eq-derivative-convergence-1}.

Second assertion. The matrix $R^{\omega}_s$ commutes with $\omega$, so $R^{\omega}_s\delta_{\omega}h_i$
equals $\delta_{\omega}(R^{\omega}_s h_i)$ (Lemma~\ref{10.7:lem-finite-cutoff-structure}(c)). Inserting
this into the definition of $c(R^{\omega}_s\cdot\vec h,\omega)$ gives the equality in
\eqref{10.7:eq-derivative-convergence-2}, namely
$c(R^{\omega}_s\cdot\vec h,\omega)=\sum_i S^T_n(R^{\omega}_s\cdot\vec f^{(i)})$.

Fix $i$. By Lemma~\ref{10.7:lem-orthogonal-continuity}, applied to the admissible tuple
$\vec f^{(i)}$, the map $R\mapsto S^T_n(R\cdot\vec f^{(i)})$ is uniformly continuous on $O(4)$. The map
$s\mapsto R^{\omega}_s=e^{s\omega}$ is continuous from $\mathbb{R}$ into $O(4)$. Their composition
$s\mapsto S^T_n(R^{\omega}_s\cdot\vec f^{(i)})$ is therefore continuous. A finite sum of continuous
functions is continuous, which proves the continuity of \eqref{10.7:eq-derivative-convergence-2}.
\end{proof}

Lemma~\ref{10.7:lem-limit-differentiability} and Lemma~\ref{10.7:lem-derivative-convergence} together
say that, for $\vec h\in\mathcal{T}^{(2)}_n(\mathcal{R},\mathsf{d})$,
\begin{equation*}
\lim_{j\to\infty}\frac{d}{dt}\Big|_{t=0}S_j(R^{\omega}_t\cdot\vec h)
=\frac{d}{dt}\Big|_{t=0}\lim_{j\to\infty}S_j(R^{\omega}_t\cdot\vec h).
\end{equation*}
The two sides are computed independently. The left side equals $c(\vec h,\omega)$ by
\eqref{10.7:eq-derivative-convergence-1}. The right side is the derivative at $t=0$ of
$S^T_n(R^{\omega}_t\cdot\vec h)$, and it equals $c(\vec h,\omega)$ by
Lemma~\ref{10.7:lem-limit-differentiability}. The interchange of $d/dt$ with $\lim_j$ is thus a
consequence of the bound of Lemma~\ref{10.7:lem-limit-differentiability}, which is uniform in $j$. It is
never an assumption.

\begin{lemma}[Surjectivity of the exponential map of the rotation group]
\label{10.7:lem-exponential-onto}
Let $\mathfrak{so}(4)$ denote the real antisymmetric $4\times 4$ matrices, let
\begin{equation*}
J=\begin{pmatrix}0&-1\\ 1&0\end{pmatrix},
\qquad
\operatorname{Rot}_\theta=\begin{pmatrix}\cos\theta&-\sin\theta\\ \sin\theta&\cos\theta\end{pmatrix}
\quad(\theta\in\mathbb{R}).
\end{equation*}
Then every $R\in SO(4)$ can be written as
$R=Q(\operatorname{Rot}_{\theta_1}\oplus\operatorname{Rot}_{\theta_2})Q^{T}$
with $Q\in O(4)$ and $\theta_1,\theta_2\in\mathbb{R}$. Moreover
\begin{equation}\label{10.7:eq-exponential-onto-1}
R=\exp\bigl(Q(\theta_1J\oplus\theta_2J)Q^{T}\bigr),
\qquad Q(\theta_1J\oplus\theta_2J)Q^{T}\in\mathfrak{so}(4),
\end{equation}
so that $\exp\colon\mathfrak{so}(4)\to SO(4)$ is onto.
\end{lemma}

\begin{proof}
Let $R\in SO(4)$. Since $R^TR=\mathbb{1}$, every eigenvalue of $R$ has modulus one.
By the real normal form of an orthogonal matrix, $\mathbb{R}^4$ is the orthogonal direct sum
of $R$-invariant subspaces of dimension one or two, of the following kinds.
Each pair $e^{\pm i\theta}$ of non-real conjugate eigenvalues gives an invariant plane on which
$R$ acts, in a suitable orthonormal basis of the plane, as $\operatorname{Rot}_\theta$.
Each real eigenvalue gives an invariant line on which $R$ acts by $+1$ or by $-1$.

Since the planes have dimension two and the total dimension is four, the number of invariant
lines is even. Each block $\operatorname{Rot}_\theta$ has determinant one, so $\det R=(-1)^{n_-}$,
where $n_-$ is the number of lines on which $R$ acts by $-1$. The condition $\det R=1$ forces
$n_-$ to be even, and hence the number of lines on which $R$ acts by $+1$ is even as well.
We pair the $-1$ lines two by two, and on each such pair $R$ acts as
$\operatorname{diag}(-1,-1)=\operatorname{Rot}_\pi$. We pair the $+1$ lines likewise, and on
each such pair $R$ acts as $\operatorname{diag}(1,1)=\operatorname{Rot}_0$.
This decomposes $\mathbb{R}^4$ into two orthogonal invariant planes, on each of which $R$ is a
rotation. Let $Q\in O(4)$ be the matrix whose columns are the resulting orthonormal basis,
ordered plane by plane. Then $R=Q(\operatorname{Rot}_{\theta_1}\oplus\operatorname{Rot}_{\theta_2})Q^{T}$.

Next, $J^2=-\mathbb{1}$. Splitting the exponential series into even and odd powers gives
\begin{equation*}
e^{\theta J}=\sum_{n\ge0}\frac{(-1)^n\theta^{2n}}{(2n)!}\,\mathbb{1}
+\sum_{n\ge0}\frac{(-1)^n\theta^{2n+1}}{(2n+1)!}\,J
=\cos\theta\,\mathbb{1}+\sin\theta\,J=\operatorname{Rot}_\theta .
\end{equation*}
The exponential of a block-diagonal matrix is the block-diagonal matrix of the exponentials, so
$e^{\theta_1J\oplus\theta_2J}=\operatorname{Rot}_{\theta_1}\oplus\operatorname{Rot}_{\theta_2}$.
Put $X=\theta_1J\oplus\theta_2J$. Since $Q^TQ=\mathbb{1}$, we have $(QXQ^T)^n=QX^nQ^T$ for every
$n\ge0$, and therefore
\begin{equation*}
\exp(QXQ^T)=Q\,e^{X}Q^T=Q(\operatorname{Rot}_{\theta_1}\oplus\operatorname{Rot}_{\theta_2})Q^T=R .
\end{equation*}
Finally, $X^T=-X$ because $J^T=-J$. Hence $(QXQ^T)^T=QX^TQ^T=-QXQ^T$, so
$QXQ^T\in\mathfrak{so}(4)$. This proves \eqref{10.7:eq-exponential-onto-1}, and with it the
surjectivity of $\exp$.
\end{proof}

\begin{lemma}[From special orthogonal to orthogonal invariance]\label{10.7:lem-orthogonal-invariance}
Let a class of tuples of test functions be given that is stable under the action of $W_4$ and
under the action of $SO(4)$, the action of $O(4)$ on test functions being the one fixed in
\ref{10.7:not-test-function-moduli}. For $R \in O(4)$ and an element $F$ of the class, write
$R\cdot F$ for the transformed tuple. Let $S$ be a function on the class that is
$W_4$-invariant and $SO(4)$-invariant, that is, $S(R\cdot F) = S(F)$ for every $F$ in the class
and every $R \in W_4$, respectively every $R \in SO(4)$. Then $S$ is $O(4)$-invariant: for every
$R\in O(4)$ and every $F$ in the class, $R\cdot F$ lies in the class and $S(R\cdot F)=S(F)$.
\end{lemma}

\begin{proof}
Let $R \in O(4)$. Then $\det R = \pm 1$. If $\det R = 1$, then $R \in SO(4)$, and the assertion
is the assumed $SO(4)$-stability of the class and $SO(4)$-invariance of $S$.

Suppose $\det R = -1$. The matrix $w_0 = \operatorname{diag}(-1,1,1,1)$ is a signed permutation
matrix, so $w_0 \in W_4$. Moreover $w_0^{-1} = w_0$ and $\det w_0 = -1$. Hence
\begin{equation*}
\det\bigl(R w_0^{-1}\bigr) = \det R \cdot \det w_0^{-1} = (-1)(-1) = 1 ,
\end{equation*}
so that $R w_0^{-1} \in SO(4)$, and
\begin{equation*}
R = \bigl(R w_0^{-1}\bigr)\, w_0 .
\end{equation*}
Since $F \mapsto R\cdot F$ is an action of $O(4)$, this factorisation gives
\begin{equation*}
R\cdot F = \bigl(R w_0^{-1}\bigr)\cdot\bigl(w_0\cdot F\bigr).
\end{equation*}
By $W_4$-stability of the class, $w_0\cdot F$ lies in the class, and by $W_4$-invariance of $S$,
$S(w_0\cdot F) = S(F)$. By $SO(4)$-stability of the class, applied to $w_0\cdot F$ and
$R w_0^{-1}\in SO(4)$, the tuple $R\cdot F$ lies in the class. By $SO(4)$-invariance of $S$,
applied to the same pair,
\begin{equation*}
S(R\cdot F) = S\bigl(\bigl(R w_0^{-1}\bigr)\cdot (w_0\cdot F)\bigr) = S(w_0\cdot F) = S(F).
\end{equation*}
This holds for every $R \in O(4)$ and every $F$ in the class, which is the assertion.
\end{proof}

\begin{lemma}[Finite-order elements of a one-parameter group]\label{10.7:lem-finite-order-elements}
Let $\omega_0\in\mathfrak{so}(4)$ with $\omega_0\neq 0$. Then the set
\begin{equation*}
\{\,t\in\mathbb{R} : e^{t\omega_0}\ \text{has finite order in } SO(4)\,\}
\end{equation*}
is countable.
\end{lemma}

\begin{proof}
Write
\begin{equation*}
J=\begin{pmatrix}0&-1\\ 1&0\end{pmatrix},
\qquad
\operatorname{Rot}\theta=\begin{pmatrix}\cos\theta&-\sin\theta\\ \sin\theta&\cos\theta\end{pmatrix},
\end{equation*}
so that $\operatorname{Rot}\theta=e^{\theta J}$, as in the proof of Lemma~\ref{10.7:lem-exponential-onto}.
By the real normal form of a real skew-symmetric matrix, there are $Q\in O(4)$ and
$\theta_1,\theta_2\in\mathbb{R}$ with $\omega_0=Q(\theta_1J\oplus\theta_2J)Q^{\mathsf T}$.
Since $\omega_0\neq 0$, at least one of $\theta_1,\theta_2$ is non-zero; say $\theta_i\neq 0$.
Because $Q^{\mathsf T}=Q^{-1}$, the exponential commutes with conjugation by $Q$, and the
exponential of a block-diagonal matrix is block diagonal, whence for every $t\in\mathbb{R}$
\begin{equation*}
e^{t\omega_0}=Q\bigl(e^{t\theta_1J}\oplus e^{t\theta_2J}\bigr)Q^{\mathsf T}
=Q\bigl(\operatorname{Rot}(t\theta_1)\oplus\operatorname{Rot}(t\theta_2)\bigr)Q^{\mathsf T}.
\end{equation*}
Suppose $e^{t\omega_0}$ has finite order, that is, $(e^{t\omega_0})^n=\mathbb{1}$ for some integer
$n\ge 1$. Since $(e^{t\omega_0})^n=e^{nt\omega_0}$, the display with $nt$ in place of $t$ gives
$\operatorname{Rot}(nt\theta_i)=\mathbb{1}$, hence $nt\theta_i\in 2\pi\mathbb{Z}$, and therefore
$t\theta_i\in\pi\mathbb{Q}$. Thus the set in the statement is contained in
$\theta_i^{-1}\pi\mathbb{Q}$, which is countable.
\end{proof}

\begin{proposition}[Cut-off independence of the boundary term]
\label{10.7:prop-boundary-independence}
Assume the hypotheses $(\mathrm{H}\text{-}\mathrm{A1})_{\mathrm{conn}}$ and
$(\mathrm{H}\text{-}\mathrm{A2}')_{\mathrm{pw}}$ of
Definition~\ref{10.7:def-hypotheses-used-form}, in the form
\eqref{10.7:eq-hypotheses-used-form-b}, with $\omega$ as there, and let $r\ge 2$.
Write $D_j$, $\mathcal{N}^{\chi}_j$ and $\mathcal{B}^{\chi}_j$ for
$D^{\omega}_{K_j,m_j}(\vec h)$, $\mathcal{N}^{\omega,\chi}_{K_j,m_j}(\vec h)$ and
$\mathcal{B}^{\omega,\chi}_{K_j,m_j}(\vec h)$. Then for every
$\vec h\in\mathcal{T}^{(r)}_n(\mathcal{R},\mathsf{d})$ and every $\chi$ admissible at
radius $\mathcal{R}$,
\begin{equation}\label{10.7:eq-boundary-independence-1}
\lim_{j\to\infty}\mathcal{B}^{\omega,\chi}_{K_j,m_j}(\vec h)=c(\vec h,\omega).
\end{equation}
In particular $\lim_j\mathcal{B}^{\chi}_j$ exists and does not depend on $\chi$, and for any
two such cut-offs $\chi,\chi'$ and all sufficiently large $j$,
\begin{equation}\label{10.7:eq-boundary-independence-2}
\mathcal{B}^{\chi}_j-\mathcal{B}^{\chi'}_j
=\mathcal{N}^{\chi'}_j-\mathcal{N}^{\chi}_j\longrightarrow 0 .
\end{equation}
\end{proposition}

\begin{proof}
Fix $\vec h\in\mathcal{T}^{(r)}_n(\mathcal{R},\mathsf{d})$ and $\chi$ admissible at
radius $\mathcal{R}$. By $(\mathrm{H}\text{-}\mathrm{A1})_{\mathrm{conn}}$ there is an
index $j_2=j_2(n,\mathcal{R},\mathsf{d},\omega,\chi)$ such that for every $j\ge j_2$,
that is, beyond the floors, the decomposition $D_j=\mathcal{N}^{\chi}_j+\mathcal{B}^{\chi}_j$
holds. Hence
\begin{equation}\label{10.7:eq-boundary-independence-3}
\mathcal{B}^{\chi}_j=D_j-\mathcal{N}^{\chi}_j\qquad(j\ge j_2).
\end{equation}
Since $r\ge 2$, every real $C^r_c$ function is $C^2_c$, so
$\mathcal{T}^{(r)}_n(\mathcal{R},\mathsf{d})\subset\mathcal{T}^{(2)}_n(\mathcal{R},\mathsf{d})$.
Therefore the convergence of the finite-cutoff rotation derivative
(Lemma~\ref{10.7:lem-derivative-convergence}) applies to $\vec h$ and gives
$D_j\to c(\vec h,\omega)$. Moreover $\mathcal{N}^{\chi}_j\to 0$ as $j\to\infty$.
Letting $j\to\infty$ in \eqref{10.7:eq-boundary-independence-3} gives
\eqref{10.7:eq-boundary-independence-1}.

The limit $c(\vec h,\omega)$ does not involve $\chi$, so $\lim_j\mathcal{B}^{\chi}_j$ is
independent of $\chi$. For a second admissible $\chi'$, with index $j_2'$, the quantity
$D_j$ is defined from $S^T_{n;K_j,m_j}$ and the rotation $R^{\omega}_t$ alone and is the
same for $\chi$ and $\chi'$. Subtracting the two instances of
\eqref{10.7:eq-boundary-independence-3} for $j\ge\max\{j_2,j_2'\}$ yields the identity in
\eqref{10.7:eq-boundary-independence-2}. Its right-hand side tends to zero because both
$\mathcal{N}^{\chi}_j$ and $\mathcal{N}^{\chi'}_j$ do.
\end{proof}

\begin{proposition}[The infrared decay hypothesis forces the anomaly to vanish]
\label{10.7:prop-anomaly-vanishes}
Assume the connected form of the rotational Ward identity and the pointwise form of the
dimension-six coefficient bound, as recorded in Definition~\ref{10.7:def-hypotheses-used-form}
and assumed in Proposition~\ref{10.7:prop-boundary-independence}. Assume in addition the infrared
hypothesis in the pointwise form recorded in Definition~\ref{10.7:def-hypotheses-used-form}
(see \eqref{10.7:eq-hypotheses-used-form-b}). Then
\begin{equation*}
c(\vec h,\omega) = 0 \qquad \text{for every } \vec h \in \mathcal{T}^{(r)}_n(\mathcal{R},\mathsf{d}).
\end{equation*}
\end{proposition}

\begin{proof}
Fix $\vec h \in \mathcal{T}^{(r)}_n(\mathcal{R},\mathsf{d})$. For an admissible $\chi$ and
$j \ge j_2(n,\mathcal{R},\mathsf{d},\omega,\chi)$, the index of
Definition~\ref{10.7:def-hypotheses-used-form}, write
\begin{equation*}
\mathcal{B}^{\chi}_j := \mathcal{B}^{\omega,\chi}_{K_j,m_j}(\vec h);
\end{equation*}
this sequence is defined for all sufficiently large $j$.
By the cut-off independence of the boundary term
(Proposition~\ref{10.7:prop-boundary-independence}), $\mathcal{B}^{\chi}_j \to c(\vec h,\omega)$ as
$j \to \infty$ for every admissible $\chi$. Hence $|\mathcal{B}^{\chi}_j| \to |c(\vec h,\omega)|$, and
\begin{equation*}
\limsup_{j} |\mathcal{B}^{\chi}_j| = |c(\vec h,\omega)|
\qquad \text{for every admissible } \chi .
\end{equation*}
The left-hand side therefore does not depend on $\chi$. By the infrared hypothesis in the pointwise
form of Definition~\ref{10.7:def-hypotheses-used-form} (see
\eqref{10.7:eq-hypotheses-used-form-b}), the infimum of the left-hand side over admissible $\chi$ is
$0$. The infimum of a quantity that does not depend on $\chi$ is that quantity itself. Hence
$|c(\vec h,\omega)| = 0$, that is, $c(\vec h,\omega) = 0$.
\end{proof}

\begin{proposition}[Constancy along rotation orbits, pointwise proof]
\label{10.7:prop-orbit-constancy-pointwise}
Let $\omega\neq 0$, and assume that the three hypotheses in the form
\eqref{10.7:eq-hypotheses-used-form-b} hold for the data
$(\omega,n,\mathcal{R},\mathsf{d})$. Let
$\vec f\in\mathcal{T}^{(r)}_n(\mathcal{R},\mathsf{d})$ and put
\begin{equation*}
\varphi(s):=S^T_n\bigl(R^{\omega}_s\cdot\vec f\,\bigr),\qquad s\in\mathbb{R}.
\end{equation*}
Then $\varphi$ is constant on $\mathbb{R}$. More precisely, $\varphi$ is differentiable on
$\mathbb{R}$ and $\varphi'\equiv 0$.
\end{proposition}

\begin{proof}
Fix $s\in\mathbb{R}$. Since $R^{\omega}_s=e^{s\omega}$ is a rotation,
Lemma~\ref{10.7:lem-rotated-generator} gives
$R^{\omega}_s\cdot\vec f\in\mathcal{T}^{(r)}_n(\mathcal{R},\mathsf{d})$.
The family $(R^{\omega}_t)_{t\in\mathbb{R}}$ is a one-parameter group, so
$R^{\omega}_{s+u}=R^{\omega}_uR^{\omega}_s$. Rotations act on tuples of test functions by a group
action, and therefore, for every $u\in\mathbb{R}$,
\begin{equation*}
\varphi(s+u)=S^T_n\bigl(R^{\omega}_u\cdot(R^{\omega}_s\cdot\vec f\,)\bigr).
\end{equation*}

We apply Lemma~\ref{10.7:lem-limit-differentiability} with base point
$\vec h:=R^{\omega}_s\cdot\vec f$. The time bound of that lemma is
$t_0=\mathsf{d}/(2\|\omega\|\mathcal{R})$. It depends only on $\mathsf{d}$, $\omega$ and
$\mathcal{R}$, and so it is the same for every $s$. For $0<|u|\le t_0$ the lemma gives
\begin{equation}
\label{10.7:eq-orbit-constancy-pointwise-1}
\Bigl|\frac{\varphi(s+u)-\varphi(s)}{u}
-c\bigl(R^{\omega}_s\cdot\vec f,\omega\bigr)\Bigr|
\le C_*\,\Sigma_s(u).
\end{equation}
Here $\Sigma_s(u)$ denotes the bracket on the right-hand side of the estimate of
Lemma~\ref{10.7:lem-limit-differentiability} for the tuple $\vec h=R^{\omega}_s\cdot\vec f$, that is,
the sum of the terms built from the rotation moduli
$\mu_{\delta_\omega h_i}(|u|\|\omega\|)$ and $\mu_{h_j}(|u|\|\omega\|)$. By that lemma,
$\Sigma_s(u)\to 0$ as $u\to 0$. It follows from \eqref{10.7:eq-orbit-constancy-pointwise-1} that
$\varphi$ is differentiable at $s$, with
\begin{equation*}
\varphi'(s)=c\bigl(R^{\omega}_s\cdot\vec f,\omega\bigr).
\end{equation*}

The tuple $R^{\omega}_s\cdot\vec f$ lies in the class, as shown above, and the hypotheses
\eqref{10.7:eq-hypotheses-used-form-b} hold. Proposition~\ref{10.7:prop-anomaly-vanishes}
applied to this tuple therefore gives $c(R^{\omega}_s\cdot\vec f,\omega)=0$. As $s$ was
arbitrary, $\varphi$ is differentiable on $\mathbb{R}$ with $\varphi'\equiv 0$.

Finally, let $s<s'$. The mean value theorem, applied to the real and imaginary parts of
$\varphi$ on $[s,s']$, gives points $\xi_1,\xi_2\in(s,s')$ such that
\begin{equation*}
\varphi(s')-\varphi(s)
=(s'-s)\bigl(\operatorname{Re}\varphi'(\xi_1)
+i\operatorname{Im}\varphi'(\xi_2)\bigr)=0.
\end{equation*}
Hence $\varphi$ is constant on $\mathbb{R}$.
\end{proof}

\begin{proposition}[Constancy along rotation orbits, integrated proof]
\label{10.7:prop-orbit-constancy-integrated}
Assume the three hypotheses \eqref{10.7:eq-hypotheses-used-form-b} of this section. Let
$\omega$ be a real antisymmetric $4\times 4$ matrix and let
$R_t = R^{\omega}_t = e^{t\omega}$, $t\in\mathbb{R}$. Let $\mathcal{R}>0$, $\mathsf{d}>0$ and
$n\ge 2$, and let
$\vec f = (f_1,\dots,f_n) \in \mathcal{T}^{(2)}_n(\mathcal{R},\mathsf{d})$, and write
$R_t\cdot\vec f = (R_t f_1,\dots,R_t f_n)$. Then, for every $t\in\mathbb{R}$,
\begin{equation}\label{10.7:eq-orbit-constancy-integrated-1}
S^T_n(R_t\cdot\vec f) = S^T_n(\vec f).
\end{equation}
\end{proposition}

These are the hypotheses and the conclusion of the preceding proposition of this section, on
constancy along rotation orbits; the proof below is independent of the proof given there and
integrates the derivative along the orbit.

\begin{proof}
For each $j$ put $\varphi_j(t) = S_j(R_t\cdot\vec f)$, where
$S_j = S^T_{n;K_j,m_j}$ along the diagonal. Let $D_j = D^{\omega}_{K_j,m_j}$ be the derivative
along the orbit. We consider $j$ so large that $L^{m_j}\ge \mathcal{R}/2$, which is what
Lemma~\ref{10.7:lem-finite-cutoff-structure}(c) requires. For such $j$, that lemma gives that
$\varphi_j$ is $C^1$ on $\mathbb{R}$ and that
\begin{align*}
\varphi_j'(s) &= D_j(R_s\cdot\vec f)\\
&= \sum_{i=1}^n S_j\bigl(R_s f_1,\dots,R_s\delta_\omega f_i,\dots,R_s f_n\bigr),
\end{align*}
with $R_s\delta_\omega f_i = \delta_\omega(R_s f_i)$.

\emph{Uniform bound.} Fix $s$. Since $\vec f \in \mathcal{T}^{(2)}_n(\mathcal{R},\mathsf{d})$ and
$R_s$ is an isometry fixing the origin, the sets $R_s(\operatorname{supp} f_k)$, $k=1,\dots,n$,
have two properties: they lie in $B(0,\mathcal{R}/2)$, and they are pairwise at distance at
least $\mathsf{d}$.

The function $\delta_\omega f_i = -(\omega x)\cdot\nabla f_i$ is real, compactly supported and
Lipschitz, because $f_i$ is $C^2$ with compact support. Moreover
$\operatorname{supp}\delta_\omega f_i \subset \operatorname{supp} f_i$. Hence the $i$-th tuple
in the sum above satisfies the hypotheses of the equicontinuity estimate,
Lemma~\ref{10.7:lem-equicontinuity}(b), with $\rho = \mathcal{R}/2$ and
$\mathsf{d}' = \mathsf{d}$.

The sup norm and the Lipschitz constant are invariant under composition with $R_s^{-1}$. Hence
$\|R_s h\|_{\sharp} = \|h\|_{\sharp}$ for every compactly supported Lipschitz $h$. The first of
the bounds in
\eqref{10.7:eq-equicontinuity-a}, applied to each term of the sum, therefore gives, for
$j\ge j_1(n,\mathcal{R}/2)$,
\begin{equation}\label{10.7:eq-orbit-constancy-integrated-2}
\bigl|D_j(R_s\cdot\vec f)\bigr|
\le C^\infty_n(\mathsf{d};\mathcal{R}/2)\sum_{i=1}^n \|\delta_\omega f_i\|_{\sharp}
\prod_{k\ne i}\|f_k\|_{\sharp}.
\end{equation}
Neither $j_1(n,\mathcal{R}/2)$ nor $C^\infty_n(\mathsf{d};\mathcal{R}/2)$ depends on the position
of the sets $R_s(\operatorname{supp} f_k)$ inside $B(0,\mathcal{R}/2)$. So the right side
of \eqref{10.7:eq-orbit-constancy-integrated-2} is a constant independent of $s$ and of $j$.

\emph{Pointwise limit.} For each $s$, the tuple $R_s\cdot\vec f$ again lies in
$\mathcal{T}^{(2)}_n(\mathcal{R},\mathsf{d})$, by the two properties of the sets
$R_s(\operatorname{supp} f_k)$ above.
Lemma~\ref{10.7:lem-derivative-convergence}, applied at $R_s\cdot\vec f$, gives
\begin{equation*}
\lim_{j\to\infty} D_j(R_s\cdot\vec f) = c(R_s\cdot\vec f,\omega).
\end{equation*}
Under the third of the hypotheses \eqref{10.7:eq-hypotheses-used-form-b} the anomaly vanishes
(Proposition~\ref{10.7:prop-anomaly-vanishes}), so $c(R_s\cdot\vec f,\omega) = 0$. Thus
$\varphi_j'(s)\to 0$ for every $s$.

\emph{Conclusion.} By the fundamental theorem of calculus,
\begin{equation*}
\varphi_j(t)-\varphi_j(0) = \int_0^t \varphi_j'(s)\,ds .
\end{equation*}
Consider the bounded interval between $0$ and $t$. For $j$ large enough that
$L^{m_j}\ge\mathcal{R}/2$ and $j\ge j_1(n,\mathcal{R}/2)$, the integrands $\varphi_j'$ are
continuous there, bounded by the constant in \eqref{10.7:eq-orbit-constancy-integrated-2}, and
tend to $0$ pointwise. By dominated
convergence the integral tends to $0$ as $j\to\infty$.

On the other hand, by the infinite-volume theorem, in the form stated in
Lemma~\ref{10.7:lem-equicontinuity}(a), we have $\varphi_j(t)\to S^T_n(R_t\cdot\vec f)$ and
$\varphi_j(0)\to S^T_n(\vec f)$. Hence
\begin{equation*}
S^T_n(R_t\cdot\vec f)-S^T_n(\vec f) = 0,
\end{equation*}
which is \eqref{10.7:eq-orbit-constancy-integrated-1}.
\end{proof}

\begin{remark}[No interchange of derivative and limit]\label{10.7:rem-no-interchange}
Let $\omega$, $n$, $\mathcal{R}$, $\mathsf{d}$ and $\vec f\in\mathcal{T}^{(r)}_n(\mathcal{R},\mathsf{d})$ be as in
Proposition~\ref{10.7:prop-orbit-constancy-pointwise}, and write $R_s=R^{\omega}_s=e^{s\omega}$. Let
$\mathcal{N}_j$ and $\mathcal{B}_j$ denote the connected ultraviolet and infrared insertions along the
diagonal sequence indexed by $j$. For $t\in\mathbb{R}$, passing to the limit $j\to\infty$ termwise in
the two integrals
\begin{equation*}
\int_0^t \mathcal{N}_j(R_s\cdot\vec f)\,ds
\qquad\text{and}\qquad
\int_0^t \mathcal{B}_j(R_s\cdot\vec f)\,ds
\end{equation*}
would require two things: first,
\begin{equation}\label{10.7:eq-no-interchange-1}
\sup_{s,j}\,\bigl|\mathcal{N}_j(R_s\cdot\vec f)\bigr|\to 0 ,
\end{equation}
and second, an integrated infrared statement. Propositions~\ref{10.7:prop-orbit-constancy-pointwise}
and~\ref{10.7:prop-orbit-constancy-integrated} require strictly less than this, and in neither of them
is the derivative $d/dt$ interchanged with the limit $\lim_j$.
\end{remark}

\begin{proof}
Requirement \eqref{10.7:eq-no-interchange-1} is automatic for the following reason. Each $R_s$ acts
on test functions through an orthogonal map of $\mathbb{R}^4$, and composition with an isometry changes
neither the supremum norm nor the Lipschitz constant of a function. Since
$\|f\|_{\sharp}=\max\{\|f\|_\infty,\,40M\operatorname{Lip} f\}$, this gives
\begin{equation*}
\|R_s f_i\|_{\sharp}=\|f_i\|_{\sharp}\qquad\text{for every } s \text{ and every } i .
\end{equation*}
The estimate for the connected ultraviolet insertion along the diagonal bounds
$|\mathcal{N}_j(\vec h)|$, for a tuple $\vec h$ of separated test functions, by a quantity that tends to
zero as $j\to\infty$ and depends on $\vec h$ only through the norms $\|h_i\|_{\sharp}$, with constants
that do not depend on the positions of the supports of the $h_i$. Applied to the tuple
$\vec h=R_s\cdot\vec f$, this bound depends on $s$ only through the norms $\|R_s f_i\|_{\sharp}$, which
by the identity just displayed equal $\|f_i\|_{\sharp}$. Hence $|\mathcal{N}_j(R_s\cdot\vec f)|$ is
bounded, uniformly in $s$, by a quantity that tends to zero as $j\to\infty$; this is requirement
\eqref{10.7:eq-no-interchange-1}. Thus the estimates used along the diagonal are uniform in the orbit
parameter $s$.
\end{proof}

\begin{proposition}[Extension of invariance from smooth to continuously differentiable tuples]
\label{10.7:prop-density-extension}
Let $S$ be either $S^T_n$ or a limit obtained in the simultaneous subsequential convergence
lemma (Lemma~\ref{10.7:lem-simultaneous-convergence}), and let $R$ be given. Suppose that
\begin{equation*}
S(R\cdot\vec g)=S(\vec g)
\qquad\text{for all }
\vec g\in\bigcup_{\mathcal{R},\mathsf{d}}\mathcal{T}^{(r)}_n(\mathcal{R},\mathsf{d}).
\end{equation*}
Then $S(R\cdot\vec f)=S(\vec f)$ holds, for the same $R$, for every admissible $C^1_c$ tuple
$\vec f=(f_1,\dots,f_n)$.
\end{proposition}

\begin{proof}
Let $\vec f$ be an admissible $C^1_c$ tuple of real functions, and let $\mathsf{d}_0>0$ be the
least distance between the supports of two distinct entries. Fix a real function
$\psi\in C^\infty_c(\mathbb{R}^4)$ with $\psi\ge 0$, $\int\psi=1$ and support in the closed unit
ball, and put $\psi_\delta(x):=\delta^{-4}\psi(x/\delta)$. For $0<\delta\le\mathsf{d}_0/4$ set
$f_i^{\delta}:=f_i*\psi_\delta$. Each $f_i^{\delta}$ is real and lies in $C^\infty_c$, hence in
$C^r_c$.

Supports. Since $\operatorname{supp}\psi_\delta$ lies in the closed ball of radius
$\delta\le\mathsf{d}_0/4$, we have $\operatorname{supp}f_i^{\delta}\subset A_i$, where
$A_i:=(\operatorname{supp}f_i)^{+\mathsf{d}_0/4}$ is the closed $\mathsf{d}_0/4$-neighbourhood.
Two of these sets are at distance at least $\mathsf{d}_0-2\cdot\mathsf{d}_0/4=\mathsf{d}_0/2$.
Choosing $\rho$ with $\operatorname{supp}f_i\subset B(0,\rho-\mathsf{d}_0/4)$ for all $i$, all
$A_i$ lie in the fixed ball $B(0,\rho)$, independent of $\delta$. Hence
$\vec f^{\,\delta}:=(f_1^{\delta},\dots,f_n^{\delta})$ belongs to
$\mathcal{T}^{(r)}_n(\mathcal{R},\mathsf{d})$ for a suitable pair $(\mathcal{R},\mathsf{d})$ not
depending on $\delta$, and the hypothesis gives $S(R\cdot\vec f^{\,\delta})=S(\vec f^{\,\delta})$.

Convergence. Since $\psi_\delta$ is a probability density supported in the ball of radius
$\delta$, $|f_i^{\delta}(x)-f_i(x)|\le m_{f_i}(\delta)$ for every $x$. Moreover
$\nabla f_i^{\delta}=(\nabla f_i)*\psi_\delta$, so
$\sup|\nabla f_i^{\delta}-\nabla f_i|\le m_{\nabla f_i}(\delta)$, and this bounds the Lipschitz
constant of the $C^1$ function $f_i^{\delta}-f_i$. Therefore
\begin{equation*}
\|f_i^{\delta}-f_i\|_{\sharp}
\le\max\{m_{f_i}(\delta),\,40M\,m_{\nabla f_i}(\delta)\}\longrightarrow 0
\qquad(\delta\to 0),
\end{equation*}
since $f_i$ and $\nabla f_i$ are continuous with compact support, hence uniformly continuous.
By the rotation stability of the hypothesis classes (Lemma~\ref{10.7:lem-class-stability}), and
since $f\mapsto R\cdot f$ is linear,
\begin{equation*}
\|R\cdot f_i^{\delta}-R\cdot f_i\|_{\sharp}
=\|R\cdot(f_i^{\delta}-f_i)\|_{\sharp}
=\|f_i^{\delta}-f_i\|_{\sharp}.
\end{equation*}

Passage to the limit. On the unrotated side, $f_i^{\delta}$ and $f_i$ are Lipschitz with support
in $A_i$, and the containers $A_1,\dots,A_n$ lie in $B(0,\rho)$ and are pairwise at distance at
least $\mathsf{d}_0/2$. The equicontinuity estimate \eqref{10.7:eq-equicontinuity-a} for $S$
(Lemma~\ref{10.7:lem-equicontinuity}) therefore gives $S(\vec f^{\,\delta})\to S(\vec f)$ as
$\delta\to 0$. On the rotated side, Lemma~\ref{10.7:lem-class-stability} shows that
$R\cdot f_i^{\delta}$ and $R\cdot f_i$ are real, Lipschitz and compactly supported, with
supports in $R(A_i)$. The containers $R(A_1),\dots,R(A_n)$ are again pairwise at distance at
least $\mathsf{d}_0/2$ and lie in $B(0,\rho)$, because these distances and this ball are
preserved. The same estimate, applied with these containers, gives
$S(R\cdot\vec f^{\,\delta})\to S(R\cdot\vec f)$. Hence
\begin{equation*}
S(R\cdot\vec f)=\lim_{\delta\to 0}S(R\cdot\vec f^{\,\delta})
=\lim_{\delta\to 0}S(\vec f^{\,\delta})=S(\vec f).
\qedhere
\end{equation*}
\end{proof}

\begin{theorem}[The conclusion step]\label{10.7:thm-conclusion-step}
Let $N=2$, so that the gauge group is $SU(2)$. Assume the following.
\begin{enumerate}
\item Hypotheses (H1)--(H7) hold with the quantifier prefix of clause (iii)(b) of the main
theorem. That is, $L\ge L_0(N)$ is odd; the auxiliary parameters $\mathfrak{p}$ and then
$\gamma$, $\gamma_H$ are chosen in the fixed order of choice; $\varrho_J>0$;
$\gamma_J(L,\mathfrak{p},\gamma;\varrho_J)\le\gamma_H$; and $g\in\,]0,\gamma_J]$.
\item The sequence $g_0^{(j)}\downarrow 0$ of clause (ii-b) is given, with
$K_j=K(g_0^{(j)},g)$. The torus exponents satisfy $m_j\to\infty$ and
$m_j\ge m_{\mathbb{T}}(g_-(g))$. All limits are taken along the joint diagonal of clause
(ii-e) and of the infinite-volume theorem.
\item The per-ball volume-free bound and the infinite-volume theorem hold.
\item The lemma on the invariances of the Wilson measure holds.
\item There is some $r\ge 2$ such that, for all $\omega\in\mathfrak{so}(4)$, $n\ge 2$,
$\mathcal{R}>0$ and $\mathsf{d}>0$, three hypotheses hold along the joint diagonal. These
are the connected rotational Ward identity; the pointwise dimension-six hypothesis, or else
either of the two forms of the dimension-six coefficient bound stated with the hypotheses of
the main theorem, from which the pointwise form follows by part (a) of
Lemma~\ref{10.7:lem-hypotheses-reduction}; and the far-sphere flux hypothesis
$(\mathrm{IR}\text{-}\mathrm{fl})_{\omega_0}$ of the main theorem, taken at $\omega_0=\omega$
and in its pointwise form: the flux of the lattice rotation current through far spheres
vanishes along one sequence of radii. A sufficient condition for it is that the connected
correlation of one far-away stress-tensor insertion with the observables decays faster than
$|x|^{-4}$, uniformly in the cutoff. The connected rotational Ward identity, the pointwise
dimension-six hypothesis and the pointwise far-sphere flux hypothesis are those of
Definition~\ref{10.7:def-hypotheses-used-form}, display~\eqref{10.7:eq-hypotheses-used-form-b}.
\item Suppose the pointwise dimension-six hypothesis is not assumed directly but is derived
from one of these two forms by part (a) of
Lemma~\ref{10.7:lem-hypotheses-reduction}. Then clause (0) of the
main theorem is assumed in addition: $K_j\to\infty$ along the sequence of clause (ii-b).
\end{enumerate}
Let $(S^T_n)_{n\ge 2}$ be any subsequential limit family of the infinite-volume theorem, and
let $n\ge 2$ and $R\in O(4)$. Let $f_1,\dots,f_n\in C^1_c(\mathbb{R}^4)$ be real functions
with pairwise disjoint supports. Then
\begin{equation}\label{10.7:eq-conclusion-step-1}
S^T_n(R\cdot f_1,\dots,R\cdot f_n)=S^T_n(f_1,\dots,f_n).
\end{equation}
Equivalently, together with the translation invariance contained in the infinite-volume
theorem, $S^T_n$ is invariant under $\mathbb{R}^4\rtimes O(4)$. The same statement holds
verbatim for any limit $S''$ along a further subsequence, as provided by
Lemma~\ref{10.7:lem-simultaneous-convergence}.

No constant enters the conclusion. The list of what the infinite-volume theorem does not
assert stands, except in so far as an item of that list is the Euclidean invariance
\eqref{10.7:eq-conclusion-step-1}, which is asserted here under hypotheses 5 and 6. In
particular, the statement concerns every subsequential limit family, and it is not asserted
that the limit family is unique.
\end{theorem}

\begin{proof}
Call a tuple $\vec f=(f_1,\dots,f_n)$ of real compactly supported functions with pairwise
disjoint supports \emph{admissible}. We first prove \eqref{10.7:eq-conclusion-step-1} for
admissible $C^r_c$ tuples and $R\in SO(4)$. We then extend it to admissible $C^1_c$ tuples,
and finally to $R\in O(4)$.

\emph{Smooth tuples in a class $\mathcal{T}^{(r)}_n(\mathcal{R},\mathsf{d})$.} Let
$\vec g=(g_1,\dots,g_n)$ be an admissible $C^r_c$ tuple. The union of the supports is
compact, so we may choose $\mathcal{R}>0$ with
$\bigcup_i\operatorname{supp}g_i\subset B(0,\mathcal{R}/2)$. Let $\mathsf{d}$ be the minimal
separation of the supports,
\begin{equation*}
\mathsf{d}:=\min_{i\neq k}\operatorname{dist}(\operatorname{supp}g_i,\operatorname{supp}g_k).
\end{equation*}
It is positive because the supports are compact and pairwise disjoint. Hence
$\vec g\in\mathcal{T}^{(r)}_n(\mathcal{R},\mathsf{d})$.

\emph{Rotations in $SO(4)$.} Let $Q\in SO(4)$. By surjectivity of the exponential map of the
rotation group (Lemma~\ref{10.7:lem-exponential-onto}) there is $\omega\in\mathfrak{so}(4)$
with $Q=e^{\omega}$.

If $\omega=0$, then $Q$ is the identity and there is nothing to prove.

If $\omega\neq 0$, the connected rotational Ward identity, the pointwise dimension-six
hypothesis and the pointwise far-sphere flux hypothesis hold for the data
$(\omega,n,\mathcal{R},\mathsf{d})$, since they are assumed for all such data. If the
dimension-six hypothesis is assumed only in one of the two forms stated with the hypotheses
of the main theorem, its pointwise form
follows by part (a) of
Lemma~\ref{10.7:lem-hypotheses-reduction}, with clause (0) available as assumed. The pointwise
proof of constancy along rotation orbits
(Proposition~\ref{10.7:prop-orbit-constancy-pointwise}) then shows that
$\varphi(s):=S^T_n(R^{\omega}_s\cdot\vec g)$ is constant in $s\in\mathbb{R}$. The independent
integrated proof (Proposition~\ref{10.7:prop-orbit-constancy-integrated}) gives the same
conclusion. Since $R^{\omega}_1=e^{\omega}=Q$ and $R^{\omega}_0$ is the identity,
\begin{equation}\label{10.7:eq-conclusion-step-2}
S^T_n(Q\cdot\vec g)=\varphi(1)=\varphi(0)=S^T_n(\vec g).
\end{equation}

\emph{Extension to $C^1_c$ tuples.} The tuple $\vec g$ was an arbitrary element of
$\bigcup_{\mathcal{R},\mathsf{d}}\mathcal{T}^{(r)}_n(\mathcal{R},\mathsf{d})$, and $Q$ was
fixed. So \eqref{10.7:eq-conclusion-step-2} holds for this $Q$ and all such tuples. The
extension of invariance from smooth to continuously differentiable tuples
(Proposition~\ref{10.7:prop-density-extension}) therefore gives
$S^T_n(Q\cdot\vec f)=S^T_n(\vec f)$ for this $Q$ and every admissible $C^1_c$ tuple $\vec f$.
As $Q\in SO(4)$ was arbitrary, $S^T_n$ is $SO(4)$-invariant on the class of admissible
$C^1_c$ tuples.

\emph{Reflections.} An orthogonal linear change of variables preserves reality,
continuous differentiability and compact support. It also maps pairwise disjoint supports to
pairwise disjoint supports. Hence the class of admissible $C^1_c$ tuples is stable under
$W_4$ and under $SO(4)$. On this class $S^T_n$ is exactly $W_4$-invariant, by part (c) of
Lemma~\ref{10.7:lem-equicontinuity}. The passage from special orthogonal to orthogonal
invariance (Lemma~\ref{10.7:lem-orthogonal-invariance}) now applies. If $\det R=-1$, then
\begin{equation*}
R=(Rw_0^{-1})\,w_0,\qquad Rw_0^{-1}\in SO(4),\qquad w_0\in W_4,
\end{equation*}
and $S^T_n$ is invariant under both factors. This proves \eqref{10.7:eq-conclusion-step-1}
for every $R\in O(4)$ and every admissible $C^1_c$ tuple.

Invariance under $\mathbb{R}^4\rtimes O(4)$ is invariance under translations together with
invariance under $O(4)$. Translation invariance is contained in the infinite-volume theorem,
which gives the equivalent formulation. The argument above applies verbatim to any limit
$S''$ along a further subsequence as in Lemma~\ref{10.7:lem-simultaneous-convergence}, with
$S''$ in place of $S^T_n$. This proves the theorem.

\emph{Order of limits.} For orientation we record the order in which the limits in the
argument above, and in the propositions it invokes, are taken. Fix $n$, $\mathcal{R}$,
$\mathsf{d}$, $\omega$,
a cut-off $\chi$ admissible at radius $\mathcal{R}$, and
$\vec h\in\mathcal{T}^{(r)}_n(\mathcal{R},\mathsf{d})$. Write
\begin{equation*}
D_j:=D^{\omega}_{K_j,m_j}(\vec h),\qquad
\mathcal{N}^{\chi}_j:=\mathcal{N}^{\omega,\chi}_{K_j,m_j}(\vec h),\qquad
\mathcal{B}^{\chi}_j:=\mathcal{B}^{\omega,\chi}_{K_j,m_j}(\vec h),
\end{equation*}
and $c=c(\vec h,\omega)$ for the rotational anomaly. The indices are restricted to
\begin{equation*}
j\ge\max\{\,j_1(n,\mathcal{R}/2),\ j_2,\ j_3\,\}.
\end{equation*}
Here $j_1(n,\mathcal{R}/2)$ is the index of Lemma~\ref{10.7:lem-equicontinuity}, and
$j_2=j_2(n,\mathcal{R},\mathsf{d},\omega,\chi)$ is the index of the connected rotational Ward
identity. The index $j_3$ is the floor of the dimension-six bound. The limits are then taken
in the following order.
\begin{enumerate}
\item First $j\to\infty$. Then $D_j\to c$ and $\mathcal{N}^{\chi}_j\to 0$. Since
$D_j=\mathcal{N}^{\chi}_j+\mathcal{B}^{\chi}_j$, it follows that $\mathcal{B}^{\chi}_j\to c$.
\item Next, the infimum over admissible $\chi$ is taken; this step gives $c=0$. By the
corollary to the vanishing of the rotational anomaly in this chapter (the equivalence of the
infimum over admissible $\chi$ with the limit $\mathcal{R}'\to\infty$), taking this
infimum is the same as letting $\mathcal{R}'\to\infty$ in the form of the hypotheses stated
with the main theorem.
\item Then the orbit parameters are released: $s\in\mathbb{R}$ in
Proposition~\ref{10.7:prop-orbit-constancy-pointwise} (respectively $t\in\mathbb{R}$ in
Proposition~\ref{10.7:prop-orbit-constancy-integrated}) is arbitrary, by constancy along
rotation orbits, and
$R\in SO(4)$ is arbitrary, by surjectivity of the exponential map. The remaining coset of
$O(4)$ is reached through $w_0$.
\item Finally, the mollification parameter $\delta\to 0$ in
Proposition~\ref{10.7:prop-density-extension}.
\end{enumerate}
\end{proof}

\begin{corollary}[The dichotomy for the far-sphere flux hypothesis]\label{10.7:cor-clustering-dichotomy}
Of the three hypotheses displayed in \eqref{10.7:eq-hypotheses-used-form-b}, assume the first two;
the third, the far-sphere flux hypothesis $(\mathrm{IR}\text{-}\mathrm{fl})_{\omega_0}$ (the flux of the
lattice rotation current through far spheres vanishes along one sequence of radii), is not assumed.
Fix the diagonal sequence $(K_j,m_j)$ and a subsequential limit
family $(S^T_n)_{n\ge2}$ along it, as in the hypotheses of
Theorem~\ref{10.7:thm-conclusion-step}: $K_j=K(g_0^{(j)},g)$ for a sequence
$g_0^{(j)}\downarrow0$, and $m_j\to\infty$ with $m_j\ge m_{\mathbb{T}}(g_-(g))$.

For each fixed $(\omega,n,\mathcal{R},\mathsf{d})$, write
$\mathcal{B}^{\chi}_j(\vec h)=\mathcal{B}^{\omega,\chi}_{K_j,m_j}(\vec h)$; with
$\mathcal{T}^{(r)}_n(\mathcal{R},\mathsf{d})$, $r\ge2$, the class of
Proposition~\ref{10.7:prop-boundary-independence}, the following are equivalent.
\begin{enumerate}
\item[(i)] The far-sphere flux hypothesis $(\mathrm{IR}\text{-}\mathrm{fl})_{\omega_0}$, in the form
displayed in \eqref{10.7:eq-hypotheses-used-form-b}, holds at the fixed
$(\omega,n,\mathcal{R},\mathsf{d})$ for every $\vec h$ of the class; here and in (ii) the hypothesis
is meant only in that displayed form.
\item[(ii)] For every $\vec h$ of the class, the double limit in the far-sphere flux hypothesis
vanishes for one (equivalently, every) admissible family of radial cut-offs $(\chi_{\mathcal{R}'})$,
in the sense of that hypothesis:
\begin{equation*}
\lim_{\mathcal{R}'}\ \limsup_{j}\ \bigl|\mathcal{B}^{\chi_{\mathcal{R}'}}_j(\vec h)\bigr| = 0 .
\end{equation*}
\item[(iii)] $\lim_j \mathcal{B}^{\chi}_j(\vec h)=0$ for every $\vec h$ of the class, for one
(equivalently, every) admissible $\chi$.
\item[(iv)] $c(\vec h,\omega)=0$ for every $\vec h$ of the class
$\mathcal{T}^{(r)}_n(\mathcal{R},\mathsf{d})$.
\end{enumerate}
Taken over all $\omega$, $n$, $\mathcal{R}$, $\mathsf{d}$, these conditions are equivalent to
\begin{enumerate}
\item[(v)] invariance of the limit family $(S^T_n)_{n\ge2}$ under $O(4)$, that is, the conclusion of
Theorem~\ref{10.7:thm-conclusion-step} for this subsequential limit.
\end{enumerate}
Without the far-sphere flux hypothesis the argument yields exactly the following, and nothing more:
the rotation
current is conserved in the limit, and its anomaly $c$ equals the $\chi$-independent flux
$\lim_j\mathcal{B}^{\chi}_j$ through any admissible shell.
\end{corollary}

\begin{proof}
Fix $(\omega,n,\mathcal{R},\mathsf{d})$ and let $\vec h\in\mathcal{T}^{(r)}_n(\mathcal{R},\mathsf{d})$.
Everything rests on the cut-off independence of the boundary term
(Proposition~\ref{10.7:prop-boundary-independence}). It uses only the first two hypotheses of
\eqref{10.7:eq-hypotheses-used-form-b}, and gives, for every admissible $\chi$,
\begin{equation}\label{10.7:eq-clustering-dichotomy-1}
\lim_{j}\mathcal{B}^{\chi}_j(\vec h)=c(\vec h,\omega),
\end{equation}
with a limit that exists and does not depend on $\chi$.

\emph{(iii)$\Leftrightarrow$(iv).} By \eqref{10.7:eq-clustering-dichotomy-1}, the limit
$\lim_j\mathcal{B}^{\chi}_j(\vec h)$ vanishes for one admissible $\chi$ if and only if
$c(\vec h,\omega)=0$. For the same reason it vanishes for one admissible $\chi$ if and only if it
vanishes for every admissible $\chi$. Letting $\vec h$ range over the class gives the equivalence.

\emph{(ii)$\Leftrightarrow$(iv).} Since the limit in \eqref{10.7:eq-clustering-dichotomy-1} exists,
\begin{equation*}
\limsup_j\bigl|\mathcal{B}^{\chi_{\mathcal{R}'}}_j(\vec h)\bigr|=|c(\vec h,\omega)|
\end{equation*}
for each member $\chi_{\mathcal{R}'}$ of any admissible family. The outer limit over $\mathcal{R}'$ is
therefore the limit of a constant sequence and equals $|c(\vec h,\omega)|$, whatever the family.
Hence the double limit vanishes for one admissible family if and only if it vanishes for every
admissible family, and this happens if and only if $c(\vec h,\omega)=0$. Letting $\vec h$ range over
the class gives the equivalence.

\emph{(i)$\Leftrightarrow$(ii).} Add the far-sphere flux hypothesis, in the form
\eqref{10.7:eq-hypotheses-used-form-b}, to the hypotheses. By
Proposition~\ref{10.7:prop-anomaly-vanishes}, $c(\vec h,\omega)=0$ on the class, so (iv) holds, and
therefore (ii) holds by the preceding paragraph. Conversely, assume (ii). By the preceding paragraph
(iv) holds, and then \eqref{10.7:eq-clustering-dichotomy-1} gives
\begin{equation*}
\limsup_j\bigl|\mathcal{B}^{\chi}_j(\vec h)\bigr|=|c(\vec h,\omega)|=0
\end{equation*}
for every admissible $\chi$ and every $\vec h$ of the class. Hence every infimum or limit over
admissible cut-offs formed from these quantities vanishes; this is (i) at the fixed
$(\omega,n,\mathcal{R},\mathsf{d})$.

\emph{(iv)$\Rightarrow$(v).} Suppose that (iv) holds for all $\omega$, $n$, $\mathcal{R}$,
$\mathsf{d}$. By the equivalence of (iv) and (i) established above for each such datum, the
far-sphere flux hypothesis holds, in the form displayed in
\eqref{10.7:eq-hypotheses-used-form-b}, for all $\omega$, $n$, $\mathcal{R}$, $\mathsf{d}$. Together
with the standing hypotheses of this corollary, the hypotheses of the conclusion step,
Theorem~\ref{10.7:thm-conclusion-step}, are then in force in the form in which that theorem uses
them, and its conclusion, applied to the subsequential limit $(S^T_n)_{n\ge2}$, is (v).

\emph{(v)$\Rightarrow$(iv).} Suppose that (v) holds, so that the limit Schwinger functions are
invariant under $O(4)$. For $\omega=0$ there is nothing to prove: $\delta_0h_i=0$ for every $i$, so
each term of $c(\vec h,0)$ vanishes by linearity of $S^T_n$ in each argument. Let $\omega\ne0$.
Since $r\ge2$, $\vec h\in\mathcal{T}^{(2)}_n(\mathcal{R},\mathsf{d})$ and
Lemma~\ref{10.7:lem-limit-differentiability} applies; let
$t_0$ be as there, and put $R_t=R^{\omega}_t=e^{t\omega}$.
Invariance makes the function
$t\mapsto S^T_n(R_t\cdot\vec h)$ constant, so the difference quotient
\begin{equation*}
\frac{S^T_n(R_t\cdot\vec h)-S^T_n(\vec h)}{t}
\end{equation*}
vanishes for $0<|t|\le t_0$. The differentiability of the limit along one-parameter groups
(Lemma~\ref{10.7:lem-limit-differentiability}) bounds the distance of this quotient from
$c(\vec h,\omega)$ by a quantity that tends to zero as $t\to0$. Hence the derivative of
$t\mapsto S^T_n(R_t\cdot\vec h)$ at $t=0$ exists and equals $c(\vec h,\omega)$. Since the function is
constant, this derivative is zero, so $c(\vec h,\omega)=0$ for every $\vec h$ of the class. This is
(iv), for every $\omega$, $n$, $\mathcal{R}$, $\mathsf{d}$.

\emph{The case without the far-sphere flux hypothesis.} The proof of
Proposition~\ref{10.7:prop-boundary-independence} uses only the first two hypotheses of
\eqref{10.7:eq-hypotheses-used-form-b}, and the diagonal. For $j$ beyond the floor
indices of that proof, it writes
\begin{equation*}
\mathcal{B}^{\chi}_j(\vec h)=D^{\omega}_{K_j,m_j}(\vec h)-\mathcal{N}^{\omega,\chi}_{K_j,m_j}(\vec h),
\end{equation*}
the derivative along the orbit minus the connected ultraviolet insertion, and shows that this
insertion tends to zero; the rotation current is conserved in the limit. The same proposition gives
\eqref{10.7:eq-clustering-dichotomy-1}, which identifies the anomaly $c(\vec h,\omega)$ with the
flux $\lim_j\mathcal{B}^{\chi}_j(\vec h)$ through an admissible shell, and shows that this flux does
not depend on $\chi$. None of the implications above produces more than this without the
far-sphere flux hypothesis.
\end{proof}

\begin{remark}[The theorem for each subsequential limit]
\label{10.7:rem-each-subsequential-limit}
The conclusion step, Theorem~\ref{10.7:thm-conclusion-step}, holds for every subsequential limit
of the joint diagonal $(K_j,m_j)$ of clause~(ii-e), in the form (b2); no two such limits are
compared, and no uniqueness statement is used. In particular, clause~(iv) of Theorem~0, in the
form of the corollary of the uniqueness theorem, is not used, and it would not remove the word
\emph{subsequential} from Theorem~\ref{10.7:thm-conclusion-step}.
\end{remark}

\begin{proof}
Fix one subsequential limit family $(S^T_n)_{n\ge 2}$ of (b2), together with its own diagonal
$(K_j,m_j)$. This family is exactly $W_4$-invariant. Moreover, the three hypotheses in the form
\eqref{10.7:eq-hypotheses-used-form-a} hold at every $(K_j,m_j)$ of this diagonal. These are the
only inputs that Theorem~\ref{10.7:thm-conclusion-step} requires of the family. Its conclusion
therefore holds for this family by itself, with no reference to any other subsequential limit.
Hence no two limits are compared, and no uniqueness statement enters.

We turn to the second assertion. Clause~(iv), in the form of the corollary of the uniqueness
theorem, would give two things: subsequence-free limits in the sense of clause~(iii)(b) on
each fixed torus $\mathbb{T}_m$, and the identification of the subsequential limits taken along
the first-passage cutoffs $K(g_0,g)$ with some $S^{(\hat w)}$. It is not used in this subsection.
Even if it were available, it would not remove the word \emph{subsequential} here. The limit in
(b2) is the joint limit $(K_j,m_j)\to\infty$ on $\mathbb{R}^4$, and the uniqueness of this joint
limit is among the statements that (b2) does not assert. This concerns uniqueness both in $m$
and in the cutoff on $\mathbb{R}^4$.
\end{proof}

\begin{corollary}[Invariance along the full sequence]\label{10.7:cor-full-sequence}
Let $m_j\to\infty$, $g_0^{(j)}\downarrow 0$ and $K_j=K(g_0^{(j)},g)$, and for $n\ge 2$ write
$S_j=S^T_{n;K_j,m_j}$. Suppose that the three hypotheses displayed in
\eqref{10.7:eq-hypotheses-used-form-b} hold as follows: the rotational Ward identity hypothesis at
every $(K_j,m_j)$ with $j$ beyond its floor indices, and the dimension-six coefficient bound and the
far-sphere flux hypothesis $(\mathrm{IR}\text{-}\mathrm{fl})_{\omega_0}$, in their pointwise form,
along the full sequence. Then, for every
$n\ge 2$, every $R\in O(4)$ and every admissible tuple $\vec f$ of $C^1_c$ test functions,
\begin{equation*}
S_j(R\cdot\vec f)-S_j(\vec f)\longrightarrow 0\qquad (j\to\infty).
\end{equation*}
\end{corollary}

\begin{proof}
Suppose the conclusion fails. Then there are $n\ge 2$, $R\in O(4)$, an admissible $C^1_c$ tuple
$\vec f$, a number $\eta>0$ and an infinite set $J'$ of indices such that
\begin{equation}\label{10.7:eq-full-sequence-1}
\bigl|S_j(R\cdot\vec f)-S_j(\vec f)\bigr|\ \ge\ \eta\qquad\text{for every } j\in J'.
\end{equation}

The sequence indexed by $J'$ is a subsequence of the given one. By the simultaneous convergence
lemma (Lemma~\ref{10.7:lem-simultaneous-convergence}) it therefore has a further subsequence,
indexed by an infinite set $J''\subset J'$, along which $S^T_{n';K_j,m_j}$ converges for every
$n'\ge 2$ on every admissible $C^1_c$ tuple; in particular $S_j$ converges on $\vec f$ and on
$R\cdot\vec f$. We write $S''$ for the limit family so obtained, one member for each arity
$n'\ge 2$, and also $S''$ for its member of arity $n$. By the same lemma, $S''$ obeys the
equicontinuity estimate \eqref{10.7:eq-equicontinuity-a} in its limit form and is $W_4$-invariant.

We now check that the three inputs of the argument leading to the conclusion step
(Theorem~\ref{10.7:thm-conclusion-step}) that refer to the approximating sequence remain
available along $J''$. Fix a generator $\omega$ of a one-parameter rotation group
$R^{\omega}_t=e^{t\omega}$. For each cut-off $\chi$ write $\mathcal{N}_j$ and
$\mathcal{B}^{\chi}_j$ for the members on the lattice $(K_j,m_j)$ of the insertions
$\mathcal{N}^{\omega,\chi}$ and $\mathcal{B}^{\omega,\chi}$; the dependence on $\omega$, and in
$\mathcal{N}_j$ also that on $\chi$, is suppressed in the notation.

First, since the rotational Ward identity hypothesis holds at every $(K_j,m_j)$ beyond the floor
indices, the orbit identity on which the proof of Theorem~\ref{10.7:thm-conclusion-step} rests,
which relates the derivative
$D^{\omega}_{K_j,m_j}$ along the orbit to the insertions $\mathcal{N}_j$ and
$\mathcal{B}^{\chi}_j$, holds at each $j\in J''$ beyond those indices.

Second, for every cut-off $\chi$, $\mathcal{N}_j\to 0$ along $J''$.

Third, a $\limsup$ along a subsequence does not exceed the $\limsup$ along the full sequence.
Hence, for every cut-off $\chi$,
\begin{equation*}
\limsup_{j\in J''}\bigl|\mathcal{B}^{\chi}_j\bigr|\ \le\ \limsup_{j\to\infty}\bigl|\mathcal{B}^{\chi}_j\bigr|.
\end{equation*}
By the far-sphere flux hypothesis $(\mathrm{IR}\text{-}\mathrm{fl})_{\omega_0}$ in its pointwise
form along the full sequence, as displayed in \eqref{10.7:eq-hypotheses-used-form-b}, the
infimum over $\chi$ of the right-hand side is zero.
Taking the infimum over $\chi$ on both sides gives
\begin{equation}\label{10.7:eq-full-sequence-2}
0\ \le\ \inf_{\chi}\,\limsup_{j\in J''}\bigl|\mathcal{B}^{\chi}_j\bigr|
\ \le\ \inf_{\chi}\,\limsup_{j\to\infty}\bigl|\mathcal{B}^{\chi}_j\bigr|\ =\ 0 .
\end{equation}
So the infimum over $\chi$ is still zero when the $\limsup$ is taken along $J''$.

The argument leading to the conclusion step needs one further input. It used the fact that the
limit family is a subsequential limit along the joint diagonal only through three statements on
the limit family, each established in the course of the argument leading to
Theorem~\ref{10.7:thm-conclusion-step}; of the first of these, in the order in which that argument
establishes them, only part (a) is used. Each of these three
is valid for $S''$, which is a limit along the subsequence $J''$ obeying the equicontinuity
estimate \eqref{10.7:eq-equicontinuity-a} in its limit form and $W_4$-invariance. Together with the
identity at each $j\in J''$, the vanishing of $\mathcal{N}_j$ along $J''$ and
\eqref{10.7:eq-full-sequence-2}, that argument therefore applies to $S''$ unchanged. It yields
\begin{equation*}
S''(R\cdot\vec f)=S''(\vec f).
\end{equation*}

Finally, along $J''$ we have $S_j(R\cdot\vec f)\to S''(R\cdot\vec f)$ and
$S_j(\vec f)\to S''(\vec f)$. Consequently
\begin{equation*}
\lim_{j\in J''}\bigl|S_j(R\cdot\vec f)-S_j(\vec f)\bigr|
=\bigl|S''(R\cdot\vec f)-S''(\vec f)\bigr|=0 .
\end{equation*}
Since $J''\subset J'$, this contradicts \eqref{10.7:eq-full-sequence-1}. Hence
$S_j(R\cdot\vec f)-S_j(\vec f)\to 0$ along the full sequence.
\end{proof}

\begin{remark}[Only vanishing of the ultraviolet term is used, not its rate]
\label{10.7:rem-rate-not-used}
Let $R_t=R^{\omega}_t=e^{t\omega}$ act on a tuple $\vec f=(f_1,\dots,f_n)$ componentwise, and let
$S_j$ denote $S^T_{n;K_j,m_j}$ and $D_j$ denote $D^{\omega}_{K_j,m_j}$, the derivative along the
orbit, along the diagonal. Of the ultraviolet insertion $\mathcal{N}^{\chi}_j$ (the insertion
$\mathcal{N}^{\omega,\chi}$ at the $j$-th member of the diagonal), the argument uses
only the statement
\begin{equation*}
\mathcal{N}^{\chi}_j(\vec h)\longrightarrow 0 \qquad\text{at each fixed tuple } \vec h .
\end{equation*}
The conclusion is therefore unchanged if the dimension-six coefficient bound is available only in its
weaker form, without a rate; what is lost is the displayed rate.

Suppose instead that both the ultraviolet and the infrared term come with rates: the ultraviolet
contribution along the orbit is bounded by $C\eta(\varepsilon_j)\prod_i\|f_i\|_{\sharp}$, where
$\eta(\varepsilon_j)$ is the rate supplied at the $j$-th member of the diagonal by the
dimension-six coefficient bound in its form with a rate, and $\beta(j;R_s\cdot\vec f)$ denotes the
infrared bound at the point $R_s\cdot\vec f$ of the orbit. Then Proposition~\ref{10.7:prop-orbit-constancy-integrated}
gives, at finite cutoff,
\begin{equation}\label{10.7:eq-rate-not-used-1}
\bigl|S_j(R_t\cdot\vec f)-S_j(\vec f)\bigr|
\le |t|\Bigl[C\eta(\varepsilon_j)\prod_{i=1}^{n}\|f_i\|_{\sharp}
+\sup_{s}\beta(j;R_s\cdot\vec f)\Bigr].
\end{equation}
\end{remark}

\begin{proof}
As in the proof of Proposition~\ref{10.7:prop-orbit-constancy-integrated}, the function
$\varphi_j(t)=S_j(R_t\cdot\vec f)$ is continuously differentiable, with
$\varphi_j'(s)=D_j(R_s\cdot\vec f)$. In that proof the ultraviolet insertion enters only through the
convergence $D_j(R_s\cdot\vec f)\to c(R_s\cdot\vec f,\omega)=0$ at each fixed $s$, and, as far as
the ultraviolet insertion is concerned, this convergence uses only that
$\mathcal{N}^{\chi}_j(\vec h)\to 0$ at each fixed tuple $\vec h$; the infrared term is treated as in
that proof and is unaffected. This proves the first assertion.

For the rate, the fundamental theorem of calculus gives
\begin{equation*}
S_j(R_t\cdot\vec f)-S_j(\vec f)=\varphi_j(t)-\varphi_j(0)=\int_0^t D_j(R_s\cdot\vec f)\,ds,
\end{equation*}
and hence
\begin{equation*}
\bigl|S_j(R_t\cdot\vec f)-S_j(\vec f)\bigr|\le |t|\,\sup_s\bigl|D_j(R_s\cdot\vec f)\bigr| ,
\end{equation*}
the supremum being over $s$ between $0$ and $t$.
In the proof of Proposition~\ref{10.7:prop-orbit-constancy-integrated}, $D_j(R_s\cdot\vec f)$ is
controlled through its ultraviolet and its infrared insertion, the anomaly
$c(R_s\cdot\vec f,\omega)$ being zero. With both rates available, the ultraviolet part of
$D_j(R_s\cdot\vec f)$ is bounded by $C\eta(\varepsilon_j)\prod_i\|f_i\|_{\sharp}$ and the infrared
part by $\beta(j;R_s\cdot\vec f)$.
Inserting these two bounds into the last display, and enlarging the supremum of
$\beta(j;R_s\cdot\vec f)$ over $s$ between $0$ and $t$ to its supremum over all $s$, gives
\eqref{10.7:eq-rate-not-used-1}.
\end{proof}

\begin{lemma}[Robustness under a renormalised rotational Ward identity]
\label{10.7:lem-renormalised-identity}
Let $\vec h\in\mathcal{T}^{(r)}_n(\mathcal{R},\mathsf{d})$ and let $\chi$ range over the admissible
cut-offs. Suppose that the rotational Ward identity, the first hypothesis of
\eqref{10.7:eq-hypotheses-used-form-b}, holds, in place of its stated form, in the renormalised form
\begin{equation}\label{10.7:eq-renormalised-identity-1}
  (1-\eta_{K,m})\,D^{\omega}_{K,m}(\vec h)
  =\mathcal{N}^{\omega,\chi}_{K,m}(\vec h)+\mathcal{B}^{\omega,\chi}_{K,m}(\vec h),
\end{equation}
where the factor $\eta_{K,m}$, the \emph{renormalisation factor}, depends on neither $\vec h$
nor $\chi$. Along the diagonal write $\eta_j=\eta_{K_j,m_j}$, $D_j$, $\mathcal{N}^{\chi}_j$ and
$\mathcal{B}^{\chi}_j$ for the quantities in \eqref{10.7:eq-renormalised-identity-1} at
$(K,m)=(K_j,m_j)$, and $c=c(\vec h,\omega)$. Assume $r\ge 2$ and the second hypothesis of
\eqref{10.7:eq-hypotheses-used-form-b}.
\begin{enumerate}
\item[(a)] Assume
\begin{equation}\label{10.7:eq-renormalised-identity-a}
  \delta:=\liminf_{j\to\infty}\,\lvert 1-\eta_j\rvert>0 .
\end{equation}
Then the second hypothesis of \eqref{10.7:eq-hypotheses-used-form-b} and the far-sphere flux
hypothesis $(\mathrm{IR}\text{-}\mathrm{fl})_{\omega_0}$ (the flux of the lattice rotation current
through far spheres vanishes along one sequence of radii), in
the pointwise form of the same display, give $c\equiv 0$ on the class, and
hence the conclusion of Theorem~\ref{10.7:thm-conclusion-step}.
\item[(b)] The boundary term remains independent of the cut-off in the limit:
$\mathcal{B}^{\chi}_j-\mathcal{B}^{\chi'}_j\to 0$ for all admissible $\chi,\chi'$.
\item[(c)] If $\eta_j\to 1$, then \eqref{10.7:eq-renormalised-identity-1} carries no information
about $c$.
\item[(d)] If the factor depends on the slot, $\eta^{(i)}_j$ for $i=1,\dots,n$, not all equal, then
the left-hand side is the derivative along a path that is not a one-parameter orbit, and the
argument yields nothing: the individual slot terms do not vanish even for $O(4)$-invariant
Schwinger functions; only their sum does.
\end{enumerate}
\end{lemma}

\begin{proof}
We use two limits. By the convergence of the derivative along the orbit to the rotational anomaly
on the class, which holds for $r\ge 2$, $D_j\to c$; by the second hypothesis of
\eqref{10.7:eq-hypotheses-used-form-b}, $\mathcal{N}^{\chi}_j\to 0$ for every admissible $\chi$. By
\eqref{10.7:eq-renormalised-identity-1},
\begin{equation}\label{10.7:eq-renormalised-identity-2}
  \mathcal{B}^{\chi}_j=(1-\eta_j)\,c+(1-\eta_j)(D_j-c)-\mathcal{N}^{\chi}_j .
\end{equation}

(a) Condition \eqref{10.7:eq-renormalised-identity-a} is stable under passage to subsequences,
since the lower limit along a subsequence is at least the lower limit along the whole sequence.
Suppose first that $\lvert 1-\eta_j\rvert$ is bounded. Then the last two terms of
\eqref{10.7:eq-renormalised-identity-2} tend to zero, and for every admissible $\chi$
\begin{equation*}
  \limsup_{j}\lvert\mathcal{B}^{\chi}_j\rvert
  =\lvert c\rvert\,\limsup_{j}\lvert 1-\eta_j\rvert
  \ge\lvert c\rvert\,\liminf_{j}\lvert 1-\eta_j\rvert=\delta\lvert c\rvert .
\end{equation*}
Suppose next that $\lvert 1-\eta_j\rvert$ is unbounded, so that $\lvert 1-\eta_j\rvert\to\infty$
along a subsequence, and that $c\ne 0$. Since $D_j\to c$,
\begin{equation*}
  \lvert\mathcal{B}^{\chi}_j\rvert
  \ge\lvert 1-\eta_j\rvert\,\lvert D_j\rvert-\lvert\mathcal{N}^{\chi}_j\rvert\to\infty
\end{equation*}
along that subsequence, so $\limsup_j\lvert\mathcal{B}^{\chi}_j\rvert=\infty$ for every admissible
$\chi$. By the far-sphere flux hypothesis $(\mathrm{IR}\text{-}\mathrm{fl})_{\omega_0}$, as in
Proposition~\ref{10.7:prop-anomaly-vanishes}, the
infimum over admissible $\chi$ of $\limsup_j\lvert\mathcal{B}^{\chi}_j\rvert$ is zero. In the
bounded case this gives $\delta\lvert c\rvert\le 0$, hence $c=0$ because $\delta>0$; in the
unbounded case it excludes $c\ne 0$. Either way $c=0$. As $\vec h$ was arbitrary in the class,
$c\equiv 0$, from which the conclusion of Theorem~\ref{10.7:thm-conclusion-step} follows as there.

(b) Since $\eta_{K,m}$ does not depend on $\chi$, subtracting
\eqref{10.7:eq-renormalised-identity-1} for $\chi'$ from the same identity for $\chi$ gives
$\mathcal{B}^{\chi}_j-\mathcal{B}^{\chi'}_j=\mathcal{N}^{\chi'}_j-\mathcal{N}^{\chi}_j$, which
tends to zero.

(c) If $\eta_j\to 1$, then $(1-\eta_j)D_j\to 0$ since $D_j$ converges, so
\eqref{10.7:eq-renormalised-identity-1} yields $\mathcal{B}^{\chi}_j\to 0$ whatever the value of
$c$; it places no constraint on $c$.

(d) Write $D^{\omega}_{K,m}(\vec h)=\sum_{i=1}^n D^{\omega,(i)}_{K,m}(\vec h)$, where the slot
term $D^{\omega,(i)}_{K,m}(\vec h)$ is the contribution in which the rotation acts on $h_i$ alone.
With slot-dependent factors the left-hand side becomes
$\sum_i(1-\eta^{(i)}_j)D^{\omega,(i)}_{K_j,m_j}(\vec h)$, the derivative at $t=0$ along the path
\begin{equation*}
  t\longmapsto\bigl(R^{\omega}_{(1-\eta^{(1)}_j)t}h_1,\dots,R^{\omega}_{(1-\eta^{(n)}_j)t}h_n\bigr),
\end{equation*}
which, the factors being unequal, is not a one-parameter orbit of the simultaneous rotation of all
slots. For $O(4)$-invariant Schwinger functions only the sum of the slot terms, the derivative
along the orbit, vanishes; the individual slot terms do not. The identity with slot-dependent
factors therefore does not determine $c$, and the argument yields nothing.
\end{proof}

\begin{lemma}[The transfer lemma for small scalar annulus terms]\label{10.7:lem-annulus-transfer}
Let $\vec h = (h_1,\dots,h_n)$ be the tuple of real Lipschitz test functions of the rotational
Ward identity, supported in $B(0,\mathcal{R}/2)$ with supports pairwise at distance at least
$\mathsf{d}$. Suppose that a term of that identity is exactly $\mathcal{O}(w_K^{(m)})$, where
$w_K$ is a real Lipschitz function with
\begin{equation*}
  \operatorname{supp} w_K \subset \bar B(0,\rho_w), \qquad
  \operatorname{dist}\bigl(\operatorname{supp} w_K, B(0,\mathcal{R}/2)\bigr) \ge \mathsf{d}_w > 0,
  \qquad \|w_K\|_{\sharp} \le C_w \varepsilon ,
\end{equation*}
for every cut-off level $K$, with $\rho_w$, $\mathsf{d}_w$ and $C_w$ independent of $K$; here
$\varepsilon = \varepsilon(K)$ is the parameter with which the rotational Ward identity is
written at cut-off level $K$, which tends to $0$ as $K \to \infty$, and
$\varepsilon_j := \varepsilon(K_j)$.
Write $w$ for $w_{K_j}$. Then for $j \ge j_1(n+1, \max\{\rho_w,\mathcal{R}/2\})$
\begin{equation}\label{10.7:eq-annulus-transfer-1}
\begin{split}
  \bigl|\kappa\bigl(\mathcal{O}(w^{(m_j)});\mathcal{O}(\vec h)\bigr)\bigr|
  &= \bigl|S^T_{n+1;K_j,m_j}(w,h_1,\dots,h_n)\bigr| \\
  &\le C^\infty_{n+1}\bigl(\min\{\mathsf{d},\mathsf{d}_w\};\max\{\rho_w,\mathcal{R}/2\}\bigr)
     \, C_w \varepsilon_j \prod_{i=1}^n \|h_i\|_{\sharp},
\end{split}
\end{equation}
and the right-hand side tends to $0$ as $j \to \infty$. Such terms satisfy the conclusion of
the pointwise form of the dimension-six coefficient bound, taken as a hypothesis, by the per-ball
volume-free bound alone, and may be counted in the inserted lattice function $N$ of the identity.
Orientation-resolved annulus pieces are not of the form $\mathcal{O}(w)$ and remain in the
inserted lattice function $A$; they are treated in the regrouping of the annulus terms in this
chapter.
\end{lemma}

\begin{proof}
The equality in \eqref{10.7:eq-annulus-transfer-1} is the definition of the connected
Schwinger function at $(K_j,m_j)$ as the joint cumulant of the observables
$\mathcal{O}(w^{(m_j)}), \mathcal{O}(h_1^{(m_j)}),\dots,\mathcal{O}(h_n^{(m_j)})$.

For the inequality we apply the first estimate in \eqref{10.7:eq-equicontinuity-a} of the
equicontinuity estimate (Lemma~\ref{10.7:lem-equicontinuity}(b)) with $n+1$
Lipschitz sources $w, h_1,\dots,h_n$, radius $\rho = \max\{\rho_w,\mathcal{R}/2\}$, the sets
$\operatorname{supp} w, \operatorname{supp} h_1,\dots,\operatorname{supp} h_n$ as the supporting
sets, and separation $\mathsf{d}' = \min\{\mathsf{d},\mathsf{d}_w\}$. Its hypotheses hold: all
the sets lie in the ball of radius $\max\{\rho_w,\mathcal{R}/2\}$, since
$\operatorname{supp} w \subset \bar B(0,\rho_w)$ and $\operatorname{supp} h_i \subset
B(0,\mathcal{R}/2)$; the sets $\operatorname{supp} h_i$ are pairwise at distance at least
$\mathsf{d}$; each $\operatorname{supp} h_i$ lies in $B(0,\mathcal{R}/2)$, whose distance from
$\operatorname{supp} w$ is at least $\mathsf{d}_w$; and
$j \ge j_1(n+1,\max\{\rho_w,\mathcal{R}/2\})$. That estimate then bounds
the modulus of the connected Schwinger function by
$C^\infty_{n+1}(\min\{\mathsf{d},\mathsf{d}_w\};\max\{\rho_w,\mathcal{R}/2\})$ times
$\|w\|_{\sharp}\prod_{i}\|h_i\|_{\sharp}$. Inserting $\|w\|_{\sharp} \le C_w\varepsilon_j$ gives
\eqref{10.7:eq-annulus-transfer-1}. The constant does not depend on $j$, since $\rho_w$,
$\mathsf{d}_w$ and $C_w$ do not depend on $K$. By clause~(0) of the
main theorem, $K_j \to \infty$ along the diagonal sequence of clause~(ii-b), so, since
$\varepsilon(K) \to 0$ as $K \to \infty$, we have $\varepsilon_j = \varepsilon(K_j) \to 0$;
hence the right-hand side tends to zero.

The bound \eqref{10.7:eq-annulus-transfer-1} gives for such a term the conclusion of the
pointwise form of the dimension-six coefficient bound, taken as a hypothesis, and it uses only
the per-ball volume-free bound; this is why the term may be counted in $N$. An orientation-resolved annulus
piece is not a single smeared observable $\mathcal{O}(w)$, so the argument above does not apply
to it, and it remains in $A$.
\end{proof}

\begin{definition}[Coordinate-plane generators and signed-permutation conjugation]
\label{10.7:not-plane-generators}
Let $(e_\mu)$ be the standard basis of $\mathbb{R}^4$. For coordinate indices $\rho\neq\sigma$ write
$E_{\rho\sigma}$ for the matrix unit, $E_{\rho\sigma}e_\tau=\delta_{\sigma\tau}e_\rho$, and put
\begin{equation*}
A_{\rho\sigma}:=E_{\sigma\rho}-E_{\rho\sigma},
\qquad\text{so that}\qquad
A_{\rho\sigma}e_\rho=e_\sigma,\quad A_{\rho\sigma}e_\sigma=-e_\rho .
\end{equation*}
Set $G_{\rho\sigma}(\theta):=\exp(\theta A_{\rho\sigma})$ for $\theta\in\mathbb{R}$; then
\begin{equation}\label{10.7:eq-plane-generators-1}
G_{\rho\sigma}(\theta)
=\mathbb{1}-(1-\cos\theta)\,(E_{\rho\rho}+E_{\sigma\sigma})+\sin\theta\,A_{\rho\sigma}.
\end{equation}
Write $SO(2)_{\rho\sigma}:=\{G_{\rho\sigma}(\theta):\theta\in\mathbb{R}\}$. In the group $W_4$ of signed
permutation matrices, $D\subset W_4$ denotes the diagonal sign matrices, $s_\tau$ the sign change of the
coordinate $\tau$, and $P_\pi$ the permutation matrices. The space $\Lambda^2\mathbb{R}^4$ carries the basis
$e_\rho\wedge e_\sigma$, $\rho<\sigma$. Let $\operatorname{Sym}_{off}$ be the space of real symmetric
$4\times4$ matrices with zero diagonal, with basis $S_{\rho\sigma}:=E_{\rho\sigma}+E_{\sigma\rho}$,
$\rho<\sigma$, on which $W_4$ acts by conjugation. For a permutation $\pi$ of the coordinate indices and
signs $\varepsilon=(\varepsilon_\mu)$, $\varepsilon_\mu\in\{\pm1\}$, let $w=w(\pi,\varepsilon)\in W_4$ be
the matrix $e_\mu\mapsto\varepsilon_\mu e_{\pi(\mu)}$. Then for all indices $a,b$ and all $\rho\neq\sigma$,
\begin{equation}\label{10.7:eq-plane-generators-a}
\begin{gathered}
wE_{ab}w^{-1}=\varepsilon_a\varepsilon_b\,E_{\pi(a)\pi(b)},\\
wA_{\rho\sigma}w^{-1}=\varepsilon_\rho\varepsilon_\sigma\,A_{\pi(\rho)\pi(\sigma)},
\qquad
wS_{\rho\sigma}w^{-1}=\varepsilon_\rho\varepsilon_\sigma\,S_{\pi(\rho)\pi(\sigma)} .
\end{gathered}
\end{equation}
\end{definition}

\begin{proof}
The two values of $A_{\rho\sigma}$ on $e_\rho,e_\sigma$ follow from $E_{\sigma\rho}e_\rho=e_\sigma$,
$E_{\rho\sigma}e_\rho=0$, $E_{\sigma\rho}e_\sigma=0$ and $E_{\rho\sigma}e_\sigma=e_\rho$.

For \eqref{10.7:eq-plane-generators-1}, use $E_{ab}E_{cd}=\delta_{bc}E_{ad}$ and $\rho\neq\sigma$:
\begin{equation*}
A_{\rho\sigma}^2
=E_{\sigma\rho}E_{\sigma\rho}-E_{\sigma\rho}E_{\rho\sigma}-E_{\rho\sigma}E_{\sigma\rho}
+E_{\rho\sigma}E_{\rho\sigma}
=-(E_{\rho\rho}+E_{\sigma\sigma}).
\end{equation*}
Multiplying by $A_{\rho\sigma}$ on the right and expanding the same way,
$(E_{\rho\rho}+E_{\sigma\sigma})A_{\rho\sigma}=-E_{\rho\sigma}+E_{\sigma\rho}=A_{\rho\sigma}$, hence
$A_{\rho\sigma}^3=-A_{\rho\sigma}$. By induction, for $j\ge1$,
$A_{\rho\sigma}^{2j}=(-1)^j(E_{\rho\rho}+E_{\sigma\sigma})$ and
$A_{\rho\sigma}^{2j+1}=(-1)^jA_{\rho\sigma}$. Inserting this into the exponential series,
\begin{equation*}
\exp(\theta A_{\rho\sigma})
=\mathbb{1}+\Bigl(\sum_{j\ge1}\frac{(-1)^j\theta^{2j}}{(2j)!}\Bigr)(E_{\rho\rho}+E_{\sigma\sigma})
+\Bigl(\sum_{j\ge0}\frac{(-1)^j\theta^{2j+1}}{(2j+1)!}\Bigr)A_{\rho\sigma},
\end{equation*}
and the two series are $\cos\theta-1$ and $\sin\theta$, which is \eqref{10.7:eq-plane-generators-1}.

For \eqref{10.7:eq-plane-generators-a}: since $\varepsilon_c^2=1$, $w^{-1}e_{\pi(c)}=\varepsilon_c e_c$ for
every index $c$. Hence
\begin{equation*}
wE_{ab}w^{-1}e_{\pi(c)}=\varepsilon_c\,wE_{ab}e_c=\varepsilon_c\delta_{bc}\,we_a
=\delta_{bc}\,\varepsilon_a\varepsilon_b\,e_{\pi(a)},
\end{equation*}
and as $c$ runs over the indices, $e_{\pi(c)}$ runs over the basis; the right-hand side is
$\varepsilon_a\varepsilon_bE_{\pi(a)\pi(b)}e_{\pi(c)}$, giving the first line. The second line follows
by linearity of conjugation from $A_{\rho\sigma}=E_{\sigma\rho}-E_{\rho\sigma}$ and
$S_{\rho\sigma}=E_{\rho\sigma}+E_{\sigma\rho}$, since the factor $\varepsilon_\rho\varepsilon_\sigma$ is
symmetric in $\rho,\sigma$.
\end{proof}

\begin{lemma}[Irreducibility of the antisymmetric square under the hyperoctahedral group]
\label{10.7:lem-antisymmetric-irreducible}
Let $e_1,\dots,e_4$ be the standard basis of $\mathbb{R}^4$, let $\mathfrak{so}(4)$ be the space of
real antisymmetric $4\times 4$ matrices, and for $g\in O(4)$ let $\Lambda^2(g)$ be the linear map
of $\Lambda^2\mathbb{R}^4$ with $\Lambda^2(g)(u\wedge v)=gu\wedge gv$. Let $D$, $s_\tau$, $P_\pi$,
$S_{\rho\sigma}$ and $\operatorname{Sym}_{off}$ be as in Definition~\ref{10.7:not-plane-generators}, and let
$W_4$ act on $\operatorname{Sym}_{off}$ by $\eta\mapsto g\eta g^{-1}$.
\begin{enumerate}
\item[(a)] The map $\Phi\colon\mathfrak{so}(4)\to\Lambda^2\mathbb{R}^4$,
\begin{equation}\label{10.7:eq-antisymmetric-irreducible-1}
\Phi(\omega)=\sum_{\rho<\sigma}\omega_{\rho\sigma}\,e_\rho\wedge e_\sigma ,
\end{equation}
is a linear isomorphism, and $\Phi(g\omega g^{-1})=\Lambda^2(g)\Phi(\omega)$ for every
$g\in O(4)$.
\item[(b)] $\Lambda^2\mathbb{R}^4$ is irreducible under $W_4$.
\item[(c)] $\operatorname{Sym}_{off}$ is irreducible under $W_4$ and is not equivalent to
$\Lambda^2\mathbb{R}^4$ as a representation of $W_4$.
\end{enumerate}
\end{lemma}

\begin{proof}
(a) An antisymmetric $\omega$ is determined by its entries $\omega_{\rho\sigma}$, $\rho<\sigma$,
and these may be prescribed freely; since $\{e_\rho\wedge e_\sigma\}_{\rho<\sigma}$ is a basis of
$\Lambda^2\mathbb{R}^4$, $\Phi$ is a linear bijection. Because $\omega_{\sigma\rho}=-\omega_{\rho\sigma}$,
$e_\sigma\wedge e_\rho=-e_\rho\wedge e_\sigma$ and $e_\rho\wedge e_\rho=0$,
\begin{equation}\label{10.7:eq-antisymmetric-irreducible-2}
\Phi(\omega)=\tfrac12\sum_{\rho,\sigma}\omega_{\rho\sigma}\,e_\rho\wedge e_\sigma
\qquad(\omega\in\mathfrak{so}(4)).
\end{equation}
For $g\in O(4)$ we have $ge_\rho=\sum_\alpha g_{\alpha\rho}e_\alpha$, hence
\begin{align*}
\Lambda^2(g)\Phi(\omega)
&=\tfrac12\sum_{\rho,\sigma}\omega_{\rho\sigma}\sum_{\alpha,\beta}g_{\alpha\rho}g_{\beta\sigma}\,
e_\alpha\wedge e_\beta\\
&=\tfrac12\sum_{\alpha,\beta}\Bigl(\sum_{\rho,\sigma}g_{\alpha\rho}\omega_{\rho\sigma}g_{\beta\sigma}\Bigr)
e_\alpha\wedge e_\beta\\
&=\tfrac12\sum_{\alpha,\beta}(g\omega g^{\mathsf T})_{\alpha\beta}\,e_\alpha\wedge e_\beta .
\end{align*}
Since $g^{\mathsf T}=g^{-1}$, the matrix $g\omega g^{\mathsf T}=g\omega g^{-1}$ is antisymmetric, and
\eqref{10.7:eq-antisymmetric-irreducible-2} applied to it identifies the last expression with
$\Phi(g\omega g^{-1})$.

(b) For $S\subset\{1,2,3,4\}$ and $\varepsilon=\operatorname{diag}(\varepsilon_1,\dots,\varepsilon_4)\in D$
put $\chi_S(\varepsilon)=\prod_{\tau\in S}\varepsilon_\tau$, and write
$\chi_{\rho\sigma}=\chi_{\{\rho,\sigma\}}$. Since $\varepsilon e_\rho=\varepsilon_\rho e_\rho$,
\begin{equation*}
\Lambda^2(\varepsilon)(e_\rho\wedge e_\sigma)=\chi_{\rho\sigma}(\varepsilon)\,e_\rho\wedge e_\sigma ,
\end{equation*}
so $D$ acts on each of the six lines $\mathbb{R}e_\rho\wedge e_\sigma$ by the character
$\chi_{\rho\sigma}$. These six characters are pairwise distinct: $\chi_{\rho\sigma}(s_\tau)=-1$ if
and only if $\tau\in\{\rho,\sigma\}$, so $\chi_{\rho\sigma}$ determines $\{\rho,\sigma\}$. For
$S\neq\emptyset$, summing first over $\varepsilon_\tau=\pm1$ for one $\tau\in S$ gives
$\sum_{\varepsilon\in D}\chi_S(\varepsilon)=0$. Define
\begin{equation*}
\Pi_{\rho\sigma}=\frac1{16}\sum_{\varepsilon\in D}\chi_{\rho\sigma}(\varepsilon)\,\Lambda^2(\varepsilon).
\end{equation*}
Then $\Pi_{\rho\sigma}(e_\alpha\wedge e_\beta)$ equals $e_\alpha\wedge e_\beta$ times
$16^{-1}\sum_\varepsilon\chi_S(\varepsilon)$, with $S$ the symmetric difference of $\{\rho,\sigma\}$
and $\{\alpha,\beta\}$, because $\chi_{\rho\sigma}\chi_{\alpha\beta}=\chi_S$. This is
$e_\alpha\wedge e_\beta$ if $\{\alpha,\beta\}=\{\rho,\sigma\}$ and $0$ otherwise. Thus
$\Pi_{\rho\sigma}$ is the projection onto $\mathbb{R}e_\rho\wedge e_\sigma$ along the other five
lines, and $\sum_{\rho<\sigma}\Pi_{\rho\sigma}$ is the identity. Let $V$ be a $D$-invariant subspace.
Each $\Pi_{\rho\sigma}$ is a linear combination of the maps $\Lambda^2(\varepsilon)$, so
$\Pi_{\rho\sigma}V\subset V$. Hence every $v\in V$ is the sum
$v=\sum_{\rho<\sigma}\Pi_{\rho\sigma}v$ of vectors of $V\cap\mathbb{R}e_\rho\wedge e_\sigma$, and $V$
is spanned by the lines $\mathbb{R}e_\rho\wedge e_\sigma$ that it contains. Now let $V\neq0$ be
$W_4$-invariant. It contains some line $\mathbb{R}e_\rho\wedge e_\sigma$. The permutation matrices
$P_\pi$, $\pi$ in the symmetric group $S_4$, lie in $W_4$; each $P_\pi$ maps every basis vector
$e_\rho$ to a basis vector $e_{\rho'}$, and the induced permutations $\rho\mapsto\rho'$ exhaust
$S_4$. Hence
\begin{equation*}
\Lambda^2(P_\pi)(e_\rho\wedge e_\sigma)=e_{\rho'}\wedge e_{\sigma'},
\end{equation*}
which is $\pm$ the basis vector of the image pair $\{\rho',\sigma'\}$. Since $S_4$ acts transitively
on two-element subsets of $\{1,2,3,4\}$, $V$ contains all six lines, and $V=\Lambda^2\mathbb{R}^4$.

(c) The argument of (b) applies to $\operatorname{Sym}_{off}$, with the conjugation formulas
\eqref{10.7:eq-plane-generators-a} in place of the action of $\Lambda^2$. Reading entrywise,
$(\varepsilon \eta\varepsilon^{-1})_{\alpha\beta}=\varepsilon_\alpha \eta_{\alpha\beta}\varepsilon_\beta$
for $\eta\in\operatorname{Sym}_{off}$, so
$\varepsilon S_{\rho\sigma}\varepsilon^{-1}=\chi_{\rho\sigma}(\varepsilon)S_{\rho\sigma}$. Moreover
$P_\pi S_{\rho\sigma}P_\pi^{-1}=S_{\rho'\sigma'}$, with $\rho',\sigma'$ as in (b) and
$S_{\sigma\rho}=S_{\rho\sigma}$. The
six lines $\mathbb{R}S_{\rho\sigma}$ therefore carry the six distinct characters $\chi_{\rho\sigma}$
of $D$. The projections formed with these characters show that a $D$-invariant subspace is spanned
by the lines $\mathbb{R}S_{\rho\sigma}$ it contains, and transitivity of $S_4$ gives irreducibility.

For inequivalence, compare traces at the transposition $t=P_{(12)}\in W_4$. On $\Lambda^2\mathbb{R}^4$:
$e_1\wedge e_2\mapsto e_2\wedge e_1=-e_1\wedge e_2$ contributes $-1$; $e_3\wedge e_4$ is fixed and
contributes $+1$; the pairs $e_1\wedge e_3\leftrightarrow e_2\wedge e_3$ and
$e_1\wedge e_4\leftrightarrow e_2\wedge e_4$ are swapped and contribute $0$. So the trace of
$\Lambda^2(t)$ on $\Lambda^2\mathbb{R}^4$ is $-1+1+0=0$. On $\operatorname{Sym}_{off}$, $S_{12}$ and
$S_{34}$ are
fixed, contributing $+1$ each, and $S_{13}\leftrightarrow S_{23}$, $S_{14}\leftrightarrow S_{24}$ are
swapped, contributing $0$. The trace of conjugation by $t$ on $\operatorname{Sym}_{off}$ is therefore
$+1+1+0=2$. Equivalent representations have equal
traces at every group element, so the two representations are not equivalent.
\end{proof}

\begin{lemma}[Coordinate-plane rotations generate the rotation group]
\label{10.7:lem-plane-rotations-generate}
Write $G_{\rho\sigma}(\theta)$, $1\le\rho<\sigma\le4$, for the coordinate-plane rotations and
$SO(2)_{\rho\sigma}=\{G_{\rho\sigma}(\theta):\theta\in\mathbb{R}\}$, and let $\|\cdot\|$ be the
operator norm of real $4\times4$ matrices.
\begin{enumerate}
\item[(a)] Every $R\in SO(4)$ is a product of at most six coordinate-plane rotations.
\item[(b)] For every $w\in W_4$, with underlying permutation $\pi$, and all $\rho<\sigma$,
\begin{equation*}
w\,SO(2)_{\rho\sigma}\,w^{-1}=SO(2)_{\pi(\rho)\pi(\sigma)}.
\end{equation*}
\item[(c)] Every infinite subgroup $\Gamma\subset SO(2)_{\rho\sigma}$ is dense in $SO(2)_{\rho\sigma}$.
\end{enumerate}
\end{lemma}

\begin{proof}
(a) For a real $4\times4$ matrix $B$, a column index $c$ and a pair $\rho<\sigma$, put
$r=(B_{\rho c}^2+B_{\sigma c}^2)^{1/2}$. If $r>0$, choose $\varphi$ with
\begin{equation*}
\cos\varphi=B_{\rho c}/r,\qquad \sin\varphi=-B_{\sigma c}/r ;
\end{equation*}
if $r=0$, put $\varphi=0$. Left multiplication by $G_{\rho\sigma}(\varphi)$ changes only rows $\rho$
and $\sigma$. Its block on the coordinates $\rho,\sigma$ has rows $(\cos\varphi,-\sin\varphi)$ and
$(\sin\varphi,\cos\varphi)$, so the new entries in column $c$ are
\begin{align*}
(G_{\rho\sigma}(\varphi)B)_{\rho c}&=\cos\varphi\,B_{\rho c}-\sin\varphi\,B_{\sigma c}
=(B_{\rho c}^2+B_{\sigma c}^2)/r=r\ge0,\\
(G_{\rho\sigma}(\varphi)B)_{\sigma c}&=\sin\varphi\,B_{\rho c}+\cos\varphi\,B_{\sigma c}
=(-B_{\sigma c}B_{\rho c}+B_{\rho c}B_{\sigma c})/r=0 .
\end{align*}
When $r=0$ both entries are already zero and $G_{\rho\sigma}(0)=\mathbb{1}$. Thus the step zeroes the
$(\sigma,c)$ entry and makes the $(\rho,c)$ entry equal to $r\ge0$.

Starting from $B=R$, apply this step to column $1$ with the planes $(1,2)$, $(1,3)$, $(1,4)$ in turn.
Each step alters only row $1$ and one other row, so a zero created in an earlier step is not
disturbed. Afterwards the entries $(2,1),(3,1),(4,1)$ vanish and $(1,1)$ is non-negative.

Next apply the step to column $2$ with the planes $(2,3)$, $(2,4)$. These steps mix only rows
$2,3,4$, whose column-$1$ entries are zero, so column $1$ is unchanged. Afterwards $(3,2)$ and $(4,2)$
vanish and $(2,2)$ is non-negative.

Finally apply the step to column $3$ with the plane $(3,4)$. It mixes only rows $3$ and $4$, whose
entries in columns $1$ and $2$ are zero and stay zero. Afterwards $(4,3)$ vanishes and $(3,3)$ is
non-negative.

The result $T=G_6\cdots G_1R$, where $G_1,\dots,G_6$ are the six rotations used, is orthogonal. It
is upper triangular with non-negative diagonal entries except possibly the last. Since $T^{-1}=T^{T}$
is lower triangular, while the inverse of an invertible upper-triangular matrix is upper triangular,
$T$ is diagonal. Orthogonality makes its diagonal entries $\pm1$, so the first three equal $1$. Since
$\det G_i=1$, we have $\det T=\det R=1$. Hence
\begin{equation*}
T=\operatorname{diag}(1,1,1,\det R)=\mathbb{1}.
\end{equation*}
Therefore $R=G_1^{-1}\cdots G_6^{-1}$, and $G_{\rho\sigma}(\varphi)^{-1}=G_{\rho\sigma}(-\varphi)$ is
again a coordinate-plane rotation. Factors equal to $\mathbb{1}$ may be omitted, which gives at most
six.

(b) This is the conjugation identity \eqref{10.7:eq-plane-generators-a}: conjugation by $w$ carries
$SO(2)_{\rho\sigma}$ onto $SO(2)_{\pi(\rho)\pi(\sigma)}$.

(c) Write $G(\theta)=G_{\rho\sigma}(\theta)$. We have
$G(\theta)-G(\theta')=G(\theta')\bigl(G(\theta-\theta')-\mathbb{1}\bigr)$, and $G(\theta')$ is
orthogonal. Put $\alpha=\theta-\theta'$. The matrix $G(\alpha)-\mathbb{1}$ vanishes off the
$(\rho,\sigma)$-plane. On that plane it equals $(2-2\cos\alpha)^{1/2}$ times a rotation matrix, since
$(\cos\alpha-1)^2+\sin^2\alpha=2-2\cos\alpha$. Hence, using $|\sin x|\le|x|$,
\begin{equation}\label{10.7:eq-plane-rotations-generate-1}
\|G(\theta)-G(\theta')\|=2\,\Bigl|\sin\frac{\theta-\theta'}{2}\Bigr|\le|\theta-\theta'| .
\end{equation}

Fix an integer $N\ge1$. Since $\Gamma$ is infinite, it contains $N+1$ distinct elements
$G(\theta_0),\dots,G(\theta_N)$ with $\theta_i\in[0,2\pi)$; the $\theta_i$ are distinct. The $N$
intervals $[2\pi j/N,2\pi(j+1)/N)$, $j=0,\dots,N-1$, cover $[0,2\pi)$. By the pigeonhole principle,
two of the angles, say $\theta_i\neq\theta_j$, lie in the same interval, so
\begin{equation*}
0<|\theta_i-\theta_j|<2\pi/N .
\end{equation*}
The group $\Gamma$ contains $G(\theta_i)G(\theta_j)^{-1}=G(\theta_i-\theta_j)$ and its inverse.
Hence it contains $G(\delta)$ with $\delta=|\theta_i-\theta_j|$, so $0<\delta<2\pi/N$.

The powers $G(\delta)^k=G(k\delta)$ lie in $\Gamma$, and every $\theta\in[0,2\pi)$ is within $\delta$
of some $k\delta$ with $0\le k\le\lceil2\pi/\delta\rceil$. By
\eqref{10.7:eq-plane-rotations-generate-1}, every $G(\theta)$ lies within distance
$\delta<2\pi/N$ of an element of $\Gamma$. As $N$ is arbitrary, $\Gamma$ is dense in
$SO(2)_{\rho\sigma}$.
\end{proof}

\begin{lemma}[Exponential estimates and the product formula]\label{10.7:lem-exponential-estimates}
All matrices below are complex square matrices of one fixed size, and $\|\cdot\|$ is the operator
norm.
\begin{enumerate}
\item[(a)] If $\mu\ge0$, $\|A\|\le\mu$, $\|B\|\le\mu$ and $k\ge1$ is an integer, then
\begin{equation}\label{10.7:eq-exponential-estimates-1}
\|A^k-B^k\|\le k\mu^{k-1}\|A-B\|.
\end{equation}
Consequently $\|e^P-e^Q\|\le e^{\mu}\|P-Q\|$ whenever $\|P\|\le\mu$ and $\|Q\|\le\mu$, and
$\|e^Z-\mathbb{1}-Z\|\le\|Z\|^2$ whenever $\|Z\|\le1$.
\item[(b)] If $W_n\to0$, $k_n\in\mathbb{N}$ and $k_nW_n\to W$, then $(\mathbb{1}+W_n)^{k_n}\to e^{W}$.
\item[(c)] \emph{(Lie product formula.)} For matrices $X,Y$ one has
$e^{X/n}e^{Y/n}=\mathbb{1}+(X+Y)/n+T_n$ with $\|T_n\|\le C/n^2$ for all $n\ge1$, where $C$ does
not depend on $n$; and $(e^{X/n}e^{Y/n})^n\to e^{X+Y}$ as $n\to\infty$.
\end{enumerate}
The hypothesis $W_n\to0$ in \textup{(b)} cannot be dropped: take matrices of size four, $k_n=1$ and
$W_n=\pi A_{12}$ for every $n$, where $A_{12}=E_{12}-E_{21}$ is the antisymmetric matrix of this
chapter's notation, $E_{\rho\sigma}$ being the matrix units.
\end{lemma}

\begin{proof}
(a) Summing the telescoping terms $A^{k-j}B^j-A^{k-1-j}B^{j+1}$ over $j=0,\dots,k-1$ gives
\begin{equation*}
A^k-B^k=\sum_{j=0}^{k-1}A^{k-1-j}(A-B)B^{j}.
\end{equation*}
By submultiplicativity each summand has norm at most
$\mu^{k-1-j}\|A-B\|\mu^{j}=\mu^{k-1}\|A-B\|$, and there are $k$ summands; this is
\eqref{10.7:eq-exponential-estimates-1}. For $\|P\|,\|Q\|\le\mu$, the series
$e^P-e^Q=\sum_{k\ge1}(P^k-Q^k)/k!$ converges absolutely, and
\eqref{10.7:eq-exponential-estimates-1} applied to each term gives
\begin{equation*}
\|e^P-e^Q\|\le\sum_{k\ge1}\frac{k\mu^{k-1}}{k!}\,\|P-Q\|=e^{\mu}\|P-Q\|.
\end{equation*}
For $\|Z\|\le1$ we have $e^Z-\mathbb{1}-Z=\sum_{k\ge2}Z^k/k!$ and $\|Z^k\|\le\|Z\|^k$, so
\begin{equation*}
\|e^Z-\mathbb{1}-Z\|\le\|Z\|^2\sum_{k\ge2}\frac{\|Z\|^{k-2}}{k!}
\le\|Z\|^2\sum_{k\ge2}\frac{1}{k!}=(e-2)\|Z\|^2\le\|Z\|^2.
\end{equation*}

(b) Put $a_n:=\|W_n\|$ and $b_n:=k_na_n=\|k_nW_n\|$. Since $W_n\to0$, $a_n\le1$ for all large $n$;
we consider only such $n$. Since $k_nW_n$ converges, $\beta:=\sup_n b_n$ is finite, and
$\|W\|=\lim_n\|k_nW_n\|\le\beta$. For indices with $k_n=0$ both $(\mathbb{1}+W_n)^{k_n}$ and
$e^{k_nW_n}$ equal $\mathbb{1}$. For $k_n\ge1$ apply \eqref{10.7:eq-exponential-estimates-1} with
$A=\mathbb{1}+W_n$, $B=e^{W_n}$, $k=k_n$ and $\mu=e^{a_n}$; the hypotheses hold because
$\|\mathbb{1}+W_n\|\le1+a_n\le e^{a_n}$ and $\|e^{W_n}\|\le e^{a_n}$. Since
$B^{k_n}=e^{k_nW_n}$, and by the second estimate of (a) with $Z=W_n$,
\begin{equation*}
\begin{split}
\|(\mathbb{1}+W_n)^{k_n}-e^{k_nW_n}\|
&\le k_ne^{(k_n-1)a_n}\|\mathbb{1}+W_n-e^{W_n}\|\\
&\le k_ne^{b_n}a_n^2=b_ne^{b_n}a_n\le\beta e^{\beta}a_n\longrightarrow0 .
\end{split}
\end{equation*}
Moreover $\|e^{k_nW_n}-e^W\|\le e^{\beta}\|k_nW_n-W\|\to0$ by the first consequence in (a) with
$\mu=\beta$. Together, $(\mathbb{1}+W_n)^{k_n}\to e^W$.

In the example after the statement, $k_nW_n=\pi A_{12}=:W$ for all $n$ and
$(\mathbb{1}+W_n)^{k_n}=\mathbb{1}+\pi A_{12}$. Since $A_{12}^2=-(E_{11}+E_{22})$, the power series
gives
\begin{equation*}
e^{\pi A_{12}}=\mathbb{1}+(\cos\pi-1)(E_{11}+E_{22})+\sin\pi\,A_{12}=\operatorname{diag}(-1,-1,1,1);
\end{equation*}
this differs from $\mathbb{1}+\pi A_{12}$, whose $(1,2)$ entry is $\pi$.

(c) Write $x:=\|X\|$, $y:=\|Y\|$ and $T_n:=e^{X/n}e^{Y/n}-\mathbb{1}-(X+Y)/n$. For
$n\ge\max\{1,x,y\}$ put $R:=e^{X/n}-\mathbb{1}-X/n$ and $R':=e^{Y/n}-\mathbb{1}-Y/n$; by (a),
$\|R\|\le x^2/n^2$ and $\|R'\|\le y^2/n^2$. Expanding the product,
\begin{equation*}
T_n=\frac{XY}{n^2}+\Bigl(\mathbb{1}+\frac{X}{n}\Bigr)R'+R\Bigl(\mathbb{1}+\frac{Y}{n}\Bigr)+RR',
\end{equation*}
and since $\|\mathbb{1}+X/n\|\le2$, $\|\mathbb{1}+Y/n\|\le2$ and $n\ge1$,
\begin{equation*}
\|T_n\|\le\frac{xy+2x^2+2y^2+x^2y^2}{n^2}.
\end{equation*}
Enlarging the numerator to cover the finitely many $n<\max\{1,x,y\}$ gives a constant $C$,
independent of $n$, with $\|T_n\|\le C/n^2$ for all $n\ge1$. Now apply (b) with
$W_n:=(X+Y)/n+T_n$ and $k_n:=n$: then $e^{X/n}e^{Y/n}=\mathbb{1}+W_n$ and
\begin{equation*}
\|W_n\|\le\frac{x+y}{n}+\frac{C}{n^2}\to0,\qquad\|nW_n-(X+Y)\|=n\|T_n\|\le\frac{C}{n}\to0,
\end{equation*}
so that (b) gives $(e^{X/n}e^{Y/n})^n=(\mathbb{1}+W_n)^n\to e^{X+Y}$.
\end{proof}

\begin{lemma}[One-parameter subgroups of infinite closed subgroups]
\label{10.7:lem-one-parameter-subgroup}
Let $T \subset O(4)$ be an infinite closed subgroup. Then there is a real antisymmetric
$4\times 4$ matrix $Y \neq 0$ such that $e^{tY} \in T$ for every $t \in \mathbb{R}$; that is,
$T$ contains the one-parameter group $\{e^{tY} : t \in \mathbb{R}\}$.
\end{lemma}

\begin{proof}
Throughout, $\|\cdot\|$ is the operator norm of real $4\times 4$ matrices acting on
$\mathbb{R}^4$ with the Euclidean norm, and $\mathbb{1}$ is the identity matrix.

Since $T$ is infinite, it contains a sequence $(g_n)$ of pairwise distinct elements. Since
$O(4)$ is compact, we may pass to a subsequence and assume that $g_n \to g$ for some
$g \in O(4)$. Put $h_n := g_{n+1} g_n^{-1}$. Then $h_n \in T$ because $T$ is a subgroup, and
$h_n \neq \mathbb{1}$ because $g_{n+1} \neq g_n$. Moreover $h_n \to g g^{-1} = \mathbb{1}$,
by continuity of multiplication and inversion. Discarding finitely many terms, we may assume
$\|h_n - \mathbb{1}\| \le 1$ for all $n$.

The matrix $h_n + \mathbb{1}$ is invertible. Indeed, if an orthogonal $h$ had $-1$ as an
eigenvalue, with $hv = -v$ and $|v| = 1$, then $|(h-\mathbb{1})v| = 2$. Since also
$\|h - \mathbb{1}\| \le \|h\| + 1 = 2$, this would give $\|h-\mathbb{1}\| = 2$, contradicting
$\|h_n - \mathbb{1}\| \le 1$. Quantitatively, write
$h_n + \mathbb{1} = 2\bigl(\mathbb{1} + \tfrac12 (h_n - \mathbb{1})\bigr)$ with
$\|\tfrac12(h_n-\mathbb{1})\| \le \tfrac12$. The Neumann series then gives
\begin{equation*}
\|(h_n+\mathbb{1})^{-1}\| \le \tfrac12 \cdot \frac{1}{1-\tfrac12} = 1 .
\end{equation*}
Define the Cayley transform
\begin{equation}\label{10.7:eq-one-parameter-subgroup-1}
S_n := (h_n - \mathbb{1})(h_n + \mathbb{1})^{-1}.
\end{equation}
All matrices built from $h_n$ by sums, products and inverses are rational functions of $h_n$,
so they commute with one another.

First, $S_n$ is antisymmetric. Using $h_n^{\mathsf T} = h_n^{-1}$ we get
\begin{equation*}
S_n^{\mathsf T} = \bigl((h_n+\mathbb{1})^{-1}\bigr)^{\mathsf T}(h_n-\mathbb{1})^{\mathsf T}
= (h_n^{-1} - \mathbb{1})(h_n^{-1}+\mathbb{1})^{-1}.
\end{equation*}
Here the two factors commute. Since $h_n^{-1} - \mathbb{1} = h_n^{-1}(\mathbb{1} - h_n)$ and
$h_n^{-1} + \mathbb{1} = h_n^{-1}(\mathbb{1} + h_n)$, this equals
\begin{equation*}
(\mathbb{1} - h_n)(\mathbb{1} + h_n)^{-1} = -S_n .
\end{equation*}

Second, $S_n \neq 0$, since $h_n - \mathbb{1} \neq 0$ and $(h_n+\mathbb{1})^{-1}$ is
invertible.

Third, $\|S_n\| \le \|h_n - \mathbb{1}\|\,\|(h_n+\mathbb{1})^{-1}\| \le \|h_n - \mathbb{1}\|$,
so $\|S_n\| \to 0$.

The matrix $\mathbb{1} - S_n$ is invertible, because a real antisymmetric matrix has purely
imaginary eigenvalues, so $1$ is not an eigenvalue of $S_n$. Multiplying
\eqref{10.7:eq-one-parameter-subgroup-1} on the right by $h_n + \mathbb{1}$ gives
$S_n h_n + S_n = h_n - \mathbb{1}$, that is, $\mathbb{1} + S_n = (\mathbb{1} - S_n) h_n$. Since
$S_n$ commutes with $(\mathbb{1}-S_n)^{-1}$, it follows that
\begin{equation}\label{10.7:eq-one-parameter-subgroup-2}
h_n = (\mathbb{1} - S_n)^{-1}(\mathbb{1} + S_n)
= (\mathbb{1}-S_n)^{-1}\bigl((\mathbb{1} - S_n) + 2S_n\bigr) = \mathbb{1} + W_n,
\end{equation}
where
\begin{equation*}
W_n := 2 S_n (\mathbb{1} - S_n)^{-1}.
\end{equation*}
By \eqref{10.7:eq-one-parameter-subgroup-2}, $W_n = h_n - \mathbb{1} \to 0$. In addition, since
$S_n \to 0$ and inversion is continuous, $(\mathbb{1} - S_n)^{-1} \to \mathbb{1}$.

The matrices $S_n/\|S_n\|$ lie in the unit sphere of the space of real antisymmetric
$4\times 4$ matrices. This sphere is compact, because that space is a closed linear subspace.
Hence, along a subsequence, which we do not relabel, $S_n/\|S_n\| \to Y$ with $Y$
antisymmetric and $\|Y\| = 1$. In particular $Y \neq 0$.

Fix $t > 0$ and set $k_n := \lfloor t/(2\|S_n\|)\rfloor$. Because $\|S_n\| \to 0$, we have
$k_n \ge 1$ for all large $n$, and we discard the finitely many other terms. From
$t/(2\|S_n\|) - 1 < k_n \le t/(2\|S_n\|)$ we obtain
\begin{equation*}
t - 2\|S_n\| < 2k_n\|S_n\| \le t ,
\end{equation*}
hence $2k_n \|S_n\| \to t$. Therefore
\begin{equation*}
k_n W_n = \bigl(2k_n\|S_n\|\bigr)\,\frac{S_n}{\|S_n\|}\,(\mathbb{1} - S_n)^{-1}
\longrightarrow t\,Y\,\mathbb{1} = tY .
\end{equation*}
We have $W_n \to 0$, $k_n \in \mathbb{N}$ and $k_n W_n \to tY$, so part~(b) of
Lemma~\ref{10.7:lem-exponential-estimates} applies and gives
\begin{equation*}
h_n^{k_n} = (\mathbb{1} + W_n)^{k_n} \longrightarrow e^{tY}.
\end{equation*}
Each $h_n^{k_n}$ lies in $T$, since $T$ is a group, and $T$ is closed. Hence
$e^{tY} \in T$ for every $t > 0$. For $t = 0$ we have $e^{0} = \mathbb{1} \in T$. For $t < 0$,
$e^{tY} = (e^{-tY})^{-1}$ with $-t > 0$, so $e^{tY} \in T$. Thus $e^{tY} \in T$ for all
$t \in \mathbb{R}$.
\end{proof}

\begin{lemma}[The tangent space of a closed subgroup]\label{10.7:lem-tangent-space}
Let $H\subset O(4)$ be a closed subgroup, and put
\begin{equation*}
\mathfrak{h}:=\{X\in\mathfrak{so}(4) : e^{tX}\in H \text{ for all } t\in\mathbb{R}\}.
\end{equation*}
Then $\mathfrak{h}$ is a linear subspace of $\mathfrak{so}(4)$, and $h\mathfrak{h}h^{-1}=\mathfrak{h}$ for every $h\in H$.
\end{lemma}

\begin{proof}
\emph{Zero and scalar multiples.} Since $e^{t0}=\mathbb{1}\in H$, we have $0\in\mathfrak{h}$. If $X\in\mathfrak{h}$ and $s\in\mathbb{R}$, then $e^{t(sX)}=e^{(ts)X}\in H$ for every $t$, so $sX\in\mathfrak{h}$.

\emph{Sums.} Let $X,Y\in\mathfrak{h}$ and $t\in\mathbb{R}$. For every $n\ge 1$, the matrices $e^{tX/n}$ and $e^{tY/n}$ lie in $H$, and $H$ is a group. Hence
\begin{equation*}
\bigl(e^{tX/n}e^{tY/n}\bigr)^n\in H .
\end{equation*}
Apply the Lie product formula, Lemma~\ref{10.7:lem-exponential-estimates}(c), to the pair $tX,tY$. It gives
\begin{equation*}
\bigl(e^{tX/n}e^{tY/n}\bigr)^n\to e^{t(X+Y)}\qquad (n\to\infty).
\end{equation*}
The limit $e^{t(X+Y)}$ lies in $O(4)$. It is a limit of elements of $H$, and $H$ is closed in $O(4)$, so $e^{t(X+Y)}\in H$. This holds for every $t$, so $X+Y\in\mathfrak{h}$.

\emph{Conjugation.} Let $h\in H$ and $X\in\mathfrak{h}$. Since $h$ is orthogonal, $h^{-1}=h^{\mathsf T}$. Hence
\begin{equation*}
(hXh^{-1})^{\mathsf T}=hX^{\mathsf T}h^{-1}=-hXh^{-1},
\end{equation*}
so $hXh^{-1}\in\mathfrak{so}(4)$. Conjugating the exponential series term by term gives
\begin{equation*}
e^{t\,hXh^{-1}}=h\,e^{tX}h^{-1}.
\end{equation*}
The right-hand side is a product of three elements of $H$, so it lies in $H$ for every $t$. Therefore $hXh^{-1}\in\mathfrak{h}$, that is, $h\mathfrak{h}h^{-1}\subset\mathfrak{h}$.

Applying the same inclusion to $h^{-1}\in H$ gives $h^{-1}\mathfrak{h}h\subset\mathfrak{h}$. Conjugating this by $h$ yields $\mathfrak{h}\subset h\mathfrak{h}h^{-1}$, and hence $h\mathfrak{h}h^{-1}=\mathfrak{h}$.
\end{proof}

\begin{theorem}[The hyperoctahedral group and one rotation of infinite order generate the orthogonal group]
\label{10.7:thm-generation-density}
Write $\mathfrak{so}(4)$ for the space of real antisymmetric $4\times 4$ matrices. Let
$W_4^+ = W_4 \cap SO(4)$ be the subgroup of signed permutation matrices of determinant one. For
$\tau\in\{1,2,3,4\}$ let $s_\tau$ be the diagonal matrix that reverses the sign of the $\tau$-th
coordinate.
\begin{enumerate}
\item[(i)] A closed subgroup $H \subset O(4)$ containing $W_4$ is either finite or equal to $O(4)$.
\item[(ii)] Consequently, for every $R_0 \in O(4)$ of infinite order, the closed subgroup generated
by $W_4$ and $R_0$ is $O(4)$.
\item[(iii)] If $R_0 = G_{\alpha\beta}(\theta_0)$ with $\theta_0/\pi \notin \mathbb{Q}$, then already
the closure of the subgroup generated by $W_4^+$ and $R_0$ is $SO(4)$.
\end{enumerate}
\end{theorem}

\begin{proof}
(i) Let $H$ be infinite, and let $\mathfrak{h}$ be the set of $X \in \mathfrak{so}(4)$ with
$e^{tX} \in H$ for every real $t$. By Lemma~\ref{10.7:lem-one-parameter-subgroup}, applied to the
infinite closed subgroup $H$, there is a non-zero antisymmetric $Y$ with $e^{tY} \in H$ for all $t$;
thus $Y \in \mathfrak{h}$ and $\mathfrak{h} \neq 0$. By Lemma~\ref{10.7:lem-tangent-space},
$\mathfrak{h}$ is a linear subspace of $\mathfrak{so}(4)$ with $h\mathfrak{h}h^{-1} = \mathfrak{h}$
for every $h \in H$, in particular for every $w \in W_4 \subset H$. By
Lemma~\ref{10.7:lem-antisymmetric-irreducible}(a), the isomorphism
$\Phi\colon \mathfrak{so}(4) \to \Lambda^2\mathbb{R}^4$ satisfies
$\Phi(wXw^{-1}) = \Lambda^2(w)\Phi(X)$; hence $\Phi(\mathfrak{h})$ is a non-zero subspace of
$\Lambda^2\mathbb{R}^4$ invariant under $\Lambda^2(w)$ for all $w \in W_4$. By
Lemma~\ref{10.7:lem-antisymmetric-irreducible}(b), $\Lambda^2\mathbb{R}^4$ is irreducible under
$W_4$, so $\Phi(\mathfrak{h}) = \Lambda^2\mathbb{R}^4$, and since $\Phi$ is an isomorphism,
$\mathfrak{h} = \mathfrak{so}(4)$.

Each coordinate-plane rotation $G_{\rho\sigma}(\theta)$ is $e^{\theta X}$ for an antisymmetric $X$
whose only non-zero entries are the $(\rho,\sigma)$ and $(\sigma,\rho)$ entries, equal to $\pm 1$;
since $X \in \mathfrak{so}(4) = \mathfrak{h}$, every $G_{\rho\sigma}(\theta)$ lies in $H$. By
Lemma~\ref{10.7:lem-plane-rotations-generate}(a), every element of $SO(4)$ is a product of
coordinate-plane rotations, so $SO(4) \subset H$. Finally, $s_1 \in W_4 \subset H$ and
$\det s_1 = -1$. If $g \in O(4)$ has $\det g = -1$, then $s_1^{-1}g \in SO(4)$; hence
\begin{equation*}
O(4) = SO(4) \cup s_1 SO(4) \subset H .
\end{equation*}

(ii) Let $H$ be the closure of the subgroup generated by $W_4$ and $R_0$. It is a closed subgroup of
$O(4)$ containing $W_4$ and the powers $R_0^n$, $n \in \mathbb{Z}$. These powers are pairwise
distinct: $R_0^n = R_0^m$ with $n \neq m$ would give $R_0^{n-m} = \mathbb{1}$, so that $R_0$
would have finite order. Hence $H$ is infinite, and $H = O(4)$ by (i).

(iii) The powers $R_0^n = G_{\alpha\beta}(n\theta_0)$ form a subgroup of $SO(2)_{\alpha\beta}$. It is
infinite, since $G_{\alpha\beta}(n\theta_0) = \mathbb{1}$ forces $n\theta_0 \in 2\pi\mathbb{Z}$, which
for $\theta_0/\pi \notin \mathbb{Q}$ happens only for $n = 0$. By
Lemma~\ref{10.7:lem-plane-rotations-generate}(c), this infinite subgroup is dense in
$SO(2)_{\alpha\beta}$. Let $\bar H^+$ be the closure of the subgroup generated by $W_4^+$ and $R_0$.
Being closed and containing a dense subset of $SO(2)_{\alpha\beta}$, it contains
$SO(2)_{\alpha\beta}$.

Let $\{\rho,\sigma\}$ be any of the six coordinate planes, and choose a permutation $\pi$ with
$\pi(\{\alpha,\beta\}) = \{\rho,\sigma\}$. Since $\det P_\pi = \pm 1$ and $\det s_1 = -1$, one of
$P_\pi$ and $P_\pi s_1$ lies in $W_4^+$. By Lemma~\ref{10.7:lem-plane-rotations-generate}(b),
conjugation by that element carries $SO(2)_{\alpha\beta}$ onto $SO(2)_{\rho\sigma}$. Since
$\bar H^+$ is a group containing this element and $SO(2)_{\alpha\beta}$, it contains
$SO(2)_{\rho\sigma}$. Thus $\bar H^+$ contains every coordinate-plane rotation, and by the Givens
factorisation of Lemma~\ref{10.7:lem-plane-rotations-generate}(a), $SO(4) \subset \bar H^+$.

Conversely, $W_4^+ \subset SO(4)$ and $R_0 \in SO(4)$, and $SO(4)$ is a closed subgroup of $O(4)$.
Hence $\bar H^+ \subset SO(4)$, and $\bar H^+ = SO(4)$.
\end{proof}

\begin{proposition}[The rotation subgroup of the hyperoctahedral group does not suffice]
\label{10.7:prop-rotation-subgroup-caveat}
Write $W_4^+ := W_4 \cap SO(4)$ for the rotation subgroup of $W_4$. Identify $\mathbb{R}^4$ with the
quaternions $\mathbb{H}$ by sending $e_1, e_2, e_3, e_4$ to $1, i, j, k$, and write
$\operatorname{Im}\mathbb{H}$ for the span of $i, j, k$. For $\xi \in \operatorname{Im}\mathbb{H}$ let
$L_\xi x := \xi x$ and $R_\xi x := x\xi$, and put
$\mathfrak{l} := \{L_\xi : \xi \in \operatorname{Im}\mathbb{H}\}$ and
$\mathfrak{r} := \{R_\xi : \xi \in \operatorname{Im}\mathbb{H}\}$. Let $q := (4+3i)/5$ and
$R_0 x := qx$. Then:
\begin{enumerate}
\item[(a)] $R_0 = G_{12}(\theta_0)\,G_{34}(\theta_0)$ with $\cos\theta_0 = 4/5$ and
$\sin\theta_0 = 3/5$; $R_0$ lies in $SO(4) \cap GL_4(\mathbb{Z}[1/5])$ and has infinite order;
\item[(b)] the closed subgroup $\overline{\langle W_4^+, R_0\rangle}$ generated by $W_4^+$ and $R_0$
is a proper subgroup of $SO(4)$;
\item[(c)] $\mathfrak{so}(4) = \mathfrak{l} \oplus \mathfrak{r}$, and $\mathfrak{l}$ is a non-zero proper
$\operatorname{Ad}(W_4^+)$-invariant subspace, so that $\mathfrak{so}(4)$ is reducible under $W_4^+$;
the element $\operatorname{diag}(1,-1,-1,-1) \in W_4 \setminus W_4^+$ maps $\mathfrak{l}$ onto
$\mathfrak{r}$ under $\operatorname{Ad}$;
\item[(d)] the closed subgroup $\overline{\langle W_4, R_0\rangle}$ generated by $W_4$ and $R_0$ is
$O(4)$.
\end{enumerate}
\end{proposition}

\begin{proof}
Write $|x|$ for the quaternion norm, $|x|^2 = x\bar x$, which on $\mathbb{R}^4 = \mathbb{H}$ is the
Euclidean norm. Write $S^3 := \{a \in \mathbb{H} : |a| = 1\}$ for the group of unit quaternions. For
$a, b \in S^3$ put $\Psi(a,b)x := a x \bar b$, and let $\ell_a := \Psi(a,1)$ and
$r_b := \Psi(1,b)$, that is, $\ell_a x = ax$ and $r_b x = x\bar b$.

\emph{The double cover.} The norm is multiplicative, $|xy| = |x|\,|y|$. Hence
$|a x \bar b| = |x|$, so each $\Psi(a,b)$ is orthogonal. Moreover
\begin{equation*}
\Psi(a,b)\Psi(a',b')x = a a' x\, \bar b' \bar b = (aa')\,x\,\overline{(bb')},
\end{equation*}
so $\Psi$ is a homomorphism from $S^3 \times S^3$ to $O(4)$. Its determinant is continuous with values
in $\{\pm 1\}$ on the connected group $S^3 \times S^3$ and equals $1$ at $(1,1)$; hence it is identically
$1$ and $\Psi$ takes values in $SO(4)$. If $\Psi(a,b) = \mathbb{1}$, then $x = 1$ gives $a\bar b = 1$,
so $a = b$, and then $a x \bar a = x$ for all $x$, so that $a$ commutes with all of $\mathbb{H}$ and is
real; thus $a = b = \pm 1$, and the kernel of $\Psi$ is $\{\pm(1,1)\}$.

For surjectivity let $g \in SO(4)$ and $u := g(1) \in S^3$. Then $\ell_{\bar u} g$ lies in $SO(4)$ and
satisfies $\ell_{\bar u} g(1) = \bar u u = 1$. It therefore preserves the orthogonal complement
$\operatorname{Im}\mathbb{H}$ of $1$ and acts on it with determinant $1$, that is, as a rotation of
$\operatorname{Im}\mathbb{H} \cong \mathbb{R}^3$. By Euler's theorem this rotation is a rotation about
a unit vector $n \in \operatorname{Im}\mathbb{H}$ by some angle $\psi$. For $\varphi \in \mathbb{R}$
let $c := \cos\varphi + \sin\varphi\, n$ and $C_c x := c x \bar c$. The map $C_c$ fixes $1$ and $n$,
since $c$ commutes with $n$. For $x \in \operatorname{Im}\mathbb{H}$ orthogonal to $n$ one has
$xn = -nx$, hence $x\bar c = cx$ and
\begin{equation*}
C_c x = c^2 x = \cos 2\varphi\; x + \sin 2\varphi\; n x ,
\end{equation*}
where $nx$ is the vector product of $n$ and $x$. So $C_c$ is the rotation about $n$ by $2\varphi$.
Choosing $\varphi = \psi/2$, the maps $\ell_{\bar u} g$ and $C_c = \Psi(c,c)$ agree on
$\operatorname{Im}\mathbb{H}$ and both fix $1$, so $\ell_{\bar u} g = \Psi(c,c)$. Therefore
$g = \ell_u\Psi(c,c) = \Psi(uc, c)$, and $\Psi$ maps onto $SO(4)$.

\emph{The subgroup $N$.} Let $N := \{\ell_a : a \in S^3\}$. It is the image of the compact group
$S^3$ under the continuous homomorphism $a \mapsto \ell_a$, hence a compact subgroup of $SO(4)$. It is
normalised by all of $SO(4)$. Indeed, for $g = \Psi(a,b)$ one has $g^{-1} = \Psi(\bar a, \bar b)$, and
\begin{equation*}
g\,\ell_c\, g^{-1} x = a\bigl(c\,\bar a x b\bigr)\bar b = (a c \bar a)\,x ,
\qquad\text{that is,}\qquad g\,\ell_c\, g^{-1} = \ell_{a c \bar a}.
\end{equation*}
Since $W_4^+ \subset SO(4)$ normalises $N$, the set $W_4^+ \cdot N$ is a subgroup of $SO(4)$. It is
the union of the cosets $wN$, $w \in W_4^+$, of which there are at most $|W_4^+| = 192$. Each is
compact, so $W_4^+ \cdot N$ is compact and therefore closed.

On the other hand, $N$ has infinitely many distinct cosets in $SO(4)$. For $b, b' \in S^3$, the
equality $r_b N = r_{b'} N$ holds if and only if $r_{b'}^{-1} r_b = \Psi(1, \bar b' b)$ equals some
$\ell_a = \Psi(a,1)$. This happens if and only if $\Psi(\bar a, \bar b' b) = \mathbb{1}$, which by the
kernel computation means $\bar b' b = \pm 1$, that is, $b = \pm b'$. Since $S^3$ is infinite, the cosets
$r_b N$ are infinitely many. Consequently
\begin{equation}
\label{10.7:eq-rotation-subgroup-caveat-1}
W_4^+ \cdot N \neq SO(4).
\end{equation}

\emph{The element $R_0$.} Since $|q|^2 = (16+9)/25 = 1$, we have $R_0 = \ell_q \in N$. Its columns are
the coordinates of $q \cdot 1$, $q \cdot i$, $q \cdot j$ and $q \cdot k$:
\begin{align*}
q \cdot 1 &= \tfrac45 + \tfrac35\, i, & q \cdot i &= -\tfrac35 + \tfrac45\, i,\\
q \cdot j &= \tfrac45\, j + \tfrac35\, k, & q \cdot k &= -\tfrac35\, j + \tfrac45\, k,
\end{align*}
using $i^2 = -1$, $ij = k$ and $ik = -j$. Hence
\begin{equation*}
R_0 =
\begin{pmatrix}
4/5 & -3/5 & 0 & 0\\
3/5 & 4/5 & 0 & 0\\
0 & 0 & 4/5 & -3/5\\
0 & 0 & 3/5 & 4/5
\end{pmatrix}
= G_{12}(\theta_0)\, G_{34}(\theta_0),
\qquad \cos\theta_0 = \tfrac45,\ \ \sin\theta_0 = \tfrac35,
\end{equation*}
in the notation of \ref{10.7:not-plane-generators}. The matrix $R_0$ lies in $SO(4)$ because
$\Psi$ takes values there. Its entries lie in $\mathbb{Z}[1/5]$, and so do those of its inverse
$R_0^{\mathsf T}$; hence $R_0 \in SO(4) \cap GL_4(\mathbb{Z}[1/5])$.

Suppose $R_0^n = \mathbb{1}$ for some $n \ge 1$. Since $R_0^n$ rotates each of the two coordinate
planes by $n\theta_0$, this gives $\cos n\theta_0 = 1$ and $\sin n\theta_0 = 0$. Hence
$n\theta_0 \in 2\pi\mathbb{Z}$, and $\theta_0/\pi \in \mathbb{Q}$. By Niven's theorem, if $\theta/\pi$
and $\cos\theta$ are both rational, then $\cos\theta \in \{0, \pm\frac12, \pm 1\}$. But
$\cos\theta_0 = 4/5$ is rational and not in this set, a contradiction. So $R_0$ has infinite order,
which completes (a).

\emph{Proof of (b).} Both $W_4^+$ and $R_0 \in N$ lie in the closed subgroup $W_4^+ \cdot N$. Hence
\begin{equation*}
\overline{\langle W_4^+, R_0\rangle} \subset W_4^+ \cdot N \subset SO(4),
\end{equation*}
and the second inclusion is strict by \eqref{10.7:eq-rotation-subgroup-caveat-1}.

\emph{Proof of (c).} For $\xi \in \operatorname{Im}\mathbb{H}$, $L_\xi$ and $R_\xi$ are the
derivatives at $t = 0$ of the one-parameter groups $x \mapsto e^{t\xi}x$ and $x \mapsto x e^{t\xi}$ of
multiplications by unit quaternions. These groups lie in $SO(4)$ by the first step, so $L_\xi$ and
$R_\xi$ lie in $\mathfrak{so}(4)$. In coordinates, $L_i = A_{12} + A_{34}$ and
$R_i = A_{12} - A_{34}$, and similarly for $j$ and $k$. The maps $\xi \mapsto L_\xi$ and
$\xi \mapsto R_\xi$ are injective, since $L_\xi 1 = R_\xi 1 = \xi$, so
$\dim\mathfrak{l} = \dim\mathfrak{r} = 3$. If $L_\xi = R_\eta$, then evaluation at $1$ gives
$\xi = \eta$, and then $\xi$ commutes with all of $\mathbb{H}$; being real and imaginary, it is $0$.
Thus $\mathfrak{l} \cap \mathfrak{r} = 0$, and since $\dim\mathfrak{so}(4) = 6$, we obtain
$\mathfrak{so}(4) = \mathfrak{l} \oplus \mathfrak{r}$.

For $g = \Psi(a,b)$ the computation made for $N$ gives
$g L_\xi g^{-1} x = a\xi\bar a\,x$, that is, $g L_\xi g^{-1} = L_{a\xi\bar a}$. Here
$a\xi\bar a \in \operatorname{Im}\mathbb{H}$, because $C_a$ preserves $\operatorname{Im}\mathbb{H}$.
Since $\Psi$ is onto $SO(4)$, $\mathfrak{l}$ is $\operatorname{Ad}(SO(4))$-invariant, hence
$\operatorname{Ad}(W_4^+)$-invariant. It is non-zero and proper, since $\mathfrak{r} \neq 0$. So
$\mathfrak{so}(4)$ is reducible under $W_4^+$ alone.

Quaternion conjugation $x \mapsto \bar x$ has the matrix $\operatorname{diag}(1,-1,-1,-1)$. This is a
signed permutation matrix of determinant $-1$, hence an element of $W_4 \setminus W_4^+$; it is its own
inverse. For $\xi \in \operatorname{Im}\mathbb{H}$ one has
\begin{equation*}
\overline{\xi\,\bar x} = x\,\bar\xi = -x\xi = -R_\xi x ,
\end{equation*}
so conjugation by this element sends $L_\xi$ to $-R_\xi$ and maps $\mathfrak{l}$ onto $\mathfrak{r}$.
This is how the full group $W_4$, unlike $W_4^+$, acts irreducibly on $\mathfrak{so}(4)$, in accordance
with Lemma~\ref{10.7:lem-antisymmetric-irreducible}.

\emph{Proof of (d).} By (a), $R_0$ is an element of $O(4)$ of infinite order. By
Theorem~\ref{10.7:thm-generation-density}, the closed subgroup generated by $W_4$ and such an element is
$O(4)$.

Finally, the hyperoctahedral symmetry that enters the $O(4)$-invariance theorem of this chapter is the
whole of $W_4$, reflections included. Therefore the failure exhibited in (b) for the rotation subgroup
$W_4^+$ does not affect that theorem.
\end{proof}

\begin{lemma}[The fixed group of a bounded separated functional]\label{10.7:lem-fixed-group}
Let $\|\cdot\|$ denote the operator norm of a real $4\times 4$ matrix and $\bar B(0,\rho)$ the closed
ball of radius $\rho$ about the origin of $\mathbb{R}^4$. The action $Q\cdot f$ of $O(4)$ on test
functions, the support radius $\rho_f$, the modulus of continuity $m_u(\delta)$ and the norm
$\|\cdot\|_{\sharp}$, as well as the admissible tuples of test functions and the admissible
combinations of such tuples in one slot, are those fixed in
Definition~\ref{10.7:not-test-function-moduli}.
\begin{enumerate}
\item[(a)] Let $Q\in O(4)$ and $f\in C^1_c$. Then
\begin{equation}\label{10.7:eq-fixed-group-1}
\begin{aligned}
\|Q\cdot f-f\|_\infty &\le \|\nabla f\|_\infty\,\|Q-\mathbb{1}\|\,\rho_f,\\
\operatorname{Lip}(Q\cdot f-f) &\le \|Q-\mathbb{1}\|\,\|\nabla f\|_\infty
+m_{\nabla f}\bigl(\|Q-\mathbb{1}\|\,\rho_f\bigr),
\end{aligned}
\end{equation}
so that $\|Q\cdot f-f\|_{\sharp}\to 0$ as $Q\to\mathbb{1}$.
\item[(b)] Let $\omega$ be a real antisymmetric $4\times4$ matrix, $R^{\omega}_t=e^{t\omega}$, and
$f\in C^2_c$. Then $\|(R^{\omega}_t\cdot f-f)/t-\delta_{\omega}f\|_{\sharp}\to0$ as $t\to0$.
\item[(c)] Let $\mathcal{S}$ be a real functional defined on admissible tuples
$\vec h=(h_1,\dots,h_n)$ of test functions of class $C^1_c$, additive in each slot on admissible
combinations, and suppose that there are constants $C(\mathsf{d},r)$ such that
\begin{equation}\label{10.7:eq-fixed-group-2}
|\mathcal{S}(\vec h)|\le C(\mathsf{d},r)\prod_{i=1}^n\|h_i\|_{\sharp}
\end{equation}
whenever the supports of the $h_i$ are pairwise at distance at least $\mathsf{d}$ and each $h_i$ is
supported in some ball of radius $r$. Then $Q\mapsto\mathcal{S}(Q\cdot\vec f)$ is continuous on
$O(4)$ for every admissible $\vec f$: precisely, if $\operatorname{supp}f_i\subset\bar B(0,\rho)$ for
all $i$, the supports are pairwise at distance at least $\mathsf{d}$, and $Q,Q_0\in O(4)$ satisfy
$\|Q-Q_0\|\le\mathsf{d}/(2\rho)$, then
\begin{equation}\label{10.7:eq-fixed-group-3}
|\mathcal{S}(Q\cdot\vec f)-\mathcal{S}(Q_0\cdot\vec f)|
\le C(\mathsf{d}/2,\rho)\sum_{i=1}^n\|(Q_0^{-1}Q)\cdot f_i-f_i\|_{\sharp}
\prod_{j\ne i}\|f_j\|_{\sharp}.
\end{equation}
The set
\begin{equation*}
\operatorname{Fix}(\mathcal{S}):=\{R\in O(4):\ \mathcal{S}(R\cdot\vec f)=\mathcal{S}(\vec f)\
\text{for all admissible }\vec f\,\}
\end{equation*}
is a closed subgroup of $O(4)$. If
$W_4\subset\operatorname{Fix}(\mathcal{S})$ and $\operatorname{Fix}(\mathcal{S})$ contains one
element of infinite order, of either determinant, then $\operatorname{Fix}(\mathcal{S})=O(4)$.
\end{enumerate}
\end{lemma}

\begin{proof}
Recall that $(Q\cdot f)(x)=f(Q^{-1}x)$. Hence, for $P,Q_0\in O(4)$,
$Q_0\cdot(P\cdot f)=(Q_0P)\cdot f$, the action is linear
in $f$, and $\|P\cdot f\|_\infty=\|f\|_\infty$, $\operatorname{Lip}(P\cdot f)=\operatorname{Lip}f$,
so that $\|P\cdot f\|_{\sharp}=\|f\|_{\sharp}$, because $P^{-1}$ is an isometry of $\mathbb{R}^4$ onto
itself.

(a) Since $Q^{-1}x-x=-Q^{-1}(Q-\mathbb{1})x$ and $Q^{-1}$ is orthogonal,
$|Q^{-1}x-x|=|(Q-\mathbb{1})x|\le\|Q-\mathbb{1}\|\,|x|$; thus $|x|\le\rho_f$ implies
$|Q^{-1}x-x|\le\|Q-\mathbb{1}\|\rho_f$. If $|x|>\rho_f$, then also $|Q^{-1}x|=|x|>\rho_f$; both
points lie in the open set $\{|y|>\rho_f\}$ on which $f$ vanishes, so $f$ and $\nabla f$ vanish at
$x$ and at $Q^{-1}x$. If $|x|\le\rho_f$, the mean value theorem on the segment from $x$ to $Q^{-1}x$
gives
\begin{equation*}
|f(Q^{-1}x)-f(x)|\le\|\nabla f\|_\infty|Q^{-1}x-x|\le\|\nabla f\|_\infty\|Q-\mathbb{1}\|\rho_f .
\end{equation*}
This is the first bound of \eqref{10.7:eq-fixed-group-1}. Next, by the chain rule
$\nabla(Q\cdot f)(x)=(Q^{-1})^{\mathsf T}\nabla f(Q^{-1}x)=Q\nabla f(Q^{-1}x)$, hence
\begin{equation*}
\nabla(Q\cdot f-f)(x)=(Q-\mathbb{1})\nabla f(Q^{-1}x)+\bigl[\nabla f(Q^{-1}x)-\nabla f(x)\bigr].
\end{equation*}
The first term has modulus at most $\|Q-\mathbb{1}\|\|\nabla f\|_\infty$. The second vanishes if
$|x|>\rho_f$, and if $|x|\le\rho_f$ it is at most $m_{\nabla f}(\|Q-\mathbb{1}\|\rho_f)$, since
$|Q^{-1}x-x|\le\|Q-\mathbb{1}\|\rho_f$. As $Q\cdot f-f$ is $C^1$ on the convex set
$\mathbb{R}^4$, its Lipschitz constant is at most the supremum of its gradient, which gives the
second bound. Finally $\nabla f$ is continuous with compact support, hence uniformly continuous, so
$m_{\nabla f}(\delta)\to0$ as $\delta\to0$; both right-hand sides of \eqref{10.7:eq-fixed-group-1}
tend to $0$ as $Q\to\mathbb{1}$, and so does $\|Q\cdot f-f\|_{\sharp}$. Only $f\in C^1_c$ was used.

(b) This is Lemma~\ref{10.7:lem-rotated-differentiability}, applied to $f\in C^2_c$; it is for
differentiability in $\|\cdot\|_{\sharp}$ that the class $C^2_c$ is used, whereas for the continuity
in (a) the class $C^1_c$ suffices.

(c) Let $A_i=\operatorname{supp}f_i$, and for $0\le i\le n$ put
$\vec f^{[i]}=(Q_0\cdot f_1,\dots,Q_0\cdot f_i,Q\cdot f_{i+1},\dots,Q\cdot f_n)$, so that
$\vec f^{[0]}=Q\cdot\vec f$ and $\vec f^{[n]}=Q_0\cdot\vec f$. The tuples $\vec f^{[i-1]}$ and
$\vec f^{[i]}$ differ only in the $i$-th slot, where
$Q\cdot f_i=(Q\cdot f_i-Q_0\cdot f_i)+Q_0\cdot f_i$. For $1\le i\le n$ put
\begin{equation*}
\vec h^{(i)}=(Q_0\cdot f_1,\dots,Q_0\cdot f_{i-1},\,Q\cdot f_i-Q_0\cdot f_i,\,Q\cdot f_{i+1},
\dots,Q\cdot f_n).
\end{equation*}
The $j$-th entry of $\vec f^{[i]}$ is supported in $Q_0A_j$ or $QA_j$; the $i$-th entry of
$\vec h^{(i)}$ is supported in $QA_i\cup Q_0A_i$, the $j$-th in $Q_0A_j$ or
$QA_j$; all lie in $\bar B(0,\rho)$, since $Q$ and $Q_0$ preserve $|x|$. For $i\ne j$, $x\in A_i$,
$y\in A_j$ and $P,P'\in\{Q,Q_0\}$: if $P=P'$, then $|Px-P'y|=|x-y|\ge\mathsf{d}$; if $P\ne P'$, then
\begin{equation*}
|Px-P'y|\ge|P'x-P'y|-|(P-P')x|\ge\mathsf{d}-\|Q-Q_0\|\rho\ge\mathsf{d}/2 .
\end{equation*}
So the supports of the entries of every $\vec f^{[i]}$ and of every $\vec h^{(i)}$ are pairwise at
distance at least $\mathsf{d}/2$ inside $\bar B(0,\rho)$; hence every $\vec f^{[i]}$ and every
$\vec h^{(i)}$ is admissible with separation at least $\mathsf{d}/2$, additivity in the $i$-th slot
applies, and telescoping gives
\begin{equation}\label{10.7:eq-fixed-group-4}
\mathcal{S}(Q\cdot\vec f)-\mathcal{S}(Q_0\cdot\vec f)=\sum_{i=1}^n\mathcal{S}(\vec h^{(i)}).
\end{equation}
By \eqref{10.7:eq-fixed-group-2} with $(\mathsf{d}/2,\rho)$, $|\mathcal{S}(\vec h^{(i)})|$ is at most
$C(\mathsf{d}/2,\rho)$ times the product of the $\|\cdot\|_{\sharp}$-norms of the entries. Now
$\|Q_0\cdot f_j\|_{\sharp}=\|Q\cdot f_j\|_{\sharp}=\|f_j\|_{\sharp}$, and
$Q\cdot f_i-Q_0\cdot f_i=Q_0\cdot\bigl((Q_0^{-1}Q)\cdot f_i-f_i\bigr)$, whose norm is
$\|(Q_0^{-1}Q)\cdot f_i-f_i\|_{\sharp}$. Summing over $i$ in \eqref{10.7:eq-fixed-group-4} gives
\eqref{10.7:eq-fixed-group-3}. Since $\|Q_0^{-1}Q-\mathbb{1}\|=\|Q_0^{-1}(Q-Q_0)\|=\|Q-Q_0\|$,
part (a) applied to $Q_0^{-1}Q$ and each $f_i\in C^1_c$ shows that every term of the finite sum in
\eqref{10.7:eq-fixed-group-3} tends to $0$ as $Q\to Q_0$; this is the continuity at $Q_0$.

For the fixed group: $\mathbb{1}\in\operatorname{Fix}(\mathcal{S})$. If
$R,R'\in\operatorname{Fix}(\mathcal{S})$ and $\vec f$ is admissible, then $R'\cdot\vec f$ and
$R^{-1}\cdot\vec f$ are admissible, and
\begin{align*}
\mathcal{S}((RR')\cdot\vec f)&=\mathcal{S}(R\cdot(R'\cdot\vec f))=\mathcal{S}(R'\cdot\vec f)
=\mathcal{S}(\vec f),\\
\mathcal{S}(R^{-1}\cdot\vec f)&=\mathcal{S}(R\cdot(R^{-1}\cdot\vec f))=\mathcal{S}(\vec f).
\end{align*}
In the first line, the first equality is the composition rule for the action, the second is
$R\in\operatorname{Fix}(\mathcal{S})$ applied to the admissible tuple $R'\cdot\vec f$, and the third
is $R'\in\operatorname{Fix}(\mathcal{S})$. In the second line, the first equality is
$R\in\operatorname{Fix}(\mathcal{S})$ applied to the admissible tuple $R^{-1}\cdot\vec f$, and the
second is the composition rule. So
$\operatorname{Fix}(\mathcal{S})$ is a subgroup. For each admissible $\vec f$ the set
$\{R:\mathcal{S}(R\cdot\vec f)=\mathcal{S}(\vec f)\}$ is closed, being the zero set of a continuous
function of $R$; $\operatorname{Fix}(\mathcal{S})$ is their intersection, hence closed. If
$W_4\subset\operatorname{Fix}(\mathcal{S})$, Theorem~\ref{10.7:thm-generation-density}(i) says that
$\operatorname{Fix}(\mathcal{S})$ is finite or equal to $O(4)$; an element of infinite order makes
it infinite, so $\operatorname{Fix}(\mathcal{S})=O(4)$.
\end{proof}

\begin{corollary}[The fixed group of a subsequential limit]\label{10.7:cor-fixed-group-limit}
Assume the per-ball volume-free bound and the infinite-volume theorem, and let $S^T_n$ be any
subsequential limit furnished by the infinite-volume theorem, taken along the diagonal as in
Lemma~\ref{10.7:lem-equicontinuity}. Then Lemma~\ref{10.7:lem-fixed-group}(c) applies to
$S^T_n$ with $C:=C^{\infty}_n$. In particular:
\begin{enumerate}
\item[(a)] for every $\rho>0$ and every admissible $C^1_c$ tuple $\vec f$ with supports in
$B(0,\rho)$ the map $R\mapsto S^T_n(R\cdot\vec f)$ is
continuous on $O(4)$. The same holds, uniformly in $j$, for the functionals $S_j$ for all
$j\ge j_1(n,\rho)$, where $j_1(n,\rho)$ is the floor index of the per-ball volume-free bound;
\item[(b)] $S^T_n$ is invariant under $W_4$, by the infinite-volume theorem; if it is moreover
invariant under one rotation of infinite order, it is invariant under $O(4)$. It suffices, for
instance, that $S^T_n$ be invariant under the full one-parameter group
$\{R^{\omega}_t\}_{t\in\mathbb{R}}$ for one $\omega\neq0$, as the Ward-identity route of this
chapter provides, or under one commensurable rotation of infinite order, as in the route through
commensurable rotations;
\item[(c)] for every admissible $C^2_c$ tuple $\vec f$ and every $\omega\neq0$ the map
$t\mapsto S^T_n(R^{\omega}_t\cdot\vec f)$ is differentiable on $\mathbb{R}$, with derivative
\begin{equation}\label{10.7:eq-fixed-group-limit-1}
\frac{d}{dt}\,S^T_n(R^{\omega}_t\cdot\vec f)
=\sum_{i=1}^n S^T_n\bigl(R^{\omega}_t\cdot f_1,\dots,R^{\omega}_t\cdot\delta_{\omega}f_i,\dots,
R^{\omega}_t\cdot f_n\bigr).
\end{equation}
\end{enumerate}
\end{corollary}

\begin{proof}
By Lemma~\ref{10.7:lem-equicontinuity}(a), $S^T_n$ exists as a limit along the diagonal on every
admissible $C^1_c$ tuple. It is real, symmetric and additive on admissible combinations.

Fix $\rho>0$. The per-ball volume-free bound holds for the functionals $S_j$, on tuples supported
in $B(0,\rho)$, with one constant for all
$j\ge j_1(n,\rho)$. Passing to the limit along the diagonal, this bound supplies the constant
$C:=C^{\infty}_n$ for $S^T_n$ in Lemma~\ref{10.7:lem-fixed-group}(c).

Since the constant does not depend on $j$, the same lemma applies with the same $C$ to each $S_j$,
$j\ge j_1(n,\rho)$. The equicontinuity estimate of Lemma~\ref{10.7:lem-equicontinuity}(b) holds
in the same range of $j$. Lemma~\ref{10.7:lem-fixed-group}(c) then gives the continuity in (a),
uniformly in $j$ at finite cutoff.

It also gives (b), once $W_4$-invariance of $S^T_n$ is known. That invariance is part of the
infinite-volume theorem.

For (c), choose $\mathcal{R}$ and $\mathsf{d}>0$ with
$\vec f\in\mathcal{T}^{(2)}_n(\mathcal{R},\mathsf{d})$; such a pair exists because $\vec f$ is a
$C^2_c$ tuple with compact, pairwise separated supports. Fix $t\in\mathbb{R}$ and put
$\vec h:=R^{\omega}_t\cdot\vec f$, where
$(Q\cdot f)(x)=f(Q^{-1}x)$. Each $R^{\omega}_t$ is an isometry of $\mathbb{R}^4$ fixing the origin.
It therefore maps $B(0,\mathcal{R}/2)$ onto itself and preserves the distances between supports.
Hence $\vec h\in\mathcal{T}^{(2)}_n(\mathcal{R},\mathsf{d})$.

The group law gives $R^{\omega}_{t+s}\cdot\vec f=R^{\omega}_s\cdot\vec h$. By the differentiability
of the limit along one-parameter groups (Lemma~\ref{10.7:lem-limit-differentiability}), applied
to $\vec h$, we obtain
\begin{equation*}
\lim_{s\to0}\frac{S^T_n(R^{\omega}_{t+s}\cdot\vec f)-S^T_n(R^{\omega}_t\cdot\vec f)}{s}
=c(\vec h,\omega)
=\sum_{i=1}^n S^T_n(h_1,\dots,\delta_{\omega}h_i,\dots,h_n).
\end{equation*}

It remains to check that $\delta_{\omega}(R^{\omega}_t\cdot f)=R^{\omega}_t\cdot\delta_{\omega}f$.
Write $y:=e^{-t\omega}x$ and use that $e^{-t\omega}$ commutes with $\omega$. Since $\omega$ is
antisymmetric, $e^{t\omega}$ is orthogonal and $(e^{-t\omega})^{\mathsf T}=e^{t\omega}$. The chain
rule gives
$\nabla(R^{\omega}_t\cdot f)(x)=(e^{-t\omega})^{\mathsf T}\nabla f(y)=e^{t\omega}\nabla f(y)$,
and then
\begin{equation*}
\delta_{\omega}(R^{\omega}_t\cdot f)(x)
=-(\omega x)\cdot e^{t\omega}\nabla f(y)
=-(e^{-t\omega}\omega x)\cdot\nabla f(y)
=-(\omega y)\cdot\nabla f(y).
\end{equation*}
The right-hand side equals $(R^{\omega}_t\cdot\delta_{\omega}f)(x)$. Substituting this identity
into the limit above gives \eqref{10.7:eq-fixed-group-limit-1}.
\end{proof}

\begin{remark}[Commensurable rotations compatible with the half-offset lattices]
\label{10.7:rem-commensurable-rotations}
Let $L$ be odd. Write $\mathbb{Z}[1/L]$ for the ring of rationals $n/L^{e}$ with $n\in\mathbb{Z}$
and $e\ge 0$, and $M_4(\mathbb{Z})$ and $M_4(\mathbb{Z}[1/L])$ for the $4\times4$ matrices with entries in
$\mathbb{Z}$, respectively in $\mathbb{Z}[1/L]$. Let $\mathbb{L}_t=L^{-t}(\mathbb{Z}+\tfrac12)^4$ ($t\ge 0$)
be the half-offset lattices.
An element of $O(4)\cap GL_4(\mathbb{Z}[1/L])$ is called a \emph{commensurable rotation}. We identify
$\mathbb{R}^4$ with the quaternions $\mathbb{H}$ through the basis $(1,i,j,k)$, write $\bar q$ for the
conjugate of $q\in\mathbb{H}$, $N(q)=q\bar q$ for its norm, $\operatorname{Im}\mathbb{H}$ for the span of
$i,j,k$, and call $q$ an integer quaternion if its four coordinates are integers. For $\rho<\sigma$ and
$\theta\in\mathbb{R}$, $G_{\rho\sigma}(\theta)$ is the rotation by $\theta$ in the plane of the $\rho$-th
and $\sigma$-th coordinate vectors, equal to the identity on the orthogonal complement of that plane.
\begin{enumerate}
\item[(a)] A matrix $A\in M_4(\mathbb{Z})$ maps $(\mathbb{Z}+\frac12)^4$ into itself if and only if every
row sum $\sum_\sigma A_{\rho\sigma}$ is odd.
\item[(a$'$)] If $s\ge 0$, $R\in O(4)$ and $L^sR\in M_4(\mathbb{Z})$, then
$R\,\mathbb{L}_t\subset\mathbb{L}_{t+s}$ for all $t\ge 0$. Moreover
\begin{equation}\label{10.7:eq-commensurable-rotations-1}
O(4)\cap GL_4(\mathbb{Z}[1/L])=O(4)\cap M_4(\mathbb{Z}[1/L])
=\bigcup_{s\ge 0}O(4)\cap L^{-s}M_4(\mathbb{Z}),
\end{equation}
so that every commensurable rotation $R$, whether or not it acts in a single coordinate plane, satisfies
$R\,\mathbb{L}_t\subset\mathbb{L}_{t+s}$ for some $s\ge 0$ and all $t\ge 0$. For even $L$ the inclusion
$R\,\mathbb{L}_t\subset\mathbb{L}_{t+s}$ can fail for $R\in O(4)\cap M_4(\mathbb{Z}[1/L])$.
\item[(b)] The matrix
\begin{equation*}
R_0^{pl}:=\frac15\begin{pmatrix}4&-3\\3&4\end{pmatrix}\oplus\begin{pmatrix}1&0\\0&1\end{pmatrix}
\end{equation*}
lies in $GL_4(\mathbb{Z}[1/L])$ if and only if $5\mid L$, and it has infinite order.
\item[(b$'$)] $O(4)\cap M_4(\mathbb{Z}[1/L])$ contains an element of infinite order of the form
$G_{\rho\sigma}(\theta)$ if and only if $L$ has a prime factor $p\equiv 1\pmod 4$.
\item[(c)] Let $q=a+v$ be an integer quaternion with $a\in\mathbb{Z}$, $v\in\operatorname{Im}\mathbb{H}$,
$N(q)=L$, $a\neq 0$ and $v\neq 0$, and put $C_q(x):=q\,x\,\bar q/L$. Then
$C_q\in SO(4)\cap L^{-1}M_4(\mathbb{Z})$; $C_q$ fixes $1$ and $v$ and rotates the plane
$v^{\perp}\cap\operatorname{Im}\mathbb{H}$ by an angle $2\alpha_q$ with
\begin{equation*}
\cos 2\alpha_q=\frac{2a^2}{L}-1\in\mathbb{Q}\setminus\{0,\pm\tfrac12,\pm1\};
\end{equation*}
$C_q$ has infinite order; and $C_q\,\mathbb{L}_t\subset\mathbb{L}_{t+1}$ for all $t\ge 0$. For $L=3$ and
$q=1+i+j$,
\begin{equation}\label{10.7:eq-commensurable-rotations-2}
3\,C_q=\begin{pmatrix}3&0&0&0\\0&1&2&2\\0&2&1&-2\\0&-2&2&-1\end{pmatrix},
\qquad \cos 2\alpha_q=-\frac13 .
\end{equation}
\item[(d)] For every odd $L\ge 3$ there is an integer quaternion $q$ as in (c).
\item[(e)] In the argument through the rotational Ward identity,
Theorem~\ref{10.7:thm-generation-density} is needed only when that identity is available for some but
not all generators $\omega$; this is the situation of Lemma~\ref{10.7:lem-one-generator}, in which the
identity is available for a single generator $\omega_0\neq0$ only.
\end{enumerate}
\end{remark}

\begin{proof}
Throughout we use the classical theorem on rational values of the cosine at rational angles:
\begin{equation}\label{10.7:eq-commensurable-rotations-3}
\theta\in\pi\mathbb{Q}\ \text{and}\ \cos\theta\in\mathbb{Q}
\quad\Longrightarrow\quad \cos\theta\in\{0,\pm\tfrac12,\pm1\}.
\end{equation}
If an orthogonal map of $\mathbb{R}^4$ leaves a plane invariant and acts on it as the rotation by $\theta$,
its $n$-th power acts there as the rotation by $n\theta$; if the map has finite order $n$, then
$n\theta\in 2\pi\mathbb{Z}$, so $\theta\in\pi\mathbb{Q}$. By
\eqref{10.7:eq-commensurable-rotations-3}, such a map has infinite order whenever
$\cos\theta\in\mathbb{Q}\setminus\{0,\pm\frac12,\pm1\}$.

(a) Let $x\in(\mathbb{Z}+\frac12)^4$. Each $2x_\sigma$ is an odd integer, so $2(Ax)_\rho$ is an integer and
\begin{equation*}
2(Ax)_\rho=\sum_\sigma A_{\rho\sigma}(2x_\sigma)\equiv\sum_\sigma A_{\rho\sigma}\pmod 2 .
\end{equation*}
A real number $y$ with $2y\in\mathbb{Z}$ lies in $\mathbb{Z}+\frac12$ exactly when $2y$ is odd. If every row
sum is odd, each $(Ax)_\rho$ therefore lies in $\mathbb{Z}+\frac12$. If the $\rho$-th row sum is even, then
for $x=(\frac12,\frac12,\frac12,\frac12)$ the number $2(Ax)_\rho$ is even, so $(Ax)_\rho\in\mathbb{Z}$ and
$Ax\notin(\mathbb{Z}+\frac12)^4$.

(a$'$) The rows of $R\in O(4)$ are orthonormal, so each row of the integer matrix $L^sR$ has squared norm
$L^{2s}$, which is odd. Since $n^2\equiv n\pmod 2$ for $n\in\mathbb{Z}$, the sum of the squares of the
entries of a row is congruent modulo $2$ to the number of its odd entries; hence every row of $L^sR$ has an
odd number of odd entries, and hence an odd row sum. By (a), $L^sR$ maps $(\mathbb{Z}+\frac12)^4$ into
itself. For $x\in\mathbb{L}_t$ put $y:=L^tx\in(\mathbb{Z}+\frac12)^4$; then
\begin{equation*}
Rx=L^{-(t+s)}\,(L^sR)\,y\in L^{-(t+s)}(\mathbb{Z}+\tfrac12)^4=\mathbb{L}_{t+s}.
\end{equation*}
For \eqref{10.7:eq-commensurable-rotations-1}: the inclusion of the first set in the second holds because
$GL_4(\mathbb{Z}[1/L])\subset M_4(\mathbb{Z}[1/L])$ by the definition of $GL_4$.
Conversely, if $R\in O(4)$ has entries in $\mathbb{Z}[1/L]$, its inverse is the transpose $R^{T}$, whose
entries also lie in $\mathbb{Z}[1/L]$, so $R\in GL_4(\mathbb{Z}[1/L])$. For the second equality, write the
sixteen entries of $R\in M_4(\mathbb{Z}[1/L])$ as $n_{\rho\sigma}/L^{e_{\rho\sigma}}$ and let $s$ be the
largest $e_{\rho\sigma}$; then $L^sR\in M_4(\mathbb{Z})$. The reverse inclusion
$L^{-s}M_4(\mathbb{Z})\subset M_4(\mathbb{Z}[1/L])$ holds because each entry $n/L^{s}$, $n\in\mathbb{Z}$,
lies in $\mathbb{Z}[1/L]$. The assertion about commensurable rotations
now follows from the first part of (a$'$).

For even $L$ we have $\frac12=(L/2)/L\in\mathbb{Z}[1/L]$, so the matrix
\begin{equation*}
R_{\mathrm{even}}:=\frac12\begin{pmatrix}1&1&1&1\\1&1&-1&-1\\1&-1&1&-1\\1&-1&-1&1\end{pmatrix},
\end{equation*}
whose rows are orthonormal, lies in $O(4)\cap M_4(\mathbb{Z}[1/L])$. It maps the point
$L^{-t}(\frac12,\frac12,\frac12,\frac12)$ of $\mathbb{L}_t$ to $L^{-t}(1,0,0,0)$, which has a vanishing
coordinate, whereas no point of any $\mathbb{L}_{t'}$ has a vanishing coordinate because
$0\notin\mathbb{Z}+\frac12$. Hence $R_{\mathrm{even}}\mathbb{L}_t\not\subset\mathbb{L}_{t'}$ for all
$t,t'\ge 0$.

(b) If $5\mid L$, then $\frac15=(L/5)/L\in\mathbb{Z}[1/L]$, so $R_0^{pl}\in M_4(\mathbb{Z}[1/L])$; it is
orthogonal, so by \eqref{10.7:eq-commensurable-rotations-1} it lies in $GL_4(\mathbb{Z}[1/L])$.
Conversely, if $\frac45=n/L^{e}$ with $n\in\mathbb{Z}$ and $e\ge 0$, then $4L^{e}=5n$, so the prime $5$
divides $4L^{e}$, hence divides $L$. The matrix $R_0^{pl}$ acts on the first coordinate plane as the
rotation by $\theta_0$ with $\cos\theta_0=\frac45\notin\{0,\pm\frac12,\pm1\}$, so it has infinite order by
the consequence of \eqref{10.7:eq-commensurable-rotations-3} stated at the beginning of the proof.

(b$'$) Suppose first that $p\mid L$ with $p$ prime and $p\equiv1\pmod 4$. We show that $p=u^2+v^2$ with
integers $u,v\ge1$. Put $c:=((p-1)/2)!$. Pairing $k$ with $p-k$ for $1\le k\le (p-1)/2$ gives
\begin{equation*}
(p-1)!\equiv\prod_{k=1}^{(p-1)/2}k\,(-k)=(-1)^{(p-1)/2}c^2\equiv c^2\pmod p,
\end{equation*}
because $(p-1)/2$ is even. By the classical congruence $(p-1)!\equiv-1\pmod p$, we get
$c^2\equiv-1\pmod p$. Let $m:=\lfloor\sqrt p\rfloor$. The $(m+1)^2>p$ pairs $(x,y)$ with
$0\le x,y\le m$ give at most $p$ residues $x-cy\bmod p$, so two distinct pairs give the same residue;
their differences $X,Y$ satisfy $|X|,|Y|\le m$, $(X,Y)\neq(0,0)$ and $X\equiv cY\pmod p$. Then
$X^2\equiv c^2Y^2\equiv-Y^2$, so $p\mid X^2+Y^2$. Since $p$ is prime it is not a square, so $m^2<p$ and
\begin{equation*}
0<X^2+Y^2\le 2m^2<2p;
\end{equation*}
therefore $X^2+Y^2=p$. Neither $X$ nor $Y$ vanishes, since $p$ is not a square;
put $u:=|X|$, $v:=|Y|$.

Now set $a:=(u^2-v^2)/p$ and $b:=2uv/p$. Since $\frac1p=(L/p)/L$, both lie in $\mathbb{Z}[1/L]$, and
\begin{equation*}
a^2+b^2=\frac{(u^2-v^2)^2+4u^2v^2}{p^2}=\frac{(u^2+v^2)^2}{p^2}=1 .
\end{equation*}
So $a=\cos\theta$, $b=\sin\theta$ for some $\theta$, and $G_{12}(\theta)\in O(4)\cap M_4(\mathbb{Z}[1/L])$.
If it had finite order, \eqref{10.7:eq-commensurable-rotations-3} would force $a\in\{0,\pm\frac12,\pm1\}$.
But $a=0$ gives $u=v$ and $p=2u^2$, even; $a=\pm\frac12$ gives $p=\pm2(u^2-v^2)$, even; $a=\pm1$ gives
$u^2-v^2=\pm(u^2+v^2)$, hence $v=0$ or $u=0$. All three contradict $p$ odd and $u,v\ge1$.

Conversely, suppose $L$ has no prime factor $\equiv1\pmod4$. Since $L$ is odd, every prime factor of $L$ is
$\equiv3\pmod4$. Let $G_{\rho\sigma}(\theta)\in M_4(\mathbb{Z}[1/L])$. Its entries include $\cos\theta$ and
$\sin\theta$, so with a common denominator $\cos\theta=A/L^s$ and $\sin\theta=B/L^s$, where
$A,B\in\mathbb{Z}$ and $s\ge0$; thus $A^2+B^2=L^{2s}$. If $AB=0$, then $\theta\in\frac{\pi}{2}\mathbb{Z}$
and $G_{\rho\sigma}(\theta)$ has order at most $4$. Suppose $AB\neq0$. Let $n_0:=\gcd(A,B)$ and write
$A=n_0A'$, $B=n_0B'$. Since $n_0^2\mid L^{2s}$,
\begin{equation*}
A'^2+B'^2=L^{2s}/n_0^2=:N'\in\mathbb{Z},\qquad \gcd(A',B')=1,
\end{equation*}
and $N'\ge2$ because $A'B'\neq0$. Let $p'$ be a prime factor of $N'$. Then $p'\mid L^{2s}$, so $p'\mid L$
and $p'\equiv3\pmod4$. If $p'\nmid A'$, let $x\equiv B'A'^{-1}\pmod{p'}$; then $p'\mid A'^2+B'^2$ gives
$x^2\equiv-1$. Since $p'$ is odd, $-1\not\equiv1$, so $x^2\not\equiv1$ and $x^4\equiv1$: $x$ has order $4$
in $(\mathbb{Z}/p'\mathbb{Z})^\times$. That group has order $p'-1\equiv2\pmod4$, and the order of an
element divides the order of the group, a contradiction. Hence $p'\mid A'$, then $p'\mid B'^2$ and
$p'\mid B'$, contradicting $\gcd(A',B')=1$. So $AB=0$, and every element of the form $G_{\rho\sigma}(\theta)$
in $O(4)\cap M_4(\mathbb{Z}[1/L])$ has finite order.

Under the same hypothesis on $L$, the descent applies block by block to elements of
$O(4)\cap M_4(\mathbb{Z}[1/L])$ preserving a splitting of $\mathbb{R}^4$ into two coordinate planes, and
shows that each such element has finite order. It does
not apply to arbitrary products of coordinate-plane rotations, which exhaust $SO(4)$; these include the
maps $C_q$ of (c), which by (c) and (d) have infinite order for every odd $L\ge3$.

(c) Since $N(q)=a^2+|v|^2=L$ and $\bar q=a-v$, we have $q\bar q=\bar qq=L$, so $q^{-1}=\bar q/L$ and
$C_q(x)=qxq^{-1}$. For $x\in\{1,i,j,k\}$ the product $qx\bar q$ is an integer quaternion, so
$LC_q\in M_4(\mathbb{Z})$. From $N(xy)=N(x)N(y)$ we get $N(qx\bar q)=L^2N(x)$, hence $|C_qx|=|x|$ and
$C_q\in O(4)$. Further $C_q(1)=q\bar q/L=1$; so $C_q$ preserves $1^{\perp}=\operatorname{Im}\mathbb{H}$. As $v$
commutes with $q=a+v$, $C_q(v)=v$.

Put $\nu:=v/|v|$ and define $\alpha_q\in(0,\pi)$ by $\cos\alpha_q=a/\sqrt L$, $\sin\alpha_q=|v|/\sqrt L$.
Then $q=\sqrt L\,\epsilon$ with $\epsilon:=\cos\alpha_q+\sin\alpha_q\,\nu$, $N(\epsilon)=1$ and
$C_q(x)=\epsilon\,x\,\bar\epsilon$. Let
$w\in\operatorname{Im}\mathbb{H}$ with $w\perp\nu$. For pure quaternions $xy=-x\cdot y+x\times y$; hence
$\nu w=\nu\times w=-w\nu$, and, since $\nu\times w\perp\nu$,
\begin{equation*}
\nu w\nu=(\nu\times w)\,\nu=(\nu\times w)\times\nu=w\,(\nu\cdot\nu)-\nu\,(w\cdot\nu)=w .
\end{equation*}
Therefore
\begin{align*}
\epsilon\,w\,\bar\epsilon
&=\cos^2\!\alpha_q\,w+\sin\alpha_q\cos\alpha_q\,(\nu w-w\nu)-\sin^2\!\alpha_q\,\nu w\nu\\
&=\cos2\alpha_q\,w+\sin2\alpha_q\,(\nu\times w).
\end{align*}
For a unit vector $w$, the vectors $w$, $\nu\times w$ form an orthonormal basis of
$v^{\perp}\cap\operatorname{Im}\mathbb{H}$, and the computation applied to both shows that $C_q$ rotates
this plane by $2\alpha_q$. In the orthonormal basis $(1,\nu,w,\nu\times w)$, $C_q$ is the identity on the
first two vectors and a rotation on the last two, so $\det C_q=1$ and
$C_q\in SO(4)\cap L^{-1}M_4(\mathbb{Z})$. Further
\begin{equation*}
\cos2\alpha_q=2\cos^2\alpha_q-1=\frac{2a^2}{L}-1\in\mathbb{Q}.
\end{equation*}
The value $0$ would give $L=2a^2$; the
value $\frac12$ would give $4a^2=3L$, so $4\mid L$; the value $-\frac12$ would give $L=4a^2$; each makes $L$
even. The value $1$ would give $a^2=L$, i.e.\ $v=0$; the value $-1$ would give $a=0$. All are excluded, so
$C_q$ has infinite order by the consequence of \eqref{10.7:eq-commensurable-rotations-3} stated at the
beginning of the proof. Finally, (a$'$) with $s=1$ gives $C_q\mathbb{L}_t\subset\mathbb{L}_{t+1}$.

For $q=1+i+j$ we have $\bar q=1-i-j$ and $N(q)=3$, and multiplying out with $ij=k$, $jk=i$, $ki=j$,
\begin{equation*}
q\,i\,\bar q=i+2j-2k,\qquad q\,j\,\bar q=2i+j+2k,\qquad q\,k\,\bar q=2i-2j-k,
\end{equation*}
while $q\,1\,\bar q=3$. These are the columns of $3C_q$, which gives
\eqref{10.7:eq-commensurable-rotations-2}. Its rows are pairwise orthogonal of squared norm $9$, with row
sums $3,5,1,-1$, all odd. The lower right $3\times3$ block has determinant $27$, so
$\det(3C_q)=81$ and $\det C_q=1$; and $\cos2\alpha_q=2\cdot1/3-1=-\frac13$.

(d) First, every odd positive integer is a sum of four integer squares. The identity $N(XY)=N(X)N(Y)$ shows
that the product of two sums of four squares is again one (the coordinates of $XY$ are integers when those
of $X,Y$ are); since $1=1^2$, it suffices to treat primes, and since $L$ is odd only odd primes $p$ occur in
its factorisation. For $0\le x\le(p-1)/2$ the residues $x^2$ are distinct (as $x^2\equiv x'^2$ forces
$x\equiv\pm x'$), and likewise the residues $-1-y^2$ for $0\le y\le(p-1)/2$; the two families have
$p+1>p$ members in total, so some $x^2\equiv-1-y^2\pmod p$. Then
\begin{equation*}
x^2+y^2+1^2+0^2=mp,\qquad 1\le m\le\frac{2\bigl((p-1)/2\bigr)^2+1}{p}<p .
\end{equation*}
Let $m_0$ be the least positive integer such that $m_0p$ is a sum of four squares; thus $m_0<p$. Suppose
$m_0>1$ and write $m_0p=x_1^2+x_2^2+x_3^2+x_4^2$.

If $m_0$ is even, an even number of the $x_i$ are odd, so after renumbering $x_1\equiv x_2$ and
$x_3\equiv x_4\pmod2$, and
\begin{equation*}
\tfrac{m_0}{2}\,p=\Bigl(\tfrac{x_1+x_2}{2}\Bigr)^2+\Bigl(\tfrac{x_1-x_2}{2}\Bigr)^2
+\Bigl(\tfrac{x_3+x_4}{2}\Bigr)^2+\Bigl(\tfrac{x_3-x_4}{2}\Bigr)^2
\end{equation*}
is a sum of four integer squares, contradicting minimality.

If $m_0\ge3$ is odd, choose $y_i\equiv x_i\pmod{m_0}$ with $|y_i|<m_0/2$. Then
$\sum y_i^2\equiv\sum x_i^2\equiv0\pmod{m_0}$, so $\sum y_i^2=m_1m_0$ with $0\le m_1<m_0$ because
$\sum y_i^2<m_0^2$. If $m_1=0$, all $x_i$ are divisible by $m_0$, so $m_0^2\mid m_0p$ and $m_0\mid p$,
impossible for $1<m_0<p$. Thus $1\le m_1<m_0$. Let $X:=x_1+x_2i+x_3j+x_4k$ and
$Y:=y_1+y_2i+y_3j+y_4k$. The coordinates of $X\bar Y$ are integer bilinear forms in the coordinates of $X$
and $Y$; since $Y\equiv X$ coordinatewise modulo $m_0$, we get
\begin{equation*}
X\bar Y\equiv X\bar X=m_0p\equiv0
\end{equation*}
coordinatewise modulo $m_0$. Hence $X\bar Y=m_0Z$ with $Z$ an integer quaternion and
\begin{equation*}
m_0^2N(Z)=N(X)N(Y)=m_0p\cdot m_0m_1,\qquad\text{so}\qquad N(Z)=m_1p ,
\end{equation*}
contradicting minimality. Hence $m_0=1$, and every odd $L$ is a sum of four squares.

We now prove (d) by strong induction on odd $L\ge3$. Write $L=\sum_{\rho}x_\rho^2$. If $L$ is not a perfect
square, at least two of the $x_\rho$ are non-zero; after renumbering, $x_1\neq0$, and
$q:=x_1+(x_2i+x_3j+x_4k)$ has $N(q)=L$, real part $a:=x_1\neq0$ and imaginary part $v\neq0$. If $L$ is a
perfect square, $\sqrt L$ is an odd integer with $3\le\sqrt L<L$, and by the induction hypothesis there is
an integer quaternion $q'=a'+v'$ with $N(q')=\sqrt L$, $a'\neq0$, $v'\neq0$. Since $v'^2=-|v'|^2$,
\begin{equation*}
q:=q'^2=(a'^2-|v'|^2)+2a'v',\qquad N(q)=N(q')^2=L .
\end{equation*}
Its imaginary part $2a'v'$ is non-zero, and its real part is non-zero, since $a'^2=|v'|^2$ would make
$\sqrt L=2a'^2$ even. The least case $L=3$ is not a square and gives $q=1+i+j$.
\end{proof}

\begin{lemma}[One generator suffices]\label{10.7:lem-one-generator}
Suppose that the hypotheses \eqref{10.7:eq-hypotheses-used-form-b} hold for one single
$\omega_0 \in \mathfrak{so}(4)$, $\omega_0 \neq 0$, and for all $n$, $\mathcal{R}$ and
$\mathsf{d}$. Then the invariance of the correlation functions under the full orthogonal group
still follows. That is, for every $n \ge 2$, every $R \in O(4)$ and every admissible tuple
$\vec f$,
\begin{equation}\label{10.7:eq-one-generator-1}
S^{T}_{n}(R\cdot\vec f) = S^{T}_{n}(\vec f).
\end{equation}
\end{lemma}

\begin{proof}
Throughout, $R\cdot\vec f$ acts componentwise by $(R\cdot f_i)(x) = f_i(R^{-1}x)$. Hence
$(RR')\cdot\vec f = R\cdot(R'\cdot\vec f)$ for $R, R' \in O(4)$. An element of $O(4)$ is an
isometry of $\mathbb{R}^4$ fixing the origin. It therefore carries admissible tuples to
admissible tuples.

\emph{Invariance along the orbit of $\omega_0$.} The hypotheses
\eqref{10.7:eq-hypotheses-used-form-b} hold for the data $(\omega_0, n, \mathcal{R}, \mathsf{d})$
for all $n$, $\mathcal{R}$, $\mathsf{d}$. Consequently Proposition~\ref{10.7:prop-anomaly-vanishes}
applies with $\omega = \omega_0$, and so does Proposition~\ref{10.7:prop-orbit-constancy-pointwise},
whose proof uses it. The latter proposition shows that for every
$\vec f \in \mathcal{T}^{(r)}_n(\mathcal{R},\mathsf{d})$ the function
$s \mapsto S^{T}_{n}(R^{\omega_0}_s\cdot\vec f)$ is constant on $\mathbb{R}$. Thus
\begin{equation*}
S^{T}_{n}(e^{t\omega_0}\cdot\vec f) = S^{T}_{n}(\vec f)
\end{equation*}
for all $t$ and all such $\vec f$, for every $\mathcal{R}$ and $\mathsf{d}$. Fix $t$. By
Proposition~\ref{10.7:prop-density-extension}, applied with $R = e^{t\omega_0}$, this identity
extends to every admissible tuple. Hence it holds for all $t$, all $n$ and all admissible
$\vec f$.

\emph{The invariance group.} Put
\begin{equation*}
G := \bigl\{R \in O(4) : S^{T}_{n}(R\cdot\vec f) = S^{T}_{n}(\vec f)
\text{ for all } n \text{ and all admissible } \vec f\bigr\}.
\end{equation*}
Clearly $\mathbb{1} \in G$.

If $R, R' \in G$ and $\vec f$ is admissible, then $R'\cdot\vec f$ is admissible, and
\begin{equation*}
S^{T}_{n}((RR')\cdot\vec f) = S^{T}_{n}(R\cdot(R'\cdot\vec f))
= S^{T}_{n}(R'\cdot\vec f) = S^{T}_{n}(\vec f).
\end{equation*}
If $R \in G$, apply the defining identity of $G$ to the admissible tuple $R^{-1}\cdot\vec f$.
This gives $S^{T}_{n}(\vec f) = S^{T}_{n}(R^{-1}\cdot\vec f)$, so $R^{-1} \in G$. Thus $G$ is a
subgroup of $O(4)$.

By Lemma~\ref{10.7:lem-orthogonal-continuity}, the map $R \mapsto S^{T}_{n}(R\cdot\vec f)$ is
continuous on $O(4)$. Hence each set
$\{R : S^{T}_{n}(R\cdot\vec f) = S^{T}_{n}(\vec f)\}$ is closed. Their intersection $G$ is
therefore closed.

By part~(c) of Lemma~\ref{10.7:lem-equicontinuity}, $W_4 \subset G$. By the first step,
$e^{t\omega_0} \in G$ for all $t \in \mathbb{R}$.

\emph{The tangent space.} Put
\begin{equation*}
\mathfrak{g} := \{X \in \mathfrak{so}(4) : e^{tX} \in G \text{ for all } t \in \mathbb{R}\}.
\end{equation*}
Three properties are needed.

\emph{It is a linear subspace.} Clearly $0 \in \mathfrak{g}$. If $X \in \mathfrak{g}$ and
$c \in \mathbb{R}$, then $e^{t(cX)} = e^{(tc)X} \in G$. Now let $X, Y \in \mathfrak{g}$. The Lie
product formula, part~(c) of Lemma~\ref{10.7:lem-exponential-estimates}, gives
\begin{equation*}
e^{t(X+Y)} = \lim_{k\to\infty}\bigl(e^{tX/k}e^{tY/k}\bigr)^k .
\end{equation*}
Each term on the right is a product of elements of $G$, hence lies in $G$. Since $G$ is closed,
the limit lies in $G$, so $X + Y \in \mathfrak{g}$.

\emph{It is invariant under conjugation by $W_4$.} Let $w \in W_4$ and $X \in \mathfrak{g}$.
Since $w$ is orthogonal, $w^{-1} = w^{\mathsf{T}}$, and
$w e^{tX} w^{-1} = e^{t\,wXw^{\mathsf{T}}}$. The left side lies in $G$ because $w$, $w^{-1}$ and
$e^{tX}$ do. Moreover $wXw^{\mathsf{T}}$ is antisymmetric. Hence $wXw^{\mathsf{T}} \in \mathfrak{g}$.

\emph{It is non-zero.} By the previous step $\omega_0 \in \mathfrak{g}$, and $\omega_0 \neq 0$.

Part~(b) of Lemma~\ref{10.7:lem-antisymmetric-irreducible}, transported to $\mathfrak{so}(4)$ by
the intertwining isomorphism $\Phi$ of part~(a) of that lemma, states that $\mathfrak{so}(4)$ is
irreducible under conjugation by $W_4$. The sign changes in $W_4$, which have determinant $-1$,
are used essentially there. A non-zero $W_4$-invariant subspace is therefore all of
$\mathfrak{so}(4)$, so $\mathfrak{g} = \mathfrak{so}(4)$.

\emph{Conclusion.} By Lemma~\ref{10.7:lem-exponential-onto}, every $R \in SO(4)$ equals $e^{X}$
for some $X \in \mathfrak{so}(4) = \mathfrak{g}$. Hence $R = e^{1\cdot X} \in G$, and
$SO(4) \subset G$.

Finally, $w_0 = \operatorname{diag}(-1,1,1,1) \in W_4 \subset G$ and
$O(4) = SO(4) \cup w_0\,SO(4)$. Since $G$ is a group, $G = O(4)$, which is
\eqref{10.7:eq-one-generator-1}.
\end{proof}

\begin{remark}[Variants of the hypotheses that change nothing]\label{10.7:rem-hypothesis-variants}
Nothing in the preceding argument changes under any of the following variants:
\begin{itemize}
\item the test-function classes are restricted to tuples of class $C^\infty_c$;
\item other norms are used in the dimension-six coefficient bound;
\item the floor indices along the diagonal are allowed to depend on the tuple;
\item the families of cut-offs $\chi$ are restricted;
\item the rotational Ward identity is delivered directly in connected form.
\end{itemize}
Write $D_j = D^{\omega}_{K_j,m_j}$ for the derivative along the orbit, taken along the diagonal.

Assume hypothesis (b2) of Corollary~\ref{10.7:cor-clustering-dichotomy}.
Then the three hypotheses together are equivalent,
by that corollary, to the statement that $\lim_j D_j \equiv 0$ on the classes. These three
hypotheses are:
\begin{itemize}
\item the rotational Ward identity;
\item the dimension-six coefficient bound;
\item the far-sphere flux hypothesis $(\mathrm{IR}\text{-}\mathrm{fl})_{\omega_0}$: the flux of the
lattice rotation current through far spheres vanishes along one sequence of radii; a sufficient
condition for it is that one far-away stress-tensor insertion decorrelates from the observables
faster than $|x|^{-4}$, uniformly in the cutoff.
\end{itemize}
The substance of the route through these three hypotheses lies therefore not in that equivalence.
It lies in their producing a splitting for which the two limit statements are provable.
\end{remark}

\begin{definition}[The dimension-four annulus functionals and the tensor-class bounds]
\label{10.7:def-annulus-functionals}
Throughout this section the rotational Ward identity in its connected form and the dimension-six
coefficient bound in its pointwise form, as stated in the hypotheses in the form used
(Definition~\ref{10.7:def-hypotheses-used-form}), are in force, together with the two further
standing conditions imposed alongside them there. We write $\varepsilon$ for the lattice spacing, $e_\mu$ for
the unit vector in the direction $\mu$, $x(p)$ for the barycentre of a plaquette $p$, and
$a_p := 1-\operatorname{Re}\operatorname{tr}U(\partial p)$. We write $U_{\rho\mu}(x)$ for the
plaquette variable of the plaquette at the lattice point $x$ spanned by the directions $\rho$ and
$\mu$. The matrix $\omega=(\omega_{\rho\sigma})$ is the generator of the one-parameter rotation
group $R^{\omega}_t=e^{t\omega}$. The tuple $\vec h=(h_1,\dots,h_n)$ is fixed, and
$\kappa_j(Y;\mathcal{O}(\vec h))$ denotes the joint cumulant of the bounded variable $Y$ with
$\mathcal{O}(h^{(m_j)}_1),\dots,\mathcal{O}(h^{(m_j)}_n)$ in the lattice measure of index $j$
of the diagonal sequence $(K_j,m_j)$. We write $g_0^{(j)}$ for the bare coupling at that index.

\begin{enumerate}
\item[(a)] \emph{The displayed forms.}
Let $\Theta^{\varepsilon}_{\nu\sigma}$ be the bounded lattice stress density of the rotational
Ward identity, with its factor $g_0^{-2}$ included. Let $\kappa_N$ be the numerical constant of
that identity; it is a number, not a cumulant. For $\chi\in C^\infty_c(\mathbb{R}^4)$ set
$\xi(y):=\chi(y)\,(\omega y)$. For each
$\mu$, with no summation over $\mu$, set
\begin{equation*}
w_\mu(y):=\frac{\xi_\mu(y+\varepsilon e_\mu)-\xi_\mu(y)}{\varepsilon}.
\end{equation*}
Further set
\begin{equation*}
P_\mu(x):=\sum_{\rho\neq\mu}a_{\rho\mu}(x),
\qquad
a_{\rho\mu}(x):=1-\operatorname{Re}\operatorname{tr}U_{\rho\mu}(x).
\end{equation*}
The dimension-four annulus functionals are
\begin{equation}
\label{10.7:eq-annulus-functionals-a}
\begin{aligned}
B^4_\chi&:=\sum_x\varepsilon^4\sum_{\rho,\sigma,\nu}\omega_{\rho\sigma}\,x_\rho\,
(\partial_\nu\chi)(x)\,\Theta^{\varepsilon}_{\nu\sigma}(x),\\
J^4_\chi&:=\kappa_N\sum_x\sum_\mu w_\mu(x)\,P_\mu(x).
\end{aligned}
\end{equation}
Here $J^4_\chi$ is the Haar-divergence annulus piece of the rotational Ward identity. It carries
no factor $g_0^{-2}$ in front.
The form of $J^4_\chi$ displayed in \eqref{10.7:eq-annulus-functionals-a} is an assumption of
this section; it is not derived here from the rotational Ward identity.

\item[(b)] \emph{The total functional and its insertion.}
Set $\mathfrak{F}_\chi:=B^4_\chi+J^4_\chi$ and
$\mathcal{B}^{\chi}_j:=\kappa_j(\mathfrak{F}_\chi;\mathcal{O}(\vec h))$.
Both formulas in \eqref{10.7:eq-annulus-functionals-a} are linear in $\chi$. For every
$\chi\in C^\infty_c(\mathbb{R}^4)$ they define bounded lattice variables.

\item[(c)] \emph{Radial cut-offs.}
We fix a function $\chi_1$ on $[0,\infty)$ with $\chi_1\equiv1$ on $[0,1]$ and
$\chi_1=0$ on $[2,\infty)$, and we choose it such that $u\mapsto\chi_1(|u|)$ lies in
$C^\infty_c(\mathbb{R}^4)$.
Set $\chi_{\mathcal{R}}:=\chi_1(|\cdot|/\mathcal{R})$. For $k\le3$ set
\begin{equation*}
A_k:=\max_{|\alpha|=k}\ \sup_u\bigl|\partial^\alpha\bigl(\chi_1(|u|)\bigr)\bigr|.
\end{equation*}

\item[(d)] \emph{The tensor-class hypotheses.}
Let $\pi$ run over the six coordinate planes. For a real Lipschitz function $g$ set
\begin{equation*}
\mathcal{O}_\pi(g):=\sum_{p\parallel\pi}g(x(p))\,a_p .
\end{equation*}
For a six-tuple $\vec g=(g_\pi)_\pi$ set
\begin{equation*}
\mathcal{T}(\vec g):=g_0^{(j)\,-2}\sum_\pi\mathcal{O}_\pi(g_\pi).
\end{equation*}
The \emph{bounded tensor-class hypothesis} for the given tuple $\vec h$ is the following. There
is $C(\vec h)$ such that
\begin{equation}
\label{10.7:eq-annulus-functionals-1}
\sup_j\bigl|\kappa_j(\mathcal{T}(\vec g);\mathcal{O}(\vec h))\bigr|
\le C(\vec h)\,\max_\pi\|g_\pi\|_{\sharp}
\end{equation}
for every closed cube $\Delta$ of side $1$ at distance at least $1$ from
$\bigcup_i\operatorname{supp}h_i$. It is required for every six-tuple $\vec g$ of real
Lipschitz functions supported in $\Delta$ with $\sum_\pi g_\pi\equiv0$.

Its weakening, the \emph{weak tensor-class hypothesis}, reads as follows. There are $C(\vec h)$
and numbers $\eta_j\to0$ such that, for every cube $\Delta$ and six-tuple $\vec g$ as above and
every $j$,
\begin{equation}
\label{10.7:eq-annulus-functionals-2}
\bigl|\kappa_j(\mathcal{T}(\vec g);\mathcal{O}(\vec h))\bigr|
\le \eta_j\,g_0^{(j)\,-2}\,C(\vec h)\,\max_\pi\|g_\pi\|_{\sharp}.
\end{equation}

\item[(e)] \emph{Decay of the stress density.}
For a Lipschitz function $g$ set
\begin{equation*}
\Theta^{\varepsilon}_{\nu\sigma}(g):=\sum_x\varepsilon^4g(x)\,\Theta^{\varepsilon}_{\nu\sigma}(x).
\end{equation*}
The \emph{stress-density decay hypothesis} is the following. There are
$C_{CL}(\vec h)$ and $\delta>0$ such that
\begin{equation}
\label{10.7:eq-annulus-functionals-3}
\sup_j\bigl|\kappa_j(\Theta^{\varepsilon}_{\nu\sigma}(g);\mathcal{O}(\vec h))\bigr|
\le C_{CL}(\vec h)\,\|g\|_{\sharp}\,D^{-4-\delta}
\end{equation}
for every cube $\Delta$ of side $1$ at distance $D\ge1$ from $\bigcup_i\operatorname{supp}h_i$.
It is required for every Lipschitz $g$ supported in $\Delta$ and every pair of indices
$\nu,\sigma$.
\end{enumerate}
\end{definition}

\begin{proof}[Proof of the assertion in \textup{(b)}]
We first prove linearity. The map $\chi\mapsto\partial_\nu\chi$ is linear. Each summand of
$B^4_\chi$ is therefore linear in $\chi$, and so is their sum.

The map $\chi\mapsto\xi=\chi\,(\omega y)$ is linear. The map $\xi\mapsto w_\mu$ is a difference
quotient, hence linear. Since $P_\mu$ does not depend on $\chi$, the functional $J^4_\chi$ is
linear in $\chi$. Consequently $\mathfrak{F}_\chi$ is linear in $\chi$. Since the joint cumulant
is multilinear in its entries, $\mathcal{B}^{\chi}_j$ is linear in $\chi$ as well.

We next prove boundedness. Let $\chi\in C^\infty_c(\mathbb{R}^4)$.

Consider first $B^4_\chi$. The summand at $x$ vanishes unless $x\in\operatorname{supp}\chi$. The
lattice of spacing $\varepsilon$ meets the compact set $\operatorname{supp}\chi$ in finitely
many points, so the sum has finitely many non-zero terms. The summand at such a point $x$ is
bounded in modulus by
\begin{equation*}
\varepsilon^4\sum_{\rho,\sigma,\nu}|\omega_{\rho\sigma}|\,|x_\rho|\,
\sup|\partial_\nu\chi|\,\sup|\Theta^{\varepsilon}_{\nu\sigma}|,
\end{equation*}
and this is finite because $\Theta^{\varepsilon}_{\nu\sigma}$ is bounded.

Consider next $J^4_\chi$. For each $\mu$, the function $\xi_\mu$ is smooth and supported in
$\operatorname{supp}\chi$. Hence $w_\mu(x)=0$ unless $x\in\operatorname{supp}\chi$ or
$x+\varepsilon e_\mu\in\operatorname{supp}\chi$. This again leaves finitely many lattice points.
By the mean value theorem, $|w_\mu(x)|\le\sup|\partial_\mu\xi_\mu|<\infty$.

Each $U_{\rho\mu}(x)$ lies in $SU(N)$, and the trace is normalised by
$\operatorname{tr}\mathbb{1}=1$. Thus $\operatorname{tr}U_{\rho\mu}(x)$ is the average of $N$
complex numbers of modulus one. It follows that $|\operatorname{Re}\operatorname{tr}U_{\rho\mu}(x)|\le1$,
and hence $0\le a_{\rho\mu}(x)\le2$. Since $P_\mu(x)$ is a sum of three such terms,
$0\le P_\mu(x)\le6$. Therefore $J^4_\chi$ is a finite sum of bounded terms.

It follows that $B^4_\chi$ and $J^4_\chi$, and hence $\mathfrak{F}_\chi$, are finite linear
combinations of bounded functions of the configuration, that is, bounded lattice variables. In
particular the joint cumulant $\mathcal{B}^{\chi}_j$ of $\mathfrak{F}_\chi$ with the bounded
observables $\mathcal{O}(h^{(m_j)}_i)$ is defined.
\end{proof}

\begin{proposition}[The rotational anomaly as a boundary flux of the annulus functional]
\label{10.7:prop-anomaly-flux}
Let $S^T_n$ denote the subsequential infinite-volume limit furnished by the infinite-volume
theorem along the diagonal sequence $(K_j,m_j)$, and write $\kappa_j(Y_1,\dots,Y_N)$ for the joint
cumulant under the lattice measure at $(K_j,m_j)$. For a bounded variable $Y$ and a tuple $\vec h$
abbreviate
\begin{equation*}
\kappa_j(Y;\mathcal{O}(\vec h))
:=\kappa_j\bigl(Y,\mathcal{O}(h_1^{(m_j)}),\dots,\mathcal{O}(h_n^{(m_j)})\bigr).
\end{equation*}
Fix $\mathcal{R}$, a real antisymmetric $4\times4$ matrix $\omega$, $n\ge2$, $\mathsf{d}$ and
$\vec h\in\mathcal{T}^{(2)}_n(\mathcal{R},\mathsf{d})$, and write
$\mathcal{B}^{\chi}_j:=\mathcal{B}^{\omega,\chi}_{K_j,m_j}(\vec h)$ and
$\mathcal{N}^{\chi}_j:=\mathcal{N}^{\omega,\chi}_{K_j,m_j}(\vec h)$. For a tuple $\vec f$ of real
$C^2_c$ functions with pairwise disjoint supports put
\begin{equation*}
c(\vec f,\omega):=\sum_{i=1}^n S^T_n(f_1,\dots,\delta_\omega f_i,\dots,f_n),
\end{equation*}
which on $\mathcal{T}^{(2)}_n(\mathcal{R},\mathsf{d})$ is the anomaly of
Lemma~\ref{10.7:lem-limit-differentiability}. For $a\in\mathbb{R}^4$ and $w\in O(4)$ put
$(\tau_a f)(x):=f(x-a)$ and $(w\cdot f)(x):=f(w^{-1}x)$, acting componentwise on tuples. We
say that $S^T_n$ is $O(4)$-invariant if $S^T_n(R\cdot\vec f)=S^T_n(\vec f)$ for every
$R\in O(4)$ and every admissible $C^1_c$ tuple $\vec f$. Taken over all $n\ge2$, this is the
conclusion of the $O(4)$-invariance theorem of this chapter for the family of limits
$(S^T_n)_{n\ge2}$ along the chosen subsequence.
\begin{enumerate}
\item[(a)] Assume the hypotheses of Proposition~\ref{10.7:prop-boundary-independence}. For
every admissible cut-off $\chi$ the limit of $\mathcal{B}^{\chi}_j$ exists and
\begin{equation}\label{10.7:eq-anomaly-flux-a}
\lim_{j\to\infty}\mathcal{B}^{\chi}_j(\vec h)=c(\vec h,\omega).
\end{equation}
In this sense the anomaly is the limiting flux of the dimension-four annulus functional
$\mathfrak{F}_\chi=B^4_\chi+J^4_\chi$ through the shell $\{\nabla\chi\ne0\}$, whichever
admissible shell is chosen.
\item[(b)] Under the same hypotheses there is no flux out of source-free regions: if $\chi$ is
admissible and $\psi\in C^\infty_c(\mathbb{R}^4)$ vanishes on $B(0,\mathcal{R})$, then
$\chi+\psi$ is admissible, $\mathfrak{F}_{\chi+\psi}=\mathfrak{F}_\chi+\mathfrak{F}_\psi$, and
\begin{equation}\label{10.7:eq-anomaly-flux-1}
\kappa_j(\mathfrak{F}_\psi;\mathcal{O}(\vec h))
=\mathcal{B}^{\chi+\psi}_j-\mathcal{B}^{\chi}_j
=\mathcal{N}^{\chi}_j-\mathcal{N}^{\chi+\psi}_j\longrightarrow0 .
\end{equation}
\item[(c)] Using only the per-ball volume-free bound and the infinite-volume theorem: on all
real $C^2_c$ tuples with pairwise disjoint supports, $c(\cdot,\omega)$ is real, multilinear,
symmetric and linear in $\omega$; it is $W_4$-equivariant,
\begin{equation}\label{10.7:eq-anomaly-flux-2}
c(w\cdot\vec f,\,w\omega w^{-1})=c(\vec f,\omega)\qquad(w\in W_4),
\end{equation}
and translation invariant, $c(\tau_a\vec f,\omega)=c(\vec f,\omega)$ for all $a\in\mathbb{R}^4$.
\item[(d)] Using only the per-ball volume-free bound and the infinite-volume theorem: $S^T_n$
is $O(4)$-invariant if and only if $c(\vec f,\omega)=0$ for every real $C^2_c$ tuple $\vec f$
with pairwise disjoint supports and every real antisymmetric $\omega$.
\item[(e)] Let $\rho_0>0$ with $\bigcup_i\operatorname{supp}h_i\subset B(0,\rho_0)$, and assume
the hypotheses of Proposition~\ref{10.7:prop-boundary-independence} for the data
$(\omega,n,\mathcal{R},\mathsf{d})$ and, for every $\mathcal{R}'\ge2\rho_0$, for the data
$(\omega,n,\mathcal{R}',\mathsf{d})$. Let
$(\chi_{\mathcal{R}'})_{\mathcal{R}'\ge2\rho_0}$ be any family with $\chi_{\mathcal{R}'}$
admissible at radius $\mathcal{R}'$. Then for every $\mathcal{R}'\ge2\rho_0$ the limit
$\lim_j|\mathcal{B}^{\chi_{\mathcal{R}'}}_j(\vec h)|$ exists and equals $|c(\vec h,\omega)|$.
Consequently the following are equivalent:
\begin{enumerate}
\item[(e1)] the double limit
\begin{equation}\label{10.7:eq-anomaly-flux-3}
\lim_{\mathcal{R}'\to\infty}\ \limsup_{j\to\infty}
\bigl|\mathcal{B}^{\chi_{\mathcal{R}'}}_j(\vec h)\bigr|=0;
\end{equation}
by the first assertion the inner limit does not depend on $\mathcal{R}'\ge2\rho_0$, so
\eqref{10.7:eq-anomaly-flux-3} holds as soon as the inner limit superior vanishes along one
sequence of radii; for $\omega\ne0$, (e1) is the analogue, for the total annulus functional
$\mathfrak{F}=B^4+J^4$ at the datum $\vec h$, of the far-sphere flux hypothesis
$(\mathrm{IR}\text{-}\mathrm{fl})_{\omega_0}$ with $\omega_0=\omega$ (the flux of the lattice
rotation current through far spheres vanishes along one sequence of radii);
\item[(e2)] $\lim_j\mathcal{B}^{\chi}_j(\vec h)=0$ for one cut-off $\chi$ admissible at radius
$\mathcal{R}$;
\item[(e3)] $\lim_j\mathcal{B}^{\chi}_j(\vec h)=0$ for every cut-off $\chi$ admissible at radius
$\mathcal{R}$;
\item[(e4)] $c(\vec h,\omega)=0$.
\end{enumerate}
If moreover the annulus terms have the forms displayed in
\eqref{10.7:eq-annulus-functionals-a}, the uniform boundedness of the tensor-class correlations
displayed there holds, and
$(\chi_{\mathcal{R}'})$ is the radial family for which \eqref{10.7:eq-divergence-vanishes-b}
is stated, then for every fixed $\mathcal{R}'\ge\max\{12,2\rho_0\}$
\begin{equation}\label{10.7:eq-anomaly-flux-4}
\lim_{j\to\infty}\kappa_j(J^4_{\chi_{\mathcal{R}'}};\mathcal{O}(\vec h))=0,
\qquad
\lim_{j\to\infty}\kappa_j(B^4_{\chi_{\mathcal{R}'}};\mathcal{O}(\vec h))=c(\vec h,\omega),
\end{equation}
and the equivalences hold with $\kappa_j(B^4_{\chi_{\mathcal{R}'}};\mathcal{O}(\vec h))$ in
place of
\begin{equation*}
\mathcal{B}^{\chi_{\mathcal{R}'}}_j(\vec h)
=\kappa_j(\mathfrak{F}_{\chi_{\mathcal{R}'}};\mathcal{O}(\vec h))
\end{equation*}
in \eqref{10.7:eq-anomaly-flux-3}; for $\omega\ne0$ this $B^4$-only double limit with radial
cut-offs is the far-sphere flux hypothesis $(\mathrm{IR}\text{-}\mathrm{fl})_{\omega_0}$, at
$\omega_0=\omega$ and the datum $\vec h$, in which the flux through far spheres is read through
the stress-tensor density $\Theta^{\varepsilon}$; by the second limit in
\eqref{10.7:eq-anomaly-flux-4} its inner limit does not depend on $\mathcal{R}'$, so its
vanishing along one sequence of radii is equivalent to the double limit.
\end{enumerate}
\end{proposition}

\begin{proof}
\emph{Part (a).} Since $\vec h\in\mathcal{T}^{(2)}_n(\mathcal{R},\mathsf{d})$,
Proposition~\ref{10.7:prop-boundary-independence} with $r=2$ applies and gives
$\mathcal{B}^{\omega,\chi}_{K_j,m_j}(\vec h)\to c(\vec h,\omega)$ for every admissible $\chi$,
where $c(\vec h,\omega)$ is formed with the limit $S^T_n$ given by the infinite-volume theorem.
This is \eqref{10.7:eq-anomaly-flux-a}.

\emph{Part (b).} The weights of the displayed forms \eqref{10.7:eq-annulus-functionals-a} are
linear in the cut-off: $B^4_\chi$ carries the factor $x_\rho(\partial_\nu\chi)(x)$, and $J^4_\chi$
carries the forward differences $w_\mu$ of $\xi=\chi\cdot(\omega y)$, both linear in $\chi$,
while the densities $\Theta^{\varepsilon}_{\nu\sigma}$ and $P_\mu$ do not depend on $\chi$.
Hence $\mathfrak{F}_{\chi+\psi}=\mathfrak{F}_\chi+\mathfrak{F}_\psi$. Since $\psi$ is smooth,
compactly supported and vanishes on $B(0,\mathcal{R})$, the cut-off $\chi+\psi$ coincides with
$\chi$ on $B(0,\mathcal{R})$ and is again admissible. The boundary term $\mathcal{B}^{\chi}_j$ is
the connected correlation
$\kappa_j(\mathfrak{F}_\chi;\mathcal{O}(\vec h))$ of the annulus functional with the
observables. The cumulant is linear in its first entry, so
\begin{equation*}
\kappa_j(\mathfrak{F}_\psi;\mathcal{O}(\vec h))
=\kappa_j(\mathfrak{F}_{\chi+\psi};\mathcal{O}(\vec h))-\kappa_j(\mathfrak{F}_\chi;\mathcal{O}(\vec h))
=\mathcal{B}^{\chi+\psi}_j-\mathcal{B}^{\chi}_j .
\end{equation*}
By Proposition~\ref{10.7:prop-boundary-independence} applied to the admissible pair $\chi$,
$\chi+\psi$, this difference equals $\mathcal{N}^{\chi}_j-\mathcal{N}^{\chi+\psi}_j$ and tends
to zero. This is \eqref{10.7:eq-anomaly-flux-1}.

Suppose the rotational Ward identity, Theorem~\ref{10.1:thm-rotational-identity}, were available
for non-admissible cut-offs, with no term
on the observable side when the cut-off vanishes near the sources. Then
\eqref{10.7:eq-anomaly-flux-1} would say literally that the Ward identity for the deformation
generated by $\psi$ has no observable term. The argument above does not use that form.

\emph{Part (c).} Let $\vec f$ be a real $C^2_c$ tuple with pairwise disjoint supports. Its
supports are compact, so $\vec f\in\mathcal{T}^{(2)}_n(\mathcal{R}_1,\mathsf{d}_1)$ for some
$\mathcal{R}_1,\mathsf{d}_1>0$. For each $i$, $\delta_\omega f_i=-(\omega x)\cdot\nabla f_i$ is
a real $C^1_c$ function with $\operatorname{supp}\delta_\omega f_i\subset\operatorname{supp}f_i$.
Hence each tuple $(f_1,\dots,\delta_\omega f_i,\dots,f_n)$ is an admissible $C^1_c$ tuple, and
$c(\vec f,\omega)$ is defined.

The functional $S^T_n$ is real on real tuples, multilinear and symmetric, being the limit of
joint cumulants
of the real bounded variables $\mathcal{O}(f_i^{(m)})$. The map $\delta_\omega f$ is linear in
$f$ and in $\omega$. Hence $c(\cdot,\omega)$ is real and multilinear, and it is linear in
$\omega$. Permuting the $f_i$ permutes the summands of $c$, so by the symmetry of $S^T_n$ the
function $c$ is symmetric.

\emph{Equivariance.} Let $w\in W_4$; it is orthogonal, so $w\omega w^{-1}$ is antisymmetric and
$\nabla(w\cdot f)(x)=w\nabla f(w^{-1}x)$. Therefore
\begin{equation*}
\delta_{w\omega w^{-1}}(w\cdot f)(x)
=-(w\omega w^{-1}x)\cdot w\nabla f(w^{-1}x)
=-(\omega w^{-1}x)\cdot\nabla f(w^{-1}x)=(w\cdot\delta_\omega f)(x).
\end{equation*}
Consequently
\begin{equation*}
c(w\cdot\vec f,w\omega w^{-1})
=\sum_i S^T_n(w\cdot f_1,\dots,w\cdot\delta_\omega f_i,\dots,w\cdot f_n)
=c(\vec f,\omega),
\end{equation*}
where the last step uses the exact $W_4$-invariance of the limit given by the infinite-volume
theorem. This is \eqref{10.7:eq-anomaly-flux-2}.

\emph{Translation invariance.} Put $v:=\omega a$. Writing $\omega x=\omega(x-a)+v$ gives
\begin{equation*}
\delta_\omega(\tau_a f)(x)
=-(\omega(x-a))\cdot\nabla f(x-a)-v\cdot\nabla f(x-a),
\end{equation*}
that is, $\delta_\omega(\tau_a f)=\tau_a(\delta_\omega f)-\tau_a(v\cdot\nabla f)$. For each $i$,
$v\cdot\nabla f_i$ is a real $C^1_c$ function with support in $\operatorname{supp}f_i$, so each
tuple with $v\cdot\nabla f_i$ in the $i$-th slot and $f_k$ in the others is admissible. Insert
this in each summand of $c(\tau_a\vec f,\omega)$ and use multilinearity together with the exact
translation invariance of the limit given by the infinite-volume theorem. This yields
\begin{equation}\label{10.7:eq-anomaly-flux-5}
c(\tau_a\vec f,\omega)
=c(\vec f,\omega)-\sum_{i=1}^n S^T_n(f_1,\dots,v\cdot\nabla f_i,\dots,f_n).
\end{equation}
It remains to prove the translation Ward identity: the last sum vanishes. The function
$s\mapsto S^T_n(\tau_{sv}\vec f)$ is constant, by translation invariance. We show that it is
differentiable at $s=0$ with derivative $-\sum_i S^T_n(f_1,\dots,v\cdot\nabla f_i,\dots,f_n)$;
the derivative of a constant function is zero, so the sum vanishes.

For $f\in C^2_c$ and $s\ne0$ set
\begin{equation*}
g_s:=\frac{\tau_{sv}f-f}{s}+v\cdot\nabla f .
\end{equation*}
Then
\begin{equation*}
g_s(x)=\int_0^1 v\cdot\bigl[\nabla f(x)-\nabla f(x-\theta sv)\bigr]\,d\theta,
\end{equation*}
and $\nabla g_s$ is the same integral with $\nabla(v\cdot\nabla f)$ in place of $v\cdot\nabla f$.
Hence
\begin{equation}\label{10.7:eq-anomaly-flux-6}
\|g_s\|_{\sharp}\le\max\bigl\{|v|\,m_{\nabla f}(|s||v|),\ 40M|v|\,m_{\nabla^2 f}(|s||v|)\bigr\}
\longrightarrow0\quad(s\to0).
\end{equation}
The triangle inequality then gives
\begin{equation*}
\|\tau_{sv}f-f\|_{\sharp}\le|s|\bigl(\|v\cdot\nabla f\|_{\sharp}+\|g_s\|_{\sharp}\bigr)
\longrightarrow0 .
\end{equation*}

Now telescope as in Lemma~\ref{10.7:lem-limit-differentiability}. By multilinearity,
\begin{equation*}
S^T_n(\tau_{sv}\vec f)-S^T_n(\vec f)
=\sum_{i=1}^n S^T_n\bigl(\tau_{sv}f_1,\dots,\tau_{sv}f_{i-1},\,\tau_{sv}f_i-f_i,\,
f_{i+1},\dots,f_n\bigr).
\end{equation*}
Write $g^{(i)}_s$ for $g_s$ built from $f_i$. Dividing by $s$ and adding
$\sum_iS^T_n(f_1,\dots,v\cdot\nabla f_i,\dots,f_n)$, we find that the sum equals
\begin{align*}
&\sum_i S^T_n\bigl(\tau_{sv}f_1,\dots,\tau_{sv}f_{i-1},\,g^{(i)}_s,\,f_{i+1},\dots,f_n\bigr)\\
&\quad-\sum_i\sum_{j<i}S^T_n\bigl(\tau_{sv}f_1,\dots,\tau_{sv}f_{j-1},\,\tau_{sv}f_j-f_j,\,
f_{j+1},\dots,f_{i-1},\\
&\qquad\qquad v\cdot\nabla f_i,\,f_{i+1},\dots,f_n\bigr).
\end{align*}
Here the second line is the telescoping of the difference between the tuples with $f_j$ and
with $\tau_{sv}f_j$ in the slots $j<i$.

Take $|s||v|\le\mathsf{d}_1/4$. Every tuple occurring above then has its entries supported in
the sets $(\operatorname{supp}f_k)^{+\mathsf{d}_1/4}$. These sets lie in
$B(0,\mathcal{R}_1/2+\mathsf{d}_1/4)$ and are pairwise at distance at least $\mathsf{d}_1/2$.
The multilinear bound for the limit in \eqref{10.7:eq-equicontinuity-a}, supplied by the
per-ball volume-free bound, with the constant
$C^{\infty}_n(\mathsf{d}_1/2;\,\mathcal{R}_1/2+\mathsf{d}_1/4)$, therefore shows that the
displayed sum is at most
\begin{equation*}
\begin{split}
&C^{\infty}_n\bigl(\tfrac{\mathsf{d}_1}{2};\tfrac{\mathcal{R}_1}{2}+\tfrac{\mathsf{d}_1}{4}\bigr)
\Bigl[\sum_i\|g^{(i)}_s\|_{\sharp}\prod_{k\ne i}\|f_k\|_{\sharp}\Bigr.\\
&\qquad\Bigl.+\sum_i\sum_{j<i}\|\tau_{sv}f_j-f_j\|_{\sharp}\,\|v\cdot\nabla f_i\|_{\sharp}
\prod_{k\ne i,j}\|f_k\|_{\sharp}\Bigr].
\end{split}
\end{equation*}
This uses $\|\tau_{sv}f_k\|_{\sharp}=\|f_k\|_{\sharp}$. By \eqref{10.7:eq-anomaly-flux-6} and
the bound following it, the bracket tends to zero as $s\to0$. This proves the claimed
derivative, hence the vanishing of the sum in \eqref{10.7:eq-anomaly-flux-5}, and
$c(\tau_a\vec f,\omega)=c(\vec f,\omega)$.

\emph{Part (d), necessity.} Assume $S^T_n$ is $O(4)$-invariant. Let $\vec f$ be as in Part (c),
so that $\vec f\in\mathcal{T}^{(2)}_n(\mathcal{R}_1,\mathsf{d}_1)$. If $\omega=0$ then $c=0$ by
linearity in $\omega$. If $\omega\ne0$, each $R^\omega_t\cdot\vec f$ is an admissible $C^1_c$
tuple, so $S^T_n(R^\omega_t\cdot\vec f)-S^T_n(\vec f)=0$ for all $t$. By
Lemma~\ref{10.7:lem-limit-differentiability} the difference quotient converges to
$c(\vec f,\omega)$ as $t\to0$. Hence $c(\vec f,\omega)=0$.

\emph{Part (d), sufficiency.} Assume $c\equiv0$ on all such tuples and all $\omega$. Let
$\vec f\in\mathcal{T}^{(2)}_n(\mathcal{R}_1,\mathsf{d}_1)$ and $\omega\ne0$, and put
$\varphi(s):=S^T_n(R^\omega_s\cdot\vec f)$. Rotations about the origin preserve
$B(0,\mathcal{R}_1/2)$, distances between supports, and the class $C^2_c$. Hence
$R^\omega_s\cdot\vec f\in\mathcal{T}^{(2)}_n(\mathcal{R}_1,\mathsf{d}_1)$ for every $s$.
Moreover
$\varphi(s+u)=S^T_n(R^\omega_u\cdot(R^\omega_s\cdot\vec f))$.

Lemma~\ref{10.7:lem-limit-differentiability} at the base point $R^\omega_s\cdot\vec f$, with the
same $t_0$ for all $s$, shows that $\varphi$ is differentiable with
$\varphi'(s)=c(R^\omega_s\cdot\vec f,\omega)$. This vanishes by hypothesis; in the argument of
Proposition~\ref{10.7:prop-orbit-constancy-pointwise} the hypothesis $c\equiv0$ here replaces
the appeal to Proposition~\ref{10.7:prop-anomaly-vanishes}. By the mean value theorem $\varphi$
is constant on $\mathbb{R}$.

By Lemma~\ref{10.7:lem-exponential-onto}, every $R\in SO(4)$ equals $e^{\omega}=R^\omega_1$ for
some antisymmetric $\omega$ (if $\omega=0$ then $R=\mathbb{1}$ and there is nothing to prove).
Hence
\begin{equation*}
S^T_n(R\cdot\vec f)=\varphi(1)=\varphi(0)=S^T_n(\vec f)
\end{equation*}
for all $R\in SO(4)$ and all $\vec f$ in the union of the classes
$\mathcal{T}^{(2)}_n(\mathcal{R}_1,\mathsf{d}_1)$ over $\mathcal{R}_1,\mathsf{d}_1>0$. By
Proposition~\ref{10.7:prop-density-extension}
the same holds for every admissible $C^1_c$ tuple.

Finally let $R\in O(4)$ with $\det R=-1$, and put $w_0=\operatorname{diag}(-1,1,1,1)\in W_4$.
Then $w_0R\in SO(4)$ and $R\cdot\vec f=w_0\cdot((w_0R)\cdot\vec f)$. The exact
$W_4$-invariance of the infinite-volume theorem, followed by the $SO(4)$-invariance just
proved, gives $S^T_n(R\cdot\vec f)=S^T_n((w_0R)\cdot\vec f)=S^T_n(\vec f)$.

\emph{Part (e).} Let $\mathcal{R}'\ge2\rho_0$. The supports of the $h_i$ lie in
$B(0,\rho_0)\subset B(0,\mathcal{R}'/2)$, and their mutual separation is unchanged. Hence
$\vec h\in\mathcal{T}^{(2)}_n(\mathcal{R}',\mathsf{d})$, and
\eqref{10.7:eq-anomaly-flux-a} applied with the data $(\omega,n,\mathcal{R}',\mathsf{d})$, for
which the hypotheses of Proposition~\ref{10.7:prop-boundary-independence} are assumed in Part (e),
gives
$\lim_j\mathcal{B}^{\chi_{\mathcal{R}'}}_j(\vec h)=c(\vec h,\omega)$. The right-hand side does
not depend on $\mathcal{R}'$, since it is defined through $S^T_n$ on fixed tuples. Therefore
$\lim_j|\mathcal{B}^{\chi_{\mathcal{R}'}}_j(\vec h)|=|c(\vec h,\omega)|$ for every
$\mathcal{R}'\ge2\rho_0$, and the outer limit in \eqref{10.7:eq-anomaly-flux-3} exists and
equals $|c(\vec h,\omega)|$.

This proves (e1)$\Leftrightarrow$(e4). By \eqref{10.7:eq-anomaly-flux-a} at radius
$\mathcal{R}$, each of (e2) and (e3) is equivalent to (e4). Thus, given the two ultraviolet
hypotheses of Proposition~\ref{10.7:prop-boundary-independence}, the limit
$\mathcal{R}'\to\infty$ in \eqref{10.7:eq-anomaly-flux-3} is superfluous.

Under the additional assumptions, the forms in \eqref{10.7:eq-annulus-functionals-a}, the
uniform boundedness of the tensor-class correlations displayed there and the per-ball
volume-free bound give, through parts (i) and (ii) of
\eqref{10.7:eq-divergence-vanishes-b}, that the $J^4$ term tends to zero at every fixed
$\mathcal{R}'\ge\max\{12,2\rho_0\}$. This is the first limit in \eqref{10.7:eq-anomaly-flux-4}.
Since $\mathfrak{F}=B^4+J^4$ and the cumulant is linear in its first entry,
\begin{equation*}
\kappa_j(B^4_{\chi_{\mathcal{R}'}};\mathcal{O}(\vec h))
=\mathcal{B}^{\chi_{\mathcal{R}'}}_j-\kappa_j(J^4_{\chi_{\mathcal{R}'}};\mathcal{O}(\vec h)).
\end{equation*}
By \eqref{10.7:eq-anomaly-flux-a} and the first limit, the right-hand side tends to
$c(\vec h,\omega)$. The equivalence argument above then runs unchanged with $B^4$ in place of
$\mathfrak{F}$.
\end{proof}

\begin{remark}
For families of cut-offs $(\chi_{\mathcal{R}'})$ that are not radial, the scalar part of $J^4$
contributes
$\kappa_N S^T_{n+1}(\tfrac12\operatorname{div}\xi^{\chi},\vec h)$, with
$\xi^{\chi}:=\chi\cdot(\omega y)$, as stated in the remark on non-radial cut-offs that
accompanies \eqref{10.7:eq-divergence-vanishes-b}. The vanishing of this term as
$\mathcal{R}'\to\infty$ is itself an infrared decay statement of the same kind. This is the
reason for displaying
the far-sphere flux hypothesis with radial cut-offs.

We note in what sense the far-sphere flux hypothesis is nevertheless weaker in
kind than the invariance to which it is equivalent.
\begin{enumerate}
\item[(i)] Given the two ultraviolet hypotheses of
Proposition~\ref{10.7:prop-boundary-independence}, Part (e) shows that the analogue for
$\mathfrak{F}$ of the far-sphere flux hypothesis
is equivalent, tuple by tuple, to $c(\vec h,\omega)=0$. Taken over all tuples and all $\omega$,
Part (d) shows that it is equivalent to the $O(4)$-invariance of $S^T_n$. It is therefore not
logically weaker.
\item[(ii)] A sufficient condition for it is its stress-tensor decay form: the clustering bound in
\eqref{10.7:eq-annulus-functionals-a} with $\delta>0$, together with the uniform boundedness of
the tensor-class correlations displayed there. This condition is of a different kind. It is a
bound, uniform in $j$, on one
family of two-cluster truncated correlations: a single component of a dimension-four
gauge-invariant tensor density in a unit cube at distance $D$ (the distance variable of the
clustering bound in \eqref{10.7:eq-annulus-functionals-a}), against $n\ge2$ separated scalar
sources. It mentions no rotation, generator, cut-off or limit. It is of the type that a cluster
expansion at the last scale, carried out in infinite volume, would give. It leaves four powers
of perturbative headroom: the perturbative rate $|x|^{-8}$ against the required
$|x|^{-4-\delta}$. It is
not implied back by $O(4)$-invariance.
\item[(iii)] As to the exponent: decay $D^{-4-\delta'}$ in the same distance variable gives,
through part (iv) of
\eqref{10.7:eq-divergence-vanishes-b}, the shell bound $C_B\mathcal{R}'^{-\delta'}$ at cut-off
radius $\mathcal{R}'$. This tends to
zero if and only if $\delta'>0$; for $\delta'=0$ it yields only $|c(\vec h,\omega)|\le C_B$.
\end{enumerate}
\end{remark}

\begin{definition}[Conventions for the covariant quasi-free models]
\label{10.7:not-quasi-free-conventions}
The model systems set up with the following conventions are not gauge theories, and nothing
about $SU(2)$ is inferred from them. The conventions are these.

\begin{enumerate}
\item \emph{Fourier transform.} For a test function $f$ on $\mathbb{R}^4$,
\begin{equation*}
\hat f(p) := \int_{\mathbb{R}^4} e^{-ip\cdot x} f(x)\,dx ,
\end{equation*}
and the inverse transform carries the factor $(2\pi)^{-4}$. We write
\begin{equation*}
G(x) := \frac{1}{4\pi^2 |x|^2}, \qquad \hat G(p) = |p|^{-2}.
\end{equation*}

\item \emph{Representation and symbol.} $(W,D)$ is a real orthogonal representation of
$O(4)$ on a space $W$ of dimension $N_W \ge 1$. The symbol is
\begin{equation*}
K(p) = \sum_{\mu,\nu} K^{\mu\nu} p_\mu p_\nu ,
\end{equation*}
with the following three properties. It is symmetric. It is covariant: for every
$R \in O(4)$,
\begin{equation*}
K(Rp) = D(R)\,K(p)\,D(R)^{T}.
\end{equation*}
It is elliptic: $K(p) \succ 0$ for $p \ne 0$. Its ellipticity constant is denoted
$\kappa_K$, to keep it apart from Ba{\l}aban's parameter $\kappa$:
\begin{equation*}
\kappa_K := \min_{|e| = 1} \lambda_{\min} K(e) > 0 ,
\end{equation*}
where $\lambda_{\min}$ denotes the smallest eigenvalue. $K(\partial)$ is the operator with
multiplier $K(p)$.

\item \emph{Time, momentum and reflection.} We write $x = (t,\vec x)$ and
$\theta(t,\vec x) = (-t,\vec x)$. Dually, $p = (p_0,\vec p)$ and $\omega_p := |\vec p|$.
The geometric reflection of the field $E$ is
\begin{equation*}
\Theta_D\bigl(E(f)\bigr) := E\bigl(D(\theta)\,(\bar f\circ\theta)\bigr),
\end{equation*}
where $\bar f$ is the complex conjugate of $f$.

\item \emph{Truncated functions.} Let $\mathbf{C} = (\mathbf{C}_{aa'})$ be the covariance
of the model, with component indices $a,a'$. Let $P_1,\dots,P_n$ be local Wick polynomials
placed at pairwise separated points $x_1,\dots,x_n$. Each leg of the vertex $P_i$ carries a
component index $a$ and a derivative multi-index $\alpha$. The truncated function
$T_{\mathbf{C}}[P_1,\dots,P_n]$ is the connected graph sum defined as follows.
\begin{itemize}
\item The sum runs over the perfect matchings of the legs.
\item No matching joins two legs of the same vertex, that is, there are no
self-contractions.
\item The graph that a matching induces on the vertices $\{1,\dots,n\}$ is connected.
\end{itemize}
A line joining the leg $(a,\alpha)$ of $P_i$ to the leg $(a',\alpha')$ of $P_j$ carries
the tensor obtained by applying $\partial^{\alpha}$ in the variable $x_i$ and
$\partial^{\alpha'}$ in the variable $x_j$ to $\mathbf{C}_{aa'}(x_i - x_j)$, that is,
\begin{equation*}
\partial^{\alpha}_{x_i}\partial^{\alpha'}_{x_j}\bigl[\mathbf{C}_{aa'}(x_i - x_j)\bigr]
= (-1)^{|\alpha'|}\,\bigl(\partial^{\alpha+\alpha'}\mathbf{C}_{aa'}\bigr)(x_i - x_j).
\end{equation*}
Since $\mathbf{C}_{a'a}(-u) = \mathbf{C}_{aa'}(u)$, this tensor does not depend on which
end of the line is called $i$.
A local Wick polynomial is called a \emph{$D$-scalar} when its vertex tensors are invariant
under $R^{\otimes}\otimes D(R)^{\otimes}$ for every $R \in O(4)$.

\item \emph{Diagonal action.} $O(4)$ acts on fields diagonally, by
\begin{equation*}
E \longmapsto D(R)\,\bigl(E\circ R^{-1}\bigr).
\end{equation*}
\end{enumerate}
\end{definition}

\begin{theorem}[Quasi-free fields positive for the geometric reflection are rotation invariant]
\label{10.7:thm-quasi-free-reflection}
Let the representation $(W,D)$ of $O(4)$, the symbol $K(p)$ with ellipticity constant $\kappa$, the
operator $K(\partial)$, the field $E(f)$, the geometric reflection $\Theta_D$ and the anomaly $c$ be as in
the conventions for the covariant quasi-free models, Definition~\ref{10.7:not-quasi-free-conventions}. Let
$W_{\mathbb{C}}$ be the complexification of $W$, and let $G$ denote the Green's function of $-\Delta$ on
$\mathbb{R}^4$.

Let $\mu$ be a Gaussian probability measure on $\mathcal{S}'(\mathbb{R}^4;W)$, write $\mathbb{E}_{\mu}$
for expectation with respect to $\mu$, and assume that $\mu$ has the following properties.
\begin{itemize}
\item $\mu$ is translation invariant.
\item $\mu$ solves the Schwinger--Dyson equations
\begin{equation}\label{10.7:eq-quasi-free-reflection-1}
\mathbb{E}_{\mu}\bigl[E(K(\partial)g)\,F\bigr]=\mathbb{E}_{\mu}\bigl[\partial_g F\bigr],
\qquad F\in\{1,\,E(f)\}.
\end{equation}
\item For every $f\in C^\infty_c(\mathbb{R}^4;W_{\mathbb{C}})$ supported in $\{t>0\}$,
\begin{equation}\label{10.7:eq-quasi-free-reflection-2}
\mathbb{E}_{\mu}\bigl[\Theta_D(E(f))\,E(f)\bigr]\ge 0 .
\end{equation}
\end{itemize}
Then $D\equiv 1$ and $K(p)=p^2\Pi^{-1}$ with $\Pi\succ 0$. Moreover $\mu=\mathcal{N}(m,\Pi G+H_0')$,
with $m\in W$ constant and $H_0'\succeq 0$ constant. Finally, $\mu$ is $O(4)$-invariant, every
$D$-scalar local Wick probe is an orbital scalar, and $c\equiv 0$.
\end{theorem}

\begin{proof}
We write $p=(p_0,\vec p)$, $p^2=p_0^2+|\vec p|^2$, $\omega_p=|\vec p|$, and $\hat f$ for the Fourier
transform $\mathcal{F}f$. The letter $\theta$ is the time reflection $(t,\vec x)\mapsto(-t,\vec x)$,
which acts on momenta by $p_0\mapsto -p_0$. Since $\theta^2=1$ and $D(\theta)$ is orthogonal, we have
$D(\theta)^{\mathsf T}=D(\theta)^{-1}=D(\theta)$.

\emph{Structure of $\mu$.} The mean $f\mapsto\mathbb{E}_{\mu}[E(f)]$ is a translation-invariant
distribution, hence a constant $m\in W$. By the Schwartz kernel theorem and the Bochner--Schwartz
theorem, the covariance $C$ of $\mu$ has the form
\begin{equation}\label{10.7:eq-quasi-free-reflection-3}
C(\bar f,g)=(2\pi)^{-4}\int \hat f(p)^{*}\,\hat C(dp)\,\hat g(p),
\end{equation}
where $\hat C$ is a positive semidefinite Hermitian matrix-valued tempered measure. We now take
$F=E(f)$ in \eqref{10.7:eq-quasi-free-reflection-1} and pass to Fourier variables. This gives the
identity of measures $K(p)\,\hat C(dp)=\mathbb{1}\,dp$.

For $p\ne 0$ the matrix $K(p)$ is invertible, $K^{-1}$ is continuous there, and
$|K^{-1}(p)|\le\kappa^{-1}|p|^{-2}$. The last bound is locally integrable in dimension four, so
$K^{-1}(p)\,dp$ is a tempered measure. On $\mathbb{R}^4\setminus\{0\}$ we therefore have
$\hat C=K^{-1}(p)\,dp$. It follows that $\hat C-K^{-1}(p)\,dp$ is carried by $\{0\}$. Since
$K^{-1}(p)\,dp$ gives no mass to $\{0\}$, this difference equals $H_0\delta_0$ with
$H_0=\hat C(\{0\})\succeq 0$.

Positivity of $\hat C$ has a further consequence. Since $\hat C$ is Hermitian and positive
semidefinite, for every Borel set $B\subset\mathbb{R}^4\setminus\{0\}$ the matrix
$\int_B K^{-1}(p)\,dp=\hat C(B)$ is Hermitian and positive semidefinite. Taking for $B$ small balls and
using continuity, we find that $K^{-1}(p)$ is Hermitian and positive semidefinite for every $p\ne 0$.
Being invertible, it is positive definite. In particular $\det K(p)>0$ for $p\neq 0$.

\emph{The determinant.} We show that $\det K(p)=c_K\,(p^2)^{N_W}$ with a constant $c_K>0$. By the
covariance of the symbol, $K(Rp)=D(R)K(p)D(R)^{\mathsf T}$, and $\det D(R)^2=1$. Hence $\det K$ is an
$O(4)$-invariant polynomial.

Such a polynomial is a polynomial in $p^2$. To see this, restrict it to a line $\mathbb{R}e$ through the
origin, with $|e|=1$. Since some $R\in O(4)$ maps $e$ to $-e$, the restriction is an even polynomial in
the line parameter, that is, a polynomial in its square. Because $O(4)$ is transitive on spheres, this
determines $\det K(p)$ as that polynomial evaluated at $p^2$.

Since $K$ is homogeneous of degree $2$ and of size $N_W\times N_W$, the polynomial $\det K$ is
homogeneous of degree $2N_W$. Therefore $\det K(p)=c_K\,(p^2)^{N_W}$, and $c_K>0$ because
$\det K(p)>0$ for $p\ne0$, as shown at the end of the structure step.

We now draw the consequences for fixed $\vec p\ne 0$. The function
$K^{-1}(p_0,\vec p)=\operatorname{adj}K/\det K$ is rational in $p_0$, with denominator
$c_K\,(p_0^2+\omega_p^2)^{N_W}$. Its poles therefore lie only at
$p_0=\pm i\omega_p$. Let $r$ be the order of the pole at $-i\omega_p$; then $r\le N_W$. We also have
$r\ge 1$. Otherwise $K^{-1}(\cdot,\vec p)$ would have no pole at $-i\omega_p$. By the covariance
relation at $R=\theta$ established below, it would then have none at $+i\omega_p$ either. It would thus be
entire and tend to zero at infinity, hence vanish by Liouville's theorem, which is impossible. Since
$\operatorname{adj}K$ has degree at most $2N_W-2$ in $p_0$, we get
$|K^{-1}(p_0,\vec p)|\le C_1|p_0|^{-2}$ for large complex $|p_0|$, with $C_1$ depending on $\vec p$.

The covariance at $R=\theta$ reads $K(-p_0,\vec p)=D(\theta)K(p_0,\vec p)D(\theta)^{\mathsf T}$ for real
$p_0$. By rationality it also holds for complex $p_0$. It gives
\begin{equation}\label{10.7:eq-quasi-free-reflection-4}
K^{-1}(-p_0,\vec p)=D(\theta)\,K^{-1}(p_0,\vec p)\,D(\theta)^{\mathsf T},
\end{equation}
so the poles at $+i\omega_p$ and $-i\omega_p$ have the same order.

\emph{Fibre form of the reflection positivity.} Let $0<t_0<t_1$, let
$a\in C^\infty_c((t_0,t_1);W_{\mathbb{C}})$, and let $g$ be a scalar function with
$\hat g\in C^\infty_c(\mathbb{R}^3\setminus\{0\})$. Put $f(t,\vec x)=a(t)g(\vec x)$. Then
$\hat f(p)=A(p_0)\hat g(\vec p)$, where
\begin{equation*}
A(p_0):=\int e^{-ip_0t}a(t)\,dt, \qquad \bar A(p_0):=\int e^{-ip_0t}\,\overline{a(t)}\,dt .
\end{equation*}
Since $\hat g(0)=0$, we have $\int f=0$. Hence the term $H_0\delta_0$ and the mean $m$ do not contribute
to the left side of \eqref{10.7:eq-quasi-free-reflection-2}, which becomes
\begin{equation}\label{10.7:eq-quasi-free-reflection-5}
\mathbb{E}_{\mu}\bigl[\Theta_D(E(f))E(f)\bigr]
=(2\pi)^{-4}\int d^3\vec p\,|\hat g(\vec p)|^2\,q_{\vec p}(a),
\end{equation}
with
\begin{equation}\label{10.7:eq-quasi-free-reflection-6}
\begin{split}
q_{\vec p}(a)
&:=\int_{\mathbb{R}}\bar A(p_0)^{\mathsf T}D(\theta)K^{-1}(p_0,\vec p)A(p_0)\,dp_0\\
&=-2\pi i\,\operatorname*{Res}_{p_0=-i\omega_p}
\Bigl[\bar A(p_0)^{\mathsf T}D(\theta)K^{-1}(p_0,\vec p)A(p_0)\Bigr].
\end{split}
\end{equation}

The second equality comes from closing the contour in the lower half-plane. For $\operatorname{Im}p_0\le 0$
and $t\ge t_0>0$ we have $|e^{-ip_0t}|\le 1$. Two integrations by parts therefore give
$|A(p_0)|\le\|a''\|_1|p_0|^{-2}$ there, and the same bound holds for $\bar A$. Together with
$|K^{-1}|\le C_1|p_0|^{-2}$, the integrand is $O(|p_0|^{-6})$ on large lower semicircles. Its only pole in
the lower half-plane is at $-i\omega_p$, and the clockwise orientation produces the factor $-2\pi i$.

The function $\vec p\mapsto q_{\vec p}(a)$ is continuous on $\vec p\ne 0$. Indeed, near a given
$\vec p_0\ne0$ the residue is an integral over one fixed small circle around $-i\omega_p$, on which the
integrand depends continuously on $\vec p$.

The function $f$ is a Schwartz function supported in $\{t>0\}$ but not compactly supported in $\vec x$.
Replace $g(\vec x)$ by $g(\vec x)$ times a cut-off equal to one on a large ball. This produces functions in
$C^\infty_c(\mathbb{R}^4;W_{\mathbb{C}})$ supported in $\{t>0\}$ that converge to $f$ in $\mathcal{S}$.
The quadratic form in \eqref{10.7:eq-quasi-free-reflection-2} is continuous on $\mathcal{S}$, because
$\hat C$ is tempered. So \eqref{10.7:eq-quasi-free-reflection-2} extends to $f$.

Now choose a sequence $\hat g_k$ with $|\hat g_k|^2$ of unit integral concentrating at a point
$\vec p_0\ne 0$. By \eqref{10.7:eq-quasi-free-reflection-5} and continuity, $q_{\vec p_0}(a)\ge 0$. Hence
$q_{\vec p}(a)\ge 0$ for every $\vec p\ne0$ and every $a$.

\emph{The pole is simple.} Fix $\vec p\ne 0$ and put $\zeta:=p_0+i\omega_p$. Expand
\begin{equation*}
K^{-1}(p_0,\vec p)=\sum_{j\ge -r}Q_j\zeta^j,\qquad Q_{-r}\ne 0 .
\end{equation*}
Since $e^{-ip_0t}=e^{-\omega_pt}e^{-i\zeta t}$, we have
\begin{equation*}
A(p_0)=\sum_{k\ge0}\frac{(-i\zeta)^k}{k!}\alpha_k,\qquad
\bar A(p_0)^{\mathsf T}=\sum_{k\ge0}\frac{(-i\zeta)^k}{k!}\alpha_k^{*},\qquad
\alpha_k:=\int t^ke^{-\omega_pt}a(t)\,dt .
\end{equation*}
The residue in \eqref{10.7:eq-quasi-free-reflection-6} is the coefficient of $\zeta^{-1}$. Only the
indices with $k+l+j=-1$ and $j\ge -r$ contribute, that is, those with $k+l\le r-1$. Hence
\begin{equation}\label{10.7:eq-quasi-free-reflection-7}
q_{\vec p}(a)=\sum_{k+l\le r-1}\alpha_k^{*}M_{kl}\alpha_l,\qquad
M_{kl}=-2\pi i\,\frac{(-i)^{k+l}}{k!\,l!}\,D(\theta)Q_{-1-k-l}.
\end{equation}

The map $a\mapsto(\alpha_0,\dots,\alpha_{r-1})$ is onto $W_{\mathbb{C}}^r$. To see this, apply it
componentwise. If its image on scalar functions were a proper subspace of $\mathbb{C}^r$, then some
nonzero combination $\sum_kc_kt^ke^{-\omega_pt}$ would integrate to zero against every test function on
$(t_0,t_1)$. It would then vanish on $(t_0,t_1)$, contradicting the linear independence of the functions
$t^ke^{-\omega_pt}$ there.

Suppose $r\ge2$. Take $\alpha_0=u$ and $\alpha_{r-1}=sv$ with $u,v\in W_{\mathbb{C}}$ and $s$ real, and
all other $\alpha_k=0$. The term with $k=l=r-1$ does not occur, since $2r-2>r-1$. By
\eqref{10.7:eq-quasi-free-reflection-7},
\begin{equation*}
q_{\vec p}(a)=u^{*}M_{00}u+s\bigl(u^{*}Yv+v^{*}Yu\bigr),
\end{equation*}
where
\begin{equation*}
Y:=M_{0,r-1}=M_{r-1,0}=-2\pi i\,\frac{(-i)^{r-1}}{(r-1)!}\,D(\theta)Q_{-r}\ne 0 .
\end{equation*}
Since $q_{\vec p}(a)\ge0$ for all real $s$, the coefficient of $s$ vanishes: $u^{*}Yv+v^{*}Yu=0$ for all
$u,v$. Replacing $v$ by $iv$ gives $u^{*}Yv-v^{*}Yu=0$. Hence $u^{*}Yv=0$ for all $u,v$, so $Y=0$, a
contradiction. Therefore $r=1$. By \eqref{10.7:eq-quasi-free-reflection-4} the pole at $+i\omega_p$ is
also simple.

Consequently $(p_0^2+\omega_p^2)K^{-1}(p_0,\vec p)$ is rational without poles, i.e. a polynomial matrix in
$p_0$. By $|K^{-1}|\le C_1|p_0|^{-2}$ it is bounded, hence independent of $p_0$.

Put $\Phi(p):=p^2K^{-1}(p)$ for $p\ne0$. This function has four properties:
\begin{itemize}
\item it is continuous;
\item it is homogeneous of degree $0$;
\item it is covariant, $\Phi(Rp)=D(R)\Phi(p)D(R)^{\mathsf T}$;
\item it is constant on every line parallel to the time axis that misses the origin.
\end{itemize}
By covariance, for every $R\in O(4)$ the function $\Phi$ is constant on every line parallel to the image
under $R$ of the time direction that misses the origin: along such a line,
\begin{equation*}
\Phi(p+sRe_0)=D(R)\,\Phi(R^{-1}p+se_0)\,D(R)^{\mathsf T},
\end{equation*}
where $e_0$ is the unit time vector. Since $O(4)$ is transitive on directions, $\Phi$ is constant on every
line missing the origin. Any two points of $\mathbb{R}^4\setminus\{0\}$ are joined by a chain of at most
two segments lying on such lines. Hence $\Phi$ is constant:
\begin{equation*}
\Phi\equiv\Pi=K(e)^{-1}\succ 0,
\end{equation*}
for any unit vector $e$ (the value does not depend on $e$ since $\Phi$ is constant), and positive
definiteness was shown above. Covariance gives
$D(R)\Pi D(R)^{\mathsf T}=\Pi$ for all $R$, and $K(p)=p^2\Pi^{-1}$.

With $K^{-1}=\Pi/\bigl((p_0-i\omega_p)(p_0+i\omega_p)\bigr)$, the residue at $-i\omega_p$ is
$\Pi/(-2i\omega_p)$. Also $A(-i\omega_p)=\alpha_0$ and $\bar A(-i\omega_p)^{\mathsf T}=\alpha_0^{*}$.
Hence \eqref{10.7:eq-quasi-free-reflection-6} becomes
\begin{equation}\label{10.7:eq-quasi-free-reflection-8}
q_{\vec p}(a)=\frac{\pi}{\omega_p}\,\alpha_0^{*}D(\theta)\Pi\,\alpha_0 .
\end{equation}

\emph{$D\equiv1$.} Covariance at $R=\theta$ gives $D(\theta)\Pi D(\theta)^{\mathsf T}=\Pi$, that is,
$D(\theta)\Pi=\Pi D(\theta)$. The matrix $P:=D(\theta)\Pi$ is therefore symmetric, because
$P^{\mathsf T}=\Pi D(\theta)=P$. By \eqref{10.7:eq-quasi-free-reflection-8}, the non-negativity of
$q_{\vec p}$ and the surjectivity of $a\mapsto\alpha_0$, we have $P\succeq 0$.

Let $W_-$ be the $(-1)$-eigenspace of the involution $D(\theta)$; it is $\Pi$-invariant since the two
matrices commute. For $0\ne v\in W_-$,
\begin{equation*}
v^{\mathsf T}Pv=(D(\theta)v)^{\mathsf T}\Pi v=-v^{\mathsf T}\Pi v<0,
\end{equation*}
which contradicts $P\succeq0$. Hence $W_-=0$ and $D(\theta)=1$.

Then $D(R\theta R^{-1})=D(R)D(\theta)D(R)^{-1}=1$ for all $R\in O(4)$. Every hyperplane reflection has
the form $R\theta R^{-1}$, because $O(4)$ is transitive on unit normals. So $D$ is trivial on every
hyperplane reflection.

By the Cartan--Dieudonn\'e theorem, every element of $O(d)$ is a product of hyperplane reflections. We
recall the induction on $d$. In $O(1)=\{\pm1\}$ the element $-1$ is a reflection. Let $R\in O(d)$. If
$Re_1\ne e_1$, the reflection $\sigma$ in the hyperplane orthogonal to $Re_1-e_1$ maps $Re_1$ to $e_1$.
Then $\sigma R$, or $R$ itself if $Re_1=e_1$, fixes $e_1$. It is therefore an element of $O(e_1^\perp)$,
a product of reflections by the inductive hypothesis; these extend to reflections of $\mathbb{R}^d$.
Hence $D\equiv 1$.

\emph{Conclusion.} With $D\equiv1$ the diagonal and the orbital actions of $O(4)$ coincide. By the
structure paragraph and $K^{-1}=\Pi/p^2$, the measure $\mu$ is $\mathcal{N}(m,\Pi G+H_0')$. Here $m$ is
constant, $G$ is radial, and $H_0'\succeq0$ is the constant kernel contributed by the atom
$H_0\delta_0$ in \eqref{10.7:eq-quasi-free-reflection-3}. A Gaussian measure is determined by its mean and
covariance, and both are $O(4)$-invariant, so $\mu$ is $O(4)$-invariant.

The truncated functions of $D$-scalar local Wick probes are contractions of the vertex tensors of the
probes with line tensors $\partial^\alpha G$. The vertex tensors are invariant, and the line tensors are
covariant. Hence the truncated functions are invariant, and every $D$-scalar local Wick probe is an
orbital scalar. Finally, let $S(\vec f)$ be the Schwinger function of the tuple $\vec f$ and
$s\mapsto e^{s\omega}$ the one-parameter rotation group, as in
Definition~\ref{10.7:not-quasi-free-conventions}. Since $s\mapsto S(e^{s\omega}\cdot\vec f)$ is constant,
\begin{equation*}
c=\frac{d}{ds}\Big|_{s=0}S(e^{s\omega}\cdot\vec f)=0 .
\end{equation*}
\end{proof}

\begin{corollary}[Mixtures of quasi-free solutions and the standard reflection]
\label{10.7:cor-quasi-free-mixtures}
Let $(W,D)$, $K$, $\Theta_D$ and the field $E(f)$ be as in
Theorem~\ref{10.7:thm-quasi-free-reflection}. Let $\theta$ be the time reflection
$(t,\vec x)\mapsto(-t,\vec x)$, and let $G$ be the kernel on $\mathbb{R}^4$ with Fourier
transform $p^{-2}$.
\begin{enumerate}
\item[(a)] Let $\lambda$ be a probability measure on a measurable space and, for each point $u$
of it, let $m_u\in W$ be constant and $H_u\succeq 0$ a constant covariance, both measurable in
$u$, such that every
\begin{equation*}
\mu_u:=\mathcal{N}(m_u,\mathcal{F}^{-1}K^{-1}+H_u)
\end{equation*}
is a translation-invariant quasi-free solution of the Schwinger--Dyson equations of
Theorem~\ref{10.7:thm-quasi-free-reflection}, for one and the same $K$. Put
$\mu:=\int\mu_u\,d\lambda(u)$. Suppose that $E_\mu[\Theta_D(E(f))E(f)]\ge 0$ for every
$f\in C^\infty_c(\mathbb{R}^4;W_{\mathbb{C}})$ supported in $\{t>0\}$ with $\int f=0$. This
holds in particular if the inequality holds for all $f$ supported in $\{t>0\}$. Then
$D\equiv 1$, and $K(p)=p^2\Pi^{-1}$ with $\Pi\succ 0$. Moreover
$\mu_u=\mathcal{N}(m_u,\Pi G+H_u)$ for every $u$; every $\mu_u$ satisfies the remaining
conclusions of Theorem~\ref{10.7:thm-quasi-free-reflection}, the statements on $D$-scalar local
Wick probes and on $c\equiv 0$ being asserted here componentwise, for each $\mu_u$; and $\mu$
is $O(4)$-invariant. Only the single reflection $\theta$ enters; no invariance of $\mu$ under
$W_4$ is assumed.
\item[(b)] Let $W=\bigoplus_i W_i$ be a decomposition into $D$-irreducible blocks, and let
$c_i>0$. Suppose that $D(\theta)$ is not the identity on some block $W_i$. Then the Gaussian
measure
$\nu':=\mathcal{N}(0,\bigoplus_i c_i\mathbb{1}_{W_i}\otimes G)$ is not $\Theta_D$-positive:
there is a real $f$ supported in $\{t>0\}$ with $E_{\nu'}[\Theta_D(E(f))E(f)]<0$. Here, for
real $f$, we take $\Theta_D(E(f))=E(D(\theta)(f\circ\theta))$.
\item[(c)] Let $\nu:=\mathcal{N}(0,\mathbb{1}_W\otimes G)$ and, in this part only, let
$\Theta_0$ denote $\Theta_D$ for the trivial representation $D\equiv\mathbb{1}$, that is,
$\Theta_0(E(f))=E(f\circ\theta)$ for real $f$. Let $f_1,\dots,f_n\in C^\infty_c(\mathbb{R}^4;W)$ be
real and supported in $\{t>0\}$. For complex $c_1,\dots,c_n$ and
\begin{equation*}
F=\sum_k c_k e^{E(f_k)}
\end{equation*}
put $\Theta_0(F)=\sum_k\bar c_k e^{E(f_k\circ\theta)}$. For a polynomial
$F=P(E(f_1),\dots,E(f_n))$ with complex coefficients put
$\Theta_0(F)=\bar P(E(f_1\circ\theta),\dots,E(f_n\circ\theta))$, where $\bar P$ is the
polynomial whose coefficients are the complex conjugates of those of $P$. Then, for $F$ of
either form,
\begin{equation}\label{10.7:eq-quasi-free-mixtures-a}
E_\nu[\Theta_0(F)F]\ge 0 .
\end{equation}
\end{enumerate}
\end{corollary}

\begin{proof}
\emph{(a).} The proof of Theorem~\ref{10.7:thm-quasi-free-reflection} uses the positivity
hypothesis only for test functions with $\int f=0$. Let $f$ be such a function. The reflected
test function $f^\theta$ is obtained from $f$ by composing with $\theta$, applying $D(\theta)$
and, for complex $f$, conjugating. Hence $\Theta_D(E(f))=E(f^\theta)$, and $\int f^\theta=0$
because $\theta$ preserves Lebesgue measure.

Under $\mu_u$ the mean of $E(h)$ is the pairing of $m_u$ with $\int h$, since $m_u$ is
constant. The contribution of the constant covariance $H_u$ to $E_{\mu_u}[E(h)E(f)]$ is the
$H_u$-pairing of $\int h$ with $\int f$. Both vanish for $h=f^\theta$. Therefore
\begin{equation*}
E_{\mu_u}[\Theta_D(E(f))E(f)]=C_K(f^\theta,f),
\end{equation*}
where $C_K$ is the covariance form of $\mathcal{F}^{-1}K^{-1}$. The right-hand side does not
depend on $u$. Integrating in $\lambda$ gives $E_\mu[\Theta_D(E(f))E(f)]=C_K(f^\theta,f)$.
Thus the inequality assumed for $\mu$ holds for every $\mu_u$ at every such $f$: the positivity
used does not see $(m_u,H_u)$.

Each $\mu_u$ is Gaussian, translation invariant and a solution of the Schwinger--Dyson
equations. The proof of Theorem~\ref{10.7:thm-quasi-free-reflection} therefore applies to
each $\mu_u$ and gives $D\equiv 1$, $K(p)=p^2\Pi^{-1}$ with $\Pi\succ 0$, and the remaining
conclusions for $\mu_u$. Since $\mathcal{F}^{-1}(\Pi p^{-2})=\Pi G$, we get
$\mu_u=\mathcal{N}(m_u,\Pi G+H_u)$. Each $\mu_u$ is $O(4)$-invariant, hence so is the mixture
$\mu$.

\emph{(c).} Use the Fourier transform $\hat f(p)=\int e^{-ip\cdot x}f(x)\,d^4x$. The covariance
of $\nu$ is
\begin{equation*}
C(f,g)=(2\pi)^{-4}\int\langle\hat f(p),\hat g(p)\rangle\,p^{-2}\,d^4p ,
\end{equation*}
where $\langle\cdot,\cdot\rangle$ is the Hermitian product on $W_{\mathbb{C}}$, antilinear in
its first argument. Put
\begin{equation*}
a_f(t,\vec p)=\int e^{-i\vec p\cdot\vec x}f(t,\vec x)\,d^3x .
\end{equation*}
Then $\widehat{f\circ\theta}(p_0,\vec p)=\int e^{ip_0t}a_f(t,\vec p)\,dt$.

For $\omega>0$ and $T>0$ consider the integral of $e^{ip_0T}(p_0^2+\omega^2)^{-1}$ over
$p_0\in\mathbb{R}$. Closing the contour in the upper half-plane, where the only pole is
$p_0=i\omega$, with residue $e^{-\omega T}/(2i\omega)$, gives
\begin{equation}\label{10.7:eq-quasi-free-mixtures-1}
\int_{\mathbb{R}}\frac{e^{ip_0T}}{p_0^2+\omega^2}\,dp_0=\frac{\pi e^{-\omega T}}{\omega}.
\end{equation}
Since $(p_0^2+\omega^2)^{-1}$ is even in $p_0$, the substitution $p_0\mapsto-p_0$ shows that
\eqref{10.7:eq-quasi-free-mixtures-1} holds with $e^{-ip_0T}$ in place of $e^{ip_0T}$.

Let $f,g$ be supported in $\{t>0\}$. For $t$ in the time support of $a_f$ and $t'$ in the time
support of $a_g$ we have $T=t+t'>0$, so \eqref{10.7:eq-quasi-free-mixtures-1} applies. Moreover
$a_f,a_g$ have compact support in time and decay rapidly in $\vec p$, and $1/\omega_p$ is
locally integrable in $d^3\vec p$. Hence Fubini's theorem and
\eqref{10.7:eq-quasi-free-mixtures-1} give
\begin{equation}\label{10.7:eq-quasi-free-mixtures-2}
\begin{split}
b(f,g)&:=C(f\circ\theta,g)\\
&=(2\pi)^{-4}\int d^3\vec p\,\frac{\pi}{\omega_p}
\Bigl\langle\int e^{-\omega_pt}a_f\,dt,\int e^{-\omega_pt}a_g\,dt\Bigr\rangle .
\end{split}
\end{equation}
For complex numbers $\eta_k$, formula \eqref{10.7:eq-quasi-free-mixtures-2} gives
\begin{equation*}
\sum_{k,l}\bar\eta_k\eta_l\,b(f_k,f_l)
=(2\pi)^{-4}\int\frac{\pi}{\omega_p}\Bigl|\sum_l\eta_l\int e^{-\omega_pt}a_{f_l}\,dt\Bigr|^2
d^3\vec p\ \ge 0 .
\end{equation*}
So $(b(f_k,f_l))_{k,l}$ is positive semidefinite.

Now take $h,h'$ real. Then $\hat h(-p)=\overline{\hat h(p)}$, so $C(h,h')$ is real and
symmetric. Since $p^{-2}$ is invariant under $p_0\mapsto-p_0$, we have
$C(h\circ\theta,h\circ\theta)=C(h,h)$. For the Gaussian $\nu$,
$E_\nu[e^{E(h)}]=e^{\frac12C(h,h)}$. Consequently
\begin{equation*}
E_\nu[\Theta_0(F)F]=\sum_{k,l}\bar c_kc_l\,e^{\frac12C(f_k,f_k)}e^{\frac12C(f_l,f_l)}
e^{b(f_k,f_l)} .
\end{equation*}
Write $B=(b(f_k,f_l))$. By the Schur product theorem every Hadamard power $B^{\circ n}$ is
positive semidefinite. The entrywise exponential $(e^{b(f_k,f_l)})_{k,l}=\sum_n B^{\circ n}/n!$
is therefore positive semidefinite. With $\eta_k=c_ke^{\frac12C(f_k,f_k)}$ the right-hand side
above equals
\begin{equation*}
\sum_{k,l}\bar\eta_k\eta_l\,e^{b(f_k,f_l)}\ge 0,
\end{equation*}
which is \eqref{10.7:eq-quasi-free-mixtures-a} on the exponential algebra.

For polynomials, the monomial $\prod_jE(f_j)^{\alpha_j}$ is $\partial_s^\alpha$ of
$e^{E(\sum_js_jf_j)}$ at $s=0$. It is a limit of difference quotients in the real parameters
$s$, pointwise for every configuration. Each difference quotient is a real combination of
exponentials $e^{E(h)}$ with $h$ real
and supported in $\{t>0\}$. Hence $F=P(E(f_1),\dots,E(f_n))$ is the limit of combinations
\begin{equation*}
F_\varepsilon=\sum_k c_{k,\varepsilon}e^{E(h_{k,\varepsilon})}
\end{equation*}
with complex $c_{k,\varepsilon}$ and real $h_{k,\varepsilon}$ supported in $\{t>0\}$, and
$\Theta_0(F_\varepsilon)$, obtained by conjugating the coefficients and replacing each
$h_{k,\varepsilon}$ by $h_{k,\varepsilon}\circ\theta$, converges in the same way to
$\Theta_0(F)$. By the exponential case, $E_\nu[\Theta_0(F_\varepsilon)F_\varepsilon]$ is a
quadratic form in a positive semidefinite matrix, and hence is $\ge 0$. For step sizes
$\varepsilon\le 1$ the mean value theorem bounds $|F_\varepsilon|$ and
$|\Theta_0(F_\varepsilon)|$ by a constant times
\begin{equation*}
\exp\Bigl(c\sum_j\bigl(|E(f_j)|+|E(f_j\circ\theta)|\bigr)\Bigr),
\end{equation*}
the constants $c$ and the prefactor depending only on $P$. The Gaussian $\nu$ has
exponential moments of all orders, so the square of this bound is $\nu$-integrable, and by
dominated convergence these quantities converge to $E_\nu[\Theta_0(F)F]$, which is therefore
$\ge 0$.

\emph{(b).} We have $\theta^2=1$, so $D(\theta)^2=\mathbb{1}$. As $D(\theta)\ne\mathbb{1}$ on
$W_i$, there is a real $v\in W_i$, $v\neq 0$, with $D(\theta)v=-v$.

Let $\psi\in C^\infty_c(\mathbb{R}^4)$ be real, $\psi\ge 0$, not identically zero, supported
in $\{t>0\}$, and
put $f=\psi v$. Then $\Theta_D(E(f))=-E((\psi\circ\theta)v)$. The covariance on $W_i$ is
$c_i\mathbb{1}_{W_i}\otimes G$, so \eqref{10.7:eq-quasi-free-mixtures-2}, applied to the
scalar $\psi$, gives
\begin{equation*}
q:=E_{\nu'}[\Theta_D(E(f))E(f)]=-c_i|v|^2\,b(\psi,\psi).
\end{equation*}
The function $\vec p\mapsto\int e^{-\omega_pt}a_\psi\,dt$ is continuous and equals
$\int\psi>0$ at $\vec p=0$. Hence $b(\psi,\psi)>0$, and $q<0$.
\end{proof}

\begin{theorem}[Quasi-free fields on periodic boxes converge to the invariant limit]
\label{10.7:thm-quasi-free-boxes}
We use the conventions for the covariant quasi-free models fixed in
Notation~\ref{10.7:not-quasi-free-conventions}. Throughout, $G$ is the Green's function of $-\Delta$ on
$\mathbb{R}^4$, so that $-\Delta G=\delta_0$ and $G$ is homogeneous of degree $-2$.

For $\ell>0$ let $\nu_\ell$ be the centred Gaussian measure on $\mathbb{R}^4/\ell\mathbb{Z}^4$ with
covariance
\begin{equation}\label{10.7:eq-quasi-free-boxes-1}
  C_\ell(x)=\ell^{-4}\sum_{p\in(2\pi/\ell)\mathbb{Z}^4\setminus\{0\}}|p|^{-2}e^{ip\cdot x},
\end{equation}
the massless scalar field with the zero mode removed. Let $n\ge 2$, let $P_1,\dots,P_n$ be local
Wick polynomials, Wick-ordered with respect to $\nu_\ell$, and let
$f_1,\dots,f_n\in C^\infty_c(\mathbb{R}^4)$ have pairwise disjoint supports. For a covariance
$C$, $S_C(\vec f)$ denotes the truncated expectation of $P_1(f_1),\dots,P_n(f_n)$ under the centred
Gaussian measure of covariance $C$, and $S_G(\vec f)$ the same expression with $C$ replaced by $G$.
Then
$S_{C_\ell}(\vec f)\to S_G(\vec f)$ as $\ell\to\infty$, along the whole family and not only along
subsequences. If the $P_i$ are $O(4)$-scalars, the limit is $O(4)$-invariant and the rotational
anomaly $c$ vanishes identically.

The same holds
\begin{enumerate}
\item[(a)] for the massive covariance with symbol $(|p|^2+M^2)^{-1}$, $M>0$, keeping all modes,
  with $\ell\to\infty$ first and then $M\downarrow 0$;
\item[(b)] for tori $\mathbb{R}^4/\ell\Lambda$ of any shape, where $\Lambda\subset\mathbb{R}^4$ is a
  lattice.
\end{enumerate}
\end{theorem}

\begin{proof}
Let $\mathfrak{g}$ be the function on the unit torus $\mathbb{R}^4/\mathbb{Z}^4$ defined by
\begin{equation*}
  \mathfrak{g}(y):=\sum_{q\in2\pi\mathbb{Z}^4\setminus\{0\}}|q|^{-2}e^{iq\cdot y}.
\end{equation*}
Then $-\Delta\mathfrak{g}=\sum_{q\ne0}e^{iq\cdot y}=\delta_0-1$ on the unit torus. The substitution
$p=q/\ell$ in \eqref{10.7:eq-quasi-free-boxes-1} gives $|p|^{-2}=\ell^2|q|^{-2}$, and hence
$C_\ell(x)=\ell^{-2}\mathfrak{g}(x/\ell)$.

Fix $\chi_0\in C^\infty_c(B(0,1/4))$ with $\chi_0=1$ on $B(0,1/8)$. Since $B(0,1/4)$ lies inside a
fundamental domain of $\mathbb{Z}^4$, $\chi_0G$ is a function on the unit torus. Put
$h:=\mathfrak{g}-\chi_0G$. By the Leibniz rule,
\begin{equation*}
  -\Delta(\chi_0G)=\chi_0(-\Delta G)+\eta=\delta_0+\eta,
  \qquad \eta:=[-\Delta,\chi_0]G=-(\Delta\chi_0)G-2\nabla\chi_0\cdot\nabla G .
\end{equation*}
Therefore
\begin{equation*}
  -\Delta h=(\delta_0-1)-\delta_0-\eta=-1-\eta .
\end{equation*}
The function $\eta$ is supported in the annulus $1/8\le|y|\le1/4$, where $\chi_0$ is not locally
constant and $G$ is smooth. Hence $\eta$ is a smooth function on the torus, and its Fourier
coefficients $\hat\eta(q)$ decrease rapidly. Comparing Fourier coefficients at $q\ne0$ gives
$|q|^2\hat h(q)=-\hat\eta(q)$, that is, $\hat h(q)=-|q|^{-2}\hat\eta(q)$. These coefficients decrease
rapidly, so $h\in C^\infty$ of the torus.

Let $B\subset\mathbb{R}^4\setminus\{0\}$ be compact and let $\ell>8\sup_{x\in B}|x|$. On the open set
$\{x\ne0:\ |x|<\ell/8\}$, which contains $B$, we have $|x/\ell|<1/8$, hence $\chi_0(x/\ell)=1$ and
$\mathfrak{g}(x/\ell)=G(x/\ell)+h(x/\ell)$. By the homogeneity of $G$,
$\ell^{-2}G(x/\ell)=G(x)$, so that on this set
\begin{equation}\label{10.7:eq-quasi-free-boxes-2}
  C_\ell(x)=G(x)+\ell^{-2}h(x/\ell).
\end{equation}
Differentiating \eqref{10.7:eq-quasi-free-boxes-2} gives, for every multi-index $\alpha$,
\begin{equation*}
  \partial^\alpha C_\ell(x)-\partial^\alpha G(x)=\ell^{-2-|\alpha|}(\partial^\alpha h)(x/\ell).
\end{equation*}
The right-hand side is bounded by $\ell^{-2-|\alpha|}\sup|\partial^\alpha h|$ on $B$, so it tends to $0$
uniformly on $B$ as $\ell\to\infty$.

The truncated functions are finite sums of finite products of such derivatives. Since the $P_i$ are
Wick-ordered with respect to the Gaussian measure, the truncated expectation is a sum over connected
graphs without self-lines. Each line joins two distinct vertices $i\ne j$ and carries a derivative
$\partial^\alpha C_\ell$ evaluated at $x_i-x_j$. The result is integrated against
$f_1(x_1)\cdots f_n(x_n)$. The differences $x_i-x_j$ range over the compact set
\begin{equation*}
  \bigcup_{i\ne j}\bigl(\operatorname{supp}f_i-\operatorname{supp}f_j\bigr),
\end{equation*}
which does not contain $0$ because the supports are pairwise disjoint. Taking this set as $B$, every
line converges uniformly, and hence $S_{C_\ell}(\vec f)\to S_G(\vec f)$. Since the limit exists for
the full family $\ell\to\infty$, no subsequence is needed.

Now let the $P_i$ be $O(4)$-scalars. In $S_G$ every line carries the tensor $\partial^\alpha G$, which
is covariant under $O(4)$ because $G$ is rotation invariant. Every vertex tensor is invariant. Hence
$S_G$ is $O(4)$-invariant, and therefore the rotational anomaly $c$ vanishes.

\emph{Massive variant.} Let $C_{\ell,M}$ be the covariance with Fourier coefficients
$\ell^{-4}(|p|^2+M^2)^{-1}$, $p\in(2\pi/\ell)\mathbb{Z}^4$, and let $G_M$ be given by
\begin{equation*}
  G_M:=\mathcal{F}^{-1}(|p|^2+M^2)^{-1}.
\end{equation*}
Then $G_M\in L^1$, and $G_M$ is smooth off $0$ with all derivatives rapidly decreasing. By the
Poisson summation formula,
\begin{equation*}
  C_{\ell,M}=\sum_{n\in\mathbb{Z}^4}G_M(\cdot+\ell n),
\end{equation*}
since both sides have the same Fourier coefficients. On a compact $B\subset\mathbb{R}^4\setminus\{0\}$,
and for $\ell\ge2\sup_B|x|$, we have $|x+\ell n|\ge\ell|n|/2$ for $n\ne0$. Hence the rapid decrease
of the derivatives of $G_M$ gives $\partial^\alpha(C_{\ell,M}-G_M)\to0$ uniformly on $B$ as
$\ell\to\infty$.

Next, $\partial^\alpha G_M\to\partial^\alpha G$ on compacts off $0$ as $M\downarrow0$. To see this,
split $p$-space into a neighbourhood of $0$ and a neighbourhood of $\infty$. Near $0$ we use dominated
convergence with the dominant $|p|^{|\alpha|-2}$, which is integrable near $0$ in four dimensions and
bounds $|p^\alpha|(|p|^2+M^2)^{-1}$. Near $\infty$ we use symbol bounds for $(|p|^2+M^2)^{-1}$ that
are uniform in $M$.

The graph argument above, applied first with $C_{\ell,M}$ and $G_M$ and then with $G_M$ and $G$, gives
\begin{equation*}
  \lim_{M\downarrow0}\lim_{\ell\to\infty}S_{C_{\ell,M}}(\vec f)=S_G(\vec f).
\end{equation*}

\emph{Other shapes.} Replace $\mathfrak{g}$ by the zero-mean Green's function of
$\mathbb{R}^4/\Lambda$, and run the argument above with it.
\end{proof}

\begin{remark}
For a general elliptic covariant operator $K$ in the sense of the conventions, with symbol $K(p)$,
the same argument applies with $\Gamma:=\mathcal{F}^{-1}K^{-1}$ in place of $G$ and the zero-mean
Green's function of $K(\partial)$ on the unit torus in place of $\mathfrak{g}$. This $\Gamma$ is a
fundamental solution of $K(\partial)$, homogeneous of degree $-2$, smooth off $0$ and covariant, and
$\det K(p)$ is a constant multiple of $|p|^{2N_W}$. This case is not used in what follows.
\end{remark}

\begin{remark}
The massive scalar box is reflection positive for the torus reflection $t\mapsto-t$. Write
$p=(p_0,\vec p)$ and put
\begin{equation*}
  \omega_M:=(|\vec p|^2+M^2)^{1/2}.
\end{equation*}
For a fixed spatial momentum, the time dependence of the covariance is given by the one-dimensional
periodic kernel, which on $(0,\ell)$ reads
\begin{equation}\label{10.7:eq-quasi-free-boxes-a}
  \begin{split}
  \ell^{-1}\sum_{p_0\in(2\pi/\ell)\mathbb{Z}}\frac{e^{ip_0s}}{p_0^2+\omega_M^2}
  &=\frac{\cosh\bigl(\omega_M(\ell/2-s)\bigr)}{2\omega_M\sinh(\omega_M\ell/2)}\\
  &=\frac{e^{\omega_M\ell/2}e^{-\omega_Ms}+e^{-\omega_M\ell/2}e^{\omega_Ms}}
        {4\omega_M\sinh(\omega_M\ell/2)} .
  \end{split}
\end{equation}
This is of the form $Ae^{-\omega_Ms}+Be^{\omega_Ms}$ with $A,B>0$.

Let $a$ be a function of time supported in $(0,\ell/2)$. Pairing $a$ with its reflection involves the
kernel at $s=t+t'$ with $t,t'\in(0,\ell/2)$, so that $s\in(0,\ell)$. Each of the two terms, regarded
as a kernel in $(t,t')$ through $s=t+t'$, is positive of rank one, namely
\begin{equation*}
  e^{-\omega_M(t+t')}=e^{-\omega_Mt}e^{-\omega_Mt'},
  \qquad
  e^{\omega_M(t+t')}=e^{\omega_Mt}e^{\omega_Mt'} .
\end{equation*}
Positivity on the exponential algebra then follows from the Schur product theorem, as in
\eqref{10.7:eq-quasi-free-mixtures-a}.

The Fourier coefficients of $C_{\ell,M}$ depend on $|p|$ only, and $(2\pi/\ell)\mathbb{Z}^4$ is
$W_4$-invariant. Hence the massive boxes supply reflection-positive, $W_4$-invariant box
approximants. By part (b) of Theorem~\ref{10.7:thm-quasi-free-boxes}, the shape of the torus plays no
role in the limit. Consequently, the quasi-free models furnish no counterexample to the hypothesis
that the infinite-volume limit does not depend on the shape of the periodic boxes.
\end{remark}

\begin{theorem}[The anomaly vanishes exactly for uniform mixtures]
\label{10.7:thm-uniform-mixtures}
Let $\mathfrak{P}^+$ be the real commutative unital algebra freely generated by symbols $P(f)$,
$f\in C^\infty_c(\mathbb{R}^4)$ real, and $T(h)$, $h\in C^\infty_c(\mathbb{R}^4;\operatorname{Sym}^2\mathbb{R}^4)$
real, subject only to linearity of $P(f)$ in $f$ and of $T(h)$ in $h$. The support of a monomial is
the union of the supports of its factors. Let $\mathfrak{P}^+_{\neq}$ be the linear span of $1$ and of
the monomials whose factors have pairwise disjoint supports, and let
$\mathfrak{P}_{\neq}\subset\mathfrak{P}^+_{\neq}$ be the span of $1$ and of those monomials that contain
no factor $T(h)$. The group $\mathbb{R}^4\rtimes O(4)$ acts on $\mathfrak{P}^+$ by the automorphisms
\begin{equation*}
\alpha_{(a,R)}P(f):=P\bigl(f(R^{-1}(\cdot-a))\bigr),\qquad
\alpha_{(a,R)}T(h):=T\bigl(R\,h(R^{-1}(\cdot-a))\,R^{\mathsf{T}}\bigr),
\end{equation*}
which map $\mathfrak{P}^+_{\neq}$ and $\mathfrak{P}_{\neq}$ onto themselves; we write $\alpha_a$, $\alpha_R$
for $\alpha_{(a,\mathbb{1})}$, $\alpha_{(0,R)}$, put $R\cdot f:=f(R^{-1}\,\cdot)$, act on tuples
$\vec f=(f_1,\dots,f_n)$ componentwise, and let $\theta=w_0=\operatorname{diag}(-1,1,1,1)$ be the time
reflection, the first coordinate $x_1$ being time. A \emph{state} is a linear functional $\nu$ on
$\mathfrak{P}^+_{\neq}$ with $\nu(1)=1$. A state $\nu$ is $\alpha_\theta$-reflection positive if
$\sum_{k,l}\bar c_k c_l\,\nu(\alpha_\theta(F_k)F_l)\ge 0$ for all finite families of complex numbers
$c_k$ and of elements $F_k\in\mathfrak{P}^+_{\neq}$ supported in $\{x\in\mathbb{R}^4: x_1>0\}$.
Joint cumulants $\langle X_1;\dots;X_n\rangle^\nu$ of factors $X_i$ of a monomial in
$\mathfrak{P}^+_{\neq}$ are defined by the moment--cumulant formula
\begin{equation}\label{10.7:eq-uniform-mixtures-1}
\nu\Bigl(\prod_{i\in B}X_i\Bigr)=\sum_{\pi\in\Pi(B)}\ \prod_{b\in\pi}\langle X_i; i\in b\rangle^{\nu},
\end{equation}
$\Pi(B)$ being the set of partitions of the finite set $B$ (the empty product and the partition of
the empty set contributing $1$); $S^\nu_n(f_1,\dots,f_n):=\langle P(f_1);\dots;P(f_n)\rangle^\nu$. For a
real antisymmetric $4\times4$ matrix $\omega$ let $\delta_\omega f:=-(\omega x)\cdot\nabla f$, the
derivative at $t=0$ of $e^{t\omega}\cdot f$, and
\begin{equation*}
c^\nu(\vec f;\omega):=\frac{d}{dt}\Big|_{t=0}S^\nu_n\bigl(e^{t\omega}\cdot\vec f\bigr).
\end{equation*}
A tuple is $\mathsf{d}$-separated if the supports of its entries are pairwise at distance at least
$\mathsf{d}>0$. Assume:
\begin{itemize}
\item[(F1)] $(\mu_u)_{u\in\mathcal{O}}$ is a family of translation-invariant states indexed by a
compact homogeneous $O(4)$-space $\mathcal{O}$, weak-$*$ continuous in $u$ (that is,
$u\mapsto\mu_u(F)$ is continuous for every $F\in\mathfrak{P}^+_{\neq}$), and covariant:
$\mu_u\circ\alpha_R=\mu_{R^{-1}u}$ for $R\in O(4)$, $u\in\mathcal{O}$.
\item[(F2)] (optional; the conclusions of (iii) that refer to it are asserted when it holds) each
$\mu_u$ satisfies a common family of identities linear in the state, of the form $\nu(\Xi)=\beta$
with fixed $\Xi\in\mathfrak{P}^+_{\neq}$ and $\beta\in\mathbb{R}$, among them the flux form
of the rotational Ward identity: for every $\omega$, every $\mathcal{R}>0$, every $\chi$ admissible at
radius $\mathcal{R}$ in the sense fixed with the rotational Ward identity, every $k\ge1$ and all real
$f_1,\dots,f_k$ with pairwise disjoint supports in $B(0,\mathcal{R}/2)$,
\begin{equation}\label{10.7:eq-uniform-mixtures-a}
\sum_{i=1}^k\nu\Bigl(P(\delta_\omega f_i)\prod_{j\neq i}P(f_j)\Bigr)
=\nu\Bigl(T(h^\omega_\chi)\prod_{j=1}^kP(f_j)\Bigr),\qquad \nu\bigl(T(h^\omega_\chi)\bigr)=0,
\end{equation}
where
\begin{equation*}
(h^\omega_\chi)_{\sigma\tau}:=\tfrac12\sum_\rho\bigl(\omega_{\rho\tau}x_\rho\partial_\sigma\chi
+\omega_{\rho\sigma}x_\rho\partial_\tau\chi\bigr).
\end{equation*}
\item[(F3)] there are $W_4$-invariant, $\alpha_\theta$-reflection positive states $\nu_j$, each
defined on the elements supported in the $j$-th cube of an exhausting sequence of cubes, with
$\nu_j(F)\to\bar\mu(F)$ for every $F\in\mathfrak{P}^+_{\neq}$, where
$\bar\mu:=\int_{\mathcal{O}}\mu_u\,d\lambda(u)$ and $\lambda$ is a Borel probability on
$\mathcal{O}$.
\item[(Reg)] for all $\mathcal{R}_0>0$, $\mathsf{d}>0$, $n\ge1$ there is $C_n(\mathcal{R}_0,\mathsf{d})$,
independent of the position of the supports, such that
\begin{equation*}
|\bar\mu(P(f_1)\cdots P(f_n))|\le C_n(\mathcal{R}_0,\mathsf{d})\prod_i\|f_i\|_{C^1}
\end{equation*}
for all $\mathsf{d}$-separated tuples supported in $\bar B(0,\mathcal{R}_0)$.
\end{itemize}
Let $m_{\mathcal{O}}$ be the $O(4)$-invariant Borel probability on $\mathcal{O}$. Then:
\begin{enumerate}
\item[(i)] $\bar\mu$ is translation- and $W_4$-invariant, and $\lambda$ may be replaced by its
$W_4$-average $\tilde\lambda:=|W_4|^{-1}\sum_{w\in W_4}w_*\lambda$: $\int\mu_u\,d\tilde\lambda=\bar\mu$.
\item[(ii)] $\bar\mu$ is $\alpha_\theta$-reflection positive, and its centred Schwinger functions are
OS-positive on separated test functions in the open half-space $\{x_1>0\}$.
\item[(iii)] $\bar\mu$ satisfies every identity of (F2), in particular
\eqref{10.7:eq-uniform-mixtures-a}, and for every antisymmetric $\omega$, every $\mathcal{R}>0$, every
$\chi$ admissible at radius $\mathcal{R}$, and real $f_1,\dots,f_n$ with pairwise disjoint supports in
$B(0,\mathcal{R}/2)$,
\begin{equation}\label{10.7:eq-uniform-mixtures-2}
c^{\bar\mu}(\vec f;\omega)=\bigl\langle T(h^\omega_\chi);P(f_1);\dots;P(f_n)\bigr\rangle^{\bar\mu},
\end{equation}
the right side being independent of $\chi$.
\item[(iv)] $\bar\mu\circ\alpha_R=\int\mu_u\,d\bigl((R^{-1})_*\lambda\bigr)(u)$ for $R\in O(4)$; hence
$\lambda=m_{\mathcal{O}}$ implies that $\bar\mu$ is $O(4)$-invariant, which implies
$c^{\bar\mu}\equiv0$.
\item[(v)] Assume (Sep): the linear span of the functions $u\mapsto\mu_u(F)$, $F\in\mathfrak{P}_{\neq}$,
is dense in $C(\mathcal{O})$. Then $c^{\bar\mu}\equiv0$ if and only if $\lambda=m_{\mathcal{O}}$. If
$\mathcal{O}$ is infinite and $\lambda$ is carried by a finite $W_4$-orbit, then
$c^{\bar\mu}\not\equiv0$.
\end{enumerate}
\end{theorem}

\begin{proof}
By weak-$*$ continuity each $u\mapsto\mu_u(F)$ is continuous on the compact space $\mathcal{O}$,
hence bounded and Borel, so $\bar\mu(F)=\int\mu_u(F)\,d\lambda(u)$ is defined; $\bar\mu$ is linear and
$\bar\mu(1)=1$. The map $R\mapsto\alpha_R$ is a homomorphism, since
$\alpha_R\alpha_SP(f)=P(R\cdot(S\cdot f))=P(RS\cdot f)$, and likewise on the $T(h)$.

(i) Let $w\in W_4$ and $F\in\mathfrak{P}^+_{\neq}$. For $j$ large, $F$ and $\alpha_wF$ are both
supported in the $j$-th cube, and by $W_4$-invariance of $\nu_j$ and (F3),
\begin{equation*}
\bar\mu(\alpha_wF)=\lim_j\nu_j(\alpha_wF)=\lim_j\nu_j(F)=\bar\mu(F).
\end{equation*}
For $a\in\mathbb{R}^4$,
\begin{equation*}
\bar\mu(\alpha_aF)=\int\mu_u(\alpha_aF)\,d\lambda=\int\mu_u(F)\,d\lambda=\bar\mu(F)
\end{equation*}
by the translation invariance of each $\mu_u$. By covariance, $\mu_{wu}=\mu_u\circ\alpha_{w^{-1}}$, so
\begin{equation*}
\int\mu_u(F)\,d(w_*\lambda)(u)=\int\mu_{wu}(F)\,d\lambda(u)=\bar\mu(\alpha_{w^{-1}}F)=\bar\mu(F),
\end{equation*}
and averaging over $w\in W_4$ gives
\begin{equation*}
\int\mu_u\,d\tilde\lambda=|W_4|^{-1}\sum_w\bar\mu\circ\alpha_{w^{-1}}=\bar\mu.
\end{equation*}

(ii) For fixed $c_k$ and $F_k$ supported in $\{x_1>0\}$, the sums
$\sum_{k,l}\bar c_kc_l\,\nu_j(\alpha_\theta(F_k)F_l)$ are non-negative for $j$ large and converge to the
same sum with $\bar\mu$; so $\bar\mu$ is $\alpha_\theta$-reflection positive. Since $\theta\in W_4$,
(i) gives $\bar\mu\circ\alpha_\theta=\bar\mu$. For $F\in\mathfrak{P}_{\neq}$ let $F^c$ be obtained by
replacing each factor $P(f)$ by $P(f)-\bar\mu(P(f))\,1$ and expanding; $F^c\in\mathfrak{P}^+_{\neq}$ has
the support of $F$, and the centred Schwinger functions are the values $\bar\mu(F^c)$. As
$\alpha_\theta(P(f)-\bar\mu(P(f))1)=P(\theta\cdot f)-\bar\mu(P(\theta\cdot f))1$, centring commutes
with $\alpha_\theta$: $\alpha_\theta(F^c)=(\alpha_\theta F)^c$, and
$\alpha_\theta(F_k^c)F_l^c=(\alpha_\theta(F_k)F_l)^c$. Reflection positivity applied to the family
$F_k^c$ gives $\sum_{k,l}\bar c_kc_l\,\bar\mu((\alpha_\theta(F_k)F_l)^c)\ge0$ for $F_k$ built on separated
test functions in $\{x_1>0\}$.

(iii) An identity $\nu(\Xi)=\beta$ holding for every $\mu_u$ holds for $\bar\mu$:
$\bar\mu(\Xi)=\int\mu_u(\Xi)\,d\lambda=\beta$. This applies to \eqref{10.7:eq-uniform-mixtures-a}.
We bring \eqref{10.7:eq-uniform-mixtures-a} to connected form. Let $\nu$ satisfy it, fix $f_1,\dots,f_n$
as there, and write $\kappa(b):=\langle P(f_i);i\in b\rangle^\nu$, $\kappa^{(i)}(b)$ for the same with
$f_i$ replaced by $\delta_\omega f_i$ ($i\in b$; $\operatorname{supp}\delta_\omega f_i\subset\operatorname{supp}f_i$),
and $\tau(A):=\langle T(h^\omega_\chi);P(f_i);i\in A\rangle^\nu$. Let $\emptyset\ne B\subset\{1,\dots,n\}$.
Expanding $\nu(P(\delta_\omega f_i)\prod_{j\in B\setminus\{i\}}P(f_j))$ by
\eqref{10.7:eq-uniform-mixtures-1}, each index lies in exactly one block; grouping by the block $A$
containing $i$ and summing over $i\in B$, the left side of \eqref{10.7:eq-uniform-mixtures-a} for the
subfamily $B$ is
\begin{equation*}
\sum_{\emptyset\neq A\subset B}\ \Bigl(\sum_{i\in A}\kappa^{(i)}(A)\Bigr)
\sum_{\pi\in\Pi(B\setminus A)}\prod_{b\in\pi}\kappa(b).
\end{equation*}
Expanding the right side by \eqref{10.7:eq-uniform-mixtures-1}, grouped by the set $A\subset B$ of
indices sharing the block of $T(h^\omega_\chi)$, gives the same expression with $\tau(A)$ in place of
$\sum_{i\in A}\kappa^{(i)}(A)$, plus the term $A=\emptyset$, which carries
$\tau(\emptyset)=\nu(T(h^\omega_\chi))=0$. With $D(A):=\sum_{i\in A}\kappa^{(i)}(A)-\tau(A)$,
\eqref{10.7:eq-uniform-mixtures-a} for every subfamily therefore yields, for every $B\neq\emptyset$,
\begin{equation*}
\sum_{\emptyset\ne A\subset B}D(A)\sum_{\pi\in\Pi(B\setminus A)}\prod_{b\in\pi}\kappa(b)=0.
\end{equation*}
The term $A=B$ is $D(B)$; by induction on $|B|$ all other terms vanish, so
$D(B)=0$. For $B=\{1,\dots,n\}$:
\begin{equation}\label{10.7:eq-uniform-mixtures-3}
\sum_{i=1}^nS^\nu_n(f_1,\dots,\delta_\omega f_i,\dots,f_n)
=\bigl\langle T(h^\omega_\chi);P(f_1);\dots;P(f_n)\bigr\rangle^\nu .
\end{equation}

It remains to show that $t\mapsto S^{\bar\mu}_n(e^{t\omega}\cdot\vec f)$ is $C^1$ with derivative at
$t=0$ equal to the left side of \eqref{10.7:eq-uniform-mixtures-3}; this is the argument of the
differentiability lemma (Lemma~\ref{10.7:lem-limit-differentiability}), with (Reg) in place of the
bound used there. Inverting \eqref{10.7:eq-uniform-mixtures-1},
\begin{equation*}
\langle X_1;\dots;X_n\rangle^{\nu}=\sum_{\pi\in\Pi(\{1,\dots,n\})}(-1)^{|\pi|-1}(|\pi|-1)!
\prod_{b\in\pi}\nu\Bigl(\prod_{i\in b}X_i\Bigr),
\end{equation*}
so (Reg) gives, with
\begin{equation*}
C'_n:=\sum_{\pi\in\Pi(\{1,\dots,n\})}(|\pi|-1)!\prod_{b\in\pi}C_{|b|}(\mathcal{R}_0,\mathsf{d}),
\end{equation*}
the bound
\begin{equation}\label{10.7:eq-uniform-mixtures-4}
|S^{\bar\mu}_n(g_1,\dots,g_n)|\le C'_n\prod_i\|g_i\|_{C^1}
\end{equation}
for $\mathsf{d}$-separated tuples in $\bar B(0,\mathcal{R}_0)$. Let $\vec f$ be $2\mathsf{d}$-separated in
$\bar B(0,\mathcal{R}_0)$ and $g^t_i:=e^{t\omega}\cdot f_i$. Rotations fix $\bar B(0,\mathcal{R}_0)$ and
$|e^{s\omega}x-x|\le(e^{|s|\|\omega\|}-1)|x|$, so for $|t|$ so small that
$(e^{|t|\|\omega\|}-1)\mathcal{R}_0\le\mathsf{d}/2$, every tuple below (entries $g^t_j$, $f_j$,
$\delta_\omega f_j$ or $g^t_j-f_j$ in place $j$) is $\mathsf{d}$-separated in $\bar B(0,\mathcal{R}_0)$.
Telescoping one entry at a time, by multilinearity and $P(g)-P(f)=P(g-f)$,
\begin{equation*}
S^{\bar\mu}_n(\vec g^{\,t})-S^{\bar\mu}_n(\vec f)
=\sum_{i=1}^nS^{\bar\mu}_n\bigl(g^t_1,\dots,g^t_{i-1},g^t_i-f_i,f_{i+1},\dots,f_n\bigr).
\end{equation*}
Dividing by $t$ and subtracting $\sum_iS^{\bar\mu}_n(f_1,\dots,\delta_\omega f_i,\dots,f_n)$, the $i$-th
difference is
\begin{align*}
&S^{\bar\mu}_n\bigl(g^t_1,\dots,g^t_{i-1},t^{-1}(g^t_i-f_i)-\delta_\omega f_i,f_{i+1},\dots,f_n\bigr)\\
&\quad+S^{\bar\mu}_n\bigl(g^t_1,\dots,g^t_{i-1},\delta_\omega f_i,f_{i+1},\dots,f_n\bigr)\\
&\quad-S^{\bar\mu}_n\bigl(f_1,\dots,f_{i-1},\delta_\omega f_i,f_{i+1},\dots,f_n\bigr),
\end{align*}
the difference of the last two terms telescoped again into terms with one entry $g^t_j-f_j$, $j<i$. By
\eqref{10.7:eq-uniform-mixtures-4} these are bounded by $C'_n$ times products of $C^1$ norms
containing $\|t^{-1}(g^t_i-f_i)-\delta_\omega f_i\|_{C^1}$ or $\|g^t_j-f_j\|_{C^1}$, and both tend to $0$
as $t\to0$, because for $f$ of class $C^2$ with compact support the rotated difference quotients
$t^{-1}(e^{t\omega}\cdot f-f)$ converge to $\delta_\omega f$ in $C^1$ and $e^{t\omega}\cdot f\to f$ in
$C^1$. Hence the derivative at $0$ exists and equals the left side of
\eqref{10.7:eq-uniform-mixtures-3}. Applied to $e^{t\omega}\cdot\vec f$, which is again admissible
since rotations are isometries fixing $0$, this gives the derivative at $t$,
\begin{equation*}
\frac{d}{dt}S^{\bar\mu}_n(\vec g^{\,t})
=\sum_{i=1}^nS^{\bar\mu}_n\bigl(g^t_1,\dots,e^{t\omega}\cdot\delta_\omega f_i,\dots,g^t_n\bigr),
\end{equation*}
since $\delta_\omega(e^{t\omega}\cdot f)=e^{t\omega}\cdot\delta_\omega f$, $\omega$ commuting with
$e^{t\omega}$; it is continuous in $t$ by \eqref{10.7:eq-uniform-mixtures-4}. Thus
$c^{\bar\mu}(\vec f;\omega)$ is the left side of \eqref{10.7:eq-uniform-mixtures-3} for $\nu=\bar\mu$,
which gives \eqref{10.7:eq-uniform-mixtures-2}; this left side does not involve $\chi$.

(iv) By covariance,
\begin{equation*}
\bar\mu(\alpha_RF)=\int\mu_{R^{-1}u}(F)\,d\lambda(u)
=\int\mu_v(F)\,d\bigl((R^{-1})_*\lambda\bigr)(v).
\end{equation*}
If $\lambda=m_{\mathcal{O}}$ then $(R^{-1})_*\lambda=\lambda$ and
$\bar\mu\circ\alpha_R=\bar\mu$ for all $R\in O(4)$. Then for every $b\subset\{1,\dots,n\}$ the joint
moments
\begin{equation*}
\bar\mu\Bigl(\prod_{i\in b}P(e^{t\omega}\cdot f_i)\Bigr)
=\bar\mu\Bigl(\alpha_{e^{t\omega}}\prod_{i\in b}P(f_i)\Bigr)
=\bar\mu\Bigl(\prod_{i\in b}P(f_i)\Bigr)
\end{equation*}
of the rotated tuple equal those of the $P(f_i)$, hence so do the cumulants (polynomials in the
moments), so $S^{\bar\mu}_n(e^{t\omega}\cdot\vec f)$ is constant in $t$ and $c^{\bar\mu}\equiv0$.

(v) The implication $\lambda=m_{\mathcal{O}}\Rightarrow c^{\bar\mu}\equiv0$ is (iv). Conversely let
$c^{\bar\mu}\equiv0$, $n\ge2$, and $\vec f$ a tuple of real test functions with pairwise disjoint
supports; these are at positive mutual distance and lie in some $B(0,\mathcal{R}/2)$. Since
$e^{t\omega}\cdot\vec f$ is again such a tuple,
\begin{equation*}
\frac{d}{dt}S^{\bar\mu}_n(e^{t\omega}\cdot\vec f)
=\frac{d}{ds}\Big|_{s=0}S^{\bar\mu}_n\bigl(e^{s\omega}\cdot(e^{t\omega}\cdot\vec f)\bigr)
=c^{\bar\mu}(e^{t\omega}\cdot\vec f;\omega)=0,
\end{equation*}
so $S^{\bar\mu}_n(e^{\omega}\cdot\vec f)=S^{\bar\mu}_n(\vec f)$ for every antisymmetric $\omega$; by the
surjectivity of the exponential map (Lemma~\ref{10.7:lem-exponential-onto}),
$S^{\bar\mu}_n(R\cdot\vec f)=S^{\bar\mu}_n(\vec f)$ for all $R\in SO(4)$. For $n=1$, (Reg) makes
$f\mapsto\bar\mu(P(f))$ a distribution of order at most one, translation invariant by (i); its first
derivatives vanish, so it equals $\kappa_1\int f\,dx$ for a constant $\kappa_1$ (the density of the
first cumulant), which is $O(4)$-invariant. Moments are polynomials in cumulants by
\eqref{10.7:eq-uniform-mixtures-1}, so
$\bar\mu\circ\alpha_R=\bar\mu$ on $\mathfrak{P}_{\neq}$ for $R\in SO(4)$. For $R\in O(4)\setminus SO(4)$,
$\theta R\in SO(4)$ and $\alpha_R=\alpha_\theta\alpha_{\theta R}$, so by (i) again
$\bar\mu\circ\alpha_R=\bar\mu$ on $\mathfrak{P}_{\neq}$. By (iv),
$\int\mu_u(F)\,d((R^{-1})_*\lambda)=\int\mu_u(F)\,d\lambda$ for all $F\in\mathfrak{P}_{\neq}$; both sides are
integrals against probabilities, hence continuous in the sup norm, so by (Sep) the two measures
agree on $C(\mathcal{O})$, and by uniqueness in the Riesz--Markov theorem $(R^{-1})_*\lambda=\lambda$.
Thus $\lambda$ is $O(4)$-invariant. Let $dR$ be the Haar probability of $O(4)$, $u_0\in\mathcal{O}$ and
$\varphi\in C(\mathcal{O})$. By invariance and Fubini,
\begin{equation*}
\int\varphi\,d\lambda=\int_{O(4)}\!\int_{\mathcal{O}}\varphi(Ru)\,d\lambda(u)\,dR
=\int_{\mathcal{O}}\!\int_{O(4)}\varphi(Ru)\,dR\,d\lambda(u),
\end{equation*}
and by transitivity every $u$ equals $R_uu_0$ for some $R_u\in O(4)$, so by right invariance of Haar
measure the inner integral equals
$\int_{O(4)}\varphi(Ru_0)\,dR$ for every $u$. Hence $\lambda$ is the image of $dR$ under
$R\mapsto Ru_0$, which is $m_{\mathcal{O}}$. Finally, if $\mathcal{O}$ is infinite, transitivity and
invariance give all points equal $m_{\mathcal{O}}$-mass, hence mass zero, so $m_{\mathcal{O}}$ has no
atoms; a $\lambda$ carried by a finite $W_4$-orbit differs from $m_{\mathcal{O}}$, and
$c^{\bar\mu}\not\equiv0$.
\end{proof}

\begin{remark}
Point separation of $\mathcal{O}$ by the functions $u\mapsto\mu_u(F)$ does not suffice in (v). For
instance, let $\mu_u(F):=\mu_0(F)+u\cdot\ell(F)$ with $\mu_0$ invariant and $\ell$ equivariant and
linear, $u\cdot\ell(F)$ denoting the pairing of $u$ with $\ell(F)$, and with $-u\in\mathcal{O}$
whenever $u\in\mathcal{O}$. Then the measure
$\lambda':=\tfrac12(\delta_u+\delta_{-u})$ gives $\bar\mu=\mu_0$, which is invariant, while
$\lambda'\ne m_{\mathcal{O}}$.
\end{remark}

\begin{remark}
The pure components cannot cluster at the stress tensor. Suppose some $\mu_u$ satisfied
\eqref{10.7:eq-uniform-mixtures-a}, (Reg), and, for some $\delta>0$, the bound
\begin{equation*}
\bigl|\langle T(h);P(f_1);\dots;P(f_n)\rangle^{\mu_u}\bigr|\le C\|h\|_{C^1}D^{-4-\delta}
\end{equation*}
when the supports of $h$ and of the $f_i$ are at distance $D$. By the argument of (iii) applied to
$\mu_u$,
\begin{equation*}
c^{\mu_u}(\vec f;\omega)=\bigl\langle T(h^\omega_\chi);P(f_1);\dots;P(f_n)\bigr\rangle^{\mu_u}
\end{equation*}
for every admissible $\chi$, so the right side does not depend on $\mathcal{R}$, and it would be
bounded by $C\mathcal{R}^4\mathcal{R}^{-4-\delta}\to0$; so $c^{\mu_u}\equiv0$
and $\mu_u$ would be $SO(4)$-invariant, which is incompatible with (Sep) unless $|\mathcal{O}|\le2$.
Hence, in any family $(\mu_u)$ satisfying (F1) and (Sep) with $|\mathcal{O}|>2$ in which each $\mu_u$
satisfies \eqref{10.7:eq-uniform-mixtures-a} and the bound (Reg) with $\mu_u$ in place of $\bar\mu$,
the correlations $\langle T;P;\dots;P\rangle^{\mu_u}$ have non-decaying shell integrals: the borderline
$|x|^{-4}$ behaviour. This is the precise sense in which the equivalence of the far-sphere flux
hypothesis $(\mathrm{IR}\text{-}\mathrm{fl})_{\omega_0}$, that the flux of the lattice rotation current
through far spheres vanishes along one sequence of radii, with $c=0$ is an infrared statement; the
decay bound displayed above, by which one far-away stress-tensor insertion decorrelates from the
observables faster than $|x|^{-4}$, is a sufficient condition for it.
\end{remark}

\begin{proposition}[A hypercubic counter-example without the flux identity]
\label{10.7:prop-flux-counterexample}
Let
\begin{equation}\label{10.7:eq-flux-counterexample-1}
v_4(p) := \sum_{\mu=1}^{4} p_\mu^4 - \tfrac12\,|p|^4 , \qquad p\in\mathbb{R}^4 .
\end{equation}
Then $v_4$ is a harmonic quartic, invariant under $W_4$ and not invariant under $O(4)$. Let
$\mathcal{H}_4$ be the real space of harmonic homogeneous quartic polynomials on $\mathbb{R}^4$, on
which $O(4)$ acts by the natural representation $D$, $(D(R)u)(p) := u(R^{-1}p)$; we write
$R\cdot u := D(R)u$, identify $u\in\mathcal{H}_4$ with the symmetric traceless tensor
$u^{\mu\nu\rho\sigma}$ for which $u(p) = u^{\mu\nu\rho\sigma}p_\mu p_\nu p_\rho p_\sigma$, and write
$u(\partial)$ for the corresponding constant-coefficient operator. Let
$\mathcal{O} := O(4)\cdot v_4\subset\mathcal{H}_4$; it is an infinite compact homogeneous space.

Consider a field $E$ with values in $\mathcal{H}_4$ and a real scalar field $\varphi$, with the local,
linear dynamics
\begin{equation}\label{10.7:eq-flux-counterexample-2}
\partial_\mu E = 0 \quad (\mu = 1,\dots,4), \qquad -\Delta\varphi = 0 ,
\end{equation}
understood in the Schwinger--Dyson sense; \eqref{10.7:eq-flux-counterexample-2} is invariant under
the diagonal action of $O(4)$ ($D$ on $E$, the scalar action on $\varphi$). Let
$G := \mathcal{F}^{-1}(|p|^{-2})$, so that $G(x) = (4\pi^2)^{-1}|x|^{-2}$ and $\Delta G = -\delta_0$.
For $u\in\mathcal{O}$ let
\begin{equation*}
\mu_u := (E\equiv u)\otimes\mathcal{N}(0,G)
\end{equation*}
be the state in which $E$ is the constant $u$ and $\varphi$ is the centred Gaussian field of
covariance $G$, and let the probe be
\begin{equation}\label{10.7:eq-flux-counterexample-3}
P(x) := \langle E(x),\nabla^4\varphi(x)\rangle
= E^{\mu\nu\rho\sigma}(x)\,\partial_\mu\partial_\nu\partial_\rho\partial_\sigma\varphi(x),
\end{equation}
a $D$-scalar; under $\mu_u$ one has $P = u(\partial)\varphi$. Then the following hold.
\begin{enumerate}
\item[(a)] The family $(\mu_u)_{u\in\mathcal{O}}$ consists of translation-invariant states, depends
weak-$\ast$ continuously on $u$, and is covariant: the two-point kernel
$G_u := u(\partial)^2G = \mathcal{F}^{-1}(u(p)^2/|p|^2)$ satisfies $G_u\circ R = G_{R^{-1}\cdot u}$
for $R\in O(4)$. The moments are given by kernels smooth off the diagonals, and every $\mu_u$
satisfies the same identities $\mu_u(P(-\Delta g)F) = 0$ for $F$ supported away from
$\operatorname{supp} g$.
\item[(b)] The state $\bar\mu := \mu_{v_4}$ (the mixture with mixing measure
$\lambda = \delta_{v_4}$, carried by the one-point $W_4$-orbit $\{v_4\}$) is $W_4$-invariant,
translation-invariant, and $\vartheta$-reflection positive on the full polynomial and exponential
algebra.
\item[(c)] The state $\bar\mu$ is the limit of the $W_4$-invariant reflection-positive states
$(E\equiv v_4)\otimes(\text{massive field on the torus } \mathbb{R}^4/\ell\mathbb{Z}^4)$ as
$\ell\to\infty$ and then the mass tends to $0$. The connected functions of $P$ under $\bar\mu$ have
every property of the following structural list: they are multilinear, symmetric and real; they are
invariant under $\mathbb{R}^4\rtimes W_4$; they obey position-free bounds on separated tuples in
balls; they are distributions of finite order off the diagonals; the centred functions are
Osterwalder--Schrader positive; and the Osterwalder--Schrader reconstruction applies to them.
\item[(d)] The rotational anomaly $c^{\bar\mu}$ of $\bar\mu$ does not vanish identically, already
for $n = 2$: there are $f_1,f_2$ and $\omega$ with $c^{\bar\mu}(f_1,f_2;\omega)\neq 0$.
\item[(e)] There is no local Wick-polynomial symmetric tensor $T$ built from $(E,\varphi)$ for which
$\bar\mu$ satisfies the flux identity \eqref{10.7:eq-uniform-mixtures-a}.
\end{enumerate}
Consequently, the conjunction of: $O(4)$-invariant local linear dynamics (Schwinger--Dyson
identities); the structural list in (c); translation invariance; geometric reflection positivity;
being a limit of $W_4$-invariant reflection-positive states on cubic periodic boxes; and a local
probe that is a $D$-scalar, does not imply $c = 0$. Among the properties through which the vanishing
of the anomaly is derived from the rotational Ward identity and the dimension-six coefficient bound,
the only one that this system violates is the flux identity, which those two hypotheses deliver.
\end{proposition}

\begin{proof}
\emph{The quartic $v_4$ and the orbit $\mathcal{O}$.} In dimension four,
$\Delta\sum_\mu p_\mu^4 = \sum_\mu 12p_\mu^2 = 12|p|^2$, and, since $\Delta|p|^2 = 8$ and
$|\nabla|p|^2|^2 = 4|p|^2$,
\begin{equation*}
\Delta|p|^4 = 2\,|\nabla|p|^2|^2 + 2|p|^2\Delta|p|^2 = 8|p|^2 + 16|p|^2 = 24|p|^2 .
\end{equation*}
Hence $\Delta v_4 = 12|p|^2 - \tfrac12\cdot 24|p|^2 = 0$, and $v_4\in\mathcal{H}_4$. A signed
permutation of the coordinates leaves $\sum_\mu p_\mu^4$ and $|p|$ unchanged, so $v_4$ is
$W_4$-invariant. On the other hand
\begin{equation*}
v_4(e_1) = 1 - \tfrac12 = \tfrac12, \qquad
v_4\bigl((e_1+e_2)/\sqrt2\bigr) = 2\cdot\tfrac14 - \tfrac12 = 0 ,
\end{equation*}
and the two arguments are unit vectors related by a rotation in $SO(4)$; so $v_4$ is not even
$SO(4)$-invariant. The orbit $\mathcal{O}$ is compact as the continuous image of $O(4)$ and is
homogeneous by construction. The orbit $SO(4)\cdot v_4\subset\mathcal{O}$ is connected, being the
continuous image of the connected group $SO(4)$, and contains at least two points; a connected
metric space with two points is infinite, so $\mathcal{O}$ is infinite.

\emph{The probe.} Under $\mu_u$ the field $E$ is the constant $u$, so
\eqref{10.7:eq-flux-counterexample-3} reads
$P = u^{\mu\nu\rho\sigma}\partial_\mu\partial_\nu\partial_\rho\partial_\sigma\varphi = u(\partial)\varphi$.
Since $u$ has even degree, $u(\partial)$ is its own formal adjoint, and for a test function $f$
\begin{equation}\label{10.7:eq-flux-counterexample-4}
P(f) = \varphi\bigl(u(\partial)f\bigr), \qquad \operatorname{supp} u(\partial)f\subset\operatorname{supp} f .
\end{equation}
The contraction in \eqref{10.7:eq-flux-counterexample-3} is invariant under the diagonal action of
$O(4)$, so $P$ is a $D$-scalar.

\emph{Proof of (a).} Each $\mu_u$ is the product of a constant field and the translation-invariant
Gaussian $\mathcal{N}(0,G)$, hence is translation-invariant. By
\eqref{10.7:eq-flux-counterexample-4},
\begin{equation*}
\mu_u\bigl(P(f)P(g)\bigr) = \int\!\!\int f(x)\,G_u(x-y)\,g(y)\,dx\,dy ,\qquad G_u = u(\partial)^2G ,
\end{equation*}
and since $u(\partial)$ acts in Fourier space as multiplication by $u(ip) = u(p)$ ($u$ is quartic),
$G_u = \mathcal{F}^{-1}(u(p)^2/|p|^2)$. For $R\in O(4)$, substituting $p = Rq$ in the Fourier
integral,
\begin{equation*}
G_u(Rx) = \int e^{ip\cdot Rx}\,\frac{u(p)^2}{|p|^2}\,\frac{dp}{(2\pi)^4}
= \int e^{iq\cdot x}\,\frac{u(Rq)^2}{|q|^2}\,\frac{dq}{(2\pi)^4} = G_{R^{-1}\cdot u}(x),
\end{equation*}
because $(R^{-1}\cdot u)(q) = u(Rq)$. By Wick's theorem the moments of the $P(f_i)$ are sums over
pairings of products of the numbers $\langle f_i,G_uf_j\rangle$, each quadratic in $u$, and
$\mu_u(e^{iP(f)}) = \exp(-\tfrac12\langle f,G_uf\rangle)$; these depend continuously on $u$, which is
the weak-$\ast$ continuity. The kernels of the moments are products of the functions
$G_u(x_i-x_j)$; as $G$ is smooth off $0$, so is $G_u = u(\partial)^2G$, and the kernels are smooth off
the diagonals. Finally let $F$ be a polynomial or exponential in probes $P(f_j)$ with
$\operatorname{supp} f_j$ disjoint from $\operatorname{supp} g$, and put $k := u(\partial)g$,
$k_j := u(\partial)f_j$. Then $P(-\Delta g) = \varphi(-\Delta k)$ and, since $-\Delta G = \delta_0$,
\begin{equation*}
\mu_u\bigl(\varphi(-\Delta k)\varphi(k_j)\bigr) = \langle -\Delta k, G k_j\rangle = \langle k,k_j\rangle = 0 ,
\end{equation*}
because $\operatorname{supp}k\subset\operatorname{supp}g$ and
$\operatorname{supp}k_j\subset\operatorname{supp}f_j$ are disjoint. Gaussian integration by parts
expresses $\mu_u(\varphi(-\Delta k)F)$ as a combination of these covariances, so
$\mu_u(P(-\Delta g)F) = 0$. These identities do not depend on $u$; together with $\partial_\mu E = 0$,
which holds since $E$ is constant, they are the dynamics \eqref{10.7:eq-flux-counterexample-2}.

\emph{Proof of (b).} Since $\{v_4\}$ is a $W_4$-orbit, (a) gives $G_{v_4}\circ w = G_{w^{-1}\cdot v_4} = G_{v_4}$
for $w\in W_4$, so $\bar\mu$ is $W_4$-invariant; it is translation-invariant by (a). By
\eqref{10.7:eq-flux-counterexample-4}, $P(f) = \varphi(v_4(\partial)f)$ with
$\operatorname{supp}v_4(\partial)f\subset\operatorname{supp}f$, and
$v_4(\partial) = \sum_\mu\partial_\mu^4 - \tfrac12\Delta^2$ is even in $\partial_t$, hence commutes with
the time reflection $\theta$: $v_4(\partial)(f\circ\theta) = (v_4(\partial)f)\circ\theta$. Thus
$\vartheta(P(f)) = \varphi((v_4(\partial)f)\circ\theta)$, and an element of the polynomial or
exponential algebra supported in $\{t>0\}$ is a function of fields $\varphi(h)$ with $h$ supported
in $\{t>0\}$. The reflection positivity of $\bar\mu$ therefore reduces to that of the Gaussian field
$\mathcal{N}(0,G)$ under the standard reflection, which is
\eqref{10.7:eq-quasi-free-mixtures-a} (Corollary~\ref{10.7:cor-quasi-free-mixtures}).

\emph{Proof of (c).} The convergence of the states
$(E\equiv v_4)\otimes(\text{massive torus field})$ to $\bar\mu$, as $\ell\to\infty$ and then the
mass tends to $0$, is Theorem~\ref{10.7:thm-quasi-free-boxes} in the form
\eqref{10.7:eq-quasi-free-boxes-a}. Under $\bar\mu$ the probes are linear in the real Gaussian field
$\varphi$; their connected functions are therefore multilinear, symmetric and real, those of order
$n\ge 3$ vanish, and the connected two-point function has kernel $G_{v_4}$. Invariance under
$\mathbb{R}^4\rtimes W_4$ is (a) and (b). The kernel $G_{v_4} = v_4(\partial)^2G$ is smooth off $0$
and homogeneous of degree $-2-8 = -10$; hence, if $\operatorname{supp}f_1$ and
$\operatorname{supp}f_2$ are at distance at least $\mathsf{d}$,
\begin{equation*}
\Bigl|\int\!\!\int f_1(x)G_{v_4}(x-y)f_2(y)\,dx\,dy\Bigr|
\le \mathsf{d}^{-10}\sup_{|z|=1}|G_{v_4}(z)|\;\|f_1\|_{L^1}\|f_2\|_{L^1},
\end{equation*}
a bound free of the positions of the supports, and off the diagonals the connected functions are
distributions of order zero. The Osterwalder--Schrader positivity of the centred functions follows
from the reflection positivity in (b), and with it the Osterwalder--Schrader reconstruction applies.

\emph{Proof of (d).} Expanding,
$v_4(\partial)^2 = (\sum_\mu\partial_\mu^4)^2 - (\sum_\mu\partial_\mu^4)\Delta^2 + \tfrac14\Delta^4$.
Constant-coefficient operators commute and $\Delta G = -\delta_0$, so every term containing $\Delta$,
applied to $G$, is a derivative of $\delta_0$ and vanishes off $0$. Hence, off the diagonal,
\begin{equation*}
G_{v_4} = \Bigl(\sum_\mu\partial_\mu^4\Bigr)^{2}G
= \frac{1}{4\pi^2}\Bigl(\sum_\mu\partial_\mu^8 + 2\sum_{\mu<\nu}\partial_\mu^4\partial_\nu^4\Bigr)|x|^{-2}.
\end{equation*}
At $x = re_1$: $\partial_1^8$ acts on $x_1^{-2}$ and gives $9!\,r^{-10} = 362880\,r^{-10}$; for
$\nu\neq 1$, the coefficient of $y^8$ in $(r^2+y^2)^{-1}$ is $r^{-10}$, so
$\partial_\nu^8|x|^{-2} = 8!\,r^{-10} = 40320\,r^{-10}$; $\partial_\nu^4(x_1^2+y^2)^{-1}$ at $y=0$
equals $4!\,x_1^{-6}$, and $\partial_1^4x_1^{-6} = 3024\,x_1^{-10}$, so
$\partial_1^4\partial_\nu^4|x|^{-2} = 72576\,r^{-10}$; for distinct $\mu,\nu\neq 1$ the coefficient of
$y^4z^4$ in $(r^2+y^2+z^2)^{-1}$ is $6r^{-10}$, so $\partial_\mu^4\partial_\nu^4|x|^{-2} = 3456\,r^{-10}$.
Therefore
\begin{equation*}
\Bigl(\sum_\mu\partial_\mu^4\Bigr)^{2}|x|^{-2}\Big|_{x=re_1}
= (362880 + 3\cdot 40320 + 6\cdot 72576 + 6\cdot 3456)\,r^{-10} = 940032\,r^{-10}.
\end{equation*}
At $x = (1,1,1,1)$: from
$(x_1^2+3)^{-1} = (2i\sqrt3)^{-1}[(x_1-i\sqrt3)^{-1} - (x_1+i\sqrt3)^{-1}]$ and
$1\pm i\sqrt3 = 2e^{\pm i\pi/3}$, whose ninth powers both equal $-2^9$, one has
$\partial_\mu^8|x|^{-2} = 0$. Next, $\partial_y^4(s+y^2)^{-1} = 24(5y^4-10y^2s+s^2)(y^2+s)^{-5}$; at
$y = 1$, $s = x_1^2+2$, and with $q := x_1^2+3$ this equals $24(q^{-3}-12q^{-4}+16q^{-5})$. Writing
$x_1 = 1+t$, $q = 4+2t+t^2$, the coefficients of $t^4$ in $q^{-3}$, $q^{-4}$, $q^{-5}$ are $-9/1024$,
$-15/4096$, $-5/4096$, so
\begin{equation*}
\partial_1^4\partial_2^4|x|^{-2}\big|_{(1,1,1,1)}
= 24\cdot 24\cdot\Bigl(-\frac{9}{1024} + \frac{180}{4096} - \frac{80}{4096}\Bigr) = 9 .
\end{equation*}
By symmetry every $\partial_\mu^4\partial_\nu^4$, $\mu\neq\nu$, gives $9$, and by homogeneity of
degree $-10$,
\begin{equation*}
\Bigl(\sum_\mu\partial_\mu^4\Bigr)^{2}|x|^{-2}\Big|_{x=\frac r2(1,1,1,1)}
= 2^{10}\cdot 12\cdot 9\;r^{-10} = 110592\,r^{-10}.
\end{equation*}
The points $re_1$ and $\tfrac r2(1,1,1,1)$ have equal norm and are related by a rotation in
$SO(4)$, and $G_{v_4}$ is continuous off $0$; taking $f_1$ concentrated near $0$ and $f_2$ near
$re_1$, and comparing with the rotated pair, the connected two-point function $S^{\bar\mu}_2$ of
$P$ under $\bar\mu$ is not $SO(4)$-invariant. By the equivalence between rotation invariance and
the vanishing of the anomaly in \eqref{10.7:eq-anomaly-flux-a}
(Proposition~\ref{10.7:prop-anomaly-flux}), which uses only bounds of the type of the regularity
condition in \eqref{10.7:eq-uniform-mixtures-a}, some $c^{\bar\mu}(f_1,f_2;\omega)\neq 0$.

Structurally: $v_4^2 = h_8 + |p|^2q_6$ with $h_8\in\mathcal{H}_8$ (harmonic homogeneous of degree
eight) and $q_6$ homogeneous of degree six. Here $h_8\neq 0$: otherwise $|p|^2$, which is prime in
$\mathbb{C}[p]$, would divide $v_4$, whereas $v_4(1,i,0,0) = 1+1-\tfrac12\cdot 0 = 2\neq 0$ on the
null quadric. Off $0$ the term $\Delta q_6(\partial)G = -q_6(\partial)\delta_0$ vanishes, so
$v_4(\partial)^2G = h_8(\partial)G$, which is proportional to $h_8(x)|x|^{-18}$; this is never
rotation-invariant, since a rotation-invariant $h_8$ would be a multiple of $|x|^8$, which is not
harmonic.

\emph{Proof of (e).} Let $T$ be a local Wick-polynomial symmetric tensor built from $(E,\varphi)$.
Under $\bar\mu$, $E\equiv v_4$ is constant and the probes $P(y_1)$, $P(y_2)$ are linear in $\varphi$;
by Wick's theorem only the part of $T$ quadratic in $\varphi$ contributes to
$\langle T(x);P(y_1);P(y_2)\rangle$, and that contribution is a sum of products of two factors
$\partial^a v_4(\partial)G(x-y_i)$, $i = 1,2$. Each factor is homogeneous of degree $-6-|a|$ off $0$,
hence $O(|x|^{-6-|a|})$ for $x$ in the shell where $\nabla\chi_{\mathcal{R}}\neq 0$ and $y_i$ in
$\operatorname{supp}f_i$. The flux weight $h^{\omega}_{\chi_{\mathcal{R}}}$ is supported in that shell,
of volume $O(\mathcal{R}^4)$, so with $C$ independent of $\mathcal{R}$
\begin{equation*}
\bigl|\langle T(h^{\omega}_{\chi_{\mathcal{R}}});P(f_1);P(f_2)\rangle\bigr|
\le C\,\mathcal{R}^{4}\mathcal{R}^{-12}\longrightarrow 0 \qquad (\mathcal{R}\to\infty).
\end{equation*}
The flux identity \eqref{10.7:eq-uniform-mixtures-a} would make the left-hand side equal to
$c^{\bar\mu}(f_1,f_2;\omega)$ for every $\mathcal{R}$, and for the pair of (d) this is non-zero.

\emph{Conclusion.} The system $(E,\varphi)$ with the state $\bar\mu$ and the probe $P$ has
$O(4)$-invariant local linear dynamics by (a), the structural list, translation invariance and
geometric reflection positivity by (b) and (c), is a limit of $W_4$-invariant reflection-positive
states on cubic periodic boxes by (c), and its probe is a local $D$-scalar; yet $c^{\bar\mu}\neq 0$
by (d). By (e) it violates the flux identity, which is the property among those named in the
statement that the rotational Ward identity and the dimension-six coefficient bound deliver.
\end{proof}

\begin{remark}[A full counter-example requires broken rotation symmetry]
\label{10.7:rem-broken-rotation}
Consider the following data.
\begin{itemize}
\item[(a)] A compact homogeneous space $\mathcal{O}$ of $O(4)$ with infinitely many points.
\item[(b)] A family $(\mu_u)_{u\in\mathcal{O}}$ of translation-invariant states, weak-$*$ continuous
in $u$ and covariant. Each $\mu_u$ satisfies
the Schwinger--Dyson identities of one local, $O(4)$-invariant, geometrically
reflection-positive dynamics in four dimensions. The probe $P$ is an $O(4)$-scalar
local field, $T$ is the symmetric stress tensor of this same dynamics, and one action of $O(4)$
is used throughout. Each $\mu_u$ also satisfies the
flux identity for all admissible $\chi$, together with the regularity and separation
conditions, in the form displayed in \eqref{10.7:eq-uniform-mixtures-a}.
\item[(c)] A probability $\lambda\neq m_{\mathcal{O}}$ on $\mathcal{O}$; for selection by cubic
boxes, $\lambda$ is carried by a finite $W_4$-orbit. In addition, one is given $W_4$-invariant,
$\vartheta$-reflection-positive states $\nu_j$ of $W_4$-invariant reflection-positive local
dynamics on the cubic periodic tori $\mathbb{R}^4/\ell_j\mathbb{Z}^4$, $\ell_j\to\infty$,
converging to $\int_{\mathcal{O}}\mu_u\,d\lambda(u)$.
\end{itemize}
The question is whether
\begin{equation}\label{10.7:eq-broken-rotation-a}
\text{there exist } \mathcal{O},\ (\mu_u)_{u\in\mathcal{O}},\ \lambda\neq m_{\mathcal{O}},\
(\nu_j)_{j\ge 1} \text{ with (a), (b), (c)}.
\end{equation}
Any system as in \eqref{10.7:eq-broken-rotation-a} has the following properties.
\begin{itemize}
\item[(1)] It satisfies every one of the following conjuncts:
\begin{itemize}
\item the infinite-volume content of the ultraviolet input, namely the exact rotational Ward
identity with vanishing bulk and observable-side terms, and a flux independent of $\chi$;
\item the structural list of properties;
\item translation invariance, including that of the stress tensor;
\item geometric reflection positivity;
\item the cubic periodic-box limit.
\end{itemize}
\item[(2)] Its rotational anomaly satisfies $c\not\equiv 0$.
\item[(3)] It has no quasi-free solution.
\item[(4)] It has no frozen-background Gaussian solution.
\end{itemize}
The existence of such a system is the existence of a reflection-positive local
four-dimensional quantum field theory with spontaneously broken rotation symmetry, selected by
cubic boxes. Its existence is not asserted: it is neither constructed nor excluded here.
\end{remark}

\begin{proof}
We prove the properties (1)--(4) of a system as in \eqref{10.7:eq-broken-rotation-a}; its
existence is not asserted.

Properties (1) and (2) follow from the anomaly theorem for uniform mixtures
(Theorem~\ref{10.7:thm-uniform-mixtures}), applied to the family $(\mu_u)_{u\in\mathcal{O}}$ of
(b), the mixing measure $\lambda\neq m_{\mathcal{O}}$ of (c), and the approximating states
$\nu_j$ of (c). That theorem gives each conjunct listed in (1); since $\lambda\neq m_{\mathcal{O}}$
by (c), the mixture is not the uniform one, and the same theorem gives $c\not\equiv 0$, which
is (2).

Property (3) follows from three results applied to the same system. The first is the theorem
that quasi-free fields positive for the geometric reflection are rotation invariant
(Theorem~\ref{10.7:thm-quasi-free-reflection}). The second is its corollary on mixtures of
quasi-free solutions (Corollary~\ref{10.7:cor-quasi-free-mixtures}). The third is the theorem
that quasi-free fields on periodic boxes converge to the invariant limit
(Theorem~\ref{10.7:thm-quasi-free-boxes}). Each of them gives $c\equiv 0$ for, respectively,
a quasi-free solution positive for the geometric reflection, a mixture of quasi-free
solutions, and a cubic periodic-box limit of quasi-free states with $O(4)$-scalar probes; by
(2) the system is none of these, so it has no quasi-free solution.

Property (4) follows from the hypercubic counter-example without the flux identity
(Proposition~\ref{10.7:prop-flux-counterexample}), which shows that the system has no
frozen-background Gaussian solution.
\end{proof}

\begin{remark}[What holds without the far-sphere flux hypothesis]\label{10.7:rem-without-clustering}
Throughout, $R_t = e^{t\omega}$ is the one-parameter rotation group generated by an antisymmetric
matrix $\omega$, acting on test functions as in Lemma~\ref{10.7:lem-limit-differentiability}, and
$c(\vec f,\omega)$ is the rotational anomaly defined there. The far-sphere flux hypothesis
$(\mathrm{IR\text{-}fl})_{\omega_0}$ of clause (vi) is not assumed. That hypothesis says that the
flux of the lattice rotation current through far spheres vanishes along one sequence of radii. Nor
is the stress-tensor decay condition assumed, which is sufficient for it: that the connected
correlation of one far-away stress-tensor insertion with the observables decays faster than
$|x|^{-4}$, uniformly in the cutoff. The following hold for every subsequential limit
$(S^T_n)_{n\ge 2}$.
\begin{enumerate}
\item[(i)] The map $c(\cdot,\omega)$ is real, multilinear, symmetric, linear in $\omega$,
translation-invariant and $W_4$-equivariant. These properties use the hypotheses (b1) and (b2) of
Proposition~\ref{10.7:prop-anomaly-flux} alone.
\item[(ii)] Assume the two further hypotheses of Proposition~\ref{10.7:prop-anomaly-flux}, the
rotational Ward identity and the dimension-six coefficient bound, in the forms stated there. Then,
for every $\vec f\in\mathcal{T}^{(2)}_n(\mathcal{R},\mathsf{d})$ and every admissible $\chi$,
\begin{equation*}
c(\vec f,\omega)
=\lim_{j}\kappa_j\bigl(\mathfrak{F}_\chi;\mathcal{O}(f_1),\dots,\mathcal{O}(f_n)\bigr).
\end{equation*}
That is, the anomaly is the limit flux through every admissible shell. If a translate
$\vec f(\cdot-a)$, $a\in\mathbb{R}^4$, lies in $\mathcal{T}^{(2)}_n(\mathcal{R},\mathsf{d})$, then
$c(\vec f,\omega)=c(\vec f(\cdot-a),\omega)$, and this is the limit flux computed for the
translate.
\item[(iii)] For $\vec f\in\mathcal{T}^{(2)}_n(\mathcal{R},\mathsf{d})$ the map
$t\mapsto S^T_n(R_t\cdot\vec f)$ is $C^1$ with derivative $c(R_t\cdot\vec f,\omega)$. Hence
\begin{equation}\label{10.7:eq-without-clustering-1}
S^T_n(e^{t\omega}\cdot\vec f)-S^T_n(\vec f)=\int_0^t c(e^{s\omega}\cdot\vec f,\omega)\,ds ,
\end{equation}
so the change under a finite rotation is the integrated flux at infinity.
\item[(iv)] Under the same two hypotheses as in (ii), and with (b1) and (b2) as in (i), the limit
$(S^T_n)_{n\ge 2}$ is invariant under $O(4)$ if and only if $c(\vec h,\omega)=0$ for every
$\omega$, every $n\ge 2$, all $\mathcal{R},\mathsf{d}>0$ and every
$\vec h\in\mathcal{T}^{(2)}_n(\mathcal{R},\mathsf{d})$.
\item[(v)] Nothing of this kind survives in finite volume. If $R\in O(4)$ maps $2L^m\mathbb{Z}^4$
into itself, then $R\in W_4$. Hence the only linear isometries of $\mathbb{R}^4$ that descend to
$\mathbb{T}_m$ are the $384$ elements of $W_4$.
\end{enumerate}
\end{remark}

\begin{proof}
Assertion (i) is the list of algebraic properties of the anomaly in
Proposition~\ref{10.7:prop-anomaly-flux}, where these properties are derived from (b1) and (b2)
only.

Assertion (ii), for tuples in $\mathcal{T}^{(2)}_n(\mathcal{R},\mathsf{d})$, is the flux
identification in Proposition~\ref{10.7:prop-anomaly-flux}, that is, the identification of the
anomaly with the limit flux through every admissible shell. Let $\vec f$ be a
tuple and $a\in\mathbb{R}^4$ a vector such that the translate $\vec f(\cdot-a)$ lies in that
class. By the translation invariance in (i), $c(\vec f,\omega)=c(\vec f(\cdot-a),\omega)$, and the
identification applies to the translate.

For (iii), fix $\vec f\in\mathcal{T}^{(2)}_n(\mathcal{R},\mathsf{d})$. If $\omega=0$, both sides
of \eqref{10.7:eq-without-clustering-1} vanish: the left because $R_t$ is the identity, the right
by linearity of $c$ in $\omega$. Let $\omega\neq 0$ and $\tau\in\mathbb{R}$. Since $R_\tau$ is a
linear isometry fixing the origin, $R_\tau\cdot\vec f$ has the same regularity as $\vec f$, its
supports lie in $B(0,\mathcal{R}/2)$, and their mutual distances are unchanged. Hence
$R_\tau\cdot\vec f\in\mathcal{T}^{(2)}_n(\mathcal{R},\mathsf{d})$.

Lemma~\ref{10.7:lem-limit-differentiability}, applied to $\vec h=R_\tau\cdot\vec f$, gives
\begin{equation*}
\lim_{t\to 0}\frac{S^T_n(R_t\cdot R_\tau\cdot\vec f)-S^T_n(R_\tau\cdot\vec f)}{t}
= c(R_\tau\cdot\vec f,\omega).
\end{equation*}
Since $R_tR_\tau=R_{t+\tau}$, the map $t\mapsto S^T_n(R_t\cdot\vec f)$ is differentiable at every
$\tau$ with derivative $c(R_\tau\cdot\vec f,\omega)$.

It remains to show that $D(\tau):=c(R_\tau\cdot\vec f,\omega)$ is continuous in $\tau$. Put
$\phi(\tau):=S^T_n(R_\tau\cdot\vec f)$; being differentiable everywhere, $\phi$ is continuous. For
$0<|t|\le t_0$, with $t_0$ as in Lemma~\ref{10.7:lem-limit-differentiability} (it depends on
$\mathsf{d}$, $\omega$, $\mathcal{R}$ only), the bound of that lemma for
$\vec h=R_\tau\cdot\vec f$ reads
\begin{equation*}
\Bigl|\frac{\phi(\tau+t)-\phi(\tau)}{t}-D(\tau)\Bigr|\le \varsigma_\tau(t),
\end{equation*}
where $\varsigma_\tau(t)$ is the right-hand side of that bound for $\vec h=R_\tau\cdot\vec f$.
This right-hand side does not depend on $\tau$. Indeed, $C_*$ depends on $n$, $\mathsf{d}$,
$\mathcal{R}$ only; $\|R_\tau\cdot\varphi\|_{\sharp}=\|\varphi\|_{\sharp}$ for every test function
$\varphi$, since the sup norm and the Lipschitz constant are unchanged under composition with an
isometry; $\delta_\omega(R_\tau\cdot\varphi)=R_\tau\cdot\delta_\omega\varphi$, since
$\delta_\omega\varphi$ is the pointwise derivative at $t=0$ of $R_t\cdot\varphi$ and
$R_tR_\tau=R_\tau R_t$; and the rotation modulus satisfies $\mu_{R_\tau\cdot\varphi}=\mu_\varphi$,
because conjugation by $R_\tau$ preserves the distance of a rotation to the identity. Hence
$\varsigma(t):=\varsigma_\tau(t)$ is independent of $\tau$, and
$\varsigma(t)\to 0$ as $t\to 0$ by the lemma. The continuous functions
$\tau\mapsto(\phi(\tau+t)-\phi(\tau))/t$ therefore converge to $D$ uniformly on $\mathbb{R}$ as
$t\to 0$, and $D$ is continuous. So $t\mapsto S^T_n(R_t\cdot\vec f)$ is $C^1$, and the fundamental
theorem of calculus on $[0,t]$ gives \eqref{10.7:eq-without-clustering-1}.

The fourth assertion is the equivalence, in Corollary~\ref{10.7:cor-clustering-dichotomy}, between
the vanishing of $c(\vec h,\omega)$ for every $\omega$ and every tuple of every class and
$O(4)$-invariance of the limit; the hypotheses of that corollary are the two hypotheses of (ii)
together with (b1) and (b2).

For (v), let $R\in O(4)$ satisfy $R(2L^m\mathbb{Z}^4)\subseteq 2L^m\mathbb{Z}^4$. Since $R$ is
linear, dividing by $2L^m$ gives $R\mathbb{Z}^4\subseteq\mathbb{Z}^4$. By part (1) of
Lemma~\ref{10.7:lem-lattice-stabiliser}, $R\in W_4$. Conversely, that lemma shows that every
element of $W_4$ maps $\mathbb{Z}^4$, and hence $2L^m\mathbb{Z}^4$, into itself.

A linear map of $\mathbb{R}^4$ induces a well-defined map of
$\mathbb{T}_m=\mathbb{R}^4/2L^m\mathbb{Z}^4$ exactly when it maps $2L^m\mathbb{Z}^4$ into itself.
So the linear isometries descending to $\mathbb{T}_m$ are the elements of $W_4$. These are the
signed permutation matrices, of which there are $2^4\cdot 4!=384$.
\end{proof}

\begin{lemma}[Regrouping the divergence term by plaquette planes]\label{10.7:lem-plane-regrouping}
Assume the displayed forms \eqref{10.7:eq-annulus-functionals-a} of the annulus functionals in
Definition~\ref{10.7:def-annulus-functionals}, with lattice spacing $\varepsilon$, and let
$e_1,\dots,e_4$ be the standard unit vectors of $\mathbb{R}^4$. For a plane
$\pi=\{\rho,\mu\}$, $\rho\neq\mu$, put
\begin{equation*}
s_\pi:=\tfrac{\varepsilon}{2}\,(e_\rho+e_\mu),\qquad
V_\pi(y):=(w_\rho+w_\mu)(y-s_\pi),
\end{equation*}
and set
\begin{equation*}
\bar V:=\frac16\sum_\pi V_\pi,\qquad V^0_\pi:=V_\pi-\bar V,
\qquad\text{so that}\quad \sum_\pi V^0_\pi=0 .
\end{equation*}
Write $\vec V^{\,0}=(V^0_\pi)_\pi$. Then
\begin{equation}\label{10.7:eq-plane-regrouping-1}
J^4=\kappa_N\sum_\pi\mathcal{O}_\pi(V_\pi)
=\kappa_N\,\mathcal{O}(\bar V)+\kappa_N\,g_0^{2}\,\mathcal{T}(\vec V^{\,0}).
\end{equation}
Here $\mathcal{O}_\pi$ is the plane-restricted observable,
$\mathcal{O}_\pi(\varphi)=\sum_{p\parallel\pi}\varphi(x(p))\,a_p$ with
$a_p:=1-\operatorname{Re}\operatorname{tr}U(\partial p)$, the sum running over the plaquettes
parallel to $\pi$, and
$\mathcal{T}(\vec g)=g_0^{-2}\sum_\pi\mathcal{O}_\pi(g_\pi)$.
\end{lemma}

\begin{proof}
By Definition~\ref{10.7:def-annulus-functionals},
\begin{equation*}
J^4=\kappa_N\sum_x\sum_\mu w_\mu(x)P_\mu(x),\qquad
P_\mu(x)=\sum_{\rho\neq\mu}a_{\rho\mu}(x),\qquad
a_{\rho\mu}(x)=1-\operatorname{Re}\operatorname{tr}U_{\rho\mu}(x).
\end{equation*}

First, $a_{\rho\mu}(x)=a_{\mu\rho}(x)$ for every $N$. Indeed $U_{\mu\rho}(x)=U_{\rho\mu}(x)^{-1}$,
and for $V\in SU(N)$ one has $V^{-1}=V^{*}$, hence
$\operatorname{tr}V^{-1}=\overline{\operatorname{tr}V}$ and
$\operatorname{Re}\operatorname{tr}V^{-1}=\operatorname{Re}\operatorname{tr}V$.
We may therefore write $a_\pi(x):=a_{\rho\mu}(x)=a_{\mu\rho}(x)$ for $\pi=\{\rho,\mu\}$.

Second, the double sum $\sum_\mu\sum_{\rho\neq\mu}w_\mu(x)a_{\rho\mu}(x)$ runs over ordered
pairs $(\rho,\mu)$ with $\rho\neq\mu$. Grouping the pair $(\rho,\mu)$ with $(\mu,\rho)$ and using
the symmetry just proved gives
\begin{equation}\label{10.7:eq-plane-regrouping-2}
\sum_\mu w_\mu(x)P_\mu(x)=\sum_{\pi=\{\rho,\mu\}}(w_\rho+w_\mu)(x)\,a_\pi(x).
\end{equation}

Third, the map sending $(x,\pi)$, $\pi=\{\rho,\mu\}$, to the plaquette $p$ with corners
$x$, $x+\varepsilon e_\rho$, $x+\varepsilon e_\mu$, $x+\varepsilon e_\rho+\varepsilon e_\mu$ is a
bijection from (sites)$\times$(planes) onto the plaquettes, both on $\mathbb{R}^4$ and on the
torus; the image $p$ is parallel to $\pi$, its barycentre is $x(p)=x+s_\pi$, and
$a_\pi(x)=a_p$, because $a_\pi(x)$ does not depend on the orientation of $\partial p$. On the
torus the values of $w_\mu$ and $V_\pi$ are read at representatives in $\mathbb{R}^4$ of the
sites and barycentres, and these representatives are unambiguous under the floor
$L^m\ge 2(2\mathcal{R}+3)$. Summing \eqref{10.7:eq-plane-regrouping-2} over $x$ and reindexing
by this bijection, with $x=x(p)-s_\pi$, we obtain
\begin{equation*}
\begin{split}
\sum_x\sum_\mu w_\mu(x)P_\mu(x)
&=\sum_\pi\sum_{p\parallel\pi}(w_\rho+w_\mu)\bigl(x(p)-s_\pi\bigr)\,a_p\\
&=\sum_\pi\sum_{p\parallel\pi}V_\pi(x(p))\,a_p
=\sum_\pi\mathcal{O}_\pi(V_\pi).
\end{split}
\end{equation*}
Multiplying by $\kappa_N$ gives the first equality in \eqref{10.7:eq-plane-regrouping-1}.

Fourth, split $V_\pi=\bar V+V^0_\pi$ and use the linearity of $\mathcal{O}_\pi$ in its argument.
Since every plaquette is parallel to exactly one plane,
$\sum_\pi\mathcal{O}_\pi(\bar V)=\mathcal{O}(\bar V)$; and by the definition of $\mathcal{T}$,
$\sum_\pi\mathcal{O}_\pi(V^0_\pi)=g_0^{2}\,\mathcal{T}(\vec V^{\,0})$. This is the second
equality in \eqref{10.7:eq-plane-regrouping-1}. Finally,
$\sum_\pi V^0_\pi=\sum_\pi V_\pi-6\bar V=0$ by the definition of $\bar V$, there being six planes.
\end{proof}

The identity \eqref{10.7:eq-plane-regrouping-1} is exact: the shift $s_\pi$ is carried inside
$V_\pi$ and is not expanded.

\begin{lemma}[Bounds on the radial divergence weights]\label{10.7:lem-divergence-weights}
Let $\omega$ be the real antisymmetric $4\times 4$ matrix of
Definition~\ref{10.7:def-annulus-functionals}, $\chi=\chi_{\mathcal{R}}$ the radial cut-off,
$\xi=\chi\cdot(\omega y)$ and $w_\mu$ the forward differences of
Definition~\ref{10.7:def-annulus-functionals}. Let $s_\pi$, $V_\pi$, $\bar V$, $V^0_\pi$ be as in
Lemma~\ref{10.7:lem-plane-regrouping}. Let $\mathcal{R}\ge 12$, $0<\varepsilon\le 1$, and let $|\omega|$
denote the operator norm of $\omega$.

Of the cut-off $\chi_{\mathcal{R}}(y)=\chi_1(|y|/\mathcal{R})$ of
Definition~\ref{10.7:def-annulus-functionals} we use only the following. First, $\chi_1$ is
of class $C^3$, constant on $[0,1]$ and zero on $[2,\infty)$. Second, there are constants
$A^{\chi}_0,\dots,A^{\chi}_3$, independent of $\mathcal{R}$, such that
$\|D^k\chi_{\mathcal{R}}\|_\infty\le A^{\chi}_k\,\mathcal{R}^{-k}$ for $k=0,1,2,3$.

Here, for a $C^k$ function $f$ on $\mathbb{R}^4$,
\begin{equation*}
\|D^kf\|_\infty:=\sup_{y}\ \sup_{|v_1|=\dots=|v_k|=1}\bigl|D^kf(y)[v_1,\dots,v_k]\bigr|,
\end{equation*}
(for $k=0$ this is $\|f\|_\infty$), so that $\operatorname{Lip}f=\|Df\|_\infty$ for $f\in C^1$. We write
$\|\partial^k\xi\|_\infty:=\max_\mu\|D^k\xi_\mu\|_\infty$,
$\|\partial^kw\|_\infty:=\max_\mu\|D^kw_\mu\|_\infty$ and $\|w\|_\infty:=\max_\mu\|w_\mu\|_\infty$.
Let $\operatorname{Ann}_{\mathcal{R}}:=\{\mathcal{R}-2\le|y|\le 2\mathcal{R}+2\}$.
\begin{enumerate}
\item[(a)] Where $\chi$ is constant on the segment $[y,y+\varepsilon e_\mu]$, one has
$w_\mu(y)=\chi(y)\omega_{\mu\mu}=0$. Hence $\operatorname{supp}w_\mu\subset\{\mathcal{R}-1\le|y|\le
2\mathcal{R}+1\}$, and $V_\pi$, $\bar V$, $V^0_\pi$ are supported in $\operatorname{Ann}_{\mathcal{R}}$.
This set lies at distance at least $\mathcal{R}/2-2\ge\mathcal{R}/4$ from $\bar B(0,\mathcal{R}/2)$.
\item[(b)] With $\xi_\mu=\chi\,(\omega y)_\mu$,
\begin{equation*}
\|\partial\xi\|_\infty\le(2A^{\chi}_1+A^{\chi}_0)|\omega|,\quad
\|\partial^2\xi\|_\infty\le\frac{2(A^{\chi}_1+A^{\chi}_2)|\omega|}{\mathcal{R}},\quad
\|\partial^3\xi\|_\infty\le\frac{(2A^{\chi}_3+3A^{\chi}_2)|\omega|}{\mathcal{R}^2},
\end{equation*}
and
\begin{equation*}
\|w\|_\infty\le\|\partial\xi\|_\infty,\qquad
\|\partial w\|_\infty\le\|\partial^2\xi\|_\infty,\qquad
\|\partial^2 w\|_\infty\le\|\partial^3\xi\|_\infty .
\end{equation*}
\item[(c)] The field $\xi$ is divergence-free, $\operatorname{div}\xi\equiv0$, and
\begin{equation*}
\Bigl\|\sum_\mu w_\mu\Bigr\|_\infty\le2\varepsilon\|\partial^2\xi\|_\infty
\le\frac{4(A^{\chi}_1+A^{\chi}_2)|\omega|\varepsilon}{\mathcal{R}},\qquad
\Bigl\|\nabla\sum_\mu w_\mu\Bigr\|_\infty\le2\varepsilon\|\partial^3\xi\|_\infty .
\end{equation*}
\item[(d)] For $\pi=\{\rho,\mu\}$, and with $|s_\pi|=\varepsilon/\sqrt2$,
\begin{equation*}
\|V_\pi-(w_\rho+w_\mu)\|_\infty\le|s_\pi|\cdot2\|\partial w\|_\infty,\qquad
\operatorname{Lip}\bigl(V_\pi-(w_\rho+w_\mu)\bigr)\le|s_\pi|\cdot2\|\partial^2w\|_\infty .
\end{equation*}
\item[(e)] One has $\sum_\pi V_\pi=3\sum_\mu w_\mu+\sum_\pi[V_\pi-(w_\rho+w_\mu)]$. Furthermore
$\|\bar V\|_\sharp\le C_{\bar V}\varepsilon/\mathcal{R}\le C_{\bar V}\varepsilon$ and
$\|V^0_\pi\|_\sharp\le C_{V^0}$, where
\begin{align}
C_{\bar V}&:=(1+\sqrt2)|\omega|\bigl[2(A^{\chi}_1+A^{\chi}_2)+40M(3A^{\chi}_2+2A^{\chi}_3)\bigr],
\label{10.7:eq-divergence-weights-1}\\
C_{V^0}&:=2(A^{\chi}_0+2A^{\chi}_1)|\omega|+80M(A^{\chi}_1+A^{\chi}_2)|\omega|+C_{\bar V}.
\label{10.7:eq-divergence-weights-2}
\end{align}
These constants depend only on $|\omega|$, $M$ and $A^{\chi}_0,\dots,A^{\chi}_3$. The bounds hold
with these same constants for every $\mathcal{R}\ge12$ and every $0<\varepsilon\le1$; the bounds
on $\|\bar V\|_\sharp$ use only $\mathcal{R}\ge1$, and the bound on $\|V^0_\pi\|_\sharp$ uses only
$\mathcal{R}\ge2$.
\end{enumerate}
\end{lemma}

\begin{proof}
(a) Suppose $\chi$ takes the constant value $c$ on $[y,y+\varepsilon e_\mu]$. Then
\begin{equation*}
\xi_\mu(y+\varepsilon e_\mu)-\xi_\mu(y)=c\bigl[(\omega(y+\varepsilon e_\mu))_\mu-(\omega y)_\mu\bigr]
=c\,\varepsilon\,\omega_{\mu\mu}.
\end{equation*}
It follows that $w_\mu(y)=\chi(y)\omega_{\mu\mu}$, which vanishes because $\omega$ is antisymmetric.

Every point $z$ of the segment satisfies $|y|-\varepsilon\le|z|\le|y|+\varepsilon$, and $\varepsilon\le1$.
If $|y|\le\mathcal{R}-1$, the segment lies in $\{|z|\le\mathcal{R}\}$, where $\chi_{\mathcal{R}}$ is
constant because $\chi_1$ is constant on $[0,1]$. If $|y|\ge2\mathcal{R}+1$, the segment lies in
$\{|z|\ge2\mathcal{R}\}$, where $\chi_{\mathcal{R}}=0$. Hence $w_\mu(y)\ne0$ only if
$\mathcal{R}-1<|y|<2\mathcal{R}+1$, and $\operatorname{supp}w_\mu$ lies in the closed annulus stated.

Now $V_\pi(y)=(w_\rho+w_\mu)(y-s_\pi)$ and $|s_\pi|=\tfrac{\varepsilon}{2}|e_\rho+e_\mu|=\varepsilon/\sqrt2<1$.
So $V_\pi(y)\ne0$ only if $\mathcal{R}-2<|y|<2\mathcal{R}+2$. The functions $\bar V$ and $V^0_\pi$ are
linear combinations of the $V_\pi$, so all three are supported in $\operatorname{Ann}_{\mathcal{R}}$.

For $|y|\ge\mathcal{R}-2$ the distance from $y$ to $\bar B(0,\mathcal{R}/2)$ is at least
$|y|-\mathcal{R}/2\ge\mathcal{R}/2-2$. This is at least $\mathcal{R}/4$, since
$\mathcal{R}/2-2\ge\mathcal{R}/4$ exactly when $\mathcal{R}\ge 8$, and $\mathcal{R}\ge12$.

(b) On $\{|y|>2\mathcal{R}\}$ the field $\xi$ vanishes identically. On $\{|y|\le2\mathcal{R}\}$ we have
$|(\omega y)_\mu|\le|\omega y|\le2|\omega|\mathcal{R}$, and $|(\omega v)_\mu|\le|\omega|$ for $|v|=1$. By
the Leibniz rule, since $y\mapsto(\omega y)_\mu$ is linear,
\begin{align*}
D\xi_\mu[v_1]&=D\chi[v_1]\,(\omega y)_\mu+\chi\,(\omega v_1)_\mu,\\
D^2\xi_\mu[v_1,v_2]&=D^2\chi[v_1,v_2]\,(\omega y)_\mu
+D\chi[v_1](\omega v_2)_\mu+D\chi[v_2](\omega v_1)_\mu,
\end{align*}
and $D^3\xi_\mu[v_1,v_2,v_3]$ is $D^3\chi[v_1,v_2,v_3](\omega y)_\mu$ plus the three terms
$D^2\chi[v_i,v_j](\omega v_k)_\mu$, $k\in\{1,2,3\}$, $\{i,j\}=\{1,2,3\}\setminus\{k\}$. Inserting
$\|D^k\chi\|_\infty\le A^{\chi}_k\mathcal{R}^{-k}$ gives the three bounds:
\begin{align*}
\|\partial\xi\|_\infty&\le A^{\chi}_1\mathcal{R}^{-1}\cdot2|\omega|\mathcal{R}+A^{\chi}_0|\omega|,\\
\|\partial^2\xi\|_\infty&\le A^{\chi}_2\mathcal{R}^{-2}\cdot2|\omega|\mathcal{R}
+2A^{\chi}_1\mathcal{R}^{-1}|\omega|,\\
\|\partial^3\xi\|_\infty&\le A^{\chi}_3\mathcal{R}^{-3}\cdot2|\omega|\mathcal{R}
+3A^{\chi}_2\mathcal{R}^{-2}|\omega|.
\end{align*}
By the fundamental theorem of calculus,
\begin{equation}\label{10.7:eq-divergence-weights-3}
w_\mu(y)=\int_0^1\partial_\mu\xi_\mu(y+u\varepsilon e_\mu)\,du
=\int_0^1D\xi_\mu(y+u\varepsilon e_\mu)[e_\mu]\,du .
\end{equation}
Hence $|w_\mu|\le\|\partial\xi\|_\infty$. Differentiating under the integral,
$Dw_\mu(y)[v]=\int_0^1D^2\xi_\mu(y+u\varepsilon e_\mu)[e_\mu,v]\,du$, so
$\|Dw_\mu\|_\infty\le\|\partial^2\xi\|_\infty$. In the same way
$\|D^2w_\mu\|_\infty\le\|\partial^3\xi\|_\infty$.

(c) Near $y=0$ the function $\chi$ is constant, so $\nabla\chi=0$ there. For $y\ne0$ the chain rule gives
$\nabla\chi_{\mathcal{R}}(y)=\chi_1'(|y|/\mathcal{R})\,y/(|y|\mathcal{R})$. Therefore
\begin{equation*}
\operatorname{div}(\chi_{\mathcal{R}}\omega y)
=\chi_1'(|y|/\mathcal{R})\frac{y\cdot\omega y}{|y|\mathcal{R}}+\chi_{\mathcal{R}}(y)\operatorname{tr}\omega
=0 .
\end{equation*}
Both terms vanish: $y\cdot\omega y=(\omega^{T}y)\cdot y=-(\omega y)\cdot y$ forces $y\cdot\omega y=0$, and
$\operatorname{tr}\omega=0$.

Subtracting $\operatorname{div}\xi(y)=\sum_\mu\partial_\mu\xi_\mu(y)=0$ from the sum over $\mu$ of
\eqref{10.7:eq-divergence-weights-3} gives
\begin{equation*}
\sum_\mu w_\mu(y)=\sum_\mu\int_0^1\bigl[\partial_\mu\xi_\mu(y+u\varepsilon e_\mu)
-\partial_\mu\xi_\mu(y)\bigr]\,du .
\end{equation*}
By the mean value theorem along $e_\mu$, the bracket is at most $u\varepsilon\|\partial^2\xi\|_\infty$ in
modulus. Since $\int_0^1u\,du=\tfrac12$, the sum of the four terms is at most
$2\varepsilon\|\partial^2\xi\|_\infty$. Part (b) bounds this by
$4(A^{\chi}_1+A^{\chi}_2)|\omega|\varepsilon/\mathcal{R}$.

Since $\operatorname{div}\xi\equiv0$, also $\nabla\operatorname{div}\xi=0$. So for $|v|=1$,
\begin{equation*}
D\Bigl(\sum_\mu w_\mu\Bigr)(y)[v]=\sum_\mu\int_0^1\bigl[D^2\xi_\mu(y+u\varepsilon e_\mu)[e_\mu,v]
-D^2\xi_\mu(y)[e_\mu,v]\bigr]\,du .
\end{equation*}
Each bracket is at most $u\varepsilon\|\partial^3\xi\|_\infty$ in modulus, and the same count gives
$\|\nabla\sum_\mu w_\mu\|_\infty\le2\varepsilon\|\partial^3\xi\|_\infty$.

(d) Put $f:=w_\rho+w_\mu$. Then $V_\pi(y)-f(y)=f(y-s_\pi)-f(y)$. Its modulus satisfies
\begin{equation*}
|V_\pi(y)-f(y)|\le|s_\pi|\operatorname{Lip}f
\le|s_\pi|\bigl(\|Dw_\rho\|_\infty+\|Dw_\mu\|_\infty\bigr)\le|s_\pi|\cdot2\|\partial w\|_\infty .
\end{equation*}
Its gradient is $\nabla f(y-s_\pi)-\nabla f(y)$, whose norm is at most
$|s_\pi|\|D^2f\|_\infty\le|s_\pi|\cdot2\|\partial^2w\|_\infty$. This bounds the Lipschitz constant.

(e) There are six planes, and each index $\mu$ lies in exactly three of them. Hence
$\sum_\pi(w_\rho+w_\mu)=3\sum_\mu w_\mu$, which is the stated identity. Dividing by six,
\begin{equation*}
\bar V=\tfrac12\sum_\mu w_\mu+\tfrac16\sum_\pi[V_\pi-(w_\rho+w_\mu)] .
\end{equation*}
By (c), (d), (b) and $|s_\pi|=\varepsilon/\sqrt2$,
\begin{equation*}
\|\bar V\|_\infty\le\varepsilon\|\partial^2\xi\|_\infty+\sqrt2\,\varepsilon\|\partial w\|_\infty
\le(1+\sqrt2)\frac{2(A^{\chi}_1+A^{\chi}_2)|\omega|\varepsilon}{\mathcal{R}} .
\end{equation*}
In the same way, using $\operatorname{Lip}(\sum_\mu w_\mu)\le\|\nabla\sum_\mu w_\mu\|_\infty$ and
$\mathcal{R}\ge1$,
\begin{equation*}
\operatorname{Lip}\bar V\le\varepsilon\|\partial^3\xi\|_\infty+\sqrt2\,\varepsilon\|\partial^2w\|_\infty
\le(1+\sqrt2)\frac{(3A^{\chi}_2+2A^{\chi}_3)|\omega|\varepsilon}{\mathcal{R}} .
\end{equation*}
Bounding the maximum in $\|\cdot\|_\sharp$ by the sum of its two entries gives
$\|\bar V\|_\sharp\le C_{\bar V}\varepsilon/\mathcal{R}$, with $C_{\bar V}$ as in
\eqref{10.7:eq-divergence-weights-1}. This is at most $C_{\bar V}\varepsilon$ since $\mathcal{R}\ge1$.

Next,
\begin{equation*}
\|V_\pi\|_\infty\le\|w_\rho\|_\infty+\|w_\mu\|_\infty\le2\|\partial\xi\|_\infty
\le2(A^{\chi}_0+2A^{\chi}_1)|\omega| .
\end{equation*}
Since $V_\pi$ is a translate of $w_\rho+w_\mu$,
\begin{equation*}
\operatorname{Lip}V_\pi\le2\|\partial w\|_\infty\le2\|\partial^2\xi\|_\infty
\le\frac{4(A^{\chi}_1+A^{\chi}_2)|\omega|}{\mathcal{R}}\le2(A^{\chi}_1+A^{\chi}_2)|\omega|,
\end{equation*}
where the last step uses $\mathcal{R}\ge12\ge2$. Hence
$\|V_\pi\|_\sharp\le2(A^{\chi}_0+2A^{\chi}_1)|\omega|+80M(A^{\chi}_1+A^{\chi}_2)|\omega|$.

Both $\|\cdot\|_\infty$ and $\operatorname{Lip}$ are subadditive, so $\|\cdot\|_\sharp$ satisfies the
triangle inequality. Therefore, since $\varepsilon\le1$ and by
\eqref{10.7:eq-divergence-weights-2},
\begin{equation*}
\|V^0_\pi\|_\sharp\le\|V_\pi\|_\sharp+\|\bar V\|_\sharp\le\|V_\pi\|_\sharp+C_{\bar V}\varepsilon
\le C_{V^0} .
\end{equation*}

None of these constants involves $\mathcal{R}$ or $\varepsilon$.
\end{proof}

\begin{proposition}[The divergence term vanishes in the limit]\label{10.7:prop-divergence-vanishes}
Assume the standing inputs of Definition~\ref{10.7:def-annulus-functionals}:
\begin{itemize}
\item the hypotheses displayed in \eqref{10.7:eq-hypotheses-used-form-b} other than the infrared
one, which is the form used in this chapter of the far-sphere flux hypothesis
$(\mathrm{IR}\text{-}\mathrm{fl})_{\omega_0}$ of clause (vi);
\item the per-ball volume-free bound on $\mathbb{R}^4$, and the remaining standing input of
Definition~\ref{10.7:def-annulus-functionals};
\item the forms of $B^4_\chi$ and $J^4_\chi$ displayed in \eqref{10.7:eq-annulus-functionals-a}.
\end{itemize}
Fix $\omega$, $n$, $\mathsf{d}$ and $\vec h=(h_1,\dots,h_n)$ as there. Let $\rho_0>0$ be such that
$\bigcup_i\operatorname{supp}h_i\subset B(0,\rho_0)$, and let $\mathcal{R}\ge\max\{12,2\rho_0\}$,
so that $\vec h\in\mathcal{T}^{(2)}_n(\mathcal{R},\mathsf{d})$ for every such $\mathcal{R}$.
By Lemma~\ref{10.7:lem-plane-regrouping},
\begin{equation}\label{10.7:eq-divergence-vanishes-1}
J^4_{\chi_{\mathcal{R}}}=\kappa_N\,\mathcal{O}(\bar V)+\kappa_N g_0^2\,\mathcal{T}(\vec V^0),
\qquad \mathcal{T}(\vec g)=g_0^{-2}\sum_\pi\mathcal{O}_\pi(g_\pi).
\end{equation}
Put $C_{\mathrm{cube}}:=2^{11}\pi^2$. For a function $g$ write
$\Theta^{\varepsilon}_{\nu\sigma}(g):=\sum_x\varepsilon^4g(x)\Theta^{\varepsilon}_{\nu\sigma}(x)$.
Here $\varepsilon_j$ and $g_0^{(j)}$ are the lattice spacing and the bare coupling of the $j$-th
member of the diagonal sequence of clause (ii-b), $C_{\bar V}$ and $C_{V^0}$ are the constants of
Lemma~\ref{10.7:lem-divergence-weights}, and $\delta$, $C(\vec h)$, $C_V$ and $C_B$ are fixed in
(ii)--(iv) below.
Consider the three bounds
\begin{equation}\label{10.7:eq-divergence-vanishes-b}
\begin{gathered}
\bigl|\kappa_j(\kappa_N\mathcal{O}(\bar V);\mathcal{O}(\vec h))\bigr|
\le\kappa_N\,C^{\infty}_{n+1}\bigl(\min\{\mathsf{d},\mathcal{R}/4\};2\mathcal{R}+3\bigr)
\,C_{\bar V}\,\varepsilon_j\prod_{i=1}^n\|h_i\|_{\sharp},\\
\bigl|\kappa_j(\kappa_N g_0^2\mathcal{T}(\vec V^0);\mathcal{O}(\vec h))\bigr|
\le\kappa_N\,g_0^{(j)2}\,C_{\mathrm{cube}}\,\mathcal{R}^4\,C(\vec h)\,C_V,\\
\sup_j\bigl|\kappa_j(B^4_{\chi_{\mathcal{R}}};\mathcal{O}(\vec h))\bigr|\le C_B\,\mathcal{R}^{-\delta}.
\end{gathered}
\end{equation}
\begin{enumerate}
\item[(i)] For every $j\ge j_1(n+1,2\mathcal{R}+3)$ the first bound of
\eqref{10.7:eq-divergence-vanishes-b} holds.
Hence $\kappa_j(\kappa_N\mathcal{O}(\bar V);\mathcal{O}(\vec h))\to0$ as $j\to\infty$.
\item[(ii)] Suppose that the tensor-class boundedness bound of
\eqref{10.7:eq-annulus-functionals-a} holds with constant $C(\vec h)$. Then the second bound
of \eqref{10.7:eq-divergence-vanishes-b} holds for every $j$, with
$C_V:=(1+160M)C_{V^0}$. Hence that cumulant tends to $0$ as $j\to\infty$ at fixed $\mathcal{R}$.
If instead the tensor-class smallness bound of \eqref{10.7:eq-annulus-functionals-a} holds,
with its sequence $\eta_j\to0$, the same cumulant is at most
$\kappa_N C_{\mathrm{cube}}\mathcal{R}^4C(\vec h)C_V\eta_j$, which tends to $0$.
\item[(iii)] Let $\vec g=(g_\pi)$ be a tuple indexed by the six planes $\pi=\{\rho,\mu\}$.
Put $(P\vec g)_\mu:=\sum_{\rho\ne\mu}g_{\rho\mu}$ and $\hat g:=\tfrac12P\vec g$.
Let $c_\Theta$ be the constant for which the leading classical coefficient of
$\mathcal{T}(\vec g)$ on the tuples $\vec g=(\hat g_\rho+\hat g_\mu)_{\{\rho,\mu\}}$ with
$\sum_\mu\hat g_\mu=0$ coincides with that of $c_\Theta\sum_\mu\Theta^{\varepsilon}_{\mu\mu}(\hat g_\mu)$;
no value of $c_\Theta$ is used. Suppose that
for some constant $C'(\vec h)$
\begin{equation}\label{10.7:eq-divergence-vanishes-a}
\sup_j\Bigl|\kappa_j\Bigl(\mathcal{T}(\vec g)-c_\Theta\sum_\mu\Theta^{\varepsilon}_{\mu\mu}(\hat g_\mu);
\mathcal{O}(\vec h)\Bigr)\Bigr|\le C'(\vec h)\max_\pi\|g_\pi\|_{\sharp}
\end{equation}
for every traceless tuple $\vec g$ with the following properties: $\sum_\pi g_\pi=0$, the $g_\pi$
are real Lipschitz functions, and they are supported in a common cube of side one at
distance at least one from $\bigcup_i\operatorname{supp}h_i$.
Suppose moreover that the infrared decay bound for the stress density $\Theta^\varepsilon$ with an
exponent $\delta>0$, in the form displayed in \eqref{10.7:eq-annulus-functionals-a}, holds with
constant $C_{CL}(\vec h)$. Then the tensor-class boundedness bound of (ii) holds with
$C(\vec h):=6|c_\Theta|C_{CL}(\vec h)+C'(\vec h)$.
\item[(iv)] Assume in addition:
\begin{itemize}
\item the infrared decay bound for $\Theta^\varepsilon$ with an exponent $\delta>0$;
\item one of the two tensor-class bounds used in (ii) (for instance through (iii)).
\end{itemize}
Put
\[
b_{\nu\sigma}(x):=\sum_\rho\omega_{\rho\sigma}x_\rho\,\partial_\nu\chi_{\mathcal{R}}(x),
\qquad C_b:=\bigl(2A_1+40M(A_1+2A_2)\bigr)|\omega| .
\]
Then the third bound of \eqref{10.7:eq-divergence-vanishes-b} holds for every
$\mathcal{R}\ge\max\{12,2\rho_0\}$, with
\[
C_B:=16\cdot4^{4+\delta}\,C_{\mathrm{cube}}\,C_{CL}(\vec h)\,(1+160M)\,C_b .
\]
Consequently $c(\vec h,\omega)=0$, and by Proposition~\ref{10.7:prop-anomaly-flux} the
subsequential limit along the sequence of clause (ii-b) is $O(4)$-invariant.
\end{enumerate}
\end{proposition}

\begin{proof}
Throughout, $\kappa_j(Y;\mathcal{O}(\vec h))$ is the joint cumulant of $Y$ and
$\mathcal{O}(h_1^{(m_j)}),\dots,\mathcal{O}(h_n^{(m_j)})$. It is taken in the $j$-th member of
the diagonal sequence of clause (ii-b), whose lattice spacing is $\varepsilon_j$ and whose bare
coupling is $g_0^{(j)}$. The cumulant is linear in its first argument.
By \eqref{10.7:eq-divergence-vanishes-1}, $J^4_{\chi_{\mathcal{R}}}$ splits into a scalar part
$\kappa_N\mathcal{O}(\bar V)$ and a traceless part $\kappa_Ng_0^2\mathcal{T}(\vec V^0)$.
The first two items treat these parts in turn.

\emph{(i) The scalar part.} We apply the transfer lemma for small scalar annulus terms
(Lemma~\ref{10.7:lem-annulus-transfer}) to the term $\mathcal{O}(\bar V)$. Its hypotheses hold
for the following reasons.
\begin{itemize}
\item By \eqref{10.7:eq-divergence-vanishes-1}, this term is exactly of the form
$\mathcal{O}(w)$ with $w:=\bar V$.
\item $\bar V$ is a real Lipschitz function.
\item By Lemma~\ref{10.7:lem-divergence-weights}, $\operatorname{supp}\bar V\subset
\mathfrak{A}_{\mathcal{R}}\subset\bar B(0,2\mathcal{R}+2)$, and this support lies at distance at
least $\mathcal{R}/2-2\ge\mathcal{R}/4>0$ from $\bar B(0,\mathcal{R}/2)$.
\item By the same lemma, $\|\bar V\|_{\sharp}\le C_{\bar V}\varepsilon_j$.
\end{itemize}
So the lemma applies with $\rho_w:=2\mathcal{R}+2$, $\mathsf{d}_w:=\mathcal{R}/4$,
$C_w:=C_{\bar V}$ and $\varepsilon=\varepsilon_j$. It applies equally with
$\rho_w:=2\mathcal{R}+3$, since $\bar B(0,2\mathcal{R}+2)\subset\bar B(0,2\mathcal{R}+3)$; we use
the latter radius. Since $2\mathcal{R}+3\ge\mathcal{R}/2$, we have
$\rho_w\vee\mathcal{R}/2=2\mathcal{R}+3$.
The lemma yields, for $j\ge j_1(n+1,2\mathcal{R}+3)$ and after multiplication by $\kappa_N$, the
first line of \eqref{10.7:eq-divergence-vanishes-b}.
The input entering here is the per-ball volume-free bound, with $n+1$ test functions,
through Lemma~\ref{10.7:lem-annulus-transfer}. No infrared input is used. The only limit needed
is $\varepsilon_j\to0$, which holds because $K_j\to\infty$ along the sequence of clause (ii-b),
by clause (0).

\emph{(ii) The traceless part.} Let $\theta_0(t):=\max\{0,1-|t|\}$, and for $z\in\mathbb{Z}^4$ put
\[
\Theta_z(y):=\prod_{i=1}^4\theta_0(2y_i-z_i).
\]
These functions form a Lipschitz partition of unity, by the following four facts.
\begin{itemize}
\item Since $\sum_{k\in\mathbb{Z}}\theta_0(t-k)=1$ for all real $t$, the sum over
$z=(z_1,\dots,z_4)$ factorises and $\sum_z\Theta_z\equiv1$.
\item The factor $\theta_0(2y_i-z_i)$ vanishes unless $|y_i-z_i/2|\le\tfrac12$. Hence
$\operatorname{supp}\Theta_z\subset Q_z:=z/2+[-\tfrac12,\tfrac12]^4$, a cube of side one.
\item Each factor takes values in $[0,1]$ and is Lipschitz in $y_i$ with constant $2$. Hence
\[
|\Theta_z(y)-\Theta_z(y')|\le2\sum_i|y_i-y'_i|\le4|y-y'|,
\]
that is, $\operatorname{Lip}\Theta_z\le4$.
\item For a Lipschitz $g$ we have $\|\Theta_zg\|_\infty\le\|g\|_\infty$ and
$\operatorname{Lip}(\Theta_zg)\le4\|g\|_\infty+\operatorname{Lip}g$. Hence
$40M\operatorname{Lip}(\Theta_zg)\le(160M+1)\|g\|_{\sharp}$, and so
\begin{equation}\label{10.7:eq-divergence-vanishes-2}
\|\Theta_zg\|_{\sharp}\le(1+160M)\|g\|_{\sharp}.
\end{equation}
\end{itemize}
Since $\mathcal{T}$ is linear in the tuple and $\sum_z\Theta_z\equiv1$,
\[
\mathcal{T}(\vec V^0)=\sum_z\mathcal{T}\bigl((\Theta_zV^0_\pi)_\pi\bigr).
\]
Each tuple $(\Theta_zV^0_\pi)_\pi$ is traceless, because
$\sum_\pi\Theta_zV^0_\pi=\Theta_z\sum_\pi V^0_\pi=0$ by Lemma~\ref{10.7:lem-plane-regrouping}.
By Lemma~\ref{10.7:lem-divergence-weights}, $\operatorname{supp}V^0_\pi\subset
\mathfrak{A}_{\mathcal{R}}=\{\mathcal{R}-2\le|y|\le2\mathcal{R}+2\}$. So only those $z$ with
$Q_z\cap\mathfrak{A}_{\mathcal{R}}\ne\emptyset$ contribute.

We count these $z$. If $y\in Q_z\cap\mathfrak{A}_{\mathcal{R}}$, then $|y|\le2\mathcal{R}+2$ and
$|y-z/2|\le1$, the half-diagonal of $Q_z$. Hence the labels satisfy $|z|\le4\mathcal{R}+6$.
The cubes $z+[0,1)^4$, $z\in\mathbb{Z}^4$, are pairwise disjoint and of volume one. Every point of
$z+[0,1)^4$ lies at distance less than $2$ from $z$, so for $|z|\le4\mathcal{R}+6$ these cubes
lie in the open ball $B(0,4\mathcal{R}+8)$, of volume $(\pi^2/2)(4\mathcal{R}+8)^4$.
As $\mathcal{R}\ge2$, we have $4\mathcal{R}+8\le8\mathcal{R}$. Therefore the number of
contributing $z$ is at most
\[
\tfrac{\pi^2}{2}(8\mathcal{R})^4=2^{11}\pi^2\mathcal{R}^4=C_{\mathrm{cube}}\mathcal{R}^4 .
\]

Each contributing tuple is supported in the side-one cube $Q_z$. Every point of $Q_z$ lies within
the diameter $2$ of a point of $\mathfrak{A}_{\mathcal{R}}$, so it has modulus at least
$\mathcal{R}-4$. Since $\bigcup_i\operatorname{supp}h_i\subset B(0,\rho_0)\subset
B(0,\mathcal{R}/2)$, the cube $Q_z$ lies at distance at least $\mathcal{R}/2-4\ge2$ from the
supports of the $h_i$, because $\mathcal{R}\ge12$. By \eqref{10.7:eq-divergence-vanishes-2} and
the bound $\|V^0_\pi\|_{\sharp}\le C_{V^0}$ of Lemma~\ref{10.7:lem-divergence-weights}, each
component satisfies
\[
\|\Theta_zV^0_\pi\|_{\sharp}\le(1+160M)C_{V^0}=C_V .
\]
Under the tensor-class boundedness bound, applied to each contributing tuple (distance
$\ge2\ge1$), every term satisfies
\[
\bigl|\kappa_j\bigl(\mathcal{T}((\Theta_zV^0_\pi)_\pi);\mathcal{O}(\vec h)\bigr)\bigr|
\le C(\vec h)C_V .
\]
Multiplying by $\kappa_Ng_0^{(j)2}$ and summing over at most $C_{\mathrm{cube}}\mathcal{R}^4$
labels gives the second line of \eqref{10.7:eq-divergence-vanishes-b}. Its right side tends to
$0$ as $j\to\infty$ at fixed $\mathcal{R}$, because $g_0^{(j)}\downarrow0$ by the definition of
the sequence in clause (ii-b).
Under the tensor-class smallness bound, the same count gives
$\kappa_NC_{\mathrm{cube}}\mathcal{R}^4C(\vec h)C_V\eta_j\to0$.
No decay hypothesis is used: the distance of the cubes from the supports of the $h_i$ enters
only through the requirement that it be at least one.

\emph{(iii) Reduction of the tensor-class bound to the infrared decay bound for
$\Theta^\varepsilon$.}
The infrared decay bound is formulated for the stress density $\Theta^\varepsilon$ and does
not literally cover $J^4$. To see why, split the traceless plane tuples. Write $\mathbb{R}^4_0$
for the $4$-tuples $\hat g$ with $\sum_\mu\hat g_\mu=0$, and put
$(E\hat g)_{\{\rho,\mu\}}:=\hat g_\rho+\hat g_\mu$. The tuple $E\hat g$ is traceless, since
$\sum_\pi(E\hat g)_\pi=3\sum_\mu\hat g_\mu=0$. Moreover
\[
(PE\hat g)_\mu=\sum_{\rho\ne\mu}(\hat g_\rho+\hat g_\mu)=2\hat g_\mu+\sum_\rho\hat g_\rho=2\hat g_\mu .
\]
For traceless $\vec g$ we also have $\sum_\mu(P\vec g)_\mu=2\sum_\pi g_\pi=0$. Hence every traceless
$\vec g$ decomposes as
\[
\vec g=E\hat g+(\vec g-E\hat g),\qquad \hat g=\tfrac12P\vec g,\quad \vec g-E\hat g\in\ker P .
\]
This splitting is $W_4$-equivariant. The first summand, the \emph{$3$-channel}
$E(\mathbb{R}^4_0)$, corresponds at the level of the classical densities to
$\sum_\mu\hat g_\mu\Theta_{\mu\mu}$. The second, the two-dimensional \emph{$2$-channel}
$\ker P$, corresponds classically to densities such as
$(F_{12}^2+F_{34}^2)-(F_{13}^2+F_{24}^2)$, which have no counterpart among the components of
$\Theta$. The $3$-channel of $\mathcal{T}$ agrees with
$c_\Theta\sum_\mu\Theta^\varepsilon_{\mu\mu}(\hat g_\mu)$ only in its leading classical
coefficient. The exact additional bound needed, uniform in the cutoff, is therefore
\eqref{10.7:eq-divergence-vanishes-a}.

Now let $\vec g$ be as in (iii). By linearity of the cumulant,
\begin{align*}
\kappa_j(\mathcal{T}(\vec g);\mathcal{O}(\vec h))
&=\kappa_j\Bigl(\mathcal{T}(\vec g)-c_\Theta\sum_\mu\Theta^\varepsilon_{\mu\mu}(\hat g_\mu);
\mathcal{O}(\vec h)\Bigr)\\
&\quad+c_\Theta\sum_\mu\kappa_j\bigl(\Theta^\varepsilon_{\mu\mu}(\hat g_\mu);\mathcal{O}(\vec h)\bigr).
\end{align*}
The first term is bounded by $C'(\vec h)\max_\pi\|g_\pi\|_{\sharp}$ by
\eqref{10.7:eq-divergence-vanishes-a}. For the second, note that
$\hat g_\mu=\tfrac12\sum_{\rho\ne\mu}g_{\rho\mu}$ is a sum of three terms and that
$\|\cdot\|_{\sharp}$ is subadditive. Hence
\[
\|\hat g_\mu\|_{\sharp}\le\tfrac32\max_\pi\|g_\pi\|_{\sharp}.
\]
Also $\hat g_\mu$ is supported in the same side-one cube, at a distance $D\ge1$ from the
supports of the $h_i$, so $D^{-4-\delta}\le1$. The infrared decay bound therefore estimates each
of the four cumulants by
\[
C_{CL}(\vec h)\cdot\tfrac32\max_\pi\|g_\pi\|_{\sharp}.
\]
Summing over the four $\mu$ gives
\[
\sup_j|\kappa_j(\mathcal{T}(\vec g);\mathcal{O}(\vec h))|
\le\bigl(6|c_\Theta|C_{CL}(\vec h)+C'(\vec h)\bigr)\max_\pi\|g_\pi\|_{\sharp}.
\]
This is the tensor-class boundedness bound with the constant stated in (iii).

For $J^4$ itself less is needed: boundedness per cube in the $3$-channel, and any growth
$o(\varepsilon_j^{-1})$ in the $2$-channel. Indeed, the $2$-channel component of $\vec V^0$ has
$\|\cdot\|_{\sharp}=O(\varepsilon_j/\mathcal{R})$ by Lemma~\ref{10.7:lem-divergence-weights}.
The reason is that the unshifted tuple $((w_\rho+w_\mu))_\pi$ lies in the image of $E$: its
traceless part is exactly $E\hat g$ with $\hat g_\mu=w_\mu-\tfrac14\sum_\rho w_\rho$.
A bound of the type of the correlation theorem for the scalar observable $\mathcal{O}$, applied
plane by plane to the restricted weights, would give only boundedness of the $3$-channel piece,
not its vanishing. The tensor-class boundedness bound asserts the genuine cancellation of order
$g_0^2$ in the traceless channel.

\emph{(iv) Conclusion.} We first locate and bound the weights $b_{\nu\sigma}$.
\begin{itemize}
\item The weight $b_{\nu\sigma}$ carries the factor $\partial_\nu\chi_{\mathcal{R}}$, so it is
supported in $\{\mathcal{R}\le|x|\le2\mathcal{R}\}$.
\item There, $|\sum_\rho\omega_{\rho\sigma}x_\rho|\le2\mathcal{R}|\omega|$, and
$|\partial_\nu\chi_{\mathcal{R}}|\le A_1/\mathcal{R}$. Hence $\|b_{\nu\sigma}\|_\infty\le2A_1|\omega|$.
\item By the product rule, the gradient of $b_{\nu\sigma}$ consists of a term bounded by
$|\omega|A_1/\mathcal{R}$ and a term bounded by $2\mathcal{R}|\omega|A_2\mathcal{R}^{-2}$. Hence
\[
\operatorname{Lip}b_{\nu\sigma}\le(A_1+2A_2)|\omega|/\mathcal{R}\le(A_1+2A_2)|\omega| .
\]
\item Combining the last two items, $\|b_{\nu\sigma}\|_{\sharp}\le C_b$, uniformly in
$\mathcal{R}\ge1$.
\end{itemize}
By the form of $B^4_\chi$ in \eqref{10.7:eq-annulus-functionals-a},
$B^4_{\chi_{\mathcal{R}}}=\sum_{\nu,\sigma}\Theta^\varepsilon_{\nu\sigma}(b_{\nu\sigma})$.
Inserting $\sum_z\Theta_z\equiv1$ gives
\[
B^4_{\chi_{\mathcal{R}}}=\sum_z\sum_{\nu,\sigma}\Theta^\varepsilon_{\nu\sigma}(\Theta_zb_{\nu\sigma}).
\]
Only labels $z$ with $Q_z$ meeting $\{\mathcal{R}\le|x|\le2\mathcal{R}\}\subset
\mathfrak{A}_{\mathcal{R}}$ contribute, and by the count in (ii) there are at most
$C_{\mathrm{cube}}\mathcal{R}^4$ of them. Each such $Q_z$ consists of points of modulus at least
$\mathcal{R}-2$. So it lies at distance at least $\mathcal{R}/2-2\ge\mathcal{R}/4\ge3$ from
$B(0,\mathcal{R}/2)\supset\bigcup_i\operatorname{supp}h_i$, because $\mathcal{R}\ge12$.
The infrared decay bound applies to each of the $16$ pairs $(\nu,\sigma)$ and each contributing
$z$, with $D\ge\mathcal{R}/4$ and, by \eqref{10.7:eq-divergence-vanishes-2},
$\|\Theta_zb_{\nu\sigma}\|_{\sharp}\le(1+160M)C_b$. It gives
\[
\sup_j|\kappa_j(B^4_{\chi_{\mathcal{R}}};\mathcal{O}(\vec h))|
\le16\,C_{\mathrm{cube}}\mathcal{R}^4C_{CL}(\vec h)(1+160M)C_b(\mathcal{R}/4)^{-4-\delta} .
\]
The right side equals $C_B\mathcal{R}^{-\delta}$, which is the third line of
\eqref{10.7:eq-divergence-vanishes-b}.

Since the total annulus functional is $\mathfrak{F}_{\chi_{\mathcal{R}}}=
B^4_{\chi_{\mathcal{R}}}+J^4_{\chi_{\mathcal{R}}}$, and since
$\mathcal{B}^{\chi_{\mathcal{R}}}_j(\vec h)=\kappa_j(\mathfrak{F}_{\chi_{\mathcal{R}}};\mathcal{O}(\vec h))$
is the cumulant of this functional with the sources, as in
Proposition~\ref{10.7:prop-anomaly-flux}, linearity of the cumulant gives
\[
\kappa_j(B^4_{\chi_{\mathcal{R}}};\mathcal{O}(\vec h))
=\mathcal{B}^{\chi_{\mathcal{R}}}_j(\vec h)-\kappa_j(J^4_{\chi_{\mathcal{R}}};\mathcal{O}(\vec h)).
\]
As $j\to\infty$ we have $\mathcal{B}^{\chi_{\mathcal{R}}}_j(\vec h)\to c(\vec h,\omega)$ by
Proposition~\ref{10.7:prop-anomaly-flux}, in the form \eqref{10.7:eq-anomaly-flux-a}. Also
$\kappa_j(J^4_{\chi_{\mathcal{R}}};\mathcal{O}(\vec h))\to0$ by
\eqref{10.7:eq-divergence-vanishes-1} together with (i) and (ii). Hence
$\kappa_j(B^4_{\chi_{\mathcal{R}}};\mathcal{O}(\vec h))\to c(\vec h,\omega)$, and therefore
\[
|c(\vec h,\omega)|\le C_B\mathcal{R}^{-\delta}\qquad\text{for every }
\mathcal{R}\ge\max\{12,2\rho_0\}.
\]
The anomaly $c(\vec h,\omega)$ does not depend on $\mathcal{R}$, and $C_B$ does not depend on
$\mathcal{R}$. Letting $\mathcal{R}\to\infty$ gives $c(\vec h,\omega)=0$. The data
$\omega,n,\mathsf{d},\vec h$ were arbitrary. Together with Proposition~\ref{10.7:prop-anomaly-flux}
(see \eqref{10.7:eq-anomaly-flux-a}), this gives the $O(4)$-invariance of this subsequential limit.
\end{proof}

\begin{remark}[Non-radial cut-offs leave a scalar limit term]\label{10.7:rem-non-radial}
Let $\chi$ be an arbitrary admissible cut-off, let $\xi^{\chi}$ be the deformation field attached
to it, and let $\bar V$ be formed from $\xi^{\chi}$ as in Lemma~\ref{10.7:lem-plane-regrouping}.
Let $\vec h=(h_1,\dots,h_n)$ be the test functions of the remaining $n$ sources. Then:
\begin{itemize}
\item the regrouping identity of Lemma~\ref{10.7:lem-plane-regrouping} holds for $\chi$;
\item $\bar V\to\tfrac12\operatorname{div}\xi^{\chi}$ in $\|\cdot\|_{\sharp}$ as $j\to\infty$;
\item the scalar part of the divergence term has the limit
\begin{equation}\label{10.7:eq-non-radial-1}
\lim_{j\to\infty}\kappa_j\bigl(\kappa_N\mathcal{O}(\bar V);\mathcal{O}(\vec h)\bigr)
=\kappa_N S^T_{n+1}\bigl(\tfrac12\operatorname{div}\xi^{\chi},\vec h\bigr).
\end{equation}
\end{itemize}
The limit in \eqref{10.7:eq-non-radial-1} is an $(n+1)$-point connected Schwinger function of
smeared curvature observables, which is the kind of quantity in which the main theorem is stated.
Hence $\lim_j\kappa_j(B^4_{\chi};\cdot)$ depends on $\chi$, whereas the total flux does not.
For the radial family $\chi_{\mathcal{R}}$ one has $\operatorname{div}\xi=0$, so this piece is
absent exactly.
\end{remark}

\begin{proof}
Since $\varphi\mapsto\mathcal{O}(\varphi)$ is linear and the joint cumulant is multilinear, the
difference between the left member of \eqref{10.7:eq-non-radial-1} before the limit and its right
member splits into two terms:
\begin{align*}
&\kappa_j\bigl(\kappa_N\mathcal{O}(\bar V);\mathcal{O}(\vec h)\bigr)
-\kappa_N S^T_{n+1}\bigl(\tfrac12\operatorname{div}\xi^{\chi},\vec h\bigr)\\
&\qquad=\kappa_N\,\kappa_j\bigl(\mathcal{O}(\bar V-\tfrac12\operatorname{div}\xi^{\chi});
\mathcal{O}(\vec h)\bigr)\\
&\qquad\quad+\kappa_N\Bigl[\kappa_j\bigl(\mathcal{O}(\tfrac12\operatorname{div}\xi^{\chi});
\mathcal{O}(\vec h)\bigr)
-S^T_{n+1}\bigl(\tfrac12\operatorname{div}\xi^{\chi},\vec h\bigr)\Bigr].
\end{align*}
The first term is controlled by the per-ball volume-free bound, applied with $n+1$ sources. It is
bounded by a constant times
$\|\bar V-\tfrac12\operatorname{div}\xi^{\chi}\|_{\sharp}\prod_i\|h_i\|_{\sharp}$, and this
tends to zero because $\bar V\to\tfrac12\operatorname{div}\xi^{\chi}$ in $\|\cdot\|_{\sharp}$.
In the second term the $(n+1)$-tuple $(\tfrac12\operatorname{div}\xi^{\chi},\vec h)$ does not
depend on $j$ and is a fixed tuple of class $C^1_c$. By the infinite-volume theorem, the cumulant
$\kappa_j$ of this tuple converges to $S^T_{n+1}$, so the bracket tends to zero.
This proves \eqref{10.7:eq-non-radial-1}.

For the radial family, $\operatorname{div}\xi=0$. The right member of
\eqref{10.7:eq-non-radial-1} is then $S^T_{n+1}$ evaluated with the zero function in its first
slot, and this vanishes by multilinearity.
\end{proof}

\begin{remark}[Minimal and sufficient infrared hypotheses]\label{10.7:rem-infrared-forms}
Throughout, the standing inputs are those of Definition~\ref{10.7:def-annulus-functionals}.
\begin{enumerate}
\item[(1)] The minimal infrared hypothesis, written
$(\mathrm{H}\text{-}\mathrm{IR})_{\mathfrak{F}}$, is the following: for each $n\ge 2$, each
$\omega$, each $\mathcal{R}$ and $\mathsf{d}$, and each tuple $\vec h$ of the class to which
Corollary~\ref{10.7:cor-clustering-dichotomy} applies for these data,
\begin{equation}\label{10.7:eq-infrared-forms-1}
\lim_{j\to\infty}\mathcal{B}^{\omega,\chi}_j(\vec h)=0
\quad\text{for one admissible (radial) cut-off }\chi .
\end{equation}
Given the standing inputs, $(\mathrm{H}\text{-}\mathrm{IR})_{\mathfrak{F}}$ is equivalent to the
$O(4)$ invariance of the subsequential limit, that is, to item~(v) of
Corollary~\ref{10.7:cor-clustering-dichotomy}.
\item[(2)] A sufficient form of the infrared hypothesis is the conjunction of the following two
conditions:
\begin{enumerate}
\item[(a)] decay of the connected correlation of one far-away insertion of
$\Theta^{\varepsilon}_{\nu\sigma}$ with the observables, as in
\eqref{10.7:eq-annulus-functionals-a}, with decay exponent $-4-\delta$ in the sense of that
display, for some $\delta>0$;
\item[(b)] the boundedness condition among the tensor-class bounds in
\eqref{10.7:eq-annulus-functionals-a};
\end{enumerate}
or, in place of the boundedness condition, the divergence bound
\eqref{10.7:eq-divergence-vanishes-a}. The condition is $\delta>0$ strictly; the borderline decay
with $\delta=0$ is not asserted to suffice. Each of these is a condition on a tensor-valued observable.
None of them lies in the class of smeared curvature Schwinger functions in which Theorem~0 is
formulated. This form enters through the expression of the rotational anomaly as a boundary flux
of the annulus functional, \eqref{10.7:eq-anomaly-flux-a}, and through
Proposition~\ref{10.7:prop-divergence-vanishes}.
\item[(3)] The hypercubically scalar plaquette density is never written with the canonical factor
$g_0^{-2}$, that is, as $g_0^{-2}\mathcal{O}(h)$. Only traceless combinations, such as the
traceless plane functional $\mathcal{T}(\vec g)=g_0^{-2}\sum_\pi\mathcal{O}_\pi(g_\pi)$, carry
the factor $g_0^{-2}$.
\item[(4)] If $J^4\equiv 0$, then $\mathfrak{F}=B^4$. In that case the divergence-term analysis, up to
and including part~(iii) of Proposition~\ref{10.7:prop-divergence-vanishes}, is void. The boundedness
condition of item~(2) then drops out of the sufficient form.
\end{enumerate}
\end{remark}

\begin{proof}
(1) Fix $\omega$, $n\ge 2$, $\mathcal{R}$ and $\mathsf{d}$. By
Corollary~\ref{10.7:cor-clustering-dichotomy}, the following two statements are
equivalent under the standing inputs:
\begin{itemize}
\item item~(iii) of that corollary, which states that
$\lim_j\mathcal{B}^{\omega,\chi}_j(\vec h)=0$ for one, equivalently every, admissible $\chi$;
\item item~(i) of that corollary, for these data.
\end{itemize}
Taken over all $\omega,n,\mathcal{R},\mathsf{d}$, the same corollary makes these items equivalent to
its item~(v), the $O(4)$ invariance of the subsequential limit. Since item~(iii) holds for one
admissible $\chi$ if and only if it holds for every admissible $\chi$, and the radial cut-offs
$\chi_{\mathcal{R}}$ of Definition~\ref{10.7:def-annulus-functionals} are admissible,
item~(iii) taken over all these data is equivalent to
\eqref{10.7:eq-infrared-forms-1}.

(2) That the conditions of item~(2), namely (a) and (b), or (a) and the divergence bound
\eqref{10.7:eq-divergence-vanishes-a}, imply \eqref{10.7:eq-infrared-forms-1} follows from the
expression of the rotational anomaly as a boundary flux, \eqref{10.7:eq-anomaly-flux-a}, together
with Proposition~\ref{10.7:prop-divergence-vanishes}, in which these conditions are used; that
argument is not repeated here.

(3) Write $g_0^{(j)}$ for the bare coupling at the $j$-th point of the diagonal.
The joint cumulant of the plaquette observables is a connected Schwinger function:
\begin{equation}\label{10.7:eq-infrared-forms-2}
\kappa_j\bigl(\mathcal{O}(h);\mathcal{O}(\vec h)\bigr)=S^T_{n+1;K_j,m_j}(h,\vec h).
\end{equation}
By the per-ball volume-free bound and the standing input (b2) of
Definition~\ref{10.7:def-annulus-functionals}, which supplies the convergence along the diagonal,
the right-hand side of
\eqref{10.7:eq-infrared-forms-2} is $O(1)$ and convergent as $j\to\infty$. Clause~(0) of Theorem~0
gives the two-sided bound $K\asymp g_0^{-2}$, and hence $\bigl(g_0^{(j)}\bigr)^{-2}\asymp K_j$.
Consequently
\begin{equation*}
\bigl(g_0^{(j)}\bigr)^{-2}\,S^T_{n+1;K_j,m_j}(h,\vec h)=O(K_j).
\end{equation*}
Along the diagonal of clause~(ii-b), clause~(0) gives $K_j\to\infty$. A scalar density written as
$g_0^{-2}\mathcal{O}(h)$ would therefore carry a prefactor of order $K_j$, which tends to infinity,
while the
Schwinger function itself stays bounded. This is why, in item~(3), the factor $g_0^{-2}$ is reserved
for traceless combinations.

(4) By definition $\mathfrak{F}=B^4+J^4$, so $J^4\equiv 0$ gives $\mathfrak{F}=B^4$. The
divergence-term analysis leading to Proposition~\ref{10.7:prop-divergence-vanishes}, up to and
including its part~(iii), concerns the functional $J^4$. With $J^4\equiv 0$ there is nothing for it
to estimate, and the boundedness condition of item~(2) is no longer required.
\end{proof}

\begin{definition}[Coordinate reflections and separated slabs]\label{10.7:not-reflected-slabs}
Fix a cutoff $(K,m)$, so that the configurations are bond variables on the lattice of spacing
$\varepsilon$ in the torus $\mathbb{T}_m$. Fix a coordinate direction $\nu\in\{1,2,3,4\}$, with
unit vector $e_\nu$, and a number $c\in\mathbb{R}$ with $2c\in\varepsilon\mathbb{Z}$. Let
$r:=r_{\nu,c}$ be the reflection
\begin{equation}\label{10.7:eq-reflected-slabs-1}
  r_{\nu,c}\colon\ x_\nu\longmapsto 2c-x_\nu ,
\end{equation}
all coordinates other than $x_\nu$ being left unchanged. Because $2c\in\varepsilon\mathbb{Z}$, the
map $r_{\nu,c}$ is an isometry of the lattice. The half-spaces $H^{\pm}_{\nu,c}$ are those of the
notation of the lemma on reflection positivity of the lattice measure, and for a function $Y$ of
the configuration, $Y^{r}$ denotes, as there, the function of the configuration transformed by
the isometry $r$.

Let $0<\ell\le L^m$, let $u\ge0$ with $2u\in\varepsilon\mathbb{Z}$, and let $\ell'>0$ with
$u+\ell'\le L^m$. The two slabs are
\begin{equation}\label{10.7:eq-reflected-slabs-2}
  \Sigma^{+}:=\{\,c<x_\nu<c+\ell\,\},\qquad
  \Sigma^{-}:=\{\,c-u-\ell'<x_\nu<c-u\,\}.
\end{equation}
Write $\tau_{a}$ for the translation $x\mapsto x+a$. The reflection $r_{\nu,c-u}$ sends $x_\nu$
to $2(c-u)-x_\nu$, and the translation $\tau_{2ue_\nu}$ then adds $2u$ to the $\nu$-th
coordinate, the result being $2c-x_\nu$; neither map changes the other coordinates. Hence
\begin{equation}\label{10.7:eq-reflected-slabs-3}
  r_{\nu,c}=\tau_{2ue_\nu}\circ r_{\nu,c-u}.
\end{equation}
Since $2u\in\varepsilon\mathbb{Z}$ and $2(c-u)\in\varepsilon\mathbb{Z}$, both factors on the
right are lattice isometries. Consequently, for $Y$ measurable on $\Sigma^{-}$, the function
$Y^{r}$ is the mirror image of $Y$ in the plane $x_\nu=c-u$, translated by $2ue_\nu$. We call
$2u$ the \emph{gap}.
\end{definition}

\begin{proposition}[Reflection-positivity bound for slab correlations]
\label{10.7:prop-slab-bound}
Fix a cutoff, and let $E$ denote expectation in the lattice measure at that cutoff. Let
$\Sigma^{\pm}$, $r_{\nu,c}$ and $H^{\pm}_{\nu,c}$ be as in
Definition~\ref{10.7:not-reflected-slabs}. For a function $F$ of the configuration write $F^{r}$ for
$F$ composed with the action of $r_{\nu,c}$ on configurations. Let $X$ be a bounded real function
that is measurable with respect to the bond variables of the bonds with both endpoints in
$\Sigma^{+}$, and let $Y$ be a bounded real function measurable in the same sense on
$\Sigma^{-}$. Then
\begin{equation}\label{10.7:eq-slab-bound-1}
\bigl|E[XY]\bigr| \le E\bigl[X\,X^{r}\bigr]^{1/2}\, E\bigl[Y\,Y^{r}\bigr]^{1/2},
\end{equation}
and both factors on the right are non-negative. If $X$ and $Y$ are centred, then all three
expectations in \eqref{10.7:eq-slab-bound-1} are covariances.
\end{proposition}

\begin{proof}
We first locate the two functions. The function $X$ is measurable on
$\Sigma^{+}\subset H^{+}_{\nu,c}$. The function $Y^{r}$ is measurable on the reflected slab, which is
\begin{equation*}
r_{\nu,c}(\Sigma^{-}) = \{\, c+u < x_\nu < c+u+\ell' \,\} \subset H^{+}_{\nu,c},
\end{equation*}
where $u$ and $\ell'$ are the offset and the width of $\Sigma^{-}$.

On bounded real $H^{+}_{\nu,c}$-measurable functions define
\begin{equation*}
b(F,G) := E\bigl[F^{r}G\bigr].
\end{equation*}
Since $(F+tG)^{r}=F^{r}+tG^{r}$ and expectation is linear, $b$ is bilinear.

We show that $b$ is symmetric. The lattice measure is invariant under the reflection
$r_{\nu,c}$, by the invariance lemma for the Wilson measure. Moreover $r_{\nu,c}^{2}$ is the
identity, so $(F^{r})^{r}=F$. Hence
\begin{equation*}
b(F,G) = E\bigl[(F^{r}G)^{r}\bigr] = E\bigl[FG^{r}\bigr] = b(G,F).
\end{equation*}

We show that $b$ is positive semidefinite. Let $\Theta_{\nu,c}$ denote the reflection acting on
functions. For real $F$ we have $\Theta_{\nu,c}F = F^{r}$, and therefore
\begin{equation*}
b(F,F) = E\bigl[(\Theta_{\nu,c}F)F\bigr] \ge 0
\end{equation*}
by reflection positivity of the lattice measure.

Cauchy--Schwarz for $b$ follows. For every real $t$, bilinearity and symmetry give
\begin{equation*}
0 \le b(F+tG,F+tG) = b(F,F) + 2t\,b(F,G) + t^{2}\,b(G,G).
\end{equation*}
A non-negative real quadratic polynomial in $t$ has non-positive discriminant. If $b(G,G)=0$, the
linear polynomial $b(F,F)+2t\,b(F,G)$ is non-negative for every $t$, which forces $b(F,G)=0$. In
either case
\begin{equation}\label{10.7:eq-slab-bound-2}
|b(F,G)| \le b(F,F)^{1/2}\, b(G,G)^{1/2}.
\end{equation}

We now apply \eqref{10.7:eq-slab-bound-2} with $F=Y^{r}$ and $G=X$; both are
$H^{+}_{\nu,c}$-measurable by the first paragraph. Using $(Y^{r})^{r}=Y$,
\begin{equation*}
E[XY] = E\bigl[(Y^{r})^{r}X\bigr] = b(Y^{r},X).
\end{equation*}
The two diagonal terms are
\begin{equation*}
b(Y^{r},Y^{r}) = E\bigl[Y\,Y^{r}\bigr],
\qquad
b(X,X) = E\bigl[X^{r}X\bigr] = E\bigl[X\,X^{r}\bigr].
\end{equation*}
Inserting these into \eqref{10.7:eq-slab-bound-2} gives \eqref{10.7:eq-slab-bound-1}. Both factors
on the right have the form $b(F,F)$, so both are non-negative.

For the centred case, suppose $E[X]=E[Y]=0$. By the invariance of the measure under
$r_{\nu,c}$, also $E[X^{r}]=E[X]=0$ and $E[Y^{r}]=E[Y]=0$. Thus each of $E[XY]$, $E[X\,X^{r}]$ and
$E[Y\,Y^{r}]$ is the expectation of a product of two centred variables, that is, a covariance.

The same argument holds with $H^{-}_{\nu,c}$ in place of $H^{+}_{\nu,c}$. For
$H^{-}_{\nu,c}$-measurable $F$, the function $G:=\Theta_{\nu,c}F$ is
$H^{+}_{\nu,c}$-measurable. Since $\Theta_{\nu,c}^{2}$ is the identity,
\begin{equation*}
E\bigl[(\Theta_{\nu,c}F)F\bigr] = E\bigl[(\Theta_{\nu,c}G)G\bigr] \ge 0,
\end{equation*}
so the positivity of $b$, and with it \eqref{10.7:eq-slab-bound-2}, carries over.
\end{proof}

\begin{lemma}[Cumulants as covariances with a cluster polynomial]\label{10.7:lem-cluster-polynomial}
Let $n\ge 1$, and let $X,\eta_1,\dots,\eta_n$ be bounded real random variables. Put
$\tilde\eta_i:=\eta_i-E[\eta_i]$. In the truncated algebra
$\mathbb{A}_n=\mathbb{C}[z]/(z_1^2,\dots,z_n^2)$ put
\begin{equation*}
M_0:=E\Bigl[\prod_{i=1}^n(1+z_i\tilde\eta_i)\Bigr]\in\mathbb{A}_n,
\end{equation*}
the expectation being taken coefficient by coefficient, and define the random variable
\begin{equation}\label{10.7:eq-cluster-polynomial-1}
Y:=[z_1\cdots z_n]\Bigl(M_0^{-1}\prod_{i=1}^n(1+z_i\tilde\eta_i)\Bigr),
\end{equation}
where $[z_1\cdots z_n]$ denotes the coefficient of the monomial $z_1\cdots z_n$. Then $M_0$ is
invertible in $\mathbb{A}_n$, and $Y$ is a real multilinear polynomial in
$\tilde\eta_1,\dots,\tilde\eta_n$ whose coefficients are polynomials in the centred moments
$E[\prod_{i\in T}\tilde\eta_i]$, $T\subset\{1,\dots,n\}$. Moreover
\begin{equation}\label{10.7:eq-cluster-polynomial-2}
\kappa(X,\eta_1,\dots,\eta_n)=E[XY],\qquad E[Y]=0,
\end{equation}
and $Y$ is measurable with respect to $\sigma(\eta_1,\dots,\eta_n)$. In particular, for
$\eta_i=\mathcal{O}(h_i^{(m)})$, $Y$ is a function of the bond variables on the bonds of the
plaquettes whose barycentres lie in $\bigcup_i\operatorname{supp}h_i^{(m)}$.
\end{lemma}

\begin{proof}
Write $\mathfrak{m}\subset\mathbb{A}_n$ for the ideal spanned by the monomials of degree at least
one; since a nonzero monomial of $\mathbb{A}_n$ is square-free, $\mathfrak{m}^{n+1}=0$. Expanding
the product,
\begin{equation*}
M_0=\sum_{T\subset\{1,\dots,n\}}z^T\,E\Bigl[\prod_{i\in T}\tilde\eta_i\Bigr],
\qquad z^T:=\prod_{i\in T}z_i .
\end{equation*}
The term $T=\emptyset$ is $1$, so $M_0-1\in\mathfrak{m}$, and
\begin{equation*}
M_0\sum_{j=0}^{n}(1-M_0)^j=1-(1-M_0)^{n+1}=1;
\end{equation*}
hence $M_0$ is invertible and
$M_0^{-1}=\sum_{j=0}^{n}(1-M_0)^j$. Its coefficients are therefore real polynomials in the
centred moments; writing $M_0^{-1}=\sum_S c_S z^S$, the definition
\eqref{10.7:eq-cluster-polynomial-1} gives
\begin{equation*}
Y=\sum_{T\subset\{1,\dots,n\}}c_{\{1,\dots,n\}\setminus T}\prod_{i\in T}\tilde\eta_i ,
\end{equation*}
which is the asserted multilinear form.

By Lemma~\ref{10.7:lem-joint-cumulants}(c) the cumulant of $N=n+1\ge 2$ variables is
multilinear and symmetric and vanishes when one argument is constant; writing
$\eta_i=\tilde\eta_i+E[\eta_i]$ and expanding one argument at a time, every term containing a
constant argument vanishes, so
\begin{equation*}
\kappa(X,\eta_1,\dots,\eta_n)=\kappa(X,\tilde\eta_1,\dots,\tilde\eta_n).
\end{equation*}

Apply Lemma~\ref{10.7:lem-joint-cumulants}(b) with $N=n+1$ to the variables
$X,\tilde\eta_1,\dots,\tilde\eta_n$, writing the generators of $\mathbb{A}_{n+1}$ as
$w,z_1,\dots,z_n$, so that $w^2=0$. Set $P_X:=E[X\prod_i(1+z_i\tilde\eta_i)]\in\mathbb{A}_n$.
Then
\begin{equation*}
M:=E\Bigl[(1+wX)\prod_{i=1}^n(1+z_i\tilde\eta_i)\Bigr]=M_0+wP_X .
\end{equation*}
Since the algebra is commutative and $w^2=0$, for $k\ge1$
\begin{equation*}
(M-1)^k=(M_0-1)^k+k\,w\,(M_0-1)^{k-1}P_X .
\end{equation*}
Summing with the weights $(-1)^{k-1}/k$ over $k\le n+1$,
\begin{equation*}
\log M=\sum_{k=1}^{n+1}\frac{(-1)^{k-1}}{k}(M_0-1)^k
+w\sum_{j=0}^{n}(1-M_0)^j\,P_X=\log M_0+w\,M_0^{-1}P_X ,
\end{equation*}
where the first sum involves no $w$. The term $k=n+1$ of the first sum vanishes because
$(M_0-1)^{n+1}\in\mathfrak{m}^{n+1}=0$, so the first sum is $\log M_0$; in the second sum the
factor $k$ cancels $1/k$, and with $j=k-1$ one has $(-1)^{k-1}(M_0-1)^{k-1}=(1-M_0)^j$, so that
the second sum is $M_0^{-1}$ by the formula for $M_0^{-1}$ above.
Hence, by Lemma~\ref{10.7:lem-joint-cumulants}(b),
\begin{align*}
\kappa(X,\tilde\eta_1,\dots,\tilde\eta_n)
&=[wz_1\cdots z_n]\bigl(\log M_0+w\,M_0^{-1}P_X\bigr)
=[z_1\cdots z_n]\bigl(M_0^{-1}P_X\bigr)\\
&=E\Bigl[X\,[z_1\cdots z_n]\Bigl(M_0^{-1}\prod_{i=1}^n(1+z_i\tilde\eta_i)\Bigr)\Bigr]
=E[XY].
\end{align*}
In the third equality the expectation was moved outside: $M_0^{-1}$ has deterministic
coefficients, the coefficient of $z_1\cdots z_n$ is a finite linear combination of products of
bounded variables, and $E$ is linear. This proves the first identity in
\eqref{10.7:eq-cluster-polynomial-2}.

Taking for $X$ the constant $1$ gives $E[Y]=\kappa(1,\eta_1,\dots,\eta_n)$, which vanishes by
Lemma~\ref{10.7:lem-joint-cumulants}(c), since $n+1\ge2$.

Finally, $Y$ is a polynomial in $\tilde\eta_i=\eta_i-E[\eta_i]$ with deterministic
coefficients, hence $\sigma(\eta_1,\dots,\eta_n)$-measurable. If
$\eta_i=\mathcal{O}(h_i^{(m)})$, that is,
\begin{equation*}
\eta_i=\sum_p h_i^{(m)}(x(p))\bigl[1-\operatorname{Re}\operatorname{tr}U(\partial p)\bigr],
\end{equation*}
only plaquettes $p$ with $x(p)\in\operatorname{supp}h_i^{(m)}$ contribute, and each such
term depends only on the bond variables of $\partial p$; so $Y$ depends only on the bonds of
plaquettes with barycentres in $\bigcup_i\operatorname{supp}h_i^{(m)}$.
\end{proof}

\begin{lemma}[Slab geometry for far stress-tensor cubes]\label{10.7:lem-far-geometry}
Let $\vec h\in\mathcal{T}^{(2)}_n(\mathcal{R},\mathsf{d})$, and let $\varepsilon$ be the lattice spacing. For a fixed pair
of indices let $\Theta^{\varepsilon}_{\nu\sigma}(x)$ be a local gauge-invariant lattice operator of fixed range
$r_\Theta$ (in lattice units), the same at every site, and for a function $g$ put
\begin{equation}\label{10.7:eq-far-geometry-a}
\Theta^{\varepsilon}_{\nu\sigma}(g):=\sum_x \varepsilon^4\, g(x)\,\Theta^{\varepsilon}_{\nu\sigma}(x).
\end{equation}
With the pair of indices fixed we write $\Theta^{\varepsilon}(g)$ for \eqref{10.7:eq-far-geometry-a}; from here on
the letter $\nu$ denotes a coordinate direction. Let $\Delta$ be a closed cube of side $2$ with centre $y$, where
$|y|\ge \mathcal{R}'/2$ and $\mathcal{R}'\ge 4\mathcal{R}+16$, and let $g$ be Lipschitz and supported in $\Delta$.
Assume
\begin{equation}\label{10.7:eq-far-geometry-1}
\varepsilon(r_\Theta+2)\le \tfrac18,\qquad |y|+3\le L^m .
\end{equation}
Choose $\nu$ with $|y_\nu|=|y|_\infty$; then $|y_\nu|\ge |y|/2$. If $y_\nu>0$, choose
\begin{equation}\label{10.7:eq-far-geometry-2}
c\in \tfrac{\varepsilon}{2}\mathbb{Z}\cap\bigl[y_\nu-\tfrac32-\tfrac18,\;y_\nu-\tfrac32\bigr],\qquad
c'\in \tfrac{\varepsilon}{2}\mathbb{Z}\cap\bigl[\tfrac{\mathcal{R}}2+\tfrac18,\;\tfrac{\mathcal{R}}2+\tfrac3{16}\bigr]
\end{equation}
(both sets are non-empty), and put $u:=c-c'$, $\ell:=3$, $\ell':=\mathcal{R}+1$,
$\Sigma^+:=\{c<x_\nu<c+3\}$ and $\Sigma^-:=\{c'-\mathcal{R}-1<x_\nu<c'\}$. If $y_\nu<0$, use the mirror-image
data, obtained from these by the substitution $x_\nu\mapsto -x_\nu$ (the planes then lie on the negative side), with
the half-space $H^-_{\nu,c}$ in place of $H^+_{\nu,c}$ in Proposition~\ref{10.7:prop-slab-bound}. Then:
\begin{enumerate}
\item[(a)] $X:=\Theta^{\varepsilon}(g)-E[\Theta^{\varepsilon}(g)]$ is measurable on $\Sigma^+$;
\item[(b)] the cluster polynomial $Y$ of Lemma~\ref{10.7:lem-cluster-polynomial} for
$O_i=\mathcal{O}(h_i^{(m)})$ is measurable on $\Sigma^-$;
\item[(c)] the data $\nu,c,u,\ell,\ell'$ are admissible for Proposition~\ref{10.7:prop-slab-bound} as soon as
$u+\mathcal{R}+1\le L^m$ and $2u\in\varepsilon\mathbb{Z}$, and the second of these conditions holds;
\item[(d)] the mirror gap of $X$, that is, the distance in the coordinate $x_\nu$ between the set of bonds on which
$X$ depends and its image under the reflection $r_{\nu,c}$ (in the notation fixed in
Definition~\ref{10.7:not-reflected-slabs}), is at least $2\cdot\tfrac38$;
\item[(e)] the separation $2u$ satisfies
\begin{equation}\label{10.7:eq-far-geometry-3}
2u\;\ge\;2\bigl(y_\nu-\tfrac{13}8-\tfrac{\mathcal{R}}2-\tfrac3{16}\bigr)\;\ge\;|y|-\mathcal{R}-4\;\ge\;
\tfrac{\mathcal{R}'}4 .
\end{equation}
\end{enumerate}
\end{lemma}

\begin{proof}
Since $d=4$, $|y|^2=\sum_{\mu}y_\mu^2\le 4|y|_\infty^2$, so $|y_\nu|=|y|_\infty\ge |y|/2$.

Since $r_\Theta\ge 0$, the first condition in \eqref{10.7:eq-far-geometry-1} gives $2\varepsilon\le\tfrac18$, that is
$\varepsilon\le\tfrac1{16}$, and hence $\varepsilon/2\le\tfrac1{16}$. The two intervals in
\eqref{10.7:eq-far-geometry-2} are closed and of lengths $\tfrac18$ and $\tfrac1{16}$, both at least $\varepsilon/2$,
so each contains a point of $\tfrac{\varepsilon}2\mathbb{Z}$; the choices of $c$ and $c'$ are possible.

We treat the case $y_\nu>0$. In the case $y_\nu<0$ the mirror-image data are the images of the data above under
$x_\nu\mapsto -x_\nu$, and every inclusion and inequality below holds for them after the same substitution,
applied to $x_\nu$, $y_\nu$, $c$, $c'$, with $H^-_{\nu,c}$ in place of $H^+_{\nu,c}$.

We note that both slabs lie in $\{|x_\nu|<L^m\}$. Indeed, by \eqref{10.7:eq-far-geometry-1},
\[
c+3\le y_\nu+\tfrac32\le |y|+\tfrac32<L^m,
\]
and, since $L^m\ge |y|+3\ge \mathcal{R}'/2+3\ge 2\mathcal{R}+11$,
\[
c'-\mathcal{R}-1\ge \tfrac{\mathcal{R}}2+\tfrac18-\mathcal{R}-1=-\tfrac{\mathcal{R}}2-\tfrac78>-L^m .
\]

(a) By \eqref{10.7:eq-far-geometry-a}, $\Theta^{\varepsilon}(g)$ involves only the operators
$\Theta^{\varepsilon}_{\nu\sigma}(x)$ with $x\in\operatorname{supp}g\subset\Delta$, and each of these depends only on
bonds within distance $\varepsilon r_\Theta$ of $x$. By \eqref{10.7:eq-far-geometry-1},
$\varepsilon r_\Theta\le\tfrac18$, so the bonds on which $\Theta^{\varepsilon}(g)$ depends lie within $\tfrac18$ of
$\Delta$. As $\Delta$ has side $2$ and centre $y$, $\Delta\subset\{|x_\nu-y_\nu|\le 1\}$, so these bonds lie in
$\{|x_\nu-y_\nu|\le\tfrac98\}$. By \eqref{10.7:eq-far-geometry-2},
\[
c\le y_\nu-\tfrac32<y_\nu-\tfrac98,\qquad c+3\ge y_\nu-\tfrac{13}8+3=y_\nu+\tfrac{11}8>y_\nu+\tfrac98,
\]
so all these bonds have both endpoints in $\Sigma^+$. Subtracting the constant $E[\Theta^{\varepsilon}(g)]$ does not
change this, and $X$ is measurable on $\Sigma^+$.

(b) By Lemma~\ref{10.7:lem-cluster-polynomial}, for $O_i=\mathcal{O}(h_i^{(m)})$ the cluster polynomial $Y$ is
measurable on the bonds of plaquettes whose barycentres lie in $\bigcup_i\operatorname{supp}h_i^{(m)}$. Since
$\vec h\in\mathcal{T}^{(2)}_n(\mathcal{R},\mathsf{d})$, each $h_i$ is supported in $B(0,\mathcal{R}/2)$, so such
barycentres satisfy $|x(p)_\nu|\le\mathcal{R}/2$; the vertices of $p$ differ from $x(p)$ by $\varepsilon/2$ in each
coordinate, so the bonds of these plaquettes lie in $\{|x_\nu|\le\mathcal{R}/2+\varepsilon\}$. By
\eqref{10.7:eq-far-geometry-2} and $\varepsilon\le\tfrac1{16}$,
\[
c'\ge\tfrac{\mathcal{R}}2+\tfrac18>\tfrac{\mathcal{R}}2+\varepsilon,\qquad
c'-\mathcal{R}-1\le-\tfrac{\mathcal{R}}2-\tfrac{13}{16}<-\tfrac{\mathcal{R}}2-\varepsilon,
\]
so these bonds have both endpoints in $\Sigma^-$, and $Y$ is measurable on $\Sigma^-$.

(c) With $\ell'=\mathcal{R}+1$ one has $\Sigma^-=\{c-u-\ell'<x_\nu<c-u\}$ and $\Sigma^+=\{c<x_\nu<c+\ell\}$, which
is the slab configuration of Proposition~\ref{10.7:prop-slab-bound}; its admissibility is asserted under the proviso
$u+\mathcal{R}+1\le L^m$ together with $2u\in\varepsilon\mathbb{Z}$. The latter holds: $c$ and $c'$ lie in
$\tfrac{\varepsilon}2\mathbb{Z}$, so $2u=2c-2c'\in\varepsilon\mathbb{Z}$.

(d) By (a) the bonds on which $X$ depends lie in $\{x_\nu\ge y_\nu-\tfrac98\}$, and by
\eqref{10.7:eq-far-geometry-2}, $y_\nu-\tfrac98-c\ge y_\nu-\tfrac98-(y_\nu-\tfrac32)=\tfrac38$. The reflection
$r_{\nu,c}$ maps $x_\nu$ to $2c-x_\nu$, so the image lies in $\{x_\nu\le c-\tfrac38\}$, and the gap is at least
$2\cdot\tfrac38$.

(e) By \eqref{10.7:eq-far-geometry-2}, $c\ge y_\nu-\tfrac{13}8$ and $c'\le\tfrac{\mathcal{R}}2+\tfrac3{16}$, which
gives the first inequality in \eqref{10.7:eq-far-geometry-3}. Its middle member equals
$2y_\nu-\mathcal{R}-\tfrac{29}8$; since $2y_\nu\ge|y|$ and $\tfrac{29}8\le4$, it is at least $|y|-\mathcal{R}-4$.
Finally, $|y|\ge\mathcal{R}'/2$ and $\mathcal{R}'\ge4\mathcal{R}+16$ give
\[
|y|-\mathcal{R}-4\ge\tfrac{\mathcal{R}'}4+\Bigl(\tfrac{\mathcal{R}'}4-\mathcal{R}-4\Bigr)\ge\tfrac{\mathcal{R}'}4 .
\]
\end{proof}

\begin{theorem}[The infrared hypothesis from scalar reflected decay]
\label{10.7:thm-infrared-reduction}
Assume the two hypotheses displayed in \eqref{10.7:eq-hypotheses-used-form-b} other than
the stress-tensor decay form of the infrared hypothesis, in the forms displayed there. Assume also the further
standing inputs listed in Definition~\ref{10.7:def-annulus-functionals}. In particular, $j$ indexes
the diagonal sequence of lattices fixed there, along which the subsequential limit family is taken.
Assume finally the structural assumption \eqref{10.7:eq-far-geometry-a} on the lattice stress
density $\Theta^{\varepsilon}_{\mu\sigma}$. Write $\mathbb{E}_j$, $\operatorname{Cov}_j$ and
$\kappa_j$ for expectation, covariance and joint cumulant under the lattice measure of the $j$-th
member of that diagonal sequence, and $\varepsilon=\varepsilon_j=L^{-K_j}$ for its lattice spacing.
For a function $g$ put
\begin{equation*}
\Theta^{\varepsilon}_{\mu\sigma}(g):=\sum_x\varepsilon^4g(x)\,\Theta^{\varepsilon}_{\mu\sigma}(x).
\end{equation*}
Fix $\omega,n,\mathcal{R},\mathsf{d}$ as in Corollary~\ref{10.7:cor-clustering-dichotomy} and
$\vec h\in\mathcal{T}^{(2)}_n(\mathcal{R},\mathsf{d})$. Write $\mathcal{O}(\vec h)$ for the tuple
$(\mathcal{O}(h_i^{(m_j)}))_{1\le i\le n}$, and $Y_j$ for its cluster polynomial
(Lemma~\ref{10.7:lem-cluster-polynomial}) under the $j$-th measure. We make two assumptions.

The first is an ultraviolet bound of tensor class. There is a constant $C_\Theta$, independent of
$j$, for which the first inequality of \eqref{10.7:eq-infrared-reduction-a} below holds for all
sufficiently large $j$, for every closed cube $\Delta$ of side $2$ together with a reflection
$r_{\nu,c}$ and plane data as produced for $\Delta$ by Lemma~\ref{10.7:lem-far-geometry}, for every
pair $(\mu,\sigma)$, and for every real Lipschitz function $g$ supported in $\Delta$.

The second is an infrared decay condition on the reflected covariances of the cluster polynomial
$Y_j$. For $u_0>0$ let
\begin{equation}\label{10.7:eq-infrared-reduction-1}
\rho(u_0):=\sup\,\limsup_{j\to\infty}
\operatorname{Cov}_j\bigl(Y_j,Y_j^{r_{\nu,c_j}}\bigr)^{1/2}.
\end{equation}
Here the supremum runs over the direction $\nu$, the sign, and the $j$-dependent lattice plane data
$(c_j,c'_j)$ that Lemma~\ref{10.7:lem-far-geometry} produces for cubes whose offset $u$ of the slab
$\Sigma^-$ from the reflecting hyperplane $\{x_\nu=c_j\}$ satisfies $u\ge u_0$. The condition is the
second relation below. The two assumptions read
\begin{equation}\label{10.7:eq-infrared-reduction-a}
\begin{aligned}
&\operatorname{Cov}_j\bigl(\Theta^{\varepsilon}_{\mu\sigma}(g),
\Theta^{\varepsilon}_{\mu\sigma}(g)^{r_{\nu,c}}\bigr)\le C_\Theta\|g\|_{\sharp}^2,\\
&\lim_{u_0\to\infty}u_0^4\,\rho(u_0)=0.
\end{aligned}
\end{equation}

Then, with an absolute constant $C_0$ and a constant $C'(\omega,\chi_1,M)$ depending only on
$\omega$, $\chi_1$ and $M$, both fixed in the proof, the radial cut-offs $\chi_{\mathcal{R}'}$
entering \eqref{10.7:eq-annulus-functionals-a} satisfy, for every $\mathcal{R}'\ge4\mathcal{R}+16$,
\begin{equation}\label{10.7:eq-infrared-reduction-2}
\limsup_{j\to\infty}\bigl|\kappa_j\bigl(B^4_{\chi_{\mathcal{R}'}};\mathcal{O}(\vec h)\bigr)\bigr|
\le16\,C'(\omega,\chi_1,M)\,C_0\,\mathcal{R}'^4\,C_\Theta^{1/2}\,\rho(\mathcal{R}'/8),
\end{equation}
and the right-hand side tends to $0$ as $\mathcal{R}'\to\infty$.

Together with Proposition~\ref{10.7:prop-divergence-vanishes}, parts (i) and (ii), this gives the
following. If the infrared condition holds for every
$\vec h\in\mathcal{T}^{(2)}_n(\mathcal{R},\mathsf{d})$, then the infrared hypothesis holds for
$(\omega,n,\mathcal{R},\mathsf{d})$ in the stress-tensor decay form of
\eqref{10.7:eq-hypotheses-used-form-b}, by which one far-away stress-tensor insertion decorrelates
from the observables, and $c(\vec h,\omega)=0$ for every $\vec h$ of the class. This decay form is a
sufficient condition for the far-sphere flux hypothesis (IR-fl)$_{\omega_0}$ of clause~(vi). If the
infrared condition holds for all $\omega,n,\mathcal{R},\mathsf{d}$ and all such $\vec h$, then the $O(4)$ invariance
asserted by the conclusion step, Theorem~\ref{10.7:thm-conclusion-step}, holds for this
subsequential limit family.
\end{theorem}

\begin{proof}
Fix $\mathcal{R}'\ge4\mathcal{R}+16$ and put
\begin{equation*}
w_{\mu\sigma}(x):=\sum_\rho\omega_{\rho\sigma}\,x_\rho\,(\partial_\mu\chi_{\mathcal{R}'})(x).
\end{equation*}
Comparing with the display of $B^4_\chi$ in \eqref{10.7:eq-annulus-functionals-a} (with the
summation index $\nu$ there renamed $\mu$), we obtain
\begin{equation*}
B^4_{\chi_{\mathcal{R}'}}=\sum_{\mu,\sigma}\Theta^{\varepsilon}_{\mu\sigma}(w_{\mu\sigma}).
\end{equation*}
The sum runs over the sixteen pairs $(\mu,\sigma)$. The weights satisfy
$\|w_{\mu\sigma}\|_{\sharp}\le C(\omega,\chi_1,M)$ uniformly in $\mathcal{R}'$. Indeed,
$\chi_{\mathcal{R}'}$ is the dilate $x\mapsto\chi_1(x/\mathcal{R}')$, and $|x|\le2\mathcal{R}'$ on the
support of $\partial_\mu\chi_{\mathcal{R}'}$. Hence, for $\mathcal{R}'\ge1$, both
$\|w_{\mu\sigma}\|_\infty$ and $\operatorname{Lip}w_{\mu\sigma}$ are bounded by constants depending
only on $\omega$ and on the derivative bounds of $\chi_1$.

Let $(\psi_\Delta)_\Delta$ be a Lipschitz partition of unity on $\mathbb{R}^4$, fixed once and for
all. It is indexed by the closed cubes $\Delta$ of side $2$ centred at the points of
$\mathbb{Z}^4$, with $0\le\psi_\Delta\le1$, $\operatorname{supp}\psi_\Delta\subset\Delta$,
$\sum_\Delta\psi_\Delta=1$ and $\sup_\Delta\operatorname{Lip}\psi_\Delta<\infty$. We bound the norm
of a localised weight. Since $0\le\psi_\Delta\le1$, we have
$\|w_{\mu\sigma}\psi_\Delta\|_\infty\le\|w_{\mu\sigma}\|_\infty$. The product rule gives
\begin{equation*}
\operatorname{Lip}(w_{\mu\sigma}\psi_\Delta)\le\operatorname{Lip}w_{\mu\sigma}
+\|w_{\mu\sigma}\|_\infty\operatorname{Lip}\psi_\Delta.
\end{equation*}
Hence
\begin{equation*}
\|w_{\mu\sigma}\psi_\Delta\|_{\sharp}\le(1+40M\operatorname{Lip}\psi_\Delta)\,
\|w_{\mu\sigma}\|_{\sharp}.
\end{equation*}
Put
\begin{equation*}
C'(\omega,\chi_1,M):=\Bigl(1+40M\sup_\Delta\operatorname{Lip}\psi_\Delta\Bigr)\,C(\omega,\chi_1,M).
\end{equation*}
Then $\|w_{\mu\sigma}\psi_\Delta\|_{\sharp}\le C'(\omega,\chi_1,M)$ for every $\Delta$.

The function $w_{\mu\sigma}\psi_\Delta$ vanishes unless $\Delta$ meets the support of
$\partial_\mu\chi_{\mathcal{R}'}$. That support lies in the shell
$\{\mathcal{R}'\le|x|\le2\mathcal{R}'\}$. At most $C_0\mathcal{R}'^4$ cubes meet this shell, with
$C_0$ absolute.

Let $\Delta$ be such a cube, with centre $y$. Every point of $\Delta$ lies within distance $2$ of
$y$, the half-diagonal of a cube of side $2$ in $\mathbb{R}^4$. Hence $|y|\ge\mathcal{R}'-2$, and
$\mathcal{R}'-2\ge\mathcal{R}'/2$ because $\mathcal{R}'\ge4$. Linearity of
$g\mapsto\Theta^{\varepsilon}_{\mu\sigma}(g)$ and multilinearity of the cumulant give
\begin{equation}\label{10.7:eq-infrared-reduction-3}
\kappa_j\bigl(B^4_{\chi_{\mathcal{R}'}};\mathcal{O}(\vec h)\bigr)
=\sum_{\mu,\sigma}\sum_\Delta\kappa_j\bigl(\Theta^{\varepsilon}_{\mu\sigma}
(w_{\mu\sigma}\psi_\Delta);\mathcal{O}(\vec h)\bigr).
\end{equation}
Here the sum runs over the cubes meeting the shell.

Fix $(\mu,\sigma)$ and such a cube $\Delta$, and put
$X:=\Theta^{\varepsilon}_{\mu\sigma}(w_{\mu\sigma}\psi_\Delta)$. This is a bounded real variable,
since $\Theta^{\varepsilon}_{\mu\sigma}$ is bounded. By Lemma~\ref{10.7:lem-cluster-polynomial},
\begin{equation*}
\kappa_j\bigl(X;\mathcal{O}(\vec h)\bigr)=\mathbb{E}_j[XY_j],
\qquad \mathbb{E}_j[Y_j]=0.
\end{equation*}
Moreover $Y_j$ is measurable with respect to the bonds of the plaquettes whose barycentres lie in
$\bigcup_i\operatorname{supp}h_i^{(m_j)}$. Put $X_\Delta:=X-\mathbb{E}_j[X]$. Since
$\mathbb{E}_j[Y_j]=0$, we have $\mathbb{E}_j[XY_j]=\mathbb{E}_j[X_\Delta Y_j]$.

Lemma~\ref{10.7:lem-far-geometry} (slab geometry for far stress-tensor cubes) applies to $\Delta$
and $g=w_{\mu\sigma}\psi_\Delta$: indeed $|y|\ge\mathcal{R}'/2$, $\mathcal{R}'\ge4\mathcal{R}+16$, and
$g$ is Lipschitz and supported in $\Delta$. Its conditions on the lattice spacing and on the torus,
\begin{equation*}
\varepsilon_j(r_\Theta+2)\le1/8,\qquad |y|+3\le L^{m_j},
\end{equation*}
hold for all sufficiently large $j$, because $\varepsilon_j\to0$ and $m_j\to\infty$ along the
diagonal. Since there are finitely many cubes, one threshold in $j$ serves all of them. The lemma
produces a direction $\nu$, a sign and lattice plane data $(c_j,c'_j)$, with slabs $\Sigma^\pm$ for
the reflection $r=r_{\nu,c_j}$ and offset $u\ge\mathcal{R}'/8$ of $\Sigma^-$ from the hyperplane
$\{x_\nu=c_j\}$. One of $X_\Delta,Y_j$ is measurable on $\Sigma^+$ and the other on $\Sigma^-$.

The reflection-positivity bound for slab correlations (Proposition~\ref{10.7:prop-slab-bound}), whose
right-hand side is symmetric in the two variables, applies to the centred variables $X_\Delta$ and
$Y_j$ and gives
\begin{equation}\label{10.7:eq-infrared-reduction-4}
\bigl|\kappa_j\bigl(X;\mathcal{O}(\vec h)\bigr)\bigr|=|\mathbb{E}_j[X_\Delta Y_j]|
\le\operatorname{Cov}_j(X_\Delta,X_\Delta^{r})^{1/2}\operatorname{Cov}_j(Y_j,Y_j^{r})^{1/2}.
\end{equation}
The covariance does not change when constants are subtracted from its arguments, so
$\operatorname{Cov}_j(X_\Delta,X_\Delta^r)=\operatorname{Cov}_j(X,X^r)$. By the first bound of
\eqref{10.7:eq-infrared-reduction-a}, for all large $j$ this is at most
\begin{equation*}
C_\Theta\|w_{\mu\sigma}\psi_\Delta\|_{\sharp}^2\le C_\Theta\,C'(\omega,\chi_1,M)^2.
\end{equation*}
The second factor in \eqref{10.7:eq-infrared-reduction-4} is taken along plane data produced by
Lemma~\ref{10.7:lem-far-geometry} for a cube with $u\ge\mathcal{R}'/8$. By
\eqref{10.7:eq-infrared-reduction-1} its upper limit is therefore at most $\rho(\mathcal{R}'/8)$.
Taking $\limsup_j$ in \eqref{10.7:eq-infrared-reduction-4} gives
\begin{equation*}
\limsup_{j\to\infty}\bigl|\kappa_j\bigl(\Theta^{\varepsilon}_{\mu\sigma}
(w_{\mu\sigma}\psi_\Delta);\mathcal{O}(\vec h)\bigr)\bigr|
\le C'(\omega,\chi_1,M)\,C_\Theta^{1/2}\,\rho(\mathcal{R}'/8).
\end{equation*}

The sum \eqref{10.7:eq-infrared-reduction-3} has at most $16\,C_0\mathcal{R}'^4$ terms. The upper
limit of a finite sum is at most the sum of the upper limits, and
\eqref{10.7:eq-infrared-reduction-2} follows. Its right-hand side contains the factor
\begin{equation*}
\mathcal{R}'^4\rho(\mathcal{R}'/8)=8^4\,(\mathcal{R}'/8)^4\rho(\mathcal{R}'/8).
\end{equation*}
By the second relation of \eqref{10.7:eq-infrared-reduction-a}, this factor tends to $0$ as
$\mathcal{R}'\to\infty$.

By Definition~\ref{10.7:def-annulus-functionals} and \eqref{10.7:eq-annulus-functionals-a},
$\mathcal{B}^{\chi}_j(\vec h)=\kappa_j(\mathfrak{F}_\chi;\mathcal{O}(\vec h))$ with
$\mathfrak{F}_\chi=B^4_\chi+J^4_\chi$. Multilinearity of the cumulant therefore writes
$\mathcal{B}^{\chi_{\mathcal{R}'}}_j(\vec h)$ as the sum of the cumulants of
$B^4_{\chi_{\mathcal{R}'}}$ and of $J^4_{\chi_{\mathcal{R}'}}$ against $\mathcal{O}(\vec h)$. By
\eqref{10.7:eq-infrared-reduction-2}, the first contributes nothing to the double limit in item (ii)
of Corollary~\ref{10.7:cor-clustering-dichotomy}. By
Proposition~\ref{10.7:prop-divergence-vanishes}, parts (i) and (ii) (see
\eqref{10.7:eq-divergence-vanishes-b}), the divergence term $J^4_{\chi_{\mathcal{R}'}}$ vanishes in
the limit. Hence the double limit
\begin{equation*}
\lim_{\mathcal{R}'\to\infty}\limsup_{j\to\infty}\bigl|\mathcal{B}^{\chi_{\mathcal{R}'}}_j\bigr|=0
\end{equation*}
holds for the family $\chi_{\mathcal{R}'}$. The tuple
$\vec h\in\mathcal{T}^{(2)}_n(\mathcal{R},\mathsf{d})$ was arbitrary, and the infrared condition is
assumed for every such $\vec h$. This is therefore item (ii) of
Corollary~\ref{10.7:cor-clustering-dichotomy} for $(\omega,n,\mathcal{R},\mathsf{d})$. Its
equivalence with items (i) and (iv) there gives the infrared hypothesis in the stress-tensor
decay form of \eqref{10.7:eq-hypotheses-used-form-b} and $c(\vec h,\omega)=0$ on the class. When this
holds for all $\omega,n,\mathcal{R},\mathsf{d}$, the equivalence with item (v) there gives the $O(4)$
invariance asserted by the conclusion step, Theorem~\ref{10.7:thm-conclusion-step}, for this
subsequential limit family.
\end{proof}

\begin{proposition}[Scalar reflected decay from clustering or a mass gap]
\label{10.7:prop-clustering-mass-gap}
Let $\vec h=(h_1,\dots,h_n)$ be a tuple of test functions as in the infrared hypothesis
\eqref{10.7:eq-infrared-reduction-a} of Theorem~\ref{10.7:thm-infrared-reduction}, supported in
$B(0,\mathcal{R}/2)$. As in Theorem~\ref{10.7:thm-infrared-reduction}, $\mathbb{E}_j$,
$\operatorname{Cov}_j$ and $\kappa_j$ denote expectation, covariance and joint cumulant in the
lattice measure at index $j$ of the diagonal sequence. For each index $j$ of the diagonal sequence let $Y_j$ be the cluster
polynomial of Lemma~\ref{10.7:lem-cluster-polynomial} for $O_i=\mathcal{O}(h_i^{(m_j)})$. For a
direction $\nu$, a sign and plane data $(c_j,c'_j)$ with gap parameter $u$, as in the definition
of $\rho(u_0)$ in \eqref{10.7:eq-infrared-reduction-a}, let $\sigma_j$ be the Euclidean reflection
$x_\nu\mapsto 2c_j-x_\nu$ (for $c\in\mathbb{R}$ we write $r_{\nu,c}$ also for the Euclidean
reflection $x_\nu\mapsto 2c-x_\nu$ of $\mathbb{R}^4$, acting on test functions by
$h\mapsto h\circ r_{\nu,c}$, so that $\sigma_j=r_{\nu,c_j}$), and let $Y_j^{r}$ be the same
polynomial evaluated at the observables
$\mathcal{O}((h_i\circ\sigma_j)^{(m_j)})$; the supports of the $h_i\circ\sigma_j$ lie at
$\nu$-gap at least $2u$ from those of the $h_i$. A sub-family $\beta$ of
\begin{equation*}
\{\mathcal{O}(h_i^{(m_j)})\}_{i\le n}\cup\{\mathcal{O}((h_i\circ\sigma_j)^{(m_j)})\}_{i\le n},
\end{equation*}
and its cumulant $\kappa_j(\beta)$, is called \emph{crossing} if it has members of both families.
Consider the two conditions (in the second, $H$ and $\Omega$ are the data of (b) below)
\begin{equation}\label{10.7:eq-clustering-mass-gap-a}
\begin{gathered}
\kappa_j(\beta)\le C(\vec h)\,(1+u)^{-8-2\delta}
\quad\text{for every crossing } \beta \text{ and all large } j,\\
\ker H=\mathbb{C}\Omega,\qquad \operatorname{spec}H\subset\{0\}\cup[m_{gap},\infty),
\quad m_{gap}>0 .
\end{gathered}
\end{equation}
\begin{enumerate}
\item[(a)] If for some $\delta>0$ the first condition in \eqref{10.7:eq-clustering-mass-gap-a}
holds for every crossing cumulant, every $\nu$, sign and plane data, then for all large $j$
one has $\operatorname{Cov}_j(Y_j,Y_j^{r})\le C'(\vec h)(1+u)^{-8-2\delta}$, and
$\rho(u_0)\le C''u_0^{-4-\delta}$; in particular \eqref{10.7:eq-infrared-reduction-a} holds.
\item[(b)] Let $(\mathcal{H},\Omega,T_t=e^{-tH})$ be the Osterwalder--Schrader data of
Definition~\ref{10.7:def-reflection-data}, Proposition~\ref{10.7:prop-reflection-form} and
Theorem~\ref{10.7:thm-rotation-representation}, built from the limit functions $S^T_n$ with the
reflection $\theta_0\colon x_\nu\mapsto -x_\nu$ in place of the time reflection (these data are
unitarily equivalent to those built with the time reflection, by an element of $W_4$). If the
second condition in \eqref{10.7:eq-clustering-mass-gap-a} holds, then
$\rho(u_0)\le C_0(\vec h)^{1/2}e^{-u_0 m_{gap}}$, where $C_0(\vec h)$ depends on $\vec h$ only
and bounds $\operatorname{Cov}_j(Y_j,Y_j^{r'})$ for all large $j$ and all plane data, $r'$
denoting the reflection at $c'_j$ (such a constant exists by the per-ball volume-free bound);
in particular \eqref{10.7:eq-infrared-reduction-a} holds.
\end{enumerate}
\end{proposition}

\begin{proof}
\emph{Expansion of the reflected covariance.} By Lemma~\ref{10.7:lem-cluster-polynomial},
$\mathbb{E}_j[Y_j]=0$, so that
$\operatorname{Cov}_j(Y_j,Y_j^{r})=\mathbb{E}_j[Y_jY_j^{r}]$. By the same lemma $Y_j$ is a real
multilinear polynomial in the centred variables $\tilde O_i$; write
$Y_j=\sum_{B}a_{B,j}\prod_{i\in B}\tilde O_i$, the sum over $B\subset\{1,\dots,n\}$, and
correspondingly $Y_j^{r}=\sum_{B}a_{B,j}\prod_{i\in B}\tilde O_i^{r}$ with $\tilde O_i^{r}$ the
centred $\mathcal{O}((h_i\circ\sigma_j)^{(m_j)})$. The coefficients $a_{B,j}$ are polynomials in
centred moments and are bounded uniformly in $j$ by the per-ball volume-free bound. By
multilinearity, $\mathbb{E}_j[Y_jY_j^{r}]$ is the finite sum over $(B,B')$ of
$a_{B,j}a_{B',j}$ times the centred moment of
$\{\tilde O_i\}_{i\in B}\cup\{\tilde O_{i'}^{r}\}_{i'\in B'}$. By the moment--cumulant formula
for centred variables, this moment is the sum, over the partitions $\pi$ of the index set
without singletons, of $\prod_{\beta\in\pi}\kappa_j(\beta)$; since every block has at least two
members, $\kappa_j(\beta)$ is the cumulant of the corresponding sub-family of the uncentred
observables. Call $\pi$ non-crossing if each block lies in one family. A non-crossing $\pi$ is
exactly a pair of partitions without singletons of $B$ and of $B'$, so the non-crossing terms
factorise and their total is
\begin{equation*}
\sum_{B,B'}a_{B,j}a_{B',j}\,
\mathbb{E}_j\Bigl[\prod_{i\in B}\tilde O_i\Bigr]\,
\mathbb{E}_j\Bigl[\prod_{i'\in B'}\tilde O_{i'}^{r}\Bigr]
=\mathbb{E}_j[Y_j]\,\mathbb{E}_j[Y_j^{r}]=0 .
\end{equation*}
Hence
\begin{equation}\label{10.7:eq-clustering-mass-gap-1}
\operatorname{Cov}_j(Y_j,Y_j^{r})
=\sum_{B,B'}a_{B,j}a_{B',j}\sum_{\pi\ \text{crossing}}\ \prod_{\beta\in\pi}\kappa_j(\beta),
\end{equation}
and every surviving term carries at least one crossing cumulant.

\emph{Proof of} (a). In each term of \eqref{10.7:eq-clustering-mass-gap-1} the coefficient is
bounded uniformly in $j$ by the per-ball volume-free bound, one crossing factor is at most
$C(\vec h)(1+u)^{-8-2\delta}$ by hypothesis for all large $j$, and each remaining factor is
either crossing, hence at most $C(\vec h)$, or a cumulant of a sub-family of one family, bounded
uniformly in $j$ by the per-ball volume-free bound. The number of terms depends on $n$ only.
Summing gives $\operatorname{Cov}_j(Y_j,Y_j^{r})\le C'(\vec h)(1+u)^{-8-2\delta}$ for all large
$j$, uniformly in $\nu$, the sign and the plane data. Taking the $\limsup_j$, the square root
and the supremum over data with $u\ge u_0$ in the definition of $\rho(u_0)$,
\begin{equation*}
\rho(u_0)\le C'(\vec h)^{1/2}(1+u_0)^{-4-\delta}\le C''u_0^{-4-\delta},
\qquad C'':=C'(\vec h)^{1/2}.
\end{equation*}

\emph{Proof of} (b). Fix $\nu$, a sign and plane data, and a subsequence along which
$\operatorname{Cov}_j(Y_j,Y_j^{r})$ tends to its $\limsup$. Pass to a further subsequence along
which $(c_j,c'_j)\to(c_*,c'_*)$, and put $\sigma_*:=r_{\nu,c_*}$. For $y\in\mathbb{R}$ let
$\tau_y$ be the translation by $y$ in the $\nu$-th coordinate direction, acting on test
functions by $(\tau_yh)(x)=h(x')$, where $x'$ agrees with $x$ except that $x'_\nu=x_\nu-y$. Then
$h_i\circ\sigma_j=\tau_{2(c_j-c_*)}(h_i\circ\sigma_*)$ with $2(c_j-c_*)\to0$, and
$\|\tau_yh-h\|_{\sharp}\to0$ as $y\to0$; by the equicontinuity estimate
\eqref{10.7:eq-equicontinuity-a} of Lemma~\ref{10.7:lem-equicontinuity}, replacing $\sigma_j$ by
$\sigma_*$ changes every cumulant in \eqref{10.7:eq-clustering-mass-gap-1} by $o(1)$. By the
infinite-volume theorem each cumulant then converges, and so do the coefficients; let $a_B$ be
the limits of the $a_{B,j}$. With $S^c$ the centred correlation functions of
Definition~\ref{10.7:def-centred-functions}, the limits of the centred moments, we obtain
\begin{equation}\label{10.7:eq-clustering-mass-gap-2}
\lim_j\operatorname{Cov}_j(Y_j,Y_j^{r})
=\sum_{B,B'}a_Ba_{B'}\,S^c\bigl((h_i)_{i\in B},(h_{i'}\circ\sigma_*)_{i'\in B'}\bigr).
\end{equation}
Define
\begin{equation*}
\hat v:=\sum_B a_B\,\iota\bigl[(h_i\circ r_{\nu,c'_*/2})_{i\in B}\bigr]\in\mathcal{H};
\end{equation*}
the supports of the $h_i\circ r_{\nu,c'_*/2}$ lie in $\{x_\nu\ge c'_*-\mathcal{R}/2\}$, with
$c'_*-\mathcal{R}/2>0$. For the plane data in question $2c_*-c'_*=c'_*+2u$, and with
the translations $\tau_y$ as above, direct substitution gives
\begin{equation}\label{10.7:eq-clustering-mass-gap-3}
\theta_0\cdot(h\circ r_{\nu,c'_*/2})=\tau_{-c'_*}h,\qquad
\tau_{2u}\cdot(h\circ r_{\nu,c'_*/2})=\tau_{-c'_*}(h\circ r_{\nu,c_*}).
\end{equation}
The inner product of $\mathcal{H}$ evaluates $S^c$ on the $\theta_0$-reflected first argument
together with the second, and $T_{2u}$ acts on symbols by $\tau_{2u}$. Hence, by
\eqref{10.7:eq-clustering-mass-gap-3}, the translation invariance of $S^c$ and
\eqref{10.7:eq-clustering-mass-gap-2},
\begin{equation*}
\langle\hat v,T_{2u}\hat v\rangle
=\sum_{B,B'}a_Ba_{B'}\,S^c\bigl((h_i)_{i\in B},(h_{i'}\circ\sigma_*)_{i'\in B'}\bigr)
=\lim_j\operatorname{Cov}_j(Y_j,Y_j^{r}).
\end{equation*}
Since $h\circ r_{\nu,c'_*/2}=\tau_{-c'_*}(h\circ r_{\nu,c'_*})$, the same reasoning at $u=0$
gives $\|\hat v\|^2=\lim_j\operatorname{Cov}_j(Y_j,Y_j^{r'})\le C_0(\vec h)$, by the per-ball
volume-free bound, and $\langle\Omega,\hat v\rangle=\lim_j\mathbb{E}_j[Y_j^{r'}]=0$.

Because $\ker H=\mathbb{C}\Omega$, the spectral projection of $H$ onto $\{0\}$ is the projection
onto $\mathbb{C}\Omega$, and $\hat v\perp\Omega$; as
$\operatorname{spec}H\subset\{0\}\cup[m_{gap},\infty)$, the spectral measure of $\hat v$ is
carried by $[m_{gap},\infty)$. By the spectral theorem
\begin{equation*}
\lim_j\operatorname{Cov}_j(Y_j,Y_j^{r})=\langle\hat v,e^{-2uH}\hat v\rangle
\le e^{-2um_{gap}}\|\hat v\|^2\le C_0(\vec h)\,e^{-2um_{gap}} .
\end{equation*}
This holds along the chosen subsequence, hence for the $\limsup$ along the full sequence.
Taking square roots and the supremum over data with $u\ge u_0$ yields
$\rho(u_0)\le C_0(\vec h)^{1/2}e^{-u_0m_{gap}}$, so \eqref{10.7:eq-infrared-reduction-a} holds,
with exponential room.

The argument is carried out on the Osterwalder--Schrader space of the limit. The corresponding
inequality on the finite torus,
$\mathbb{E}_j[F^{r}F^{\tau_{2t}}]\le e^{-2m_{gap}t}\,\mathbb{E}_j[F^{r}F]$
for all $t\ge0$ and every function $F$ of the bonds of a half-torus, with $F^{r}$ the reflected
function, is false, since the reflection of the torus fixes two hyperplanes; it is not used.
\end{proof}

\begin{corollary}[A non-zero anomaly forces slow decay]\label{10.7:cor-slow-decay}
Assume all hypotheses of Theorem~\ref{10.7:thm-infrared-reduction} other than its scalar
reflected-decay hypothesis (the decay $\rho(u_0)=o(u_0^{-4})$ displayed in
\eqref{10.7:eq-infrared-reduction-a}). Among the hypotheses assumed are the slab geometry
\eqref{10.7:eq-far-geometry-a} for far stress-tensor cubes and the ultraviolet tensor-class bound
in \eqref{10.7:eq-infrared-reduction-a}. Assume, together with these, the standing inputs listed in
Definition~\ref{10.7:def-annulus-functionals} and the boundedness estimate in
\eqref{10.7:eq-annulus-functionals-a}. Suppose that
$c(\vec h,\omega)\neq 0$ for some tuple $\vec h$ of test functions as in
Theorem~\ref{10.7:thm-infrared-reduction}. Then it is not the case that
\begin{equation}\label{10.7:eq-slow-decay-1}
\sup_{\substack{\nu,\ \pm,\ (c_j,c'_j)\\ u\ge u_0}}\ \limsup_{j}
\operatorname{Cov}_j\bigl(Y_j,Y_j^{r_{\nu,c_j}}\bigr)=o(u_0^{-8})
\qquad(u_0\to\infty),
\end{equation}
where the supremum is taken over the direction $\nu$, the sign and the lattice plane data
$(c_j,c'_j)$ of Theorem~\ref{10.7:thm-infrared-reduction} for this tuple, for cubes whose gap
parameter $u$ in that data satisfies $u\ge u_0$. Here, in the notation of
Theorem~\ref{10.7:thm-infrared-reduction}, $Y_j$ is the cluster polynomial of $\vec h$ at
index~$j$, $Y_j^{r_{\nu,c_j}}$ is its image under the lattice reflection $r_{\nu,c_j}$ in the
plane $x_\nu=c_j$, whose reflected supports lie at $\nu$-gap at least $2u$ from the original
ones, and $\operatorname{Cov}_j$ is the covariance in the lattice measure at index~$j$ along the
diagonal.
Two-cluster correlations of the scalar curvature observables therefore fail to decay at this rate.
This is an infrared pathology of the smeared curvature Schwinger functions themselves; compare
clustering, a mass gap and uniqueness of the vacuum, none of which is asserted anywhere in this
book.
\end{corollary}

\begin{proof}
We argue by contraposition for the fixed tuple $\vec h$. Suppose that \eqref{10.7:eq-slow-decay-1}
holds. The covariances $\operatorname{Cov}_j(Y_j,Y_j^{r_{\nu,c_j}})$ are non-negative by
reflection positivity of the lattice measure with respect to the plane $x_\nu=c_j$, since
$Y_j^{r_{\nu,c_j}}$ is the reflection of $Y_j$ in that plane; this is the fact on which the
definition of $\rho(u_0)$ through square roots in Theorem~\ref{10.7:thm-infrared-reduction} rests.
Since the square root is increasing and continuous on $[0,\infty)$, it commutes with $\limsup_j$
and with the supremum, and taking square roots gives
\begin{equation*}
\rho(u_0)=\sup_{\substack{\nu,\ \pm,\ (c_j,c'_j)\\ u\ge u_0}}\ \limsup_j
\operatorname{Cov}_j\bigl(Y_j,Y_j^{r_{\nu,c_j}}\bigr)^{1/2}
=o(u_0^{-4}).
\end{equation*}
This is the scalar reflected-decay hypothesis of Theorem~\ref{10.7:thm-infrared-reduction},
displayed in \eqref{10.7:eq-infrared-reduction-a}.

Every other hypothesis of that theorem is assumed in the corollary, so every hypothesis of
Theorem~\ref{10.7:thm-infrared-reduction} is then in force, and the theorem applies to $\vec h$.
Its conclusion is, for $\vec h$, the stress-tensor decay condition: one far-away stress-tensor
insertion decorrelates from the observables faster than $|x|^{-4}$, uniformly in the cutoff.
This condition is sufficient for the far-sphere flux hypothesis $(\mathrm{IR\text{-}fl})_{\omega_0}$
of clause (vi), the vanishing of the flux of the lattice rotation current through far spheres along
one sequence of radii, which therefore holds for $\vec h$. The conclusion step is
applied with hypothesis $(\mathrm{IR\text{-}fl})_{\omega_0}$ for $\vec h$, with the boundedness
estimate in \eqref{10.7:eq-annulus-functionals-a}, and with the standing inputs listed in
Definition~\ref{10.7:def-annulus-functionals}; the first has just been obtained and the other two
are assumed in the corollary, so the conclusion step gives $c(\vec h,\omega)=0$. This proves the
contrapositive.
\end{proof}

\emph{Scope.} The proof of Theorem~\ref{10.7:thm-infrared-reduction} delivers
the decay bound in \eqref{10.7:eq-annulus-functionals-a} for the stress-tensor density
only in the far zone. There the cube
lies far outside a ball containing all the sources, and this is all that the stress-tensor decay
condition, the sufficient condition for the far-sphere flux hypothesis
$(\mathrm{IR\text{-}fl})_{\omega_0}$, requires. The bound
for cubes in general position is not asserted.

\begin{definition}[Truncated lattice expectations and the joint diagonal]
\label{10.8:not-joint-diagonal}
\begin{enumerate}
\item[(a)] \emph{Lattice, torus, measure.} The blocking factor $L \ge L_0$ is a fixed odd integer.
For $m \ge 1$ and $K \ge 0$ the torus is
$\mathbb{T}_m := \mathbb{R}^4 / 2L^m\mathbb{Z}^4$ with its quotient metric $\operatorname{dist}$.
The lattice spacing is $\varepsilon := L^{-K}$; the letter $a$ is also used for this spacing.
The bare sites are the points of $\varepsilon(\mathbb{Z}+\tfrac12)^4$ modulo $2L^m\mathbb{Z}^4$.
A plaquette $p$ is an elementary square of this lattice, with barycentre $x(p)$ and
action density
\begin{equation*}
a_p := 1 - \operatorname{Re}\operatorname{tr} U(\partial p) \in [0,2],
\end{equation*}
where the gauge group is $SU(2)$ and $\operatorname{tr}\mathbf{1} = 1$.
The measure $\mu_{K,m}$ is the Wilson probability measure with density
$\rho_0 = \exp\bigl[-g_0^{-2}\sum_p a_p\bigr]$, and $\langle\cdot\rangle_{K,m}$ is its
expectation. For bounded real random variables $Y_1,\dots,Y_N$ the \emph{joint cumulant} is
\begin{equation}\label{10.8:eq-joint-diagonal-1}
\kappa(Y_1,\dots,Y_N) := \partial_{z_1}\cdots\partial_{z_N}
\log\bigl\langle e^{\sum_{i=1}^N z_i Y_i}\bigr\rangle_{K,m}\Big|_{z=0},
\end{equation}
and the \emph{truncated expectation} is
$\langle A;B_1;\dots;B_n\rangle_{K,m} := \kappa(A,B_1,\dots,B_n)$;
truncation is always taken with respect to $\mu_{K,m}$.
For a function $f$ we write $\|f\|_{\sharp} := \max\{\|f\|_\infty,\ 40M\cdot\operatorname{Lip} f\}$,
where $M$ is the cube side among Ba{\l}aban's auxiliary parameters.
\item[(b)] \emph{Observables and connected correlations.} For a real function $\varphi$ on
$\mathbb{T}_m$ we set $\mathcal{O}(\varphi) := \sum_p \varphi(x(p))\,a_p$, the sum running over
all plaquettes in all six orientations. For a compactly supported function $h$ on $\mathbb{R}^4$,
$h^{(m)}$ denotes its periodic lift to $\mathbb{T}_m$. The connected $n$-point curvature
correlations are
\begin{equation}\label{10.8:eq-joint-diagonal-2}
S^T_{n;K,m}(\vec h) := \kappa\bigl(\mathcal{O}(h_1^{(m)}),\dots,\mathcal{O}(h_n^{(m)})\bigr).
\end{equation}
\item[(c)] \emph{The joint diagonal.} Along the sequence $g_0^{(j)} \downarrow 0$ of
clause (ii-b) one has $K_j = K(g_0^{(j)},g) \to \infty$ by clause (0), hence
$\varepsilon_j := L^{-K_j} \to 0$, and $m_j \to \infty$; the sequence $(K_j,m_j)$ is the
\emph{joint diagonal} of clause (ii-e). Along it the following holds: for every $n \ge 2$ and
all real $f_1,\dots,f_n \in C^1_c(\mathbb{R}^4)$ with pairwise disjoint supports,
\begin{equation}\label{10.8:eq-joint-diagonal-3}
S^T_{n;K_j,m_j}(\vec f) \longrightarrow S^T_{n;\infty}(\vec f)
\qquad (j \to \infty),
\end{equation}
and the limit is multilinear, symmetric, real, exactly invariant under
$\mathbb{R}^4 \rtimes W_4$, where $W_4$ is the group of signed permutation matrices, and it
obeys the position-uniform bound of clause (iii-b).
In what follows $\limsup_{j\to\infty}$ always runs along this diagonal, and
$\langle\cdot\,;\dots;\cdot\rangle_j := \langle\cdot\,;\dots;\cdot\rangle_{K_j,m_j}$.
\end{enumerate}
\end{definition}

The convergence \eqref{10.8:eq-joint-diagonal-3} and the listed properties of its limit are
the $\mathbb{R}^4$ part of clause (ii) of Theorem~0, and they are used in this chapter
without re-proof.

\begin{definition}[The rotation generator and the radial cut-offs]\label{10.8:def-radial-cutoff}
Let $\omega_0$ be one fixed real antisymmetric $4\times 4$ matrix with $\omega_0\neq 0$, that is, a
non-zero element $\omega_0\in\mathfrak{so}(4)$. The corresponding linear rotation field on
$\mathbb{R}^4$ is
\begin{equation}\label{10.8:eq-radial-cutoff-1}
\xi(x):=\omega_0 x,\qquad
\xi^{\rho}(x)=\sum_{\tau}(\omega_0)_{\rho\tau}\,x_{\tau}\qquad(x\in\mathbb{R}^4).
\end{equation}
Fix a profile $\chi_1\in C^{\infty}(\mathbb{R})$ which is non-increasing, equal to $1$ on
$]-\infty,1]$ and equal to $0$ on $[2,\infty[$. For a radius $\mathcal{R}>0$ the radial cut-off
at radius $\mathcal{R}$ is
\begin{equation}\label{10.8:eq-radial-cutoff-2}
\chi_{\mathcal{R}}(x):=\chi_1\bigl(|x|/\mathcal{R}\bigr)\qquad(x\in\mathbb{R}^4).
\end{equation}
It is a radial function in $C^{\infty}_c(\mathbb{R}^4)$, equal to $1$ on $B(0,\mathcal{R})$ and
equal to $0$ off $B(0,2\mathcal{R})$; it is smooth at the origin because $\chi_1$ is constant on
$]-\infty,1]$, so that $\chi_{\mathcal{R}}$ is constant on $B(0,\mathcal{R})$.

Along the joint diagonal $(K_j,m_j)$ of Notation~\ref{10.8:not-joint-diagonal}, with lattice
spacing $\varepsilon_j=L^{-K_j}$, for $j$ large, namely for those $j$ with
\begin{equation*}
2\mathcal{R}+\varepsilon_j<L^{m_j},
\end{equation*}
the function $\chi_{\mathcal{R}}$ and the vector field $\chi_{\mathcal{R}}\xi$ are read on the
torus $\mathbb{T}_{m_j}$ through representatives: a point of $\mathbb{T}_{m_j}$ is represented
by its unique representative $x\in[-L^{m_j},L^{m_j})^4$, and $\chi_{\mathcal{R}}$ and
$\chi_{\mathcal{R}}\xi$ are evaluated at that $x$.
\end{definition}

\begin{remark}[Necessity of the local identity: anisotropic free fields]
\label{10.8:rem-anisotropic-fields}
The ultraviolet half of the argument, the local identity, cannot be dropped either. Fix $m>0$
and numbers $a_1,\dots,a_4>0$, not all equal. For each of the twenty-four permutations $\pi$ of
$\{1,2,3,4\}$ (we write $\pi\in S_4$) let $\phi_\pi$ be
the centred Gaussian field on $\mathbb{R}^4$ with covariance
\begin{equation*}
\Bigl(-\sum_{\mu=1}^4 a_{\pi(\mu)}\,\partial_\mu^2+m^2\Bigr)^{-1},
\end{equation*}
the twenty-four fields $\phi_\pi$ being independent. For $f\in C_c(\mathbb{R}^4)$ put
\begin{equation*}
\mathcal{O}(f):=\int_{\mathbb{R}^4} f(x)\sum_{\pi\in S_4}{:}\phi_\pi(x)^2{:}\,dx ,
\end{equation*}
each $\phi_\pi$ Wick ordered with respect to its own law. Let $S_n(f_1,\dots,f_n)$, $n\ge 2$, be the
cumulants of $\mathcal{O}(f_1),\dots,\mathcal{O}(f_n)$ for test functions with pairwise disjoint
supports, and $G_\pi$ the covariance kernel of $\phi_\pi$. Then the $S_n$ are invariant under
$\mathbb{R}^4\rtimes W_4$, reflection positive at every axis hyperplane $\{x_\mu=c\}$, bounded uniformly
in the positions of the supports, and clustering exponentially, with a mass gap. Moreover
\begin{equation*}
S_2(f_1,f_2)=2\sum_{\pi\in S_4}\iint f_1(x)f_2(y)\,G_\pi(x-y)^2\,dx\,dy ,
\end{equation*}
and $S_2$ is not $O(4)$-invariant.

Every infrared hypothesis therefore holds, and what fails is the local identity. With the
conserved stress tensor $\sum_\pi T^\pi$ of the model (written out in the proof below), the
antisymmetric part of $\sum_\pi T^\pi_{\nu\rho}$ does not vanish. It is a dimension-four operator
in the $\Lambda^2$ channel, and such an operator exists because the action of $W_4$ permutes the
species.

The model also fails reflection positivity at the twelve diagonal mirrors
$\{x_\mu\pm x_\nu=c\}$.
The Wilson weight has that positivity on $\mathbb{Z}^4$: with $\theta$ the reflection in the
mirror, the plaquettes
crossing $\{x_\mu=x_\nu\}$ enter it as $\operatorname{Re}\operatorname{tr}(\mathsf{A}\,(\theta\mathsf{A})^{\dagger})$,
where $\mathsf{A}$ is the product of the two links on one side of the mirror. A diagonal hyperplane,
however, does not separate the periodic torus on which the lattice measures of Theorem~0 live.
This positivity is therefore not asserted for
the limits of Theorem~0 and is not used in this book.

By Proposition~\ref{10.8:prop-threshold-sharpness}, even reflection positivity at all sixteen
mirrors of $W_4$ does not suffice without an infrared or growth input of its own.
\end{remark}

\begin{proof}
For $z\in\mathbb{R}^4$ put
\begin{equation*}
\mathsf{q}_\pi(z):=\sum_\mu z_\mu^2/a_{\pi(\mu)} .
\end{equation*}
For $u>0$ put
\begin{equation*}
\mathsf{G}(u):=(a_1a_2a_3a_4)^{-1/2}\int_0^\infty(4\pi t)^{-2}e^{-m^2t-u/(4t)}\,dt .
\end{equation*}
The change of variables $y_\mu=z_\mu/\sqrt{a_{\pi(\mu)}}$ turns $-\sum_\mu a_{\pi(\mu)}\partial_\mu^2+m^2$
into $-\Delta_y+m^2$, whose kernel is the same integral without the prefactor. The Jacobian
supplies the factor $(a_1a_2a_3a_4)^{-1/2}$, which does not depend on $\pi$. Hence
$G_\pi(z)=\mathsf{G}(\mathsf{q}_\pi(z))$ for $z\ne 0$. Differentiating under the integral gives
$\mathsf{G}>0$, $\mathsf{G}'<0$ and $\mathsf{G}''>0$. Consequently
\begin{equation*}
(\mathsf{G}^2)''=2(\mathsf{G}')^2+2\mathsf{G}\,\mathsf{G}''>0 .
\end{equation*}

\emph{Invariance.} The covariances depend only on $x-y$, which gives translation invariance. Let
$w\in W_4$ act by $(wz)_\mu=s_\mu z_{\sigma^{-1}(\mu)}$, with signs $s_\mu\in\{\pm1\}$ and a
permutation $\sigma$. Then
\begin{equation*}
\mathsf{q}_\pi(wz)=\sum_\nu z_\nu^2/a_{\pi(\sigma(\nu))}=\mathsf{q}_{\pi\circ\sigma}(z).
\end{equation*}
So the family $(\phi_\pi\circ w)_\pi$ has the law of $(\phi_{\pi\circ\sigma})_\pi$. Relabelling the sum
over $\pi$ shows that $\sum_\pi{:}\phi_\pi^2{:}\circ w$ has the law of $\sum_\pi{:}\phi_\pi^2{:}$
jointly in all points, and hence the $S_n$ are $W_4$-invariant.

\emph{Reflection positivity at axis hyperplanes.} The same change of variables maps $\phi_\pi$ to
the isotropic free field of mass $m$. It maps $\{x_\mu=c\}$ to
$\{y_\mu=c/\sqrt{a_{\pi(\mu)}}\}$, and the reflection in the one to the reflection in the other.
Each $\phi_\pi$ is therefore reflection positive at $\{x_\mu=c\}$. Let $\theta$ denote the
reflection in $\{x_\mu=c\}$. For products
$F=\prod_\pi F_\pi$, with $F_\pi$ a functional of $\phi_\pi$ on one side, independence gives
$\langle\theta(\bar F)F\rangle=\prod_\pi\langle\theta(\bar F_\pi)F_\pi\rangle$. The Gram matrix of a
finite family of such products is the entrywise product of positive semidefinite matrices, which
is positive semidefinite. The moments of the $\mathcal{O}(f)$ with $f$ supported on one side are
limits of linear combinations of such products, so they are reflection positive, and so are the
centred moments.

\emph{The two-point function.} Mixed cumulants of independent families vanish. Wick's theorem
gives $\langle{:}\phi_\pi(x)^2{:}\,{:}\phi_\pi(y)^2{:}\rangle=2G_\pi(x-y)^2$, and the integral
converges because the supports are compact and at positive distance. Take $x=re_1$. Then
$\mathsf{q}_\pi(x)=r^2/a_{\pi(1)}$, and each value $a_j$ occurs for six permutations, so
$2\sum_\pi G_\pi(x)^2=12\sum_j\mathsf{G}(r^2/a_j)^2$. Take next $x'=\tfrac r2(1,1,1,1)$, which has the
same norm. Then $\mathsf{q}_\pi(x')=\tfrac{r^2}4\sum_j a_j^{-1}$ for every $\pi$, and
\begin{equation*}
2\sum_\pi G_\pi(x')^2=48\,\mathsf{G}\Bigl(\frac{r^2}{4}\sum_j a_j^{-1}\Bigr)^2 .
\end{equation*}
The function
$\mathsf{G}^2$ is strictly convex and the $a_j^{-1}$ are not all equal, so Jensen's inequality with
weights $\tfrac14$ gives
\begin{equation*}
12\sum_j\mathsf{G}(r^2/a_j)^2>48\,\mathsf{G}\Bigl(\frac{r^2}{4}\sum_j a_j^{-1}\Bigr)^2
\qquad(r>0).
\end{equation*}
The kernel of $S_2$ is continuous off the origin and takes different values at two points of the
same norm. Testing it with $f_1,f_2$ concentrated near $0$ and near $x$, respectively near $x'$,
shows that $S_2$ is not $O(4)$-invariant.

\emph{Bounds and clustering.} Let $\mathsf{d}>0$ be the least distance between two of the supports
of $f_1,\dots,f_n$. For $n\ge2$, Wick's theorem writes the cumulant of
${:}\phi_\pi(x_1)^2{:},\dots,{:}\phi_\pi(x_n)^2{:}$ as $2^{n-1}$ times the sum, over the $(n-1)!$
cyclic permutations $\sigma$, of $\prod_iG_\pi(x_i-x_{\sigma(i)})$. Put $a_+:=\max_ja_j$. Then
$\mathsf{q}_\pi(z)\ge|z|^2/a_+$, and $\mathsf{G}$ is decreasing. Since $m^2t+u/(4t)\ge m\sqrt u$,
\begin{equation*}
m^2t+\frac u{4t}\ge\frac{m\sqrt u}2+\frac12\Bigl(m^2t+\frac u{4t}\Bigr).
\end{equation*}
Hence $\mathsf{G}(u)\le e^{-m\sqrt u/2}\,\mathsf{G}_{1/2}(u)$, where $\mathsf{G}_{1/2}(u)$ is the
finite integral, decreasing in $u$, obtained from the definition of $\mathsf{G}$ by replacing
$e^{-m^2t-u/(4t)}$ with $e^{-m^2t/2-u/(8t)}$. Thus $|z|\ge\mathsf{d}$ implies
$G_\pi(z)\le C\,e^{-m|z|/(2\sqrt{a_+})}$ with $C=C(m,a_1,\dots,a_4,\mathsf{d})$. It follows that
\begin{equation*}
|S_n|\le 24\cdot2^{n-1}(n-1)!\,C^n\prod_i\|f_i\|_{L^1},
\end{equation*}
uniformly in the positions of the
supports. If the indices split into two non-empty groups whose supports are at distance
$R\ge\mathsf{d}$, every cycle through all $n$ points has at least two edges joining the groups, and
\begin{equation*}
|S_n|\le 24\cdot2^{n-1}(n-1)!\,C^n e^{-mR/\sqrt{a_+}}\prod_i\|f_i\|_{L^1}.
\end{equation*}
For the mass gap, take the direction $\mu$ as time. The Fourier-transformed covariance of
$\phi_\pi$ has its poles in $p_\mu$ at $\pm iE_\pi(p)$, where
\begin{equation*}
E_\pi(p)=\Bigl(\bigl(\textstyle\sum_{\nu\ne\mu}a_{\pi(\nu)}p_\nu^2+m^2\bigr)/a_{\pi(\mu)}\Bigr)^{1/2}
\ge m/\sqrt{a_+},
\end{equation*}
so every state orthogonal to the vacuum has energy at least $m/\sqrt{a_+}$.

\emph{The stress tensor.} Put
\begin{equation*}
T^\pi_{\nu\rho}:=a_{\pi(\nu)}\partial_\nu\phi_\pi\,\partial_\rho\phi_\pi
-\delta_{\nu\rho}\Bigl(\tfrac12\sum_\sigma a_{\pi(\sigma)}(\partial_\sigma\phi_\pi)^2
+\tfrac12m^2\phi_\pi^2\Bigr).
\end{equation*}
The field equation is $\sum_\nu a_{\pi(\nu)}\partial_\nu^2\phi_\pi=m^2\phi_\pi$. Using it,
$\sum_\nu\partial_\nu\bigl(a_{\pi(\nu)}\partial_\nu\phi_\pi\partial_\rho\phi_\pi\bigr)$ equals
\begin{equation*}
m^2\phi_\pi\partial_\rho\phi_\pi
+\sum_\nu a_{\pi(\nu)}\partial_\nu\phi_\pi\,\partial_\nu\partial_\rho\phi_\pi
=\partial_\rho\Bigl(\tfrac12m^2\phi_\pi^2+\tfrac12\sum_\nu a_{\pi(\nu)}(\partial_\nu\phi_\pi)^2\Bigr).
\end{equation*}
This is the divergence of the $\delta_{\nu\rho}$-term, so $\sum_\nu\partial_\nu T^\pi_{\nu\rho}=0$. The
antisymmetric part is
\begin{equation*}
\sum_\pi\bigl(T^\pi_{\nu\rho}-T^\pi_{\rho\nu}\bigr)
=\sum_\pi\bigl(a_{\pi(\nu)}-a_{\pi(\rho)}\bigr)\partial_\nu\phi_\pi\,\partial_\rho\phi_\pi .
\end{equation*}
For $\nu\ne\rho$ some $\pi$ has $a_{\pi(\nu)}\ne a_{\pi(\rho)}$, because the $a_j$ are not all
equal. The species are independent, so this dimension-four operator is not zero.
\end{proof}

\begin{definition}[The shell functional of the rotational Ward identity]\label{10.8:def-shell-functional}
Let $\omega_0\in\mathfrak{so}(4)$, $\omega_0\neq 0$, the rotation field $\xi(x)=\omega_0 x$ and the radial
cut-offs $\chi_{\mathcal{R}}$ be as in Definition~\ref{10.8:def-radial-cutoff}. Let $a=L^{-K}$ be the
lattice spacing, and let $a_p=1-\operatorname{Re}\operatorname{tr}U(\partial p)$ be the plaquette action
density. Let $x(p)$ be the barycentre of the plaquette $p$, and let $x_\ell$ be the point attached to
the bond $\ell$ in the densities $\mathfrak{d}_{\ell,p,\rho}$ of the rotational Ward identity,
Theorem~\ref{10.1:thm-rotational-identity}. For a compactly supported smooth
function $\chi$ on $\mathbb{R}^4$ (below, $\chi=\chi_{\mathcal{R}}$), read on $\mathbb{T}_m$ together
with $\chi\xi$ through representatives as in Definition~\ref{10.8:def-radial-cutoff}, and with
$\sum_p$ running over the plaquettes of the lattice of spacing $a$ in $\mathbb{T}_m$, the
\emph{shell functional} of the rotational Ward identity for the deformation field $\chi\xi$ is
\begin{equation}\label{10.8:eq-shell-functional-1}
\mathfrak{B}^{\mathrm{sh}}_{\chi}
:=g_0^{-2}\sum_{p}a_p\,\bigl(\mathsf{L}^{a}_{\xi}\chi\bigr)(x(p))
-g_0^{-2}\sum_{p}\mathfrak{h}_{\chi}(p)+\operatorname{div}V_{\chi\xi}.
\end{equation}
Here $\mathsf{L}^{a}_{\xi}$ is the lattice Lie derivative along $\xi$ at spacing $a$, and
$\mathfrak{h}_{\chi}(p)$ is the plaquette shell density. Further, $V_{\chi\xi}$ is the vector field whose
lattice divergence $\operatorname{div}V_{\chi\xi}$ is the Jacobian term. These three objects are those
of the rotational Ward identity, Theorem~\ref{10.1:thm-rotational-identity}, and are not redefined here.

The following properties of $\mathfrak{B}^{\mathrm{sh}}_{\chi}$ are not part of the definition and are
used below: $(\alpha)$ is part of the rotational Ward identity, and $(\beta)$ and $(\gamma)$ are the
shell-class form of the shell functional.
\begin{itemize}
\item[$(\alpha)$] $\mathfrak{B}^{\mathrm{sh}}_{\chi}$ is the explicit shell term of the exact lattice
rotational Ward identity for the deformation field $\chi\xi$.
\item[$(\beta)$] For a family of
weights $k=(k_{\ell,p,\rho},k^0_p)$, the \emph{shell class} is
\begin{equation}\label{10.8:eq-shell-functional-2}
\mathfrak{T}^{\mathrm{sh}}(k)
:=g_0^{-2}\sum_{p,\ \ell\subset\partial p,\ \rho}k_{\ell,p,\rho}\,\mathfrak{d}_{\ell,p,\rho}
+g_0^{-2}\sum_{p}k^0_p\,a_p ,
\end{equation}
where $\mathfrak{d}_{\ell,p,\rho}$ are the lattice dimension-four tensor densities of the rotational
Ward identity. Thus
$\mathfrak{T}^{\mathrm{sh}}(k)$ is a smeared combination of components of these densities; its
classical counterpart is, schematically, $\int k\cdot g_0^{-2}\vec F_{\rho\mu}\cdot\vec F_{\sigma\mu}$,
with $\vec F_{\rho\mu}$ the curvature written as the vector of its colour components.
The weights $k_{\chi}$ of the deformation field $\chi\xi$ are
\begin{equation}\label{10.8:eq-shell-functional-3}
\begin{aligned}
k_{\ell,p,\rho}&=-a^{-1}\bigl(\chi(x_\ell)-\chi(x(p))\bigr)\,\xi^{\rho}(x_\ell),\\
k^0_p&=\bigl(\mathsf{L}^{a}_{\xi}\chi\bigr)(x(p))-\tfrac{3}{2}\,g_0^{2}\bigl[\xi^{\rho}\nabla^{c}_{\rho}\chi
+\xi^{\mu}\nabla^{c}_{\mu}\chi\bigr](x(p)),
\end{aligned}
\end{equation}
where $\nabla^{c}$ is the centred lattice gradient, and
$\mathfrak{B}^{\mathrm{sh}}_{\chi}=\mathfrak{T}^{\mathrm{sh}}(k_{\chi})$.
\item[$(\gamma)$] Let $\chi=\chi_{\mathcal{R}}$. Then the weights $k_{\chi_{\mathcal{R}}}$ are supported
in the shell
\begin{equation*}
\{\,x:\ \mathcal{R}-a\le|x|\le 2\mathcal{R}+a\,\}.
\end{equation*}
With a constant $C$ independent of $\mathcal{R}\ge 1$, they satisfy $|k|\le C|\omega_0|$ and
$\operatorname{Lip}k\le C|\omega_0|/\mathcal{R}$.
Hence their $\|\cdot\|_{\sharp}$-norm is $O(|\omega_0|)$, uniformly in $\mathcal{R}\ge 1$.
\end{itemize}

At each cutoff of the joint diagonal (see Notation~\ref{10.8:not-joint-diagonal}), suppose the shell fits,
that is, $2\mathcal{R}+\varepsilon_j<L^{m_j}$, where $\varepsilon_j=L^{-K_j}$ is the spacing at the
$j$-th cutoff (so that $a=\varepsilon_j$ there). Then the fields $\chi_{\mathcal{R}}$ and
$\chi_{\mathcal{R}}\xi$ are
read on $\mathbb{T}_{m_j}$ as in Definition~\ref{10.8:def-radial-cutoff}, and
$\mathfrak{B}^{\mathrm{sh}}_{\chi_{\mathcal{R}}}$ is a bounded gauge-invariant function of the bare
field on $\mathbb{T}_{m_j}$.
\end{definition}

The functional $\mathfrak{B}^{\mathrm{sh}}_{\chi}$ coincides with the total dimension-four annulus
functional $\mathfrak{F}_{\chi}=B^4_{\chi}+J^4_{\chi}$ of the conclusion step at cutoff $K$, less its
bulk defect under the full form of the dimension-six coefficient bound for the tensor densities
$\mathfrak{d}_{\ell,p,\rho}$. The bulk defect and this identification are given with the conclusion
step, and the full form of the dimension-six coefficient bound is the one stated with the conclusion
step.

\begin{definition}[The coupling-weighted stress tensor and its admixtures]
\label{10.8:def-stress-admixtures}
Let $x_0 = g_0^{-2}$ denote the inverse-coupling weight, that is, the factor $g_0^{-2}$ standing in
front of the shell class $\mathfrak{T}^{\mathrm{sh}}(k)$ in the shell functional of the rotational
Ward identity (Definition~\ref{10.8:def-shell-functional}).
\begin{enumerate}
\item For $\nu,\sigma \in \{1,\dots,4\}$, $\Theta_{\nu\sigma}$ denotes the lattice stress-tensor
density of dimension four, transforming as a tensor under $W_4$, which carries the weight $x_0$ and
is traceless at tree level. Tracelessness at tree level means that the combination carrying the
weight $x_0$ has vanishing $W_4$-trace classically. The pure trace $\sum_{\nu}\Theta_{\nu\nu}$
never enters with the weight $x_0$.
\item The \emph{admixtures} of $\Theta$ are the symmetric tensor densities of dimension four which
carry no factor $x_0$ and enter the shell functional with coefficients of order one. In the shell
functional they are the following two terms. First, the term
\begin{equation*}
-\tfrac{3}{2}\,g_0^{2}\bigl[\xi^{\rho}\nabla^{c}_{\rho}\chi + \xi^{\mu}\nabla^{c}_{\mu}\chi\bigr]
\end{equation*}
of the weight $k^0_p$ of Definition~\ref{10.8:def-shell-functional}, multiplied by the factor
$g_0^{-2}$ in front of $\mathfrak{T}^{\mathrm{sh}}$;
the product is the bracket above times $-\tfrac{3}{2}$, a coefficient of $a_p$ of order one.
Second, the diagonal tensor of the Jacobian term
$\operatorname{div} V_{\chi\xi}$.
\item Whenever decay is required of the components of $\Theta$, it is required also of its
admixtures.
\item At lattice spacing $\varepsilon$, the density $\Theta^{\varepsilon}_{\nu\sigma}(x)$
realising $\Theta_{\nu\sigma}$ is a real, bounded, gauge-invariant polynomial in the link variables
lying within distance $r_{\Theta}\varepsilon$ of the site $x$, where $r_{\Theta}$ is a fixed range
measured in lattice spacings. It is the same polynomial at every site and
contains the factor $g_0^{-2}$. For a test function $g$ its smearing is
\begin{equation}
\label{10.8:eq-stress-admixtures-1}
\Theta^{\varepsilon}_{\nu\sigma}(g) := \sum_{x} \varepsilon^{4}\, g(x)\,
\Theta^{\varepsilon}_{\nu\sigma}(x),
\end{equation}
where the sum runs over the sites of the lattice of spacing $\varepsilon$.
\end{enumerate}
The identification of $\Theta^{\varepsilon}_{\nu\sigma}$ with the corresponding combinations of
the densities $\mathfrak{d}_{\ell,p,\rho}$ and $a_p$ of the shell class, and of the annulus
functional $\mathfrak{F}_{\chi}$ of the conclusion step with $\mathfrak{B}^{\mathrm{sh}}_{\chi}$,
is used beyond leading coefficients only as an explicit hypothesis, stated in the conclusion step.
\end{definition}

\begin{definition}[the far-sphere flux hypothesis $(\mathrm{IR}\text{-}\mathrm{fl})_{\omega_0}$]
\label{10.8:def-clustering-hypothesis}
Let the joint diagonal $(K_j,m_j)_{j}$ of clause (ii-$\mathbb{R}^4$) be fixed, as in
Notation~\ref{10.8:not-joint-diagonal}. Let
$\chi_1$ be the profile of Definition~\ref{10.8:def-radial-cutoff}, and put
$\chi_{\mathcal{R}}(x) := \chi_1(|x|/\mathcal{R})$ for $\mathcal{R} > 0$. Write
$\varepsilon_j = L^{-K_j}$ for the lattice spacing along the diagonal. For
$h \in C^\infty_c(\mathbb{R}^4)$, $h^{(m)}$ denotes the periodic lift of $h$ to
$\mathbb{T}_m$, and $\mathcal{O}(h^{(m)})$ is the smeared curvature observable built from it.
As in Notation~\ref{10.8:not-joint-diagonal},
$\langle A;B_1;\dots;B_n\rangle_{K,m}$ denotes the truncated expectation, that is, the joint
cumulant of $A,B_1,\dots,B_n$ with respect to the Wilson measure $\mu_{K,m}$ at
cutoff $K$ on $\mathbb{T}_m$.

\emph{Hypothesis} $(\mathrm{IR}\text{-}\mathrm{fl})_{\omega_0}$, the far-sphere flux
hypothesis of clause (vi) (the flux of the lattice rotation current through far spheres
vanishes along one sequence of radii), is the following
assertion. There is one real antisymmetric $4\times 4$ matrix $\omega_0 \neq 0$ with the
following property. Let the shell functional $\mathfrak{B}^{\mathrm{sh}}_{\chi}$ of
Definition~\ref{10.8:def-shell-functional} be formed with the rotation field
$\xi(x) := \omega_0 x$. Then for every $n \ge 2$ and every tuple
$\vec f = (f_1,\dots,f_n) \in \bigl(C^\infty_c(\mathbb{R}^4)\bigr)^n$
whose members have pairwise disjoint supports,
\begin{equation}\label{10.8:eq-clustering-hypothesis-a}
\liminf_{\mathcal{R}\to\infty}\;\limsup_{j\to\infty}\;
\Bigl|\bigl\langle \mathfrak{B}^{\mathrm{sh}}_{\chi_{\mathcal{R}}};
\mathcal{O}\bigl(f_1^{(m_j)}\bigr);\dots;\mathcal{O}\bigl(f_n^{(m_j)}\bigr)
\bigr\rangle_{K_j,m_j}\Bigr| = 0 .
\end{equation}

The inner expression in \eqref{10.8:eq-clustering-hypothesis-a} is defined for every $j$
with
\begin{equation*}
2\mathcal{R} + \varepsilon_j < L^{m_j},
\end{equation*}
that is, for each fixed $\mathcal{R}$, for all sufficiently large $j$; so the
limit superior over $j$ is meaningful.

The order of the quantifiers is part of the hypothesis. The matrix $\omega_0$ is
existentially quantified and is chosen first; it is then one and the same for all $n$ and all
tuples. The
assertion is required for all $n \ge 2$ and for all tuples of $C^\infty_c$ functions with
pairwise disjoint supports. The cut-offs $\chi_{\mathcal{R}}$ form the radial family
generated by the fixed profile $\chi_1$. In \eqref{10.8:eq-clustering-hypothesis-a} the limit
superior over $j$ is taken inside, at fixed $\mathcal{R}$, and the limit inferior over
$\mathcal{R}$ is taken outside.
\end{definition}

Beside $(\mathrm{IR}\text{-}\mathrm{fl})_{\omega_0}$ we display its stress-tensor decay form,
a condition sufficient for it: the connected correlation of one insertion of the
stress-tensor density $\Theta_{\nu\sigma}$ at a point $x$ of $\mathbb{T}_{m_j}$ far from the
supports of $f_1,\dots,f_n$ (with $|x|$ read through the representative of $x$ in
$[-L^{m_j},L^{m_j})^4$) with the observables
$\mathcal{O}(f_1^{(m_j)}),\dots,\mathcal{O}(f_n^{(m_j)})$ is bounded, for all $j$ and all
such $x$, by one function of $|x|$, independent of $j$, which is $o(|x|^{-4})$ as
$|x|\to\infty$. The smearing and normalisation of the insertion, and the proof that this
condition implies \eqref{10.8:eq-clustering-hypothesis-a}, are given in the statement of this
chapter that derives the far-sphere flux hypothesis from the stress-tensor decay form.

\begin{remark}[One generator suffices; necessity given ultraviolet bounds]
\label{10.8:rem-one-generator}
The far-sphere flux hypothesis $(\mathrm{IR}\text{-}\mathrm{fl})_{\omega_0}$ of clause (vi)
(Definition~\ref{10.8:def-clustering-hypothesis}), which says that the flux of the lattice
rotation current through far spheres vanishes along one sequence of radii, has the following
three features.
\begin{enumerate}
\item[(a)] \emph{Nature of the hypothesis.} It is an infrared hypothesis on the limit
state. For each fixed $\mathcal{R}$ the shell term in \eqref{10.8:eq-clustering-hypothesis-a}
is a smeared insertion of dimension four, placed at distance of order $\mathcal{R}$ from the
sources, and its limit superior in $j$ is taken first. No uniformity in $K$, or in $m$, is
asked for beyond the iterated order of the two limits: first in $j$, then in
$\mathcal{R}$. The natural setting of the hypothesis is the effective theory on the lattice
of the last renormalisation step, beyond the scale at which the renormalisation-group flow
stops.
\item[(b)] \emph{One generator suffices.} Under the hypotheses of the conclusion step, namely
clauses (ii-$\mathbb{R}^4$) and (iii-b) and the ultraviolet bounds entering that step, the
rotation defect $D_{\omega}$ is linear in the generator $\omega\in\mathfrak{so}(4)$ and
$W_4$-covariant, and the dichotomy of the conclusion step says that the
stabiliser in $O(4)$ of every limit on $\mathbb{R}^4$ is either a finite group containing
$W_4$ or all of $O(4)$. Vanishing of the defect for the single generator $\omega_0$ makes the
stabiliser infinite; by the dichotomy it is then $O(4)$.
\item[(c)] \emph{Given the ultraviolet bounds, the hypothesis is necessary as well as
sufficient.} Assume clauses (ii-$\mathbb{R}^4$) and (iii-b) and the ultraviolet bounds
entering the conclusion step. Then, as shown in the conclusion step, for every tuple
$\vec f$ as in Definition~\ref{10.8:def-clustering-hypothesis} the defect
$D_{\omega}(\vec f)$ equals the limiting flux of the shell insertion
$\mathfrak{T}^{\mathrm{sh}}(k_{\chi_{\mathcal{R}}})$, the smeared dimension-four density with
the weight $k_{\chi_{\mathcal{R}}}$ of the deformation field $\chi_{\mathcal{R}}\xi_{\omega}$,
through any shell surrounding the supports of $\vec f$: for every $\mathcal{R}$,
\begin{equation}\label{10.8:eq-one-generator-1}
\lim_{j\to\infty}
\bigl\langle \mathfrak{T}^{\mathrm{sh}}(k_{\chi_{\mathcal{R}}});
\mathcal{O}(\vec f)\bigr\rangle_{j}
= D_{\omega}(\vec f),
\end{equation}
where $\mathcal{O}(\vec f)$ abbreviates the $n$ slots
$\mathcal{O}(f_1^{(m_j)});\dots;\mathcal{O}(f_n^{(m_j)})$ and $\langle\cdot\rangle_j$ is the
truncated expectation along the diagonal $(K_j,m_j)$ of
Definition~\ref{10.8:def-clustering-hypothesis}.
Consequently, if the limit is $O(4)$-invariant, then $D_{\omega}=0$ for every $\omega$, and
the limit inferior in \eqref{10.8:eq-clustering-hypothesis-a} vanishes. If the limit is not
$O(4)$-invariant, then for every generator $\omega\neq 0$ the hypothesis with $\omega$ in
place of $\omega_0$ in the shell functional fails: there is a tuple $\vec f$ with
$D_{\omega}(\vec f)\neq 0$, and for this tuple the limit inferior in
\eqref{10.8:eq-clustering-hypothesis-a} equals $|D_{\omega}(\vec f)|>0$.
\end{enumerate}
The radial shell functional $\mathfrak{B}^{\mathrm{sh}}_{\chi_{\mathcal{R}}}$, in which the
hypothesis is formulated, suffices for clause (vi) once the one-loop anisotropy inequality
holds with exponent $\kappa=1$ (Corollary~\ref{10.6:cor-anisotropy-inequality}), the radial
form requiring only $\kappa>\tfrac12$ by the
ultraviolet thresholds displayed with that inequality.
\end{remark}

\begin{definition}[Pointwise stress decay faster than $|x|^{-4}$]\label{10.8:def-pointwise-decay}
Let $\Theta_{\nu\sigma}$, $\nu,\sigma\in\{1,\dots,4\}$, be the lattice stress-tensor density of
dimension four weighted by $x_0=g_0^{-2}$, and let its $x_0$-free symmetric dimension-four
admixtures be as in Definition~\ref{10.8:def-stress-admixtures}. Define $\mathcal{K}_{\Theta}$ to be
the class of components consisting of the following.
\begin{itemize}
\item The off-diagonal components $\Theta_{\nu\sigma}$, $\nu\neq\sigma$, of the $x_0$-weighted
stress tensor. Each of them is traceless by itself.
\item The traceless diagonal combinations $\sum_{\nu}c_{\nu}\Theta_{\nu\nu}$, where
$c_1,\dots,c_4$ are real constants with $\sum_{\nu}c_{\nu}=0$ and $\max_{\nu}|c_{\nu}|\le1$; they
are spanned by the diagonal differences $\Theta_{\nu\nu}-\Theta_{44}$, $\nu\le3$. The
$W_4$-trace $\sum_{\nu}\Theta_{\nu\nu}$, carrying the weight $x_0$, is not itself a member of the
class.
\item The $x_0$-free symmetric dimension-four admixtures of
Definition~\ref{10.8:def-stress-admixtures}.
\end{itemize}
For a member $\Phi$ of $\mathcal{K}_{\Theta}$ and a real weight $h$ on $\mathbb{R}^4$, we write
$\Phi(h)$ for that member smeared with $h$, as in
Definition~\ref{10.8:def-stress-admixtures}; for a traceless diagonal combination this is
$\sum_{\nu}c_{\nu}\Theta_{\nu\nu}(h)$, with $\Theta_{\nu\nu}(\cdot)$ the smeared diagonal
component. We write
$\langle\,\cdot\,;\dots;\,\cdot\,\rangle_{K,m}$
for the truncated expectation (joint cumulant) with respect to $\mu_{K,m}$, as fixed
in Notation~\ref{10.8:not-joint-diagonal}.

The family $(\mu_{K_j,m_j})_j$ of measures along the joint diagonal has \emph{pointwise stress
decay faster than $|x|^{-4}$} if the
following holds. For every $n\ge 2$ and every tuple $\vec f=(f_1,\dots,f_n)$ of test functions
with pairwise disjoint supports, there are $\delta>0$ and $C(\vec f)<\infty$ with the property
below. Let $\Delta$ be any closed axis-parallel cube of side $1$ at distance
$D\ge 1$ from $\bigcup_i\operatorname{supp}f_i$. Let $h$ be any real Lipschitz weight supported
in $\Delta$, and let $\Phi$ be any member of $\mathcal{K}_{\Theta}$. Then
\begin{equation}\label{10.8:eq-pointwise-decay-a}
\sup_{j}\,\bigl|\langle\Phi(h);\mathcal{O}(f_1);\dots;\mathcal{O}(f_n)\rangle_{K_j,m_j}\bigr|
\le C(\vec f)\,\|h\|_{\sharp}\,D^{-4-\delta}.
\end{equation}
Here $(K_j,m_j)_j$ runs along the joint diagonal of Notation~\ref{10.8:not-joint-diagonal}, so that the
supremum over the cutoff $K=K_j$ is a supremum over $j$, and the bound holds uniformly in $K$ and
in $m$ along it. At the $j$-th member of the diagonal, $D$ is the distance in the quotient metric of the
torus $\mathbb{T}_{m_j}$. The norm $\|h\|_{\sharp}$ is $\max\{\|h\|_{\infty},40M\operatorname{Lip}h\}$.

The condition can also be written with the shell class. For a real Lipschitz weight $k$ supported
in $\Delta$ (here $k$ is a weight, not a level index), let $\mathfrak{T}^{\mathrm{sh}}(k)$ denote a
lattice dimension-four tensor density of the shell class, that is, of the class of densities
entering the shell functional of Definition~\ref{10.8:def-shell-functional}, smeared with the
single weight $k$, so that $\|k\|_{\sharp}$ is the sharp norm of that one weight. Here the
densities of the shell class are read as the members of $\mathcal{K}_{\Theta}$. Write
$\langle\,\cdot\,\rangle_j=\langle\,\cdot\,\rangle_{K_j,m_j}$ along the joint diagonal. Then
\eqref{10.8:eq-pointwise-decay-a} reads
\begin{equation}\label{10.8:eq-pointwise-decay-1}
\sup_{j}\,\bigl|\langle\mathfrak{T}^{\mathrm{sh}}(k);\mathcal{O}(f_1);\dots;\mathcal{O}(f_n)\rangle_{j}\bigr|
\le C(\vec f)\,\|k\|_{\sharp}\,D^{-4-\delta},
\end{equation}
for every such $k$, with the same $\delta$, $C(\vec f)$, $\Delta$ and $D$.
\end{definition}

Condition \eqref{10.8:eq-pointwise-decay-a} is a hypothesis, not an assertion. Clause (vi) of
Theorem~0 carries as its hypothesis the far-sphere flux hypothesis
$(\mathrm{IR}\text{-}\mathrm{fl})_{\omega_0}$: the
flux of the lattice rotation current through far spheres vanishes along one sequence of radii.
In that hypothesis, at each fixed radius the limit superior over the cutoff along the joint
diagonal is taken first, and only then is the radius sent to infinity along one sequence; no
bound uniform in the cutoff is part of the hypothesis itself.
Beside it stands a sufficient condition in stress-tensor form: the connected correlation of one
far-away stress-tensor insertion with the observables decays faster than $|x|^{-4}$, uniformly in
the cutoff. Condition \eqref{10.8:eq-pointwise-decay-a} is this sufficient condition read
pointwise and restricted to the class $\mathcal{K}_{\Theta}$.

We now explain why the radial shell requires decay for $\mathcal{K}_{\Theta}$ rather than for all
sixteen components $\Theta_{\nu\sigma}$. Let $\xi(x)=\omega_0x$ be the linear rotation vector
field of the generator $\omega_0\in\mathfrak{so}(4)$, so that
$\xi^{\sigma}(x)=\sum_{\rho}(\omega_0)_{\sigma\rho}x_{\rho}$. Let $\chi_1$ be the fixed radial
profile, and let $\chi_{\mathcal{R}}(x)=\chi_1(|x|/\mathcal{R})$ be the radial cut-off at radius
$\mathcal{R}$.

The deformation field is $\chi_{\mathcal{R}}\xi$. In the transcription of the annulus functional
$\mathfrak{F}_{\chi}$ of the conclusion step, its strain weight is
\begin{equation}\label{10.8:eq-pointwise-decay-2}
b^{\chi_{\mathcal{R}}}_{\nu\sigma}(x)
=\sum_{\rho}(\omega_0)_{\rho\sigma}\,x_{\rho}\,(\partial_{\nu}\chi_{\mathcal{R}})(x)
=-\xi^{\sigma}(x)\,\partial_{\nu}\chi_{\mathcal{R}}(x).
\end{equation}
The second equality in \eqref{10.8:eq-pointwise-decay-2} uses the antisymmetry
$(\omega_0)_{\rho\sigma}=-(\omega_0)_{\sigma\rho}$, which gives
$\sum_{\rho}(\omega_0)_{\rho\sigma}x_{\rho}=-\xi^{\sigma}(x)$. For $x\neq0$ the chain rule gives
\begin{equation*}
\partial_{\nu}\chi_{\mathcal{R}}(x)=\chi_1'\bigl(|x|/\mathcal{R}\bigr)\,\frac{x_{\nu}}{\mathcal{R}|x|}.
\end{equation*}
Hence each diagonal weight equals
\begin{equation*}
b^{\chi_{\mathcal{R}}}_{\nu\nu}(x)=-\xi^{\nu}(x)\,x_{\nu}\,\frac{\chi_1'(|x|/\mathcal{R})}{\mathcal{R}|x|}.
\end{equation*}
These weights do not vanish individually. Their sum, however, is
\begin{equation}\label{10.8:eq-pointwise-decay-3}
\begin{aligned}
\sum_{\nu}b^{\chi_{\mathcal{R}}}_{\nu\nu}(x)
&=-\frac{\chi_1'(|x|/\mathcal{R})}{\mathcal{R}|x|}\sum_{\nu}x_{\nu}\,\xi^{\nu}(x)\\
&=-\frac{\chi_1'(|x|/\mathcal{R})}{\mathcal{R}|x|}\;x\cdot\omega_0x \;=\;0.
\end{aligned}
\end{equation}
The last equality holds because
$x\cdot\omega_0x=\sum_{\nu,\rho}x_{\nu}(\omega_0)_{\nu\rho}x_{\rho}$ is a quadratic form with
antisymmetric matrix, and so vanishes. Thus the radial weights are diagonal-traceless pointwise.

At each point, the diagonal part of the strain weight therefore pairs with the diagonal
components of the stress tensor as
$\sum_{\nu}b^{\chi_{\mathcal{R}}}_{\nu\nu}\Theta_{\nu\nu}$. By \eqref{10.8:eq-pointwise-decay-3},
$b^{\chi_{\mathcal{R}}}_{44}=-\sum_{\nu\le3}b^{\chi_{\mathcal{R}}}_{\nu\nu}$, and hence
\begin{equation*}
\sum_{\nu=1}^{4}b^{\chi_{\mathcal{R}}}_{\nu\nu}\Theta_{\nu\nu}
=\sum_{\nu=1}^{3}b^{\chi_{\mathcal{R}}}_{\nu\nu}\,\bigl(\Theta_{\nu\nu}-\Theta_{44}\bigr).
\end{equation*}
This is a sum of three diagonal differences, each a traceless diagonal combination in the sense
of the second item above, smeared with the weight $b^{\chi_{\mathcal{R}}}_{\nu\nu}$. The
$W_4$-trace with the large weight $x_0$ never occurs.

The conclusion differs between the radial shell and the axis tubes.
\begin{itemize}
\item On the radial shell, decay of the off-diagonal components, of the traceless diagonal
combinations (the diagonal differences) and of the $x_0$-free admixtures suffices. This is the
class $\mathcal{K}_{\Theta}$, and it is what \eqref{10.8:eq-pointwise-decay-a} asks.
\item On the axis tubes, decay of the off-diagonal components alone suffices, with no $x_0$-free
admixture. There the axis-tube deformation field $\eta^{\natural}_{\mathcal{R}}$ of
Theorem~\ref{10.5:thm-reduction-without-shell} has identically vanishing diagonal strain (compare
Lemma~\ref{10.5:lem-no-diagonal-strain}) and no Jacobian term.
\end{itemize}

\begin{lemma}[Pointwise decay implies the far-sphere flux hypothesis]\label{10.8:lem-pointwise-implies}
Let the joint diagonal $(K_j,m_j)$ and the profile $\chi_1$, with
$\chi_{\mathcal{R}}(x)=\chi_1(|x|/\mathcal{R})$, be as in the far-sphere flux hypothesis
$(\mathrm{IR}\text{-}\mathrm{fl})_{\omega_0}$ of Definition~\ref{10.8:def-clustering-hypothesis},
which states that the flux of the lattice rotation current through far spheres vanishes along one
sequence of radii. Let $\omega_0\in\mathfrak{so}(4)$ be
any non-zero real antisymmetric matrix, $\xi(x)=\omega_0x$, and write
$\varepsilon_j=L^{-K_j}$. Assume the following.
\begin{enumerate}
\item[(a)] Pointwise stress decay faster than $|x|^{-4}$
(Definition~\ref{10.8:def-pointwise-decay}), in its shell-class form: for every $n\ge 2$ and
every tuple $\vec f=(f_1,\dots,f_n)\in(C^\infty_c(\mathbb{R}^4))^n$ with pairwise disjoint
supports there are $\delta>0$ and $C(\vec f)<\infty$, independent of $j$, such that for every $j$,
every closed axis-parallel unit cube $\Delta$ at distance $D\ge 1$ on $\mathbb{T}_{m_j}$ from
$\bigcup_i\operatorname{supp}f_i^{(m_j)}$, and every real Lipschitz weight $k$ supported in
$\Delta$ for which $\mathfrak{T}^{\mathrm{sh}}(k)$ is a sum of components of the class
$\mathcal{K}_{\Theta}$ of Definition~\ref{10.8:def-pointwise-decay},
\begin{equation}\label{10.8:eq-pointwise-implies-1}
\bigl|\langle\mathfrak{T}^{\mathrm{sh}}(k);\mathcal{O}(f_1^{(m_j)});\dots;
\mathcal{O}(f_n^{(m_j)})\rangle_{K_j,m_j}\bigr|
\le C(\vec f)\,\|k\|_{\sharp}\,D^{-4-\delta}.
\end{equation}
\item[(b)] The shell-class form of the shell functional
(Definition~\ref{10.8:def-shell-functional}): $\mathfrak{B}^{\mathrm{sh}}_{\chi}
=\mathfrak{T}^{\mathrm{sh}}(k_{\chi})$, and for $\chi=\chi_{\mathcal{R}}$ the weights
$k_{\chi_{\mathcal{R}}}$ are supported in the shell
$\{\mathcal{R}-\varepsilon_j\le|x|\le 2\mathcal{R}+\varepsilon_j\}$ and have $\sharp$-norm at most
$C_{\mathrm{sh}}|\omega_0|$, with $C_{\mathrm{sh}}$ independent of $\mathcal{R}\ge 1$ and of $j$;
moreover $\mathfrak{T}^{\mathrm{sh}}(k_{\chi_{\mathcal{R}}})$ is a sum of components of the class
$\mathcal{K}_{\Theta}$, its diagonal strain weights being traceless at every point.
\item[(c)] There is a Lipschitz partition of unity $(\varphi_{\Delta})_{\Delta}$ subordinate to
unit cubes, indexed by a family of closed unit cubes $\Delta$, with
$\operatorname{supp}\varphi_{\Delta}\subset\Delta$ and $\sum_{\Delta}\varphi_{\Delta}=1$, such that
$\|\varphi_{\Delta}k\|_{\sharp}\le(1+160M)\|k\|_{\sharp}$ for every weight $k$, and such that every
ball of radius $r\ge1$ in $\mathbb{R}^4$ meets at most $C_{\mathrm{cnt}}r^4$ cubes of the index
family, with $C_{\mathrm{cnt}}$ independent of $r$.
\end{enumerate}
Then for every $n\ge2$, every $\vec f\in(C^\infty_c(\mathbb{R}^4))^n$ with pairwise disjoint
supports, and every $R_f>0$ with $\bigcup_i\operatorname{supp}f_i\subset B(0,R_f)$, with
$\delta=\delta(n,\vec f)$ given
by (a), there is $C'=C'(n,\vec f)<\infty$ such that for every $\mathcal{R}\ge 2R_f+6$
\begin{equation}\label{10.8:eq-pointwise-implies-2}
\limsup_{j\to\infty}\bigl|\langle\mathfrak{B}^{\mathrm{sh}}_{\chi_{\mathcal{R}}};
\mathcal{O}(f_1^{(m_j)});\dots;\mathcal{O}(f_n^{(m_j)})\rangle_{K_j,m_j}\bigr|
\le C'|\omega_0|\,\mathcal{R}^{-\delta},
\end{equation}
the $\limsup$ taken over the $j$ with $2\mathcal{R}+\varepsilon_j<L^{m_j}$. In particular the
left side of \eqref{10.8:eq-pointwise-implies-2} tends to zero as $\mathcal{R}\to\infty$, and the
far-sphere flux hypothesis $(\mathrm{IR}\text{-}\mathrm{fl})_{\omega_0}$, display
\eqref{10.8:eq-clustering-hypothesis-a}, holds with this $\omega_0$
and with the $\liminf_{\mathcal{R}\to\infty}$ replaced by a full limit.
\end{lemma}

\begin{proof}
Fix $n\ge2$ and $\vec f$ as in the statement, let $\delta$ and $C(\vec f)$ be given by (a), and
let $\mathcal{R}\ge 2R_f+6$. Consider only the $j$ with $2\mathcal{R}+\varepsilon_j<L^{m_j}$;
these are all $j$ from some index on, and for them the shell fits in $\mathbb{T}_{m_j}$ and
$\mathfrak{B}^{\mathrm{sh}}_{\chi_{\mathcal{R}}}$ is a bounded gauge-invariant function of the
bare field, so the cumulant in \eqref{10.8:eq-pointwise-implies-2} is defined. Note that
$\varepsilon_j\le 1$.

Let $(\varphi_{\Delta})_{\Delta}$ be the partition of unity of (c), indexed by unit cubes, with
$\operatorname{supp}\varphi_{\Delta}\subset\Delta$ and $\sum_{\Delta}\varphi_{\Delta}=1$. By (b),
$\mathfrak{B}^{\mathrm{sh}}_{\chi_{\mathcal{R}}}=\mathfrak{T}^{\mathrm{sh}}(k_{\chi_{\mathcal{R}}})$,
and since $\mathfrak{T}^{\mathrm{sh}}(k)$ is by its definition linear in the weights $k$
(Definition~\ref{10.8:def-shell-functional}),
\begin{equation*}
\mathfrak{B}^{\mathrm{sh}}_{\chi_{\mathcal{R}}}
=\sum_{\Delta}\mathfrak{T}^{\mathrm{sh}}(\varphi_{\Delta}k_{\chi_{\mathcal{R}}}),
\end{equation*}
where only cubes $\Delta$ meeting the shell contribute. The shell lies in the ball
$B(0,2\mathcal{R}+1)\subset B(0,3\mathcal{R})$, since $\mathcal{R}\ge1$, so by (c), applied with
$r=3\mathcal{R}$, it meets at most $C_{\mathrm{cnt}}(3\mathcal{R})^4=81\,C_{\mathrm{cnt}}\mathcal{R}^4$
cubes of the partition; thus
$\mathfrak{B}^{\mathrm{sh}}_{\chi_{\mathcal{R}}}$ splits into at most
$81\,C_{\mathrm{cnt}}\mathcal{R}^4$
pieces. Each piece $\mathfrak{T}^{\mathrm{sh}}(\varphi_{\Delta}k_{\chi_{\mathcal{R}}})$ is of the
form $\mathfrak{T}^{\mathrm{sh}}(k)$ with $k=\varphi_{\Delta}k_{\chi_{\mathcal{R}}}$ a real Lipschitz
weight supported in the single cube $\Delta$, the form to which (a) applies
(Definition~\ref{10.8:def-pointwise-decay}). It is a sum of components of the class
$\mathcal{K}_{\Theta}$: by (b) this holds for $k_{\chi_{\mathcal{R}}}$, and multiplication of the
weights by the real scalar $\varphi_{\Delta}(x)$ maps each component to a component of the same
kind, while a diagonal strain weight $b$ traceless at every point stays so, since
$\sum_{\nu}\varphi_{\Delta}(x)b_{\nu\nu}(x)=\varphi_{\Delta}(x)\sum_{\nu}b_{\nu\nu}(x)=0$.
The cube $\Delta$ meets the shell and, by (c) and (b), the weight satisfies
\begin{equation*}
\|\varphi_{\Delta}k_{\chi_{\mathcal{R}}}\|_{\sharp}\le(1+160M)\,C_{\mathrm{sh}}|\omega_0|.
\end{equation*}
A closed unit cube in $\mathbb{R}^4$ has diameter $2$, and a cube $\Delta$ meeting the shell
contains a point $y$ with $|y|\ge\mathcal{R}-\varepsilon_j$, so every $x\in\Delta$ has
$|x|\ge\mathcal{R}-\varepsilon_j-2\ge\mathcal{R}-3$. Since
$\bigcup_i\operatorname{supp}f_i\subset B(0,R_f)$, every such cube is therefore at distance
\begin{equation*}
D\ge\mathcal{R}-3-R_f\ge\mathcal{R}/2\ge 3
\end{equation*}
from $B(0,R_f)$, the middle inequality being $\mathcal{R}\ge 2R_f+6$. The same lower bound holds
for the distance on $\mathbb{T}_{m_j}$ of (a). Indeed, a non-zero period of $\mathbb{T}_{m_j}$ has
norm at least $2L^{m_j}$, so every point of the translate of $B(0,R_f)$ by it has norm at least
$2L^{m_j}-R_f$, while every point of a cube meeting the shell has norm at most
$2\mathcal{R}+\varepsilon_j+2$; since $2\mathcal{R}+\varepsilon_j<L^{m_j}$ and
$R_f\le\mathcal{R}/2-3$, the distance between them is at least
\begin{equation*}
2L^{m_j}-R_f-(2\mathcal{R}+\varepsilon_j+2)
>2\mathcal{R}+\varepsilon_j-2-R_f
\ge\tfrac32\mathcal{R}+1>\mathcal{R}/2 .
\end{equation*}
In particular $D\ge 1$, so (a) applies to each piece.

The joint cumulant is multilinear, in particular linear in its first slot. Hence, for every
admissible $j$,
\begin{align*}
&\bigl|\langle\mathfrak{B}^{\mathrm{sh}}_{\chi_{\mathcal{R}}};
\mathcal{O}(f_1^{(m_j)});\dots;\mathcal{O}(f_n^{(m_j)})\rangle_{K_j,m_j}\bigr|\\
&\qquad\le\sum_{\Delta}\bigl|\langle\mathfrak{T}^{\mathrm{sh}}(\varphi_{\Delta}k_{\chi_{\mathcal{R}}});
\mathcal{O}(f_1^{(m_j)});\dots;\mathcal{O}(f_n^{(m_j)})\rangle_{K_j,m_j}\bigr|\\
&\qquad\le 81\,C_{\mathrm{cnt}}\mathcal{R}^4\cdot C(\vec f)\,(1+160M)\,C_{\mathrm{sh}}|\omega_0|
\,(\mathcal{R}/2)^{-4-\delta}\\
&\qquad=C'|\omega_0|\,\mathcal{R}^{-\delta},
\end{align*}
with $C'=81\cdot2^{4+\delta}C_{\mathrm{cnt}}\,C(\vec f)\,(1+160M)\,C_{\mathrm{sh}}$, where the third
line is \eqref{10.8:eq-pointwise-implies-1} applied to each piece with $D\ge\mathcal{R}/2$. The
bound is uniform in $j$, since none of $C(\vec f)$, $C_{\mathrm{cnt}}$, $C_{\mathrm{sh}}$ depends
on $j$, so taking $\limsup_{j\to\infty}$ gives
\eqref{10.8:eq-pointwise-implies-2}. Because $\delta>0$, the right side of
\eqref{10.8:eq-pointwise-implies-2} tends to zero as $\mathcal{R}\to\infty$; the left side is
non-negative, so its limit as $\mathcal{R}\to\infty$ exists and equals zero. A fortiori its
$\liminf$ is zero, which is \eqref{10.8:eq-clustering-hypothesis-a}, for every $n\ge 2$ and every
$\vec f$ with pairwise disjoint supports.
\end{proof}

\begin{remark}
With $\delta=0$ the same count gives only a bound independent of $\mathcal{R}$: the threshold of
the argument is decay strictly faster than $|x|^{-4}$.

Two features of hypothesis (a) bear on how the lemma is used. First, the uniformity in $j$ along
the joint diagonal of the constant in \eqref{10.8:eq-pointwise-implies-1}
contains the ultraviolet boundedness statement for the stress-tensor insertions of the class
$\mathcal{K}_{\Theta}$, namely the torus bound for the tensor insertion
(Proposition~\ref{10.4:prf-tensor-torus-bound}) of the stress-tensor-insertion correlation
theorem (Theorem~\ref{10.4:thm-insertion-correlation}) restricted to that class; assuming (a)
includes assuming that statement. Second, applied to cubes
near the sources the same inequality would bound the contact terms and so would place into the
hypothesis an ultraviolet statement not proved in this book; the proof above uses
\eqref{10.8:eq-pointwise-implies-1} only for cubes at distance $D\ge\mathcal{R}/2$ from the
sources, with $\mathcal{R}\to\infty$.
\end{remark}

\begin{remark}[Further sufficient conditions for the hypothesis]\label{10.8:rem-sufficient-conditions}
Let the joint diagonal $(K_j,m_j)$, the radial cut-off $\chi_{\mathcal{R}}$, the generator
$\omega_0\in\mathfrak{so}(4)$ and the shell functional $\mathfrak{B}^{\mathrm{sh}}_{\chi_{\mathcal{R}}}$
be as in the far-sphere flux hypothesis $(\mathrm{IR}\text{-}\mathrm{fl})_{\omega_0}$ of clause (vi),
which says that the flux of the lattice rotation current through far spheres vanishes along one
sequence of radii (Definition~\ref{10.8:def-clustering-hypothesis}), and let
$\vec f=(f_1,\dots,f_n)\in(C^\infty_c(\mathbb{R}^4))^n$, $n\ge 2$, have pairwise disjoint supports.
Each of the following conditions implies
\eqref{10.8:eq-clustering-hypothesis-a}; none of them is asserted in this book.
\begin{enumerate}
\item The shell-averaged (log-Ces\`aro) form. Split the shell functional by a Lipschitz partition
of unity $\{\varphi_{\Delta}\}$ subordinate to the unit cubes $\Delta$, as in the proof of
Lemma~\ref{10.8:lem-pointwise-implies}: with $k_{\Delta}:=\varphi_{\Delta}\,k_{\chi_{\mathcal{R}}}$,
\begin{equation*}
\mathfrak{B}^{\mathrm{sh}}_{\chi_{\mathcal{R}}}=\mathfrak{T}^{\mathrm{sh}}(k_{\chi_{\mathcal{R}}})
=\sum_{\Delta}\mathfrak{T}^{\mathrm{sh}}(k_{\Delta}),
\end{equation*}
the sum running over the at most $C\mathcal{R}^{4}$ unit cubes $\Delta$ that meet the shell. The
condition is: there is a sequence $\mathcal{R}_i\to\infty$ such that, with $\mathcal{R}=\mathcal{R}_i$,
\begin{equation*}
\limsup_{j\to\infty}\sum_{\Delta}\bigl|\langle\mathfrak{T}^{\mathrm{sh}}(k_{\Delta});
\mathcal{O}(f_1^{(m_j)});\dots;\mathcal{O}(f_n^{(m_j)})\rangle_{K_j,m_j}\bigr|
\longrightarrow 0 \qquad (i\to\infty);
\end{equation*}
equivalently, the average of these terms over the unit cubes of the shell is
$o(\mathcal{R}_i^{-4})$. This condition is not asserted.
\item Pointwise stress decay faster than $|x|^{-4}$ for the components of the class
$\mathcal{K}_{\Theta}$, that is, \eqref{10.8:eq-pointwise-decay-a} of
Definition~\ref{10.8:def-pointwise-decay}; that it implies the far-sphere flux hypothesis is
Lemma~\ref{10.8:lem-pointwise-implies}. This condition is not asserted.
\item A mass gap uniform in $K$. This condition is not asserted.
\item This condition and the two following ones are formulated in terms of the Hilbert space
$\mathcal{H}^{\mathrm{curv}}$, the vacuum $\Omega$, its projection $P_{\Omega}$ and the Hamiltonian
$H$ of the limit theory, which reflection positivity of the lattice measure furnishes. The
condition is: the ultraviolet boundedness, uniformly along the diagonal $(K_j,m_j)$, of one
$\mathcal{K}_{\Theta}$-class stress-tensor insertion truncated against the observables
$\mathcal{O}(f_1),\dots,\mathcal{O}(f_n)$, together with
\begin{equation*}
\bigl\|e^{-\tau H}(1-P_{\Omega})\Psi\bigr\| = o(\tau^{-4}) \qquad (\tau\to\infty)
\end{equation*}
for every $\Psi$ in the linear span, in $\mathcal{H}^{\mathrm{curv}}$, of the vectors given by
products of smeared curvature observables. This condition is not asserted.
\item The same ultraviolet boundedness as in the preceding condition (or
any polynomial bound on the shell variance), together with
\begin{equation*}
\inf \operatorname{spec}\bigl(H\restriction \Omega^{\perp}\bigr) > 0 .
\end{equation*}
This condition is not asserted.
\item Uniqueness of the vacuum $\Omega$, together with a bounded shell variance: the shell variance
at radius $\mathcal{R}$ and inner scale $\mathcal{R}/8$, taken in the limit superior along the
diagonal, is bounded uniformly in $\mathcal{R}$. This condition is not asserted.
\end{enumerate}
\end{remark}

\begin{proposition}[Sharpness of the $|x|^{-4}$ threshold: a free-field counter-model. Proof outstanding]
\label{10.8:prop-threshold-sharpness}
Let $\phi$ be the massless Gaussian free field on $\mathbb{R}^4$, with covariance
$G(x) = (4\pi^2 |x|^2)^{-1}$. Put
\begin{equation}\label{10.8:eq-threshold-sharpness-1}
h_4(x) := \sum_{\mu=1}^{4} x_\mu^4 - \tfrac12 |x|^4 ,
\end{equation}
a harmonic polynomial that is $W_4$-invariant and not $O(4)$-invariant. For an ordered
$4$-tuple $I = (\mu_1,\mu_2,\mu_3,\mu_4)$ of indices in $\{1,\dots,4\}$ write
$\partial_I = \partial_{\mu_1}\partial_{\mu_2}\partial_{\mu_3}\partial_{\mu_4}$, let
$H_I := \partial_I h_4$ (these are constants) and $\psi_I := \partial_I \phi$. Fix
$\lambda \neq 0$, and for $f \in C_c(\mathbb{R}^4)$ let
\begin{equation}\label{10.8:eq-threshold-sharpness-2}
\mathcal{O}^{\lambda}(f) := \int_{\mathbb{R}^4} f(x)\,
{:}\sum_{I} \bigl(\psi_I(x) + \lambda H_I\bigr)^2{:}\; dx ,
\end{equation}
the sum running over all ordered $4$-tuples $I$ and the Wick ordering being taken with respect
to $\phi$. Thus $\mathcal{O}^{\lambda}(f)$ is the smeared $O(4)$-scalar composite
${:}(\partial^4\phi)^2{:}$ of the free field shifted by $\lambda h_4$, since
$\psi_I + \lambda H_I = \partial_I(\phi + \lambda h_4)$. Let $S^{\lambda}_n(f_1,\dots,f_n)$
be the joint cumulant of $\mathcal{O}^{\lambda}(f_1),\dots,\mathcal{O}^{\lambda}(f_n)$.

Then for every $n \ge 2$ and all $f_1,\dots,f_n \in C_c(\mathbb{R}^4)$ with pairwise
disjoint supports, $\mathsf{d}$ denoting the separation of the supports:
\begin{enumerate}
\item[(1)] $S^{\lambda}_n$ is multilinear, symmetric and $\mathbb{R}^4 \rtimes W_4$-invariant;
\item[(2)] the model, that is the family of Schwinger functions of the observables
$\mathcal{O}^{\lambda}(f)$, is reflection positive at every hyperplane normal to any of the
sixteen mirror normals of $W_4$;
\item[(3)] there is a constant $C_n(\lambda) < \infty$, depending only on $n$ and $\lambda$,
such that $S^{\lambda}_n$ obeys the position-uniform bound
\begin{equation}\label{10.8:eq-threshold-sharpness-3}
\bigl|S^{\lambda}_n(f_1,\dots,f_n)\bigr|
\le C_n(\lambda)\, \mathsf{d}^{-10}\bigl(1 + \mathsf{d}^{-10(n-1)}\bigr)
\prod_{i=1}^{n} \|f_i\|_{L^1},
\end{equation}
whose right-hand side is non-increasing in $\mathsf{d}$;
\item[(4)] the $S^{\lambda}_n$ cluster with a power law, at least like the inverse tenth power
of the distance, and the vacuum reconstructed from them is unique;
\item[(5)] the two-point function is
\begin{equation}\label{10.8:eq-threshold-sharpness-4}
S^{\lambda}_2(f_1,f_2) = S^{0}_2(f_1,f_2)
+ 4\cdot 24^2\,\lambda^2 \iint f_1(x)\, f_2(y)\,
\Bigl[\Bigl(\sum_{\mu=1}^{4}\partial_\mu^4\Bigr)^{2} G\Bigr](x-y)\, dx\, dy ,
\end{equation}
$S^{0}_2$ denoting the two-point cumulant of the unshifted observable $\mathcal{O}^{0}$, the
case $\lambda = 0$ of \eqref{10.8:eq-threshold-sharpness-2}; here
\begin{equation}\label{10.8:eq-threshold-sharpness-5}
\Bigl(\sum_{\mu=1}^{4}\partial_\mu^4\Bigr)^{2} G(x) = P_8(x)\, |x|^{-18}
\end{equation}
with $P_8$ a non-zero harmonic polynomial of degree $8$; in particular $S^{\lambda}_2$ is
not $O(4)$-invariant;
\item[(6)] the short-distance germ of $S^{\lambda}_n$ is rotation invariant.
\end{enumerate}
Moreover, the local Ward identity for the model's stress tensor holds, and the correlation of the
model's stress tensor with the observables $\mathcal{O}^{\lambda}(f_i)$ decays at distance $|y|$
exactly like $|y|^{3}\cdot|y|^{-7} = |y|^{-4}$. Thus the threshold $|x|^{-4}$ of the
pointwise decay condition \eqref{10.8:eq-pointwise-decay-a} of
Definition~\ref{10.8:def-pointwise-decay} is attained and is sharp; what the model violates is
the translation-covariance requirement on the stress-tensor density used in the derivation of
$O(4)$-invariance in this chapter: the model's stress tensor is not translation covariant.
Replacing $h_4$ by a $W_4$-invariant harmonic polynomial of higher
degree, and the
fourth-order derivatives $\partial^4$ by derivatives of the corresponding higher even order,
gives power-law clustering with any exponent.

Consequently there exists a reflection-positive, $\mathbb{R}^4 \rtimes W_4$-invariant,
position-uniformly bounded free-field model, with power-law clustering and a unique vacuum, that
is not rotation invariant and whose stress-tensor correlation decays exactly at the rate
$|x|^{-4}$. Reflection positivity, position-uniform bounds and clustering at that rate
(indeed at any power-law rate) cannot replace the far-sphere flux hypothesis
$(\mathrm{IR\text{-}fl})_{\omega_0}$ of Definition~\ref{10.8:def-clustering-hypothesis}.
No argument that uses only reflection positivity, $\mathbb{R}^4 \rtimes W_4$-invariance and
position-uniform bounds can prove $(\mathrm{IR\text{-}fl})_{\omega_0}$.
\end{proposition}

Proof outstanding.

\emph{Route.} A proof has to establish, for the Gaussian field $\phi$ with covariance $G$ and the
deterministic shift $\lambda h_4$, the following. First, the elementary facts about $h_4$ in
\eqref{10.8:eq-threshold-sharpness-1}: that it is harmonic, invariant under $W_4$ and not under
$O(4)$, so that the $H_I$ are constants forming a $W_4$-invariant tensor. Second, the
expansion of the joint cumulants of the Wick-ordered quadratic functionals
\eqref{10.8:eq-threshold-sharpness-2} of $\phi$ into terms built from derivatives of order eight
of $G$ and from the constants $H_I$; from this expansion, the invariance of $G$ and the
$W_4$-invariance of the $H_I$ have to give (1); the reflection positivity of the free field at
the mirror hyperplanes of $W_4$, together with the behaviour of $h_4$ under the corresponding
reflections, has to give (2); the decay of the eighth derivatives of $G$ like $|x|^{-10}$ has
to give \eqref{10.8:eq-threshold-sharpness-3} and the clustering in (4), and the uniqueness
of the reconstructed vacuum has to be derived from that clustering. Third, the explicit
two-point formula \eqref{10.8:eq-threshold-sharpness-4} with the constant $4\cdot 24^2$, and
the computation \eqref{10.8:eq-threshold-sharpness-5} exhibiting $P_8$ as a non-zero harmonic
polynomial of degree eight which is not $O(4)$-invariant, which together give the failure of
$O(4)$-invariance in (5), while the leading short-distance part, carried by the unshifted
field, gives (6). Fourth, a stress tensor for the model for which the local Ward identity
holds, the computation of its connected correlation with the $\mathcal{O}^{\lambda}(f_i)$ as
the product of a factor growing like $|y|^3$ and a factor decaying like $|y|^{-7}$, and the
failure of its translation covariance; this places the model exactly at the threshold
$|x|^{-4}$ of \eqref{10.8:eq-pointwise-decay-a}, whereas
Lemma~\ref{10.8:lem-pointwise-implies} obtains the far-sphere flux hypothesis
$(\mathrm{IR\text{-}fl})_{\omega_0}$ from decay strictly faster
than $|x|^{-4}$. Fifth, the same computations with $h_4$ replaced by a $W_4$-invariant harmonic
polynomial of higher degree and $\partial^4$ by derivatives of the corresponding higher even
order, giving clustering at any prescribed power-law rate. The consequences stated at the end of
the proposition then follow from the existence of the model, since the model has every property
listed there and is not rotation invariant.

\begin{remark}
The far-sphere flux hypothesis $(\mathrm{IR\text{-}fl})_{\omega_0}$ of clause (vi)
(Definition~\ref{10.8:def-clustering-hypothesis}) says that the flux of the lattice rotation
current through far spheres vanishes along one sequence of radii; decay of one far-away
stress-tensor insertion faster than $|x|^{-4}$, uniformly in the cutoff, as in
\eqref{10.8:eq-pointwise-decay-a}, is a sufficient condition for it
(Lemma~\ref{10.8:lem-pointwise-implies}).
Neither the monotonicity in $\mathsf{d}$ of the bound in clause~(iii-b) of Theorem~0 nor
ergodicity under translations suffices for the far-sphere flux hypothesis
$(\mathrm{IR\text{-}fl})_{\omega_0}$.
Conversely, the ultraviolet input of the derivation of $O(4)$-invariance in this chapter, the
local Ward identity, cannot be dispensed with either: there is a massive,
exponentially clustering, reflection-positive, $\mathbb{R}^4 \rtimes W_4$-invariant model in
which every infrared hypothesis, the far-sphere flux hypothesis included, holds and
$O(4)$-invariance fails because the local Ward identity fails.
\end{remark}

\begin{definition}[The rate-free pointwise form on $C^r$ tuples]\label{10.8:def-rate-free}
Fix one joint diagonal $(K_j,m_j)_{j\ge 1}$ and use the truncated expectations
$\langle\cdot\,;\dots;\cdot\rangle_{K,m}$ with respect to $\mu_{K,m}$, in the sense of
Notation~\ref{10.8:not-joint-diagonal}.
Fix an integer $r\ge 2$.
For $\omega\in\mathfrak{so}(4)$, $n\ge 2$, $\mathcal{R}>0$ and $\mathsf{d}>0$ define the following.

\begin{enumerate}
\item The class of \emph{$C^r$ tuples} is
\begin{equation}\label{10.8:eq-rate-free-1}
\begin{split}
\mathcal{T}^{(r)}_n(\mathcal{R},\mathsf{d})
:=\Bigl\{\vec h=(h_1,\dots,h_n) :\ & h_i\in C^r_c(\mathbb{R}^4)\ \text{real},\
\operatorname{supp}h_i\subset B(0,\mathcal{R}/2),\\
&\operatorname{dist}(\operatorname{supp}h_i,\operatorname{supp}h_k)\ge\mathsf{d}
\ \text{for}\ i\ne k\Bigr\}.
\end{split}
\end{equation}
\item A function $\chi$ is \emph{admissible at radius $\mathcal{R}$} if
$\chi\in C^\infty_c(\mathbb{R}^4)$ and $\chi\equiv 1$ on $B(0,\mathcal{R})$.
\item For $\chi$ admissible at radius $\mathcal{R}$, for every cutoff $K$ and every torus
exponent $m$, and for
$\vec h\in\mathcal{T}^{(r)}_n(\mathcal{R},\mathsf{d})$, put
\begin{equation}\label{10.8:eq-rate-free-2}
\mathcal{B}^{\omega,\chi}_{K,m}(\vec h)
:=\bigl\langle \mathfrak{F}_{\chi};\mathcal{O}(h_1^{(m)});\dots;
\mathcal{O}(h_n^{(m)})\bigr\rangle_{K,m}.
\end{equation}
Here $h_i^{(m)}$ is the periodic lift of $h_i$ to $\mathbb{T}_m$.
The functional $\mathfrak{F}_{\chi}$ is the total dimension-four annulus functional of the
conclusion step; it depends on $\omega$ through the deformation field and on $K$ and $m$
through the spacing $\varepsilon=L^{-K}$ and the torus $\mathbb{T}_m$, and it is given by
\begin{equation}\label{10.8:eq-rate-free-3}
\begin{gathered}
\mathfrak{F}_{\chi}=B^4_{\chi}+J^4_{\chi},\\
B^4_{\chi}=\sum_{\nu,\sigma}\Theta^{\varepsilon}_{\nu\sigma}(b^{\chi}_{\nu\sigma}),\qquad
J^4_{\chi}=\kappa_N\sum_x\sum_{\mu}w^{\chi}_{\mu}(x)\,P_{\mu}(x).
\end{gathered}
\end{equation}
In this formula:
\begin{itemize}
\item $\Theta^{\varepsilon}_{\nu\sigma}(\cdot)$ is the smeared coupling-weighted
stress-tensor density of Definition~\ref{10.8:def-stress-admixtures};
\item $b^{\chi}_{\nu\sigma}$ are the strain weights, and $w^{\chi}_{\mu}$ the weights of
the Jacobian piece, of the deformation field $y\mapsto\chi(y)\,\omega y$, as in the display
of the annulus functional in the conclusion step;
\item $\kappa_N$ and $P_{\mu}$ are as in that display of the conclusion step.
\end{itemize}
\end{enumerate}

We say that the \emph{rate-free pointwise form} holds along $(K_j,m_j)$ with regularity $r$
if, for every $\omega\in\mathfrak{so}(4)$, every $n\ge 2$, every $\mathcal{R}>0$, every
$\mathsf{d}>0$ and every $\vec h\in\mathcal{T}^{(r)}_n(\mathcal{R},\mathsf{d})$, one has
\begin{equation}\label{10.8:eq-rate-free-4}
\inf\Bigl\{\limsup_{j\to\infty}\bigl|\mathcal{B}^{\omega,\chi}_{K_j,m_j}(\vec h)\bigr| :
\ \chi\ \text{admissible at radius}\ \mathcal{R}\Bigr\}=0 .
\end{equation}

No rate of decay is required in \eqref{10.8:eq-rate-free-4}. The infimum runs over
admissible scalar cut-offs $\chi$ only. The deformation field is always
$\chi\cdot(\omega y)$, that is, $y\mapsto\chi(y)\,\omega y$.
\end{definition}

\begin{lemma}[Identification of the two shell functionals in the limit]
\label{10.8:lem-shell-identification}
Fix a diagonal sequence $(K_j,m_j)$, an integer $r\ge 2$, an integer $n\ge 2$, a radius
$\mathcal{R}>0$, a separation $\mathsf{d}>0$ and $\omega\in\mathfrak{so}(4)$. Write
$\langle\cdot;\cdots;\cdot\rangle_j$ for the truncated expectation with respect to
$\mu_{K_j,m_j}$, write $\mathcal{B}^{\omega,\chi}_j(\vec h)$ for
$\mathcal{B}^{\omega,\chi}_{K_j,m_j}(\vec h)$ as in Definition~\ref{10.8:def-rate-free}, and let
$\mathcal{O}(\vec h)$ stand for the $n$ slots
$\mathcal{O}(h_1^{(m_j)}),\dots,\mathcal{O}(h_n^{(m_j)})$. Let $\chi$ be admissible at radius
$\mathcal{R}$, with deformation field $\chi\,\xi_{\omega}$, where $\xi_{\omega}(y)=\omega y$,
and let $\mathfrak{B}^{\mathrm{sh}}_{\chi}$
be the shell functional of the rotational Ward identity
(Definition~\ref{10.8:def-shell-functional}). Assume:
\begin{itemize}
\item[(a)] the agreement of shell functionals: at each cutoff $K$,
$\mathfrak{F}_K(\chi)=\mathfrak{B}^{\mathrm{sh}}_{\chi}+\mathfrak{E}_K(\chi)$, where
$\mathfrak{F}_K(\chi)$ denotes the functional at cutoff $K$ of the rotational Ward identity for the
deformation field $\chi\,\xi_{\omega}$, of which $\mathfrak{B}^{\mathrm{sh}}_{\chi}$ is the explicit
shell term, and $\mathfrak{E}_K(\chi)$ denotes the remainder, its bulk defect, the term controlled
by the dimension-six coefficient bound for the lattice tensor densities;
\item[(b)] the formula for the annulus functional: by the connected form
of the rotational Ward identity assumed in the conclusion step, the functional
$\mathfrak{F}_{\chi}$ entering $\mathcal{B}^{\omega,\chi}_{K,m}$ is the total dimension-four
annulus functional of that identity; its explicit form
$\mathfrak{F}_{\chi}=B^4_{\chi}+J^4_{\chi}$ and, with it, the equality
$\mathfrak{F}_{\chi}=\mathfrak{F}_K(\chi)$ at each cutoff $K$ are assumed, not derived;
\item[(c)] the dimension-six coefficient bound for the lattice tensor densities in its full
form or, alternatively, in its pointwise form as a hypothesis of the conclusion step: the bulk
defect tends to $0$ in $j$, in the sense that for every
$\vec h\in\mathcal{T}^{(r)}_n(\mathcal{R},\mathsf{d})$ one has
$\langle\mathfrak{E}_{K_j}(\chi);\mathcal{O}(\vec h)\rangle_j\to 0$ as $j\to\infty$; this is
the form in which the bound enters here.
\end{itemize}
Then, for every $\vec h\in\mathcal{T}^{(r)}_n(\mathcal{R},\mathsf{d})$,
\begin{equation}\label{10.8:eq-shell-identification-1}
\limsup_{j\to\infty}\bigl|\mathcal{B}^{\omega,\chi}_j(\vec h)\bigr|
=\limsup_{j\to\infty}\bigl|\langle\mathfrak{B}^{\mathrm{sh}}_{\chi};\mathcal{O}(\vec h)\rangle_j\bigr|.
\end{equation}
The conclusion already uses hypothesis (b) and the vanishing of the bulk defect (c);
a later appeal to this identification with no further hypotheses named uses nothing beyond
(a)--(c).
\end{lemma}

\begin{proof}
The identification of $\mathfrak{F}_{\chi}$ with $\mathfrak{F}_K(\chi)$ is not derived here: it is
the content of hypothesis (b). Fix $\vec h\in\mathcal{T}^{(r)}_n(\mathcal{R},\mathsf{d})$ and
write $u_j:=\mathcal{B}^{\omega,\chi}_j(\vec h)$ and
$v_j:=\langle\mathfrak{B}^{\mathrm{sh}}_{\chi};\mathcal{O}(\vec h)\rangle_j$.

By Definition~\ref{10.8:def-rate-free}, $u_j$ is the joint cumulant of $\mathfrak{F}_{\chi}$ and the
$n$ observables with respect to $\mu_{K_j,m_j}$, that is
$u_j=\langle\mathfrak{F}_{\chi};\mathcal{O}(\vec h)\rangle_j$. By (b) this equals
$\langle\mathfrak{F}_{K_j}(\chi);\mathcal{O}(\vec h)\rangle_j$. By (a),
$\mathfrak{F}_{K_j}(\chi)=\mathfrak{B}^{\mathrm{sh}}_{\chi}+\mathfrak{E}_{K_j}(\chi)$. The joint
cumulant is linear in each argument, in particular in the first, so
\begin{equation*}
u_j=\langle\mathfrak{B}^{\mathrm{sh}}_{\chi};\mathcal{O}(\vec h)\rangle_j
+\langle\mathfrak{E}_{K_j}(\chi);\mathcal{O}(\vec h)\rangle_j
=v_j+\langle\mathfrak{E}_{K_j}(\chi);\mathcal{O}(\vec h)\rangle_j .
\end{equation*}
By (c), with the bound in its full form, or in its pointwise form as a hypothesis of the
conclusion step, $u_j-v_j\to 0$.

Two sequences that differ by a null sequence have the same limit superior of absolute values.
Indeed, by the triangle inequality $\bigl||u_j|-|v_j|\bigr|\le|u_j-v_j|$, so
\begin{equation*}
|u_j|\le|v_j|+|u_j-v_j|,
\qquad
|v_j|\le|u_j|+|u_j-v_j|.
\end{equation*}
Taking $\limsup_{j\to\infty}$ and using $|u_j-v_j|\to 0$ gives
$\limsup_j|u_j|\le\limsup_j|v_j|$ and $\limsup_j|v_j|\le\limsup_j|u_j|$. These inequalities hold
also when either side is $+\infty$. This is \eqref{10.8:eq-shell-identification-1}.
\end{proof}

\begin{remark}
The far-sphere flux hypothesis
(Definition~\ref{10.8:def-clustering-hypothesis}), that the flux of the lattice rotation current through
far spheres vanishes along one sequence of radii, is formulated for tuples
$\vec f \in (C^\infty_c(\mathbb{R}^4))^n$ with pairwise disjoint supports, whereas the rate-free pointwise form
(Definition~\ref{10.8:def-rate-free}) is formulated for tuples with $h_i \in C^r_c$ for some fixed $r \ge 2$,
the class $\mathcal{T}^{(r)}_n(\mathcal{R},\mathsf{d})$ being defined for every $r \in \{1,2,\dots,\infty\}$.
By part (a) of Lemma~\ref{10.8:lem-test-classes} below, the real-valued tuples of the first kind are exactly
the $C^\infty_c$ members of the classes $\mathcal{T}^{(r)}_n(\mathcal{R},\mathsf{d})$. The difference is not
one of formulation only. The cumulant bound for the scalar observables controls the joint cumulants
$\langle\mathcal{O}(f_1);\dots;\mathcal{O}(f_{n'})\rangle_j$ of any number $n'$ of scalar observables, the
subscript $j$ referring to the diagonal $(K_j,m_j)$ fixed in
Definition~\ref{10.8:def-clustering-hypothesis}, and it gives no continuity, uniform in $j$, of the map
\begin{equation*}
\vec h \longmapsto
\langle \mathfrak{B}^{\mathrm{sh}}_{\chi};\mathcal{O}(h_1);\dots;\mathcal{O}(h_n)\rangle_j ,
\end{equation*}
a cumulant with one insertion of the shell functional $\mathfrak{B}^{\mathrm{sh}}_{\chi}$, which is of
tensor class rather than a scalar
observable. Passing from $C^\infty_c$ tuples to $C^r_c$ tuples therefore requires a further input: either a
continuity statement of this kind, which is not proved here, or the argument through
ultraviolet bounds given later in this section.
In Lemma~\ref{10.8:lem-test-classes} below the rotations act on functions by
$(R\cdot h)(x) := h(R^{-1}x)$; with the opposite convention $(R\cdot h)(x) := h(Rx)$ the infinitesimal
action $\delta_{\omega}$ defined there becomes $\delta_{\omega} h = +\,\xi_{\omega}\cdot\nabla h$.
Every later use of $\delta_{\omega}$ either depends on $\omega$ only through $|\omega|$
or holds for all $\omega \in \mathfrak{so}(4)$, so nothing depends on this choice of sign.
\end{remark}

\begin{lemma}[Test classes and the sharp norm of a rotation derivative]
\label{10.8:lem-test-classes}
Let $n \ge 2$ and $r \in \{1,2,\dots,\infty\}$, and, for $\mathcal{R},\mathsf{d} > 0$, let
$\mathcal{T}^{(r)}_n(\mathcal{R},\mathsf{d})$ denote the class of Definition~\ref{10.8:def-rate-free}.
\begin{enumerate}
\item[(a)] Every tuple $\vec f = (f_1,\dots,f_n)$ of real-valued functions in $(C^\infty_c(\mathbb{R}^4))^n$
whose members have pairwise disjoint supports lies in $\mathcal{T}^{(r)}_n(\mathcal{R},\mathsf{d})$ for some
$\mathcal{R} > 0$ and some $\mathsf{d} > 0$. Conversely, every member $\vec h$ of
$\mathcal{T}^{(r)}_n(\mathcal{R},\mathsf{d})$ with
every $h_i \in C^\infty_c(\mathbb{R}^4)$ is a tuple in $(C^\infty_c(\mathbb{R}^4))^n$ with pairwise disjoint
supports.
\item[(b)] For $\omega \in \mathfrak{so}(4)$ put $R^{\omega}_t := e^{t\omega} \in \mathrm{SO}(4)$, let
$\mathrm{SO}(4)$ act on functions on $\mathbb{R}^4$ by $(R\cdot h)(x) := h(R^{-1}x)$, put
$\xi_{\omega}(x) := \omega x$, and, for a function $h$ on $\mathbb{R}^4$, put
$\delta_{\omega} h := \frac{d}{dt}\big|_{t=0}(R^{\omega}_t\cdot h)$. Let $\rho > 0$ and let
$h \in C^2_c(\mathbb{R}^4)$ with $\operatorname{supp} h \subset B(0,\rho)$. Then
\begin{equation}
\label{10.8:eq-test-classes-1}
\delta_{\omega} h = -\,\xi_{\omega}\cdot\nabla h ,
\end{equation}
$\delta_{\omega} h \in C^1_c(\mathbb{R}^4)$, $\operatorname{supp}\delta_{\omega} h \subset \operatorname{supp} h$, and
\begin{equation}
\label{10.8:eq-test-classes-2}
\begin{aligned}
\|\delta_{\omega} h\|_{\infty} &\le 2|\omega|\rho\,\|h\|_{C^2}, \qquad
\operatorname{Lip}\delta_{\omega} h \le |\omega|(2+4\rho)\,\|h\|_{C^2},\\
\|\delta_{\omega} h\|_{\sharp} &\le C_{\sharp}(\rho)\,|\omega|\,\|h\|_{C^2}, \qquad
C_{\sharp}(\rho) := 2\rho + 40M(2+4\rho),
\end{aligned}
\end{equation}
and moreover
\begin{equation}
\label{10.8:eq-test-classes-3}
\|h\|_{\sharp} \le (1+80M)\,\|h\|_{C^2}.
\end{equation}
\end{enumerate}
Here $\|h\|_{C^2} := \max_{|\alpha|\le 2}\|\partial^{\alpha} h\|_{\infty}$, $|\omega|$ is the operator norm
of $\omega$ on $\mathbb{R}^4$, and $|\cdot|$ on vectors is the Euclidean norm.
\end{lemma}

\begin{proof}
(a) Let $\vec f \in (C^\infty_c(\mathbb{R}^4))^n$ have pairwise disjoint supports. For $i \ne k$ the sets
$\operatorname{supp} f_i$ and $\operatorname{supp} f_k$ are compact and disjoint, hence at positive distance;
put
\begin{equation*}
\mathsf{d} := \min_{i\ne k}\operatorname{dist}(\operatorname{supp} f_i,\operatorname{supp} f_k) > 0 .
\end{equation*}
The union $\bigcup_i \operatorname{supp} f_i$ is compact, hence bounded,
so it lies in $B(0,\mathcal{R}/2)$ for some $\mathcal{R} > 0$. Since each $f_i$ is real-valued and
$C^\infty_c(\mathbb{R}^4) \subset C^r_c(\mathbb{R}^4)$, we have
$\vec f \in \mathcal{T}^{(r)}_n(\mathcal{R},\mathsf{d})$.
Conversely, if $\vec h \in \mathcal{T}^{(r)}_n(\mathcal{R},\mathsf{d})$ with every $h_i \in C^\infty_c$, then
for $i \ne k$ the distance of $\operatorname{supp} h_i$ and $\operatorname{supp} h_k$ is at least $\mathsf{d} > 0$,
so these supports are disjoint.

(b) Since $h \in C^1$, the chain rule applied to $t \mapsto h(e^{-t\omega}x)$ gives at $t = 0$ the value
$\nabla h(x)\cdot(-\omega x)$, which is \eqref{10.8:eq-test-classes-1}. As $\xi_{\omega}$ is linear and
$\nabla h \in C^1$, the function $\delta_{\omega} h$ is $C^1$. On the open set
$\mathbb{R}^4 \setminus \operatorname{supp} h$
the function $h$ vanishes identically, so $\nabla h = 0$ and $\delta_{\omega} h = 0$ there; hence
$\operatorname{supp}\delta_{\omega} h \subset \operatorname{supp} h$, which is compact, and
$\delta_{\omega} h \in C^1_c$.

Each component $\partial_{\mu} h$ is bounded by $\|h\|_{C^2}$, so $|\nabla h| \le \sqrt{4}\,\|h\|_{C^2}
= 2\|h\|_{C^2}$. On $\operatorname{supp} h \subset B(0,\rho)$ we have $|\xi_{\omega}(x)| \le |\omega|\,|x| \le
|\omega|\rho$. Hence, at points of $\operatorname{supp} h$, $|\delta_{\omega} h| \le |\xi_{\omega}|\,|\nabla h|
\le 2|\omega|\rho\|h\|_{C^2}$, and $\delta_{\omega} h = 0$ elsewhere; this is the first bound in
\eqref{10.8:eq-test-classes-2}.

Write $\xi^{\nu}_{\omega}(x) = \sum_{\lambda}\omega_{\nu\lambda}x_{\lambda}$, so that
$\partial_{\mu}\xi^{\nu}_{\omega} = \omega_{\nu\mu}$. By the product rule, with summation over $\nu$,
\begin{equation*}
\partial_{\mu}\bigl(\xi^{\nu}_{\omega}\partial_{\nu} h\bigr)
= \omega_{\nu\mu}\partial_{\nu} h + \xi^{\nu}_{\omega}\partial_{\mu}\partial_{\nu} h ,
\end{equation*}
that is, $\nabla(\xi_{\omega}\cdot\nabla h) = \omega^{T}\nabla h + (\nabla^2 h)\xi_{\omega}$. The operator norm of
$\omega^{T}$ equals $|\omega|$, so $|\omega^{T}\nabla h| \le 2|\omega|\,\|h\|_{C^2}$. The Hessian $\nabla^2 h$ has
sixteen entries each bounded by $\|h\|_{C^2}$, so its operator norm is at most its Hilbert--Schmidt norm,
$\sqrt{16}\,\|h\|_{C^2} = 4\|h\|_{C^2}$. The Hessian vanishes off $\operatorname{supp} h$, and on
$\operatorname{supp} h$ we have $|\xi_{\omega}| \le |\omega|\rho$, so everywhere
$|(\nabla^2 h)\xi_{\omega}| \le 4\|h\|_{C^2}|\omega|\rho$. Therefore
\begin{equation*}
\sup_x|\nabla\delta_{\omega} h(x)| \le 2|\omega|\,\|h\|_{C^2} + 4|\omega|\rho\|h\|_{C^2}
= |\omega|(2+4\rho)\|h\|_{C^2}.
\end{equation*}
Since $\delta_{\omega} h$ is $C^1$, the mean value theorem along the segment joining two points bounds
$\operatorname{Lip}\delta_{\omega} h$ by this supremum, which is the second bound in
\eqref{10.8:eq-test-classes-2}.

For the sharp norm, $\max\{a,b\} \le a + b$ for $a,b \ge 0$ gives
\begin{align*}
\|\delta_{\omega} h\|_{\sharp}
&\le \|\delta_{\omega} h\|_{\infty} + 40M\operatorname{Lip}\delta_{\omega} h\\
&\le 2|\omega|\rho\|h\|_{C^2} + 40M|\omega|(2+4\rho)\|h\|_{C^2}
= C_{\sharp}(\rho)\,|\omega|\,\|h\|_{C^2}.
\end{align*}
Finally, $\|h\|_{\infty} \le \|h\|_{C^2}$, and, by the mean value theorem again,
$\operatorname{Lip} h \le \sup|\nabla h| \le 2\|h\|_{C^2}$, so that
\begin{equation*}
\|h\|_{\sharp} \le \|h\|_{\infty} + 40M\operatorname{Lip} h \le (1+80M)\|h\|_{C^2},
\end{equation*}
which is \eqref{10.8:eq-test-classes-3}.
\end{proof}

\begin{lemma}[The far-sphere flux hypothesis implies the rate-free form on smooth tuples]\label{10.8:lem-smooth-tuples}
Fix the joint diagonal $(K_j,m_j)$ and the radial cut-offs $\chi_{\mathcal{R}}(x)=\chi_1(|x|/\mathcal{R})$ of
the far-sphere flux hypothesis $(\mathrm{IR}\text{-}\mathrm{fl})_{\omega_0}$ of clause (vi), which states that
the flux of the lattice rotation current through far spheres vanishes along one sequence of radii
(Definition~\ref{10.8:def-clustering-hypothesis}; the decay \eqref{10.8:eq-pointwise-decay-a} of the
correlation of one far-away stress-tensor insertion with the observables, faster than $|x|^{-4}$ and
uniformly in the cut-off, is a sufficient condition for it). Let
$r\in\{1,2,\dots,\infty\}$, $n\ge 2$, $\mathcal{R}>0$ and $\mathsf{d}>0$, and let the class
$\mathcal{T}^{(r)}_n(\mathcal{R},\mathsf{d})$, the
cut-offs admissible at radius $\mathcal{R}$ and the cumulants $\mathcal{B}^{\omega,\chi}_{K,m}(\vec h)$ be as in
the rate-free pointwise form on $C^r$ tuples (Definition~\ref{10.8:def-rate-free}), and write
$\mathcal{B}^{\omega,\chi}_j:=\mathcal{B}^{\omega,\chi}_{K_j,m_j}$. For $\omega\in\mathfrak{so}(4)$, a cut-off
$\chi$ and a tuple $\vec h=(h_1,\dots,h_n)$ put
\begin{equation*}
\langle\mathfrak{B}^{\mathrm{sh}}_{\chi};\mathcal{O}(\vec h)\rangle_j
:=\langle\mathfrak{B}^{\mathrm{sh}}_{\chi};\mathcal{O}(h_1^{(m_j)});\dots;\mathcal{O}(h_n^{(m_j)})
\rangle_{K_j,m_j},
\end{equation*}
where the shell functional $\mathfrak{B}^{\mathrm{sh}}_{\chi}$ is formed with the vector field $\xi(x)=\omega x$;
for $\chi=\chi_{\mathcal{R}'}$ this is defined for all $j$ with $2\mathcal{R}'+\varepsilon_j<L^{m_j}$. Assume:
\begin{enumerate}
\item[(I)] the identification of the two shell functionals in the limit
(Lemma~\ref{10.8:lem-shell-identification}): for every $\omega\in\mathfrak{so}(4)$, every $\vec h\in
\mathcal{T}^{(r)}_n(\mathcal{R},\mathsf{d})$ and every $\chi$ admissible at radius $\mathcal{R}$, with
deformation field $x\mapsto\chi(x)\,\omega x$,
\begin{equation*}
\limsup_{j}\bigl|\mathcal{B}^{\omega,\chi}_j(\vec h)\bigr|
=\limsup_{j}\bigl|\langle\mathfrak{B}^{\mathrm{sh}}_{\chi};\mathcal{O}(\vec h)\rangle_j\bigr| ;
\end{equation*}
\item[(C)] the far-sphere flux hypothesis at every generator: for every $\omega\in\mathfrak{so}(4)$ the display
\eqref{10.8:eq-clustering-hypothesis-a} of the far-sphere flux hypothesis
$(\mathrm{IR}\text{-}\mathrm{fl})_{\omega_0}$ holds with $\omega$ in place
of $\omega_0$ (that
is, with $\xi(x)=\omega x$ in the shell functional), for every $n\ge 2$ and every tuple
$\vec f\in(C^\infty_c(\mathbb{R}^4))^n$ with pairwise disjoint supports.
\end{enumerate}
Then:
\begin{enumerate}
\item[(a)] under (I) alone, for every $\omega\in\mathfrak{so}(4)$ and every $\vec h\in
\mathcal{T}^{(r)}_n(\mathcal{R},\mathsf{d})$,
\begin{equation}\label{10.8:eq-smooth-tuples-1}
\begin{split}
\inf\Bigl\{\limsup_{j}\bigl|\mathcal{B}^{\omega,\chi}_j(\vec h)\bigr| :
\chi\text{ admissible at radius }\mathcal{R}\Bigr\}
&\le\inf_{\mathcal{R}'\ge\mathcal{R}}\,\limsup_{j}
\bigl|\langle\mathfrak{B}^{\mathrm{sh}}_{\chi_{\mathcal{R}'}};\mathcal{O}(\vec h)\rangle_j\bigr|\\
&\le\liminf_{\mathcal{R}'\to\infty}\,\limsup_{j}
\bigl|\langle\mathfrak{B}^{\mathrm{sh}}_{\chi_{\mathcal{R}'}};\mathcal{O}(\vec h)\rangle_j\bigr| ;
\end{split}
\end{equation}
\item[(b)] under (I) and (C), for every $\omega\in\mathfrak{so}(4)$ and every
$\vec h\in\mathcal{T}^{(r)}_n(\mathcal{R},\mathsf{d})$ with $h_i\in C^\infty_c(\mathbb{R}^4)$ for all $i$, the
left-hand side of \eqref{10.8:eq-smooth-tuples-1} vanishes, so that the display of the rate-free pointwise form
(Definition~\ref{10.8:def-rate-free}) holds for $\vec h$. If $r=\infty$, this holds for every member of
$\mathcal{T}^{(r)}_n(\mathcal{R},\mathsf{d})$.
\end{enumerate}
\end{lemma}

\begin{proof}
Fix $\omega\in\mathfrak{so}(4)$ and $\vec h\in\mathcal{T}^{(r)}_n(\mathcal{R},\mathsf{d})$, and put
\begin{equation*}
a(\mathcal{R}'):=\limsup_{j}
\bigl|\langle\mathfrak{B}^{\mathrm{sh}}_{\chi_{\mathcal{R}'}};\mathcal{O}(\vec h)\rangle_j\bigr|
\in[0,\infty],\qquad \mathcal{R}'>0 .
\end{equation*}

First, for $\mathcal{R}'\ge\mathcal{R}$ the radial cut-off $\chi_{\mathcal{R}'}$ is admissible at radius
$\mathcal{R}$: the function $x\mapsto\chi_1(|x|/\mathcal{R}')$ belongs to $C^\infty_c(\mathbb{R}^4)$ and is
identically $1$ on $B(0,\mathcal{R}')$, which contains $B(0,\mathcal{R})$. Hence the family
$\{\chi_{\mathcal{R}'}:\mathcal{R}'\ge\mathcal{R}\}$ is contained in the set of cut-offs admissible at radius
$\mathcal{R}$, and the infimum over the larger set is at most the infimum over this family. For each
$\mathcal{R}'\ge\mathcal{R}$, the identification (I), applied to the admissible cut-off $\chi=\chi_{\mathcal{R}'}$,
gives
\begin{equation*}
\limsup_{j}\bigl|\mathcal{B}^{\omega,\chi_{\mathcal{R}'}}_j(\vec h)\bigr|=a(\mathcal{R}') .
\end{equation*}
Together these give the first inequality in \eqref{10.8:eq-smooth-tuples-1}.

Second, an infimum over a tail is at most the limit inferior. Indeed, for every $T\ge\mathcal{R}$ the set
$\{\mathcal{R}'\ge T\}$ is contained in $\{\mathcal{R}'\ge\mathcal{R}\}$, so
$\inf_{\mathcal{R}'\ge\mathcal{R}}a(\mathcal{R}')\le\inf_{\mathcal{R}'\ge T}a(\mathcal{R}')$; taking the supremum
over $T\ge\mathcal{R}$,
\begin{equation*}
\inf_{\mathcal{R}'\ge\mathcal{R}}a(\mathcal{R}')
\le\sup_{T\ge\mathcal{R}}\,\inf_{\mathcal{R}'\ge T}a(\mathcal{R}')
=\liminf_{\mathcal{R}'\to\infty}a(\mathcal{R}') .
\end{equation*}
This is the second inequality in \eqref{10.8:eq-smooth-tuples-1}, and (a) is proved.

For (b), let in addition $h_i\in C^\infty_c(\mathbb{R}^4)$ for every $i$. By the test-class lemma
(Lemma~\ref{10.8:lem-test-classes}), a member of $\mathcal{T}^{(r)}_n(\mathcal{R},\mathsf{d})$ all of whose
entries are in $C^\infty_c$ is a tuple in $(C^\infty_c(\mathbb{R}^4))^n$ with pairwise disjoint supports (its
supports are pairwise at distance at least $\mathsf{d}>0$). Hypothesis (C) for this $\omega$ therefore applies
to $\vec h$, and it states exactly that $\liminf_{\mathcal{R}'\to\infty}a(\mathcal{R}')=0$. By
\eqref{10.8:eq-smooth-tuples-1} the left-hand side, which is non-negative, vanishes; this is the display of the
rate-free pointwise form for $\vec h$. Since $\omega\in\mathfrak{so}(4)$ was arbitrary, the rate-free form holds
for $\vec h$ for every $\omega$. If $r=\infty$, every member of $\mathcal{T}^{(r)}_n(\mathcal{R},\mathsf{d})$ has
all entries in $C^\infty_c$, and the last assertion follows.
\end{proof}

\begin{remark}
An inequality of the same kind along the tail of the radial family, for the annulus functional
$\mathfrak{F}_{\chi}$ of the conclusion step on that step's own test class, appears in the conclusion step; it
does not bear on the difference between the test classes. The rate-free form is required for some fixed
$r\ge 2$, not necessarily $r=\infty$. For $\vec h\in\mathcal{T}^{(r)}_n(\mathcal{R},\mathsf{d})$ with some
$h_i\in C^r_c\setminus C^\infty_c$ the last step of the proof above would need (C) for $\vec h$, and (C) is
formulated only for $C^\infty_c$ tuples; Lemma~\ref{10.8:lem-smooth-tuples} makes no claim for such tuples, and
from (I) alone none is made. Such tuples are treated later in this section under the ultraviolet hypotheses of
the conclusion step.

A bound that would give the extension without those hypotheses is the following; it is stated here, it is not
asserted in this book, and it is not used. For every $n\ge 2$, $R'>0$ and $\mathsf{d}'>0$ there is a constant
$C^{\mathrm{sh}}_n(\mathsf{d}';R')<\infty$, independent of $j$ and of $\mathcal{R}'$, such that for every
$\omega\in\mathfrak{so}(4)$ (the shell functional being formed with $\xi(x)=\omega x$), every
$\mathcal{R}'\ge 4$, every $j$ with $2\mathcal{R}'+\varepsilon_j<L^{m_j}$ (and $j$ beyond the volume and coupling
thresholds of the bound on cumulants of curvature observables used in this section), and all real Lipschitz
functions $g_1,\dots,g_n$ on $\mathbb{R}^4$, each supported in a ball of radius $R'$ inside
$B(0,\mathcal{R}'/2)$, with supports pairwise at distance at least $\mathsf{d}'$,
\begin{equation}\label{10.8:eq-smooth-tuples-a}
\bigl|\langle\mathfrak{B}^{\mathrm{sh}}_{\chi_{\mathcal{R}'}};\mathcal{O}(g_1^{(m_j)});\dots;
\mathcal{O}(g_n^{(m_j)})\rangle_{K_j,m_j}\bigr|
\le C^{\mathrm{sh}}_n(\mathsf{d}';R')\,|\omega|\prod_{i=1}^n\|g_i\|_{\sharp} .
\end{equation}
This is a bound, uniform in the cut-off, on a cumulant with one insertion of a smeared tensor density of the
shell class and $n$ smeared curvature observables, a bound of the same type as a cut-off-uniform bound on
cumulants with one insertion from the class $\mathcal{K}_{\Theta}$. A bound of this type is implicit in the
sufficient decay form of the far-sphere flux hypothesis, in which the correlation of one far-away
stress-tensor insertion with the observables is required to decay uniformly in the cut-off $K$; no bound of this
type, with
curvature observables in the remaining slots, is asserted in this book.
It is what the far-zone count yields from the pointwise stress decay \eqref{10.8:eq-pointwise-decay-a} with the
exponent $4+\delta$ replaced by $4$ (for $\delta=0$ the count gives a constant, uniform in $\mathcal{R}'$); it
thus sits at the threshold, below \eqref{10.8:eq-pointwise-decay-a}, and like that decay it is not asserted in
this book.

Under \eqref{10.8:eq-smooth-tuples-a} the extension would read as follows. Let $r\ge 1$,
$\vec h\in\mathcal{T}^{(r)}_n(\mathcal{R},\mathsf{d})$ and $0<\eta\le\mathsf{d}/4$. Take $C^\infty_c$ tuples
$\vec g^{(\eta)}$ with $\operatorname{supp}g^{(\eta)}_i$ in the $\eta$-neighbourhood of
$\operatorname{supp}h_i$ and $\|h_i-g^{(\eta)}_i\|_{\sharp}\to 0$ as $\eta\downarrow 0$; for instance
$g^{(\eta)}_i:=h_i*\phi_{\eta}$, with $\phi_{\eta}\ge 0$ smooth, supported in $B(0,\eta)$ and of integral $1$.
For these, $h_i*\phi_{\eta}\to h_i$ in $C^1$ since $h_i\in C^1_c$; hence
$\|h_i-g^{(\eta)}_i\|_{\infty}\to 0$ and, by the mean value theorem,
$\operatorname{Lip}(h_i-g^{(\eta)}_i)\le\|\nabla(h_i-g^{(\eta)}_i)\|_{\infty}\to 0$, so the $\sharp$-norm of the
difference tends to $0$. Moreover $\|g^{(\eta)}_i\|_{\infty}\le\|h_i\|_{\infty}$ and
$\operatorname{Lip}g^{(\eta)}_i\le\operatorname{Lip}h_i$, since $\phi_{\eta}$ is a probability density, so
$\|g^{(\eta)}_i\|_{\sharp}\le\|h_i\|_{\sharp}$. Put $\Gamma:=\max\{1,\max_i\|h_i\|_{\sharp}\}$, so that every
$h_k$ and every $g^{(\eta)}_k$ has $\sharp$-norm at most $\Gamma$, and put
$R':=\mathcal{R}/2+\mathsf{d}/4$. Each of $h_i$, $g^{(\eta)}_i$, $h_i-g^{(\eta)}_i$ is supported in the
$\eta$-neighbourhood of $\operatorname{supp}h_i\subset B(0,\mathcal{R}/2)$, hence in the ball of radius $R'$
about the origin, which lies inside $B(0,\mathcal{R}'/2)$ once $\mathcal{R}'\ge\mathcal{R}+\mathsf{d}/2$; and
for $i\ne k$ these neighbourhoods are at distance at least $\mathsf{d}-2\eta\ge\mathsf{d}/2$. The observable
$\mathcal{O}(\cdot)$ and the lift $h\mapsto h^{(m)}$ are linear, the cumulant is multilinear in the observable
slots, and $\|(h-g)^{(m_j)}\|_{\sharp}=\|h-g\|_{\sharp}$ by the isometry of the lift once $L^{m_j}/2>R'$.
Writing the difference of the cumulants of $\vec h$ and $\vec g^{(\eta)}$ as the telescoping sum over $i$ of the
cumulants of the mixed tuples
\begin{equation*}
\bigl(g^{(\eta)}_1,\dots,g^{(\eta)}_{i-1},\,h_i-g^{(\eta)}_i,\,h_{i+1},\dots,h_n\bigr),
\qquad i=1,\dots,n,
\end{equation*}
whose members other than $h_i-g^{(\eta)}_i$ have $\sharp$-norm at most $\Gamma$, and applying
\eqref{10.8:eq-smooth-tuples-a} with $\mathsf{d}'=\mathsf{d}/2$ to each of them, one obtains for every
$\mathcal{R}'\ge\max\{4,\mathcal{R}+\mathsf{d}/2\}$
\begin{equation}\label{10.8:eq-smooth-tuples-2}
\begin{split}
\limsup_{j}\bigl|\langle\mathfrak{B}^{\mathrm{sh}}_{\chi_{\mathcal{R}'}};\mathcal{O}(\vec h)\rangle_j\bigr|
&\le\limsup_{j}\bigl|\langle\mathfrak{B}^{\mathrm{sh}}_{\chi_{\mathcal{R}'}};
\mathcal{O}(\vec g^{(\eta)})\rangle_j\bigr|\\
&\quad+n\,C^{\mathrm{sh}}_n(\mathsf{d}/2;R')\,|\omega|\,\Gamma^{n-1}
\max_i\|h_i-g^{(\eta)}_i\|_{\sharp} .
\end{split}
\end{equation}
Write $b(\mathcal{R}')$ for the left-hand side of \eqref{10.8:eq-smooth-tuples-2}.
The tuple $\vec g^{(\eta)}$ is in $(C^\infty_c(\mathbb{R}^4))^n$ with pairwise disjoint supports, so (C) gives
that the limit inferior as $\mathcal{R}'\to\infty$ of the first term on the right of
\eqref{10.8:eq-smooth-tuples-2} is $0$. Since the second term does not depend on $\mathcal{R}'$, taking the limit
inferior in $\mathcal{R}'$ in \eqref{10.8:eq-smooth-tuples-2} bounds
$\liminf_{\mathcal{R}'\to\infty}b(\mathcal{R}')$ by that second term, for every $\eta\in(0,\mathsf{d}/4]$; as
$R'$ does not depend on $\eta$, letting $\eta\downarrow 0$ gives
\begin{equation*}
\liminf_{\mathcal{R}'\to\infty}\,\limsup_{j}
\bigl|\langle\mathfrak{B}^{\mathrm{sh}}_{\chi_{\mathcal{R}'}};\mathcal{O}(\vec h)\rangle_j\bigr|=0 ,
\end{equation*}
and the two inequalities \eqref{10.8:eq-smooth-tuples-1} give the rate-free pointwise form for $\vec h$. The
uniformity of $C^{\mathrm{sh}}_n$ in $\mathcal{R}'$ is what lets the outer limit inferior pass; a constant
depending on the cut-off would not suffice, since the cut-offs at which the tuples $\vec g^{(\eta)}$ approach the
infimum may vary with $\eta$.

Finally, the far-sphere flux hypothesis $(\mathrm{IR}\text{-}\mathrm{fl})_{\omega_0}$ asserts the display
\eqref{10.8:eq-clustering-hypothesis-a} for
one generator
$\omega_0$; hypothesis (C), for every $\omega\in\mathfrak{so}(4)$, follows from it only under the ultraviolet
hypotheses of the conclusion step.
\end{remark}

\begin{lemma}[Cut-off independence of the limiting flux]\label{10.8:lem-limiting-flux}
Fix $r\ge 2$, $\omega\in\mathfrak{so}(4)$, $n\ge 2$ and $\mathcal{R},\mathsf{d}>0$. We use the
diagonal $(K_j,m_j)$, the expectations $\langle\cdot\rangle_j$, the truncated expectations
$\langle\cdot;\dots;\cdot\rangle_j$, the joint cumulant $\kappa_j$ with respect to
$\mu_{K_j,m_j}$ and the abbreviation $S^T_{n;j}:=S^T_{n;K_j,m_j}$ fixed in
Notation~\ref{10.8:not-joint-diagonal}. We also use the classes
$\mathcal{T}^{(r')}_n(\mathcal{R},\mathsf{d})$, $r'\ge 1$, admissibility of a cut-off $\chi$ at radius
$\mathcal{R}$, and the cumulant $\mathcal{B}^{\omega,\chi}_j:=\mathcal{B}^{\omega,\chi}_{K_j,m_j}$
of Definition~\ref{10.8:def-rate-free}. Finally, we use the rotations $R^{\omega}_t=e^{t\omega}$,
their action $(R\cdot h)(x)=h(R^{-1}x)$ and the infinitesimal action $\delta_{\omega}h$ of
Lemma~\ref{10.8:lem-test-classes}. For $\vec h\in\mathcal{T}^{(r)}_n(\mathcal{R},\mathsf{d})$ put
\begin{equation*}
D^{\omega}_j(\vec h):=\frac{d}{dt}\Big|_{t=0}S^T_{n;j}(R^{\omega}_t\cdot\vec h),
\qquad R^{\omega}_t\cdot\vec h:=(R^{\omega}_t\cdot h_1,\dots,R^{\omega}_t\cdot h_n).
\end{equation*}
Assume, for every cut-off $\chi$ admissible at radius $\mathcal{R}$:
\begin{enumerate}
\item[(a)] the connected form of the rotational Ward identity: for every
$\vec h\in\mathcal{T}^{(r)}_n(\mathcal{R},\mathsf{d})$ and every $j\ge j_2$, $j_2$ being the
threshold index supplied with that identity,
\begin{equation}\label{10.8:eq-limiting-flux-1}
\mathcal{B}^{\omega,\chi}_j(\vec h)=D^{\omega}_j(\vec h)-\mathcal{N}^{\omega,\chi}_j(\vec h),
\end{equation}
where $\mathcal{N}^{\omega,\chi}_j(\vec h)$ is the remainder term of that identity;
\item[(b)] the pointwise form of the dimension-six coefficient bound, which yields: the remainder
term $\mathcal{N}^{\omega,\chi}_j(\vec h)$ of~(a) tends to $0$ as $j\to\infty$ for every
$\vec h\in\mathcal{T}^{(r)}_n(\mathcal{R},\mathsf{d})$;
\item[(c)] the bound on the connected Schwinger functions $S^T_{n';j}$ of the scalar observables
along the diagonal, that is, on the cumulants
$\kappa_j(\mathcal{O}(f_1^{(m_j)}),\dots,\mathcal{O}(f_{n'}^{(m_j)}))$, $n'\ge 2$,
of the smeared curvature observables, for tuples $(f_1,\dots,f_{n'})$ of test functions; this
hypothesis is not used in the proof below;
\item[(d)] the convergence statement on $\mathbb{R}^4$ (clause (ii) of Theorem~0), in the form:
for every $\vec f\in\mathcal{T}^{(1)}_n(\mathcal{R},\mathsf{d})$ the limit
$S^T_{n;\infty}(\vec f)=\lim_{j}S^T_{n;j}(\vec f)$ exists.
\end{enumerate}
Then, for every $\vec h\in\mathcal{T}^{(r)}_n(\mathcal{R},\mathsf{d})$:
\begin{enumerate}
\item[(1)] at every $j$ with $L^{m_j}>\mathcal{R}$ the derivative $D^{\omega}_j(\vec h)$ exists and
\begin{equation}\label{10.8:eq-limiting-flux-2}
D^{\omega}_j(\vec h)=\sum_{i=1}^{n}S^T_{n;j}(h_1,\dots,\delta_{\omega}h_i,\dots,h_n);
\end{equation}
\item[(2)] $D^{\omega}_j(\vec h)\to c(\vec h,\omega)$ as $j\to\infty$, where
\begin{equation}\label{10.8:eq-limiting-flux-3}
c(\vec h,\omega):=\sum_{i=1}^{n}S^T_{n;\infty}(h_1,\dots,\delta_{\omega}h_i,\dots,h_n);
\end{equation}
\item[(3)] for every $\chi$ admissible at radius $\mathcal{R}$ the limit
$\lim_j\mathcal{B}^{\omega,\chi}_j(\vec h)$ exists and equals $c(\vec h,\omega)$, a number
independent of $\chi$, which is the limit as $j\to\infty$ of the derivatives at $t=0$ of the
functions $t\mapsto S^T_{n;j}(R^{\omega}_t\cdot\vec h)$;
\item[(4)] the rate-free pointwise form on $C^r$ tuples of the far-sphere flux hypothesis
(Definition~\ref{10.8:def-rate-free}) holds for $\vec h$ if and only if $c(\vec h,\omega)=0$.
\end{enumerate}
\end{lemma}

\begin{proof}
Fix $\vec h\in\mathcal{T}^{(r)}_n(\mathcal{R},\mathsf{d})$ and a cut-off $\chi$ admissible at
radius $\mathcal{R}$.

\emph{The observables along the rotation.} Let $j$ satisfy $L^{m_j}>\mathcal{R}$. Since $R^{\omega}_t$
is orthogonal, $\operatorname{supp}(R^{\omega}_t\cdot h_i)=R^{\omega}_t(\operatorname{supp}h_i)$ lies
in $B(0,\mathcal{R}/2)$ for every $t$. This ball has radius $\mathcal{R}/2<L^{m_j}/2$. Hence two
distinct translates of it by the period lattice $2L^{m_j}\mathbb{Z}^4$ have centres at distance at
least $2L^{m_j}>\mathcal{R}$ and are therefore disjoint. The periodic lift
$(R^{\omega}_t\cdot h_i)^{(m_j)}$ is the sum of the translates of $R^{\omega}_t\cdot h_i$ by the
period lattice. Thus at each plaquette barycentre $x(p)$ at most one translate contributes, and it
is the same for all $t$. Precisely, if the class $x(p)+2L^{m_j}\mathbb{Z}^4$ meets
$B(0,\mathcal{R}/2)$, it does so in exactly one point $y_p$, and then
$(R^{\omega}_t\cdot h_i)^{(m_j)}(x(p))=h_i(e^{-t\omega}y_p)$ for all $t$. Otherwise this value is
$0$ for all $t$. Consequently
\begin{equation*}
\mathcal{O}\big((R^{\omega}_t\cdot h_i)^{(m_j)}\big)=\sum_p(R^{\omega}_t\cdot h_i)^{(m_j)}(x(p))\,a_p
\end{equation*}
is a finite sum, over the plaquettes of the lattice of spacing $\varepsilon_j$ on
$\mathbb{T}_{m_j}$, of the products of the numbers $u_{i,p}(t):=h_i(e^{-t\omega}y_p)$ (or $0$) with
the functions $a_p$ of the configuration. Here $0\le a_p\le 2$, since
$|\operatorname{Re}\operatorname{tr}U(\partial p)|\le 1$.

Since $h_i\in C^r_c\subset C^1_c$, each $u_{i,p}$ is continuously differentiable in $t$ with
\begin{equation*}
\dot u_{i,p}(t)=-\nabla h_i(e^{-t\omega}y_p)\cdot\omega e^{-t\omega}y_p .
\end{equation*}
Its modulus is at most $|\omega|\,(\mathcal{R}/2)\,\|\nabla h_i\|_\infty$, because the point
$e^{-t\omega}y_p$ has norm $|y_p|$, and $\nabla h_i$ vanishes outside $B(0,\mathcal{R}/2)$. Also
$|u_{i,p}(t)|\le\|h_i\|_\infty$. At $t=0$ one has $\dot u_{i,p}(0)=-(\omega y_p)\cdot\nabla
h_i(y_p)=(\delta_{\omega}h_i)(y_p)$. Moreover $\operatorname{supp}\delta_{\omega}h_i\subset
\operatorname{supp}h_i\subset B(0,\mathcal{R}/2)$, so by the same one-translate argument
$(\delta_{\omega}h_i)^{(m_j)}(x(p))=(\delta_{\omega}h_i)(y_p)=\dot u_{i,p}(0)$ when $y_p$ exists,
and $(\delta_{\omega}h_i)^{(m_j)}(x(p))=0$ otherwise, in which case the coefficient of $a_p$ is $0$
for all $t$. Writing $Y_i(t):=\mathcal{O}((R^{\omega}_t\cdot h_i)^{(m_j)})$, we conclude that
$t\mapsto Y_i(t)$ is continuously differentiable. It is bounded, together with its $t$-derivative,
uniformly in $t$ and in the configuration, and
$\dot Y_i(0)=\mathcal{O}((\delta_{\omega}h_i)^{(m_j)})$.

\emph{Differentiation of the cumulant.} The cumulant is multilinear in its bounded arguments and
is the polynomial
\begin{equation*}
\kappa_j(Y_1,\dots,Y_n)=\sum_{\pi}(-1)^{|\pi|-1}(|\pi|-1)!\prod_{\beta\in\pi}
\Big\langle\prod_{i\in\beta}Y_i\Big\rangle_j
\end{equation*}
in the joint moments, the sum running over the partitions $\pi$ of $\{1,\dots,n\}$ into $|\pi|$
blocks $\beta$. Each integrand
$\prod_{i\in\beta}Y_i(t)$ is continuously differentiable in $t$, and its $t$-derivative is bounded
uniformly in $t$ and in the configuration. By the mean value theorem, its difference quotients at
$t=0$ are therefore dominated by a constant, which is integrable because the configuration space is
compact. By dominated convergence,
\begin{equation*}
\frac{d}{dt}\Big|_{t=0}\Big\langle\prod_{i\in\beta}Y_i(t)\Big\rangle_j
=\sum_{i\in\beta}\Big\langle\dot Y_i(0)\prod_{\ell\in\beta\setminus\{i\}}Y_{\ell}(0)\Big\rangle_j .
\end{equation*}
Apply the product rule to each term of the partition formula. For a given $i$, exactly one block of
$\pi$ contains $i$. Collecting the terms in which $Y_i$ has been replaced by $\dot Y_i(0)$ therefore
reproduces the partition formula for $\kappa_j(Y_1,\dots,\dot Y_i(0),\dots,Y_n)$. Hence the
cumulant may be differentiated slot by slot. Since
$S^T_{n;j}(R^{\omega}_t\cdot\vec h)=\kappa_j(Y_1(t),\dots,Y_n(t))$, this proves
\eqref{10.8:eq-limiting-flux-2} and assertion~(1).

\emph{The limit.} Because $r\ge 2$, each $\delta_{\omega}h_i$ lies in $C^{r-1}_c\subset C^1_c$. It
is real, and $\operatorname{supp}\delta_{\omega}h_i\subset\operatorname{supp}h_i$. Hence the tuple
$(h_1,\dots,\delta_{\omega}h_i,\dots,h_n)$ has real $C^1_c$ components supported in
$B(0,\mathcal{R}/2)$, with supports pairwise at distance at least $\mathsf{d}$. That is, it belongs
to $\mathcal{T}^{(1)}_n(\mathcal{R},\mathsf{d})$. Hypothesis~(d) gives
\begin{equation*}
S^T_{n;j}(h_1,\dots,\delta_{\omega}h_i,\dots,h_n)\to
S^T_{n;\infty}(h_1,\dots,\delta_{\omega}h_i,\dots,h_n)
\qquad (j\to\infty)
\end{equation*}
for each $i$. By \eqref{10.8:eq-limiting-flux-2}, applied at every $j$ with $L^{m_j}>\mathcal{R}$,
this yields $D^{\omega}_j(\vec h)\to c(\vec h,\omega)$, with $c(\vec h,\omega)$ as in
\eqref{10.8:eq-limiting-flux-3}. Here we use that the torus side $2L^{m_j}$ tends to infinity along
the diagonal (Notation~\ref{10.8:not-joint-diagonal}), so that \eqref{10.8:eq-limiting-flux-2}
holds for all sufficiently large $j$. This is assertion~(2).

\emph{Independence of the cut-off.} By hypothesis~(a), for $j\ge j_2$ the identity
\eqref{10.8:eq-limiting-flux-1} holds. By hypothesis~(b) its last term tends to $0$. Together with
assertion~(2) this gives
\begin{equation*}
\lim_{j\to\infty}\mathcal{B}^{\omega,\chi}_j(\vec h)=c(\vec h,\omega)-0=c(\vec h,\omega).
\end{equation*}
The right-hand side does not involve $\chi$. By the definition of $D^{\omega}_j(\vec h)$ and
assertion~(2),
\begin{equation*}
c(\vec h,\omega)=\lim_{j\to\infty}\frac{d}{dt}\Big|_{t=0}S^T_{n;j}(R^{\omega}_t\cdot\vec h).
\end{equation*}
This is assertion~(3).

\emph{The pointwise form of the far-sphere flux hypothesis.} By
Definition~\ref{10.8:def-rate-free}, this form holds for $\vec h$ exactly when
\begin{equation*}
\inf\Big\{\limsup_{j}|\mathcal{B}^{\omega,\chi}_j(\vec h)|:\chi\text{ admissible at radius }
\mathcal{R}\Big\}=0 .
\end{equation*}
By assertion~(3), $\limsup_j|\mathcal{B}^{\omega,\chi}_j(\vec h)|=|c(\vec h,\omega)|$ for every
admissible $\chi$. The set of admissible cut-offs is non-empty, since a smooth compactly supported
function equal to $1$ on $B(0,\mathcal{R})$ exists. The infimum is therefore the
infimum of the constant $|c(\vec h,\omega)|$, that is, $|c(\vec h,\omega)|$ itself. It vanishes if
and only if $c(\vec h,\omega)=0$, which proves assertion~(4).
\end{proof}

\begin{lemma}[$C^2$-continuity of the limiting flux]\label{10.8:lem-flux-continuity}
Let $n\ge 2$, $\mathcal{R}>0$, $\mathsf{d}>0$, $r\ge 2$ and $\omega\in\mathfrak{so}(4)$.
Fix $\vec h\in\mathcal{T}^{(r)}_n(\mathcal{R},\mathsf{d})$ and put $\eta_0:=\mathsf{d}/4$,
\begin{equation*}
U_k:=\{x\in\mathbb{R}^4:\ \operatorname{dist}(x,\operatorname{supp}h_k)\le\eta_0\},
\qquad R':=\mathcal{R}/2+\eta_0,\qquad \mathsf{d}':=\mathsf{d}/2 .
\end{equation*}
Assume the following two statements, for tuples of real test functions on $\mathbb{R}^4$.
\begin{enumerate}
\item[(a)] \emph{Clause (ii-$\mathbb{R}^4$) of Theorem~0: existence and multilinearity of the
limit.} For every $n$-tuple $\vec u=(u_1,\dots,u_n)$ of real functions in $C^1_c(\mathbb{R}^4)$
with pairwise disjoint supports, the limit $S^T_{n;\infty}(\vec u)$ of the connected curvature
correlations along the diagonal sequence of cutoffs $(K_j,m_j)_j$ fixed in this chapter exists.
Moreover, for the sets $U_1,\dots,U_n$ above, the map
$\vec u\mapsto S^T_{n;\infty}(\vec u)$ is multilinear on the tuples with
$\operatorname{supp}u_k\subset U_k$ for every $k$.
\item[(b)] \emph{The sharp-norm bound for the limit correlations, with the constant
$C^\infty_n$.} If $\vec u$ is a tuple as in (a) whose supports lie in a ball of
radius $R'$ and are pairwise at distance at least $\mathsf{d}'$, then
\begin{equation*}
\bigl|S^T_{n;\infty}(\vec u)\bigr|\le C^\infty_n(\mathsf{d}';R')\prod_{k=1}^n\|u_k\|_{\sharp}.
\end{equation*}
\end{enumerate}
Let $\delta_\omega$, the norm $\|\cdot\|_{C^2}$ and the constant $C_{\sharp}(\cdot)$ be as in
Lemma~\ref{10.8:lem-test-classes}.
For a tuple $\vec u$ of real $C^2_c$ functions put
\begin{equation*}
c(\vec u,\omega):=\sum_{i=1}^n S^T_{n;\infty}(u_1,\dots,u_{i-1},\delta_\omega u_i,u_{i+1},\dots,u_n)
\end{equation*}
whenever the right-hand side is defined; for $\vec u=\vec h$ this is the limiting flux of
Lemma~\ref{10.8:lem-limiting-flux}. Let $\vec g=(g_1,\dots,g_n)$ with $g_k\in C^2_c(\mathbb{R}^4)$
real and $\operatorname{supp}g_k\subset U_k$, and put
\begin{equation*}
\hat H:=\max\Bigl\{1,\ \max_{1\le k\le n}\max\bigl\{\|h_k\|_{\sharp},\|g_k\|_{\sharp},
\|\delta_\omega h_k\|_{\sharp},\|\delta_\omega g_k\|_{\sharp}\bigr\}\Bigr\}.
\end{equation*}
Then $c(\vec h,\omega)$ and $c(\vec g,\omega)$ are defined, and
\begin{equation}\label{10.8:eq-flux-continuity-a}
\begin{split}
|c(\vec h,\omega)-c(\vec g,\omega)|
&\le n^2\,C^\infty_n(\mathsf{d}';R')\,\hat H^{\,n-1}\bigl(1+80M+C_{\sharp}(R')|\omega|\bigr)\\
&\qquad\times\max_{1\le k\le n}\|h_k-g_k\|_{C^2}.
\end{split}
\end{equation}
Thus $c(\cdot,\omega)$ is locally Lipschitz in the tuple, with respect to the $C^2$ norms of the
$h_k$, uniformly on sets bounded in the sharp norms of $(h_k,\delta_\omega h_k)$, the constant
being $C^\infty_n(\mathsf{d}';R')$ up to the explicit factors in \eqref{10.8:eq-flux-continuity-a}.
Said exactly, \eqref{10.8:eq-flux-continuity-a} is continuity in the sharp norms of the slots $h_k$
and $\delta_\omega h_k$, as in hypothesis (b); since $\delta_\omega$ costs one derivative, this is
continuity in the $C^2$ norm of the $h_k$ themselves, which is what approximation of tuples in
$C^r_c$, $r\ge 2$, by smooth tuples requires, and it is all that the subsequent density argument on
$\mathcal{T}^{(r)}_n(\mathcal{R},\mathsf{d})$ uses.
\end{lemma}

\begin{proof}
We begin with the geometry. Each $U_k$ is closed and bounded, hence compact. Since
$\vec h\in\mathcal{T}^{(r)}_n(\mathcal{R},\mathsf{d})$, the support of $h_k$ lies in
$B(0,\mathcal{R}/2)$, so every point of $U_k$ lies within $\eta_0$ of a point of
$B(0,\mathcal{R}/2)$, hence in $B(0,\mathcal{R}/2+\eta_0)=B(0,R')$; thus
$U_k\subset B(0,R')$. For $i\ne k$ the supports of $h_i$ and $h_k$ are at distance at least
$\mathsf{d}$, and every point of $U_i$, respectively $U_k$, lies within $\eta_0$ of
$\operatorname{supp}h_i$, respectively $\operatorname{supp}h_k$; hence
\begin{equation*}
\operatorname{dist}(U_i,U_k)\ge\mathsf{d}-2\eta_0=\mathsf{d}/2=\mathsf{d}'.
\end{equation*}

Next, the functions involved. Since $r\ge 2$, each $h_k$ lies in $C^2_c$, as does each $g_k$ by
assumption. By Lemma~\ref{10.8:lem-test-classes}, for $u\in C^2_c(\mathbb{R}^4)$ one has
$\delta_\omega u\in C^1_c$ and $\operatorname{supp}\delta_\omega u\subset\operatorname{supp}u$, and
$\delta_\omega u$ is real when $u$ is. Hence $h_k$, $g_k$, $h_k-g_k$ and their images under
$\delta_\omega$ are real $C^1_c$ functions supported in $U_k$. Consequently every $n$-tuple whose
$k$-th member is one of these functions, for each $k$, is a tuple of real $C^1_c$ functions with
pairwise disjoint supports; by (a), $S^T_{n;\infty}$ is defined on it, in particular
$c(\vec h,\omega)$ and $c(\vec g,\omega)$ are defined. By (a), $S^T_{n;\infty}$ is multilinear on
the tuples with $k$-th member supported in $U_k$. (This multilinearity is inherited through the limit
from the finite cutoffs: there $S^T_{n;K,m}$ is a joint cumulant, multilinear in the observables,
and the observable $\mathcal{O}$ of a periodic lift is linear in the test function.) Finally all these
tuples have supports in the ball $B(0,R')$, pairwise at distance at least $\mathsf{d}'$, so (b)
gives
\begin{equation}\label{10.8:eq-flux-continuity-1}
\bigl|S^T_{n;\infty}(\vec u)\bigr|\le C^\infty_n(\mathsf{d}';R')\prod_{k=1}^n\|u_k\|_{\sharp}
\end{equation}
for each of them.

Fix $i\in\{1,\dots,n\}$ and define tuples $\vec H^{(i)}$, $\vec G^{(i)}$ by
$H^{(i)}_k:=h_k$, $G^{(i)}_k:=g_k$ for $k\ne i$, and $H^{(i)}_i:=\delta_\omega h_i$,
$G^{(i)}_i:=\delta_\omega g_i$. By the definition of $c$,
\begin{equation}\label{10.8:eq-flux-continuity-2}
c(\vec h,\omega)-c(\vec g,\omega)
=\sum_{i=1}^n\Bigl(S^T_{n;\infty}(\vec H^{(i)})-S^T_{n;\infty}(\vec G^{(i)})\Bigr).
\end{equation}
Replacing the slots one at a time, from the first to the last, and using multilinearity in each slot
(all members involved being supported in the corresponding $U_k$), we obtain the telescoping identity
\begin{equation*}
\begin{split}
&S^T_{n;\infty}(\vec H^{(i)})-S^T_{n;\infty}(\vec G^{(i)})\\
&\qquad=\sum_{k=1}^n S^T_{n;\infty}\bigl(G^{(i)}_1,\dots,G^{(i)}_{k-1},
H^{(i)}_k-G^{(i)}_k,H^{(i)}_{k+1},\dots,H^{(i)}_n\bigr).
\end{split}
\end{equation*}
In the $k$-th term, each of the $n-1$ slots other than the $k$-th carries one of $h_l$, $g_l$,
$\delta_\omega h_l$, $\delta_\omega g_l$, whose sharp norm is at most $\hat H$. Hence, by
\eqref{10.8:eq-flux-continuity-1}, the $k$-th term is bounded in absolute value by
\begin{equation*}
C^\infty_n(\mathsf{d}';R')\,\hat H^{\,n-1}\,\bigl\|H^{(i)}_k-G^{(i)}_k\bigr\|_{\sharp}.
\end{equation*}
It remains to bound the differences. For $k\ne i$ one has $H^{(i)}_k-G^{(i)}_k=h_k-g_k$, and by the
conventions of Lemma~\ref{10.8:lem-test-classes},
\begin{equation*}
\|h_k-g_k\|_{\sharp}\le(1+80M)\|h_k-g_k\|_{C^2}.
\end{equation*}
For $k=i$, since $\delta_\omega$ is linear,
$H^{(i)}_i-G^{(i)}_i=\delta_\omega(h_i-g_i)$, and $h_i-g_i$ is a $C^2_c$ function supported in
$U_i\subset B(0,R')$; the same lemma therefore gives
\begin{equation*}
\|\delta_\omega(h_i-g_i)\|_{\sharp}\le C_{\sharp}(R')|\omega|\,\|h_i-g_i\|_{C^2}.
\end{equation*}
Both right-hand sides are at most
\begin{equation*}
\bigl(1+80M+C_{\sharp}(R')|\omega|\bigr)\max_{l}\|h_l-g_l\|_{C^2}.
\end{equation*}
Summing the $n$ terms over
$k$ shows that, for each $i$,
\begin{equation*}
\begin{split}
\bigl|S^T_{n;\infty}(\vec H^{(i)})-S^T_{n;\infty}(\vec G^{(i)})\bigr|
&\le n\,C^\infty_n(\mathsf{d}';R')\,\hat H^{\,n-1}\bigl(1+80M+C_{\sharp}(R')|\omega|\bigr)\\
&\qquad\times\max_{l}\|h_l-g_l\|_{C^2}.
\end{split}
\end{equation*}
Summing over the $n$ values of $i$ in \eqref{10.8:eq-flux-continuity-2} gives
\eqref{10.8:eq-flux-continuity-a}.

For the final assertion, fix a bound on the sharp norms of the $h_k$, $g_k$, $\delta_\omega h_k$,
$\delta_\omega g_k$; then $\hat H$ is bounded, and \eqref{10.8:eq-flux-continuity-a} is a Lipschitz
bound in $\max_k\|h_k-g_k\|_{C^2}$ for all $\vec g$ with $\operatorname{supp}g_k\subset U_k$, with
constant $C^\infty_n(\mathsf{d}';R')$ times the explicit factors displayed.
\end{proof}

\begin{proposition}[Extension to $C^r$ tuples by mollification]\label{10.8:prop-mollified-extension}
Fix the diagonal $(K_j,m_j)$ of the far-sphere flux hypothesis
$(\mathrm{IR}\text{-}\mathrm{fl})_{\omega_0}$
(Definition~\ref{10.8:def-clustering-hypothesis}), and $r\ge 2$ (an integer or $\infty$); write
$\mathcal{B}^{\omega,\chi}_j$ for $\mathcal{B}^{\omega,\chi}_{K_j,m_j}$ and
$\langle\cdot\,;\dots;\cdot\rangle_j$ for $\langle\cdot\,;\dots;\cdot\rangle_{K_j,m_j}$.
Assume, along this diagonal, for every $\omega\in\mathfrak{so}(4)$, every $n\ge 2$, all
$\mathcal{R}>0$, $\mathsf{d}>0$, every $\vec h\in\mathcal{T}^{(r)}_n(\mathcal{R},\mathsf{d})$,
every cut-off $\chi$ admissible at radius $\mathcal{R}$ and all sufficiently large $j$:
\begin{enumerate}
\item the identification of the two shell functionals in the limit,
Lemma~\ref{10.8:lem-shell-identification};
\item the connected rotational Ward identity, which writes $\mathcal{B}^{\omega,\chi}_j(\vec h)$
as the finite-cutoff rotation derivative $D^{\omega}_j(\vec h)$ minus a remainder
$\mathcal{N}^{\omega,\chi}_j(\vec h)$;
\item the pointwise vanishing of the remainder $\mathcal{N}^{\omega,\chi}_j(\vec h)$ as $j\to\infty$;
\item the volume-free bound on $S^T_{n;\infty}$ with the constants $C^\infty_n(\mathsf{d};R)$;
\item the existence of the limit $S^T_{n;\infty}$ on $\mathbb{R}^4$ along the diagonal,
clause (ii-$\mathbb{R}^4$) of Theorem~0.
\end{enumerate}
Suppose that for every $\omega\in\mathfrak{so}(4)$ the display
\eqref{10.8:eq-clustering-hypothesis-a} holds with $\omega$ in place of $\omega_0$, that is, with
$\xi(x)=\omega x$ in the shell functional $\mathfrak{B}^{\mathrm{sh}}_{\chi}$. Then for every
$\omega\in\mathfrak{so}(4)$, every $n\ge 2$, all $\mathcal{R},\mathsf{d}>0$ and every
$\vec h\in\mathcal{T}^{(r)}_n(\mathcal{R},\mathsf{d})$ the rate-free pointwise form
(Definition~\ref{10.8:def-rate-free}) holds:
\begin{equation}\label{10.8:eq-mollified-extension-1}
\inf\Bigl\{\limsup_{j\to\infty}\bigl|\mathcal{B}^{\omega,\chi}_j(\vec h)\bigr| :
\chi \text{ admissible at radius } \mathcal{R}\Bigr\}=0 .
\end{equation}
\end{proposition}

\begin{proof}
Fix $\omega$, $n$, $\mathcal{R}$, $\mathsf{d}$ and
$\vec h=(h_1,\dots,h_n)\in\mathcal{T}^{(r)}_n(\mathcal{R},\mathsf{d})$. As in the setting of the
$C^2$-continuity of the limiting flux (Lemma~\ref{10.8:lem-flux-continuity}) put
$\eta_0:=\mathsf{d}/4$,
\begin{equation*}
U_k:=\{x : \operatorname{dist}(x,\operatorname{supp}h_k)\le \eta_0\},\qquad
R':=\mathcal{R}/2+\eta_0,\qquad \mathsf{d}':=\mathsf{d}/2 ;
\end{equation*}
each $U_k$ is compact, $U_k\subset B(0,R')$, and $\operatorname{dist}(U_i,U_k)\ge\mathsf{d}'$ for
$i\ne k$, since the supports of $h_i$ and $h_k$ are at distance at least $\mathsf{d}=4\eta_0$.

\emph{The mollified tuples.} Let $\phi\in C^\infty_c(B(0,1))$ with $\phi\ge 0$ and $\int\phi=1$,
put $\phi_\eta(x):=\eta^{-4}\phi(x/\eta)$, and for $0<\eta\le\eta_0$ set
$g^{(\eta)}_k:=h_k\ast\phi_\eta$, so that
\begin{equation*}
g^{(\eta)}_k(x)=\int_{\mathbb{R}^4}h_k(x-y)\,\phi_\eta(y)\,dy .
\end{equation*}
Since $\phi_\eta\in C^\infty_c$ and $h_k$ is real, $g^{(\eta)}_k\in C^\infty_c(\mathbb{R}^4)$ is real;
since $\operatorname{supp}\phi_\eta\subset \overline{B(0,\eta)}$ and $\eta\le\eta_0$, the support of
$g^{(\eta)}_k$ lies in the $\eta$-neighbourhood of $\operatorname{supp}h_k$, hence in $U_k$.
Because $\phi_\eta$ is a probability density, $|g^{(\eta)}_k(x)|\le\|h_k\|_\infty$ and
\begin{equation*}
|g^{(\eta)}_k(x)-g^{(\eta)}_k(x')|\le\int|h_k(x-y)-h_k(x'-y)|\,\phi_\eta(y)\,dy
\le \operatorname{Lip}(h_k)\,|x-x'| ;
\end{equation*}
as $\|f\|_\sharp=\max\{\|f\|_\infty,40M\operatorname{Lip}f\}$, this gives
$\|g^{(\eta)}_k\|_\sharp\le\|h_k\|_\sharp$. For $|\alpha|\le 2\le r$ one has
$\partial^\alpha g^{(\eta)}_k=(\partial^\alpha h_k)\ast\phi_\eta$, and the same bound on sup-norms
gives $\|g^{(\eta)}_k\|_{C^2}\le\|h_k\|_{C^2}$. Consequently the constant $\hat H$ of
Lemma~\ref{10.8:lem-flux-continuity}, formed for the pair $(\vec h,\vec g^{(\eta)})$,
\begin{equation*}
\hat H:=\max\Bigl\{1,\ \max_k\max\bigl\{\|h_k\|_\sharp,\|g^{(\eta)}_k\|_\sharp,
\|\delta_\omega h_k\|_\sharp,\|\delta_\omega g^{(\eta)}_k\|_\sharp\bigr\}\Bigr\},
\end{equation*}
is bounded uniformly in $\eta$: by the bound on the sharp norm of a rotation derivative in
Lemma~\ref{10.8:lem-test-classes}, applied to $g^{(\eta)}_k$, which is supported in
$U_k\subset B(0,R')$,
\begin{equation*}
\|\delta_\omega g^{(\eta)}_k\|_\sharp\le C_\sharp(R')\,|\omega|\,\|g^{(\eta)}_k\|_{C^2}
\le C_\sharp(R')\,|\omega|\,\|h_k\|_{C^2},
\end{equation*}
and the remaining entries are bounded by $\|h_k\|_\sharp$ and $\|\delta_\omega h_k\|_\sharp$, which do
not depend on $\eta$. Finally, $\|h_k-g^{(\eta)}_k\|_{C^2}\to 0$ as $\eta\downarrow 0$: for
$|\alpha|\le 2$,
\begin{equation*}
\bigl|\partial^\alpha h_k(x)-\partial^\alpha g^{(\eta)}_k(x)\bigr|
\le\int\bigl|\partial^\alpha h_k(x)-\partial^\alpha h_k(x-y)\bigr|\,\phi_\eta(y)\,dy
\le\sup_{|y|\le\eta}\bigl|\partial^\alpha h_k(x)-\partial^\alpha h_k(x-y)\bigr| ,
\end{equation*}
and the right-hand side tends to $0$ uniformly in $x$ because $\partial^\alpha h_k$ is continuous
with compact support, hence uniformly continuous.

\emph{The flux of the mollified tuples vanishes.} Fix $0<\eta\le\eta_0$. The tuple
$\vec g^{(\eta)}=(g^{(\eta)}_1,\dots,g^{(\eta)}_n)$ consists of real $C^\infty_c$ functions with
supports in $U_k\subset B(0,R')$, where $R'=(\mathcal{R}+2\eta_0)/2$, pairwise at distance at least
$\mathsf{d}'$; it is therefore a pairwise-disjoint $C^\infty_c$ tuple and a member of
\begin{equation*}
\mathcal{T}^{(\infty)}_n(\mathcal{R}+2\eta_0,\mathsf{d}')\subset
\mathcal{T}^{(r)}_n(\mathcal{R}+2\eta_0,\mathsf{d}') .
\end{equation*}
The hypotheses are assumed for all radii and all distances, so they hold at radius
$\mathcal{R}+2\eta_0$ and distance $\mathsf{d}'$. For $\mathcal{R}'\ge\mathcal{R}+2\eta_0$ the radial
cut-off $\chi_{\mathcal{R}'}$ is admissible at radius $\mathcal{R}+2\eta_0$, since
$\chi_1(|x|/\mathcal{R}')$ is $C^\infty_c$ and identically $1$ on
$B(0,\mathcal{R}')\supseteq B(0,\mathcal{R}+2\eta_0)$, as in the proof of
Lemma~\ref{10.8:lem-smooth-tuples}. The cut-off independence of the limiting flux
(Lemma~\ref{10.8:lem-limiting-flux}), which rests on the connected rotational Ward identity,
the vanishing of the remainder and the existence of the limit $S^T_{n;\infty}$ among the
hypotheses above, therefore gives
\begin{equation*}
\lim_{j\to\infty}\mathcal{B}^{\omega,\chi_{\mathcal{R}'}}_j(\vec g^{(\eta)})
=c(\vec g^{(\eta)},\omega)\qquad\text{for every } \mathcal{R}'\ge\mathcal{R}+2\eta_0 ,
\end{equation*}
and combined with the identification of the two shell functionals
(Lemma~\ref{10.8:lem-shell-identification}) this yields
\begin{equation}\label{10.8:eq-mollified-extension-2}
\limsup_{j\to\infty}\bigl|\langle\mathfrak{B}^{\mathrm{sh}}_{\chi_{\mathcal{R}'}};
\mathcal{O}(\vec g^{(\eta)})\rangle_j\bigr|=\bigl|c(\vec g^{(\eta)},\omega)\bigr|
\qquad\text{for every } \mathcal{R}'\ge\mathcal{R}+2\eta_0 .
\end{equation}
Since $\vec g^{(\eta)}$ is a pairwise-disjoint $C^\infty_c$ tuple, the display
\eqref{10.8:eq-clustering-hypothesis-a} with $\omega$ in place of $\omega_0$ applies to it and says
that the limit inferior as $\mathcal{R}'\to\infty$ of the left-hand side of
\eqref{10.8:eq-mollified-extension-2} is $0$. By \eqref{10.8:eq-mollified-extension-2} this
function of $\mathcal{R}'$ is constant, equal to $|c(\vec g^{(\eta)},\omega)|$, for all
$\mathcal{R}'\ge\mathcal{R}+2\eta_0$, so its limit inferior is that constant. Hence
$c(\vec g^{(\eta)},\omega)=0$ for every $0<\eta\le\eta_0$.

\emph{Passage to $\vec h$.} Each $g^{(\eta)}_k$ is a real $C^2_c$ function supported in $U_k$, so
the $C^2$-continuity of the limiting flux, \eqref{10.8:eq-flux-continuity-a}, in which the
volume-free bound among the hypotheses above enters, applies with
$\vec g=\vec g^{(\eta)}$:
\begin{equation*}
\bigl|c(\vec h,\omega)-c(\vec g^{(\eta)},\omega)\bigr|
\le n^2C^\infty_n(\mathsf{d}';R')\,\hat H^{\,n-1}\bigl(1+80M+C_\sharp(R')|\omega|\bigr)
\max_k\|h_k-g^{(\eta)}_k\|_{C^2} .
\end{equation*}
The constant $\hat H$ is bounded uniformly in $\eta$ and
$\max_k\|h_k-g^{(\eta)}_k\|_{C^2}\to0$ as $\eta\downarrow0$, by the first part of the proof;
hence
\begin{equation*}
c(\vec h,\omega)=\lim_{\eta\downarrow 0}c(\vec g^{(\eta)},\omega)=0 .
\end{equation*}
By Lemma~\ref{10.8:lem-limiting-flux}, for every $\chi$ admissible at radius $\mathcal{R}$ the limit
of $\mathcal{B}^{\omega,\chi}_j(\vec h)$ exists and equals the $\chi$-independent number
$c(\vec h,\omega)$, so the infimum in \eqref{10.8:eq-mollified-extension-1} is the infimum of the
constant $|c(\vec h,\omega)|$, and \eqref{10.8:eq-mollified-extension-1} holds if and only if
$c(\vec h,\omega)=0$. Thus \eqref{10.8:eq-mollified-extension-1} holds for $\vec h$. Since $\omega$,
$n$, $\mathcal{R}$, $\mathsf{d}$ and $\vec h\in\mathcal{T}^{(r)}_n(\mathcal{R},\mathsf{d})$ were
arbitrary, the proposition follows.

For the members of $\mathcal{T}^{(r)}_n(\mathcal{R},\mathsf{d})$ all of whose entries lie in
$C^\infty_c$, the conclusion is Lemma~\ref{10.8:lem-smooth-tuples}, which uses, among the
hypotheses above, only the identification of the two shell functionals.
\end{proof}

\begin{corollary}[Equivalence of the two forms under ultraviolet bounds]
\label{10.8:cor-forms-equivalence}
Fix the diagonal $(K_j,m_j)$ and the integer $r\ge 2$ of Definition~\ref{10.8:def-rate-free}, and
write $\langle\cdot\,;\dots;\cdot\rangle_j$ for $\langle\cdot\,;\dots;\cdot\rangle_{K_j,m_j}$ and
$\mathcal{B}^{\omega,\chi}_j$ for $\mathcal{B}^{\omega,\chi}_{K_j,m_j}$. For
$\omega\in\mathfrak{so}(4)$ the \emph{clustering condition at $\omega$} is
\eqref{10.8:eq-clustering-hypothesis-a} with $\xi(x)=\omega x$ in place of $\xi(x)=\omega_0 x$, for every
$n\ge 2$ and every $\vec f\in(C^\infty_c(\mathbb{R}^4))^n$ with pairwise disjoint supports; hypothesis
(H-IR) of Definition~\ref{10.8:def-clustering-hypothesis} is this condition at one $\omega_0\neq 0$.

Assume the identification of the two shell functionals in the limit
(Lemma~\ref{10.8:lem-shell-identification}); assume, for all $\omega\in\mathfrak{so}(4)$, all
$n\ge 2$, all $\mathcal{R},\mathsf{d}>0$, all $\vec h\in\mathcal{T}^{(r)}_n(\mathcal{R},\mathsf{d})$
and all $\chi$ admissible at radius $\mathcal{R}$, hypotheses (i) and (ii) below; and assume (iii)
and (iv):
\begin{itemize}
\item[(i)] the first of the two hypotheses from which, together with (iii) below,
Lemma~\ref{10.8:lem-limiting-flux} is proved, as stated in that lemma;
\item[(ii)] the second of these hypotheses, as stated there, under which the remainder
$\mathcal{N}^{\omega,\chi}_j(\vec h)$ of that lemma tends to $0$ as $j\to\infty$;
\item[(iii)] convergence of the connected correlations along the diagonal: for pairwise disjoint real
$C^1_c$ tuples $\vec u$ the limit
\begin{equation*}
S^T_{n;\infty}(\vec u):=\lim_{j\to\infty}S^T_{n;K_j,m_j}\bigl(u_1^{(m_j)},\dots,u_n^{(m_j)}\bigr)
\end{equation*}
exists and is multilinear in the slots;
\item[(iv)] the uniform sharp-norm bound on the limit correlations: for the tuples $\vec u$ of (iii)
whose supports lie in a ball of radius $R$ and are pairwise at distance at least $\mathsf{d}$,
\begin{equation*}
|S^T_{n;\infty}(\vec u)|\le C^\infty_n(\mathsf{d};R)\prod_k\|u_k\|_{\sharp}.
\end{equation*}
\end{itemize}
Then the following three statements are equivalent:
\begin{itemize}
\item[(a)] the clustering condition holds at every $\omega\in\mathfrak{so}(4)$;
\item[(b)] the rate-free pointwise form of Definition~\ref{10.8:def-rate-free} holds;
\item[(c)] $c(\vec h,\omega)=0$ for every $\omega\in\mathfrak{so}(4)$, every $n\ge 2$, every
$\mathcal{R},\mathsf{d}>0$ and every $\vec h\in\mathcal{T}^{(r)}_n(\mathcal{R},\mathsf{d})$, where
$c(\vec h,\omega)$ is the limiting flux of Lemma~\ref{10.8:lem-limiting-flux}.
\end{itemize}
Statement (a) reads (c) on the $C^\infty_c$ tuples, statement (b) on the $C^r_c$ tuples.
\end{corollary}

\begin{proof}
\emph{(b) is equivalent to (c).} Fix $\omega$, $n\ge 2$, $\mathcal{R},\mathsf{d}>0$ and
$\vec h\in\mathcal{T}^{(r)}_n(\mathcal{R},\mathsf{d})$. By Lemma~\ref{10.8:lem-limiting-flux}, which
is proved from hypotheses (i), (ii) and (iii), for every $\chi$ admissible at radius $\mathcal{R}$ the
limit $\lim_j\mathcal{B}^{\omega,\chi}_j(\vec h)$ exists and equals the number $c(\vec h,\omega)$,
which does not depend on $\chi$. Hence $\limsup_j|\mathcal{B}^{\omega,\chi}_j(\vec h)|=|c(\vec h,\omega)|$
for every admissible $\chi$, and the infimum over admissible $\chi$ in
Definition~\ref{10.8:def-rate-free} is the infimum of the constant $|c(\vec h,\omega)|$, namely
$|c(\vec h,\omega)|$ itself. The rate-free pointwise form therefore holds at $\vec h$ if and only if
$c(\vec h,\omega)=0$. Quantifying over $\omega$, $n$, $\mathcal{R}$, $\mathsf{d}$ and $\vec h$ gives the
equivalence of (b) and (c).

\emph{(c) implies (a).} Fix $\omega\in\mathfrak{so}(4)$, $n\ge 2$ and a tuple
$\vec f\in(C^\infty_c(\mathbb{R}^4))^n$ with pairwise disjoint supports. By
Lemma~\ref{10.8:lem-test-classes}, $\vec f$ lies in $\mathcal{T}^{(r)}_n(\mathcal{R},\mathsf{d})$ for
some $\mathcal{R},\mathsf{d}>0$. Let $\mathcal{R}'$ be any radius at which $\chi_{\mathcal{R}'}$ is
admissible for $\vec f$, that is, admissible at the class radius $\mathcal{R}$; this is the case for
$\mathcal{R}'\ge\mathcal{R}$. By Lemma~\ref{10.8:lem-limiting-flux} and (c),
\begin{equation*}
\lim_{j\to\infty}\mathcal{B}^{\omega,\chi_{\mathcal{R}'}}_j(\vec f)=c(\vec f,\omega)=0 .
\end{equation*}
By the identification of the two shell functionals in the limit
(Lemma~\ref{10.8:lem-shell-identification}), the limsup over $j$ of the absolute value of the cumulant
of the shell functional $\mathfrak{B}^{\mathrm{sh}}_{\chi_{\mathcal{R}'}}$ with the $n$ observables
equals the absolute value of the limit just computed:
\begin{equation*}
\limsup_{j\to\infty}\bigl|\langle\mathfrak{B}^{\mathrm{sh}}_{\chi_{\mathcal{R}'}};
\mathcal{O}(\vec f)\rangle_j\bigr|
=\Bigl|\lim_{j\to\infty}\mathcal{B}^{\omega,\chi_{\mathcal{R}'}}_j(\vec f)\Bigr|=0
\end{equation*}
for every such $\mathcal{R}'$, where $\mathcal{O}(\vec f)$ abbreviates the $n$ slots
$\mathcal{O}(f_1^{(m_j)});\dots;\mathcal{O}(f_n^{(m_j)})$. In particular the liminf over
$\mathcal{R}'\to\infty$ in \eqref{10.8:eq-clustering-hypothesis-a}, with $\omega$ in place of
$\omega_0$, vanishes; indeed the limsup in $j$ vanishes at every radius $\mathcal{R}'\ge\mathcal{R}$,
so the limit over the radius is superfluous. As $\omega$, $n$ and $\vec f$ were arbitrary, (a) holds.

\emph{(a) implies (b).} This is the extension to $C^r$ tuples by mollification
(Proposition~\ref{10.8:prop-mollified-extension}), which is proved as an implication from the
identification of Lemma~\ref{10.8:lem-shell-identification} together with hypotheses (i)--(iv); thus
the $C^\infty_c$ tuples and the $C^r_c$ tuples carry the same statement.

The three implications together prove the equivalence of (a), (b) and (c).
\end{proof}

\begin{remark}
In terms of the conclusion step, the role of hypotheses (i)--(iv) and of
Lemma~\ref{10.8:lem-shell-identification} is played by clauses (ii-$\mathbb{R}^4$) and (iii-b)
of the main theorem together with the ultraviolet hypothesis of the conclusion step. Granting these, the
conclusion step identifies the rotation defect $D_{\omega}(\vec f)$ with the limiting flux through any
shell surrounding the supports; $c(\vec h,\omega)$ and $D_{\omega}(\vec f)$ are two notations for the
derivative of the limit state under the rotations generated by $\omega$. Hence, granting these clauses
and that hypothesis, the clustering condition at every $\omega$ is equivalent to $D_{\omega}(\vec f)=0$
for all $\omega$ and all pairwise disjoint $C^\infty_c$ tuples $\vec f$, and by the dichotomy of the
conclusion step (the stabiliser of an $\mathbb{R}^4$ limit in $O(4)$ is finite or all of $O(4)$) this
is in turn equivalent to the clustering condition at a single $\omega_0\neq 0$, that is, to
hypothesis (H-IR).

From Lemma~\ref{10.8:lem-shell-identification} alone, that is, granting only the hypotheses from which
that lemma is proved and none of (i)--(iv), the clustering condition at every
$\omega$ gives the rate-free pointwise form on the $C^\infty_c$ members of the classes
$\mathcal{T}^{(r)}_n(\mathcal{R},\mathsf{d})$, by the statement that clustering implies the rate-free
form on smooth tuples, and is in that sense the stronger of the two forms; no converse, and no
extension to tuples in $C^r_c\setminus C^\infty_c$, is available without hypotheses (i)--(iv) or the
bound \eqref{10.8:eq-smooth-tuples-a}.
\end{remark}

\begin{remark}[Two-insertion tensor clustering at the canonical rate]\label{10.8:rem-two-insertion}
Clause (vi) carries the far-sphere flux hypothesis $(\mathrm{IR}\text{-}\mathrm{fl})_{\omega_0}$, which
states that the flux of the lattice rotation current through far spheres vanishes along one sequence of
radii. A sufficient condition for it is displayed in \eqref{10.8:eq-clustering-hypothesis-a}, in which one
far-away stress-tensor insertion decorrelates from the observables. Besides these and the rate-free
pointwise form of Definition~\ref{10.8:def-rate-free}, there are two-point hypotheses from which the
rate-free pointwise form follows.
They are stronger than the rate-free pointwise form and, under the inputs granted below,
than the far-sphere flux hypothesis; they are of a different kind, and they are neither hypotheses of
Theorem~0 nor proved in this book. They concern truncated correlations of two tensor-class
insertions, localised in separated unit cubes and taken in the finite-volume measures
$\mu_{K_j,m_j}$ themselves, along the diagonal.

Write $\mathfrak{F}_{\chi} = B^4_{\chi} + J^4_{\chi}$ for the
total dimension-four annulus functional that enters $\mathcal{B}^{\omega,\chi}_{K,m}$ in
Definition~\ref{10.8:def-rate-free}, with stress-tensor part $B^4_{\chi}$ and Jacobian part $J^4_{\chi}$.
For the part $B^4_{\chi}$ an insertion is a traceless-weighted
combination of the components of the smeared stress-tensor density $\Theta^{\varepsilon}_{\nu\sigma}$.
For the Jacobian part $J^4_{\chi}$ it is the insertion carried by a tuple of plaquettes lying in a single
plane.

The bounds in question estimate the truncated expectation $\langle\cdot\,;\cdot\rangle_{K_j,m_j}$ of
two such insertions, localised in unit cubes $\Delta_1,\Delta_2$ at mutual distance $D$ (in this remark
$D$ denotes this distance), at the canonical rate $D^{-8}$, with no surplus in the exponent. They hold
locally uniformly in the cutoff, that is, along the diagonal for all indices $j$ beyond a floor index
depending on a parameter of the bound.

These bounds come in two forms, the first used together with uniqueness of the vacuum $\Omega$, the
second without it.
\begin{itemize}
\item Uniqueness of $\Omega$ together with the first form implies the rate-free pointwise form of
Definition~\ref{10.8:def-rate-free}. This implication is the statement of this book deriving the
rate-free pointwise form from uniqueness of $\Omega$ and the first form; it is proved there as an
implication from six statements: the transcription of $\mathfrak{F}_{\chi}$ into the two parts above,
the structure of the stress-tensor density, reflection positivity of the lattice measure, and the three
further statements listed among its hypotheses, one of which has uniqueness of $\Omega$ on its list of
statements not asserted and another of which is a lemma of this book.
\item The second form alone, without uniqueness of $\Omega$, implies the rate-free pointwise form. This
implication is the statement of this book deriving the rate-free pointwise form from the second form;
it is proved there as an implication from four of the six statements just listed: the transcription of
$\mathfrak{F}_{\chi}$, the structure of the stress-tensor density, reflection positivity of the lattice
measure, and the remaining one of the three further statements of the preceding implication.
\end{itemize}

The rate-free pointwise form concerns a single insertion of $\mathfrak{F}_{\chi}$ and carries no rate: it
is an infimum over admissible cut-offs $\chi$. The two-insertion bounds supply it through reflection
positivity.

Grant the hypotheses of Corollary~\ref{10.8:cor-forms-equivalence}: the ultraviolet bounds under which it
is stated and the identification, along the diagonal and in the form in which that corollary uses it, of
$\mathcal{B}^{\omega,\chi}_{K_j,m_j}$ with the truncated correlation of the shell functional
$\mathfrak{B}^{\mathrm{sh}}_{\chi}$ with the observables. Then each of the two conditions in the list
above (uniqueness of $\Omega$ with the first form, or the second form alone) implies the rate-free
pointwise form and hence, by that corollary, the far-sphere flux hypothesis
$(\mathrm{IR}\text{-}\mathrm{fl})_{\omega}$, that is, $(\mathrm{IR}\text{-}\mathrm{fl})_{\omega_0}$ with
$\omega_0$ replaced by an arbitrary $\omega\in\mathfrak{so}(4)$.

The converse implications are not available. Neither two-point condition is implied by the
rate-free pointwise form or by the far-sphere flux hypothesis.

Among the sufficient conditions for the far-sphere flux hypothesis collected in the remark of this book
listing such conditions is
uniqueness of $\Omega$ together with a bound on the variance of the shell functional. It is of the same
kind as uniqueness of $\Omega$ together with the first form of the two-insertion bounds. We do not
identify the two conditions.
\end{remark}

\begin{remark}[The axis-tube variant of the far-sphere flux hypothesis]\label{10.8:rem-axis-tube}
Let $\omega_0 \in \mathfrak{so}(4)$, $\omega_0 \neq 0$, the rotation field $\xi(x) = \omega_0 x$ and the profile
$\chi_1$ be those of Definition~\ref{10.8:def-radial-cutoff}. Let $(K_j, m_j)$ be the joint diagonal of
clause (ii-$\mathbb{R}^4$) fixed in Definition~\ref{10.8:def-clustering-hypothesis}. This is the diagonal
along which the far-sphere flux hypothesis $(\mathrm{IR}\text{-}\mathrm{fl})_{\omega_0}$ of clause (vi) is
stated (Definition~\ref{10.8:def-clustering-hypothesis}): the flux of the lattice rotation current through
far spheres, in truncated correlation with every tuple of curvature observables with pairwise disjoint
supports, vanishes along one sequence of radii. The pointwise stress decay faster than $|x|^{-4}$
(Definition~\ref{10.8:def-pointwise-decay}) is a sufficient condition for it
(Lemma~\ref{10.8:lem-pointwise-implies}). Let
$\langle \cdot\,; \dots ; \cdot \rangle_j$ denote the truncated expectation with respect to
$\mu_{K_j, m_j}$. For $\mathcal{R} > 0$ put $\theta_{\mathcal{R}}(s) := \chi_1(|s|/\mathcal{R})$,
$s \in \mathbb{R}$. The \emph{axis-tube deformation field} $\eta^{\natural}_{\mathcal{R}}$ is defined
componentwise by
\begin{equation}\label{10.8:eq-axis-tube-1}
(\eta^{\natural}_{\mathcal{R}})^{\rho}(x) := \xi^{\rho}(x) \prod_{\tau \neq \rho} \theta_{\mathcal{R}}(x_\tau),
\qquad \rho = 1, \dots, 4,
\end{equation}
and is read on $\mathbb{T}_m$ whenever $4\mathcal{R} < L^m$. This is the field of
Theorem~\ref{10.5:thm-reduction-without-shell}. It is $C^\infty$ and periodic, and it is own-coordinate-free
in the sense of Definition~\ref{10.5:def-own-coordinate-cutoffs}: its
component $(\eta^{\natural}_{\mathcal{R}})^{\rho}$ does not depend on the coordinate $x_\rho$, since the
product omits the factor $\theta_{\mathcal{R}}(x_\rho)$, and $\xi^{\rho}(x) = \sum_{\tau} (\omega_0)_{\rho\tau} x_\tau$
does not involve $x_\rho$ because $(\omega_0)_{\rho\rho} = 0$. The field equals $\xi$ on the box
\begin{equation*}
B_{\mathcal{R}} := \{ x : |x_\mu| \le \mathcal{R} \text{ for every } \mu \}.
\end{equation*}

Let $\mathfrak{B}^{\natural}_{\mathcal{R}}$ be the tube shell functional: $-g_0^{-2}$ times the smearing,
against the strain weight of $\eta^{\natural}_{\mathcal{R}}$, of a lattice tensor density of the channel
$\mathbf{6}$ of Theorem~\ref{10.4:thm-insertion-correlation}. Only the following
properties of it are used. Its weight is bounded by $C|\omega_0|$ and has Lipschitz constant at most
$C|\omega_0|/\mathcal{R}$, with a constant $C$ independent of $\mathcal{R}$, so that for $\mathcal{R} \ge 1$
its $\sharp$-norm is bounded by a multiple of $|\omega_0|$ independent of $\mathcal{R}$; it vanishes on
$B_{\mathcal{R}}$.
It carries the large weight $x_0 = g_0^{-2}$, the bare inverse coupling, on an
off-diagonal, hence traceless, tensor density, and it has no admixture free of $x_0$.

The \emph{axis-tube variant} of the far-sphere flux hypothesis asserts the following. For every $n \ge 2$ and
every tuple $\vec f = (f_1, \dots, f_n) \in (C^\infty_c(\mathbb{R}^4))^n$ with pairwise disjoint supports,
\begin{equation}\label{10.8:eq-axis-tube-2}
\liminf_{\mathcal{R} \to \infty} \limsup_{j \to \infty}
\bigl| \langle \mathfrak{B}^{\natural}_{\mathcal{R}}; \mathcal{O}(f_1); \dots; \mathcal{O}(f_n) \rangle_j \bigr| = 0 .
\end{equation}
Here the inner expression is defined eventually in $j$, namely for all $j$ with $4\mathcal{R} < L^{m_j}$ and
$\bigcup_i \operatorname{supp} f_i \subset \,]-L^{m_j}, L^{m_j}[^4$, the $f_i$ being read on
$\mathbb{T}_{m_j}$ through their periodic lifts $f_i^{(m_j)}$ as in
Definition~\ref{10.8:def-clustering-hypothesis}. It stands beside
the far-sphere flux hypothesis $(\mathrm{IR}\text{-}\mathrm{fl})_{\omega_0}$ as follows.

(a) Suppose that the tensor density of which $\mathfrak{B}^{\natural}_{\mathcal{R}}$ is the smearing is
identified with the off-diagonal components $\Theta_{\nu\sigma}$, $\nu \neq \sigma$, of the $x_0$-weighted
stress-tensor density of \eqref{10.8:eq-pointwise-decay-a}. Under this identification,
\eqref{10.8:eq-axis-tube-2} is implied by the pointwise stress decay faster than $|x|^{-4}$
(Definition~\ref{10.8:def-pointwise-decay}) in the following reading: only its instances for the
off-diagonal components are used, and only at unit cubes far from the sources (at quotient distance at
least $\mathcal{R}/2$, as arranged below), and the bound is taken on each torus $\mathbb{T}_{m_j}$ with one
constant $C(\vec f)$ for all unit cubes at quotient distance $D \ge 1$ from
$\bigcup_i \operatorname{supp} f_i$ and with no lower bound on $j$ depending on $D$.
The implication goes through the tube count. Split $\mathfrak{B}^{\natural}_{\mathcal{R}}$, as in the proof
of Lemma~\ref{10.8:lem-pointwise-implies}, by a Lipschitz partition of unity subordinate to unit cubes, at
the cost of a fixed factor in the $\sharp$-norm of the weight; each piece is then bounded by
\eqref{10.8:eq-pointwise-decay-a}. Since $\chi_1$ vanishes on $[2,\infty[$, the weight is supported within
one lattice spacing of the support of $\eta^{\natural}_{\mathcal{R}}$, hence in the set where
$|x_\tau| \le 2\mathcal{R} + 1$ for all coordinates $\tau$ but at most one, which is contained in the union of
the box $B_{2\mathcal{R}+1}$ and the eight half-tubes along the
positive and negative coordinate axes. For a unit cube $\Delta$ call its height the integer part of
$\min_{x \in \Delta} \max_\mu |x_\mu|$. Let $\mathcal{R}$ be so large, depending on $\vec f$, that each unit
cube meeting the support of the weight lies at quotient distance at least $\max\{1, \mathcal{R}/2\}$ from
$\bigcup_i \operatorname{supp} f_i$, and at least $t/2$ if its height $t$ is at least $2\mathcal{R}$.
The cubes of height below $2\mathcal{R}$ meet $B_{2\mathcal{R}+1}$ and number at most $C\mathcal{R}^4$; for
each $t \ge 2\mathcal{R}$, at most $C\mathcal{R}^3$ cubes of height $t$ meet the eight half-tubes.
Hence, with $\delta > 0$ the exponent of Definition~\ref{10.8:def-pointwise-decay} and $C$
independent of $\mathcal{R}$ and $j$,
\begin{equation}\label{10.8:eq-axis-tube-3}
\begin{split}
\limsup_{j \to \infty}
\bigl| \langle \mathfrak{B}^{\natural}_{\mathcal{R}}; \mathcal{O}(f_1); \dots; \mathcal{O}(f_n) \rangle_j \bigr|
&\le C\mathcal{R}^4 (\mathcal{R}/2)^{-4-\delta}
+ C\mathcal{R}^3 \sum_{t \ge 2\mathcal{R}} t^{-4-\delta} \\
&\le C \mathcal{R}^{-\delta}.
\end{split}
\end{equation}
The first term counts the unit cubes of height below $2\mathcal{R}$. The second counts the unit cubes along
the axis tubes, at height $t$. The right side tends to zero as $\mathcal{R} \to \infty$, which gives
\eqref{10.8:eq-axis-tube-2}. The decay threshold is $|x|^{-4}$, the same as for the radial shell in
Lemma~\ref{10.8:lem-pointwise-implies}.

(b) As stated, the rate-free pointwise form on $C^r$ tuples (Definition~\ref{10.8:def-rate-free}) does not by
itself yield \eqref{10.8:eq-axis-tube-2}. The field $\eta^{\natural}_{\mathcal{R}}$ is not of the form $\chi\,\xi$ with
$\chi$ a scalar cut-off admissible at some radius: it is not a scalar function times $\xi$, and it is not
compactly supported. Hence the infimum in
that form does not range over it.

(c) Two further bounds of the kind in Definition~\ref{10.8:def-pointwise-decay}, for unit cubes $\Delta$ at
distance $D$ from the sources, are asserted uniformly in $j$ only locally, and neither, as stated, implies
\eqref{10.8:eq-axis-tube-2}. The first is uniform in $j$ only for unit cubes $\Delta$ which,
together with a prescribed set of cubes attached to $\Delta$, lie in a fixed ball $B(0,R_1)$, and only from a
threshold $j \ge j_1(R_1)$ on, where $j_1(R_1)$ depends on $R_1$. The second is uniform in $j$ only for
distances $1 \le D \le D^*$, with $D^*$ a fixed upper bound, and only from a threshold $j \ge j_2(D^*)$ on,
where $j_2(D^*)$ depends on $D^*$. The tube functional
at cutoff $j$ reaches quotient distance $L^{m_j}$, beyond either restriction.

(d) Tuple by tuple and at the generator $\omega_0$, \eqref{10.8:eq-axis-tube-2} is equivalent to the instance
at $\omega = \omega_0$ of the rate-free pointwise form once both
sides' ultraviolet inputs are granted. On the side of the rate-free form these are the ultraviolet hypotheses
of Corollary~\ref{10.8:cor-forms-equivalence}, taken for admissible cut-offs $\chi$. On the side of the
axis-tube variant they are:
\begin{itemize}
\item the versions for $\eta^{\natural}_{\mathcal{R}}$ of the rotational Ward identity
(Theorem~\ref{10.1:thm-rotational-identity});
\item the inputs for $\eta^{\natural}_{\mathcal{R}}$ of the reduction without a shell term
(Theorem~\ref{10.5:thm-reduction-without-shell}): the version for $\eta^{\natural}_{\mathcal{R}}$ of the
vector package in output form (Definition~\ref{10.5:def-vector-package-output}), the hypotheses of the
reduction theorem (Theorem~\ref{10.5:thm-reduction}, whose proof is incomplete at one step) other than its
vector-package and anisotropy inputs, and an anisotropy input with exponent greater than $-\tfrac12$, met by
the one-loop anisotropy inequality with exponent $\kappa = 1$
(Corollary~\ref{10.6:cor-anisotropy-inequality}) under that corollary's hypotheses, since $1 > -\tfrac12$;
\item an ultraviolet bound taken for the deformation field $\eta^{\natural}_{\mathcal{R}}$.
\end{itemize}
Under these inputs both forms then say that the limiting flux $c(\vec f, \omega_0)$ of
Corollary~\ref{10.8:cor-forms-equivalence} vanishes, for every tuple
$\vec f \in (C^\infty_c(\mathbb{R}^4))^n$ with pairwise disjoint supports.

(e) The far-sphere flux hypothesis $(\mathrm{IR}\text{-}\mathrm{fl})_{\omega_0}$ in the form of
Definition~\ref{10.8:def-clustering-hypothesis}, that is, along one sequence of radii and with the radial
shell functional, does not by itself yield \eqref{10.8:eq-axis-tube-2}.

(f) With the exponent $\kappa = 1$ in the one-loop anisotropy inequality, the axis-tube variant is not
needed for clause (vi).
\end{remark}

We now assemble clause (vi) of Theorem~\ref{3.1:thm-main}. The clause concerns $N = 2$ in the
regime of clause (iii), and only the $\mathbb{R}^4$ limits of (ii-$\mathbb{R}^4$). Under hypothesis
(H-IR(A)) of Theorem~\ref{3.1:thm-main} it asserts the following for every such limit
$S^T_{n;\infty}$, every $n \ge 2$, every $R \in O(4)$ and every tuple $\vec f = (f_1, \dots, f_n)$
of real functions in $C^1_c(\mathbb{R}^4)$ with pairwise disjoint supports:
$S^T_{n;\infty}(R \cdot \vec f\,) = S^T_{n;\infty}(\vec f\,)$, where $R \cdot \vec f$ is the
tuple of the functions $f_i \circ R^{-1}$.
The deduction uses five groups of inputs.
First, it uses the rotational Ward identity, with the split of its $O(4)$-breaking term into spin,
orbital and shell parts, and the dimension-six classification. Both are proved earlier in this
chapter.
Second, it uses the vanishing in the limit of the spin-channel insertion. This is proved as an
implication from the correlation theorem (stated; proof incomplete), the one-lattice Lipschitz step
and the statements that put their inputs in the form these two theorems require.
Third, it uses the vector-channel statement, Proposition~\ref{10.2:prop-vector-channel}. This is
proved as an implication from the vector package (proved as an implication); the half of the
stress-tensor-insertion correlation theorem that transports the insertion along the flow (proved
as an implication); the half of that theorem that puts the insertion in the form to which the
correlation theorem applies (stated; proof incomplete); the reduction theorem,
Theorem~\ref{10.5:thm-reduction} (stated; proof incomplete); and the one-loop anisotropy inequality
with exponent one (proved as an implication from the hypotheses listed in
Corollary~\ref{10.6:cor-anisotropy-inequality}).
Fourth, it uses clause (ii-$\mathbb{R}^4$), the per-ball bound (iii-b) and $K_j \to \infty$ from
clause (0).
Fifth, it uses hypothesis (H-IR(A)), the infrared decay hypothesis displayed in
Theorem~\ref{3.1:thm-main}. It is read on the traceless lattice stress tensor $\Theta_{\nu\sigma}$,
weighted by the coordinate $x_0$ of the point of insertion, together with its $x_0$-free symmetric
admixtures of dimension four. For every $n \ge 2$ and every tuple $(f_1, \dots, f_n)$ as above it
asks that the connected correlation of $\Theta$ at distance $\mathcal{R}$ from the supports with
$\mathcal{O}(f_1), \dots, \mathcal{O}(f_n)$ be $o(\mathcal{R}^{-4})$, uniformly in the cutoff (and
in $m$ along the joint diagonal), in the sense that the integral over the shell
$\mathcal{R} \le |y| \le 2\mathcal{R}$ of the absolute value of this correlation, after the limit
in the cutoff, tends to zero along one sequence of radii $\mathcal{R} \to \infty$.
In this form (H-IR(A)) is sufficient for the one-generator shell statement used in the deduction:
for one real antisymmetric $\omega_0 \neq 0$, every $n \ge 2$ and every such tuple of
$C^\infty_c$ functions, the connected correlation with $\mathcal{O}(f_1), \dots, \mathcal{O}(f_n)$
of the shell functional at radius $\mathcal{R}$ built on $\omega_0$ has, in absolute value, a
$\limsup$ along the diagonal of (ii-$\mathbb{R}^4$) whose $\liminf$ as $\mathcal{R} \to \infty$ is
zero; no uniformity in the cutoff is asked beyond this order of limits.
The first-order change of a limit $S^T_{n;\infty}(\vec f\,)$ under the rotations generated by an
antisymmetric matrix $\omega$ is linear in $\omega$ and $W_4$-covariant. Hence the stabiliser in
$O(4)$ of a limit is either a finite group containing $W_4$ or all of $O(4)$. Since $W_4$ and one
one-parameter subgroup generate $O(4)$, vanishing of this change for $\omega_0$ gives $O(4)$. Call
the first three groups of inputs the ultraviolet half. Given it, the one-generator shell statement
is necessary as well as sufficient. Both facts are proved as an implication from
(ii-$\mathbb{R}^4$), (iii-b) and the ultraviolet half.

Clause (vi) is not proved. It is a complete deduction from hypothesis (H-IR(A)), about whose truth
nothing is asserted, and from the statements of the first, second and fourth groups, the
vector-channel statement and the two facts just stated, each proved or proved as an implication.
Among the statements on which these depend are the half of the stress-tensor-insertion
correlation theorem that puts the insertion in the required form and the reduction theorem (both
stated; proof incomplete), the statements used to prove that half (stated; proof incomplete), and
the reading of the correlation theorem's collar clauses on the inner collar that is used by the
vector package (stated; proof incomplete). Three further ones are
shared with clauses (i), (iii) and (v): the correlation theorem (stated; proof incomplete), the
cluster bound (stated; proof incomplete) and the carry lemma (proved as an implication).

The following decay bound on the far stress-tensor insertion is sufficient for (H-IR(A)). There must
be $\delta > 0$ and $C(\vec f\,)$ with this property: for every unit cube $\Delta$ at distance
$D \ge 1$ from the supports of the $f_i$, and every Lipschitz $h$ supported in $\Delta$, the
connected correlation of $\Theta_{\nu\sigma}(h)$ with $\mathcal{O}(f_1), \dots, \mathcal{O}(f_n)$
is bounded in absolute value, uniformly in the cutoff, by
$C(\vec f\,) \|h\|_\sharp D^{-4-\delta}$. This is a decay strictly faster than $|x|^{-4}$. A decay
slower than $|x|^{-4}$, even uniform in the cutoff, does not suffice. At tree level the decay is
$|x|^{-8}$. No decay of this insertion is asserted.

Clause (vi) asserts none of the following: anything about tensor multiplets; rotation invariance
of block probes or of Wilson loops; anything on a torus beyond the finite-volume covariance of
commensurate rotated towers; that (H-IR(A)) holds; clause (vi) along the tuned family of clause (v)
on $\mathbb{R}^4$; uniqueness of $\Omega$, or clustering of the correlations of the curvature
observables themselves, from which (H-IR(A)) is a different statement and which it is not asserted
to imply. A second route, which compares two commensurate orientations directly and needs
independence of the torus shape in place of (H-IR(A)), is not adopted.

\subsection*{Inputs used by statement}

Of the clauses of Theorem~\ref{3.1:thm-main}, only clause~(vi) is treated in this chapter.
It is stated under the far-sphere flux hypothesis $(\mathrm{IR}\text{-}\mathrm{fl})_{\omega_0}$
displayed in the statement of clause~(vi): the flux of the lattice rotation current through far
spheres vanishes along one sequence of radii. The infrared decay condition displayed beside it,
that one far-away stress-tensor insertion decorrelates from the observables faster than the
inverse fourth power of the distance, uniformly in the cutoff, is sufficient for it. This
hypothesis is not listed below. Two kinds of input
used by statement in the proofs of this chapter are
listed below. The first are statements published by Ba{\l}aban, each with its reference. The
second are statements of this book whose own proof here is incomplete or outstanding, each with
that standing. They are arranged by the clauses of Theorem~\ref{3.1:thm-main}, in the order of its
statement.

\begin{description}
\item[(0)] Not treated in this chapter.
\item[(i)] Not treated in this chapter.
\item[(iii)] Not treated in this chapter.
\item[(ii)] Not treated in this chapter.
\item[(iv)] Not treated in this chapter.
\item[(v)] Not treated in this chapter.
\item[(vi)] Published statements cited by statement:\par
\begin{itemize}
\item \cite[Proposition~1.1, (1.89)--(1.90)]{B-CMP95}, in
Lemma~\ref{10.6:lem-representation-covariance}, Proposition~\ref{10.6:prop-regulator-admissible}
and Theorem~\ref{10.6:thm-one-loop-drift};
\item \cite[(1.42)--(1.45)]{B-CMP95}, in Lemmas~\ref{10.6:lem-background-generators}
and~\ref{10.6:lem-representation-covariance}, Proposition~\ref{10.6:prop-regulator-admissible}
and Theorem~\ref{10.6:thm-one-loop-drift};
\item \cite[(1.95), (1.97)]{B-CMP95}, in Lemma~\ref{10.6:lem-representation-covariance};
\item \cite[(1.99)--(1.101)]{B-CMP95}, in Lemma~\ref{10.6:lem-representation-covariance},
Proposition~\ref{10.6:prop-regulator-admissible} and Theorem~\ref{10.6:thm-one-loop-drift};
\item \cite[Proposition~3]{B-CMP98} and \cite[(139), (143)]{B-CMP98}, in
Lemma~\ref{10.6:lem-uniform-majorant};
\item \cite[(3.3)--(3.10)]{B-CMP99b}, in Theorem~\ref{10.6:thm-one-loop-drift}; (3.10) also in
Lemma~\ref{10.6:lem-background-generators};
\item \cite[(3.19)--(3.26)]{B-CMP99b}, in Theorem~\ref{10.6:thm-one-loop-drift}; (3.19) also
in Lemma~\ref{10.6:lem-background-generators};
\item \cite[(3.25)--(3.27)]{B-CMP99b}, in Lemma~\ref{10.6:lem-representation-covariance};
\item \cite[(0.1)]{B-CMP109}, in the conventions on lattices, constant curvatures and affine
fields preceding Lemma~\ref{10.6:lem-constant-fields};
\item \cite[(0.5)--(0.16)]{B-CMP109}, in Theorem~\ref{10.6:thm-one-loop-drift}; among these,
(0.5)--(0.12) also in Lemma~\ref{10.6:lem-uniform-majorant}, (0.10)--(0.12) also in
Lemmas~\ref{10.6:lem-constant-fields} and~\ref{10.6:lem-background-generators}, (0.14)--(0.16)
also in Lemma~\ref{10.6:lem-half-trace}, and (0.15) also in
Proposition~\ref{10.6:prop-bare-trace};
\item \cite[(1.4)--(1.5)]{B-CMP109}, in Theorem~\ref{10.6:thm-one-loop-drift}; (1.5) also in
Lemmas~\ref{10.6:lem-half-trace} and~\ref{10.6:lem-one-step-identities} and
Proposition~\ref{10.6:prop-truncated-ring};
\item \cite[(2.1)--(2.12)]{B-CMP109}, in Theorem~\ref{10.6:thm-one-loop-drift}; (2.2)--(2.5)
also in Lemma~\ref{10.6:lem-one-step-identities}, (2.2)--(2.3) and (2.5) also in
Lemma~\ref{10.6:lem-half-trace}, and (2.4)--(2.5) also in
Lemma~\ref{10.6:lem-background-generators};
\item \cite[(2.16)]{B-CMP109}, in Lemma~\ref{10.6:lem-half-trace};
\item \cite[(3.31)]{B-CMP109}, in Definition~\ref{10.3:def-exponents-rates};
\item \cite[(4.14)]{B-CMP109}, in Lemmas~\ref{10.3:lem-kernel-decay},
\ref{10.6:lem-second-moment} and~\ref{10.6:lem-half-trace}, and in the remark that inner maps
enter only through their linearisations;
\item \cite[(4.33)]{B-CMP109}, in Definition~\ref{10.3:def-kernel-comb-gauge};
\item \cite[(4.35)--(4.37), (5.12)--(5.15), (5.30)--(5.42), \S5]{B-CMP109}, in
Lemma~\ref{10.4:lem-typed-slot-kernel};
\item \cite[(4.40)--(4.41), (5.37)--(5.38)]{B-CMP109}, in
Lemma~\ref{10.4:lem-coefficient-bounds};
\item \cite[clause~(iv)]{B-CMP109}, in Lemma~\ref{10.3:lem-real-values-unchanged};
\item \cite[(1.36), (1.42), (1.43)]{B-CMP116} and \cite[(2.1)--(2.14)]{B-CMP116}, in
Proposition~\ref{10.4:prop-vertex-response};
\item \cite[(2.3), (2.22)]{B-CMP116}, in the statement on the cut-free one-loop numbers that
precedes Lemma~\ref{10.4:lem-terminal-matching};
\item \cite[(2.26)]{B-CMP119}, in Lemma~\ref{10.3:lem-vector-rim};
\item \cite[(2.28)]{B-CMP119}, in Notations~\ref{10.3:not-units-insertion}
and~\ref{10.4:not-standing-notation} and Lemma~\ref{10.3:lem-relocalisation};
\item \cite[(2.29)]{B-CMP119}, in Lemma~\ref{10.3:lem-relocalisation};
\item \cite[(2.32)]{B-CMP119}, in Definition~\ref{10.3:def-vector-production-bounds} and
Lemma~\ref{10.3:lem-vector-refiling};
\item \cite[(3.66)]{B-CMP119}, in Lemma~\ref{10.4:lem-typed-irrelevance};
\item \cite[(3.67)]{B-CMP119}, in Definition~\ref{10.3:def-exponents-rates};
\item \cite[(1.63)]{B-CMP122a}, in Lemma~\ref{10.3:lem-consumer-domain};
\item \cite[(1.16)--(1.17), (1.36), (1.46)--(1.47)]{B-CMP122b}, in
Proposition~\ref{10.4:prf-sourced-format};
\item \cite[(1.59), (1.61)]{B-CMP122b}, in Lemma~\ref{10.3:lem-vector-rim};
\item \cite[(1.66)]{B-CMP122b}, in Lemmas~\ref{10.3:lem-relocalisation}
and~\ref{10.3:lem-vector-rim} and in the discussion of the decay rates of sources, sharp
families and defects.
\end{itemize}
Statements of this book whose proofs are incomplete or outstanding:\par
\begin{itemize}
\item the domination estimate, with its abelian case and its form on the torus, stated after
Theorem~\ref{10.1:thm-rotational-identity} --- stated; proof incomplete;
\item the production bounds for the vector part under the split rule,
Definition~\ref{10.3:def-vector-production-bounds} --- proof outstanding;
\item the vector package (bulk, freezing and rim bounds for the vector part of the run),
Proposition~\ref{10.3:prop-vector-package} --- stated; proof incomplete;
\item the inner-collar cores of the vector part under the extended correlation theorem,
Proposition~\ref{10.3:prop-cores-extended} --- stated; proof incomplete;
\item the transfer of the scalar lemmas to a tensor slot, stated before
Proposition~\ref{10.4:prop-vertex-response} --- stated; proof incomplete;
\item the response of one step to the tensor vertex,
Proposition~\ref{10.4:prop-vertex-response} --- stated; proof incomplete;
\item the typed step with a marginal tensor slot, Proposition~\ref{10.4:prop-typed-step}
--- stated; proof incomplete;
\item the re-filing of typed marginal families, Lemma~\ref{10.4:lem-typed-refiling}
--- stated; proof incomplete;
\item the sourced format of the local tangent run, Proposition~\ref{10.4:prf-sourced-format}
--- stated; proof incomplete;
\item the statement on the cut-free one-loop numbers that precedes
Lemma~\ref{10.4:lem-terminal-matching} --- stated; proof incomplete;
\item the vector package in output form,
Definition~\ref{10.5:def-vector-package-output} --- proof outstanding;
\item the reduction theorem, Theorem~\ref{10.5:thm-reduction} --- stated; proof incomplete;
\item the two-sided decay of the tensor normalisation,
Corollary~\ref{10.6:cor-normalisation-decay} --- stated; proof incomplete;
\item the statement on rotations in the planes containing the time axis, following
Theorem~\ref{10.7:thm-rotation-representation} --- stated; proof incomplete.
\end{itemize}
\end{description}

No statement listed above is deduced from clause (vi), or from any statement whose proof uses
clause (vi). The published statements are cited from the papers as stated there, and none of
them depends on a statement of this book. The dependence of clause (vi) on the statements
listed therefore contains no cycle.

%% ---- end inlined file: chapters/c10 ----

\clearpage
\setcounter{section}{10}
\refstepcounter{section}
\section*{Chapter~\thesection\quad Assembly of the proof of Theorem~0 from its clauses\addcontentsline{toc}{section}{\protect\numberline{\thesection}Assembly of the proof of Theorem~0 from its clauses}\label{ch:11}}
%% ---- begin inlined file: chapters/c11 ----
Chapters~\ref{ch:4}--\ref{ch:10} have treated the seven clauses of Theorem~\ref{3.1:thm-main} one at a time; this
chapter records, once, which clause rests on which, derives the standing of each clause from its
chapter, and counts the statements that are not proved outright.

\begin{description}
\item[(0)] The statement on the coupling flow and clause~(i) form a single induction on the level
$k$. At each step, the coupling $g_{k+1}$ and the large-field threshold of level $k+1$ are fixed
from data of index at most $k$, and this is done before the large-field half of step $k+1$ is
carried out. The induction is therefore not circular (Chapter~\ref{ch:4}).
\item[(i)] Clause~(i) belongs to the same induction on the level as (0) and takes from it the
coupling of each level (Chapter~\ref{ch:4}).
\item[(iii)] The cutoff-uniform bounds on the connected Schwinger functions rest on (0) and (i)
(Chapters~\ref{ch:5}, \ref{ch:6}).
\item[(ii)] Clause (ii) is read in three limit regimes (Chapter~\ref{ch:7}).
\begin{itemize}
\item (ii-$\mathbb{T}$-a), the torus part along tuned couplings, has exactly the inputs of (v)
and is obtained from (v).
\item (ii-$\mathbb{T}$-b), the torus part along the first-passage cutoffs $K(g_0,g)$, rests on
(iii-a) and on reflection positivity of the lattice measure.
\item (ii-$\mathbb{R}^4$), the part on $\mathbb{R}^4$, rests on clause (iii), through the
correlation theorem and the per-ball volume-free bound (iii-b), on reflection positivity of the
lattice measure, and on (0).
\end{itemize}
\item[(iv)] Clause (iv), the curvature two-point witness, is stated inside the regime of (iii),
uniformly in the cutoff and in the volume, under the further one-sided condition
$g \le \gamma_W$ with $\gamma_W \le \gamma_J$, and rests on three inputs:
\begin{itemize}
\item the correlation theorem, read at the \emph{witness depth} $s(\mathsf{d})$, the number of
levels below the cutoff that the separation $\mathsf{d}$ of the two supports determines;
\item clause (0), including its torus-independence lemma, which enters when the two-point
function is compared on tori of different sizes; that comparison also uses reflection positivity
of the lattice measure and the lower stability bound of clause~(i);
\item its own statements on large-field regions, formulated in Chapter~\ref{ch:8}.
\end{itemize}
Its consequences in the limits of (ii), those on $\mathbb{R}^4$ included, take clause (ii) and
reflection positivity of the lattice measure as further inputs.
\item[(v)] Uniqueness up to one real parameter rests on (iii), applied along the prescribed
scaffold of couplings and cutoffs of Chapter~\ref{ch:9}, and on the statements proper to that
chapter. The corollary of the uniqueness theorem also takes (0) as an input.
\item[(vi)] The statement on $\mathrm{O}(4)$ invariance rests on the following inputs
(Chapter~\ref{ch:10}):
\begin{itemize}
\item clauses (0), (i), (ii) and (iii);
\item the one-lattice Lipschitz step of (v-e), of which only the statement is used;
\item its displayed infrared hypothesis, the far-sphere flux hypothesis (IR-fl)$_{\omega_0}$:
the flux of the lattice rotation current through far spheres vanishes along one sequence of
radii. A sufficient form of it, displayed beside it, is the infrared decay condition that one
far-away stress-tensor insertion decorrelates from the observables faster than the inverse fourth
power of the distance, uniformly in the cutoff. It is a hypothesis of clause~(vi) and is derived
nowhere in this book.
\end{itemize}
\end{description}

The statements on Ba{\l}aban's block laws collected in Remark~\ref{3.2:rem-block-laws} are an
input of no clause. The dependency relation just listed has no cycle. Among the clauses it orders the
pair (0), (i) before (iii), (iii) before each of (ii), (iv) and (v), and (ii) and (v) before
(vi); (iv) is read by no other clause.
Inside the group (ii), (iv), (v) the only dependences are that of (ii-$\mathbb{T}$-a) on (v) and
that of the consequences of (iv) in the limits on (ii); neither (ii) nor (v) takes (iv) as an input.
Inside the pair (0), (i) there is no cycle because the two are one induction on the level, that
of Chapter~\ref{ch:4}. The constants of (H2)--(H3), which every clause uses, are chosen in this
order: first $L$, odd with $L \ge L_0(N)$; then Ba{\l}aban's auxiliary parameters, in one
fixed order, each of which is determined by earlier ones, or bounded on one side by a positive
function of earlier ones, or confined to a window that earlier conditions make non-empty; then
$\gamma$; and then $\gamma_H$. The further constants of (H7), used by (iii)--(vi), are chosen
after $\gamma_H$.
That this order contains no cycle is Lemma~\ref{2.4:lem-no-cycle}, which is proved. That the
one-sided conditions of (H3) can be met simultaneously is
Proposition~\ref{2.4:prop-order-consistent}, which is proved as an implication from the
independence from $M$ of the random-walk constants, Proposition~\ref{2.4:prop-m-freeness}; its
proof also uses Lemma~\ref{2.4:lem-no-cycle}. Proposition~\ref{2.4:prop-m-freeness} is in turn
proved as an implication from the operator theorem at free current and the statements named in
its section. These constants are those of Ba{\l}aban's bounds for
old terms beside a new region, whose standing is ``printed without proof; cited as printed''
\cite[p.~375]{B-CMP122b}.

\paragraph{The standing of the clauses.}
Every statement of this book carries one of the phrases ``proved'', ``proved as an implication
from'' the statements named, ``stated; proof incomplete'', ``not proved'' and ``printed without
proof; cited as printed''. No clause of Theorem~\ref{3.1:thm-main} is proved outright.
\begin{description}
\item[(0)] Proved as an implication, jointly with (i) in one induction on the level
(Chapter~\ref{ch:4}). It rests on its own four steps, each proved as an implication. The proofs
of these use statements printed without proof and cited as printed, statements of status
``stated; proof incomplete'', and statements not proved.
\item[(i)] Proved as an implication, jointly with (0) (Chapter~\ref{ch:4}). It rests on the same
statements and, in addition, on:
the positivity dichotomy of the large-field operations, which is of status ``stated; proof
incomplete'';
Ba{\l}aban's transfer of the analysis of one renormalization step in his small-field
approximation to a cover by cubes with a subordinate partition of unity, which is printed without
proof in \cite[p.~276]{B-CMP119} and cited as printed;
further steps of status ``stated; proof incomplete'', among them the steps carrying the case
$N\ge3$; and statements that are not proved, on which several of the steps of (i) rest.
\item[(iii)] Proved as an implication from (0), (i) and the correlation theorem
(Chapter~\ref{ch:5}). The correlation theorem is itself proved as an implication, by a complete
deduction, from Ba{\l}aban's statements as cited, namely theorems cited by statement and
sentences printed without proof and cited as printed, among them his bounds for old terms beside
a new region \cite[p.~375]{B-CMP122b}, with the reading chosen stated at each place where his
construction admits more than one, and with two sentences (the sentence of \cite[p.~269]{B-CMP119} on the
other characteristic functions on the strip, and the remark of \cite[p.~191]{B-CMP122a}) used
here as stated in Chapter~\ref{ch:5}, and from statements of status
``stated; proof incomplete'', among them the weight lemma and the steps of the ultraviolet
induction on which (i) rests, whose other inputs have the standings given under (i); and
the case $N\ge3$ rests on a step of status ``stated; proof incomplete''. Item (d) of (iii), the
cumulative one-loop law of the running coupling, is proved as an implication from the statement of
the correlation theorem on the flow and a lemma on the one-loop coefficient, under the
identification of the increment of the coupling in the correlation theorem with that of clause
(0), both being Ba{\l}aban's coefficient of the complete regular part, with no further
unproved input. The per-ball volume-free
bound (Chapter~\ref{ch:6}) is proved as an implication from the correlation theorem, one
inequality of Ba{\l}aban cited by statement \cite[(1.26)]{B-CMP116}, reflection positivity of
the lattice measure and the uniformity of the last-scale constants in the torus
exponent. The last of these is of status ``stated; proof incomplete''.
\item[(ii)] On each torus, along tuned bare couplings and without subsequences, the limits, their
invariance and their reflection positivity are proved as an implication from the statements on
which (v) rests. Along subsequences of first-passage cutoffs they are proved as an implication
from (iii) and reflection positivity of the lattice measure, which is proved (Chapter~\ref{ch:6}).
On $\mathbb{R}^4$ the limits along one joint subsequence, their invariance under translations and
the hypercubic group, their reflection positivity and the reconstruction from the span of centred
products of curvature fields with pairwise separated supports are proved as an implication from
the per-ball volume-free bound, reflection positivity of the lattice measure and (0)
(Chapter~\ref{ch:7}).
\item[(iv)] Clause (iv) is stated on every torus of the family, uniformly in the cutoff, the
volume and the position, together with a bound on the difference of the two-point functions on two
tori of the family; its consequences hold in every limit of (ii), on the tori and on $\mathbb{R}^4$.
It is proved as an implication (Chapter~\ref{ch:8}): a complete deduction at the standing of the
correlation bounds (iii), and, through the normaliser statement of the passage to all volumes, of
the torus-independence lemma of clause (0), itself proved as an implication from earlier lemmas
of this book and Ba{\l}aban's definitions, under one further one-sided lower bound on $M$, the
convergence and analytic extension of the expansions on the joined localisation domains being
cited as printed from \cite[p.~192, (1.70)]{B-CMP122a} and
\cite[p.~377, (1.68)--(1.69)]{B-CMP122b}.

Here a large-field region created at step $j$ has \emph{age} $k-j$ at level $k$; it is
\emph{young} or \emph{old} as this age lies below or above the threshold displayed in
Chapter~\ref{ch:8}, and \emph{healed} once a later step has absorbed its large-field factor back
into the small-field format and the region has been re-filed into the regular expansion.

On every torus of the family the two-point law of the clause is deduced from: the finite-volume
closure of the large-field age series, the Gaussian extraction, the localisation of the sourced
step, the healed old-region bound, the age lemma, the rarity of young large-field blocks, and the
last-lattice stability bounds of the correlation theorem. It holds there at depths above a floor
that depends on the coupling alone on the reference torus of hypothesis (H5) and on every torus
below a volume threshold depending on the coupling (the weak-coupling horizon of
Chapter~\ref{ch:8}), and on the coupling and the volume beyond that threshold. Large-field
regions still carrying their
large-field factor at the level read enter according to their age: the young ones through the
rarity of young large-field blocks, the old ones healed by then through the healed old-region
bound, and the old ones not healed by then through the finite-volume closure, which bounds their
contribution by the un-normalised age weight at the last lattice of the run, under the depth floor
just described; the weight of the joint
presence of two such regions is bounded by that of the presence of one of them, so that the
weight of a single region enters, and enters linearly; for the old regions not healed by then no
statement on their rarity in the normalised measure is used, in particular not the estimate,
uniform in the volume, named below. Of these and the
statements beneath
them, one part of the small-piece statement beneath the localisation of the sourced step is of
status ``stated; proof incomplete''; the second-order part of the Gaussian extraction is proved as
an implication from the correlation theorem, a cluster-expansion lemma, which is proved, and
Ba{\l}aban's regularity theorem for the variational problem \cite[Theorem~1]{B-CMP102b}, under
one further one-sided condition tying two of his free integer exponents, $q\ge 21p_0+1$, with $q$
the exponent of \cite[(2.28)]{B-CMP119} and $p_0$ the fixed exponent of the degree function
$p_0(\cdot)$; the
last-lattice bounds of the correlation theorem have the
standing of (iii) and (i); the others are proved or proved as implications. The proof uses in
addition
the source lemma and the flat shell bound, both stated in Chapter~\ref{ch:8} with their proofs
outstanding, and the continuum
kernel, whose statement and proof are not written out there. The passage to all volumes, which carries the two-point law with
depth floors depending on $g$ alone to every torus of the family, uses the volume-transfer
theorem, proved as an implication, a transfer lemma and an averaging lemma, both proved, and the
normaliser statement, which rests on the torus-independence lemma, whose standing is given above,
and on a lemma on the normalising constant, which is proved.
Of the large volume the passage
uses only reflection positivity and the lower stability bound. Not asserted: convergence or
uniqueness in the volume,
uniqueness of the vacuum, clustering, a mass gap, and invariance under $\mathrm{O}(4)$ outside
clause (vi); nor the rarity, uniformly in the volume, of old large-field regions not healed, which
is neither used nor proved; nor the
transfer for more than two sources, or
the sharp power of the scale in the transfer bound.
\item[(v)] Proved as an implication (Chapter~\ref{ch:9}) from the following:
\begin{itemize}
\item clause (iii) along the runs compared; the re-filing lemmas, each proved as an implication
from statements on which (iii) rests; the bi-Lipschitz lemma away from the endpoints of its
range, proved as an implication from its displayed inputs, among them a re-centred reading of the
running-shift statement (its clause at the endpoints is not proved); the one-lattice
Lipschitz step, itself proved as an implication; and the limit-member identification, proved as
an implication from the change-of-units identity, clause (0) and the other parts of the clause,
with no further input. The carry
lemma, on which the proof of the one-lattice Lipschitz step rests, is of status ``stated; proof
incomplete''.
\item the counting of the large-field regions of different sizes met in the comparison of the
two runs, which is proved, under two further one-sided lower bounds on $M$; and the comparison of
the ambient couplings of the two runs, in the form stated in Chapter~\ref{ch:9}, proved as an
implication.
\item the one-run expansions on labels, one step of the argument for the exponential rate, and
two statements of the induction of (iii) rewritten in the re-filed coordinates, each of status
``stated; proof incomplete''; and a third such statement, proved as an implication in the
variant form stated in Chapter~\ref{ch:9}, one of whose auxiliary clauses, supplying an input the
variant needs, is not proved.
\item a bound, specific to the comparison of two cutoffs, on the factors kept beside large-field
regions, one further statement of that induction in the re-filed coordinates, a statement on the
dependences of the clause's constants, and the extension to $N\ge3$, each not proved; beneath the
first of these lie a lemma on the fields enclosed by those regions, of status ``stated; proof
incomplete'', and two statements comparing two runs whose proofs are outstanding: the difference
between two runs
of the remainder terms whose details are omitted in \cite{B-CMP122a}, and the corresponding
difference for the nested interrupted Gaussian integrations; the two-cutoff response statement at
zeroth order is proved as an
implication (Chapter~\ref{ch:5}), and the bookkeeping by which clause (v) reads it at the finer
zones of its own two-cutoff comparison is stated in Chapter~\ref{ch:9} and not proved.
\end{itemize}
Its sub-clause, stated for the constructions at two coprime blocking factors, on the convergence
of the increments of the inverse coupling, the exact covariance under a change of the reference
length by a power of $L$ and the injectivity of the renormalized coupling
is proved as an implication at the standing of (iii), (iv) and (v), under one further one-sided
lower bound on the radius $\varrho_2$; the corollary on uniqueness at short distances stated after
the clause is proved as an implication at the same standing.
\item[(vi)] Not proved (Chapter~\ref{ch:10}). The clause, invariance under $\mathrm{O}(4)$ of
every limit of (ii) on $\mathbb{R}^4$, is a complete deduction from the far-sphere flux hypothesis
$(\mathrm{IR}\text{-}\mathrm{fl})_{\omega_0}$: for one nonzero antisymmetric $\omega_0$, every
$n\ge2$ and all test functions with pairwise disjoint supports, the upper limit, as the cutoff is
removed, of the modulus of the connected correlation of the flux, through the sphere of radius
$\mathcal{R}$, of the lattice rotation current generated by $\omega_0$ with the $n$ smeared
curvature fields tends to zero along one sequence of radii $\mathcal{R}\to\infty$, the cutoff being
removed first and the radius second. The decay condition displayed beside this hypothesis in
Chapter~\ref{ch:10} (one far stress-tensor insertion
decorrelates from the observables faster than the inverse fourth power of the distance, uniformly
in the cutoff) is sufficient for this hypothesis. The clause is deduced from the hypothesis and
from:
\begin{itemize}
\item the part of (ii) on $\mathbb{R}^4$, itself proved as an implication.
\item the rotational Ward identity and the dimension-six classification, which are proved, the
decomposition of the identity's $\mathrm{O}(4)$-breaking part being of status ``stated; proof
incomplete''.
\item statements each proved as an implication: (0), (i), (iii), the vector package, the
transport of stress-tensor insertions through the renormalization steps, and the one-loop
anisotropy inequality with exponent one.
\end{itemize}
Beneath these implications lie the following statements, none of status ``proved''. The reduction
theorem, in the sharp form
of its coefficients used here, is of status ``stated; proof incomplete''. Among the inputs of the
vector package, its bounds on the rim terms of the vector channel, those carrying the position
weight, rest on a statement about healed regions of status
``stated; proof incomplete'', and its bound on the bulk term of that channel, in the form in which
the clause reads it, is proved as an implication for a fixed gap between the scales compared and is
stated in
Chapter~\ref{ch:10}, with its proof outstanding, in the form weighted over all gaps. The proofs of
these implications also use statements interpreting the form in which Ba{\l}aban writes his
renormalization step for the terms of his class, among them the renormalization step and the
re-filing written in that form. One of them, the reading of that form at the collars of the vector
package, is proved as an
implication from an extension of the statement of the correlation theorem, declared in
Chapter~\ref{ch:10} and not proved, together with its other named inputs and under one further
one-sided ceiling on the coupling; the renormalization step and the re-filing written in that form
are proved as implications in part, are of status ``stated; proof incomplete'' in part, and in one
further part are not proved.
\end{description}
Every clause after (i) inherits the standing of (0), (i) and (iii), so none can stand better than
the correlation theorem and the ultraviolet steps on which that theorem rests. As a statement
about the running coupling, not about the Schwinger functions, and at the standing of (0) and
(i): the running
coupling exists at every scale and is strictly positive; its inverse square obeys two-sided linear
flow bounds uniformly in the cutoff; and at every scale the effective action is the Yang--Mills
action at that coupling plus terms irrelevant in Ba{\l}aban's norms.

\paragraph{The count.}
Theorem~\ref{3.1:thm-main} rests on twenty-seven intermediate statements of Ba{\l}aban's
programme that this book does not prove outright. They fall into the following three groups.
\begin{itemize}
\item Seven are counted as implications from named statements. Six are proved as such in
Part~II: the
independence of the localised terms from the small-field region; the operator bounds at free
current; the representation and bound of the large-field terms; the large-field part of the
coupling increment $\beta_{k+1}$; the one-power hereditary class; and the operator theorem at
free current. The seventh, the positivity dichotomy of the large-field operations, whose only
displayed inputs are the representation of the density established in clause~(i) and
clause~(0), is printed in Part~II with the standing ``stated; proof incomplete''.
\item Four concern assertions of Ba{\l}aban's papers. Three of them are sentences printed there
without proof, with the standing ``printed without proof; cited as printed'': the transfer of
the analysis of \cite[\S3]{B-CMP109} to the cover used in \cite{B-CMP119}, stated in
\cite[p.~276]{B-CMP119}; Ba{\l}aban's bounds for old terms beside a new large-field region,
\cite[p.~375]{B-CMP122b}, used together with the $M$-freeness of the random-walk constants
(Proposition~\ref{2.4:prop-m-freeness}), which is proved as an implication from the operator
theorem at free current and the further inputs displayed with it; and
the admissibility of the prepared large-field data,
\cite[pp.~178--179]{B-CMP122a}. The fourth, that the initial data lie in the
regularity classes of \cite[(3.35), (3.37)]{B-CMP99b} at the two lowest tower levels on the
cells on which the operator theorem at free current uses them, is proved in Part~II on every
such cell as an implication from the inputs displayed with it there, under one further
one-sided lower bound on an enlarged block size, chosen after $L$ and $M_1$ in the order of
choice.
\item Sixteen are counted under the standing ``stated; proof incomplete'', each being printed
in Part~II with its own standing: the ten statements of the
representation and the induction of clause~(i); the step for complex directions, an input of
clause~(ii); the correlation theorem, whose deduction in Part~II is an implication from
Ba{\l}aban's statements as cited, each used in the reading stated at its point of use, in
particular at \cite[p.~269]{B-CMP119} and \cite[p.~191]{B-CMP122a}, and from those ten
statements, and which is therefore counted with them; the
weight lemma; and, needed only for $N \ge 3$, three statements: the
$\mathrm{SU}(N)$ form of the constants of clauses~(i) and~(ii), the
passage of the arguments of the induction to $\mathrm{SU}(N)$, and the extension of the
correlation theorem to $\mathrm{SU}(N)$.
\end{itemize}
None of the twenty-seven is of the kind ``not proved''. At $N = 2$ the three statements needed
only for $N \ge 3$ are not used, and the counted statements number
twenty-four: the seven and the four of the first two items, thirteen stated with proof
incomplete, and none not proved.

By clause, the twenty-seven are used as follows: four are clause~(0)'s own steps, the first
four of the first group; two enter every clause, the one-power hereditary
class and the statement on the initial data at the two lowest tower levels; fourteen enter
clause~(i): the transfer to the cover, the positivity dichotomy, the ten statements of
clause~(i) and the two $\mathrm{SU}(N)$ statements that concern it; one, the step for complex
directions, is an input of clause~(ii) and, through the Gaussian extraction, of
clause~(iv); four enter clause~(iii): the correlation theorem, the weight lemma, Ba{\l}aban's
bounds beside a new region and the extension to $\mathrm{SU}(N)$; and two are inputs of both
clauses~(0) and~(i): the operator theorem at free current and the admissibility of the
prepared large-field data.
Item~(d) of clause~(iii), the exact one-loop law of the construction's running coupling, adds
nothing to this count: it is proved as an implication from the flow statement of the
correlation theorem (its bounds on the increments of the inverse coupling) and the one-loop
lemma on the constants of Definition~\ref{1.2:def-flow-constants}, under the identification of
the correlation theorem's coupling increment with the flow's increment $\beta_{k+1}$, both
being Ba{\l}aban's coefficient of the complete regular part, with the same standing as
clause~(iii).

The block-law statements of Remark~\ref{3.2:rem-block-laws} rest on twenty-nine further
intermediate statements, disjoint from the twenty-seven, which no clause of
Theorem~\ref{3.1:thm-main} uses and which are not counted here.

Clauses (ii)--(vi) of Theorem~\ref{3.1:thm-main} also depend on further inputs that are
statements about the construction's measures and their curvature correlation functions. Among
them are the second-derivative room
lemma of the Gaussian extraction, the source lemma and the continuum kernel.
Neither these inputs nor the further sentences of Ba{\l}aban cited as printed at their points
of use are among the twenty-seven; each is printed in place in Part~II with its standing, and
those whose proof is outstanding or incomplete, or which are printed without proof, are listed
in the index of Chapter~\ref{ch:13}, with which this count agrees.

\paragraph{The case $N\ge 3$.}
Hypothesis (H1) fixes the range of $N$ for each clause. Theorem~0 states clauses (0) and (i) for
every $N\ge 2$, and states and argues clauses (ii)--(vi) at $N=2$ only. Clauses (0) and (i) are
proved as implications from the statements on which they rest.

For $N\ge 3$ those statements include two that concern the extension of the argument from
$N=2$ to $\mathrm{SU}(N)$:
\begin{itemize}
\item the form, for $\mathrm{SU}(N)$, of the stability constants $E_\pm$ of clause (i);
\item the statement that every argument beneath clauses (0) and (i) given at $N=2$ carries over to
$\mathrm{SU}(N)$.
\end{itemize}
Each of the two is stated; proof incomplete.

For clause (iii), the extension to $N\ge 3$ is the extension of the correlation theorem to
$\mathrm{SU}(N)$. It is stated; proof incomplete. Clause (iv) at $N\ge 3$ would rest on it through
clause (iii). Clause (v) at $N\ge 3$ would need, in addition, a uniqueness theorem for
$\mathrm{SU}(N)$, which is not proved. Clause (vi) is argued at $N=2$ only.

The rank $N$ enters in four places:
\begin{itemize}
\item the one-loop constant $c_0$, through the one-sided lower bound $\lambda_\sharp\ge c_0/4$ on the
slope of the coupling flow;
\item the single existential constant $L_0=L_0(N)$;
\item the constants of the Lie algebra $\mathfrak{su}(N)$ and the constants built from them, among
them the stability constants $E_\pm$ of clause (i), and hence the thresholds on the coupling;
\item the factor $(N^2-1)/2$ of clause (iv).
\end{itemize}

Accordingly, nothing is asserted at $N\ge 3$ beyond clauses (0) and (i). At $N\ge 3$ these two
clauses are proved as implications whose antecedents include the two statements about
$\mathrm{SU}(N)$ listed above, each of which is stated; proof incomplete.

%% ---- end inlined file: chapters/c11 ----

\part{Statements awaiting proof}
\clearpage
\setcounter{section}{11}
\refstepcounter{section}
\section*{Chapter~\thesection\quad Statements used without proof in this account\addcontentsline{toc}{section}{\protect\numberline{\thesection}Statements used without proof in this account}\label{ch:13}}
%% ---- begin inlined file: chapters/c13 ----
Every statement of Part~II whose proof in this book is outstanding or incomplete is listed below
once, and so is every sentence of Ba{\l}aban's papers that is used here as printed there without
proof. A statement that carries a number is listed by its number; otherwise it is listed by its
content. Each entry gives the printed name, the chapter and the form. The entries fall under five
heads: (1) statements for which no complete argument is known; (2) statements used as stated whose
proofs are not written out here; (3) sentences of Ba{\l}aban's papers
printed there without proof and used here as printed; (4) statements whose proof is incomplete;
(5) model theorems. Any count of the statements of Part~II that lack a complete proof is a count of
this list. A statement proved as an implication from named statements appears here only through
those of its antecedents that are listed. Where a route is given, it says in mathematical terms from
what the statement is to be deduced.

Large-field regions, their age and their healing are as defined in Chapter~\ref{ch:8}. A region
created at step $j$ has age $k-j$ at level $k$; it is live at level $k$ if it still carries its
large-field factor at that level; it is healed when a later step absorbs its large-field factor
back into the small-field format. The never-healed old regions are those created more than the
prescribed number of steps before the witness depth and not healed by then; the witness depth is
the depth below the cutoff at which the two-point witness theorem is applied, and a depth floor is
a lower bound on that depth. Islands, too, are as defined in Chapter~\ref{ch:8}.

\paragraph{(1) Proof outstanding: statements for which no complete argument is known.}
\begin{enumerate}
\item Proposition~\ref{8.12:prop-uniform-old-rarity} $\cdot$ uniform rarity of never-healed old
large-field regions $\cdot$ Chapter~\ref{ch:8} $\cdot$ proof outstanding
\item the exit bound $\cdot$ Chapter~\ref{ch:8}
$\cdot$ proof outstanding
\item the half-skin bound at the canonical
comparison point $\cdot$ Chapter~\ref{ch:8} $\cdot$ proof outstanding
\item the unbounded-union depth statement $\cdot$ Chapter~\ref{ch:8} $\cdot$ proof outstanding
$\cdot$ route: one model rate lemma and its transfer to Ba{\l}aban's step
\item the collar transfer $\cdot$ Chapter~\ref{ch:8} $\cdot$ proof outstanding
\item Proposition~\ref{5.13:prop-sun-correlation} $\cdot$ the correlation theorem for
$\mathrm{SU}(N)$, $N\ge 3$ $\cdot$ Chapter~\ref{ch:5} $\cdot$ proof outstanding
\item Proposition~\ref{9.14:prop-sun-uniqueness} $\cdot$ the uniqueness theorem for
$\mathrm{SU}(N)$, $N\ge 3$ $\cdot$ Chapter~\ref{ch:9} $\cdot$ proof outstanding
\end{enumerate}

\paragraph{(2) Proof outstanding: statements used as stated whose proofs are not written out
here.}
\begin{enumerate}
\item the re-blocking step outside admissible boxes $\cdot$ Chapter~\ref{ch:4} $\cdot$ proof
outstanding
\item uniqueness and locality of the constrained minimiser $\cdot$ Chapter~\ref{ch:4} $\cdot$
proof outstanding
\item the cube gauge condition for the small pairs and the real-background pairs $\cdot$
Chapter~\ref{ch:4} $\cdot$ proof outstanding
\item positivity dichotomy of the large-field operations $\cdot$ Chapter~\ref{ch:4} $\cdot$ proof
outstanding
\item the volume counts of the correlation theorem, as used in the relative small-field machine
$\cdot$ Chapter~\ref{ch:4} $\cdot$ proof outstanding
\item two bounds on the data entering lemmas of Chapter~\ref{ch:5}, displayed as hypotheses of
the relative small-field machine $\cdot$ Chapter~\ref{ch:4} $\cdot$ proof outstanding
\item the kernel-difference bound displayed as a hypothesis of the relative small-field machine
$\cdot$ Chapter~\ref{ch:4} $\cdot$ proof outstanding
\item geometry and counting of live large-field regions $\cdot$ Chapter~\ref{ch:4} $\cdot$ proof
outstanding
\item a kept term as a functional of the data: its label, norm and law sites $\cdot$
Chapter~\ref{ch:4} $\cdot$ proof outstanding
\item kept terms add one summand to the probe functional $\cdot$ Chapter~\ref{ch:4} $\cdot$ proof
outstanding
\item kept terms at a large-field operation on a foreign component $\cdot$ Chapter~\ref{ch:4}
$\cdot$ proof outstanding
\item Combes--Thomas decay of the block covariance $\cdot$ Chapter~\ref{ch:4} $\cdot$ proof
outstanding
\item conversion of per-rate bounds to the stage norm $\cdot$ Chapter~\ref{ch:4} $\cdot$ proof
outstanding
\item holonomy bound for walks across several cubes $\cdot$ Chapter~\ref{ch:4} $\cdot$ proof
outstanding
\item Proposition~\ref{5.6:prop-polydisc-holomorphy} $\cdot$ holomorphy of the localised terms on
the decoupling polydisc $\cdot$ Chapter~\ref{ch:5} $\cdot$ proof outstanding
\item Lemma~\ref{5.11:lem-sixth-family} $\cdot$ the sixth-family lemma $\cdot$ Chapter~\ref{ch:5}
$\cdot$ proof outstanding
\item Lemma~\ref{7.3:lem-first-passage-offsets} $\cdot$ non-convergence of the first-passage
offsets $\cdot$ Chapter~\ref{ch:7} $\cdot$ proof outstanding
\item identity of representations for the shifted frozen run $\cdot$ Chapter~\ref{ch:8} $\cdot$
proof outstanding
\item Lemma~\ref{8.21:lem-source} $\cdot$ the source lemma $\cdot$ Chapter~\ref{ch:8} $\cdot$ proof
outstanding
\item Lemma~\ref{8.21:lem-continuum-kernel} $\cdot$ the continuum kernel $\cdot$ Chapter~\ref{ch:8}
$\cdot$ proof outstanding $\cdot$ route: a classical computation for the free field
\item Definition~\ref{9.1:not-radius-exponent-floors} $\cdot$ the assertions on the radius
exponents, the floors on $\kappa_1$ and the fat core $\cdot$ Chapter~\ref{ch:9} $\cdot$ proof
outstanding
\item Definition~\ref{9.1:not-shell-floor-chain} $\cdot$ the assertions on the floor chain on
shells and the ratio $A_0/A_1$ $\cdot$ Chapter~\ref{ch:9} $\cdot$ proof outstanding
\item Definition~\ref{9.1:not-levelwise-rate} $\cdot$ the assertions on the level-wise rate, the
number of stages and the strip width $\cdot$ Chapter~\ref{ch:9} $\cdot$ proof outstanding
\item Proposition~\ref{9.1:prop-blind-site-expansions} $\cdot$ site expansions with island-blind
inputs $\cdot$ Chapter~\ref{ch:9} $\cdot$ proof outstanding
\item Proposition~\ref{9.2:prop-modulus-majorant} $\cdot$ the history-wise modulus majorant $\cdot$
Chapter~\ref{ch:9} $\cdot$ proof outstanding
\item the nesting clause of the majorant hypothesis used in the proof of
Theorem~\ref{9.2:thm-geometric-forgetting} $\cdot$ Chapter~\ref{ch:9} $\cdot$ proof outstanding
\item Proposition~\ref{9.3:prop-propagator-stationarity} $\cdot$ weak stationarity of the
fine-lattice propagators $\cdot$ Chapter~\ref{ch:9} $\cdot$ proof outstanding
\item the two-cutoff response statement at zeroth order carried over to the shells used by
clauses~(v) and~(ii-$\mathbb{T}$-a), which are not those on which it is proved $\cdot$
Chapter~\ref{ch:9} $\cdot$ proof outstanding
\item the residual clauses of the increment statement $\cdot$ Chapter~\ref{ch:9} $\cdot$ proof
outstanding
\item the bound on the difference, between the two runs, of the remainder terms of case~(ii) of
\cite{B-CMP122a}, whose bounds that paper states without details $\cdot$ Chapter~\ref{ch:9} $\cdot$
proof outstanding
\item the bound on the difference, between the two runs, of the nested Gaussian integrals
interrupted by large-field regions $\cdot$ Chapter~\ref{ch:9} $\cdot$ proof outstanding
\item the bound at the current depth for the two factors kept un-integrated inside a large-field
region $\cdot$ Chapter~\ref{ch:9} $\cdot$ proof outstanding $\cdot$ depends on the two preceding
items
\item Lemma~\ref{9.6:lem-maximum-principle} $\cdot$ a maximum principle for coupled scale-indexed
inequalities $\cdot$ Chapter~\ref{ch:9} $\cdot$ proof outstanding
\item Definition~\ref{9.6:def-comparison-chain} $\cdot$ the assertions made in the definition of
the comparison chain $\cdot$ Chapter~\ref{ch:9} $\cdot$ proof outstanding
\item Lemma~\ref{9.6:lem-chain-structure} $\cdot$ well-definedness and structure of the comparison
chain $\cdot$ Chapter~\ref{ch:9} $\cdot$ proof outstanding
\item Corollary~\ref{9.8:cor-subsequential-limits} $\cdot$ first-passage offsets and the
identification of subsequential limits $\cdot$ Chapter~\ref{ch:9} $\cdot$ proof outstanding
\item Proposition~\ref{9.9:prop-remainder-derivative} $\cdot$ bare-offset derivative of the witness
remainder $\cdot$ Chapter~\ref{ch:9} $\cdot$ proof outstanding
\item Gaussian skeleton and in-box precision $\cdot$ Chapter~\ref{ch:9} $\cdot$ proof outstanding
\item the regrouped two-run identity at the real centres for the re-filed boundary terms $\cdot$
Chapter~\ref{ch:9} $\cdot$ proof outstanding
\item the bound on the own-history part of the boundary terms produced at a large-field step
$\cdot$ Chapter~\ref{ch:9} $\cdot$ proof outstanding $\cdot$ route: from the corresponding clause
for the history part
\item the endpoint case of Corollary~\ref{9.12:cor-same-limit} $\cdot$ Chapter~\ref{ch:9} $\cdot$
proof outstanding
\item Lemma~\ref{9.13:lem-limit-member} $\cdot$ the limit-member identification $\cdot$
Chapter~\ref{ch:9} $\cdot$ proof outstanding
\item Definition~\ref{10.3:def-vector-production-bounds} $\cdot$ the production bounds for the
vector part under the split rule $\cdot$ Chapter~\ref{ch:10} $\cdot$ proof outstanding
\item Definition~\ref{10.5:def-vector-package-output} $\cdot$ the bounds for the vector part of
the run in output form $\cdot$ Chapter~\ref{ch:10} $\cdot$ proof outstanding
\end{enumerate}

\paragraph{(3) Printed without proof; cited as printed.}
\begin{enumerate}
\item the transfer of the analysis of \cite[\S3, Lemmas~5 and~6]{B-CMP109} to the cover of
\cite[\S3]{B-CMP119} $\cdot$ Chapter~\ref{ch:4} $\cdot$ printed without proof in
\cite[p.~276]{B-CMP119}; cited as printed
\item Remark~\ref{5.10:rem-old-terms-beside} $\cdot$ Ba{\l}aban's bounds for old terms beside a new
region $\cdot$ Chapter~\ref{ch:5} $\cdot$ printed without proof in \cite[p.~375]{B-CMP122b}; cited as
printed
\item admissibility of the prepared large-field data $\cdot$ Chapter~\ref{ch:4} $\cdot$ printed
without proof in \cite[pp.~178--179]{B-CMP122a}; cited as printed
\end{enumerate}

\paragraph{(3, continued) Further sentences cited as printed at their use sites, not counted among
those above.}
\begin{enumerate}
\item the regularity conditions at the two lowest tower levels,
\cite[p.~396, (3.35), (3.37)]{B-CMP99b} $\cdot$ Chapter~\ref{ch:4} $\cdot$ cited as printed at its
use site
\item the extension of the bounds \cite[(3.26)--(3.28)]{B-CMP109} on block-averaged curvature to
near data, \cite[p.~277]{B-CMP109} $\cdot$ Chapter~\ref{ch:4} $\cdot$ cited as printed at its use
site
\item locality of the interaction terms, \cite[p.~379]{B-CMP122b} $\cdot$ Chapter~\ref{ch:4} $\cdot$
cited as printed at its use site
\item the inclusion sentence \cite[(1.17)]{B-CMP109}, read at complex pairs $\cdot$
Chapter~\ref{ch:4} $\cdot$ cited as printed at its use site
\item the bounds and walk expansion of the inverse of the constrained covariance,
\cite[\S3, the sentence preceding (3.132)]{B-CMP99b} $\cdot$ Chapter~\ref{ch:4} $\cdot$ cited as
printed at its use site
\item the walk expansions of the further operators, \cite[\S3, after Theorem~3.9]{B-CMP99b}
$\cdot$ Chapter~\ref{ch:4} $\cdot$ cited as printed at its use site
\item analyticity of the derivative bounds for the averaging maps, \cite[\S E, p.~43]{B-CMP98}
$\cdot$ Chapter~\ref{ch:4} $\cdot$ cited as printed at its use site
\item the second inequality of the bound on the current of the far family,
\cite[(3.49), (3.69)]{B-CMP99b} $\cdot$ Chapter~\ref{ch:4} $\cdot$ cited as printed at its use site
\item the polymer-gas representation of the bracket, \cite[(1.90)]{B-CMP122b} $\cdot$
Chapter~\ref{ch:4} $\cdot$ cited as printed at its use site
\item the bounds of the large-field terms under the remaining modifications,
\cite[p.~386]{B-CMP122b} $\cdot$ Chapter~\ref{ch:4} $\cdot$ cited as printed at its use site
\item the layerwise gauge clauses of the large-field spaces,
\cite[(1.64)--(1.69), pp.~190--191]{B-CMP122a} $\cdot$ Chapter~\ref{ch:4} $\cdot$ cited as printed
at its use site
\item the layerwise ceiling as a condition of small coupling, \cite[p.~192]{B-CMP122a} $\cdot$
Chapter~\ref{ch:4} $\cdot$ cited as printed at its use site
\item the absolute bound and gauge invariance of the boundary terms,
\cite[(1.11), (1.23)--(1.48), p.~369]{B-CMP122b} $\cdot$ Chapter~\ref{ch:4} $\cdot$ cited as printed
at its use site
\item the analytic solvability of the equation \cite[p.~367, (1.40)--(1.41)]{B-CMP122b} $\cdot$
Chapter~\ref{ch:4} $\cdot$ cited as printed at its use site
\item smallness of the sum of the localised remainders of the localisation and exponentiation
\cite[(1.11), (1.23)--(1.48)]{B-CMP122b}, stated in \cite[p.~368, after (1.46)--(1.48)]{B-CMP122b}
$\cdot$ Chapter~\ref{ch:5} $\cdot$ cited as printed at its use site
\item the identification of the flat quadratic form of the small-field step with that of the free
field, \cite[\S1]{B-CMP95} $\cdot$ Chapter~\ref{ch:8} $\cdot$ cited as printed at its use site
\item positivity and bound of the normalised island operation, \cite[\S1]{B-CMP122b} $\cdot$
Chapter~\ref{ch:8} $\cdot$ cited as printed at its use site
\item convergence and analytic extension, with their bounds, of the expansions on the joined
localisation domains, \cite[p.~192, (1.70)]{B-CMP122a} and
\cite[p.~377, (1.68)--(1.69)]{B-CMP122b} $\cdot$ Chapters~\ref{ch:4}, \ref{ch:8} and~\ref{ch:9}
$\cdot$ cited as printed at its use sites
\item the passage between \cite[(189)]{B-CMP102b} and \cite[(190), p.~308]{B-CMP102b} $\cdot$
Chapter~\ref{ch:9} $\cdot$ cited as printed at its use site
\end{enumerate}

\paragraph{(4) Stated; proof incomplete.}
\begin{enumerate}
\item Lemma~\ref{4.2:lem-unmarked} $\cdot$ Chapter~\ref{ch:4} $\cdot$ stated; proof incomplete
\item the re-blocking step inside admissible boxes $\cdot$ Chapter~\ref{ch:4} $\cdot$ stated; proof
incomplete
\item real-factor slot bounds inside admissible boxes $\cdot$ Chapter~\ref{ch:4} $\cdot$ stated;
proof incomplete
\item contribution of the kept terms to the stage totals $\cdot$ Chapter~\ref{ch:4} $\cdot$ stated;
proof incomplete
\item field dependence of the densities in cube gauges $\cdot$ Chapter~\ref{ch:4} $\cdot$ stated;
proof incomplete
\item oscillation of a Gaussian piece under small moves $\cdot$ Chapter~\ref{ch:4} $\cdot$ stated;
proof incomplete
\item dissolution of the arrested region $\cdot$ Chapter~\ref{ch:4} $\cdot$ stated; proof
incomplete
\item the $\mathbf{T}$-step layer in a background field $\cdot$ Chapter~\ref{ch:4} $\cdot$ stated;
proof incomplete
\item weights, the dated size parameter and heredity $\cdot$ Chapter~\ref{ch:4} $\cdot$ stated;
proof incomplete
\item cube families under a stage move $\cdot$ Chapter~\ref{ch:4} $\cdot$ stated; proof incomplete
\item current thresholds across an increase of density $\cdot$ Chapter~\ref{ch:4} $\cdot$ stated;
proof incomplete
\item the one-point current probe $\cdot$ Chapter~\ref{ch:4} $\cdot$ stated; proof incomplete
\item no index-zero margin on source domains $\cdot$ Chapter~\ref{ch:4} $\cdot$ stated; proof
incomplete
\item inclusion of the ambient domain at size $\gamma_0^{(k)}$, the dated radius of the arrest
theorem $\cdot$ Chapter~\ref{ch:4} $\cdot$ stated; proof incomplete
\item the arrest theorem $\cdot$ Chapter~\ref{ch:4} $\cdot$ stated; proof incomplete
\item newborn terms, space inclusions and continuation of stored domains $\cdot$
Chapter~\ref{ch:4} $\cdot$ stated; proof incomplete
\item bounds of later steps on live arrested pieces $\cdot$ Chapter~\ref{ch:4} $\cdot$ stated;
proof incomplete
\item the current radius read at raw points of live pieces $\cdot$ Chapter~\ref{ch:4} $\cdot$
stated; proof incomplete
\item index-zero members under maps, at the real point and at completion $\cdot$
Chapter~\ref{ch:4} $\cdot$ stated; proof incomplete
\item the core-freezing step $\cdot$ Chapter~\ref{ch:4} $\cdot$ stated; proof incomplete
\item no lattice point on both boundaries of a stratum $\cdot$ Chapter~\ref{ch:4} $\cdot$ stated;
proof incomplete
\item the cell-counting bound on the stratified frames of the induction $\cdot$
Chapter~\ref{ch:4} $\cdot$ stated; proof incomplete
\item the stability bound on the near part of the effective action at $N\ge 3$, displayed among
the standing hypotheses of the constants for $\mathrm{SU}(N)$ $\cdot$ Chapter~\ref{ch:4} $\cdot$
stated; proof incomplete
\item Proposition~\ref{5.10:prop-m-freeness} $\cdot$ $M$-freeness of the random-walk constants $\cdot$
Chapter~\ref{ch:5} $\cdot$ stated; proof incomplete at the step marked $(\ast)$ $\cdot$
independence of one of the random-walk constants from $M$
\item all two-cutoff differences on one lattice vanish $\cdot$ Chapter~\ref{ch:5} $\cdot$ stated;
proof incomplete
\item the small-piece statement $\cdot$ Chapter~\ref{ch:8} $\cdot$ stated; proof incomplete
\item radius tails of the majorant masses $\cdot$ Chapter~\ref{ch:8} $\cdot$ stated; proof
incomplete
\item the second-derivative room lemma $\cdot$ Chapter~\ref{ch:8} $\cdot$ stated; proof incomplete
\item covariance of the remainder column, small in the running coupling $\cdot$
Chapter~\ref{ch:8} $\cdot$ stated; proof incomplete
\item monotonicity of the common prefactor under island removal $\cdot$ Chapter~\ref{ch:8} $\cdot$
stated; proof incomplete
\item later steps do not read earlier healed islands $\cdot$ Chapter~\ref{ch:8} $\cdot$ stated;
proof incomplete
\item annealing and conditioning of the removal bound $\cdot$ Chapter~\ref{ch:8} $\cdot$ stated;
proof incomplete
\item translation covariance of the label histories $\cdot$ Chapter~\ref{ch:8} $\cdot$ stated;
proof incomplete
\item independence of the steps from the target coupling $\cdot$ Chapter~\ref{ch:8} $\cdot$
stated; proof incomplete
\item the age lemma $\cdot$ Chapter~\ref{ch:8} $\cdot$ stated; proof incomplete
\item rarity of never-healed old regions above a volume-dependent depth floor $\cdot$
Chapter~\ref{ch:8} $\cdot$ stated; proof incomplete
\item Theorem~\ref{9.2:thm-carry} $\cdot$ the carry lemma $\cdot$ Chapter~\ref{ch:9} $\cdot$
stated; proof incomplete
\item the extension proposition $\cdot$ Chapter~\ref{ch:9} $\cdot$ stated; proof incomplete
\item zone-resolved variation of the first-group boundary terms $\cdot$ Chapter~\ref{ch:9} $\cdot$
stated; proof incomplete
\item the tail of the rest terms above a threshold zone $\cdot$ Chapter~\ref{ch:9} $\cdot$ stated;
proof incomplete
\item aligned parts of nested leaves and the unmatched term $\cdot$ Chapter~\ref{ch:9} $\cdot$
stated; proof incomplete
\item the first index over the full history $\cdot$ Chapter~\ref{ch:9} $\cdot$ stated; proof
incomplete
\item the per-unit response factor at the step $\cdot$ Chapter~\ref{ch:9} $\cdot$ stated; proof
incomplete
\item Lemma~\ref{9.2:lem-read-set} $\cdot$ read sets of cells meeting the domain $\cdot$
Chapter~\ref{ch:9} $\cdot$ stated; proof incomplete
\item the boundary-family entry of the Lipschitz step $\cdot$ Chapter~\ref{ch:9} $\cdot$ stated;
proof incomplete
\item linear size at least one for constructed regions $\cdot$ Chapter~\ref{ch:9} $\cdot$ stated;
proof incomplete
\item the one-run expansions on labels $\cdot$ Chapter~\ref{ch:9} $\cdot$ stated; proof incomplete
\item Lemma~\ref{9.1:lem-glued-minimisers} $\cdot$ ambient regularity of the two runs, in the
form imposing the stronger regularity conditions only where a proof reads them $\cdot$
Chapter~\ref{ch:9} $\cdot$ stated; proof incomplete
\item Theorem~\ref{9.10:thm-refiled-induction} $\cdot$ the induction theorem for the re-filed run
$\cdot$ Chapter~\ref{ch:9} $\cdot$ stated; proof incomplete
\item Theorem~\ref{9.10:thm-refiled-lipschitz} $\cdot$ re-filed one-lattice Lipschitz step $\cdot$
Chapter~\ref{ch:9} $\cdot$ stated; proof incomplete
\item kernel form of the re-filed Lipschitz step $\cdot$ Chapter~\ref{ch:9} $\cdot$ stated; proof
incomplete
\item geometric source bounds for the re-filed run $\cdot$ Chapter~\ref{ch:9} $\cdot$ stated;
proof incomplete
\item finite-range envelope of last-lattice loads $\cdot$ Chapter~\ref{ch:9} $\cdot$ stated; proof
incomplete
\item geometric last-lattice comparison for the re-filed run $\cdot$ Chapter~\ref{ch:9} $\cdot$
stated; proof incomplete
\item hatted source families with geometric envelope $\cdot$ Chapter~\ref{ch:9} $\cdot$ stated;
proof incomplete
\item the domination estimate accompanying the rotational Ward identity of
Theorem~\ref{10.1:thm-rotational-identity}, with its abelian case and its form on the torus $\cdot$
Chapter~\ref{ch:10} $\cdot$ stated; proof incomplete
\item Proposition~\ref{10.3:prop-vector-package} $\cdot$ the bulk, freezing and rim bounds for the
vector part of the run $\cdot$ Chapter~\ref{ch:10} $\cdot$ stated; proof incomplete
\item Proposition~\ref{10.3:prop-cores-extended} $\cdot$ the inner-collar cores of the vector part
under the extended correlation theorem $\cdot$ Chapter~\ref{ch:10} $\cdot$ stated; proof incomplete
\item the transfer of the scalar lemmas to a tensor slot $\cdot$ Chapter~\ref{ch:10} $\cdot$
stated; proof incomplete
\item Proposition~\ref{10.4:prop-vertex-response} $\cdot$ the response of one step to the tensor
vertex $\cdot$ Chapter~\ref{ch:10} $\cdot$ stated; proof incomplete
\item Proposition~\ref{10.4:prop-typed-step} $\cdot$ the typed step with a marginal tensor slot
$\cdot$ Chapter~\ref{ch:10} $\cdot$ stated; proof incomplete
\item Lemma~\ref{10.4:lem-typed-refiling} $\cdot$ the re-filing of typed marginal families $\cdot$
Chapter~\ref{ch:10} $\cdot$ stated; proof incomplete
\item Proposition~\ref{10.4:prf-sourced-format} $\cdot$ the sourced format of the local tangent
run $\cdot$ Chapter~\ref{ch:10} $\cdot$ stated; proof incomplete
\item the statement on the cut-free one-loop numbers $\cdot$ Chapter~\ref{ch:10} $\cdot$ stated;
proof incomplete
\item Theorem~\ref{10.5:thm-reduction} $\cdot$ the reduction theorem $\cdot$ Chapter~\ref{ch:10}
$\cdot$ stated; proof incomplete
\item Corollary~\ref{10.6:cor-normalisation-decay} $\cdot$ the two-sided decay of the tensor
normalisation $\cdot$ Chapter~\ref{ch:10} $\cdot$ stated; proof incomplete
\item the statement on rotations in the planes containing the time axis $\cdot$
Chapter~\ref{ch:10} $\cdot$ stated; proof incomplete
\end{enumerate}

\paragraph{(5) Model theorems: proved in their models; not statements about Ba{\l}aban's step.}
\begin{enumerate}
\item the model theorem on inactivity at full band width $\cdot$ Chapter~\ref{ch:8} $\cdot$
model; proved in its model
\item the model theorem on the exit value $\cdot$ Chapter~\ref{ch:8} $\cdot$ model; proved in its
model
\item the model theorem on the exit rate $\cdot$ Chapter~\ref{ch:8} $\cdot$ model; proved in its
model
\item the model theorem on value separation $\cdot$ Chapter~\ref{ch:8} $\cdot$ model; proved
in its model
\end{enumerate}

%% ---- end inlined file: chapters/c13 ----

%% ---- end inlined file: body ----

\clearpage
\addcontentsline{toc}{section}{References}

\begin{thebibliography}{99}

\bibitem{B-CMP95}
T.~Ba{\l}aban, \emph{Propagators and renormalization transformations for lattice gauge theories.~I}, Comm. Math. Phys. \textbf{95} (1984), 17--40.

\bibitem{B-CMP96}
T.~Ba{\l}aban, \emph{Propagators and renormalization transformations for lattice gauge theories.~II}, Comm. Math. Phys. \textbf{96} (1984), 223--250.

\bibitem{B-CMP98}
T.~Ba{\l}aban, \emph{Averaging operations for lattice gauge theories}, Comm. Math. Phys. \textbf{98} (1985), 17--51.

\bibitem{B-CMP99a}
T.~Ba{\l}aban, \emph{Spaces of regular gauge field configurations on a lattice and gauge fixing conditions}, Comm. Math. Phys. \textbf{99} (1985), 75--102.

\bibitem{B-CMP99b}
T.~Ba{\l}aban, \emph{Propagators for lattice gauge theories in a background field}, Comm. Math. Phys. \textbf{99} (1985), 389--434.

\bibitem{B-CMP102a}
T.~Ba{\l}aban, \emph{Ultraviolet stability of three-dimensional lattice pure gauge field theories}, Comm. Math. Phys. \textbf{102} (1985), 255--275.

\bibitem{B-CMP102b}
T.~Ba{\l}aban, \emph{The variational problem and background fields in renormalization group method for lattice gauge theories}, Comm. Math. Phys. \textbf{102} (1985), 277--309.

\bibitem{B-CMP109}
T.~Ba{\l}aban, \emph{Renormalization group approach to lattice gauge field theories.~I. Generation of effective actions in a small field approximation and a coupling constant renormalization in four dimensions}, Comm. Math. Phys. \textbf{109} (1987), 249--301.

\bibitem{B-CMP116}
T.~Ba{\l}aban, \emph{Renormalization group approach to lattice gauge field theories.~II. Cluster expansions}, Comm. Math. Phys. \textbf{116} (1988), 1--22.

\bibitem{B-CMP119}
T.~Ba{\l}aban, \emph{Convergent renormalization expansions for lattice gauge theories}, Comm. Math. Phys. \textbf{119} (1988), 243--285.

\bibitem{B-CMP122a}
T.~Ba{\l}aban, \emph{Large field renormalization.~I. The basic step of the R operation}, Comm. Math. Phys. \textbf{122} (1989), 175--202.

\bibitem{B-CMP122b}
T.~Ba{\l}aban, \emph{Large field renormalization.~II. Localization, exponentiation, and bounds for the R operation}, Comm. Math. Phys. \textbf{122} (1989), 355--392.

\bibitem{OS73}
K.~Osterwalder and R.~Schrader, \emph{Axioms for Euclidean Green's functions}, Comm. Math. Phys. \textbf{31} (1973), 83--112.

\bibitem{OS75}
K.~Osterwalder and R.~Schrader, \emph{Axioms for Euclidean Green's functions.~II}, Comm. Math. Phys. \textbf{42} (1975), 281--305.

\bibitem{OSe78}
K.~Osterwalder and E.~Seiler, \emph{Gauge field theories on a lattice}, Ann. Physics \textbf{110} (1978), 440--471.

\bibitem{Sei82}
E.~Seiler, \emph{Gauge Theories as a Problem of Constructive Quantum Field Theory and Statistical Mechanics}, Lecture Notes in Physics \textbf{159}, Springer, Berlin, 1982.

\bibitem{Sid67}
Z.~{\v{S}}id{\'a}k, \emph{Rectangular confidence regions for the means of multivariate normal distributions}, J. Amer. Statist. Assoc. \textbf{62} (1967), 626--633.

\bibitem{CT73}
J.~M.~Combes and L.~Thomas, \emph{Asymptotic behaviour of eigenfunctions for multiparticle Schr{\"o}dinger operators}, Comm. Math. Phys. \textbf{34} (1973), 251--270.

\bibitem{GJ87}
J.~Glimm and A.~Jaffe, \emph{Quantum Physics. A Functional Integral Point of View}, 2nd ed., Springer, Berlin, 1987.

\bibitem{KP86}
R.~Koteck\'y and D.~Preiss, \emph{Cluster expansion for abstract polymer models}, Comm. Math. Phys. \textbf{103} (1986), 491--498.

\bibitem{FILS78}
J.~Fr\"ohlich, R.~Israel, E.~H.~Lieb and B.~Simon, \emph{Phase transitions and reflection positivity. I. General theory and long range lattice models}, Comm. Math. Phys. \textbf{62} (1978), 1--34.

\end{thebibliography}

